\documentclass{book}
%\usepackage[letterpaper, portrait, margin=1cm]{geometry}
%\usepackage[letterpaper, bindingoffset=0.2in, left=1in,right=1in,top=.5in,bottom=.5in,footskip=.25in,marginparwidth=5em]{geometry}
\usepackage[letterpaper, left=1in,right=1in,top=.5in,bottom=.5in,footskip=.25in,marginparwidth=1cm]{geometry}
% ---------------------
% mini-table-of-content
% ---------------------
\usepackage{minitoc}
\setcounter{minitocdepth}{1}
\setlength{\mtcindent}{24pt}
\setcounter{secnumdepth}{-2}
%\renewcommand{\mtcfont}{\small\rm}
%\renewcommand{\mtcSfont}{\small\bf}
%\usepackage{setspace}
%\usepackage{tocloft}
%\setlength\cftparskip{-1.2pt}
%\setlength\cftbeforesecskip{1.3pt}
%\setlength\cftaftertoctitleskip{2pt}
%\renewcommand{\cftsecafterpnum}{\hspace*{02.0em}}
%\renewcommand{\cftsubsecafterpnum}{\hspace*{02.0em}}

% ---------------------------
% Chinese Characters Packages
% ---------------------------
\usepackage{fontspec} 
\usepackage{xeCJK}
\setmainfont{Times}
\setCJKmainfont{BiauKai}
\newfontfamily\sblgoodhebrew{SBL BibLit}[Script=Hebrew,Contextuals=Alternate]
\newfontfamily\sblgoodgreek{SBL BibLit}[Script=Greek,Contextuals=Alternate]

\usepackage{ifpdf,cite,algorithmic,url,tikz}
\usepackage[cmex10]{amsmath}

% ---------------------------
% Hebrew Characters Packages
% ---------------------------
\usepackage{polyglossia}
\setmainfont{Times New Roman}

% -------
% General
% -------
\usepackage{multicol}
\usepackage{multirow}
\usepackage{color,colortbl}
\usepackage{xparse}
\usepackage{pbox}
\usepackage{stackengine}
\usepackage{titlesec}% http://ctan.org/pkg/titlesec
\usepackage{tabularx}
\usepackage{xltabular}
\usepackage{titlesec}
\usepackage{makecell}
\newcommand{\sectionbreak}{\clearpage}

\author{
  Editor, Michael Chan\\
  \texttt{michaelchan\_wahyan@yahoo.com.hk}
}
\usepackage{tocloft}

\usepackage{hyperref}
\hypersetup{
    colorlinks=true, % set true if you want colored links
    linktoc   =all , % set to all if you want both sections and subsections linked
    linkcolor =blue, % choose some color if you want links to stand out
}

% ----------
% Afterword
% ----------
\usepackage{marginnote}
\usepackage{sectsty}
\usepackage{ragged2e}
\usepackage{lineno}
\usepackage{xcolor}
\usepackage{paracol}

\begin{document}

\clearpage
%% temporary titles
% command to provide stretchy vertical space in proportion
\newcommand\nbvspace[1][3]{\vspace*{\stretch{#1}}}
% allow some slack to avoid under/overfull boxes
\newcommand\nbstretchyspace{\spaceskip0.5em plus 0.25em minus 0.25em}
% To improve spacing on titlepages
\newcommand{\nbtitlestretch}{\spaceskip0.6em}
\pagestyle{empty}
\begin{center}
\bfseries
\nbvspace[1]
\Huge
{%\nbtitlestretch
\Large
\textbf{城聯社區教會 粵語講道逐字稿 2018-present \\
       Youtube Channel: Urban Voice Community Church
       }}

\nbvspace[1]

{\large
Editor: Michael\\
\texttt{michaelchan\_wahyan@yahoo.com.hk}
}

\nbvspace[1]

{\large
Revision: \texttt{v1.3}\\
Last Update: \today
}


\vfill
\begin{tikzpicture}
    %remove comment for OT cover%\node (0,0) [opacity=0.03]{\includegraphics[width=15cm]{../bible_out/ot_frontcover.png}} ;
    %remove comment for NT cover%\node (0,0) [opacity=0.03]{\includegraphics[width=15cm]{../bible_out/christ_on_cross.png}} ;
    %remove comment for Bible cover%\node (0,0) [xshift=0.8cm, yshift=+2cm, opacity=0.03]{\includegraphics[width=10cm]{./christ_on_cross.png}} ;
    %remove comment for Bible cover%\node (0,0) [              yshift=-2cm, opacity=0.03]{\includegraphics[width=14cm]{./ot_frontcover.png}} ;
\end{tikzpicture}
\vfill

\end{center}

\newpage

\setcounter{tocdepth}{0}
\dominitoc
\begin{multicols}{3}
\addtocontents{toc}{\protect\hypertarget{toc}{}}
\tableofcontents
\end{multicols}

\large
%\twocolumn

% the color definition syntax is as follow:
% \definecolor{name}{system}{definition}
% example: a mono-channel color can be defined as
%          \definecolor{Gray}{gray}{0.9}
% example: an rgb-3-channel color can be defined as
%          \definecolor{LightCyan}{rgb}{0.88,1,1}
%          \definecolor{pink}{rgb}{0.68,0,0.68}

\definecolor{CUV1LightRed}{rgb}{1,0.75,0.75}     % for CUV1
\definecolor{LZZVLightGray}{rgb}{0.9,0.9,0.9}    % for LZZ
\definecolor{KJVVLightGreen}{rgb}{0.75,1,0.85}   % for KJV
\definecolor{CUV2LightYellow}{rgb}{1,1,0.75}     % for CUV2
\definecolor{CNVVLightBrown}{rgb}{1,0.85,0.7}    % for CNV
\definecolor{NRSVLightBlue}{rgb}{0.75,1,1}       % for NRSV
\definecolor{WENLLightPurple}{rgb}{0.95,0.85,0.9}% for WENL
\definecolor{TCV19PaleGreen}{rgb}{0.85,1,0.95}   % for TCV19
\definecolor{MSGVLightWhite}{rgb}{0.98,0.98,0.98}% for MSGV
\definecolor{NETSLightRed}{rgb}{1,0.75,0.75}     % for NETS
\definecolor{JPS1917LightYellow}{rgb}{1,1,0.75}  % for JPS1917
\definecolor{SBLGNTPaleRed}{rgb}{1,0.85,0.80}    % for SBLGNT

\section{目錄\small{(順時)}}
\label{sec:index_chronic}
{ \scriptsize


\begin{xltabular}{\textwidth}{|p{0.15\textwidth} p{0.6\textwidth}|p{0.07\textwidth} p{0.1\textwidth}|}
\hline
  & \hyperref[sec:Kg3cX9_FXZI]{3/30/18 受苦節崇拜  Good Friday Service} & 2018-04-01 & \href{https://youtube.com/watch?v=Kg3cX9_FXZI}{\texttt{Kg3cX9\_FXZI}} \\
  & \hyperref[sec:Z07LUlzvgG0]{4/8/18 主日講道: 基督徒成長12步--\#4 探索真理 (陳崇基牧師)} & 2018-04-10 & \href{https://youtube.com/watch?v=Z07LUlzvgG0}{\texttt{Z07LUlzvgG0}} \\
  & \hyperref[sec:a492NMYDqm0]{4/15/18 主日講道: 兩種禱告與生命(吳君才牧師)} & 2018-04-17 & \href{https://youtube.com/watch?v=a492NMYDqm0}{\texttt{a492NMYDqm0}} \\
  & \hyperref[sec:OhwVkVNV6k4]{4/29/18 主日講道: 基督徒成長12步--\#5 小組團契} & 2018-04-30 & \href{https://youtube.com/watch?v=OhwVkVNV6k4}{\texttt{OhwVkVNV6k4}} \\
  & \hyperref[sec:Y6D4skNOieM]{6/3/18 主日講道: 基督徒成長12步\#7 恩賜事奉 (陳崇基牧師)} & 2018-06-11 & \href{https://youtube.com/watch?v=Y6D4skNOieM}{\texttt{Y6D4skNOieM}} \\
以賽亞書 41:10 & \hyperref[sec:qSAJBUcq9PQ]{1/13/2019 2018 年最多人搜尋的聖經 - 2 - 以賽亞書 41\_10} & 2019-01-15 & \href{https://youtube.com/watch?v=qSAJBUcq9PQ}{\texttt{qSAJBUcq9PQ}} \\
  & \hyperref[sec:_FsWKr6IpKs]{7/1/2018 主日講道: 宣教--定意研讀並遵行神的話 (Pastor Kyle Stevens)} & 2019-07-03 & \href{https://youtube.com/watch?v=-FsWKr6IpKs}{\texttt{-FsWKr6IpKs}} \\
  & \hyperref[sec:MAxHEDY7zDA]{6/30/19 主日講道:屬靈的人 (吳君才牧師)} & 2019-07-06 & \href{https://youtube.com/watch?v=MAxHEDY7zDA}{\texttt{MAxHEDY7zDA}} \\
  & \hyperref[sec:1dt83c6RWRk]{我屋企ge 主人} & 2019-07-09 & \href{https://youtube.com/watch?v=1dt83c6RWRk}{\texttt{1dt83c6RWRk}} \\
詩篇  & \hyperref[sec:ip7wGoUuLe8]{詩篇人生 (1) - 詩篇 73 講道 (陳崇基牧師)} & 2019-07-10 & \href{https://youtube.com/watch?v=ip7wGoUuLe8}{\texttt{ip7wGoUuLe8}} \\
  & \hyperref[sec:1lBwTIY05FM]{當被壓太重時...} & 2019-07-13 & \href{https://youtube.com/watch?v=1lBwTIY05FM}{\texttt{1lBwTIY05FM}} \\
  & \hyperref[sec:Qfkef2FSGUk]{3/17/19 主日講道:智取人生 (二) - 陳崇基牧師} & 2019-07-14 & \href{https://youtube.com/watch?v=Qfkef2FSGUk}{\texttt{Qfkef2FSGUk}} \\
  & \hyperref[sec:qc2nb_Nb9qc]{8/19/18 主日講道: 基督徒的教會觀(吳君才牧師)} & 2019-07-14 & \href{https://youtube.com/watch?v=qc2nb-Nb9qc}{\texttt{qc2nb-Nb9qc}} \\
  & \hyperref[sec:TEi_5oJPKrM]{11/25/18 主日講道 真誠親密的感恩 (張麗仙傳道)} & 2019-07-16 & \href{https://youtube.com/watch?v=TEi-5oJPKrM}{\texttt{TEi-5oJPKrM}} \\
  & \hyperref[sec:NpmOHXaC8DQ]{11/4/18 主日講道:  女子組加油 \#2: 顛覆男人世界的聖經女子(陳崇基牧師)} & 2019-07-16 & \href{https://youtube.com/watch?v=NpmOHXaC8DQ}{\texttt{NpmOHXaC8DQ}} \\
  & \hyperref[sec:YsIxBX5tTR0]{6/17/18 主日講道: 信靠主 (吳君才牧師)} & 2019-07-16 & \href{https://youtube.com/watch?v=YsIxBX5tTR0}{\texttt{YsIxBX5tTR0}} \\
  & \hyperref[sec:o0LouZaecN0]{7/15/2018 主日講道: When Should We Pray (Steve Tan)} & 2019-07-16 & \href{https://youtube.com/watch?v=o0LouZaecN0}{\texttt{o0LouZaecN0}} \\
詩篇  & \hyperref[sec:lgP5_8Xevo8]{詩篇人生 (2) - 詩篇 90 講道 (陳崇基牧師)} & 2019-07-16 & \href{https://youtube.com/watch?v=lgP5-8Xevo8}{\texttt{lgP5-8Xevo8}} \\
  & \hyperref[sec:5SLOPVCSwvk]{12/16/18 主日講道: The Meaning of Christmas (Sam Johnson)} & 2019-07-17 & \href{https://youtube.com/watch?v=5SLOPVCSwvk}{\texttt{5SLOPVCSwvk}} \\
  & \hyperref[sec:J7N4kVhPD7E]{6/10/18 主日講道: 這裡敬拜的是怎樣的人 (張麗仙傳道)} & 2019-07-18 & \href{https://youtube.com/watch?v=J7N4kVhPD7E}{\texttt{J7N4kVhPD7E}} \\
  & \hyperref[sec:3CLlvNWl34w]{4/21/19 主日講道:初熟的果子(陳崇基牧師)} & 2019-07-19 & \href{https://youtube.com/watch?v=3CLlvNWl34w}{\texttt{3CLlvNWl34w}} \\
  & \hyperref[sec:OTRUSy4_3CA]{4/14/2019 主日講道:棕樹節主日-「榮進耶路撒冷」(陳崇基牧師)} & 2019-07-20 & \href{https://youtube.com/watch?v=OTRUSy4_3CA}{\texttt{OTRUSy4\_3CA}} \\
  & \hyperref[sec:zcyVrDXZup0]{7/22/2018 主日講道: 基督徒成長12步\#8 委身傳揚(陳崇基牧師)} & 2019-07-21 & \href{https://youtube.com/watch?v=zcyVrDXZup0}{\texttt{zcyVrDXZup0}} \\
  & \hyperref[sec:QRMnvNFd6QE]{2/24/19 主日講道: 事奉的迷思 (Antonio Luo)} & 2019-07-22 & \href{https://youtube.com/watch?v=QRMnvNFd6QE}{\texttt{QRMnvNFd6QE}} \\
  & \hyperref[sec:UcXT0Snl5mA]{5/19/19 主日講道\_耶穌天國的小故事-3\_麵酵 (陳崇基牧師)} & 2019-07-22 & \href{https://youtube.com/watch?v=UcXT0Snl5mA}{\texttt{UcXT0Snl5mA}} \\
  & \hyperref[sec:Kxw1fL4mQTc]{7/8/2018 主日講道: 基督徒的婚姻 (吳君才牧師)} & 2019-07-23 & \href{https://youtube.com/watch?v=Kxw1fL4mQTc}{\texttt{Kxw1fL4mQTc}} \\
  & \hyperref[sec:me3DczkMU3E]{9/23/18 主日講道: 基督徒成長12步--\#10 領導事工 (陳崇基牧師)} & 2019-07-23 & \href{https://youtube.com/watch?v=me3DczkMU3E}{\texttt{me3DczkMU3E}} \\
  & \hyperref[sec:lPh_PhzFDQI]{9/2/18 主日講道: 教會藍圖\#2 -- 教會需要問責 (陳崇基牧師)} & 2019-07-24 & \href{https://youtube.com/watch?v=lPh-PhzFDQI}{\texttt{lPh-PhzFDQI}} \\
  & \hyperref[sec:KZEdWxG87Rg]{9/9/18 主日講道: 教會藍圖\#3 -- 教會領袖的素質 (陳崇基牧師)} & 2019-07-24 & \href{https://youtube.com/watch?v=KZEdWxG87Rg}{\texttt{KZEdWxG87Rg}} \\
提摩太前書 2:8-15 & \hyperref[sec:tDpoW3ZdL20]{12/2/18 主日講道: 女子組加油! \#4 提摩太前書2\_8-15  (陳崇基牧師)} & 2019-07-25 & \href{https://youtube.com/watch?v=tDpoW3ZdL20}{\texttt{tDpoW3ZdL20}} \\
耶利米書 29:11 & \hyperref[sec:zzrp0tNV2iI]{1/6/19 主日講道: 2018 最多人搜尋的聖經 (陳崇基牧師) - 耶利米書 29\_11} & 2019-07-25 & \href{https://youtube.com/watch?v=zzrp0tNV2iI}{\texttt{zzrp0tNV2iI}} \\
詩篇  & \hyperref[sec:y_ka_wkoG0A]{詩篇23 篇 (陳崇基牧師) 2018-06-24 主日講道} & 2019-07-25 & \href{https://youtube.com/watch?v=y-ka_wkoG0A}{\texttt{y-ka\_wkoG0A}} \\
以賽亞書 14:1-4 & \hyperref[sec:nzhjFy_6Umo]{3/24/2019 主日講道 - 以賽亞書14\_1-4 (Tony Wong 牧師)} & 2019-07-26 & \href{https://youtube.com/watch?v=nzhjFy-6Umo}{\texttt{nzhjFy-6Umo}} \\
  & \hyperref[sec:dNN3GBdDj2k]{7/21/19 主日講道: 懇切需要主耶穌 {Pastor John Yi}} & 2019-07-28 & \href{https://youtube.com/watch?v=dNN3GBdDj2k}{\texttt{dNN3GBdDj2k}} \\
  & \hyperref[sec:_5Jd2EPsTDg]{8/5/18 主日講道: 基督徒成長12步--\#9 門訓關係 (陳崇基牧師)} & 2019-07-29 & \href{https://youtube.com/watch?v=_5Jd2EPsTDg}{\texttt{\_5Jd2EPsTDg}} \\
  & \hyperref[sec:1yfYPGx_z_4]{10/7/18 主日講道: 基督徒成長12步 \#12 植堂擴展 (陳崇基牧師)} & 2019-07-30 & \href{https://youtube.com/watch?v=1yfYPGx-z-4}{\texttt{1yfYPGx-z-4}} \\
以弗所書  & \hyperref[sec:6IhSTx8c43Q]{10/14/18 主日講道: 以弗所書與教會選舉 (張麗仙傳道)} & 2019-07-31 & \href{https://youtube.com/watch?v=6IhSTx8c43Q}{\texttt{6IhSTx8c43Q}} \\
  & \hyperref[sec:dClKoGWQ8rs]{10/28/18 主日講道: 女子組加油 \#1: 顛覆男人世界的聖經女子(陳崇基牧師)} & 2019-08-01 & \href{https://youtube.com/watch?v=dClKoGWQ8rs}{\texttt{dClKoGWQ8rs}} \\
  & \hyperref[sec:9yvFXOIg3So]{8/26/18 主日講道: 教會藍圖\#1 -- 在醜聞中如何重建教會 (陳崇基牧師)} & 2019-08-01 & \href{https://youtube.com/watch?v=9yvFXOIg3So}{\texttt{9yvFXOIg3So}} \\
  & \hyperref[sec:ALcJ_99Kiw8]{安息日為人設立 - 張國棟博士 (全篇講道)} & 2019-08-01 & \href{https://youtube.com/watch?v=ALcJ-99Kiw8}{\texttt{ALcJ-99Kiw8}} \\
  & \hyperref[sec:PGikCuC4cZg]{1/27/19 主日講道: 創造與敬拜 (張麗仙傳道)} & 2019-08-02 & \href{https://youtube.com/watch?v=PGikCuC4cZg}{\texttt{PGikCuC4cZg}} \\
  & \hyperref[sec:wT2Sf41UjYU]{10/21/18 主日講道: 三件要事 (吳君才牧師)} & 2019-08-03 & \href{https://youtube.com/watch?v=wT2Sf41UjYU}{\texttt{wT2Sf41UjYU}} \\
  & \hyperref[sec:NJ2Erozw8Ss]{11/18/18 主日講道: 女子組加油! \#3 顛覆男人世界的聖經女子 (陳崇基牧師)} & 2019-08-03 & \href{https://youtube.com/watch?v=NJ2Erozw8Ss}{\texttt{NJ2Erozw8Ss}} \\
  & \hyperref[sec:QpmQmJgiMg8]{4/7/2019 主日講道: 智取人生(三) 陳崇基牧師} & 2019-08-03 & \href{https://youtube.com/watch?v=QpmQmJgiMg8}{\texttt{QpmQmJgiMg8}} \\
  & \hyperref[sec:JN7FdDAB0F8]{4/28/19 主日講道\_耶穌天國的小故事-1\_撒種(陳崇基牧師)} & 2019-08-05 & \href{https://youtube.com/watch?v=JN7FdDAB0F8}{\texttt{JN7FdDAB0F8}} \\
 4:6-8 & \hyperref[sec:tJHkj9NzTVE]{7/29/18 主日講道: Running The Race - 2 Tim 4\_6-8 (Rev. John Yi)} & 2019-08-05 & \href{https://youtube.com/watch?v=tJHkj9NzTVE}{\texttt{tJHkj9NzTVE}} \\
  & \hyperref[sec:pK11m4fqws8]{8/12/18 主日講道:  邊個係惡人 (Joyce Tam)} & 2019-08-05 & \href{https://youtube.com/watch?v=pK11m4fqws8}{\texttt{pK11m4fqws8}} \\
  & \hyperref[sec:qkxKhCtF31k]{5/12/19 主日講道\_耶穌天國的小故事-2\_恩典 (陳崇基牧師)} & 2019-08-06 & \href{https://youtube.com/watch?v=qkxKhCtF31k}{\texttt{qkxKhCtF31k}} \\
  & \hyperref[sec:KWfIxDVxl3w]{12/30/18 主日講道: 齊心作主工 (吳君才牧師)} & 2019-08-07 & \href{https://youtube.com/watch?v=KWfIxDVxl3w}{\texttt{KWfIxDVxl3w}} \\
詩篇  & \hyperref[sec:J6o6vqGpqR4]{詩篇人生 - 詩篇 2 - 全篇講道 (陳崇基牧師)} & 2019-08-08 & \href{https://youtube.com/watch?v=J6o6vqGpqR4}{\texttt{J6o6vqGpqR4}} \\
  & \hyperref[sec:6oKVLkh536w]{3/10/2019 主日講道:神愛世人 (Marcus Yang)} & 2019-08-09 & \href{https://youtube.com/watch?v=6oKVLkh536w}{\texttt{6oKVLkh536w}} \\
  & \hyperref[sec:cPh4mXwlq7I]{1/20/19 主日講道: 信靠主 (吳君才牧師)} & 2019-08-13 & \href{https://youtube.com/watch?v=cPh4mXwlq7I}{\texttt{cPh4mXwlq7I}} \\
箴言 4:23 & \hyperref[sec:H3Nf68xn8Zg]{2/3/19 主日講道: 2018年最多人搜尋的聖經 3: 箴言4\_23(陳崇基牧師)} & 2019-08-13 & \href{https://youtube.com/watch?v=H3Nf68xn8Zg}{\texttt{H3Nf68xn8Zg}} \\
  & \hyperref[sec:C4RfQvzd3mA]{2/10/19 主日講道: 施比受更為有福 (陳崇基牧師)} & 2019-08-14 & \href{https://youtube.com/watch?v=C4RfQvzd3mA}{\texttt{C4RfQvzd3mA}} \\
  & \hyperref[sec:NaL7AxYsROQ]{3/31/19 主日講道:重獲生命樹 (張麗仙傳道)} & 2019-08-14 & \href{https://youtube.com/watch?v=NaL7AxYsROQ}{\texttt{NaL7AxYsROQ}} \\
  & \hyperref[sec:iiRocZadlgU]{9/30/18 主日講道: 基督徒成長12步\#11 核心帶領 (陳崇基牧師)} & 2019-08-14 & \href{https://youtube.com/watch?v=iiRocZadlgU}{\texttt{iiRocZadlgU}} \\
  & \hyperref[sec:xXuB_PHmaoM]{3/3/19 主日講道: 智取人生 - 學會生命的智慧(陳崇基牧師)} & 2019-08-15 & \href{https://youtube.com/watch?v=xXuB-PHmaoM}{\texttt{xXuB-PHmaoM}} \\
  & \hyperref[sec:8I6GZqi_kx0]{9/16/18 主日講道:我所擁有的(張麗仙傳道)} & 2019-08-15 & \href{https://youtube.com/watch?v=8I6GZqi_kx0}{\texttt{8I6GZqi\_kx0}} \\
  & \hyperref[sec:sixId3mzKXA]{12/9/18 主日講道: 女子組加油! \#5 馬利亞(陳崇基牧師)} & 2019-08-18 & \href{https://youtube.com/watch?v=sixId3mzKXA}{\texttt{sixId3mzKXA}} \\
詩篇  & \hyperref[sec:AG5ZPTEPotw]{詩篇人生 - 上行之詩 - 詩篇 120-134 - 全篇講道 (陳崇基牧師)} & 2019-08-27 & \href{https://youtube.com/watch?v=AG5ZPTEPotw}{\texttt{AG5ZPTEPotw}} \\
詩篇  & \hyperref[sec:JPzN1z7hdRY]{8/25/19 主日講道: 詩篇人生 - 詛咒詩歌 (陳崇基牧師)} & 2019-09-01 & \href{https://youtube.com/watch?v=JPzN1z7hdRY}{\texttt{JPzN1z7hdRY}} \\
詩篇  & \hyperref[sec:C3L6GM4kyjo]{9-1-19 主日講道: 詩篇人生 -哀歌 (陳崇基牧師)} & 2019-09-08 & \href{https://youtube.com/watch?v=C3L6GM4kyjo}{\texttt{C3L6GM4kyjo}} \\
詩篇  & \hyperref[sec:EogeiLrQX3g]{9/8/19 主日講道: 詩篇人生 -有福的人(陳崇基牧師)} & 2019-09-12 & \href{https://youtube.com/watch?v=EogeiLrQX3g}{\texttt{EogeiLrQX3g}} \\
  & \hyperref[sec:7VnGbL9RHgg]{9/22/19 主日講道: 在敬拜中遇見你 (Brenda Li 傳道)} & 2019-09-29 & \href{https://youtube.com/watch?v=7VnGbL9RHgg}{\texttt{7VnGbL9RHgg}} \\
創世記  & \hyperref[sec:FaXJPQMXp6c]{與人同在的神 \#1 - 創世記 1 - 20190929 主日講道 (陳崇基牧師)} & 2019-10-05 & \href{https://youtube.com/watch?v=FaXJPQMXp6c}{\texttt{FaXJPQMXp6c}} \\
  & \hyperref[sec:LfJIKNZbdWE]{10/06/19 主日講道: 參透生命的智慧 (吳君才牧師)} & 2019-10-10 & \href{https://youtube.com/watch?v=LfJIKNZbdWE}{\texttt{LfJIKNZbdWE}} \\
創世記 1:26-27 & \hyperref[sec:1bmqQECjSXg]{與人同在的神 \#2 - 創世記 1\_26-27 - 2019/10/13 (陳崇基牧師)} & 2019-10-17 & \href{https://youtube.com/watch?v=1bmqQECjSXg}{\texttt{1bmqQECjSXg}} \\
創世記  & \hyperref[sec:0WFluTmE3pc]{與人同在的神 \#3 - 創世記 3 - 2019/10/20 (陳崇基牧師)} & 2019-10-23 & \href{https://youtube.com/watch?v=0WFluTmE3pc}{\texttt{0WFluTmE3pc}} \\
創世記 11:1-9 & \hyperref[sec:Gl7QFa_KpBw]{10/27/19 主日講道:與人同在的神 IV - 創世記 11\_1-9 巴別塔 (陳崇基牧師)} & 2019-11-01 & \href{https://youtube.com/watch?v=Gl7QFa-KpBw}{\texttt{Gl7QFa-KpBw}} \\
  & \hyperref[sec:xwcSwwS07wQ]{11/3/19 主日講道: 圍繞父母親的見證人 (萬楊婉貞博士)} & 2019-11-07 & \href{https://youtube.com/watch?v=xwcSwwS07wQ}{\texttt{xwcSwwS07wQ}} \\
  & \hyperref[sec:z_G8mnyC7Vg]{11/17/19 主日講道: 與人同在的神 5 - 律法(陳崇基牧師)} & 2019-11-24 & \href{https://youtube.com/watch?v=z-G8mnyC7Vg}{\texttt{z-G8mnyC7Vg}} \\
  & \hyperref[sec:pqDLnSQO_8g]{11/24/19 主日講道:感恩,愛,赦免(吳君才牧師)} & 2019-11-30 & \href{https://youtube.com/watch?v=pqDLnSQO-8g}{\texttt{pqDLnSQO-8g}} \\
  & \hyperref[sec:SGo8UblJv4Q]{Premiere Pro video editing (part 1)} & 2019-12-05 & \href{https://youtube.com/watch?v=SGo8UblJv4Q}{\texttt{SGo8UblJv4Q}} \\
  & \hyperref[sec:joskqj5_Na8]{Premiere Pro video editing (part 2)} & 2019-12-05 & \href{https://youtube.com/watch?v=joskqj5_Na8}{\texttt{joskqj5\_Na8}} \\
  & \hyperref[sec:2Ci7sbMUq14]{12/1/19 主日講道:Connect! 密切連繫 (張麗仙傳道)} & 2019-12-06 & \href{https://youtube.com/watch?v=2Ci7sbMUq14}{\texttt{2Ci7sbMUq14}} \\
 2:1-7 & \hyperref[sec:zYSzdaPaiDI]{12/8/19 主日講道: 與人同在的神 6 - 路2\_1-7 (陳崇基牧師)} & 2019-12-15 & \href{https://youtube.com/watch?v=zYSzdaPaiDI}{\texttt{zYSzdaPaiDI}} \\
馬太福音 1:1-17 & \hyperref[sec:99dN28SUzkQ]{12/15/19 主日講道: 與人同在的神 7, 馬太福音1\_1-17 (陳崇基牧師)} & 2019-12-20 & \href{https://youtube.com/watch?v=99dN28SUzkQ}{\texttt{99dN28SUzkQ}} \\
馬太福音 2:13-23 & \hyperref[sec:4ZoSLfbOWng]{12/29/19 主日講道: 聖誕故事 - 馬太福音2\_13-23 (陳崇基牧師)} & 2020-01-05 & \href{https://youtube.com/watch?v=4ZoSLfbOWng}{\texttt{4ZoSLfbOWng}} \\
腓立比書 2:6-11 & \hyperref[sec:HLQkwHUhsgc]{1/12/20 主日講道: 天國公民地上行 2 (腓立比書2\_6-11) - 陳崇基牧師} & 2020-01-19 & \href{https://youtube.com/watch?v=HLQkwHUhsgc}{\texttt{HLQkwHUhsgc}} \\
腓立比書 2:6-11 & \hyperref[sec:6ypOdtEuNl8]{1/5/20 主日講道: 天國公民地上行 1 (腓立比書2\_6-11) - 陳崇基牧師} & 2020-01-19 & \href{https://youtube.com/watch?v=6ypOdtEuNl8}{\texttt{6ypOdtEuNl8}} \\
提摩太前書 2:1-24 & \hyperref[sec:Pe1TNggJhos]{主日講道: 屬靈生命的模式 提摩太前書2\_1-24 (吳君才牧師)} & 2020-01-26 & \href{https://youtube.com/watch?v=Pe1TNggJhos}{\texttt{Pe1TNggJhos}} \\
利未記 13:1-6 & \hyperref[sec:ez9_0uQY1v4]{2020/2/2 主日講道: 疫境前行 (利未記13\_1-6) - 陳崇基牧師} & 2020-02-06 & \href{https://youtube.com/watch?v=ez9_0uQY1v4}{\texttt{ez9\_0uQY1v4}} \\
利未記 11:1-19 & \hyperref[sec:imYjaG_VAbM]{2020/2/9 主日講道: 疫境前行 2 (利未記11\_1-19) - 陳崇基牧師} & 2020-02-14 & \href{https://youtube.com/watch?v=imYjaG_VAbM}{\texttt{imYjaG\_VAbM}} \\
歌羅西書 1:25-29 & \hyperref[sec:CWGXYDFB8bU]{2020/2/16 主日講道: 全面的佈道事工 (歌羅西書1\_25-29) - 黃茂林傳道} & 2020-02-21 & \href{https://youtube.com/watch?v=CWGXYDFB8bU}{\texttt{CWGXYDFB8bU}} \\
使徒行傳 10:1-48 & \hyperref[sec:zCN_cR7XoxQ]{2020/2/23 主日講道: 疫境前行 3 (使徒行傳 10\_1-48) - 陳崇基牧師} & 2020-03-01 & \href{https://youtube.com/watch?v=zCN-cR7XoxQ}{\texttt{zCN-cR7XoxQ}} \\
詩篇 27:1-14 & \hyperref[sec:d0PXI2JCtZU]{2020/2/9 主日講道: 恐懼中, 記著這一件事 (詩篇27\_1-14) - 張麗仙傳道} & 2020-03-10 & \href{https://youtube.com/watch?v=d0PXI2JCtZU}{\texttt{d0PXI2JCtZU}} \\
傳道書  & \hyperref[sec:EN7WJcnMP5Y]{2020/3/22 崇拜直播 vlog \#1 疫境前行 (傳道書 3) - 陳崇基牧師} & 2020-03-23 & \href{https://youtube.com/watch?v=EN7WJcnMP5Y}{\texttt{EN7WJcnMP5Y}} \\
詩篇  & \hyperref[sec:o3fsHQqXKjc]{2020/3/29 崇拜直播 vlog \#2 疫境前行 (詩篇 139) - 陳崇基牧師} & 2020-03-30 & \href{https://youtube.com/watch?v=o3fsHQqXKjc}{\texttt{o3fsHQqXKjc}} \\
耶利米哀歌  & \hyperref[sec:G9uX1Fwy_5w]{2020/4/5 崇拜直播 vlog \#3 疫境前行 (耶利米哀歌) - 陳崇基牧師} & 2020-04-06 & \href{https://youtube.com/watch?v=G9uX1Fwy_5w}{\texttt{G9uX1Fwy\_5w}} \\
帖撒羅尼迦前書使徒行傳  & \hyperref[sec:kAntZ1DnVvY]{2020/4/12 復活節主日崇拜直播 vlog \#4 疫境前行 (帖撒羅尼迦前書 4; 使徒行傳 17) - 陳崇基牧師} & 2020-04-13 & \href{https://youtube.com/watch?v=kAntZ1DnVvY}{\texttt{kAntZ1DnVvY}} \\
雅各書  & \hyperref[sec:HEMq3Mncf_A]{2020/4/19 智慧廚房 vlog 信息 \#1 (雅各書 1) - 陳崇基牧師} & 2020-04-18 & \href{https://youtube.com/watch?v=HEMq3Mncf-A}{\texttt{HEMq3Mncf-A}} \\
  & \hyperref[sec:_qut0JfQmyc]{狼人殺(聖經版) round 1 精華版} & 2020-04-22 & \href{https://youtube.com/watch?v=_qut0JfQmyc}{\texttt{\_qut0JfQmyc}} \\
雅各書  & \hyperref[sec:CDVnSLlj71Y]{2020/4/26 智慧廚房 vlog 信息 \#2 (雅各書 2) - 陳崇基牧師} & 2020-04-27 & \href{https://youtube.com/watch?v=CDVnSLlj71Y}{\texttt{CDVnSLlj71Y}} \\
  & \hyperref[sec:o_g_uZyyIHo]{1 十一少被賣記原身家庭(一)} & 2020-05-04 & \href{https://youtube.com/watch?v=o-g-uZyyIHo}{\texttt{o-g-uZyyIHo}} \\
  & \hyperref[sec:qk8r6sElCdc]{Kingdom Values III - Sermon on the Mount - Pastor Sam Johnson} & 2020-05-04 & \href{https://youtube.com/watch?v=qk8r6sElCdc}{\texttt{qk8r6sElCdc}} \\
  & \hyperref[sec:VkJzbPZVpM0]{2 十一少被賣記原生家庭(二)} & 2020-05-05 & \href{https://youtube.com/watch?v=VkJzbPZVpM0}{\texttt{VkJzbPZVpM0}} \\
  & \hyperref[sec:3aPT0T6_ZVI]{3 十一少被賣記:原生家庭(三)} & 2020-05-05 & \href{https://youtube.com/watch?v=3aPT0T6-ZVI}{\texttt{3aPT0T6-ZVI}} \\
雅各書  & \hyperref[sec:I_xWBBrkSzQ]{2020/5/3 智慧廚房 vlog 信息 \#3 (雅各書 3) - 陳崇基牧師 (updated)} & 2020-05-09 & \href{https://youtube.com/watch?v=I-xWBBrkSzQ}{\texttt{I-xWBBrkSzQ}} \\
雅各書  & \hyperref[sec:v2v5pAP5WjA]{2020/5/10 智慧廚房 vlog 信息 \#4 (雅各書 4) - 陳崇基牧師} & 2020-05-12 & \href{https://youtube.com/watch?v=v2v5pAP5WjA}{\texttt{v2v5pAP5WjA}} \\
雅各書  & \hyperref[sec:teCFMiZwZp0]{2020/5/17 智慧廚房 vlog 信息 \#5 (雅各書 5) - 陳崇基牧師} & 2020-05-18 & \href{https://youtube.com/watch?v=teCFMiZwZp0}{\texttt{teCFMiZwZp0}} \\
  & \hyperref[sec:N2MGZTvagIA]{4 十一少被賣記:原生家庭(四)} & 2020-05-26 & \href{https://youtube.com/watch?v=N2MGZTvagIA}{\texttt{N2MGZTvagIA}} \\
  & \hyperref[sec:12OrdZOBpSQ]{5 十一少被賣記:有神同在(一)} & 2020-05-26 & \href{https://youtube.com/watch?v=12OrdZOBpSQ}{\texttt{12OrdZOBpSQ}} \\
  & \hyperref[sec:uQ6L96b0_II]{6 十一少被賣記:有神同在(二)} & 2020-06-16 & \href{https://youtube.com/watch?v=uQ6L96b0_II}{\texttt{uQ6L96b0\_II}} \\
  & \hyperref[sec:sSMekSYd36k]{7 十一少被賣記:有神同在(三)} & 2020-06-16 & \href{https://youtube.com/watch?v=sSMekSYd36k}{\texttt{sSMekSYd36k}} \\
  & \hyperref[sec:y63dV5dKyc4]{10 十一少被賣記:因禍得福－下(一)} & 2020-06-18 & \href{https://youtube.com/watch?v=y63dV5dKyc4}{\texttt{y63dV5dKyc4}} \\
  & \hyperref[sec:LqhzJpn7Fyo]{8 十一少被賣記:因禍得福－上(一)} & 2020-06-18 & \href{https://youtube.com/watch?v=LqhzJpn7Fyo}{\texttt{LqhzJpn7Fyo}} \\
  & \hyperref[sec:YbuM_vsj2Sc]{9 十一少被賣記:因禍得福－上(二)} & 2020-06-18 & \href{https://youtube.com/watch?v=YbuM_vsj2Sc}{\texttt{YbuM\_vsj2Sc}} \\
  & \hyperref[sec:dNRqxYjeS0o]{Kingdom Values 10 - Sam Johnson - A Conversation on Race} & 2020-06-22 & \href{https://youtube.com/watch?v=dNRqxYjeS0o}{\texttt{dNRqxYjeS0o}} \\
  & \hyperref[sec:qRHTClh4NmQ]{12 十一少被賣記大團圓－上(一)} & 2020-07-05 & \href{https://youtube.com/watch?v=qRHTClh4NmQ}{\texttt{qRHTClh4NmQ}} \\
  & \hyperref[sec:hwLDlGYTpto]{13 十一少被賣記大團圓－上(二)} & 2020-07-06 & \href{https://youtube.com/watch?v=hwLDlGYTpto}{\texttt{hwLDlGYTpto}} \\
  & \hyperref[sec:BmdnpSdGr5s]{15十一少被賣記大團圓－上(四)} & 2020-07-06 & \href{https://youtube.com/watch?v=BmdnpSdGr5s}{\texttt{BmdnpSdGr5s}} \\
  & \hyperref[sec:F_RgXPWQknI]{16十一少被賣記大團圓－上(五)} & 2020-07-06 & \href{https://youtube.com/watch?v=F_RgXPWQknI}{\texttt{F\_RgXPWQknI}} \\
  & \hyperref[sec:WV2ED6LTYMM]{17 十一少被賣記大團圓－中(一)} & 2020-07-06 & \href{https://youtube.com/watch?v=WV2ED6LTYMM}{\texttt{WV2ED6LTYMM}} \\
  & \hyperref[sec:egka2gEIyik]{20 十一少被賣記大團圓－中(四)} & 2020-07-06 & \href{https://youtube.com/watch?v=egka2gEIyik}{\texttt{egka2gEIyik}} \\
  & \hyperref[sec:fHYowIViTao]{22 十一少被賣記大團圓－下(二)} & 2020-07-06 & \href{https://youtube.com/watch?v=fHYowIViTao}{\texttt{fHYowIViTao}} \\
  & \hyperref[sec:6gEpaWbkTWE]{23 十一少被賣記大團圓－下(三)} & 2020-07-06 & \href{https://youtube.com/watch?v=6gEpaWbkTWE}{\texttt{6gEpaWbkTWE}} \\
  & \hyperref[sec:oFXXJCLRQwo]{14十一少被賣記大團圓－上(三)} & 2020-07-07 & \href{https://youtube.com/watch?v=oFXXJCLRQwo}{\texttt{oFXXJCLRQwo}} \\
  & \hyperref[sec:2kNd6PvPzMU]{18 十一少被賣記大團圓－中(二)} & 2020-07-07 & \href{https://youtube.com/watch?v=2kNd6PvPzMU}{\texttt{2kNd6PvPzMU}} \\
  & \hyperref[sec:QNtUzTDF7W4]{11十一少被賣記因禍得福－下(二)} & 2020-07-08 & \href{https://youtube.com/watch?v=QNtUzTDF7W4}{\texttt{QNtUzTDF7W4}} \\
  & \hyperref[sec:X74nEQWZ90I]{21 十一少被賣記大團圓－下(一)} & 2020-08-02 & \href{https://youtube.com/watch?v=X74nEQWZ90I}{\texttt{X74nEQWZ90I}} \\
  & \hyperref[sec:xOb08azcrmI]{19 十一少被賣記大團圓－中(三)} & 2020-08-05 & \href{https://youtube.com/watch?v=xOb08azcrmI}{\texttt{xOb08azcrmI}} \\
  & \hyperref[sec:EEqL82eh5Po]{聖地網上行 - 步主腳蹤 [第三集] 約旦河 - 曠野} & 2020-08-22 & \href{https://youtube.com/watch?v=EEqL82eh5Po}{\texttt{EEqL82eh5Po}} \\
  & \hyperref[sec:QHoIKk7ZHEQ]{聖地網上行 - 步主腳蹤 [第二集] 拿撒勒} & 2020-08-26 & \href{https://youtube.com/watch?v=QHoIKk7ZHEQ}{\texttt{QHoIKk7ZHEQ}} \\
  & \hyperref[sec:KnqUXPbp8Kk]{聖地網上行 - 步主腳蹤 [第八集] 登山變像} & 2020-08-26 & \href{https://youtube.com/watch?v=KnqUXPbp8Kk}{\texttt{KnqUXPbp8Kk}} \\
  & \hyperref[sec:vLHlHGB1MwM]{聖地網上行 - 步主腳蹤 [第七集] 凱撒利亞腓立比} & 2020-08-27 & \href{https://youtube.com/watch?v=vLHlHGB1MwM}{\texttt{vLHlHGB1MwM}} \\
  & \hyperref[sec:jXnsElF6xnM]{聖地網上行 - 步主腳蹤 [第九集] 榮進耶路撒冷} & 2020-08-27 & \href{https://youtube.com/watch?v=jXnsElF6xnM}{\texttt{jXnsElF6xnM}} \\
  & \hyperref[sec:PFNNLzN6LA8]{聖地網上行 - 步主腳蹤 [第十一集] 最後晚餐} & 2020-08-28 & \href{https://youtube.com/watch?v=PFNNLzN6LA8}{\texttt{PFNNLzN6LA8}} \\
  & \hyperref[sec:c1irNOxf2YQ]{聖地網上行 - 步主腳蹤 [第四集] 加利利 - 迦百農} & 2020-08-29 & \href{https://youtube.com/watch?v=c1irNOxf2YQ}{\texttt{c1irNOxf2YQ}} \\
  & \hyperref[sec:aXI0__JmVi0]{和你約定 [1] - 傾o下國安法} & 2020-09-03 & \href{https://youtube.com/watch?v=aXI0__JmVi0}{\texttt{aXI0\_\_JmVi0}} \\
  & \hyperref[sec:u3NKvcUYcmI]{聖地網上行 - 步主腳蹤 [第五集] 加利利 - 撒種者湖灣/八福山/抹大拉/希布斯} & 2020-09-06 & \href{https://youtube.com/watch?v=u3NKvcUYcmI}{\texttt{u3NKvcUYcmI}} \\
  & \hyperref[sec:m3nKOXb4VNI]{聖地網上行 - 步主腳蹤 [第十三集] 十字架} & 2020-09-06 & \href{https://youtube.com/watch?v=m3nKOXb4VNI}{\texttt{m3nKOXb4VNI}} \\
  & \hyperref[sec:o6x_Tz0d6Ww]{浩劫重生 【第1集:挪亞洪水】 陳崇基牧師 - 主日講道視頻系列} & 2020-09-29 & \href{https://youtube.com/watch?v=o6x-Tz0d6Ww}{\texttt{o6x-Tz0d6Ww}} \\
  & \hyperref[sec:9fgDv4ygGRU]{神的名字1\_Elohim (靈思牧禱)} & 2020-09-30 & \href{https://youtube.com/watch?v=9fgDv4ygGRU}{\texttt{9fgDv4ygGRU}} \\
  & \hyperref[sec:_8xftiiwoRs]{浩劫重生 【第2集:洪水之後】 陳崇基牧師 - 主日講道視頻系列} & 2020-10-05 & \href{https://youtube.com/watch?v=-8xftiiwoRs}{\texttt{-8xftiiwoRs}} \\
  & \hyperref[sec:fE8jDxIIYZM]{Kingdom Values XXIV: Matthew Chapters 5 - Pastor Sam Johnson} & 2020-10-06 & \href{https://youtube.com/watch?v=fE8jDxIIYZM}{\texttt{fE8jDxIIYZM}} \\
  & \hyperref[sec:x6dSk6R1byI]{聖地網上行 - 步主腳蹤 [第一集] 伯利恆} & 2020-10-07 & \href{https://youtube.com/watch?v=x6dSk6R1byI}{\texttt{x6dSk6R1byI}} \\
  & \hyperref[sec:vkydkBsADbY]{浩劫重生 【第3集:出埃及十災】 陳崇基牧師 - 主日講道視頻系列} & 2020-10-12 & \href{https://youtube.com/watch?v=vkydkBsADbY}{\texttt{vkydkBsADbY}} \\
  & \hyperref[sec:O6HbgHZahfI]{神的名字2\_Yahweh (靈思牧禱)} & 2020-10-14 & \href{https://youtube.com/watch?v=O6HbgHZahfI}{\texttt{O6HbgHZahfI}} \\
  & \hyperref[sec:Lh5hqFxfj84]{Kingdom Values XXVI: Matthew Chapter 6 Continued - Pastor Sam Johnson} & 2020-10-19 & \href{https://youtube.com/watch?v=Lh5hqFxfj84}{\texttt{Lh5hqFxfj84}} \\
  & \hyperref[sec:XYSdVdRz0wM]{浩劫重生 【第4集:出埃及路線】 陳崇基牧師 - 主日講道視頻系列} & 2020-10-20 & \href{https://youtube.com/watch?v=XYSdVdRz0wM}{\texttt{XYSdVdRz0wM}} \\
 6:5-6 & \hyperref[sec:c2n9YZYOe80]{Kingdom Values XXVII: Matthew Chapter 6\_5-6 Continued - Pastor Sam Johnson} & 2020-10-28 & \href{https://youtube.com/watch?v=c2n9YZYOe80}{\texttt{c2n9YZYOe80}} \\
約拿書  & \hyperref[sec:y_7ZwDcfiBI]{經緯線 [釋經視頻] EP1 約拿書 - 介紹,文學體裁} & 2020-10-30 & \href{https://youtube.com/watch?v=y-7ZwDcfiBI}{\texttt{y-7ZwDcfiBI}} \\
約拿書 1:1-6 & \hyperref[sec:WmRxFI5hs7M]{經緯線 [釋經視頻] EP2 約拿書 1\_1-6} & 2020-10-30 & \href{https://youtube.com/watch?v=WmRxFI5hs7M}{\texttt{WmRxFI5hs7M}} \\
 6:16-18 & \hyperref[sec:lgq1HtuVP6w]{Kingdom Values XXVIII: Matthew Chapter 6\_16-18 - Pastor Sam Johnson} & 2020-11-01 & \href{https://youtube.com/watch?v=lgq1HtuVP6w}{\texttt{lgq1HtuVP6w}} \\
  & \hyperref[sec:LPCHL652i_w]{浩劫重生 【第5集:亂局中的以利亞】 陳崇基牧師 - 主日講道視頻系列} & 2020-11-10 & \href{https://youtube.com/watch?v=LPCHL652i_w}{\texttt{LPCHL652i\_w}} \\
  & \hyperref[sec:rfX8QtfxB9g]{Kingdom Values XXX: Matthew Chapter 6 Prayer - Pastor Sam Johnson} & 2020-11-14 & \href{https://youtube.com/watch?v=rfX8QtfxB9g}{\texttt{rfX8QtfxB9g}} \\
約拿書 1:7-16 & \hyperref[sec:MzNtNdJRGgQ]{經緯線 [釋經視頻] EP3 約拿書 1\_7-16} & 2020-11-14 & \href{https://youtube.com/watch?v=MzNtNdJRGgQ}{\texttt{MzNtNdJRGgQ}} \\
約拿書 1:17-2:10 & \hyperref[sec:_ZlpJCFpNOg]{經緯線 [釋經視頻] EP4 約拿書 1\_17-2\_10} & 2020-11-14 & \href{https://youtube.com/watch?v=-ZlpJCFpNOg}{\texttt{-ZlpJCFpNOg}} \\
  & \hyperref[sec:ksA3W2OSofQ]{浩劫重生 【第6集:急難中的希西家】 陳崇基牧師 - 主日講道視頻系列} & 2020-11-15 & \href{https://youtube.com/watch?v=ksA3W2OSofQ}{\texttt{ksA3W2OSofQ}} \\
  & \hyperref[sec:UwZdh9361sc]{神的名字4\_Yahweh Yireh 耶和華預備 (靈思牧禱)} & 2020-11-17 & \href{https://youtube.com/watch?v=UwZdh9361sc}{\texttt{UwZdh9361sc}} \\
約拿書 3:1-10 & \hyperref[sec:kB2Ko_XJaY8]{經緯線 [釋經視頻] EP5 約拿書 3\_1-10} & 2020-11-20 & \href{https://youtube.com/watch?v=kB2Ko_XJaY8}{\texttt{kB2Ko\_XJaY8}} \\
約拿書 4:1-11 & \hyperref[sec:EWFiHkhUrkU]{經緯線 [釋經視頻] EP6 約拿書 4\_1-11} & 2020-11-25 & \href{https://youtube.com/watch?v=EWFiHkhUrkU}{\texttt{EWFiHkhUrkU}} \\
約拿書  & \hyperref[sec:vU00DNbRAss]{經緯線 [釋經視頻] EP7 約拿書 - 研經資料推介} & 2020-11-27 & \href{https://youtube.com/watch?v=vU00DNbRAss}{\texttt{vU00DNbRAss}} \\
  & \hyperref[sec:DlruT1gd8EY]{浩劫重生 【第7集:流亡被擄的以西結】 陳崇基牧師 - 主日講道視頻系列 7} & 2020-12-02 & \href{https://youtube.com/watch?v=DlruT1gd8EY}{\texttt{DlruT1gd8EY}} \\
  & \hyperref[sec:h_zZvugRB1g]{浩劫重生 【第8集:亂世中的耶利米】 陳崇基牧師 - 主日講道視頻系列 8} & 2020-12-07 & \href{https://youtube.com/watch?v=h_zZvugRB1g}{\texttt{h\_zZvugRB1g}} \\
路加福音  & \hyperref[sec:wsAyaNl6nbo]{疫年聖誕信息 (路加福音二章) - 陳崇基牧師} & 2020-12-26 & \href{https://youtube.com/watch?v=wsAyaNl6nbo}{\texttt{wsAyaNl6nbo}} \\
  & \hyperref[sec:eJV9zLgjq90]{作神的子民 - 再讀歷代志 (直播重溫) - 陳崇基牧師} & 2020-12-29 & \href{https://youtube.com/watch?v=eJV9zLgjq90}{\texttt{eJV9zLgjq90}} \\
  & \hyperref[sec:n3sfAz_OcWA]{舊約網上行 - 上帝的足跡:EP1 起初} & 2021-01-12 & \href{https://youtube.com/watch?v=n3sfAz_OcWA}{\texttt{n3sfAz\_OcWA}} \\
  & \hyperref[sec:lVpUws1JssE]{神的名字6\_耶和華賜平安 Yahweh Shalom (靈思牧禱)} & 2021-01-13 & \href{https://youtube.com/watch?v=lVpUws1JssE}{\texttt{lVpUws1JssE}} \\
  & \hyperref[sec:u79BpJLNaZg]{經緯線 - 耶穌的比喻 (EP1 撒種的比喻)} & 2021-01-13 & \href{https://youtube.com/watch?v=u79BpJLNaZg}{\texttt{u79BpJLNaZg}} \\
  & \hyperref[sec:xKw__Td71oI]{舊約網上行 - 上帝的足跡:EP2 叛變} & 2021-01-19 & \href{https://youtube.com/watch?v=xKw--Td71oI}{\texttt{xKw--Td71oI}} \\
  & \hyperref[sec:ke055dTE6Oc]{經緯線 - 耶穌的比喻 (EP2 麥田稗子的比喻)} & 2021-01-29 & \href{https://youtube.com/watch?v=ke055dTE6Oc}{\texttt{ke055dTE6Oc}} \\
  & \hyperref[sec:1qsacx0Vy78]{舊約網上行 - 上帝的足跡:EP3 文明} & 2021-02-03 & \href{https://youtube.com/watch?v=1qsacx0Vy78}{\texttt{1qsacx0Vy78}} \\
  & \hyperref[sec:tF7ESjGDOjk]{經緯線 - 耶穌的比喻 (EP3 葡萄園工人的比喻)} & 2021-02-12 & \href{https://youtube.com/watch?v=tF7ESjGDOjk}{\texttt{tF7ESjGDOjk}} \\
 6:30-34 & \hyperref[sec:3j_OvRyfEuc]{Kingdom Values XL Matthew 6\_30-34} & 2021-02-15 & \href{https://youtube.com/watch?v=3j-OvRyfEuc}{\texttt{3j-OvRyfEuc}} \\
  & \hyperref[sec:OwmCWGNsDAQ]{經緯線 - 耶穌的比喻 (EP4 十個童女的比喻)} & 2021-02-16 & \href{https://youtube.com/watch?v=OwmCWGNsDAQ}{\texttt{OwmCWGNsDAQ}} \\
  & \hyperref[sec:UecWhHGM39A]{Kingdom Values XLI - Judgement} & 2021-02-21 & \href{https://youtube.com/watch?v=UecWhHGM39A}{\texttt{UecWhHGM39A}} \\
  & \hyperref[sec:bJj9ij22kAs]{經緯線 - 耶穌的比喻 (EP5 才幹的比喻)} & 2021-02-24 & \href{https://youtube.com/watch?v=bJj9ij22kAs}{\texttt{bJj9ij22kAs}} \\
 7:1-5 & \hyperref[sec:jf7dAWwLNe0]{Kingdom Values XLII: Matthew 7\_1-5} & 2021-03-01 & \href{https://youtube.com/watch?v=jf7dAWwLNe0}{\texttt{jf7dAWwLNe0}} \\
  & \hyperref[sec:QRUeCLAvrYk]{舊約網上行 - 上帝的足跡:EP4 方舟} & 2021-03-02 & \href{https://youtube.com/watch?v=QRUeCLAvrYk}{\texttt{QRUeCLAvrYk}} \\
  & \hyperref[sec:rLZ_I4dgAlo]{經緯線 - 耶穌的比喻( EP6 研經資料推介)} & 2021-03-06 & \href{https://youtube.com/watch?v=rLZ-I4dgAlo}{\texttt{rLZ-I4dgAlo}} \\
  & \hyperref[sec:w_y0M2H_ndU]{舊約網上行 - 上帝的足跡:EP5 巴別} & 2021-03-09 & \href{https://youtube.com/watch?v=w_y0M2H_ndU}{\texttt{w\_y0M2H\_ndU}} \\
  & \hyperref[sec:hPoAv0AxRzM]{舊約網上行 - 上帝的足跡:EP6 離別吾珥} & 2021-03-16 & \href{https://youtube.com/watch?v=hPoAv0AxRzM}{\texttt{hPoAv0AxRzM}} \\
  & \hyperref[sec:uhU9sAnkAzo]{考古探秘 (發現新死海古卷) -  陳崇基牧師 Mark Sir (2021.03.21)} & 2021-03-22 & \href{https://youtube.com/watch?v=uhU9sAnkAzo}{\texttt{uhU9sAnkAzo}} \\
  & \hyperref[sec:neXNn__inUY]{舊約網上行 - 上帝的足跡:EP7 試煉} & 2021-03-23 & \href{https://youtube.com/watch?v=neXNn--inUY}{\texttt{neXNn--inUY}} \\
  & \hyperref[sec:MgtgOAT6lXs]{考古探秘 (聖墓堂) -  陳崇基牧師 Mark Sir (2021.03.31)} & 2021-04-01 & \href{https://youtube.com/watch?v=MgtgOAT6lXs}{\texttt{MgtgOAT6lXs}} \\
出埃及記  & \hyperref[sec:_gtM1xmdWw0]{考古探秘 (蘇伊士運河與出埃及記) -  陳崇基牧師 Mark Sir (2021.04.10)} & 2021-04-11 & \href{https://youtube.com/watch?v=_gtM1xmdWw0}{\texttt{\_gtM1xmdWw0}} \\
  & \hyperref[sec:FvI2UPC3nfc]{Kingdom Values XLVII: Sermon on the Mount} & 2021-04-17 & \href{https://youtube.com/watch?v=FvI2UPC3nfc}{\texttt{FvI2UPC3nfc}} \\
  & \hyperref[sec:U5GUh9ATRrs]{考古探秘 (法老黃金遊行) - 陳崇基牧師 Mark Sir (2021.04.18)} & 2021-04-18 & \href{https://youtube.com/watch?v=U5GUh9ATRrs}{\texttt{U5GUh9ATRrs}} \\
  & \hyperref[sec:Ks1q0nCc4X4]{舊約網上行 - 上帝的足跡:EP8 亞伯拉罕與羅得} & 2021-04-20 & \href{https://youtube.com/watch?v=Ks1q0nCc4X4}{\texttt{Ks1q0nCc4X4}} \\
  & \hyperref[sec:QTJ66JH0gPw]{舊約網上行 - 上帝的足跡: EP9 凱旋回歸} & 2021-05-03 & \href{https://youtube.com/watch?v=QTJ66JH0gPw}{\texttt{QTJ66JH0gPw}} \\
 7:13-14 & \hyperref[sec:afSlbD_al4Y]{Kingdom Values: Sermon on the Mount L Matthew 7\_13-14} & 2021-05-18 & \href{https://youtube.com/watch?v=afSlbD_al4Y}{\texttt{afSlbD\_al4Y}} \\
  & \hyperref[sec:_taiMeC2wK0]{舊約網上行 - 上帝的足跡: EP10 等待之約} & 2021-05-30 & \href{https://youtube.com/watch?v=_taiMeC2wK0}{\texttt{\_taiMeC2wK0}} \\
  & \hyperref[sec:17rNRRMvNT8]{舊約網上行 - 上帝的足跡: EP11 聖約的神} & 2021-06-21 & \href{https://youtube.com/watch?v=17rNRRMvNT8}{\texttt{17rNRRMvNT8}} \\
  & \hyperref[sec:hPfz5WlNiPQ]{Sam Johnson Kingdom Values LV} & 2021-06-26 & \href{https://youtube.com/watch?v=hPfz5WlNiPQ}{\texttt{hPfz5WlNiPQ}} \\
  & \hyperref[sec:p69k7rhhWPQ]{God's Goodness Episode 1 Gratefulness Towards God} & 2021-07-04 & \href{https://youtube.com/watch?v=p69k7rhhWPQ}{\texttt{p69k7rhhWPQ}} \\
  & \hyperref[sec:hd_nz1Oh6bM]{考古探秘 (耶路巴力銘文) - 陳崇基牧師 (2021.07.14)} & 2021-07-17 & \href{https://youtube.com/watch?v=hd-nz1Oh6bM}{\texttt{hd-nz1Oh6bM}} \\
  & \hyperref[sec:yCNlGPEoTik]{舊約網上行 - 上帝的足跡: EP13 所多瑪 - 審判與恩典} & 2021-07-31 & \href{https://youtube.com/watch?v=yCNlGPEoTik}{\texttt{yCNlGPEoTik}} \\
  & \hyperref[sec:rJBi5QVSDGc]{舊約網上行 - 上帝的足跡: EP14 守約者} & 2021-08-09 & \href{https://youtube.com/watch?v=rJBi5QVSDGc}{\texttt{rJBi5QVSDGc}} \\
  & \hyperref[sec:PyoQq9QIPJA]{Mapping the Bible  - 1 Geography of Faith (Pastor Mark Chan)} & 2021-08-13 & \href{https://youtube.com/watch?v=PyoQq9QIPJA}{\texttt{PyoQq9QIPJA}} \\
  & \hyperref[sec:x_0xBEFerwg]{舊約網上行 - 上帝的足跡: EP15 最後試煉} & 2021-09-01 & \href{https://youtube.com/watch?v=x_0xBEFerwg}{\texttt{x\_0xBEFerwg}} \\
  & \hyperref[sec:bgfQNbxQms4]{2021.08.29  God's Goodness - Pastor Sam Johnson} & 2021-09-06 & \href{https://youtube.com/watch?v=bgfQNbxQms4}{\texttt{bgfQNbxQms4}} \\
 25:31-36 & \hyperref[sec:XIq2yvbK8oo]{Welcoming Jesus in Disguise (Matthew 25\_31-36) - Keith Draper} & 2021-09-13 & \href{https://youtube.com/watch?v=XIq2yvbK8oo}{\texttt{XIq2yvbK8oo}} \\
  & \hyperref[sec:9svuWXYcy2M]{舊約網上行 - 上帝的足跡: EP16 撒拉的埋葬} & 2021-09-13 & \href{https://youtube.com/watch?v=9svuWXYcy2M}{\texttt{9svuWXYcy2M}} \\
  & \hyperref[sec:pfpxXxJmbDw]{Anointed by God - Part 1 (1 Samuel 17) - Tommy Lee} & 2021-09-21 & \href{https://youtube.com/watch?v=pfpxXxJmbDw}{\texttt{pfpxXxJmbDw}} \\
  & \hyperref[sec:JJV1ojc4_5E]{舊約網上行 - 上帝的足跡: EP17 以撒與利百加} & 2021-09-25 & \href{https://youtube.com/watch?v=JJV1ojc4_5E}{\texttt{JJV1ojc4\_5E}} \\
 3:1-21 & \hyperref[sec:YtsGA4E1iho]{Rejoice in Knowing Christ (Philippians 3\_1-21, 4\_4) - Jonathan Barney} & 2021-11-01 & \href{https://youtube.com/watch?v=YtsGA4E1iho}{\texttt{YtsGA4E1iho}} \\
  & \hyperref[sec:MkL9xiQNZJM]{We are Jonah (Jonah 4) - Brock Lockenour} & 2021-11-08 & \href{https://youtube.com/watch?v=MkL9xiQNZJM}{\texttt{MkL9xiQNZJM}} \\
  & \hyperref[sec:GChzvOXTH_M]{Anointed by God - Part 7 (1 Samuel 17) - Tommy Lee} & 2021-11-16 & \href{https://youtube.com/watch?v=GChzvOXTH-M}{\texttt{GChzvOXTH-M}} \\
詩篇 46:1-11 & \hyperref[sec:eQLpWVai5F0]{What's up (詩篇46\_1-11) - 陳德晶傳道 [HeArea心靈手作室同工]} & 2021-11-17 & \href{https://youtube.com/watch?v=eQLpWVai5F0}{\texttt{eQLpWVai5F0}} \\
路得記  & \hyperref[sec:I4XJBn_dGH4]{亂世生存手冊 (路得記) - 朱莉傳道 [基督中心堂佐敦堂 ]} & 2021-11-17 & \href{https://youtube.com/watch?v=I4XJBn_dGH4}{\texttt{I4XJBn\_dGH4}} \\
  & \hyperref[sec:oiXSurvMOKc]{舊約網上行 - 上帝的足跡: EP19 以掃與雅各} & 2021-11-18 & \href{https://youtube.com/watch?v=oiXSurvMOKc}{\texttt{oiXSurvMOKc}} \\
 10:19-25 & \hyperref[sec:cKHN4UF2tsA]{The Church is the gathered (Hebrews 10\_19-25) - Rev. John Yi} & 2021-11-23 & \href{https://youtube.com/watch?v=cKHN4UF2tsA}{\texttt{cKHN4UF2tsA}} \\
 7:1-23 & \hyperref[sec:BxAa7Ok6gjk]{Keep your heart from the evil more than from COVID (Mark 7\_1-23) - Rev. Jaeuk Jeong} & 2021-11-29 & \href{https://youtube.com/watch?v=BxAa7Ok6gjk}{\texttt{BxAa7Ok6gjk}} \\
以賽亞書 7:13-17 & \hyperref[sec:CNrSX78PCg8]{以馬內利的兆頭 (以賽亞書 7\_13-17) - 陳崇基牧師} & 2021-12-07 & \href{https://youtube.com/watch?v=CNrSX78PCg8}{\texttt{CNrSX78PCg8}} \\
 1:18-25 & \hyperref[sec:e3EJDUy_J3k]{The birth of Jesus Christ (Matthew 1\_18-25) - Rev. Mark Chan} & 2021-12-13 & \href{https://youtube.com/watch?v=e3EJDUy-J3k}{\texttt{e3EJDUy-J3k}} \\
 1:18-25 & \hyperref[sec:lTIqtT1KU5c]{耶穌的誕生 (太 1\_18-25) - 陳崇基牧師} & 2021-12-13 & \href{https://youtube.com/watch?v=lTIqtT1KU5c}{\texttt{lTIqtT1KU5c}} \\
  & \hyperref[sec:PUn4RKNOgHU]{2021 年十大聖經考古發現 (陳崇基牧師)} & 2022-01-02 & \href{https://youtube.com/watch?v=PUn4RKNOgHU}{\texttt{PUn4RKNOgHU}} \\
 1:6-12 & \hyperref[sec:6HyiM6YA71c]{Following God in Tough Times (Job 1\_6-12) - Rev. Mark Chan} & 2022-01-07 & \href{https://youtube.com/watch?v=6HyiM6YA71c}{\texttt{6HyiM6YA71c}} \\
傳道書  & \hyperref[sec:deow9FxV6oI]{逆要作門徒 (2) — 再思傳道書 (陳崇基牧師) 2022.01.09} & 2022-01-11 & \href{https://youtube.com/watch?v=deow9FxV6oI}{\texttt{deow9FxV6oI}} \\
 1:12-14 & \hyperref[sec:qJYnzOjf_Ec]{Pray extraordinarily to be filled with the Spirit (Acts 1\_12-14, 2\_1-4) - Rev. Jaeuk Jeong} & 2022-01-12 & \href{https://youtube.com/watch?v=qJYnzOjf_Ec}{\texttt{qJYnzOjf\_Ec}} \\
創世記  & \hyperref[sec:SlU_7mFzyiw]{創世記第2章 (2.1) - 背景釋經 (陳崇基牧師)} & 2022-01-16 & \href{https://youtube.com/watch?v=SlU-7mFzyiw}{\texttt{SlU-7mFzyiw}} \\
創世記  & \hyperref[sec:l4PeWdLBiWQ]{創世記第1章 - 背景釋經 (陳崇基牧師)} & 2022-01-18 & \href{https://youtube.com/watch?v=l4PeWdLBiWQ}{\texttt{l4PeWdLBiWQ}} \\
創世記  & \hyperref[sec:rm7wwoQR_Zw]{創世記第2章 (2.2) - 背景釋經 (陳崇基牧師)} & 2022-01-25 & \href{https://youtube.com/watch?v=rm7wwoQR-Zw}{\texttt{rm7wwoQR-Zw}} \\
創世記  & \hyperref[sec:yR3tIvy_FRE]{創世記第2章 (2.3) - 背景釋經 (陳崇基牧師)} & 2022-01-25 & \href{https://youtube.com/watch?v=yR3tIvy-FRE}{\texttt{yR3tIvy-FRE}} \\
詩篇  & \hyperref[sec:oVRXAcV7BoQ]{逆要作門徒 (3) — 再思詩篇一二 (陳崇基牧師) 2022.01.23} & 2022-01-25 & \href{https://youtube.com/watch?v=oVRXAcV7BoQ}{\texttt{oVRXAcV7BoQ}} \\
 2:1-3 & \hyperref[sec:gQBhlqNH6hM]{Bring the Christ in the World (1 Kings 2\_1-3) - Tommy Lee} & 2022-01-28 & \href{https://youtube.com/watch?v=gQBhlqNH6hM}{\texttt{gQBhlqNH6hM}} \\
創世記  & \hyperref[sec:wRAUIVuDJgw]{創世記第3章 (3.1) - 背景釋經 (陳崇基牧師)} & 2022-02-01 & \href{https://youtube.com/watch?v=wRAUIVuDJgw}{\texttt{wRAUIVuDJgw}} \\
創世記  & \hyperref[sec:rx44uhcmvMM]{創世記第3章 (3.2) - 背景釋經 (陳崇基牧師)} & 2022-02-01 & \href{https://youtube.com/watch?v=rx44uhcmvMM}{\texttt{rx44uhcmvMM}} \\
箴言  & \hyperref[sec:Q90kQe4p0_k]{逆要作門徒 (4) — 再思箴言 (陳崇基牧師) 2022.01.30} & 2022-02-01 & \href{https://youtube.com/watch?v=Q90kQe4p0_k}{\texttt{Q90kQe4p0\_k}} \\
創世記  & \hyperref[sec:q8ZgN14Imcg]{創世記第4章 (4.1) - 背景釋經 [比對研究的圖解,近東宇宙觀] (陳崇基牧師)} & 2022-02-02 & \href{https://youtube.com/watch?v=q8ZgN14Imcg}{\texttt{q8ZgN14Imcg}} \\
出埃及記 12:1-14 & \hyperref[sec:xeNKKwLqznQ]{逾越年月 (出埃及記12\_1-14) -  朱莉傳道} & 2022-02-07 & \href{https://youtube.com/watch?v=xeNKKwLqznQ}{\texttt{xeNKKwLqznQ}} \\
 21:1-14 & \hyperref[sec:INsxhCoyb9U]{我想和你無所事事 (約21\_1-14) -  陳德晶傳道} & 2022-02-15 & \href{https://youtube.com/watch?v=INsxhCoyb9U}{\texttt{INsxhCoyb9U}} \\
  & \hyperref[sec:b3sjDQwWdlg]{Family Sunday School - EP 1 (Between Christmas and Easter)} & 2022-02-22 & \href{https://youtube.com/watch?v=b3sjDQwWdlg}{\texttt{b3sjDQwWdlg}} \\
創世記  & \hyperref[sec:bkEucgq83Sg]{創世記第4章 (4.2) - 背景釋經 [亞達哈斯史詩 vs 創世記] (陳崇基牧師)} & 2022-02-22 & \href{https://youtube.com/watch?v=bkEucgq83Sg}{\texttt{bkEucgq83Sg}} \\
 6:1-2 & \hyperref[sec:1enAW_IL3Aw]{First of the Seven Seals (Rev 6\_1-2) - Rev. Mark Chan} & 2022-02-28 & \href{https://youtube.com/watch?v=1enAW-IL3Aw}{\texttt{1enAW-IL3Aw}} \\
  & \hyperref[sec:X5Jubq7tdoc]{Family Sunday School - EP 2 (The Last Supper)} & 2022-03-03 & \href{https://youtube.com/watch?v=X5Jubq7tdoc}{\texttt{X5Jubq7tdoc}} \\
啟示錄 6:1-2 & \hyperref[sec:FcE9snFLM6s]{亂世啟示 (1) 第一印 白馬  (啟示錄6\_1-2) - 陳崇基牧師} & 2022-03-07 & \href{https://youtube.com/watch?v=FcE9snFLM6s}{\texttt{FcE9snFLM6s}} \\
啟示錄 6:3-4 & \hyperref[sec:jFHDnibTFuo]{亂世啟示 (2) 第二印 紅馬 (啟示錄6\_3-4) - 陳崇基牧師} & 2022-03-07 & \href{https://youtube.com/watch?v=jFHDnibTFuo}{\texttt{jFHDnibTFuo}} \\
創世記  & \hyperref[sec:shYH8FbN1Gk]{創世記第4章 (4.3) - 背景釋經 [近東爭辯詩,創四章,該隱的後代] (陳崇基牧師)} & 2022-03-08 & \href{https://youtube.com/watch?v=shYH8FbN1Gk}{\texttt{shYH8FbN1Gk}} \\
  & \hyperref[sec:hkH9QTCOPtw]{Family Sunday School - EP 3 (Jesus Gets Arrested)} & 2022-03-11 & \href{https://youtube.com/watch?v=hkH9QTCOPtw}{\texttt{hkH9QTCOPtw}} \\
 10:38-42 & \hyperref[sec:yncumdAd8iU]{Choosing What is Better (Luke 10\_38-42) - Rev. Sam Johnson} & 2022-03-14 & \href{https://youtube.com/watch?v=yncumdAd8iU}{\texttt{yncumdAd8iU}} \\
創世記  & \hyperref[sec:NyL1j_DJIlI]{創世記第4章 (4.4) - 背景釋經 (陳崇基牧師)} & 2022-03-14 & \href{https://youtube.com/watch?v=NyL1j_DJIlI}{\texttt{NyL1j\_DJIlI}} \\
馬太福音 24:1-34 & \hyperref[sec:5qb8PRS0GKE]{亂世啟示: 橄欖山講道 - 終末論 (馬太福音 24\_1-34) - 陳崇基牧師} & 2022-03-15 & \href{https://youtube.com/watch?v=5qb8PRS0GKE}{\texttt{5qb8PRS0GKE}} \\
創世記  & \hyperref[sec:RE3zuT2jZmI]{創世記第5章 (5.1) - 背景釋經 [比對研究,亞當家譜] (陳崇基牧師)} & 2022-03-15 & \href{https://youtube.com/watch?v=RE3zuT2jZmI}{\texttt{RE3zuT2jZmI}} \\
創世記  & \hyperref[sec:mEHavpOJC9o]{創世記第5章 (5.2) - 背景釋經 (陳崇基牧師)} & 2022-03-15 & \href{https://youtube.com/watch?v=mEHavpOJC9o}{\texttt{mEHavpOJC9o}} \\
創世記  & \hyperref[sec:mFPkQ1HAFUA]{創世記第5章 (5.3) - 背景釋經 (陳崇基牧師)} & 2022-03-15 & \href{https://youtube.com/watch?v=mFPkQ1HAFUA}{\texttt{mFPkQ1HAFUA}} \\
  & \hyperref[sec:ng1CRwPs7JE]{從「信仰百川」到「港風之谷」(上):網絡的言論自由與不自由} & 2022-03-16 & \href{https://youtube.com/watch?v=ng1CRwPs7JE}{\texttt{ng1CRwPs7JE}} \\
  & \hyperref[sec:B4MkD7Z7rb0]{從「信仰百川」到「港風之谷」(下):建立健康網絡信仰文化的努力} & 2022-03-16 & \href{https://youtube.com/watch?v=B4MkD7Z7rb0}{\texttt{B4MkD7Z7rb0}} \\
  & \hyperref[sec:IQh3_NuFQqI]{Family Sunday School - EP 5 (Jesus Goes to Trial)} & 2022-03-21 & \href{https://youtube.com/watch?v=IQh3_NuFQqI}{\texttt{IQh3\_NuFQqI}} \\
 6:1-17 & \hyperref[sec:XtJhkYaLq8U]{Ordinary Me (Judges 6\_1-17, 25-26) - Tommy Lee} & 2022-03-28 & \href{https://youtube.com/watch?v=XtJhkYaLq8U}{\texttt{XtJhkYaLq8U}} \\
  & \hyperref[sec:pinIq8rUt5E]{放下自己的接待(路14:1–14)- 黃嘉漢牧師} & 2022-03-28 & \href{https://youtube.com/watch?v=pinIq8rUt5E}{\texttt{pinIq8rUt5E}} \\
  & \hyperref[sec:CgJV5or0kgk]{Family Sunday School - EP 6 (The King of Jews)} & 2022-03-30 & \href{https://youtube.com/watch?v=CgJV5or0kgk}{\texttt{CgJV5or0kgk}} \\
  & \hyperref[sec:deMiI9YWwUk]{Family Sunday School - EP 7 (Good Friday)} & 2022-04-04 & \href{https://youtube.com/watch?v=deMiI9YWwUk}{\texttt{deMiI9YWwUk}} \\
 24:36-44 & \hyperref[sec:i4txNEY8_xE]{The Last Days According to Jesus 2 (Matthew 24\_36-44) - Rev. Mark Chan} & 2022-04-04 & \href{https://youtube.com/watch?v=i4txNEY8_xE}{\texttt{i4txNEY8\_xE}} \\
以西結書  & \hyperref[sec:sGka_a0cCzQ]{亂世啟示 : 歌革的預言 (以西結書 38-39 章) - 陳崇基牧師} & 2022-04-04 & \href{https://youtube.com/watch?v=sGka-a0cCzQ}{\texttt{sGka-a0cCzQ}} \\
創世記  & \hyperref[sec:jwup3nNhxBw]{創世記第6章 (6.1) [神的兒子的叛變] - 背景釋經 (陳崇基牧師)} & 2022-04-12 & \href{https://youtube.com/watch?v=jwup3nNhxBw}{\texttt{jwup3nNhxBw}} \\
  & \hyperref[sec:Gv0nluM2dc8]{「流散/流放」神學初探 (下):華人教會實況(2022-4-16)} & 2022-04-16 & \href{https://youtube.com/watch?v=Gv0nluM2dc8}{\texttt{Gv0nluM2dc8}} \\
  & \hyperref[sec:H5eDGDUsh5o]{「流散/流放」神學初探 (上):從聖經神學定義(2022-4-15)} & 2022-04-17 & \href{https://youtube.com/watch?v=H5eDGDUsh5o}{\texttt{H5eDGDUsh5o}} \\
創世記  & \hyperref[sec:17_oWOqNjfw]{創世記第6章 (6.2) - 背景釋經 [一百二十年] (陳崇基牧師)} & 2022-04-25 & \href{https://youtube.com/watch?v=17_oWOqNjfw}{\texttt{17\_oWOqNjfw}} \\
創世記  & \hyperref[sec:ezSMl3tcxvs]{創世記第6章 (6.3) - 背景釋經 [耶和華後悔?] (陳崇基牧師)} & 2022-04-25 & \href{https://youtube.com/watch?v=ezSMl3tcxvs}{\texttt{ezSMl3tcxvs}} \\
馬太福音 28:18-20 & \hyperref[sec:fBoYPwedTkU]{艱險我奮進 (馬太福音28\_18-20) - 朱莉傳道} & 2022-04-25 & \href{https://youtube.com/watch?v=fBoYPwedTkU}{\texttt{fBoYPwedTkU}} \\
創世記  & \hyperref[sec:RE1JRWJg7GI]{聖經考古學 - 創世記,洪水,方舟 (2)  -  陳崇基牧師} & 2022-04-29 & \href{https://youtube.com/watch?v=RE1JRWJg7GI}{\texttt{RE1JRWJg7GI}} \\
 4:13-18 & \hyperref[sec:2oEQOesAhNs]{亂世啟示: 空中會見主 (帖撒羅尼加前書4\_13-18) - 陳崇基牧師} & 2022-05-05 & \href{https://youtube.com/watch?v=2oEQOesAhNs}{\texttt{2oEQOesAhNs}} \\
  & \hyperref[sec:eLoBtjnY0Kg]{比對研究 2: 必要, 發展, 方法   -  陳崇基牧師} & 2022-05-05 & \href{https://youtube.com/watch?v=eLoBtjnY0Kg}{\texttt{eLoBtjnY0Kg}} \\
 1:1 & \hyperref[sec:Wp2QtRm4GIw]{Apostles' Creed \#1 (Genesis 1\_1) - Rev. Mark Chan} & 2022-05-09 & \href{https://youtube.com/watch?v=Wp2QtRm4GIw}{\texttt{Wp2QtRm4GIw}} \\
  & \hyperref[sec:io1AlESJGhY]{Apostles' Creed \#2 (Psalm 19) - Rev. Mark Chan} & 2022-05-16 & \href{https://youtube.com/watch?v=io1AlESJGhY}{\texttt{io1AlESJGhY}} \\
詩篇  & \hyperref[sec:2oUtiGUEP5Q]{防禦性悲觀性樂觀 (詩篇 27 篇) - 朱莉傳道} & 2022-05-16 & \href{https://youtube.com/watch?v=2oUtiGUEP5Q}{\texttt{2oUtiGUEP5Q}} \\
 11:1-44 & \hyperref[sec:V2NOXpx01xg]{An Emotive Lord (John 11\_1-44) - Brian Ng} & 2022-05-24 & \href{https://youtube.com/watch?v=V2NOXpx01xg}{\texttt{V2NOXpx01xg}} \\
馬太福音 1:1-17 & \hyperref[sec:ytVA464zDA8]{耶穌的騎呢家譜 (馬太福音 1\_1-17) - 何建邦傳道} & 2022-05-24 & \href{https://youtube.com/watch?v=ytVA464zDA8}{\texttt{ytVA464zDA8}} \\
創世記  & \hyperref[sec:lNtIgECCCGM]{創世記第6章 (6.4) - 背景釋經  [洪水問題 \#1] (陳崇基牧師)} & 2022-05-28 & \href{https://youtube.com/watch?v=lNtIgECCCGM}{\texttt{lNtIgECCCGM}} \\
創世記  & \hyperref[sec:0vMb7YIM6HI]{創世記第6章 (6.5) - 背景釋經  [洪水問題 \#2] (陳崇基牧師)} & 2022-05-28 & \href{https://youtube.com/watch?v=0vMb7YIM6HI}{\texttt{0vMb7YIM6HI}} \\
  & \hyperref[sec:sMa4o_PUcpQ]{比對研究 3: 必要, 發展, 方法   -  陳崇基牧師} & 2022-05-28 & \href{https://youtube.com/watch?v=sMa4o_PUcpQ}{\texttt{sMa4o\_PUcpQ}} \\
 19:7-14 & \hyperref[sec:JkOQgNW0fh8]{Apostles' Creed \#3  (Psalm 19\_7-14) - Rev. Mark Chan} & 2022-06-01 & \href{https://youtube.com/watch?v=JkOQgNW0fh8}{\texttt{JkOQgNW0fh8}} \\
啟示錄 6:5-11 & \hyperref[sec:s7F8M3fuVJY]{亂世啟示:第三至第五印 (啟示錄 6\_5-11) - 陳崇基牧師} & 2022-06-01 & \href{https://youtube.com/watch?v=s7F8M3fuVJY}{\texttt{s7F8M3fuVJY}} \\
  & \hyperref[sec:vmXTNOZZzYc]{Apostles' Creed \#4 (Jesus Christ) - Rev. Mark Chan} & 2022-06-06 & \href{https://youtube.com/watch?v=vmXTNOZZzYc}{\texttt{vmXTNOZZzYc}} \\
  & \hyperref[sec:Bxy4FSvUGwA]{Apostles' Creed \#5 (Virgin Birth) - Rev. Mark Chan} & 2022-06-21 & \href{https://youtube.com/watch?v=Bxy4FSvUGwA}{\texttt{Bxy4FSvUGwA}} \\
  & \hyperref[sec:1_MORjvbY98]{父親雅各的轉變 (創 32 章) - 陳崇基牧師} & 2022-06-21 & \href{https://youtube.com/watch?v=1-MORjvbY98}{\texttt{1-MORjvbY98}} \\
啟示錄 6:12-8:5 & \hyperref[sec:5Ocf3zsSFnw]{亂世啟示:第六至第七印 (啟示錄 6\_12-8\_5) - 陳崇基牧師} & 2022-06-24 & \href{https://youtube.com/watch?v=5Ocf3zsSFnw}{\texttt{5Ocf3zsSFnw}} \\
創世記  & \hyperref[sec:X2LCQ0cGbnk]{創世記第6章 (6.6) - 背景釋經 [方舟,再創造 1] (陳崇基牧師)} & 2022-06-24 & \href{https://youtube.com/watch?v=X2LCQ0cGbnk}{\texttt{X2LCQ0cGbnk}} \\
  & \hyperref[sec:zuvpXyLO2aE]{比對研究 4: 必要, 發展, 方法   -  陳崇基牧師} & 2022-06-24 & \href{https://youtube.com/watch?v=zuvpXyLO2aE}{\texttt{zuvpXyLO2aE}} \\
  & \hyperref[sec:HciKujZi6ec]{比對研究 5: 必要, 發展, 方法   -  陳崇基牧師} & 2022-06-24 & \href{https://youtube.com/watch?v=HciKujZi6ec}{\texttt{HciKujZi6ec}} \\
 4:25-5:2 & \hyperref[sec:v_Mfq956F90]{Relationship in the Church (Ephesians 4\_25-5\_2) - Steve Tan} & 2022-06-27 & \href{https://youtube.com/watch?v=v-Mfq956F90}{\texttt{v-Mfq956F90}} \\
 3:16-19 & \hyperref[sec:w05Jl15bCEk]{Why are we here? (John 3\_16-19) - Rev. Sam Johnson} & 2022-06-28 & \href{https://youtube.com/watch?v=w05Jl15bCEk}{\texttt{w05Jl15bCEk}} \\
約翰福音  & \hyperref[sec:GttbrJ8u_y4]{酒喇, 係咁先喇 (約翰福音 2: 1-11) - 陳德晶傳道} & 2022-06-28 & \href{https://youtube.com/watch?v=GttbrJ8u_y4}{\texttt{GttbrJ8u\_y4}} \\
希伯來書  & \hyperref[sec:ufhLD6E1n5U]{帶著信靠的順服 (希伯來書11:8-10) - 黃嘉漢牧師} & 2022-07-06 & \href{https://youtube.com/watch?v=ufhLD6E1n5U}{\texttt{ufhLD6E1n5U}} \\
  & \hyperref[sec:ndrseOs8_m4]{Apostles' Creed \#7 (Holy Spirit) - Rev. Mark Chan} & 2022-07-19 & \href{https://youtube.com/watch?v=ndrseOs8-m4}{\texttt{ndrseOs8-m4}} \\
詩篇  & \hyperref[sec:8rUfl6H7c2c]{創造之歌 (詩篇19) - 陳崇基牧師} & 2022-07-19 & \href{https://youtube.com/watch?v=8rUfl6H7c2c}{\texttt{8rUfl6H7c2c}} \\
  & \hyperref[sec:KZ1AIqUYQds]{比對研究 6: 基迦美斯史詩 \#1  -  陳崇基牧師} & 2022-07-19 & \href{https://youtube.com/watch?v=KZ1AIqUYQds}{\texttt{KZ1AIqUYQds}} \\
  & \hyperref[sec:ExG_dJbvHck]{比對研究 7: 基迦美斯史詩 \#2  -  陳崇基牧師} & 2022-07-19 & \href{https://youtube.com/watch?v=ExG_dJbvHck}{\texttt{ExG\_dJbvHck}} \\
詩篇  & \hyperref[sec:bjjiDZD_qQI]{詩前想後 - 詩篇一 (研經) - 陳崇基牧師} & 2022-07-19 & \href{https://youtube.com/watch?v=bjjiDZD_qQI}{\texttt{bjjiDZD\_qQI}} \\
 18:16-33 & \hyperref[sec:B4uBevTGH_s]{我們與義的距離 (創18\_16-33) - 朱莉傳道} & 2022-07-22 & \href{https://youtube.com/watch?v=B4uBevTGH_s}{\texttt{B4uBevTGH\_s}} \\
  & \hyperref[sec:QzV41v7JIro]{Apostles' Creed \#8 (Church and Eternal Life) - Rev. Mark Chan} & 2022-08-02 & \href{https://youtube.com/watch?v=QzV41v7JIro}{\texttt{QzV41v7JIro}} \\
  & \hyperref[sec:BAdEYU5b_F0]{群體之美 (詩 133) - 陳崇基牧師} & 2022-08-02 & \href{https://youtube.com/watch?v=BAdEYU5b_F0}{\texttt{BAdEYU5b\_F0}} \\
約翰福音 4:43-54 & \hyperref[sec:pMRJMgqJITI]{信你就奇 (約翰福音 4\_43-54) - 陳德晶傳道} & 2022-08-10 & \href{https://youtube.com/watch?v=pMRJMgqJITI}{\texttt{pMRJMgqJITI}} \\
  & \hyperref[sec:KfYBQZvmZ4Y]{Apostles' Creed \#9 (Life Eternal) - Rev. Mark Chan} & 2022-08-16 & \href{https://youtube.com/watch?v=KfYBQZvmZ4Y}{\texttt{KfYBQZvmZ4Y}} \\
  & \hyperref[sec:AlQ26GIFgJs]{Psalms for Every Season (Psalm 63) - Rev. Mark Chan} & 2022-08-23 & \href{https://youtube.com/watch?v=AlQ26GIFgJs}{\texttt{AlQ26GIFgJs}} \\
  & \hyperref[sec:DvMyIFbhzjs]{曠野的記憶 (詩 63) - 陳崇基牧師} & 2022-08-23 & \href{https://youtube.com/watch?v=DvMyIFbhzjs}{\texttt{DvMyIFbhzjs}} \\
詩篇  & \hyperref[sec:X4wQCaypggM]{聖經地理 EP2 - 地理與釋經 - 詩篇125-1 (陳崇基牧師)} & 2022-08-23 & \href{https://youtube.com/watch?v=X4wQCaypggM}{\texttt{X4wQCaypggM}} \\
詩篇  & \hyperref[sec:b63slc4ANbo]{聖經地理 EP4 - 地理與釋經 - 詩篇125-3 (陳崇基牧師)} & 2022-08-23 & \href{https://youtube.com/watch?v=b63slc4ANbo}{\texttt{b63slc4ANbo}} \\
  & \hyperref[sec:USIUOn0X4jM]{聖經地理 EP5 - 地理與釋經 - 中間之地-1 (陳崇基牧師)} & 2022-08-25 & \href{https://youtube.com/watch?v=USIUOn0X4jM}{\texttt{USIUOn0X4jM}} \\
  & \hyperref[sec:dVqF8V7wdDo]{聖經地理 EP6 - 地理與釋經 - 中間之地-2 (陳崇基牧師)} & 2022-08-25 & \href{https://youtube.com/watch?v=dVqF8V7wdDo}{\texttt{dVqF8V7wdDo}} \\
 11:19-26 & \hyperref[sec:PjaM3O4lbt8]{The beauty of ordinary faithfulness (Acts 11\_19-26) - Rev. Sam Johnson} & 2022-08-30 & \href{https://youtube.com/watch?v=PjaM3O4lbt8}{\texttt{PjaM3O4lbt8}} \\
  & \hyperref[sec:f_Inn31K7Tc]{上行之詩 (詩 122) - 陳崇基牧師} & 2022-08-30 & \href{https://youtube.com/watch?v=f-Inn31K7Tc}{\texttt{f-Inn31K7Tc}} \\
 3:16 & \hyperref[sec:FkGSJO7nqgY]{Hang Tight (John 3\_16) - Joyce Tam} & 2022-09-05 & \href{https://youtube.com/watch?v=FkGSJO7nqgY}{\texttt{FkGSJO7nqgY}} \\
何西阿書 5:15-6:6 & \hyperref[sec:ahiCSxfoOC8]{主愛又如何? (何西阿書 5\_15-6\_6) - 吳君才牧師} & 2022-09-05 & \href{https://youtube.com/watch?v=ahiCSxfoOC8}{\texttt{ahiCSxfoOC8}} \\
 20:20 & \hyperref[sec:T9LnkydmX7E]{What Does God Really Care? (Reflection on Hezekiah's Achievements - 2 Kings 20\_20)} & 2022-09-13 & \href{https://youtube.com/watch?v=T9LnkydmX7E}{\texttt{T9LnkydmX7E}} \\
  & \hyperref[sec:aMelPey4yl8]{佳節倍思恩(一): 住棚節(上) - 陳崇基牧師} & 2022-09-24 & \href{https://youtube.com/watch?v=aMelPey4yl8}{\texttt{aMelPey4yl8}} \\
  & \hyperref[sec:CndBkmupUxY]{The role of Nations \ and  Neighbors in a growing church (Luke 10) - Keith Draper} & 2022-09-30 & \href{https://youtube.com/watch?v=CndBkmupUxY}{\texttt{CndBkmupUxY}} \\
  & \hyperref[sec:IDFYoAmcEPE]{Faith and time: Feast of Tabernacle 1 (Rev. Mark Chan)} & 2022-10-08 & \href{https://youtube.com/watch?v=IDFYoAmcEPE}{\texttt{IDFYoAmcEPE}} \\
 5:1-5 & \hyperref[sec:eW87A4hZtZA]{算得上的人生 (創5\_1-5, 21-32) - 朱莉傳道} & 2022-10-10 & \href{https://youtube.com/watch?v=eW87A4hZtZA}{\texttt{eW87A4hZtZA}} \\
馬太福音 4:1-11 & \hyperref[sec:W7L4nTwMQOU]{曠野的吶喊 (馬太福音4\_1-11) - 張麗仙傳道} & 2022-10-12 & \href{https://youtube.com/watch?v=W7L4nTwMQOU}{\texttt{W7L4nTwMQOU}} \\
 15:1-19 & \hyperref[sec:yX0mME__Hqg]{Lightening Our Yokes (Acts 15\_1-19) - Brian Ng} & 2022-10-19 & \href{https://youtube.com/watch?v=yX0mME--Hqg}{\texttt{yX0mME--Hqg}} \\
  & \hyperref[sec:R_BYGFCmbDs]{佳節倍思恩(三): 安息日  - 陳崇基牧師} & 2022-10-19 & \href{https://youtube.com/watch?v=R_BYGFCmbDs}{\texttt{R\_BYGFCmbDs}} \\
  & \hyperref[sec:DueqCDH4G_Q]{聖經地理 EP11 - 地理與釋經 - 中央山地 - 1 (陳崇基牧師)} & 2022-10-21 & \href{https://youtube.com/watch?v=DueqCDH4G-Q}{\texttt{DueqCDH4G-Q}} \\
  & \hyperref[sec:Pgm9dzwLc3c]{聖經地理 EP12 - 地理與釋經 - 中央山地 - 2 (陳崇基牧師)} & 2022-10-21 & \href{https://youtube.com/watch?v=Pgm9dzwLc3c}{\texttt{Pgm9dzwLc3c}} \\
  & \hyperref[sec:26WK5R9cMYM]{宣教路 - Steve Tan} & 2022-10-25 & \href{https://youtube.com/watch?v=26WK5R9cMYM}{\texttt{26WK5R9cMYM}} \\
  & \hyperref[sec:mXwPLGGYFwA]{聖經地理 EP7 - 地理與釋經 - 交通-1 (陳崇基牧師)} & 2022-11-01 & \href{https://youtube.com/watch?v=mXwPLGGYFwA}{\texttt{mXwPLGGYFwA}} \\
  & \hyperref[sec:PWN_xBip_GQ]{聖經地理 EP8 - 地理與釋經 - 交通-2 (陳崇基牧師)} & 2022-11-08 & \href{https://youtube.com/watch?v=PWN-xBip-GQ}{\texttt{PWN-xBip-GQ}} \\
  & \hyperref[sec:wlVkY6Ioxpo]{佳節倍思恩(四): 逾越節  - 陳崇基牧師} & 2022-11-10 & \href{https://youtube.com/watch?v=wlVkY6Ioxpo}{\texttt{wlVkY6Ioxpo}} \\
  & \hyperref[sec:PVQ3X0QwXNk]{Faith and time: Passover - Rev. Mark Chan} & 2022-11-13 & \href{https://youtube.com/watch?v=PVQ3X0QwXNk}{\texttt{PVQ3X0QwXNk}} \\
  & \hyperref[sec:E_n46N8Pd_8]{The Greatest Commandment (Matthew 22: 34-40) - Brock Lockenour} & 2022-11-13 & \href{https://youtube.com/watch?v=E_n46N8Pd_8}{\texttt{E\_n46N8Pd\_8}} \\
 5:1-9 & \hyperref[sec:042xzVMibbM]{一個人在 (約 5\_1-9) - 陳德晶傳道} & 2022-11-14 & \href{https://youtube.com/watch?v=042xzVMibbM}{\texttt{042xzVMibbM}} \\
  & \hyperref[sec:i2GG1qfPho8]{聖經地理 EP9 - 地理與釋經 - 地景 (陳崇基牧師)} & 2022-11-15 & \href{https://youtube.com/watch?v=i2GG1qfPho8}{\texttt{i2GG1qfPho8}} \\
  & \hyperref[sec:JhyniyhyQUs]{Cultivating Gratitude (Psalm 100) - Rev. Sam Johnson} & 2022-11-16 & \href{https://youtube.com/watch?v=JhyniyhyQUs}{\texttt{JhyniyhyQUs}} \\
  & \hyperref[sec:xHuMbJI2NdE]{Faith and Time: Shavuot - Rev. Mark Chan} & 2022-11-28 & \href{https://youtube.com/watch?v=xHuMbJI2NdE}{\texttt{xHuMbJI2NdE}} \\
  & \hyperref[sec:LC2O0QPKB1U]{佳節倍思恩(五): 五旬節  - 陳崇基牧師} & 2022-11-28 & \href{https://youtube.com/watch?v=LC2O0QPKB1U}{\texttt{LC2O0QPKB1U}} \\
  & \hyperref[sec:SrXfS_omV5U]{佳節倍思恩(四): 逾越節  - 陳崇基牧師} & 2022-12-06 & \href{https://youtube.com/watch?v=SrXfS_omV5U}{\texttt{SrXfS\_omV5U}} \\
  & \hyperref[sec:kSJ3Fo7YGUs]{Reason For The Season - Rev. Mark Chan} & 2022-12-20 & \href{https://youtube.com/watch?v=kSJ3Fo7YGUs}{\texttt{kSJ3Fo7YGUs}} \\
  & \hyperref[sec:du2XrTgHsNA]{聖誕的原因 - 陳崇基牧師} & 2022-12-20 & \href{https://youtube.com/watch?v=du2XrTgHsNA}{\texttt{du2XrTgHsNA}} \\
  & \hyperref[sec:LoAxQ1cr1f0]{A Child is Born - Rev. Mark Chan  有一嬰孩 - 陳崇基牧師} & 2022-12-26 & \href{https://youtube.com/watch?v=LoAxQ1cr1f0}{\texttt{LoAxQ1cr1f0}} \\
  & \hyperref[sec:lBh3rgPDiTo]{6/2/19 主日講道\_耶穌天國的小故事-4\_寶貝與珍珠(陳崇基牧師)} & 2023-01-02 & \href{https://youtube.com/watch?v=lBh3rgPDiTo}{\texttt{lBh3rgPDiTo}} \\
撒母耳記下 24:1-25 & \hyperref[sec:j1mkZMKAJog]{蓋棺定論(撒母耳記下24\_1-25) - 朱莉傳道} & 2023-01-02 & \href{https://youtube.com/watch?v=j1mkZMKAJog}{\texttt{j1mkZMKAJog}} \\
 13:3-5 & \hyperref[sec:ZVIv_jPkXmQ]{Stewarding the Gift of Power (John 13\_3-5, 12-16) - Tommy Lee} & 2023-01-10 & \href{https://youtube.com/watch?v=ZVIv_jPkXmQ}{\texttt{ZVIv\_jPkXmQ}} \\
雅各書 4:13-17 & \hyperref[sec:DsrtebXqSNI]{計劃明天(雅各書4\_13-17)- 陳崇基牧師} & 2023-01-17 & \href{https://youtube.com/watch?v=DsrtebXqSNI}{\texttt{DsrtebXqSNI}} \\
  & \hyperref[sec:daNqcFTl4KQ]{Becoming A Healthier Church 1 (Ephesians) - Rev. Mark Chan} & 2023-01-24 & \href{https://youtube.com/watch?v=daNqcFTl4KQ}{\texttt{daNqcFTl4KQ}} \\
以弗所書  & \hyperref[sec:hnLQJ3HSDsE]{同建健康的教會 1 (以弗所書 5) - 陳崇基牧師} & 2023-01-26 & \href{https://youtube.com/watch?v=hnLQJ3HSDsE}{\texttt{hnLQJ3HSDsE}} \\
彌迦書 4:1-5 & \hyperref[sec:5LAc82ueC_Y]{假使世界原來不像你預期(彌迦書4\_1-5)- 陳智祥傳道} & 2023-01-30 & \href{https://youtube.com/watch?v=5LAc82ueC-Y}{\texttt{5LAc82ueC-Y}} \\
  & \hyperref[sec:NqYX_zjUCDg]{Building A Healthier Church 3: United for Mission - Rev. Mark Chan} & 2023-02-06 & \href{https://youtube.com/watch?v=NqYX-zjUCDg}{\texttt{NqYX-zjUCDg}} \\
約翰福音 2:1-11 & \hyperref[sec:Cv_wNlk__oQ]{你可以活得更豐富 (約翰福音2\_1-11)- 程子信牧師} & 2023-02-07 & \href{https://youtube.com/watch?v=Cv_wNlk-_oQ}{\texttt{Cv\_wNlk-\_oQ}} \\
 17:1-7 & \hyperref[sec:qMFmuMv64OQ]{Before Whom I Stand (1 King 17\_1-7) - Tommy Lee} & 2023-02-13 & \href{https://youtube.com/watch?v=qMFmuMv64OQ}{\texttt{qMFmuMv64OQ}} \\
箴言 3:5-12 & \hyperref[sec:8B7OdFQev9A]{我兒... (箴言3\_5-12)- 譚文倩莊傳道} & 2023-02-13 & \href{https://youtube.com/watch?v=8B7OdFQev9A}{\texttt{8B7OdFQev9A}} \\
  & \hyperref[sec:vHt7tPIvb_w]{Building A Healthier Church 4: Spiritual Gifts - Rev. Mark Chan} & 2023-02-20 & \href{https://youtube.com/watch?v=vHt7tPIvb-w}{\texttt{vHt7tPIvb-w}} \\
  & \hyperref[sec:u_InczBESQY]{同建健康教會 4:恩賜 - 陳崇基牧師} & 2023-02-20 & \href{https://youtube.com/watch?v=u_InczBESQY}{\texttt{u\_InczBESQY}} \\
  & \hyperref[sec:zKpR0b97ynU]{比對研究 8: 巴力宗教 1 - 陳崇基牧師} & 2023-02-24 & \href{https://youtube.com/watch?v=zKpR0b97ynU}{\texttt{zKpR0b97ynU}} \\
  & \hyperref[sec:9YLRCR_mn34]{Art, Beauty, and God: Gen 1-2 - Rebecca Wu} & 2023-02-28 & \href{https://youtube.com/watch?v=9YLRCR-mn34}{\texttt{9YLRCR-mn34}} \\
  & \hyperref[sec:5SRk2JnEBBc]{同建健康教會 5: 愛與恩賜 - 陳崇基牧師} & 2023-02-28 & \href{https://youtube.com/watch?v=5SRk2JnEBBc}{\texttt{5SRk2JnEBBc}} \\
  & \hyperref[sec:CFdO0Ddf8oE]{Becoming A Healthier Church 5: Way of Love - Rev. Mark Chan} & 2023-03-07 & \href{https://youtube.com/watch?v=CFdO0Ddf8oE}{\texttt{CFdO0Ddf8oE}} \\
使徒行傳 26:15-20 & \hyperref[sec:z2UTnfWH0W8]{非常任務 (使徒行傳26\_15-20) - 黃茂林傳道} & 2023-03-07 & \href{https://youtube.com/watch?v=z2UTnfWH0W8}{\texttt{z2UTnfWH0W8}} \\
  & \hyperref[sec:bEexjAddGwM]{比對研究 9: 巴力宗教 2 - 陳崇基牧師} & 2023-03-11 & \href{https://youtube.com/watch?v=bEexjAddGwM}{\texttt{bEexjAddGwM}} \\
  & \hyperref[sec:2q17d1G4vqo]{Becoming A Healthier Church 6: 1 Cor 14 - Sound Teaching  - Rev. Mark Chan} & 2023-03-14 & \href{https://youtube.com/watch?v=2q17d1G4vqo}{\texttt{2q17d1G4vqo}} \\
 12:1-11 & \hyperref[sec:lL2aGhSKVWQ]{Final Days of Jesus: Bethany (John 12\_1-11) - Rev. Mark Chan} & 2023-03-20 & \href{https://youtube.com/watch?v=lL2aGhSKVWQ}{\texttt{lL2aGhSKVWQ}} \\
約翰福音 12:1-11 & \hyperref[sec:zbYaP2QF_GY]{耶穌最後的日子:伯大尼  (約翰福音12\_1-11) - 陳崇基牧師} & 2023-03-20 & \href{https://youtube.com/watch?v=zbYaP2QF-GY}{\texttt{zbYaP2QF-GY}} \\
  & \hyperref[sec:5_8ZAI7O6Mc]{比對研究 10: 巴力宗教 3 - 陳崇基牧師} & 2023-03-25 & \href{https://youtube.com/watch?v=5-8ZAI7O6Mc}{\texttt{5-8ZAI7O6Mc}} \\
 13:1-15 & \hyperref[sec:WPZJt_sBFNI]{Final Days of Jesus: Example of The Cross (John 13\_1-15) - Rev. Mark Chan} & 2023-03-27 & \href{https://youtube.com/watch?v=WPZJt_sBFNI}{\texttt{WPZJt\_sBFNI}} \\
以弗所書 4:11-16 & \hyperref[sec:9Bv9h74L7iM]{我健康,成長嗎?(以弗所書4\_11-16) - 吳君才牧師} & 2023-03-27 & \href{https://youtube.com/watch?v=9Bv9h74L7iM}{\texttt{9Bv9h74L7iM}} \\
  & \hyperref[sec:k6VIyucQ7dI]{UV 匯客室- 朱莉傳道} & 2023-03-30 & \href{https://youtube.com/watch?v=k6VIyucQ7dI}{\texttt{k6VIyucQ7dI}} \\
 17:8-16 & \hyperref[sec:hADCTqthLUc]{Spent And Empty (1 King 17\_8-16) - Tommy Lee} & 2023-04-03 & \href{https://youtube.com/watch?v=hADCTqthLUc}{\texttt{hADCTqthLUc}} \\
申命記 8:1-20 & \hyperref[sec:QvSDeFIblPQ]{當紀念主 (申命記 8\_1-20) - 陳智祥傳道} & 2023-04-03 & \href{https://youtube.com/watch?v=QvSDeFIblPQ}{\texttt{QvSDeFIblPQ}} \\
  & \hyperref[sec:aVFFH_WArz0]{比對研究 11: 巴力宗教 4 - 陳崇基牧師} & 2023-04-08 & \href{https://youtube.com/watch?v=aVFFH_WArz0}{\texttt{aVFFH\_WArz0}} \\
  & \hyperref[sec:_Z2VGiey2J8]{Resurrection Gospel: John 20 - Rev. Mark Chan} & 2023-04-18 & \href{https://youtube.com/watch?v=_Z2VGiey2J8}{\texttt{\_Z2VGiey2J8}} \\
約翰福音  & \hyperref[sec:PhsM03F2OIM]{復活的福音:約翰福音 20 章 - 陳崇基牧師} & 2023-04-18 & \href{https://youtube.com/watch?v=PhsM03F2OIM}{\texttt{PhsM03F2OIM}} \\
  & \hyperref[sec:SK0VEvCIu6I]{比對研究 12: 巴力宗教 5 - 陳崇基牧師} & 2023-04-22 & \href{https://youtube.com/watch?v=SK0VEvCIu6I}{\texttt{SK0VEvCIu6I}} \\
 6:33-34 & \hyperref[sec:9WOLJvd45fU]{First Things First: Matthew 6\_33-34 - John Yi} & 2023-04-24 & \href{https://youtube.com/watch?v=9WOLJvd45fU}{\texttt{9WOLJvd45fU}} \\
哥林多前書  & \hyperref[sec:CsnEAjBlXuk]{復活的福音:哥林多前書 15 章 - 陳崇基牧師} & 2023-04-24 & \href{https://youtube.com/watch?v=CsnEAjBlXuk}{\texttt{CsnEAjBlXuk}} \\
 9:1-9 & \hyperref[sec:wxXYrlF2bG8]{The Gospel of Resurrection: Act 9\_1-9 - Rev. Mark Chan} & 2023-05-02 & \href{https://youtube.com/watch?v=wxXYrlF2bG8}{\texttt{wxXYrlF2bG8}} \\
使徒行傳 9:1-9 & \hyperref[sec:upXqQwuuS0M]{復活的福音:使徒行傳 9\_1-9 - 陳崇基牧師} & 2023-05-02 & \href{https://youtube.com/watch?v=upXqQwuuS0M}{\texttt{upXqQwuuS0M}} \\
  & \hyperref[sec:ixAcYGJ3Y7Y]{比對研究 13: 巴力宗教 6 - 陳崇基牧師} & 2023-05-06 & \href{https://youtube.com/watch?v=ixAcYGJ3Y7Y}{\texttt{ixAcYGJ3Y7Y}} \\
 17:13 & \hyperref[sec:hS4BoW8e_54]{Go and Do: 1 Kings 17\_13 - Tommy Lee} & 2023-05-09 & \href{https://youtube.com/watch?v=hS4BoW8e_54}{\texttt{hS4BoW8e\_54}} \\
 1:26-38 & \hyperref[sec:UpJWRYxxqyc]{Unseen Service (Luke 1\_26-38) - Sam Johnson} & 2023-05-16 & \href{https://youtube.com/watch?v=UpJWRYxxqyc}{\texttt{UpJWRYxxqyc}} \\
申命記 4:9 & \hyperref[sec:_7dHuWNR_kw]{謹記這一口氣 (申命記 4\_9, 23) - 張麗仙傳道} & 2023-05-16 & \href{https://youtube.com/watch?v=_7dHuWNR-kw}{\texttt{\_7dHuWNR-kw}} \\
  & \hyperref[sec:GOoaQbj9lhw]{比對研究 14: 巴力宗教 7 - 陳崇基牧師} & 2023-05-20 & \href{https://youtube.com/watch?v=GOoaQbj9lhw}{\texttt{GOoaQbj9lhw}} \\
  & \hyperref[sec:i8guCKlYag0]{Beholding Your Glory - Rev. Mark Chan} & 2023-05-22 & \href{https://youtube.com/watch?v=i8guCKlYag0}{\texttt{i8guCKlYag0}} \\
  & \hyperref[sec:OQpeQNO5TLA]{看見...你的榮耀 - 陳崇基牧師} & 2023-05-22 & \href{https://youtube.com/watch?v=OQpeQNO5TLA}{\texttt{OQpeQNO5TLA}} \\
  & \hyperref[sec:jG3lPKnsQeA]{People Who Worship God (Psalms 95) - Rev. Mark Chan} & 2023-05-29 & \href{https://youtube.com/watch?v=jG3lPKnsQeA}{\texttt{jG3lPKnsQeA}} \\
詩篇  & \hyperref[sec:wLYSrSGINOw]{敬拜神的子民:詩篇95 - 陳崇基牧師} & 2023-05-29 & \href{https://youtube.com/watch?v=wLYSrSGINOw}{\texttt{wLYSrSGINOw}} \\
  & \hyperref[sec:4x5CA4Q65_w]{比對研究 15: 巴力宗教 8 - 陳崇基牧師} & 2023-06-03 & \href{https://youtube.com/watch?v=4x5CA4Q65_w}{\texttt{4x5CA4Q65\_w}} \\
 17:17-24 & \hyperref[sec:kj0zBwcuLqE]{Death in The Family: 1 King 17\_17-24 - Tommy Lee} & 2023-06-05 & \href{https://youtube.com/watch?v=kj0zBwcuLqE}{\texttt{kj0zBwcuLqE}} \\
 8:21-22 & \hyperref[sec:7Sx7wB7NbaI]{Cost of Following Jesus (Matthew 8\_21-22) - Rev. Mark Chan} & 2023-06-13 & \href{https://youtube.com/watch?v=7Sx7wB7NbaI}{\texttt{7Sx7wB7NbaI}} \\
羅馬書 12:1-2 & \hyperref[sec:TVT11XSJIWY]{生命的祭 (羅馬書 12\_1-2) - 譚文倩莊傳道} & 2023-06-13 & \href{https://youtube.com/watch?v=TVT11XSJIWY}{\texttt{TVT11XSJIWY}} \\
  & \hyperref[sec:DoxRaGGvrWQ]{比對研究 16: 巴力宗教 9 - 列王記上 18 - 陳崇基牧師} & 2023-06-17 & \href{https://youtube.com/watch?v=DoxRaGGvrWQ}{\texttt{DoxRaGGvrWQ}} \\
 1:7 & \hyperref[sec:2hbMIJp1DE8]{Wisdom and Fear of God (Proverbs 1\_7) - Rev. Mark Chan} & 2023-06-20 & \href{https://youtube.com/watch?v=2hbMIJp1DE8}{\texttt{2hbMIJp1DE8}} \\
箴言 1:7 & \hyperref[sec:F8HOtvMKDlc]{智慧與敬畏(箴言 1\_7) - 陳崇基牧師} & 2023-06-20 & \href{https://youtube.com/watch?v=F8HOtvMKDlc}{\texttt{F8HOtvMKDlc}} \\
 8:31-32 & \hyperref[sec:K5EeMHs5F9w]{Dwelling in The Truth (John 8\_31-32) - Sam Johnson} & 2023-06-27 & \href{https://youtube.com/watch?v=K5EeMHs5F9w}{\texttt{K5EeMHs5F9w}} \\
 18:1-4 & \hyperref[sec:hjFzQBCexic]{Are You Limping Through Your Faith (1King 18\_1-4, 17-40) - Tommy Lee} & 2023-07-08 & \href{https://youtube.com/watch?v=hjFzQBCexic}{\texttt{hjFzQBCexic}} \\
 12:1-2 & \hyperref[sec:NGEVacpuPjM]{The Perfect Sacrifice (Romans 12\_1-2) - Joyce Tam} & 2023-07-11 & \href{https://youtube.com/watch?v=NGEVacpuPjM}{\texttt{NGEVacpuPjM}} \\
羅馬書 12:3-8 & \hyperref[sec:bp9nD3hO7O0]{生命的賜(羅馬書 12\_3-8)- 譚文倩莊傳道} & 2023-07-11 & \href{https://youtube.com/watch?v=bp9nD3hO7O0}{\texttt{bp9nD3hO7O0}} \\
  & \hyperref[sec:hrYJaJP5Fzk]{4/1/18 復活節崇拜話劇--轉機,轉機 Easter Sunday Skit} & 2023-07-13 & \href{https://youtube.com/watch?v=hrYJaJP5Fzk}{\texttt{hrYJaJP5Fzk}} \\
 20:20 & \hyperref[sec:mbI7RGfwLmY]{Exploring the Holy Land: Hezekiah Tunnun (2 Kings 20\_20, 2 Chronicles 32\_30) -  Rev. Mark Chan} & 2023-07-18 & \href{https://youtube.com/watch?v=mbI7RGfwLmY}{\texttt{mbI7RGfwLmY}} \\
 1:1-18 & \hyperref[sec:LITCDKMXv4k]{Count It All Joy (James 1\_1-18) - Caleb Lee} & 2023-07-25 & \href{https://youtube.com/watch?v=LITCDKMXv4k}{\texttt{LITCDKMXv4k}} \\
羅馬書 12:9-21 & \hyperref[sec:2mwrzSe3J14]{生命的致:致力愛 (羅馬書 12\_9-21) - 譚文倩莊傳道} & 2023-07-25 & \href{https://youtube.com/watch?v=2mwrzSe3J14}{\texttt{2mwrzSe3J14}} \\
 1:19-27 & \hyperref[sec:qy6pJ6fRtpk]{Being A Doer, Not A Hearer (James 1\_19-27) - George Crawley} & 2023-08-01 & \href{https://youtube.com/watch?v=qy6pJ6fRtpk}{\texttt{qy6pJ6fRtpk}} \\
 6:1-5 & \hyperref[sec:LDb_e4fz_Ck]{Exploring The Land: Little Town Jericho (Joshua 6\_1-5)- Rev. Mark Chan} & 2023-08-08 & \href{https://youtube.com/watch?v=LDb_e4fz-Ck}{\texttt{LDb\_e4fz-Ck}} \\
  & \hyperref[sec:NkubOgp9wRQ]{6/9/19 主日講道:神竟然比我和祂講數 (張麗仙傳道)} & 2023-08-13 & \href{https://youtube.com/watch?v=NkubOgp9wRQ}{\texttt{NkubOgp9wRQ}} \\
 2:7 & \hyperref[sec:T08JyDUvISQ]{Learning to Love Dependence (Genesis 2\_7) - Sam Johnson} & 2023-08-14 & \href{https://youtube.com/watch?v=T08JyDUvISQ}{\texttt{T08JyDUvISQ}} \\
約翰福音 14:1-6 & \hyperref[sec:1rUpXzwI0NI]{新天新地:人生終末的再思 1 將來的父家 (約翰福音14\_1-6) - 陳崇基牧師} & 2023-08-14 & \href{https://youtube.com/watch?v=1rUpXzwI0NI}{\texttt{1rUpXzwI0NI}} \\
 2:1-13 & \hyperref[sec:p_FHxvDNGe8]{Show No Favoritism (James 2\_1-13) - Caleb Lee} & 2023-08-21 & \href{https://youtube.com/watch?v=p-FHxvDNGe8}{\texttt{p-FHxvDNGe8}} \\
啟示錄 21:1-4 & \hyperref[sec:S12GkifPujs]{新天新地:人生終末的再思 2 將來新天地 (啟示錄21\_1-4) - 陳崇基牧師} & 2023-08-21 & \href{https://youtube.com/watch?v=S12GkifPujs}{\texttt{S12GkifPujs}} \\
 14:1-7 & \hyperref[sec:mdfK7VGG0u4]{New Heaven and New Earth 1 - (John 14\_1-7) - Rev. Mark Chan} & 2023-08-29 & \href{https://youtube.com/watch?v=mdfK7VGG0u4}{\texttt{mdfK7VGG0u4}} \\
帖撒羅尼迦前書 4:13-18 & \hyperref[sec:OzeSIq8OBsY]{新天新地:人生終末的再思 3 將來的君王 (帖撒羅尼迦前書 4\_13-18) - 陳崇基牧師} & 2023-08-29 & \href{https://youtube.com/watch?v=OzeSIq8OBsY}{\texttt{OzeSIq8OBsY}} \\
 22:1-14 & \hyperref[sec:d5kiiPL31Z0]{New Heaven and New Earth: The Coming Banquet (Matthew 22\_1-14) - Rev. Mark Chan} & 2023-09-05 & \href{https://youtube.com/watch?v=d5kiiPL31Z0}{\texttt{d5kiiPL31Z0}} \\
馬太福音 22:1-14 & \hyperref[sec:ovpilkoc6io]{新天新地:人生終末的再思 4 將來的婚宴 (馬太福音 22\_1-14) - 陳崇基牧師} & 2023-09-05 & \href{https://youtube.com/watch?v=ovpilkoc6io}{\texttt{ovpilkoc6io}} \\
  & \hyperref[sec:D2gsbWPv_jE]{比對研究 17: 巴力宗教 10 - 列王記上 18 - 陳崇基牧師} & 2023-09-09 & \href{https://youtube.com/watch?v=D2gsbWPv-jE}{\texttt{D2gsbWPv-jE}} \\
 2:14-26 & \hyperref[sec:EH56cMFj_rc]{A Living Faith (James 2\_14-26) - Caleb Lee} & 2023-09-13 & \href{https://youtube.com/watch?v=EH56cMFj-rc}{\texttt{EH56cMFj-rc}} \\
尼希米記 1:1-11 & \hyperref[sec:p__HCvfoaaM]{靈命省察:四個重要的問題 (尼希米記 1\_1-11) - 張麗仙傳道} & 2023-09-13 & \href{https://youtube.com/watch?v=p_-HCvfoaaM}{\texttt{p\_-HCvfoaaM}} \\
 1:3-14 & \hyperref[sec:1Q8_M16V0ZI]{We Are Worshippers (Ephesians 1\_3-14) - Sam Johnson} & 2023-09-19 & \href{https://youtube.com/watch?v=1Q8_M16V0ZI}{\texttt{1Q8\_M16V0ZI}} \\
創世記  & \hyperref[sec:Z86XRROS_74]{上主的護理 (1) 細想約瑟的故事 1 (創世記 37) - 陳崇基牧師} & 2023-09-19 & \href{https://youtube.com/watch?v=Z86XRROS-74}{\texttt{Z86XRROS-74}} \\
  & \hyperref[sec:_iX85FoHSZc]{比對研究 18: 巴力宗教 11 - 列王記上 18 - 陳崇基牧師} & 2023-09-23 & \href{https://youtube.com/watch?v=_iX85FoHSZc}{\texttt{\_iX85FoHSZc}} \\
  & \hyperref[sec:mJHLbAHSJBg]{God's Providence 1 (Genesis 37) - Rev. Mark Chan} & 2023-09-26 & \href{https://youtube.com/watch?v=mJHLbAHSJBg}{\texttt{mJHLbAHSJBg}} \\
創世記  & \hyperref[sec:Fs4lhTNde4E]{上主的護理 (2) 細想約瑟的故事 2 (創世記 39) - 陳崇基牧師} & 2023-09-26 & \href{https://youtube.com/watch?v=Fs4lhTNde4E}{\texttt{Fs4lhTNde4E}} \\
 3:1-12 & \hyperref[sec:Jzl6PyqrviE]{Words Of Life (James 3\_1-12) - Caleb Lee} & 2023-10-03 & \href{https://youtube.com/watch?v=Jzl6PyqrviE}{\texttt{Jzl6PyqrviE}} \\
  & \hyperref[sec:VkqqimXzBU4]{弟兄中的一個最小 - 譚文倩莊傳道} & 2023-10-03 & \href{https://youtube.com/watch?v=VkqqimXzBU4}{\texttt{VkqqimXzBU4}} \\
  & \hyperref[sec:lg0Kj2q_1qU]{方舟考古問題逐一破障 (2023-10-2 鄺牧道頻,城匯網台聯播)} & 2023-10-04 & \href{https://youtube.com/watch?v=lg0Kj2q_1qU}{\texttt{lg0Kj2q\_1qU}} \\
  & \hyperref[sec:55qaAhCWIcM]{God's Providence 2 (Genesis 39) - Rev. Mark Chan} & 2023-10-11 & \href{https://youtube.com/watch?v=55qaAhCWIcM}{\texttt{55qaAhCWIcM}} \\
創世記  & \hyperref[sec:ZcJI3G9ScsI]{上主的護理 (3)細想約瑟的故事 3 (創世記 40-41) - 陳崇基牧師} & 2023-10-11 & \href{https://youtube.com/watch?v=ZcJI3G9ScsI}{\texttt{ZcJI3G9ScsI}} \\
  & \hyperref[sec:s6E5zSGVp5c]{比對研究 19 巴力宗教 12 - 列王記上 19 - 陳崇基牧師} & 2023-10-13 & \href{https://youtube.com/watch?v=s6E5zSGVp5c}{\texttt{s6E5zSGVp5c}} \\
 15:1-7 & \hyperref[sec:iNlk2JOLf1Q]{The Power of Storytelling (Luke 15\_1-7, Matthew 18\_1-6,10-13) - Rebecca Wu} & 2023-10-17 & \href{https://youtube.com/watch?v=iNlk2JOLf1Q}{\texttt{iNlk2JOLf1Q}} \\
  & \hyperref[sec:RamdTJjsZh8]{以巴衝突的反思 (1) - 陳崇基牧師} & 2023-10-17 & \href{https://youtube.com/watch?v=RamdTJjsZh8}{\texttt{RamdTJjsZh8}} \\
 13:24-30 & \hyperref[sec:Xbkzaf9FOTY]{The Power of Analogy: Weed, Wheat, or What? (Matthew 13\_24-30) - Rebecca Wu} & 2023-10-24 & \href{https://youtube.com/watch?v=Xbkzaf9FOTY}{\texttt{Xbkzaf9FOTY}} \\
 19:1-19 & \hyperref[sec:fRha9TL_FUQ]{我在這裡作什麼 (列王記上19\_1-19) - 羅懿信牧師} & 2023-10-24 & \href{https://youtube.com/watch?v=fRha9TL_FUQ}{\texttt{fRha9TL\_FUQ}} \\
  & \hyperref[sec:_t1VkAvsN7o]{聖經地理 EP1 - 地理神學 (陳崇基牧師)} & 2023-10-30 & \href{https://youtube.com/watch?v=_t1VkAvsN7o}{\texttt{\_t1VkAvsN7o}} \\
  & \hyperref[sec:H8JKxV9_GHo]{Israel-Palestine Conflict - A Reflection (1) - Rev. Mark Chan} & 2023-10-31 & \href{https://youtube.com/watch?v=H8JKxV9_GHo}{\texttt{H8JKxV9\_GHo}} \\
  & \hyperref[sec:giSEqlRniIQ]{以巴衝突的反思 (2) - 陳崇基牧師} & 2023-10-31 & \href{https://youtube.com/watch?v=giSEqlRniIQ}{\texttt{giSEqlRniIQ}} \\
  & \hyperref[sec:XjRMb_PqUwc]{Israel-Palestine Conflict - A Reflection (2) - Rev. Mark Chan} & 2023-11-07 & \href{https://youtube.com/watch?v=XjRMb-PqUwc}{\texttt{XjRMb-PqUwc}} \\
  & \hyperref[sec:g4cKcCqr1bs]{以巴衝突的反思 (3) - 陳崇基牧師} & 2023-11-07 & \href{https://youtube.com/watch?v=g4cKcCqr1bs}{\texttt{g4cKcCqr1bs}} \\
  & \hyperref[sec:_21dlkrR_Vw]{Do You Even Lift, Bruh? - Brian Ng} & 2023-11-14 & \href{https://youtube.com/watch?v=-21dlkrR_Vw}{\texttt{-21dlkrR\_Vw}} \\
  & \hyperref[sec:z0bvgIPiO7c]{認識三一真神 - 陳漢儀姊妹} & 2023-11-14 & \href{https://youtube.com/watch?v=z0bvgIPiO7c}{\texttt{z0bvgIPiO7c}} \\
 5:16-22 & \hyperref[sec:__GpGaclvJI]{Daily Reminders from Jesus (1 Thess 5\_16-22)- Rev. Mark Chan} & 2023-11-21 & \href{https://youtube.com/watch?v=_-GpGaclvJI}{\texttt{\_-GpGaclvJI}} \\
  & \hyperref[sec:gNNugJCYC2g]{以巴衝突的反思 (4) - 陳崇基牧師} & 2023-11-21 & \href{https://youtube.com/watch?v=gNNugJCYC2g}{\texttt{gNNugJCYC2g}} \\
  & \hyperref[sec:UIZxTjRVxE4]{Israel-Palestine Conflict - A Reflection (3) - Rev. Mark Chan} & 2023-11-28 & \href{https://youtube.com/watch?v=UIZxTjRVxE4}{\texttt{UIZxTjRVxE4}} \\
 18:13-19:37 & \hyperref[sec:wrvuGGxg0R0]{最終話事權 (列王記下18\_13-19\_37) - 譚文倩莊傳道} & 2023-11-28 & \href{https://youtube.com/watch?v=wrvuGGxg0R0}{\texttt{wrvuGGxg0R0}} \\
 3:15 & \hyperref[sec:a0q_gDDhWfE]{The Advent of Hope (Genesis 3\_15) - George Crawley} & 2023-12-05 & \href{https://youtube.com/watch?v=a0q-gDDhWfE}{\texttt{a0q-gDDhWfE}} \\
路加福音  & \hyperref[sec:NrY2oymR5Fc]{將臨節的挑戰 - 施洗約翰出生的報喜 (路加福音1-2) \#1 - 陳崇基牧師} & 2023-12-05 & \href{https://youtube.com/watch?v=NrY2oymR5Fc}{\texttt{NrY2oymR5Fc}} \\
  & \hyperref[sec:SsOhyCfDa_I]{回應方舟問題續集:詳解局部洪水說的理據,回應抱此看法是不信聖經的質疑,處理聖經無誤,原文釋經,文學釋經及末世時代論等相關課題(2023-12-6 鄺牧道頻.城匯網台聯播)} & 2023-12-06 & \href{https://youtube.com/watch?v=SsOhyCfDa-I}{\texttt{SsOhyCfDa-I}} \\
 2:1-12 & \hyperref[sec:aFC7A7v5zzs]{The King Is Born  (Matthew 2\_1-12) - Sam Johnson} & 2023-12-12 & \href{https://youtube.com/watch?v=aFC7A7v5zzs}{\texttt{aFC7A7v5zzs}} \\
路加福音  & \hyperref[sec:iKrjLPh_hzM]{將臨節的挑戰 - 耶穌出生的報喜 (路加福音1-2) \#2 - 陳崇基牧師} & 2023-12-12 & \href{https://youtube.com/watch?v=iKrjLPh-hzM}{\texttt{iKrjLPh-hzM}} \\
  & \hyperref[sec:JrlUbBSqsPA]{Reflection on Advent (Luke 1-2) - Rev. Mark Chan} & 2023-12-19 & \href{https://youtube.com/watch?v=JrlUbBSqsPA}{\texttt{JrlUbBSqsPA}} \\
路加福音  & \hyperref[sec:vk3Vfk1fF_M]{將臨節的挑戰 - 神子降臨 (路加福音1-2) \#3 - 陳崇基牧師} & 2023-12-19 & \href{https://youtube.com/watch?v=vk3Vfk1fF-M}{\texttt{vk3Vfk1fF-M}} \\
  & \hyperref[sec:SCf1z8t8nOA]{Atomic Wisdom 1 - Rev. Mark Chan} & 2024-01-02 & \href{https://youtube.com/watch?v=SCf1z8t8nOA}{\texttt{SCf1z8t8nOA}} \\
  & \hyperref[sec:8jU_KxWNams]{原子智慧 1 - 陳崇基牧師} & 2024-01-02 & \href{https://youtube.com/watch?v=8jU-KxWNams}{\texttt{8jU-KxWNams}} \\
  & \hyperref[sec:YZjpAlGepXw]{Atomic Wisdom 2 - Rev. Mark Chan} & 2024-01-09 & \href{https://youtube.com/watch?v=YZjpAlGepXw}{\texttt{YZjpAlGepXw}} \\
  & \hyperref[sec:C09znJHrIGU]{原子智慧 2 - 陳崇基牧師} & 2024-01-09 & \href{https://youtube.com/watch?v=C09znJHrIGU}{\texttt{C09znJHrIGU}} \\
 1:15-23 & \hyperref[sec:RlCj0uSBVME]{Know Your Hope (Ephesians 1\_15-23) - Sam Johnson} & 2024-01-16 & \href{https://youtube.com/watch?v=RlCj0uSBVME}{\texttt{RlCj0uSBVME}} \\
  & \hyperref[sec:__gIiUwiPLI]{原子智慧 3 - 陳崇基牧師} & 2024-01-16 & \href{https://youtube.com/watch?v=--gIiUwiPLI}{\texttt{--gIiUwiPLI}} \\
 18:1-4 & \hyperref[sec:G8td0HdJn8w]{Are You Limping Through Your Faith? (1 King 18\_1-4, 17-40) - Tommy Lee} & 2024-01-24 & \href{https://youtube.com/watch?v=G8td0HdJn8w}{\texttt{G8td0HdJn8w}} \\
創世記  & \hyperref[sec:9LbyKyl1jKw]{上主的護理 (4) 細想約瑟的故事 4 (創世記 42-43) - 陳崇基牧師} & 2024-01-24 & \href{https://youtube.com/watch?v=9LbyKyl1jKw}{\texttt{9LbyKyl1jKw}} \\
  & \hyperref[sec:TqQFqNVa814]{Atomic Wisdom 3 - Rev. Mark Chan} & 2024-01-30 & \href{https://youtube.com/watch?v=TqQFqNVa814}{\texttt{TqQFqNVa814}} \\
  & \hyperref[sec:WYvLrPDJi5g]{The Chosen (Sermon Series) \#1 - Rev. Mark Chan} & 2024-02-06 & \href{https://youtube.com/watch?v=WYvLrPDJi5g}{\texttt{WYvLrPDJi5g}} \\
創世記  & \hyperref[sec:QExp8HlGgEo]{上主的護理 (5) 細想約瑟的故事 5 (創世記 44-50) - 陳崇基牧師} & 2024-02-06 & \href{https://youtube.com/watch?v=QExp8HlGgEo}{\texttt{QExp8HlGgEo}} \\
 2:1-3 & \hyperref[sec:eEU2FEv4C5I]{The Chosen - Sermon Series \#2 (Genesis 2\_1-3) - Rev. Mark Chan} & 2024-02-13 & \href{https://youtube.com/watch?v=eEU2FEv4C5I}{\texttt{eEU2FEv4C5I}} \\
路加福音 8:1-3 & \hyperref[sec:aIlHRCRsv_k]{The Chosen - 講道系列 \#1 (路加福音 8\_1-3) - 陳崇基牧師} & 2024-02-13 & \href{https://youtube.com/watch?v=aIlHRCRsv_k}{\texttt{aIlHRCRsv\_k}} \\
 19:1-8 & \hyperref[sec:LTR4cqb1AYI]{What Are You Doing Here? (1 Kings 19\_1-8) - Tommy Lee} & 2024-02-20 & \href{https://youtube.com/watch?v=LTR4cqb1AYI}{\texttt{LTR4cqb1AYI}} \\
箴言 1:1-7 & \hyperref[sec:Lczwsjcf0tU]{用盡一生的腦袋 (箴言 1\_1-7) - 譚文倩莊傳道} & 2024-02-20 & \href{https://youtube.com/watch?v=Lczwsjcf0tU}{\texttt{Lczwsjcf0tU}} \\
 18:21-35 & \hyperref[sec:ON4fMvKp6vs]{Forfeiting Forgiveness (Matthew 18\_21-35) - George Crawley} & 2024-02-28 & \href{https://youtube.com/watch?v=ON4fMvKp6vs}{\texttt{ON4fMvKp6vs}} \\
創世記 2:1-3 & \hyperref[sec:LB9q6sOwuKc]{The Chosen - 講道系列 \#2 (創世記 2\_1-3) - 陳崇基牧師} & 2024-02-28 & \href{https://youtube.com/watch?v=LB9q6sOwuKc}{\texttt{LB9q6sOwuKc}} \\
 6:5-8 & \hyperref[sec:_33GO1favQo]{Remember: You are speaking with your father (Matt 6\_5-8) - John Yi} & 2024-02-29 & \href{https://youtube.com/watch?v=-33GO1favQo}{\texttt{-33GO1favQo}} \\
  & \hyperref[sec:_2c5tRWbVyU]{The Road Unknown - Joshua 3 (Speaker: Tommy Lee)} & 2024-02-29 & \href{https://youtube.com/watch?v=-2c5tRWbVyU}{\texttt{-2c5tRWbVyU}} \\
 2:1-11 & \hyperref[sec:q3J7I2JVwVE]{The Chosen - Sermon Series \#3 (John 2\_1-11) - Rev. Mark Chan} & 2024-03-05 & \href{https://youtube.com/watch?v=q3J7I2JVwVE}{\texttt{q3J7I2JVwVE}} \\
 11:13-16 & \hyperref[sec:LtoT2yWzuEc]{客旅人生 (來 11\_13-16, 彼前 2\_10-11) -  黃茂林傳道} & 2024-03-05 & \href{https://youtube.com/watch?v=LtoT2yWzuEc}{\texttt{LtoT2yWzuEc}} \\
  & \hyperref[sec:01fKINSJbMU]{The Last Days According to Jesus 1 (Matthew 24) - Rev. Mark Chan} & 2024-03-09 & \href{https://youtube.com/watch?v=01fKINSJbMU}{\texttt{01fKINSJbMU}} \\
 2:1-10 & \hyperref[sec:o3vdukdwyRI]{Know Your Identity (Ephesians 2\_1-10) - Sam Johnson} & 2024-03-12 & \href{https://youtube.com/watch?v=o3vdukdwyRI}{\texttt{o3vdukdwyRI}} \\
馬可福音  & \hyperref[sec:Ep4OFfIa5EM]{耶穌最後的一週 \#1 (馬可福音 11) - 陳崇基牧師} & 2024-03-12 & \href{https://youtube.com/watch?v=Ep4OFfIa5EM}{\texttt{Ep4OFfIa5EM}} \\
  & \hyperref[sec:1C8pDQHXwCU]{Family Sunday School - EP 8 (Yes, He Is Risen!)} & 2024-03-13 & \href{https://youtube.com/watch?v=1C8pDQHXwCU}{\texttt{1C8pDQHXwCU}} \\
創世記  & \hyperref[sec:2mbeDPTg12s]{聖經考古學 - 創世記,洪水,方舟  (1) -  陳崇基牧師} & 2024-03-18 & \href{https://youtube.com/watch?v=2mbeDPTg12s}{\texttt{2mbeDPTg12s}} \\
  & \hyperref[sec:3W39B_RRLi0]{Apostles' Creed \#6 (Death \ and  Resurrection) - Rev. Mark Chan} & 2024-03-19 & \href{https://youtube.com/watch?v=3W39B-RRLi0}{\texttt{3W39B-RRLi0}} \\
  & \hyperref[sec:EpFmOvtCyQg]{God's Design and Calling (Joshua 6) - Tommy Lee} & 2024-03-19 & \href{https://youtube.com/watch?v=EpFmOvtCyQg}{\texttt{EpFmOvtCyQg}} \\
 4:13-17 & \hyperref[sec:3t6kNRmRV_s]{Planning For Tomorrow (James 4\_13-17) - Rev. Mark Chan} & 2024-03-19 & \href{https://youtube.com/watch?v=3t6kNRmRV-s}{\texttt{3t6kNRmRV-s}} \\
 4:13-18 & \hyperref[sec:3kSkaNorqlI]{Ready for Christ's Return (1 Thess 4\_13-18) - Rev. Mark Chan} & 2024-03-19 & \href{https://youtube.com/watch?v=3kSkaNorqlI}{\texttt{3kSkaNorqlI}} \\
馬可福音 12:38-44 & \hyperref[sec:dON_DJyLmhA]{耶穌最後的一週 \#2 - (馬可福音 12\_38-44) - 陳崇基牧師} & 2024-03-19 & \href{https://youtube.com/watch?v=dON_DJyLmhA}{\texttt{dON\_DJyLmhA}} \\
  & \hyperref[sec:55R6ROPsg74]{Kingdom Values XXV: Matthew Chapter 6 - Pastor Sam Johnson} & 2024-03-22 & \href{https://youtube.com/watch?v=55R6ROPsg74}{\texttt{55R6ROPsg74}} \\
  & \hyperref[sec:5fl01IlnFNQ]{聖經地理 EP10 - 地理與釋經 - 地景/氣候 (陳崇基牧師)} & 2024-03-23 & \href{https://youtube.com/watch?v=5fl01IlnFNQ}{\texttt{5fl01IlnFNQ}} \\
  & \hyperref[sec:6DyYwyvwR5E]{Faith and time: Sabbath - Rev. Mark Chan} & 2024-03-25 & \href{https://youtube.com/watch?v=6DyYwyvwR5E}{\texttt{6DyYwyvwR5E}} \\
  & \hyperref[sec:6uSoGUnudqk]{神的名字5\_公義的神 El Tsaddik (靈思牧禱)} & 2024-03-26 & \href{https://youtube.com/watch?v=6uSoGUnudqk}{\texttt{6uSoGUnudqk}} \\
  & \hyperref[sec:lvZuhPuvg84]{Final Week of Jesus (Mark 14-15) - Rev. Mark Chan} & 2024-03-27 & \href{https://youtube.com/watch?v=lvZuhPuvg84}{\texttt{lvZuhPuvg84}} \\
創世記  & \hyperref[sec:7tBg4MSn2jE]{創世記第6章 (6.7) - 背景釋經 [方舟,再創造 2] (陳崇基牧師)} & 2024-03-27 & \href{https://youtube.com/watch?v=7tBg4MSn2jE}{\texttt{7tBg4MSn2jE}} \\
馬可福音  & \hyperref[sec:IxOjYdkRpUg]{耶穌最後的一週  \#3 (馬可福音 14-15 ) - 陳崇基牧師} & 2024-03-27 & \href{https://youtube.com/watch?v=IxOjYdkRpUg}{\texttt{IxOjYdkRpUg}} \\
 6:25-34 & \hyperref[sec:_seeVdxX33U]{Seek Him (Matthew 6\_25-34)- Rev. Mark Chan} & 2024-03-29 & \href{https://youtube.com/watch?v=-seeVdxX33U}{\texttt{-seeVdxX33U}} \\
馬可福音 4:35-41 & \hyperref[sec:87VD_Eud_tU]{漫步風雨中 (馬可福音4\_35-41) - 張文龍牧師} & 2024-03-29 & \href{https://youtube.com/watch?v=87VD-Eud_tU}{\texttt{87VD-Eud\_tU}} \\
  & \hyperref[sec:9_6WJzTToTc]{聖地網上行 - 步主腳蹤 [第十集] 潔淨聖殿} & 2024-03-31 & \href{https://youtube.com/watch?v=9-6WJzTToTc}{\texttt{9-6WJzTToTc}} \\
  & \hyperref[sec:AndEm9fTjLE]{UV 匯客室 - 鮑維均博士} & 2024-04-01 & \href{https://youtube.com/watch?v=AndEm9fTjLE}{\texttt{AndEm9fTjLE}} \\
 7:21-23 & \hyperref[sec:BAxyx20fKwY]{Kingdom Values Matthew 7\_21-23 Sermon on the Mount} & 2024-04-02 & \href{https://youtube.com/watch?v=BAxyx20fKwY}{\texttt{BAxyx20fKwY}} \\
  & \hyperref[sec:BG6FWzPxpgI]{神的名字3\_Yahweh Rapha 耶和華醫治(靈思牧禱)} & 2024-04-02 & \href{https://youtube.com/watch?v=BG6FWzPxpgI}{\texttt{BG6FWzPxpgI}} \\
  & \hyperref[sec:D8rFOM5_n_k]{聖地網上行 - 步主腳蹤 [第十二集] 出賣與受審} & 2024-04-05 & \href{https://youtube.com/watch?v=D8rFOM5_n-k}{\texttt{D8rFOM5\_n-k}} \\
  & \hyperref[sec:DKQrjJlzvbM]{舊約網上行 - 上帝的足跡: EP18 揀選} & 2024-04-05 & \href{https://youtube.com/watch?v=DKQrjJlzvbM}{\texttt{DKQrjJlzvbM}} \\
  & \hyperref[sec:DbCFhByjju0]{Sing Hallelujah to Lord 詩歌意義 (黃嘉漢牧師)} & 2024-04-06 & \href{https://youtube.com/watch?v=DbCFhByjju0}{\texttt{DbCFhByjju0}} \\
  & \hyperref[sec:Dw3oWngj5UA]{Faith and time: Feast of Tabernacle 2 (Rev. Mark Chan)} & 2024-04-07 & \href{https://youtube.com/watch?v=Dw3oWngj5UA}{\texttt{Dw3oWngj5UA}} \\
 10:11-18 & \hyperref[sec:DrODNpzG4PQ]{I am a good shepherd (John 10\_11-18; Psalm 23) - Samantha Lockenour} & 2024-04-07 & \href{https://youtube.com/watch?v=DrODNpzG4PQ}{\texttt{DrODNpzG4PQ}} \\
詩篇  & \hyperref[sec:_tKyeazYj4I]{聖經地理 EP3 - 地理與釋經 - 詩篇125-2 (陳崇基牧師)} & 2024-04-08 & \href{https://youtube.com/watch?v=-tKyeazYj4I}{\texttt{-tKyeazYj4I}} \\
 39:1-13 & \hyperref[sec:vWnZGyYbGoI]{How Do We Suffer Differently as God's People (Psalms 39\_1-13) - Chee Yie Jong} & 2024-04-09 & \href{https://youtube.com/watch?v=vWnZGyYbGoI}{\texttt{vWnZGyYbGoI}} \\
約翰福音 4:1-42 & \hyperref[sec:u_YoyAgPSkE]{打破性別定型的女人 (約翰福音 4\_1-42) - 何建邦傳道} & 2024-04-09 & \href{https://youtube.com/watch?v=u_YoyAgPSkE}{\texttt{u\_YoyAgPSkE}} \\
  & \hyperref[sec:HYnhtr6gM20]{Everlasting Praise (Psalms 8) - Rev. Mark Chan} & 2024-04-15 & \href{https://youtube.com/watch?v=HYnhtr6gM20}{\texttt{HYnhtr6gM20}} \\
詩篇  & \hyperref[sec:DfbxT6XIBuk]{永恆的敬拜 1 (詩篇 8) - 陳崇基牧師} & 2024-04-15 & \href{https://youtube.com/watch?v=DfbxT6XIBuk}{\texttt{DfbxT6XIBuk}} \\
  & \hyperref[sec:yghFb3uoZ_c]{11/11/18 主日講道 (Pastor Tony Wong)} & 2024-04-19 & \href{https://youtube.com/watch?v=yghFb3uoZ_c}{\texttt{yghFb3uoZ\_c}} \\
  & \hyperref[sec:GRZOoQIm4Fg]{5/13/18 主日講道: 告訴子女兩句重要的話 (張麗仙傳道)} & 2024-04-19 & \href{https://youtube.com/watch?v=GRZOoQIm4Fg}{\texttt{GRZOoQIm4Fg}} \\
  & \hyperref[sec:GXyrjafkb9Y]{Gaining a New Perspective Can Change Your Life - Keith Draper} & 2024-04-20 & \href{https://youtube.com/watch?v=GXyrjafkb9Y}{\texttt{GXyrjafkb9Y}} \\
  & \hyperref[sec:KF6pQDFyEx4]{聖地網上行 - 步主腳蹤 [第六集] 加利利湖} & 2024-04-20 & \href{https://youtube.com/watch?v=KF6pQDFyEx4}{\texttt{KF6pQDFyEx4}} \\
  & \hyperref[sec:cEtF7vpVnrE]{Everlasting Praise 2 (Genesis 2) - Rev. Mark Chan} & 2024-04-22 & \href{https://youtube.com/watch?v=cEtF7vpVnrE}{\texttt{cEtF7vpVnrE}} \\
士師記  & \hyperref[sec:H_MgtrWeR24]{從懼怕到召命的覺醒 (士師記 6) - 莫家承牧師} & 2024-04-22 & \href{https://youtube.com/watch?v=H-MgtrWeR24}{\texttt{H-MgtrWeR24}} \\
創世記  & \hyperref[sec:g90s_0fTe_c]{永恆的敬拜 2 (創世記 2) - 陳崇基牧師} & 2024-04-22 & \href{https://youtube.com/watch?v=g90s_0fTe-c}{\texttt{g90s\_0fTe-c}} \\
 10:25-37 & \hyperref[sec:Hq5No0HLb30]{Kingdom Values XLV: Luke 10\_25-37} & 2024-04-26 & \href{https://youtube.com/watch?v=Hq5No0HLb30}{\texttt{Hq5No0HLb30}} \\
約伯記  & \hyperref[sec:Gigyg3Yfrs0]{逆要作門徒 (1) – 約伯記 (陳崇基牧師) 2022.01.02} & 2024-04-27 & \href{https://youtube.com/watch?v=Gigyg3Yfrs0}{\texttt{Gigyg3Yfrs0}} \\
 8:18-30 & \hyperref[sec:ANI7oTOQdWo]{Hope in Heaven for Suffering on Earth (Romans 8\_18-30) - George Crawley} & 2024-04-30 & \href{https://youtube.com/watch?v=ANI7oTOQdWo}{\texttt{ANI7oTOQdWo}} \\
 4:7-12 & \hyperref[sec:J6PHRTa1EUI]{Love, the Greatest Gift, the Greatest Power, and the Greatest Purpose (1 John 4\_7-12) - Rev. Jeong} & 2024-04-30 & \href{https://youtube.com/watch?v=J6PHRTa1EUI}{\texttt{J6PHRTa1EUI}} \\
出埃及記  & \hyperref[sec:vMWYJSKl8Ig]{永遠的頌讚 3 (出埃及記 25-31) -  陳崇基牧師} & 2024-04-30 & \href{https://youtube.com/watch?v=vMWYJSKl8Ig}{\texttt{vMWYJSKl8Ig}} \\
  & \hyperref[sec:JpQUnDljfo8]{Family Sunday School - EP 4 (It Could Have Been You)} & 2024-05-04 & \href{https://youtube.com/watch?v=JpQUnDljfo8}{\texttt{JpQUnDljfo8}} \\
  & \hyperref[sec:K7tbzr00wo0]{佳節倍思恩(二): 住棚節(下) - 陳崇基牧師} & 2024-05-04 & \href{https://youtube.com/watch?v=K7tbzr00wo0}{\texttt{K7tbzr00wo0}} \\
  & \hyperref[sec:7G7VT1O0Fmo]{Everlasting Praise 3 (Exodus 25-31) - Rev. Mark Chan} & 2024-05-07 & \href{https://youtube.com/watch?v=7G7VT1O0Fmo}{\texttt{7G7VT1O0Fmo}} \\
  & \hyperref[sec:yIgKl3jW8_8]{永恆的敬拜 3 (列王記上 6-7 ) - 陳崇基牧師} & 2024-05-07 & \href{https://youtube.com/watch?v=yIgKl3jW8-8}{\texttt{yIgKl3jW8-8}} \\
 24:1-12 & \hyperref[sec:EISv2Eie8Pg]{Easter: Luke 24\_1-12 He Conquered Sin and Death} & 2024-05-08 & \href{https://youtube.com/watch?v=EISv2Eie8Pg}{\texttt{EISv2Eie8Pg}} \\
  & \hyperref[sec:6bRdtM__Fr0]{愈靠近,愈堅強 - 張麗仙傳道} & 2024-05-09 & \href{https://youtube.com/watch?v=6bRdtM__Fr0}{\texttt{6bRdtM\_\_Fr0}} \\
  & \hyperref[sec:D_DiJk8_Dhc]{Everlasting Praise 4 (1 Kings 6-7) - Rev. Mark Chan--} & 2024-05-14 & \href{https://youtube.com/watch?v=D_DiJk8_Dhc}{\texttt{D\_DiJk8\_Dhc}} \\
  & \hyperref[sec:NNhx4LQjZTc]{比對研究 1: 必要, 發展, 方法   -  陳崇基牧師} & 2024-05-22 & \href{https://youtube.com/watch?v=NNhx4LQjZTc}{\texttt{NNhx4LQjZTc}} \\
民數記 6:24-26 & \hyperref[sec:NOP7HqFawO8]{神賜的福   民數記 6\_24-26  (陳崇基牧師)} & 2024-05-22 & \href{https://youtube.com/watch?v=NOP7HqFawO8}{\texttt{NOP7HqFawO8}} \\
  & \hyperref[sec:NP_1P0Hd_Zc]{考古探秘 (發現死海古卷的傳奇故事) - 陳崇基牧師 Mark Sir (2021.05.30)} & 2024-05-22 & \href{https://youtube.com/watch?v=NP-1P0Hd_Zc}{\texttt{NP-1P0Hd\_Zc}} \\
 2:11-22 & \hyperref[sec:OP86jLVzWZ4]{Jesus Established the Chruch (Eph 2\_11-22) - Steve Tan} & 2024-05-27 & \href{https://youtube.com/watch?v=OP86jLVzWZ4}{\texttt{OP86jLVzWZ4}} \\
  & \hyperref[sec:OVsK1efMh70]{聖地網上行 - 步主腳蹤 [第十四集] 死裡復活} & 2024-05-27 & \href{https://youtube.com/watch?v=OVsK1efMh70}{\texttt{OVsK1efMh70}} \\
 2:11-22 & \hyperref[sec:qLV2rgTxxTw]{Reconciled in Christ (Ephesians 2\_11-22) - Sam Johnson} & 2024-05-28 & \href{https://youtube.com/watch?v=qLV2rgTxxTw}{\texttt{qLV2rgTxxTw}} \\
腓利門書  & \hyperref[sec:FBZ55r0dhNA]{升呢: 福音重新定義人與人之間的關係 (腓利門書) - 譚文倩莊傳道} & 2024-05-28 & \href{https://youtube.com/watch?v=FBZ55r0dhNA}{\texttt{FBZ55r0dhNA}} \\
  & \hyperref[sec:fTMtbscgTAE]{Everlasting Praise 5 (Corinthians 14) - Rev. Mark Chan -} & 2024-06-05 & \href{https://youtube.com/watch?v=fTMtbscgTAE}{\texttt{fTMtbscgTAE}} \\
哥林多前書  & \hyperref[sec:mwhshrf5O44]{永恆的敬拜 4 (哥林多前書 14) - 陳崇基牧師} & 2024-06-05 & \href{https://youtube.com/watch?v=mwhshrf5O44}{\texttt{mwhshrf5O44}} \\
  & \hyperref[sec:Gu87k6iNMK0]{Everlasting Praise 6 (Colossian 3) - Rev. Mark Chan} & 2024-06-13 & \href{https://youtube.com/watch?v=Gu87k6iNMK0}{\texttt{Gu87k6iNMK0}} \\
歌羅西書  & \hyperref[sec:3ck_KAakvyM]{永恆的敬拜 5 (歌羅西書 3) - 陳崇基牧師} & 2024-06-13 & \href{https://youtube.com/watch?v=3ck_KAakvyM}{\texttt{3ck\_KAakvyM}} \\
  & \hyperref[sec:xTyYHQxRMCA]{Everlasting Praise 7 (Rev 4) - Rev. Mark Chan} & 2024-06-19 & \href{https://youtube.com/watch?v=xTyYHQxRMCA}{\texttt{xTyYHQxRMCA}} \\
啟示錄 4:1-11 & \hyperref[sec:VHqi_cqG4KM]{永恆的敬拜 6 (啟示錄 4\_1-11) - 陳崇基牧師} & 2024-06-19 & \href{https://youtube.com/watch?v=VHqi_cqG4KM}{\texttt{VHqi\_cqG4KM}} \\
 5:1-16 & \hyperref[sec:cMH1ZiMWOm8]{The Upside-Down Kingdom (Matthew 5\_1-16) - Chee Yie Jong} & 2024-06-25 & \href{https://youtube.com/watch?v=cMH1ZiMWOm8}{\texttt{cMH1ZiMWOm8}} \\
腓利門書 1:4-7 & \hyperref[sec:hDHhZ8zCYYE]{愛回家:教會群體的功課  (腓利門書 1\_4-7) - 譚文倩莊傳道} & 2024-06-25 & \href{https://youtube.com/watch?v=hDHhZ8zCYYE}{\texttt{hDHhZ8zCYYE}} \\
  & \hyperref[sec:VPCYBjM4gao]{Back to Basics (Psalms 119) - Rev. Mark Chan} & 2024-07-03 & \href{https://youtube.com/watch?v=VPCYBjM4gao}{\texttt{VPCYBjM4gao}} \\
 3:1-9 & \hyperref[sec:KO8k0_vDZOY]{Back to Basics \#2 (2 Timothy 3\_1-9) - Rev. Mark Chan} & 2024-07-09 & \href{https://youtube.com/watch?v=KO8k0_vDZOY}{\texttt{KO8k0\_vDZOY}} \\
提摩太後書 3:1-9 & \hyperref[sec:x0e0tuTJkmQ]{回到基礎 \#1 (提摩太後書 3\_1-9) - 陳崇基牧師} & 2024-07-09 & \href{https://youtube.com/watch?v=x0e0tuTJkmQ}{\texttt{x0e0tuTJkmQ}} \\
腓利門書  & \hyperref[sec:eiPsrBdDRJk]{請勿超越黃線 (腓利門書) - 譚文倩莊傳道} & 2024-07-17 & \href{https://youtube.com/watch?v=eiPsrBdDRJk}{\texttt{eiPsrBdDRJk}} \\
 3:10-15 & \hyperref[sec:QOZ3Ws7CSlE]{Back to Basics \#3 (2 Timothy 3\_10-15) - Rev. Mark Chan-} & 2024-07-18 & \href{https://youtube.com/watch?v=QOZ3Ws7CSlE}{\texttt{QOZ3Ws7CSlE}} \\
  & \hyperref[sec:QovpodRHlhw]{流散與福音廣傳 - 劉漢杰醫生} & 2024-07-21 & \href{https://youtube.com/watch?v=QovpodRHlhw}{\texttt{QovpodRHlhw}} \\
 3:16-17 & \hyperref[sec:umenrD4C3Dw]{Back to Basics \#4 (2 Timothy 3\_16-17) - Rev. Mark Chan} & 2024-07-23 & \href{https://youtube.com/watch?v=umenrD4C3Dw}{\texttt{umenrD4C3Dw}} \\
提摩太後書 3:10-17 & \hyperref[sec:koqOwadJaEc]{回到基礎 \#2 (提摩太後書 3\_10-17) - 陳崇基牧師} & 2024-07-23 & \href{https://youtube.com/watch?v=koqOwadJaEc}{\texttt{koqOwadJaEc}} \\
  & \hyperref[sec:9TgjTCjHTHQ]{Back to Basics \#5 (Philippians 4) - Rev. Mark Chan} & 2024-08-01 & \href{https://youtube.com/watch?v=9TgjTCjHTHQ}{\texttt{9TgjTCjHTHQ}} \\
腓利比書  & \hyperref[sec:gEFzH4C2B2Y]{回到基礎 \#3 (腓利比書 4) - 陳崇基牧師} & 2024-08-01 & \href{https://youtube.com/watch?v=gEFzH4C2B2Y}{\texttt{gEFzH4C2B2Y}} \\
 12:1-8 & \hyperref[sec:WzsPGK70Ghc]{Back to Basics \#6 (Romans 12\_1-8) - Rev. Mark Chan} & 2024-08-06 & \href{https://youtube.com/watch?v=WzsPGK70Ghc}{\texttt{WzsPGK70Ghc}} \\
  & \hyperref[sec:l_mjl7q5DSI]{Back to Basics \#8 (Corinthians 12) - Rev. Mark Chan} & 2024-08-23 & \href{https://youtube.com/watch?v=l-mjl7q5DSI}{\texttt{l-mjl7q5DSI}} \\
  & \hyperref[sec:nNDuojAF0eE]{開往水深之處: 信仰歷程中不同階段的禱告生命 - 余淇煥傳道} & 2024-08-23 & \href{https://youtube.com/watch?v=nNDuojAF0eE}{\texttt{nNDuojAF0eE}} \\
 15:16 & \hyperref[sec:jtr0kNAa9Hw]{Fruit that Lasts (John 15\_16) - Sam Johnson} & 2024-08-27 & \href{https://youtube.com/watch?v=jtr0kNAa9Hw}{\texttt{jtr0kNAa9Hw}} \\
哥林多後書 8:1-15 & \hyperref[sec:7j_P4EWgCCI]{恩典的慷慨 (哥林多後書 8\_1-15, 9\_6-15) - 譚文倩莊傳道} & 2024-08-27 & \href{https://youtube.com/watch?v=7j_P4EWgCCI}{\texttt{7j\_P4EWgCCI}} \\
  & \hyperref[sec:pxG7cdEbvuY]{What Christ Thinks of the Church (Revelation 2-3) - Rev. Mark Chan} & 2024-09-05 & \href{https://youtube.com/watch?v=pxG7cdEbvuY}{\texttt{pxG7cdEbvuY}} \\
  & \hyperref[sec:YG4nAX61j6k]{基督眼中的教會 (啟示錄 2-3) - 陳崇基牧師} & 2024-09-05 & \href{https://youtube.com/watch?v=YG4nAX61j6k}{\texttt{YG4nAX61j6k}} \\
  & \hyperref[sec:U3A7_wArCTw]{What Christ Thinks of the Church \#2  (Revelation 2-3) - Rev.  Mark Chan} & 2024-09-17 & \href{https://youtube.com/watch?v=U3A7-wArCTw}{\texttt{U3A7-wArCTw}} \\
  & \hyperref[sec:h0WrV1zwMlM]{What Christ Thinks of the Church \#3 (Revelation 2-3) - Rev. Mark Chan} & 2024-09-17 & \href{https://youtube.com/watch?v=h0WrV1zwMlM}{\texttt{h0WrV1zwMlM}} \\
歌羅西書 1:28-29 & \hyperref[sec:FGEt94QjF9Q]{在基督裡的勉導同行 (歌羅西書1\_28-29) - 蘇妙嫻傳道} & 2024-09-17 & \href{https://youtube.com/watch?v=FGEt94QjF9Q}{\texttt{FGEt94QjF9Q}} \\
啟示錄 2:1-3:22 & \hyperref[sec:S2QbUdfuEWU]{基督眼中的教會 \#2 (啟示錄 2\_1-3\_22) - 陳崇基牧師} & 2024-09-17 & \href{https://youtube.com/watch?v=S2QbUdfuEWU}{\texttt{S2QbUdfuEWU}} \\
  & \hyperref[sec:KVIuGXk07jc]{What Christ Thinks of the Church \#4 (Revelation 2-3) - Rev. Mark Chan} & 2024-10-01 & \href{https://youtube.com/watch?v=KVIuGXk07jc}{\texttt{KVIuGXk07jc}} \\
  & \hyperref[sec:pG_vfjaVL_8]{What Christ Thinks of the Church \#5 (Revelation 2-3) - Rev. Mark Chan} & 2024-10-01 & \href{https://youtube.com/watch?v=pG-vfjaVL_8}{\texttt{pG-vfjaVL\_8}} \\
使徒行傳 7:2-60 & \hyperref[sec:9bZktUw13B0]{史上最撩交打ge 自辯 (使徒行傳 7\_2-60) - 何建邦傳道} & 2024-10-01 & \href{https://youtube.com/watch?v=9bZktUw13B0}{\texttt{9bZktUw13B0}} \\
啟示錄 2:1-3:22 & \hyperref[sec:1ry__RW4dIE]{基督眼中的教會 \#3 (啟示錄 2\_1-3\_22) - 陳崇基牧師} & 2024-10-02 & \href{https://youtube.com/watch?v=1ry_-RW4dIE}{\texttt{1ry\_-RW4dIE}} \\
  & \hyperref[sec:yfuGrCAshKk]{What Christ Thinks of the Church \#6 (Revelation 2-3) - Rev. Mark Chan} & 2024-10-08 & \href{https://youtube.com/watch?v=yfuGrCAshKk}{\texttt{yfuGrCAshKk}} \\
 5:17-18 & \hyperref[sec:vjkb_KdvSz0]{工人的工價 (提摩太前 5\_17-18) - 陳崇基牧師} & 2024-10-08 & \href{https://youtube.com/watch?v=vjkb_KdvSz0}{\texttt{vjkb\_KdvSz0}} \\
  & \hyperref[sec:hfmoR2RDphM]{What Christ Thinks of the Church \#7 (Revelation 2-3) - Rev. Mark Chan} & 2024-10-15 & \href{https://youtube.com/watch?v=hfmoR2RDphM}{\texttt{hfmoR2RDphM}} \\
啟示錄 2:1-3:22 & \hyperref[sec:sCrh8yIfIpY]{基督眼中的教會 \#4 (啟示錄 2\_1-3\_22) - 陳崇基牧師} & 2024-10-15 & \href{https://youtube.com/watch?v=sCrh8yIfIpY}{\texttt{sCrh8yIfIpY}} \\
 24:30-34 & \hyperref[sec:yEZdtKfic4o]{Way of Wisdom \#1 (Proverbs 24\_30-34) - Rev. Mark Chan} & 2024-10-22 & \href{https://youtube.com/watch?v=yEZdtKfic4o}{\texttt{yEZdtKfic4o}} \\
啟示錄 2:1-3:22 & \hyperref[sec:VjJ7v3bZQ7A]{基督眼中的教會 \#5 (啟示錄 2\_1-3\_22) - 陳崇基牧師} & 2024-10-22 & \href{https://youtube.com/watch?v=VjJ7v3bZQ7A}{\texttt{VjJ7v3bZQ7A}} \\
 16:18 & \hyperref[sec:uDsOylL_mRw]{Way of Wisdom \#2 - Proverbs 16\_18 (Rev. Mark Chan)} & 2024-11-02 & \href{https://youtube.com/watch?v=uDsOylL_mRw}{\texttt{uDsOylL\_mRw}} \\
約拿書  & \hyperref[sec:v8IvKYM0B_8]{約拿一小步 (約拿書 23) - 鍾志愷弟兄} & 2024-11-02 & \href{https://youtube.com/watch?v=v8IvKYM0B-8}{\texttt{v8IvKYM0B-8}} \\
 5:17-20 & \hyperref[sec:T8J2ny3cosE]{Truly Blessed Life According to Jesus (Matthew 5\_17-20) - Chee Yie Jong} & 2024-11-06 & \href{https://youtube.com/watch?v=T8J2ny3cosE}{\texttt{T8J2ny3cosE}} \\
詩篇  & \hyperref[sec:Pk_kHNdAhpo]{這人真是有福 (詩篇 1) - 黃國維博士} & 2024-11-06 & \href{https://youtube.com/watch?v=Pk-kHNdAhpo}{\texttt{Pk-kHNdAhpo}} \\
  & \hyperref[sec:ywevukHYUvY]{Who is Your Lord? \#2 Belshazzar (Daniel 5) - Rev. Mark Chan} & 2024-11-19 & \href{https://youtube.com/watch?v=ywevukHYUvY}{\texttt{ywevukHYUvY}} \\
馬可福音 3:1-6 & \hyperref[sec:KOpSYq8TJvA]{新約中的情緒 (1): 憤怒 (馬可福音 3\_1-6) - 何建邦傳道} & 2024-11-19 & \href{https://youtube.com/watch?v=KOpSYq8TJvA}{\texttt{KOpSYq8TJvA}} \\
  & \hyperref[sec:dQy1lcfV2Eg]{Who is Your Lord? \#1 Cyrus (Isaiah 45) - Rev. Mark Chan} & 2024-11-26 & \href{https://youtube.com/watch?v=dQy1lcfV2Eg}{\texttt{dQy1lcfV2Eg}} \\
  & \hyperref[sec:fC6e7drn7H4]{感恩,因為.感恩,縱然.Give Thanks because... even though... (陳崇基牧師 Rev. Mark Chan)} & 2024-11-26 & \href{https://youtube.com/watch?v=fC6e7drn7H4}{\texttt{fC6e7drn7H4}} \\
以賽亞書  & \hyperref[sec:HjvG0k680po]{誰是你主?\#1 居魯士 (以賽亞書 45) - 陳崇基牧師} & 2024-11-26 & \href{https://youtube.com/watch?v=HjvG0k680po}{\texttt{HjvG0k680po}} \\
 5:21-26 & \hyperref[sec:sCP4BPS1_XA]{Righteousness as Redefined by Jesus: Anger (Matthew 5\_21-26) - Chee Yie Jong} & 2024-12-03 & \href{https://youtube.com/watch?v=sCP4BPS1_XA}{\texttt{sCP4BPS1\_XA}} \\
但以理書  & \hyperref[sec:FaphUTtEgYU]{誰是你主?\#2 伯沙撒王 (但以理書 5) - 陳崇基牧師} & 2024-12-03 & \href{https://youtube.com/watch?v=FaphUTtEgYU}{\texttt{FaphUTtEgYU}} \\
  & \hyperref[sec:_TmbaOYl2cg]{Who is Your Lord \#3 Augustus (Luke 2) - Rev. Mark Chan} & 2024-12-11 & \href{https://youtube.com/watch?v=-TmbaOYl2cg}{\texttt{-TmbaOYl2cg}} \\
路加福音 19:37-48 & \hyperref[sec:WhzVz61wqtY]{耶穌:復興的聖殿 (路加福音 19\_37-48) - 梁紫健傳道} & 2024-12-11 & \href{https://youtube.com/watch?v=WhzVz61wqtY}{\texttt{WhzVz61wqtY}} \\
 20:1-17 & \hyperref[sec:UlWyJIvtS2w]{New Principles for New People (Exodus 20\_1-17) - Rev. Mark Chan} & 2025-01-04 & \href{https://youtube.com/watch?v=UlWyJIvtS2w}{\texttt{UlWyJIvtS2w}} \\
路加福音  & \hyperref[sec:3GTzRC3gFVM]{聖誕的平安 - 路加福音 2 (陳崇基牧師)  Peace of Christmas - Luke 2 (Rev. Mark Chan)} & 2025-01-04 & \href{https://youtube.com/watch?v=3GTzRC3gFVM}{\texttt{3GTzRC3gFVM}} \\
路加福音 2:8-20 & \hyperref[sec:Lw531TtsFSc]{隱藏的聖誕 (路加福音 2\_8-20) - 譚文倩莊傳道} & 2025-01-04 & \href{https://youtube.com/watch?v=Lw531TtsFSc}{\texttt{Lw531TtsFSc}} \\
  & \hyperref[sec:pzPtCyCLwfM]{The Wonder Of Christmas (Luke 2) - Rev. Mark Chan} & 2025-01-05 & \href{https://youtube.com/watch?v=pzPtCyCLwfM}{\texttt{pzPtCyCLwfM}} \\
路加福音  & \hyperref[sec:KuhthfKDgTw]{誰是你主?\#2 奧古斯都 (路加福音 2) - 陳崇基牧師} & 2025-01-05 & \href{https://youtube.com/watch?v=KuhthfKDgTw}{\texttt{KuhthfKDgTw}} \\
出埃及記 20:1-17 & \hyperref[sec:x5SxAthpTtg]{新子民,新原則 \#1 (出埃及記 20\_1-17) - 陳崇基牧師} & 2025-01-07 & \href{https://youtube.com/watch?v=x5SxAthpTtg}{\texttt{x5SxAthpTtg}} \\
 20:1-17 & \hyperref[sec:kxR54r6RFz8]{New Principals, New People \#2 (Exodus 20\_1-17) - Rev. Mark Chan} & 2025-01-14 & \href{https://youtube.com/watch?v=kxR54r6RFz8}{\texttt{kxR54r6RFz8}} \\
出埃及記 20:1-17 & \hyperref[sec:9TO87He9Dbc]{新子民,新原則 \#2 (出埃及記 20\_1-17) - 陳崇基牧師} & 2025-01-14 & \href{https://youtube.com/watch?v=9TO87He9Dbc}{\texttt{9TO87He9Dbc}} \\
羅馬書  & \hyperref[sec:kf1XXedwCsU]{神愛大地 (羅馬書 8) - 陳崇基牧師} & 2025-01-28 & \href{https://youtube.com/watch?v=kf1XXedwCsU}{\texttt{kf1XXedwCsU}} \\
 5:31-32 & \hyperref[sec:CwJ8DZXjXd8]{Righteousness as Redefined by Jesus: Marriage \ and  Divorce (Matthew 5\_31-32) - Chee Yie Jong} & 2025-02-05 & \href{https://youtube.com/watch?v=CwJ8DZXjXd8}{\texttt{CwJ8DZXjXd8}} \\
約翰福音 3:14-15 & \hyperref[sec:_BgBlG2kkQM]{耶穌——賜生命的蛇 (約翰福音 3\_14-15) - 何建邦傳道} & 2025-02-05 & \href{https://youtube.com/watch?v=-BgBlG2kkQM}{\texttt{-BgBlG2kkQM}} \\
  & \hyperref[sec:cIQZdCWcAV4]{神愛大地 (舊約篇) - 陳崇基牧師} & 2025-02-11 & \href{https://youtube.com/watch?v=cIQZdCWcAV4}{\texttt{cIQZdCWcAV4}} \\
馬太福音 25:14-30 & \hyperref[sec:UraagaffkTo]{神國藍圖 (馬太福音 25\_14-30) - 余淇煥傳道} & 2025-02-18 & \href{https://youtube.com/watch?v=UraagaffkTo}{\texttt{UraagaffkTo}} \\
  & \hyperref[sec:iar_apDpNcQ]{神愛大地 (新約篇) - 陳崇基牧師} & 2025-02-25 & \href{https://youtube.com/watch?v=iar-apDpNcQ}{\texttt{iar-apDpNcQ}} \\
 5:33-42 & \hyperref[sec:IH4xN4bYlR0]{Righteousness as Redefined by Jesus: Word and Deed (Matthew 5\_33-42) - Chee Yie Jong} & 2025-03-05 & \href{https://youtube.com/watch?v=IH4xN4bYlR0}{\texttt{IH4xN4bYlR0}} \\
約翰福音  & \hyperref[sec:LHY953jFc_g]{生命頻道 \#1 (約翰福音 10) - 陳崇基牧師} & 2025-03-06 & \href{https://youtube.com/watch?v=LHY953jFc_g}{\texttt{LHY953jFc\_g}} \\
馬可福音 14:32-36 & \hyperref[sec:KKxcZYP9oLY]{新約中的情緒 (2): 恐懼 (馬可福音 14\_32-36) - 何建邦傳道} & 2025-03-11 & \href{https://youtube.com/watch?v=KKxcZYP9oLY}{\texttt{KKxcZYP9oLY}} \\
約翰福音  & \hyperref[sec:gEwUekxM30c]{生命頻道 \#2:聆聽上主的聲音 (約翰福音 15) - 陳崇基牧師} & 2025-03-18 & \href{https://youtube.com/watch?v=gEwUekxM30c}{\texttt{gEwUekxM30c}} \\
約翰福音  & \hyperref[sec:1U5fivtmb3s]{生命頻道 \#3:聆聽上主的聲音 (約翰福音 14-17) - 陳崇基牧師} & 2025-03-25 & \href{https://youtube.com/watch?v=1U5fivtmb3s}{\texttt{1U5fivtmb3s}} \\
約翰福音 6:52-60 & \hyperref[sec:ETczAl2v1zo]{生命麵包,你今日食jor 未?(約翰福音 6\_52-60) - 梁紫健傳道} & 2025-04-01 & \href{https://youtube.com/watch?v=ETczAl2v1zo}{\texttt{ETczAl2v1zo}} \\
 5:43-48 & \hyperref[sec:95H_Df1nGvs]{Righteousness as Redefined by Jesus: Love Your Enemies (Matthew 5\_43-48) - Chee Yie Jong} & 2025-04-10 & \href{https://youtube.com/watch?v=95H_Df1nGvs}{\texttt{95H\_Df1nGvs}} \\
約伯記  & \hyperref[sec:FaEgprmwTI0]{在問題中遇見神 (約伯記) - 余淇煥傳道} & 2025-04-10 & \href{https://youtube.com/watch?v=FaEgprmwTI0}{\texttt{FaEgprmwTI0}} \\
使徒行傳 20:17-38 & \hyperref[sec:i_UUX_QnitY]{明辨真理 (使徒行傳 20\_17-38) - 陳崇基牧師} & 2025-04-16 & \href{https://youtube.com/watch?v=i_UUX-QnitY}{\texttt{i\_UUX-QnitY}} \\
 5:27-30 & \hyperref[sec:TstKn_BEGfA]{Righteousness as Redefined by Jesus: Lust (Matthew 5\_27-30) - Chee Yie Jong} & 2025-04-20 & \href{https://youtube.com/watch?v=TstKn_BEGfA}{\texttt{TstKn\_BEGfA}} \\
  & \hyperref[sec:as4vlJLVesI]{5/6/18 主日講道: 我如何從回教徒成為基督徒 (Pastor Wissam Youssef )} & 2025-04-25 & \href{https://youtube.com/watch?v=as4vlJLVesI}{\texttt{as4vlJLVesI}} \\
 1:8 & \hyperref[sec:Ml1uo7s2nGU]{Cross Over (Acts 1\_8) - Rev. Mark Chan} & 2025-04-29 & \href{https://youtube.com/watch?v=Ml1uo7s2nGU}{\texttt{Ml1uo7s2nGU}} \\
 6:1-1 & \hyperref[sec:BPEn5KuFa4c]{Pursuing Kingdom Righteousness (Matthew 6\_1-1) - Chee Yie Jong} & 2025-05-07 & \href{https://youtube.com/watch?v=BPEn5KuFa4c}{\texttt{BPEn5KuFa4c}} \\
創世記 22:1 & \hyperref[sec:Wr2X99fNMqk]{「空白」的三天 (創世記 22\_1–19) - 楊靖詩博士} & 2025-05-07 & \href{https://youtube.com/watch?v=Wr2X99fNMqk}{\texttt{Wr2X99fNMqk}} \\
使徒行傳  & \hyperref[sec:rAFDbYbX6HA]{天國的啟動 (使徒行傳 1) - 陳崇基牧師} & 2025-05-13 & \href{https://youtube.com/watch?v=rAFDbYbX6HA}{\texttt{rAFDbYbX6HA}} \\
  & \hyperref[sec:bewfsUoYndE]{對兒女的最大期望:智慧 - 黃國棟醫生} & 2025-05-21 & \href{https://youtube.com/watch?v=bewfsUoYndE}{\texttt{bewfsUoYndE}} \\
使徒行傳  & \hyperref[sec:MAWjFwIPHJE]{神國的模樣 (使徒行傳 2) - 陳崇基牧師} & 2025-05-28 & \href{https://youtube.com/watch?v=MAWjFwIPHJE}{\texttt{MAWjFwIPHJE}} \\
 3:1 & \hyperref[sec:34yXkIAS4Hk]{對上帝喜悅的祈求 (列王記上 3\_1–15)  - 江芷敏傳道} & 2025-06-05 & \href{https://youtube.com/watch?v=34yXkIAS4Hk}{\texttt{34yXkIAS4Hk}} \\
馬可福音  & \hyperref[sec:EgZkScLaImE]{新約中的情緒 (3): 哀傷 (馬可福音 14:32-36) - 何建邦傳道} & 2025-06-11 & \href{https://youtube.com/watch?v=EgZkScLaImE}{\texttt{EgZkScLaImE}} \\
箴言 22:6 & \hyperref[sec:lPvFkn_YDDw]{父母的職事 (箴言 22\_6) - 陳崇基牧師} & 2025-06-19 & \href{https://youtube.com/watch?v=lPvFkn-YDDw}{\texttt{lPvFkn-YDDw}} \\
使徒行傳 3:1-10 & \hyperref[sec:cZT_1gxTxfc]{福音的突破 (使徒行傳 3\_1-10) - 陳崇基牧師} & 2025-06-24 & \href{https://youtube.com/watch?v=cZT-1gxTxfc}{\texttt{cZT-1gxTxfc}} \\
  & \hyperref[sec:1fX6w4FAuCw]{舊約網上行 - 上帝的足跡: EP12 款待與報喜} & 2025-07-01 & \href{https://youtube.com/watch?v=1fX6w4FAuCw}{\texttt{1fX6w4FAuCw}} \\
路加福音 10:25-37 & \hyperref[sec:ot0FwdbiPyw]{跨越界線的愛 (路加福音 10\_25-37) - 梁紫健傳道} & 2025-07-01 & \href{https://youtube.com/watch?v=ot0FwdbiPyw}{\texttt{ot0FwdbiPyw}} \\
使徒行傳 28:1-10 & \hyperref[sec:EEmGj96kwiU]{接待的福音(Be Our Guests) (使徒行傳28\_1-10, 30-31) - 朱莉傳道} & 2025-07-09 & \href{https://youtube.com/watch?v=EEmGj96kwiU}{\texttt{EEmGj96kwiU}} \\
使徒行傳 4:32-37 & \hyperref[sec:_p6p1yxlKsI]{成為聖靈的群體 (使徒行傳4\_32-37) - 陳崇基牧師} & 2025-07-14 & \href{https://youtube.com/watch?v=-p6p1yxlKsI}{\texttt{-p6p1yxlKsI}} \\
 6:9-15 & \hyperref[sec:6C_6zjSlvhM]{A Model for Prayer (Matthew 6\_9-15) - Chee Yie Jong} & 2025-08-24 & \href{https://youtube.com/watch?v=6C_6zjSlvhM}{\texttt{6C\_6zjSlvhM}} \\
使徒行傳 2:1-47 & \hyperref[sec:HQs8juXt_Kg]{Contagious Church (Acts 2) - Rev. Mark Chan 有感染力的教會 (使徒行傳2\_1-47) - 陳崇基牧師} & 2025-08-24 & \href{https://youtube.com/watch?v=HQs8juXt-Kg}{\texttt{HQs8juXt-Kg}} \\
使徒行傳 2:42-47 & \hyperref[sec:J7Rxq8y7uj4]{Cross Over - Be Committed ＆ Connected (Acts 2\_42-47) - Rev. Mark Chan 跨越 - 委身與連結  (使徒行傳2\_42-47) - 陳崇基牧師} & 2025-08-24 & \href{https://youtube.com/watch?v=J7Rxq8y7uj4}{\texttt{J7Rxq8y7uj4}} \\
使徒行傳 2:38 & \hyperref[sec:TtUFBp1ccf4]{Only One Way to Be Saved (Acts 2\_38; 4\_12) - Rev. Mark Chan 唯一的拯救 (使徒行傳2\_38,4\_12) - 陳崇基牧師} & 2025-08-24 & \href{https://youtube.com/watch?v=TtUFBp1ccf4}{\texttt{TtUFBp1ccf4}} \\
路加福音 22:15 & \hyperref[sec:3tR0PVmebsI]{彌賽亞 - 耶穌的筵席: 從排他到包容的救恩 (路加福音22\_15 -23) - 梁紫健傳道} & 2025-08-24 & \href{https://youtube.com/watch?v=3tR0PVmebsI}{\texttt{3tR0PVmebsI}} \\
傳道書 1:1-11 & \hyperref[sec:Xoz1f5_xHTQ]{日光之下當思考 (傳道書 1\_1-11) - 譚文倩莊傳道} & 2025-08-24 & \href{https://youtube.com/watch?v=Xoz1f5-xHTQ}{\texttt{Xoz1f5-xHTQ}} \\
使徒行傳 8:1-40 & \hyperref[sec:rZsT_71UjNw]{流散到地極 (使徒行傳8\_1-40) - 陳崇基牧師} & 2025-08-24 & \href{https://youtube.com/watch?v=rZsT_71UjNw}{\texttt{rZsT\_71UjNw}} \\
使徒行傳 6:1-15 & \hyperref[sec:GgAMYWfpmPk]{從磨擦到拓展 (使徒行傳6\_1-15) - 陳崇基牧師} & 2025-08-30 & \href{https://youtube.com/watch?v=GgAMYWfpmPk}{\texttt{GgAMYWfpmPk}} \\
使徒行傳 4:31-5:16 & \hyperref[sec:xrB_pn8UZeQ]{聖靈二次降臨 (使徒行傳4\_31-5\_16) - 陳崇基牧師} & 2025-08-30 & \href{https://youtube.com/watch?v=xrB-pn8UZeQ}{\texttt{xrB-pn8UZeQ}} \\
詩篇 19:1-14 & \hyperref[sec:EMLJsLtjt_w]{The God Who Speaks (Ps 19) - Rev. Mark Chan 對你說話的神 (詩篇19\_1-14) - 陳崇基牧師} & 2025-09-01 & \href{https://youtube.com/watch?v=EMLJsLtjt_w}{\texttt{EMLJsLtjt\_w}} \\
路加福音 4:1-13 & \hyperref[sec:8R1ydESgnbc]{屬靈的謊言 (路加福音 4\_1-13) - 梁紫健傳道} & 2025-09-01 & \href{https://youtube.com/watch?v=8R1ydESgnbc}{\texttt{8R1ydESgnbc}} \\
馬太福音 7:13 & \hyperref[sec:Iq3SYwSRmiI]{What Would You Choose? (Matthew 7\_13 - 27) - Chee Yie Jong 你會怎樣選擇?(馬太福音 7\_13-27) - 楊志溢傳道} & 2025-09-08 & \href{https://youtube.com/watch?v=Iq3SYwSRmiI}{\texttt{Iq3SYwSRmiI}} \\
傳道書 1:12-18 & \hyperref[sec:0No_EkcZeRE]{日光之下再思考 (傳道書 1\_12-18) - 譚文倩莊傳道} & 2025-09-08 & \href{https://youtube.com/watch?v=0No_EkcZeRE}{\texttt{0No\_EkcZeRE}} \\
使徒行傳 9:1-43 & \hyperref[sec:dx4StTt36lE]{不可預期的福音 (使徒行傳 9\_1-43) - 陳崇基牧師} & 2025-09-20 & \href{https://youtube.com/watch?v=dx4StTt36lE}{\texttt{dx4StTt36lE}} \\
路加福音 8:26-39 & \hyperref[sec:lW1Rt3jzEUM]{Who is your enemy (Luke 8\_26-39) - Tsz Kin Leung 誰是敵人 (路加福音 8: 26-39) - 梁紫健傳道} & 2025-09-23 & \href{https://youtube.com/watch?v=lW1Rt3jzEUM}{\texttt{lW1Rt3jzEUM}} \\
使徒行傳 9:10-19 & \hyperref[sec:6VvjU9_F7kQ]{歡迎從不易 (使徒行傳 9\_10-19) - 陳崇基牧師} & 2025-09-23 & \href{https://youtube.com/watch?v=6VvjU9_F7kQ}{\texttt{6VvjU9\_F7kQ}} \\
以弗所書 1:3-14 & \hyperref[sec:7m42sBmfTkI]{You Are Blessed (Ephesians 1\_3-14) - Paul Zhang 你是有福的人 (以弗所書 1\_3-14) - 張威傳道} & 2025-10-01 & \href{https://youtube.com/watch?v=7m42sBmfTkI}{\texttt{7m42sBmfTkI}} \\
傳道書 2:1-11 & \hyperref[sec:9E6MnSM4U0Q]{抗疫虛空 (傳道書 2\_1-11) - 譚文倩莊傳道} & 2025-10-01 & \href{https://youtube.com/watch?v=9E6MnSM4U0Q}{\texttt{9E6MnSM4U0Q}} \\
詩篇 63:1-11 & \hyperref[sec:Dh5pU8Y4jPw]{Seeking God (Psalm 63) - Rev. Mark Chan 尋求神 (詩篇63\_1-11) - 陳祟基牧師} & 2025-10-14 & \href{https://youtube.com/watch?v=Dh5pU8Y4jPw}{\texttt{Dh5pU8Y4jPw}} \\
使徒行傳 13:1-3 & \hyperref[sec:jhQYPPMz5JE]{向外的教會 (使徒行傳13\_1-3) - 陳崇基牧師} & 2025-10-14 & \href{https://youtube.com/watch?v=jhQYPPMz5JE}{\texttt{jhQYPPMz5JE}} \\
使徒行傳 16:6-10 & \hyperref[sec:ZCODt6kOCdM]{不似預期的福音 (使徒行傳 16\_6-10) - 陳崇基牧師} & 2025-10-29 & \href{https://youtube.com/watch?v=ZCODt6kOCdM}{\texttt{ZCODt6kOCdM}} \\
士師記 1:1-2:23 & \hyperref[sec:Qv1tOecnO_s]{Are We Becoming More Like the World? (Judges 1-2) - Chee Yie Jong 我們是否越來越像這個世界 (士師記 1\_1-2\_23) - 楊志溢傳道} & 2025-11-03 & \href{https://youtube.com/watch?v=Qv1tOecnO-s}{\texttt{Qv1tOecnO-s}} \\
以弗所書 2:1-10 & \hyperref[sec:0vNwJDwewOw]{From Death to Life (Eph 2\_1-10) - Paul Zhang 出死入生 (以弗所書2\_1-10) - 張威傳道} & 2025-11-03 & \href{https://youtube.com/watch?v=0vNwJDwewOw}{\texttt{0vNwJDwewOw}} \\
出埃及記 1:8-22 & \hyperref[sec:_S6mqFk4KdE]{上帝的應許與保守(出埃及記1\_8-22) - 江芷敏傳道} & 2025-11-03 & \href{https://youtube.com/watch?v=_S6mqFk4KdE}{\texttt{\_S6mqFk4KdE}} \\
路加福音 8:26-39 & \hyperref[sec:19Ss941ItL4]{誰是敵人? (路加福音8\_26-39) - 梁紫健傳道} & 2025-11-03 & \href{https://youtube.com/watch?v=19Ss941ItL4}{\texttt{19Ss941ItL4}} \\
詩篇 73:1-28 & \hyperref[sec:SzBowSpLGuQ]{Why do the Wicked Prosper? (Psalm 73) - Rev Mark Chan 惡人亨通時 (詩篇73\_1-28) - 陳崇基牧師} & 2025-11-10 & \href{https://youtube.com/watch?v=SzBowSpLGuQ}{\texttt{SzBowSpLGuQ}} \\
傳道書 2:12-26 & \hyperref[sec:I3sM6YZywrk]{感謝天父賜我飲食 (傳道書2\_12-26) - 譚文倩莊傳道} & 2025-11-10 & \href{https://youtube.com/watch?v=I3sM6YZywrk}{\texttt{I3sM6YZywrk}} \\
詩篇 8:1-9 & \hyperref[sec:lX_7jxrhzjs]{True Worship (Psalm 8) - Rev. Mark Chan 真實的敬拜 (詩篇8\_1-9) - 陳崇基牧師} & 2025-11-18 & \href{https://youtube.com/watch?v=lX-7jxrhzjs}{\texttt{lX-7jxrhzjs}} \\
使徒行傳 17:1-34 & \hyperref[sec:mJ4c0qfZcx0]{在異地作見證(使徒行傳17\_1-34) - 陳崇基牧師} & 2025-11-18 & \href{https://youtube.com/watch?v=mJ4c0qfZcx0}{\texttt{mJ4c0qfZcx0}} \\
以弗所書 2:11-22 & \hyperref[sec:G1IwEMMFCQQ]{One in Christ (Eph 2\_11-22) - Paul Zhang 在基督裡合一 (以弗所書2\_11-22) - 張威傳道} & 2025-12-01 & \href{https://youtube.com/watch?v=G1IwEMMFCQQ}{\texttt{G1IwEMMFCQQ}} \\
路加福音 4:16-30 & \hyperref[sec:45ohI_TBLkg]{逆轉的人生(路加福音4\_16-30) - 梁紫健傳道} & 2025-12-01 & \href{https://youtube.com/watch?v=45ohI-TBLkg}{\texttt{45ohI-TBLkg}} \\
士師記 2:6-23 & \hyperref[sec:VcgHfBIcHMQ]{Enduring Faith (Judges 2\_6-23) - Chee Yie Jong 堅韌的信仰(士師記2\_6-23) - 楊志溢傳道} & 2025-12-08 & \href{https://youtube.com/watch?v=VcgHfBIcHMQ}{\texttt{VcgHfBIcHMQ}} \\
出埃及記 2:23-25 & \hyperref[sec:8eaa7ARgWZg]{願意靠近我們的神(出埃及記2\_23-25) - 江芷敏傳道} & 2025-12-08 & \href{https://youtube.com/watch?v=8eaa7ARgWZg}{\texttt{8eaa7ARgWZg}} \\
馬太福音 2:1-23 & \hyperref[sec:3dbOIx2ZjUI]{The Worship of the Magi (Matt 2) - Rev. Mark Chan 東方博士朝拜主(馬太福音2\_1-23) - 陳崇基牧師} & 2025-12-15 & \href{https://youtube.com/watch?v=3dbOIx2ZjUI}{\texttt{3dbOIx2ZjUI}} \\
馬太福音 2:1-23 & \hyperref[sec:f1NDkt3u9eY]{東方博士朝拜主 (馬太福音2\_1-23) - 陳崇基牧師} & 2025-12-15 & \href{https://youtube.com/watch?v=f1NDkt3u9eY}{\texttt{f1NDkt3u9eY}} \\
詩篇 119:1-176 & \hyperref[sec:4msYGc1nvL4]{Loving God's Word (Psalm 119) - Rev. Mark Chan 愛神的話語 (詩篇119\_1-176) - 陳崇基牧師} & 2025-12-29 & \href{https://youtube.com/watch?v=4msYGc1nvL4}{\texttt{4msYGc1nvL4}} \\
約翰福音 3:16 & \hyperref[sec:wPcDqaSy0yg]{神愛世人 (約翰福音3\_16) - 陳崇基牧師} & 2025-12-29 & \href{https://youtube.com/watch?v=wPcDqaSy0yg}{\texttt{wPcDqaSy0yg}} \\
路加福音 2:1-14 & \hyperref[sec:VP0gKdg8cek]{誰是我王( 路加福音 2\_1-14) - 梁紫健傳道} & 2025-12-29 & \href{https://youtube.com/watch?v=VP0gKdg8cek}{\texttt{VP0gKdg8cek}} \\
士師記 3:1-31 & \hyperref[sec:kGtKq4doSFU]{God’s Holy Love (Judges 3 ) - Chee Yie Jong 神聖的愛 (士師記3\_1-31) - 楊志溢傳道} & 2026-01-14 & \href{https://youtube.com/watch?v=kGtKq4doSFU}{\texttt{kGtKq4doSFU}} \\
詩篇 133:1-3 & \hyperref[sec:6JWiBdm7MIc]{Pilgrim's Community (Psalm 133) - Rev. Mark Chan 上行的群體(詩篇133\_1-3) - 陳崇基牧師} & 2026-01-14 & \href{https://youtube.com/watch?v=6JWiBdm7MIc}{\texttt{6JWiBdm7MIc}} \\
詩篇 133:1-3 & \hyperref[sec:4c2vxrJygVs]{上行的群體(詩篇133\_1-3) - 陳崇基牧師} & 2026-01-14 & \href{https://youtube.com/watch?v=4c2vxrJygVs}{\texttt{4c2vxrJygVs}} \\
詩篇 8:1-9 & \hyperref[sec:LP_MCyZeclI]{人算什麼?(詩篇 8\_1-9) - 陳崇基牧師} & 2026-01-14 & \href{https://youtube.com/watch?v=LP_MCyZeclI}{\texttt{LP\_MCyZeclI}} \\
約翰福音 1:1-5 & \hyperref[sec:n8Rs4GhmcQo]{現身的上帝(約翰福音1\_1-5) - 梁紫健傳道} & 2026-01-24 & \href{https://youtube.com/watch?v=n8Rs4GhmcQo}{\texttt{n8Rs4GhmcQo}} \\
 3:1-15 & \hyperref[sec:_Xk1Tbpr5MU]{難先要你做 (傳道書3\_1-15) - 譚文倩莊傳道} & 2026-01-24 & \href{https://youtube.com/watch?v=_Xk1Tbpr5MU}{\texttt{\_Xk1Tbpr5MU}} \\
  & \hyperref[sec:Ly_O4TyukTE]{Seen and Known (Psalm 139) - Rev. Mark Chan 被神認識 (詩139) - 陳崇基牧師} & 2026-02-02 & \href{https://youtube.com/watch?v=Ly-O4TyukTE}{\texttt{Ly-O4TyukTE}} \\
約翰福音 1:1-5 & \hyperref[sec:L_CPDfFPb5k]{現身的上帝(約翰福音1\_1-5) - 梁紫健傳道} & 2026-02-02 & \href{https://youtube.com/watch?v=L-CPDfFPb5k}{\texttt{L-CPDfFPb5k}} \\
詩篇  & \hyperref[sec:XrYioYwg95Q]{Be Still (Psalm 46)- Rev. Mark Chan 學習安靜(詩篇46) - 陳崇基牧師} & 2026-02-09 & \href{https://youtube.com/watch?v=XrYioYwg95Q}{\texttt{XrYioYwg95Q}} \\
詩篇  & \hyperref[sec:QVz_a2O_21A]{學習安靜(詩篇46) - 陳崇基牧師} & 2026-02-09 & \href{https://youtube.com/watch?v=QVz_a2O-21A}{\texttt{QVz\_a2O-21A}} \\
士師記  & \hyperref[sec:espYGx8NsQM]{The Trustworthy God (Judges 4) - Chee Yie Jong  信實的上帝 (士師記 4) - 楊志溢傳道} & 2026-02-16 & \href{https://youtube.com/watch?v=espYGx8NsQM}{\texttt{espYGx8NsQM}} \\
約翰福音 1:14-18 & \hyperref[sec:Ulr7EZgFyUo]{道成肉身 (約翰福音1\_14-18) - 梁紫健傳道} & 2026-02-16 & \href{https://youtube.com/watch?v=Ulr7EZgFyUo}{\texttt{Ulr7EZgFyUo}} \\
 19:11-16 & \hyperref[sec:xg_qWk2v3jg]{The King on the White Horse (Rev 19\_11-16) - Rev. Mark Chan 白馬之王 (啟19\_11-16) - 陳崇基牧師} & 2026-02-22 & \href{https://youtube.com/watch?v=xg_qWk2v3jg}{\texttt{xg\_qWk2v3jg}} \\
啟示錄 19:11-16 & \hyperref[sec:KdIPOrLMIS0]{白馬之王 (啟示錄19\_11-16) - 陳崇基牧師} & 2026-02-22 & \href{https://youtube.com/watch?v=KdIPOrLMIS0}{\texttt{KdIPOrLMIS0}} \\
詩篇 1:1-6 & \hyperref[sec:am5WL2bo5gk]{The Righteous Person (Psalm 1) - Rev. Mark Chan 義人的生命 (詩篇1\_1-6) - 陳崇基牧師} & 2026-03-02 & \href{https://youtube.com/watch?v=am5WL2bo5gk}{\texttt{am5WL2bo5gk}} \\
以斯帖記 8:3-9:16 & \hyperref[sec:mjg2Vzt9lP0]{神的救贖 人的責任 (以斯帖記 8\_3-9\_16) - 黃國棟醫生} & 2026-03-02 & \href{https://youtube.com/watch?v=mjg2Vzt9lP0}{\texttt{mjg2Vzt9lP0}} \\
詩篇  & \hyperref[sec:_lnAtpZRm3Q]{Holistic Worship (Psalm 150) - Rev. Mark Chan 全人的敬拜 (詩篇 150) - 陳崇基牧師} & 2026-03-09 & \href{https://youtube.com/watch?v=_lnAtpZRm3Q}{\texttt{\_lnAtpZRm3Q}} \\
詩篇  & \hyperref[sec:y65Q_8wRaJ8]{全人的敬拜 (詩篇 150) - 陳崇基牧師} & 2026-03-09 & \href{https://youtube.com/watch?v=y65Q_8wRaJ8}{\texttt{y65Q\_8wRaJ8}} \\
出埃及記 3:1-22 & \hyperref[sec:CVOZX9nsoIc]{呼召的關係 (出埃及記3\_1-22) - 江芷敏傳道} & 2026-03-16 & \href{https://youtube.com/watch?v=CVOZX9nsoIc}{\texttt{CVOZX9nsoIc}} \\
士師記 6:1-40 & \hyperref[sec:6DmNaa6pfms]{真實的敬拜 (士師記6\_1-40) - 楊志溢傳道} & 2026-03-16 & \href{https://youtube.com/watch?v=6DmNaa6pfms}{\texttt{6DmNaa6pfms}} \\
約翰福音 2:1-12 & \hyperref[sec:TQN008eNyL0]{Suprising Transformation - (John 2\_1-12) - Tsz Kin Leung 驚喜的轉化 (約翰福音2\_1-12) - 梁紫健傳道} & 2026-03-24 & \href{https://youtube.com/watch?v=TQN008eNyL0}{\texttt{TQN008eNyL0}} \\
約翰福音 5:1-18 & \hyperref[sec:a3FPlfJ6MYc]{屬靈的力量 (約翰福音 5\_1-18) - 鍾智愷傳道} & 2026-03-24 & \href{https://youtube.com/watch?v=a3FPlfJ6MYc}{\texttt{a3FPlfJ6MYc}} \\
使徒行傳 13:1-52 & \hyperref[sec:_FBi_OEMziA]{How to Keep Going? (Acts 13) - Rev. Mark Chan 如何繼續走下去? (使徒行傳13\_1-52) - 陳崇基牧師} & 2026-03-30 & \href{https://youtube.com/watch?v=_FBi_OEMziA}{\texttt{\_FBi\_OEMziA}} \\
約翰福音 2:1-12 & \hyperref[sec:ZqMpnK5BxrA]{行奇事的神 (約翰福音2\_1-12) - 梁紫健傳道} & 2026-03-30 & \href{https://youtube.com/watch?v=ZqMpnK5BxrA}{\texttt{ZqMpnK5BxrA}} \\
 3:12-4:3 & \hyperref[sec:a8Vgzt3GCx8]{公義缺席 (傳3\_12-4\_3) - 譚文倩莊傳道} & 2026-04-05 & \href{https://youtube.com/watch?v=a8Vgzt3GCx8}{\texttt{a8Vgzt3GCx8}} \\
 3:14-22 & \hyperref[sec:FQREQHip6Ag]{成為合乎主用的人 (啟3\_14-22) - 陳崇基牧師 Becoming Useful for God's Purpose (Rev 3\_14-22) Rev. Mark Chan} & 2026-04-05 & \href{https://youtube.com/watch?v=FQREQHip6Ag}{\texttt{FQREQHip6Ag}} \\
以弗所書 3:1-21 & \hyperref[sec:mR4JE8s_dsw]{復活節主日講道 (以弗所書3\_1-21) - 陳崇基牧師 Easter Sunday Sermon 2026 (Eph 3) - Rev. Mark Chan} & 2026-04-09 & \href{https://youtube.com/watch?v=mR4JE8s-dsw}{\texttt{mR4JE8s-dsw}} \\
創世記 22:15-19 & \hyperref[sec:FC_bOIHzaL8]{Sand (Gen 22\_15-19) - Rev. Mark Chan 海沙 (創世記22\_15-19) - 陳崇基牧師} & 2026-04-18 & \href{https://youtube.com/watch?v=FC_bOIHzaL8}{\texttt{FC\_bOIHzaL8}} \\
創世記 22:15-19 & \hyperref[sec:Lyb7YfDNUgY]{海沙 (創世記22\_15-19) - 陳崇基牧師} & 2026-04-18 & \href{https://youtube.com/watch?v=Lyb7YfDNUgY}{\texttt{Lyb7YfDNUgY}} \\
約翰福音 2:12-25 & \hyperref[sec:FHgJwAMW_uQ]{Cleanse your hearts of its "Cattle and Sheep" (Jn 2\_12-25) - Tsz Kin Leung 趕走心中的牛羊(約翰福音2\_12-25) - 梁紫健傳道} & 2026-04-28 & \href{https://youtube.com/watch?v=FHgJwAMW_uQ}{\texttt{FHgJwAMW\_uQ}} \\
約翰福音 2:12-25 & \hyperref[sec:YUWGLt45Wz8]{趕走心中的「牛羊」(約翰福音2\_12-25) - 梁紫健傳道} & 2026-04-28 & \href{https://youtube.com/watch?v=YUWGLt45Wz8}{\texttt{YUWGLt45Wz8}} \\
約翰福音 4:43-54 & \hyperref[sec:bMtCjPvaMwM]{你是耶穌的粉絲嗎 (約翰福音4\_43-54) - 譚偉明宣教士} & 2026-05-07 & \href{https://youtube.com/watch?v=bMtCjPvaMwM}{\texttt{bMtCjPvaMwM}} \\
以弗所書 3:1-6 & \hyperref[sec:3zbQTQ6NDLE]{你是誰?當福音重新定義你的人生 (以弗所書3\_1-6)Who Are You? When the Gospel redefines your life ( Eph 3\_1-6)} & 2026-05-07 & \href{https://youtube.com/watch?v=3zbQTQ6NDLE}{\texttt{3zbQTQ6NDLE}} \\
  & \hyperref[sec:6UvylJcofPU]{Craftsman of Broken Pieces (Gen 12-50) - Rev. Mark Chan 碎片的工匠 (創 12-50) - 陳崇基牧師} & 2026-05-10 & \href{https://youtube.com/watch?v=6UvylJcofPU}{\texttt{6UvylJcofPU}} \\
  & \hyperref[sec:9gxDgFyKyk4]{碎片的工匠 (創12-50) - 陳崇基牧師} & 2026-05-10 & \href{https://youtube.com/watch?v=9gxDgFyKyk4}{\texttt{9gxDgFyKyk4}} \\
  & \hyperref[sec:KUqZwRWtyog]{遊歷土耳其 - 老底嘉 (Laodicea) - 陳崇基牧師} & 2026-05-15 & \href{https://youtube.com/watch?v=KUqZwRWtyog}{\texttt{KUqZwRWtyog}} \\
士師記 6:11-16 & \hyperref[sec:jnGxkuaDf2M]{Faith over Fear (Judges 6\_11-16) - Chee Yie Jong 信心勝過恐懼 (士師記6\_11-16) - 楊志溢傳道} & 2026-05-18 & \href{https://youtube.com/watch?v=jnGxkuaDf2M}{\texttt{jnGxkuaDf2M}} \\
約翰福音 19:25-27 & \hyperref[sec:HpXR4dbZvbk]{母親 這是你的兒子 (約翰福音19\_25-27) - 梁紫健傳道} & 2026-05-18 & \href{https://youtube.com/watch?v=HpXR4dbZvbk}{\texttt{HpXR4dbZvbk}} \\
出埃及記 20:1-26 & \hyperref[sec:y_bUxa4hgT8]{Stone (Exod 20) - Rev. Mark Chan 石板 (出埃及記20\_1-26) - 陳崇基牧師} & 2026-05-24 & \href{https://youtube.com/watch?v=y_bUxa4hgT8}{\texttt{y\_bUxa4hgT8}} \\
出埃及記 20:1-26 & \hyperref[sec:OKPQ6OtPkus]{石板 (出埃及記20\_1-26) - 陳崇基牧師} & 2026-05-24 & \href{https://youtube.com/watch?v=OKPQ6OtPkus}{\texttt{OKPQ6OtPkus}} \\
約翰福音 4:1-26 & \hyperref[sec:wv60KUfih4o]{井旁相遇 (約翰福音4\_1-26) - 梁紫健傳道 An Encounter at the Well (John 4\_1-26) - Tsz Kin Leung} & 2026-05-30 & \href{https://youtube.com/watch?v=wv60KUfih4o}{\texttt{wv60KUfih4o}} \\
傳道書 4:4-5:7 & \hyperref[sec:OAEc6rF7Lr4]{關係的失衡(傳道書4\_4-5\_7) - 譚文倩莊傳道} & 2026-05-30 & \href{https://youtube.com/watch?v=OAEc6rF7Lr4}{\texttt{OAEc6rF7Lr4}} \\
利未記 23:39-43 & \hyperref[sec:augzRkvXpyY]{Booths (Lev 23\_39-43) Rev. Mark Chan 住棚 (利未記23\_39-43) - 陳崇基牧師} & 2026-07-12 & \href{https://youtube.com/watch?v=augzRkvXpyY}{\texttt{augzRkvXpyY}} \\
出埃及記  & \hyperref[sec:pCh24jxVqNk]{The Tent (Exod 25-40) - Rev. Mark Chan 會幕 (出埃及記25-40) - 陳崇基牧師} & 2026-07-12 & \href{https://youtube.com/watch?v=pCh24jxVqNk}{\texttt{pCh24jxVqNk}} \\
士師記  & \hyperref[sec:Cofy9lMF2CI]{When Less is More (Judge 7) - Chee Yie Jong 少即是少 (士師記 7) - 楊志溢傳道} & 2026-07-12 & \href{https://youtube.com/watch?v=Cofy9lMF2CI}{\texttt{Cofy9lMF2CI}} \\
約翰福音 10:27-30 & \hyperref[sec:qC9_GTWJyK4]{一雙永不鬆手的手 (約翰福音10\_27-30) - 梁紫健傳道 The Hands that Never Let Go (John 10\_27-30) - Tsz Kin Leung} & 2026-07-12 & \href{https://youtube.com/watch?v=qC9_GTWJyK4}{\texttt{qC9\_GTWJyK4}} \\
利未記 23:39-43 & \hyperref[sec:xpOSI0B_Qx0]{住棚 (利未記 23\_39-43) - 陳崇基牧師} & 2026-07-12 & \href{https://youtube.com/watch?v=xpOSI0B-Qx0}{\texttt{xpOSI0B-Qx0}} \\
  & \hyperref[sec:IBF8VFRUG4s]{再思女性的事奉 - 陳崇基牧師} & 2026-07-12 & \href{https://youtube.com/watch?v=IBF8VFRUG4s}{\texttt{IBF8VFRUG4s}} \\
 3:7-13 & \hyperref[sec:MlRDdW70h9M]{在基督裡坦然無懼(弗3\_7-13)張威傳道 Boldness in Christ (Eph 3\_7-13) Paul Zhang} & 2026-07-12 & \href{https://youtube.com/watch?v=MlRDdW70h9M}{\texttt{MlRDdW70h9M}} \\
約翰福音 4:3-14 & \hyperref[sec:YT4UObZ7V44]{放下水桶 (約翰福音4\_3-14) - 梁紫健傳道} & 2026-07-12 & \href{https://youtube.com/watch?v=YT4UObZ7V44}{\texttt{YT4UObZ7V44}} \\
出埃及記  & \hyperref[sec:rfQ3ZwOlzM4]{會幕 (出埃及記25-40) - 陳崇基牧師} & 2026-07-12 & \href{https://youtube.com/watch?v=rfQ3ZwOlzM4}{\texttt{rfQ3ZwOlzM4}} \\
出埃及記  & \hyperref[sec:jsiErX4TZXs]{生命的定位 (出埃及記 6) - 江芷敏傳道} & 2026-07-12 & \href{https://youtube.com/watch?v=jsiErX4TZXs}{\texttt{jsiErX4TZXs}} \\
以弗所書 3:14-21 & \hyperref[sec:DRGVYhPEXTM]{被基督的愛充滿 (以弗所書3\_14-21) - 張威傳道 Filled with the Love of Christ (Eph3\_14-21) - Paul Zhang} & 2026-07-12 & \href{https://youtube.com/watch?v=DRGVYhPEXTM}{\texttt{DRGVYhPEXTM}} \\
傳道書 5:8-6:8 & \hyperref[sec:2jdjy7oOq1s]{錢值幾錢?(傳道書5\_8-6\_8) - 譚文倩莊傳道} & 2026-07-12 & \href{https://youtube.com/watch?v=2jdjy7oOq1s}{\texttt{2jdjy7oOq1s}} \\
以弗所書 4:1-7 & \hyperref[sec:3LIoE3b6SXA]{不只是名字,更是生命 (以弗所書 4\_1-7) - 張威傳道 Not Just a Name But Life - Paul Zhang} & 2026-07-20 & \href{https://youtube.com/watch?v=3LIoE3b6SXA}{\texttt{3LIoE3b6SXA}} \\
哥林多前書 11:2-16 & \hyperref[sec:8hcsYSy4IhI]{再思女性的事奉 (哥林多前書11\_2-16) - 陳崇基牧師} & 2026-07-20 & \href{https://youtube.com/watch?v=8hcsYSy4IhI}{\texttt{8hcsYSy4IhI}} \\
\end{xltabular}
}
\newpage



\section{}
\label{sec:Kg3cX9_FXZI}
\textbf{3/30/18 受苦節崇拜  Good Friday Service}
\newline
\newline
連結: \href{https://youtube.com/watch?v=Kg3cX9_FXZI}{\texttt{https://youtube.com/watch?v=Kg3cX9\_FXZI}} ~~~~ 語音日期: 2018-04-01
\newline
\newline
\hyperref[sec:code]{\small{< < < PREV SERMON < < <}}
~
\hyperref[sec:index_chronic]{\small{[返順時目]}}
~
\hyperref[sec:index_scriptual]{\small{[返順卷目]}}
~
\hyperref[sec:Z07LUlzvgG0]{\small{> > > NEXT SERMON > > >}}
\newline
\newline
$^{1}$受苦受死的情況.
聖經記載.
一到早晨.
祭司長和長老.
問事全公會的人.
大家相議.
就把耶穌捆綁.
解去交給彼拉多.
每逢這節期.
巡撫照眾人所求的.
釋放一個囚犯給他們.
指示祭司長.
亦挑唆眾人.
寧可釋放巴拉巴給他們.
把耶穌鞭打了.
交給人釘十字架.
冰釘把耶穌帶進衙門院裡.
叫齊了全營的兵.
他們給他穿上紙袍.
又用荊棘篇做冠冕.
給他戴上.
就慶賀他說.
恭喜猶太人的王.
又拿一根圍紙.
打他的頭.
吐汗抹在他面上.
然後把他扣上.
吐汗抹在他面上.
屈膝拜他.
戲弄完了.
就給他脫了紙袍.
仍穿上了他自己的衣服.
帶他出去.
要釘十字架.
從那裡經過的人.
辱罵他要穿頭雪.
哈!你這拆毀聖殿.
三日又建造起來的.
可以夠自己.
從十字架上哈來吧.

$^{41}$祭司長和文士.
也試圖戲弄他.
彼此說.
他救了別人不能救自己.
從悟證至新初.
遍地都是黑暗了.
新初的時候.
耶穌大聲喊著說.
以羅爾拉瑪撒巴各帶來.
翻出來就說.
我的神我的神.
為什麼離棄我.
旁邊站著的人.
有的聽見就說.
看吶!他叫以利亞呢.
有一個人跑去.
把海洋沾滿了醋.
綁在位置上送給他.
說且等著.
看以利亞來不來把他取下.
耶穌大聲喊叫.
戲就斷了.
電影的萬紙從上到下.
裂開兩半.
對面站著的白夫長.
看見耶穌這樣喊叫斷氣.
就說.
這人真是神的兒子.
我怎樣扮他呢.
這句話.
是來自經文第十二節末句.
是由一個羅馬高官所說.
他就是彼拉多.
他管轄猶大地.
他是總督.
他的權力相當大.
可以指揮軍隊.
管理財政.
徵收稅項.
甚至設立大祭司.

$^{81}$審斷案件.
他更操有生殺之訣.
我怎樣扮他呢.
他.
是指耶穌基督.
扮.
就是處置對待.
今晚我們這個受難節的聚會.
我們就從剛才所看的經文.
再重溫一下.
耶穌受苦的情形.
給我們一些反省.
提醒.
和提醒.
我們要如何去承受.
受苦的痛苦.
我們要如何去承受.
有些反省.
有些提醒.
讓我們愛主更多.
我們先有個簡單的祈禱.
天父安靜我們的心.
我們要對準你.
思想主耶穌.
在十字架上的大愛.
還有我們.
是怎樣扮你.
對待你.
昔日的人這樣對待你.
今天我們也有類似的情況.
求你光照.
又提醒我們.
將以後的時間交託.
用你自己的話語再一次提醒.
和.
勉勵我們.
祈禱奉主耶穌基督的名.
阿們.
我怎樣扮他呢.
有很多人對耶穌.

$^{121}$有不同的態度和表現.
今天我們找其中五項的情況來看看.
第一樣.
我們要如何去承受.
第一樣.
就是大家相依.
你看看.
第一節說了.
按當時猶太人的律法.
需要先有人告發.
來作見證.
那件案才可以成立.
才可以開審.
但你看耶穌這件案.
犯人說是耶穌.
但未有人正式告發.
又沒有見證人出庭.
使這件案可以成立.
所以耶穌這件審訊.
可以當是非法審訊.
非法被捉拿.
而公會實際上.
在當時是沒有.
定人死罪的權柄.
定死罪的權柄.
唯有羅馬政府發出.
而且由羅馬政府.
來作出審訊.
而公會.
在當時是沒有.
定人死罪的權柄.
而公會.
在當時是沒有.
定人死罪的權柄.
而公會.
在當時是沒有.
定人死罪的權柄.
而公會.
在當時是沒有.
定人死罪的權柄.

$^{161}$而公會.
在當時是沒有.
定人死罪的權柄.
而公會.
在當時是沒有.
定人死罪的權柄.
而公會.
在當時是沒有.
定人死罪的權柄.
而公會.
在當時是沒有.
定人死罪的權柄.
而公會.
在當時是沒有.
定人死罪的權柄.
而公會.
在當時是沒有.
定人死罪的權柄.
公會為了.
所謂按章辦事.
所謂急不及待.
一大清早就召開會議.
要確定.
昨晚所出的決議審訊.
馬可帶著一個.
沉痛的心情.
講出.
全公會的人.
是一同相議.
是大家相議.
這次所謂正式的會議.
原本應該全體.
有71位會員.
都不需要全部出席.
因為我們很熟悉的.
亞利馬太約瑟.
和尼哥底姆.
我們在這裡都明顯看到.
他們有在場的證據.
這次的會議.

$^{201}$必須將耶穌的案件.
有所判決.
在四福音的.
另一個作者.
馬太就是這樣說.
他用了一個很清楚的描述.
他說大家相議.
要致死耶穌.
我們特別要注意.
公會裡的成員.
是有些什麼人呢.
你看第一節的時候.
你就看到有什麼人.
有祭司長.
有長老和文士.
這一班的人.
正是會中的領袖.
是屬靈的長者.
他們是通了聖經律法的人.
是信仰的忠奸分子.
但是他們.
公然的橫行無忌.
不理良心的公正.
竟然被邪惡和狠毒所蒙蔽.
立心不良.
要致死耶穌.
謀害耶穌.
他們不容許耶穌的存在.
這實在是可悲.
以及可憐的光景.
當時的猶太人.
當時的猶太國.
的黑暗.
可以說是到了一個極點.
連那些有信仰的人.
那些是民間的領袖.
屬靈長者.
竟然落到這樣的地步.
聖經其實是有記載.
這班人.

$^{241}$心裡面是充滿惡毒和仇恨的.
他們由起初的控訴.
只是很單單.
一項褻瀆的罪名.
即是侮辱上帝.
開始.
但後來一路演變.
變成了幾項嚴重的罪名.
包括.
煽動民眾.
不准人民納稅.
被蓋殺.
還自稱為皇.
盧伽福音23章1到2.
只有說.
他們企圖渲染成一個.
很嚴重的政治理由.
唯恐彼拉多總督.
充耳不聞.
一個公正.
無私的工會.
竟然變成了陰謀權術.
迫害義人的機構.
一批信仰的忠正分子.
竟然敗壞到.
是害人陰毒.
陷害忠良.
敵黨基督.
關鍵原因在哪裡呢.
你看第十節.
這句話.
披露有關的重要線索.
他.
就是指彼拉多.
願曉得濟師長.
是因為嫉妒.
才把耶穌解了來.
嫉妒.
原來是嫉妒.
曾元說.

$^{281}$嫉妒是骨中的扭爛.
亞國書3章16節.
何處有人.
不知不覺.
被人.
嫉妒.
在16節說.
何處有嫉妒紛爭.
就在何處有擾亂.
和各樣的壞事.
嫉妒原來可以.
蠶食人.
可以腐化人.
可以擾亂.
和成就很多壞事.
嫉妒本身.
就是一個自我中心.
嫉妒其實.
就是心中只有私意.
不用納到嫉妒.
所以.
當神沒有到.
弟兄姊妹.
今天我們心裡.
有沒有被一些嫉妒.
自我中心.
所影響.
甚至不經意地.
把嫉妒推到一邊.
偏行幾路.
很多時候.
會出現一個情況.
就是變得.
任意妄為.
近代猜拳.
運動之父.
威廉克里.
他自己看聖經.
從聖經裡面領受到.
要向普世宣教的教訓.

$^{321}$在當時.
正是18世紀.
教會裡面一班.
屬靈長者和領袖.
卻是沒有突破的體會.
在1789年.
浸信會的同工.
開大會.
大家商議.
教會未來要做些甚麼.
發展是怎樣.
這位年輕的同工.
威廉克里.
就想提出.
教會應該對異教徒.
有一個宣教的服事.
當時不少的教會領袖.
很多的長者.
都是看低這個年輕人.
覺得這個人.
懂得甚麼.
不看低他.
結果當時的大會主席.
經過大家商議之後.
公然的對威廉克里說.
年輕人.
坐下吧.
當時異教徒想主的時候.
不著你和我替他操心.
這班沒冷水的小朋友.
這班沒冷水的說話.
就是從這個同工大會.
是大家一起商議.
所得出的結果.
他們根本是對主耶穌的大使命.
注之不理.
對神的旨意都不徹底的進行.
原因是基於他們.
自以為是.
他們自我中心的執著.

$^{361}$不抵得這個年輕人.
提出這種的真知灼見.
大大的失去了海外宣教的良知.
所以我們要留意.
當我們心中有一些執著.
有一些私意.
有一些看低人.
甚至嫉妒的時候.
就會被自我中心.
自以為是這個偏執的思想.
所盤據著.
很容易.
將一些正直.
真理和神.
推到一邊.
走自己的路.
走自己的道路.
走自己認為要走的路.
漠視了神的旨意.
所以我怎麼.
窩扮他呢.
第一樣.
就是大家商議.
我們看第二樣.
是寧可.
第十一和十五節.
這是第二樣.
對耶穌的態度.
這個寧可.
就是要棄絕耶穌.
寧可要釋放.
窮凶極惡的強盜.
巴拉伯.
不選擇耶穌.
關於巴拉伯的事.
除了福音書的記載之外.
資料真的不是很多.
可能由於他是一個.
薄煞空巷.
而且是犯了.

$^{401}$王法的人.
但是他闖出了一個名堂.
據說在當時的社會裡面.
有一班.
很熱血的民族主義份子.
隨時找機會.
煽動群眾暴動.
這一班人.
經常準備一些短刀.
放在外衣裡面.
見機行事.
可以說是當時的.
熱血極端的民族革命者.
他們企圖用不擇的手段.
用兇殺.
或者暗殺的方法.
達到他們的目的.
他們的目的.
暗殺的方法.
達到他們的目的.
當時的巴拉伯.
或者是一個這樣的勇士.
專以暗殺出名.
他雖然是一個匪徒.
但是他滿腔熱血.
為解救民族.
這種請命.
所以有些人反過來覺得.
他很愛國.
用烈士.
愛國之士去稱呼他.
有些群眾.
對他也有好感.
我們要想一想.
一個星期前.
群眾夾道歡迎.
騎著小驢.
的耶穌入城.
入耶路撒冷.
為什麼一個星期之後.

$^{441}$群眾那麼善忘.
又要呼喊要釘耶穌.
在十字架上.
原因很簡單.
因為一個星期前.
那班在路上.
歡迎耶穌的人.
一個星期之後.
已經分散在聖地.
各處.
而呼喊要釘在十字架上的那班人.
是另外的一班群眾.
你想想當時.
工會不依法行事.
在夜間.
匆匆地審訊耶穌.
當時有些門徒已經嚇到.
逃跑.
所以當時.
這些呼喊的民眾.
裡面.
很少很少.
是耶穌的擁躉支持者.
那這班民眾成員.
是哪些人呢.
大部分大部分.
是巴拉巴的擁護者.
他們全心前來.
要求釋放巴拉巴.
所以當他們.
看到另一個犯人.
耶穌出現的時候.
而彼拉多似乎暗示.
要釋放耶穌.
這樣的心意的時候.
他們大失所望.
立刻高聲呼喊.
釘耶穌在十字架.
而祭司看見這個千載難逢的機會.
更加煽動.

$^{481}$一些群眾.
拼命呼喊.
寧可要巴拉巴.
你看這節經文裡.
提及群眾.
寧可要巴拉巴.
這個強度而不揀選耶穌.
其中一個原因.
原來是被.
祭司長挑唆.
挑唆.
原文有勸導.
或者說服的意思.
今天有些人.
不肯說耶穌.
不肯揀選耶穌.
原來正正是受了外界.
不少人的勸導說服.
甚至挑唆.
有沒有聽過.
類似的說話.
有些人想勸導一些人.
信耶穌.
有些人反過來叫他們不要信.
你信什麼耶穌.
信耶穌有什麼好.
條條大路都可以.
何須要信耶穌.
有些人說.
信耶穌.
不好,綁手又綁腳.
這樣不可以.
那樣又不可以.
沒有歌聲,沒有自由.
有些人說.
信耶穌又要.
參加聚會,又要讀經.
祈禱,聽到滿極.
你又聽過類似的說話.
信耶穌.

$^{521}$你可以信.
不過不要信得太投入.
偶爾回教會就可以了.
偶爾隨便奉獻.
主日崇拜.
都不用每個星期去.
聖餐的時候才去都可以.
這些所謂的勸導.
這些所謂的說服.
無靈.
其實是跳梭.
使不少人走歪了.
失去了豐盛喜樂積極的.
屬靈生命.
今天有些人說.
信了耶穌.
可能被世界上很多東西吸引了.
為了事業.
為了愛情.
為了家庭子女而孫.
更加被一些情慾.
私意.
跳梭.
擺弄.
寧可要這個世界.
寧可要這個物質享樂.
地位事業.
都不選擇耶穌.
寧可睡多點.
寧可出去吃喝玩樂.
寧可去多點街.
寧可.
主日崇拜的時候.
約見親友去喝茶或者活動.
都不想回教會.
寧可參與外面的一些事.
都不肯去事奉.
其實這是基督徒.
一個價值觀的問題.
是你和我的人生優先次序的問題.

$^{561}$是眾者過.
輕那過.
這個問題.
第十五節說到.
彼拉多.
他有一個優先次序.
他說為要叫人喜悅.
他為了要這樣.
他就要.
為了要這樣.
今天我們是討神喜悅.
還是討人歡心.
求主幫助你和我.
不要被這個世界.
不要被撒旦老鵝.
肉體絲肉.
所跳梭.
我們仍然是寧可要耶穌.
我們寧可要愛主更深.
我們要有警覺.
去分別一些人說的話.
有沒有一些跳梭.
令到我們慢慢.
走來主.
第三我們看另一樣.
有些人對耶穌的態度.
就是戲弄.
雞肖和肉麻.
這是第三個態度.
在肉麻戲弄.
耶穌.
也都帶來鞭打耶穌.
羅馬人的鞭打.
是一個很可怕的號形.
受鞭打的人.
要脫下衣服.
將身體綁在柱子上.
或者是苦服.
受困.
受綁.

$^{601}$甚至有些人按在地上.
露出背部.
然後任人打.
用來鞭打的.
是一束皮帶的短辮.
幾條皮帶.
是傷了.
有鐵釘.
鉛鐵.
骨片.
或者骨頭.
所以.
你聽過俗稱叫做九尾鞭.
一發下去.
皮肉都會扯開.
皮開肉爛.
以及有些勁的時候露出骨.
其實抽打的時候.
最初重實的皮條.
割破皮的表皮.
但你再打幾下.
抽打就切進了皮下的組織.
的深層.
造成皮膚.
微絲血管和小靜脈的破裂.
的破裂.
而流血.
如果你再繼續打下去.
肌肉再深入一點.
連較大的動脈.
受抽打影響.
破裂.
一路流血.
曾經有記載.
有些人打高一點.
不小心將眼球也抽出來.
扯開.
更是有些人.
頂不住.
當場死去.

$^{641}$被打.
被打完之後.
有些人搞到殘廢.
或者精神崩潰.
神經失常.
很少很少人被鞭打之後.
可以若無其事.
仍然很輕鬆.
甚至神智仍然很清醒.
兄姐妹.
我們的主耶穌.
正是在當時.
被人這樣抽打.
這個鞭刑.
這個仇父的待遇.
就因為他要擔當.
世人的罪.
要為你和我的過法.
而受這個折磨和刑罰.
我們看不看到.
主耶穌的負心.
丁丁當時奉命.
將耶穌帶到衙門.
衙門即是指.
總督府.
他們當時不能夠.
隨便走來走去.
按《日行福音》十九章四節說.
吩咐丁丁.
帶耶穌進廷院.
之前彼拉多已經進去了.
然後他們在戲弄.
之後.
彼拉多又出來了.
所以鞭打和戲弄.
都是彼拉多的主意.
鞭打的時候.
彼得和工會眾人.
一定看得清清楚楚.
彼拉多鞭打戲弄耶穌的用意.

$^{681}$原本是想著.
讓猶太人和會員.
可以在這裡.
讓猶太人和會員.
看到這個悲戰可憐的人.
軟弱的耶穌.
根本做不出什麼大事.
他怎麼可以做王呢.
他是一個這樣軟弱的人.
無損無害.
沒什麼大作為.
何必要釘死他呢.
彼拉多其實是想.
這樣的可憐的情景.
使猶太人改變他們的主意.
而釋放他.
但是這個如意算盤.
發生不了效力.
因為在約翰福音19章1到7節.
類似是這樣說.
鞭打戲弄.
只是代表一種對待耶穌的態度.
今天我們看到一些.
未信耶穌的人.
極盡詆毀.
還有喜樂嘲笑.
淺達.
很多糟質.
信耶穌的人.
其實都是變相的戲弄.
和一個的.
技巧.
特別.
大家都知道.
現在是末世的日子.
未來的日子.
基督徒有沒有心準備為主受苦呢.
這個需要更多的裝備.
鞭釘戲弄.
耶穌之後.

$^{721}$其中一樣要穿上紙袍.
耶穌原來穿的衣服要脫光.
然後再穿紙袍.
但是你記得他打完之後.
他已經被皮開肉爛.
現在那個紙袍是有硬.
是有粗的紙袍.
血又被皮開肉爛.
這樣的粗.
這樣的壓在他身上.
更加好痛.
你可以想像得到.
還有皇帝有皇冠荊棘.
他們就找了那個荊棘.
打圈.
小小的圈.
然後用力壓在耶穌的頭上.
那些刺的荊棘.
將耶穌的頭劃破.
血就流出來.
你看到一些名畫.
有時候看到就是這樣.
在耶穌很痛苦的時候.
那些鞭釘.
卻是在笑嘻嘻的.
嬉皮笑臉的來取笑耶穌.
是落耶穌.
恭喜猶太人都會黃啊.
他們不曉得耶穌真是黃.
真是黃.
萬王之黃萬主之主.
他是上帝的兒子.
現在主耶穌落得這樣的下場.
真是令人心痛.
鞭釘又拿了一個位置.
一些很長又硬的植物給他.
象徵這個是皇杖皇權.
起初就放在耶穌的手裡.
作為扮演.
但是鞭釘在耶穌面前.

$^{761}$氣流的時候.
竟然把這個位置.
從耶穌的手上奪去.
然後用這個長長硬硬的位置.
來打他的有經濟的頭.
一打下去又再壓一壓.
又再深入那個創傷.
經濟刺入得更深.
疼痛更加難忍.
他們拿這個皇杖.
象徵這個位置來打頭.
其實是想嘲弄耶穌.
根本沒有權柄.
任何人都可以拿你的皇杖來打你的頭.
你還有什麼作為呢.
然後他們又吐唾沫.
即是口水在他面上.
當時來說是極盡羞辱的能事.
《約翰福音》十九章三節更加說.
冰釘用手掌打耶穌.
右一巴,左一巴.
這樣刮過去.
折磨了一段時間之後.
耶穌被釘在十字架上.
耶穌的神智體力達到崩潰的地步.
在十字架經過的人.
即是那些從城裡面出來的人.
他們在耶穌受審的時候都看過.
現在又加入戲弄耶穌的行列.
聖經說他們是辱罵他.
原文是褻瀆.
即是咒罵上帝.
上帝因為兒子的緣故被人辱罵.
有些人搖頭.
意思是表示否認和輕看.
甚至大祭司和文士.
這些有信仰的忠堅分子.
屬靈的領袖都降低人格.
加入這個褻瀆,譏笑,戲弄耶穌的陣線.
一同攻擊,侮辱主耶穌.

$^{801}$你看看這些屬靈的領袖.
這些人.
他們的眼睛可以說是瞎了.
他們的心可以說是變質,敗壞了.
如果一個人.
就算他掛有什麼屬靈的名堂身份.
而他的眼睛瞎了.
心變質的時候.
可以說他是敗壞到+01.
他的言行表現簡直和身份不相稱.
這個值得我們今天做基督徒的.
求主幫助我們警醒.
這兩個和耶穌一起釘的強盜.
起初都同樣在譏笑,嘲弄.
不過其中一個我們知道.
他信了耶穌.
這裡各種的人.
他們所辱罵,戲弄,譏笑的背後.
其實有一個共通的挑戰.
是什麼呢.
就是從十字架上走下來.
我們就會相信你了.
這些人以為看到才相信.
今天有很多人都是這樣的光景.
你說耶穌是不是真神呢.
你以為主耶穌不可以用神力使自己從十字架下來嗎.
當然可以,而且祂能夠.
但是主耶穌卻不會這樣被人激起.
因為祂要貫徹始終的.
擔當世人的罪.
要救世人.
這個重大的艱鉅任務.
難怪以前救世君的創辦人.
卜威廉就這麼說.
正因為主耶穌甘願受十字架之苦.
祂沒有從十字架走下來.
我們才相信祂.
第四項我們看看.
我怎樣扮演祂的態度.
是誤解.

$^{841}$33到36節.
那裡有說.
今天一般人的印象裡.
十字架是美麗的飾物.
但是大家知道在耶穌的時代.
當日羅馬人用十字架.
專門對付一些謀殺犯.
造反的奴隸和一些罪大惡極的罪犯.
釘上十字架.
是為他們而設.
你也聽過很多釘十字架.
是一個很長很大的釘.
就不像一些人看見的手板.
不是,而是手掌和碗之間.
大大的穿過去.
有些人記載.
根本這是很刺痛的事.
因為壓碎了手碗和手掌之間.
那裡有一個最大的神經.
你掛在十字架上的時候.
身體的重量.
由上而下墜.
就像閃電的劇痛.
從你的手掌穿過你的手臂.
一直到你的頭.
從頭到手.
痛得像要爆炸一樣.
而那個被釘的人.
是需要因為用肌肉和橫隔膜來呼吸.
他同樣受到很大的壓力.
胸部你經常處於一個吸氣的狀態.
他為了要呼氣.
兩腳,因為耶穌釘,兩隻腳疊起來.
兩腳必須要硬高,向上推動.
卻使得被釘的雙腳和背部產生另一類的劇痛.
被釘的人.
每一次一口氣的呼吸都非常吃力.
直到口喘,筋疲力竭.
不能再呼氣為止.
我們主耶穌釘十字架.

$^{881}$前後有六個小時.
飽受這樣的折磨和慘痛.
耶穌被釘的那一刻有一個驚人的現象.
從午靜,即是十二點.
到新初,即是下午三點.
全地都是黑暗.
這三個小時的時間.
太陽的光線起了極大的變化.
聖經說是遍地黑暗.
這完全是神跡.
因為是神使得全地黑暗.
黑暗也是審判的奇特的作為.
因為在舊約,約如書2章1到2節.
以塞俠書5章30節.
甚至在馬太福音24章29節.
以及路加福音21章25節.
都提及黑暗和審判有相關的事.
這次的黑暗正是父神.
為主耶穌擔當聖人的罪.
而顯明出來的審判.
新初的時候,即是三點鐘的時候.
黑暗要過去的時候.
耶穌在這個象徵審判最高峰的時候.
大聲地喊著.
「伊羅伊爾,伊羅伊爾,拉瑪撒,巴國大尼」.
整句意思就是亞蘭話.
翻譯出來的意思就是.
「我的神,我的神,為何離棄我?」.
因為主耶穌擔當了聖人的罪孽.
全世界的人類.
包括你和我的罪惡和咒詛.
都在祂身上.
以致我們的神.
我們的上帝離開了祂.
關係也變得疏離.
耶穌的呼喊是「我的神」.
而不是比較親切的「我的天父」.
原因也在這裡.
耶穌成為擔當罪孽的羔羊.
目的要救贖世人.

$^{921}$當時戲弄他們的人.
只聽到耶穌所說的前面兩個字.
抓住其他的就不理會.
斷章取義.
「伊羅伊」是亞蘭話.
馬太福音就寫成「伊利」.
是希伯來話.
當時的猶太人是懂得希伯來話.
自己又說亞蘭話.
所以他們必定明白「伊利」或「伊羅伊」的意思.
應該不會搞錯是「以利亞」.
因為希伯來和「伊利」是大大不同的.
但是那些人有偏見又心腸惡毒的人.
卻是用這句話來說.
「他叫以利亞」.
「看以利亞會不會下來救他」.
這些猶太人對耶穌的態度是誤解.
這個態度起碼有兩種表現.
第一種表現是不弄清楚.
是有意歪曲.
聽到一點不徹底.
近似立刻很偏見地下了訣解.
是有意扭曲真實的情況.
第二是以訛傳訛.
一個傳一個這樣錯解下去.
這兩種表現都是變得不求爭議.
而是屈往爭議.
同樣我們今天有些未信耶穌的人.
對基督教都是有誤解.
正是類似這種情況.
未搞清楚人說他又說.
錯誤的了解真理.
甚至有些基督徒都是.
聽到聽一點又不明白一點.
或者聽到那時候又不專心.
或者是吸收又不知道怎樣.
不求甚解.
錯誤解了聖經.
或者聽的時候已經有些東西投入了.
然後拿喜歡聽的東西.

$^{961}$甚至有時我們聽見一些東西.
將它不分辨的亂說亂講.
有些人教人喜歡教人.
但是教錯了.
又不知道教錯.
這樣才慘.
這三件事的問題就是.
我們有沒有辨識的能力.
清清楚楚.
所以我們需要好好讀聖經.
好好靈修.
好好禁戒自己和一些論斷.
這是求主幫助我們.
第五樣的態度.
我怎樣辦他.
就是真誠信靠.
37至39節.
我們見到當時耶穌為了成就救贖大功.
氣絕身亡.
聖殿當中聖所和至聖所之間的.
布障萬子重上而下.
列為兩半.
好像有一雙無形的手將他拉低.
至聖所顯露出來.
布萬裂開.
宣告猶太人的大祭司.
一年一次大祭司的血贖罪這件事終結了.
大祭司的職份從此也結束了.
因為神聖的大祭司耶穌.
已經用自己的性命和寶血贖罪了.
得盡天上至聖的居所了.
當時站在遠處的白夫長.
非常感動.
這位白夫長.
根本是能征冠戰的羅馬長官.
等級相等於今天的準衛.
很多戰役當中.
他見過不少人臨終怎樣死都見過.
但從未見過像耶穌這個人.
這麼勇敢面對死亡.

$^{1001}$我相信這位白夫長.
也知道耶穌被釘的原因.
看到耶穌在受寬當中所表現的態度.
和其他人很不同.
他又看到猶太人怎樣戲弄耶穌.
包括自己的馬仔兵丁怎樣加入.
耶穌死的時候.
馬太夫音更記載地震.
盆石裂開.
墳墓也裂開了.
這一切我相信.
必然令他有很大的感動.
觸動甚至生發敬畏的心.
根據遺傳的說法.
這個白夫長叫做朗吉路.
他在十字架底下脫置信了耶穌.
宣告說:這真是上帝的兒子.
白夫長是心悅誠服.
親自耳聞目睹.
他第一身的親身經歷.
充滿信心宣告.
白夫長真是真心信靠耶穌的態度.
值得你和我今天又再想想.
我們聽的道理很多.
我們有沒有一個親歷其境和神.
這位耶穌有親身經歷的關係.
我們有沒有謙卑下來經歷.
這位耶穌對我自己的意義.
我需要耳聞目睹.
親自進入內實.
和這位神有一個直接的關係.
讓我們學習這位白夫長一樣.
讓今晚彼拉多說.
我怎麼扮他呢.
今天我們怎樣對耶穌呢.
在耶穌時代的五類.
第一類是大家相依.
但偏行幾路.
第二,寧可要其他都不要耶穌.
第三,是戲弄,譏笑,辱罵.

$^{1041}$表明是輕項主任.
不乾不淨的表現.
第四,是錯誤解說.
不肯認真.
第五,是真誠信教.
親身經歷這個信仰的實在.
不知道這五類.
我們比較接近哪一類呢.
我們低頭誠心祈禱.
多謝主.
主,你愛我們.
藉著主耶穌的犧牲擺上.
求你在十誡的大愛.
抹滑在我們心板上.
讓我們被你的愛感動.
也愛你更多.
愛人更多.
愛教會更多.
主,你監察我們,幫助我們.
祈禱奉主耶穌基督的名.
Amen.
多謝主給我們好好提醒.
主耶穌在十字架上為我們承受.
告訴我們我們的罪孽深重.
每逢受難節的時候.
我們知道我們需要被提醒.
主對我們恩典.
我們不要做一個忘恩的人.
人很善忘.
特別是基督徒.
我們有時候為了事奉,為了工作.
我們忘記了主的恩典.
也有很多時候我們離開了教會.
離開了上帝.
我們少了祈禱.
少了去默想主的恩典的時候.
我們就忘記了.
耶穌基督為我們所拜上的有多少.
受難節給我們一個操練.
讓我們能夠記起主的恩典.

$^{1081}$耶穌在上十字架之前.
他設立了一個操練機會.
讓每一個跟隨他的人.
我們每逢吃這個餅.
喝這個杯的時候.
記起主對我們恩典.
當我們領受這個餅,這個杯的時候.
我們在身心靈上有這個經歷.
拿起這個餅.
喝這個杯的時候.
是我們肉眼能夠看到.
是我們的手,我們的身體能夠去領受的.
這個聖餐是為了我們每一個跟隨主的門徒所設.
這個是每一年猶太人.
他們當時會守的御節晚餐.
你記起《追加俠記》的時候.
那個故事.
當他們離開埃及的時候.
那一晚耶和華神跟他們說.
我要設立這個御節晚餐.
要他們很趕快的去領受.
穿著衣服,穿著鞋.
他們要很快的離開這個遺老之地.
不要再眷戀在那個地方.
埃及要離開.
不要放在心裡面.
不要回頭.
上帝以恩典去拯救他們的子民的時候.
這一隻羔羊.
他們那一晚來到領受了之後.
他們每一年重複的這樣做.
去告訴他們.
心裡面的埃及是需要殘缺的.
他們也不要回頭.
要回到以前埃及的地方.
這是一個寓意.
告訴我們每一個離開埃及的人.
就是離開罪惡的生活.
沒錯,我們得救得贖的人.
我們仍然是百分之一百的罪人.

$^{1121}$不過既然離開了埃及.
我們也是百分之一百成為恩典的子民.
不要眷戀,不要回頭.
向前看.
去上帝要我們的目的地.
去那裡進化.
預節的晚餐就是做這件事.
就是告訴以色列民聽.
告訴神所恩愛的一班子民.
包括我們在內.
上帝對我們的恩典,祂大能的拯救.
讓我們不要再在黑暗,罪惡,埃及這個地方.
所以這個預節晚餐.
當耶穌在十字架之前.
祂賦予一個全新的意義.
祂說羊羔其實就是祂自己.
祂就是完完整整,完美的每一個羊羔.
祂所流出的寶血.
在十字架上為我們捨去的生命.
就是他們喝的葡萄酒.
這是一個新的藥.
一個永恆的藥.
我們同樣在宿難節裡.
領受這個聖餐的時候.
我們想在耶穌去拯救我們.
真的出真正的埃及.
進入真正的加拿大.
真正的新天神.
不要忘記主的恩典.
所以耶穌在最後一次同門徒.
來受這個預節晚餐的時候.
馬可福音第十四章說.
他們吃飯的時候.
耶穌拿起這個無烤餅.
祂捉住之後.
祂就張開,帶給他們.
祂將這個無烤餅打碎的時候.
然後碎片就分給門徒.
祂說你們拿去領受.
這個是代表我的身體.

$^{1161}$今天我們同樣拿起主的餅.
紀念祂給我們的恩典.
我請三位弟兄姊妹上來.
(剪接中).
我們同心靈受主的餅.
照樣的馬可福音告訴我們.
當我們吃了餅之後.
耶穌又拿起他的杯.
他祝謝了之後.
他就遞給他的門徒.
他說你們去靈受.
耶穌說這是我納藥的血.
為多人流出來的.
我實在告訴你們.
我不再喝這葡萄汁.
直到我在上帝的國裡面.
喝新的那一支.
猶太人的禹節.
是有四個杯.
這四個杯是取自於.
《出亞給記大六章》.
上帝對他所恩愛的以色列人說了四句話.
第一句話是.
我會將你們脫離埃及這個地方.
遺老之地.
罪惡所捆綁的地方.
第二句話.
上帝對他的子民說.
我會救贖你.
第三句話.
我會買贖你.
從罪惡從黑暗買贖你回來.
第四句話.
你們會成為我耶和華的百姓.
所以猶太人的禹節.
有四個葡萄汁的杯.
耶穌拿起這個杯.
祂說你們拿來喝.
這是我納藥的血.
去買贖你們.

$^{1201}$去賦予你們.
祂說從此之後.
我不再喝這葡萄汁.
就是說祂沒有拿起第四個杯.
祂沒有完成.
整個猶太人.
PASSA.
禹節的晚餐.
因為.
逢領受這個杯.
我們就是.
他們所買贖的子民.
到新的天地.
耶穌才會再拿第四個杯.
給我們一個promise.
一個應許.
無論你現在經過的是什麼.
只要上帝已經救贖了我們.
耶穌說我在天國裡.
會學新的那一字.
會和我們一起完成.
這個禹節的晚餐.
我們今天拿起的杯.
等於就是第三個杯.
記起耶穌和我們說.
祂說我不再喝葡萄汁.
第四個杯.
直到新天新地.
上帝的國那裡.
我們會學新的.
耶穌的死.
不是故事的完結.
因為三天之後.
祂復活.
並且祂將來再回來.
祂會完完整整地救贖我們.
不要忘記上帝的恩典.
是一個永恆的恩典.
一個陪伴我們一生.
直到世界終結的時候.

$^{1241}$這個恩典.
用來拿起這個杯.
第三個杯.
我們跟生命受益者的杯.
上帝我們多謝你.
因為你對我們的恩典.
你去很具體的.
你都可以提醒我們.
你在十字架上.
撕裂所打碎的身體.
是去拯救我們的罪.
去洗劫我們一切的不義.
你所流盡的所有的寶血.
為了成就這個救恩.
要我們這個救恩.
要我們這個盼望.
有一個永恆天國裡面的盼望.
實在求你讓我們一班.
屬你的子民.
我們明白到.
主你的恩典是有多深.
同樣我們的生活.
我們的生命.
能夠去反映出主你的恩典.
祈禱奉主耶穌的名字.
\newpage



\section{}
\label{sec:Z07LUlzvgG0}
\textbf{4/8/18 主日講道: 基督徒成長12步--\#4 探索真理 (陳崇基牧師)}
\newline
\newline
連結: \href{https://youtube.com/watch?v=Z07LUlzvgG0}{\texttt{https://youtube.com/watch?v=Z07LUlzvgG0}} ~~~~ 語音日期: 2018-04-10
\newline
\newline
\hyperref[sec:Kg3cX9_FXZI]{\small{< < < PREV SERMON < < <}}
~
\hyperref[sec:index_chronic]{\small{[返順時目]}}
~
\hyperref[sec:index_scriptual]{\small{[返順卷目]}}
~
\hyperref[sec:a492NMYDqm0]{\small{> > > NEXT SERMON > > >}}
\newline
\newline
$^{1}$(主席發言).
大家好,歡迎大家回到城會教會.
特別歡迎新來的朋友.
我們會有一張聯絡卡.
鼓勵大家留下名字,聯絡方法.
或者有待討試項.
都歡迎你和我們聯絡.
接下來是十一奉獻的時間.
這個奉獻提醒我們要在金錢上做一個好的管家.
有請我們的歡迎團隊幫我們存這個奉獻袋.
UV的YouTube頻道.
30天的詩篇零修.
明天是第一集開始.
每個影片大概三分鐘.
如果還未訂閱我們UV頻道的姊妹.
或者弟兄都歡迎.
歡迎大家看完之後留言.
有什麼意見可以給我們.
歡迎大家讚好和分享我們的影片.
四月份開始.
姊妹旅行團的聚會.
會轉在第一個星期五晚的七點半.
男子組會在第三個星期五的七點半舉行.
請大家留意時間.
蔣麗萍姊妹的報道會.
會在四月十五日星期五五點在UV舉行.
鼓勵大家邀請身邊的朋友和親戚.
參與蔣麗萍一生最美祝福的音樂報道會.
除此之外.
還有關心賢小姐的報道會.
會在五月十六日星期三七點半舉行.
鼓勵大家邀請朋友和親戚參與聚會.
很感恩UV越來越多BB.
為了不影響身邊人的崇拜.
歡迎父母將六個月大的BB放在紓壓所.
大家和父母可以好好享受崇拜的時間.
專心崇拜.
今天會有一個家事大會.
會在崇拜後十一點十分開始.
大約一小時.

$^{41}$到時會有分組討論成會教會的財務.
教會事工和選舉方法等等.
大家完了聚會都歡迎參加家事大會.
時間交給麥sir.
麥sir:早安.
我們有家事大會請各位留步.
雖然我們只是短短一個小時的時間.
但有很多重要的事需要大家合作建立教會.
上星期復活節主日我們說得很清楚.
復活節不只是讓我們拿護照等上天堂.
復活節告訴我們我們有工作要做.
因為上帝在耶穌基督死而復活的身上.
已經開始了救贖的工作.
迪斯力袋的信徒教會也在做這件事.
所以我們有一個任務要做.
所以我們有一個家事大會提醒大家.
我們有很多不同的角色.
有很多不同的事奉.
有很多不同的工作.
我們需要大家一起去建立上帝的教會.
和去興旺福音.
所以希望大家能夠留步.
我們在今年開始的時候.
我們每個月一次.
我們說到一個系列.
是十二個成長的方面.
每一個基督徒的教會.
個人信徒的生活.
都需要有起碼這十二步.
我們說到從福音開始.
我們來跟隨耶穌.
以耶穌為我們的主.
然後第二月的時候說到.
我們主日的崇拜.
為什麼我們需要來敬拜上帝.
來跟祂建立關係.
以祂為我們的君王.
這是一個最基本的成長的一方面.
第三方面就是說到.
我們要認識信仰的基本.

$^{81}$一些基礎我們需要.
然後今天我們說到.
我們要繼續去探索上帝的真理.
如何去整合真理.
以至我們的生活每一方面.
我們的職業.
我們的工作等等.
如何在上帝的話語的基礎上.
我們能夠活出一個屬靈的人.
所以今天我們看到第四步.
福音書裡面有很多篇幅.
大家要知道.
無論是《馬太》《馬可羅》加《約翰》.
有很多時候都告訴我們.
不是耶穌基督的死和復活這麼簡單.
沒錯.
福音書最後四福音最後的時候.
是說到死而復活.
這是最高峰的時候.
耶穌來為我們的罪死.
然後救贖我們.
以我們能夠在十字架底下.
成為上帝的子民.
這個當然是最重要的.
但是有很多篇幅.
不只是說耶穌的死和復活.
有很多是說到耶穌的生活.
耶穌和門徒的交往.
耶穌傳道.
耶穌如何接觸其他人.
為什麼呢.
因為主要是要告訴我們.
我們今天已經接受了耶穌基督的人.
我們應該如何去活出這個生命.
而福音書有很多篇幅.
就是告訴我們.
介紹這位耶穌.
介紹這位耶穌帶著的這群門徒.
如何能夠在不同的地方.
不同的處境.

$^{121}$他們來到.
有工作要做.
在天國裡面如何去延伸上帝自己的工作.
今天我們來看《馬可福音》.
《馬可福音》有一段聖經.
由第一章至第三章.
就是介紹耶穌的時候.
他是一個風塵撲撲的僕人.
特別說到他接觸了很多其他人.
然後吸引了很多人.
去聽道.
去到耶穌的面前.
然後第四章.
《馬可福音》就說.
耶穌有到海邊教訓人.
有很多人來到他的面前去聚集.
他只能坐在船上.
船在海裡面.
眾人靠近海.
然後站在岸邊.
耶穌就在坐.
坐或者站.
在船上坐.
他就用比喻來教訓這些人.
教訓他們許多道理.
在教訓之間對他們說.
就是說耶穌用很多方法.
來跟這些湧在他面前的人.
跟隨他們的門徒.
這些群眾告訴他們.
天國是一回事.
這裡說到海邊.
說到加利利海.
或者加利利湖.
因為這是一個淡水湖.
不過我們叫它海.
因為希伯來文沒什麼分淡水和咸水.
但是這是加利利湖.
加利利湖的湖邊.
加伯隆等等我們聽過的地方.

$^{161}$就是耶穌傳道開始公開事奉的基地.
這個海邊.
是耶穌去到呼召門徒的時間.
的地方.
他說這裡很多人容許他在他面前.
他只能在船上坐.
太多人了.
他在岸邊的時候.
不如租一艘船或者划一艘船出去一點.
然後向著岸邊去港島.
因為地方很不足夠.
大家知道加利利湖.
旁邊是很傾斜的.
到今天還有一個地方叫做Sawers Cove.
一個叫做殺種者的海灣.
就是半圓形的地方.
為什麼呢?.
因為湖邊四周.
地勢非常傾斜.
所以他划船出去的時候.
向著岸邊去港島.
有很多群眾坐在半山之上.
其實就是一個天然的劇場.
所以有些人說這個就是那個地方.
Sawers Cove.
無論如何.
他就用當時的環境.
耶穌就用比喻來教導他們.
當時猶太人的拉比.
是特別多用比喻的.
這個絕對不是耶穌發明的教學方法.
parable 比喻.
其實就是用地上的事.
來代表或者表達天上的事.
地上的事我們都很熟悉.
我們接觸過.
我們碰過的東西.
他就用這些事物.
來告訴我們.
天國上帝神的國的道理是怎樣的.

$^{201}$他用很多的比喻來教導他們.
對他們說什麼呢?.
其中一個比喻馬克方就說.
他說你們聽.
有一個殺種的出去殺種.
殺的時候.
有一個在路旁的.
飛了你吃盡了.
也有落在淺石頭地上的.
土技不深.
發芽最快.
白天出來一曬.
因為沒有根就枯乾了.
也有落在荊棘裡面的.
荊棘長起來的時候.
就把它擠住了.
結果就不能結實.
也有一些落在好的土上面.
就發生長大.
結果有三十倍.
六十倍.
一百倍.
耶穌說有耳可聽的.
就應當聽.
這裡說到.
有一個出去.
好像很笨拙的翻譯.
其實直譯出來就是.
有一個殺種的出去殺種.
特別說出去殺種的字.
就是耶穌在頭三章.
說他出去不同的地方.
來公開傳道和事奉.
所以我很清楚殺種的那個農夫.
一定是主自己.
一定是耶穌自己.
種子就是天國福音的真理.
他將真理的道殺出去的時候.
就有不同的接收的方法.
或者不同的反應.

$^{241}$他說了幾種.
大家看看有多少種.
在巴勒斯坦的土地.
山多平地少.
我不知道中間有沒有試過耕田.
有沒有.
有吧.
或者這裡可能也算是耕田.
你在口一那裡種東西.
蔥太貴自己種蔥.
或者也可以有一點點的連繫.
種東西.
有多少種呢.
巴勒斯坦地是很特別的地方.
山多平地少.
正如我所說.
所以沒有一塊很大的喉一.
沒有一塊很大的麥田.
讓你去.
好像美國的conville那樣.
有個大車在那裡駛過就搞定了.
山多平地少的地方.
是怎樣去耕種呢.
他就要靠梯田.
梯田的意思就是說在半山斜坡.
然後你就很辛苦地作.
然後就作一級一級出來.
然後那些作出來的大石.
你就放在梯田的邊邊那裡.
然後雨或者侵蝕等等.
射了一些泥土下來的時候.
你就可以裝那些泥土.
承載它一些肥沃的泥土.
你就可以裝在那些不同的梯級那裡.
那些梯級就成為了一點點的平地.
讓你可以在那裡耕種.
耶穌就用這個比喻.
特別是加里里的地方.
耶路撒冷等等的地方.
都是有很多這樣的梯田.

$^{281}$一級一級一級的.
到今時今日都是.
種那些果樹.
小的麥田.
有很多都是在一些terrace裡面去栽種.
他說有一個殺種的出去殺種.
就殺在什麼地方呢.
他說殺在路旁.
所以是一個很地道的比喻.
當時的農民.
當時的平民百姓一聽就明白.
殺種.
殺出去.
當然有一些就落在路旁.
那些路旁是讓你在梯田那裡.
上上落落.
然後去那裡挑灌.
掘水等等.
有些種子落在路旁.
這個比喻之後有解釋.
耶穌在暗中和門徒來解釋的時候.
他說殺在路旁是什麼意思呢.
那些種子落在路旁的時候.
就有飛鳥來吃那些種子.
那些種子就代表撒旦.
魔鬼撒旦.
就是等人聽了道之後.
他就在那裡等著.
來將這些道種.
在人的心裡面.
如果你的道是落在路旁的話.
他就立刻會將它拿走.
也有一些殺在淺石地上.
你看到梯田有很多石頭.
可能是大塊的.
不能完整地做成一個梯田的時候.
就有些淺的石地.
他說發苗是非常快的.
因為土非常淺.
這個比喻其實在說什麼呢.

$^{321}$耶穌解釋的時候.
就是有些人聽了福音聽了真理.
他立刻很有反應.
很有感覺.
很溫馨.
很有感觸等等.
OK.
甚至可能你哭三天三夜.
可以是這樣.
但那些是發苗非常快.
但是他說日頭一曬的時候.
因為這些苗是沒有心的根.
所以很快就會枯乾.
這些全部都是很地道的言語.
日頭一曬.
那些跟我去過以色列的人.
就知道日頭一曬是什麼意思.
曬到你真的快燒焦了.
巴勒斯坦地.
那些淺土上面的苗.
因為沒有根的時候.
就不能夠吸到水分.
就不能夠去抗衡太陽的熱力.
或者水分蒸發得太快.
就會很多時候枯乾.
那些快的苗.
很多時候就是這樣的收場.
這是第二種.
第三種.
他說是落在荊棘裡面的.
有很多地方是有野草長出來的.
摸起來很像手的野草.
耶穌解釋是什麼呢.
那些荊棘就是世上的屍類.
錢財的迷幻.
別樣的屍肉湧進來的時候.
就好像那些荊棘將那些發苗的植物.
來去束住它們.
或者吸收它們的養分.
讓它們不會可以結果子.

$^{361}$這些所謂屍肉.
世上的屍類或者錢財的迷幻.
很多時候都是一些好的東西.
可以是我們的工作.
可以是我們的家庭.
可以是任何的好東西.
不一定是錯的.
不一定是罪的.
只不過很多時候就是讓那些東西.
來去屍類誘惑的時候.
就把我們塞住了.
吸收了我們的養分和精力.
我們就不能夠去結果.
這是第三種.
第四種就是落在好土上面.
就是有人聽了道.
他去領受.
然後結出好的果子.
有四樣的土壤.
那個比喻就是告訴我們.
到底我們是哪樣.
我們是那些路旁的人.
還是淺土上面.
淺石地上面的一種心境.
還是我們很多時候.
我們的心境.
我們聽道我們聽福音的時候.
就是一種在荊棘裡面的土壤.
聽了道.
但是有太多東西搞了.
生活裡面太多東西擠擁我.
於是上帝的道.
天國的道理.
就不能夠在我們的生命裡面.
來去化瓦結果.
還是我們是一種好的土呢.
馬可夫林描述耶穌.
去作天國的講道的時候.
他首先用了三章.
去講到主去到全道去殺種的時候.

$^{401}$人有不同的反應.
如果我們查一查.
馬可夫林一至三章的時候.
他也不解釋那麼多.
耶穌一開剝.
他就出來去講道.
很多人有不同的反應.
有一些群眾.
他就來到耶穌的面前.
因為他可以行神蹟.
可以治病.
眾人就擁擠他.
擠過去.
好像舉個例子.
這裡沒有醫院.
沒有醫生.
個個有病的人.
就擁去醫生那裡排隊.
排八個小時都沒問題.
排五個小時都沒問題.
十個小時都沒問題.
治得好就行了.
這些人來到耶穌面前.
都是同一個心態.
他們希望能夠得到治病.
得到神蹟.
這些是沒錯的.
不過.
很多這樣的群眾.
這些擠擁的人.
就是想在主上得到好處.
或者他聽到上帝的道.
聽到耶穌講道的時候.
他教學的時候.
他覺得有道理.
很有感覺.
不過他自己覺得很驚訝.
這個人形式是不同的.
這些人.
就在這個殺種比喻裡面出現.

$^{441}$這些群眾.
就是想在主上得到好處.
And that's all.
我只是在上帝面前.
在上帝的身上.
我只是想拿一些對我有用的東西.
你如果使我升官發財的話.
我為什麼不來聽呢.
這些就在殺種比喻裡面出現.
還有另外一些群眾.
是什麼群眾呢.
就是法利賽人.
法利賽人就是當時第一世紀猶太教的權貴.
權貴人士.
這些是宗教的領袖.
他們來到耶穌面前.
他們覺得他很特別.
不過.
他們一聽他講話.
或者聽到他做神蹟奇事的時候.
他們就說這個人是被鬼附的.
這個人是瘋的.
還有什麼呢.
因為這個人做的事情.
跟我們的規矩是不同的.
因為他特別是在安息日裡面.
他自病行不公然去挑戰當時的權貴的紅線.
屬於耶路撒冷的宗教領袖.
法利賽人.
這些人來到耶穌面前.
他們就去反對.
甚至馬夫人說.
他們想除滅耶穌.
因為有耶穌就沒有我們.
所以我們就要將他致死.
設不同的陷害計謀.
想要他致命.
這些就是反對的一群人.
路旁的殺種.
下去路旁的.

$^{481}$就是這種人.
還有另外一種.
就是門徒.
被呼召跟隨耶穌的時候.
他們就算是一種典範的.
一種天國的生活.
他們聽了.
接受領袖的.
所以有人說這個殺種的比喻.
其實不應該叫殺種的比喻.
應該叫他做.
禮土的比喻.
因為主要說的不是說殺種.
殺種的當時平民百姓農夫就是這樣殺種的.
一定是這樣的.
你說是不是浪費了.
沒辦法的.
你在巴勒斯坦地.
你去種田這樣做農夫的時候.
你一定會這樣.
地勢實際環境.
是逼使你這樣.
所以最主要不是說殺種的農夫.
主要是說收種子的禮土土樣.
是有哪幾種.
這裡說了有四種.
殺種的比喻開始說什麼呢.
他就是告訴我們.
他想我們成為一種怎樣的土樣.
他說你們聽.
有一個殺種的出去殺種.
結束的時候也是這樣說.
耶穌說有耳可聽的.
就應當聽.
to listen.
去聆聽專心的.
來聆聽上帝天國的道理.
聖經裡面的話語.
但你不要覺得馬可福音.
說這麼簡單的道理.

$^{521}$這麼簡單的故事.
馬可不是這樣的.
馬可是很精細的.
馬可很多時候.
他去記載耶穌一些事蹟的時候.
他用主題去編排.
有很多所謂三文字出現.
三文字意思就是說.
說完一件事 A.
然後他就說 B.
然後他就說回 A.
有點像我們唱歌一樣.
唱完第一字然後唱副歌.
副歌又唱第一字.
最重要在哪裡.
副歌就是中間.
三文字.
馬可特別多這樣的手法.
來記載耶穌的事蹟.
是一種文學的手法.
所以四章第一節說.
有一個殺種的出去殺種.
就是耶穌說的殺種的比喻.
或者是圍圖的比喻.
然後最後的時候他就去解釋.
門徒私底下來到主人面前.
他就去解釋這個殺種的比喻.
到底是說什麼.
地上的事物.
怎樣連結天上天國的道理.
夾住中間的四章十二節.
那個就是最重要.
最中間最中心的東西.
所以真的要所謂要你命的.
真的要你專注的是什麼.
不是那個很有趣的故事.
而是中間那一段.
就是比喻教導的解釋.
中間他說什麼.
他說沒有人來的時候.

$^{561}$群眾散開了.
拿了好處走了.
跟隨耶穌的人.
和十二門徒.
那些真的緊緊跟隨耶穌的人.
他就問耶穌.
這個比喻的意思是什麼.
耶穌對他說.
他說天國的奧秘.
只叫你們知道.
就是說你們這群十二門徒.
和緊緊跟隨耶穌的人.
如果對外人說的話.
我就說比喻.
為什麼會這樣.
為什麼只有門徒.
真的跟隨耶穌那群人.
才能夠明白這個奧秘.
因為眾人散開之後.
還在耶穌身邊的是什麼人.
那些就不只是拿好處.
那些真的明白.
去耶穌面前跟隨的時候.
十二門徒跟隨他的人.
那些是什麼人.
馬可就說希望是我們.
我們是跟隨耶穌和十二門徒的那群人.
呼召了跟隨他的時候.
特別說到十二門徒的時候.
第三章之前那章.
耶穌說呼召十二門徒.
就是呼召十二門徒.
常與他同在.
事上的跟著他.
事上的和他建立關係的那群人.
這一群人和耶穌真正建立關係的時候.
他們就明白天國比喻的真相.
他們就不是那些淺土.
不是那些荊棘裡面的.
他們肯去聽.

$^{601}$肯去思想.
肯去考問.
到底耶穌在說什麼呢.
這一群人就是好土的人.
你是不是這樣.
去聽.
在希伯來人的文化裡面聽到.
不只是聽.
是研究.
是傳播.
還要教導其他人.
如果你看看.
新面基特律章.
以色列人最出名每天都會背誦的經文.
每天都會背誦的祈禱.
他說以色列你聽.
耶和華是獨一的主.
然後繼續說下去的時候.
你要思想去探論.
思想去研究.
思想去探索.
然後在你家裡.
你要釘在門上.
教導你的兒女.
新教,現教的.
讓他們明白上帝自己的華語.
聽上帝的道.
所以思想有研究.
有整合的意思.
和我心裡面有什麼關係呢.
和我的行業有什麼關係呢.
和我的行業,我的職業.
和我的身份有什麼關係呢.
這些全部都是聽的意思.
聽不是只是塞你的時間.
而是我們真的用心去領受之後.
我們肯去思考.
肯去研究.
肯去辯論.
肯去問與答.

$^{641}$這個其實就是我們一直以來.
我們所謂教會裡面.
有一個很重要的.
很需要的傳統.
就是叫做主一行.
有幾百年的歷史.
教會舉行主一行.
不代表幾百年之前沒有主一行.
不過叫做Sunday School.
只是幾百年.
十八世紀左右.
開始有這樣的運動.
就是說我們需要去聽上帝的華語.
去研究,去查考.
然後去整合它出來.
我們的探索課程都是一樣.
我們的探索課程都是.
一個小時的時間.
我們就是希望.
我們大家去整合.
去研究一個社群的.
去建立和主的關係.
不是那些淺土的走去聽完.
我們就罷了.
而是我們真是Sing in.
吸收它.
去整合它.
我們年頭的時候提過大家.
我們有一個222的行動.
今時四月了.
再提一提你.
我們兩個小時建立.
在團契,小組做什麼呢?.
不是塞我們的時間表.
大家都已經很忙了.
很多荊棘.
大把荊棘在我們身邊.
而是建立我們屬靈的友誼.
是一個Community.
一起來跟隨上帝.

$^{681}$兩個小時就是我們主日敬拜的生活.
一個小時我們來敬拜.
我們會聽道當然.
我們會聽分享.
甚至有時候有話劇,有歌劇.
上個星期不是這樣嗎?.
但另一個小時.
就試探索課程的時間.
第三個我們想提醒你.
剛才報告已經有提及.
我們30天的靈周.
我們又再開播.
不是塞你的時間.
而是操練我們個人的屬靈健康.
222.
我們知道現代人的注意力越來越短.
對嗎?.
大家有聽過什麼叫TED Talk.
TD.
他們的研究說.
今天成人的注意力是18分鐘.
18分鐘.
所以每一個TED Talk上台講話.
都只是要求他18分鐘.
我們UV.
有時候我也有這樣的想法.
講18分鐘就不可以了.
當然我會比18分鐘長.
但最主要的重點是.
我們知道我們的注意力越來越短.
那怎麼辦呢?.
只要有兩個出路.
第一.
叫講的人講短一點.
對嗎?.
第二.
就去鍛鍊你的注意力.
鍛鍊你自己怎樣去聽.
怎樣去聲音.
鍛鍊你自己不要那麼分心.

$^{721}$因為有太多根基在我們身上.
然後檢查電話.
電話又震又響.
太多東西令我們有分歧.
我們要鍛鍊我們怎樣去專注.
在一件事上.
特別是做父母的.
就更加需要新教.
我可以坐定.
打開聖經.
真的鍛鍊我的注意力.
聽到.
希臘人的文化.
聽除了研究探索之外.
就是走出來.
去應用出來.
猶太人的文化.
由舊約到新約.
聽與行不是分割的.
不是我們現代人那樣.
聽純粹是智慧的練習.
一個理性的作業.
我們做完聽完就罷了.
聽完再走出來.
整件事就叫做聽.
猶太人有一個拉比的傳統.
老師教門生做什麼.
甚至每天早上之前.
有拉比甚至這樣教.
早上前.
早上前去上帝面前祈禱之前.
你要做什麼.
你出去奉獻.
你走到街上.
看到有需要的東西.
看到上帝國外面有什麼需要.
有什麼工作.
是你需要去奉獻.
需要施捨的.
是施羽的.

$^{761}$你就先做那件事.
為什麼猶太人拉比會教這件事.
你記得我們講過主禱文.
主禱文六章第九節.
我們經常說要看上下文.
上文講什麼.
主禱文上文.
我們只會背主禱文.
主禱文之前講什麼.
就是拉比的傳統.
就是講施捨.
耶穌說你們施捨.
你們奉獻的時候.
就不像宗教領袖權貴那樣.
要別人知道.
你暗中去奉獻.
掌上就在天上.
他講了一段是講施捨.
先講奉獻講施捨.
然後才講主禱文.
猶太人拉比為什麼會這樣教.
他說.
so that.
以致什麼.
以致你祈禱的時候.
你不會只是為你自己自私的祈禱.
因為你明白世界的需要.
錢都拿出來了.
奉獻了.
先奉獻了.
我就明白.
上帝要我為世界的需要.
我需要奉獻什麼.
需要拿什麼出來.
我才為自己祈禱.
以致我的祈禱不是為了自己安心.
為了自己平安.
為了自己發大.
為了自己好處才禱告.
你的祈禱就會轉化過來.

$^{801}$就是為了上帝天國去禱告.
願人尊你的名為.
成員你的國降臨.
願你指引在地上.
願你降臨在天上.
It's not me, me, me all the time.
先做那件事.
這個跟我們傳統的教導很不同.
第一件事當然是祈禱.
但怎知猶太人的拉比就將他反轉.
他說你千萬不要這樣祈禱.
你先去行道.
然後才在上帝面前.
你的生命就會轉化.
反過來.
很多時候只是我自己.
越祈就越衰.
對不對.
你試試.
耶穌說了一件事.
聽完我們去行.
或者行與聽.
沒有分開一和二.
沒有說我做完這件事.
聽多一點.
吃多一點.
這樣我才走出來.
No.
行與聽是不分開的.
我想鼓勵你.
Two, two and two.
我們探索真理.
去明白上帝天國的生命.
是需要下功夫的.
我們詩篇三十天又來了.
明天.
Again.
我們盡量將它減到三分鐘.
零收三秒鐘好嗎.
沒理由這麼低.

$^{841}$如果上帝只是三秒鐘.
我想你不會相信.
你不要信罷了.
你不要告訴別人你相信.
減了三分鐘.
並且我們不是跟你說廢話.
有一些很重要的字眼.
有一些很重要的句子.
在上帝的道裡面.
你帶著一個好的心態.
一個好的土壤的心態.
不是當交差.
所以我們不是搞什麼.
三十天之後誰看完三十天.
我們送一些東西給他.
我們不搞這些的.
搞什麼.
送你又不是送很多.
I don't know what's that for.
Certificate都不做了.
我不知道為什麼華人教會.
很喜歡搞什麼Certificate course.
上完之後就給你一張紙.
封塵.
What's that for.
因為我們真正來到上帝面前.
來尋求祂探索祂的時候.
我們不是去交差.
不是做功課.
而是和上帝同在.
正如耶穌所說.
十二門徒不是跟隨他的人.
私底下就來問耶穌.
天國的奧秘.
那個奧秘就屬於你.
那些只是來拿好處的群眾.
他們就沒有份.
上與上帝的同眾.
來聽他自己的話.
甘心樂意.

$^{881}$明天開始.
我們繼續探索上帝的道.
我們預備聖餐之前.
我們同心低頭祈禱.
上主我們今天來領受你的恩典.
求你讓我們成為好的土壤.
讓你撒出的種子.
真理天國的道.
能夠真正進入我們的生命當中.
主讓你自己成為我們人生的主宰.
讓我們時刻與你同在.
讓我們Sing in.
接受你的道的種子的時候.
發芽生長三十倍六十倍一百倍.
讓你的道去引領我們的生活.
結合我們的不同的角色.
社會裡面不同的職責.
讓你的道來指引我們.
使我們真正能夠與耶穌.
是上的同在.
我們要領受這個聖餐.
為了我們每一個.
我們能夠得蒙救贖.
你死而復活的恩典.
臨到我們人生的.
我們獻上感恩讚美.
讓這個餅這個杯.
讓我們領受的時候.
想起你對我們的恩典.
並且為我們生活.
與這個恩典上情.
祈禱奉主耶穌的名字.
這個聖餐是為每一個.
基督徒的門徒而設的.
常與耶穌同在的人.
你今天願意去跟隨他的話.
無論你十年之前做的決定.
或者今天做的決定.
起碼拿這個餅這個杯.
我們回到座位那裡.

$^{921}$我們一起等候.
其他人來對你領受.
我們三張聖餐桌.
如果你準備好的話請你出來.
(聖餐禮儀).
(主持人:你好,你們好).
(主持人:你們好,我是Anna).
(主持人:我想介紹給你們聽).
(主持人:這是世界聲明會的活動).
(主持人:他們一直在做一個活動,叫做6K for Water).
(主持人:是怎樣的呢?是甚麼來的呢?).
(主持人:我們可以看下一張嗎?).
(主持人:在一些未發展好的地方).
(主持人:有些小朋友和大人沒有乾淨的水).
(主持人:他們所喝的水有機會令他們感染或有病).
(主持人:他們平均每天走的路去拿水的時候).
(主持人:例如要走6K,即是大約6公里,大約3.3公里).
(主持人:很多人會拿著一個很大的水桶走).
(主持人:但這裡不是平均的,很多時候都是不平均的).
(主持人:或者有很多石頭).
(主持人:我們現在做這個活動的時候).
(主持人:每次一個人報名,大約50元).
(主持人:可以幫到一個小朋友去幫他們製造機器).
(主持人:令他們有乾淨的水).
(主持人:我們可以看看影片嗎?).
(主持人:可以看之前的影片嗎?).
(主持人:看不到嗎?).
(主持人:不要緊,你們回去自己看).
(主持人:你們可以自己Google 6K for water).
(主持人:我們先請Cecilia上來).
(主持人:這位是我的侄女).
(主持人:她是我們的人版).
(主持人:Cecilia,我先問你).
(主持人:請你到下一個頁面).
(主持人:我問Cecilia的問題).
(主持人:你為什麼想用6K for water?).
(Cecilia:我姨姨問我).
(Cecilia:想不想用6K for water).
(Cecilia:我回答是).
(Cecilia:我第一次用6K for water).

$^{961}$(Cecilia:我已經六歲了).
(主持人:她的侄女是主人).
(主持人:和她一起跑的,是主人).
(主持人:這個活動也是主人介紹我們參與的).
(主持人:她第一次跑,你看我的衣服).
(主持人:是2014年,這件就是2014年的衣服).
(主持人:Cecilia,你已經做了多少次?).
(Cecilia:我已經三次用6K for water).
(主持人:她第一次跑是六歲).
(主持人:之後她繼續跟著主人跑了幾次).
(主持人:這是第二次跑嗎?).
(Cecilia:第二次跑).
(主持人:你看她第一章剛開始的時候).
(主持人:是很有力量的).
(主持人:然後Cecilia發生什麼事?).
(Cecilia:我累了).
(主持人:但她繼續跑,完成了).
(主持人:最後她拿到什麼?).
(Cecilia:一塊金牌).
(主持人:她拿到一塊金牌).
(主持人:這次也有金牌和T恤).
(主持人:我等下給你看看).
(主持人:你從這次經歷中得到什麼?).
(Cecilia:我認為如果我把話語傳遍).
(Cecilia:我可以讓人加入6K for water).
(Cecilia:並為水提供資金).
(主持人:其實第一次Cecilia加入的時候).
(主持人:我突然說:Cecilia,你也加入吧).
(主持人:我錄了一個影片給她).
(主持人:她現在六歲).
(主持人:我跟她說:你先看這個YouTube).
(主持人:我們可以給人看).
(主持人:然後叫人支持你).
(主持人:我們在那晚錄影).
(主持人:她就看了我們影片的連結).
(主持人:看完後她說:我要加入了).
(主持人:但那時候小朋友只有50元).
(主持人:我跟她說:那怎麼辦?).
(主持人:我錄了一個影片).
(主持人:叫她去找人支持她).

$^{1001}$(主持人:她就錄了一個影片).
(主持人:我們現在也找不到影片).
(主持人:但應該可以找到).
(主持人:我跟她說:你先錄了一個影片).
(主持人:我們在那晚錄影後).
(主持人:馬上找到姑婆).
(主持人:是姑婆嗎?).
(主持人:是).
(主持人:兩個人支持了她).
(主持人:她就可以加入了6K for Water).
(主持人:請我們看下一張).
(主持人:其實我看到她今年).
(主持人:她的小朋友已經變成25元).
(主持人:其實兩個小朋友).
(主持人:可以幫到一個大人).
(主持人:不是,一個小朋友去拿清水).
(主持人:如果大人報名,就是50元).
(主持人:我們現在是).
(主持人:開了一個團隊叫Team UV).
(主持人:如果你報名的時候).
(主持人:你就可以說我要加入這個團隊).
(主持人:跟UV大人小朋友一起跑).
(主持人:其實不一定要跑,走都可以).
(主持人:所以我決定拿著小朋友推他走).
(主持人:這個就是她今年的T恤).
(主持人:這是一個樣本給大家看).
(主持人:所有T恤上都有一個BIP).
(主持人:這個BIP就是).
(主持人:她的小朋友的名字).
(主持人:這個BIP就是).
(主持人:有張相片告訴你).
(主持人:你是為了這個小朋友而跑的).
(主持人:每張都不一樣).
(主持人:每張都有自己的密碼).
(主持人:還有,跑完的時候).
(主持人:就會有一個獎牌).
(主持人:會很有成功感的).
(主持人:如果大家想跑的話).
(主持人:你可以自己在網上登記).
(主持人:如果需要幫忙,你可以找我).

$^{1041}$(主持人:UV的目標就是要提高一千元).
(主持人:其實我覺得應該很容易達到).
(主持人:我們現在已經有四百元).
(主持人:不是,四百二十元).
(主持人:我們現在已經有四百二十元).
(主持人:這個就是條碼).
(主持人:你們可以去按Join Team).
(主持人:就可以回答她的問題).
(主持人:如果有什麼問題).
(主持人:你們就可以來找我,謝謝).
(主持人:好,謝謝Anna).
(主持人:謝謝Cecile).
(主持人:6K這麼小可以跑,厲害).
好,我們今天就結束了.
\newpage



\section{}
\label{sec:a492NMYDqm0}
\textbf{4/15/18 主日講道: 兩種禱告與生命(吳君才牧師)}
\newline
\newline
連結: \href{https://youtube.com/watch?v=a492NMYDqm0}{\texttt{https://youtube.com/watch?v=a492NMYDqm0}} ~~~~ 語音日期: 2018-04-17
\newline
\newline
\hyperref[sec:Z07LUlzvgG0]{\small{< < < PREV SERMON < < <}}
~
\hyperref[sec:index_chronic]{\small{[返順時目]}}
~
\hyperref[sec:index_scriptual]{\small{[返順卷目]}}
~
\hyperref[sec:OhwVkVNV6k4]{\small{> > > NEXT SERMON > > >}}
\newline
\newline
$^{1}$弟兄姊妹平安.
有幾次我邀請一些肢體帶領大家祈禱.
但對方拒絕.
理由是自己祈禱不行不好.
確實有些基督徒信主若干年.
不過很遺憾的是他們怕祈禱.
尤其是在很多人面前或小組聚會中領禱.
不知你有沒有類似的情況.
其實不用怕.
禱告是靈命的操練.
多禱告多操練你的靈命就會慢慢提升.
進一步說.
聽人的祈禱.
從你的祈禱.
其實可以反映到人或自己的屬靈狀況.
求主幫助我們提升自己禱告和靈命.
我們先簡單的有個祈禱.
天父將以下時間交託提醒我們.
我們的禱告和靈命究竟應該怎樣.
討你的喜悅.
願主你幫助我們.
祈禱奉主耶穌基督的名禱.
阿們.
今天跟大家分享的經文就是.
路加福音十八章九到十四節.
看看經文.
稍後解的時候會再抽.
不過這段時間我們很快地拿一拿.
我不讀出來了.
再另一張.
我們看到經文的比喻.
其實是法利賽人和瑞利的祈禱.
這兩個人的祈禱正正反映了他們兩個的靈命.
或者基督徒的靈命狀況.
我們藉著這兩個人的祈禱來看看自己的靈命和祈禱的生活.
我們先看法利賽人的祈禱.
他的祈禱和靈命的第一樣很特色的.
就是自我中心.
耶穌向那些仗著自己是二人藐視別人的人.
他有三樣很特別的表現.

$^{41}$第一樣提了.
就是自我中心.
第九節.
你看到.
耶穌根本是針對那些仗著自己是二人的人.
用這個比喻.
現代中文翻譯為.
我做的事都對,我行.
其實有時這些也是自我中心的典型.
第十一節.
你看看第十一節.
法利賽人是站著自言自語的禱告.
劈頭就說.
聽起來你看.
多屬靈.
多懂得感恩.
但其實他是自言自語.
他其實是對自己禱告.
向自我炫耀.
還有他實際上並不是真的在向神傾心吐意的祈禱.
表面上他是聖敬神.
口說是上帝.
神啊.
但骨子裡其實心裡想到的都是自己.
樣樣都是標榜自己.
你看看這個法利賽人一直沒有辦法擺脫自我.
他根本沒有提到神.
你看看第十一節到十二節.
你數一數有多少個我字.
我,我,我.
充分流露去標榜自我.
滿有自我中心的特色.
有人說今天是一個自我中心的世代.
有很多微信的人都會說.
你也經常聽人說.
人不為己.
下一句天誅地滅.
所以他們在生活上是很怕吃虧的.
很多時候做什麼都爭先恐後.
千萬不要執輸.

$^{81}$唯恐不及.
不知道你有沒有見過一些視像.
說一些中國大陸的旅行團.
好像去到馬來西亞還是哪裡.
吃自助餐.
吃自助餐夾東西.
用夾子和匙羹.
他又不是.
是用碟子.
看到碟子.
看到蝦.
啪啪啪啪啪.
接著都不夠.
整盤拿走.
疊到高高.
在其他地方.
失禮我們中國人.
那種自私,自我.
你會發現我們其實.
基督徒.
我們何嘗不是自私自我.
當你的生活.
或者你的禱告.
只顧自己.
只集中自己.
或者是專注自己的家人.
專注自己的好處.
而漠視了其他的人.
包括外人.
或者社會.
或者教會.
其他人的時候.
會不會都有一些自我的成份呢.
你記得三幾個月前.
在佛羅里達州.
發生一個中學槍擊事件.
有17個師生被槍殺.
在這麼危難逃難當中.
有一個15歲的華裔少年人.
他是很鎮定的.

$^{121}$和很有條理的為同學.
開門疏散一批又一批.
驚恐的人群.
他忘我為人到一個地步.
只顧照顧同學逃生.
結果被槍手瞄準.
砰砰.
結果他死了.
原來他是西點軍校.
預備錄取的學生.
在他喪禮的那天.
西點軍校正式宣告.
他是未來的學生.
並為他不畏自我.
能夠捨己救人.
在棺木上蓋上了國旗.
向他行軍禮.
我更記得好幾年前.
是10月2日.
發生在賓州.
亞米希社區的校園.
槍擊事件.
造成了5死5傷.
在生還者口中.
透露當時一些很感人的情景.
彰顯了一些不自我中心的人性光輝.
和那種正直的道德勇氣.
更加反映了原來從小到大.
基督教教育對兒童青少年.
原來可以產生很正面的作用.
當槍手試圖殘殺10個女生之前.
其中比較年齡大的那一個女生.
是13歲的瑪麗安.
她挺身而出.
她捨己救人.
不顧自己的命.
和槍手談判.
她說你殺我一個就好了.
就放了其他人.
小小的年紀不會自保求存.

$^{161}$反而要求槍手.
你放過他們.
先朝我開槍.
結果那位空手真的.
將瑪麗安打死.
仍然不罷休還想.
瑪麗安11歲的妹妹巴比.
同樣跟著起來大言無慚地說.
你對著我再開槍.
下一個就對著我開槍.
她挺身而出捨己繼續要救人.
結果那槍手真的開槍兩槍.
一槍對她肩膊.
一個對她的電骨.
幸好大能不死.
只是受重傷.
這兩個小姐妹.
未殺之前曾經跟槍手說.
為什麼你這樣做.
真是想不到13,11歲.
竟然這麼大義無畏.
冷言正直的.
不自私.
不自我中心的表現.
其實正正就是家庭.
或者那些教育社區.
那種感染教導的成果.
所以弟兄姊妹.
你現在有嬰兒.
有小朋友很重要.
從小就要怎樣培訓.
能夠好好栽培.
其實我們會看到.
今天一個現象.
很多基督徒.
也會慨嘆.
今天為什麼福音的威力.
不能夠震盪這個世界.
為什麼我們仍然.
不能夠好好積極傳福音.

$^{201}$阿斯利神學院的院長.
就這樣說.
各樣就是因為.
作為基督徒的你和我.
每一個人.
雖然我們會不惜身外物.
但卻是我們未曾好好地.
將我們那個自我拿出來.
讓我們靠賴上帝的恩典.
不要死抓住.
騙住那個自我.
保持那個自我.
而是願意放手.
然後對神說.
隨你的意思走.
要讓基督在你和我身上.
和生命裡面.
做主做王.
弟兄姊妹.
我們真的要醒悟和警醒.
那個自我中心.
所引發的一連串影響我們.
這種生命屬靈的光景的病態.
就是自我中心的徵狀.
求主光照幫助.
第二樣.
我們看到這個法利賽人.
第二樣的表現.
他的禱告和靈命的表現.
就是論斷比較.
第九節那裡說.
藐視別人.
其實現代中文翻譯為輕視別人.
你看看第十一節.
他說我不像別人.
也不像這瑞麗.
看不看到.
這個法利賽人的批評論斷出現了.
你請注意.
他是藉著一個很神聖的舉動行為.

$^{241}$是甚麼.
是祈禱.
祈禱是否很神聖.
他用祈禱的方法.
這個屬靈活動.
從事他自己不良的目的.
公然夠膽在父神面前批評論斷人和.
漠視只有神才是審判主的真理.
似乎真的當神沒有到.
因為論斷批評是沒有愛心的事.
是不屬靈的表現.
所以弟兄姊妹.
今天如果你聽到一些任何弟兄姊妹.
向你批評某一個人.
或者論斷某一件事.
不妨你靠主用愛心說誠實話跟他說.
你這樣說這樣做.
是不屬靈的表現.
是不合神心意的行為.
你要為這個人和這件事.
好好交托守望禱告.
這個反而是一個正確的態度.
因為神不喜悅人.
批評論斷.
神知道論斷是會啟爭斷.
是留地步給魔鬼和人的私怨.
撒旦會乘虛而入.
所以弟兄姊妹我們要高度提防和警醒.
特別在教會當中.
我做了二十多年的牧者.
看到批評論斷是很蠶食和殺傷.
我們是新興的教會.
我們更需要警覺這件事.
你又看法利賽.
他很熟讀律例典章.
精通舊約魔洗五經.
他知道什麼叫盡心盡意盡性盡力的愛主.
他還知道怎樣要愛人如己.
但是他沒有用.
知多多頭腦多聰明.

$^{281}$走不出來.
都是沒有用.
這件事你看到他.
在批評人的過錯.
控訴人的不是.
舉報人的不對.
比較人的軟弱.
這些正是法利賽屬靈的大問題.
聖經學者巴赫萊博士驗證指出.
沒有一個看不起同伴的人.
能夠好好的禱告.
在禱告中我們從不會將自己的名字高舉在別人之上.
我們要記住我們自己都是犯罪受苦憂傷人群中的一份子.
都是跪在父神面前.
慈悲的寶座面前.
真正的禱告是必然發自我們自己的生命.
是放在屬靈的上帝生命旁邊作一個比較.
而不是跟人作一個比較.
如果我們能夠將自己的生命放在奇妙的耶穌的生命旁邊.
和放在神聖這位神的面前作一個比較的時候.
你和我只能夠說一句話.
上帝啊.
開恩可憐我這個罪人.
巴赫萊博士這番話是一個很精準的提醒.
好幾年前在馬來西亞舉行第六屆華福大會.
青年士工組的工作坊.
就21世紀的領袖的質素作一個研討.
結果他們結論得出來.
你做一個教會的領袖或者做一個好好事奉神的人.
你有四個C這個質素.
第一個C是Confident即是信心.
第二個C是Committed即是委神.
第三個C是Competent即是有能力.
第四個C是Compassionate即是愛心.
愛心佔其中的重要位置.
當你不批評不論斷比較.
正是具體彰顯愛心的其中一個重要的表現.
讓我們留意我們在生命當中.
有沒有這樣不足的地方.
求主加能賜力.

$^{321}$幫助我們.
勒住自己的舌頭.
不隨意批評論斷.
法利素學人第三樣的禱告.
或者他的靈命表現.
是甚麼呢.
我們看看.
就是驕傲自義.
第九節說將著自己的義人.
然後第十二節.
他說我一個禮拜禁食兩次.
凡我所得的都捐上十分之一.
猶太人的律法只規定一天.
絕對必須禁食的就是救贖日那一天.
但有些人想多蒙神的恩典.
所以在禮拜一,禮拜四.
他們都禁食.
當時利未人.
以以色列人出產的十分之一作為酬勞.
但當時法利素人.
將所得的都捐出來.
十分之一.
即使不需要捐上十分之一的東西.
他都捐上.
他整個態度正是當時典型的.
法利素人的作風.
我們對照一下我們今天的愛神.
愛到好像這個法利素人.
有很好的禱告.
甚至禁食禱告.
為神的家守望.
一個禮拜有一次.
他不止一次.
這個法利素人是兩次.
我們有為愛主.
關顧神的家這個教會.
成為教會的需要.
我們有沒有甘心主動.
熱切的奉獻十分之一呢.
你對教會金錢奉獻.

$^{361}$父神有沒有一些特別的感動.
要你可以突破進一步做呢.
這個法利素人即使不需要捐上十分之一的.
他都捐上.
但是會不會今天我們.
連起碼的十分之一.
都很難為的去做.
或者吝惜的去做.
弟兄姊妹.
新約時代要求不是十分之一這麼少.
而是全然奉獻.
今天我們做父母的.
做長輩的.
不妨教一些小的.
從小.
從孩童時期.
在在學的青少年.
兒童.
挑戰教導他們從小養成起碼奉獻的.
十分之一的習慣.
如果他一直這樣操作.
我擔保.
他出來社會工作的時候.
他會覺得是理所當然.
他不會感到為難.
從小培訓他.
大把做父母的.
你有沒有教你的小朋友.
回教會要奉獻.
試一下這樣.
甚至有人計算過.
如果會有裡面20到25個.
真真正正實行十一的奉獻.
就差不多可以養得起一個傳道人.
聖會教會人人如果這樣做.
我想我們現在的全職.
勵仙傳道.
肯定教會就不用憂心經費的問題.
是不是.
再說回法利賽人.

$^{401}$他為了自己所謂虔誠愛神.
而自滿驕傲.
其實在教會歷史上.
曾經記載一個叫西門的拉比.
他居然大言不慚說.
整個世界上.
兩個二人.
嘩!真是厲害.
這兩個二人.
一個是我.
一個是我的兒子.
但如果只是選一個.
就已經目熟了.
你看這個人多牙刷.
這個拉比和這個法利賽人.
都是一擔擔.
在上帝面前.
他們真的不是好好禱告.
而是向父神炫耀.
告訴神知.
我多好多聰明多聰明多厲害.
這正是屬靈生命的致命點.
驕傲自義.
你不知道有沒有看一本.
另一個雜誌.
除了中信雜誌之外.
一個姐妹的雜誌.
叫做《傳書》.
他記載一個見證.
有一個華人企業家.
叫做李學勉.
他做一個見證.
他說我對那些銀行家.
律師開宗明義講明.
我很多東西都不懂.
我是門外漢.
但我不怕學.
我也樂於虛心請教.
他說最重要的一點.
為什麼我會請教呢.

$^{441}$因為信仰使我變得很謙卑.
信仰使我的價值觀有所改變.
信仰我是很重視.
我也不會只看外在的形象問題.
而是怎樣從心裡學到虛心請教.
這樣我才學到東西.
一位屬靈長者更加講.
提醒人.
特別是在教會生活很久的人.
他說一個領袖.
或者一個事奉神的人.
最大的危險就是驕傲.
驕傲到神沒有辦法再用你.
如果你想繼續被神使用.
你就要常常生活言行裡面.
流露出謙卑.
學著幫助我們學習謙卑.
以致靈命更長進.
真的.
剛才說的那三樣東西.
自我中心.
論斷比較.
驕傲自如.
都是法理賽人的靈命.
我們要引以為鑑.
說一個很正面的人物.
就是瑞麗的禱告和她的靈命.
我們看看.
她有兩樣顯著的表現.
第一.
她是謙卑虛己.
第十三節.
常.
瑞麗遠遠站著.
連舉目望天也不敢.
這個瑞麗是遠遠的站著.
感到自己不配走前一點去親近神.
她那種謙恭虔誠的低頭.
舉目望天也不敢的表現.
充滿了謙卑柔順和虛己.

$^{481}$這個正是她蒙福.
得著恩典的秘訣和人生態度.
剛才我說的華人成功企業家李學勉.
他說我做生意.
我的商場原則.
就要一定.
不同外面其他人.
我要做一個謙和實在.
不斷學習新知的人.
他披露自己的心路歷程.
他說照著人的本性.
很少人願意承認自己不行.
但是信仰卻是影響了我.
使得我變得很謙卑.
他又說.
信仰使我學習將事情一一交託給神.
我也看到今天我所有的一切.
如果不是神的特別恩典.
我什麼都不可以擁有.
是的.
今天不是神的恩典.
你的健康.
你所擁有的機會.
你的家庭子女如孫.
你怎麼可以誇呢.
上兩任的華福總幹事.
叫高雲康.
一個牧師.
我曾經跟一些人說過.
他生前已經被主接了去.
他生前有腎病.
有一段時間.
嚴重到每天要洗腎五次.
人看來半條命就快死了.
但後來有一個姐妹.
將腎臟移植給他.
很成功.
移植腎臟.
然後健康好轉.
不過每天仍然要吃三十顆藥丸.

$^{521}$這麼軟弱的人.
你也想不到.
他卻是擁有兩個博士學位.
直接或間接建立教會有八間.
還有五所禮拜堂.
還有派出不少的孫教士.
從自己的堂會出去.
這麼有本事的人.
卻認為要他擔任華福總幹事.
他心感不配.
他說我有四筆.
哪四筆呢.
他說我年紀不輕.
身體不壯.
恩慈不大.
經驗不多.
你看看.
謙虛.
新加坡神學院前兩任院長.
沈納德牧師.
他是我以前在香港浸會大學.
另一個學系的學長.
有一年我們基督徒團契.
舉辦夏令會.
他擔任大會主席.
我是大會的靈修.
當時要他們查羅馬書.
他是一位很重視禱告和讀經的弟兄.
是一個很謙和的弟兄.
更是神極重用的一個人.
他在幾年前已經從新加坡神學院正式退休.
但是他服務過那間神學院.
有些學生為他做見證.
其中一個說.
我想我永遠不會忘記沈院長.
他說有一次我病得很嚴重的時候.
床都起不來的時候.
他竟然親自走到宿舍.
帶了一些東西來餵我吃.
向我表達一個溫情.

$^{561}$這件事對他來說是舉手之勞.
但是他這樣的舉動.
我一生銘記感激.
沈納德牧師當時不是院長.
他只是當時神學院的訓導主任.
想一想一個嚴肅紀律.
鐵面無私的訓導主任.
竟然好像藍丁格爾那樣.
溫情慈愛的樣式流露出來.
這是一個謙卑服侍的表現.
虛己不以身份地位的表現.
成為教會弟兄姊妹.
這是挑戰我們.
我們在教會裡面.
我們願不願意謙卑虛己去愛人.
就由你對身邊的人做起.
瑞麗的禱告和靈明弟弱.
我們要看到的就是.
她吻吾同悔.
第十三節下半節.
她說只是隨著空說.
上帝啊開恩可憐我這個罪人.
法利賽人和瑞麗同樣呼求神.
上帝啊.
但只是想到自己的義.
根本沒有顧及神.
但瑞麗卻一心想到自己.
我是一個罪人我不配.
她只是全心仰賴神的慈悲.
舉目望天她都不敢.
瑞麗真是有一個很親切.
吻吾和痛苦的心.
我相信如果你是做了基督徒很久.
你就要從實這個吻吾痛悔的心志.
否則你麻木了或者慣性了.
你又想想我自己心靈空間.
有留多少空間給神呢.
我是否真是全心仰賴神.
純一真智的.
檢視到自己的屬靈光景呢.

$^{601}$剛才我們唱歌的歌詞都說我們很忙.
很迷失.
有時真是的.
你說人在江湖身不由己很多時候.
在這些處境裡.
我們真的求主給我們一個開放.
切切親近神的心.
用一個樂於回應神的心態去行事為人.
好幾年前有個神學家說到明.
今天基督徒最致命的罪.
就是對犯罪無動於衷.
不覺得一回事.
如果你想心靈安息寧靜.
就必須先得著神的赦免.
第一步就必須直接的.
自己去正視這個醜惡的罪.
神啊求你賜給我一個清潔的心.
正直的靈.
求主監察.
你為所做的能夠痛悔.
然後回轉歸向天父.
正如路加福音十五章.
那個浪子回頭的故事一樣.
今天我們聽的東西不少.
但這些對我們好像很疏離.
沒有切身的感動和推動.
這一番話是值得我們提醒.
有一個老益思C.S. Lewis.
他說一般有兩種人.
一種是自視過高.
很驕傲自大.
傲慢專橫.
另一種人是妄自菲薄.
很低看自己.
自我形象很破碎很低落.
樣樣退縮怕事.
其實這些都是人的劣根性和軟弱.
今天我們說要做一個土神喜悅的基督徒.
我們真的要靠主去吻悟.
讓聖靈光照我們.

$^{641}$我們能夠正視面對.
針對自己人性的軟弱.
和一些劣根的罪性.
然後有回應.
深切的痛悔.
然後行動的願意去對付.
好好地回轉歸向父神.
這樣才有幸福喜樂平安.
正如《耳塞書》55章第7節說.
惡人當離棄自己的道路.
不義的人當除掉自己的意念歸向耶和華.
耶和華就必連戌他.
當歸向我們的神.
因為神必廣行赦免.
《耳塞書》57章第18節.
更加說到神會對悶悶痛悔罪的人是怎樣.
那裡說.
我看見他所行的道.
也要醫治他.
又要引導他.
使他和那一同傷心的人.
再得安慰.
原來神的醫治.
引導.
安慰.
會臨到悶悶痛悔的人身上.
我真的求神給我們每一個人.
有一個無偽的心.
有一個正直的靈.
這是真切痛悔.
祈求主幫助的正確態度.
更是個人和教會復興之道.
但願成會教會人人都是這樣.
請主妹.
讓我們的禱告反映出我們好的靈明.
好的信心.
馬丁路德曾經將信心界定為.
是對上帝活潑的不問情由的信任.
很有意思.
祈求神幫助你和我多多操練禱告.

$^{681}$是更多的禱告.
更多貼近神的禱告.
以及對神有一個活潑.
不問情由的信靠.
一個信任依靠.
若果是這樣.
父神應許說.
「這人回家去,比那人」.
即是發離錯人.
「倒算為二了」.
我們低頭,祈禱.
求主你幫助監察我們.
我們不單是聽道.
我們更被你的靈感動.
要做一個行道,守道,榮耀你的人.
神你加能,賜力.
祈禱奉主耶穌基督的名.
阿們.
冬瓜.
\newpage



\section{}
\label{sec:OhwVkVNV6k4}
\textbf{4/29/18 主日講道: 基督徒成長12步--\#5 小組團契}
\newline
\newline
連結: \href{https://youtube.com/watch?v=OhwVkVNV6k4}{\texttt{https://youtube.com/watch?v=OhwVkVNV6k4}} ~~~~ 語音日期: 2018-04-30
\newline
\newline
\hyperref[sec:a492NMYDqm0]{\small{< < < PREV SERMON < < <}}
~
\hyperref[sec:index_chronic]{\small{[返順時目]}}
~
\hyperref[sec:index_scriptual]{\small{[返順卷目]}}
~
\hyperref[sec:Y6D4skNOieM]{\small{> > > NEXT SERMON > > >}}
\newline
\newline
$^{1}$早晨各位.
今天很高興.
各位看到這盞燈.
在這幾個月.
我們有很多.
精彩的演出.
希望能夠在這些.
消息和訊息當中.
提供幫助.
如果看我們今天的節目的話.
每一節會有不同的演員.
每週節目.
可能不到下星期.
我們請了一個本身是.
伊斯蘭教徒.
信奉耶穌.
本身是伊拉克人.
伊拉克巴格達本地人.
他們經歷過.
在那裡洗禮禮.
在那裡讀大學.
在那裡福音.
在那裡修煉.
是一個很激烈的演出.
下星期他們會親自.
和我們分享.
為什麼呢?因為下星期的週末.
周末.
在大西洋區的教會.
有一個叫愛沃的教徒.
愛沃的教徒.
愛沃是一個穆斯林.
伊斯蘭的信徒.
他是其中一個.
他來自美國.
來到美國.
現在在俄羅斯.
來到紐西蘭.
在俄羅斯.
有很多阿拉伯人.

$^{41}$很多在中東.
維班爾.
因為以前美國的.
汽車廠請了很多錢給他們.
所以有很多.
阿拉伯人.
這可以說是.
在中國.
有很多.
在中國的穆斯林.
伊斯蘭教會.
很多人.
上星期沒有錯過.
當然這個星期.
沒有錯過.
但是.
一個很重要的數字.
因為現在我們經常.
很難說.
就是認識伊斯蘭.
有一個朋友.
在伊斯蘭.
有一個朋友.
有在俄羅斯.
有在印度.
在一個很難說的地方認識.
這個宗教.
這個文化.
因為我們是最好的.
所以千萬不要錯過.
我們有教過很多.
上次有拉伊斯蘭教會.
有拉伊斯蘭教會.
上次有拉伊斯蘭教會.
我們有教會.
UV成長.
我們的生命.
我們的質素.
我們能夠.
保證你.

$^{81}$你自己親自來到.
我們印度.
我們每個人都要保守.
我們都是法輪功的.
你將我們.
在耶穌基督的恩典當中.
你去寫在我們的罪.
給我們全新的生命.
你叫你的子民.
你叫你的弟子.
都能夠成為一個新的.
Humanity.
新的一部Community.
讓我們有Community.
能夠彼此的信任.
能夠在最好上得一點.
自己在這個中間.
能夠將其自己保護.
最後多謝你.
在這一年開始之後.
我們有4個月.
有GA.
調整.
希望我們能夠.
繼續努力去.
賣下主養.
在下星期.
有拉伊斯蘭的.
本身.
出世的長大.
威信耶穌基督.
我們會.
使這麼多人認識.
特別是阿拉伯人.
伊斯蘭人.
認識耶穌基督的前輩.
信仰基督的反對.
一個都是我們的期望.
因為我們身邊有很多人.
要成全我們.

$^{121}$我們由保護會.
剛進入的.
都是保護會.
我們在殺種的時候.
我們見到這麼多人.
他們普通的種子.
是大家保養成長.
成長.
我們在這幾年前.
我們都想將自己的生命再次發達.
因為你制止了我們生活的環境.
保住了我們的品食.
讓我們更不再.
受到影響.
你是一個在我們生活當中.
讓我們受害的人.
在我們生活當中.
讓我們做英明的人.
主要是.
身邊有需要的.
在苦難裡面.
在病毒裡面.
我們有很多人主要努力.
去維持自己的健康.
在我們專家辦公室.
我們有很多.
方便前人的訓練.
或者是工作.
主要幫我們需要.
的工作時可以知道我們在做什麼.
在我們訓練室裡.
我們準備.
三十年來的工作.
在他在下層樓.
我們有很多人在裡面.
我們每一句.
十二次的當下.
譬如說一個政權的教會.
一個政權的基督徒.
應該由邊界方面.

$^{161}$我們已經講了四到五十四次.
我們從福音開始.
然後我們說.
我們為何要崇拜.
崇拜在一件什麼事.
例如今日崇拜在一件什麼事.
是以我們的神為主.
亦都認識我們信仰的基甸.
和今天.
為何不崇拜.
為何這件事不是古老.
過時的事.
都要教信仰到今天.
很多人類的問題.
我們自己面對的問題.
都是講明天的事.
在講明天.
我們探索到一個真理.
我們就不在這裡.
在這裡成長.
第五方面就是小組長.
大部份人.
都需要近看生活.
這些事全部都很坦白.
不過聖經裡面.
講了很多這些.
講到近看.
不屬於我們.
神不是只關心.
我們的法子.
祂會關心.
我們的思維.
我們的信念.
我們裡面.
你這樣走是不是很乾淨.
很乾淨.
但可能你裡面.
有些事是很困難.
有些事可以幫助我們.
我們能夠安定.

$^{201}$有些事可以目標我們.
以我們能夠向前走.
就是同一生活.
聖經裡面經常講彼此.
one another.
care for one another.
彼此相愛.
彼此相護.
彼此互相擁抱.
所以同一生活.
在聖經裡面.
這個教信仰.
最重要的地方.
最得精的地方.
就是同一人.
信仰的宗教.
很多人以為是你和我的神.
或者是一個很抽象的.
being,other,other.
是和他交情.
這個教信仰.
本來都是一個.
原來是一個關係.
是一組人的班人.
一起去.
認識了信仰.
沒有了國民體實就沒有了.
這個教信仰.
這個教信仰就是廢物.
就是殘棄在暫停.
信仰和民體一定是爛泥.
我不知道你們.
誰喜歡打暴力.
或者打過暴力.
有嗎?.
(笑聲).
看過我的書或是聽過我的研究.
都是十年多前.
很著名的研究.
現在還在做的研究.

$^{241}$這個研究十多年前的時候.
我才公佈出來.
我的研究員.
主要負責的.
就是柯克頓.
柯克頓在我們現在美國.
是一個很著名的scientist.
他專門研究社會生態.
他就說.
在我的研究裡面.
他就說這個研究的名字.
叫做Bowling Dome.
他就說.
一條人打個Bowling.
(笑聲).
一條人自己打Bowling.
好不好玩?.
(笑聲).
他說,美國是一個現代化的社會.
不是單單美國.
有很多其他地方.
甚至很多已經.
所謂發展中的國家.
都比較像這樣.
就是Bowling Dome什麼意思呢?.
就是打Bowling.
就是說,他研究出來的數據.
就是說,單單說美國.
美國的社會裡面.
有些人打Bowling.
幸好有很多人.
不過原來有很多人是自己去打Bowling.
意思就是說,他不再去參與.
那些什麼.
團隊啊.
去比賽啊.
打了Bowling.
自己去打.
他其實想說什麼呢?.
他說,我們社會面對一個很重要的問題.

$^{281}$Health issue.
一個公共問題.
維生公共健康問題.
就是說,人類原來是精神上的問題.
我們好像認識很多人.
但其實我們沒有什麼朋友.
沒有什麼人去支持我們.
那個叫Social Capital.
就是說.
社會的資產.
社交的資產.
越來越少.
越來越難.
為什麼說是Health issue呢?.
因為公眾的心理.
很多的問題.
很多心理的問題引申出來.
不單止心理上.
心理上.
就是那個new system.
受命的問題會受影響.
所以他就發出了一個警告.
他說,現代化的社會.
那些不存在的事實.
越來越變成一個很難的.
就是人不在一起.
打保齡,以前是在一起打的.
打疫情,以前是在一起打的.
有很多一起去做的.
社交的工作.
已經不再是一起做的.
而是一條流通.
社群.
Community.
聖經就是整個社群.
人是離不開社群.
得的是一個順傳的.
一個稱自己.
是跟隨耶穌基督的.
但不可能是自己.

$^{321}$所以永遠都是講彼此的關係.
和公眾的關係.
今天我們來看一本聖經.
是新界共和其中.
一本書信,那本書信.
不知道是誰寫的.
是一個他墓者.
用他的墓者的心腸寫給.
當時的一班.
會眾,那班基督徒.
這本書信是由希特勒書家.
Learned to the Hebrews.
我們不知道.
作者的名字.
因為這本書.
全部都是他的講道.
他的教導.
他就是告訴當中的.
信徒.
讓他們知道.
彼此的關係,不是關係.
是不可以.
散掉的.
一定要堅持.
希特勒寫給當時.
第一世紀猶太的基督徒.
所以我叫他Hebrews.
因為有很多古董就有背景.
但不可以形容.
所以我們所謂的外邦人.
或者有太教背景的人.
有時看來看去都看不明白.
其實那幫墓者.
就是用他們理解的語言.
來告訴他們.
每一方面的信仰.
每一方面的重要性.
為何特別寫給他們呢.
因為在當時.
應該是第一世紀的60年代.

$^{361}$60至68AD.
那時在社會.
有很多變化.
很多分裂的時期.
這些初期教會的信徒.
特別是猶太教.
信祖主的基督徒.
那些叫做文明的主室.
那些主室是猶太人.
猶太教的.
他們改信了.
信了耶穌.
相信耶穌是救主的御師.
教會受到很多的迫害.
和很多的打壓和威脅.
唯一最大的威脅.
對他們來說.
就是猶太人本身的意圖.
以前在教會長大.
當他們信了耶穌時.
他們就被隔絕了.
被disintegrate了.
然後他們成了一個很小的教會.
基督教的教會.
敬拜耶穌的教會.
為何有這麼大的威脅呢.
因為當時的羅馬帝國.
是宗教.
猶太教.
但是.
當開教會.
幾十年代結尾的時候.
基督教和猶太教.
開始分步了.
一者不相信耶穌是救主.
不相信耶穌是聖人.
另外一者.
教會相信天主.
和耶穌是聖父的主人.
是神的耶穌.

$^{401}$所以很多時候.
基督教就變成了不合法的組織.
在地下的民眾.
很多的壓力.
很多的威脅.
所以有很多人很心癱.
很洩氣.
就叫他們回到以前的地位.
回到以前的猶太.
回到以前的.
基督地位.
那麼.
這封書.
是給我們.
不要這麼做.
因為如果你這麼做的話.
你不知道什麼是錯.
你不知道錯是最好.
最好.
什麼叫猶太教.
猶太教是三個字來講.
三個字.
Torah.
就是九月的.
我們講的六月.
第二天.
聖殿.
GPS.
很重要.
越近就越好.
因為你不用它.
就是靠GPS.
靠你越來越近.
神就聽你的.
第三個T.
就是哲學.
他們有的Jobar.
大哲學.
所以留書的作者.
用頭十章.

$^{441}$告訴他們不要回到這個.
不要回到這個.
猶太教的舊有的生活.
為什麼呢?.
因為原來基督來的時候.
耶穌已經滿足了這三個.
耶穌就是律法的實現.
耶穌就是.
那個Temple.
你要去真正上帝.
不是藉著GPS去填.
而是藉著我們.
相信耶穌基督的那一生.
不需要藉著大祭司.
不需要藉著那些祭禮.
也不需要去守.
那些文化的尖銳.
那些本能背景的人.
其實只要是一個.
實體在.
就是耶穌基督已經到了.
所以說不要回去.
因為回去.
最真最重要的.
就是耶穌基督自己.
就是那個Fulfillment of all your desires.
整個舊的世界.
天國和聖經其實只要有耶穌自己.
你抓住耶穌就不會.
如果你回到那裡.
你就不會失去.
這句話是在.
十章十九書.
「各位大兄弟.
我們既因耶穌的血.
得而產而進入至聖上.
是藉著耶穌基督.
給我們該令一體.
又新又奇異.
從萬子經或者萬子教.

$^{481}$到其他的神明.
又有一位大祭司.
是上帝的教.
並幫我們收縮天良的.
虧欠我們的罪性.
已經灑毀了已經清潔了.
身體已經用清水來洗滌.
就當全誠心.
充足的信心.
我們就來到.
耶穌基督的聖體.
在老子的好心.
又新又活的路.
意思就是說.
不單止是Brand New的意思.
從前未有的.
舊的靈靈是靠著祭司.
祭禮,建制,聖殿.
才能來到神明.
現在耶穌基督已經進入了.
他的至聖所.
這是又新又奇異.
耶穌在神之家上所做的.
讓我們能夠直接看到神的位置.
留意地面.
看到至聖所.
看到萬子當時的聖殿.
是至聖所,即是上帝坐立的地方.
寶座的地方.
是有一條萬子的牆.
遮住.
因為聖殿的上帝.
你有罪,我有罪.
我們是不可以留任一人.
我們要留著祭禮.
但是耶穌基督在神之家上已經死了.
他的身體已經像是被擦開了.
他就是光環子.
他就是廢去左邊的光環子.
所以任何接觸耶穌基督的人.

$^{521}$都會有罪的感覺.
所以不是一個.
不是一個公正的事.
而是一群人去祈禱.
神的家.
The family of God.
我們接住耶穌的神的延遲.
我們就成為神的兒女.
在這裡.
我們沒有那麼多時間去講.
但是如果你有時間的話.
可以去看開頭的第一節.
他在講什麼呢?.
在教會大教.
他有一個節日.
每年一次.
就是那一天.
就帶著一隻牛.
帶著一隻羊.
經過大祭司獻祭.
然後大祭司進到他的寺廷所裡.
然後去跪祝祂的贖罪.
經過光環子.
進到祂的上帝的寺廟.
但是三個同日.
聖禮禮一天是不可以動的.
那個沒有做什麼呢?.
那個在聖經裡面的祭禮.
叫做Yom Kippur.
Yom就是Day的意思.
Kippur就是Cover.
就是Account.
贖罪的意思.
Yom Kippur.
是整年三個同日裡面.
猶太教最重要的High Holy Day.
是非常重要的.
而這次也一樣.
這幾個道理.
我會講一回.

$^{561}$是老天的祭祀.
一會到的時候.
有一個就是Yom Kippur.
就是贖罪祭.
贖罪的日子.
雖然他們不再有變樣.
全部都是在圖畫裡面.
Yom Kippur發生什麼事呢?.
哇!那天都是靜不住.
平時上班都不靜.
街上都不停地打.
街上都在打.
全部人都在家裡.
不做任何事.
Yom Kippur那麼厲害.
那麼厲害.
死成這樣.
我家裡都不做任何事.
因為在那裡猶太人的圖畫.
在定義基礎.
將我們廣東的同樣.
要克服基礎.
你要感情.
你要感.
因為這個人子.
就是大祭司宗子聖所.
我們叫他親人聖體.
哇!程度不得了.
我可以這樣.
為什麼?.
因為他都不開口.
在我過來之前.
他走到超級市場.
和那些公仔店.
那些公仔店不是販賣布筋.
不是販賣布筋.
是那裡的布筋.
哇!.
我真的不記得了.
哇!.

$^{601}$很危險.
我一直都在那裡.
因為我們贖罪的日子.
所以.
雖然.
藉著猶太教.
可以來到神的面前.
但是只有一天.
再藉著大祭司.
他自己都有罪.
他自己都要藉著祭禮.
來到神的面前.
希伯學書就告訴你.
Don't miss the whole point.
我問耶穌基督在十字架上.
成為一個完美的祭.
你藉著相信耶穌基督的時候.
你就直接來到神的面前.
成為真人.
你就來到神的面前.
所以《貫穿整度》.
其實希伯學書就告訴你.
一開頭第一章文字是怎樣的.
說耶穌是真人.
神的天使.
不是天使.
不是魔神.
不是約書林.
不是太平洋.
真名天使.
真名神秘.
只能夠帶我們去神的家.
就是耶穌的面前.
所以我們.
藉著相信耶穌基督的時候.
我們已經能夠成為.
神的兒女.
永遠屬於神的家.
舉個例子.
你一個人去我家.

$^{641}$有些人不肯.
有些人去我家.
有些人不肯.
有些人想試試.
我可以允許他們.
我帶他們去我家.
讓他們試試.
另一個方法.
就是做我自己.
如果你不跟我一起去.
他們就開始大人.
他們自己有自己的家.
他們自己的家.
他們就去我家.
他們不用跟我去神的家.
他們說我有時間.
我可以陪他們去神的家.
師父說好.
我可以陪他們去神的家.
好像藉著他們的感覺.
就能夠有很多的.
這就是我們神的家.
藉著神子耶穌基督.
我們就能夠成為神的家.
請你們記住.
因為說得邪惡.
他說弟兄和弟姐.
耶穌基督.
是一群人.
很多人.
不是一致政策.
是一群人.
所以在教會裡面.
我們就要學習相傳.
所以在《學術》的作者.
第十章.
他講到已經來到.
神的家中.
藉著耶穌基督.
坦然無懼.

$^{681}$我們就要.
在這個教會裡面.
學習如何建立神的家.
教會.
我們就要學習.
如何建立神的家.
教會.
我們就要學習.
如何建立神的家.
教會.
我們就要學習.
如何建立神的家.
教會.
我們就要學習.
如何建立神的家.
教會.
我們就要學習.
如何建立神的家.
教會.
我們就要學習.
如何建立神的家.
教會.
我們就要學習.
如何建立神的家.
教會.
我們就要學習.
如何建立神的家.
教會.
我們就要學習.
如何建立神的家.
教會.
我們就要學習.
如何建立神的家.
教會.
我們就要學習.
如何建立神的家.
教會.
我們就要學習.
如何建立神的家.
教會.

$^{721}$我們就要學習.
如何建立神的家.
教會.
我們就要學習.
如何建立神的家.
教會.
我們就要學習.
如何建立神的家.
教會.
我們就要學習.
如何建立神的家.
教會.
我們就要學習.
如何建立神的家.
教會.
我們就要學習.
如何建立神的家.
教會.
我們就要學習.
如何建立神的家.
教會.
我們就要學習.
如何建立神的家.
教會.
我們就要學習.
如何建立神的家.
教會.
我們就要學習.
如何建立神的家.
教會.
我們就要學習.
如何建立神的家.
教會.
我們就要學習.
如何建立神的家.
教會.
我們就要學習.
如何建立神的家.
教會.
我們就要學習.

$^{761}$如何建立神的家.
教會.
我們就要學習.
如何建立神的家.
教會.
我們就要學習.
如何建立神的家.
教會.
我們就要學習.
如何建立神的家.
教會.
我們就要學習.
如何建立神的家.
教會.
我們就要學習.
如何建立神的家.
教會.
我們就要學習.
如何建立神的家.
教會.
我們就要學習.
如何建立神的家.
教會.
我們就要學習.
如何建立神的家.
教會.
我們就要學習.
如何建立神的家.
教會.
我們就要學習.
如何建立神的家.
教會.
我們就要學習.
如何建立神的家.
教會.
我們就要學習.
如何建立神的家.
教會.
我們就要學習.
如何建立神的家.

$^{801}$教會.
我們就要學習.
如何建立神的家.
教會.
我們就要學習.
如何建立神的家.
教會.
我們就要學習.
如何建立神的家.
教會.
我們就要學習.
如何建立神的家.
教會.
我們就要學習.
如何建立神的家.
教會.
我們就要學習.
如何建立神的家.
教會.
我們就要學習.
如何建立神的家.
教會.
我們就要學習.
如何建立神的家.
社交語言.
英文有句說話叫做「lock the door」.
即是說你在平台上.
生活中有很多時間.
有很多機會.
你可能不能「lock the door」.
即是見到很多人.
或者很多人打交道.
這叫做社交語言.
即是在那裡有粉絲.
或者在那裡做公關.
或者在那裡做一些事.
一起去做一些事.
這就是最重要的因素.
我們訪問那裡的老人家.
那裡的老人家全部都是有親人的.

$^{841}$親戚.
伴侶.
孫子.
更加容易.
年輕人.
覺得榮幸可以和老人家.
共度時光.
換句話說.
這條路沒有人是冠軍.
大部分都是由face to face的.
互動.
英文就說.
比較.
chill out.
merge.
大家最好不要去這裡.
因為有社交關係.
這些絕對不是新的事.
大家可能覺得.
不是很驚訝.
可能覺得.
怎麼會有人這樣說.
沒有人這樣說.
什麼意思.
人生命是一樣的.
如果不懂就不懂.
有一個大的數據.
就是研究出.
到底有什麼人不得好.
應該是很重要的.
和這些大數據的研究.
都有很多可能.
都是一些孤心的.
半得長不可人見.
另外一個島又是很出名.
但是人不得長壽.
並且不得好.
不得風俗.
就是這裡.
亦都是一個很強的.

$^{881}$一個風俗.
對面的研究指出.
每個月開價4億.
你的壽命可以平均延長.
4至14年.
或者就算不是真的延長.
你都不覺得長了.
廣州城兩小時多.
三百多小時.
感覺上都長了.
這個真的很重要.
社會的重要性.
人們的觀念.
我最近聽到一個新聞.
在Dot.com的新聞.
有記者訪問一些無價錢的.
當地的家庭.
可能是房價值.
他們的房價值.
他們訪問人們.
特別是當天暖和.
十度以下.
聽到他們的訪問.
訪問人們的無價值的住宅.
他們說很喜歡住大型的.
他們不算.
很喜歡住那裡.
他們用很多電腦.
在Facebook上有很多人.
很多人知道他們.
住在這裡.
真的很感動.
心很痛.
有很多同行.
在這裡訪問.
我們online.
地鐵和網絡.
交易成功.
蘇俊彬說.
如果不知原來在你同屋的時候.

$^{921}$你住在這裡.
你會發揮很多東西.
原來你有design.
你有工藝.
原來你完全不感動.
是會令你感動.
他的意思是.
我們這班人和我們的show.
是變化出來的.
因為我們真是有不同的.
構圖.
和工藝.
在大陸同一個地方.
發揮出來.
我不是說.
我們住在這裡.
我不是這樣想的.
人的數據.
我們不會猜到.
是沒有實際的.
正經就是.
香港人.
對我們一班.
電光電光的.
社會的貢獻.
更不捨得.
the dispersion.
你不可以.
是Lone.
宿命上.
驚恐孤獨的話.
更加是大事.
希望的輸在.
就是要Heal你.
堅守我們自王的道理.
不容許我們的都會.
是Amen的信神.
不容許.
彼此相互激發內心.
勉勵彼此.

$^{961}$24小時.
看到這篇文章.
Let's proceed.
中文翻譯我才翻出來.
我仔細查過.
只有在國內現代的.
當代的.
要想辦法.
彼此相互激發內心.
勉勵彼此.
即是說.
我們要想辦法.
想清楚我們如何.
可以彼此激發內心.
君子的意思是甚麼.
君子的意思好像是.
今早起床時化妝.
很intent.
去check自己.
pay attention.
留意.
這個意思是甚麼.
我們彼此相愛.
我們相互看待.
我們需要.
Let's proceed.
Let's think about.
其實想一想方法.
我們如何彼此相互激發.
然後才能夠.
為內心.
Social, spiritual, human.
激發內心.
Stir up.
精明點可以幫人打氣.
但這絕對不代表負面.
負面上就是.
人與人之間激氣.
吵架也是.
level 激發.

$^{1001}$所以有同志與人吵架.
可能會長期激發.
五樂五會.
很煩惱.
不認識他們.
可能會被人激發.
可能會互相打麻將.
我準備的時候.
我覺得五樂五會.
會開始很早.
因為五樂五會.
都差一點時間.
最後突然間.
心臟有事.
但他.
很壯.
好像很瘦.
但他很壯.
精英是很高的.
但他有一樣東西.
自己不做,別人也做.
他喝咖啡買杯茶.
然後做.
我做什麼.
他就做.
但現在.
原來不是.
那個action.
是interaction.
我自古以來都沒有玩遊戲.
沒有人可以有那麼多interaction.
激大家也好.
但interaction doesn't matter.
let's go.
很多人都說.
英國有間遊客的安老所.
他搞麻將網絡會.
黑迪爆.
麻將的會.
打死圈.

$^{1041}$那時候我們用那些.
安老那些會.
打電話.
大家去聊聊天.
一人分出來.
有些人說.
我不做這件事.
那我不做這件事.
我想讓我回憶一下.
讓我回憶一下.
ok.
我明白那時候.
那幾分鐘的小學.
好像很無聊.
但是.
social integration就像小學一樣.
對嗎.
有很多人來教會.
你對嗎.
為甚麼.
因為他很想你.
但教會本身就有.
你走這一步.
你會相識你會很相處.
然後認識不到大家.
我知道我們很難遇上.
我們要說一些很深的東西.
我問你.
你和你的知己.
有誰第一見面的時候.
你馬上說.
沒有啊.
你要很多時間去擦身.
去擦肩膀才能夠去見到那個社會.
有嗎.
不可停止進步.
保障產生一定之破爛的東西.
但你留意這裡.
這本書的作者.
他沒有說一些現在好像很時尚的.

$^{1081}$一些神話.
甚麼是時尚神話.
我自己上網上班就多了.
我都是聽這裡的.
你都不明白.
好.
大班人講了很多東西.
中文有幾個.
英文有幾個.
現在我上網的youtube.
有很多人在看.
重點是人體的心.
不是只是眼.
是we.
一起做一些事.
我上網的網絡.
給大家收的資料.
都是we.
是很大的輔助.
是輔助.
都說不夠.
我聽了很多.
有啊.
不是寫部份.
是we.
我想講講.
Jack的一個故事.
他是一個.
現在因為看聖經.
就信了主.
的基督徒.
他十二年來.
他沒有試過去教.
他十二年想找人幫他守則.
都找不到.
他是一個會教.
理想教的聖人.
他想找一個故事.
他第一次去基督徒教會.
就上網.

$^{1121}$其中一個.
一個英國人.
美國.
一個人.
上網查基督徒.
第一次問他.
他問你什麼故事.
他說.
他說很想找一個地方.
教會.
他想找一個地方.
他就問了一個地方.
他就說.
我會自己來幫.
我會幫他考慮.
我會practice.
我會實行我的.
基督徒的信念.
基督的名義.
教會.
我會幫他考慮.
然後我會幫他找聖經.
他就問.
他想找.
一個主題.
基督的信念.
彼此的信念.
我們講講二元的行為.
一個人.
這個人大學的弟弟.
Edward.
二元.
有一個宣傳.
我們建立了屬靈的.
友情.
這個群體就形容了.
為什麼我會想去大學.
為什麼會有integration.
為什麼會有兩樣東西.
我有幾個空間.

$^{1161}$想去大學.
我去大學.
這就是我的大學.
我可以.
完成我的.
我會做特務的義務.
我會做特務.
關於這個.
這就是我的大學.
思考二元.
是.
Fuselsmart.
友情就是.
Safehouse.
就是你自己的安全地帶.
或者找一個地方.
找一個團體.
小組.
我會小組.
我會新人.
我會.
我不會說.
不過.
我會小組.
我們開一個小組.
每個人都會有一個大法.
一起.
全部都是我們的東西.
這就是我的.
這樣.
我會小組.
沒有的.
想想.
其他.
其他人的影子.
建立一個.
很激進的團體.
沒有人可以理會.
上次我們將會有教會.
每一個地方的教會.

$^{1201}$我們都.
會有一個頭來.
讓我們在神經的不同.
處理的時候.
可以跟進.
很多學生的想像.
學生團體的心.
幫忙教會.
我們大家有警鐘.
讓我們在十字架底.
能夠互相的認識.
大家不要分開去.
讓我們在十字架底.
去知道.
我們之間的長遠不相干.
在我們這個.
教會.
靠近的.
一個所謂的.
一起有神意.
互相激發.
互相勉勵.
謝謝.
\newpage



\section{}
\label{sec:Y6D4skNOieM}
\textbf{6/3/18 主日講道: 基督徒成長12步\#7 恩賜事奉 (陳崇基牧師)}
\newline
\newline
連結: \href{https://youtube.com/watch?v=Y6D4skNOieM}{\texttt{https://youtube.com/watch?v=Y6D4skNOieM}} ~~~~ 語音日期: 2018-06-11
\newline
\newline
\hyperref[sec:OhwVkVNV6k4]{\small{< < < PREV SERMON < < <}}
~
\hyperref[sec:index_chronic]{\small{[返順時目]}}
~
\hyperref[sec:index_scriptual]{\small{[返順卷目]}}
~
\hyperref[sec:qSAJBUcq9PQ]{\small{> > > NEXT SERMON > > >}}
\newline
\newline
$^{1}$我們由一月開始,每一個月的主日.
我們說一個整傳的屬靈的生命有十二部.
我們眨眼已經過了六次了.
上次說了受洗和見證.
這次我們說到恩賜是逢.
我不知道你覺得.
說到恩賜是逢的時候.
是一回什麼事呢?.
你想起一回什麼事呢?.
由我第一次開始說的時候.
我們說每一個信了主的基督徒.
我們知道信了福音之後.
我們不是坐著等待將來我們就去天家.
其實我們有工作要做.
因為基督去到完成了救恩.
他死而復活.
升天作為王的時候.
他賜給我們很多很多的恩賜.
然後這個王國就要來到地上.
所以我們有一個工作要去做.
所以我們想起是逢的時候.
我們就需要知道.
我們要為國度,為神的命來做些什麼呢?.
建立教會.
這就是恩賜是逢的這一環.
我們需要去明白.
很多時候我們當作是逢.
只不過是幫教會的忙.
誰誰出聲了.
我們不好意思推他.
Mark Sir 出了聲.
祈禱也不用了.
我幫幫他吧.
因為當作是我們做家務.
當作教會裡面也是我們的家.
沒有人洗碗.
我去洗一下.
純粹幫一下家務.
這個是不是就是恩賜是逢呢?.
今天我就想跟你說.

$^{41}$這個主題.
雖然可能你聽過.
但是是一個很重要的一環.
討論到恩賜是逢的時候.
聖經裡面特別是新約.
最多討論的就是一本書.
叫做《哥林多前書》.
如果你第一次翻開聖經.
或者你不知道聖經的來龍去脈.
那些新約裡面有一些書卷.
如果這本書叫做《哥林多前書》.
就是寫給在哥林多教會.
哥林多的弟兄姊妹.
基督徒的一些教導.
哥林多前書講了很多這方面.
恩賜與是逢.
十二章第一節.
那裡說什麼呢?.
他說.
弟兄們.
論到屬靈的恩賜.
我不願你們不明白.
恩賜原有分別.
聖靈卻是一位.
即是也有分別.
主卻是一位.
公用也有分別.
上帝卻是一位.
在眾人裡面運行一切的工作.
這本書叫做《哥林多前書》.
由十二章到十四章.
都是講到屬靈恩賜的.
講到上帝給我們一些技能.
一些不同的能力.
或者一些不同的態度.
去做成主的工作.
十二至十四章都是講一件事.
整封書信裡面.
十六章.
十二至十四章都是講一件事.

$^{81}$有五分一的篇幅.
保羅寫信的人.
那位牧者.
都是講到屬靈的恩賜.
什麼叫做屬靈的恩賜呢?.
Spiritual gifts.
英文叫做.
留意那個字.
有兩塊.
第一塊是形容詞.
第二塊是名詞.
第一塊是叫做Spiritual.
就是屬靈的.
什麼意思呢?.
什麼叫做屬靈的.
這個形容詞.
其實就是屬於聖靈的.
就是由聖靈來產生.
由聖靈神自己.
來工作所產生出來的效果.
這個是Spiritual.
後面那塊就是名詞.
現代英文.
有時候翻做Charisma.
就是說魅力.
現代人翻譯就叫做魅力.
通常英文翻譯就是Spiritual gifts.
禮物.
就是上帝賜予.
在聖靈裡面.
祂親自來給你一些東西.
我們作為現代人.
我們有時候搞錯了.
我們以為Spiritual gifts.
就好像聖誕禮物.
就是上帝給了你.
這個就是屬於你的.
你就可以自己去.
私有化自己去使用.
其實Spiritual gifts.

$^{121}$不是這個意思.
其實可能有一個更好的翻譯.
應該叫做Spiritual manifestation.
就是屬靈.
聖靈所產生的表現.
聖靈所產生的工作.
怎麼看呢?.
在第七節.
他說聖靈顯在各人身上.
是要叫眾人得到益處.
這裡就說到聖靈所產生的工作.
Spiritual的工作.
保羅這裡特別用了一個字叫做顯在.
或者顯明.
顯示出來.
Manifestation.
就是說上帝的靈在我們心靈當中.
只要我們信了主之後.
我們心口真心承認.
耶穌為我們主的時候.
聖靈就在我們的內心.
在我們的心靈裡面.
其實是在工作的.
你不要管祂的工作.
是怎麼樣.
是要你哭三天才止眼淚.
然後你就悔改過來.
你有改變.
還是你立刻很平靜.
去經歷得救的過程.
每個人都不一樣.
但是這裡說得很肯定.
就是聖靈在我們身上.
有不同的Manifestation.
有不同的顯明.
信了主之後.
原來上帝的工作就在我們身上.
來更新和改變.
我們慢慢慢慢改變過來.
舉個例子.

$^{161}$信主之前.
可能你有一些興趣.
有一些能力.
有一些技能.
你只是拿來自我陶醉.
就是自我發展.
但是信了主之後.
上帝去改變你的時候.
你又讓上帝來改變你.
那件事就不是純粹自我陶醉.
就變成去服侍其他人.
要叫別人來得到益處.
信主之前.
音樂可能是你自我陶醉的方法.
信了主之後.
慢慢上帝改變你的時候.
你就拿來服侍別人.
信主之前.
可能你煮飯已經很厲害了.
色香味什麼都可以.
你可以是自我陶醉的.
但是信了主之後.
上帝繼續改變你的時候.
你就找到位置來服侍別人.
用你自己的技能.
用你自己的態度.
甚至你的嗜好.
甚至你的品性.
就用來洗其他人.
洗教會.
洗神國裡面的人.
能夠得到益處.
所以這個就叫做恩賜.
就是聖靈在你生命當中.
工作的彰顯.
的manifestation.
active manifestation.
你想一想.
你信了主.
到現在.

$^{201}$是不是有所改變呢.
是不是有些事情.
上帝去到擴闊了你的視野.
或者擴闊了你的角色.
弄大了你的崗位.
或者影響的範圍呢.
這就是聖靈所做的工作.
聖靈的恩賜.
第八節繼續說.
保羅在哥林多前書裡面說.
原來聖靈在教會當中.
賜下很多各種多元化的恩賜.
不同的能力.
他說:這人蒙聖靈賜他智慧的言語.
那人也蒙這位聖靈.
是同一位聖靈.
賜給他智慧的言語.
知識的言語.
又有一個.
蒙這位聖靈賜他信心.
還有一位蒙這位聖靈賜給他醫病的恩賜.
這裡說的不是做醫生.
可能是一些人.
上帝特別經過他.
去做一些醫治的工作.
又叫一人能行二能.
又叫一人能作先知.
又叫一人能辨別諸靈.
分辨是非.
一些真與假的道.
又叫一人能說方言.
一些天國裡面的奧秘.
用一些言語表達出來.
又叫一人能夠翻方言.
這裡我們不能夠逐個逐個說.
但這是其中一個新約聖經裡面說的.
屬靈恩賜的列表之一.
整本新約聖經裡面有好幾次.
有五次是有這樣的列表.
我只是想將這個列表告訴你.

$^{241}$讓你有種感覺.
到底上帝在教會裡面做過什麼.
以下這個列表.
就是將五個屬靈恩賜的列表.
我將它放出來.
剛才我們看第一個 column A.
第一個最左的 column.
就是哥林多前書十二章.
列了好幾個的恩賜.
然後另外那四個 column.
另外那四段聖經.
就說出不同的恩賜.
就是神原來在教會當中.
在我們基督徒的生命裡面.
有不同的顯現.
這些顯現你的感覺就知道.
是很 diverse(多元)的.
是很多元化的.
沒有一個是一樣的.
用一個現代的比喻來說.
其實就像一個工具箱.
教會就是上帝的工具箱.
你知道我們做這件東西.
這件三角的東西.
你知道用了多少種工具嗎.
我數一數.
買了二十種工具.
你猜是不是要二十種工具才做得到.
要尺子.
然後要三角形的那個角尺.
就是鑿著九十度的那個.
然後要一把鋸.
鋸要用電的.
要用電又怎樣.
就要充電.
我們做事的那班人.
我們要戴手套.
已經六樣了.
要釘子.
要鑽釘又要有pre-drill.

$^{281}$不然木頭就會裂開.
那就要鑽嘴.
要鑽嘴就要鑽那把東西.
鑽那把東西要上螺絲.
又要另一把又要上螺絲.
數一數 買了十幾樣.
原來教會就是上帝的工具箱.
每個人都好像不同的工具一樣.
你都去做成上帝自己的工作.
上帝原來做我們.
帶領我們都不一樣.
你有你的獨特之處.
我也有我的獨特之處.
但是我們的能力.
我們的技能.
我們的品性.
我們不同的興趣.
加起來的時候.
就能夠配合到.
做成上帝要的工作.
所以你看一看.
我們就有不同的工具在這裡.
有些人好像士巴拿.
有些人好像尖尖的.
可能是螺絲批.
有些人是很粗重的.
他是鋸.
每個人都是上帝的工具.
所以今天在這個時候.
我想你做一件事.
跟著那張紗有三個人在這裡.
那三個人中間那個是你.
你就是上帝的工具.
左邊一個人右邊一個人.
你先跟右邊的人說一件事.
說什麼呢.
幸好我不像你.
現在你跟右邊的人說.
幸好我不像你.
因為每個人都是獨特的.

$^{321}$現在你跟左邊的人說什麼呢.
你不要說錯.
不是的.
一個而不多.
他很重要.
你右邊那個也很重要.
一個就夠了.
沒有人是多過一個.
因為每個人都是獨特的人.
每個人都是獨特的工具.
十二章第十一節說什麼呢.
他說這一切.
原來我們每個人不同.
每個人都是獨特.
每個人都是獨特的工具.
以至於去做成上帝的工作.
他說這一切都是這位獨一的聖靈.
一位的聖靈所運行.
隨己而隨主的意思.
隨神的意思.
來分給我們的.
就如身子是一個.
它有許多的肢體.
而且肢體雖多.
仍是一個身子.
基督也是這樣.
他用一個很有機的概念.
例子就是我們的身體.
看看我們的身體.
這個身體就像基督復活的榮耀的身體一樣.
教會一樣.
有不同的器官.
有不同的部分.
每個部分都不同.
手和腳不同.
眼和耳是不同的.
但大家不需要去比較.
因為這個比方就是告訴我們.
每個人都需要彼此的幫助.
大家試過病的人都知道.

$^{361}$有一樣東西壞了.
很嚴重的事情.
是不是.
你什麼都生.
突然生腎石就很煩.
這是我們這幾個月的佳話.
不是開玩笑的.
(笑).
不是開玩笑的.
不是關天的.
就是你有一樣東西壞了.
其他東西還不行.
你知道我們上星期電腦壞了.
我們有些很有技能的IT專業人士.
看過之後.
原來什麼事.
原來風扇壞了.
就是CPU的晶片.
上面有風扇.
我想是這個意思.
我不懂那些東西.
風扇散熱.
壞了電腦就常常關掉.
就不做事.
就是因為風扇.
早就出聲了.
我家裡有很多風扇.
我拿來放在上面就OK了.
教會也是一樣.
如果某一個部分不行.
或者不願意投身的話.
就會常常失靈.
教會很多時候會關掉.
然後再重新開機.
然後就壞機.
所以我們每一個有屬靈恩賜的人.
上帝原來就告訴我們.
我們有工作要做.
說到屬靈的不同方面.
很多元化之外.

$^{401}$這段聖經也說到.
如何去運用我們屬靈的恩賜.
上帝給了你.
你的崗位是其中一個工具.
你就要知道你如何去運用.
如何去發展你自己的才幹.
以至能夠繼續作成上帝的工作.
在哥林多前書.
這本書信裡面.
這幾章我講過了.
三章講到屬靈的恩賜.
其實他很詳盡地講到.
十二章.
他就解釋什麼是屬靈恩賜.
然後十三章就講到.
最大的恩賜是什麼.
就是愛.
就是為了其他人著想.
不是為了我自己著想.
大家很熟悉.
愛篇經常在婚禮裡面用.
其實就是在屬靈恩賜上下文裡面用.
不是在婚禮裡面用.
OK 你婚禮用.
都好都沒問題.
但最主要是講屬靈恩賜如何運用.
在愛的中間運用.
愛不是feel good這麼簡單.
愛是一種commitment to others.
不是commitment to我自己.
不只是自己開心.
self happiness.
self fulfillment.
而是對於對方.
對於所謂大我的群體.
我有什麼貢獻.
有任何屬靈恩賜都好.
如果沒有了愛.
保羅說一切都沒用.
你只不過是一個好響的鼓.

$^{441}$好響的叉.
是表現你自己.
因為上帝給我們屬靈的恩賜.
不同的才幹.
是用來造福給教會.
然後十四章就講到屬靈恩賜運用的一些regulation.
對於哥林多教會來說.
最大的問題就是他們有一些很表現性的恩賜.
譬如講方言.
我不解釋這麼多什麼是講方言.
什麼是speaking in tongues.
但那種很表現性.
很想自high的屬靈恩賜.
保羅就說.
有些regulation.
principle.
一些原則.
告訴你怎樣去使用這些恩賜.
對於我們來說.
可能不一定是有講方言的問題.
但我們同樣可能都有哥林多教會的問題.
講到屬靈恩賜的時候.
譬如十二章第一節.
他講到大英姐妹.
他說論到屬靈的恩賜.
我願意你們不願意你們不明白.
其實就是在質問他.
不要ignorant.
因為哥林多教會的人覺得自己很有學識.
很有智慧.
個個都是PhD.
我們有什麼不懂.
但保羅說.
I don't want you to be ignorant.
因為其實你很foolish.
其實你有很多東西你不知道.
但七節他說.
聖靈顯在個人身上.
是叫人得到益處.
聖靈顯在我身上.

$^{481}$我的恩賜我的才幹.
就是要為著common good.
為著大家的好處在愛裡面來運用.
不是表現自己.
而是為著其他人的好處.
今天總結一下.
幾個屬靈恩賜運用的原則.
第一個可以總結出來.
就是你是一個基督徒.
你是有屬靈恩賜的.
剛才列表裡面列了很多很多.
現代人去整理過之後.
我覺得大概有七種.
七種不同的才幹與恩賜.
第一種是leading.
管理或者領導.
organization組織.
有些人特別厲害.
可能他的家庭背景.
他的工作等等.
第二種是show mercy.
是憐憫人的.
去關心人的.
人與人之間的關係.
那種人很通的.
EQ很高的.
第三種是服事.
服事的不要小看.
那些是躲在後面工作的.
不出聲的.
最近跟一些浸會的弟兄姊妹.
其他教會.
有些短訓隊來的.
有些是南部來的人.
有人提醒我.
南部來的人不喜歡說話.
他們寧願用手工作.
我說是喔.
上年有一班是南部來的.
Texas來的.

$^{521}$他們不是跟你說很多話的.
談天說地他們不懂.
不過你叫他們竇木.
他們就很專心地做.
寧願用手做.
有些人就是服事的因子.
坐在背後擔椅做後勤工作.
第四種是教導.
有些人很喜歡去研究聖經.
明白一些熟人的資訊.
然後跟別人分享.
這就是教導.
第五種是encouraging.
有時候叫做faith.
信心.
當教會.
當大家去到灰心或危機的時候.
出來鼓勵大家.
這些新舊約裡面有很多故事.
撒迦利亞等等先知.
都是在做這件事.
他們就是提醒神的子民.
我們不要灰心.
要去信教上帝.
這就是勸勉.
exhortation或者encouragement.
第六是傳福音.
即是傳道.
使徒保祿.
去建立教會.
去做一些新的教會工作.
去外展的工作.
要出去告訴別人教會在這裡.
要去邀請人來.
要去跟別人傳福音.
要去分享上帝的好處.
要更多人.
未信的人能夠去明白.
真理明白福音.
這些就是傳道恩賜.

$^{561}$最後一樣就是宣講.
或者先知的恩賜.
那種責備.
或者辯別是非.
判辯誅靈.
宣講指正等等.
這種恩賜.
我想問你.
你是哪一種.
或者你不是哪一種.
第一個.
你是有屬靈恩賜的.
第二.
是什麼呢.
第二.
是你的恩賜.
原來是上帝按照他的紀義去安排的.
神的聖靈.
信主之後.
他怎樣使用你.
給你不同的特別的才幹和恩賜.
你是不需要跟別人比較的.
因為上帝有他自己的主義.
是按照上帝的心意.
你不能夠說.
我很謙卑的.
我沒有恩賜.
如果你這樣說.
就是不聽上帝的話.
因為上帝說.
起碼你都會有一種屬靈的才幹.
去作成他的功.
去建立他的教會.
所以你的恩賜.
是上帝賜給你的.
第三.
你的恩賜要在愛裡面.
去建立教會.
為了眾人的益處.
為了共同的好處.

$^{601}$不是為了你自己的表現.
第四.
是你去發掘.
你去發現.
你去發展和使用你的恩賜.
你雖然有那種恩賜.
教導恩賜才算對.
比較明顯一點.
不代表你不需要去發展.
你需要多學一些東西.
然後明白怎樣去教導.
譬如有些人有音樂恩賜.
當然他需要來事奉.
但不代表你有那種才幹.
你就不需要去發展.
有些東西是需要學的.
音樂就是這樣.
其他東西都是一樣.
在這裡我想用一個很簡單的方法告訴你.
你的路線.
將來的路線.
應該朝著這個方向去看.
想像一下有兩個圈.
一個圈叫做恩賜.
即是上帝給我的獨特的才幹.
另一個圈就是我心裡面.
我們基督徒用語術語就是負擔.
有什麼負擔.
其實我最關懷的是什麼.
我最關心的是什麼.
然後這兩樣東西加起來.
很多時候就是你屬靈的崗位.
就是上帝要在你身上.
要用你來顯明他的工作.
我舉幾個例子.
第一個例子.
這個列表就說.
所有的關懷都是一樣.
我很關心小朋友的生命.
很關心他們屬靈生命的成長.

$^{641}$但恩賜有不同.
恩賜各有不同.
如果那個恩賜.
我調轉了兩個圖標.
恩賜是關心小朋友之餘.
我是想教導他們的.
我是一個教導者.
一個教師的恩賜.
我的顯明就很明顯.
我就是教主一學.
這個很清楚.
很簡單.
但如果我的關懷是小朋友.
但我的恩賜是管理.
那我就可能在背後.
做策劃去組織.
你知不知道教會約人開會很煩.
約人開會,約來約去都約不到.
所以為什麼很多教會牛龜這麼慢.
約來約去,約了半年才開一次會.
開一次會又沒有效率.
又約過.
年尾又再開一次.
說來說去,說不到任何東西.
就是需要一些策略的人.
一些組織領導的人.
來做一些這樣的聯絡.
所以你要關心小朋友.
但你的恩賜是管理的話.
你就可以做一些領導或者組織的工作.
背後可能算一下錢.
背後可能做一下registration.
做一些後勤的管理的工作.
組織的工作.
但如果你的關懷是小朋友.
你的恩賜是服侍.
那可能出現的顯現方法又不同.
我就可能去託管小朋友.
我去服侍他.
做一些服務的工作.

$^{681}$所以關懷和你的恩賜拼在一起的時候.
你就找到你的崗位.
第二個例子.
如果你關懷的是.
我們教會不是很歡迎.
剛才我不是這樣問嗎.
是真的不太歡迎.
那你的恩賜是什麼呢.
如果你恩賜是服侍的.
在背後做後勤的工作.
那可能加入我們歡迎的團隊.
但如果你的恩賜是勸勉的.
那你可能個人就需要做關懷和跟進的工作.
但這個我們有的.
我們UV Care那班弟兄姊妹.
是有做這件事的.
所有新朋友我們會有一個寫卡給他.
或者跟他接觸.
但如果你的恩賜是全副音.
你說我想教會歡迎新的人進來.
那我就需要去接觸這些微信主的人.
所以關懷和恩賜拼在一起的時候.
你就找到你的崗位.
第三個實例.
如果你的關懷是青年.
但你的恩賜是服侍.
那你可能去服侍一些.
真是實際需要的東西.
我舉個例子.
以前我在溫哥華教會.
有些弟兄姊妹.
他們是幫那些年輕工作人員.
特別是有小孩的年輕工作人員.
他們星期五因為有年輕學員會.
他們要去事奉的時候.
他們扔不下小孩.
他們就上門去幫他們事奉.
這是照顧那些事奉者的需要.
我不是親手做青年工作.
不是親手接觸青年.

$^{721}$不過我在背後做服侍.
以致我們能夠彼此配搭.
那些年輕工作人員.
父母他們能夠去到教會.
他們就去青年團契服侍.
如果關懷是長者.
你的恩賜又是服侍.
那就是做後勤工作.
有很多東西要設置.
一個團契發生一個聚會.
有很多東西.
擔椅子也有人擔.
椅子不會自己去那個地方.
你怎麼祈禱他也不會去.
總之有人擔過椅子.
你跪玻璃祈禱也沒用.
是有人擔椅子有人收椅子.
有人去清潔.
有人去預備.
關懷如果是崇拜.
崇拜的聚會.
你想做得更加好.
如果你恩賜是服侍.
那你就做後台.
不是每個人都站台.
這就好了.
有不同的角色.
大家不同的工具.
但我們有不同的角色.
我想鼓勵你.
去到什麼叫做發掘和發展.
就是你要去嘗試.
我知道有什麼恩賜測驗.
那些可以做.
但不是很能幫到.
真的落手落腳.
你試一下的時候.
你就明白.
這是第四個原則.
最後一個原則.

$^{761}$就是追求最大的屬靈恩賜.
就是在愛裡面.
我們去成長.
去到同工.
去到彼此.
你都要合作.
這是上帝要我們教會有的心意.
同樣在哥林多前書.
亦都很強調這件事.
我們大家都有不同的屬靈恩賜.
不同的才幹.
但必須要去配合.
事奉不是給家用.
奉獻亦不是給家用.
事奉不是幫下手.
不好意思.
不如我幫下.
事奉原來是上帝的靈.
在你心靈裡面啟動他的工作.
當你事奉的時候.
你同樣會發現.
神在你的靈裡面去到工作.
不是只是幫下人手.
讓人不要那麼辛苦.
而是行出聖靈.
在你心靈當中有的心意.
在愛裡面建立我們教會.
哥林多前書亦都講到.
怎樣建立.
怎樣提到我們教會是一個愛的教會.
就是聖餐.
在《屬聖經》第十一章.
未講到屬靈恩賜之前.
它講到耶穌基督去成就舊人.
我們每逢這樣行的時候.
就要紀念主對我們的恩典.
如果恩典不是繼.
你厲害嗎.
你認信耶穌基督的話.
我們就同樣在十字架的恩典裡面.

$^{801}$可以成為一個屬靈的教.
我歡迎你如果已經認信耶穌基督.
你相信主了.
整理一下你心裡面最近的狀況.
來到主的面前.
我們拿起餅的杯子.
紀念他對我們的恩典.
我們三張聖餐桌.
拿一個杯子拿一個餅.
我們回到座位.
記住我們今天守聖餐.
我想強調的是我們是一個屬靈的家庭.
在十字架底下我們無分彼此.
只要你認信耶穌基督.
接受了主的恩典.
你就是同代吃飯的.
上帝恩典.
我們祈禱.
去工作.
你為了的是建立教會.
建立你自己的國度.
主我們感恩.
不單止我們能夠去認識這個恩典.
也都聖靈在我們生命當中.
啟動了他的工作.
給我們恩賜給我們不同的才幹.
以致我們能夠.
為著眾人的逆轉.
以致這個教會成為一個更加有愛的教會.
繼續建立我們這個教會.
建立我們的生命.
當我們拿起這個餅的杯子.
我們紀念耶穌基督.
你在恩典當中.
你受苦受難.
受刑罰.
你都可以買贖我們這個生命.
在十字架上面你.
捨棄你的身體.
流出你的血.

$^{841}$成為一個最美的獻祭.
以致我們最得赦免.
我們能夠成為你的兒女.
成為你屬靈家靈的人.
你已經受聖靈在我們身上的工作.
讓我們今天拿起這個餅的杯子.
上帝你極厚的恩典.
也都讓我們知道我們生活中.
怎樣去回應.
祈禱奉主名禁.
準備好的話.
我們就拿起餅的杯子.
領主沒人同婚一個餅.
同心靈愁.
亦同領主的病.
我們這樣祈禱奉主的名.
Amen.
十二秒鐘希望見到你.
我們會講到同性戀這個問題.
願上帝祝福各位.
\newpage



\section{以賽亞書 41:10}
\label{sec:qSAJBUcq9PQ}
\textbf{1/13/2019 2018 年最多人搜尋的聖經 - 2 - 以賽亞書 41\_10}
\newline
\newline
連結: \href{https://youtube.com/watch?v=qSAJBUcq9PQ}{\texttt{https://youtube.com/watch?v=qSAJBUcq9PQ}} ~~~~ 語音日期: 2019-01-15
\newline
\newline
\hyperref[sec:Y6D4skNOieM]{\small{< < < PREV SERMON < < <}}
~
\hyperref[sec:index_chronic]{\small{[返順時目]}}
~
\hyperref[sec:index_scriptual]{\small{[返順卷目]}}
~
\hyperref[sec:_FsWKr6IpKs]{\small{> > > NEXT SERMON > > >}}
\newline
\newline
以賽亞書 41:10
\newline
\begin{longtable}{cl}
\hline
\hline
章節 & 經文 (和合本修訂版)\\
\hline
41:10 & \begin{tabularx}{0.7\textwidth}{X} 你不要害怕，因為我與你同在；不要驚惶，因為我是你的神。我必堅固你，幫助你，用我公義的右手扶持你。 \end{tabularx} \\ \\
[1ex]
\hline
\hline
\end{longtable}
$^{1}$不知道你最近在聖經或Google上在搜尋甚麼.
聖經是值得我們去搜索的.
有很多人生中的大小課題.
值得我們在神的話語中尋找.
上帝會給我們指引.
有時一些大原則或生命中的難題.
很多時會給我們一些導向.
不一定是要我們選A或B.
不一定是這樣.
但對我們的生命和心靈起了作用.
告訴我們動機是為了甚麼.
很多時上帝的話語就是給我們一種調教.
我們對準我們的生命是否對準著一些永恆的事物.
所以我們想鼓勵你去看聖經.
當然不是看聖經多就一定是屬靈一點.
或親近神多一點.
不一定的.
所以你不單要看聖經.
我們的心靈都要向上帝謙卑.
讓他來對我們說話.
我們肯去聽他講話.
這樣讀經的生活才可以幫到我們.
如果我們硬著頸汗硬著心去看聖經的話語.
是不會對我們有幫助的.
所以不妨你試一試.
無論你參與我們的讀經計劃也好.
不參與也好.
試試打開聖經的時候我們去祈禱.
祈禱上帝才開我們的心.
然後我們才進去聖經.
以至聖經我們用信心來讀.
上帝此時此地對我們說的話給我們提醒.
這樣我們的生命就會繼續向前.
我知道不是每個人都用電話.
不是每個人都喜歡用APP.
我知道新一代可能比較沒問題.
我們都鼓勵你.
不一定要加入APP和我們一起讀經.
以往我試過有一個素靈的習作.
我覺得對我有幫助.

$^{41}$就是每天我看一章的箴言.
每天看一章的箴言.
箴言有31章.
所以第一天.
今天幾號?.
13號才算.
看13章.
第13章看完了.
到明天我就看第14章.
同樣的我們都是用神的話語來親近上帝自己.
在聖經裡面有很多值得我們留意和搜索的東西.
讓我們這幾個星期來更加清楚.
到底人在找什麼呢?.
為什麼我們會在聖經裡面找這些東西.
而你在聖經裡面是想找什麼呢?.
我們就看去年2018年.
到底人們在找什麼聖經呢?.
上個星期我們講到耶利米蘇29章.
那裡講到上帝對我們的計劃.
對我們的意念.
是一些祝福的意念.
降和的意念.
我們解釋了那段聖經.
想問問你到底在講什麼.
今天我們就看另一段是去年特別.
其實這個是第一章.
根據新版本.
是最多人分享和記錄的一段聖經.
就是以賽亞書第41章第10節.
這個也很驚訝.
因為這個不是教會裡面這麼多年我們會背的聖經.
不過居然在新版本裡面是最多人分享和傳遞的一段聖經.
這段聖經說什麼呢?.
他說你不要害怕.
因為我與你同在.
不要驚慌因為我是你的神.
我必兼顧你.
我必幫助你.
我必用我公義的右手扶持你.
不知道這段聖經能不能給你一些提醒.

$^{81}$今天你可能遇見一些難阻.
或者一些迷失的時候.
這段聖經可能就是因為這個緣故.
起碼看表面.
你都知道上帝肯與我同在.
無論我遇見什麼困難都好.
祂用右手來扶持我.
正如我們所說.
我們想用一個大一點的圖畫來看.
到底這段聖經.
這些人搜索的聖經.
到底想說什麼.
大家知道《以塞瓦書》.
有些人說是一本很難看的書.
我們剛剛開始讀經計劃的時候.
我們三十多人大概都能拍得上.
都追得上.
剛才吳牧師說了一句什麼話.
「虎頭蛇尾」.
開始的時候當然容易.
到你看到《以塞瓦書》的時候就不是了.
是一本很難看的書.
因為有很長.
當然不是最長的《以塞瓦書》.
你不要誤會.
是《耶利米蘇》才是最長的.
《以塞瓦書》都很長.
有66章.
裡面說的內容很多.
所以是很難看的.
有些人說參加一些查經班.
去看《以塞瓦書》的時候.
就覺得很困難.
因為說來說去不知道想說什麼.
所以我們要了解一下.
到底《以塞瓦書》是說什麼.
只有20-30分鐘.
到底《以塞瓦書》是想說什麼.
大概我想告訴你.
《以塞瓦書》又叫做.

$^{121}$為什麼這麼重要.
因為它又叫做Fifth Gospel.
第五本的福音書.
大家知道.
就算今天第一次來我們教會.
你是不離開教會.
不看聖經的.
你知道聖經分類.
可以分成兩大部分.
第一部分是前一部分.
就是舊約傳書.
39本書.
舊約傳書.
或者希伯來聖經.
Hebrew Bible.
舊約傳書就是第一部分.
第二部分就是新約傳書.
到底舊約和新約有什麼分別.
其實很簡單.
舊約就好像一隊音樂隊.
他們在唱一首歌.
就是說將來的救主會來了.
他很快就會來了.
舊約的師班就是唱這首歌.
新約的師班就不是.
新約的師班就是唱.
救主已經來了.
他曾經來了.
並且他將來會再回來.
所以舊約和新約有什麼分別.
你知道新約傳書裡面有四本福音書.
開頭的四本就是.
《馬太》《馬可》《路加》和《約翰福音》.
那裡就介紹耶穌的生平.
祂的傳道.
死而復活.
升天等等.
所以初信的弟兄姊妹.
剛剛想認識信仰的.
你就應該去看這四本福音書.

$^{161}$新約.
舊約的福音書.
可以叫做《爾賽亞書》.
為什麼呢.
因為那六十六章的聖經.
《爾賽亞書》.
就將整個舊約的福音.
那個是盼望.
就是說師班將.
唱那首歌.
就是說救主將會來到的時候.
他是怎麼樣的呢.
上帝怎樣去成就他的救恩.
救我們離開死亡.
離開罪惡呢.
其實在《爾賽亞書》就夠了.
就是說你將舊約所有書卷燒掉.
如果你只留下一本的話.
你都可以掌握整個舊約裡面的.
上帝想說的話.
就留在《爾賽亞書》就夠了.
這本就是《艾賽亞》.
提到《爾賽亞書》.
去年發生什麼事呢.
我想告訴你.
因為你在教會一定要知道.
這些很少人會說的.
只有我會說的.
去年發生什麼事呢.
去年有一宗新聞.
就是說在《爾賽亞》那裡.
找到一張很小的.
手甲那麼小.
一個印章.
那個印章是一個泥印章.
你知道以前的人.
他們寫完東西.
在羊皮上寫完之後.
那些經文.
他們就會捲起來.

$^{201}$然後用繩子紮住.
然後就會弄一塊泥在那裡.
然後用一塊戒指.
就會在泥上封上經書.
就是authorize(允許)它.
封上它.
或者給它一個權柄.
去年就在《爾賽亞》那裡.
就說到一宗很頭條的新聞.
就是找到這個印章.
這個印章是說什麼的呢.
這個印章那麼小.
你看到它是分上中下三部分的.
頭那部分的左上角已經脫落了.
沒有了.
就是爛了.
你看到略略是一隻鹿仔的尾巴.
和它們後面兩隻腳.
看到嗎.
看到吧.
OK吧.
是鹿仔的尾巴.
你不覺得是嗎.
當然這些是比較很多其他找到的.
人都刁鹿.
這隻當然是鹿.
難道是狗嗎.
中間那部分.
你把它翻譯出來.
就是Le-Yeshayahu.
Le是什麼意思呢.
Le就是在白俄文裡的意思.
就是belonging to.
屬於.
Yeshayahu像什麼名字呢.
就是Isaiah.
爾賽亞.
它是屬於爾賽亞.
belonging to Isaiah.
然後下面有三個或者四個字母.

$^{241}$你重構出來大概應該有四個字母.
就是Navi.
就是Prophet的意思.
Navi就是希伯來文Prophet的意思.
就是先知.
所以整句就是Le-Yeshayahu Navi.
就是屬於爾賽亞先知.
這個就很大新聞了.
不單在以色列世界裡面很大新聞.
聽說最近這東西還去了美國.
不知道是Kentucky還是哪裡.
在那裡展覽.
不是在那裡嗎.
Anyway在美國.
你把它翻譯出來之後.
你就發現為什麼這麼大件事.
因為從來沒有出過一件東西.
是跟爾賽亞有關的.
那你猜.
爾賽亞這個人.
雖然是一個希伯來文.
希伯來文裡面很普通的一個名字.
但是又是Navi的話.
又是先知的話.
似乎只有一個人.
所以這個就是第一次出現.
有聖經裡面.
爾賽亞這個人.
一件東西.
一件可以摸到的東西.
正如你看聖經要看上下文.
其實找這些東西都要看上下文.
意思就是找找它的context.
到底是從哪裡找出來的.
當然這東西不是去年發現的.
是幾十年前發現的.
不過放在博物館的儲存室裡面很久.
然後才清洗.
然後因為現在的鏡頭越來越厲害.
那些camera越來越厲害.

$^{281}$你想想你帶著那個camera.
一千塊的電話.
厲害到不得了.
所以有新的photography的technique.
才能夠將這些東西再一次呈現出來.
但無論如何.
這個東西到底在哪裡找到的.
奇怪的地方就是這樣.
這東西附近大概十呎左右.
就是我跟你前面那一行那個.
大概十呎左右.
又找到另外一顆.
那顆是什麼呢.
那顆就是這顆.
那是什麼呢.
你把它翻譯出來的話.
它裡面說什麼呢.
它就說.
Hesychia Ahaz.
Malek Yahudi.
就是說屬於Hesychia.
希西加.
Ahaz就是阿哈斯.
聖經講得很清楚.
阿哈斯的兒子希西加作王.
然後下面那一行.
就是King of Judah.
Malek Yahudi.
猶大的王.
你想想十呎以外.
這麼近的.
就找到這一顆.
中間那些是什麼呢.
那兩隻東西.
然後有個太陽.
然後有些太陽光射出來.
這樣可以吧.
明白了.
換句話說.
換句話說這兩顆東西.

$^{321}$就是用一隻戒指做出來的.
這些是很個人的東西.
這些所謂印象所產生出來的東西.
因為那隻戒指可能只有一隻.
正如你這隻結婚戒指.
你不會讓我戴吧.
我們不會說.
我們交換戴一下.
這樣很漂亮.
因為這些是很個人的東西.
所以這些印象都是很個人的東西.
這裡呢.
這顆東西是幾年前.
都是一個頭條新聞.
但問題就是這樣.
為什麼這麼近可以找到.
那那個耶穌一號.
或者屬於以賽亞的那個印象.
到底是不是先知以賽亞呢.
那個機會率就大很多.
因為聖經裡面.
以賽亞先知.
和希西加王.
大家是很熟的.
同台吃飯.
你看看以賽亞書.
37,38,39章.
或者聶王記.
你就知道他們兩個.
大家是很熟的.
因為以賽亞.
先知有很多種.
有些先知.
好像現在的總統的內閣成員.
他是那些advisor.
以賽亞就是其中一種.
叫做court prophet.
就是在宮廷裡面.
做先知的.
傳遞上帝的話語.

$^{361}$無論如何.
這個讓我們看到.
聖經裡面有很多歷史的可靠.
歷史的可靠性.
是千真萬確的.
聖經說的話.
是對在時空底下的人說的.
是很具體.
很清楚.
所以為什麼我們看聖經呢.
因為看聖經不是純粹看語言故事.
而是上帝對每一代的人說的話.
如果我們解得通.
解得清楚的時候.
我們就知道.
聖經到底現在對我們說什麼.
以賽亞書.
分兩部分.
頭39章.
就是被擄之前.
我們上個星期開始說了.
說猶太人亡國.
他們流亡70年於巴比倫的時候.
以賽亞書頭那39章.
其實就是對那些將會被擄.
將會亡國的人.
他們已經預言了.
以賽亞已經說給他們聽.
如果你們不聽話.
不信靠上帝的話.
你們就會流亡.
如果你一心不願意歸向上帝的話.
你轉向那些外邦的政權.
以為那些外邦的政權能夠幫到你.
你轉向那些外邦的政權.
外邦的國家.
等於就是拜他們的神明.
叫他們來救你.
無論是埃及也好.
是其他國家也好.

$^{401}$你們就一定會流亡.
以賽亞說將會流亡.
所以在586年BC的時候.
他們結果是流亡.
上次我們看耶利米茲29章.
就是對那些剛剛流亡於巴比倫.
散居於巴比倫的猶太人.
他們也會說話.
結果在這段時間.
當然以色列人沒有聽以賽亞所說的話.
北國以色列.
十個支派.
Lost Ten Tribes.
在亞述國之下流亡.
然後到最後巴比倫流亡.
在70年裡面.
他們就住在一個異鄉.
40至66章就是下一半.
以賽亞說的下一半.
就是對那些將會流亡於巴比倫的人.
說預言.
一早就說了給他們聽.
告訴猶太人.
你在70年流亡的結局的時候.
快要完結的時候.
上帝會跟你們說什麼話呢.
大家想想.
如果你流亡70年.
我上次說過.
我們說到exile.
說到流亡的時候.
是跟我們現在的生活.
我們的遭遇.
有一些相似的地方.
你記得我說過流亡.
不只是homeless.
沒有家的意思.
你有家.
一個你所屬的地方.
你覺得是最親的地方.

$^{441}$你不能夠回去.
這個叫流亡.
你想想70年裡面流亡.
流亡到最後的時候.
你的感覺會是怎樣.
但神的子民到這個關頭的時候.
他們就開始對上帝有很多懷疑.
或者他們有很多的anxiety.
很多的焦慮.
因為他們看到現實生活沒有什麼改變.
仍然是巴比倫的帝國作王.
或者可能現實生活裡面.
他的情況越來越壞.
前景有很多的焦慮.
有很多的anxiety.
這個起碼不是我們現代人經歷的問題.
在去年2018年.
美國的心理學會做了一個調查.
他說被訪者裡面評估一下焦慮的程度.
其實生活裡面令我們焦慮的東西只有五樣.
你說什麼焦慮.
五樣.
讓你選五樣是什麼.
這五樣是什麼呢.
健康.
安全.
財務.
關係.
和政治.
五樣已經殺光了.
我想我們沒有第六樣需要焦慮的了.
健康,財務,安全,關係.
裡面有的焦慮.
政治上有的焦慮.
這個成為一個公共衛生的問題.
美國如果你是做藥廠的就知道.
我們到了一個地步.
有些pain killers已經上癮.
藥公司當然是做藥的賺很多錢.
因為我們的pain.

$^{481}$所以引致他們可以賺很多錢.
但是焦慮不安和工作壓力.
是一個流亡的生態.
我們每一個人都在面對.
兩三個星期前.
我去一間中學做義工.
因為有些connection.
在日常生活中.
有誰可以flexible去那些中學做翻譯呢.
因為那裡有一些華人的家庭.
他們和下一代關係上有很多張力.
那個媽媽不會說英文.
要見校長.
因為那個女生不肯上學.
又說謊.
每逢早上她就說她生病.
不舒服.
所以就要告假.
越告假就越多.
或者上學又覺得自己有病.
然後就要家裡來載她.
那個媽媽不知道怎麼做.
就要見校長和校務處的主任.
我就得了毒.
我就做翻譯.
嘗試去理解那個問題.
兩邊都有壓力.
媽媽當然有壓力.
媽媽來到這裡流亡.
來這裡做什麼呢.
是為了兒女.
在中國大陸來這裡.
不是因為這裡有金汁.
而是因為希望下一代有一個好一點的教育.
她又很焦慮.
那個女兒也很焦慮.
因為新移民來到的時候.
十多歲.
學業的壓力很大.
有很多事情追不上.

$^{521}$越病得多就追不上.
越追不上就更加有病.
再更加有病的時候又追不上.
那個校長也很好.
很盡力去幫一個這樣的家庭.
我也覺得是很好的一個服侍.
我坐在這裡.
我也介紹我自己.
我做一個翻譯.
無論如何我們現代人面對很多的焦慮.
因為焦慮和將來.
和你將來的未知之數很有關.
如果你肯定到你將來的話.
你應該沒有太多的焦慮.
這是我們現代人面對的問題.
所以為什麼今年3,4月開始.
我們會請了一個Whitten的博士生.
她是在讀哲學的.
心理學的博士.
她是一個很好的姐妹.
她會來做一連串的心理健康的講座.
我覺得我們教會,聖經,講信仰的時候.
如何去關注這件事.
無論我們成人也好,小朋友也好.
如何去處理我們心理上的健康.
以賽亞書其實就是在講將來.
流亡的人到底將來會是怎樣.
講到流亡70年很多不確定的人民.
他們在等上帝來拯救他們的時候.
所以為什麼這句話會出現呢.
就是叫你們不要害怕.
因為上帝與你同在.
不需要驚惶.
因為我是你的神.
我會堅固你.
我的右手會扶持你.
聽完這句話.
你說是不是一切OK呢.
那是不是我就坐在沙發上.
我就在等呢.

$^{561}$這就是最大問題.
就是以賽亞書下一半要處理的問題.
不是這樣就OK.
不是聽完上帝與我同在我就搞定.
因為起碼我的不確定仍然在.
還有一個最大的問題就是什麼呢.
就是聽說話的人.
我們是三心兩意.
我們經常對上帝的忠心不一致.
我們這裡又看那裡又看.
我們說去倚靠上帝.
很多時候我們在生活中.
我們找很多這樣的偶像.
以為他們可以幫到我們.
最大的問題就是我們.
聽這句話的人.
我們是不是對上帝真心真意.
就好像亞國書所說.
不是心懷異議.
以賽亞書下一半有很多重點.
我抽幾個.
上帝說你倚靠我的話.
你歸向我的話.
你一心跟隨我的話.
你就有將來.
有什麼重點呢.
第一個重點.
上帝應許說他會興起另一個政權.
而這個政權會打敗巴比倫.
即是巴比倫欺負了你這麼久.
到時候我會在波斯這個地方.
興起古列.
Cyrus.
他就會打敗巴比倫.
這個是政治的問題.
解決了政治的問題.
但你一定要一心歸向我.
第二.
上帝說我會再帶我的子民.
帶你們回歸故土.

$^{601}$好像出埃及一樣.
以往在出埃及的.
怎樣離開慰勞之地.
現在你都一樣.
可以離開巴比倫.
解決了政治的問題.
解決了土地的問題.
但棘手的地方就在這裡.
猶太人縱使回歸故土也好.
你要知道.
不是成為天國.
巴比倫完了就波斯.
波斯完了是誰.
希臘.
希臘完了就到羅馬.
連猶太人都會自己鬼打鬼.
回歸之後他們仍然覺得沒有歸屬感.
仍然是homeless.
仍然是exile.
最大的問題是什麼.
上帝要解決的.
就是上帝要差派他的僕人.
他的救主.
受苦的救主.
受苦的僕人.
這個就是以塞爾蘇後半段所說的.
他要應許他的servant會來.
解決我們最根本的問題.
以致我們不應該再三心兩意.
你知道一個好的醫生.
不是只解決你symptom的問題.
是解決你根本的問題.
因為根本我們不能夠回家.
不能夠回歸的原因就是我們有罪.
我們在屬靈上和上帝破裂了關係.
所以就算你解決健康.
解決政治.
解決你土地的問題.
上帝說都不會完全解決.
唯有他的救主來到.

$^{641}$才可以一切解決.
所以以塞爾蘇後半段不是說.
我會和你和神會同在.
你儘管回歸.
以塞爾蘇後半段說清楚.
你需要離棄你自己以往的道路.
你要歸向上帝.
你才能夠真正回歸.
這段聖經我們很熟悉.
說什麼呢.
當趁耶和華可以尋找的時候就尋找.
相近的時候就求過他.
惡人當離棄自己的道路.
不義的人當除去你自己的意念.
歸向你的主人.
又說我的意念是我的意念.
我的道路是我的道路.
你先這樣離棄地.
照樣我的道路是你的道路.
這樣就高過你的意念.
我們經常將八九字拿出來.
我們覺得上帝帶領很奇妙.
是這樣的.
他的道路不止你的道路.
其實上文下你講什麼.
你不要再跟著你自己的意念.
自己好像蠻頭羊一樣周圍走.
你要歸向上帝.
因為耶和華的意念比你的意念更加高.
你要信任得過上帝.
不要信任單單你自己.
信仰就是真正的回歸.
真正的回歸是回歸上帝.
不是在生活裡面解決.
純粹我們健康平安等等的問題.
我們做牧師的經常都探病.
是不是.
探病.
我想如果我病的話.
有牧師來探病.

$^{681}$我希望牧師說我做什麼.
我希望牧師灑兩滴水.
祈個靚肚.
我馬上彈起來.
很自然的.
我一定是這樣.
如果我有困難的時候.
可能我失業.
跟牧師聊天.
我希望牧師做什麼.
我想想如果我失業的話.
我想找牧師來做什麼.
可能灑兩滴水.
祈個靚肚.
我就找到工作.
在過去這麼多年.
我跟很多人寫推薦信.
我們這個牧師行業已經不再有人尊敬.
我不知道為什麼有些人還找我寫推薦信.
我照寫.
你想找工作 想申請學校.
有什麼待遇.
你想想如果我想找一個人寫推薦信.
我希望他寫什麼.
寫一些對我有好處的東西.
寫到我有多好.
天上有地下沒有.
我也有這樣做過.
但很多時候拿了好處.
很多時候人就不見了.
這個不是真正回歸.
神說你要離棄你的意念.
不要只有一些短線的東西.
不要以為短線解決了.
上帝保佑我.
祈禱順利了.
順利了 上帝不見了.
你就不見了.
這是短線的盼望.
上帝說你一定要將你的信心.

$^{721}$你的盼望要放在神的身上.
整本《耳塞書》後一段.
都是說你將你的hope.
你的trust你的faith.
要放在上帝那裡.
我們說很多盼望.
什麼叫盼望 什麼叫hope.
《耳塞書》怎麼說.
什麼叫盼望.
盼望不是一種感覺.
你不要搞錯.
盼望也不是說積極面對.
positive thinking我們就叫盼望.
可能你想都沒想過盼望.
聖經是怎麼定義的.
聖經怎麼定義盼望或者信.
聖經定義它為一種美德.
virtue.
一種德性 一種本性.
一種character.
意思是什麼呢.
如果盼望或者信心.
在生命裡面的virtue.
是一種質素的時候.
這種美德.
無論是希伯來人的.
猶太人信仰的傳統.
或者古希臘的傳統.
新舊約也好.
說到virtue的時候.
說到美德的時候.
就不是一種感覺.
不是一種想好一點就算了.
美德是一種需要你出力.
你要去鍛鍊才得到的技能.
你有沒有學過樂器.
彈琴 彈結他 打鼓.
等等不同的樂器.
是需要下苦功的.
你需要每一天去練.

$^{761}$little and often.
你不斷來操練才能得到那種技能.
你有沒有踢過球.
你踢那個球要踢中的話.
你要不斷again and again去practice.
以致到落場的時候.
你需要踢那個feel高.
你才能踢過去.
你要訓練.
如果你訓練是沒有的話.
你就不用指望可以踢到一場好球.
新的語言也是一樣.
學一個new language.
Spanish Hebrew.
任何的新語言都是.
little and often你要去practice.
這個就叫做virtue.
virtue就是這個意思.
你要take time.
你要去每一天鍛鍊出來.
不是一朝一夕就能出來.
這個叫做信心 這個叫做盼望.
在2012年的時候.
有一個美國人叫做Colin O'Brady.
他是第一個美國人有史以來.
去到南極洲Antarctica.
你記得這個新聞嗎.
他是步行走過.
不需要任何人幫助的.
跨過南極洲.
步行跨過南極洲.
每一天要走12個小時.
每一天要走20里.
他走了932里.
總共走了54天.
他就是拖著他的食物.
拖著他的營.
然後走了900多里.
誰知道932英里有多長.
就是你現在這裡.

$^{801}$你去多不多.
然後再回來.
這個就差不多900差不多1000里.
是走的.
每一天他吃東西.
需要吃8000卡路里的食物.
所以他每一天要走12個小時.
你知道8000卡路里有多少嗎.
剛才說減肥.
我找過.
你知道一餐麻辣火鍋.
一餐麻辣火鍋大概1500卡路里.
就是說你每天要吃五餐的火鍋.
你才能夠有這樣的力量可以走.
第34天的時候.
Colin O'Brady.
在他的電話上.
推特了一件事.
他說chop wood and carry water.
他說這個是他以前訓練的時候.
在Triathlon的時候.
經常有一句這樣的術語.
chop wood and carry water.
就是說斬柴拿水.
就是他最孤單的時候.
最辛苦的時候.
他就經常記住這個術語.
就是他訓練的時候得到的.
就是恆常的一種習慣.
you stick to it.
你跟著它.
每一天當你快要瘋了的時候.
你就是chop wood and carry water.
chop wood and carry water.
做同一樣東西.
繼續走.
one step at a time.
這個就叫做virtue.
是練習出來的.
是你每一天累積出來的一種生命的character.

$^{841}$以致你能夠去克服這些困難.
這個叫美德.
Anti-White在一本書.
After You Believe, Why Christian Characters Matter.
裡面define美德是什麼.
美德是怎樣產生的.
Virtue是怎樣產生的.
就是你要做一千八百個微選擇.
專心一意地做出一些美善的事.
或者公義的事.
而這些事情不是你自然可以做的.
但是你不斷地做一千八百個這些micro-decision.
choices.
你的生命才會慢慢產生美德.
有virtue出來.
然後你一千零一次的時候.
你真的需要那種virtue發揮出來的時候.
你就自然會發揮出來.
當勇氣和智慧的選擇.
成為你自己second nature.
第二個天性的時候.
這個就是virtue.
為什麼有些人可以這麼有信心.
為什麼有些人可以這麼有盼望.
為什麼有人可以有節制.
你以為他突然睡醒就會有嗎.
是因為他生命裡面不斷地作出千萬個.
很micro 很小的small choices.
很小的選擇.
他才能累積到生命裡面的一種本性.
這是virtue.
這個猶太信仰裡面根本是一個基本的生命的韻律.
所以保羅寫羅馬書的時候.
講到我們信著之後.
他自然就留了這句出來.
他說因信稱義我們與神相和的時候.
進入了這個恩典裡面歡歡喜喜的盼望上帝的榮耀.
我們相信上帝.
盼望上帝.
不但如此患難也是歡歡喜喜的.

$^{881}$患難生忍耐忍耐生老年.
就是成熟.
老年再生盼望.
盼望不至於羞愧.
這個就是上帝的愛去瞭解我們的心.
這個就是生命裡面有的virtue.
生命裡面有的美德.
不是一次兩次出來.
今天上帝提醒我們什麼.
如果我們要歸向上帝.
要將我們的信心將我們盼望歸向上帝的時候.
你不要想一些很美很大的事.
就想今天.
有沒有一些小的選擇你需要做.
而且你是不斷拖延不肯做的.
今天就是時候你做一些所謂的mini choices.
千百個這樣的choices.
以至需要的時候.
你就可以流露出來.
最近Facebook有一個朋友.
他是做院務的.
他在Facebook上面.
他就說到一個專門研究children literature的人.
兒童故事,兒童文學的人.
他說小孩子一定要看這些兒童故事書.
因為這些兒童故事書.
如果小孩子可以入心的話.
他將來就可以develop他的virtue.
他就成為一個生命裡面有韌力的人.
你有沒有聽過一個兒童故事.
叫做大象與小豬.
Elephant and Piggy.
裡面有一個故事.
講到小孩子人格培訓的故事.
那個故事的名字叫做.
Waiting is not easy.
就是等候非容易.
就是教小朋友盼望的品格忍耐.
我想借這個時間.
我和你go through這個story.

$^{921}$然後我們結束.
有一個小孩子在youtube上.
他做了一個兒童故事.
我就拿了他來.
我用口述.
Waiting is not easy.
等候非容易.
他說Joe.
我有一個surprise給你.
他說什麼來的surprise.
什麼驚喜.
這個驚喜就是一個驚喜.
哦.
是不是很大.
他說是很大.
是不是很美.
當然美.
我們可不可以分享.
可以.
哎呀我等不及了.
但是一定要等.
要等為什麼.
因為驚喜還沒有到.
那我是不是真的一定要等.
是.
哎呀.
那我唯有等.
我會等的.
我等著呢.
等真的很不容易.
但是我很想現在看到那個驚喜.
對不起因為你一定要等.
哎呀有沒有搞錯.
我等完了我覺得.
我不覺得值得等.
我不想再等了.
好了我繼續等.
是值得的.
哎呀有沒有搞錯等這麼久.
豬仔呀我們等了很久了.

$^{961}$那我們越來越黑了.
越來越暗了.
我們看不到大家.
我們什麼都看不到.
我們浪費了一整天還在等.
我們等等等又再等.
等什麼呢.
就是等這個.
值不值得等.
相信上帝盼望上帝.
等候不容易.
但是在上帝的大令裡面絕對值得.
但是你必須要今天開始.
做千百個這些未決定.
千萬個在分秒裡面需要等候的東西.
其實你每一個主日回來.
我們去崇拜我們去仰望我們上帝.
都是一個未決定.
你去團計你去小組.
不會立刻會這麼漂亮.
是一個微決定.
你每天祈禱讀經.
你去傳道去跟別人分享福音.
你去事奉.
都是要恆心去等候.
絕對不自然.
不過值得.
願我們歸空上帝.
去動心低流祈禱.
上主我們多謝你.
因為你是信實的神.
主我們學習去信靠你.
也不只是坐下被動.
我們願意在我們生命裡面.
每一個微小的決定決策裡面.
我們學習去信靠你.
學習去仰望你.
主在我們生命當中.
你自己不斷去培養我們對你的盼望.
對你的信心.

$^{1001}$讓我們信望愛每一方面.
藉著我們生命裡面每一個微決定.
能夠去培養出來.
使我們生命更加壯烈.
主讓我們繼續在我們人生當中.
每一個階段.
我們仍然去仰望和信靠的時候.
求你賜給我們信心.
我們承認我們很多事不順.
或者對你小心.
求你再一次幫助我們去歸向你.
全心全意的去等候和信任你.
祈禱放轉美感.
拜拜!.
\newpage



\section{}
\label{sec:_FsWKr6IpKs}
\textbf{7/1/2018 主日講道: 宣教--定意研讀並遵行神的話 (Pastor Kyle Stevens)}
\newline
\newline
連結: \href{https://youtube.com/watch?v=-FsWKr6IpKs}{\texttt{https://youtube.com/watch?v=-FsWKr6IpKs}} ~~~~ 語音日期: 2019-07-03
\newline
\newline
\hyperref[sec:qSAJBUcq9PQ]{\small{< < < PREV SERMON < < <}}
~
\hyperref[sec:index_chronic]{\small{[返順時目]}}
~
\hyperref[sec:index_scriptual]{\small{[返順卷目]}}
~
\hyperref[sec:MAxHEDY7zDA]{\small{> > > NEXT SERMON > > >}}
\newline
\newline
$^{1}$.
首先想多謝各位讓我們成員教會能夠去接待他們到我們中間來服侍我們中間的青年人.
.
今年的課程叫做Game On.
.
主題是我們在宣教的遊戲裡面不要再做旁觀者.
我們走到場裡面一起玩這個遊戲.
.
所以今天我想請教一個人.
這個人就是從一個旁觀者變成一個比賽裡面的選手.
.
我們祈禱.
.
天父上帝.
.
就是這一天你所定的.
.
我們走在一起的時候無論我們是西人還是華人.
是任何的移民我們都能夠在你的主裡面聚集.
.
我們可能說的語言不一樣.
.
我們可能的背景不一樣.
.
我們過去的背景也不一樣.
.
但是在主的名下讓我們能夠一起聚集.
.
我們讓你自己的說話進入我們心靈.
.
讓我們能夠在芝加哥和世界每一個地方來繼續福音工作.
.
奉耶穌基督的名求.
.
今天我們來看的是以色列記.
.
首先在看聖經之前給一點背景.
.
在這個故事開始之前.
150年之前.

$^{41}$尼泊爾撒巴比倫王將猶太人俘虜到異地.
.
只是在耶路撒冷有一個很小的群的猶太人.
.
一代一代他們都沒有一個領袖.
.
也沒有敬拜.
.
波斯帝王讓一些猶太人從異地回到他們的家鄉.
.
不單止讓他們回到他們的家鄉.
也給他們錢來建聖殿.
.
所以容許一些猶太人.
他們懂得律法和敬拜.
回到耶路撒冷帶領子民.
.
所以我們今天看這段聖經的時候.
那個聖殿已經被建築了.
.
有一個年輕人叫以斯拉回到耶路撒冷.
.
我們看看聖經.
.
是以斯拉記第七章六至第十節.
.
我先用英文然後中文.
.
以斯拉記第七章是一個故事.
以斯拉從巴比倫上來.
他是敏捷的文士.
通達耶路撒冷以色列神所賜的摩西律法書.
波斯王准許他一切所求.
是因為耶路撒冷的神幫助他.
亞達斯西王第七年.
.
所以以斯拉回到耶路撒冷.
他不單單一個人去.
他帶領著一群同工.
他們熟悉聖殿敬拜.

$^{81}$那些唱歌的守門的利未人.
他不單單以斯拉來到耶路撒冷.
他還住在那裡.
.
他將耶和華的道帶回到這個城市.
.
今天我們在聖經這個故事.
看到兩點.
我們將它帶回到我們現在這個城市.
.
剛才你聽到一句說話兩次重覆了.
就是上帝的手一路幫助他.
當以斯拉上路的時候.
他不單單一個人去做.
他不是用他自己的能力去做.
他這樣做是因為神推動他這樣做.
而是在事奉當中有上帝的能力與他同在.
所以以斯拉有一個很大的使命.
就是將敬拜帶回耶路撒冷.
而耶路撒冷在六代裡面都沒有敬拜.
沒有人留在耶路撒冷.
是記起以往聖殿敬拜是怎樣的.
沒有人留在耶路撒冷的人.
明白上帝對他的子民說的話.
所以以斯拉就在一個這樣的情況下回到.
重新將敬拜和律法帶給他的子民.
無論你住在芝加哥.
還是在阿根廷的熱水洞.
我們作為基督的門徒.
神給我們同一個話語.
因為神創造所有人是為了敬拜他.
我們心裡面都有一種慾望去敬拜神.
但世界文化就告訴我們我們要敬拜其他東西.
我們可能敬拜金錢.
可能是權力.
可能是你的樣子是否漂亮.
所以每個人都在敬拜一些東西.
當以斯拉回到耶路撒冷的時候.
他們就不懂得去敬拜.
他們絕對沒有敬拜上帝這個心.

$^{121}$這些就是我們平日在城市裡接觸到的人.
他們心裡面有一個空隙我們需要去敬拜神.
我們的心很希望能夠敬拜這位真神.
但我們不知道怎麼做.
上帝就在他的教導中告訴我們怎樣去回答這個問題.
我們怎樣去敬拜神.
所以我們作為基督的門徒.
我們就要告訴別人神是誰.
到底耶穌是誰.
到底耶穌為我們做了什麼.
讓我們能夠進入敬拜上帝的關係.
這就是以斯拉回到耶路撒冷要做的事.
但他不是靠自己的能力去做的.
他是靠著神的靈去做的.
而我們同樣有這個能力.
我們作為基督的門徒有聖靈.
在上帝的聖靈裡我們就有能力去教導其他人去敬拜神.
但這個故事中以斯拉告訴我們還有另外一個元素.
我們看第十節.
第十節說以斯拉定義來考究遵行耶路撒冷的律法.
又將典章律例教導以色列人.
他不是回到耶路撒冷的時候自己構想一些東西出來.
他不是在耶路撒冷的時候想到時怎麼辦.
他有上帝的靈.
他有計劃.
他定義去研究耶和華的律法.
明白上帝的典章.
去學習然後去教導其他人.
去教導其他人到底耶和華的神說了什麼.
我今早去考慮講章的時候.
有一樣東西我學到.
傳統上說以斯拉記住了耶和華的律法.
以斯拉記住了聖經的頭五卷書.
摩西所寫下的所謂的陀拉.
所以裡面記住了很多東西.
所以他回到耶路撒冷的時候.
他不是在空想猜想到底要怎麼做.
到底教了什麼.
他有了,他知道了.
同樣應用在我們身上.

$^{161}$當我們告訴別人耶穌是誰的時候.
我們真的要知道耶穌到底是誰.
我們要知道耶穌說過什麼.
如果我們不定義去考究祂的話語.
讓聖經來充滿我們和挑戰我們的時候.
我們就不能告訴別人真正的敬拜是怎樣.
所以他不只是上路之際的時候才想想要怎麼做.
有清楚的計劃去研究律法.
去遵守它.
並且去教導.
所以當我們研究,學了耶和華的律法的時候.
我們同樣要遵守聖經.
在我們日常生活裡面遵守神的話.
當別人看到我們接觸我們的時候.
他們看到我們是不一樣的.
我們說話不一樣.
我們反應不一樣.
我們所言所行的裡面,我們告訴別人我們遵守神的話.
昨天他去了Potillo排隊.
每年來芝加哥一定要去Potillo一次.
所以一定要吃意大利的牛肉.
那些熱狗.
排隊的時候,後面的人.
我看到他有點特別,因為他不像群體.
他很熱,流了很多汗.
我們在排隊的時候,他有很多的改變.
他在等待的時候,他在炒零碎的零錢.
我開始和他有對話.
我問他今天怎樣.
他說他今天很熱,在外面很熱.
他很辛苦.
失業了,所以.
所以今天希望能夠賺到錢,買到東西吃.
然後找到一個比較冷的地方,能夠在那裡睡.
他就去歡迎他,一起吃午飯.
一起在同一張桌子上坐.
開始問他信仰的問題.
第一件事,他說我相信神.
我相信耶穌是我的救主.
我有點失望,因為我想和他傳福音.

$^{201}$但我很為他興奮.
他說我知道神是誰,因為以往的困難,神都帶我經過.
我說你今天有點困難.
你介不介意我現在立刻為你祈禱.
為你找到工作而祈禱.
希望你找到一個比較涼快的地方來睡.
我們一起祈禱.
我們一起分享我們的故事,說我們的家庭.
他要走了.
但他離開之前,他說他一生都住在芝加哥.
他說他從未見過一個陌生人,和我說神.
但他願意聆聽我,並且為我祈禱.
我們需要遵從神的要求.
因為我們不知道我們會站在一個人身上.
一個人在尋找正確的事情來慮拜.
我們不知道火車上的乘客是否受傷,是否有需要.
我們研究神的話語.
我們看見上帝的話語來遵從.
我們找一個時機,將神的話語告訴別人.
我們看見以色列.
我們找到一個時機來遵從神的話語.
不僅是以色列人去研究和遵從神的話語.
他們也會去教導別人.
因為以色列人回到一個地方,所有人都不知道如何真正的拜神.
很多人會以為他們知道自己如何拜神.
但其實他們不知道.
如果我們翻轉到尾端的DNI,我們會看到他們在法律上找到一個圖案.
他們會花很多時間,大群人去讀法律.
整個會眾會第一次去看神的話語.
因為他們之前不知道.
我們知道基督的道.
我們作為門徒.
不只是告訴他們耶穌是誰,讓他們自己解決.
我們會和他們同行,解釋基督告訴我們的命令.
以色列給我們一個很好的榜樣,如何行這件事.
我們第一要知道,我們不是靠我們的能力去做.
我們是藉著神的能力去做.
因為神施恩的手去幫助你.
我們要去研究和學習,誰是耶穌.
我們要去遵循.

$^{241}$我們要去教導其他人.
我們處於一個很好的位置.
因為我們有神的靈在我們心,我們是基督的徒弟.
如果我們敏感的話.
我們會在社區找到我們需要的人.
無論在芝加哥,或是在熱水湖.
我們就能夠教導這些人,如何真正敬拜神.
我現在想你做的.
你一起來低頭祈禱.
我們有時間去回應,神在我們面前的一切.
也許你不是基督的門徒.
也許你每個星期天都會來.
也許有人邀請你來.
你需要遵循,認定耶穌是誰.
你今天需要去認定耶穌是你的主.
你去回應,告訴神,我是一個罪人.
而耶穌基督就是唯一的道路.
祂為了你的罪,死在神的架上.
祂今天仍然活著.
祂是唯一的道路.
祂是唯一的道路帶領你去到神的面前.
如果你今天有這樣的反應,這樣的回應.
請找我們教會同工來聊聊天.
但可能在我們大部分的宗教.
我們認識耶穌.
我們知道祂的說話.
我們明白祂是我們的救主.
但可能我們沒有去學習祂的說話.
也許我們沒有教導祂的吩咐.
也許我們有些不順服的地方.
如果今天是你的反應.
如果你不舒服.
我們會挑戰你.
挑戰你在心裡靜靜的.
跟神說,我想你去察合你的話語.
進行你的話語,即是教導你的話語.
我想你想一想社區身邊的人的面.
也許是學生在你旁邊,在你學校裡.
也許是你同學旁邊的人.
也許是你每天購物的人.

$^{281}$你如何能夠遵從和跟他們分享基督的道理.
親愛的神父.
上帝,我們來到你面前.
主,我們來到你面前禱告.
你已經將耶穌基督的救恩給我們.
但你居然在耶穌基督這個人身上來到我們面前.
我們今天來到你面前來認罪.
我們重新立志來研究你的律法.
就像以色列一樣.
我們立志來考究你的律法.
正如以色列如何遵守一樣.
正如以色列一樣.
我們想進入別人的生命.
我們想教導別人你的道理.
主,我特別想祈求.
我們來到城會教會來禱告.
主,你給了Mark和他的團隊一個很好的視野.
去接觸這個社區.
主,你給了他們很好的資源.
去做這個.
能夠這樣做.
主,我們求你繼續來祝福我們這個教會.
讓城會教會知道.
不是靠紀靈來傳你的福音.
我們每一個的生命.
我們所接觸到的.
每一個人我們能夠去接洽的.
主,我們知道因為你的聖靈在他們的生命裡工作.
主,當我們找一些新的方法.
來跟這個社區接觸的時候.
主,我們求你.
主,我們的聖靈在這裡工作.
你給我們一個地方.
在這裡我們能夠一起來敬拜你.
主,我們無論在芝加哥.
在Hot Springs.
我們會為你給我們的資源.
你給我們神的話語的知識.
在我們社區中成為你的光.
在耶穌的名下,Amen.

$^{321}$耶穌基督明求 阿們.
(音樂結束).
\newpage



\section{}
\label{sec:MAxHEDY7zDA}
\textbf{6/30/19 主日講道:屬靈的人 (吳君才牧師)}
\newline
\newline
連結: \href{https://youtube.com/watch?v=MAxHEDY7zDA}{\texttt{https://youtube.com/watch?v=MAxHEDY7zDA}} ~~~~ 語音日期: 2019-07-06
\newline
\newline
\hyperref[sec:_FsWKr6IpKs]{\small{< < < PREV SERMON < < <}}
~
\hyperref[sec:index_chronic]{\small{[返順時目]}}
~
\hyperref[sec:index_scriptual]{\small{[返順卷目]}}
~
\hyperref[sec:1dt83c6RWRk]{\small{> > > NEXT SERMON > > >}}
\newline
\newline
$^{1}$香港最近幾個星期的局勢有很大的挑戰,.
因著修改逃犯條例的迴響,反送中大遊行,民間聲音,市民,教育界,甚至媽媽,律師,大律師,中學,大學,小學,信徒等等,他們都紛紛反對..
但特首林鄭月娥對有關要求和回應,雖然曾經以書面和親自致歉,都被泛民,民陣及青年人不收服..
但林鄭月娥的回應反映了她不會辭職,也不會設立獨立委員會調查,仍然不會撤回..
所以很多人都認為她太官僚,太官腔,甚至有人說她正在撥火,明天7月1日香港還有一個更大的遊行,問題會繼續延續下去..
有人認為林鄭月娥和政府實在太離地,今天作為基督徒,我們對時事,對社會世界局勢是否都很抽離呢?.
我們的信仰,信心,愛心以及生活表現是否能夠不離地呢?要做一個合神心意的屬靈的人,仍然是今時今日的挑戰..
我要做一個不離地的基督徒,去見證榮耀神,所以今天我用《迦勒頓大書》六章一到五節,釋經解說,如何活出一個不離地的屬靈人,我們先簡單的祈禱..
天父安靜我們的心,讓我們生活不是不吃人間煙火,而是入世又出世,如何好好地不離地地做一個土里喜悅的屬靈的人..
我們以下時間奉主耶穌基督的名禱.我們看看一些經文,《迦勒頓大書》六章一到二節..
你們若有人偶然被過犯所成你們屬靈的人,就當用溫柔的心把他挽回過來,又當自己小心,恐怕也被引誘,你們各人的重擔要互相擔當,如此就完全了基督的律法..
人有無有,自己幻以為有,就是自欺了,個人應當審查自己的行為,這樣他所誇的,就專在自己,不在別人了,因為個人必擔當自己的擔子..
這段寶貴的經文提醒我們有六方面要注意,才是真真實實可以做到一個不離地的屬靈人..
第一,能夠用溫柔的心挽回犯過錯的人..
你們看到經文,爭取時間,不讀了..
不知道對犯過錯或得罪你的人,你的態度會是怎樣呢?.
會跟他捏過,不放過,還是會另一個態度呢?.
你是不是很激動,很生氣,就罵他呢?.
還是用一個溫柔沉靜的態度去回應呢?.
哪樣更有效呢?.
曾經有一位傳道人,報道家施暴真,面對一個姐妹..
這個姐妹是一個很挑剔,講是非,詆毀人的人,但是她的性格又很敏感..
她做了這些事後,施暴真以為會向蔣則部施行教會紀律,所以他先發制人..
他向施暴真先下馬威,在情緒高漲,殺氣騰騰地跟施暴真說話..
施暴真卻是很安靜,溫柔的,細聲地回應這位姐妹..
他說目前教會還沒有任何決定會行教會紀律..
施暴真仔細地聆聽這位姐妹的感受和想法,.
為了她的敏感和誤會,施暴真竟然可以向這位姐妹反而表達一個歉意..
請施暴真以耳主的愛去包容,甚至有些不足的地方請施暴真原諒..
原來施暴真用一個溫柔,忍耐和疏導的方法令這位姐妹安靜下來,令她反省..
很多年前我曾經附和一位姐妹,她後來因一些誤會而大聲地指著我來說話..
罵完之後她很爽,我當然不爽,但她仍然不忿..
隔了一段日子,不知道,一人做事一人當,她連師母都生氣,不理睬我們..
我們跟她打招呼,想跟她聊天,她當我們透明,不聽不看,不理不睬..
幾次都令我們難受很難堪..
但我的領受是,越對這些人,我們更要加倍有愛心和忍耐..
結果在大半年時間,我學習了一些寶貴的功夫..
仍然用愛和溫柔和忍耐去挽回..
「姐妹,你是否一個溫柔的人?」.
溫柔在這個原文的意思是,不用一些語言很暴戾,殺氣騰騰的態度,.

$^{41}$或者盛氣凌人的對待,而是很柔和謙卑的語氣和和平的態度去表達..
經文裡面說的過範,其實是有跌倒,偏離,甚至是走歪了聖靈所帶領的路線的意思..
而這裡挽回,就像是駁骨折骨,更引伸有裝備,承傳,保忘的意思..
挽回過範的意思是指使一個跌倒,軟弱,走歪了的人,.
回到正途,保回走應有的路,走回原有的地位和路線..
好比哥林多前書十二章二十三節所提醒的,.
身上肢體我們看為不體面的,就越發體面,不盡美的,越發得著盡美..
甚至二十五節所提醒,免得肢體分門別類,總要肢體彼此相顧..
肢體彼此相顧是一個很重要的道理..
我深信在座的弟兄姊妹,你們一定可以明白,以及靠主做得到..
這個表現正是用溫柔的心,沒有因為持愛,忍耐,而不是激憤,不滿,怨恨,甚至對立的態度..
對一些真的有軟弱,得失理的人,有過範的人,求主幫助我們應該怎樣做,.
這樣才可以做到一個不離地的屬靈人..
第二,能警醒,自己小心,恐怕也被引誘..
有本著識書就是這樣提醒,強調,因為我們都有犯罪的傾向,可能有一天也需要別人來憐憫的..
有一項調查顯示今時今日,甚至多年來,美國的福音信仰和信徒比例,.
是從1900年總人口42.2\%降到2005年14.7\%,是逐年下降的.近年也有下降的趨勢..
美國社會正處於信仰的危機,不單止是這樣,美國的一些宗派,或者因為教會的問題,.
或者分裂,甚至有些教睦以及領袖的醜聞,社會上也有非議,大大影響了社會形象和時空..
這是值得我們每一個基督徒去守望和代禱的事..
美國教會的牧者,其實很多人都以為牧者個個很堅強很淑齡,.
其實正如經文提醒,真的要警醒,要小心,恐怕也會被引誘,誘軟弱..
其實有些令人很擔憂的現象和統計資料,令人不得不為牧者禱告和守望..
我舉幾項其中的資料,第一,美國全國每個月有1500位牧師,.
因為道德問題,或者因為工作過勞,分裂,或者教會內部的不和,而離開他們的崗位..
第二,美國全國每年雖然增加4000間教會,但同時有7000間教會關門..
第三,80\%的牧師和84\%的師母對事奉感到乏力失望..
第四,70\%的牧師經常和抑鬱交戰..
不過以上的資料數字是好幾年前的,有待繼續研究..
主要是指美國一般的西人教會,.
但美國的華人教會或少數族裔教會,未有正式的統計資料調查..
你找到的就告訴我們..
不過如果看這些數字,這些情況,當然會喚起我們每一個人要好好為教會為教牧祈禱..
有幾個令人發人深省的問題,可以幫助我們思考,值得留意..
當我們為生活奔波時,生命在哪裡失落呢?.
當我們追尋知識時,智慧在哪裡失落呢?.
當我們尋找資訊時,知識在哪裡失落呢?.
人在人生路或信仰路途上,不時很容易失落,迷失..
經歷一些事件更令人迷失..
這就是Bonnie真誠的一些剖白,大家都可以感受到..

$^{81}$但更重要的是,這個世界上有很多邪情,私慾和世界的惡者,有惑,迷惑,很多的人..
如何防止自己被引誘,如何可以小心警醒呢?.
我相信唯有你和我時時刻刻同神有一個保持緊密的關係,有警醒的心..
這樣你就可以做得到..
曾經有一個家庭主婦,她說:.
我在洗碗時,我的心對神說,願我的心不會被這些油脂所蒙蔽,能夠清潔乾淨..
當我在洗衣機洗衣服時,她又感覺和神對話,祈禱說:願我的心內的污穢被主潔淨..
當她拿著掃帚掃地時,她在心裡對主說:願我的心沒有污穢,讓我除去不義和我在神眼中看為垃圾的事..
你說這個姐妹傻不傻呢?.
其實她赤誠警醒靠著的心,我相信主必定願立..
另一位姐妹面對一些排山倒海的工作,因為除了要管家務外,她在外面也要工作..
還有一些子女的教養問題,但她仍然這樣說:.
無論我做什麼,都要好像是為神而做..
神啊,求你幫助我做世上的事,雖然人體是世俗的事,也要我做得非常神聖..
不知道你有沒有這樣的態度..
求主幫助我們對罪有敏感的心,在聖潔上多做功夫..
我不知道你個人的靈修生活,以及你聽道的時候的心態是怎樣..
是否竭力去學習呢?這樣你才有能量做一個不離地的熟人..
第三樣,我們看,能夠互相擔當重擔..
在這裡的時候,我們聽過一個英文字,Family,F-A-M-I-L-Y..
一個快樂的家庭是一個能夠互相擔當重擔的家庭..
教會是一個大家庭,也是這樣..
有些人說Family,F是代表Father,A是代表Anne,M是代表Mother,媽媽,I是代表我,L是代表Love,Y是代表You..
所以家庭就是Family,F是代表Father,M是代表Mother,I是代表Love,Y是代表You..
即是爸爸媽媽我愛你..
家庭其實是強調愛和關係,教會也是,強調愛和關係..
有人將Love再細分,L是Listen,Obligate,Value,E是Excuse,.
在彼此分擔重擔,在教會這個淑齡家庭裡面,我們要學習去聆聽,去感恩,去彼此尊重,去學習饒恕..
是不是?不容易..
為什麼在教會或者家庭裡面經常有問題發生呢?.
就是因為分擔不到,彼此不能夠相愛,當大家不和的時候會吵架..
而吵架為什麼通常很大聲呢?你有沒有聽過人吵架很小聲,陰聲細氣?.
因為當兩個人不和的時候,互相的憤怒怨恨,他們之間的心距離越來越遠..
為了填補這個距離關係,他們必須呼喊,有這麼大聲,罵這麼大聲,這樣才可以令對方聽到..
這樣他們的憤恨怨恨心與心之間,卻是造成越來越遠..
夫妻關係如是,那些階層變成懸嬰通常也有這樣的過程..
是不能夠好好彼此分擔,互相關懷,相愛..
就算同床都是異夢,甚至分床而睡,甚至同一屋簷下,冷戰..
請問你能不能夠為主而活呢?.
活就是Live,如果不能夠好好為主而活,相反或者倒轉,Live是什麼?.

$^{121}$就是Evil,活在罪惡當中,不討神的喜悅..
今日,弟兄姊妹,基督的律法總綱就是愛,能夠互相擔當重擔,愛就得以完全..
重擔在原文的意思就是為罪憂傷難過的重擔..
但其實我們都知道,你有你的重擔,我有我的重擔..
你有你的壓力,我有我的壓力,很多人都這樣說..
但是我們不要單顧自己的事,也不要顧及別人的事..
人生路上難免有困難和迷惑,當你學到不是自我中心,能夠去關顧別人,分擔別人的時候,.
你會有一些感受輕散,有人幫你分擔一些東西..
好的說話,好的行動,實在是扶人一把,讓別人的生命能夠奮起飛揚,那又何樂而不為呢?.
我們經常都要依靠主才能做到,靠自己沒有可能..
所以希望我們弟兄姊妹多說使人歡喜的話和行動,多說鼓勵的話和行動,多說肯定的話和行動,.
甚至多讚嘆的話和行動,而不是挑剔這樣那樣,多依靠主說好話行好事,少說壞話是非和消極破壞的行為..
其實一句關懷人的話和行動,能夠使那些沮喪的人有生命活下去的勇氣,.
是建立生命的工程,是生活化,是不離地的表現,所以我們都求主繼續幫助我們..
我還記得一件事發生在十八年前,我印象很深刻,那時候我們一家移民來美國,.
芝加哥的Aurora是我姐姐申請我和接待我,但她住在埠仔,Aurora你知道有多遠,.
很多時候她都愛莫能助,幸好我認識一個在香港認識的牧師,.
他比我早一年來到芝加哥,後來我找回他,他都有很好的照顧,那時候沒有車,.
他載我去消基市場,給我和師母一個小時去買東西,載我們回去,諸如此類,.
很多的照應,我們都很感激,但他住在很遠的地方,特意來,我很感動..
後來有一段時間我等候工場,這段時間可以說是沒有收入,只好支出,.
你知道,好像倒水一樣,有時候你經歷過這一點.後來我在紐約有工場的呼聲,.
我答應了,於是我又由芝加哥搬去紐約,臨走前我約這位牧師大家見面喝茶,.
臨走前這位牧師塞了信封給我,知道你要去紐約牧會,要搬運一切物品前去,.
我有一點心意,作為支持你的搬運費,你收了它吧.其實這位牧師和我不算是深交,.
他所有牧羊的教會,我都知道是一間小的教會,人數不多,他的收入都是有限的,.
所以他遞給我,我知道是什麼回事,我不收,但他硬塞在我手裡說,.
是我和太太經過禱告,覺得你有這樣的需要,是神提醒我,感動我而來的,.
你這些小小心意你收了吧,堅持了一段時間,再見,拿了信封回家拆開,.
發現裡面有一疊美金,其實數目不少,這位牧師擔當我那時候的一些需要和預約,.
我沒有跟他說,但他付出很多,十幾年之後,我現在還印象猶新,.
肢體之愛是很窩心,是令人鼓舞感動感激的..
弟兄姊妹,你又怎樣可以做到擔當人的重擔呢?有沒有關心周邊的人的需要呢?.
求主幫助我們,這不是只有口說,總要在行為和誠實上,不是在舌頭上,.
求主讓我們怎樣做一個這樣的人.第四樣,能夠誠實無偽..
有個電腦動畫,埃及王子的影片,有些人看過,.
有一幕是摩西從水中被拉出之後,他和法魯的兒子拉美西斯在王宮一同長大,.
一天這兩位埃及王子拉著車比賽,以極快的速度,追追續續,.
不理會其他人的安全,推翻了周邊的一些檔攤,建築工程的支架,.
沿途留下很多很多混亂和殘破不堪的局面..

$^{161}$後來法魯王發現他們的魯莽行徑的時候,就嚴厲責備拉美西斯,.
擔心埃及王朝岌岌可危..
但這位小王子卻是無知的回答,.
埃及王朝已經立國好幾百年了,一個表面無雙大雅的小事件應該不會影響的..
但法魯很有智慧的回答說,只要一個脆弱的環節,.
就足以瓦解整個帝國,國家如是,教會如是,家庭如是,個人如是..
小小的問題,小小的軟弱,小小的詭詐不誠實,小小的罪過,.
足以蠶食,甚至摧毀整個的生命..
誠實無偽,就是不會自欺,不會以無辯為友,而自以為友..
第三節說,神其實是參透我們裡面的心思意念的,.
你的內心世界是怎樣,你的口說得好聽,其實神是知道,騙不到..
進一步說,今天的世界的風氣其實是入侵教會,.
使教會逐漸失去獨特的聖潔和單純,.
以及誠實無偽的特質..
你回教會多,聽教會的東西多,你就知道..
有位神學家分析,當代的教會潛在一些危機,影響基督徒的,包括什麼呢?.
第一,教會中,瀰漫著普遍的反權並的氛圍,.
人與人之間有彼此不信任的傾向,.
以及一些從屬關係,他不喜歡的,.
更加會恨惡任何人告訴我應該怎樣做,.
有些團體教你怎樣做,他不喜歡聽..
第二,以人的感覺超越思想,.
今天人所關心的,不是到底人應該怎樣思考,.
而是人的感覺如何,.
最重要是合適感覺,最重要是我喜歡,.
導致變相抬舉個人的經驗和感受,.
過於神的道和聖經的真理..
第三,宗教商品化,.
舉凡一切都以數字,尺寸或預算收入和出席人數作為衡量標準,.
很多時候,外面工商界的模式成為衡量成功與否的準繩..
第四,活動取代了聖潔,.
教會時舉辦不少活動或會議等同教會的活靈,.
以忙碌多事工遮蓋了靈修,退修,安靜,.
醒察,軟弱和罪過等方面..
第五,教會在性關係議題上妥協,.
教會中的男女問題,家庭貞潔,離婚,婚外情,.
色情圖片和性戀等事項,.
教會在這方面有堅定的立場和教導,.
以致上述這一切,令教會屬靈的見證慢慢下降..
弟兄姊妹,讓我們按照聖經要求神的心意,.

$^{201}$我們不會自欺欺控自己,.
我們要以人的言行真誠無遺,誠實不詭詐..
幾年前在紐約,我牧羊教會時,.
曾經為一個施洗弟兄通電話聯絡,.
看看他的近況如何,.
那次在電話中,他跟我說,牧師,我對不起,.
我撒謊騙你,我說我來教會已經一個月,.
其實不是,我是差不多來了超過大半年,.
我騙你是因為我想搞身份而撒謊,.
這位弟兄很好,最後都能夠誠實面對一些事..
今日,我們為著弟兄姊妹能夠真誠無遺,.
對神的誠實,正直感恩,.
讓你們這份美德繼續發揚光大下去,.
正如提摩太被保羅所勸導的一樣,.
就是應該怎樣在清潔的心,無愧的良心,.
無遺的信心這三方面表現出特質..
當我們誠實無遺面對人生時,.
我們會有很多學習,.
你聽過有些很精準的提醒,.
因為他磨練了你的心智,.
因為他強化了你的雙腿,.
因為他增加了你的智慧,.
因為他覺醒了你的自尊,.
因為他教會你獨立,.
求主幫助我們誠實無遺,.
正面積極面對自己,.
以致成為一個不離地的淑女..
我們看第五樣,.
能夠審查自己的行為,.
審查,測驗就是審查,.
我們每天都要測驗自己的行為,.
有位老師教小朋友,.
上課時問,為何今天這麼黑,.
為何這學生的臉這麼黑,.
為何你們的衣服都很深色,.
那些學生不如問,.
老師,你戴了深色的太陽眼鏡,.
為何還未脫下來,.
今天我們會否戴著有色的眼鏡,.
去看人,看事,.

$^{241}$以致失真呢?.
先脫掉自己有色眼鏡,.
先脫掉自己眼中的良目,.
才可以看得清,自我反省..
要檢視自己的內心世界,.
裡面的深層思想是怎樣,.
前新加坡神學院院長陳世協博士,.
提出一個秘訣,.
就是先跟自己的眼睛納入藥,.
因為我眼養我心..
今天你所看的東西,.
無論是影碟,互聯網,.
手機,報紙,刊文,.
你要校正一下,.
這是眼睛跟你納入藥,.
是一個好的開始..
怎樣說,.
他所誇的就專在自己,.
不在別人呢?.
意思就是說,.
他不在乎跟別人比較,.
他在於自己所做的,.
是有價值的工作而自誇..
求主幫助我們所做的事,.
雖然有些人說,.
我做普通家庭主婦,.
我是退休人士,.
我家務買菜煮飯,.
或者照顧子女兒孫,.
這麼苦悶單調乏味的,.
但是都好像沒什麼價值,.
但是對準神,.
就一定一定是金銀寶石,.
是很有價值的工作..
有一個人笑著對太太說,.
成功的男士背後有兩個女人,.
哇,他太太聽到立時反應,.
哇,她說,.
一個是我母親,.
一個是太太來囉,.

$^{281}$這樣才發出會心微笑..
作為一個好的母親,.
好的太太,.
真的不簡單..
求主幫助我們眾姐妹,.
我們多測驗自己的行為,.
能夠榮神益人與否..
最後快快說完第六樣,.
能擔當自己的擔子,.
擔子是只有責任的擔子,.
是不能夠卸給別人,.
不可以做A字膊..
今天我們一起共勉,.
你自己的家庭,.
你自己的本份,.
是否全心全力去做,.
專心的致力擺上,.
最好的呢?.
神巴利在不同崗位一定有意思,.
沒有人可以代替..
有位患末期癌病的人,.
對他的好朋友說,.
從前我以為自己很聰明,.
現在才知道我自己是一個大笨蛋..
朋友問什麼意思呢?.
他很懊悔地說,.
以往我將大量的精力時間.
花在事業上,.
親人,家庭,.
我只是配丁點的時間,.
到現在這一刻,.
我癌病末期,.
我才體會人生什麼最重要..
其實類似的故事每天都發生,.
只是有沒有人.
能夠從這些人的故事裡.
吸取到教訓..
今天的教會,.
神是我的家,.
今天我們有個家事大會,.

$^{321}$求主都幫助我們,.
好好的,好好的,.
能夠忠心擺上自己,.
在家庭,在崗位,.
否則我們到主交障的那一天,.
我們會覺得羞愧..
既然神是我們的牧者,.
我們有神的同在,.
同行,同供,.
主恩典救我們用..
求主提醒我們今天三幾樣東西,.
六樣,.
第一,善用溫柔的心,.
以愛,以包容去挽回犯過或者軟弱的人,.
第二,好好的警醒,.
自己都小心恐怕被軟弱所引誘,.
第三,學習互相搭擋重擔,.
彼此以愛和行動互相守望支持,.
第三,求主幫助你和我做一個誠實無畏的人,.
時常謹慎自己脆弱的環節,.
信仰堅守立場,.
第五,時刻審查自己的行為,.
審查自己內心,.
和自己的眼睛立約,.
保守自己所看的東西..
第六,操練擔當自己擔子,.
作一個有承擔有責任感的人,.
恰如其分的作一個土著喜悅的人,.
以致可以成為一個不離地的熟人的人,.
我們誠心祈禱..
多謝主帶領我們,.
藉著加拿大書給我們一些提示,.
加能賜力讓我們走未來的日子,.
祈禱奉主耶穌基督的名,.
阿們.
\newpage



\section{}
\label{sec:1dt83c6RWRk}
\textbf{我屋企ge 主人}
\newline
\newline
連結: \href{https://youtube.com/watch?v=1dt83c6RWRk}{\texttt{https://youtube.com/watch?v=1dt83c6RWRk}} ~~~~ 語音日期: 2019-07-09
\newline
\newline
\hyperref[sec:MAxHEDY7zDA]{\small{< < < PREV SERMON < < <}}
~
\hyperref[sec:index_chronic]{\small{[返順時目]}}
~
\hyperref[sec:index_scriptual]{\small{[返順卷目]}}
~
\hyperref[sec:ip7wGoUuLe8]{\small{> > > NEXT SERMON > > >}}
\newline
\newline
$^{1}$在某一夜,我竟然會做出這樣的決定.
就是向神投降.
而且打開了道門.
和耶穌和我一起去見祂們.
在那天晚上,我看到耶穌的聖體在我身邊.
祂們的聖體是一顆聖骸,是一顆聖骸.
祂們的聖骸是一顆聖骸,是一顆聖骸.
祂們的聖骸是一顆聖骸.
和耶穌進入了一個我最隱密的地方.
我的心.
想想,基督教信仰真是不可思議.
人人的家庭都會盡量找一些最美最堅固的材料來建造.
但神選擇了他的家庭是用破碎的心靈來建造.
而是一些相信他和接受他救恩的人的心.
但這是甚麼道理呢?.
耶穌的確說過.
誰真心愛我,聽我的話去做人.
我就住在他那裡.
哇!耶穌進入了我的心裡住?.
我隨時可以和他談心事,甚至談到通宵?.
想起都覺得很棒.
而他搬進來住的那天.
雖然沒有開甚麼歡迎會,慶祝他進來住.
但我真的感受到並不是一時衝動.
就是這樣,耶穌搬進來我的心裡.
他的光照亮了全屋每一個角落.
甚至那些最隱密,從來沒有人去過的房間.
他點起壁爐的火,很溫暖,很溫馨.
他拿起那些我從來沒有認真用過的樂器彈奏起來.
我感受到每個房間都充滿了愛.
我認為那晚我開門給他.
是我人生中最明智的一個選擇.
而當耶穌進來住後.
我感受到和他的關係親密了很多.
我個人開心了,樂觀了.
而且價值觀變得很正面.
主啊,你就當這裡是你的家.
我的東西就是你的東西.
現在,不如我帶你到處看看吧.
我帶耶穌進了第一間房,就是書房.

$^{41}$也就是我的思想大本營.
這間房雖然不算很大,但這裡其實很重要.
我有很多文件,大小事情,我都在這裡決定.
而耶穌一進來這間房.
他的眼睛就橫掃了我所有書和電腦的APP.
他越看,我就越感到渾身不自在.
說起來很奇怪.
因為我一向都不覺得有什麼問題.
但現在他那雙什麼都看穿的眼睛.
令我感到很尷尬.
因為問心嘛.
我現在知道,其實我有很多書和APP.
根本就是垃圾,不見得人.
但你的東西已經放在這個思想大本營很久了.
不說是有自可.
說著說著,都覺得自己很丟臉.
我回頭對耶穌說.
主啊,我知道這個思想大本營需要翻新.
不如你幫我選好書.
順便重整我的電腦.
以後由你做我這間房的老頭好不好.
耶穌很爽快地說.
非常好啊Ben.
我來的目的就是要幫你.
你首先刪除那些不見得人的書.
包括色情或者沒益的書和APP.
買那些有關聖經的好書.
這個思想大本營要放多些聖經的道理.
就能幫你日夜專注神的話.
還有啊,掛在牆上的畫都不是對你有益.
不如你掛這一幅吧.
它會幫到你的.
於是耶穌送了一幅祂自己的油畫給我.
還叫我將這幅畫掛在這房間最中心的位置.
就是這樣.
這段日子.
我思想集中在耶穌這幅畫上的時候.
祂的純潔和威力.
就幫助我不再想那些骯髒的事.
我和耶穌從書房來到廚房.

$^{81}$因為我對食一向都很有興趣.
而且我對廚藝很有心得.
所以我留在這裡的時間都很多.
時時刻刻都想下一餐炮製些什麼吃.
今晚還煮了幾道菜給耶穌吃.
那我就和耶穌說.
這個廚房都幾寬敞,設備齊全.
我想你會很喜歡我們今晚吃的東西.
耶穌就說.
今晚吃什麼?.
今晚我親自下廚.
有我最拿手的XO醬炒神戶牛柳.
紅燒石斑球.
瑤柱豆苗黃.
吃完飯還有冰糖紅棗燉燕窩.
我想這些家常小菜.
應該都不失禮的.
反正耶穌說過.
基督徒的生命不應該單靠食物來滿足.
好了,菜都煮好了,預備好了.
是時候開飯了.
接著我和耶穌出去飯廳吃飯.
菜放上桌已經有一段時間.
奇怪耶穌仍然很沉默,沒什麼夾菜.
我忍不住就問.
耶穌說.
是不是不合口味?.
耶穌開口了.
沒有.
其實我自己預備了食物,是你看不到的.
我的食物就是中行天父的旨意.
耶穌看到我不是很明白的表情.
他很認真地看著我說.
阿Ben.
如果你想得到令你真正滿足的食物.
就要好好求天父的旨意.
不是單單滿足你的食慾.
只要你求天父的旨意.
一心討他的歡喜.
這樣就是令你飽足的食物了.

$^{121}$說完之後.
耶穌給我嘗一口聆聽天父旨意的食物.
嘩,真的很好吃.
我吃了全世界沒有一樣東西能夠彼得上的食物.
令我心靈得到很滿足.
不久之後,我們走進客廳.
這地方環境很寧靜舒適.
一大張意大利米白色的沙發.
名設計師的書架放了很多新潮.
很有型的擺設.
還有這個溫馨的壁爐.
我很喜歡在這裡享受.
看來耶穌對這裡也有好感.
祂說.
這個客廳很舒服.
不如我們多點來吧.
因為這裡夠寧靜.
能使你的心遠離生活繁瑣的事.
你和我就可以多點時間見面聊天了.
耶穌這樣說.
使我感到喜出望外.
我心想.
有什麼比單獨和耶穌見面聊天更好呢?.
耶穌約定我.
我每天早上就在這裡等你.
我們就可以一起開始新的一天.
就是這樣.
每天早上我來到客廳.
耶穌在書架上拿出聖經.
細心向我講解其中的寶貴真理.
我的心感到很溫暖.
在講解當中.
我領會到祂對祂兒女的愛和恩典.
而福音又如何可以改變人的生命.
這段時間實在太美妙了.
我們開始叫這個客廳改了名.
叫做靈修室.
很可惜.
因為生活有很多責任和壓力.
這段時間越來越馬虎了.

$^{161}$唉 我不是故意的.
但就是這樣.
每天家庭事務,朋友約會,工作壓力.
或者這些都是藉口.
總之不單時間縮短了.
而且次數也越來越少了.
有一天早上我又匆匆忙忙趕著出門.
經過客廳.
看到壁爐的火點著了.
耶穌一個人坐著.
我突然感到很不好意思.
耶穌是客人.
是我請祂來住的.
我竟然這樣對待祂.
於是我立即走進客廳.
充滿歉意地對祂說.
主啊 物理.
每天早上都一個人在這裡?.
耶穌這樣回應我.
當然了.
我不是答應了你.
每天早上都在這裡和你約會嗎?.
耶穌這樣一說.
我更加感到無地自容.
縱使我這樣不守承諾.
但祂始終很守信用地等著我.
於是我很慚愧地求祂原諒我.
想不到祂一點也不介意.
很爽快地原諒了我.
祂還很有耐性地和我解釋.
阿Ben.
你的煩惱在於你一直把靈修當作單單上課學習.
從來沒有想到是我們親密相處的時間.
我和你的關係.
我是很珍惜的.
我希望你也重視我.
無論日子如何.
也不要放棄這段我們親密相處的時間.
過了一會.
耶穌問我.

$^{201}$你家裡有沒有工作室?.
有的.
在地牢是有一個工作室的.
那裡確實有一些器材和工具.
但其實我很少用.
有空的時候.
我也會下去搞搞小玩意出來.
但從來沒有認真地做過一些有意義的事.
當我帶耶穌到工作室.
祂看到那張工作桌.
由細心觀察我懂得的技能.
祂就說.
這裡的設備也不錯.
你有沒有想過你能為神的國做些什麼出來?.
而當祂看到一兩件.
我隨便放在桌上的小製成品.
祂就問我.
在你的一生中.
你就是只做這些嗎?.
我這樣回答祂.
主啊.
我有試過盡力.
我知道不算什麼.
我也很想做多些.
只是有心無力.
奇妙地耶穌這樣回應我.
好 阿Ben.
你伸出雙手給我.
我會請聖靈透過你的雙手來工作.
如果聖靈得到你的雙手和你的心.
就一定會像度身訂做.
幫助你為神的國做一些有意義的事.
然後耶穌走到我後面.
祂神聖的雙手很溫柔地拍拍我的肩膀.
祂的意思是鼓勵我繼續為神做些事.
我肯定自己還有很多東西要學習.
而且我對自己過去做出來的事不是很滿意.
但現在我明白到.
無論為神做什麼.
都是靠著神賜給我的能力而做出來的.

$^{241}$還記得當耶穌問起我家的那間打機房的時候.
我真的情願去沒有問過.
因為裡面擺滿了打打殺殺色情又暴力的遊戲.
一定過不了那關.
但是我一向都是這樣玩慣的.
是我生活的一部分.
很多時候打機打到三更半夜.
搞到第二天上班一味釣魚.
還很多時候起不了床回主的嬰兒墳.
所以當耶穌問起的時候.
我還不扯開話題嗎?.
還記得有一晚.
不是很想留在家裡.
我就收聖約朋友去他們家玩到深夜.
之後還去了酒吧追女孩.
跟了人吃丸子.
其實不留在家裡更好.
可以更自由地煮煙,喝啤酒,賭錢.
耶穌見我要出街就跟我說.
你出街嗎?我跟你一起去好嗎?.
唉,他這樣問一問.
很想馬上拒絕他.
但是我都要尊重他.
我去的這個地方你不會有興趣的.
那好吧,明天我參加祈禱會.
到時一起出去吧.
耶穌非常失望.
好,不過我還以為搬進你家之後.
你去哪裡我們都一起去,一起生活.
就是那晚,竟然玩到天亮.
雖然整晚看起來玩得很盡興.
但是回家的路上.
感到非常孤單,沮喪.
心情很混亂,越想越難過.
回到家裡,見耶穌的房間燈還亮了.
我的心覺得很內疚,就跟他說.
主啊,對不起.
現在我終於明白了.
我人生沒有了你.
我的心無法有真正的平安,真正的快樂.

$^{281}$其實越玩到天昏地暗.
我的心就越空虛.
我決定了,從現在開始.
去哪裡都跟你一起.
耶穌聽到之後,很開心.
他還陪伴我,幫助我.
戒掉那些大大小小的不良嗜好.
也因為這樣.
我的人生體會到真正的喜樂和滿足.
從此之後,我的家充滿了歡樂的笑聲.
耶穌的家,我會永遠的愛著.
祝福你們.
耶穌搬進來住後這段日子.
我洗心革面.
放棄以前放蕩不羈的生活習慣.
時時和耶穌在一起.
現在我家裡所有的房間.
都被耶穌的愛和光充滿了.
不過,有一天我回到家裡.
耶穌在大門前等我.
臉色有點不對勁.
耶穌很擔心地說.
Ben,你回來就好了.
樓上有一陣很難聞的味道.
好像有些東西過期發臭.
應該是從屋頂那間小房間傳出來的.
唉,其實他一說,我已經心知肚明.
是,在家裡確實有間鎖著的秘密小房間.
因為我藏起了一些不讓人知道的東西在裡面.
我明知道裡面的東西已經發臭了.
但我不想扔掉.
我想是滿足我心底的一些情慾吧.
所以我很怕被人發現.
而當我們一步一步走上去的時候.
那陣臭味越來越勁.
耶穌指著這間小房間說.
就是,就是,就是這裡裡面了.
這個時候,我快忍不住發脾氣.
我心想,我已經給你面子了.
我的書房,廚房,客廳,工作室,打機房.

$^{321}$全屋每一個角落,我都聽你的話.
但你連這間四呎乘四呎那麼大把的小房間.
都不放過?.
夠了,這間房間我一定不會給你鑰匙進去的.
但耶穌卻強硬地說.
我是不會容忍和這些爛爛又臭的東西住在一起的.
我寧願睡車房,我不會遷就的.
沒錯,自從耶穌搬進來,我們越來越親密.
祂不單翻新了我的家,祂還翻新了我的生命.
我的心靈感受到從未有的釋放和快樂.
我很珍惜這個關係.
我認真地想了一會兒,就跟耶穌說.
啊,那我決定,不如我把我家的鑰匙交給你.
拜託你按你的意思,扔掉那些不乾淨和對我有壞影響的東西.
其實,我也不想這些東西挾持我的人生.
耶穌一聽到,就很爽快地在我手裡拿了一條鑰匙.
逐個房間進去,把所有不乾淨甚至腐爛發臭的東西都放在垃圾袋裡.
拿出屋外全扔掉.
然後,祂很認真地吸塵,拖地,把牆重新塗油.
轉眼間,整個房間翻新了.
這個時候,我感到整個人好像被放牢一樣.
覺得我的心靈很釋放,自由.
我一向都很無知.
以為我這間屋可以靠自己的決心保持得乾淨,整潔.
但後來我才知道,我是不可能做得到的.
因為當我清潔完第一間房,第二間房又髒了.
清潔完第二間房,第三間房又髒了.
結果,好像人的罪一樣,弄得自己殘破,瘦削,還是那麼骯髒.
所以,我鼓起勇氣跟耶穌說.
主啊,有沒有可能,我整間屋都被你管.
由你來清潔和翻新,就好像你對付的房間一樣.
我還未說完我想說的話,耶穌已經笑容滿面了.
Ben,你懂得想就好了.
如果你單靠自己,是不可能做得到的.
這間屋要保持乾淨,才可以靠我.
不過,我不是這間屋的業主.
有很多事情,我沒有權力去做.
這時候,我抬起頭,看一看這個所謂屬於我的物業.
我就衷心地說.
主啊,這間屋,雖然你過去是客人,我是主人.

$^{361}$但我現在決定將這個物業權轉交給你.
說完之後,我就走進書房,拿出這張屋契.
將業主的名字改成耶穌.
以上是一個比喻卻是真實的故事.
房屋比喻我的生命.
說的是自從我決置信耶穌.
邀請耶穌進入我心的那一刻的心路歷程.
我已經決定將整個生命都交給耶穌.
由祂做我生命的主人.
透過這個故事,我很想告訴你.
沒有任何其他方法比完全信靠和交託給主.
感到更自由,更快樂.
我祝願你也像我一樣.
邀請耶穌入住你的心.
讓祂做你生命的主人.
當你這樣做,你會發現.
你走的是一條神祝福的道路.
看到的是不一樣的天空.
\newpage



\section{詩篇}
\label{sec:ip7wGoUuLe8}
\textbf{詩篇人生 (1) - 詩篇 73 講道 (陳崇基牧師)}
\newline
\newline
連結: \href{https://youtube.com/watch?v=ip7wGoUuLe8}{\texttt{https://youtube.com/watch?v=ip7wGoUuLe8}} ~~~~ 語音日期: 2019-07-10
\newline
\newline
\hyperref[sec:1dt83c6RWRk]{\small{< < < PREV SERMON < < <}}
~
\hyperref[sec:index_chronic]{\small{[返順時目]}}
~
\hyperref[sec:index_scriptual]{\small{[返順卷目]}}
~
\hyperref[sec:1lBwTIY05FM]{\small{> > > NEXT SERMON > > >}}
\newline
\newline
$^{1}$當我們看這些詩篇的時候,有一個問題可以問的.
就是我們在讀這些詩詞歌賦.
詩篇的時候.
我們就問一問到底.
哪一個人會發出一種這樣的詩歌.
誰會唱這些詩歌.
另一個問題可以問的就是.
到底詩歌裡面所說的.
到底是什麼情況呢.
什麼人呢.
對象是在描述什麼情境呢.
套入我們今天的現代的境況裡面.
到底誰可以唱這首詩歌.
或者這首詩歌到底是想說什麼.
73篇其實很清楚分開三節的.
有三節,好像我們平時唱詩歌有好幾節一樣.
73篇總共有三節.
一至十二節.
就是一段,我們先看頭六節.
他說:神實在善待以色列.
就是神的子民.
善待那些內心清潔的人.
就是上帝和祂的子民立約.
無論是舊約裡面.
藉著摩西過紅海,西乃山等等的歷史.
或者到新約,在基督裡面.
我們是那一群和上帝立約的人.
有一個這樣的信念.
只要我們守信約的話.
神的子民守信約的時候.
神就會恩待他的百姓.
會祝福他的百姓.
God is good,神實在善待.
good to those.
那些是righteous的.
那些是願意守約的人.
比如第一篇,大家應該很熟悉.
不從來人的計謀等等.
那個是有福的人.
他只要跟著上帝的律法.

$^{41}$這個就是一個有福的人.
有良知,行得正,站得正.
行公義,好憐憫的這群人.
上帝就是祝福這群人.
其實我們教子女也是這樣的.
誠誠實實在你的學校.
在你成長的過程裡面.
你盡力做好,做好你自己.
不要做壞事.
這個就是舊約新約.
一直有的conviction.
就是他相信.
宣佈耶和華神一定會恩待那些守約的人.
但是他宣佈完之後.
他就對聖殿裡面一起assembly的那群人.
God is good.
但是第二節開始.
他說世情很多時候不是這樣的.
世情往往是相反的.
他見證了什麼呢?.
他說他自己曾經差不多跌倒.
差不多滑倒.
他說惡人狂傲人.
惡人狂傲人不是那些常常說髒話.
很不禮貌的人.
不一定是那些人.
惡人狂傲人在詩篇裡面.
舊約裡面.
就是那些不相信上帝.
不守他的律法.
不信神.
甚至是有權有勢.
在社會裡面有尊利的人.
那些人反而見到什麼呢?.
他們得到平安.
得到Shalom.
明明是Shalom.
平安.
就是給那些守約的人.
他說世情卻是相反的.

$^{81}$那些反而大享平安.
有Shalom的那些人.
他們就是那些惡人和狂傲人.
沒有疼痛.
沒有苦.
沒有災.
這些人反而有平安.
代入我們今天的情況.
就是他們有健康照顧.
還是很好的健康照顧.
他們不需要擔心任何事情.
甚至說到他們身上的打扮.
戴的項鍊.
或者穿的衣服.
都是一些強暴的.
第七節.
他說他們的眼睛.
因體叛而突出.
你不要以為聖經是針對那些胖的人.
沒有這樣的事情.
其實為什麼呢?.
你知道古代的壁畫.
比如埃及的壁畫.
畫的皇帝比較胖.
有肚腩的人.
都是做君皇的人.
或者做皇后的人.
為什麼呢?.
因為古代裡面.
很少會夠吃的.
真的很少會夠吃的.
我們平常就是十常八九.
我們吃剩的東西很多.
丟的東西很多.
特別是在美國.
我們很大浪費.
古代是正正相反的.
吃不夠是常事.
吃得夠.
甚至有餘的人.

$^{121}$通常都是一些有權有勢的人.
所以人肥.
在古代.
在聖經的世界.
這是非常的事情.
絕對不是常規.
為什麼可以有餘呢?.
或者夠吃夠住呢?.
就是因為這些人.
會欺壓其他人.
搶奪其他人.
我之前也說過.
我們要明白.
聖經裡面的經濟的話.
Economy.
我們就要明白.
聖經古代的世界.
資源永遠都是有限的.
就是說這麼多人.
我們就分一個的Pizza.
一個的蛋糕.
但是那些有權有勢的那幾個人.
1\%裡面的population.
他們就拿了95\%的蛋糕.
那幾個拿了蛋糕就走了.
然後剩下的那些.
剩下的那些鈣.
就是給我們眾人去分的.
這個就是狂傲的人.
或者有權有勢.
有專利的人.
在聖經古代的社會裡面就是這樣.
如果你這樣看的時候.
其實今天也有一點跟我們相似的.
無端端就說加稅了.
上個星期是吧.
燃油稅是不是加.
Gasoline.
有人提醒我.
有了.

$^{161}$我又忘記了.
第二天.
早上一起來.
因為我們要開車.
我老婆就說.
昨天你不進.
現在多給幾毫.
忘記了.
無端端為什麼要加呢.
又加.
那些貪得無厭的.
貪完又要再貪的.
那些就是.
眼睛因體叛而突出.
肥到眼睛都飆出來了.
第八節.
他們憑惡意來說.
欺壓人的說話.
說話是至高.
這裡不是說粗言穢語的.
可能有些人說話是很斯文.
但是卻是說一些欺壓至高的說話.
你看今天在世界裡面.
任何一個地方.
一些掌權的人.
他們說的話未必一定是粗俗的.
但是很多時候.
比粗俗的說話更難聽.
這些就是欺壓的說話.
第九節.
他們的口.
狂惡人的口.
是褻瀆上天.
他們的舌頭是毀謗全地.
你看這些詩歌的時候.
你要知道.
希伯來的詩歌特別.
就是經常有並聯的.
有A句,然後有B句.
很多時候甚至有C句.

$^{201}$以前我們以為.
A句和B句說的話是完全一樣的.
現在開始我們明白.
原來不是A句和B句說的話是一樣的.
不是redundant的東西.
而是B句很多時候.
去到elaborate.
或者去擴充.
或者補充A句所說的話.
所以B是大於A的.
所以你這樣細細去看的時候.
第二句很多時候是expand.
或者擴大上一句.
第一句所說的話.
他說他們的口是褻瀆上天.
褻瀆上天.
就是他們是褻瀆上帝的.
B句說什麼呢?.
他說他們的舌頭是毀謗全地.
你可以想像到什麼呢?.
我就想像到.
就是那隻口是很大的.
上唇碰到上天.
然後下唇就在地上.
他要吞噬的是那些平民百姓的人.
就是那隻口很大很大的.
獅子張開大口.
上唇接天.
然後下唇就碰地.
就吞噬了所有的東西.
貪得無厭的一班人.
第十節.
因此他的民就歸回那裡.
這個和合本翻譯是神的民.
其實應該是他的民.
就是那些暴君的人民.
就歸去他那裡.
然後喝盡滿杯的苦水.
這個苦水翻譯得很對.
或者喝了很多的水.

$^{241}$受了很多的苦的意思.
他們說.
就是那些狂傲的人.
神怎能曉得.
至高者豈有自識.
看下者就是惡人.
詩人說什麼呢.
他們常享安人.
財富也是家宅.
這一段是73篇.
頭的三分一.
詩人就告訴我們他所見所聞.
在他的人生當中.
他遇見很多的惡人.
看見很多的惡人橫行霸道.
這個就是那個世情.
很多時候社會世界就是這樣.
古代就是這樣.
現代也可能是這樣.
上個星期我有機會看回.
2014年我寫過一篇文.
大概二三萬字的一篇文.
那時候為什麼寫這篇文呢.
就是剛從以色列回來.
那一年2014年就是以巴衝突.
又另一個回合.
我們計劃了去那裡做的事情.
全部都打調.
因為經常都有炮彈的警號.
經常都響.
整天在那個地方到處都響.
在街上又響.
在酒店裡面又響.
然後又去防空洞.
為什麼回來要寫這篇文呢.
就是很多人去指責那些巴勒斯坦人.
說他們很暴力.
說他們譴責他們.
為什麼要傷害平民百姓.
他們在加沙那裡.

$^{281}$也有寫過很多火箭炮.
過去以色列那邊.
但是為什麼要寫這篇文呢.
我覺得很多時候我們不明白.
就是我覺得沒有那個context.
去understand為什麼人會有violence.
在以巴這個情況.
是一個官逼民反的情況.
他們射火箭炮當然不對.
但為什麼他們要射火箭炮呢.
為什麼呢.
因為這麼多年以來.
已經70年了.
他們就是被趕盡殺絕的一群人.
開頭建國的時候.
有很多村落裡面的巴勒斯坦人被殺.
被殺戮.
然後有很多人被趕絕.
回不了自己家.
然後去了難民營.
一直到現在.
有兩代人.
70年都是這樣被趕絕.
引導絕望的時候.
你就不可以叫他們冷靜.
他們有沒有所謂和平的運動.
和平的示威呢.
有,絕對有.
但是我很贊成一個.
在當地以色列的一個學者說.
他說如果政府只是懂得一種語言.
而那種語言.
是暴力的語言.
Language of violence.
對方一定會跟你說同一種語言.
對嗎.
如果政權只會說language of violence.
在政權底下的人.
同樣也會說同一種語言.
因為似乎這種語言.

$^{321}$是唯一一種大家能夠明白的語言.
你就捲入了一個以暴易暴的情況底下.
不是以巴是這樣.
很多世界各地不同的地方其實都是一樣.
你仔細想想.
很多地方很多社會很多的政局.
其實是否被綁在一個死局裡面.
上個星期我去了多倫多.
度假了幾天.
昨晚才回來.
上個星期去了多倫多.
正在吃粥.
那裡有電視.
我先看電視.
在說香港的情況.
就說年青人佔據了立法會.
然後做了很多很暴力的事情.
看完那個報道之後.
為什麼沒有提到那個政權拖延.
或者政權無能呢.
為什麼沒有提到政權霸道等等方面呢.
因為制度也可以很暴力的.
這種以暴易暴的情況.
為什麼沒有提到.
原來右下角寫TVB.
怪不得.
如果大家只不過是在一種暴力的情勢底下.
說不對就打的.
沒有對話的.
沒有大家可以謀求的understanding.
這個一定是惡性循環.
以巴是這樣.
香港也是走向這個方向.
這就是C編在73篇所說.
狂傲的人.
去打壓下面的人.
你不要只說普通的政局.
教會也很多時候被打壓.
特別現在在中國大陸.
不需要留意很多新聞我們都明白.

$^{361}$最近我跟大家提過.
有一個牧師在中國大陸.
最近被指控.
說他有顛覆國家的罪.
甚至捉了教會裡面的其他信徒.
餵他們吃一些不知道是什麼的藥.
然後逼他們下一些假的證供.
然後指證那個牧師.
指證那個教會的領袖.
來將他們入罪.
或者有女的信徒被性侵犯等等.
這些是一些很不公義的事情.
這就是世情.
C編73篇所留意到的世情.
我想我提過的.
我認識一個講國語的傳統人.
他的親戚在新疆也做牧師.
不知道有沒有跟大家提過.
我記得去年一起吃飯的時候.
吃完飯我們一起祈禱.
他一個祈禱的事項是.
他的親戚在新疆不見了.
我問他不見了多久.
不見了九個月.
有沒有他的音訊.
沒有音訊.
WeChat也關了.
每樣東西都沒有了.
我問這位傳統人.
我問你親戚怎麼樣.
一直圍著他祈禱.
他就放了出來.
九個月十個月.
受了很多苦.
被人抓進了所謂的再教育營.
新疆裡面的.
他受了很多苦.
如果你有機會去看BBC.
去採訪新疆裡面的問題.
新疆裡面的再教育營.

$^{401}$真的看到你無骨悚然.
問題不是那些人.
那些人為了什麼.
他們不可以說什麼.
問題是他們不可以說什麼.
不可以說什麼.
這就是詩篇.
七十三篇所提到的情況.
往往西程就不是這樣的.
很多的詩篇.
很多的聖經.
很多時候我們要很斟酌翻譯.
那些字.
到底現代人怎麼用中文英文.
等等現代語言.
應該怎麼翻譯它.
怎麼掌握到那個意義.
詩篇七十三篇大概都不用的.
怎麼翻譯都好.
準不準都好.
大家都很明白.
這個世界.
很多時候就是這樣.
到第二段.
十三節.
詩人見到這樣的世情的時候.
他說我緊守.
我的心.
順存.
實在是徒然.
我洗手.
表明清白也是枉然.
因為我終日受傷害.
連他自己都受苦被欺壓.
每天早晨受懲罰.
如果我心裡面說.
我要說這樣的話.
我就是對你的這一代.
眾兒女神的兒女不重要.
我思想這件事的時候.

$^{441}$就覺得很煩惱.
英文翻譯.
oppressive.
不單止政權.
不單止那些狂妄的人.
令他很oppressive.
連他掙扎這個問題.
他都覺得很oppressive.
覺得很累.
詩人說.
他見到很多不公義的事.
甚至令到他.
甚至會失腳.
甚至會離開神殿.
整首詩篇由第一節到現在.
走下坡.
走下坡的意思就是說.
詩人見到世情.
社會裡面發生的事.
狂妄的人.
張牙舞爪.
然後他心懷不平.
他心裡面覺得很不舒服.
甚至到現在.
他說.
一切都是陶冶.
一切.
都是empty.
守住上帝的約.
原來是陶冶的.
他去到這個地步.
最低最低的地步.
我們人的傾向都是這樣.
We define God.
by our problems.
我們見到的問題.
我們會將上帝定義.
以我們見到的問題.
來定義上帝.
很少會相反.

$^{481}$讓上帝來定義我們的問題.
我們見到的問題.
我們會將上帝定義.
以我們見到的問題.
來定義我們的問題.
詩人似乎是想.
提出一個這樣的角度.
眼見惡人享福.
很多時候我們會定義上帝.
with our problems.
就是說.
見到這麼多不公義的事.
一是神不公義.
一是上帝無能.
這是很典型的.
problem of evil.
就是苦罪的原委.
底線就是.
如果上帝不阻止這些事.
上帝就無能.
要麼上帝能夠阻止這些事.
但是他本性不是善的.
要麼上帝能夠阻止這些事.
但是他本性不是善的.
所以他容許一些不義的事來發生.
最後他掙扎的是什麼呢.
就是說如果我這樣說話的話.
我就會對不起.
你的兒女.
下一代的兒女.
下一代的教會會是怎樣呢.
如果我繼續這樣想下去.
但是這樣想下去的時候.
又很oppressive.
又很辛苦.
怎樣去守住約.
但是又見到.
世情是這麼不公義.
到第十七字.
直到.

$^{521}$我進了神的聖所.
才明白他們的結果.
他一直感覺到很oppressive.
直到他來到.
聖殿裡面去敬拜.
來到教會.
或者聽到一些的詩歌.
他才想起有.
另外一個的角度.
有些什麼他看到.
或者他進了上帝的聖所.
他明白到的.
會不會像.
先知以在亞.
去到.
聖殿敬拜的時候.
他看到高高在上有上帝.
坐在寶座上面.
所以他明白原來上帝仍然是.
掌權.
還是像以西結去被擄的時候.
流亡的時候.
他在河邊祈禱的時候.
他看到上帝的戰車.
有四活物.
有火有雷電的戰車.
上帝的座駕原來去到他們的那裡.
不是被困於聖殿.
我不知道他看到什麼.
但他明白了什麼.
他明白到有幾樣東西.
第一.
他在敬拜的裡面.
他明白到.
惡人的結局.
惡人將來.
會有什麼結局.
不是他現在看到的結局.
而是他明白將來.
會有一個結局.

$^{561}$第十八至第二十節.
他說你上帝.
實在將這些惡人放在.
滑地.
使他們倒下滅亡.
他們忽然之間.
就會成為多麼荒謬.
被突然的驚鴻.
完全消滅.
人睡醒了.
怎樣看夢.
主日你睡醒也照樣.
興翰他們如幻形.
他看到惡人的結局.
我覺得這個幻形.
這個我翻譯的.
我跟ESV翻譯的.
ESV翻譯叫做Phantoms.
即是說這些惡人.
他們多囂張也好.
終有一天他們會滅亡.
終有一天.
好像這些Phantoms一樣.
這些鬼影一樣.
一過了.
就不再有.
他在聖所裡面了解這件事.
就是上帝.
現在不施行審判.
不等於他永遠不施行審判.
就是上帝將來.
一定會施行審判.
上個星期度假的時候.
我們住的那間屋.
我們就有TV.
那間大屋.
Airbnb.
我們有時候看電影.
我記得看一套.
就是Mission Impossible.

$^{601}$Fallout.
即是最近那套.
劇情當然很緊張.
但因為我看過這套電影.
所以我看到中間的時候.
那些令人很害怕的東西.
所謂的.
劇情很緊湊.
當我明白結局是怎樣的時候.
令我看這一段.
是完全不同的.
私人似乎就是在說這件事.
就是看一些.
譬如長篇劇.
其中一集.
有一個人.
其中一集.
那個惡人是令你很不舒服的.
但你留意.
你只是在看其中一集.
大結局還未到.
私人就是說.
現在很生氣.
但他知道編劇是誰.
將來會有大結局.
是直至到他入聖所的時候.
他才看到.
就是這位永恆的主.
原來一定會有.
將來私人審判.
所以他了解到惡人的結局.
並且.
他了解到自己是無知的.
二十一節.
我心中酸苦我的肺腑.
刺痛的時候.
我是愚昧無知的.
我在你面前.
就好像畜類一樣.
畜類一樣.

$^{641}$他自知什麼呢.
就自己.
自知自己是無知的.
自知什麼呢.
就自己好像畜生一樣.
英文有個翻譯.
Dumbass.
我就是Dumbass.
我在酸苦的當中.
原來我是一個.
Dumbass.
我是一個畜類一樣.
大家知道.
井底之蛙.
在井底.
看上去的時候.
你只會看到一個很狹窄的視角.
你的視野.
是非常狹窄的.
你會看到.
一些狂傲的人.
他們是在欺壓.
日光之下的平民百姓.
但只不過你在井底看.
在進到聖所的時候.
他忽然好像.
從井底再上來.
他看到有永恆.
有永恆.
如果只有井底的話.
這個世界就沒有公義.
一切都是徒然.
守約也是徒然.
走得正站得正也是徒然.
但施人說.
他從井底上來的時候.
他就看到永恆的角度.
永恆的視角.
讓他看到上帝原來是有.
公義的.

$^{681}$將來一定會有結果.
現在他就知道.
我是無知的.
因為我只不過是.
such a dumb ass.
我只是一個井底.
只能看到的東西.
是好笑的.
23節.
但是我仍與你同在.
你緊握著我的右手.
你以你的奮然.
引領我.
以後還要接我到榮耀裡去.
除你以外在天上.
我還有誰.
除你以外在地上.
我又沒所愛慕的.
我的肉身.
我的內心雖然漸漸衰弱.
神卻永遠.
是我心裡的盤石.
是我的業份.
看那遠離你的.
必定滅亡.
凡是對你不徵的.
你都要滅絕.
對我來說.
親近神是好的.
我以主耶和華.
為我的避難所.
我要述說.
你一切的行為.
整篇詩篇.
你畫個V字.
你就明白整篇詩篇想說什麼.
他由開始.
他看到狂傲人.
然後他心懷不平.
覺得每樣東西都是陶然.

$^{721}$他非常酸苦的時候.
直到他.
進了上帝的聖所.
他就明白惡人的結局.
明白自己無知.
明白上帝的同在.
他又再回到上帝.
28節最後.
他又用同樣的字.
第一節就是這樣說.
他說神善待以色列.
God is good.
28節.
他說親近上帝是美好的.
Good.
回到上帝.
是一個信心的階梯.
下樓梯.
然後再上樓梯.
整個的轉捩點.
那個轉折在哪裡.
就是他肯來到上帝的聖所.
很多時候我們.
過我們人生.
我們可以很酸苦的.
我們可以很苦澀.
但我們又不願意.
讓上帝來對我們說話.
但我們又不願意.
讓上帝來對我們說話.
我們只有.
C邊73邊的頭那一段.
總之我們就很酸苦.
每樣東西都覺得很徒然.
沒有的,什麼都沒有.
詩人說.
不要做井底子話.
進到上帝的聖所.
讓上帝.
給我們一個新的角度.

$^{761}$舊約裡面的聖殿官.
Temple.
是有兩方面的.
第一,上帝的Temple.
在耶路撒冷的Temple.
就是一個Copy.
Copy of what?.
Copy of Heaven.
就是天上的寶座.
給一個Copy給你.
讓你在地下能夠看到.
原來上帝在天上仍然掌權.
所以他進到上帝的聖所的時候.
他就明白.
原來上帝仍然在天上掌權.
另一方面.
聖殿的Temple.
也是一個路標.
你開車的時候.
你會看到.
How many miles away.
就有Exit.
或者目的地是.
多少里以外.
那個Temple.
就是那個路標.
就是告訴你將來.
原來天國會降臨.
你從舊約的整個歷史.
從阿巴拉罕的時候.
在祭壇.
到摩西大連的意識人.
出埃及.
在西大山底下.
他們建造會幕.
是一個portable的Temple.
帶著他們在曠野裡面.
牽著他們的右手一路走.
到他們入主了迦南.
就在耶路撒冷建一個聖殿.

$^{801}$一個石造的聖殿.
是一個所謂Permanent.
的標記.
然後到了耶穌的時候.
這個聖殿會過去.
真正的聖殿在哪裡.
在我這裡.
在我死而復活.
一個活著的主.
一個信仰的裡面.
在福音裡面.
所以現在不需要一定要去地圖的那一點.
去敬拜神.
在基督裡面.
不再在地圖的某一點裡面.
某一個Building裡面.
而在基督的信仰裡面.
基督的福音裡面.
天王並且將來.
也會有新天新地.
聖經沒有用一些很哲學的.
抽象的.
1234第一點第二點第三點.
用Logic來解決這個.
Problem of evil.
聖經不用這個方法.
聖經用一個故事的方法.
整個大歷史的敘事.
就應該是這樣.
基督已經第一次來了.
並且他會再回來.
我們現在就在中間.
原來這些Evil.
這些惡事.
狂傲的人.
有一天將會有結果.
上帝在天上仍然是掌權.
你就要明白.
你信耶穌.
不是要上帝Fit在你的故事裡.

$^{841}$我人生裡面在搞什麼.
上帝一定要幫我.
而是相反的.
上帝有一個大的故事.
是一個Narrative.
是我怎樣Fit在那個Narrative.
我已經被上帝所愛.
得到上帝的恩典.
我自己現在的人生.
我怎樣改寫.
應該Fit在上帝這個大故事裡.
可以和上帝同工.
上帝是公義的.
我仍然為公義說話.
上帝是真理的我仍然傳真理.
上帝是一個有盼望的神.
我因為基於耶穌基督死而復活.
這個盼望我能夠活著.
是我怎樣Fit在上帝的故事裡.
不是上帝怎樣Fit在我的故事裡.
我們結束的時候會守聖餐.
我不知道私人.
他說進到聖所看到什麼.
可能是一句詩歌.
給他一個新的體會.
可能是.
其他.
兄弟姐妹對他說的話.
給他一些鼓勵.
或是可能他宣佈的時候.
他來到上帝面前祈禱.
他發現上帝仍然聽他禱告.
可能他拿起一些.
實物.
一些標記.
他了解原來天上上帝仍然掌權.
並且一個大的.
Narrative告訴他.
原來我現在正在看的.
只不過很多集裡.

$^{881}$長篇劇集裡其中一集.
結局是上帝掌管.
可能他看到聖所裡面的一些雕刻.
可能看到祭壇.
他可能聞到.
上帝.
給上帝的祭物.
那些味道.
可能看到.
聖所裡面事奉的祭司.
等等.
讓他了解到原來.
上帝仍然掌權.
今天我們來到教會也是一樣.
可能你聞到是咖啡味.
可能你聽到的是.
弟兄姊妹講的話.
可能你.
一個人的祈禱.
或者我們聖餐.
拿起這個杯子.
你要知道原來上帝仍然掌權.
我進到上帝的聖所.
我先了解到.
上帝在做什麼.
然後上帝在我的生命裡面講說話.
多謝.
\newpage



\section{}
\label{sec:1lBwTIY05FM}
\textbf{當被壓太重時...}
\newline
\newline
連結: \href{https://youtube.com/watch?v=1lBwTIY05FM}{\texttt{https://youtube.com/watch?v=1lBwTIY05FM}} ~~~~ 語音日期: 2019-07-13
\newline
\newline
\hyperref[sec:ip7wGoUuLe8]{\small{< < < PREV SERMON < < <}}
~
\hyperref[sec:index_chronic]{\small{[返順時目]}}
~
\hyperref[sec:index_scriptual]{\small{[返順卷目]}}
~
\hyperref[sec:Qfkef2FSGUk]{\small{> > > NEXT SERMON > > >}}
\newline
\newline
$^{1}$(悲傷的音樂).
如果我說.
沒有胃口,頭痛,胃痛,失眠.
你會聯想什麼呢?.
或者你想都不用想.
馬上就有答案了.
就是生活壓力.
或者你正在這個情況裡面.
實在太無奈了.
你和我在生活上.
難免承受不同的壓力.
來自家庭,工作,疾病,經濟.
人際關係.
有時壓力還一下子四方八面地來.
我們真是叫苦連天.
沒錯,生活壓力太沉重.
是可以弄壞我們的.
我們會做出一些.
連自己都很不喜歡的事.
或者是對別人說話很不客氣.
或者是為了一些小事發很大的脾氣.
又或者是心情很憂鬱.
好像跌落深谷一樣.
自然地壓力.
是會使我們的眼目離開了神.
單單注目在自己的問題上.
不如我們稍微停一下.
聽一下保羅的一句話.
在哥林德號書第一章那裡他說.
「我們被壓大重,力不能勝.
叫我們不靠自己,只靠神」.
我估計我們大部分人.
都沒有像保羅那麼淒涼.
因為他曾經試過.
坐著的船都要快下沉.
他又試過被鞭打差不多要死.
又試過被一群反對他的人追殺.
甚至要躲在籃子裡.
深夜被人從城牆那裡跳下來逃命.
以上這些情況.

$^{41}$都是保羅所說的被壓太重.
力不能勝.
也是他為福音所承受的.
非常沉重的壓力.
但是他在這些壓力裡得到的秘訣就是.
「叫我不靠自己,只靠神」.
即使在面對重重的壓力裡.
我們必須要知道.
如果我們不小心處理的話.
它會傷害我們和人的關係.
而且後果可大可小.
而在承受沉重的壓力裡.
讓我們都像保羅一樣.
「不靠自己,只靠神」.
我們可以這樣祈禱.
「主啊,求你幫助我們.
我很需要從你而來的智慧和力量」.
而當我們看到情況不受我們控制的時候.
我們更加要學習堅定相信神.
我們要數算祂過去在我們身上的恩典.
我們要信得過.
在當前這些看起來好像很難.
很失敗的情況裡.
神會繼續賜下祂的恩典.
使我們得到勝利.
如果我們這樣做.
我們不僅體會到從神而來的智慧和力量.
我們也會體會到.
原來在壓力裡.
我們仍然可以做一個.
冷靜面對令人喜愛的人.
「我愛你」 詞:陳曦 曲:陳曦.
「我愛你」 曲:陳曦.
\newpage



\section{}
\label{sec:Qfkef2FSGUk}
\textbf{3/17/19 主日講道:智取人生 (二) - 陳崇基牧師}
\newline
\newline
連結: \href{https://youtube.com/watch?v=Qfkef2FSGUk}{\texttt{https://youtube.com/watch?v=Qfkef2FSGUk}} ~~~~ 語音日期: 2019-07-14
\newline
\newline
\hyperref[sec:1lBwTIY05FM]{\small{< < < PREV SERMON < < <}}
~
\hyperref[sec:index_chronic]{\small{[返順時目]}}
~
\hyperref[sec:index_scriptual]{\small{[返順卷目]}}
~
\hyperref[sec:qc2nb_Nb9qc]{\small{> > > NEXT SERMON > > >}}
\newline
\newline
$^{1}$.
\newpage



\section{}
\label{sec:qc2nb_Nb9qc}
\textbf{8/19/18 主日講道: 基督徒的教會觀(吳君才牧師)}
\newline
\newline
連結: \href{https://youtube.com/watch?v=qc2nb-Nb9qc}{\texttt{https://youtube.com/watch?v=qc2nb-Nb9qc}} ~~~~ 語音日期: 2019-07-14
\newline
\newline
\hyperref[sec:Qfkef2FSGUk]{\small{< < < PREV SERMON < < <}}
~
\hyperref[sec:index_chronic]{\small{[返順時目]}}
~
\hyperref[sec:index_scriptual]{\small{[返順卷目]}}
~
\hyperref[sec:TEi_5oJPKrM]{\small{> > > NEXT SERMON > > >}}
\newline
\newline
$^{1}$.
恭喜七位今天受浸的弟兄姊妹.
時間有限.
原本三十幾分鐘的講道.
現在剩下十分鐘都不夠.
我很快看一輪.
我的powerpoint就差不多了.
我想提醒今天受洗的弟兄姊妹.
和在座的基督徒.
我們信了耶穌之後.
我們應當有什麼表現.
我用耶福所書四章一到六節.
跟大家這樣說.
第一件事.
這裡就是經文.
你拿一次.
看到了.
看完,英文.
看完,好.
四到六節中文.
講到合一.
有七個要素.
看英文.
接著.
看完就很快不記得.
所以我又提兩件事.
今天你一定要記得.
我們要實行相稱的生活.
第二件事.
我們需要有一個合一的心智.
受洗的七位弟兄姊妹.
更加要明白多一點.
好,看回頭.
第一件事,我們要有相稱的生活.
相稱,你看到背景.
好像一把秤.
一面倒.
有些人信耶穌.
但和他的表現不相稱.
相稱原文的意思.

$^{41}$就是配合.
合而相配.
有些譯去做.
就是應該.
一個信了耶穌的人.
表現不太好.
就真的不太好.
其中有兩樣質素.
我們說要有相稱的生活.
第一樣,我們要追求屬靈的美德.
謙卑.
溫柔,忍耐.
謙卑.
其實.
是不容易學的功課.
教父古斯丁說.
人生的靈情.
有三樣東西.
第一樣是學謙卑.
第二樣又是學謙卑.
第三樣.
仍然是學謙卑.
這裡說凡事.
謙虛,溫柔,忍耐.
謙虛是第一樣東西.
我們要學.
有些聖經學者.
教大家謙卑.
講到一個情況.
我不自誇.
我看別人比自己強.
比自己厲害.
我仍然能夠謙卑.
這是一個不容易的功課.
因為你覺得.
有一個人.
比我厲害.
我仍然能夠和他和諧相處.
如果我們有時合作上.
跟別人出現問題.

$^{81}$你要看自己是否謙卑.
第二樣是溫柔.
溫柔是耶穌的榜樣.
耶穌說:我心裡柔和謙卑.
柔和和謙卑一起.
有人用溫柔.
柔和.
很比海綿.
吸吸吸.
但你怎樣屈他都不會.
他可以承受壓力.
今天我們說.
我心裡柔和謙卑.
他的謙卑.
不是怕死.
而是真的能夠.
好好面對一些處境.
另外一件事.
我們說忍耐.
忍耐在這裡專注.
特別有些人無理.
對你一些不好的.
你怎樣回應.
忍耐的意思.
就是對一些不好的東西.
能夠自我約束.
而不會再有衝突.
這是靠主.
靠主才可以.
一個夢照的人.
你會發現.
我們都需要經歷這些功課.
第二件事.
我們說有相情的生活.
第二張片.
要重視.
相交的關係.
這裡說到用愛心互相寬容.
用和平彼此聯絡.
你見到有互相有彼此.

$^{121}$是相情的.
我們基督教信仰是群體的信仰.
不是自我的信仰.
這麼簡單.
我們是有關係.
有連結的.
不是疏離的.
如果一個人回到教會.
好像屬靈的獨行俠.
或者屬靈的吉督賽人.
這是一個畸形的表現.
我們需要靠主.
留意.
我們是用愛心互相寬容.
用和平彼此聯絡.
寬容.
是用愛心.
你用愛心有多少.
寬容又有多大.
我們有時會覺得.
信耶穌教的人一定很好.
但剛才有見證.
有些人比不信的人更壞.
但你的信仰不是建基於人.
而是建基於神.
還有一件事.
聖經沒有說.
信了耶穌的人不會有爭執.
不會有意見磨擦.
不是.
一個寬容.
就是寬大的包容.
有一個饒恕的心.
你就可以消除和諧和衝突.
寬容並不表示.
我需要別人的寬容.
我去寬容別人.
但其實你自己也需要別人的寬容.
因為很多人已經包容了你.
寬容的意思是寬大的包容.

$^{161}$接納.
和接納.
接納.
和.
能夠.
在一些事上用愛心.
來代替你.
求主幫助我們也可以這樣做.
如果有一個人對你不好.
弟兄姊妹.
令你很難受.
你真的要靠主去面對.
怎樣用愛心寬容.
用和平彼此的聯絡.
不單這樣.
和平彼此聯絡.
和平的意思.
是一個很和藹.
和善溫文.
和睦友好的態度.
在教會裡我們需要.
求主有這樣的表現.
我們亦需要.
互相守望.
互相扶持.
如果教會大家一盤散沙.
很容易魔鬼.
從中在這裡.
搞擾.
我們很難合一.
再看下去.
我們要持守的是.
合一的心智.
有七件事.
我們需要理解.
有時很多人.
回到教會.
若當年看到教會一些令人心up 的事.
教會分裂.
或者有些事.

$^{201}$令到你很難接受.
要離開教會的事件.
這些就很難.
教會的合一.
是否只靠制度或甚麼合一呢?.
其實不是.
是生命的合一.
你信了耶穌之後.
你已經有神的生命.
有聖命在這裡.
有基督拆毀所有的隔閡.
我們在主裡合一.
最近Villow Creek社區教會.
柳溪社區教會.
那個標.
海波斯牧師.
因為性騷擾的事件.
擾攘多時.
不單止海波斯牧師.
要提早離職.
連兩位原本想要.
找他做繼任人的.
都請辭.
所有原有的執事.
全部下台.
新的執事.
董事再上任.
這間教會.
受到這些事件的衝擊.
他們依然在這裡.
好像聖經說.
在這裡.
好像聖經說.
是竭力.
盡力保守自己的合一.
要重建.
仍然為了整體.
教會的合一和重建.
當然他們的前途.
還是很艱難.

$^{241}$已漫長.
但他們是竭力盡力保守教會合一.
所以我們.
在成會的第一節.
我們都有這七件事.
今天你的認訊當中.
有這七件事.
你也可以想一想.
你也可以拿這七件事.
作為一個指標.
因為異端.
邪說.
邪教.
不會有這七件事.
一生是指一個身體.
那裡是指耶穌基督的身體.
今天我們看到.
教會是基督的身體.
只是一個.
你已經加入了基督名下.
所以我們人人尊主偉大.
高舉基督.
人人在一生.
在基督底下好好配合.
我們就可以合一.
有些人說.
一生另一個涵義.
就是你.
堅心一致的.
就是毀身在這間教會.
不輕易離開.
好好地.
用你的恩賜才覺.
在教會生活.
這七位今天受浸的弟兄姊妹.
希望你們也是這樣.
第二是一靈.
一個靈是指聖靈.
你和我信主之後有聖靈.
貫通我們各人.

$^{281}$一個人如果沒有聖靈.
他是不是基督徒.
這樣都值得疑問.
但是有一個說法.
這個說法.
一靈也可以指靈性的生活.
你有沒有靈性的生活呢.
教會有聖靈.
作供就會生機蓬勃.
有生命的活力.
為什麼信徒失去靈明呢.
搞得好像薩迪教會一樣.
在啟示錄的薩迪教會.
神責備他.
安你一命是活的.
其實是死的.
因為你沒有靈性.
你沒有靈性.
神責備他.
安你一命是活的.
其實是死的.
我們千萬不要走這條路.
我們需要靠主恩典.
恩亡是盼望.
這個很豐富.
包括永生的盼望.
包括主的再來.
包括天堂的基業.
將來的復活.
甚至新天新地.
很榮耀的賞賜.
所以我們很有盼望.
你仍然抓緊神有盼望.
第四.
是一神.
一主.
一主就是救主.
是耶穌基督.
也是我們人生生命萬有的主.
是我們的主宰.

$^{321}$今天我們信耶穌.
很多人是接受耶穌做救主.
但是不肯讓耶穌.
做一生的主宰.
將耶穌推到一邊.
這是值得我們反省.
我們信這個信.
是信神.
信耶穌基督.
除去諸外沒有別的救恩.
我們這個信.
是因信稱義.
其他的邪端.
以教邪說.
未必是這樣.
現在有全能神的教.
其實是東方閃電教.
另外一個女的基督.
諸如此類的邪說.
他們不是這樣的信.
第六是一洗.
剛才是用浸泥.
有些教會是用灑水泥.
無論什麼形式.
最重要是你.
像陳牧師提醒.
你向罪死.
被埋葬.
然後向神活.
有些人信了耶穌.
受了洗之後.
受了浸之後.
生命依舊.
不是洗泥浸泥有問題.
所以我們都要求主幫助.
最後一樣是一神.
基督教是一神觀.
除他以外.
不可有別的神.
可惜很多時候.

$^{361}$我們有很多別的神.
別的偶像.
別的一些東西.
纏累我們.
我們被這些分心分神.
在這裡斯圖波羅.
特別提到這位神.
是超乎眾人之上.
觀乎眾人之中.
也住在眾人之內.
超乎就是神的權能.
智慧,威嚴.
超乎一切.
值得你和我相信和信服.
說到他是觀乎.
是意思是.
貫穿,通過.
古往今來的人.
說到在眾人之內.
其實主神的同在.
今天神和我們同在.
特別以今天.
一浸的弟兄姊妹同在.
我們經歷神的同在.
就不怕去面對明天.
我們低頭結束祈禱.
求主你自己話語.
提醒我們.
我們有七個要素.
七個重要的.
信仰,心智,持守.
可以令到教會合一.
亦都讓我們生命與相稱.
的生活.
信仰,相稱.
活出你的樣式.
以幫助人道,奉主耶穌基督明卡.阿們..
\newpage



\section{}
\label{sec:TEi_5oJPKrM}
\textbf{11/25/18 主日講道 真誠親密的感恩 (張麗仙傳道)}
\newline
\newline
連結: \href{https://youtube.com/watch?v=TEi-5oJPKrM}{\texttt{https://youtube.com/watch?v=TEi-5oJPKrM}} ~~~~ 語音日期: 2019-07-16
\newline
\newline
\hyperref[sec:qc2nb_Nb9qc]{\small{< < < PREV SERMON < < <}}
~
\hyperref[sec:index_chronic]{\small{[返順時目]}}
~
\hyperref[sec:index_scriptual]{\small{[返順卷目]}}
~
\hyperref[sec:NpmOHXaC8DQ]{\small{> > > NEXT SERMON > > >}}
\newline
\newline
$^{1}$感恩節剛過去了,我想問你,回想起來,哪一餐吃得最好吃呢?.
我想問你,吃了這麼多,歡聚了這麼多,你有沒有把握機會表達衷心的感恩呢?.
我知道的是對神和對人,如果有的話,這個感恩節就過得有意思了.
哎呀,我又想問你,你有沒有把握機會表達衷心的感恩呢?.
有些人在Facebook上看我發過這句話,.
不是開心的人才會感恩,而是感恩的人才會開心.
不過,說到感恩,有時不是不想,但很多人都不知道怎樣去表達感恩.
所以感恩節就是吃一大餐,是好事,不要讓它停.
不如我們問一下,我們感恩的人,像我們的父母,配偶,親愛的人.
這些感恩大餐,不能缺乏的一樣是什麼呢?.
我相信他們會告訴你,就是你真心在這裡,跟他們一起過一個親密的時候.
真的,如果人在,心不在,比如坐在一起吃飯,吃感恩大餐.
你一邊在這裡點著手機,一邊看電視,多豐富的感恩大餐,其實都沒有意思.
因為你還在自己的世界裡感恩,對方感到很沒趣就真的.
一切感恩大餐都是表面的,如果你的心不在,是沒有意思的.
真心在一起,對方是會感受到的,感受到你的珍惜和重視.
這個就是最好的感恩.
我們在教會裡時時施奉神,奉獻很多錢.
但是如果我們沒有真心想跟神在一起,這些行動恐怕都是空殼,是沒有什麼意思的.
經歷了聖經,很多次都說,神所體重的不是我們為祂做任何的事.
因為神其實是不欠缺的.
我們對神最好的感恩就是真心跟祂在一起,一種真誠親密的時刻.
你有沒有試過這樣做呢?你又懂不懂得這樣做呢?.
今天我就很想和你分享一下我在這方面的成長和體會.
不過在說之前,我很想用詩篇139篇的四幅圖畫提醒一下我們和神之間的關係是怎樣的.
第一幅畫是《你這麼愛我》.
看看聖經裡的畫,它說:神啊,你監察我,你認識我.
小心看,監察這個字其實是指神關注我們的一切.
不單只是表面的活動,更加是我們內心深處的感受,祂也監察,祂也知道.
接著說:你認識我.
認識這個字,傳授詩出現了四次.
一方面是強調神是無所不知的,祂深深知道我們的一切.
而這個字也強調神和他的納若子民有一個真誠親密的關係.
第二句詩篇說:我坐下,我起來,我行路,我躺臥,你都曉得,你都世察.
想想,我坐下,可以是為生活的事想到頭都爆了的時侯.
我起來,可以說我起身去接受一些我接受不到的挑戰時侯,我很害怕的時侯.
我行路,可以是我很迷失,不知道哪一條才是正路的時侯.
而躺臥更加是我擔憂時侯,點床點席,整晚睡不著的時侯.
這些時侯,神都曉得和世察.
而且詩篇第四節那裡,詩人更加說:連我沒有說出來的說話,神你都知道.

$^{41}$而且,你在前後環繞著我,按手在我身上.
這些動作都是表示神的愛和保護的動作.
神就是這麼愛我們,這是第一幅圖畫.
第二幅圖畫就是:我躲避你,多麼諷刺.
為何這麼愛我們的神,我們都躲避他呢?.
不過說真的,心裡都明白的,大概就好像父母對我們這樣.
我們怕了父母,經常告訴我們,這又不對,那又不對.
而我們越錯,就越不想見他們,我們對神就是這樣.
不過詩篇說:無論如何,你走到上天嗎?神都在那裡.
你走到地獄嗎?神都在那裡.
你展開青神翅膀,飛到海極嗎?意思是說,你由極東飛到極西嗎?神都在那裡.
我們問:為何我們要刻意地躲避神呢?.
或者我們是「wishful thinking」.
以為我們在這些黑暗的地方,神就看不到我們.
因為黑暗遮蔽我們,但小心我們.
「黑暗」這個字,其實是指「罪」.
如果我們被黑暗遮蔽我們,想神不見我們.
其實我們是非常錯,我們也看錯了.
因為詩篇第十節告訴我們.
原來就在那裡,你的手必引導我,你的右手必扶持我.
原來神對我們的追蹤,不是想審判我們,而是愛我們.
原來即使我們在罪的黑暗裡,神一直都愛我們.
我們再來看第三幅圖畫.
就是「你創造我」.
「你創造我」,我們終於可以了解為何神這麼愛我們.
因為祂是創造我們生命的神.
我們一起讀讀詩篇第十三至十八節.
那裡說:我的肺腑是你所造的,我在母腹中,你已覆蔽我.
我要稱謝你,因我受造奇妙可謂.
你的作為奇妙,這是我心心知道的.
我在暗中受造,在地的深處被聯絡.
那時我的形體並不向你隱藏.
我未成形的體質,你的眼早已看見了.
你所定的日子,我尚未到一日.
你都寫在你的冊上了.
神啊,你的意念向我何等寶貴,奇數何等眾多.
我若數點,比海沙更多.
這裡說:神創造我們.
你試過創造東西沒有?.
很有感情的.

$^{81}$譬如問問Mark Sir 竇木或Irene設計.
或你創造過某些東西,或你喜歡煮食.
你都知道這種創造的心血.
你明白到這種感情.
不過,說到底,我們的創造都是物質上的.
但神創造的是有血有肉有靈的生命,是我們.
這就是另一回事.
如果最靠近的,可以說是父母.
他將生命帶到這個世界上,是他的骨肉.
不過聖經又清楚說明.
神創造我們的生命.
父神比父母的愛更加深切.
因為父母是帶我們到這裡,神創造我們的生命.
所以當我們在罪裡面墮落滅亡的時候.
神會盡力拼命拯救我們.
甚至他猜他的愛子耶穌基督為我們的罪而死生.
在十字架上,他要挽救我們的生命.
這就表明了神對我們的愛.
難怪詩人大衛說.
神啊,你的意念向我何等寶貴,其數何等眾多.
因為你是創造我生命的神.
你的心思對我是多麼寶貴,是多麼多.
但意想不到的是下一幅圖畫.
因為第四幅畫是堅決受敵.
很奇怪,前面詩篇的內容是多麼溫馨的情節.
講到神多麼重視我們,創造我們.
為何讀到這裡突然殺出這些情節.
好像和上面的內容很不和諧.
所以我以前讀詩篇的時候.
很多時候有意無意地略過這段不讀,我不知你呢.
但對大衛來說,原來這個局面是很真實的.
因為當他寫這詩的時候,他正被仇敵攻擊誣陷他.
他知道這些仇敵是抵擋耶和華旨意的人.
他亦深信耶和華必定會勝過他們.
對我們來說,其實這段是有很重大的提醒.
我們和大衛都是一樣.
雖然我們和神有一個很美好的關係.
但我們的仇敵是想動搖我們對神的信心.
他時時想摧毀我們和神之間的關係.
記不記得上個月我和大家講到.

$^{121}$《耳窟所書》的時候,講到屬靈的戰爭.
一場屬靈的戰爭是真真實實.
到今時今日,現在,此時此刻.
都在打仗,撒旦要時時摧毀我們和神之間的關係.
難怪大衛會這樣說.
他說:恨惡你的,我豈不恨惡他們嗎?.
攻擊你的,我豈不憎厭他們嗎?.
弟兄姊妹,我們也要學好像大衛這樣.
有一個堅決仇敵的心.
最後能夠和大衛一樣.
講出以下兩句說話.
再一次認定我們和神之間寶貴的關係.
這兩句說話就是.
神啊,求你監察我.
知道我的心思.
思念我,知道我的意念.
看在我裡面,有什麼惡行沒有.
引導我走永生的道路.
在幾分鐘之前,我才和大家講到.
我們對神最好的感恩,就是真心和祂在一起.
我也和大家說過.
我很想分享一下我在這方面的成長和體會.
為什麼我又要花一大段時間讀詩篇139篇呢?.
是因為這詩篇一再提醒我們.
我們和神之間的關係原來如此寶貴.
祂愛我們,祂做我們.
甚至祂付出很大的代價.
去戰勝摧毀這個關係的仇敵.
那麼我們要怎樣珍惜我們和神的關係呢?.
我認為就是.
我們要時時和神有真誠親密的時刻.
就是說,我們時時說的是.
靈修禱告的時刻,即是靜止時間.
我們都知道,退離人群,單獨和神在一起.
我們禱告不難的,我們時時都做.
但是我們禱告要真誠親密.
就相當不容易了.
我們都體會到很多掙扎,主要有兩個.
第一就是,不知道要說什麼.
特別是面對這個世界,突然間亂槍掃射.

$^{161}$突然間山林大火,突然間地震海嘯.
很多無辜的人喪生.
我們祈禱真的不知道要說什麼.
再貼身一點.
我們看到身邊很愛的人.
他正在忍受癌症的折磨.
我們看到本來幸福的家庭.
關係破裂,不可以修補.
這個時候,我們祈禱更加不知道要說什麼.
第二個掙扎就是,我們精神不集中.
因為我們生活實在太多掛念.
或者是太多吸引.
其實我和你都是一樣,這些都是我的掙扎.
以致我來到神的面前.
很多時候祈禱,就像交功課一樣.
很難做到真誠親密.
不過,以下我就很想和你分享.
兩個改變了我和神同在的方法.
第一個方法就是.
馬丁路德的A Simple Way of Prayer.
其實年初講道的時候,我已經和大家分享過.
不過今天值得和大家再refresh一下這個方法.
是一個以神的話為禱告內容的方法.
我們平時都會讀一段聖經或者一節聖經.
作為我們的禱告開始.
可以說,例如我們的靈修三十天.
最緊的.
只不過我們時常懷著的是交課的心態.
三分鐘讀完或者聽完.
我們就閉上眼睛,祈禱,逐樣逐樣.
很快很快地說完,就阿們.
不過,其實我們的禱告可以好很多.
深入很多的.
又例如,我講回我們人際關係的靈修.
在9月14日的靈修裡面.
Mark Sir引用了《箴言》18章13節.
那裡說,未曾聽完就回答的.
便是他的愚昧和羞辱.
Mark Sir慢慢地和大家解釋清楚.
這節聖經是甚麼意思.

$^{201}$當你聽完的時候.
第一樣.
你開始祈禱的時候.
你可以反思.
你對神說,神啊,我讀這節聖經.
你讓我明白了.
我終於明白了甚麼甚麼了.
就例如,我終於明白了.
聆聽原來是這麼大的智慧.
不聽,其實就是很大的愚蠢.
這節聖經教了你甚麼.
你就跟神說,我學了,我明白了甚麼.
這是第一樣,反思.
即是說,你讀完這節聖經之後.
神啊,你讓我更加認識.
你是一個怎樣怎樣的神.
例如,神啊,你是這麼忍耐.
這麼願意聆聽我的神.
第三,就是認罪.
即是說,這節聖經幫助我.
審查到自己某方面的罪.
例如說,我讀了這節聖經之後.
我明白到原來我自己在人際溝通那裡.
就是事事犯了這個毛病.
我們不聽就說.
以致引起很多很多的矛盾衝突.
傷害.
在這個時候,你向神認罪.
第四.
你現在開始祈求.
你盡情為生活的事和身邊的人祈禱交托.
盡情說.
因為經過先前的三步.
我可以包裹你現在的禱告交托是不同了.
因為你是以神為主.
不再以自己為主.
沒錯了.
馬丁路德的這個Simple Way of Prayer.
其實很簡單.
是一個建基於神的話上.

$^{241}$可以包裹你這樣禱告的話.
你發覺你的禱告有深度了很多.
這個對我自己來說.
是有很深遠的影響.
因為它矯正了我.
不再只看渺小的自己.
而轉向看這位.
掌管我們今生和永遠的神.
我讓他的說話啟動我每天的生活.
以他為主.
我發覺禱告不會再是沒話說了.
不過.
如果說到精神不集中.
那又怎樣呢?.
禱告精神不集中.
可以說信了主很多年.
甚至我在教會做童工,導師.
我一直這個禱告精神不集中.
都是我的問題.
的確,人生很多吸引.
亦有很多掛慮.
所以我禱告很多時候就像游魂一樣.
直至幾年前.
我讀了一篇靈修的文章.
講到一個很好的方法.
可以說它影響到我和神的禱告.
影響我和神的關係.
直至到今時今日.
這個方法.
我給它一個名字.
五分鐘奇妙的起飛.
這是我真實的經歷.
只是五分鐘.
它是平凡而有效的.
就是我們在祈禱開始的時候.
什麼都放下.
特別是我們的手機.
這是當然的.
不過即使我們放下很多東西.
放下一切東西.

$^{281}$但我們的心思掛慮.
往往如潮湧來.
當我們祈禱的時候.
我們發覺我們沒辦法專注.
像有一樣東西湧過來.
在這個時候.
你只需要說一句話.
就是說.
神啊,你是創造我,救贖我.
引導我生命的神.
此刻,什麼都不重要.
除了你和我的關係是最重要.
或者你會馬上想到你的子女.
你要買樓.
或者你失業.
或者你有一些關係的衝突.
甚至是你想到面對癌症.
你都可以說.
此刻,什麼都不重要.
神啊,這一切掛慮都不重要.
唯獨是你和我的關係.
在這一刻是最重要.
或者你到了這個時候.
你又湧起很多很多的思念.
例如說今天生日派對還沒買東西.
或者有一個項目還沒完成.
但是你又要記住說.
神啊,你是創造我,救贖我.
和引導我的神.
此刻,什麼都不重要.
只有你和我的關係最重要.
開始的時候,你可能有點不習慣.
但是你這樣再做再做.
來回來回.
我可以保證,幾分鐘之內.
你的心靈就漸漸進入一個空間.
在那裡,你安靜,集中.
說個比喻就像飛機飛到跑道上.
慢慢熱身.
熱身到,就一下起飛.

$^{321}$這五分鐘,就是能夠讓我們專注熱身.
以致我們能夠一下起飛.
來到一個與神同在的空間.
在這裡,只有你和神親密的說話.
在這一刻.
你可以說出剛剛開始禱告的時候.
佔據著你腦海的每一樣掛慮.
因為你已經可以專注在神面前.
你就一樣一樣交託給神.
跟祂說你的感受.
不知不覺,你會發覺你說得很長.
不知不覺,你會發覺你和神的關係已經很深.
弟兄姊妹,特別是在傷痛,恐懼,沒智慧的日子.
你更加會感受到這個時刻的寶貴.
因為祂叫我們的信心堅固起來.
就是這麼簡單的一個方法.
我們只要願意付出這額外的五分鐘熱身.
幫助我們進入一個與神安靜,專注,同在的空間.
我們就能夠與神培養到想不到的親密關係.
我可以告訴你.
以後你會發覺,慢慢不用五分鐘這麼長.
甚至兩分鐘,甚至一分鐘.
你就可以進入這個專注的精神狀態裡面.
與神有一個空間,對祂傾心吐耳.
這個方法是我真實的體驗.
我能夠在這五分鐘期待的期間,影響我的禱告.
沒錯,就是這兩個方法.
第一個是馬丁路德的Simple Way of Prayer.
是以神的話語作為反思,讚美,認罪,祈求的禱告方法.
第二個就是五分鐘奇妙的起飛.
就是以五分鐘專注禱告的時間.
作為我們與神有的空間.
在那裡我們能夠對神傾心吐耳的禱告.
這兩個方法,我有時用這個,有時用那個.
有時兩個一起用.
可以說,這兩個方法是大大改變我與神同在的時間.
現在我的禱告不會再說我不知道說什麼.
又或者是精神不集中.
相反,我是能夠享受到真誠親密的與神同在.
我感到很開心.

$^{361}$我相信神也會很開心.
不過,也許你也會像我這樣馬上想到第三個問題.
就是,如果沒有時間怎麼辦呢?.
但是,我覺得我們可能都問錯了這個問題.
我們也許應該這樣問.
與這位那麼愛我們,創造我們.
甚至受苦,受死,救贖我們的永生神同在.
我們能不能找到時間呢?.
弟兄姊妹,如果你一向與神有一個真誠親密的時刻.
我為你感到高興.
但是如果你一直覺得自己與神的關係很游離.
我很鼓勵你開始嘗試.
也歡迎你跟我分享.
你嘗試之後的成功和失敗經驗.
無論如何,感恩節很快就過去.
但是感恩這件事是永遠都不會過去的.
特別是對愛我們的家人,好朋友.
更加特別的就是對於這位愛我們的神.
讓我們一起努力.
開墾一個單獨與神的真誠親密的空間.
我很相信,這是神最喜悅的感恩.
願神的話,祝福你.
(音量注意).
\newpage



\section{}
\label{sec:NpmOHXaC8DQ}
\textbf{11/4/18 主日講道:  女子組加油 \#2: 顛覆男人世界的聖經女子(陳崇基牧師)}
\newline
\newline
連結: \href{https://youtube.com/watch?v=NpmOHXaC8DQ}{\texttt{https://youtube.com/watch?v=NpmOHXaC8DQ}} ~~~~ 語音日期: 2019-07-16
\newline
\newline
\hyperref[sec:TEi_5oJPKrM]{\small{< < < PREV SERMON < < <}}
~
\hyperref[sec:index_chronic]{\small{[返順時目]}}
~
\hyperref[sec:index_scriptual]{\small{[返順卷目]}}
~
\hyperref[sec:YsIxBX5tTR0]{\small{> > > NEXT SERMON > > >}}
\newline
\newline
$^{1}$今天我們繼續一個女子組加油的系列.
我們上個星期第一次開始.
我們講到聖經裡面怎麼看待女性.
怎麼去 correct(正確).
或者改正我們怎麼去看待女性的方法.
聖經絕對是對我們二十一世紀的人.
很大的提醒.
雖然一本很古老的書.
但是我們今天就要講.
聖經裡面的女性.
我們今天就要講.
聖經裡面的女性.
我們今天就要講.
很大的提醒.
雖然是一本很古老的書.
但是是一本很劃時代的信息的書.
我們就 go into the Bible(去到聖經).
去看看有些什麼.
在聖經裡面怎麼看待女性.
聖經是不是責成男女平等呢?.
上次我告訴大家.
我先跟大家 present(展示)幾個命題.
這些命題.
大家談一談.
想一想.
到底是對還是錯呢?.
第一個命題就是.
聖經是父權.
就是重男輕女的社會裡面的產品.
這是對還是錯呢?.
是的.
因為以往古代的社會.
全部都是重男輕女.
其實現在也是.
你看到別人生育都是女的.
你覺得他很慘.
其實你也是重男輕女.
但是他不斷生育子女.
你又不會覺得他慘.
當然我不是這麼想.

$^{41}$我有一個女兒一個兒子.
(笑).
但是無論如何.
當時父權的社會就是這樣.
神要跟他們說話的時候.
就是 assume(假設).
或者假設預設他們當時有的價值觀或者語言.
不代表聖經是維護父權主義.
但是他會假設了.
有一些重男輕女的觀念.
所以第二個命題就是說.
聖經是不是挑戰父權主義的觀念呢?.
patriarchal values(父權價值觀).
對還是錯呢?.
是對的.
聖經裡面.
女人雖然比較少出現.
說話比較少一點.
但是每逢她們說話或者做事的時候.
很多時候都是物以罕為貴.
去挑戰當時男人的觀念.
或者當時一些傳統的觀念.
金庸死了是不是?.
有些人用女性主義的角度來看金庸小說.
你有沒有看過這樣的文章?.
那些 literary criticism(文學批評).
用這樣的角度來看金庸小說.
有些人也說.
金庸小說裡面的女性.
往往在關鍵的時候.
她們的勇氣和見識都比男人多一點.
想不想起有什麼主角?.
待會再聊好不好?.
所以聖經的女性也是一樣的.
出現的時候很多時候.
去顛倒或者去挑戰.
一些以男為上的觀念.
不是說神可以用男人做事.
更加可以用女性來做事.
在聖經的歷史裡面.

$^{81}$女性是扮演著策略性的角色的.
Yes or No?.
True or False?.
是對的.
是不是?.
所以女人不是可有可無的.
是扮演很多時候一些關鍵的大角色.
她們正在推進上帝自己的工作.
上一次我們留意到.
寫聖經的歷史.
或者留意聖經的歷史裡面.
一些關鍵時刻.
有什麼發生的時候.
很多時候我們自己的假設就是.
一定是男人.
其實你可以這樣寫的.
凡是一些關鍵的時候.
就是女人出現的時候.
譬如出埃及記.
為什麼會有摩西.
就是因為有些女人來保護他.
為什麼可以進入迦南.
以色列人可以進入迦南.
不是因為約書亞.
他有勇字在胸口.
不是因為他好勇鬥狠.
或者他戰術好.
而是一個愛邦的女子叫做拉合.
是當時的一個妓女.
她去投靠以色列一方來幫助他們.
等等等等.
直到耶穌基督來的時候.
雖然她選的是十二宮門徒.
但是生的當然是.
生耶穌基督出來的當然是媽媽.
誰不知道媽媽是女人嗎?.
當她復活的時候.
第一個來迎接復活主的時候.
又是一群婦女.
第一個將耶穌基督從死裡面復活.

$^{121}$這個信息傳給那群男人的.
也是一群弱質女子所為.
第四個命題.
今天這個命題是新的.
剛剛那些是 review(評論)而已.
解讀聖經的時候.
我們需要極度敏感的背景.
對不對?.
先打個問號.
我想今天給你看幾段聖經.
特別是針對女人所說的一些聖經.
有沒有當時的背景呢?.
如果明白了當時的背景的時候.
會不會給我們一些新的體見呢?.
路加福音第十章.
這是一個很熟悉的故事.
馬大瑪利亞.
我們看看這個故事.
也說到一個女子的故事.
他說她們走路的時候.
耶穌進入了一個村莊.
有一個女人叫做馬大Mother.
接她回自己的家裡面.
她有一個妹妹叫做瑪利亞.
坐在耶穌的腳前.
聽她講道.
馬大因為事後的事情多.
心裡面很忙亂.
她就進前來跟耶穌講.
她說主啊.
我妹妹留下我一個來事後.
你不在意嗎?.
請你吩咐她來幫助我.
耶穌就回答她.
馬大.
你為許多事情來思慮煩擾.
但是不可少的只有一件.
瑪利亞已經選上了上好的.
是沒有人能夠奪去的.
這是上好的福分.

$^{161}$傳統教會長大後.
我一直看聖經這個故事的時候.
都會覺得.
馬大就是一種忙碌型的信徒.
瑪利亞就是一種靜修靈修型的信徒.
耶穌就著重.
他要靈修.
要靜態.
不要太忙.
要安靜.
來操練自己.
無疑這是故事的其中一方面.
但是我有時候想.
第一世紀.
以前古代的.
他們聽到這個故事的時候.
知道這個情節的時候.
到底他們覺得最震撼的地方是什麼呢?.
有一些背景可以解釋一下.
背景是什麼呢?.
第一世紀的世界.
古代的世界.
我們很多華人都知道.
以前的傳統社會都是這樣的.
就是男女的社會空間.
social space(社會空間).
是畫得很清楚的.
Right?.
男主外.
是不是?.
他去見客人.
接待客人.
女主內.
他就是煮飯.
餵奶.
女主內.
男女之間的社會空間是河水不犯井水的.
其實今天甚至還是這樣的.
如果你有機會去耶路撒冷.
西牆.

$^{201}$那裡叫做曲牆.
現在叫西牆.
Western Wall.
那裡其實是一個露天的猶太人會堂.
synagogue(會堂).
也分開男女的.
中間有一個圍欄.
上面那一處.
左邊那一處.
就是男人的.
那裡是比較好的地方.
那裡還可以.
下雨的時候可以走進洞裡面.
那裡有一個比較漂亮的會堂.
女人在哪裡呢?.
女人在右邊.
右邊有很多白色的椅子.
西牆的廣場就是分為男女的.
很清晰的男女的空間.
這個其實也沒有什麼很奇怪的.
對我來說.
我人生裡面經常都會有女人跟我說.
你不要進來.
這裡是廚房.
你出去.
馬大當然是因為工作多而有怨氣.
但是我想馬大最不喜歡的是什麼呢?.
或者最詫異的是什麼呢?.
就是瑪麗亞在這個故事裡面.
她進入了一個男人的空間.
因為學習律法.
Torah(律法).
是猶太人男人的專利.
對馬大來說.
這是打破了當時社會的慣例.
他說他的妹妹叫做瑪麗亞.
她在耶穌的腳前坐下來.
聽他說道.
我們現代人經常想.
一個畫像.

$^{241}$就是那個女孩.
坐在一個很美好的老師的下面.
她托著頭.
然後很仰望他.
說這個真的是很厲害的高中老師.
又帥又什麼的.
我們想起的.
就是一種這樣的畫像.
但是問題是.
是不是真的有這個意思呢?.
坐在拉比的腳前.
其實是猶太人的專用語.
保羅有這麼用過.
保羅自己說什麼呢?.
他說我原本是猶太人.
在大數城裡面長大的.
在加瑪烈的門下.
我查過很多中文的翻譯本.
都是叫做門下.
都有它的意思.
因為門生的意思.
但其實原文是腳下的.
At the feet.
坐在加瑪烈的門下.
坐在拉比的門下是什麼意思呢?.
坐在拉比的腳下.
不是查經BSF這麼簡單.
不只是去看聖經.
不只是神學進修.
不只是拿神學學位.
坐在拉比的腳下.
就是說你要學東西.
然後將來你成為一個拉比.
將來你成為一個 teacher(老師).
一個有權威的律法.
的教師.
一個群體的領袖.
耶穌就說.
瑪利亞有權.
也有福分.

$^{281}$能夠成為拉比教師和 leader(領袖).
如果這樣看 social space(社會空間)的時候.
你就會發現.
原來仔細地看聖經.
你會覺得有些很有趣的地方.
譬如呢.
耶穌被釘十字架的時候.
那些門徒去了哪裡呢?.
藍門徒.
不是全躲起來了嗎?.
然後那些女人就像自出自入一樣.
對不對?.
為什麼會這樣呢?.
或者在空墳墓的時候.
有羅馬兵兵來守衛.
那時候教會被打壓.
跟隨耶穌的運動的人.
全部被打壓.
那些男的全部躲起來.
女的就自出自入.
就好像耶路撒冷當局不放他們在眼內一樣.
就因為男女的 social space(社會空間).
完全是不同的.
完全是河水不犯井水的.
有一個在中東宣教.
還有做了學者幾十年的.
叫做 Kenneth Bailey.
他呢.
幾年前他過世了.
他在黎巴嫩啊.
以色列等等的地方.
他做了幾十年宣教士.
他說這種情況.
男女分得很清楚.
現在還在中東發生.
譬如黎巴嫩有戰爭的時候.
最高峰的時候.
那些男人.
打仗的那些.
全部都躲在家裡.

$^{321}$或者出入的時候非常謹慎.
但是女人就很奇怪.
已經炸得飛起來了.
女人都是自出自入去買東西.
買菜煮飯.
帶著小孩出門.
這個表示.
我想這個背景是非常貼合.
福音書.
新約的背景.
為什麼被打壓的時候.
那些女人還可以自出自入.
走到十字架之前.
你還記得有一個叫耶穌所愛的門徒.
來到十字架之前.
耶穌說將母親交給你.
那個我們想應該是小約翰.
就是寫約翰福音的約翰.
我們猜到.
他那時候應該是比較年輕一點.
他跟那些婦女在一起.
這種社會的空間是畫得很清楚的.
如果我們明白的時候.
我們就看到原來很有趣.
但是到你呢.
看完福音書到你看史唐人傳的時候.
史唐人傳在司祭梵第一個被殉道的門徒之後.
教會就被大規模地打壓.
史唐人傳怎麼記載呢?.
九章第一節.
那個掃羅打壓教會的人叫掃羅.
他仍然向主的門徒口託威嚇封殺的話.
就去見大祭司.
求文書.
給大馬士革另外一個 town(鎮)的各個會堂.
若有信奉這道的人.
無論男女都捆綁帶回去.
然後受審.
唯一一個可能性.
為什麼無論男與女呢?.

$^{361}$唯一一個可能性就是當時初期教會.
他們的 call team member 不只是男人.
是有女人的.
所以女人也被抓了.
這個很清楚.
看到男與女都是同等.
能夠在一個屬靈的群體裡面.
教會的群體裡面.
來到去做領導.
被裝備了.
然後他們去帶領教會.
這是第一段經文.
第二段經文我想給你看的就是哥林多前書.
第十四章三十三節.
你可能會問.
為什麼有些地方又叫女人閉嘴呢?.
真的有啊.
十四章三十三節.
他說因為上帝不是叫人去混亂.
是叫人安靜.
這是保羅寫信給哥林多教會.
他們崇拜的時候.
有很多混亂的情況出現.
三十四節.
他說婦女在會中要閉口不言.
像在聖徒的宗教會一樣.
因為不准他們說話.
他們總要信服.
正如律法所說的.
這句我們男人聽到會覺得很爽.
這樣說真的很好.
他們要學什麼.
可以在家裡面問自己的丈夫.
因為婦女在會中而說話呢?.
原是可恥的.
首先第一件事.
這裡不是意思什麼.
意思一定不是.
有些人一聲不響.
一定不是的.

$^{401}$因為上文十一章的時候.
已經說到女人可以公道.
還可以講道.
但是如果你看社會空間.
社會的空間.
來看這段經文的時候.
就有些幫助.
今天的中東.
有些教會也是這樣的.
男女是分開坐的.
就是你進來的時候.
師事會告訴你.
你是男還是女.
你不可以選擇 other(其他).
那時候沒有這樣的東西.
男就坐這邊.
女就坐這邊.
牧師或者講道的人.
就是用所謂 classical Arabic.
就是那種古典的阿拉伯語.
來講道的.
男人在阿拉伯的社會裡面.
每個人都讀一下書.
每個人都會聽的.
但是女人呢.
就會講鄉下話.
但是就不會聽 classical Arabic.
剛剛我講過的學者叫做 Kenneth Bailey.
他這麼多年在中東裡面打滾.
去做宣教.
他就觀察到什麼呢?.
每當牧師用 classical Arabic 講道的時候.
男人就在聽.
女人又如何呢?.
因為他們聽不明白.
他們就在自言自語.
自言自語.
然後自己在講話.
直到牧師說.
女人請你安靜好不好?.

$^{441}$他們就靜了幾分鐘.
然後又開始自言自語.
又吵起來.
之後牧師或者就會說.
如果女人想講話.
想問話的話.
請你回去問你老公.
我覺得這麼多的解釋裡面.
這段經文.
是最切合當時的背景.
為什麼叫女人不要講話呢?.
因為講道的.
男人才聽得懂的語言.
或者 liturgical language.
其實跟今天大致上都相似的.
你不要覺得這些是很奇怪的東西.
你去戲院看戲.
有東西放出來叫你做什麼?.
叫你關掉手機.
叫你不要吵.
因為這樣會影響到其他人來聚會.
尊重其他人.
不就這麼簡單了嗎?.
其實教會也是一樣的.
你準時到.
是不是?.
小朋友.
It's okay.
你不放心放在那裡.
你可以帶他進來.
但是當他吵的時候.
你自然會帶他出去.
因為你不想影響其他人.
絕對保羅沒有針對女人.
或者看低她們.
不讓她們講話.
他關心的是什麼呢?.
他說上帝聚會的時候.
我們聚會的時候.
上帝不是叫我們混亂.

$^{481}$而是叫我們安靜.
其實叫做有 order(命令)的意思.
就是要井井有條.
因為很多時候聚會裡面.
有些微信主人進來的時候.
他會覺得.
為什麼這裡會這樣呢?.
好像街市一樣.
紙層層的女人在那裡.
保羅所關心的.
就是一個這樣的 situation(情況).
所以有背景的時候.
我們就明白多一點.
《多瑪書》第十六章.
這是最後一段聖經我想講的.
為什麼會想起這段聖經呢?.
因為今天我們開始.
我們教會的獎執.
舉行一次.
我們會選舉.
我終於用了大概五天時間.
追趕每一個被提名的人.
他們回答我.
到底他們會不會 run(舉行)這個 election(選舉).
有九位被提名出來的候選人.
出現了.
我們有六位女的.
有三位男的.
好不好?.
女子組加油!.
我還沒介紹他們之前.
我想告訴你.
我就想起這段聖經.
就是《羅馬書》最後一章.
我經常覺得《羅馬書》最後一章都是.
一些問安語.
最後好像寫信那樣.
sincerely.
regard(尊重).
nobody cares(沒有人在乎).

$^{521}$right?.
但是《羅馬書》第十六章不是這樣的.
《羅馬書》第十六章這裡.
保羅特別寫一封信收尾的時候.
就向羅馬教會.
或者對羅馬宣教有共識的二十六個人.
他提名來說出來.
其中有九位是女士.
受將第一節.
他說我對你們羅馬教會要舉薦我們的姐妹.
Phoebe.
她是堅格理教會的女執事.
他說她真的很堅.
堅格理教會的女執事.
請你為主來接待她.
合乎聖徒的體統.
她在何事上要你們幫助.
你們就幫助她.
因為她素來幫助很多人.
也幫助過我.
這是教會歷史裡面第一個女執事.
Phoebe.
也可能是帶著這封羅馬書.
這封信.
在堅格理的海港.
哥林多的海港.
搭船送到羅馬的宗教會.
所以是一個很縱容的人.
是英文的servant.
我們很多時候英文翻譯為servant.
其實是一個女的執事.
她說又問安了.
問誰安呢?.
伯基拉與阿居拉.
哪個男哪個女.
Priscilla and Aquila.
Priscilla.
就是伯基拉是一個女孩的名字.
他們在基督耶穌裡面與我同工.
也為我的命.

$^{561}$將自己的頸項自之度外.
伯基拉與阿居拉是兩夫婦.
這裡特別的地方就是.
女人的名字是先行的.
過了陳三陳太.
我們介紹這個couple的時候.
有四次都是說到陳太陳三.
陳太先到陳三.
這裡本來也用了一個專用語來形容他們.
或者形容Priscilla.
Fellow workers.
Co-workers.
是一個title.
保羅尊稱教會領袖.
來做co-workers.
或者同工.
十六章第六節.
她說我又問馬利亞安.
馬利亞當然就是女性的名字.
她為你們多受勞苦.
又問我的親屬.
這裡很少說.
原來保羅有親戚.
與我同坐牢的安多尼古和尤利亞安.
他們在使徒裡面有名望.
也使比我先在基督裡面.
這裡就是陳三陳太.
因為陳太太是尤利亞安.
這裡說到他們是一個couple.
男與女.
是跟十二使徒來到同列的.
所以你這樣看的時候.
他們其實稱這個人叫做女使徒.
按照他們的名字.
當時人的名字不是很多變化的.
不是我們今天這樣.
我們生了孩子想名字想好一點.
想一些很丟臉的名字.
當時人人都叫陳大文就是陳大文.
整條街都是陳大文.

$^{601}$你問一下就知道他的鄉下.
是多少人.
或者他屬於什麼人.
他的名字尤利亞安.
應該是巴勒斯坦的猶太人.
他們.
不知道什麼原因.
他們就去了羅馬或者與保羅一起.
去了地中海的世界.
他們跟過耶穌.
作為使徒.
他們的使徒.
全部都是曾經親自跟過耶穌.
或者親自見證過耶穌基督的復活.
這裡是教會的第一個女使徒.
Female Apostle.
如果你這麼看的話.
你沒有時間看完.
這麼多個名字.
有男也有很多的女.
其實不只是一些問安語這麼簡單.
就是保羅的教會名人榜.
他的 Honor Roll.
在耶穌和保羅的眼中.
很可以肯定.
女性在神國的使宮裡面.
是有尊重的地位.
甚至是具有尊重.
我們有什麼可以 take away(拿走)呢?.
第一.
就是我上次提過的.
重男輕女的思想.
某一種重男輕女的思想.
我說過了.
不只是男人才會重男輕女.
各位女士可能會.
你們覺得我只不過是弱質的女子.
我能做什麼呢?.
是不是?.
但是請你用上帝的角度去看待自己.

$^{641}$特別是教會神的工作.
不是 gender-based(性別基礎).
是 gift-based(禮物基礎).
不是按照性別去劃分.
按照你的恩賜,你的才幹去劃分.
不是 gender-based(性別基礎).
很清楚地讓我們看到.
從九曰到新曰.
一個這樣的角度.
你不要只顧自己.
不要把自己當成弱質女流.
上帝當然可以用男人.
但是上帝也可以同樣用女性.
我覺得新時代的教會就應該是這樣.
是不是?.
重男輕女.
男與女一起去配搭.
嗯.
我們很虛心地.
一起來服侍.
願意地一起去配搭.
為了每一個我們能夠.
為了這九位.
最後.
能夠出來做候選人.
主,我們求你在他們的生命當中.
在他們的信心裡面.
來成就你們的工作.
無論是姐妹還是弟兄.
讓我們經歷你的恩典.
在主的恩典當中我們去承載.
去領導,去事奉,去學習.
就是主耶穌基督你怎樣去服侍我們.
讓這個教會能夠成為一個.
有主理形象的教會.
我們感激老方昌.
多謝各位收看,再見!.
\newpage



\section{}
\label{sec:YsIxBX5tTR0}
\textbf{6/17/18 主日講道: 信靠主 (吳君才牧師)}
\newline
\newline
連結: \href{https://youtube.com/watch?v=YsIxBX5tTR0}{\texttt{https://youtube.com/watch?v=YsIxBX5tTR0}} ~~~~ 語音日期: 2019-07-16
\newline
\newline
\hyperref[sec:NpmOHXaC8DQ]{\small{< < < PREV SERMON < < <}}
~
\hyperref[sec:index_chronic]{\small{[返順時目]}}
~
\hyperref[sec:index_scriptual]{\small{[返順卷目]}}
~
\hyperref[sec:o0LouZaecN0]{\small{> > > NEXT SERMON > > >}}
\newline
\newline
$^{1}$足夠法定人數.
主席.
主的話必永不落空.
我們看第十和十一節的經文.
我用現代中文譯本讀給大家參照.
我的話像雪和雨水降下不再返回.
卻灌溉大地使五穀生長.
讓殺種的有種子.
讓吃的人有食糧.
我的話就是這樣.
神的話從不落空.
卻要完成我的計劃.
要執行我的旨意.
一旦神的話預備解開善明.
必定有奇妙的感召力.
令人為罪為義為審判.
會自己責備自己.
神的話更是你我腳前的燈路上的光.
會指引我們.
走在一條正確美好的道路.
神的話真的不會落空.
是有成效的.
依從祂的旨意和計劃而成就.
執行祂神聖的任務.
問題不是神的話語無力會落空.
而關鍵的就是人對神的話語是怎樣回應.
怎樣去珍惜.
怎樣去善用呢.
我記得有一次我在香港的時候.
我陪同一個未信主的人上法庭.
因為之前發生了一些事.
他拿不住自己的怒氣出手.
不幸傷害了人.
要告上法庭.
這是可大可小的事.
特別他是一個專業的人士.
如果有些案底.
就成世前途盡毀.
在這件事件當中我介入.
我以詩篇23篇送給他.

$^{41}$安慰他 堅諾他.
我和他一起祈禱.
我希望他抓緊神和主的應許.
其中詩篇23篇第4節你很熟悉.
我送給他其中一篇第4節.
是金句.
當日上庭的早上我約了他.
一起和他祈禱.
提醒他.
第4節怎樣說.
主是和你同在 祂的像 祂的光都安慰你.
果然很奇妙.
那天負責案件的法官.
真的很仁慈.
輕判了他.
罰款和賠償一些湯藥就了事.
當時剛剛有條新的例.
叫做改過自新的條例.
犯了這些輕微的罪.
三年之後會自動取消你的案底.
這又是這麼奇妙.
神真是這麼奇妙.
神的同在應許是何等的真相.
主的話是不會落空的.
聽姐妹今天你和我究竟說了什麼神的話.
以賽亞先知告訴你和我.
神的話是不會落空的.
只要人以赤誠使求的心.
必然有更真實和寶貴的經歷和體會.
主的話真的不會徒然返回.
必定執行祂的旨意.
完成祂的計劃.
但問題是我們有時人太畏縮.
太小信或沒有信心.
忘記了神的權能和忘記主的應許會兌現.
以致主的話語的福樂我們未嘗經歷到.
我們求主幫助我們.
神的話語是真金不怕紅爐火.
你嘗試抓住祂.
好好地善用.

$^{81}$你必然經歷更多的福氣.
我記得某一年黃昏時分.
我和教會一班職業團契的青年人.
即是已經出來工作的一些人.
大家一起出去公園報道.
派單張.
那次我們用四律.
很多人當時的團友都半信半疑.
可不可以用四律這麼簡單的東西.
說說聽聽就會決自信耶穌.
但四律所用的聖經的神的話語是滿有能力的.
神的話不會落空的.
他們見到一個大概四十歲的婦人.
在公園的燈光下.
有三幾個人圍著她.
用四律跟她報道傳福音.
這位婦女有學識有教養的.
竟然被這麼簡單的四律福音折服了.
而且居然肯決自祈禱.
連那幾個半信半疑的團友都大得激勵.
弟兄姊妹.
我們要不以福音為恥.
不以神的道為恥.
福音是神的大能.
要救一切相信的人.
主的話必永不落空.
主的道不會廢去.
天地會廢去.
但主的道是永遠長存.
進一步剖析.
究竟我們愛慕神的話.
神的道程度是怎樣的呢?.
你每天花多少時間讀經.
看神的話呢?.
還是忙到放在一旁呢?.
家裡的聖經是封塵的呢?.
有一次.
施暴真去到蘇格蘭.
路經一條小鎮.
他在一間小旅館裡度宿.

$^{121}$他在房間裡找找找.
找到一本舊的聖經.
拿出來一看.
為什麼這本聖經是穿洞的呢?.
原來是被蟲蛀過.
一打開燈光底下.
哇!這本聖經的透光.
是由創世記一直透到主到啟示錄.
施暴真祈禱說.
但願我能像那些蟲一樣.
將聖經完全鑽研通透.
另一位屬靈長者.
A.C. Dixon 就這件事回應說.
但願我們能夠這樣.
吃神的話語.
一句句的吃.
一張張的吃.
一卷卷的吃.
我們就一定受益不淺.
如果我們這樣默想神的話語.
我們就可以每天在其中.
找到豐富的寶藏.
請問我們是不是這樣.
好像詩篇的作者說.
我趁天未亮呼求.
我仰望你的話語.
我趁夜更未換.
將眼睜開.
為要思想你的話語呢?.
詩篇有十九篇.
一百四十八字所提醒.
究竟我們愛慕神的話.
讀聖經的心智有多少呢?.
如果我們真的愛慕神的話語.
我們就肯付代價去閱讀.
多累都會閱讀.
多忙都會閱讀.
會抽時間默想去靈修.
而且我們還會持之以恆.
有些人讀經,祈禱,靈修.

$^{161}$有一段時間得薄.
接著做很多事.
爛尾,正式是虎頭牌.
即是虎頭蛇尾.
這樣的時候.
難怪神的話語.
在我們生命和經驗裡面.
是落空的.
但其實是不落空的.
神的話語是不會徒然返回的.
我們看第三點.
主的恩典必引導賜福.
我們看第十二,十三節的經文.
你望一望.
我用原來中文譯本讀給大家聽.
你們將歡歡喜喜地離開巴比倫.
平平安安地歸回故鄉.
大山小山要一齊歌唱.
樹木要拍手歡呼.
長荊棘的地方將有松柏.
長疾泥的地方將有番石榴.
這是永遠的記號.
要紀念我上主的作為.
這一幅是喜樂無遺的圖畫.
是描述那些被父神親手救贖的人.
來到神所預備的居所.
大自然都和他一同歡樂.
呈現一番新的氣象.
作為神救贖永遠的一個證據.
這是印證主的恩典.
必引導賜福給倚靠父神的人.
各位姐妹.
不知道你們現在的心情.
或者你們面對一些什麼.
如果你真的走在一個低谷苦惱.
似乎你覺得神好像離你很遠很遠.
正如昔日以色列人.
被擄割破家亡一樣.
那段痛苦的日子.
那種滋味感受是不容易的.

$^{201}$但我請你抬起頭來.
仍然仰望倚靠主.
適合父神.
如何在他的子民歷盡艱困苦痛之後.
仍然以慈愛恩典引導和賜福.
我希望今天這位色在金在永在的大能主宰.
仍然會拖大你.
仍然會顧念你.
所以請你用信心的眼光.
等候和仰望主.
倚靠他的必不設羞愧.
因為神不會撇棄你和我.
尤其是當你的心是專注於神.
在你尊敬神的國和神的意為念的時候.
你會更加配的體會神的恩典和福氣.
美國石油大王約翰·洛克菲勒.
曾經是世界的首席富翁.
不過他只顧著賺錢.
只顧著自己.
他變成一個很可憐的人.
因為他晚晚失業.
他疑心又很重.
所以他需要特別請一些保鏢來保護自己的安危.
到了他53歲那一年.
他患了一種怪病.
醫生都不知道是怎樣一回事.
他的頭髮還脫得比我少.
還有身體又萎縮.
醫生說估計這樣的情況只有一年的壽命.
洛克菲勒思前想後.
想到自己的生命.
想到自己的健康.
想到永恆的問題.
他開始改變.
我不能再自我中心的生活.
我要想想其他.
我要顧念其他有需要的人.
關心教會和天國的事.
自此開始捐獻金錢.
資助教會和有需要的窮人.

$^{241}$他更成立了洛克菲勒基金.
負擔稀有疾病的研究經費.
就是這樣著力去做的時候.
全心投入的時候.
他的健康慢慢好轉過來.
結果和醫生的預測相反.
因為洛克菲勒.
竟然活到98歲的高齡才去世.
請姐妹.
當我們人在心智上,情感上,觀念上都聚焦在神的時候.
當我們能夠有個心是關顧人.
關顧神的國,神的義,神的示唯念的時候.
神不獨不虧待他.
反而有加倍的恩典和福氣臨到.
主的恩典確實是與倚靠他,敬愛他,敬畏他的人同在同行.
父神以信實為量.
以慈悲仁愛引導賜福給敬畏他的人.
以色列人被擄歸國的歷史事實.
正是一個很好的提醒.
身邊弟兄姊妹.
每一個成為的弟兄姊妹.
我們要操練更多的倚靠神.
更多的經歷主.
今日那三樣寶貴的學習功課就是.
第一,我們要認定神的安排勝於世人.
神的意念和道路總比人更高超,更可信賴.
第二,神的話語必永不落空.
不會徒然返回的,定然會成就的.
我們更加要以愛慕和肯宣揚主的話語真道的心.
去行事慧入.
第三,主的恩典必定引導賜福.
但是我們需要先求神的國和神的義.
我們低頭有個祈禱,結束.
神的簡單的道.
祂是深藏很寶貴的提醒.
讓我們知道你的安排勝過我們人所想.
你的意念道路比人更高超可靠.
神下你的話是永不落空的.
讓我們重視你的話語.
揣摩你的話語.

$^{281}$靈修,閱讀你的話語.
你的話語不會徒然返回,必定成就.
讓我們有個愛慕你的話語.
甚至宣揚你的話語的心.
神我們更求你幫助我們.
信得過你,有恩典會引導我們.
但是我們要先求神你的國和你的義.
加能,賜力給每個弟兄姊妹.
祈禱奉主耶穌基督的名.
阿們.
\newpage



\section{}
\label{sec:o0LouZaecN0}
\textbf{7/15/2018 主日講道: When Should We Pray (Steve Tan)}
\newline
\newline
連結: \href{https://youtube.com/watch?v=o0LouZaecN0}{\texttt{https://youtube.com/watch?v=o0LouZaecN0}} ~~~~ 語音日期: 2019-07-16
\newline
\newline
\hyperref[sec:YsIxBX5tTR0]{\small{< < < PREV SERMON < < <}}
~
\hyperref[sec:index_chronic]{\small{[返順時目]}}
~
\hyperref[sec:index_scriptual]{\small{[返順卷目]}}
~
\hyperref[sec:lgP5_8Xevo8]{\small{> > > NEXT SERMON > > >}}
\newline
\newline
$^{1}$[Applause].
[Cantonese].
\newpage



\section{詩篇}
\label{sec:lgP5_8Xevo8}
\textbf{詩篇人生 (2) - 詩篇 90 講道 (陳崇基牧師)}
\newline
\newline
連結: \href{https://youtube.com/watch?v=lgP5-8Xevo8}{\texttt{https://youtube.com/watch?v=lgP5-8Xevo8}} ~~~~ 語音日期: 2019-07-16
\newline
\newline
\hyperref[sec:o0LouZaecN0]{\small{< < < PREV SERMON < < <}}
~
\hyperref[sec:index_chronic]{\small{[返順時目]}}
~
\hyperref[sec:index_scriptual]{\small{[返順卷目]}}
~
\hyperref[sec:5SLOPVCSwvk]{\small{> > > NEXT SERMON > > >}}
\newline
\newline
$^{1}$你知道聖經裡面有很多不同的內容.
有些是故事,有些是律法,有些是書信.
其中一些是詩歌.
就是說詩人寫了的時候,本身是用音樂來唱出來.
當然那些音樂已經失傳了.
但是歌詞仍然留給我們.
所以看這些詩詞歌賦的時候.
我記得上一次我們開始的時候.
我就說應該問兩個問題.
起碼問兩個問題就是什麼呢?.
你看這些詩詞歌賦的時候.
你問問到底作者,詩人是誰.
還有他唱出來的時候.
或者寫出來的這些詩歌.
他是在什麼情況,背景下寫出來的.
他想述說些什麼呢?.
以致我們在詩歌裡面能夠看多一點這些內容.
我們首先看第一兩節.
第一兩節是這麼說的.
他說:神人摩西的祈禱.
整首詩歌都是一個禱告.
主呀,你世世代代在我們居所.
群山尚未生出.
大地世界尚未造成的時候.
從永遠到永遠,你是神.
只是這兩節打開之後.
你就會發現.
似乎就帶我們回到聖經頭五本書.
我們知道聖經有兩舊.
舊約傳書和新約傳書.
舊約傳書就是耶穌基督未來之前所寫的.
新約傳書就是耶穌基督來了之後所寫的.
但是舊約呢.
頭五本書我們有時候叫做摩西五經.
摩西五經有五本書的.
《創世記》,《出埃及記》,《利未記》,《民數記》和《申命記》.
為什麼說這兩節開頭.
就帶我們回到摩西五經頭五本書呢.
因為這裡說到群山尚未出.
世界尚未造成的時候.

$^{41}$上帝已經在這裡.
上帝是永遠到永遠的.
即是上帝是無始無終的.
我們每個人都有開始的.
我們地上的生命也會有結終的.
但是上帝是沒有的.
祂是由永恆到永恆.
祂就是在這裡.
並且這裡說到神人摩西.
形容摩西是man of God 神的人.
屬神的人.
或者神所尊貴的人.
man of God.
在舊約聖經裡面.
摩西五經頭五本書.
只有一個地方提到摩西.
是說他是man of God 神人.
就是《申命記》第三十三章.
《申命記》你記得就是第五本書.
摩西五經的結束.
最後那兩章.
他就說到摩西臨終的時候.
有一個祝福和祈禱.
他說以下就是神人摩西.
未死之前.
先為以色列人所祝的福.
所以告訴我們.
回到聖經的開頭.
提到我們知道上帝.
是從亙古到永恆.
很多時候我們生命裡面.
我們遇見的東西.
我們很多時候會將那件事放大.
無論我們生活裡面的焦慮.
我們的疾病.
我們的前途.
這些問題我們經常將它放在眼前.
放大到一個地步.
我們只看見這些問題.
神人摩西的禱告一說.

$^{81}$原來是有永恆的.
將我們所注目的前面的問題.
我們無限放大的問題.
就將它拿開一點.
然後讓我們看到整個局面.
整幅圖畫.
從亙古到永恆.
耶和華是神.
你世世代代作我們的居所.
上帝所創造給我們的家居.
住在這個世界.
我們就是一個這樣的地方.
原來是從這位永恆的神而來的.
然後第三節開始.
這個祈禱就帶我們回頭.
再想想我們現在的人生.
上帝是永恆的.
但是我們呢.
祂說.
神,你使人歸於塵土.
說你們世人要歸回.
在你看來千年如同已過的昨日.
猶如夜更的一更.
你叫他們如水沖去.
如水一較.
早晨如生長的草.
早晨發出的野草.
然後晚上就割下來去枯乾.
歸回塵土.
這裡就告訴我們.
我們生命是有限的.
我們的人生永遠都是有啟發的.
有什麼好處呢?.
這個詩篇所提到的東西.
就是告訴我們我們有限期.
我們每個人都有一個警鐘.
無論你今天是年輕也好.
你是年長也好.
你是中年也好.
我們每一個都會歸回.

$^{121}$歸回塵土.
就是說我們生命裡面.
只有這個機會.
只有一個機會.
如何再思考生命如何過去.
摩西的祈禱就試探索這個問題.
人是有限的.
我們很快就會過去.
上帝是永恆的.
但是這兩者中間.
神和我們人.
這位永恆的主.
和我們這些有限的人.
很快我們的生命就會過去.
我們的關係是如何的.
dynamics.
如何去相處呢?.
你想像一下.
有一張很大的白色的dry erase board.
你在白板上面.
你點一點.
然後你數1,2,3.
你馬上抹掉.
這個就是你的生命.
我們就是這麼快.
或者我們數三聲.
我們拍一下手指.
可以嗎?.
1,2,3.
這個就是你和我的生命.
就是可以這麼快.
我們如水沖去.
如早晨生出的草.
然後晚上就會不見.
但是永恆的上帝看時間.
是不同的.
因為神是outside of time.
就是他在時間以外的.
我們都不知道如何去描述這件事.
因為我們全部都在時間以內的.

$^{161}$時空底下生活的人.
上帝不是的.
上帝是在時空的以外的存在者.
我們只能夠說.
他是outside of time.
就是在時間以外的.
我們過一千年.
等於上帝過一天.
我們過一更.
夜更.
一更大概是四個小時.
三個小時的時間.
上帝是這麼快的看時間.
我們過了一千年.
對上帝來說.
只不過是夜裡面的一更.
昨天你睡了幾更.
已經幾千年了.
上帝是永恆的.
所以我們有限的人.
我們看時間和上帝看時間是不同的.
或者可以這樣說.
打個比喻.
我的小朋友小的時候.
幾歲三四歲的時候.
我想我們中間很多家庭.
都是在經歷這個生活的階段.
你開車去某個地方.
然後小朋友坐在後面.
那時候小的時候.
我們無論去suburb.
四十分鐘車.
三十分鐘車我們覺得很快會到.
但是小朋友會覺得很久.
他每分鐘就問你.
How long does it take?.
到底還有多久?.
問完之後回答他.
要再等多三十分鐘.
過多一分鐘之後.

$^{201}$他又說還要等多久才到?.
How long?.
他對時間的觀感.
和我對時間的觀感是完全不同的.
你將這個分別再誠意無限.
就是永恆的上帝.
和我們有限的人.
那個時間感.
是完全不同的.
你記不記得你年輕的時候.
你很年輕的時候.
你等生日呢.
等很久.
怎麼這麼久都還沒到生日?.
剛剛生日結束了.
下次生日什麼時候?.
我記得我女兒幾歲的時候.
她等她生日.
就是預備她生日party的時候.
她是幾個月前預備的.
我心想用不著這麼早.
你覺不覺得你年紀大了的時候.
生日來得很快.
眨眼.
原來生日了.
如水沖去.
一眨眼我們就會過.
上帝經歷時間.
和我們經歷時間.
是完全兩回事.
不單止這樣.
上帝看我們人生裡面的過煩.
我們的破碎.
我們的罪.
都是有所不同.
我們過每一天.
我們做錯事.
做錯了之後可能say sorry.
然後如水沖去.
我們就繼續過生活.

$^{241}$做錯了一件事.
第二天都會繼續.
Life goes on.
但對永恆的主就不同.
因為永恆上帝去看我們每一天生活的時候.
因為時間的過去.
和將來都常在他面前.
因為他永恆到永恆.
所以整天的事都在他面前.
第七節.
我們做人的.
因為你上帝的怒氣而消滅.
因你的憤怒而驚惶.
你將我們的罪孽常擺在你的面前.
將我們的隱惡擺在你面光之中.
我們經過的日子都在你震怒之下.
我們進度的年歲好像人一聲嘆息.
我們一生的年日是七十歲.
若強壯的話可以有八十歲.
但其中所驚誇的.
不過是勞苦受煩.
轉眼成空.
我們便如飛而去.
誰曉得你怒氣的權勢.
誰安著你該受的敬畏.
曉得你的憤怒.
我們的罪孽常在這位永恆的主.
時間以外的主的面前.
這也是一個很難明白的概念.
正如我所說.
神是外在的時間.
我們在用時空之下的語言來形容他.
到底他怎樣看我們破碎的人生.
我們的過犯.
他說常在他的面前.
常在他的面光之中.
我們的隱惡都在上帝面前.
可能可以這樣說.
我打個比喻.
又說回年輕的家長.

$^{281}$小孩子小的時候.
我們不是經常跟他們拍照嗎.
經常跟他們拍照.
拍完有些漂亮的我們便曬出來.
印出來.
我們便經常放在家裡.
所以我們看小孩子.
永遠都是看他小時候的樣子.
我的女兒十八歲了.
最近我跟她說笑.
我叫她妹妹.
因為小時候我都叫她妹妹.
她說我不是妹妹.
除了她生活中有個阿姨還是叫她妹妹.
沒有人再叫她妹妹.
但在我心目中.
她都是常時妹妹.
我永遠都想起她幾歲的時候.
想起她以前的樣子.
想起那些照片.
這就是時間觀感差別的問題.
上帝看我們其實都是一樣.
上帝是永恆的.
他是無所不知的.
過去現在將來都在他的面前.
這位聖潔的上帝.
這位公義的上帝.
如果按照他的公義的話.
我們每個人應該得到的.
就是審判.
所以他在這裡說.
嘆息我們整個人生.
在這位永恆的主裡面.
無論有多少年歲都好.
都是在他的憤怒和審判當中.
他說人的生命.
長可以七十歲.
強壯的話可以八十歲.
在我們中間有很多長者可能過了八十歲.
這些全部都是bonus.

$^{321}$但人的生命大概都是這幾十年.
這篇詩篇的祈禱就是掙扎的問題.
他就向上帝禱告.
他說上帝如果你這樣看我們的話.
你是不是真的想這樣繼續紀念我們.
在你永恆的主的眼中.
我們的罪事上在你的面前.
如果你這樣定義我和你之間的關係.
我們只會憤怒.
還有什麼罪呢.
可不可以用另一個方法去定義我們的關係.
第十二節.
他說求你教我們怎樣數算自己的日子.
好使我們得著智慧的心.
在這位永恆的主的目光下.
怎樣去過得有智慧呢.
數算不是只是計算有多少日子.
數算是有計劃的意思.
數算是有accounting的意思.
即是evaluation的意思.
核數的意思.
即是核數一下我們的生命裡面.
評估一下.
到底我們有限的年日.
我們是否懂得用一個智慧的心去活出來.
這是神人摩西的禱告.
如果你想起摩西的話.
詩篇是一本書.
裡面有150首詩歌是集成的.
有些是大衛寫的.
有些是其他的作者寫的.
只有一篇是摩西所寫的.
就是這篇90篇.
如果你是熟悉摩西的人生的話.
在摩西五經裡面記載摩西的生平.
你可能會想起某一些字眼.
在這個禱告裡面.
在詩篇90篇裡面.
經常在摩西的生平的時間出現.
或者那個episode出現.

$^{361}$我想起什麼episode會經常有.
要上帝轉回等等.
這些字眼在哪裡出現呢?.
你知道摩西的人生.
很簡單的說.
他是在埃及出生.
然後在埃及的宮廷長大.
但他沒有忘記自己是希伯來人的血統.
然後他又流亡到別處.
因為不同的事故.
然後上帝在曠野裡呼召他回到埃及.
救他自己希伯來的人民.
帶領以色列人出埃及.
過紅海.
在曠野裡他們去了一個山.
叫西乃山.
在那裡上帝和他的人民.
藉著摩西來到立約.
拿了十戒,拿了法板,拿了律法.
然後在曠野裡留殮了40年.
快要去到上帝所認許的迦南地獄.
他就進不去.
他只能夠看到.
這就是摩西說一聲的人生.
如水沖去就過了.
但有件很大的事.
聖經有很多的細節說到.
就是在摩西帶領以色列人出埃及的時候.
我們說過了.
我們看聖經的時候.
我們不可以只看聖經的內容.
我們要看聖經作者怎樣寫出內容.
因為他怎樣寫出內容呢?.
很能告訴我們他想說些什麼.
這裡已經是摩西帶領以色列人出埃及.
經過了一段曠野的路程.
然後去到西乃山的山腳底下.
神告訴摩西.
你上來.
上到山頂.

$^{401}$他和神面對面.
他將上帝將律法書交給摩西.
再帶回去給人民.
山腳底下發生什麼事呢?.
那些人等得久的時候.
他們發現摩西會不會壽終正寢呢?.
不見了人呢?.
他們就說.
不如我們弄一隻金牛出來.
他們就將耳環,首飾等等.
金銀寶石.
然後將它溶化.
然後弄一隻金牛出來.
以前古代的人就是這樣.
他們拜的是什麼呢?.
拜的是最有力氣的東西.
就是家畜裡面.
人養的動物.
能力最大的是什麼呢?.
就是牛.
就好像你今天在家裡.
你最疼愛的可能是你的車.
最大力的東西.
或者最有能量的東西.
最有產力的東西.
就是牛.
他們就弄了一隻金牛.
然後他們說.
這隻金牛就是我們的神.
他們做了一個這樣的牛像.
然後就去下拜.
這就是24章至40章所說的.
首先24章說什麼呢?.
你看到A-24.
神的榮耀首先來到西大山.
然後將這些律法書.
或者是會幕裡面的藍圖.
告訴摩西.
你要準備了.
要怎樣弄祭壇,約櫃.

$^{441}$會幕裡面的布.
等等不同的金燈台.
這就是一個所謂的Portable的聖殿.
就是說神願意.
我願意.
永恆的主願意住在你們中間.
你就要建造這個會幕.
他拿了藍圖之後.
32章就交代什麼事呢?.
C-32.
就是人民背若.
他們就做了一隻金牛.
原來摩西受到政懲.
不如弄一隻金牛吧.
摩西還沒下山就聽到他們.
下面搞什麼鬼呢?.
然後呢.
上帝就發怒.
因為我剛剛才拯救了你們出來.
你們就背著我.
然後就弄了一隻金牛.
拜一隻金牛,拜假神.
拜偶像.
然後呢.
這裡很奇妙.
就是摩西.
在上帝發異怒的時候.
想說殲滅他的人民的時候.
摩西就為了他的人民來禱告.
告訴上帝求你赦免他們.
並且回傳.
上帝真的有這樣做.
然後上帝就仍然願意.
用他的榮耀帶領他的子民.
一路走.
曠野他們.
這就是X-33章.
這個地方.
34章又說.
雖然他們背若.

$^{481}$但上帝願意再立若.
摩西不是很生氣嗎?.
你記得那個故事.
那個法板打爛了.
但上帝說你再上來.
我願意再跟他們立若.
再雕法板.
然後再帶回來.
到B所謂一點.
B'.
35至39章.
就真的可以建造回望.
這個故事結束.
出埃及的結束.
最後那章40章.
A'.
神的榮耀.
仍然在回望.
這個就是那個故事的背景.
我相信就是.
詩篇90篇.
背後的故事.
就是他怎樣為人民來土共.
我們來看看.
出埃及的第32章.
然後對摩西說.
你看看這些百姓是硬著脖子的百姓.
你扯油著我.
你誘我.
我要向我的百姓發烈怒.
憤怒.
將他們滅絕.
使你的後裔成為大國.
這句才是過癮的.
他說我要滅絕我所拯救的這班子民.
因為他們背弱.
但是呢.
如果將他們滅絕的話.
他說摩西.
不如從你身上.

$^{521}$我將你成為大國好不好.
如果是你的你好不好.
都挺好的.
因為我的名聲就會顯大.
但摩西居然做了一件事.
他說摩西便懇求耶和華他力神說.
耶和華.
你為什麼向你的百姓來發烈怒.
求你轉回.
求你轉移不發你的烈怒.
回心轉意.
不要降和於你的百姓.
求你紀念你僕人阿拉罕爾薩.
以色列亞國.
第14節.
於是耶和華神就回心轉意.
回心轉意.
發烈怒等等.
這些我們已經highlight了的字.
正正是在詩篇第90篇.
經常出現.
他說求你轉回.
13節.
他說耶和華.
我們要等到什麼時候.
求你轉回.
為你的僕人回心轉意.
求你使我們早早飽得你的慈愛.
好叫我們一生一世可以歡呼喜樂.
求你照著你使我們受苦的日子.
糟爛的年歲叫我們有喜樂.
願你的作為向你的僕人顯現.
願你的榮耀向他們子孫.
他們下一代來顯明.
願主我們的神榮美歸於我們的身上.
願你建立我們手足做的功.
我們手足做的功.
願你建立.
這篇詩篇是一首哀歌.
他說如果永恆的主常常見到我們的罪.

$^{561}$常常擺在他面前的話.
我們就不得了.
下一半他說.
可不可以用另一個定義來跟我們相處.
你這位上帝.
可不可以你回心轉意.
在你的慈愛.
在你的恩典當中.
來跟我們建立關係.
下一半這裡說.
請求上帝你.
根據你的慈愛.
你都決定神與人之間的關係.
你要轉會.
因為沒有了恩典的話.
我們只有審判.
我們過這一生只有憤怒.
上帝的耳濡.
我們只有苦難糟難的日子.
所以為什麼詩篇第九十篇結論的時候這麼說.
他說願主我們的神.
祂的慈愛臨到我們身上.
願你為我們的緣故.
建立我們所做的功.
我們所做的功願你建立.
你再回到剛才回顧摩西這一段的時候.
你就知道工作這個字在哪裡出現.
就是在拜金牛之後.
上帝說我回心轉意.
我願意赦免他們.
再跟他們納藥.
然後他們興建會幕的時候.
就經常提到有工作.
出現工作這個字.
很多很多次.
就是興建會幕的工作.
興建會幕就是說讓上帝的榮耀能夠住在我們中間.
上帝不依我們的罪來跟我們相處.
上帝以祂的慈愛.
祂的恩典臨到我們身上.

$^{601}$跟我們一起同行.
這是摩西在CB90篇的禱告.
願我所做的工作.
你能夠建立.
那些工作是什麼工作.
就是我們可以參與在神的會幕興建的工作.
就是說上帝用祂的恩典赦免了我們之後.
我們得到赦罪之恩.
我們能夠參與祂的工作.
這是很美的事.
我們是有罪.
我們是很脆弱.
我們是如水沖去.
我們是人生很快會過.
但是上帝居然肯去施恩.
我們蒙恩得救赦免的罪人.
我們能夠去參與祂榮耀的工作.
你有沒有想過你可以參與上帝的工作.
不是因為你自己碰.
不是因為你被人好.
我們很容易在教會跟別人比較.
在我們生命跟別人去比較.
你看我不是比他好嗎.
為什麼上帝不選擇我做.
你有沒有想過這裡說的.
你能夠參與上帝的工作.
不是因為你有多少才幹.
不是你覺得自己有多好.
而是因為上帝赦罪的恩典.
不要輕看上帝的恩典.
不要看你自己那麼高.
因為你可驚誇的.
只不過是老虎受煩.
不是你自己有多少.
你要誇的就是誇上帝的恩典.
祂願意讓我能夠參與祂的工作.
有生之年能夠去做到上帝的聖供.
我們多麼好.
剛才進來崇拜之前.
有些弟兄姊妹早到.

$^{641}$有個會友給我一張照片.
他說有朋友在別的地方.
都算年長了.
剛剛到完神學科程.
裝備自己.
能夠參與上帝的工作.
有生之年我們能夠有這樣的機會.
能夠去討好上帝.
真是很大的恩典.
我們生命是有限期.
並且有很多事情是我們猜不到的.
我媽現在71歲.
前兩星期她去醫院.
突然間去醫院.
我們不是剛剛去完嗎.
為什麼突然去醫院.
因為急症她就進了醫院.
因為血壓高.
一直都下不來.
她完全掌握不到.
為什麼突然間會這樣.
我們人繼續大下去.
當我們發現生日越來越快的時候.
就好像車一樣.
有些事情會故障.
上星期我以為我可以專心預備教會的事.
誰知道我又病了.
我以為我老虎打死幾隻.
剛好上星期.
天氣又好.
誰知道我又病了.
我又躺下了一天半.
我們沒有什麼可以經過.
但既然上帝給我們有赦罪的恩典.
我們能夠參與他的工作.
這是我們應該把握的機會.
你能夠參與上帝的工作不會多.
救主給我們.
建立我們守守祖國.
\newpage



\section{}
\label{sec:5SLOPVCSwvk}
\textbf{12/16/18 主日講道: The Meaning of Christmas (Sam Johnson)}
\newline
\newline
連結: \href{https://youtube.com/watch?v=5SLOPVCSwvk}{\texttt{https://youtube.com/watch?v=5SLOPVCSwvk}} ~~~~ 語音日期: 2019-07-17
\newline
\newline
\hyperref[sec:lgP5_8Xevo8]{\small{< < < PREV SERMON < < <}}
~
\hyperref[sec:index_chronic]{\small{[返順時目]}}
~
\hyperref[sec:index_scriptual]{\small{[返順卷目]}}
~
\hyperref[sec:J7N4kVhPD7E]{\small{> > > NEXT SERMON > > >}}
\newline
\newline
$^{1}$[no audio].
Well, good morning. I just want to start by saying how grateful I am to be here.
and how grateful I am to you all for having me and my family here..
[Cantonese].
We love Urban Voice and love the time that we've spent here, and being here today is really special for us..
It feels like a little reunion seeing a lot of good friends again, so it's a blessing..
[Cantonese].
And this time of year is a special time of year, time of year that I enjoy when we celebrate Christmas..
[Cantonese].
I enjoy the get-togethers like today, the food, the decorations, all of those things..
[Cantonese].
I even like the Christmas movies, all the good ones and the cheesy ones that are on TV all the time this year..
[Cantonese].
But this morning I want us to think a little bit about the meaning of Christmas and what the heart of the meaning is..
[Cantonese].
Because Christmas is much more than just the gifts and the decorations and the lights and the time with family and friends..
Those things are wonderful, but the meaning is deeper than that..
[Cantonese].
So let's look at a few verses in the book of Luke, chapter 2..
[Cantonese].
We're going to read verses 8 to 14..
[Cantonese].
These verses come from right after the birth of Jesus..
[Cantonese].
Jesus has just been born in Bethlehem in the stable..
[Cantonese].
And the Bible tells us that there were shepherds living out in the fields nearby, keeping watch over their flocks at night..
[Cantonese].
An angel of the Lord appeared to them, and the glory of the Lord shone around them, and they were terrified..
[Cantonese].
But the angel said to them, "Do not be afraid..
[Cantonese].
I bring you good news of great joy that will be for all the people..
[Cantonese].
Today, in the town of David, a Savior has been born to you..
[Cantonese].
He is Christ, the Lord..
This will be a sign to you..
You will find a baby wrapped in cloths and lying in a manger..
[Cantonese].

$^{41}$Suddenly, a great company of the heavenly host appeared with the angel..
[Cantonese].
Praising God and saying, "Glory to God in the highest.".
[Cantonese].
And on earth, peace to men, among whom his favor rests..
[Cantonese].
Now these verses tell us that the birth of Jesus was good news, really good news..
[Cantonese].
In fact, verse 10 specifically tells us it was good news of great joy to all people..
[Cantonese].
But why? Why was it such good news?.
[Cantonese].
Why was this such a special event?.
[Cantonese].
To begin to answer that question,.
[Cantonese].
we have to zoom out a little bit and consider the place of the Christmas story in human history..
[Cantonese].
The Bible details to us God's work throughout history..
[Cantonese].
And although the Bible contains many different stories,.
[Cantonese].
it's really one big story..
[Cantonese].
One grand narrative that runs from beginning to end..
[Cantonese].
It has a lot of different climatic moments..
[Cantonese].
And the Christmas story is a climatic moment in the Bible..
[Cantonese].
But it's not the beginning of the story..
[Cantonese].
And it's not the end of the story either..
[Cantonese].
To start to grasp the importance of the birth of Jesus,.
[Cantonese].
we need to go back to the beginning of the story.
when God began to work with humans..
[Cantonese].
The story of the Bible begins with the act of creation..

$^{81}$[Cantonese].
Now before creation there was God..
[Cantonese].
The scripture tells us that he has always existed..
[Cantonese].
He has no beginning or end. He is uncreated. He has always been..
[Cantonese].
Earth did have a beginning and so did we..
[Cantonese].
God spoke the earth into existence..
[Cantonese].
He created out of nothing from the force of his power and wisdom..
[Cantonese].
And he made everything on the earth including us..
[Cantonese].
God created us and made us uniquely in his image..
[Cantonese].
Nothing else that he made was made in his image..
[Cantonese].
Humans are very special in that regard..
[Cantonese].
Now being made in his image means many beautiful things for us..
[Cantonese].
But one thing that it means is that we are relational.
beings..
[Cantonese].
God created us for relationship..
[Cantonese].
He created us for a right and healthy relationship to him..
[Cantonese].
And for right and healthy relationships with one another..
[Cantonese].
But not long after the moment of creation,.
[Cantonese].
there was a problem..
And it was us..
[Cantonese].
See, we as humans have always thought we're a little bit smarter than God..
[Cantonese].
And that we could really improve on the world he made if only we were in charge..

$^{121}$[Cantonese].
And the very first man and woman decided they didn't want to let God be in control..
[Cantonese].
They didn't want to follow his instructions..
[Cantonese].
They felt like they could do better making their own choices..
[Cantonese].
And so they chose to reject God..
[Cantonese].
And to reject his instructions for life..
[Cantonese].
And that was the first broken relationship..
[Cantonese].
The first breakup that happened..
[Cantonese].
And it wasn't God, it was us..
[Cantonese].
We're the ones who said, "God, it's not me, it's you. It's your fault.".
[Cantonese].
"I want to go my own way and do things on my own.".
[Cantonese].
And ever since that moment throughout time,.
[Cantonese].
all of us have made that same choice..
[Cantonese].
All of us at some point have chosen to reject God..
[Cantonese].
And to pursue our own comfort and our own ideas and what seems best to us..
[Cantonese].
And that has led to all of the pain and the heartache of the human condition..
[Cantonese].
Instead of choosing right relationship with God and each other,.
[Cantonese].
we choose to please ourselves..
[Cantonese].
And out of this comes selfishness,.
and greed and lust and oppression, hatred, wars, jealousy, racism, lies, betrayal, theft,.
the list goes on and on, all these things..
[Cantonese].
And the reality is there's not any one of us who hasn't done that..

$^{161}$[Cantonese].
Everyone in this room has at some point acted selfishly..
[Cantonese].
Or done something to hurt another human..
[Cantonese].
We've all done that, I certainly have..
[Cantonese].
All of us have sinned and fallen short of what God wants us to be..
[Cantonese].
But there's a glimmer of hope here..
[Cantonese].
And that is that God wasn't okay with a broken relationship..
[Cantonese].
God didn't want to leave it that way..
[Cantonese].
So God had a plan from the very beginning..
[Cantonese].
His plan was to restore and heal our relationship..
[Cantonese].
To make things right..
[Cantonese].
And to make a way for us to be forgiven..
[Cantonese].
God's plan was to send us something..
[Cantonese].
Or rather someone..
[Cantonese].
God didn't choose to send us a letter or a text or a voicemail..
[Cantonese].
Not an email or a DM or any other kind of message..
[Cantonese].
God sent his only son..
[Cantonese].
And that's who Jesus was..
[Cantonese].
The son of God sent on a rescue mission..
[Cantonese].
To heal and restore and ultimately save us..
[Cantonese].
That's why in verse 10 the angel says the birth of Jesus is good news of great joy..

$^{201}$[Cantonese].
Jesus was the first gift of the first Christmas..
[Cantonese].
And God gave this gift to all of us..
[Cantonese].
Every human that has ever lived or ever will live on the earth..
[Cantonese].
God gave us this gift of his son..
[Cantonese].
And through his son..
[Cantonese].
The gift of relationship with our creator..
[Cantonese].
But how is it that Jesus could rescue us?.
[Cantonese].
And what was it about his birth that was so different than all the other births?.
[Cantonese].
The Bible tells us that Jesus was and is someone very special..
[Cantonese].
Because he was the son of God..
[Cantonese].
He was fully God..
[Cantonese].
And yet also fully human..
[Cantonese].
The book of John chapter 1 tells us that Jesus is the word of God made flesh..
[Cantonese].
The revealing of all that God is..
[Cantonese].
That means all of God's justice, God's love, God's peace, and infinitely more than that.
brought together into human flesh..
[Cantonese].
Full of grace and truth..
[Cantonese].
Jesus was Emmanuel..
[Cantonese].
God with us..
[Cantonese].
Notice in verse 14, the original language doesn't just communicate peace to us..
[Cantonese].

$^{241}$It communicates peace among us..
[Cantonese].
Jesus was the manifestation of God..
[Cantonese].
Not just to show us who God is..
[Cantonese].
But to be with us..
[Cantonese].
To live among us..
[Cantonese].
And to make us able to live with God..
[Cantonese].
Jesus left the comfort of heaven to do this..
[Cantonese].
He made himself subject to all the limitations and brokenness of humanity..
[Cantonese].
So that he could rescue us..
[Cantonese].
Because he was human..
[Cantonese].
Jesus dealt with every struggle that we deal with..
[Cantonese].
He became tired..
[Cantonese].
He experienced hunger..
[Cantonese].
He got cold..
[Cantonese].
All of those things that happened to us..
[Cantonese].
He was also tempted to reject God and sin just like we are..
[Cantonese].
He was tempted to be angry..
[Cantonese].
To be selfish..
[Cantonese].
To lust..
[Cantonese].
To seek revenge on those who had wronged him..
[Cantonese].

$^{281}$He was tempted to be greedy and pursue his own comfort..
[Cantonese].
He was tempted to be lazy and to be fearful and to give in to his fears..
[Cantonese].
He was tempted to be consumed with anxiety and to be stressed and worried..
[Cantonese].
All the things that challenge us and hurt us..
[Cantonese].
He dealt with those same things..
[Cantonese].
But in all those temptations he didn't sin..
[Cantonese].
He didn't give in to any of those things..
[Cantonese].
He loved and served and healed those that he met here on earth..
[Cantonese].
And at the end of his life, he gave himself as a sacrifice for us..
[Cantonese].
He died on a cross..
[Cantonese].
Murdered in a way that he never deserved..
[Cantonese].
So that he could win life for us..
[Cantonese].
In a way that we never deserved..
[Cantonese].
He traded his perfection for our brokenness..
[Cantonese].
So that through him, we might find perfection and have our brokenness restored..
[Cantonese].
His death in our place gave us an opportunity to know our creator..
[Cantonese].
But he didn't stay dead..
[Cantonese].
After three days he was gloriously resurrected..
[Cantonese].
And he defeated death once for all..
[Cantonese].
He opened a doorway to freedom..
[Cantonese].

$^{321}$From sin and death..
[Cantonese].
He opened a doorway to relationship with our creator..
[Cantonese].
And he opened a doorway to heaven and eternal life for those who believe in him..
[Cantonese].
That's why his birth was good news..
[Cantonese].
It was the beginning of God's great rescue plan to restore us..
[Cantonese].
The arrival of Jesus was the first step in our healing..
[Cantonese].
And that brings us to another unique and powerful truth about the birth of Jesus..
[Cantonese].
And that is the intentional nature of Jesus stepping into our world..
[Cantonese].
He did it on purpose..
[Cantonese].
Have you ever thought about the fact that Jesus actually chose where he would be born?.
[Cantonese].
He's the only person in history who's been able to do that..
[Cantonese].
If you could go back and choose your birth,.
[Cantonese].
your family, where you would be born,.
[Cantonese].
maybe there's something that you would change..
[Cantonese].
But Jesus did have the opportunity to choose..
[Cantonese].
He could have been born anywhere to any parents..
[Cantonese].
But he chose to be born to an unknown, poor, and hard-working couple.
from a tiny, struggling town in an oppressed country..
[Cantonese].
When Jesus was born,.
[Cantonese].
Rome ruled his homeland..
[Cantonese].
A Caesar sat at the pinnacle of human power in Europe..

$^{361}$[Cantonese].
And the Han Dynasty sat at the pinnacle of human power in Asia..
[Cantonese].
There were kings throughout the world..
[Cantonese].
Influential families,.
[Cantonese].
wealthy clans,.
[Cantonese].
people with palaces and grand estates..
[Cantonese].
And riches and servants and animals..
[Cantonese].
But Jesus didn't choose to be born there..
[Cantonese].
He didn't choose to be born in a capital city..
[Cantonese].
Or in a palace..
[Cantonese].
He chose to be born in a country that was suffering under heavy oppression..
[Cantonese].
He was born in a small town so full of travelers,.
[Cantonese].
that his parents couldn't even find a decent room..
[Cantonese].
He was born in a place where animals were kept..
[Cantonese].
And his first bed was literally a feeding trough..
[Cantonese].
This shows us clearly that Jesus.
[Cantonese].
purposely entered into our brokenness..
[Cantonese].
He didn't just barely approach us and hold his nose to get away from the stench..
[Cantonese].
He didn't just come around the corners of our brokenness and mostly avoid it..
[Cantonese].
He dove straight into the middle of our human mess..
[Cantonese].
He saw us at our very worst..

$^{401}$[Cantonese].
And he loved us there..
[Cantonese].
And he knowingly walked straight into all of our chaos..
[Cantonese].
And that's good news..
[Cantonese].
That is such good news..
[Cantonese].
Because we humans are still a mess..
[Cantonese].
Our lives are still full of brokenness and chaos..
[Cantonese].
We have struggle and hurt and pain..
[Cantonese].
But Jesus showed us beyond a shadow of a doubt.
[Cantonese].
that our brokenness doesn't keep him away..
[Cantonese].
He's here to fix the worst that we could have to offer..
[Cantonese].
There's nothing too broken that he can't fix..
[Cantonese].
There's nothing or no one too dirty that he can't cleanse..
[Cantonese].
There's no pain or hurt too deep that he can't heal..
[Cantonese].
And there's nothing so dark that he cannot forgive it..
[Cantonese].
And that is good news..
[Cantonese].
News of great joy to all people..
[Cantonese].
That's why the birth of Jesus was such a big deal..
[Cantonese].
That's why heaven was celebrating that night..
[Cantonese].
Stars were moving in the heavens..
[Cantonese].
Angels were singing..

$^{441}$[Cantonese].
That's why we still celebrate..
[Cantonese].
Over 2,000 years later,.
[Cantonese].
we still talk about that night and we sing about it..
[Cantonese].
Because of what it meant..
[Cantonese].
Jesus stepped into our brokenness to bring healing..
[Cantonese].
Now there are just three points of application.
[Cantonese].
I want us to think about as we close..
[Cantonese].
First of all, perhaps there are some people here today.
who don't have relationship with Jesus..
[Cantonese].
Maybe you've never heard much of this story before..
[Cantonese].
Maybe you always thought Christmas was just about Santa.
and presents and ugly sweaters..
[Cantonese].
[Laughter].
Maybe even Christmas is a difficult time for you,.
something that you don't really like..
[Cantonese].
That's okay..
[Cantonese].
Maybe your life is messy..
[Cantonese].
The good news for you.
[Cantonese].
is that's why Jesus came..
[Cantonese].
If you're experiencing brokenness.
[Cantonese].
if there's difficulty in your relationship with family and friends.
[Cantonese].
or in your marriage or your job.

$^{481}$[Cantonese].
or your emotional or mental health.
[Cantonese].
or your finances or you're trapped in some type of addiction or pain.
[Cantonese].
the good news for you.
[Cantonese].
is that Jesus has chosen to show up.
[Cantonese].
right where you are.
[Cantonese].
in the middle of all the pain and brokenness.
[Cantonese].
and there's nothing he can't fix..
[Cantonese].
If you're willing to believe on him.
[Cantonese].
and follow him.
[Cantonese].
then he'll begin to heal your mess.
[Cantonese].
and bring you back into that healthy and right relationship with God.
[Cantonese].
that you were created for..
[Cantonese].
This is the good news for you..
[Cantonese].
Jesus has seen you at your worst.
[Cantonese].
and loved you..
[Cantonese].
Call out to him for help.
[Cantonese].
and if you're open to relationship with him.
[Cantonese].
he's waiting with open arms for you..
[Cantonese].
Second.
[Cantonese].
maybe you're here today and you already have relationship with Jesus.

$^{521}$[Cantonese].
maybe you've had that relationship for a short time or a long time.
[Cantonese].
but there's still some mess in your life.
[Cantonese].
just because we're a Christian, a follower of Jesus.
[Cantonese].
doesn't mean everything is easy.
[Cantonese].
we still struggle.
and have pain and hurt.
[Cantonese].
the good news for you.
[Cantonese].
is that the same Jesus who saved you.
[Cantonese].
wants to be at work in your life right now.
[Cantonese].
redeeming the mess.
[Cantonese].
fixing those broken things.
[Cantonese].
healing you.
[Cantonese].
turning the ashes of your failure and my failure into the beauties of his grace.
[Cantonese].
if you're in that position this morning.
[Cantonese].
maybe you're trying to hide some of those messy corners of your life in your heart.
[Cantonese].
from Jesus.
[Cantonese].
but if you're burdened and struggling.
[Cantonese].
cast your cares on him.
[Cantonese].
surrender to his work in your life.
[Cantonese].
and ask him for healing.
[Cantonese].

$^{561}$open every door of your heart to him.
[Cantonese].
and invite him into those places of hurt and brokenness.
[Cantonese].
and see what miracles he can work.
[Cantonese].
with your heart and with your life.
[Cantonese].
the third thing I want us to consider.
[Cantonese].
is if you are a person who's a follower of Jesus.
[Cantonese].
then I want to remind you.
[Cantonese].
that Jesus has given us an example to follow.
[Cantonese].
he left a place.
[Cantonese].
of perfect comfort.
[Cantonese].
and peace.
[Cantonese].
and tranquility.
[Cantonese].
to come here.
[Cantonese].
to come to this.
[Cantonese].
our brokenness.
[Cantonese].
and he did it out of love.
[Cantonese].
as his followers.
[Cantonese].
he's called us to do the same.
[Cantonese].
to carry his love.
[Cantonese].
into the messy places around us.
[Cantonese].

$^{601}$into the messy lives around us.
[Cantonese].
into the places that are broken and chaotic and hurting.
[Cantonese].
Jesus calls us to follow him in this way.
[Cantonese].
he calls us.
[Cantonese].
to help others.
[Cantonese].
come into that right and healthy relationship with God.
[Cantonese].
and that's also good news for us.
[Cantonese].
as his followers.
[Cantonese].
not only has he given us right relationship to God.
[Cantonese].
but he's given us purpose.
[Cantonese].
and a pathway to serve him.
[Cantonese].
and to love others.
[Cantonese].
so may each one of us be encouraged.
[Cantonese].
this joyful good news.
[Cantonese].
may we hear and receive the words that the angel shared.
[Cantonese].
do not be afraid.
[Cantonese].
I bring you good news.
[Cantonese].
great joy.
[Cantonese].
that's for all the people.
[Cantonese].
good news for you.
[Cantonese].

$^{641}$good news for me.
[Cantonese].
good news for our neighbors.
[Cantonese].
and may our voices rejoice with the angels.
[Cantonese].
as we look at what God has done for us.
[Cantonese].
and we say glory to God in the highest.
[Cantonese].
and on earth peace to men.
[Cantonese].
on whom his favor rests.
[Cantonese].
let's pray.
[Cantonese].
God I just thank you this morning.
[Cantonese].
for what you did.
[Cantonese].
when you left the comfort of heaven.
[Cantonese].
to step into our broken and messy world.
[Cantonese].
you did it out of love for us God.
[Cantonese].
God I pray that we would encounter your love this morning.
[Cantonese].
wherever we're at in life.
[Cantonese].
all of us have hurt.
[Cantonese].
all of us have fears.
[Cantonese].
all of us have brokenness.
[Cantonese].
may we encounter your love.
[Cantonese].
in those places this morning.
[Cantonese].

$^{681}$and may from this place.
[Cantonese].
may we carry your love.
[Cantonese].
to the hurting world around us.
[Cantonese].
in Jesus name we pray.
[Cantonese].
Amen.
[BLANK AUDIO].
\newpage



\section{}
\label{sec:J7N4kVhPD7E}
\textbf{6/10/18 主日講道: 這裡敬拜的是怎樣的人 (張麗仙傳道)}
\newline
\newline
連結: \href{https://youtube.com/watch?v=J7N4kVhPD7E}{\texttt{https://youtube.com/watch?v=J7N4kVhPD7E}} ~~~~ 語音日期: 2019-07-18
\newline
\newline
\hyperref[sec:5SLOPVCSwvk]{\small{< < < PREV SERMON < < <}}
~
\hyperref[sec:index_chronic]{\small{[返順時目]}}
~
\hyperref[sec:index_scriptual]{\small{[返順卷目]}}
~
\hyperref[sec:3CLlvNWl34w]{\small{> > > NEXT SERMON > > >}}
\newline
\newline
$^{1}$想問問大家.
中間有誰聽過小明的故事?.
聽過的話請舉手.
哎!.
弄得JIMBO嗎?.
那我是不是跟JIMBO這麼老呢?.
因為我的年代小明的故事是非常流行的.
因為我們的年代是很純情的.
沒有現在這麼多玩具.
老師一說故事就說.
從前有個小明.
就說說說說很多故事.
已經很久很久沒有聽過小明的故事.
但是最近WhatsApp.
竟然有人WhatsApp一些小明的故事給我.
而且是在一個中文課室堂那裡.
說的小明的故事.
非常之.
你告訴我怎樣.
我介紹幾個小明的故事給你聽.
中文老師說.
小明不是小明.
小明和所有同學.
中國字是這樣的.
凡是有三點水.
凡是跟水有關.
都是有三點水的.
所以江河湖海.
小明就說.
老師那麼沙漠呢?.
老師就說.
滾!.
小明就說.
老師我滾出去的時候都沒有看到有水.
然後老師和小明一班就說.
現在就是中文成語的時間.
你們說一下.
找一下成語就是.
如果你怎樣形容那個人的樣子.
很醜的樣子.

$^{41}$有人就說奇貌不揚.
這個就答得更加好了.
奇醜無比.
然後小明就說.
有點像你.
老師看著他.
OK 接著中文數據.
各位.
如果每一滴水都能代表一份祝福的話.
那我就送給你一片海洋.
大家請作句.
照這個句式來作句.
可以嗎?.
照這個句式來作句.
小明就說.
如果每一朵花都代表一份祝福.
我會送給你一塊花圈.
究竟小明的故事和我這次說的有什麼直接的關係呢?.
我告訴你其實是沒有的.
只不過是.
我接下來要和大家分享的聖經.
是會比較嚴肅一點.
所以開始的時候.
我也想你和我一起輕鬆一下.
那今天想和大家說哪一段的聖經呢?.
一段大家都很熟悉的.
我們一起讀一下好嗎?.
他們來到耶路撒冷.
進入聖殿.
說出殿內做買賣的人.
推倒兌換銀錢之人的桌子.
和賣甲子之人的底子.
也不許人拿著器材從殿內經過.
便教訓他們說.
經上不時記著說.
我的殿必稱為萬國禱告的殿吧.
你們倒視他便為賊窩了.
祭司長和文士聽見這話.
就想發誓要除滅耶穌.
卻又怕他.

$^{81}$因為眾人都欺其他的教訓.
今天耶穌在做什麼呢?.
祂在做的是將聖殿的桌子推倒.
將聖殿的椅子推倒.
叫人們不許走過.
祂要趕走買賣的人.
我想如果耶穌今天在我們城會教會做.
就知道嚴重性有多大.
把東西撥走.
把桌子上的東西全部放下.
究竟耶穌為什麼這麼激動呢?.
究竟我們知不知道.
為什麼耶穌這麼激動呢?.
我們知道.
祂一定是有很重要的原因在背後.
無論我們摸索的原因是否最正確.
我們知道背後一定有一個很嚴肅的原因.
耶穌會這麼激動.
所以今天我和大家一起重溫這個事跡.
看看耶穌為什麼這麼激動呢?.
因為耶穌這個行動有很多的解經家.
會引申這個理論.
是玷污聖殿的理論.
Impurity Theory.
是什麼意思呢?.
就是說因為這些買賣的人.
在聖殿買賣.
買到又叫又叫又叫.
所以搞到聖潔的聖殿.
搞到很不正常.
很不聖潔.
搞到聖殿好像一個街市一樣.
所以因為這個原因.
耶穌要這麼激動.
在聖殿裡結聖聖殿.
但是當我們看.
這個是聖殿.
當時的買賣是在聖殿的入口.
看到我圈住的那些.
好像橫著一排的房子嗎?.

$^{121}$那個是皇室處郎.
當時的買賣就在這個地方舉行.
當時的買賣就是在這個地方舉行.
並不是說整個聖殿都是充滿買賣.
所以我很難想像.
耶穌當時會這麼激動.
為了擺在前排房屋的買賣.
祂會這麼激動要趕走這些買賣.
還要結聖聖殿.
就像我們去威斯康辛看球賽.
進去後.
類似吧.
進去後有買東西吃的地方.
那些買東西吃的地方不是在球場裡舉行.
只不過是前面入口的地方舉行.
為什麼耶穌會這麼激動.
為了這些事而結聖聖殿.
這樣的行為呢?.
而且我們知道.
聖殿這些買賣祭物的活動.
一直都存在的.
是有它的需要的.
因為很多遠方來到的聖殿.
牽制的猶太人.
他們是不可能帶一個遠方帶來的祭物.
帶到聖殿的那個祭物是沒有病,沒有痛.
是無瑕疵的.
不可能.
所以他們一定要在聖殿買祭物.
還有聖殿兌換銀錢也是很需要的.
因為他們當時是需要聖殿特定的錢幣才可以買到祭物.
所以這個理論我覺得是有問題的.
第二.
很多時候就說.
耶穌這麼激動.
是因為當時聖殿的買賣很貪婪.
貪婪得利.
這個是Greed Theory.
耶穌會激動.
因為裡面的人的買賣很貪心.

$^{161}$兌換銀錢的人將兌換率提得很高.
買祭物的人將買祭物的價錢提得很高.
更不用說聖殿裡面的祭師.
管理聖殿的宗教領袖.
他們從這些買賣得利.
所以一切都漲價.
使遠道而來的牽制的人負擔很重.
所以不滿是極之有可能的.
耶穌是不滿這些貪婪買賣是極之有可能的.
再加上我們看到最後那段有一個說話.
耶穌說你們將聖殿變成賊窩了.
很多時候我們就會總結.
耶穌這麼激動.
是因為聖殿的買賣太過份了.
貪得太過份了.
不過這個理論也有它的漏洞.
我們看到耶穌當時激動趕走的人.
不單是買東西的人.
他連買東西的人都趕走.
我們看到他不單是推倒買東西的桌子和椅子.
他連在聖殿裡面拿著器具.
拿著這些牽制的東西走過的人.
他都要趕走.
似乎耶穌的激動是因為那裡的人貪心.
似乎耶穌是有一些東西要說.
就是似乎他要停止當時聖殿的牽制.
不准買不准賣不准牽制.
他似乎要freeze當時的牽制.
聖殿裡面的敬拜的人.
所以這個理論也讓我們看到它的漏洞.
你如何去解釋到耶穌要趕走買東西的人.
和阻止這些牽制的活動呢.
如果只是為了貪心的話.
所以在我們未探討下去以前.
我們要注意馬可這段記載.
耶穌的動作是有很多的.
耶穌是趕走聖殿裡面的人.
推倒兌換銀錢的桌子和賣甲子人的椅子.
不讓人拿器具從殿裡走過.
但是耶穌說的話就只有兩句.

$^{201}$經上五世記載說.
我的殿必要稱為萬國土國的殿嗎.
這是第一句.
第二句是你們卻把他變成賊窩了.
萬國土國的殿.
耶穌的激動是因為.
聖殿的活動佔據了愛邦人院的地方.
是不是呢.
很多聖經學者都認為是.
但是好像我剛才一開始已經說過.
當時的買賣是在聖殿的入口的牆廊裡面.
並沒有佔據整個愛邦人的院.
耶穌會不會因為這樣而這麼激動呢.
所以這個愛邦人的理論.
Gentile Theory.
我覺得也有它的問題.
不過當我們細心去想這個萬國土國的殿的時候.
我們發現原來也有問題的.
我們也要知道.
神是很體重猶太人.
他也很體重世界上所有的人.
愛邦人.
所以在二次亞書第五十六章.
神直接先知說過.
還有那些和我聯合的愛邦人.
要事奉我的話.
要愛我的名.
要作我的僕人.
就是凡事安息日.
不干犯又持守我約的人.
我必令他們到我的聖山.
使他們在禱告我的殿中喜樂.
他們的繁祭和平安祭.
在我壇上必蒙月立.
因我的殿必稱為萬民禱告的殿.
這是耶和華所說的.
耶穌所說.
我的殿必稱為萬國禱告的殿.
就是出自二次亞書第五十六章.
我們可以看到.

$^{241}$神是很重視愛邦人對他的敬拜.
也接納愛邦人對他的繁祭和平安祭.
他們對神的敬拜.
但是當時的聖殿是怎樣的呢?.
當時的聖殿有愛邦人院.
一進門口買賣的地方.
有一個長廊.
有愛邦人院.
愛邦人院前面有一個聖殿圍欄.
圍著聖殿的地方.
那些地方只有哪些人可以去呢?.
所有猶太人都可以去.
男女老少都可以去.
但是愛邦人就不可以去.
愛邦人可以說.
他們在開放區域.
他們的範圍就是去到那裡.
你從那裡走到那裡.
裡面這裡是封了的.
不可以進去的.
是不讓他們進去的.
不僅不讓他們進去.
還有一個牌寫明.
愛邦人止步.
進去的必要處死.
原來當時的聖殿有愛邦人院.
但是那個愛邦人院.
是愛邦人在那裡.
也是他遊客止步的地方.
他們不可以進到聖殿裡面.
敬拜獻祭.
神說.
我的殿必稱為萬國禱告的殿.
但是聖殿的宗教領袖.
他是違抗了神的心意.
禁止愛邦人進到獻祭和禱告的地方.
因此耶穌結證聖殿.
有可能是這個原因.
但是絕對不是只有這個原因.
因為記不記得耶穌說了兩句話.

$^{281}$我的殿必成為萬國禱告的殿.
你們卻使它變為賊窩了.
最後說接下來的一個原因.
我想提一下你對舊約聖經的認識.
你都會記得在舊約的時代.
當神要跟以色列說話.
要給他們一些警告的時候.
他是loud and clear的.
他不會含蓄.
他是會很直接很嚴肅地.
藉著先知去跟他們說.
而且有時候不單是用語言說.
還有時候用行動做.
所以這些叫做先知的行動.
Prophetic action.
舉個例子.
在以徐亞書20章.
神藉著以徐亞先知赤裸裸.
不穿衣服不穿鞋子三年的時間.
做什麼呢?.
是象徵埃及和古實的災難.
這個記載在以徐亞書20章.
另外在耶利米斯19章.
神藉著耶利米斯先知在城門打碎啞瓶.
象徵猶大.
如果你不聽從神.
你們敬拜偶像.
你們就好像啞瓶一樣的命運.
這些叫做Prophetic action.
是先知的行動.
當神要很嚴厲地將一個訊息告訴他們的時候.
有時候他會藉著先知很戲劇化.
好像很極端的行動.
來講這個訊息.
所以當我們去了解耶穌.
潔淨聖殿這個激動的行為的時候.
有些聖經學者就說.
這是一個先知的行動.
這是一個Prophetic action的行動.
因為耶穌要表達神一個非常之不滿.

$^{321}$所以祂要做這個行動出來.
這個行動是代表神什麼的不滿呢?.
我們今天2000年之後.
可能對於這個行動.
為什麼耶穌會這麼激動呢?.
會很不明白.
想了很多很多原因都不明白.
但是當時在場的猶太人一看就明白.
因為他們對舊約的聖經很熟悉.
當耶穌說你們將聖殿變成賊窩了.
他們就明白耶穌當時為什麼這麼激動.
所以這段聖經最後說.
祭司和文士聽到耶穌說.
看到耶穌的行動.
同時想法子除滅耶穌.
但是又怕祂.
因為眾人都欺負祂的教訓.
證明當時的人一看就明白.
耶穌這個潔淨聖殿是代表什麼的意思.
但是到了我們今天.
是有些困難去明白.
但是以下一段聖經可以幫我們去明白.
耶利米書第七章一至十一節.
那裡記載說.
耶華說臨到耶利米才知道.
祂說你是站在耶華的殿裡面.
在那裡宣傳我這一番話.
你要進到那裡.
對敬拜耶華的一切猶大人.
湯吞耶華的話.
究竟耶華的話是什麼呢.
為什麼要這麼鄭重地叫耶利米.
才知道去到聖殿的門.
要跟他們說這些話呢.
耶華的話就是.
萬貫之耶華,以色列的神如此說.
你們要改正行動作為.
我就使你們在這地方安然居住.
你們不要依靠虛方的話說.
這是耶華的殿,這是耶華的殿.

$^{361}$這是耶華的殿.
你們若實在改正行動作為.
在人和靈社中間.
盛宴施行公平.
不欺壓寄居的孤兒寡婦.
在這地方不流無辜人的血.
也不隨從別臣陷害自己.
我就使你們在這地方安然居住.
就是我古時所賜給你們滅灶的地.
直到永遠.
這是耶華的話.
你也看得很清楚.
當時的猶太人是怎樣.
耶華說兩番跟他們說.
你們要改正你們的行動作為.
你們的行動作為是甚麼呢.
你們時時欺壓孤兒寡婦.
你們時時敬拜偶像.
你們時時在靈社中間行不公平的事.
這些等等是你們當時的行為.
是激怒神的行為.
如果你們改變這些.
我會使你們在這地安然居住.
但是繼續.
耶華說.
但是你們偷渡,殺害,姦淫,起假誓.
向霸力燒香並隨從數不認識的別臣.
但是你們且來到這個稱為我名下的殿.
在我面前敬拜.
有說我們可以得到自由了.
你們這舉動是行那些可憎的事嗎.
這稱為我名下的殿在你們眼中豈可看為賊窩嗎.
這句說話原來出自耶利米書七章一讀.
藉著這一段說話.
我們會清楚一點耶穌當時的憤怒是甚麼.
我們也知道當時猶太人怎樣一聽就明白.
因為耶利米書已經是清清楚楚的.
是陳列出當時猶太人怎樣在聖殿裡面激怒耶和華.
他們在外面犯了很多壞事.
已經列出來了.

$^{401}$他們以為回到聖殿就若無其事.
只要我獻祭,敬拜神.
就甚麼都不計較了.
他們是這樣以為.
他們在外面做盡傷天害理的事.
但他們回到聖殿.
聖殿是他們的hideout.
是他們的安樂窩.
或者我們一向對賊窩這個字有少少誤解.
我們以為賊窩就是在聖殿裡面做買賣的人.
在那裡賊購.
賊窩的意思.
原來我們看清楚.
不是指正在做壞事的地方.
是指賊做完壞事.
回家休息去享受賊贓的安樂窩.
這個才是賊窩.
Den of Roberts.
是賊的溫巢.
原來當時的猶太人.
他們在外面做盡很多不好的事.
違反神旨意的事.
是一些很傷天害理的事.
但他們以為回到聖殿裡面.
就甚麼問題都沒有.
只要他們獻祭和敬拜.
聖殿就是他們的安樂窩.
神卻說.
你看錯了.
你以為錯了.
這稱為我名下的殿.
在你們眼中豈可變為賊窩呢.
神非常氣憤.
當日要藉著耶利米先知去到聖殿門口.
用先知行動.
Loud and clear.
警告以色列人.
他們不可以再這樣做.
他們要快點悔改.
他們不要想錯.

$^{441}$以為你在外面犯罪.
聖殿可以成為你的避難所.
成為你的安樂窩.
你是錯的.
難怪.
我們看到.
耶穌當日的激動.
潔淨聖殿.
祂當時並不是大發脾氣.
我們不要看錯.
耶穌不是發脾氣.
Lose temper 失控制.
耶穌這個行動.
是神設計的一個先知的行動.
為了指責當日的信徒.
和宗教領袖.
因為他們犯了當日.
舊約時代.
猶太人所做的事.
他們在外面犯罪.
以為回到聖殿.
在這裡敬拜獻祭.
就什麼都沒有發生.
在這裡.
耶穌要宣告神的審判.
祂要趕走買賣的人.
祂連拿著獻祭器具的人.
在聖殿準備獻祭.
祂也要趕走.
全部停止.
我要停頓了聖殿的獻祭.
我不再接受這種獻祭.
因為你們毀了它.
我要廢除這種獻祭.
我甚至要廢除這個聖殿.
所以有些聖經學者.
稱耶穌潔淨聖殿.
為「廢除聖殿」.
而當時耶穌這個震驚的先知行動.
就像電影一樣.

$^{481}$但它又不像電影一樣.
因為電影.
The End.
就是什麼都完結.
但耶穌這個先知行動.
這個審判.
才是開始.
今天跟大家重溫了.
這一段耶穌潔淨聖殿的事.
究竟對我們城會教會.
對在座每一位弟兄姊妹.
你有什麼提醒呢.
對我們大家一起.
有什麼的省察呢.
讓我們一起在神的面前.
真誠地去省察.
究竟你有沒有像當日的宗教領袖.
和猶太人一樣.
禮拜一至禮拜六.
在教堂以外.
隨從私慾地去活.
浸淫在一些違背神旨意的生活中呢.
究竟有沒有在家.
在工作的地方.
甚至在你的朋友圈.
做一些欺壓別人的事呢.
有沒有自私自利地.
偷取了別人的東西.
那些不是屬於你的東西呢.
有沒有在名利.
色這些事上犯罪呢.
有沒有縱容自己在某些玩意.
享樂中.
虛度了神給我們的恩賜.
給我們的光陰.
同時也傷害了我身邊愛你的人呢.
更重要的就是.
你有沒有敬拜偶像呢.
在你的生活中.
有沒有.

$^{521}$你已經敬拜了很多偶像.
但你的位置是高於神的.
我們是不是來到教堂這裡.
每個星期日.
照正常的奉獻.
獻祭.
敬拜.
就日無事了.
如果我們是這樣的話.
恐怕我們不比當日的宗教領袖.
當日的猶太人好.
因為我們這樣的話.
我們也是將神的電變為賊窩.
變為我們犯罪後的安樂窩.
上星期Mark Sir在這裡.
強調到恩賜事奉.
今天的節奏上也提到恩賜事奉.
沒錯.
當看到弟兄姊妹.
很樂意地參加事奉.
使用你的恩賜.
其實我們是很開心的.
但是我們更加要認定.
我們的神看重的不是多少人事奉.
祂也不是看重事奉的本身.
祂是看重我們每一個事奉的人.
究竟我們有沒有審查.
我們生活的行為.
以致我們來到神的電裡.
敬拜事奉祂.
沒錯.
耶穌潔淨聖殿是一個審判.
但是我們千萬不要被這個審判的訊息.
嚇得我們不敢再來到聖殿敬拜和事奉.
因為這不是耶穌的心意.
不是耶穌審判的心意.
耶穌的審判.
接下來帶來的是拯救.
是福音.
我們要相信我們的神是慈悲憐憫的神.

$^{561}$只要我們真誠承認我們的過犯.
盡力的改善.
祂的怒氣不過是轉眼之間.
祂的恩典卻是一生之久.
事實上在耶穌的審判之後.
過幾天馬可福音記載.
就說耶穌在幾天之後.
晚上和門徒坐在一起.
和他們吃午餐晚餐的時候.
耶穌拿起杯子來祝福了.
就張開了給門徒說.
你們要拿著吃.
這是我的身體.
又拿起杯子來說.
這是我的血.
這是很多人流出來的.
因為耶穌知道我們每一個人的罪性.
舊的憲制雖然無效.
但是祂因為愛我們.
祂為我們附上一個新的憲制.
就是祂自己.
我們什麼時候願意獻上一顆真誠懺悔的心.
我們什麼時候就得到.
耶穌這個新的憲制帶來的赦免.
一開始的時候.
我跟大家說小明的故事.
還記得嗎.
搞笑該死.
輕鬆一下.
有時候是很需要的.
但是耶穌潔淨聖殿的故事.
不是搞笑該死.
相反很嚴肅.
但是對我們是更加需要的.
藉著今天的訊息.
我勸勉你.
我們做信徒.
不只是在星期日.
而是更加在星期一至星期六.
我們要懷著愛神愛人的心好好地去活.

$^{601}$我們的星期日回來這個敬拜.
才是真誠的敬拜.
我也勸勉你.
當發覺自己拼命都做得不夠好的時候.
你要懷著一個謙卑的心.
依靠主耶穌在十字架上.
用寵教為你所附上的新的憲制.
這個最純潔最有效的憲制.
是屬於你的.
我們簡單祈禱.
潔淨聖殿這個行動.
讓我們醒察我們今天對你的信心.
我們今天生活的行動.
藉著你自己的憲制.
給我們有新的機會.
雖然我們今天做得不夠好.
但永遠在你裡面.
只要我們懷著真誠悔改的心.
我們永遠是有得救.
以我們城會教會每一位弟兄姊妹.
用恩主的話來醒察我們的信心.
我們做一個更合神心意的兒女.
多謝你聽我們禱告.
奉救主耶穌的名.
阿們.
\newpage



\section{}
\label{sec:3CLlvNWl34w}
\textbf{4/21/19 主日講道:初熟的果子(陳崇基牧師)}
\newline
\newline
連結: \href{https://youtube.com/watch?v=3CLlvNWl34w}{\texttt{https://youtube.com/watch?v=3CLlvNWl34w}} ~~~~ 語音日期: 2019-07-19
\newline
\newline
\hyperref[sec:J7N4kVhPD7E]{\small{< < < PREV SERMON < < <}}
~
\hyperref[sec:index_chronic]{\small{[返順時目]}}
~
\hyperref[sec:index_scriptual]{\small{[返順卷目]}}
~
\hyperref[sec:OTRUSy4_3CA]{\small{> > > NEXT SERMON > > >}}
\newline
\newline
$^{1}$Happy Easter.
各位 福祉快樂.
我們可否給桌上的人.
給他們一個High Five.
Happy Easter.
如果你們記得在星期五.
星期五.
OK.
兩三天前.
耶穌死亡的最可怕的死亡.
是基督的死亡.
祂被懲罰.
被判刑.
祂的身體被困在黑暗的寶寶墓.
所有的希望都失去了.
而是 邪惡的 死亡 都獲得了勝利.
在今天第三天.
早上非常早.
女性 女性的徒弟們.
他們去到耶穌的寶寶墓.
他們想參觀祂的身體.
他們來到寶寶墓.
他們發現寶寶墓已經打開.
寶寶墓已經空了.
耶穌已經不在了.
別害怕 聖經說.
祂不在了.
祂已經從死亡中出生了.
在這時.
女性們記得.
耶穌已經告訴她們.
很久以前 重複地說.
祂會被祭祀.
祂會從死亡中出生.
在第三天.
記得這一點.
女性徒弟們去告訴.
在藏身的徒弟們.
他們告訴他們好消息.
耶穌出現了.

$^{41}$祂的徒弟們 祂的家人.
並且顯示了 祂已經從死亡中出生了.
耶穌呼喚了祂的徒弟們.
以及所有的世代的教會.
來祈求生命的希望.
耶穌已經擊敗了死亡.
耶穌已經擊敗了邪惡.
耶穌已經擊敗了一切.
祂就是勝利.
這就是所有的世代的信息.
他們都在祈求.
在聖誕節我們也在慶祝這個事實.
我們也在活著這個呼喚.
我們被神呼喚.
去到外面 去告訴一個完全不同的故事.
去看到一件完全不同的事.
去看到現實 去看到我們的生活.
上週 巴黎發生了大火.
納德達教堂也被火燒了.
這令人心痛.
整個世界都在悲傷.
這令人特別痛苦.
因為上週 本來是聖誕節.
聖誕節是在聖誕節中.
聖誕節是在聖誕節中 讓整個基督教世界都記得受苦.
和耶穌的痛苦.
看到大火燒掉了教堂.
我們都感到憂慮.
失去了文明.
失去了所有的文明和建築.
失去了歷史.
我們都看到大火和毀滅.
我們都感到不朽.
我們都看到大火 和毀滅了教堂.
我們都感到不朽.
我們都看到大火 和毀滅了教堂.
我們都感到不朽.
我們都看到大火 和毀滅了教堂.
我們都感到不朽.
我們都看到大火 和毀滅了教堂.

$^{81}$我們都感到不朽.
我們都看到大火 和毀滅了教堂.
我們都感到不朽.
Everything is gone and destructed. All you see is smoke and ash of destruction. But behind this is a golden cross..
This is how a good photographer can tell a totally different story, a totally different perspective to see reality..
High above the order, high above the ash of destruction is a powerful reminder of Christian hope..
This is the hope that we have with the empty tomb of Christ..
And this is the hope we can see and experience in the midst of destruction, in the midst of ash of fire..
Our hope is on the empty cross. Christ is not there. He died on Friday. He was buried. He was in the tomb. And His tomb is empty..
Our hope is on the empty cross to raise in Christ. And please don't forward this picture and say this is a miracle..
So gold has a much higher melting point than wood. This is not a miracle..
But however, this is a totally different way to see reality..
Easter Sunday calls us to see reality with a different lens..
Put on a different eye glasses. Put on a different lens, a spiritual lens to see life..
If life is just here and now, death and destruction have triumphed..
This is what you get. Ash of destruction, forces of evil, death and Satan..
But it shows us a higher reality, a different way, a totally alternative way to see reality..
And that is the eternal lens. With Jesus' resurrection, the forces of evil, evil authorities, death and Satan have been defeated..
With Jesus' resurrection, new creation, the Bible says, is set in motion..
All in Christ, all people who put their trust in Christ, will share the same victory of Jesus..
1 Corinthians 15.20. It says, "But in fact, Christ has been raised from the dead..
He is the first of great harvest of those who have died in Christ..
Christ's resurrection is the first fruits.".
That's an interesting term. It's a farming term..
It means joy and hope with the first harvest..
And Jesus' resurrection is the first harvest, filled with joy and gratefulness..
You know, you look around our world, it's filled with sin and suffering and death..
You probably caught the news story in Sri Lanka this morning..
There was a series of bombing of Christian churches as they celebrated Easter Sunday..
More than 100 people died. Hundreds of people are wounded..
This is our world. As of destructions..
But Easter Sunday marks a totally different way to see the world..
It celebrates a new beginning..
The first fruits of the harvest will guarantee the rest of the field..
There will be harvest. There will be a new beginning..
So you come to Christ today. You put your trust in Him..
With the resurrection, you have new hope..
And you are new. New beginning and new creation..
Our church launched three years ago on Easter Sunday..
With the same conviction, same hope that we have in the raising Christ..

$^{121}$We launched this church and this new mission..
Why we do that? Because we know the tomb of Jesus was empty..
There is hope in resurrection. There is a totally different new way to see our lives..
And because of the hope of resurrection, we have a new mission..
We are new people..
So I pray that may the resurrection of Jesus encourage us all to continue to build the church..
To plant the church. To extend the mission of God. To continue in our faith..
He said in the book of Acts 15, verse 20, "But Christ is already in death.".
And Christ's resurrection is the beginning of the people who have been sown..
Who said the people who have been sown?.
It must be the church at that time. A new church..
Some people are suffering. Some families have left this world..
But they are the people in the tomb. They have left the world..
What are they doing? They are in despair..
Paul wrote this letter to teach them..
We still have despair. Because we have a desire for resurrection in Christ..
Paul said a lot of things about the fact of the resurrection of Christ..
He came to this chapter and said, "The resurrection of Christ is like the fruit of the first ripening.".
It is a word of the agricultural field..
You know, at this time, in spring, the weather in Barcelona,.
This is the time of the ripening of wheat..
It is the first batch of ripening. The first batch of fruit will not be produced..
As the first batch of ripening, the first ripe fruit,.
There must be a lot of joy and gratitude..
Because after the rainy season, there are several months of very cold rainy season..
At this time, it starts to be messy, just like Chicago..
It was snowing last week..
This is the time of the first batch of ripening..
People will be full of hope and gratitude..
And what does the first ripe fruit say?.
It is the reason for despair..
Why can there be a reason for despair?.
Because without the first batch of ripening, the rest of your crops will be ripened..
If there is no first batch of ripening, there is no guarantee that the rest of your crops will be ripened..
Jesus Christ is the first real ripening..
The first ripe fruit..
Really resurrected from death and will not die again..
He has completely won over death..
Jesus's resurrection is also a guarantee that all of us who believe in Jesus and rely on Jesus to the end..
We will have eternal hope in the future..

$^{161}$The main purpose of the second batch of ripening is to ask us to look at the world with a completely different perspective..
Look at our life..
Because of the resurrection, God will show a more real existence..
This is heaven..
This is God's kingdom..
The picture you just saw is the one from the St. Martin's Cathedral in Paris last week..
People all over the world are very sad..
Lost a lot of history and heritage..
Of course, it doesn't matter to some people..
The most important thing is that they are rich..
They don't appreciate culture, art, or history..
But there are many reasons for the destruction of heritage..
What everyone sees is just ashes and fire..
We think it's helpless..
We think there's only some destruction..
But a good photographer,.
When using a different lens to see the truth,.
There is a completely new concept..
This picture is after the fire..
A photographer went to the scene and took the first picture..
A good photographer uses a different lens to see the truth..
What he sees is not what ordinary people see..
What ordinary people see is just ashes and fire..
But he saw a cross behind..
If you can understand English, I have already mentioned it..
Please don't say this is a miracle..
Or you will be laughed at..
Because the melting temperature of gold is higher than that of wood..
So it's not hard to burn the gold cross..
It's very reasonable..
But this picture shows that we can see our real life from a new perspective..
What does the resurrection ceremony do?.
The empty tomb is actually to shake us..
To re-examine how we see the world..
Change a superhuman eye to a superhuman lens..
Change another pair of lenses to let you see God and God..
A spiritual lens..
The 15th chapter of the Gospel of John.
Proposes a lot of evidence of resurrection..
Paul said that if there is no resurrection of the dead,.

$^{201}$If there is no resurrection of Jesus Christ,.
Our life is just a closed system..
There will be no external things..
It's a dead end..
The last thing we will see is the end..
Because the last thing we will see is the final enemy..
That is death..
He said we only have this life of hope..
But we are so poor..
This is the special thing about the Christian faith..
If you pay attention to the Christian faith,.
The special thing is that it is built in a historical event..
It's not a theory of abstract..
It's not a very solid philosophical thought..
The Bible never provides a very solid philosophical thought..
But the Bible tells us that the Christian faith.
Is built in a historical event..
The death of Jesus Christ on Friday..
The resurrection of Jesus Christ today.
Is a historical event..
This is the most unique thing..
Because of the resurrection of Jesus Christ,.
We can see this world clearly..
It's not just a closed system..
The resurrection day forces us to realize that our life needs to be changed..
There is a higher reality..
God really exists..
He created the world..
And loves us..
He entered our world and saved us..
These sinners..
Because the resurrection of Jesus Christ,.
Think about it..
Life, human life, the last one who speaks in the final stage of salvation,.
Is no longer death..
It's resurrection..
It's despair..
It's love..
It's God..
It's not an evil regime..

$^{241}$It's eternal life..
So the resurrection of Jesus Christ from death.
Brings a lot of changes to many people's lives..
From the first generation of believers,.
Those followers were greedy and afraid of death..
Then they were willing to dare to come out and spread the resurrection of Jesus Christ..
All of them believed to death..
Another follower,.
Changed 180 degrees,.
Is Paul..
He was originally a person who persecuted the church..
Now he is willing to.
Because he saw the resurrection of Jesus Christ..
He understood this and he is no longer afraid of authority..
Then he went to all over the world to build a church and spread the gospel..
Today, in your life,.
You may encounter different difficulties..
I want you to bring the desire of the resurrection of Jesus Christ.
Into your life..
God can direct you to show his own work..
You are the one who brings the resurrection of Jesus Christ into your life..
You can tell people.
This world and life can be completely different..
We will pray later..
We will pray for the brothers and sisters who are in trouble..
We will pray for our friends and brothers and sisters..
We will pray for each of us..
We have a mission..
We have the desire to bring the resurrection of Jesus Christ..
We go back to our home..
Go back to our life..
Go back to the place where you work..
You can see the place where you don't know Jesus Christ..
You can tell people that Jesus Christ is really resurrected..
How to be resurrected?.
Not only a historical fact, but also an experience in your life..
Your attitude is completely different..
This is your mission..
We also pray for the church..
Our church says that today is the third year of the baptism..

$^{281}$We are still very young..
We still have a lot of needs..
We still have a lot of places to be filled..
But the desire to be resurrected pushes us forward..
Let us know that what we have done is of eternal value..
We will not hate this..
We are going to pray together..
Pray for those who are weak..
We know there are many friends around us..
Brothers and sisters in Christ..
They are facing some struggles..
We are going to remember them..
In our prayers..
We also pray for the church..
This church is found on the hope and faith of Jesus' resurrection..
We want to pray for the church because we want the same hope to motivate us..
So that we can go back to our home, our workplaces, our neighborhood..
To witness and to testify the living hope of Jesus Christ..
Let us all pray together..
我們請司見上來帶領我們祈禱..
\newpage



\section{}
\label{sec:OTRUSy4_3CA}
\textbf{4/14/2019 主日講道:棕樹節主日-「榮進耶路撒冷」(陳崇基牧師)}
\newline
\newline
連結: \href{https://youtube.com/watch?v=OTRUSy4_3CA}{\texttt{https://youtube.com/watch?v=OTRUSy4\_3CA}} ~~~~ 語音日期: 2019-07-20
\newline
\newline
\hyperref[sec:3CLlvNWl34w]{\small{< < < PREV SERMON < < <}}
~
\hyperref[sec:index_chronic]{\small{[返順時目]}}
~
\hyperref[sec:index_scriptual]{\small{[返順卷目]}}
~
\hyperref[sec:zcyVrDXZup0]{\small{> > > NEXT SERMON > > >}}
\newline
\newline
$^{1}$我們可以聽到一些這樣的心靈見證啊!.
我們這位主呢,是一位創天造地的神!.
聖經裡面講得很清楚,上帝從起初創造這個世界,.
雖然我們被創造,我們是人類啊!.
聖經裡面的記載,說我們呢,作為人,.
我們是被創造的,我們是人類啊!.
聖經裡面的記載,說我們呢,作為人,.
我們很多時候都不理會上帝,.
那個去創造我們,去供應我們每一天的生命,.
每一天我們的需要,這一位的上帝,很多時候我們不理會他,.
甚至我們去反對他,不理會他,.
但是這位的主,就是一位那麼愛我們的神,.
不單止他有權能去創造萬有,.
他也有愛來接近我們,來救贖我們,.
我們一班罪人,在我們中間每一位,.
我們生命裡面有很多破碎,有很多過犯,.
有很多我們心思行為所作的,.
我們可能連自己,打開自己的心,.
去看清楚我們裡面,生命裡面,.
每一個細節的時候,我們都未必能夠去接受到我們自己,.
但是這位愛我們的主,就是用一個恩典,.
一個無條件的愛來愛我們,.
剛才我說過,這是基督教裡面最重要的一個禮拜,.
因為這個禮拜,是去記載耶穌基督在歷史裡面,.
人類的歷史裡面,他為我們罪人做了什麼,.
今天所謂的中樹節主日,.
可能你聽到的時候會有一點陌生,.
什麼叫做Palm Sunday,.
什麼叫做Palm Branches Sunday,.
或者中樹節的主日,.
其實很簡單,.
聖經說到耶穌基督在地上最後的七天,.
這一天就是所謂的進入耶路撒冷的日子,.
這一天叫做中樹節主日,為什麼呢?.
因為在這一天裡面,.
很多時候,連今天,.
每一次到了這一天,Palm Sunday的時候,.
在耶路撒冷這個地方,以色列這個地方,.
人們就會在一個橄欖山上聚集,.
有很多基督徒,還有很多朝聖者,.

$^{41}$他們就會拿著這些中樹枝,.
在猶太人的世界裡面,中樹枝是代表勝利的意思,.
聖經是有這樣記載,.
說耶穌在最後的一個星期,.
主日裡面,隨著很多朝聖的人,.
還有他的門徒,.
有很多的群眾,他們進入耶路撒冷,.
這個就叫做Palm Sunday,.
聖經是怎麼記載的呢?.
在馬可福音第十一章第一節,.
耶穌和門徒將近耶路撒冷的時候,.
到了伯伐奇和伯大利的村落,.
在橄欖山那裡,.
耶穌就打發兩個門徒對他們說,.
他說你們去對面的村莊裡面,.
一進去的時候,就會看到一匹的驢驢,.
綁在那裡,.
從來沒有人騎過的,.
一隻驢子,可以解離,然後牽來,.
若有人問你們為什麼要做這些事呢?.
你就對那個所謂的驢主說,.
驢子的主人說,主要用他,.
那人必立時讓你們牽來,.
那些門徒去了之後,.
他們就看到一匹驢驢,.
在門外街道上,.
就將它解開了,.
在那裡站著的人,.
那裡有幾個就說,.
他說你解開驢驢幹什麼呢?.
門徒就照著耶穌的吩咐去回答他,.
那些人就被他解開,.
拿了這隻驢子,.
他們把驢驢牽到耶穌那裡,.
把自己的衣服搭在上面,.
耶穌就騎上這隻驢子,.
有很多人就將衣服鋪在地上,.
也有人將田間的樹枝,.
那些棕樹枝撼下來,.
然後鋪在路上,.

$^{81}$前行後退的人都喊著說,.
和煞那,和煞那,.
奉主名來的,是應當稱重的,.
那將來來的,.
我祖大位之國,.
是應當稱重,.
高高在上,和煞那..
這就是那個棕樹節主日的記載,.
《馬可福音》是其中一段,.
四本福音書都有講這個事蹟,.
你可以想像一下,.
在耶路撒冷這個地方,.
是一個很小的地方,.
每逢這個時候,.
每年按月曆來計算,.
這個時候就是猶太人的御節的時候,.
他們的Passover,.
Passover的時候,.
就是壽難節和壽苦節的時間,.
每逢這個時候,.
就會有很多朝聖的人,.
會去到耶路撒冷,.
他們會擠在城裡面,.
然後住在裡面的客店,.
然後在城外面,.
他們可以紮營,.
然後在這裡等候,.
去守這個最大的節日,.
有一點像我們華人的新年團聚,.
這樣的意思,.
他們在Passover的時候,.
他們會發現,.
曾經在歷史裡面,.
神拯救他們離開了埃及,.
一個他們做奴隸的地方,.
被人強權壓迫的地方,.
上帝藉著摩西拯救他們,.
然後歷史歷代他們就去慶祝這個所謂的Passover,.
就是因為神為他們所作的事情,.
救他們出來埃及,.

$^{121}$他們就有自由,.
所以在這個時間,.
中樹節主日,.
這個星期之前,.
就有很多人擠擁在耶路撒冷當中,.
是很熱鬧,.
而且是很擠擁的情況,.
那些人看到耶穌進來城的時候,.
他們就知道,.
這個耶穌不單止有名聲,.
很多人知道祂,.
有很多人跟隨祂,.
並且他們想認定,.
這個就是他們等了很久要來的國王,.
這個國王會幫助他們,.
因為當時你知道,.
在耶穌的時代,那些猶太人,.
他們是被羅馬帝國的強權佔領著,.
一直被人欺負的時候,.
他們就很希望,.
上帝在舊約當中,.
在聖經裡面曾經預言,.
幾百年前預言過,.
有一位的王會來到,.
來給他們自由,.
再出一次埃及,.
所以耶穌到耶路撒冷的時候,.
他們就發現,.
這個就是我們的國王,.
所以他們就高呼,.
來講什麼呢?.
和撒拿,和撒拿是什麼意思呢?.
在希伯來文,和撒拿就是Hoshe Anar,.
就是說,.
Please save me,.
Please save us,.
就是求上帝來拯救,.
是一首家傳戶曉,.
在猶太人當中,.
歷史裡面經常唱的一首詩歌,.

$^{161}$詩篇118篇,.
每逢Passover,.
每逢這個時間,.
他們就會唱起這首歌,.
因為他們去請求上帝來拯救他們,.
從一種困境,.
從一種壓迫裡面去拯救他們的時候,.
因為上帝是信實,.
一定會履行祂的諾言,.
所以這個和撒拿的意思,.
就慢慢變成讚美上帝的意思,.
因為上帝是值得讚美,.
因為在創天造地的主是會拯救我們的,.
所以他們就唱出這首詩歌,.
當耶穌進入耶路撒冷的時候,.
他說這個國王就是上帝預言會差派來到我們中間,.
他們現在在哪裡呢?.
在這個故事裡面,.
他們就在鑒欖山上面,.
我們看這張照片,.
就是在鑒欖山上拍的,.
你見到拍到對面,.
有城牆的地方就是耶路撒冷,.
一個很小的地方,.
一公里乘一公里,.
一個正方形的城市,.
這個就是歷史歷代上帝做事的地方,.
你見到對面有一個金頂,.
見到嗎?.
有一個比較小的,.
有一個圓形的金頂,.
那個叫做Dune of the Rock,.
盤石清真寺,.
那個是主耶穌之後,.
七百年才建的,.
一個回教的清真寺,.
然後你看到一個城牆,.
清真寺外面有一個長方形的城牆,.
那個城牆就是耶穌的時候,.
見到的聖殿山,.

$^{201}$是一個很大的,.
用石做的一個平台,.
大到二十個美式足球場那麼大,.
那個就是世界歷史的中心,.
就是上帝要在人類歷史裡面成就的東西,.
就是由這裡開始,.
也是當時古代世界最中心的一個地方,.
你看到這個圖畫,.
你看過去的時候,.
你就會發現,.
我們現在在橄欖山上面,.
看過去對面那座東西,.
那裡就是當時的聖殿耶路撒冷,.
其實如果你這樣下山,.
再上去那個城那裡,.
其實只需要十幾分鐘,.
你相信我,.
我經常帶隊去以色列,.
我曾經在山腳跑上去找人,.
因為走失了一些人,.
所以又跑回山腳找人,.
我十幾分鐘就搞定了,.
你看看我,.
我平時也不跑步,.
那些跑步慣的,.
可能十分鐘,.
七分鐘就搞定了,.
其實是非常近的,.
在這一處,.
耶穌就騎上路進入耶路撒冷,.
你想一下,.
耶穌在北面加利利一個地方來,.
他已經走了百多里路,.
還有那些跟隨他的人,.
那些朝聖者,.
都走了百多里路,.
走到最後,.
剩下十分鐘,.
十五分鐘,.
去對面就到耶路撒冷,.

$^{241}$為什麼在這個時候就要騎路呢?.
你明白我的意思嗎?.
今天你來教會,.
你沒理由去到街口那裡,.
然後你去叫Uber,.
是不是?.
你不可能在這個時候,.
你就會叫車,.
然後上車,.
就來到教會門口,.
是不是?.
所以他做的一定是有他的意思在這裡,.
為什麼最後的時候,.
他要安排一隻Donkey,.
一隻驢,.
然後騎著驢進入耶路撒冷,.
有四樣東西值得留意,.
第一樣東西,.
就是有中樹之人,.
耶穌的時代大概第一世紀,.
百多年前,.
猶太人曾經得過一次自由,.
他們打退了很多的希臘帝國的軍隊,.
當他們擊倒這些軍隊之後,.
他們就開始一個叫馬加比的時代,.
大概有一百年,.
他們是獨立,.
能夠自主的,.
那個元帥叫猶太馬加比,.
當他擊退了希臘的將軍,.
軍隊之後,.
他就騎著戰馬進入耶路撒冷,.
潔淨聖殿,.
然後他說,.
神的國就到了,.
誰知,.
建立了他的國度,朝代之後,.
一百年之後,.
他們就自己腐敗,.
然後就鬼打鬼,.

$^{281}$一百年之後就羅馬帝國去征服他們,.
曾經猶太人就想起這個Good old days,.
如果好像馬加比的時候就好了,.
所以當這個馬加比元帥進入耶路撒冷,.
擊退了希臘帝國的軍隊的時候,.
那些群眾就拿起中樹枝,.
來歡迎他,.
你看那些猶太人的硬幣,.
那時候丟上去的記號,.
就是中樹枝的記號,.
是以色列一種本地的土產樹,.
那些Palm Tree,.
他們就舉著,.
仰著,.
放在地上,.
來歡迎騎著戰馬進來的猶太人馬加比,.
但是在這個故事裡面值得留意的,.
就是耶穌騎上的不是戰馬,.
他騎上的是驢,.
是驢仔,.
這個就讓那些認識聖經的人,.
想起聖經裡面預言的人,.
他就一定會想起一章聖經,.
叫做撒迦利亞書第九章,.
那裡說什麼呢?.
那裡是幾百年之前,.
曾經預言,.
神的王會到,.
會統治我們,.
他說錫安的民,.
就是耶路撒冷的民,.
他說應當大大歡喜,.
耶路撒冷的民應該要歡呼,.
看你的王來到你們的樹立,.
他是公義的,.
是會私人拯救的,.
但是他是謙謙和和的,.
騎著驢,.
也會除滅戰車和戰馬,.
所以這個神預言會來的君王,.

$^{321}$就不是騎戰馬進聖城,.
他是騎著一隻謙謙和和的驢,.
進去耶路撒冷,.
我們現在繼續看下去,.
你就知道在這七天裡面,.
耶穌發生在他身上的事情,.
絕對不是一個軍隊指揮官會做的事情,.
他沒有刀,沒有槍,沒有戰馬,.
他是謙謙和和的,.
來到去教導人,.
然後被人釘在十字架上,.
受很大的痛苦,.
是他自己願意這樣做的,.
絕對不是耶穌進去以為群眾都會接納他,.
誰知道原來搞錯了,.
然後人們就將他釘死了,.
原來耶穌的死是一個意外的死,.
絕對沒有這樣的事情,.
聖經的記載告訴我們,.
他出奇的地方就在這裡,.
騎著驢進去耶路撒冷的時候,.
原來是他自己親自願意去捨棄他的生命,.
這個對比是一個很強烈的對比,.
現在我們看一個模型,.
是一比五十的模型,.
是當時耶路撒冷的一比五十的模型,.
我們從南面看上去,.
上面是北面,.
右邊就是東面,.
就是耶穌騎著驢進城,.
在艦艦山進入耶路撒冷的方向,.
他是騎著驢進來的,.
很多時候在猶太人的禦節的時候,.
因為有很多群眾來到耶路撒冷,.
羅馬帝國的王就很緊張,.
他們的總督彼拉多,.
就會騎著戰馬進入城裡,.
因為要去鎮壓他們,.
用強權去打壓他們,.
以十字架,.

$^{361}$以這些刑罰去壓住那些猶太人,.
這是一種帝國式的進城,.
一種帝國式的佔領,.
一種帝國式的和平,.
這就是所謂的羅馬和平,.
但是耶穌進來的,.
他帶來的,.
就不是以強權帶來的平安,.
而是以愛帶來的平安,.
以他的寫明,.
在十字架上的受苦,.
使我們這群的罪人能夠得以赦免,.
能夠得到平安,.
這是一個非帝國,.
或者反帝國的一種進城的方法,.
這個王不是我們眼中很多時候.
assume的那個王,.
沒有刀沒有槍,.
是以愛,以恩典進入我們心靈的一個王,.
你繼續看這個故事,.
你就知道不單止有鐘樹枝,.
不單止有騎著驢的君王耶穌進來,.
連他那個驢子也是他自己去預備的,.
你記得他說他插兩個門徒進去村莊裡面,.
他說你記得拿那個驢子過來,.
等於今天他插兩個人去你的車房拿車,.
那你會有什麼反應?.
你為什麼要拿我的車?.
你是不是想偷車?.
他就給了他們一個密碼,.
他說主要用他,.
換句話說,老祖耶穌已經預備了他想騎的驢,.
他那個Uber一早已經叫了,.
一早已經預訂了,.
繼續去記載這個星期發生的事情,.
有很多的地方是耶穌親自預備自己的受苦,.
自己預備自己的死,.
絕對不是一種意外發生的事情,.
你記得最後晚餐,.
耶穌說我們要預備這個最後晚餐,.

$^{401}$預備最後晚餐好像我們華人,.
預備我們的團年飯一樣,.
是要做很多功夫的,.
誰去預備?.
他就告訴那兩個門徒,.
你們為我最後預設的晚餐預備,.
他說你們進城的時候,.
必要有人拿著一瓶水迎面而來,.
你們就知道要跟他,.
然後在上樓房一個客廳的地方,.
你就為我們預備這個晚餐,.
這有什麼特別?.
有一個男人拿著一瓶水,.
你知道在古代,.
打水的是女人的工作,.
男人是不會拿著水瓶的,.
男人永遠都是坐梳化,.
在那裡享受,.
是女人做這些辛苦的工作,.
當門徒進城的時候,.
看到一個男人拿著水瓶,.
你就知道有境過,.
是有不同的,.
等於你今天去到公眾的地方,.
你看到一個男人夾著一個LV手袋,.
他一定是有不同的,.
這個就是一個密碼,.
爐子是他預備的,.
晚餐也是他自己預備的,.
這裡想告訴你一件事,.
就是耶穌基督,.
一位創天造地的主,.
他甘願來到我們人間,.
他拿了人的肉身,.
他百分之百是神,.
也百分之百是人,.
他預備這些東西,.
他親自在神的身上為我們寫記,.
不是意外發生的事情,.
而是他甘願來拯救我們,.

$^{441}$約翰福音第十章的說話,.
耶穌自己親自說的,.
道賊來,.
這裡說到是魔鬼撒旦,.
道賊來,.
是要殺害,是要毀壞我們的生命,.
耶穌說我來了,.
是要叫楊德森明,.
他用了一個比喻,.
他就是牧羊人,.
我們罪人就是羊,.
他是要我們楊德森明,.
並且得到更豐盛,.
他說我是好牧人,.
好牧人是為羊寫命,.
沒有人要他的命,.
只有他願意寫下,放下,.
最重要的東西,.
我們記念耶穌基督的受苦受難,.
是他親自來到世上,.
成為一個牧人,.
願意以無罪之身代替我們這班有罪的人,.
以致我們的罪能得到赦免,.
我們能成為一個新造的人,.
所以為什麼今天叫做中樹節,.
然後星期五晚上,.
我們在樹上有受難節的崇拜,.
我們會看看耶穌基督為我們所作,.
但同樣他不單止在十字架上死亡,.
但三天復活,.
他的墳墓是空的,.
這麼多信仰的系統,.
這麼多世界觀,.
只有耶穌基督的信仰,.
他自己的墳墓是空的,.
教主,運動的發起人,.
他們都是會死去的,.
唯獨是耶穌基督的墳墓是空的,.
這就是復活節,.
所以為什麼這個星期發生的事情是多麼重要,.

$^{481}$你剛才聽到一些見證,.
或者你認識一些基督徒,.
你就知道他們認識了主之後,.
不單止主幫助他們的難關,.
去醫好他們的病,.
最重要的就是我們基督裡面,.
我們能夠成為一個新造的人,.
你記不記得上帝的恩典,.
什麼時候進入你的生命,.
就是當你跟上帝說,.
我願意重聚悔改,.
我願意歸向你,.
我願意接受耶穌基督,.
你成為我的教主,.
那一刻開始,.
上帝的恩典就進入我們的生命裡面,.
我們的生命,我們的心態,.
我們對人生,對永恆的看法,.
就完全改觀了,.
這就是封城的生命,.
就是耶穌基督親自來,.
為我們寫明,.
帶給我們的祝福,.
我希望你今天,.
看完這些觀禮之後,.
看完這些朋友,.
他們接受了耶穌基督,.
作了建成,.
接受了水禮,接受了浸禮,.
我希望你看到,.
上帝的改變,.
上帝自己十字架的愛,.
他的犧牲,.
不單止改變他們的生命,.
也可以改變你的生命,.
如果你今天願意相信耶穌基督的話,.
你在基督裡面,.
就是一個全新的人,.
你以往發生的事,.
當然不可以.

$^{521}$奮不自禁地埋葬,.
但是在上帝眼中,.
你就是一個全新,.
新創造的英雄,.
一個新的生命,.
上帝看你,.
就像看他自己完美的獨生子耶穌基督一樣,.
這就是豐盛的生命,.
也有復活節,.
我們生命裡面一定會有困難,.
甚至我們會面對死亡,.
但是耶穌基督的墳墓是空的,.
祂給我們永遠的盼望,.
我希望這個恩典能夠成為你自己生命裡面的實在,.
你接受耶穌基督的時候,.
我們每一個都在主的面前成為一新造的人,.
我們同心低頭祈禱,.
主愛多謝你,.
我們能夠去禱告,.
主愛你感恩,.
因為你親自的來到,.
為我們赦命,.
我們是一班不配受你的愛,.
不配受你恩典的罪人,.
我們都不是完美的,.
但是你自己願意來到十字架上,.
為我們作真正的君王,.
以愛以犧牲來救贖我們,.
去改變我們,.
主愛你願意在我們中間,.
每一位的朋友,.
每一位的弟兄姊妹,.
我們重新明白到,.
在主的恩典當中,.
你給我們的十字架的救恩,.
使我們能夠成為新造的人,.
幫助我們能夠認識這個救恩,.
認識你成為我們的救主,.
我們最能夠得到赦免,.
我們在永恆裡面能夠有盼望,.

$^{561}$為我們犧牲,.
為我們復活的主,.
你給我們生命裡面重新的有盼望,.
神愛多謝你,.
我們繼續求你帶領我們弟兄姊妹,.
我們教會,.
讓我們在新的一年當中,.
我們能夠更加努力的去事奉你,.
去見證你,.
祈禱奉主耶穌基督的名字..
\newpage



\section{}
\label{sec:zcyVrDXZup0}
\textbf{7/22/2018 主日講道: 基督徒成長12步\#8 委身傳揚(陳崇基牧師)}
\newline
\newline
連結: \href{https://youtube.com/watch?v=zcyVrDXZup0}{\texttt{https://youtube.com/watch?v=zcyVrDXZup0}} ~~~~ 語音日期: 2019-07-21
\newline
\newline
\hyperref[sec:OTRUSy4_3CA]{\small{< < < PREV SERMON < < <}}
~
\hyperref[sec:index_chronic]{\small{[返順時目]}}
~
\hyperref[sec:index_scriptual]{\small{[返順卷目]}}
~
\hyperref[sec:QRMnvNFd6QE]{\small{> > > NEXT SERMON > > >}}
\newline
\newline
$^{1}$今天我們想來強調一個主題.
可能這個主題你們聽過.
我們每個月一次會講.
由今年年頭開始到現在.
我們會講一個主題.
跟我們基督徒成長有關的.
每一個信耶穌的人.
由第一步開始跟隨耶穌信徒的福音之後.
我們主要崇拜的是認識上帝.
然後認識信仰的基本.
繼續探索真理等等.
一直受洗見證等等.
我們已經到了第八個了.
第八個剛剛就講到宣教.
毀身傳若.
剛才聽完May May分享.
我不知道你有沒有想過.
宣教是一個這樣的title.
Missionary.
一個這樣的title.
你想的時候你會覺得.
會不會是你自己呢.
或者你聽見missionary的時候.
可能是站在台上分享.
或者那些被差派出去的人.
那些才算是missionary.
算是宣教士.
我自己就不是.
我覺得那些宣教士是有光環的.
無論我去那裡做事都好.
但是聖經是不是這樣講呢.
聖經原來說你和我.
每一個信耶穌的人.
每一個都認信耶穌基督的人.
其實我們就是宣教士.
Missionary.
Missionary的意思就是說.
我們帶著耶穌基督給我們的message.
我們走到我們自己的地方.
我們的崗位.

$^{41}$我們就是去宣講這個訊息.
是不是.
聖經裡面有很多這樣的故事.
今天我想和你提一段很重要的.
記載到第一世紀的教會.
2000年之前的教會.
第一世紀的教會.
這個教會在一本書裡面.
文件裡面記載的.
叫做《使徒行傳》.
這本叫做Book of Acts的書.
記載了一個很重要的發展的突破.
當時教會耶穌基督死而復活之後.
大概AD主後30年左右的時候.
去到這個時候大概AD40多的時候.
是一個宣教上的突破.
11章第19節.
那些因私提反Steven這個人.
的事而遭患難四散的門徒.
直走到菲尼基一個地方.
另外一個地方就是塞浦路斯.
並且去到更遠的地方叫做安提奧Antioch.
他們不向其他人講道.
不去跟其他人講福音.
只是向猶太人講福音.
但是內宗有塞浦路斯人或者古利內人.
他們來到安提奧就向希臘人.
就是非猶太人.
所謂猶太人眼中的外邦人.
傳講主耶穌.
主與他們同在,信義歸主.
就很多了.
如果你看一看地圖.
大概這個《使徒行傳》.
這本Book of Acts.
你當它是歷史文件.
就算你沒看過聖經.
你不去教會你不信耶穌.
這本文件記載一件事.
就是記載當時在耶路撒冷開始的運動.

$^{81}$耶穌基督死而復活之後.
他們得到這個訊息被差派出去.
在耶路撒冷開始.
大部分都是猶太人的基督徒.
他們在猶太的地區.
耶路撒冷來到建立教會.
傳福音,有人信住.
建立了教會,Church planting.
然後就有一件事發生.
《使徒行傳》第七章說到.
有一個叫司提反的人.
他就殉道.
在教會又開始.
幾年之後.
就有第一個宣教士殉道.
Steven這個宣教士死了之後.
因為教會受到大規模的迫害.
被人欺凌,被人趕走.
所以門徒就四散.
《使徒行傳》到今天我們看《元年經文》.
有一章這個故事就說到.
他們四散到北上.
上到塞浦路斯的地方.
坐船過去.
或者去到Fenicia這個地方.
這些全部都是當時最重要的trading route.
最重要的商貿路線.
他們就去到這些地方.
然後他們帶著福音.
他們就去到那裡.
就向自己的人,猶太人的基督徒.
向猶太人去港島.
甚至有些去到安提柯.
另一個大的商貿的城市.
他們去到這裡又是帶著福音.
然後將福音傳開.
留意裡面一些的智女行間.
是需要我們明白當時發生了什麼事.
他在這裡說什麼呢.
他說他們只是向猶太人港島.

$^{121}$其他人他們都不說.
為什麼呢.
原來安提柯這個地方.
如果你查一下當時的新聞的歷史.
你就知道幾年前.
這件事發生幾年前.
就有猶太人和希臘人或者羅馬人.
有一個很大的衝突.
死了很多的猶太人.
因為種族歧視的問題.
耶路撒冷總部的大祭司.
就派了三千個男丁上去安提柯.
然後就將那裡的愛邦人又殺了很多.
這些叫做劈友.
這些真正的劈友.
就是愛邦人和猶太人.
希臘人和猶太人是河水不犯井水.
所以那些猶太人基督徒.
之前上到安提柯的時候.
只是和自己同聲同氣的人來港島.
自己種族的人就說.
但是這裡第一個突破是什麼呢.
居然有這樣的信徒.
來到安提柯的時候.
他們不單止向自己的種族的人來港島.
甚至向希臘人.
即是愛邦人來傳講耶穌基督.
這個是第一個革命.
打破了這個僵局.
然後有一些愛邦人來信署.
我想問你福音是為什麼人而設的.
是不是只是為和我們同聲同氣的人而設.
我們講廣東話的.
我們大粵語民族的.
只是為我們而設的.
當然不是.
是為所有人而設的.
無論他講什麼話.
他是什麼種族都好.
我們其實都有這樣的機會.

$^{161}$或者有這樣的推動力.
我們應該要向他們傳講耶穌基督.
教會應該是這樣.
因為我們教會也是這樣.
不只是我們同聲同氣的人走在一起就算了.
安提柯的信徒就是有這樣的心智.
原來帶著這個福音是一個普世的福音.
任何人都適合聽.
任何人都應該要知道.
他裡面講到.
那些因屍體犯的人遭患難四散的門徒.
那些人沒有名字.
可能當中沒有誰出名.
屍體犯殉道之後.
這些信徒被四散到不同地方.
他們其實帶著福音到處去.
其實你看小行傳這本書.
那些因迫害而四散的門徒.
其實這是第二次出現.
第一次出現是什麼時候呢?第八章.
第八章發生什麼事呢?.
之前發生什麼事呢?.
屍體犯被害之後.
門徒四散.
然後教會的領袖也有被四散.
肥力是其中一個.
或者Peter.
這兩個人就去了有意外的地方傳福音.
然後有人信了主.
甚至有埃及的人信了主.
現在小行傳十一章十九節.
又講門徒又四散.
他們散到安提柯的地方.
然後他們又將福音傳到更遠的地方.
其實兩件事是平行的.
就是告訴我們.
縱使我們不一定是教會的領袖.
什麼牧師什麼有title的宣教士.
不是肥力我們不是彼得.
但是我們四散到周圍的地方的時候.

$^{201}$我們就等於他們一樣.
並且我們將福音傳到更遠的地方.
兩個故事平行就是告訴我們.
所謂普通的信徒.
我們覺得我們是Late Christians.
或者我們是平信徒.
其實我們和教會的領袖宣教士.
其實在做同一件事.
就是我們將福音傳到更遠的地方.
我們教會很多時候.
有一個迷.
就是明星迷.
我覺得如果有些藝人.
是基督徒的藝人.
請他來這裡傳福音就好了.
因為他賣座.
如果印海報.
Mark Sir有個分享會.
分享他自己的見證和經歷.
有沒有人來?沒有人來.
關心賢就不同了.
如果關心賢來就幾百人來.
號召力夠.
但這裡想告訴你什麼呢?.
原來我們每一個信徒.
我們都是宣教士.
他們是什麼人.
他們只不過是普通的基督徒.
所以每一個人都是宣教士.
你去歐洲玩的時候.
有沒有想過當地教會的需要.
你去日本玩的時候.
有沒有想過當地教會的需要.
有很多宣教士需要我們支持.
需要我們祈禱.
需要我們和他們同工.
我記得有一次我自己一個人.
去到埃及的開羅.
我擠了很久的巴士和的士.
我才去到那個地方.

$^{241}$因為有一個牧師我認識的.
那裡有一個宣教士.
不如去探望他.
我約了他去餐廳等.
等了一個小時因為他有遲到.
大家談了一晚.
他一個人在那裡做阿拉伯的服務工作.
很孤單.
所以有人為他祈禱.
所以去支持他.
我們每一個人去到哪一處.
其實我們都是被四散的信徒.
第二點我想和你說的.
就是我們每一個基督徒都是牧者.
11章22字.
他說這個風聲傳到耶路撒冷教會的耳中的時候.
總部耶路撒冷教會的總部.
他打發一個人去探索一下到底發生什麼事.
他打發了一個人叫巴拿巴.
巴拿巴去到安提河.
他去到那裡的時候.
他見到上帝所賜的恩就歡喜.
並且勸勉眾人立定心志.
恆久靠主.
這個巴拿巴原本是一個好人.
被聖靈充滿大有信心.
於是有很多人來去歸服了主.
巴拿巴作為一個猶太人的信徒.
去到安提河一個比較多愛幫人的教會.
他做什麼?.
他不是潑冷水.
他去支持他們.
去鼓勵他們.
不是隨便你們自己做.
他是幫助他們然後牧養他們.
見到新的一代開始.
見到新的教會開始.
見到新的宣教的工作.
是值得歡喜的.
值得去鼓勵的.

$^{281}$事實上這個巴拿巴.
由小林傳開始介紹他的時候.
他就是願意奉獻.
願意出錢出力.
來推動教會方面的工作.
就是因為這個緣故.
所以安提河教會就能夠被建立出來.
於是我們其實就是巴拿巴.
我們就是帶著我們的使命來到教會.
我們去到堅強的教會.
使教會能夠得到更多好處.
短孫也好.
留下的我們在社區裡面.
不同的工作崗位不同的地方.
其實我們就帶著這個使命.
使教會能夠得以建立.
25節.
還有一樣東西巴拿巴做了.
就是什麼呢.
他又往大掃城來找蘇羅.
找到了就帶他到安提河去.
他們足有一年的功夫.
和教會一同聚集.
教訓許多人.
門徒稱為基督徒是從安提河起手.
安提河是羅馬帝國.
第三大城市.
不小的.
那裡有很多的外邦人.
很多的多元的文化.
巴拿巴作為一個前輩.
他就知道如果他繼續建立教會.
繼續建立福音的教會.
繼續建立宣教的工作的話.
他就需要有幫手.
他就找了一個人叫蘇羅.
蘇羅就是保羅.
猶太人的名字或者羅馬的名字.
同一個人.
找了他 為什麼呢.

$^{321}$找蘇羅就最適合了.
因為大掃城在安提河附近.
大掃城也是一些金融的中心.
商貿的中心.
保羅在那裡從小到大.
他就聽了很多羅馬希臘的字號.
所以這個人是三大文化都熟悉的.
保羅熟悉什麼文化.
熟悉猶太人自己的文化.
因為他自己本身是希臘.
希伯來人 是猶太人.
他也讀過律法書.
保羅也知道希臘人的字號的思想.
所以他的雅典去到愛邦人的地方.
沒有什麼問題.
保羅也是羅馬的公民.
所以按照他的文化來說.
按照哲學來說.
按照政治來說.
保羅就最適合了.
所以巴拿巴就去找了蘇羅.
以致能夠將福音做得更加好和更加大.
更多人來認識.
教會就需要有新舊的體會.
安提安教會就是第一個打破這個僵局.
不只是舊人做.
而是有新的人來做.
繼續宣教 繼續有福音的工作.
教會其實也是一樣.
香港建築學院那裡畢業.
梁家倫院長.
今年他退休.
退休的致詞裡面說什麼呢.
退休的講道.
他說不阻人上位.
是我最大的貢獻.
梁家倫.
我從小到大受他影響很多.
因為他著書等身.
一個這樣的學者.

$^{361}$他有很好的識經書.
很好的基督教的書籍.
他退休的時候說什麼呢.
不阻人上位是他最大的貢獻.
很坦然的說.
我們這些老餅.
這些死老鬼.
是年輕人不會接受到的方法.
讓年輕人用他們自己的方法.
去表揚耶穌基督的愛.
不用我們老餅的方法.
用年輕人的方法.
使教會.
使這樣的工作.
使神學院可以年輕化.
其實我們教會也是一樣.
我們成立了兩年.
今年年底我們要改選.
所謂老餅.
有些要走.
我不是說只有年輕人才做.
不是.
但我們應該要給年輕人多一點空間.
這麼多代跨代的.
我們一起來領導教會.
最重要的是新一代.
安提拉教會就是做了這個.
史瑭珠寫下了一個故事.
其實就是挑戰每一代的教會.
甚至今天的教會.
到底我們是不是這樣.
他說門徒稱為基督徒是從安提拉起手.
你知道新約聖經裡面.
叫做基督徒的人.
都不會用這個Christian.
Christian基督徒這個名字.
在新約的文件裡面你查一下.
大部分的名字都是叫弟兄姊妹的.
家人.
信徒.

$^{401}$門徒.
當時沒有人叫基督徒.
基督徒是哪些人叫他們的呢.
就是那些不信主的人.
看到那些人.
他們說他們就是Christian.
為什麼呢.
因為他們經常提的就是基督.
是Christ.
所以加個人就叫基督人.
其實原文就是這個意思.
原文希臘文就是基督的人.
講基督的人.
或者身上有基督名字的人.
到第二世紀的時候.
就是AD 200的時候.
教會.
開始被茁壯建立起來.
他們就用了這個名字.
不如我們就叫我們自己做基督徒.
其實這樣是最適合.
因為我們每一個說自己信耶穌的人.
其實就是基督的人.
我們就是帶著耶穌基督在我們身上.
別人會看到.
你怎麼生活.
然後他會從你的生活中.
認識到你的基督是怎樣的.
別人會看到你怎樣和別人相處.
你怎樣和別人相處.
就告訴別人耶穌基督的福音是怎樣的.
我們就是Christian.
是一個missionary的church.
其實我們UV的vision.
都是一樣的.
我們每一個都應該做一個missionary.
我們有三個核心價值.
Be real and be bold and be relevant.
我們要勇毅真誠.
還有緊貼善待.

$^{441}$走到我們不同的崗位.
無論你今天.
無論你為什麼去日本.
你去中東.
你去台灣.
你去哪裡都好.
你回到自己家裡都好.
我們就是Christian.
我們一起祈禱希望教會成為這樣.
好嗎.
在我祈禱當中我想問你幾個問題.
第一.
你是不是一個宣教士.
上帝在你心目中.
給你一個什麼人或者哪一類人.
是你要向他傳福音.
第二.
你怎樣去成為一個牧者.
可以鼓勵到教會.
可以同心聽到.
主我們多謝你.
因為有其他的人.
是你自己猜拳的宣教士.
的信徒.
能夠將福音帶給我們.
以致我們能夠成為你自己的弟兄姊妹.
我們成為基督徒.
主讓我們.
帶著這個新的身份.
我們就是missionary.
我們去到社區.
城市世界不同的角落.
我們就是帶著這個missionary.
讓更多人認識主你.
我們也帶著這個牧者的身份.
在我們的教會.
去到不同地方的教會.
能夠使教會來茁壯.
使更多人認識主你.
讓你自己藉著我們的生活.

$^{481}$能夠告訴別人.
耶穌基督你是值得相信.
值得跟隨.
值得去敬拜的神.
主幫助我們的生活.
我們的一言一行.
讓我們能夠映照出.
我們就是基督的人.
是基督徒.
Christian.
讓我們生命帶著你的message去到不同的地方.
使其他人能夠得到這個福氣.
祈禱給主明.
\newpage



\section{}
\label{sec:QRMnvNFd6QE}
\textbf{2/24/19 主日講道: 事奉的迷思 (Antonio Luo)}
\newline
\newline
連結: \href{https://youtube.com/watch?v=QRMnvNFd6QE}{\texttt{https://youtube.com/watch?v=QRMnvNFd6QE}} ~~~~ 語音日期: 2019-07-22
\newline
\newline
\hyperref[sec:zcyVrDXZup0]{\small{< < < PREV SERMON < < <}}
~
\hyperref[sec:index_chronic]{\small{[返順時目]}}
~
\hyperref[sec:index_scriptual]{\small{[返順卷目]}}
~
\hyperref[sec:UcXT0Snl5mA]{\small{> > > NEXT SERMON > > >}}
\newline
\newline
$^{1}$各位弟兄姊妹平安.
我這幾天有點不舒服.
所以我的聲音有點奇怪.
所以希望大家諒解.
2月底,剛剛我們過了中國人的農曆新年.
我發現來美國久了.
我們在美國就不會再像以前在香港或中國大陸慶祝新年那種氣氛.
我們不會再做一些以前會做的事.
例如我們不會逛花街,不會去做什麼的,我已經開始忘記了.
派利是,我們已經不再做了.
我們不再去吃湯圓,或者我們不會再去燒豹掌.
很多這些傳統的東西,我們中國人來到美國之後就開始慢慢流失了.
但我們也開始接納一些西方的傳統.
來到新年的時候,我們會做一些新年的新年決議.
大家有沒有做過這些?.
都有的吧?.
我自己會有.
通常每年的大計就是今年真的要減肥.
或者今年要做一些,上班要勤力一點,讀書要認真一點.
這些通常在新年大計都是走兩邊的.
一邊是很模糊的,今年做好一點,上年不夠,今年做好一點.
另一邊可能會說,不是的,真的落實一個計劃去做好一點.
真的去做好這件事好不好.
那時候就會說,好像後者會比較容易成功.
譬如說,五月前我們跑馬拉松,那你會開始計劃的.
你要開始鍛鍊,開始去找資料,去報名.
如果你說今年真的要減肥,真的要做了.
那你也開始節食,開始看看營養.
開始去準備自己的鍛鍊計劃.
通常有計劃就比較容易成功.
但是有多少人真的試過做好這些research,做好這些計劃.
但是也沒有做到呢?.
我又是其中一個,這個我太太經常提醒我.
為什麼呢?為什麼我們不會做到呢?.
有時候可能我們頭腦上,知識上,知道要做些什麼.
但實際行動,我們都缺乏一些動力,缺乏一些決心去做到這些東西.
會不會我們其實事奉上面都會是這樣,有時.
我們事奉上面都是用我們的頭腦上的事奉神.
不是用我們的身體行動.
我們缺乏一些動力,缺乏一些知識,還是缺乏一些什麼呢?.

$^{41}$今天我跟大家分享的這段經文是約書亞記二十四章.
這裡我們會用以色列人的一個反例,一個反面的例子.
去幫助我們思考事奉究竟是一樣什麼.
事奉究竟是不是我們要做一些事情呢?.
事奉究竟是不是我們要付出很多呢?.
是幫教會做事?幫牧師做事?.
事奉究竟是什麼一回事?.
我不知道大家對約書亞記熟不熟.
約書亞記二十四章,大家通常最熟悉的那一節就是第十五節.
至於我和我的家,必定會事奉耶和華.
當年我也是用這句話向我太太求婚.
她也說OK,那就OK吧.
但約書亞記二十四章不是在說這件事.
不是說我們要怎麼事奉,是說另外一些東西.
二十四章開始,一至十三節說的是約書亞再重申一次.
召集以色人在事件的地方跟他們再說一次.
你們快要住進這個地方了,你們要做些什麼?.
約書亞記二十四章第一至十三節重申一次.
神在你們當中做了些什麼事.
這裡說到…不好意思,忘記拿了.
不好意思.
這裡約書亞重申跟以色列人說的.
重申一次神對他們做了些什麼.
一至十三節,在複述神在他們當中的作為.
神怎麼將他們的列祖,亞巴拉罕,雅各,以撒,以撒那各.
從他們本身的地方帶出來,來到他們迦南的地方.
後來也將他們帶到埃及,也在埃及將以色列一族帶出來.
當中經歷過的所有的歧視,約書亞重申的再跟他們說.
但不單只是以他們列祖的事.
還有他們上一代已經經歷的事.
上一代經歷的是什麼呢?就是在抗夜的四十年,他們經歷過的.
他們知道的,看到上一代因為叛逆神他們進不去.
但神仍然帶領他們經歷的所有的對他們的祝福和恩惠.
來到這裡了.
他們不單止快要進入迦南地.
甚至是進入迦南地的過程中.
他們將耶利哥城倒塌,完全拆毀.
讓他們準備分好地,準備住進去.
準備住進神應他們流瀨與物的地方.
這就是約書亞一至十三節的背景.

$^{81}$十四節二十八節大概分為三段對話.
是約書亞和以色列人的對語答.
大概分為三段.
我們首先看第一段對話.
十四節說:所以你們要敬畏耶和華,誠心誠意地事奉他.
摒棄你們祖先在友法拉底河那邊和在埃及所拜的神明,專心事奉耶和華.
如果你們不願意事奉耶和華,今天就選擇你們的神明吧.
或在大河那邊你們祖先事奉的神明.
或者你們在這裡阿摩利人的神明.
至於我和我的全家,我們必事奉耶和華.
第一句說話,約書亞對他們說:你們要選擇.
你們要選擇你們要怎樣選擇.
你們要事奉他還是不事奉他.
你們要事奉你們以前列祖的神.
你們要事奉在埃及那邊的神.
事奉這個地方的人,你們快要住進這個地方了.
你們和他們一起去住吧,怎樣呢.
約書亞也對他們說,他們自己的答案.
他和他的全家會事奉耶和華.
但為什麼約書亞要他們選擇呢?.
他們不是一直都事奉耶和華嗎?.
為什麼約書亞要跟他們說,你們要專心事奉耶和華.
你們要選擇你們的神明.
你們不是已經選擇了嗎?他們不是一直都跟著這裡嗎?.
我們看看以色列人怎樣回答.
他說我們決不會背棄耶和華去事奉別的神明.
因為我們上帝耶和華曾經令我們和我們的祖先脫離埃及人的奴役.
在我們眼前行大神跡,保護我們安然來到列邦.
把這個地方賜給我們.
把這個地方的阿摩利人和外族人在我們前面趕走.
我們必定要事奉耶和華.
因為祂是我們的上帝.
我們聽上去其實以色列人答得很對.
他沒有說錯話.
他說,他說了約書亞剛才說的話.
神已經做了那麼多事,我看到的.
我們會做的.
我們會選擇的,因為這是我們的上帝.
但似乎約書亞不是很相信他們.
我們之後會再看到,約書亞再挑戰他們兩次.

$^{121}$為什麼呢?.
我們待會再回答這個問題.
但他們明白的,他們頭腦上是明白的.
就像我們這樣,我們很多人都明白.
我們明白神在我們心中做了什麼.
我們知道神怎樣拯救了我們.
我們知道耶穌基督曾經在十字架上為我們寫命.
赦免我們的罪.
也知道祂在死裡復活.
然後升天.
我們知道這些東西,但代不代表我們知道.
就是可以事奉上帝呢?.
我們再看.
約書亞第二段對話開始就說.
你們不能夠事奉耶和華.
因為祂是一個聖潔的上帝.
祂痛恨不正.
必不赦免你們的過犯和罪惡.
事奉神,剛才約書亞給了他們一個選擇.
叫他們選擇.
現在約書亞就說.
你們根本不可以事奉上帝.
讓你們選擇,但你們不可以.
因為祂是一個聖潔的上帝.
為什麼呢?.
為什麼要描述這個神是聖潔的神呢?.
叫他們選擇事奉神.
但也同時說.
你是一個聖潔的神.
你們不可以隨便地接觸祂.
這個對神的描述.
有很多地方都在說聖潔這件事.
特別是在《利未記》裡面.
整本書,頭十一章都說.
祭司如何需要聖潔.
讓他可以進入聖殿裡面.
去至聖所,去敬拜神.
但是之後.
十二章打後到尾尾.
都是在說整個以色列人.

$^{161}$整個群體.
是同神同在的.
不斷重複地說.
我,耶和華,你的神.
是聖潔的神.
十九章第二節說.
你們要聖潔,因為我.
耶和華,你們的上帝,是聖潔的上帝.
是聖潔的.
後來《申命記》.
之前的.
這裡之前的說.
你們屬於你們的上帝耶和華.
是聖潔的民族.
天下萬族中選擇了你.
做祂的子民,祂寶貴的產業.
為什麼要說聖潔這件事.
如果說.
神說的聖潔.
就是說,我要.
傳言.
你們是來不了我那裡的.
我是聖潔的.
為什麼要提出聖潔這件事.
如果說,你們不能接觸我.
根本你們不能事奉我.
為什麼要提出聖潔這件事.
我想可能.
或者下一句.
下一句.
說的是,祂是一個聖潔的上帝.
因為什麼呢?因為祂痛恨不正.
在這個和合本.
我們通常會記得.
和合本翻譯就會說.
是嫉妒的,忌邪的神.
和合本翻譯通常是.
忌邪的神.
如果英文聖經就會翻譯.
He is a jealous God.

$^{201}$是嫉妒的.
但這個.
嫉妒,Jealousy這個字.
是一個充滿感情的字.
是一個關係的詞語.
神不會形容.
祂自己說,我是一個很嫉妒的神.
我傳言,你們不可以.
來接近我,我是很嫉妒的.
不是只是說對自己的特性.
這是一個relational term.
是一個和我們有關係的詞語.
所以當我們想到神的聖潔的時候.
其實是神想說.
我和你有一個關係.
這個嫉妒這個詞通常都形容.
是我們夫妻之間的那種關係.
不可以有第三者.
如果中間有第三者的時候.
我們的嫉妒會出現.
這是一個充滿情感,充滿感情的用字.
如果神將這個字用在自己的時候.
其實他是很重視.
神和以色列人.
之間的那個關係.
不單止是他們.
要做的事,不是他們要做什麼.
是他和以色列人之間的關係.
他說.
必不赦免你們的過犯和罪惡.
為什麼要不赦免他們呢?.
因為他是聖潔.
因為他是公義的神.
他不可以.
他不是一種.
會溺愛的.
做錯了不要緊,不要緊,再來一次.
有時候我們父母.
都可能會溺愛我們自己的兒女.
說不要緊,不要緊,你做錯了不要緊.

$^{241}$但是神不可以這樣.
因為他是聖潔,因為他是公義的.
在二十節那裡說.
你們如果背棄耶和華.
去拜外族的神明.
儘管他曾經恩代你們.
賦禍給你們,毀滅你們.
這裡說的是.
剛才我們第一節到十三節.
若書雅已經說.
曾經怎樣待你們的先祖.
你們的先祖.
甚至你們的父輩.
還有你們.
神對你們的恩典已經顯明了.
他已經恩代了你們.
但是如果你們繼續.
這樣去拜外邦的神.
他必會轉.
他必會轉向你們.
然後賦禍給你們.
我們聽的時候就會說.
這個神好像很凶惡.
這個神好像.
真的不是.
不是那麼容易去做.
好像很難.
為什麼這麼嚴謹.
但是這個.
我們可以想成.
神不想我們走錯.
好像.
我們如果去.
這段時間沒有什麼山地.
如果我們去到.
加州的Yosemite.
你去山上的話.
你會看到山路旁邊.
就是懸崖.
他就一定要放一個護衛的路.

$^{281}$路邊的護衛.
讓你不要走進去山裡.
神是真的有一個.
真正的警告.
真的警告跟你說.
不要走錯.
不要掉下去.
神想的.
不是去嚇他們.
神想的是.
讓他們知道.
你專一事奉我.
專一事奉神.
好讓你們得福.
承受這個地方.
文中在第二段對話裡.
他說不是我們一定要事奉耶和華.
事奉我們的上帝.
來到第三次的對話的時候.
約書亞就說.
你們現在自己作證.
你們發誓.
你們已經選擇了耶和華.
選擇事奉耶和華.
文中說.
我們願意為自己作證.
我們願意發誓.
作證就像我們在公證人前面.
sign notary的時候.
你會舉手發誓.
所說的一切是真實的.
白紙.
黑字上面.
所寫的東西都是真實的.
都是真確的.
如果真的這樣的話.
你們現在.
就摒棄你們中間的.
外族實名.
就是說.

$^{321}$這班以色列人.
準備進入.
他們迦南的美地的時候.
他們抱著很多神.
說要事奉耶和華.
明白嗎.
一開始第一次對話.
約書亞跟他們說.
你們選擇吧.
你們列祖的愛邦神.
還是這裡的阿摩利人的神明.
你們選擇你們事奉誰還是事奉耶和華.
其實他們已經抱著.
一些.
我們說的菩薩佛祖.
在說我們要.
我們要事奉.
我們要事奉上帝.
我們要事奉主耶穌.
但不是.
約書亞說你們不可以這樣做.
會不會我們其實.
我們事奉的時候.
我們也帶著一些其他的偶像.
去事奉上帝呢.
會不會我們都.
忘記了.
我們的上帝其實.
不是認識我們的上帝.
其實是一個聖潔.
一個.
質度的上帝呢.
這次是第三次.
約書亞跟他們說.
叫他們選擇.
叫他們.
第一次叫他們.
當中有偶像.
但第三次提到他們當中有偶像這個問題.
民眾第三次.

$^{361}$也回答我們必事奉.
我們的上帝耶和華.
聽從祂的話.
第三次的回答.
這次多了一句.
不是說我們一定會事奉上帝的.
我們會聽話.
但聽話這裡.
也不是說從此以後.
我們就專心.
去事奉上帝.
宏觀地看其實.
以色列人進了迦南地.
下一章就是.
士司祭.
我們會發現士司祭.
基本上是一個悲劇.
這句話是什麼意思呢.
這句話必聽從他的話是什麼意思呢.
其實是.
其實是一個.
怎麼說呢.
Anticipatory.
他們是.
我們會事奉上帝的.
你說什麼我們就說什麼.
這是一個回答.
所以我們看到.
約書亞就跟他們立約.
重新再說一次.
神的律例章典是什麼.
跟他們立約.
神需要的是什麼.
重新再說一次.
神接著摩西.
抗夜跟他們說的.
所有的律例章典.
包括昌克及記.
利未記.
還有申命記等等.

$^{401}$重新再說一次.
神究竟需要他們.
怎樣的.
對他們有什麼要求.
他們怎樣的群體與神.
去相交.
享受這個關係.
他們在事件立約.
在這個.
像樹下面納了一個大石頭.
作為一個見證.
這些見證就是.
我們會有的.
一些紀念碑.
類似這些物體.
這個見證著你們納的約.
如果你們背棄.
耶和華.
上帝的話他們必會作證.
然後約書亞就.
叫人們回到自己的地方.
我們從這三段對話中.
我們去想.
究竟.
事奉是什麼.
或者是以色列人.
他們對事奉有什麼迷思.
他們以為事奉.
上帝好像和.
其他外邦的神一樣.
抱得越多就越有好處.
或者是.
我們拜黃大仙.
這個上帝很厲害.
有能力.
應驗.
我們去拜祂.
我們去跟祂.
我們準備住進一些.
我們不用建造的城市.

$^{441}$不用栽種的葡萄園.
我們可以拿來吃.
我們做的田地.
我們就可以享受這些一切.
是不是因為這些東西.
令他們覺得.
這個神真的可以.
可以跟著.
他們可能有這樣的想法.
或者他們.
一直以來.
神根本是.
容忍著他們.
帶領他們來到這裡.
令他們覺得這個神真的好.
我們可以.
隨便吧,他無所謂.
但是約書亞提醒他們.
你將要進入這個地方.
你們對神興萬不得.
不可以這樣.
跟我們有什麼關係?.
事奉上帝.
有時我們對事奉的想法.
就是有些奇怪的想法.
奇怪的想法.
可能事奉上帝有好處.
在教會做事有好處.
或者做傳道人有好處.
雖然我不覺得這是一個事實.
所以我都.
我覺得.
事奉上帝其實是有代價.
不是有好處.
但有些人會覺得事奉上帝有好處.
可能在教會裡.
拿到權利.
拿到別人的尊重.
因為自己的事奉而驕傲.
或者我覺得.

$^{481}$我的茶經帶得很好.
我的小組又搞得不錯.
我會覺得.
我自己的事業.
小組又搞得不錯.
有好處.
因為自己的驕傲有成功感.
或者覺得事奉.
令我們熟悉一點.
事奉說.
我有事奉.
最近有事奉真的覺得不錯.
最近沒有事奉我覺得.
跟神人的關係有點差.
或者事奉.
令我們有人稱讚.
得到肯定.
所以約書亞和以色列人一直說的.
不是事奉.
不是這些東西.
事奉是什麼呢.
我用兩點.
我覺得.
或者可以從我們今天的經裡面去看到.
其實事奉是一種敬拜.
約書亞一直和以色列人.
針對他們敬拜偶像這件事.
來跟我們說.
是一種敬拜.
不是你做一些什麼.
是你敬拜的態度.
你究竟是不是敬拜.
專一地去敬拜這個獨一的神.
羅馬書要說到.
保羅說.
所以弟兄姊妹.
我憑上帝的憐憫.
勸你們獻上自己的身體.
作聖潔蒙上帝月立的活祭.
你們理所當然這樣事奉.

$^{521}$如果事奉是敬拜的話.
事奉這件事.
不只是牧師.
或者神學生.
或者執事長老做.
而是每一個願意來到神面前去敬拜他的人.
都應該參與的一件事.
事奉是每一個願意來到神面前去敬拜他的人.
都可以參與.
都應該參與的一件事.
所以事奉其實是一種關係.
神不斷重複的.
神不斷重複的說他的聖潔.
說他對以色列人的要求.
其實他著重的是他們之間的關係.
他著重的是.
和以色列人之間沒有第三者出現.
不希望他們抱著其他的人.
而來到上帝說主啊我愛你.
你們還沒記住第19章第二節.
我們重新讀一次.
你們要聖潔.
因為你們的上帝耶和華是聖潔的.
從來神說到他自己的聖潔.
都是說和你們有什麼關係.
和我們有什麼關係.
我們事奉的時候.
其實就是進入這一段的關係.
願意被神使用.
容許我用一個例子來說明這件事.
兩個星期前應該是情人節.
我小朋友回到學校.
他們都會慶祝情人節的活動.
拿了很多糖果回來.
其中我二兒子Eliot就拿了很多糖果.
然後他就拿一個出來.
他說爸爸這個給你.
我很開心.
是不是因為他給了我一顆糖果.
我想吃那顆糖果而開心呢.

$^{561}$而不是.
當然不是.
是因為他知道.
我和他有一個很親密的關係.
他願意用那顆糖果.
去彰顯他對我的愛.
神同樣的.
希望我們每一個人.
去事奉他.
無論是金錢上.
用我們在教會的服侍.
甚至是在你自己生命裡面.
不同的崗位上.
我們做的時候.
不是為了做一些事給神.
而是為了回應神.
已經給了最好的禮物給我們.
讓我們進入這個.
和他修復.
進入這個父與子.
父與女兒這種親密的關係.
我們可以做一些回應.
所以我鼓勵大家.
願意我們在想事奉的時候.
不是在想我們要為教會做些什麼.
幫Mark Sir做些什麼.
幫這個機構做些什麼.
幫神做些什麼.
反而是.
在思想上是.
我們是一段關係.
和神的一種親密關係的一種回應.
讓我們再次去享受這段關係.
通過我們做的事.
我們的生命的活出來.
無論你是家庭的主婦.
你在工作上.
怎樣去待人處事.
如果你是學生.
在你的學業.

$^{601}$在你的.
尋求知識的過程上.
怎樣去彰顯去榮耀神.
怎樣去事奉神呢.
這個我留給你們自己去思想.
但我想說的.
事奉其實是我們的敬拜.
我們每一個尋求.
我們來到主耶穌面前.
願意稱祂為主.
願意藉著祂的名得救.
進入這一段新的.
和神修復這種新的關係.
的時候我們.
一個敬拜的行動.
事奉是我們的關係.
是我們享受經歷和神.
這個關係的一個過程.
也希望大家可以.
在你們的事奉當中.
繼續去經歷上帝.
經歷祂為你做的神蹟奇事.
繼續經歷祂給你們的恩典.
我們一起低頭祈禱.
慈悲的父母我們多謝你.
多謝你給你的意志我們.
讓我們能夠重新的.
可以和你有一個新的關係.
我們多謝你給很多很多的東西.
讓我們豐足不知缺乏.
我們多謝你給這個教會.
這個地方.
給我們身邊有弟兄姊妹.
一個屬靈群體.
一起學習去支持你.
一起學習去認識你.
去事奉你.
父母我們求你幫助我們.
幫助我們重新的去思想事奉.
這件事.

$^{641}$讓我們明白這個是.
我們可以經歷你.
經歷你的同在經歷你的祝福.
經歷你的帶領的一個回應.
求你幫助我們讓我們事奉裡面有一個喜樂.
幫助我們事奉裡面可以見到你自己的真實.
幫助我們在事奉當中.
去更加的認識你.
和認識我們自己.
求你幫助我們撇去我們心裡面的偶像.
求你幫助我們專一的.
單單的.
單純的去.
去到你面前將我們有的交上.
無論是金錢,時間.
甚至我們的生命.
好讓我們.
在這個世上.
經歷到你給我們的豐盛.
神父我們多謝你.
感恩祈求,奉主的聖名.
謝謝大家.
\newpage



\section{}
\label{sec:UcXT0Snl5mA}
\textbf{5/19/19 主日講道\_耶穌天國的小故事-3\_麵酵 (陳崇基牧師)}
\newline
\newline
連結: \href{https://youtube.com/watch?v=UcXT0Snl5mA}{\texttt{https://youtube.com/watch?v=UcXT0Snl5mA}} ~~~~ 語音日期: 2019-07-22
\newline
\newline
\hyperref[sec:QRMnvNFd6QE]{\small{< < < PREV SERMON < < <}}
~
\hyperref[sec:index_chronic]{\small{[返順時目]}}
~
\hyperref[sec:index_scriptual]{\small{[返順卷目]}}
~
\hyperref[sec:Kxw1fL4mQTc]{\small{> > > NEXT SERMON > > >}}
\newline
\newline
$^{1}$我們現在打開聖經.
去到《馬太福音》第十三章三十三節.
《馬太福音》第十三章三十三節.
我們繼續去看耶穌基督講的比喻.
以及一些小的故事.
這裡只是一節的經文.
這裡他說.
天國好像面孝.
有婦人那來.
藏在三斗面裡.
直等到全團都發起來.
很簡單的一個比喻.
我們講過看耶穌的比喻的時候.
不是單單聽一個衣索如燃的故事.
而是要留意故事裡面.
到底要挑戰.
以及讓我們不舒服的是什麼地方.
那種比喻的挑釁性.
這裡講到上帝工作.
是好像那些面孝.
我不知道中間有誰做過麵包.
親手做過麵包.
就是真的要搓麵粉那些.
那些面孝就是讓麵粉發大.
因為我自己也獻身做了一次麵包.
所以我想大家去看看.
到底我學了什麼.
大家去看看.
(笑聲).
做麵包真的不容易.
燒焦了是吧.
不想的.
我也是第一次算是這樣.
我覺得很有意思.
上帝給我們一些啟示.
耶穌講的比喻.
很多是有機的比喻.
你記得我上次講過芥菜種的比喻.
這次我們講麵孝的比喻.
其實也是一樣.

$^{41}$講到麵孝的時候.
就是在廚房裡面常常會看到的東西.
大概可能猶太人或者耶穌講的麵孝.
不是我們今天在Whole Foods裡面買到一包的麵孝.
而是每次搓麵粉留下一塊.
明天又再搓要用回那塊.
那叫做Leaven.
就是那種麵孝.
但是做的功能其實也是差不多.
就是要用麵團來發起.
如果你說到麵孝的時候.
在猶太人的世界裡面.
一提起這件事.
他們會想起一個故事.
就是想起出埃及記的故事.
你記得有一個叫做無孝餅節.
或者是御節之後的七天.
這個就是猶太人常常會想起.
一想起麵孝的時候就想起無孝餅.
無孝餅就是平時我們吃一頓飯的時候.
我們拿的那個餅.
無孝就是好像那些Cracker.
那些餅乾那樣.
因為沒有發過.
為什麼會想起這個故事呢.
因為當時以色列人他們在埃及記那裡.
他們是受了苦很多年.
做了奴隸很多年.
然後上帝用恩典來拯救他們出來.
用十災來擊打法老的暴權.
將這個政權來擊打之後.
就去….
以致最後法老都說不如你們走吧.
因為我們受不了是不是.
那麼你看到這裡.
譬如說第34節.
百姓就拿著沒有孝的生麵.
將那些盤子包在衣服裡面.
然後就趕快托著肩膀就走了.
在夜更裡面就起行離開.

$^{81}$所以為什麼他們會守無孝節呢.
無孝餅呢.
吃無孝餅.
因為你沒有時間再等.
你知道整個麵包.
你要等那些孝來發起的時候.
很多時候要等很多個小時.
很長的時間.
他們出埃及的時候就要趕著來抓人.
就是沒有任何時間讓他們來等.
所以每逢去到御節的時候.
就是我們大概受難節和復節的時候.
猶太人的家庭就會大掃除.
就像我們華人新年.
我們大掃除一樣.
將所有這些麵包就拿走.
除去它.
因為這些麵包就是給予腐敗.
不潔的東西.
或者是擔言我們的東西等等.
所以在福音書裡面.
耶穌都用麵孝 "yeast" 或者 "leaven".
來形容一些腐敗的東西.
你記得耶穌跟門徒的一些對話裡面.
他曾經說過要防備法利賽人.
殺刀該人的孝.
這些宗教領袖的孝.
就是他們說的話.
他們的教訓.
你要小心.
因為很多時候都是雙重標準.
或者自以為是.
是一種對律法的選擇性.
方便自己.
合自己心意的就要走.
不方便的.
或者福音書經常提醒我們.
這些宗教領袖.
他們是貪錢貪權的.
他們就將這些東西.

$^{121}$就當上帝舞蹈.
行公義,好憐憫等等的事情.
他們就當上帝舞蹈.
耶穌就說你要防備這些孝.
防備這些滲入.
或者這種令我們腐敗的東西.
但是在這一處.
耶穌說天國上帝的工作.
他的比喻.
他是用面孝來比喻.
怎樣看也好.
在猶太人的心目中.
都是有一點點顛覆性的.
有一點點很奇怪.
就等於今天我們可以說.
上帝做的工作.
好像是小強.
打也打不死的.
或者上帝的工作.
天國就好像電腦的病毒.
你怎樣剷它也剷不走.
還有散播得很快的.
有一點點這樣的感覺.
你記得我前一次說.
這個芥菜種的比喻的時候.
其實就是這個比喻.
之前的比喻.
所以這兩個比喻.
都是大家合起來看.
他說天國就像芥菜種.
百種裡面最少的有人拿來.
然後長成之後變成很大的.
甚至好像樹木一樣.
飛鳥可以停在樹枝上面.
第二個比喻就是今天這個.
就好像面孝.
有婦人拿來.
藏在三豆麵的裡面.
然後等到全團的麵團來發大.
這些比喻我們可以叫做對立的比喻.

$^{161}$小變成大.
或者見不到的東西.
變成肉眼很容易很明顯看得到的.
面孝藏在麵團裡面.
當你搓一塊麵粉的時候.
那個麵粉有孝和沒有孝.
開始的時候你是看不到的.
你給時間去等待的時候.
你就能夠看得到那個分別.
開始的時候是無法看的.
但是經過這些孝素來滲透的時候.
你的肉眼就能夠分得出一個已經發了孝的麵團.
他說藏在三豆麵裡面.
三豆麵英文叫做three measures.
其實大概是四十至六十磅的麵粉.
如果你有做過這些搓麵粉或者做麵包的.
你會知道四十至六十磅的麵粉其實是很多的.
餵飽可能有一百人.
會成為一個很大的現值.
皆宰種由小變成大.
由謙卑開始.
變成一個擴張到四周都是黃花的景象.
面孝的比喻.
是由一個不顯眼.
但變成一個影響全團發大的功效.
我想比喻最重要的地方.
就是講到面孝的影響力.
那種滲透力.
上帝的工作.
上帝的影響.
很多時候是在我們人的內心裡面發生的.
不是外界強迫.
以致有改變.
是在裡面發生的.
在裡面發生面孝的作用.
以致整團發大.
滲透到整團.
福音或者天國的工作.
其實都是從我們內心開始的.
當我們接受了上帝的恩典.

$^{201}$由我們的生命裡面.
從裡面出來.
在上個星期.
在美國有一個很爭議性的法律的新聞.
大家可能有聽過.
在阿拉伯州.
通過了有史以來最保守的反墮台法案.
我想最終目的是挑戰1973年的時候.
有一個很出名的.
Supreme Court的一個決定.
一個裁定是使墮胎合法.
法國叫做Grover Way.
這個新聞是一個很大的爭議.
也是令人很激動的.
因為你說到墮胎.
在我們中間.
我們有不少.
我們是做媽媽的.
我們是做父母的.
我們一定會覺得這是一個很個人的.
很感情的一件事.
但我想說說.
有一點歷史的來由.
是我們不能夠不知道的.
這種墮胎的法案.
是由美國裡面一派.
叫做宗教右派.
來興起.
或者religious right,Christian right.
這一班人來興起.
為什麼需要去解釋一下背景呢.
如果我們不知道內容的去脈.
有時候我們不是很掌握得到整幅圖畫.
我相信對這方面有研究或者解釋的.
就是這本書叫做《大Kingdom Come》.
是一個福音派的信徒.
叫做Randall Balmer.
他是一個學者.
也是神學院的教授.
也是一個牧師.

$^{241}$他寫了不少近代的歷史政治倫理的書.
或者一些背景的東西.
他寫這本書是為什麼呢.
他想去解釋一下.
到底他有什麼哀歌.
他作為一個福音派的信徒.
他這本書的主旨就是說.
這些宗教的右派.
極保守的右派.
是怎樣不知不覺地扭曲了信仰.
或者將信仰重新重建.
然後威脅了社會和政治之間的關係.
其實他追查.
為什麼這個宗教的右派.
會和現在的所謂Republican.
有這麼密切的關係呢.
話說在七零年代的時候.
保守派的人.
他們想去強化他們在社會的政治地位或者權力.
曾經用過不少的.
不同的一些倫理議題.
或者社會的議題去號召民眾.
號召很多的選民.
拿一些vote(投票).
發現最後找到一個議題.
是可以號召到很多選票的.
那個就是反墮胎.
只要出一種呼聲.
很多人就會出來.
有很多人會用一個很單一的議題.
來支持某一個議題.
你要知道這些背景的原因.
因為很多時候的爭議.
就是回到到底真正想做的是什麼.
到底是想拿political power(政治權力).
還是真是在信仰裡面.
我們有一個很整體的.
想有一個見證.
很多時候這些情況.
你追溯那個來龍去脈的時候.

$^{281}$給外界的人看.
就會覺得很多時候.
這些右派.
就會很容易是以偏概全.
你要知道.
我想告訴你.
早在一九八零年代的時候.
七零年代開始.
你知道那個Wall-E-Way(反墮胎).
是七三年的時候裁定的.
墮胎合法化.
七零年代的時候.
各自有那個Wall-E-Way之後.
起碼五至七年.
如果你看當時所發的文章.
《牧者法》的文章.
或者倫理.
甚至是右派裡面發表的文章.
有很多都是贊成墮胎的.
有部分是反對.
但是有不少是贊成墮胎.
在一九八零年代初的時候.
當時的Christianity Today.
今日基督教.
這個雜誌.
其實他已經出了一篇文章.
Christianity Today是一個很保守.
偏右派的基督教雜誌.
這個雜誌出了一篇文章.
他說什麼呢.
他說教會.
你要行公義和好憐憫.
你要反貪婪.
反貪婪金錢.
反對富者越富.
貧者越貧的現象.
要幫助弱勢的人.
等等.
他說這些的命令在聖經裡面.
差不多每一章都有.

$^{321}$但是墮胎反墮胎呢.
不是那麼多.
上個星期.
這個Alabama的.
推出這個法案的時候.
你要知道.
基督教裡面.
一個很保守的代表人.
叫做Pat Robinson.
他也出來說.
這個行動太激進了.
這個推出的法案.
太過restrictive.
但你要知道.
Pat Robinson在基督教裡面.
保守派裡面.
是表表者.
如果老人家都說這個法案有問題.
我們就需要想清楚一點.
當然.
我不會反對.
聖經裡面大體上是pro-life.
聖經裡面大體上.
是擁護生命權.
但這本書指出.
為什麼呢.
他說保守派的信徒.
往往有一種叫做.
insufficiently pro-life.
就是不足夠地.
或者選擇性地.
去擁護生命權.
你問一問自己.
反對墮胎是維護生命.
是不是.
但同時可不可以.
也去維護.
有死刑有好刑呢.
歧視.
歧視一些小眾.

$^{361}$或者移民.
或者擁護白人至上的.
種族主義.
或者甚至發動戰爭.
如果你這樣想的時候.
我們就不可以純粹.
或反墮胎.
而我們呢.
發動戰爭就OK.
上兩個.
三個禮拜.
我就.
被邀請.
我和一班的牧者.
去了見一個.
本地的一個議員.
我們走到議員的辦公室.
我們就說出.
我們為什麼對這個.
墮胎合法化的法案.
有一些保留的關注.
我經常都問的.
有時候.
要出席這樣的活動.
我自己都會問自己問題.
為什麼反墮胎的時候.
就號召很多人出來.
但是當美國入侵其他國家.
發生戰禍.
令到整個地區散了的時候.
就沒有人出聲了.
這個我覺得是我們需要問的.
我們不能夠這樣閉上眼.
我們只能選擇性地.
去看一個議題.
到底我們POLIFE.
有多POLIFE.
是不是我需要用個handheld mic.
.
那麼.

$^{401}$巴馬說什麼呢.
我覺得有一句.
幾句說話我想引用的.
他說在這本書的結論裡.
他說他認為.
減少墮胎.
最有效的方法是通過教育.
好像推廣一些教育的運動.
或者類似.
去禁吸煙.
禁吸毒.
或者瘋酒配偶虐待等等.
這些公共服務.
來改變這個纏繞問題的整個大氣候.
這種策略.
比較侵犯個人權利.
無效的法律.
就是由愛來壓迫.
由愛來影響.
是更加.
讓我們進一步去.
廢除墮胎.
這個就是在說滲透力.
基督徒花了多少時間和精力.
真的要做這件事.
在一個反墮胎的網頁上.
我看到有一個人留言.
他問了一個問題.
他說在這裡留言的人.
有多少人曾經.
助養過一個小孩.
或者肯助養一個小孩.
肯幫助一個.
在一個不受權利的地方出來的人.
有多少人願意幫助一個人.
還是我們只是口說反墮胎.
這樣就夠了嗎.
和這個麵包的比喻有什麼關係.
我覺得有關係.
巴黎在最後的結論.

$^{441}$他說當信仰獨立於政治的權勢的時候.
才能夠發揮最大的作用.
初期教會的信徒.
是反文化.
他們是反文化.
是因為他們當時和當時的政權是分開的.
正是因為分開了.
沒有任何利益的關係.
就能夠去批判強權者.
因為他們是連結的.
disinterest.
沒有投身在一個權貴的架構裡面.
事實上歷史裡面.
告訴我最有效最有活力的.
宗教信仰的運動.
都是和弱勢認同.
將自己定位在社會的邊緣.
而不是在政權的中心.
你想一想.
耶穌說的比喻.
很多都是這樣.
是社會的franchise或者margin.
不是去靠攏一些地上的權貴或者政黨.
這是對我們一個當頭棒喝.
到底什麼時候我們專門靠一些政黨.
或者靠國家的權貴.
去做神的工作.
上帝的工作.
耶穌說什麼呢.
是等於面孝.
是不明顯的.
是從裡面滲透的.
沒有政治的強勢.
是隱藏的.
是由人的內心開始發出來的.
天國就是面孝.
我們就是面孝.
在各自不同的地方.
我們將我們信仰給我們的能力.
愛和恩典.

$^{481}$都去影響.
滲透在這個世界當中.
這就是神做事的方法.
不是以外向內.
而是以內向外.
說到政治倫理等等.
我一定不是專家.
我們中間有一個專家.
張國棟博士.
在我們中間.
張國棟博士今天飛回香港.
他會去神學院和不同的教會講學.
他七月尾就會回來.
經過芝加哥.
我們會特別邀請他講講這些專題.
這些是值得我們去問的問題.
我不是想今天講一些.
只是yes或者no.
而是我想擴大討論的領域.
讓我們多問一些應該問的倫理問題.
天國的問題.
到底神做事是怎樣做.
我們最精會神要做的事是甚麼.
這是第一點.
面孝是一種滲透力.
是內之外.
不是靠一些強權或勢力去做事.
而是我們自己各自.
在我們社區.
在世界裡面去做事.
最後一樣我相信面孝可以告訴我們的是甚麼.
就是需要時間.
我第一次做麵包.
那個video.
我以為兩個小時就搞定.
原來這樣看太幼稚了.
因為那個星期六的中午我們開始.
當然我們set up camera等等.
其實也不是浪費很多時間.
因為很低製作.

$^{521}$我以為要兩個小時就搞定.
原來要幾個小時.
誰知一搞就搞到我吃晚飯.
搞了五六個小時.
六七個小時.
就是等那些孝嫂去做他們的事.
急不來的.
我不可以用風筒吹他們.
然後叫他們快點去做事.
現在我明白了.
為甚麼那些餅店的工作人員.
要早上三四點起床.
去搓麵包.
現在我明白了.
原來那一元的麵包其實也不是那麼貴.
天未亮真的要做.
你要知道上帝的工作也是這樣.
不是很壯觀放煙花.
不是我們人可以去控制的.
我記得我女兒我們開始這個church plan的時候.
她星期五晚開始了一個青年團契.
現在說回來的時候.
她開始的時候有點失望.
因為開始的時候只有十個人.
八個人.
她以為是一個很大的turn out.
但學習了甚麼呢.
是需要很多耐性.
很多時間.
不是放煙花.
麵包是需要時間來工作.
神的角度其實也是一樣.
往往都是看不見的.
是一些無名氏在背後來工作.
我們就是麵包.
我跟你們說過.
我最近一年在西南面一間中學做義工.
做翻譯的義工.
那裡有一些華人的家庭.
一些新移民的家庭.

$^{561}$你說辛苦嗎幫人家做這些翻譯.
做的時候不辛苦.
但時間是需要的.
因為他們有report card day的時候.
我要在那裡等.
這個就大半天的時間.
你不要以為我們監獄沒事做.
我兒女以前就以為我沒事做.
你只是做一個上半天的工作.
不錯啊.
平時在家裡喝下咖啡.
不是那麼簡單.
一花了你大半天的時間和精力.
但我覺得是有教無類.
特別是一些新移民的家庭來的時候.
是需要幫助他們.
這些就是麵包的工作.
上帝就說天國好像麵包.
你有沒有做一些這樣的工作.
可能是一些隱藏的工作.
你有沒有讓上帝的能力.
在你內心裡面開始發起.
以致影響你自己的生命.
你的生命可以影響其他人的生命.
我們同心低頭祈禱.
上主我們多謝你.
因為你是我們的神.
主我們求你幫助我們.
讓我們每一個小信徒.
在我們這個世界.
在我們這個社區.
能夠成為面孝.
以致更多人被你感染.
你先用你的福音.
你的恩典來感染我們.
以致我們的生命.
我們的家庭.
我們的婚姻.
我們和兒女的關係.
我們和家人的關係.

$^{601}$我們和朋友的關係.
我們的工作.
我們的工場.
我們的學校.
能夠使更多人認識主.
讓我們的教會.
能夠成為一個面孝的教會.
是你自己所用的.
能夠去影響這個社區.
影響這個世界.
祈禱奉主耶穌基督的名字.
原上祝福各位.
我們11點半鐘.
我們也會有主日學時間.
我們希望能夠見到你.
\newpage



\section{}
\label{sec:Kxw1fL4mQTc}
\textbf{7/8/2018 主日講道: 基督徒的婚姻 (吳君才牧師)}
\newline
\newline
連結: \href{https://youtube.com/watch?v=Kxw1fL4mQTc}{\texttt{https://youtube.com/watch?v=Kxw1fL4mQTc}} ~~~~ 語音日期: 2019-07-23
\newline
\newline
\hyperref[sec:UcXT0Snl5mA]{\small{< < < PREV SERMON < < <}}
~
\hyperref[sec:index_chronic]{\small{[返順時目]}}
~
\hyperref[sec:index_scriptual]{\small{[返順卷目]}}
~
\hyperref[sec:me3DczkMU3E]{\small{> > > NEXT SERMON > > >}}
\newline
\newline
$^{1}$我們先有個簡單的祈禱.
多謝主讓我們再重思.
你自己的婚姻的藍圖.
家庭的藍圖.
讓我們再想一想.
我們有什麼需要更走在你的旨意裡.
以下時間我們分別出來.
聖靈你提醒人道.
奉主耶穌基督明求,阿們.
今日的經文是《馬學福音》十章二到十二節.
我們很快地看一看.
說到有些法理賽人來問耶穌.
優妻可以嗎?.
然後摩西怎樣回答.
其中摩西說因為你們心硬.
繼續第七節.
神說婚姻有一個很重要的訊息.
就是人要離開父母,與妻子聯合.
二人成為一體.
所以就有兩個.
耶穌又面對一些人的挑戰.
祂又說優妻另娶的凡奸人.
姑父是妻子的.
如果妻子離棄丈夫另嫁.
亦都是凡奸人.
再看下去.
我要先說婚姻的另類.
獨身或單身的問題.
以前在中神任教的劉達方博士.
寫過一篇文章.
對獨身觀念的兩個常見誤解.
在這篇文章裡面提到.
另一半這個名稱都不知道從何而來.
或者因為聖經所說.
人要成為一體.
所以未結婚的人似乎只有半個人.
或者因為在人的潛意識當中.
未婚的人是很不完全.
這幾句說話點出現在一些人的典型想法.
所以有些長輩或不少父母.

$^{41}$對已經適婚年齡的子女或親人.
耳濡目染.
有時甚至單單打打.
最慘的是過時過節的時候.
令未婚的人避之則吉.
不想見太多親友.
有一本書叫《結伴同行》.
有些人稱它是單身一族的錦囊.
它羅列出很多單身人士所遇到的掙扎.
或者在在的單身人士曾經經歷過.
或者你結了婚.
以前做單身人士亦都曾經歷過.
包括從實況裡面將傳說分開.
第二從壞的勸告裡面要過濾一些好的東西.
第三應付父母和同伴的壓力.
第四從愛在的期望和個人優先的次序要作出選擇.
另一位在美國多年研究單身問題的專家.
他又列出單身人士其實有很多好處.
或者一些理由.
包括有很多個人的時間.
可以自由去追尋自己的興趣或者理想.
以及自己的專業.
又因為不急需要籌備結婚家庭的問題.
可以很自由自在地去旅遊.
或者自由自在地去生活.
又可以隨意選擇自己的生活方式.
有更充裕的時間.
可以做什麼?.
看書,尋索和發揮個人的潛質.
有更多的機會甚至去探險.
有更多的機會去認識其他的朋友.
相對來說壓力沒有家庭的負累是少一點的.
所以可以有更多的安靜時間.
這些都是好處.
但是其實如果我們研究看一看聖經和神的心意這個藍圖.
我們有時候都不會那麼困擾.
反而得到一些釋放安慰.
父神的恩典其實是以不同的方式臨到你和我的.
對一些人來說神賜給你婚姻.
讓你享受婚姻的好處.

$^{81}$但是對另外一些人來說.
祂就給他有獨身的經歷.
哇!那神就很不公平.
不是.
神是知道你受得.
才賜給你.
獨身的好處.
其實都是一個恩典.
你從另一個角度來看.
有些未婚的人被取笑是單身貴族.
聽起來很瀟灑.
很自在.
無拘無束.
但是如果我們用屬靈的角度去看.
獨身是一個恩賜.
就是父神在他的大愛和恩典底下.
安排人獨身.
或者一個單身的生活.
這種體會.
一方面讓你和我明白.
單身,獨身.
不是表示少了一種福分.
其實是和已經結婚的人一樣有福.
另一方面.
既然獨身是從神來的恩典.
我們就應該要學習用感恩的心.
去領受和接納.
學習不要這麼容易發怨言.
怪責神.
第三方面.
人也不需要再苦苦說說.
自己有沒有長期獨身生活的本領.
其實只需要若過一日.
父神沒有改變我這個單身的境況.
我就是蒙了神賜給我這個單身的恩典.
這是一個表記.
直到他日.
神要改變我這個境況.
賜給我一份人.
這也是另一個恩典的表象.

$^{121}$我個人來看.
人只要生活得愉快喜樂.
心中和諧安穩.
要能夠積極發揮潛能.
和群體有好好的相處.
守住崗位.
做好本份.
這樣結婚也不是大問題.
也沒有妨礙神在你的生命中定下的旨意.
請本聖經要你很小心留意.
神並沒有白紙黑字列明人一定要結婚.
神並沒有說.
人二十歲以後就該住二.
三十歲以後就不要結婚.
無論得時不得時.
總要專心找個配偶.
並用百般的辦法.
吸引人.
取悅人.
務求發財.
聖經沒有這樣說.
聖經也沒有這樣說.
人若不結婚.
就不能作我的門徒.
聖經只要求我們盡心盡意盡性盡力愛主.
我們的神.
又要愛人如己.
所以說.
神並不是要求我們每個人盡一切力量去找一個歸屬.
可以說.
結婚和獨身.
同樣都是有價值的.
同樣都是有意義的.
所以未婚的人.
或者現在還是單身的人.
不要為此覺得羞愧不安.
或者覺得難過.
因為在聖經上看.
神性的法則.
那個precaptive view of God.

$^{161}$這其實是對神.
擺我們有不同的帶領.
唯有神的法則的旨意.
就像剛才所說的.
盡心盡意盡力愛神.
愛人如己.
這就是屬於神的法則性的旨意.
人人都要遵守.
你很熟悉的十戒.
只屬於法則性的旨意.
其他是神的帶領的旨意.
明白神的旨意.
這是另一個課題.
我們要探討.
神對你對我.
是會把最好的給你給我的.
今天你信不信這件事呢.
神的帶領是因人而異的.
無法比的.
不能與旁邊的人.
或與你一起成長的人相比.
為何大家同樣的資歷.
這樣的條件.
他有我又沒有呢.
我們要學習.
安心.
守住自己的身份.
不要覺得自己是不幸或不圓滿.
要自要緊.
明白神在你身上的旨意.
我們經常說.
萬事互相效力.
叫愛神的人得益處.
一切的好處不會神意外.
究竟神的心意是怎樣呢.
只要我們能夠做新造的人.
能夠效法神而止.
耶穌基督的榜樣.
我們全然成聖.
這反而是我們人生要追求的大前提.

$^{201}$我們要信靠耶穌.
而邁向一個完整.
今天我們再想想甚麼是信靠.
你說你信耶穌.
你靠耶穌.
是甚麼意思呢.
信靠其實是包羅你.
信靠神有絕對的主權.
包括你現在的婚姻或生活狀態.
信靠是包括你相信神有大能.
信靠是包括你.
信得過神最好的智慧部署.
信靠是包括你仍然信得過.
神在你現在此一刻.
生命仍然是有關懷和慈愛.
所以我們能夠學到.
在神面前好好的禱告.
反省自己心裡有甚麼疾度.
憎恨甚至不滿.
以及種種的疑惑.
甚至對未來那種憂慮.
那種孤獨.
或被人誤解的憤怒.
你和我都要一一向神陳寗傾訴.
神一定會垂聽.
神一定有最美好的安排和引導.
我很同意從觀察所得.
有些專家說.
為何有些獨身的人.
單身的人.
他們的人生好像支離破碎.
或很不開心.
他們不時有埋怨.
甚至會疾度.
又會自我中心.
他們裡面充滿憤恨和怒氣.
因為他們的生命被苦劫所腐蝕.
但同樣的,為何有些人.
同樣是單身的.
仍然覺得活得很喜樂.

$^{241}$很精彩.
生命很充實.
因為他有目標和人生的方向.
因為他也只得自己去幫助人.
關心人.
分別在哪裡呢.
一個有甚麼.
而另一個就沒有甚麼.
不是,他們的條件和各樣東西都差不多.
只不過因為一個肯信服.
肯依靠神.
肯面對神的帶領.
他不生神氣.
不向神發脾氣.
另一個真的信得過神.
無論如何.
在我生命上.
總有美好的帶領.
能夠仍然以神為樂.
所以未婚的弟兄姊妹.
這些對你和我有甚麼啟發呢.
獨身的問題.
婚姻的另類.
我到這裡點到即止.
主題是講第二項.
有關美滿的婚姻.
我已婚多年.
我和師母不斷在學習.
我做牧師.
我也有很多婚姻上生活上的軟弱和挑戰.
有時真的覺得自己是壞人.
但問題是你和我都有虧欠.
但去到神面前.
可以靠著恩典去面對.
今日我們再重溫一下.
反省一下.
無論你結婚多少年.
你也要再想一想.
神給我們有三個美滿婚姻的要素.
第一,要離開父母.

$^{281}$大家都明白離開父母不是拋棄父母.
不是不理自己的父母.
不是這個意思.
也不是說自己有無成家立室.
就特意搬到老遠的地方居住.
離開父母.
那是指你自己在心智上和行為上.
能夠自立,自主,自住,自治.
管治這個「自」.
這個幸福是由你和另一半掌握.
意思是你離開父母.
不再只依靠父母的任避,呵護.
或甚麼都只依靠自己以前的父母.
而是要在你和另一半新的家庭中做抉擇.
種種生活細節當中.
離開父母其實是指身份,角色和人際關係的轉移.
開始以自己和你的配偶的感受.
作為一個很重要的考慮.
你不應該像以前未嫁或做少爺時在家中的表現.
你也不是拐彎抹角不理家務的心態.
離開父母是表示你能夠自治和承擔更大的責任.
當然你仍然要做一個好的兒子,好女兒.
但你仍然要在新組合的家庭中盡心做一個好丈夫或好妻子.
如果你有了子女後你更要學做一個好的父母.
很多婚姻問題出現.
就是不少人仍然受著我們輔導有的原生家庭的影響.
你未改變和能夠脫身.
以致有些扭曲或負面的看法太強.
你的感情和作風仍然跟父母那一套.
被他牽著走亦不自覺.
我抄一段有代表性的東西.
因為他的母親太完美.
以致他將過高的期望加諸妻子的身上.
他心理上對母親的依戀.
使他難以接納與母親截然不同的妻子.
還有很多東西,先讀到這裡.
無論是男是女.
如果你的心態有一種閉塞或不開放,不調教.
你過高未來開放的期望.
會構成你潛伏的壓力和婚姻的難處.

$^{321}$不過我又講另一個角度.
如果你是做父母的.
你要學習接納與被子女離開的事實.
雖然我們知道天下父母心.
確實是難捨難離.
自己的子女結了婚離開.
但千萬不要霸佔或遙控.
想子女這一生一世黏著你.
在你的附屬底下.
我曾經知道或遇到一些做長輩的父母.
他亦因為離不開父母這個聖經給我們的觀念.
令到婆媳和子女極度難堪的局面.
特別是一些父母和子女一起住.
很多時候很容易引起婆媳的問題.
確實是聽過一個典型的例子.
他是一個作母親的.
和子女一起住.
子女又有結婚.
新婆又來了.
但離譜到要求已婚的兒子.
你們的房間不要鎖門.
因為我習慣像以前那樣.
到時候晚上看兒子有沒有蓋.
被子蓋不好.
還有吃完飯.
我們的家規就是要坐在一起.
吃水果,看電視,聊天.
才好生活.
但做媳婦的.
根本不喜歡看電視.
很浪費時間.
我只喜歡看書.
你想想這種價值取向.
有時我們做父母,做子女.
有時我們真的要學習這一點.
如果我們想子女更有自知自願的能力.
要獨立.
在心態上,在性格上.
要離開父母.
我們就要多給一些.

$^{361}$Let go and let God.
即是說放心和放手給神.
第二樣我們要看.
與妻子聯合.
這裡時間多一點.
說到聯合方面.
基督徒的婚姻應該在主內.
甚至和信仰相同.
以及差不多靈命的人結合.
是最理想不過.
不過世事沒有那麼容易.
有一篇講到與不信主的人結婚的文章.
是楊木谷牧師所寫的.
他在結論時寫出這樣的文字.
他說血戀中的人的世界.
只有一個人.
不是兩個人.
因為他連自己都忘掉了.
只有對方.
羅密歐與朱麗葉.
若果是真實人物.
他們能夠以結婚而終.
也必定是走著同一的道路.
二人若有相同的宗教背景.
婚姻關係可以自然地發展下去.
不然的話.
一方若是以宗教信仰為他人生的期盼和價值.
而另一方則對宗教冷漠甚至反對.
他就會覺得對方是反對他這個人.
更悲哀的是.
反對他的是他最愛的人.
與不信主結婚的苦果.
就越加苦澀.
正如說來.
信與不信不能同父一家.
不騙就不僅是一個命令.
也是一個結論.
因此基督徒的婚姻聯合.
是很重視神.
重視神的心意和信仰.

$^{401}$是有共同的生命結合.
而這種與妻子聯合.
是不應該在婚前或婚外.
今天拍攝的基督徒.
或已婚的弟兄姊妹.
我們要警惕和克制.
否則我們很容易產生淫亂的行為.
也不把父神的旨意放在眼內.
記得有一次.
在我牧羊和關貴一個有教會的男士的過程中.
有一天.
忽然有人按我家門鈴.
打開門一看.
是一個中年婦人.
這位婦人怒氣沖沖.
原來是找上門來聲討我.
你不要管那麼多.
高過人的事.
婚姻問題你不要管那麼多.
你不要插手.
臨走之前.
這位中年女人拋下一句.
我自己也有老公.
每個人都出來玩.
我當場呆了一會.
也覺得很痛心.
你有妻兒.
她有丈夫.
仍然是恬不知恥.
惡向膽中生.
沉溺在罪之中.
讓自己.
甚至自己的配偶.
自己的家庭婚姻.
被情慾和罪來蠶食.
破壞甚至徹底撕裂.
與妻子聯合.
或與丈夫聯合.
這是聖經的真理.
很慘.

$^{441}$其實我們明白到一個很真實的事.
進一步說.
與妻子聯合.
就是作為丈夫.
就是真的要好好和自己合法的妻子聯合.
而不是另外和另一些女人聯合.
不少出現問題的時候.
就是男或女.
另外與其他人的聯合.
我們要認定這裡的聯合.
真的需要付出決心毅力.
是以堅誠,忠誠,愛.
全程投入毀身對方.
其實在香港,紐約,芝加哥.
我曾經輔導不少和關顧很多的人.
其中有些已婚人士.
她明知自己有妻子,丈夫.
卻是另外有外遇.
她發展一些很畸形的關係.
與婚姻以外的異性結合.
令到帶來更多的痛苦和悲愴的事.
結果婚姻,家庭破裂,離婚收場.
聽起來令人痛心.
其中一些是透過上網.
現在看到很容易.
在網絡世代.
科技很先進.
交友的網站很多.
原來這些魔鬼的伎倆.
蠶食很多人.
所以發展成網戀.
網戀的關係.
可以約,何時見.
你在外地,我去找你見.
結果慌張自己的私慾感情.
導致不可挽回的慘痛後果.
其實婚姻這種聯合.
按神的心意.
是一個不可輕易改變的盟誓.
夫婦都受到這個規範.

$^{481}$所以無論任何的境況.
都要忠於對方.
不論什麼環境.
互托終生.
是一個義無反顧的毀身.
所以很多人在婚姻的誓詞裡都很熟悉.
這個不是隨口說的.
是一生的承諾.
誓詞說無論健康或疾病.
貧富或富裕.
快樂或痛苦.
喜樂或悲傷.
順境或逆境.
意見相同或分歧的時候.
仍然終身結合.
互相隸屬.
終身結合.
這個togetherness.
連結在一起.
在婚姻裡真的很重要.
很多發生問題.
導致離婚.
甚至不和的夫婦.
你會發現一個很顯著的徵象.
你可以留意一下.
甚至自己反省一下.
就是少了togetherness.
很少在一起.
而是各顧各.
各做自己的事.
各去自己的地方.
亦不想知對方.
不理對方.
聽起來各自各精彩.
少了togetherness.
這些活動和二人世界.
其實你的婚姻已經亮了紅燈.
接著就沒什麼溝通.
沒什麼對話.
同一屋簷下.

$^{521}$大家可以不秋不徐.
是,都是日常生活的小節.
有togetherness.
這是一個很重要的功課.
其中一樣.
涉及的就是forgiveness.
饒恕.
對方曾經做過一些不好的事.
或者犯錯.
你學習怎樣饒恕.
怎樣去接納.
給對方更多空間和信任去重建.
以免仍然懷恨的心.
或者找回以往的不幸.
去挑剔.
引起傷痛不和.
饒恕這件事在夫妻的關係上.
是一個很重要的功課.
而且可以聯合得更好的功課.
有機會再深入來說.
不過其實我們說.
要彼此聯合.
互相隸屬.
就要學到有什麼問題.
大家要有商有量.
有什麼關係困難的.
要跟著大家談.
而不是像一些人發展到.
發生一些事.
就大禍臨頭.
各自飛.
這麼極端.
基督徒的婚姻結合.
是按神的心意.
是今天明天.
是一生一世.
都要聯合在一起.
所以我們真的要求主幫助我們.
第三樣.
嘩!時間這麼快.

$^{561}$還有一些重要的事未說.
就是二人成為一體.
神所設立的婚姻.
是男女之間最完美的關係.
他的親密和緊密的情況.
只可以用二人成為一體來形容.
神所設立的婚姻.
其實是一個很令人感受到美好的事.
因為這是最完美.
最恆久.
最穩定的聯繫.
所以為了這個關係.
神定立了一個制度.
就是一夫一妻.
是二人成為一體.
不是三人.
不是有第三者.
是男與女.
不是男與男.
不是女與女.
而成為一體.
成為一體最基本的意思.
就是有性關係和肉體上的結合.
為了在正式的婚姻裡.
性關係是神聖,美好,美麗.
否則的話是醜陋,卑賤和有罪.
多年前.
香港突破中心.
有一個很資深的輔導.
跟我們教母同工說.
現在的潮流.
現在美國也是.
都是流行旅行.
流行同居.
但是很遺憾.
很危險.
不少未婚的熱戀中的男女基督徒.
趁著假期.
又一起去外地.
這樣二人同住一室.

$^{601}$就必有損失.
還有情到濃時.
剛柴烈火.
二人就發生了婚前的性行為.
提早了二人成為一體.
結果引致日後.
心理上,感情上出現不少問題.
而需要接受婚姻輔導.
他說這種個案是越來越多.
這位資深的輔導員.
刻意請求你們教母同工.
麻煩你們回到自己的堂會.
多多提醒教導青年的弟兄姊妹.
避免一起去旅行時.
共住一室.
這是一個很堂醫毒藥的試探.
要成為一體.
其實還有一個更深層的意思.
就是完全毀生.
完全相熟.
因此夫婦可以分享一切的所有.
無論身體上所擁有的東西.
所以那句話好像說笑.
其實是真的.
你的是我的.
我的是.
很多人說我的是我的.
我的也是你的.
這樣是完全的.
這樣分享一切.
不單止物質.
甚至思想.
意見.
你的因此才剛.
你的困難.
你的失敗.
你的痛苦.
都可以彼此契合交流.
然後有行動的關注.
同心.

$^{641}$靠賴神的力量.
今天家庭瓦解.
婚姻不美滿.
就是因為任何一方面.
太自我中心.
現在的大男人做不到.
就變成小男人.
大女人做不到就變成小女人.
結果引致很多的傾責和不開心的事.
要成為這一樣.
神的心意.
二人成為一體.
我們真的要好好學習.
養大父神給我們的功課.
其實還有第三項.
不過時間有限.
有機會再說下集大結局.
我們低頭有個祈禱.
天父你帶領我們再重新思想.
我們在婚姻的含義和心意.
我們知道有時做得很甩漏.
或者不合理的心意.
你讓我們能夠明白.
能夠回轉.
靠你的恩典.
好好改過.
向神你自己加能,賜力.
讓你明白我們的掙扎和苦惱.
特別是為了單身的弟兄姊妹.
你加倍的恩典臨到.
讓他們能夠做個開心快活人.
在神面前仍然能夠很釋放的.
發展自己的恩賜才剛.
為主而活.
已婚的他們家庭各遭遇有難處.
我們求主你力上加力.
恩上加恩.
祝福引導他們每一個人.
我們仰望你.
求你個別施恩.

$^{681}$祈禱奉主耶穌基督名求.
阿們.
加番給陳牧師.
阿們.
加番給陳牧師.
\newpage



\section{}
\label{sec:me3DczkMU3E}
\textbf{9/23/18 主日講道: 基督徒成長12步--\#10 領導事工 (陳崇基牧師)}
\newline
\newline
連結: \href{https://youtube.com/watch?v=me3DczkMU3E}{\texttt{https://youtube.com/watch?v=me3DczkMU3E}} ~~~~ 語音日期: 2019-07-23
\newline
\newline
\hyperref[sec:Kxw1fL4mQTc]{\small{< < < PREV SERMON < < <}}
~
\hyperref[sec:index_chronic]{\small{[返順時目]}}
~
\hyperref[sec:index_scriptual]{\small{[返順卷目]}}
~
\hyperref[sec:lPh_PhzFDQI]{\small{> > > NEXT SERMON > > >}}
\newline
\newline
$^{1}$今年我們有一次的《屬靈成長》十二部系列.
我們由第一次開始.
我們說我們作為基督徒.
我們跟隨耶穌的時候.
有什麼十二方面我們需要去成長呢?.
今天我們說到已經是領導侍宮了.
我們用我們的恩賜來事奉之餘.
我們學習做一個領袖.
學習做一個領袖.
每一個基督徒在一種普遍的意義上.
我們都是領袖,都是領袖.
所以我們沒有選擇的.
因為你是一個基督徒的時候.
你就帶給別人影響.
我思考了一下今天要怎麼說呢?.
因為又要做這麼多報告.
剩下一點點時間.
但我想告訴你.
事奉可以用很多方法去說.
我們在聖經裡面說過.
每一個基督徒有他屬靈的恩賜.
Spiritual Gift.
我們可以用我們的屬靈的恩賜.
我們的Gift來事奉我們的神.
有些人行政很厲害.
有些人關懷很厲害.
有些人可能唱歌很厲害.
他們有不同的恩賜.
我們大家放在一起的時候.
我們就需要去服侍.
但說到領導的時候.
我怎麼說呢?.
我思考了一下.
其實最重要的分別在哪裡?.
作為一個領袖.
其實是一個生命裡面基本的態度有所不同.
一個領袖和一個非領袖.
其實最基本我們說的是一種面對生命的態度.
面對生命的態度.
我覺得最好就是找聖經裡面的所謂領袖.

$^{41}$看看他們.
到底上帝怎麼容他們和怎麼呼召他們.
如果回到聖經.
從頭回到尾.
我想沒有誰會說.
這個人不是領袖.
這個人叫做摩西.
Moses.
在《摧哀局記》裡面.
去介紹這個民族的英雄.
所謂的以色列國民的領導.
我想他在這麼多個領導裡面.
除了軍王,大衛等等.
他應該是別標者.
所以在他的刻畫裡面.
或者在他的技術裡面.
能夠告訴我們.
作為神的子民.
到底我們應該有什麼心態.
有什麼態度.
我們要帶入我們的生命裡面.
以致我們成為一個領袖.
摩西他自己一生.
我知道在《摧哀局記》.
他有好幾個所謂的定點時刻.
定點時刻.
在人生裡面.
基督徒的人生裡面.
是一些所謂很定義性的時刻.
領袖的時刻.
你想一想你自己.
信了耶穌這麼久.
有什麼定點時刻.
是你覺得我要站出來.
來事奉的呢.
你想一想什麼時候.
還有發生了什麼事.
然後我們看看摩西.
摩西所謂第一次出現的時候.
在《摧哀局記》第二章.

$^{81}$那裡說什麼呢.
這裡第一次摩西長大.
所謂成人的時候.
他首次出現在聖經歷史的舞台.
他說後來摩西長大.
他出去到他弟兄那裡.
看他們的重擔.
他看到一個埃及人.
打希伯來人的一個弟兄.
就把埃及人打死了.
藏在沙土裡面.
這個是整本聖經裡面.
最大的摩西民族英雄領袖.
第一次出現在舞台上.
之前沒有怎麼交代過.
摩西為什麼突然間長大成人之後.
他突然間醒覺到自己希伯來人的身份.
為什麼他是一個以色列人.
他不是一個埃及人.
這段聖經很清楚.
說得很清楚.
他知道自己的身份.
不是以色列人.
而是一個希伯人.
到底是他父母教他的.
還是他自己在成長的影響裡面.
以致他能夠理解到.
我不是埃及人而是一個希伯來人.
這裡提到弟兄兩次.
他說他的弟兄.
我們假設他應該是在皇宮裡面出來.
然後他看到他的弟兄.
那些以色列人希伯來人.
他們在做苦工.
然後有一個被人打.
到底是什麼意思呢.
他左看右看見到沒有人.
其中一個看法可能是.
他左右看見沒有人.
看到的時候.

$^{121}$他就去砌下那個.
KO了那個埃及人.
可以是這樣.
有三個人.
第一個是誰.
就是被埃及人打的希伯來人.
第二個呢.
打希伯來人的埃及人.
第三個是誰.
就是他自己.
所以只有三個人.
有一個是被埃及人打的希伯來人.
一個呢.
被埃及人打的希伯來人.
打希伯來人的埃及人.
然後就是他自己.
他身上看起來好像是埃及人.
但其實裡面是希伯來人.
在這樣的環境.
他見到沒有目擊證人.
他就KO了打人的埃及人.
但過兩節之後.
你見到你熟悉聖經的話.
你就知道其實有人見到.
可以是這樣解釋.
意思就是說.
在埃及的工場上.
其實不會是經常開放的.
埃及的工場.
砌磚的工場.
一定有很多工人一起在開工.
所以不可能沒有人在場地附近.
一定是有很多人在身邊.
似乎意思就是說.
沒有人來幫忙的時候.
他就發現不如我KO他.
既然都沒有人幫忙.
其他希伯來人都照樣砌磚.
他們不敢去冒險.
四圍看的時候沒有人肯介入.

$^{161}$他自己去介入.
這是第一次這樣介紹.
成人的摩西.
似乎第二個解法是挺合理的.
這叫做Bystander Effect.
大家可能聽過一些新聞.
其實好幾年前在紐約.
有一個女人在街上被人打.
到處都有人在.
但沒有人幫她.
什麼叫Bystander Effect呢?.
有些科學家告訴我們.
做過很多實驗.
人的心理就是這樣.
人見到有其他旁觀者在.
因為沒有人出手.
所以他自己的結論.
就會覺得這個不是緊急情況.
所以我又不出手.
亦都可能有人認為.
有其他旁觀者在.
假設其他人一定會幫忙.
我就不需要幫忙.
這個叫Bystander Effect.
所以就會變成責任的推卸.
這就是摩西第一次出場的時候.
他的情況.
他見到附近沒有人出手.
有人被欺壓.
有人被攻擊.
他就正義凜然走出來.
去幫助被欺壓的人.
這似乎是那個經民的處境.
他夠膽去冒險.
他夠膽去take risk.
有些人是守株待兔.
等其他人來做.
有些人覺得這裡很黑.
就在那裡就坐.
有些人在黑暗中.

$^{201}$他就去點著眼燈.
就不是在那裡就坐.
做領袖做事奉的人.
就是一種這樣的moment.
是一種這樣的警覺.
就不是作為一個Bystander.
而是見到有需要.
我不會作為一個Bystander.
我覺得這個不是緊急情況.
或者可能有其他人會加入.
這個就是摩西不同的地方.
就是他有那種leadership的drive.
有那種leadership的動力.
等我出手.
繼續看下去.
介紹摩西再講兩個故事.
都是這樣的.
第二天摩西出去.
見到兩個希伯來人.
同胞來的.
他們在互相鬥.
他對那個希夫人說.
為什麼你打同族的人.
又是走出來.
他看不過眼.
那人就說.
誰立你作為我們的首領和審判官.
難道你要殺我嗎.
好像殺你昨天那個埃及人.
然後他就害怕.
摩西就知道這件事.
當然是鬧到法老身上.
法老聽見這件事就想殺摩西.
摩西就躲避法老.
然後逃到米田的地方居住.
第三個故事.
他去到米田.
有一天他在井旁坐下.
米田的祭司七個女.
出來打水.

$^{241}$打滿了.
要喝父親的群洋.
有其他的牧羊人來到.
把他們趕走了.
摩西起來幫他們.
然後去餵群羊.
讓他們可以喝水.
這個路見不停拔刀相助.
這個就是摩西自己本身有的驅動.
這樣去介紹摩西.
用這麼簡單的三個故事.
其實就是告訴我們.
這個人是一個以色列信仰.
聖經信仰的化身.
聖經的信仰.
Biblical Faith.
是應該怎樣的呢.
你看這些人物的傳記.
我以為他只是在說一個人物的經歷.
其實他是想介紹一個人出來.
這個人就是我們追求的目的.
這個就是我們信仰應該有的生命.
他就是這樣介紹摩西.
就是說這個人是embody整個聖經的vision.
整本舊約新約聖經說什麼.
那個人是有freedom.
但是從freedom而來.
就是會有責任.
Responsibility of freedom.
聖經告訴我們.
我們每一個人.
上帝給我們自由的時候.
我們就必然有責任.
有些人很喜歡認命.
或者將聖經的信仰.
化成一種只會認命的信仰.
世界是這麼破爛的.
教會是有問題的.
家庭是會破裂的.
教會是不夠人用的.

$^{281}$福音工作是不夠人力和財力的.
他們就會認命.
或者我們說得好聽一點.
我們就說不如祈禱吧.
祈禱是可以的.
上帝會搞定.
這個不是聖經的信仰.
這個只不過是我們從freedom而來有的.
不想有的責任.
聖經從頭到尾告訴我們.
是God acts through us.
not for us.
我再說一次.
God acts through us.
not for us.
上帝是藉著你去工作.
上帝不是代替你去工作.
上帝是藉著你的事奉.
藉著你的領導.
你去工作.
不是我們一言不發.
然後舉高雙手.
我們投降.
這個不是聖經的信仰.
雖然有很多宗教信仰.
有很多人生觀是教你做這件事.
世界是這麼苦的.
這是苦海.
我們只需要逃避.
逃出來就行了.
你不要感覺到它.
那你就不用苦.
這個不是聖經的信仰.
聖經的信仰告訴我們.
是God acts through us.
最近網上有一些教會的新聞.
不知道你們有沒有留意.
就是中國大陸的教會被打壓.
很清楚.
我們不用分析.

$^{321}$這個政權是打壓教會的.
你愛不愛國是另一回事.
這個政權確實是在打壓教會.
新聞說什麼呢.
海外的華人教會一言不發.
不出聲.
於是有人發起聯署.
說我們需要去指責.
如果你是認命的人.
你會覺得無所謂.
或者我們會解經解到一個地步.
有打壓是好.
小龍珠打壓越大越厲害.
像蟑螂一樣.
越打越多.
但你想想小龍珠的信息不是這樣.
小龍珠不是說繼續鼓勵有打壓.
小龍珠只是告訴你事實是這樣.
雖然是這樣.
但上帝仍然可以工作.
不代表我們要鼓勵有打壓.
或者隨便可以讓打壓的事情發生.
所以我們應該出聲.
應該去到指責.
才去做這件事.
是不是.
不可以一言不發.
如果要聯署的話.
應該去聯署.
應該.
起碼你出了聲.
因為上帝是藉著我們工作.
不是代替我們工作.
我們不要隨便找上帝去做招牌.
然後我們就隻手不動.
我們就不走出來.
摩西不是這樣.
希伯來文很有趣.
變一點點的意思是很有趣的.
希伯來文.

$^{361}$責任這個字.
responsibility.
叫做什麼呢.
叫做akheryo.
akheryo是從哪個字出來的呢.
就是從akher出來的.
所以akher變一變就變成akheryo.
就是責任.
由akher他者這個字出來.
other.
就是說.
其他人的關注.
就變成我的關注.
這個叫做責任.
社會的責任.
他者的需要.
變成我的需要.
other變成我的responsibility.
在聖經的信仰裡面告訴我們.
有兩個他者.
是需要在我們心靈裡面出現的.
第一個.
那個other就是上帝自己.
divine other.
他的關懷要成為我的關懷.
所以上帝所關懷的東西.
就自然成為我的責任.
第一個最大的誡命.
你要盡心盡力的愛主你的神.
主我的神.
他的關懷他的關注.
就成為我的關懷.
這個叫做responsibility.
第二個他者就是愛人如己.
就是我身邊的鄰舍.
身邊的人.
身邊人的需要.
就成為我的責任.
成為我的關懷.
這個叫responsibility.

$^{401}$摩西第一個的leadership moment.
就是一個這樣的moment.
Acher, Aquarius.
其他人有的關懷.
變成我自己有的關懷.
這個人就是民族英雄.
就是leader.
也是每一個基督徒有的生命.
這個是我們追求的ideal.
你聽聽今天上帝跟你說什麼.
Who is the other?.
上帝關懷的是什麼?.
今天成為你的關懷.
其他人的需要.
今天成為你的需要.
教會如果需要有新的leader出來.
Well, 上帝呼召你.
上帝感動你的.
這個變成你的responsibility.
不要做bystander.
做一個leader.
因為上帝是藉著你工作.
不是代替你工作.
我們一起低頭祈禱.
讓上帝藉著我們工作.
再一次將我們的生命交給上帝.
主, 我們的生命是屬於你的.
對我們的呼召.
我們成為你自己的子民.
坐在我們身邊的other people.
其他的人.
就是我們的兄弟姐妹.
主讓他們的關懷成為我們的關懷.
上帝所關心的.
在世界上要工作的.
成為今天我們要的concern.
主, 我們教會進入到一個改選的季度.
主, 我們求你讓教會的關懷成為我們自己的關懷.
讓我們切身的去參與事奉.
因為當我們事奉.

$^{441}$當我們工作的時候.
你就藉著我們去工作.
神求你讓我們盡上我們的責任.
盡心盡力的去愛你.
盡心盡力的去愛我們身邊的人.
愛你的教會.
我們感恩祈禱.
奉主耶穌基督的名字.
所以我今天就講到這裡,謝謝大家!.
\newpage



\section{}
\label{sec:lPh_PhzFDQI}
\textbf{9/2/18 主日講道: 教會藍圖\#2 -- 教會需要問責 (陳崇基牧師)}
\newline
\newline
連結: \href{https://youtube.com/watch?v=lPh-PhzFDQI}{\texttt{https://youtube.com/watch?v=lPh-PhzFDQI}} ~~~~ 語音日期: 2019-07-24
\newline
\newline
\hyperref[sec:me3DczkMU3E]{\small{< < < PREV SERMON < < <}}
~
\hyperref[sec:index_chronic]{\small{[返順時目]}}
~
\hyperref[sec:index_scriptual]{\small{[返順卷目]}}
~
\hyperref[sec:KZEdWxG87Rg]{\small{> > > NEXT SERMON > > >}}
\newline
\newline
$^{1}$這三個星期,今天是第二個星期.
我們講到教會在一個很不利的環境裡面.
特別是現在很多聽到基督徒的領袖.
教會的領袖.
他們在聖上跌倒.
在這樣的環境.
在這樣的世代當中.
我們應該如何去面對呢?.
我們回到聖經裡面一套書.
三本的.
叫做提摩太前書.
提摩太后書.
和提多書.
就在教穆的書信.
Pastoral Levels裡面.
我們回頭看.
到底這些新約的書信.
給我們有什麼啟迪呢?.
上一次我講到這三本書.
其實跟我們寫的時候.
跟我們現在這個年代的環境其實差不多.
因為我們提過.
教穆書信Pastoral Levels.
是主考60年代的時候寫的.
60年代發生什麼事呢?.
就是基督教本身是猶太教其中一個支派.
跟其他猶太教都是一樣.
所以羅馬帝國裡面有法律的保護.
都算是合法.
但60年代.
第一世紀60年代的時候.
就開始大家分道揚鑣了.
猶太教仍然有它在羅馬帝國裡面的法律保護.
不過基督教開始分道揚鑣的時候.
就沒有了這種環境.
就變成了非法.
變成了被人歧視.
是一種對福音很不利的環境.
對教會很不利的環境裡面.
教會是如何在世教裡面立足呢?.

$^{41}$這三本書信經常提到一件事.
就是教會請你搞好你自己.
Put it in order.
然後上帝的道不會有人毀旁.
我們會在教穆書信裡面去看一看不同的經文.
譬如這一段.
提多書二章四節.
提多是一個人.
他是在地中海一個大的海島上.
去到穆會.
他是保羅的一個徒弟之一.
他接到這封信.
這封信二章四節說什麼呢?.
他剛剛說到.
告訴提多這位穆者.
他在教會裡面.
在眾多小組的教會裡面.
他要怎麼處理不同的.
各式各樣的人.
以及不同的年齡層.
這裡剛剛提到少年婦女.
就是提到一些.
I guess young ladies(年輕女士).
他說好指教少年婦女.
她們要愛丈夫.
愛兒女.
謹守.
Self-control(自我控制).
要有貞潔.
要料理家務.
待人要有恩.
馴服自己的丈夫.
免得什麼呢?.
上帝的毀謗.
你可以看到當時的教會.
遇見一個這樣的環境.
有什麼行差踏錯呢?.
新聞就會看到.
人們就會對教會.
嘲笑教會也好.

$^{81}$或者反對教會也好.
或者抨擊教會也好.
所以你看到.
跟我們今天的時代其實也是一樣.
教會裡面出現一些性醜聞.
基督教裡面一些性醜聞.
對教會是非常不利的.
對福音是非常不利的.
人已經對福音有些反感了.
加上教會的見證做得不是很好的時候.
特別有性失德的時候.
有些性醜聞的時候.
更加使教會很難立足.
提多書講的事情.
環境也是一樣.
今天針對的其實也是我們這個環境.
他說免得上帝的道理.
他的福音被人毀謗.
我們可以怎麼講呢?.
上個星期講了一點點.
今天我可以從天主教開始講.
雖然我們跟天主教不是一樣.
不過基本上.
基本上的教義其實都是相關的.
天主教爆出的醜聞越來越大件事.
那個《大紀報》.
上次我提了.
可能你回去有興趣.
有好奇心.
你可能看過一下.
近幾天.
有人說教宗要下台.
為什麼呢?.
因為Francis本身知道.
有一個主教.
或者一些神職人員.
他有性醜聞.
或者他搞過一些小孩.
卻隱瞞真相.
這件事越搞越大.

$^{121}$甚至有很多天主教教徒離開教會.
或者仍然留在教會.
不過他們覺得.
我們不要聽教宗講話.
我們自立門戶.
也有人說.
天主教是時候要改制度了.
神職人員應該可以給他結婚.
是不是?.
閒話家常的人很多時候都會聊這件事的.
為什麼神職人員經常有性失德呢?.
結婚會不會好一點呢?.
你知道天主教的傳統.
差不多有上千年的傳統.
神職人員,主教,神父.
是不容許他們結婚的.
這個大家都知道.
為什麼神職人員要單身呢?.
你知道由第一世紀開始.
聖經是有記載的.
直到教會的歷史中間.
主教的一千年.
那頭一千年.
其實很多神職人員.
都是有自己的家庭的.
有自己的妻子.
有自己的異女.
不如做個聖經混合比賽吧.
彼得有沒有結婚呢?.
有啊.
怎麼看彼得有結婚呢?.
對了.
有外母嘛.
彼得有外母.
有外母你說他結婚了嗎?.
如果沒結婚就是女朋友的媽媽.
巴拿巴呢?.
巴拿巴有啊.
因為保羅說.
你以為我不想像彼得和巴拿巴一樣.

$^{161}$到處去傳道的時候.
有太太在身邊.
但我放棄了權利.
那保羅呢?.
保羅有沒有結婚?.
應該有.
但按文化來說.
circumstantial(紀錄性).
可能會有.
但按傳統來說是沒有的.
anyway.
至少他寫過那麼多前書的時候是沒有的.
所以有神職人員是單身的.
也有神職人員是專業的.
只不過後期.
就是.
主後大概1100年左右.
當時的天主教.
就加上了這件事情上去.
就是覺得.
因為這種是神學的推測.
既然神職人員是代表耶穌.
代表主.
那你說主耶穌有沒有結婚?.
沒有啊,他沒有伴侶.
他是單身的.
還是活得最美的人生.
既然是這樣的話.
那神職人員應該prefer(更好)他單身.
但這個沒有聖經根據的.
這個沒有聖經根據.
所以很多人有不同的議論.
因為這個Catholic Church有這種性失德的問題.
當然.
我們作為福音派Evangelicals.
不可以五十步笑八步.
因為我們近期也有很多這種性醜聞.
這個問題呢.
不在於神職人員是否結婚的問題.
就在於神職人員.

$^{201}$有沒有一種所謂的accountability(負責).
就是有沒有問責.
在教會事奉在他們工作裡面.
有沒有被問責.
上一次我介紹了一本書.
是Stanley Grant寫的.
這本書講到第一本在90年代出版的.
全面來討論處理性失德的問題.
就是牧者性失德的問題.
Stanley Grant在這本書裡面提過.
他說今天我們最大的問題是什麼呢?.
他不只是講教牧.
是講到所有的人.
包括教會裡面的基督徒.
他說最大的問題是什麼呢?.
就是我們很individualistic(個人性質).
很個人.
我們自己一個人.
我又不去團契.
我又不去小組.
我又不會跟人建立關係.
我不會被人問我是否accountable(負責).
就是mind your own business(管自己的生意).
我來教會只是看一場戲而已.
是不是?.
看一場戲你不會跟其他觀眾打交道.
完了就走了.
這就是我們最大的問題.
individualism(個人性質).
也是西方文化裡面最大的問題.
因為我們很喜歡孤獨.
孤獨有什麼好處呢?.
就是因為沒有accountability(負責).
沒有人hold我accountable(負責).
沒有人speak into my life(說進我的生命).
I like it(我喜歡).
Right?.
為什麼呢?.
因為我可以隨心所欲.
這就是最大的問題了.

$^{241}$特別是講到性失德的時候.
很多時候教育.
都是怎麼說呢?.
就是.
人在高處.
就是有點孤獨.
沒有什麼朋友.
是不是?.
他整天去幫人.
就沒有人去幫他.
所以在很多的處境裡面他就會跌倒.
如果你打開教育書信三本.
我發現.
是有四重的問責的關係的.
有四重問責關係.
第一.
保羅.
向他的徒弟那些牧者問責.
提摩太.
提多.
無論他哪個是童工.
他向他們問責.
這三本的教育書信其實就是一種問責的書信.
那.
譬如大家看.
提摩太前書第五章第一節.
這裡又是講到.
提摩太怎麼處理教會裡面更多不同層面的人.
他說不要嚴責老年人.
只要勸他如同父親.
勸少年人如同弟兄.
勸老年的婦女如同母親.
這裡講到父親和母親.
就是給他年長的父親和母親.
講回猶太教的倫理.
因為第五屆.
十屆的第五屆.
就是要孝敬父母.
就是說.
以他為尊敬他.

$^{281}$給他有價值.
來尊重他.
不需要嚴責他.
尊重他.
他說遇見一些比他大的.
比他大的弟兄姊妹.
就當他爸爸或者媽媽.
同年紀的.
或者可能比他年輕的.
勸少年的婦人如姐妹.
他自己家裡的sister.
他說總要清清潔潔.
也要尊敬那些真為寡婦.
這裡不解釋那麼多了.
什麼叫真為寡婦.
就是有很多假為寡婦.
就是扮可憐.
去教會拿錢也好.
貪錢也好.
Anyway.
他特別講到.
勸少年的婦女如姐妹.
總要清清潔潔.
他加了"總要"這個字.
In all purity.
或者absolute purity.
就是說任何的境況當中.
你要處理這些教會的信徒的時候.
那個保羅穆者就說.
他說徒弟.
你不要輕看聖的事談.
我也想不到有什麼.
令福音更不利的醜聞.
是比聖事談.
或者聖事談.
更加大破壞力.
教會可以與很多事情交易.
可能有失誤.
有些事情可能應該說沒說.
不應該說說了.

$^{321}$有些立場可能剔了.
那些固然有它破壞性.
第一個信徒.
特別是教會的領袖.
在聖的上面跌倒.
破壞力是最大的.
是最丟架的.
最多可以被人笑笑.
你不理裡面的事.
不理外面的事.
或者沒有為哪個人發聲.
但是一有聖失德的時候.
丟架是最大的.
我們美國南方壽信會.
我們屬於他的宗派.
Southern Baptists.
有這些指令.
有這些指示.
對教會領袖.
特別是牧者.
就是男女不可以共處於一個密室裡面.
一個男孩一個女孩.
要逃避這種環境.
即使你見一個人.
第二.
在你辦公室裡面見.
你的門一定要有玻璃進去.
就是人家可以裝進去.
你可以看得出去.
第三.
應該有一個轉介的表.
有很多個案不應該由教務處理.
應該轉介給其他人.
無論是輔導員也好.
什麼也好.
可能你覺得這些很老土.
我已經是十一世紀了.
男男女女這些東西.
我們很自由的.
但我想告訴你.

$^{361}$我們開始直堂的時候.
我們要簽一些約.
這就是其中一些約要簽的.
可能你會覺得.
這麼老土.
這些幾十年前才做的.
現在很開放.
你要知道約的精神在哪裡.
約不是故意限制我們某些東西.
或者覺得這個信仰很局促.
要做這些事事奉我們作為信徒.
為什麼這麼局促.
搞出這樣的事.
你要知道.
你要知道有些電影的名字.
叫做《光柴烈火》.
我當然沒有看過.
我沒有看過.
但你知道那些名字叫《光柴烈火》.
不只是電影名字那麼簡單.
其實它在刻畫人性.
人的墮落.
性上的墮落.
我告訴你.
我們每一個人.
無論你是老還是亂也好.
你信了主.
多久也好.
你有多少經驗也好.
我們每一個人.
在一個對的時間.
對的地方.
對的機會裡面.
我們都可以做大衛.
大衛和別士巴的故事.
大家意識到是吧.
但你知道那個故事.
你詳細回去看.
昨晚我再看一次.
我發現.

$^{401}$整個故事裡面是沒有文集的.
只有他一個人.
是吧.
他在皇宮頂走.
一個人走.
他要決定一些事情.
他一個人決定.
他要處理一些問題.
他一個人來處理.
沒有人在身旁來問責.
每個人.
我們每一個人.
都可以做大衛.
只要是對的時間.
對的地方.
對的機會.
我們每個人都可以跌倒.
所以有一些這樣的指令.
其實是幫助我們的.
以至我們不會.
去到.
整天在試探裡面.
在這裡轉來轉去.
總要清清潔潔.
這個很好提醒.
提醒我們.
無論年長或者年輕的.
男男女女的關係.
你不要看輕.
人性的墮落.
上個星期我們提到.
我們都是罪人中的罪魁.
如果保羅也是這麼說的話.
我們沒有一個人能夠垮口.
所以保羅向.
他的徒弟問責.
提摩爾前書第四章第十三節.
我們只是看十六節.
他說.
你要謹慎自己.

$^{441}$跟自己的教訓.
就是說你要做運動.
自我控制.
關心.
小心你自己.
還有你自己的教導.
在這些事上.
教導的事上.
你要恆心.
這樣恆的話.
你能夠救你自己.
最終可以救你自己.
也可以救到聽你道的人.
這是一個問責的關係.
我每個月.
我跟你說過.
我都要見一個.
我們所謂的Sending Church的牧師.
我是很不喜歡見他.
我很不喜歡見他.
因為很遠.
在北面.
你問我喜不喜歡.
我當然不喜歡.
但我想告訴你.
即使我不想要.
也許我需要.
他問我的問題.
主要很多時候不是問這些問題.
他問你個人的生活.
你家庭生活怎麼樣.
即使我是.
一敗塗地也好.
他問完我之後.
我就發現.
原來我有些事要去留心.
這就叫Accountability.
我想問你們生活裡面.
基督徒的靈性生活裡面.
有沒有這些人在你身邊.

$^{481}$可能有些你很不喜歡的人.
他就專門說一些你不喜歡聽的事.
那你可能.
你就當頭棒喝.
可能他.
你最不喜歡聽的事.
就是你最需要聽的事.
Accountability.
保羅對提摩太有Accountability.
提摩太也對他的長質有Accountability.
三章第一節講到長質的.
揀選的時候.
他說人若要得到監督的職份.
這是善無成功.
則或是可信的.
作監督的.
即是作教會領袖的.
長質的必須無可指責.
初入教的不要作監督.
恐怕他自高自大.
就落在魔鬼所受的刑罰裡面.
監督也必須在教外面有好的名聲.
恐怕又是被人毀謗.
落在魔鬼的網絡裡面.
所以提摩太對他的長質也有要求.
無可指責不是完美的意思.
教我們不是完美.
教會領袖我們絕對不是完美.
無可指責是什麼意思.
不是完美的意思.
而是說沒有多少可以公開抨擊的地方.
即是說這個人的操守.
他的生活.
沒有一些可以開放地攻擊他的東西.
這就是對監督的要求.
提摩太就要向他的長質作問責.
我們講過了,接下來9月底開始.
我們會有選舉.
我們教會UV是時候有新的領袖出來.
Accountability.

$^{521}$我們一起事奉的時候就需要向大家問責.
保羅向他的徒弟提摩太問責.
提摩太向長質問責.
長質向會眾問責.
要事實上去教導他們.
事實上去引領他們.
第四重.
你聽到嗎?保羅向提摩太.
提摩太向長質.
長質向會眾.
我們不花那麼多時間看.
但最後會眾向領袖問責.
這是第四重.
就是說那個Congregation.
會對他的領袖來問責.
我們講過.
保羅寫的教目書信.
不只是寫給個人提摩太和提多.
他們看完之後就收在櫃桶.
是給眾人看的.
就是說提摩太提多收到保羅的書信之後.
他就要在教會崇拜里公讀解讀出來.
就是說會眾知道保羅對提摩太有什麼要求.
有什麼accountability(負擔)的items(物品).
會眾是要知道的.
所以每個人都看到.
是transparent(透明)的.
提摩太前書五章十九節.
那裡說什麼呢?.
他說控告長老的情子.
必有兩三個見證就不可收.
非有兩三個見證就不可收.
犯罪的人.
當在眾人面前來襲避他.
使其餘的人可以懼怕.
後面那句就叫做church discipline(教會律師律).
church不是讓我們亂來的地方.
如果有一些明顯的罪.
是需要來到去懲責的話.
就要在眾人面前來襲罵.

$^{561}$這個對我們現代人來說.
很聽不進去.
因為我們很value(價值)什麼.
privacy(私隱).
是不是?.
我私人問題.
my own business(我的業務).
但是他講教會的order(命令)的時候就不是這樣.
因為是一個群體的order(命令).
就不能夠說你隨便點.
你自己生活怎麼樣.
沒有影響的.
是有影響的.
所以是有一個問責discipline(教會律師律).
另外會眾對長老.
對長子有什麼要求呢?.
就都是有問責的.
因為可以寫信去charge(責罰)他.
譬如他生活裡面有什麼不妥當的地方.
但是這裡保羅講的不是那種亂那樣控告.
而是真的有真憑實據的時候.
就需要來到去提醒.
或者去keep(保留)那個長老來到accountable(負責).
我們結束的時候我想問你幾個問題.
是講到accountability(負責).
right?.
我相信如果我們有更多accountability(負責)的時候.
正如教務書信所講.
我們就能夠在一個這樣的環境裡面.
雖然對福音對教會不利.
但仍然可以立足.
因為我們活給世人看.
福音,真理,上帝到底是怎樣的.
如果有更多當面問責.
我相信會減少很多這些性失德的問題.
如果會眾為教會的領袖去祈禱.
知道他也是軟弱者.
只不過他的呼召不夠.
我相信也會減少很多這些問題.
不會全部迎刃而解.

$^{601}$不過起碼會少很多這些事情.
神的道不會那麼容易被毀謗.
如果我們每個人都有一個accountability(負責)的partner.
或者group(群體)的時候.
我們跌倒的機會自然會少.
我在想我自己人生裡面.
有好幾個弟兄.
是我做得不妥當的時候.
或者是有一點不妥當.
他們願意撕破臉跟我說.
不怕傷害我們友情.
他說Mark你這樣不行.
請你這樣做好嗎.
是不是難聽?當然是難聽.
那種confrontation(爭吵)是不是難過?.
當然是難過.
但問責的精神就在這裡.
你作為一個基督徒.
在你人生當中有沒有這些人.
可能是你最恨的人.
那些就是你最需要的人.
可能有些話你最怕聽.
你巴不得避開他.
但我想告訴你那些話可能是你最需要聽的.
如果你想起一個這樣的人.
或者多過一個這樣的人.
為他感恩.
因為有這樣的人hold you accountable(負責).
為你祈禱去守望你.
那你不會孤軍作戰.
因為你孤軍作戰是死定的.
不久之前我看了這樣的圖.
他說到每一個基督徒.
特別是教會領袖.
他需要在所謂的關懷圈.
就是一個question(問題).
一個leader(領袖)的circles of concern(團體).
就是我們應該要put it into order(安排)的時候.
我們應該有什麼次序.
第一是個人.

$^{641}$個人清清潔潔.
規規矩矩.
按照神的心意.
不是按照我肉體的心意.
第二.
我家庭的關係.
無論是夫妻也好 兒女也好.
每一方面.
put it back to order(安排好).
有沒有一些我漏了的東西.
第三 教會.
和弟兄姊妹的關係.
可以keep我accountable(負責).
spiritually accountable(精神上負責)的.
我沒有關懷他們.
沒有這個圈.
另外就是教愛.
我對教愛 我的明星是怎樣.
我想請你.
最後我想請你留意.
主要的目的.
主要目的剛才所講的.
這兩個禮拜 三個禮拜所講的.
不是說要限制我們什麼.
不是要規限我們的自由.
而是最主要目的.
是要想上帝的名不會被毀謗.
讓福音 讓上帝的道.
能夠表揚出來.
that's the whole point(這就是重點).
是missionary(任務).
不是personal(個人).
不是針對我 針對他.
為什麼要這樣 鉗制我人生.
最主要的目的不是這樣.
是missionary(任務).
是我們作為一群基督徒的會眾.
我們怎樣去見證神.
讓上帝的福音.
能夠表揚出愛慕.

$^{681}$神的名不會被毀謗.
我們結束的時候.
在上帝的恩典裡面去成長.
我們會守聖餐.
聖餐告訴我們.
作為罪人中的罪魁.
我們站在十字架底下.
由恩典開始.
也由恩典裡面去長進.
讓我們來到神的面前.
我們一起去反省.
上主我們.
來到禰面前我們認罪.
因為我們每一個都是千瘡百孔.
作為人中的罪魁.
我們有很多虧欠的地方.
在我們的心思意念.
我們的行為.
我們的操守上.
我們有很多不是無可指責的.
主要是提醒我們.
讓我們謹守.
總之在我們這個世代裡面.
對福音 對你的名.
有很多的諉謗 很多的抨擊.
主要是讓我們站起來.
在這個世代當中.
做你的見證.
我們身邊的弟兄姊妹.
教會裡面的領袖.
教會裡面的導師.
弟兄姊妹.
他們遠以來都可以提醒我們.
多謝禰賜這些生命給我們.
讓我們走我們的這幾路人生的時候.
不至於跌倒.
上主禰紀念我們.
讓我們在你的恩典當中成長.
我們拿起這個餅 這個杯.
知道你對我們的恩典是多麼深厚.

$^{721}$你去買熟我們.
去更新我們.
以至我們的生活.
我們的生命.
我們的品格.
能夠見證你.
祈禱奉主明禱.
\newpage



\section{}
\label{sec:KZEdWxG87Rg}
\textbf{9/9/18 主日講道: 教會藍圖\#3 -- 教會領袖的素質 (陳崇基牧師)}
\newline
\newline
連結: \href{https://youtube.com/watch?v=KZEdWxG87Rg}{\texttt{https://youtube.com/watch?v=KZEdWxG87Rg}} ~~~~ 語音日期: 2019-07-24
\newline
\newline
\hyperref[sec:lPh_PhzFDQI]{\small{< < < PREV SERMON < < <}}
~
\hyperref[sec:index_chronic]{\small{[返順時目]}}
~
\hyperref[sec:index_scriptual]{\small{[返順卷目]}}
~
\hyperref[sec:tDpoW3ZdL20]{\small{> > > NEXT SERMON > > >}}
\newline
\newline
$^{1}$今天我們繼續我們三個星期的一個 topic(話題).
就是叫做教會藍圖.
教會藍圖講什麼呢?.
教會藍圖是講.
我不知道為什麼有 Abraham(阿博)出來.
這是英文課我講到的講道.
在 folder(文件夾)裡面有其他的 slide(slide).
Anyway 先關掉它.
我可以沒有 powerpoint(電子顯示)也講得到.
我們講教會藍圖是做什麼呢?.
就是我們看看新約的文件.
新約的聖經裡面.
告訴我們到底.
教會在一個不利的環境裡面.
要怎麼去生存.
怎麼去立足.
特別是教會做得不是很好.
很多時候教會有醜聞.
教會有一些失見症的地方.
我們跌倒的地方.
其實這些是很不奇怪的事情.
只不過很多時候新聞的媒體.
他們把它放大而已.
教會也是一群聚人走在一起的地方.
不是因為我們碰到所以走在一起.
而是因為我們碰不到.
我們需要神給我們的恩典和赦免.
我們要學做一個像基督的人.
所以很多時候跌跌碰碰我們會跌倒.
在不利的環境裡面.
怎麼也好.
怎麼可以在教會裡面立足呢?.
新約有一堆書信.
叫做教穆書信.
有三卷書.
是提摩太前書.
提摩太後書.
以及提多書.
這三卷書信是講什麼的呢?.
這三卷書信就是.

$^{41}$在第一世紀.
新約聖經的時代.
教會的時代.
有一個牧者叫做保羅.
保羅是一個很有經驗的牧者.
他有兩個徒弟.
至少兩個徒弟.
一個徒弟叫做提摩太.
一個徒弟叫做提多.
一個叫做 Timothy.
一個叫做 Titus.
他在各自不同建立了一個教會的城市裡面.
在牧養教會.
這三卷書信.
提摩太前書和後書.
以及提多書.
就是到師父保羅.
去教他的徒弟.
怎麼去牧養教會.
特別是在一個不利的環境裡面.
怎麼去牧養教會.
雖然以前歷史裡面發生的事情.
跟我們今天發生的事情有不相似的地方.
但是其實有很多都盡相同的.
當時他們對於信仰.
是很不利的社會.
同樣今天也是一樣.
有不同的原因.
但是大體上其實都是一樣的.
那麼.
怎麼去重建神自己的名.
上帝自己的聲譽呢?.
特別是對我們的教會.
城會教會.
其實也很切身.
因為我們這個月9月底.
我們就會開始進行改選.
我們CORE TEAM的會員.
有部分會step down.
他們會step down之後.

$^{81}$然後有一群新的領袖.
我們希望能夠在選舉裡面產生.
我們會有三部分會做.
我們每一個聚會的弟兄姊妹.
要重新登記做選民.
然後會有提名.
是由弟兄姊妹會來提名.
一些他們認為合資格.
希望代表他們來帶領教會.
最後我們就會有選舉.
我們很民主的.
所以我們希望能夠.
產生新一批的領袖來帶領教會.
希望這個新的CORE TEAM.
能夠在未來三年裡面.
三年一次我們有選舉.
去帶領我們教會.
聖經裡面也有說很多教會領袖.
所以今天我們說CORE TEAM的資格.
提出這些資格的時候.
教務書信.
為什麼這麼多說這件事.
因為你要去到一個不利福音的環境裡面.
重新振作到立足的話.
你就需要去到對領袖有要求.
因為領袖就是教會的代表.
所以教務書信經常提到.
那些領袖或者所謂監督.
有什麼資格.
監督有什麼資格.
有什麼要求.
等一下我們會看一段這樣的經文.
但是我想告訴你.
教務書信不是寫給牧者這麼簡單.
我們兩個星期前提過了.
是寫給眾人.
整個教會裡面的會眾.
有一段話想你留意.
你聽清楚.
那句話說.

$^{121}$Leadership is an extension of discipleship.
就是領袖是作基督門徒的延伸.
領導是基督徒.
做基督徒門徒的延伸.
就是說你已經成為一個基督徒了.
神父每一個願意去毀身.
去爬這個基督徒.
我們都需要去增長.
我們的領導的能力和質素.
都需要去深化我們的心靈.
以致我們能夠成為一個更好的領袖.
This is for everybody.
不只是某一些人.
牧者傳道人他們做好一點就可以了.
原來每一個基督的門徒.
都需要變成一個成熟的領袖.
如果有聖經的話.
跟我打開去提摩太前書第三章第一節.
這裡說人想要得監督的職責.
就是善無聖功.
這話是可信的.
聖經說如果人想做一個教會領袖.
能夠得到監督的職責.
是可信的.
就是說可以佩服的.
所以聖經很欣賞一個人.
或者應該鼓勵每一個人.
每一個基督徒.
能夠得到監督的職責.
一個好的意願.
是一件很好的事來的.
想要得到監督的職責.
是一件很好的事.
這裡說到想得到監督的職責.
那個動詞.
有一個小朋友.
很想在曲奇盒裡面拿曲奇.
大家有小朋友知道的.
或者觀察過小朋友.
想要糖果,零食的時候.

$^{161}$他很想伸手就拿.
人若想得到監督的職責.
是善無聖功.
這句話是值得去信服的.
值得我們去確認的一句話.
是一件美事.
可以事奉上帝.
為什麼不是美事.
不是為了自己舒服這麼簡單.
是為了教會.
為了神國的事功.
是一件美事.
做監督有什麼要求呢.
三件事.
第一.
成熟的生命.
第二.
完整的關係.
第三.
實際的技能.
記不記得.
成熟的生命.
成熟的性格.
完整的關係.
證明的技能.
三件事.
這三方面在教務書信裡面提到監督的質素的時候.
一定有提到.
三章第二節.
之前說什麼.
人若想得到監督的職責.
善無聖功.
這句話是可信的.
三章第二節.
作監督的.
必要無可指責.
只作一個婦人的丈夫.
有節制.
自首.
端正.

$^{201}$樂意接待軟人.
善於教導.
不因酒之事.
不打人.
只要溫和.
不憎勁.
不貪財.
我想問你們剛才那三方面.
生命的成熟.
然後完整的關係.
還有一個證明的技能.
一個實際可以應用到的技能.
三方面.
等下我們看這些資格的時候.
我想問你們是投入哪一個組呢.
哪三個.
是生命.
是技能.
技能還是關係.
必須無可指責.
這裡絕對不是代表完美.
Blameless的意思.
Above reproach.
就是說她生命裡面沒有一些東西是open to attack.
沒有一些東西去到.
別人懷疑或者會攻擊的地方.
這裡說到生命.
無可指責.
指了一個婦人的丈夫.
就是說他最基本的.
家.
他的婚姻一定要完整.
要美滿.
這裡當然不是說.
結了婚才可以做監督.
不是這個意思.
也不是說只有男人才可以做監督.
女人就不可以做監督.
因為有當時的社會背景.
我們以前的華人.

$^{241}$傳統很久很久之前.
很多時候女人都不是很識字.
因為她們沒機會讀書這麼簡單.
因為.
通常男人才識字.
以前的以色列社會都是一樣.
羅馬帝國的社會也是一樣.
男人有education.
所以他們很多時候做leaders.
這是出自那時候的背景.
我們不需要去到.
照辦主管.
聖經不是要我們照辦主管.
回到以前古代.
而是讓我們當時的文化.
看到一些大原則.
現在可以應用到.
最重要的是什麼呢.
他要作一個婦人的丈夫.
就是說他婚姻一定需要完整.
剛才那三個東西.
就是關係.
Holistic relationship.
你不可以說我不理我家人.
我只要有技能就行了.
No.
因為上帝看你.
是看了你的關係.
他說要有節制.
Temporary.
就是說有self control.
要自守.
這裡說到emotional的東西.
就是情緒上我們需要去管理.
跟一些大情大性的人一起.
成功是很難的.
很難的.
動不動就發脾氣.
動不動就罵情緒.
這樣就搞不定.

$^{281}$教會需要一起向前來事奉主.
同宮就需要有節制.
特別說到emotion.
就是情緒上.
這個就是說到自己實在的生命.
端正.
什麼意思呢.
端正.
Respectable.
就是說他走出來的時候.
他的生命.
他的生活整體上.
都是受別人尊敬的.
這個說到關係.
樂意接待遠人.
很多時候初期教會.
有很多傳道人不是駐守在一個堂會.
因為人手不夠.
所以要四處去.
ABCD四處去廣道.
去到不同城市的時候.
就需要長老.
長輩監督.
去接待他們.
就是樂意開放自己的家.
這個也是說到他對外的關係.
善於教導終於有技能.
不只是教導這麼簡單.
是善於教導.
Able to teach.
他自己要愛上帝的話.
他要追求上帝的話.
他要有很好的根基.
並且不是自樂這麼簡單.
是要教其他人.
Able to teach.
這個就是技能.
不因酒之事.
就是不要醉酒.
OK.

$^{321}$不要醉酒.
其實他主要說什麼呢.
當時當然現在有他的社會背景.
當時的水大部分都是酒.
去歐洲玩.
酒比水便宜.
當時一定有一點酒精.
水沒有.
清清的水.
很少.
但很多時候有一些所謂substance abuse.
就是事物濫用.
這裡又說回我們的生命.
一個監督應該有的生命.
不打人.
我覺得這句和諧本也很有趣.
他不可以打人.
其實新的中文譯本怎麼譯呢.
他不喜歡鬥氣.
這個人不喜歡經常和別人鬥氣.
鬥誰贏誰輸.
開會的時候.
到底誰的意見好一點.
誰贏誰輸.
他說只要溫和.
不堅持永遠自己是對的.
溫和.
溫和不是lying.
所以為什麼要解經呢.
gentle.
其實就是說.
希利文有一個意思就是.
You are not right all the time.
I am not right all the time.
就是溫和的意思.
其他人補足我.
不爭競.
peaceable.
就是有親和力.
你明白我的意思嗎.

$^{361}$親和力是什麼意思.
教會永遠都是很多不同的人背景走在一起.
親和力就是說.
大家可以相親可以相處.
不只是一舊的人.
我才能接納他.
其他人我就排擠他.
做監督或者我們應該有的目標就是應該要這樣.
其實這裡說什麼.
不打人要溫和.
不要堅持自己永遠是對的.
不要爭競其實是什麼.
不要濫用權力.
不要abuse或者misuse.
我們的power.
講到關係.
不貪財.
也是講到關係.
為什麼呢.
這裡不是說知足自己夠吃夠住.
你知道古代的社會.
古代社會.
可能中間比較年長的可能也有一點認同.
古代的社會.
經濟是怎樣的呢.
就是大家這麼多人.
我們分一個蛋糕.
在我們中間1\%的人.
那些top 1\%富有的人.
他就分了蛋糕的99\%.
他已經拿走了.
剩下的那一塊.
就由誰分呢.
99\%的人去分.
所以很多窮人.
所以古代社會裡面.
大家家常說要賺錢.
expand.
business要擴闊一點.
其實就是貪的意思.

$^{401}$不貪財是什麼意思.
不只是知足這麼簡單.
而是幫助其他人.
幫助教會.
慷慨generosity.
肯去幫其他人.
如果你是那位比較富有的.
資源多的.
就自然應該拿出來.
如果你不拿出來.
其實就是貪財.
就是守財.
肯去慷慨generous.
所以說了這麼多.
暫時說這麼多.
有多少是技能呢.
我數來數去都只有一樣.
大部分都是和我們內在生命.
和我們的關係有關.
特別是和權力有關係.
就是不去濫用權力.
或者不去堅持自己有的權位.
我不知道你有沒有聽過這樣的情況.
就是大家一齊走過來.
商議的時候.
開會的時候.
解決不了問題.
其中一個的同工會發火.
或者臉黑.
其實就是emotional manipulation.
在那種環境裡.
大家不要傷和氣.
算了給他吧.
或者那個人會說.
如果你不同意的話.
我今天立刻走.
或者我立刻辭職.
這是情緒的威脅.
這個是abuse of power.
有牧者會是這樣.

$^{441}$有長老是這樣.
有團體的職員會是這樣.
這些不是同工.
這些是自己的ego太大.
自己的ego太大.
如果一個大ego的人.
自己很sit on power.
抓住權的話.
身邊的人.
他就不想身邊的人來看見.
或者他會斷絕.
他想人做什麼呢.
是奉承他.
喜歡去稱讚他.
心胸狹窄.
我們經常說方丈很小器.
就是類似這樣的意思.
如果你有big ego的話.
這種人是很難跟他同工.
不僅如此.
他的生命裡也可能有很多未解決的問題.
心裡面有些big ego的人有什麼問題.
你聽清楚.
用英文說.
就是有一種need to be needed.
有一種被需要的需要.
想爆你.
need to be needed.
是很微妙的.
在一些關係群組裡面.
一個教會也是一樣.
好像一個組織.
如果有一群人跟著我.
有一組人跟著我.
我就覺得人很需要我.
所以很微妙的.
我望鏡的時候.
我多厲害.
有一群人跟著我.
需要我.

$^{481}$自我良好的感覺.
所以在教牧的書中說到.
領袖.
他的qualification.
一定是生命和關係走在前面.
他不會去爭勁.
不會去爭power.
謙卑的來服侍.
生命先成熟.
關係先走.
我希望我講這些.
不僅是應用在教會裡面.
想一想監督.
leaders.
想一想工場裡面.
職場裡面.
是不是其實也有很多時候.
用這些原則.
是走得通的呢.
其實21世紀裡面.
有很多講領導學的研究.
都是用聖經的原則.
所謂servant leadership.
大行其道.
經常說servant leadership.
或者character leadership.
virtual leadership.
這些字眼從哪裡出來的呢.
是聖經裡面的倫理.
異曲同工.
好幾年前.
有一個研究商業的專家.
叫做Jim Collins.
Jim Collins寫了一本書叫做Good to Great.
可能如果你有認識的話.
你應該有看過這本書.
他花了20年的時間.
研究一些很久歷史的公司.
他說有什麼分別呢.
一家好的公司.

$^{521}$和一家頂級一流的公司.
Good to Great.
到底有什麼分別.
裡面有很多很多值得我們留意的東西.
我只講一件事.
遲些有機會我們再講多一點.
他在這本書裡面說.
Leadership分五個層面.
OK?五個層面.
最高層面.
Level 5的Leaders.
有什麼特別的地方呢.
我引用他其中一句話.
Level 5 Leaders是做什麼的呢.
Level 5 Leaders.
五星級的領袖.
是將自己的需求.
從自己轉移到.
去建立群體.
建立一家偉大公司.
一個更大的目標.
不代表五星級的Leaders.
沒有自我或者沒有自身的利益.
事實上他們是有的.
是雄心勃勃的.
Ambitious.
但是他們的抱負放在哪裡呢.
不是自己.
而是放在群體其他人身上.
就是他願意放下自己的社企.
去建立一個更加偉大的Institution.
或者一個Group.
這個叫Level 5 Leaders.
其實保羅的書信.
教務書信.
其實都是在講這個.
我們都要有一個Goal.
每一個Commit的Question.
去到這個Level 5的.
一個這樣的生命.

$^{561}$所以這些的Qualifications.
全部都是去鼓勵我們.
或者給我們一個遠一點的指標.
要你和我都可以達到.
三章第四節.
他說你要好好管理自己的家.
使兒女凡事端莊順服.
人若不知道管理自己的家.
焉能照管上帝的教會.
初入教的不可以作監督.
恐怕他自高自大.
就落在魔鬼所受的刑罰裡面.
監督亦都必須在教外有好的名聲.
恐怕被人毀旁.
落在魔鬼的網羅裡面.
問你.
這裡這麼多句.
有哪幾句是講生命.
有哪幾句是講關係.
或者講技能.
好好照管自己的家.
使兒女凡事端莊順服.
家很重要.
因為當時.
特別是當時的教會.
第一世紀的教會.
是家庭教會.
就是說你要有教會聚會.
就要進入你家.
家裡看到什麼.
就等於你教會是怎樣的.
所以上上下下都應該要去敬畏上帝.
因為人先需要去管理自己很好的家.
然後才可以延伸到上帝的家.
教會.
是不是.
這裡又講到完整的關係.
生命呢.
初入教的不可以作監督.
恐怕他自高自大.

$^{601}$落在魔鬼所守的刑罰.
魔鬼因為自高自大.
然後墮落.
如果初入教還沒經過實際考驗根基.
還沒夠深的.
做了領導.
很多時候就撞板.
講到內在生命.
然後教外要有好的明星.
免得被人毀謗.
外在關係.
講了這麼多.
我真的找不到一個.
講很多技能.
的描述.
我去過很多領導的會議.
很多時候都是講技能.
我不是說技能不重要.
技能當然重要.
是不是.
例如.
舉個例子.
唱歌才算是順位.
你不會彈就不會彈.
是不是.
那些半途出家的已經學不到.
很聰明的就學到.
但這些十居其九都不行.
你是五音不全就五音不全.
可能一輩子都五音不全.
當然有技能.
但如果你站在台上.
你沒有一個生命.
你對主的毀身是不夠的.
空殼.
因為沒有生命的關係.
過去.
星期四星期五.
我很忙.
因為總會有會議要開.

$^{641}$因為UV都兩年了.
我們可能都成熟.
可以變成其他的直行牧者.
開新教會.
Church planter.
所謂sending church.
我們可以差派人出去.
有些手續要做.
就說不如你來吧.
你學學怎麼做教練.
怎麼做好coach.
好吧.
開了兩天會議.
其實就是一個assessment.
來自美國有很多不同的church planters.
我們就走在一起.
剛剛進入教會模樣.
教會建立的工作.
這些church planters.
我們走在一起.
做一個assessment.
我們就有份做一點.
很少部分.
有很多的.
我都很欣賞.
我們教會和總會.
浸會.
有的系統.
那些人apply.
申請.
想做直行.
他有一個很full.
很仔細的assessment給他.
等於你的公司.
你剛做新人.
有一個所謂pre-qualification.
看清楚你到底對不對.
有哪方面需要有improvement.
我聽到不少故事.
我聽到一個我認識的牧者.

$^{681}$他就說他幾年前.
他在芝加哥開了一間新的教會.
有一班call team的同工和他走在一起.
開了.
幾十人之後.
結果同工.
大家就吵架.
無法在一起.
他沒有親和力.
不知道怎樣和人相處.
social skill.
他自己有很多的issues沒有處理.
那些need to be needed.
結果.
崩潰了.
婚姻就成為很大危機.
心理上.
他就很不舒服.
婚姻就成為很大危機.
心理也成為很大危機.
兩年後這東西就沒有了.
誰知道.
幸好有一個牧者的朋友就把他抓回來.
就和他說一些很深入的東西.
那些need to be needed.
那些.
你知道事奉.
是可以揭露一個人.
最美的.
一方面.
同時.
可以打開最脆弱的人.
It's a trial.
It's a test.
問題是你夠不夠膽.
go through this test.
結果這位.
朋友.
這位弟兄.
他真的在上帝面前重新再整理過.

$^{721}$認罪再悔改.
他現在再做回牧者.
不是他自己開的教會.
是幫人家.
幫人家.
起到模樣.
你留意這段經文.
我們看了這麼久.
魔鬼出現了兩次.
你在事奉.
你在做很多事情.
其實魔鬼也在做.
魔鬼可能.
比你做的更多.
教務書信裡面說的.
這些qualification.
是生命與慣性.
比你所做的事情重要.
Being.
和Doing.
上帝所關懷的.
是我們的Being.
不是我們的Doing.
聖經強調的是你個人生命.
你的品格,你的關係.
不是舞台上的表現.
不是一個話劇.
這個不是run script的地方.
教會是上帝的家.
神的心意不是我們做的事情.
給他看.
神的心意是我們生命.
可以成長.
可以更加完備.
看完教會的.
頂袖的所謂履歷.
監督的.
qualification,我不知道你有什麼感覺.
我的感覺就是.
What a list.

$^{761}$我以往都是看.
這些經文.
水過鴨背.
當我想清楚.
每個字,每句說話.
說什麼的時候.
我才知道上帝.
他緊張的是什麼.
緊張你的生命.
緊張你的Doing.
緊張我的Being.
三方面,我們希望有.
在事奉當中.
我們有mature character.
有holistic relationship.
有proven skills.
正如我開頭說.
我們9月會開始.
改選,這個實在是很興奮的事.
我覺得有新一代的.
領袖.
新一代的領袖出來.
但同時也是很嚴肅.
因為我們知道我們不是單單.
做事,而是生命要經歷.
要經歷歷年.
我們需要去成長.
上帝可能在我們事奉當中.
會揭露一些.
我們需要面對的事.
我們要去.
重整的事,無論是我們的關係.
也好,我們的生命也好.
在我們的技能當中也好.
是興奮.
是非常嚴肅.
我希望你每一個.
當我們進入9月尾的.
改選的時候.
我們為自己祈禱.

$^{801}$我們在神的.
話語當中我們去反省.
讓上帝在我們.
生命當中的工作.
讓我們更加像.
耶穌基督.
我們同心一流祈禱.
上主我們再次多謝你.
因為你緊張我們的生命.
多於一切.
主你要我們更加像你的.
獨生子耶穌基督.
在我們的品格,我們的心態.
我們內在的生命.
我們各種的關係.
主你緊張.
讓我們帶領一個聚會.
帶領一個茶經.
做很多的program.
唱詩,音樂.
話劇等等,所有東西.
這些都是次要.
我要在你面前.
被你訴訟.
成為一個更加.
討你喜悅的人.
主你多謝你對我的愛.
你的恩典.
求你繼續幫助我們教會成長.
我們快要改選的時候.
主,我們求你.
引導整個過程.
幫助我們不單止選.
某些的弟兄姊妹出來.
帶領教會.
我們自己被提醒.
我們要成為一個更加完備的生命.
討你喜悅.
祈禱奉祝耶穌的名字.
\newpage



\section{提摩太前書 2:8-15}
\label{sec:tDpoW3ZdL20}
\textbf{12/2/18 主日講道: 女子組加油! \#4 提摩太前書2\_8-15  (陳崇基牧師)}
\newline
\newline
連結: \href{https://youtube.com/watch?v=tDpoW3ZdL20}{\texttt{https://youtube.com/watch?v=tDpoW3ZdL20}} ~~~~ 語音日期: 2019-07-25
\newline
\newline
\hyperref[sec:KZEdWxG87Rg]{\small{< < < PREV SERMON < < <}}
~
\hyperref[sec:index_chronic]{\small{[返順時目]}}
~
\hyperref[sec:index_scriptual]{\small{[返順卷目]}}
~
\hyperref[sec:zzrp0tNV2iI]{\small{> > > NEXT SERMON > > >}}
\newline
\newline
提摩太前書 2:8-15
\newline
\begin{longtable}{cl}
\hline
\hline
章節 & 經文 (和合本修訂版)\\
\hline
2:8 & \begin{tabularx}{0.7\textwidth}{X} 我希望男人舉起聖潔的手隨處禱告，不發怒，不爭論。 \end{tabularx} \\ \\ \relax
2:9 & \begin{tabularx}{0.7\textwidth}{X} 我也希望女人以端正、克制和合乎體統的服裝打扮自己，不以編髮、金飾、珍珠和名貴衣裳來打扮。 \end{tabularx} \\ \\ \relax
2:10 & \begin{tabularx}{0.7\textwidth}{X} 要有善行，這才與自稱為敬畏神的女人相稱。 \end{tabularx} \\ \\ \relax
2:11 & \begin{tabularx}{0.7\textwidth}{X} 女人要事事順服地安靜學習。 \end{tabularx} \\ \\ \relax
2:12 & \begin{tabularx}{0.7\textwidth}{X} 我不許女人教導，也不許她管轄男人，只要安靜。 \end{tabularx} \\ \\ \relax
2:13 & \begin{tabularx}{0.7\textwidth}{X} 因為亞當先被造，然後才是夏娃； \end{tabularx} \\ \\ \relax
2:14 & \begin{tabularx}{0.7\textwidth}{X} 亞當並沒有受騙，而是女人受騙，陷在過犯裡。 \end{tabularx} \\ \\ \relax
2:15 & \begin{tabularx}{0.7\textwidth}{X} 然而，女人若持守信心、愛心，又聖潔克制，就必藉著生產而得救。 \end{tabularx} \\ \\
[1ex]
\hline
\hline
\end{longtable}
$^{1}$今天我們繼續講女子做家友的系列.
為什麼呢?.
講到現在.
因為下個星期就會結束了.
下個星期無論如何都要講一下瑪麗亞.
因為聖誕節到了.
無論如何都要講一下Mary.
Mary當然是女人.
所以我們用女子做家友.
Mary作為一個結論.
如果你今天第一次來.
也不要緊的.
我們會review一下.
我們每一次頭三分鐘.
review一下我們到底講過什麼.
到底過去三幾個星期我們講過什麼呢?.
我們說過我們知道聖經拿在手上的時候.
它是一個以前的社會.
父權社會的文學作品.
當時的社會是重男輕女.
或者可能今天的社會也是有一點點.
但是當時的聖經的社會.
確實是一個父權的patriarchal society.
所以裡面有一些記載.
一些描述.
都是在溝通當時的一些文化價值.
這是沒錯的.
另外.
聖經是挑戰父權主義的.
挑戰重男輕女的一種思想.
所以我們找到很多地方.
其實聖經雖然有時候是假設.
但是也會去挑戰或者發出疑問.
到底這些重男輕女的思想.
是不是不對的呢?.
第三我們講到在聖經歷史裡面.
女性扮演很多重要的角色.
這個大家有留意.
不只是男人做事.
上帝也用很多女人,女士去做事.

$^{41}$第四我們講到解經或者看聖經.
要了解明白的時候.
我們需要去明白背景.
盡量去明白背景.
上一次或者前一次.
我們講到一段聖經.
是講到女人要蒙頭的.
你不清楚聖經的內容.
不要緊.
這裡講到有一段聖經.
講到女人敬拜的時候.
講到的時候需要蒙頭的.
上次我們略為有些關鍵字去解釋之後.
我們發現可能女人要蒙頭.
是過不多大男人的壞鬼神學.
不是保羅自己的神學.
保羅寫的那段聖經.
很可能是去矯正他們的一些錯誤.
我們講過了.
我們講到聖經女人的角色.
有一些很困難的經文一定要處理.
那些是不能避開的.
上一次我們講到.
《江南多前書特首》一章.
那個不容易解決.
今天我們也要講一段很難解的經文.
我還沒看這段經文之前.
我想告訴你.
今年年初.
有一個很出名的牧師.
寫了一篇文.
放上網了.
他說女人最好留在家裡.
煮飯,洗衣服,照顧孩子.
這是聖經的教導.
他們不要在教會裡面做很多領導的角色.
如果在教會裡面你要事奉.
你要做事的話.
女性最適合做什麼呢?.
就是去處理部門.

$^{81}$或者去托兒所.
去做其他的工作.
不要去做一些公開的教導.
因為他說這是聖經所教的.
一出這篇文.
就引起很多很多的反應.
Twitter裡面有很多很多人表達他的意見.
其中一個女作家.
她說.
我不去爭辯這些問題了.
不過我想歡迎所有的男士.
你在Twitter裡面回答.
告訴我.
你生命裡面有沒有一些女性.
或者女牧師.
女教師.
女屬靈的長者.
是幫過你,影響過你的.
一出Twitter.
就有很多人.
特別是男人.
他就放了很多的意見上去.
他說我人生裡面有很多很多女性.
是祝福過我的.
我想告訴你.
這篇文.
就是叫女生.
不要做這麼多公開的事情.
不要領導.
他的聖經的根據是來自哪裡呢?.
就是今天看這段聖經.
提摩太前書第二章.
第八節開始.
你知道聖經提摩太前書是什麼嗎?.
我先解釋一下.
就是保羅當時是一位牧師.
他寫信給在二佛所打理教會的一個徒弟.
他的名字叫做Timothy提摩太.
在書信裡面很多時候.
就是告訴他怎樣迎接這些挑戰.

$^{121}$怎樣處理教會裡面的問題.
這個叫做提摩太前書.
就是第一封書信.
第二封書信是提摩太後書.
就講到他如何去教會.
提摩太如何去處理教會裡面的問題.
這裡第二章第八節.
他說什麼呢?.
他說所以我願意.
就是他叫提摩太你在二佛所裡面.
教那些男人做什麼呢?.
在每一個場合裡面.
都舉起聖潔的手來禱告.
不要發牢.
不要爭論.
當然這裡是男人有問題.
所以他就說.
不如教那些男人.
不要經常爭論.
不要經常爭什麼.
舉起聖潔的手.
不要發牢.
並且不要爭論.
第九節他就講到女人.
他說同樣我願意女人.
要以端莊的衣著.
樸實和自律來裝飾自己.
不以編髮,金飾,珍珠,名貴的衣裳.
而以善行裝飾.
而且才語自稱為敬畏神的女人相稱.
女人要安靜學習.
完全信服.
我不容許女人教導人.
管轄男人.
女人都要安靜.
因為是亞當先被造.
然後才是亞娃.
而且不是亞當受騙.
卻是女人受騙.
陷在過犯裡.

$^{161}$女人若以自律,慈守.
著信心,愛心,聖潔.
就必藉著生產而得救.
這段經文.
幾乎每個字都有困難.
特別這段.
什麼以生產得救.
生產不是工廠裡的.
貼貼東西.
做線的生產.
生產是生孩子的意思.
如果女人藉著生孩子得救.
很奇怪.
沒理由生孩子就得救.
我們是信耶穌得救.
但你看看.
不只是這一節有問題.
幾乎每一個字都有困難.
而且這段聖經很奇怪.
就是保羅用很多字.
只是在這裡用.
保羅的字典裡.
這段經文裡用了很多字.
只出現過一次,兩次.
譬如管轄男人.
男人的字.
整本聖經裡只用了一次.
就是這裡.
到底怎樣了解這段聖經呢.
剛才我說過.
那篇文章.
那位很出名的牧師.
他當然是照字面來看.
他說不許女人講道.
當然不讓女人講道.
或者做教導,領導.
不讓她管轄男人.
當然我們是男人.
我們無所謂.
沒人管我們是很好的.

$^{201}$你看看結婚的老婆.
出差她很自由.
昨晚我女兒.
我差點想說我女人.
我太太.
我太太出去跟姐妹吃飯.
子女很開心.
為什麼.
我們晚餐吃什麼.
我們出去買東西吃.
媽媽在就不會.
媽媽會煮健康的東西給她.
我們就去買雞肉雞蛋.
買快餐.
大家很開心.
我們男人當然沒問題.
但如果是女人.
可能會有一點疑問.
不許女人講道.
是不是有些文化背景我們不太明白呢.
或者當時可能有些學者覺得.
當時正是羅馬帝國裡面.
有一種新的潮流.
就是女人有新的髮色.
特別是跟第一女郎.
羅馬皇帝的太太.
有很多雕像找到.
她的髮色扮得很漂亮.
其中一個老婆叫Livia.
第一夫人Livia是這樣穿衣服.
這樣去裝飾.
成為一個當時的潮流.
所以他就寫了一段聖經說.
女人不要跟這些.
所謂新羅馬女士的潮流.
可能是這樣.
譬如今天我們總統第一夫人.
穿什麼衣服.
也有很多宣傳.
你會跟她嗎.

$^{241}$OK的.
不會吧.
會跟第一夫人的服飾等等.
會不會是保羅說.
不要叫教會盲目跟潮流呢.
有沒有新的指引.
讓我們去了解這段聖經呢.
是有的.
是近這幾年.
為什麼近幾年這麼奇怪呢.
讓我告訴你到底發生什麼事.
近幾年有一個人在美國.
他做古代社會研究.
他專門研究當時社會.
以弗所這個地方的有錢人是怎樣生活的.
他做博士論文的時候.
他看了很多希臘羅馬的文學小說.
剛才Emily提到有小說.
他突然間看到一本小說.
是一個希臘的哲學家寫的.
那本小說叫做.
Ephesiaca.
這個Zenophil.
正是以弗所裡面的一個作家.
他寫了一本小說.
大概有幾冊.
講到什麼呢.
講到一個以弗所的男女的愛情故事.
有什麼大不了呢.
大不了的地方就在這裡.
就是以前那些人以為那本小說.
是主後二百多年才寫的.
然後十幾年前之後.
有人再重新寫過.
重新擺放它應該屬於哪個時期.
正是屬於第一世紀五十年代.
六十年代的時候.
有什麼大不了呢.
皮摩泰前書和保羅皮摩泰.
所遇見的以弗所的處境.

$^{281}$正是第一世紀五六十年代的時候.
所以將文學著作放回這個地方.
就很厲害了.
你看這本著作.
看這本小說的時候.
你就知道當時的人在想什麼.
然後你再扣回去.
到底聖經有沒有跟這件事有相似的地方呢.
我想告訴你.
文學著作寫作的時間是很重要的.
你知不知道人們說魏康為什麼這麼厲害.
因為魏康幾十年前.
這裡有沒有魏康迷呢.
應該有吧.
有不少吧.
幾十年前他寫的科幻小說.
他已經想像到現在生活是怎麼樣的.
以前沒有互聯網.
他已經想到會有互聯網這件事.
所以如果你說魏康現在寫.
那就沒有什麼大不了.
已經出現了.
但幾十年前他已經想到這件事.
就知道他想像的空間有多厲害.
所以寫作的著作是很重要的.
這個愛情故事是說什麼呢.
就是說一個十幾歲的帥哥和美女.
男神和女神.
他們都是異異所的人.
異異所有一個很大的廟.
叫做阿底.
Artemis.
每年會有巡遊.
新年的時候有巡遊.
他們這對男女.
在巡遊中一見鍾情.
但是又不夠膽向對方表白.
然後就單戀成病了.
如果你是父母你怎麼辦.
看到你的子女各自單戀成病.

$^{321}$他們就去一個廟.
找一個神靈.
去問一下姻緣.
那個人就說.
如果這對男女走在一起.
一定不會有好日子過.
他們一定會被俘虜.
然後有海盜,水險,火險等等.
但是父母你會怎麼做.
因為他們單戀成病.
男神女神好像有段姻緣.
不想他們受苦.
所以兩方面的父母.
就讓這對男女閃電結婚.
那些不好的.
那些會有災禍發生的怎麼辦.
閃電結婚之後.
父母就送他們走.
搭船送他們去埃及.
送他們去埃及的時候.
誰知道就真的有沉船.
誰知道就真的有海盜.
誰知道就真的有火險,水險等等.
兩個人愛人就被迫分開.
一到那女生想服毒自殺.
但是又不行.
死不去.
死了幾個小時又再起身.
然後兩個人又被賣到不同的地方.
男的就決定萬里尋妻記.
結局是怎樣呢.
兩個人又再重遇.
就回到異佛所.
在亞底米女神的廟裡奉獻.
是可歌可泣的劇情.
這些劇情是不是很熟悉.
是不是.
有點像韓劇.
就是這樣.
但是我想告訴你最重要的是什麼.

$^{361}$那個小說不只是愛情小說那麼簡單.
那個小說是講到異佛所的宗教文化.
那個女生為什麼最後有一個好的姻緣.
原因是因為她很敬畏亞底米女神.
所以你扣下來就知道.
異佛所當時是一個很重要的女神.
那個女神叫做亞底米.
Orthomus 或者羅馬叫她Diana.
她就是守護異佛所的一個大女神.
當時我們大家都會去旅遊.
旅遊就一定會打卡.
有些地方一定會去.
古代的人有七個地方一定會去.
古代世界的文明.
Seven Wonders of Ancient World.
其中一個就是亞底米女神的廟.
那個大的廟.
那裡有如今天的拉斯維加斯.
那裡的廟是一個女神的廟.
每年會巡遊.
也有很多女性在當中做祭司.
做先知.
做領導.
特別是一些貴族的女子.
每逢巡遊的時候.
她就需要打扮到某一個方法.
某一種的髮色.
然後她就模仿那個女神.
當她模仿那個女神的時候.
那個女神就會祝福她.
如果她不模仿那個女神的時候.
那個女神就會.
亞底米女神你看到她的樣子就知道她很兇.
她就會降禍給你.
所以為了討好她.
貴族的女子.
就會在這個神廟裡面來事奉.
甚至大聲的來頌土.
甚至宣揚神話.
這是故事的背景.

$^{401}$但是很奇怪.
當你將這個小說.
和提摩塔前書第二章.
這段聖經放在一起平衡的看的時候.
你就發現起碼有八個字.
專有的字.
是兩邊都有用的.
有什麼字呢?.
例如女人.
字律.
裝飾.
編髮.
金飾.
名貴等等.
善行.
或者敬畏神.
或者敬畏神明.
所以裝飾才算.
女孩子越模仿越像那個女神.
就像那個小說裡面的女主角.
她就越得神的保護.
編髮也是一樣.
當時貴族的女生就是一種編髮.
就表示她對亞底米女神的敬畏.
她那種頭髮的打扮梳髮.
就好像那個女神的髮色一樣.
所以不單單是一種興朝流的問題.
而是一種方法.
來敬畏亞底米女神.
以弗所的女人.
你想想如果從小到大.
她都受著這種異教的影響的時候.
你可以想像.
當她們信了耶穌.
到了以弗所教會裡面.
提摩太就是她們的牧者.
就是一群女人她們信了主.
你可以預期.
她們以前的生活.
在亞底米女神廟裡面的那種價值.

$^{441}$那種做法.
她就帶進了以弗所的教會.
以前是怎樣敬拜神.
假神.
現在進到以弗所的教會裡面.
她就用回同樣的方法.
你可以想像在以弗所的提摩太.
剛剛開始默會的時候.
就令她有很大的問題.
會不會令人混雜了.
基督教的教會是不同的.
跟外面是一樣的.
你可以想像到她就寫信求救.
問保羅她的師父.
到底怎樣算數.
所以如果你這樣看的時候.
你就會明白.
第九.
第十節.
那裡說到打扮.
他就叫那些女人.
不要再像以前在神廟裡面那種打扮.
因為上帝不是這樣讓她得到喜悅的.
那時候是你以前在亞底米女神裡面.
你得到她歡心的方法.
現在你來到教會.
不不不.
以前以為.
所謂拜得神多自有神庇佑.
你越像女神就越有女神庇佑.
有一種舊有的習慣.
亂帶進來教會的時候.
保羅就說不是.
在教會在基督裡面.
就掃了一種新的文化.
那種文化不是看你的外表.
而是看你內在.
生命.
品格的培養.
就是人看見你信主之後.

$^{481}$有的改變.
在裡面發出來的改變.
這個才是福音.
這個才是敬拜.
這個才是基督徒的生活.
其實說這些不是很抽離的.
我記得很多年前我去以色列的時候.
我認識了一個女孩.
這個姐妹.
她的先生是猶太人.
猶太人知道她不信耶穌.
大部分都不信.
並且反對.
她一直都很反對這個姐妹去教會.
但是這個姐妹很堅持.
我要繼續回去.
起初她的老公非常反對.
大家經常吵架.
但是過了一段時間.
她的老公開始不吵了.
因為她的老公覺得很奇怪.
因為他老婆回來教會之後.
她的心情特別好一點.
老公就覺得心情好.
Happy wife, happy life.
Why not?.
雖然她還沒有信耶穌.
但是她覺得教會做了一些事.
挺好的.
起碼對她的生命有些改變.
不單止她沒有反對.
她甚至鼓勵她回去.
這種就是內在生命的改變.
就是那種呼吸的感染力.
保羅在《提摩太前書》第二章.
講到這些姐妹.
這些女人.
不要像以前那樣.
去跟神相處.
要重新整理過.

$^{521}$然後用自己內在生命的改變去表達出來.
11至14節.
那裡講到安靜學習.
講到不許女人講道.
不許管轄男人.
然後又講到先做的是阿當.
後做的是夏娃.
這個也可能有這個背景.
因為當時亞底米女神廟.
裡面的女孩子.
每逢節日的時候.
她們就要在廟裡面.
很大聲的唸經.
歌頌女神.
因為女神就是眾生之母.
連男人也是她生出來的.
神話裡面.
講到亞底米女神幫她媽媽生她的兄弟.
如果沒有了女神.
她的兄弟就會死了.
所以女人就大一點.
各地的旅客去到伊佛所的時候.
去亞底米女神朝聖打卡的時候.
他們就是聽這些神話.
如果你這樣背景看的時候.
你就會了解為什麼他說女人.
在教會裡面.
是需要安靜來學習的.
你想想那些在異教裡面.
convert了.
信了耶穌的女士.
她們以為我們好像以前一樣.
以前也是這樣.
祈禱這樣去教.
為什麼來到教會不行呢?.
保羅不是說不行.
而是說他們需要安靜才去學習.
裝備好自己.
然後做一個合格的老師.
你記得我們幾個星期前.

$^{561}$《瑪達瑪利亞》這個故事.
坐在耶穌的腳前.
其實就是這個意思.
他就可以變成.
Torah scholar.
安靜學習.
就是將來他能夠成為一個合格的老師.
我不許女人講道.
這個也是有時候我帶兩個小孩出去.
我女兒突然跟我說.
不如我開車吧.
我說不.
因為我沒有試過開車的課程.
我說我不許你開車.
不代表我永久地不許你開車.
裝備好你還沒學好.
其實也是.
不許管轄男人.
這個是亞歷米女神宗教.
女人當道.
似乎跟這個宗教的背景是有關的.
他們覺得.
就像以前一樣.
我們是在開演唱會.
保羅說不是.
是男女一起來做帶領.
亞歷米女神宗教.
我講過.
也是女人當道.
所以他們解釋創世記的時候.
可能會調轉來解讀.
所以保羅在這裡來矯正.
先做的是亞當.
然後做的是夏娃.
第十五節就要守的人.
女人若常存信心,愛心,聖潔至守.
就必在生產上得救.
這個真的很麻煩.
上個星期的主學就知道.
生產可以是一個比喻的手法.

$^{601}$正如聖經裡面說.
基督徒的生命要結果子.
不是說我們一起過基督生活成長的時候.
我們家庭就多了一件東西.
不是這個意思.
是一個比喻的手法.
結果子就是說我們有一個信得的生命的改變.
出來了.
有喜樂,和平,節制等等.
聖靈的果子.
可以是一個比喻的手法.
還有沒有其他背景呢?.
也可能有.
因為亞當與什麼女神裡面的神話呢?.
正如我所說.
是他幫他媽媽去生那個孿生兄弟.
所以亞底米女神是什麼呢?.
是二佛所的婦科聖手.
就是說你一定要尊重.
一定要敬奉亞底米.
然後每一個女生.
當時古代的社會就一定會生孩子.
大部分都會.
99.9\%當時的女人都沒有選擇.
她一定要做媽媽.
所以她們就要去敬畏亞底米女神.
否則你生孩子的時候不止痛.
一定會死的.
你知道當時去產房.
不是今天那麼安全.
今天比較安全一點.
當時是很危險的.
所以他們就很怕得罪亞底米女神.
保羅說不是.
你已經信了主了.
不是那個女神保護你.
而是耶穌基督去保守你.
所以必在生產上得救.
你不需要怕那位以前的女神.
你要怕的,你要敬畏的.

$^{641}$就是真神耶穌基督.
這個對我們來說.
應用其實是很切身的.
我們華人很多時候.
我們有一些以前的包袱.
我們經常說.
不少大兄姐妹都有這樣的經歷.
我信了主.
我以前是要裝香的.
如果信了主就不能裝香.
那不是很大件事嗎.
可能會不會有遭殃呢.
信了主會不會有所改變呢.
還可以怎麼應用呢.
以前可能你覺得.
信耶穌只不過是上天堂.
我拿個護照.
就在意萬一有什麼事.
起碼可以去天堂.
這是舊有的觀念.
現在新要有的觀念.
保羅就是要等上天堂.
信了耶穌就有一個交托.
有一個方向.
可能你以前信主之前.
或者你以往舊有的觀念.
只不過是信耶穌.
我只是拿東西.
我拿到什麼好處.
現在你就需要180度改變.
改變過來原來不是get get get.
而是give give and give.
以前可能你的信仰生活.
只不過是間中有時間.
我就去敬拜.
我就去聚會.
沒時間就把它放在一邊.
可有可無做一個註腳.
今天就不是了.
上帝的生命要改變你.

$^{681}$你要完完全全地改變.
跟這些女人都一樣.
我們被福音,被上帝恩典改變.
今天正是章林節開始.
我想提醒你.
不只是點一下蠟燭.
拍一下照片.
這麼簡單.
你迎接上帝進入你的生命.
就讓上帝在你整個生命裡面.
作改變.
作更新.
改變你的價值.
改變你追求的方向.
在advent當中想一想.
到底有什麼你還是扣住.
不想上帝來改變.
可能你是看時間的價值.
可能你是看金錢的價值.
可能你是看一些關係.
有些事你是不是不想let go.
上帝臨到你人生就需要去改變.
我們同心低頭祈禱.
上主我們願意你來改變我們.
願意你進入到我們生命當中.
去更新我們.
用你的愛,恩典來幫助我們.
讓我們畏著這個信仰.
畏著你的恩典.
我們作出回應.
更多事奉你,更多愛你.
將你自己的真理帶進我們人生當中.
改變我們的談吐.
改變我們的所作所為.
改變我們一舉一動.
讓人見到耶穌基督在我們身上.
讓我們內在生命能夠發放耶穌基督那種恩典.
我們這樣祈禱.
奉主耶穌基督的名字.
\newpage



\section{耶利米書 29:11}
\label{sec:zzrp0tNV2iI}
\textbf{1/6/19 主日講道: 2018 最多人搜尋的聖經 (陳崇基牧師) - 耶利米書 29\_11}
\newline
\newline
連結: \href{https://youtube.com/watch?v=zzrp0tNV2iI}{\texttt{https://youtube.com/watch?v=zzrp0tNV2iI}} ~~~~ 語音日期: 2019-07-25
\newline
\newline
\hyperref[sec:tDpoW3ZdL20]{\small{< < < PREV SERMON < < <}}
~
\hyperref[sec:index_chronic]{\small{[返順時目]}}
~
\hyperref[sec:index_scriptual]{\small{[返順卷目]}}
~
\hyperref[sec:y_ka_wkoG0A]{\small{> > > NEXT SERMON > > >}}
\newline
\newline
耶利米書 29:11
\newline
\begin{longtable}{cl}
\hline
\hline
章節 & 經文 (和合本修訂版)\\
\hline
29:11 & \begin{tabularx}{0.7\textwidth}{X} 我知道我向你們所懷的意念是賜平安的意念，不是降災禍的意念，要叫你們末後有指望。這是耶和華說的。 \end{tabularx} \\ \\
[1ex]
\hline
\hline
\end{longtable}
$^{1}$(寶寶叫聲).
不知道你去年在Google或電話等地方.
你搜尋了甚麼.
可能是找餐館或是找地方.
今年是2019年.
去年年尾2018年的時候.
YouVersion和Bible Gateway.
這兩個算是最大的.
專門搞網路聖經.
或者是網路聖經.
或者是手機聖經的App.
他們就發佈了一些大數據的回收.
分析過之後.
他們發現原來去年.
在網路上看聖經的人.
或者是找聖經的東西的人.
是顯著有增加的.
很有趣吧.
顯著有增加.
全球顯著有增加.
所以這幾個主日.
我們想看一下他們發佈的數據.
他們的數據告訴我們.
有些聖經的章節或經文.
是有很多人來分享.
或者寫筆記.
到底這些這麼受歡迎.
或者這麼多人搜索的聖經經文到底是什麼呢.
和我們有什麼生活上的關係呢.
所以我們想看看這些經文.
每一個星期我們看一節.
今天我們看一個最多人搜索的.
去年最多人搜索的一段經文.
就是耶利米蘇第29章.
29章第11節.
不知道你們會不會背呢.
耶利米蘇第29章第11節.
或者記不記得這段聖經是說什麼的呢.
這段聖經是很有名的.
不只是去年才是最受歡迎的.

$^{41}$其實這麼多年裡面.
都是首屈一指的.
很多人都是分享這句話.
特別是失戀或是有些難處.
他就會發這句經文給你.
可能會是這樣.
所以收到這段經文不一定是好事.
但是這段經文給了很多人鼓勵.
給了很多人提醒.
這段經文說什麼呢.
這段經文是.
耶和華說.
我知道我向你所懷的意念是賜平安的意念.
不是降災禍的意念.
要叫你在末後有子夢.
失戀的人合不合聽.
或者那個市倒了合不合聽.
都合聽.
撞車又合聽.
病又合聽.
所以這節的聖經.
很多人都覺得很受歡迎.
因為人之常情.
我們覺得到新年的時候.
當然我們希望明天會更好.
為什麼這麼多人會去倒數.
倒數有什麼不好玩的.
這麼累.
這麼睏.
有誰沒有去倒數的.
很多人睡覺.
但是人之常情.
我們內心的深處.
我們很希望將來會更加好.
無論你個人.
或者你的家庭.
或者你的職業.
你的學業.
不同的地方.
這個很真實.

$^{81}$看完聖經.
起碼就當它是揮春來看.
都給我們一種盼望.
但是我們想知道這段經文到底是說什麼.
我們經常說.
看完聖經.
一定要有上下文.
因為你看到它是第十一節.
你就知道不是這一節開頭.
到底上下文說的是什麼.
或者看看上下文.
我們會更加明白這段聖經更加深的意思.
到底什麼叫沒後有子亡.
二十九章第一節開始的時候.
耶利米書是這樣說的.
他說先知耶利米.
從耶路撒冷.
被避路的.
流亡的.
祭司先知和眾民.
並生存的.
就是死剩的長老.
就是尼泊爾格尼撒巴比倫王.
從耶路撒冷.
俘虜到巴比倫那一幫人.
耶利米就留在耶路撒冷.
這個是主善五百八十六年.
當巴比倫軍隊去到耶路撒冷的時候.
他們就洗劫整個耶路撒冷.
死了很多人.
然後將那些國家的精英分子.
那些祭司先知識字的人.
有些教育的人.
或者有很多其他.
等等的.
長老啊.
眾民啊.
他們就將他們俘虜到一個異鄉.
就是巴比倫.
聖殿就燒毀了.

$^{121}$然後有大量的猶太人.
神的子民.
就因為這樣被迫去到巴比倫.
或者你看這種事蹟.
這種背景的時候.
可能有一些畫面可以幫到你.
什麼叫做被擄或者流亡呢.
XR.
可能這個畫面可以幫到你.
就是二次大戰的時候.
德國納粹黨.
去到屠殺猶太人的時候.
有很多的家庭.
有很多的人被迫.
然後要面對很多的恐怖.
死亡.
去到集中營.
或者可能這個畫面可以幫到你.
就是今天我們看到全球有很多不同的地方.
因為戰亂因為動盪.
有很多的難民.
如果你用這套畫來看.
耶利米蘇29章的時候.
你就可能會更加逼真一點.
我們經常說什麼被擄啊.
XR.
或者流亡啊.
到底是什麼意思呢.
我想.
我也很少說正確的.
詳細的去定義.
到底什麼叫做XR.
什麼叫做流亡呢.
上個星期吳牧師講道的時候.
就引用一個聖經學者.
那個聖經學者就專門去研究.
這段時期的文學作品.
就是這段時期到底背景發生什麼事.
這位女神學家.
她就這樣去將流亡去定義.

$^{161}$就算你Google什麼叫做被擄.
什麼叫做流亡.
你都很難找到一個這麼詳盡的定義.
定義是什麼呢.
她說流亡不是僅僅無家可歸.
相反你本身有家.
但是你的家已經被敵人強佔了.
流亡不是僅僅無根.
相反你本身根深蒂固.
現在呢.
他被連根拔起.
被迫活在異地裡面.
你的根就被暴露無遺.
然後你很痛苦去掙扎.
寒冬嘲笑裡面.
敵人嘲笑裡面.
你渴望回到你原根的土壤.
但是你作為一個流亡者.
固土所屬.
卻深知你無能回家.
仍然未能.
或者永遠都不能.
這個就叫做Exile.
這個就叫做流亡.
流亡不僅僅是地理上的遷移.
不是地理上的遷移.
其實可能是心理上的遷移.
可能是心理上的遷移.
或者你心裡面精神上有一些流失.
一些失所迷失了.
或者在事業上都可以有很多流亡.
可以在你前途上有很多流亡.
心底裡面你總之感覺不到你所屬的地方.
你就是一個流亡的狀態.
對於舊約的神的子民來說.
更加慘的是什麼呢?.
因為進入到信仰的層面.
因為當你居住在異地.
被巴比倫殲滅了的一群民族的時候.
你就知道到底誰贏了.

$^{201}$和誰輸了.
我本身是敬拜耶和華.
屬於耶和華的.
屬於天地萬物的真神.
現在我被流亡了.
就是說我的上帝失敗了.
那誰贏了呢?.
就是巴比倫的眾神贏了.
那些偶像贏了.
所以糾纏和掙扎和無力是很深的.
我講到流亡的定義的時候.
我相信有不少人都有心同感受.
或者未必去到那個地步.
但是我們心底裡.
或者有一點點覺得.
我也可能可以去認識到.
在過年前後.
因為我預備一些東西.
我就去YouTube看了兩個節目.
兩個節目都是紀錄2018年的香港.
就是說香港這個社會.
有多少的改變.
這兩個節目說到有很多的港人.
正在考慮移民.
或者想離開香港.
總之也有人覺得只是想走.
我已經離開了香港很久了.
我已經很脫節了.
但我看完這兩個節目之後.
我真的有一點點失落感.
有一點無力感.
人在一個社會裡面.
本身是他根直的地方.
但開始失去他的語言或者文化.
或者他自己的身份.
節目裡面也有提過.
甚至台灣一個講普通話的地方.
就拍了一個廣告.
用廣東話拍了一個廣告.
就拍給香港人.

$^{241}$就叫香港人不如移民來台灣.
我看到這裡覺得很奇怪.
也有很多的.
你可以叫他做流亡感覺.
我們知道我們美國對上加拿大.
回到加拿大的人.
是有上升的趨勢.
上一年底的時候.
去多倫多玩玩.
認識了一對年輕的牧者.
他們在多倫多的牧會.
大家在閒談當中.
他跟我說.
內行人如果知道那個形勢的時候.
他說如果只是在多倫多.
只是做這一班人的事工.
只是做這一班回到加拿大的人的教會事工.
也可以做幾十年.
你說厲不厲害.
所以我們就經常想去多倫多.
有很多無力感.
或者有很多的悲觀.
對現實生活.
或者我們不一定是香港人.
我們是講廣東話.
懂得聽廣東話.
說真的我們坐在這裡.
來到這裡.
我們都是所謂流亡於美國.
如果有人問你.
你死的時候想葬在哪裡.
如果價錢都是一樣的話.
你或者不會選這裡.
這裡冰天雪地.
很冷.
當然是去一些暖的地方.
或者你的故鄉.
東南亞也好.
台灣也好.
或者廣東其他地方.

$^{281}$或者有不少我們華人來到美國.
我們都有種流亡的感覺.
以往我們在故鄉的時候.
我們可能是專業人士.
來到美國我們就要洗碗.
要在餐館裡面捱.
這個就是20世紀的exile.
displacement.
是社會或者政治的狀況有不穩.
所以我們離開了某些東西.
也知道那裡是我們的家.
但是我們連根拔起.
我們也知道我們回不去.
起碼現在還不能夠回去.
甚至今天在美國.
有很多人都想搬去加拿大.
你不要說華人.
近兩年因為政治很多的分化.
社會很不友善.
很不歡迎外來者.
其實美國公民移民加拿大.
也有上升.
雖然不是很多.
也有些去了之後後悔.
但是怎樣也好.
總是一種exile的感覺.
我們一年的.
全年的讀經計劃開始.
通常讀經計劃.
從新年開始的時候.
看創世記頭兩章.
就一定是家常便飯.
看完兩章之後.
你也不用繼續看下去.
很快你就會看到人類在流亡.
本身人類被造的時候.
就是和上帝有團契有關係.
是一個美滿的關係.
那個叫伊甸園.
到第三章的時候.

$^{321}$人類就被丟出伊甸園.
這個就是一個exile.
一個流亡.
是回不去家的.
只不過上帝沒有放棄我們.
上帝就將他的居所建在人間.
有會幕有聖殿.
然後基督親自來到.
到最後啟示了上帝將天家.
將天上的耶路撒冷帶到地上.
在基督裡面在主裡面.
就有永恆的永久的家鄉.
這個就是流亡.
還有我們作為一些exile的人.
在基督的恩典裡面.
在基督的救恩裡面.
我們有一個永恆的身份.
這個就是耶利米面對的處境.
神就叫耶利米做什麼呢.
叫他寫信給流亡於巴比倫的神的子民.
信上面寫什麼呢.
他說萬官之猶說.
以色列上帝對一切被擄流亡的.
就是我使他們從耶路撒冷被擄巴比倫的.
如此說你們要蓋造房屋住在其上.
栽種田園吃其中所產的.
娶妻生兒女.
為你們的兒女也娶妻.
使你們的女兒嫁人生兒養女.
在那裡就是在巴比倫.
你們要生養眾多不至減少.
我所使你們被擄流亡到那城.
你們要為那城求平安.
為那城禱告耶和華.
因為那城得平安你們有一脆著可以得平安.
萬官之猶說以色列上帝如此說.
不要被你們中間的先知和占卜的引誘你們.
也不要聽信自己所做的夢.
因為他們託我的名對你們說的是假預言.
我並沒有差遣他們.

$^{361}$這事耶和華說的.
這很奇怪.
被流亡的人當然想回家.
特別對猶太人來說.
在一個外邦的地方.
一個不潔淨的地方.
他們當然很想夢寐以求每一天想回家.
在外國中間.
你來了美國沒多久的人.
你剛剛來那幾年是不是想回家.
只不過我們來的時間久了.
我們老殘 美殘.
我們就覺得這裡就是我們的家.
他們當然很想回家.
但是很奇怪.
信中說什麼.
你不是告訴神的子民.
他說不用想回家.
還要做什麼.
你們留在他鄉.
在那裡耕種然後建立家庭.
然後有日常的生活.
這個很有趣.
你不喜歡一個地方你想走.
偏偏神就說你留下.
你喜歡芝加哥嗎.
一般吧.
假設你不喜歡芝加哥.
你很想離開.
找個你自己心中夢寐以求的地方.
你不喜歡芝加哥的城市.
想搬出去.
或者搬去別的地方.
搬到西岸也好.
你不喜歡這裡的天氣.
說真的來了這麼久.
我來了差不多20年了.
這裡好像亂了.
沒以前那麼冷.
但怎樣也好.

$^{401}$你有沒有想過這樣祈禱的時候.
上帝偏偏不回答你.
或者跟你說你留下.
因為你留下的時候.
你原來有一個見證要做.
原來你有一個使命要做.
原來你人生的使命不是找回你自己的故鄉.
人生的使命是以上帝為你故鄉.
這裡就是耶利米對神的子民說.
你還要留下.
並且留三代.
其實耶利米有預言說過.
他們流亡要多久.
起碼要70年.
起碼70年.
如果你是被流亡的那個.
你會不會等到70年完結了.
然後回歸.
沒有吧.
你不會見到的.
就算讓你看到你也走不動.
70年之後才可以回歸.
有生之年你未必一定見到上帝應許的實現.
不如你不要說.
那是OK的.
因為此時此地.
此刻.
你這個人生.
原來是有見證可做.
他說要為那個地方求平安.
平安說了幾次.
為那個城市求平安.
猶太人在愛邦的地方.
愛邦的世界.
就正是見證神的時間.
你比較一下第一節我們剛才說的.
你記得的話.
第一節說為什麼猶太人會流亡於巴比倫.
就是因為巴比倫王尼布格尼薩打敗了他們.
然後抓了他們去巴比倫.

$^{441}$主體是巴比倫王.
是一個愛邦不認識神.
不敬畏神的君王.
到第四節.
這裡就將這件事重新定義.
流亡是誰做主角.
就是耶和華神自己.
原來是神去差派他們.
使他們被擄到巴比倫.
流亡到這個地方.
所以他們在巴比倫.
在異鄉.
不是一件偶然的事.
而是上帝親自安排帶他們來.
你試一下應用在你自己身上.
你可能會找到一個嶄新的人生觀.
作為流亡的人.
你不會再日日夜夜在想.
到底我何時可以回到過去.
你最多只是想著.
偶爾可能去多倫多.
吃一下那裡的東西.
那些是度假.
你不會夢寐以求回到馬來西亞.
我回到新加坡.
我回到香港.
我回到中國.
你反而要想的是.
在這個流亡的地方.
你沒有地上的歸屬感.
但是神仍然跟你同在.
神仍然對你說話.
或者會容你.
在流亡的時候.
就在逆勢裡面.
在受苦當中.
在困難裡面.
偏偏很多時候聖經一次又一次告訴我們.
原來在我們的苦難.
在我們的流亡當中.

$^{481}$不是結束我們跟神的關係.
反而是我們跟神的關係.
屬靈生命一個新的開始.
It's not the end.
but a new beginning.
上帝更加真實.
沒有了以往那種歸屬.
沒有了以往那種靠山.
現在上帝屬靈的事.
可以更加真實.
你不要寫信給他們.
告訴他們.
你們就要留下.
為那城求平安.
你有道可傳.
有見證可以做.
第二.
方信說什麼.
他繼續說.
耶路撒冷說.
為巴比倫所定的七十年滿了之後.
我就會眷顧你們.
使你們成就我的恩賢.
使你們迎回此地.
即是回到耶路撒冷.
他說我知道我向你們所懷的意念.
是治平安的意念.
不是降災禍的意念.
要叫你們沒後有子亡.
這句就是那句金句.
你們要呼求我.
禱告我.
我就應運你們.
你們要在巴比倫尋求我.
若專心一意尋求我.
就必尋見上帝.
他說我必被你們尋見.
我也必使你們被擄流亡的人歸回.
將你們從各國中.
所趕你們到各處的焦聚.

$^{521}$返回來.
又將你們帶回.
我使你們被擄離開的故鄉.
這是耶和華所說的話.
最受歡迎的那段經文.
就是保留在這裡.
偏偏最受歡迎.
很多時候我們就不理上上文下理.
我們就把它亂解.
就當上帝一定會使我明天會更好.
使我一定會風生水起.
耶利米偏偏就在這裡說.
和什麼人對話呢.
對抗什麼人呢.
就是對抗在巴比倫裡面的假先知.
那些假先知就是報喜不報憂.
他們在說什麼呢.
剛剛去到被擄的時候.
他們就跟那些被擄的猶太人說.
我們很快就可以走了.
一兩年之後.
我們就會回到猶太人那裡.
因為上帝一定會眷顧我們.
耶利米偏偏就說.
這是謊言.
你們是不會走的.
你們不如就滴起心肝.
在這裡落根.
上帝信仰他.
信仰基督.
不一定會脫苦海的.
不一定會脫苦海.
但是在流亡裡面.
上帝可以和我們同行.
這個才是真理.
這個才是真的先知.
雖然聽起來很硬.
很不舒服.
因為我是想回去.
但是上帝說你就要留下.

$^{561}$這很難說.
很難吞下.
但是很多時候.
忠言就是逆耳.
真理就是這樣.
你知不知道去年2018年.
時代雜誌每年都會選.
Person of the Year.
就是全年的風雲人物.
最有影響的人物.
上年2018年.
時代雜誌選的風雲人物是什麼.
Person of the Year.
是Journalist.
是新聞界的工作者.
那些被打壓.
或者會受害的.
仍然勇於說話的.
世界各地的新聞工作者.
這些就是真相的守護者.
Guardians.
如果你說到舊約裡面.
聖經裡面.
就像那些真的先知.
真的先知就是上帝真理的代言人.
或者上帝真理的Guardians.
他們就是守護著真相.
守護著真理.
縱使有困難.
仍然要說出來.
這些就是先知.
就是Prophet.
如果說到這裡的時候.
我覺得我們起碼要加上一個.
我起碼也會加上一個.
就是在信仰上的Guardian.
去年大家知道中國大陸.
打壓教會.
非常厲害.
在中國大陸成都的一個教會牧師叫做王怡.

$^{601}$在一個叫做秋雨教會.
的一個牧者.
我在一個華人北美的微信組群裡面.
去年年底不斷收到一些代討的通知.
為了王怡.
為了他的教會.
為了他的家庭.
兩夫妻為了師母來祈禱.
因為他被控煽動顛覆國家的罪名.
然後入獄.
太太也是一樣.
王怡牧師.
他說過一番話.
在華人牧者的組群裡面經常出現.
就是年底的時候有人發出他的其中一段話.
他說人生活著是為三件事.
第一是跪著禱告.
站著講道.
和跑著宣教.
這個就是我們人生屬神的生命.
我們最重要做的三件事.
跪下我們去祈禱.
我們是依靠上帝的能力.
並且我們站起來.
站得直我們去講真理.
第三是跑著趕著去宣教.
趕著去使人認識主.
人生所做的最基本最永恆價值的.
就是這三件事.
不是吃喝玩樂.
也不是搬屋.
也不是買車.
就是這三件事.
跪著禱告.
站著.
站得直去祈禱.
站得直去講真理.
跑著的人去宣教.
主比感激.
三個的wisdom.

$^{641}$就是隨時搬家.
隨時坐牢.
隨時回天家.
這個就是天國的生命.
當頭棒喝.
你想想你自己.
你安書了多久.
你想想你最近祈禱.
或者你去年祈禱.
你為了什麼祈禱.
或者你可能為了一點點.
你的portfolio幾點上落去禱告.
可能你在想20年後30年後.
我的退休金能不能留在這裡.
我可不可以安安心心過完我晚年.
可能你就是在想這件事.
沒錯.
這些是尾時.
但不是最重要的事.
要分清楚priority.
有什麼是永恆.
有什麼是暫時.
如果看得清楚的時候.
我們生命裡面有多少隨時.
其實都ok的.
隨時我不回故鄉.
It's fine.
但我必須跪著禱告.
站著講真理.
並且跑出去宣教.
新的一年.
剛才我給了三個挑戰給大家.
去活出天國的生命.
我們連於地上的教會.
我們去關愛身邊其他人.
讓他們認識.
三個就是我們的目標.
你記住.
人生裡面.
要想清楚.

$^{681}$真是有永恆價值的東西.
真是能帶走的東西.
搞清楚我們的優先次序.
我們的生命就會向著主.
我們用心低頭祈禱.
主要我們低頭來禱告.
因為我們在天國當中.
在你的救恩裡面.
我們實在不配.
我們有很多愧欠的地方.
神我們求你來憐憫我們.
主叫我們轉變.
悔改.
再一次迎向你.
在我們人生當中.
我們搞好我們三件事.
禱告.
站得直的.
你都可以講真理.
講公義.
並且趕著的.
來宣教.
主我們求你來憐憫我們.
憐憫我們這裡每一位.
憐憫我們整個教會.
叫我們隨時的.
進行你自己的命.
無論去哪一處.
都是你自己的安排.
讓我們就在這個地方.
在這個角色裡面.
我們來到去見證你.
求你繼續帶領我們.
讓我們教會今年當中.
我們更加經歷你自己的恩典.
祈禱奉主名禱.
\newpage



\section{詩篇}
\label{sec:y_ka_wkoG0A}
\textbf{詩篇23 篇 (陳崇基牧師) 2018-06-24 主日講道}
\newline
\newline
連結: \href{https://youtube.com/watch?v=y-ka_wkoG0A}{\texttt{https://youtube.com/watch?v=y-ka\_wkoG0A}} ~~~~ 語音日期: 2019-07-25
\newline
\newline
\hyperref[sec:zzrp0tNV2iI]{\small{< < < PREV SERMON < < <}}
~
\hyperref[sec:index_chronic]{\small{[返順時目]}}
~
\hyperref[sec:index_scriptual]{\small{[返順卷目]}}
~
\hyperref[sec:nzhjFy_6Umo]{\small{> > > NEXT SERMON > > >}}
\newline
\newline
$^{1}$Cindy 和我們的報告.
今天他們說要走火警.
所以講到要短一點.
很簡單的,我就講短一點好嗎?.
等下有多一點時間,大家可以慢慢去做防火演習.
或者你今天可以這樣.
可以當今天是一個靈修.
有時候講到不一定是長才是重要的.
長有時候是逞氣的.
有時候短可能更加精簡.
所以今天他們說要走火警.
我們少一點時間,不如我們就講短一點.
我就選了一段聖經,一段詩篇.
在過去幾個星期,經常在我腦裡思想.
上星期我也有提過一點,詩篇90篇.
今天我不是講詩篇90篇.
有好幾段是我最近幾個星期.
經常掛在我心裡.
這篇詩篇叫做詩篇23篇.
是很熟悉的一段聖經.
有很多人可能會背.
如果你還未懂得背,我勸你背了它.
不知為何偏偏詩篇23篇.
就是每樣東西都適合用.
生日會也適合.
等下可能出去慶祝生日的時候又可以講詩篇23篇.
喪禮也適合.
又可以講詩篇23篇.
有時候畢業,學校畢業又可以講詩篇23篇.
結婚生子,什麼都可以,詩篇23篇.
但為什麼我想起23篇呢?.
就是最近跟我們社區或世界發生很多事.
令我們心裡很不安.
或者令我們心裡有些很不舒服的感覺.
或者你可能是親受其害的人.
你更加覺得最近你的心靈.
你的心靈好像整個圖畫粉碎了.
我們講詩篇23篇的時候.
你想一想到底最近有什麼令你心裡有點不安.
很多時候我們華人很喜歡將那件事壓了它.

$^{41}$掃它到地氈底.
我們就算了,不理它.
你想一想最近有什麼令你很不安呢?.
是什麼事呢?.
我們才來到神的話語面前.
有什麼呢?.
可能是一些死訊.
親人過去,或者一些朋友的親人.
或者是一些家人可能有重病.
社區裡面可能發生一些不安的事.
想一想那件事是什麼呢?.
或者可能是阿根廷書波.
令你心裡不安.
到底這篇詩篇想跟我們講些什麼呢?.
很多時候我發現我們現代人.
我們很多時候將我們的生活.
將我們的生命分成不同的軌統.
我們覺得有些事是專家才可以講的.
這件事是專家才講的.
這件事是另外一個專家講的.
沒錯的.
不過很多時候我們分得太細.
我們忘記了其實上帝說過.
有時候對我們生命,對我們心靈.
是有治療的作用的.
我不是說我們任何事.
我們翻開聖經就印印印而解.
我不是這個意思.
我們要看醫生,我們要看精神科醫生.
我們要看心理醫生.
是需要的.
不過我們不要去underestimate.
低估了上帝自己說話.
對我們的心靈.
對我們心靈的醫治的能力.
這就回到我們每一個基督徒.
為了我們信神的.
我們需要有的習慣.
我們翻開聖經.
不是一個屬靈的習俗.

$^{81}$而是我們用信心來領受上帝的話.
是生命之道.
我們不可以,我不知道怎樣再強調.
如果上帝是不說話的神.
那你大概可以燒一本聖經,沒所謂的.
但如果上帝真的好像聖經所說.
他自己在說話裡面所說.
他是一個會說話的上帝.
他是在每一代的信徒的心裡面.
對你會說話的.
你就要豎起耳朵聽清楚.
這篇23篇想跟我們說什麼.
我們先用信心來領受.
大概20分鐘.
我想跟你說一下.
這篇23篇想說什麼.
有一個很好的方法.
我覺得用圖片來說.
Picture is worth a thousand words.
百文不如一見.
我就選了,有六節聖經.
就選了六張相片.
我有的相片.
我覺得這些相片可以點出.
這篇23篇裡面.
一些今天對我們.
生命,心靈裡面.
有的不同的方面.
不同的意義.
23篇第一節.
耶和華是我的牧者.
我必不自缺乏.
看聖經的時候.
我們知道背景是中東.
中東,特別是古代.
只有兩樣.
經濟裡面只有兩件事.
頂住整個生活.
整個經濟人的生活.
只有兩件事.

$^{121}$你可以想像一下.
芝加哥可能有IT.
有醫療.
有不同的工業.
來發展.
頂住經濟.
在中東只有兩件事.
第一,耕田.
第二,芳陽.
牛,羊,馬,奴仔等等.
畜牧業.
這篇詩篇是說牧羊.
用一個羊仔的角度來看人生.
羊我們知道是很軟弱的.
不能夠自保的.
牠的生存是靠牧羊人.
Shepherd帶著牠.
才能夠生存.
所以他說神.
耶和華是我的牧者.
我必不自缺乏.
其實原文是將來式.
就是說我將來不會缺乏.
現在牠是一個缺乏的狀況.
是一個不太安全.
或者缺水,缺食物的時間.
但這篇詩篇宣告牠去信靠一位牧者.
因為這位牧者是牠的牧者.
牧羊人就像今天的天文台.
牠懂得去預測天氣.
牠不可以呼風喚雨.
但牠知道何時會有雨落.
何時會有風來.
牧羊人應該帶羊去哪裡呢?.
哪個山頭最好呢?.
牧羊人就是做一個這樣的角色.
就是供應和保護那隻羊.
第二節.
他說牠必使我.
刊我在青草地上.

$^{161}$令我在寒安傑的水邊.
大家可以看到這張照片.
不是那些一片綠油油.
一望無際的大草原.
在巴勒斯坦以色列.
牧羊呢.
這個叫青草地.
這個叫溪水棚.
巴勒斯坦地的季節,天氣.
分為兩頭.
寒季,夏季.
四,五月至九月.
雨季,冬天.
秋天,冬天等等.
即是九月,十月至三,四月.
只有兩頭.
春夏秋都不太分明.
像我們芝加哥一樣.
只有夏季.
大概冬季.
秋季冬季.
五月至九月.
即是我們現在這段時間.
我們芝加哥很多雨.
我不知道你特別濕濁.
各位大兄姐妹.
可能水浸.
雨多到不得了.
有時候走走雷.
下下暴雨.
在巴勒斯坦這個時間.
是最乾旱的時間.
詩篇裡面經常有這樣的說話.
乾旱疲乏無水之地.
詩篇裡面有三篇.
沒有雨下的.
整個天都是藍色.
藍到好像很假的藍.
一塊雲都沒有.
是非常乾燥的.

$^{201}$他說青草地.
青草.
這裡希伯來文的意思.
是說八月,九月的時候.
第一次下的雨.
一下在那些曠野的地方.
本身寸草不生的.
它就長出一些很小的綠草.
那個叫做青草地.
又或者這個時間.
五月至九月寒季的時候.
早上天朦朧的時候.
這些曠野.
是霧土.
有很多石頭.
石頭與石頭中間.
因為早上有霧水吹過.
石頭中間有少許水分.
然後就出少許草.
所以羊過去的時候.
撩開那些石頭.
它就有青草.
剛剛嫩芽發出來的草.
過了下午就沒有了.
因為太陽一曬它就死了.
這個叫做青草地.
剛剛足夠我們的需用.
神就去供應我們.
這裡說.
他使我躺我在青草地.
使我在可安歇的水邊.
羊很少.
現在看清楚羊的影片.
羊很少坐在樹上.
因為羊很害怕.
四隻腳經常站在那裡.
什麼時候可以躺我.
什麼時候可以坐在水邊.
是它感覺到最安全的時候.
最安心的時候.

$^{241}$那種狀況.
或者最飽.
不再口渴.
不再去找其他東西的時候.
羊才會坐下.
四隻腳跪下.
在水邊.
上帝.
我們必須要的飲食.
那牧者賜給我們.
第三節.
他使我靈魂蘇醒.
為自己的名引導我走二路.
他說靈魂蘇醒其實.
廣東話說得貼切一點.
就是死過翻生的意思.
差不多死了.
渴死也好餓死也好.
或者可能有猛獸來攻擊我的時候.
我差不多沒命的時候.
他使我的靈魂蘇醒.
救回我的命仔.
死過翻生的意思.
一隻羊可以很固執.
這張照片是約旦南部的曠野拍的.
那時候我跟著一些牧羊人.
他們早上就帶羊出去放牧.
我們隨後就跟著.
他帶著一群羊.
其中一隻羊.
是一隻黑山羊.
你留意牠的前後腳是綁著的.
所以牠走的時候就卡卡卡.
快要跌倒了.
我們很好奇.
就找個翻譯跟阿拉伯的牧羊人說.
怎麼那麼多羊綁著牠.
只綁著這隻.
牧羊人就經過翻譯跟我們說.
他說這隻羊很麻煩.

$^{281}$整天自己走失.
走失了很多次.
偏行幾路就是這個意思.
整天都不想跟著牧羊人.
整天自己就走到不知去哪.
然後又要走去找牠.
結果牧羊人受不了.
就把兩隻腳綁起來.
就算偏行幾路.
就算自己走失了.
也不會走那麼快.
可以找牠回來.
死過翻生.
引導我走二路.
是這個意思.
上帝在我們生命裡.
可能陳設了不同的結構.
或者陳設了不同的.
你叫做左頭左勢的東西.
有沒有想過這個圖畫呢.
引導我走二路.
引導我走一條對的道路.
其實我們看了這麼多節.
兩三節.
走曠野.
曠野裡面有食物.
有水喝.
最後會去到目的地.
其實這個舊約裡面.
經常出現的故事.
整個舊約的大故事.
The Big Story.
就是講這件事.
因為神在舊約裡面.
在歷史裡面.
做到最大的事.
就是救以色列人出埃及.
然後他們就經過曠野.
最後才到目的地.
就在曠野裡面.

$^{321}$給他們水喝 給他們草吃.
人當然不吃草.
但意思是給他們飲食.
去供應他們一切的需要.
引導他們走二路.
其實以色列人.
出埃及 去迦南地.
其實就算古代的時候.
沒有馬騎.
就這樣走.
馬甚麼都會到.
就這樣行走.
但為甚麼用了四十年呢.
原來上帝有另外一個意思.
就是為自己的名引導他們走.
不是為你的名.
不是為了你的名聲.
而是要人在你身上.
在我身上.
看到上帝的名聲.
所以他們這樣大意色的人.
走四十年.
走曠野.
去訓練他們.
去操練他們.
要他們知道.
他們唯一一個的神.
是上帝.
唯一一個可以靠得住的 應該是上帝.
為自己的名 引導我走二路.
詩篇23年將整個的大故事.
就濃縮在一個人.
自己 以色列人.
個人的生命裡.
你想想 你個人的生命都是一樣.
上帝要我們走一條路.
可能不一定.
是最快的那條路.
不一定是最順利的那條路.
不過.

$^{361}$如果是為神的名引導.
你走那條路.
那條就是正路.
就是二路.
第四節.
我雖然行過死刃的幽谷.
也不怕遭害.
因為你與我同在.
你的將 你的肝.
都安慰我.
死刃的幽谷.
芝加哥看不到的.
芝加哥是一片.
大的平地.
巴斯坦地 以色列就看到.
山多平地少.
很多的山谷.
死刃的幽谷.
就是一些黑到不得了.
原文意思是absolute darkness.
我不知道你記不記得.
你曾經去過的地方.
是伸手不見五指的.
好不好.
我想起這段經文的時候.
我想起以前小學的時候.
小學的時候 在香港長大.
去一個.
Field trip.
學校去Field trip.
去一個島 叫長洲.
那裡有個洞.
叫張寶仔洞.
張寶仔洞.
是以前一個盜賊的.
藏身之處.
第一次進去的時候.
都可以 這樣滑進去.
肥一點進不去.
因為真的進去的時候.

$^{401}$我第一次這麼大個人.
伸手不見五指.
完全自己.
Lost.
不知道方向在哪裡.
死刃的幽谷就是這個意思.
伸手不見五指.
其實最重要的.
不是迷失.
幽谷最.
最考人的是什麼呢.
因為你經過一些迂迴曲折的.
山路的時候.
幽谷最大的威脅就是.
你不知道.
What is going on around the corner.
轉右或轉左.
看不到的時候 轉過去.
到底會有什麼呢.
是不是有隻狼等著你 或者有隻豹等著你.
對一隻羊來說.
死刃的幽谷就是一個.
會很害怕的地方.
就是一個這樣的意境.
他相信上帝不單止.
會公認他.
走過伸手不見五指的行路.
下一個彎是怎樣.
不知道 不過跟著那個.
的牧者.
他就不需要怕遭寒.
也說得很現實.
是一定會經過幽谷的.
人生就是這樣.
傳道書志偉文學告訴我們.
生有時死有時.
哀痛有時喜樂有時.
人生一定會高呼.
不同的境況.
只不過有時.

$^{441}$很obsessed by我們的happiness.
我們很.
著迷於我們現在有的快樂.
我們現在快樂就好了.
上帝永遠都.
在一個高山裡面.
帶我們就好了.
死刃的幽谷.
我一定會經過.
詩篇也可以.
不單止是安慰.
不單止是信心的詩篇.
也是一個智慧的文學.
智慧文學告訴我們.
在上帝的眼中我們應該怎樣看人生.
有時我們.
看人生.
我們很truncated 展開了.
有些現實.
我們不想面對的時候就不看.
死刃的幽谷就是一個很現實的.
一個智慧的文學.
告訴我們生有時.
死亦有時.
有時可能面臨死亡的時候.
面臨一些絕症的時候.
我們先會想到.
詩篇23篇或是想到愛和神的關係.
永恆的事.
聽說.
Steve Jobs臨死的時候.
他一番話說什麼呢.
很長的一番說話.
他說.
他人生去到最高峰.
是不是.
現在他有絕症了.
他其中一句說話.
最令我留意到的人.
他人生最高峰的時候.

$^{481}$現在我有癌症.
我快要面臨死亡.
在病床上.
在醫院上.
我才明白.
可能人吃飯.
吃一餐300元的.
和你吃一餐3元的.
其實都是一樣.
沒有分別.
智慧文學就是想教你這件事.
不是.
牧者特意想害你.
而是牧者想鍛鍊你.
想告訴你.
走他的路才是對的.
他說你的象你的肝都安慰我.
這是一個指導.
不需要去解釋.
很直接.
因為是詩.
不需要解釋很直接.
說一定有兩支棒.
沒有兩支棒就不是牧者.
你看到這個牧羊人.
他沒有棒.
他手裡有什麼?.
石頭.
為什麼拿著石頭?.
我看他拿著石頭.
扔了,然後再撿一塊.
他沒有棒.
原來.
帶著群羊的時候.
你又沒有牧羊狗.
又沒有車,沒有馬.
羊在偏行幾路.
牧羊群羊遠了的時候.
自己走到那邊.
要在崖邊吃草.

$^{521}$他就用石頭再扔遠一點.
然後擲地有聲的時候.
那隻羊就哭了.
然後他又回來.
不是用石頭扔那隻羊.
是扔羊遠一點的地方.
讓牠回到牠的家.
上主對人生的安排.
就是要我們這樣做.
你與我同在.
你的將領的國安都安穩.
那種同在.
With me, Emmanuel.
以馬來尼就是這個意思.
無論經過什麼地方.
他都和我們一起.
最近打世界盃.
有一個很令我們感動的場面.
就是南韓隊波輸了.
球員進到更衣室的時候.
他看到南韓的總統.
來拍拍他.
安慰他.
然後他就哭了.
流淚而淚.
那種同在.
就是這個意思.
你的同在都安慰我.
上帝就是經歷.
和我們一起經歷.
和我們同行.
在我們人生裡不同的高山.
低谷特別.
就是去拍拍我們.
第五字.
在我敵人面前你為我擺設現場.
你用油高了我的頭.
洗我的福杯滿人.
這裡就轉了過來.
不是牧羊人.

$^{561}$而是主人家.
阿拉伯人都很好客.
每逢有陌生人來的時候.
他就當你是朋友.
就帶著牌子.
上面有兩張.
左邊就是阿拉伯人.
或者巴勒斯坦人.
他們很流行煮一些東西.
就是雞飯.
磨一些麵粉.
然後就包起.
就像Taco那樣.
然後就放一些飯上去.
然後就放其他的染料.
就放上去.
煮好的時候就拿出來.
整群人就在帳幕裡.
在沙漠裡的帳幕.
在曠野裡的帳幕.
就圍著我們在吃.
用手吃的.
不要用錯手.
要捏著.
捏成一團團的.
然後就放進口裡.
就像做壽司一樣.
坐在敵人面前.
你為我擺設筵席.
可能意境有一點.
可以想像到.
帶羊到曠野.
吃完草喝完水.
回到家裡的時候.
他就為他們擺設筵席.
說不定在放羊的那幾個小時.
八個十個小時裡面.
可能有獅子老虎.
但是那個牧羊人.
就把這些猛獸打敗了.

$^{601}$然後就覺得.
對了,有野味.
就帶他們回來.
在我敵人面前.
你為我擺設筵席.
就是把敵人吊起的時候.
敵人對你無可奈何的時候.
因為你跟著這個牧者.
你有安全.
你就可以享受那個筵席.
你擺設筵席.
你看到這些東西好不好吃.
當然不好吃.
對我們來說的口味當然不好吃.
但是你要想想.
在那種意境的氣氛裡面吃.
你就覺得真的是大排筵席.
因為外面是沙漠.
吃風沙就有.
但是當你回到帳幕的時候.
你就發現.
這個真的是滿日的.
筵席的豐富.
你用油膏撩我的頭.
那些羊,你記得.
撩開那些石頭.
然後才吃到釘到的草.
他們走過放牧的時間.
他們會受傷的.
面會被.
鼻子,眼睛會擦傷.
或者那些草.
草裡面草.
你撩不動的.
你捐頭的時候.
整個面都會擦傷.
擦傷.
油膏撩我的頭.
橄欖油.
是醫治.

$^{641}$那些羊去體貼的.
來醫治我們.
使我們能夠痊癒.
為我們去療傷.
福杯滿人.
喝的東西.
那個杯子.
你去喝茶的時候.
那個杯子是滿到.
水面都會凸出來.
什麼時候.
有兩個可能性.
第一個可能性,你和人喝茶的時候.
那個杯子是滿到爆的時候.
第一.
那個主人家不停地倒茶給你喝.
喝一點點他又倒給你.
又倒給你.
已經倒不滿了,再倒.
這是第一個可能性.
什麼時候會是福杯滿人.
就是你不口渴的時候.
是不是.
還想著可能要上廁所.
他還倒茶給你.
就是說你不再渴不再餓的時候.
這個就是這樣的情況.
那天羊回到主人家裡面的時候.
那個福杯滿人的情景.
我一生一世必有恩惠.
慈愛隨著我.
我且要住在耶和華的天空.
直到永遠.
恩惠慈愛.
是神和人性約中間的.
總結.
就是說你在上帝的約裡面.
上帝既然是守約既然是信實.
如果你也聽他.
說你不背約的話.

$^{681}$你一生守約的人.
就會有恩惠慈愛隨著你.
這個就是這樣的意思.
他說隨著我.
隨字.
跟著的意思.
其實也很特別的.
因為在希伯來的聖經裡面.
用到隨著這個字.
通常都是負面的.
就是說那些禍患會.
跟著你.
如果你背約的話.
如果你離開上帝的話.
那些咒詛.
申命記經常說.
那些咒詛就會隨著你.
就是說會追趕你會追殺你.
這樣的意思.
你將這件事逆轉過來.
如果他跟著上帝.
跟著這個牧人.
就算都會變成咒詛.
那些咒詛會.
好像仇人那樣追殺你.
就是這樣的意思.
如果你真的守約.
你明白牧者是供應你.
一切所需要的東西.
不是想要和不想要的東西.
會帶你走一條易路的時候.
你又繼續跟著他的話.
不偏左右.
那些咒詛就會追趕你.
好像仇人那樣追殺你.
這個就是一個這樣的意境.
或者好像今天這樣.
你晚上回家.
特別是Chinatown.
治安不是很好.

$^{721}$你回家的時候下車.
然後你去門口.
你一路走的時候.
後面有一腳步聲跟著你.
你越走得快的時候.
他就越走得再快一點.
你準備開門的時候.
他就拿著頭過來.
原來是你老爸.
他追趕你因為他想保護你.
一生一世.
必有恩惠之愛除了我.
就是這樣的意思.
我請你住在夜晚的殿中.
直到永遠.
希伯來文是用字母串出來的.
很有趣的.
一個字母會變化出不同的意思.
住在就是一個這樣的字.
如果差一個字母.
都可以變化出來.
同樣串法.
但原本的串法.
是有些不同.
變化出來是一樣的.
但本身是兩個字.
兩個字的意思.
是有所不同.
但可以是相輔相成.
住在就是一個這樣的字.
另一個.
變化出來的字.
就是回歸.
其實這個更加貼切.
用放羊的比喻.
就是我且要回到.
耶和華的家.
放羊放了十幾個小時.
晚上回來羊棚的時候.
牧羊人打開羊圈.

$^{761}$歡迎羊回來.
為牠們擺飾.
用油膏.
來頭.
其實這個意境.
就是詩篇23篇.
告訴我們.
每一個守羊的.
我們都能夠經歷.
耶和華這個字.
又出現了.
第一節出現了.
現在最後一節.
又出現了.
就是說牠是有頭有尾的上帝.
創始成中的耶穌.
牧羊人不會.
帶羊去抗疫.
然後扔下牠們.
那是故宮.
不會理羊的生死.
但牧羊人在羊圈.
不會理羊的生死.
但牧羊人就不會.
帶了牠們又會帶回牠們.
歸回耶和華的殿.
回到牠們的目的地.
這就是舊約的大故事.
神大耶穌的人就是這樣.
經過抗疫去迦南地.
問題就是.
到底牠們肯不肯跟神.
如果不肯跟的話.
牠們就自尋死路.
肯跟的牠們就會.
回到耶和華的家.
肯跟的話.
肯信服的話.
肯信任到底的話.
這個神同樣我們經歷到的.

$^{801}$就是由頭到尾的恩典.
一生一世必有恩為前為主.
我們教牧.
探一個弟兄.
探了好幾次.
過去病了好多月.
特別是吳牧師.
探了他好多次.
很多個月.
去ICU,去Hospice.
好慘,真是死任的憂愁.
為祈禱.
安慰.
弟兄說得很好.
他說在病床上.
依靠他的安慰.
是上帝的說話.
想起上帝的說話.
梁先生今天坐在這裡.
坐在後面.
走過死任的憂愁.
但不需要.
懼怕,因為有上帝.
跟我們同在.
我們再讀多一次.
詩篇13篇,然後我們唱回應詩歌.
讓這些.
詩篇神的話語.
成為你現在縱使不安的時候.
成為你的鼓勵和安慰.
一起.
(唸詩).
我必不自拒佛.
祂使我躺在青草地上.
令我在寒安寂的水邊.
祂使我的靈魂蘇醒.
為自己的命.
引導我走義路.
我雖然行過死任的誘局.
也不怕遭害.

$^{841}$因為你與我同在.
你的像,你的光.
都安慰我.
在我的人面前.
你為我擺設現證.
你擁有高了我的頭.
使我的福輩滿人.
我一生一世.
必有恩惠慈愛隨著我.
我且要住在.
耶和華的殿中.
直到永遠.
(演唱).
耶和華是我的不折.
我必不自拒佛.
祂使我躺在青草地上.
令我在寒安寂的水邊.
祂使我的靈魂蘇醒.
為自己的命.
引導我走義路.
耶和華是我的不折.
我必不自拒佛.
祂使我躺在青草地上.
令我在寒安寂的水邊.
祂使我的靈魂蘇醒.
為自己的命.
引導我走義路.
我雖然試過失憶的勇敢.
也不肯招牌.
因為你與我同在.
你的槍 你的槓.
都安慰我.
在我敵人的面前.
你為我扮演戲.
你擁有了我的囚.
是我的 是我的.
腐敗滿衣.
我一生 一世.
你為誰來 誰則我.
我這樣 輸在.

$^{881}$耶和華的點中.
我一生 一世.
你為誰來 誰則我.
我這樣 輸在.
耶和華的點中.
我們跟隨你的時候 必定不知缺乏.
只給我們這個信心 只給我們這個信服.
無論我們生命中現在遭遇的.
可能是不安 可能是渴水的時候.
我們深信跟隨你的時候.
我們必不知缺乏.
為了我們人生所經歷的.
儘管我們信心的說一句.
因為你帶領我們走的路.
是一條易路.
教導我們如何信靠你.
而且要永永遠遠的.
回到你的家裡.
的歡迎與引領.
\newpage



\section{以賽亞書 14:1-4}
\label{sec:nzhjFy_6Umo}
\textbf{3/24/2019 主日講道 - 以賽亞書14\_1-4 (Tony Wong 牧師)}
\newline
\newline
連結: \href{https://youtube.com/watch?v=nzhjFy-6Umo}{\texttt{https://youtube.com/watch?v=nzhjFy-6Umo}} ~~~~ 語音日期: 2019-07-26
\newline
\newline
\hyperref[sec:y_ka_wkoG0A]{\small{< < < PREV SERMON < < <}}
~
\hyperref[sec:index_chronic]{\small{[返順時目]}}
~
\hyperref[sec:index_scriptual]{\small{[返順卷目]}}
~
\hyperref[sec:dNN3GBdDj2k]{\small{> > > NEXT SERMON > > >}}
\newline
\newline
以賽亞書 14:1-4
\newline
\begin{longtable}{cl}
\hline
\hline
章節 & 經文 (和合本修訂版)\\
\hline
14:1 & \begin{tabularx}{0.7\textwidth}{X} 耶和華要憐憫雅各，再度揀選以色列，將他們安頓在本地。寄居的必與他們聯合，加入雅各家。 \end{tabularx} \\ \\ \relax
14:2 & \begin{tabularx}{0.7\textwidth}{X} 外邦人要將他們帶回本地。以色列家必在耶和華的地上得外邦人為僕婢，也要擄掠先前擄掠他們的，轄制先前欺壓他們的。 \end{tabularx} \\ \\ \relax
14:3 & \begin{tabularx}{0.7\textwidth}{X} 當耶和華使你得享安息，脫離愁苦、煩惱，和被迫做苦工的日子， \end{tabularx} \\ \\ \relax
14:4 & \begin{tabularx}{0.7\textwidth}{X} 你必唱這詩歌嘲諷巴比倫王說：「欺壓人的竟然滅亡！他的兇暴竟然止息！ \end{tabularx} \\ \\
[1ex]
\hline
\hline
\end{longtable}
$^{1}$.
我們請Pastor Tony去歡迎他.
.
其實想一想,問一問自己,又問一下在座當中的朋友和弟兄姊妹.
有哪一個人喜歡憂傷痛苦,常常緊貼自己的生命當中?.
又有哪一個人很喜歡罪惡所帶來的痛苦,常常在我們眼前出現?.
又有哪一個人很希望看到所身處的社會和環境,常常受到強權的壓制,眼見邪惡在我們當中出現?.
慈愛和公義漸漸消失在人群當中.
我相信這些問題的答案,你們和我都知道.
沒有人希望看到不好的事情,邪惡的事情,傷心的事情,常常在我們眼前出現.
人在罪惡當中被捆綁,社群被罪惡的勢力所欺壓的時候,人怎樣可以得到釋放?.
什麼可以使我們得到安息,得到釋放呢?.
如果可以的話,請大家一起打開聖經,看回我們今天所讀的經文,就是以賽亞書第十四章一至四節.
這段經文其實就是以賽亞書的十三章到二十四章,一個很大段落的其中一段.
如果整段這十章的大段落的經文,其實就是以賽亞先知向著昔日十個國家,圍繞著以色列的十個國家.
所發出要受審判的預言,而今天的經文是上帝要向巴比倫發出審判之前.
對以色列的國民所說的一番安慰的說話,是安慰的說話來的.
為什麼要安慰呢?因為以色列人將要被巴比倫人和亞述人奴許.
讓我們透過今天的經文,看看究竟哪一個是可以給到真正的安息和釋放.
首先我們先留意經文的第一節和第二節.
當我們看這兩節經文的時候,我們發現有關於上帝和他的子民,以色列民,一些安慰的訊息,安慰的預言.
在這個預言當中,耶和華要像母親一樣用濃情厚愛來連戌亞國.
連戌這個字是有母親一樣的濃情厚愛的意思.
耶和華再一次揀選他們,將他們安置在昔日他們離開的土地.
這表明昔日,如果你熟悉舊約聖經的話,以色列民原本是一個民族,後來分為兩個國家,南與北.
之後一個國家,北面先被亞述滅了.
然後在這個經文的背景,預言著,巴比倫要滅掉南面的猶太國.
在這個預言中,他告訴他們,這個昔日分裂的國家將來不單止要再連合.
而且他們能夠被釋放在這個地方當中,回到上帝賜給他們的一個流奶與麥之地.
因著上帝這份憐憫的愛,臨到以色列人身上的時候,就安置在他們當中.
而當地所居住的人,因為很羨慕上帝與以色列百姓的關係.
以至在經文中有一句話,是緊貼雅各家,是願意與以色列人一起定居.
還有看第二節,在第二節的時候,我們發現這個預言不單止是外邦人.
不單止將以色列人帶回到當地,而且歸回之後出現了一個很奇怪的現象.
在第二節說到,以色列人會得著外邦人成為他們的瀑被.
要將再一次擄掠他們的人,還要核製以前欺壓過他們的人.
是不是真的這樣呢?其實有些學者說,這段經文其實是用一個象徵的意思.
象徵就是好似,象徵的意思,他想表達一樣東西,就是列邦的驕傲要消逝.
上帝的榮耀卻是無止境,作為上帝的子民.
他們在那一刻,本來將要受到世界其他列國的核製,他們將要成為世上的王國.

$^{41}$昔日他們是被擄掠,被欺壓的民,他朝他們要成為擄掠欺壓者的統治者.
想想,我剛才說過,這一群猶大國的人將要面對巴比倫的帝國,要令他們國破家亡.
他們將會成為奴僕,將要被一個強國擄到他方,這班以色列人來說.
剛才我們說這兩節經文,是一個安慰,是一個將來等候的安慰.
但同一時間,對於要擄掠他們那班巴比倫帝國的勢力來說.
亦都告訴那班人聽,這個喪鐘的訊息.
表明,巴比倫,巴比倫,你不是一個無堅不摧的強國.
上帝在這裡,其實告訴當時的以色列人聽,他憐憫和釋放屬他的子民.
還有,上帝不單止憐憫和釋放他的子民.
如果你再看經文第三第四節的時候.
你就會發現,上帝使他們能夠脫離受苦煩惱,和被別人強加在他們身上的苦功.
而且能夠經歷到從他所帶來的全然安息.
即是說,他原本在第三節說,他受很多的壓抑.
但都要消逝,成為安息.
反過來,在第四節,經文那裡.
上帝輪到這個權傾中東的巴比倫帝國.
巴比倫帝國在哪裡呢?.
大家知道現在在什麼位置嗎?.
先考考大家.
伊拉克那裡,沒錯了.
現在伊拉克的位置,就是昔日巴比倫強國的大約位置.
這個昔日權傾中東的帝國.
他以為自己無堅不摧.
他又欺壓著很多小小國家的強國.
卻是要被熄滅.
他所作的強暴要被止息.
這個表明上帝臨到的日子.
他一方面使到卑微受欺壓的人得安息.
同時他也要伸手攻擊那些恃著強勢的惡人.
上帝必會出手干預惡人的行為.
實際上那個藐視他的國.
的確,上帝使得屬他的子民得安息.
而且同時會止息世上的邪惡.
你留意第十三章和第二十三章整頓經文.
在第十三章開端的時候說明.
耶和華向這十個國家發出宣告是一個原因.
因為他們藐視耶和華.
說到這裡,如果你看聖經的歷史發展的時候.
上帝是很守信的.
一方面他將以色列人.

$^{81}$他被擄去巴比倫七十年之後.
透過將巴比倫打敗的波斯王.
將他們釋放出來.
而且將他們帶回巴勒斯坦地.
如果你讀聖經的話.
你有時間可以去看以色列的《以斯拉記》.
還有《尼希米記》.
已經有提及到.
將他們釋放.
但問題是釋放了是建立了一個以色列國.
但這些以色列民仍然活在波斯帝國的權勢下.
以色列國好像沒有成為世上強國.
仍然活在很多強國之下.
你在這裡看歷史發覺也是.
那是不是聖經錯了呢?.
是不是聖經搞錯了呢?.
不是.
有些學者提到.
可能巴比倫的傾倒.
其實是上帝要來臨的日子的一個象徵.
是要體驗這個事實的一個象徵圖畫.
要知道先知書.
先知站在這裡.
將來的世界在這裡.
他看在這裡遠遠地看著將來的世界.
他未必看得很清楚.
他看得到輪廓.
就好像如果.
我行山的時候.
在行山當中.
我站在這裡.
看不看得到?.
如果在我的角度.
我覺得好像是一個山.
因為你當你的手是山.
但你看到有多少個山?.
兩個的,是不是?.
我從我的視野看得到大約只有一個.
先知那時候.
他朦朧地看到一點點東西.

$^{121}$我們有著新舊約聖經.
我們可以看到什麼?.
兩個世代.
看得更加長遠.
如果從這個角度去想的時候.
其實學者提出.
巴比倫是象徵著在世界上.
所有強大暴君的模式.
在世界上各種黑暗勢力.
邪惡的權勢.
邪惡的權勢?.
你是指某個國家?.
是不是?.
哪個強國?.
是,當然它可以包含著國家政權.
但也可以包攬著.
世上一些擁有權力欺壓他人的人.
什麼意思?抽象一點.
說些實際一點的.
大家都打工的.
用一個例子說給大家聽.
舉例在一間公司裡面.
有兩個主管.
主管A.
為人很公正的.
他用上級給他的權力.
好好做他的工作.
用公平公正的態度對待他的下屬.
主管A.
坦白說.
如果你是下屬的.
你也想跟隨主管A.
因為你覺得他公平公正.
做到事.
但在同一間公司裡面.
有一個主管B.
主管B是什麼?.
他同時擁有同樣的權力.
但他濫用那份權力.
只是為自己而送.

$^{161}$甚至迫迫他的下屬.
或者.
我們可能對這種人.
會給很多花名.
我不想在台上說什麼花名.
可能你心中已經想到.
有個人是這樣的.
有些人是這樣的.
我給你安個花名.
你可以叫他巴比倫.
明白我的意思嗎?.
不一定是一個國家.
一個國家有一個權勢.
一個力量.
他濫用了那個權力.
欺壓一些手上沒什麼權力的人.
可以在你公司.
可以在你家裡.
可以在政權.
可以在社會各方的勢力.
世上的確有很多黑暗的勢力.
邪惡的權勢.
他們使我們信主的人.
因為信主的緣故.
而面對壓力.
也有些時候.
他們因著我們信主而迫迫我們.
更甚的就是.
有時這些權勢所做的事情.
不單止只是傷害基督徒.
傷害更多的人.
而他們所行的.
最重要的是什麼?.
是抵擋著上帝捍衛公正.
和慈愛的事.
所以你這樣說.
巴比倫是可以預言著.
抵擋著上帝的權勢.
巴比倫的傾福.
傾倒.

$^{201}$其實是要指出上帝掌管歷史的原則.
就是雖然世上的權力.
權勢.
這些權勢是真實的.
但他們必定要為自己的暴行.
在神面前較仗.
而且上帝最終極.
要處理這些邪惡的方法.
在哪裡出現呢?.
聖經有說.
我剛才跟他說.
兩個角度.
先知站在這裡.
他看到一些東西.
但這個新約有沒有聖經說過一些相似的東西呢?.
有的.
如果你有機會揭開啟示錄.
第十九章.
十九節到二十一節的時候.
聖經這樣說.
我看見那獸和地上的君王.
並他們兩軍都聚集.
要與騎白馬的.
即耶穌基督.
並他的軍隊爭戰.
那獸被擒拿.
那在獸面前曾行歧視.
迷惑受獸恩基.
和拜獸匠之人的假先知.
也與獸同被擒拿.
他們兩個就活活地被冰在燒著流亡的火壺裡.
其餘的人都被騎白馬者口中的劍殺死.
飛了都吃飽他們的肉.
說什麼呢?.
又是象徵語言.
在啟示錄裡.
他想說什麼呢?.
說什麼呢?.
就是說耶穌基督.
是要將敵黨他的邪惡勢力完全消滅.

$^{241}$他要永遠作王.
直到永遠.
這就是巴比倫傾覆所顯現出來的最終事實.
所以先知其實是在看一個很遠的遠像.
很遠的遠像告訴我們.
今天我們可能看到很多邪惡.
但最終消滅.
那你可能說.
黃牧師我們已經說了.
現在10點47分.
你10點半上來.
說了13分鐘.
差不多10幾分鐘的聖經.
那跟我今天的生活有什麼關係呢?.
趁大家還沒打哭鬧之前.
或者睡著之前.
我講一下現代意義好不好?.
經文就這樣說.
有沒有現代意義呢?.
是有現代意義的.
我們要承認.
在我們的眼中.
眼前.
在我們這個世代裡面.
是看到很多巴比倫出現的.
可能你一起來.
你一踏入公司.
你已經看到你的老闆.
巴比倫老闆走過了.
然後看到你的同事.
這個也是巴比倫來的.
救命啊.
好了.
打開新聞.
看到某些人.
某些勢力.
某些政權.
這個也是巴比倫來的.
大家笑.
我相信你明白我說什麼.

$^{281}$來到心領神會.
在世界上.
你和我總會發現.
有這些人存在.
有這樣的權力存在.
我們會體會到罪.
出現的時候.
我講罪了.
和罪有關?有.
因為當我們犯罪的時候.
我們不單止和神的關係破裂這麼簡單.
不止這件事.
我很喜歡十字架.
十字架告訴我們.
如果上帝.
是和人復和的話.
不單止和他復和.
打棟.
打橫那條.
也是和人復和.
我們信主的人.
理應相信耶穌之後.
我們是因著信仰的力量.
和其他人復和.
所以反過來.
為什麼上帝要這樣做.
因為罪的破壞能力.
不單止和神人關係破壞.
也會將人人關係扭曲.
試想想.
有多少家庭的問題.
不是因為人的自私.
人的仇視.
而產生出來的嗎?.
有多少婚姻的不和.
不是因為.
只看到自己而看不到別人.
社會上.
有多少.
不公平不公義的事情.

$^{321}$是因為.
我們覺得.
我有權.
他沒有權.
我用權用到盡.
可以這樣說.
當人與人之間的關係.
被扭曲的時候.
使人會運用.
手上的權力.
來傷害其他相對.
沒有權力的人.
不公平不公義的.
暴力的事情.
會出現.
家裡會不會出現.
會.
社會會不會出現?.
會.
工作環境會不會出現?.
會.
很多地方都會出現.
國家裡面會不會出現?.
都會.
我們不單止.
被罪所影響.
被罪壓制.
失去了自由.
同時.
真是真實的.
有時候我想想.
都很容易會犯了這個錯誤.
同時.
我們每個人都會因著受到罪的影響.
我們都成為了欺壓.
其他手上.
無權的人.
很容易的.
我是主管.
我見這個夥計.

$^{361}$不順眼.
當然不是叫大家.
他做錯事.
不叫他改.
當我們有權的時候.
我們可以做一些事情.
對這個下屬.
給多一些事情來做.
是不勝利的.
明白嗎?.
是發生的.
是真實的.
所以如果你這樣想的話.
這是一個現實.
我們會面對.
但是.
上帝來的目的是什麼?.
他來是要釋放他的子民.
當我們接受了耶穌基督的救恩的時候.
我們不單止因著信.
得到和神之間的關係復和.
十字架.
和神的關係復和.
同時因著.
和神的關係復和.
和人的關係.
開始修和.
復和和修和.
同一時間.
上帝也透過信主的人.
在地上的工作.
使其他人能夠真正體會.
從神而來的憐憫.
安息.
和釋放.
我們透過耶穌基督的救恩.
我們固然能夠完全體會到.
上帝的憐憫.
安息和釋放.
英文有一首詩歌.

$^{401}$Just as I am.
大家如果想聽主樂.
可以聽一聽.
在這首詩歌裡.
我們就是這樣.
可以坦然無懼.
去到上帝面前.
上帝將我們從罪裡釋放.
從我們過往的經歷裡.
釋放出來.
哪管我們是一個.
可以這樣說.
是一個平平白白的普通人.
又或者是一個窮凶極惡的殺人犯.
又或者是一個權傾天下的財主.
當我們來到耶穌基督當中.
我們能夠得釋放.
但是不單止這樣.
上帝也幫助我們.
透過我們使其他人得釋放.
世上的確出現很多巴比倫在我們面前.
可以是一個人.
可以是一個組織.
可以是一個制度.
可以是一個國家.
有時會給基督徒很大的壓力.
甚至希望推倒上帝.
和祂所傳揚的公義和慈愛.
有時這些巴比倫的勢力真的很大.
可能給很多信主的人壓力.
但是上帝必然會使得這些黑暗的勢力消滅.
他們必須要向上帝交仗.
的確上帝作王.
使得我們能夠接受救恩的人.
能夠得釋放.
弟兄姊妹和朋友.
這是什麼意思?.
意思其實不是很難明白.
意思就是說.
當你和我領受了這個救恩的時候.

$^{441}$第一方面.
我們是因著上帝的緣故已經得釋放.
可能你未信主之前.
你有很多的重擔.
背負著很多人生的歷史.
甚至乎可能你現在心底裡.
就算缺了自信耶穌也好.
你心底暗暗底裡面都在問一個問題.
我配上帝愛我嗎?.
上帝愛我們.
上帝是要釋放我們.
從那些重擔.
枷鎖當中.
從罪當中釋放.
但同一時間.
他這個釋放更加讓我們看到.
我們面對這個社會裡面有更多的人.
心靈上被重擔壓抑.
或者面對在一個社會制度上被壓抑的時候.
作為教會在這個地上.
也要做一個釋放的工作.
釋放者就是.
將釋放和好的福音帶給他們.
這個不僅僅.
不僅僅是我們去傳福音這麼簡單.
那個重要的.
一定要做.
但在傳福音的過程.
You have good words.
好的說話.
好的真理.
但是我們都需要good work.
好的關懷.
對社區的工作.
將那個真理.
導成了肉身帶入社區當中.
我自己的教會.
不是一間很大的教會.
我自己的教會.
其實我想.

$^{481}$我想加起來.
你們教會的整體人數.
現在這個數字.
不是很大.
算起來.
最多一百人送中英文課.
你自己想想.
很家庭式.
但當我初初來這間教會服事的時候.
有一樣東西.
我是很感動.
我到現在都很感動.
我們教會.
每隔一個星期.
就會在社區裡面.
去食物銀行拿食物.
帶到我們教會當中.
隔一個星期服事.
我們社區當中四十多個家庭.
因為我們教會的社區環境結構變了.
有很多很有需要.
經濟很困難的人.
再加上現在本城的經濟.
真的一般般.
失業率很高.
真實隔一個星期.
曾經有人問過我.
我在上個星期最後說過.
其實四十多個家庭.
可能最多.
我現在算過.
最多只有一個人回來教會崇拜.
好像不是英文所說的.
cost effective.
好像不是這麼有成本效益.
但我在想.
究竟是成本效益.
還是將釋放的訊息帶出去呢?.
上帝的福音是沒錯.
那裡的人.

$^{521}$未必會回到我的教會.
我知道有些基督徒.
未必會回到我的教會.
但我肯定現在.
心裡面有一個肯定.
就是他走進這間教會.
他覺得有愛.
走進這個教會裡面.
是有關懷.
當然我們怕忘記信耶穌.
我們繼續有些英文課.
弟兄姊妹繼續去做.
但我覺得.
很實際說.
作為一個教會的社區當中.
我們是要將釋放的恩典.
帶到我們社群裡面.
有哪一個人不想被釋放.
我們很多人在罪裡面被捆綁.
有哪一個人是希望.
在社會制度被壓抑下.
不被拯救.
沒有的.
但作為教會.
我們願不願意.
因著領受了這份福音的好處.
領受了這個救恩之後.
因著同神的關係打通了.
而我們去學習.
跟其他人的關係復活.
我來了美國一段時間.
上上落落.
有時跟美國的同學聊天的時候.
發覺原來.
不同種族.
不同文化.
不同社會階層.
都有很多的張力.
加拿大也有.
不過有些東西沒那麼厲害.

$^{561}$但存在.
但問題是作為教會.
我們願不願意.
將這個使人釋放安息的福音.
帶到他們身上.
我再說.
世界上是有巴比倫.
你可以想.
有政權.
某些人.
某個體制.
但問題是.
作為一個被上帝將會.
願意釋放我們的人.
好像那班以色列民.
我們會不會帶著盼望來過生活.
而且將盼望的福音帶給身邊的人.
也將盼望的福音.
成為一個行動.
讓社區的人看到.
每個社區不同.
你不要抄襲我.
我的社區和你的社區不同.
但成為你這個社區.
究竟上帝等你們在這裡.
怎樣將一個釋放的福音.
帶到這個社區當中.
讓我們一起祈禱.
因為你的福音.
就是使人復和的福音.
釋放人從罪裡面的福音.
也將釋放的福音.
不單止在我們心靈當中被釋放.
也透過我們的生活.
透過我們的見證.
在社區社群裡面.
將釋放的福音.
讓人知道.
讓人看到.
讓人經歷感受到.

$^{601}$主要是求你幫助我們.
既然我們領受了這麼好的福音的時候.
已經從罪裡面被釋放.
已經可以回家.
有一個牧人帶領我們的時候.
求你幫助我們.
就引領別人.
去到這個陽圈當中.
得安慰和得釋放.
父神啊 求你幫助我們.
帶領我們.
讓我們的心.
吹向你的使命.
在這個社區裡面.
做好上帝要成為教會.
要做的服事.
多謝你.
隨聽我們禱告.
奉主耶穌基督的名祈求.
阿們.
\newpage



\section{}
\label{sec:dNN3GBdDj2k}
\textbf{7/21/19 主日講道: 懇切需要主耶穌 {Pastor John Yi}}
\newline
\newline
連結: \href{https://youtube.com/watch?v=dNN3GBdDj2k}{\texttt{https://youtube.com/watch?v=dNN3GBdDj2k}} ~~~~ 語音日期: 2019-07-28
\newline
\newline
\hyperref[sec:nzhjFy_6Umo]{\small{< < < PREV SERMON < < <}}
~
\hyperref[sec:index_chronic]{\small{[返順時目]}}
~
\hyperref[sec:index_scriptual]{\small{[返順卷目]}}
~
\hyperref[sec:_5Jd2EPsTDg]{\small{> > > NEXT SERMON > > >}}
\newline
\newline
$^{1}$.
張美上來講道之前.
我會和大家讀一讀今天的聖經.
在路加福音第五章十五到二十六節.
請大家聽.
.
但耶穌的明星越發傳揚出去.
有極多的人聚集來聽道.
也指望醫治他們的病.
耶穌卻退到曠野去禱告.
有一天.
耶穌教訓人有法利賽人和教師.
在旁邊坐著.
他們是從加利利各鄉村和猶太並耶路撒冷來的.
主的能力與耶穌同在.
使他能醫治病人.
有人用玉紙抬著一個毯子.
要抬進去放在耶穌面前.
卻因人多尋不出髮紙抬進去.
就上了房頂.
從眼間把他連玉紙聚到當中.
正在耶穌面前.
耶穌見他們的信心.
就對毯子說.
你的罪赦了.
民事和法利賽人就議論說.
這說前往華的是誰.
除了神以外誰能赦罪呢.
耶穌知道他們所議論的.
就說你們心裡議論的是什麼呢.
話說你的罪赦了.
話說你起來行走哪一樣容易呢.
但要叫你們知道.
人子在地上有赦罪的權柄.
就對毯子說.
我吩咐你.
起來拿起你的玉紙回家去吧.
那人當眾人面前立刻起來.
拿著他所討愛的玉紙回家去.
歸榮耀與神.

$^{41}$眾人都驚奇.
也歸榮耀與神.
並且滿心懼怕.
說我們今天看到非常地事了.
歡迎John.
.
大家知道了.
.
讓我們一起祈禱.
.
今天你的僕人可以聽你的話.
願大家一起傾聽.
讓你向我們說話.
奉主耶穌的名求.
你可以看螢幕上的經文.
也可以打開你的電話.
我們一起看今天的聖經.
今天是路加福音第十五章十五到二十六節.
因為翻譯所以今天的講道會長一點.
所以我肯定今天是John的責任.
今天的講道是短的.
我嘗試把整個講道的內容都剪掉.
我們來專注於神的話.
我們馬上開始.
雖然現在是早期的耶穌的教會.
但他已經成為了這麼著名的.
因為人們都會在他身邊聚集.
每當他進入新的地方.
所以路加福音第十五章說.
人們想聽耶穌教導.
他們也想得到治癒他們的疾病.
路加福音提到.
人們很想聽耶穌教導和得到治癒.
他是這麼著名.
甚至很多宗教領袖都從各地來到這裡.
我們想像一下這個場景.
耶穌正在教導.
當時的群眾當中有法利賽人.
有拉比甚至是有文士.
第十八節說.

$^{81}$當時有很多人來聽耶穌的說話.
他們很想親近耶穌.
但是也有四個人.
他們帶著他們的癱子朋友.
同樣是想親近耶穌.
當我們讀到第十八和第十九節的時候.
我們發覺有一個句子是重覆了兩次.
這可能是一個很好的時間.
我們要記住.
每當我們看到一些文字重覆.
這是一個文學的工具.
用來給予人們重視.
讓我們注意到.
這一點是重要的.
這句話有什麼意思呢.
究竟他想強調什麼呢.
原來在第十八和十九節.
這四位朋友.
他們是很想將他的癱子朋友.
帶到耶穌的面前.
這是他重覆的.
而這四位朋友是很堅決地做的.
在這個故事裡面我們看到緊急性.
我感覺到他們很緊張,很拼命.
讓我問一問你.
今天在你當中.
在你的朋友裡面.
有哪一位你是這麼著緊拼命地.
帶到耶穌的面前呢.
你的朋友可能就是有身體殘疾.
好像今天這個癱子一樣.
可能你的朋友是有感情上.
或者是心理上的一些疾病.
或者你的朋友是毀壞了.
自己的家庭和婚姻.
或者是你的親戚朋友.
他正在面對一些經濟結局.
或者是在他的情感上.
在道德上發生了一些問題.
你有沒有很緊張,很拼命地.

$^{121}$將他們帶到耶穌的面前呢.
這四位朋友.
他們是在你當中.
做了什麼呢.
他們是在你當中.
做了什麼呢.
他們是在你當中.
做了什麼呢.
他們無法在耶穌附近.
找到床.
因為有很多人.
也有同樣的想法.
那四個人做了什麼呢.
他們放棄了嗎.
他們回家了嗎.
不.
他們決定了.
無論如何.
他們都要在耶穌面前.
找到自己的朋友.
所以在第十九節說.
正如第十九節那裡說.
很多人.
沉不出髮紙抬進去.
就上了房頂.
從瓦間把他連玉紙.
聚到當中.
正在耶穌面前.
他們竟然拆散了屋頂.
為的就是要將他這個.
癱子朋友.
聚到耶穌的面前.
你知道當我準備這個訊息的時候.
這個屋頂的毀滅.
對我來說是有點私人的.
最近我們要把屋頂.
換了一個房子.
因為颱風把屋頂損壞了.
總共要15000元.
這算是很多錢吧.

$^{161}$你可以買.
一個便宜的新車.
15000元.
或者一個很好的舊車.
我可以給每個人.
100元.
每人.
15000元.
其實呢.
這15000元.
張華每人派.
100元都夠的.
可惜今天他不會這樣做.
因為他已經用光了你的錢.
這個是他今天的第一點.
第一個重點就是.
當我們拼命地.
為了自己和別人的需要.
向耶穌求助的時候.
其實是充分地表現了.
我們對耶穌真誠的信心.
我們想想.
如果我們很常為我們的朋友祈禱.
為我們的朋友祈禱.
但是我們.
我們很常去幫助朋友.
但是我們沒有向耶穌求助去幫助他.
我們問問自己.
為什麼我們會這樣做.
為什麼我們不會來到耶穌面前.
求祂幫助呢.
是不是我們覺得.
耶穌不能夠幫到我們呢.
我們是不是真的相信.
耶穌會想幫助我們呢.
我們想想.
我們有多少次說過.
我們會祈禱給某人.
然後我們會忘記.
為什麼我們不勉強地.

$^{201}$向耶穌求助.
為別人的物質和身體的需要.
我們是不是覺得.
耶穌可能不在乎我們身體的需要呢.
是不是因為我們覺得.
耶穌不在乎.
不重視這些東西呢.
我們相不相信.
耶穌真正是我們的幫助呢.
因為如果我們相信的時候.
我們真的會向耶穌呼求.
去求祂的幫助.
而這裡有一個重要的部分.
不只是他能夠幫助我們.
我們真的相信.
他想幫助我們嗎.
我想說的很清楚.
聖經說得很清楚.
神是能夠幫助我們.
祂有能力去幫助.
祂也想幫助.
祂想要幫助我們.
祂也想幫助我們.
我很希望在這裡.
大聲清楚地跟大家說.
耶穌是會幫助我們.
祂是想幫助我們.
祂有能力幫助我們.
因為耶穌.
祂就是福音.
祂有拯救的大能.
祂一定會拯救凡相信的人.
在《羅馬書》第1章16字裡說.
祂說.
首先是猶太人.
然後是希臘人.
當然.
也是對中國人.
那些相信的.
耶穌是神的愛.

$^{241}$和祂的恩典.
對於我和你.
的充滿愛和恩典.
祂想拯救我們.
在2彼得3:9.
耶穌說.
主不慢於祂的保證.
如某人說的.
祂不慢.
但祂耐心向你們.
不希望任何人死去.
但希望所有人來求復.
我再說一次.
我們真誠地為耶穌.
求助.
是我們.
非常認真地.
向耶穌求助.
並向其他人求助.
這不是我們在.
在我們在.
在《聖經》裡所讀的.
這個不是我們剛才.
看到的經文所說嗎.
看看第20節那裡說.
耶穌見他們的信心.
因為他們這麼肉緊.
要拆了屋頂.
將他的朋友帶到.
耶穌的面前.
《聖經》說.
耶穌看到他們的信心.
以下是第二個.
我想很快說的要點.
當我們將人.
看得重要過物的時候.
也都表現了我們.
對基督的信心.
我們對基督的信心.
就表現在我們對.

$^{281}$看人重要過物.
當他們拆了屋頂的時候.
告訴我.
他們的朋友.
比屋頂更重要.
告訴我.
他們的朋友的健康.
比屋頂的損害.
更重要.
他們的朋友的救救.
比他們的名譽.
更重要.
他們的朋友的健康.
更值得.
更值得損害.
他們的朋友的健康.
比他們面對冒險.
更重要.
他們必須將他們的朋友.
帶到耶穌面前.
不論什麼價錢.
因為人比物更重要.
當我們表現.
人比物更重要的時候.
我們也在向別人表現.
這也是神的看法.
神也認為.
人比物更重要.
所以不奇怪嗎.
即使是基督徒.
我們也覺得.
我們的物品比人更重要.
是不是很奇怪.
因為我們基督徒.
也時時將事物.
看得比人更重要.
有時甚至是.
我們比我們的家庭重要.
有時我們看物件.
比我們的教會.

$^{321}$大家庭重要.
讓我問大家一個問題.
你的人生.
你的教會.
會不會看得不一樣.
如果人比物更重要.
我們一起想想這個問題.
如果人比你的時間更重要.
你會如何對待他們.
如何不同地對待他們.
如果人比你的錢更重要的時候.
我們想想這個教會.
如果人比你的錢更重要.
你會如何對待他們.
如果人比你的錢更重要的時候.
我們要想想.
如果我們成為教會.
看待人是很重要的時候.
我們每一個星期的早上.
我們這個教會.
是如何不同的呢.
我們的僑港區.
我們周圍的旅社.
又如何看待我們的教會呢.
如果我們看待人.
是這麼重要的時候.
可能性是很多很多的.
回到我們.
在這間屋裡面.
教訓耶穌.
突然間他們聽到屋頂有聲音.
有好像拆屋的聲音.
有些塵掉下來.
一片片瓦片掉下來.
你可以想像到這個情況嗎.
大家推來推去.
為了要大家人身安全.
他們望上屋頂.
有一個大洞.
光照進來.

$^{361}$最後他們看到一張床.
慢慢慢慢地.
落到地上.
耶穌第一句說話.
是說什麼呢.
他有沒有說.
發生了什麼事.
他有沒有說.
誰敢干擾我的講道.
又或者他有沒有說.
你們有沒有事.
在第20節說.
看到他們的信心.
他說.
耶穌沒有說那些.
但是在第20節告訴我們.
耶穌看到他們的信心.
就對譚子說.
是不是很奇怪呢.
耶穌竟然這樣說.
我們可能期望耶穌.
說出我們提到的所有事.
但是說出.
對一個完全不認識的人.
在這個情況下.
你的罪是被赦免的.
這是我第三點.
即使我們可能覺得.
我們知道我們需要什麼.
或我們朋友需要什麼.
即使我們以為.
我們知道自己需要什麼.
但是耶穌更加知道.
我們真正的需要是什麼.
耶穌比我們更加知道.
我們的需要.
這四位朋友.
其實他們是想著.
他們需要身體的醫治.
但是耶穌知道.

$^{401}$這個譚子更加大的需要.
他需要和神修補關係.
我們都是一樣.
沒有一樣比這一樣重要.
我們需要和神修補關係.
和我們的神.
和我們的創造主.
和我們的審判主.
因為聖經說到.
我們大家都犯了罪.
虧缺了神的榮耀.
因為我們大家都是罪人.
朋友,你的罪被赦免了.
你看到一個治癒了.
健康的身體.
有一天會再次感染.
即使一個人.
被治癒了.
也不會再次感染.
他們肯定會體驗到.
一件事.
有一天.
每個人都會死.
他們都會站在神的面前.
那一天.
健康的身體.
對你來說.
都沒有意義了.
健壯的腿.
對死亡者來說.
都沒有意義.
我們的錢.
和我們世界上的財富.
都沒有意義.
即使我們在神面前.
很在乎.
你開什麼車.
他不會在乎.
你的學位和名額.
他不會在乎.

$^{441}$你在這地球上的榮譽.
但你知道我們會怎樣令他感到驚訝嗎?.
他會想知道.
你和我的兒子.
是怎樣的.
他想知道.
他想重視的是.
你怎樣回應我的兒子.
他會問你認不認識我的兒子.
你有沒有接受我這位兒子.
我給你的禮物.
我的兒子認不認識你.
因為父親命令.
只有透過耶穌.
他的兒子.
才能讓罪人.
與神的關係正確.
因此.
知道這病毒的真正需要.
耶穌對他說.
朋友,你的罪.
已經被賦予了.
但耶穌知道.
這病毒的真正需要.
所以他說.
朋友,你的罪已經賦予了.
耶穌也賦予了.
這個賦予.
給我,你和每一位.
這個賦予.
給你的朋友,你的家庭.
每一個人.
你會不斷地.
來面對他.
接受你和每一個人.
真正需要的.
唯一的東西.
是解決我們的罪問題.
這是問題.
讓我們與神分開.

$^{481}$因為這個罪的問題.
使我們與神分離.
而耶穌給我們的解決方法.
不是我們可以用錢買得到.
沒有肥皂或洗潔液.
可以洗淨我們的罪.
也沒有足夠的好行為.
我們可以實行.
沒有任何的好行為.
可以遮蓋我們的罪.
只有耶穌.
耶穌是唯一的解決.
我們所有的罪問題.
而你知道.
保祿教的教師.
他本人是法律的教師.
他最初是恨恨耶穌.
他甚至是迫害了教會.
你知道他說了什麼嗎.
你應該很熟悉保祿.
他自己是一個很出名的法律人.
他很憎恨耶穌.
首先他很憎恨耶穌.
而且迫害了教會.
但後來他怎麼做呢.
保祿竟然這樣說.
他說.
保祿將我們脫離黑暗的權勢.
把我們遷到他愛子的國內.
我們在愛子裡得蒙救贖.
罪過得以赦免.
保祿是一個殺人犯.
當時他最大的喜樂.
就是去殺基督徒.
但他卻說.
我得到赦免.
今天你有沒有任何的罪.
你認為耶穌是不可以赦免的.
沒有的.
沒有罪是不能夠被赦免的.

$^{521}$我想請你和我們一起祈禱.
請大家閉上眼睛.
不要看周圍.
我們不需要看到周圍的人.
這是你和神.
我想問你.
誰在這裡.
渴望了解.
神的美妙赦免.
有人想體驗.
重重的罪.
被脫離他們的肩膀.
有沒有任何一位.
你是願意.
你罪的重擔.
能夠被釋放.
朋友.
你所做的事.
是不得到神的赦免的.
但是你一定需要.
來到耶穌面前.
相信祂.
看到你的信心.
祂會對你說.
朋友,你的罪是被赦免的.
來,相信神.
已經把他的兒子.
賠償了你的罪和我的罪.
請你來.
相信這個大能的神.
已經在猜探他的兒子.
來赦免你和我的罪.
請你來接受.
相信耶穌基督.
得到的好處.
祂會使你.
得到赦免.
使你得到永遠的生命.
當你閉上眼睛.
如果你需要.

$^{561}$有人來為你祈禱.
請你把手舉起來.
有人會來為你祈禱.
請你把手舉起來.
還有沒有其他人.
需要有人為他們祈禱.
還有沒有其他人.
你是想有人為你祈禱的呢.
有時我們不懂說什麼.
我們不知道怎麼說.
但是我們有聖靈的幫助.
即使我們的朋友幫不到我們.
聖靈會幫到.
讓我們來到聖靈的面前.
看到你的信心.
祂會說你的罪.
已經被赦免了.
來接受你真正需要的.
你真的需要耶穌嗎.
我們說「阿們」.
在大街上.
我們說到成為社區教會.
我們是不是.
實在很需要耶穌呢.
讓我們相信耶穌.
讓我們一起祈禱.
首先我向兩位姐妹們祈禱.
在房間後面.
我向你們祈禱.
我知道你們看到他們.
把手舉起來.
而這只是在說.
他們心裡的心意.
這是一個外在的證明.
「是的,主啊,我需要你們」.
我們也會說.
「神啊,我們需要你」.
「我們極之需要你」.
「我們需要來到你的面前」.
「接受你的赦免」.

$^{601}$「我相信耶穌能夠赦免我」.
「我相信耶穌能夠修補我和神的關係」.
「我將我們所有的罪放在你面前」.
「主啊,赦免我和清潔我」.
「讓我更清晰,比雪更白」.
「主啊,取走所有的罪」.
「讓我重新開始」.
「因為我相信耶穌」.
「憑著從神而來的權柄」.
「我想告訴你,如果你相信耶穌的話」.
「朋友,你的罪已經得到赦免」.
「你的罪已經得到赦免」.
「我們禱告」.
「Amen」.
「謝謝,John」.
「再一次和我們分享」.
「我們的願神的話語祝福大家」.
「願大家帶著今天這個赦免的訊息」.
「進入你的生活」.
「也和你身邊的朋友分享」.
「我們今天頌拜到此為止」.
「願神的話祝福大家」.
\newpage



\section{}
\label{sec:_5Jd2EPsTDg}
\textbf{8/5/18 主日講道: 基督徒成長12步--\#9 門訓關係 (陳崇基牧師)}
\newline
\newline
連結: \href{https://youtube.com/watch?v=_5Jd2EPsTDg}{\texttt{https://youtube.com/watch?v=\_5Jd2EPsTDg}} ~~~~ 語音日期: 2019-07-29
\newline
\newline
\hyperref[sec:dNN3GBdDj2k]{\small{< < < PREV SERMON < < <}}
~
\hyperref[sec:index_chronic]{\small{[返順時目]}}
~
\hyperref[sec:index_scriptual]{\small{[返順卷目]}}
~
\hyperref[sec:1yfYPGx_z_4]{\small{> > > NEXT SERMON > > >}}
\newline
\newline
$^{1}$我們每個月講一個專題叫做十二步.
我們說十二方面.
我們在一個基督徒的生命裡面成長.
在我們中間信主的.
你應該知道或者你再一次被提醒.
我們信主之後.
我們的生命會慢慢改變.
基本上有很多東西我們會重新整理.
為什麼呢?.
因為耶穌就好像我們平時開車.
我們信了主.
現在由祂開車.
我們就幫祂開車.
我們坐在旁邊.
輔助而已.
最主要是祂帶我們去一個地方.
或者祂帶我們經歷一些事情.
是祂的心意.
所以這十二步就是我們慢慢慢慢在各方面.
去調整和改變.
你信了主之後到現在.
無論你昨天信也好.
十年前,二十年前,五十年前信也好.
你慢慢會有不同的改變.
不同的更新.
在你人生不同的階段裡面.
神可能有不同的方法來提醒你.
我們的興趣會改變.
我們的嗜好.
我們的價值觀會改變.
我們的優先次序.
都應該會有改變.
所以我們這十二步就是我們希望去整理一下.
我們生命裡面.
走基督徒這條路程裡面.
有什麼需要去整理.
今天我們去到第三十五六七八九.
第九了.
第九就是門訓的關係.
門訓的關係.

$^{41}$如果你看我們海報的話.
你知道不只是基督徒.
這麼簡單.
是spiritual direction.
我們每一個基督徒都不可以做一個孤島.
因為神在我們生命裡面的工作.
是靠著其他人來工作的.
有其他人來做一些人道.
其實你想一下.
一個門訓的關係或者一個人際關係裡面.
有人來影響我或者帶領我的時候.
幾乎我們人生裡面接觸任何一方面都是這樣的.
我們的學業.
有counselor.
在我們身邊給我們指引.
或者他給我們一些經驗的分享.
我們的事業.
無論你是做哪一行都好.
我們都有些所謂的師父.
每樣東西都需要別人來幫助指導.
而且在那個關係裡面.
不單止別人去幫我獲益.
我自己有經驗之後.
我也幫到其他人.
我女兒現在17歲.
她已經升高中第四年.
你知道升高中第四年就很嚴重.
她要找大學.
你知道年輕人想的東西很不實際.
很不實際.
你想讀哪一行.
她沒有想過.
很短視的.
只想我喜歡做什麼.
你知道你喜歡做什麼.
不一定能吃飯.
這個世界很實際.
你寧願吃飯.
然後你吃飽飯.
你可以找一些你喜歡做的事情去做.

$^{81}$比你找一些你喜歡做的事情.
但又吃不到飯好.
很實際的.
你不要覺得基督徒說的話很離地.
他們說的很理想.
不一定.
不是這樣的.
你知道.
爸爸媽媽說的話.
他們不會聽.
這個夏季.
因為放夏季.
他去某個地方做義工.
是一個基督徒的後輩學習程.
他很喜歡做那件事.
他一向對小朋友都很好.
他很有智慧.
很懂得怎樣處理小朋友.
管理他們.
他那個後輩學習程.
有一個所謂導師.
就是他的老闆.
也是一個很好的基督徒.
回到教會.
他跟他聊天.
給他一些提示.
開導他.
他開始明白.
應該怎樣選科.
將來怎樣選擇.
這個就是別人對我們的影響.
有一個所謂門訓的關係.
有一點像.
一個指導.
無論是一些職業的指導.
或者一些人生裡.
我們下一步需要有的指導.
無論是事業也好.
學業也好.
我們都需要其他人來幫助我們.

$^{121}$告訴他們.
去打理教會.
這個就是提摩太前後書.
有一本叫提多書.
寫給他另外一個徒弟叫Titus.
告訴他.
他在他當地的教會.
應該怎樣處理的問題.
中間一定是有來往書信.
只不過你現在有的.
所謂提摩太前後書.
就是他寫的回答.
你叫保羅做師父.
你叫做提摩太做徒弟.
巴拿巴是甚麼呢.
師父的師父是誰.
士工.
他可以叫中師.
一代中師.
葉問就是這樣.
葉問都有徒弟.
他看葉問的電影.
最後都是說.
詠春怎樣發揚光大.
這是重點.
你不理他電影.
這個內容是否真是事實.
但他最重要做的事.
就是有這麼多人學詠春.
最後那一幕都是一班阿公.
Uncle.
年輕人.
年輕女孩.
一起打詠春.
所以你看到巴拿巴.
一代中師.
三代在新約的文件中展現出來.
所以今天想說的.
是巴拿巴的.
新約聖經記載這個人.

$^{161}$他做了甚麼.
以至告訴我們.
我們應該怎樣去走.
我們基督徒的旅程.
巴拿巴我們知道.
他不是十二使徒之一.
但他在《少林寺行傳》出現了好幾次.
都是在一些很關鍵的情況出現.
關鍵的時間出現.
剛剛初期教會在耶路撒冷被建立的時候.
他就出現了.
在這裡怎樣說呢.
Acts 4.
《二使堂傳》第四章.
他說當時耶路撒冷的教會剛剛開始.
真是第一代的直黨運動.
第一代plan的church.
他說許多人都是一心一意的.
沒有人說他的東西.
有一樣是自己的.
到時大家共用.
十二使徒大有能力.
見證主耶穌復活.
眾人也蒙大恩.
內宗也沒有一個缺乏的.
因為人人都將他的田產和房屋都賣了.
把所賣的家銀拿來.
放在使徒的腳前.
照各人所需要的分給各人.
這裡有一句話是很特別的.
內宗沒有一個缺乏.
你翻覽舊約聖經.
差不多只有一處地方看到內宗沒有缺乏.
就是神的子民裡面沒有窮人.
沒有一個needy person.
在哪裡出現呢.
在《申命記》其中一章.
前後大概說到.
耶和華神對他的子民說.
你們進入以色列應許之地的時候.

$^{201}$生活一個社群的時候.
你們需要大家不要欠債等等.
以至你們中間沒有窮人.
沒有一個needy.
整本舊約聖經只有那處地方說到.
就是上帝對他的子民說.
如果你真的愛我盡心來事奉我.
你們應該內宗沒有缺乏.
史唐人傳似乎告訴我們.
初期教會聖靈降臨的時候.
那群信耶穌的人.
他們就taste到.
到底什麼叫做上帝子民的生活.
就是內宗沒有缺乏.
以色列好像在初期教會taste到.
實現到.
有猶太人有第二代講希臘語的猶太人.
也開始有外邦人進來.
沒有一個缺乏的.
就是教會神的子民.
他們有一種分享的精神.
即是每個人不只是靠自己.
看著自己.
原來每一個人都是artist center.
去到照顧其他人需要.
照顧大家的需要.
就在這個時候巴拿巴就出現了.
四章三十六節.
剛剛剛才說完分享之後.
路加在史唐人傳說.
有一個利未人.
他生在塞浦路斯的塞浦路斯.
他的希伯來名字就是約瑟.喬治.
試圖給他一個綽號叫做巴拿巴.
巴拿巴這個綽號翻出來的意思.
大概是encouragement的意思.
即是他很懂得去鼓勵.
為人打氣來鼓勵教會.
他做了什麼呢.
他有田地也買了.

$^{241}$把稼銀拿來放在使徒的腳前.
這裡給了很多資料.
一些所謂的sociological information.
告訴我們這個人是怎樣的.
很多單面的人物是不會介紹這些detail.
巴拿巴就不是.
他說他一個利未人.
即是一個以色列人.
他就在外邦的地方.
異鄉的地方長大.
生在塞浦路斯.
塞浦路斯是東地中海的一個很大的島.
在所謂以色列的以外的地方.
所以他算是第二代.
或者散居周圍的猶太人.
到第一世紀的時候.
他們大概有很多猶太人.
他們就在異地搬回耶路撒冷.
去等候上帝的應許臨到.
上帝的國.
以賽亞那裡說的上帝的國臨到.
到底會怎樣呢.
很多人就在外地回到自己的故鄉.
然後他們有田地.
他自己有田有地.
在耶路撒冷那裡有田有地.
他就賣了.
他信了主.
信了耶穌.
然後他就將這些價錢放在使徒的面前.
得到初期教會很大的鼓勵.
大家知道.
看《真聖經》有上下文.
下文我們不說.
但下文他也說了另外一個奉獻的故事.
巴拿巴他鼓勵初期教會.
以至沒有人是有需要的.
沒有一個窮人在中間.
他賣了他的田地.
然後放在使徒的面前.

$^{281}$下文說了什麼呢.
就是另一對夫婦.
叫做阿拿里亞和薩菲拉.
如果你們有機會回去看《使徒行之第五章》.
是另一個奉獻的故事.
他們做了什麼呢.
他們都是有田地.
賣了田地.
也有價錢拿了.
他們放在使徒奉獻之前.
他們就收起自己一點.
然後才交給使徒奉獻給教會.
大家知道收藏是怎樣.
因為他們是棄用聖靈.
所以死了.
很簡單.
所以兩個奉獻的故事.
做一個對照.
巴拿巴巴納布斯.
就是一個願意不會有回報.
自己有沒有一點點好處也好.
他也願意奉獻出來.
跟下面那對夫婦.
就是阿拿里亞那對夫婦就不同了.
想奉獻.
但又想自己有一點點好處.
有一個很大的對比.
巴拿巴巴納布斯.
有田有地.
按道理說家境應該挺富裕的.
如果耶路撒冷有地的話.
特別今天也是一樣.
如果耶路撒冷有塊地.
是很大的一件事.
可能你的地比他那塊地大.
但你耶路撒冷沒有地.
那你就差一點.
很多猶太人他們在異鄉回到耶路撒冷的時候.
有地為甚麼.
因為有墳地.

$^{321}$他們想重新在耶路撒冷.
重新在耶路撒冷.
到神的國來到的時候.
他們的盼望就是復活.
就是在耶路撒冷復活.
首先復活.
這個是他們的希望.
巴拿巴納布斯的地.
可能包括他的墳地.
就賣出來了.
即是我們今天叫做4-1-K.
即是他人生裡面的安全.
他願意放出來.
還有一樣東西要明白.
當時羅馬的社會.
每次說到奉獻捐獻的人.
都有一種這樣的文化.
我們中國人叫做禮尚往來.
以前的羅馬社會都是這樣的.
我們叫做benefactor.
即是說我是比較有錢一點.
我就捐錢出來給窮人.
幫他們建路也好.
幫他們建劇院也好.
是為了大家的.
是公共建築.
但我建了之後就禮尚往來.
即是我有錢.
我建一些東西給你.
一班窮人的時候.
你這班窮人要為我做甚麼.
就是要你回報我.
是雙方的.
你要為我立個石碑也好.
你要給我建廟也好.
或者你要建一些東西來紀念我也好.
即是我要有名堂.
我要有威望.
這種reciprocal relationship.
即是雙方禮尚往來的文化.

$^{361}$在羅馬的社會裡面特別深.
因為捐獻者和受惠者.
就有一種依附.
我是有錢的.
我捐給你.
你又回報給我.
你又不會造反.
社會就會安定.
你不會亂出來遊行去抗爭.
因為我對你好.
你又對我好.
這個故事在文化情況就完全不同.
因為它不是reciprocal.
它放在使徒的腳前.
即是全盤交給使徒去看管.
不是他自己拿著數.
不是他自己所謂的string attached.
我奉獻了一大筆.
不過我有些條件.
你要給我做執事.
你要給我做上位.
類似這些東西.
其實這些文化經常都有.
當時羅馬的文化更加厲害.
有一種無條件的思慾.
為什麼呢.
因為他有願見.
巴拿巴是看見將來.
教會的將來.
他看見這個初生的教會.
他沒有問.
他五年後,十年後,二十年後.
是不是值得我投資呢.
他看見上帝在他們工作.
他就願意奉獻出來.
這就是他為什麼有個綽號叫encouragement.
這不是因為巴拿巴捐獻的特別大份.
所以使徒行傳或路加夫特別提醒他.
路加夫和使徒行傳都是同一個人寫的.
一二集.

$^{401}$裡面講了很多捐獻的故事.
寡婦捐獻等等.
不是講那個amount.
而是講他的心態.
他的那種體液的角度.
那種價值觀.
是和他不同的.
是被耶穌改變了.
是因為他的心態.
所以特別提到巴拿巴這個人.
我們UV開了兩年.
有一個卡.
或者一個奉獻我收到.
是很感動我的.
兩年以來.
久不久他會寄一張卡給我.
那是一個寡婦寄來的.
她認識我.
在郊區很遠的地方.
我記得她去年寫了一張卡給我.
她說牧師我是一個寡婦.
我給的不可以很多.
不過我盡量給.
她以為很少數目.
但其實我覺得對她來說.
數目很大.
不少錢.
但她說我只是一個窮寡婦.
所以我只能夠奉獻這麼多給你.
去支持你新的教會.
新的工作.
她看到的是什麼呢.
就好像巴拿巴所看到的.
就是建立上帝的工作.
有新的發展.
有神的工作在當中.
她就願意來做一個鼓勵.
不是數目.
而是體液的角度.
她的價值觀.

$^{441}$所以巴拿巴第一次出現的時候.
就是一個捐獻者.
第二次出現.
她就是一個宗師.
怎麼看到呢.
九章第二十六節.
這裡說到.
保羅這個人.
他本身是一個猶太的拉比.
一個很有能力的人.
他去打壓教會.
去迫害教會.
也抓那些信耶穌的.
拉他入監.
然後他去到一個地方.
叫大馬色.
在路上的時候.
主耶穌向他顯現.
他就降伏在主的面前.
他說原來這個才是真正的神.
然後他去到主之後.
他就去到耶路撒冷.
因為當時教會只有耶路撒冷.
在九章二十六節.
他就說到蘇羅或者保羅.
他到了耶路撒冷.
他就想用門徒來結交.
因為那些第一代的信徒.
或者十二使徒.
都是在耶路撒冷.
他們卻怕他.
不信他是門徒.
因為不信保羅是門徒.
為什麼呢.
因為突然間.
為什麼一百八十度改變.
我信不過你.
二十七節唯有.
獨有巴拿巴來接待他.
就領導他去見十二使徒.

$^{481}$然後將他在路上.
看到主的異象.
神對他說的話.
還有大馬色.
如何奉耶穌的名來放膽傳道.
都告訴使徒.
即是說做他擔保的人.
擔保人.
用他的名字.
你知道以前古代的社會.
是一種榮辱的文化.
Honor and shame.
你有的currency.
不一定是你的錢包.
你有的currency.
最重要的就是你的名字.
你有一個有美名的人.
有名譽的人.
你做人擔保.
就是用你的名字來做擔保.
等於我們中間有不少.
是否扮移民.
現在都是要人擔保.
那個人就拐個頭出來.
去擔保你.
是不是.
如果你有什麼事.
你走到這裡.
瞞稅.
然後騙錢.
然後斬死幾個人.
那個人的名字就有問題.
做一個sponsor.
巴拿巴就是做他的sponsor.
他見到這個人.
他居然歸主.
然後見到他在大馬式那裡.
歸主之後.
他就去傳道.
然後被人迫害.

$^{521}$被人追斬.
他照樣去傳道.
他知道這個是真的.
他唯有巴拿巴親自去接待他.
以他的姓名.
以他的名聲.
他擺上.
然後擔保這個人.
他說這個保羅是真的.
是可信的.
然後就帶他出道.
十一章.
巴拿巴又在關鍵的時候出現.
十一章這裡發生什麼事呢.
上個星期.
我記得是前個星期還是上個星期.
前個星期.
我說過安提拉教會.
在巴拿巴出現一個很關鍵的時候.
他說這個巴拿巴原本是一個好人.
被聖靈充滿.
大有信心.
並且有很多人歸服了.
他也帶掃來找掃羅.
找到了.
他就帶他去安提拉.
然後兩個人.
巴拿巴和保羅就足有一年的功夫.
很實在的在那裡和他一起聚會.
去教訓教會裡面很多人.
基督徒是從安提拉開始的.
這裡怎麼可以看到他是一個宗師.
一個道師呢.
就是他特意去找他.
在安提拉教會.
有猶太人信.
有外邦人信.
巴拿巴大概可以自己一個人做.
但是他特意去找一個新人保羅.
我們知道保羅的時間線.

$^{561}$他大馬式信主.
然後去耶路撒冷之後.
見了使徒之後.
他就回到自己的家鄉大掃城.
他應該去那裡做什麼呢.
他應該去處理一下他的神學.
到底這個信仰是什麼呢.
還有在大掃城附近的地區傳道.
巴拿巴就特意去找掃羅.
找這個字.
是很特別的字.
因為路加夫曾經用過這個字.
就是你記得十二歲耶穌的時候.
耶穌上聖殿.
然後回去家的途中.
父母不見了他.
然後四處找他.
那個找字就是這個找字.
就是這個找字.
不是我去請人.
我想請一個牧師.
就當是廣告.
誰來者不拒.
我就去見一見他.
不是的,他是特意去找他.
等於你不見了錢.
不見了錢包.
四處找他.
巴拿巴就是這樣.
離開安提安.
然後去到大掃城.
他知道保羅在那裡.
但是他逐家逐戶去找他.
不知道他去了哪裡.
因為他知道這個人是可用之材.
四處找.
四處找,發散的人去找.
拼命地找.
等於像是有些竹澤.
或者有些失物.

$^{601}$認命那樣.
特意去找保羅.
因為他賞識這個人.
他知道這個人是一腳踏三船的.
兩個禮拜前.
你記得我說過.
保羅.
為什麼一定要找保羅去安提安.
因為他是有三個文化背景的人.
猶太的拉比.
希臘哲學思想.
和羅馬的公民.
所以他什麼都有.
他最重要的是教會的將來.
他看的阿根廷都不是他自己的地位.
聲望.
而是下一代.
下一代可不可以被建立出來.
這個和猶太的傳統文化.
或者甚至我們傳統的華人文化很不同.
我們傳統華人文化就是.
教會的徒弟沒有師父.
所以他教徒弟的時候.
師父留下兩招絕招.
讓徒弟作反的時候可以打死他.
華人文化就是這樣.
猶太人的傳統都是一樣.
名字.
榮譽.
他要保持著榮譽.
其實這些文化.
教會也不少有.
教會有很多牧者很忌賽.
真的.
我過來告訴你.
我們會忌賽的.
咦 他說得這麼厲害.
不要讓他說那麼多.
一年說幾次就好了.
為什麼.

$^{641}$他這樣他霸佔了我的講道.
我不用說的.
其實為什麼不好呢.
是不是.
講少一點不好.
你有空的話.
或者做其他更重要的東西.
我們很忌賽.
或者我們不想交給下一代.
總是想留下幾招.
我們有一種戒心.
有一種妒忌.
有一種文化.
巴拿巴就沒有了.
你看到小男傳的記載.
從這裡開始.
都是說巴拿巴和索羅.
巴拿巴和保羅.
巴拿巴的名字是走線.
然後是保羅.
這個很重要.
因為到之後那幾章之後.
就是保羅和巴拿巴.
或者是保羅自己.
甚至去到一些外邦人的地方.
那些人的反應.
對保羅和巴拿巴的傳統反應.
就是覺得保羅是說話多一點.
巴拿巴是脆口的.
所以有一個歐洲的學者.
他寫了一本書.
他只是說巴拿巴.
把他背後的傳統都炒掉.
整理之後.
他說巴拿巴這個人.
是中間人.
他永遠不是走頭的.
是走中間.
因為他是扶持新一代出來.
這個就叫做Discipleship.

$^{681}$想起Discipleship.
那些番內教的.
很多時候我們的經驗就是.
死記東西.
記很多金句.
記很多聖經的東西.
全部都是內容基礎.
但不是關係基礎.
巴拿巴和保羅.
作為宗師.
他就去教導保羅.
Mentor他出來.
大家知道.
如果你看到有隻龜.
在柱子上.
你就很肯定.
那隻龜是有人幫他.
是吧.
有個哲學家說了這句話.
What the heck.
你看看是不是有意思.
如果你看到這樣的圖畫.
Somebody has helped him.
他不會自己爬上去.
有人幫他上去的.
那個就是宗師或者導師.
Mentor.
每個基督徒都需要有個宗師或者導師.
我看到有人皺眉頭.
不太明白.
Turtle on the pole.
Somebody has helped him.
什麼叫導師.
什麼叫coach.
什麼叫Mentor.
很多時候都是同一樣東西.
有些人說是long term.
有些人說是short term.
其實我覺得最重要不是一樣.
我覺得有一個很好的定義.

$^{721}$什麼叫導師.
什麼叫教練.
有一個好的導師.
會對你的好問題提出好答案.
有一個好的coach.
就會對你的答案提出好問題.
想爆你.
Mentor has some great answers for your questions.
Coach.
Some great questions for your answers.
這兩年我們背後.
譬如上星期Johnny來過.
有些所謂的church planting staff.
他們幫我們這些新的教會建立.
我上過他們的課.
我覺得真的學很多東西.
因為學東西的時候.
我們經常覺得.
是不是有考試.
我要記什麼.
那些session.
我上過一個連續的session.
是12個session.
是cover教會12方面的.
他們不是給你很多內容.
不過他們問你問題.
真的有一個好的coach.
懂得去問你問題.
問你怎麼看.
他不是純粹問了問題之後.
你覺得呢.
他不是這樣.
他幫你引導你.
然後去問一些更加好的問題.
或者去深化你的答案.
這些就是disciples裡面應該要做的事.
我每個月都見一個.
所謂淑玲的導師.
spiritual coach.
他都會問我一列的問題.

$^{761}$最近你在聖經裡面有什麼新的insight.
有什麼新的發現.
其實他就不是純粹問你有沒有靈修.
而是你有沒有想起一些新的insight幫到你.
其實他就鼓勵你去study.
最近你的家庭怎樣.
你的子女你老婆怎樣.
他關心你個人.
當你問到這個問題的時候.
當然你平時都有留意.
但他面對面問你的時候.
你就覺得.
我真的要想一想.
就算我隨便回答他.
回到家的途中.
我一個人開車的時候.
我都會想一想.
有些事情可能要整理一下.
要留意一下.
他問你最近有沒有運動.
連這些都問.
我又會想一想.
我有沒有注意我整體的生命.
discipleship就是這個意思.
其實每一個基督徒我們在座的.
我們都需要有導師.
你不要想到導師是那些很老套.
一定要兩個人坐在一起.
要談一些很formal的事情.
不是的.
在一種對話裡面.
有人給你feedback.
有人給你答案.
亦都有人問你問題.
我記得在那些coaching session裡面.
其中一個整天問的問題就是.
我就說我教會遇見這樣的問題.
他又不給我答案.
當然他給不到什麼答案.
答案都是罐頭答案.

$^{801}$但是那個coach又問我.
你試過做什麼.
做過這個做過那個.
你有沒有想過有其他方法.
去tackle這個問題.
有人問你的時候.
你就會想.
是喔.
再回到開頭我所說的.
人生裡面每一個方面其實都是這樣.
我們都需要counseling coach.
你想一下你的職業.
一個好的coach.
為什麼你今天事業有成.
都是應該有一個coach去幫你的.
無私的奉獻去幫你.
屬靈心靈其實都是一樣.
今天開始.
我想你祈禱.
你說我沒有一個mentor沒有一個coach.
Well, be open.
找一個導師.
團體的導師.
小組裡面的組長.
跟他傾談對話.
然後在對話的關係裡面.
你才會成長.
是relationship來的.
神就是藉著其他relationship來訴訟我們.
你為了這個人去祈禱.
你自己是不是可以做一個更加好的導師.
更加好的教練.
怎樣去適應問題.
能夠引導其他人去成長.
我們低頭幾度.
主要我們多謝你.
因為我們生命當中有不同的人來祝福我們.
你使用這些弟兄姊妹.
導師,教練,mentor,coach.
讓我們開竅.

$^{841}$讓我們看到人生裡面更多不同的盲點.
就是為了這些人來獻上感恩.
求你讓我們能夠成為其他人的恩師.
其他人的教練,coach.
幫助我們這方面來成長.
幫助我們和其他弟兄姊妹一起到one on one談談的時候.
小組裡面談談的時候.
我們能夠帶領人來成長.
歸榮要給你.
祈禱奉傳明教.
\newpage



\section{}
\label{sec:1yfYPGx_z_4}
\textbf{10/7/18 主日講道: 基督徒成長12步 \#12 植堂擴展 (陳崇基牧師)}
\newline
\newline
連結: \href{https://youtube.com/watch?v=1yfYPGx-z-4}{\texttt{https://youtube.com/watch?v=1yfYPGx-z-4}} ~~~~ 語音日期: 2019-07-30
\newline
\newline
\hyperref[sec:_5Jd2EPsTDg]{\small{< < < PREV SERMON < < <}}
~
\hyperref[sec:index_chronic]{\small{[返順時目]}}
~
\hyperref[sec:index_scriptual]{\small{[返順卷目]}}
~
\hyperref[sec:6IhSTx8c43Q]{\small{> > > NEXT SERMON > > >}}
\newline
\newline
$^{1}$由年頭開始到現在.
我們在講這個系列.
這次是第十二次了.
我覺得這次是最重要的.
因為這次是大結局.
講這個系列的時候.
中間當然有其他的系列.
我們講過了.
但是到了這裡十二個.
我們想提醒你和挑戰你.
我們希望你能夠.
在這次的系列中.
有一個更好的發展.
希望你能夠有一個更好的發展.
希望你和我們.
在我們的所謂.
spiritual journey.
我們跟隨主的旅程.
我們可以去成長.
所以我們這次最後一次講到.
力堂和擴展.
教會需要繼續去成長.
不可以停滯不前.
如果教會要繼續成長.
不可以停滯不前的話.
換句話說.
我們每一個作為教會的弟兄姊妹.
我們就需要.
繼續向前走.
怎樣講呢?.
我們講到力堂和擴展的時候.
可以有很多方法去講.
我們可以講一下最近的趨勢.
最近的人口的變動.
最近新的波動.
新的運動是什麼.
為什麼要有church planting.
這件事出現.
可以這樣講一下.
我們都可以留意這些東西.

$^{41}$但我想最重要的就是.
到底我們有沒有那種spiritual strength.
有沒有那種屬靈的力量.
來推動教會.
在教會事奉.
來跟隨我們的主.
去見證他.
這些全部都是講spiritual strength.
不是講我們的physical strength.
沒錯我們需要.
教會需要.
奉獻需要有資源.
做工作需要有resource.
這個一定是.
但這些全部都是physical resource.
最重要的是.
我們的spiritual strength.
到底我們心裡面有沒有力量.
如果這樣看的時候.
我們發現聖經裡面.
時常提醒我們.
和教訓我們.
到底我們所謂的信主.
那個信心到底是什麼.
到底本質應該是怎樣的呢.
所以這三個禮拜.
我們看了一些人物.
摩西.
他的leadership moments.
他曾經經歷過.
一些怎樣的神的呼召.
以致他能夠.
踏出來事奉.
今天我們也想講這件事.
就是.
那個生命.
在我們人生裡面.
有主的恩典到了.
他在一些moment裡面.
來挑戰我們.

$^{81}$我不知道你心裡面.
有沒有這樣的時機.
可能是一個邀請.
可能是一些危機.
來到的時候.
突然間就令到你.
好像很.
面對.
真正的生命是怎樣的.
面對神.
面對上帝是怎樣的.
有一個人.
他人生裡面大概經歷了.
兩個很重要的.
所謂calling.
兩個很重要的.
defining moments.
在他生命裡面.
信心的路程上.
給他一個很大的挑戰.
今天這個人就是我們想看的.
就是亞伯拉罕.
我們看一段.
聖經記載亞伯拉罕.
第22章.
《創世記》.
他說:這時候神要.
試驗亞伯拉罕對他說.
亞伯拉罕.
他回答我在這裡.
神說:帶著你的兒子.
就是你所愛的.
獨生子爾撒.
到摩萊亞地.
在一座山上.
獻為梵濟.
第30章.
亞伯拉罕清早起來.
預備好你的奴.
帶著兩個同僕.

$^{121}$和自己的兒子爾撒.
又劈好你獻梵濟用的柴.
就起程到神指示他的地方去.
第三天.
亞伯拉罕.
遠遠地看見了那地方.
亞伯拉罕對他的同僕說.
你們和奴留在這裡.
我與孩子.
要到那裡.
那邊去敬拜.
然後回到你們這裡.
亞伯拉罕把.
獻梵濟的柴.
放在他兒子爾撒的身上.
自己手裡拿著火.
與刀.
他們二人一起.
往前走的時候.
爾撒問他父親.
亞伯拉罕說.
爸爸.
亞伯拉罕回答.
我兒有什麼事.
爾撒說.
你看火與柴都有了.
可是獻梵濟用的羊羔在哪裡呢.
亞伯拉罕回答.
我兒神自己.
必親自預備.
獻梵濟用的羊羔.
於是二人繼續一起前行.
他們到了神指示他的地方.
亞伯拉罕就在那裡.
捉了一座祭壇.
擺好了柴.
捆綁了自己的兒子爾撒.
就把他放在祭壇的柴上.
亞伯拉罕伸手拿刀.
要宰殺自己的兒子的時候.

$^{161}$耶和華的使者從天上呼叫.
亞伯拉罕說.
我兒子已經死了.
天上呼叫亞伯拉罕說.
亞伯拉罕說.
亞伯拉罕回答我在這裡.
天使說.
不可在這孩子身上下手.
一點也不可害他.
現在我知道.
你是敬畏神的.
因為你沒有留下你的兒子.
你的獨生子不及我.
看見後面有一隻公綿羊.
兩角扣在.
稠密的小樹叢中.
亞伯拉罕就去取了那隻公羊羔.
把他獻為凡濟.
代替自己的兒子.
亞伯拉罕給那地方起名叫耶和華以南.
直到今天.
人還說.
在耶和華的山上.
神必有預備.
這是很難處理的經文.
有人說.
如果你有這樣的經文.
如果你有這樣的權利.
你重新再編輯聖經的話.
把一些聖經刪除.
有人就說.
這一章就把它拿走就好了.
因為有很多人拿著.
這個故事.
來反抗上帝也好.
或者反抗聖經裡說的話也好.
確實是很難去啃下去的.
我們通常的做法就是.
把他靈化.
沒有這樣的事情發生的,怎麼會這樣呢?.

$^{201}$神怎麼會賜兒子給亞伯拉罕呢?.
然後又要了他的兒子的性命.
可能有些背景可以幫助我們.
其實這幾年.
特別多這些背景的爭論.
到底亞伯拉罕的時代.
究竟發生了什麼事?.
究竟亞伯拉罕的時代.
究竟發生了什麼事?.
當時人的宗教的practice是怎麼樣的?.
如果你有留意的話.
英國有好幾個歷史學家.
他們說到底背後發生了什麼事.
然後有其他人就說沒有這樣的事.
爭論了好一段時間.
明白了.
原來當時發生了這樣的事.
看完經文的時候.
首先你可能會問.
神叫亞伯拉罕獻上耳朵.
是不是做梵祭?.
你知道什麼叫做梵祭嗎?.
就是燒,將它燒了.
將它獻為梵祭.
為什麼亞伯拉罕的反應會是.
所謂英文叫做meta effect.
你叫我這樣做,我就這樣做.
好像沒有什麼反抗.
沒有什麼異議.
有沒有背景可以衡量到.
為什麼他會有這樣的反應呢.
正如我剛才所說.
近幾年特別多這樣的update.
發現加林的風俗的習慣.
加林宗教的風俗的習慣.
有一件事叫做獻英祭.
獻英不是飛的那些鷹.
獻人就是human sacrifice.
或者是infant sacrifice.
在北非一個地方.

$^{241}$叫做加泰基.
在舊約的時代.
在這個地方他們挖掘一些東西出來的時候.
發現是一個嬰孩的.
不是嬰孩的墳場.
而是那些嬰孩.
被獻祭之後.
然後獻給神靈.
獻給神靈.
獻給假的神靈.
然後火化之後.
放在一些骨灰罐裡面.
放在地上.
埋葬.
然後上面就立石.
做一個記號.
所以有些碑在那裡.
那些碑很多時候都會說.
這樣獻上.
上帝會繼續祝福我.
繼續去祝福我.
這個就是所謂的cultural logic.
當時的文化邏輯.
文化宗教的邏輯是什麼呢?.
特別是迦南宗教的邏輯.
就是上帝給你一個兒子.
我們中間有不少抱著的.
上帝給你一個兒子.
上帝.
神靈自然就有.
權利去要求.
要求拿回這個兒子.
特別是在生命裡面.
最激.
最激的時候.
或者最危難的時候.
如果你將兒子獻上.
的時候.
有一點像.
跟神靈.

$^{281}$強迫他也好.
賄賂他.
叫他介入.
有一幅畫在埃及的宮廷裡面.
或者是迦南的廟裡面.
就是描述.
那個法老.
攻打迦南的一個城市.
那個城市叫.
阿薩基倫.
你將線條畫出來的時候.
你發現法老.
帶著埃及的軍隊.
打迦南的城市的時候.
那個戰況.
是非常激烈的.
但是還沒分勝負.
為什麼你知道還沒分勝負呢?.
因為城門還沒開.
城門還沒開.
是還不知道誰贏.
和誰輸的.
這個法老當然是.
帶著軍隊壓頂.
來攻打這個城市.
你看到.
城牆上面.
就有城市裡面的王.
還有一堆長老.
他們在做什麼呢?.
他們就舉手.
不是向法老求情.
而是向上天.
那個神靈.
向他來幫助.
右邊.
城牆的右邊.
是吊著一個人下來的.
那個是嬰兒.
你看他的髮色.

$^{321}$就知道有條辮子.
是一個小孩.
左邊也有.
就是緊急的時候.
如果將你的兒子.
獻上給神靈的時候.
就是希望神靈做什麼呢?.
就是去.
加入去幫助你.
事實上在舊約裡面.
有這樣的描述.
在黃記下第三章.
那裡說的是什麼呢?.
就是以色列人打摩訶人的時候.
摩訶人有一個.
他的宗教裡面有一個假神.
叫做摩洛.
那麼.
在戰況裡面怎麼去描述呢?.
他說摩訶王見到陣勢甚大.
難以對敵的時候.
他就率領七百個拿刀的兵.
就要衝過陣.
然後去到以東王那裡避難.
卻是不能.
27節.
便將那應當.
接續他作王的長子.
在城上獻為飯食.
那時發生什麼事呢?.
以色列人.
遭遇耶和華的大怒.
其實不知道為什麼中文譯成.
耶和華的大怒.
遭遇大怒.
於是三王就離開了摩訶王.
他就各自.
回到本地去.
大家知道.
戰爭是心理戰術.

$^{361}$他這樣獻上的時候.
不僅求神靈介入.
也告訴敵軍.
這件事很殘酷.
所以以色列人見到的時候.
他說不要打.
太過.
太過殘酷.
所以在舊約聖經裡面.
先知書可很明文.
去批判.
這些野蠻的風俗.
《申命記》就是其中一節.
第18章第9節.
他說不可以有人.
把自己的耳濡腺蒂.
火祭.
這些是迦南.
野蠻的風俗.
所謂的.
獻應祭.
或者.
sacrifice.
有這樣的背景.
我們可以了解到.
阿巴罕會有這樣的反應.
就像是要求的.
他就照做.
因為在那種環境.
在那種背景下.
但我想.
問你一個問題.
整個故事裡面.
以撒有沒有死.
沒有吧.
整個故事就是告訴我們.
並且有袋屬的羔羊出現.
所以上帝.
給一個機制.
一個substitution的機制.

$^{401}$就是不需要人.
來獻上.
自己的耳濡腺蒂.
這個不是.
舊約的信仰.
也不是聖經的信仰.
因為上帝不需要你這樣做.
迦南的宗教的假神就會這樣做.
所以和迦南的宗教.
劃清楚界線.
我們的信仰絕對不會這樣.
有一個羔羊.
來作代替.
如果整本聖經這樣看.
其實就是在預表將來.
有一個真正的羊羔.
就是耶穌基督.
他願意愛我們.
來獻上自己.
所以有這樣的背景.
我們大概可以了解到.
原來這個故事寫下來.
起碼有一個這樣的.
一重的意思.
就是想和.
別教.
別神的宗教.
這些practice.
來劃清楚界線.
但是.
很清楚這個故事的目的.
不是主要做這件事.
因為如果要主要做這件事.
的話.
一句就可以說.
就像申命記那樣說.
你就是不可以獻你的兒子作繁祭.
故事的目的是什麼呢.
和我們信心的成長有什麼關係呢.
我們說過.

$^{441}$亞伯拉罕在他人生當中.
有兩個很重要的關頭.
是上帝呼召他.
讓他能夠前進的.
第一個關頭是什麼呢.
就是第十二章.
他剛開始的時候.
上帝呼召他離開本地本族本格.
叫他撇下.
他自己生命裡面.
有的security(安全).
然後去到迦林地.
以至萬國能夠得到祝福.
這是第一個關頭.
第二個關頭是什麼呢.
就是這個故事.
就是他差不多結束的時候.
上帝已經賜了.
兒子給他.
給了他.
然後他就要去試驗他.
二十二章第一節.
他說這事以後.
這事以後就是說.
兒子已經出世了.
你知道亞伯拉罕按照.
聖經裡面的情節.
他七十五歲離開.
米索不達米亞.
然後去迦林.
創世記律十二章.
七十五歲的時候.
那時候他沒有兒子.
結果等了二十五年.
兒子終於出現了.
出世了.
他一百歲的時候出世.
這個故事告訴我們大概.
以則有多大.
背得起柴的.

$^{481}$我想不應該是手譜.
問到一些有意義的問題.
不應該是硬硬學語的時候.
所以大概應該是.
年少的時候.
是不是.
一十歲多一點.
還能走三天路.
肯定不是手譜.
就在這段時候.
這些事以後.
這個問題就是這樣.
亞伯拉罕有了二十.
他就要去.
亞伯拉罕有了二十.
二十又長大了.
是不是.
差不多快成年了.
這個其實應該是亞伯拉罕.
最風光的時候.
就是說是他最.
春風得意的時候.
最順利的時候.
有兒子就萬事足.
是不是.
這個是上帝.
祝福他最多最多的時候.
你可以大概在.
二號華對他說的話裡面.
你能夠理解一點.
他是多麼疼愛這個兒子.
我記得以前也說過.
我們想起兒子的時候.
我們不應該想起今天.
我們抱著兒子的時候.
或者有人生了兒子.
恭喜你啊,很可愛啊.
有理有道.
兒子在古代的世界裡面.
代表什麼呢.

$^{521}$就是你的411k.
你的生命裡面的.
security.
古代的人老了之後.
沒有兒子很慘,橫屍街頭.
沒人照顧.
有個兒子其實就是.
養個兒子來.
退休.
他去照顧你.
你看到神怎麼跟他說.
中文又是.
搞錯了的詞綴.
第一句應該是.
你的兒子.
第二句.
你的獨生子.
第三句.
你所愛的.
那個兒子.
三句裡面.
神似乎跟他說.
帶著你的兒子.
神跟他說.
我知道你很愛兒子.
第二句.
你的獨生子.
第三句.
你所愛的.
那個兒子.
你太愛他了.
整個故事裡面,兒子這個字我不知道你們有沒有留意.
出現了10次.
那種依戀和.
那種情感的強烈.
我們應該感受得到.
你很鍾情那件事.
你很喜歡那件事.
你朝思暮想.
你很抱著.

$^{561}$那件事.
你的安全.
神就跟你說.
我要去試驗你.
試驗是什麼意思.
試驗就是說.
到底你去跟隨神.
你去信神.
到底是為了.
神給你的好處.
還是為了神自己.
是不是因為上帝給你祝福.
有不成.
所以你去試驗.
有不成.
所以你去跟隨主.
還是你跟隨主的時候.
是不計祝福.
我們有孩子的.
我們知道.
你獎勵他的時候.
他很乖.
這個不是真的乖.
但如果你不獎勵他,他也會乖.
這個叫做.
真事.
你明不明我意思.
如果你跟他中間.
有所謂贈福.
贈福.
獎勵.
有好處的時候.
他先去跟隨你的話.
這個不是很純正.
但到他長大了.
他發現不需要.
去賄賂他的時候.
不需要給他獎勵的時候.
他也會照樣做.
你喜歡的事.

$^{601}$照你的說話去做.
這個就叫做成長.
亞伯拉罕信心的成長.
都是一樣.
12章開始.
創世的12章.
上帝呼召你.
他說你要出來.
你要成為大國.
我會祝福你的.
藉著你,你要祝福其他人.
有贈福在.
你走出來.
你冒險.
你冒一冒險.
你來事奉我,你跟我.
然後你就成為大國.
有兒子.
有一個家族.
然後你就可以祝福其他人.
這是一個不錯的贈福.
我們申請信用卡的時候.
我們經常說有沒有贈福.
有沒有好處.
多少現金.
我才做會員.
亞伯拉罕第一次.
亞伯拉罕第一次接受呼召的時候.
就是這樣.
但到這裡的時候.
神就要試驗他.
要測試他.
到底你是不是因為有贈福.
你才信耶穌.
你才跟耶穌.
你才去事奉他.
這個測試這個字的目的.
逢親.
聖經裡面說.
測試.

$^{641}$如果這個主題是上帝.
是神的時候.
目的永遠都是好的.
就像煉鋼一樣.
要將雜質消除.
這個就是這個意思.
二十二章第三節.
這裡.
他接受這個呼召的時候.
亞伯拉罕做了什麼.
他說清早起來.
然後預備了他的奴隸.
帶著兩個同步.
然後自己的兒子.
就劈他的柴.
朝上.
你的次序.
應該是怎樣.
你想想.
你應該怎樣.
你就起來.
劈柴.
然後最後遇上奴.
你出門上班都是這樣.
刷牙.
吃宵夜.
叫Uber.
你沒理由先叫Uber.
然後才來.
忘記吃早餐.
我不知道為什麼有這個意思.
有很多拉比.
猶太人拉比挺有趣.
看聖經.
去咀嚼裡面的用詞.
為什麼最後才劈柴.
到底是不是思覺失調.
還是.
劈柴才是最決定性的動作.
因為你可以.

$^{681}$拿走所有東西.
然後去到那個地方.
你要去獻祭的時候.
你還要去劈柴.
不好意思了,我先回家.
最重要的事.
最後才做.
還要走三天路程.
總之就是告訴他.
這不是一時.
是一時衝動.
你要承諾.
你要去盡.
你要想清楚.
跟隨耶穌基督.
是一回什麼事.
你計算成本.
然後你拿走.
有一個猶太人的神學家.
掙扎了一個問題.
在他的書裡面這樣說.
他說回憶.
猶太人在.
猶太人最大的事件.
歷史事件.
在大屠殺.
在德國的大屠殺.
他們在集中營.
他說很奇怪.
在死亡營的時候.
仍然守七日.
他們削不照樣.
有他們的節禮.
照樣去敬拜.
照樣去讀經,去祈禱.
割禮照做.
在裡面生孩子照做割禮.
在裡面有什麼節期.
他們照樣守.
那個神學家就很奇怪.

$^{721}$為什麼呢?.
因為好像上帝已經放棄了這群人.
他們仍然去堅持.
為什麼會這樣.
去問那些猶太人.
那些猶太人怎麼回答呢?.
他說我不知道.
為什麼現在會這樣.
我不知道為什麼上帝搞什麼.
我們受這麼的難.
但是我們這樣做.
是因為上帝.
是我們的上帝.
不是因為上帝給我benefit.
給我風山水地.
給我順利.
我就跟隨他.
這就是故事的精髓.
也是故事最獨特的地方.
就是這個故事.
不可以今天照樣主管.
為什麼呢?.
因為在上帝和亞伯拉罕的約中間.
成與敗的約.
是在乎以實.
只有亞伯拉罕和神之間的關係.
才是這麼獨一.
唯一一個爛jor .
他有信心繼續跟隨神.
就是以實.
就是他所愛的獨一.
他所愛的獨生子.
所以上帝要去試煉他的時候.
就是針對以實.
所以不可以照樣主管.
今天我獻上我兒子.
沒有的.
因為你和上帝的關係.
不是在於兒子.
但是你也可以問一個問題.

$^{761}$到底我們和神之間.
是為了好處.
去信祂.
去敬拜祂.
去事奉祂.
還是我們敬拜祂.
去事奉祂.
因為祂就是上帝.
沒有其他東西在中間.
很多時候我們開頭的時候.
我們總會.
找一些好處.
相信得很清楚.
我們人是有緣約.
我們人是有.
我們的.
所謂初步的時候.
我們有很多祝福.
有很多的供應.
等等.
但我們總要去到一個地步.
你問我們自己.
到底你去事奉上帝.
你作為一個基督徒.
都是什麼原因才可以.
是不是因為上帝.
讓我順順利利.
去事奉祂.
如果上帝不讓我順利的時候.
我就走了.
在我人生裡面.
我認識了很多很多基督徒.
有不少人.
都是因著.
著數才來見上帝.
有很多弟兄姊妹.
是因為有好處.
有些benefit.
可以事奉神.
如果你要.

$^{801}$你的信心繼續去成長的話.
你就要confront(抗拒)你自己的問題.
我去教會.
是不是因為有著數.
才去.
很值得.
所以我去.
我去打開聖經.
我去看神的話語.
有時候我們都可以是.
講著數.
有得著.
我就去看.
沒得著.
我就想去拜祭.
我認識很多人.
事奉是為了有著數.
是不是feel good.
如果我做.
那個責任.
那個事奉.
是讓我能夠feel good.
給我一點fulfillment(滿足).
做完之後很過癮很爽.
of course.
但如果拿走那件事.
你還做不做.
奉獻也是一樣.
我可以奉獻.
但如果有incentive(勵志)有1\%cashback(現金).
yeah maybe i'll do more.
我們有incentive.
但你要信心繼續成長的時候.
你就要將那件事拿走.
可能那就是你的eason(依賽).
那就是你的security(安全).
你放錯了位.
神才是你的security(安全).
我們可以講直堂.
我們講教會發展.

$^{841}$每樣東西都可以.
我們可以開頭的時候很過癮.
很爽.
很exciting.
但到我要付出的時候.
你想想不容易.
我要去建堂.
我要去.
讓教會發展到直新堂的時候.
我要付出很多東西.
有很多東西我要重新去建立.
我人際的網絡.
我時空的環境.
我要重新的習慣重新去認識.
多辛苦.
我們不是十幾歲.
你多年輕都好你都沒有十幾歲.
我們人大了.
我們很不喜歡.
那些不familiar(熟悉).
不熟悉的環境.
但上帝說.
就是要去試問.
這就是我們的spiritual strength(靈性).
如果我們沒有這個spiritual strength(靈性).
我們講什麼都沒有用.
我們可以開不同的團體.
開不同的小組.
做不同的事工.
一路加上去.
我們可以繼續加上去.
但什麼是重點.
如果你的faith(信念)沒有增長.
我們不會屠殺你.
因為我們還在講有沒有著數.
如果.
我們今天跟著的話.
我們就要清清楚楚問我們自己.
都是為了什麼benefit(好處).
而我們願不願意去give up(放棄).

$^{881}$我覺得這是一個很大的paradox(幻想).
spiritual growth(靈性成長)就是一個paradox(幻想).
你捨去你的異常.
你會得回異常.
並且會得更多.
這起碼就是耶穌對門徒說.
人若為我捨棄.
背起十字架.
跟著你就可以得著真正的生命.
要得著生命的就捨棄生命.
如果你肯捨棄生命的話.
你就得著一個真正真正的生命.
這就是那個反含.
就是信仰前進.
信仰成長的反含.
我跟著我去事奉主.
不是因為順順利利.
不是因為有benefit(好處).
不是因為個人.
而是因為祂是我的主.
祂愛我.
甚至為我捨棄.
我值得去盡.
值得的.
我們結束的時候會受聖讚.
提醒一下我們.
聖經裡面說.
上帝既然愛惜自己的獨生子.
願意他的獨生子耶穌基督為我們死.
有什麼祂不會給我們.
意思不是要我們再看看有什麼好處.
而是告訴我們.
我們將來的未來.
一個這樣的主.
這個是我的上帝.
我跟著祂.
就是因為祂是我的神.
預備我們心海和心西側的祈禱.
主,你叫我們的生命繼續成長.
我們的信心繼續來練正.

$^{921}$教導我們.
看見你自己.
不是看見你給我們的禮物.
焦點在你自己身上.
不是你所給的祝福.
主,警醒我們,提醒我們.
到底我們和你之間.
有什麼難在.
主,我們求你建立我們的信心.
幫助我們今天.
拿起這個餅的杯子.
因為你是我們的主.
就是因為你是真正的主.
我們相信你的時候.
你有不同的福利給我們.
永生最得赦免.
救人.
但我們跟隨你的時候.
不是因為你的福利.
而是因為你是那位.
值得我們敬拜師父的神.
主,讓我們心靈反映這件事.
讓我們信心.
去反映這個事實.
你是我們的神.
祈禱奉傳明禁.
\newpage



\section{以弗所書}
\label{sec:6IhSTx8c43Q}
\textbf{10/14/18 主日講道: 以弗所書與教會選舉 (張麗仙傳道)}
\newline
\newline
連結: \href{https://youtube.com/watch?v=6IhSTx8c43Q}{\texttt{https://youtube.com/watch?v=6IhSTx8c43Q}} ~~~~ 語音日期: 2019-07-31
\newline
\newline
\hyperref[sec:1yfYPGx_z_4]{\small{< < < PREV SERMON < < <}}
~
\hyperref[sec:index_chronic]{\small{[返順時目]}}
~
\hyperref[sec:index_scriptual]{\small{[返順卷目]}}
~
\hyperref[sec:dClKoGWQ8rs]{\small{> > > NEXT SERMON > > >}}
\newline
\newline
$^{1}$各位聽眾好.
我記得上一個月講道的時候.
跟大家虛構了一個阿Ben的故事.
還記得嗎?.
今天也有一個虛構的人物.
天外來客.
不要以為他是有兩隻翅膀的.
也不要以為他像UFO一樣.
他只不過是像你和我一樣.
穿得旗旗利利.
這個天外來客在城匯那裡走來走去.
想去了解一下城匯的情況.
我用三幅圖畫告訴你.
他做了些什麼.
第一幅就是.
他想看看最近城匯有什麼最熱的事發生.
泊車?.
好像不是.
泊車也是.
不過不是最熱的話題.
最熱的話題就是Mark Sir花了很多時間站在那裡說的.
就是會有選舉和會有登記和領袖選舉.
天外來客還在訪問了這些會友.
知道他們的感受.
有些人就說.
Mark Sir搞這些東西很煩的.
選舉.
你看看有多少規矩.
有多少規矩.
有些人就說.
好啊好啊選舉好.
有些人就說.
一想起這些選舉投票.
好像想起了多些政治.
OK.
不過沒有人像Mark Sir這樣說的.
是想拉票的.
好像沒有.
因為很多人都想領選他.
接著第二幅畫就是.

$^{41}$天外來客又看看會有什麼事發生.
我們也知道跨世代團體過去大半年時間.
在查這個與不所書.
而且有一班很忠心的組員.
不是 組長.
他們在這大半年時間.
不單每個月帶查經.
而且帶查經之前還要很聽勵聲傳道.
一起去預查.
所以這班組長是很忠心的.
過了大半年時間.
終於查完與不所書了.
他們大家都很有滿足感.
組員也很有滿足感.
天外來客就訪問一下.
究竟你們讀與不所書讀了些什麼.
每個人的解釋是不同的.
每個人的回答也不同.
總之就是.
越勤力的就越答得開心.
這是第二幅圖畫.
第三幅圖畫就是.
我們每個星期二上午在這裡.
我們四位教務同工都有會議.
這個天外來客得了下來.
他還問了一個問題.
他說與不所書和現在教會的會友登記領袖選舉有什麼關係.
讓他問一問.
我心想有什麼關係呢.
我覺得這個問題也很有意思.
所以我花了幾天時間.
想想如何回答他.
今天就在這裡回答他.
大家也不少看過與不所書.
或者看一張.
或者看幾節.
但如果你是曾經.
全國的與不所書都看過的話.
你會很清楚知道.
與不所書有一個主題.

$^{81}$就是教會.
教會.
與不所書的主題是教會.
而如果說到教會的時候.
我們成為教會.
當然要很重視.
究竟與不所書對我們說了些什麼呢.
今天在這裡我想和大家說.
特別兩件事.
第一就是.
與不所書很清楚地說教會是什麼.
教會是什麼呢.
教會不是一塊東西.
教會不是一個建築物.
教會是人.
說得更加清楚就是.
教會是神創造的一個新人類.
聽清楚.
教會是神所創造的一個新人類.
由古今中外.
凡真心悔改.
相信耶穌為主的人組成的.
包括你和我.
如果包括你和我.
我們就要問.
我要問你.
或者你也要問我.
究竟我們清不清楚.
神為我們做了些什麼呢.
或者可以這麼說.
究竟我們清不清楚.
我們信什麼呢.
究竟我們清不清楚呢.
與不所書第二章幫助我們.
因為他講得非常清楚.
他列得非常清楚.
一個基督徒是信什麼的.
教會裡面每一個人.
你是信什麼的.
與不所書第二章.

$^{121}$第一章特別是.
他說我們每一位神所創造的新人類.
我們基督徒.
我們是被神揀選的.
在基督耶穌裡面揀選我們.
使我們成為聖潔.
沒有瑕疵的.
然後又說.
我們是被神的愛子的血救贖.
使得我們最得意赦免.
然後又說.
神是特意揀選我們.
使我們預定我們有兒子的名份.
就是說以前我們犯罪離開了神.
但是祂現在拉我們回來.
使我們得到兒子的名份.
兒子的名份是什麼.
我們是會承受遺產的.
我們是會承受神在天上的基業的.
而且神不單止說了這些.
祂還非常認真.
你可以肯定.
非常認真是什麼.
祂打了一個印下來.
這個印是由什麼造成的.
就是聖靈作為印記.
各位弟兄姊妹.
你知不知道.
只要你真心相信耶穌為你的救主.
你已經是一個基督徒的時候.
你是有聖靈在你身上作為印記.
告訴你你是神的兒女.
這件事是打風下雨.
什麼都吹不冷.
天崩地裂都吹不冷.
神非常認真.
難怪保羅在《以弗所書》第一章.
說了這句話.
他說「願眾贊歸於我們主耶穌基督的父神」.
因為祂在基督耶穌那裡.

$^{161}$神賜給我們天上各種屬靈的福氣.
就是我以前所說的.
祂使我們成為祂的業.
能夠得到天上的基業.
我說.
新人類.
意味著舊的我們.
是因罪而敗壞.
是沒得救的了.
所以神要創造這個新的.
《以弗所書》給這個新人類的名字.
就是教會.
我剛才說的.
你可以想.
新人類等於教會.
神所創造的新人類就是教會.
我們知不知道呢.
我們成為了新人類.
我們知不知道.
我們是神所創造的新人類.
我們這個教會.
和天下普世的教會.
我們是神所創造的新人類.
而這個教會的主是誰呢.
主任牧師就是陳崇基.
但是這個教會的主是耶穌基督.
這個耶穌基督是怎樣的呢.
就是神照在基督身上所做的大能大力.
祂曾經死而復活過.
現在耶穌基督坐在神的右邊.
祂是這個一切執政的掌權的靈界一切的勢力.
而且不是今世的.
連來世都超過.
這個就是耶穌基督從死裡復活.
能夠得到一個拯救的大能.
所以神能夠做了這個新的人類.
而耶穌基督是今天教會的主.
教會是祂的身體.
耶穌基督是教會的主.
我可以用《爾伏所書》第二章第一句說話.

$^{201}$來總結這個教會最大的特徵.
我們這群人都是死在過犯罪惡裡.
神叫我們活過來的.
我記得我以前講過這句話.
我們是死在過犯罪惡裡.
神叫我們活過來.
使我們得到二子的名分.
使我們得到天上的基業.
還有保羅提出一樣非常非常重要的.
我們真的要打開耳朵聽清楚.
就是我們得到這件事.
我們得到成為神的兒女.
教會的一份.
我們完全不是因為自己的功勞.
第二章第八九節說.
你們得救是本乎因.
也是因著信.
你們得救是本乎因.
不是因為自己的行為.
免得有人自誇.
這是要很清楚的.
一定不要看錯.
我們能夠成為神教會裡.
我們能夠成為新人類的一份子.
是恩典.
不是行為.
不是任何功德.
我曾經很多次問弟兄姊妹.
如果你要將基督教的信仰.
和世界上這麼多宗教信仰比較.
你以為你能說出一樣最不一樣的地方嗎?.
有很多弟兄姊妹說.
是平安.
是愛.
我可以告訴你.
很多宗教都有平安.
都有愛.
但唯獨基督教.
這件事是其他宗教沒有的.
就是恩典.

$^{241}$一切宗教都說.
你要做什麼好事.
平時已得到上天堂.
已得到神的喜悅.
但唯獨基督教告訴我們.
是靠恩典.
不是靠行為.
所以教會是一個什麼地方?.
教會是一個.
一群受過耶穌基督.
百百恩典的人.
神要求我們.
在我們的生命裡.
在我們的教會裡.
活出恩典.
所以去到第二點.
教會應該怎樣活?.
活出恩典.
有幾樣東西值得我們留意.
保羅在《爾弗所書》.
保羅是《爾弗所書》的作者.
他在《爾弗所書》裡清楚說明.
第一樣就是.
教會不能夠排斥非我族類.
在第二章裡.
很有趣的.
你看到保羅經常說.
我們.
你們.
究竟我們你們是什麼呢?.
原來他指我們是猶太人.
你們是外邦人.
原來當時在《爾弗所書》的背景裡.
保羅身為猶太人.
代表了很多猶太人的自尊.
一種自己的誇耀.
就是我們是神的選民.
其他那些人都不是.
他們是恩典以外的.
好像這句話.

$^{281}$你們和基督無關.
你們是局外人.
你們在神的恩典約裡.
沒有希望.
你們沒有神.
這是猶太人怎麼看待外邦人.
他們是很抗拒外邦人.
覺得他們不是像我們這樣.
神聖的選民.
但是.
但是在第二章裡.
保羅繼續說.
你們死在過犯罪惡之中.
夠壞了吧.
死在過犯罪惡之中.
接著他說.
我們也是.
我們從前也在他們中間.
放縱肉體的私慾.
這是什麼意思.
原來保羅已經有一個深深的領會.
耶穌基督光照了他.
他知道.
原來你們是這麼壞.
我們也是這麼壞.
我們完全不是比你們好多少.
我們猶太人和你們外邦人一樣.
都是罪人.
這裡說得很清楚.
保羅明白神的恩典之後.
他知道不能排斥非我族類.
所以他能夠說出這樣一句話.
他說.
是神拆毀了我們中間的圍牆.
以前我們分得很清楚.
現在是神拆開.
現在你們外邦人不再作為愛人和客類.
與聖徒同國.
在神家裡的人和我們一樣.
對猶太人來說.

$^{321}$這是一個很震撼的改變.
很震撼的接受.
很多猶太人都不接受.
所以你看到保羅這麼慘.
經常被人要命.
要逃跑.
是因為保羅傳這種道理而追殺他.
但是保羅卻要用生命來拼.
因為他知道神的教會.
是一個不會排斥非我族類的教會.
這是伊弗所說的第一樣.
讓我們停一停.
反省一下.
我們成會教會今天怎麼樣.
你和我有沒有排斥非我族類呢?.
我們不是猶太人.
我們很多是香港人,中國人.
我們有沒有排斥不同膚色的人呢?.
引申來說.
我們會排斥.
你是香港人.
我是越南人.
你是中國人.
你是中國哪裡來的?.
你會不會這樣互相排斥呢?.
神說不可以.
神的教會是不可以排斥非我族類的.
我們不同學歷.
不同背景.
我們都不能夠以一種優越感來排斥別人.
記住.
教會不是排斥非我族類的地方.
是恩典的地方.
這是第一樣.
第二,伊弗所說繼續告訴我們.
教會是一個怎樣的地方.
是一個不可以排斥我不喜歡的地方.
或者我們從今早到現在.
進來成為教會.
我們都有多少人是看不順眼的.

$^{361}$心裡面不說出來.
因為我們大家性格很不同.
有些人感性.
有些人理性.
有些人主動.
有些人被動.
有些人快.
有些人慢.
有些人爽捷.
有些人含蓄.
有些人前衛.
有些人保守.
總而言之就是看不順眼對方.
所以就有很多的衝突.
除了這個不同之外.
我們還有很多比較的心.
我們會比較外表.
他漂亮我漂亮.
他多錢我多錢.
他厲害我厲害.
他有什麼成就.
甚至我們會比較大家住什麼車.
開什麼車.
吃什麼.
旅遊去什麼地方.
連狗都比較慢.
所以這些比較使我們很多時候抗拒.
抗拒我不喜歡那些人.
還不要再說.
如果說政見你更慘.
因為今天美國的社會.
是很兩極的政見.
所以一坐在一起說.
一不對真的可以火紅火綠.
但是.
神說教會不是一個.
不應該是一個排斥非我喜歡的教會.
衝突的時候我們要怎樣.
保羅說衝突的時候.
我們要謙虛溫柔忍耐.

$^{401}$用愛心互相寬容.
用和平彼此聯絡.
目的是竭力保守.
聖靈所賜合而為一的心.
教會是合一的.
不要被這些不喜歡的我不喜歡的.
搞到教會分裂.
我們都經常說到謙卑溫柔忍耐.
你不要看錯.
這裡說的不是叫你個人收身這麼簡單.
沒錯,謙卑溫柔忍耐.
是聖經裡面其他地方說出的.
聖靈的果子.
是我們應該有的行為.
但是在這個上下文之下.
這裡是特別指教會.
在教會裡面.
我們怎樣去對付一些我們大家不喜歡的東西呢.
謙卑溫柔忍耐.
用愛心互相寬容.
用和平彼此聯絡.
這是我們需要做的.
教會是合一的地方.
是恩典的地方.
第三.
教會是怎樣呢.
是需要勇於面對當代的問題.
什麼是當代問題呢.
在爾弗所書裡面.
當代的問題就是淫亂.
sexual immorality.
在今天的時代.
2000年之後.
很妙.
一樣.
今天我們的社會.
今天我們面對的不是很多淫亂的事嗎.
今天的道德標準已經低到嚇死人了.
保羅指出淫亂是一個很大的問題.
在當時.

$^{441}$他提到說淫亂是污穢的.
是貪心.
記不記得我上次提到.
貪是在神的眼裡.
貪是覺得這些是我的.
我喜歡怎樣就怎樣.
這個在神的眼裡是貪.
當你這麼貪的時候.
你為自己建立了偶像.
這些是我喜歡的.
我要這些.
我沒有他就不行.
你是用這個去取代神.
當日.
保羅寫信給爾弗所書.
他們是外邦的信徒.
他們在一個很淫亂的環境之下.
在文化和宗教之下.
受了這些宗教的影響.
他們成為了基督徒.
保羅要強調告訴他們.
而且用整張.
爾弗所書整張跟他們說.
你們不可以再有這件事.
今天我們的教會.
仍然面對世風日下的問題.
淫亂.
性男教婚外情.
夫妻教婚.
還有很多很多很恐怖的事.
是極之嚴重的問題.
結婚的人要注意.
未結婚的人也要注意.
在學校裡有這種問題.
在公司裡有這種問題.
在教會裡有這種問題.
在教會外面有這種問題.
對於淫亂這個問題.
保羅在這裡說了一句話.
他說從前你們是暗媒的.

$^{481}$但是如今在主裡.
你們是光明的.
你們行事為人.
要做光明的子女.
不可以再有這些淫亂的黑暗的事.
你們不可以再這樣.
你們不要做愚昧人.
要做智慧人.
不要做糊塗人.
不要醉酒.
你以為爾弗所書在說很多不同.
第五章在說很多不同的東西.
原來他只有一個中心.
就是淫亂.
就是提到不要醉酒.
不要做糊塗人.
要做智慧的人.
但是雖然保羅在這裡.
這麼深的境界.
他給了一個正面的指引.
他激勵教會做一件事.
就是你們要用詩章.
頌詞,靈歌彼此對唱.
口唱心和的讚美神.
要換句話說.
這些是什麼呢.
保羅說教會我們要一起著重大家的敬拜.
而且不只是在教會敬拜.
就是我們這些教會的人.
我們在生命裡面也要活出敬拜的生命.
以致我們可以遠離這種淫亂的生活.
這是一個很重要的敬拜.
時間關係.
爾弗所書其實還有很多東西要說.
但是我要說到最後.
很重要的一樣.
就是保羅提到一個很重要的現實.
屬靈戰爭.
在爾弗所書第六章那裡.
當保羅說完以上重要的東西之後.

$^{521}$他提出一樣我們不可以不知道的東西.
就是神所創造的這個新人類.
原來是時時刻刻受受敵的攻擊.
一場屬靈的戰爭是不會停止的.
不會停止的.
是時時刻刻都會發生的.
在你的心和我的心.
在我們教會裡面和外面.
新人類我們要留意.
攻擊什麼呢.
這個戰爭攻擊什麼呢.
就好像我們以上所說的.
我們的受敵會攻擊我們.
不相信耶穌基督的救贖.
寶血的救贖.
會使我們輕視這個新人類的身份.
使我們不相信恩典.
只相信行為.
使我們的生命也不會出恩典.
我們的受敵也會攻擊我們.
使我們抗拒一些不同我們膚色.
不同我們學歷.
不同我們背景的人.
使我們人與人處理很多的衝突.
在教會很多的紛爭.
我們不會用溫柔牽扯來.
我們只是互相的譴責和傷害.
我們的受敵也會攻擊我們.
當面對一些像淫亂的時代問題.
他使我們沒有勇氣去面對.
沒有勇氣去指正.
翻臉的,假裝看不到.
或者甚至是隨波逐流.
我們做了其中一個.
我們的受敵是時時的攻擊我們.
保羅知道這個嚴重性.
他說了這句話.
我們面對的戰爭是一個巨大的戰爭.
我們不是肉眼所見的.
而是天空屬靈的掌權者去打.

$^{561}$這一場是一場很大的屬靈戰爭.
你注意,保羅說了兩件事.
他說我們不是屬血氣的征戰.
寫著摔跤.
原來保羅寫這封書的時候.
是在羅馬的文化背景.
摔跤是一個很受歡迎的運動.
在競技場上摔跤.
拳打肉搏.
打得很狠,有時還會打死人.
但大家看得很開心.
保羅說這一場被摔跤.
還沒有那麼可怕.
因為一拳打來一拳打去.
是可以看到的.
但我們面對的屬靈戰爭是看不到的.
但是攻擊我們.
你看不到敵人在你面前.
打死你也不知道.
死了也不知道.
不過保羅也用了一個羅馬的.
很受歡迎的比喻.
就是軍裝.
當時羅馬的軍事勢力非常強.
保羅用了一個盔甲.
羅馬人的盔甲作為一個比喻.
他說你想像這個盔甲.
想像神給你一個全副的屬靈軍裝.
他說要拿起這個神所賜的軍裝.
我們就能夠在惡魔來到的日子.
能夠成就一切.
永遠責罵得住.
因為今天的時間關係.
我不會詳細講這個屬靈的軍裝.
你也不想我講.
因為現在十一點.
但是這個屬靈的軍裝.
有幾樣東西我特別要強調的.
我們要有真理.
彰顯真理.

$^{601}$維護公義.
我們要傳揚福音.
我們要用信德守護著.
我們要用救恩抵擋受敵的攻擊.
我們要重視學習和遵守神的話語.
以致我們能夠抵擋撒旦的攻擊.
然後他還說.
我們要靠著聖靈.
時時都祈禱.
祈禱不是屬靈軍裝的其中一件.
但是祈禱是屬靈軍裝的一個根基.
我們要時時祈禱.
不要輕看祈禱在教會的能力.
這裡講得很清楚.
講到這裡.
我要回到姐姐的問題.
就是這個天外來客問過我的.
究竟已不鎖書和會有登記領袖選舉有什麼關係呢.
天外來客不是一定是男人.
我弄了一個女的在這裡.
究竟有什麼關係呢.
我們知道領袖選舉的第一步.
Mark Sir 也講過.
就是登記為會友.
會友們我們要明白的就是.
我們既然是神家裡.
我們是神所創造的新人類.
這個是我家.
這裡是我的地方.
我是這裡的一份子.
我不是旁觀者.
我不是過路客.
我要登記做會員.
也表示責任.
這個家要選淑玲的領袖.
我們每一個家人都要一起選出來.
領袖選舉的第二步.
就是提名這三個星期.
然後就是選舉.
接下來的三個星期.

$^{641}$我們需要為這個家選出成熟和委身的信徒.
作為領袖.
這些人不一定本身就很有領袖才能.
但是他一定要有愛主和委身的心.
就好像今天以弗所說.
我們需要選一些委身愛主的領袖出來.
帶領教會.
幫助信徒清楚自己的信仰.
使教會在合一的精神之下.
能夠彼此關顧.
互相建立.
使大家能夠聯手一起面對.
我們當代一些很大的問題.
可能不單只是淫亂的問題.
還有很多其他的問題.
我們也需要選舉淑玲的領袖.
能夠帶領大家一起去抵擋真正的敵人.
就是天空淑玲掌權者.
當然我們教會有牧師有傳道.
但是牧師傳道是不會做得完的.
我們的能力是有限的.
有些東西我們是不懂的.
我們是要靠整個的團隊.
整個中心的團隊.
一起起來建立這個教會.
最後就是.
究竟這個天外來客.
上去的時候.
他會為教會打一個什麼樣的報告呢?.
我但願他的報告是.
這個教會有委身的會員.
有委身的領袖.
他們大家一起都合一.
能夠做到行事為人.
與他們蒙召的恩相承.
他們能夠在這個教會裡面.
使到耶穌基督得榮耀.
他們會繼續祝福這個社區.
在這個社區作榮神益人的一個教會.
這個是我們的盼望.

$^{681}$這個是我們需要共同努力的.
我們一起祈禱.
作為新人類.
你所創造的新人類.
教會.
Local Church.
就是城會教會.
求你教我們.
從你的說話裡面.
我們再一次的明白.
教會是不能夠讓我們傾看.
是需要讓我們.
忠心,委身.
因為我們怎樣回應教會.
就是我們怎樣回應你.
多謝你的話語的提醒.
奉救主耶穌的名.
\newpage



\section{}
\label{sec:dClKoGWQ8rs}
\textbf{10/28/18 主日講道: 女子組加油 \#1: 顛覆男人世界的聖經女子(陳崇基牧師)}
\newline
\newline
連結: \href{https://youtube.com/watch?v=dClKoGWQ8rs}{\texttt{https://youtube.com/watch?v=dClKoGWQ8rs}} ~~~~ 語音日期: 2019-08-01
\newline
\newline
\hyperref[sec:6IhSTx8c43Q]{\small{< < < PREV SERMON < < <}}
~
\hyperref[sec:index_chronic]{\small{[返順時目]}}
~
\hyperref[sec:index_scriptual]{\small{[返順卷目]}}
~
\hyperref[sec:9yvFXOIg3So]{\small{> > > NEXT SERMON > > >}}
\newline
\newline
$^{1}$今天我們講一個新的題目.
新的系列.
這個系列就叫做.
厲害了.
加油.
女子組加油.
這句加油真的挺厲害的.
最近大家知道.
新聞裡面經常說.
Add oil 這句話.
就加進了牛津的字典裡面.
我們真的挺厲害的.
廣東話這句話.
Add oil.
加油.
那麼.
我們想為什麼人加油呢?.
為我們身邊的姐妹.
我們身邊的女性.
來加油.
我不知道你有沒有想過.
這個題目其實也是很切身的.
今時今日經常都會講這類問題.
Gender issue.
或者 Gender equality.
就是性別的平權等等這些問題.
因為有Me too 又殺出來了.
所以經常都要提倡.
男女要平等.
有很多人覺得勝經是很古老的.
很古董的.
跟這個時代沒有什麼話好說的了.
只不過放在書架上封塵而已.
其實不是的.
聖經雖然是一本很古老的書.
一本很古老的文件.
但是裡面所說的話.
其實是很劃時代的.
所以今天我們開始的時候.
在接下來幾個星期.

$^{41}$我們看看聖經怎麼看女性.
聖經怎麼看女性這個角色.
所以今天想先打一點基礎.
有一點 framework.
讓大家知道怎麼去看.
如果你是女性的話.
你怎麼看自己.
是不是?.
神怎麼看你.
如果我們是男的.
不要走.
看看神怎麼看女性.
是不是?.
我們用神的角度來看他們.
所以打打基礎.
我會告訴你們.
我 present(展示)出來.
有幾句話.
OK?.
你就要告訴我.
或者你跟旁邊的傾聽.
到底這句話是對還是錯呢?.
第一句話.
聖經是父權社會的產品.
你認為對還是錯?.
對還是錯?.
Product of patriarchal society.
什麼叫父權社會呢?.
父權社會就是重男輕女.
對嗎?.
重男輕女.
你生兩個女兒.
Oh no!.
起碼也要生一個兒子.
重男輕女就是這樣的角度.
或者以男人為上的文化.
以男人為主.
聖經是不是社會的文學產品呢?.
是不是?.
是的.

$^{81}$OK.
是的.
為什麼說是呢?.
你看聖經很清楚.
文化背景,社會背景.
全部的舊約也好,新約也好.
都是在假設一個父權的文化.
或者以男為上的背景.
你光是看語言就知道了.
但是神是沒有性別的.
神是零.
God is spirit.
神是沒有性別的.
但是你知道聖經裡面從頭到尾都是說什麼?.
神是He的.
是不是?.
我知道有很多女性的形象形容神.
這是挺驚訝,挺有趣的.
不是今天說.
不過無論如何.
By and large.
大部分來說.
其實都是以男人為主的.
這是聖經的背景.
這是不能質疑的.
其實甚至今天.
在中東一些山區的地方.
都是一樣的.
都是以父權文化為主的.
都是很守舊的地帶.
一個女人是不能出門的.
如果一個女人要出門的話.
就一定要有一個男人陪著.
無論那個是你爸爸.
是你的先生.
是你的兄弟.
或者是你的叔叔.
都要有一個男人陪你出去逛街.
因為如果一個女人在街上出現.
就代表什麼呢?.

$^{121}$她可以侵犯你.
這就是父權文化的背景.
所以.
這是聖經的背景.
沒錯的.
但是很可惜的就是很多時候我們拿著聖經.
就是說聖經背景既然是這樣.
所以就用聖經來做男兄女.
對嗎?.
就會用聖經來支持做男兄女.
或者是抑壓看低那些女人.
我聽過不少故事.
有些人去到教會.
某些教會.
他看到講道的是女人.
或者去到一個主學學.
聖經課程裡面.
他看到教的是一個姐妹的時候.
她跑了.
這麼有趣.
她馬上急著跑了.
謝了頓.
因為她的心裡面想的是.
既然聖經背後的文化.
是一個 patriarchal 的社會的話.
既然聖經是這樣的話.
我們現在應該是這樣的.
是不是?.
這班主管.
這是錯誤的觀念.
所以第一句話.
我們先認定.
聖經是那個時代背景出來的.
沒錯的.
第二句.
聖經是創造或者維護父權主義的.
是不是?.
對還是錯?.
很大聲.
我聽到是女生說的.

$^{161}$是錯的.
因為他沒有創造出來.
他只是假設預設當時的社會.
是不是?.
聖經是上帝對人說的話.
上帝對人說的話.
我們叫 revelation.
叫做啟示.
啟示其實是一種溝通.
要溝通的話.
就一定要用當時人的語言.
用當時人的文化.
用當時人的價值.
來通過他們所用的語言.
來跟他們說話.
是不是?.
今天我跟你說話.
我不可以用希伯來文說話的.
因為你 understand(明白)的語言.
是我們現在說的.
所以我們用你所明白的.
或者所預設的價值語言等等文化.
來跟你說話.
所以聖經沒有創造父權文化.
也沒有特意去維護以男為上這種主義.
我們知道.
聖經裡面是以父權為上.
對.
一家之主是男人.
是不是?.
女人是要馴服.
你看看我的家譜就知道了.
家譜是跟男人的.
很少有女人的.
對嗎?.
家譜是跟男人的.
你如果回到鄉下.
你去祠堂看你的族譜.
跟男人是跟女人的嗎?.
是跟男丁的.

$^{201}$對嗎?.
跟男丁的.
如果你在鄉下.
有塊地怎麼辦?.
他們還按照男丁傳授下來的.
就要有男孩子才可以拿到塊地.
對嗎?.
我沒有錯.
沒有錯吧?.
你是女人.
你是沒有地位的.
聖經的文化背景是這樣的.
你想想.
聖經裡面的英雄.
heroes(英雄).
大部分都是男人.
對嗎?.
我記得有姐妹跟我說過.
我看聖經的時候好像不是很認同.
為什麼不能夠identify(認出).
因為要我學大衛.
大衛是男人.
摩西是很好的屬靈領袖.
摩西是男還是女?.
摩西是男人.
亞伯拉罕是男人.
所以有時候不知道你覺不覺得.
看聖經你很難認同.
我講的是我男人就沒事了.
女的看聖經會不會有些奇怪呢?.
不知道怎麼去認同.
怎麼identify(認出).
我給你數據.
93 49 14056.
在聖經裡面有名的女人.
只有93個.
93個.
什麼我不說.
真的提到上了聖經舞台的.
93個.

$^{241}$有名的.
真的說出她的名字.
瑪利亞.
莫大拉瑪利亞.
底波拉Deborah.
撒拉.
這些有名的給她的.
49個而已.
49個.
這些人.
我那時候聖經的故事.
講到這些女子的時候.
有時候是直接引用她說話的.
直接direct speech(直接說話).
直接引用她有說過多少樣東西.
這些人加起來.
這些女子加起來.
說過的字.
說過的字數.
只是14000多個字.
14000多個字.
多不多?.
你是女子.
你女子.
你是女子.
你一直都不止說14000多個字.
所以.
其實你這麼說.
總數來算.
所有的direct speech(直接說話).
所有說.
誰誰說過這句話.
引用她.
其實佔聖經裡面1.1\%而已.
1.1.
就是99\%.
都是男人開口說話的.
不是女人開口說話的.
但是我想你留意.
少不代表不重要.

$^{281}$物以罕為貴.
對嗎?.
不是說得多才重要的.
是不是?.
有時候說少.
才是重要的.
所以每逢.
一有女人說話.
或者有女人.
出現在聖經舞台.
就代表你 pay attention(注意).
留意.
留心清楚.
可能有些重要的事情發生了.
是不是?.
有些突出的地方.
我不知道大家有沒有看過一部電影.
叫做《女人40》.
有沒有?.
沒有啊?.
那些看過的應該都大部分.
《女人40》是施芳芳做主角的.
講女人.
中年的女人.
講她生活上很多挑戰和壓力.
喬宏就做她的老爺.
那女人的家屬要跟老爺住已經很辛苦了.
對嗎?.
那.
這部電影你記不記得?.
那個劇情.
就是老爺患了老人癡呆症.
是不是?.
全部人都不認得.
只認得家屬.
那就麻煩了.
是不是?.
不認得家屬就不用理他了.
誰知道只認得他.
那.

$^{321}$他就要.
就是.
肩擔一個很重的責任.
然後.
在他的辦公室.
做工的地方.
那.
施芳芳就.
在辦公室裡面.
出現了一個很年輕的女孩.
那個人做事又快.
又會用電腦.
她就受到威脅.
有很多壓力和威脅.
做女人很不容易.
這部電影為什麼這麼厲害呢?.
因為這部電影是第一部.
是.
第一部香港的電影.
是拿到.
叫做大滿貫.
什麼叫大滿貫呢?.
就是拿到最.
先看看.
最佳電影.
最佳導演.
最佳編劇.
最佳男主角.
和最佳女主角.
突然間有這部電影出來.
你就知道突出了.
就是說.
你注意一下.
原來有這樣的故事是需要聽的.
需要知道的.
這就是聖經的女子.
出現的.
特色.
每逢她們突然間出現的時候.
你就應該豎起兩隻耳.

$^{361}$豎起你那兩根天線.
你就想一想.
到底.
有什麼重要的東西.
我是漏掉了呢?.
第三句話.
聖經是挑戰父權主義的觀念的.
Yes 還是No呢?.
對還是錯?.
Patriarchal values.
Challenge.
是的.
是.
挑戰.
很多時候.
聖經是特意.
來挑戰.
當時假設的.
父權為上的社會.
我們的假設就會是.
如果聖經是出自那個父權主義的文化的產品.
就應該對女人有貶低的味道.
或者甚至一字不提.
你看看那些.
古代的.
跟聖經同期的文學作品.
無論是巴比倫也好.
埃及也好.
很少提到女人的.
提到女人都是貶低她們.
或者當她們是二等公民.
聖經很多時候不是這樣的.
這就是最 surprising(驚訝)的地方.
甚至有時候聖經是質疑那些重男輕女的.
價值觀.
或者有時候聖經.
去到.
表達出一個男主角的故事的時候.
就是.
去貶低他們.

$^{401}$貶低那個男人.
說那個男人.
不是一個真正的男人.
所以聖經.
如果你看清楚的時候.
在顛倒當時.
很多的文化背景和價值.
有時候提到聖經的女子是受害者.
是受害者.
她們是 victim(受害者)的時候.
就是說她們是沒有地位.
以及是跟那些卑微的人去認同.
譬如那些女人很多時候.
聖經裡面.
跟那些孤兒.
外來者.
移民.
奴婢.
是跟她們是同一群人.
上帝就是去親近那一群人.
那一群被打壓.
或者被輔代的那一群人.
上帝就特別來跟他們親近.
如果你想看看聖經裡面.
怎麼 present(展示)這些 story(故事).
聖經裡面有一段話.
有幾章.
你不妨回去看看.
士師記最後幾章.
就是很殘酷.
就是說當時的男人怎麼對待當時的女人.
其實聖經的士師記寫出來就是在做什麼.
他寫的故事.
就是要我們看.
我們覺得很殘酷.
就是一個在父權主義文化裡面.
社會的狀況.
是一個這麼黑白不分明.
去到殘害一些無權的人.
殘害一些女性.

$^{441}$是很嚴厲的來譴責當時的以色列社會.
男人的社會.
所以聖經裡面.
經常都會挑戰一些父權主義.
這個可能你沒聽過.
如果我們男人看聖經.
看到一個地步我們覺得.
哇!我們男人很威風很厲害.
可能你看錯聖經了.
但如果我們看聖經是能夠挑戰我們的話.
我們會突然間想起.
咦!可能我真的有些事情需要改變.
這個就是聖經的用意.
第四句.
在聖經歷史裡面.
女性很多時候扮演一些策略的角色.
對還是錯?.
對的.
我給你看看.
這裡我嘗試很簡單.
告訴你整個聖經的歷史是怎麼樣.
由創世開始.
去到族長時期.
Patriarch.
然後到Church of Eccles.
入主迦南.
然後定居在迦南.
然後王國分裂.
然後被擄.
他們又回歸.
然後等候基督來.
耶穌就到了.
是不是?.
那就是第一世紀.
就是這麼簡單.
你想知道整本聖經說什麼.
That's it.
Right?.
但是在這裡.
所有的.

$^{481}$叫做什麼?.
那些策略的人.
Strategic people.
那些故事.
在哪裡出現呢?.
族長時期當然是阿巴拉罕.
阿巴拉罕以賽亞國.
出埃及.
誰是民族英雄?.
摩西.
Moses.
入主迦南.
當然是約書亞.
他帶著軍隊進去.
然後.
Settle down.
在迦南地.
就是事實.
王國時期就是誰?.
素羅大衛.
所羅門.
王當然是男人.
然後回歸的時候.
以斯拉和尼希米.
十二叛徒.
耶穌呼召的.
有沒有想過.
這個其實就是我們平時.
Study Bible.
或者.
參考書.
告訴我們.
聖經的歷史就是這樣.
點出那些重要的.
全部.
你有沒有想到.
這個 timeline.
少了什麼?.
全部都是男人.
女子在哪裡出現呢?.

$^{521}$今天我給你看看整個聖經故事.
我隨意昨晚想一想.
每逢有女人上舞台.
是什麼時間?.
是不是關鍵時刻?.
第一個.
出埃及為什麼要先拿摩西?.
因為當時法老要殺所有.
希伯來的男嬰.
有誰留住希伯來的男嬰.
以至摩西人的出生?.
就是有兩個收生婆.
那些 midwives.
幫人生孩子的.
自己沒有孩子.
但是幫人生孩子的.
有名的.
你回去看看.
出埃及的第一章.
然後就是摩西的姐姐.
摩西的子子.
Miriam.
摩西不是放在籃子裡.
然後在河裡流嗎?.
那麼她就跟著他.
姐姐跟著他.
然後就叫法老的女兒.
不如你 adopt 她吧.
於是摩西就被他自己的親生母親.
在皇宮裡面被養大了.
是不是?.
突然間女性是 pop up 的.
所以為什麼摩西會留住.
摩西其實是死症種.
但是因為女人的緣故.
就生了摩西出來.
以至能夠出埃及.
入主加林.
應該是.
怎麼樣呢?.

$^{561}$揮軍直入.
因為是正統加林.
誰出現了呢?.
第一個出現的主角.
是一個迦南的女子.
妓女拉合.
她去幫那兩個探子.
然後還教探子.
你進來幹什麼?.
神已經給你了地了.
還有什麼好探的.
所以就出現拉合了.
就叫他們.
你們信得過神的話.
是時候你們應該入主加林了.
女子的故事.
定居加林呢?.
settlement(定居).
沒錯,大部分事司都是男人.
但是突然間有一個出現.
是誰呢?.
是 Deborah.
如果 Deborah 你知道那個故事的話.
你就知道是男人不敢出戰.
然後她自己出戰.
這就是希伯來的花木蘭.
不是代父從軍.
而是代男從軍.
當然是在譴責當時的男人不夠僵.
男人都沒有僵.
那怎麼行呢?.
那 Deborah 出來吧.
事司的時代.
各自為政的.
當時所謂的胡亂華.
你看看時史記你就知道了.
有什麼女子.
告訴我們上帝的約繼續守下去.
上帝的工作繼續堅持下去呢?.
就是路德.

$^{601}$就是摩訶女子.
一個外邦的女子.
她願意來信守耶和華的約.
是不是?.
堅定地來跟隨耶和華.
然後就差不多到王國時期了.
就是薩姆爾基開始的時候.
是誰開始呢?.
是哈拿開始.
第一章.
薩姆爾基上.
是一個生不了孩子的婦女.
她求上帝給她一個兒子.
上帝真的給了她一個兒子.
生了之後.
養了之後.
她就交給聖殿.
奉獻了薩姆爾.
她就成為最後一個的事司.
也是第一個的先至.
第一個的領導.
誰的功勞呢?.
是哈拿的功勞.
然後到了王國時期.
所有事情都所謂國泰民安了.
整頓了.
就是 settle down 了.
然後有了政府了.
就沒有什麼女人出現了.
為什麼呢?.
因為父權文化嘛.
所以在第一和七王國.
找不到多少女人.
因為當時已經是 patriarchal society.
逢出混亂的當中.
逢出危機的當中.
這些女子的故事就出現了.
是吧?.
國泰民安就定了.
她們就回到自己家.

$^{641}$廚房裡面.
被擄波斯時期.
就是以色列出現.
以色列進去做皇后.
然後波斯有一個奸臣.
就想殲滅所有的猶太人.
就是他肯冒險.
以色列來介入.
以至保守了他神的子民.
他們才能夠回歸.
才能夠回到重建聖殿.
所以又在一個危機的時候.
就出現了民族的救星.
這些民族的救星就是女人.
耶穌基督的降生當然不用說了.
就是瑪利亞.
她願意去相信耶和華.
蒙恩來到她的身邊.
大救主來到地上.
耶穌基督復活.
第一個去迎接.
在空墳墓.
看到空墳墓迎接復活基督的.
又是女人.
其中一個是莫大拉·瑪利亞.
那些門徒在哪裡?.
門徒躲藏了.
害怕藏起來.
躲在房子裡面.
所以每一次都是重要關頭.
無論是以色列的女子.
或者是外邦的女子.
她們出來拯救神的子民.
是在推動上帝的工作繼續向前.
到目前為止只不過是很簡單的說.
你就知道這些女英雄的故事.
似乎告訴我們.
聖經.
沒錯是66本書拼在一起.
似乎拼在一起編輯的人.

$^{681}$或者那群人.
或者背後的上帝.
是有用意.
去精心安排這些故事的.
不是亂打亂撞放在一起的.
好像你聽一首歌一樣.
在一個專業的人編曲之後.
你知道有些歌.
你聽聽看.
有一段solo出來的.
譬如電結他solo.
五分鐘的歌.
它的solo有多長?.
不會四分半鐘吧.
可能三十秒而已.
但就在那裡是精髓的地方.
所以麥爾漢遺跡就在這裡.
我們這樣打打基礎.
我想問你.
到底聖經是不是我們要重新再看呢?.
你想一想.
聖經怎麼看女性的角色?.
我想你帶走幾件事.
就是我們盡量除去那些.
就是重男輕女的思想.
上帝可以用男人做事.
更加可以用女人做事.
重男輕女不一定是男人有的觀念.
可能各位女士也會有重男輕女的觀念.
讓男人做吧.
男人是可以的.
其實有時候男人是不行的.
你不要有重男輕女的觀念.
覺得我也是輕視我自己的地位.
聖經告訴我們.
上帝要推進他的工作.
很多時候是用姐妹的.
怎麼除去重男輕女的思想呢?.
幾點而已.
第一.

$^{721}$我們想起領袖.
想起長輩.
想起團長.
領袖的人.
你會想起什麼人呢?.
你profile一個人出來.
如果你經常都profile一個男人出來.
那應該是時候要改了.
有男也有女.
做教師的.
教聖經的.
你怎麼profile一個港島的人.
你心中的印象.
他是一個怎樣的人呢?.
有男的也有女的.
我們跟姐妹要同工.
一起去事奉.
姐妹要告訴男的.
我們應該怎麼跟你同工.
怎麼去明白你.
不是說.
你領導的方法不是一個領導的方法.
而是大家領導的方法不同而已.
我們怎麼在異同當中.
我們能夠一起同工.
在家庭也是一樣.
做父母的.
男人的.
請你做男人.
男女也是平等的.
大家是同等.
同權.
一起做partner.
在教會.
在我們事奉當中.
讓我們知道.
咦! 聖經原來有這麼劃時代的訊息給我們.
我們同心低頭祈禱.
做很多的工作.
他們很熱心地來服侍.

$^{761}$主我們求你.
來到親切的祝福使用他們.
主我們知道你自己的工作.
你自己的.
地上神國的福音的運動.
是要藉著男.
藉著女.
來向前的.
主你求你使用我們每一個.
讓我們除去一些舊有的思想.
讓你自己的聖經來挑戰.
來更新我們的觀念.
以致我們同心同意的.
來服侍你.
彰顯你自己的國度.
我們這樣祈禱奉主耶穌的名字.
\newpage



\section{}
\label{sec:9yvFXOIg3So}
\textbf{8/26/18 主日講道: 教會藍圖\#1 -- 在醜聞中如何重建教會 (陳崇基牧師)}
\newline
\newline
連結: \href{https://youtube.com/watch?v=9yvFXOIg3So}{\texttt{https://youtube.com/watch?v=9yvFXOIg3So}} ~~~~ 語音日期: 2019-08-01
\newline
\newline
\hyperref[sec:dClKoGWQ8rs]{\small{< < < PREV SERMON < < <}}
~
\hyperref[sec:index_chronic]{\small{[返順時目]}}
~
\hyperref[sec:index_scriptual]{\small{[返順卷目]}}
~
\hyperref[sec:ALcJ_99Kiw8]{\small{> > > NEXT SERMON > > >}}
\newline
\newline
$^{1}$.
提摩太前書第四章第二節,這裡又提到良心,說到一些全假道的人,他們是假冒的,保羅說這等人的良心如同被熱鐵下灌尿一般,意思是說他們有一個意象,我們的良心好像那種,你試過有些傷口,戒損了之後,.
然後你去撓它,撓著撓著,它又流血,又再損,然後它又再戒,你又再撓它的時候,它就會變成疤痕,有疤痕的地方,或者你身上也有不少,你的感覺度就會漸漸降低,.
保羅說如果我們今天在一個不利的環境裡面,我們第一件事就是需要去取回我們生命裡面的良心,一個good conscience,讓上主洗滌了我們之後,我們的良心我們的conscious再一次活躍上來,.
並且主的恩格外豐富,使我在基督耶穌裡面有信心和愛心..
第十五節,基督耶穌降世,為要拯救罪人,這話是可信的,是十分可佩服的,在罪人中可視為罪魁,而我蒙了憐憫,是因為耶穌基督要在我這罪魁身上,顯明他一切的忍耐,給後來信他得永生的人作榜樣,但願尊貴榮耀,貴給那不能扭壞,不能按鍵永世的君王,獨一的上帝,直到永遠完,阿們..
看了這段聖經,你有什麼感覺呢?.
一個很有經驗的牧者,他是第一代的神學家,你想想,保羅寫的東西,今天還有人在研究,他是一個頗好漢學問的人,一個頗有見處的人,他居然寫給他的徒弟講什麼呢?.
他的口氣,你留意,是沒有自大自滿的口氣,完全沒有一丁點說我已經上了岸,我是很有經驗的牧者,你要聽我說..
他說什麼呢?他說他今天,他的雙腳還站不在上帝的恩典裡面,他是罪人裡面的罪人..
我們聽到這些醜聞的時候,我們會很失望,原因有時候是這樣的,牧職這件事是神聖的,職位的職位,牧職,我們叫神職人員,是很神聖的,那個職位是很神聖的..
但我想告訴你,那個職位裡面事奉的人,就不是神聖的,你明白嗎?那個職位是神聖的,但那個在職位裡面事奉的人,絕對不是神聖的..
牧師,牧者,他不是神聖的,傳道人,長執,他們不是,聽清楚,他們不是90\%的罪人,他們也不是99\%的罪人,他們是100\%的罪人,任何人,包括我在內,男與女,我們都是百分之百的罪人..
所以不要將牧者,雖然他在聖堂裡面,他去事奉,他接受了一個呼召去那裡事奉,就把他推到神位,我們有時候不知不覺也會的,因為一起祈禱,誰去祈禱?當然是找牧師祈禱,牧師祈禱比較靈,不是的,我們祈禱是奉耶穌的名,.
不是奉牧者的名,無論你多麼尊敬他也好,你對他一般般也好,牧者不是神聖的,但牧職是神聖的,我們應該要尊重..
所以,很多時候聽到一些性騷擾,性失德的問題,我們經常都會問,為什麼會發生這樣的事情?有什麼因素?我相信有很多因素,Stanley Grant 也在這裡講了很多很多的因素,但最重要的因素是什麼?.
就是我們每個人都是百分之一百的罪人,保羅說我是罪人之中的罪魁,他只不過是一個蒙恩基督教的罪人,到他有見樹的時候,保羅說我仍然是一個罪魁,.
我不是以名目去自居,不是一個有名的人,明星,我們叫做celebrity pastor,來自居,只不過是一個蒙恩基督教的罪人..
所以,有聽到這些消息的時候,我們會感到很失望,不過同樣,這是一個現實,罪人之中,我是一個罪魁..
如果連保羅也這樣說,今天有誰可以誇獎自己有多厲害?.
在過去幾個月裡面,我們講到這個教會,雖然有些處理不當,但有一樣東西我覺得可以肯定和可以欣賞,.
是什麼呢?起碼最後肯認錯,全體的領導甚至願意承擔然後下台,這個真的很難做的,.
勇敢面對自己的失敗,成功的時候我們會拍掌,但面對失敗的時候,這真正是一個試煉,願意謙卑,讓外界其他的第三者介入處理這個問題..
這個可以值得欣賞,雖然很痛苦,但我相信這是一個教會的新開始,願意下台,承擔責任..
香港教會更新協會的一個牧師胡志偉說到這件事的時候,.
他說如果換上是華人的堂會,大可能都是死不認錯,死撐,結果同性留下,不同意的可以離開,.
這個傾向就是要保守自己的權力,保守自己的聖餘,而不是保守上帝的聖餘..
在教會裡面,我們成功的時候我們會學習,但其實失敗的時候我們都應該學習..
你聽到一些牧師跌倒的時候,聽到一些基督徒失見證的時候,你的反應可能會很不同,可能你會覺得這整件事很傻,你可能會覺得這個信仰是白痴的..
但請你不要忘記,請你不要忘記有幾樣東西,請你不要將你的信仰放在一個人,或者一個制度身上..
無論那個制度外表有多美麗,無論那個人外表有多值得你去欣賞,你的信仰,你的信心不是放在人或制度上,你的信心是放在上帝的信仰..
基督耶穌降世,為要拯救罪人,這話是可信,是十分可佩服..
這句話是初期教會認信的confession,其中一個,一個宣言,trustworthy saying,是要我們full acceptance,就是將我們的信仰,我們的信心放在救我們的主..
不是某某的人,或者某某的機構制度,某某的名字,而是放在上帝的名字上..
所以你有什麼反應都好,請你放正你自己的信仰,放在上帝的信仰..
第二,請你都小心,我們不要覺得我們自己是隔岸觀火..
保羅在另一封書信說,所謂自以為站得穩的,需要謹慎,免得跌倒..
我們見到別人跌倒,我們自然會覺得,咦,我站得很穩..
保羅在這裡警惕我們,不要自告,因為我們都是百分之一百罪的,我們不能夠跟失德的人說,咦,我比你好..
請你小心,我們每一個自以為站得穩的人要謹慎,免得跌倒..

$^{41}$第三,我想告訴你,有很多人仍然是很努力的,很有誠信的,去服侍教會的..
請你不要忘記,有很多這樣的人..
新聞裡面特別強調某幾單,這些失德的問題..
但是我想告訴你,有很多人正在努力的,在建立教會,很有誠信..
願意奉上來服侍我們的上帝..
但是不是因為我們厲害,不是因為人跌倒,而是神的福音是可信的..
保羅也是一樣,保羅沒有說,咦,我繼續去做,我繼續去事奉,繼續去信..
是因為我沒有跌倒過,他說什麼呢?.
可信的是基督耶穌說了,要拯救我們罪人..
恩典裡面重新建立教會..
我們之後兩個星期,我們繼續掙扎這個問題,可能不是很舒服的問題,.
不過是需要我們去正視,需要我們去面對的..
最後的,我想告訴你,提醒我們每一個蒙恩得救的罪人,我們需要立志,.
在我們生活裡面任何的一個方位,盡心去事奉主..
可以放心低頭祈禱..
主我們今天的祈禱,是不叫我們如見是談,救我們脫離凶惡..
主願你自己在國度,你自己的主權,臨到我們的生命,.
臨到我們生命裡面每一方面,無論是性的方面,.
我們個人的生活,我們的操守,我們的家庭,我們很多生活裡面的枝節..
主讓你給我們一個無藝,清潔的良心,來愛你,事奉你..
主我們為了其他的教會,可能是遭遇一些性醜聞的影響,主我們為了祈禱,.
為了Wheel of Quick這個教會,我們去禱告,求你親自與他們同在,.
讓他們經過悔改,在這個時機裡面,你自己做新的工作..
主我們求你讓當中的領袖,長老,牧師,他們有一個重新更新的機會,.
因為你是有恩典的主.主我們在你面前,沒有一個人能夠垮口,.
因為我們每一個都是罪人當中的罪魁,.
是蒙你的拯救,蒙你的恩典,一個你不值得接受的恩典,.
我們重整過來的,幫助我們在恩典當中,我們要長進,.
幫助我們教會能夠成為你自己更好的見證..
縱使環境是惡劣,是不利,我們仍然堅信耶穌基督的降世,.
是可信,可佩服.主你拯救我們,你要拯救更多的人,.
藉著我們這個這麼破碎的生命,不完整不完美的生命,來服侍你..
不是因為我們厲害,而是因為神,你是值得我們相信和佩服..
我們這樣祈禱,奉主耶穌基督的名字..
希望大家有所得益..
\newpage



\section{}
\label{sec:ALcJ_99Kiw8}
\textbf{安息日為人設立 - 張國棟博士 (全篇講道)}
\newline
\newline
連結: \href{https://youtube.com/watch?v=ALcJ-99Kiw8}{\texttt{https://youtube.com/watch?v=ALcJ-99Kiw8}} ~~~~ 語音日期: 2019-08-01
\newline
\newline
\hyperref[sec:9yvFXOIg3So]{\small{< < < PREV SERMON < < <}}
~
\hyperref[sec:index_chronic]{\small{[返順時目]}}
~
\hyperref[sec:index_scriptual]{\small{[返順卷目]}}
~
\hyperref[sec:PGikCuC4cZg]{\small{> > > NEXT SERMON > > >}}
\newline
\newline
$^{1}$我們今天有一位客席講完,很難得.
張國棟博士Daniel來到芝加哥.
跟我們開兩個星期的課程.
我都請了他來分享神話語.
張博士是在俄亥俄州來的.
他在Dayton University教哲學.
他一向的專科都是哲學,商業倫理等等.
他除了寫書之外,他也在香港神學院教一些特別的課程.
剛剛Summer教完,他就回到美國來芝加哥.
除了這個,Daniel也有一個網頁叫信仰百川的發起人.
其實都是他一個人做,大部分時間都是你一個人做.
如果你不知道什麼是信仰百川,請你去信仰百川看一些好的文章.
是一個很好的基督教平台.
有很多不同的聲音在那裡寫文章.
發表一些觀察或者神學的信念等等.
信仰百川小弟也是其中一個很懶的作者.
偶爾都會寫東西,只是很偶爾.
但是有很多其他作者寫了很多文章.
你能夠藉著這個這麼好的平台.
你可以裝備自己去認識社會,認識信仰.
今天張博士講的正道經文是馬可福音第二章二十三節.
有聖經的話請你打開,或者APP打開.
馬可福音第二章二十三節.
耶穌當安息日從墨地經過,他們途行路的時候,喝了墨水.
巴利賽人對耶穌說:看啊,他們在安息日為什麼做不可做的事呢?.
耶穌對他們說:經上記著大衛和跟隨他的人缺乏飢餓之時所做的事.
你們沒有念過嗎?.
他當亞比亞他作大祭司的時候,怎樣盡了聖殿神的殿.
吃了陳設餅,又給跟他的人吃.
這餅除了祭司以外,人都不可吃.
又對他們說:安息日是為人設立的,人不是為安息日設立的.
所以人子也是安息日的主.
我們歡迎丹尼爾.
在剛才大家讀的經文裡面有兩組人.
第一組人就是耶穌和他的跟隨者.
他們路經這個麥田,有些肚餓,就吃了麥水.
第二組人就是耶穌提到大衛和跟隨他的人.
這兩組人,大衛就進去拿陳設餅吃.
這兩組人都犯了一些,我們可以這樣說,是宗教規條.
但是我們所相信的主耶穌基督,他又認為違反這些規條並不是大問題.

$^{41}$這是很奇怪的.
耶穌主動告訴大家,你不用守的.
這樣也可以,究竟是什麼意思呢?.
我想在今天短短的時間,和大家一起思考聖經裡面這方面的道理.
我想提出三個要點,然後有些綜合的討論.
第一個要點就是,究竟他們違反了什麼呢?.
在座當中,如果看得仔細的朋友,可能你也會有個想法.
在安息日喝麥水,未必真的違反舊約聖經的律法.
因為那很可能是法利賽人想出來的人間遺傳.
這個說法,耶穌也有提過.
如果大家仔細找出埃及記第34章21節的話.
那裡提到,須在耕種收割的時候也要安息,那是十誡之一.
可能就是因為這條誡命,令到後來的法利賽人覺得.
有很多事情都不准做了,包括你喝麥水,充飢都不可以.
如果這樣說的話,可能耶穌也未必真的有犯到律法.
不過我們就不可以說戴偉沒有犯,因為戴偉犯的那件事.
是舊約裡面清楚說明不可以做的事.
不過這裡也需要解釋一點.
因為按照撒慕爾記上21章3至6節.
當時的大祭司對戴偉講的那番說話和耶穌這條福術有一點出入.
當時的大祭司說,只要男子沒有親近婦人就可以吃.
就是他加了一個要求,有一個condition.
如果你那些跟隨你的人,或者你近來沒有親近婦人,就可以吃那些餅.
但是其實當時的大祭司講出這番說話給一個condition.
戴偉和他的跟隨者,其實是他自己作出來的.
因為如果你看回利未記22章,那裡有提過.
只有亞倫的後裔要在一些結證了的條件之下才可以吃陳詞餅.
其他人不可以吃.
那位的祭司沒有理會戴偉和他的跟隨者是否亞倫後裔.
大部份都應該不是.
然後還將結證的要求重新詮釋.
只要你們的男子這幾天都沒有親近女性就可以了.
戴祭司是自己作條條例出來的,在那個場合裡面.
挺奇怪的,他可以這樣說,是因時制宜走過方便幫助戴偉和他的跟隨者.
這裡可以舉個例子,我初初來美國的時候,傻傻地開車到處走.
我都忘記了甚麼事.
然後去到一個地方突然肚子痛,想找個洗手間.
然後就見到有個pizza的餅店.
衝進去,最多吃完出來買個薄餅吃.
誰知原來那裡是delivery的.

$^{81}$根本不是公開讓人進去.
但店員就很好心說,總之你保持乾淨就行,你用吧.
其實都是走過方便,規矩上不是這樣做.
但他覺得你有需要就給你.
說起洗手間,大家可以回想香港人的遊行.
去到某些地方,去到一些教會.
有些教會就說不准你來用我們的洗手間.
但有很多教會現在都愈來愈願意打開門口.
歡迎遊行人士進去休息,用洗手間,充電等等.
這是第一點.
第二點就是,焦點要放在安息日上.
如果不是安息日的話,其實舊約有很多的明文的說法.
是說一般人都可以喝墨水食.
例如申命記23章25節.
他說道,你們盡了淪捨暫著的和家,可以用手摘水子.
只是不可以用鐮刀割取和家.
即是你吃少少都可以,但不要用鐮刀割,吃人家那麼大份.
因為都是人家好心給你吃,不要讓你那麼慘,那麼肚餓.
利未記第十九章第九至十節又這樣說.
在你們的地收割莊稼,不可割進田割,也不可拾取所遺落的.
不可摘進葡萄園的果子,也不可拾取葡萄園所吊的果子.
為什麼呢?因為要留給窮人和寄居的.
所以舊約的經文,舊約的律法本身是有它恩慈的一面.
即是社會有需要的人可以去你那裡拿些東西吃.
你不可以太絕情.
換言之其實下墨水本身不是很大問題.
但為什麼來到耶穌這個經驗的時候,反而觸怒了法利賽人呢?.
就因為那個日子是安息日,焦點就放在安息日上.
而且安息日是十誡裡面其中一個命令.
但奇怪在耶穌認為這個也不是很大問題.
雖然當時的法利賽人和猶太人不是這樣想.
舉個假設性的例子.
假如你是一個蛋糕店的老闆.
你平日都會將一些多出來的食物捐給有需要的人士.
不過有一天來拿那些多餘食物的運貨.
那個人就問可不可以星期日來拿呢?.
因為這個星期不方便過來.
假設星期六不方便過來.
你會不會因為那天是星期日.
你要回教會所以你不喜歡.

$^{121}$然後就不開門給他拿呢?.
還是你會盡量遷就到.
令到他可以做到一個送貨呢?.
究竟主日是否重要到一個地步很神聖.
然後我們不可以改時間.
一定要關門.
我回了教會我就不幫你.
是不是這樣呢?.
這個也可以思考一下.
第三個重點就是.
他犯了什麼規呢?原因在哪裡?.
令到耶穌一早有人犯規.
違反規條的.
是因為他們有衝擊的需要.
而不是因為他們說要安息.
所以我們不可以將這件事.
詮釋為有些宗教泛民辱節.
令人在安息日無法安息.
因此那些律法就不能夠達到他本身的目的.
那段經文就沒有這個意思.
那段經文只不過是說.
如果那些人有這個.
就是.
就是討福的需要.
就是很餓很累.
那你就吃點東西.
所以重點就應該放在那一句經文.
安息日是為人設立.
人就不是為安息日設立.
如果大家掀開聖經的話.
會發覺.
今天我們讀的是《馬學福音》第二章.
但其實在《馬太福音》有同一個記載.
也都說到耶穌和他的門徒去喝墨水.
不過在那邊的記載就沒有第27節.
就立即跳到第28節.
就說人子是安息日的主.
但如果你仔細想想背後的原因.
似乎人子是安息日的主.
不是最主要的原因.

$^{161}$如果你將焦點放在.
人子是安息日的主的話.
那個是耶穌.
他是天上地下的主宰.
他想怎樣都可以.
他要吃東西就吃東西.
安息日無所謂.
如果你將焦點放在第28節的話.
就很容易這樣想.
但其實第27節是很有意思的.
那裡說安息日是為人設立.
人就不是為安息日設立.
而且你想想大衛的故事.
你不可以說大衛是安息日的主.
或者是陳設餅的主.
如果我們要將大衛的故事.
拿來作為一個對比.
就說耶穌這樣做都沒有問題的話.
或者是做相似的事.
那個焦點就不是放在主角.
是否陳設餅的主.
或者是安息日的主.
而是安息日或者陳設餅等等.
是為人設立.
不是相反的.
這裡說的安息日為人設立.
是為.
就是單單是為人.
或者我們說是人性.
Humanity.
安息日為人設立.
人不是為安息日設立.
這個精神.
和我們平日對舊約律法的理解.
可能是有些出入的.
一般來說我們印象中.
對舊約律法的想法.
都是覺得是絕對黑白分明的.
甚至因為舊約裡面有些經文提到.
如果你不遵守這些的話.

$^{201}$你就會被石頭扔死.
多嚴重.
你當然絕對不敢越雷持半步.
但是一方面舊約裡面是有一些教導.
但是另一方面.
其實我們看清楚.
聖經裡面也有提到一些.
好像很模糊.
很有彈性的教訓.
例如這一節.
《箴言26章》.
這兩節很有趣.
「不要照愚昧人的愚妄說回答他.
恐怕你與他一樣」.
接著下一節.
「要照愚昧人的愚妄說回答他.
免得他以為有智慧」.
即是照不照愚昧人的愚妄說話呢?.
如果你很自免.
當是一個規條來遵守的話.
其實你看完這兩節.
你真的不知道該怎麼做.
為什麼聖經會有這樣的經文呢?.
究竟這兩節.
或者這一類經文的意思是什麼?.
這個可能值得我們想一想.
當然這個也不是很難解釋.
一般人的想法就是.
我們行事為人要有智慧.
要因時制宜.
不可以一成不變.
只是一條規則.
然後你就遵守.
這麼簡單.
另外也可以看看《傳道書》第七章.
第十六節.
「不要行義過分.
也不要過於自請智慧.
何必自取滅亡呢?」.
接著「不要行惡過分.

$^{241}$也不要為人愚昧.
何必不到其而死呢?」.
第十八節也很奇怪.
「你持守者為美.
哪個也不要鬆手.
因為敬畏神的人.
必從者兩樣出來」.
這幾節經文又有什麼意思呢?.
大概兩個星期前.
我在哥倫布斯那邊的華人教會.
回到團契的時候.
他們也討論過一個類似的問題.
就是《傳道書》.
《傳道書》好像很喜歡搞事.
總是說一些古古怪怪的東西.
《傳道書》裡面有提過.
你在神面前不要亂說話.
說話要謹慎.
否則就很糟糕.
那段查經的材料.
裡面有提到.
耶穌也有說.
你祈禱的時候不要重重複複.
因為你不應該以為重複.
就像愛邦人一樣.
重複多一點.
神就會順聽.
這是錯的.
但是那時候我就反問了一些問題.
我說聖經另外也有些教導.
就是你要向神傾心吐意.
將你所想的都告訴神.
而另外我們在教會的練習.
根本都是重複祈禱.
大家看看你手上的程序表的代禱事項.
你每個禮拜都重複為那些事祈禱.
究竟重複還是不重複呢?.
究竟應該多說話.
覺得像朋友一樣.
說什麼都可以.

$^{281}$還是要很謹慎.
好像你見一個大人物.
想好要說什麼.
說完兩句都不敢再說.
要閉嘴馬上走開.
什麼態度才是正確.
為什麼聖經有時候會教你這樣.
為什麼有時候會教你這樣呢?.
所以就算說到祈禱.
應該有什麼心態.
什麼祈禱的方式.
其實聖經裡面的說法.
都好像不是那麼清楚.
如果你單單覺得要遵守一些規條.
你覺得勝經有規條.
所以你要遵守.
用這樣的心態去看的話.
你就完全不明白那些經文.
想表達什麼意思.
是不是?.
他們那裡可能有些特別的意思.
當我們遵守.
看回《傳道書》第七章.
當我們遵守一些我們以為很重要.
很絕對的生活原則的時候.
我們有沒有想過.
會否行義過分呢?.
這樣說很容易令我們很驚恐.
是不是真的不用遵守呢?.
這個世界可以做什麼都可以.
做基督徒也可以隨便享受最終之樂.
是不是這個意思呢?.
那倒不是.
因為你如果享受最終之樂.
你就行義過分了.
下一節已經說了不行了.
那是怎樣呢?.
很煩惱.
做得太盡又不行.
你不做又不行.

$^{321}$第十八節也很有趣.
你持守者為美.
哪個也不如鬆手.
因為敬畏神的人在兩樣中間出來.
我們一般的想法就是.
守得越好,守得越嚴的人.
就越敬畏神.
但怎知道這裡說.
你守得這麼嚴的話.
你就行義過分.
真正敬畏神的人.
其實在兩樣中間出來.
這裡說到敬畏神的人.
如果大家想起的話.
曾言提到敬畏耶和華是智慧的開端.
我們可否將敬畏神.
那幾個字改為有智慧的人.
有智慧的人.
就會在這兩樣出來.
舉個例子.
我們覺得每個基督徒.
每一天都要靈修,祈禱.
每個星期都要回教會.
熱心遵守一切教訓.
這些不是壞事.
大家回教會這麼久都明白.
但是.
當你很努力遵守這一切的時候.
會否落入一種行義過分呢?.
最後和大家思考一段經文.
就是《馬太福音》第五章.
這個大家都很熟悉了.
「此事我告訴你們不要與惡人作對.
有人打你的右臉.
連左臉也轉過來由他打.
有人想要告你.
要拿你的女衣.
連外衣也由他拿去.
有人強迫你走一里路.
你就和他走二里.

$^{361}$有人求你的就給他.
有向你借貸的不可推辭.
你們聽見有話說.
幫愛你的鄰舍.
恨你的仇敵.
此事我告訴你們.
要愛你的仇敵.
為那逼迫你們的禱告.
這樣就可以作你們天父的兒子.
因為他叫日頭召好人也照歹人.
降雨給異人也給不異的人.
你們若單愛那愛你們的人.
有甚麼賞賜呢?.
就是遂理不也是這樣行嗎?.
你們若單請你弟兄的安.
給人有甚麼長處呢?.
就是愛幫人不也是這樣行嗎?.
所以你們要完全.
像你們的天父完全一樣」.
這段經文大家應該非常熟悉.
尤其是打左臉打右臉的.
我們經常說的.
但其實你有沒有可能.
全面遵守這段經文說的話呢?.
沒甚麼可能的.
如果你回想香港抗爭的場面.
如果警察打你.
你還要站著讓他打多兩拳嗎?.
打你兩拳就死定了.
前陣子有段影片.
在一條隧道上.
有個胖子在打.
另一個人站著讓他打.
雖然你覺得他很有偉大的犧牲精神.
不過一般的小朋友.
在家裡也不要亂學.
這樣打法是很危險的.
這樣打法真的會死人的.
就算不死也會有永久傷害.
後來有人說被打的人也是練過武的.

$^{401}$這樣他才能夠撐那麼多拳.
不然他也會暈倒很久.
剛才讀的那段經文也提到.
有人問你借錢.
你就借給他.
一有人問就要借.
不是應該多想一點嗎?.
如果你想深一層的話.
其實這段經文說的.
你要怎樣做.
是很難按照字面這樣做的.
那個難的原因不是因為你不想跟隨神.
不想做一個好基督徒.
而是本身好像不太合理.
究竟這裡應該怎樣去思考.
應該怎樣做呢.
有人強迫你走一里路.
你就跟他走兩里.
求你的就急急他.
借太的不可推遲.
究竟這裡耶穌想表達什麼.
認為我們應該做一個怎樣的基督徒.
會不會這麼簡單在字面上所說的.
我們做完就對了呢.
還是我們應該想想背後有什麼意思.
來到這裡的時候.
開始大家可能有點擔憂.
我們平日讀信的一些生活行為原則.
是不是都要放棄呢.
是不是應該間中守一下就可以呢.
不過我想重點不是你可以放棄或不遵守.
而是我們用什麼心態去遵守.
我們是否明白和願意.
有時候應該怎樣做怎樣做.
這個就回到剛才那段經文的主題.
安息日是為人設立.
人就不是為安息日設立.
這裡給我們一些啟迪.
耶穌沒有建議我們胡亂遵守.
你守也行不守也行.

$^{441}$看什麼情況也行.
有時候也行.
他的意思不是這樣.
他將重點放在為人設立這個概念上.
換言之有很多的生活行為原則或規條.
其實是為了我們成長得更加好.
做一個更好的人.
做一個合神心意的人.
所以才設立.
正因如此如果有些規條或原則.
慢慢發展成為一些壓抑了人性的.
壓抑了你生命成長.
靈性成長的一些事情的時候.
我們有時候要學習不要死守.
要將焦點放在人的生命成長上.
我這裡舉些例子.
直到今時今日香港仍然有少部分的教會或中學.
會說中學生時代談戀愛是不對的.
你要留意.
他不僅不建議你這樣做.
他甚至說你是做錯事.
錯到一個地步.
有些中學仍然會當你犯校規.
記你一個小過.
中學生談戀愛你會犯小過.
被學校懲罰.
有些教會因為你拍拖.
然後又會懲罰你.
這些會不會就是行義過分呢?.
將一些原意美好的建議.
變成一些規條.
然後太過黑白分明地實施.
然後有些人遵守不到.
雖然不像有什麼大錯.
但你都要照樣罰他.
這個會不會就是行義過分呢?.
又或者有些信徒對於不可停止聚會.
這個說法是很強烈地遵守.
是否真的有必要呢?.
例如去年在香港山竹的颱風.

$^{481}$吹得很嚴重.
你看這些影片都看到.
擺得很嚴重.
杏花村的海水都撞進了地牢和一樓的房間.
在這樣的情況下.
是否應該回教會呢?.
是星期日.
不可停止聚會.
是否我們死都要回去呢?.
還是大家為了生命安全.
可以跳過這一天.
再說遠一點的就是.
政治制度.
基督徒都會嘗試討論一些政治倫理.
社會的願景.
究竟基督徒.
究竟一個社會怎樣才合神心意.
怎樣才比較好.
有些人會討論民主制度有多重要.
有很多不同的制度.
哪些制度.
但其中一個我們應該留意的焦點.
哪些制度才是對人有足夠的尊重.
對人的尊嚴有足夠的保障呢?.
如果我們將焦點放在這裡的時候.
就比較容易看得到.
哪些比較好.
哪些比較差.
如果純粹想想經濟發展.
賺到很多錢.
喜歡吃什麼.
有很多物質享受.
可能你最後的判斷.
哪個制度好.
哪個制度不好呢?.
結果就不是很相同了.
現在到最後.
讓我去總結.
今天這半個小時.
跟大家說了什麼.

$^{521}$我想提出的就是.
很多時候.
我們的生活行為原則或規條.
我們遵守是好事.
我們在裡面有成長.
但是我們要留意.
人不是為那些規條去設立.
人不是為安息日設立.
而是相反.
安息日為人設立.
所以如果有些原則.
有些規條.
發展到一個地步.
窒息了我們的成長.
我們的人性的話.
可能我們就不需要遵守得那麼死板.
這個說法好像很危險.
但是如果你明白背後的道理.
其實是很合理的.
就算安息日那麼重要.
在猶太人來說.
在舊約聖經來說.
這是十誡之一.
耶穌也說.
其實安息日是為人設立.
人不是為安息日設立.
讓我最後用一個比喻來總結.
聖經裡面有很多教導和督察.
詩篇提到.
你的話是我腳前的燈.
是我腦上的光.
沒錯的.
但是如果你在一個很黑暗的地方.
例如你去行山.
天開始黑.
你開燈.
開著電筒來照一照.
你的焦點是看向哪裡呢?.
你不是看著燈.
不是看著電筒.

$^{561}$你是看著路.
你是四處看.
你最後的目的地在哪個方向.
然後你按照燈所照亮的路.
一直走向那個方向.
我這裡想提出的是.
明燈本身不是我們終極要去的地方.
很多的教導,規矩.
甚至很多制度.
其實本身的價值.
是指向我們走向一個神要我們走向的方向.
我們的焦點不是單單看著那盞燈.
覺得好像找到真理.
我今天說的就這麼多.
多謝各位.
多謝收看成匯網台.
記住like,comment或訂閱我們.
\newpage



\section{}
\label{sec:PGikCuC4cZg}
\textbf{1/27/19 主日講道: 創造與敬拜 (張麗仙傳道)}
\newline
\newline
連結: \href{https://youtube.com/watch?v=PGikCuC4cZg}{\texttt{https://youtube.com/watch?v=PGikCuC4cZg}} ~~~~ 語音日期: 2019-08-02
\newline
\newline
\hyperref[sec:ALcJ_99Kiw8]{\small{< < < PREV SERMON < < <}}
~
\hyperref[sec:index_chronic]{\small{[返順時目]}}
~
\hyperref[sec:index_scriptual]{\small{[返順卷目]}}
~
\hyperref[sec:wT2Sf41UjYU]{\small{> > > NEXT SERMON > > >}}
\newline
\newline
$^{1}$這樣好看一點吧.
各位頂尖的朋友.
看到你們都很開心.
我有兩個問題想問你們.
就是你今早起床的時候.
回到教會的時候.
你會想到什麼呢?.
有什麼在你腦海裡出現呢?.
這是第一個問題.
第二個問題就是.
你進到禮堂這裡崇拜的時候.
你又會想到什麼呢?.
在你腦海裡又在流行著什麼呢?.
為什麼我要問這兩個問題呢?.
我暫時不說.
不過我一定會告訴你們為什麼.
但是我想跟大家說.
今天我們要先看經文.
我今天要跟大家看的聖經.
是初代世紀第一章一到二節.
我想很多弟兄姊妹.
馬上會心微笑.
不是吧?聽了這麼多年道.
讀了這麼多經.
又重頭來過?.
神創造天地.
沒錯.
各位弟兄姊妹.
我們不可以失去眼光.
就是我們的信仰的根基.
是神的創造.
祂是我們的創造主.
這是我們信仰的根基.
而創世記第一章一到二節.
這一節聖經.
我們以為我們非常熟練.
倒轉頭也會背.
但是.
或許我們不是想像中這麼熟悉這節聖經.
起初神創造天地.

$^{41}$是空虛混沌.
冤冤黑暗.
神的靈運行在水面上.
這一節聖經.
這兩節聖經.
歷史歷代很多聖經學者.
有很多不同的研究.
不同的看法.
有很多的討論.
很多的爭辯.
可以說很久很久.
不過暫時我們不要說學者.
不要說聖經學者.
我告訴你.
我小時候怎麼看這節聖經呢.
我看到這兩節聖經的時候.
我就想像.
有一天.
神突然間心血來潮.
來到一個慘不忍睹.
黑漆漆的爛地上.
神就像一隻鷹一樣.
又莊嚴又優雅.
運行在水面上.
然後呢.
神開口.
就一樣一樣一樣.
好像迪士尼卡通片一樣.
好像神仙棒一樣一樣.
天地都做了.
整個漆黑的爛地.
變成彩色繽紛的世界.
這個是我小時候.
了解的創世記.
了解的這句話.
我不知道跟你對於創世記的了解.
版本像不像呢.
無論像不像也好.
我相信我和大家一樣.
我們都知道.

$^{81}$或者我們都這樣了解.
神在一片荒蕪那裡.
由無創造有.
這個是無可置疑的.
但是.
究竟.
是不是由無創造有.
但是這節聖經.
有沒有說多一點的東西給我們知道呢.
我們要承認一件事.
聖經對神六天的創造.
都有仔細的描述.
但是.
對於創造的背景.
的記載是很有限的.
我們所知的.
也是非常有限的.
就說.
起初.
起初是什麼時候呢.
你有沒有想過是多少年前.
神創造世界的呢.
你想不想知道呢.
很多聖經學者都想知道.
盡量去研究.
但是其實我們是不知道的.
創世之前.
發生什麼事呢.
創世的時候又發生什麼事呢.
老實說因為聖經沒有記載.
我們不要太勉強去掘.
因為可能會越掘越錯.
但是我們又再看回.
創世的第一第二節那裡.
簡單的幾句.
地是空虛混沌.
冤冤黑暗.
神的靈運行在水面上.
這裡.
其實有沒有些東西是我們可以知道多一點的呢.

$^{121}$我最近讀到一本.
一個聖經學者的解釋.
他的研究.
他是從近東文化和聖經的用詞去解釋.
特別是第二節這裡.
我覺得他的解釋.
對我對創世的時候.
有什麼事發生.
是有很大的幫助.
我們看.
地是空虛混沌.
冤冤黑暗.
神的靈運行在水面上.
幾筆很簡單的去形容創世的時候是怎樣的.
我們可以了解.
神當時面對的是三個問題.
面對的是空虛混沌.
空虛混沌是什麼意思呢.
就是沒有生命.
和非常非常混亂.
超混亂.
怎樣混亂我也不知道.
總而言之就是很混亂.
第二個問題就是冤冤.
冤冤的意思其實就是大水包圍著.
大水在近東文化的了解.
是指邪惡的意思.
我們都知道.
大水是指邪惡的意思.
大水包圍就是說.
當時神面對的就是.
創世的時候有很多邪惡包圍著.
然後黑暗.
同樣近東文化的了解.
這個黑暗的字.
不是只是關了燈就黑了的意思.
而是黑暗代表邪惡的勢力存在.
可以說.
空虛混沌冤冤黑暗.
是意味著創世的時候.

$^{161}$是超級混亂.
沒有生命.
卻有邪惡的勢力存在.
但是.
神的靈運行在水面上.
改變了一切.
神的靈運行在水面上.
改變了一切.
我們看看神的靈這個字.
神的靈這個靈字.
其實在聖經裡跟風字是同一個字.
而且這個字.
不是指微風.
溫柔的風.
而是狂風巨風.
一個最好的參考.
就是以西絕書第37章.
那裡說到.
神將他的靈像風一樣吹進骷髏骨頭裡.
叫他活起來成為極大的軍隊.
去幫助他的子民作戰.
那個風就給我們一個很好的參考.
到底創世記這裡的風.
是什麼意思呢.
是狂風巨風.
是神跟邪惡的勢力作戰的意思.
如果我們再看下去的時候.
我們會更加覺得.
這個可能是很有可能的.
接著我們看運行這個字.
神的靈運行在水面上.
運行這個字.
這個希伯來文.
在舊約聖經裡出現過兩次.
創世記這裡一次.
另外一度.
就是在新靈記第32篇11節.
那裡記載到神對他的子民.
就好像魚鷹攪動茶鍋.
在初鷹以上兩翼展線.

$^{201}$接取初鷹.
背在兩翼之上.
什麼意思呢.
原來在這次聖經裡.
是描述神對他的子民.
母鷹攪動茶鍋.
到底是一回什麼事呢.
原來母鷹很生氣.
因為她的小鷹在茶鍋裡.
懶惰,吃慣,坐慣,不學飛.
因此母鷹就要.
你不飛,怎麼做鷹啊.
你不做事,你怎麼做人啊.
母鷹就要幫小鷹飛.
她就要攪動茶鍋.
首先我們要知道.
鷹巢往往在懸崖那裡.
母鷹攪動茶鍋的意思.
是因為想你飛.
所以鷹的翼很大.
好像直升機一樣.
在你家裡.
你猜你家裡會怎樣.
如果有直升機在你家裡.
會把東西吹亂.
小鷹當時就這樣.
因為母鷹在挑戰他們.
把茶的東西吹亂.
他們就像地震一樣.
怎麼辦.
正在想怎麼辦.
都未止.
母鷹飛到懸崖上.
做一個像自殺式的動作.
放下來.
向鷹巢襲擊.
這個是什麼意思呢.
就是讓小鷹覺得.
好像被母鷹殺了一樣.
嚇得他們脫毛脫翼.

$^{241}$這個時候他們會怎樣做呢.
會逃跑嗎.
讓他們自殺嗎.
所以小鷹就會飛出茶外.
學飛.
懶得死也要飛.
在這個時候.
母鷹就在他們底下.
伸開他們的翼.
接住他們.
如果他們飛不動.
跌倒的話.
接住他們.
讓他們繼續飛.
這個比喻是什麼意思呢.
母鷹對小鷹的期望.
希望他們學飛.
也代表了神對他們的指紋.
你們的懶惰.
你們的叛逆.
我要像母鷹攪動茶窩一樣.
來教訓你們.
讓你們生生勝勝.
這個就是這裡的意思.
說回頭.
這個攪動.
原來是和創世記.
第一章二節的運行.
是同一個字.
那你怎樣看神運行在水面上呢.
是不是像我小時候這樣想.
像鷹一樣很優雅.
很莊嚴地飛.
這麼簡單呢.
我們可以看到.
原來神運行在水面上.
就像母鷹攪動茶窩一樣.
在這裡撲擊.
甚至在這裡作戰.
原來這裡給我們看到的.

$^{281}$是一幅神的靈與邪惡.
爭鬥的一幅圖畫.
也可以這麼說.
原來在創世記未記載神六天的創造之先.
他簡單地說了一句話.
地是空虛.
雲頓.
淵面黑暗.
神的靈運行在水面上.
簡單的幾筆.
告訴我們.
原來創造的時候.
神是面對一場.
與邪惡勢力的爭戰.
這是我最近讀到的一個了解.
我覺得這個解釋.
對我來說是有點新鮮.
或者對你來說.
是有點大膽.
但是當我們繼續看創世記.
神一二三四五六一的創造的時候.
我們會發覺原來這裡的解法.
神在作戰破擊.
不是沒有根據的.
我們首先看.
我們說神當時面對三個局面.
第一就是黑暗.
針對黑暗.
神做了什麼呢.
我們知道第一天.
神做了光.
光衝破了黑暗.
這個時候.
有光有暗.
分為晝和夜.
而晝夜.
也就是有白天和黑夜.
但是這個夜晚.
與之前的黑暗是不同的.
希伯來文有一個很特別的地方.

$^{321}$其實中國字也有點像.
中國人知道形聲字.
水流流.
什麼濕濕的.
這些形聲字.
原來希伯來文.
黑暗這個字.
讀成粗錫.
讀起來覺得怎麼樣.
很刺激.
很刺激的.
代表什麼.
代表了邪惡的勢力.
希伯來文是用粗錫.
在創世的意思是.
冤冤黑暗就是這個字.
但是當神做了光.
有晝夜.
有光和暗.
夜晚的時候.
用了另外一個字.
不是粗錫.
是拉拿.
拉拿讀起來怎麼樣.
很溫柔.
很通順.
很舒服.
由粗錫變成拉拿.
這個黑暗.
神變成了一個夜晚.
安詳的夜晚.
這是神的創造.
第一天.
什麼意思.
神的創造.
徹底改變了.
當時的黑暗局面.
接著我們又知道.
第四天.
神創造光體.

$^{361}$大陽,月亮.
還有眾星神.
一個什麼樣的局面.
就是黑暗的權勢.
已經被控制.
一切變得有秩序.
不再受黑暗的控制.
接著.
冤冤.
我之前也說過.
在近東文化.
冤冤.
這裡是指大水的意思.
近東文化.
大水是指邪惡勢力的意思.
第二天神做了什麼.
針對冤冤.
做了空氣.
使得以上是天.
下面是海.
不提醒創世記也不知道.
原來是這樣.
神用空氣.
來改變了.
大水包圍的邪惡勢力.
空氣以上為天.
就是我們了解的窮倉.
空氣以下的.
是我們現在看到的海.
神的創造.
使得有天有海.
意味著大水邪惡的.
充滿的局面.
神又徹底改變了.
然後第五天.
我們看到.
神做水中的活物.
和天空的飛鳥.
有天了,所以做天空的飛鳥.
有水了,做水中的活物.

$^{401}$一切各從其類.
變得有秩序.
接著針對第三天.
空虛混沌.
第三天.
神將水最後一處.
洗漢地露出來.
原本只有大海.
現在有地了.
還有神的計劃繼續進行.
神在地上.
做出各樣的植物.
地有生命.
也可以說.
這些生命是一個伏筆.
因為在第六天.
神創造更多的活物.
地上有活物.
最後在神創造的活物裡.
神創造了人.
神創造了活物.
神創造了人.
這是神要解決空虛混沌的局面.
將它徹底改變.
空虛變成有生命.
混沌變成一切都有秩序.
各從其類.
而且神創造了人.
是要人管理萬物.
就是一切有秩序.
說到這裡我要再重複一下.
就是我要強調.
創世的背景是怎樣.
因為聖經沒有記載.
如果我們太勉強去解釋.
有時候會越解越錯.
但是我覺得以上的解釋.
是有它的根據.
也是用創造的時候那種.
淤棉黑暗空虛混沌的情景.

$^{441}$是對應了神.
接著所創造的六天.
怎樣改變了這個情景.
而在這個改變裡面.
我們看到.
神說一切都是好的.
因為神是將邪惡的勢力征服了.
神作了戰爭.
打了一場仗.
將邪惡的勢力.
即是黑暗和大水控制住.
這是當時神所做的第一件事.
第二,祂創造了生命和秩序.
空虛混沌不再有.
淤棉黑暗不再有.
我們的神勝利了.
這是創造.
是我們的創造主所成就的.
當然.
我們也很熟悉接下來的聖經.
因為創世記告訴我們.
接下來有很曲折的發展.
這個發展可以說是.
關鍵在第六天創造了人.
我相信創造人對邪惡的勢力.
是最大的打擊.
想想,魔鬼為什麼不去試探那些馬.
獅子,老虎.
跑來試探人呢?.
因為聖經告訴我們很清楚.
神安著他的形象造人.
而且他要派人有一個很榮耀的使命.
去管理他所創造的宇宙.
所以魔鬼不可以放過.
去試探人.
他要令他們墮落.
要他們灌溉.
要他們返回神所創造的美善.
而我們知道很不幸的是.
人,阿當,夏娃.

$^{481}$在這個試探中失敗了.
墮落了.
整個人類又陷入空虛混沌.
冤冤黑暗的局面.
是不是?.
各位弟兄姊妹.
你過去一個星期過得怎麼樣?.
在你的生命裡.
有沒有被疾病恐懼.
窒到欺詐.
或者背叛所傷害?.
我們的世界又怎麼樣呢?.
貧窮,貪污,強姦,醉酒,毒品,暴政,廝殺.
什麼時候停過呢?.
你自己又怎麼樣呢?.
在醉的試探裡.
你有多少跌滑過呢?.
你不想做的事,你做了.
你想做的好事,你做不到.
再看我們的城市.
芝加哥,多少槍殺.
多少搶劫.
搶錢包,搶車.
連加拿大羽絨也搶了.
這些是冤冤黑暗,空虛混沌.
是因為人在試探中墮落了.
但是各位弟兄姊妹.
不要忘記.
這些空虛混沌,冤冤黑暗是暫時的.
因為在人墮落的一刻.
神,上帝已經立刻預言拯救.
而在歷史上.
我們的確看到.
神是不斷地為我們作戰.
完成了拯救.
為什麼耶穌基督的福音是這麼重要?.
為什麼對你和我是這麼重要?.
對這個世界是這麼重要?.
因為藉著耶穌.
神戰勝了邪惡的勢力.

$^{521}$恢復了祂創造的美善.
因為藉著耶穌.
神已經得到永遠的勝利.
開始的時候我說.
我們的信仰不能夠失去一個眼光.
就是神的創造.
是我們信仰的根基.
我們這麼重視耶穌基督的福音.
是因為我們知道.
藉著耶穌基督的福音.
可以將我們殘缺墮落的世界.
帶回創造主那裡.
我們也希望.
有一天不會再淵源黑暗空虛混沌.
我要問大家兩個問題.
我要告訴大家為什麼要問你們兩個問題.
就是今早準備來會議的時候.
你們心裡在想什麼?.
今早這麼冷.
當然有很多事情要想.
想想能不能出來.
車子能不能出.
穿什麼衣服.
吃什麼早餐.
下午吃什麼.
去哪裡吃東西.
還有現在快過年了.
我們要準備什麼.
下雪下成這樣.
很麻煩.
很多事情要想.
我想問.
在你們的腦海裡.
有沒有想到.
原來你們是來教會敬拜神?.
神有沒有出現在你們的想像裡?.
這是第一.
第二就是.
你進入這個禮堂的時候.
你想到什麼?.

$^{561}$當你聽到敬拜李真偉大爺的時候.
你有沒有全情投入去唱.
去歌頌這位創造我們的神?.
我們有沒有將當得的榮耀給祂?.
在這裡我想分享一個.
我小時候很誇張的事情.
給大家聽.
十一,二歲就回教會.
我的做法就是.
唱完歌之後.
台上唱完歌.
儀式完結之後.
坐下來.
牧師一開口.
就像我今天一開口.
我就說話了.
我就聊天了.
和我旁邊的人聊天.
不管是誰.
我總是有很多話要說.
不過我講話都很小聲.
小聲到有時候.
我覺得不會影響到別人.
所以我一直說.
可以這樣說.
牧師說多久.
我就說多久.
OK.
講到我前面的大人.
看過來.
有些就笑一笑.
我說.
過了一分鐘.
我又講又講又講.
我都不知道要說什麼.
接著.
差不多十幾分鐘.
大人又來.
噓.
再兇一點.

$^{601}$接著我又停一會.
又繼續講又講.
接著大人又看過來.
不是很溫柔.
眉頭皺了.
又沒有說話.
就是這樣.
我不知道為什麼.
我是這樣.
牧師說話.
我就不停說.
我自己覺得.
我都不明白.
為什麼他這麼不濟.
為什麼他這麼不服氣.
因為我講話都很小聲.
影響不到別人.
無法吵到別人.
牧師都聽不到.
牧師都沒空聽我說話.
我心裡是這樣想的.
但是.
後來我才發覺.
原來不是.
原來我一說話.
多小聲都好.
是吵到別人的.
是影響到別人的.
為什麼我要講這件事.
我昨天才想起這些.
打算揭開給大家看.
看看我有什麼樣的人.
為什麼要這樣講.
因為各位電視節目朋友.
我以下想給大家一支提醒.
提醒之前.
我都想告訴你.
你不用太feel bad.
當我提醒你.
因為我都衰過很多.

$^{641}$我想提醒大家兩件事.
第一件事就是.
當我們敬拜的時候.
我們不要被別的東西.
做了我們的主.
我看到今天我們有兩個很大的時代問題.
第一就是.
當然大家都明白.
手機的問題.
不是聊天的問題.
是手機的問題.
因為你在手機上.
老實說可以慢聊天.
我們用手機很方便.
手機實在太可愛.
但手機帶到崇拜裡面.
你都忍受不住用.
這就不是一個好的現象.
如果你一直聽到.
唱歌,祈禱的時候.
你都可以走開.
因為手機.
你不在網上.
或者社交媒體.
這絕對是一個問題.
這絕對是一個影響.
影響你自己.
影響你身邊的人.
影響你的敬拜.
我甚至曾經見過.
我派主餐去某些座位的時候.
我不是說現在.
在領主餐.
有人在按手機.
我心想.
你知不知道你現在在做什麼.
你正在領主的身體和步行.
你還在按手機上網.
有時候有些事情.
你做了之後.

$^{681}$你自己是不知道的.
我只有一句話.
奉勸大家戒掉.
戒掉這件事.
如果我說中了.
你千萬不要以為我真的針對你.
而是這些事情.
是我們大家互相提醒.
我實在不是針對任何一個人.
我奉勸你要戒掉.
其實很簡單.
如果你真的需要手機.
拿出來讀聖經.
你按飛行模式.
這麼簡單.
你也知道.
但是如果你又發覺.
你拿手機出來不是看聖經.
但你又忍不住上網.
那你就不要拿出來.
收起來可以嗎.
收起來.
崇拜之後拿出來.
歸根究底.
我知道這不是手機的事.
這是人的事.
是用手機的人忘記了.
尊榮我們的神.
忘記了這是一個崇拜的時間.
我們來敬拜祖我們的主.
我願意在這裡.
向大家提醒.
勸勉.
第二個時代問題就是.
我覺得.
現代的崇拜.
有很多科技.
很精彩.
很多影音.
有時會令我們覺得.

$^{721}$來崇拜好像看演唱會.
你有興趣的部分.
你看演唱會是怎樣.
有興趣的部分.
你會留心或留在這裡.
沒興趣的部分.
你會走出去.
或者人在心不在.
這是一個時代問題.
現在崇拜太精彩了.
但我們要明白.
原來崇拜的每一部分.
都有意思的.
崇拜的每一部分都有意思的.
我們在崇拜.
我們怎樣呢.
崇拜的時候.
唱詩我們稱謝神的作為.
禱告我們向祂讚美呼求.
聽到的時候我們更加認識祂.
奉獻我們是值金錢.
甚至說是尊婆在這裡報告.
原來都是崇拜的一部分.
是要我們關心教會.
和弟兄姊妹.
這是敬拜.
每一個環節.
其實都重要.
都是值得.
都是需要你的全心全意在這裡.
各位弟兄姊妹.
願意我們.
再一次被今天.
創世記創世的故事提起.
提醒我們.
我們的神是一個怎樣為我們作戰.
拯救我們創造我們的神.
讓我們省察.
我們怎樣去敬拜.
這位偉大又慈愛的神.

$^{761}$願意我們在這裡.
坐在這裡.
我們成會教會.
去做一個尊榮神的敬拜.
我們願意神會覺得.
I am proud of you.
You are the church.
因為你們是很尊崇我.
在這裡敬拜.
願意在這個禮堂裡.
從左到右.
從前一排到最後一排.
我們大家.
都是全心全意的尊崇.
全心的敬拜.
將榮耀歸給我們的主.
我們的創造主.
我們的救贖主.
我們的神.
願意我們坐在這裡.
使每一位都看到.
原來我們的敬拜.
是這麼有意思.
因為我們是將榮耀.
歸給這位創造的主.
最近.
我正在讀一本Tim Keller.
有關於真言靈修的書.
我都知道真言是講智慧的.
Tim Keller 講了兩個.
其中很大的智慧.
我很想在今天.
崇拜結束的時候.
送給大家.
特別是.
如果你是今天.
或者是最近第一次來教會.
是最近才開始認識神.
我以下兩個智慧.
是特別送給你的.

$^{801}$第一個智慧就是.
Tim Keller 提到說.
我們要明白.
我們正在敬拜的.
不是任何一位菩薩或神明.
我們敬拜的是一位.
創造我們的神.
沒錯.
今天創世記告訴我們.
神是戰勝邪惡.
創造萬物和利我的主.
我們需要謙卑地敬拜祂.
我們就必能夠經歷到.
祂帶領我們進入.
面對星光.
我們人生一切的空虛混沌.
冤冤黑暗.
始於我們敬畏.
我們的敬拜創造主器.
我們敬拜的不是任何一位菩薩.
是創造我們的神.
第二.
我們在這個世界上.
有很多關係.
Relations.
但是沒有一種關係.
能夠取代.
我們和創造主的關係.
想想這是很大的智慧.
我們這個世界上.
有很多關係.
父母兄弟姊妹.
夫妻子女.
朋友.
甚至是教會.
請記住.
這一切關係.
不能夠取代.
你和創造主的關係.
相反.

$^{841}$只有你重視.
你和創造主的關係.
你敬畏祂.
你才有智慧.
去珍惜和擁抱.
你人生一切的關係.
這兩樣是我.
在結束的時候.
特別送給大家.
我們一起來到神面前祈禱.
我們的創造主.
我們的神.
我們不知道.
因為很多話都不能夠表達.
但是容讓我們說一句.
感激你做了我們.
我們的生命.
我們的氣息.
如果不是在你的恩手裡.
根本不會有存在我們.
主願意今天.
我們讀《創世記》聖經的時候.
再一次地明白你.
再一次地將我們的心.
擺在你的面前.
獻上給你.
願意我們每一次在這崇拜的時候.
我們給你的是.
你當得的榮耀.
你當得的敬拜.
幫助我們成為教會.
成為一個.
我們每一個人都會.
Be proud of.
你更加會Be proud of.
我們的敬拜.
幫助我們.
我知道我們不會一天就做到perfect.
但是求主你幫助我們.
有這樣的心去追求.

$^{881}$去建立一個這樣的敬拜.
一切的關係.
不能夠取代你和我們的關係.
求你教我們珍惜.
你和我們的關係.
你給我們一切的恩典.
你不是普通的神明.
你是創造天地宇宙萬物的真神.
奉救主耶穌的聖名.
阿們.
\newpage



\section{}
\label{sec:wT2Sf41UjYU}
\textbf{10/21/18 主日講道: 三件要事 (吳君才牧師)}
\newline
\newline
連結: \href{https://youtube.com/watch?v=wT2Sf41UjYU}{\texttt{https://youtube.com/watch?v=wT2Sf41UjYU}} ~~~~ 語音日期: 2019-08-03
\newline
\newline
\hyperref[sec:PGikCuC4cZg]{\small{< < < PREV SERMON < < <}}
~
\hyperref[sec:index_chronic]{\small{[返順時目]}}
~
\hyperref[sec:index_scriptual]{\small{[返順卷目]}}
~
\hyperref[sec:NJ2Erozw8Ss]{\small{> > > NEXT SERMON > > >}}
\newline
\newline
$^{1}$主席.
今日的主要崇拜,現在剩下大約二十多分鐘,所以我會很簡單地提出三件事,特別是針對做核心領袖的,或者預備選核心領袖的,甚至每一個基督徒都應該要著意的三件事..
我選的經文是希伯來書一章五到十一節,我們一起讀一讀..
弟兄讀單數,姊妹讀雙數..
所有的天使,神從來對那一個說,你是我的兒子,我今日生你,又指著哪一個說,我要作他的父,他要作我的子..
(主席再讀).
弟兄,論到使者有說,神以風為使者,以火焰為伏翼.姊妹..
(主席再讀).
弟兄,你喜愛公義,恨惡所以神,就是神,用喜內柔告你,勝過告你的同伴..
(主席再讀).
弟兄,天地要長存,你卻要長存,天地都要仗衣服,漸漸夠了..
(主席再讀).
我們低頭有個簡單的劇主,但你仍是誰顧我們每一個的人.主,你提醒我們在你面前的,如何更多依靠你,更多榮耀你,更多經歷神你自己的恩典..
我們今日聽道的時間雖短,但求你幫助我們有一個謙卑,信服和願意受教的心..
祈禱奉主耶穌基督的名,阿們..
我提出三件要事,第一件就是要尊敬主..
在第五到第六節,你會見到有關的經文..
現代中文,日本語都最清楚,不過時間關係不講..
我們可以看到,天使雖然被稱為神的眾子,但父神唯獨對耶穌基督說:你是我的兒子,我今日生你..
我們見到耶穌基督的位格和神聖,他與父神同樣的尊貴,是超過天使等級的..
我們又看到天使都要拜他,敬畏他..
在第六節提到長子,你會見到,史圖保羅用手生的,源始的,再許多弟兄中作長子來形容耶穌基督..
一切信耶穌的人都得到神兒子的名份,是嗎?.
耶穌基督不單是神的獨生子,更成為長子,成為一切相信耶穌救贖的人的手生,源始或長子..
而這位父神的長子,愛子,連天上所有的使者,天使都要敬他,拜他..
所以耶穌基督的超越性和獨特性,是有尊貴的得勝者..
所以第六節提到,第七節,連風和火焰都要聽他的話..
弟兄姊妹,我們又檢視一下,這位天上的父神,直接被稱為神的愛子..
他連萬有的父,都看重他..
特別欽點他的耶穌基督,我們抱著什麼態度對他呢?.
連上帝所有的天使都要拜他,尊貴他,榮耀他..
但是相對來說,今時今日我們對耶穌的表現是怎樣呢?.
何時我們當主耶穌是主呢?.
在我們日常生活當中,我們有沒有真的尊崇主,耶穌是主呢?.
我們拜他,敬他,那種心態表現又是怎樣呢?.
我們常說要用心靈誠實來拜他,敬畏他..
但是在日常生活中,甚至在事奉中,我們敬畏,尊主為大的程度又是怎樣呢?.
這正是我們值得再想一想..
特別我們說我們要作教會的核心領袖,.
我們將要選出事奉主的人,.

$^{41}$甚至你說我自己只是一個很平凡的平信徒,.
但是這裡提醒我們,敬畏主,敬虔的去愛神,.
是一個屬靈生命很基本的態度和心智..
我們一起彼此面對..
曾經有一個牧師面對一個教友,.
被他所說和表現的東西搞得啼笑皆非..
這位教友以前是一個好賭的人,是一個賭徒,.
日夜沉迷賭博,但後來他信了耶穌,.
他說我信了耶穌,我要變,我要尊敬主,我要放棄賭博了..
但是他不斷說,他不是很徹底,.
他仍然留地步給魔鬼,反而狡詐的對牧師說,.
這副啤牌跟了我很多年,我一定要保留它,.
其實這副啤牌可以見到神的恩典,.
當我見到Ace的時候,我記起神是最大,全能獨一的真神,.
二字牌的時候,我想到神起初做亞當,夏娃,建立二人世界,.
三牌令我想起給了我三次不認主,我要更加警醒,.
四牌使徒保羅念同去羅馬宣教的路程,.
有大大話話四次,提起我要將福音遍傳耶路撒冷,.
猶大傳遞撒瑪利亞直到地極,我要做一個傳福音的使者,為主作見證..
五牌,令我想起利未記,那五種制字,他很厲害,.
哪五種制字?利未記的五種制字,.
凡制,數制,平安制,贖罪制,贖軒制,很多人都不懂,.
但他講到,這五種制字提醒我要聖潔的過日子,.
啤牌的J,思洗藥巷是主耶穌的開路先鋒,.
Queen,想起瑪利亞的信服,King,想起萬王之王,萬主之主,和平之君,.
啤牌第一隻是Ace,耶穌是首先的,最後一隻是King,耶穌是末後的,.
心,Heart,提醒我要愛人,用心,.
愧是,好像一把銼,一個鏟,提醒我為主努力作功,.
階磚是鑽石,提醒我的信心要像鑽石一樣堅硬,.
他解釋得好像是頭頭是道,這位完全未解道的教友,.
滔滔不絕發表偉論,似是而非,滿口屬靈的名詞,.
其實是將聽過的道理,用私心,私意,仍然順著自己好賭的情慾,.
美化甚至堂然王之的論調表達出來,.
根本與蒙召的因未有相稱,.
其實當主耶穌仍然是無道,生命未徹底改變,.
不是真正的敬主,愛主的人,.
他的屬靈表現是這樣的貨色,.
但我們回想一下自己,今時今日,我何嘗不是這樣呢?.
在神面前好像知很多事,講得很堂皇,.
不知是怎樣一回事呢?.

$^{81}$弟兄姊妹,你在我日常生活中,是否真真實實的尊崇主,敬畏主呢?.
這實在是我們需要求主幫助我們..
第二件事,我們看第八,第九節,就是要彰顯主..
「夫臣的兒子耶穌基督,他的寶座是世世長存,.
他的國權是正直,守公義,恨邪惡,賜予喜樂,.
每一位信徒,每一位作主門徒的人.」.
我們生命和生活當然是要彰顯主,要榮耀神,見證基督,.
但實際情況是否這樣呢?.
我們特別信了主若干年的人,.
我們似乎被一種基督教文化或氣氛所感染,.
很順理成章地表達一些東西,但實際上裡面的餡是如何呢?.
英國著名牧師默遜死後,.
有人在書桌的抽屜裡找到一封信,.
他生前沒有提及,有人看了這封信,.
原來寫這封信的人是因為聽過牧師默遜的講道而信了耶穌..
他在信裡對牧師說:.
「我最初對你所傳的道理有興趣,不是因為你所講的道,.
而是因為你講道時面上那種聖潔的美麗和喜樂的光彩,.
吸引我,要知道你所傳講的,究竟這位耶穌是怎樣的..
請展回真是,你嘴上這樣說,.
但你的生命和生活的行動是更加實際,更有影響力,.
道言之言不變之變比講話更有力..
本身生命的外顯和感染力,程度是怎樣,.
其實是反映你裡面的生命的質素和活力是怎樣的..
今日我們每一位被提醒要好好培養我們裡面的屬靈質素和生命是怎樣的..
請問你花了多少時間好好靈修讀經,好好去到神的面前呢?.
不要說每一天,每個星期,計算總數是多少,當然不是用量來計算,.
而是質素是怎樣的呢?.
有沒有好好地經歷在神那裡支取到力量,.
以至外面流露出來基督的那種氣香,能夠令人羨慕呢?.
按經文所講的,這裡提及幾個質素,包括有正直公平,.
包括有喜愛公義,恨惡罪惡,包括有喜樂的表演等等..
究竟我們能夠彰顯這幾方面有多少呢?.
這個不同,我們真的需要靠主流露基督的美事..
英國大文豪查理士·狄更斯和當時一位很出名的朋友戴格里發生爭執,.
大家不太妥當,在1863年的聖誕節,他們剛剛在倫敦這個地方碰頭,.
大家碰到面,視若陌路人,雙方不理睬,.
但戴格里覺得良心不安,轉身,主動走向狄更斯,.
他舉起了雙手,表示不能忍受我們之間的隔膜,不能忍受疏離的關係..
狄更斯深受感動,於是大家對話聊天,前言盡釋,和好如初..

$^{121}$沒多久之後,戴格里突然離世..
一段日子之後,有人見到狄更斯站在戴格里的墓前,.
我相信他心裡一定很釋然,問心無愧,再無遺憾,.
因為他們曾經不和,有爭執,有心病,.
但在戴格里死前幾日,狄更斯終於和他握手言和,.
已經拆毀了雙方之間彼此的嚴規..
弟兄姊妹,主耶穌基督,因著祂的愛,.
因著祂自己的恩典,已經拆毀,接受了我們..
我們和神之間的牆已經拆毀,直的關係沒問題,.
但橫的關係,我們和人與人之間的關係,那種嚴規有沒有拆毀呢?.
這經文告訴我們,神喜愛公義,恨惡罪惡,.
所以耶穌是和平之君,.
祂來到這個世間,要拆毀人與人之間的厚厚的牆..
我不知道你和某些人,或者弟兄姊妹,或者你的家人,.
或者親人,有沒有一些厚厚的牆隔著?.
有些人和某些親人,可以不說話,甚至同一屋簷下,.
我真的見過聽過這些例子,其實很痛苦,但仍然當沒有回事..
有沒有這些難處呢?.
但我們知道,當我們讓基督掌權的時候,.
神的寶藏在我們心裡面的時候,這裡說是世世代代,是永遠的國權,是正直的..
我們這裡人與人之間的厚牆,那種嚴規,我們需要拆毀..
你對某些人的不滿,不和,或者對弟兄姊妹,有一些不接受,不相往來,.
如果是同一個教會,同一個團契,我們需要靠主去對付,去處理..
特別作為一個教會的領袖,預備作主聖公的人,.
他必須與人和善相親,他能夠謙卑愛人,他不能夠冷漠,很孤傲..
他不是這樣才容易與人相處,和配搭,同工..
求主幫助我們,藉著耶穌的大能大力,幫助我們,.
深信神今天願助著為王,加我們力量,賜我們智慧..
神必定祝福一些按他心意行事為人的人..
第三,我們看看第十到十二節,提出要注意的要事,就是要依靠主..
三一真實,包括耶穌基督,去創天造地,創造萬物,.
萬物都會改變,會消逝,天地都會捲起來..
但主耶穌是常存的,永不改變的,主的年歲是無窮盡的,永在不變..
大能的主甚至能夠改變天地和環境,.
但我們經常聽過,或經歷過,主的意念不同人的意念,主的道路非同人的道路,.
天怎樣高過地,照樣主的道路高過人的道路,主的意念也高過人的意念..
《耳塞書》55章8-9提醒我們,我們的主是永恆的主,不服在時間以下,.
他是昭然永存,無始無終,主處理的事件不會依附我們的時間表,.
因為他是萬有的主,掌管一切的神,時間和一切都在神的掌握下,.
命立就立,命成就就成就,沒有什麼可以攔阻或左右他的指引和安排..

$^{161}$所以弟兄姊妹,我們要明白,人有時並非全部都能幹,人始終有限,.
我們與這位無限奇畏的神相比之下,我們只是滄海的一叔,.
我們的一知半解以不知為知,我們有時用直覺,用自己的想法,自以為是,其實會變得很無端..
我們自以為認識很多東西,其實會留於偏差,我們自以為聰明,.
這些辦法以為行,其實可能行不通..
當行不通的時候,我們會誤解神,你不聽我祈禱,我們會埋怨神,不照我的方法辦事,.
神啊,你真是無能,你也不是這樣的,你不理我,你也不肯照應我..
或者你曾聽過一個故事,爸爸帶女兒去一個地方玩具店買洋娃娃,.
進去店舖,小女孩一見到洋娃娃,在面前說:我要這個..
其實爸爸心目中有一個更好,更大,更美的送給她..
爸爸隨手拿走女兒想要的次好的洋娃娃,女兒當場發了冷顫,.
抱著不放,大哭大讓,沒辦法,唯有按她的心意,給她一個次好的..
那個更好的,更美的,因為這個女兒自以為是,自恃自主,.
以有限的短小的眼光和思想,拒絕那個更上好的,更美好的福分,結果走寶..
同樣我們基督徒也是,有時我們以為一些事是應該這樣,是很好的,但神有更好的..
我們不記得經文,羅馬書八章三十二節說:.
「不愛惜自己的兒子,也為我們眾人寫明:喜不也把萬物和他人一同白白賜給我們.」.
神會將最好的給我們,特別今天我們面對一些處境,.
我們想都不明白為何神會讓這些事情發生,但神總有最美善的旨意,可能是一個試煉,.
可能經歷了這些的時候,我們會認識神更多,自己和神的關係會進深..
有時我們的信心和眼光真的太短小,主耶穌基督今天提醒我們,.
祂今天仍然是永不改變的主,祂今天仍然坐著為王的,.
掌管一切包括你人生大小的事,你的疾病,你的健康,你的經濟,你的家庭,你的兒女,子孫,.
你信不信神會將最好的給你呢?為了我們更多的依靠主,你會經歷更多..
最近我和師母回香港,有機會見到一個姐妹,她在香港某大學做教授,.
她見證神的信實和奇妙,大家都知道,在大學,很多大學都要爭表現,.
現在香港更標榜要創新科研發展,要爭取政府的金錢資助,.
這個姐妹又要教學,又要研究,又要提交提案,壓力大到她吃不消,結果她病了..
偏偏病了的那段日子,她要安排見政府的資助評審委員會,要提出她的計劃方案..
她說當時她完全在病裡面,不是狀態,但卻將這件事切切的交給神,.
她很高的祈禱依靠主,神是掌管她當時的論述和講話的應對,.
很稀奇,過了一段日子,她聽聞局方決定全面資助她的計劃,.
她自己當時覺得自己說得不好,或者表達得很差,但因為她依靠神,.
她就把軟弱當成剛強,放一些不可能的為神奇..
姐妹很感恩,是聖靈帶領她的解說和決策的人的決議,.
一切都是因為她依靠主而得的成果..
真的弟兄姊妹,今天我們要學習更多的依靠主,.
我們事奉神的人,我們更需要很好的依靠主..
在北極附近有一個海峽,見到一堆一堆巨大的浮冰,.
最大的浮冰有三四百尺高,當時吹著很強的南風,.

$^{201}$但很奇怪,浮冰卻是冒著大浪的衝擊,逆風向南飄浮,而不是向北飄浮..
原來冰山露出的水面只佔全體的體積百分之一,.
然後還有百分之七在下面,當時的水流是向南流,.
雖然海面上吹著很大的強風,但冰山的流向主要是受水流的影響,.
冰山大部分在水的下面,再大的風都不能改變冰山的流向..
同樣道理,弟兄姊妹,當我們的信心紮根在主耶穌基督那裡,.
上面吹著很強大的風,.
怎樣吹襲,怎樣打擊,甚至狂風巨浪,.
我們因為仰望主,依靠神,受著主的激流影響而推動,.
我們才能夠逆風而上,前進,前進,而不會被風向改變..
今天,作教會領袖的,要作中流砥柱的人,.
要能夠面對逆流而上,圍繞土柱的喜悅..
弟兄姊妹,盼望你和我在我們的日常生活裡面,.
尤其是善無聖功,要作核心領袖的肢體,.
以及我們將要選出教會領袖的弟兄姊妹,.
今天我們重溫這三件重要的事,.
就是要尊敬主,要彰顯主,要信靠主,.
願神加力賜恩典給我們,來低頭去結束祈禱..
願主你自己綏顧我們每一個人,.
我們都是很軟弱,我們都有很多不合你心意的地方,.
求神你赦免,用你的大能寶血,潔淨我們,.
也更新我們,繼續成為聖潔合乎主用的人,.
讓我們更多依靠你,更多在生活上彰顯榮耀你,.
我們更好好的敬畏你,祈禱奉主耶穌基督的名,阿們..
阿們..
我每天都這樣想著你.
我希望你能夠明白我認真地愛你.
我一直都這麼認真地愛你.
我一直都這樣想著你.
我希望你能夠明白我認真地愛你.
我一直都這麼認真地愛你.
我一直都這樣想著你.
我希望你能夠明白我認真地愛你.
我一直都這麼認真地愛你.
我一直都這樣想著你.
我一直都這麼認真地愛你.
\newpage



\section{}
\label{sec:NJ2Erozw8Ss}
\textbf{11/18/18 主日講道: 女子組加油! \#3 顛覆男人世界的聖經女子 (陳崇基牧師)}
\newline
\newline
連結: \href{https://youtube.com/watch?v=NJ2Erozw8Ss}{\texttt{https://youtube.com/watch?v=NJ2Erozw8Ss}} ~~~~ 語音日期: 2019-08-03
\newline
\newline
\hyperref[sec:wT2Sf41UjYU]{\small{< < < PREV SERMON < < <}}
~
\hyperref[sec:index_chronic]{\small{[返順時目]}}
~
\hyperref[sec:index_scriptual]{\small{[返順卷目]}}
~
\hyperref[sec:QpmQmJgiMg8]{\small{> > > NEXT SERMON > > >}}
\newline
\newline
$^{1}$我們今天繼續幾個星期的女子組加油系列.
我們說到用另一個角度來看聖經.
為了什麼人加油呢?.
為了女性.
為了姐妹.
我們去加油.
一個很切身的題目.
因為現在女性的地位.
男女平等.
是一個很切身的題目.
有很多不同的運動.
都是跟這個gender有關.
Me Too的運動.
聖經裡面有沒有一些提醒給我們呢?.
或者聖經雖然是一本這麼古老的書.
有沒有一些劃時代的信息值得我們留意呢?.
我們這幾個星期在討論這個題目.
聖經雖然很舊.
但是.
裡面的信息有很多是劃時代的.
是我們今時今日21世紀的人.
我們都覺得可以形容得到.
所以我們由第一個星期開始到今天第三次.
我們就說有一些命題.
我們會提出來.
然後我們大家去想想.
到底是對還是錯呢?.
如果你今天第一次來的話.
我知道我們中間有一些新朋友.
我們目前說的東西是這些.
第一個命題就是.
聖經是一個父權社會裡面的文學作品.
這是對的.
是一個patriarchal literature.
裡面有很多背景.
是當時古代的社會.
重男輕女的價值浮現出來的.
但是不代表聖經是維護重男輕女的思想.
或者父權主義.
所以第二個命題就是.

$^{41}$聖經很多時候是挑戰.
顛覆父權主義的觀念.
特別是聖經提到一些女人出現的時候.
雖然不是那麼多.
但是物以罕為貴.
每逢女士出現的時候.
有很多時候都是一些很重要的角色.
很關鍵的角色.
就叫我們看聖經的時候.
留意一下.
原來上帝也是藉著這些女性.
來做很多事情.
做她自己的工作.
第三就是聖經歷史裡面.
由舊約到新約.
女性是扮演一些很關鍵很策略的角色.
我記得我們看了這個 chart(表)很多次.
下面那一部分就是很多時候我們在一些參考書裡面提到.
聖經歷史的發展.
上帝用什麼人來帶領自己的工作.
讓上帝的天國能夠臨到地上.
下面那些全部都是男人.
我們平時的參考書都是這麼說的.
沒有錯的.
不過從另一個角度去看的時候.
其實上面那些女性.
很多時候就是做先鋒或者做最關鍵的角色.
所以我們不可以不留意這些事情.
女人不是可有可無.
在聖經裡面.
上帝不僅是用男人.
也用女人.
用女人做很多重要的工作.
第四就是解讀聖經的時候.
我們需要極度敏感.
聖經的背景.
我想告訴你.
這些事情不是很深的事情.
你第一次來教會的時候.
第一次看聖經或者沒看過聖經.

$^{81}$你不要覺得這些是很深奧的事情.
我舉個例子.
上年.
大概十月的時候.
你記得我們教會不是有一個話劇組.
不是有一些修女在做話劇的嗎?.
記不記得?.
我老婆是其中一個演員.
她是其中一個修女逃走.
參加宗教改革的運動.
根據真實的故事.
她是其中一個演員.
她排戲的時候.
因為她排戲要記很多台詞.
我經常在家裡.
我們在做事的時候.
家裡很安靜.
她突然就會爆一句出來.
她說有道理不是一切.
你想想.
不只是這樣.
她在煮飯的時候就會說.
或者無緣無故.
我們很平靜地在做事.
很集中精神在做事的時候.
她可能在廁所裡面大叫一句.
有道理不是一切.
如果你沒有背景.
沒有上文下理.
你就不明白她在說什麼.
你就很害怕.
以為她在罵你.
其實我一向說話都很有道理的.
但是.
但是她那句話不是跟我說的.
那句話是她台詞的其中一句.
所以你明白背景.
明白前文後理的時候.
你就明白那句話的意思.
其實《體性經》也是這樣的.

$^{121}$有一些經文令我們非常困惑.
因為很多時候.
我們只聽到對話的某一方面.
另外那方面我們看不到的.
就好像你只收到一個回覆的電郵.
原本的電郵就沒有了.
你就要明白當時的背景有什麼處境.
引致那本《性經》.
或者那本書信的經文是這樣寫的.
就是這麼簡單.
我們要學習聆聽.
作為一個基督徒.
甚至作為一個人.
都要學習去聆聽.
聆聽配偶說的話.
或者聆聽家人說的話.
其實《體性經》也是有這樣的需要的.
今天我們講到女性的時候.
我們就無法避免.
有一些比較棘手的經文.
有些聽了之後.
好像突然間.
有道理不是一大堆的.
我們知道一點處境.
知道一點背景的時候.
可能我們會多一點了解.
我們一直.
我想告訴你.
我們不需要怕.
去對《性經》問一些很尖銳的問題.
神可以處理的.
所以我們那些青年人.
英文青年人開始有一個 Discipleship Group.
有一個 Biostudy.
有一個導師帶他們查經.
我就跟我的兒女說.
他們現在十五十七歲.
我就說麻煩你.
你去查經的時候.
因為那些青年人經常覺得查經很悶.

$^{161}$我就說麻煩你.
麻煩你去查經的時候.
我請你不要問一些典型的主要學的問題.
你問一些 Critical questions.
最好是問到老師回答不了.
我真的這麼說.
他們就搔頭.
是啊,先問一下.
每次吃飯的時候我就問.
上個星期五你們查經的時候問了什麼問題呢?.
那個帶查經的能不能回答呢?.
他說回答不了.
我說好啊.
回答不了就好.
我們應該用腦子的.
所以其實我們一向都是希望大家.
作為基督徒我們都要吃腦子的.
並且我們最近的 Sunday School.
或者過去這麼多個 Sunday School 的題目.
其實都是在說這件事情.
更加今天有一個這麼特別的講座.
李華弟兄跟我們講聖經翻譯.
其實很難得的.
我們應該 Critical 一點.
不是要去批評什麼.
而是去問一些尖銳的問題.
不需要怕.
所以今天我們處理一段很激受的經文.
我們試一下用一個這樣的角度去看.
到底這段聖經想講些什麼呢?.
特別講到女性在教會裡面的角色.
這段聖經是哥林多前書第十一章.
我知道在我們中間有一些剛剛初來的.
未信主的.
我不想將這件事情講得太複雜.
因為覺得這麼複雜.
不知道你在說什麼.
但是我想藉著以下的時間.
給大家一點指引.
然後大體上怎麼去應用.

$^{201}$這段聖經跟我們有什麼關係.
哥林多前書第十一章.
這段出了名.
是最棘手的經文.
有很多不同的解法.
有很多不同的體會.
但是不代表我們不看.
我們看看這段聖經講些什麼.
十一章第二節.
保羅是一個牧者.
當時第一世紀初期教會的一個牧者.
他寫信給一間很有問題的教會.
叫做哥林多教會.
這本書叫做哥林多前書.
是一本書信.
寫到第十一章的時候.
他就處理到女人在教會.
在崇拜裡面的衣著.
髮色等等的問題.
就是頭髮等等的問題.
十一章第二節.
這是新漢語譯本.
幫李懿華弟兄買了廣告.
新漢語譯本.
是一個很好的譯本.
十一章第二節.
他說.
我稱讚你們.
因為你們事事紀念我.
又堅守我所傳給你們的傳統.
但我希望你們知道.
基督是每個男人的頭.
男人是女人的頭.
神是基督的頭.
每個男人禱告或者宣講神的信息.
如果蒙著頭.
就是羞辱自己的頭.
OK.
每個女人.
如果禱告或者宣告神的信息.

$^{241}$如果不蒙著頭的話.
就是羞辱自己的頭.
因為這就等於剃了頭一樣.
女人如果不蒙著頭.
就應該把頭髮剪掉.
女人如果以剪髮剃頭為可恥的.
就應當蒙著頭.
就是宣講信息.
講到禱告的時候.
就需要蒙頭.
女人需要蒙頭.
第七節.
男人不應該蒙著頭.
因為他是神的形象和榮耀.
女人卻是男人的榮耀.
因為男人是從女人而出.
女人卻是從男人而出.
並且男人不是為女人而做.
女人卻為男人而做.
因此為天使的緣故.
女人頭上應當有權柄的記號.
十一節.
但在主裡面.
在主耶穌基督裡面.
女人不可沒有男人.
男人不可沒有女人.
因為正如女人是從男人而出.
就像男人也是藉著女人生出來的.
但萬有都事出於神.
你們自己判斷.
女人向神禱告.
不蒙著頭是恰當的嗎?.
難道你們本性沒有教導你們.
男人若留長頭髮是他的羞辱嗎?.
女人若留長頭髮.
卻是他的榮耀嗎?.
因為賜給他的長頭髮.
就是跟蒙頭的外衣寬領一樣.
如果有人想要爭辯.
我們.

$^{281}$使徒保祿說.
我們可有這樣的習俗.
神的宗教會也沒有.
聽完之後.
還說什麼呢?.
什麼事?.
正如我所說.
這是一個很棘手的經文.
到底說什麼呢?.
好像聽電話一樣.
聽別人說話一樣.
只聽到一邊.
不知道對方到底在問什麼問題.
他在回答什麼.
這段聖經看的時候.
有一點像這樣的情況.
但是我想告訴你.
不僅是棘手.
也成為很多時候教會裡面的爭議.
幾個月前我去過加州港島.
就聽了這樣的故事.
有一個教會.
有一個女傳道人.
她要港島.
是不是編排她去港島.
這是第一個問題.
但是如果她去港島的話.
讓她去港.
那應該怎麼去港呢?.
結果那個教會.
結果那個長執開完會之後.
就說可以去港.
那個女傳道可以去港島.
不過她需要蒙頭.
蒙頭是什麼意思呢?.
就是有條紗巾的意思.
I guess(我猜測).
就是當時就是這個意思.
所以那個女傳道一定要蒙著頭去港島.
他們的根據就在哪裡呢?.

$^{321}$就是所謂這段聖經.
似乎這段聖經是在說這件事情.
那麼.
這樣應用對不對呢?.
那就不知道了.
就是.
不同教會有不同傳統.
一處鄉村一處泥.
不過我想問一個問題.
就是這樣去應用.
就是說現在女人出來港島的時候.
或者教主學這樣才算.
她就需要用條紗巾.
蒙著頭.
是蒙頭這樣.
這樣.
我想問一個問題就是.
這樣應用.
是不是不是很敏銳當時的文化背景呢?.
一般的理解.
說保羅在說什麼呢?.
似乎是在針對一個問題.
就是當時哥林多的教會.
有一些女人.
信了主.
她信了保羅所傳的福音.
信了耶穌.
那麼他們發現在耶穌裡面.
他們有一個新的 identity(身份).
新的身份.
那麼教會就成為一個新的人類.
新的 humanity(人類).
男女是平等的.
原來在基督裡面男女可以平等的.
保羅是這麼傳道的.
那麼這個福音就顛倒了.
當時羅馬社會的那種所謂 honor and shame(榮耀和羞辱).
就是很注重榮耀和羞辱的那種文化.
重男輕女的社會.
那麼你可以想象.

$^{361}$如果你是當時的女人.
你從小到大都被人歧視的.
被男人歧視.
就是女子無才便是德.
你有沒有聽過這句話?.
讀這麼多書做什麼的女人?.
就是你經常都是在那種環境裡面長大的.
是男人做頭的.
女人是跟尾的.
那麼你是女人.
你想象一下.
如果你是女人.
你在那種文化的分途裡面.
你就一定要.
像這樣的樣子.
有這樣的身材或者有這樣的美貌.
你才會得到尊重.
才會得到接納的.
你突然間聽到保羅所傳的福音.
在耶穌裡面原來.
我們是受神尊重.
受神的愛.
受神的恩典的.
男女都是一樣的.
那麼這個是教會.
一個新的社會.
一個新的群體.
在群體裡面.
在基督裡面.
就沒有分高低貴賤的.
你想一下是多麼吸引.
是一個很吸引的信息.
所以他們進來教會是信了的.
我舉個例子.
我想八幾年前.
八幾年.
八幾年前.
就是西教士進去中國.
去傳福音的時候.
有很多女西教士.

$^{401}$就在中國.
她們想很多的辦法.
就是去廢除紮腳的習俗.
在我們中間有誰是紮過腳的?.
沒有.
但是其實這件事情不是很遠的.
你媽媽或者你的祖母.
Great-grandma.
可能就是受到一些這樣的對待.
所以西教士就盡量去廢除這件事情.
你想一下你在那種情況下.
出來的女性聽到.
上帝的福音.
是一件多麼革命性的事情.
來到主的面前就不分高低貴賤的.
紮腳也不用紮.
是多麼吸引.
By the way.
很多在中國八幾年前.
建立的女子中學.
都是在西教士建立的.
我在香港讀中學的.
有誰是在真光女書院讀書的?.
那時候我們讀中學.
我的男生很喜歡認識真光女書院的女生.
因為他們穿的校服很漂亮.
說來說去.
但是那些女子書院很多時候都是在西教士建立的.
因為他們看到那些女生.
他們沒有機會得到教育.
是一個很革命性的信息.
福音這個Gospel.
可能.
這個情況就是這段話背後發生什麼事情.
就是那些女人信了主之後.
她們進了光明堂的教會.
她們覺得我們不需要受制於當時羅馬的社會的蒙頭風俗.
因為我們看到一些石雕像.
那些女士.
到了憲制的時候.

$^{441}$那些假神宗教的女士.
她們蒙頭.
所以她們進到教會的時候.
她們索性就不蒙頭了.
頭巾都不要了.
那麼.
也可能當時的女性.
進到教會.
她們說我們崇拜或者敬拜的時候.
就是禱告或者.
講道的時候.
我們就可以披頭散髮.
對當時的社會來說.
披頭散髮的女人.
是一個壞女人.
所以.
可能這個就是當時的處境.
保羅寫了一封信.
是去關注一個問題.
就是他想當時的敬拜的聚會.
教會.
的聲譽不會受損.
在我們西方現代的社會裡面.
男女的髮型.
或者頭飾.
我們是很自由的社會.
我們覺得沒什麼所謂的.
有些男人是長頭髮的.
也有女的.
可以是短頭髮.
但是在東方的社會.
或者在.
武守的社會.
這件事情是很緊張的.
很緊張.
譬如.
我以色列.
我每年都會帶一些弟兄姊妹去.
我經常去到某些地方.
我是被人罵的.

$^{481}$因為我是帶著帽子的.
你知道去猶太人的會堂.
就要帶帽子的.
你不可以光著頭.
男人不可以光著頭.
但是到了基督教教堂.
你不可以帶帽子進去的.
我經常被基督教的修士罵我.
脫掉帽子.
你這麼不尊重.
這是一個神聖的地方.
這些社會的文化就會很緊張.
這些所謂衣著的事情.
你外表的事情.
或者女士.
我們一團去.
去不同的地方.
去遊覽的時候.
女人不可以露這麼多肉的.
是不是?.
裙子不可以太短.
肩膀不可以露出來.
你需要套一件巾.
這樣你才可以進去的.
你可以想象.
那個情況有一點像今天.
如果那些女士.
來教會崇拜的時候.
她穿少布一點.
因為西方社會講自由.
但如果每個人都是這麼少布的.
你覺得教會給人的印象好不好?.
正如我所說,我在香港長大的.
中學的時候.
讀書的時候我們穿校服.
久不久就要塗校服了.
我記得那些女生要跪在桌子上.
然後用間尺量度.
她有沒有收短一條裙子.
你想想如果社會也是這樣的話.

$^{521}$教會是不是應該要留意呢?.
因為不想去妨礙我們的見證.
似乎這個就是當時的背景.
所以你不要以為.
只是保羅去關心.
那個女人要穿什麼.
或者戴什麼上頭.
保羅所緊張的.
那個大圖畫是什麼呢?.
就是基督徒的崇拜.
教會的聚會.
是一個公開的見證.
public testimony.
我們是一個見證.
坐在一起.
我們原來今天就是一個見證.
另一處很多的地方.
保羅講到女士的時候.
他就說那些女士進來的時候.
特別是一些所謂貴族身份的女士.
進來之後.
就不要穿金戴銀了.
因為在教會裡面.
一個新的社群裡面.
就不應該貧富有落差.
所以不是dress up to go to church.
dress down.
是嗎?.
dress down.
因為不要在這裡分辨到底是哪一個呢?.
是富是貧.
有一點像這種背景.
這是第一個指引可以給大家的.
或者.
那邊電話是在講這件事情.
而保羅呢.
說了不多前書.
就是在處理.
關注這個問題.
所以他說女人講到要蒙頭.

$^{561}$但這當然是當時文化的背景.
還有什麼指引可以給大家呢?.
一些pointer(指引)呢?.
正如我所說的這段聖經.
很多不同的說法.
很多不同的解法.
那麼.
也在這裡講到女人要蒙頭.
似乎也很奇怪.
因為保羅在其他地方沒有怎麼講過.
就算你看背景.
有些女人.
外邦的女人.
有蒙頭.
有些也是沒有蒙頭.
那麼近年有什麼新的方向呢?.
或者什麼pointer(指引)呢?.
那麼有一個女色經家.
正如我所說的這幾次.
我們特別看了很多女色經家寫的書.
到底他們有什麼留意到.
是我們男人未必留意到的呢?.
那麼有些新的說法就是這麼說的.
就是說.
可能.
哥林多前書11章.
這裡講到女人要蒙頭的時候.
女人講到公土.
要蒙頭的時候.
不是保羅自己的意思.
可能是他引用.
哥林多教會裡面的人講的話.
他引用他們的說法.
然後駁斥他們的說法.
那麼這個看法呢.
就180度不同了.
那麼你這麼看的時候.
保羅不是自己在講女人要蒙頭.
而是哥林多教會有這樣的 practice(練習).
他就引用他們的說法做法.

$^{601}$然後駁斥.
或者是發出疑問.
這些不是先例的.
特別是在哥林多前書.
當然有很多其他書信.
保羅都會做這些事情.
比如《勃羅馬書》.
有些所謂假想敵.
哥林多前書也有的.
比如《六章》第十二節.
現在新的《聖經翻譯》也有很多.
已經 point out(指出)告訴你了.
這裡保羅是直接引用當時哥林多教會.
自以為好.
自以為對的話.
他們的說法.
然後他就駁斥他.
黃色的那些就是保羅講的.
白色的那些就是哥林多教會當時流行.
他們教會裡面經常講的.
所謂的 slogan(口號).
他說凡事都可行.
其實是在引用他.
你看到 NIV 就是 quote(引用)著.
Everything is permissible for me.
但是 but.
保羅講什麼呢.
but.
不是每樣東西都有益的.
他就在駁斥他們.
不要以為順了我們有新的自由了.
我們什麼都做得方中正確都可以.
但是凡事都可行.
但是不盡都有益處.
凡事都可行.
但是每一件我不受他的核制.
他就是在駁斥.
當時哥林多教會一些錯誤的思想.
如果.
講到女人要蒙頭.

$^{641}$那段聖經.
其實保羅是在發出疑問.
哥林多教會有這樣的習俗的時候.
這樣看來又有些不同了.
我就引用這個女聖經家.
Lucy Piper.
她的解法.
我們再讀一讀這段聖經.
聽一聽到底保羅想講什麼.
首先第二節.
黃色的就是保羅講的話.
白色的就是他所謂引用哥林多教會當時的風俗.
或者他們的講法和做法.
首先第二節.
我稱讚你們.
因為你們事事都紀念我.
又堅守我所傳給你們的傳統.
但是我希望你們知道.
基督是每個男人的頭.
男人是女人的頭.
神是基督的頭.
你們說.
這些廣東話我加上去.
當然了.
新學院日本沒有的.
這個我是根據Papir的解釋.
You say.
你們說.
每個男人禱告或者宣講神的信息的時候.
如果蒙著頭就是收了自己的頭.
如果每個女人要公禱.
要去宣講神的信息.
如果不蒙著頭的話.
就是收了自己的頭.
因為.
就如剃了頭一樣.
這樣的話.
So it means.
保羅就說了.
你的意思就是說.

$^{681}$女人如果不蒙著頭的話.
就應該將頭髮剪掉.
女人如果以剪髮剃頭為可恥的.
就應該蒙著頭嗎?.
你們這樣解釋.
男人不應該蒙著頭.
因為他是神的形象和榮耀.
女人卻是男人的榮耀.
因為男人不是從女人而出.
女人卻是從男人而出的.
男人才有女人.
並且男人不是為女人而做的.
你們又說女人卻是為男人而做的.
因此為天使的緣故.
女人頭上應該有權柄的記號.
我想矯正你們.
OK.
Now I correct you.
但在主耶穌基督裡面.
女人不可以沒有男人.
男人不可以沒有女人.
什麼意思呢?.
就是男女是.
大家dependent.
interdependent.
因為正如女人是從男人而出.
照樣男人也是藉著女人而出的.
是不是?.
有媽媽才能生我們.
但萬有是出於神.
那.
你自己判斷一下.
女人如果要向上帝禱告.
不蒙著頭是不是恰當呢?.
當然是恰當的.
因為男女平等.
Of course.
不用蒙著頭.
是不是?.
難道你們的本性.

$^{721}$就是你們的文化不是這樣教導你們嗎?.
男人如果留長頭髮就是他的羞辱.
女人如果留長頭髮就是他的榮耀嗎?.
因為賜給女人的長頭髮.
就是跟蒙頭外衣歡慶一樣.
如果你們中間有誰呢?.
想去爭辯.
跟我爭辯的話.
我想告訴你.
我們沒有你們這種習俗.
神的宗教也沒有.
就是說女人蒙頭.
只是哥林多教會的習俗而已.
這個解法有些通.
有些通是什麼呢?.
如果你這麼看的話.
他是在引用哥林多教會錯誤的壞鬼神學.
他似乎整段聖經.
整段話我們聽到的時候.
保羅就沒有前言不對後語.
沒有九不搭八.
因為他不是贊成哥林多教會的做法.
而是他發出疑問.
也解釋了最後這一句.
他說宗教會沒有這種做法.
沒有這種習俗.
女人蒙頭的習俗.
如果這麼看的話.
似乎保羅說的這句話.
他在電話裡面說的就是.
他正在針對哥林多教會裡面的大男人.
他們就視女人為他們的二等公民.
就是說他們的女人要講道的話.
可以講.
不過呢.
就要求他們要蒙頭.
女人要蒙頭.
這個是哥林多教會大男人的神學.
不是保羅自己的神學.
保羅本身是不贊成女人蒙頭的.

$^{761}$針對的.
這段聖經是針對的.
我們似乎就以為是針對那些很自由的女人.
你這麼看的話.
保羅所針對的就不是女人.
就是大男人.
教會裡面.
一些有霸權的男人.
其實這麼看的話.
又配合整本哥林多前書.
整封書信.
因為整封書信就是在針對.
保羅作為牧者就是在針對當時的教會有大男人.
他們分黨.
他們愛權.
他們自大.
自為事.
他們放縱情慾.
彼此爭議甚至搞到上了法庭.
甚至他們認為自己某一些才華.
然後看低其他人.
等等等等.
似乎這樣的脈絡.
又多麼純.
這跟我們有什麼關係呢.
我想我講了這麼多篇.
第一篇.
這篇唯一一篇我經常講.
似乎.
可能.
我們其實也在掙扎這個問題.
這段經文到底背後是怎麼樣的.
但是我想你留意一件事情.
千萬不要只見樹木而不見森林.
無論為什麼也好.
OK.
無論怎麼去看也好.
哥林多前書第十一章這裡處理的是什麼呢.
你不要見不到.
見到樹木就見不到森林.

$^{801}$第十一章處理的是什麼呢.
就是教會公開敬拜聚會的問題.
有什麼問題呢.
從十一章開始.
直到十四章.
都是在處理這個問題.
就是教會敬拜的行為.
有沒有 jeopardize(妨礙)教會的見證.
第十一章講到女人的蒙塗.
無論是本來贊成也好.
反對也好.
其實都是重點在處理.
整個大圖畫就是.
我們教會敬拜一起聚會的時候.
給人一種什麼印象.
然後他就處理到.
十一章之下的.
就處理到聖餐.
community.
就是他們哥林多的大男人.
他們就分.
我們這群有錢的就有東西吃.
你那群遲到的沒錢的.
剛剛下班.
來到.
吃東西沒東西吃.
他就是在處理這種問題.
他說這是一個失見症.
因為我們聚會的時候.
被人看到這種情況.
這個不可以的.
第十二章講到才幹和恩賜的時候.
都是一樣.
因為哥林多的大男人.
是口齒薄皮的.
然後十三章的時候.
他就講到外騙.
是不是?.
這段話就是經常昏迷胡容的.
但其實他就是叫教會.

$^{841}$不要爭權奪利.
要彼此的幫助.
在外面來建立.
然後第十四章他講到方言的問題.
其實他重點講什麼呢?.
要你們教會的聚會崇拜的時候.
要讓人看出神跟我們同在.
神是真實的.
所以十四章最後的時候.
講什麼呢?.
他說所以全教會在一處聚會的時候.
若說方言.
偶然有一個不通方言的人.
或者有未信主的人.
進來.
如果你每個人都講一些東西.
是我們人們不明白的話.
豈不是說你們是癲狂呢?.
二十四節.
但若都是先知講道.
就是說一些東西是大家明白的.
甚至未信主的人進來是明白的.
偶然有不信的.
或者不通方言的進來.
就被眾人勸醒.
被眾人審明.
他心裡面的隱情顯露出來.
他就將面服於地.
然後敬拜神說.
神實在在你們中間.
保羅整個著眼的地方就是這樣.
就是教會在基督裡面.
是新的人類.
是新的社會.
是新的一個群體.
一個新的聖殿.
在這裡我們能夠跟神相通.
見到上帝.
能夠遇上上帝.
這是一個新的創造.

$^{881}$所以你留意.
保羅真的好像一個有智慧的家長.
小孩子開始長大的時候.
他就不會只跟他說 do and don't.
他給他一個新的 vision.
原來教會我們走在一起的時候.
就是一個 new creation.
上帝住在我們中間.
所以為什麼我們聚會的時候.
我們需要明白.
要以愛來建立.
要去給予彼此的幫助.
要去苦就卑微的人.
要有愛有恩典在我們中間流露出來.
這就是一個 new creation.
我們就是一群新的人類.
他不只是單單跟你說.
我們會聚得舍.
搞定.
是一個 new people.
一個 new creation.
我們教會走在一起就是這樣.
結束的時候我想有幾件事.
可以帶走.
無論你覺得剛才那段聖經.
真的很難搞.
很棘手.
不要緊.
但最重要.
有幾件事值得應用的就是.
敬拜.
我們主日走在一起聚會的時候.
我們如何可以讓人看到.
神真的在我們中間.
如果按照剛才我第二個解釋的話.
就是說.
這段聖經保羅所處理所針對的人.
不是那些自由的女人.
而是那些濫權的大男人.
他逼那些女人去蒙頭.

$^{921}$如果是這樣的看法的話.
就是說那些濫權的人.
在教會裡面限制了女人.
不准她.
不准女人來實踐上帝對她的呼召.
是不是?.
你儘管可以講道做教道.
不過我想你蒙頭.
信服.
在男人面前.
如果按照這個解釋.
這就是思路.
在幾天前.
我跟一個姐妹又聯絡了.
認識了她很多年.
她就讀完神學.
她就去一個教會事奉.
我跟她再聯絡.
她說不如約出來聊聊天.
她在布吉博特附近的一個教會.
她說我以前是去A教會的.
現在我去了B教會.
她說這麼多年沒見面.
我們沒有聯絡.
我以為你一直都在A教會.
她說不是.
我讀完神學.
原來在A教會裡面.
我最多只可以做助理.
助理的助理.
她說她跟她的先生再祈禱之後.
她覺得還是move on.
就轉了去現在B教會.
也是布吉博特.
她說那裡的空間可能大一點.
比她事奉的空間大一點.
我說.
I guess.
It's time to move on, right?.
你看到那個意思嗎?.

$^{961}$如果按照剛才第二個解法的話.
就不是單是那些女性.
而是那些愛權,濫權的大男人.
上帝要顛覆他們的價值.
因為在基督裡面.
男女按照恩賜去事奉.
不要限制他們.
我第一間神學院讀M.Div的時候.
是.
我們讀M.Div的.
就是準備做牧師的那班人.
我們有女同學.
那時候.
我記得那些教授教女生.
你可以在M.Div program.
Fine.
但我想告訴你.
這個是for your own personal enrichment.
你將來不會去港島.
因為我們的理念.
就是你不可以去港島.
男人說.
你可以教.
教小孩可以.
大人就不可以教.
I don't know.
就是.
有傳統我們應該尊重.
是不是?.
但我就想.
我們作為基督徒.
二十一世紀的人.
我們想一想.
聖經裡面.
到底怎樣去掙扎這些經文的意思.
怎樣去應用它們.
如果這樣看法的話.
那個大圖畫就是.
我們一起走.
一起事奉.

$^{1001}$一起敬拜的時候.
我們要問.
到底是不是.
我們有時候錯誤應用.
純粹問一個問題.
還是.
我們需要繼續問.
我們有沒有一些愛權濫權的時候.
我們就限制其他人.
來服侍.
我們有九位的.
長執的.
候選人.
六個是女的.
三個是男的.
我覺得.
That's the way to go.
就是男與女我們能夠彼此的.
去配搭.
並且我想最後帶走的.
就是我們.
保羅你記得那個大圖畫是什麼嗎?.
在敬拜的時候我們給人.
社會.
一個的見證.
教會你有沒有想過.
走在一起.
主日崇拜的時候一舉一動.
都是一個公開的見證.
Public testimony.
敬拜不是純粹一個個人的事情.
信主信耶穌都不是一個個人的事情.
雖然我們經常.
口頭禪說個人救主.
是一個群體.
在基督裡面.
新的創造.
我們有沒有.
很好的敬拜.
留意.

$^{1041}$我們整個的敬拜.
我們呢.
是一個.
要是一個好的見證讓人知道.
上帝在我們中間.
有沒有一些留習.
是我們需要去改變的.
可能是一些很小的事情.
遲來早走.
我們不專心.
不專注.
這些都是.
我們給人的感覺.
我想說.
有很多.
來參觀教會.
第一次來教會的人.
他們很聰明.
他們豎起天線.
一進來就豎起天線.
是應該的.
他們就在看我們基督徒.
到底我們是不是假冒為善的人.
我們口唱心和.
就說上帝是我們主.
又這個又那個.
你說祂多麼重要.
多麼尊貴.
如果我們的.
聚會.
以致我們的生活.
不是反映出這個價值.
人們就會發現到.
有一個女作家說.
今時今日.
不信主的人.
通常不去教會的人.
他們有一種叫做.
BS.
Detector.

$^{1081}$BS Detector.
你明白我的意思吧.
不用翻譯了.
BS Detector.
就是.
他對教會有一種.
很強烈的BS Detection.
他一進來就問到.
他就知道你.
到底是不是真的.
基督徒去敬拜.
基督徒去聚會.
我們大家走在一起.
作為教會.
神的子民.
我們給人家.
任何一個.
強烈的證據.
告訴別人.
上帝真的在我們中間.
我們同心投地去祈禱.
主要我們今天來敬拜你.
我們知道你和我們同在.
主要讓我們.
恭敬.
在你面前.
尊你為我們的主宰.
尊你為大.
主赦免我們很多時候.
很隨便地來敬拜.
赦免我們很多時候.
我們沒有留意.
我們對你的敬拜和事奉.
主讓你在我們教會.
能夠同心.
不分男女.
我們不分恩切.
我們一起.
來見證.
去敬拜你.

$^{1121}$讓你在我們中間.
能夠和我們同在.
也讓其他人看到.
我們.
去敬拜你的時候.
去尊你為我們王的時候.
神你是真的實在的.
使我們這個城會教會.
能夠繼續在愛裡面.
得以建立.
我們這樣祈禱.
奉主命.
願上帝祝福各位.
Happy Thanksgiving.
多謝收看城會網台.
記住Like, Comment.
或者訂閱, Subscribe我們.
\newpage



\section{}
\label{sec:QpmQmJgiMg8}
\textbf{4/7/2019 主日講道: 智取人生(三) 陳崇基牧師}
\newline
\newline
連結: \href{https://youtube.com/watch?v=QpmQmJgiMg8}{\texttt{https://youtube.com/watch?v=QpmQmJgiMg8}} ~~~~ 語音日期: 2019-08-03
\newline
\newline
\hyperref[sec:NJ2Erozw8Ss]{\small{< < < PREV SERMON < < <}}
~
\hyperref[sec:index_chronic]{\small{[返順時目]}}
~
\hyperref[sec:index_scriptual]{\small{[返順卷目]}}
~
\hyperref[sec:JN7FdDAB0F8]{\small{> > > NEXT SERMON > > >}}
\newline
\newline
$^{1}$今天我們繼續看這個系列.
我們一直在說一個叫做智慧文學的系列.
我們開始的時候我問過大家一個問題.
你記不記得或者你領悟到你人生裡面曾經有一個智者出現過.
而那個智者有什麼特別的地方.
記不記得?.
還記得嗎?.
有沒有這麼多個星期之後你想多幾個?.
那些智者在你生命裡面.
可能是一些長者或者一些朋友.
給你一些生命的引導.
這些智者.
能夠幫助你在人生裡面有更多的智慧.
所以我們今天也會繼續看聖經裡面怎麼說智慧.
你記得我們說過在智慧文學裡面.
其中一本書叫做《傳道書》.
這本書是說到人生裡面最終極的關懷是什麼.
有什麼最終極的關懷.
智者說是需要我們知道的呢?.
我們今天就來看看.
一章第一節.
他說在耶路撒冷作王大位的兒子.
傳道者的言語.
傳道者說.
虛空的虛空.
凡事都是虛空.
人一切的牢獄.
就是他在日光之下的牢獄.
有什麼益處呢?.
答案就是沒有益處.
一切都是empty.
meaningless.
虛空的.
這裡說說傳道者是什麼人.
希伯來文叫做cohelic.
cohelic的意思就是說.
有一個像樣的哲學家.
或者一個人生的哲理專家.
或者是一個聚集一群人.
assembly.

$^{41}$去跟他們說話的一個人.
可以叫做今天的傳道者.
傳道書就說有一個這樣的人.
有一些言語.
他想告訴我們.
虛空是虛空.
凡事都是虛空.
希伯來文.
就是聖經舊約的原文.
很多時候是跟那些聲音.
聲音的色彩.
是很有關的.
是一個很colorful的language.
什麼叫虛空呢?.
就是heavy.
你試一下把手放在你的嘴唇前面.
你這樣.
heavy.
就是這樣的意思.
什麼叫虛空呢?.
虛空就好像一口氣.
你一口氣有多長?.
很短的.
你這樣走完之後.
你的人生就是這樣.
傳道者就是要我們想最終極的東西.
既然人生是這樣的話.
他說如果在日光之下.
life under the sun.
就樣樣東西都是empty.
和沒有意義的.
你這麼消極.
為什麼呢?.
因為有兩個原因.
整本書解釋了兩個原因.
第一個原因就是.
在日光之下發生什麼事呢?.
你以為人生短促.
我們盡量把握時間.
做一些有意義的事.

$^{81}$做一些好的事.
以致有些價值.
但他說在日光之下發生什麼事呢?.
他說有異人行異.
反戰滅亡.
有惡人行惡.
倒響長壽.
這都是我在日光之下.
虛度之年日所見到的.
就是說世界裡面有很多的不公.
使得這個生命更加虛空.
我們看報紙.
看網上的新聞.
幾乎單單都是和這些有關的.
犯錯的人.
罪大惡極的人.
可能逍遙法外.
但很多時候中宗直直的人.
就要受苦.
受迫害.
所以他說第一個原因.
這個生命裡面是有意義的.
是沒有意義的.
為什麼呢?.
因為有不公.
當然這些不公的事.
都會發生.
但他說在光明之下.
在虛度之年日所見到的.
所以人生就沒有意義.
第二個原因.
為什麼我們人生是有意義的.
是好像一口氣.
一吹就消失了.
就是因為有死亡.
每個人都要經歷死亡.
每個人都在邁向死亡.
戴吉勵士說過.
死定了.
我們每個人都要面對人生的終極.

$^{121}$所以很簡單.
整本書他講什麼呢?.
用了很多篇幅.
他告訴我們.
人生在世如果只有日光之下.
只有here and now.
人死如燈滅的話.
每樣東西.
都是虛空的.
用英文講有句說話挺好的.
英文叫Life sucks.
想想中文怎麼翻譯.
Life sucks.
傳道者就說.
Life sucks.
and then you die.
Right?.
中文怎麼翻譯呢?.
廣東話可以這樣翻譯.
做人很實際.
然後就死了.
是不是很實際?.
做人就是這樣.
在日光之下.
我們就是這樣.
傳道者就告訴我們.
用了很多篇幅.
在今生今世裡面.
這位傳道者.
他嘗試在人生裡面.
在日光之下找意義.
增加自己的學識.
多讀幾個學位.
去找多一點快樂.
無論是情慾上的快樂.
或者不同的.
消遣裡面的快樂.
工作裡面的成就.
或者我們儲很多的財富.
有很多的地位.

$^{161}$但最後一切都是虛空的.
因為Life sucks.
and then you die.
每樣東西你儲蓄起來.
能夠建立的就乘以零.
一百萬乘以零.
就是零.
一億萬乘以零.
所以每個人都要面對一個這樣的終局.
所以這本書開頭的時候.
他用一個人.
來做一個化身.
他說.
看看這個人的生命.
想一下這個人的生命.
你就會知道是一個虛空.
再虛空的生命.
耶路撒冷作王大衛的兒子.
他就是傳道者.
我們知道這個人是誰.
馮檢舊的書裡面.
講到大衛的兒子.
就是所羅門.
所羅門.
為什麼會跟智慧有關係呢?.
或者跟學識有關係.
跟權力也有關係.
因為在所羅門的人生裡.
他是最有.
最有學識.
最有財富.
最多派對.
你記得勝經怎樣介紹他.
他有飛奔.
七百三百個.
一千個.
所羅門就好像今天的Bill Gates.
或者好像Playboy雜誌的Hefner.
你要的女人他全部都有.
或者好像今天的美國總統.

$^{201}$或者好像今天.
最有學識的人.
他有這些東西.
但他人生裡面的終結是怎樣呢?.
你記得舊約的歷史.
就是告訴我們一個這樣的事實.
他所建立的國度.
死了之後就分裂.
然後就.
北國沒有了,南國也沒有了.
他就失去了一切.
這叫做敗家.
全部東西都沒有了.
然後他自己也都歷經背道.
甚至他的信仰也都可能失去了.
我小時候.
我們常常都會帶他去.
Nick Shaw那裡.
在湖邊那裡他們就很喜歡玩沙.
你們小時候.
都是做這些事.
放假的時候就帶他出去玩沙.
他們在岸邊.
在水邊那裡.
就做一些sandcastle.
做完之後就做得很漂亮.
花了很多時間.
然後就一個浪蓋過來.
全部東西都沒有了.
他們就哭了.
又要再做.
又要做完之後.
又一個浪蓋過來.
這個就叫做虛空的虛空.
人生很多時候就是這樣.
是不是很濕濕的.
傳道者在傳道書裡面說的就是這樣.
他說如果我們有更多錢.
你想想.
你想想你自己人生.

$^{241}$如果我們有更多錢.
好不好.
好,沒有人嫌錢聲.
再多一點學識.
再多一點人尊重.
或者再多一點權力,再多一點快樂.
你是不是真的可以絕對滿足呢.
一浪蓋過來的時候.
就什麼都成零了.
沒有任何的.
中央.
傳道者就怕.
convince不到我們.
到最後那一章.
最後那一章他講什麼呢.
他說.
就算不想死亡的終結.
不要想.
這麼慘的事情.
你也要想想人生之後.
的時間.
他說你趁著年幼的時候.
衰敗的日子.
就在想衰敗的日子.
尚未來到就是你所說的毫無喜樂.
那些年日.
未曾臨近之先.
你當紀念你的主.
光明月亮.
星星變成黑暗.
雨後雲彩.
返回.
這個日子.
是什麼日子.
你想想你什麼時候.
會等到白天.
月亮星星變成黑暗.
就是老眼昏花的時候.
對不對.
老眼昏花.

$^{281}$我也有很多.
弟兄姊妹身邊的.
也多了幾副眼鏡.
又有老花.
看東西不清楚.
然後他講什麼呢.
他說他看守的房屋發顛.
看守的房屋發顛.
下一節他講.
第三節.
他說看守房屋的發顛.
有力的屈身.
推磨的稀少就止息.
從窗戶往外看的月光.
也都昏暗.
街門關閉.
推磨的響聲微小.
雀鳥一叫.
人就起來.
唱歌的女子也都睡眠.
在這裡講什麼.
房屋會顫抖.
其實是一個比喻.
他說我們的身體.
就好像房屋.
就是說我們的身體會顫抖.
我們身體會顫抖.
就是說我們的身體會顫抖.
你覺不覺得.
芝加哥越來越冷.
起碼主觀性地.
你覺得越來越冷.
如果寒冬來到的時候.
你覺得很慘.
有力的屈身.
以前很大力的.
以前是很多力量的青年.
到你老去的時候.
你都要面對一個日子.
就是你屈身.

$^{321}$你不會繼續高下去.
你只會矮.
因為你的脊椎會寒背.
推磨的稀少就紫色.
推磨你知道是什麼意思.
今天我們吃麵包.
大麥小麥有兩塊石.
然後磨它變成麵粉.
就做成麵包.
但這也是一個比喻.
什麼叫推磨.
推磨就是我們的牙齒.
是不是磨來磨去就越來越少.
對嗎.
你年紀大了就知道了.
要用牙齒做牙齒.
傳道者就是說這件事.
從窗往外看.
就昏暗了.
又說到我們的視野越來越差.
我們看的東西.
我都發現我工作的時候.
我的電腦.
我的電腦的螢幕.
我現在要的越來越大.
不是因為我想工作.
方便一點而是我需要.
你明白我的意思嗎.
不是我想要.
我兒子說你的螢幕這麼大.
我心想也不是我想要的.
只不過我開始眼睛.
不是那麼有力.
我就要越來越大.
你家裡的電視是不是越來越大.
哈哈.
昏暗.
街門要關閉.
不想出街.
聽到快樂的聲音.

$^{361}$推磨的聲音就微笑.
人.
什麼時候會被吵醒.
睡不著.
就著了一覺的時候.
三四點就醒了.
譬如說晚上.
我明明十二點前睡.
乖乖的我就去睡.
我預備今天事奉.
誰知道睡了幾個小時.
OK 睡得不錯.
看了一看時鐘.
三點鐘.
哈哈 幹嘛.
跟誰聊天.
一叫就起來.
唱歌的女子.
睡了沒.
快樂的聲音全部都沒有了.
我有一個親戚.
退休老了.
身體又開始軟弱.
每天.
十個多小時.
他就在家裡梳化.
看著電視.
然後.
拿著電話.
他就在按.
他就在轉台.
這就是.
衰老的日子.
唯一可以做到的動作.
就是按電話.
你一叫人就起來.
我們睡不著.
這些不是去神魂.
這些是避神魂.
你沒辦法.

$^{401}$你一定要去神魂.
第五節.
人不怕高處.
路上有驚慌.
下樓梯都怕.
有高處的時候.
你就怕.
杏樹開花很有意思.
杏樹是巴基斯坦地.
第一種樹開花.
全年頭的時候會開花.
白色的花.
但他在說的比喻是什麼.
就是你的頭髮.
有白頭髮.
你開花.
杏樹開花.
越來越多.
渣猛成為重擔.
明明很輕的東西.
現在手震.
捉一隻渣猛.
手震.
人所願的也是廢掉.
情慾也廢掉.
因為人歸自己.
永遠的家.
吊喪的.
在街上往來.
淫土的凝歸塵土.
靈的凝歸次靈的上帝.
我曾經看過一篇.
醫生寫的文章.
醫生用一個笑話的方法寫.
他說人的生命.
由出生到他.
十八歲的時候.
十八歲就是他人生成長.
最高峰的時候.
十八歲之後.

$^{441}$就一路走下坡.
你的細胞越來越小.
你就有衰老的跡象.
所以你過了十八歲.
不要那麼得意.
已經在走下坡.
不過只是慢慢走.
有些人可能走得快一點.
上年我去加州的時候.
探望以前一個導師.
他現在七十多歲.
開始有一隻耳朵聽不到東西.
不單止這樣.
他父母九十多歲.
他每一天.
要花五六個小時.
去照顧他的母親.
七十多歲.
有很多疾病.
有很多小的毛病.
又要每一天去照顧母親.
他覺得很辛苦.
有一天他在街上.
他遇到一個教會的朋友.
大家聊聊天的時候.
那個朋友就知道.
他父母九十多歲長壽.
我們華人會怎麼說.
長壽真的有福.
那位弟兄就很生氣.
他就立刻罵人.
你們這些老人家.
他立刻開口就罵人.
他說.
你怎麼知道長壽一定有福.
我每天去照顧他.
又抹屎抹尿.
要扶他去哪裡.
要去看醫生.
然後就破口大罵.

$^{481}$那兩天之後.
他就覺得很不好意思.
他就打電話給別人.
打電話給別人.
打電話給別人做什麼.
又罵人.
他道歉.
心裡面很煩.
說了這些話.
你想想.
我們是pro-life的.
聖經叫我們pro-life.
你想想什麼叫pro-life.
有沒有想過quality of life是什麼.
傳道者就要我們想這件事.
人生的終局.
如果只有子得今世的話.
老到一定是福.
前陣子.
有一個年輕的姐妹.
跟我說.
她說你的肚腩這麼大.
你什麼時候知道自己老了.
就是一個年輕貌美的女孩.
去評論你的身體的時候.
你是完全不覺得傷害的.
誰在乎.
我不在乎.
你知道你過了中年的時候.
你就會發福.
我現在都ok.
但我想告訴你.
我以前更加好的.
我兒子十五歲.
他不用運動已經六塊腹肌.
我怎樣運動都只是一塊.
但你知道你老的時候.
你覺得這些說話完全不傷害你.
你已經百毒不侵.
你幾乎不在乎.

$^{521}$這個就是傳道者說.
日光之下的生命.
如果是這樣.
就是虛空的虛空.
很多人會覺得.
He who dies with most toys win.
如果只有今世的話.
你應該怎樣過生活.
你希望你離場的時候.
你能夠有最多的玩具.
這個就是很多人的人生觀.
我有什麼享受.
我有什麼物慾的擁有.
那我就贏了.
這個就是中文叫.
離場的時候有最多的玩具.
那就是贏家.
但這個也是虛空的.
傳道者說.
你最後想清楚.
始終也是成零.
做人就是很濕.
Life sucks and then you die.
整本書好像很消極.
但傳道者要我們想一個最邏輯的結論.
如果人生只有here and now的話.
一切都是empty.
就像一口氣.
你呼完之後就沒有了.
現在清明時節.
我不是很check到夠力.
但我知道清明快到了.
你知道我怎麼知道嗎.
過了.
我不是很熟悉這些.
過了什麼時候.
剛好過了禮拜五.
我不是過很多.
你知不知道我怎麼知道.
因為幾個禮拜前收到一張卡.

$^{561}$當然有些人說.
你死了.
但你帶著最多的玩具還是會死.
誰會帶著最多的玩具.
還是會死.
我怎麼知道什麼時候是清明.
我收到這張卡.
這張卡很漂亮.
我開的時候.
這麼奇怪有人寄這張卡給我.
是從清明時節寄來的.
他說清明時節.
有些優惠派對.
有個邀請卡請你去.
我真的沒收過這些東西.
這個殯儀館和墓地都很關心我.
我們這些熟客.
因為我經常跟他們做喪禮.
我是牧師.
他們有我的資料.
真的寄這張卡給我.
真是大吉利是.
無端端寄這張卡給我.
不知道是生日卡還是新年卡.
原來是清明的卡.
宣傳單.
是時候買墓地了.
趁便宜就買了.
預備好.
這就是我們需要想的東西.
傳道者.
就是要跟我們說這件事.
你解傳道書.
這本智慧書的時候.
你要明白是有兩個的聲音在.
中間的傳道者.
就是告訴我們.
人生的終結只是.
現在的時候全部都是空的.
包著的是什麼呢.

$^{601}$是父親.
父親的智者.
智者跟孩子說.
你聽不聽到傳道者說什麼.
他說的話是100\%正確.
100\%對.
你要聽清楚.
所以對他有叮囑.
有叮囑.
有叮鈴.
傳道者跟你說真心說話.
日光之下做人就很濕濕的.
智者就要我們看.
如果我們看日光之下的話.
什麼都是虛假的.
但最後的時候他說什麼呢.
他說這件事已經聽見了.
那個種意就是這樣.
不用再聽了.
請慰神.
謹守祂的誡命.
這是人在世當盡的本分.
因為人在世所做的.
連一切隱藏的事.
無論是善是惡.
神將來.
必定會有審問.
他說有永遠.
我們不要只想著生活在陽光下.
只是在日光之下過活.
要有一個新的角度去看我們的人生.
是在陽光之上過活.
要有一個永恆.
上主的角度去看我們.
看我們今生今世要過的生活.
他說敬畏神.
第一句說話.
第二是謹守祂的誡命.
第三是知道將來會有審問.
敬畏神.

$^{641}$這個一定是智慧文學經常提及的說話.
其實是在說我們和上帝有一個正確的關係.
一個正確永約的關係.
敬畏神不是敬而怕.
然後避開祂.
而是我們敬畏祂的時候.
我們來到祂面前.
來事奉祂.
去尊敬祂.
這個叫做fear God.
第二件事就是謹守祂的誡命.
忠於祂的約.
忠心地去跟隨祂.
有審問就是以永恆將來.
Live in light of eternity.
來看我們今天的人生.
有一天我們都要結帳.
要見上主.
上主怎樣看我們的人生.
做了多少事.
是有永恆價值.
我覺得這三方面.
正正是整本舊約聖經裡面所說.
當然耶穌基督包含了新約.
舊約裡面說什麼呢.
敬畏神就是智慧文學.
還有那些叫做writing.
希伯來聖經有三部分.
第一部分就是writing.
裡面有很多智慧文學.
傳道書是其中一部分.
敬畏神.
第二方面.
謹守祂的誡命.
就是律法書.
頭五本摩西五經.
講到上帝的律法.
然後將來有審問.
這裡是講將來的摩西.
就是先知書.

$^{681}$先知和歷史書都是神的代言人.
跟我們說話.
智者就說不要只想.
人生只有life under the sun.
如果是這樣的話.
每樣東西都是虛度.
要想life above the sun.
就是壽苦節和復活節的意思.
不是一年一次我們去一次教會.
給面子.
有些人真的這樣.
Easter去給面子.
穿西裝.
去一下sunrise service.
你知道.
壽苦節和復活節.
就是告訴我們人生有另一個更高的現實.
這個更高的現實.
不是life under the sun.
是life above the sun.
復活節告訴我們.
死亡不是人生的必然句號.
不一定死亡是我們最後一句說話.
永生才是我們最後一句說話.
因為上帝.
藉著他的兒子.
在墳墓裡面復活.
空的墳墓就告訴我們.
有更高的存在的現實.
這就是壽苦節和復活節的意思.
我們最近的老朋友.
在我們的chat group裡面.
我們收到一個消息.
就是以前.
所謂看著我們長大的一位女傳統.
她現在輾轉之下.
她去了香港.
她有二臟癌.
大家聽到二臟癌就知道是甚麼事.
是這個意思.

$^{721}$她要做一些手術.
我們一班大英姐妹就幫她籌款.
如果只有life under the sun的話.
每樣東西都是空的.
為甚麼?.
但如果有life above the sun.
有一個eternity的話.
她所做的事.
以往做的事.
就算是這樣終結的話.
都是有價值的.
分別就大在這裡.
一個乘以零.
一個乘以百倍,千倍甚至萬倍.
因為有eternal value.
傳道者說敬畏上帝.
這是你當盡的本份.
這才是生命裡最重要的東西.
先放好這件事.
其他東西就會come in place.
當你和上帝有一個正確的關係的時候.
其他生命裡的優次,次序.
全部東西都會fall into place.
我們同心低落.
主我們多謝你.
因為你對我們的愛.
縱使我們人生裡面日光之下.
很多時候我們會受到很多苦楚.
有很多難處.
但我們知道你是我們的盼望.
在耶穌基督的十字架,九印.
他的空墳墓裡面給過我們.
神愛多謝.
多謝你對我們的恩典.
也讓我們現在在生命當中.
我們以敬畏你.
守你誡命.
為我們的優次.
我們優先要做的.
就是和你建立一個正確的關係.

$^{761}$主我們.
求你讓我們的生命來事奉你.
按照你的話語而行.
讓我們能夠成為一個敬畏神.
以致我們生命裡面.
雖然短暫.
我們所作的一切.
是能夠有永恆的價值.
能夠永耀神利.
我們感恩,讚美,祈禱奉主美金.
\newpage



\section{}
\label{sec:JN7FdDAB0F8}
\textbf{4/28/19 主日講道\_耶穌天國的小故事-1\_撒種(陳崇基牧師)}
\newline
\newline
連結: \href{https://youtube.com/watch?v=JN7FdDAB0F8}{\texttt{https://youtube.com/watch?v=JN7FdDAB0F8}} ~~~~ 語音日期: 2019-08-05
\newline
\newline
\hyperref[sec:QpmQmJgiMg8]{\small{< < < PREV SERMON < < <}}
~
\hyperref[sec:index_chronic]{\small{[返順時目]}}
~
\hyperref[sec:index_scriptual]{\small{[返順卷目]}}
~
\hyperref[sec:tJHkj9NzTVE]{\small{> > > NEXT SERMON > > >}}
\newline
\newline
$^{1}$剛才那一下是想大家醒的.
今天我們開始一個新的topic.
叫做天國的小故事.
我不知道大家成長當中有沒有聽過一些short story.
或者一些很簡單的故事.
可能是衣索寓言等等.
有沒有?.
有吧?.
今天開始這個topic.
我們希望能夠看看聖經裡面的小故事.
一些short stories.
我們叫它小故事.
其實聖經裡面很多時候叫它做比喻.
parable.
什麼是parable呢?.
就是一些很小的.
耶穌講的一些比方.
如果你記得的話.
聖經裡面有不少這些比喻.
我們先看看聖經.
馬太福音第十三章.
耶穌又設一個比喻對他的門徒說.
天國好像一粒芥菜種.
有人拿去種在田裡.
這原是百種裡面最少的.
等到長起來.
卻比各樣的菜蔬都大.
且成了樹.
天上的飛鳥也都縮在他的肢上.
就是這麼簡單.
今天我們來看一個short story.
一個小故事.
一個比喻.
大家知道.
你打開新約聖經.
馬太,馬可,盧迦.
三本的福音書記載了耶穌很多的比喻.
因為沒有什麼比喻.
頭三本才有這麼多.
有這些比喻.

$^{41}$我想考考你.
我們先做個測驗.
等下我給你一些比喻的名字.
或者一些小故事的名字.
你告訴我哪個故事是耶穌有說過.
哪個故事是耶穌沒有說過的.
可以嗎?.
準備好了嗎?.
我們試一試.
先給一個text.
第一.
殺種的比喻.
第二,布恩的老鼠.
第三,麥田敗子的比喻.
第四,是麵烤的比喻.
第五,葡萄園裡面工人的比喻.
第六,迷陽的比喻.
狼來了的比喻.
故事.
十個童女的故事.
好殺馬尼亞人.
以及寡婦和法官的比喻.
哪個short story不是耶穌說的?.
哪一個?.
有好幾個.
哪一個?.
老鼠.
老鼠不是耶穌說的.
布恩老鼠,有沒有聽過?.
不記得?.
不記得?你們成長這麼差的?.
成長背景.
那個獅子想吃老鼠.
放了牠.
獅子遇難.
被網綁住了.
老鼠經過看到牠.
咬開牠,幫牠.
類似這個意思.
記不記得?.

$^{81}$不記得?.
難怪你不幫人.
那個故事是教我們幫人.
布恩老鼠不是耶穌說的.
OK?.
It's a good story,但不是耶穌說的.
還有哪個不是耶穌說的比喻?.
狼來了,大家有聽過?.
有誰沒聽過?.
沒聽過,你童年少了很多樂趣.
還有沒有一些耶穌沒說的比喻?.
其實只有這兩個.
當然.
你看聖經,看經文記載.
你就知道那個故事是什麼.
那個title叫什麼.
其實都有斟酌的.
因為這些是後者.
我們聖經加上去.
寡婦和法官的比喻是有的.
不過你未必叫他做寡婦和法官的比喻.
但總之.
這些比喻都是常見的.
什麼叫做比喻?.
比喻是一些人間的故事.
特別是耶穌說的比喻.
是人間發生或人間能夠看見的事.
用這些來述說天國上神的事情.
神的工作.
神的道理.
還有很多時候比喻是間接的.
是 indirect communication.
為什麼呢?.
因為聖經常說的是 indirect communication.
不是直接告訴你你要這樣做.
而是用一個故事形式來告訴大家.
天國是怎樣的.
神的道理是怎樣的.
耶穌也是這樣.
你知道耶穌很少會直接回答問題.

$^{121}$有人說在福音書記載.
耶穌說的答案只有兩次.
很多時候人家問了問題.
或者他遇到一個情況或處境.
他就會間接說出.
我講一個比喻給你聽.
為什麼要這樣做呢?.
因為除了當時第一世紀的拉比.
猶太人的老師叫拉比 Rabbi.
他們會這樣回答一些問題之外.
也是用間接的方法來叫醒我們.
因為很多時候直接說是沒用的.
一定要繞個彎來踢你一腳.
這樣你才會醒.
比喻的意思就是這樣.
就是要把你抽進那個故事裡面.
然後你想一想那個故事裡面有什麼道理.
原來這麼多年我都不記得.
這麼多年我都沒有想起.
你想想耶穌有多久.
可能你告訴別人很久.
但你可能一直睡覺.
比喻就是用間接的方法來叫醒你.
大家記得舊約裡面有一個故事.
大衛作為君王的時候.
平定四景.
他下午起床.
他見到一個女人叫做拔士巴.
然後就和她犯姦淫.
是一件很大的罪.
之後誰去提醒她.
她做錯了事.
不是突然晴天霹靂.
是有一個仙子叫做拿丹.
他就去見大衛王.
也不是直接跟他說你做錯了事.
他跟她說了一個比喻.
記不記得你回去看看.
說完比喻之後大衛王就中招.
他就說那個人怎麼可以這樣貪心.

$^{161}$拿丹就直接一句跟他說那個就是你.
所以比喻的目的就是這樣.
間接地令我們進入一個故事當中.
讓我們去改變.
如果不是直接跟我們說是沒有用的.
還有很多時候比喻不是給我們一些娛樂.
比喻 耶穌說的比喻.
是去挑釁我們.
廣東話叫做串你.
是故意去串你.
我們特別是上了教會很久的人.
看了聖經很久的人.
信了耶穌很久的人.
我們常常覺得耶穌好像被馴服了的寵物.
很乖很紳士.
但其實很多時候不是這樣.
耶穌說的比喻是有很多挑釁性.
是去引發我們一些反應.
聽完之後是很不舒服的.
其實是要你醒起來.
睡覺的時候要醒.
聽完之後很不舒服.
然後你才會有一個新的開始.
這個就是比喻的目的.
我們會說很多的比喻.
我想舉個例子.
有一個比喻叫做浪子的比喻.
有沒有聽過.
有一個人有兩個兒子.
小兒子跟他說.
不如我們現在分家產.
爸爸就跟他們分家產.
分給兩個兒子一個大一個小.
小兒子就去花天酒地.
然後很窮.
浪費了錢.
他就在那裡餵豬.
很慘的時候他就跑過來了.
這個小兒子.
不如回家吧.

$^{201}$我做我爸爸的傭人.
好過在那裡捱苦.
在那裡吃豬哨.
於是他就真的回家了.
原來他父親在那裡等他.
接他回家.
就去接他.
整個比喻有沒有留意.
沒有用過愛字.
love.
從來沒有出現過整個比喻.
很多經文.
沒有說過愛.
但是裡面說的.
就是天父.
沒有說過愛字.
等於你跟你太太沒有說過愛字.
你做了很多事.
說了.
比喻就是這樣.
就是跟我們說天父的愛.
但是故事沒有停止.
最重要的高峰在哪裡.
就是大兒子.
原來大兒子的兄弟走了很久.
他一直都在田裡工作.
很勤力的工作.
回來為什麼大家這麼開心.
歡迎浪子的小兒子回來.
他就很生氣.
然後父親跟他說了一大輪的事情.
就告訴他.
這是失而復得的一個兒子.
你不明白我在想什麼嗎.
你跟我相處了這麼久.
你不知道你父親在想什麼.
那個就是浪子的比喻.
你有沒有留意.
這個比喻在路加福音第十五章.
到底想挑釁我們什麼.

$^{241}$其實這是路加福音第十五章.
三個比喻裡面第三個.
我們可以叫做失子的比喻.
就是有個兒子不見了.
現在找回來.
第一個比喻是.
是故意將那三個比喻放在一起.
失羊的比喻.
你記得吧.
有個牧羊人.
有一百隻羊.
他不見了一隻.
那他怎麼做.
他就留下那九十九隻羊在羊圈裡.
周山去找他.
找到了.
他就回來跟其他鄰居和朋友慶祝.
因為找到這隻.
Loss and found.
這就是失羊的比喻.
然後接著耶穌又說到一個比喻.
失錢的比喻.
有個女人.
她有十個金幣或十個銅錢.
十個裡面不見了一個.
她就怎樣.
反轉房子到處找.
找了很久.
然後真的找到第十個不見了的金幣.
她就叫她的朋友.
身邊的鄰居來.
說我們慶祝一下.
十個金幣裡面找到一個.
然後才說這個失子或浪子的比喻.
有沒有留意三個的比喻.
都是跟號碼有關.
大家很數口.
大家是壞人.
一百個裡面不見了一個.
有沒有試過數一百.

$^{281}$由一數到一百.
睡不著的時候有沒有試過數羊.
數著數著.
二十幾就忘記了數了多少.
又重新數.
數一百.
就是牧羊人去數羊的時候.
一二三四一隻.
回羊欄的時候.
五六七八九十三十二.
去到最後好像少了一隻.
又回頭再數.
是不是真的數錯了呢.
又再數.
真的不見了一隻.
那他就出去找她.
是跟號碼有關的.
跟數字有關.
失錢的比喻呢.
有沒有試過數自己的錢.
錢包的錢.
或者銀行的號碼.
以前我們有很多現金.
很多神沙的.
有些零錢的時候.
我們不是經常數嗎.
一個兩個三個四個五個.
那個女人就是在數這個.
真的好像少了一個.
以前的錢幣很小.
好像手指尾的指甲那麼小.
數著數著.
總是差一個.
原來真的不見了.
就反轉房子.
就找回那個不見了的.
第十個的錢幣.
失紙的比喻.
或者浪紙的比喻.
兩個.

$^{321}$有沒有試過數二呀.
一和二呀.
咦少了一個.
一百個少了一個.
和十個少了一個.
到最後是兩個少了一個.
其實在說什麼.
三個都是跟數字有關.
並且三個越來越容易數的.
一百個就很難數.
十個就容易一點.
但是到兩個呢.
沒理由不會數的.
除非你白痴.
你會數一個一.
下一個數字是三.
沒理由不會數.
這樣你都不會數.
這是最簡單的數目.
兩個裡面不見一個.
你應該知道.
你知道在故事失紙的比喻裡面.
誰不會數.
不是爸爸不會數.
我想不應該是爸爸.
爸爸在等他回來.
是大兒子不懂得數.
家裡少了一個人.
你不知道嗎.
你回家少了一個人.
你知不知道.
你快點報警吧.
為什麼大兒子居然不會數.
想挑釁什麼.
比喻不只是說一個醫學語言那麼簡單.
一定是有常文下的.
因為耶穌答問題是用比喻來答.
這三個比喻答什麼問題呢.
第十五章第一節就開始的時候.
為什麼會說這三個這麼好聽的比喻.

$^{361}$這樣的故事.
第一節.
他說縱需你和罪人都靠近耶穌.
要聽他講道.
很多社會唾棄的人.
他就是靠近耶穌.
想聽耶穌講道.
法利錯人和文士就私下議論.
這個人.
接待罪人.
又和他們一起吃飯.
心裡面在說.
或者自己自尋在議論.
說耶穌這個人不行.
接待這些罪人.
社會唾棄的人.
耶穌用比喻跟他們說.
說什麼.
一百個少了一個的失陽比喻.
十個錢銀少了一個的比喻.
然後就是失子的比喻.
兩個兒子少了一個兒子的比喻.
就是想挑釁什麼.
就是說你.
你會數東西.
牧羊人都會數.
女人都會數.
你做大兒子居然不會數.
大兒子是誰.
就是那些所謂宗教的人士.
以為自己信了主很久.
自己都懵然不知.
連數最簡單的數目都不會數.
一加一是二.
二減一是一.
連這些你都搞錯.
不明白上帝的心意.
上帝的心意是要一個人悔改也好.
要一個人迴轉也好.
你居然不懂上帝的父親的心意.

$^{401}$枉你一直在家裡面.
心不在.
你來了這麼多次教會.
你都不明白天父的心意.
你翻了這麼多聖經.
記了這麼多聖經.
你居然不明白上帝的心意.
一個人的迴轉對上帝是多大的歡樂.
這就是挑釁.
就是搖一搖醒你.
到底你關心的是什麼.
我來到這裡.
我的朋友在這裡就不可以了.
我認識的人在這裡就不可以了.
你有沒有想過那些沒有在這裡的人.
一個都好.
你都要掛念他.
你去為他祈禱.
你去找他.
這就是重點.
講道理.
不是要我們聽完之後.
不,上帝說你瘋了嗎.
一.
再數,二減一是多少.
是一,少了一個.
對上帝來說是重要的.
所以每一次我們看比喻的時候.
想一想,到底想挑釁你什麼.
想挑釁你什麼.
想弄醒你什麼.
今天我們看這個比喻.
叫做「戒菜種」的比喻.
我們再看,剛才我們讀了.
耶穌有一次比喻對他的門徒說.
他說天國好像一粒的戒菜.
有人拿去種在田裡.
這原是百種裡最少的.
等到它長起來的時候.
他說比各樣的菜更大.

$^{441}$且成了樹.
天上的飛鳥來.
屬在他的枝上.
他說什麼像戒菜種.
天國像戒菜種.
Kingdom.
你一講到國這件事.
在當時耶穌的年代.
或者可能今天也是.
國這件事是一個很敏感的詞語.
Kingdom.
因為誰最常用Kingdom這個詞語.
就是羅馬帝國.
羅馬帝國是用強權.
用飛機大炮去統治的.
用戰馬.
用槍去統治.
去打壓那些小民.
耶穌的門徒.
耶穌的群眾就是被打壓的人.
耶穌的門徒.
他會立很多惡的法.
來去統治他們.
來去打壓他們.
這個就是國.
Kingdom.
耶穌說的是天國.
天國有什麼不同呢.
天國就是.
自然界的東西.
是一粒的.
戒菜種.
神的國.
上帝的統治.
上帝的工作.
就不同羅馬帝國.
羅馬帝國會打鐵,打劍,打刀.
打盾牌.
但天國是什麼呢.
天國是很自然的.

$^{481}$天國是很自然的東西.
他突然就帶你去好闊.
你看看天國是怎樣的.
他就拿戒菜種.
我就去了馬里安諾.
買了一瓶戒菜種.
有沒有人買過.
有沒有人需要.
有需要我給你一粒.
很小的戒菜種.
戒菜種有黑有黃的.
他就跟你去市場那裡看看.
是很小的.
我今天想.
既然講了比喻.
我想跟你玩個遊戲.
我們這裡這麼多人.
大家玩個遊戲好不好.
我跌了一粒.
很貴的.
但我想跟你玩個遊戲.
我們這裡分四楷.
這楷就.
給你幾粒.
我給頭一個人.
我給他一兩粒.
串到後面.
看誰最快.
好嗎.
我先給Jelissa.
不要串著我.
不要跌.
Jimmy.
一隻手都可以.
George.
不要跌我.
不要串著.
幾多粒.
你先數數.
公平一點.

$^{521}$好.
開始串了.
最後那個阿Ron那裡有16粒.
好來.
是一致的.
串過去.
你們串到哪裡.
(沒有十多粒).
你一看就知道沒有十多粒.
果然是.
那天買慣送.
沒有十多粒.
串過去.
好像.
這裡串到哪裡.
聽不到你說話.
這裡串到哪裡.
跌了在地上.
不要撿.
這裡呢.
去哪裡.
什麼.
冷靜一點.
這裡去哪裡.
OK.
你們解釋少一點容易一點.
幾多粒.
看看這裡幾多粒.
(原本只有15粒).
原本15粒變16粒.
OK.
多了.
這裡那粒去哪裡.
在哪裡斷了.
為什麼會跌了.
不知道.
太小了.
這一戒呢.
哪裡開始不見.
在哪裡.

$^{561}$沒有了.
這邊呢.
串到最後.
這是神石.
神石.
串串串.
先弄一粒神石出來.
他在說這件事.
那粒戒菜種很小.
小到.
大家知道.
種子很小.
像什麼.
像鼻屎.
真的.
我不知道串那粒是不是戒菜種.
中間變了.
不奇怪.
他說天國像什麼.
天國像戒菜種.
當然不是.
百種裡面最少.
不是全世界的種子裡面最少.
有很多種子比戒菜種更小.
肯定.
當然講比喻.
有些誇大.
告訴你那個小的程度.
有多少.
或者等於我們今天講廣東話.
蚊子那麼小.
他說天國就像一粒戒菜種.
但我都很奇怪.
為什麼耶穌會用.
這一種植物.
來代表天國.
代表上帝的國.
你去巴勒斯坦地.
或者以色列.
你經常看到.

$^{601}$到處都有的戒菜種.
的所謂樹.
由一粒那麼小的.
黃色或者黑色.
的種子.
生出一棵大概可以.
五六尺高的.
戒菜種.
的菜.
也可以去到六尺至八尺高.
是.
經常會出現的.
你去以色列.
你經常留意那些景色的時候.
就是遍地開花.
黃花.
特別是現在春天的時候.
起初是很小.
但慢慢會出現黃色的花.
到最後可以有.
一個人那麼高.
這個是還沒開花的.
戒菜種.
最後天上的飛鳥.
是可以.
在上面.
去棲息.
不是築巢.
是棲息.
那隻牛.
你見到後面那隻牛.
牠就躲在那些戒菜種黃花的裡面.
小鳥可以在這裡.
來棲息.
他說天國.
好像戒菜種.
將來大了.
就可以變成這樣.
並且遍地會有黃花.
但你要留意.

$^{641}$他最重要.
要你留意的.
是將來.
怎樣遍地黃花.
也不是將來戒菜種有多大.
他要留意的就是.
天國好像一粒戒菜種.
就是說天國上帝做事.
是由小開始的.
小到不顯眼.
小到你不覺得.
他有做過事.
但這個就是.
上帝國度奇妙的地方.
你想想.
天國.
福音.
教會.
公義.
真理.
他的散播.
是由一粒戒菜種開始.
一個很小的.
開始.
一個你完全不覺得有意義.
但是他是.
這樣開始.
下星期我們有一個客席講完.
我們.
也很榮幸認識到他.
Feruzur Krasin.
是一個姐妹.
她是在中亞.
一個回教.
或者伊斯蘭.
國家裡面長大.
在那裡出生在那裡長大.
17歲的時候.
她就信主.
是一個很.

$^{681}$戒菜種式的信主的方法.
就是這麼簡單.
那個英文班.
是用聖經學英文的.
就認識了一些基督徒.
有些人.
跟她建立了關係.
她就被基督教的信仰.
耶穌吸引.
你聽她的故事的時候.
你幾乎想不到.
有什麼大事發生.
她現在就.
成為了一個.
她現在是麻省.
她就成為了一個.
在聖經翻譯的.
事工裡面.
她幫人翻譯.
她認識六種的語文.
她就慢慢地.
幫忙翻譯.
使那裡的人能夠看到聖經.
是一個很小的動作.
認識一下她.
我也沒有什麼把握.
她會不會相信呢.
誰知道她就相信了.
天國的工作就是這樣.
不是一開始.
就放煙花.
而且很多時候沒有煙花放.
只不過是一粒一粒的.
芥菜種.
她來工作.
並且你留意一下.
芥菜種.
她是很粗生的.
最近我也種了一些東西.
可能年紀大了.

$^{721}$其實不是年紀大.
是因為我覺得蔥很貴.
你覺不覺得蔥很貴.
為什麼這麼貴.
搞個蔥來.
貴到多少.
一元一束.
神經病.
你放在麵裡.
就把蔥蔥了.
我就去youtube學一下.
原來這麼容易.
把蔥頭浸在水裡.
然後生.
你看到後面飄出來的東西嗎.
我就是吃那些蔥.
芥菜種.
就是很粗生的.
我剛才說過.
在以色列裡面.
有很多很多田園的景色.
芥菜種就是經常遍滿路邊.
大家知道耶穌的比喻.
有很多都是跟好闊有關的.
是有機的.
是跟植物有關.
跟菜蔬有關.
葡萄園.
葡萄園的工人的比喻.
我經常說.
種葡萄樹是很浪費功夫的.
因為你要修剪那些樹枝.
讓那些葡萄出到.
好的葡萄.
然後好的葡萄才能做到好的葡萄酒.
釀酒就是靠那些葡萄.
也要靠那些天氣.
薩種的比喻.
麥田裡面的比喻.
是有耕田的.

$^{761}$耕田是要很多的人來做事.
要下很多的功夫.
或者應該這麼說.
很容易拿到功勞.
你看這塊田是我種的.
他出的菜.
裡面的麥子.
都是因為我的.
努力而成的.
但是.
芥菜種有什麼不同.
沒有人會去種一個.
芥菜種的田.
起碼古代沒有的,現在有.
現代有.
他就是.
被風吹散.
然後去到不同的地方.
就有芥菜種出來.
就是這麼粗生.
你看看比喻的第一句說什麼.
三十字.
他說天國很像一粒芥菜種.
有人拿去種在田裡.
田裡面不是那些菜田.
田的意思是野地.
路邊.
那些很粗生的地方.
他說有人.
拿去種在田裡面.
野地裡面.
誰啊.
就有個人.
什麼人.
任何人.
Anybody.
或者他應該是一個沒有.
留名的一個人.
有人拿去種在田裡面.
這個人就是一個任何人.

$^{801}$或者叫做老人甲.
就是說.
天國的擴張.
他的擴散.
你能夠看到滿地遍地黃花.
看到天國擴張.
誰的功勞.
沒有人有功勞.
沒有人可以說拍攝幕後.
因為我努力.
我種了多少田.
才可以遍地黃花.
沒有人可以做.
因為.
芥菜種的成長.
是不需要你去搞它.
不需要你去.
下手去做事.
人.
是沒有什麼可以拿獎賞.
天國的擴張.
其實就是這樣.
你有沒有留意.
教會的發展.
團計的發展.
天國福音的發展.
我們阻礙地球轉是什麼.
就是人.
就有個人不種.
是不是.
有個人不肯去做事.
那些叫做阻礙地球轉.
為什麼我沒份上台.
為什麼我不可以做主席.
為什麼要你做.
不是我做.
為什麼.
是你有權我沒權.
這些就是阻礙地球轉.
我們很渴求.

$^{841}$人的共識.
共認.
我們自己是舉足輕重的.
耶穌.
就在好福那裡.
打下來.
你願意做個.
路人甲嗎.
任何人.
只要你種的話.
將來就會遍地開花.
但是功勞不在你那裡.
你只是一個人.
沒命的.
在紐約有一個交響樂團.
有一個以前很出名的指揮家.
那個指揮家.
有一次訪問裡面.
有人問他.
他說你覺得在交響樂團裡面.
這麼多件樂器.
最難的是哪一件樂器.
你覺得最難的是哪一件樂器.
他說最難的樂器.
就是第二件樂器.
就是第二件小提琴手.
你知道有一個.
是第一主席的小提琴手.
那個就拿了所有spotlight.
但那個最難的是什麼.
就是輔助他那個.
第二小提琴手.
這個是沒有人會留意的.
就是有人.
拿下賽種就在田裡.
就是那個人.
他就在田裡.
沒有人留意他.
但這個他最難.
天國裡面的事奉其實都是一樣.

$^{881}$是不是.
有人說.
在聖經裡面你翻覽聖經.
領袖leader這個字.
出現過六次.
但僕人呢.
出現過九百次以上.
事奉就是second foot.
我不會拿credit.
我會推人家上台.
我會推上帝上台.
second foot.
有人拿去還在田裡.
誰.
誰care.
只要上帝的國.
是遍地黃花開.
誰拿credit.
不是最重要.
你有沒有想過.
戒賽種的信心就是這樣.
有一個你很關心的人.
你很希望他能夠信主.
是不是一定要你帶他信主.
可能你是為他祈禱.
誰知道你祈禱.
沒有人知道.
但為什麼可以有人得贖.
因為你做過小小的事.
不明顯的.
不顯眼的.
你的禱告.
是很歡迎.
或者你和人建立關係.
有想過我們去實行上帝的旨意.
其實很多時候都是戒賽種的方法.
你去愛一個人.
你去關心一個人.
你去建立一個家庭.
很多都是一些很瑣碎的事.

$^{921}$但卻是戒賽種的事.
為什麼世界可以有很多地方.
社會又公平又公義.
又自由.
是不是每個人都很出名的.
可以推動某些事.
很多時候不是.
只是一個很小的種子.
一朵很小的黃花.
但放在神的手裡.
上帝就使他遍地掌門.
讓我們今天的事奉.
帶著戒賽種的心態.
我們請勁八隊上來.
我們唱一首回應詩歌.
讓這首回應詩歌.
去鼓勵我們.
在上帝面前我們做一個戒賽種.
我就是那個.
某人.
將戒賽種放在田裡.
然後我希望.
上帝去使用.
假以自然.
上帝會使他的國度.
遍地黃花.
但你想不想做那朵黃花.
我們唱一首.
\newpage



\section{4:6-8}
\label{sec:tJHkj9NzTVE}
\textbf{7/29/18 主日講道: Running The Race - 2 Tim 4\_6-8 (Rev. John Yi)}
\newline
\newline
連結: \href{https://youtube.com/watch?v=tJHkj9NzTVE}{\texttt{https://youtube.com/watch?v=tJHkj9NzTVE}} ~~~~ 語音日期: 2019-08-05
\newline
\newline
\hyperref[sec:JN7FdDAB0F8]{\small{< < < PREV SERMON < < <}}
~
\hyperref[sec:index_chronic]{\small{[返順時目]}}
~
\hyperref[sec:index_scriptual]{\small{[返順卷目]}}
~
\hyperref[sec:pK11m4fqws8]{\small{> > > NEXT SERMON > > >}}
\newline
\newline
$^{1}$Well, before we go into God's Word, let me pray for us..
Father, it is your Word we are going to read and study..
So we trust you will give us your Word and it will return to you having accomplished.
your will and your purposes..
Be glorified as your Word is preached and proclaimed..
We pray in Jesus' name..
Amen..
So we are in 2 Timothy chapter 4, verse 6-3..
So 2 Timothy is believed to be one of the last letters that Paul wrote..
And it was written to his son and his apprentice, his disciple, Timothy..
So you know, it's been said that when people are at the edge of death, that their life.
flashes before their eyes..
So what do you imagine are the final thoughts you're going to have when you come to that.
day?.
Will you remember a life that has lived well?.
Maybe you'll have dreams that you remember, dreams that have been forgotten, plans that.
are not fulfilled..
Perhaps you'll have no regrets at all..
Or maybe you'll have many regrets..
And then what about what's on the other side of death?.
Do you think about what will be waiting for you?.
The Apostle Paul knows that he is about to face the end of his life..
So he tells Timothy these final words..
And it's from his second letter in chapter 4, verse 6..
And he says, "I am already being poured out like a drink offering, and the time has come.
for my departure..
I have fought the good fight, I have finished the race, I have kept the faith..
Now there is in store for me the crown of righteousness, which the Lord, the righteous.
judge, will award to me on that day..
And not only to me, but also to all who have longed for His appearance.".
So Paul says, "I am already being poured out like a drink offering.".
And I believe he is talking about a drink offering, referencing a drink offering that.
is talked about in the Old Testament..
And this drink offering involves wine being poured on a hot altar..
So when the alcohol and the sugars burn up, there's this wonderful smell in the air..
An offering that goes to God..
And the idea is that God would smell this wonderful smell, and he would be worshipped..
So yes, Paul is talking about that he is going to die soon..
Kind of like pouring out a cup and saying, "It's over.".
But he's also referring to, "My life is being poured out like worship to God.".

$^{41}$So he's not just saying, "My life is over, I'm done.".
But he's saying, "Look, my life is worship to God.".
And Paul talks and uses language that we would hear in athletics, in sports, and he's.
done this in other places too, but he does it here..
So he says, "I have fought the good fight like a boxer.".
I have finished the race like a runner..
And then he makes it clear, he's not just talking about physical activity, and he says,.
"I have kept the faith.".
In other words, Paul is saying, "Look, my life is filled with so much evidence that.
my life was lived for Jesus.".
This life is a wonderful scent to God..
But I ask, what about me?.
What about you?.
Is your life, is my life so filled with proof that we have lived it well for Jesus?.
That you and me, we could both say confidently, "Look, my life has been lived for worship.
of God.".
God would smell the scent of my life and be, "Ah, please.".
Think about what would people who know you say about your life?.
How about a coworker, or a neighbor, or a friend?.
Maybe a family member..
Would they say, "Are you a Christian?".
What if they came and surprised you, they visited you at this church, didn't know you.
go to church here, and they find you and they say, "Wow, you're a Christian?".
Would they say, "You know, I thought you might be because of the way you act.".
I thought there was something different about you because the way you are always so kind.
to me and other people..
I thought maybe you were a follower of Jesus because the way you live, it reminds me of.
Jesus..
Or would somebody find you at church and say, "What?.
You're a Christian?".
The way you're always complaining about the boss at work..
How mean you are to the neighbors..
Your language, the way you curse all the time and say evil words..
I had no idea you were a Christian..
What would people say of us?.
See our lives, our speech, our attitude, our behaviors, these can tell people that something.
has happened to you..
When Jesus has taken over our lives, people might ask, "What's going on?".
And what an opportunity to tell them, "It's because of Jesus.".
Let me tell you about Jesus..

$^{81}$He's loved me..
He's changed me..
He's made me new..
I hope that our lives would invite people to ask..
Now let's not forget as we look at this passage that Paul knows that he's going to die too..
He says, "The time of my departure has come.".
And he's not just going to a new city..
He's leaving this world and he will be in the presence of God..
But Paul is not worrying about it..
He's not fearful of it..
I think he's excited about it..
That Paul sees this as a release..
He's being untied..
From this body of death..
In another letter he writes to the Romans in chapter 7, he says, "What a wretched man.
that I am.".
This is one of the holiest men in the world..
And he says, "What a wretched man that I am.".
He says, "Who is going to save me from this body of death?".
And what does he say?.
"Thanks be to God through Jesus Christ our Lord because Jesus has delivered me from this.
body of death.".
So most of us know somebody who seems to be trapped in their body..
Maybe they're paralyzed or handicapped in some way..
Maybe they have some kind of physical weakness..
And their body has become like a prison..
And can you imagine that as we want to live, we also want to be free..
It's like somebody trapped in a burning house or in a car that's been in an accident and.
you're trapped inside the car and you're saying, "Somebody please help me get out.".
"Rescue me.".
Paul recognizes that Jesus has rescued him..
Has Jesus rescued you?.
Have you put your faith in the one who says, "I can save you?".
The one who says, "Not only I can save you, but I want to save you..
I want you to be free..
I want you to have a new life.".
You see, what we really need is not a doctor..
We need the great physician..
What we need is not someone who cares about us for a little while and then forgets about.
us once we're gone..

$^{121}$As much as you might have a great relationship with your doctor, do you think he thinks about.
you after you've left the office?.
Jesus is a doctor who never stops caring about you..
He never stops thinking about you..
He never stops working on you..
He is our creator..
He is the one who has made us..
So we don't need somebody who will just, you know, take a cut and fix it..
Or somebody who will give us temporary comfort or temporary pleasure..
We need someone who has conquered death and hell..
Jesus alone has conquered death and hell..
Death cannot touch him..
And he offers us eternal life..
This great physician is Jesus..
And he is the creator who loves us and cares for us..
He is the one who conquered death and Hades..
You see, Paul was this very religious man..
He thought he had done everything possible, everything necessary so that he could also.
have eternal life..
But then one day he met Jesus..
He was actually on the way to go and arrest and harm the people that love Jesus..
And Paul too believed and trusted Jesus..
He had lived such a perfect life that in Philippians, another letter he writes, he says, "If anyone.
else thinks he has reason to put confidence in the flesh," meaning the things he has done.
with his body, he says, "I have more confidence.".
And basically he says, "I have obeyed all the laws..
I come from the right kind of family..
And I'm from the right nation..
And the way I live has been with great passion for God.".
So passionate I even persecuted his church..
He says, "As far as the law is concerned, no one could blame me.".
What does he say?.
For Jesus, all of this is loss..
To all of Paul's Jewish peers, they would have said, "Oh, I wish I had his life..
I wish I could be righteous like him.".
But Paul continues and says, "Whatever was my profit, I now consider loss because of.
Jesus.".
And then he says, "What is more, I consider everything loss compared to the surpassing.
knowledge of knowing Christ Jesus, my Lord.".
He says, "I consider it all rubbish, garbage, junk, worthless compared to knowing Jesus.".

$^{161}$Wow..
Can you say that?.
That compared to knowing Jesus, everything else is just garbage?.
That without Jesus, everything is meaningless?.
That's why Paul says in verse 10 of chapter 3 of Philippians, "I want to know Christ and.
the power of his resurrection and the fellowship of sharing in his sufferings.".
Do you get that?.
Paul doesn't just say, "I want all the benefits of knowing Jesus.".
He says, "I even want to share in his suffering..
Jesus is so valuable..
I even want to know what it feels like to suffer like him.".
And then even becoming like him in his death so that he could have the resurrection..
I mean this sounds like crazy talk, doesn't it?.
But see this is not false humility..
Paul gave up everything because he saw that Jesus was such a treasure..
Everything else was junk..
You see Paul met Jesus on the road to Damascus and his life was changed forever..
That's when his race began..
That's when his fight began..
And Paul says, "I have fought the good fight.".
He says about the race, "I have already finished the race and I have kept the faith.".
It is believed that in ancient Rome and in ancient Greece that the athletes would make.
a promise..
They would devote their race or their performance to their God..
So if it was Zeus, that is their God, they would say, "For Zeus, I will run this race.
and I will win.".
And today we see this tradition carried out in the Olympics..
Now I won't read their oath but you can go online and look up, Google the Olympic oath..
What did they do?.
They promised that they are going to play fair..
They are not going to cheat..
They are going to compete honorably..
And there is no question in my mind that Paul is thinking of this..
All those who make promises to false gods, Paul says, "I have lived my life..
I have run my race..
I have fought my fight for the true God..
I have kept the faith.".
So having completed what he was charged to do, Paul reveals his confidence..
So look at 2 Timothy chapter 4, verse 8..
He says, "Now there is in store for me the crown of righteousness which the Lord, the.

$^{201}$righteous judge, will award me on that day, and not only to me but also to all who have.
longed for His appearing.".
Do you know what's waiting for you at the finish line?.
When the fight is done and the race is over, do you know what's waiting for you?.
When the ancient athletes won, they would receive a crown made of leaves..
And while they might look beautiful and fresh for a moment, in just a couple days they will.
dry up and blow away..
But Paul says, "I'm going to receive the crown of righteousness.".
This is an imperishable crown..
This is a crown that the Lord delivers..
And this is the crown that Paul says is not only for him, but it's for all of us who long.
for His appearing..
So what does it mean to long for His appearing?.
Have you ever missed somebody?.
Maybe your husband or your wife had to go away on a long business trip..
I figure in this room there are many immigrants..
And it's a common story among immigrants that one husband or wife comes to America first..
And sometimes it's years before they are reunited..
That's what it's like to long for someone's appearing..
That everything is focused on that day when you will see them again..
And Paul says, "If you are this way with Jesus, you cannot wait to see Him again.".
The same reward is there for us..
The crown of righteousness..
And why is that?.
Because in Jesus is all of our hope..
He's the one who holds our eternal future..
Through Jesus, all the blessings of God come to us..
But see it's not just what He gives us..
But it's because He's the one we love..
He is our treasure..
So what's waiting for you when the race is over and the fight is done?.
I hope that at the finish line, Jesus will say, "Well done, good and faithful servant.".
You see some of us are going to finish so fast and so strong, we're going to race across.
the line..
Others, we're going to think we're dying because we're so exhausted and we're going.
to cross the finish line..
I have a feeling some of us, we're going to cross the finish line bloody, beat up..
With our last breath, we'll cross that line..
But for all of those who long to see Jesus coming, He will welcome us..
He will reward us with the crown of righteousness..

$^{241}$He'll say to us, "Well done, good and faithful servant.".
I hope every person in this room will hear those words..
Will you be one of those who is shocked at the end?.
Jesus talks about some who will come to Him on the last day and they will say to Him,.
"Lord, Lord.".
And they're going to be shocked because Jesus will say to them, "Who are you?.
I don't know you..
Get away from me..
You're an evildoer.".
And they'll even plead with Him..
"But Lord, didn't we do all kinds of amazing things for you?.
Even in your name?".
But Jesus will say, "Get away from me.".
And what is it that they're missing?.
See Paul says, "I have done everything I need to do..
I've obeyed all the laws..
But His salvation was not sure until He knew Jesus..
Until He said, "I don't trust me..
I trust Jesus.".
I don't put my faith in me..
I put my faith in Jesus..
Where do you put your faith?.
You know, Paul reasoned God's solution to the sin problem that we all have..
And it's best summarized, I think, in Romans chapter 5..
And he says, "You see, at just the right time, when we were still powerless, Christ died.
for the ungodly..
Very rarely will anyone die for a righteous man, though for a good man someone might possibly.
dare to die..
But God demonstrates His love for us in this..
While we were sinners, still sinners, Christ died for us.".
While we were still sinners, Jesus died for us..
He took the punishment for our sin..
Do you believe this?.
Urban Voice, do you believe this?.
Pastor, you must have a church that doesn't like to speak out during the service..
It's very quiet..
Don't be ashamed..
Do you believe this?.
Yes..
I hope so..

$^{281}$This is our hope, that Jesus died for us..
My friends, if you have not trusted Jesus, could today be the day of your salvation?.
Is Jesus saying to you, "What are you waiting for?".
I'm going to ask everybody to close their eyes right now..
Let's keep our eyes closed..
If there is someone here who has never trusted Jesus for the forgiveness of their sins,.
for eternal life, would you just raise your hand?.
No one's looking..
It's just you and God..
And would you tell Him in your own words?.
And it's okay..
You can whisper it because He can hear you..
Would you tell Him, "Jesus, I need you..
I believe that you died for me because I am a sinner..
I'm nothing like Paul..
And if somebody like Paul had no hope in his own righteousness, how could I hope in.
my righteousness?.
I need your righteousness..
Forgive me for my sin..
Wash me and make me clean..
I believe Jesus can do this for me..
I trust Him..
I believe He can give me eternal life..
I believe Him..
And now I also believe He can teach me how to live for Him..
Because I want to live for Jesus..
The one who loved me from the beginning and will always love me..
The one who doesn't just say it, but the one who even laid down His life for me..
Jesus, I trust you..
Welcome me at the finish line..
And say of me, "Well done, good and faithful servant..
Enter into the joy of your master.".
Lord Jesus, I believe you will do this..
Thank you..
We pray in Jesus' name..
Amen..
So my brothers and sisters, for those of us who have already put our faith in Jesus Christ,.
let's not be foolish runners who don't know where the finish line is..
Who don't know that there is a finish line..
We are not running just in any direction..

$^{321}$We are running to our Savior..
So I ask you again, can you say confidently like Paul, "There's so much evidence that.
my life has been lived well..
There's so much proof that Jesus is in my life.".
If there isn't, perhaps this is a day when you might think, "What do I need to do?.
How do I need to change?.
What do I need to pray and ask God to help me with?.
Maybe it's time to ask, "What are the things that are keeping me from living that good.
life?".
And asking God for help..
Do you think he will help you?.
He does want to help you..
Do you think he wants you to live and experience the good life?.
Of course he does..
So when we ask God something he wants to do, do you think he'll say no?.
Of course not..
He wants to give it to you..
He's waiting for us to ask so he can give it to us..
Urban Voice, God loves you..
God loves this church..
He has given birth to this church..
And he will finish the work he started here..
And I hope we'll all run across the finish line together and hear those wonderful words..
Urban Voice, well done..
My good and faithful servant..
Join me in saying amen if you agree..
Amen..
Thank you, John..
Thank you for the challenge..
[APPLAUSE].
\newpage



\section{}
\label{sec:pK11m4fqws8}
\textbf{8/12/18 主日講道:  邊個係惡人 (Joyce Tam)}
\newline
\newline
連結: \href{https://youtube.com/watch?v=pK11m4fqws8}{\texttt{https://youtube.com/watch?v=pK11m4fqws8}} ~~~~ 語音日期: 2019-08-05
\newline
\newline
\hyperref[sec:tJHkj9NzTVE]{\small{< < < PREV SERMON < < <}}
~
\hyperref[sec:index_chronic]{\small{[返順時目]}}
~
\hyperref[sec:index_scriptual]{\small{[返順卷目]}}
~
\hyperref[sec:qkxKhCtF31k]{\small{> > > NEXT SERMON > > >}}
\newline
\newline
$^{1}$我都不知道我原來是新一代.
(笑聲).
Hello 你們好.
在預備今天的講章的時候.
我就想去你們UV的網頁看.
聽聽你們最近在看什麼.
在聽什麼.
誰知道我看到Ming 雞講到.
(笑聲).
她講到之前她說.
在座很多人認識她很多年.
知道她是什麼人.
她是個什麼人.
然後我想說.
我不是更大隻嗎.
就像剛才Mark Sir說的.
我哥 我嫂.
我媽媽的Best Friend 燕姨.
從小學已經認識我的朋友.
或者以前我在CCUC做中學生導師的團友都在當中.
我老公就跟我說.
如果我是你.
就會很緊張.
(笑聲).
或者你Host 我的時候.
可能你都在想.
其實我都在等你緊張.
既然大家都這麼緊張.
不如我們首先一起做個祈禱.
除了幫助我.
也幫助你去.
把焦點.
校正在我們今天要聽的訊息中.
我們一起低頭祈禱.
親愛的天父.
多謝你讓我們再一次能夠在.
今天聚集在一起去敬拜,讚美你.
主啊!來到你面前的時候.
我們一起去思想,研讀你的華語.
求你讓我們有聆聽的耳朵.

$^{41}$受教的心.
你的說話教訓,督責,校正我們.
叫我們得益,得以完全.
奉耶穌名求.
今天我說的訊息.
是在詩篇第一篇那裡.
可能你們很多人都很熟悉.
但我想和你們讀一次.
詩篇第一篇.
不從惡人的計謀.
不暫聚人的道路.
不助洩萬人的座位.
為喜愛耶和華律法.
就耶思想.
這人便為有福.
他要像一棵樹栽在溪水旁.
安時後結果子.
葉子也不枯乾.
凡他所做的盡都順利.
惡人並不是這樣.
乃像康丕被風吹散.
因此當審判的時候.
惡人必站立不住.
罪人在二人的會中也是如此.
因為耶和華知道二人的道路.
惡人的道路.
卻必滅亡.
根據聖經學者的分析.
詩篇第一篇.
就像一個宣言.
其實包含了詩篇150篇裡.
主要的訊息.
不從惡人的計謀.
這一節裡.
在原文裡.
之前還有一個.
在原文裡之前還有幾個字.
那幾個字是說.
有福的人.
意思是說一個有福的人.

$^{81}$不跟從一些不認識神.
不遵行耶和華律法.
不敬畏神的人.
也不容許一些遠離神的思想.
將我們從點A帶到點B.
更加不會留戀.
或者繼續在罪裡面.
誰是惡人?.
在聖經舊約薩姆爾記下.
記載了一個故事.
以色列王大衛.
有一天睡完午覺.
起床.
在平台走來走去的時候.
他就遠遠看到一個女人洗澡.
他心裡動了淫念.
而這個畫面就停留在他腦海裡.
然後他到處打聽.
究竟這個女人是誰呢?.
有人就告訴他.
她叫別士巴.
她是以蓮的女兒.
烏尼亞的老婆.
我想今時今日如果.
有一個男生要去追一個女生.
就去打聽一下她是誰.
我想問原來她有老公.
她說「算了」.
但大衛的反應不是這樣.
他還派人去找別士巴.
接了她回來.
在皇宮裡跟她發生了關係.
之後把她送回家.
但沒多久.
別士巴就派人去告訴大衛.
「大西了,我有了」.
很明顯他知道大衛就是那個經手人.
因為她的老公出差了.
他在戰場打仗.
根本他不可能帶頭.

$^{121}$邪惡的思想進入邪惡的行為.
帶來罪的後果.
大衛在這個時候還沒停下來.
他想想.
有什麼辦法可以掩蓋這件事.
不要讓別人知道.
他就想不如這樣.
我叫烏先生回來.
我就問候他.
打仗打得怎樣.
你最近好嗎.
你真是辛苦了.
不如這樣吧.
我放你大假.
你回家休息一下吧.
你也辛苦了很久.
送他走的時候還給他一些禮物.
讓他回去享受一下.
誰不知烏利亞是一個很正直的人.
他就在想.
我的兄弟在外面還在打仗.
我有什麼理由自己回家休息呢.
他就不肯.
被大衛知道之後當然氣死他.
他心想.
糟了他不肯回去怎麼辦呢.
他就再想另一條辦法.
再一次邀請烏利亞去皇宮.
跟他吃頓好的飯.
任他吃.
但在吃的當中.
他又給了很多酒去喝.
希望灌醉他.
讓他去嗅一嗅.
嗅一嗅就回家了.
但烏利亞就算吃飽了.
喝完酒也好.
結果他也不肯回家.
我想到了這個時候.
大衛真的越來越急.

$^{161}$因為想到想到.
別士巴的肚子就越來越大了.
那怎麼辦呢.
更邪惡的思想.
更邪惡的行動.
借刀殺人.
大衛在想.
你這麼喜歡打仗.
好,我就派你去打仗.
我還要把你放在前線.
然後大衛就吩咐他的軍隊.
我們去打仗.
見到敵人的時候.
你們一群人退下.
剩下烏利亞一個在前面.
讓他被敵人殺死.
結果烏利亞真的戰死沙場.
而大衛也名正言順地將別士巴接到皇宮.
成為他其中一個老婆.
很明顯大衛在這件事上是一個惡人.
是一個罪人.
是一個褻瀆人.
他有壞的思想.
也將這些壞的思想付諸行動.
越踩越深.
聖經記載大衛所做的事.
在耶和華眼中看為惡.
我相信我們在座很多人.
如果你來教會一段日子.
已經不是第一次聽這個故事.
你也不需要我告訴你.
大衛在這件事上犯了罪.
在Facebook上我看到一句說話.
.
意思是我們熟悉聖經.
甚至我們背誦聖經也是很重要.
但除非這些頭腦上的知識.
去到我們的心中.
我們才能夠真正活出.
生命根深義變化.

$^{201}$第二節.
.
今年八月一日至十日.
在香港有港九培靈研經會.
其中有一個講員說了一句說話.
我覺得和這一句是同一個意思.
意思是如果我們認識神.
又信服神的話.
神的同在會影響我們做事的方法.
我們走的每一步.
做的每一個決定.
你有沒有留意大衛和別士巴這件事上.
神好像不知道去了哪裡.
在傳道書裡面.
所羅門提醒我們.
生命是神給我們的恩典.
他鼓勵我們去享受神給我們的東西.
吃喝活樂玩樂所擁有的東西.
我們的家庭.
我們的親戚朋友.
但他也提醒我們.
我們所做的一切.
連一些沒有人知道的事.
有一天到審判的時候.
我們都要和神交仗.
因為敬畏神.
遵守神的誡命.
這是你和我都當盡的本份.
第三節.
他要像一棵樹栽在溪水旁.
按時候結果指.
葉子也不枯乾.
凡他所做的盡都順利.
我特別在今年經常在想.
我們經常去教會.
乖乖地去團契.
去Sunday school.
去崇拜.
但為什麼有時候.
不太看到很多生命的改變.

$^{241}$哪裡出了問題呢?.
我自己在想.
原來我們只是經常聽道.
甚至看聖經.
其實都沒有去到.
生命轉化的這個位置.
生命轉化是靠著神的恩典.
我們要信得過神的話.
祂會給我們最好的.
以神和我的同在.
就成為了我人生的價值觀.
我生命的方向.
我處事的方法.
例如在大衛這件事上.
聖經已經很清楚提醒我們.
我們要逃避少年人的恥辱.
我們每個人都知道這個提醒.
只差在這個temptation來臨的時候.
我有沒有去想起.
這句聖經的教導.
然後我真的要走出來.
我們都知道.
但我們知道之後.
我們有沒有去走出來.
第三句這裡提到.
「凡他所做的盡都順利」.
這裡不是說我們遵守神的話.
或者怎樣.
我們生命就順順利利.
沒有問題.
這裡提到的是一個有福的人.
一個認識神的人.
無論我們的境況好與壞.
神都和我們同在.
一個認識神.
和神同行的人.
真的可以面對風吹雨打的時候.
都不怕.
我們可以有一個超越常理.
平時人都不能夠經歷到.

$^{281}$那個抗難的能力.
因為你知道世界上沒有一樣東西.
沒有一個人可以勁得我們所信的神.
詩篇一篇第四至六節.
繼續提醒我們.
如果人的思想所作所為.
都是隨心所欲的時候.
審判的時候來到.
只有滅亡.
相反二人的道路.
就是我們否定我們自己.
逃離恣慾.
放下我的話事權.
放下我覺得做這個人一定要這樣.
It has to be this way.
放下.
Everything has to be under my control.
因為我們以為這樣好像是.
會令到我們感覺良好一點.
我實在很想挑戰弟兄姊妹.
特別是你已經信了一段時間的時候.
我們真的需要思想去挑戰我們自己.
用聖經的說話去挑戰我們.
去改變我們一些做事的方法.
怎樣去過我們的生活.
選擇神在你的生命裡面.
去更新你去改變你.
我們常常說信靠神.
但其實很多時候我們不是很信得過神.
但你再想深一層.
其實我們信主那一刻.
其實我們不是已經做了決定.
想這樣做嗎?.
但很多時候其實信了主之後.
我們就是沒有好好地去挑戰我們自己.
去練習聖經裡面的說話.
直至我們其實常常走回那惡人的道路.
在預備講章的時候.
我不斷在思想這個經文.
祈禱的時候.

$^{321}$突然間我腦海裡出現了一件兩年前發生的事.
我在神學院讀了四年半.
以前很熟悉我的人.
就知道我不是一個讀書人.
我也不是很能坐穩.
所以讀得很辛苦.
特別是去到第二年.
有一個科目是很辛苦.
讀極了都考得不好.
讀極了的文章都做得不好.
我們二十四個人開始讀那一課.
去到中段的時候我們只剩下十二個.
其中有三個還要是落選神學院.
我在想.
我怎樣可以做一些事情去釋放我的壓力.
令到我自己沒有那麼辛苦.
因為那段時間辛苦得來.
我睡不著覺.
我脾氣很差.
我老公叫了很多次.
我晚上還經常做噩夢.
我看到有些教授在跟我考試.
我就想不如這樣.
我買一支結他.
唱一首詩歌.
可能會幫助我自己放鬆一下.
但我又不敢跟老公說.
想買一支結他.
因為我已經不是工作.
交學費也很多錢.
但那支結他經常在我腦袋裡.
我經常上網看.
我煮飯也好.
做家務也好.
我當然會跟老公提.
那支結他多好.
外面賣多少錢.
現在賣這麼多錢.
說來說去他沒有什麼反應.
他很心急.

$^{361}$那不行了.
其實有點生氣.
我再想辦法.
怎樣可以買到一支結他.
想到另一條辦法.
我先叫我朋友買.
寄來我家.
然後我每個月會退一點錢給他.
但當我老公發現.
為什麼桌子旁邊有一支結他.
我就會告訴他.
這支結他是誰買的.
但我不會告訴他.
其實我每個月都在退錢給他.
但我能夠想像.
我講這些話的時候.
我心裡非常火熱地在罵他.
因為我心裡在想.
如果你聰明一點.
說OK賣給我.
就不用我想這些辦法了.
其實我真的發現我老公很無辜.
我又想到.
我七月底生日.
不如叫我哥和我弟弟.
分份送給我.
他們兩個是很疼我.
我知道如果我出聲.
他們一定會賣給我.
Thank you.
真的.
他們真的會賣給我.
雖然當時是三,四月.
還有幾個月才到我生日.
但我想告訴他們.
我知道你們會預備賣東西給我生日.
所以你們不要想了.
我想了.
不如這樣賣.
早幾個月賣.

$^{401}$我就將這些辦法和我的朋友說.
我的朋友又沒有發聲.
不知道你們有沒有發覺.
我們平時沒有什麼創作力.
但在這些時候.
創作總監上身.
創作力非常澎湃.
我的朋友又沒有發聲.
最後他就說.
不如你正正經經和你老公說吧.
(笑).
惡人的計謀.
罪人的道路.
涉患人的座位.
當然現在說起來很好笑.
你不知道我當時有多麼的掙扎.
想了這些詭計.
連事奉也出了.
因為我跟老公說.
可能我買一支結他.
彈一首歌.
我會帶事奉.
神很恩待我.
我身邊有一個這樣的朋友.
好像拿丹一樣.
在殺母義記下記載.
大衛借刀殺人之後.
神差派拿丹去承認大衛.
因為大衛藐視耶和華的命令.
做神眼中看為惡的事.
可能你在想.
Joyce你做的事很小事.
你又不是借刀殺人.
又不是犯奸淫.
你只是買一支結他.
我在思想的時候.
我在想.
我想不是大事.
還是小事的問題.
而是在這件事上.

$^{441}$我的所思所想.
都是很隨心所欲.
跟一個不認識神的人.
不一樣.
一模一樣.
更羞愧的時候.
發生這件事的時候.
我是在神學院的.
給我自己也有一個提醒.
其實我真的需要時時警醒.
不是說我是讀神學院.
甚至今時今日我已經畢業.
我就不會再犯罪得罪神.
忘記神的律法.
在神的眼中.
是沒有分別的.
在這件事上.
我都是一個惡人.
說得好像很慘.
但我們也不要忘記.
大衛雖然做神眼中看為惡的事.
但當他願意悔改.
神仍然寬恕他.
仍然重用他.
同樣神也寬恕我.
當我再去分析.
我讀神學院.
分析我自己的心想什麼.
有什麼想法的時候.
我亦看到一件事.
很有趣.
是我以前沒有留意到.
整件事裡面.
神是沒有在.
沒有牽涉.
不像我平時.
我平時很喜歡出去走路.
走路的時候我就跟神聊天.
但在這件事上.
我是沒有跟神聊天.

$^{481}$不像我平時的習慣.
我在想為什麼會這樣.
其實我自己想.
可能是我心虛.
我知道.
我在想的這件事.
在做的這件事.
是神不喜悅.
我自己知道.
只不過我自己不去面對.
求神幫助我們.
每一個認識祂的人.
我們不是只是片面.
去聽一篇道.
去看聖經.
去知道一些經文.
就算.
就好像是足夠.
但其實一直都停留在這個位置.
聖靈其實經常有跟我們說話.
但如果我們習慣了.
將它放在一邊.
我繼續做.
繼續用我的方法去處理很多事的時候.
其實你永遠都不會經歷到.
神同在的恩典.
我真的已經一次又一次.
很多次去經歷了.
有些事我覺得這樣做似乎是好一點.
我自己舒服一點.
安全一點.
但神有聲音跟我說.
不是這樣做.
而在聖經裡面亦都很清楚.
給了我一些指引.
似乎不是應該這樣做.
當我願意信靠神去經歷的時候.
接著出來的結果.
第一是我真的猜都猜不到.
第二.

$^{521}$我真的猜不到.
絕對是比我自己所想的好.
然後我覺得藉著那件事.
神建立了我對祂的信心.
其實我真的覺得我經常很沒有信心.
我們人其實是沒有信心.
那個信心是神給我們的.
但藉著這些困難.
藉著這些事.
神就是去建立我們的信心.
以致下一次再面對不同的挑戰.
困難的時候.
我的信心比上次大一點.
我知道原來我信神.
是比信我自己好.
甚至是人給我的意見.
好像一棵樹栽在溪水旁.
將神的話.
栽在你的心裡.
用聖靈跟你說話.
甚至如果神真的在你身邊.
擺一個兩個拿單的時候.
你要好好地去學習去聆聽.
讓神的話時刻刻影響我們的心思意念.
我們所走的每一步.
所做的每一個決定.
不做一個惡人.
有一天.
能夠在主的面前.
站立得穩穩.
我們一起祈禱.
親愛的主耶穌.
求你赦免我們很多時候.
都看不到其實我們自己就是那個惡人.
我們偏行己路.
主啊求你幫助我們.
不要去看輕你的律法.
你的教導.
提醒我們真的要晝夜思想.
你的話成為我們腳前的燈.

$^{561}$路上的光.
幫助我們做一個蒙福的義人.
願眾讚榮耀都歸還給你.
阿們.
謝謝.
\newpage



\section{}
\label{sec:qkxKhCtF31k}
\textbf{5/12/19 主日講道\_耶穌天國的小故事-2\_恩典 (陳崇基牧師)}
\newline
\newline
連結: \href{https://youtube.com/watch?v=qkxKhCtF31k}{\texttt{https://youtube.com/watch?v=qkxKhCtF31k}} ~~~~ 語音日期: 2019-08-06
\newline
\newline
\hyperref[sec:pK11m4fqws8]{\small{< < < PREV SERMON < < <}}
~
\hyperref[sec:index_chronic]{\small{[返順時目]}}
~
\hyperref[sec:index_scriptual]{\small{[返順卷目]}}
~
\hyperref[sec:KWfIxDVxl3w]{\small{> > > NEXT SERMON > > >}}
\newline
\newline
$^{1}$各位早安!.
今天我們繼續講到的系列.
叫做耶穌的小故事.
一些小故事.
即是聖經裡面.
特別是三本科目書.
馬太,馬可,路加.
頭三本的科目書是講到耶穌的一些比喻的故事.
上一次我們開始的時候.
講到這些比喻的故事.
不單止是來娛樂一下我們.
或者是一些很容易記的故事.
也都是講到耶穌的比喻的時候.
這些比喻很多時候.
最重要的事情就是去挑釁我們.
挑釁我們一些平時有的思想.
然後把我們搖一搖醒.
以致我們能夠掌握到.
在天國裡面的運作原則.
在天國裡面做一個天主教的運作.
問他們為什麼整天站在這裡無所事事呢?.
他們說因為沒有人顧我們.
他對他們說.
你們也進葡萄園去吧.
傍晚.
葡萄園的主人對公頭說.
把工人叫來.
給他們工錢.
從最後來的人開始.
直到最先來的人.
那些在十一時.
坐五時的公頭.
左右才僱用的工人來了.
每個人領取一個特拿利.
最先僱用的人也來了.
他們以為自己會多領一些.
但他們每個人.
都是領取一個特拿利.
他們領取之後.
對園主.

$^{41}$咕咕隆隆.
法院言說.
那些是最後來的.
只做了一個小時.
而我們整天勞苦受熱.
你卻使他們跟我們.
有同等的待遇?.
園主回答他們其中一個人.
朋友,我們有虧待你.
你不是跟我協定好.
一個特拿利嗎?.
那屬於你的那份酒吧.
我願意給那.
最後來的人.
跟給你們一樣.
難道我按自己的意思.
用自己的東西.
也是不合意嗎?.
還是因為.
我有善心.
你就眼紅了?.
第十六節.
這樣.
在後的將要在前.
在前的將要在後.
那考考你吧.
葡萄園的主人.
出了市集請工人.
請了多少次?.
請了多少批?.
清早一批.
清早剛開工的時候.
六點鐘.
接著呢?.
九點鐘,十二點鐘.
三點鐘,四批.
然後.
五點鐘又一批.
後來下班一個小時.
又再請多批.

$^{81}$總共請了五批.
當時的人沒有手錶,也沒有手提電話.
他怎麼看時間的?.
看太陽.
但是有些大戶的人家.
或者有錢一點的人.
好像園主.
他很多時候用那些叫日規.
有沒有見過什麼叫日規?.
Sun dial.
有沒有?.
你知道這個叫Sun dial.
有沒有用過?.
沒有用過?.
你那些大戶人家.
那些South Loop.
你也有用過在你的Paleo.
我想問你.
那裡有一支針.
有一條針.
然後有一些marking.
有一些標記在樹上.
如果你要安裝這個東西.
哪邊是東,哪邊是西?.
就是太陽射到.
雪的針那裡.
就會有一個影在那些標記上面.
因為太陽.
你知道太陽會動的.
從我們的角度看.
太陽是從東邊出來的.
從我們的角度看.
西邊下山.
知道嗎?.
知道吧?.
自然科學.
你要怎麼安裝這個東西?.
哪邊是東,哪邊是西?.
大哥.
東邊是右邊還是左邊?.

$^{121}$右邊,對了.
右邊,東邊是右邊.
左邊呢?.
就是西邊.
太陽會動的時候.
現在大概幾點鐘?.
我剛才沒有跟你揭曉.
拍這張照的時候.
大概是下午四點半.
看到嗎?.
四點半,五點鐘左右.
行不行?.
這樣動的.
東邊在右邊.
那支針就自然.
放在第一.
有十二個標記.
然後到中午的時候就在上面.
再到左邊的時候.
就是日落的時候.
西邊.
那支針就放在.
大概十二個格.
第一時,第二時,第三時就是這個意思.
什麼時候是午靜,什麼時候是新初.
這些是舊的term.
所以以前的人就是這樣.
tell time.
那個葡萄園的園主.
就是在這五個時間.
出去請人.
太陽出來的時候.
清早,就是第一時.
第一時.
當然他們當時看時間.
就不是說我們今天這麼準.
現在是十點四十一分.
他們是看時辰的.
好像我們華人傳統.
誰時辰.

$^{161}$你記得耶穌經常說.
事後到了.
或者事後還未到.
The hour has come.
其實就是在說這個hour.
不是exactly在那個point.
而是在那段時間.
他說事後到了就是這個意思.
那個hour已經來到.
這個園主清早出去一次.
然後三時的時候.
又出去一次.
第六個的時辰.
中午午靜的時候.
又出去一次.
九點的時候也是出去一次.
最後的時候.
臨放工.
日落西山.
之前的一個小時.
來請一些人.
我們要了解當時的背景.
我們經常說聖經一定要了解.
當時的背景,怎樣聽那個故事.
人家當時怎樣聽那個故事.
然後我們知道今天我們應該怎樣聽那個故事.
當時的工人.
特別是在葡萄園.
或者農田上工作的工人.
那個情況.
是這樣的.
他們會一早.
他們的工人.
特別是做苦力的.
他們會到輪流聚集.
那時候有人請他們的時候.
他們做一天的工.
第二天又是輪流這樣去做.
沒有什麼permanent.
有4-1-K.

$^{201}$那些這樣的東西,沒有的.
是一天起來.
去到市集那裡去等.
特別是在以色列巴勒斯坦地.
大約九月中.
或者九月底的時候.
有十幾天,兩個星期.
就是葡萄園.
要去到被採摘.
或者收割的時候.
那些葡萄.
成熟了之後.
就一定要在那兩個星期裡面.
盡量去採摘.
採摘了之後就去釀酒.
或者去做其他東西等等.
所以那段時間是非常忙的.
所以市集裡面就有很多的人.
因為你不趕時間去採摘的時候.
九月底過了之後.
到十月就會第一陣雨.
就會到.
到第二陣雨到的時候.
這些葡萄如果是掉了.
或者爛了,就不再值錢.
就浪費了.
你要釀酒,你要賺錢,就在這段時間一定要出去.
所以主人.
就會去市集那裡.
找很多很多的人.
來幫他,快點去摘這些葡萄.
以致能夠拿到.
這些收成.
等於我們今天,你的fiscal year.
你要結帳.
你要做一些quality report.
那段時間是很辛苦的.
我們中間.
有些弟兄姊妹是幫人報稅的.
前排就很辛苦.

$^{241}$因為每個人都last minute才做.
沒有人一早做的.
所以那兩個禮拜.
或者三個禮拜.
就會很辛苦.
當時就是這樣.
當時那些工人就去市集那裡聚集.
走走.
然後主人就會選擇.
那些.
呼籲那些人.
去到工場裡面做.
今天其實也是這樣,有些地方.
譬如home depot.
那些停車場.
是不是很多人在等.
等著開工.
我記得.
前兩三年的時候.
我們這裡開始教會.
我經常去home depot.
我去買東西.
我就一個人搬那些4x8.
就是那些很大塊的.
搬上車.
自然就有人來幫我.
那些Latino那些人.
因為他們都去搵食.
當然我不請他們.
我一個人就搞定.
有很多情況.
都好像今天這樣.
或者果欄.
一些搬運.
操縱工作.
一天就做一天的臨時工.
然後.
希望有人請你.
你就在那裡等.
等十個鐘十二個鐘.

$^{281}$然後回到家裡.
你整家就要叫餓了.
如果幸運地有人請你.
那一晚就有東西吃.
這個就是我們現代人最難明白的地方.
就是這個比喻.
這個比喻是按當時的情況.
而去.
講出來的.
沒有了當時的背景我們就不明白.
我以為這個比喻是教我們.
對勞工不公平.
其實不是那個比喻的意思.
勞工是需要公平的.
你知道上個星期有很多人罷工.
Uber.
很多全世界主要的城市.
這些Uber的司機.
Uber上司.
那些CEO.
每年有百多個million.
每年百多個million.
而那些司機都吃都沒得吃.
所以他們罷工.
你看這個比喻你不要以為.
這個比喻是贊成或者容許.
這些事情發生.
Economic justice都是聖經裡面所講的.
耶穌是不贊成.
這些不公平的對待.
勞工勞資的問題.
聖經是concern的.
神是concern的.
但是講回這個比喻的時候.
你就是要明白當時的背景.
你才知道他想講什麼.
不是在講我們可以容許.
勞資不公平的情況.
是不是.
但是耶穌用當時的實況.

$^{321}$來道出一種天國的道理.
什麼道理呢.
第一個道理.
就是天國裡面.
恩典運作的原則.
一個人進天國.
到底恩典是怎麼運作的呢.
就是按照家主.
來呼召人進天國.
不是按你有多少好處.
或者有多少才幹.
因為家主.
原主去到市集的時候.
有很多人.
他就選了一些人.
進到天國.
這個故事的實景.
就是這樣.
原主冒險的工人.
或者市集裡面等的人.
任何人.
進到葡萄園.
是天國的恩典.
他能夠進到這個工場.
是天國的恩典.
你留意20章第一節.
他說因為.
換句話說.
有上文.
比喻講的故事.
其實就是間接的.
去回答上面的一些問題.
上面發生什麼事呢.
上面有一個故事.
就是一個很有錢的年輕的官.
他來問耶穌.
他說老師.
你要做什麼.
才可以得到永生.
耶穌就跟他說.

$^{361}$一大堆東西.
你要變賣你一切的東西.
不要再貪了.
不要再做金錢的奴隸.
然後你跟隨我.
你就可以進天國.
就是一個年輕的官.
像今天的財政司長.
很有錢的.
聽見耶穌這樣說的時候.
他就很憂愁.
有些門徒在身邊.
聽見一個這樣的對答的時候.
彼得就問一個問題.
彼得就說漢.
我們已經撇下所有.
來跟隨你.
將來我們可以得到什麼.
這個問題.
耶穌就這樣回答.
那個實在告訴你.
你們這一班跟隨我的人.
到神國復興的時候.
就是末日的時候.
人子在榮耀的寶座上.
你們也要在十二個寶座上.
審判以色列十二個支派.
凡為我明.
撇下房屋.
弟兄姊妹父母.
兒女田地的.
必會得著百倍.
並且會承受永生.
請你不要按字面解釋.
這裡不是叫我們拋妻棄子.
這裡是說神知道你.
跟隨耶穌的時候.
你的付出神是不會忘記的.
神是會紀念的.
並且會得百倍.

$^{401}$百倍也不是我們今天.
百倍是美好的生命.
就是你跟隨耶穌的時候.
已經有豐盛的生命.
並且會得一個更加豐盛的生命.
耶穌這樣回答他.
他接著就說了這個比喻.
這個比喻.
我們不可以說.
不如我等到最後一小時.
我才跟隨耶穌.
那我就不用做了.
最後的時候.
大半個小時.
我就不想吃虧了.
不想吃虧.
那我們就搞錯了.
因為我們不是按照當時.
巴勒斯坦地的工場的情況.
如果按照當時的.
實際的情況.
你就知道.
被選進去葡萄園.
是完全恩典.
是All of Grace.
是所有.
因為那個家主.
去選你.
被選的人.
就會歡天喜地.
然後做了幸運兒.
他就會進去這個工場.
那個恩典的運作就是這樣.
就是你能夠得到永生.
你能夠進到天國.
原因是什麼呢?.
請你不要痴心妄想.
你以為.
你自己有丁點的本能或資格.
沒有一個人.

$^{441}$具備有任何的必然.
是神要選我們的.
你被選.
不進到上帝的葡萄園.
天國.
完全是恩典.
不單止你能夠進到天國.
是一個恩典.
你能夠進到天國.
事奉也是一個恩典.
整個比喻裡面.
我覺得比較有挑戰性的說話.
很可能就是最後那句.
大家看看.
他說.
然而有許多在前的.
將要在後.
在後的.
將要在前.
你記得.
比喻裡面.
為什麼說這句話.
因為那些先被請的.
一早上六點鐘.
被請進去園裡工作的人.
他們就發怨言.
然後他們就囉囉嗦嗦.
或者嘰嘰噠噠.
為什麼你這麼不公平.
我們做到汗流浹背.
做了這麼多個小時.
我們得的仍然是那一樣.
而那些剛剛做了一個小時的.
他們居然.
可以跟我們得的.
工價是一模一樣.
園主回答他們其中一個人.
他說朋友我沒有虧待過你.
因為你進完紙的時候.
我們已經協定了.

$^{481}$你得這麼多.
拿你的走.
我願意.
最後的.
得到的跟你一樣.
是按我自己的心意.
按我自己的恩典.
去做的.
那不合宜嗎.
你是不是因為這樣而眼紅呢.
你留意這個經文.
裡面的脈絡.
之前解答了一個.
年輕的財主.
發出的問題.
我可以得到什麼.
我可以承受永生.
彼得問了一個問題.
我撇下所有.
我到底有什麼好處.
耶穌就說請你小心.
你跟隨我的話.
你必定會得百倍.
我一定會紀念你.
然而有許多在前的.
將要在後.
我講個故事給你聽.
因為天國就好像葡萄園.
的家主.
出去不同的時間.
請不同的工人.
這個故事就說了出來.
故事的結尾.
這樣在後的將要在前.
在前的將要在後.
你留意當時的實況.
第一批進去工作.
去事奉的那一班.
就代表第一批.
有上帝恩典揀選的人.

$^{521}$最後的一批.
亦都是.
按照上帝的恩典.
亦都被揀選進去天國.
並且進去天國裡面事奉.
你知道最後的一批.
可能就是最感激的一批.
因為他們以為.
整天就會捱餓.
但最後的一小時.
那個園主.
去到工場試習.
就去選照他們.
你說是不是.
最後的一小時被請的那一班人.
是最感激的那一班人.
按照當時的實況.
他們就是那些知道.
什麼是恩典.
我們不配被揀選.
但我們最後的一小時居然被揀選.
我們就努力的.
甘心的.
或者很感激的.
來到工作.
第一批的.
就是那些身在恩典裡面.
然後不知道恩典的.
你記得耶穌曾經跟法利賽人說過一句話.
他說你知道嗎.
上帝可以從石頭那裡.
興起阿巴拉罕的後裔.
你不要以為.
你是猶太人.
你是那些所謂被揀選的子民.
我們揀選了這麼久.
這麼多的時間.
我們一定是坐定笠六.
耶穌說.
上帝可以從石頭裡面.

$^{561}$爆出來.
一些天國裡面的子民.
其實就是在說他們.
猶太人那種.
被揀選.
能夠去事奉.
那種驕傲.
當頭棒喝.
以為自己了不起.
竟然敢跟上帝.
的恩典.
來討價還價.
或者去斤斤計較.
你進入天國.
你能夠事奉.
全部都是上帝自由的揀選.
上帝可以沒有你.
上帝可以不揀你.
但如果今天揀了你.
你應該用什麼心態.
來事奉.
所以為什麼說在後的會在前.
在後被揀選的會知道上帝的恩典.
我真的努力去工作.
在前的.
信了很久的.
被揀選了很久的.
那些就在發牢騷.
上帝可以從石頭裡面去興起人.
你如果不懂恩典的時候.
身在恩中不知恩.
你可能就會在前的反而在後.
你可能會失去恩典.
失去你在葡萄園裡面的地位.
你知道有人.
是天生五音不全.
你知道嗎.
有些人真的五音不全.
沒有一句.
說得不客氣的話.

$^{601}$你知道有些人不生出來音樂是比較厲害的.
對不對.
這個世界是這樣的.
但你知道.
上帝有時候.
他可能會聽.
喜歡聽一些五音不全的人.
去讚美他.
如果你不知道你有的才幹.
是上帝給你的話.
上帝可以.
我不聽.
今天母親節.
媽媽最喜歡聽誰唱歌.
是不是American.
Gertrude Talon.
唱得很好聽.
最喜歡聽什麼.
是小朋友的寶寶.
所唱的歌.
唱得多不好聽.
都說你唱得很好聽.
上帝很多時候.
對我們也是這樣.
沒有一個人.
有權去選擇我們怎樣出生.
全部都是上帝的恩典.
沒有一個人能夠先選擇.
我出生的時候.
我要有什麼套餐.
我要有多高.
我要有多多.
技能.
沒有一個人有權選擇.
是上帝的恩典.
在前的反而在後.
後者是心知本來沒有機會.
但現在可以得到上帝的恩典.
在最重要的關頭.
我能夠進到葡萄園裡去事奉.

$^{641}$最後一個鐘.
你留意整個的比喻.
有些很有趣的地方.
你真的留意.
有些故事.
他故意不說一些東西.
不是因為那件事不重要.
他故意不說一些東西.
是因為不說那件事才是最重要.
你明白我的意思嗎.
想爆你.
他錯過了不說的東西.
那些叫做gap.
不是blank.
這個故事裡面.
第五批被請的工人.
你有沒有留意沒有說什麼.
就是家主沒有跟那班工人.
說那個工價.
頭四批都有說的.
頭四批都有說.
只是第五批最後一個中國批.
你們為什麼在這裡無所事事.
進來一起工作.
他們就湧進去田園裡面.
來工作.
沒有說過多少錢.
因為最後那批被請的.
沒有說明工價也好.
他們知道他們是最後一班.
是完全不會理會.
有多少工錢.
只要今天有機會.
他就帶著感激的心.
進去工場裡面事奉.
進去天國裡面事奉.
就是因為知恩.
反而得到最好的.
再後的反而再賤.
再賤的是跟著計較.

$^{681}$你想想.
今天你在天國裡面.
你檢查一下你對自己說的話.
最近這幾個星期.
你有沒有跟著計較.
有沒有在上帝面前覺得.
好像上帝欠了你一件事.
上帝沒有欠過你一件事.
因為你能進天國.
是完全上帝恩典.
你能夠事奉.
也是宣佈上帝的恩典.
你應該知恩圖報.
你想想你身邊的人.
有很多可能被你所謂早信.
他們今天失去了救恩.
不見了.
我認識一個弟兄.
最近他.
突然間有一天.
發現心律不靜.
很害怕.
跟我差不多靈機.
突然間才恍然大悟.
我最基本的生命氣息.
原來都是靠上帝的恩典.
你跟我都是一樣.
進到天國.
能夠在天國裡面事奉.
不要記憶的.
是上帝的恩典.
給我們一個警醒.
當頭棒學.
讓我們全心感激的心.
來事奉.
我們同心協力祈禱.
上主我們多謝你.
因為你賜給我們生命的氣息.
我們今天能夠呼吸能夠存活.
是因為你對我們的恩典.

$^{721}$多謝你這個小故事.
來提醒我們.
能夠向你祈禱.
能夠在天國裡面.
能夠成為你的子民.
能夠事奉你.
主,這個是我們當然要去感恩的.
因為我們本身就不配受.
完全的簡說.
這個完全的恩典.
我們不配受,卻淋到我們人生.
主讓我們帶著感激.
帶著感恩.
來去服侍你.
叫我們.
能夠永遠都在前.
幫助我們用一個感恩的心.
來去回應你.
對我們的恩典.
我們這樣祈禱,奉主名禱.
\newpage



\section{}
\label{sec:KWfIxDVxl3w}
\textbf{12/30/18 主日講道: 齊心作主工 (吳君才牧師)}
\newline
\newline
連結: \href{https://youtube.com/watch?v=KWfIxDVxl3w}{\texttt{https://youtube.com/watch?v=KWfIxDVxl3w}} ~~~~ 語音日期: 2019-08-07
\newline
\newline
\hyperref[sec:qkxKhCtF31k]{\small{< < < PREV SERMON < < <}}
~
\hyperref[sec:index_chronic]{\small{[返順時目]}}
~
\hyperref[sec:index_scriptual]{\small{[返順卷目]}}
~
\hyperref[sec:J6o6vqGpqR4]{\small{> > > NEXT SERMON > > >}}
\newline
\newline
$^{1}$.
當然對這些仇苦擔心和一些令你很懼怕的事.
你會怎樣去面對呢?.
很多時候我們會很擔心,很受煩,無所適從.
但是我們有沒有想到去交託,依靠神呢?.
令自己仍然有一個剛強,喜樂,走前面的日子呢?.
有時候我們不單要為自己,我們也要為神的國度,神的家.
你有沒有想過要有什麼準備可以做呢?.
已經選出新的核心領袖.
他們已經在10月8日有退休.
每個人都預備在主的面前,很好的大展拳腳,很好的事奉.
但是所有其他在座的弟兄姊妹,或者今天不在的弟兄姊妹.
我們都不要放慢,或者放軟手腳.
你和我都要繼續同心合意的齊心作主功.
所以今天我的講道崇拜經題,我們找尼希米記第三和第四章.
當中我們有學習和被提醒,我們先有個簡單祈禱.
求主接著尼希米記第三和第四章.
讓我們每個人在你面前有領受和學習.
願你的靈在我們中間,將我們憂慮,重擔,受煩,各樣的擺在一邊.
專注在你身上,在你的話語身上.
求你向我們說話,個別施恩,奉主耶穌基督的名求,阿們.
昔日尼希米重建城牆,策動當時的百姓成為修造的人.
今天我們用尼希米記第三和第四章.
怎樣可以成為新的時代修造的人,作為齊心作主功的人.
有些表現和一些東西我們要掌握和運用.
我提出三大方面的東西.
第一,我們要各盡所能,分工合作.
可不可以打出來.
那裡的經文有三章四章,諸如此類,我們會一直找下去.
再到標題.
我們會看到很多的情況.
當代一位神學家,Escenius,伊健拿斯.
分析以色列記和尼希米記,當時的時代意義.
據他的觀察,他提出三大立論,我抽其中兩項.
第一,他認為以色列記和尼希米記當中用了四分之一的篇幅.
記錄了九個大小的名單.
成書的人著眼,顯然是群眾而不是個別的英雄人物.
所以明顯地表示,父神需要的不單是卓越的個別英雄領袖.
整個神子民的群體也有一起參與.
換言之,每一個成員在神的角度也有份.

$^{41}$可以成為一個修造的人,齊心作主公的人.
修建成像是這樣的,修造教會,建立教會的事實上也如是.
所以第二項,尼希米在修完成像後.
特別是尼希米記七章一節和二後後.
他重覆了第一批回歸耶路撒冷的人的名單.
回顧了當時的古列,即是主前五三百年.
到達薩西王,即是主前四百三十三年.
超過一百年的重建歷史.
由當時起初所羅巴伯建聖殿起.
到以斯拉重建聖文,尼希米重建聖城成像.
整個過程計算起來有四五代.
換句話說,這個重建的大業,修造的事工.
絕對不是個別領袖,窮一生之力來完成.
而是需要個別會眾,和一代與一代,再另一代.
投身參與,各盡其力,合作又分工才能成就.
成為社區教會,我們要發展未來的事工.
我們需要有個觀念,我們要超越一代.
伸延到下一代,甚至以後的一代.
所以面對未來,我們仍然需要全體弟兄姊妹身體力行,投身.
其中,當我們看第三四章的時候.
我們會見到第一個情況,請下.
就是領袖的以身作則.
第一,回到第一章.
這裡你會見到經文,第一節我們讀一讀.
(文字).
你會見到記載大祭司,在陽門.
大祭司以利亞實,原來與反對造成的多比亞和參巴拉.
都有親色的關係.
(文字).
第十七到二十六節記載.
當時還有利未人和祭司的名單.
所以你可以見到任何事工的核心領袖的認同.
以及很熱切的投入為身.
要帶領一個工作是極重要的事.
我們看第二章,下一章.
(文字).
如果看聖經的記載.
(文字).
第二到七節,時間關係,你看一看,我不讀出.
你計算一下,第二到七節有什麼人呢?.

$^{81}$(文字).
還有,欽利的兒子薩克.
亞西娜的子孫.
阿哈斯的孫.
烏利亞的兒子米利末.
他們都是祭司之一.
巴拿的兒子薩都.
提哥亞人是阿摩斯故鄉的人.
巴西亞的兩個兒子等等.
他們則在北邊城牆建築.
而在西牆的修造.
就有基變人,米斯巴人.
他們位於耶路撒冷的西北或北邊.
另外南邊牆修造的人.
在第十三十四節有記載.
你開後一點,再下一章.
然後當時.
修造東牆的就在第十五到三十二節.
又有記載.
當時波斯人統治猶大地區.
似乎是以區作為他們的行政單位.
每區又再劃一些分區.
這一章提到行政區包括有耶路撒冷.
在伯克基林.
米斯巴.
白宿.
基拉等等.
這些地區的人都總動員起來修造.
積極投身.
這給我們今天成會的弟兄姊妹.
東南西北各邊.
神給你分佈在不同的地區.
你可以怎樣有回應.
第三,看下去.
專業人士也毀身參與.
一起讀第八節.
其次是銀匠阿開雅的兒子烏舍修造.
其次是作鄉的阿哈尼雅修造.
這些人修建耶路撒冷直到寬牆.
看到這裡.

$^{121}$這裡出現了銀匠和作鄉.
銀匠是製造珠寶打造金銀飾物的人.
作鄉是製造香膏.
供宗教儀式清潔身體用途.
銀匠和作鄉都是專門技術人士.
你看到這裡第31節.
提到一位銀匠.
叫馬基雅.
第32節提到商人.
他們在城的閣樓和陽門中間修造.
弟兄姊妹,你有沒有留意到.
原來你和我的背景.
或者我們所從事的專業工作或職業.
經過不少訓練神的造就.
原來可以像昔日的銀匠,作鄉,商人等等.
有什麼聖器可以獻上為神所用呢?.
神都會預納我們.
神都希望我們好好為教會所用.
作成修造的人.
齊心作工,竭盡所能.
這是實際的貢獻.
對分工合作和教會的建立.
未來會有積極的影響意義.
接著我們看第四項.
是婦女姐妹更佔一重要的職位.
一起讀第十二節.
其次是管理耶路撒冷的那一半.
哈羅哈的兒子沙龍和他的女兒修造.
你看到這裡.
哈羅哈的兒子沙龍和他的女兒們修造.
特別提及女兒們.
或者沙龍沒有兒子.
所以他找女兒們來承擔修造的事工.
但其實整個建築修牆的事業裡.
婦女都有份參與.
投身在內.
在座的弟兄姊妹.
你要很欣賞一下姐妹.
這些婦女很重要的.
千萬不要被忽略.

$^{161}$她們很有價值.
修造城牆的事.
教會的建立也需要弟兄姊妹.
甚至社區的服務關懷.
也不能沒有姐妹,婦女.
從剛才的分析.
我們可以看到.
當姐妹,婦女竭盡所能的時候.
分工合作的時候.
積極參與投入的時候.
是很重要的.
同樣也是今天教會的呼聲.
事奉神的呼聲.
十九世紀.
美國著名詩人.
愛麥申.
曾經這樣說過.
自私自利的人.
不願意參與.
不願意施濟.
他所損失的.
是比受者所得的更多.
你自私.
你不肯付出.
你不肯投身.
原來最終的失去.
是更多的福氣和恩典.
齊心作主公.
第二個指引.
下一章.
就是要常常禱告.
是有信心的行動.
你收造城牆這件事是一件美事.
但也是一件屬靈的事.
撒旦是不會放過.
所以就會有屬靈的爭戰.
有愛患.
有來憂.
不單止以前是這樣.
今時今日也有這樣的情況.

$^{201}$當我們想完成神的事宮的時候.
我們會有屬靈的爭戰.
有困難的打擊和種種的人事問題.
這些是無可避免的.
尼熙米收造城牆的工作.
面對內憂和愛患.
但他能夠用禱告仰望神的心去得成.
他能夠持守禱告的能力.
當時尼熙米的愛患是來自.
北面的和倫人.
的參巴拉和薩瑪利亞人.
東南面的阿滿人.
多比亞.
南面的阿拉伯人.
和西面的阿薩德人.
所以你看到敵人的第一招.
是用激將法.
是動口不動手.
甚至是扭曲事實的伎倆.
第四章一到三節有講.
這裡列了出來.
你望一望.
一到三節不讀出.
但我用現代中文譯本的譯法讀給你聽.
神說:這些虛弱的猶太人在做什麼呢?.
他們想要重新建造城牆嗎?.
他們以為藉著憲制就能在一天之內完工嗎?.
他們意想從燒過的碎瓦堆中.
找石塊來建造呢?.
第三節.
阿滿人多比亞站在旁邊附和說.
如果能做出這樣的牆.
恐怕一隻狐狸都能踩到它呢?.
他們嘲弄猶太人.
以為靠虔誠就可以有所作為.
這裡火燒的石頭.
其實是經過大火燒的石頭.
是指特別在尼布加利沙旺攻城的時候.
石牆焚燒了.
堡塔遺留下來.

$^{241}$一大堆一大堆的殘破石頭.
其實有不少的石頭可以重新用.
這與敵人所稱的歪曲並非事實.
猶太人更不是想藉著憲制一天之內完工.
他們用激將法.
甚至揶揄,嘲笑.
令人心裡搖動.
甚至會有負面的效應.
令人很軟弱,很自卑.
當這些行不通的時候.
這群人就聯合起來再搞其他的陰謀.
擾亂.
第七,第八節.
下一章,麻煩你.
我們一起讀一讀七八這兩節.
(朗讀).
你看到他們是這樣做的.
但是他們不敢那麼猖狂.
因為這個修建城牆的事.
是得著皇帝的詔書所批准.
所以他只能夠盡量周邊搞擾和威脅.
第二章,下一章,麻煩你.
他們面對的方法是怎樣呢.
就是繼續嘗嘗禱告.
面對這些難阻,離希米的招式.
值得我們學校.
繼續嘗嘗禱告.
並盡自己的本分.
第四,第六節.
那裡有說.
你看一看四到六節.
我用現代中文譯本讀給大家聽.
(朗讀).
求你不要赦免他們所做的壞事.
不要忘了他們的罪惡.
因為他們侮辱了我們這些正在修建的人.
於是我們繼續重建城牆.
不久就建到一半高度.
因為人民都熱心工作.
你看到這裡.

$^{281}$離希米每當遇到困難的時候.
面對挑戰的時候.
往往立刻禱告.
這個現象經常在離希米記可以找到.
因為離希米深深知道.
禱告可以將危機轉化成為傳記.
一個修造城牆的人.
我們想做主公的人.
是要追求上上的禱告.
在台灣陽明大學醫學系畢業生.
有一個叫做連嘉欣的弟兄.
他以外交醫療服務.
來代替台灣服兵役.
他用這個去到西非很貧瘠的地區.
不知道你有沒有聽過這個地區.
這個叫做布吉納法索.
聽過沒有.
布吉納法索很小的地方.
他眼見當地人沒什麼衣服穿.
於是向武會和弟兄姊妹發出呼籲.
叫他們穿上衣服.
而且這個勸喻捐物的行動.
牽動了其他世界各地的人支持.
於是有1500箱的衣物寄來.
不過入境要付關稅.
關稅計算出來是700萬西非的法郎.
一個這樣的年輕人.
一個這樣的宗教士.
何來這麼大筆的錢呢.
他有一個辦法.
就是跪下向神祈禱.
所以他的跪墊是紙皮.
跪下的時候.
他連頭也放進去.
於是形成一個頭和膝.
三角形.
他寫了一些記錄的書.
叫做紙板上的三個洞.
跪下的地方.
就是因為迫切禱告.

$^{321}$日子有功.
造成一個記號.
奇妙的是當人很願意切禱告.
又竭盡本份的時候.
神動功了.
當地的省長.
他的爸爸病了.
連嘉欣是醫生.
請他來診治.
趁著爸爸的病.
趁著這個機會.
連嘉欣就請省長.
你可不可以幫我寫一個公文.
去到首都幫我申請減稅.
結果很奇妙.
那個減稅的費用.
結果只需要十分之一.
而且在最後的期限完成.
一個曉得常常禱告.
不住依靠神的人.
必定經歷不少神奇妙的作為.
我們要好好交託禱告.
甚至我們會經歷更多神奇妙的作為.
成會社區教會.
也需要眾人齊心好好禱告.
會經歷更多的作為.
連嘉欣弟兄這個年輕的醫生.
他做了禱告.
動員了布吉納傑索的地區.
幫他開水井.
起孤魚丸.
買地.
建宣教士總部.
人看來很難的.
不容易的.
他是一一的成就.
我們再看第九節.
第九節.
那裡說到禱告.
然而我們禱告我們的神.

$^{361}$又因他們的緣故.
派人看守.
晝夜防備.
充分地.
再前一章.
充分地顯示出.
希米常常禱告.
依靠神的經歷是怎樣.
他盡自己的本分.
盡自己的能力.
應該做的事.
有行動.
這不單止禱告.
另一面我們要作為.
行動.
一個想建立神國的人.
是極需要留意.
有一個屬靈長者說.
信徒事奉.
成功或失敗.
是在於禱告的功夫這一環.
因為這是一個屬靈的戰爭.
我們除了要讀經之外.
我挑戰弟兄姊妹.
要好好地個人地禱告.
有沒有禱告的生活呢?.
我們教會都鼓勵弟兄姊妹禱告.
所以我們每個主日.
九時在1050室.
都有一個神禱會.
歡迎你們來.
和我們家母講.
第三.
警醒作供.
守望黑己.
當時有外環.
也有內憂.
這裡有十到十二節.
有講.
你又看一看.

$^{401}$我用現代中文的日本讀給你聽.
有大的居民有一首歌說.
工人的力氣越來越衰弱.
有那麼多的稅額要搬.
我們今天怎能收完城牆.
我們的敵人以為我們不會看到他們.
不會知道發生了什麼事.
要等到他們到這裡殺害我們.
阻止我們的工作的時候.
我們才知道.
那時住在敵人當中的猶太人.
已經一次又一次地來警告我們.
告訴我們敵人的陰謀計劃.
第十一節講到這裡.
第十一節講到敵人狂妄自大的話.
第十二節.
是靠近猶太居住的猶太人.
他們通風報信.
講一些不利的消息.
害怕到居民不得了.
動搖他們的鬥志.
但我們看下一章.
尼熙米卻沒有被他恐嚇.
他越是這樣警醒守望.
艱難的局勢.
發動禱告.
守望.
奮勇起來更大.
你看看.
這是十三到十四節.
那裡有講的.
你看到這裡.
尼熙米以家族為單位.
組織武裝起來.
一些措施.
成為守望戒備的防範.
你看到每個家庭.
每個家族.
為了神的家.
為了當時的情況.

$^{441}$怎樣努力.
今天成為的每一個弟兄姊妹.
我們也可以這樣做.
警醒關顧守望.
這個行動是有效的.
因為敵人的詭計.
心理攻勢都瓦解了.
你看看第十五節.
第十五節.
一起讀一讀.
道的聽見我們知道他們的心意.
見神也破壞他們的計謀.
就不來了.
我們都回到城牆那裡.
各作各的功.
你看到.
尼熙米並沒有鬆懈.
因為他知道這是長期的戰爭.
不容許片刻的放鬆.
繼續警醒.
由早到晚.
由天亮到星秀出現.
即是入黑的時候.
他都有很好的精心部署.
我請弟兄讀第十六到第十八節.
弟兄讀.
從那日起.
我的僕人一半作弓一半拿槍.
拿盾牌拿弓穿盔甲.
官長都站在猶太眾的後邊.
修造城牆的人.
共台材料的.
都一手作弓一手拿兵器.
修造的人.
都要間配刀修造.
摧割的人在我旁邊.
我們對貴就官長或其他人說.
這工程好大.
我們在城牆上相離甚遠.
你看到是一個圖畫.

$^{481}$尼熙米說我的僕人是指他的侍從護衛.
他們一半做事.
一半手拿著弓盾兵器穿盔甲.
作為防範和戒備.
他們防守的人.
是安定周圍的人的心.
這個也是所有修造的人.
通通武裝起來.
怎麼說呢.
工台的人是用肩膀.
他們的手還可以拿武器.
另一些人.
修造的人.
他們是武裝.
腰間有刀劍.
所以有需要的時候.
拔出刀劍.
因為城牆很遼闊.
所以分佈的很廣闊.
他們建立一個通訊網.
靠吹號角.
哪裡有問題.
那裡吹號角.
大家都知道.
所以通傳.
一切都是警醒的措施.
你是不是細意吩咐下去.
19到21節我想請姐妹讀一讀.
下一張.
19.
我對貴就官長和其他人說.
這工程浩大.
我們在城牆上相離甚遠.
就聚集到我們那裡去.
我們神必為我們爭戰.
於是我們作弓.
一本拿兵器.
從天亮直到星宿出現的時候.
好,看到這裡.
你會看到警醒有戒備的心.

$^{521}$神會為我們作弓爭戰.
這個心態很重要.
不是警醒是靠人的努力.
靠人精密的部署安排.
而且是在乎神的作弓.
這個是依靠神.
肅靈的爭戰的正義.
所以我們都要好好依靠神.
有這樣的心智的人.
必定有具體的行動.
22到23節.
我們一起讀一讀.
那時我又對百姓說.
神和他的僕人當在耶路撒冷住宿.
好在耶間保守我們.
伯奏作弓.
這樣我和弟兄僕人.
並跟從我的護兵.
都不脫衣服.
出去打水也帶兵器.
講得很具體.
所以我們可以看到一幅圖畫.
怎樣好禱告.
警醒守望.
另一邊好好靠主.
作弓.
這班人不計較麻煩.
辛苦的付出.
和代價多大.
今天我們說.
我們年紀有一定程度.
或者說我年紀尚輕.
很多事都不懂.
但是我們有沒有想.
神給我們有生命氣息.
原來我們的動作.
全都在乎神.
我們在有生之年.
有生命氣息健康活力的時候.
我們會不會立志.

$^{561}$2019年來了.
好好為主而活.
我們要為愛人.
愛教會而活.
甘願為主付代價.
好好行動.
在某一年.
在香港舉行福音會議.
的培靈言經會.
請到李思敬博士來講道.
他在結語裡面.
用一個精準的呼聲.
他說.
信徒特別是領袖.
很容易混淆個人的得失.
和使命的成效.
上帝重建耶路撒冷.
歷時一個世紀.
誰敢肯定.
在他有生之年.
就可以目睹教會的復興.
國家的迴轉呢.
我們的眼光.
如果只局限於一己生命的長短.
我們只會力求自保.
而不懂得捨己.
跟從基督.
在今天這個重視自我成長的氣候裡.
小心.
不要讓只有今生的人本假設.
而取代了我信復活的真理.
在不輕率.
耗盡的嚴謹當中.
我們仍然要堅持.
為主消盡成灰.
這種壯志和決心.
我想李思敬博士這番話.
亦提醒我們.
如何好好地警醒為主作供.
守望黑己的嚴謹表現.

$^{601}$今天.
我們在2018年.
最後一個的主日.
求主再提醒我們.
激勵我們.
齊心地作主供.
第一,各盡所能.
好好地為新分工合作.
第二,事上好好地禱告.
信心行動上要有好的操練.
第三,在警醒作供.
守望黑己上.
我們要有落實的回應.
願主祝福每一位.
成為一個蒙神喜悅.
修造的人.
齊心作主供.
我們低頭結束祈禱.
求主你直著利器米技提醒我們.
如何各盡所能.
好好地為新投入神理的家.
更幫助我們事上切切禱告.
信心行動上要有更多的操練.
我們更需要警醒作供.
在守望教會,守望弟兄姊妹的事上.
能夠付代價.
願主你幫助引導.
祈禱奉主耶穌基督的名.
阿們.
\newpage



\section{詩篇}
\label{sec:J6o6vqGpqR4}
\textbf{詩篇人生 - 詩篇 2 - 全篇講道 (陳崇基牧師)}
\newline
\newline
連結: \href{https://youtube.com/watch?v=J6o6vqGpqR4}{\texttt{https://youtube.com/watch?v=J6o6vqGpqR4}} ~~~~ 語音日期: 2019-08-08
\newline
\newline
\hyperref[sec:KWfIxDVxl3w]{\small{< < < PREV SERMON < < <}}
~
\hyperref[sec:index_chronic]{\small{[返順時目]}}
~
\hyperref[sec:index_scriptual]{\small{[返順卷目]}}
~
\hyperref[sec:6oKVLkh536w]{\small{> > > NEXT SERMON > > >}}
\newline
\newline
$^{1}$打開聖經到詩篇.
大概在《新舊約聖經》的中間.
詩篇是一本詩集.
裡面有多少首歌.
有150首歌.
裡面的歌也有不同的類型.
這裡有些數字告訴你.
類似有不同的格式.
或者不同類型的歌.
有讚美詩.
有信心的詩歌.
有敬拜的詩歌.
例如我們剛剛唱的一些詩歌.
屬於敬拜的詩歌.
是向上帝唱的.
智慧的詩歌等等.
哀歌也很多.
也有一些小部分.
但是很重要的.
叫做Royal Songs.
這些叫做君王的詩歌.
什麼叫做君王的詩歌呢.
你要明白什麼叫做Royal Songs的時候.
你要想一想.
當君王在職的時候.
登基的時候.
很多時候用一些這樣的詩詞歌賦.
來慶祝.
好像我們美國總統登基的時候.
就是Inauguration的時候.
一定會有國歌等等.
這些Royal Songs就有一些類似的用處.
就是一些以色列國.
他們登基的時候.
來慶祝.
這樣的情況來用.
但是你一定要明白.
古以色列國.
舊約的政體.
到底是怎麼樣.

$^{41}$去安排.
你就會明白為什麼會有這樣的Royal Songs.
這裡有一個很簡單的圖.
告訴你.
在舊約以色列國.
是一種神治的社會.
叫做Diocese.
神治社會就是說上帝是真正作王的.
他是唯一一個王.
King of Kings.
King, Capital King.
他去管治他的子民.
直處是他所揀選的君王.
所以他揀選了.
譬如大衛.
所羅門等等不同的君王.
來治理他自己的國民.
然後讓他們有律法.
他們活在法治的底下.
他們這樣生活.
其實就是要告訴周邊的鄰國.
告訴周邊的外邦國家.
上帝的律法是有智慧的.
以及上帝治理國度是截然不同的.
以往我們講過.
有人權.
有一些保護.
有一些跟其他近東世界的律法.
是有不同的地方.
當然有相似.
但是也有很多不同的地方.
這個就是舊約的政體.
舊約的政體就是上帝作為The King.
然後他揀選Kings.
這些Kings有時候叫他們做兒子.
長子.
Son of God.
Sons of God.
叫他們來治理國民.
有時候用的用語.

$^{81}$就用彌賽亞.
彌賽亞是什麼意思呢?.
彌賽亞希伯文文就是Machia.
其實意思就是說.
安立一個皇位登基的時候.
就用一些以前古代世界很貴的高人.
來高在他的頭上.
來安立他.
所以就叫受高者.
Machia就是Anointed的意思.
就是將他受高的意思.
接受他的恩高的意思.
所以為什麼有時候我們會唱彌賽亞.
Machia.
這個就是上帝所選的君王.
然後他會代表上帝來管治他的國民.
譬如以色列民.
或者之後猶大國的國民.
這些就是People of God.
Kingdom of God.
這些子民有時候也叫做上帝的兒子.
因為這些神國的子民.
就活在這個君王底下.
而這個君王就活在上帝這位君王底下.
The people in the King.
And the Kings are in the King.
Capital King.
這個就是以色列的正體.
所以每當新王登基的時候.
他們就會唱一些這樣的歌.
叫做Royal Song.
所以今天我們來看一首叫做Royal Song.
這首詩篇第二篇.
排名第二.
我們知道150首的詩歌.
第一篇和第二篇.
是好像兩扇的大門一樣.
就是給每一個來敬拜上帝的.
先看這兩篇的詩篇.
就好像打開一道雙門.

$^{121}$第一篇和第二篇.
預備好他們的生命.
才能夠進來神的面前來敬拜.
這首Royal Song是說什麼呢?.
第二篇.
愛國為什麼更難.
萬民為什麼謀殉虛妄的事.
世上的君王一起起來.
宰成一同相議.
說我們要睜開他門.
就是以色列國的捆綁.
脫去他們的成熟.
頭三節這裡說到.
外邦有爭論.
或者外族這些的萬民.
就是以色列以外這些的國民.
他們一起起來.
來相議或者去謀算.
去謀反.
反誰呢?.
就是反上帝的受告者.
他們說什麼呢?.
他們說要脫去他們的捆綁等等.
這是當時以色列國.
周一朝會發生的事情.
你乘著時光機回到大衛所在的時間.
你就知道當時是一些.
很多時候會有戰亂的時候.
每當一個國家有新的王登基的時候.
其他的國民就會.
重重複雜.
想作一些謀反等等.
每當一個王死了之後.
然後另一個王再立他.
來到他登基.
接續王位的時間.
是一個很動盪的時候.
各邦各族就會來到他抵擋.
以色列的仇.
以色列的君主.

$^{161}$想推翻他們.
圖謀作反.
就是去攻打他們等等.
這些在舊約的歷史裡面.
很多時候都會發生的.
譬如一些以色列周邊的國家.
以色列的死對頭.
譬如一些叫做菲利斯人.
以東人.
摩艾人.
阿摩尼人.
等等不同的.
在周邊不同的國家.
他們就是想去敵對.
上帝的律法.
以及上帝的統治等等.
這個歷史.
在靈性上.
這些國民或者這些所謂外邦.
爭鬧到謀反的這些國民.
似乎是在說一些政權.
說一些勢力.
好像我說兩師門.
好像《詩篇第一篇》最後說的.
他說是惡人.
這些惡人就是一些.
反對上帝.
敵黨上帝律法的這些勢力.
他們是肆無忌憚的.
無恥無義的.
他們去反對上帝.
這些世上的君王.
就是這一首王莊首先開始來描述.
描述是一種敵黨的勢力出現.
萬國來對上帝的和平.
或者以色列的和平.
這些好像是幾千年前.
社會的情況.
其實在我們今天.
都會發生的.

$^{201}$並且是陸續繼續會有發生.
就是上帝的子民被打壓.
或者被迫害.
上星期我和Johnny.
也和這位在利日尼亞的牧者.
這位阿爾牧者.
他一起和他見面.
去開會.
你記得前陣子新聞說些什麼呢.
就是在Nigeria.
有二百多個的女學生.
被綁架.
記得嗎.
他就在那裡出現.
就是在Nigeria.
利日尼亞被一些所謂.
Islamic militants.
伊斯蘭教的武裝分子.
武裝分子來到去.
搞得很亂的.
就是一種Islamic fundamentals.
這位牧者曾經在利日尼亞那裡.
他去到發表一些言論.
或者一些書.
然後他藉著這些言論去批評.
當地的.
福克華人的政權.
所以他就被打壓.
然後他就逃亡離開了他的國家.
兩年前的時候.
他就來到美國.
逃亡了來美國.
他就在芝加哥的南邊住.
現在就是在拿政治庇護.
在利日尼亞很多的基督徒.
很多的教會受到史無前例的迫害.
這是詩篇第二篇.
這首歌是Ronald Saunders所說的.
外邦爭霸或者萬民要謀樹一些虛妄的時候.
來去抵擋神的指紋.

$^{241}$抵擋上帝的權勢.
不僅是非洲是這樣.
不僅是利日尼亞.
中國更加不需要說了.
已經說了好幾次了.
最近林友買了一本新書給我.
這本新書叫做《新時代中國的宗教秩序與基督教》.
是嚴福珍教授所寫的.
最近有沒有看書?.
現在看書就好了.
很好看這本書.
因為真的是講事實和講資料的.
一本大概都是四百多本的書.
是新書出爐的.
兩個月前出爐的.
剛剛朋友經過芝加哥.
我叫他買了一本給我.
開頭這本書講什麼呢?.
導演說請用事實來說服我.
他說再想一想中共的宗教政策.
以及六招所謂為中共辯解的套路.
有六種的說法.
導演的時候主要想寫給誰呢?.
就是寫給在中國以外.
或者甚至海外的.
基督教律師.
因為事實上有一種說法.
流傳在坊間之內.
就是說.
中國官方的政策來作辯解.
譬如什麼呢?.
譬如說中國的教會受打壓.
這不是中央的政策.
只不過是地方某些部門發生的事情.
易富珍說不.
不要轉移事實.
第四他說什麼呢?.
他說與從前.
就是文革.
已經有很大的進步.

$^{281}$這就是很多時候為官方說話的一種流傳的說法.
就是說現在宗教自由將會越來越開放.
易富珍說請你用事實來說服我.
因為事實說不是.
因為宗教的自由是越來越差的.
等等等等.
寫給我們.
我再看一遍.
他去分析中國宗教最近的政策的時候.
他說中國已經發展成為一個成熟的警察國家.
全天候的視頻監控.
到社會信用制度.
還有各項國安和維穩的事項.
黨是無所不在的.
全面的去領導.
他引用了一位挪威的政治社會學家所說的一個形容詞.
就是叫做controloquism.
你聽過一個字.
democracy.
theocracy.
這是controloquism.
最後一句.
他說如果我們看不見這個事實的時候.
我們就無視中國存在很多的矛盾與問題.
更加不會承認我們所多數的所謂中國的夢.
中國的進步.
只不過是藉著強權高壓來維穩.
而犧牲的是公義.
你就可以支撐著一些所謂偽裝成神.
再下去他有一篇文章.
一個chapter.
他是回應一位學者.
這位學者一直以來為中國的宗教政策來護航.
他說什麼呢?.
我覺得很有道理.
他說宗教自由並非孤立.
而是涉及到結社自由.
出版自由.
表達自由.
言論自由等等公民的基本權利.

$^{321}$跟整個中國社會的現實息息相反.
如果你是關心中國教會的話.
就不應該將宗教自由的問題.
抽離於一個宏觀的處境.
中國對宗教的強化和管轄.
只不過是打壓基本人權.
以及公民社會的質實的寬容.
一本很好的書.
如果你真的想知道.
想明白這麼多.
請你借這本書.
不過先等我看完.
我應該很快看.
這些說話都是告訴我們.
就是詩篇第二篇的一個成果.
外方爭立.
要抵擋上帝.
這些勢力.
要去打壓屬於上帝的子民.
《佛誓講脈》就是詩篇第二篇所說的.
如何去面對這些溫和的詩篇呢?.
很多時候就是告訴我們.
到底我們應該如何去有所活然.
初代的教會.
第一世紀的教會.
可以回到初始.
第一世紀的文件的時候.
其中一本叫做《使徒行傳》.
Book of Acts.
這本文件記載初代教會的處境和歷史.
翻開《使徒行傳》的時候.
你就知道他們都受嚴重的打壓.
當時的希拉多.
大力地.
直接地間接地去迫害耶路撒冷的教會.
在他們的講道裡面.
Book of Acts.
《佛誓講脈》的詩章都會引用.
就是詩篇第二篇.
就是說外邦政權來打壓教會的時候.

$^{361}$正在發生在他們的眼前.
但是他們是有盼望的.
為什麼?.
我想詩篇第二篇這裡.
最重要的一個字.
就是剛才第二章第三節所說的.
虛妄的事.
萬民要去打壓或者反對上帝是Invade.
是一種很愚蠢的做法.
英文有一個譯本叫做New Living Translation NLT.
翻譯得不錯.
他說.
They waste their time.
to plot invade.
浪費了他們的時間精力.
來打壓教會或者去敵黨上帝.
為什麼呢?.
因為神仍然掌權.
他一定會受到.
審判一定會落實.
而且教會.
你看Book of Acts.
《詩章傳》你就知道.
仍然在打壓當中.
在困難的裡面.
是勢不可擋的.
第四至第九節.
剛才頭三節就說了地上發生什麼事.
第四節至第九節.
四五六七八九.
就帶領我們上到天上.
在天上發生什麼事.
「那坐在天上的必發笑.
主必恥笑他們.
那時神要在路中.
責備他們.
在路中驚嚇他們說.
我已經立我的君在石安(耶路撒冷).
在我的聖山之上.
受高者說」.

$^{401}$這個被立為王的大衛或之後的君.
「我要傳聖旨.
耶路撒冷對我說.
你是我的兒子.
我今日受你.
你求我.
我就將列國賜給你為基業.
將地極賜你為田寨.
你必用鐵杖打破他們.
你必將他們如同要杖的瓦器來摔破」.
我們現代人.
我們經常想像上帝.
Imagine God.
就是一個白鬚風.
他拿著拐杖.
很齒長的.
你戳他一下他就死了.
或者我們現代人會想像上帝.
我們Imagine上帝會是怎樣.
就是我們的好朋友.
他的工作是做什麼的.
就是要開解我的心靈.
我的心靈不要不開心的時候.
他會保持我們開心.
我們Imagine的上帝就是這樣的一種上帝.
或者我們Imagine上帝是一個菩薩.
這個上帝的工作.
就是要使我安居樂業.
每件事都是順順利利的.
這個就是上帝的工作描述.
詩篇第二篇也告訴我們.
真正的上帝是怎樣.
完全沒有我們剛才所想像的用意.
他描述的上帝是一位絕對有主權的上帝.
坐在天上.
他是不受威嚇的.
不會受到你威脅.
他說他笑他們.
笑地上的勢力來阻擋上帝的國度.
笑這個字.

$^{441}$其實是用來.
去恥笑一些無能者反抗的這種情況.
主要是去笑他們.
並且不單止有主權.
他們也會有權能的.
在路中去執行.
在烈怒之中去審判.
這些是我們不能夠想像到的.
一個這樣的上帝.
如果你說.
這些東西全部是人造出來的.
如果真的要造出來的話.
要哄人信的話.
那個人應該不做這些事情.
上帝去啟示他自己的時候.
他說他是一位有公義之路的上帝.
面對著這些敵黨勢力的時候.
神在哪裡.
神仍然在錫安.
Zion.
或者.
Jerusalem.
是同一件事.
Zion和Jerusalem.
他仍然立他自己的君王.
他說你是我的兒子.
我今天生你.
這些是立王登基的時候的一些用語.
不是真的說到那個王.
真的變成神或者半個神.
你是我的兒子.
我今天生你.
就是說我今天立你為我的兒子.
為我的君王.
就如他立大位為他的兒子一樣.
他的君王.
長子.
Phae Solon.
他立.
Solon.

$^{481}$Phae Adolia.
等等.
上帝是掌管一切.
並且好像堯將將這些瓦器打爛.
在形容一個圖像.
你見到一個K形的碗或者碟.
是不能夠再用的.
他就用猛烈的用棍勾過去.
然後就將它打爛.
這是公義上主的心.
以中第七詞.
你有沒有想到英文的翻譯.
是跟那個LORD.
世界.
接著ORD.
跟這個LORD.
四個都是LORD.
翻譯是不同的.
因為背後的字是不同的.
中文翻譯.
是這個意思.
為什麼呢?.
因為我們用音來翻譯.
希伯文.
大概那個字的音是Yahweh.
或者以前King James的翻譯是Jehovah.
所以中文流傳下來的.
就變成了耶和華.
在整本聖標裡面.
什麼時候首次出現?.
第一次真正出現的時候.
你記得出埃及記高考試.
就是摩西被警察帶回去埃及.
見打壓臣子民的法老.
這個暴君法老.
摩西就問上帝.
如果我真的回來了.
如果你們的子民問我.
Who is he?.
Who is your God?.

$^{521}$你這個上帝是哪位呢?.
請問.
他的名字叫什麼?.
當我們很多時候問.
Who am I?的時候.
或者Who are you?的時候.
我們通常.
我們入隊.
一定會回答.
答覆一定是有其他的參考.
Right?.
我們去指我們自己的時候.
通常我們都是用其他的參考去指我們.
其他的人士或者什麼.
也都會指出我們的身份.
所以我們是某某的異類.
某某人的異類.
我們是某某公司的員工.
我們是某某學校的學生等等.
但是上帝他介紹自己的時候.
他就不是這樣介紹.
3月14日.
神就對摩西說.
我是自由榮耀的.
你要對聖人說.
那自由者.
打發我.
到你們打我.
什麼意思呢?.
這裡翻譯就很笨.
因為I am who I am.
Who are you?.
上帝你是誰?.
上帝說I am.
I.
Right?.
我是誰?.
我是.
我.
這樣就不會回答問題.

$^{561}$我就是我.
但是這裡說的是什麼呢?.
你知道在基督教的神學裡面.
有一門學問叫做系統神學.
系統神學想做什麼呢?.
Systemic Theology.
其實就是想去總結.
整本聖經.
如何去討論某某的題目.
譬如整本聖經如何去討論.
基督.
耶穌基督是怎樣的.
系統神學就會出來說.
基督是這樣的.
或者整本聖經如何說神的子民.
教會.
系統神學就是去做這件事.
系統神學裡面說到神的時候.
用了一個字來表達上帝的本性.
自存式.
就是他是自有榮耀.
他沒有開始也沒有結束.
上帝沒有出世旨.
你有出世旨.
多少年,多少月,多少日,多少時.
你是出世.
你是開始.
上帝沒有這個開始.
上帝也沒有任何的Cause.
沒有任何的因.
因為他的存在不是依賴他者的.
這就是「我是自有榮耀的意思」.
就是「也我」的意思.
也我.
I am.
自由者.
差派你回去對這個暴權.
去講說話.
去拯救我的子民.
自由者.

$^{601}$意思就是說.
他不會怕一些有出世旨的人.
不會怕地上的惡權.
並且最後他會心痛.
這個詩篇第二篇.
給我們一個完全不能想像的生活.
不是人自己想出來的.
而是上帝啟示自己.
是一位自由者.
一位坐在天上作王的神.
大家知道詩歌.
歌這東西.
在不同的時間長有不同的意思.
我十多歲的時候.
十年的時候聽過這首歌.
當時很有感覺.
現在的感覺是完全不同.
以前你聽過歌.
到你今天再唱的時候.
又有些不同.
君王的詩篇都是一樣.
在以色列的歷史裡面.
曾經他們有國度的時候.
有以色列有猶太國的時候.
他們立王登基的時候.
他們用這些皇帝服裝.
就是慶賀上帝仍然在他們中間的位置上.
但是當他們流亡的時候.
流亡於巴比倫的時候.
以色列國.
因為在寶座上再沒有一個君王.
大衛已經不在船上.
大衛的子孫也不再在皇位上.
那怎麼辦呢?.
我們流亡亡國的時候.
這些舊約的神的子民.
他們會再唱這些皇帝的詩.
但是他們的意義有些不同.
因為他們是盼望著明天黎明.
現在他們在黑暗的當中.

$^{641}$不過他們唱君王的詩篇.
他們是以真相唱的.
不是以實相.
不是說看到有王坐在那裡.
他們才相信上帝.
而是見不到他們依然唱這些歌.
等等待將來.
上帝所揀選的王者一定會歸來.
這個就是舊約裡面一直有的.
所謂的以色列國.
將來王者.
彌賽亞會來的盼望的根源.
不是現在的事情.
而我們唱這首詩歌的時候.
我們在盼望上帝曾經對我們的不祖.
在歷史裡面.
他曾經留下的承諾.
對大衛說過.
他的國位是永遠不會結束的.
所以當他流亡的時候.
唱的時候.
他就是在尋找未來的王者.
將來的王者會.
你記得這個圖片.
剛才我們看的只是左邊.
現在我們看右邊.
當基督來到的時候.
他就說.
新約不死的天皇.
他已經滿足了.
第二邊軍王的所有盼望.
他就是神所揀選的軍王.
是K.
他是神的兒子.
是S.
他是 Messiah.
不是在舊約裡面的Messiah.
他是The Messiah.
The Anointed.
唯獨一位的宗保.

$^{681}$受保者.
他建立這個角度.
是超乎時間和地理的.
所有神的子民.
我們就是教會.
猶太人.
希臘人.
所有其他人.
到這篇詩篇最後的時候.
他就是邀請.
既然這位的王已經到.
他將來都會回來.
你現在做什麼.
他說你們軍王應該醒悟.
世上的審判官應該受到警教.
當全威懼.
極威懼的事件.
當全戰驚而快樂.
當以咀親自.
恐怕他發怒.
你們在道上面就要亡.
因為他的好戲快要發作.
凡投靠上帝的人.
都是不夠福氣.
以咀親自.
這是古墾通的一種語句.
即是當一個僕人來到一位軍王面前.
他下巴來親他的腳.
是代表效忠.
代表順服的意思.
代表他親近關係的意思.
最大的忠誠在這裡邀請我們.
在地上最大的Allegiance.
不是任何一個地上的黨或國.
而是天上的天上統祖的熱愛.
或者他所立的受告之位.
主耶穌基督.
我們歸屬是他.
不是地上任何的統神.
所以靠主耶穌基督.

$^{721}$所以為什麼這首詩篇.
你要讀了.
然後你才可以來到敬拜神.
敬拜神不只是唱歌.
而是你願意來到祂面前作祂的指紋.
來到祂的墓碑.
來到祂的荒廬.
來到祂的福事.
最後我想留下幾件事給你.
你信了耶穌基督給你.
你接納耶穌基督作你心靈的主的時候.
你就是以咀親自.
你在祂面前來祂的降服.
然後去敬拜.
去效忠祂.
你清不清楚認識神.
如果你不那麼清楚.
你需要認識多一點.
上帝的話語.
上帝的本性.
我們教會就要開.
第二輪.
第三第二輪的主教法.
基礎的信仰法.
我們就會告訴你.
到底聖經裡面所啟示,解釋的上帝.
是否是委任的上帝.
為受迫害的信徒.
教會來到.
我們這些這麼安殊的人在美國.
是不是?.
我們為了世界裡面不同的地方.
信徒.
我們來禱告.
去認識下他們現在到底發生什麼事.
他們面對著什麼處境.
為公義,為真理而去信徒.
我們有恩賜的.
我們有才幹的.
我們要去事奉上帝.

$^{761}$五件事.
請你拿著這件事離開.
在你心靈裡面.
應付.
下來.
跟隨耶穌.
事奉祂.
傳祂福音.
為天國而講話.
為其他人而有需要.
更深的去認識上帝.
我用心解讀.
多謝收看城匯網台.
記住Like,Comment.
或者訂閱,Subscribe我們.
\newpage



\section{}
\label{sec:6oKVLkh536w}
\textbf{3/10/2019 主日講道:神愛世人 (Marcus Yang)}
\newline
\newline
連結: \href{https://youtube.com/watch?v=6oKVLkh536w}{\texttt{https://youtube.com/watch?v=6oKVLkh536w}} ~~~~ 語音日期: 2019-08-09
\newline
\newline
\hyperref[sec:J6o6vqGpqR4]{\small{< < < PREV SERMON < < <}}
~
\hyperref[sec:index_chronic]{\small{[返順時目]}}
~
\hyperref[sec:index_scriptual]{\small{[返順卷目]}}
~
\hyperref[sec:cPh4mXwlq7I]{\small{> > > NEXT SERMON > > >}}
\newline
\newline
$^{1}$.
各位有聖經的話.
和我打開到約翰福音第三章十六節開始.
今日正道的經文是約翰福音第三章十六節開始.
神愛世人.
甚至將他的獨生子賜給他們.
叫一切信他的不至滅亡,反得永生.
因為神差他的兒子降世.
不是要定世人的罪.
乃是要叫世人因他得救.
信他的人不被定罪.
不信的人罪已經定了.
因為他不信神獨生子的名.
光來到世間.
世人因自己的行為是惡的.
不愛光,倒愛黑暗.
定他們的罪就是在此.
凡作惡的別恨光.
並不來就光.
恐怕他的行為受責備.
但恆真理的.
必來就光.
要顯明他所行的是靠神而行.
大家好.
非常感恩可以來到UV來分享聖的說話.
我是Marcus.
我是一個台灣人.
長大時也是說國語長大的.
在2017年的時候.
神開始給我一個呼籲.
是服侍說廣東話的群體.
從那一年開始.
神就開始訓練我怎樣說廣東話.
到現在應該都有兩年的時間了.
其實與其說是我服侍說廣東話的群體.
其實我都覺得是這個群體的人.
是栽培我,是幫助我的.
因為在我生命有難關的時候.
都是神栽培一些說廣東話的人來幫助我的.
在我2016年初初來到美國的時候.

$^{41}$我是第一次知道自己的生命是很軟弱的.
在這個陌生的大環境裡面.
我都覺得自己是一個很渺小的人.
但都是因為這些經驗.
讓我明白到聖是很大很偉大的.
聖是非常有愛的.
我在美國讀聖學院已經讀了三年的時間了.
神經常栽培一些說廣東話的長者來幫助我.
走這條服侍神的道路.
在我聽這些長者分享他們的二萬個故事的時候.
我都發現到其實很多人都經歷到相同的孤單.
相同的掙扎.
好像每一個人都經歷過一個人生的階段.
就是要面對一個選擇.
到底是要放棄回到自己的國家.
或者是要留在美國然後生活下去.
我相信每一個人背後都有一個故事.
每一個基督徒.
就是他在生命裡面都可以見到耶和華.
見到耶穌的榮耀.
每次聽完弟兄姊妹分享的見證之後.
我都會思想兩個問題.
第一個就是為什麼神要愛我們.
第二個問題就是為什麼神愛我們之後.
我們要怎樣回應他.
希望今天分享完神的話之後.
可以回答到這兩個問題.
我們剛剛有讀到經文裡面有提到.
神栽培他的兒子到世界上來.
不是要定死人的罪.
是要使死人因他耶穌基督來得到救恩.
也就是說信他的人就不會被定罪了.
這是一個很好的消息.
因為我們可以被神愛.
我們可以得到神的幫助.
都是因為我們得到耶穌基督為我們的罪而死.
但這也是一個非常不好的消息.
因為在第十八節後.
有限是說到不信的人就已經被定罪了.
譬如我的爺爺或者我的朋友.

$^{81}$我的家人他們沒有信主的話.
他們已經被定罪了.
他們不需要等到耶穌回到世界那天才會被定罪.
他們現在就已經被定罪了.
這番話其實對猶太人是一個更大的衝擊.
因為他們信了主是幾千年的時間.
但有限到這一刻才對他們說.
你們是得不到救恩的.
因為你們是沒有信了主的.
一直以來猶太人都有一個思想.
就是說他們是一個被神挑選的民族.
他們一直以為救恩已經臨到他們之間了.
但有限在這裡直接挑戰他們信仰的學生.
有限跟他們說救恩還沒有臨到你們之間.
但是審判已經臨到你們之間了.
幾千年的信仰到這一刻.
猶太人才知道自己信的不完整.
幾千年的信仰到這一刻.
才知道自己完全不知道聖罪是什麼.
就好像有很多人在美國居住了很多年的時間.
一直以來都以為自己是明白到美國的律法.
但是很多人出去美國旅行的時候.
我們知道最近有很多的案例.
他們出去旅行之後就回不來了美國了.
之前有一個Mudie的學生.
他回到大陸跟他的家人.
但是當他想回到美國的時候.
他就回不來了.
美國政府是有絕對的華語權.
決定誰可以回來誰不可以.
美國政府是有決定的.
所以同樣意思.
約翰這個說話對於猶太人來說是一個很大的震撼.
因為他們信了主幾千年.
但是到這一刻他們是得不到救恩的.
所以猶太人可能感覺.
哇有沒有搞錯啊約翰.
為什麼說這句話呢.
為什麼華邦人可以信了主可以得到救恩.
但是為什麼我們不行呢.

$^{121}$他們是非常不開心的.
你是不是搞錯了約翰.
他們可能會這樣感覺.
如果我們是猶太人的話.
我們是怎麼想呢.
我們有什麼感覺呢.
我們可以再繼續讀下去.
如果我們用審判的角度來看.
讀第18節和第16節的經文.
應該可以找到一個很出奇的地方.
因為第16節是講神的愛.
神對世人的愛.
但是第18節是講到神要審判那些不信的人.
這樣好像有點古怪.
因為好像神的審判裡面是沒有愛的.
這樣是不是約翰的邏輯有點不對呢.
我們可以再讀多一次這兩節的經文.
(閩).
原來這兩節經文其實有對比的關係.
神愛世人不是因為我們自己有多好.
神愛的世人在順了主之前就已經被定罪了.
意思就是說神愛的世人每一個人都是罪人.
所以神的愛是跟我們沒有關係的.
我們所做的事是沒有關係的.
我們是不需要去努力營造神的愛.
不需要去做一些特別的事情去營造神的憐憫.
或者是我們的救恩.
神拯救我們是因為祂給我們恩典.
是因為祂自己的特性.
祂給我們這個無條件的愛.
在我學習廣東話的時候經常發生一些很好笑的事情.
比如說我不知道怎麼講Snow的廣東話.
因為在台灣我是看不到Snow的.
所以每次看到Snow的時候就很開心.
有幾次我看到Snow的時候就很開心.
我就跟我的朋友說外面有綠帶血.
大家聽到我說這句話的時候就覺得很奇怪.
什麼綠帶血.
我看到他的眼神的時候就覺得.
我明白到我自己講錯了.

$^{161}$我是要講綠帶雪不是綠帶血.
所以我自己回家想像的時候.
都覺得很恐怖.
綠帶血是一灘血在地上.
是一個非常恐怖的事情.
所以他們因為一個字就明白不到我到底想講什麼.
另外一件事就是.
另外一個例子就是.
我之前在一個教會裡面.
有一個牧師要分享.
我們打開心扉神想認識我們.
但是因為我是分不清認識的識和吃吃這個字的分別.
所以我一直在那個教會裡面思想.
為什麼我信了基督之後.
神要吃我.
想吃我.
所以就覺得很古怪.
原來神是想認識我們.
不是想吃我.
所以都是同樣的意思.
猶太人他們信了主幾千年的時間.
但是他們完全不知道神到底在講什麼.
所以約翰也很明白到這個道理.
如果神沒有給我們恩典的話.
所有的死人都會在黑暗之中被定罪.
所以約翰在第十九節裡面解釋.
為什麼人會被定罪.
為什麼猶太人都會被定罪.
我們再讀多一下第十九節到第二十一節.
光來到世間,世人因自己的行為是我的.
不外光,道外黑暗,定他們的罪就是在此.
第二十一節.
凡作我的並行光,並不來就光.
恐怕他的行為受責備.
第二十一節.
但行真理的必來就光.
與顯明他所行的是靠神而行.
約翰一開始解釋的時候就講到光來到世間.
約翰講的光到底是什麼光.
我們教會裡面的光可不可以拯救我們.

$^{201}$或者是電燈站裡面的光可不可以拯救我們.
如果我們要明白到約翰到底是講什麼光.
我們就要讀回第一章的第六節到第十節.
有一個人是從神那裡出來的.
名叫約翰.
這人來為宇宙見證.
為光做見證.
叫眾人因他可以信.
他不是那光.
乃是要為光做見證.
那光是真光.
照亮一切身在世上的人.
他在世界.
世界也許是他做的世界決定.
不認識他.
所以約翰在這裡是想挑戰猶太人的信仰學生.
想挑戰他們可不可以看完聖經歷史.
因為在舊約裡面的一切的記錄.
都是指向一個終極的問題.
就是誰是理柴亞.
誰是救世主.
約翰在這篇經文裡面是說.
耶穌基督就是那個理柴亞.
就是世界的真光.
世界就是耶穌基督的創造.
猶太人希望以色列人可以明白到這個道理.
同樣的我們在讀這段經文的時候.
都去思想一下我們是否相信耶穌基督.
就是我們生命裡面的真光.
我們是生命裡面的救主.
之後明白到這個道理之後.
問題就來了.
我們相信這個真光到底有什麼用呢.
要明白到到底有什麼用.
我們就要再退一步去看整本福音書.
寫書的目的.
我們再讀一下第20章.
從第30到第31節.
耶穌在門徒面前另外行了許多的神蹟.
沒有在這書上.

$^{241}$但這些事是要叫你們信耶穌是基督.
是神的兒子.
並且叫你們信了他.
就可以因他的名得生命.
在我們學校上學圈子裡面.
大家都一直認同這兩段經文.
就是整本約翰福音寫書的目的.
在這21節裡面有一個字.
在我們學校裡面有特別的討論.
就是這個「信」字.
因為文化的關係.
這個「信」字有兩種不同的解釋方法.
一個是你們因此可以信.
或者是第二個.
你們可以認足你們的信心.
You may continue in faith.
所以我自己覺得是第二個意思是比較對的.
因為整本福音書是要挑戰猶太人的信仰學生.
是希望他們可以得到救恩.
在我們信了主之後.
我們知道耶穌基督為我們死之後.
之後要做什麼.
是不是信了主之後就什麼事情都不用做了.
我們再讀第三章第二十節到二十一節.
「凡作我的並行光,並不來走光,恐怕他的行為受責備,但行真理的必來走光,而顯明他所行的是靠神而行.」.
在很多翻譯本裡面.
第二十節的後面.
翻譯是「恐怕他的行為受責備」.
但是在原文或者是英文的ESV裡面.
那個責備的翻譯其實是有其他的意思.
可能可以翻譯成「暴露」或者是「揭發」.
「暴露」或者是「揭發」.
這個字是很有意思的.
因為意思就是說光可以揭發我們的罪.
就好像我們如求感光.
神就可以用光來暴露我們的罪.
如果我們如求感光.
我們的行為就會被警察揭發出來.
如果作我的人就會感到羞恥.
就好像有些政治家.

$^{281}$他們做一些不對的事情.
但這些不好的事情.
一場新聞就見光死了.
所以他們的功能可能因此就沒有了.
所以在讀到這段經文的時候.
可以思想一下.
我們的生命裡面.
有沒有一些東西是不可以見光的.
有沒有一些東西是一見光之後.
你就覺得很羞恥的.
當我們在讀這篇經文的時候.
如果加上第20節和第21節這兩段經文.
加上一起讀的時候.
就會發現到一個很奇怪的事情.
有些古古怪怪的事情.
約翰是用「做我的」和「行真理」這兩件事來做對比.
他不是用「做我的」和「做善事」來做對比.
為什麼?.
為什麼要這樣做呢?.
為什麼不這樣做呢?.
為什麼要這樣做呢?.
因為約翰很希望猶太人可以明白到.
耶穌基督的教導不是教你怎樣做一個好人.
因為基督教不是要教你怎樣做一個好人.
當然不是教你怎樣做壞人.
基督教要教我們的事情就是.
我們是罪人.
這是事實.
我們是自己拯救不了自己的.
我們可以思想一下.
如果耶穌基督是我們的保佑.
耶穌基督有沒有教我們怎樣做好人呢?.
或者耶穌基督是教我們怎樣用行真理呢?.
我們可以思想一下.
耶穌基督的生命裡面.
其實有很多人是反對他的.
最後一刻.
當你讀到福音書最後一刻的時候.
耶穌基督其實是以一個罪犯的身份死在十字架上.
他不是一個社會認可的好人.

$^{321}$所以約翰用俄式和行真理這兩件事來做對比.
就是因為想跟猶太人和我們表達.
行真理的時候有時候是得不到頒獎的.
不像做善事的時候.
可以有人為你歡呼.
為你頒獎.
為你開心.
約翰想表達的就是.
順了耶穌基督之後.
要行真理.
不是做善事這麼簡單.
因為有時候是不會有人為你頒獎的.
行真理的意思就是要跟父聖的旨意一起行.
在神呼喚我當牧師之前.
我的夢想是當一個飛機師.
我第一次接觸廣東話.
就是因為看一部片.
TVB拍過一部奇襲.
叫《沖上運輸》.
不知道有沒有人看過.
都是因為這個奇襲.
《沖上運輸》讓我開始喜歡廣東話.
就是我在大學的時候.
成都都安排了很多人進來我的生命裡面.
我們都很喜歡看TVB的奇襲.
還有成都安排.
甚至乎我的室友他們都是說廣東話的.
所以都是因為他們.
所以我都對廣東話非常喜歡廣東話.
我對香港都有很好的好感.
但是每一個故事.
每一個奇襲都有些高低起伏.
如果有高低起伏.
奇襲才會好看.
所以我的故事都是一樣的.
就是八年前.
我自己從台灣開飛機.
自己去到香港參加香港大學的一個研討會.
我永遠記得有一晚.
我自己去到香港的.

$^{361}$微敦道附近的一間大吉樂的餐廳.
我永遠記得那晚我拿到自己的餐之後.
發現到自己是沒有.
他們沒有給我筷子.
所以我自己走回拿餐的地方.
想跟侍應拿一雙筷子.
但是那個時候我跟侍應是說國語的.
當我跟他一說國語之後.
他就非常不開心.
他就拿了一雙筷子來叮我.
所以覺得我非常不開心.
因為真的非常沒有禮貌.
所以那件事之後.
我就跟自己說不想在香港說國語了.
我就不再說國語了.
但是三年之前.
神就開始翻轉我對香港的看法.
因為我初初來到美國的時候.
我的生命裡面有很多的需要.
所以神就派了一對之前去香港事奉的家庭來幫助我.
就是Brian和Rebecca.
第一年的時候他們煮菜給我吃.
帶我出去逛街鼓勵我.
Brian煮的東西真的非常好吃.
如果大家有機會可以試一下.
Rebecca也非常鼓勵我.
她很有智慧.
經常鼓勵我怎樣能夠堅持事奉神.
他們畢業之後.
神也很純實地.
差別地對另外一個家庭來幫助我追隨神.
那個家庭就是幾個月之前.
有來UV分享神經翻譯的李如華先生.
我想有些人應該也記得他.
也是因為李如華先生和Brian和Rebecca.
我對香港很有愛,很有負擔.
想服侍講廣東話的群體.
甚至乎之後我回台灣.
我可能也會在台灣的教會裡成立一個廣東堂.
事奉講廣東話的人.

$^{401}$或者是從香港移民過來的香港人.
但是在學習廣東話的時候.
也不是一直很順利.
我學習廣東話的時候.
心情也有些起伏.
因為事奉廣東話不是我自己的語言.
有時候事奉廣東話的群體的時候.
我也忘記了為什麼我要事奉神.
我忘記了自己是因為想回應神給我的愛.
所以才想事奉廣東話的群體.
我很常不知道自己究竟是要行真理.
想事奉神或者是想證明自己.
因為我的廣東話不好.
很常被其他人罵.
或者是有些人說一些不中聽的話.
所以每次經歷到這些的時候.
我就跟自己說.
一個大男人怎麼可以被其他人罵.
我一定要學到廣東話.
我一定要學到.
但是每一次有這樣的想法.
這樣的感受的時候.
神就提醒我.
要我讀福音書去看看.
讀耶穌基督的生平.
是要我明白到.
基督徒是要像耶穌基督一樣.
追隨神.
不是要證明自己.
耶穌基督的生命的故事.
耶穌基督拯救我們這個事實.
對我們來說是最好的鼓勵.
因為祂是用祂自己的生命來拯救我們.
祂是愛我們這群罪人.
昨天我收到一通電話.
我爸爸打電話給我說.
我爺爺行到他生命的終點.
我和家人一直以來都想把握.
最後的機會傳福音給他.
我也很希望我爺爺最後一下.

$^{441}$可以順主,可以上天堂.
什麼叫順主?什麼叫永生?.
這兩個問題一直在我腦海裡轉.
因為真的很希望我爺爺可以順主.
如果我爺爺順了主的話.
就是他最後一次吐氣之後.
他就可以上天堂了.
神就會帶他上天堂.
但如果他沒有順主的話.
就是他最後一次吐氣之後.
神就會帶他到天堂以外的地方.
因為他沒有順主,他就是一個罪人.
我們都不知道爺爺最後有沒有順主.
但是我的家人,未順家人最後.
決定用一個佛教的方法幫他辦一個禮儀.
我一直希望爺爺的救恩.
傳福音給我爺爺的責任.
不只是我自己一個人身上.
或者我家裡有順主的人身上.
因為我一直希望神會掌權.
會親自傳福音給我爺爺.
因為我真的很希望爺爺可以順主.
可以上天堂見我的奶奶.
如果我們相信死亡之後還有生命還有天堂.
那我更要明白到.
耶穌基督為我們的罪死在十字架上.
這是事實.
那我們更要遵守耶穌基督的教導.
神愛世人,神子將他的獨生子賜給我們.
叫一切順他的人,不自滅亡,奮德換生.
當我們順了耶穌基督之後.
我們要行真理,要跟隨父聖的旨意.
我們要傳福音給家人.
要傳福音給我們可能不喜歡的人.
因為我們還是罪人的時候.
就愛我們了,拯救我們了.
我想問問在你們心中有沒有一個人.
或者一個群體.
是你一直以來都知道你想傳福音給他們的.
但是你也明白到跟他們傳福音是很辛苦的事情.

$^{481}$你也知道跟他們傳福音是得不到獎的.
但是我希望,也想鼓勵你.
希望你可以因著耶穌基督拯救你的事實.
去傳福音,特別是傳給那些你不喜歡的人.
因為耶穌基督是因為你的罪而死的.
耶穌基督是給了我們無條件的愛.
給了我們的恩典,才可以讓我們事奉祂.
讓我們可以得到榮幸.
希望大家也去思想一下.
我們心中有沒有一個人.
是希望,是神要我們去傳福音的.
讓我們一起祈禱.
為我們的罪死在十字架上.
我們也很感恩你可以讓我們進入永生.
很多時候,生命裡也有很多的難關.
但是你也順勢地幫助我們走生命的道路.
也請你給我們勇氣和智慧.
可以去服侍你,事奉你.
請你打開我們的眼睛,見到我們的榮譽.
讓我們可以傳福音給那些我們可能不喜歡的人.
在我們覺得軟弱的時候.
也可以有盼望,可以見到永生,見到你的榮耀.
我們把我們的生命奉獻給你.
希望你可以帶我們去走這條行真理的道路.
謝謝收看 下次見 拜拜.
香辣椒.
\newpage



\section{}
\label{sec:cPh4mXwlq7I}
\textbf{1/20/19 主日講道: 信靠主 (吳君才牧師)}
\newline
\newline
連結: \href{https://youtube.com/watch?v=cPh4mXwlq7I}{\texttt{https://youtube.com/watch?v=cPh4mXwlq7I}} ~~~~ 語音日期: 2019-08-13
\newline
\newline
\hyperref[sec:6oKVLkh536w]{\small{< < < PREV SERMON < < <}}
~
\hyperref[sec:index_chronic]{\small{[返順時目]}}
~
\hyperref[sec:index_scriptual]{\small{[返順卷目]}}
~
\hyperref[sec:H3Nf68xn8Zg]{\small{> > > NEXT SERMON > > >}}
\newline
\newline
$^{1}$.
\newpage



\section{箴言 4:23}
\label{sec:H3Nf68xn8Zg}
\textbf{2/3/19 主日講道: 2018年最多人搜尋的聖經 3: 箴言4\_23(陳崇基牧師)}
\newline
\newline
連結: \href{https://youtube.com/watch?v=H3Nf68xn8Zg}{\texttt{https://youtube.com/watch?v=H3Nf68xn8Zg}} ~~~~ 語音日期: 2019-08-13
\newline
\newline
\hyperref[sec:cPh4mXwlq7I]{\small{< < < PREV SERMON < < <}}
~
\hyperref[sec:index_chronic]{\small{[返順時目]}}
~
\hyperref[sec:index_scriptual]{\small{[返順卷目]}}
~
\hyperref[sec:C4RfQvzd3mA]{\small{> > > NEXT SERMON > > >}}
\newline
\newline
箴言 4:23
\newline
\begin{longtable}{cl}
\hline
\hline
章節 & 經文 (和合本修訂版)\\
\hline
4:23 & \begin{tabularx}{0.7\textwidth}{X} 你要保守你心，勝過保守一切，因為生命的泉源由心發出。 \end{tabularx} \\ \\
[1ex]
\hline
\hline
\end{longtable}
$^{1}$很精彩的分享.
我覺得有時這些分享幫助我們了解這個世界.
我們基督徒應該有一個特別的世界觀.
一個和其他人不同的人生的觀念.
你可以去吃泰國菜,去泰國玩,去其他地方玩.
但我們不會忘記原來上帝在這個地上繼續工作.
祂在拯救人,讓人認識祂.
這件事真的不要忘記.
因為這唯一一件基督徒有的分別.
我們每逢打開聖經的時候.
我們就讓聖經來挑戰我們,提醒我們.
因為有時我們太過狹窄,只看到我們自己.
我們在生活中被上帝的話語提醒的時候.
讓我們知道神最終想做些什麼.
不只是討我們的喜歡.
很多時候我們去教會,我們參與一些所謂宗教的節目.
我們都是說我們自己喜不喜歡.
說我們自己有沒有東西落袋.
其實我們忘記了上帝最重要是要差遣我們去別的地方.
今天你來這裡敬拜完之後.
你認識主多一點.
你就被差遣到你的崗位,你自己的地方.
你就要想盡辦法用一些智慧,用一用腦.
祈禱去依靠主.
在你的生命的崗位裡面,怎樣可以再進一步.
能夠讓人看得更加清楚神是怎樣的.
這個才是最終的目的.
所以我們每一次講道,每一次打開聖經.
我們都希望能夠挑戰你去做這件事.
請你不要忘記.
今天這片道,和每一次的講道都不是讓你落袋的.
這個不是一個顧客的地方.
這是一個差遣的地方.
你讓主提醒你多一點,你認識主多一點的時候.
你今天離開這個地方,回去你的家.
回去你的婚姻,回到你的家庭,回到你的職場.
你就有一個更加新的目的.
我們這幾個星期都在講一些去年很受歡迎.
經常被人搜索的聖經.
這次是第三次.

$^{41}$我們講一段聖經,叫做《箴言》.
如果你有聖經的話,你可以打開.
因為銀幕上的,是給一些新來的朋友.
他們不是很熟悉聖經的.
如果你可以的話,打開你的電話.
去到這一節的聖經,是四章第23節.
《箴言》.
這裡說什麼呢?.
保守你心,勝過保守一切.
因為一生的果效,由心發出.
正如我剛才所講.
上帝去處理我們的時候.
是應該從心裡面去處理的.
《箴言》是講什麼呢?.
《箴言》是講人生的智慧.
讓神的子民活得有智慧的時候.
活得有上帝的生命的時候.
他的生活有所不同.
以至這種智慧能夠去感染到其他人.
這就是最終的目的.
不是讓我們舒服舒服地去生活.
成成功功地去生活.
而是讓我們最終在生命裡面見到.
讓人見到神在我們中間.
保守你心,勝過保守一切.
因為一生的果效,由心發出.
我們說到智慧.
我想問你第一件事.
你自己如果是0至10.
你會給自己智慧有多少分呢?.
0至10.
10是最有智慧的.
不要搞錯.
0就是最沒有智慧.
不一定是蠢的.
可能你拿了很多PhD.
可能你的智慧是0.
你會是0還是10呢?.
你在哪裡呢?.
在這個規模上.

$^{81}$我最近看到一個這樣的分類.
將人類的知識.
或者將人類的數據知識分類.
怎樣分類呢?.
他說最底就是數據.
Data.
Facts.
事實發生了什麼事.
這些是最原始的東西.
接著再上一層就是連結這些事實.
成為Information,即是資料.
我想你公司需要的就是資料.
無論你是做accountant也好.
還是做一個財務的管理員.
Data不是最重要.
你需要Information.
更加需要的是什麼呢?.
就是知識.
處理了,整合了之後.
怎樣去做一個決定.
這個是要有知識才行的.
但最重要的是要有智慧.
這個就是要經歷去累積的.
並且能夠去教授到其他人.
成為其他人的榜樣.
能夠傳授給其他人.
這個叫做Wisdom.
剛才你給自己多少分呢?.
你是不是在想數據很多時候?.
拿了多少學位?.
學校多少分?.
是不是在說Wisdom呢?.
智慧在聖經裡面是很practical的東西.
是要走出來,活出來的.
我不管你房裡面有多少張證書.
我不管你有多少學位.
你都可能是零智慧.
什麼叫智慧?.
什麼叫Wisdom?.
是一些很practical的東西.

$^{121}$我女兒開始學車.
昨天真真正正第一次.
我和她上路去學車.
都幾驚嚇.
因為我的車是MDX.
我只有一部手波車.
一部自動波.
那你先學手波車還是自動波?.
手波車.
真是有少少驚嚇,有少少流汗.
因為我的MDX沒有手煞車.
手煞車在腳上,在左腳上.
好像沒有什麼渣拿.
上車就跟她說.
這個tie 盤就在這裡.
倒後鏡就在這裡.
這些只是數據.
什麼叫Information?.
什麼叫資料?.
功能是什麼?.
油門在哪裡?.
油門和煞車.
你不要搞錯.
我都讓我兒子試一試.
因為他下星期要拿permit.
在parking lot讓他試一試.
我看一看他.
他右腳踩油門.
左腳又踩煞車.
我心想,你坐車嗎?.
幸好看一看.
這些就是Information.
什麼叫知識?.
就是路上有的處境.
你知道怎樣面對.
什麼叫智慧?.
一個有智慧的開車的人.
很practical.
就是.
走快了.

$^{161}$不單止你要navigate.
你還要有經驗去處理.
不同的緊急的事情.
爆胎怎麼辦?.
車撐不住怎麼辦?.
兩天前車撐不住.
那你做什麼?.
這個就是叫Wisdom.
或者你學音樂也是一樣.
數據的東西只是音符琴鍵.
到智慧就是叫做Improvise.
上到場的時候.
你能夠和別人合奏.
這個叫智慧.
不單止打捷機.
人生也是一樣.
我們剛剛信主的時候.
有很多數據.
聖經是怎樣組織的.
敬拜是什麼意思?.
幾點鐘到.
這些全部都是數據.
到最終的時候就是要智慧.
你經過很多不同的層面.
去認識的時候.
我受了洗.
公開見證了.
然後我懂得怎樣去看聖經.
明白敬拜的意思.
明白做基督徒的意義.
事奉等等.
到最後你就是教導.
和能夠做到其他人的Mentor.
這個叫做Wisdom的Level.
在聖經裡面.
就是告訴我們.
有生命裡面的智慧.
我們需要去操的.
不會自然而然的.
在舊約聖經裡面有五卷書.

$^{201}$那五卷書是什麼呢?.
我想考考你.
我以下就列出這些書卷.
哪五卷是智慧文學?.
哪一卷是智慧文學?.
《九音真經》是智慧文學.
哪五卷?.
下一張我告訴你哪五卷.
下一張.
不是這五卷.
再下一張才是真正的五卷.
《約伯記》.
《箴言詩篇》.
《雅各傳道書》.
《伯施針傳曲》.
這些全部都是講我們人生裡面的智慧.
可能你沒想過.
這五卷東西到底大體上是講什麼.
我們很喜歡箴言.
講一兩句像是揮春的東西.
還可以.
但你整體去看.
到底聖經裡面怎樣講智慧呢?.
人生的practical living.
什麼叫智慧很practical?.
其實就是skillful living的意思.
就是你心靈裡面.
由你的心態,你的說話,你的作為.
你所走出來的東西.
都是skillful的.
是有技巧的.
能夠讓人看到神有分別.
留了一個圖給你看清楚.
到底整體是講什麼.
人生什麼叫有智慧呢?.
這個圖像是一個pizza.
你沒想過詩篇原來是有智慧文學的味道.
詩篇第一篇.
記得嗎?.
詩篇十九篇.

$^{241}$第八十九篇.
講到我們跟神之間的關係.
所以詩篇打開的時候.
有人生的智慧.
逆境的時候.
你要繼續去敬拜神.
你可以祈禱.
你可以唱哀歌.
這是跟神的關係.
你在好境的時候.
不要忘記去感恩.
但在生命裡面.
曾經有逆勢轉順的時候.
上帝救贖過你之後.
你不要忘記.
上帝的恩典.
你可以去學感恩.
上帝醫好你.
你不要離開.
不要在病房裡面.
才去祈禱.
求神拜佛.
當逆勢轉順的時候.
你學習去感恩.
這就是人生的最中心.
神在我們最中心.
箴言講什麼呢?.
箴言講很多人生關係.
人際的關係.
你跟配偶怎麼相處.
朋友怎麼相處.
同事怎麼相處.
怎麼管理你的金錢.
怎麼管理你人生裡面.
一切的生活關係.
怎麼有智慧去生活.
溝通有沒有藝術.
這個不是學位學到的.
這個不是書本可以學到的.
你需要上帝的話語去提醒你.

$^{281}$你的情緒怎麼去管理.
雅歌也是智慧的文學.
因為雅歌是講愛情.
男女之間的關係.
告訴我們愛情的威力很大.
所以要小心處理.
傳道書有一句說話很出名的.
虛空的虛空.
什麼都是虛空.
什麼意思?.
他說日光之下的生活.
一切都是虛空的.
因為一切都會死.
你死了什麼都沒有.
你存了多少錢.
你存了多少科文K沒用的.
日光之下一切都會虛空的.
他想逼你想最終的東西.
但最重要老師告訴你什麼.
要生活得是日光以上的生活.
不要住在太陽下.
住在太陽下所有東西都沒有意義.
但住在太陽上.
在日光之上.
就是要去敬畏上帝.
認識上帝的時候有永恆的角度.
你今生做的事情才會有意義.
約百記.
讓我多花一點時間.
因為今年我們有些弟兄姊妹.
他們全年讀經.
現在又在讀約百記.
一天讀幾章很難受.
我一向都覺得約百記好像三文治.
頭兩章好像第一塊方包.
最後幾章就是最後那塊.
中間那些就很難受.
中間那些很難受.
我都有很多東西很不明白.
但我想多花一點時間.

$^{321}$告訴你約百記說什麼.
我們知道他很多時候是說苦難.
我們知道故事大概是這樣.
約百是一個無可指責的二人.
但他受苦.
受了很多苦.
然後那些朋友去辯論.
朋友最主要辯論是什麼呢.
長篇大論.
給你的感覺就是真的很長篇大論.
可不可以結尾很快地說.
不可以.
因為很多時候人都是這樣想.
那些朋友怎麼想呢.
約百二人受苦是因為他做錯了.
還是他得罪了上帝.
既然他得罪了上帝的話.
他就應該認罪.
如果認罪的話.
上帝就會將他扭轉.
但約百就不肯認.
他就說他是無辜的.
他不應該不值得受這些苦.
最後幾章.
就是上帝真的出現.
他沒有給約百任何受苦的原因.
他沒有跟他解釋背後的原因是什麼.
但上帝就問約百什麼呢.
他說你可不可以猜透我管理世界的智慧.
你受苦的時候.
你願不願意相信我這個猜不透的智慧.
你願意相信上帝的智慧嗎.
所以整本約百記是在說上帝的智慧.
猜不透的智慧.
後面那段有幾節我很喜歡.
他在告訴約百.
他的智慧.
上帝的智慧是在說猜不透.
不是像人一樣.
38章第25節.

$^{361}$他說誰為大禹分水道.
誰為雷電定道路.
誰為雨降在無人之地.
落在無人居住的曠野.
誰為荒廢的地方得到資源.
清草在旱地裡得長出來.
這不是人想的東西.
人想的雨在那裡做什麼.
沒有用.
我們常常想好人就有好報.
壞人就有壞報.
這就是我們的智慧.
但上帝的智慧不是這樣.
雖然無人耕作的地方.
無人吃東西的地方.
有菜蔬出來.
有雨降下來.
上帝的智慧是猜不透的.
但你願不願意相信上帝.
所以整本約百記似乎不是在說苦難.
不是解釋苦難.
而是說你在苦難裡面.
怎樣去思想神.
怎樣去認識神.
不是解釋你為什麼會病.
為什麼會受苦.
而是在病的當中你怎樣去認識上帝.
怎樣去認知上帝.
上個星期我再一次看一本.
很久之前看過的書.
叫做《Lament for a Son》.
這是一個很有名的基督教神學家寫的書.
他為什麼寫這本書呢.
因為他兒子25歲的時候.
在一次登山意外裡喪生死了.
這本書就是他喪子之後哀悼的日記.
是一本很難謹的書.
整本書從頭到尾再看一次.
不是很長.
不過是很難謹.

$^{401}$因為裡面是很辛苦.
你看到他經歷的事情是多麼坎坷.
今年一月的時候.
這位神學家.
他再一次被一個基督教的雜誌訪問.
訪問最後的時候問到這個問題.
他說如果你沒有辦法解釋苦難的話.
為什麼你不放棄神.
他就這樣回答.
正如他書裡面這樣回答.
他說今時今日都是.
他說如果我無法解釋苦難.
為什麼不放棄神.
他說我不能夠.
因為當我世察宇宙那種驚人的複雜性.
巨大的複雜性的時候.
他見到人類他自己的自我意識.
和智慧的奇蹟的時候.
他說他無法相信一切是偶然而然的.
必定是有一個難以測透的智慧.
想像力或者能力來創造這一切.
或者更加正確地說.
不是創造了這一切.
而是不斷地創造了這些東西出來.
他開始認識到神其實在表演.
Performance.
他在密西根的西部書房.
望出去秋天的景色的時候.
深藍的天空.
絢麗的樹葉色彩.
他就見到上帝的表現.
上帝的演出.
雖然只是很短暫.
但是一個很完美的一刻.
不是一切前而言.
沒有解答苦難的恩由.
但他問在苦難裡面如何去認識神.
信不信得上帝的智慧.
你回到那個圖畫.
你就知道整個圖畫是教授我們一個基督徒的人生觀.

$^{441}$這五本書就是在說以上帝為我們的中心.
上帝為中心才能夠活到一個有智慧不同的人生.
《箴言第一章》第七節.
很多人都很喜歡的.
也是其中一節的《箴言經》.
他說敬畏耶和華是智慧的開端.
你要開始有智慧的人生.
你的開端你的beginning.
就一定要先去敬畏耶和華.
敬畏不要以為是嚇怕的意思.
而是經而生畏.
是和上帝有一個正確的關係.
proper relationship.
你回到這一段聖經.
很多人去搜尋的四章二十三節.
他說你要保守你的心勝過保守一切.
因為一生的果效由心發出.
如果你的心對上帝是不正的話.
沒有以他為你的開端的話.
一生的果效就不能發出上帝對你的審美.
回到最後.
我想問你.
你的智慧有多少.
如果要你去判斷一下.
給自己分.
或者不要這樣.
你和你隔壁的那個人比分.
你覺得他在關係上智慧有多少.
十是最有智慧.
不要搞錯.
零是最沒有智慧.
不是說學位.
不是說你拿了多少PhD.
溝通有沒有智慧.
領導包括領導你自己.
self leadership.
或者領導其他人.
財務.
你的奉獻你對上帝理財的觀念.
你的分配是怎樣.

$^{481}$情緒的管理.
苦難裡面.
我們怎樣去管理我們自己.
這些一切我們都需要以耶和華為我們開端.
如果你在這方面能夠定一些目標.
我奉勸你.
如果在關係上你需要一些改善.
溝通上需要一些改善.
你不妨打開《針言說》.
裡面講了很多說話處世的藝術.
情緒也一樣在《針言說》裡面有很多.
約百記.
如果你正在看的話.
繼續看.
建立一個完全不同的人生的觀念.
可以同心協力地祈禱.
上主我們多謝你.
因為你賜給我們人生.
是一個完全嶄新的觀念.
我們所看的不是一個顧客.
只是在受你的福氣.
神.
我們求你.
讓我們的人生活出智慧.
在我們每一方面.
我們的溝通.
我們的關係.
財務.
我們的情緒等等.
能夠見到神的生命.
願上帝幫助我們成為一個智慧的人.
來見證自己的恩典.
祈禱奉祝耶穌的名字.
願上帝祝福各位.
下星期我們希望能夠和你一起去敬拜神.
\newpage



\section{}
\label{sec:C4RfQvzd3mA}
\textbf{2/10/19 主日講道: 施比受更為有福 (陳崇基牧師)}
\newline
\newline
連結: \href{https://youtube.com/watch?v=C4RfQvzd3mA}{\texttt{https://youtube.com/watch?v=C4RfQvzd3mA}} ~~~~ 語音日期: 2019-08-14
\newline
\newline
\hyperref[sec:H3Nf68xn8Zg]{\small{< < < PREV SERMON < < <}}
~
\hyperref[sec:index_chronic]{\small{[返順時目]}}
~
\hyperref[sec:index_scriptual]{\small{[返順卷目]}}
~
\hyperref[sec:NaL7AxYsROQ]{\small{> > > NEXT SERMON > > >}}
\newline
\newline
$^{1}$.
OK 恭喜發財各位.
希望各位豬年快樂.
Happy New Year to all of you.
It's the year of pig.
We're gonna have wish you piggy.
OK.
So I'm gonna preach the sermon.
I'm gonna share my message in Chinglish.
How's that OK.
講英文混中文.
今年豬年.
希望各位都能夠像豬一樣發財.
OK.
I wanna share a couple of images with you.
我想和你分享兩張相片.
第一張相片.
有隻豬去找一個紋身的.
那個紋身師就問他.
想紋些什麼.
Piggy, what kind of tattoo do you want?.
那隻豬就說.
I want to look rich.
我想紋一些東西.
令到人家看起來很有錢.
Alright.
如果你是那個紋身師.
你會紋些什麼.
What would you do if you were the artist?.
What tattoo would you put on the pig?.
That would make the pig look rich.
Dollar sign OK.
What else?.
Yeah, this is what he did.
You wanna be a piggy bank, right?.
You wanna be a piggy bank.
That you look rich.
在新年的時候.
我們常常祝福大家.
我們說恭喜發財.

$^{41}$What does it mean?.
恭喜發財.
What is it?.
Wish you well, right?.
Or have a lot of money.
This is all kinds of blessings that we say to our friends, our families, our neighbors.
In Chinese New Year.
We want them, wish them well.
Right? A lot of material prosperity.
A lot of blessings.
More money, more.
Maybe bigger house.
Bigger car.
多些錢.
多些.
層樓大些,多些食物.
I want to think about.
How does God expects us to handle our blessings.
We know that God is God of blessings.
Material blessings is part of His blessings.
神是一個很豐富,很多祝福的神.
聖經是這樣說的.
其實上帝給我們的祝福和供應.
是有物質上的供應.
物質上的,吃的東西,住的東西,我們享受的東西.
This is all part of God's blessings to all of us.
However, how do we do with it?.
What do we do with it?.
How do we handle them?.
其實耶穌經常說錢,財富的教導.
耶穌經常說錢.
所以我們教會也會說錢.
Jesus talks a lot about money and wealth.
That's why we in the church, we also talk about it.
Especially in year of pig.
We all want to be piggy bank.
So I want to share a very simple story.
Told by Jesus.
It has to do with our perspective on wealth and prosperity and blessings.
這個故事我想今天和你分享.

$^{81}$就是在路加福音第十二章.
這裡耶穌講了一個很簡單的故事.
就是和我們地上的祝福是有關係的.
我先讀一讀.
Gospel of Luke chapter 12.
Starting from verse 13.
Someone in the crowd said to Jesus.
Teacher, tell my brother to divide the inheritance with me.
Jesus replied, man, who appointed me as a judge or an operator between you?.
Then the man said to, Jesus said to the disciples, said to them.
Watch out, be on your guard against all kinds of greed.
A man's life does not consist in abundance of his possessions.
路加福音第十二章十三節.
他說眾人之中有一個對耶穌說.
他說夫子老師請你吩咐我的兄長和我分我的家產.
分家產.
耶穌說你這個人誰立我作你段氏的官.
給你們分家產呢.
於是對眾人說.
你們要謹慎自守免去一切貪心.
因為人的生命不在乎家道豐富.
And the story continues.
And Jesus told them this parable, this simple story.
The ground, the field, the farmland of a certain rich man.
produced a good harvest, a good crop.
He thought to himself, what shall I do?.
I have no place to store my crops or my harvest.
Then he said, this is what I'll do.
I will tear down my barns and build bigger ones.
And there I will store all my grain and all my goods.
And I will say to myself.
You have plenty of good things laid up for many years.
Take life easy, chill, eat, drink and be merry.
But God said to him, you fool.
This very night your life will be demanded from you.
Then who will get what you have prepared for yourself?.
This is how it will be with anyone who stores up treasures thanks for himself.
but is not rich toward God.
耶穌就用比喻對他們講.
他說有一個財主田產豐富.

$^{121}$自己心裡面思想說.
我的田產沒有地方收藏,怎麼辦?.
又說我要這樣辦.
要將我的倉房拆了.
另蓋更大的,在那裡好收藏我一切的糧食和財物.
然後要對我的靈魂講.
靈魂啊,你有許多的財物積存.
可作多年的費用.
只管安安日日地吃喝快樂吧.
神卻對他說.
無知的人啊,今夜必要你的靈魂.
你所預備的要歸誰呢?.
凡為自己積財,在神面前卻不富足.
也是這樣.
Very simple story, ok?.
This story is called the parable of rich fool.
叫做無知財主的比喻.
這個故事很簡單.
他有一個很有錢的人.
他有很多錢,他越來越有錢.
然後他就想,那我就可以安安日日.
神就對他說.
Guess what? Your time is up.
And you don't even know it.
Talking about rich man.
你看看這個故事.
你知道這個人沒有做錯任何事.
你知道他田產豐富.
This rich guy didn't do anything immoral.
He got a good harvest.
He got a good produce.
Quantity earning exceed expectation.
Exceed forecast.
或者他可能去年做得很辛苦.
Working really hard.
這些是他值得有的東西.
越來越多錢.
他有個good year.
他有個很好的回報.
Most probably he work overtime I guess.

$^{161}$And he deserve all this good things.
All this blessing.
And nothing immoral that he has done.
沒有做過不道德的事.
他值得得到這些.
But what's wrong with the picture.
Now, a surface of plain meaning of the story.
Will have us believe.
That Jesus is against all rich people.
Jesus is against wealth.
表面上看,你會覺得.
耶穌或者聖經是很爭錢的.
耶穌覺得人都要窮.
有錢他就最憎.
有錢他就要拿你的錢.
But this is a very superficial understanding of the story.
但不可以這樣去解.
不可以說這個故事.
就是說耶穌很憎我們有錢.
我們有錢就是有罪.
沒有這個意思.
You have to remember.
This story was told in a Jewish context.
A Jewish culture.
We all know in the Jewish society.
Financial management, financial gain, material blessing.
Is the backbone of Jewish culture.
All the Jews, most of the Jews are rich I think.
Many of the Jews.
Although there are many poor Jews in the world.
但我們知道這個故事.
是說這個故事的時候.
是對猶太人的文化,猶太人的社會說的.
我們都知道很多猶太人都有錢.
猶太人很會做生意.
耶穌就是跟這個背景.
跟這些人來說話.
而我們知道猶太人的文化.
管錢,賺錢.
是他們文化的軸心.

$^{201}$是他們文化的一部分.
Now we turn to the Old Testament.
The scripture of Jewish people.
It talks about the book of prophets.
Especially it talks about working hard.
Diligence, saving, right.
Managing your money well.
And managing your wealth wisely.
Instead of letting nature take its course.
All kinds of things that has to do with wealth.
And material blessings.
So Jewish people.
If we understand the Old Testament.
The Torah.
We need to save, we need to manage our money.
We need to gain and save.
如果你看看猶太人的典籍.
特別是舊約,有一本書叫做《箴言》.
你就知道裡面教很多東西都跟錢有關.
我們要去勤力做事.
我們要努力去賺錢.
然後我們要儲錢.
我們要好好管理我們的錢財.
任由它荒蕩地花.
這些全部都是猶太人的文化.
猶太人的教導.
So we know that Jesus is not against having wealth.
Jesus is not against having money.
Jesus is not against having material blessings.
所以你不要想耶穌說.
耶穌是受苦的.
或者神很不喜歡我們有錢有這些祝福.
Now go back to this parable.
I find this something very interesting.
Let me tell you how to read the Bible.
How to read the story of Jesus.
There is something you need to pay attention.
When you come to reading some of the stories or parables told by Jesus.
你看看這個故事的時候.
我想你留意一件事.

$^{241}$我想告訴你怎樣去看這些故事.
這些故事很有趣的地方就是.
不是比喻說了什麼才有趣.
而是比喻沒有說什麼才有趣.
So this is how you read the story.
What is most interesting is not what the parable is saying.
What is interesting is what the parable does not say.
到底那個故事沒有說什麼.
而它應該說的就不說.
So let's think about that.
Let's think about it for a moment.
What is not being said in the story is the most interesting part.
這個比喻沒有說的才是最有趣最值得留意的.
Now remember Jesus is speaking to a Jewish society, Jewish audience.
And all Jewish men believe in God.
Believe in the Creator.
Believe in the God of Blessing.
Every Jewish man begins every morning, each morning with one prayer.
And that prayer is called the Thanksgiving prayer.
阿尼莫德.
I give thanks.
I give thanks to the God of Blessing.
Every Jewish man begins every single morning when he wakes up.
He prays the same prayer.
I give thanks to the God of Blessing.
你要知道這個故事.
再提一次.
就是和一個猶太人的背景聽眾說的.
而我們知道每一個猶太的男人都是信創造主.
救贖主和上帝是一個滿有祝福的上帝.
其實每一個猶太人的男人.
But if you look at this story.
Alright, if you look at this story.
A Jewish story.
Talking about this rich guy.
Never a word is attributed to God.
Never a thought.
think about God.
So what's wrong with this guy.
This guy is suffering from some sort of.

$^{281}$You know, I guess spiritual Alzheimer.
He forgets about God.
Never talk about God.
Never thought about Him.
With all the material blessings.
But somewhere up there.
There's a higher power who grant me all these things.
這個就是最有趣的地方.
就是那個故事.
說到這個男人的時候.
他想都沒有想過上帝.
他提都沒有提過上帝.
他有很多很多的東西.
越來越多.
他的田產越來越多.
他賺到的錢.
但他從來沒有想過.
原來在生命裡面.
有一個更高的存在.
一個上帝.
在上面祝福他.
來供應他一些這樣的東西.
So Jesus is not against wealth.
Jesus is not against having money.
He is not against wise management of our resources.
especially our financial resources.
Jesus is poking fun of something.
When you receive tons of blessing.
With extra money in your pocket.
And then you have no thought of.
There has to be a being up there.
Granting us all this.
Providing us.
And sustaining us.
所以耶穌沒有說受苦.
你千萬不要覺得整本聖經.
就是不喜歡我們有錢.
不喜歡我們賺錢.
耶穌絕對告訴我們.
我們應該要好好做管家.

$^{321}$好好管理我們的錢.
但財主給予故事.
有趣的地方就在這裡.
就是有錢的沒有想過.
原來上面有一個上帝.
我的生命是由他來sustain.
我的生命是由他來供應.
而不是.
你從17-19的故事來看.
你看看17-19.
這個有錢的人對自己說什麼.
這是富家子弟對自己說.
他說了五次.
我會.
他說了五次.
我會做這個.
我會把我的田割下來.
我會建造更大的.
我會存放我的寶物.
我會開設銀行帳戶.
我會開設投資帳戶.
我會告訴自己.
我種了很多好東西.
你看看他說了什麼.
他說我的田產多豐富.
我要做什麼.
我要拆我的倉庫.
我要開多個帳戶.
我要開多個投資戶口.
我要收藏我所有的東西.
然後對我的靈魂說.
你不需要擔心.
他說了我的農產.
我的果子.
我的農場.
我的財產.
甚至我的靈魂.
我的靈魂.
我的倉庫.
我的生命.

$^{361}$每樣東西都是我.
我 我 我 還是我呢.
我想這就是聖經的主題.
主要的問題不是富有.
不是恭喜發財.
當然神想讓你發財.
神想祝福你.
主要的問題不是富有.
主要的問題是我們的態度.
對於富有.
我們的態度對於錢.
而耶穌想要的是這個.
最主要的問題在哪裡.
不是財富問題.
而是我們對財富的心態.
你拿著錢.
我們儲著錢.
到底我們的心態是怎樣去看它呢.
這是第十五節.
耶穌說什麼呢.
耶穌說.
祂說.
注意.
對於財富.
我出生在香港.
我十歲的時候.
我跟隨我的父母.
去見一個家庭.
一個大家庭.
所有的叔叔阿姨一起.
所以中國新年最好的事情.
對於一個十歲的小孩來說.
最好的事情是什麼呢.
就是你得到錢.
你得到紅包.
你得到很多現金.
我們有很多現金.
我非常開心.
我記得我得到五,六百元現金.
從豆禮士.

$^{401}$恭喜發財.
當我說恭喜發財的時候.
有錢進入我的口袋.
這是十歲的最好的事情.
我們去見他.
我們一起吃飯.
我們吃得很豐盛.
然後叔叔阿姨.
特別是叔叔.
叔叔阿姨叫了所有小孩.
家裡有很多小孩.
來玩個遊戲.
遊戲叫做魚,蝦和蝦.
你玩過這個遊戲吧.
我們全部豆禮士.
十歲的小孩當然很開心.
我記得那次.
我自己有五,六百元.
然後叔叔阿姨叫我們玩魚蝦蟹.
我們小孩當然是玩.
我們這麼多錢都玩魚蝦蟹.
誰做莊.
那個叔叔做莊.
我們在玩的時候.
我們很開心.
有輸有贏.
玩這個遊戲很有趣.
但這遊戲很糟糕.
為什麼.
我們玩了一場魚蝦蟹.
我們很享受.
但最後我們輸光了所有錢.
我們輸光了所有錢.
我們輸光了所有錢包.
他們都去哪了.
他們都回到叔叔那裡.
那是一個邪惡的叔叔.
那個叔叔原來是哄我們玩.
搶光我們的錢.
我從來沒有試過在人生裡.

$^{441}$有一個這樣的時刻.
我們覺得人真的很邪惡.
你給我們錢.
沒有條件給我們錢.
我們會發財.
但他把所有錢都拿回來了.
但在那時候.
我才明白了義氣和慈悲的分別.
還有我的看法.
在於富貴.
耶穌說什麼.
他說.
你要謹慎自守.
免得一切的貪心.
因為人的生命不在乎.
家道豐富.
耶穌警告他的徒弟們.
要注意.
要超越神.
對付所有的義氣.
因為人的生命不在乎.
家道豐富.
那是義氣和慈悲的分別呢.
貪和慷慨到底有什麼分別呢.
我想給你們一些想法.
我想給你們兩句話.
讓你們去想一想.
貪和慷慨到底是什麼呢.
我想我們在座中間.
沒有誰說自己會貪的.
我們不是貪.
我們只是省錢.
義氣和謙虛的分別是什麼呢.
我覺得耶穌定義義氣和謙虛的分別是.
義氣是關於為自己而獲得的.
而謙虛和謙虛是關於為別人而獲得的.
貪婪是什麼呢.
就是只為自己團職.
懂得感恩或者慷慨.
懂得去祝福其他人.

$^{481}$這就是義氣和謙虛和謙虛的巨大分別.
有些事要想一想.
義氣是關於提高自己的生活標準.
而謙虛是關於提高自己的獻祭標準.
貪婪就是只想提高自己的生活水準.
感恩和慷慨就是提高自己的施與標準.
生活標準和獻祭標準.
有些事要想一想.
神給我們很多東西.
有意圖讓我們分享和祝福別人.
特別是那些真的需要.
特別是他的國王工作.
上帝祝福我們很多東西.
其實想我們做什麼呢.
不是想我們團職.
而是想我們分享和祝福其他人.
特別是祝福他的工作.
我兩年前就讀了一本新研究.
《學術學院》.
它講述了父母.
把孩子養大.
把孩子養大有感恩心.
有感恩心.
有慷慨心.
這本研究的題目是.
"感恩父母養大感恩孩子".
這引起了我的注意.
是一本研究關於如何把六到九歲的孩子養大.
如何幫助他們發展感恩心.
這是兩年前報告出來的科學研究.
科學研究是說.
家長如何養大孩子.
特別是六到九歲的孩子.
如何培養孩子有種德性.
有種本性.
他們懂得感恩.
懂得慷慨.
裡面有一樣東西是他講到的.
他說.
"Greater parent gratitude and family religious involvement".

$^{521}$"were the only significant predictors".
"of more frequent parent-reported child gratitude behavior".
父母有感恩的父母.
懂得感恩的父母.
還有他們對信仰的投入.
生活的投入.
整家生活的投入.
是唯一能夠預測到.
就是孩子會觀察到他有感恩的行為.
我們知道孩子.
他們可以說謝謝.
但不代表.
他們說謝謝.
當他們收到某些東西.
但他們不代表.
感恩不是說謝謝.
感恩是比喻好態度更重要.
你得到某些東西.
你說謝謝.
感恩是更重要的.
感恩是對你的心.
而不是只是外在的行為.
說謝謝.
去感恩.
不是純粹說謝謝.
我們知道感恩是從心裡發出來.
孩子懂得說謝謝.
但他未必心裡真的在說謝謝.
我發現這份研究很有趣.
它談到"Role Model".
"Spiritual involvement of parents, adults".
"We are happy in simple things".
"We are happy to volunteer".
"We are happy to give, to share".
"All these have great impact on our next generation".
"There is a family culture of giving instead of getting".
"That heavily impacts our children".
我覺得很有意思.
你想想.
當我們成為成人.

$^{561}$當我們成為父母.
我們懂得去感恩.
我們懂得去慷慨.
懂得去施與的時候.
懂得去跟別人分享的時候.
這就是我們家庭的文化.
我們的家庭文化.
是跟我們下一代有很大的影響.
這就帶來耶穌在《事實》中說的.
他說.
The Lord Jesus said.
It is more blessed to give than to receive.
聖經記載了耶穌說了一句話.
他說"施比受更為有福".
真的有Blessing的人.
真的有祝福的人.
是那些願意施與多過願意去接受.
這是我第一段信息.
第二段.
應用部分.
就在午餐後.
我想在這裡暫停.
我先說一句話.
我們的教會正在建造和種植.
是因為所有的人.
他們貢獻.
他們捐獻.
他們捐錢.
他們幫助.
這是一家人的慈悲.
這是一家人的感激.
我們想繼續這樣.
我們已經三年了.
我們看到了很多的福音.
神邀請我們成為更好的管理員.
祂所給予的東西.
所有的物質和福音.
是神給予的.
我們如何看待它們.
如何處理它們.

$^{601}$我們會否只為自己而獲得.
還是我們學會了給予.
就像神給予我們很多.
我們成會已經三年了.
在我們開始的時候到現在.
我聽到很多這樣的故事.
就是一些感恩.
願意施比受更為有福的故事.
我們教會.
神的工作可以繼續延伸.
完全是因為這些故事.
新年流流.
恭喜發財之餘.
我們希望大家能夠多一點祝福.
但最重要的是我們如何在神面前.
活得更加豐盛.
施比受更為有福.
在上帝眼中做一個有福的人.
我想展示給大家一個影片.
這個影片是為了介紹.
2019年的新的管理員.
等一下我們有一個短片.
我想用這個短片來介紹.
我們新的管理員.
我們教會的領袖.
我們選出來的領袖.
我想請大家想像這個影片.
是一個邀請.
並不是只關於知道誰是誰.
但這個影片是一個邀請.
給所有人.
每一個人.
我們想想如何建立教會.
我們想想如何一起服務.
我們想想如何活得更加豐盛.
更加豐盛的禮物.
彼得到的更豐盛.
看過影片的時候想一想.
到底我們除了認識這些叫團成員.
這些舊的領袖之外.

$^{641}$我們想一想.
我們如何可以在上帝給我們這麼多祝福的時候.
我們做一個好的管家.
施比受更為有福.
\newpage



\section{}
\label{sec:NaL7AxYsROQ}
\textbf{3/31/19 主日講道:重獲生命樹 (張麗仙傳道)}
\newline
\newline
連結: \href{https://youtube.com/watch?v=NaL7AxYsROQ}{\texttt{https://youtube.com/watch?v=NaL7AxYsROQ}} ~~~~ 語音日期: 2019-08-14
\newline
\newline
\hyperref[sec:C4RfQvzd3mA]{\small{< < < PREV SERMON < < <}}
~
\hyperref[sec:index_chronic]{\small{[返順時目]}}
~
\hyperref[sec:index_scriptual]{\small{[返順卷目]}}
~
\hyperref[sec:iiRocZadlgU]{\small{> > > NEXT SERMON > > >}}
\newline
\newline
$^{1}$每年大概是在這個日子.
有幾天我是特別開心.
我不知道你們知不知道我說什麼.
我不知道你們有沒有像我一樣的感受.
是因為看到樹木突然間全綠了.
今天還沒有綠.
只是在發芽.
但是可以保證過幾天天氣暖一點.
就很快全綠了.
想起當天的樹變成這樣.
樹葉全下了.
因為冬天都會淋寒.
所以樹葉全下了.
其實很難看.
不單難看.
還有一種很淒涼的感覺.
很蒼涼.
很灰暗.
感覺很沒有希望的感覺.
但是當春天一來.
大地回春的時候.
我們發覺.
你和我都看到.
一下子就全綠了.
好像全都變成了.
你覺得嗎.
只要天氣暖兩天.
好像全都變成了.
全部都綠起來了.
這是我來了美國三十年.
我覺得很震撼的地方.
為什麼這麼厲害.
真的過兩天.
全部好像都變成了.
我們一起綠了.
真的從冬天到春天.
我們看到枯萎的樹木.
全部都變綠了.
給我們很多人生的啟示.
各位電視姐妹.

$^{41}$回望你過去一年.
你有沒有一些嚴寒冬天的日子.
有沒有感覺到在這些日子裡面.
你覺得很沒有希望.
很灰暗.
可能是一個心愛的人.
他在病裡面掙扎.
你看到很傷心.
甚至看到他死.
可能是一個關係.
一個夫妻的關係.
一個朋友的關係.
一個父母子女之間的關係.
因為各種的原因破裂了.
也可能是身體衰老.
無法走回頭.
問題就一天比一天多.
讓你覺得很沮喪.
也有可能是前途的問題.
你不知道未來幾個月會是怎樣.
你不知道你前途會是怎樣.
留在這裡或回香港或去哪裡.
你不知道怎樣的時候.
你很迷茫.
當我們面對這些人生的艱難.
無奈的時候.
有時候我們回來城會.
說真的.
都是因為見朋友.
因為我們心底裡覺得.
神好像對我們沒什麼幫助.
有時候我們覺得在痛苦裡面.
神是靜默的.
神是不是真的靜默呢.
我們再看.
由冬天枯萎的樹.
直到春天恢復生機.
你聽不聽到神的聲音.
每一年都是這樣演奏出來.
你聽不聽到神的聲音.

$^{81}$事實上我們的神不單止是.
每一年枯萎又復生的樹.
對我們說話.
你看清楚原來貫穿聖經.
創造我們的神.
祂藉著一棵樹.
對艱難的我們說話.
這棵是生命樹.
神要對我們所答應的.
神藉著這棵生命樹對我們答應的.
是沒有人可以答應得到的.
無論那些是多厲害的人.
多有錢的人.
多有地位的人.
甚至是一些很受歡迎的神明.
他們是不可以答應.
神藉著生命樹答應我們的.
所以各位弟兄姊妹.
今天我很想用三幅圖畫.
和你重溫聖經裡面.
講到生命樹的故事.
和神在裡面答應我們什麼.
第一幅圖畫.
是生命樹失落了.
生命樹失落了.
記得上一次我講道的時候.
我是講道創世記.
講道地是空虛混沌.
冤冤黑暗竟未答.
的確創世記的頭幾章.
有很多精彩和濃縮的創世故事.
告訴我們.
神創造人因為神的命令而失落.
你讀的時候.
有沒有留意到創世記.
在這幾章.
這一棵生命樹呢.
我們看這一節聖經.
創世記第二章第九節說.
耶和華神使各樣的樹.

$^{121}$從地裡長出來.
可以為人眼目.
其上的果子好作食物.
原子當中又有生命樹和分辨善惡的樹.
生命樹在這裡出現.
我們很清楚.
接下來的故事.
就是亞當及瑪夏娃被蛇引誘.
他們違背了神的命令.
所以他們吃了分辨善惡的果子.
結果就是被港族醫電完.
在這個時候.
在這個港族裡面.
我們不要忽略了.
不要看漏了.
聖經記載給我們的一個細節.
就在創世記第三章二十四節說.
於是把他趕出去了.
又在醫電員的東邊.
安設基路伯和四面轉動發火焰的劍.
要把守生命樹的道路.
原來不單止趕走.
神還要天使基路伯.
在門那裡把守著生命樹.
用火焰的劍把守著生命樹.
其實是非常嚴密的保安措施.
不讓人回來.
不讓人回來再接近生命樹.
神很嚴厲地說.
你們不可以接近生命樹.
你們是失落了這個生命樹.
不再關你們的事.
切斷了.
人失去生命樹.
就好像人失落天堂一樣.
Paradise Lost.
失去了之後怎樣.
結果我們看到的就是.
人類落在一個因罪而殘缺敗壞的局面.
每一天我們都體會到很多很多的問題.

$^{161}$人不斷在罪的惡果裡面受到傷害.
而且災害是沒完沒了的.
遠一點我們看到戰爭.
政權的壓迫.
飢餓.
逼迫.
廝殺.
強姦.
色情.
毒品.
我說得有沒有過分了.
沒有.
因為你每天看新聞都有這些事.
貼身一點.
關係破碎.
疾病.
衰老.
人在試探引誘之下那種掙扎和跌滑.
而承受到其中的惡果.
這些一切一切都是因為人失去了生命樹之後的病徵.
而我沒有誇張.
因為你和我都絕對體會到.
而且是時時體會到.
即使不是每天體會到.
都經常體會到.
受傷醫不到.
無奈很多時候衝不破.
難怪你和我都會感受到.
生命的冬天.
樹木的枯萎.
蒼涼無助.
這個是人在生命樹失落的景象.
但是聖經告訴我們.
不是沒有盼望的.
我帶大家看第二幅圖畫.
就是生命樹重獲.
怎樣重獲呢.
怎樣可以重獲這棵生命樹呢.
我先不說.
我想和你談一談一部電影.

$^{201}$又是叫做生命樹.
在2011年曾經提名奧斯卡最佳影片.
我想問哪一位看過.
舉手給我看.
舉高一點.
Tree of Life.
有幾個.
這套電影叫做甚麼呢.
Tree of Life.
Brad Pitt.
Sean Penn.
Jessica Chastain.
這些都是很出名的明星.
一部很出名的電影.
有人看了十分鐘就想離開.
告訴你都是這樣的電影.
但有些人從頭到尾都拍手掌.
感覺到很精彩.
這部電影很明顯是另類的電影.
頭十幾分鐘說甚麼呢.
你坐下來看吧.
星球體系的出現.
生命的起源.
嬰孩出生.
魚,爬蟲,恐龍的出現.
之後隕石墜落.
宇宙大爆炸.
生命怎樣進化.
經歷千數百年之後.
分子變成今天的你和我.
這套電影的頭十幾分鐘都是說這些.
一句說話都沒有說.
全部都是經典音樂.
或者是一些很特別的音效.
還有很精彩,很美的畫面.
一幕又一幕.
一個接一個.
很壯觀,很藝術.
其實是很好看的.
看得你呼吸不到.

$^{241}$很好看.
但你也可以覺得很沉悶.
為甚麼呢.
我不是打算看.
我看電影當然是看劇情.
十幾十分鐘都是做這些.
有些人覺得這十分鐘.
二十分鐘很享受.
有些人已經起床走了.
不過我告訴你.
這套電影不是沒有劇情的.
是有劇情的.
是一個在Waco,Texas的家庭.
有爸爸,媽媽和他的子女.
不過我告訴你.
這套電影幾乎像默片.
不是很多句說話.
數起來不知道有沒有幾十句.
是50,50年代,60年代的背景電影.
爸爸是非常重規矩.
是一個很嚴厲的人.
他不是壞人.
規矩很重要.
媽媽很重愛,很重恩典.
爸爸是象徵大自然的紀律.
你不做事,不勤力就餓死了.
你做錯事就受罰了.
媽媽是象徵恩典.
做錯事都應該以愛相待.
就是這樣,不同的角色.
好像很多家庭一樣.
這個家庭是經歷很緊張的關係.
爸爸因為這麼嚴厲.
跟子女沒有話說.
幾個兒子.
所以關係很僵.
沒有話說.
我長子特別生爸爸的氣.
爸爸很嚴厲,無情.
罵就打.

$^{281}$當他們受傷的時候.
去到媽媽那裡得到安慰.
又好像這個家庭一樣.
又好像很多家庭一樣.
這個家庭是經歷了痛苦和打擊.
特別是爸爸原本有一份很安定,很收入的工作.
突然間中年的時候失業.
家人有點徬徨.
爸爸在這個時候有很多生命的反省.
另外更加是一件大的慘劇.
就是小兒子去當兵.
誰知在戰場裡面死了.
由這個時候全家人陷入極度的痛苦和失落.
雖然是一套劇情不是很激昂的戲.
但是這個家庭所面對的問題.
其實很多家庭都會面對.
跳到這個電影的最後的部分.
也是最重要的部分.
就是大兒子人到中年.
是一個很成功的建築師.
但是他每一天的心裡面.
都是很悲痛的懷念在戰場裡死去的小孩.
終於有一天他有一個幻象.
有一個幻象.
他看到在一個海邊.
他和爸爸媽媽.
還有很多其他失去家人的人.
他們在這個海邊.
和死去的家人重逢.
是一個很安靜,安詳的景象.
家人重逢.
是好的.
是好,應該是好.
是喜樂的事.
很安靜,安詳.
這部戲在這裡結束.
而這個景象.
這個盼望.
就是這部戲所呈現給世人的盼望.
原來到最後我們可以和家人見面.

$^{321}$但是.
不見得他們非常的喜樂.
因為似乎.
在人的生裡面面對的很多問題.
他們還沒有得到解決.
他們並不是有一種.
可以克服這些問題的盼望.
家人再走在一起.
但是沒有再進一步的盼望.
說了這麼久.
我不知道你們想不想看這部戲.
如果你們想看.
可以在YouTube那裡看.
好像可以給幾塊錢看.
如果你看了之後.
請你們回來跟我分享.
但是你們喜不喜歡看都好.
我們都體會到這個家庭的人性的緊張.
張力.
體會到這個家庭的高低起跌.
也有時候會體會到這個家庭遇到的慘劇.
我們很希望.
有一天如果是到世界末日的時候.
或者是再來生見到家人的時候.
我們希望那個時候.
過去的一切破碎.
不復再有.
不再只是這麼低調.
這麼安詳.
那就算了.
我們實在希望.
這個世界有新天新地.
不再有眼淚.
不再有悲傷破碎.
但是這部戲沒有提供給我們這個盼望.
這部戲《Tree of Life》沒有提供給我們這個安慰.
雖然在這部電影開始的時候.
都引用了約伯記的一句話.
就是神對約伯說.
我納大地根基的時候.

$^{361}$你在哪裡呢.
你若有聰明你說吧.
那個時候神聲一同歌唱.
神的眾子也都歡呼.
好像是一個由神為開始的電影.
但是電影一直發展下去.
神並沒有角色.
而且神並沒有從神而帶來的盼望.
事實上不單止是這部電影.
有很多電影或電視劇.
很創意.
很新鮮地探討人死了之後.
來生是怎樣的呢.
末日是怎樣的呢.
可惜越探就越亂.
有時你會覺得.
如果這樣來生.
我寧願沒有來生就算了.
甚至有很多宗教所宣告的天堂.
都沒有提供給我們人生最初的盼望.
就是希望這個世界不再有罪.
不再有破碎.
我們能夠恢復當初創造美滿的樂園.
沒有電影.
沒有宗教.
給我們真正的盼望.
但是和電影相反.
聖經一直針對這個問題.
當人們在創世上失落的時候.
我們去到聖經最後那一段.
啟示錄那一段.
給我們看到一個重新而來的盼望.
清清楚楚地告訴我們.
和電影相反.
聖經告訴我們的生命樹.
是帶給我們絕大的盼望.
或者我們一起來看.
啟示錄第22章1到6節.
我們一起讀好嗎.
天使由子時樂.

$^{401}$在城內街道當中一道生命水的河.
明亮如水晶.
從神和羔羊的寶座流出來.
在河這邊和那邊有生命樹.
結十二月講子.
每月都結果子.
樹上的葉子乃為以至萬民.
以後再沒有就早.
在城裡有神和羔羊的寶座.
他的僕人都要事奉他.
也要見他的面.
他的名字必寫在他們的額上.
不再有黑夜.
他們也不用燈光,日光.
因為主神要光照他們.
他們要作王直到永永遠遠.
天使又對我說.
這些話是真實可信的.
我相信在座的兄弟姐妹.
很多都不是第一次讀這段聖經.
你也知道這段聖經是誰寫的.
是耶穌的使徒約翰所寫的.
約翰見到一個異象.
那時候耶穌已經死了.
使徒約翰見到異象.
他見到天使.
正如第一次所說.
天使帶我去到一條生命河那裡.
生命水的河那裡.
什麼叫做生命水的河呢?.
這個很重要.
什麼叫做生命水的河呢?.
事實上在舊約和新約聖經裡.
水時時都象徵著神的拯救.
而在這裡所說.
生命水的河是從神的寶.
生命水的河明亮如水晶.
從神和羔羊的寶座那裡流出來.
這條河是從神和羔羊的寶座那裡流出來.
這裡象徵著耶穌來到這個世界上.

$^{441}$成為人的代罪羔羊.
聖經前面清清楚楚的記述這件事.
而耶穌在十字架上流的寶血.
成為了生命河.
要將人帶回去當日失落的樂園那裡.
生命水的河是耶穌基督的寶血.
在這個河那裡我們看到什麼呢?.
這個河帶我們去看到什麼呢?.
看到有三樣東西.
第一,是生命樹依次一切受傷破碎的.
我們在電影裡看不到的東西,在這裡看到.
在其他宗教裡看不到的東西,在這裡看到.
這個樹結十二樣的桿子.
當說十二的時候.
很多熟悉聖經的弟兄姊妹.
馬上想到,連串起來的就是十二個支派.
象徵了十二個支派.
也象徵了所有被神救贖的子民.
生命河帶我們來到生命樹那裡.
看到這棵樹很豐盛,結果子.
這棵樹結的果子是覆蓋所有被救贖的子民.
而這些葉子是依次萬民的.
這棵生命樹的葉子.
是依次萬民先前所受的創傷和破碎.
這是第一樣.
生命河帶我們來到.
看到生命樹帶來的依次.
第二樣,生命河帶我們來到的是.
再無咒詛了.
我們看看聖經怎麼說.
在城裡面有神和羔羊的寶座.
他的僕人要事奉他,要見他的面.
他的名字必寫在他們的額頭上.
我們可以再事奉神,見他的面.
這意思是我們可以再接近神.
之前因為生命樹沒了,切斷了.
現在我們可以回到神那裡.
我們有尊榮事奉他.
而且神的名字會刻在我們的額頭上.
什麼意思呢?.

$^{481}$神的名字刻在我們的額頭上.
就是說我們是屬於神的.
在舊約聖經裡面,刻著黑名字在額頭上.
說我們是屬於某人的.
如果今天在額頭上寫UV兩個字.
就是屬於UV的.
神說,你們以前失落的人.
你們現在的額頭上有我的名字.
就是說你們是屬於我的了.
以前因為罪,我和你斷絕了.
現在你們回到我那裡,屬於我的了.
在我的家裡,就像爸爸迎接浪子回頭一樣.
其實這就是浪子的比喻.
爸爸家裡所有的東西全部都是你的.
父子的關係恢復了.
家裡的東西是你的了.
我們不再分開了.
是這裡所說的.
再沒有就坐了.
以前就坐帶來的分離.
以前就坐帶來的懺決敗壞.
眼淚相同.
不會再有了.
因為我們不再分離.
你在父家裡面.
第三樣,生命何必得情.
不再有黑夜的意思是.
罪的權勢不可以得情.
聖經說,不再有黑夜.
他們不用用燈光,日光.
因為主神要光照他們.
他們要作王,永永遠遠.
我們不再受罪惡的打擊.
我們回到神那裡.
神的光芒照亮我們.
我們會和神一起同在.
我們有榮耀的日子.
這是啟樹律所告訴我們的.
最後,和剛才說的生命樹電影相反.
這個不是幻象.

$^{521}$是神切切實實給我們的眼光.
請我們留意一下.
這裡所說的是很認真,很鄭重的.
看到這節聖經最後說.
有些人說,我不單看見這個異象.
天使還告訴我們.
這些話是真實,可信的.
生命河帶我們回來.
看到生命樹一切受傷破碎.
我們不再受咒了.
我們不再有黑暗了.
祂讓我們重獲生命樹.
不過我們停一停.
又想想,剛才說第一幅圖畫的時候.
說得很絕.
說到生命樹失落的時候.
神是要用天使去把守住.
不讓人碰到.
為什麼我們可以突然重獲生命樹呢?.
這個讓我們去看第三幅圖畫.
是全因為那個為我們被掛上在樹上的.
我們看一節聖經.
在泰書第三章十三節說.
基督既為我們受咒了.
屬我們脫離律法的咒詛.
因為經上記著.
凡掛在木頭上的,都是被咒了的.
基督為我們受咒了.
屬我們脫離律法的咒詛.
正常來說,人犯罪是要受刑罰的.
所謂律法的咒詛.
是指律法指控你的罪.
你犯錯,你犯罪會帶來的刑罰.
但這裡清楚地告訴我們.
耶穌基督為我們的罪受了這些刑罰.
受了這些咒詛.
使我們不用再受這些咒詛.
令我們想起另一節的聖經.
以上是書第五十三章那裡說.
我知道祂為我們的過犯受害.

$^{561}$為我們的罪孽壓傷.
因祂受的刑罰,我們得醫治.
因祂受的變相,我們得醫治.
大家都很熟悉這節聖經.
沒錯,讓我們再一次重溫.
是耶穌基督將我們的罪當成是祂的罪.
耶穌基督在十字架上為我們受苦受死.
為我們承擔了罪的刑罰.
為我們承擔咒詛.
微妙的是當日的十字架.
耶穌在上面釘死的十字架.
這個十字架不是用塑膠做.
不是用金屬做,不是用法普膠做.
是用什麼做的?.
用樹木做的.
十字架,耶穌被釘在木頭上.
這個十字架是樹做的.
什麼意思?.
十字架是樹上砍下來的木頭做的.
這個象徵是告訴我們.
耶穌釘在樹上為我們爭取.
為我們換取生命樹.
釘在樹上的基督.
為我們幫助我們重獲生命樹.
其實這是一個很有意思的象徵.
使我們不能不相信.
神的拯救是多麼的認真.
神對我們的愛是多麼的深切.
說到耶穌被釘在木頭上.
我想和大家再看多一次聖經.
這是一節很重要的聖經.
約翰福音第三章十四到十五節.
寫這個的同樣是寫啟示錄的.
是使徒約翰所寫的.
約翰說:摩西在曠野如何舉蛇?.
人子也必照樣被舉起來.
叫一切順他的都得永生.
摩西在曠野如何舉蛇?.
這段歷史故事.
我相信很多弟兄姊妹都記得.

$^{601}$當時神興起摩西.
帶領以色列民逃離埃及奴隸的生活.
可惜這些被神的恩典帶領.
逃離奴隸生活的子民.
他們在曠野裡面有很多的怨言.
遇到困難,遇到沒有東西吃.
他們多番的不信神.
他們多番的對神懷疑.
多番的埋怨.
因此神激怒到一個地方.
說要用蛇毀滅他們.
但是在這個時候.
神同時的伸出一個拯救.
他叫摩西做一條火蛇.
掛在樹幹上面.
凡舉頭望這個蛇的.
就必得活.
是這一段的歷史故事.
而約翰引用這段歷史故事來講耶穌.
他說摩西在曠野怎樣舉蛇.
人子也照樣被舉起來.
叫一切信他的都得到永生.
意思就是說.
耶穌基督也是被舉起來的樹.
舉起來的木頭.
叫凡謙卑.
好像當日在曠野裡面.
願意懺悔承認自己不信得罪神的.
只要你舉目.
你就可以得活.
一望耶穌.
你就可以得活.
實在這句話.
是人人都應該聽的.
在我們當中.
或者有一些是剛剛聽這個道理的人.
或者以前從來沒聽過.
原來聖經有這樣講.
如果你今天真的是第一次聽的話.
我鼓勵你.

$^{641}$你認真的繼續思量.
這個釘在十字架上的耶穌.
被木頭掛在木頭上的耶穌.
跟你人生有什麼關係.
你是不是都要謙卑的回到他那裡.
你就要拭目仰望得到拯救了.
好,讓我們再來總結今天講的三幅圖畫.
第一幅是生命樹的失落.
在人類遲早犯罪的一刻.
生命樹失落了.
我們難免因為罪而帶來殘缺敗壞.
痛苦破碎.
我們的人生就是這樣.
第二幅是生命樹重獲了.
神給我們看到異象.
有一天我們重獲生命樹.
我們可以回到樂園.
而我們有這個盼望.
我們可以重獲生命樹.
是因為那一位被我掛在木頭上的耶穌.
那一位被掛在木頭上的耶穌.
祂來到這個世界.
殘酷地被釘在木頭上.
而祂要帶來的拯救並不是短暫的.
而是永遠的.
弟兄姊妹,在這個世界.
如果醫好的會再病.
修補的會再破裂.
甚至當日耶穌叫人復活.
那些人都會再死.
所以要告訴你.
耶穌是釘在十字架上.
釘在這個木頭上.
要告訴你.
祂要來拯救的.
不是短暫的.
是一個永恆.
是我們永恆的生命.
我們永恆的命運.
就是我們回到生命樹那裡.

$^{681}$一開始的時候.
我和大家回望.
過去一年.
或者是甚至今天.
你有沒有經歷冬天這種枯萎蒼涼.
你是不是在這些事情裡面.
你感覺到你聽不到神的聲音呢.
但是神藉著生命樹.
再一次對我們說.
祂在乎.
而且祂一直伸手來拯救我們.
在耶穌基督裡面.
神告訴我們.
祂已經救了我們.
今天我們所受的苦楚是短暫的.
靠著祂對我們的愛.
我們是可以撐得住.
絕對撐得住.
而我們仰望那一天.
祂再來的.
祂將我們帶回生命樹的日子.
各位弟兄姊妹.
特別在這個紀念耶穌基督.
受苦受死的大齋期.
如苦期的日子.
讓我們再一次的仰望生命樹.
仰望這位為我們.
帶我們回到生命樹的主耶穌基督.
對我們的愛是這麼大.
是這麼不可思議的.
讓我們一起祈禱.
我們看到你為我們所成就的.
我們是無言以對.
因為我們是這麼不配的人.
我們其實和當日在抗耶的.
以色列人一樣.
我們時時都是以不信.
以不感恩.
以埋怨來對待你.
我們的生命裡面.

$^{721}$我們很多時候沒有想到.
你是我們的主.
我們抗拒你.
我們承受最意大利的殘缺敗壞.
就好像當日亞當和夏威一樣.
你的聖經再一次的提醒.
再一次的光照我們.
讓我們看到從你以來.
你所答應我們的事.
就是當我們舉目仰望.
我們必回到生命樹那裡.
我們能夠與你同在.
不再有晝坐.
不再有黑夜.
我們深進這個盼望.
這個應許.
是會帶著我們對今天難過的日子.
一切的艱難.
是給我們力量.
給我們安慰.
因為我們知道我們的神是得勝的.
我們的神是在乎的.
奉救主耶穌的名.
阿們.
願神的話語祝福大家.
主席.
\newpage



\section{}
\label{sec:iiRocZadlgU}
\textbf{9/30/18 主日講道: 基督徒成長12步\#11 核心帶領 (陳崇基牧師)}
\newline
\newline
連結: \href{https://youtube.com/watch?v=iiRocZadlgU}{\texttt{https://youtube.com/watch?v=iiRocZadlgU}} ~~~~ 語音日期: 2019-08-14
\newline
\newline
\hyperref[sec:NaL7AxYsROQ]{\small{< < < PREV SERMON < < <}}
~
\hyperref[sec:index_chronic]{\small{[返順時目]}}
~
\hyperref[sec:index_scriptual]{\small{[返順卷目]}}
~
\hyperref[sec:xXuB_PHmaoM]{\small{> > > NEXT SERMON > > >}}
\newline
\newline
$^{1}$主,我們在祈禱當中.
我們去認定你是那位聽我們禱告的神.
我們在生活裡面.
我們絕對不孤單.
因為有你自己同床.
主,你聽見我們的心聲.
聽見我們的訴苦.
甚至我們每一天經歷的.
主,讓我們今天來到這裡.
你是我們每一天的主人.
主,讓我們今天來到你面前.
我們讓你.
在我們生命當中.
心靈裡面去打掃清楚.
潔淨我們的罪.
主,讓我們今天在這個.
敬拜安息日的當中.
我們再一次被你自己的愛和恩典來重圓.
起使我們的生命和我們的生活.
能夠再重整過來.
在耶穌基督的恩典裡面.
我們不單止知道我們是罪得赦免.
我們知道我們有新的生命.
你會重新在恩典裡面創造我們.
以致我們的生活,我們的品格.
每一天更加長烈.
主,求你繼續在恩典當中使我們成長.
讓我們教會,我們身邊弟兄姊妹.
我們有一些生命裡面的重擔.
主,我們為他們祈禱.
無論是一些身體的疾病.
一些職業上,前途上.
遇見很多不同的難處.
很多未知之數.
主,幫助我們.
不單止是去憂慮.
而是我們願意去到你面前.
求你來帶領我們.
主,讓這個教會我們同心合意的.
來一起建築.

$^{41}$不用這麼多人來認識.
我們為下個星期會搬進來的.
Park Community Church.
為這個教會我們要感恩.
因為在芝加哥.
這個Park Community Church.
做走了很多.
祝福了很多.
主,他們Bridgeport的Campus.
來到這裡的時候.
主,我們也繼續祈禱.
祝福我們兩邊的教會.
讓我們能夠同心合意的.
在你們的恩典當中來成長.
讓我們盡力的去洗人.
來去歸住.
盡力的洗教會.
能夠在你們的道上.
能夠在你們的道上成長.
主,叫我們教會.
站立在地上的這個地方.
我們是代表你的.
我們不是一些很抽離的一班人.
我們為了世界上任何的苦難.
我們都會連向你們呼求.
我們為了剛剛在印度尼西亞.
印尼的海嘯.
地震,海嘯.
死亡傷的人數.
數以百計千計.
主,我們求你憐憫這個地方.
主,讓你自己的平安.
你自己的安慰.
臨到一些受苦難.
受災情影響的家庭.
主,我們為了中國大陸的教會.
被政權來迫害打壓.
有很多地下的教會.
家庭的教會.
忠於所託的不斷的來見證.

$^{81}$但是被封殘.
傳道人可能會失蹤.
主,我們求你親自來保守你的義.
在這個土地當中.
他們繼續來見證你.
主,我們求你親自來讓各地的教會.
在這個時代.
我們無論哪一個崗位.
盡心竭力的來為你作見證.
繼續讓我們聽你的道.
聽你的說話.
讓你的說話來挑戰.
來提醒我們.
我們感激到,奉主名.
我們由年頭開始.
到現在我們有12個系列.
每月一次.
上次我們說到.
剛剛上星期.
領導侍工.
我們希望在上帝的恩典中.
我們長進成為領袖.
成為一個領袖.
這次去到尾二.
就是核心領袖.
這就是我們剛剛開始了選舉的季節.
我們希望鼓勵弟兄姊妹.
我們儘管踏上來.
來做一個領袖.
甚至是一個教會的領袖.
事奉和做一個領袖.
是一個很好的經歷.
上帝在我們事奉當中.
去鋪排我們一些弱點.
一些盲點.
弟兄姊妹我們同心一切來事奉的時候.
我們生命裡面能夠成長.
所以我們能夠在我們生命裡面.
繼續去成長.
我們需要去挑戰我們自己.

$^{121}$所以這次尾二.
我們說到怎樣做一個領導.
怎樣做一個在教會裡面.
核心去帶領.
好像我們中間由開始到現在.
差不多三年了.
有很多的call的members.
和我們一起來同工.
以致我們繼續去發展我們的教會.
這是一個很好的經歷.
我們上個星期開始.
最好是怎樣說呢.
就是找一個聖經裡面的所謂領袖.
一個英雄.
讓他在生命裡面的領袖時刻.
他被上帝呼召.
或者他自己踏上去事奉的時候.
他在那個時刻有什麼讓我們學習到的功課.
上個星期開始我們看摩西這個人.
我們說摩西.
由聖經第一章到最後那章.
沒有什麼人可以跟他比較.
他是一個很重要的領袖.
一個被上帝用的人.
所以我們用摩西.
用摩西另一個領袖的時刻.
上一次我們說他是一個很有責任感的人.
他那些路見不平不都相助.
那種的人.
上帝呼召他.
他成為一個事奉主的人.
今天我們再看摩西另一個.
在他事奉的生命裡面.
另一個所謂leadership的moment.
我們重溫一下.
到底摩西的人生是怎樣的.
摩西的人生大概是一個.
整天去到.
你感覺上.
好像經過很多冤枉路的人.

$^{161}$很多detour.
我不知道你有沒有想過.
其實聖經裡面有很多人物.
跟我們人生現在處境.
我們有的經歷其實都有點像.
很多的detour.
好像經過很多冤枉路.
或者很多阻滯的路.
我們重溫一下摩西頭三分二的人生.
用地圖來看就很清楚了.
上一本書就是加拿大.
就是現在的以色列巴勒斯坦.
最後的時候交代什麼呢.
他說阿布拉罕爾什的後裔.
帶走了一個很大的家族.
他們因為那裡有饑荒.
不夠雨水.
不夠出產.
他們就去埃及避難.
所以就住在了埃及的尼羅河三角洲.
當時是埃及帝國的一段時間.
歡迎這些人在那裡居住.
所以這個所謂摩西的祖宗.
希伯來人.
大部分都搬到了埃及的地方.
他們去避難.
上一次我們知道摩西.
就是其中一個以色列人的孩子.
他是希伯來人家庭裡的一個兒子.
因為法路來打壓那些以色列人.
殺了很多這些孩子.
他是所謂的死刑種.
我們知道他被一個埃及的公主收養了.
救了他.
這個死刑種應該是在埃及皇宮長大的.
以前的習慣就是這樣.
埃及帝國會抓很多外族的人.
收留他們進來皇宮.
然後再跟他們重新教育.
然後派他們回去.

$^{201}$如果他們守迦南.
就派他們去迦南不同的城邦.
幫埃及做事.
以致發展埃及的土地.
摩西似乎就是這個時代的背景.
上次我們說到他在皇宮出來.
他看到自己的同胞弟兄希伯來人被欺壓的時候.
他終於受不了.
他就殺了一個埃及的管公.
然後法路發現了之後.
就去追殺他.
所以他就逃走.
去了一個地方叫做米田.
米田就是現在地圖右邊.
就是現在北面的沙地阿拉伯的地方.
是一些游牧民族.
在那裡被難平之後.
移民到那裡.
跟當地的一個祭師.
叫做葉剔娥.
在那裡暫居.
寄人籬下.
因為幫過他的女兒.
連女兒也給了他.
在那裡結婚生小孩.
在那裡流離失所.
有幾十年的時間.
結果有一次.
摩西帶著他的羊去牧羊.
去到西乃的曠野.
近西乃山這個地方.
Mount Sinai.
一個地方三角形的樹.
他就見到上帝跟他顯現.
上帝呼召他.
他說是時候了.
現在我要去拯救我的子民.
我的百姓.
他們還在埃及受苦.
就要派摩西回去.

$^{241}$糾纏了很久.
摩西最終面對自己的過去.
然後他順應這個呼召.
他就回到埃及.
跟法老對抗.
大家知道了.
有很多的.
十災.
過紅海.
他們經歷過曠野.
摩西就帶著以色列人.
差不多去到西乃山.
發生什麼事呢?.
有一群外族的游牧民族.
叫做Amalekites.
叫做阿瑪力人.
他想去偷襲這些以色列人.
然後有一場仗打了.
如果大家記得的話.
神有一個名號.
叫利西.
他得勝了.
快要去到西乃山.
要去領受上帝的約和律法的時候.
就是這個故事發生的時候.
看到這麼多箭咀.
摩西的人生.
雖然他在空間上.
他真的走來走去.
但其實好像.
反映我們自己現在所經歷的人生.
都一樣.
我們不一定是到處走.
或者我們其實也是.
離開了家鄉.
走了很多冤枉路.
但很多時候上帝就是.
要我們離開我們自己的地方.
我們的崗位.
其實是有比我們更好的計劃.

$^{281}$摩西也是一樣.
去到這裡.
快要去到西乃山了.
就有一個這樣的領導性的故事.
就發生了.
十八章第一節.
那裡發生什麼事呢?.
那裡說摩西的岳父出現了.
岳父來的.
為什麼說岳父嫁到那麼奇怪呢?.
他說米殿的祭司葉提羅.
聽見神為摩西和神的百姓.
以色列所行的一切事.
就帶著摩西的妻子西波拉.
就是摩西從此打發了回去.
不知道他什麼時候打發.
可能剛剛過完紅海打發了.
又帶著西波拉的兩個兒子.
來到神的山.
第七節.
摩西迎接他的岳父.
向他下拜與他親嘴彼此的文案.
進了他的帳篷.
摩西就將.
耶和華為以色列人的緣故.
向法老和埃及人所行的一切事.
以及路上所遭遇的一切艱難.
並耶和華怎樣答救他們.
都述說給他的岳父聽.
跟他說故事.
葉提羅因為耶和華.
代以色列一切的好處變甚歡喜.
葉提羅說.
耶和華是應當稱頌的.
他救了你們脫離埃及人和法老.
將這百姓從埃及人的手下救出來.
我現今在埃及人.
向這百姓發狂澳的時得志.
耶和華比萬神都大.
聽完這些古錄.

$^{321}$出埃及,過紅海等等.
上帝怎樣救他們的子民.
這是一個很驚喜的一件事.
如果你熟悉.
舊約聖經,希伯來聖經的話.
突然有一個人說了一句這樣的說話.
還是一個外邦人.
還是一個外邦的祭司.
他說他敬拜的本身不是真神耶和華.
他說原來耶和華比萬神都大.
他只聽到這些故事.
他已經很驚訝.
他去讚歎上帝所做的事.
發現這位上帝才是真一位上帝.
這很驚訝,為什麼呢?.
因為如果你熟悉上文.
發生什麼事呢?.
你會知道那些以色列人被耶和華救出來之後.
過了紅海.
他們在曠野走了幾個月.
他們對神有什麼反應呢?.
他們說沒有水,沒有東西吃.
然後大大發怒.
不如收拾包袱回埃及吧.
埋怨摩西,埋怨上帝.
連自己的子民都不會去感恩.
但居然一個外邦的葉剔羅.
一個外邦的祭司.
他只是聽見而已.
他沒有親眼見過,沒有親眼經歷過.
他就知道這個神是一個真神.
耶和華比萬神都大.
祂才是最大的.
我相信作者是故意想將一個外邦人.
祂有一個反應.
和祂本族的希伯來人.
做一個對立.
對立什麼呢?.
就是原來一個外邦人都懂得去感恩.
一個自己信神的人.

$^{361}$可能很多時候我們都只是在埋怨.
上帝為什麼不給我們這樣.
我們繼續這個故事.
這個故事告訴我們.
米殿人葉剔羅.
就是零碩不同的外邦人.
聽見故事,然後他說耶和華是最大的.
跟之前的那幫外邦人.
阿瑪利人和偷襲以色列人完全不同.
更加不像自己的子民.
希伯來人.
所以他出現了.
探摩西.
我們繼續這個故事.
第十八章第十三節.
第二天,探了他之後,吃完東西了.
第二天.
摩西坐在那裡審判百姓.
百姓從早到晚都站在摩西的左右.
摩西的岳父漢建他向百姓.
所做的一切事就問他.
你向百姓做的是什麼呢?.
你為什麼獨自坐著?.
摩西對岳父說.
這是因為百姓.
到我這裡來要去求問上帝.
我便在兩座之間來施行審判.
我又叫他們知道上帝的律例和法道.
一個領袖.
他帶領著很多子民.
很多時候有很多的爭拗.
或者一些的案件.
就帶到摩西面前.
摩西做一個審判.
但從早到晚他都在那裡做審判.
岳父見到女婿這樣.
岳父會有什麼反應?.
如果你是岳父.
看到女婿這麼辛苦.
你會怎樣?.

$^{401}$不做.
不做,叫他不做.
他就說quit了.
摩西對岳父說.
你做的不好.
not good.
你和這些百姓都會疲憊的.
你會exhaust(累積).
你會burn out(燒出).
因為這件事太重了.
你獨自一個人扮你.
你知道在聖經裡面.
從頭到現在.
從第七誡到第十八章.
只出現過not good這個字.
出現過兩次.
這一次.
not good.
你一個人not good.
頭一次在哪裡?.
就是上帝做人的時候.
他說那人.
亞當.
一個人獨居.
not good.
不好.
只有他自己一個去承擔.
做上帝的工作是做不到的.
摩西在這裡.
岳父跟他說.
岳父真的很疼他.
他說你這樣不行.
你這樣會死的.
不僅你死.
所有整個國民都會死.
這個是leadership(領導)的moment.
突然間好像亮了燈.
到底應該怎樣呢?.
葉蒂羅怎麼說呢?.
他說現在請你聽我說.

$^{441}$我幫你出個主意.
願上帝與你同在.
你要替百姓在上帝面前.
將案件就告上帝.
要將律例和法道教導他們.
指示他們在當行的路上所做的事.
並要從百姓裡面.
去到揀選有才能的人.
就是敬畏神若.
誠實無妄的.
恨不義之才的人.
派他們去做什麼?.
千夫長,百夫長,五十夫長,十夫長.
管理百姓.
叫他們可以隨時審判百姓.
大事才呈到你面前.
小事他們自己可以審判.
只要你輕散一些.
他們也可以同當此刑.
留意一下文理.
外府當然希望保住女婿.
就算他事奉也好.
也不要burn out.
但我想你留意一下上文下理.
十八章發生什麼事?.
他說葉蒂羅帶著老婆去見摩西的時候.
在什麼地方?.
在神的山.
西乃山.
如果熟悉聖經.
就知道十九章他們才去到西乃山.
二十章有十誡.
然後有其他律法.
即是說作者是不順時寫東西的.
這一章發生的事.
應該是頒佈十誡和律法之後才發生.
You get my point.
他故意將這件事扼在頒佈十誡和律法之前.
其實有什麼用意呢?.
他的用意是.

$^{481}$他想告訴我們.
如果我們要做leader.
要lead一個nation的時候.
lead一個people的時候.
人出的智慧和神的律法不一定對立.
有時候人想的獻計.
葉蒂羅的獻計.
怎樣去管理.
怎樣去服侍.
和神自己的律法.
神的智慧.
是不一定對立的.
所以神國的建立.
或者教會的建立.
是需要人的智慧.
也是需要上帝的智慧.
這個很明顯是作者故意想做的.
因為在這裡葉蒂羅獻計.
外邦人獻計.
然後十九章他們頒佈了律法.
就有上帝給他們的guideline.
他們怎樣去生活.
十戒等等.
人的智慧和上帝的智慧.
不一定對立.
昨天我去了一個conference.
整天去了.
由早到晚.
九點到四點.
我去了一個conference.
那個conference叫做fulfill.
是集合了一班在大自教區裡面的.
直堂的人.
就是我們這些牧者.
特別是一班種族教會.
ethnic church的領袖.
我們集合的時候.
有一些講員.
來分享.
有一個很出色的.

$^{521}$印度籍的宣教研究者.
Sam George.
他分享了什麼呢.
他說.
他專門研究現在世界大趨勢.
人口的趨勢等等.
他說ethnic church.
種族的教會.
是去幫助北美的基督教會.
科林派的基督教會.
不會下降的一個趨勢.
主流的教會開始下降.
應該這樣說.
百人的教會開始下降.
但種族.
multi-ethnic的church.
就是一個新的growth.
所以我們一班.
我們有不同的人.
我是做華人的.
有蒙古人的教會的人.
有中東人的教會的人.
有東歐人的教會的人.
有韓國人.
有東南亞的教會的church planter.
我們一起去明白.
到底現在神在做什麼.
他裡面說了什麼呢.
他說.
他引用一個在UIC的社會學家.
叫做Steven Warner.
UIC的.
在我們附近.
UIC的sociologist.
Steven Warner.
他說什麼呢.
他說21世紀.
現在這一代發生的事.
他說有件事發生比500年前.
宗教改革更加大的事件.

$^{561}$是什麼呢.
他說是migration.
是人口的流動.
人口的活動.
immigration 移民.
或者難民.
這種最大的趨勢.
是在影響21世紀.
比500年前宗教改革.
你知道宗教改革之後.
整個歐洲福音的運動.
開展了.
然後派了很多宣教士出去.
我沒有聽錯的話.
應該是20個million的人.
在各地建立教會.
現在有一個新的趨勢.
就是migration.
人們到處走.
到處去不同的地方.
無論他們被人displace(移民)也好.
什麼也好.
總之就是人的研究.
如何配合我們神的教導.
在這樣的浪潮裡.
我們應該要做什麼呢.
這就是leadership moment.
就是人的智慧.
人的觀察.
與上帝的誡明.
如何配合.
剛才說的東西.
下星期再說.
因為下星期第12節.
我們說什麼.
直堂.
我們說教會的planting.
我們繼續這個故事.
第14節.
於是摩西聽從他岳父的話.

$^{601}$按照他說的去行.
摩西從以色列人當中.
選擇了有才能的.
立他們為百姓的首領.
作千夫長,百夫長,五十夫長,十夫長.
隨時來審判百姓.
有爛斷的案件就呈交摩西那裡.
各樣小事他們自己去審判.
之後摩西讓他岳父回去.
他就回到自己本地.
沒有一個人能夠孤獨地來領導.
無論那個人.
十張刀,十把刀來.
沒有一個人可以這樣.
你知道以色列人去到這個情況.
他們發現.
原來那個領袖.
就是民眾成長的瓶頸.
摩西就是那個bottleneck.
不是要將他剷除.
他是一個很好的領導.
很有責任的領導.
但是如果只有他一個人孤獨去導航的話.
不會對大家有好處.
所以他需要立其他領袖.
千夫長,百夫長,五十夫長,十夫長.
這是第一次.
在聖經的歷史裡面.
有所謂delegation of leadership.
就是將他分享出去.
讓其他人能夠有領導的責任.
民主化就是這樣開始.
有千夫長,有百夫長,有五十夫長,有十夫長.
就是說每一千人裡面立一個領袖.
另外每一百人裡面立一個領袖.
每五十個人立一個領袖.
每十個人立一個領袖.
我想問你有多少領袖?.
在一千人裡面有多少領袖?.
如果你這樣算的話.

$^{641}$有131個領袖.
如果你算一千人的話.
131個就是13\%.
13\%.
每八個人裡面就有一個領袖.
每七個八個裡面就應該有一個領袖.
這就是突破瓶頸的方法.
就是大家share leadership.
你要知道《出家及記》裡面給一個這樣的策略.
給一個這樣的體現的時候.
是對當時的民眾.
對當時的社會來說.
是一個很新的idea.
因為你知道領袖.
以前古代的民族.
通常只有一個.
葉蒂羅就是那個領袖.
他在外邦米殿人當中是祭司.
只有他幾個.
但如果你繼續追溯《出家及記》裡面所說的.
這裡說要有131個領袖出來.
每八個人裡面應該有一個領袖被按立.
然後到第十九章之後.
神到立約頒佈律法的時候.
他說以色列人要做什麼呢.
他說你要歸我成為祭司的國度.
聖潔的國民.
Kingdom of.
不是singular priest.
是pro.
是眾數的priest.
就是所有的國民.
所有的屬上帝的子民.
信上帝的子民.
他們就是一個祭司.
這是一個劃時代的idea.
就是人人皆祭司就在這裡出來.
每一個神的子民.
都是有責任來好好管理自己.
打理好自己的家.

$^{681}$然後管理整個教會.
整個人民的靈明.
每一個人都是領袖.
每一個人都是牧者.
每一個人都不是羊.
你知道猶太教到今天.
都有些拉匹在說.
No Israeli is a sheep.
沒有一個以色列人是羊.
每一個人都是什麼.
Shepard.
是牧者.
所以今天我們三十靈州.
剛剛好.
就說到每一個人都是牧者.
你回去聽聽.
如果可以的話.
我們現在播放了.
我們一起聽.
我們每一個人都是牧者.
我們每一個人都是神的子民.
我們就是領袖.
所以摩西在領導什麼呢.
是領導一個nation of servant leaders.
每一個人都是祭司.
每一個人都要明白上帝的道.
明白上帝的guideline.
然後去到有責任.
去帶領教會.
我想問.
我們成會教會什麼時候才上軌道.
我覺得起碼有幾樣東西.
我們要做到.
第一我們要自足.
我們奉獻我們的經濟.
我們需要自足.
不單止自足.
應該還要多20\%.
不是剛剛好.
每一個人的責任.

$^{721}$你不要翹雙手.
好像我上次所說.
那些by standard effect.
人家都會.
No 你是一個leader.
你應該contribute.
第一要自足.
第二要自存.
我們不是靠別人來幫我們.
然後我們可以自己自存.
書記證券班別人幫我們做的.
青年項目別人幫我們做的.
應該是我們自己有leader.
出來承擔.
我們自存.
另外最後應該要自建.
我們自己有能力.
有潛力.
有potential去plan另外一個church.
這樣才叫成熟.
才叫上軌道.
我想我們距離這個有一點點.
還不夠.
自足自存自建.
這三字.
但不是中共那種三字.
每一個教會我們需要去起身.
然後去擔任我們領導的角色.
今年我們開始選舉.
這是一個很serious的matter.
聖經裡面告訴我們.
每一個領袖.
是在帶領屬靈的群眾.
你是代表上帝來帶領的.
我已經說過很多次.
所謂合不合聖經.
看聖經的時候.
不可以照辦主換.
聖經做什麼我們就照做.
我們應該解清楚裡面的精神是什麼.

$^{761}$然後應用在我們身上.
以前的聖經所寫的.
所有領袖都是按立的.
人選出來的.
保羅選出來的.
是不是可以這樣照辦主換呢.
我覺得不是.
我覺得應該.
現在當中的訊息是給我們的.
告訴我們我們的領袖需要負責.
成熟.
提名的選舉出來的都是一樣.
所以在我們的Quality Team裡面.
我大概break down給大家看.
我們要負責的東西有多少.
其實這個已經很簡化了.
還有很多的.
我們有財務,有教育.
有兒童不同階段的事工.
有敬拜.
為什麼今天你可以這麼enjoy坐在這裡.
因為有人負責.
有leader在這裡.
不會突然抽一根筋出來帶的.
No.
他們星期三晚上在練歌.
背後有很多人在做策劃和schedule.
還有recruit人.
阿貴在求你了.
出來事奉.
真的要命.
你知不知道事奉是一個尊貴.
是上帝給你的.
你不要到病床上的時候.
想事奉上帝,上帝不會給你的.
是個尊貴.
你明白嗎?.
是個尊貴.
不是給你的.
上帝如果叫你去事奉的話.

$^{801}$是尊重你.
是給你尊榮.
小組.
外展的事工,傳科的事工.
財務等等.
所以我今天想work up.
告訴你.
你人生裡面.
是不是應該有一個leadership moment.
由上個星期開始到這個星期.
我有一個領袖.
人人在UV都應該參與事務.
並且我們選出合適的領袖.
做管理有才幹的.
敬畏神的.
誠實的人.
做眾人榜樣的.
我們去提名,去選舉.
產生一個領袖.
讓我們教會繼續前進.
leadership moment.
今天上帝跟你說什麼.
我們同心,低頭,祈禱.
上主,我們多謝你.
因為你給我們事奉的福分.
事奉的機會.
我們在北美的教會.
我們不受迫害.
我們有很多的空間.
我們能夠事奉你.
主要讓你來容我們的生命.
容我們的才幹.
主要讓我們事奉.
能夠亡你的起源.
讓我們教會.
我們去選舉的當中.
我們產生的執事領袖.
能夠繼續的同心.
讓這個教會來長進.
使更多人來得到你的福氣.

$^{841}$祈禱奉主明禱.
\newpage



\section{}
\label{sec:xXuB_PHmaoM}
\textbf{3/3/19 主日講道: 智取人生 - 學會生命的智慧(陳崇基牧師)}
\newline
\newline
連結: \href{https://youtube.com/watch?v=xXuB-PHmaoM}{\texttt{https://youtube.com/watch?v=xXuB-PHmaoM}} ~~~~ 語音日期: 2019-08-15
\newline
\newline
\hyperref[sec:iiRocZadlgU]{\small{< < < PREV SERMON < < <}}
~
\hyperref[sec:index_chronic]{\small{[返順時目]}}
~
\hyperref[sec:index_scriptual]{\small{[返順卷目]}}
~
\hyperref[sec:8I6GZqi_kx0]{\small{> > > NEXT SERMON > > >}}
\newline
\newline
$^{1}$什麼叫做智慧呢?.
聖經裡面說的智慧.
不是說學位有多少.
也不是說你年紀有多大.
更加不是說你智慧.
智慧是什麼呢?.
我想問你.
你想想你人生當中.
你誰覺得是有智慧的人呢?.
你想一想.
你想想他的名字或是他的樣子.
在你人生當中.
不用告訴我.
OK?.
你覺得他有什麼特點呢?.
那個人是很有智慧的.
聖經裡面說的智慧是什麼呢?.
是一種藝術.
是技巧.
Art or skill.
Skill of living.
是你平常生活的智慧.
聖經說的智慧.
就是怎樣活出上帝給我們的生命.
有技巧.
有智慧.
有技巧.
有藝術.
有美感的活出來.
你人生當中看到他的時候.
不單止在他身上可以學到東西.
是一個很漂亮的生命.
在聖經裡面有很多的經文.
都是教我們.
怎樣去做一個有智慧的人.
有上帝智慧的人生.
我想解釋清楚一點.
在聖經裡面有不同的說話.
是上帝給我們的說話.
你有聽過什麼叫做先知書嗎?.

$^{41}$先知不是明天下不下雪的那種.
未卜先知不是什麼意思.
先知是代言人.
是上帝的代言人.
上帝直接跟他說話.
然後先知就將他神的話.
帶給其他人.
這個是從上而下.
這個叫做Prophetic literature.
或者是先知所說的話.
是直接由上.
上帝的說話.
將說話帶給我們.
有很多這樣的經文.
都是直接神跟我們說話.
但是智慧文學.
Wisdom literature.
或者是Wisdom的passage.
就不是.
Wisdom的passage就是.
智者.
屬靈的智者.
他去觀察人生.
觀察世界裡面有的定律.
觀察了.
然後他就.
領略了一些屬天的真理.
所以是有一點.
下而上的.
就是我們.
打開我們的眼睛.
來觀察一下人生當中有什麼.
值得我們學習.
然後就將他帶到上帝面前.
看上帝有什麼道理.
真理來教導我們.
聖經的智慧文學有五本.
百思針傳歌.
若百詩篇.
箴言傳道書.

$^{81}$還有雅歌.
這五本就好像.
Wikipedia那樣.
就將所有的.
聖經裡面告訴我們.
智慧人生.
智慧的.
生活的技巧.
就濃縮在這五本.
書卷當中.
今天我們想說.
傳道書.
並且傳道書裡面夾住中間.
最重要的那一段.
傳道書第七章.
這些有一點.
中國人的格言文化.
但是.
我想告訴你.
不是純粹格言那麼簡單.
而是智者去觀察人生之後.
領略了一些東西的時候.
將我們帶到上帝的面前.
這些很寶貴.
第一句說什麼呢.
第一字.
美名優於美好的高尤.
死亡的日子.
勝過誕生的日子.
往復喪之家.
比往燕落之家還好.
因為死是人人的結局.
活人要把這事放在心上.
憂愁勝於稀少.
因為面帶受用.
能使內心得到好處.
智慧人的心在復喪之家.
愚昧人的心在歡樂之家.
這裡有幾句.
經常將一些東西比較.

$^{121}$什麼比什麼好.
智者第一件事要告訴我們.
就是人生裡面的抉擇.
很多時候.
是最重要.
和重要.
我再說一次.
他告訴我們分清楚.
是誰的.
他告訴我們分清楚.
在人生裡面的抉擇當中.
很多時候不是單單好與壞.
而是重要.
和最重要.
你要選擇.
最重要的東西.
很多時候我們選擇的東西.
或者我們作的人生抉擇.
不是好或壞.
而是好.
和最好.
智者就是想引導你.
如何拿到最好.
最好的東西.
他說最重要的是什麼.
比膏油更重要.
在古代的世界.
橄欖油那些東西是很重要的.
那些是最好的商品.
特別是.
最好的.
第一次炸出來的橄欖油.
最好的膏油.
就是那些財產.
都不及什麼.
不及美名.
就是一個人好的聲譽.
那個名字.
是比一切世界上.
有的物質.

$^{161}$更加重要.
如果你有看新聞的話.
你都知道最近有很多不同的新聞.
令我們跌眼鏡.
真的.
就算有些所謂的熟齡的巨人.
熟齡的偉人.
他們都會有跌倒的時候.
權力會使人腐敗.
會使人做一些.
愚蠢的.
決定.
看完這些故事的時候.
智者就說.
我們需要警醒反省一下.
他說活人.
如果我們還活著的話.
我們就要想清楚.
智者說什麼呢.
美名的秘訣.
秘訣在哪裡呢.
怎樣一個人.
到他結束人生的時候.
好的聲譽.
中文有幾個字說得很對.
垂流青史.
或者什麼.
為臭萬年.
你的名字.
真的可以為臭萬年.
秘訣在哪裡呢.
就是了解我們人.
我們作為一個人.
我們的聲譽.
我們的名字.
其實是go beyond death.
你死了之後.
人還記得你的名字.
所以.
這裡他說.

$^{201}$你要留住你的美名.
你的美好的生命.
一定要好好栽培.
比一切更加重要.
雖然這裡說了一些.
我們華人很不舒服的東西.
說死我.
如果你今天第一次來教會.
大吉利是.
剛剛不是新年過了.
還沒過.
還差幾天.
差幾天.
OK.
差幾天.
你明白智者想做什麼.
智者想逼我們想一些.
最重要的東西.
很多時候我們搞來搞去.
我們人生都是搞一些不是很重要的東西.
不是最重要的東西.
他要留意.
我們人生裡面的結局.
比我們的開始.
更加重要.
當你留意到你的結局的時候.
你就好好地小心.
有智慧地過你的人生.
他說煙樂之家.
裡面是愚蠢的.
煙樂之家就是有喜事.
有派對.
有婚禮.
每個人都喜歡去.
有演唱會.
開心的東西.
我們每個人都喜歡去.
但智者想逼我們想.
這些場合.
這樣的聚會.

$^{241}$都是愚昧人.
沒有人生智慧.
在裡面.
上個星期.
你知道我的女兒.
剛好走開了.
我的女兒.
她是一對粉絲.
她是一對韓國的.
跳舞唱歌的.
叫做BTS的粉絲.
有沒有聽過BTS.
年紀大一點的可能沒有聽過.
你回去想.
你回去找一下韓國的BTS.
你就知道他們多厲害.
五月.
BTS就是來.
芝加哥Social field.
那裡有音樂會.
世界巡迴.
來到芝加哥這一站.
Social field.
有六萬多人.
一對韓國的樂隊.
要買票.
我幫我女兒買票.
她開始買的時候.
我們三部機在那裡.
她一部,我一部.
她媽媽一部.
什麼時候進場.
先進場買.
在上網裡面先進場買.
我們三個都點進去.
那.
原來這麼搶手.
三個小時已經買了.
最近的消息.
原來還再開多一場.

$^{281}$再開多一場.
星期六開,星期日又再開.
要加場.
我就幫我女兒買票.
我一輩子都沒做過這件事.
現在這麼麻煩.
要再買票.
買多少錢.
是不是.
買多少錢.
其實都OK的.
我們必.
爭位上網去買票.
他才買了兩張.
每張175元.
不是很貴.
因為我想.
BTS這麼厲害.
四五百元.
在哪裡賺錢.
原來175元.
還算前.
100的位置.
他才買了.
他告訴我.
你可以offline了,不用買了.
等了這麼久.
我也進去看看.
沒多久就到我了.
一進去就到我.
全世界都在買票.
我當然就點進去.
我剛剛找到兩張.
貼在一起.
上少少300位.
他說買了.
我都想了一會.
我真的買了.
我買多兩張.
因為你可以再賣出去.

$^{321}$Ticketmaster會讓你這樣做.
他說這麼多人要.
人們搶著要.
特別是少男少女搶著要.
沒所謂,買兩張.
無論如何.
誰知道話口未完.
兩個小時後,他朋友買不到.
我說賣給他.
就賣我那張給他.
我買150元一張給他.
我買回來100元.
加上稅120元.
150元賣給他.
那個人就願意.
他兩個人.
他的同學和女朋友一起去.
多謝我.
你可以拍拖.
我150元便宜.
我不是那些炒票的人.
你想想,我幫朋友.
這些事,每個人都喜歡去.
每個人都喜歡去.
這是燕樂之家.
燕樂之家.
你去聽淑天的道理,就沒有人搶著去.
因為我們每個人都沒有智慧.
前陣子我做婚禮的時候.
幫人家做婚禮.
婚禮之前.
負責拍照和做DJ的姐姐.
跟我說.
牧師,你的衣服皺了.
皺了,不太好看.
我跟她說.
你冷靜一點.
沒有人會在意我.
牧師做婚禮.
沒有人會理我.

$^{361}$結果婚禮結束之後.
他們分享了照片.
我看回照片.
一張都沒有我.
要不被新郎遮住.
要不被新娘遮住.
要不被伴娘遮住.
我就是這個意思.
這是經驗,我做了很多次.
我都不在意自己有沒有梳頭.
無所謂.
沒有人在意.
結束之後一定有人說.
What a beautiful wedding.
What a beautiful ceremony.
多謝你啊,牧師.
但我想告訴你.
不久之前,做完一次喪禮.
做完一次婚禮.
結束之後,有個中年男人.
走來.
經過,來我面前.
他跟我握手.
Thank you.
I needed this.
我正是需要這件事.
正是需要聽.
你剛才所說的.
這就是傳道書.
智者裡面所說的.
他說喪事比喜事好.
因為喪事.
會教我們人生的智慧.
所以為什麼智者說.
死比生更重要.
你開生日派對.
不是很重要.
但你要想清楚.
活人要將這件事.
放在心上.

$^{401}$因為死可以提醒我們的結局.
七章第三節.
他說憂愁勝過喜笑.
因為愁容能夠使內心得到好處.
希伯來文的意思是.
有個better heart.
我們有一個更加mature的heart.
這些是怎樣來的呢.
是從憂愁.
復喪.
或者一些.
喪事裡面.
來教導我們.
智者說話.
我們就聽清楚.
然後第五節.
他說如果聽智慧人的叱tak.
勝過愚昧人的歌唱.
愚昧人的笑聲.
好像窩底下的鴿子.
好像窩底下的荊棘爆聲一樣.
是虛空的.
勒索使智慧人.
變為愚昧.
賄賂能夠敗壞人的心.
責備和歌唱是作一個對比.
責備當然不易受.
人家去教你或者去矯正你.
罵你.
歌唱當然是.
動聽一點.
你留意他們兩個對比.
智慧人是單數.
你看看英文.
Songs of fools.
是中數.
有很多人對你歌唱當然好聽.
我們當然想聽.
但忠言是逆耳.
智慧人就是告訴你.

$^{441}$一句真心責備你的說話.
勝過千萬句動聽的歌聲.
所以.
當有人責備你的時候.
你豎起耳朵.
聽清楚.
是為了你人生有智慧的好處.
七章大六字是原末人的笑聲.
好像窩底下的.
燒荊棘的爆聲.
你以前用過一些.
不是那麼高科技的煮飯的方法.
燒柴那些.
那些.
腐枝燒起來.
別人就說很過癮.
但不會有溫暖的.
很快就會消失.
就是窩底下的爆聲.
就是那些派對的聲音.
就是那些很悅耳的聲音.
但智慧人告訴你.
不要聽那麼多.
一句責備的說話.
一次的喪禮.
勝過你百次的派對.
勝過你.
一百首好聽的歌曲.
第八節.
事情的結局.
勝於開端.
全心忍耐的.
勝過居心驕傲.
第九節.
你心裡面.
不要輕易的動腦.
因為腦袋.
是留在愚昧人胸懷裡面.
要忍耐.
忍耐.

$^{481}$吞了一口氣.
看事情看長一點.
忍耐是需要時間.
因為你要時間.
因為你要看一件事情的結局.
比它現在開始.
是更加重要.
換句話說.
很多時候我們太快的蓋棺定律.
這個就是愚昧人的怒氣.
太快的.
我們反應過敏.
我覺得這個.
也是我其中一個挑戰.
我自己做人.
不是很有耐性.
這是一個很大的挑戰.
怎樣輕易的.
怎樣看長一點.
神仍然在工作.
不要這麼老奴.
不要動了你的怒氣.
七章第八節.
全心忍耐的.
勝過居心驕傲的.
看看原文是什麼意思.
希伯來的文字.
有時候很有詩意.
他說全心忍耐的.
勝過居心驕傲的.
其實他說全心忍耐是什麼意思.
希伯來文的原文有意思就是.
你的心是長的話.
long time.
你的heart.
是看long term的話.
比你居心驕傲的好.
希伯來文的意思是.
高的.
tall spirit.

$^{521}$很高的.
在上的.
現在你錯了.
我很生氣.
我要砌你.
居心驕傲.
tall spirit.
智者說你寧願全心忍耐long spirit.
看長一點的時間.
take it easy.
你的生命.
你的心靈.
你可以用兩個方法去拉.
拉長它.
一是.
居心驕傲的.
拉高自己.
謙卑地忍耐學習等候神.
這就是人生的智慧.
愚昧人.
就是將自己的心靈不斷拉高.
心高氣傲.
去鞭屍也好.
去.
覺得自己比別人好.
這是愚蠢.
第十節.
他說你不要說.
有一樣你千萬不要說.
就是為什麼.
從前的日子勝過現在.
因為你這樣問的時候.
你就不是出於智慧.
這個就是.
不單只是人生的領會.
也是智者告訴我們.
神是怎樣看.
萬事.
用神的角度看.
有什麼不同.

$^{561}$意思就是說.
不要常說.
以往多好.
以前多好.
很依戀或很懷念過去.
如果你這樣說.
這個不是出於智慧.
因為你忘記了.
上帝仍然是掌舵.
你忘記了.
上帝仍然在將來工作.
不只是.
The good old days.
是一個生命的教練.
這樣教人.
我覺得也挺有道理.
跟這裡有一點像.
Make peace with your past.
So it won't disturb your future.
他說.
跟你的過去和解.
以致.
不會搗亂你的將來.
你現在有的生活.
你將來有的生活.
很多時候是懷念著過去.
All the good old days.
是不是.
無論是教會還是屬靈生命.
什麼都好.
不要只懷念過去.
唱回以前的那首歌就好了.
因為時光不會停留.
智者說.
最重要的是什麼.
最重要就是不要忘記.
神仍然在工作.
神仍然在你的生命裡面.
在你的將來裡面.
仍然要工作 要改變.

$^{601}$無論你今天年紀多大.
年長 年輕.
上帝仍然在你生命裡面.
繼續工作.
That's the future.
是上帝仍然會做新的事情.
不要只想著.
以往是多麼好.
想回到過去.
回不到去.
有智慧的不如就洞察一下.
上帝要我將來做什麼.
上帝要我怎樣去事奉他.
上帝要我怎樣去見證他.
十一節.
智慧和產業都是好的.
兩樣東西都重要.
不是說產業不重要.
對得見天逆的人.
是有益的.
因為受智慧的庇護.
如同受銀紙的庇護.
唯有智慧能夠保存智慧人的生命.
這個就是智識的益處.
上帝說.
智慧是一種能量.
是智識的益處.
他說兩樣東西都重要.
不過最重要的是什麼.
能夠保存你.
能夠去保存你的.
就是上帝給我們的智慧.
我們說過.
有很多人.
可能萬箭不保.
無論是政治家或者官員.
或者甚至.
一些熟人的巨人.
你回想一下.
如果他們當時.

$^{641}$是有一些智慧人.
來對他們說話.
或者旁人夠膽來說一句.
責備的話.
結果會差多遠.
這個就是我們需要去警醒.
我們想一想.
到底我們生命裡面.
是不是有缺乏智慧的地方.
七章第十三節.
他說你要察看神的作為.
因為神.
是為曲的.
誰能變為直.
遇坑通的日子.
你當喜樂.
遭患難的日子.
你當思想.
因為神使這兩樣東西.
並列.
為的事叫理.
查不出身後.
有什麼事.
剛剛頭說什麼呢.
我們要察看.
留意上帝在做什麼.
因為只有上帝.
才有主權.
上帝變成曲的東西.
不會變回直.
直的就沒有人變回橫.
因為這是上帝的.
sovereignty.
這是上帝的主權.
正是因為這樣的緣故.
上帝有主權.
他是完全in control的時候.
當你現在有.
坑通的日子.
你就應該感恩要喜樂.

$^{681}$不要忘記上帝.
原意.
原本的意思.
希伯來文是什麼意思呢.
坑通的日子你當喜樂.
其實就是你今天坑通好的日子.
如果你現在有good days的時候.
你就要活得好.
你活得好.
不要讓勝利.
充分投入.
你現在享受的時候.
每樣東西都差不多.
挺好.
得得點.
有好日子你會活得好.
你必須要活得好.
live a good life.
如果你有好日子.
live a good life.
有工作有健康.
live a good life.
敬畏上帝.
想想如何事奉他.
如何有智慧.
但如果你有患難的時候.
bad days.
你就要去察看上帝的作為.
你當思想.
想想上帝有什麼可以教導你.
如何得到更多智慧.
上帝才是主.
在我們坑通或患難的日子.
是生是死.
都是有上帝自己的心意.
我們是查不出的.
這是way beyond.
我們的comprehension.
way beyond.
我們的理解.

$^{721}$我們是完全不理解的.
這是我們的comprehension.
way beyond我們能夠插頭的東西.
人生裡面有很多謎語.
死那天都可能解不清.
為什麼會有這樣的際遇.
為什麼會有這樣的經歷.
你查不出.
不過最重要的是.
當你現在有的坑通的日子.
live a good life.
如果你現在患難的時候.
思想.
察看一下.
上帝有什麼作為.
做一個智慧人.
最後我想你做一件事.
剛才不是有兩單新聞.
你給分數.
現在你給自己.
給自己.
十或零.
你給自己多少.
智慧 十是最有智慧.
零是最愚蠢的.
你給自己多少.
你想清楚.
在你的生活.
你有沒有live a good life.
當你現在暢通無阻的時候.
當你現在可能有些阻滯.
有些患難的時候.
你有沒有領會到.
領略到上帝的智慧.
零至十.
你給自己多少.
你低頭多少.
你低頭.
主我們今天來到你面前.
我們去禱告.

$^{761}$因為我們需要你.
給我們人生的智慧.
讓我們活得更加像你自己.
讓我們活出你自己天國的生命.
當我們現在有患難的時候.
或者我們有不同的阻滯當中.
主讓我們有智慧.
來察看你自己的主權.
主幫助我們.
思想到我們人生的時候.
教導我們.
教導我們怎樣去過一個有智慧的生命.
好像詩篇裡面教導我們.
這樣祈禱.
求你指教我們怎樣數算自己的日子.
好叫我們得著智慧的心.
給我們的生命當中.
我們全心的忍耐.
學習你都等候你.
也都.
逃避那些愚蠢的試探.
叫我們不要.
去到居心驕傲.
叫我們去察看你自己.
有作為.
信任你.
繼續去跟隨你.
我們這樣祈禱.
防止耶穌基督的名字.
希望下個禮拜能夠見到你.
下個星期能夠見到你.
法定人數不足.
\newpage



\section{}
\label{sec:8I6GZqi_kx0}
\textbf{9/16/18 主日講道:我所擁有的(張麗仙傳道)}
\newline
\newline
連結: \href{https://youtube.com/watch?v=8I6GZqi_kx0}{\texttt{https://youtube.com/watch?v=8I6GZqi\_kx0}} ~~~~ 語音日期: 2019-08-15
\newline
\newline
\hyperref[sec:xXuB_PHmaoM]{\small{< < < PREV SERMON < < <}}
~
\hyperref[sec:index_chronic]{\small{[返順時目]}}
~
\hyperref[sec:index_scriptual]{\small{[返順卷目]}}
~
\hyperref[sec:sixId3mzKXA]{\small{> > > NEXT SERMON > > >}}
\newline
\newline
$^{1}$今天不說小敏的故事了,因為最精彩的都說完了.今天和你們說Ben的故事.Ben是香港人,50歲..
不過回頭看,1990年的時候,22歲,他來到美國,去了這裡最有名的芝加哥大學讀MBA..
到了25歲,他畢業了.因為他太厲害了,頭腦太靈活了,所以很快就被有名的公司聘用了..
一直都做得非常好,做金融.說一下他1998年,30到40歲這段期間,他事業非常有成,賺了很多錢..
不過我還沒說他的樣子,類似這些樣子.不過這個太年輕了,如果Ben50歲的話,可能像這個..
真的這麼帥,所以他讀書的時候就有很多女生追求他..
他也拍拍手,30歲在一個生日派對看到一位很漂亮的女生,是香港來的,也是讀MBA的..
所以一見鍾情就約會了.拍了半年拖,一晚Ben請女朋友去一家很高檔的意大利餐廳,很浪漫的意大利餐廳吃飯..
叫了一瓶很漂亮的紅酒,喝一喝,有侍應推著一大車玫瑰花來到這張桌子..
有幾個穿著小提琴的,拉著小提琴的,大家知道為什麼了..
然後Ben在西裝拿出一個盒子,一打開就發光了,大家知道是什麼了..
然後就跪在地上,女朋友就……大家都知道接下來會怎樣了..
所以五年之間Ben結婚,生了一個小孩,一對雙胞胎,很可愛,還有一個兒子,三個..
所以一家五口富裕家庭,外人看起來很羨慕,但也有暗湧的..
因為Ben應酬很多,所以子女經常見不到他..
太太睡著睡著,半夜也不見他在身邊..
再說Ben到了45歲,子女開始變青年了,大家都知道了,沒什麼可以說的了..
太太經常說話,但Ben又不想聽..
說不說也好,反正Ben生活很豐富,因為他雄心壯志..
當時他公司想請他做CEO,但他賺錢不夠,所以他自己出來做生意,做物流,Logistics..
在芝加哥的單攤旁邊買了一間舊舊的辦公室,買了一棟來做辦公室..
大家也知道2013年,早五年左右,美國經濟也不是很好..
但他做物流就好了,因為中國經濟開放,所以打通中國的那條線,Ben真的有很多東西賺..
他生意做得非常好,有幾十個手下.再過兩年,特朗普上台後,還將很多規矩減少了,所以他的生意更加好..
所以真的多到做不完..
他想想,要搬辦公室了,不過不好了,芝加哥的地點這麼好,不如把舊的辦公室拆掉,建一個新的..
再請多些人,所以在一兩年之間,又做了這件事..
而且新的辦公室建了之後,生意非常好,新公司建起來了..
BEN還拿了很多錢投資房子,所以錢真的財源滾來滾去,擋也擋不住..
今年八月,適逢Ben五十歲,所以他就舉辦了一個慶祝會..
請很多人恭喜他,恭喜之餘,Ben很開心,他自己就閃到露台,左手一個雪茄,右手一個布蘭地,放下摸摸自己的肚腩..
自言自語說,我真的挺令人羨慕的,一個美麗的老婆,一對可愛的雙胞胎,一個可愛的兒子,生意投資順利,財源滾來滾去,.
生意不如交給手下,我不用做,不如天天吃盡山珍海味,每個月環遊世界,哈哈哈哈,很厲害..
誰知道,突然間Ben的心狀狀變得很痛,人生什麼都能預料,但這一切他是預料不到的,所以突然間就這樣了..
大家知道了吧?大家睜大眼睛看著我,這個故事是不是真的?.
不過我還沒回答你之前,我想問你,你聽這個故事的時候,你想想聖經裡面耶穌說過什麼比喻最像這個故事呢?.
司見說的,財主的比喻,在路加福音第十二章,財主的比喻,不如我們看看,我們一起讀吧..
字很小,看到就讀吧..
有一個財主田產豐盛,自己心裡思想說,我的出產沒有地方收藏,怎麼辦呢?.
又說,我要怎樣辦?要把我的創房拆了,令它更大的,在那裡好收藏我一切的糧食和財物,然後對我的靈魂說,.

$^{41}$靈魂啊,你有很多藏的積存,可作多年的費用,只管安安逸逸地吃喝快樂吧..
神對他說,無知的人啊,今夜必要你的靈魂,你所預備的要歸誰呢?.
真的跟Ben的故事有點像,如果我們看這個財主的話,我們可以用四句話來說他..
第一,他頭腦很靈活,財產多到不夠放置,他拆了再建一個,不用搬,但又可以建一個大的創房,擺放新的財寶,新的收成,這是財主的比喻..
他成功富有,這裡說得很清楚..
第三,他知滿自得,我們看得很清楚,當他很有錢的時候,他對自己的靈魂說,你有很多積存,可作多年的費用,只管安安逸逸地吃喝快樂吧..
知滿自得..
最後一個就是,一鳴護符..
Ben的故事不過是現代家園家造版的財主的故事,純屬虛構,但內容並非憑空道讚,當然作者就是我..
財主的故事,如果你們讀英文,你們就記得是甚麼?.
Parable of Rich Fool,是愚蠢財主的故事..
這個愚蠢財主是神對有錢人的境界,究竟這個財主的比喻是神對有錢人的境界呢?.
所以我們這些不是財主的,可以在旁邊沾沾自喜地看show..
究竟財主的比喻是不是真的有錢人的問題呢?.
要回答這個問題,我們要細心偵查一下這段聖經的內容,有沒有給我們答案..
一開始我們看到財主說,我的出產,我要這樣做,我的倉房,好收藏,我的一切,你覺得這個人怎樣?.
很自大,他心目中有自己這個成功的人,但完全沒有半點感恩的心..
可以說在這個人的人生字典裡面是沒有神的..
他有自己,我的出產,我要這樣做,我的倉房,好收藏,我的一切..
再看下去,他自滿自得的時候,他對自己的靈魂說,你有很多的積藏,可以作多年的費用,只管安安逸逸,吃喝快樂..
這個人不單覺得什麼都是他自己,我所擁有的是我的,而且他覺得我想怎樣做就怎樣做..
可以說他就是他自己的神,他的結局是怎樣呢?.
我們知道如果神是真的話,不是你說沒有神就沒有神,結局就是神出來了,神對財主說,無知的人啊,今夜我要你的靈魂,你所預備的要歸誰呢?.
我們想像一下,財主慶祝的時候,他一生的財源滾滾,他有很多的慶祝,玩的時候是彩色繽紛的,但是突然間變成一片漆黑..
我們可以想像一下,他的派對的狂歡聲,喝酒碰杯聲,突然間停頓..
當神出來說要收回一切的時候,一切都停頓.財主在這個時候突然間覺悟到,原來我所擁有的不是我的,原來拿走就拿走了,不過太遲了..
說到這裡,耶穌給了一個輕描淡寫的評論,他說,凡為自己積財,在神面前卻不負足的也是這人..
口語化就是說,凡是那些富有而目中無神的人,他的結局就是這樣,一名戶父卻無永恆的福分..
回到剛才的問題,這個財主的比喻是不是只講給有錢人聽的呢?.
如果我們真的要找到答案的話,我們就要像看電影一樣,倒敘一下,看回頭..
我們看兩個倒敘,第一個倒敘就是,耶穌講這個比喻的源由是什麼呢?.
我們看一看,原來在財主的比喻之前有幾句是我們不可以忽視的..
我們看第十三節,眾人當中有一個人對耶穌說,夫子,請你吩咐我的兄長和我分開家業..
這個人來到耶穌面前,耶穌說,你幫我和我的兄長分開家業,我們很難搞定,我們家產都咬得臉紅臉綠,你們幫我們分開家業..
耶穌的反應有兩個,第一個,他說,你這個人,哪一個納我做你段氏的法官,讓你分家業的?.
這句話是什麼意思呢?.
如果我們想,耶穌當時在這個世界上的時候,他很努力很忙碌很盡心地傳講天國的道理,叫罪人悔改,又醫病趕鬼,彰顯天國的榮耀,耶穌正在做這些事情..
而這個弟兄,他好像不是,他好像是明白了,他來到耶穌面前,耶穌很忙碌的時候,他說,耶穌,你幫我分家產..
所以耶穌跟他講了這句話,第二個反應才厲害,耶穌說,於是對眾人說,你們要謹慎自守,免對一切的貪心,因為人的生命不在乎家道豐富..
你聽下去,如果你叫耶穌幫你分家產,耶穌這樣跟你說,我沒空幫你做這件事,然後跟別人說,貪,不要貪,那不就是在說他,好像很狠..

$^{81}$為什麼這個人求耶穌分家產,耶穌不幫他不要緊,還要這樣作弄他,說他貪呢?.
其實我們不知道這個仁兄是不是真的貪,因為聖經沒有說,但是我們知道的就是,耶穌說完貪之後,然後就說這個財主的比喻..
可以說四個字,就是耶穌要藉著這個弟兄求分家產的事,借題發揮,說一件事,就是貪..
然後他就說了財主的故事..
我們知道,我們一般來說,貪是拿了別人的東西,那就貪,拿了屬於別人的東西,一般來說,貪是對人的..
但是耶穌在這裡很明顯地給了一個新的意思,他說,貪是拿了不屬於自己的東西,是屬於神的東西,是對神的..
所以他就用了財主的比喻,用這個財主,他很多的財富,自滿自得的時候,他認為我所擁有的,全部是我的..
我喜歡怎樣就怎樣,這個人生態度,耶穌要用這個人生態度來點名,各位,你要小心,這個就是貪..
在神面前,這個就是貪..
我們不要忘記,耶穌當時是從那個分家產的弟兄轉面向眾人說,這個轉面是很重要的..
換句話說,耶穌說的貪不是警戒那個弟兄這麼簡單,而是要警戒眾人,當時聽他講道,跟隨他每一個人,.
還有今天我們坐在這裡,看著這段聖經的人,你們要謹慎自肅,免去一切的貪心..
耶穌所講的財主的比喻,這個警戒,對財主的警戒,不是只給有錢人,而是給我們每一個人..
所以我們看這個比喻的時候,我們絕對不能抽空來看,只說有一個人很有錢,我們要看回前幾句,為什麼耶穌講這個比喻?.
我們有了這個倒序之後,我們就明白,耶穌要警戒我們每一個人,我們要小心我們的人生態度是怎樣..
我們每個人都擁有很多東西,你要知道你擁有的,你的字典裡面有沒有神,你有沒有敬畏神,這個是耶穌在這裡所講的..
這是第一個對象,但是我們不要看到這段聖經,看到這裡就完..
因為當我們看回《路加福音》的時候,在弟兄求婚生家這件事之前,耶穌所面對的是什麼呢?.
《路加福音》裡面講了很多很多東西,就是一直都在講當耶穌傳講天國的道理,很受到人擁戴,很多人去聽他所講的時候,.
但是有另外一班人,他們是很針對耶穌,這班是嫉妒的宗教領袖,因為耶穌這麼受歡迎,所以他們是很嫉妒,.
而且是多了一些人想試探耶穌,想置耶穌於死地,想做一些事情,想試探耶穌,下他陷阱,置他於死地..
這個時候,這個氣氛是相當緊張的,所以我們要看回這個氣氛,.
以致到耶穌在講財主的比喻這件事之前,他曾經對他身邊的人講了這句話,那些殺身體之後不能做什麼,.
另外一條聖經說,那些殺身體不能殺靈魂的,你不要怕他們,耶穌對跟隨他們的門徒說,因為耶穌已經聞到宗教領袖對他的殺機,.
然後耶穌對跟隨的人說,人帶你到會堂,並官府有權柄的人面前,你不用思考如何分數,講什麼,聖靈會教你的..
大家看到嗎?耶穌已經面對著宗教領袖的殺機,他的門徒也面對著這個很大的威脅,耶穌要講這兩件事..
如果你看回這兩句話,你會看到當時的對抗是很緊張的..
耶穌就在這個背景下,當有錢人兄來叫他分家產的時候,他就講出這個比喻..
因此,我相信耶穌講財主的比喻,不單只是講給有錢人聽,不單只是講給跟隨他的人聽,而且也講給這些敵對他的宗教領袖聽..
這些宗教領袖表面虔誠,但他們的生活態度就是,我所擁有的是我的,我喜歡怎樣就怎樣..
耶穌警戒他們,你們這些不見不敬的人,你們這些不見不敬卻抱著權位的人,其實你們就是偷取了神的權位..
但是你要知道,神要拿走,隨時可以拿走..
耶穌對這些宗教領袖用財主的比喻警戒他們..
所以,耶穌這個財主的比喻可以這樣說,第一是警戒宗教領袖,告訴他們,不要以為你們擁有權位就以為了不起,其實一切都是屬於神的..
神要什麼時候拿走,就拿走,但是耶穌同時要提醒跟隨他的人,你們不需要害怕這些宗教領袖,因為真正擁有一切,掌管一切的是神,不是他們..
你們也不要像這些宗教領袖那樣,假冒為善,相反你們要敬畏神,你們要在一切擁有的東西上,在你們所擁有的一切要懷著尊崇神的心去使用..
如果你把神的榮耀當作你自己的,這就是貪婪..
我認為財主的比喻是這樣的意思..
弟兄姊妹,聽到耶穌講財主的比喻的時候,對你有什麼提醒呢?.
我們有沒有像當日的宗教領袖那樣,以神賜給我們的東西來抵擋神呢?.

$^{121}$我希望我們不會,但是我們有沒有像財主或Ben那樣,擁有很多東西,但漠視一切賜給我們的神呢?.
我們有沒有以為我所擁有的是我的,我想怎樣就是怎樣去過活呢?.
不要以為我們每個星期來教堂,我們就沒有這種態度,我們會有的..
所以我們真的很需要神的話語提醒我們,不如我們反省一件事..
反省一件我們現代人最容易的陷阱,就是沉溺,indulge..
有幾樣東西,先說吃吧,人人都喜歡吃,我也很喜歡吃..
美食,人人都嚮往,但是我們有沒有花太多的時間,金錢和精力,把位置放得太高,高過愛神和愛人呢?.
說說玩遊戲機,雖然我一般般喜歡玩,但是我知道很多人喜歡玩,我老公也很喜歡玩..
你玩遊戲機沒關係,但是你有沒有玩到你放下更加重要的事,或者你玩到不關心你身邊的人,你只顧著玩,玩,玩,玩,玩,玩..
或者你玩到連崇拜都回不了,重要的事沒有做,甚至喪失了你堅持前途的機會..
購物,我也很喜歡,但是你有沒有購到很多債,有沒有購到上癮呢?.
或者是應該用在一些更加重要的事上的錢,你用了來購物,因為你上癮了..
又說性了,你會不會追求性到一個地步,就是你做一些要掩著五下給別人聽的事呢?.
各位弟兄姊妹,我們要審查我們今天所擁有的東西,我們是不是懷著一個敬畏神的心去做?.
我們是不是貪,以為這些東西是我的,我喜歡怎樣就怎樣..
如果我們下次在這種陷阱中發覺,我不知道怎麼辦,我已經習慣了,我請你記住這句話,.
我所擁有的不是我的,神想我怎樣做,我就怎樣做..
我們今天未必是大富大貴,但是我們實在擁有很多,我們擁有才幹,.
我很喜歡鼓勵年輕人,我很喜歡鼓勵他們,如果你爸爸媽媽曾經供過你去讀什麼,.
例如彈琴,畫畫,打功夫,你不要浪費了它,不要隨便說一句拍拍手,.
說我不喜歡這些東西,算了,不要隨便這樣做,為什麼呢?.
因為你爸爸媽媽花了很多錢,很多時間,載你們去安排很多東西,其實他們很辛苦,付出過的,.
千萬不要隨便說不喜歡,不如好好地用它來報答你的父母和祝福別人的生命,.
這樣做的時候,其實我們就是珍惜父神給我們的才幹,我們不會埋沒不用,.
學了就不用放在這裡,相反我們好好地培養,來報答我們生命的神,.
我們用這些來祝福別人的生命,我們用這些來建立教會,我們的生命一定會更加有意思..
金錢,我們有沒有明智地使用我們有的錢呢?.
我們知道揮霍無度不是神的意思,用金錢來建造自己的生命,.
幫別人建立自己的生命,這個就是神的意思..
更加,如果神給你足夠的,你有沒有真誠慷慨地奉獻給教會?.
敬畏神,是要使神的家不缺糧,是要我們一起做,不是某些人做,.
是我們一起做,要神的家不缺糧..
我們的家庭,神賜給我們配偶,子女,父母,.
你有沒有好好地珍惜這些?.
有沒有好好地珍惜跟他們在一起的時間,好好栽種大家的感情呢?.
如果我要說的話,還有很多很多要說下去,.
但總括來說,今天耶穌說財主的比喻,這個比喻再一次提醒我們,.
其實我們每一個人都擁有很多,我們都是財主,我們都是阿Ben,.
但是我們絕對不可以像財主和阿Ben的結局,.
我們要記住耶穌以財主的比喻向我們發出的勸勉,.

$^{161}$就是我們懷著一個尊崇神的心,去使用神給我們的一切..
我們要深信一件事,我們敬畏神去過活,神必祝福我們的人生..
如果我們發覺,我們聽到這裡,能夠反省到很多自己做得不足的地方,並不緊要,.
因為我們的神是滿有恩典的神,.
不如就好像剛剛Praise Tim跟我們唱那首回家那樣,.
我們的心是要想著回到神那裡,我們做錯完全不是問題,.
我們要時時回到我們的神那裡,我們要在祂那裡建立一個祂喜歡的態度,.
一個敬虔的人生..
因為下個星期又是回家音樂報道會,鼓勵每位弟兄姊妹,.
好像我剛剛說的,我們一起參與,就算你不來也要拼命為這個活動祈禱,.
讓其他人能夠聽聞福音,能夠回轉歸神..
這就是你使用你擁有的時間,一個正確的方法..
讓我們再一次去唱這首回家,如果我們做錯,做得不好的地方,.
我們回到神那裡,我們要將人帶回神那裡..
而且再一次說的是,我們要珍惜靈修三十天,.
教會誠懇邀請大家用靈修三十天,.
繼續在神的話語裡面去建造你的人際關係.請Praise Tim..
求主你教我們不單止坐在這裡聽,.
我們回到家的時候,我們行出你的話,.
我們做一個更愛你的人,我們做一個更愛你的教會,.
讓我們一切所擁有的成為你的榮耀,成為別人的祝福..
多謝你奉求主耶穌的聖名,阿們..
我們接下來有探索班.
如果大家有空的話可以來.
\newpage



\section{}
\label{sec:sixId3mzKXA}
\textbf{12/9/18 主日講道: 女子組加油! \#5 馬利亞(陳崇基牧師)}
\newline
\newline
連結: \href{https://youtube.com/watch?v=sixId3mzKXA}{\texttt{https://youtube.com/watch?v=sixId3mzKXA}} ~~~~ 語音日期: 2019-08-18
\newline
\newline
\hyperref[sec:8I6GZqi_kx0]{\small{< < < PREV SERMON < < <}}
~
\hyperref[sec:index_chronic]{\small{[返順時目]}}
~
\hyperref[sec:index_scriptual]{\small{[返順卷目]}}
~
\hyperref[sec:AG5ZPTEPotw]{\small{> > > NEXT SERMON > > >}}
\newline
\newline
$^{1}$聖誕節一定要講一下聖誕的故事.
我們做這行的習慣了.
每逢聖誕節什麼節期我們都找經文來講.
正在講女子組加油這個系列.
我們一定要講一個女性.
就是瑪麗亞.
Mary這個女性是一個很關鍵的角色.
如果你去歐洲旅遊或者去以色列旅遊.
你會看到很多很多的教堂或者景點.
都是跟Mary有關的.
是不是?.
其實你開車去Bridgeport.
到處去.
很多古教堂,舊教堂都是跟Mary有關的.
聖經裡面的Mary是怎樣的呢?.
我們作為基督徒.
我們叫做Protestant.
我們是在Reformation開始的.
馬丁路德等等.
有時候看到Mary的時候就不太自在.
因為我們覺得天主教.
為什麼講Mary講得這麼過火.
或者將她的尊崇放得太高.
所以很多時候我們有些反彈.
隻字不提.
其實是不需要的.
我想告訴你.
在聖經裡面記載的autobiographical data.
就是每一個人物.
你當今天是一個聖經人物的研究.
看看不同的聖經人物,不同的角色.
到底有什麼值得我們提醒的呢?.
這樣的autobiographical data.
就是一個自傳的或者個人的資料.
最多的是誰?.
當然是耶穌.
你要知道耶穌是最多的.
第二個呢?.
大衛.
大衛是最多的.

$^{41}$你聽一下我們的YouTube program就知道了.
保羅第三個就是.
第四個可能是彼得.
剛才那些全部都是男人.
第五個應該是瑪利亞.
對他的描述或者對他的心境裡面的經歷等等.
瑪利亞是很多的.
那麼.
你要知道教會自耶穌基督升天作王.
由耶路撒冷開始到現在大概2000年左右.
有一個古代的信經.
一個文件.
叫做Apostles' Creed.
我知道最近的新聞Apostles' Creed經常上頭條.
因為我們的總統就翹起雙手.
不讀.
過去這麼多屆總統都會按照傳統讀.
他又不讀.
所以就上了頭條.
Apostles' Creed是什麼呢?.
其實是一個早期的教會.
算是最初代的教會留下來的信仰的總結.
這些總結是用來防止異端.
讓我們能夠信的東西是清楚正統的.
那麼這些信經就是給初期教會在崇拜裡面使用的.
或者一個人信了之後.
作為一個教導.
或者一個人去受洗之前.
他需要明白Apostles' Creed.
或者這些不同的初代的信經.
這些初期神學的總結.
有一句是講到瑪利亞的.
那你可想而知瑪利亞的角色是多麼的重要.
第二段講到我信的時候.
他說我信耶穌基督上帝的獨生兒子.
我們的主因著聖靈成人.
從同類瑪利亞所生.
特別要提到瑪利亞同精類懷孕.
為什麼這個教義這麼重要呢?.
神學上可以這樣推理.

$^{81}$如果耶穌要去頂替我們的罪.
他當然要成為一個完全的人.
當然要有母親所生.
但同時他也不能夠本身是罪惡之身.
我們都是罪人.
我們需要被別人救.
需要被別人贖罪.
但耶穌就不能夠這樣.
所以所謂原罪.
就不可以在耶穌的生命裡面出現.
所以他是沒有錯百分之一百的人.
但同時也不能夠有我們的罪性.
原罪是什麼意思呢?.
就是那個罪惡的狀況裡面.
我們生活出生的.
但耶穌卻不是這樣.
所以為什麼這個同精類瑪利亞的教義.
是一個很重要的教義.
同時聖經也告訴我們.
每逢有奇蹟的誕生.
每逢有神蹟出生的嬰孩出現的時候.
無論是亞伯拉罕的兒子.
以撒等等不同的奇蹟誕生出現的時候.
聖經裡面告訴我們.
只有上帝才能夠成事.
不是人能夠做.
是上帝藉著他的恩典和權能去成就這些事.
所以為什麼這個同精類瑪利亞的教義這麼重要.
當然翻開聖經.
看瑪利亞的故事或記載的時候.
他就不是抽象地講這些神學理論.
主要是講門徒.
一個跟隨耶穌基督的人.
或者跟天國有關的女子.
她到底有什麼經歷.
是一個有血有肉的歷史故事.
所以今天結束這個女子組的系列的時候.
用瑪利亞來結束.
抽幾段瑪利亞對於瑪利亞的描述.
到底對我們有什麼鼓勵.

$^{121}$這個人很重要.
這個人的信心歷程.
也是我們每一個基督徒.
需要被提醒的不同地方.
今天我們看聖經的時候.
有一樣東西大家一定要做.
看聖經.
就是一種神學的想像.
Illogical imagination.
不是亂想的,不是幻想.
而是有了聖經的資料.
有了背景的資料.
再加上一些我們可能從經歷.
或者人性裡面的.
一些的illogical imagination.
就想一想到底他經歷的是什麼.
所以這幾樣東西.
我們今天在這幾個字裡面一起做.
瑪利亞的記載有好幾處.
第一處當然是.
盧家福音第一章.
這個女子出現的時候.
第26節.
他說到了第六個月.
天使加伯列奉上帝的差遣.
往加利尼的一座城去.
這個城名叫做拿撒勒.
是一個很小很小的地方.
按統計當時應該大概幾百人最多的人口.
一個很小很小的town.
到一個童女那裡.
是已經許配給大衛加的一個人名叫約瑟.
童女的名字叫做瑪利亞.
天使進去對她說.
蒙大恩的女子我問你安.
主和你同在了.
瑪利亞因為這話就很驚慌.
反覆思想.
這樣問安是什麼意思.
天使對她說.

$^{161}$瑪利亞不要怕.
你在上帝面前已經蒙恩了.
你要懷孕生子.
吸他在名叫耶穌.
他要為大.
稱為至高者的兒子.
主上帝要把他祖大位的位給他.
他要作亞國家的皇席到永遠.
他的國也沒有窮盡.
這個所謂天使報喜.
就是annunciation.
通常呢.
人有新人.
就是有媽媽.
他有了.
我們應該會說什麼.
恭喜.
希望是這樣吧.
應該是這樣吧.
他有嬰兒.
有恭喜.
所以叫做報喜.
但如果我們想一想當時的父權文化.
會不會就是這樣呢.
當時瑪利亞已經訂婚了.
或者所謂許配了給約瑟.
當時猶太人的規矩.
訂婚許配已經是正式的了.
不能chicken out了.
只不過是還沒有正式出嫁.
還有一年就會正式出嫁.
如果這個時候有新人.
在一個父權的文化.
重男輕女的文化裡面.
會有多大的衝擊.
會有多大的羞辱.
那個壓力會有多大呢.
問一問我們中間的男人.
如果你是約瑟.
你的未婚妻懷孕了.

$^{201}$你會有什麼反應呢.
以往我們一直以為.
這個故事報喜之後.
約瑟就有一個掙扎.
他的抉擇就是可以迎娶她.
按照同樣計劃迎娶她.
或者另一個選擇就是離婚.
所以馬太福記載這件事的時候.
約瑟是一個異人.
他不想明明羞辱瑪麗亞.
所以他就想靜靜的離婚.
不想明明羞辱她.
你可以想想.
這對男女報喜之後.
到底他們的對話會是怎樣.
大家聊天的時候會是怎樣.
特別是在一個父權社會的陰影下.
男人有什麼抉擇.
我們以往一直以為.
唯一一個給約瑟的抉擇就是離婚.
但其實也可以有其他的抉擇.
按照一些文化的平衡.
我們知道他的抉擇.
可以是一件事叫做honor killing.
大家可能聽過.
現在在一些很保守的社會裡.
都有這樣的事情發生.
2014年華盛頓郵報有一個這樣的報導.
在一些巴勒斯坦的地區.
一些保守的地區.
有越來越多的honor killings.
什麼叫做honor killings.
你要想一想.
一個大男人的文化裡.
如果家裡有一個女子.
被侵犯,被污辱.
或者有姦淫等等.
在家族裡的男人.
需要去維護家族的光榮.
很多時候會將被羞辱的.

$^{241}$或者是shame了我們family name的人.
姐妹.
將他們殺死.
這件新聞發生什麼事呢.
就是一個家族裡的女子.
她跟一個家族不批准她跟他一起的男人.
她走了去找那個男人.
那個兄弟就去那個男人的家.
他就拍門.
他說給我妹妹.
我們不會對她做什麼.
你儘管給她帶回家.
他就將妹妹帶回家.
他們就在食物裡下毒.
然後就將那個女子毒死了.
這個叫做honor killings.
當然這個聖經不批准.
這是大男人的問題.
所以你想想.
收到這個信息的時候.
報喜的時候.
是不是純粹我們今天在Facebook說.
恭喜.
這麼簡單.
是非常複雜的.
瑪利亞也非常奇怪.
一章第34節.
瑪利亞對天使說.
我都沒有出嫁.
怎麼會有這事呢.
天使回答說.
聖靈要臨在你的身上.
至高者能力要任避.
你要保護你.
這個時候你就需要protection.
因此你所生的聖者必稱為上帝的兒子.
況且你的親戚以利沙白.
是一個老年的親戚.
在年老的時候也懷了男胎.
是數來稱為畢生育的.

$^{281}$現在有人六個月了.
因為出於上帝的話.
沒有一句不大能力.
瑪利亞說.
我是主的使者.
請願照你的話.
成就在我身上.
天使就離開他.
走了.
請願照你的話.
成就在我身上.
你想想.
在這種社會的背景.
你想像一下.
theological imagination.
他受的壓力會是怎樣.
而說出這句話.
做出這樣的決定.
他的faith是怎樣的faith.
是一個很risk taking的faith.
是一個很不簡單的決定.
在眾人眼中是看為羞恥的.
但他願意去聽上帝的話.
他沒有這樣的一個欠.
你想想.
生命可能會有危險.
by the way.
瑪利亞多少歲.
這個也需要我們一點imagination.
當時的少女出嫁.
大概十二三歲就會出嫁.
不是我們今天在等.
在等.
先買一間房子.
先買一輛好車.
不是這樣的.
以前是沒有選擇的.
十三四歲她就會出嫁.
你想想.
這種情況.

$^{321}$一個junior high student.
十多歲的一個猶太人女子.
她可以有一個這樣的信心表現.
告訴我們.
信心.
對上帝的信仰.
有時候真的要take risk.
信心的相反是什麼?.
不信.
有時候我們覺得信心相反就是問題.
我們有doubt.
有掙扎.
有很多疑問.
有很多困難.
這裡不是信心相反.
信心相反是什麼?.
是certainty.
是我們有把握.
所以人們不信主.
沒有把握.
先給我看天堂.
我只需要有100\%的certainty.
這就是和信心相反.
瑪利亞就是在這種uncertainty裡面.
去投靠上帝.
當然我不是說.
someone跳過湧字就可以.
最近有篇大熱的新聞.
大家聽過了.
John Aaron Chau.
一個美國的宣教士.
他去了印度附近的一個孤島.
擅自闖入孤島.
然後去那裡宣教.
然後被土人殺死.
你有聽過吧.
當然我沒有看完所有的東西.
但我想問的問題就是.
這樣做會不會只有勇字.
而缺乏常識呢?.

$^{361}$common sense.
但無論你怎麼說.
信心的底線都是要冒險的.
這就是瑪利亞的信心的典範.
一個中學生.
在我們中間沒有一個是中學生.
沒有一個是十二三歲的.
外面有的.
他有這樣的faith.
是不簡單的.
在聖經的描述裡面.
瑪利亞是經常去.
聽到主的話.
他就會去ponder.
心裡面反覆思想.
反覆思想不只是分析思考這麼簡單.
也不是安安靜靜的去默想.
不是這樣的意思.
是一種掙扎.
是一種struggle.
是很多晚上可能睡不著的.
就是struggle到底上帝的計劃是怎樣.
如果你這樣看的時候.
你會再看瑪利亞的故事.
是一個很浪漫的故事.
聖誕卡那一隻.
是一個很大的struggle.
很大的掙扎裡面.
他要去完成上帝的計劃.
所以說到瑪利亞的信心.
也說到瑪利亞有的掙扎.
在掙扎的當中.
他自己也在信心的旅程裡面.
有跌有碰.
有成長.
第二段的聖經說到瑪利亞.
就是耶穌已經成人了.
他出來傳道.
馬可福音也有這樣的記載.
耶穌進了一間房屋.

$^{401}$眾人就聚集了.
甚至他連飯也顧不得吃.
吃飯也忘記了.
21節.
耶穌的親屬聽見.
當然一定是瑪利亞在裡面.
我想路約瑟已經過世了.
聖經沒有怎麼說他.
耶穌的親屬聽見這件事的時候.
他就出來拉住他.
或者應該這麼翻譯.
就是想強拉住他.
因為他們說他癲狂了.
為上帝做事做到癲了.
31節當下.
耶穌的母親和弟兄就站在屋外.
打發人來叫他.
有很多人在耶穌周圍坐著的時候.
他們就告訴他.
看啊,你的母親和弟兄在外面找你.
耶穌回答.
媽媽來找你.
耶穌回答說.
誰是我的母親.
誰是我的弟兄.
他說看啊,我的母親.
我的弟兄.
凡遵行上帝旨意的人.
就是我的弟兄姊妹和母親.
你怎麼解讀聖經也好.
也有rebuke,責備的成分.
就是告訴瑪利亞.
告訴媽媽.
到底什麼才是最重要.
是上帝天國的東西.
還是在地上的家庭的東西.
你又要做一點想像.
當時他出來加勒里傳道的時候.
耶穌的傳道不是這麼浪漫的.
他引起很多人不滿.

$^{441}$特別是什麼人不滿呢.
特別是社會裡面的權貴不滿.
就是耶路撒冷的文士和法利才人.
所以古時裡面是一個很重honor and shame的society.
一個社會著重榮辱的文化.
特別是長子.
長子的責任.
天經地義的東西.
就是要保住家裡的明星.
跟耶路撒冷的權貴有衝突.
就有損家族的明星.
長子應該要照顧家庭.
違反孝道的話.
這是不能要的.
你這樣想像的時候.
你就發現.
這段聖經到底說什麼.
那個張力在哪裡.
耶穌就告訴瑪利亞.
放優先次序.
是神的過.
間接改正了他的家人.
包括他母親的錯誤.
是神國最重要.
不是地上的東西.
你想想.
我在教會很多年.
做父母的.
自己的子女去宣教.
自己的子女要去奉獻做傳道.
我教會這麼多年.
十該其九都是父母反對的.
很少父母會舉腳贊成.
很多事情要去攔阻.
特別是教會的執事.
那些領袖.
他們可能知道做牧師不是很好.
所以他們大部分都反對.
誰是我的母親誰是我的弟兄.
凡遵行上帝旨意的.

$^{481}$就是我的弟兄姊妹就是我的母親.
天國裡面的家庭才是最重要.
不是地上的家庭.
不是叫我們不孝.
而是告訴我們盡了能力.
是時候你離開了.
天國裡面的事.
永恆的事.
是更加重要.
這就是瑪利亞有的掙扎.
有的跌碰.
他需要學的東西.
是他生的兒子.
但是這個兒子.
是為上帝工作的.
主要不是為自己工作.
第三.
瑪利亞和耶穌的關係.
在他的信心歷程裡面.
他也受過很多苦.
這裡需要多一點想像力.
而有一個景點是每個人都去的.
就是中世紀有一群人.
想像了中世紀的一個苦路.
苦傷路.
Villa Dolorosa.
這個神學的想像是什麼呢?.
就是耶穌在他被判釘十字架.
到他在刑場上被釘死的中間.
那個過程到底是怎樣.
Villa Dolorosa就有14個站.
14個地點.
第四個地點.
在這些想像當中.
這些想像裡面.
就是會想.
一定會有媽媽在.
況且約翰福林也有說.
在十架面前.
瑪利亞在場.

$^{521}$所以怎會不在苦路上呢?.
你想想你的兒子.
你帶他去打針.
你也痛得不得了.
我記得我第一次帶我弟弟去.
兩個都是.
他要打針的時候.
他哭了.
我哭還可以吧.
我弟弟受痛苦的時候.
你多麼的肉緊.
對吧.
這就是theological imagination.
你看到你的兒子這樣受苦的時候.
你會不會不在場呢?.
你自己心裡面會是怎樣呢?.
這第四個站.
就是講到瑪利亞.
在苦路上遇見背著十字架的耶穌.
是一種很人性的想像.
一種想像.
一種interpretation.
教堂裡面有很多畫.
很多畫像.
或者經文.
是講到瑪利亞遇見她的兒子.
知道她的兒子快要.
上刑場.
死在十字架上.
這種痛苦.
你要去想像才能夠經歷得到.
在天主教的傳統裡面.
另一個神學的想像也說到.
瑪利亞的七個sorrow.
七個悲傷.
由耶穌出生.
到她離開世界.
這七個是什麼呢?.
第一在聖殿裡面.
還在手抱的時候.

$^{561}$抱著耶穌的時候.
有一個老人家叫西冕.
他講了一句寓言.
就是說.
兒子會成為彌賽亞.
但最後他加了一句.
西冕和瑪利亞說什麼呢?.
你的心會被刀刺頭.
將來這個兒子會令你受很多苦難.
心就像被刀刺頭一樣.
然後希律王來追殺他.
他要逃亡去埃及.
這個也是需要想像的.
拿東西走.
然後在聖殿裡面.
十二歲的耶穌失蹤了.
心裡面很擔心.
你有沒有試過子女不見了.
有沒有.
這麼好嗎?.
我試過一次.
我們去Marshall買東西.
買完之後忘了兒子在哪裡.
原來那個PA就響了.
那個喇叭就響了.
你們有個小孩在門口.
原來兒子跑去那裡.
不見了兒子.
你想一下多麼困擾.
到找到兒子的時候.
兒子還指著他兩句.
我不是應該要在父家裡面.
突然間就一聲叮.
這個兒子有點特別.
很不同.
在苦路上相遇.
第四個sorrow.
到耶穌被釘十字架而死.
到他身體被除下的時候.
他就想起.

$^{601}$他曾經有一個兒子.
他身體被除下的時候.
我們說白頭人送黑頭人.
是一個多大的悲傷.
到最後被賣了.
這個就是我們有的interpretation.
如果我們說.
信心是需要冒險.
信心是需要去掙扎的話.
第三就是告訴我們.
我們的信心.
我們的信仰歷程裡面.
suffering is part of the package.
我們是需要受苦的.
是會痛的.
無論是上帝拿走你生命裡面.
一些你很痛愛的東西.
或者.
有些困難.
你遭遇一些低谷.
it's part of the package.
sorrow is part of the package.
只不過我們在現代的生活裡面.
我們經常都覺得.
我們是有entitlement.
我們是鳳子.
信了主他就要給我飛黃騰達.
給我步步青雲.
是不是?.
不是這樣的.
suffering is part of the package.
到最後瑪利亞.
我們看到他的事奉.
新的瑪利亞出現的時候.
是什麼時候呢?.
史唐人傳第一章第十四節.
聖靈降臨的時候.
這個五旬節.
120個人來到聚會.
在當中有什麼人呢?.

$^{641}$他說.
這些人同著幾個婦女.
即是十二個門徒.
和初期教會的領袖.
和耶穌的母親瑪利亞.
並耶穌的弟兄.
都同心合意的來行設禱告.
等候聖靈降臨.
特別提到瑪利亞的時候.
其實是和路加福音.
頭兩章是呼應的.
因為路加寫第二集史唐人傳的時候.
就拿出瑪利亞這個人物.
換句話說.
瑪利亞和這些婦女.
並非初期教會的註腳.
不是一個footnote.
不是可有可無的人.
而是有特異的認證.
證明他們在教會中.
是重心的角色.
瑪利亞你留意.
整個故事.
是救恩的開始.
報喜.
願意生耶穌出來.
亦都出現於教會的開始.
史唐人傳第一章.
是愛門徒的典範.
你信了主那麼久.
有多少冒險.
有多少掙扎.
願不願受苦.
有沒有一個新的你.
願意事奉的你.
一個新的瑪利亞.
看清楚到底信心是甚麼.
我們結束的時候.
會讀.
《史道信經》.

$^{681}$讀的時候.
不要覺得這是一些很抽離頭腦上的.
statement.
proposition.
讀的時候.
要知道神在做甚麼.
神藉著女子在做甚麼.
神藉著瑪利亞在做甚麼.
她有的信心.
她有的擺相.
是一回事.
我們認信的時候.
讀《信經》.
就是我們認信這些東西.
同時到底我們的信心.
有沒有在上帝面前.
說我真的信.
如果上帝是這樣.
我願意.
照上帝的話.
和祂的旨意.
去承受.
傳統教會讀《信經》的時候.
是要起立的.
我們結束的時候一起起立.
我們一起來看.
一段《信經》.
一起來讀.
讓這些《信經》成為我們心靈裡.
信心的鼓勵.
《史道信經》.
《基督信經》.
我信上帝.
全能的父.
創造天地的主.
我信耶穌基督.
上帝的獨生子.
我們的主.
因著聖靈聖人.
從婦女瑪利亞所生.

$^{721}$在本雕彼拉多手下雨南.
被釘在十字架上.
死了 撞了 下到陰間.
第三天從死裡復活.
後升天.
在在無所不能的父上帝的右邊.
將來要從那裡降臨.
審判活人死人.
我信聖靈.
一聖基督教會.
聖徒相同.
聚得舍滅.
欲生復活.
並且永生.
我們誦讀.
上主我們信.
我們信上帝無所不能的上帝.
他願意拯救我們.
在恩典當中.
他的獨生子.
藉著童貞女瑪利亞降生.
無罪之身.
來成為.
一個.
罪惡的頂題.
成為罪惡的高人.
去贖我們.
去赦免我們.
去拯救我們.
主讓我們的信心.
不單用我們的口.
來說我信.
也讓我們.
為著你的國.
冒險.
去冒險.
去嘗試.
去掙扎.
甚至去受苦.
主讓我們.

$^{761}$用瑪利亞的信心來事奉你.
由瑪利亞這個女子.
她的信仰歷程.
她的典範.
成為我們眾人的鼓勵.
讓我們像瑪利亞跟你說.
我們是.
願你的旨意.
按照你的旨意來行.
在我們身上.
願以你自己的天國.
臨到.
我們這樣祈禱奉祝耶穌基督的名字.
各位請做我們.
蜜桃散會.
各位請做我們蜜桃散會.
\newpage



\section{詩篇}
\label{sec:AG5ZPTEPotw}
\textbf{詩篇人生 - 上行之詩 - 詩篇 120-134 - 全篇講道 (陳崇基牧師)}
\newline
\newline
連結: \href{https://youtube.com/watch?v=AG5ZPTEPotw}{\texttt{https://youtube.com/watch?v=AG5ZPTEPotw}} ~~~~ 語音日期: 2019-08-27
\newline
\newline
\hyperref[sec:sixId3mzKXA]{\small{< < < PREV SERMON < < <}}
~
\hyperref[sec:index_chronic]{\small{[返順時目]}}
~
\hyperref[sec:index_scriptual]{\small{[返順卷目]}}
~
\hyperref[sec:JPzN1z7hdRY]{\small{> > > NEXT SERMON > > >}}
\newline
\newline
$^{1}$我們繼續看詩篇.
我們在人生當中有一些詩篇.
是正是告訴我們我們如何做人和如何生活.
我們先看122篇.
大衛上行之師人對我說.
我們往耶和華的殿去我就歡喜.
耶路撒冷我們的腳站在你的門內.
耶路撒冷被建造如同連絡整齊的一座城.
眾之派就是耶和華的派.
上那裡去.
按以色列的常例稱讚耶和華的名.
因為在那裡設立審判的寶座.
就是大衛家的寶座.
你們要為耶路撒冷求平安.
耶路撒冷.
愛你的人必然興旺.
願你城中平安.
願你宮內興旺.
因我的弟兄和同伴的緣故.
我要說.
願平安在你中間.
因耶和華我們上帝殿的緣故.
我要為你求福.
我不知道你們有沒有想過.
我們基督徒可以用什麼形象來代表或形容我們.
我們基督徒可以用門徒.
門徒的意思其實是學生的意思.
門生的意思.
就是說你拜一個人為師.
然後跟著他學習.
就是一個終身學習的經過.
這個就叫做門徒.
所以基督徒就叫做基督的門徒.
或者很多時候聖經裡面說我們是兄弟姐妹.
著重點這個形象.
就是說我們是一班的家人.
我們是一個家庭裡面的成員.
就是我不見到你的時候.
我會掛念你.
或者你有難的時候.

$^{41}$我們需要去關心你等等.
就是大家齊上齊落.
就是兄弟姐妹.
另外在聖經裡面也說到follower.
說我們基督徒是一班的跟隨者.
follower是什麼意思呢?.
就是說我們跟著一個不同的生活方式.
是不是?.
基督給我們的生活的習慣.
生活的價值.
生活的模式是不同的.
所以我們是followers of Jesus.
就是說我們是基督的跟隨者.
跟隨他的方式.
不是跟隨世界的方式.
我們跟隨另外一個方式.
我不知道你們有沒有想過.
基督徒也可以叫做朝聖者.
pilgrims.
我想我們現代的基督徒很少會說去朝聖的.
是不是?.
我們最多去朝聖的可能是多倫多.
去一些我們喜歡去的地方等等.
在我們的傳統裡面.
很少會說到我們去某個地方.
去朝聖.
去做pilgrim.
或者去go on一個pilgrimage.
譬如有時候我會帶隊去以色列.
你會看到有其他的宗教.
或者其他基督教的傳統.
他們有他們的朝聖的practice.
東正教或者天主教.
在這方面就會比較強一些.
就去到那個地方.
去到聖地或者一些遊覽景點的地方.
跟聖經是有關係的.
跟他們宗教信仰生命是有關係的.
他們就會特別注重去那個地方.
然後去那個地方不是單單拍照這麼簡單.

$^{81}$他們是對他的信仰或者對他的生命.
是會有幫助的.
給他生命裡面一種經歷.
這種就叫做pilgrim.
什麼叫做朝聖呢?.
朝聖就是說我們永遠都是在一個路程上.
是不會停下來的.
去某一個目的地.
這個就是每一個pilgrim會做的事情.
就是說明白我們真正的家鄉.
真正所屬的地方.
歸屬的地方.
就不是在這裡.
而我們每一個都是上路者.
朝聖就是一路上路.
由這個地方.
你會去到一個終點一個目的地.
而那個地方就是屬神的地方.
或者神自己.
聖經裡面真的有這麼說的.
特別是詩篇.
說到有15篇的詩篇.
就是說我們的生命.
或者邀請我們的生命.
做一個朝聖者.
這15篇的詩篇.
剛剛排在第120至134.
這15篇的詩篇裡面.
每一篇的詩篇.
這15篇裡面.
它的title都是一樣的.
就叫做《上行之詩》.
Song of Essence.
就是上去的意思.
有上升的意思.
就是攀上的意思.
在摩西的律法裡面.
吩咐神的子民.
他們一年有三次.
他們要去到所謂朝聖的.

$^{121}$要去到神所揀選的地方.
《心明記》第16章.
一切的藍丁.
當然是當時的社會是藍士行.
他就帶著其他的家人.
然後一年三次.
在三個節日裡面.
除孝節.
七七節.
還有鑄棒節.
除孝節大概是三四月的時候.
七七節就是festival week.
就是說除孝節之後的.
七四十九.
五十日之後.
就是五旬節.
那個時間大概是四月尾.
五月頭的時候.
鑄棒節.
九月左右.
這三個大的節日.
這些以色列人.
屬神的子民.
他們就要上去神所揀選的地方.
就是耶路撒冷.
或者石安山.
石安山.
耶路撒冷.
都是同一個地方.
大衛城也是同一個地方.
他們要去那裡朝見上帝.
並且不能夠空手來朝見.
各人要帶著自己的.
的供物.
的奉獻.
上去見耶和華自己.
這個地方就是他們的目的地.
而這些上行之師.
就是去邀請他們.
在上行的當中.

$^{161}$他們就去唱這些詩篇.
所以你要這樣看這些詩篇的時候.
你才明白到.
到底他們想怎樣.
告訴我們.
我們怎樣過我們人生.
上行的人生.
一個朝聖的人生.
我們不是這裡.
是我們的目的地.
目的地是將來.
是另一個地方.
為什麼15篇詩篇呢.
為什麼剛剛好就是15篇上行之詩.
沒有多也沒有少.
因為初期的猶太教的.
一些文件告訴我們.
原來一個朝聖的經歷.
由他自己的家鄉離開.
去到耶路撒冷.
去到耶路撒冷裡面的聖殿.
進去聖殿裡面的庭院.
中間是有15級的樓梯.
在這15級樓梯.
這些初期的文件告訴我們.
就是當時的師班.
就是尼美人 濟斯等等.
他們就會在這15級上面.
去唱詩歌.
如果你這樣看的話.
就是唱那15首的詩歌.
就是說這15首的詩歌.
是一個陪同我們.
邀請我們.
去做上行的一些詩篇.
Songs of Essence.
就是邀請你和我.
去做一個上行者.
就是說什麼叫上行者.
你可以想一下.

$^{201}$那些遊客.
你可能去過不同的地方去玩.
你去做Tourist.
Traveler.
你有幾個條件.
第一個條件.
你就是不會在那個地方逗留太久.
你也不會在那個地方逗留得太舒適.
不會太自滿.
你在那裡落地生根.
你只不過是經過Transition.
在那個地方.
你永遠Move on to另外一個地方.
這個就是一個朝聖者.
或者Traveler有的心態.
朝聖者或者Pilgrim.
就永遠不會依附在這個世界上.
或者依附在地上某一點.
因為他永遠都是On the go.
他要去某一個地方.
如果你Travel慣的話.
你就知道你的行李一定是很輕便的.
我的女兒這個星期四.
她就去升大學了.
我們就幫她預備好一些東西.
搬東西.
我剛剛就幫她買了一個很小的20元.
煮麵的Pot.
插電的.
我還特意買了一個可以自動關掉的.
就是Overheat了不會燒到別人的Dorm就沒有了.
很小的,剛剛煮到一個麵.
Travel Size.
又幫她買了一個熱水壺.
也是Travel Size的,一個Liter.
一定是輕便的.
無論我的車有多大.
或者可以裝到多少東西都好.
她都是行李有限.
或者她帶的東西也是非常有限.

$^{241}$這是一個Program,一個Traveler一定會預備的.
或者一定會有的條件.
在路上當然也有很多危險.
有很多的阻滯.
你一定要預備好.
或者有這樣的所謂吃得鹹魚抵得寒.
你上得路的話,就可能要預期有這樣的事情發生.
早幾個月,我媽媽去了歐洲玩.
她玩完之後,我就打電話給她.
我說玩得開心嗎?我看到你的照片很開心.
她說開心你的頭.
我去到的時候,行李就被人偷了.
真是麻煩.
哪裡好玩?那裡一點都不好玩.
總之不去.
早點說吧,我們去夾錢.
給你去玩.
朝聖者或者Traveler就永遠都有這樣的準備.
或者你去到某個地方,可能吃的東西不夠.
或者耗盡你的資源.
你就需要去補給等等.
之後你要繼續前進,繼續前往.
聖經確實這樣告訴我們.
我們就是朝聖者,就是客人.
希伯來書第十一章十三節.
他說,這些人都是存著信心死的.
並沒有得著所認許的.
是從遠處他們望見,且歡喜的迎接.
又承認自己在世上是寄居,是客裡.
說這樣的話,人們表明自己是找另一個家鄉.
他們若想念所離開的家鄉,還可以回去的機會.
他們卻羨慕一個更美,將來的家鄉在天上.
所以神被稱為他們的神,不以為恥.
因為他們已經給他們預備一座城.
這個就是新的耶路撒冷,將來的耶路撒冷.
就是說,每一個基督徒都是一個客裡.
每一個都是在地上暫居,寄居.
我們是在去某一個目的地.
對於舊約的人來說.
他們真正的目的地就是耶路撒冷.

$^{281}$那是他們真正的家.
所以他們很多時候都希望葬在那裡.
耶路撒冷的墓地是最貴的.
為什麼呢?因為他們覺得那是他們的家鄉.
無論他們住在歐洲也好,非洲也好,亞洲也好.
世界各地不同的地方.
他們始終都是歸向那個國家.
對於新約的人來說,我們已經知道耶穌基督已經來了.
我們的家就是將來的新耶路撒冷.
新耶路撒冷.
新天新地來到,耶穌基督再次來到的時候.
那個就是我們的家鄉.
所以上行之師就是告訴我們.
真正我們上路做一個旅客的時候.
我們應該懷著什麼心態.
以及一個怎樣的心情來上路.
所以如果看這些上行之師的時候.
這樣看是一個上行的詩歌.
你就會明白多一點.
我們今天看一點點這些詩篇.
120篇,第一篇的上行之師.
這首是哀歌.
就是說這個詩人哀嘆什麼呢?.
哀嘆他現在住的地方是遠離耶路撒冷很遠的地方.
並且他遇見很多困難,困境,危險.
第四節,他說我在急難中求告耶和華.
他就應允我.
他回憶起以往在急難中神曾經應允了他的祈禱.
好像你生活中經歷過一些事情.
上帝答應了你,應允了你的禱告.
這個詩人哀記起在急難中求告過耶和華.
然後他講了很多他遇見的困境.
第五節,他說我寄居在米切.
住在基達帳棚之中.
有禍了.
我願那恨惡和睦平安的人許久痛存.
我願和睦但我發熱的時候他們就要振前.
很奇怪的一段詩篇.
第一篇上行之師.
詩人就抱怨什麼呢?.

$^{321}$你可以想像他行經一些異鄉的時候.
正在上耶路撒冷的途中.
他經過了暫住在某一些地方.
這個地方叫做米切或者是基達.
這兩個名字都是在舊約裡面的所謂恐怖分子.
米切就在以色列很遠的北面的地方.
那裡的人很出名是很暴戾的.
基達就在以色列的東面沙漠的地方.
這些是一些幼默的民族.
是一些好戰的勇鬥狠的戰士.
他說他住在他們中間的時候.
他感覺非常不平安.
這些人就代表著一些壓迫者.
這個詩人上路的時候.
他行經這些地方有很多的衝突.
有很多的危險.
他渴望有和平有和睦.
但他找不到.
唯獨是找到和睦.
Shalom.
就是要去到耶路撒冷才找到.
但是他在這些危地當中.
在危險的地方裡面.
他去信靠上帝.
上帝救他脫離了這些險境.
所以他繼續上路的時候.
遠離耶路撒冷.
這個詩篇第120篇.
第一篇上行之詩.
第二篇.
他開始上路了.
上路的時候他做什麼呢.
第121篇.
他說我要向眾山舉目.
我的幫助從何而來.
我的幫助從造天地的耶和華而來.
他看到很多的山.
因為我經常覺得芝加哥是一個很悶的地方.
對於旅行是一個很悶的地方.
他說去Wilson Park.

$^{361}$其實每個地方都一樣.
因為四景都是平的.
沒有山沒有水.
去哪裡旅行.
或者去哪裡露營.
其實到處都是一樣.
因為是一個大平原.
我們的景色都差不多.
有山有水是以色列的景色.
以色列很多的山和谷.
有些溪澗.
有很多的山谷等等.
如果你看詩篇第23篇.
就是講述這些情景.
你想像到這些情景的時候.
這個詩人上行去耶路撒冷.
他就說向高山來舉目.
當然他的目的地是耶路撒冷.
耶路撒冷是海拔以上的2500呎的地方.
在群山當中.
耶路撒冷在群山的裡面.
所以他一定是往上望的.
上行之詩就是這個意思.
升上攀上的意思.
永遠都是上去的.
這種意境.
到現在以色列人都有這樣的概念.
因為以色列人在那些散居世界不同國地的.
以色列人猶太人.
他們如果要移民回到現在的以色列.
Modern Israel.
那個移民的過程叫什麼呢.
叫做Aliyah.
Aliyah的意思就是上去的意思.
上行的意思.
就是說他們上去耶路撒冷.
所以無論哪一個方位.
你去耶路撒冷都要爬山的.
一定是上行上去的.
當然在古晉東裡面.

$^{401}$很多時候那些山頂都是一些假神的廟宇所在的地方.
這個山是誰的神的.
這個山是屬於誰的神明的.
譬如在舊約裡面一個很出名的神明叫做巴黎.
Bel.
Bel的基地.
他的山就是Mount Savon.
在北面一個高山的那裡.
詩人說他望見這些這麼重的山的時候.
他就想起這裡有很多的假神.
在路上有很多景象.
他看到很多的山.
但是這些山的神是不能夠幫助他的.
唯獨是一個是誰呢.
就是創造天地的耶路撒冷.
他就上路去耶路撒冷.
來敬拜這位天地創造的主.
他是唯一的幫助.
上路的時候是唯一的幫助.
所以第21篇第二篇的上行詩篇.
就是他上路的時候看到的景色和他想起的禱告.
上行,昭聖的時候.
他依靠神.
神給他幫助.
就等於我們人生一路過程的時候.
我們經歷不同的迂迴曲折.
或者上山下山.
就是上帝幫助我們.
我們去依靠上帝.
再下一篇122篇.
就是剛才我們開頭的時候讀的那篇.
122篇.
他說人對我說我往耶路撒冷的殿墟我就歡喜.
耶路撒冷啊.
我不的確站在你的門內.
耶路撒冷被建造如同連絡整齊的一座城.
122篇.
第三篇就是這篇.
上行之詩.
他就到了哪裡呢?.

$^{441}$大概到了耶路撒冷.
因為他說我們的腳站在你的門內.
你可以想像有很多其他的朝聖者一路上行.
一路來到耶路撒冷的時候.
他說我們的腳.
就是我們一群人來到你的城門裡面.
是一件很歡喜快樂的事情.
因為你到達了目的地.
以前的古人.
他們上路的時候不像我們這樣.
我們可能坐飛機最多都是坐十幾個小時的飛機.
他們上路要去預備一個朝聖的行程.
很多時候是要預備幾天的時間.
要走幾個星期.
或者要幾個月.
他們才到達.
我記起有一次在巴勒斯坦地的時候.
我在West Bank.
我在一個營裡面幫一個巴勒斯坦的基督徒.
他有訪客.
剛好有一天他的訪客是一群歐洲人.
他們在歐洲步行去以色列.
那個人已經全身臭了.
沒有洗澡.
步行了幾十天.
從歐洲好像是德國.
跟一群朋友他們就背著背囊一路走.
經過土耳其.
然後經過敘利亞.
經過黎巴嫩.
然後到達以色列.
他就揭開他的護照.
他就告訴我他在哪裡丟過人.
我經過多少個地方.
多少大國小國.
他要預備很多很多東西.
他才能夠到達以色列.
對他來說是聖地.
你看到這些朝聖者終於來到耶路撒冷.
他就站在城門裡面.

$^{481}$他看到耶路撒冷是什麼呢?.
是一個很可愛的地方.
是一個連絡整齊的一座城.
今天去到耶路撒冷.
其實都是保持著近冬的風俗.
那些店鋪和房子.
全部都是很密集的.
車不能進去.
一定要走進去.
那些古城現在都是保留著那種風俗.
如果你去過歐洲玩.
大概都是這樣.
比較舊式的生活習慣.
那些街道迂迴曲折和很雜.
其實就是在形容這種情況.
連絡整齊.
其實意思就是說團結在一起.
大家聚在一起就是神的子民.
聯絡在一起.
來建立上帝的國度.
建立上帝的聖城.
第四節.
眾支派就是耶和華的支派.
十二支派.
上到耶路撒冷去.
按以色列的常例去稱讚耶和華的名.
因為在那裡設立審判的寶座.
就是大衛家的寶座.
很奇怪.
為什麼去到這裡居然說審判的寶座.
你去耶路撒冷讚美神.
當然是因為你去到聖殿讚美耶和華的名.
為什麼第五節他特別提到審判的寶座.
Judgment.
公義的意思.
在耶路撒冷設有公義的寶座.
在舊約裡面.
很多時候人去耶路撒冷.
不是純粹去敬拜不簡單.
很多時候他們在自己的鄉村.

$^{521}$他們去耶路撒冷其實是尋求公義.
特別是那些本地的法院.
那些當地的法院.
如果有些不義的事情發生.
或者他的案件沒有一個合理的賠償.
或者一個合理的審議的話.
或者他當時被當地的權貴欺壓的話.
他就會上訴.
平民百姓就會上路去耶路撒冷.
希望去耶路撒冷做什麼呢.
就是將他的冤案交給那個君王.
那個君王就是大衛家的君王.
耶路撒冷就是最後上訴的地方.
他想尋求公義.
所以在賞賀之詩裡面.
他說他上耶路撒冷是希望見到有公義的彰顯.
就是說他希望去到君王寶座審判的面前.
他能夠討回他的公道.
公義.
在舊約裡面會不會講到公義的時候呢.
Justice的時候.
通常會出現幾個人.
摩西的律法叫那些人.
你要行公義,你要好憐憫.
永遠提到公義的時候都會出現幾種人.
什麼人呢.
孤兒,寡婦,寄居的.
或者真正貧窮的人.
不是假裝貧窮的人.
是真正貧窮的人.
換句話說,在社會裡面.
這四個人都在社會裡面是沒有權勢的.
他們是低下階層的人.
他們是很脆弱的一群人.
特別是有打仗,有饑荒等等.
他們就是第一個受害的.
換句話說,今天包括什麼人呢.
例如難民.
或者一些無助的移民.
無家者.

$^{561}$Homeless.
被虐待,被打壓的人.
這些就是往往舊約裡面提到過.
公義的時候.
這些無權無勢的人.
所以大部分時間.
在舊約的聖經裡面.
提到要去行公義的時候.
大體上都是叫什麼人行公義呢.
就是社會裡面有權有勢的人去行公義.
因為行公義就是由上而下組的.
上面有權有勢的.
他就可以行公義.
以公義來代替下面的人.
有權有勢的人.
就可以幫助下面被打壓的人.
沒有資格去反抗的那群人.
所以大部分時間講到行公義.
都是那些有資源,有勢力的人.
做君王的.
那時候就是被圈滅.
要去行公義,好憐憫.
我中文是一般般的.
但是昨天我就認識了一些新的東西.
有句成語叫做.
黃台之瓜.
何堪再擇.
就是李嘉誠在《大公報》刊登的.
我找找典故的時候.
我覺得挺有趣的.
大家聽過這句話吧.
看文人才能寫得出來.
典故是講到武則天.
他打壓他的兒子.
殺他的兒子.
把他們趕盡殺絕.
兒子還沒死之前就有個絕命詩寫給他.
黃台之瓜.
就是唐代的宗室.
何堪再擇.

$^{601}$就是你再一味再擇下去的時候.
他就沒有了.
不要趕盡殺絕的意思.
這是我的解釋.
或者我找到的解釋.
我知道可能有其他人有不同的解釋.
但是無論如何.
這句話是對誰說的呢?.
是對有權有勢的人說的.
那些叫自己做媽媽的人.
就是你不要對子女趕盡殺絕.
做一點點合理的事.
行一點點憐憫,行一點點公義.
《舊約》裡面就是這樣.
每次講到行公義,好憐憫的.
都是針對什麼人呢?.
是有權的.
在《舊約》裡面.
每次稱讚軍王.
都是說他們肯去行公義.
就會稱讚他們.
因為這個是本身《舊約》裡面世界.
軍王的責任.
因為他們有權.
所以他們可以秉行公義.
《撒母耳記》下第八章.
這是第一次講到大衛的王朝.
當他們克勝所有的仇敵之後.
四景平定.
他說大衛作以色列眾人的王.
又向下面的眾民秉公行義.
Justice, exercise justice.
有權有勢的就是要去行公義.
很少會講到平民百姓要行公義.
因為我們哪裡有權有勢去做這些事.
當然是上面有權有勢的人才可以彰顯公義.
並且這裡講到上帝.
祂才是真義軍王.
而上帝就是那位愛連率.
以及去行公義.

$^{641}$為著那些社會裡面無權無勢的人.
孤兒寡婦,無家寄居等等的人.
去到行公義.
無權者,在下面的人.
只能尋求公義.
或者祈禱希望有公義.
或者希望伸冤.
在上帝面前伸冤.
他可以求助尋求公義.
但是上面做軍王的就要秉行公義.
要有一點念敏的心.
所以講到公義的時候.
講到審判的補助的時候.
就是對一些有權勢的人.
告訴他們如果對公義疏忽的話.
不單單是沒有念敏.
不單單是沒有念敏之心.
而是你干犯了上帝公義的標準.
所以有能力者.
就應該要去行公義.
為著下面比較貧苦的人去行公義.
聖經告訴我們行公義好念敏就是這個意思.
前幾個星期我朋友在網上寫了一段話.
是講到我們美國.
怎樣對非法移民.
大概他講到的情況就是這樣.
我引述給你聽.
他說他看到一些新聞.
講到一些ICE.
那些人去抓一些非法移民的家庭的父母.
拆散了他們的家庭.
拆散之餘就沒有去幫助那些小朋友.
那些小朋友來到的時候.
他們就被拆散了家庭.
當然國家安全.
國家邊界是需要的.
聖經也有講.
是需要有國家安全的關注.
但問題就是這樣.
這些人已經到了.

$^{681}$那硬被拆散的時候.
我朋友為什麼很sad呢.
為什麼他說很悲傷呢.
就是他看到一些基督徒.
講一些言論.
就說這些基督徒我們是愛上帝的.
但卻同時對這些寄居者.
說他們是小強.
說他們是入侵者.
或者用一些言語去侮辱這些人.
我覺得他說得很對.
基督是吩咐我們.
做什麼呢.
要去愛他們的靈體.
我們不可以一口說愛上帝.
但是對其他人.
對我們的靈體.
我們不關心或者落井下石.
這是國家社會最脆弱的.
當然邊界是重要的.
安全是重要的.
但是如果我們一口用這些話去形容人.
我們是不是背道而馳.
是不是行公義好憐憫.
這篇詩篇最後繼續說什麼呢.
他說要有公義才會有平安.
他說你們要為耶路撒冷求平安.
耶路撒冷愛你的人必然興旺.
願你城中平安.
願你宮中興旺.
因愛弟兄和同伴的緣故.
我要說願平安在你中間.
因耶路撒冷說我們上帝奠的緣故.
我為你求福.
這裡提到平安.
平安不是單單安居樂業.
平安不是單單無風無浪.
沒有人阻止我平時生活的日常生活.
平安不是單單沒有戰爭.
平安不是單單心裡面平靜.

$^{721}$Salom是well being.
wellness的意思就是說有公義有真理.
就是社會裡面有公義有真理.
耶路撒冷是最高法院願意去做公義的工作.
審議的工作.
還人一個的公道.
這個就叫做Salom.
這個叫平安.
上行之師為什麼上到耶路撒冷.
需要做這件事.
就是這是神國子民追求的東西.
社會裡面人的關係裡面.
是需要有Salom.
要有justice.
不只是心裡面我覺得.
我起居飲食搞定就行了.
如果你是一個上行者.
你關心的是什麼.
如果你是上行者的話.
你關心的絕對不會只是今時今日.
現在我能夠吃餐安樂茶飯.
這個不是你最終關心的東西.
你最終關心的是終點.
去到新耶路撒冷的時候.
這個神的國.
它的價值.
如果上帝的國是要行公義好憐憫的話.
今天我的起居飲食就是次要.
我今天可不可以有安樂茶飯吃.
也是次要.
因為我始終是上行的.
公義的東西我們關心多少.
或者耶路撒冷的價值我們關心多少.
神國的價值我們關心多少.
沒有權沒有制的人在我們中間大有人在.
我們是一些移民的community.
我想不少人比我們沒有那麼幸福.
我們關注多少.
還是我們只是忙於我們世上的attachment.
我有好一點的生活就是這麼簡單.

$^{761}$可以比上行早.
可以同心低流幾度.
多謝收看城會網台.
記得like comment 或者訂閱 subscribe我們.
快播的呀!.
\newpage



\section{詩篇}
\label{sec:JPzN1z7hdRY}
\textbf{8/25/19 主日講道: 詩篇人生 - 詛咒詩歌 (陳崇基牧師)}
\newline
\newline
連結: \href{https://youtube.com/watch?v=JPzN1z7hdRY}{\texttt{https://youtube.com/watch?v=JPzN1z7hdRY}} ~~~~ 語音日期: 2019-09-01
\newline
\newline
\hyperref[sec:AG5ZPTEPotw]{\small{< < < PREV SERMON < < <}}
~
\hyperref[sec:index_chronic]{\small{[返順時目]}}
~
\hyperref[sec:index_scriptual]{\small{[返順卷目]}}
~
\hyperref[sec:C3L6GM4kyjo]{\small{> > > NEXT SERMON > > >}}
\newline
\newline
$^{1}$我們繼續看詩篇,我們知道在人生當中有不同的階段,有不同的經歷,在150首詩歌裡面,總有一首是講到我們人生裡面不同的經歷,不同的情況..
所以我們會看不同的詩篇,以致詩歌能夠幫助我們,面對我們人生不同的階段..
我們今天會看詩篇第58篇,如果你有聖經的話請你現在打開,電話的聖經也好,手上拿著的聖經也好..
詩篇第58篇,這篇的詩篇有11節,我們先看這段的詩篇..
我用的是新的譯本,所以我們現在開始讀..
如果你看不到畫面,如果你還沒有聖經的話,可以跟旁邊的分享..
或者你聽一聽這段的聖經的詩人的情緒是怎樣..
詩篇第58篇第一節..
(閩南語).
\newpage



\section{詩篇}
\label{sec:C3L6GM4kyjo}
\textbf{9-1-19 主日講道: 詩篇人生 -哀歌 (陳崇基牧師)}
\newline
\newline
連結: \href{https://youtube.com/watch?v=C3L6GM4kyjo}{\texttt{https://youtube.com/watch?v=C3L6GM4kyjo}} ~~~~ 語音日期: 2019-09-08
\newline
\newline
\hyperref[sec:JPzN1z7hdRY]{\small{< < < PREV SERMON < < <}}
~
\hyperref[sec:index_chronic]{\small{[返順時目]}}
~
\hyperref[sec:index_scriptual]{\small{[返順卷目]}}
~
\hyperref[sec:EogeiLrQX3g]{\small{> > > NEXT SERMON > > >}}
\newline
\newline
$^{1}$今天我們繼續看詩篇.
我們說過詩篇裡三分之一的歌都是一些哀歌.
一些lament.
一些protest.
一些在上帝面前申訴的詩詞.
不是我們心裡想.
來到教會,來到神面前一定是很rosy(浪漫)的.
一定是很浪漫的.
一定是很正面的.
很多時候這些詩篇告訴我們一個有信心.
對上帝有信仰的人.
很多時候就會在上帝面前作哀歌.
以及申訴.
我們也告訴大家.
在這些哀歌當中.
既然佔了三四十個百分比都是這一類的禱告.
詩篇或者聖經告訴我們.
我們如何處理我們心裡面一些負面的情緒.
很多時候我們華人的文化很喜歡將這些東西壓抑在我們心底裡.
我們說過了.
如果我們這樣做的時候.
很多時候就會在其他地方出來.
然後傷害其他人.
或者我們會讓這些情緒去hijack(拆解)我們.
騎劫了我們的生命.
以致我們做一些事情.
可能會隨意發洩.
或者會傷害自己.
或者我們會孤立我們自己等等.
似乎聖經裡面的詩篇這些哀歌.
這些禱文.
就是給我們有一條中間的路.
就是藉著去祈禱.
我們能夠來到上帝面前去處理.
我們心裡面看到一些比較難堪的事情.
我們今天要看詩篇22篇.
我的神我的神為什麼離棄我.
為什麼遠離不救我.
不聽我哀慳的聖的言語.
我的神啊.

$^{41}$我白日呼求.
你不應運.
夜間呼求.
並不住聖.
你看到這是一個很personal的.
很個人的一個表達.
他說我的神.
My God.
我自己的神.
和上帝有一個很個人的關係.
第三個字.
但你神是聖潔的.
是用以色列的讚美為寶座的.
我們的祖宗倚靠你.
他們倚靠你.
你便解救他們.
他們哀求你.
便蒙解救.
他們倚靠你.
就不羞愧.
這個詩人在申訴裡面.
complaint(投訴)的裡面.
他說上帝離開他.
他感覺到上帝離開了他.
上帝不在他身邊.
並且他提醒.
去挑戰上帝.
提醒他.
提醒上帝什麼呢.
他說他遠古的時候曾經做過一些事情.
叫上帝不要忘記.
每逢這些舊約的人.
他講到遠古的時候.
祖宗的時候做過的事情.
解救過他們.
拯救過他們.
通常舊約的人都會知道.
舊約裡面其中一個最重要的.
所謂拯救,得救.
就是出埃及.

$^{81}$在埃及的時候.
那些以色列人.
他們受苦.
他們被欺壓.
上帝聽見他們受苦的時候.
所有的哀求.
然後上帝下來打倒暴政,法老.
然後將以色列人救出來.
這就是他講的祖宗的時候.
我們的祖宗已靠過你.
也得到你的拯救.
記起這件事.
叫上帝記起這件事.
叫上帝現在去做他以前遠古的時候做過的事情.
他繼續去到新數的時候.
他說:但我是一條蟲,不是人.
被眾人羞辱.
被百姓藐視.
他在人間裡面.
別人看他.
他自己覺得自己是一個subhuman.
或者一個不像人的人.
凡看見我的都瘋了.
他們滅咀.
搖頭說.
他把自己交託耶和華.
耶和華說可以救他.
耶和華既喜悅他.
可以答救他吧.
十一節.
求你不要遠離我.
因為危難臨近我.
沒有人幫助我.
有許多公牛圍繞我.
巴山大力的公牛四面困住我.
這些都是當時的用語.
詩詞的用語.
那種牛不是說牠在動物園裡面.
也不是在非洲.
而是一些詩詞.

$^{121}$說一些很強力的敵人來欺壓他.
十三節.
他們向我張口.
好像在嘲笑的詩詞.
我如水被倒出來.
我的骨頭都脫了節.
我的心在我裡面.
如蠟融化.
我的精力護骨.
好像牙片.
我的舌頭貼在我的牙床上.
你將我安置在死地的塵土之中.
他就像死人一樣.
你也可以呼求上帝.
犬類圍著我.
惡黨環繞我.
他們擇戀我的手和我的腳.
馮崢說到狗的時候.
不是你養我的狗.
是說野狗.
或者是一些野狼.
我的骨頭.
我都能數過.
他們瞪著眼看我.
他們分我的內衣.
為我的內衣占溝.
看了已經十八節.
詩人說到一種情況.
是他在一個不公義的欺壓中.
到底是怎樣的呢?不知道.
當然這些詩篇很多時候都是比較.
普遍化.
以至後代的人能夠想起這些詩歌.
用這些圖文的時候.
能夠應用在他自己現在的情況.
或者你的情況也好像是這樣.
可能他在醫院的病床上.
可能這位詩人在很多壓迫的情況下.
有很多人糟蹋他等等.
但是你要知道.

$^{161}$看了十八節.
這只是這篇很長的詩篇.
二十二篇的頭一半.
他用了十八節.
他都沒有求上帝.
很多要求.
他只是告訴上帝.
他現在有的抗議.
他現在有的投訴.
他現在有的申訴.
這些哀歌背後有一個很強的意念.
那個意念就是.
上帝很在乎我們裡面每一句的投訴.
我們很多時候都會覺得上帝不理我們這些事.
如果是神的話.
天上的神.
他怎會在乎我這個人.
裡面現在經歷的是什麼.
但是這些哀歌.
這篇二十二篇就告訴我們.
他很在乎我們的抗議.
每一句話他都聽見.
詩人沒有強求或急於來到上帝面前.
求這個求那個.
要這個要那個.
他只不過在上帝面前表達他的不滿.
表達他現在的情況.
大家知道我女兒剛剛去了大學.
過幾個星期.
我半空巢.
有些人問有什麼不習慣.
我最不習慣的是什麼.
就是沒有人投訴.
少了一半的人投訴.
這麼安靜.
為什麼好像沒有人投訴.
平時很多人申訴.
電腦又太慢.
或者WIFI又太慢.
或者這個東西不好用.

$^{201}$那個東西又太舊.
又要換新的.
有沒有想過上帝很在乎.
很想聽.
我們心裡的牢訴.
我們的抗議.
或者我們人生裡面遇見的一些難題.
我們有的抗議.
或者我們在一些觀察當中.
在世界裡面你見到的新聞.
或者一些故事.
你心裡面很不舒服.
那些不舒服原來上帝.
每一句都很在乎.
有一個釋經家叫做.
Kathleen O'Connor.
她說到哀歌的時候.
她說哀歌叫做一件事.
哀歌就是要去Name it.
Name what is wrong in our world.
就要道出世間上有什麼是錯的.
在這個屬於上帝的世界裡面.
有什麼是失序.
一些很坦率的哀悼.
就是要將這些的錯處揭露出來.
去說明出來.
以致它能夠被看見.
最終我們才能夠得到醫治.
我記得有一次跟一個牧師的朋友聊天.
他有很多正能量.
他經常說我們有什麼投訴都好.
譬如人家受苦才算.
人家受苦,那我就很不舒服.
那我就去投訴.
我在上帝面前投訴.
他總是要最後一句.
起碼你不是受苦.
我就不是很頑固.
有一點格外的觀火.
詩歌或者哀歌.

$^{241}$沒錯我們是需要感恩,需要喜樂.
這些是聖靈的果子.
或者聖靈的生命裡面的工作.
但不代表沒有哀歌.
正如我所說,詩篇裡面的40\%都是哀歌.
一個有生命的人才會有哀歌.
一個屬靈的人才會有哀歌.
而這些哀歌背後表達著一個上帝.
這個上帝可能不是我們意想當中.
生活中有的上帝.
有一個天主教的神學家叫魯云.
他在一本書裡說過.
他是一個在絕望或者困惑裡面.
默然的和我們同在.
在悲傷的時候.
喪失的當中和我們一起.
可以容忍不知道.
可以容忍不確定.
沒有答案.
又不強行去醫治.
不草率下藥.
這個才是真正關懷我們的朋友.
當你心悶的時候.
當你有辛訴的時候.
對方坐在你身邊.
和你在一起.
不草率下藥.
不去強行醫治.
這個才是真正關懷我們的朋友.
神就是這樣.
這些哀歌就是告訴我們.
真正的上帝應該是這樣的上帝.
所以說了十八節之後.
他才來到神的面前.
有三節.
他去求.
他說:求你不要遠離我.
我為救主求你快來幫助我.
求你救我的靈魂脫離刀劍.
救我的生命脫離犬類.

$^{281}$救我脫離獅子的口.
你已經應允我.
使我脫離野牛的國.
最後的時候.
他才把他的請求.
帶到上帝的面前.
然後到二十二節.
才開始有轉變.
二十二節他說什麼呢.
他說:我要將你的名傳於我的弟兄.
在會中我要讚美你.
你們敬畏.
也就是說那些人要讚美他.
雅各的後裔都要榮耀他.
以色列的後裔都要懼怕他.
因為他沒有藐視恨惡受苦的人.
神也沒有向他掩面.
那受苦的人呼籲的時候.
他就垂停.
然後有一個語氣的轉換.
語氣的轉換.
哀歌很多時候都有這樣的格式.
就是哀求完之後.
那個詩人就來到上帝面前.
他有一些語氣上的轉折.
他不知道發生了什麼事.
可能是他去到聖殿敬拜的時候.
他聽到一些說話.
聽到一些歌.
或者那個牧者.
或者那個祭司.
跟他說了一些話.
以致他突然間蘇醒過來.
想清楚.
想起那些東西.
然後他有一些轉變.
他就說我要將你的名傳在我的弟兄.
在會中我要讚美你.
這些全部都是哀歌常有的格式.
那個格式就是.

$^{321}$開頭有一些申訴訴苦.
然後最後他會有信心的表達.
去讚美.
上帝可能已經找了他禱告.
或者他再一次堅定相信.
上帝一定會聽他的禱告.
你見到這些詩人在哀歌裡面.
他們不會純粹說.
上帝聽了我的禱告.
然後我就說多謝.
然後就走.
詩人會做什麼呢.
他會轉向在會中見證上帝.
在那個敬拜的會中.
他會邀請其他人.
來聽聽他的經歷.
我們的感恩.
永遠都是一種community.
社群性的一種反應.
不是我自己跟上帝說.
thank you 就完.
他會帶其他人.
進入他的生命裡面.
他會出來做見證.
鼓勵其他人.
或者可能是因為他的經歷.
可以幫到其他人.
或者他的經過可以幫到.
鼓勵其他人面對相似情景的人.
所以很多時候他們就會出來.
會在會中.
見證去讚美上帝.
我相信這些就是我們教會要做的事.
我們有過的經歷.
上帝聽過你的祈禱.
或者舒緩過你某一些的辛酥.
在會中去讚美.
以致後代其他人.
能夠更加深去體會.
到底這個上帝是怎樣的上帝.

$^{361}$不是純粹你自己收起來孤立你自己.
這種的靈性.
這種的生命.
是一種群體的生命.
是男女躲起來.
自己搞定.
他是在社會裡面去做見證.
以致那個社會.
能夠幫到.
鼓勵到其他人.
我們看了大概二十多次.
我們不會看完這麼多次.
不過.
很多時候我們會問一個問題.
我們就說未必會有下一半.
我們生命裡面很多時候未必有下一半.
你想一下你最近的投訴是什麼.
或者你曾經試過的困境.
未必真的會有下一半.
二十二節開始.
突然間反彈上來.
在谷底上來.
可能有這樣的經歷.
可能也沒有.
如果沒有的話會怎樣呢.
我相信這篇詩篇最特別的地方.
詩篇二十二篇.
在詩篇二十二篇.
在新約的聖經裡面.
是最多被引用的.
在哪裡被引用呢.
就是在四福音書裡面.
記載耶穌在十字架上受苦.
受刑罰的時候.
被打.
以致骨頭都斷了的時候.
在這種刑罰當中.
福音書的作者.
很多時候會引用這篇詩篇.
馬可福音第十五章三十三節.

$^{401}$這裡記載耶穌在十字架上發生什麼事.
在正午到下午三點鐘.
遍地坑暗了.
下午三點的時候.
耶穌大聖的呼號.
以來伊 以來伊.
拉瑪撒巴 葛大利.
其實是亞蘭文.
翻譯出來.
這句話.
馬可福音說什麼呢.
翻譯出來是什麼意思.
就是我的神 你為什麼離棄我.
這就是在引用二十二篇.
你記得第一節.
二十二篇的哀歌.
就是說我的上帝.
你為什麼離棄我.
他就是在說這件事.
他在引用.
二十二篇來祈禱.
這是我們的主耶穌在十字架上.
他到了最黑暗的時候.
受苦最深的時候.
他就是引用二十二篇.
來祈禱.
嗯.
二十二篇不是預言.
你不要覺得.
二十二篇是預言在十字架上的情況.
但是二十二篇說什麼呢.
就是在十字架上的君王.
他受苦的時候.
他就用一個這樣的祈禱.
二十二篇的祈禱.
就是居然一位天上的主.
願意成為受苦的人.
就是和我們這群受苦的人.
受難的人.
很多時候我們在世界上.

$^{441}$有很多辛酸的.
他來到和我們一起.
去認同.
這個是最大的轉折.
上帝願意成為被拋棄者.
使我們不再被拋棄.
十字架.
就是說在受苦的裡面.
原來主來到.
我們說二十二篇下一半.
很多時候我們未必親身可以經歷.
但是耶穌在十字架上經歷了.
經歷了上一半.
也經歷了下一半.
為什麼呢?因為下一半就是上帝為他申冤.
從死裡面復活.
一個二人受苦.
最後他得到公義.
死而復活.
能夠得成.
上帝為他申冤.
上帝聽他的哀歌.
最後為他伸張公義.
雖然我們人生裡未必一定會經歷到.
四篇二十二篇下一半.
但在基督裡面我們能夠分享.
我們每一個信耶穌的人.
我們就是看著十字架.
縱使耶穌在十字架上受刑罰.
是不義受鞭打.
他死,但卻復活.
上帝為他申冤.
整篇二十二篇就是他的人生.
我們未必眼見到.
二十二篇下一半.
或者未必親身經歷到.
但在基督裡面我們能夠分享.
我認識一些香港的教會.
最近有很多聯署.
這些聯署都是嘗試去.

$^{481}$把一些不公義的事情說出來.
強烈譴責香港的警方去襲擊市民.
無差別的襲擊.
市民又打.
市民又毆打.
社工又抓.
救護人員不讓進去幫忙.
不讓去救人.
坐輪椅的又打.
外國人又抓.
這就是我們的世界.
這就是詩篇二十二篇上一半.
在我們網台上有一個中國大陸的弟兄留言.
求主幫香港的弟兄姊妹保守.
在大陸的我們除了禱告之外.
做不到什麼.
這就是我們今天的世界.
公義未能夠見到可以伸張.
我記得有一次.
上兩個星期.
有其他外洲的教會來.
去探我們.
我們大家有開會.
有很多都是白人.
我記得我帶了他們去中國城坐椅吃東西.
他們吃得很開心之餘.
那些白人就一個一個問.
到底發生什麼事.
香港到底發生什麼事.
在同台的旁邊那位.
是一個在北京來的傳統人.
他們一問這些問題的時候.
我就想.
怎麼回答呢.
因為他是在北京來的.
我跟他從來都沒有談過這個問題.
不過他就開口說.
他說這個世界不可以這樣.
我們要支持香港.
你居然說他這樣.

$^{521}$真是.
打開了話題我才跟他談.
他說我在北京來的.
我很清楚北京裡面.
或者中共的政策.
在想什麼.
他首先說什麼呢.
他說教會被打壓.
他說北京三間最大的教會被封殘.
中共滲透在不同的地方.
去製造假新聞.
這位弟兄.
我認識的朋友.
他跟我說.
連他媽媽.
她媽媽當然是北京人.
她媽媽都不明示你.
我就跟他談這個話題的時候.
你怎麼幫你的媽媽.
明白這件事.
我就是將事實告訴她.
希望她能夠聽.
因為她在裡面.
聽的新聞.
我的資料.
全部都相反的.
他跟她解釋.
我跟她突然間好像.
感情好很多.
以前不是這麼熟.
但他這樣說的時候.
我發現.
我們不能夠活在這個世界當武功.
我們信上帝.
我們更加需要自首功利.
我們更加需要明白.
到底什麼是真新聞.
什麼是偽新聞.
什麼是公義的事.
什麼是不公義的事.

$^{561}$我們要懂得去辯認.
正因為我們相信有上帝.
這個世界才會有功利.
我問一下你.
如果這個世界沒有上帝.
沒有永恆的話.
就是說人生過了之後就是完結.
不需要來到審判處交帳.
是沒有公平公義的審判.
這個世界有什麼所謂.
誰的拳頭大.
誰就有功利.
誰在乎.
正因為我們相信有這位審判處.
這些私人哀歌背後的那位公義之主.
我們更加需要去關心.
更加需要去明白.
不要只說香港.
上個星期.
如果你有聽過新聞的一些故事的話.
你就知道.
在非洲有一個聖經翻譯員.
他是Wikliffe.
Bible Translation這個組織的翻譯員.
就在非洲一個國家裡面.
就是被武裝分子用刀刺殺.
他太太被武裝分子從家裡拉出來.
然後斬了她的手.
因為他們不想他翻譯聖經.
這些人都只是看到詩篇.
22篇.
上一半.
下一半在哪裡.
未必我們親身可以體會.
不過.
在基督裡面.
我們就看到詩篇22篇.
因為耶穌已經經歷過.
有上帝為他申冤.
有公義為他長年.

$^{601}$而且這個公義.
讓我們這群罪人可以分享.
這就是詩篇22篇最奇妙的地方.
沒錯.
耶穌是經歷過受苦.
經歷過不義的對待.
但我們知道.
死後祂復活.
最終有公義的申冤.
因為背後有一位永恆的主宰.
在耶穌裡面.
我們就有耶穌給我們的公義.
耶穌給我們最後的聖地.
我們與此相信.
\newpage



\section{詩篇}
\label{sec:EogeiLrQX3g}
\textbf{9/8/19 主日講道: 詩篇人生 -有福的人(陳崇基牧師)}
\newline
\newline
連結: \href{https://youtube.com/watch?v=EogeiLrQX3g}{\texttt{https://youtube.com/watch?v=EogeiLrQX3g}} ~~~~ 語音日期: 2019-09-12
\newline
\newline
\hyperref[sec:C3L6GM4kyjo]{\small{< < < PREV SERMON < < <}}
~
\hyperref[sec:index_chronic]{\small{[返順時目]}}
~
\hyperref[sec:index_scriptual]{\small{[返順卷目]}}
~
\hyperref[sec:7VnGbL9RHgg]{\small{> > > NEXT SERMON > > >}}
\newline
\newline
$^{1}$好,我們一起來讀聖經好嗎?.
我們回到詩篇第一篇.
第一篇,我們一起來讀.
不從惡人的計謀.
不暫聚人的道路.
不助洩漫人的座位.
遺棄愛耶和華的律法.
晝夜思賞,者人便為有福.
他要像一棵樹栽在溪水旁.
安事後結果子.
葉子也不枯幹.
凡他所做的盡都順利.
第四節,惡人並不是這樣.
乃像康丕被風吹散.
因此,當審判的時候.
惡人必站立不住.
罪人在二人的會中也是如此.
因為耶和華知道二人的道路.
惡人的道路卻必滅亡.
好,不如用回這個吧.
好,差很遠.
我們今天結束詩篇的系列.
我們最後看詩篇第一篇.
有沒有想過詩篇是一百五十首詩歌來組成的?.
所以我們有時候想想.
你買的CD或唱片.
可能有十首,二十首的歌.
有沒有想過.
擺放的人或設計的人.
他將第一首歌放在開頭.
是有他的榮譽.
你想想,如果我們在說敬拜的時候.
我們開頭敬拜的時候.
當然有音樂和唱詩.
但哪一個在耶穌的時代之前.
將一百五十首詩歌編排的時候.
他是將詩篇第一和第二篇放在開頭.
第一篇和第二篇不是唱出來的.
是一種教導的詩歌.
是一種智慧的詩篇.

$^{41}$就像你來到崇拜的時候.
第一件事做的是講道.
你一坐在那裡十點鐘的時候.
我們就講道.
講完道我們才唱歌.
我們會覺得有點不自然.
但其實也有他的榮譽.
第一篇詩篇是一篇教導的詩篇.
第二篇也是.
就像告訴我們.
我們首先要明白.
怎樣來敬拜神.
怎樣來認識我們的主.
我們要先聽一聽.
要受教.
詩篇第一篇就是有這樣的用意.
第一個字其實是有福.
中文的翻譯當然不能將福字放在最頭.
因為不會一開始就說有福的.
那個人就是這樣的.
因為文法的緣故.
所以我們就放在第二句.
什麼什麼什麼.
他不做什麼他做什麼.
這個人為有福.
整本詩集150首詩歌.
第一個字其實是有福.
在聖經裡面有兩個主要的字是講到福氣的.
第一個字就是在創世記裡面出現的.
上帝創造了人之後.
祂就去祝福他們.
祝福萬有.
叫他們生養眾多.
叫他們多生.
那個字就叫做Barak.
意思就是說你得到上帝的祝福.
就是上帝給你生命的能力.
生命的potential.
productivity.
multiplication.

$^{81}$給你那種發揚光大發展下去.
這樣的意思.
所以當耶穌復活之後.
在升天之前祂就祝福門徒.
祝福初期的教會.
叫初期的教會什麼.
他們要去傳福音建立教會.
祂就祝福他們.
這個就是同一個字.
就是叫他們去multiply.
去他們的地方.
然後將上帝的福音傳開.
然後去發展去develop.
去multiply他們.
那個是第一個字.
Barak.
祝福.
不是詩篇第一篇開頭那個字.
第二個祝福的字就是.
這個所謂有福這個字.
這個字是什麼意思.
這個字是講到將來的.
就是說你將來一定會得到上次.
但是你將來得到上次.
是在乎於你現在和上帝的關係.
所以那個祝福是有些不同的.
那個祝福不是一定講你現在.
是肯定講你將來.
就是說將來你一定會得到reward.
in the future.
你一定會得到上帝的blessing.
但是你現在要和上帝搞清楚你的關係.
所以這個祝福就有些不同.
這個祝福很多時候.
你現在的時候可能是受困苦.
你看這個字怎麼用.
在約伯記有一節.
是第五章第十七節.
他說Blessed is the person.
whom God corrects.

$^{121}$就是那個有福的人是怎樣呢.
就是上帝在管教懲治的人.
如果有人在管教你的時候.
你不相信你會很開心.
你可能在一些很困難的時候.
一些困境.
有很多軟弱等等.
這個是你現在可能在受苦.
但是為什麼會有福呢.
因為你將來.
ultimately 最終末的時候.
你就會得到上帝的祝福.
這個就是詩篇第一篇.
第一個字告訴我們.
有福的人將來會是怎樣的.
現在是怎樣的.
等於譬如你年頭的時候.
老闆跟你說.
年尾的時候你會有花紅.
會有bonus.
你現在努力工作.
現在努力工作好不好.
當然不好.
但是你要知道將來.
你就會有bonus.
或者你現在開學了.
努力讀書.
老師告訴你.
你要努力讀書.
你一定會有好的成績.
你現在努力讀書的時候.
當然不是很爽.
但是將來你會結尾的時候.
見到那個結果.
你去找一個人.
教你.
訓練你做運動.
他說你努力工作.
努力做你的運動.
努力去exercise.

$^{161}$將來你一定會收到你要的成效.
正在做的時候當然不爽.
但是將來會有.
這個就是詩篇第一篇所說的.
有福的人就是這樣的人.
耶穌其實也是這樣用福字.
不是說現在很幸福的意思.
你記得登山寶訓.
說到門徒的特質的時候.
他說虛心的人有福了.
they will be blessed.
為義受迫害的.
是有福的.
人辱罵你的時候.
迫害你的時候.
是有福的.
應當歡喜快樂.
因為什麼.
future.
天上的獎賞是大的.
以前的先知人都是這樣對他們.
so you are blessed.
不是因為現在你很有錢.
而是你將來一定會有reward.
一定會有獎賞或者賞賜.
或者第一節.
詩篇第一篇第一節.
可以這樣去翻譯.
或者必得神賞賜的人.
正是一些不從惡人計謀.
不暫罪人徒勞.
不坐視萬人座位的人.
為喜愛耶護和律法的人.
這個人就是必得上帝的賞賜.
詩篇一二篇就是這樣開始的.
為什麼一二篇是一起的呢?.
因為有很多字眼是一樣的.
特別是最後的字眼.
第二篇最後最後的字眼.
第二篇大家聽過了.

$^{201}$我解釋過了.
第二篇講到耶和華已經高立了他的君王.
最終就是耶穌基督.
他就是治理萬國的這位君王.
我們當以咀親子.
恐怕他來到審判的時候.
他發怒.
我們應該去投靠他.
如果這樣做的話.
就是有福.
最終得到神的賞賜.
所以一二篇就像兩扇大門.
去介紹給我們.
這位值得我們敬拜的神.
我們首先要怎樣去準備自己.
才能去敬拜他.
意思就是說.
有福的話.
如果你真的想有福的話.
最終得到上帝賞賜的話.
你就要做一個承諾.
做一個未來的人.
做一個將來的人.
未來人.
什麼叫未來人.
意思就是說我們的目標是放在將來.
將來的天國.
將來的新天新地.
不是放在現在.
不是說我現在有福了.
上帝我求他什麼.
他就給我.
這個不是福.
真正的福是將來.
你知道最近有什麼新聞.
在基督教裡面.
有一個人叫做Benny Hinn.
Benny Hinn在過去二十多年.
他建立了一個超大的教會.
這是一個所謂靈因派裡面.

$^{241}$一個很極端的位置.
他經常說一些叫做.
Prosperity Gospel.
就是一些叫做幸福音.
就是說你信主的話.
你就有幸福了.
你現在信主的話.
你去相信主的話.
你給奉獻的話.
上帝就會雙倍給你.
現在就會給你.
所以有很多人去那些教會.
有些人去那些教會.
不代表那些教會在傳福音.
他只是在傳幸福音.
二十多年都可以這樣做.
Benny Hinn.
很多人去那些教會.
去他的教會.
但是最近有什麼新聞.
他居然在自己的公開聚會.
他說我要悔改.
我要改回我的神學.
因為原來我一直以來所傳的福音.
是一種幸福音.
突然間他這樣扭轉.
所以這是一個很大的新聞.
一個人歸回.
明白到底上帝所說的福.
或者上帝所說的福音.
到底是什麼.
這就是第一扇門我們要打開的.
我們來見上帝之前.
想清楚到底我們所說的福.
Blessing.
上帝給我們Blessing.
是不是好?當然好.
但這種Blessing的本質是怎樣的?.
是叫我們做Future People.
必得上帝賞賜的人.

$^{281}$Future People.
你留意詩篇第一篇.
這裡的流程.
他說不從惡人的計謀.
不站罪人的道路.
不坐褻瀆萬人的座位.
有三對東西.
三對的東西.
是有進程的.
第一.
計謀.
Thinking.
Council.
思想.
道路.
走出來的.
就是生命行為.
生活行為.
座位.
是地位.
身份.
所以不只是想東西.
不只是生活.
更加是他坐的座位.
他有的身份.
特別掌權的身份.
不從.
不站.
不坐.
不從就是不跟隨的話.
的意思.
幫兇的意思.
不站.
是不支持.
不附和.
惡人的做法.
立場.
不坐.
就是不捐卵.
他的身份.

$^{321}$惡人.
是一些會暗算人的.
罪人.
是一些會做事去害人的.
褻瀆萬人.
是一些用言語.
來行暴的.
所謂言語藝術就是這個意思.
不是那種藝術.
是假的.
Fake的.
藝術.
把歪理說成真的.
或者張開眼睛說謊.
異人和惡人有什麼分別呢?.
詩篇第一篇開始就告訴我們.
在上帝的眼中.
和上帝有關係和沒有關係的人.
只有兩種.
一種是異人.
一種是惡人.
一種是罪人.
在詩篇整篇裡面.
或者一些智慧文學.
經常提到異人和惡人有什麼分別.
很簡單去解釋定義.
其實異人是什麼人呢?.
就是因著上帝和你的關係.
你會.
disadvantage yourself.
to advantage others.
損己.
利人.
Wicked.
就是一些漠視上帝.
去 disadvantage others.
to advantage oneself.
one self.
異人與惡人.
不只是做出來.

$^{361}$不簡單.
也不是說異人是完美的人.
而是將別人.
擺在前面.
惡人呢.
就將自己擺在前面.
這裡說.
你敬拜上帝的時候.
因著上帝.
因著你和主的關係.
你做一個righteous的人.
做一個異人的.
心態,價值觀.
因著上帝.
我和他的關係.
我advantage其他人.
寧願.
很多時我disadvantage我自己.
這是一個ethical commitment.
一個倫理上的.
毀身.
先行這一步.
這一個就最終.
會有福.
最終會得到賞賜.
這個人怎麼去形容他呢.
為喜愛而華律法.
晝夜思想者人便為有福.
最終會得賞賜.
他要長一棵樹.
栽在溪水旁.
安事後結果自.
業子也不富幹.
凡他所做的一切盡到順利.
這裡說為喜愛而華律法.
他不是說做了.
行了律法.
這是一個心智上的方向.
一個喜愛是心裡的東西.
不是說行足律法.

$^{401}$就會得福.
而是.
心裡喜愛的東西.
這裡假設了什麼呢.
就是一個新的生命.
我們本身不喜歡上帝的律法.
但我們因為信了主上帝的福音.
上帝的生命進入了我們.
改變了我們.
我們就會開始喜愛.
新的心.
新的生命.
你有沒有想過.
信主前信主後有基本上的改變.
你對上帝的態度.
你對身邊事物的態度.
會慢慢漸漸有改變.
你不會變成正人.
但你基本上.
你會有改變.
你的想法會有改變.
你喜歡的東西會有改變.
有誰喜歡喝黑咖啡呢.
你記得你第一次喝黑咖啡的時候.
你喜不喜歡.
是獲得的味道.
有些是後天習慣.
養成的.
一些味道.
那些喜好.
所以你喜歡喝黑咖啡.
你是真的懂得喝咖啡的人.
因為你獲得了那個味道.
有誰喜歡吃苦瓜呢.
今天有個姐妹.
給了我一個大的苦瓜.
那個是獲得的味道.
有誰喜歡吃榴槤呢.
這些是獲得的味道.
你第一次吃的時候.

$^{441}$你記得有多辛苦.
不辛苦嗎.
這些是天生的.
上帝的說話都是一樣.
是一種獲得的味道.
當你第一次看聖經.
可能真的不太明白.
上一下課程.
聽一聽.
啟開一下我的心靈.
學一些新的東西.
慢慢你就會喜愛.
更加喜愛.
這就是詩篇第一篇說.
你見上帝之前.
你需要去後天有的習慣.
為喜愛而獲得法.
是晝夜思想.
當然不是說.
聖經裡面說的晝夜思想.
不是說什麼都不做.
不務正業,我們只是去靈修.
不是這個意思.
而是習慣,什麼叫習慣.
是我們意識上的習慣.
當你很忙的時候.
你當然不會去思想上帝的律法.
當然不會.
但當你騰空時間的時候.
生命的韻律.
有一些優子乎的時候.
你第一件事會想起什麼.
你應該想起上帝的律法.
想起上帝自己.
想起上帝的話語.
晝夜思想.
也不是說騰空我們的.
思想.
也不是說騰空我們的.
心懷意念.

$^{481}$任何東西都倒空了.
不是這個意思.
而是.
我們勤於去思考.
勤於去研習.
去學習.
去懂得去.
發問,懂得去思考.
懂得去翻查,甚至.
去到自我批判.
變成自我批判.
想清楚到底我自己的想法.
到底有沒有需要改變的地方.
等於我剛才所說.
有人改變他自己的神學.
去想清楚.
到底以往的理解可能太初步.
現在是需要去進一步.
耶和華的律法.
Torah.
代表它的指引.
它的指引.
我們需要來到喜愛上帝給我們的.
指南.
在我們生活當中生命裡面有的指南.
聖經裡面說到.
Torah的時候特別夠弱.
都是說到整體的.
不能切割開來的.
一個整體的教導.
每字每句.
每篇每章都是重要的.
這個叫做耶和華的律法.
The Torah.
不是只選一些我喜歡的.
而是整個的Torah.
我記得.
那時候.
去準備.
會議的時候.

$^{521}$我去了一個神學院.
他們的口號是什麼呢.
他們的口號是.
The whole council of God.
就是說上帝.
整全的心意.
我們需要學習和祂交代.
是整全.
不可以選擇性的去做.
不是說喜歡的東西我就拿.
不喜歡的東西我就不理它.
Torah.
The Torah.
上帝整全的律法.
在華盛頓DC.
有一間博物館.
叫做Museum of the Bible.
是2018年開的.
應該是2018或2017.
總之是最近開的.
一間博物館.
它是專門展覽一些.
舊的聖經.
最近在展館裡面.
展示了一個文物.
這個文物不是那些.
Ancient history.
只不過是200年左右的文物.
是1807年所印製的一本聖經.
那本聖經可以叫做.
Slave Bible.
奴隸聖經.
為什麼叫奴隸聖經呢.
是當時.
19世紀初.
1807年.
1800年代的時候.
印製出來的一本聖經.
這本聖經是一種.
叫做.

$^{561}$抽了一些.
聖經的經文.
來印出來的.
不只是印整本.
只是印一部分.
專門給當時的.
所謂的黑奴來看的一本聖經.
這本聖經有什麼有趣呢.
就是舊約.
有90\%沒有了.
新約呢.
有50\%是不見了.
整本的聖經有多少張.
1189張.
1200張.
這本聖經呢.
只有232張.
最有趣的地方.
不是他留了什麼在這本聖經裡面.
而是沒有留下.
撇走的.
那些到底是什麼的經文.
有什麼經文.
是沒有了呢.
加泰書三章.
沒有了.
在基督裡面是不分猶太人.
希臘人自主的.
圍爐的.
或男或女的.
這本聖經就沒有了.
加泰書三章.
28節可能我都講過.
是人權平等的宣言.
Magna Carta.
Humanity.
這本的.
Slave Bible就沒有將這本聖經.
放在上面.
如果你想看的話.

$^{601}$這本聖經你不一定要去華盛頓看.
網上可以看.
因為已經200年了.
但你想看網上.
到底他留了什麼.
我給你看看.
這裡是.
創世記.
第45章.
你知道創世記到第45章.
發生什麼事.
就是亞國的其中一個兒子.
約瑟被賣到埃及.
做什麼的.
做奴隸.
做了很多好事.
神就祝福他.
故事就在這裡.
創世記有50章.
他停在第45章.
你繼續看下去.
第45章一跳就跳到哪裡.
跳到出埃及的Exodus.
19章.
如果你略為對聖經有些認識的話.
你就知道他跳過了什麼.
就是跳過了Exodus.
出埃及的故事.
出埃及的故事是說.
上帝打暴權法老.
對他一點都不客氣.
然後領以色列人.
本身在埃及做什麼的.
奴隸,沒錯了.
領我們出埃及.
去到西乃山然後賜艾明給他.
整段故事就沒有了.
就是說故事的訊息.
其實就是上帝對這些暴權一點都不會客氣.
他會打殘他們.

$^{641}$然後領他們的子民出埃及.
他們就得到自由.
Liberation.
那本Slave Bible.
就沒有放這件事上.
還有少了什麼呢.
很多舊約五經的律法.
是說到奴隸制度.
就是說去制衡奴隸制度.
當然你要知道.
聖經時代的那種奴隸.
又不和二百多年前那種.
黑奴的奴隸.
有些不同.
地位上當然階層上有相似的地方.
但是五經裡面有很多.
是制衡奴隸的主人.
叫他們不要苦待他們.
叫他們.
譬如到第五十年的時候.
有希年有Jubilee.
你就要釋放那些奴隸.
除非他想留下來服侍你.
等等這些制衡奴隸主人的經文.
全部被剔除.
還有剔除了什麼呢.
整本啟示錄都沒有了.
你知道啟示錄說什麼嗎.
啟示錄請你不要以為.
是說中東現在新聞發生的事.
No.
那些apocalyptic literature.
天啟的文學.
啟示錄是其中一半.
是什麼文學呢.
英文有個字可以形容這種文學.
天啟文學是一種resistance literature.
就是那些被打壓受打壓的.
受欺壓的人.
很多時候猶太人.

$^{681}$他們就會經天啟文學.
去批判在上掌權.
這個是apocalyptic literature.
你不要打開apocalyptic literature.
找到啟示錄的時候.
嘩 現在中東這麼亂.
發生這麼多事.
一定是上帝預言了.
不是預言.
是resistance.
是讓受打壓的人.
能夠上到天上的時候.
知道上帝是延長權.
他們是不需要懼怕.
他們可以繼續忠心跟隨上帝.
這個叫apocalyptic literature.
這本奴隸的聖經.
整本啟示錄就沒有了.
當然沒有了這段.
這是最終啟示錄最後一句.
一切在書上預言的作見證.
如果有人在預言上加添了一些東西.
上帝會將寫上去的災害.
加在他身上.
剔走了的.
上帝會在書上所說的應許.
都會剔走.
是一個警告.
留下了的.
有些是什麼經文呢.
以弗所書錄的五字就留下了.
你們作僕人的要懼怕戰經.
要用一輩子的心理去服侍你肉身的主人.
好像服侍基督一樣.
誰印這本聖經呢.
很明顯不是.
哎呀,忘記了,不好意思.
因為他是abridged version.
是故意這樣選出來的.
很可能是當時的英國.

$^{721}$大英帝國的商家.
去到殖民地的時候.
他們操控殖民地.
加上當時的城坐.
帝國主義.
去的宣教士.
他們就印了這些聖經.
當然這本不是當時唯一一本聖經.
是其中一本.
是專門給黑奴看的.
就是那些獨裁者.
他們刪改聖經.
以聖經的名義去做一些欺壓的事.
斷章取義等等.
最近有一個海外的牧者.
去提醒我們.
就算現在中國.
都想自己翻譯聖經.
這是最危險的.
十誡裡面第一誡剔除.
其實我們看這些東西.
我們不需要太驚訝.
為什麼呢.
因為你和我都可能寫得到.
我們喜歡的東西就留下.
可能你心目中.
Abridged version的Bible.
就可能是這本.
原來聖經是在維護我現在的地位.
我現在舒適的生活.
上帝是在保障這件事.
可能你的Abridged version就是這樣.
可能你一直信主.
你所看的聖經原來就是一本.
你自己按照你的喜好去抽出的聖經.
而不是聖經的全部.
不是The Torah.
是一種選擇性的去跟隨.
我的意思就是.
我們不需要太驚訝.

$^{761}$因為我們很多時候都會做.
為什麼我們會有門徒訓練呢.
為什麼呢.
為什麼需要每一季.
我們四個禮拜.
我們坐下來.
專門去研究一些聖經.
讓我們有整全的去理解.
本聖經到底說什麼.
我這個門徒訓練的課程.
叫你記的金句.
都不是你喜歡記的金句.
你記得我們都說過.
我們喜歡記的金句.
都是令我們溫馨.
保障著我現在舒服的生活.
那些金句我們就是不停記.
有很多其他重要的經文.
我們就失去.
我們是選擇性的.
當然沒有人可以一坐下來.
就這樣看聖經.
永遠都是有些選擇性的.
但問題就是不能斷章取義.
這就是喜愛耶和華律法要做的事.
他說如果你這樣做的話.
就好像樹栽在溪水旁.
是一種有生命根深蒂固的.
安事後會結果的.
你知道近東的社會.
聖經的世界.
每逢看到樹的時候.
有什麼地方有樹?.
伊甸園有樹.
聖殿裡面有樹.
有植物.
將來新天新地耶路撒冷.
都是會有樹的.
有生命的樹.
意思就是說.

$^{801}$如果我們這樣喜愛耶和華律法.
去到思想去的話.
我們就好像樹栽在溪水旁.
就是說我們會來到神的同在.
你可以享受到上帝的presence.
是有requirement的.
不是只是等上帝的同在.
是我們進到上帝的同在.
並且我們要就夜思想.
很奇怪的.
詩篇開始的時候.
他不是給我們一首歌來唱.
而是給我們一個教導.
警告告訴我們.
你要最終得福.
最終得賞賜的話.
你就必須喜愛耶和華律法.
不從惡人的計謀.
你要立志來到上帝的面前.
在他的律法底下生活.
惡人就不一樣了.
惡人不是這樣.
乃象康丕被風吹散.
因此當審判的時候.
惡人必站立不住.
罪人在二人的會中也是如此.
因為耶和華知道.
看顧二人的路.
惡人的道路.
卻必滅亡.
康丕是什麼?.
康丕.
康丕是那些農夫.
在那些高地.
那個flushing floor.
打麥的高地.
在和場那裡.
採摘了大小麥的時候.
就放在那裡.
然後硬碎或打出來.

$^{841}$將麥子打出來.
然後康丕就是那些殼.
你知希伯來文.
樹和康丕的音是很相似的.
所以詩篇第一篇就是選擇.
惡人象康丕.
因為一定會被風吹散.
但是不會有將來.
一吹就走了.
你的生命.
來到上帝的面前.
上帝要你好像一棵樹.
來到上帝面前就要認識他.
這是生命的決定.
這是需要下功夫.
需要去研讀.
需要用心.
這就是一個最重特賞賜的生命.
求主幫助我們教會.
讓我們教會成為一個.
詩篇第一篇的教會.
可以同心協力.
主讓我們聽見你的說話.
聽見你的提醒.
主讓我們離開這個地方.
我們去探索課程.
去到團契.
主讓我們不要忘記.
我們的生命不能像康丕.
要像一棵根深蒂固的樹.
有生命的.
要喜愛你的律法.
求主你親自讓我們教會弟兄姊妹.
去同心.
去讀你的話語.
去認識你自己的心.
讓我們行在你自己的心意當中.
讓我們成為一個義人.
願意去建立其他人.
願意去傳你的福音.

$^{881}$去建立你自己的國度.
幫助我們.
讓我們今天做一個生命的決定.
去明白你的心意.
我們這樣祈禱 奉傳明教.
\newpage



\section{}
\label{sec:7VnGbL9RHgg}
\textbf{9/22/19 主日講道: 在敬拜中遇見你 (Brenda Li 傳道)}
\newline
\newline
連結: \href{https://youtube.com/watch?v=7VnGbL9RHgg}{\texttt{https://youtube.com/watch?v=7VnGbL9RHgg}} ~~~~ 語音日期: 2019-09-29
\newline
\newline
\hyperref[sec:EogeiLrQX3g]{\small{< < < PREV SERMON < < <}}
~
\hyperref[sec:index_chronic]{\small{[返順時目]}}
~
\hyperref[sec:index_scriptual]{\small{[返順卷目]}}
~
\hyperref[sec:FaXJPQMXp6c]{\small{> > > NEXT SERMON > > >}}
\newline
\newline
$^{1}$弟兄姊妹.
(笑).
那麼.
不好意思 因為.
經過兩天很長的.
教導 唱歌.
大家都聽到我的聲音.
應該只剩下一半.
那麼 不過呢.
我發覺每次在這些位置.
就說 哎呀 怎麼辦.
因為我們靠聲音.
很緊張的時候.
神的恩典就會出現.
我記得有一次.
我最後是比現在的狀況更差.
但是我教了三個小時.
歌唱的教室.
是完成了的.
但是那一晚就失聲了.
之後就失聲了一個星期.
那麼我都憑信心宣告.
這件事是今天不會發生的.
如果會的話.
我會把看醫生的錶.
發給你們的教會.
(笑).
給我機會來到Chicago.
更加感恩來到你們當中.
因為我發覺這個教會.
就像它的名字一樣.
Urban Voice.
就是在社區裡.
發出宣告神的恩典.
和愛的一把聲音.
既然你們教會的名字.
已經說你們是一個Voice.
我相信神的心意.
更加是希望這個成為一個.
在敬拜裡.

$^{41}$向祂和社區發聲的教會.
今天的訊息主題.
在敬拜中預見神.
稍後會說更多.
我們如何預見神.
不如開始之前.
我們一起做一個簡單的禱告.
給我們一群弟兄姊妹.
也有當中未認識你.
但學悟去經歷你更多的朋友.
我們來到.
願意你都藉著今天.
我們剛才的敬拜.
藉著我們領受的訊息.
藉著我們彼此的相交.
我們能夠更深去體會.
耶穌基督你這份愛和恩典.
祝福我們今天的時間.
給我們每一個.
離開的時候.
我們裡面都很飽足.
因為我們能夠預見你.
奉耶穌基督的聖名祈求.
阿們.
可能我們都覺得.
我們對敬拜不陌生.
可能我們很多人.
就算我們現在.
在教會每個星期的聚會.
無論是主日或是.
我們的團契小組.
我們都會有一些敬拜的部分.
在當中.
甚至你回想你未信耶穌之前.
你未正式.
很毀身在教會之前.
你應該都有接觸過.
詩歌或是敬拜.
或許這份熟悉感.
會給我們一個感覺.

$^{81}$其實我們都知道什麼是敬拜.
我們每個星期都在敬拜.
我們都在做敬拜這件事.
一個小時的聚會.
20分鐘的敬拜.
真的很多.
我們覺得敬拜應該很熟悉.
每個人都在.
應該都有開口唱.
其實我們都懂得.
應該不需要怎樣學.
學敬拜的應該留給敬拜隊.
敬拜的氣氛如何.
敬拜的外貌.
是敬拜隊的責任.
所以有人來訓練.
訓練敬拜隊就夠了.
因為這個範疇屬於敬拜隊.
但原來這個理解.
不是太正確.
剛才的報告也有提及.
我們都要學讀經.
聖經也給了你.
手機裡也有.
我們說明我們要看.
你看就行了.
為何要學呢.
因為我們在當中有方法.
讀經不是只看字.
我們就會明白.
讀經裡有很多不同的途徑.
可以讓我們在經文裡.
吸收更多神話語.
明白更多.
以致我們能應用在生命當中.
其實敬拜也是一樣.
不是有幾首歌出來.
我們跟著歌詞唱了.
就等於我們敬拜了.
原來敬拜也有.

$^{121}$我們需要學的東西.
如果我們首先從敬拜是什麼.
我們先理解敬拜的意思.
敬拜的原文是.
是在說事奉神.
這個事奉神.
是在說我們每一天的生活.
我們的生命.
我們每一天的擺盤.
其實都是一種敬拜.
但敬拜同樣有另一個意義.
是說到好像一個.
你可以看成一個活動.
即是一個特定的時間.
一種的行動.
是我們將這個時間.
將這個行動分別為性.
我們特別去專注於神.
就是為了讚美祂.
而我們今天用的方式.
是一種知歌音樂的方式.
我們用這樣的方式.
來做讚美神的行動.
所以如果我們這樣去理解敬拜.
我們會知道敬拜不是說唱歌.
雖然我們用唱歌的方式.
我承認如果一聽到一些人說.
「哦」 即是進來教會.
「哦」 正在唱歌.
我的心就很不舒服.
甚至有一點被冒犯.
因為我覺得不是在唱歌.
我們是在讚美.
其實唱歌是一個動作.
或者是我們的表現.
但唱歌的背後.
其實裡面是有它的意義.
亦都有神在我們當中的心意.
待會我們可以說得更多.
既然我們說敬拜不是唱歌.

$^{161}$既然不是唱歌.
就不是說我們的心情.
有人說心情好就唱歌.
心情不好又唱歌.
在這裡宣洩一下情緒.
或者會覺得如果是唱歌.
唱得好的就唱.
唱得不好的就應該唱少一點.
不要吵到別人.
但如果敬拜不是唱歌的時候.
唱歌好不好聽.
喜歡不喜歡唱歌.
唱歌好不好聽.
這些完全不應該是考慮.
不應該影響我們敬拜的心情.
因為敬拜並不是唱歌.
敬拜是讓我們重新將焦點放在神的身上.
讓我們定精在祂的恩典.
在祂的美善之中.
如果說到這樣.
敬拜是一件很認真的事.
原來神對我們的敬拜都是有要求.
為什麼呢?.
其實我們今天的敬拜.
就相等於昔日舊約時代的獻祭.
人為什麼要獻祭?.
因為贖罪.
如果你看回聖經裡很多關於獻祭.
其實有很多規矩.
有班部長十屆之後出家求妓.
之後一直說的.
都是說要怎樣去獻祭.
祭司有什麼條件.
有什麼要求.
有什麼規矩.
你還沒記住.
基本上全書卷都是說一些關於獻祭的規矩.
相比之下.
我們今天應該幸福很多.
我們都不需要學到這麼複雜.

$^{201}$我們都不需要記每一個祭的名字.
什麼時間應該獻什麼祭.
祭司要做什麼.
如果祭司做得不對又會死的.
今天敬拜人都不好.
唱一奏音應該不會有什麼事.
所以相比之下.
我們今天的敬拜應該容易很多.
我們也不需要帶著那個生存.
因為那個羔羊已經是耶穌基督自己.
成為了那個羔羊.
成為了那個祭物.
他一次的獻上.
又成為了我們的祭司.
以致我們可以毫無阻隔地去到神的面前.
去敬拜神.
去肩負上帝和好.
以致我們能夠重新成為神的兒女的身份.
既然今天的敬拜的要求.
和我們需要做的事.
相比以前容易很多的時候.
學應該不難.
不過我們一定要學.
我想我們記得一件事.
我們的敬拜.
其實是基於.
神對我們的敬拜的根基.
是耶穌基督的救恩.
這件事很重要.
這件事也是我們每一次回到神的面前.
無論我們聚在一個聚會裡的群體.
或者我們自己去敬拜神.
我們要知道這個敬拜的根基.
是神的救恩.
為什麼這樣說呢.
因為你知道現在很多異端.
很多不是基督教的.
不想說邪教.
宗教的群體.
他們也有詩歌.

$^{241}$有時候如果你一不小心.
你會不經意被他的旋律.
和某些詞彙欺騙了.
因為在佛教的詩歌裡.
你也不難找到「恩典」這個字.
「慈愛」這個字.
「感恩」這個字.
其實這些都是屬於我們基督徒.
常用或我們聖經裡常用的詞彙.
他的旋律也舒服,優美.
但你會知道.
你在那首歌裡是找不到耶穌的.
另外.
一個最近很出名.
大家一聽到就知道.
要避之則吉的異端.
全能神教會.
即是東方閃電.
他們很厲害.
他們已經發展到有不同語言的詩歌.
我聽過國語.
我聽過韓文.
聽過英文.
應該有其他語言存在.
但我不懂.
你看下去.
因為他是一個異端.
他也有說耶穌.
他也有說神.
但如果你是一個認識信仰的人.
認識真理的人.
你聽聽他的歌詞.
那些東西怪怪的.
不太對勁.
甚至連唱歌的人的樣子.
我不是猛猴怕.
我真的聽到那首歌.
看到他們的樣子.
我為何好像在表演.
好像很搞笑.

$^{281}$很老套.
那種假的.
很老套.
很裝模作樣.
裝模作樣.
很興奮的那種.
很假.
然後再看歌詞.
為何那麼奇怪.
後來再查看歌詞的名字.
就看到.
原來這是全能神教會的歌.
所以我們要知道.
我們的信仰.
我們的敬拜.
是基於救恩.
因為這件事.
在其他假宗教裡.
是找不到的.
他們可以用我們的詞彙.
但他們裡面卻沒有.
耶穌基督在我們生命中.
所做過的事情.
因為耶穌基督.
是屬於神的兒子.
唯有從神而來.
我們才能得到這份救恩.
所以.
亦因為這樣.
敬拜也無形中成為了我們.
屬神而來的一個身份的象徵.
我們知道.
這裡有不少香港人.
或是熟悉香港的文化.
我們知道前一陣子.
很多教會的團體.
在香港的街頭去唱詩等等.
其實當一唱詩的時候.
無論是詩歌的內容本身.
或是唱詩敬拜的表現動作.

$^{321}$其實已經讓人知道.
這群人是基督徒.
這群人是信神.
這群人是信耶穌.
原來我們敬拜和唱詩.
這是我們身份的象徵.
我們是一群得蒙救贖.
已經被神的救恩買贖.
亦成為了他們的兒女.
一群人的身份象徵.
所以原來敬拜當中.
有這麼豐富的含義的時候.
我也很想弟兄姊妹.
去看重我們的敬拜.
原來我們的敬拜.
在對外是一個見證.
我們的敬拜對神來說.
亦是一樣祂很重視和喜悅的事情.
剛才我也有提及.
敬拜這麼重要.
就好像我們獻祭一樣.
當然神也會對我們的敬拜有要求.
今天的敬拜.
神對我們的要求又是甚麼呢?.
其實聖經裡有不少地方.
是有提及敬拜.
但有一段經文是我特別喜歡的.
我相信我們也很熟悉.
經常聽到.
不足兩年.
可能一年多兩年.
我重新再看這段經文.
並且神也讓我有機會去看一本書.
令我更深入體會這段經文.
整個故事背後的意思的時候.
我發覺.
回想自己的敬拜經歷.
我真的體會到.
原來耶穌也是一個非常重視敬拜的人.
這段經文.

$^{361}$在約翰福音第四章.
就是耶穌與撒瑪利亞父人.
井旁婦人的對話.
我也問過教會的弟兄姊妹.
我知道教會裡.
很多弟兄姊妹.
都很熟悉聖經.
不用擔心,我不會考你們.
所以我不打算全段讀出來.
因為我相信你們對這個經文的背景.
是有一些認識.
但我還是想我們集中在這幾節.
耶穌在撒瑪利亞的井中.
與這個婦人的對話.
耶穌從問她要水.
作為一個切入點展開了他們的對話.
到耶穌對她說.
很有趣,耶穌問人要水.
然後對她說.
你喝過我的活水,永遠不會渴.
就像別人拿水,然後你又說自己有更好的.
耶穌真的很特別.
但就用這個活水的話題.
打開了.
讓這個婦人.
在他們的對話裡.
讓這個婦人知道.
耶穌好像不是一個普通人.
他好像知道一些一般人不會知道的東西.
他好像是一個先知.
不知道婦人是為了這個話題還是有其他原因.
其實這可能是反映了她內心的一種很真實的渴望.
她就說到敬拜這個話題.
婦人就說.
你們猶太人說要在山上.
要在耶路撒冷去敬拜.
但我們撒瑪利亞人就說要在山上.
說起這個敬拜的話題.
但耶穌就跟她說.
婦人,你要相信我.

$^{401}$時候將到.
你們拜父的不是在這個山上.
也不是在耶路撒冷.
你們所拜的你們不知道.
我們所拜的我們知道.
因為救恩是從猶太人那裡出來的.
這也再一次證明了我剛才和大家分享.
我們的敬拜根基是救恩.
我們再看下一段.
耶穌接著說.
那些真正去敬拜父的人.
是要用心靈和誠實去拜祂.
因為父要這樣的人來拜祂.
神是一個靈.
所以拜祂的人.
是必須要用心靈和誠實來拜祂.
我們常常都會用這段經文來講敬拜.
或者我們有時候祈禱.
我們也會用這段經文來講敬拜.
我們要用心靈和誠實來敬拜.
但其實什麼是用心靈和誠實呢?.
如果我們查考聖經.
坦白說這個和合本已經譯了很多年.
有時候它的翻譯未必最能表達到原先的意思.
如果我們看原文.
或者看英文的翻譯.
其實這裡所說的心靈和誠實.
我們會看為用心.
和誠實.
不會虛假.
不會騙人.
可能是這個意思.
很誠實.
但原來不是這個意思.
原來這裡是說.
我們是要用心靈.
按著真理來敬拜神.
英文是說.
In spirit and in truth.
Truth 是代表真理.

$^{441}$按著真理來敬拜.
可能對我們來說比較容易一點.
因為我們都會看聖經.
我們都有背經文.
我們對聖經都有一種追求.
我們都會認識.
還有每個星期都有.
牧者分享聖經.
我們應該在真理方面.
我們應該是抓得很穩.
但反而可能在心靈方面.
發覺原來我們很多時候.
我們不是很知道.
或者不是很明白.
怎樣用心靈來敬拜.
我也想補充一點.
在真理的敬拜.
除了我們可能會看到.
明白聖經.
或者不會像我剛才舉的那些例子.
假的詩歌不知道在拜什麼.
讚什麼.
但以為自己在按著真理敬拜.
除了這些情況之外.
其實有時候都會代表著.
我們是否真的將敬拜的焦點.
放在神和祂自己的話語裡面.
什麼情況不是將焦點.
放在神和祂的話語裡面呢?.
例如我們會留意.
究竟這個星期誰大.
如果唱得比較好的.
那個團隊的音質就開心.
如果唱得一般般.
有些難聽和悶的話.
我們就覺得他大了.
例如什麼不是專注在神.
例如音樂.
例如那首歌.
可能我們也不只是弟兄姊妹.

$^{481}$敬拜隊裡面也有.
我也聽過不少人說過.
不是,那首歌很難聽.
但我不是很想唱.
那首歌很老套.
我彈下去覺得好像90年代.
很多這些因素.
會影響我們敬拜的心情.
這些不是將焦點放在神.
不是按著真理.
就算那首歌如何.
除非歌詞有問題.
這又是另作別論.
就算那首歌的曲風如何.
就算他彈的人彈得好聽.
難聽唱的人好聽.
這些其實都不會影響神自己.
神是不變的.
神的恩典是不會改變的.
神的愛是不會改變.
神的信神不會改變.
神的美善不會改變.
神從來都不會改變.
所以我們也不應該被這些東西影響.
但有一樣東西會影響我們敬拜.
就是我們的心.
我們的心情.
我們的心靈.
我們如何去看自己和神的關係.
這些會影響我們的心.
來去敬拜.
但其實如何打開心敬拜.
打開心敬拜.
那如何打開心.
是不是要用刀來打開.
那叫打開心嗎.
很抽象嗎.
其實打開心是一個很直接.
但不簡單的一件事.
就是我們要放下自己.

$^{521}$放下自己本身在想的東西.
放下自己的執著.
放下自己的看法.
放下自己的掛慮.
放下自己那一刻的情緒.
但這個放下.
不是掩飾扮沒事.
因為神也不是要我們用這樣的方式去敬拜祂.
這個放下.
也代表著我們願意將這些東西交給神.
更重要的是我們願意用神自己的話語.
用神的真理.
用神的愛.
用神的應許.
來取代這些東西.
我們要學習如何用心靈來敬拜.
這件事真的要學習.
首先.
我們有時候聽到一些歌.
如果你懂得唱就跟著唱.
但有些時候我們也不需要很心急地唱.
但也不是叫大家不要唱.
或者說Brenda說完後.
大家合上嘴巴不出聲.
不是這個意思.
我們有時候也可以咀嚼一下歌詞.
我們理解一下.
我們感受一下歌詞.
然後我們理解了.
或者聽到這首歌.
我也很熟悉.
或者我想跟著去唱.
想跟著去唱的時候.
你真的要用心去唱.
你真的要盡心去唱.
你知道嗎?.
無論我們說話.
我們讀一些字.
或者我們唱歌.
我們唸口往式地.

$^{561}$讀在唇上.
跟你很認真地.
用心去讀是很不一樣的.
如果你不相信,你試一下吧.
今天沒有時間去做這些練習.
但原來我們真的盡心和用心去唱一些歌.
是需要我們的思想.
需要我們的專注.
需要我們的身體.
你唱歌也要用力.
你發聲的聲帶在這裡.
但你也要用力.
如果你再用你的身體語言.
向神表達你內心的感受.
其實這絕對是全身,全人的行動.
另外我們要找一些空間去聆聽.
我忘記昨天敬拜訓練的時候有沒有提到.
其實如果我們時間許可.
有時候我們也可以有一些音樂的空間.
讓人有機會去聆聽一下神.
因為當我們唱唱唱唱完一些歌的時候.
其實我們內心也有些感受.
有時候神一剎那也會跟我們說一些話.
我們也可以有些空間來安靜一下.
來聽一下,來感受一下.
甚至如果我們覺得我們已經聽到了.
我們可以用那些小小的空間.
輕聲和神簡單地回應.
這個回應可能很簡單.
「主啊,我知道了」.
很簡單的.
但我們不是單單地唱,單單地獻上.
我們也有期待.
我們可以從神那裡去接收一些東西.
剛才我提到關於身體語言.
時間有限也不會說很多.
不過很快速地去說一說.
昨天的訓練就說得比較多.
其實我們華人是一個比較含蓄的群體.
就從我們的教育制度就知道.

$^{601}$小朋友坐好,不要動,不要說話.
那些就是乖的小朋友.
到處走,又說話又多嘴那些.
是頑皮的小朋友.
就好像我小時候那樣,頑皮.
我們在這種塑造下.
我們不知不覺間就覺得.
我們的沉默,或者我們的不要亂來.
就是一種安分守己和最恰當的表現.
但我想說,其實這樣是不正常的.
小朋友怎會不叫,小朋友怎會不笑.
小朋友怎會不說話,怎會不跑.
是我們要他不要動,你明不明.
不正常,如果你的小朋友是這樣.
可能你要擔心.
同樣地,其實神創造我們.
是有很多反應.
是很自然地在我們的身體表現出來.
例如你很開心,我們就會笑.
剛才我可能說了一些冷笑話.
你也覺得好笑,你會笑.
或者我們.
舉個例子,就好像你看球賽.
你很喜歡那個球隊入分.
你會叫,如果他做得不好,你會叫交叉.
我們是一定會有反應.
情緒裡面,你的感動裡面.
你開心不開心,其實是會有反應.
例如我們見到人的時候.
我們表示歡迎.
我們會不會這樣.
歡迎你.
不會的,我們真的歡迎一個人.
我們很自然就會.
笑是最基本的.
歡迎你,或者.
我們做得多一點的,就是握手.
或者如果打開手,請進來.
或者同性別的.
通常見面就不至於這樣.

$^{641}$我們中國人可能走的時候.
同性別的就擁抱.
這些其實是很自然的反應.
感動的時候,傷感的時候.
可能我們會沉默,我們會流淚.
我們會想事情.
你真的覺得很卑微.
可能你會低頭.
甚至有些情況你會跪下等等.
其實這些都是神創造我們的時候.
給了我們一些很自然的情緒和一些反應.
但是呢.
就在我們這樣的文化的塑造底下.
就令到我們好像.
反而不能夠用我們最自然的方式.
來向神表達.
說到這裡,我想大家站起來.
站起來.
我想大家可不可以.
舉起你的手,先舉一隻.
好了,舉起兩隻.
我們一起拍手,好嗎?.
很有趣的,一邊拍手的時候.
樣子是笑著的.
我到目前為止還沒有見到一個很黑口黑面的樣子拍手.
如果你身邊有一個熟悉的人.
或者同性別的人.
你可不可以給他一個「曬黑」之類的.
(笑聲).
好了,這個示範完畢了.
大家先坐下來.
那你有沒有發現呢.
原來我們的肢體語言.
是會影響我們的情緒的.
因為原來我們剛才這樣拍手的時候.
不期然就會笑著.
其實你唱歌的時候.
如果那首歌本來是很興奮的.
你的心情也不是很差.
也會有一點開心.

$^{681}$但你跟著一起拍,會特別開心.
譬如有些歌很厲害.
譬如剛才那兩首就很厲害.
不過我不會將兩首小歌放在同一套.
很厲害的聲音.
譬如那些什麼.
「你公義降臨」.
或者提神什麼.
「願賦你的榮耀,權兵降下」.
「你是人,教授萬國萬民」.
那些歌詞.
當你唱得越來越激昂的時候.
一直重複.
重複的時候不是沒有歌詞可以唱.
其實只是在重複.
因為表示這句是很重要.
還有重要的是.
「We want it」.
就好像那些小朋友.
你答應了給我雪糕.
我就問十次.
「我有冰淇淋嗎」.
其實是「We want it」.
所以我們不斷重複.
We declare it.
但如果你一直唱.
你會不會覺得.
再唱.
「願賦你的榮耀,權兵降下」.
「你是人,教授萬國萬民」.
再重複.
它越來越澎湃,越來越厲害的時候.
我不知道大家.
我自己.
我也覺得自己是一個正常的人.
身體是有一點忍不住.
因為你在求.
你學求那件事的時候.
不會沒有反應.
不會十次都是一樣.

$^{721}$「願賦你私人國度,權兵降下」.
不會這樣的.
你會激動地說.
「榮耀降下」.
是很自然的,是會的.
因為我比較豁出去.
雖然對著你們.
我也可以很自然流動.
我們教經班的人是會這樣.
你們就不需要像我這麼誇張.
但這對我來說是很自然的.
不過我想說.
假如我們內心真的有這些感受.
我們是可以表達出來的.
我相信拍手.
譬如剛才第一首歌這麼厲害.
「萬主坐在寶座上」.
然後「萬主之主,萬王之王」.
如果是拍手.
我相信我們幾位牧者是不會走上來說.
「剛才是誰在那首歌之後拍手」.
我相信應該不會的.
或者如果我們真的有感動.
我們真的舉手.
我相信應該沒有人說.
「你的手阻礙我,請你退開,我看不到」.
應該不會的.
昨天也舉了一個例子.
你頭癢,也會舉手叫.
你想舉手,不用怕阻礙旁邊的人.
但你發生意外,我就不敢說.
這裡的空間應該足夠.
但其實原來上錄影廠做我們的時候.
他已經預備了我們的身心靈.
是對周邊的事物.
特別是對他.
是有反應,是有感受.
但我們要學習一件事.
就是我們要願意放下自己的執著.
而執著這個字的範圍很廣.

$^{761}$可能包括我剛才提到的思想,情緒,擔憂.
執著也可以包括.
這個不是我的習慣.
這個不是我的喜好.
但我想說.
其實你想想.
我們說我們來尋找神.
我們來到殿裡尋找神.
其實這個說法是不完全.
因為神已經在這裡.
神不單在這裡.
祂更加在我們的心裡.
既然這位這麼聖潔,這麼完美的主.
祂也願意來到我們當中.
甚至住在我們裡面的時候.
這份恩典是多麼大的時候.
會有什麼原因.
令我們對祂沒有反應,沒有感覺呢?.
其實耶穌在祂和撒瑪利亞婦人的對話裡.
整個聖經記載裡.
祂不單說給他們聽什麼是敬拜.
祂也展示給他們看.
也同樣給我們看什麼是敬拜.
原來敬拜就是耶穌基督願意去到人的當中.
去到撒瑪利亞.
去到一個不應該是猶太人踏足的地方.
去等待這個婦人.
去到井旁等待這個婦人.
為了要跟她說話,聊天.
為了要告訴她.
我知道你的事.
我知道現在跟你在一起的那個.
也不是你的丈夫.
我知道你過往有五個丈夫.
我知道你的羞恥.
我知道為什麼你選擇一個中午熱的時候來這裡打水.
因為你不想遇見其他人.
我知道你一直在這裡隱藏.
但是我這一刻來到.
我來到去尋找你.

$^{801}$原來當我們在敬拜裡面.
當我們願意去定位自己.
我們願意將自己放在一個位置.
就是願意去讓神進入我們的內心.
讓神來觸摸我們的內心的時候.
我們可以好像這個撒瑪利亞婦人一樣.
來遇見耶穌.
我們怎樣知道去遇見耶穌.
就從她之後所說的話.
當耶穌告訴她.
你一直在等候的那個彌賽亞.
現在跟你在說話的那個就是她.
當這個婦人跟耶穌說完這番對話.
她發現耶穌知道她的羞恥.
知道她的過去.
她發現耶穌是一個猶太人.
不應該跟她有任何的對話.
還要跟一個女人有任何的對話的時候.
但是耶穌竟然願意這樣做.
並且她也因著耶穌所做的事.
確認了耶穌就是一個彌賽亞的身份.
她相信.
她相信這個經歷改變了她的生命.
由原先是一個連拿水這麼基本的生活需要.
她一定要拿水的.
她也要選擇在一個沒有人見到她的時間去做.
而變成她可以衝出去城外.
到處跟別人說.
你知不知道我遇到一個人.
他能夠將我過去的事都說出來.
其實可能很多人都知道她的事.
所以她說她能夠將我過去的事都說出來.
大家都知道她的事.
但是她不再害怕.
她不再羞恥.
因為她遇見一位又真又活.
並且主動來尋求她的主.
所以當她遇到耶穌以後.
不單止她被主耶穌那份愛和接納.
醫治了她的內心.

$^{841}$同樣的耶穌也改變了她.
改變了她對自己的看法.
甚至讓她成為一個要去傳福音的人.
當她走出去外面到處跟別人說.
她遇見一位可能是彌賽亞的人的時候.
其實她就是在傳福音.
她不再介意將自己的過去告訴別人.
原來我們的敬拜是可以這樣.
原來我們的敬拜.
當我們願意打開我們的心.
當我們願意在敬拜當中.
將我們的焦點對準神的時候.
我們願意被主捉摸的時候.
我們可以在敬拜當中來遇見祂.
但這不是一次的經歷.
其實這會在我們心靈裡不斷循環.
我們去下一個範疇.
當我們在恩著神的救恩.
我們來到敬拜讚美.
當我們在敬拜裡用心靈來按真理來讚美.
當我們在敬拜當中遇見神的時候.
我們的生命被改變被恢復的時候.
其實這個經歷會鼓勵我們更多來敬拜.
因為你品嘗敬拜當中能與神相遇的好處.
那份甘甜的時候.
這會吸引你繼續在神的話語.
在跟神的交往裡再次尋求祂.
這會是我們生命當中不斷的循環和不斷的需要.
因為我們的生命.
就算我們信了主有多久.
我們的生命其實是一個熟練的歷程.
有高有低 有起有跌.
就是每一刻我們感覺到我們有軟弱的時候.
我們特別的尋求神.
但當我們經歷了神的醫治.
神的捉摸 神的改變的時候.
我們重新有力量的時候.
我們也同時間繼續尋求神.
讓我們發現我們不可以沒有祂.
所以第一件事.

$^{881}$讓我們去學習敬拜的時候.
敬拜的時候我們去體會.
原來敬拜是神給我們一個很大的祝福.
是神給我們一個能夠直接跟祂相交.
能夠去遇見祂一個很寶貴的時刻.
敬拜並不是那二十分鐘.
敬拜是在說我們的生命.
我們的生命時時刻刻都可以去與主相連.
並且當我們願意打開我們的心.
當我們願意去揚升 去讚美祂的時候.
我們一次又一次.
透過詩歌將神的話語.
神的應許 神的恩典.
宣告 宣告我們自己的時候.
這些真理是可以改變我們的生命.
我們敬拜並不是因為.
神要我們讚祂.
神有多好 神有多厲害.
神有多奇妙 祂不知道嗎.
神當然知道.
我們每一個都是祂創造.
原來神要我們去敬拜 去讚美祂 去高舉祂.
不是因為祂的不安全.
是因為祂知道當我們這樣做的時候.
其實我們是不斷去提醒自己.
這位主有多愛我.
這位主對我有多重要.
我們一起起來吧.
我們也用這首詩歌來回應.
高天深海數樹.
你是配得轉為.
天空掘了高牆.
贊助你的永遠.
深知你所喜愛.
卻是以心靈獻奉給你.
成為你最欣賞喜悅的敬拜.
高天深海數樹.
你是配得轉為.
天空掘了高牆.
贊助你的永遠.

$^{921}$深知你所喜愛.
卻是以心靈獻奉給你.
成為你最欣賞喜悅的敬拜.
高天深海數樹.
你是配得轉為.
天空掘了高牆.
贊助你的永遠.
深知你所喜愛.
卻是以心靈獻奉給你.
成為你最欣賞喜悅的敬拜.
高天深海數樹.
你是配得轉為.
天空掘了高牆.
贊助你的永遠.
深知你所喜愛.
成為你最欣賞喜悅的敬拜.
主啊!你根本不需要我們給你什麼.
因為天上地下的萬物包括我們.
都是你親手所做.
然而你要的.
是我們願意打破我們自己.
來完全獻上.
打開我們的心靈.
來讚美.
來歌頌.
來救贈你.
是恩典的上帝.
從出世的一世.
分手掌管所有.
並沒有所缺少.
因此你最悲哀.
是我沒有保留.
賞心悲哀的.
才是你最欣賞喜悅的敬拜.
是恩典的上帝.
創造世間一切.
親手掌管所有.
並沒有所缺少.
因此你最悲哀.
是我沒有保留.

$^{961}$賞心悲哀的.
才是你最欣賞喜悅的敬拜.
因此你最悲哀.
是我沒有保留.
賞心悲哀的.
才是你最欣賞喜悅的敬拜.
因此你最悲哀.
是我沒有保留.
賞心悲哀的.
才是你最欣賞喜悅的敬拜.
\newpage



\section{創世記}
\label{sec:FaXJPQMXp6c}
\textbf{與人同在的神 \#1 - 創世記 1 - 20190929 主日講道 (陳崇基牧師)}
\newline
\newline
連結: \href{https://youtube.com/watch?v=FaXJPQMXp6c}{\texttt{https://youtube.com/watch?v=FaXJPQMXp6c}} ~~~~ 語音日期: 2019-10-05
\newline
\newline
\hyperref[sec:7VnGbL9RHgg]{\small{< < < PREV SERMON < < <}}
~
\hyperref[sec:index_chronic]{\small{[返順時目]}}
~
\hyperref[sec:index_scriptual]{\small{[返順卷目]}}
~
\hyperref[sec:LfJIKNZbdWE]{\small{> > > NEXT SERMON > > >}}
\newline
\newline
$^{1}$今天我們開始一個新的系列.
這個系列叫做.
與人同在的神.
我們可能第一次接觸聖經.
你不知道聖經是甚麼.
或者你可能信了一段時間.
你不知道聖經是甚麼.
你可能信了一段時間.
你不知道聖經是甚麼.
你可能信了一段時間.
或者你可能信了一段時間.
你到底問一下.
到底聖經這麼厚.
到底是一本怎樣的書呢.
所以這個系列是為.
如果你是第一次接觸信仰.
你不是很清楚.
有時候你聽到的時候.
聽到牧師說.
不知道在哪裡找那段經文.
你找都未找到他已經說完了.
我們這個系列.
我們會特別為你而設.
但如果你信了主很久.
你不要覺得這些只是基本功.
可能你都夢察察.
你裡面下載了不同的.
聖經的章節,金句.
或者一些故事.
你都不知道如何擺放這些東西.
大家彼此是如何連繫的.
所以這個系列我們會幫你走過.
一些聖經裡面.
由第一章到最後.
一些重要的段落.
讓你知道整個上帝的故事.
那個藍圖是怎樣的.
上帝在世界裡面在做什麼.
在人類的社會裡面在做什麼.
然後找出我們到底站在哪一個地方.

$^{41}$以致我們能夠明白.
神對我們的角色.
我想做一些新的東西.
我想做一些現場解答的環節.
如果你有電話的話.
你有QR Code的掃描器.
請你掃描它.
我們有張很簡單的上網的表格.
你可以隨時問問題.
你問了問題的時候.
就在我手機出現了.
我有時間的話.
最後的時候.
我們就會和你做一些解答.
我們想有些互動.
當然有答案的問題就會回答.
沒有答案的.
可能在門口和你聊一下.
但怎樣也好.
我們希望能夠有問有答.
所以請你掃描它.
或者你在我們Facebook群組那裡.
你也可以找到這條鏈接.
所以如果你有問題的話.
你隨時可以馬上問.
在我這裡彈出來的時候.
到底我有沒有時間回答.
就看我控制時間控制得怎樣.
但有問題的話.
我們都會累積下去.
下一次會再回答.
怎樣也好.
我想有好幾個方法去認識.
到底基督教信仰是一回事.
譬如我們有時候去一些聚會.
我們會聽一下別人的見證.
譬如藝人報道.
或者有些人的報道會.
我們去聽一下他信主的經歷等等.
或者我們可以去做一個study.

$^{81}$看看這個教會過去二千年.
裡面的教義有什麼發展等等.
我們都可以這樣做.
但有一樣東西不可以不做.
就是要去看一看.
基督教信仰一些所謂的foundation document.
一些很基本的文件.
這個一定不可以跳過這一步.
那個foundation document就是聖經.
聖經有六十六卷.
聖經的字是bible.
其實是book的意思.
聖經本身就是book of books.
裡面你打開目錄的時候.
你會找到有六十六本小一點的書.
這些書的長短各有不同.
有些可以去到幾十章.
有些只有一章這麼短.
一板紙這麼短.
六十六卷書.
三十九卷我們叫做舊約.
是耶穌基督來之前寫的.
二十七卷就是耶穌基督來之後所寫的.
你看到這些書卷的名字.
有些是地名.
有些是人名.
有些是書卷本身的主題的名字.
你看這些書的時候.
同樣你知道.
裡面所謂文學的題材.
type of literature.
genre.
是各有不同的.
有一些好像是詩歌.
詩詞歌賦.
有一些是歷史書.
有一些是我們現在的.
二十一世紀的人.
都不會再有的一些寫作的方法.
譬如最後那本《猶下國》.

$^{121}$就是啟示錄.
那些叫做天啟文學.
apocalyptic literature.
那些是很富有象徵性的語言.
意境寫出來的一本書.
所以你打開這本書的時候.
你會發現這本聖經.
裡面好像一本圖書館.
一個圖書館.
由第一本寫完.
到最後那本寫完.
大概中間有一千五百年左右的時間.
一路的搜集起來的時候.
就成為了這本所謂正典的聖經.
舊約三十九本.
新約有二十九本.
今天我們第一次開始的時候.
我想和大家看看.
或者仔細一點想一想.
到底第一本聖經.
第一章.
創世記第一章.
到底想講些什麼.
我先讀好嗎.
創世記第一章.
如果你有聖經的話.
我鼓勵你打開.
你的Bible App也好.
或者你有一本列印出來的聖經.
這個時候你拿出來.
我們PowerPoint當然會打出來.
創世記第一章第一節.
他說起初神創造天地.
地是空虛混沌.
與棉行.
靈運行在水面上.
神說要有光.
就有了光.
上帝看光是好的.
就把光暗分開了.

$^{161}$上帝稱光為晝.
稱暗為夜.
有晚上.
有早晨.
這是頭一天.
當我們打開第一章的時候.
這幾句話.
我相信我們需要做一點說明.
為什麼呢.
二十一世紀的人.
我們拿著手機.
有QR Code.
科技很發達的人.
我們很多時候.
立刻會想起.
很多的爭論.
就是宗教信仰和科學之間.
到底兩個是不是像婆媳的關係.
永遠都是爭吵死的.
因為我們二十一世紀現代的人.
我們是預設了.
或者心裡面.
我們在想一件事.
就是所謂現代的科學.
和聖經一本這麼古舊的書.
一定是不合的.
一定是彼此的對立.
但是我想告訴你.
很多時候我看一些.
特別是非信徒.
不是信主的.
不信聖經的.
特別是科學家.
他們寫一些書.
是講到宇宙的浩瀚.
講到宇宙有多大的時候.
很多時候.
最後的時候.
我發現他們用一些語言.
或者一些語言.

$^{201}$幾乎去到我們基督教教會裡面.
唱詩歌那種語言.
一些worshipful的語言.
一些讚嘆的語言.
一些去到覺得很奇妙的語言.
這本書我記得我曾經提過.
叫做The Six Numbers.
這位科學家叫Martin Rees.
他說.
這本書很薄的.
不是很難看.
就算你沒有什麼科學背景.
好像我這樣.
你都應該可以看得明白.
他說.
這個整個宇宙.
原來是靠賴著六個數字.
六個常數constance.
這些constance不可以太大.
也不可以太小.
只要稍微有一點點轉變.
今天我們所有的生活.
我們有的宇宙.
四季氣候等等.
是不會產生出來的.
我不知道怎樣去.
我不知道這樣是不是恰當.
是不是可以用這六個常數.
來去證明.
去證明出.
上帝是存在的.
可不可以用這六個常數.
這麼奇妙的常數.
連作者都覺得很奇妙.
連學者都覺得很奇妙.
差一點點.
大一點點.
或者小一點點.
我們是不能夠生活的.
一個這麼fine-tune universe.

$^{241}$一個fine-tune的宇宙.
可不可以用這六個常數.
然後我們就說.
一步一步用邏輯.
用哲理.
來去有一個conclusion.
有一個結論.
然後我們就說.
上帝創造主是存在的.
我不知道這樣是不是很恰當.
為什麼呢.
因為我相信.
如果上帝真是上帝的話.
神真是神.
創造主本身.
就不應該是我們有限的邏輯.
來去辯證出來的結論.
如果神真是神的話.
就不應該是我們的邏輯.
辯論的步驟.
最後的conclusion.
如果神真是神的話.
應該是相反的.
上帝應該是我們一切的starting point.
是開始.
不是我們想法的結論.
正如聖經第一句所說.
起初上帝創造天地.
由神開始.
似乎聖經第一章第二章.
說到創世的時候.
就不是去解答科學的問題.
而是解答一些人生.
宇宙觀或世界觀的問題.
什麼叫世界觀呢.
世界觀我們每個人都有.
無論你是刻意地表明出來.
或者你背後有一點價值.
世界觀需要問什麼問題呢.
這本書不是很新的.

$^{281}$九幾年寫的.
但他寫了一些很有意思的東西.
這本書叫做.
Truth is stranger than it used to be.
他就說到world view到底是什麼.
世界觀人生觀是什麼.
所以你有沒有人生觀.
你有沒有世界觀.
請你點頭.
有的.
只不過你未必能夠用語言表達出來.
一個世界觀一個人生觀.
怎樣看世界宇宙和人生.
其實是回答五條問題.
這五條問題就是.
我在哪裡.
Where are we.
Where am I.
我是誰.
或者我們是誰.
Who are we.
第三.
世界有什麼問題.
今天我們看見網上的新聞等等.
一定我們會留意到.
Something is wrong.
當然有問題.
問題是什麼.
第四.
對於這個世界的問題.
有沒有終極的答案.
What is the solution.
The ultimate solution.
第五個世界觀要回答的問題.
What time is it.
現在是什麼時間.
我們處於這個世界宇宙歷史裡面.
是哪一點.
What time is it.
我們到底在哪裡.

$^{321}$我們人人都會回答這些問題.
無論你是很能夠articulate.
很能夠表明出來.
或者你心裡面其實都會有一點假設.
這個就是世界觀.
創世記頭幾章似乎就是要回答這些問題.
不是去講科學.
因為你看看創世記第一章.
你打開的時候.
你立刻應該領會到.
他用的用詞和語句.
不是科學報告.
起初上帝創造天地.
現在我們再讀一次.
我想用一點imagination.
當你聽到這些說話的時候.
聽到這些言辭的時候.
你會想起什麼圖畫.
起初神創造天地.
地是空虛混沌.
淵綿黑暗.
神的靈蘊行在水面上.
神說要有光就有了光.
神看光是好地.
就把光分開.
神稱光為晝.
稱暗為夜.
有晚上有早晨.
這是頭一天.
你想到什麼畫面.
地是有光出來.
還有呢.
有水.
有誰想到有個地球.
在黑暗裡面出現.
總之一塊好像je 喱一樣.
地是空虛混沌.
有誰想到這件事.
大概.
我想我們大部份的人都會這樣想.

$^{361}$大家想想.
我們想起地是空虛混沌.
想起創造天地.
淵綿黑暗.
有些水.
有東西在那裡.
空虛混沌.
然後有光出現.
我們自然會想起太空.
太空看到地球的一些景象.
但我想告訴你.
這種景象.
到底歷史有多悠久.
其實不是很久.
因為60年代開始.
才有電視可以看到.
相片可以看到.
就是一個這樣的畫面.
對嗎.
這是一個很近很近.
相對來說其實是很近的現象.
但你知道聖經.
是3000年之前寫的.
當時他們沒有火箭.
沒有穿梭機.
或者沒有太空的探索.
他們想起上帝創造天地的時候.
他們會想起什麼呢.
一定不會是這個景象.
一定不會是我們想的大爆炸等等.
因為上帝要和人溝通的時候.
就如他創世記這麼多年前.
要與人溝通的時候.
是要用當時的人明白的語言.
或者一些觀念.
來與他們溝通.
否則就不可以溝通.
所以與以前的人溝通的時候.
你知道聖經是當時寫的.
甚至不是用現代語言寫的.

$^{401}$是用古希伯來文.
古希伯來文來寫的.
對於第一代的作者來說.
他們不會想起太空.
他們會想起什麼呢.
上帝創造天地的時候.
他們會想起什麼呢.
其實就是想起一樣.
天空和地下.
上帝創造天地.
踩著的地.
這張照片漂亮嗎.
漂亮嗎.
你知道是誰拍的嗎.
是我拍的.
我去加利利拍的.
加利利有一座很高的山.
上了山上面拍的.
上帝創造天地就是天空.
藍色的天空.
好像有水的天空.
還有土地.
Land.
Farmland.
所以你留意一下.
那個言詞的時候.
你就知道.
回到第一代的讀者.
他們是怎樣了解.
其實是一些很普通人的經歷.
有晚上有早晨.
換句話說.
有Sunset.
有Sunrise.
他們不會想到.
有個地球圍著太陽公轉.
等等這些科學的問題.
他們只知道太陽從東邊升起.
然後西邊下山.
然後有月亮出來.

$^{441}$黑暗裡面會有光.
那個光是什麼光呢.
不是我們想的那些光子.
就是太陽光.
也不是這些燈泡的光.
就是每天早上有太陽.
他們能夠生活.
能夠工作.
日落的時候他們就作息.
起初.
神創造天地.
起初也不是我們想.
一定是大爆炸那一點.
十幾個Billion years ago.
有那一點.
那一點是很Dense的那一點.
然後爆出來.
我們現代人就會這樣想.
起初不是說時間裡面的一點.
起初原來這個字.
希伯來聖經或舊約聖經第一個字.
起初.
其實是形容一個王.
一個君王.
他沒有登基儀式之前.
他已經上了去管理一個國度.
就是那段時間.
由他開始坐立王之後.
到他正式有登基儀式的那段時間.
列王記也是經常出現的字.
就是說他去登基之後.
他去管理這個國度.
所以那段起初的時間.
上帝去管理這個世界的時間.
神創造天地.
地是空虛雲遁.
冤冤黑黑.
也不是我們想著.
當然是怎樣大爆炸.
當時物理的定律還沒有.

$^{481}$那時候就好像je 喱一樣.
或者是很混亂的.
空虛雲遁是什麼呢.
空虛雲遁是.
形容曠野的.
是一個地理.
以色列人很清楚的地理的一個語句.
希伯來文你應該試一下讀.
你就會知道為什麼是空虛雲遁.
因為是Tofu Fafou.
你試一下讀Tofu Fafou.
你都不知道說什麼.
就是這樣.
就是說曠野.
沙漠這個地方是不能夠居住的.
Tofu Fafou.
是空虛雲遁就是這個意思.
不能夠定居.
不能夠成家的一個地方.
居住在以色列的人.
他們很清楚什麼叫做Tofu Fafou.
又很清楚.
在這個世界裡面.
有Tofu Fafou.
也有水.
深海的水.
Ocean.
試探索.
不能夠探索.
不能夠在那裡居住.
為什麼呢.
因為這兩個地理的terms.
就正是在迦林地裡面的人.
遇見的兩個景象.
在他們東邊就是北邊的Arabian Desert.
在他們西邊就是地中海.
一邊是Tofu Fafou.
一邊是空虛雲遁.
是不能夠居住.
在他們另一邊就是一個深海.

$^{521}$可不可以在海上住.
以前是不行的.
有誰是水上人.
在船上長大出生的那些.
現在可以.
但以前是不行的.
他們剛剛就在中間.
他們能夠居住.
可以定居.
是誰make it possible.
是創造主make it possible.
在一邊地是空虛雲遁.
淵面黑暗.
另一邊有深的水.
誰make it possible.
能夠在中間這個地方.
我們可以耕種.
可以定居.
可以作為一個家.
誰可以呢.
神的靈運行在水面上.
希伯來文是一些很有意味的語言.
剛才我說過Tofu Fafou.
你聽到聲音就知道很混沌.
很不妥.
神的靈是什麼意思.
神的靈其實是神的風的意思.
神的氣的意思.
希伯來文怎麼讀呢.
就是Wuha.
你試一下.
Wuha.
你就一定感覺到有氣.
確實就是這樣.
Wuha.
所以要有上帝的靈.
上帝的氣.
祂說的話.
才能夠make life possible.
人才能夠有生命.

$^{561}$我們才能夠有家.
有生活.
有生命的氣息.
所以第一件事.
之後上帝所做的.
就是上帝說.
上帝就說話.
神的Wuha運行在水面上.
神用祂的說話.
祂的道.
來創造這個世界.
要有光就要有光.
要有什麼就要有什麼.
上帝用祂的呼氣.
使這個世界能夠有生命.
不再是空虛雲遁.
不再是冤冤黑暗的水.
在創世的第一章.
出現過很多次上帝說.
總共有十次.
上帝說了十次說話.
就令到這個物質的世界存在.
並且讓人能夠居住.
你要知道上帝.
我們之後會說.
上帝再說十次說話的時候.
是什麼時候.
是西大山.
十屆.
十屆原文是ten words.
是上帝說出生命的原則.
十句的說話.
就是給我們生命.
上帝的道就是給我們生命.
所以上帝創造這個世界.
用祂的說話來創造.
並且在第一章裡面.
我說過了.
不是一個科學報告.
而是一個平民百姓的語言.

$^{601}$並且是第一章很有趣的平衡.
第一天和第四天是平衡的.
第二天所做的事情.
是一個天與海.
然後第五天就是做天與海裡面的活物.
第三天是做地出來.
有乾地出來.
第六天就是做乾地上面的活物.
然後最後最後最後最頂峰的.
當然是最精彩的.
就是做人類.
第七天上帝就rest.
休息不是因為他累.
上帝安息其實就是在他的聖殿.
在他的寶座上.
他能夠登基.
他真的作王.
開始管理.
去統治這個世界.
所以不是科學的報告.
是很有詩意的.
等於你去博物館.
看那些很偉大的藝術家畫的畫.
你可以說他的油畫.
他的油彩有什麼化學成分等等.
其實你錯過了整個重點.
如果你只是談這件事.
你就要站起來.
然後你就看見一幅很大的油畫.
哇,這個真的很偉大.
我剛剛從加州回來.
兩個星期前.
有誰在加州來過?.
去過加州的?.
去過一個地方叫做Yosemite.
休山美地.
有沒有去過?.
你人生未死之前.
應該一定要去一次.
休山美地.

$^{641}$你開車上到那些遼望的地方.
你就會看見.
這個真的很偉大的創作.
很優美,很浩瀚的一個創作.
《創世記》第一章.
其實就是用一個這樣的藝術語言.
去告訴我們.
並且第一章裡面有什麼呢?.
是經常有七件事出現.
所以我知道有七天.
但是知道第一章第一節.
「死初神創造天地」.
是有多少字?.
中文有七個字.
原來希伯來文都是有七個字.
第一節有七個字.
第二節呢?.
「地球空虛運動」等等.
十四個字.
還有什麼七呢?.
第一章第一節.
你要數一下那些字母.
因為希伯來文的語言.
是Alphabetical.
好像英文一樣.
是要串出來的.
有多少個字母呢?.
28個字母.
即是七乘四.
第二章第一節到第三節.
即是最後第七天的時候.
上帝安息統治的時候.
到底有多少個字呢?.
35個字.
為什麼35個字呢?.
因為七乘五就是35.
「神」這個字出現了多少次呢?.
第一章.
35次.
「天」和「地」.

$^{681}$Heavens and Earth.
出現了21次.
「事就這樣成了」出現了7次.
還有一句話.
就是「神看是好的」.
你猜猜有多少次呢?.
一定是和七有關.
就是七次.
所以是一個很有詩意.
很精密的文件.
還有就是古人.
逢甄說到建造一些東西.
用了七天,七個星期,七個月.
或者七年的時候.
他們在說一些很特別的創造.
只有某一些的建築物.
才會用七年,七月,七星期,七日做的人.
就是聖殿.
所羅門王.
之後的故事.
做聖殿做了七年.
換句話說上帝創造一切.
不單只是功能是好.
而且整個宇宙是神聖的.
我們走在大自然裡面.
走在這個世界裡面.
看到天,看到地.
古代的人也好,現代人也好.
我們走在聖所當中.
這個世界.
《創世記》第一章說是一個神聖的世界.
是在敬拜上帝.
是在讚美上帝.
神創造最高峰當然是26節.
26節說到神創造人.
神說我們要照著我們的形象.
照著我們的樣式去做人.
使他們管理海裡面的魚,空中的飛鳥.
地上的牲畜和全地上的一切昆蟲.
神就照著自己的形象去做人.

$^{721}$那是照著他的形象去做男,做女.
之事就這樣成了.
上帝看一切都是好的.
有晚上,有早晨.
這是第六天.
一切都是美好的.
我想加一句.
連我們人體都是美好.
包括我們的性別也是美好的.
有男有女.
這是上帝自己精密的創造.
我知道現在社會的風氣已經不同以前.
我最近和一個小學老師聊天.
他說很麻煩.
他在加州教書.
那間學校要我們去跟小朋友讀故事書.
那些故事書說到男女不是生出來就是男女.
是你會變成男或女.
你生出來的時候.
性別是不知道的.
是你自己選擇.
你變成女性.
你不是變成女性.
你明白我的意思嗎?.
當然他是信主的.
連同事都覺得有問題.
未信主都覺得有問題.
今時不同往日.
所以我們需要去回應這件事.
創世第一章就是講這件事.
我們的性別其實都是一個美好的創造.
人生出來你是男或女.
是值得去欣賞的.
因為這是上帝的美好的創造.
你想想小時候.
怎樣可以讓他選擇性別.
他連大學都未選擇.
將來我去哪一間.
他都會很迷惘.
開車轉左轉右他都未知.

$^{761}$怎樣去選擇性別呢?.
I don't know.
最近在一個基督教作者寫了一本書.
叫做God Made Boys and Girls.
所謂第一本.
用聖經的角度去告訴小朋友.
你是boy or girl.
不是因為你喜歡跳舞所以一定是girl.
不是因為你喜歡修車一定是男.
不是這個意思.
而是你的性別就是上帝的創造.
一本很好的書我買了.
我跟太太說.
我出了這本新書.
今年九月頭出的.
所以新鮮在媒體上出現.
我買了回來先review.
我聽過一次我自己看一次.
很好.
我覺得我自己也想看.
我是一個man.
good man.
你是woman.
good woman.
上帝創造.
這本書有講DNA等等.
很適合小朋友來看.
我知道你心裡在想什麼.
如果你是父母特別.
你就說怎樣買到這本書.
我想告訴你.
我已經會幫你order.
等下完結後請你去後面.
有個sign up書15元一本.
你去amazon買15個交易扣.
五年shipping.
還有為什麼不去amazon買呢?.
因為我特意去new growth press買.
那些小的出版社.
我們應該去support.

$^{801}$15元收便宜.
人工也不算.
我會跟你一批訂回來.
如果你有興趣的話.
今天請你去後面.
你去sign up.
不單止你自己看.
也可以讀給你的小朋友.
由今天開始.
我們的children ministry.
因為它是一個series的書.
其他東西要說到respect our body等等.
我們會用這些story book.
去讀給我們的story land.
那些小朋友.
告訴他們.
起碼我們要先懂得appreciate.
我們的gender.
上帝創造是美好.
並且是神聖.
有一次我聽一個人.
講他信主的見證.
一個人叫Francis Collins.
Francis Collins是世界最頂尖的美國基因學家.
你說那個基因的藍圖.
怎麼map out出來呢?.
他就是那個人.
他本身是醫生.
有一次他去探病的時候.
他見到一個垂死的病人.
他見到這個病人很平安.
快要面對生命最後一日.
他見到很平安.
他跟他聊天的時候.
那個女病人婆婆問他.
她說我知道我去哪裡.
他問醫生.
你知不知道你去哪裡.
然後Francis Collins突然醒了.
他自己都不懂回答.

$^{841}$他回到醫院的時候.
他就看了一本書.
如果你有看過一些classic Christian的書.
有一本叫做Mere Christian Entity.
是派給C.S. Lewis的.
他看了之後就相信.
這個宇宙不是我們單單捉到摸到那麼簡單.
這個只是現實裡面其中一半.
原來還有一個屬靈的境界.
那個就是聖經所打開的境界.
所以他信了.
他也是一個很頂尖的一個科學家.
他說這個世界上帝就是創造基因的上帝.
我們在大教堂可以來去敬拜他.
但是在實驗室,大自然.
在奇妙的創造裡面都可以敬拜他.
是很複雜並且是絕美的.
聖經第一章其實就是打開這個真理.
叫我們相信這位造物主.
我們知道當然有些東西.
聖經本身說是美好的.
但是我們知道有很多問題.
這個就是聖經的藍圖想打開.
就是因為有罪有邪惡進入了世界.
所以現在有暴力有不義等等很多混亂的東西出現.
而聖經不是停在創世記第一章.
最後基督來.
耶穌基督來拯救我們離開罪和死亡.
現在我們站在的地方就是基督第一次來.
使我們有新天新地一個完美的宇宙.
再去到更新的時候.
我們就站在這裡.
我們不是在等待上天堂.
而是我們將天堂帶來這裡.
並且我們盼望將來新天新地會到.
這個就是基督徒有的角色.
我們的定位就是在這裡.
我結束的時候想說上星期我探望了一個姐妹.
星期四探望了她.
她病得很嚴重.

$^{881}$因為她過去這麼多年都是因為一種很奇怪的基因病.
令她受了很多苦.
但每一次她好的時候.
她都去教會敬拜神.
去聽道 去主教 去茶經.
每一次她不行的時候她就會去醫院.
星期四探望了她之後.
星期五她就離開了世界.
我陪著她的家人.
我這麼大個人.
我未試過坐在一個已經離開了的人的軀體.
坐在旁邊坐了幾個小時.
你想想如果我們沒有將來 沒有新天新地的話.
人生是多麼荒謬.
但這位姐妹一向都很相信.
正如我所說她好的時候就會去敬拜神.
沒有藉口的.
因為她很相信將來會有新天新地.
這個盼望是上帝給我們的.
這個也是我們有的世界觀.
正是因為我們有這樣的世界觀.
所以我們就能夠在今時今日都去慰注 作見證.
告訴別人這是何等好的一個信仰.
這位神是何等好的一位神.
我們結束低頭祈禱.
然後我們才回答一些問題.
上主我們多謝你.
因為你賜給我們原創的天地.
是何等美好絕美的一個世界.
我們的生命也是你所賜給我們的.
主要是多謝你 你愛我們.
神我們盼望將來新天新地.
同樣我們今天見到有很多問題.
我們願意代表你這個角度.
你都去說話 行事.
求你讓我們教會.
我們整個群體能夠去見證你.
讓我們去盼望 去飲祈禱 奉主明教.
這本書我太太在小朋友的班讀完.
如果你想要的話請你去sign up 付款.

$^{921}$下星期你可以pick up.
我想回答一個問題.
這裡有兩三個問題.
第一個問題是.
基督教可不可以說髒話.
因為時間關係我只回答一句.
我發現有很多說話比髒話更難聽.
可能說得很斯文.
我聽過很多不公義 打壓的說話.
比髒話更難聽.
下次再說 基督教可不可以說髒話.
無論如何我們會繼續run這個story.
我希望上帝祝福各位.
謝謝大家.
\newpage



\section{}
\label{sec:LfJIKNZbdWE}
\textbf{10/06/19 主日講道: 參透生命的智慧 (吳君才牧師)}
\newline
\newline
連結: \href{https://youtube.com/watch?v=LfJIKNZbdWE}{\texttt{https://youtube.com/watch?v=LfJIKNZbdWE}} ~~~~ 語音日期: 2019-10-10
\newline
\newline
\hyperref[sec:FaXJPQMXp6c]{\small{< < < PREV SERMON < < <}}
~
\hyperref[sec:index_chronic]{\small{[返順時目]}}
~
\hyperref[sec:index_scriptual]{\small{[返順卷目]}}
~
\hyperref[sec:1bmqQECjSXg]{\small{> > > NEXT SERMON > > >}}
\newline
\newline
$^{1}$(香港政局).
(香港政府).
各位朋友.
這個世界是否很不平安.
特別你關心香港政局的時候.
你會發現.
真的很憂心.
亦都很痛心.
亦都很傷心.
不過蒙面法.
反蒙面法通過之後.
你會看到更多貧乏的局面.
有人說將會如何發展.
我們並不知道.
但是將會如何.
我們真的需要直接為.
這個香港局勢保護圖告.
要求主幫助那些暴力.
或者那些仇恨.
不會惡法.
我們相信一樣東西.
很多東西我們不知道.
未來我們都不知道.
但我們知道萬事萬物都有平息.
掌握在手裡.
神坐在天上為我們.
大人一層.
長老一層的人.
所以我們真的要有參透人心.
細思智慧.
有一個這樣的事.
今日的崇拜.
現在剩下時間有限.
其實我那邊講章都幾長.
現在要跳了.
我們可以看兩樣.
或者三樣.
以及四樣.
看時間.
一會兒我講成怎樣.

$^{41}$我們會用全圖書.
三章一到十五節.
作為今日的講題經文.
大家給前面來.
我們先簡單一個祈禱.
天父讓我們在你的面前安靜.
等候你自己的話語.
給我們得著你自己的亮光.
幫助我們在你的面前.
有一個倚靠和信服的心.
帶領我們以下一些.
查考你華語聽到的時間.
奉主耶穌基督的名禱.
阿們.
全圖書在一卷格言集.
或者言語錄.
和《熟智慧文學》的針言很相似.
裡面除了講到一些作者的自述.
和反省人生之外.
其實裡面有很多格言.
訓誨,忠告,語及,吩咐,禁令等等.
全圖書你經常看的時候.
會出現多次的.
虛空的虛空.
歸納起來全圖書講了七方面的虛空.
由第一章到十二章.
第一是努力和享樂的虛空.
第二是生命有缺陷和限制的虛空.
第三是世上不公平的虛空.
第四是追求財富的虛空.
第五是惡人不一定即時得到懲罰的虛空.
第六是人命運類同的虛空.
第七是生命短暫的虛空.
每一樣都可以一大篇講章.
但是全圖書作者在結束這七方面的虛空的時候.
是教導讀者每一個人應該怎樣去參透和超越這個虛空.
以追求更有意義和快樂的生命.
所以時間所限我提出幾點.
第一.
我們要接受生命有缺陷.

$^{81}$時間關係我們不讀.
但是你看到這些.
傳道者全部的心意要提醒鼓勵我們一件事.
就是要接受我們的生命是有缺陷的.
他在頭八節用了十四對的對比.
二十八個時.
他指出人生遭遇的都是按既定的時候發生的.
十四和二十八是七的倍數.
七在聖經來說是代表完全的數字.
對比也是希伯來文學用來表達所有的文學技巧.
所以這十四對對比是暗示人生所有的一切經歷.
都有固定的時間限制.
人是很無奈的.
所以傳道者用這個對比描述人生的經歷.
他反映出另一個很重要的真理.
這個真理就是.
生命好像一個錢幣的兩面.
不是永遠是正面的也有負面.
不會永遠是積極的也有消極的.
原來生命不是我們所要求並且盼望的那麼完美.
是有缺陷的.
生命不是永遠是生的也會有死.
不會永遠是栽種的也會有拔出所栽種的.
不會永遠有哭也會有笑.
不會永遠是哀悼也會有跳舞.
不會永遠是靜默也會有語言.
不會永遠是爭戰也會有和好.
第一個對比是生死.
那裡是你和我每一個人必須經歷.
而最後一個對比就是爭戰和和好.
卻是社團或者國家相關的.
但是更重要的是這十四對的對比.
表明原來我們人生的生命是會有缺陷的.
並不是完美的.
我們有時候真的會歡笑.
但是有時候我們也會流淚.
可惜很多人誤會了基督教的信仰.
以為信耶穌之後就一帆風順.
天色常藍無風無險又過十二點.
但是沒錯有歡笑.

$^{121}$但也有眼淚.
不是嗎?很多時候.
你和我聽過類似的話.
快快信耶穌.
信耶穌實在是好的.
信了祂神會幫你的.
沒有問題的.
有些人說耶穌很厲害.
有求必應.
這些是「成功神學」.
「付諸神學」.
「信理神學」.
這樣的信仰或者信念.
只是說到信仰的一部分.
並不是全部.
這樣的信念信仰.
有一天當人面對一些流淚傷心挫敗艱苦的事.
他就會產生懷疑.
會埋怨甚至放棄.
事實上不少人放棄信仰.
是因為一些自以為不會發生的事.
竟然臨到.
根本聖經是.
未許過信主的人.
永遠會一帆風順.
不只是有歡笑.
沒有眼淚.
聖經所有的偉人包括主耶穌.
他都是道成肉身之後.
經歷世間各樣的勞苦受煩.
有歡笑.
但亦都經歷很多眼淚苦楚.
但你和我有一個盼望.
就是在新天新地的時候.
世人不會再有悲哀.
那裡說沒有悲哀哭號.
那時候信徒的生命才是完全沒有缺陷.
真的.
今天你和我在這個有限的時空裡面.
面對這個變幻無常的世界世事.

$^{161}$我們只是像聖經說的.
勞苦受煩.
勞碌和受經煉.
但如果你和我能夠接受.
生命裡面不圓滿的.
是會有缺陷的這個現實.
我們就會學到一樣的功課.
就是當遇到痛苦不幸.
逆境難處的時候.
我們不會那麼容易沮喪.
抑鬱.
自閉.
埋怨或者咒詛.
很多時候我們不快樂.
我們不喜樂.
我們對生命或者人生.
會不會有太多太多不切實際的期望和要求呢?.
我們有時太多的執著.
太多的所謂理想意願期望.
要求.
以為人生一切是按你自己的藍圖和期望出現的.
以為一生不會遇到不開心的事.
不會逆自己的意願的事.
或者因為自己根本不能夠接受不完美.
以及生命中有缺陷的人或者事.
而這個張力在這裡.
有時我們又不知不覺以為基督徒.
信了耶穌.
念受受苦.
沒有缺陷的這個特權.
其實這個不是聖經.
解釋.
這個是你的錯覺或者偏差.
我們看第二樣.
要承認自己不是全知.
我們一起讀一讀這一節.
神做萬物.
各按其時成為美好.
又將永生安置在世人心裡.
然而神從始至終的作為.

$^{201}$人不能參透.
無疑從《創世記》一章所顯示.
神所做的一切都是美好的.
甚至神將永生安置在人的心裡.
不過結果.
神始終的作為.
人卻不能參透.
永生你看到這個詞.
其實可以解作隱藏.
也可以解作永恆的意念.
兩個解釋都配合上下文.
如果你解作隱藏.
就是指人心裡因為.
被神所放的奧秘.
遮蓋了.
不能夠明白神從始至終的作為.
如果你接納永生或者永恆的觀念.
這個說法.
即是說.
人雖然有永恆的觀念.
但仍然為時間所限拘禁.
好像坐井觀天一樣.
不能夠完全明白.
神從始至終的作為.
這個正如保羅所說的.
我如今所知道的有限一樣.
只有獨一的神.
從始至終.
由頭到尾.
祂可以一目了然.
參透一切.
傳道者在這裡強調.
你和我是人.
你不是神.
所以你一定知道的東西是有限的.
你不能夠像神那麼厲害.
全知全能.
有這樣的屬性.
人是被造之物.
不是創造主.

$^{241}$所以你不可能甚麼都知道.
甚麼都懂得.
人沒有可能參透明白.
世事一切一切.
就算今天你目前所發生的事.
所面對的困難甚至不幸.
為甚麼你會問.
偏偏選中我.
發生在我自己或我家人身上呢?.
就算你怎樣證明.
怎樣去分析.
總是想不通.
或猜不透.
因為人始終不是全知.
怎樣能夠明白一切呢?.
我記得我曾經探望一個家庭.
他們未信主.
他們拜菩薩觀音.
丈夫剛失業.
太太跌傷右手.
兒子失業半年.
女兒患長期病患.
需要長期的照料治療.
當晚我探望他的時候.
做爸爸的大男人.
在兒女面前.
在我面前.
掩不住真情的壓抑.
流出眼淚.
他跟他的妻子問我.
為甚麼我們這麼倒楣.
我這麼倒楣.
我一家這麼不幸.
我回答說.
我真的不知道.
但我知道.
神愛你們.
關心你們.
神特別會照顧.
困苦的人.

$^{281}$當晚.
我邀請他.
接受耶穌.
這個未信的家庭.
為他一起祈禱.
求主帶領和顧念.
頂尖會今時今日.
當人的年紀越大的時候.
所經歷的事越多的時候.
痛苦的事.
不幸的事.
又經歷多的時候.
你明白到原來生命呢.
有很多很多的事.
是你作為人.
不可以完全操控.
是不可以完全明白透的.
你作為人.
你並沒有全知全然明白一切的權利.
如果我們遇到困難.
不明白的事的時候.
你要對自己說.
我是有限制的.
我所知的事是有限制.
我知識有限制.
你可以將這個焦點.
由問為甚麼.
而轉移到在這樣的情況下.
神啊.
我應該怎樣做.
我應該做甚麼呢.
這個會更加務實.
更具體幫助我們去面對人生當前的問題.
第三點.
我們要分享的.
要立志終身喜樂行善.
如果傳道者說.
剛才上述兩個事實就停.
我們就會覺得.
人生真是太灰暗.

$^{321}$很無奈.
不值得繼續活下去.
傳道者在這裡並不只是鼓勵人.
接受上述這兩個事實.
他繼續在十二十三節說.
他說我知道世人莫強於終身喜樂行善.
並且人人吃喝.
在他一切勞碌中享福.
這也是神的恩賜.
沒錯.
傳道者承認生命裡面有缺陷.
和知識有限制.
但是這個並不等於說.
我們以後就要永遠憂愁下去.
我們是全無快樂可言的.
相反的.
傳道者指出.
我們不單止要接受這兩件事實.
或者是殘忍的事實.
但是我們可以.
而且應該超越他們.
而為此快樂.
甚至終身喜樂行善.
行善在這裡並不是指.
我們一般人所說的.
修橋補路.
做一些好的倫理道德的善行.
善在這裡和十三節所說的福字.
是同一個字.
行善的意思在這裡.
是指製造開心快樂.
享受生命.
即是主動追求美好和快樂的生命.
呂鎮忠譯本將終身喜樂行善翻譯作.
是常歡喜一生享樂.
傳道者認為人生很重要.
很美好的事.
所以他說.
我知道世人莫強於終身喜樂行善.
莫強於圓文直譯.

$^{361}$即是說沒有比這件事更好.
傳道者認為人最好的.
莫過於你終身快快樂樂.
如何可以做到時常歡喜一生享福呢?.
他提出三個要點.
第一就是勞碌和享受勞碌.
幫我們弄出來.
第十三節那裡傳道者說.
人人吃喝在他一切勞碌中享福.
當然這不是指我們糊里糊塗的去享樂.
荒唐頹廢的只顧吃喝玩樂.
什麼都不做.
不是因為他提出勞碌.
勞碌這個意思正是傳道者意思.
我們要努力工作的意思.
並且享受辛勞工作的過程和果效的意思.
傳道者第二章二十四節同樣有說.
人莫強於吃喝.
且在勞碌中享福.
我看這也是出於神的手.
意思是神在人所從事的工作上.
得著滿足和福樂.
諾貝爾物理學得獎者當中也有華人.
其中你聽過的.
是已故的前任香港中文大學的高崑校長.
他被稱為剛簽之父.
他早期發現這件事.
但後期才辭去的罪刑.
而另一個諾貝爾化學得獎的人是華人.
李遠哲博士.
他在1998年與香港浸會大學演講時說.
我們科學家事實上是費寢忘餐.
通宵達旦工作.
因為我們被毛主席的好奇心驅使.
落在迷堆的世界裡.
我們是辛勞不息.
享受勞碌的過程.
由另一位在1998年.
諾貝爾物理學得主崔奇.
他與大學生分享他的心得.

$^{401}$他說研究難就難在不知道甚麼是問題.
到了你發現問題後.
也不知道應該怎樣問怎樣做.
因此他向大學生建議.
你們千萬不要放棄.
不要氣餒.
他又說研究好像探索.
好像探險.
欣賞就好像觀光.
無需用力.
但研究就要用力.
雖然研究要出力用力.
但每一樣東西.
崔奇覺得要享受當中的過程.
研究會做到覺得是好玩的東西.
好像去樂園玩一樣.
崔奇這番說話代表了.
他享受他的勞碌.
他的心態和觀念.
正如愛迪生所說.
我一生從未工作過.
我只是每天享受自己喜歡做的事.
一樣.
這個只是享受勞碌.
我不知道弟兄姊妹你今天這麼勞碌.
你的工作甚至家務.
有沒有這樣的心態呢?.
第二.
我們要活在此時此刻.
傳道者在第三章一開始就用十四句對比.
來表明凡事都有定期.
天下萬物都有定時.
這是人面對的限制.
但我們可以借用這個對比.
參透人生的另一個角度.
那個角度是什麼?.
我們要活在此時此刻.
意思是在災中的時候專心災中.
拔出的時候專心拔出.
拆毀的時候專心拆毀.

$^{441}$建造的時候專心建造.
哀動的時候.
你就好好哀動.
全人的投入.
跳舞的時候.
你就全人投入去跳.
活在此時此刻.
換句話說.
你享受現在這一刻.
英文說here and now.
其實這個態度就是說.
你不會擔心將來.
也不會為了過去而後悔.
如果你有看列王記上十七章記載.
以利亞先知去到西頓.
撒勒伐窮寡婦家的時候.
窮寡婦窮到很難過.
她沒有餅.
只有潭內有一些麵粉.
瓶內有一點點油.
不過只可以供應她一兩餐.
吃完就等死.
但是以利亞先知告訴她.
你日之如何力量也必如何.
你要窮寡婦.
你學信心的功課.
活在此時此刻.
只要面對現在.
不要擔心將來.
抓住神.
所以她要將一個經歷.
不要怕.
以利亞先知說.
你先為我做一個小餅.
我吃了之後.
然後為你和你的兒子做餅.
因為神說.
潭內的麵不會減少.
瓶內的油不會短缺.
窮寡婦真的這樣做.

$^{481}$憑信心活在此時此刻.
照先知的話去做.
結果奇蹟發生.
潭內的麵沒有減少.
瓶內的油沒有短缺.
正如神先知所說.
窮寡婦用這樣的方式.
食用了很多油.
神供應了很多油.
請我們不要為明天憂慮.
一天的難處.
一天當就夠.
要活在此時此刻.
全心去依靠神.
相信主.
祂是大能.
有慈愛的.
是信實的.
祂可以行神蹟奇事.
祂的恩典夠你和我用.
活在此時此刻.
經歷神的供應.
第三樣.
多思想神的恩賜.
傳道者說.
人能夠享受生命.
是神的恩賜.
這是基督徒的體驗.
神的恩賜提醒我們.
要多思想神的恩典.
而不是單單集中在.
神給予生命的缺陷不足.
或限制.
當你在痛苦.
有掙扎的時候.
有眼淚和懷疑.
迷惑的時候.
請你轉一轉.
你要仰望神.
抬起頭來.

$^{521}$思想神的恩賜.
正如使徒雅各吩咐.
你在百般試煉當中.
不要只想到所受的苦.
而是要看到苦難帶給你好處.
甚至以受苦為大喜樂.
雅各書1章2節這樣說.
所以你在痛苦裡面.
思想神的恩賜.
就好像人用手.
抹去遮住眼睛的淚水一樣.
自然看到前面.
柳暗花明柔一村.
這是很重要的東西.
因為人在軟弱的時候.
往往看到不好的東西.
看到負面不幸的東西.
你在痛苦裡面.
只看到消極的東西.
你看不到背後.
這位神.
看不到神.
繼續有恩典.
你聽過台灣有兩個.
很著名很出名的.
基督徒女作家.
杏林子和三毛.
他們兩個都是好朋友.
他們兩個都去世了.
杏林子12歲患了.
類風濕骨關節炎.
全身90\%的關節都損壞.
腳不能走.
肩不能動.
手活動有限制.
前面病癱.
幾十年才去世.
進出醫院無數之多.
記都記不到.
杏林子分享.

$^{561}$三毛有一次探她的經過.
三毛看到杏林子病得這麼難過.
她熱心過度.
杏林子我為你祈禱.
於是祈禱.
求神早些接去杏林子的生命.
讓她好好幸福.
嚇到杏林子馬上阻止她.
她說天堂雖然好.
但別人欠我的債尚未償還之前.
我不甘心走.
我欠人的情未完.
我也不能夠安心走.
我就是這麼死心眼的人.
我向神祈禱.
希望我有90歲.
不過神打不打折扣是神的事.
這樣計算.
我還有50年的生命可以做很多事.
等到那一天.
我可以在結算單上蓋印.
銀貨兩行.
即銀貨交收.
我就含笑而去.
三無所見的只是杏林子的痛苦.
眼睛被淚水所遮蓋.
看不到她在杏林子身上的恩賜.
而杏林子是抹走眼淚.
看到神的恩典.
讓自己繼續做未來要做的事.
而且覺得還有很多事等著我去做.
杏林子在她的病痛裡面.
落在其中.
不願離世.
反而三無.
你知道她後期會怎樣死.
會自殺死.
或者因為三無只是看到生命裡面有些缺陷.
有限制.
還有很多悲哀.

$^{601}$消極的事.
忘記了神的恩惠.
還有慈愛的福樂.
有些人說她好像兩個人都困在人生的監獄裡面.
一個只是看到窗戶的鐵枝.
即是巴.
另一個看到窗外的星星.
Star.
杏林子.
你和我在苦難裡面.
你有沒有超越巴這個鐵枝.
而看到外面的星星.
Star.
超越痛苦.
而看到神的恩典.
這個可以令你從憂愁變為喜樂.
第三樣時間關係我不說了.
但是大標題一定要你知道.
麻煩你弄出來.
一起讀一讀.
要以靠康復我們的父神.
裡面是說到.
神是一位怎樣的神.
祂是一個全知.
祂是永恆.
亦都完全無所增添無所減少.
亦都配不你和我敬畏的神.
而且後面那兩句.
現今的事將來的事等等.
其實張天神是色在.
金在.
永在.
將來都是掌管一切的神.
我想到有一個機構.
你以前聽過內地會.
現在改名叫OMF.
海外基督事團.
是嗎.
他們內地會以前有兩句金句.
兩句金句就是耶和華已立.

$^{641}$和已變已者.
耶和華已立是什麼意思呢.
就是主耶和華上帝必恭應與備.
已變已者是什麼意思呢.
就是到如今神都是幫助我們.
當我們敬畏神.
以靠神.
覺得這位神是恭應我們.
直到如今都幫助我們的主.
我們的心會安穩起來.
我們低頭有個祈禱.
結束.
求主你自己話語.
堅立幫助我們.
能夠參透人生.
願你顧念.
帶領.
提醒.
祈禱.
共主耶穌基督明教.
阿們.
\newpage



\section{創世記 1:26-27}
\label{sec:1bmqQECjSXg}
\textbf{與人同在的神 \#2 - 創世記 1\_26-27 - 2019/10/13 (陳崇基牧師)}
\newline
\newline
連結: \href{https://youtube.com/watch?v=1bmqQECjSXg}{\texttt{https://youtube.com/watch?v=1bmqQECjSXg}} ~~~~ 語音日期: 2019-10-17
\newline
\newline
\hyperref[sec:LfJIKNZbdWE]{\small{< < < PREV SERMON < < <}}
~
\hyperref[sec:index_chronic]{\small{[返順時目]}}
~
\hyperref[sec:index_scriptual]{\small{[返順卷目]}}
~
\hyperref[sec:0WFluTmE3pc]{\small{> > > NEXT SERMON > > >}}
\newline
\newline
創世記 1:26-27
\newline
\begin{longtable}{cl}
\hline
\hline
章節 & 經文 (和合本修訂版)\\
\hline
1:26 & \begin{tabularx}{0.7\textwidth}{X} 神說：「我們要照著我們的形像，按著我們的樣式造人，使他們管理海裡的魚、天空的鳥、地上的牲畜和全地，以及地上爬的一切爬行動物。」 \end{tabularx} \\ \\ \relax
1:27 & \begin{tabularx}{0.7\textwidth}{X} 神就照著他的形像創造人，照著神的形像創造他們；他創造了他們，有男有女。 \end{tabularx} \\ \\
[1ex]
\hline
\hline
\end{longtable}
$^{1}$今天我們繼續與人同在的神的系列.
今次是第二講.
我們說過在網上可以問一些問題.
我知道上次我說的時候那個QR Code有些問題.
今次這個QR Code應該沒有問題.
如果你掃描了的話.
你現在可以掃描了.
掃描了之後你可以submit問題.
然後我們最後的時候會做一些解答.
希望不會太長期.
最後大概有五至十分鐘可以解答中間的問題.
我已經收到有些不同的問題.
最後的時候會解答.
今次第二講我們說過在聖經裡面.
我們由第一章開始看一些重要的經節的時候.
我們知道聖經最主要要做什麼呢.
就是給我們一個新的世界觀.
什麼叫世界觀.
什麼叫world view.
就是一套我們的信念.
以至我們能夠支撐著我們的生活.
這些問題很多時候我們平時不會刻意去問.
但我們需要去問.
上一次我引用過一本書叫Truth is stranger than it used to be.
1990年代末的時候寫的一本書.
兩位基督徒哲學家說.
聖經要去解答什麼問題呢.
就是我們到底是誰.
世界有什麼問題.
這個問題有什麼最終極的答案.
而且現在我們應該怎樣去生活下去.
等等這些問題很多時候我們潛意識裡面已經有答案.
但聖經卻要告訴我們一個刻意的和一個更加明顯的.
告訴我們在聖經裡面我們需要一個怎樣的生活觀念.
一個人生的觀念.
到底我們自己是誰這個世界是怎樣的和神有什麼關係.
這些所謂世界觀的問題.
世界觀的概念.
其實不是一些很陌生的概念.
你不要覺得這些哲學的東西很抽象.

$^{41}$其實我們每個人都會有這樣的預設.
在我預備的時候.
我想起有一本書在90年代的時候我看過.
有一位科學家叫Carl Sager.
他寫了一本書叫做Pale Blue Dark.
話說美國在1977年.
太空總署射了一件東西上去.
外太空.
那件東西叫做Voyage 1.
這件Voyage 1在太空裡面射出去之後.
就一路離開地球一路進發.
現在還在進發中.
這件東西已經服務了我們42年.
它一路向外太空一路走.
現在的速度還是每小時4萬個mile per hour.
一個這樣的速度.
它繼續在外太空去搜集資料.
然後就變回來地球.
這個Voyage 1離開了太陽系.
再過了另外兩個星體.
最後一次能夠看到地球的時候.
太空總署就叫他拍最後一張照片.
就是拍地球最後一張照片.
這張照片是在1990年2月14日情人節那天拍的.
他拍下來的時候.
就發現我們這個地球.
原來只是在太陽系裡面的其中一點.
剩下一點.
作者就說.
就是在這麼小的一點裡面.
我們人類歷史過去現在將來.
所有人就是在這個點裡面生活.
這是一個世界觀.
我們在一個這麼浩瀚的宇宙裡面.
我們是多麼孤獨.
你知不知道為什麼你在這裡.
為什麼我存在.
這個世界觀就是告訴你.
這本書或這觀念告訴你.
總之你就是在這裡.

$^{81}$突然間你就在這裡.
所以以前小時候媽媽罵我.
說一句話是很對的.
她說石頭是爆出來的.
或者生一塊叉燒就生你.
所以你去餐廳買一塊叉燒.
可能價值和我的心是差不多的.
其實意思就是石頭爆出來的意思.
為什麼你會在這裡.
我是誰.
我去哪裡.
我為什麼在這裡.
如果你沒有創造沒有一個設計的話.
這個浩瀚的宇宙只不過是一個意外的出現.
換句話說我們人的存在.
我們今天活生生在這裡.
是沒有意思的.
這是一個世界觀.
因為正在回答剛才我所列出的幾個問題.
我是誰.
這個世界是怎樣的.
我去哪裡.
等等這些問題.
這個系列我們就是去看一些經文.
告訴我們聖經.
如何啟示出一個新的觀念.
一個新的世界觀.
到底你和我存在有什麼意思.
我們說過聖經是在迦南地或以色列這個地方寫的.
在古代世界這個地方.
所謂迦納這個地方.
或者現在的以色列巴勒斯坦地.
就是古代世界裡面.
古代人類文明的世界裡面.
那時候還沒有美洲.
還沒有發現美洲.
這個是當時世界的中心.
因為它是聯絡非洲亞洲和小馬紹爾亞洲的一個地橋.
在這個連接裡面就是人來人往的地方.
所以在聖經看的時候.

$^{121}$你需要明白聖經是一個生活文件.
或者是一個活著的文件.
因為這個文件就是和世界各地不同的觀念或思潮去做辯論.
或者去做一些對話.
這些思潮如何講人類.
所謂我們為何會在這裡.
聖經就去做一個回覆.
或者做一個對話.
說他們對或者很多時候說他們這樣看是不足夠是錯的.
所以你看聖經一定要知道是一個生活文件.
所以為什麼很多時候你要看背景.
背後的人到底在說什麼故事.
而為什麼聖經要講出這個故事.
去做對話或者做一個對話.
今天我想和你說一個聖經背後的故事.
背後或者一些其他帝國.
他們很流行的講一些故事是講人類為什麼會在這裡.
這個故事我就用七張的slide去表明出來.
這個故事叫做Elumar Elish.
這個故事是一些人發掘在伊拉克古文明裡面找到的一些泥板.
這些泥板講出一些古代的故事.
或者當時耳熟能詳的家傳戶曉的一些世界的故事.
這個故事就是一些創造的故事.
為什麼要看這個故事呢?.
因為我們在看創世記第頭兩章.
我們在看頭兩章的時候.
我們知道是一個創世的story.
而這些帝國四周圍的帝國特別是巴比倫的帝國.
他們怎樣去講這些故事呢?.
他們都有創世故事.
但他們創世故事的版本或者他們故事的內容到底是講什麼.
這個Elumar Elish就是其中一個很出名的故事.
當時在街上行走的人每個人都會聽到這個故事.
或者每個人都會明白這個故事所講的意思.
請你不要介意我們用一些很低成本的方法.
一些不是動畫來表達這個故事.
因為我們知道如果你由頭到尾你會悶死.
所以我用七張的圖畫來告訴你這個故事講什麼.
這個故事開頭發生什麼事呢?.
這個故事就是說.

$^{161}$泰初在天.
他說有兩個神叫做Absu和Tiamat.
是男和女.
他們是中神的父母.
所以他們因為是父母.
他們就生了很多代的小孩.
這些中神就出現了.
這是巴比倫多神論裡面有的一個開始.
這兩個父母生了很多小孩的時候.
生小孩好不好?.
生小孩好啊.
但問題就是這樣.
這個故事裡面說這些小孩生得多的時候.
就令到你睡不著.
因為他們很吵.
確實他真的這麼說.
中神的兒女發出太多的噪音.
令到那個父母Absu不能夠睡覺.
擾人清夢.
所以Absu生氣了.
就決定將這些眾子女全部殺掉.
我希望不會發生在你的家庭裡.
但故事就是這樣.
因為他失眠所以要殺死他們.
但誰知道其中一個的兒子.
第四代的兒子叫做Ian.
他就發現了這個詭計.
他就先將他的爺爺殺掉.
所以他就用咒語殺死了Absu.
然後他就做了王.
然後他就生了另外一個兒子.
叫做馬獨克.
叫做Mardu.
這個Mardu就是巴比倫的京城的大帝.
這個是最大的神.
就叫做Mardu.
你知道一開始Absu和Tiamat是夫妻.
Absu死了之後老婆怎麼樣.
有人殺了你老公你會怎樣.
他就要報仇.

$^{201}$所以Tiamat就興起他的叛軍.
然後他就決心來復仇.
這些眾神靈聽到之後就非常害怕.
然後不知道怎麼辦.
馬獨克就出來.
Mardu這個巴比倫的神.
他就出來說他願意和他決一死戰.
但他說勝利之後他就要所有的眾神來拜他.
以致他能夠成為至高者.
所以Mardu就答應了幫大家出戰.
大家就說好Mardu你去出戰.
於是他就用了很多的武器.
用了很多怪獸的武器.
文件裡面都是這樣說.
大家都出了很多不同的武器.
然後就將他們的叛軍Tiamat殺了.
然後他就創造天地.
他就將Tiamat本身是什麼神.
是水神.
將水分成兩半.
一個放在天上一個放在下面.
然後他就建立萬物.
最重要當然就是建立巴比倫.
和巴比倫的神廟.
這個就是世界的重心.
不要以為完結了.
還有的.
創造萬有的時候還有一件事還沒創造.
故事發生是什麼呢.
就是說那些眾神因為他們公司工作很忙.
他們要服侍那些大神.
他們經常要供奉他們.
服侍他們.
大神會肚餓的.
所以經常要做東西給他們吃.
他們覺得很辛苦.
所以他們就在打仗.
好像我們CPS.
我們的老師在打仗.
說威脅要打仗.

$^{241}$結果最大的神.
現在這個就叫做Marduk.
Marduk就說不如這樣吧.
我們做一班人出來.
以致你們不用做了.
是不是.
我們做一班TA出來.
那你們那些老師就比較好做了.
結果他就做了人類出來.
這些人類就成為了奴隸.
怎樣做一個人出來呢.
就將叛軍裡面其中一個神.
叫做Kingu.
就將他的血製造了人出來.
所以人類永遠是下面做奴隸的.
所以整個故事就是這樣發展.
清楚嗎.
好看嗎.
到底這些故事應該怎樣去解釋.
有些人會說是神話.
等於呂渥保天.
等等這些古代的神話.
你要知道這些所謂的myth.
其實是political story.
意思就是說.
誰會作一個這樣的故事出來.
去告訴社會的人聽.
百姓聽.
你們生存的意義是這樣的.
因為是一個政治宣傳的故事.
這個是巴比倫帝國有的世界觀.
換句話說.
那個所謂的Martin Luther.
巴比倫最大的神.
其實是代表巴比倫帝王.
皇帝.
因為當時是獨裁的社會.
上面那個就是最大的.
他就是代表巴比倫神.
坐在皇位上面.

$^{281}$而他皇室裡面的成員.
所謂的白宮裡面的內閣.
那些貴族.
就是代表那些比較小的神.
最下面最下面的.
就是那些草民和百姓.
那故事說出什麼世界觀呢.
就是.
暴力是這個世界固有的秩序.
因為你留意那個故事.
暴力是自有永有的.
永遠都是打來打去.
永遠都是互相廝殺.
誰殺得最厲害那個.
那個拳頭最大的就坐在上面.
你知道Martin Luther起來的時候.
本身已經有暴力.
老爸殺兒子.
兒子又殺老爸.
所以大家互相為仇.
政權的暴力.
是世界秩序的固有.
實際上這個故事.
Eluma Elish每年都起碼一次.
要在巴比倫的神廟裡面讀一次.
等於他們農曆新年的時候.
他們就要講出這樣的神話.
就來鞏固誰的權力呢.
就是Martin.
或者是代表神靈的帝王.
所以就是坐在上面.
來管治所有的人.
如果你想做叛軍的話.
故事就告訴你.
你就一定會被殲滅.
因為他不會手軟.
這個是一個世界觀.
等於.
可能你的背景是這樣.
或者你曾經經歷過這樣的教導.

$^{321}$等於有人說這個世界是唯物的.
只有物質.
由始至終.
人類歷史社會都是階級鬥爭.
等於暴力.
是互相鬥爭的成果.
鬥爭會繼續開始.
繼續延續下去.
這是一個世界觀.
或者等於我剛才所說的Pale Blue Dot.
這個世界是沒有意義的.
意義是怎樣來的呢.
是你自己創造出來的.
往往很多時候就是用武力去創造.
這個故事是一個政治故事.
也告訴那些草民百姓.
你們本身就是要在下面當奴隸.
而且你本身的人性.
其實是邪惡的人性.
為什麼呢?因為你怎樣做出來的.
你是叛軍的神靈的血所做出來的.
所以人是沒有尊嚴的.
沒有固有的尊嚴.
沒有固有的美好的人性.
這是一個世界觀.
所以我這樣說的時候.
請你不要以為.
這只是很老遠的事.
或者跟我沒有什麼關係的事.
所有人類的社會.
都自然有它自己的世界觀.
無論你叫它階級鬥爭也好.
叫它完全沒有意義的唯物的.
這是世界觀.
或者巴比倫神話的世界觀.
創世記一開始.
就給我們一個截然不同的世界觀.
這個世界觀說什麼呢?.
你要留意創世記.
不是一個政治故事.

$^{361}$因為一開始就說普世.
沒有說以色列.
大部分的國家主義故事.
全部都是國家主義的故事.
我們的民族是最大的.
我們的民族是最好的.
巴比倫民族是最好的.
其他反叛我們的.
那些就是非巴比倫的人.
那些是草民.
那些是要打壓的人.
或者是我們需要殲滅的人.
創世記很奇怪.
是以色列人本土出來的故事.
但頭十一章全部都是普世的.
你問阿當夏娃她是什麼國籍呢?.
她沒有國籍.
什麼國籍?.
她就是代表所有人類.
阿當夏娃沒有護照.
到以色列什麼時候才有呢?.
大概如果你說一個國家的成立.
要到出埃及記中間.
他們出埃及取得上帝的律法.
才成立成為一個希伯來民族.
所以創世記一開始的時候.
是一個普世的觀點.
說到神怎樣看待整個人類.
阿當夏娃是代表人類.
並且世界是美好的.
上個禮拜第一次我們說的時候.
我們說.
事就這人成了.
出現過多少次?.
七次.
上帝看他所做的為好的.
good.
甚至最後是very good.
都是七次.
只不過因為人叛逆上帝.

$^{401}$他想遠離上帝.
罪惡入了世界.
然後才會有神人之間的關係破裂.
男與女之間人間的權力破裂.
然後到第四章.
有一個故事叫做該隱與阿爸.
說到兄弟之間的爭鬥.
第四章尾.
說到所謂institutional violence.
該隱之後的一代叫做拉民.
他說他要報仇七次.
這是一個完全是人類的罪惡.
進入了人間之後.
破壞所有的關係.
或者將暴力的罪行.
越來越升.
越來越增加.
創世的第一章說到什麼呢?.
他說我們是誰呢?.
第一章第26,27節.
如果你有聖經打開的話.
我們看看這段聖經.
他說神說我們要照著我們的形象.
按著我們的樣式做人類.
使他們管理海裡面的魚.
空中的飛鳥.
地上的牲畜和泉地.
並地上所爬的一切昆蟲.
神就照著自己形象做人.
乃是照著他的形象做男做女.
這個中文是我自己的反應.
你見到聖經裡面說是做人的.
我回去做人類.
回去做人類.
這裡有幾樣東西需要去想一想.
或者我想告訴你這裡想說什麼.
神說我們要照著我們的形象做.
這個是眾數.
這個就很奇怪了.
為什麼神是眾數.

$^{441}$因為神是獨一的神.
這個可能會是這樣.
以前的皇帝.
古代世界的君王.
他是有朝廷的.
所以君王很多時候說話.
他都是會用自我的眾數說話.
等於今天我們聽美國總統說話.
他會說We.
我們.
我們要撤兵.
我們要立法案等等.
我們要制裁那個國家或者那個國家.
他是說我們.
所以帝王說話.
或者在上掌權的說話.
很多時候就是用我們.
眾數.
是poor.
或者一個royal poor.
所以說我們.
照著我們的形象去做人類.
人類是我的反應.
我覺得NLT的反應是不錯.
Human beings.
因為那個原文是什麼呢.
那個原文就是Hada.
我們叫阿當的意思.
第二張中文翻譯就是那人.
什麼叫Hada.
Hada就是humanity.
人類的意思.
不是一定有的意思.
而是可能是一群人的意思.
或者應該這樣說.
我們想阿當的時候.
我們不要只想一個人.
而是想一個代表全人類的人.
所以說出阿當.
或者夏娃的故事的時候.

$^{481}$聖經是在說每個人的故事.
我們以往看很多兒童的.
聖經的故事書.
由開始神創造的時候.
六日創造.
然後那個畫面圖畫一定是說.
在花園森林那裡.
有一男一女沒有穿衣服.
一定是這樣.
只有一男一女.
但第一章26-27講什麼呢.
那是人類.
Hada.
不是做一個人.
也不是做一男一女.
而是做整個人類.
而整個人類的代表.
或者化身.
就是阿當自己.
如果你這樣看的時候.
你解決了很多問題.
第四章那個人殺了阿伯之後.
他有妻子.
妻子是誰.
以往解釋可能就是.
他和他妹妹結婚.
兄妹結婚.
我們就有很多困難.
或者他殺了阿伯之後.
那個人說了一句話.
他很怕別人來報仇.
誰會報仇.
父母當然很生氣.
但他離開了他們.
他到底怕誰報仇.
所以如果你這樣看的時候.
阿當是代表全人類.
正如上帝第一章所創造.
飛鳥.
地上的活物.

$^{521}$海裡面的魚等等.
都是一群做出來的.
但你不要忘記.
阿當其實就是每個人.
所有人的代表.
阿當做的事.
就是每個人的代表.
等於我今天跟你講一個故事.
很久很久之前.
有一個人.
他整天吐痰.
很不衛生.
也很沒禮貌.
那個人叫什麼名字.
他姓鍾.
叫做鍾國仁.
你就知道我說的不是某某鍾國仁.
而是所有的鍾國仁.
你明白我的意思嗎.
所以創世記第一二三章說的是.
整個人類到底發生什麼事.
上帝創造人類.
讓他能夠管理萬物.
是一個很榮耀的責任.
能夠跟上帝一起去管治.
上帝美好的創造.
並且人是按照神的形象所做.
你有沒有想過.
你是在反映上帝.
創造主有什麼性情.
有什麼本性.
在你的身上.
或多或少.
會有被反映出來.
知情義.
我們人類有的本性.
也是上帝有的本性.
並且最重要.
這個image of God.
或者神的形象.

$^{561}$是一個皇室的用語.
代表上帝在地上擔任作管家.
作經理manager.
來管治這個世界.
並且能夠跟上帝建立關係.
所以有一個猶太的聖經家.
他說得很對.
他說上帝創造人類的時候.
給人類最大最大的禮物是什麼.
是自由.
freedom.
you use it or lose it.
你用上帝給你的自由.
去選擇去敬畏去事奉他.
自由.
人權.
就是上帝給我們最大的禮物.
所以為什麼呢.
由這一章我們會找到很多.
在先知書以後舊約的經文裡面.
講到一些人權或者社會公義問題.
這個就是一個基石.
因為每一個人都是上帝所創造.
也是帶著上帝的形象.
所以以後講所有justice的issues.
全部都是按照這裡開始.
你想想如果不是上帝.
他自己的創造.
美好的創造.
賦予人有智慧.
有人權.
有自由.
知情義等等.
這些美好的本性.
跟上帝能夠建立關係的話.
如果不是由創造開始的話.
人存在幾乎是沒有意義的.
因為你沒有一個外在的.
一個客觀的進層.
去告訴你什麼叫做意義.

$^{601}$今時今日.
這段聖經就經常被人用.
並且是應該我們思考怎麼去用.
在上個星期五.
我們美國南方浸信會的總會.
我們有一個叫做ethics and religious liberty commission.
一個這樣的委員會.
專門去關注一些倫理.
公義或者宗教自由的問題.
這位負責人叫做Russell Moore.
Russell Moore經常為一些公義的issues來說話.
我很佩服他.
他很願意說話.
很敢說話.
上個星期五他發表了什麼講話呢.
他就說到這個世界裡面.
有很多人去漠視人權.
有很多的政權也是漠視人權.
有很多的生意.
有些商業的組織.
也只是顧著賺錢.
而漠視人權的.
包括NBA.
大家知道發生什麼事了.
他寫了一封信.
這封信是給美國國會的.
他說關注在新疆人權的問題.
維吾爾族人權的問題.
也關注美國在北敘利亞撤兵.
那些所謂庫爾德族的人.
本身是美國的盟友.
幫忙打ISIS的.
因為撤兵的緣故.
他們現在面臨一個很大的人類的危機.
就會被土耳其軍可能會滅族.
這是值得關注的人權的問題.
他最後說了一句話.
他說現在這個世界.
看人越來越低賤.
越來越賤價.

$^{641}$到底人有什麼價值.
有什麼自由.
有什麼權利.
少數的民族.
各地不同的地方.
基督徒是受打壓的.
這就是image of God給我們的基礎.
我們能夠去說這些話.
我想問你們有沒有關注這些.
如果上帝的世界觀告訴我們.
人類是神的形象.
我們就應該去說這些話.
我們就應該去關注這些事情.
上帝創造人.
就是要人在大地上.
為他一起管理這個世界.
並且和上帝建立一個關係.
幾百年前有一個信經.
什麼叫信經呢.
就是教會總結出來的一些神學信仰.
那個信經叫做.
Westminster Catechism.
有一個所謂小冊子.
叫做Solid Set Catechism.
叫做小冊子.
就是小要你問答.
那個Catechism是做什麼呢.
就是給剛剛信主的人.
他們有一個很好的真理的基礎.
用一個方法就是問與答的方法.
一問一答.
第一條問題是問什麼呢.
第一條問題是.
人類主要生存的目的是什麼.
What is the main goal of humanity.
人主要的目的答案是.
就是要永要神.
永遠以神為樂.
Glorify God.
and enjoy Him forever.

$^{681}$我想你留意一件事.
什麼叫做enjoy Him forever.
你有沒有想過你和上帝之間的關係.
是可以enjoy他的痛楚.
當你翻開聖經.
當你敬拜.
當你禱告的時候.
和上帝傾訴的時候.
就是以神為樂.
你不知道有沒有想過.
我們被造的目的.
有一個這麼永要的目的.
一個這麼有尊貴的目的.
這個就是我們有的世界.
求主幫助我們和上帝之間建立一個更緊密的關係.
並且在地上.
我們按照他的形象.
按照他的真理.
去管理.
去發聲.
去行事為人.
按照上帝的聖情.
上帝的本性.
我們同心協力.
多謝收睇城匯網台.
記住like comment或者訂閱subscribe我地呀.
\newpage



\section{創世記}
\label{sec:0WFluTmE3pc}
\textbf{與人同在的神 \#3 - 創世記 3 - 2019/10/20 (陳崇基牧師)}
\newline
\newline
連結: \href{https://youtube.com/watch?v=0WFluTmE3pc}{\texttt{https://youtube.com/watch?v=0WFluTmE3pc}} ~~~~ 語音日期: 2019-10-23
\newline
\newline
\hyperref[sec:1bmqQECjSXg]{\small{< < < PREV SERMON < < <}}
~
\hyperref[sec:index_chronic]{\small{[返順時目]}}
~
\hyperref[sec:index_scriptual]{\small{[返順卷目]}}
~
\hyperref[sec:Gl7QFa_KpBw]{\small{> > > NEXT SERMON > > >}}
\newline
\newline
$^{1}$我們看過《創世記》第一.
第二章左右.
我們講過聖經啟示的這位上帝.
是一位創造萬有的神.
這是我們第一集講到的.
第二次我們講到.
就是聖經這位上帝.
是去創造人類.
按照他的形象.
有他的尊貴.
有他的榮耀去創造人類出來.
這是第二次講的.
這次我們第三集就講.
創世記第三章.
讓我一口氣讀.
創世記第三章.
(2)神所造的,惟有蛇比填冶一切的活物更狡猾..
蛇對女人說:.
「神起事親說:.
不許你們吃圓中所有樹上的果子嗎?」.
女人對蛇說:.
「圓中樹上的果子我們可以吃,.
惟有圓中那火樹上的果子.」.
神曾說:.
「你們不可吃也不可摸,免得你們死.」.
蛇對女人說:.
「你們不一定死,.
因為神知道你們吃的日子眼睛就明亮了,.
你們便如神能知道時兒.」.
於是女人見那火樹上的果子.
好作食物也悅人眼目,.
且是可喜愛的,.
能使人有智慧就摘下果子來吃了,.
又給他丈夫,他丈夫也吃了,.
他們二人的眼睛就明亮了,.
才知道自己是赤身露體,.
便拿無花果樹的葉子為自己編作裙子..
天起了涼風,.
耶和華在神的園中行走,.
那人和他妻子聽見神的聲音,.

$^{41}$就藏在園裡的樹木中,.
躲避耶和華神的面,.
耶和華神呼喚那人,.
對他說:.
「你在哪裡?」.
他說:.
「我在園中聽見你的聲音,.
我就害怕,.
因為我赤身露體,.
我便藏了.」.
耶和華說:.
「誰告訴你赤身露體呢?.
莫非你吃了我所吩咐你.
不可吃那樹上的果子嗎?」.
那人說:.
「你所賜給我,.
與我同住的女人,.
她把那樹上的果子給我,.
我就吃了.」.
耶和華神對女人說:.
「你做了什麼事?」.
女人說:.
「那蛇引誘我,.
我就吃了.」.
耶和華神對蛇說:.
「你既作了這事,.
就必受咒詛,.
比一切的生畜野獸更甚,.
你必用兔子行走,.
終生吃兔,.
我又要叫你和女人彼此為仇,.
你的後裔和女人的後裔.
也彼此為仇,.
女人的後裔要傷你的頭,.
你要傷她的腳筋.」.
又對女人說:.
「我必多多加增你懷孕的苦楚,.
你生產兒女必多受苦楚,.
你必戀慕你的丈夫,.
你的丈夫必管轄你.」.

$^{81}$又對亞當說:.
「你既聽從妻子的話吃了我所吩咐,.
你不可吃那樹上的果子,.
地必因你的緣故受咒詛,.
你必終生勞苦才能從地裡得吃的,.
地必給你長出荊棘和窒泥來,.
你也要吃田間的菜蔬,.
你必汗流滿面才得糊口,.
直到你歸了土,.
因為你是從土而出,.
你本是塵土,.
仍要歸於塵土.」.
亞當給他妻子起名叫夏娃,.
因為她是眾生之母,.
又說神為亞當和他妻子用被子做衣服,.
給她們穿,.
又說神說:.
「哪人已經與我們相似,.
能知道善惡?.
現在恐怕她伸手遊摘生命樹上的果子,.
就永遠活了.」.
又說神便打發她出伊甸園去,.
耕種她所自出之土,.
是把她趕出去了,.
又在伊甸園的東面,.
安設基路柏和四面轉動發火焰的劍,.
要把守生命樹的道路..
這個故事絕對不是陌生,.
大家可能聽過,.
如何去看這段故事呢?.
或者我們在上次的兩集,.
我們說過在《創世記》一二三章,.
到底是一種怎樣的題材,.
到底它想溝通甚麼,.
告訴我們甚麼,.
很肯定不是想說科學的東西,.
很肯定亦不是告訴我們.
一個人類的百科傳書,.
或者人類早期歷史的百科傳書,.
聖經最主要想溝通的是一個世界觀,.

$^{121}$告訴我們如何了解到底我們是誰,.
現在我們是怎樣,.
所以用一些故事的形式去表達出來..
如何解讀聖經呢?.
我想先做一個例子,.
大家有聽過一個故事叫浪子的比喻,.
有吧?.
耶穌講的這個故事的比喻,.
這個比喻如果沒有聽過很簡單,.
就是有一個好的爸爸,.
他有兩個兒子,.
一個大一個小,.
小兒子很壞的,.
他爸爸未死之前就說要跟他分家產,.
他有一天來到他爸爸面前,.
他說不如我們先分家產,.
其實是咒爸爸死,.
當爸爸是死的,.
然後他拿了那些家產,.
就出去外面花天酒地,.
花了那些錢之後,.
他慘到一個地步,.
他就去放豬,.
直到有一天他突然醒過來,.
醒過來他說,.
爸爸一向對我很好,.
提供了我所有的東西,.
為什麼我要在這裡捱窮呢?.
要在這裡放豬呢?.
不如我認錯,.
我回到家裡,.
寧願在富家裡面做一個工人好了,.
於是他就起來,.
然後回去途中遠遠望見,.
在門口,.
他爸爸居然在等他,.
他就跑出來,.
老父就跑出來,.
然後去迎接他,擁抱他,.
去接納他,.

$^{161}$講到這裡的故事一半是這樣,.
你可以問,.
浪子的故事到底有什麼意思,.
你可以有字面的意思,.
字面的意思就是說,.
有一個父親和兒子之間的關係,.
但很肯定不只是表面的意思,.
是一個象徵的意思,.
所以你不要問,.
到底是不是當時巴勒斯坦地區,.
有一些家暴的問題,.
家庭的問題,.
兒子是不是需要看心理醫生等等,.
這不是重點,.
重點不是表面的意思,.
重點是象徵的意思,.
是那個象徵的符號裡面所代表的是什麼,.
所以那個家父就是代表上帝,.
神,.
怎樣看待我們,.
那個小孩花天酒地,.
然後悔改回來父的面前,.
那個就是罪人,.
小孩就是代表罪人,.
這個才是重點,.
那個故事的威力,.
不在於歷史裡面,.
到底細節發生了什麼事,.
而最重要是他所象徵的是什麼,.
一個叛逆的兒子回到家裡,.
就代表我們一群叛逆上帝的人,.
我們回到家裡見天父,.
我們罪人悔改,.
能夠得及,.
創世記第三章其實都是一樣,.
你看一個故事的時候,.
你不需要想一下到底當時發生了什麼事,.
如果你Google搜尋一下,.
搜尋一下譬如阿當夏娃如果是華人的話,.
這個故事會發生什麼事,.

$^{201}$所以有些人畫一些卡通,.
他說如果阿當夏娃本身是華人的話,.
早就不用談了,.
就當了蛇羹,.
這是一個服務的意義,.
但到底真正的意義是什麼呢?.
到底那條蛇是不是真的會說話,.
是不是真的沒有穿衣服,.
然後在樹的下面,.
那個夏娃,那個女人,.
跟那條蛇在聊天,.
是不是那條蛇本身是有腳的,.
做錯事之後,.
引誘人犯罪,.
離棄神之後,.
她就沒有了腳呢?.
That's not the point,.
那些只是服務的意義,.
所以我們要留意到底進入這個故事裡面,.
有些什麼symbol,.
這些symbol背後要表達的,.
代表的是什麼?.
剛才說了蛇,.
蛇是最狡猾的,.
這個我想整本聖經,.
當時第一代讀者未必會明白,.
他可能會覺得,.
第一代讀者可能會覺得,.
這是一個邪惡的標記,.
因為它是引誘人類離開上帝,.
翻轉了上帝所創立的世界,.
可能會是這樣,.
但是我們整本聖經知道,.
耶穌基督都來了,.
有更完備的啟示的時候,.
我們就知道這個是代表撒旦,.
聖經裡面有很多不同形容撒旦,.
魔鬼的一些用語,.
他說他是光明的天使,.
光明天使就是說,.

$^{241}$你看出去不是,.
你看下去不是很恐怖,.
但是看下去反而是light,.
是光明的,.
不過他所說的話是半真半假,.
去引誘人的,.
我們知道啟示錄裡面告訴我們,.
這條蛇是代表撒旦,.
也是反叛上帝的那幫人,.
那幫墮落的天使的首領,.
在這一說,.
他只是說到那條蛇說了一些話,.
其實他不是說了話,.
其實只是問問題,.
那個神啟示真說,.
神是不是真的有說過這些話,.
你可以看到那個試探的mechanism,.
很多時候不是叫你做壞事,.
而是去提出一點點的懷疑,.
這位神,.
原來是一個這麼過份嚴厲的主人,.
給你一點點的樂趣都不行,.
是不是真的會這樣,.
一點點享受都不給,.
上帝是很mean,.
園裡面樹上的果子是不是都不行,.
是不是真的不給你吃掉所有樹的果子,.
上帝是不是真的這麼mean,.
然後女人就回答他,.
你知道聖經裡面由第一章第一字到現在,.
只是三章的聖經,.
第一個要否定的對上帝啟示的教導,.
否定的教義,.
就是審判的教義,.
Doctrine of Judgment,.
第一個要否定的,.
就是否定上帝會問責,.
否定我們做了事情之後是有consequence,.
留意一下社所說的話,.
或者對答問的問題,.

$^{281}$是部份是真的,.
但全部都是錯誤的,.
partially true,.
totally wrong,.
部份是真確的,.
但完全是錯誤的,.
完全是謊話,.
是社所問的問題,.
說到果子,.
分辨善惡的果子,.
果子絕對不是蘋果,.
拉丁文後期的時候,.
翻譯拉丁文的語言,.
蘋果就是Malu,.
邪惡就是Malus,.
所以人們就有一個傳奇,.
說禁果就是蘋果,.
所以很明顯蘋果,.
不是蘋果,.
果子背後所代表的才是最真實的,.
所以不是摘一個果,.
不是去超級市場買蘋果橙這麼簡單,.
不是說阿東阿蛙真的吃一個蘋果,.
這麼表面,.
而是背後代表著一些東西,.
分辨善惡,.
在聖經舊約裡面提到分辨善惡的時候,.
決定是非的時候,.
很多時候是君王才會做,.
或者是上帝才會做,.
上帝才有絕對的主權做一個分辨,.
決定善與惡,.
決定好與壞,.
good and evil,.
這個能力和角色,.
你記得第一章我們看過,.
上帝創造萬有的時候,.
上帝是看光為好,.
他看到這件事為good,.
然後他就說這件事是好的,.

$^{321}$這個就是上帝自己有的角色,.
所以摘那個果,.
拿那個分辨善惡樹上的果子,.
意思就是說我要做上帝,.
那個創造我不是我的上帝,.
所以不是說吃一個蘋果這麼簡單,.
所以有個神學家說,.
是the guardian of God,.
本身他是上帝,.
我們將他de-guard,.
他不再是上帝,.
所以決定善惡的是誰呢?.
就是我,.
我自己來決定善惡,.
自己來決定好與壞,.
這個就是吃分辨善惡樹果子的本質,.
本身創造我們的是上帝,.
但他不是我的上帝,.
我自己才是上帝,.
不是他所決定的是與非,.
我要聽,.
而是我現在自己作為一個little god,.
我可以決定任何事物的是與非,.
對與錯..
三章第六節,.
於是那女人看到樹上的果子是好作食物,.
與人眼目是可喜愛的,.
能使人有智慧就摘下來吃,.
又給了她的丈夫,.
她的丈夫也吃了,.
很明顯不只是女人類事,.
男人本身在,.
故事就是說人類的代表,.
亞當夏娃,.
她代表整個人類,.
他們在這個關頭,.
他們吃了那個果子,.
也就是說他們自稱為帝,.
或自稱為神,.
他們以人的眼睛是明亮的,.

$^{361}$就知道自己是赤身老體,.
就拿無花果樹葉來為自己編作裙子,.
這裡已經去到一個地步,.
拿了那個果子,.
人與賜生命的主割蓆,.
斷交,.
這就是聖經所說的死亡,.
就是斷絕了關係,.
拿無花果樹葉來為自己編裙子,.
這其實是聖經裡面的故事,.
是有一些幽默感在,.
因為如果你看過無花果樹的果子的話,.
是遮不到什麼的,.
當然我們看到的圖畫,.
很多時候都是遮最重要的部位,.
就是那幾點,.
但你知道無花果樹葉,.
遮到外面,.
但遮掩不了我們內心的羞恥,.
如果你犯罪,.
你偷了東西,.
被人抓了,.
你可以將你偷了的東西物歸原主,.
就當沒事了,.
但你的良心受虧損,.
是不能夠撤銷的,.
不可以undo,.
如果你犯姦淫,.
在一個你應該忠於的關係裡面,.
你要破壞它,.
你是不能undo良心上的譴責,.
你可以掩著它,.
你良心受傷,.
或者關係的破裂,.
是不可以撤銷的,不可以undo,.
犯了罪,.
就一定會有後果,.
原本他們是享受神跟他們同行同在,.
一個很好的關係,.
一個很好的fellowship,.

$^{401}$但現在他們要躲起來,.
因為他們有羞恥,.
如果人跟上帝割席的話,.
得來的結果就是死亡,.
所以神與人之間,.
本身有的fellowship,.
現在就斷絕了,.
他們避開上帝,.
人與人之間亦彼此有隔膜,.
所以你見到這個故事是彼此抵賴,.
或者彼此控訴對方,.
如果故事繼續發展下去,.
你就知道人與人之間的關係有多糟糕,.
人與地之間的關係,.
本身有的和諧亦被破壞了,.
所以有雜草叢生,.
本身耕地是有利的,.
但現在不是不可以有利,.
不過賺錢就艱難,.
我不知道你有沒有留意這整段聖經,.
那裡有一個次序,.
首先神與男人說話,.
然後神與女人說話,.
然後神對蛇說話,.
然後又跟女人說話,.
然後又跟男人說話,.
中間那段,.
如果你來過這個教會就明白了,.
那個修辭就是想你留意最中間那段,.
有沒有留意神與蛇的說話,.
不是對話,沒有問沒有答,.
上帝只是pronounce judgment,.
宣判他的審判,.
他對蛇說,.
你既作了這事,就必受咒詛,.
比一切的生畜野獸更甚,.
你用兔子行走終生吃兔,.
其實就是說撒旦一定會被打敗,.
一定會失敗的意思,.
我又要叫你為女人彼此為神,.

$^{441}$你的後裔繼續下去有的generation,.
你的後裔與女人的generation也要彼此為神,.
女人的後裔要傷你的頭,.
你要傷他的腳,.
這是一個很黑暗的畫面,.
換句話說,.
將來會有一位女人的後裔,.
會跟撒旦有一個肉搏戰,.
一個close combat,.
而那個決定或決鬥的結果,.
都已經暗示了出來,.
女人的後裔將來有一位要你的,.
要傷你的頭,.
傷蛇的頭,.
你會傷他的腳筋,.
所以這很明顯不是在說生物學的問題,.
撒旦將來會有一位女人,.
跟他之間有一個搏鬥,.
十幾年前你可能看過一套電影叫做Passion of Christ,.
有沒有看過?.
基督受難曲,.
那套電影又好又不好,.
Mel Gibson拍的,.
你記得他開場時,.
有一個這樣的畫面,.
這個就是Passion of Christ,.
開頭的開場,.
你會看到那些畫面,.
或者他用的記號或代號,.
就是出自創世記第三章,.
蛇,即是撒旦,.
與一位女人的後裔,.
即是最後的理賽亞,.
他要擊敗撒旦,.
就是一場肉搏戰,.
那個畫面我們知道是在克西瑪利園,.
耶穌決定要去犧牲,.
按照父的心意,.
他要上十字架,.
為人的罪去死,.

$^{481}$然後死後父,.
一個肉搏戰,.
就是這裡所說的,.
這裡伊甸園,.
人犯罪之後,.
上帝去給人一個認許,.
就是說這個後裔將來會去打敗這條蛇,.
我們知道這些不只是記號,.
或者一些故事這麼簡單,.
我們每個人都經歷到,.
或者體驗到一些黑暗,.
或者邪惡的事情,.
或者一些真實,.
我們每個人都會看到有很多不義,.
在這個世界裡面,.
有很多罪惡很煩,.
也有苦難,.
我們生命裡面可能會有一些大病臨到,.
有些很大的挑戰臨到,.
可能會生癌症等等,.
有很多事情我們不能夠解釋,.
確實我們不能夠解釋,.
聖經也不是主要去解釋苦難的緣由,.
因為解釋了苦難的緣由,.
那個因素的時候都幫不到你,.
聖經告訴我們一個認知,.
上帝會戰勝一切的邪惡,.
從這裡開始,.
整本聖經就跟著這個脈絡去發展,.
所有之後的後代,.
以色列,波西,大衛,.
到最後耶穌基督,.
都是跟著這個第三章的脈絡,.
三章第二十節這裡說,.
他說,阿當被他的妻子起名叫夏娃,.
她就是眾生之母,.
然後耶和華神為阿當和他的妻子用皮紙,.
and musket來做衣服穿上,.
本身他們是用無花果樹的,.
現在上帝用皮紙來為他們穿上,.

$^{521}$這個你想一下,.
就是上帝給他們一個更好的遮蓋,.
如果有皮紙的話,.
換句話說,.
一定有一個生畜,.
一隻動物要犧牲,.
然後才可以有皮紙給他們穿上,.
對不對?.
這個就是整個救恩的故事,.
整個聖經大的故事,.
裡面所追溯的一個主題,.
就是上帝為罪人預備一個罪的遮蓋,.
一個救恩,.
如果你有聽過交響樂的話,.
你會知道交響樂很多時候,.
剛剛開始那一段,.
opening那一段,.
很多時候夾雜了一些主題,.
一些用語,一些句語,.
然後在下面的所有movement裡面,.
然後他會再play out出來,.
很多時候都會這樣,.
聖經的故事其實都是這樣,.
有一種伏線在這裡,.
告訴我們,.
上帝為人,你都要預備,.
我覺得最非比尋常的就是,.
人們叛逆了上帝說自己是,.
我設要稱臣,.
臣不是上帝的時候,.
他們被趕出伊甸園的時候,.
上帝沒有留在伊甸園,.
上帝居然出了伊甸園,.
跟他們同在,.
之後整個救若的故事,.
就是告訴我們,.
上帝怎樣用不同的方法,.
讓人能夠跟祂建立關係,.
有會幕的設立,.
有聖殿的設立,.

$^{561}$然後最後有耶穌基督裡面,.
十字架的救恩,.
讓我們能夠跟祂重修這個關係,.
正如我開頭所說,.
那個故事,.
好像浪子比喻的故事,.
不是要看表面的意思,.
而是看背後所表達的意思,.
出走的兒子,.
他可以很慘,.
遭遇很慘,很孤獨,很孤單,.
或者很抑鬱,.
有什麼答案給他?.
我們人想的,.
可能就會改善他的心情,.
這是心理學上的答案,.
或者他很餓,.
他去放豬,.
他被人欺負,.
這個可以是,.
你給的答案可以是一個經濟的答案,.
或者一個政治的答案,.
給他公義等等,.
這些所有的答案都重要,.
公義,有公平,.
或者有心理健康等等,.
但是你想想,.
這些答案全部都是橫面學生,.
聖經所給的答案,.
是vertical的答案,.
這是最終人需要的東西,.
不是說那些東西不重要,.
而是最終我們需要的,.
就是那個離家的半億的兒子,.
罪人,.
願意回到家裡,.
願意見到他的天父,.
跟他復和的關係,.
創世的第三章都是一樣,.
人最終需要的,.

$^{601}$就是我們跟創造主重修關係,.
今天你們來教會,.
可能你們來過一排教會,.
或者認識過信仰,.
聽過聖經的故事,.
你要知道,.
那個大的故事,.
就是上帝為我們預備的復修的關係,.
我們願意悔改,.
願意去接受主的,.
我們就跟他重建這個fellowship,.
我們同心共悔,.
主愛多謝你,.
因為你是那位愛我們的神,.
你創造萬有,.
但是你卻沒有丟下我們,.
縱使我們叛逆,.
縱使我們不認你,.
從創世以來到現在,.
你不斷的去拯救,.
你去修復我們罪人和你之間的關係,.
在耶穌基督的十字架上,.
祂是完美的羔羊,.
祂的犧牲,.
以至我們的生命被賣熟,.
能夠跟你重修關係,.
主讓我們跟你重新建立這個關係,.
在這個關係當中,.
我們在這裡跟你同行,.
我們這樣祈禱,.
奉主的名禱..
好,我在這裡回答一個問題,.
這個問題是說,.
如果亞當是代表整個人類,.
那麼到底神有沒有創造過多一個人?.
這個細節我不知道,.
但是如果你說亞當是整個人類的代表,.
聖經第一章,.
上帝創造人,.
按照神的形象所做,.

$^{641}$做人類出來,.
一群人出來,.
亞當是代表,.
我相信伊甸園裡面發生的事情,.
第二章發生的事情,.
就是把鏡頭放到這個代表的身上,.
可能有其他人在,.
但是那個人就是被亞當和夏娃來代表的,.
所以我們要知道,.
亞當其實就是人類的意思,.
不是一個人的意思,.
而是整個人類的代表,.
如果你這樣看的時候,.
我相信你會更加清楚聖經裡面以後發展的舊因,.
羅馬書第五章,.
在亞當裡面,.
我們全部都是死亡,.
與上帝相隔,.
但是在基督裡面,.
基督是我們的代表,.
我們就得到引領,.
我們就能夠成為神的子民,.
希望下一次,.
下個星期我們繼續..
多謝收看城匯網台.
記住like comment或者訂閱subscribe我地呀.
\newpage



\section{創世記 11:1-9}
\label{sec:Gl7QFa_KpBw}
\textbf{10/27/19 主日講道:與人同在的神 IV - 創世記 11\_1-9 巴別塔 (陳崇基牧師)}
\newline
\newline
連結: \href{https://youtube.com/watch?v=Gl7QFa-KpBw}{\texttt{https://youtube.com/watch?v=Gl7QFa-KpBw}} ~~~~ 語音日期: 2019-11-01
\newline
\newline
\hyperref[sec:0WFluTmE3pc]{\small{< < < PREV SERMON < < <}}
~
\hyperref[sec:index_chronic]{\small{[返順時目]}}
~
\hyperref[sec:index_scriptual]{\small{[返順卷目]}}
~
\hyperref[sec:xwcSwwS07wQ]{\small{> > > NEXT SERMON > > >}}
\newline
\newline
創世記 11:1-9
\newline
\begin{longtable}{cl}
\hline
\hline
章節 & 經文 (和合本修訂版)\\
\hline
11:1 & \begin{tabularx}{0.7\textwidth}{X} 那時，全地只有一種語言，都說一樣的話。 \end{tabularx} \\ \\ \relax
11:2 & \begin{tabularx}{0.7\textwidth}{X} 他們向東遷移的時候，在示拿地找到一片平原，就住在那裡。 \end{tabularx} \\ \\ \relax
11:3 & \begin{tabularx}{0.7\textwidth}{X} 他們彼此商量說：「來，讓我們來做磚，把磚燒透了。」他們就拿磚當石頭，又拿柏油當泥漿。 \end{tabularx} \\ \\ \relax
11:4 & \begin{tabularx}{0.7\textwidth}{X} 他們說：「來，讓我們建造一座城和一座塔，塔頂通天。我們要為自己立名，免得我們分散在全地面上。」 \end{tabularx} \\ \\ \relax
11:5 & \begin{tabularx}{0.7\textwidth}{X} 耶和華降臨，要看世人所建造的城和塔。 \end{tabularx} \\ \\ \relax
11:6 & \begin{tabularx}{0.7\textwidth}{X} 耶和華說：「看哪，他們成了同一個民族，都有一樣的語言。這只是他們開始做的事，現在他們想要做的任何事，就沒有甚麼可攔阻他們了。 \end{tabularx} \\ \\ \relax
11:7 & \begin{tabularx}{0.7\textwidth}{X} 來，我們下去，在那裡變亂他們的語言，使他們彼此語言不通。」 \end{tabularx} \\ \\ \relax
11:8 & \begin{tabularx}{0.7\textwidth}{X} 於是耶和華使他們從那裡分散在全地面上；他們就停止建造那城了。 \end{tabularx} \\ \\ \relax
11:9 & \begin{tabularx}{0.7\textwidth}{X} 因為耶和華在那裡變亂了全地的語言，把人從那裡分散在全地面上，所以那城名叫巴別。 \end{tabularx} \\ \\
[1ex]
\hline
\hline
\end{longtable}
$^{1}$各位觀眾,如果有聖經的話.
請你打開到《創世記》第十一章.
我們繼續這個系列.
《與人同在的神》.
我們先看聖經.
第十一章.
那時天下人的口音.
言語都是一樣.
他們往東邊遷移的時候.
在士那地遇見一片平原.
就住在那裡.
他們彼此商量.
內伯,我們要造磚,把磚燒透了.
他們就拿磚當石頭.
又拿石漆當灰泥.
他們說內伯,我們要建造一座城,一座塔.
塔頂要通天.
位要全讓我們的名.
免得我們分散在全地上.
又說江南要看看世人所建造的城和塔.
又說看下他們成為一樣的人民.
一樣的言語.
如今既作了者是來.
以後他們所要做的事.
就沒有不成就的了.
我們下去在那裡變亂他們的口音.
使他們的言語彼此不通.
於是又說使他們在那裡分散在全地上.
他們就停工,不做那成了.
又說在那裡變亂天下人的言語.
使眾人分散在全地上.
所以那成名叫巴別.
我們繼續這個系列.
我們今天去到創世記第十一章.
我想告訴大家.
我們由第一章看到現在.
很清楚,由第一章看到第十一章.
聖經的頭十一章.
就是想告訴你這個世界.
當罪進入這個世界的時候.

$^{41}$當罪破壞了上帝的秩序的時候.
會一直這樣.
一直在螺旋的.
走下坡,越來越混亂.
越來越混沌.
你記得第一章第一節.
上帝創造天地的時候.
他由混沌之中帶來的order.
由混沌裡面帶來的秩序.
是整整有罪的.
但是當罪進入了這個世界的時候.
有很多的關係就破壞了.
首先是神和人之間的關係.
人和人之間的關係.
還有人和地的關係.
我們華人就說天倫.
我們和上帝的關係.
人倫就是人與人之間的關係.
地倫就是我們和這個環境.
和這個世界創造的關係.
全部都被破壞了.
你見到一直下去的時候.
人因為有罪有邪惡進入了這個世界的時候.
他們就自相的.
彼此的敗壞.
有個self destruction.
或者自我的毀壞.
到那個地步上帝就要去.
藉著洪水去滅當時的世界.
然後再重整這個世界.
但是人仍然是繼續走下坡.
到最後最後的時候.
十一章就是我們今天看的聖經.
所謂巴別塔的聖經.
人就開始去挑戰上帝.
可以這麼說.
所以一路十一章裡面.
都是告訴我們有罪惡進入這個世界的時候.
我們不會認為會越來越好的.
好像頭十一章告訴我們.

$^{81}$如果上帝不介入這個世界裡面.
人只會繼續滅亡.
或者繼續崩壞.
暴力會繼續增加.
由家庭暴力開始.
然後到制度上的暴力.
政權上的暴力.
等等等等.
一直下去到第十一章.
這十一章其實不是一些很遠古的經文.
或者一些文件.
告訴我們當時世界是這樣的.
其實你看看現在.
我們的世界也是這樣的.
你環顧一下世界裡面有哪個地方.
你聽一下新聞.
聽一下社會裡面的一些訊息.
很多的地方都是繼續去崩壞下去.
創世的頭十一章就是告訴我們.
我們要去這樣看待這個世界.
若非上帝有恩典.
你都去介入的話.
這個世界永遠都會繼續走下坡.
上個星期在英國有一個很大的新聞.
三十九條屍體在貨櫃裡面找到.
大部分是中國的公民.
聽到這樣的新聞的時候.
你真的覺得這個世界真的有問題.
世上的國度.
世上的國家都很多有問題.
我看了其中一個的遇難者.
一個女孩子.
她應該是越南華僑.
越南籍的.
她臨死之前.
她短信給她父母.
跟他們說對不起.
因為想偷渡去另外一個地方.
有更好的生活.
有一個這樣的美夢.

$^{121}$或者有一個更好的前途.
就這樣斷送了.
你看到這個世界.
就是創世記第十一章裡面所說的.
就是會越來越走差.
我們看看這段聖經怎麼說.
它說那時的人的口音.
言語都是一樣的.
他們就往東遷移了.
在創世記頭十一章.
凡是往東遷移都不是好事.
你記得亞當夏娃.
人類的代表.
他們被踢出伊甸園之後.
他們背叛上帝.
被踢出伊甸園之後.
是往東走.
然後他們有家庭暴力.
該人殺他們的兄弟.
被耶和華趕去流亡.
也是往東的.
到第十一章這裡往東走.
他們在士那地遇見一個平原.
就住在那裡.
他們彼此商量.
就說我們要做磚.
用一些燒透的磚.
然後當成石頭.
用石漆當成灰泥.
然後他們就說建一座塔.
或者建一座城.
然後這個塔艇要通天.
我們要傳揚我們的名.
免得地上我們分散.
如果你看見這些字眼的時候.
你就很明白.
是當時巴比倫帝國的字眼.
士那地.
到南邊.
米索不達米亞南邊.

$^{161}$即是現在南邊的伊拉克.
或者科威特那個地方.
這些全部都是巴比倫帝國.
或巴比倫文化科技的字眼.
燒透的磚.
這些不是巴勒斯坦地會做的事.
因為巴勒斯坦地是用石.
到處撿石來建屋.
但是在那個地方.
南非密蘇古塔米亞.
就一定要做磚.
因為那裡沒有什麼石頭.
這個地方有一點像我們今天.
中西部.
芝加哥.
是平原一帶.
這個就是南非密蘇古塔米亞.
在科威特.
或者南伊拉克的地方.
很平坦.
但是他們做一些城的時候.
他們去建一些城鎮.
一些城市.
他們就會做一些.
所謂我們今天叫摩天大廈.
或者上一代叫摩天大廈.
那些高樓大廈.
你今天開車去芝加哥.
看著下層的時候.
那些摩天大廈.
會不會霧起.
你可能去別的地方.
但你都被它吸引.
就去下層.
在大平原.
然後有一些高的建築物.
這個也是巴比倫帝國出名的地方.
在巴比倫帝國裡面.
無論是古巴比倫帝國.
或者所謂的紐巴布隆年.

$^{201}$當時猶太人被流亡的地方.
叫做紐巴布隆年帝國.
新的巴比倫帝國.
都好.
他們很出名的就是一些建築.
一些建築叫做Ziggurat.
這些Ziggurat.
其實可以譯作廟塔.
他們在那些大平原上.
就會建起來.
我昨天去搜尋這些畫面的時候.
有個人在Minecraft那裡.
建了一個廟塔.
挺厲害的.
可以走進去.
幾乎是一比一.
不知一比幾.
他真的建了一個很真實的廟塔.
叫做Ziggurat.
現在可以看到有幾十個.
這樣的塔在當時的古城那裡.
這些廟塔是做什麼的呢.
有些高起來有十層樓那麼高.
是一些巴比倫市中心.
他們downtown的建築.
廟塔的功能是什麼呢.
就不是讓人上去上面.
去所謂上香油.
不是上去上面去拜的.
而是這個廟塔下面.
就會有那個temple.
有那個神廟在那裡.
那麼廟塔是做什麼的呢.
廟塔其實就好像樓梯那樣.
好像給總統下飛機的那個樓梯.
其實就是建一個那麼高的廟塔.
就是給上面的神靈.
讓他能夠下來的.
所以廟塔上面.
不是很多人上去.

$^{241}$間中才有人上去.
那些祭司會上去.
因為廟塔上面有一個房.
有些小的可能一百多呎.
有些大的可能五百多六百呎.
裡面有傢俬的.
有桌子.
有床.
也有一些食物會放在那裡.
其實就是那個神靈的airbnb.
每逢需要他下來的時候.
希望建一個這樣的廟塔.
一個這樣的樓梯.
以致那個神靈方便他.
能夠下來地上.
然後他下來廟塔下面.
那個廟那裡.
他能夠去接受人的拜祭.
或者人的祭祀.
所以是他的房間.
或者廟塔通常的名字.
都是天與地之間的connection.
那個大門.
那個聯絡.
這個就是我們一路.
從第一章開始到現在.
我們說當時的創世記寫成的時候.
那些巴比倫或者外邦的神話.
很多時候都是描述那些神靈.
是有需要的.
好像我們人有需要.
所以我們人會累.
我們需要休息.
你不要擾我清眠.
我們人會肚餓或者會口渴.
我們就需要有飲有食.
跟人的形象是一樣.
那些神靈的神話裡面.
類似是.
如果你是拜這些神靈的話.

$^{281}$其實就是.
英文叫做.
You scratch my back, I scratch your back.
就是你握我的背.
我又握你的背.
就是我去服侍你.
你又會服侍我.
所以是一個.
互相護理的一個關係.
我給你吃東西.
然後你就會祝福我.
你就會待我好.
創世的時候一張巴比塔就是一個這樣的背景.
所以就弄了一條梯子.
就想神下來.
誰知道第五節.
耶和華真的下來.
耶和華降臨.
要看世人所見的城和塔.
耶和華說.
看他們成為一樣的人民.
有一樣的言語.
如今所做起來的事.
以後他們所做的事.
就沒有不成就了.
他們繼續這樣敗壞下去的話.
沒有一樣東西是不會繼續敗壞的.
就是這個意思.
我們下去.
在那裡變亂他們的口音.
使他們的言語彼此不通.
於是耶和華從那裡分散傳遞.
他們就停工了.
就不做這個城.
因為耶和華在那裡變亂了天下人的言語.
將眾人分散在傳遞.
那個地方就叫做巴比.
你知道聖經的故事.
有些幽默感.
他們說塔頂要通天.

$^{321}$但耶和華要下來看的時候.
他也要下來.
不是很通天.
樓梯不夠長.
或者耶和華下來的時候.
他們以為這樣可以.
讓樓梯的神靈下來.
然後接受他們的供奉.
但耶和華真神下來的時候.
他就不接受他們的供奉.
他說這是不好的.
所以他下來的時候就變亂他們的口音.
然後讓他們彼此說不通.
變亂他們的世界.
變亂他們的帝國.
以致他們的帝國瓦解了.
這個就是巴別的意思.
在巴比論文裡面.
巴別就是他們的名字.
但希伯來文聽起來就是混亂的意思.
所以有一個諧音在這裡.
按照聖經這樣看的時候.
他們以為這樣合一.
以為這樣可以建立他們的帝國.
誰知道其實就是瓦解的.
是自己來去作破壞.
所以巴別塔是象徵什麼呢.
巴別塔代表人到了一個地步.
他想操控上帝.
你建立一個樓梯.
其實就是想控制神.
因為我這樣去服侍他的話.
他就可以這樣祝福我.
神人的關係就完全顛倒了.
因為是我操控他.
我操控著上帝要怎麼做.
我是操控著上帝.
上帝成為我的奴隸.
所以誰是上帝.
你看巴別塔所象徵的.

$^{361}$就是人就變成上帝.
因為人是在操控著神.
這個你回想起創世記第三章.
所謂人的墮落那段故事的時候.
你就知道亞當夏娃.
人類的代表他們犯的罪就是這個罪.
就是顛覆上帝和人之間的關係.
人就變成神.
神就變成人的奴隸.
這個就是最大的走下坡.
神人之間關係上的顛倒.
或者應該這樣說.
這個也是創世記批判當時帝國.
自以為是的觀念.
我們這樣操控上帝的話.
我們自己就會得到利益.
只要我們有一個這麼高的塔.
我們就能夠控制天上的神靈.
這個就是最大的叛逆.
或者應該是當時社會裡面最大的問題.
你以為你控制得到.
你以為這樣我們就能夠成為上帝.
我們就控制了所有東西.
但其實就是混亂.
是巴別的一件事.
不久之前.
有一個人在社交媒體上.
他發了一段說話.
在清末的一個神宰.
叫做曾國藩.
曾國藩說什麼呢.
他說社會大亂之前.
一定有三個先兆.
第一.
就是無論何時都是黑白不分.
第二.
善良的人越來越客氣.
越來越謙卑.
而沒用的人就越來越猖狂.
第三.

$^{401}$就是當問題嚴重到一個程度的時候.
偏偏都是將凡事合理化.
一切都是默認.
莫名其妙地去應對一番.
當是沒有看到.
社會的制度崩壞.
社會文化的崩壞.
有任何破爛的東西.
我們當看不見.
社會大亂就有三個前兆.
這個其實就是巴別塔的故事.
就是巴別倫帝國未覆亡之前.
寫了一個這樣的故事.
你叫他去影射也可以.
你叫他去一個社會的批判也可以.
就說巴別倫帝國是不是真的這麼厲害.
以為自己已經塔頂通天.
我們可以控制上帝.
想做什麼.
巴別塔就是說人是為了自己的利益.
我們來操控上帝.
其實骨子裡面就是我以為我就是上帝.
我以為我自己就是神.
巴別塔代表什麼.
不單止是社會的一個問題.
不單止是一個國家的問題.
也是我們人生命裡面的問題.
我們和神的關係.
巴別塔想說什麼.
是不是我們為了自己的利益.
我們不知不覺.
或者刻意的.
我們想去控制上帝.
我們未必是實物.
我們想做一條梯出來.
但我們心裡面已經做了一條梯.
我是牧師.
我有時和人淨分.
有時也是一些不去教會的人找我淨分.
其實我會很拗頭.

$^{441}$你又不信 你又沒有信仰.
為什麼你需要一個神職人員.
來祝福.
你當我是道士.
我去做法事就好一點.
有一點這樣的感覺.
但我明白文化的傳統是這樣.
有時真的可能會是這樣.
但要想深一層.
為什麼需要有一個上帝的祝福.
在骨子裡面很多時候是這樣的.
就是我要在婚禮裡面有上帝.
但我的婚姻.
上帝你不要作主.
最好我的婚禮搞得好看一點.
找個好看一點的牧師.
或者說話好聽一點.
那我有個happy wedding.
但婚禮結束後 上帝在哪裡.
上帝不要指望我.
將上帝變成我家庭的中神.
不不不.
我只是希望上帝來祝福一下我.
像不像芭比塔.
我叫我的子女去教堂.
捉他去教堂.
或者我叫我子女去基督教的學校讀書.
因為那裡是教他們得育的.
教他們做一個好人.
教到一個好的兒子或者好女兒.
這就是我的指望.
但你不要主意在這裡改變我生活的價值.
這就是芭比塔.
就是我想去操縱上帝.
我祈求上帝讓我家臣能夠升職.
讓我能夠飛黃騰達.
但不要指望上帝改變我的習慣.
他應該加薪.
應該讓我能夠飛黃騰達.
但如果上帝要求我做任何事.

$^{481}$無論是我奉獻或是我工作的職場上.
我需要去服侍著.
我不會期望上帝會改變我.
你想想我們很多時候祈禱.
上帝讓我健康.
你治好我好不好.
現在我大能臨頭.
我有重病.
你治好我.
但治好我之後.
你不要意圖改變我對你的看法.
好像你去看醫生.
看完醫生你會不會和醫生做個更好的朋友.
未必的.
看完醫生之後你就離開.
我們和上帝祈禱很多時候都是這樣.
上帝你救我.
我們一味叫上帝救我.
上帝救了我.
上帝你不要意圖改變我.
你給我健康.
我自己的這些.
我做我自己喜歡做的事.
上帝你給我男朋友或女朋友.
但給我之後你不要意圖要我順服你.
我奉獻給上帝.
我甘情願奉獻給上帝.
上帝應該保佑我.
上帝你不要作主.
是我作主.
我事奉上帝.
我想交換什麼.
上帝你給我態度.
你給我一些好處.
但你不要意圖改變我對自己的看法.
這些全部都是巴別行.
所以不是說社會不是說帝國這麼簡單.
是說我們的帝國.
我們的人生.
我們的世界觀.

$^{521}$如果我們這樣看的時候.
難怪耶穌曾經在他登山寶訓裡面.
結束的時候這樣說.
他說凡稱呼我主呀主呀的.
不能都進天國.
唯獨進行我天父子兒的才能進去.
當那一天必有許多人對我說主呀主呀.
我不是奉你的名來傳道.
奉你的名來趕鬼和行異能.
耶穌說我明明是高少丹.
我從來不認識你.
你們這些作惡的人離開我去吧.
這段說話不能再強調.
如果真的有永恆的話.
這是最強烈的可能性.
他告訴我們.
你稱呼主呀主呀的時候.
你不要想著你一定是入硬去.
或者你一定是屬於天國.
反過來說.
這班人所謂作惡的人.
並且奉過主的名來行事.
做了許多事奉的.
這班人他們不就是當自己是主.
或者他們和上帝的關係是.
你不是說恩典嗎.
我們拿出特朗普卡.
我們說的是恩典.
什麼都是恩典.
不關我做過什麼事.
你們看清楚耶穌說什麼.
凡是唯獨進行我父子才能進去.
不要動不動.
我們拿出底牌.
你講恩典.
所以我生活裡面.
上帝改變我的習慣.
我的態度.
所謂上帝是天國要臨到我身上的時候.
這些我們不算.

$^{561}$我們算恩典.
這裡好像不是這麼說.
這是最強烈的可能性.
你這作惡的人.
避開我.
不要搞錯.
我們去跟隨主.
我們信上帝.
我們事奉上帝.
不是上帝非要你不可.
沒有這樣的事.
由創世第一章告訴我們.
上帝是不需要我們的.
上帝也沒有欠過你一丁點的.
上帝欠過你什麼.
是否欠你健康.
非給你不可.
上帝沒有欠過你任何東西.
事奉是因為上帝救了你.
你應該感恩回應.
你應該巴不得去事奉他.
上帝供應我們一切需要.
救我們脫離死亡.
給我們有永生.
無限好的東西給你.
都不是因為上帝欠你的.
是上帝用恩典施恩給你.
要講債的話.
是你欠他債.
上帝欠我債.
巴別塔就是這個意思.
你欠我的.
現在我供奉了你的東西.
我上了香油錢.
你這麼有恩典.
你應該祝福我.
這就是操控上帝.
想操控上帝的宗教.
這是完全顛覆主和人的關係的宗教.
一種假的宗教.

$^{601}$只會叫你滅亡的宗教.
這就是創世記一至十一章.
有的伏線.
你清清楚楚要知道.
是人去摧毀上帝的世界.
是人和上帝割席.
我們才會弄成這樣.
但是上帝既然有恩典.
你就好自為之.
想清想楚.
世界走下坡的時候.
有一個伏線.
上帝既然知道罪人.
來摧毀他的世界.
他沒有放棄.
那些叛逆者.
他願意去施恩.
叛逆者是欠他的人.
他願意去施恩給他.
亞當犯罪.
上帝用皮紙.
用一個動物的犧牲.
來遮蓋他的罪.
人在這段外面流亡.
沒有上帝和他的關係.
但是上帝出來.
這些人和他們來同行.
人背叛上帝.
去引入撒旦.
是顛覆上帝的秩序.
上帝應許.
說我會戰勝邪惡.
但是你要站在我那邊.
因為女人的後裔會打敗.
射撒旦,邪惡.
你看看創世記頭十一章的家譜.
其實就是告訴我們.
上帝一直一直有條線索.
就是他願意用恩典.
用信神來保護.

$^{641}$創世記十二章.
就是上帝主動來救贖所有的罪人.
叫亞伯拉罕出來.
他說你離開本地本族本家.
往我指示你的地方.
我要使你成為大國.
我要賜福給你.
最後萬國要因你來得福.
你要令到上帝的名出名.
不是你的名出名.
巴比塔就是要令我的名出名.
因為我操縱上帝.
真正上帝的國.
就是要上帝的名來得出名.
他說你要離開.
其實意思就是說.
不要猶豫不決.
Go.
意思就是這樣.
Go.
You.
Go.
離開你本地本族本家.
當時社會,政治,個人.
我們的靠山.
亞伯拉罕要離開那裡.
使萬國能夠得福.
離開哪裡?.
就是巴比塔.
因為你見到他呼召他的地方.
就正是巴比塔的地方.
亞伯拉罕要離開那裡.
離開他最安恕的地方.
永不回頭的他要跟著上帝.
上帝是他的神.
他是他的老婆.
這個就是聖經的信仰.
不要弄錯.
有人說在英文的文學裡面.
有兩種故事.

$^{681}$有一種故事叫探險的故事.
一種叫adventure.
有一種故事叫探索的故事.
叫quest.
什麼叫adventure?.
什麼叫quest?.
adventure就是我們小時候看的探險的故事.
發發勇勇走了去別的地方.
那裡很有趣的.
有森林,野獸等等.
跟他玩完一輪之後我們回來.
回來後媽媽說吃飯了.
我們就去吃飯.
這些就是兒童漫畫裡面的adventure.
有什麼特色呢?.
你去了一個地方.
你經歷了一些事情.
然後你回到原位.
你是完全沒有改變過.
這個叫adventure.
我不管你哪個世代.
以往看的那些.
我的世代什麼Pyotr Ilyin.
什麼IQ博士.
我們去完玩完.
然後我們回來的時候.
看完就算了,很過癮.
但我的生命沒有改變.
但在文學裡面有一樣叫做quest.
quest有什麼不同呢?.
就是你去一個旅程.
你上路之後.
你是不會再回來.
或者甚至.
就算你回來之後.
你都不認得你自己.
因為你的生命會因為這個旅程.
而完全改變了.
The Law of the Ring就是這種故事.
基督教信仰不是adventure.

$^{721}$不是去完回來我自己沒有改變過.
基督教信仰是quest.
是你願意上路去跟上帝.
而在這個旅程裡面.
你是完全改變了.
縱使你回到家裡.
你個人已經改變了.
基督教信仰是quest.
是你願意讓上帝改變.
這個才能夠使萬能得福.
這個才是上帝自己的計劃.
聖經信仰不試探險,試探索.
是一個人生永不回頭的旅程.
這個才是真的信仰.
巴比塔就是假的信仰.
今天讓我們提醒.
我們不是為了自己的利益.
不是想去操縱上帝.
而是上帝要改變.
這個才是真的信仰,真的宗教.
真理,我們同心協力地祈禱.
\newpage



\section{}
\label{sec:xwcSwwS07wQ}
\textbf{11/3/19 主日講道: 圍繞父母親的見證人 (萬楊婉貞博士)}
\newline
\newline
連結: \href{https://youtube.com/watch?v=xwcSwwS07wQ}{\texttt{https://youtube.com/watch?v=xwcSwwS07wQ}} ~~~~ 語音日期: 2019-11-07
\newline
\newline
\hyperref[sec:Gl7QFa_KpBw]{\small{< < < PREV SERMON < < <}}
~
\hyperref[sec:index_chronic]{\small{[返順時目]}}
~
\hyperref[sec:index_scriptual]{\small{[返順卷目]}}
~
\hyperref[sec:z_G8mnyC7Vg]{\small{> > > NEXT SERMON > > >}}
\newline
\newline
$^{1}$《希伯文》第十二章 第一章 第十二章.
今天正道的經文是在詩篇127篇第三節和《希伯文》第十二章一至二節.
我們先看聖經.
(書面文件).
我們歡迎張老師.
讓我先收集一下氣氛和環境.
各位大兄姐妹 各位朋友早.
今天想和大家分享的經文是穆斯所說的.
題目是「圍繞在父母親周圍見情人」.
起初是少許的澄清.
因為很多時候我們去崇拜聽講道的時候.
一聽到父母親的三個字便會覺得與我無關.
因為我又不是父母親 我又沒有子女 我是單身.
所以特別在此提醒 圍繞在父母親一班見證人.
即是你和我也有份.
希望大家把時間放在心中聽一聽神的教導.
我們在詩篇127篇第三節.
是一個我們比較熟悉的經文.
大家也知道是再一次提醒我們.
神把子女作為產業給我們每一個父母.
現在是禮物 其實說得很清楚.
神是擁有權在神身上.
祂給我們每一個父母 祂只希望我們去看管 去照顧.
最終是我們回應神的交代.
所以神再次提醒我們.
既然是這樣的時候 神不是要求我們操控.
也不是要求我們強權去管理.
也不是說我們只要保護 要眼不離 手不停地去黏著他們.
那不是神給我們產業的上次作用.
主要是告訴我們 如果我給你的時候.
希望父母怎樣去疼愛 去教導 去放手.
所以放手在我們中國人的文化裡面是比較困難了解的.
我忘記了 是否應該回到我的PowerPoint.
按一下就可以.
太快了 明白.
因為我發給Mark Sir的PowerPoint 底色不是這樣的.
所以我就問這是不是我的PowerPoint.
他竟然把我喜歡的紫色換了.
Mark Sir第一個禍就闖在這裡.
既然說到神的產業.

$^{41}$也告訴我們 每一個人的產業次數是不同的.
有時會有很多產業次數給我們.
有時會遲些才給 有時會早些給.
有時會跟一些不是親身的路徑去給.
可能是營養或其他.
所以很多時候產業次數的渠道不是很劃一.
但我們仍然要求 仍然是.
我們今天聽神的話語的時候.
就是有一個心態 產業是神的寶物.
是神的恩典 是神的賜予.
我們作為父母在跑道上 繼續做我們的角色的時候.
我們有什麼要留意.
今天就用這個大前題來看看.
我們看希伯來西第十二章第一節.
我們有很多很多的見證人像雲彩圍著我們.
其實大家看經文的時候.
可能會留意到最近馬拉松全馬剛結束.
當時有些人去參加.
當時的經文給我們一個畫像.
我們當時希伯來西第信徒他們在跑道上跑.
不斷地跑的時候 周圍就有很多人在.
我們看全馬的時候.
在旁邊加油 給水 叫他們有些醫療設備照顧他們.
或者在終點的時候 怎樣跟他們歡呼.
所以這個畫像是這樣給出來的.
今天我們嘗試將父母親每一個人放在畫像裡面.
我們在奔跑這條路.
當我們奔跑的時候.
我們周圍的見證人對我們有什麼提示.
怎樣幫我們繼續做父母親的角色.
所以我們嘗試放的就是.
當時是我們父母.
父母的時候我們向著哪個目標去跑呢?.
就是我們向著主耶穌.
所以我們現在有個畫像.
就是我們一班父母我們向著主耶穌去跑的時候.
我們周圍最外圈的見證人是些什麼人.
在當時來說 希伯來西的作者告訴當時的信徒.
你們一定要很忠心地跑的時候.
我們周圍的見證人其實就是希伯來西的十一章的見證人.

$^{81}$舊約裡面的先賢.
所以我們簡單來說.
十一章的先賢有亞伯拉罕,雅各,約瑟,拉合,羅亞.
一直有很多舊約上值得我們去學的榜樣.
這些先賢我們也概括一些.
他們有什麼我們要學.
拉罕因為信離開不易.
他對神的信心.
約瑟在主墓的試探.
他很清楚說我不可以做這些事.
不可以得罪神.
他對神的心智.
我們也看到.
例如基里拉合怎樣去收藏兩個探子.
因為他經歷過這個信是真的.
所以他可以很勇敢去抗拒.
大義在獅子坑口那種至死不改的忠誠.
所以十一章是告訴我們.
十二章的信徒.
當你們繼續做信徒的時候.
你看著十一章這群所提的先賢的榜樣.
如果我們放了我們父母親在這個跑步上.
其實也提醒我們.
過往的先賢他們有很多榜樣給我們做.
我們怎樣去學他們.
至死忠心.
怎樣學他們忍耐.
學他們對神的信.
學他們的態度.
學他們不要做錯事得罪神.
但也告訴我們.
今天我們父母親在跑步走的時候.
我們周圍的外圈除了先賢之外.
其實還有甚麼人呢?.
就是我們每一位弟兄姊妹.
因為其實是你們每一位.
你們跟神的關係是怎樣.
你們可不可以在你們事奉神的時候.
有做到先賢的榜樣.
那個盡忠那個盡心.

$^{121}$你們可不可以作為一班父母親.
當他們教子女的時候.
他們會看著你們怎樣做.
如果用回我們城會的現場的述職.
當你們做父母教子的時候.
你可不可以告訴他們.
你看看鎮惑敬拜小組的哥哥姐姐.
很勤力去練師的.
因為他們想帶我們父母去崇拜神.
讓我們更加親近神.
你們可以在長大的時候.
你們可不可以學這些哥哥姐姐.
這樣去親近神去崇拜.
剛剛說有音樂的火堂.
你們用這些才能去事奉神.
或者我們會說.
其實你們看我們的歡迎小組.
教會的小組.
你們進來的時候.
很多小食準備好.
他們很早就準備這個地方.
讓我們可以安靜地去事奉神.
所以你們要珍視這個教會.
有時候不要弄得地方亂.
因為有一班叔叔姨姨.
很辛苦地準備這個.
我們的事奉.
可不可以令到一班父母.
教他們的子女的時候.
會用這些說話去教他們.
甚至可能當中你們說.
有些人沒有很喜眼的事奉.
但是你的態度.
你的心智仍然是做了一班.
父母的教子女社的見證人.
你是不是可以很準時地.
進來教會去崇拜.
很安靜地坐在那裡.
或者你們參加教會.
除了禮拜天之外的活動.

$^{161}$這樣的聚餐.
Outreach聚餐.
你們會不會邀請人回來.
你們每一位弟兄姊妹所做的事.
其實是告訴當中的父母親.
當你教子女的時候.
你們不是孤單的.
你們有同路人的.
你們可以告訴他們.
看看其他人也是這樣做.
其他人也是這樣去親近神.
所以這件事是.
第一個提醒就是.
其實我們每一個人是有份的.
我們作為弟兄姊妹.
怎樣去做一班父母親的見證人.
然後我們再看.
除了外圈之外.
接著那個圈.
是甚麼可以成為父母親的見證人.
當然很清楚.
內圈是我們的配偶.
無論你是爸爸媽媽.
就是你的太太或先生.
因為我們很多時候.
我們夫婦之間.
父母之間互相的相處.
是很影響我們怎樣教導子女.
子女是看著你們怎樣相處.
看著你們之間和神的關係.
看著你信不信這個神.
看著你們很多新教就在那裡出來.
我們都很清楚知道.
夫婦之間的相處.
就是七天見面.
二十四小時.
不可能做到完美.
不會做錯事的處境.
但我們想想.
希伯留書十一章.

$^{201}$上了信心龍虎榜的舊約的先賢.
不是因為他們多完美.
他們不是因為完美才上龍虎榜.
他們純粹是因為他們對神的忠.
對神的信.
所以同樣在配偶之間.
你們怎樣互相做見證人.
去鼓勵我們做好父母親的角色.
就是你們之間.
你自己的新教和神的關係.
一個比較遠的陳年畫像.
第一就是在我的心裡.
我的朋友就是那時候.
我和我先生.
我們在美國旅行.
很喜歡開車.
開車的時候就會住酒店.
當時我還記得.
我早上起床睜開眼.
我先生就坐在圓台上.
自己做靈修.
其實那件事一直對我提醒.
其實重點不是我先生.
比我早起床做靈修的畫面.
重點是他的新教.
他告訴我他和神的關係.
重點不是說.
你為什麼這麼懶.
去旅行而已.
為什麼你和神的關係不太好.
重點不是他告訴我的.
重點是他自己怎樣做.
而我就觀察到.
那樣影響我怎樣教導子女.
所以實際來說.
作為父母的提醒.
當我們教小朋友的時候.
新教是什麼.
新教不是說.
你怎麼這麼不懂得教.

$^{241}$你自己也不和神的關係搞好.
你怎樣教導子女.
你怎樣做團體的導師.
你怎樣做主要學的教師.
新教不是這樣.
特別是在配偶之間.
新教是當我們和神有關係的時候.
我們會分享.
今天我在神裡學了什麼.
或者分享.
我今天學不到這些東西.
我經常都失敗.
我經常都做錯了.
那是將自己的脆弱.
在你配偶之間分享.
而那個分享.
你就是作為一個父母親的見證人.
見證你和神的關係.
告訴他應該要這樣教.
才可以做到一個好的父母親.
所以作為父母親.
當我們一路.
原來畫像就是在希伯來信.
告訴他.
你不斷奔跑的時候.
你周圍的見證人.
有舊日的先賢.
有今天的弟兄姊妹.
去幫助你們.
和你同行.
有你的配偶不斷提醒你們.
還有什麼可以做我們的見證人呢?.
這裡是不繼續講.
再一次提醒.
就是在配偶之間.
其實那個說話.
那個新教.
如果在《聖經》中.
是曾賢告訴我們.
他說一句話說得合而.

$^{281}$就好像金蘋果在銀網子那裡.
意思就是你講的話.
你懂不懂得去鼓勵.
你懂不懂得去提醒.
你懂不懂得去刺激他們.
去振奮你的配偶.
讓他們當他們做父母親的時候.
是做得很舒服.
很被肯定的.
另外在《真言三十一章》也告訴我們.
他家丈夫心裡倚靠他.
必不缺少利益.
很多時候倚靠在做父母之間.
不是說你做完.
你教完.
你管完.
我在外面上班.
家裡的事.
教兒子的.
你的責任.
那個倚靠不是關人.
那個倚靠其實是對你的配偶是信任的.
無論你是爸爸.
或者你是媽媽.
你都是信任那個.
那個配偶可以做到.
你都是肯定他.
用他的方法去管教.
所以那一次再次告訴你.
在配偶之間.
你的身教是可以影響大家.
去做好父母親的角色.
所以你是互相之間的見證人.
我們再退下去.
我們說的就是.
兒女是我們的見證人.
我們看希伯來書的經文.
在第一節說得很清楚.
就是告訴父母.
或者當時的信徒.

$^{321}$你在跑路的時候.
你要怎樣去.
第一樣的提醒是什麼呢?.
你要身教.
就是你要放下各樣的重擔.
在經文裡.
什麼叫做放下各樣的重擔呢?.
其實在教兒女的時候.
兒女是我們的見證人.
為什麼他可以做到我們的見證人?.
其實因為他學了什麼?.
他學了父母親怎樣教.
他學了父母親和神的關係是怎樣.
他學了父母親對神的信是怎樣.
所以其實是一個相輔相成的.
我們怎樣做好我們和神的關係.
我們教導我們的兒女.
而他們的表現.
也是見證我們.
告訴我們.
這樣做是對的.
這樣做是OK的.
這樣做是合神心意.
是土神喜悅的.
所以第一樣的提醒就是.
周圍父母親.
我們要學的是.
要放下各樣的重擔.
我想父母親的重擔是什麼呢?.
因為我們現在是.
很多時候父母親的重擔.
父母親一定會做的事是什麼呢?.
你會關心.
作為一個好的父母.
會關心你的兒女.
其實我們有很多事情關心的.
我們會關心他的學業.
關心他的性格.
他的健康.
他的信仰.

$^{361}$他的人際關係.
我們會關心.
他上學讀書.
他的老師懂得教嗎?.
他教得好不好?.
他教的東西對不對?.
有沒有教錯?.
我們甚至有些父母會踩過界.
會教那個老師.
你應該要這樣選擇.
應該要這樣說.
但是我們作為父母.
我們會關心.
我們另外會關心什麼呢?.
我們會關心兒女的人際關係.
我們會關心他們是否合群.
我們中國人很喜歡關心的就是.
見到人的時候有沒有禮貌.
你懂不懂得叫人.
這件事是很重要的.
我也聽過有些父母.
指責其他小朋友.
這麼沒禮貌.
不懂得叫人.
那些可能不是小朋友沒禮貌.
每個小朋友的個性不同.
有些比較主導.
社交很好.
有些比較內向.
但是這件事我們作為父母會關心.
我們關心其他人.
怎樣看待我們的子女.
我們關心他的性格.
他和一群人相處的時候.
是否沒所謂的.
因為他比較和順.
朋友說這樣就這樣.
我們關心他會否被人欺負.
我們關心他這麼不主動.
不主動說話.

$^{401}$你的要求會否吃虧.
我們又關心他如果太兇.
會否欺凌別人.
會否變成小霸王.
我們很關心他們和別人相處.
我們關心他的健康.
睡眠.
有沒有喝湯水.
如果他們睡得多.
我們又會關心他.
因為病嗎?.
懶惰嗎?.
我們很多事情關心.
我們關心他的物質夠不夠.
他的電話.
是否保持他同輩.
如果他的電話不夠好.
會被人取笑.
所以我們很多事情關心父母.
這些事情是對的.
但很多時候.
聖經告訴我們.
要放下重擔.
意思是我們關心.
但如果我們關心過火的時候.
變成擔心.
擔心會壓著我們的重擔.
我們一方面教導子女.
要信神.
甚麼都要交託.
但另一方面.
如果我們的子女看到我們的信.
是否放手的信.
還是懷疑的信.
看到父母要信神.
是靠自己的能力解決問題的信.
還是真的拉著神的手的信.
所以我們要放下各樣重擔.
我們真的信的時候.
教導子女.

$^{441}$用自己的身價教導他們.
怎樣去信.
彼得前書也提醒我們.
要把一切憂慮卸給神.
因為祂顧念你們.
所以這件事告訴我們.
要放下我們的重擔.
見不見到剛才說過.
在希伯來世第十一章.
也有提到.
其中一個先賢.
信了龍父榜.
就是亞伯拉罕.
他怎樣憑信離開吾爾.
他在老年的時候沒有兒子.
但他信甚麼?.
信神給他一個大國.
天上的星和地下的沙.
所以這件事告訴我們.
作為父母.
有先賢教導我們.
做我們的見證人.
教導我們怎樣放下去信.
我們要周圍的弟兄姊妹提醒我們.
怎樣去信.
我們自己教導子女的時候.
是否可以教導他們.
怎樣放下重擔.
所以這是第一件事.
作為父母親要學的.
要怎樣放下重擔.
另外在經文中.
希伯來世第一二章也說.
我們要學習脫下容易纏累我們的罪.
其實容易纏累我們的罪.
一定不是大罪.
一定不是十家都.
不要做不要貪.
不要殺人.
不要作假戰證.

$^{481}$一定不是大罪.
因為大罪我們很清楚知道.
我們很清楚知道就很容易改.
所以纏累我們的罪是甚麼罪呢.
是一些不是很清楚.
但和我們糾纏中的.
而這些糾纏令我們不小心.
就會犯罪.
就會做錯事.
所以在父母親的立場來說.
當我們做父母親的時候.
那些是甚麼罪呢.
我記得說過.
我們要學放下個人的重擔.
因為我們的關心變成擔心.
我們的擔心一直做下去的時候.
就出現了一些焦慮.
焦慮令我們很不舒服.
很緊張.
焦慮有甚麼行為會出現呢.
因為我們焦慮.
我們關心,擔心.
所以很多時候我們會出現甚麼呢.
我們會埋怨.
我們會指責.
我們會憤怒.
父母親很多時候會埋怨.
埋怨自己的配偶.
無論是丈夫或太太.
都會埋怨配偶.
你這麼懶.
你不懂教導.
你不做多一點.
不看兒女的功課.
你不陪他們多一點.
我們很多時候埋怨配偶不懂得做.
最近有個弟兄跟我說.
你知道為甚麼我們弟兄.
進去洗手間的時候.
一定要帶著手機進去嗎.

$^{521}$他說雖然環境上.
玩手機好像不太刺激.
但你知道當我們在外面玩手機的時候.
只玩一分鐘.
你的配偶就會說.
你去拿那些東西出來吧.
去搬這些東西吧.
去做那些吧.
去看兒子的功課吧.
但如果我帶著手機進去.
他不知道我做甚麼.
就算他明知我玩.
他也會假設我做一些更加需要的大事.
所以我們的焦慮.
我們會怪我們的配偶.
我們會埋怨他.
我們會要求.
我們另外一些焦慮產生的是甚麼.
我們會埋怨我們的弟弟.
我們的兒女.
不懂得自動自覺.
更重要的是我們會埋怨他們.
你不懂得.
你會將他們和別人比較.
你沒有誰那麼厲害.
沒有誰那麼懂得計劃安排自己.
一比較之下.
我們的矛眼就會出現很多憤怒.
很多指責.
我們也會埋怨.
只有我一個人做.
我們經常說.
全屋只有我一個人做.
我不做就沒有人做.
我不看頭看尾就沒有人看頭看尾.
或者我不看那些功課就沒有人做.
很多時候我們埋怨.
是覺得自己很辛苦.
就算我們很多時候不說出來.
那個埋怨其實在臉色上都出來了.

$^{561}$所以在父母親的時候.
我們要學的就是.
這些其實是在纏累我們的罪.
焦慮不是錯.
我們一定會擔心.
一定會關心.
一定會有緊張的焦慮.
但是焦慮引致那個關係破壞了.
引致了一些負面的行為.
引致了我們的憤怒出來.
一些過分的指責.
那些就變成一些纏累我們的罪.
因為我們靠我們的兒女.
做我們的見證人.
鼓勵我們做好父母的角色.
如果我們做好父母.
我們教的時候.
我們這樣教的時候.
我們出現了一些.
令關係破裂的行為.
我們的兒女學的是什麼呢?.
他們見證到的就是這些.
他們做我們的見證人的時候.
怎樣可以去見證我們.
去鼓勵我們繼續做好父母的角色.
所以在基民有兩個提醒.
放下重擔.
跳起容易纏累我們的罪.
接著另外.
基民也看到.
我們要全心忍耐.
奔擺在我們前面的路.
我們做父母的角色.
不斷地在跑步上跑.
向著主耶穌的終點跑下去.
全心忍耐.
忍耐在父母身上.
是我們經常挑戰的事.
以前有一個長輩跟我說過.
你知不知道照顧小孩最辛苦的是什麼呢?.

$^{601}$是他們很多時候.
無心之失說的話.
但很傷害我們.
令我們很沮喪.
煮好飯.
坐在一起吃的時候.
他們不小心說了一句.
為什麼這麼鹹?.
我們已經覺得.
這麼不體諒.
這麼不欣賞.
這麼在指責.
所以很多時候我們會學.
有些事情我們是.
控制自己的情緒.
我們要忍耐他們無心之失的說話.
或者很多時候.
子女之間的指控.
明顯地關心他們.
父母為什麼這麼操控.
為什麼這麼強權.
這麼獨裁.
這些字眼令父母很不舒服.
或者有些父母和子女之間的關係.
會面對一些更加爭議性的問題.
他們的行為更加令你覺得.
不應該這樣做.
在男女之間的關係.
或者性別認同之間的表態.
或者對於毒品的用.
或者是酗酒.
所以很多時候.
父母見到子女之間的行為.
我們會叫他們改.
會提醒他們.
但是沒有後果.
叫他們怎麼也不改.
或者他們一意孤行.
聖經提醒我們的時候.
其實我們是需要忍耐.

$^{641}$我們跟他們.
因為他們不接受.
不懂得去愛的時候.
我們就會把關係破裂.
忍耐不是一件容易做的事.
在羅.
在當時你可以想像.
神叫他做那艘船.
但是要這麼多空間.
這麼多準備工作.
帶一些家人.
帶一些動物上船的時候.
那段時間因為有很多人看著.
他們可能有很多回應.
很多白眼.
但是他要忍耐.
因為他知道他要向著一個目標去做.
在船上下了40天的大雨.
然後接著150多天.
水退了.
然後要再等40天.
才叫烏鴉和白鴿出去.
那段時間是很考驗一個人的忍耐.
所以當天的先賢告訴我們.
其中一個跟神的關係要學的.
就是一些逆境我們要學忍耐.
一些結果不是我們希望做到的.
做父母我們要學怎樣去忍耐.
所以這是第三.
這是基文對我們的提醒.
因為提醒是告訴我們.
產業是神給的.
神一定會負責任.
你擔心的事神會擔心.
你憂慮的事神會更加憂慮.
所以我們只可以去學.
去忍受.
去忍受一個不同的後果.
最後一點就是.
基文也提到父母.

$^{681}$當我們做父母的角色的時候.
我們要仰望誰呢?.
仰望創始成中的耶穌.
這個比起那個狠親.
也有些異曲同工的感覺.
為什麼我們要仰望創始成中的耶穌呢?.
因為我們做父母.
我們看周圍人的眼光.
看他們的子女是否出色.
看著他們就覺得我們成功做了一個好的父母.
其實這樣只會令壓力更加加深.
我們的期望越高.
我們的要求越高.
他們承受的壓力越多.
我們的壓力會將我們的關係拉扯得更遠.
所以基文告訴我們.
其實我們如果知道.
做父母的角色是神給的話.
是神的賜予.
是神的禮物.
是神的產業.
我們看著就是創始成中的耶穌.
講得很清楚.
基文說的是.
耶穌是輕看.
他忍受.
因為他在十字架的道.
受苦的時候.
是看著神奔跑.
同樣地告訴我們.
如果我們知道做父母的角色是神的賜予.
我們將目光放在神的時候.
我們不是祈求我們的子女怎樣出色.
而令我們驕傲.
我們不是祈求他們讀書好一點.
你這麼厲害.
你這麼乖.
你這麼孝順.
這麼聽話.
這麼懂得跟人相處.

$^{721}$我們不是祈求這些.
令我們覺得我們做好一個好父母的角色.
而我們祈求的是.
我們的子女.
他們是否可以做到我們.
在聖經裡說自己的學習.
他們是否可以放下重擔.
信得過神.
他們是否可以知道.
這是什麼潛累他們的罪.
而他們是不會被那些罪.
潛累他們有些負面的行為出來.
他們是否全心忍耐.
他們是否可以仰望主耶穌.
我們子女在我們身上學到這些.
成為我們的見證人.
是令我們做父母的角色更加舒服.
更加出色.
所以在結論來說.
對於大家的提示就是.
既然很清楚.
詩篇告訴我們.
產業是神給的.
是神的禮物.
我們作父母的時候.
我們所要問的就是.
我和神有沒有這個關係.
我和神.
我是否真的信靠神.
倚靠神.
我是否真的以神為我們的中心.
我們是否願意和神的生活.
以主為大.
這是我們作為父母親.
要和神之間的對答.
去學習和努力.
也告訴我們.
我們周圍的見證人.
我們的弟兄姊妹.
我們的配偶.

$^{761}$我們的家人.
你和神的關係是怎樣.
你的事奉是否盡忠.
是否完全盡心去事奉.
你可否成為一班父母的見證人.
因為你的事奉.
因為你和神的關係.
你對神的信.
而成為他們作為父母親的角色的同行.
所以希望是那些.
基文提醒.
放在我們心裡.
在我們的生活裡可以行得出來.
謝謝.
我會盡力的來為你報仇.
\newpage



\section{}
\label{sec:z_G8mnyC7Vg}
\textbf{11/17/19 主日講道: 與人同在的神 5 - 律法(陳崇基牧師)}
\newline
\newline
連結: \href{https://youtube.com/watch?v=z-G8mnyC7Vg}{\texttt{https://youtube.com/watch?v=z-G8mnyC7Vg}} ~~~~ 語音日期: 2019-11-24
\newline
\newline
\hyperref[sec:xwcSwwS07wQ]{\small{< < < PREV SERMON < < <}}
~
\hyperref[sec:index_chronic]{\small{[返順時目]}}
~
\hyperref[sec:index_scriptual]{\small{[返順卷目]}}
~
\hyperref[sec:pqDLnSQO_8g]{\small{> > > NEXT SERMON > > >}}
\newline
\newline
$^{1}$今天我們繼續這個系列.
叫做與人同在的神.
我們由開始第一次.
到現在有幾次.
我們說到聖經由第一章開始.
上帝是創造萬有和創造人.
也願意與人同在.
與我們同在.
不過人犯罪.
或者人是叛變.
來背叛上帝.
所以沒有了上帝的同在.
但是神仍然去救贖人.
仍然在伊連以外.
與人同在同行.
所以我們看到創世記第十二章.
那裡說到亞伯拉罕.
上帝特別呼召他.
由一個別國的地方.
去到迦南地.
當時世界的中心.
上帝藉著他.
以至他能夠成為一個家.
成為一個國.
就是所謂舊約的以色列國.
讓這個國家,國民,天國的子民.
能夠去作上帝自己的見證.
以致萬國.
能夠與上帝同在.
最後我們知道.
耶穌基督來到的時候.
所謂舊約和新約中間.
耶穌來到的時候.
他成就了這個救恩.
在十字架的救恩上面.
讓我們能夠藉著基督.
與神同在.
所以今天我們看到.
過了亞伯拉罕.
上帝呼召亞伯拉罕.

$^{41}$要使萬國去得福.
重歸上帝自己的關係裡面.
我們知道這個故事就是.
亞伯拉罕之後的後代.
雅各家.
十二個兒子.
因為饑荒的緣故.
就落了埃及.
在埃及那裡.
他們是.
eventually(後來)他們也是.
做了奴隸.
被法老一個強權的帝王.
來打壓他們.
然後上帝去呼召他們出來.
藉著摩西來帶領.
他們自己的子民出埃及.
然後出了埃及.
他們就去到.
經過了曠野.
去到一個地方叫做西乃山.
這個山就是舊約裡面.
最中心的一個山.
因為在這個山那裡.
上帝就出現.
來跟他們自己的子民立約.
並且跟他們去立法.
所以就設下.
他們自己的誡命和他們的律法.
來創造一個新的人類.
新的社會.
以色列國.
以色列民.
天國的子民.
他們就是上帝自己在地上的代表.
隱身自就是我們舊約新約都是一樣.
我們都是一群神國的子民.
所以我們看這些律法的時候.
我們就想去問一個問題.
藉著這些律法誡命.

$^{81}$上帝要彰顯他自己什麼樣的本性.
他對人間的社會.
或者對以色列國.
他自己的子民.
他的社會.
他的群體.
是有什麼樣的願景.
有什麼樣的vision.
特別是有什麼樣的social vision.
在這些律法裡面.
我們或多或少會看到.
上帝自己的心意.
上帝對人類社會的心意.
在我們今天這麼動盪.
這麼不安的時代.
包括香港.
包括世界各地不同的地方.
我們自然心裡面很嚮往.
所謂有law and order.
我們希望有律法.
我們希望有秩序.
希望社會能夠井井有條.
也更加我們很嚮往.
希望這個社會裡面的政權的人.
掌執政權的人.
執政者.
他們是有公義.
所以我們看到上帝的律法的時候.
其實是一些很色切性的東西.
就是上帝的律法.
他頒佈的時候.
他律法的條文背後.
要去表達一個什麼樣的意象.
一個什麼樣的vision.
所以我們就用一個這樣的角度來看.
這些舊約的律法.
上帝要創造一個新的群體.
他的律法到底要怎樣去維繫這個社會.
或者這些律法.
給我們今天身邊我們看到的局面.

$^{121}$看到的社會有什麼啟迪.
當我們看這些律法的時候.
我們可以分開兩方面來看.
因為這些律法的格式.
或者這些文法的形式.
都是有所不同的.
如果你有問題的話.
你可以掃描這個QR Code.
然後我們最後的時候.
我們有時間可以去解答一些.
你對信仰基督教有關的問題.
我們講到律法的時候.
我們講到律法的體材可以有兩樣.
你記得看律法書的時候.
當然我們知道.
這些平時在靈修的時候.
你不會覺得很溫馨的東西.
但是你要明白.
那些條文是有兩大類的.
一種就是所謂判例個案的律法.
判例個案的律法是什麼呢.
Causalistic Law.
它有一個格式.
你一看律法書你就會留意到.
它一定會有如果發生什麼事.
你就要做什麼事.
如果有這樣的情況發生.
你就要有什麼回應.
這叫做Case Law.
所謂判例.
就是有一個例子給你.
以致你引申出什麼叫做公義的判決.
另外一種就是Apodictic Law.
Apodictic就是明文的法律.
就是說你不可作這個.
或者你需要作哪一個.
Thou shalt not or Thou shalt.
我知道大家很聰明.
我給個考試給大家.
看看這個律法.

$^{161}$在初夏禁忌二十一章.
有一個這樣的律法條文.
他說牛如果捉死男或女.
總要用石頭來打死這隻牛.
總不可以吃他的肉.
牛的主人可恕人無罪.
糖藥.
那隻牛是數來都喜歡撞人的.
有人也報告給牛主.
他竟然不把這隻牛綁住.
以致把男或女撞死了.
就要用石頭打死這隻牛.
牛主都必須要自死.
這是判例的案.
還是明文的法律.
判例沒錯.
因為是約過發生什麼事.
然後你應該有什麼回應.
如果你的車救了.
你不知道的.
你的車失靈.
然後撞死了一個人.
那車就要怎樣.
毀了他.
燒了他.
你不可以吃他.
OK.
那主人車主可恕人無罪.
因為不知者不罪.
但是如果你的車一向都壞的.
一向都會失靈的.
你又知道的.
有人告訴你了.
你就不理會.
繼續開.
然後撞死了男或女.
那又怎樣.
怎樣處置車.
毀了他.
那你呢.

$^{201}$你就有罪了.
這個就是所謂判例的律法.
我們看一些這樣的律法經文.
我們留意一下.
我們會抽一些不同的經文.
到底這些背後.
上帝的心意是什麼.
人命關天.
人是很重要的.
因為人有上帝的形象.
有意外發生當然不好.
但是上帝的律法.
去表明上帝對生命神聖的心意.
每一條命都是重要的.
所以要做好所有的事情.
以致人命不會受影響.
不會有危險.
保存著神聖的生命.
這個就是律法背後.
上帝的本性.
我們看另外一段.
就是很有名的.
就是舊約裡面所謂十誡.
《春秋及經》第20章.
我一口氣讀完這十誡.
神吩咐這一切話說.
我是耶和華的神.
曾將你從埃及地為奴之地領出來.
除了我以外.
你不可有別的神.
這是第一誡.
第二誡.
不可為自己雕刻偶像.
也不可做什麼形象.
彷彿上天下地.
地底下水中的白麥.
不可以跪拜那些像人.
不可以事奉他.
因為我是耶和華的神.
是忌邪的神.

$^{241}$恨我的人.
我必追討他的罪.
直腐及止直到三四代.
愛我的守我誡明的人.
我必向他發慈愛直到千代.
這是第二個誡明.
不可為自己.
為上帝來雕刻一些偶像.
不可妄稱神的名.
神的名.
因為妄稱耶和華名.
耶和華必不以他為無罪.
這是第三誡.
不可妄稱上帝的名.
第四誡.
當紀念安息日.
守為聖日.
六天你要勞碌你一切的功.
但第七天.
是向耶和華你神當守的安息日.
這一天你和你的兒女.
瀑被生畜.
並你城裡面寄居的客人.
滿河宮都不可作.
因為六天之內.
耶和華做天地海其中的萬物.
第七天便安息.
所以耶和華賜福與安息日.
定為聖日.
守安息日.
敬拜.
七天裡面有一天是敬拜.
第五誡.
當孝敬父母洗禮的日子.
在耶和華你神所賜禮的地上.
得以長久.
第六誡.
不可殺人.
第七誡不可姦淫.
第八誡不可偷渡.

$^{281}$第九誡不可作假見證陷害人.
第十誡.
不可貪戀人的房屋.
也不可貪戀人的妻子.
瀑被牛露和他一切所有的.
這個舊約裡面.
可算是歷代裡面.
最出名的一段經文.
或者最有影響力的經文.
十誡.
頭四誡.
是說到.
我們和上帝之間的關係.
接著那六誡.
就是說到人倫的問題.
人與人之間的關係.
所以不可殺人.
不可偷竊.
不可作假見證.
陷害人等等.
我今天想問你.
我們想問.
十誡到底頒佈給誰.
十誡.
直接的受眾.
intended audience.
到底是誰.
當然我們就要說神的子民.
所有神的子民.
但是你看看.
這些條文和字眼的背後.
到底直接的受眾是誰.
direct audience.
到底是誰.
你將這些資料.
整理之後.
我大概整理到.
有這十個元素.
這個人.
或者這一群人.

$^{321}$就是十誡律法的直接受眾.
第一.
那個一定是男人.
第二.
那個是做父親的.
因為他有兒女.
第三.
他有物業.
第四.
他是已婚的.
第五.
他是年紀大到一個地步.
他的子女是在跟他打工的.
因為你看看安息的那個律法.
你跟你的兒女.
你的瀑被.
你的生畜.
都不可作供.
在安息人裡面.
所以他的子女跟他打工.
但是他們年紀也是年輕到一個地步.
還有.
在生的父母要照顧的.
當孝敬父母.
在你神所賜給你的地裡面.
你的日子可以長久.
他們住在城市裡面的人.
urbanize(擁有地區).
他們都是生活挺富足的.
怎麼看得到.
有牛.
有奴.
有工人.
他們社會的權位.
是比較重要的.
並且重要到一個地步.
是可以在城門口.
或者在法庭裡面.
可以做證人.
不可作假見證.

$^{361}$陷害其他人.
一定要在社會裡面有一點地位.
他才被 call to the court.
to testify.
right?.
所以他們在法庭上.
可以能夠.
訴訟上.
是能夠作證供的.
社會裡面的領袖.
這個是.
十誡的直接的受訟.
我知道十誡當然是給每一個以色列人.
平民百姓任何的人.
但是很明顯.
那個字眼背後.
是寫給.
當時社會的權貴.
有權的人.
或者.
leaders.
做領袖的人.
或者用今天的話來說.
這些的界名.
這些的律法.
直接的受訟.
是寫給那些管治者.
或者.
有權威的人士.
社會裡面中流砥柱.
或者所謂我們上了岸的.
到了收成期的人.
有 privilege(權利)的.
律法.
第一個.
要去到說話的.
或者頒布十大國的大的原則.
好像國家的憲法一樣.
寫給的.
頒布給的.

$^{401}$就是社會裡面權貴的.
第一個.
要跟他們說話的.
不是 the governed(政府).
而是 the government(政府).
你明白我的意思嗎?.
直接的受訟.
不是 the governed(政府).
不是被管治者.
不是下面那些的草蜢.
而是管治者.
the governing agents(政府).
政府裡面的人.
或者.
執法者.
法庭裡面的人.
城門口.
可以坐在城門口.
去到談訴訟.
法庭裡面的人.
似乎告訴我們.
上帝去頒布律法的時候.
他第一個想去到對他們說話的.
就是那一群.
在社會裡面.
高位的.
上了岸的.
所謂修成企業的人.
是跟這一群人先去到說話.
並且嘗試去約束.
去限制.
他們某一些的權力.
這個是我們今天現代法治的精神的一部分.
Law and order(法治和秩序).
Rule by law(法律的規則).
執政者.
執政者有權在手.
他是有界線的.
因為律法說.
有界線.

$^{441}$你不可以繼續去濫用你的界線.
以致你拿了別人的東西.
Right(對).
你是有一個範圍.
第二方面我們可以想到的就是.
這個十屆頒布的時候.
你想想對什麼人有好處.
或者對什麼人是沒有好處.
剛才我們問的問題就是.
Direct audience(直接的觀眾).
直接的受眾.
直接跟他說話的.
就是有位在權的人.
在位有權的人.
到底誰是有好處呢.
在保障誰呢.
這個十屆的大原則.
到底在保障誰呢.
我們說頭四屆就是神與人的關係.
所以頭四屆可以這麼說.
是保障上帝的明星.
除了我以外你不可有別的神.
不可以丟偶像.
不可以濫用上帝的權威.
不可以在上帝的名下亂做事.
你要守上帝的聖人.
安息人.
這個是保障他人.
保障上帝自己的地位和他的名聲.
這是律法永遠有的方向.
是保障其他人.
然後六屆人倫方面.
就是保障什麼人呢.
保障做勞部的人.
因為七天裡面他工作.
他不可以讓主人濫用他的權力.
七天起碼他要放假放一天.
那些工人.
或者寄居者.
在當時以色列社會裡面的.

$^{481}$那些sojourners.
travelers.
或者foreigners.
那些是沒有政治權柄的人.
沒有政治權利的人.
沒有political right的人.
你要保障他們.
因為那些寄居的.
他都不可以被凌辱的.
保障鄰舍.
不可以貪戀別人的東西.
別人的妻子.
別人的房屋土地等等.
不可以作假見證.
保障誰呢.
不是保障我.
是保障別人的名聲.
我不可以在這裡說三道四.
講是講非.
毀壞別人的名聲.
更加不可以在法庭上.
作假見證.
這裡不是純粹講大話這麼簡單.
是作一個假見證.
來損毀別人的 reputation(名聲).
他的名聲.
總括來說.
如果這麼說的話.
就是耶穌所說的.
解明的精神是什麼呢.
最大就是用盡心盡力.
愛主你的神.
透四界.
然後就是愛人如己.
愛你的鄰舍如己.
這個新的社會.
這個律法下.
這個社會才會有障礙.
如果每一個在位在權的.
只是保障自己的權柄的話.

$^{521}$這個社會永遠不會有障礙.
因為只會是一個肉強食的社會.
你看到這個律法背後.
就是彰顯了上帝自己的心意.
我們再看另外一些經文.
你美記第十九章.
我想問你.
剛才那個問題.
這段的律法.
那個的直接受眾是誰.
第二.
到底誰有好處.
或者要保障哪些人.
你美記十九章第九節.
在你們的田地收割莊稼的時候.
不可以割盡田國.
也不可以拾取所遺落的.
不可以摘盡葡萄園的果子.
也不可以拾取葡萄園所吊的果子.
要留給窮人和寄居的.
我是耶和華.
你們的上帝.
誰是直接受眾.
是什麼人.
聽到這個律法的時候.
他會覺得有意思.
當然一定是無地無產業的人.
不會覺得有意思.
因為他不會有地.
你們的地.
一定是所謂地主.
有產業有物業的人.
他們才會有田地.
收割的時候他們會去工作.
這裡有四個字.
這八個字是專有名詞.
第一.
田國.
什麼叫田國呢.
就是你那塊地.

$^{561}$你的farmland.
周邊的地方.
你可以說田國到底有多寬有多窄.
這些就是猶太拉比很喜歡辯論的東西.
有一個辯論就是說.
你的田地的六十分之一的界.
就是你的田國.
所以你不可以在那裡收割.
這個就是田國.
所為樂.
這裡說收割莊稼是說大麥小麥.
我們耕田的人就知道.
如果你以前耕過田的話.
收割的時候你拿著一手的麥水.
然後你有鐮刀在手.
然後你就斬.
你就拔出來.
斬斷了的你就拔出來了.
所為樂的.
就是肯定你拿著的時候.
你再斬的時候.
就一定有一些斬不完的.
是不是.
所以你一掠上來的時候.
你就有一些剩下的.
那些就是所為樂的.
Grain product.
Barley或者wheat.
大麥或者小麥.
不可以拾取這些為樂.
就是你斬了之後就只有這麼多了.
不可以有些剩下的.
我又再斬.
律法就是說你不要這樣.
第三個.
專有的詞.
就是葡萄園的果子.
中文翻譯應該譯作.
是未熟透的葡萄.
不可採摘未熟透的葡萄.

$^{601}$這裡說的是另外一種的植物.
另外一種的農產品.
就是你去到葡萄園那裡.
收割這些葡萄和葡萄的葡萄的時候.
有一些未熟透的.
熟透的你會摘下來.
未熟透的可能是莊稼.
Harvest時期來到的時候.
還未熟透.
還要可能多等一兩個星期.
它才會熟透.
那些你不要摘.
未熟透的果子.
然後所掉的果子.
就是那些可能熟透了.
或者你摘的時候.
剩下掉下來在地上的.
那些你不要摘.
有多少你都不要摘.
四個專有的名詞.
其實是說出什麼呢.
就是說這些東西是屬於誰的.
屬於寄居的人.
那些在你下面.
無權無勢的人.
窮人.
在經濟上比你差的.
其實主要的目的是什麼呢.
就是彰顯上帝的公義.
上帝的本性.
就是一個人愛的上帝.
特別愛什麼人呢.
就是有他形象.
並且在社會裡面.
他可能是在捱苦.
他不是那麼多權力的那一群人.
或者這樣說.
反過來說.
那個律法的精神.
就是要向那些.

$^{641}$所謂有物業的人.
上了岸的人.
修成旗的人.
就告訴他請你.
要學上帝有那種人愛.
和公義的本性.
他是愛那群下面的人.
你這樣做的時候.
你就是在彰顯上帝的本性.
in a way 就是說.
慷慨一點.
特別對下面無權的人.
或者比你少權的人.
你儘管generous(慷慨)一點.
這個很多應用的地方.
我們經常咬文嚼字.
好像那些猶太人的拉比.
田國到底有多寬啊.
我們一定要是非黑白講清楚.
60\% 60分之一.
50分之一還是70分之一.
所謂什麼叫田國呢.
其實它的總義就是.
practice generosity(謙虛).
特別對那些無權無勢的人.
你講說話.
你所行所言所作的事情.
請你去憐憫他們.
我們很多時候奉獻的時候.
我們都會想想金錢奉獻.
到底田國有多少.
到底是扣了稅才十分之一.
還是還沒扣稅之前十分之一.
不好意思.
法書沒怎麼講這些.
That's not the point(這不是重點).
重點就是上帝的本性是一個慷慨的神.
我想問你.
如果是慷慨的神.
你去反映他的本性的時候.

$^{681}$你會怎麼奉獻呢.
講到法治.
講到律法.
如何去維繫一個社會.
或者是如何去形成一個社會.
塑造一個社會.
這個社會要有上帝的形象在裡面.
你這樣看的時候就很有意思.
並且在我們今天.
我們看到的政局.
看到的社會.
有很多東西我們值得去想.
戴耀廷教授.
他曾經寫過一本書.
講法治.
我覺得他講得很有意思.
他用十六個字.
去表達到底什麼叫法治.
他說法治有四個層面.
他說第一個層面就是有法可依.
意思是有明文規定.
寫出來.
無論你是講基本法也好.
美國的憲法也好.
任何的法律.
都是有明文去寫出來的.
不是說掌權者說怎麼就怎麼.
所以第一個層次.
是有法可依.
第二個層次.
就是有法必依.
犯了法的就要受刑法.
這個很簡單.
大家都知道.
爸爸是男人.
不用說了.
犯了法.
你就要去到.
受刑法.
違例者.

$^{721}$就需要受刑.
你燒爛了東西.
是不對的.
你應該要抓他.
但是這個只是.
所謂法治的第二個層面.
因為你想想.
如果執政者.
經常開口閉口.
都只是叫人民你們要守法.
我說什麼你們就要照做.
這樣就很容易.
令到法治的精神.
會受破損.
因為如果執政者只是會說.
我們要你守法.
我們要你做這個做那個.
你聽聽話話.
那麼法律很多時候會成為.
欺壓人的工具.
所以第三個層面.
就是以法限權.
以法限權是什麼意思呢.
就是我們剛才看十屆.
或者不同的.
舊約的律法.
就是去到制衡.
keep them accountable.
執法者不可以濫權.
你想想.
如果一邊人犯法.
他沒有衝擊到法治.
犯法就要抓.
有法必依.
但如果你說執法者自己犯法.
執法者自己越權濫權.
那麼這個對法治的破壞就是最大的.
對嗎.
所以以法限權是第三個層次.
因為如果沒有限權.

$^{761}$沒有check and balance(檢查和平衡)的時候.
這個對社會造成的危機.
就會更加可怕.
如果你想想聖經裡面的例子.
不難找到.
君王大位.
他要奪人的妻.
殺人的丈夫.
他沒有犯法.
基本上他沒有犯法.
他只是貪人.
貪是不能抓的.
所有事情都是在他的權力範圍裡面.
他能夠作的.
他可以寫一個軍令.
然後害死人的丈夫.
這些例子.
多的是.
聶王記上21章.
可能不是很多人知道.
我們會聽到這樣的故事.
不過我告訴你可能會有些印象.
阿哈王.
King Ahab.
就是以利亞先知去抗衡阿哈王.
阿哈王有一個什麼故事呢.
聖經裡面有29節.
整章都記載這個故事.
就是阿哈王很喜歡.
有一個平民叫做拿伯.
他的葡萄園地.
因為葡萄園地在皇宮旁邊.
他就很想要那塊田地.
他就跟拿伯說.
拿伯我想要你的地.
不如我賠另一塊地給你.
或者我給你錢買.
拿伯回答他.
不行,這是我們的祖地.
祖宗的產業.

$^{801}$我不能給你的.
所以我不會給你的.
阿哈王就很不舒服.
很傷心.
吃不下飯.
最毒的婦人心.
皇后夜洗別就來問他.
為什麼你不吃飯.
阿哈王跟他說拿不到那塊地.
夜洗別就說.
你快點吃飯.
我明天跟你解決.
夜洗別做什麼呢.
皇后就用阿哈王的玉璽.
寫了封公函.
來告拿伯.
說他褻瀆上帝的名.
然後那封公函.
就叫城市的長老.
明天早上去城門口開會.
開法庭會的時候.
拿著這封公函.
然後叫長老來.
拿伯也在.
就讀這封公函.
然後就去告.
用莫須有的罪名告拿伯.
然後就請幾個土匪.
就是當時黑社會.
去打死他.
結果真的就是這樣發生.
夜洗別就用一個這樣的方法.
就把拿伯害死了.
有沒有犯法呢.
沒有犯法.
因為是公函.
有玉璽嘛.
他就拿了那塊地.
這件事令到耶和華非常不喜悅.
所以以利亞就去批評他.

$^{841}$再多一次批評他.
這些這樣的例子.
多的是.
就是以法限權.
就是這個意思.
但還有一個更加深的層次.
就是以法達義.
就是要建立一個人愛公義的群體.
一個社會.
這個就是上帝自己的心意.
有沒有例子.
多的是.
你看福音書.
大把.
耶穌和法理賽人之間的碰撞.
很多時候法理賽人都在說什麼呢.
就是有法必依.
因為安息日就不可以作供.
安息日就不要醫人.
醫人就要在第一天才醫.
你等一天.
耶穌偏偏不肯.
耶穌偏偏在安息日就要治病.
法理賽人開口閉口就說.
你就是破壞我們的律例.
但耶穌說什麼呢.
耶穌就說.
你根本不明白.
律法最終的目的是什麼.
最終的目的是達義.
是讓社會群體行公義.
好憐憫.
這些思考.
是很值得我們今時今日去想的.
因為我們今天太多.
這些社會的動盪.
法制上的問題.
一個社會的問題.
這些怎麼能給我們啟迪.
值得我們去思考.

$^{881}$幾年前我認識一個學者.
他在美國讀PhD.
我去載他去幫他的時候.
我就問他PhD讀什麼.
他說博士研究憲法.
constitutional law.
這麼奇怪這麼有趣.
大家談談當中.
他講到為什麼.
當然每個PhD都會justify.
為什麼他讀的東西有意思.
他說我覺得21世紀.
是一個威權統治的時代.
越來越多霸權或者威權出現.
所以為什麼需要有憲法.
建立一個社會.
維繫一個社會的時候.
是需要有法治有憲法的精神.
所以他就讀這個.
我記起有一本書.
時常被人引用.
這本書是一個耶魯大學的歷史學者.
他專門研究20世紀裡面東歐.
中歐或者蘇共.
很多時候是共產的國家.
或者威權的國家社會裡面.
一些歷史.
Snyder他寫一本書On Tyranny.
他說看20世紀裡面.
這些威權的社會當中.
有什麼lessons.
或者功課.
我們需要學的.
譬如.
他說在威權的社會裡面.
很多時候.
政權想人民有分化.
想人民彼此敵對.
彼此割席.
彼此對立.

$^{921}$所以他說.
有一個lesson我們要學的.
這裡有很多個.
其中一個他說.
Believe in truth.
我們一定要相信真相.
世間是有真理有公理的.
另外一個很有意思的第十二個.
他說Make eye contact.
Small talk.
他說因為在威權的統治下.
很多時候.
人民會被分化.
而那個氣氛就會是恐懼懷疑.
人與人之間會被拆散.
或者人與人之間會是有你無我.
有我就無你.
你偏我這邊權的話.
你就屬於我.
不偏我的權的.
他就是敵人.
他說Make eye contact.
Small talk.
就是說我們不要被威權擺弄我們.
人與人之間的關係.
是怎樣建立出來的.
眼目的接觸.
小小的對話.
所以大家彼此一個擁抱.
一個握手.
一聲早晨.
你想想你的鄰居會變成怎樣.
如果是這樣的話.
他說很多時候在二十世紀.
這些極權的社會裡面.
很多時候人與人之間.
就是不想看到大家.
他看到你然後就在那邊過馬路.
人與人之間就會分散.
有很多的分化.

$^{961}$怎樣去抵抗暴權.
Small talk.
Make eye contact.
教會就是在做這件事.
剛才說大家彼此擁抱.
我想整個星期都沒有人擁抱過你.
過去幾個星期.
可能人都沒有碰過你.
有沒有人跟你握過手.
很好笑.
有沒有人跟你Make eye contact.
很好笑.
教會在地上代表耶穌基督.
我們其實就是在做這件事.
我們很多時候會很理想化.
或者這些small talk很煩.
但你要知道人與人之間的關係.
community.
國家社會就是這樣造成出來的.
沒有small talk就不用說有什麼big talk.
教會新的群體在基督裡面建立的.
就是在做這件事.
我想起一本書叫做.
The Gospel Comes With A House Key.
特別是寫給一個在後現代.
後基督教的世人.
The Gospel Comes With A House Key.
即是我們信福音的時候.
其實我們家門的鎖匙.
是需要交給主的.
意思是我們需要去款待身邊的人.
開放我們的家庭.
讓人進入我們的私有空間.
去make community.
在後現代社會.
就是一個部落主義的世界.
有分化有對立.
如何可以締造上帝的心意.
有community 有family.
不只是你自己的家庭.

$^{1001}$讓人被邀請進入你的家.
Comes With A House Key.
我們信耶穌的時候.
把這個House Key交給主.
我們的私有空間可以跟別人分.
這就是抗衡現在的氣氛.
我想最後這段時間.
我們花一點時間去禱告.
心裡面看到一些新聞或故事.
無論世界各地不安的地方.
或香港.
我知道現在理工大學受圍困.
情況很緊張.
我們為那裡祈禱.
為你心裡想到的不安處境去禱告.
我們一起在這裡祈禱.
我先給大家一點時間去禱告.
主你是人愛.
你是公義的神.
我們知道萬有都是屬於你的.
你更加是歷史每一刻的主宰.
主你為了這個世代有很多動盪不安.
可以來到你面前.
我們知道你已經勝過黑暗的世界.
勝過掌權的撒旦.
你賜給我們平安.
不是世界能夠給到.
主要是讓平安能夠臨到不安的處境.
我們特別紀念在香港.
已經持續很多個月.
甚至很多年的社會問題.
那種張力.
主你求我們去到靈魂.
特別在這一刻.
在理工大學裡面的學生.
主你求我們去保守他們每一位.
因為他們其實就是這個社會的將來.
主你求我們使警方克制.
使他們的權力不會被濫用.
主你也為在上掌權躲在後面的執政者.

$^{1041}$主你求我們改變他們.
求我們親自去對抗他們.
制衡他們.
主你我們看不到人間有什麼很好的答案.
我們只求你幫助.
在不安的處境裡面.
我們私隱憐憫.
只求平安.
你作保護.
彰顯你的公義.
主要是讓你兒女在不同的地方.
包括我們成為教會.
我們能夠成為一個新的群體.
讓這個群體有在你自己的形象.
彰顯你自己的美善.
在這個動盪不安威權的氣氛裡面.
教導我們如何在基督裡面.
用恩典.
用你的愛.
去建立新的社會.
讓人來到教會的時候.
感受到神是自己的實在.
我們這樣祈禱.
奉主耶穌基督的名字.
\newpage



\section{}
\label{sec:pqDLnSQO_8g}
\textbf{11/24/19 主日講道:感恩,愛,赦免(吳君才牧師)}
\newline
\newline
連結: \href{https://youtube.com/watch?v=pqDLnSQO-8g}{\texttt{https://youtube.com/watch?v=pqDLnSQO-8g}} ~~~~ 語音日期: 2019-11-30
\newline
\newline
\hyperref[sec:z_G8mnyC7Vg]{\small{< < < PREV SERMON < < <}}
~
\hyperref[sec:index_chronic]{\small{[返順時目]}}
~
\hyperref[sec:index_scriptual]{\small{[返順卷目]}}
~
\hyperref[sec:SGo8UblJv4Q]{\small{> > > NEXT SERMON > > >}}
\newline
\newline
$^{1}$我11天前才從香港探親回來.
有一段時間在香港.
真的感受到他們的感受.
以及經歷他們的經歷.
那種不安 那種憤怒.
那種裡面的一些困惑.
甚至有一種很無奈的感受.
時局不斷在變化.
但我們深信神仍然是掌權.
我們都是看到他們很多的市民.
在這麼艱難的局面裡面.
仍然是怎樣去面對.
神有恩典 保守他們.
在巴西 我去探一個人 去關一個人.
這個弟兄跌倒 腦出血.
住了ICU一段時期.
出來的時候要再學校.
很多事情要重新再來.
幸好他的語言能力和表達能力都更好.
都是去加拿大 都感恩.
他康復得不錯.
我要說的就是我們今天可以坐在這裡.
有健康 有生命氣息.
可以自由崇拜.
真是恩典.
我們要好好珍惜.
今天我要和大家分享的題目就是.
感恩 愛 聖義.
我們去看一段經文.
就是在《路加福音》七章三十六到五十節.
我們看看經文.
我很快讀.
有一個法利賽人請耶穌和他吃飯.
耶穌就到法利賽人家裡去坐席.
那裡有一個女人 試過醉人.
知道耶穌在法利賽人家裡坐席.
就拿著盛香膏的玉瓶站在耶穌背後.
挨著她的骨 眼淚濕了耶穌的腳.
繼續下一章.
就用自己的頭髮擦乾.

$^{41}$又用嘴臉親她的腳.
把香膏抹上.
請耶穌帶回法利賽人家身邊者.
心裡說.
這人若是先知.
必知道我他的是誰.
是個醉人的女人.
乃是個醉人.
耶穌對她說.
西門 我有句話要對你說.
西門說 夫子請說.
耶穌說.
一個債主有兩個人欠他的債.
一個欠五十兩銀子.
一個欠五兩銀子.
因為他們無力償還.
債主就開恩免了他們兩個人的債.
這兩個人.
哪一個更愛他呢.
西門答說.
我想是那多得恩免的人.
耶穌說 你斷的不錯.
於是轉過來.
向那女人面對西門說.
你看見這女人嗎.
我還要離家.
你沒有吸過水洗腳.
但這女人.
用眼淚濕了我的腳.
用頭髮擦乾.
你沒有要我親嘴.
但這女人從我進來的時候.
就不停地用嘴親我的腳.
你沒有用油光我的頭.
但這女人.
用香膏抹我的腳.
所以我告訴你.
她許多的罪都赦免了.
因為她的愛多.
但她赦免少的.

$^{81}$她的愛就少.
於是對那女人說.
你的罪赦免了.
我們低頭向她說祈禱.
求主用你的話.
再一次提醒我們.
怎樣好好地.
趁著這個感恩節的假期.
我們思想我們和你的關係.
以及和人的關係.
主你的靈在我們中間.
說話.
祈禱紅主耶穌.
阿們.
這段故事很多人都聽過.
看到一個重要的人物.
很明顯.
就是一個人.
這個人就是西門.
他是做東.
請其中吃飯.
其實一會兒還有其他客人.
而可能這些都是.
法利賽人.
或者宗教人士.
第49節有說.
另一個就是一個妓女.
37節稱她為罪人.
這個詞的意思.
不單止指她是一個不敬畏神的人.
或者是一個不體面的人.
更加在當時的社會來說.
這個稱呼.
是指一個妓女.
而她的身份.
很多人都知悟.
整個故事發展開有幾個段落.
就是你看到.
西門是法利賽人.
這個人請耶穌吃飯.

$^{121}$一會兒有一個聲名狼藉的女人.
來到耶穌身邊.
在他後面.
表達對耶穌的感恩.
和愛.
36到38節有說.
西門對這個妓女的表現.
很反感.
心裡面暗暗的不滿.
第39節.
耶穌用兩個負債的人的比喻.
來教導西門.
將他和西門.
作一個對照比較.
然後你又見到.
耶穌向那個婦人說.
和肯定.
她的罪經寫言.
我們進入.
《釋經》部分.
看看西門請客的動機.
你會見到.
西門是一個法利賽人.
很奇怪.
如果你看《福音書》.
法利賽人和耶穌的關係.
是很緊張的.
為什麼這一次請耶穌吃飯.
這麼好人呢.
如果由經文的推測.
我們可以看到.
他的表現.
不是友好.
和崇敬主的.
如果不是的話.
為什麼他不會做一些.
由太人請客的規矩和禮貌.
就是用水餵洗腳.
甚至和客人親聚.
第44到45節有說.

$^{161}$這次西門請客.
他心裡面有一個目的.
他請了其他的同道.
第49節.
當時耶穌剛剛開始出道.
一個剛露出來的宗教領袖.
我們可以推想.
西門的動機.
可能要請耶穌來測試他.
看一看他是什麼家伙.
用我們現在的宗教術語.
可能請他做一個宗教對話.
或者一個信仰的座談會.
測一測.
耶穌.
你是真還是假.
是一個先知.
堅定立.
我們看第二天.
這個婦人的背景和表現.
37到38節.
英文一直沒有記錄.
這個婦人的名字.
是不是?你翻譯了嗎?.
37節.
在介紹當中說.
她是一個罪人.
一個罪人.
然後英文一直用.
這女人.
耶穌稱呼她四次.
那女人.
作者稱呼她三次.
你形容她.
這個女人沒有理由沒有名字.
應該有的.
但是刻意地不提.
一個聲名狼藉的人.
為什麼會這麼無媒.
在宗教領袖的宴會裡面出現.

$^{201}$在這樣的場合.
去沽沒耶穌.
你看她懷著.
一個心.
是一個感恩的心.
和熱愛之情來沽沒耶穌.
所以她向著.
耶穌獻上最高的崇敬.
用一個行動來看.
我們可以推測.
她先前.
一定見過耶穌.
接觸過耶穌.
可能聽過耶穌一些教訓.
赦罪的道理.
感受很深刻.
對她來說是一件很寶貴的事.
這個赦罪之恩.
叫她能夠經歷到身心.
她體會得很深刻.
所以她才會.
這麼大膽無媒.
來到耶穌面前.
用香膏沽沒耶穌.
獻上自己的感恩.
為什麼她用這樣的場合.
出現呢.
在這樣的場合.
是一個好的場合呢.
我們來看看.
猶太人.
請客會請一些.
人來教誨.
或者拉比來講座.
通常有些是.
大戶人皆富有的人.
往往他們將會在自己家的廷院開放.
歡迎其他外來的人.
自由進入.
以致他們都可以.

$^{241}$分享,聆聽到拉比的教誨.
而當時的飲宴.
就不像那些中國人.
圍著桌子.
西方人是坐在長桌子.
不是.
是向著飲宴桌子.
但是左傾.
側側地.
然後腳向後.
這個可以解釋.
為什麼這個婦人.
能夠在後面沽沒耶穌的腳.
這個婦人當時情不自禁.
哭起來.
眼淚一滴滴滴在耶穌的腳上.
38字有說.
她來到耶穌面前.
其實是帶有很強烈的感激之情.
感恩的心.
深深體會到自己昔日的污穢敗壞的罪.
耶穌今天已經赦免我了.
這位神的羔羊赦罪的主.
完全潔淨了我.
寬恕了我.
在主面前.
她經歷了什麼叫做無條件的赦免.
和憐憫.
被主完全的接納.
她可以重新做人.
重新抬起頭來.
面對未來.
她的過犯.
原來與被主耶穌同沒主義.
也不再紀念她以前的那種罪惡.
這個這麼深刻的經歷.
我相信在座的弟兄姊妹.
你信耶穌的時候.
會不會有一樣的感受呢?.
所以她看到.

$^{281}$知道耶穌的時候.
她完全是情不自禁的.
所以她的哭泣是真的.
而這個哭泣也是始料不及的.
克制不來的.
因此當她的眼淚滴一滴滴滴滴.
滴在耶穌的腳的時候.
她很狼狽 她措手不及.
在這樣的場合.
她更加做出一個.
當時令人很譁然的行動.
她渾然忘卻了自己當時社會的禮教和規範.
把自己一頭的秀髮打散.
然後用來擦乾耶穌的腳.
帶來的香膏.
然後擦抹上.
然後連連親吻耶穌的腳.
在當時的封建的社會.
或者一個女人披著頭髮.
是不雅的.
是不合體統的事.
當時有些人甚至可以作為優妻的其中一個理由.
不過這個女人對主耶穌的愛.
使她一再敢於打破當時的禮教和限制.
渾然忘我到一個地步.
有人說你親耶穌的腳.
其實當時的社會裡面.
親人的腳是表達一個最高的敬意.
一個感恩.
很崇高的表現.
我們看看第三.
當時耶穌的回應是怎樣呢?.
40到48節和50節.
耶穌默默地接受這個婦人對祂的一切敬愛和尊崇的行動.
祂看到這個婦女對祂的表現.
但是西門也看在眼內.
她心中有所批評.
西門的想法充分反映了當時的人.
或者法利賽人和一些人的假設.
如果你是一個先知.

$^{321}$你沒有理由不知道這個人有罪.
或者不潔來到你面前.
你摸他是一個怎樣的人.
她認為可能這個女人用來抹耶穌香膏的.
那些香膏是用她賣淫得利的錢買的.
這種錢是不為神所喜悅的.
新領記23章18節說.
如果耶穌是先知.
祂應該不讓這個女人用香膏來抹他.
在西門眼中這個女人只是一個罪人.
所以西門採取了一種輕描不捨的態度.
但是我們很感恩.
我們這位洞察人心的神.
我們的耶穌特意地說出兩個負債人的比喻.
追問耶穌.
兩個負債人當中哪一個更愛那個債主呢?.
西門給了一個很明顯和合理的答案.
耶穌肯定西門的判斷.
透過這個比喻.
直接指出一個很重要的原則.
這個原則是甚麼?.
多得恩典得赦免就多愛.
愛和恩典赦免原來是成正比的.
然後耶穌用這個原則來解釋.
這個女人對他的表現.
把他的表現和西門的表現作一個比較對照.
影襯出這個女人對耶穌非常謙卑.
極深的敬重和極熱切的愛.
接著47次耶穌提出一個很重要的真理.
這個真理就是愛是由於赦免.
赦免多就愛多.
赦免少愛就少.
相反的越少赦免反映你愛越少.
如果你沒有赦免.
即是說你沒有愛.
這個很重要的真理.
不知道你今天掌握到嗎?.
這段經文對我們今時今日有甚麼應用意義或提醒.
我會花多些時間來分享.
第一.

$^{361}$熱切的感恩敬拜.
抽離旁觀.
西門和那個女人接近耶穌的心態和動機完全不同.
兩個的表現態度完全不一樣.
西門是當時公認的宗教領袖.
是一個上流社會身份的人.
一個有地位的人.
他認識耶穌.
很大程度上是由於好奇.
他認為耶穌可能是眾先知當中其中的一位.
而且他有些同伴發力才能.
一起為耶穌獻命正身.
檢定他的資格.
但是他在招呼耶穌的時候.
用基本的禮貌態度.
款項之多的禮節.
他都怠慢了.
他跟同道的邀請耶穌.
其實態度只是輕言淡語.
只是淡淡說說.
為了有關於耶穌的事.
知道多一點的事.
完全看不到耶穌原來是那位創天造地.
萬失都向他跪拜感恩的時候.
是值得我們還讚尊崇的主.
所以他雖然跟耶穌在一起.
但是他的心態可以說是抽離冷淡.
他是不投入,不賣身的.
淡淡然的冷漠.
相反的那位曾經滄海被人鄙視的女人.
她冒著很大的阻力.
和別人的冷嘲滅諷.
不恥毅然的來接近耶穌.
在他心靈裡面.
他對耶穌的態度只有寸言的敬拜.
很投入的感激和愛護.
弟子們.
今天我們回教會.
參加聖拜保住什麼心態?.
每個星期日我們回來是怎樣呢?.

$^{401}$有些人說我一到六天在外面.
世俗凡塵打滾.
我星期日來教一些熟悉的東西.
有些人回教會不是為了見神.
而是要見人.
多過親近神.
有些人說很悶.
回來教會打發時間.
與其在家.
不如來教會襯托一下人氣.
再說.
今天什麼東西難了我們.
很好的親近神.
不能夠很好的熱切的感恩和敬拜.
如果你和我關注崇拜.
聽道的時候.
唱歌的時候.
心態.
人在心不在.
這樣又怎樣解釋呢?.
其實我們是否可以再投入一點.
我們是否可以再熱切地學毛多一點呢?.
究竟.
說回頭.
你的心有沒有空處為主.
在你的心靈的深處.
你是否真的尊崇耶穌.
信服神呢?.
以受教的心態來敬拜我們的主呢?.
是願意真的稱頌神.
數算自己的恩典嗎?.
還是主耶穌只不過是可有可無.
一時有一時無.
是我的一個忠義研究的對象.
而我是旁觀者.
這樣的心態.
我就算.
只是回教會.
關注神和教會的事.
我只是聊聊天.

$^{441}$我不是新人.
我不知道.
在教會講道的時候.
你的心態是怎樣的.
有些人是抽離的.
有些眾籌.
是抽籤的.
究竟是怎樣呢?.
如果我們失落了對主的傳言為心.
難怪我們會這樣.
我們沒有一個受教.
至愛尊崇主的回憶.
又或者.
可能我們回教會時間久了.
信主的時日久了.
我們已經失去了起初.
那種愛主的衝勁熱誠感恩的態度.
失去了虔誠的特質.
有些人對教會的毀身.
變得大不如前.
對教會的事或者問題.
都抽離.
都不想理那麼多.
不知道你是否這樣.
但願今日的主人崇拜我們重溫一下.
主對你的大愛.
好像那個女人一樣.
對主的愛再激勵你們.
讓我們的愛慕.
渴慕親近主.
而不再是一個狂歡抽離的心態.
我們要請賊的弟弱.
是過去的虧欠與現今的錯誤.
西門和他的同道當中.
認為那個女人是一個罪人.
其實當年我們說你是人渣垃圾.
有些人在中國說.
童幼平日裝童幼.
臭坑出臭草.
不會好得去變.

$^{481}$永遠都是這樣.
不能變的.
但主耶穌我們的主.
他是用另一個角度來看人.
他講了一個比喻.
證明兩個欠債者被免債的時候.
那種感激的心情的差異.
然後又論及這個女人剛才對他所做的每一個行動.
甚至指出這一晚.
公道主西門沒有做的.
這個女人一天都做了.
並且在程度上或者形式上.
發揮得更深切.
更熱切.
更深刻.
所以最後主耶穌的結論就是說.
他做那麼多.
做得那麼深.
不是因為他本來是怎樣好.
怎樣聖潔.
而是因為這個女人.
深刻體會到主耶穌赦罪的恩典.
臨到他身上.
我們的主耶穌看人看事.
有一個熟的原則.
是和世人和西門不一樣.
兄姐們.
西門和我們很多人一樣.
往往我們看人看事.
經常著眼於他的過去.
他的成就.
或者什麼缺陷.
我們看人評價人.
常常會因為他過去的表然而作為一個定律.
中文有說.
一次不忠.
八次不容.
是有些道理.
不過主耶穌基督的恩典.
福音的訊息.

$^{521}$有時卻是令到我們感到.
有違常理或者有違常情.
在這位聖潔兒身為萬有之主的耶穌的眼裡.
向來對這些人看人很敬虔.
遵守律法.
很了解經典的法利賽人.
原來同樣被視為欠債的人.
反而那位聲名狼藉.
人皆鄙視的女人.
卻被主認為是愛主和尊崇主的典範.
為什麼呢?.
這是什麼道理呢?.
原來這是福音與恩典的道理.
福利賽人西門與我們今天很多人一樣.
看完的時候.
只能看到他的過去.
尤其是他的過去的失敗,劣跡,壞事.
在人情淡薄,功利主義抬頭的社會.
尤其是這樣.
我們只會抽斥人.
有些不對,有些軟弱,犯錯的地方.
我們放大來看來說.
如果你記得以前有位歌星叫做蔡楓華.
他說了一句話.
他說什麼呢?.
「殺下光輝不是永恆」.
這句話其實很有哲理.
但是誰知道這句話.
迫得他在吳樂園歌唱界很難再立足.
很多年前.
不知記得在哪個電視節目.
不是《今日睇真D 》還是《星史追擊》.
就是陳年的電視節目.
訪問過他.
他說錯話的電視節目.
蔡楓華有很多的講座.
好幾個月前.
有黃心穎,許志安的一些安心事件.
發生之後其實很不安心.
因為令到黃心穎要離開吳樂園.

$^{561}$許志安連要舉辦演唱會都要取消.
黃心穎曾經拍過一些電視劇.
都要被雪藏.
不過好幾個月之後.
那個《留下淚高音》.
就成為無線的台慶劇.
可以現在.
是不是我們的社會.
是常常不容許一些人犯錯,講錯的呢?.
或者做錯的呢?.
就很難再翻身.
很難再接納.
很難再重建.
這些事令到很多人.
神經衰弱,精神飽受自我.
很抑鬱,很困惑.
社會是這樣對人的.
但教會是不是這樣呢?.
只是看人的過去.
只是看人的缺失和犯錯.
只是評論有哪些不對.
做了不好的事呢?.
放大來看一些事呢?.
但我們萬有的主耶穌.
不是這樣看人的.
這正是我們所信福音.
所宣示的訊息.
也是恩典的內涵.
當我們小時候.
41到46節的經文的時候.
你更加會發現.
耶穌所論及的那個女人的說話.
都是關於她當下.
主耶穌對她的立事回應.
要她怎樣的行動.
也是那個女人找到當時的機會.
奉獻和禮身.
表達愛主的那種表現.
這對耶穌來說.
他是這樣看人的.

$^{601}$看事處事.
就是我們主耶穌.
看每一個真誠悔改的罪人.
願意趁現今當下的機會.
找緊這個態度是很重要的.
主耶穌對人的看法.
是存有包容,接納和開放的態度.
主耶穌要問的.
不是這個人過去是怎樣不堪.
怎樣壞.
怎樣表現是一塌糊塗.
或者怎樣不理想,不合格.
主耶穌所看的.
是究竟這個人現今.
此時此刻.
他肯不肯對主開放和托靠.
只要人不將自己的吉利受.
自主的恩典.
他就有機會翻身.
可以重見.
主耶穌所關注的.
不是人以前過去怎樣壞.
怎樣失敗.
而是他現在此時此刻.
今天.
你肯不肯再依靠神.
好好的毀身.
好好的再重見.
主耶穌所看到的.
不是劣跡敗行.
而是在恩典當中.
你和我每一個人的生命.
是有新的機會,新的可能,新的盼望的.
這個就是我們所信的福音.
和恩典的真義.
靈姐妹.
讓你和我都有好像主耶穌那樣的心.
那樣的眼光.
有那樣的角度.
對自己的過去.

$^{641}$或者對他人的過去表現.
我們都能夠不再數算一些軟弱失敗.
不再被控訴.
不被內疚再驅趕.
而是抓緊現今重複神明的機會和恩典.
與求福音的大能再重見自己.
站立起來.
抬起頭來.
毀身給神.
再踏上重見之路.
讓我們不要再被過去的軟弱或者失敗.
或者是是非非所捆綁.
因為恩典正是向我們宣告.
主耶穌在主耶穌裡面.
祂給我們新的力量,新的機會可以重新再起力.
我們可以再重見.
忘記背後努力面前不是一個口號.
而是需要我們靠著恩典踏踏實實地.
抓緊走出來.
沒有人可以再控訴或者舊事重提.
因為主的恩典已經接納了.
已經赦免了.
這種不氣餒的態度.
可以幫助我們重新再上路.
因為這是我們重見的道理.
是有神的.
第三天.
我們看到.
愛和赦免是成正比.
主耶穌說兩個欠贈人的比喻.
讚賞完高末的女子之後.
作一個結論.
這個結論就是赦免多.
愛就多.
赦免少.
愛就少.
47字這樣說.
而主耶穌竟然對一個向祂不懷孔怯的法利賽人西門.
說愛.
人對你不好你還跟他說愛.

$^{681}$主耶穌在這個機會場合.
不正面的責備西門.
只是在第47字指出西門.
西門啊你愛神愛得太少了.
而愛得太少的原因.
原來是因為他體驗到赦免太少.
西門不像那個女人一樣.
深刻的經歷或者體會.
那裡虧欠.
而在這個法利賽人西門.
一向的信仰表現裡面.
他覺得自己很棒.
聖經的經典我很熟悉.
我不錯.
我多威風.
我很虔誠.
我很有實力.
人家都不及得我.
他很自大很自我.
其實這個是令裡面的驕傲自大.
不過主耶穌很巧妙的說話.
他並沒有公開的責備西門.
沒有譴責他.
主耶穌很寬文的.
很厚道的來開解他.
讓西門有時間有空間.
可以作出反省.
真是人算什麼.
你和我算什麼.
始終我們可以論辯批評什麼.
連主耶穌都完全沒有條件的來愛我們.
接納我們.
赦免了我們.
這個是他的恩典.
主耶穌是我們的主.
我們的基督.
他今天迎向你.
向我講說話.
提醒你和我.
赦免人多少.

$^{721}$反映我們愛有多少.
讓我們有更多的愛.
讓我們有更多的包容.
赦免對人.
對肢體.
甚至對自己不滿的人.
有更多的愛.
更多的包容.
赦免.
證明我們自己在生命當中.
是有愛的.
靖姐妹.
今天你對人究竟是怎麼樣.
是不是有些嫌隙.
你不放下.
有些是是非非令你很煩惱.
你仍然咬著不放.
是不是因為你赦免不夠.
其實是因為你的愛不夠.
你的包容不夠.
所以赦免幫助我們.
讓我們有愛.
有赦免.
在社會舊約裡面.
在這個宿命的家裡面.
我們出現有愛.
有赦免.
有包容.
因為只有愛.
只有赦免.
才是主耶穌的心.
最後一樣我要說的.
第四.
要對付生命中的西門.
在現實生活裡面.
大家會不會發覺.
有時候我們又是西門.
有時候又是那個有罪的女人.
很可惜.
是我們的聖潔謹守.

$^{761}$不及得法利賽和西門.
但我們也缺乏那個有罪的女人的謙卑.
和熟悉的洞悉.
我們在生活表現上.
有哪些態度好像西門呢.
譬如有論斷.
批評.
道場論斷.
觀龍感受體罰.
當中有沒有一些輕視.
或者不接納我們的弟兄姊妹呢.
如果不經意的.
在神面前.
自為是.
品評人的看法和做法.
其實是心裡面.
自我的驕傲和渴望作怪.
以致對人對神.
也欠缺那種恭敬.
謙卑的心.
如果有的話.
我們真的需要在神面前悔改.
另一方面.
如果我們對任何人.
包括對教會的肢體.
只是以眼前所看的片面的現實.
瀆暴地了解這個人.
判斷這個人.
而不是以一個期待.
以一個接納包容的心.
你和我都要求神.
承認.
唯有以欽點的角度.
會發生很多可能的角度.
用這個期望.
可能的角度去看人看事.
而不是釘死人.
批死人.
像那個女人一樣.
用謙卑倒空.

$^{801}$被主破碎的心態.
體會每一個人都需要神的恩典.
在恩典裡面.
學習不以以前或現在的她來行行去.
而以她將來在神面前.
可能的發生,發展.
來鼓勵她,認識她,支持她.
我相信唯有這樣.
這個愛的期待.
才有突破的可能.
和更新的空間.
如果用這個原則.
來看香港的局勢.
不只是黃絲,藍絲,黑衣,白衣,紅衣.
不是這樣分.
在香港彼此不同的意見爭拗.
可以撕裂,突破的可能.
投破血網.
但在教會是不可以的.
更加殘忍的.
用液體點火來燒人.
我們更加不可以.
我們是以神的道,聖經的真名.
作為我們的準繩.
我們是以愛,以善意.
以生命的真實來行事為由.
我們求神讓我們不要再挑剔.
更加求香港的情況.
警暴,示威人士的暴力,憤怒.
求主幫我們止息.
拆解那些冤仇.
止息一切的怒火和十恐.
停止一些強暴.
制止毀壞.
用主的愛,用主的大能.
來平靜這些躁動的心.
我們擺在禱告裡面.
紀念香港.
也紀念我們自己.
我們深知神是掌管一切.

$^{841}$總有祂美好的時間安排.
求主幫助我們.
能夠不學不用.
能夠以愛,以自拔,以包容,以赦免.
重建人的心和社會.
我們低頭祈禱.
主你自己恩代,鑑察,幫助.
我們在你面前是一個罪人.
但當我們倚靠你的時候.
我們可以更生成為.
在蒙恩成為一個聖人.
求主幫助我們.
感恩心態,示上,頌讚主在我們身上所做.
讓我們能夠有一個從你的愛,赦免,包容,接納,期待.
求你的真理和能力能夠到我們眾人身上.
求主你眷顧身心受創的人.
求你幫助醫治我們創傷的心血.
求你幫助有些軟弱.
有些需要你同在同行的人.
求主你的平安與我們同行同在.
祈禱奉主耶穌基督的名.
阿們.
\newpage



\section{}
\label{sec:SGo8UblJv4Q}
\textbf{Premiere Pro video editing (part 1)}
\newline
\newline
連結: \href{https://youtube.com/watch?v=SGo8UblJv4Q}{\texttt{https://youtube.com/watch?v=SGo8UblJv4Q}} ~~~~ 語音日期: 2019-12-05
\newline
\newline
\hyperref[sec:pqDLnSQO_8g]{\small{< < < PREV SERMON < < <}}
~
\hyperref[sec:index_chronic]{\small{[返順時目]}}
~
\hyperref[sec:index_scriptual]{\small{[返順卷目]}}
~
\hyperref[sec:joskqj5_Na8]{\small{> > > NEXT SERMON > > >}}
\newline
\newline
$^{1}$I am going to teach you how to use Premiere Pro..
If you already know the basics, please skip to Part 2..
This is Part 1 for those who have not used it yet..
First, you open the program and click New Project..
I usually don't change all the settings..
What you need to change is to make sure your location..
Make sure you open a new folder..
Make sure you set up everything..
Give it a title..
This time I am using Premiere Pro Video Editing Part 1..
After opening it, you will see what it is..
You don't have to worry about a lot of things..
You can adjust everything..
Basically, these four boxes can be moved freely..
You can move a lot of things..
For example, if you press the tool box, you can drag it to another place..
Yes, you can drag it up..
But everyone has their own preference..
So don't move things around..
For example, go to Tools..
You can drag it down and put it in your preferred position..
The most important thing is to import your files, audio, video, and pictures..
So you will see the "Import Media to Start" at the bottom left of my screen..
You can right-click here and you will see "Import"..
But if you don't have this, you may not see it when you first do it..
Go to the file above and find "Import"..
It will be the same..
First, you can choose a video or a picture..
But sometimes, if there are some files that are not clear in the name, you may not remember what it is..
There is a way to see it..
Just highlight the name and double-click the icon..
You will see it pop up in a corner..
That window is your preview window..
My setting is in the upper left corner..
You can see the video by clicking the "Play" button..
It's a preview, right?.
If you are a picture, you can see it..
But you can't play it..
You can see the timeline is moving, but it won't move the other days..
It's a picture..

$^{41}$As for the audio, you can play it..
There is no topic, but there is sound..
That's it..
How to start making videos?.
You can see the timeline in the bottom..
It's like your workspace..
For example, I want to put a video in the timeline..
There are two ways..
First, you can drag it directly..
You hold and drag it..
Or you can click the "Left" button on the cursor on the screen..
You can drag it to the timeline..
That's it..
You can start your timeline..
You can see the video is about 10 to 15 seconds..
The video can be extended or shortened by clicking the "Short" button..
You can see the video is getting shorter and shorter..
You are looking at the video getting longer and longer..
You can see the video is 1 minute and 30 seconds..
This is the function of zoom in and zoom out..
The above function is called V123..
You can play more than one video..
You can adjust the height..
You can also add audio..
If you have a lot of audio or videos,.
you can right-click on the space next to the eye and click "Add Track"..
If you want to add more tracks, you can click "Delete Track"..
It will ask you if you want to delete the track..
You can highlight and uncheck the tracks..
I will tell you the difference later..
You can mute the audio tracks..
For example, I want to add a picture..
You can click "Add" and click "Picture"..
When you play the video, the picture will be changed..
I will add one or two more videos to show you the difference..
I will keep clicking "Video", "Picture", "Video", "Picture"..
I will add Tony's audio..
You can see that I can't add audio because there is a "Video" button..
I will add the audio..
If you make a video, you will see that Tony's audio is very long..

$^{81}$I will add a lot of audio to fill up the space..
I will make the audio short this time..
You can see that the audio is not long..
We need to adjust the time..
I have a music player, so I will play the music..
This is the final product..
I only have a few pictures..
I want to make the audio longer..
What can I do?.
First, the picture is long and short..
When you highlight the picture, the cursor will be on both ends..
When the cursor is on the end of the picture, the red bracket will be changed..
This is to shorten the picture..
The picture will not move, so it will stay in the same position..
You can drag the picture as long as you want, but it will not move..
If I want to undo, I will press "Ctrl + Z"..
I will play the music..
I want to make the audio longer, so I will play the music..
The video can be long or short..
I will adjust the duration..
Now the duration is 100\%..
I will adjust the duration to 50\%..
Why is the duration not changed?.
Actually, the duration is not changed..
I can play the music longer..
The video was long, but now it is longer..
When I play the music, the video will be slower..
Can you see it?.
The video is much slower..
I will shorten the video..
The video is longer, but the video is shorter..
If the video is shorter, the duration is 150\%..
Now the video is only 16 seconds..
The video is shorter..
The video will be faster when I play the music..
This is how you play with video..
What else can be changed?.
For example, I need this picture to be 10 seconds long..
But I am bored..
There is a trick..

$^{121}$Go to the "At the beginning" of the timeline..
Go to the "Back 1 frame" of the timeline..
This is a very useful trick..
This is still a video of a girl..
This is the picture..
How can I adjust the duration of this picture?.
How can I make this picture not boring?.
There is a way..
Zoom in and zoom out..
First, highlight the picture..
Then, click on the "Preview" box..
Then, click on "Effect Controls"..
Then, click on "Highlight Picture"..
This is a little complicated..
The program reads each frame..
At the beginning of the frame, the frame rate is 100\%..
The position is 960x540..
There is a stopwatch..
First, click on the stopwatch..
Then, you will see two points..
What is strange?.
I will zoom in a little..
You will see three small points..
At this time, the program will know the frame rate is 100\%..
I want to make it 200\% now..
You can see the frame rate is bigger..
I want to zoom in the picture..
But I want to zoom out after 5 seconds..
So, you have to click on the "Add Keyframe" button..
You have to add a command here..
Click on this frame..
You can see two points..
The computer knows what to do at this moment..
You can choose the frame rate you want to change..
Click on the stopwatch..
At this time, the frame rate is 200\%..
But I want to change it to 100\% at this time..
What should I do?.
I will zoom out now..
But you don't have to do every frame in between..

$^{161}$You just need to tell the program that you want to change the frame rate to 100\% at this time..
Then the program will play out the zoom function smoothly..
I will play it for you..
Don't just look at the bottom..
You have to look at the top..
I will show you again..
It is zooming out..
Now your video is not so boring..
Because you have something moving..
Although you are playing the picture..
Zoom in and zoom out..
Just move the scale is one of the ways..
But do you remember that I asked you to put the stopwatch in the position?.
In addition to the scale, you can also play with the position..
Okay, let's go back here..
This is the final product..
I want them to be like this..
But at the beginning, I want the woman to be a little higher..
What should I do?.
960 is like a vertical axis..
Horizontal and vertical..
When you put the cursor on the 960, you will see two arrows..
You can hold it and move it left and right..
You will see that the picture is actually moving..
For example, I want to focus on a black man..
Or maybe these two black people..
At this time, focus on these two black people..
But when I go outside, I want it to go back to the middle..
I just need to move one time frame..
Now when you play, there are different effects..
You can play with this setting in many places..
But remember the original position..
When we put the picture in the middle, the settings of Premiere Pro are usually 986 and 540..
540 is the top and bottom..
If you don't remember, press 540..
The number in the middle is 960 and 540..
You can see that the screen is a little slow..
I set it to full resolution..
If you see the screen is moving, I usually set it to 1/4..
Then the screen will be smoother..

$^{201}$If you notice, I added a little more..
Because I adjusted the settings..
If you don't want to do anything, you can add more..
You can highlight it individually..
After highlighting, press delete to delete..
Of course, you can press Ctrl + Z to undo..
You can play with the picture to make it more vivid..
Zoom in and zoom out..
Zoom in and zoom out from the left and right..
You can also adjust the time..
You can also zoom in and zoom out..
You can also do this function..
For example, you can zoom in on this woman at the beginning..
Then slowly zoom out..
You can also see the background on the side..
You can see that the transition between each video is quite boring..
Because when you play, it suddenly changes to another one..
How to make a more interesting transition in the middle?.
This is called dissolve..
It is not called transition in this program..
How can I do this?.
You can go back to the import box..
You can see the top of each box..
If there is a double arrow, you can press it..
Then look at the effects..
It's not effects control..
It's effects control..
Then it tells you that there are many presets, audio transitions..
You can slowly cut it..
But if you know your function, you can type it directly..
These transitions are called dissolve..
You haven't finished playing, but it's usually out..
It's out..
There are many different dissolves..
It's actually a transition..
For example, cross dissolve..
Highlight it and drag it down..
Just put it at the end of the box..
I want to show you something..
I just went to the effects..

$^{241}$Where did I go?.
You just have to press the project button..
You will be back..
I found something..
If you have two pictures..
I have two pictures on the side..
Put a transition in the middle of the two pictures..
Go back to the effects..
When you put cross dissolve here,.
You can see that I can put the dissolve in the middle of the two pictures..
You can see that the two pictures are highlighted..
You can put it on the left..
On the right..
But it's also in the middle..
This is the dissolve function..
Put it in the middle of the two pictures..
When it transitions, it's called cross dissolve..
You can see it slowly..
Cross because both will appear at the same time..
Then dissolve to another..
A very smooth transition..
If you like this style..
If you don't like it, you can change to another..
For example, if I don't want it, I can highlight it and press delete..
I can dip it to black..
You can see that the glow will be black and come out..
Or dip to white..
There are many different dissolve methods..
There is a method called film dissolve..
I think it's the same as cross dissolve..
What I want to say is that no matter what you are,.
This dissolve transition can be very long and short..
First highlight it..
Then you will see that when the control box is opened,.
You will see that something pops out..
I usually don't need to use the top one..
I just need to do it here..
Because after highlighting, you can see that the cursor is in the film dissolve box..
There will be a red bracket..
Press it to lengthen it..

$^{281}$Just do it like this..
You can go to the preview box..
When you mouse over, you will see many different buttons..
First, you can press the lengthen button..
You can see that it is only lengthened one side..
It is different from the previous one..
Both sides are lengthened evenly..
If you see this symbol,.
You can move the box forward and backward..
There are different effects..
Because when you move it forward,.
The left hand picture will disappear very quickly..
The right hand picture will come out very quickly..
There are different effects..
I usually don't press the transition..
I usually put the mouse in the middle and press the left hand button..
It will be fine..
If you have the same nature,.
It is easy to do the same thing..
But if you want to do the same thing,.
Between a video,.
This is a video..
This is a picture..
I found that sometimes the cursor can't drop exactly in the middle..
You see, I can't drop the cursor in the middle..
What should I do?.
I can highlight the cursor..
Remember that we have this function..
Go to the effects control..
There is this function..
I can do it here..
I put the cursor in the middle..
I want the dissolve transition to be longer..
I can drag it back..
Sometimes you will find that the cursor can't be placed in the middle..
You can try to move it to the effects control..
You can play around with it..
This is one of the transition methods..
How to dissolve, how to transition..
You can cross dissolve, dip to black, dip to white..

$^{321}$These are very basic..
You can find it here..
It is what we need to do for now..
If you think it is clean cut,.
It will be better if you jump to another place..
You don't need to put anything between the clips..
After talking about video, let's talk about audio..
I have muted two audio clips..
Let's listen to the audio..
Let's listen to it..
The background music is very good..
But when Tony talks,.
I can't hear him at all..
What should I do?.
I want to keep the background sound..
The level is really good..
What should I do?.
You can do one thing..
Highlight the background..
How to adjust the volume?.
If you want to lower the volume of the audio at the same time,.
There are two ways..
First, highlight it..
Right click, go to audio gain..
You can see that it only changes the volume..
You can increase the volume..
Press positive number,.
Press 5, press 10..
You can see that the volume is reduced..
You can increase the volume..
Press 5, press 10..
You can see that the volume is reduced..
You can increase the volume..
If you want to decrease the volume,.
you can press negative..
For example, I think it's better to increase the volume..
You can decrease the volume..
You can see that the volume is reduced..
Let's listen to it again..
It seems better..

$^{361}$But Tony's voice is still very low..
So I will increase the volume..
I will increase the volume..
Press 5..
Let's listen to it again..
It's better..
Maybe....
When Tony speaks,.
the background should be quieter..
But the beginning of the background is very good..
I don't want to increase the volume..
I don't want to reduce the volume..
What should I do?.
Can I just increase the volume?.
Yes, you can..
But if you increase the volume,.
the sound will be too loud..
You should reduce the volume..
There are two ways..
First,.
you can turn the audio.
into two audio..
You can turn the audio into two audio..
You can keep the same sound in the beginning..
You can reduce the volume in the end..
But the sound will not be as beautiful..
Because the sound will be very awkward..
The sound will be very awkward..
How can you slow down the sound?.
First,.
put the timeline cursor.
a few frames before Tony speaks..
put the timeline cursor a few frames before Tony speaks..
Highlight the audio..
Preview the box..
If you don't turn on the effect controls,.
turn it on..
If you don't turn on the effect controls,.
turn it on..
If you don't turn on the effect controls,.

$^{401}$turn it on..
This is called "Levels"..
Remember,.
you are asking the computer.
to do something at this time..
So you need to set a time..
Press the stopwatch button..
You can see the diamond..
It means the computer knows what to do..
But what should you tell it?.
You want the sound to be.
quieter from here..
You want the sound to be quieter from here..
First, the audio and video are cut in the same way..
Use a razor blade to cut..
You can separate the whole clip..
You can cut it off independently..
You can separate the whole clip..
When you cut it off, you can delete it..
What if the big hole is left?.
There are several ways to cut..
Video and audio processing methods are the same..
I will show you the above audio..
There are several ways to cut..
If the clip is at the beginning or end of the clip,.
You can simply click the red bracket to cut..
\newpage



\section{}
\label{sec:joskqj5_Na8}
\textbf{Premiere Pro video editing (part 2)}
\newline
\newline
連結: \href{https://youtube.com/watch?v=joskqj5_Na8}{\texttt{https://youtube.com/watch?v=joskqj5\_Na8}} ~~~~ 語音日期: 2019-12-05
\newline
\newline
\hyperref[sec:SGo8UblJv4Q]{\small{< < < PREV SERMON < < <}}
~
\hyperref[sec:index_chronic]{\small{[返順時目]}}
~
\hyperref[sec:index_scriptual]{\small{[返順卷目]}}
~
\hyperref[sec:2Ci7sbMUq14]{\small{> > > NEXT SERMON > > >}}
\newline
\newline
$^{1}$First, I want to record a quick and easy video for you guys.
This video will give you a basic idea of how to do Premiere Pro.
I will start from the beginning.
First, you need to import your video.
It's best to sort it out.
As you can see, I have written down all the camera.
Camera 1, Camera 2, Take 1, Take 2.
It will be much easier.
First, I will do the opening.
You only need one angle for the opening.
So you won't need a lot of things.
You can see the import is done.
Basically, you click and drag it to the timeline.
You will see it.
Like this.
But we have a lot of takes.
You can edit it here.
For example, you can play it here.
(There is a Christmas, there is a Christmas, it's warm).
That's right.
But if you don't want to edit it here, there is another way.
Let's talk about how to edit it here.
For example, I want to do this last take.
You can put the cursor at the beginning or at the end.
You will see a red bracket.
First, press the cursor at the beginning.
You will see the video is cut when you drag it.
Like this.
Let me restore it first.
If I don't want to cut it like this.
What will I do?.
First, go back to this.
Press the video.
When I double click it.
You will see a small window here.
It will open a preview window.
I can play it here.
(It's so big).
I can play it until I get to this third place.
I want that take.

$^{41}$(There is a Christmas, there is a Christmas, it's warm).
This is what I want.
Let's go back.
I open a bracket here.
Mark it in.
Start from here.
(There is a Christmas).
It's not right, move it forward a little.
Go back here.
About here.
(There is a Christmas, there is a Christmas, it's warm).
Stop.
Then I mark in.
Can you see this? Mark in and out.
This is what I want.
First, I put the cursor at the video.
Left click, hold.
When I drag it down.
You will see.
The part I dragged is the part I bracketed just now.
Not the whole video.
This is another way you can do.
You can see the bar below.
Zoom in and zoom out.
This word.
You see I just cracked that video down.
Very good.
Usually.
If I just need one camera.
I will immediately get the audio down.
Import the audio first.
You have to organize it yourself.
Label all files.
Otherwise, it will be troublesome when you do it.
You must label it.
Because all files are photos.
The number will be very messy.
I have labeled my audio.
Audio.
Put it in.

$^{81}$You see my audio is also very long.
Because the audio should be so many takes.
I know it should be the third take.
So I can pull it.
Pull it to the last episode.
Okay, I need to synchronize it.
Not synchronize by myself.
Must use all programs.
So highlight both.
Right click.
Synchronize.
Must synchronize by audio.
When you press.
You see it lined up by itself.
Why is it lined up here.
Oh, that means tell me.
What happened.
I thought I got the second take.
Actually, I should have got the third take.
That means it can't line up.
Although I cut it.
But in fact, it can see.
In fact, the program did not delete.
It knows.
I am the whole video.
That means I did something wrong.
I actually want to get the third take.
I thought I was the third take.
So now let me do my wrong demonstration.
Okay, I actually need the third take.
Actually, if I want this part.
I will highlight both at once.
Trim.
Two videos at the same time.
This should be the third take.
Audio.
I play it once.
When you play.
Both tracks are on.
I don't want to.

$^{121}$I want to mute the first track.
I want to mute the original one.
Because I want to use the beautiful sound we recorded.
Now look at it.
Christmas is coming.
This is perfectly matched.
Because it's not manual.
Must use the program to synchronize.
When you play.
Why do you feel a little lacking.
Because you see.
The speed when I play.
It's just a quarter of the speed.
So when you play it for him.
There should be no problem.
But usually preview.
Don't turn it on full.
Because when you have too many running.
Your machine usually can't handle it.
And it's very slow.
You think things are not synchronized.
Usually I am one quarter or one eighth.
Depends on how big my video is.
Okay, the hard thing is coming.
Next step.
When you want to use multi-camera.
What to do when three cameras are doing it at the same time.
Do you remember every video.
We are all from three different angles.
First of all.
What to do is import this video.
I suggest you don't import all at once.
One take at a time.
So I put camera 1, camera 2.
And camera 3 take 1.
Import it first.
Then these highlights.
These cards are named.
So it doesn't have to be stuck together.
First, select them all.

$^{161}$Three different angles take 1.
Then right click.
Go to this one.
Called create multi-camera source sequence.
So what do you synchronize it by.
It must be by audio.
Then it will think about it.
Put the three videos together.
If you are lucky, it will be fine.
I have tried error message.
Maybe I can show you later.
Will also error message.
If it is an error message.
First, the first thing is to look at.
Will it suddenly show up here.
A called camera 2 or camera whatever.
Your take 1 but more.
Is multi-camera.
Let me see what the full name is.
See if I can show you.
Multi-camera.
If you have done this process.
A video like this is played.
Actually it works.
Even if it says you have a mistake.
You don't care.
First, highlight this video.
Drag down.
If you see there is a video above.
V1 is actually video.
There are videos below A1.
There are two voices.
You are already in good shape.
You have done the first part.
I deleted it.
Sometimes there is an error message.
There is this video.
There is a multi-camera video.
I drag down and only have a voice.
Only have a voice or only have a picture.

$^{201}$What do you do.
You delete it.
Highlight it, delete it.
Do it again.
I don't know why.
But it happens to me so many times.
As long as you keep doing it.
Make sure.
Of course, it can't synchronize that video.
There is a reason.
You highlighted the wrong video.
There are two videos that are not the same.
It can't synchronize.
But if you are pretty sure.
Double check.
Those three videos are the same.
Same video, different angle.
But it can synchronize.
I just keep doing it.
The second one will be take two.
Push it like this.
Take two.
Do it again.
Let's see if there is a failure this time.
I am also lucky.
Sometimes failure.
Say error message.
Sometimes.
Can do it.
It has a take two multi-camera.
Ah, I will do it for you right away.
Why is this one?.
If you are in this situation.
I basically delete it.
Say delete, yes.
I do it again.
Import the video.
Do it again.
Do it until it works.
Usually only two or three takes.

$^{241}$The worst case.
Basically, you keep importing all your videos.
Then put the audio below.
Synchronize.
This is the first step.
First look at the video.
I have put all the videos of the multi-camera.
So the opening is the start.
Then take one.
Take two, take three, take four.
See different lengths.
Below here.
I have started to put the audio.
Synchronize the audio.
It's definitely not handmade.
It's done by a machine.
So I continue to highlight three.
Actually not three videos.
It's the video and the audio.
You see it looks like three boxes of highlights.
But actually because the original video is with the audio.
Then continue to do it.
Synchronize.
Always like this.
Audio.
Then I think about it.
You see it doesn't seem to move much.
Because these two videos match very well.
The audio and the video.
But sometimes there will be a problem.
For example, I will put the audio of take two down now.
You see that sometimes our hands.
Everyone's camera and video.
There is a slight difference between opening and closing.
Sometimes it may not be perfectly matched.
Sometimes the audio is very long.
I will do a demonstration now.
First, if I put the videos in a seamless way.
Which is nice.
And you have to do this.

$^{281}$But sometimes in the early morning.
If the audio is not synced with it.
You will find that this happens.
First, I will put the audio and video of take two.
Line up.
Synchronize.
Because when it is synchronized.
It usually puts.
See.
The video will move around to adjust the audio.
I don't know if it has anything to do with the order of highlights.
But you will see when they want to align something.
It will move the video.
When you see it move to the right.
It will actually eat into the next video.
So these are things you don't want to see.
It won't push the video below.
It overlaps and eats it.
So first, press Ctrl + C to undo.
So don't usually line up the videos so quickly.
Leave a line in the middle.
In case they will move around.
Shift to the next video.
Yes.
After that, you can line up.
How to line up?.
I will zoom in for you to see.
Usually, there is no problem if the audio is not lined up below.
I basically want the video to line up.
So you drag it over.
Make sure the video is seamlessly lined up.
If they are usually seamless.
You will see a triangle here.
That means it is perfect when it is pinned together.
But if you eat it like this.
You will see a triangle on one side.
Not two sides.
That means you ate into another video.
So when you drag it down, you have to touch it.
Oops, wait a minute.

$^{321}$I forgot to highlight.
You see.
Make sure to highlight the three videos together.
Otherwise, you just move the top.
There is also a way.
If you want to group them together.
You group.
Here.
You move it, everyone moves together.
No problem.
But sometimes you remember.
When you want to highlight the video.
Ungroup it.
Otherwise, you move it all.
Move the audio.
Okay, let's continue.
The other two videos are the same.
After synchronizing it.
Line up.
You see me.
I'll do it for you quickly.
OK.
Drag it over.
Make sure the video is lined up.
Do the last one.
Take 4.
You see.
Actually before you officially add the video.
Actually there are a lot of these.
Line up.
Synchronize.
A lot of trimming work.
Have to do it first.
OK.
Synchronized.
Also all lined up.
So you can.
Start your first step.
If you are clear-headed.
You will see that our videos have different colors.

$^{361}$If you remember.
When we opened.
We said.
This.
Opening camera 2.
Used one camera.
So you see the video is blue.
And the other ones are green.
That is multi-camera.
Then.
You will see why it seems so random.
The first video is written Cam3.
The second video is written Cam2.
This one.
Highlight them with you.
What do I want to say.
Highlight them with you.
There is a difference when doing multi-camera.
I didn't remember.
When you highlight the first one.
When you click the first video.
It will take that angle.
Actually this is not important.
How to change the angle is the most important.
So everyone.
Need to know a button.
Is this.
There is a row of buttons under this big light.
You guys slowly familiarize yourself with how to use it.
Basically.
This is the marker.
This is Mark In.
Mark Out.
This color.
This Shift to In.
I rarely use.
This is.
Backtrack 1 frame.
This is Play.
This is Forward 1 frame.

$^{401}$This is Go To Out.
Go to Mark In, Mark Out.
I rarely use.
I don't know how to use these.
Export Frame.
Sometimes you can take a screenshot.
You can use this button.
Comparison View.
I don't seem to have used.
This button is what I want you to know.
This is called Taggle Multi-Camera View.
You can say.
I don't have it.
Why can't I find it.
It doesn't matter.
You go to the plus sign next to this.
There are actually many hidden buttons inside.
You haven't used it.
Maybe your button hasn't appeared yet.
So you can find the button you want to use.
For example, I drag one out.
This one.
You press it and then drag it down.
You can put it in this Toolbox.
Then click OK.
Then you can do it.
But I don't want this.
I'll cancel it.
This is very important.
I'll press it down for you to see.
After pressing.
My machine is a bit slow.
So I just saw it flashed.
It's called Media Blending.
Your machine is fast, it's okay.
There will be more cameras.
How do you see it.
First of all, it's the same as this icon.
The big camera is on this side.
The one on the right is the angle your audience sees.

$^{441}$So now the audience sees this video.
It will be in this camera 3.
Anyway, camera 1, 2, 3.
You label it yourself.
I have my own system.
Usually the camera closest to Mark Chan.
The so-called left-hand side of our audience.
Mark Chan's side.
Camera 1.
The front one is called Camera 2.
On the right side of our audience.
Camera 3.
I usually label it like this.
If you have your own system.
That's okay too.
Now the audience is watching in Camera 3.
You see here.
I divided it into smaller sections.
There are 1, 2, 3, 3 sections.
Because I only have 3 cameras.
I haven't tried it yet.
But I absolutely believe.
If you have 4 cameras.
Then there will be 1, 2, 3, 4 fill-ups.
If you have 5, then it will continue to split.
The computer also tells you.
You are watching in Camera 3.
You are using the camera 3 angle.
So it highlights the yellow box.
So the left side of this screen.
Just to let you know.
How many camera angles are there in this video.
We have 3.
But the big box on the right.
Is the angle that the audience is watching.
Why do you have to watch this.
Why do you have to know the split.
It's very important.
Because it will be very confusing later.
So everyone listen.

$^{481}$OK.
When I was watching this video.
I was watching and playing.
Play very slowly.
Because if your machine.
Is not that fast.
Actually doing this editing.
Is quite difficult.
Because it will keep splitting.
So I usually do it.
I will flip it.
I will turn it off.
Because it wastes a lot of memory.
When I usually watch the preview.
I will keep playing.
For example, I think I want to change the angle.
What should I do.
I will press this button.
Waiting again.
Yes.
There are two ways.
If your machine is playing fast.
I will set it to 1/8 or 1/16.
This is very important.
Especially because the more you open.
The slower it plays.
I have the 1/16 option.
Use it.
The screen seems to be very blurry.
But that's ok.
Because you are using 1/16 to watch.
Two ways to edit.
The first is to watch and edit.
When you keep watching.
The bar will move.
You can see a gap here.
You won't see any editing.
Until you call the machine to stop.
You keep watching and editing.
The clips will be made for you to see.

$^{521}$What is right for me.
You watch me do it now.
First, if you keep watching and editing.
First, you must press the play button.
While you are watching.
The machine is talking.
It is very slow.
It should be moving.
For example, I want to change the angle.
I will change to full screen.
You press this button.
Full screen.
Press.
It seems nothing happened.
You can see the screen has changed.
Keep watching.
For example, I want to change the angle.
I will change to Mark Chang's angle 1.
Press this button.
Camera 1.
Then.
It is working.
You can see it is not working.
It is working.
It will remember what you want to cut.
I have some delay.
You can see it has changed.
Ok.
For example, I want to stop here.
Stop.
You can see.
It is not cut yet.
Now.
After I stop.
The program will remember.
When I changed the angle.
It will cut the scene.
So.
Let me see.
I need to go back.

$^{561}$My machine is too slow.
It is hard to play.
While I am watching.
While I am editing.
It will help me change the angle.
(Laughter).
It has changed the angle.
It will help me change the angle.
My machine is too slow.
It has delayed.
(Background noise).
It will change the angle.
Next.
I will change to Mark Chang's angle.
It will change.
So.
This is a fast way to do.
Not fast.
It is called.
When you are watching the video.
You can cut.
For example, I am watching here.
Mark Chang suddenly said something.
I zoom to Mark Chang's side.
Look at Mark Chang.
It will be very smooth.
Very beautiful.
This is one of the ways.
The second way.
Is manual.
Not playing.
It is.
When you decide to stop.
You can use the razor blade to cut.
I delete what I did.
Now the video is beautiful.
I didn't cut.
The way to save memory.
Don't use the multi-angle.
Just play.

$^{601}$(Laughter).
Just play.
Okay, everyone is happy.
I decide to cut here.
You divide the video into two parts.
The first part is still angle 3.
But I want to start from here.
Go to the front.
Because I want to see everyone laughing.
So.
I will open this button.
Here.
Manual.
Press camera 2.
You see this video.
Immediately change to camera 2.
The whole video will change.
Then you continue to go down.
Go down to this point.
I want to go back to the second camera.
So you manually cut it.
So from this video.
You want to change the angle of the new video.
Then the video will turn back to camera 3.
See it changed immediately.
This is manual.
Both ways are fine.
Your machine is fast.
You just watch and cut.
If your machine is slow.
Just like me.
Turn around and cut.
Not much difference.
It will be easier if you are familiar with it.
But what I want to say is.
First of all, I regret it.
After I did it.
I found that this video should not be camera 2.
What should I do?.
It's very simple, right?.

$^{641}$Go back to this multi-camera setting.
Just change.
Basically, the angle has been changed.
Anyway, you have to make sure.
When you change.
You highlight the video.
Don't highlight the other video.
You highlight the other video.
You press it.
When you find it doesn't work.
Because you highlighted it wrong.
You have to highlight the other video.
This is one of the ways to change the angle.
Very easy.
But another thing is.
Oh, I think this video.
Stay at camera 3.
This angle is too long.
Actually, I should have changed earlier.
Change earlier.
Let's say.
For example, my next video is camera 2.
Change it first.
Usually my machine is slow.
Just like this.
Okay, now you see.
This video is camera 2.
This video is camera 3.
Camera 2.
I found that I have to change to camera 2 earlier.
This video is too long.
What should I do?.
This is the third trick.
You see.
The so-called clip trick.
The first is red, right?.
Just shorten the video.
But leave a gap.
This red bracket.
Can only be used at the beginning or end of the video.

$^{681}$If you are in the middle of the video.
Remember which button to press.
Just put the mouse here.
Press ctrl to make it yellow, right?.
Then drag it.
The video will be seamless.
There will be no gap.
But what's the problem here?.
It will make the video shorter.
Your neck is not beautiful.
Ctrl + Z to come out.
How to shorten this video.
That is, the time of camera 3 is shortened.
To change to camera 2.
The time of camera 3 is shortened.
To change to camera 2.
First put the red bracket in the middle.
Not in the middle.
Put it in the middle of the two videos.
When you press ctrl.
You see this is a new button.
You have never seen.
This has two arrow buttons.
Then you press the left button to hold it.
Then you shift it.
The advantage of this button is that the length of the video will not change.
What you change is to shorten the time of the camera.
I usually use this.
If I find that the length of the video is too long.
I want to adjust.
I don't want to affect other things.
I will use this button.
You see the video won't change.
It just changes the angle.
This is how I change the angle.
And adjust the length.
That's about it.
In the next video, I will continue to talk about how to trim.
Because I usually don't do camera.
This is just a demonstration.

$^{721}$I won't change the camera angle.
After synchronizing the video and audio.
I first trim the unwanted things.
So the correct order is.
Remember the first time I imported the video.
The video has to be multi-camera.
Then put the beautiful audio in.
Synchronize the audio.
After everything is in sync.
What you have to do is start watching the video.
Then you see something that doesn't fit.
Or some jokes, just cut it.
How to cut.
I usually put a bookmark here.
Because the bookmark may not want to cut off.
I still think about some things.
Maybe cut off.
Maybe not cut off.
Some of course, the NG points must be cut off at the first time.
But some may be the topic or the sentences are not good.
I don't know if it will cut off.
I will use the above marker.
Mark where it is.
Then cut them.
The first video should have talked about the basics.
I won't talk about it anymore.
But there are some tricks that I want to tell you.
Re-tune.
Talk later.
Give you a little more example.
See how I cut.
For example, I am here now.
I found a place to stop.
How to do it.
I will play it for you to see.
I found that I can stop here.
I added a bookmark.
Actually, it's the above.
Add bookmark.
I found that the middle paragraph can be cut.

$^{761}$I added a bookmark.
Usually, after adding a bookmark.
You press it.
It will line up.
Then I will use the razor blade to cut.
I just happened to.
Because we remember that our machine failed.
So I will have a video.
I have a beautiful soundtrack.
Then add MP3.
So I have two soundtracks now.
One video.
And I zoom in for you to see.
Actually, the soundtrack that follows this video is here.
But I don't want to see it.
I just pulled it up.
Then cut it.
Then you can do it.
Put it back in your marker.
Do the same thing again.
Usually if you are in the words.
In the words.
These questions.
I usually zoom in a little.
Make sure I cut it beautifully.
Use the razor blade to cut.
Cut.
Cut it all.
Cut all I want.
Video and those two sound clips.
Remember to go back to Arrow.
Choose your tools.
There are two ways.
First, I decided not to do this kind of thing in the middle.
Highlight it all.
Remember not to press delete immediately.
You press delete.
There will be a big gap.
Like just now.
You have to drag all the.

$^{801}$Pull it all back.
Very troublesome.
So how to do it.
If you delete.
I press Ctrl + C.
Remember to highlight.
What trick to do this time.
Just press Shift.
Then delete.
See it.
The middle piece.
Delete all.
But all below.
Will automatically line up.
That's one of the ways.
it is good.
There is another way.
Just.
I.
I will press Ctrl first.
I just cut one side.
Although these two sections in the middle.
Now we are cracking this section.
Is what I don't want.
I just need to cut one side.
How to do it.
After cutting.
I cut that side.
Then now put this.
Line up.
I want to go to.
We just said we want to go to this place.
Don't want the middle piece.
Highlight.
Remember I have a trick.
Put an arrow here.
Then press Ctrl.
Then I pull it.
Pull it to the place I want to stop.
Let go.

$^{841}$Same reason.
You can do the same effect.
Just don't want the middle piece.
Instead of.
Cut two pieces like cutting a cake.
Delete the middle one.
I just cut one side.
Then pull the whole video.
Pull it to the place I want to let go.
But remember.
After you make these adjustments.
The markers above.
Will all move.
So.
How do I delete these markers.
Go to the marker.
Right click.
You will see.
Clear all markers.
You can clear it all.
So usually these markers.
I do a paragraph.
Then do it all.
I have to delete it all.
Clear it all.
Otherwise I will be very lucky.
Sometimes in the middle of these cuts.
You will feel.
Coincidentally, their positions are different.
Will it be very awkward.
Actually, you see my usual videos.
Sometimes it will jump.
But how to make the jump smoother.
Is in the middle of these jumps.
Remember we can change the camera angle.
When you do this.
The jump will be less obvious.
And you will feel like.
You seem to be suddenly focusing on talking.
Talking about angles.

$^{881}$For example, Mark Sir said.
He will talk a lot in the middle.
I decided to jump the angle.
When you play this video.
It will be very smooth.
Because you seem to be.
I changed the angle to focus on.
The person talking.
People will not notice.
You are jumping.
Not so obvious.
Is it like a lot better.
So usually.
First make a video.
How to do it.
Remember I just told you.
Line up all the audio to video.
After that, don't care about the angle.
Don't do anything.
Angle is basically done at the end.
Just watch the video all the time.
Put some bookmarks.
Where do you think you should cut.
Where not to cut.
How.
Usually watched.
About two minutes.
Or if you feel very smooth.
Continue to watch.
After a few minutes.
You put all the bookmarks.
Just go back.
Do your cutting things.
After cutting.
When you cut the whole video.
What will you find.
I found that the sound will be loud and quiet.
For example.
If you can watch Peaks.
You will know how to watch Peaks.

$^{921}$You will see.
Here.
Is a little quieter.
But when everyone suddenly gets up.
It will be very noisy.
Wow, what should I do?.
There is a way.
First is.
For example.
I found this video.
When someone is talking.
Very quiet.
This paragraph.
Crack it up.
Now you cut the sound.
After.
You just control this paragraph.
For example, maybe Cindy is talking.
Maybe Leo is talking.
Very quiet.
Click on this video.
Then right click.
Go to Audio Gain.
It will ask you.
How much gain.
Here is actually gain and reduction.
Add and subtract can also be.
Now it is too quiet.
So I want it louder.
So you are positive number.
I usually start with 5.
You will see.
These peaks suddenly got higher.
Or I do it again for you to see.
Exaggerated.
You can see the difference.
First do 10.
Or 15.
Press.
You see the peak is much bigger.

$^{961}$Just adjust like this.
First do this adjustment.
Because this adjustment.
There won't be so much echo.
But if you say.
I still want to be higher.
There is another way.
When you click the video.
You will see.
The effect control here.
Will come out.
Change levels.
This can also change the volume.
But somehow this volume.
Change.
It will be a little different from the control gain.
You have to play around.
Usually if I want to change the overall.
Change the sound.
Add a little.
I will use control gain.
I found that the echo is not so strong.
If you just do level.
For example, from 0 to 15.
Usually the sound will be worse.
This is how you can control the volume.
But specifically sometimes.
For example, someone said a sentence.
Suddenly it's very explosive.
I can also change.
Change levels.
I will talk about this later.
Basically it's like this.
Do the whole video.
Cut all the boring details.
After that, listen again.
Make sure all the sounds are consistent.
If it's too loud, adjust it.
After adjusting all the sounds.
Then go back to the last angle.

$^{1001}$Because by the time you cut the video.
You will find that this video.
Has been cut countless times.
Then you slowly adjust the angle.
Sometimes it will jump to Mark Chan's angle.
Sometimes jump to Hale's side angle.
Sometimes go back to the middle angle.
That's about it.
The order is always doing camera angle.
Do the last last.
Will do this step.
That's all for now..
\newpage



\section{}
\label{sec:2Ci7sbMUq14}
\textbf{12/1/19 主日講道:Connect! 密切連繫 (張麗仙傳道)}
\newline
\newline
連結: \href{https://youtube.com/watch?v=2Ci7sbMUq14}{\texttt{https://youtube.com/watch?v=2Ci7sbMUq14}} ~~~~ 語音日期: 2019-12-06
\newline
\newline
\hyperref[sec:joskqj5_Na8]{\small{< < < PREV SERMON < < <}}
~
\hyperref[sec:index_chronic]{\small{[返順時目]}}
~
\hyperref[sec:index_scriptual]{\small{[返順卷目]}}
~
\hyperref[sec:zYSzdaPaiDI]{\small{> > > NEXT SERMON > > >}}
\newline
\newline
$^{1}$我都同Irene同樣感覺.
係想到你同各位Dance妹.
你地講聖誕節快樂呢.
就好似遲左D 囉.
但係同你講聖誕節快樂.
又好似太快喎.
所以我不如想同大家講一樣.
講一個日子.
係今年ge 7月10號.
呢個對我黎講係一個重要ge 日子.
係gwo 度呢.
gwo 一日夜晚我非常之gui .
想著訓覺.
但係一陣間就俾gwo D .
救傷車聲同埋警車聲呢.
Sirenge 嘈醒我.
我紮醒jor .
聽到D 回音.
啊.
我就同Patrick講.
我話我地ge 生命.
再得一樣ge .
我地晚晚都會聽到呢D 聲ge .
點解呢.
係因為7月10號今年ge 時候呢.
我同Patrick由埠仔搬到黎新加哥呢度住.
係埠仔住左二十幾年啦.
好多朋友都唔明點解佢地會搬出黎呢.
我知埠仔地方大D .
屋平D .
啊.
泊車容易D .
Shopping一切都方便D .
點解會搬出黎呢.
係因為當初我黎到成域呢度事奉ge 時候呢.
我覺得北景都係好格式.
每一次呢我都要zar .
七八個字車返返屋企.
咁所以我自己已經缺jor 心呢.
就係我同Patrick都缺心就係.

$^{41}$我一轉睇成會係full timege 時候呢.
我地就會搬出黎住架啦.
目的只係有一個.
就好似螢幕呢度所講ge .
就係.
好講.
就係一個字.
Connect.
我希望同大家有更多ge 密切聯繫.
呢個係我地搬出黎ge 盼望.
希望大家多D 機會近D .
一齊食飯 傾計 祈禱.
無論係個人腳或者係一班人ge .
多D 聚 係多D 認識 多D 關懷支持ge .
我當然啦 人好有限制 我只係得一個.
所以呢 如果要同大家密切聯繫呢.
都係要時間ge .
所以如果係短期做唔到.
請大家都睇啦.
講jor 出黎呢.
搬屋係聖經裡面.
時時.
我地做到ge .
有時係出於一D 呼召.
有時係一D 特別ge 原因.
就好似亞伯拉罕離開吾爾.
去到一個佢唔知ge 地方.
呢個係神秘ge 呼召.
又好似猶太人 猶太人佢地.
因為饑荒從猶大搬到去埃及.
跟住又峰迴路轉ge 命運.
呢個係特別ge 原因.
就講我地坐係度ge .
大部分都係移民.
我地其實都係搬屋.
不過搬得遠D .
七千幾里由香港或者中國搬到黎呢度.
但係最遠ge 其實都唔係我地.
最遠ge 係我地ge 神.
諗o下佢由天搬到黎地.

$^{81}$佢ge 目的是咩呢.
又係呢一個字.
就係想同我地.
connect 密切聯繫.
所以今朝係呢度.
我地想同佢去睇兩個神.
佢要主動黎connect我地ge 故事.
不過講呢個故事之前.
想同大家講一個序言.
講一個序幕.
呢個序幕就係話說有一日.
神創造天地.
當時並唔係陽光捧照天下氣徵.
因為當時都冇太陽.
但係按《聖經》所記載.
當時空虛混沌.
淵源黑暗.
講真就係邪惡勢力ge 存在.
當時係完全失控ge 局面.
就係呢個時候.
神一靈同呢D 黑暗.
邪惡ge 勢力搏鬥.
如果我地今日用3D或者4K.
黎睇呢D 片影真係精彩絕倫.
但係我地冇得睇.
因為我地都未存在.
而神就係咁樣制服jor 呢個.
熱水同埋黑暗象徵ge 邪惡世界.
唔單止係咁.
佢要創造呢個天地.
帶黎一個新ge 結局.
所以我地都好熟悉.
在世界一個記載話.
神係六日裡面.
佢做光 做天 做地.
同埋做係上面一切ge ye .
每一樣都係獨特美好.
每一樣神都話.
係佢眼看你係好ge .
而神做最後ge 一樣.

$^{121}$可以講係最獨特.
就係你做jor 人.
而且係按住佢ge 形象去做人.
而且係讓呢D 人.
同佢一齊住係呢片園裡面.
而且俾佢地管理一切ge 組織.
但係我地唔明.
點解神唔徹底消滅邪惡ge 勢力呢?.
因為仲有條蛇走出黎睇光陣.
呢個蛇係代表jor 邪惡ge 魔鬼.
現現ge 存在係度.
而佢心中裡面最忌ge 一樣係乜呢?.
就係gwo D .
佢唔驚獅子.
唔驚牛.
唔驚大象.
唔驚鱷魚.
唔驚恐龍.
佢驚ge 就係人.
佢驚人.
因為人係按住神ge 形象做ge .
所以奸險ge 蛇.
邪惡ge 勢力.
佢要忍有刑.
背叛神.
呢個故事我地到轉頭都識到.
我地知道好不幸.
俾蛇忍jor 之下.
我地人類ge 始祖亞當夏娃.
真係墮落jor 佢ge 行徑裡面.
佢地背叛jor 神ge 命運.
雖然我地明明知道係魔鬼靠害.
但係人唔可以話.
我係忌犯罪ge .
因為係佢地自己決定背叛神ge .
而佢地咁樣ge 背叛.
成熟ge 經驗.
好可怕.
咁點解關我地事呢.
羅馬書呢句說話講得好清楚.

$^{161}$佢話.
這就如罪是從一人入jor 世界.
死又係從罪而黎ge .
於是死就淋到罪人.
因為眾人都犯jor 罪.
因為亞當夏娃犯罪.
於是我地整個人類都有罪.
而事實上呢種罪.
唔係講下下ge .
你同我都體會得到.
你想唔想都好.
我地心內就有呢D 邪惡ge 罪孽係度.
而且我地心內唔使惡.
就會識得做好多罪ge 事.
呢一樣ye 係好急到會搞.
你同我都體會到.
因罪而帶黎ge .
撐脫敗壞破碎.
最後係死亡.
呢D 命運係你同我都揮之不去.
想我都唔得ge .
不過係呢度.
我地要稍為回答一下.
就係係呢段年度.
雖然神好似好無情咁宣佈.
因為你地犯jor 罪.
你地要受到好大ge 懲罰.
但係係同時.
佢好有情咁樣.
答應jor 一樣ye .
呢一樣ye 係唔少.
一D 都唔少.
就係佢答應將來.
係會徹底消滅邪惡.
將人救回.
而靠ge .
係一個女人ge 後裔.
係一個人.
只不過.
神好講呢件事幾時發生.

$^{201}$又只不過.
當神宣佈jor 呢種懲罰之後.
佢將阿當夏娃趕絕jor 出伊甸園.
冇得返屋企.
似乎人同神係完全隔絕jor .
似乎係咁.
好長ge 序幕.
我即刻講咩.
我話今日好想同大家講.
神搬到黎呢個世界.
佢點樣要同人去connect.
我想同大家講神同人connectge 兩個故事.
首先就係第一個.
神ge connect.
係從一個民族開始.
而呢個民族我地都好熟悉.
係猶太民族.
呢個民族後來.
神更加為佢地國家改jor 個名.
叫做以色列.
而神向佢地曉喻自己ge 名.
係耶和華.
獨一ge 真神.
創天造地ge 耶和華.
我地都知道.
但係你未必知道一樣ye .
就係原來京拜耶和華.
唔係由猶太民族開始.
近年ge 考古學發現.
對耶和華ge 京拜.
係始於迦南南面.
gwo 一D ge 猶無民族.
所以係《舊約聖經》裡面.
我地都讀到呢D 內容.
我講下.
哈巴谷書裡面.
神從提曼遙.
呢個就係迦南ge 南面.
跟住.
我地見到.

$^{241}$耶和華.
你從西耳出來.
呢D 都係南面ge 地方.
原來神早已係南面被京拜.
而係舊年考古學裡面.
更加發現.
呢D 迦南南面猶無民族.
佢地最叻係咩呢?.
就係打金.
燒銅燒鐵係佢地最本事ge 地方.
所以佢地當時好鍾意.
京拜一個咩ge 神呢?.
就係叫烈火ge 神.
而呢個烈火神ge 形象.
就正正設立係《聖經》裡面.
所形容ge 神.
好似龍光咁.
閃電,烈火,雲柱,火柱.
呢D 我地都睇過係《舊約聖經》裡面.
形容我地ge 神.
係迦南南面猶無民族.
早已京拜呢一位.
同以色列人一齊京拜ge 神.
只不過呢一D 南面ge 民族.
佢地唔單止係京拜耶和華.
佢地係京拜埋好多好多其他ge 神明.
咁你諗呢件事講明咩ye 呢?.
講明ge 就係.
雖然耶和華係列邦gwo 度已經被人京拜.
但係係佢ge 心裡面.
一直有一個好想救人ge .
救人出黎追ge 營份.
ge 一個策略.
就係佢會接觸一個民族.
從佢地中間興起救世主.
講哂一句就係.
神係好serious about佢ge 拯救.
佢唔係講呀.
然後我地又讀到.
神點樣向猶太ge 一位族長亞伯拉罕顯現.

$^{281}$話比佢聽.
你要將你全族.
由吾以大一個你都唔知咩ge 地方.
你跟住我去行就得啦.
我承諾話比你聽.
人類ge 救世主.
會係你呢一族出生.
我承諾你地.
萬國必因你ge 後裔得福.
我承諾話比你聽.
你地會得福.
萬國亦都會.
萬族亦都會因你地得福.
而你睇我地見到.
神點樣去接觸呢個民族呢.
佢不斷向佢地顯現.
同呢一個民族.
佢地ge 人生.
同佢地經歷高高低低.
請容我同你再分享一個.
好有意思ge 詞語.
係《釋經》冇出現過.
但係卻係時時ge 係發生ge .
Shikain Glory.
Shikain係一個希伯來字.
意思係高高在上ge 神.
向人顯現.
向人藉住塵世ge 物件.
向人顯現.
神剛成為塵世ge 人或者物.
向人顯現.
呢個就叫做Shikain Glory.
咩意思呢就係.
我地睇到係《舊約聖經》裡面.
神揀選jor 呢個民族.
猶太民族之後.
佢就不斷ge 向佢地顯現.
幫我等下.
佢化成jor 為人.
同阿巴拉罕對話.

$^{321}$呢個我地都好熟悉.
神化成荊棘裡面ge 火焰.
同人同摩西講說話.
神成為運柱火柱.
帶領以色列人出埃及.
神係會幕gwo 度一路同佢同在.
神係沙瀨山上面.
運中降臨.
神親臨.
所羅門王所見ge 聖殿.
神同人立約.
呢一切都係話.
個屋在上ge 神.
佢竟然黎到呢個世界.
佢用人間可以接觸ge ye .
去向人顯現.
話比佢聽.
我與你同在.
呢個就係第一.
我今日要同大家講佢ge 故事.
第一個就係.
神要connect一個民族.
作為佢拯救呢個世界ge 開始.
只不過係呢一個民族.
身份就好特別.
但係佢地ge 性格就冇乜特別.
同你我都一樣.
佢地係罪裡面.
佢地死死夾夾敬拜其他國家ge 神明.
佢地唔忠於呢一位創天造地.
帶領佢地.
一路帶領佢地.
呼召佢地ge 神.
ge 耶和華.
而且佢地唔單止一代係咁.
差唔多代去代都係咁.
最後佢地過份到一個地方.
激怒到神.
神話我要離開佢地.
我要離開聖殿.

$^{361}$睇呢一齣黎.
神係要做一個決定就係.
我要同你地disconnect.
神同人disconnect.
總結以上.
我講ge 係.
神雖然揀jor 呢一個民族.
想藉住佢地.
去拯救世人.
神藉住好多繁盛ge 人與物.
黎到向呢一個民族顯現.
黎到實現佢ge 拯救.
可惜呢一班民族係好低死.
佢地不忠於耶和華.
對佢地ge 愛.
佢地激怒到神要離開佢地ge 聖殿.
神好似同佢地disconnect.
但係微妙ge 就係.
係幾百年之後.
當冇人諗到ge 時候.
神又對於呢一班離棄佢ge 選民.
做jor 一件不可思議ge ye .
就俾佢地睇到.
神原來從來冇同佢地disconnect過.
相反.
而家神仲要行多一步.
黎到去connect你地.
雖然係你ge 罪.
你都封鎖我唔可以幫你地.
但係我做一樣ye .
就係去到我地講ge 第二個故事.
神又connect.
由一個翻牆開始.
翻牆熟悉電腦ge 人.
呢度大把熟悉電腦ge 人多.
我扮專家講幾句.
翻牆ge 意思係.
電腦裡面有好多時候有D 叫火牆.
呢個火牆firewall.
為jor 呢D program.

$^{401}$避免你比一D 病毒入jor 你ge 電腦.
呢個係一樣.
又或者係有D 極權ge 政府.
佢係要用一D 火牆.
係防止市民ge 人民.
唔可以去到國外ge 一D 網站.
封鎖佢地ge 資訊.
就係咁.
呢D 係火牆.
只不過係有一D 電腦ge 高手.
佢寫jor 一D program.
使人能夠突破jor 呢一種ge 網絡ge 封鎖.
去到gwo D 去唔到ge 網站.
可以翻過火牆.
呢個叫翻牆.
點解我講呢D ye 呢.
因為呢個係一個好好ge 比喻.
雖然神曾經主動connect人.
但係人因為罪性.
仍然有一度同神隔開jor 路過ge 牆.
我地可以講話呢個係一個火牆.
人因為係咁翻唔到神為佢地創造ge 樂園.
但係神做jor 一樣不可思議ge ye .
佢要翻越呢個火牆.
係藉住咩ye 呢.
不可思議ge 就係呢句ge 說話.
耶穌親密ge 門徒.
佢親眼見過.
佢寫jor 呢句說話出黎.
佢寫jor 呢句聖經出黎.
道成肉身住在我門中間.
充充足足有欣點有真美.
道成肉身.
呢四個字相當之震撼.
震撼jor 我地整個人類.
係神佢竟然成為jor 有血有肉ge 人.
佢搬到黎呢個世界同我地一齊住.
為jor ge 係乜ye 呢.
為jor ge 係要拯救我地.
佢翻越呢一個牆.

$^{441}$為jor ge 係要拯救我地.
翻越一道我地同神.
隔絕ge 火牆.
道成肉身呢.
黎到我地中間.
住係我地中間.
充充足足ge 有欣點有真美.
開始ge 時候Mark Sir都提過.
今日我地慶祝聖誕係咩呢.
不如我地將我地目光.
倒返轉 吞返D 離佢去睇.
我地慶祝張能權.
我地人類ge 救主快將降臨.
但係聖誕節係咪Party 禮物 裝飾係最重要呢.
重要 你同我都鍾意.
但係千祈唔好將呢D 係掩蓋jor .
聖誕節真正重要ge .
就係我地慶祝人類ge 救主.
佢翻牆ge 黎到我地呢個世界.
道成肉身黎到我地呢度.
佢搬到呢個世界.
佢衣食住行唔再係天街ge .
係繁盛ge .
而且佢揀選擇黎到呢個世界.
唔係去gwo D 高貴ge 地方.
佢揀ge 就係好平民ge 地區.
好似我地呢度McKinney Park.
Gridgeport Chinatown.
係D 好平民ge 地區.
佢搵ge 餐廳亦都係去.
如風尾呀 美妮芝ge 呀.
Basement 呀呢D 咁樣.
同我地好貼身ge .
佢選擇與我地一齊住係我地中間.
佢更加選擇ge .
唔係優厚ge 對待.
相反ge 係好惡劣ge 對待.
佢由得人咁樣對待佢.
佢對人妒忌 憎恨 誤解.
甚至係陷害.

$^{481}$呢D 人唔單止係當時ge 聖殿ge 領袖.
亦都有好多同佢親密ge 人.
有D 係佢ge 同鄉.
咁係咪無人對佢好呢?.
當然有 只不過係當佢危難ge 時候呢.
呢D 人都走jor .
但係注意 值得我地注意ge 係.
呢D 唔係意外黎ge .
我地唔使淨係計耶穌慘ge .
呢個係佢ge 計劃黎ge .
係佢計劃黎到呢個世界上面.
受到呢D ge 對待.
而至尊ge 目的係為jor 拯救你同埋我.
係罪裡面ge 困境.
使我地唔再殘缺 敗壞 破碎.
以致後死亡.
講返呢班神connectge 小民.
佢地當時完全get唔到 明唔到.
因為係佢地ge 歷史裡面.
係佢地一味睇到.
舊日ge 先知寓言ge 呢位舊世主.
佢地以為一位高貴ge 君王.
使佢地以色列國能夠光輝燦爛.
但係佢地冇諗到.
呢個舊世主.
以咁低下ge 身份黎到呢個世界上.
對佢地黎講好難接受.
但係後黎好多人都大徹大悟.
佢地到底明白.
被釘係十字架上面ge 耶穌.
原來就係神徹底connect.
為ge 係要實現佢當日係伊甸園.
答應過我地要救返我地.
道成肉身住係我地中間.
充充足足ge 有真理.
唔係咩.
佢道成肉身黎到世界上.
同我地罪人connect.
以至代我地稱秀嚴罰.
呢個係恩典.

$^{521}$而佢死jor 之後.
第三日佢復活jor .
表明佢能夠為我地人類.
打jor 一場徹底ge 仗.
亦都表明jor .
神係伊甸園裡面ge 承諾.
唔係空頭支票.
呢個係真理.
好可惜ge 就係.
等我地繼續睇呢一個.
約翰福音第一章ge 時候.
佢前文講ge ye 令到我地好尷尬.
好心疼.
佢前文講到話.
佢係世界.
世界亦都係藉住佢做ge .
但係世界卻唔認識佢.
佢黎到自己ge 地方.
世人都不接待佢.
但係雖然係咁樣.
佢仍然ge 黎connect我地.
為jor 懲戒我地.
佢仲要話比我地聽一個.
好寶貴ge 訊息.
就係凡願意接待佢ge .
就係信佢名ge 人.
佢就賜比佢地權臂.
作神ge 兒女.
意思就係話.
如果你聽到呢個故事.
你謙卑ge .
用你最大ge 真誠.
黎到接受呢一位ge 救主.
呢位降生為人ge 神.
耶穌ge 翻牆.
亦都可以使到你翻牆.
問題係.
你今日已經謙卑接受jor 佢咩?.
以上我好似淨係講一個故仔.
都未開始講到.

$^{561}$但其實講完啦.
因為今日.
踏入十二月ge 聖誕ge season.
令我覺得.
我地需要去重溫返.
呢D 咁重要寶貴ge 故事.
唔係因為呢D 故事.
令到我地覺得好窩心.
係令到我地覺得好溫馨.
而係呢個故事.
對我地人類係最重要ge 故事.
讓我地認識.
我地ge 神係幾咁愛我地.
我地又咁諗.
我地ge 神既要咁樣connect我地.
黎拯救我地.
如果我地今日.
係人生ge 孤單困惑裡面.
佢點會唔理我地呢.
各位弟兄姊妹.
或者你今日面對你人生ge 困難.
又或者係係香港今日ge 局勢裡面.
你會感覺到好窄聲.
或者我地都喊叫過.
神呀你幾時出手呀.
但係其實我地問錯jor .
我地應該問.
我地幾時出手呀.
係咁耐.
我地作為神ge 子民.
我地舉過幾多次頭部ge 手呢.
弟兄姊妹.
我地肯定一樣ye .
我地人希望ge 拯救方法.
係短暫ge .
因為係呢個有罪ge 世界.
今日我地忠言得到公義.
明日不義忌事又黎.
甚至係話.
今日你得到希治.

$^{601}$明日你ge 病痛又會黎.
所以我地要深信.
我地ge 神黎到呢個世界.
唔係拯救呢個短暫ge 世界.
佢係要拯救我地ge 永恆.
問題係我地有冇lor 出一個謙卑.
真誠.
降服ge 信心.
黎到接受佢比我地呢個靚得.
佢對我地ge connect呢.
係呢個聖誕節快黎ge 日子.
我想比大家.
同大家去操練兩個功課.
第一個係.
無論今日你ge 日子係點樣.
你要對呢位高天般樂美.
願意同我地重新connectge 神.
對佢有信心.
無論你日子點樣.
第二就係.
無論係幾唔容易.
我地都同身邊ge 人.
分享呢個故事.
讓佢地亦都能夠.
同神重新connect.
得到好似我地得到ge .
永恆ge 福氣.
願意我地一齊去做呢個功課.
我地一齊祈禱.
創造我地.
獨一ge 真誠.
當我地失散ge 時候.
你connect我地.
你黎到呢個世界.
而且.
你係付上好大ge 代價黎到.
使命我地ge 人生反轉過黎.
我地今日被光照著.
我地永恆更加得到呢個光去帶領.
讓我地唔單止.

$^{641}$而且同人去分享你呢個美妙ge 故事.
願神的話,也祝福大家.
\newpage



\section{2:1-7}
\label{sec:zYSzdaPaiDI}
\textbf{12/8/19 主日講道: 與人同在的神 6 - 路2\_1-7 (陳崇基牧師)}
\newline
\newline
連結: \href{https://youtube.com/watch?v=zYSzdaPaiDI}{\texttt{https://youtube.com/watch?v=zYSzdaPaiDI}} ~~~~ 語音日期: 2019-12-15
\newline
\newline
\hyperref[sec:2Ci7sbMUq14]{\small{< < < PREV SERMON < < <}}
~
\hyperref[sec:index_chronic]{\small{[返順時目]}}
~
\hyperref[sec:index_scriptual]{\small{[返順卷目]}}
~
\hyperref[sec:99dN28SUzkQ]{\small{> > > NEXT SERMON > > >}}
\newline
\newline
$^{1}$又到聖誕了.
我們做傳道文牧師的行業.
一定要講到聖誕節的經文.
大家知道經文不是很多.
其實如果真的要講也很多.
不過有很多我們不能解釋那麼多.
譬如舊約有很多經文.
所以有很多新約的故事.
我們一定會提一提.
一定會講一講.
所以有些人會說.
牧師你又翻炒了.
但我覺得勝經的故事.
聖經的經文是歷久常新的.
你再看細微的時候.
可能你現在現代的意思又會很不同.
你的心靈的不同經歷.
又會因著不同的際遇.
你又再看上帝的話語.
上帝對你說的話又有所不同.
所以就算到聖誕節.
或者明年的福祉.
或是後來的聖誕節.
我們都會看一些可能會很熟悉的經文.
或者會聽過耳熟能詳的經文.
這些故事或者記載.
都是希望能夠在不同的處境裡面.
對我們有不同的聲音.
我相信聖靈不斷在我們不同的.
人生的階段裡面都會繼續傳話.
只要我們願意去聽就行.
雖然我們講到聖誕節的故事.
是一些很耳熟能詳的.
但是我和一些弟兄姊妹去聊天.
分享和談話的時候.
我發現其實我們不是很熟悉.
你熟不熟悉聖誕節的故事?.
我想問你.
你問一問你旁邊那個.
我問你幾條問題.

$^{41}$看你旁邊那個是否懂得回答.
第一條問題.
到底東方有多少個博士去找耶穌?.
你問一問.
有多少個?.
有誰說你旁邊那個答了三個?.
有沒有?真的有吧?.
不是三個.
聖經不是講三個.
聖經講幾個.
當然不是三個,是幾.
我再問你.
東方的博士.
幾個博士,所謂三個博士.
幾個博士去找耶穌的時候.
到底耶穌當時大概是多大呢?.
你問一問你旁邊那個.
大概當時耶穌應該是多大?.
OK,有沒有答案?.
有沒有人答?大概兩歲左右吧?.
有吧?.
我看見幾個人.
你覺得耶穌多大?.
十五歲?.
大概應該是兩歲以下吧.
為什麼呢?你回去看看就知道了.
馬達荒第二章.
因為大希律的暴君聽到有猶太人生下來的王.
會搶他的位置,他很害怕.
然後他就計算過東方博士來的時候.
然後就去伯利林殺光那裡的男嬰.
兩歲以下都殺.
他計算到大概是兩年前.
一至兩年前.
那個星出現了.
所以我就發現其實大家不是很熟悉.
所以應該要很仔細去想一想.
到底聖達治的故事.
想跟我們說什麼.
我們很多時候想起聖誕節故事.

$^{81}$想起聖誕卡.
想起電視或電影裡面的一些情節.
我們經常覺得勝誕故事只是一個小孩的故事.
但你想一想聖經如何記載耶穌的誕生情節的時候.
你就會發現這個絕對不是給小朋友聽的.
待會我就告訴你為什麼.
我們記住這是一個系列.
我們其中最後的兩集.
這個星期和下個星期.
我們就會結束這個系列.
就是《與人同在的神》.
在上帝自己的聖經裡面.
告訴我們人類歷史應該怎麼寫.
人類歷史應該怎麼走.
它的大故事到底是怎麼樣.
而我們現在又在哪裡呢.
耶穌已經到了.
祂會再回來.
我們就在祂中間的時間.
所以我想review一下.
我們曾經說過的是什麼.
我們說過大概是.
由《創世記》第一二三章開始.
我們說上帝創造萬有創造人.
我們每一個人都有上帝的形象和他的本性.
有他的尊貴.
但是人背叛了上帝.
亞當·夏娃就是人類的代表.
就是告訴我們.
我們每一個人都背叛上帝.
但是上帝沒有wipe out(清除)我們.
上帝可以按他的公義.
去剷除所有背叛祂的人.
但是上帝去救贖他們.
然後就設立了整個的計劃.
去將人帶回祂的面前.
我們知道聖經的故事說.
神去救贖人.
是要藉著一個人.
一個家族.

$^{121}$一個國家開始的.
那個就是去到以色列的阿巴拉罕.
然後那個故事由阿巴拉罕.
到下一代以撒雅各.
到以色列.
他們因為饑荒而故落在埃及.
然後上帝用恩典來拯救他們.
離開埃及.
然後在西乃山跟他們立了約.
而聖經裡面說到立約的條文.
那些的law and regulations(法律和規範).
是教導一些已經得救了.
被上帝拯救了的人民.
他們應該怎樣去生活.
所以你看看摩西五經.
律法書裡面所說的.
那些律例.
全部都是告訴我們一個新的民族.
一個新的identity(身份).
我們已經屬於上帝了.
應該怎樣去生活.
然後上帝就帶領他們進入迦南地.
就是現在的以色列巴勒斯坦地.
在當時他們建立了古以色列國.
這個國家也是告訴萬國.
上帝自己的國度應該是怎樣的.
當然如果你熟悉聖經或者有些印象的話.
你知道舊約的故事.
就說這個國家其實都是灌滅他的.
所以以色列國.
由大國就各自滅亡了.
但是上帝恩著他的恩典的認去.
他沒有扔下他們.
也不wipe out(清除)他們.
就帶回他們回去.
以色列地.
他們被島回歸之後.
有幾百年的時間.
等到尼賽亞救世主來的中間.
上帝沒有拋棄他們.

$^{161}$直到第一世紀.
大概2000年之前.
就真正的救主尼賽亞降世.
就來到我們中間.
這個就是聖誕的意思.
這個review(評測)就告訴我們.
這是一個grand story(偉大的故事).
上帝一直在推進人類的歷史.
並且我們要看我們自己人類的歷史.
是用上帝的角度去看.
縱使我們生活裡面.
或者在我們人世之間.
我們所observe(觀察)到的.
是一些很混亂的東西.
但我們一定要明白上帝.
如何在當中繼續拯救.
所以我們今天.
既然聖誕節來到.
我們就講到基督的降臣.
對於一群已經背叛了他的國民.
一群心不歸他的國民.
上帝想做些什麼.
藉著救主.
想去到如何帶領他們.
回到神自己的面前.
我們講過這是聖誕節的故事.
絕對不是小朋友的故事.
我先讀頭三十.
路加夫第二章.
當那些日子.
凱撒奧古士都有子而後.
叫天下人民都報名上策.
這就是居尼努作敘利亞省巡撫的時候.
頭一次行報名上策的時候.
眾人各歸各城報名上策.
我為什麼說這不是小朋友的故事呢.
因為如果講小朋友的故事的話.
我不會跟你說.
當那些日子.
習近平做中國的領導人.

$^{201}$林鄭做香港特首的時候.
他就發生了這樣的事.
我不會提這些政權的問題.
我不會提這些政治社會的環境.
因為小朋友的故事是不需要的.
小朋友的故事通常是怎樣開始的.
Once upon a time.
很久很久之前.
有一件這樣的事情發生.
然後情節是怎樣.
很少會交代歷史.
特別是政治的氣氛.
所以打開路加夫第二章的時候.
看到聖誕節的故事.
就不像我們聖誕卡上看到的那些.
他先提的是什麼呢.
先提兩個的和.
一個叫做Augustus.
一個叫做Giulio.
如果你從路加夫第一章開始看.
他還提到一個本地的王烈.
那個叫做Helium.
所以一提到基督的誕生的時候.
耶穌的降生的時候.
他就先提一些政治的人物.
或者政治的環境.
這些當然是當時的讀者很清楚.
正如我今天跟你說.
習近平做領導人的時候.
或者林鄭治港的時候.
你就想起很多不同的事件.
或者新聞.
或者一些的動向.
你就明白那時候的時代是怎樣.
當時的讀者.
當時的聽眾.
他們聽見凱撒.
Augustus.
或者做這個.
敘利亞省的巡撫的.

$^{241}$Governor的Giulio的時候.
他們就明白到一個政治的開放.
這個Augustus.
就是Roman Empire裡面第一個的帝王.
他是第一位.
殺了所有他的政敵.
然後他去統治天下.
殺了很多人.
還有打了很多場仗.
最後沒有仗打了,為什麼?.
因為他殺了太多人.
就沒有仗打了.
這個就叫做天下太平.
你就想起一個這樣的時代.
Giulio都不是很多好邊.
他做敘利亞的總督.
他殺了很多不同的政敵.
然後將軍變成一個巡撫.
或者一個governor.
一個省長.
所以你從第一章開始看.
聖誕故事的時候.
你就想起一個這樣的政治氣氛.
並且當這個Augustus.
第一個的羅馬帝王.
你聽一下他們當時的所謂官媒.
就是他們國家的宣傳的新聞台.
他們說什麼呢?.
當時國家的宣傳新聞台當然沒有電子.
但是有很多這樣的神像.
為他自己雕鑽的神像.
放在不同的他所統治.
所打壓的一些城鄉城鎮裡面.
那些市中心就有他的神像在那裡.
也都在他的貨幣上面.
他去到講清楚他自己的身份是什麼.
這些這樣的金幣或者銅幣.
每一隻.
每一個的銀幣.
都是當時的官媒.

$^{281}$propaganda會講的東西.
如果你看到一點英文的話.
你會看到divine.
然後有個F.
其實拉丁文意思就是說.
Son of God.
Augustus.
Son of God.
他說他自己就是神的兒子.
他是統領天下的.
並且每逢羅馬帝國.
有兒子班過來的時候.
他們用的字眼都是說.
我們有福音帶給你.
有gospel.
good tidings.
來帶給你.
那個叫大喜的信息.
Augustus也是第一個去宣傳自己.
說他是Lord.
主.
他說他帶給天下羅馬的人民.
有救恩.
他就是savior.
在他很多銅幣上.
也都是告訴人民.
他就是神的兒子.
他就是主宰.
有他的地方.
羅馬統治的地方.
就是有和平.
有peace on earth.
並且他所管理的地方.
羅馬國度的南途.
就是他的國度.
他的kingdom.
這些全部都是當時羅馬帝國.
有的官媒.
平民百姓一定會接觸到.
所以這個故事告訴我們.

$^{321}$這個人作王的時候.
他是這樣去自稱自己.
並且是一個這樣的.
社會的政治狀況.
在經文裡面.
四次交代了.
我不知道你有沒有留意.
報名上冊這件事.
短短七節.
路加福音第二章一至七節.
講到耶穌降生.
是第一幕.
四次去講到registration.
報名上冊.
羅馬帝國每次報名上冊.
就不像我們美國.
我們人口統計.
人口統計到底有多少人.
demographic是怎樣等等.
不單單是這樣.
每次講到報名上冊的時候.
就有兩件事要做.
第一件事就是要計算你的人口稅.
到底有多少人.
男女女有多少head.
他會計算你的head的tax.
就是你的農產品.
你工廠裡面出現的東西.
你所紡織的物業等等.
你所produce的東西.
有多少是應該歸給羅馬王的.
head tax.
第二件事.
講到報名上冊的時候.
就是去徵召軍隊.
所以他去不同的鄉鎮的官員.
就會叫人去報名上冊.
然後就說你家裡面有多少個男丁.
然後有多少男丁需要.
來去當評.

$^{361}$當然.
都是猶太人來幫.
猶太人就這樣.
猶太人在羅馬帝國就不會叫他們當評.
因為他們知道他們很麻煩.
經常有動員.
但是每一次有報名上冊的時候.
就代表要加稅.
加稅就像牙痛一樣.
就像我們芝加哥.
整個人要加稅.
我們芝加哥出名的是什麼.
稅.
第二名出名的是什麼.
腐敗.
第三名出名的是什麼.
pop host.
徐佳麗.
加稅就是牙痛的意思.
但是對猶太人來說.
不僅加稅是牙痛的意思.
加稅也是對他們的identity.
一個很大的挑戰.
因為猶太人只會拜一個主.
一個神.
如果你強迫他們加稅的話.
就代表他們要對他們的身份.
say no.
要斷絕他們的身份.
要清洗他們.
所以每逢有加稅的時候.
當時的.
第一世紀歷史的記載.
就一定會有動員.
有好幾場大的動員.
都是在公元前和公元後的十幾二十年發生.
每逢加稅的時候.
社會政治的形態就會很緊張.
因為會有人出來反抗.
來做一些武力上的暴力.

$^{401}$所以說到報名上策的時候.
在當時的人來看.
是一種制度的暴力.
institutional violence.
你明白.
特別是今時今日.
你明白什麼叫做violence.
violence不只是搞爛東西.
塗鴉的東西.
去傷人.
或者破壞公物.
那些當然可以叫做violence.
但是有一種深層的violence.
不是你能夠看到的.
並且是很真實的.
那些叫做制度上的暴力.
就是說那個制度裡面.
去到壓迫你.
使你沒有氣頭.
那叫做institutional violence.
那些violence可能比那些看得見的violence會更加暴力.
在羅馬帝國裡面.
加稅會有報名上策的聲音.
這些就叫做institutional violence.
所以你看路加坊記載聖誕故事的時候.
你就知道.
他用了很多當時官媒用的字眼.
或者官媒用的ideology.
那些政府新聞台裡面經常用的字眼.
字句.
去形容主基督的國恆.
我們繼續看下去.
二章四至七節.
約瑟從加里尼的拿撒勒城.
上猶太省.
到了大衛的城.
叫做伯利恆.
因他本是大衛一族一家的.
要和他所聘之在馬尼亞一同報名上策.
那時馬尼亞的新人已經重了.

$^{441}$他們在那裡的時候.
馬尼亞的產期到了.
就心裡要投胎的意志用布包起來.
放在馬場裡.
因為黑店裡面沒有地方.
我先搞清楚裡面的內容.
約瑟和馬尼亞.
他們的本鄉在哪裡呢.
鄉下在哪裡呢.
伯利恆沒錯.
但他們身在哪裡呢.
就在拿撒勒.
拿撒勒就在北邊.
大概九十里英里以北.
他們在拿撒勒那裡.
大致上我們知道現在的背景.
大致上我們知道可能是.
可能是約瑟為了謀生.
而他要北遷上去.
北邊的拿撒勒那裡找工作.
因為我們知道上面那裡的省.
拿撒勒附近的嘉利的地方.
當時是大興圖木.
約瑟是做哪一行的呢.
木匠或者石匠的.
當時不分木與石.
總之匠就是匠.
木又丟石又丟.
所以他就上了去拿撒勒這個地方.
拿撒勒附近有一個所謂的省府.
聖經沒有提過.
不過在那裡正是這個時間.
他們大興圖木.
可能是這個緣故.
所以他們就上了去拿撒勒.
誰知道報名上冊.
一下了所謂聖旨的時候.
他就要回去他的本鄉.
來報名.
來登記.

$^{481}$當時猶太人就要逼他們回去他們自己的鄉下.
然後就要報所謂的人和稅.
所以他們兩夫婦.
瑪利亞已經從聖命.
懷孕了.
就要趕去巴里海這個地方.
這裡就是最棘手的地方.
我們通常因為預備團契的聚會.
我看了好幾齣聖誕電影.
因為上一版.
我們團契.
團契會.
他們有個電影週.
就看聖誕的電影.
我看了好幾齣都是同一個描述.
就是約瑟瑪利亞.
瑪利亞就坐著龍仔.
可能因為他懷孕了.
所以死無港.
他坐龍仔.
我猜可能應該是這樣.
你懷孕了.
你就叫Uber.
趕路去巴里海.
昨晚我還看了一齣.
Nativity Story.
那裡還說到.
約瑟帶著瑪利亞.
當時有一隻驢.
還下大雨.
然後瑪利亞穿水袋.
我想他穿水袋.
總之是作動求離要生.
約瑟就帶著瑪利亞.
拖著驢仔.
下大雨.
傾盆大雨.
到處去叫.
去找人.
有沒有地方.

$^{521}$產房有沒有床位.
到處去找地方.
找來找去都找不到.
結果就跑到山洞裡.
然後就生了出來.
通常都是有這樣的情景.
因為我們在說.
客店會有地方去.
客店會有地方.
No room for them in the inn.
我們想起inn的時候.
想起客店.
我們就一定想起holiday inn.
或者hotel 6等等.
原來報名上集的時候是這麼擠擁的.
所有的酒店房間.
床位都沒有了.
這裡有一點點棘手.
為什麼呢.
中英的翻譯.
都是給了我們西方現代人.
有這種錯覺.
我們覺得當時當然有很多hotel 6.
或者holiday inn.
他們check in找不到.
好像沒有vacancy等等.
其實路加用這個字.
是別有用心.
路加用這個字.
不是平常他們用的.
所謂hotel inn的酒店的字.
如果你想在路加方向去找.
真的有酒店這個字.
客店這個字.
你去到後面的一章聖經.
說到很撒亞的比喻.
那個比喻的故事.
那裡出現一個字.
那個字就真的是客店的意思.
你記得那個故事.

$^{561}$有一個人在坊間.
被人打了半死.
沒有人理他.
祭司也不理他.
平常我們期望他會理他的人都沒有理他.
然後有一個撒亞人.
猶太人最憎恨的種族的人.
他就經過就理了他.
包紮了.
然後送他去客店.
給錢他.
照顧他.
那個叫客店.
但這裡翻譯做客店.
其實又不是很對.
應該最好的翻譯應該是客廳.
或者英文.
英文某些的翻譯本.
就叫做guest room.
就是說.
他們去到巴黎人的時候.
英文也說清楚了.
他們在罵他們的時候.
不是突然間說到.
突然間下雨很慘.
然後又作動又穿水袋.
然後又什麼都沒有帶.
差不多打他的時候.
預備好了.
產旗到來.
就關了頭盔的耳墊.
布包放在馬場.
因為客廳.
是沒有地方的.
如果你這樣看.
那個畫面就完全不同了.
當時以色列猶太人住的房樓的格式.
我們知道大概是這樣.
兩層.
上面是住人的.

$^{601}$上面也可能會有所謂的guest room.
客廳在上面.
因為上面是睡人的.
你爬上去就可以睡人.
然後下面是做什麼的呢?.
下面是房樓的一樓.
就是你的所謂的place of productivity.
就是你家庭或者工業.
家庭工業的生產的地方.
因為下面是什麼呢?.
中間是kitchen.
然後左邊是一些storage的地方.
放雜物的地方.
然後可能另一邊.
就是放一些牛羊等等.
特別是牧羊的人.
或者耕田的人.
他們一定會有些牛羊.
你明白嗎?.
首先這些房樓.
不是你今天想的那些south of town house.
沒有這麼大間.
這些房樓大概百多呎而已.
幾十呎.
有兩層.
為什麼會這樣設計呢?.
你要知道.
人的建築是很實際的.
住著住著.
他們演變出來的建築.
是很practical的.
因為在下面放這些牛羊.
有什麼好處呢?.
就可以保溫.
你養狗養貓的.
睡覺有沒有跟貓狗睡?.
是不是很亂?.
你想想你抱著牛睡.
有多亂?.
你每晚都這樣.

$^{641}$放在下面其實是for heating.
下面也是煮飯的地方.
經常會生火等等.
所以上面睡就會亂過.
所以很practical.
按照當時巴勒斯坦的氣候.
來設計的房樓.
聖經說.
客廳會有地方.
或者樓上房間.
客房會有地方.
guest room會有地方.
於是他們就放在下面.
在下面馬槽的地方.
來生客.
所以你這樣的圖畫.
你就可以聯想到.
當時因為人口統產的問題.
被迫很多人回去巴黎.
那些同鄉巴拉人的時候.
他們很擠迫.
因為突然間回去.
人流多了.
就變成這些小的房間.
房子是沒有地方的.
很擠迫地他們就放在下面.
這個courtyard.
就是庭院的地方.
來生下耶穌.
如果你這樣看.
你就會發現羅家芬的故事.
不是說人們不要耶穌.
就是說人們其實是迎接耶穌.
客房就沒有地方.
不過讓他去下面馬槽.
無論你怎樣解釋.
都是一個很平凡的故事.
如果你說一個極主要講.
一個君王的講.
如果你說一個神的兒子出生.

$^{681}$沒理由會是這樣的境況.
羅馬帝王神的兒子出生的故事.
這些所謂的birth story.
nativity story.
如果對羅馬人來說.
這是一個家庭的表現.
怎會在一個這麼平凡的地方出生.
但如果你這樣看.
你就會發現聖誕的故事.
就是我們的家庭.
我們的心.
是不是有接耶穌的.
你要知道.
我知道中間有很多是.
花教會開不久.
可能形式信仰不是很深.
但我知道基督教信仰.
聖經信仰.
不是一個東西.
基督教信仰是一個關係.
Christianity is not a thing.
不是說我現在要了.
我現在不需要了.
我就丟了.
Christianity is not a thing.
Christianity is a relationship.
聖誕是一個關係.
是和上帝建立這個關係.
當你邀請主進入你的生命當中.
就是一個關係的開始.
並且這個關係是永遠的關係.
我們在Advent的時候.
我們可能會想到.
聖誕節有很多很美的氣氛.
大家知道聖誕節最重要的.
Advent是什麼?.
就是我們想.
就是上帝要光臨在你的生命當中.
也不是說一些美麗的氣氛是重點.
重點在於我們的生命.

$^{721}$我們願意在上帝面前去營造.
我們承認我們是脆弱的.
我們需要上帝的恩典.
這才是真正的Advent.
家庭接受耶穌.
正如我們需要去接受耶穌的恩典.
路加從頭到尾都告訴我們.
這位耶穌降生的.
死在十字架上.
復活的.
以至我們所傳的都是一位主人.
你有聽過什麼叫做頭尾呼應嗎?.
在文學上叫做頭尾呼應.
就是說開頭說的東西到最後他就作狂犯.
這首《使徒行傳》都是路加所寫的.
第一集第二集.
開頭聖誕的故事.
一二三四章.
講到耶穌的起頭的時候.
都是講到他是主.
他來到建立一個的王國.
到最後《使徒行傳》最後那兩節.
保羅為了除風音已經到了羅馬.
他在那裡有兩年的時間.
他做什麼呢?.
他就是傳講上帝的王國.
傳講主基督.
主耶穌基督.
Lord.
我邀請這位主進入我們的生命當中.
就是以他為我們真正的主.
有一首聖誕的詩歌我很喜歡.
小孩子才會唱.
大人很難唱到.
他說什麼呢?.
What can I give?.
到底這位主耶穌來到我的生命當中.
我可以給些什麼?.
如果我是一個shepherd.
我可能給一隻羊.

$^{761}$如果我是一個博士.
我可能會獻上最好的東西.
但到底我應該給些什麼?.
I will give him my heart.
就是我的生命的主權.
來交給他.
這才叫做聖誕.
Ever.
最近在Facebook有個人.
發了這篇文章.
我覺得很有意思.
12月的時候我們教小孩子做什麼呢?.
聖誕節的時候.
就是教他們每一天做一個所謂的.
Act of kindness.
就是做一些事情.
來服侍身邊的人.
一些很小的事情.
其實我們.
信了主的.
到了12月.
是不是我們需要到12月?.
有些什麼小的事情我們可以做.
是能夠祝福到身邊的人.
奉獻.
而不是我買東西給我自己.
奉獻去幫助別人.
幫助上帝的福音的事局.
去關心一個朋友.
為他做一些祈禱.
人在聖誕節很孤單的.
特別是在中間.
父母未必跟你在一起.
家庭未必跟你在一起.
人不可以很孤單.
在聖誕節裡面我們不妨多做一些.
Open.
你自己的心靈.
今年年初的時候.
英國BBC的新聞台.

$^{801}$他們有一個報告.
他說日本政府發現了一個社會的問題.
就是老人家.
經常犯法.
犯法他就.
要入獄.
為什麼呢?.
因為老人家面對很多社會的問題.
就是退休金不夠.
又負債.
又不想成為兒女的負擔.
他們想到一個辦法.
不如偷偷東西.
然後犯法.
坐牢.
因為監獄裡面有什麼?.
有三餐.
因為這樣的社會.
轉變的問題.
日本的監獄都開始改變了.
加扶手.
因為跌倒.
就成為一個很大的循環.
既然這麼好.
不如進去吧.
有一個老人家偷了我的單車.
他直接去警察局.
他說這單車是我偷的.
你抓我吧.
他說監獄裡面.
原來日本的監獄有些監獄.
有卡OK.
所以他寧願坐牢.
但其實最深層的問題是什麼?.
最深層的問題.
為什麼不要自由.
而進去坐牢?.
是因為孤單.
孤獨.
如果我們有神童座.

$^{841}$我們就明白.
原來我們不是孤單.
與未必有孤單.
與人不孤單.
這就是基督教信仰的精髓地方.
就是關係.
一個人的關係.
在十二月裡面.
Evangelist快到了.
我不知道你可以主宰.
我也不知道你生命裡面.
願不願意接受這個邀請.
就讓主耶穌進入你的生命當中.
與你建立關係.
我們同心一致.
就這樣.
上主我們.
無論正國社會世界發生什麼事.
你願意但守.
願意進入我們的生命當中.
主讓我們.
邀請你進入我們的生命當中.
來到我們的家庭.
來到我們的領受.
你自己親自的一個座位.
我們將我們的主權.
再次移交給你.
讓你來管理我們的生命.
以致我們的生命能夠成為.
其他人的祝福.
主我們教會.
我們不只是迫迫求求.
我們也做很多不同的事.
讓我們的主權交給你.
讓上帝認識我們.
祝福他和我們.
使人能夠與你建立關係.
我們這樣祈禱奉祝.
祝福各位.
我們繼續.

$^{881}$拜拜!.
\newpage



\section{馬太福音 1:1-17}
\label{sec:99dN28SUzkQ}
\textbf{12/15/19 主日講道: 與人同在的神 7, 馬太福音1\_1-17 (陳崇基牧師)}
\newline
\newline
連結: \href{https://youtube.com/watch?v=99dN28SUzkQ}{\texttt{https://youtube.com/watch?v=99dN28SUzkQ}} ~~~~ 語音日期: 2019-12-20
\newline
\newline
\hyperref[sec:zYSzdaPaiDI]{\small{< < < PREV SERMON < < <}}
~
\hyperref[sec:index_chronic]{\small{[返順時目]}}
~
\hyperref[sec:index_scriptual]{\small{[返順卷目]}}
~
\hyperref[sec:4ZoSLfbOWng]{\small{> > > NEXT SERMON > > >}}
\newline
\newline
馬太福音 1:1-17
\newline
\begin{longtable}{cl}
\hline
\hline
章節 & 經文 (和合本修訂版)\\
\hline
1:1 & \begin{tabularx}{0.7\textwidth}{X} 亞伯拉罕的後裔、大衛的子孫耶穌基督的家譜： \end{tabularx} \\ \\ \relax
1:2 & \begin{tabularx}{0.7\textwidth}{X} 亞伯拉罕生以撒，以撒生雅各，雅各生猶大和他的兄弟， \end{tabularx} \\ \\ \relax
1:3 & \begin{tabularx}{0.7\textwidth}{X} 猶大從她瑪氏生法勒斯和謝拉，法勒斯生希斯崙，希斯崙生亞蘭， \end{tabularx} \\ \\ \relax
1:4 & \begin{tabularx}{0.7\textwidth}{X} 亞蘭生亞米拿達，亞米拿達生拿順，拿順生撒門， \end{tabularx} \\ \\ \relax
1:5 & \begin{tabularx}{0.7\textwidth}{X} 撒門從喇合氏生波阿斯，波阿斯從路得氏生俄備得，俄備得生耶西， \end{tabularx} \\ \\ \relax
1:6 & \begin{tabularx}{0.7\textwidth}{X} 耶西生大衛王。大衛從烏利亞的妻子生所羅門， \end{tabularx} \\ \\ \relax
1:7 & \begin{tabularx}{0.7\textwidth}{X} 所羅門生羅波安，羅波安生亞比雅，亞比雅生亞撒， \end{tabularx} \\ \\ \relax
1:8 & \begin{tabularx}{0.7\textwidth}{X} 亞撒生約沙法，約沙法生約蘭，約蘭生烏西雅， \end{tabularx} \\ \\ \relax
1:9 & \begin{tabularx}{0.7\textwidth}{X} 烏西雅生約坦，約坦生亞哈斯，亞哈斯生希西家， \end{tabularx} \\ \\ \relax
1:10 & \begin{tabularx}{0.7\textwidth}{X} 希西家生瑪拿西，瑪拿西生亞們，亞們生約西亞， \end{tabularx} \\ \\ \relax
1:11 & \begin{tabularx}{0.7\textwidth}{X} 百姓被遷到巴比倫的時候，約西亞生耶哥尼雅和他的兄弟。 \end{tabularx} \\ \\ \relax
1:12 & \begin{tabularx}{0.7\textwidth}{X} 遷到巴比倫之後，耶哥尼雅生撒拉鐵，撒拉鐵生所羅巴伯， \end{tabularx} \\ \\ \relax
1:13 & \begin{tabularx}{0.7\textwidth}{X} 所羅巴伯生亞比玉，亞比玉生以利亞敬，以利亞敬生亞所， \end{tabularx} \\ \\ \relax
1:14 & \begin{tabularx}{0.7\textwidth}{X} 亞所生撒督，撒督生亞金，亞金生以律， \end{tabularx} \\ \\ \relax
1:15 & \begin{tabularx}{0.7\textwidth}{X} 以律生以利亞撒，以利亞撒生馬但，馬但生雅各， \end{tabularx} \\ \\ \relax
1:16 & \begin{tabularx}{0.7\textwidth}{X} 雅各生約瑟，就是馬利亞的丈夫；那稱為基督的耶穌是從馬利亞生的。 \end{tabularx} \\ \\ \relax
1:17 & \begin{tabularx}{0.7\textwidth}{X} 這樣，從亞伯拉罕到大衛共有十四代，從大衛到遷至巴比倫的時候也有十四代，從遷至巴比倫的時候到基督又有十四代。 \end{tabularx} \\ \\
[1ex]
\hline
\hline
\end{longtable}
$^{1}$昨晚給我們一個很好的分享會.
昨晚我聽到很多他們的見證.
我很感動.
因為我見到分享的人不是一些很完美的人.
很多破碎的地方,很多軟弱的地方.
或者可能這種分享,這種生命才是最讓我們能夠認同到.
我們很多時候有一種錯覺.
我們覺得上來分享.
屬於主的人,那些應該是一些完美的人.
我們有一種所謂德性的主義.
就是Triumphanism.
如果那個人,他的信主的經歷.
是一個很完美的經歷.
很Prosperous的,很棒的.
我們覺得這些經歷才能夠擺上桌子.
但這種經歷,偏偏不是我們大眾能夠認同.
我們能夠認同的,就是那些名字,那些生命.
是和我們一樣脆弱.
和我們一樣軟弱.
所以我們今天想藉著他們兩個的分享.
藉著我們這個weekend所做的分享會.
我們都會結束我們這個series.
我們這個series叫做「與人同在的神」.
我們由創世記第一章開始.
我們講講上帝整個的story.
上帝整個的大故事.
到底祂是從何而來.
就是拯救我們這班軟弱脆弱的罪人.
破碎的生命.
祂將我們拯救出來,和我們同行.
由創造開始,直至整個的舊約的故事.
到耶穌基督親自來到聖誕節.
去想耶穌基督的誕生.
到祂在十字架上為我們聚死.
在墳墓裡為我們復活,被我們盼望.
整個的grand story,到底我們自己.
我們這條命,我們這個生命.
如何fit入上帝自己的grand story.
上帝仍然在拯救人.
仍然去和罪人同在,同行.

$^{41}$這個grand story到底我們自己如何fit入.
第一次進來教會.
你是完全不認識耶穌基督這回事.
你是在摸索中.
這個grand story就是告訴我們.
原來在整個世界上帶領著的時候.
我們自己的人生可以如何fit入.
所以今天我想結束的時候.
我們看一段很長的經文.
這段經文是耶穌基督的家譜.
新約裡面第一本書.
叫做《馬太福音》.
他開始介紹耶穌的時候.
就介紹耶穌的家人.
耶穌的先祖.
之前的那一代人.
那一條line是怎樣的.
由阿班罕到大衛.
到耶穌的降生.
耶穌的家譜就在這裡.
當然我不會讀.
就算今天沒有時間也好.
我也不會讀.
我想給大家一個大的picture.
講到家譜,我們華人就有一點點能夠認同.
因為我們中間有些人回過鄉下.
去祠堂看祖譜.
有沒有誰試過.
有吧,找回自己的名字.
哪個祖宗等等.
家譜就是有這樣的意思.
通常猶太人的家譜.
就是將最顯赫的人放上去.
譬如阿班是民族英雄.
所謂安國英雄等等.
這些人人都會認識的.
這些名字,這些生命.
就是去介紹一個很重要的人.
你想想,如果你介紹耶穌基督的時候.
你會用什麼人去介紹祂.

$^{81}$在馬太福音這裡.
他介紹耶穌的降生的時候.
他就用了這一堆的生命.
這一堆的聖靈.
很簡單地告訴你.
這裡的大picture是怎樣的.
馬太福音就將整個舊約的故事.
整個grand story就分成三段.
第一段就是阿班是民族英雄.
直到大衛那個民族的最大君王.
最有名的君王.
那裡有十四代.
然後大衛這個君王.
再到他們舊約裡面最重要的一個災難.
就是他們流亡巴比奴.
那裡又有十四代.
然後流亡巴比奴他們回來.
回到故土的時候.
直到耶穌的降生.
有十四代.
如果你是熟悉聖經文化.
七是很重要的.
因為七是完美的number.
一個星期有七天.
所以用了七天來創造這個世界.
所以馬太似乎告訴我們.
十四十四十四.
就有多少個七.
六個七.
六個七沒錯.
這個question不是很trick的.
不用怕.
我們很簡單.
六個七.
所以耶穌基督來的時候就帶來多少.
第幾個七.
第七個七.
當你在猶太人的民族裡面.
大衛這個名字.
David.

$^{121}$不要響音就是DVD.
算希伯來的數目.
就是十四.
所以十四也是一個iron number.
他就告訴我們.
當我們水深火熱的裡面.
他們流亡.
回到故土的時候.
神的子民還在等救主.
很慘的等救主.
被人欺負.
生命裡面很破碎.
很黑暗的神.
所以第七個七.
是誰帶來的.
馬太就是這樣介紹耶穌基督救主的出現.
但是他做了一樣東西.
是很特別的.
因為你知道.
猶太人的族譜.
或者甚至是我們華人的族譜.
是不會有女人的命的.
因為以男為上.
很多時候女生外向.
女生是嫁走的.
所以只會放男人的命.
但是馬太.
在這裡介紹耶穌的時候.
就特意放了四個女人的命.
不單止放了四個女人的命.
這四個女人.
在猶太人的歷史當中.
不是出名的.
說起四個女人.
我們中國人有.
西大美人.
記不記得.
猶太人也有西大美人.
猶太人西大美人是誰呢.
受害.

$^{161}$利伯加.
威伯卡.
利亞和拉傑.
四個都是創世記裡面.
阿拉罕爾薩雅國的太太.
他們在兩個中間.
是讚口不絕的.
從頭讚到尾.
他們的外表.
他們的德行.
讚到他們不得了.
如果你要選四個女人的話.
我們如果和馬太一起去看.
寫耶穌的家譜的時候.
為什麼不自己寫呢.
他就偏選.
另外那四個.
而這四個呢.
都是有共通點的.
就是他們在道德上.
都是破相的.
你回去聽聽.
看一下創世記怎麼說他媽.
約書亞記怎麼說拉合.
在聖經裡面舊約新約.
父親提到拉合的時候.
大部分時間都說是妓女拉合.
她是一個妓女.
她歸順了耶和華.
她仍然都是一個叫妓女.
因為每個人都給她一個名字.
就是妓女.
是一個道德上不能夠站高的.
人來的.
路德可能好一點.
接著就是烏里亞的妻子.
佛事伯.
有一個共通點.
就是四個都是令我們不舒服的名字.
為什麼你要介紹一個這麼重要的人.

$^{201}$你要用這四個人來介紹.
有更加好的人為什麼你不選.
我不知道你去一些分享會的時候.
你會看一下那些貼文.
你會想想.
那個誰分享呢.
這個是誰呢.
這個藝人是誰呢.
不是出名的.
我不去了.
你總會是這樣.
如果你和馬太一起寫耶穌的家譜的時候.
不如放那四個四大美人來吧.
馬太說.
我們不放.
為什麼要選這四個呢.
四個都是道德上有婆娑.
四個都是外邦人.
這個告訴我們耶穌基督來到的時候.
是要恩典裡面拯救我們.
不是因為我們是天.
我們很多時候在教會.
我在教會很久很多年.
我經常聽到有個外國人在外面.
和別人談教會的事情.
有些人不上教會.
他說我不是那麼好.
我不想上教會.
他就當教會是一個聖人的博物館.
我們就是那些蠟像.
特別是懷舊的.
我們不是很動的.
好像每個人都是蠟像.
每個人都是那些英雄.
每個人都是很棒的人.
錯不是.
有一個朋友跟我說.
我不是那麼好的.
我不去教會.
我說千萬不要這樣說.

$^{241}$我們每個人都是那麼壞的.
包括我在內.
可能我比你還壞.
我們都是因為上帝的恩典.
我就跟他說.
你看清楚.
聖經不是叫我們做好人.
他就突然間說.
不是的,聖經不是叫我們做好人.
聖經告訴我們我們做不到好人.
我們是罪人.
我們是破碎的人.
我們是依靠上帝的恩典.
所以馬太特意將這四個所謂罪人.
放上台.
叫他們來這裡顯出.
原來不是因為我們.
而是上帝的恩典.
這也是馬太福音開始的時候.
他介紹的.
用中文就譯為耶穌基督的家譜.
其實不是很能譯出那個意思.
其實應該說.
The book of Genesis in Jesus Christ.
就是說在基督耶穌的恩典裡面.
有一個創世記.
拜這個時間的時候.
大概這個時間兩小時.
猶太人已經說.
創世記就是九月裡面第五書.
他們會叫他做什麼呢.
叫做Genesis.
用希臘文來叫做Genesis.
馬太就說.
在基督裡面就有Genesis.
就是我們生命裡面能夠被重整.
我們破碎的生命.
擺在上帝的面前的時候.
在他的恩典當中.
我們就被重新創造.

$^{281}$我們就是New Creation.
我們生命裡面.
總有一些東西.
是不見得光的.
你試試把所有東西都倒出來.
可能連自己都接受不了.
但是這位上帝就能夠接受.
這位上帝在恩典裡賜給我們.
我想結束的時候.
用一個祈禱來結束.
這是我生命的主結.
你沒有因為我的罪.
沒有因為我過去而遠離我.
相反.
因為你為我的罪死在上帝的手上.
生命我的罪.
從死裡復活.
賜給我新生命新的開始.
New Genesis.
我願意你成為我生命的主.
但你過一世.
我不知為何.
今晚如果你沒有來的話.
你聽到這些故事.
你可以聽到再聽.
你有來的話.
你應該很清楚.
我們人就是一個這麼破碎的東西.
不是因為我們可以.
而是耶穌基督的恩典.
我請詩琴上來.
彈一彈.
我們做一個結束祈禱.
原來聖經告訴我們.
聖誕節不是和耶穌慶祝生日.
聖誕節是慶祝我們的生日.
是耶穌臨到我們的生命.
給我們恩典.
讓我們能夠重生.
這是Spiritual Rebirth.

$^{321}$New Genesis.
新的創造.
聖誕節不是我們和耶穌唱Happy Birth.
而是祂的生命來到我們的時候.
我們有一個New Birth.
在恩典裡面能夠有新的生日.
去和聖誕節.
在心裡面如果你有同感的話.
和我一起來祈禱.
生命的主.
主,你沒有因為我的罪.
沒有因為我的過去而遠離我.
相反.
你為我的罪死在神之間.
赦免我.
從死裡復活.
你賜給我新的生命.
我願意你成為我生命的主.
從今.
你去帶領我一生.
奉主耶穌基督的名字而鑑.
Amen.
\newpage



\section{馬太福音 2:13-23}
\label{sec:4ZoSLfbOWng}
\textbf{12/29/19 主日講道: 聖誕故事 - 馬太福音2\_13-23 (陳崇基牧師)}
\newline
\newline
連結: \href{https://youtube.com/watch?v=4ZoSLfbOWng}{\texttt{https://youtube.com/watch?v=4ZoSLfbOWng}} ~~~~ 語音日期: 2020-01-05
\newline
\newline
\hyperref[sec:99dN28SUzkQ]{\small{< < < PREV SERMON < < <}}
~
\hyperref[sec:index_chronic]{\small{[返順時目]}}
~
\hyperref[sec:index_scriptual]{\small{[返順卷目]}}
~
\hyperref[sec:HLQkwHUhsgc]{\small{> > > NEXT SERMON > > >}}
\newline
\newline
馬太福音 2:13-23
\newline
\begin{longtable}{cl}
\hline
\hline
章節 & 經文 (和合本修訂版)\\
\hline
2:13 & \begin{tabularx}{0.7\textwidth}{X} 他們走後，忽然主的使者在約瑟夢中向他顯現，說：「起來！帶著小孩子和他母親逃往埃及，住在那裡，等我的指示；因為希律要搜尋那小孩子來殺害他。」 \end{tabularx} \\ \\ \relax
2:14 & \begin{tabularx}{0.7\textwidth}{X} 約瑟就起來，連夜帶著小孩子和他母親往埃及去， \end{tabularx} \\ \\ \relax
2:15 & \begin{tabularx}{0.7\textwidth}{X} 住在那裡，直到希律死了。這是要應驗主藉先知所說的話：「我從埃及召我的兒子出來。」 \end{tabularx} \\ \\ \relax
2:16 & \begin{tabularx}{0.7\textwidth}{X} 希律見自己被博學之士愚弄，極其憤怒，差人將伯利恆城裡和四境所有的男孩，根據他向博學之士仔細查問到的時間，凡兩歲以內的，都殺盡了。 \end{tabularx} \\ \\ \relax
2:17 & \begin{tabularx}{0.7\textwidth}{X} 這就應驗了耶利米先知所說的話： \end{tabularx} \\ \\ \relax
2:18 & \begin{tabularx}{0.7\textwidth}{X} 「在拉瑪聽見號咷大哭的聲音，是拉結哭她兒女；她不肯受安慰，因為他們都不在了。」 \end{tabularx} \\ \\ \relax
2:19 & \begin{tabularx}{0.7\textwidth}{X} 希律死了以後，在埃及，忽然主的使者在約瑟夢中向他顯現， \end{tabularx} \\ \\ \relax
2:20 & \begin{tabularx}{0.7\textwidth}{X} 說：「起來，帶著小孩子和他母親回以色列地去！因為要殺害這小孩子的人已經死了。」 \end{tabularx} \\ \\ \relax
2:21 & \begin{tabularx}{0.7\textwidth}{X} 約瑟就起來，帶著小孩子和他母親進入以色列地去。 \end{tabularx} \\ \\ \relax
2:22 & \begin{tabularx}{0.7\textwidth}{X} 但是他因聽見亞基老繼承他父親希律作了猶太王，怕到那裡去；又在夢中得到主的指示，就往加利利境內去了。 \end{tabularx} \\ \\ \relax
2:23 & \begin{tabularx}{0.7\textwidth}{X} 他們到了一座城，名叫拿撒勒，就住在那裡。這是要應驗先知所說的話：「他將稱為拿撒勒人。」 \end{tabularx} \\ \\
[1ex]
\hline
\hline
\end{longtable}
$^{1}$我們同心得道祈禱.
讓我們能夠有主耶穌的生命進入我們的生命當中.
主耶穌讓我們繼續歡迎和迎接你.
來到我們生命當中作主和作王.
你引導我們.
今天我們是二零一九年最後一個的主日.
你不斷地去到保守.
去到給我們有足夠的恩典.
主我們求你讓我們不是單單做一群受恩的人.
讓我們懂得在你自己給我們的角色.
在我們的崗位當中.
我們去盡心盡力的去見證去服侍你.
今天我們的會友大會.
我們求你與我們同在.
讓我們去仰望你在來年.
你的恩典足夠我們使用.
讓我們同心一致去建立我們自己的家.
同樣地我們將到我們身邊的弟兄姊妹來交托.
特別是一些有重病.
或者需要照顧重病的家人.
主我們求你讓他們心裡能夠有安穩.
在基督裡面的平安.
賜給他們力量.
賜給他們需要及時的幫助.
主我們求你讓我們這個家.
能夠繼續有主你給我們的愛在當中.
我們學習彼此的扶持.
彼此的鼓勵.
去建立你這個屬靈的家.
讓你在我們中間能夠得著榮耀.
我們這樣祈禱奉主名及.
今天是聖誕節過後的一個主日.
我們想繼續看看聖誕節的故事.
我們經常停留在十月二十五日.
降生了.
什麼事情都沒有發生.
聖經裡面有記載.
耶穌基督降生之後.
到底發生了什麼事情.
特別是馬太福音第二章.

$^{41}$如果我們有聖經的話.
請你們翻開去.
馬太福音第二章.
第二章第十三節開始.
這裡剛剛說完什麼呢?.
就是那個博士的故事了.
就是東方三個博士.
來朝拜新生王耶穌基督.
然後他們就離開了.
這裡說到他們就是那幾個博士了.
十三節.
他們去後有主的使者.
向約瑟夢中顯然說.
起來帶著小孩子和他母親.
逃往埃及.
住在那裡.
等我吩咐你.
因為希律必尋找小孩子.
要除滅他.
約瑟就起來.
夜間帶著小孩子和他的母親.
往埃及去.
住在那裡直到希律死了.
這是要應驗主席先知所說的話說.
我從埃及召出我的兒子來.
希律見自己被博士愚弄.
就大大發怒.
差人往百里行城裡.
並四警所有的男孩.
照著他向博士仔細查問的時間.
凡兩歲以來的都殺盡了.
這是應驗先知耶利米的話說.
在拉瑪聽見號嘯大哭的聲音.
是拉傑哭他的兒女.
不肯受安慰.
因為他們都不在了.
希律死了以後.
有主的使者在埃及向約瑟夢中顯現說.
起來帶著小孩子和他的母親.
往以色列地去.

$^{81}$因為要害小孩子性命的人已經死了.
約瑟就起來.
把小孩子和他母親帶到以色列地去.
只因聽見阿基努接著他父親希律.
作了猶太王.
就怕往那裡去.
又在夢中被主指示.
要往加利利境內去了.
到了一座城名叫拿撒勒.
就住在那裡.
這是應驗先知所說.
他將稱為拿撒勒人的王.
我們看一段聖經.
主要講的是大希律.
他去到撒鹿.
在伯利恆裡面.
兩歲以下的男嬰.
這個故事就是 Massacre.
的故事就是剛剛在所謂基災.
耶穌基督降生之後.
我們知道那幾個東方博士來朝拜耶穌的時候.
當時耶穌大概應該是兩歲左右.
如果了解那個故事的小情節就知道.
大概當時已經是兩歲左右.
今天你看到一些.
我們教會那些.
一歲兩歲那些已經跳來跳去走了.
那時候耶穌應該是差不多長大的時候.
不是三個博士.
幾個博士.
幾個博士來到朝拜耶穌之後.
大希律就聽到新生王.
覺得對他的政權有一個很大的威脅.
這個其實有一個歷史的.
所謂大的圖畫.
為什麼是東方來的博士呢?.
因為當時的東方.
所謂米索不達米亞這個地方.
就是今天的伊拉克等等的地方.
還沒有受到羅馬帝國的直接的統治.

$^{121}$在那裡東方的人.
是不是很承認羅馬帝國的政權的.
有一段時間大概是主耶穌降生的時候.
那時候他們很嚮往.
有一個新的王來到中間.
我們不只是在說猶太王.
只是在說普通的近東的民族.
他們都覺得有一個新的王來取締奧古斯都.
所以是一個比較是.
政治上比較不是很平安.
不是很平穩的一段時間.
東方幾個博士來到的時候.
其實他們就是帶著東方的人.
他們近東的人.
他們那種 expectation(期望).
他們希望有一個新的王.
而這個新的王是他本地的王.
是他本地的近東的王.
所以他們來到去.
敬拜新生王耶穌基督的時候.
大希律就是羅馬帝國任命去做所謂特首的.
他就是管住猶太省這個地方.
所以他聽到有一個新生王出來的時候.
他就覺得很大衛脅.
很多時候看完聖經的時候.
我們就會想起一個歷史的問題.
就是因為經外的文獻.
從來沒有提過這件事的.
沒有提過大希律.
去到殺戮伯利恆裡面兩歲以下的男嬰.
有說大希律很多其他的事蹟.
特別是他做的壞事.
譬如第一世紀有一個歷史學家.
叫做Josephus.
他寫了很多有關Hero the Great.
Historical background(歷史背景)的東西.
但是沒有提到這件事.
所以我們很多時候有一個這樣的.
Historical question(歷史問題).
就是到底這件事有沒有發生.

$^{161}$只是聖經記載夠不夠壓力呢?.
其實呢.
你要知道當時的文獻.
或者當時的新聞.
去記載一些軍王做過的事情.
很多時候這些這麼小的事情.
他們是不會記載的.
伯利恆呢.
只不過是當時一個很小很小的小鎮.
大概也是幾百人左右而已.
最多只有幾百人的一個鄉村的地方.
所以要說去殺戮這些男嬰.
可能只有十幾二十人左右的數目.
所以當時的經外的文獻.
大概應該也不會怎麼記載.
不過呢.
正如我所說.
有很多人記載大希律的事蹟的時候.
就吻合這個故事裡面所說的大希律的性格.
因為Hero the Great.
我們知道是一個很殘酷的暴君.
他又殺了他自己的老婆.
因為怕他老婆奪他的位置.
又殺了他所有的兒子.
又怕他的兒子來搶他的政權.
也殺了大祭司.
當時的大祭司就是猶太地的第二把交椅.
要殺他為什麼呢?.
因為又怕他奪取他的位置.
所以呢.
你說到大希律的時候.
就說到一個是暴君一個是昏君.
所以呢.
這些數據告訴我們.
大希律也是一個這麼殘酷的人.
所以有一個這樣的故事.
其實是很寫實的.
雖然我們只有 circumstantial evidence(事實).
這些所謂環境上的證供.
但是你看這個故事的時候.

$^{201}$我相信聖經特別要去記載.
耶穌降生之後.
還有一些這樣的下文.
有他的意思.
今天我們就來看看這個故事到底是什麼意思.
在教會的傳統裡面.
聖誕節過了.
四天到五天.
有一個特別的日子.
大公教會裡面.
所謂 Catholic Church.
裡面有一個這樣的傳統.
就是25號之後.
四到五天.
就是28號昨天.
或者29號.
視乎你是西邊的教會還是東邊的教會.
28,29號就是這兩天.
就是一個叫做 Feast of Holy Innocence 的日子.
這個日子就是去紀念.
那些所謂 saints(聖徒).
就是在上帝的國度裡面有被紀念的人.
那麼.
今天.
或者昨天.
28,29號.
就是所謂這個 Feast of Holy Innocence.
或者不是 Christmas(聖誕節).
是 Childerness(兒童節).
就是記載什麼呢?.
就是紀念那些百里行的兩歲以下.
被殺戮的這些嬰孩.
那麼這個日子.
就是特別去講到這一段經文所記載的.
正如我所說.
剛才我們看那段聖經.
聖經裡面的記載.
當時耶穌大概是兩歲多.
然後主去保守他.
然後他就離開了百里行.

$^{241}$逃難到了埃及這個地方.
那麼.
就逃往埃及.
然後主的使者就帶他回到以色列.
這個就叫做 Feast of Holy Innocence.
去紀念這些被殺害的無辜的.
可以這麼說的殉道者.
如果你有去過西班牙的話.
有誰去過西班牙玩過呢?.
如果你這兩天去過西班牙玩的話.
他們就會紀念這一天.
那麼他們的方法.
就好像愚人節一樣.
就是小孩子就將家裡的糖換成鹽.
那麼媽媽起床做早餐的時候就放錯了.
那麼放糖就變成放鹽.
或者趁著爸爸媽媽睡覺的時候.
小孩子就會把時間調到另一個時間.
爸爸媽媽起床的時候就會覺得晚了.
晚了就上班等等.
愚人節就好像 April Fools 一樣.
那麼我不知道為什麼會有這樣的傳統.
大概我們也不是很清楚為什麼.
會不會是因為這個故事裡面.
是講到幾個博士.
是愚弄暴君.
會不會是因為這樣的緣故.
所以變成了 April Fools.
就是西班牙的 April Fools.
無論如何.
這個節日.
讓我們想到這個經文裡面所講的一些事蹟.
更加讓我們去想一下.
到底這個經文.
到底想講些什麼.
熟悉舊約聖經的人.
一聽到剛剛所讀的經文的時候.
他就會自然聯想起一些.
在舊約歷史裡面.
一些很重要的所謂 Grand Narratives.

$^{281}$一些大的故事.
就是每一個民族都有他自己的一些回憶.
一些大的日子或者大的故事.
在舊約裡面最大的故事是什麼呢?.
當然就是 Exodus.
是不是?.
就是出埃及.
你記得出埃及記.
裡面一些的事蹟.
大概就好像殺六百里行這個故事的事蹟差不多.
你想想出埃及記頭幾章.
發生什麼事?.
就是有暴君出現是不是?.
那個就是法老.
然後.
有殺男嬰的故事.
就是法老想控制喜伯來人的男嬰.
將他們淹死.
然後呢.
也有一個主上帝的歷史裡面的.
掌管.
然後就留下了某某.
他就成為一個帶他的子民出埃及的所謂英雄.
那個就是摩西.
所以現在聽到剛剛那個殺六百里行男嬰的故事的時候.
自然屬於舊約的人就會一定想起這件事情.
就是他們想起似曾相識.
曾經法老暴君.
想去殲滅上帝的子民.
但是有神的子民出埃及記.
就是兩個修生婆.
去到逆權.
或者去愚弄法老.
就是騙了法老.
上帝最後打擊法老.
法老死了.
然後這個暴君是失敗的.
那個民族的英雄.
摩西.
還有以色列人.

$^{321}$整個國民.
就安全地出埃及了.
所以大希律在這個故事裡面.
去殺六百里行的男嬰.
暴君殺百里行的男嬰.
這個就是新約的Exodus.
或者應該是最大的Exodus.
就是一個暴君.
大希律.
他想殺這個男嬰.
但是卻被愚弄.
然後.
耶穌.
或者他所代表的上帝的子民.
能夠平安地出埃及.
很明顯這個故事.
想說的.
就是跟政權有關.
跟暴君有關.
跟暴政有關.
暴政.
是上帝會愚弄的.
上帝可以用他的方法來抗衡或者去愚弄.
這個其實就是聖誕節.
故事的味道.
我記得我講過兩個星期前開始講聖誕故事的時候.
我說聖誕故事不是聖誕卡小朋友的故事.
只不過我們有很多聖誕節的tradition.
燈飾禮物等等.
我們覺得只不過是一個children's story.
所以聖誕節的時候一定有children's choir的.
是不是?.
這個是好的.
是不是?.
就是扮成小朋友.
然後他們唱歌唱得不好聽也好.
我們也拍一頓夢.
我們.
這個是好的.
不過.

$^{361}$你要知道.
Christmas story不是children's story.
It's a political story.
有人說去重述.
重述.
聖誕的故事.
就是一個.
抗衡暴政的方法.
We're telling the Christmas story.
It's an act.
of resistance.
就是有聖誕就代表.
一些極權.
一些暴政.
是不會坐得穩的.
所以你要知道.
為什麼有很多極權.
共產的國家.
是不准人.
來慶祝聖誕節.
為什麼?不是因為不喜歡燈飾.
不是因為不喜歡禮物.
不喜歡聖誕老人.
而是他知道這個故事.
是有一個.
政治上.
政權上.
一個顛覆的.
味道在裡面.
所以為什麼《馬太福音》.
不只是說到.
耶穌降生這麼簡單.
而降生之後會令到當時的暴政.
當時的暴權.
有些什麼的影響.
你再細心看看這個故事.
除了上帝愚弄.
那個暴政之外.
也都說到什麼呢?.
也都說到上帝去拯救.

$^{401}$他自己的子民.
去擦去他們的眼淚.
二章十八節.
在拉瑪聽見.
號嘯大哭的聲音.
就是拉傑.
哭他的兒女.
不肯受安慰.
因為他們都不在了.
這裡要搞清楚.
因為我們.
不太熟悉舊約聖經的時候.
我們有幾樣東西不太知道.
拉瑪是什麼?.
拉瑪不是一個人.
拉瑪是一個地方.
羅馬.
拉瑪在哪裡呢?.
拉瑪在伯利恆.
耶路撒冷附近的一個鄉鎮.
在舊約的回憶裡面.
集體回憶裡面.
猶太人會想起.
拉瑪會想起什麼呢?.
就好像想起二次大戰的時候.
在德國那些集中營的地方一樣.
因為呢.
當猶大國.
最後覆亡的時候.
是主前586年的時候.
巴比倫軍.
攻入了京城.
耶路撒冷的時候.
就將這些猶太人.
俘虜到去.
帶到北面.
拉瑪這個地方.
有很多是被殺.
也有很多是被俘虜.
流亡到巴比倫的.

$^{441}$所以在拉瑪這個地方.
是最慘痛的時候.
是一個大哭.
大大哭號的一個地方.
猶太人想起拉瑪羅馬.
就想起他們亡國的時候.
他們流的眼淚.
有很多人死了.
有很多人被殺.
這裡是引用耶利米先知所說的話.
耶利米呢.
就說到拉瑪的情況.
因為耶利米是流亡的.
就是被擄之間的.
做先知的一個人.
他說到拉瑪的情況的時候.
他就將拉傑也拉進來了.
拉傑是什麼呢?.
拉傑就不是一個地方.
拉傑是一個人.
拉傑是誰呢?.
英文是Rachel.
你記得了.
在創世記裡面.
以色列國之父.
雅各.
Jacob.
God of Jacob.
Generation of Jacob.
雅各最喜愛的老婆.
就是拉傑.
而拉傑生最後最後胎的時候.
就是卞雅敏.
就是第十二個兒子.
卞雅敏生Benjamin的時候.
在哪裡生的呢?.
就是在拉瑪附近.
如果你翻查一下舊約裡面的.
記載一些事跡的時候.
你就知道拉傑呢.

$^{481}$的墳墓.
他生最後一胎的時候.
是他生了之後就死的.
他就是葬在.
伯利恆.
拉瑪附近的一個地帶.
所以耶利米說到拉瑪.
流亡慘痛的經歷的時候.
他就聯想到.
拉傑.
去生他的兒女的時候.
他要死.
好像他在墳墓裡面再重新出來.
他去到哭號一樣.
他見到他的兒女被.
拉被殺被斬被.
流亡.
到異鄉.
那種慘痛的經歷的時候.
被擄.
那種經歷.
所以拉傑就算死在墳墓也起來了.
去到為他的兒女哭號.
但是你要知道.
新約的作者.
每逢引用舊約的經文的時候.
是引用了上下文的.
他們不是引用一句.
他們是引用上下文的其中一句.
下文.
耶利米說什麼呢?.
耶利米三十一章最後.
他說這個.
的眼淚.
一定會上帝會擦去的.
因為他們有回歸的盼望.
因為上帝是信神.
上帝是不會忘記.
上帝的咨文.
上帝是不會忘記他的兒女的.

$^{521}$正如拉傑就算死了.
也不會忘記他的兒女.
那些被殺戮或者流亡的人.
上帝一定會紀念他和大衛.
甚至亞伯拉罕的約.
上帝不會丟下他們.
其實這個馬太福音.
一直第一章第一節.
就是想說這件事.
就是說上帝.
怎麼都是非福.
無論你今天落在什麼處境裡面.
當時的猶太人當然是落在一個暴君的底下.
是很苦的一個情況.
由第一章第一節.
說到耶穌的家譜的時候.
馬太也提醒我們.
這個不是在說遠古的歷史.
而是在說今天.
上帝依然守約.
上帝依然帶領著他的子民.
你記得那個家譜說過什麼.
兩個.
是不是兩個星期.
但是那兩個星期前我說過.
就是耶穌基督的家譜.
是亞伯拉罕的子孫.
是大衛的子孫.
然後.
亞伯拉罕.
傳遞到大衛有十四代.
由大衛傳遞到.
就是流亡巴比倫有十四代.
流亡巴比倫.
到基督降生.
是有十四代.
十四十四十四.
他就是告訴你.
就算我們是流亡.
就算我們是在一個慘痛的經歷裡面.

$^{561}$一個困境當中.
上帝是不會丟下他自己的子民.
所以被擄的眼淚.
這裡說應驗先知耶利米的話.
就將會被拆去.
因為真正的大衛的君主.
已經到了.
已經誕生.
上帝不會忘記他的兒女.
內華這兩天是那個.
這些小孩子.
被殺的.
故事似乎就是提醒我們.
告訴我們暴政往往.
是想保障自己的地位.
他不惜要去殺害一些無辜者.
就算是兩歲的小孩.
在人看來.
好像是政權在負責.
是真的.
你看看我們身邊.
國家,社會,世界.
似乎是那個暴政.
是他說了算.
他說要抓就抓.
他說要殺就殺.
他說要你流亡就要你流亡.
那些貧苦的百姓.
好像是很無助的.
正如百里行兩歲的嬰孩一樣.
在馬太福音小朋友這個字.
不是單單小朋友.
馬太福音小朋友,小孩子.
the least among us.
是說到初代的教會.
他們是代表初代的教會.
他們是受打壓.
是受迫害的.
在人看來是一群無助的人.
但那個故事相反.

$^{601}$就是說似乎是暴政在負責.
似乎我們是無助的.
但是呢.
上帝的工作就是藉著小孩子.
藉著耶穌基督.
藉著一群無助的人來推薦.
我不知道你2019年尾的時候.
可能也是一個好的時間我們去review一下.
到年尾的時候經常有些很無聊的statistics出來.
其中一個我這兩天碰到的一個很無聊的statistics.
就是在2019年最多人用的emoji是什麼.
我不知道你今年有沒有用過emoji.
emoji就是這些.
我們打微信.
打whatsapp這些text的時候.
我們就會有一些公仔.
來表達我們的情緒.
這個無聊的statistics是說什麼呢.
最多人用今年2019年.
就是第一個了.
就是face with tears of joy.
最多人用.
第二個就是red heart.
就是拍拖的那些特別多人用.
接著一直數下去.
我不知道你最多人用的emoji是什麼.
今年.
我最多人用的emoji.
是anger.
真的.
我看到很多新聞.
看到很多headline 看到很多直播的video.
無論你說香港.
說世界哪一個地方.
我最多用的我留意到的.
就是anger.
我用到一個地步.
我覺得anger都表達不了我.
心裡面.
對一些政權或者政局一些濫捕等等.

$^{641}$那種情緒.
我不知道你有沒有這樣的同感.
或者我看到一些身邊的朋友患重病.
突然間患重病.
我覺得.
很多問號.
我都覺得很無助.
我看到教會在不同的地方.
包括中國大陸.
特別是中國大陸.
受打壓.
那種情緒我都不知道怎麼去表達出來.
那種helplessness.
可以sum up我.
今年.
2019年.
那種的.
心境.
我就沒有什麼red heart.
我也沒有什麼.
很浪漫的東西.
但是很多anger.
在祈禱裡面.
我留意更加多.
在詩篇裡面一些.
申訴的圖文.
2019年.
我不知道你2019年怎麼過.
或者你可能都是很無助.
你會覺得.
in charge的不是你.
是一些暴政或者很evil的東西.
或者一些.
很邪惡的一些勢力.
無論是身體的病痛.
或者社會裡面的一些狀況.
但是這個故事告訴我們.
上帝工作.
上帝的鏡頭.
不是一些正要的人物.

$^{681}$上帝在歷史裡面那個camera.
不是那個.
更加不會是那個暴政.
更加不會是那個昏君.
而是那些無助的嬰孩.
在馬達福音第二章記載.
一個故事.
裡面.
我不知道你有沒有留意.
我剛才讀經的時候.
經常提到一個字.
小孩子.
和他母親.
那個是最無助的小孩子.
和他媽媽.
無辜的受害的.
其實.
就好像那些殉道者.
那個小鎮.
伯利恆.
或者耶穌基督.
最後他長大的那個地方.
Nazareth.
拿撒勒.
一個很寂寂無名的一個地方.
這個就是上帝推展他的計劃.
推展歷史.
的重心.
是一群似乎是無助無能的人.
留意二章裡面.
去記載.
耶穌基督降生的事.
包括第一章最後.
講到約瑟.
被報夢說要娶瑪利亞過來.
要應驗.
以賽亞先知第七章.
和女兒要為人生子.
有五次.
記載耶穌基督降生的事.

$^{721}$五次.
馬太說.
是應驗.
Fulfill.
上帝.
在人世間裡面.
要成就的歷史計劃.
他要帶來救恩.
他要帶來.
真正的主宰.
真正的王.
是沒有東西能夠阻止到他.
無論有昏君.
或者有個暴政也好.
上帝仍然可以.
仍然可以推展他自己.
的計劃.
沒有一件事能夠阻止到他.
在馬太福音.
不單止這五次講到應驗.
整本書十二次講到應驗.
就是在耶穌基督.
最無助的時候.
在十字架上面.
都是要應驗.
上帝.
所講的說話.
要應驗上帝在歷史裡面.
所推展的.
救贖的工作.
你不要忘記.
今年你再無助也好.
今年你覺得.
好像被.
世界的巨輪.
去到.
去到糟蹋著你.
或者你對你的生活.
好像沒有什麼控制的時候.
上帝仍然是工作.

$^{761}$他的鏡頭.
看清楚.
在耶穌基督身上.
還有屬於耶穌基督的人.
在人的眼中.
不是 in charge.
不過在上帝眼中.
這個就是他的 focus.
年尾的時候.
我們將我們的生命交給主.
好嗎?.
又一望新的一年.
我告訴你新的一年.
一定不會.
更加好的.
我想不會有明年會更好.
明天會更好.
不會有這次國際場合.
但我們知道上帝仍然掌管.
我們同心協力祈禱.
主我們是熟悉的子民.
我們知道你.
仍然是掌管歷史.
掌管我們生命的.
無論是我們的生活,我們的生命.
我們看到很多的混沌.
很多的混亂.
很多的不安.
主有多麼動蕩也好.
我們知道你仍然是推展著自己的工作.
獨生子耶穌基督.
我們的救主.
祂的降生.
進入到我們這麼動蕩,這麼不安的世界.
在十字架上為我們所作.
我們能夠成為你的兒女.
主讓我們的眼睛.
看見你自己仍然做.
你仍然是作主.
讓我們投靠的.

$^{801}$是單單投靠你.
不是世上的重要人物.
不是世上的權柄.
而是你的權柄.
主我們是屬於你的兒女.
我們將我們今年所經歷的.
無論是高山或者是低谷.
我們深信你仍然是掌管.
去推展著你的計劃.
而我們就是被你呼召.
來參與你的工作.
主叫我們仍然去堅信.
去投靠,去跟隨你.
去深信你仍然作供.
我們就跟你同供.
用我們的生命去見證你.
讓我們的生命能夠在.
你自己掌權的這個事實上.
能夠得到平安.
主讓我們的生命帶到這個世間.
我們的家庭.
我們的不同的崗位當中.
讓我們讓人知道.
神你仍然是主宰.
我們這樣祈禱奉祝耶穌基督的名字.
\newpage



\section{腓立比書 2:6-11}
\label{sec:HLQkwHUhsgc}
\textbf{1/12/20 主日講道: 天國公民地上行 2 (腓立比書2\_6-11) - 陳崇基牧師}
\newline
\newline
連結: \href{https://youtube.com/watch?v=HLQkwHUhsgc}{\texttt{https://youtube.com/watch?v=HLQkwHUhsgc}} ~~~~ 語音日期: 2020-01-19
\newline
\newline
\hyperref[sec:4ZoSLfbOWng]{\small{< < < PREV SERMON < < <}}
~
\hyperref[sec:index_chronic]{\small{[返順時目]}}
~
\hyperref[sec:index_scriptual]{\small{[返順卷目]}}
~
\hyperref[sec:6ypOdtEuNl8]{\small{> > > NEXT SERMON > > >}}
\newline
\newline
腓立比書 2:6-11
\newline
\begin{longtable}{cl}
\hline
\hline
章節 & 經文 (和合本修訂版)\\
\hline
2:6 & \begin{tabularx}{0.7\textwidth}{X} 他本有神的形像，卻不堅持自己與神同等； \end{tabularx} \\ \\ \relax
2:7 & \begin{tabularx}{0.7\textwidth}{X} 反倒虛己，取了奴僕的形像，成為人的樣式；既有人的樣子， \end{tabularx} \\ \\ \relax
2:8 & \begin{tabularx}{0.7\textwidth}{X} 就謙卑自己，存心順服，以至於死，且死在十字架上。 \end{tabularx} \\ \\ \relax
2:9 & \begin{tabularx}{0.7\textwidth}{X} 所以神把他升為至高，又賜給他超乎萬名之上的名， \end{tabularx} \\ \\ \relax
2:10 & \begin{tabularx}{0.7\textwidth}{X} 使一切在天上的、地上的和地底下的，因耶穌的名，眾膝都要跪下， \end{tabularx} \\ \\ \relax
2:11 & \begin{tabularx}{0.7\textwidth}{X} 眾口都要宣認：耶穌基督是主，歸榮耀給父神。 \end{tabularx} \\ \\
[1ex]
\hline
\hline
\end{longtable}
$^{1}$我們一起打開去.
嗯.
肥立秘書第二章.
六到十一節.
他本有神的形象.
不以自己與神同等.
為強奪的.
反倒虛己.
不以自己與神同等.
為強奪的.
反倒虛己.
不以自己與神同等.
為強奪的.
反倒虛己.
取了奴僕的形象.
成為人的樣式.
既有人的樣子.
就自己卑微.
全心馴服.
以致於死.
且死在十字架上.
所以神將他升為至高.
由此給他那超乎萬名之上的名.
叫一切在天上的.
地上的.
和地底下的.
因耶穌的名.
無不屈實.
以耶穌基督為主.
此榮耀歸於父神.
我們上星期.
開始這個系列的時候.
我們就說.
到底一個基督徒.
我們信耶穌的.
在地上生活的時候.
到底我們的身份是什麼.
有什麼是最重要的身份.
我們知道我們在世裡.
我們的身份經常會變.

$^{41}$Right.
你的工作.
你的家庭.
你個人的生活的狀況.
事實上都會改變.
但是在聖經裡告訴我們.
基督徒有些很獨特的.
identity.
這些不同的身份.
告訴我們我們應該怎樣去處世.
應該怎樣去生活.
應該怎樣去.
過地的生命.
所以我們提到在基督徒.
的.
聖經裡.
到底有什麼提醒我們.
我們的身份.
我們就嘗試回到最初的信徒.
所謂大概.
2000年之前,第一世紀的時候.
當耶穌死而復活之後.
他.
留下的教會.
初代的教會,到底他們.
在耶穌的經歷體會裡.
到底怎樣去.
inform他們的identity.
所以我們去看看新約聖經.
裡面的文件.
這些書信.
或者這些記載.
怎樣去指導我們.
告訴我們我們是一個怎樣的人.
一個信耶穌的人.
我們的身份是怎樣.
當你翻開新約聖經的時候.
你會明白不是.
放在頭的那幾本書.
就是先寫的.

$^{81}$你知道新約聖經裡面有27本書.
是不是.
有頭四本是叫做福音書.
馬太,馬可,路加和約翰.
這四本.
沒錯是記載耶穌的事跡.
他的神跡.
他的死而復活等等.
但這些絕對不是.
那27本書裡面.
開頭寫的那四本.
我們所知道這四本福音書.
是在初代的教會.
之後的第二代.
第一代就快離開的時候.
大概.
是公元後.
主後.
60年代的時候.
才開始執筆去寫.
這很合理的,為什麼呢.
因為第一代的目擊證人.
那些的眼見證人.
那些跟隨耶穌的門徒.
他們就快要離開這個世界了.
大概主後AD.
60年代的時候.
很合理的,他們就要將這些.
眼見證人的帳號.
這些見證人,目擊證人的帳號.
就要寫下.
記載下.
以致第二代第三代以後.
所有的教會他們能夠明白.
到底耶穌是一個怎樣的人.
所以很合理的.
就好像今天那些猶太人.
他們現在.
到2020年.
那些.

$^{121}$所謂經歷過二次大戰.
在納粹黨.
底下他們受到大屠殺.
的殘害.
他們是生還者.
他們很多做一些.
所謂眼見證人的.
面試,留下一些.
重要的.
記載,或者記錄.
告訴下一代.
到底曾經.
發生過什麼事.
所以我們打開福音書的時候.
我們發現不是開頭寫的.
而最開頭寫,最早期寫.
對於基督徒.
或者教會裡面.
所看的聖經,所看的經文.
到底是什麼呢.
肯定不是福音書裡面的內容.
肯定是.
一些保羅的書信.
保羅是當時的一個牧者.
一個仕徒.
他去周圍羅馬帝國.
不同的城市裡面建立教會的時候.
他就會寫一些書信.
來指導教會.
或者去教導他們.
或者去改正他們.
某一些的思想,某一些的概念.
告訴他們福音.
是什麼.
我們在心靈裡面.
去運用福音.
去應用福音的時候.
到底應該是怎麼樣的呢.
所以在這些最早期的文件裡面.
我們大概知道.

$^{161}$幾乎追溯到.
第一代的信徒,到底他們在想什麼.
他們怎樣體會.
到底耶穌是怎樣的.
我們上個星期開始看.
這個肥立比書.
肥立比書就是保羅.
寫的一些書信.
我們說這本書信.
是說很多的所謂榜樣.
一些.
跟隨耶穌基督應有的.
生命的榜樣.
那些例子.
所以這本書信很短,也很簡單.
他就告訴肥立比.
的信徒.
是他自己生的教會.
是他自己所建立的教會.
他現在正在坐牢.
為了傳道的緣故,他受打壓.
然後他現在坐牢.
去跟他們同心.
去供應他們一切的需要.
在監獄裡面去看他們.
或者去供應他們所缺乏的東西.
這是第一章.
第二章呢,保羅就說.
我們需要學校的榜樣.
所以他就說了兩個跟他一起.
致死忠心的同工.
一個叫提摩泰,一個叫.
爾巴弗提.
然後最後呢,他就說到.
一個他自己的榜樣.
那個榜樣就是.
他好像耶穌基督.
的僕人的身份.
他去服侍.
最後就是結語.

$^{201}$呼籲教會謙卑的.
來同心.
在整本書信中間夾著的餡.
就說到.
基督是怎樣偉大.
這一首的詩歌.
大概應該不是.
保羅自己親手寫的.
因為我們知道那些字眼.
或者他用的一些用語.
不會是保羅平時.
一些書信所用的.
並且這首.
的詩歌.
是保羅引用的.
換句話說.
他寫這一段的.
書信的時候.
他是在引用.
初期教會,初代教會.
在敬拜裡面.
在用的一首詩歌.
或者他們.
初期教會走在一起的時候.
他們去教導大家.
去講述耶穌基督的時候.
他們就會提到.
這一首的詩歌叫.
Christ Him.
這一首的詩歌.
大概就是去.
告訴基督徒.
告訴教會.
我們的身份.
是怎樣的.
上個禮拜你記得我說過.
在寫這本書的時候.
有一處,兩處的地方.
保羅說我們的身份.
就是天國的公民.

$^{241}$是用羅馬帝國的.
用語去告訴我們.
我們是天國的公民.
這個比我們羅馬帝國.
比我們在地上.
帝國的身份.
更加重要.
這裡大致上應該是告訴我們.
我們作為一個僕人的身份.
如果你問保羅的話.
當他寫信的時候.
很多時候會自己介紹自己.
你覺得他最多時.
或者最喜歡.
介紹自己.
是用什麼名銜去介紹自己.
保羅有時會用.
使徒.
這個稱呼來介紹自己.
但我相信.
最多時.
或者他最喜歡用的.
就是僕人,servant.
所以在腓立比書第一章第一節.
如果你有聖經.
你自己翻開看看.
他也是這樣介紹自己.
他說我是基督耶穌的僕人.
現在我寫信給.
菲律賓的教會.
servant.
或者servant.
我們聽見servant的時候.
我們想起什麼呢.
我們想起可能是菲傭.
我們想起可能是nanny.
我們想起可能是我們請的人.
清潔工人等等.
在羅馬帝國的社會階層裡面.
這是最低的社會階層.

$^{281}$在羅馬帝國的聽眾.
聽見servant.
或者bond servant.
僕人的時候.
他們會想起什麼呢.
就是想起社會學家所說的.
social death.
就是說在社會裡面.
你去到.
自殺的一個階段.
或者是.
你去到.
自殺的一個稱號.
因為每一個在羅馬帝國裡面的人.
他們希望攀爬.
那個.
corporate ladder.
那個樓梯的時候.
他們希望能夠脫離.
這個servant的身份.
成為一個freeman.
或者甚至到上面.
成為一些的貴族.
最大的那個當然就是.
羅馬帝王.
沒有人喜歡.
做一個servant.
他們希望能夠攀爬到.
到帝國上面.
令到他自己引以為傲的.
一個社會的地位.
所以.
你用一個第一世紀.
羅馬人的.
耳朵.
去聽一聽.
這些書信所說的話.
你就知道那種震撼是多大.
居然一個.
一個.

$^{321}$仕途.
他最喜歡去介紹自己.
他說我就是一個.
servant.
我們說過.
羅馬帝國有很多的殖民地.
這些殖民地.
就有著一種這樣的文化.
菲律賓就是其中一個這樣的city.
在每一個殖民地裡面.
都會有一些的.
政治的宣傳.
當時那些羅馬帝國的官媒.
那些中央的電視台.
他們就會用不同的方法.
在不同的殖民地的城市裡面.
告訴大家.
到底什麼才是最大.
哪一個才是最大.
哪一個是值得你們去景仰.
在每一個的羅馬的.
這些.
殖民地的地方.
就是要去宣揚.
羅馬帝國的好處.
和羅馬帝國的文明.
你告訴別人.
多麼仁義為傲.
那個統治者叫做凱撒.
Caesar.
Caesar會在不同的地方.
不同的羅馬帝國.
所征服的城市.
殖民地那裡.
他就會建立他的官媒.
當時當然是沒有電視.
也沒有smartphone.
官媒在哪裡呢?.
在每一個城市的中間的市中心.
他一定會樹立很多的神像.

$^{361}$這個神像就是要告訴.
那裡的平民百姓.
或者那個戰敗的民族.
聽什麼.
哪一個.
是主宰的.
哪一個是lord.
哪一個是統治.
天下的.
哪一個是萬口要.
敬拜他.
或者萬口要稱他為主的.
這個是阿古士都.
Augustus.
是一個羅馬帝王的神像.
特別在地中海的東邊.
很多的羅馬的.
殖民地的城市.
是找到這一類的神像.
你看看這些神像.
你就知道不單單是.
博物館裡面的雕塑.
很漂亮的東西.
其實是一些政治宣傳的手段.
你看看這個.
人像.
是一個當時希羅的人像的雕塑.
是最完美的身材.
各位看看.
你有沒有這樣的身材.
這是一些很標準的.
Greco-Roman的雕塑.
他就是要告訴你.
這個Augustus.
是一個多麼的威風的人.
並且他雕塑這個神像的時候.
Augustus已經不是一個青春的人.
但是他雕出來的形象.
就是一個很青春的身材.
還有他的高度等等.

$^{401}$三維.
是當時Greco-Roman最標準.
最完美的一個比例.
最標準最完美的一個比例.
這東西有多高呢?兩米高.
六尺六吋.
你有沒有六尺六吋高?.
我也希望我有六尺高.
我知道我永遠不會有.
六尺六吋高.
並且有一個戰士的風格.
你看到他舉起右手.
是他指揮他的羅馬軍隊.
來征服天下.
還有.
在他右腳的底下.
如果你近鏡去看的話.
就有一隻.
海豚.
海豚象徵什麼?在羅馬帝國裡面.
他們的海軍.
他們能夠.
有最好的戰艦.
最強的戰艦.
並且打遍天下無敵手.
很多時候都是在海戰上贏.
在Dauphin上面.
有一個女神.
那個叫Venus.
為什麼呢?因為Augustus領養.
那個之前的皇帝.
叫做Julius Caesar.
他說他自己是神.
Julius Caesar說.
他是女神生他出來的.
那個就是Venus.
所以他就將這個Venus.
放在他腳下.
他說他自己也是.
神聖的後代.

$^{441}$在胸口上面的.
所謂護心牌.
護心鏡.
你記得在.
以弗所書第六章.
說到軍隊.
屬靈的軍隊的軍裝.
Breastplate.
那個護心板.
雕了很多.
這些人像.
是代表羅馬帝國軍隊.
征服過的一些民族.
在這個人身上.
就是萬口萬實.
都要稱他為Lord.
那個人像也是赤腳的.
在希羅文化裡面.
這些人像.
每個人都要穿鞋.
除了某一些人不用穿鞋.
那些就是神人就不用穿鞋.
只會神像.
才不用穿鞋.
所以這個Augustus就是說.
我就是.
Son of God.
我就是.
天子.
就是去統領萬有的.
並且Augustus.
最重要的或最多時用的.
一個自我的稱號.
就是與神同等.
Isothero.
Isothero就是說我和上帝.
是同等的.
是God equal.
是和上帝是同樣的地位.
這個就是他有的所謂特權.

$^{481}$與生俱來的特權.
他統治帝國的時候.
他就說他自己.
就等於上帝.
在這篇的基督的讚美詩.
Christ Him.
如果聽了我剛才.
類似的介紹.
你就會發現這個.
Christ Him.
有很多的詞彙.
其實就是直接用羅馬帝國.
政權用的語言.
來描述基督.
他說他有神的形象.
與神同等.
耶穌基督是與神同等.
Isothero.
羅馬人一聽見的時候.
他就發現這句話.
就是.
凱撒羅馬王.
去自稱的一個方法.
現在就用在了耶穌基督身上.
但他做了一些顛覆性的事情.
不以自己與神同等.
Isothero.
為強奪的.
強奪這個字.
也是羅馬帝國的政治用語.
什麼叫強奪.
強奪就是說.
去建立自己的power base.
去征服不同的人.
然後將他們打壓打殘的時候.
他們跪在他面前.
然後成為他power base的一部分.
強奪就是這個意思.
藉著不同的方法.
軍事的方法也好.

$^{521}$或者以暴制服平民的方法.
以壯大自己的.
權力的底細.
這個就是羅馬帝國.
有的政權.
這個就是羅馬帝國所說的.
LORD.
這個才叫做主.
你想想在一個殖民地裡面.
羅馬殖民地裡面.
他們一直受這些宣傳.
或者這種意識形態.
來影響的時候.
當他們去.
相信耶穌之後.
你想想對他們的世界.
有多upside down.
有多完全反轉的概念.
真正的主.
這首詩歌.
是一個不以.
自己與神同等.
為強奪的.
有多大的震撼.
當耶穌基督一到的時候.
就沒有人.
可以稱自己為主.
連羅馬帝國的帝王.
都要說.
我只不過是一個僕人.
政權成為僕人.
當我以為.
我自己是一家之主的時候.
什麼都要聽我說的時候.
當耶穌基督.
到了你家裡.
你就成為一個工人.
你就成為一個servant.
不單止這樣.
耶穌基督自己本身.

$^{561}$也是一個servant.
很多時候我們看這些詩歌的時候.
我們可以欣賞它的詞彙.
但很多時候我們沒有欣賞.
的地方就是.
這首詩歌裡面有一些.
辯駁的性質.
一些prolimic.
其實聖經從頭到尾.
都有很多這種prolimic的nature.
這些辯駁的.
性質辯駁的元素在裡面.
辯駁的性質辯駁的元素在裡面.
顛覆當時的帝國.
所說的東西.
就將它反轉來說.
其實這個在猶太人的文學裡面.
一點都不陌生.
因為猶太人從開國到.
第一世紀的時候.
他們都是面對很多不同的帝國.
你記得在舊約裡面的故事.
他們是出埃及.
他們面對的是.
埃及帝國.
然後有一點獨立的時間.
大衛,所羅門等等的時間.
然後就有亞述.
有巴比倫.
有波斯.
有希臘.
然後到最後有羅馬帝國.
他們很熟習.
在帝國的.
陰影底下.
來堅持著自己的身份.
所以他們就將.
這些語言.
帝國認為最.
厲害的東西.

$^{601}$或者帝國認為最大榮譽的東西.
他們就將它拿出來.
然後去反轉.
保羅引用這段的.
Christine.
就是將耶穌反轉過來.
他本身是有神的形象.
不是他自稱的.
不是他自己拿回來的.
耶穌基督本身.
就是與神同等.
但不去到強奪.
他肯去放手.
肯去自我.
卑微,情願放手.
下來,甚至死在十字架上.
十字架在.
羅馬帝國.
的政治犯.
的下場.
是一個最失敗.
的人.
一個政治的囚犯.
一個願意.
與神同等.
下來.
放下他自己所有的title.
放下他自己所有的榮譽.
他說他寧願.
用一個最低微的方法.
來去.
服侍我們.
他說他反倒虛奇.
取.
露骨的形象.
當我們一聽到虛奇的時候.
英文說.
empty himself.
是什麼意思?我們以為任人魚肉.
其實是不是一定是這樣.

$^{641}$前幾天.
我女兒去了一個.
他們大學生.
的conference.
大學生的conference.
就是在Whitting College舉行.
現在他們這些年輕人.
去conference.
去所謂的退休會.
跟我們那一代不同,他們住酒店.
真是厲害.
載了他們.
去register的地方,載他們去.
酒店.
我在city上個星期三.
載他和我女兒.
和他們兩個朋友.
好,我載你去.
我們去Whitting College.
去register.
我就問他,去到了.
到底你們在哪裡register.
他們就找他們的資料.
原來.
就是一個叫.
Todd Beamer.
的student center.
他說我知道在哪裡.
我也在Whitting.
出入過.
我就載了他們去.
Todd Beamer student center.
他們很快就register了,拿了自己的牌.
上車,我們載他們去酒店.
載他們去酒店途中.
我問他們三個.
我說你知不知道誰叫.
Todd Beamer.
三個人.
當然說不知道.

$^{681}$我說不如你google一下.
到底Todd Beamer是一個怎樣的人.
為什麼那個building.
就是Whitting College裡面的building.
會用他的名字.
起名.
Todd Beamer到底是誰.
Todd Beamer.
是.
十幾年前2001年的時候.
2001年.
大家在世的時候,記得2001年發生什麼事.
9月11號.
是September 11的時候.
Todd Beamer就上了飛機.
那個飛機本身.
飛去三民市.
上了飛機沒多久.
那個飛機叫做聯合航空93.
United Flight 93.
他在飛機那裡.
他就知道有劫機者.
劫機者就控制了那個飛機.
40分鐘之後.
他和後面的一些乘客.
就去談.
談一下怎樣去.
破壞.
這些hijackers的計劃.
Todd Beamer就打電話.
去給一個operator.
電話就和那個女人說.
我們幾個乘客.
我們決定這樣做.
決定去衝擊他.
破壞他的計劃.
然後他們在電話那裡.
他們.
應該按照那個operator說.
他們就讀了主圖文.

$^{721}$詩篇23篇.
然後Todd Beamer就和那個女士說.
我們可不可以一起祈禱.
祈禱之後.
Todd Beamer所謂的last word.
最後所知道的non-last word.
就是說Let's roll.
我們就出去.
於是乎這班的班機.
就在Pennsylvania一個郊區的空地.
那裡就聚會.
整部機裡面.
每一個44個乘客就全部喪生.
有很多人猜測.
到底這班機本身想飛哪裡.
有些人說是威口士.
有些人說是國會大廈.
但你要知道.
Todd Beamer.
他這樣做的時候.
他有很多牽掛.
他太太有第三個.
的小孩.
有新人.
他可以有很多的牽掛.
或者有很多的保留.
可以叫做.
或者用.
Christian的.
詞語來講.
就是可以是強奪.
就是.
我possess的東西我不想放手.
是一種自願,甘願.
等低自己.
或者忘我.
這叫做虛己.
虛己不是說我們做個dormant.
做個地氈任人魚肉.
任人踩.

$^{761}$虛己相反在聖經裡面講.
其實就是我限制我自己.
放棄我自己某些的權利.
以致.
祝福其他的人.
這叫做虛己.
這叫做empty himself.
這裡講一講文法.
我知道大家不是很喜歡這些grammar.
英文文法裡面.
有一樣東西叫做pronoun.
right.
代名詞.
就是himself, herself.
right, 這叫做pronoun.
這裡用了一個字虛己.
是一個叫做reflexive pronoun.
或者.
中文可以叫做反新.
的代名詞.
什麼叫做反新代名詞呢.
就是中文譯法呢.
中文譯法叫做.
什麼什麼自己.
就是你做了一些事.
是你自己選擇.
主動選擇.
是做在你自己身上.
不是被逼.
不是被人強行.
而是你主動的放棄.
某一些東西.
這首詩篇.
講到的就是.
虛己.
什麼叫做虛己.
他願意.
自己主動的放棄.
自己的title.
自己的地位.

$^{801}$一種忘我的表現.
你想想如果教會.
如果基督徒.
有一個這樣的虛己的.
僕人的身份的時候.
你回到第一世紀.
會對羅馬帝國.
一個多大的震撼.
震盪.
在不同的城市當中.
雖然有一群人.
他們是不跟著.
羅馬帝國的corporate ladder.
他們的生活.
大家彼此是同心.
彼此能夠謙卑自己.
自我卑微.
empty himself.
願意.
是反身的.
去說我願意放下我自己.
就不跟帝國.
那種虛榮.
不跟帝國世界.
那種社會的.
動作.
去做.
跟著耶穌基督自己的榜樣.
這一種另類的.
politics.
對一個帝國的震撼是多大.
他們所稱的主.
不是Caesar.
所稱的主.
是主耶穌基督自己.
一個最低微的人.
本有上帝的形象.
但他願意降卑.
降下來.
你想想如果我們帶著.

$^{841}$這樣的身份進入我們的.
家庭.
進入我們的社會.
會是一個多麼的震撼.
的一回事.
基督.
成為我家庭的主宰.
right.
我們家庭我們的夫婦.
能夠在這個主宰下.
彼此作僕人.
基督.
是教會的主.
我們大家一起同工.
一起放棄.
我們自己本有的.
名.
title.
本有的地位.
本有的所謂privilege.
這個社會.
另類的運作.
就是跟世界說.
真正的主.
是耶穌基督.
真正有價值的東西.
是那位願意寫在十字架上的主.
\newpage



\section{腓立比書 2:6-11}
\label{sec:6ypOdtEuNl8}
\textbf{1/5/20 主日講道: 天國公民地上行 1 (腓立比書2\_6-11) - 陳崇基牧師}
\newline
\newline
連結: \href{https://youtube.com/watch?v=6ypOdtEuNl8}{\texttt{https://youtube.com/watch?v=6ypOdtEuNl8}} ~~~~ 語音日期: 2020-01-19
\newline
\newline
\hyperref[sec:HLQkwHUhsgc]{\small{< < < PREV SERMON < < <}}
~
\hyperref[sec:index_chronic]{\small{[返順時目]}}
~
\hyperref[sec:index_scriptual]{\small{[返順卷目]}}
~
\hyperref[sec:Pe1TNggJhos]{\small{> > > NEXT SERMON > > >}}
\newline
\newline
腓立比書 2:6-11
\newline
\begin{longtable}{cl}
\hline
\hline
章節 & 經文 (和合本修訂版)\\
\hline
2:6 & \begin{tabularx}{0.7\textwidth}{X} 他本有神的形像，卻不堅持自己與神同等； \end{tabularx} \\ \\ \relax
2:7 & \begin{tabularx}{0.7\textwidth}{X} 反倒虛己，取了奴僕的形像，成為人的樣式；既有人的樣子， \end{tabularx} \\ \\ \relax
2:8 & \begin{tabularx}{0.7\textwidth}{X} 就謙卑自己，存心順服，以至於死，且死在十字架上。 \end{tabularx} \\ \\ \relax
2:9 & \begin{tabularx}{0.7\textwidth}{X} 所以神把他升為至高，又賜給他超乎萬名之上的名， \end{tabularx} \\ \\ \relax
2:10 & \begin{tabularx}{0.7\textwidth}{X} 使一切在天上的、地上的和地底下的，因耶穌的名，眾膝都要跪下， \end{tabularx} \\ \\ \relax
2:11 & \begin{tabularx}{0.7\textwidth}{X} 眾口都要宣認：耶穌基督是主，歸榮耀給父神。 \end{tabularx} \\ \\
[1ex]
\hline
\hline
\end{longtable}
$^{1}$如果有聖經的話 大家可以去到.
肥納比書第三章第20節.
三章第20節.
我們卻是天上的國民 並且等候救主 就是主耶穌基督從天上降臨.
他要按著那能叫萬有 歸服自己的大能.
將我們這卑賤的身體改變形狀 和他自己榮耀的身體相似.
這裡其中一節 肥納比書三章第20節.
講到我們基督徒的身份 我們這個系列會看看新約裡面.
我們知道新約裡面有多少卷書 27卷書.
這27卷書裡面有很多都是當時的教會領袖 使徒 牧者.
他們寫給當時初代的教會 當時的教會已經在不同的城市裡面.
已經藉了堂 在不同的地區.
不同的城市裡面 也有很多不同的所謂的小組教會.
不同的家庭聚會的教會.
在這些教會散佈在羅馬帝國不同的地方的時候.
有很多時候領袖的使徒 不但會去不同的城市各鄉做教導和栽培.
他們也會寫信 給他們一些指導.
我們會翻開聖經裡面一些新約的書信.
一些所謂教會的書信.
Church letters 或者epistles.
我們看看裡面如何教導初代基督徒.
他們在世界各地不同的地方他們有的身份是什麼.
三章二十節這裡很清楚講.
就是說我們到底是一個怎樣的身份 什麼的identity.
我們是什麼.
我們翻開這本書的時候.
我們知道是叫做《肥納比書》.
如果你不是很清楚聖經的格局或者背景.
你就知道肥納比書.
肥納比是一個地方.
是當時羅馬帝國的一個城市.
裡面有一些初代的第一世紀的信徒 教會建立了.
肥納比書我們很簡單地告訴你.
到底肥納比書是講什麼的.
是一封信來的.
當時的使徒 牧者保羅叫做Paul.
他寫信給在肥納比居住的信徒 教會.
他寫這封信的時候他在坐牢.
因為傳道 傳福音的緣故.
他被迫害.

$^{41}$然後他在監獄裡面寫這封書信.
寫給肥納比的教會.
有四章的聖經 有四章的經文.
因為肥納比教會就像他的子女一樣.
是他自己建立的教會.
由第一天開始他見到這個教會 信主.
建立起來在肥納比這個城市裡面能夠開始聚會.
是他自己生的兒女.
所以他跟他們有一個很個人的關係.
正如我們今天看 別人的子女不孝 我們就不這麼不開心.
自己的子女不孝 我們就很不開心.
別人的子女乖 我們就不這麼開心.
或者也開心的.
但是自己的子女乖 特別開心.
是自己生的子女.
保羅見到自己生的教會也是一樣.
他們一直支持他的工作 一直關心他.
並且告訴他他在監獄裡面的經歷.
你記得裡面有一段他提過.
他死都無所謂.
因為死更加是一個大的祝福.
因為他可以跟主同在.
但是如果不死的話.
他就留在這裡可以更多工作 更多事奉.
所以他就用他自己的經歷來勉勵肥納比的教會.
第二章他就去鼓勵肥納比教會.
他們要同心 要謙卑的事奉.
然後他繼續說下去.
他就用了兩個例子來告訴肥納比教會.
為什麼你需要這樣做.
因為你們的目的是這樣.
提摩泰是其中一個保羅的得力助手.
他是一個很忠心的僕人.
另外一個就是爾巴弗提勒.
爾巴弗提勒雖然我們不是常聽這個名字.
但是他是肥納比教會其中一個的牧者.
他為了去供應保羅在監獄裡面的需要.
差不多他病到要死.
所以保羅就提到這種基督的榜樣.
第三章他就說到他自己作為一個榜樣.

$^{81}$告訴教會他自己如何去捨棄他以往猶太人家族的尊貴.
那種地位 社會裡面的身份.
他拿回上帝給他基督裡面的身份來事奉他.
第四章最後他就呼籲教會要同心.
呼籲教會特別一些可能大家走不在一起去同宮的人.
呼籲他們要同心.
在整本肥納比書中間的一個地方.
就是一個保羅所寫或者保羅所引用的.
當時初代教會的一個詩章.
那個詩篇就說到基督的偉大.
而那個基督的偉大不是說到羅馬帝國那種偉大.
羅馬帝國那種偉大就是要打敗所有的異己然後提高自己.
但是基督的偉大就是從天上虛己降下來.
並且作路步 信服上帝的心意.
甚至死在十字架上 上帝才將他提升至高.
他就成為萬主之主.
這個是最偉大的詩篇.
所以保羅說看這封書信的時候.
你不要忘記最重要就是要記得我們所跟隨的耶穌基督.
他自己的榜樣是怎樣.
他的榜樣是這樣.
我們就照樣事奉在我們生活裡面.
去學孝主耶穌基督自己的榜樣.
這一章的聖經二章六至十一節.
是值得你去背的.
如果你說今年我沒有想過要整年讀一次聖經.
可能你去年說要全年讀聖經.
誰知道你讀了兩個星期之後就停止了.
你放棄了 失敗了.
你覺得很洩氣 今年又繼續讀 又為了什麼呢.
如果你覺得全年讀一次聖經太多了.
我會建議你這個月不如去學學去記聖經 去讀聖經.
背了它.
背什麼呢 肥下比書二章六至十一節.
基督的心為心.
到底基督我們的主是一個怎樣的主.
怎樣偉大化.
記了它.
如果你覺得整年讀聖經太多.
不如你記了這一段經文.

$^{121}$說到肥下比書提到基督徒的身份identity的時候.
我覺得有一段聖經或兩節聖經特別提到基督徒的identity.
剛才我們讀了三章二十節.
他說我們是天上的國民.
同一個字在第一章二十七節也提到.
你們行事為人 要與基督的福音相稱.
你不用懂英文或希臘文.
你也知道大概兩個字都是一樣.
大概兩個字都是當時羅馬帝國特別的字眼.
有點像politics.
所以我們說.
現代的翻譯比較難去掌握那個味道.
他說我們卻是天上的國民.
或者你們行事為人.
國民也不知道怎樣翻譯.
英文翻譯也好.
就是citizenship.
或者citizens.
我們就是天上的heavenly citizens.
其實這個字是一個很罕有的字.
整本聖經只用了三次.
兩次就是這裡.
另外一次就是保羅自己說自己的經歷.
說他行事為人是一個好的公民.
其實這些都是羅馬帝國一些很特別的字眼.
或者更好的翻譯應該是國籍的意思.
我們卻是天上國的.
天國裡面的公民.
有天上公民的國籍.
有heavenly citizenship.
或者可能這個字眼比較容易把握.
這個特別的地方.
當時的羅馬帝國.
他們去到佔領不同的地方的時候.
他們就會將那個地方變成一個羅馬帝國的殖民地.
Colonial City是做什麼的呢.
就是將人派到那個地方.
那個所謂被佔領的地方.
然後叫那些人去宣揚羅馬帝國的好處.
是一些很國家主義很愛國的字眼.

$^{161}$他們說你們就是天上的國民.
就是說你們不是羅馬帝國的殖民.
你們是天上天國的殖民.
就是天國殖下來我們人間裡面.
一個很帝國的字眼.
但是用在了上帝的國度裡面.
另外他講到只要你們行事為人.
行事為人這個字也不是普通猶太人講到生活這麼簡單.
猶太人講到生活的時候不會用這個字.
通常都是講走路的意思.
中文也掌握到行事為.
你走路出來.
平時走路的時候就是你的生活.
但這個不是普通典型猶太人用的字眼.
這個是羅馬人用的字眼.
行事為人是什麼意思.
就是說羅馬公民的生活方式.
這個跟菲律賓的背景是有關的.
我們講過菲律賓是一個城市.
腓立比書就是寫給這個城市的教會和信徒.
菲律賓這個地方是保羅在第二次.
去到傳道宣教的路程當中.
他停留的一個地方.
他由巴基斯坦地區出發.
經過今天的土耳其.
當時的小馬西亞.
去到一個地方叫做特羅亞.
那裡是一個海港.
去到亞洲最端端的地方.
他見到一個意象叫做馬其頓的意象.
在《少林寺傳》講到馬其頓這個地方.
就是菲律賓附近一帶的人的種族叫做馬其頓人.
因為他見到這個意象.
他見到對面岸有一個福音的需要.
所以保羅就說我們坐船過去.
所以他坐船過去經過不同的海港.
最後去到菲律賓.
菲律賓是歐洲大陸最東邊.
其中一個最重要的關口.
是聯絡東與西中間的會點.

$^{201}$在這個地方他就帶了人來信主.
是一個有希臘西方羅馬帝國味道的城市.
為什麼叫做菲律賓呢.
幾百年前希臘人佔領了這個地方.
有一個很有名的將軍叫做菲利普II.
佔領了這個城市的時候.
他就把這個城市改了自己的名字.
來給自己榮耀.
叫做菲律賓.
或者叫做菲利普.
去到這個地方保羅就帶了一個女人.
叫做女底亞來信主.
然後帶了她一家來信主.
所以你看小行村第十六章.
看到菲律賓是一個很特別的地方.
那裡說什麼呢.
他說他們從德羅亞的海港開船.
一直走到另一個地方.
第二天到了另一個地方.
在哪裡就來到近海的.
內陸一點的菲律賓.
那裡是馬其頓這一方的頭一個城.
也是羅馬的駐防城.
我們就在那個城駐了幾天.
這些字眼都是很特別.
和當時菲律賓城的處境.
菲律賓是那個城的名字.
是馬其頓的一帶這一方頭一個城.
當羅馬帝國從希臘人那裡.
拿到這個城回來的時候.
羅馬帝國侵佔了這個城市.
他就將這個城市或者這一帶.
分成四楷.
我們今天說的四邑.
將變成四邑.
最大的邑 第一的城就是菲律賓.
就將它成為駐防城.
這是和合本舊的翻譯.
新的翻譯就不再叫駐防城.
因為想起駐防城就有軍隊在.

$^{241}$但最重要的不是有軍隊在.
最重要的是那是殖民地.
是羅馬帝國的殖民地.
就是說那些羅馬帝國.
幫他打生打死的羅馬將軍.
他們打完仗.
羅馬王就會派他們去駐防城或者殖民地.
那些Veterans.
派他們去殖民地的城市.
他就要去那裡住.
然後給他物業 給他福利.
為國家打生打死.
然後就給他國民的身份.
羅馬帝國的國民的身份.
那些人住在那個地方.
他們就宣揚羅馬帝國的好處.
或者繼續去推廣羅馬帝國的版圖.
當然最重要的一點.
就是拿到羅馬帝國的Citizenship.
如果拿到羅馬公民的國籍的時候.
就會有很多福利.
稅可以減.
法律和政治也有很多的保障.
所以當時平民百姓拿不到的東西.
只有少數的人才有榮譽拿到Citizenship.
拿到那個好處.
就是Roman Citizenship.
其實有點像今天.
美國國籍其實也很輕香.
我們經常有些人說.
美國也不是很好 有些人很老.
其實拿到美國國籍是很輕香的.
我記得有一次我在中東.
以色列.
我要搭長途巴士去埃及.
我就去那裡買票.
跟我一起買票有其他旅客.
我去世界不同的地方.
我記得買車票的那個姐姐.
就跟我說.

$^{281}$我們中間有幾個美國人.
拿著美國的護照.
我是其中一個.
她說你們都不知道.
我們拿著不是美國的護照.
或者是一些中東國家的護照.
你們拿著美國的護照是多麼方便.
多麼好.
多麼有保障.
那些說美國不好的人.
通常都是拿著美國的護照才說不好.
身體是很誠實的.
不是啊.
那些中國大陸的高官.
為什麼每個子女都跑去美國.
沒有那麼好的話.
為什麼每個都要跑去美國.
叫自己的人民留下.
然後他的子女全部都有美國國籍.
或者有外國國籍.
這個就是他的意思.
有了羅馬公民的國籍.
或者有羅馬公民的護照.
是多麼大的榮譽.
或者那個身份是多麼的榮耀.
菲律賓作為一個殖民地的城市.
裡面的人民.
特別是羅馬的公民.
是多麼大的榮耀.
他們在推廣宗主國.
那種愛國.
那種羅馬的好處.
所以去到菲律賓這些城市的時候.
你就會看到.
他們有很多羅馬的好處.
羅馬生活的方式.
就是將經過羅馬的東西搬去菲律賓.
那裡有什麼呢.
有劇場.
這些全部都是羅馬的東西.

$^{321}$你看到那些建築物.
不只是建築物那麼簡單.
其實是一種生活的文化.
看到羅馬的市場.
很大的平面的市場.
那些叫做Roman Forum.
有街市 有市場.
有交易商.
然後上面有羅馬的廁所.
這些都是羅馬帝國的生活方式.
很榮耀的地方.
去推崇羅馬帝國的好處.
但是保羅說什麼呢.
保羅說我們就是天國的國民.
Heavenly Citizens.
我們是神國裡面的公民殖民.
我們的行事為人.
是要與基督的福音來相稱.
最重要的identity.
不是在地上的羅馬帝國的公民.
我們最重要的是天上給我們的.
一個公民的身份.
我們代表所推廣的.
不是地上的羅馬帝國.
不是地上的國家.
是天上的國家.
是屬於耶穌基督 屬於上帝的天國.
我們的行事為人.
每一天我們所做的公民.
不是單單做一個羅馬公民.
更重要的identity.
是做一個天國的公民.
是與基督的福音來相稱.
我想這樣說.
在菲律賓的教會的人.
他們立刻catch到那個意思.
在我們現代人看這些文字的時候.
未必catch到.
他們當時在政治的氣氛.
或者社會的光景裡面.

$^{361}$他們會特別catch到.
哦 我明白了.
原來我最重要的identity.
不是地上的政治identity.
而是天上 天國裡面.
給我一個新的identity.
我不知道你們有沒有這樣想過.
做基督徒.
是一個這樣的身份.
當你去到不同的地方.
你帶著一個天上 天國的殖民身份.
去到那裡.
或者現代的比喻可以怎麼說.
我想現代的比喻可以叫.
一個比較合適的比喻就是.
embassy.
就是國家的大使館.
美國在全世界各地.
有很多很多的embassy.
很多的大使館.
大使館是做什麼的呢.
大使館就是在別國的首都裡面.
的一塊地.
他在那裡工作.
他代表著他所猜派他的國家.
他要討月的.
不是他所在的地方.
的國家.
他討月的是猜派他國家的元首.
所以大使館是什麼呢.
就是要對自己本國來效忠.
縱使他所在的地方不是他自己的地方.
美國的大使被派去不同的國家.
不同的地方.
他就要去到慰藉美國的利益.
慰藉美國所謂美國帝國的權利.
來去效忠.
唱的國歌不是當地的國歌.
而是派他那個國家的國歌.
緊張的事情是他自己國家派他國家的政務.

$^{401}$不是當地國家的政務.
所以為什麼大使館這麼重要.
最近幾天前.
在伊拉克美國的大使館被攻擊.
攻擊大使館等於攻擊國家.
巴格達剛剛被攻擊的大使館.
美國的大使館.
是全世界所有美國的大使館最大的大使館.
伊拉克之戰之後.
加大了很多.
對美國來說是一個最重要的地方.
所以為什麼情勢這麼緊張.
每個人都說快要打第三次世界大戰.
因為攻擊大使館等於攻擊國家.
大使館就是民族之間的國家.
你就是被殖民到那個地方.
掛的國旗不是當地國家的國旗.
而是派駐你的國家的國旗.
你所推廣的價值.
也不是當地所在國的價值.
而是你宗主國的價值.
在異地他鄉見證另一個國度的存在.
另一個國度的美好.
我們基督徒就是有這樣的身份.
我們是天上的國民.
天上的殖民到我們這裡.
我們教會就是embassy.
就是天上天國神國的embassy.
為什麼我們行事為人要與基督的主相稱.
因為我們所代表所推廣.
所崇高的價值.
就是天國裡面有的價值.
我們所傳的福音就是說.
耶穌基督才是唯一的主.
不是地上的帝王.
任何一個領導人都不是我們的主.
只有耶穌基督是我們的主.
剛剛過年.
一段很大的新聞.
或者很多人傳閱的消息.

$^{441}$就是在中國四川省成都地下教會的.
王怡牧師的講話.
王怡牧師是領導在成都一間叫秋雨城若教會的牧者.
在他入獄之前.
他有一個這樣的聲明.
如果你可以的話.
你回去google一下.
去看看他整個的聲明.
他很長的.
我只在這裡揭露一點.
他這個聲明的題目叫做.
我的聲明信仰上的狂明.
這間教會被中共政權打壓了很久.
最近判了王怡牧師九年監禁.
判他的罪就是莫須有的罪.
煽動國家政權.
或者不正當的非法活動.
在這段很長的文字裡面.
這段文章裡面.
我想揭露幾段.
我認為最震撼的信仰講話.
他說我對福音的篤信.
對中人的教導.
和對一切罪惡的責備.
都是出於基督在福音裡面的命令.
在這裡帶領和呼籲任何願意悔改的人.
向上帝悔改.
這是教會在中國的一切工作的目的.
就是要向世界見證基督.
向中國見證天國.
向屬地短暫生活.
來見證天上永恆的生活.
這也是我本人.
夢照作牧者的呼召.
我同時相信.
中共政權對教會的迫害.
是極其邪惡的犯罪行為.
作為基督教的牧者.
我必須對這樣的罪惡.
發出嚴厲和公開的責備.

$^{481}$我所夢的呼召.
也是要求我以一種非暴力的形式.
在和平和忍耐裡面.
去違背.
那些違背了聖經和上帝一切的人間法律.
作為牧者.
我唯一關心的.
那是信仰上的抗命.
所帶來對罪惡人性的震動.
和對基督十字架的見證.
你有沒有想過我們的福音.
我們的信仰.
當我們口稱基督是我們的主的時候.
那種震動有多大.
當我們說Jesus is Lord的時候.
祂是唯一一個主.
是最大最大的主.
這個主會震動我們生命裡面.
社會裡面.
世界裡面.
任何一種權力的設定.
當耶穌來到的時候.
祂就是最大最大的主.
我們對十字架福音的見證.
就是見證著一種這樣的國度.
這個天國裡面的國度.
是叫我們來抗衡.
世界裡面任何一個帝國的國度.
要叫每一個人來悔改.
去歸在這個主之下.
我們每一年到年頭的時候.
我們就說會讀經.
請你不要覺得讀經只不過是一些.
很黑板.
很普通的熟能的練習.
因為我們拿到聖經的時候.
我們知道這個唯一的主是對我們說話.
而我們指定名於這位主.
我們指定名於這位對我們說話的神.
這位天國的主宰.

$^{521}$當我們拿到聖經.
我們說要讀聖經的時候.
這個天上的君主.
祂所說的.
是最終極.
最有權威的旨意.
我們應當去全心謙卑地去聽.
我們的行事為人.
應該跟著祂的心意去過一套.
不是地上的國度.
我不管你是哪一個國度.
當耶穌基督要說話的時候.
祂說的話就比任何一個領導人.
任何一個黨.
任何一個政權說的話都要大.
任何一個在生命裡面.
以為在持有我們.
在支配著我們.
的主人的說話.
耶穌基督的主權應該是最大的.
祂給我們的價值.
往往跟這個世界是相反的.
很多年前我去了一個基督徒的組織.
他們專門去中東一些戰地.
去做一些和平工作.
去訪問和探訪他們一些辦公室的時候.
我記得有一次去了他們的辦公室.
他們的門口的海報是這樣的.
你看不看到這是什麼.
是一個反轉了的世界.
平時我們看世界是反轉的.
我們基督的國度.
天上的國度.
我們看世界是反轉的.
什麼叫做榮耀.
就好像耶穌基督.
做僕人.
寧願死在十字架上.
這一種才叫榮耀.
世界不會這樣看的.

$^{561}$我們偏偏是反過來看.
世界看十字架是一個羞辱.
是一個愚蠢的記號.
但是我們看十字架.
是一個寶貴的十字架.
世界看救贖主是一個弱者.
但是我們說真正的強者.
就是一個不帶刀不帶槍.
願意以愛來犧牲.
在世界上十架上死去.
這個才值得我們敬拜.
值得我們去愛.
天上的公民.
這個就是我們最重要的身份.
因為如果你的身份.
真的形成了天國殖民公民的身份.
如果是這樣的話.
你的人際關係.
在社會裡面的角色.
你自己個人的操守.
將會有多大的震動.
原來信耶穌不是我有什麼好處.
這麼簡單.
我等上天堂.
就是上帝在天房裡將我植下來.
我要推崇的是天上的國度.
耶穌基督自己的榜樣.
我說菲勒比斯最忠心的就是這篇詩章.
上下文講到什麼?.
教會如何能夠同心.
教會如何能夠成為真真正正.
行事為人與基督的福音主權相承.
就是要看我們基督.
跟我們基督的榜樣.
我們今年會領第一次聖餐.
領在這個餅這個杯的時候.
你記得嗎?.
這個就是我們天國君主的榜樣.
我們的基督.
我們這個餅這個杯所紀念的.

$^{601}$救贖我們的主宰.
他就是這樣心.
這個就是天國的價值.
我們作為殖民,作為公民.
我們就是以基督的心為我們的心.
我們讀這條聖經.
然後我們才領受聖餐.
領受聖餐的時候.
我們不是兒兒戲戲.
輕輕率率的.
讓上帝的福音給我們心靈的震動.
我們虛心的跟隨耶穌基督的榜樣.
我們同心一起讀.
肥羅比書第二章,六節開始.
你們當以基督耶穌的心為心.
他本有神的形象.
不以自己與神同等為強奪的.
反倒虛己.
取了奴僕的形象.
成為人的樣式.
既有人的樣子.
就自己卑微.
全心信服.
以至於死.
且死在十字架上.
所以神將他升為至高.
有次給他那超乎萬名之上的名.
叫一切在天上的,地上的和地底下的.
因耶穌的名無不屈身.
無不口稱耶穌基督為主.
使榮耀歸於父神.
\newpage



\section{提摩太前書 2:1-24}
\label{sec:Pe1TNggJhos}
\textbf{主日講道: 屬靈生命的模式 提摩太前書2\_1-24 (吳君才牧師)}
\newline
\newline
連結: \href{https://youtube.com/watch?v=Pe1TNggJhos}{\texttt{https://youtube.com/watch?v=Pe1TNggJhos}} ~~~~ 語音日期: 2020-01-26
\newline
\newline
\hyperref[sec:6ypOdtEuNl8]{\small{< < < PREV SERMON < < <}}
~
\hyperref[sec:index_chronic]{\small{[返順時目]}}
~
\hyperref[sec:index_scriptual]{\small{[返順卷目]}}
~
\hyperref[sec:ez9_0uQY1v4]{\small{> > > NEXT SERMON > > >}}
\newline
\newline
提摩太前書 2:1-24
\newline
\begin{longtable}{cl}
\hline
\hline
章節 & 經文 (和合本修訂版)\\
\hline
2:1 & \begin{tabularx}{0.7\textwidth}{X} 所以，我勸你，首先要為人人祈求、禱告、代求、感謝； \end{tabularx} \\ \\
[1ex]
\hline
\hline
\end{longtable}
$^{1}$.
請姐妹平安.
2020年已經開始了.
但我又不知道你在來近這一年.
有什麼計劃幫助自己.
在愛神,愛人,愛教會這方面呢?.
究竟在聖經所說的關於一些屬靈生活.
或者生命的模式.
在今時今日是否行得通呢?.
今天我會以提摩太後書二章.
一到廿四節彼此分享.
提摩太後書.
就是小桃保羅自己知道自己命不久矣.
時日無多.
對屬靈的兒子提摩太.
一些肺腑的說話.
對一位年青的傳道人.
他體質又不好.
又幾次患病.
他胃口又不清.
所以小桃保羅提醒他.
你可以稍微用一些補酒調理一下.
其實提摩太是一個性格柔弱內向的人.
普通到像你和我一樣.
甚至按肉體條件來說.
他可能不合格.
不過神卻使用了他.
呼召了他.
保羅又委以重任給他.
鼓勵他內心要剛強起來.
依靠耶穌的基督恩典.
做一個剛強勇往直前好好事奉的人.
不怕挑戰.
神恩典夠你用.
保羅在他的書信裡面提出幾種生命的模式.
或者都是一個事奉的模式.
今時今日我們可以有一些指標.
我們看看經文.
1到24節.
實在是很長.

$^{41}$我跳來讀.
我們看好了.
你看1到5節.
再另一張.
看到第8節.
好,接著跳.
好,接著14到24.
再另一張.
另一張.
再另一張.
我一會兒會抽幾道經文再讀.
再講解.
我們先有個簡單的祈禱.
切開的父神.
雖然天氣外邊寒冷.
但求主給我們心裡火熱.
愛慕你的話語.
也愛慕和你親近.
更加愛慕聆聽你自己的真道.
一會兒時間我們分別為聖.
讓我們的思慮受煩擺低.
聽你自己聖靈引導我們的心意.
求你帶領引導.
奉主名求.
阿們.
第一個.
我們看到屬靈生命的模式.
或者路向或者指標.
下一張,謝謝.
就是要作無愧的工人.
我們讀一讀經文.
你當竭力在神面前得蒙喜悅.
作無愧的工人.
按著正義分解真理的道.
意思是要努力在神面前.
做一個經得起考驗.
問心無愧的工人.
而且要好好正確的講解上帝真理的信息.
你和我.
陳牧師上幾次都說.

$^{81}$我們是天國的公民.
我們是天國公民.
我們是神的子民.
我們是祂的僕人.
是作主的工人.
如果你是基督徒.
這個身份.
你走不掉的.
丟不下的.
不過你不單止做主的工人.
你要做一個無愧的工人.
這裡說.
你當竭力.
意思是盡力.
努力.
拼命去做.
而不是懶懶行.
死死氣地.
可惜今天.
基督徒中間.
有些現象.
就是寧寧懶懶行行.
死死氣地.
毫無生命活力和積極.
很多時候.
不能夠盡力.
努力拼力在神面前.
我曾經在街上.
遇見一個.
多年前.
很熱心追求信仰的青年人.
但是他自從出來社會做事之後.
整個人都改變了.
他的信仰經不起衝擊.
經不起考驗.
變得不順.
我和他聊天發現.
他變得很使賄.
為利是圖.
為錢至上.

$^{121}$他說.
現在趁我年輕.
最重要是要賺錢.
賺錢.
不理用什麼手段.
都無所謂.
我聽了之後.
心裡面都很寒.
很心鬱.
為什麼以前這麼積極.
這麼追求.
現在變成這樣.
他人生的生命立場.
是否經得起考驗.
是否問心無愧.
我們不說其他人.
我們說回我們自己.
究竟我們自己在神面前.
是否經得起考驗.
究竟我們是否得到神的喜悅.
我行事為人.
是否問心無愧.
在這裡.
經文甚至.
看以後的章節.
是包括你.
說話都能夠無愧.
因為史圖保羅在第15節.
下半節.
接著就提及.
要按著正義.
分解真理的道.
是要回應第14節.
所提及的.
那些什麼.
不要為言詞爭辯.
和第16節所提及的.
虛談.
偏離真道.
敗壞人心的言談.

$^{161}$一個無愧的工人.
一個屬神的人.
是要留意這些.
就是說.
說話.
是要謹慎.
要分別為性.
要有靈有性的.
可以建立人.
和去分解真理的道.
簡單來說.
就是說.
要正確的講解.
甚至活出.
上帝真理的信息.
難度.
頗高的.
很多年前.
不知道你有沒有聽過.
發生一件事.
在香港.
錫安教會.
梁逸華牧師的.
雙仰水事件.
他是不以醫學病理的專業角度.
只是穿鑿附會.
聖經中所說的活水.
其實是指雨水.
是有礦物質的液體.
這類的東西.
所以雙仰水.
他認為是可以內服的.
他牧羊教會那間挺大的.
一千個會友.
結果有七成的人聽他說.
真的喝雙仰水.
結果部分的人出現問題.
腸臟出現毛病.
這是否按正義.
分解真理呢?.

$^{201}$分解神的道.
並不是要分割.
拆散神的話語.
而是好像在田園中間.
路途裡面開一條小路.
引領人直達到神的面前.
其實唯有神真正的純正的話語.
才是生命豐盛的傳言.
當然你說看聖經.
看神的真理.
可以有個人的理解.
領受不同.
是可以不同.
但你仍然需要.
持守那個識經的原則.
當你和我向人分享神的道的時候.
或者栽培人的時候.
甚至你出去報道.
和人用聖經來說耶穌.
除了依靠聖靈和信服聖靈之外.
你都要了解那個受眾.
那個對象的情況和需要.
用神的道可以透過你的講解.
用神的話語.
令他明白有適切的活用.
和傳遞在他們的宗教.
引領他們去到神面前.
得著更多的恩典.
這樣做.
你問心無愧.
這樣做.
你真是一個無愧的工人.
是合乎神的心意.
按正義分解神的道.
不過.
和這裡有一個很重要的關鍵問題.
就是問.
我們自己究竟對讀經.
對查經.
對靈修掌握有多少.

$^{241}$肯不肯付代價.
花時間去操練去看呢?.
假如你和我對聖經明白一知半解.
或者蒙查查.
怎能夠很有系統.
很有條理的和人分享交流呢?.
又可以講解得清清楚楚呢?.
求主幫助我們.
先在基本讀經的生活.
在你靈修的生活.
在你想主要學探索課程的生活.
甚至在聽道的態度.
有改變.
盼望這個也是你2020年.
新的計劃.
能夠有更新.
立志靠主去實行.
我們要看第二樣.
要成為潔淨的器皿.
我們一起讀.
人若自潔脫離卑賤的事.
就必作貴重的器皿.
成為聖潔.
合乎主用.
預備行各樣的善事.
意思是說.
如果你和我自潔脫離一切邪惡的事.
就能夠給主人器重.
因為那些已經獻給主人的人.
是主人所縱容的.
預備行各樣的善事.
大戶的人家好比神的家.
是指那些可見的教會.
或者是事奉神,屬神的人.
器皿主要是指那些屬神.
或者是事奉神的人.
要作貴重的金銀器皿.
還是做一些較卑賤用途的木器,瓦器.
是不同的.
這裡代表了兩種事奉的靈明的本質和生命表現.

$^{281}$弟兄姊妹.
你覺得你在成會教會.
或者你的生命.
你的事奉本質在神眼裡.
是不是屬於寶貴的金器,銀器這些器皿呢?.
還是只是普普通通.
甚至卑賤的瓦器,木器呢?.
每一個屬神的人.
或者事奉神的人.
都可以是很尊貴的.
也都可能可以很卑賤的.
主要有一個關鍵.
那個關鍵就是.
人若自覺脫離卑賤的事.
就可以必定作貴重的器皿.
結果有三方面.
第一,成為聖傑.
第二,合乎主容.
第三,可以預備行各樣的善事.
怎樣自結呢?.
很多人問這個問題.
波羅開宗明義說了一樣東西.
你第一樣要逃避少年的私慾.
提摩太,這很重要.
不單止當時的提摩太要注意.
今時今日每一個屬神的人.
或事奉神的人.
都是要逃避私慾.
我們看教會歷史.
看現代一些教會的發展.
有些很明目.
事奉神的領袖.
都因為私慾的問題.
無論錢,財,色,利.
被打垮.
有些人說少年人.
少年人是否像現在十多歲的年輕人.
都要時刻稱少年.
其實如果看當時的背景.
希臘人的年齡.

$^{321}$他只分兩類.
就是老年人和少年人.
40歲以上的就是老年人.
40歲以下的一概為少年人.
哇,用這樣的分法.
在座坐在那裡的那些.
40歲那些.
這樣你就已經變成老年人了.
其實提摩太不是十多歲的少年人.
而是像今天.
所在你附近的三十多歲的壯年.
少年的私慾.
不單止是情慾,私慾.
邪情,私慾.
其實也包括很多和德行有關的表現.
所以,小圖保羅接著說.
要同那清心禱告主的人.
追求公義,信德,仁愛,和平.
很明顯的.
保羅勸提摩太.
要除了逃避這些情慾,私慾的事.
那些不合神心意的事之外.
不要被人家克制之外.
你要逃避或者甚至除掉.
不單止是屬於這些.
還要包括聖經所說的.
其他的自私,淫恨,淫辭.
妄語,名利,虛諷.
甚至那些驕傲,倔強,苦毒.
怨恨,嫉妒這些毛病.
因為這些東西.
是和公義,信德,仁愛,和平有違背的.
一個不潔淨的器皿.
是充滿了污穢,淫恨.
和人的劣根性.
神不會容許.
我記得很多年前.
在香港吃蛇羹.
都很好.
誰知道吃到差不多.

$^{361}$發現了一隻金金的烏鷹.
嘩!你吃了那些東西都要吐出來.
於是跟它說.
它就換了另一些出來.
我心想會不會又是同一個源頭呢?.
同樣道理.
你想想.
倘若你裡面有些污穢.
有些邪情,私慾的事.
你居心不良的走來事奉.
可以說.
好像一個污穢的管子.
必定是堵塞和敗壞了神的話語.
是失去了豐盛的生命力.
結果他的事奉.
往往變成靠人的血氣.
或者靠肉體而行事.
2020年開始.
很多團契或者小組有新的職員.
我們盼望這些事奉的人.
先求主的寶血.
再一次潔淨我們.
每一個事奉人.
以致你和我可以成為聖潔合用的器皿.
不過我又要提醒弟兄姊妹.
人和事奉最大的假想敵是誰?.
其實就是你自己.
就是自己.
你裡面的自我.
所以人要攻克己身.
叫身服我.
很嚴格和問自己有沒有罪.
有沒有劣根性.
所以私人祈禱神啊.
求你監察我.
看我裡面有什麼惡行沒有.
不單止「吾日三省吾身」.
隨時都有對聖靈的同在和罪的敏感.
一個人我想你聽過的.
那個人就是泰臣.

$^{401}$他19歲很厲害.
很年輕成為拳王.
但是他的童年很淒慘.
他1966年生在紐約卜倫.
那個黑人的貧民區.
出生不久就被父親遺棄.
不久媽媽年幼的時候都死了.
就好像一個浪子孤兒在掙扎.
天生天養.
12歲因為偷竊.
被送入少年犯宿舍.
當時宿舍內有一個當過拳手的守衛.
教少年犯打拳.
你會打拳可以宣洩你的淚氣.
他學會了之後.
拿拳的技能經常去傷人.
他長得很大.
甚至侵犯女性.
結果在1992年.
他因為強暴一名美國黑人選美小姐.
而入獄判刑6年.
但後期在獄中因為行為良好.
3年半之後就被釋放出來.
在獄中當時他遇到一名回教的導師.
經常和他聊天關心.
引導到泰臣成為回教的跟隨者.
出獄後泰臣重返拳壇.
他打得很奮力.
爭取名聲.
打垮很多拳手成為拳王.
但經過一段日子之後.
他又因為涉嫌騷擾女侍應.
這樣被警方審查.
更因為藏毒或其他刑事原因.
又被拘捕.
兜兜轉轉的一個人生.
泰臣不能自結.
不能約束自己.
不能攻克己身.
結果好事不出門.

$^{441}$惡事纏你.
不單止泰臣.
其實你和我每個人.
我們作基督徒的.
都是要小心.
人一旦本質沒有好好依靠神.
徹底的改變.
即是你回教會多久.
即是你事奉綱位.
做得好像很有生氣.
有聲有色.
但你裡面的生命未真正的更新.
難怪你裡面仍然有不潔不義的殘餘.
經常會蠢蠢欲動.
絲肉懷胎.
嘔嘔嘔.
唯一的出路.
就是養賴主耶穌基督的寶血.
塗抹我們.
潔淨我們.
再一次靠主成為一個潔淨.
聖潔的器皿.
才可以好好地被神使用.
但願大兄姐妹.
你靠主的心智.
就要預備潔淨自己.
成為一個乾淨.
聖潔的器皿.
繼續在2020年.
好好地來事奉神.
恭克己身.
第三樣.
屬靈領袖或屬靈生命模式的路向.
第三樣.
有僕人的心態.
第24節我們一起讀.
「也如主的僕人不可增敬.
只要溫溫和和地待眾人.
善於教導.
全心忍耐」.

$^{481}$這句金句是我事奉的座右銘.
幫助我很大.
不是跟人爭.
你事奉好像出來經常跟人爭.
擺權威擺款不溫和.
待人好像不知道什麼.
不是.
你要善於教導.
全心忍耐.
裡面有很多經歷.
這裡說到僕人心態.
提到不要增敬.
要溫和.
善於教導.
又忍耐.
馬丁路德說到基督徒的自由.
曾經說.
基督徒全然是自由的眾人之主.
不受任何人管轄.
但是基督徒.
亦都全然信服眾人之僕.
受任何人管轄.
很弔詭的說話.
其實保羅也有類似的說話.
保羅說.
「我雖是自由.
然而我甘心作了眾人的奴僕」.
意思一樣.
但是你看到現實卻是另一回事.
景象就是.
人人都想坐高位.
個個都想表現自己.
想出風頭.
想做大哥.
或者做大家姐.
想說事.
多少時候想做領導人.
想指揮人.
就不想被人指.
不肯屈就.

$^{521}$因為沒有什麼話事權.
不爽的.
今天如果我們盼望成為一個好的領袖.
就必須學一個功課.
要做一個好的僕人.
作為起點.
放下尊嚴.
或者男子或者女子的氣概.
能夠好好地投入.
積極地與人配搭.
以及對你當前的肢體生活.
教會事奉.
能夠有多一些正面的看法.
這個可以作為起點.
英文有一句說話.
A good leader is a good follower.
意思就是說.
好的領袖原來是一個好的跟隨者.
你不能夠做一個好的跟隨者.
你很難成為一個好的領袖.
因為你的自我驕傲.
心態阻塞了你.
你能夠學習去互相欣賞.
互相樂意與人配搭.
你又謙卑.
能夠又包容地去服侍.
那是基於你學校主耶穌的榜樣.
你能夠知道主耶穌的心意.
以主耶穌基督的心為心.
當你面對一個不肯卑微做事的人.
不屑做一些很瑣碎的事的時候.
這個人你可以說.
他的自我都很強.
他只顧著出風頭.
根本只是好高騖遠.
好像很有野心.
但是他屬靈的基本功都做不好.
很難有很大的作為和發展.
怎樣操練這樣呢?.
就需要有一個僕人的心態.

$^{561}$不妨時刻提醒自己.
你叻極點叻都是主的僕人.
你要守住一個心態.
就要不斷地學習.
和氣待人.
能夠包容,能夠接納人.
能夠好像這裡說的很殷勤的.
善於教導人.
保羅更挑戰我們.
要以溫柔勸誡一些敵對你的人.
我們做教務的曾經面對一些很難搞的人.
但是經驗告訴我們.
你不用和他鬥.
你繼續為他禱告.
你好好的恆久忍耐.
始終神會幫助他醒悟過來.
你這樣的表現反而建立了他.
今天我們每一位事奉主的人.
我們作為土主喜悅的僕人的弟兄姊妹.
我們要明白原來一切的果效.
事奉的成效.
完全不在你手中.
是在神手裡面.
始終人算不得甚麼.
我們都是主的僕人.
主的侍者.
主的工人.
唯有自己多一點依靠主.
聽命於主.
信服主的安排和引導.
避免太小器.
太執著.
太意氣用事.
如果應用僕人的心態在你日常生活裡面.
很多東西可以說.
不過我提出其中一樣.
就是你在家裡當中.
家庭裡面的表現.
怎樣好好運用出來.
因為有些人在家裡面.

$^{601}$他是大爺來的.
他是女王來的.
神聖不可侵犯.
很多時候可能翹起雙手.
不是怎樣做事.
反而有諸多挑剔,意見要求.
又不肯謙卑.
不肯參與家務.
關顧家人等等的付出.
有專家說.
家庭主婦在家裡有三分時間忙於家務.
而作為丈夫的.
有半數以上.
每個星期做家務的時間.
只有五個小時或以下.
而百分之十五的男士.
更是不動一根指頭.
專家這樣說的.
不知道在座的你屬於哪一類.
或者你是另類.
還有我曾經看過一本楊木菊所寫的書.
《廚房哲學》這本書很厚.
裡面講很多人生的大道理.
但在生活小節上顯示出來.
其中提到中國男人.
始終就算這個世代.
仍然有男尊女卑的思想.
雖然我們說有平權的觀念.
一切平等.
但這種舊有的觀念.
影響了一些男士做家務的問題.
我想這個也跟獨人心態有關.
楊木菊牧師指出.
一個不懂家務或懂得家務.
卻不肯做家務的男人.
是個有殘障的人.
聽到我吧.
他說不管對他的伴侶.
他的家庭或自己也是一樣.
他說日本新的這一代.

$^{641}$不再受上一代的男尊女卑的思想影響.
事家務會失去了男子氣概.
與所遂的事都不想做.
但今日日本的年青爸爸.
精通帶小孩.
換尿片很熟手.
他不介意自己妻子的地位.
所從事的職位比他高.
他的工資比他高.
他也不介意.
兩天前.
日本環境大臣.
小柱晉次郎.
人們說他是明日政壇之星.
他是高官.
但他為了一件事.
他太太懷孕要生小孩.
他竟然.
以前的保守日本人不會去做.
他特別申請放育兒假.
育兒假.
即是帶小孩假的意思.
兩個星期.
他說我要陪太太生小孩.
又陪兒子.
不過如果有些內閣會議他會出席.
甚至有些視像會議要處理公務.
他願意.
他這個行徑在日本電視.
國內得到很多人的議論.
甚至他自己說.
我這件事引致國際有些讚賞.
有些人說他.
這樣是令到保守的日本帶來改變.
不過.
今日日本年青男士思想心態有改變.
反觀我們華人.
心態.
就算華人基督徒.
心態.

$^{681}$都是很閉塞.
仍然.
固步自封.
根深蒂固的觀念不是一朝一夕可以除.
心態問題.
究竟你有沒有一個僕人的心態呢.
你那種.
放下自己.
不理那個面子.
那種思想.
是不是仍然都是一個阻打.
很大的阻力呢.
各位姐妹.
無論在家.
是你工作的場所.
甚至在教會.
事奉裡面.
僕人的心態.
是我們.
繼續要學習.
和追求.
2020年.
我但願你彼此提醒.
我們要有僕人的心態去事奉.
其實我有六樣的模式.
看一個鐘時間都差不多.
還有很多東西要說.
太喜歡講道德了.
留待下集.
要停留下來.
今天.
使徒保羅提醒我們.
三樣東西要記住.
看回前面.
第一.
讀個大標題.
一.
要作無愧的工人.
第二.
要成為潔淨的器皿.

$^{721}$第三.
要有僕人的心態.
是很重要.
但願你能夠自己在神面前成為一個潔淨的器皿.
脫離卑賤和私慾的事.
你能夠以僕人的心態.
不憎敬.
不嫉妒.
很溫柔忍耐的謙卑苦咒.
好好地和人配搭.
無論在家庭.
教會裡面.
在工作場所.
都能夠榮耀神.
向我們低頭祈禱.
神你自己與我們同在.
加能賜力給我們.
讓我們在你面前.
能夠成為一個貴重的器皿.
潔淨的器皿.
更加在你面前.
能夠作為一個無愧的工人.
好好地將來交帳的時候.
都無愧.
今日.
事奉.
對人.
對神都無愧.
我們更求主.
你讓我們有一個僕人的心態.
不憎敬.
不嫉妒.
是溫柔忍耐的謙卑苦咒.
懂得和人配搭.
使教會.
家庭.
和工作.
能夠見證榮耀.
願你加能賜力.
奉主耶穌基督的名求.

$^{761}$阿們.
好.
天父祝福我們每一個人.
做一個聽道又行道的人.
我們默討之後可以散會.
\newpage



\section{利未記 13:1-6}
\label{sec:ez9_0uQY1v4}
\textbf{2020/2/2 主日講道: 疫境前行 (利未記13\_1-6) - 陳崇基牧師}
\newline
\newline
連結: \href{https://youtube.com/watch?v=ez9_0uQY1v4}{\texttt{https://youtube.com/watch?v=ez9\_0uQY1v4}} ~~~~ 語音日期: 2020-02-06
\newline
\newline
\hyperref[sec:Pe1TNggJhos]{\small{< < < PREV SERMON < < <}}
~
\hyperref[sec:index_chronic]{\small{[返順時目]}}
~
\hyperref[sec:index_scriptual]{\small{[返順卷目]}}
~
\hyperref[sec:imYjaG_VAbM]{\small{> > > NEXT SERMON > > >}}
\newline
\newline
利未記 13:1-6
\newline
\begin{longtable}{cl}
\hline
\hline
章節 & 經文 (和合本修訂版)\\
\hline
13:1 & \begin{tabularx}{0.7\textwidth}{X} 耶和華吩咐摩西和亞倫說： \end{tabularx} \\ \\ \relax
13:2 & \begin{tabularx}{0.7\textwidth}{X} 「人身上的皮膚若腫脹，或發疹，或有斑點，可能成為痲瘋 的災病，就要把他帶到亞倫祭司或亞倫的一個作祭司的子孫那裡。 \end{tabularx} \\ \\ \relax
13:3 & \begin{tabularx}{0.7\textwidth}{X} 祭司要檢查他身上皮膚的患處，若患處的毛已經變白，災病的現象深入身上皮膚內，這就是痲瘋的災病。祭司檢查後，要宣佈他為不潔淨。 \end{tabularx} \\ \\ \relax
13:4 & \begin{tabularx}{0.7\textwidth}{X} 若這人身上的皮膚有白斑，看起來並沒有深入皮膚內，其上的毛也沒有變白，祭司就要將這病人隔離七天。 \end{tabularx} \\ \\ \relax
13:5 & \begin{tabularx}{0.7\textwidth}{X} 第七天，祭司要檢查他，看哪，若災病在祭司眼前止住了，沒有在皮膚上擴散，要將他再隔離七天。 \end{tabularx} \\ \\ \relax
13:6 & \begin{tabularx}{0.7\textwidth}{X} 到了第七天，祭司要再檢查他。看哪，若災病減輕，沒有在皮膚上擴散，祭司就要宣佈他為潔淨，因為他患的不過是疹子。那人要洗自己的衣服，就潔淨了。 \end{tabularx} \\ \\
[1ex]
\hline
\hline
\end{longtable}
$^{1}$平安啊.
你知道大家能夠聚會的時候.
我們很多時候講平安啊.
這個其實是初期教會的一個習慣.
因為初期教會,第一世紀教會在羅馬帝國的打壓之下.
很難聚會的.
能夠一件 one piece(一件)來聚會的時候.
就是很平安地去聚會的.
所以他們講Shalom(祝福).
講平安的時候.
我們就說平安啊.
平安是一件很重要的事情.
因為平安是一個信仰的信念.
所以他們講Shalom(祝福).
講平安.
是一件很不簡單的事情.
今天也是一樣.
在這樣的疫情裡面.
大家應該很熟悉了.
那些新聞是不是?.
應該看新聞看得比聖經多是不是?.
是不是?.
現在真的太多太多這些資訊了.
所以.
剛剛提醒了大家.
我們的組群.
我們盡量給一些.
又不會爆炸性的給你.
但是盡量給一些.
是對的 fact check(事實查證)了的一些.
的程勢給你.
讓你夠知道.
那麼這兩個星期.
我們不可以不講一些跟疫症有關的事情.
是不是?.
因為我們是要貼地一點.
以及要知道社會裡面我們怎麼去回應.
我也嘗試怎麼去看聖經.
聖經裡面怎麼去教導我們.
你知道聖經裡面有一本書叫做.

$^{41}$《利未記》.
Leviticus(利未記).
《利未記》是摩西五經的律法書裡面.
中間那一本.
那裡講到很多我們怎麼去.
一些的律例.
是舊約神的子民的一些律例.
譬如教導他們怎麼去敬拜.
一些講到梵祭,平安祭等等.
以前的敬拜的方法跟今天是不同的.
因為耶穌基督來了的原因.
所以他們要怎麼去獻這些牛,羊.
是完美的祭物.
他們要獻上.
教導他們怎麼去親近和敬拜神.
也有一些祭禮或者一些律法.
講到祭司要做什麼的.
更加很多時候.
《利未記》.
是講到一個神國的子民.
就是舊約裡面的教會.
他們應該大家怎麼去相處.
以致這個社群.
這個 community(社群).
是有神自己同在的.
那個處境就是當以色列人出了埃及之後.
他們在曠野這麼多年.
就是上帝接著摩西來賜給他們一些律例,一些條例.
以致能夠去 regulate(管理).
以致他們的 community(社群).
是有神與他們同在的.
那麼我就選了一段聖經.
我們一起來反省一下.
或者去掙扎一下.
到底這段聖經.
對我們今天.
是不是完全無關的呢?.
如果是有關的話.
是哪一方面是有關的呢?.
如果有聖經的話.

$^{81}$你不妨打開去到《利未記》第十三章.
《利未記》第十三章.
這裡是講到跟一些傳染病有關的.
所以我就選了這段聖經.
十三章第一節.
耶和華對摩西和亞倫說.
如果有人在皮肉上生了瘤.
或腺.
或火斑.
後來皮肉變成了大麻風病症.
就要把他帶到一個作祭司的亞倫子孫那裡.
祭司要察看皮肉上面的病症.
如果患處的毛已經變白.
病症的現象.
又滲透皮下的肉.
就是大麻風病症.
祭司發言者就要宣佈他為不潔症.
第四節.
火斑若是在他的皮肉上發白.
但沒有滲透皮膚的表象.
上面的毛也沒有變白.
祭司就要把患者隔離七天.
到了第七天.
祭司要察看他.
如果看見他病情止住了.
皮上的患處沒有蔓延.
祭司就要把他再隔離七天.
到了第七天.
祭司要再察看他.
如果看見患處色澤變淡.
也沒有在皮上蔓延.
祭司就要宣佈他為潔症.
這不過是癬.
他洗淨衣服.
就得潔症了.
那麼.
這些聖經.
是我們查經不會查的.
這些也是我們平時自己讀經.
是不會感覺到很安慰或者很溫馨的.

$^{121}$這些條例你覺得好像很沒用.
那麼.
首先我們要梳理一下這些問題.
就是看這段聖經的時候.
他說大麻風.
應該不是今天現代人所說的Hansen's disease(哈根斯病).
因為我們知道在近東.
巴勒斯坦地區.
這個時候是沒有Hansen's disease(哈根斯病)的.
就是說真正我們今天醫學所說的大麻風病.
大概大麻風病這個字.
是一個很籠統很廣泛的字.
其實是形容一些會蔓延.
或者會傳播會傳染的一些問題.
一些病症.
所以可以在皮膚上長出來的.
也可以等一會兒我們看另外一些經文.
就知道.
是在牆上也可以有的.
就是石牆上也可以有的.
所以歷做大麻風.
大家就知道不是現代人所說的那種麻風.
還有為什麼這麼多的病症.
偏偏這些的律例又特別要針對.
特別要針對特別要說.
這種所謂的大麻風.
或者皮膚的病.
我想大概的原因就很簡單.
因為當時還沒有現代的醫學.
是一個很古代的一個時候.
所有的病症看得到的.
就是最明顯的.
所以有什麼病是最明顯的呢?.
你說.
暗病是不明顯的嘛.
講明是暗病.
是不是?.
所謂明病就一定是皮膚的病.
所以很多時候他們 deal with(處理)的.
就是一些比較表面.

$^{161}$現象.
phenomena(現象).
observable(可觀察).
就是可以觀察到的一些病症.
所以這一段聖經.
十三十四章.
都是處理這種病症的.
另外呢.
要處理梳理的一個問題.
就是潔淨與不潔淨.
我們一提起潔淨與不潔淨.
其實也很吵架的.
因為現代人呢.
我們覺得潔淨與不潔淨.
就是衛生的問題.
Right?.
就是你潔不潔淨.
乾不乾淨.
我們就想起.
hygiene(衛生).
但是呢.
我們知道.
在舊約.
特別是在舊約的社會裡面.
講到潔淨與不潔淨.
就不是講到 hygiene(衛生)的問題.
沒錯.
可能hygiene(衛生)的問題有部分.
但是大致上不是講到 hygiene(衛生)的問題.
其實是講到什麼呢?.
是講到 order(命令)的問題.
潔淨的東西.
就代表.
他這件東西.
或者這件事情.
是應該在他現在的位置的.
這個就為潔淨了.
Pure(純).
Clean(乾淨).
如果有一件事情發生.

$^{201}$無論是人事或物.
發生在一些不應該有的位置.
那些就叫做不潔淨了.
Unclean(不乾淨).
OK?.
所以就算是食物也是的.
下個星期我們就會講食物了.
食物的條例.
潔淨與不潔淨是在講 order(命令)的問題.
如果那件事情不應該在那裡出現.
而出現了的話.
那就是不潔淨了.
是不是?.
因為是 disorder(不潔淨).
就是擺錯了地方.
因為舊約神學的觀念.
就是說上帝有一個 creational order(創造命令).
這個 order(命令)是一定要 maintain(維持)的.
一定要維繫著這個 order(命令).
如果這個 order(命令)錯了的話.
那就會使得.
無論是社群.
或者是世界.
就會有秩序上的混亂.
所以講到潔淨與不潔淨.
不是講衛生.
而是講.
是 order(命令).
或者是 disorder(不潔淨).
其實這些每個人的社會.
人類的社會歷世歷代到今天.
都是有這樣的觀念的.
order(命令)和 disorder(不潔淨).
就是說有些事情不應該在那裡發生.
但是又發生了.
那這個就不潔淨了.
或者是每件事情應該有條有理的.
我問你.
你看到一雙鞋子.
放在地上的時候.

$^{241}$乾不乾淨?.
OK嗎?.
是OK的.
是 order(命令)的.
但是如果吃飯的時候.
有一雙鞋子放在飯桌上.
那你覺得是不是 disorder(不潔淨)呢?.
是不是?.
但是如果我想問.
如果那雙鞋子是新的鞋子呢?.
剛剛在 mall(商場)買的.
是沒穿過的.
新得不得了.
可能比你的手乾淨.
你覺得.
是 disorder(不潔淨)還是 order(命令)呢?.
我們都會覺得是 disorder(不潔淨)的.
因為鞋子應該在哪裡?.
鞋子應該在地上.
這就是我們每一個文化.
我們認為什麼叫做 dirt(髒).
什麼叫做骯髒的東西.
或者 disorder(不潔淨)的東西.
所以鞋子放在地上.
就沒有人會去投訴的.
但是.
鞋子如果放在你枕頭邊.
那你就覺得是.
不是很對的.
是不是?.
那麼這個是物件.
上面有 disorder(不潔淨)和 order(命令)的分別.
我們每一個文化都會有的.
或者這麼說.
我們做牧者的.
可能我也提過了.
我們有時候在 weekend(週末).
我們會辦喪禮的.
辦完喪禮.
喪禮很多時候都是白天的.

$^{281}$很少晚上喪禮的.
我們上完山.
很多時候同一天.
就會有婚宴舉行.
就是另外一家會有婚宴舉行.
可能也是我負責的.
那麼我又會去那個婚宴.
是不是?.
因為.
其實我們很方便.
我們不用換衣服的.
因為我們穿什麼都是黑色的.
所以我們去到婚禮.
去到婚宴 banquet(晚宴)的時候.
我就會坐下來跟大家有紅事.
就是有喜事的.
那麼.
那麼你覺得是 order(命令)還是 disorder(不順)呢?.
我記得我也曾經有這樣的情況.
但是我不記得是誰.
我不記得是誰.
但是我去到那個婚禮的時候.
我就不會跟對方說.
或者跟那兩家說.
哦!剛剛去完喪禮.
我剛剛來.
講到是不是?.
我當然不講出來是不是?.
是不是?.
我不會講出來.
因為這種是所謂 disorder(不順).
是不是?.
我們呢.
就是你知道中文.
廣東話是很傳神的.
這些是廣東話才能夠表達的.
不老禮是什麼意思呢?.
不是因為迷信的問題.
也不是因為衛生的問題.
是不是?.

$^{321}$我很乾淨的.
我去完喪禮.
我又跟著去你那個婚禮.
是沒有問題的.
那麼.
但是總之我們就覺得是不老禮了是不是?.
所以.
不好意思.
今天新年嘍嘍.
講這些東西也是不太老禮是不是?.
但是我想解釋梳理的一個問題就是.
什麼叫 purity(純粹).
什麼叫 impurity(不純粹).
在舊約裡面往往都是在講一個.
創造界社會裡面.
世界裡面.
一些 order(規矩)的問題.
一些規律的問題.
一些的.
秩序的問題.
所以講到這段經文的時候.
往往不是在講罪.
犯罪.
所以有病的問題.
不是因為犯罪所以有疫症的問題.
往往講起這些事情的時候.
就是講到 order(規矩)和 disorder(不順).
當然我們要明白.
這本聖經.
是一本古代的文獻.
它不是呢.
正如它不是現代的科學.
你不要拿著這本聖經來探索太空的奧秘.
聖經主要目的不是講那件事情.
你更加不應該拿著這本聖經.
當它是一個醫學的百科全書.
因為那裡不是講治病的.
沒有什麼所謂現代醫學的知識.
但是這麼說.
他們以前這些文獻.

$^{361}$來處理這些疫症的時候.
或者處理不同的生活習慣模式的時候.
都是有一些間接的教導.
這個就是我的出發點.
正如有些人說.
現在經常研究.
一個星期有一天是安息日.
你去敬拜.
去休息.
六天的工.
你做六天.
然後第七天你安息.
你休息.
有很多心理學家或者一些醫生.
都說.
對人的身心靈.
都是有健康.
會給你康泰的.
應該是這麼做的.
當然你去聖經裡面找.
為什麼會有安息日的條例呢?.
他不是講這一點的.
他不是說因為你想在身心靈上有益處.
所以你要守安息日.
不是.
而是你需要去敬拜這位創造天地萬物的主.
那是神學的原因.
不是醫學的原因.
但是確實是有一些好處的.
有些人去研究說.
祈禱是讓你心靈裡面得到醫治或者得到安慰的.
也是一樣的.
聖經沒有講過.
只要你祈禱.
你就會得到不同的好處的.
很多時候都不會講一些醫學的好處.
但是我們要知道.
就是那個的.
律法的條例.
很多時候不是照辦主管下到今天.

$^{401}$應用到在現在的社會的.
我們一定要通過某一些的.
所謂方法論.
我們才能夠明白.
那個律法到底想講些什麼.
Right(對).
我舉個例子.
有一個的律法的條例.
是講到說如果你建造新的房屋的話.
你就要在房頂上面弄欄杆.
免得有人在屋頂那裡掉下來.
是不是?.
然後流血的罪.
就會歸到你家了.
為什麼會這麼講呢?.
芝加哥的屋頂不是平的.
大部分都是斜的.
但是巴勒斯坦地的屋子.
都是平的.
因為上面可以儲水.
或者上面可以曬乾一些農作物.
或者可以做不同耕種的一些地方.
所以他們很多時候有平頂的房屋.
大概兩層樓高.
兩層樓高.
所以很多時候就會上到屋頂.
所以律法書就說你弄房子的時候.
一定要弄好欄杆.
要有 fence(牆).
當然你不可以說照辦主管今天.
我們就說我們屋頂就要弄欄杆了.
是不是?.
因為那個社會的處境.
建築的處境都完全不同了.
是不是?.
但是他背後有些什麼原則.
有些什麼價值呢?.
有什麼宣告呢?.
就是說生命是神聖的.
所以我們要做任何事情.

$^{441}$來保障著無辜的生命.
以至他不會死亡.
所以我們要問心一層.
問的不是照辦主管下來這裡怎麼應用.
而是問問他背後到底是什麼價值.
是屬於上帝國度的價值.
屬於神自己的價值.
所以看到這些病的時候.
我們也是需要這樣子去 approach(接觸)的.
其實這些脈絡也很簡單.
這裡講到說.
如果有人生病的話.
皮膚上面有一些惡性的傳染病的話.
就需要去找一個祭司.
然後祭司就會去做一些的診斷.
去觀察.
如果真的生病的話.
那就算為不潔症了.
Disorder(不潔症).
然後呢.
就將那個人.
病人就需要去隔離了.
然後到某一個時間.
如果他真的康復的話.
就會成為潔症.
就是 orderly(順序).
他就回來了.
這件事.
這裡講到一個社會.
一個神國.
社群.
他們要建立的時候.
要維繫的時候.
要做些什麼呢?.
有生病.
很多時候經過診斷.
他們就將有傳染病的人.
就將他去 quarantine(隔離).
今時今日這個字應該不是很陌生.
什麼叫做隔離.

$^{481}$什麼叫做 quarantine(隔離).
In fact(事實上)這個字.
其實是在中世紀裡面引申出來的.
Quarantine(隔離)這個字本身是 Italian(意大利語).
是意大利文.
當時其中一個的自圓的原因.
就是在當時有很多時候的疫症.
譬如一些 Black Death(黑死病).
黑死病的疫症.
在蔓延歐洲的時候.
那些歐洲人.
特別是意大利人.
就觀察到.
在他們城市裡面.
猶太人的族裔.
他們的群體.
他們染病的機會率.
是特別小一點的.
所以他們就很好奇.
到底為什麼他們會這樣子呢?.
然後就發現.
他們的猶太人.
因為守住五經的律例.
他們經常會做一些洗手.
或者隔離.
一些事情.
有很清楚的分結症與不結症.
所以他們就明白.
原來他們有一個這樣的原則.
Quarantine(隔離).
其實自根的意思就是.
四十的意思.
就是四十天.
因為有一些這麼摩西的律例說.
要去隔離一個病人的時候.
就要隔離他四十天.
隔離的原因很簡單.
就是不會傳染到其他人.
不會影響到整個的社群.
以致整個社群會不結症.

$^{521}$隨你怎麼看.
有很多現代的醫學倫理學者.
一些的 medical ethicist(醫學學者).
他們去研究這些的時候.
他們甚至說.
這些摩西五經的隔離的律例.
是公共衛生法案的第一個 model(模式).
Again.
我不是說舊約裡面說.
是一些 public hygiene(公共衛生)的問題.
他們只不過是由一個創造的神學.
以致他們能夠分辨到有些什麼叫 order(命令).
有些什麼叫 disorder(不和).
但是你要想想.
那時候寫的時候.
或者描述的社會的狀況.
是公元前一千多年的時候.
是.
距離我們現在.
大概有三千多年.
是.
有些人甚至說.
當時的條例.
現在這些條例.
當時的條例.
可以去到進步的地方.
那個程度.
可以去到現代社會十九世紀.
公共衛生的程度.
所以你就要去想到底.
背後有什麼神學的原因.
表達一些什麼神學的思想.
其實就是利美基的忠心.
利美基的忠心其中一樣就是愛人如己.
Right?.
Love your neighbor as yourself(愛自己的鄰居為自己).
你的 neighbor(鄰居)是怎麼樣?.
你需要去 protect(保護)他.
建立一個神的社會.
有神同在的一個社會.

$^{561}$教會.
其實就是一個這樣的原則.
每個人看著的不是自己.
而是看著隔離的人.
以致隔離的人.
不會得到傳染也好.
或者以致他的明星得到保障也好.
愛人如己.
就是這樣子衍生出來的.
不只是說今天教會.
如果你引伸到整個的社會.
甚至是一些 secular(社會)的 society(社會).
其實如果真的是愛人如己的時候.
大家的生活會有多大的分別.
如果每個人不是這麼自私.
如果每個人都是看著隔離.
那個鄰舍的需要.
那個社會會有多不同.
在1665年的時候.
英國有一個小村.
叫做伊木.
在伊木發生了一個很大的疫症.
疫症的原因是怎麼發生的呢?.
就是一個裁縫收到一疊衣服.
那疊衣服是濕濕的.
是發霉的.
當時收進去的時候.
衛生或者醫學的常識不是很多.
所以一收了的時候.
有很多的濕跑出來了.
然後就將這個疫症擴散.
在當時這條村.
有很多的村民.
因為看到疫症越來越擴散的時候.
他們就提議說.
不如我們跑到另外一個城市那裡.
去走攔.
但是他們還沒做這件事情之前.
有一位牧師就站出來.
那位牧師就說.

$^{601}$就勸大家留下.
我們整條村隔離自我的 quarantine(隔離).
以致我們不會將疫症帶到其他的村.
去傳染其他人.
他在教會這麼說的時候.
就有很大的爭議.
但是他用了很大的努力.
就是去說服村裡面的村民.
我們不如留下.
結果他們真的整條村留下了.
整條村本身是三百五十人的.
在疫症完的時候.
一年多之後.
就死剩九十人.
但是因為這樣的緣故.
因為他們自我隔離.
他們犧牲自己.
其實就挽救了很多其他村的村民.
這個就叫做愛人如己.
這部電影其實就是想說.
如何去活出那種愛人如己.
不是自私自利的.
而是去 protect(保護)其他人.
那種的原則.
這個就是天國的原則.
天國的價值.
今天無論在教會.
你在我們的社會裡面.
其實都是需要一個這樣的原則.
特別是在疫症的時期.
你是不是看到很多這樣的.
真的很瘋狂的表現.
其實就是代表我們人不僅沒有平安.
就是我們根本掌握不了這樣的原則.
愛人如己.
第二樣呢.
講到什麼呢?.
我們看看十四章第三十三節.
又說對摩西和阿倫說.
你們進到我賜給你們作產業的家林地的時候.

$^{641}$我給你們所作產業之地的房屋.
進到房子裡面的時候.
如果是發霉的話.
房主就要去告訴祭司.
說他聽什麼呢?.
據我看來.
房屋裡面似乎是染了病症.
祭司就要去那裡做檢查.
查看房屋.
如果發現房屋的牆上面.
是發霉或者發紅.
有腐蝕的斑紋的話.
或者.
而且深入了牆的現象的話.
祭司就要將房子清理了等等.
那些details(細節)我不講了.
但是最後就是將房子.
來到 quarantine(隔離).
for seven days(七天).
將他隔離七天.
這裡處理另外一個問題.
就不是皮膚病的問題.
again 又是那種大麻風.
就是上到房子的人住的牆上.
那麼.
大家可能也知道過了.
在 basement(地下室).
那些 dry wall(乾屋).
發霉.
其實就是差不多這樣的事情.
他就要找祭司.
然後叫祭司來檢查一下.
祭司是當時的專業人員.
他就要去做防疫.
然後去做隔離.
但是我想這裡最重要的是什麼呢?.
這裡最重要的就是房主要親自告訴祭司.
Right?.
你買了房子.
是不是?.

$^{681}$房子是屬於你的.
如果有大麻風病.
有毛菌的話.
房子隨時都會被人收起來.
是一個經濟的損失.
但是這裡的律例講什麼呢?.
他說房主就要去告訴祭司.
Right?.
是要那個人自己肯去透明地說出真相.
我們已經很多人在講.
防疫抗疫.
最重要是什麼呢?.
就是要講真理.
就是不可以隱瞞.
Right?.
你想想他在假設什麼呢?.
假設在社會裡面.
是有言論.
有新聞的價值.
是不是?.
Otherwise(否則)就不會說.
就是要告訴別人.
要說出真相.
武漢肺炎爆發之前.
有八位醫生.
中國的醫生.
他們早就已經留意到.
但是就變成了.
一講出來的時候.
就不是.
就是被封口.
是不是?.
12月中.
早知道那個病毒已經可以人傳人了.
很明顯是有.
遲遲不報或者.
一些隱瞞的事情.
以至引致到現在更大的擴散.
這裡講到一個社會的基礎.
如果你只是單講公共衛生的話.

$^{721}$其實是一個政治問題.
我們經常說我們不理政治.
政治不關我們的事.
你戴著口罩就是政治.
我們現在要處理一個這樣的情勢.
其實就是一個政治的問題.
一個社會沒有了一個所謂言論的價值.
言論的自由.
沒有一個新聞的自由的話.
會影響多大呢?.
這些全部都是政治問題.
不可以我們不關心的.
每一天的生活.
吃的東西.
喝的東西.
衣食住行.
其實每一樣都是很實際的.
也是跟政治有關係的.
所以那個社群的基礎是什麼呢?.
是要說出真理.
說出真相.
我們沒有時間看其他的.
彼此應該怎麼做.
那些愛人如己.
弟兄姊妹如果有肝飯來的.
你應該親自道出真理.
用愛心說誠實話.
就是這樣的意思.
要將真相透明的說出來.
你還沒記十三章十四章.
如果你一口氣讀的話.
你就會經常留意什麼呢?.
你會留意那個祭司經常出現.
幾乎每一節.
都起碼會提一次祭司.
甚至有時候兩次祭司.
兩章聖經十三十四章還沒記.
處理這些疫症的時候.
就提到祭司九十五次.
因為只是兩章聖經而已.

$^{761}$幾乎每一節都提祭司.
如果你按整個的.
我們經常說 big data(大數據).
整個的舊約或者甚至整本的 Leviticus(利比特).
祭司出現的地方.
最多就在這裡.
最密就在這裡.
幾乎是不成比例的.
就是處理疫症為什麼九十五次.
不成比例地.
去提出.
說要祭司在這裡.
背景原因.
我們明白了.
我們知道為什麼經常提祭司.
因為近東的社會裡面.
他們處理這些所謂的皮膚病.
這些病症的時候.
你就知道他們很多時候.
就是要去.
有很多迷信的 practice(練習).
他們可能會用一些咒語.
或者一些法術.
或者一些通靈的方法.
護身符的方法.
來處理這些病症.
你有牙痛.
你就跟牙神去禱告.
就希望他弄走你牙裡面的蟲.
是不是?.
當然他們沒有這樣的.
現代醫學的知識.
但是他們很多時候就是去.
用一些符咒的方法.
你知道以色列的社會.
舊約的律例.
是不讓他們這麼做的.
所以一味就是說.
你一定要 go to the priest(去禮尊).
是上帝所按立的那個人.

$^{801}$來處理這些事情.
重複又重複.
強調一些虛假的神靈.
是幫不了你,治不了你的.
是要經過祭司的診斷.
是不是?.
你用那些迷信的方法是不行的.
這些就是當時的假新聞.
今天.
疫症蔓延的時候.
假新聞也有很多.
蔓延是不是?.
在中國大陸因為封城的緣故.
中國的票房.
就大受損失是不是?.
所以很多人在家裡 stream(直播)他的 movie(電影)進來.
在家裡自己看.
最多人看的其中一部電影是什麼呢?.
就是 Contagion(感染).
Contagion(感染)是幾年前拍的.
好幾年前拍的.
就是講到病毒的情節.
那個情節是怎麼樣呢?.
就是說女主角去了香港出差.
然後她接觸了一些病毒.
就帶了回來美國.
然後疫症就爆發了.
由她的家人開始.
一直爆發.
那個電影的 subtitle(字幕).
講中了一種現象.
Nothing spreads like fear(恐懼無比).
沒有一件事情.
比恐懼傳染得更加快.
你想想今天.
居然搞到.
人們打劫可以去搶口罩.
打劫不是去金舖打劫.
打劫可以去.
去貨倉搶口罩.

$^{841}$就是那些人.
一罩難求.
我沒想過.
輝春可以寫這四個字.
簡直令人很….
真的苦中作樂這些是不是?.
你想想那些人沒有口罩.
可以有多少的創意是不是?.
你沒想過你可以….
我們覺得好笑是不是?.
一定是覺得好笑的.
Of course(當然).
但是.
但是你真的覺得.
啊?.
這個就是恐懼是不是?.
右下角那個不是中國大陸的.
我想是昨天看到在 Houston(休斯頓).
有人上飛機這樣子上.
那你說他恐怖分子還是.
想防 virus(病毒)呢?.
香港甚至有一個高官說.
可以清蒸口罩.
就是可以翻用.
清蒸翻用口罩.
就是這些假消息.
清蒸炒煮炸.
都是錯的.
是不可以的.
這些事情真的會害死人的.
你明白我的意思嗎?.
口罩只能夠用一次是不是?.
其實要解決沒有口罩的.
也是一個政治問題是不是?.
是不是?.
政府可以解決的.
是不是?.
所以有很多假新聞.
我們真的需要去分辨一下.
以往的假新聞就是那些迷信的東西.

$^{881}$今天的假新聞就是.
其中一些就是這些.
是不是?.
這張很厲害.
這張說蒸馬拉糕是終點是不是?.
蒸燒賣是大點.
蒸蝦餃是特點.
蒸口罩就是老點.
你說是多少東西老點我們是不是?.
為什麼這麼多東西可以老點我們呢?.
就是因為我們.
我們很多恐懼你明白嗎?.
很多恐懼.
不知道怎麼做.
不知道怎麼辦.
是不是?.
所以.
這些真的是苦中作樂是不是?.
在那部電影.
Contagion.
有不少的情節強調什麼呢?.
因為整部電影都是在講抗疫.
放疫.
就是跟報導真相息息相關.
CDC需要將真正的真相透明地.
向大眾交代.
是不是?.
不辭而言的交代.
說出他的真相是不是?.
所以一個社會的建立.
神國社會的建立.
神的價值就是那個真相是不是?.
第三方面.
就是一個社群的基礎是什麼呢?.
就是.
有盼望.
經常提到祭司.
就是代表一個嚴謹的監察機制.
因為祭司是什麼呢?.
祭司是中間人.

$^{921}$祭司是做上帝.
跟人群的中間人.
他要做教導.
他要做安慰.
他要為患者祈禱.
當然律例上沒有說.
但是你可以看到很多詩歌.
詩篇裡面有很多詩歌.
有一些叫做哀歌.
是一些病人發出的哀求.
是他們的土文.
很多時候哀歌有一個格式.
後一半.
他由哀求變成信心去仰望轉向神.
相信因為有祭司.
他去到聖殿的時候.
他看到祭司.
祭司去教導他.
引導他.
帶他去到上帝面前.
這個就是盼望.
全人的醫治的盼望.
就是醫治你.
不是單單搞好你的 symptom(症狀).
而是你的心靈也得到完整.
有一個醫學倫理的學者.
聖經學者.
也是醫學倫理家.
他寫了一本書叫做.
Reading the Bible in the Strange World of Medicine.
裡面有一段話很有意思.
他說醫學很多時候是幫助我們辨識.
理解.
以及控制病痛.
辨識是將病痛客觀化的第一步.
為病情提出一個標籤.
一個合理的解釋.
創建醫學的一種語言.
但這個就是最諷刺的地方.
因為很多時候.

$^{961}$這些醫學所創造的語言.
並非病人的語言.
而是病人的語言.
當這些醫學的語言.
變成官方的語言.
或者唯一獲准使用的語言的時候.
病人就會怎麼受苦呢?.
他就會自己覺得.
我是一個陌生人.
無言以對.
對於世界.
對於他自己痛苦.
就沒有盼望.
因為他找不到一種 language(語言).
認同他現在的感受.
這個就是所謂醫學世界裡面.
現代醫學世界的怪異錄.
一個病人去看醫生.
你聽到的東西.
得到的診斷.
只不過是一些 official(官方)的 medical language(醫學語言).
而是對你內心深處的東西.
是不能夠認同的.
大概你應該試過.
在我們中間.
很多大英姐妹都可能試過.
當你患重病的時候.
你進去見醫生.
你進去手術室.
你進去門診部的時候.
你只聽到一些醫學的專用語言.
一些很難明白的 terms(語言).
但是我想問你.
當時.
你真正的需要是什麼呢?.
很肯定不只是這些 medical terms(醫學語言).
而是你需要心裡面的安慰.
是不是?.
你更加需要有人去認同.
你心裡面的那種痛苦.

$^{1001}$要有人跟你同行.
這就是之所以為什麼.
醫院裡面有院目.
有 chapter(章節).
為什麼我們牧者會去探訪.
為什麼我們大英姐妹會去探望其他人.
會去為他祈禱.
會跟他同行呢?.
這就是現代醫學做不到的事情.
所以為什麼會有 counselor(顧問).
有輔導員.
因為輔導員應該會幫助你.
用你的語言.
你明白的語言.
你能夠認同的語言.
使你心裡面能夠得到安慰.
有那種安穩.
結束的時候.
讓我們提醒自己.
上帝自己天國的價值.
抑正.
有些人說.
往往就是揭示社會在政治文化上的問題.
這些是我們需要關注的.
我們作為天國公民.
我們需要關注.
疫症當中.
讓我們也幫助一下身邊的人.
你不要囤積你的口罩.
有很多人沒有的.
甚至護士,醫護人員沒有的.
去幫助一下其他人.
如果你現在在你的處境裡面.
你有些恐懼.
有些憂慮.
讓你來到上帝面前.
我們有一個最好的祭祀.
耶穌基督.
完美的大祭祀.
軟弱的時候能夠給予我們力量.

$^{1041}$能夠安慰我們.
開童心大道基督.
因為你是我們的平安.
主我們求你將這個平安.
讓我們穩穩妥妥地抓住.
讓我們知道我們在基督的十誡之下.
我們有永遠的盼望.
你是那位醫治者.
求你醫治我們心靈裡面的病痛.
我們的憂慮,我們的恐懼.
主你臨到.
讓我們安穩.
也讓我們能夠將自己的平安.
帶到這個世界.
沒有平安的世界.
求主你藉著我們所言所行.
能夠祝福更多的人.
繼續為疫症蔓延的.
受很多影響的.
中國大陸不同的城市.
香港.
不同的地方.
主我們求你來幫助.
求你來憐憫.
我們仰望你.
祈禱奉主耶穌基督的名字.
\newpage



\section{利未記 11:1-19}
\label{sec:imYjaG_VAbM}
\textbf{2020/2/9 主日講道: 疫境前行 2 (利未記11\_1-19) - 陳崇基牧師}
\newline
\newline
連結: \href{https://youtube.com/watch?v=imYjaG_VAbM}{\texttt{https://youtube.com/watch?v=imYjaG\_VAbM}} ~~~~ 語音日期: 2020-02-14
\newline
\newline
\hyperref[sec:ez9_0uQY1v4]{\small{< < < PREV SERMON < < <}}
~
\hyperref[sec:index_chronic]{\small{[返順時目]}}
~
\hyperref[sec:index_scriptual]{\small{[返順卷目]}}
~
\hyperref[sec:CWGXYDFB8bU]{\small{> > > NEXT SERMON > > >}}
\newline
\newline
利未記 11:1-19
\newline
\begin{longtable}{cl}
\hline
\hline
章節 & 經文 (和合本修訂版)\\
\hline
11:1 & \begin{tabularx}{0.7\textwidth}{X} 耶和華吩咐摩西和亞倫，對他們說： \end{tabularx} \\ \\ \relax
11:2 & \begin{tabularx}{0.7\textwidth}{X} 「你們要吩咐以色列人說，地上一切的走獸中可吃的動物是這些： \end{tabularx} \\ \\ \relax
11:3 & \begin{tabularx}{0.7\textwidth}{X} 凡蹄分兩瓣，分趾蹄而又反芻食物的走獸，你們都可以吃。 \end{tabularx} \\ \\ \relax
11:4 & \begin{tabularx}{0.7\textwidth}{X} 但那反芻或分蹄之中不可吃的是：駱駝，反芻卻不分蹄，對你們是不潔淨的； \end{tabularx} \\ \\ \relax
11:5 & \begin{tabularx}{0.7\textwidth}{X} 石獾，反芻卻不分蹄，對你們是不潔淨的； \end{tabularx} \\ \\ \relax
11:6 & \begin{tabularx}{0.7\textwidth}{X} 兔子，反芻卻不分蹄，對你們是不潔淨的； \end{tabularx} \\ \\ \relax
11:7 & \begin{tabularx}{0.7\textwidth}{X} 豬，蹄分兩瓣，分趾蹄卻不反芻，對你們是不潔淨的。 \end{tabularx} \\ \\ \relax
11:8 & \begin{tabularx}{0.7\textwidth}{X} 這些獸的肉，你們不可吃；牠們的屍體，你們也不可摸，對你們都是不潔淨的。 \end{tabularx} \\ \\ \relax
11:9 & \begin{tabularx}{0.7\textwidth}{X} 「水中可吃的是這些：凡在水裡，無論是海或河，有鰭有鱗的，都可以吃。 \end{tabularx} \\ \\ \relax
11:10 & \begin{tabularx}{0.7\textwidth}{X} 凡在海裡、河裡和水裡滋生的動物，就是在水裡所有的動物，無鰭無鱗的，對你們是可憎的。 \end{tabularx} \\ \\ \relax
11:11 & \begin{tabularx}{0.7\textwidth}{X} 牠們對你們都是可憎的。你們不可吃牠們的肉；牠們的屍體，也當以為可憎。 \end{tabularx} \\ \\ \relax
11:12 & \begin{tabularx}{0.7\textwidth}{X} 凡在水裡無鰭無鱗的，對你們是可憎的。 \end{tabularx} \\ \\ \relax
11:13 & \begin{tabularx}{0.7\textwidth}{X} 「飛鳥中你們當以為可憎，不可吃且可憎的是：鵰、狗頭鵰、紅頭鵰， \end{tabularx} \\ \\ \relax
11:14 & \begin{tabularx}{0.7\textwidth}{X} 鷂鷹、小鷹的類群， \end{tabularx} \\ \\ \relax
11:15 & \begin{tabularx}{0.7\textwidth}{X} 所有烏鴉的類群， \end{tabularx} \\ \\ \relax
11:16 & \begin{tabularx}{0.7\textwidth}{X} 鴕鳥、夜鷹、魚鷹、鷹的類群， \end{tabularx} \\ \\ \relax
11:17 & \begin{tabularx}{0.7\textwidth}{X} 鴞鳥、鸕鶿、貓頭鷹， \end{tabularx} \\ \\ \relax
11:18 & \begin{tabularx}{0.7\textwidth}{X} 角鴟、鵜鶘、禿鵰， \end{tabularx} \\ \\ \relax
11:19 & \begin{tabularx}{0.7\textwidth}{X} 鸛、鷺鷥的類群，戴鵀與蝙蝠。 \end{tabularx} \\ \\
[1ex]
\hline
\hline
\end{longtable}
$^{1}$各位早安!我們繼續這個系列叫做逆境前行.上個星期開始我們看一些不是平時我們停修的時候看的聖經,或者我們很少會留意的聖經,看看裡面有什麼神學的教訓,信仰,和生命裡面的一些提醒..
你記得上次我們打開去舊約一本書叫做《利未記》,今天就意猶未盡,我們仍然是在看《利未記》的.有聖經的話請你幫我打開..
去到《利未記》第十一章一節開始..
(1)對摩西和亞倫說:你們曉如以色列人說,在地上一切的酒受中可吃的乃是這些.凡提婚兩婦,杜雀的酒受,你們都可以吃..
但那杜雀或分堤之中不可吃的乃是,駱駝因為杜雀不分堤,就與你們不結證;沙蕃因為杜雀不分堤,就與你們不結證;兔子因為杜雀不分堤,就與你們不結證;豬因為堤分兩反,卻不杜雀,就與你們不結證;.
這些瘦的肉,你們不可吃;死的,你們不可摸,都與你們不結證.水中可吃的乃是這些.凡在水裡,海裡,河裡,有刺,有鱗的,都可以吃.凡在海裡,河裡,並一切水裡游動的活物,無刺無鱗的,你們都當以為可憎..
這些無刺無鱗以為可憎的,你們不可吃他的肉;死的,也當以為可憎.凡水裡無刺無鱗的,你們都當以為可憎.雀鳥當中,你們當以為可憎的,不可吃的乃是:鵰,狗頭鵰,熊頭鵰,妖鷹,小鷹與其類,烏鴉與其類,鴕鳥,夜鷹,魚鷹,鷹與其類..
這裡有17字..
鵰鵰,鵰鵰,貓頭鷹,鴿齒,鶘鵰,貫鵰,鵰鵰與其類,大吟與蝙蝠..
這些經文的內容都是與舊約的神的子民猶太人的食譜有關..
有些食物可以吃,有些不可以吃,有些是潔淨的,可以放進口中,有些是不潔淨的,不可以放進口中..
其實經文的內容提到有三類的活物,在世間當中,有地上的,有水裡面的,也有天空上的活物..
所以你不要理會他列出的那些什麼什麼的,我弄了很久才知道廣東話怎麼說..
地上不可吃的是什麼呢?.
他說找到的那些活物是分體的和加上杜著的,即是反初的,那些你才可以吃..
水裡面的有刺有鱗的,其他就是可憎的不潔的..
天上的飛的,沒有列出可以吃什麼,不過列了一堆不可以吃的,有十幾種..
上次我們提到舊約聖經裡面神學,說到潔淨與不潔淨的時候,很多時候主要說的不是罪與義的問題..
不是說這些狗頭雕有罪,不是moral的問題,也不是說衛生hygiene的問題,也不是說科學的問題..
潔淨與不潔淨是關乎在我們不同的文化,這些經文當然是說到舊約時候文化,他們有不同的category..
我們每一個人,每一個文化都有我們的category,上次我們舉了一個例子,鞋放在哪裡,我們認為它應該belong,屬於哪個地方..
鞋放在地上,我們就不會覺得是dirt來的,但是鞋放在餐桌上,我們就會覺得是污糟的..
或者我們亞洲人,我們進屋的時候,我們要做什麼?.
這是我們的ritual,我們的禮儀,你必須要脫鞋,你不能說我剛剛買的鞋是新的,我可以穿上去到處踩,因為是不屬於裡面的..
這個就是我們文化裡面有的order和disorder,order就是說屬於那裡的東西應該放在那裡,disorder就是說不應該放在那裡的東西放在那裡,放錯了就變成了混沌,類別不對..
所以我們說到潔淨與不潔淨,應該是這樣想,什麼叫做category..
說到食物的category,我們不同的文化有不同的類別所分的,有些屬於食物,有些不屬於食物..
中國的華人,亞洲人其實都是有這樣的觀念的..
我和大家玩個遊戲好不好?快問快答..
我們說到食物的時候,有些什麼是可以吃的,有些什麼是不可以吃的..
白雲鳳爪你覺得可以吃嗎?.
有人認同,對你來說如果你認同的話,就是不潔了..
我太太不吃白雲鳳爪,但是豉汁的鳳爪她也吃,所以白雲鳳爪她就不潔了..
這個是她的category,我再問你,豬紅可以吃嗎?.
有誰覺得可以吃的?.
那些沒有舉手的可能是不吃的,對你來說就是不潔了..
不潔不是因為不衛生,不潔就是不在食物的category裡面..
大腸呢?吃不吃?有人不吃,不吃就覺得你不在那個category,無論放在哪裡都好,你都不會吃的..
但是我想問你,如果炸得很香你吃不吃?都不吃,所以就不是煮法問題,總之大腸炸是不可以吃的..
我沒有照片,免得弄一張照片出來,狗肉呢?.

$^{41}$這些不在我們的category裡面..
或者我說一個水果的時候,你覺得水果是可以吃的嗎?有什麼水果是絕對不可以吃的?.
有誰在這裡不吃榴槤的?死都不吃的..
你學吃呢?有些人學吃就喜歡吃的,會不會?不會吧?永遠都不會?那你就走寶了,水果王..
所以我們每個人都有不同的category,所以這個不是hygiene的問題,而是我們的組織和我們的所謂類別..
如果你問猶太人的食物的話,他們都是這樣的,每一個文化的食譜都有它的傳統,可能是歷世歷代遺留下來的傳統..
華人的文化,當然說到食物的文化,我們很多東西都可以吃的,對於一些白人來說,可能有很多東西對他們來說是厭惡的,可憎的,disgusting或者是unclean..
我記得我讀研究院的時候,我有一個白人的老師,那個教授,他說他小時候有機會去Chinatown,.
他小時候在Chinatown看到有些店鋪有一個水缸,他看到魚在游來游去,他覺得那個是一個寵物店,.
誰知道進去的時候發現他罵那些魚,然後就砍了它,然後就賣給別人,他覺得很disgusting,他覺得很噁心,.
他說為什麼要這樣對待魚,這個是他的文化和我們的文化的差異.猶太人都是一樣,或者這些食物的文化都是基於一些生物的特性,.
有些可能是可憎的食物,一些可能是雜食性的食物,鷹,一些比較殘暴獵食性的雜食性的一些生物,他覺得烏鴉不好吃,.
這些是文化傳統遺留下來的一些價值,一些category,或者可能是一些宗教的category,在近東裡面,.
以色列人為什麼這麼強調不吃豬肉呢?不是因為public hygiene的問題,當然可能有些public hygiene的原因,.
以前古代他們煮東西可能煮不熟,所以不衛生等等,最主要是他們不喜歡吃豬肉,到今天還是不吃豬肉,原因很多時候都是宗教的問題,.
因為在近東的文化裡面有很多異教色彩的社會,他們是將豬獻給一些邪靈的,一些陰間的邪靈,所以他們覺得是一種不潔的東西,.
你怎麼煮都好,他們都不會吃的,所以有很多不同的因素,以至特定的文化,有他自己的食物的category,.
這段經文的分段其實很清楚,你看頭十九字,他分地上的活物,然後水裡面的活物,天上的活物,.
無論你覺得他們的食譜為什麼有些東西可以吃,有些東西不可以吃,原因在哪裡,我相信最大的原因,或者可以捕捉到這段聖經到底想教什麼,想講什麼,.
最主要其實就是他們想起天地海,他們的宇宙觀是分開三個level,天地與海,當你想到天地與海的時候,你會想起哪一段聖經,猶太人很熟的,.
創世記第一章,因為上帝創造萬有,整理萬有的時候,分開黑夜,白晝,全部東西分開的時候,就是在那六天裡面,創世記的第五和第六天,.
就將天地與海的活物創立出來,所以他們想起的就是他們對創造的神學,theology of creation,他們認為上帝創立這個世界的時候,有他的category,有他的類別,.
另外你看十一章的話,你就知道整個字會出現,除了那些類別的動物之外,就是類字,kind,category,這個字在整章十一章出現了很多次,.
最多出現這個字是什麼呢?就是經文,創世記第一章,創世記第一章,就是各從其類,上帝創立整頓這個世界的時候,去到分配不同種類的時候,就是各從其類,天上的飛鳥是各從其類,.
海裡面的魚也是各從其類,就是上帝有秩序,整整有條的去創立order,跟這個order的東西就是潔淨,不跟這個order的東西,我們就不知道他放在哪裡,那些我們就認為他們是不潔淨,所以這個是猶太人舊約裡面的神學,.
譬如說水裡面的魚,有什麼可以吃的呢?有什麼潔淨的呢?魚就是有刺有鱗的,一定是有刺有鱗的,你想一下,在水裡面游的東西,就應該要有刺有鱗,不然怎麼游呢?.
但偏偏有些東西在水裡面不是有刺有鱗的,但在水裡面游,那些叫鱔,有些其他蝦蟹等等,那些算不潔淨,為什麼?因為不屬於那個category,怎麼define呢?.
沒有刺沒有鱗,但又在水裡面住的,所以這個又是講到廣東話才有這麼全神的表達,其他語言都沒有,叫不頂不吊是什麼意思呢?.
什麼叫不頂不吊?就是你fit不進去那個category,所以就成為所謂不潔的東西,天上的飛鳥,你會發現他們有不同的雜性,鷹有不同的雜性,或者蝙蝠可能是很古怪的東西,很disgusting的東西,或者樣子很噁心的東西,所以就變成不潔淨,.
但是請你好像上個禮拜一樣,利未記去處理一些疫症的時候,我們不可以照字搬紙的,我們要想深一層,到底有些什麼神學的概念,有些什麼conviction,背後底下有些什麼深一層的ideology,是引發他們有這樣的看法,那個才是最重要的訊息,.
食物的條例我們不只是表面去看這些字,我們就說照辦主觀,今天不能吃的東西全部在這裡,那不就不行了,那你就不可以吃焗豬扒飯,不可以這樣照字搬紙,你要想一下他背後有什麼world view,而那個dietary law,那個飲食的條例,就是去表述一個世界觀,一個live view,怎樣去看世界,以往舊約的人就是這樣看世界,.
原因就是出自那個live view,那個世界觀是什麼呢,就是上帝是創造主,上帝是創立這個世界,真正生命的主權,終極生命的主權,像飲食那個,是上帝,不是你努力就可以拿回來,是上帝有主權去分配給你的,.
泥土,這個是一個很重要的宣告,就是說雖然我們人以為我們是萬物之靈,但是我們是沒有最大權的,創造主才有絕對的主權,.
你想一下你今天呼吸,都是創造主有的主權,大概兩個星期前,Coby Bryant,就一個空難就去世了,你上都沒有上過,人生就是這樣,一個的病症或者一個的疫症,就揭發了我們人生命其實是很脆弱的,.
我們很無力,在湖北有一個女孩子,她最近在網上發表了她的遺言,她睡在醫院裡,一個很年輕的女孩子,她很不甘心,.
父母或者親戚以為她是得救,聽從醫院的安排,其實她自己知道她快要死,沒有醫生來,沒有人問她,沒有人管她,心裡面很不忿,太多話已經不知道怎麼說了,.
我知道還有我這個人發病的人,然後死在這裡,肯定是不少,我們如此相信政府,現在落到這個結局,.
你看看生命是多麼脆弱,可能外在你會覺得有些事情很富強,可能有一個國家,一個大國,可能是可以看上去是戰無不勝,.
但一個疫症來臨時,一些你看不到的東西,一些病毒,就可以封鎖成千上萬的人,那種恐慌是你不能理解的,確實因為我們知道自己很脆弱,.
我們知道自己是不行的,不要亂靠一些東西在世上,無論那東西表面上粉飾得多麼強大也好,真正幫到你,真正值得你倚靠,應該是生命的主,應該是創造的主,.

$^{81}$食譜就是想說這件事,在我們吃之前,見到大自然,見到世界,我們能夠有資源,有基本的需要,供應,全部都是出自上帝自己,.
在2002至2003年的時候,我有機會和那時候教會有個報道會,沈祖堯醫生經過,他講完見證,我就做呼召,報道,福音,訊息,一個很好的醫生,很懂得鼓勵人的醫生,.
在他面對沙士爆發時期的時候,他說每天早上他找醫院裡面的同事集合他們,原來有很多個基督徒,就和他們一起祈禱,祈禱講什麼呢?.
祈禱就是說,跟神說我們真的沒有辦法,但求上帝不要因為我們無知,令到病人的生命犧牲,因為上帝才是掌管生命的神,人真的很渺小,能力是很有限,即使我們覺得自己很厲害,我們自己覺得很厲害,.
但其實我們很渺小,因為生命真的不是我們可以手去操縱,所以我們看到利未記,你可能覺得一些很無關痛癢的東西,今天我們不會照著利未記的食譜去吃東西,.
原文的意思也不是這樣的意思,最主要是告訴我們背後就是操縱生命創立生命的主,每一次你把飯放在你口中,每一次你能夠起床,可以去有口氣,你能夠去讀書,能夠去工作,這些都是上帝的安排,上帝的恩典,生命的主權是上帝,.
另外,如果你看這些食譜的時候,你總是會覺得有些隔硬的感覺,你一定會覺得這樣,無論你怎樣解釋也好,我已經盡量去解釋用人類學的方法,什麼category等等,我們去解釋到底怎樣了解什麼是clean,什麼是unclean,.
無論如何我們都覺得有些隔硬,譬如為什麼一定要反錯又要分題呢,為什麼一定要反錯又要分題呢,你可以說那個category裡面地上的東西應該是這樣的,如果不符合那個規格,他們就覺得是不潔的,是unclean的,無論如何都是有些隔硬,隔硬是什麼意思呢,.
我想大家很少會想過食物是跟我們的ethics倫理有拉上關係的,我們吃飯的時候好像會想到food跟ethics息息相關,我們真的沒有,.
但有一個全世界最高的研究利美技學者,是一個猶太人的拉比,Jacob Milgram他已經去世了,他花了他一輩子的人生就是去看這本書,就是去看Latificus,.
他說這個dietary law其實是一個ethical system,是一個倫理的系統,就是說那些食譜你看的時候其實是在看倫理的表述,一些ethical statement,什麼意思呢,就是說食物是跟倫理息息相關的,猶太人這樣吃東西代表他們是上帝聖結的子民,.
是在效法上帝,效法上帝做的事情,吃東西不吃什麼就是去到親近上帝,是他靈性的一部分,所以你看到這個食譜很強調某一樣東西,感覺你大概應該get到,.
雖然讀那些食譜很笨很煩很細節,但你應該大概get到感覺就是很強調是限制,不是說你找到那東西動的就可以殺了它然後煮來吃,有很多的限制,換句話說拿什麼吃去滿足你的食慾呢,不到你說了算,.
是創造主華,什麼可以放在飯桌上是有限制的,是有limitation的,你還沒記住飲食的條例,你想一下列出的不可以吃的,篇幅是多一點的,換句話說給我們一些強調我們飲食上是要節制的,.
人不可以說我看到那東西動我就拿來吃,怎樣煮來煮去都不要就炸它,炸香它就好吃,你不可以這樣,因為創造界不是給人去濫用,沒錯是給我們來享用,但卻非給我們濫用,不是給我們去explore,.
不是每樣東西都可以動的都可以放進口,這個again不只是衛生問題,而是上帝給創造界裡面住的人,他們去管理,他們可以去享有,但最終的主權是creator,在這個疫症裡面一個市場大家很熟,.
經常聽到這個市場,武漢的華南海鮮市場,他們上面賣的東西,有個UV有個弟兄跟我說,你看看那裡,任何你想到的東西都可以賣的,想到的東西都可以吃的,除了人肉之外,.
未記過的食譜,講到creation theology的時候,不是每樣東西都可以放進口,因為不是給你濫用,上帝創造在大自然裡面某些東西是可以讓你去研究,可以讓你去欣賞,但不是給你濫用,不是給你consume,.
在中世紀,大公教會,他們一個神學的整合,整合了聖經裡面講很多我們人犯錯,我們的罪,我們的破碎等等,這個也是系統神學,整合了人類有七大宗罪,seven deadly sins,.
seven deadly sins是什麼呢?有人甚至把現在的seven deadly sins放在online上面,seven deadly sins是什麼呢?七宗是致命的罪,第一是窒道,人家有的東西我沒有,我好窒道,第二是懶惰,不去being productive,不去生產,不去貢獻,.
第三是傲慢,pride,第四是憤怒,那些是不義之怒,第五是暴食,蜥蜴和貪婪,這裡就是講到暴食,為什麼講到暴食,是one of the seven deadly sins,暴飲暴食,.
就是說不管這麼多,三七二十一,總之我治了我的胃口,滿足我的食慾,那就算數,放縱我們的食慾,無限制地放縱,這個也是one of the seven deadly sins,.
你有沒有想過,你吃東西是你生命靈性的indicator,是表達你靈裡面的一個,你和上帝關係的一個指標,你試一下這樣想,I think我們就get到到底那個message在哪裡,那個食譜想講什麼,.
前幾年我認識一個白人的年青牧師,二十多歲的,剛剛上學院畢業,開始,牧會開始,港獨開始去帶人家,教導人家,有一次跟他吃飯,坐下來聊天,幾年前,.
大家說原來他很努力地在減肥,他真的很肥,他那種肥不是我們華人那種肥,那種肥是白人那種肥,你知道白人那種肥是,那些噸位很厲害,那些mini cooper他進不去,要買一些大的美國車才行,.
跟他聊天的時候,他說他很努力減肥,我跟他聊天的時候就問為什麼,我想他見到我的時候應該是很質度,很多牧師都很肥,我跟他聊天的時候我就問他為什麼,.
他就說出他心裡面在想什麼,他說我一定要減肥,我立定決心要減肥,為什麼呢,他說我作為一個牧者,如果我要教我的教會的人,要去discipline,要去自律的話,你想一下我這樣,.
怎麼可以教到別人自律,如果不是我就假務為事,所以他就很努力地減肥,他說這個我一定要先克服,結果幾個月他真的減了很多,他真的學了很多東西,真的佩服佩服佩服,.
我記得十幾年前我也是這樣,我都曾經肥過,一天可以吃兩個大的飯盒,做了飯桶都不知道自己,但是不吃不行,覺得自己,因為吃開了,.
吃到自己身體有點差,負荷很大,當然負荷大,身體裡面的機能很大的負荷,overload,你美肌教人做什麼,.
你美肌我正如上一次所說,整本notificus,不只是說獻祭獻什麼,吃東西吃什麼,然後去聖所的時候要預備什麼等等,最主要說什麼呢,就是要教我們,.
如何去親近上帝,教我們如何去生活,以致我們能夠和上帝建立一個好的關係,人是有罪,人是會濫用美好的創造,我們的食慾是一定要揭紙的,一定要牽制的,.
律例就是這樣的意思,就是要我們知道我們要去牽制我們有罪的慾望,叫我們去尊重上帝的創造,要appreciate上帝的creation,而不是去exploit,去consume,.
一味去滿足我自己一己的食慾,如果你這樣想的時候,那種牽制,那種restriction,discipline,其實不只是食慾這麼簡單,而是我們如何看待這個creation都是一樣,在我們中間做研究的不少人,科學研究,科學研究都可以有很多禍闖出來的,.
現在已經有另外一些人說不是野味出來的,病症,是實驗室可能洩漏某一些所謂進攻性的病毒,很高級很厲害的病毒,.
這令人很擔憂,因為我們去研究這個世界的時候,科技固然是好,沒有人會說科技不好,但是我們會過界,會濫用,或者會大義,不顧安全,不顧生命神聖,.
做出來的時候很多時候就是不堪設想,你想一下整個的律法,Torah,整本聖經其實就是教我們如何生活,其實不是單單說飲食,食慾這麼簡單,情慾,我們的言語,我們的關係,全部律法都有講,.
你如果想對你的仇人,你的仇家,他不見了一隻牛,牛走了去你那裡,你應該怎樣去看顧他,你的仇家牛去了你家,你會做什麼?你宰了他,然後我吃了他,你應該還給他,這個就是人際人倫的關係,.
時間怎樣運用?時間也是上帝給我們的美好的東西,我們就不要去濫用它,所以七天裡面有一天是安息日,是屬於上帝的,我們的物質的享受,我們財富都是discipline,奉獻都是一個discipline,.
就不只是滿足我consumerism,我一味去到拜物,一味去到滿足我自己的私慾,有錢我就去滿足我自己的私慾,這個不是體現上帝的天國的生命,這個是我們的罪性的表現,十幾年前,你記得有一本書叫做Purpose Driven Life,是Wick Warren寫的,那本書你不知道為什麼他一寫,是一本很簡單的書,一寫就很暢銷,.

$^{121}$然後他就變成了百萬富翁,賺很多錢,幾乎翹起雙手,書不停賣,有一次他在TED Talk,你知道什麼是TED Talk,TED Talk就是現在有這樣的平台,有些人去share一些idea,他有一次講TED Talk,他就跟觀眾講,.
他說現在我們很多錢,不單止奉獻給教會,甚至叫做reverse tithing,不是十一奉獻,而是調轉的十一奉獻,什麼意思呢?.
就是他每年從十開始就一直加上去,十一,十二,十三,十四,幾年後他就給二十,再幾年後他就給二十五\%,那些人問他,現在你發達了,你的生活方式有沒有改,他說他跟他太太一點都沒有改,仍然開舊的Ford,仍然住那間本身的樓房,.
他說了一件很有趣的事情,很值得我們去深入去想,他說他為什麼要這樣做呢?他說要break the chain of materialism,他說要將我們物質主義,享受消費主義,在那個綑綁中釋放出來,.
你的食慾滿足之後,你不竭致地去滿足的時候,其實食慾就控制了你,食慾,envy,傲慢等等,這些全部的desire其實會回來控制我們,所以一定要去discipline我們自己,.
這些在自由的社會裡面,很多時候都不想講,因為我們全部都是講freedom,但是freedom不是我想做什麼就做什麼,freedom是我應做什麼我就做什麼,而我不應做的我不做,這個叫freedom,最後讓我們結束的時候,十一章四十五節,那個教導最終是想講什麼?.
就是告訴世人聽,我耶和華是將你們從埃及遺老之地領出來的,要作你們的神,所以你們要聖潔,正如我為聖潔,求主幫助我們,我們同心低頭祈禱,主我們都是一班不配受恩典的,在他裡面的人,.
我們很多破碎,我們很多私慾,所以我們求你來讓我們認定你就是我們生命的主,認定你是我們生命當中,主讓我們來到你面前進行你的心意,讓你的道量,你生命的道,來引領我們,.
我們需要悔改的地方,我們需要去重新去整理我們生命的地方,求你幫助我們,讓我們在神國的生命裡面,我們有節制,有自律,讓我們生命能夠告訴別人,神是永恆的人,求主你幫助,祈禱奉主名,.
\newpage



\section{歌羅西書 1:25-29}
\label{sec:CWGXYDFB8bU}
\textbf{2020/2/16 主日講道: 全面的佈道事工 (歌羅西書1\_25-29) - 黃茂林傳道}
\newline
\newline
連結: \href{https://youtube.com/watch?v=CWGXYDFB8bU}{\texttt{https://youtube.com/watch?v=CWGXYDFB8bU}} ~~~~ 語音日期: 2020-02-21
\newline
\newline
\hyperref[sec:imYjaG_VAbM]{\small{< < < PREV SERMON < < <}}
~
\hyperref[sec:index_chronic]{\small{[返順時目]}}
~
\hyperref[sec:index_scriptual]{\small{[返順卷目]}}
~
\hyperref[sec:zCN_cR7XoxQ]{\small{> > > NEXT SERMON > > >}}
\newline
\newline
歌羅西書 1:25-29
\newline
\begin{longtable}{cl}
\hline
\hline
章節 & 經文 (和合本修訂版)\\
\hline
1:25 & \begin{tabularx}{0.7\textwidth}{X} 我照神為你們所賜我的職分作了教會的僕役，要把神的道傳得完滿； \end{tabularx} \\ \\ \relax
1:26 & \begin{tabularx}{0.7\textwidth}{X} 這道就是歷世歷代所隱藏的奧祕，但如今向他的聖徒顯明了。 \end{tabularx} \\ \\ \relax
1:27 & \begin{tabularx}{0.7\textwidth}{X} 神要讓他們知道，這奧祕在外邦人中有何等豐盛的榮耀；就是基督在你們心裡成了得榮耀的盼望。 \end{tabularx} \\ \\ \relax
1:28 & \begin{tabularx}{0.7\textwidth}{X} 我們傳揚他，是用諸般的智慧，勸戒各人，教導各人，要把各人在基督裡完完全全地獻上。 \end{tabularx} \\ \\ \relax
1:29 & \begin{tabularx}{0.7\textwidth}{X} 我也為此勞苦，照著他在我裡面運用的大能盡心竭力。 \end{tabularx} \\ \\
[1ex]
\hline
\hline
\end{longtable}
$^{1}$.
我召神為你們所賜我的積分.
作了教會的執事.
要把神的道理傳得完備.
這道理就是歷世歷代所隱藏的墓碑.
但如今向他的聖徒顯明了.
神願意叫他們知道.
這墓碑在外邦人中有何等豐盛的榮耀.
就是基督在你們心裡承了有榮耀的法律.
我們傳揚祂使用諸夫的智慧.
勸戒各人 教導各人.
要把各人在基督裡.
完完全全地引到神面前.
我也為此勞苦.
照著祂在我裡面運用的大能.
盡心竭力.
我們歡迎穆斯林.
各位弟兄姊妹.
早安.
很開心也很期待.
去年.
你們叫Mark Sir.
我叫陳穆斯.
一樣的 別說錯了.
很期待能夠和大家碰面.
聊天 分享.
分享聖經.
分享一些我生命裡覺得很重要的事.
我感動.
甚至可能是.
在報道的人生.
報道的工作裡.
都是離不開我的家庭.
我的子女 我的太太.
很多人的關係.
但這一切都是連接報道的生命.
首先再一次.
跪到代表短宣中心.
芝加哥的短宣中心.
知道你們也有參加我們耶穌愛之城.

$^{41}$我知道上一年你們有些困難.
但我知道你們今天也很積極.
所以為你們感恩.
其實能夠和大家分享這個訊息.
今天的感動.
很想說全面的報道事工.
因為我覺得.
我想每個人都不同意.
傳福音重要嗎?.
是重要的.
小的可能五六七歲.
大的可能有八十九十歲.
可能去教會有這種文化分途.
他們會覺得重要.
不過.
重要的是他們的事.
對我們來說.
我要問幾個問題.
就是這件事.
福音或者報道.
其實我們自己怎麼看.
我們怎麼去看待這件事.
福音的本質究竟是什麼.
真相是什麼.
福音和我們有什麼關係.
有些很好的和我們沒有關係.
有些地方很好吃的東西.
對我沒有關係.
因為我去不到.
或者我付不起那些錢.
福音的內容是什麼.
很多人都是來聽教會的一些訊息.
知道一點點.
無論他是否扮演代表.
可能他不太知道福音的真相.
全面是怎麼樣的.
報道的方法.
也是很重要的.
因為有時候你衝出去.
走過去和那個人說話.

$^{81}$那個人可能是怕你打劫.
你的熱誠對他來說是一個壓力.
所以有一定的方法.
應該怎麼做都很重要.
然後最常有的.
屬靈的棄嬰.
我自己信耶穌.
也有半個可能性是屬靈的棄嬰.
帶我信耶穌的人.
帶我信耶穌.
我很開心.
你們也會握手拍手掌.
為什麼下次.
他好像不覺得我存在.
其實我有一點傷心.
不過說的是差不多40年前的事.
不要怪他.
不過我反思.
如果有一個人.
你和他談得好好的.
甚至帶他去信耶穌.
下次你見到他會怎麼做呢.
這些令我很大的膠戰.
我不要做一個這樣的基督徒.
我曾經立過這樣的志.
希望大家也是這樣.
如果神在你生命裡面.
藉著你將一些好東西.
特別是福音的真理帶給人的時候.
神要繼續使用你.
下次見到他.
你沒有時間怎麼辦呢.
很簡單.
記住他的名字.
說一兩句問候他.
如果你真的能夠這樣做.
每次你見到那個人.
為他祈禱幾句.
下次見到他.
你就不會像說謊話.

$^{121}$或者內疚地說.
我會為你祈禱.
這是很容易有的方法.
你知不知道.
你知不知道我有一個罪.
曾經為了這件事認.
我是怎樣的呢.
做牧者.
我答應了為他祈禱.
在可能範圍裡面.
我真的要坐下.
蹲下.
跪下.
為他祈禱.
否則我在神面前.
就好像說了一個謊話一樣.
好.
這些是引言.
我想說的今天這段聖經.
我自己有思想.
為什麼保羅會說這段話.
其實在二十四節.
保羅說他是為了福音受苦.
覺得自己受苦.
彷彿就好像和耶穌基督.
一同的為了福音受苦.
滿足神的意思.
回到一個問題.
為什麼事受那麼多.
可以這麼說.
苦難或者不如意的事.
而去傳福音呢.
這個會不會是你的困難呢.
我們都很忙的.
有沒有誰不忙的.
請舉手.
我想認識一下你.
好.
你不忙.
那我一定要認識一下你.

$^{161}$你不忙的.
是很特別的.
是不是.
很多人都很忙.
還要在這麼忙的時間裡.
加上這麼偉大的工作.
傳福音.
或者報道.
又要我請假.
又要我回來教會.
或者我要出去社區.
或者去到天橋底.
不用去天橋底.
那麼多東西.
回到一個問題.
有什麼因素這麼有價值呢.
原來保羅.
接著在二十五節那裡.
彷彿就是用他的方式.
去回答了我們上面的問題.
他知道這件事.
是存在在我們的內心.
以至於去決定.
我們會不會覺得想去繼續.
做這個我們看為重要.
不過好像很艱鉅的工作.
如果這個是你心裡面的問題.
我答應你.
今天我們透過保羅.
嘗試拿一些激勵和答案.
首先我們可以看看.
在保羅第二十五節裡.
其實是說什麼呢.
我的領會就是.
其實他告訴大家.
他是因為神的一個恩典.
是為了那些對象.
領受了一個積分.
我看這個是一個使命.
你有沒有使命.

$^{201}$你有什麼使命.
我記得我有些親戚.
很喜歡說.
當然是賺錢.
這是不是你的使命.
其實賺錢不是很好的一個使命.
當然可能是你生命裡.
很重要的工作.
保羅他的使命是什麼呢.
保羅就是說.
他有一個積分.
是神給他的.
為的對象是為了他所服侍的人.
他稱為教會的執事.
我要解釋一點.
你們教會有沒有執事的.
哪個是教會執事.
不用怕.
沒有後果的.
你放心.
這裡說的教會執事.
不一定是你.
為什麼呢.
原來這個教會的執事的上下文.
是告訴我們.
這個執事是說福音的使命的那種服侍者.
執事的文字是服侍人.
可以是servant.
也可以用這個字說執事.
這個意思就是.
神給你有一個積分.
有一個使命.
去將福音傳開.
所以你說.
我不是這個執事.
也可能是.
如果我們不覺得自己有領受.
所以我們想到.
其實有時候我們會不會做一件事.
會牽涉到.

$^{241}$我覺不覺得自己想有這個工作.
這麼重的負擔.
保羅就覺得自己是很重要.
他有很重要的形式.
自己有這個負擔.
要去做這件事.
不要阻我.
沒有東西可以攔阻他.
在神的面前.
這個使命.
有些人怎麼看使命呢.
有些人看使命就是.
哦.
那我就要開執事會了.
如果以前是浸信會事奉.
開執事會.
開會又會決定.
今年這個猜拳.
這個報道工作會不會做.
在英國開始.
大概幾百年前.
他們有一個叫貴格會.
貴格會.
貴格物片那個字.
和物片無關.
他很著重內心的亮光.
所以他平時不傳福音也好.
有一天突然間亮光.
今天要傳福音.
就衝出去.
拿著一大單張.
可能就給那個人.
希望他信.
誰知道他又沒有反應.
還將那單張扔走.
於是他很灰心.
他就會覺得.
這個傳福音為什麼會這樣.
不知道會不會是.
我們那種貴格式的報道心態.

$^{281}$保羅看到人.
他自己的生命.
為什麼會這麼看重.
這個工作是一個使命.
所以弟兄姊妹.
有時我們不覺得那件事重要.
我們也不覺得那件事.
有一個這麼清晰的使命.
交給我們.
但是.
毋庸置疑.
耶穌基督升天之前.
已經將這個大使命給了我們.
無論我們知道也好.
或者不知道也好.
這件事.
你也要做.
我來到芝加哥.
我自己就很掙扎.
你知不知道掙扎什麼.
報稅.
我這個星期還在搏鬥.
什麼叫做.
Self-employment tax.
什麼叫做Social Security.
你不知道我說什麼.
你和我所面對的困難一樣.
無論我理不理解那條法例.
但是說真的.
我要不要給稅.
要.
要吧.
沒有人說我不用.
如果有些高層.
或者有權力的人士告訴我.
原來我不用.
那我很開心.
我是要的.
主耶穌給了我們使命.
不論我們承認.

$^{321}$已經給了我們.
在這段經文裡.
主給保羅的使命.
保羅特別說.
他是要將神的話語.
來到.
說神的教訓.
傳德傳弊.
那裡有兩個可能性.
我以往一直都認為.
傳德傳弊.
傳得齊齊全全.
傳得好一點.
不要搞到你們睡著了.
這樣應該是挺好的.
後來我再反覆去查.
傳德傳弊這個字.
是用一個完成.
或者說fulfill的意思.
原來它的重點.
可能第二個理解.
是那個範圍.
因為要fulfill這個神的話.
或者是完成到神的話.
這裡上下文是說福音.
是要怎樣呢?.
不可以保羅只是集中.
某個小地區的人的福音的需要.
好比今天有很多教會.
我行了.
我這裡這麼多人.
我向他們傳福音就夠了.
我哪有時間理這麼遠的那些.
吃咖喱的那些.
關我什麼事.
他們可能身體好一點.
不怕新冠狀病毒的.
弟兄姊妹.
原來這個意思.
是保羅看到範圍.

$^{361}$要傳避.
所以以保羅的人生.
他應受這個使命.
他不只是關心某個小地區的人.
我知道你們教會.
都很有負擔.
除了這個地方.
可能你們有不同的社區.
甚至你們也會去.
再遠一些去做差傳的工作.
我覺得這個是很寶貴的.
為什麼我們會這樣做呢?.
這個是第一個大點.
我們有個使命.
弟兄姊妹.
希望今年2020年.
除了你脫口罩之外.
我挑戰你.
在神面前有個使命.
就是為了神的國度.
來服侍.
我禮拜一.
很多時候會教一個福音的訓練班.
那些參加者大部分是.
五十幾到七十幾歲.
很多都覺得.
我們這麼大的年紀.
沒什麼可以做了.
我挑戰他們.
我說今年.
你會不會在神面前有個祈禱.
一個靈魂.
為神來拯救一個靈魂.
當然是神的大能.
但是神通過你.
可不可以呢?.
是不是很困難呢?.
如果你做銷售員.
你的老闆說.
這個月一單生意好不好呀?.

$^{401}$我想你也不會拒絕吧?.
那些我很容易理解.
但是神給我們的使命.
我們好像感受不到.
這是第25節.
不單止這樣.
保羅他多說一點.
這是我覺得很寶貴的26和27節.
我以前都不明白.
我覺得這裡好像和上面說到.
或者28節說的沒有聯繫.
但是我發覺這個是很重要.
26和27節是說一種裝備.
這個裝備裡面.
並沒有說報道法.
保羅說什麼呢?.
首先是說福音的本質.
有人說賣花要讚花香.
其實你要了解.
你在賣的是什麼花.
有些花會令人很容易敏感.
有些花其實是有股臭味.
或者是沒有味道.
我們真的要知道.
自己在說的是什麼.
有時候我們宣講的東西.
自己都不清楚.
甚至我們自己都搞到.
真的是狐狸單刀.
所以傳福音傳完之後.
你沒有幫到他.
你是令到那個人更混亂.
我不是詆毀.
有些人傳完福音之後.
我作為組長.
或者我作為另外一組的組長.
我都要再補妃.
當然有原因.
有些視頻就說到.
有些人身處高位.

$^{441}$想教人怎樣抗疫.
但是他出來的時候.
是調轉了口罩戴的.
你就會覺得.
他知不知道自己在做什麼.
我二十年前左右.
我就開始上課.
上那個叫希伯來文.
去過神學院讀課程.
或者全時間.
兼職都好.
不知道你對希伯來文的感覺是什麼.
聖經公會很好.
他知道我們修讀希伯來文.
他就派一個牧師.
送每個學生.
送一本希伯來文的聖經.
來課室送給我們.
我們很開心.
這裡沒有慢慢地出.
本來答案不給大家.
我的同學拿著一本.
希伯來文聖經去看.
牧師走過來的時候.
就說同學.
這樣看吧.
其實他都很愚蠢.
因為不單止他和牧師知道.
周圍的同學都看到.
因為那時候是第一課.
我們還沒有怎麼讀.
有時你是神學生.
你搞到自己的希伯來文.
都不知道哪一邊看得起.
事實上是很困難的.
你看看.
這是哪一段呢.
是Try及記.
我們有些翻譯了的聖經.
不用和希伯來文搏鬥.

$^{481}$但是你會發覺.
你傳福音不要搞到.
自己都不知道福音是怎麼回事.
你要負點責任好不好.
舉個例子.
你講福音的時候.
你有沒有試過被人問一個問題.
地獄和陰間是不是一樣的.
異端和異教怎麼分呢.
兩個都異的.
差不多.
法輪功是異教還是異端.
耶和華見證是異端還是異教.
等等的問題你要搞得很清晰.
如果不是.
你傳福音就有很大問題.
但是保羅更加想說的是.
這個福音本身的本質再深入一點.
我看它為有三方面.
在這個本質裡面.
他提到這個福音是榮耀的.
福音是榮耀的.
集中所說的介紹的耶穌基督是榮耀的.
他帶來永恆的意義是榮耀的.
榮耀我覺得很難去解釋.
或者用舊約的想法就是.
舊約當榮耀這個字出現的時候.
意思就是heavy.
很重.
有些東西很重你覺不覺得.
你不覺得嗎.
你不覺得跟一些很喜歡做運動的去健身.
給個一百磅的你去拉一下.
再不是給個二百磅的拉一下.
有些東西有重量的感覺.
榮耀是什麼呢.
榮耀是一種有重量的東西.
我們好像不太留意.
什麼是榮耀呢.
舊約這個榮耀曾經出現過一次.

$^{521}$在哪裡出現呢.
就是在《以斯帖記》第一章那裡.
就是當這個阿俠隋老覺得自己一百二十七個省這麼威.
在那裡統管一切.
於是就搞了個大會.
用了一百八十天開大餐.
就在席間裡面將他國家的那些美麗的東西show給那些人看.
你看我們多榮耀這個波斯國.
榮耀是什麼呢.
那個榮耀是不是每一件的物件呢.
不是.
那個榮耀是這一切物件的總和.
或者加上去帶出的那種整個的體會.
什麼叫榮耀.
所以有時候傳福音.
你如果體會耶穌基督的福音.
體會這位的主很少.
體會這個的盼望很少.
其實我們出去講福音呢.
不多吸引的.
你有沒有試過講福音告訴別人.
人家說哦這樣啊.
走開.
但是有些人講福音呢.
講完真的彷彿在流口水一樣.
為什麼.
他體會得到福音那種的寶貴.
在他生命裡面.
他又講到這位主的榮耀.
你知不知道這個就是吸引.
所以弟兄姊妹.
傳福音.
我還沒講完.
傳福音其中一個竅門就是.
你看不體會到福音那種的榮耀.
那個是寶來的.
很多人來到認識基督教.
甚至受浸之後.
他的感覺是什麼.
我已經通過了.

$^{561}$我已經pass了.
到了4月12號浸完之後呢.
我就可以遲些少些回來.
又不用去浸禮班了等等.
錯.
如果你有條鑰匙.
進到一個寶藏裡面.
進去之後並不是你的終結.
是不是.
你應該進去寶藏裡面拼命拿東西.
美國有個節目.
給一分鐘他衝入去超級市場拿東西.
多蠢呢.
那個小朋友多蠢.
他不斷拿廁紙.
廁紙沒用.
但是他的大人就醒目了.
嘩 當然是拿那時候.
早些播的.
拿iPod.
拿DVD player.
拿電視機.
這個就是會不會寶.
神的榮耀我們看到多少呢.
就激勵我們有多大的緊張去傳福音.
我曾經有一次去點貨.
以前我是做核數師的.
就去到勞力士錶廠的倉那裡.
嘩 我人生裡面第一次拿著一隻錶.
港幣 說的是很多年前.
30年前.
那隻錶當時值50萬.
如果幾十年前的50萬.
現在可能數百萬.
嘩 一隻錶.
嘩 覺得很寶.
沒有機會摸.
要買都幾難.
什麼叫榮耀呢.
不只是好像有些.

$^{601}$我有平安.
或者我有什麼好處.
頂紫門.
這個榮耀是超過這件事.
所以當仕途寫啟示錄的時候.
看到天上的榮耀的時候.
不知道你會不會覺得.
不是吧 仕途約翰.
什麼你都不明白.
只有敬畏一位神嗎.
你竟然會企圖想拜那個使者.
大家不明白.
其實那一刻他看到神永恆的榮耀.
他很大的感動.
他真的好像傻傻的.
突然間做了一樣.
他平時不會做的錯.
是他看到榮耀.
保羅認識到.
不斷想講解.
在26 27節就是.
神都很想基督徒知道這個榮耀.
特別是外邦人知道.
有耶穌基督是多麼好.
這個可能都是我們很想他們.
所以就帶進一樣東西.
為什麼我們要.
除了有這個使命.
我們要有些裝備.
最怕頂紫門有時候.
不斷裝備.
但是不做第三樣.
是什麼呢.
報道.
傳播音.
用保羅的話.
傳揚基督.
傳揚基督很難的.
這麼多東西要講.
這麼多東西要學.

$^{641}$但是其實學一樣.
基本功.
以基督做中心.
舉一些例子.
就是說.
你想不同的東西都好.
問一問耶穌基督.
主啊.
這樣好不好呢.
如果昨天.
我很累的時候.
大家說主啊.
明天來不來城會好呢.
芝加哥都有兩單.
這樣.
好不好出街呢.
你問了主之後.
主就說.
不是這樣.
你事奉來的.
你祈禱了很多個月.
都不來.
你知道以基督為中心.
這個東西是要做的.
就不同.
我記得我做青年工作.
有一天不見他回來.
我就call他.
我都很緊貼.
他很怕我.
但是他們很愛我.
知道我關心他.
我是那些令他又愛又恨的人.
我為什麼不見你.
這個morning.
I don't feel like going to church.
不知道這是不是你.
其實你又不可以怪他.
他並沒有以基督為中心.
他覺得.

$^{681}$我這個身體狀態.
昨晚五點鐘才睡到床.
或者今早五點鐘才睡到床.
我怎麼可以在八點鐘又醒呢.
他為了愛惜自己的身體.
又不想自己又感冒.
抵抗力低.
於是他就不想醒了.
他就沒有回來.
所以你見不到他.
你都知道明白.
原因.
又沒有以基督為中心.
弟兄姊妹我們都是.
我們用這個尺度的話.
很容易發覺.
在報道裡面我們差什麼呢.
其實就是要以耶穌基督為中心.
我們有這樣的心.
我們去傳耶穌基督.
我們傳得有力有勁.
在這個內容裡面.
補充兩件事.
我們會面對要勸戒那些人.
我的報道工作就是.
聖經講的那些我都跟著他.
很忌諱.
很難搞.
勸那些人.
不是的.
你都要告訴他.
有些事是他在神面前要面對的.
就是他的罪.
他的背役.
我做小朋友青年工作.
但我亦都要做長者工作.
長者最難搞的就是這件事.
他有幾十年的功力.
犯罪的功力.
但是你要告訴他.

$^{721}$他有罪.
他又不想承認.
那怎麼辦呢.
很多時候都要祈禱.
有些就是.
我幫他查一些聖經.
我祈禱就是神給他感應到一些.
他好像在困惑的地方.
在那一點裡面.
我就會問他.
那你自己會怎樣看.
你有沒有這樣做過.
有時候真的很奇妙.
他就會認.
舉起手.
是的我是這樣.
很奇妙.
勸戒各人.
教導各人.
要將一些聖經的真理告訴他.
弟兄姊妹.
有些真理你不告訴他.
就算他不斷舉手.
不斷點頭.
是不真的相信的.
我舉個例子.
如果有個人不相信天堂地獄.
而你又不解決天堂地獄的真實性.
神的應許是真的.
他就算老老莽莽地舉起手.
他這樣相信.
是不清楚的.
我們叫短裝.
什麼叫短裝.
打樁本應是打50尺.
80尺.
他省錢省脾.
打了30尺算了.
將來會倒的.
信仰很嚴重的.

$^{761}$他不清晰.
那個真理.
就算他表示友好.
信了.
他不清晰.
是一個危險.
在這些裡面.
既然這麼清晰要做.
但是呢.
每個人有不同的狀況.
有一個是學者.
那你用的款式.
你用的approach就不同了.
有些人呢.
他是比較是.
一般的平民老百姓.
那你不要說這麼複雜.
的東西.
你說罪就罪了.
你就不要說.
其他複雜的形容詞了.
所以呢.
保羅加一個豬腳.
這個途徑要怎麼做呢.
豬般的智慧.
不是豬這樣的智慧.
是各種各樣的智慧.
In all wisdom.
所以我提一提.
這個就是訓練.
會幫到我們.
拿回正式的方向.
那個方法.
怎麼去形容.
還有學習對應不同的人.
有不同方式的.
來幫助他.
事實上教導.
這件事是很重要的.
在新約裡面.

$^{801}$用教導這個詞.
用了九十多次.
反而呢.
如果就這樣講道.
好像我們現在.
講道訊息.
這個講道這個詞呢.
是用了三十多次.
在福音書用了十幾次.
似乎呢.
耶穌基督的訊息裡面.
固然是講道.
但是呢.
那個教導的成分.
是非常之重.
是要將那個真理.
不要混淆地告訴他.
讓他清楚.
明白.
你說報道.
挺好的.
其實一個人有使命.
又要做裝備.
又要出去報道.
做不太好.
但是傳遞的報道.
應該是怎樣呢.
全面的應該是怎樣呢.
不可以少了這部分.
就是跟進.
28節下部分.
你留意.
就是他說.
傳揚了之後.
他是要將他.
引導到神的面前.
在基督裡面.
完完全全的引導到神的面前.
講解一下.
首先就是.

$^{841}$這個完完全全.
是指哪樣東西呢.
我的查考.
我發覺這個完完全全.
是在講.
這個被你所帶領.
已經信了主的人.
他應該去到什麼呢.
是一個成熟.
完全的地步.
有些翻譯做perfect.
但是perfect的意思.
不是指到他完美.
因為那個真正的完美.
要到永恆天家.
才是那種最終的完美.
所以現在他那種.
指成熟的完美.
就是他那種指成熟.
mature.
mature和不mature的分別是什麼呢.
就是看小朋友和大人.
不mature的人會怎樣呢.
我碰到不少.
我那時候剛開始工作.
排隊搭車.
突然間有一個大哥哥.
在我以前的教會.
在香港.
走過來.
問候我.
說著說著就說.
你多少錢工資.
已經和他不太熟了.
那他這個是不是不mature.
就是說他在不適當的情況.
不適當的環境.
關係裡面.
問一個不適當的問題.
是什麼問題呢.

$^{881}$不mature.
不成熟.
在我們生命裡面.
如果在靈情.
沒有盡心裡面.
就是不成熟.
那怎樣令他成熟呢.
他已經信了主耶穌了.
他已經有這個永生了.
就是需要跟進.
弟兄姊妹.
這個是很重要的.
今天的教會的增長.
能不能夠真的增長.
能不能夠真的有實際的果效.
就是跟進工作.
所以這個東西.
大進今年.
芝加哥短孫.
我們有一個負擔.
就是開這個叫.
盡心跟進課程.
是很盡心的.
所以每個教會.
都是派幾個.
真的覺得值得栽培的.
才推薦給我們.
去栽培他.
這樣是分階段的.
為什麼呢.
哪裡能夠教這麼多重要的東西.
先教第一階段.
我們現在訂五堂.
這樣就真的.
將最好的東西都儘量教他.
答應你.
精挑細選.
忠心去好好地去教.
讓他在跟進方面突破.
突破那些坑坑.

$^{921}$突破那些困難.
特別是看到有果效的.
你想不想.
你祈禱好不好.
你祈禱想.
主啊,你會不會將感動給我.
我去參與這個工作.
不要等這麼多信了耶穌的人.
就等著跟我們說拜拜.
應該信了的耶穌.
不斷跟我們說Hi才好.
不要整天跟我們說拜拜.
保羅最後我覺得最有力的.
可能你.
有時候我看了.
跟28節有什麼關係呢.
前面保羅是在說我們.
我們這班人.
但是到最後29節.
他在說我自己.
他說自己是奴婦.
為了這個目標.
就是帶那些人.
完完全全去到神的面前.
成熟的地步去到神的面前.
不是帶到他自己的面前.
有時候我們都很喜歡幫自己.
好像做些事情.
但是保羅說不是.
不是為自己.
但是他去拼盡他的生命去做.
盡心竭力.
這個字就是.
也都形容過人.
去爭勝去參加比賽.
出汗了.
沒命的.
就是那個盡心竭力.
有些人真的跑步跑到.
會心臟病發的.

$^{961}$會死的.
當然他應該檢查身體.
沒有亂來.
不要隨便跑半馬拉松.
但是弟兄姊妹.
保羅的生命見證.
我覺得很寶貴.
保羅他看到的.
不是那個行動.
或者參加了多少個報道法.
我想問大家.
有沒有人當中已經參加過.
兩個三個報道法的課程以上.
有沒有.
牧師傳道我知道你們會參加.
但是其他的有沒有.
有我不敢舉手.
不是這麼複雜.
我回想自己以前也參加過報道法.
但是我覺得自己最大的失敗是什麼呢.
就是我不去報道.
就是我這個生命裡面.
很多事情不可告人.
我以前真的有不可告人.
明明在報紙檔.
拿了別人的東西來看.
當然是沒有給錢.
有這些你不對付.
你怎麼跟人講福音.
你去報紙檔告訴他.
你怎麼說.
很困難.
所以在我們生命裡面.
對付.
生命裡面突破.
我們的報道就不同了.
我沒有經常講那一套.
但是.
這個圖我不知道你能不能理解.
我想表達的就是.

$^{1001}$領受即是可能是.
橙色就是你心裡面.
神給過你的負擔.
給你的體見.
你裡面發生的.
基督在你裡面運用的大能.
工作改變你.
使你生命也不是容易看到.
不過.
唯有你有這種裡面生命的改變.
就是一個真正的大的.
這個大的Error.
或者大的箭咀是什麼.
就是一種動力.
你就能夠為這個福音.
可以勞苦拼搏.
就是那個所謂.
可以翻譯做拼搏.
就是好像拼命一樣.
你有沒有覺得.
服侍神你要感覺到.
是要拼命才開心.
才是快樂的.
我就是這樣.
所以有些人說我瘋狂.
但我發覺.
原來你體會到.
服侍神.
依奉耶穌基督是多麼好的時候.
你真的不會覺得累的.
很多時候我發覺.
有一個反定律.
我去報道之前.
心靈掙扎.
有一點點好像很累.
甚至近乎想呼.
但呼不起來.
當去報道.
去和人分享福音的時候.
嘩裡面很踴躍.

$^{1041}$那兩個箭咀就不斷在動.
跟著結果怎樣呢.
病就好火了.
想呼都醫好了.
而且還有多餘力量.
再做多點其他東西.
我發覺事奉主真是寶貴.
我不知道你會不會這樣.
因為我們看到.
我們人生做了些很有價值的事.
姑且這樣說.
幫人.
我用一個.
一兩個分享做結束.
在2000年的時候.
有人做過一個統計.
很有趣的.
他說在2000.
或者說應該是幾年前.
他就拿了2000份.
在報紙裡登出來的那些賦文.
即是說死了人之後.
登了一段東西出來.
這個人統計完之後.
發覺.
奇怪.
最多出現的是什麼字呢.
最多出現就是.
介紹這個人在世的時候.
有多少身家.
不是.
住多大的地方.
又不是.
什麼高級的職位.
又不是.
主要不是說.
最多提到.
差不多必然會提.
就是這個人的一生裡面.
有些什麼貢獻.

$^{1081}$特別有個字詞出現的.
幫助.
弟兄姊妹原來有些東西.
在你離開這個世界.
或者我們離開這個世界.
人家會容易想起的是什麼呢.
你對人的幫助.
你對人真的有些什麼貢獻.
在永恆的角度.
我在想.
我們如果有些東西去列出來.
會不會更加有價值.
在神的眼中.
我們圍到這個國度.
出力.
出錢.
出生命.
那樣東西.
就是神覺得最有價值的東西.
在我去不同地方.
事奉的時候.
有時候我都試過有一次.
很灰心.
在香港.
也都不開心的時候.
又見到以前自己.
二十幾三十年前.
教主的學生.
說著說著.
他就反過來說.
老師.
不要這麼不開心.
然後又請我吃東西.
又對我很好.
突然覺得.
為什麼你這麼好呢.
他還加多一句.
他說老師.
不如你想一想.
你就當你以前.

$^{1121}$經常這麼關心我們.
愛惜我們.
今天我給你一點回報.
是他安慰我.
我突然在想.
我那時候沒有這樣想過.
但原來我們真的.
做在主裡面.
好像保羅說.
是不陶冶.
頂指我不知道你有什麼思想.
總之.
在我總結的時候.
我就用這幾個的關係來總結.
神.
怎樣令我們有一個全面的事奉.
或者在教會裡面.
我們怎樣去看一個全面的事奉.
我相信我們每一個人.
教會性的一個使命.
或者我們個別應該有領袖的使命.
我們需要裝備.
我們不要每個人都覺得自己.
生出來就知道.
報道就是一個報道人.
或者什麼都懂.
我們需要一些裝備.
我們要謙卑學習.
但我們不單止學東西.
我們要出去.
我們要報道.
生了這些屬靈的果子.
我們不要令它成為棄嬰.
我們去學習去跟進.
所以我們需要多些訓練.
去一起去事奉.
但是能夠承托住這四點.
是什麼呢.
是生命.
是我們這個的生命.

$^{1161}$我記得有一個夏令會.
我們去為一個傳道人慶祝生日.
那個傳道人覺得很奇怪.
我幾十年都沒有為自己慶祝生日了.
那個傳道人.
當然了.
他為主坐牢坐了很多年.
他去講福音信息.
他去分享主的慈愛.
他去分享主的信實.
他每一句.
都不是那0.01秒的意思.
他每一秒彷彿是用很多年的時間去孕育.
去告訴你.
今天我就是用這些這麼寶貴的生命.
去跟你分享這個福音.
分享這位耶穌基督.
所以我們那時候年輕人.
講了三十多年前.
我們很感動.
不是他那個神學的解說.
甚至那些是遲早.
而是他這個生命.
所以在他這個生命所承托下的使命.
裝備.
報道跟進.
是無與倫比的.
是你抵擋不到的.
因為他的生命很強.
他要講榮耀什麼.
他講給你很形象化.
很真實的.
神的真實.
神是令我很吸引.
甚至到幾十年的今天.
我想起他去分享的.
我都好像覺得是昨天聽的.
這個是他的生命.
願主使用我們.
我們一起去祈禱.

$^{1201}$親愛的主.
願你再一次.
以你的生命來吸引.
激勵我們.
主啊.
多麼的美好.
我們能夠蒙了這個恩典.
我們又有這個動力.
使命.
奉你的名.
跟眾弟兄姊妹去出去.
去見證你.
當他們願意歸向你的時候.
我們又可以跟他們並肩的成長.
提攜.
讓他們都在主的面前成長.
可以這樣獻醜在主的面前.
因主啊.
願你這樣施恩.
豐富.
強化我們這個生命.
誰聽我們的禱告.
向我們眾人施恩.
無論在我們中間幾年青的.
到幾成熟年長的.
主啊我們都相信.
你願意使用我們.
透過我們的生命.
去見證你.
誰聽我們這樣誠心的禱告.
直立耶穌基督的聖名.
\newpage



\section{使徒行傳 10:1-48}
\label{sec:zCN_cR7XoxQ}
\textbf{2020/2/23 主日講道: 疫境前行 3 (使徒行傳 10\_1-48) - 陳崇基牧師}
\newline
\newline
連結: \href{https://youtube.com/watch?v=zCN-cR7XoxQ}{\texttt{https://youtube.com/watch?v=zCN-cR7XoxQ}} ~~~~ 語音日期: 2020-03-01
\newline
\newline
\hyperref[sec:CWGXYDFB8bU]{\small{< < < PREV SERMON < < <}}
~
\hyperref[sec:index_chronic]{\small{[返順時目]}}
~
\hyperref[sec:index_scriptual]{\small{[返順卷目]}}
~
\hyperref[sec:d0PXI2JCtZU]{\small{> > > NEXT SERMON > > >}}
\newline
\newline
使徒行傳 10:1-48
\newline
\begin{longtable}{cl}
\hline
\hline
章節 & 經文 (和合本修訂版)\\
\hline
10:1 & \begin{tabularx}{0.7\textwidth}{X} 在凱撒利亞有一個人名叫哥尼流，是意大利營的百夫長。 \end{tabularx} \\ \\ \relax
10:2 & \begin{tabularx}{0.7\textwidth}{X} 他是個虔誠人，他和全家都敬畏神。他多多賙濟百姓，常常向神禱告。 \end{tabularx} \\ \\ \relax
10:3 & \begin{tabularx}{0.7\textwidth}{X} 有一天，約在下午三點鐘，他在異象中清楚看見神的一個使者進來，到他那裡，對他說：「哥尼流。」 \end{tabularx} \\ \\ \relax
10:4 & \begin{tabularx}{0.7\textwidth}{X} 哥尼流定睛看他，驚惶地說：「主啊，甚麼事？」天使對他說：「你的禱告和你的賙濟已達到神面前，蒙記念了。 \end{tabularx} \\ \\ \relax
10:5 & \begin{tabularx}{0.7\textwidth}{X} 現在你要派人往約帕去，請一位稱為彼得的西門來。 \end{tabularx} \\ \\ \relax
10:6 & \begin{tabularx}{0.7\textwidth}{X} 他住在一個皮革匠西門的家裡，房子就在海邊。」 \end{tabularx} \\ \\ \relax
10:7 & \begin{tabularx}{0.7\textwidth}{X} 向他說話的天使離開後，哥尼流叫了兩個僕人和常伺候他的一個虔誠的兵來， \end{tabularx} \\ \\ \relax
10:8 & \begin{tabularx}{0.7\textwidth}{X} 把一切的事都講給他們聽，然後就派他們往約帕去。 \end{tabularx} \\ \\ \relax
10:9 & \begin{tabularx}{0.7\textwidth}{X} 第二天，他們走路將近那城，約在正午，彼得上房頂去禱告。 \end{tabularx} \\ \\ \relax
10:10 & \begin{tabularx}{0.7\textwidth}{X} 他覺得餓了，想要吃。那家的人正預備飯的時候，彼得魂遊象外， \end{tabularx} \\ \\ \relax
10:11 & \begin{tabularx}{0.7\textwidth}{X} 看見天開了，有一塊好像大布的東西降下，四角吊著縋在地上， \end{tabularx} \\ \\ \relax
10:12 & \begin{tabularx}{0.7\textwidth}{X} 裡面有地上各樣四腳的走獸、爬蟲和天上的飛鳥。 \end{tabularx} \\ \\ \relax
10:13 & \begin{tabularx}{0.7\textwidth}{X} 又有聲音對他說：「彼得，起來！宰了吃。」 \end{tabularx} \\ \\ \relax
10:14 & \begin{tabularx}{0.7\textwidth}{X} 彼得卻說：「主啊，絕對不可！凡污俗和不潔淨的東西，我從來沒有吃過。」 \end{tabularx} \\ \\ \relax
10:15 & \begin{tabularx}{0.7\textwidth}{X} 第二次有聲音再對他說：「神所潔淨的，你不可當作污俗的。」 \end{tabularx} \\ \\ \relax
10:16 & \begin{tabularx}{0.7\textwidth}{X} 這樣一連三次，那東西隨即收回天上去了。 \end{tabularx} \\ \\ \relax
10:17 & \begin{tabularx}{0.7\textwidth}{X} 正當彼得心裡困惑，不知所看見的異象是甚麼意思時，哥尼流所差來的人已經找到了西門的家，站在門外， \end{tabularx} \\ \\ \relax
10:18 & \begin{tabularx}{0.7\textwidth}{X} 喊著問有沒有一位稱為彼得的西門住在這裡。 \end{tabularx} \\ \\ \relax
10:19 & \begin{tabularx}{0.7\textwidth}{X} 彼得還在思考那異象的時候，聖靈對他說：「有三個人來找你。 \end{tabularx} \\ \\ \relax
10:20 & \begin{tabularx}{0.7\textwidth}{X} 起來，下去，跟他們同去，不要疑惑，因為是我差他們來的。」 \end{tabularx} \\ \\ \relax
10:21 & \begin{tabularx}{0.7\textwidth}{X} 於是彼得下去見那些人，說：「我就是你們要找的人，你們是為了甚麼緣故在這裡？」 \end{tabularx} \\ \\ \relax
10:22 & \begin{tabularx}{0.7\textwidth}{X} 他們說：「百夫長哥尼流是個義人，敬畏神，為猶太全民族所稱讚。他蒙一位聖天使指示，叫他請你到他家裡去，要聽你講話。」 \end{tabularx} \\ \\ \relax
10:23 & \begin{tabularx}{0.7\textwidth}{X} 彼得就請他們進去住宿。次日，他起身和他們同去，還有約帕的幾個弟兄跟他一起去。 \end{tabularx} \\ \\ \relax
10:24 & \begin{tabularx}{0.7\textwidth}{X} 又次日，他進入凱撒利亞，哥尼流已經請了他的親朋好友在等候他們。 \end{tabularx} \\ \\ \relax
10:25 & \begin{tabularx}{0.7\textwidth}{X} 彼得一進去，哥尼流就迎接他，俯伏在他腳前拜他。 \end{tabularx} \\ \\ \relax
10:26 & \begin{tabularx}{0.7\textwidth}{X} 但是彼得拉他起來，說：「你起來，我自己也不過是人。」 \end{tabularx} \\ \\ \relax
10:27 & \begin{tabularx}{0.7\textwidth}{X} 彼得和他一邊說話一邊進去，見有好些人聚集， \end{tabularx} \\ \\ \relax
10:28 & \begin{tabularx}{0.7\textwidth}{X} 就對他們說：「你們知道，猶太人和別國的人結交來往本是不合規矩的，但神已經指示我，無論甚麼人都不可看作污俗或不潔淨的。 \end{tabularx} \\ \\ \relax
10:29 & \begin{tabularx}{0.7\textwidth}{X} 所以，我一被邀請，沒有推辭就來了。現在請問，你們為甚麼叫我來呢？」 \end{tabularx} \\ \\ \relax
10:30 & \begin{tabularx}{0.7\textwidth}{X} 哥尼流說：「四天前，這個時候，我在家中守著下午三點鐘的禱告，忽然有一個人穿著明亮的衣裳站在我面前， \end{tabularx} \\ \\ \relax
10:31 & \begin{tabularx}{0.7\textwidth}{X} 說：『哥尼流，你的禱告已蒙垂聽，你的賙濟在神面前已蒙記念了。 \end{tabularx} \\ \\ \relax
10:32 & \begin{tabularx}{0.7\textwidth}{X} 你要派人往約帕去，請那稱為彼得的西門來，他住在海邊一個皮革匠西門的家裡。』 \end{tabularx} \\ \\ \relax
10:33 & \begin{tabularx}{0.7\textwidth}{X} 所以我立刻派人去請你。你來了真好。現在我們都在神面前，要聽主吩咐你的一切話。」 \end{tabularx} \\ \\ \relax
10:34 & \begin{tabularx}{0.7\textwidth}{X} 彼得開口說：「我真的看出神是不偏待人的。 \end{tabularx} \\ \\ \relax
10:35 & \begin{tabularx}{0.7\textwidth}{X} 不但如此，在各國中那敬畏他而行義的人都為他所悅納。 \end{tabularx} \\ \\ \relax
10:36 & \begin{tabularx}{0.7\textwidth}{X} 神藉著耶穌基督—他是萬有的主—傳和平的福音，把這道傳給以色列人。 \end{tabularx} \\ \\ \relax
10:37 & \begin{tabularx}{0.7\textwidth}{X} 這話在約翰傳揚洗禮以後，從加利利起，傳遍了猶太。 \end{tabularx} \\ \\ \relax
10:38 & \begin{tabularx}{0.7\textwidth}{X} 神怎樣以聖靈和能力膏了拿撒勒人耶穌，這都是你們知道的。他到處奔波，行善事，醫好凡被魔鬼壓制的人，因為神與他同在。 \end{tabularx} \\ \\ \relax
10:39 & \begin{tabularx}{0.7\textwidth}{X} 他在猶太人之地和耶路撒冷所行的一切事，有我們作見證人。他們竟把他掛在木頭上殺了。 \end{tabularx} \\ \\ \relax
10:40 & \begin{tabularx}{0.7\textwidth}{X} 第三天，神使他復活，使他顯現出來； \end{tabularx} \\ \\ \relax
10:41 & \begin{tabularx}{0.7\textwidth}{X} 不是顯現給所有的人看，而是顯現給神預先所揀選為他作見證的人看，就是我們這些在他從死人中復活以後和他同吃同喝的人。 \end{tabularx} \\ \\ \relax
10:42 & \begin{tabularx}{0.7\textwidth}{X} 他吩咐我們傳道給眾人，證明他是神所立定，要作審判活人、死人的審判者。 \end{tabularx} \\ \\ \relax
10:43 & \begin{tabularx}{0.7\textwidth}{X} 眾先知也為這人作見證：凡信他的人，必藉著他的名得蒙赦罪。」 \end{tabularx} \\ \\ \relax
10:44 & \begin{tabularx}{0.7\textwidth}{X} 彼得還在說這些話的時候，聖靈降在一切聽道的人身上。 \end{tabularx} \\ \\ \relax
10:45 & \begin{tabularx}{0.7\textwidth}{X} 那些奉割禮的信徒和彼得同來，見聖靈的恩賜也澆在外邦人身上，就都驚奇； \end{tabularx} \\ \\ \relax
10:46 & \begin{tabularx}{0.7\textwidth}{X} 因聽見他們說方言，稱讚神為大。於是彼得回答： \end{tabularx} \\ \\ \relax
10:47 & \begin{tabularx}{0.7\textwidth}{X} 「這些人既受了聖靈，跟我們一樣，誰能阻止用水給他們施洗呢？」 \end{tabularx} \\ \\ \relax
10:48 & \begin{tabularx}{0.7\textwidth}{X} 他就吩咐奉耶穌基督的名給他們施洗。於是他們請彼得住了幾天。 \end{tabularx} \\ \\
[1ex]
\hline
\hline
\end{longtable}
$^{1}$那麼我們回到聖經,好嗎?.
我們繼續我們這個系列叫做《逆境前行》.
我們去《史萊姆傳》第十章.
第一節開始.
這裡聖經這麼記載.
他說在該撒尼亞有一個人叫做哥尼留.
是意大利營的白夫長.
他在那裡做了一件事.
就是他在那裡做了一件事.
他在那裡做了一件事.
他是虔誠的人.
他和全家敬畏神.
對人民行過許多善事.
常常向神禱告.
有一天大約下午三點鐘.
在異象中.
他清清楚楚看見神的一位天使.
來到他那裡.
對他說:哥尼留.
他聽見一看害怕.
就說:主呀,什麼事?.
天使說:你的禱告和善行.
已經達到神面前.
蒙他紀念了.
現在你要派人到約柏去.
請那名叫彼得的西門來.
他在一個祭陪工人西門的家裡作客.
房子就在海邊.
和他說話的天使走了之後.
他就叫了兩個家僕.
事後他的一個虔誠士兵來.
把一切事向他們講明.
然後派他們到約柏去.
我們讀這八節.
這個故事是使徒函傳第十章.
我們知道新約聖經有開頭四卷書.
馬太,馬可,路加和約翰福音.
這四本約翰福音都可以叫做歷史的題材.
是介紹耶穌的出生和生平.
他死而復活甚至升天.

$^{41}$四本福音書.
第五本書就是我們今天看的使徒函傳.
使徒函傳裡面講什麼呢.
是講到初期教會的報道.
宣教的歷程.
如果你看下一張的slide的話.
就是一張地圖.
我想在這張地圖找三個地方.
該撒尼亞,約柏和耶路撒冷.
我想在這張地圖找三個地方.
找到吧.
很容易吧.
找到這三個地方就知道.
大概這個故事是環繞這三個地方.
當時初期教會開始之後.
他們在耶路撒冷開始.
然後到其他地方擴散.
In fact在聖靈降臨之後.
聖靈降臨之後.
他就empower當時的教會.
去到猶太傳遞.
就是耶路撒冷北面的地帶.
然後去到撒馬利亞.
去到其他地方.
聖靈就要他們.
將福音,教會.
一路直傳,發展.
所以這個地圖讓我們看到.
大概這三個的三點.
就是我們今天看的故事.
有一個羅馬的將軍.
叫做 Cornelius.
他信主的經歷.
或者那個故事是平衡的.
又有另一個故事.
就是彼得教會的領袖.
他看到一個異象.
這個故事就是環繞這三個地方.
如果你看到這個地圖的時候.
不僅找到這三點.

$^{81}$你也找到以色列這個地方.
其實沒有什麼天然港口.
你看到那個海岸線.
其實也是很直.
不是很深.
所以在以色列的這個地方.
沒有天然港口.
一切的港口都是比較小的.
約柏.
其實就是一個比較小的.
從舊約到新約.
一個小的港口.
到了羅馬的時期就有一點不同.
公元前大概22年的時候.
就是22 B.C.的時候.
有一個當時幫羅馬帝國.
去守這塊地的.
那個叫做大希倫.
Helder Great.
22 B.C..
他就找了一個地方.
一個小的漁港.
然後將它擴大變成一個很大的港口.
叫做.
蓋薩尼亞石屎利亞.
為什麼他可以將一個.
天然的小海港.
然後變成一個很大很大的地中海的海港呢.
因為當時.
羅馬帝國他們發明了.
什麼叫做石屎.
什麼叫做concrete.
有沒有用過concrete.
你坐的地方其實就是concrete.
站的地方是嗎.
就是當時剛剛開始.
發明出來的.
就是有沙和泥混在一起.
放水進去會變硬的.
所以大希倫就用了.

$^{121}$這個concrete的.
建築的方法.
就做了一個.
很大的人工的海港.
叫做Seserea.
就好像.
香港那樣.
當它是地中海的香港那樣.
又可以有貨櫃碼頭或者大的.
漁船,大的貨船.
可以賣到岸的.
所以他就將這個海港變成一個.
大的人工海港.
就將那些錢塞進來.
將那些貿易的錢帶進來.
以色列他管的那個地方.
就將這個海港.
叫做Seserea.
他老闆就叫Seserea.
Caesar.
你看到一個很羅馬的.
帝國風格的一個地方.
有賽馬場.
有劇院.
有人工海港.
有一些遺跡.
在海旁上.
有一個宮殿.
等等.
這個故事就發生在這個地方.
在Seserea.
就是說有一個意大利的.
伯夫長.
伯夫長是什麼意思呢?.
就是說他管一百個軍人.
或者有多.
或者有少.
總之是一個大的將軍.
他就駐守在這個.
所謂愛邦人為主的.

$^{161}$Seserea的海港.
也提到.
這個人經常去祈禱.
以及奉獻.
並且去敬畏神.
看到他.
早上三點鐘.
他就按照猶太人.
下午的祈禱時間.
做了典型猶太人.
會做的事.
不過他自己本身當然是愛邦人.
猶太人怎麼看人呢?.
我記得.
我以前跟你講過.
猶太人看人.
好像我們今天看的舊的.
純種的特別貴.
雜種的就便宜一點.
人就是這樣.
猶太人眼中的人.
如果他是.
完全猶太人.
純種的.
Jewish people.
他們是守律法的.
他們是信奉耶和華的.
獨一的神.
這些就是純種的人.
猶太人.
在他們眼中猶太人以外.
分幾樣的.
第一樣就是一些.
所謂的歸信者.
他們是依了猶太教的.
那些是非猶太人血統的人.
也是.
他們信耶和華.
雖然他們本身的血統不是猶太人.
他們也守律法.

$^{201}$行割禮.
食物的條例等等.
上一次我們講了.
飲食的culture law.
或者一些節期.
特別是Shabbat.
再這個以外.
就是那些.
沒有受割禮.
或者沒有傳授律法的.
都是非猶太人的血統.
但是他覺得.
這個神是獨一的神.
所以他來認識他.
只不過他們沒有.
傳授律法.
這些就叫做虔誠人.
再外面那些.
就是他們不接受.
猶太教的信仰.
也不是猶太教血統的人.
我想問你.
God fears 是哪一種人.
他是Jews.
Converts.
God fears 或者.
Gentiles.
是哪一種呢.
大家很聰明吧.
就是虔誠人.
God fears.
那些是信守.
猶太教的信仰.
只不過他還沒有.
去守一些.
沒有一些律法的移民.
就是一些條文.
他有一點像.
今天我們一些.
所謂seekers.

$^{241}$一些在教會裡面.
他可能覺得.
這個聖經裡面的神.
才是那個神.
才是那個主.
但是他還沒有去到一個地步.
去到相信.
不過他又會來教會.
他會參加這裡的活動.
他心裡面可能似乎覺得.
會不會可能我以前的背景.
是無神也好,不可知論者也好.
或者我是.
其他信仰的也好.
這些人可能都有一個共通點.
共通點可能就是.
What if I am wrong.
會不會是我看錯了.
看人生.
會不會這個神才是真神.
是一個創造主.
以耶穌所說.
他自己是誰的時候.
可能耶穌真的不是這麼簡單.
不是純粹一個好的教師.
很有智慧的人.
可能他有一些很特別的地方.
他吸引了他們.
哥尼留大概就是一個這樣的人.
God fears.
In fact,在當時的猶太的世界裡面.
有一類這樣的人.
他們有時候會這樣形容他們.
一些所謂虔誠人.
或者righteous gentiles.
他們是愛邦人的血統.
不是猶太人的血統.
他們來尋求神.
在第十章的故事.
在該撒尼亞的地方.

$^{281}$哥尼留就在禱告裡面.
他就見到天使.
天使就告訴他.
喂,你已經得到遇納了.
請你派人去約柏.
你知道他在該撒尼亞.
他派人南下.
去約柏找一個人.
就是找Peter.
他在什麼地方呢?.
就在一個製皮工人的家裡面.
那個人的家.
大概.
accidentally意外地.
又是叫西門.
彼得又叫西門.
所以他就在那裡作客.
他的airbnb就在那裡.
他就留在那裡.
由該撒尼亞南下.
去到約柏.
約柏.
大概30英里.
現在開車大概一個小時.
我們知道當時的約柏港.
也有一些猶太人在那裡.
但這裡很特別的地方.
就是他去到一個人的家.
那個人的家.
是一個workshop.
他做什麼生意呢?.
是做.
可能你會覺得是皮草.
其實是皮的工廠.
製皮的工廠.
那我想問製皮需要什麼呢?.
需要皮,ok.
sure.
怎麼樣?.
你需要很多東西.

$^{321}$製皮是一定需要的.
以前古代的時候製皮.
就是水.
因為很臭.
所以在海邊.
有一個製皮的工廠.
有一個公司.
所以就派人去找他.
已經給了地址.
你聞到臭味就知道那裡就是.
或者近海的地方.
但這裡特別要留意的就是.
第十章這裡特別說到.
這是製皮工人的家.
也是近海的.
by the sea.
在猶太人的地理觀念裡面.
約柏港應該是海港.
海港的一個城鎮.
製皮工人的工廠就在海邊.
在猶太人的地理觀來講.
這個就真的去到地極了.
因為從耶路撒冷出發.
去到約柏港.
去到約柏港的海邊.
這個就去到地的極處.
對於猶太人來說.
所以這個故事.
就去到在彼得眼中.
或者猶太人眼中的一個地極的地方.
第九節.
第二天大約正午的時候.
他們走近這個城.
被哥利奈派去的工人.
走近這個城的時候.
他們看到彼得上了房頂去祈禱.
這個就是他們的地極.
他們的房頂去祈禱.
這是另外一個故事.
他說那家人呢.

$^{361}$西門製皮的家人正在預備飯的時候.
當你聞到飯香你會覺得.
你會覺得肚餓.
彼得就覺得肚餓了.
很想吃飯.
誰知道這個時候他祈禱的時候.
就雲遊障礙.
他看到天上有一件東西.
好像一塊大布綁著四角.
然後放在地上.
有地上各種的四足生畜.
還有昆蟲和天空的飛鳥.
然後有聲音對他說.
他說彼得起來宰鳥吃.
然後煮千萬不可.
他這麼說.
我從來沒有吃過這些不潔的東西.
上次你記得我說過.
你沒記錯的一張.
猶太人有的飲食裡面.
有一些不潔的飲食.
現在他看到的是什麼.
是一些混雜了的食物.
當然對彼得來說.
這些不是不潔的東西.
他是不會吃的.
第二次有聲音.
他就對他說.
神所潔淨的.
你不可當作植物.
這樣一連三次.
這個東西立刻就收回天上.
一個很奇怪的現象.
三次看到一些.
潔與不潔的一切.
放在一張桌布上.
跌下來叫他吃.
第二次.
又第三次.
在約伯港那裡.

$^{401}$今天約伯港.
有一個.
聖彼得教堂.
一個天主教教堂.
裡面有一幅畫.
祭壇上有一幅畫.
這幅畫可能幫到大家.
去想像到底發生了什麼事.
彼得,天使和上面什麼人.
有一塊布.
包著一些.
可能是魚蝦蟹.
一些猶太人不會吃的東西.
你記得上次我們提過.
猶太人的culture(文化)的law(法律).
dietary law(食物法律).
是他們的identity marker(標誌).
就是猶太人.
吃什麼不吃什麼.
是他們身份的象徵.
今天的猶太人也是一樣.
我沒多久之前認識了一個猶太的基督徒.
那些miss Anna Jew.
她開始去了一個教會.
她說那個教會有一次.
教potluck(聚餐).
那裡的教會有些亞洲人在.
也有些華人在.
誰知道他們.
弄了一隻燒豬來.
大家很開心.
但她整家人就沒有去.
那個聚會.
說我們不去.
不要叫她吃.
聞都不想聞.
這個就是他們的identity marker(標誌).
所以你可以想像.
對彼得來說是很大的震撼.
他去到這個地方.

$^{441}$約伯港.
一個最邊緣最邊緣的地方.
最地極的地方.
不可以再出去.
這個已經突破了.
他很多的.
所謂種族上的.
泛地.
種族上的.
一些barrier(障礙).
其實那個故事.
大概的重點也是一樣.
就是上帝.
要彼得.
去到突破他.
種族的界限.
因為上帝.
叫潔淨的.
你不可以當為族民.
留意那句話.
神所潔淨的.
你.
不可以為族.
就是上帝要彼得.
以他的角度.
他普世的角度.
來看事物.
上帝要你.
用他的方法.
他的perspective(角度).
來看人事與物.
我不知道你有沒有留意.
有時候我們教會.
走在一起,我們崇拜的時候.
我們音樂唱歌.
有些歌詞.
有一個是這樣說的,他說要依靠神.
不要依靠我們自己.
我們.
有時候會想.

$^{481}$可能會錯覺認為.
依靠神不用依靠自己.
上班就依靠神就行.
就不用工作了.
讀書就依靠神就行.
就不用工作了.
當然不是.
什麼叫依靠上帝.
什麼叫depend on God.
很多時候去以色列的時候.
我們帶團.
跟弟兄姊妹去看.
不同的地方.
我們一定會上山.
每逢去一個地方,你盡量找一個高的地方.
那個可能是一個小的山丘.
或者是一些.
瞭望到.
更加大的.
景色的地方.
我那時候就要爬上山.
在盆石上面.
盆石不是我們玩石頂水那些.
一粒粒的春石.
盆石是那些.
的foundation rock.
是那些石山.
山頂,那個叫盆石.
你在盆石上面才能夠.
看到東西更加清楚.
就要用一個這樣的角度.
去看事物.
依靠上帝其實是什麼意思.
就是用上帝的角度來看.
人事與物.
很多時候上帝.
就是要擴大我們的框架.
我們每個人都要有.
一點點框架.
我們會限制上帝想做的事情.

$^{521}$我們認為上帝一定要這樣做.
一定要這樣聽我祈禱.
一定要這樣答應我,這樣才算是.
神的,是不是.
有我們的框架,猶太人當然他的框架.
很多時候就是種族上面的框架.
我不知道你的框架是什麼.
但是這裡很肯定.
上帝就是要讓.
去突破他自己的框架.
說起種族的框架.
其實我們看到現在.
這個新冠狀武漢肺炎.
其實種族這件事.
又再一次.
就放上台了.
又成為一些很爭論性的.
一些的事情.
在我們現在這麼多.
疫情的裡面.
疫症的裡面.
其實真的有些困難.
一個人在哪裡來.
如果他在疫區來.
我們會覺得.
可能有一些的框架.
我們會另一些眼光.
來看他.
有人弄了這張這樣的東西.
我在Facebook看到.
就是兩班那些小混混.
小孩子你哪裡來的.
他說我是.
什麼.
我是鴻興的.
他說我剛剛武漢回來的.
他們馬上都走了.
這些都是.
我們需要去.
想清楚在疫症裡面.

$^{561}$或者疫症之後.
我們怎麼去看其他的人.
近日在紐約的地鐵站.
大家聽過這件事的新聞.
有一個黑人去襲擊一個.
戴口罩的華裔婦女.
罵她為什麼將這個病毒帶進來.
在現在美國的時代.
特別近這幾年.
很多這些很極端的.
種族的問題.
就出現.
最近唐人街的.
劫殺案.
都是.
令我們想起種族的問題.
所以為什麼我們鼓勵大家.
華人.
在這裡不要躲起來.
我們需要被控制.
我們需要知道我們在這裡.
其他人要知道我們在這裡.
也不要只顧自己.
所以當彼得.
去到Cornelius的家裡的時候.
他就聽見.
Cornelius原來有上帝.
給他的訊息.
原來上帝也接納外邦人.
彼得就說什麼呢.
他說原來神不偏待人.
在各民各族裡面.
凡敬畏他的人.
都蒙上帝自己的月亮.
這個對彼得來說.
是很稀奇的事.
聖靈居然降臨在他們身上.
一些God fearer.
一些Gentiles.
一些沒有受到國禮.

$^{601}$沒有按照猶太人的規矩的人.
所以他們非常稀奇.
換句話說.
上帝要做事.
上帝要出手的時候.
就不一定要按照你的框架來做.
出乎你意料之外.
你就需要用這樣的角度.
來看事物.
到第十一章.
不單止彼得要學新的功課.
其實耶路撒冷的教會也要.
十章就說到.
哥利留信主.
然後受洗.
他歸主之後.
一個外邦人歸主.
能夠成為神的子民.
回到耶路撒冷匯報的時候.
就很奇怪.
為什麼那裡的弟兄姊妹.
居然不太接受.
十一章第一節.
使徒和猶太的弟兄.
聽見外族人.
也接受了神的道.
有上帝的聖靈.
彼得上去耶路撒冷的時候.
那些守國禮的與他爭論.
他說你竟然.
跟未受國禮的人.
在那裡跟他們一起吃飯.
Again.
Again,這個就是猶太人的框架.
跟他們吃飯.
就是說.
令到你自己不潔.
令到我們不潔.
於是彼得就需要跟他們解釋.
你想想猶太人框架就是.

$^{641}$猶太人有一種潔癖.
他不想跟一些非猶太裔的人一起.
是.
嗯.
這個是上帝要他們去突破的.
初期的教會.
要見到上帝.
自己的工作.
其實這個也是史上傳.
上文夏利的.
大的圖畫.
下一張你就會看到.
上文夏利發生什麼事呢.
第九章,保羅就信主.
外邦人的信徒.
Right.
然後我們就講到今天.
這個哥尼留信主彼得傳道.
第十章至十一章開頭.
到十一章第十九節.
就講到安提柯教會.
安提柯教會是北面.
另一個外邦人的一個大的海港.
那裡就建立了第一個多元化.
多文化的一個教會.
教會是需要學功課.
初期教會.
今天的教會.
都是一樣.
要突破我們的框架.
嗯.
那裡的教會.
就有很多不同種族的人走在一起.
結果.
結束的時候.
我們看看安提柯教會.
安提柯.
他說有些被迫害的信徒.
他們四散的時候.
走到不同的地方.

$^{681}$有些叫做菲利基,有些叫做塞浦路斯.
並且去到安提柯.
這些猶太人的信徒就不跟其他人講道.
只是跟猶太人講道.
但是.
裡面有一些人.
也跟一些希臘人.
傳耶穌的福音.
主跟他們同座.
信主.
歸主的人就多.
嗯.
你想想.
如果教會每個人都跟我們同聲同氣.
好不好.
好.
I think it's good.
但是不是great嗎.
It's not great.
It's good.
but not great.
我們教會是不是一個good的教會.
是不是一個good church.
I think so.
Right.
但我們是不是great church.
就是看到上帝所看到的.
看到上帝去突破我們很多的框架.
包括種族的框架.
語言的框架.
我們的興趣的框架.
讓我們看遠一點,看到地極.
希望你能夠繼續為教會祈禱.
甚至有些弟兄姊妹.
參與我們英文的.
英文課的試學.
Right.
其他語言為什麼不行.
我們這裡兩文三語.
我們每個人都可以.

$^{721}$就是我們突破了.
我們才能夠繼續向前.
走到地極.
讓上帝突破你的框架.
以致我們教會能夠成為一個.
更加土俱起源的教會.
好.
\newpage



\section{詩篇 27:1-14}
\label{sec:d0PXI2JCtZU}
\textbf{2020/2/9 主日講道: 恐懼中, 記著這一件事 (詩篇27\_1-14) - 張麗仙傳道}
\newline
\newline
連結: \href{https://youtube.com/watch?v=d0PXI2JCtZU}{\texttt{https://youtube.com/watch?v=d0PXI2JCtZU}} ~~~~ 語音日期: 2020-03-10
\newline
\newline
\hyperref[sec:zCN_cR7XoxQ]{\small{< < < PREV SERMON < < <}}
~
\hyperref[sec:index_chronic]{\small{[返順時目]}}
~
\hyperref[sec:index_scriptual]{\small{[返順卷目]}}
~
\hyperref[sec:EN7WJcnMP5Y]{\small{> > > NEXT SERMON > > >}}
\newline
\newline
詩篇 27:1-14
\newline
\begin{longtable}{cl}
\hline
\hline
章節 & 經文 (和合本修訂版)\\
\hline
27:1 & \begin{tabularx}{0.7\textwidth}{X} 耶和華是我的亮光，是我的拯救，我還怕誰呢？耶和華是我生命的保障，我還懼誰呢？ \end{tabularx} \\ \\ \relax
27:2 & \begin{tabularx}{0.7\textwidth}{X} 那作惡的就是我的仇敵，前來吃我肉的時候就絆跌仆倒。 \end{tabularx} \\ \\ \relax
27:3 & \begin{tabularx}{0.7\textwidth}{X} 雖有軍隊安營攻擊我，我的心也不害怕；雖然興起戰爭攻擊我，我仍舊安穩。 \end{tabularx} \\ \\ \relax
27:4 & \begin{tabularx}{0.7\textwidth}{X} 有一件事，我曾求耶和華，我仍要尋求，就是一生一世住在耶和華的殿中，瞻仰他的榮美，在他的殿宇裡求問。 \end{tabularx} \\ \\ \relax
27:5 & \begin{tabularx}{0.7\textwidth}{X} 因為我遭遇患難，他必將我隱藏在他的帳棚裡，把我藏在他帳幕的隱密處，將我高舉在磐石上。 \end{tabularx} \\ \\ \relax
27:6 & \begin{tabularx}{0.7\textwidth}{X} 現在我得以昂首，高過四面的仇敵。我要在他的帳幕裡歡然獻祭，我要唱詩歌頌耶和華。 \end{tabularx} \\ \\ \relax
27:7 & \begin{tabularx}{0.7\textwidth}{X} 耶和華啊，我呼求的時候，求你垂聽我的聲音；求你憐憫我，應允我。 \end{tabularx} \\ \\ \relax
27:8 & \begin{tabularx}{0.7\textwidth}{X} 你說：「你們當尋求我的面。」那時我的心向你說：「耶和華啊，你的面我正要尋求。」 \end{tabularx} \\ \\ \relax
27:9 & \begin{tabularx}{0.7\textwidth}{X} 求你不要轉臉不顧我，不要發怒趕逐你的僕人，你向來是幫助我的。救我的神啊，不要離開我，也不要撇棄我。 \end{tabularx} \\ \\ \relax
27:10 & \begin{tabularx}{0.7\textwidth}{X} 即使我的父母撇棄我，耶和華終必收留我。 \end{tabularx} \\ \\ \relax
27:11 & \begin{tabularx}{0.7\textwidth}{X} 耶和華啊，求你將你的道指教我，因我仇敵的緣故引導我走平坦的路。 \end{tabularx} \\ \\ \relax
27:12 & \begin{tabularx}{0.7\textwidth}{X} 求你不要把我交給敵人，遂其所願；因為妄作見證的和口吐凶言的都起來攻擊我。 \end{tabularx} \\ \\ \relax
27:13 & \begin{tabularx}{0.7\textwidth}{X} 我深信在活人之地必得見耶和華的恩惠。 \end{tabularx} \\ \\ \relax
27:14 & \begin{tabularx}{0.7\textwidth}{X} 要等候耶和華，當壯膽，堅固你的心，要等候耶和華！ \end{tabularx} \\ \\
[1ex]
\hline
\hline
\end{longtable}
$^{1}$各位觀眾 今天是3月8日婦女節.
大家以為我今天是講聖經中的女性.
但今天我想跟大家講另一個主題.
思考另一個題目.
是跟男女老少有關.
就是一個字 驚.
當我預備今天的講道的時候.
我看到新聞 第三宗 第四宗 第五宗 第六宗.
而且是越來越近的 在我們伊利諾州這裡.
所以老實說其實都有點害怕.
說真的 人生本身就已經有很多害怕.
我們會怕出門忘記關火 怕撞車 坐飛機怕撞機.
還有就是怕失業 不但沒有了薪金 而且沒有了保險才會拿命.
還有怕打劫 怕患了重症 怕關係破裂.
有很多很多事情我們真的怕都怕不完.
但真的怎樣都想不到來這個驚就是一場瘟疫.
真的嚇死了.
但在這個驚的時候我們有很多不同的反應.
我們會頭暈 會失眠 會胃痛.
很多時候都是因為害怕的反應.
害怕其實本身不是壞事.
因為害怕是令我們提高警覺.
但最壞的就是我們被害怕操縱了.
例如做了這些事情出來.
真的很驚天.
又例如這樣.
當我們不理性地害怕的時候.
會越來越害怕 會製造很多恐慌.
我就想起一個很著名的美國作家Max Lucado.
他說了一句這樣的說話.
他說恐懼會探訪每一個人.
你可以讓他做一個訪客.
但不要讓他在你的心裡做一個住客.
這句說話我真的覺得很對.
我聽說恐懼的時候.
我們的身體腎上腺分泌會令我們變得很軟.
而且恐懼不合而處理.
會令我們更加害怕 失去理智.
說到認識恐懼和處理恐懼的情緒.
我想留給Zoe來說吧.

$^{41}$作為牧者 我今天想跟大家說的是一件很重要的事.
就是在恐懼裡面有一件很重要的事.
以致我們知道這件事.
無論什麼事情都會發生.
我們都不會活得像一個沒有盼望的人.
很想跟大家一起讀這篇詩篇.
也是很熟悉的詩篇 我們一起讀吧.
14節不是很長.
耶和華是我的亮光 是我的拯救 我還怕誰呢.
耶和華是我性命的保障 我還懼誰呢.
那作惡的就是我的仇敵.
前來吃我肉的時候 就半跌倒.
雖有軍兵安然攻擊我 我的心也不害怕.
雖然興起刀兵攻擊我 我必仍舊安穩.
有一件事 我曾求告耶和華.
我仍要尋求.
就是一生一世住在耶和華的殿中.
沾揚他的榮美 在他的殿內求民.
因為我遭遇患難 他必暗暗地保守我.
在他亭子裡 把我藏在他帳幕的隱密處.
將我高舉在盤石上.
現在我得意昂首 高過四面的仇敵.
我要在他的帳幕裡 歡言獻祭.
我要唱詩歌頌耶和華.
耶和華啊 我用聲音呼籲的時候 求你誰聽.
求你連恤我 應運我.
你說 你們當尋求我的面.
那時我心向你說.
耶和華啊 你的面我正要尋求.
不要向我掩面 不要發怒趕逐僕人.
你向來是幫助我的 救我的神啊.
不要丟棄我 也不要離棄我.
我父母離棄我 耶和華必收留我.
耶和華啊 求你將你的道指教我.
因我仇敵的緣故 引導我走平坦的路.
求你不要把我 交給敵人 遂其所願.
因為妄作見證 和口吐窮言的 起來攻擊我.
我又不信在活人之地 得見耶和華的恩惠.
就早已喪膽了 要等候耶和華.
當喪膽 堅固你的心 我在說 要等候耶和華.

$^{81}$我相信這個詩篇 弟兄姊妹你們有很多都讀過.
但你們讀的時候 有沒有留意到這篇詩篇很特別.
就是裡面有很多很多 害怕的 很多恐懼.
有沒有留意到那些字眼呢.
我們一起來看看 我highlight了黃色的.
有怕字 有懼字 又說到我仇敵.
仇敵說了很多次 仇敵來踐踏吃我肉.
軍兵還扎型來攻擊我.
但我又說到害怕 攻擊我.
然後在11節繼續說 仇敵敵人.
口吐狂言的 起來攻擊我 喪膽.
看清這篇詩篇 其實有很多很害怕的東西.
但很妙的是 這篇詩篇 也被稱為信心的詩篇.
聖經很多時候是很吊詭的.
例如說 說到黑暗的時候 就有光明.
說到狂風巨浪的時候 就有平安.
說到死亡的時候 就有生命.
說到恐懼的時候 就有信心.
事實上 詩篇27篇是一個很好的心路歷程.
一開始的時候 第一句就已經說 耶和華是我的亮光.
說到耶和華是我的亮光的時候.
我想起前幾天在網上看到一段說話.
一個禱告 是一個武漢 在水深火熱的城市.
一位武漢的牧師所說的話.
他說 在有光明的日子.
我體會不到自己的有限和軟弱.
但是在黑暗的日子.
我突然間發覺到 原來人是金毛座的.
我們不敢向神呼喊求祂幫助.
這個牧師說 在黑暗的日子.
我呼喊求神幫助.
我相信他這個心情 十足是大衛當時的心情.
他說 這個是這個牧師所說的.
大衛當時所說的就是說.
耶和華是我的亮光 是我的拯救.
因為在《懼怕》裡面 大衛他經歷了這件事.
他說 耶和華是我的亮光 是我的拯救.
這兩句說話其實是同一件事.
原文是 光和拯救是等於一樣.
大衛體會到 在黑暗的日子裡面.

$^{121}$就好像這個武漢的牧師所經歷到的.
今日日子真的很黑暗.
我真的摸不到去處.
我只是感受到敵人攻擊我.
我甚至經歷到好像死陰幽谷一樣.
但是在這個時候 大衛說.
我體會到有光照進來.
那個光是神給我的拯救.
大衛不單止感覺到.
神是我性命的保障.
英文是Stronghold.
原文的意思是 是一個城堡.
神是我的城堡是什麼意思呢.
就是說敵人攻擊的時候.
我有一個堅固的城牆圍繞著 我不怕.
這個是大衛當時面對他驚懼的情況.
受敵人攻擊的情況 他所感受到的.
神是我的光和拯救.
神是我性命的保障.
我還怕誰呢.
你繼續看他這些說話的時候.
他說 我的受敵前來吃我的肉.
我的受敵前來吃我的肉.
我的受敵還扎型攻擊我.
這些都是非常時期.
真的敵人打過來.
我的輕兵 輕起他的刀劍來攻擊我.
但是我仍舊安穩.
為什麼大衛會經歷到這些呢.
他一直說 因為是耶和華.
耶和華 在詩篇裡面時時都說到耶和華.
耶和華是神.
但是他不是普通的神明.
他是與以色列 與大衛的先祖立約的神.
大衛回顧這位耶和華.
他過去歷世歷代的跟他們立約.
說我是你們的神 你是我的子民.
你忠於我 我就保佑你.
而且歷世歷代 他完全沒有失過約.
因為這個緣故.

$^{161}$大衛再一次經歷到.
在他仇敵前來吃他的肉的日子.
他體會到耶和華的幫助.
耶和華在黑暗的日子裡.
是我的亮光 是我的拯救.
在敵人打過來的時候.
他是我堅固的城堡.
這是大衛所經歷到的.
弟兄姊妹 你有沒有經歷到呢.
詩篇我們讀了很多.
你有沒有經歷到這件事呢.
或者是 你今天有沒有經歷到.
仇敵吃我的肉.
我的仇敵責刑攻擊我這件事呢.
仇敵吃我的肉.
我相信 我們不是自命的意思.
就是人際之間很多時候互相攻擊.
或者是你受到一些不公義的對待.
甚至他們責刑一齊地攻擊你.
一群一群地攻擊你.
你試過這些日子沒有呢.
我相信試過的弟兄姊妹.
你感覺到這些攻擊比病毒更加厲害.
但是 你相不相信.
你有沒有試過在這些日子裡.
你看到神的光 祂的拯救 祂的山寨呢.
大衛這麼有信心.
是不是說就有呢.
我們自己經常想到.
我真的沒有大衛這麼有信心.
我來的時候 仇敵打過來的時候.
我還是腳軟.
病毒打過來 我一樣是腳軟.
我沒有他的信心.
但是原來大衛也是沒有這個信心.
在這個詩篇裡.
大衛跟我們分享到一樣很真實的東西.
原來他是沒有信心的人.
但是有一件事改變了他.
我們看這次的《聖經》.

$^{201}$第四節那裡說.
有一件事我尋求耶和華.
我仍要尋求.
就是一生一世在耶和華的殿中.
沾釀他的榮美.
在他的殿裡面求民.
什麼是有一件事呢.
這裡說得很清楚.
大衛說是住在耶和華的殿裡面.
說真的 其實當時大衛的時候.
都沒有見聖殿.
當然就不是指耶路撒冷的聖殿.
可以說是敬拜的地方.
如果對我們來說.
如果我們說一生一世住在聖殿裡.
是不是說我們一起來住在教會UV這裡呢.
如果是的話 這裡的老闆就很開心了.
不過就租給我們.
我們住不完給我們.
但是聖殿不是指我們真的住在這裡.
什麼叫做住在耶和華的殿裡面呢.
昨天我們在這裡有一個音樂團隊的聚會.
我們統計過這三年裡面.
在教會裡面唱得最多的那首歌.
就是我一生一世住在你聖殿.
這首歌是很受歡迎.
這個詩篇也是很耳熟能詳.
但是什麼叫做住在你聖殿呢.
住在你聖殿.
如果我們看詩篇23篇到30篇.
原來這個是一組的詩篇.
裡面時時提到住在耶和華的聖殿.
意思就是住在耶和華裡面.
意思就是說我們和耶和華有一個親密的關係.
說完就是這一句.
住在耶和華的聖殿.
就是我們和耶和華有一個親密的關係.
我們有沒有住在耶和華的聖殿裡面呢.
這個是很關鍵 很重要的.
我們有沒有和祂保持一個親密的關係.

$^{241}$不單止是這樣.
在這個親密的關係裡面.
我們沾釀祂的榮美.
沾釀祂的榮美又是什麼意思呢.
大衛所指的就是.
我會回顧過去日子裡面.
神多麼信實待我.
待我的先祖.
待我們的民族.
這位神祂說過的話.
還有祂在我身上所行的恩惠慈愛是多麼的大.
這個就是大衛所說的沾釀神的榮美.
弟兄姊妹你們有沒有沾釀神的榮美.
當我們回想起我們今天的信仰的時候.
我們有沒有只看到前面的東西.
不如我們回顧神過去在我們身上所做過什麼.
我們不要忘記.
我們一樣一樣去思想.
我們會看到.
這個就是神的榮美.
當我們住在神的殿裡面.
沾釀祂的榮美.
這件事改變了我們的生命.
以致當我們面對恐懼的時候.
我們會將這個恐懼變為力量.
我們不再害怕.
因為我們體會到耶和華和我們同在.
新冠狀病毒這個疫情發展到現在.
真的可以說情況越來越嚴重.
首先我們最希望的是我們不會得到病.
我們親愛的人不會得到病.
不會受到感染.
所以我們會積極做很多事.
但是原來這個病毒帶來的.
這個疫情帶來的問題.
不單只是身體的問題.
還是有一個整個全世界.
不單只是受到健康的影響.
而且是經濟的影響.
這些影響我們現在還沒有貼身.

$^{281}$在中國大陸就已經開始貼身了.
但是在海外.
我們慢慢都會體會到這個影響.
這個經濟帶來的衝擊.
這個海嘯真的很大很大.
可以說是空前的.
所以跟別人說著說著的時候.
有時候我都忍不住.
給了這個疫情一個綽號.
就是「世紀大鑊」.
因為真的不知道這個鑊.
有多大很大的影響.
慢慢慢慢會來到.
不過說真的.
有哪一個鑊.
是比伊甸園那個鑊更加大呢?.
當亞當夏娃受到試探.
墮入在罪裡面的時候.
那個鑊真的大到.
我們全人類今天還在承受.
但是大家都深深知道.
這個鑊神已經補了.
而且是在及至時候滿足的時候.
耶穌基督為女子所生.
將這個人間大鑊完全補了.
凡相信的人都知道.
都有這個信念.
就是神已經為我們補了鑊.
我們不再需要怕.
所以既然這個鑊都補了.
有什麼人生的大鑊.
有什麼世紀的大鑊我們怕呢?.
我們要住在耶和華的殿中.
沾釀它的榮美.
而且等候它.
向它求敏.
這個是我們今天所需要做的.
這個是會改變我們今天的經據.
由經據變為平安.
以致我們繼續看聖經的時候.

$^{321}$會看到大衛說.
在這個經據裡面.
因為神的緣故.
我能夠遭遇患難的時候.
祂暗暗保守我.
祂在亭子裡面.
將我藏在祂的帳幕隱密處.
將我高舉在盤削上.
這個是完全的勝利.
一個經據的人完全得到保護.
這個是大衛所經歷到的.
完全的勝利.
不單止這樣.
祂知道祂會繼續的勝利.
所以在第六節說.
現在我得以昂首.
高過四面的首敵.
我要在祂的帳幕裡面.
歡言獻祭.
我要嘗試歌頌耀華.
第六節說.
現在我得以昂首.
其實如果你看英文的聖經.
會更加清楚.
它是一個將來式.
祂說現在我深信一件事.
就是我將來會受到.
更加大的災難的時候.
我都會高過四面的首敵.
我會繼續在神的帳幕裡面.
歡言獻祭.
我會歌頌耀華.
即使有未來的災難.
將來會發生在我身上.
我仍然有這個信心.
這個是大衛在經據裡面.
因為在祂能夠沾養耀華.
在神的殿裡面沾養他的榮美.
而帶來的信心.
其實這篇詩篇.

$^{361}$讀到這裡可以結束了.
因為這裡已經是一個總結.
已經是完全的勝利.
但很有趣.
我讀詩篇27篇的時候.
往往都覺得.
為什麼接著又有這段話.
已經是完全的勝利.
為什麼還在說.
我要用聲音呼籲你的時候.
求你垂聽.
求你連續說.
好像現在很不濟.
求你不要向我掩面.
不要丟棄我.
我父母離棄我.
你也不要丟棄我.
為什麼會有這樣的事呢.
為什麼鄭華已經是.
完全的勝利.
突然間又這麼大.
好像很負面.
走回頭來呢.
很簡單.
信心並不是沒有試探.
信心並不是沒有掙扎.
這裡只是詩篇27篇.
大衛說完一個總結之後.
他繼續所說的.
信心是有掙扎的.
等於今天我們的信心也有掙扎.
雖然我們知道.
我們已經得到完全的勝利.
特別在復活節的時候.
我們唱復活的詩歌.
勝利的詩歌.
唱得很興奮.
但是當我們回到日常生活裡面.
我們仍然知道.
我們的信心會有很多掙扎.

$^{401}$但是大衛繼續跟我們分享.
在掙扎的時候.
要呼求.
我要呼求.
我要呼籲.
我要尋求神的面.
我繼續要尋求.
這是大衛跟我們第一樣的分享.
第二就是.
即使有時候我覺得神好像不理會.
默默無聲.
我呼籲了很多次.
我祈求了很多次.
祂也沒有反應.
你繼續要死掉.
好像神不願意向我們掩面.
大衛說不要向我掩面.
我要死在這裡.
你不要丟棄我.
不要離棄我.
在這個信心的掙扎裡面.
這些是很重要的.
還有.
大衛繼續祈禱說.
神你不要將我交給敵人.
求你將正確的道路指示我.
祂不斷地去祈求.
總括來說.
信心是面對很多的掙扎.
我們今天.
我們時時都面對這些掙扎.
各位大英姐妹.
你有沒有想起.
你過去這個星期的掙扎.
或者是過去一個月.
甚至是過去一年裡面的掙扎.
很多時候.
落在大能裡面.
我們會很害怕.
我們會覺得.

$^{441}$神離開了我們.
祂不再聽我們.
但是讓我們再一次看到大衛這裡.
他的經歷說.
我呼求.
我死走.
我不結求.
因此我經歷到神的同在.
祂沒有離開過我.
以致到祂可以說到.
接著.
十三十四絕的說話.
我若不信在活人之地.
得見耶和華的恩惠.
就早已喪膽了.
要等候耶和華.
當壯膽.
堅固你的心.
我再說.
要等候耶和華.
第十三節那裡.
其實原文的意思是.
我很相信.
在活人之地.
我會見到耶和華的恩惠.
不過我們的和合本.
的翻譯.
可圈可點地加多了.
將它譯為.
我不喪膽.
我若不信在活人之地.
得見耶和華的恩惠.
就早已喪膽了.
我覺得和合本這個譯得很全神.
因為實在在恐懼的日子.
在受敵吃我的肉.
攻擊我的日子.
我真的會嚇到瘋了.
嚇到腳軟.
嚇到我腎上腺分泌.

$^{481}$控制了我.
嚇到我要飛車去Costco.
買廁紙.
買口罩.
買米.
當我們很害怕的時候.
如果我們不是見到.
耶和華的同在.
如果我們不是回顧.
祂過去如何帶領我們.
祂的恩典詞.
我們如何不離不棄.
無論我們多軟弱的日子.
無論我們多不配.
祂都是和我們一起的時候.
如果我們不是.
回顧祂的過去.
我們真的會灰心喪膽.
但是讓我們記起.
神過去如何帶領我們.
這個神是如何信實地帶領我們.
甚至祂在十字架上.
如何將我們這個人類的大鑊.
完全補救了.
我們不自喪膽.
無論今天多害怕.
無論疫情是如何.
老實說真的害怕.
最高的數字是美國.
今天有五個字嗎.
真的會繼續上升.
當檢測越多的時候.
數字會越上升.
我們伊利諾州也會繼續上升.
你等著吧.
一定會有的.
但是讓我們不要因為這樣.
害怕得完全軟弱.
我們若不信在活人之地.
得見耶和華的恩惠.

$^{521}$我們是會喪膽的.
但我們已經見過耶和華的恩惠.
讓我們一起堅信這件事.
讓我們死抓住祂.
讓我們不停地信靠祂.
依靠祂.
最後大家有沒有看到.
如果你看第十四節.
你覺得大衛這個時候.
說到這一個詩篇最後的時候.
他的心情是怎樣的.
其實是挺振奮的.
因為不是振奮的人.
他不會說出這些話.
他說當壯膽.
堅固你的心.
我再說要等候耶和華.
只有一個經歷過信心的掙扎.
他仍然看到神沒有離開過他.
而且是不斷地去保守他的人.
他會說這些話.
你要等候.
當壯膽堅固你的心.
其實就是當日.
摩西和約斯亞所說的話.
你當剛強壯膽.
完全是那些字.
你當剛強壯膽.
等候耶和華.
無論你面前的大敵是多麼大.
無論你面前是多麼不可知.
你進入迦南地.
是多麼像如臨大敵一樣.
但是你仍然剛強壯膽.
就好像我們今天.
無論我們面前.
我們根本不知道.
這個疫情會繼續發生到什麼時候.
我們仍然剛強壯膽.
想起耶和華.

$^{561}$而且是要等候他.
知道他是行事.
有他的時間.
不是我們的時間.
雖然黑暗的日子.
我們是會見到光.
在差不多結束這本詩篇27篇的時候.
我很想和大家看一個很重要的東西.
就是詩篇27篇.
是很短短的14節的聖經.
但是這14節的聖經.
是有一個很特別的結構在.
就是一到三節.
在恐懼裡面.
我們深信神的保護.
就正如祂是我們的光.
我們的拯救.
我們大衛是深信神的保護.
到13,14節的時候.
也是重複這件事.
在恐懼裡面.
他不信在活人之地得見神.
因為他灰心喪膽.
但是他現在相信.
所以他說要等候耶和華.
剛強壯膽.
他是在恐懼裡面.
深信神的保護.
而中間4至12節.
原來就是有一件事改變了我.
等於我們讀聖經的時候.
我們可以從結構上.
更加看到這段聖經.
帶給我們神的心意.
和這段話的力量.
這種結構我們可以說是首尾呼應.
也可以說是一個三文治的結構.
我們在主義學上也說了很多.
往往首尾講的那件事是重要的.
不過中間那個更加重要.

$^{601}$那個可以說是一個訊息的精華所在.
詩篇27篇是講恐懼裡面.
我深信神的保佑.
但是在中間這一段.
是告訴你.
你怎樣能夠得到這個保佑呢?.
有一件事改變了我.
住在耶和華的殿裡面.
沾釀它的榮美.
即使在信心的掙扎裡面.
你仍然去求告,呼求.
去圍著它繼續地求.
這就是詩篇27篇.
願意你回到家裡.
願意你今晚.
願意你接著面對疫情的日子.
你時時都記住這個詩篇.
神給你的應許.
但是你要做一件事.
就是你願意沾釀它的榮美.
在它的聖殿裡面求民.
沒錯.
在面對困難,面對疫情.
我們是要很小心,謹慎.
這是真的.
我們不能夠經常說.
我們不怕不怕.
有耶和華同在.
我們就什麼都不做.
不是不是.
不是這個意思.
我們仍然是要謹慎.
我們仍然是要小心.
盡我們的力.
做應該的防護工作.
但是我們要記住的是.
當事情未那麼嚴重的時候.
我們的水只是淹了腳.
就好像掉了進大海一樣.
這樣對我們其實是沒什麼用的.

$^{641}$因為這樣會害怕.
上加害怕.
帶來很多恐慌.
但是相反.
我們去依靠耶和華.
我們不要好像一個.
沒有神的人一樣.
當我們死亡的時候.
其實我們就是說.
神啊,你搞不定.
我自己來吧.
如果是這樣的時候.
你永遠經歷不到神的大能.
讓我們再一次記起.
有一件事.
一生一世住在耶和華的殿裡面.
沾釀它的榮美.
在它的殿裡面求民.
我再回到第一句開始的時候.
說到Max Lucado的這句話.
恐懼是會探訪每一個人.
但是你可以讓它做一個訪客.
不要成為你的住客.
不要讓恐懼住在這裡.
我再補充一句.
在恐懼的日子裡面.
我們特別記住.
請你讓你成為神的住客.
不是訪客.
我們住在神裡面.
我們住在祂那裡.
我們不試探祂一下就走了.
讓我們因為這樣.
而經歷到神的同在.
神的能力.
神在恐懼的日子.
給我們的力量和安慰.
我們一起祈禱.
\newpage



\section{傳道書}
\label{sec:EN7WJcnMP5Y}
\textbf{2020/3/22 崇拜直播 vlog \#1 疫境前行 (傳道書 3) - 陳崇基牧師}
\newline
\newline
連結: \href{https://youtube.com/watch?v=EN7WJcnMP5Y}{\texttt{https://youtube.com/watch?v=EN7WJcnMP5Y}} ~~~~ 語音日期: 2020-03-23
\newline
\newline
\hyperref[sec:d0PXI2JCtZU]{\small{< < < PREV SERMON < < <}}
~
\hyperref[sec:index_chronic]{\small{[返順時目]}}
~
\hyperref[sec:index_scriptual]{\small{[返順卷目]}}
~
\hyperref[sec:o3fsHQqXKjc]{\small{> > > NEXT SERMON > > >}}
\newline
\newline
$^{1}$各位好.
這是芝加哥時間星期六的下午四點.
我們正在等候州政府的家居隔離的命令開始.
命令是星期六下午今天五點.
所以還有一個小時.
所謂 stay at home(留在家中).
家居隔離的命令就會開始了.
因為這個疫症的緣故.
所以我們有很多芝加哥的市民.
有很多人都回家了.
所以你看到街道上也是沒有什麼人的.
但是我們正在準備.
每一個星期日的崇拜是怎麼做的呢?.
我們有很多的弟兄姊妹幫忙.
以及在我家的 basement(地下室).
在土庫那裡進行崇拜的直播.
我帶你進來.
這裡就是我家.
那麼.
我們就在這裡.
我的 basement(地下室).
來做崇拜的直播.
你看到我們已經 set up(準備)好所有東西了.
因為現在隔離令的緣故.
我們是不能夠去教會的地方做直播.
所以我們將所有東西搬到家裡這裡.
但是我們照舊崇拜.
中文英文的唐崇拜.
我們都會正常地進行的.
只不過你在手機或者家裡的電視.
你的平板電腦那裡看.
那麼我希望你能夠和我們一起來去敬拜主.
無論疫情是怎麼樣也好.
我們認為上帝仍然是值得我們事奉.
值得我們敬拜去讚美的.
現在是星期六的晚上八點鐘.
芝加哥的時間.
我們出去了拿一些東西.
因為留了一些東西在教會.
待會回家.

$^{41}$現在在 downtown(市中心)兜一兜.
明顯地車是少了的.
在芝加哥很少會有這樣的情況.
沒有塞車.
沒有任何的交通.
很明顯很多人都趕著回家.
或者留在家裡.
讓我想起很多時候去以色列.
那個 Shabbat(安息日).
安息日的星期五晚上.
Shabbat 一開始.
人們全部都回家了.
街上的車都很少.
從來沒有見過美國.
芝加哥會有這麼 Shabbat(安息日)的情況.
我們現在來到 Chinatown(中國城).
這裡平時星期六晚上是很熱鬧的.
因為餐館都開了.
車都停不下來.
近來特別是今天有一個家居隔離令.
就更加少人了.
整個停車場.
還有一些貨車停在這裡.
都沒有其他車了.
(車停下來).
再一次歡迎各位.
來到我們城會教會主日崇拜的直播.
這次是第一次在我自己家的 basement(地下室).
做一個家居版 Home Edition(家居版).
我不知道你過去一段時間.
無論在芝加哥或全球哪一個地方.
我相信都是一些很失控的時候.
令到我們情緒或生命生活.
都很多挫折和打擊.
我們芝加哥這個地方.
昨天星期六下午五點.
我們就開始了這個所謂的 Stay at Home Order(家居自我隔離).
就是在家居自我隔離的命令.
州長在伊利諾州就說.
每一個人都要留在自己家裡.

$^{81}$嘗試去減慢疫症的蔓延.
這個其實遲了做.
我們應該一早要去控制疫症的蔓延.
因為在中國武漢開始爆發武漢肺炎的時候.
其實事情已經非常之大.
現在你知道國際不同的地方.
都已經大爆發.
在美國都開始很快的蔓延.
昨天在一個朋友的Facebook上.
他就發佈了一個意大利醫院的情況.
大家可能看到.
你見到那個情況真的好像那個headline(頭銜).
有緣末日一樣.
真的很需要美國.
我現在身處美國這個地方的社會的人.
要去明白這個新冠病毒.
這個病毒不是普通的流行性感冒.
那個殺傷力和會令到整個醫療系統癱瘓.
那種的染例.
在美國我估計還沒有感染過.
不像香港等等的地方.
可能十多年前有過SARS.
他們的人民都非常警覺性很高.
在美國好像很輕鬆.
很輕鬆的.
這是一個很大的教訓.
我有一個朋友是基督徒.
在疫症在中國開始爆發的時候.
我已經跟他說一定要小心.
一定要take it seriously.
要很嚴肅的去看待.
他總是跟我說.
我們很有信心沒事的.
我們是信神.
只是嚴重一點的感冒而已.
現在肯定不是.
那時候一早說的時候.
我心裡面都覺得很不舒服.
為什麼你經常將信心和智慧.
我們人的責任.

$^{121}$你都將它對立.
我們信神of course(當然).
我們信神.
我們有盼望.
但不等於我們不需要去.
用我們的智慧和社會的責任.
特別在這個疫情當中.
真的有需要的人.
未必一定是我們年輕力壯.
我們能夠走得走得.
我們可能感染到.
但我們未必會有事.
未必會死.
但有很多被我們脆弱的人.
老人家等等.
如果受到這個疫症的影響.
他們的生命會有很大的危機.
希望這個社會.
希望全地球所有的人.
他們都能夠明白到.
這個疫情的嚴重性.
我們真的要take it seriously(認真地看待).
我們的faith(信仰).
我們的wisdom(智慧).
信仰和我們的智慧.
一定需要有一個很好的平衡.
我有一個朋友是白人.
他在美國deep south(深南).
美國很南邊的州.
在那裡募會.
他在Facebook.
昨天星期六.
他發了一個訊息.
他說在deep south(深南)有一間教會.
他們堅持要聚會.
結果他們在上星期的聚會當中.
有四十個人受感染.
他說得很對.
這個時候雖然對教會來說.
是一個阻滯.

$^{161}$我們要上網.
在網上直播崇拜.
每個人都不喜歡.
每個人都覺得這不是很理想的事.
實體崇拜始終是最理想的.
但是這樣的疫情的緣故.
我們必須要轉去上網.
否則我們很容易連累到其他人.
所以我們教會都需要做這方面的措施.
我們不會在教會直播.
在家裡直播.
無論如何都是一種生命失控的感覺.
上星期我的女兒已經去大學.
她有一個星期的春假.
學校宣佈了停止實體的教學.
全部人都會上網.
她都會上網.
所以她放完春假就不用回去.
她就回家.
在這裡一直到學期尾.
都是用網上的方法上課.
這也是一個好消息.
但是對我們來說.
突然間有的時間表.
或者安排.
全部打斷了.
我們全部都要改變.
要作一些調節.
相信你也是一樣.
教會當然是.
全球的教會.
你的教會.
我的教會.
上星期我們還在教會直播.
今個星期我們就要把東西搬回來.
在我的地下室.
在家裡直播.
但無論如何我們相信.
神有他自己的時間.
在疫情當中.

$^{201}$我特別想起一段聖經.
在過去一個星期.
經常在我腦海裡出現.
就是傳道書第三章.
傳道書第三章說什麼呢?.
我們先讀這段經文.
傳道書第三章第一節.
凡事都有定期.
天下萬物都有定時.
生有時死有時.
栽種有時.
拔出所栽種的也有時.
殺戮有時.
醫治有時.
擇偉有時.
建造有時.
哭有時.
笑有時.
哀動有時.
跳舞有時.
拋擲石頭有時.
追罪石頭有時.
懷抱有時.
不懷抱有時.
尋找有時.
失落有時.
保守有時.
捨棄有時.
撕裂有時.
逢補有時.
靜默有時.
言語有時.
喜愛有時.
恨惡有時.
爭戰有時.
這樣看來.
做事的人在他勞碌上.
有什麼益處呢?.
我見上帝叫世人勞苦.
使他們在其中所經歷.

$^{241}$上帝做萬物.
割愛其事.
成為美好.
有將永生安置在世人心裡.
然而上帝從始至終的作為.
人不能參透.
這一段聖經.
經常提到時間.
什麼都有時.
時間不單單是時間列表上的一點.
某年某月某日.
某時某分.
時間也有處境的意思.
我們受制於很多.
我們不能控制的情況.
一些 occasion(情況).
一些警況當中.
所以這裡提到.
我們每一個人在日光之下.
都是受制於時間空間.
我們愛在的處境.
為什麼我經常想起這段聖經呢?.
因為傳道書很貼身.
講的東西.
特別是現在這種情況當中.
我們人很脆弱.
傳道書你記得.
它經常說.
虛空的虛空.
虛空其實就是我們一口氣.
我們戴口罩的時候.
我們感受到.
我們真的有一口氣.
虛空的意思就是這樣.
我們虛空的虛空.
我們日光之下.
人生真的很短小.
而且很脆弱.
他就思想到人生到底是怎麼回事.
剛才那段經文第三章.

$^{281}$傳道者講出他怎麼看人生.
他列出了14對人生經歷的對比.
在第二至第八節.
每一句都講到我們人生經歷的不同方面.
第二節他講生有時死有時.
沒有人可以控制到你什麼時候生什麼時候死.
或者在什麼情況當中.
什麼背景裡面出生和死亡.
也講到栽種有時.
拔出栽種的東西也有時.
這裡講到草.
花草.
菜梳.
連菜梳的生命.
什麼時候栽種他都開始.
到他被拔出的時候.
都不到他去掌管.
殺戮有時.
醫治也有時.
我們生老病死不到我們去控制.
拆毀有時.
建造也有時.
這裡講到我們的工作.
我們的建築.
特別提到建築的人.
因為是最明顯的.
什麼時候拆毀什麼時候建造.
有時候真的不到我們去掌管.
在這個逆境當中.
我已經聽到一些朋友.
弟兄姊妹.
甚至教會.
經濟上受到很大打擊.
我們真的不知道將來會怎樣.
也有些已經失業.
或者一些特別餐飲業.
服務業的行業的人.
他們都面對很多困境.
這個又是不到他們去控制.
哭有時笑有時哀動.

$^{321}$跳舞也有時.
這裡就講到我們人生裡面的情緒.
什麼時候哭什麼時候笑.
什麼時候可以哭.
什麼時候想笑但是笑不到.
在疫情裡面我們經歷很多.
第五節說到懷抱有時.
不懷抱也有時.
這裡講到人際之間的分離.
那種聚合.
什麼時候懷抱.
什麼時候能聚在一起.
可以吃飯可以見面.
有它的處境限制.
不懷抱也有時.
特別我們現在在疫情當中.
我們需要去社交距離.
社交要保持六尺的距離.
要在家居裡面隔離.
這個都不到我們去作主.
尋找失落有時.
保守捨棄有時.
有時不見面.
第七節說到撕裂有時.
馮保有時.
希伯來文撕裂的時候.
就是他們哀動.
譬如去喪禮等等.
他們撕爛衣服.
代表他們哀悼或懺悔.
撕裂有時完結之後.
他們哀悼完結後.
他們要補衣服.
也有時.
靜默言語也有時.
特別在最近這麼多混亂不安的當中.
有很多訊息.
有很多說話.
不一定是你說.
有時你也避不了.

$^{361}$有很多人排山倒海的.
說出一些訊息.
這些言語有時.
又不到我們去控制.
喜愛恨戶.
戰爭和好也有時.
甚至說到國際的關係.
戰爭和和好.
很多時候.
都不到我們控制的處境.
這裡說到.
神將人經歷的萬無.
他都去掌管著.
但很多時候.
不到我們去參透.
第十一第十二節.
他說.
神造萬物.
國安其事成為美好.
將永生安置在世人的心裡.
然而.
神從始至終.
他所做的事.
他的作為.
人不能夠參透.
偏偏就是人.
我們這麼有限的人.
在日光之下.
我們不能夠參透.
我們不能夠控制.
不能夠去知曉.
你問疫情何時會過去.
沒有人會知道.
甚至專家都不能夠.
很肯定地估計得到.
這也是我們控制的.
上帝就要我們.
脆弱的人.
我們要明白到.
我們就是虛空的虛空.

$^{401}$真的不到我們去作主.
而我們需要去仰望他.
所以他這裡說到.
神造萬物.
國安其事.
安他的時間.
成為美好.
在這裡.
希伯來這個字.
很有趣.
美好是Yahoo.
Yahoo是什麼呢.
是美麗的意思.
在以色列有一個小的海港.
在Tel Aviv下面.
靠著Tel Aviv.
叫做約柏港.
一個很小的漁港.
舊約裡面也經常看過.
新約裡面也看過.
一個很小小.
風景很美麗的海港.
很宜人的風景.
他就用了這個字Yahoo.
他說上帝去編排.
萬物的時候.
按照他的時間.
成為美好.
我知道現在是一個很困難的時間.
對很多人來說.
有很多的不確定性.
傳道書其實就是要告訴我們.
我們不可以在日光之下.
Under the sun.
用世人的眼光去審視萬物.
我們需要Above the sun.
在日光之上.
用上帝自己永恆的角度去看萬事.
這樣的時候我們知道.
無論發生什麼事.

$^{441}$上帝總會有他自己.
Work out(做到)到的美好.
不一定在我們人類有限的視角裡面能夠看到.
也未必在我們今生今世能夠看到那個美好.
但我們知道是在各安其事.
安神自己的時間.
所以我們可以去信靠神.
可以去信任得過他.
無論多麼動盪多麼不安.
一切在上帝自己的時間.
不是按照我的時間.
會成為美好.
在過去一個星期.
我經常想起一首詩歌.
一首比較舊的詩歌.
這首詩歌很能夠導出我們很多人的心聲.
這首歌叫做「神你在掌管」.
這首歌可以看到我們很多時候.
在很多不明朗的時候.
我們不能夠去掌管一切.
但我們知道神在他的手中.
一切都是有他的時間.
我們希望這首歌能夠成為一個回應.
好嗎?.
我望著前路.
是迷亂張望一片.
幾多挫折打擊.
未明白做甚麼指引.
我踏著細沙.
在海邊走過.
晚風中心裡聽你聲.
聽灑灑海風.
在寧靜細雨低霜.
眺望滔滔海風.
在無言默默宣告.
諸浩大能力.
海天皆響應.
每一刻變幻.
亦在你的指裡.
神你在掌管.

$^{481}$時間萬有空間.
在你的手裡一切定有時.
神我的一生.
神你定有預備.
莫計得失只想你來到人.
神你在掌管.
時間萬有空間.
在你的手裡一切定有時.
神我的一生.
神你定有預備.
莫計得失只想你來到人.
希望這首歌能成為大家的鼓勵.
幫助大家親近神.
很多時候是一些詩歌.
不是人說的敘述.
能夠帶領我們親近神.
所以在聖經裡.
很多時候有很多詩詞,歌賦,詩篇.
這些詩歌.
能夠讓我們去到神的面前.
在我們最後結束的時候.
希望大家不要忘記.
我們不只是自己安全就算.
我們要去顧及其他人.
特別是社區裡面的.
特別在疫情當中.
有很多人的生命受到很大的打擊.
他們的心裡面都很煩亂.
昨天和太太下午去了超市.
去超級市場.
在那裡買東西.
走的時候就有一點風波發生.
剛才我們去超級市場買東西的時候.
有一個事件發生.
在收銀處.
有一個黑人女人和一個白人女人.
有一點口角.
我不清楚發生什麼事.
但似乎那個黑人女人說.
那個白人女人拿了她的東西.

$^{521}$於是兩個人吵架.
吵著吵著的時候.
那個黑人女人想動手.
那些售貨員嘗試去調解.
那個黑人女人居然拿一個胡椒噴霧.
噴在那個白人女人身上.
噴到我們去買東西的物件上.
就有一個小的事件.
小的風波出現.
職員立刻報警.
警方就到了調停.
所以在這樣的情況下.
人的情緒都很波動.
所以很容易去吵架.
所以我們很需要去找個地方.
找個方法去調和不安的情緒.
我相信不只是社區裡.
可能你家庭也是這樣.
可能你自己的心境也是很混亂.
在這個時候.
其實真的需要我們去神面前求神給我們那種屬靈的平安.
讓我們不會那麼嘈.
讓我們有一種屬天的平安.
能夠去面對這個困難的時候.
所以我們需要去社區來禱告.
還有在行動上其實你可以做到很多事.
我們可以去探望一些有需要的人.
打電話給他們.
看看他們有什麼特別的實際上的需要.
我們可以幫到他們.
有很多比較脆弱的年老的人.
甚至一些醫護.
在美國的醫院很多的醫護人員都是不夠口罩等等.
如果有一些的眾籌.
來幫到他們.
請你伸出援手.
我們不要只是囤積我們的東西.
讓我們自己安全就算了.
同時也會有不同的情況去守望.
我們希望能夠繼續在網上去事奉神去崇拜.

$^{561}$雖然實體崇拜未必做到.
也不知道什麼時候可以做到.
但我們在網上依然可以去到親近住.
在世界各地不同的地方我們能夠去看崇拜直播.
所以結束的時候再一次多謝各位.
在不同的地方去收看我們.
參加我們網上的崇拜直播.
希望上帝能夠祝福你們.
我們結束祈禱.
上主我們求你在這個時候給我們屬天的平安.
讓我們知道在主耶穌基督裡面.
我們有的倚求.
主祈求你給我們應許.
你對我們的信實.
讓我們在全球各地不同的地方.
我們一起來去敬拜弟兄姊妹.
我們知道真正的平安是從耶穌基督而來.
我們在基督裡面已經有了.
主讓這個平安能夠成為我們.
落實在我們生命當中.
並且在我們的社區.
我們的生活圈子裡面.
能夠去彰顯出來.
主祈求你在每一個教會.
每一個弟兄姊妹.
我們散落在不同的地方.
世界各地.
我們能夠將你的平安.
帶到這個不安動盪的世代當中.
使人能夠見到你自己.
是那位掌管萬物.
國安其事成為美好的神.
我們這樣祈禱奉祖耶穌基督的名字.
Amen.
再一次多謝你來收看我們崇拜直播.
我們也希望你能夠支持我們.
成會教會成會網台的事工.
我們除了做網上崇拜直播之外.
我們也做了很多熟人的訊息節目.
希望能夠在不同的地方成為更多人的祝福.

$^{601}$所以希望你繼續努力.
我們下個星期能夠在網上的直播.
崇拜裡面.
芝加哥時間星期日十五早上.
我們能夠見到你.
\newpage



\section{詩篇}
\label{sec:o3fsHQqXKjc}
\textbf{2020/3/29 崇拜直播 vlog \#2 疫境前行 (詩篇 139) - 陳崇基牧師}
\newline
\newline
連結: \href{https://youtube.com/watch?v=o3fsHQqXKjc}{\texttt{https://youtube.com/watch?v=o3fsHQqXKjc}} ~~~~ 語音日期: 2020-03-30
\newline
\newline
\hyperref[sec:EN7WJcnMP5Y]{\small{< < < PREV SERMON < < <}}
~
\hyperref[sec:index_chronic]{\small{[返順時目]}}
~
\hyperref[sec:index_scriptual]{\small{[返順卷目]}}
~
\hyperref[sec:G9uX1Fwy_5w]{\small{> > > NEXT SERMON > > >}}
\newline
\newline
$^{1}$我們都會覺得是一篇讚美詩.
因為我們都很熟悉.
我們都很熟悉.
我們都會覺得.
這篇詩是一篇讚美詩.
我們都會覺得.
這篇詩是一篇讚美詩.
我們都會覺得.
這篇詩是一篇讚美詩.
我們都會覺得.
這篇詩是一篇讚美詩.
我們都會覺得.
這篇詩是一篇讚美詩.
我們都會覺得.
這篇詩是一篇讚美詩.
我們都會覺得.
這篇詩是一篇讚美詩.
我們都會覺得.
這篇詩是一篇讚美詩.
因為裡面的詩詞歌賦.
讓我們好像很感激上帝.
很感恩等等.
大家看看這篇詩篇.
很多時候的用語.
都不是一篇典型的讚美詩會用的.
他向耶和華呼求.
他去求主.
去為他伸冤.
有很多的敵人來壓迫他.
這些都不是普通讚美詩會寫的東西.
反而在詩篇裡面.
很多時候會提到.
叫上帝伸冤.
這些都是哀歌經常會提到的.
哀歌是什麼呢?.
就是詩人面對很多重重的壓迫.
重重的不安.
他會去求上帝為他伸冤.
他面對敵人的時候.
面對欺壓的時候.

$^{41}$他去求主幫助他.
他的心來到神的面前.
求上帝給他力量.
給他能夠有那種智慧.
去面對重重的壓迫.
今天在疫情當中.
這篇詩篇就是對我們說話的.
不如我們看看這篇詩篇到底說什麼.
詩篇139篇.
耶和華你監察了我.
你認識我.
我坐下我起來.
你都知道.
你在遠處就明白我的意念.
我行路我躺臥.
你都細察.
我的一切行為你都熟悉.
耶和華我的舌頭還沒有發炎.
你已經完全知道了.
你在我前後環繞著我.
你的手按在我身上.
這樣的知識奇妙.
是我不能理解的.
高超是我不能達到的.
這首頭六節告訴我們.
這篇詩篇經常出現一個字.
叫做知道熟悉或者深知.
這裡說到神的全知.
說到上帝知道我們一切的行動.
知道我們心裡面的思想.
幾天前在我的飯廳.
我突然說了一句話.
我說Oh my god.
然後突然電話的Google Assistant.
它就說How can I help you.
我就覺得很奇怪.
原來它一直在聽我說話.
我說Oh my god的時候.
好像在求救一樣.
Google Assistant馬上來問.

$^{81}$到底我有什麼需要.
我不是很喜歡這樣的東西.
但是你會發現.
今天很多時候電話不斷聽我們說話.
或者周圍很多的監察.
監察著我們.
但是無論如何.
我說的話它都可以聽到.
我的行動鏡頭可能可以看得到.
可以偵查得到.
但是一樣東西.
這些鏡頭是不能夠知道的.
就是我心裡面所想的東西.
我心裡面的思想.
我心裡面每一天煩擾的東西.
每一天我想起的東西.
這些是不能夠見到.
不能夠說不出來.
但是這篇詩篇就說.
神是很清楚.
是很熟知的.
知道我心裡面最深深處.
在想什麼.
這些全部都不能夠逃避.
這位全知的上帝對我的熟知.
詩篇139篇.
這個詩人就在生活裡面.
很多壓迫.
很多不安的當中.
他就想起神學的信念.
他說原來上帝什麼都知道.
什麼都看察.
我經過什麼.
或者現在體會的是什麼.
神其實是很清楚.
並且很緊張想知道.
在現在隔離的日子.
有時候真的有辣有不辣.
沒錯我們多了時間和家人一起.
但是同樣我們也多了很多機會.

$^{121}$去到有爭執.
因為在隔離的日子.
大家困在家裡.
生活的空間又少了很多.
當然我們在美國.
芝加哥這個地方.
我們怎樣也好.
我們的生活空間.
其實也比很多其他城市.
其他地方更加廣闊一點.
不過大家困在一起的時候.
難免有很多爭執.
意見不同.
所以有辣有不辣.
但是在這些情況當中.
詩人說原來上帝很熟悉.
還有上帝很想知道.
並且不斷地來觀察.
我們坐在這裡.
或者我們走路.
我們睡覺躺在這裡.
在我們家居的日子當中.
原來上帝一直和我們同在.
詩人不單單說.
在他生活的壓迫和不安裡面.
認定神是全知.
是很清楚的.
他在七至十二節怎麼說呢.
他說我往哪裡去逃避你的靈.
我往哪裡躲避你的念.
我若升到天上你在哪裡.
我若在陰間下榻你也在哪裡.
我若展開清晨的翅膀.
飛到海極居住.
就是在哪裡你的手必引導我.
你的右手也必扶持我.
我若說黑暗必定遮蔽我.
我周圍的亮光必成為黑夜.
黑暗也不能遮蔽我.
使你不見.

$^{161}$黑夜卻如白晝發亮.
黑暗和光明在你看到時一樣.
他說就算他展開清晨的翅膀.
這個神.
全知的神都會不斷來和他同在.
他說直到海極之處.
去到海邊.
這個上帝都會來看見他所經歷的.
展開清晨的翅膀對於一個古代的以色列人來說.
是東邊的意思.
就是日出之處.
到海極的地方.
對於一個住在巴勒斯坦地區的人.
海極的地方就是西邊.
就是由東到西.
無論他在哪一個地方.
整個地面上.
神都會來監察著.
認識著他.
知道他一切所行所經歷.
在古代的神.
很多時候有一個是管上晝.
一個是管夜間.
日出到日落.
是一個神靈來管理.
另外夜間當中就是另外一個神靈.
這裡告訴我們.
耶和華24小時.
都是來看著我們.
監察著我們.
每一天留在家裡.
活動的空間越來越小的時候.
心裡面的心情可能都會有很多很煩擾的事.
特別當你覺得有些事你想做但做不到.
有些人你想見但見不到的時候.
我不知道你的心思意念會有多混亂.
我相信這個詩人都是同樣.
有一種這樣的感覺.
被遺棄或者被壓迫.
或者被一種很無形的壓力.

$^{201}$使他喘不過氣.
但是他去到仰望上帝.
他說這個上帝原來是全知.
是很明白的.
這個上帝沒有背面來利器他.
當他發現有很多事情.
令他失控的時候.
原來這個上帝一直來清楚知道他所受的是什麼.
第十三節.
他說我的肺腑是你所做的.
我在模糊之中你已經覆蔽我.
然後他說我要稱謝你.
因為我受造.
奇妙可為.
你的作為奇妙.
是我深深知道的.
我就暗中受造.
在地的深處被聯絡.
那時我的形體並不向你隱藏.
我未成形的體質.
你的眼早已經看見.
你所定的日子.
我尚未到一日你都寫在你的冊上.
上帝啊你的意念是何等寶貴.
奇數何等眾多.
我若數點比海沙更多.
我睡醒的時候.
仍和你同在.
這裡說到我們的肺腑.
我們的身體.
所有的我們的存有.
神都已經知道.
這些都是一些很典型的近東的詩歌會說到的.
有什麼人會說自己是神去保守著他.
他未出生之前已經得到上帝的福備.
就是那些亞述王啊.
埃及王啊.
那些在當時古代的社會裡面.
他們是有權有勢的人.
他們的詩歌經常說.

$^{241}$讚嘆那位神去到不斷的保守著他.
就算他未出生之前.
神已經和他同在.
讓我們想起在疫情當中.
很多時候有些有權有勢的人.
他們就有很多福利.
他們可以首先得到醫療的照料.
然後有時候他們會有更好的照顧.
你想一下現在這個疫情當中.
有些有權有勢的人.
他們有特別好的醫療服務.
平民百姓就沒有了.
平民百姓很多時候我們甚至想去測試.
想診斷一下都未必有.
近這兩天我們 McKinney Park 有一個家庭.
其中一個成員確診了.
然後有很多病徵出現.
他就帶著整個家庭去醫院.
想說可不可以讓我們做一點確診.
誰知道醫院就關掉了他們.
就說不好意思.
因為你暫時先留在家裡.
我們做不到確診.
當然了因為測試套件都不夠.
使得這個家庭很大掙扎.
很無助.
這篇詩篇就將這件事反過來.
他就說不只是有權有勢的軍王官員才能夠有神來保守.
他說原來全部的人.
我們平民百姓每一個人.
我們的體質.
我們的肺腑.
我們的受造.
都是奇妙可為.
都是有上帝與我們同在.
這不是給我們很大的安慰嗎?.
這位上帝創造了我們的神.
這位醫治了我們的主.
帶領了我們心靈的主.
原來每一個細節都在祂的保守下.

$^{281}$第十九節.
祂說:神啊,你必要殺六惡人.
所以你們好流人血的,你開我去吧.
因為他們說惡言頂撞你.
你的仇敵也罔稱你的名.
耶和華說:恨惡你的,我豈不恨惡他們嗎?.
攻擊你的,我豈不憎厭他們嗎?.
我切切的恨惡他們,以他們為仇敵.
這幾段話不是我們平時崇拜.
或者平時茶經的時候經常會看的.
但是正如我所說.
這篇詩篇是一篇愛國.
上帝在很多不安,壓迫下.
有很多仇敵.
去頂撞上帝,也頂撞他.
他就去求上帝伸冤.
在逆境裡面.
我們很容易就會想起.
什麼是我們的仇敵.
還有一些什麼是面對著的欺壓.
一些不滿.
有很多我們覺得是仇敵的東西.
病毒是無處不在.
有很多時候.
各難當前.
很多人就親火打劫.
或者說一些風涼話等等.
正如詩篇這個詩人.
作為一個義人.
他受到很多無情的壓迫.
無情的待遇的時候.
他去上帝的面前伸訴.
在他生命裡面覺得很失控的時候.
他就去仰望上帝.
求上帝伸張正義.
在逆境當中我們很多時候會很生氣.
我們在很多疫情裡面會覺得.
有沒有搞錯.
國家居然這樣做.
或者在上掌權者說一些不負責任的話.

$^{321}$這些都會令我們心煩意亂.
就像這個詩人這樣說.
他面對著很多亂說話.
頂撞上帝.
不說真理的人.
令他非常心煩.
他求上帝來他身怨.
他也深信上帝不單止是全知.
也是一個公義的審判處.
所以他可以以信心來仰望他.
縱使在生活裡面他很多事情不能控制.
他知道神會伸張公義.
最後這個詩歌這個祈禱.
這個伸訴是這麼結束的.
神啊求你監察我.
知道我的心思.
思念我知道我的意念.
看在我裡面有什麼惡行沒有.
引導我走永生的道路.
這篇詩篇最後是一個詩人伸訴的時候.
他去到仰望上帝.
他知道上帝是會伸張公義的.
並且他願意立志.
去告訴上帝.
我一定會跟他.
跟到底.
縱使現在有多少不安的遭遇.
縱使他面對多少的仇敵也好.
他相信這位上帝會繼續去引導他.
也都決意來跟隨這位上帝.
在這個逆境當中.
我們的生活裡面有很多失控.
有很多本身有的計劃.
不能夠盡前盡尾來殺人出來.
或者今天不知道明天的事.
但是這個禱告這個立志.
也是應該我們的立志.
就像這個詩人所說.
無論生活有多不穩定.
有多少未知之數也好.

$^{361}$我們知道這位上帝會繼續引導我們.
但是我們需要立志.
我們需要告訴這位主.
你是全知的.
我們要在你面前.
來求你引導.
在我們心思意念當中我們去跟隨上帝.
在我們的.
在我們心思意念當中.
我們也要討上帝的喜悅.
求上帝來引導我們走前面的道路.
在詩集詩篇.
150首的詩歌當中.
有三分之一都是哀歌.
這些哀歌都是一些詩詞.
是聖經的作者詩人.
他們遇見一些苦難.
遇見一些逆勢當中.
他們向上帝申訴.
向上帝的呼求.
很多時候這些Lament.
這些哀歌就能夠帶領我們來到上帝的面前.
很多時候我們.
每件事都很順利.
每件事都很平穩的時候.
我們未必聽到上帝跟我們說話.
很多時候只是一些逆勢裡面.
逆境當中.
上帝去到向我們呼喊提醒我們.
人的脆弱和上帝自己的榮幸.
有一個作家.
已故的一個基督教學者.
叫做C.S. Lewis.
他曾經說過一段說話很有意思.
他說在我們平靜安穩的時候.
在我們快樂的時候.
上帝向我們在耳邊.
施施細語.
去到告訴我們.
我在這裡.

$^{401}$但是在逆勢裡面.
在我們痛苦當中.
上帝向我們呼喊.
去到告訴我們提醒我們.
他說在這個已經對他耳聾的世界當中.
神藉著苦難.
好像一個大聲公.
來告訴我們.
神自己的永恆.
神的真實.
神真的同在.
我們有苦難和軟弱的人.
我們才會體會到原來上帝.
真的需要去依靠他.
真的需要去親近這位創造的主.
剛才提到詩篇139篇的時候.
我就一直想起一首詩歌.
叫做《展開清晨的翅膀》.
這首國語的詩歌.
都是出自詩篇139篇.
我希望這首詩歌能夠帶領我們.
來到神的面前.
我們繼續去敬拜他.
繼續去依靠他.
《展開清晨的翅膀》.
初夜過後 你已經檢查了我.
我坐下 我起來 你都已曉得.
我醒路 我躺臥 你都稀缠.
你也深知 我一切所惜.
初夜過後 你已經檢查了我.
我坐下 我起來 你都已曉得.
我醒路 我躺臥 你都稀缠.
你也深知 我一切所惜.
我舌頭上的瓜 你沒有一句不知道.
你在我前後環繞著我 暗守在我身上.
這樣的氣貌 是我不能擦透.
你的自告 你的珍貴 是我永遠不能說盡.
我可以往哪裡去躲避你的意義.
我可以往哪裡去逃避你的念.
我若展開清晨的翅膀飛到地底.

$^{441}$就在那裡 你的雙手也披不到我.
初夜過後 你已經檢查了我.
我坐下 我起來 你都已曉得.
我醒路 我躺臥 你都稀缠.
你也深知 我一切所惜.
好 一齊低頭祈禱.
上帝我們多謝你.
縱使我們在被隔離的當中.
我們在逆境的旅途.
我們深知你是全知道一切.
你是那位全知的神.
你對我們的意念.
你對我們的想念.
是多如海沙.
是不能夠去到數算的.
主愛多謝你.
我們坐在這裡.
我們出街.
我們在順境.
甚至在逆境的當中.
你自己都不斷與我們同在.
主耶和華.
你就是我們的神.
主讓我們在順境.
甚至在逆境當中.
在逆勢的裡面.
我們都依靠你 去親近你.
願上帝你自己繼續引導我們.
正如詩人所說.
願意去監察我們的心.
在知道我們的心思意念.
讓我們有你自己的手去引導我們.
走永生的道路.
因為我們知道每一步都有你同在.
求你繼續幫助我們.
無論什麼情況.
我們都會與你同行.
願你自己一命得主永耀.
祈禱奉尊名求.
在過去幾天.

$^{481}$在網上的一個祈禱.
我覺得是非常有意思的.
他盼望我們多為一些有需要的人去禱告.
我在beach上發了文.
告訴大家.
希望大家留意.
因為我們大部分的人.
可能只是面對生活上.
一些不太方便的時刻.
但我們需要紀念一些在逆勢當中.
生活裡面有危機的人.
特別是一些可能免疫系統上.
比較脆弱的人.
我們需要去留意.
需要去關顧他們.
同樣的.
有些人他們沒有其他選擇.
他們不能夠留在家裡工作.
在外教會裡面也有不少這樣的人.
這些就是醫護人員.
護士或者醫生等等.
特別我們禱告裡面一定要去紀念一些.
前線的醫護人員.
他們去照顧一些受病毒感染的病人.
他們都是屬於高危的人.
我們在祈禱裡面一定要去紀念他們.
同樣我們紀念一些被迫在家裡.
他們要去home school的parents.
一些家長.
他們也是很辛苦的.
在家裡幫助小朋友學習.
他們有很多耐性.
這些都是因為現在沒有實體的班房.
實體的學校能夠去.
所以這些家長我們在禱告裡面要紀念他們.
同樣我們也要紀念一些.
根本在錢財上沒有積蓄沒有依靠的人.
我們有很多的人可能在儲蓄上受到一些虧損.
這些都是真的虧損.
但是要紀念在我們人群當中.

$^{521}$社會裡面有很多需要的人.
還有為了一些無家可歸的人.
我們隔離在家中的.
我們起碼也有瓦遮頭.
但是有很多homeless的人.
他們很需要很多的援助.
另外我們需要憑著這個時間.
這個逆勢的時間.
我們多一點去留意.
多一點去守望身邊我們的鄰舍.
我們的朋友他們的需要.
下一個星期我們一些UV的弟兄姊妹.
他們可能會去Chicago Food Depository.
來做一些義工.
我的女兒說想去參加.
如果弟兄姊妹或者朋友.
你是希望能夠出一點點的面力來幫助的話.
請你來通知我們教木頭工.
我們會安排去Food Bank去事奉.
去做義工的時間.
另外我們知道.
我們一些弟兄姊妹.
他們也會織一些口罩.
能夠幫助到其他人.
或者我們一些弟兄姊妹.
正在做一些收集一些醫護的用具.
保護的用具.
口罩等等.
來幫助一些前線的醫護人員.
這些都是非常有意思的東西.
請你不要忘記.
讓我們不要獨善其身.
在這個逆勢當中.
我們盡我們能力把握每一個機會來事奉神.
再一次多謝你來參加我們網上的崇拜.
我們每個星期日都會播出崇拜的環節.
希望能夠藉著這些訊息.
我們的詩歌.
我們所做的一切.
能夠去宣揚神自己的名.

$^{561}$在這個逆勢當中.
我們去說明我們在基督裡面有的盼望.
如果你今天第一次收看.
我們希望你能夠繼續支持我們.
繼續以奉獻支持我們.
或者以事奉來支持我們.
去分享我們的連結.
去留言.
去讚好.
以致我們成為教會的網台.
我們能夠去服侍更多的人.
願意主自己的話語.
給我們平安和安穩.
我們能夠在這個時候.
向別人宣講.
並且宣講得更加遠.
希望下星期能夠在同一個時間.
在我們網上崇拜裡面.
能夠和你見面.
和一起去事奉.
去敬拜我們的主.
多謝收看成會網台.
記住like comment 或者訂閱 subscribe我們.
拜拜.
\newpage



\section{耶利米哀歌}
\label{sec:G9uX1Fwy_5w}
\textbf{2020/4/5 崇拜直播 vlog \#3 疫境前行 (耶利米哀歌) - 陳崇基牧師}
\newline
\newline
連結: \href{https://youtube.com/watch?v=G9uX1Fwy_5w}{\texttt{https://youtube.com/watch?v=G9uX1Fwy\_5w}} ~~~~ 語音日期: 2020-04-06
\newline
\newline
\hyperref[sec:o3fsHQqXKjc]{\small{< < < PREV SERMON < < <}}
~
\hyperref[sec:index_chronic]{\small{[返順時目]}}
~
\hyperref[sec:index_scriptual]{\small{[返順卷目]}}
~
\hyperref[sec:kAntZ1DnVvY]{\small{> > > NEXT SERMON > > >}}
\newline
\newline
$^{1}$Hello 又是我.
今天是3月30號.
現在是早上大概11點.
去教會拿一些東西.
我經過想去Drop off一些洗手液.
洗手液是不是很想要呢?.
我Drop off去一個UV的家庭.
路上的車少了很多.
都沒有人搭車.
剛剛經過這個火車站.
CTA沒有人搭.
但是還有巴士.
譬如我剛剛看到這架巴士.
但是巴士都沒有人.
只有司機在.
無論如何我們隔離了一段時間.
已經過了兩星期.
在家裡坐太久會發霉.
所以出來走走都好.
我來到Barry的家了.
現在我將這些重要的洗手液交給他.
喂Barry.
(喂好啊).
(不要擋著我洗手液).
(你自己交給他).
(我們不要走太近).
(不要走太近).
(過來).
(跑好一點).
(再來).
(哇).
(兩支就只有).
(有沒有搞錯啊).
(三個人來教會只有兩支).
(哇只有兩支).
(不是吧).
(不是開玩笑的).
(哇這麼多).
(教會真是好啊).
不如說說新聞吧.

$^{41}$今天早上看到一篇新聞.
看了新聞真的氣死了.
那個教會居然有一個師班.
違反了家居隔離令.
照樣幾十人一起去練歌.
誰知道練完之後就發現有人確診.
有時候我不知道這些人在想什麼.
很多時候我們人們有一句術語.
說要有信心.
信心的解釋就是.
總之是靠著心口上.
不管這麼多.
不管三七二十一.
我對上帝有信心.
靠著心口我就衝上去.
上帝一定會保守.
這些好像是試探神.
多過信心.
因為聖經裡面說信心是什麼呢.
聖經裡面說信心.
大部分時間都是強調什麼呢.
不是我們有多強.
不是我們有多有信心.
聖經大部分時間說信心.
是說信心的對象.
不是信心的本體.
信心是什麼呢.
聖經裡面說的信心.
你留心聽.
雖然現在開車.
我們要看路面情況.
也要留意安全.
英文怎麼說呢.
信心.
Faith.
是Focus什麼呢.
不是Subject of Faith.
而是Object of Faith.
中文大概怎麼翻譯呢.
不是信心的本體.

$^{81}$而是信心的對象.
不是Subject.
我就是Subject.
Subject of Faith.
不是說我有多厲害.
我有多大膽.
而是Object of Faith.
信心的對象.
信心的對象是值得我們信任的.
Object of Faith就是上帝自己.
所以聖經裡面經常說.
神是信實的.
上帝是Faithful.
不是我們Faithful.
相反聖經說我們是Faithless.
我們是小信.
我們是不信.
但最重要的是.
神自己是信得過的.
在過去一個多星期的家居隔離裡面.
不知道大家想起什麼聖經.
經常想起神對你說什麼話.
如果可以的話.
在留言那裡.
現在可以分享.
告訴我們最近神提醒你什麼經文.
過去一個星期我經常想起.
耶利米哀歌.
Book of Lamentations.
這本書.
Book of Lamentations.
耶利米哀歌.
我們很少會去說.
為什麼會想起呢.
因為耶利米哀歌.
告訴我們痛苦是很真實的.
是人生必經的階段.
當然我們很多時候.
現代人.
我們經常覺得痛苦是不對的.

$^{121}$或者是一些不好.
不想要的情緒.
所以我們特別在美國.
在現代的生活概念裡面.
我們覺得.
這些情緒在我們生活裡面.
不佔任何地位.
最好不要哭.
最好不要有哀動.
Life is good.
有時候我們對基督教信仰生活的詮釋.
我們都會覺得.
只不過是Life is good.
我們經常都這麼開心.
上帝經常要我們常常喜樂.
常常去感恩.
也沒有錯的.
聖經確實有這麼說.
不過你要想想.
常常喜樂,常常感恩.
這些都是使徒保羅在監獄裡面.
受苦的時候說的.
所以我們不可以經常想.
只不過是Life is good就夠了.
因為聖經裡面很真實地告訴我們.
痛苦是人生必經的階段.
因為我們在一個罪惡破壞了的世界裡面.
我們很多時候因為自己的罪.
或者是一些群體的罪.
或者是一些結構性的罪.
我們陷入了一個低谷的雪.
或者是在混亂的當中.
我們就承受了那個罪的結果.
有一個舊約的學者叫做Walter Buegermann.
Walter Buegermann將所有聖經的詩歌,詩篇.
分成三類.
第一類叫做Songs of Orientation.
就是正確方向的詩篇.
這一類的詩篇就是Life is good.
我們現在在山頂的時候.

$^{161}$每樣東西都好像很順利.
這一個就叫做Songs of Orientation.
那些讚美詩,敬拜詩.
很多時候都是屬於這一種的詩歌.
Songs of Orientation.
第二類他認為叫做Songs of Disorientation.
就是一種迷失方向的詩篇.
這一種的詩歌就是詩人在那種苦難裡面.
他向上帝來哀求.
向上帝來申訴.
或者甚至要上帝來為他申冤.
這些叫做Songs of Disorientation.
第三類的詩篇叫做Songs of Reorientation.
就是重新找到方向的詩歌.
不知道你有沒有試過.
開車或者走路的時候迷失方向.
那些感覺不是很熟悉.
但是當你迷失方向.
你再找回方向的時候.
你的心就會很安定.
你就會很安穩.
這一種的詩篇就叫做Reorientation.
就是重新找到方向.
譬如什麼詩歌是屬於這一類的經歷呢?.
就是一些感恩的詩歌.
譬如說曾經有一個大病.
或者有些苦難.
上帝拯救了他.
上帝去醫治了他.
所以他就找回他的方向.
他就來讚美,敬拜上帝.
耶利敏哥的歌就算是第二類了.
就是Songs of Disorientation.
一種迷失方向的詩歌.
耶利敏哥的歌就有五張經文.
其實就是五段的詩歌拼在一起.
每一段的詩歌.
每一張的經文.
都有二十二段或者二十二節.
因為希伯來文的字母是有二十二個的.

$^{201}$所以每一段或者每一節.
它是由一個字母開始.
從頭到尾.
有二十二段或者二十二節.
如果你查一下希伯來文的聖經.
你就知道第一,第二和第四章.
都有二十二節.
每一節都是由頭到尾.
希伯來文的字母排列出來.
中間那一章第三章.
就有六十六節這麼長.
因為每三節或者每三段.
就是一個希伯來文的字母開始.
你看到它這樣排列的時候.
就好像在耶路撒冷淪陷的.
慘況,苦難當中.
那種混亂中嘗試找出方向.
好像生活被打亂.
很多東西失控的時候.
到國難當前的時候.
很多人民受到打壓,被擄.
流亡到其他地方.
那種慘淡的當中.
他嘗試用詩歌去整理.
他自己心中的情緒.
所以就用這樣的方法.
去表達內心深處的痛苦.
就像當我們很抑鬱,很難過的時候.
很多時候我們都會藉著詩歌,詩詞.
去表達我們心中的哀嘆.
第一章.
從前人煙稠密的城中.
現在怎麼竟然孤獨地坐在那裡.
從前在列國中偉大的.
現在怎麼竟然像個寡婦.
從前在眾省份中為王侯的.
現在怎麼竟然成為奴隸.
耶路撒冷在困苦,漂流的日子.
就想念他昔日的一切珍寶.
他的人民落在敵人手裡的時候.

$^{241}$沒有人幫助他.
敵人看見他就譏笑他的滅亡.
上星期去了Chinatown.
拿著鏡頭去.
以前真是人煙稠密的Chinatown.
現在居然變成一個人影都沒有的地方.
然後去了Downtown那裡逛逛.
情景都是一樣慘淡.
你可以想像一下當時耶路撒冷淪陷的時候.
那種淒慘的情況.
那些婦孺,平民百姓.
甚至官員都被拉被斬.
帶到異鄉.
以前耶路撒冷是世界的中心.
很多遊客,很多朝聖者.
聚集在耶路撒冷守節.
人們被遷徙和流亡到另一個地方.
而剩下的只有幾隻小貓.
耶路撒冷特別的地方就是.
嘗試在苦難中亂中尋序.
用字幕的思維去找出秩序.
但到第五章的時候.
發現它是完全亂了.
不再是跟著A to Z.
全部的字幕都亂了.
結果在第五章最後的時候是沒有答案的.
第五章最後說什麼呢?.
他說:但耶和華,你永遠坐著為王.
你的寶座代代長存.
你為什麼永遠忘記我們?.
為什麼長久離棄我們?.
耶和華,除非你真的棄絕我們.
你向我們所發的怒氣難收.
不然求你使我們歸向你.
好使我們回轉過來.
求你更新我們的日子.
像古時一樣.
整本書到最後的時候.
給我們一種不是答案.
而是苦難中的張力.

$^{281}$張力就是現在現實.
我們看到很多的疫情,苦難,失控.
而他仰望到將來.
將來才能夠有秩序.
當上帝歸回歸向我們.
憐憫我們的時候.
我們很多時候就在這種張力裡面.
不一定現在有答案.
但是耶利米在哀嘆裡面.
他來到神的面前.
這就是耶利米哀歌寶貴的地方.
他不去告訴我們.
你受傷,連塊膠布就算了.
反而他說你受傷.
你去痛苦的時候.
你可以來到上帝的面前.
因為上帝不會覺得.
這些事情是我不會聽的.
上帝不會覺得你受苦的時候.
我會覺得很煩.
你儘管來到我的面前.
我會覺得很煩.
你儘管來到我的面前.
我會覺得很煩.
你儘管來到我的面前.
我會覺得很煩.
你儘管來到我的面前.
我會覺得很煩.
你儘管來到我的面前.
我會覺得很煩.
你儘管來到我的面前.
我會覺得很煩.
你儘管來到我的面前.
我會覺得很煩.
你儘管來到我的面前.
他所說的哀嘆.
不是單單是個人的哀歌.
我們知道詩歌裡面有兩種不同的哀歌.
第一種叫做 Individual Lament.
是一種個人的哀歌.

$^{321}$我們經常在詩篇裡面提到.
就是詩人大衛受苦難的時候.
他被追殺的時候.
他個人來到上帝的面前哀求.
那種叫做個人的哀歌.
第二種哀歌也有很多.
不是個人的.
就是群體性的.
Communal 或者是 Collective.
或者是 National.
或者是國家的群體一起哀嘆.
這些詩歌有很多.
耶利米哀歌就是這種 National Lament.
就是國家性的群體性的一種哀歌.
就是說我們現在受的苦難.
不單單是我們自己一個人去受.
而是我們整個群體,整個國家.
甚至整個全球的人去受這種苦難.
這裡其實提醒我們.
我們的信仰和信仰的經歷.
我們不單單是得種個人的經歷.
我們需要被提醒.
我們不可以成為一個宗教信仰的個人消費者.
我們只和我們神之間有個人的關係.
這是必然的.
但是同樣更加重要的.
就是我們整個群體在上帝面前有一個關係.
而在苦難裡面我們整個群體哀求上帝.
在上帝面前去 lament.
停下來看看第一章十一節.
他說他所有的人民都在哀嚎都諸尋覓食物.
他們為了維持生命拿自己的珍寶去換取糧食.
他說要我求你垂看求你監察.
因為我被人藐視.
這裡說到他是女性的他.
很多時候詩人,詩人就會將耶路撒冷.
或者整個以色列文去形容成一個女子.
當她很無助的時候.
她就去到深淵.
在上帝面前作哀過.

$^{361}$去買東西.
耶律勒的寫書這一節很寫實.
當耶路撒冷淪陷的時候.
那些人民就會用金銀珠寶來換取糧食.
平時我們覺得一些理所當然的.
糧食,食水甚至廁紙.
你會覺得這些東西理所當然必定會有.
但當逆勢的當中居然比金銀珠寶更加罕有.
在逆勢當中我們要了解.
每一個人都受到打擊.
每一個人的生活都會受到很大的影響.
所以在逆勢當中我們真的需要去思考一下.
人人都有責任.
對社會有責任.
對我們身邊的鄰居有責任.
因為我們每個人受影響的時候.
我們就需要彼此守望相助.
能夠彼此去鼓勵.
(廣播).
其實很多東西都是有貨的.
其實不需要去擔心那麼多.
我們這裡是很豐足的社會.
只不過恐慌的時候.
很多時候人們就不理會那麼多.
每樣東西都立即先拿走.
你看到有很多東西.
甚至剛才廁紙也有不少.
不過限制了.
每人只買一包.
今天出來看到市集市場.
其實多了一些人戴口罩.
其實這樣就好了.
早就應該要戴了.
這就是一個National Lament.
一個群體的哀歌.
寫實的地方.
那種痛苦不單止是個人.
也是整個群體需要去分擔.
在這個疫情當中.
提醒我們和其他人一樣.

$^{401}$要去守望相助.
和其他鄰舍.
無論我們認識他們也好.
不認識也好.
我們其實每個人都是在同一頁面.
大家都是受到這個逆勢來打擊.
耶利敏哀歌不單止讓我們看到.
個人的受苦是很真實.
並且是一個群體受苦也是很真實.
最後我相信他也說到.
是上帝與人受苦.
在第二章那裡.
他有一句說話.
我覺得是今時今日很切身的.
就是他在哀求的時候.
他告訴上帝.
原來上帝正在聽他的哀求.
我停下來讀一讀這段聖經.
第二章第十八節.
石安,即耶路撒冷,居民的城牆.
你要一心向主哀求.
願你的眼淚像江河般湧流.
晝夜不息.
願你得不著歇息.
願你眼中的同人不能休息.
第十九節.
夜裡每度交更的時分.
你要起來呼喊.
在主面前你要傾心如水.
你的孩童在各街頭因飢餓而昏倒.
你要為他們的性命向主舉手.
禱告.
這是最詫異的地方.
上帝很想人受苦的時候.
來到祂面前傾心如水.
和我們一起來受苦.
快要到受苦節的時候.
快要紀念耶穌基督十字架上.
和我們一起來受苦.
這一位的神,這一位的主.

$^{441}$不是一個高在上的教案官房.
給我們一些答案.
而是親自到道成肉身來到我們中間.
和我們一起來經歷這個痛苦.
並且他經歷的痛苦是最大的痛苦.
是十字架的羞辱,十字架的痛苦.
首先我們說到整本愛國最後最後第五章結束的時候.
有一種張力.
他在求問上帝.
為什麼上帝要你離棄我.
為什麼上帝好像看不到我現在受苦.
正在經歷的事.
真正的答案在哪裡呢?.
真正的答案原來不是在愛國.
不是在整本書最後最後那幾節.
真正的答案原來是在整本書中間第三章.
我們說的第三章就是最長的那一章.
六十六段,六十六節.
那裡說到一段經文.
就是耶利米先知.
他在苦難當中.
在國家淪亡的時候.
他來到仰望上帝.
那裡就是整本書最精髓和最中間的地方.
第三章第十九節.
當回憶我的困苦漂流.
就像苦梗和毒草.
每逢愛心想起往事.
我的心就消沉.
但是我的心一想起下面的這件事.
我就有指望,我就有哭.
第十二節,耶和華的慈愛永不斷存.
他的憐憫永不止息.
每天早晨都是新的.
你的信實多麼廣大.
我心裡面說,耶和華是我的業份.
所以我必仰望他.
就像耶利米在哀歌的精髓最中間的地方.
他就說想起上帝的信實.
上帝以往在很多災難和困苦當中.

$^{481}$他沒有丟棄過他的子民.
在歷史裡面,他證明了.
這位神從來由始至終都會守約.
他就用這個作為基礎的信念.
來仰望著.
在逆勢當中,我們可能有些很難受的時候.
但是讓我們知道神和我們在一起.
甚至神願意和我們一起去受苦.
接下來的環節就是一個很好的時機.
就是讓我們知道,縱使我們在逆勢裡面.
我們知道神沒有丟棄我們.
耶穌基督在十字架上和我們一起去受苦.
在苦難當中,祂給我們一個答案.
就是祂願意和我們一起去受苦.
願意拯救我們.
願意給我們一個永生的盼望.
記不記得我開始的時候問大家一個問題.
就是過去一個星期,神有沒有提醒你一些經文.
祂自己的說話能夠將你帶領到祂自己面前.
神自己的華語能力就在這裡.
有很多時候不安在我們逆勢裡面.
我們可以藉著上帝的華語.
祂的那種力量來親近祂,靠近祂.
第二個問題,或者可以問大家.
就是過去一個星期有沒有藉著一些詩歌.
上帝將你帶領到祂自己面前.
無論在不安的當中給你安慰.
或者在煩擾的裡面.
這些詩歌教導你去禱告,去親近祂.
過去這幾天在社交媒體.
經常都有朋友去分享.
這位已故樂壇的巨星張國榮.
他的金曲和一些很動人心弦的歌曲.
都是張國榮去演繹的.
我不知道你記不記得這隻白膠碟.
如果你曾經見過或者handle過.
你的年齡已經被透露出來了.
我還記得這隻白膠碟剛面世的時候.
在唱片簿我去等它.
然後立刻買.

$^{521}$因為裡面的音樂很能夠去說服我們.
能夠打動我們的情緒.
這幾天想起一首詩歌叫做.
只有你永恆的主.
這首詩歌在不安的時候.
成為我的禱告.
在這個逆勢當中.
希望這首詩歌能夠幫助你.
來去教近神.
誰知心痛苦 誰解心裡憂.
誰可於這刻伴我身邊數.
只有你永恆的主.
每一刻等待靜聽心聲.
誰犧牲苦中 誰寫專講悲.
誰於漆黑中讓每點光照遍.
只有你永恆的主.
帶出光輝 令我溫暖.
人海中常失落.
每天感到壓迫寬鬆.
還願你每天也紀念.
求恩處溪底.
是我能知道.
你是永恆.
你是盼望常在我心.
容許我靠近你.
永遠沒掛慮.
容許我愛慕你.
因你是我力量.
明天縱有幻變.
求你引導.
願此生全為你走過.
誰犧牲苦中 誰寫專講悲.
誰於漆黑中讓每點光照遍.
只有你永恆的主.
帶出光輝 令我溫暖.
希望你能夠繼續支持我們成會社區教會的事工.
和成會網台的事工.
其中一個支持我們的方法當然就是奉獻.
你可以用不同的電子奉獻方法.
來支持我們.

$^{561}$以致我們的器材.
或者我們的訊息能夠做得更加好.
另外一個支持方法就是.
來like share comment.
和subscribe去訂閱我們成會網台.
在youtube那裡請你按subscribe訂閱的字.
然後再按那個叮叮.
你就可以看到我們每一個新的影片.
希望能夠和你繼續上路.
我們都想到.
(遊戲畫面).
\newpage



\section{帖撒羅尼迦前書使徒行傳}
\label{sec:kAntZ1DnVvY}
\textbf{2020/4/12 復活節主日崇拜直播 vlog \#4 疫境前行 (帖撒羅尼迦前書 4; 使徒行傳 17) - 陳崇基牧師}
\newline
\newline
連結: \href{https://youtube.com/watch?v=kAntZ1DnVvY}{\texttt{https://youtube.com/watch?v=kAntZ1DnVvY}} ~~~~ 語音日期: 2020-04-13
\newline
\newline
\hyperref[sec:G9uX1Fwy_5w]{\small{< < < PREV SERMON < < <}}
~
\hyperref[sec:index_chronic]{\small{[返順時目]}}
~
\hyperref[sec:index_scriptual]{\small{[返順卷目]}}
~
\hyperref[sec:HEMq3Mncf_A]{\small{> > > NEXT SERMON > > >}}
\newline
\newline
$^{1}$[音樂].
今天是復活節的主日.
我們來慶祝耶穌基督從死裡面復活.
耶穌我們的主.
祂的墳墓是空的.
這個慶祝是基督教的教會最重要的一天.
復活節的主日.
Easter Sunday.
是全年裡面我們的日曆最重要的一個主日.
不過今年這個主日就非常特別了.
因為全世界每一個地方.
很多的教會都不能夠實體地在一個地方聚會.
來崇拜和慶祝耶穌的復活.
你想清楚.
其實也不要緊的.
無論是實體崇拜.
或者是網上的崇拜.
我們最重要知道的.
是耶穌基督從死裡面復活.
祂是得勝的.
祂是勝利的.
祂是打敗了死亡和陰間的權柄.
所以勝利是屬於我們的主.
王者耶穌基督.
所以用什麼方法來崇拜呢?.
其實沒有什麼所謂.
最重要的是我們知道.
這一個是在基督的信仰裡面.
我們有的永恆的盼望.
想起復活節的時候.
我們就想起在歷史真實的時空實體裡面發生的事情.
就是耶穌在三天前.
祂在十字架上死去.
然後被埋葬.
然後在主日今天復活.
在基督教的信仰.
聖經的信仰.
最特別的地方就是這裡了.
就是信仰的核心.
不是一些抽象的理論.

$^{41}$也不是一些學說.
也不是一些單單講到人生怎麼做人的道理.
核心的信仰.
是放在一個真實時空裡面發生的歷史事實.
歷世歷代的人.
不只是現代的人才會這樣.
古來的人也是一樣的.
就想起死而復活的時候.
會不會有一點點荒誕呢?.
會不會有一點點奇怪呢?.
死人會復活.
到底墳墓是不是真的空的?.
每一年我都會想起寫得很好的一本書.
這本書是由一個以色列的歷史考古學家所寫的.
他叫做Simon Gibson.
他在他的那本書叫做.
《最後的日子》.
Archaeological Evidence.
那裡做了一個研究.
研究什麼呢?.
就是研究初代教會第一世紀的文獻.
以及研究教會以外有的文獻.
他將這些收集在一起的時候.
分析了他們的資料等等.
發現原來這些文件所記載的東西.
是非常確實的.
他的研究結論是這麼說的.
他說有很多人嘗試將耶穌空墳墓這件事情.
去解釋掉.
就是將它用不同的方法去解釋.
沒有這些事情發生的.
有些人會說.
可能耶穌根本沒有死.
可能祂只是暈倒在神賜卡上面.
然後被人堆在墳墓裡面.
然後他又再被人以為.
他是從死裡面復活.
其實他沒有死.
這是其中一個理論.
另外一個理論可能是一些門徒.

$^{81}$偷了耶穌的屍體.
等等.
Simon Gibson.
他本身是一個以色列的猶太人.
他也不是信耶穌的.
他做了一個很中肯的結論.
結論是什麼呢?.
他說在這麼多理論當中.
其實他說這些都不能夠解釋.
歷史裡面曾經耶穌的墳墓是空的.
如果你要去解釋.
耶穌的墳墓為什麼是空的.
這是超越了歷史科學的範疇.
唯獨是.
除非你去接受聖經的神學信仰解釋.
這個神學信仰解釋就是.
神使耶穌從墳墓裡面復活.
說到耶穌從死裡面復活.
一個最好去解釋這件事情的地方.
就是這裡了.
墳場.
這裡是墳場.
這個墳場叫做Mount Auburn Cemetery.
Mount Auburn墓園.
在大芝加哥區這個地方有幾個大的墓地.
幾個大的墳場.
這是其中一個.
來這個地方很多次.
很多時候都是為朋友,論社,教會會友.
做安息禮拜,做喪禮.
每逢來到這些地方.
我相信大家的心情都一定是很複雜的.
首先複雜的地方就是.
我們不喜歡來.
我們盡量想避開.
不來墳場這些地方.
因為這些是生人物近的地方.
另外我們又覺得不能夠避免.
為什麼不能夠避免呢?.
無論是你的親友,你自己自身的經歷.

$^{121}$應該可能都不少.
一些親友要離開這個世界.
我們有時被迫或者必須要來到這個地方.
我們做牧師的當然.
我們來這個地方事奉.
所以來到這個墓地墳場的地方.
我們總會覺得有點不自然.
我曾經在這裡做過很多安息禮拜,落葬禮.
在這裡流過很多眼淚.
這是一個很令人悲傷的地方.
但來到這個地方你就想到人是多麼脆弱.
生與死真的不到我們作主.
剛才在這個墓地隨意散散步.
看到很多林林種種的墓碑.
有大的,有小的.
有土葬的,有火葬的.
我留意到有很多不同的年份.
有些是1910年代出生的.
我看到一個是1946年.
也看到一個是1971年.
也看到一個是1989年.
更加令我痛心的是看到一個是2014年.
(悲傷的音樂).
特別現在在疫情當中.
全球的人都受到影響.
無論你是什麼背景.
無論你是社會地位高或低.
你是什麼膚色.
我們每一個都被捲入這個疫勢當中.
我們真的感覺到人很多東西.
我們不能控制,控制到.
生與死,病,健康.
有沒有工作.
我們的經濟完全不能在我們手中把握.
前幾天我去了唐人街.
經過的時候我看到很多情景.
這些是我從來沒有見過的.
我現在在Chinatown.
這裡水靜河飛.
我想知道哪一間餐館是開的.

$^{161}$這裡是開的.
不是,這裡已經關了.
有一間叫江湖烤肉.
關了.
老四川還開.
我想可以買外賣的.
明軒.
明軒已經關了.
先看看.
裡面真的有人打掃.
是不是一些禮物店.
全部都不開.
關了.
這裡還請人.
這裡是Cheese Cafe,即是旺角,也是關了.
李榮華,關了.
但是皇上王呢.
皇上王在開.
還可以有外賣的.
開了.
411是開的,是喝的.
明軒也關了.
NCCB也關了.
全部都關了.
有些人在喝酒.
在坐椅.
強記關了.
長城Great Wall Restaurant關了.
長城燒臘,沒有燒臘吃了.
人影都不見了.
知義城還開,可以買東西的.
Essentials,這些都關了.
左一條圈.
(風聲).
剛才在芝加哥的Chinatown逛了幾分鐘.
我告訴你我個人的感受是什麼.
我在芝加哥這裡都有20年多的時間.
我沒試過在Chinatown白天1點,2點鐘.
水靜河飛.
我從來沒試過找到位置,閉上眼也找到.

$^{201}$冠狀病毒對飲食業很大打擊.
我們教會也有飲食業,酒店業等等.
這些行業的影響很大.
真的要為這些大英姐妹去守望,祈禱.
是一個很抑鬱的場面.
在今天復活節主日.
我想起很多聖經裡面說到.
一個復活的信仰.
我們相信耶穌基督復活的時候.
給我們一個怎樣超越死亡的盼望.
其中一段聖經是貼在我們家前書第四章.
他說「論到誰了的人,我們不願意你們不知道.
免得你們憂傷,就像那些沒有盼望的人一樣.
我們若信耶穌死了又復活了.
照樣也應該相信那些靠著耶穌已經睡了的人.
神必定把他們和耶穌一同帶來」.
這裡形容一個死而被埋葬的人.
他是一個睡了的人.
即是一個睡了覺的人.
睡了覺的意思不是一個最終的終結.
而是一個等候.
他在等待耶穌再回來的時候.
他會把新天新地帶到我們中間.
並且會帶一些在基督裡面已經睡了覺.
在基督裡面去世.
安息主懷的人一起帶回來.
然後把新天新地帶到地上.
這就是死而復活的耶穌.
可以給予我們的盼望.
因為我們知道死亡不是一個句號.
死亡是一個優質符號.
我們死了之後會與主相圓.
但是有一天耶穌會再回來.
然後給我們一個復活的身體.
好像祂自己榮耀的復活身體.
當我們想到復活節的時候.
我們不要只想只是一個歷史時空發生的事.
這個當然沒錯.
但同時我們要想起.
這個復活的主耶穌基督.

$^{241}$祂是一個活著的神.
祂是每一天和我們同在的神.
並且祂與死復活的時候.
祂就帶給我們一個超越疫病,超越死亡的盼望.
《貼撒萊加前書》第四章十五節.
「我們現在照著主的話告訴你們」.
這是初代教會的信約.
一個信條.
我們照著主的話.
曾經這樣教導我們.
現在再一次向你重申.
「我們這些活著存留到主再來的人.
絕不能在那些碎了的人而先.
因為主必親自從天降臨.
那時有法靈的聲音.
有天使長的呼聲.
還有神的號角聲.
那些在基督裡面死了的人.
必先復活.
然後我們還活著存留的人.
必和他們一同被提到魂裡.
在空中與主相會.
這樣我們就要與主常常同在.
所以你們應該用這些話彼此勸勉」.
《貼撒萊加前書》第四章十五節.
「我們現在照著主的話告訴你們」.
這是初代教會的信約.
一個信條.
我們照著主的話.
曾經這樣教導我們.
現在再一次向你重申.
「我們這些活著存留到主再來的人.
絕不能在那些碎了的人而先.
因為主必親自從天降臨.
那時有法靈的聲音.
還有神的號角聲.
這段聖經講到有歡呼聲.
有號角聲.
一些很大聲的歡迎聲音.
這些是當時羅馬帝國的人民很清楚.

$^{281}$很熟知的一些情景.
每逢羅馬帝王出巡.
或者一些大官貴人.
去訪問一些城市的時候.
城中的百姓長老就會很興奮.
他們就會去組織一些歡迎的列隊.
他們就會去到城外.
當大官或帝王還沒到之前.
他們就會在路上.
用號角聲.
用歡呼聲來歡迎他.
在途中與他相會.
其實相會這個字也是一個.
羅馬帝國的政治的.
專有動詞.
meet together.
城中的百姓長老.
就會走出去與帝王.
帝王.
相會.
然後就將他escort.
帶他進城裡.
然後就會有很多歡呼.
很多慶祝.
對於羅馬人來說.
這個很重要.
因為一個帝王出巡.
能夠訪問你的城市.
或者一個大官貴人.
去到你的城市訪問的時候.
一定會帶來很多好處.
即是說看得起你.
去到你的城裡給你很多福利.
很多權利.
所以你可以想像當時的.
羅馬人.
很希望這些皇者.
羅馬帝國的皇者能夠駕臨.
在這幾個星期芝加哥的疫情當中.
你會看到很多的公路.

$^{321}$很多的馬路.
甚至市中心.
平時是很多的車輛.
以及很多很繁忙的交通.
現在突然間空了.
不過讓我想起.
其實也曾經有過這樣的情況出現.
就是每當美國總統.
去到訪問芝加哥.
來到這個城市裡面.
有一個大會的時候.
那些保安.
或者警方.
就會將這些路全部封了.
我曾經記得.
那個公路突然間空了.
發生什麼事呢?.
原來美國總統來到芝加哥.
有訪問.
這種的情景.
就好像一個羅馬帝王的皇者.
駕臨的時候.
那些歡迎會議.
就走出去.
去城外面.
來迎接皇者的駕臨.
因為這個皇者會帶給他們.
很多的好處.
這個也是城裡面.
每一個居民他們有的.
盼望.
你知道貼索雷加前書.
就是寫給貼索雷加這個城市.
一個希臘羅馬城市.
裡面的.
教會.
在貼索雷加這個城市裡面的居民.
很清楚這些羅馬人的政治意識.
如果你翻開去到.
《士徒行傳》第17章.

$^{361}$那裡講到保羅和他的同工.
他們去到.
貼索雷加傳道的時候.
發生了一件.
很妖嚴的事.
第17章第一節講什麼呢?.
他說保羅和西蘭.
經過不同的地方來到貼索雷加.
他就按照他的習慣.
一年三個的安息日.
就去到猶太人的會堂裡面.
來到廣道.
用聖經.
舊約的聖經來證明.
這位耶穌一定是要先.
受害受死被埋葬.
然後復活.
這個才是真正的基督.
就是猶太人一直所盼望的救世主.
也是西人.
希望能夠見到的一位救世主.
保羅就這樣用聖經.
跟他們辯論.
然後證明耶穌就是基督.
中間有些人.
特別是一些比較顯耀的人.
或者一些官長.
或者一些社會地位比較高的人.
就相信了耶穌.
但是中間有一些猶太人.
特別是一些有權力的猶太人.
他們就質問.
然後他們就找了.
在城市裡面的一些流氓.
他們就找了.
在城市裡面的一些流氓.
他們想拘捕保羅.
想趕走保羅.
找不到的時候.
就找了保羅的一些同工.

$^{401}$然後就抓了他們.
去到一些地方的長官那裡.
然後就去指控他們.
那個指控是什麼呢?.
是一個很政治性的指控.
他說這些搞亂天下的人.
來到我們這裡.
他們違反了.
該殺的法令.
他們去傳有另外一個王.
耶穌.
就是說保羅在傳道的時候.
說到這個死而復活的耶穌.
他們就違反了羅馬帝國王者.
那個凱撒斯的法令.
因為他們在說的是另外一個王.
大家可以想像為什麼那些貼薩尼加人.
會大動肝火.
指控為什麼是一種這麼煽動性.
政治性的指控.
他們違背了.
該殺的法令.
羅馬王者的法令.
傳另外一個王.
正如我所說.
如果你是一個貼薩尼加城裡的市民.
你一輩子.
你最大的指望.
最大的寄望.
就是王者.
羅馬帝王或者他的官員.
來到你的城市裡.
你心裡面只有這位主.
而這位主能夠給你安穩.
給你太平.
給你經濟好.
給你健康.
所以一個城市.
親羅馬王是很應該的.
但是保羅.

$^{441}$去到這個城市的時候.
他就說起死而復活的耶穌.
他說這個死而復活的耶穌.
是生死的.
他才是我們的王者.
這個羅馬王.
或者我們地上的所有的權柄.
所有的政權.
其實都是和我們一樣.
都是這麼脆弱.
沒有人能夠去掌握生與死.
換句話說.
這個福音就告訴我們.
耶穌從死裡復活.
就證明他才是真正的王者.
這個政權.
地上的人.
我們都不能夠掌握到生與死.
如果我們.
只是寄望在.
我們自己身上.
寄望在這些政權.
政客.
或者人的身上.
我們的信心投資一定會破產.
所以.
死而復活的耶穌.
這個信仰.
就告訴我們.
與其我們去依賴.
和我們一樣脆弱的勢力.
不如我們去養賴.
這個死而復活的主耶穌基督.
我們養賴主耶穌基督.
不代表我們要離地.
因為.
在空中與主相遇.
我們就好像那些.
歡迎的群眾.
我們走到城外.

$^{481}$在空中與主相遇.
其實最重要的目的.
不是我們上天堂.
而是我們將耶穌基督帶到地上.
去到escort.
去到引領王者.
來到我們生命的地上.
Thy kingdom come.
Thy will be done.
願你的國光臨.
行在地上如同行在天上.
在空中與主相遇.
其實就是想去迎接.
這位主耶穌基督.
王者.
駕臨到我們的生命當中.
在我們的生活,在我們城市裡面.
能夠成為其他人的福氣.
城匯社區教會的弟兄姊妹.
他們很好.
他們就自發.
自己車衣,自己織東西.
就做一些.
自製的口罩.
裡面可以放一個.
濾棉.
給一些有需要的.
鄰舍,因為這裡已經.
一早,這些口罩.
美國早就掃清了.
所以有些弟兄姊妹.
他們自願的,24/7.
不停地,不停地做.
就是做一些.
自家的口罩,自己衣車做.
自己手織,我們已經做了.
幾百個了,已經.
他們這麼熱心,來服侍.
我們的鄰舍,所以今天.
做了第一批,這一批.

$^{521}$我們會先送去.
第一家的老人大廈.
然後我們就會逐批逐批.
送去一些有需要的鄰舍.
包括一些.
家護助理.
等等,他們都.
是一個很高危的情況.
所以今天是非常興奮的.
我們有一個Mission.
叫做Show me your love.
是不是很想看看第一批的.
自製的口罩的製成品.
是怎樣呢?我也很興奮,很想看.
你跟我來吧.
好,來吧,我們有第一批的.
自製的口罩.
製成品,嘩,多漂亮.
我打開來看看.
先看看.
這裡.
這個口罩,嘩,做得挺好.
很漂亮,先看看.
哦,這樣.
兩邊有繩子可以綁住.
我們有這個成會教會的說明書.
這裡說清楚.
這個口罩是怎樣用的.
以及設計是來自哪裡.
當然這個不是工業專業的程度.
但是我們送出去的時候.
說清楚.
這個是一種Hong Kong made.
DIY的.
收了之後有一個聲明.
需要用戶自己去清潔.
中間的濾芯要自己去供應.
有中文,有英文.
希望這些DIY的口罩.
能夠幫助到一些鄰舍.

$^{561}$好,我們來到這個唐人街的知英大廈.
我們是成會社區教會.
你們可不可以戴上口罩.
Hello,你不介意我拍一點東西嗎?.
我們放在這裡可以嗎?.
我們把口罩放在這個盒子裡.
不要緊,我也知道的.
我們也有距離.
這位姑娘很好人.
我們是成會社區教會UV的.
我們將每一個.
這裡有大概一百多個口罩.
每一個都分開包的.
裡面有說明.
中英文.
希望能夠對你有幫助.
OK.
好,謝謝你.
我現在來到.
這裡其中一層.
有很多這些.
預防的指示.
以免人們感染到病毒.
這裡每一個單位.
都有一個.
Hello,Hello,Hello.
我們現在進去.
我們送一個.
自製的口罩給你.
早上防毒工作.
我做完.
很好.
很聰明.
OK,謝謝.
每一個家戶.
都有一個列表.
可以檢查他們的情況.
每一天都有.
會跟他們探熱.
檢查後.

$^{601}$他們會檢查.
任務完成.
Mission accomplished.
我將口罩.
送給裡面的用戶.
送了150個.
我覺得這個非常有意義.
一些大英姐妹.
義工他們自己.
夜車製作了一百多個口罩.
不止一百多個.
我們第一批製作了幾百個.
剛剛送了一百多個.
這些非常有意義.
首先多謝每一位同工.
我們希望所做的事情.
作為一個教會.
作為神的子民.
能夠代表基督的愛.
能夠關顧到我們身邊的鄰舍.
我相信在疫情當中.
怎樣嚴峻都好.
我們作為天國的子民.
當基督進入我們心靈裡的時候.
我們總會有事奉主的機會.
可能用心留意一下.
在我們社區.
教會.
在我們的城市裡.
可能有很多人,很多鄰舍.
他們比我們更加脆弱.
我們不妨用一些實際的行動.
來關心一下他們.
去事奉一下他們.
將我們的王者帶到地上.
將我們耶穌基督復活的生命.
帶到人群裡.
去引滾那裡.
不要看著天爆炸.
耶穌隨時會回來.

$^{641}$不過耶穌回來之前.
記得訂閱我們城會網台.
按訂閱鍵和鈴鐺.
下次就會通知你.
什麼時候有新片出.
\newpage



\section{雅各書}
\label{sec:HEMq3Mncf_A}
\textbf{2020/4/19 智慧廚房 vlog 信息 \#1 (雅各書 1) - 陳崇基牧師}
\newline
\newline
連結: \href{https://youtube.com/watch?v=HEMq3Mncf-A}{\texttt{https://youtube.com/watch?v=HEMq3Mncf-A}} ~~~~ 語音日期: 2020-04-18
\newline
\newline
\hyperref[sec:kAntZ1DnVvY]{\small{< < < PREV SERMON < < <}}
~
\hyperref[sec:index_chronic]{\small{[返順時目]}}
~
\hyperref[sec:index_scriptual]{\small{[返順卷目]}}
~
\hyperref[sec:_qut0JfQmyc]{\small{> > > NEXT SERMON > > >}}
\newline
\newline
$^{1}$各位好.
我們今天有一個新的訊息系列.
叫做「智慧廚房邊煮邊講」.
為什麼呢?.
因為現在家居隔離.
我們每個人都躲在家裡.
簡單簡單地煮一頓飯.
其實我看很多Facebook.
他們的訊息一點都不簡單.
但是無論如何.
我們今天會煮一些簡單一點的東西.
我告訴你什麼叫做簡單簡單地煮一頓飯.
不過我們邊煮邊講.
講什麼呢?.
就是講《亞國書》.
Book of James.
非常重要.
我告訴你.
我們煮什麼.
今天煮什麼呢?.
我們煮焗火腿.
為什麼煮焗火腿呢?.
因為我太太買了這個很大塊的火腿.
她說很便宜.
5.89元.
我說好像很便宜.
我不知道.
我吃米不知道價錢的.
我們就嘗試去做這件事.
我已經在網上找到一個.
糖漿焗火腿的recipe.
預備時間是30分鐘.
煮的時間是一個半小時.
其實我從來沒有做過這些事情.
其實我很佩服那些.
無論是爸爸或者媽媽.
總之你在家裡.
要煮三餐.
一味想.
又不可以repeat.

$^{41}$又不可以經常煮來煮去.
都是那幾件事.
所以我很佩服他們.
每天去買菜.
煮飯.
其實有點像我們牧師.
我們牧師說到.
我們都不可以repeat.
我們不可以經常炒那幾道菜.
起碼我們這些駐堂牧師.
我們不可以一曲走天涯.
所以我們.
一於邊煮邊說.
我很少下廚.
所以做這個東西都很大壓力.
無論如何.
食譜就是這樣說的.
所以當然一塊火腿.
另外有一些cloves.
有honey.
橙.
因為拿了橙汁.
就倒進火腿裡.
才放進去焗的.
說起來很容易.
但是好像一點都不簡單.
煮一餐的話.
我希望待會我們完成了這個火腿的時候.
起碼今晚有東西吃.
因為外面隔離令.
很多餐館都不開了.
現在唯有吃自己.
第一步.
我們將這個火腿.
放在room temperature.
室內的氣溫.
就拿出來了.
我已經一早拿出來了.
所以準備了.
然後就要切開它.

$^{81}$先開了它.
說起這個下廚.
其實有點像做人.
因為我們傳統來說.
我們經常說.
每天都要吃.
每天都有廚房.
做人就好像在廚房一樣.
為什麼這個訊息是你呢?.
我們說到廚房.
你知道大家現在因為疫情的緣故.
全球的人都要留守在家.
所以經常都要想煮什麼.
想想自古以來.
廚房是人生的中心.
我們沒有了吃就很慘.
每天都要吃三餐.
所以廚房裡面有不少.
所謂人生的道理.
都是由我們吃.
由我們煮開始的.
所以這個地方廚房.
其實在我們生命裡面有很多智慧.
我們需要去學.
特別是現在在逆勢當中.
神似乎叫我們回到最基本的東西.
煮和吃.
簡簡單單.
基本的一餐.
是上帝供應給我們.
生命的基本當中.
逆勢裡面我們都需要去想這個東西.
現在我在做什麼呢?.
我將這個 ham 切成十字.
每個地方都要切成十字.
然後將丁香塞進去.
以致焗的時候就會有味道.
說回廚房確實是真的.
有很多人生道理在當中.
聖經也說了很多廚房.

$^{121}$平時人怎麼生活.
一些很實際的智慧.
所以為什麼我選了《雅各書》.
《雅各書》是說什麼呢?.
如果你有印象的話.
《雅各書》只有幾章.
五章.
《雅各書》是說一些很實際的基督徒生活.
一些基本的基督徒的生活.
《雅各書》是寫給當時的教會.
信徒面對很多困難.
他們就要重整他們自己的生命.
其實有一點像我們今天這樣.
我們也需要去重整一下自己的生命.
剛才我將火腿切成十字.
他說只是切皮和裡面的肥.
不要切裡面的肉.
用丁香.
他說要先塞進去.
然後才放醬汁.
我們先找碗.
《雅各書》也是一本智慧的書.
裡面講了很多的題材.
是跟舊約裡面的智慧文學很相似的.
譬如講到人遇見苦難的時候.
他需要去忍耐.
在我們的日常生活裡面.
我們需要在言語上面有智慧等等.
不如讀一下《雅各書》第一章.
「神和主耶穌基督的僕人雅各.
向散居各地的十二支派問安.
我弟兄們你們遭遇各種試煉的時候.
都要捍衛喜樂.
因為知道你們的信心經過考驗.
就產生忍耐.
但忍耐要堅持到底.
使你們可以完全毫無缺乏」.
開頭他說他是雅各.
是耶穌的兄弟.
他信了耶穌是基督.

$^{161}$是這位救世主.
在十字架上死亡和墳墓裡復活.
這位耶穌的弟弟.
他成為了耶路撒冷教會的領袖之一.
他很有名的.
不單是聖經有提他.
在聖經以外的文獻也有提到這個James.
很多時候人們會叫他James the Just.
就是二人雅各.
因為他有很多講道和教導.
包括雅各書.
都是去斥責一些權貴.
人們貪心或者說話非常直接.
好像先知那種宣告.
所以很有那種正氣.
很有那種正義感.
所以很多時候叫他James the Just.
雅各在第一代教會當領袖.
所以這封書特別寫給猶太背景的基督徒.
他們遇見很多苦難.
遇見很多情勢他們不能控制.
鼓勵他們在基督的生命裡堅持.
在基督的生命裡成長.
有一點像我們今天這個時候.
很多時候我們都遇見一些逆勢.
我聽到很多弟兄姊妹很擔心.
生活裡面的前途.
無論是去雨流.
或者在讀書的擔心學校.
在工作的擔心失業.
有一點像我們煮東西.
有很多時候材料是很生的.
還沒有預備好.
試煉來到的時候就好像來預備我們的生命.
有時候不是很好受.
但是使他更加成熟.
其實這裡也是這麼說的.
這裡的雅各書第一章第二節.
他說你們遭遇各種試煉的時候.
你們要捍衛喜樂.

$^{201}$這裡的原文很有意思.
遭遇或者和合本.
也就是落在百般試煉當中.
都要以為大喜樂.
這個意思就是說.
這些試煉來到我們中間的時候.
是一種我們感覺到好像很被動的時候.
我們好像被大風大雨.
很多的風暴淋到我們人生的時候.
我們完全朝霞不離.
就是這個意思.
他說落在百般試煉當中.
要以為大喜樂.
什麼叫喜樂呢?.
我很喜歡一個英文的翻譯.
MSG的翻譯.
MSG不是味精.
是Message.
Message的簡稱.
一個現代的英文的譯本.
他說什麼呢?.
他說要以為大喜樂.
其實是傳言的恩典.
傳言的因詞.
Share gift.
意思就是說.
這些試煉來到我們中間的時候.
我們需要去想清楚最後的結果.
是一個Share gift.
一個完全的.
上帝給我們的恩賜.
恩賜是什麼呢?.
就是在我們生命裡面能夠被建立.
這裡的動詞.
認為捍衛喜樂.
或者英文是Consider.
我們需要去鑽研.
用另一個的perspective(觀點).
去看那件事情.
特別是當我們很失控地.

$^{241}$遇見很多的危難的時候.
第三節這裡說什麼呢?.
他說因為知道你們的信心.
經過考驗.
就產生忍耐.
但是忍耐要堅持到底.
使你們可以完全毫無缺乏.
這裡說了一個過程.
好像我們主宋有一個很長很長的過程.
就是由我們的信心受到考驗.
然後就說會生出.
產生忍耐.
然後這個Patience(耐心).
就會要我們去堅持到底.
一直成為一個完成的過程.
使我們的生命成熟.
以及無缺陷.
以前的翻譯本.
和學本就譯作完全.
完全當然不是說我們整個人.
我們無罪.
變成聖人.
完全的意思其實就是.
預備妥當.
就是好像一頓飯.
你由一個原料.
一些原始的材料.
然後經過一個過程.
你之後做起一件事.
是色香味俱全.
這是一個完全的意思.
[音樂].
按照這個食譜.
將這些香丁放進去.
大概有50粒.
插進去.
以致你焗的時候.
味道就可以走進去.
我真的沒做過.
不過儘管跟著食譜去做.

$^{281}$其實有些新的事物要去嘗試.
其實都是一種試煉.
我們落在一些疫情當中.
其實這些真的從來沒有人會經歷過.
整個世界受到影響.
不妨我們去想一想.
到底我們有什麼可以學習呢.
當然我們希望這個逆勢很快過去.
我們都是這樣祈禱求上帝恩典.
但是既然我們在這個所謂百般試煉當中.
百般的考驗當中.
我們就要想.
到底我們有什麼人生的智慧.
在我們的生命當中.
在逆勢當中我們能夠成長.
然後我們將這些香丁插進去之後.
總共有50粒.
插進去火腿.
然後我們就可以將焗爐開了.
325度.
現在大概300度.
我就想看一看亞哥書第一章.
十二節這裡怎麼說呢.
他說能忍受試探的人士有福地.
因為他經過考驗之後.
必得著生命的冠冕.
這冠冕就是主應許給愛他之人的.
人被試探不可說.
我被神試探.
因為神不能被惡試探.
他也不試探任何人.
每一個人受試探.
都是被自己的私欲所勾引誘惑的.
私欲懷了胎.
就生出罪來.
罪長成了.
就產生死亡.
剛才說過亞哥書很像智慧文學.
智慧文學經常警告我們.
現在有一個選擇給你.

$^{321}$要麼你就選生命.
要麼你就選死亡.
這裡第一章一開始就說.
在我們這個逆勢裡面.
很多的困難,試探,試煉裡面.
我們就需要做一個選擇.
這裡也說受到試探的時候.
到底怎麼辦呢.
他在這裡警告我們.
原來我們受試探的時候.
很多時候我們會賴上帝.
我們會覺得上帝在作弄我們.
他在這裡說.
如果我們在試煉,試探當中.
我們跌倒.
我們stumble(倒下).
我們不可以以為上帝在作弄我們.
為什麼呢.
因為他說最主要的原因.
原來是我們自己的私欲.
讓我們跌倒.
或者讓我們離開主.
這裡第一章.
試煉和試探兩個字.
雖然翻譯是不同的.
但原文其實是一樣的.
似乎是告訴我們.
每當我們遇到考驗逆境的時候.
其實可以有兩個結果.
第一個結果.
就是一個會幫助我們成長.
使我們信心能夠更加堅固.
這個結果.
另一個結果.
可能會使我們離開主.
甚至使我們生命破產.
我們信心會破產.
到底你選擇哪一個呢.
這個就是智慧的所在.
每每聖經都說.

$^{361}$人是自己的責任.
我們跌倒.
在試探當中.
我們離開主.
或者我們信心破產.
是因為我們自己有的私欲.
其實阿當夏娃.
一個很典型的每個人的故事.
其實都是因為我們自己眼目的情慾.
我們的驕傲等等.
是我們自己的私欲引致我們離開主.
這裡說了一個比喻很有意思.
私欲會懷胎.
懷胎十月就是這個意思.
就是我們私裡面的desire.
我們的evil desire.
懷孕.
然後慢慢慢慢會長成.
長到最後的時候.
就會有死亡.
雅各作為一個智者.
就說你要想清楚.
在這個逆勢當中.
我們就是一個兩者的抉擇.
如果選擇生命的時候.
我們就明白上帝原來他美善的心意.
是去建造我們的生命.
如果我們選擇死亡的時候.
其實我們就誤會了上帝.
以為他作弄我們.
但上帝是不會這樣.
所以你選擇什麼.
原來要reheat的時候.
我們要去做上面的醬.
看看怎麼做.
死火真的沒做過.
做醬我就需要一個measuring cup.
在哪裡呢.
我在廚房是找不到東西的.
原來要先放進去.

$^{401}$我先放進去.
要煮2小時20分鐘.
timer已經set了.
現在我們可以準備甜醬.
四分三的honey.
要很大力.
不如先打開這個.
倒進去.
快點弄出來.
哇 斷了.
快點斷了 慘了.
死火弄成這樣.
慘了.
不如剪開.
剪不開.
先把筷子拿出來.
看到嗎 筷子斷了.
怎麼剪都剪不出來.
出來了.
終於剩下半個筷子.
這些筷子不太耐用.
我都很有信心這個可以的.
一半的honey.
弄錯了.
只需要一半.
半個cup的brown sugar.
這個有.
quarter cup of orange juice.
我有橙.
等一下可以擠點橙汁.
煮東西真的不容易.
沒想過弄到這麼搞笑.
我平時都是煮公仔麵.
煮到這些的時候.
壓力很大.
剛剛好.
現在還有一個多小時.
還有很多時間.
還可以趕賭.
現在最精彩的部分.

$^{441}$就是弄這個橙汁.
我有這個神器.
弄橙汁.
很久沒弄過橙汁.
需要quarter cup.
這樣真的要五個橙.
說起來.
雅各書.
說到試煉試探的時候.
就好像金和銀.
在金礦銀礦裡挖出來的時候.
就需要被試煉.
需要被approve.
然後才可以煉出那些真金真銀.
和我現在榨橙汁都一樣.
想起一句話.
葡萄被壓才釀出新酒.
它才會成為可以有用的東西.
exactly.
我們的生命就是這樣.
如果沒有一些這樣的難處.
沒有一些這樣的難關.
我們的信心永遠都不會成長.
好像小朋友一樣.
在新約裡面.
其他的經文經常會說到.
譬如教目書信就會說到.
是合用的工人.
approved worker.
一個人要比上帝更加重.
使用的話.
更加重用的話.
那個人要been approved.
去試煉過.
好像你去過考試.
有那種qualification.
好.
這個橙汁已經預備好了.
我們把它混合.
火腿吃得多.

$^{481}$我沒試過真的下手去做.
通常都是老婆做.
或者別人做.
我們吃那個製成品.
但過程一點都不過癮.
試煉試探都是一樣.
不過我們需要with the end in mind.
我們需要有那個.
信心得到考驗.
上帝能夠去拆毀重建.
以及使我們的生命能夠成長.
按照高人的指點.
就是那些甜醬.
就需要一層一層塗上去.
現在雖然在焗.
但是拿出來塗一層.
讓那個cam不會那麼乾.
好.
拿出來.
嘩你看.
好.
這個.
嘩你看.
我覺得ok.
好好吃.
關掉.
現在我們要.
塗上一層.
讓那個cam.
不會那麼乾.
說真的.
這些廚房的東西.
我真的不太熟悉.
現在我當作是塗一塗.
這些.
對我來說是完全新的.
這些經歷.
前幾天上網的時候.
就和一些.
教務的朋友.

$^{521}$都是北美的.
大家就聊天.
在這個疫情當中我們教會.
事奉啊, 牧會啊.
我們有什麼新的適應.
大家都.
不約而同.
這個情況.
因為疫情.
因為禁聚令等等.
我們都會有很多的適應.
和很多新的學習.
其實.
這個就是雅各書的精神.
當我們有試煉.
試探來到的時候.
我們怎樣去consider it all joy.
好, 好像可以了.
塗了一層.
又放回去.
ok.
等這個火鍋的時候.
我們就要讀完雅各書的第一章.
其他經文.
這裡第五節說什麼呢.
你們中間若有人缺少智慧.
就當向那厚似眾人.
而且不斥人的神祈求.
他就必得著.
可是他應該憑著信心來求.
不要有疑惑.
因為疑惑的人.
就好像被風吹蕩翻騰的海浪.
那樣的人.
不必想從主得到什麼.
因為三心兩意的人.
在一切的道路上.
他都是搖擺不定.
剛才說到智慧, 什麼叫做wisdom.
廚房裡面有wisdom.

$^{561}$但是人生裡面的wisdom, 雅各書就說.
在逆勢當中我們需要去做一個判斷.
一個好的outcome.
在試煉裡面能夠信心得到考驗.
我們信心得到成長.
經歷神更多.
還是我們選擇一個不好的outcome.
那個不好的outcome就是最後引致我們.
死亡, 甚至離開上帝.
這些說法都是.
很有智慧文學的味道.
你想起智慧文學的時候.
我們不單止想起.
針綿, 傳道書等等.
都想起很多舊約裡面的故事.
譬如創世記裡面的.
藥室.
藥室被賣到埃及的時候.
他被波提法的.
妻子去引誘他.
在那個盛世探訪當中.
他能夠去得勝.
譬如尼希米.
他在波斯帝國.
強權帝國的底下.
居然做了一個.
很高級的官員.
以致他能夠用那個機會.
來事奉神.
包括被露的猶太人.
很多時候這些.
智慧文學色彩的記載.
都是說到一些神的子民.
他落在一個.
所謂百般試煉當中.
他們用他們自己的身份.
用他們職場上的崗位.
能夠去到.
慰藏神的國.
慰藏神的價值.

$^{601}$他們活出來.
都是有一種態度.
就是做出最好的事.
在一個這樣的難處裡面.
如何能夠去利用它.
如何去到轉變的角度.
去看事情.
我們可以.
怨天尤人.
同樣我們都可以.
在局勢裡面.
那種試煉當中.
可能你覺得很辛苦.
很侷促或很吃力.
但是.
上帝的本性.
上帝自己的美善的本質.
以他.
為考慮.
現在如何面對.
現在的逆勢的態度.
不知道你在.
這個逆勢當中.
你有沒有把握這個抉擇.
我們在上帝面前.
我們寧願更多.
去求上帝.
你在我們心靈當中.
現在你想做些什麼.
當我們還是在.
等候的當中.
要考考我們的 patience.
不如我們看第一章最後的一段.
受獄節.
他說我親愛的弟兄不要寒鋤寥.
各樣美好的賞賜.
各樣完備的恩賜.
都是從上面從縱光之附降下來的.
他本身並沒有改變.
也沒有轉動的影移.

$^{641}$他憑著自己的旨意藉著真道.
生了我們.
使我們作他所做萬物中.
初熟的過子.
這裡說到上帝.
是一個不變的上帝.
他在他身上沒有.
轉動的影移.
平時走路的時候你會發現.
我們是隨著太陽的方向.
隨著外界的世界.
給我們的光線.
我們的影就會經常變來變去.
他沒有任何的改變.
就是說上帝自己的本性.
如果是美善的話.
什麼時候都是美善的.
順境的時候是美善.
逆境的時候.
上帝依然是美善.
在聖經裡面的理解.
萬物是被罪所破壞.
但是因為耶穌基督來到的時候.
他就將整個萬物.
開始展開一個更新的工作.
由耶穌的死而復活.
到他升天作王.
到他再回來的途中.
一直向新天新地邁進.
就是在夏國書.
最後這18節所說.
上帝的旨意.
是要使我們作他所做.
萬物裡面的初熟.
的過子.
整個的創造上帝.
不斷地在更新不斷地在改變.
走向新天新地.
就是神的國會臨到地上.
那個目的地.

$^{681}$而第一個.
要去成就的就是.
在真道裡面在福音裡面.
來生我們出來.
成為了基督徒之後我們就有一個.
這樣的新的生命.
一直仰望著.
一直往這個新天新地來走.
我們不單止有這樣的盼望.
也告訴我們我們的身份.
是初熟的果子.
但是初熟的果子.
當我們想起的時候.
是整個的 harvest(收成)裡面的第一份.
的農作物收成.
他就是 guarantee(保證)整個的.
田地會有.
豐富的田產.
所以那個 first harvest(第一種收成).
就是獻給上帝的.
那個 first harvest(第一種收成)也是一個保證.
原來我們每一個在.
基督裡面我們有一個新的.
生命的人我們不單止.
仰望著上帝將來的新天新地.
同樣我們.
是那種的保證.
保證上帝的工作.
一定會實現.
那麼在試煉試探裡面.
有什麼智慧呢.
我記得這個.
的圖文叫 serenity prayer(平靜的祈禱).
我覺得這個圖文很有智慧.
裡面說什麼呢?.
祂說神啊求你賜給我平靜的心.
去接受我無法改變的事.
賜我勇氣.
去做我能改變的事.
賜我智慧去分辨者兩者的不同.

$^{721}$一天一天的過日子.
一時一刻的去欣賞生命.
接納艱苦.
是通往平安的路.
能夠效法耶穌.
接納這個罪孽深重的世界的本相.
而非自己所期望的那個世界.
堅信你必定會把萬事變為美好.
若我順服你的旨意.
好叫我此生能喜樂地生活.
並在永世當中得享與你無關的生活.
你是我的天使.
你是永世當中得享與你永遠同在的大喜樂.
好了,我想應該搞定了.
好像OK,已經聞到很香.
Let's see.
Wow.
先探探熱.
Let's see.
110度,正確.
120度.
Oh, it's ready.
好了,上去.
Yeah.
哇,好熱.
剛剛好,130度.
Yeah.
It's good.
先放上來.
兩分鐘.
先放上來.
兩塊.
Oh.
Hey, I think it tastes good.
嗯.
Very good.
吃完東西當然要洗碗.
看完成會網台要做什麼呢.
有兩樣東西你可以做.
第一樣就是電子封面.

$^{761}$下面的description box有不同的電子封面方法.
你支持我們的工作,我們可以做更多節目.
更多好的訊息.
祝福更多人,好不好.
第二樣就是like, share, comment,訂閱.
按個subscribe和叮叮.
下一次你就知道我們煮什麼菜.
現在煮完又講完.
我想有個問題問你.
這個星期你可以思考一下.
最近在你心中有什麼考驗.
而上帝如何給你智慧力量.
來面對呢.
記住這個星期來思考一下.
希望下一次能夠見到你.
我們邊煮邊講.
智慧廚房.
最後希望這首歌成為我們的祈禱.
就是塗灶在你手.
(主 正深受眾安排).
(塗灶在我的一生).
(必給指引在不安穩的時代).
(主 願我在你手).
(能改變跟著你).
(塗灶著我).
(都欣著你的光).
(主 正深受眾安排).
(塗灶著我的一生).
(必給指引在不安穩的時代).
(主 願我在你手).
(能改變跟著你).
(祈求願意將生命再交上).
(就在你手中).
(一分一刻都安排在你手).
(你要叫我跟主前往去).
(你要永遠牽著我的手).
(我是你手工).
(你說我正是你手足工作).
(要我靠你在悠暗地方看).
(必得著你的恩惠).

$^{801}$(主 正深受眾安排).
(塗灶著我的一生).
(必給指引在不安穩的時代).
(主 願我在你手).
(能改變跟著你).
(祈求願意將生命再交上).
(就在你手中).
(一分一刻都安排在你手).
(你要叫我跟主前往去).
(你要永遠牽著我的心).
(我是你手工).
(你說我正是你手足工作).
(要我靠你在悠暗地方看).
(必得著你的恩惠).
(就在你手中).
(一分一刻都安排在你手).
(你要叫我跟主前往去).
(你要永遠牽著我的心).
(我是你手工).
(你說我正是你手足工作).
(要我靠你在悠暗地方看).
(必得著你的恩惠).
(主 正深受眾安排).
(塗灶著我的一生).
(必給指引在不安穩的時代).
(主 願我在你手能改變跟著你).
(祈求願意將生命再交上).
(能隨著你 正深被你掃蕩).
(祈求願意將生命再交上).
\newpage



\section{}
\label{sec:_qut0JfQmyc}
\textbf{狼人殺(聖經版) round 1 精華版}
\newline
\newline
連結: \href{https://youtube.com/watch?v=_qut0JfQmyc}{\texttt{https://youtube.com/watch?v=\_qut0JfQmyc}} ~~~~ 語音日期: 2020-04-22
\newline
\newline
\hyperref[sec:HEMq3Mncf_A]{\small{< < < PREV SERMON < < <}}
~
\hyperref[sec:index_chronic]{\small{[返順時目]}}
~
\hyperref[sec:index_scriptual]{\small{[返順卷目]}}
~
\hyperref[sec:CDVnSLlj71Y]{\small{> > > NEXT SERMON > > >}}
\newline
\newline
$^{1}$(魔鬼叫聲).
起初神創造天地.
有日頭,有夜晚.
有光明,有黑暗.
魔鬼就是這樣.
混入了在人的中間.
將人一個,又一個地打死.
在這個十二人的模擬世界.
將會分成三組.
在壞人的陣營有四隻魔鬼.
目標就是要混淆師兄.
和要殺光好人陣營其中一組的所有人.
而其中的首領撒旦.
如果投票被殺的話,還可以帶多一個人離開.
好人的陣營有兩組.
第一組是四個的亞當和夏娃.
即是沒有特別技能的平民.
好人的另一組是四個有特別技能的光明使者.
第一個是先知以才雅.
每晚可以用意象查一個人的身份.
第二個是耶穌.
他夜晚可以叫死人復活.
亦都可以審判殺死一人.
但每個技能只可以用一次.
第三個是光明使者所樂無望.
早上討論的時間,他可以公開查證一個人的身份.
對方是壞人,就壞人死.
是好人的,所樂無望就要以死謝罪.
第四個是參孫.
當他被殺的時候,可以找多一個人和他同歸於盡.
當三組其中一組全軍覆沒的時候,遊戲就會結束.
在這場光明和黑暗之戰,究竟誰旗開得勝,誰旗差一著呢?.
大家是不是都收到自己的角色了?.
第一晚了.
第一晚的時候,大家請閉上眼.
四個魔鬼,開眼.
現在麻煩把撒旦首領都示意.
大家知道是誰了.
你們今晚想殺誰呢?.
是吧?每個人都選這個了.

$^{41}$現在魔鬼可以合眼.
然後耶穌請開眼.
今晚死的人是他.
你有復活的技能,你用不用呢?.
你也有最後審判的技能,你用不用呢?.
好,那耶穌麻煩合眼.
現在到先知起身了.
好,先知你今晚想查哪一個身份呢?.
(魔鬼的身份是誰?).
好,現在這個環節只給觀眾知道.
鎖羅門王麻煩揮手.
好,知道了,你可以合眼.
現在參孫麻煩揮手.
你不揮手的嗎?.
可以合眼了.
天亮了,大家可以起身了.
今晚死的人是1號.
1號因為是第一晚死,所以你可以說遺言.
之後夜晚死的人就不能再說遺言了.
你們今晚這麼快叫我回去做夜展演.
因為你得罪人不可能多.
你指票皇后經常得罪人.
大家說得多又被人壞壞.
說話又被人壞壞.
我真的沒有想過.
你死了沒有說過這麼多話.
第一個可以說的.
大家好人加油.
好人加油,我關燈了,再見.
我們現在可以輪流發言了.
我覺得Emily死得很慘,我會幫你報仇的.
我會幫你報仇的,OK?.
11號.
今天Barry這麼安靜,平時有很多話說.
10號,今天Jason這麼多話說.
這個遊戲大家多說點話,否則會令我很奇怪.
我現在暫時看不到任何話.
下一位.
Emily死得慘.
一開始就被人這樣殺死.

$^{81}$我會幫你找到你的仇人.
下一位是不是我?.
是.
過.
下一位就是我了.
我是一個好人.
沒有很多的資訊給大家.
我是5號.
多謝你的意見,Presley.
完,過.
4號.
Emily我會幫你尋尋不捨.
3號.
我也是一個好人.
但我覺得有疑惑就是今天Jason很多話說.
是不是有先知?先知可以查一下.
這兩個先知.
OK,過.
到你了,Moshe.
我一定會支持公道的,我告訴你.
那個殺手一定會橫屍街頭.
好,說完,現在有一個自由討論的時間.
我覺得Hale的分析都非常合理.
合理得來就有些掩飾.
剛好有三四個人說了Emily死得很慘.
為何這四個人會剛好就是壞人呢?.
對,會不會無端端Amanda第一次來玩就說要為她報仇呢?.
但都可以不要騙Persi的.
因為Persi平時不會想這些這麼黑暗的事.
不知為何她今天特別勇武.
她太正經得太厲害了.
我覺得要小心Persi這個人.
我也覺得要小心Persi.
不要,不要,不要殺死我.
我是壞人,不是,我是好人.
哈哈哈哈哈哈.
現在就是投票的時候了.
現在請舉手.
我看到6號有三票,所以她是最多票數的.
所以6號今天就要表達.

$^{121}$6號因為你早上死,所以你可以說遺言.
我覺得Tim應該是一個好人.
但我很肯定Amanda一定是壞人.
所以我可以說遺言,你可以關燈.
剛才6號的全部都是壞人.
大家閉上眼睛不要理他.
魔鬼可以開眼.
魔鬼你今晚想殺的是誰呢?.
OK,看到了.
Jesu可以開眼了.
剛剛死的是他.
你選不選擇用你的復活技能呢?.
你也有最後審判的技能,你選不選擇用呢?.
你可以閉上眼睛了.
現在到仙姿起身了.
仙姿你今晚想查誰的身份呢?.
OK,她的身份是.
仙姿閉上眼睛了.
那就天亮了.
有一個人要死,那個人是8號.
你是夜晚死,所以你不能說遺言,現在可以關燈了.
討論的是5號先,然後順著下去.
跟大家十不相瞞,其實我是耶穌.
我剛剛殺了其中一個殺手.
那就下去了,6號死了,所以到7號了.
沒有,過了.
下一個,8號死了,9號.
剛剛聽5號說的話,基本上是不合理的.
因為如果他殺了人的話,應該昨晚是有兩個人死的.
我懷疑他是壞人.
5號,他說他是耶穌,為什麼他說他殺了一個殺手呢?.
他為什麼是殺手呢?.
很複雜,下一個.
我也覺得Clinton的東西是有些問題.
好像在掩飾自己.
我很想相信這個Clinton說的話.
然後.
為什麼你不相信?.
很多人都反對Clinton.
那些反對Clinton的人是不是同一黨的人呢?我不知道.

$^{161}$我其實不太清楚是發生什麼事的.
我覺得有可能Clinton平時說話都是這樣,所以我習慣了.
但是呢.
我覺得有可能如果他是耶穌的話,所以死了一個是對的.
因為殺手一定會殺人的.
可能他就救了那個,然後他就再毒死了其他人.
我覺得一起說Clinton是有問題的那些很有嫌疑.
我覺得有功能,有身份的人應該要多做事.
好,就是這樣.
其實我也不太清楚發生什麼事.
不過剛剛Percy在說我是壞人.
我突然在想,如果我是壞人的話.
其實應該魔鬼是不會想殺我.
因為如果他殺我的話,那不就是知道我是好人嗎?.
好,Amanda你也是很有嫌疑的.
OK.
我覺得Harold很有可疑.
為什麼呢?.
他想我們好人投錯票,白天又死一個好人.
晚上又死一個好人,肯定的了.
因為耶穌只能夠救一次.
我們每晚都死人的,OK.
我就知道他是一個好人.
下一個你檢查的是2號,4號,7號和12號.
好,是不是討論完了?討論完就去投票的時間.
大家就躲起來,今晚投死誰呢?.
投完了沒有?.
就有3way tie,4號,10號和12號.
其實我是好人,我知道Jason投我.
但我覺得有點古怪.
因為他之前也沒有提及過任何關於我的事情.
但他現在就馬上輪到我.
所以我覺得Jason是挺有可疑的.
10號.
唉,投了我死就糟糕了.
我是有功能的好人,OK.
12號.
我覺得那些投了票,例如11號很可疑.
Jason和Mark Sir和Harold.
這三個人都很可疑.

$^{201}$大家相信我是好人.
好,現在就再投票,可以再閉上眼.
我看到4,4,4,12,12和10.
所以4號就是要死了.
4號你可以講位置.
總之我覺得投我的都是有嫌疑的.
即是3號和5號.
而且我也是好人,所以對不起各位.
我也要承認了,祝你好運.
好人萬歲,再見.
是不是每個人都閉上眼了?.
等等,等等.
好,現在魔鬼可以起身了.
你今晚想殺誰呢?.
是不是這個?.
好,知道了.
耶穌,請你開眼.
剛剛死的人是他.
你有兩個技能,你選不選擇用死裡復活呢?.
是,OK.
那你選不選擇用第二個功能,最後審判呢?.
OK,那耶穌可以閉上眼了.
先知可以起身.
你今晚想查誰呢?.
OK,他的身份是….
先知閉上眼.
天亮了.
昨晚死的人是12號.
那12號你選不選擇發動你的功能呢?.
選擇,我可以拉一個人死嗎?.
你可以帶一個人走.
暫時.
我覺得這個Jackson很可疑.
因為他每次投票都要有頭人.
然後很想大家快點死.
所以我要拉他一起死.
11號你選不選擇發動你的技能呢?.
有什麼技能呢?.
你為什麼拉我一起死呢?.
12號拉了11號一起死.

$^{241}$所以11號你選不選擇發動你的技能呢?.
算了,我沒有技能.
你沒有技能,所以你現在就關燈了.
因為你們兩個都是晚上死的.
所以你們都沒有遺言.
9號現在可以順序發言.
因為本來以為Jackson是先知.
因為他投了4號.
無可奈何地他就死了.
不如先猜猜.
我覺得Jackson應該是最壞的壞人.
只不過他不發動他的技能.
Katrina又可以發動技能.
她又可以發動技能.
一是魔鬼大佬.
另一個是參選.
Very快點解釋.
什麼解釋呢.
我就是楚內門王.
我不想出技能.
不可以亂出.
我覺得有身份的人是時候發動.
要不然贏就輸了.
有身份的人發動了.
你知道他是.
然後他死,我們贏.
我就是楚內門王.
願意嗎?.
好人,要不要出技能.
當然要出.
現在這個時候.
不如2號.
好,2號.
是一個壞人.
所以不需要死.
但是2號因為你早上死.
所以你可以說是遺言.
快點說,你同黨是誰.
我真的不知道發生什麼事.
好,多謝你的遺言.

$^{281}$因為楚內門用了他的技能.
所以今晚是不能投票.
現在可以立即去下一晚.
魔鬼,你今晚選擇殺誰呢.
魔鬼可以合眼.
耶穌請開眼.
剛剛死的是這個.
你可以令人死裡覆沒.
你用不用你的技能.
你也可以用最後審判,你用不用.
耶穌請閉上眼.
先知,起身.
你今晚想查誰的身份呢.
他的身份是.
先知閉上眼.
就到天光了.
昨晚是有一個人要死.
那個人就是10號.
10號因為是夜晚死.
所以你是沒有遺言的.
7號說話才是.
我相信就通過了.
如果這個階段只剩下4個人.
應該只剩下一個殺手.
我覺得5號的嫌疑比較大.
如果照朝神所說.
Krinson其實是一個高玩.
我相信還有一個有身份的好人.
如果Krinson說的是耶穌.
現在就合理了.
應該沒有先知了.
因為先知應該死了.
應該一早死了.
我覺得早就沒有先知了.
我們投票吧.
我看到9975.
所以9號是要死的.
所以你可以說是遺言.
我是好人.
你殺錯人了.

$^{321}$這樣一定輸了.
暫時是完結了.
因為只剩下一個殺手.
殺其中一個人.
都會贏.
所以殺手贏了.
魔鬼贏了.
大家可以回來了.
第一局.
第一局玩的遊戲.
竟然是魔鬼贏了.
耶.
我是先知.
你第一局殺死了我.
一定沒有先知.
沒有人在說.
是不是Krinson.
是不是7號.
唉.
都說這個遊戲.
不說話的人都很有意見.
好.
第一局.
我們現在跳去第二局.
交給Calvin.
謝謝.
\newpage



\section{雅各書}
\label{sec:CDVnSLlj71Y}
\textbf{2020/4/26 智慧廚房 vlog 信息 \#2 (雅各書 2) - 陳崇基牧師}
\newline
\newline
連結: \href{https://youtube.com/watch?v=CDVnSLlj71Y}{\texttt{https://youtube.com/watch?v=CDVnSLlj71Y}} ~~~~ 語音日期: 2020-04-27
\newline
\newline
\hyperref[sec:_qut0JfQmyc]{\small{< < < PREV SERMON < < <}}
~
\hyperref[sec:index_chronic]{\small{[返順時目]}}
~
\hyperref[sec:index_scriptual]{\small{[返順卷目]}}
~
\hyperref[sec:o_g_uZyyIHo]{\small{> > > NEXT SERMON > > >}}
\newline
\newline
$^{1}$各位好,又到智慧廚房邊煮邊講.
今次我們繼續煮,繼續講亞國書到第二章.
今天想煮什麼呢?就是港式豬腸粉.
就是自製豬腸粉,大家可能吃過了,是不是?.
回味無窮的.
按照我在YouTube和網上的資料搜集.
就有不同的做法和不同的配料.
大致上那個粉漿很重要.
就需要粘米粉,粟粉和燉麵粉.
還有蔥,蝦米等等.
還有其他醬料.
這些東西大部分家裡都有的.
不過好像不太夠.
所以你要跟我們去超級市場買一些東西.
現在來到市場,很多人.
木薯粉,鎮粉,木薯粉.
看看有沒有.
花工粉.
哇.
有,粘米粉一包.
這個粘米粉.
我要燉粉和米粉.
番薯粉.
這個是什麼?.
生粉.
或者西米粉,是不是一樣的?.
真的不懂這些東西.
原來白粉也有這麼多種.
每種都白色,但每種都不一樣.
這個還是太白粉.
有沒有一些不太白粉的?.
這些很像西米粉.
甜豉油,不可以沒有的.
甜豉油一定要有的.
煲仔飯豉油.
下次可能可以做煲仔飯.
真的少見多怪,有臭食粉.
臭食粉是什麼東西?.
臭食粉,真的.
看到也不想要.

$^{41}$放回去.
真的見識多了.
原來有這麼多種粉,現在才知道.
但找不到燉粉.
一定要有芝麻醬.
這個芝麻醬就是.
看清楚哪裡出的.
台中市的,台灣產品.
先試一下,芝麻醬.
這些粉我都已經買齊了.
醬料也買齊了.
沒有的燉粉,家裡有一點.
不然可能要問人借.
不過應該夠用的.
現在找齊東西可以走了.
不過我在超級市場找東西很慢.
你叫home depot找,我想我會比較厲害.
好像廚房一樣,超級市場我經常找不到東西.
怎樣也好,找完了.
現在找不到的是什麼?.
找不到跟我一起來的老婆.
我現在在找她,她在哪裡?.
(找不到).
找不到她,不過找到芝麻.
這個很重要,借不到的.
好,芝麻.
其他東西家裡都有,燉粉應該夠用的.
好,東西都買齊了.
我們現在可以開工了.
怎樣去做豬腸粉呢?.
第一樣東西當然是配料.
簡單的配料,首先做蝦米腸.
當然要蝦米,已經在暖水裡浸泡了.
然後有蔥.
我們先把蝦米切了.
做這個腸粉有一個秘訣.
雖然我從來沒有做過蔥,只是第一次.
有一樣東西,我觀察到的.
不同的食譜,不同的做法.
就是那個combination.

$^{81}$那個配搭,那個配合是很重要的.
如果有些東西多了,有些東西少了.
那個配搭不對.
就很容易做到那個腸粉不好吃.
無論是粉漿,還是配料等等.
所以先做蝦米,這是最簡單的步驟.
首先做蝦米腸.
當然我們不喜歡做蝦米腸.
可以做其他配料.
不過我們等一下有時間,可能試一試.
譬如叉燒腸,蝦腸等等.
這些是隨後的畫.
我們先成功地研製了.
適合的配料,適合的combination.
我們就可以把它弄得複雜一點.
蝦米很小,但要它更加小.
現在浸了,我不知道你聞不聞到.
已經有些蝦米的很香的香味.
說到不同的配搭.
就正是《雅各書》裡面所說的.
《雅各書》裡面說的是我們基督徒的生命.
信了耶穌的人.
我們的生命需要有一個很好的配搭.
而這個生命沒有一個很好的配搭.
只是偏門得一樣.
缺少了其他,其實就不是一個完備的生命.
這是整個《雅各書》會提到的.
所以你看《雅各書》,它整天說.
上帝的律法,上帝的道是完備的.
我們的生命照樣是需要完備的.
我們切了這些蝦米.
然後切一點蔥.
蔥其實暫時切一點就夠了.
就是在加料的時候,蒸的時候.
我們才能夠盡量去到.
加多一點.
蔥就不怕多,因為切多了.
我們就可以拿去做其他東西.
蔥就切好了.
我們就把它放在碗裡.

$^{121}$蔥就切好了.
不如也切這些.
快點切掉.
因為我們主要不是煮東西.
當然煮東西也很重要.
我也很想吃豬腸粉.
但是主要是邊煮邊說.
在說《雅各書》.
OK.
現在我們把這些蝦米切粒.
蔥就切粒了.
第二部分是什麼呢.
就是做豬腸粉的皮.
你知道豬腸粉的皮有多重要.
特別是港式腸粉一定要滑和彈牙.
按照我的研究來說.
有三樣東西需要.
第一樣就是主角.
就是粘米粉.
按照我有限的認識.
粘米粉沒有糯米粉那麼黏.
那些煮點心,包東西的人.
弟兄姊妹應該很清楚.
另外我們需要加什麼呢.
就是加Cornstarch.
就是平時我們做湖水的粟粉.
聽說這個粟粉.
就是令到粘米粉不會太黏.
比較不會那麼硬.
所以加一點點.
然後另外一樣就叫做澄粉.
澄粉是什麼呢.
澄粉就是令到腸粉比較彈牙一點.
所以有三樣東西.
粘米粉我們需要60克.
60克.
然後有一點點的Cornstarch.
我們就要多20克.
做腸粉真的很不簡單.
你吃就以為很簡單.

$^{161}$你買這些東西的時候.
你就知道原來有很多東西要準備.
才可以做得到.
現在都未必做得到.
看看製成品是怎樣.
現在我們需要澄粉.
澄粉是做什麼的呢.
彈牙.
澄粉.
就要10克.
問題是越多就越彈牙.
你覺得是不是呢.
我覺得不是.
最重要的元素就是.
幾樣東西全部加起來.
是一定需要有一個好的配搭.
有一個好的配合.
鹽就要少少.
然後就要水.
水就要260毫升.
以前讀化學的時候.
要觀察那條線.
是底部.
你看到底部.
260毫升.
隨便吧.
攪拌均勻.
這個時候我們攪拌的時候.
不如一邊攪拌一邊看聖經.
上次我們講了第一章.
現在到第一章尾.
記得第一章上次講到什麼呢.
就是講到一個基督徒.
我們遇見試煉試探的時候.
我們需要知道上帝的美意.
就是使我們的心靈能夠成長.
所以有考驗來到的時候.
我們就需要做出最好的表現.
在逆境裡面我們如何鍛鍊我們的信心.
鍛鍊我們的性情.

$^{201}$以致我們更加像神自己.
神的心靈能夠在我們的心靈當中表揚出來.
第一章尾的時候說什麼呢.
我親愛的弟兄,你要知道人人應該快快的聽慢慢的講.
慢一點動腦.
因為人的怒氣並不能夠承傳神的意.
所以你們應當擺脫一切污穢和所有的邪惡.
以溫柔的心領受神所栽種的道.
這道能救你們的靈魂.
你們應該做行道的人.
不要單單去聽道.
自己欺騙自己.
因為人指著聽道的人不去行道.
他就好像一個人對著鏡子.
看見自己本來的面貌.
過後走開之後馬上就忘記自己的樣子.
唯有詳細察看那使人自由的傳遞的律法.
事實上遵守的人.
他不是聽了就忘記.
而是實行出來.
就必因自己所作的夢復.
這裡講到說.
上帝的道首先是種了在我們的生命當中.
記得上次我說過.
我們是初實之果.
上帝在真道裡面生了我們出來.
然後這裡雅谷書說什麼呢.
他說道是栽種在我們的生命當中.
受神栽種的道.
這個道能夠救我們的生命.
並且這個道是要聽.
更重要的是要行出來.
首先第一件事.
上帝的道進入了我們的生命當中的時候.
其實會被我們的性情改變.
這裡提出一個例子.
很多時候我們有怒氣.
怒氣的根源是什麼呢.
很多時候我們很自我中心.
我們覺得整個世界的中心是我們.

$^{241}$有一樣東西不順利.
有一樣東西不work out.
或者不是我們心目中的計劃進行的時候.
我們就會發怒.
他提到上帝的道進入了我們的生命當中的時候.
怒氣慢慢就應該要退到最後.
我很喜歡英文的翻譯叫Message.
上次說過了叫MSG.
不是美精.
他說要lead with your ear.
就是用你的耳朵帶領自己的生命.
你的性情.
以上帝的道聽他的道.
所以那個生命的道進入了我們生命當中.
我們就應該有一種性情上慢慢的改變.
基本上我們會改變了.
我們應該成為一個謙虛.
來聽上帝的話的人.
而不是聽我們自己的性情.
怒氣或者其他的感覺.
所以這裡提到上帝的道在我們生命當中.
在性情上有的效果.
有一個問題要問大家.
就是過去一個星期.
上帝的道有沒有使你更加謙虛.
使你不會越來越自我中心.
有沒有將你去自我化.
在你生命當中你覺得.
我自己的想法.
我自己的行徑.
不一定是最重要.
最重要的是我聽上帝的道.
好,我想我已經搞好運救了.
搞了很久了.
因為要趕到.
現在我們就要去蒸東西了.
蒸東西我們需要一個蒸的鍋.
我拿了這個打火鍋的鍋.
我們盡量用一些家居裡面找到的東西.
不需要特別買一些特別的工具.

$^{281}$因為有些人買一些蒸腸粉的抽屜.
蒸腸粉的器皿.
那些東西不是每個人都有的.
我們盡量邊煮邊說.
這個智慧廚房就是用智慧.
怎樣用現在現成有的東西.
來製造出我們好的.
能吃到的,滑的豬腸粉.
蒸腸粉皮.
我首先試一下用玻璃的鍋.
平的.
倒進去.
試一下這樣行不行.
這個剛剛好.
好,火就開了.
等水滾.
我們就試一下用不同的方法.
看看怎樣的工具可以製造出好的豬腸粉的皮.
怎樣蒸得到,漂亮.
要試的.
要不斷不斷去嘗試.
這是第一個方法.
其實那些下廚的朋友就知道了.
煮了這麼多年.
你已經試了很多不同的東西.
失敗過不同的東西.
無數次失敗過.
然後又再試驗過.
然後又再加強版.
然後又用其他的方法.
這些是經驗.
經驗是怎樣來的呢.
就是不斷失敗.
不斷成功.
然後去到累積出來的.
我記得有一次去一個訓練.
那裡有句說話挺好的.
他說最好的教師不是經驗.
Best teacher is not experience.
他說最好的教師是什麼呢.

$^{321}$就是評估過的經驗.
Best teacher.
是Evaluate experience.
就是你失敗過之後你再想想.
不用同一個方法.
用其他的方法.
你就可以有一個生命裡面最好的教師.
其實這個也是《雅各書》裡面所說的.
我們要有傳遞的信仰.
傳遞的基督徒生活.
就是不斷去實踐出來.
所以不要單單做聽道的人.
要做行道的人.
水滾了.
我先試試塗點油上去.
少少.
好熱啊.
好 這樣吧.
各位.
我們就塗了少少油.
現在就應該準備好了.
一手拿著.
然後一邊倒下去.
其實煮東西不是最難的.
煮東西和拍攝才是最難的.
你也想不到.
我預備材料的時間.
是比我預備鏡頭的時間少.
現在滾了.
我們就將它加進去.
然後放在一邊.
我們就加蝦米進去.
這個要放平.
OK 整回平均的.
OK 行.
先放蝦米進去.
然後有些蔥.
因為少少夠了.
看看行不行.
放進去.

$^{361}$爭取時間.
你看我多流暴.
又要拍攝又要煮東西.
這個其實也是一個學習.
大概要蒸多久呢.
我想三分鐘四分鐘就行.
先看看.
現在我不可以分心.
我要看清楚.
因為這是第一次.
然後才可以邊煮邊說.
其實現在我覺得有少少時間.
我們不如看看這裡.
亞國書是一個很實用的書.
所以他經常不放過我們.
不像保羅的書信.
比如羅馬書.
他講了十一章.
都是比較教義性或神學性的內容.
那些是信念的東西.
就是Right Belief的東西.
那些都是重要的.
但他講幾章.
解釋了很多章.
然後才進入實用的東西.
比如以弗所書你記得.
頭三章都是比較教義性的.
教會神學的東西.
然後四五六章是講家庭生活.
一個基督徒的生活.
在世裡面應該怎樣去行動.
生活行為怎樣.
亞國書是很不同的.
亞國書的語氣.
經常都是講實用的東西.
他不肯放過我們.
只是頭腦上的知識就夠了.
他一味逼我們想想怎樣去行動.
所以二十六字那裡講什麼呢.
他說如果有人自為虔誠.

$^{401}$自為自己是屬靈的話.
如果不約束自己的舌頭.
反而自己欺騙自己.
這個人的虔誠就沒有用了.
其實是假的.
.
所以是非常實用.
這裡好像已經搞定了.
現在有點厚有點薄.
因為這個東西斜了.
現在先關掉.
我們把它拿出來.
然後看看行不行.
現在準備好了.
先把它拿出來.
然後把它刮出來.
其實一隻手就搞定了.
弄了第一個皮.
這個皮不知道行不行.
我們把它.
好像黏住了.
有點硬.
好像油不夠.
好.
可以上碟了.
這裡一條.
搞定.
OK.
我不知道好不好吃.
不過我覺得很OK.
彈一下彈一下.
先弄多一條.
等的時候.
三分鐘.
不如看看亞國書第二章.
正如我所說亞國書是一些.
很實用.
逼我們要去行道.
不要只是坐在這裡聽.
頭腦上有很多很多知識.

$^{441}$認識很多聖經.
但是我們行道的時候就沒有力.
他說神給我們虔誠的生命.
spiritual life(靈性).
spirituality(靈性).
其實需要什麼呢?.
首先不要被世俗的東西所玷污.
不要被世俗的環境.
以及價值來玷污.
另外他就說到.
社會公義.
說到什麼呢?.
是照顧患難中的人.
那些孤兒寡婦.
在社會當中.
在群體裡面.
他沒有什麼勢力.
能力是很有限的.
那群人.
我們有沒有下手下腳.
來照顧他們呢?.
你觀察一下這個世界.
大把的人.
現在就算是21世紀.
就算我們生活在一些.
現代化的社會當中.
有很多人是很需要.
人去照顧的.
喂!.
好像還可以.
這是第一條.
這是第二條.
好像這個呢.
是.
但怎樣可以令到它不會爛呢?.
剛才沒有爛.
現在為何會爛呢?.
這個是不行了.
為什麼呢?.
剛才那個OK的.

$^{481}$你看這個爛了.
這樣都可以.
這些是不能播出的.
好,我們不要緊.
再接再厲.
好,我們再來一次.
到第二章的時候.
他給了一個非常實際的例子.
那個例子是什麼呢?.
就是有人來到教會的時候.
你怎樣去款待他.
他說什麼呢?.
他說你們既然對我們榮耀的主耶穌基督是有信心的.
就不應該憑外貌來待人.
如果有人手帶金戒指穿華麗的衣服進來會堂的時候.
另外有一個穿衣衫藍柳的窮人進來的時候.
你怎樣去款待他呢?.
他就指責他們去看重那些穿華麗衣服的人.
給他一個好位子坐.
窮人就說不如你站在這裡.
坐在腳椅下面.
這不是你們對人有歧視.
成了心懷惡意的審判官.
怎樣去待人切密呢?.
其實都是我們的靈性的一部分.
我們跟上帝的生命是很有關的.
在疫情當中.
經常有人說笑.
Who cares?.
這個 World Health Organization.
有很多社會的評論都說.
Who?.
WHO.
世衛.
經常說一些東西是沒有幫助疫情的.
只不過是去做一些宣傳而已.
所以Who cares?.
這裡雅各書就說.
Who cares?.
那些無權無勢的人.

$^{521}$到底誰應該去 care for 他們呢?.
雅各書一直想跟猶太背景的基督徒去強調一件事情.
他說.
你們很喜歡聽道.
你們對道可能很認識.
不過你們這樣以外貌待人.
眾父輕貧.
其實就是違反了律法書裡面所說的東西.
第八節說什麼呢?.
他說你們要照著聖經所說.
你們懂得很多.
要愛人如己者華.
要去完成者至尊的律法.
你們就要做對了.
如果你們要按外貌待人的話.
就是犯罪了.
律法就要裁定你們是犯法的人.
因為凡是遵守全部律法的.
只要在一條上面是失足的話.
就是違反所有的律法.
他這裡說到.
上帝的律法是整全的.
是 wholeness(完整).
是整全的.
不能夠把它切斷.
就是你守一點.
然後又不守一點.
你留一點.
然後就輕忽另外一點.
你就會成為一個不願避的一個人.
因為上帝的律法是不可以斬開你的.
照樣我們的生命.
去跟隨律法的話.
就不應該斬開你.
這裡第幾條.
一二三四.
二三四都是失敗的.
四就更加不用說了.
第五就好很多了.
你看看這個.

$^{561}$這個很不錯.
比第一條反而更加好.
好,那你看到這裡好像第五條了.
第五條好很多.
因為做了一些改善.
那麼.
經驗就是這樣來的.
evaluate(審核)了.
我們再學過.
那麼經常都是這樣跌跌碰碰.
那麼我想雅各的意思也是這樣.
雅各不是說我們一步能夠登天.
他講到我們要進到完全的地步.
也不是說我們完美了.
我們不再會犯罪了.
其實就不是這樣的意思.
他真的說.
那個完備的信仰.
其實就是wholeness(完整)的意思.
就是說你經常有整幅.
整幅更大幅的圖畫.
in mind(心中).
那麼他就去到對比了.
魔鬼的信仰.
以及真正的信仰會是怎樣呢.
那麼他就給了兩個例子.
在希伯來舊約聖經裡面.
兩個信心的英雄.
第一個就是叫做亞伯拉罕.
那麼亞伯拉罕個個都認識.
特別是猶太背景的基督徒.
那麼亞伯拉罕就是所謂猶太人的國父.
那麼你看看創世記裡面.
記載亞伯拉罕信心的故事.
由上帝呼召他.
他信主之後.
他有的成長.
當然他不是完美的人.
但是肯定有一樣東西是值得讚賞他的.
就是他的信心.

$^{601}$不是單單頭腦上的信仰.
說是真實的.
他是實行出來的.
第二個例子就是妓女拉罕.
那麼這個女子就是迦南人.
在以色列人眼中.
她是一個滅亡的人.
那麼你記得那個故事.
約聖下記第二章.
就是以色列人快要進入應許之地.
他們首先要進入迦南.
要攻打耶利哥城.
那麼拉罕這個妓女.
就在那裡生活.
那麼她就去迎接探子.
並且投奔於耶和華.
相信耶和華.
兩個例子.
亞伯拉罕和拉罕.
都是有信心並且有行為的.
那麼你看到第一條是OK的.
第一次.
然後第二,二,三,四.
就越來越爛.
然後就改善了那個tray.
make sure它是乾淨的.
make sure油是even的.
如果沒有油真的不行.
好,那麼現在進入第三部分了.
這一部分也是非常重要的.
大家吃過豬腸粉的人.
就知道那個醬汁也很重要.
我們可以買現成的.
這些一枝過你敢記.
買芝麻醬.
如果不行就用花生醬.
那麼另外我們也可以自己做.
配料就是一份水,一份豉油和三份糖.
其實我覺得不需要怎麼量.
然後一份水.

$^{641}$差不多.
三份一糖.
這樣就可以了.
不要吃太多了.
即是糖尿.
攪一攪它.
然後有一件事漏了.
就是芝麻了.
現在慢慢將這些芝麻.
由白芝麻變成金黃色.
好,人生第一次的自製的豬腸粉.
終於面世了.
先試一試.
這個雖然不帥.
不過呢.
也能吃的.
賣相不行.
現在又不是開舖.
沒所謂的.
還可以的.
這次是第一次.
人生第一次.
嘩.
真的不簡單.
那麼.
我開頭的時候不是說過.
港式的豬腸粉.
或者豬腸粉一般來說.
有分三樣東西.
那麼你想想.
其實亞哥斯也是這樣.
有三方面.
那麼我想起就是.
這三方面是神學家將它總結出來的.
一個基督徒的生命.
一個基督徒的生活.
是有三方面.
第一方面是什麼呢.
就是 orthodoxy.
就是正統的信仰.

$^{681}$就是信仰的內容.
當然不能偏差.
要根據聖經.
第二就是.
orthopraxy.
什麼叫 orthopraxy.
你猜也猜到是什麼意思.
就是正統的 living, practice.
就是實踐出來.
那一定要按照我們所認識的.
正確地 practice 出來.
這個就是亞哥斯經常去到.
強調的一件事.
就是你不可以沒有實踐.
你的信仰一定要有實踐.
第三呢.
就是 orthopathy.
什麼叫 orthopathy.
意思就是說.
那個 emotion.
或者我們這麼說 affection.
就是我們的情意.
我們的感覺.
那麼這個意思不是說.
我們每一個基督徒都要某種感覺.
某種.
你當他什麼都好.
好像很 high 或者哭等等.
不是單單說那種.
感覺的表露.
而是說那種的.
價值觀.
我們生命當中的價值觀.
我們生命裡面的感受.
我們生命裡面感覺到.
identify 到.
認同到上帝的感受.
這個叫 orthopathy.
這個就是他說到的.
基督徒的品性.

$^{721}$基督徒的品性應該是怎樣呢.
不可以.
有些不義之路.
我們要慢慢慢慢的.
說話.
說話要有智慧.
我們的生命當中.
第一件事我們要聽上帝的道.
跟隨他的心意.
慢慢地聽.
這件事少了也是不可以的.
因為不可以只有.
我們一些很冷的知識.
然後一些純粹是行為上的東西.
而是我們內心的性情也會改變.
這是一章最後的時候.
亞國書所說的.
所以這三件事.
orthodoxy.
正確的信念.
然後 orthopraxy.
正確的生活行為.
第三是 orthopathy.
我們 identify 認同到上帝所感受到的東西.
我想問你.
在你的基督徒的經驗.
在你的基督徒的旅程當中.
你少了哪一樣呢.
少一樣也不可以.
三樣也需要.
讓我們信聖經裡面所信.
讓我們行聖經裡面所吩咐.
也讓我們去認同基督的心意.
親愛的朋友.
多謝你收看城匯網台.
城匯網台是屬於城匯社區教會.
我們位於美國芝加哥.
我們網台的目的很簡單.
希望在網上做到更多好的節目.
好的資訊.

$^{761}$祝福更多人.
使更多人認識我們的主耶穌基督.
這個時宮希望你能夠和我們同工同行.
下面description box有不同電子奉獻方法.
你看我們的節目背後製作.
軟件硬件都是非常簡單的.
你的電子奉獻絕對能夠幫助我們繼續和發展這個時宮.
最後不要忘記.
like share comment和subscribe.
請你按那個subscribe按鈕和叮叮.
下一次你就知道我們有什麼新的動向.
願上帝祝福你.
\newpage



\section{}
\label{sec:o_g_uZyyIHo}
\textbf{1 十一少被賣記原身家庭(一)}
\newline
\newline
連結: \href{https://youtube.com/watch?v=o-g-uZyyIHo}{\texttt{https://youtube.com/watch?v=o-g-uZyyIHo}} ~~~~ 語音日期: 2020-05-04
\newline
\newline
\hyperref[sec:CDVnSLlj71Y]{\small{< < < PREV SERMON < < <}}
~
\hyperref[sec:index_chronic]{\small{[返順時目]}}
~
\hyperref[sec:index_scriptual]{\small{[返順卷目]}}
~
\hyperref[sec:qk8r6sElCdc]{\small{> > > NEXT SERMON > > >}}
\newline
\newline
$^{1}$有一首詩歌的歌詞說道.
常見到人生像演戲.
如果人的一生真的像一部戲.
我們每一個人都會是戲中的男主角,女主角.
不過主角所扮演的角色,性格,家庭背景和外形等等.
包括劇情的發展和結局.
都不會是做主角可以自編自導.
劇中每一個角色都必須按著劇本一言一動.
都要聽導演的指導.
如果人生真的像世人所說的一部戲.
這樣人的一生由開始到結束.
似乎是很無奈亦毫無意義.
因為人的一生竟然由開始到結局.
都只可以按著劇本.
在劇中的主角既沒有自由選擇自己喜歡的劇本.
亦沒有能力改寫自己的一生.
不過對重生,有新生命的基督徒來說.
基督徒的人生卻充滿著希望和意義.
為什麼我會這樣說呢?.
事實上基督徒重生後.
我們會有新的生命.
但一樣是沒有辦法改變一些事實.
不過我們基督徒的人生的確充滿希望.
因為重生後的基督徒.
生命已經不再一樣.
我們重生後的基督徒有新的生命.
新生命是屬於靈性的生命.
人未信耶穌前,靈性也是死的.
但當人信耶穌後,靈性得到復生.
而且復生後的靈性生命遠超過人的肉體生命.
為什麼我會這樣認為.
靈性的生命遠超過人的肉體生命呢?.
地卿姐妹.
有很多事在人出於母胎前已經定了不能再改.
舉個例子.
我們不能選擇自己出生時的父母是誰.
出生時我們排行第幾.
出生時我們是男生還是女生.
出生時在什麼年代.
出生時家庭背景又是怎樣.

$^{41}$每一個重生後的基督徒.
我們的肉身都是一樣.
所以我們沒有自由去選擇.
也沒有能力去改.
不過,我們重生後的基督徒.
靈性的生命當得到復生後.
因為有聖靈.
在聖靈的帶領和掌管下.
基督徒的新生命就不再一樣.
即使在現實的生活中.
我們依然要面對著不能改變的狀況.
例如,天生有缺陷.
孤兒在單親家庭中長大.
天生患有絕症.
親人遇意外,失望等等.
但重生得救後的基督徒的靈性生命.
是遠超過我們肉體的生命.
加上有聖靈帶領和掌管下.
使我們重生後的基督徒.
可以勝過現實不能改變的咒語.
今次我打算和大家看一個聖經的故事.
故事當中的主角.
他的一生可以說是非常戲劇性.
故事的主角出於一個大家庭.
這個大家庭在故事中.
跟電視劇的情景,小說的情節很相似.
內容是主角約瑟.
約瑟被賣的原因.
他被賣後的遭遇.
和他被賣後的結果.
希望大家透過約瑟被賣的原因.
他被賣後的遭遇.
和他被賣後的結局.
讓我們可以從約瑟的經歷中.
學習一些熟悉的功課.
下次我會和大家分享.
主角約瑟被賣的前因.
完.
\newpage



\section{}
\label{sec:qk8r6sElCdc}
\textbf{Kingdom Values III - Sermon on the Mount - Pastor Sam Johnson}
\newline
\newline
連結: \href{https://youtube.com/watch?v=qk8r6sElCdc}{\texttt{https://youtube.com/watch?v=qk8r6sElCdc}} ~~~~ 語音日期: 2020-05-04
\newline
\newline
\hyperref[sec:o_g_uZyyIHo]{\small{< < < PREV SERMON < < <}}
~
\hyperref[sec:index_chronic]{\small{[返順時目]}}
~
\hyperref[sec:index_scriptual]{\small{[返順卷目]}}
~
\hyperref[sec:VkJzbPZVpM0]{\small{> > > NEXT SERMON > > >}}
\newline
\newline
$^{1}$Hey family, it's good to be with you again..
I hope that you are well and that you're safe and that you're experiencing God's presence.
and His goodness at this time..
I know I say that every week, but I mean it..
I pray those things for you, that God will protect you, keep you healthy and safe, and.
also just keep you encouraged and keep you in His peace and in His comfort and in His.
joy..
Today, we're continuing on with this study we've been doing, looking at kingdom values.
that Jesus taught us in the Beatitudes, which is a part of the Sermon on the Mount in Matthew.
5..
And specifically, we're looking at verse 4 today, Matthew 5, verse 4, and it says this,.
"Blessed are those who mourn, for they shall be comforted.".
Blessed are those who mourn, for they shall be comforted..
Now this verse has a lot of really deep meaning for us and deep meaning for what it means.
to be a follower of Jesus, a disciple, for what it means to have a right relationship.
to God..
And in a really deep way, it tells us a lot about what it means to have a right relationship.
to the people around us and to the world around us..
We're going to get into some of those things today..
But to start, I want us to just think about what it means..
What does it mean when it says that blessed are those who mourn and that they will experience.
comfort?.
What do those words mean?.
Well, the word for mourning in the original language is a lot like what you might think..
It's a word that communicates grief or lament..
And so it's communicating a sense of sadness or grief, but it's a very strong word..
It's communicating a strong sense of sadness or grief..
It was a word that was used often when someone had died and people were mourning and really.
grieving deeply over that loss..
And the word for comfort in the original language communicates a sense of receiving encouragement.
or consolation or strength being strengthened..
And so again, it means kind of what we might think..
It means receiving some comfort, some care, some encouragement, some consolation..
That's what it's referring to..
But this verse is more than just the words..
It's the ideas that are in this verse that are really significant for us..
So we're going to start out just talking about a few of the things that it doesn't mean,.
a few of the things that we might assume about this verse initially that aren't necessarily.
what this verse is talking about..
And first of all, it doesn't mean that we're supposed to be unhappy or kind of sad people.

$^{41}$as Christians..
That's not what this verse is communicating..
It's not saying that it's a kingdom value for the followers of Jesus just to be sad.
all the time or to be in despair all the time..
That's not who we're supposed to be..
And we know from other parts of the Bible, it's clear the Bible talks about the fact.
that we have the joy of the Lord as our strength, that we're supposed to be rejoicing in what.
God has done for us..
We're supposed to be joyful people..
And this verse reminds us that that joy that we have as believers in Christ is not a fake.
joy..
It's not a false joy..
It's not the type of pretentious thing we put on where we pretend to be happy when we're.
not or where we're kind of out of touch with reality and when there's tragedy or hurt going.
on around us, we're just kind of got our head in the clouds and we don't even care and we're.
like, "Oh, I'm still great..
Everything's fine," even when things are not fine..
We're not supposed to be like that as believers, but we can and should have joy that is beyond.
our circumstances because it's resting in our security in God, not in our circumstances.
themselves..
We can have joy beyond the things we experience..
And we've talked about that some in the past..
But even though we know that the Bible communicates that we're supposed to be joyful people, still.
this verse tells us that it's a kingdom value to mourn and to experience God's comfort..
So what does it mean when it says that we should mourn?.
And again, this is not just talking about simply grieving or being sad over loss or.
tragedy..
I think sometimes that's one of the things we assume about this verse..
When we just look at it initially, we think, "Well, it's just saying that when we go through.
loss or tragedy and we mourn and we're sad and we're grieving, that God cares and He.
comforts us and He takes care of us.".
And that is completely true..
The Bible makes that clear in so many other places, but I don't think that's what this.
verse specifically is talking about..
But I want to stop there for a second and be real clear about that..
God does care about our grief..
God cares about the pain in our lives..
God does want to comfort us in those things..
And that's clear in so many places in the Scripture..
Hebrews 4, verses 14 and 15 talk about how Jesus is our high priest and how He has been.

$^{81}$touched with the feelings of everything that we go through because He went through it Himself.
when He was here on the earth living as a human..
And that verse tells us crystal clear that Jesus cares about what we feel and what hurts.
us and our pains and our grief..
He cares about those things and He wants to meet us there and comfort us and encourage.
us..
Psalms 34, verse 15, we talked about this a few weeks ago, such a beautiful verse that.
says, "The eyes of the Lord are toward the righteous and His ears are listening to their.
cries.".
What that verse is saying, His ears are toward their cries..
It's saying that God is looking and listening to our need and He cares about our hurt and.
our pain and our grief..
Second Corinthians, chapter 1, verses 3-7 talks about how God is the God of all comfort.
and He comforts us in all of our affliction, all of our pain..
So I want to be clear about that..
It is a biblical truth that God cares about your pain and your grief and your tragedy.
and your mourning..
And when you experience that, God wants to meet you in that and bring you comfort and.
peace..
That is a biblical truth..
And if you're experiencing that right now, grief and mourning, then God wants to meet.
you in that and bring you comfort and peace..
But while that is a biblical truth, again, I don't think that's necessarily what this.
verse is specifically talking about..
So we're going to dig a little bit deeper into this and we're going to go for a little.
drive and talk more about what it means to mourn, what it does mean to the followers.
of Jesus, to disciples of Jesus, to mourn, and what it does mean to receive God's comfort..
As we talk today about what it means to mourn and what it means to be comforted, I want.
us to look out..
I'm sure you're tired of looking at me every week..
I'm tired of looking at me every week..
So today we're going to go for a drive and since we're in my neighborhood, we're going.
to look at my neighborhood..
But more than that, we're going to look at ourselves..
We're going to look at Chicago..
We're going to look at us..
We're going to talk about what it means to mourn over our brokenness and we're going.
to talk about the comfort that God has for us in that mourning and the hope that we have.
for God's solution, God's redemption, God's rescue..
Last week, as we discussed Matthew 5.3, I said that the Beatitudes build on one another..

$^{121}$So to understand the meaning of verse four, we're actually going to start with verse three..
Last week we talked about what it means to be poor in spirit..
The fact that it means we are aware of our need for God..
We're aware of our spiritually bankrupt state without Him..
Our need for Him is profound..
It's inescapable..
And we discussed how poverty of spirit is an essential kingdom value because it's necessary.
for entrance in the kingdom..
Entrance into the kingdom of God comes by God's grace through faith in the forgiveness.
of our sins through the sacrifice of Jesus on our behalf..
God gives the grace..
We're supposed to exercise our faith in Him and that's how we enter into God's kingdom..
But it can't even happen if we don't know we need Him..
If we're not aware of our need for God, if we're not aware of our own sinfulness and.
the bankrupt state of our own virtue apart from God, then we'll never ask forgiveness..
We'll never seek repentance..
We'll never accept God's grace and we'll never experience the all-encompassing and.
completely sufficient solution of God's redemptive plan..
None of that will happen unless we're aware of the problem of our own brokenness..
So you have to become poor in spirit to even enter the kingdom..
However, being poor in spirit is also a primary value in the kingdom because even once you've.
entered the kingdom, once you have relationship with God, you won't grow in that relationship..
You won't become more like Jesus..
You won't grow in love and joy and peace and the other fruits of the spirit unless you're.
aware of your continuing need for God..
Unless you're acknowledging your need for Him and asking for God's activity and work.
in your life and asking for God's activity and work in your life to increase..
This is when we first begin a relationship with God..
As we continue following Jesus, we still need to be aware of our own brokenness and need.
in order to experience the solution of God's grace, God's strength, and God's transformation.
in our lives..
We have to remain poor in spirit to grow and progress in our relationship with Him and.
in the kingdom of God..
Now, all this awareness of our own brokenness can be a little bit discouraging..
And that brings us back to the meaning of Matthew 5, 4, the verse we're looking at today..
See when this verse talks about mourning, I believe that it's talking about a mourning.
over sin, a mourning over brokenness, a mourning over the reality that we as humans have colossally.
wrecked our lives and our world..
Remember that God created this world as a perfect place and He created us as perfect.
beings in it..

$^{161}$Originally, there was no sin..
There was no suffering..
There was no disease..
There was no hatred..
There was no murder..
There was no racism..
There was no poverty..
There was no hunger..
There was nothing painful or hurtful in that original plan..
But we humans wrecked it with our own selfishness and we've been hurting ourselves and hurting.
one another ever since..
And the more we begin to grasp the reality of this brokenness, the more grief we'll feel,.
the more lament we'll feel..
And just like this verse describes it as mourning, it's a deep grief..
It's a deep lament as we really grasp the depth of our need for God, the depth of our.
bankruptcy without Him..
But this isn't just some type of hopeless despair because there is hope..
And we'll get to that in a minute..
But it is a very real assessment of the reality of our brokenness..
We need to have that..
And when we get that real, when we get that up close with our brokenness, it's ugly..
It makes us sad..
Humans don't like to mourn..
We don't like to be sad or uncomfortable..
And a lot of times our first immediate temptation is just to look away, just to ignore the brokenness,.
just to act like it isn't there and try to find something happier to look at..
But Jesus is telling us in this verse that that behavior isn't acceptable for His followers..
We can't just ignore the brokenness..
We can't just imagine it away..
We need to be honest about it and we need to allow ourselves to feel the sting of its.
reality..
That's true in our own lives..
We can't just ignore our brokenness or imagine it away..
But it's also true as we look at the world around us..
As we're looking at Chicago today and more specifically at my neighborhood, you can see.
with your eyes some physical manifestations of brokenness..
There's boarded up houses that have been left to deteriorate over many years as property.
values have declined precipitously..
There's empty storefronts that used to hold bustling businesses and they're now sitting.
silent..

$^{201}$A few of those have been boarded up in the last few months because of the COVID-19 pandemic,.
but the truth is most of them have been shut for years, even decades..
And the few small businesses that were here when the pandemic began, sadly, probably will.
not survive the economic downturn..
That's a physical manifestation of brokenness in my community, in our city..
There are boarded up churches with no worshippers..
That's a physical manifestation of brokenness in my community, in our city..
Pay attention when you see the police barricade with officers standing around it and manning.
it..
That's the block in front of Roseland Community Hospital..
It's a small neighborhood hospital with only 134 beds..
And they've been so overwhelmed with the COVID-19 patients that they had to set up triage tents.
in the parking lot to process people who are walking up and they had to use their ER as.
an overflow for their ICU..
That's a physical manifestation of brokenness in my community and in our city..
Keep looking and you'll also see the human toll..
That's the most devastating part of it all..
People whose lives have been wounded by the pain of poverty, the bondage of addictions,.
the trauma of abuse or violence, the sting of abandonment..
You can see those things written in people's faces..
And it hurts to see those things..
But don't just look away..
Remember, Jesus ruled that out in verse 4..
We need to see, we need to mourn..
That's what he said..
"Blessed are those who mourn.".
It's a kingdom value to see, to mourn..
But if the first temptation of seeing our brokenness is just to look away, then the.
second is to shift blame from ourselves, to explain away our involvement and absolve ourselves.
of any guilt or responsibility..
It's easy to do that..
It's easy to do that as we look at Chicago, as we look at my neighborhood..
It's easy to tell ourselves that none of it's my fault and of course I don't want any of.
these things to happen..
I don't wish for these physical manifestations of brokenness to be here, so it doesn't have.
anything to do with me..
But that's too easy..
And I believe Jesus rules out that type of thinking in verse 4 as well..
The mourning that Jesus is talking about here always starts with my own poverty, my own.
need, my own sin, my own guilt..

$^{241}$You know, when we look at a neighborhood that has a lot of needs, we often just see the.
superficial issues..
We see the crime, we see issues of poverty, we see issues of homelessness or boarded up.
buildings, things like that..
And a lot of times we can just simply assume that the only source of those issues is in.
the community itself..
Now clearly there are issues of brokenness and sin in every community that need to be.
addressed by that community..
And that's certainly true here in my neighborhood..
But there are also many external issues involved as well..
This neighborhood is an excellent example..
It went from thriving to struggling by means of a multi-decade process that started with.
racial prejudice and fear, discriminatory and predatory housing practices that were.
motivated by racial prejudice, disinvestment in community businesses and institutions that.
followed discrimination, even a loss of jobs due to changes in the global economy..
And along the way there were breakdowns in the relational health of families, collapses.
in the moral health of churches and religious leaders, plus a heaping helping of personal.
and political selfishness and greed and corruption..
That's the Chicago way, right?.
Well maybe so, but it's more aptly the human way..
That's really all of us..
And it's so easy for us to look at a specific community or a specific person or place, especially.
if that community or that person or that place is different than we are, and just assume.
that it's their problem, not mine..
But that's not true..
It's hard to see all of those other things when we look at a boarded up house, but it's.
all there..
All that human brokenness wrapped up and reflecting us, ourselves..
Because no one actually tries to make these broken things happen..
No one wishes for them or wants them to happen, but they do happen all the time..
And they happen while a lot of us are so busy rushing in pursuit of our own success and.
comfort and happiness that we end up running away from the places of need as fast as we.
can..
Even many of us who are followers of Jesus and who know better and who should be involved.
in serving Jesus in those places run away to pursue other things..
And in that way, we're all a part of the brokenness..
Because honestly, we're all wrapped up in selfishness..
And as we spoke about a few weeks ago, we're all wrapped up in building our own little.
kingdom where we're the monarch and everything is custom made for our own happiness..
So we can't just shift the blame and absolve ourselves..

$^{281}$We can't just say, "It's not my fault..
It has nothing to do with me.".
We need to be able to see our own brokenness reflected in the brokenness of the world around.
us..
That's what Jesus calls us to do..
Jesus calls us to see, to be aware of our need for Him, to be aware of the need around.
us, to mourn, not to blame others, but to mourn our own need, our own poverty, our own.
brokenness..
Now, at this point, you might be thinking, "Wow, this is kind of a heavy message..
I think I'm ready to turn it off and go watch some cartoons or something.".
But stay with me..
This is heavy..
Jesus was challenging us with some deep and honestly offensive truth here..
But don't give up yet because there is hope..
See we're only halfway through this verse..
And in the second half of the verse, Jesus promises that those who mourn will experience.
comfort..
Now, what does that mean?.
What does this phrase, "They shall be comforted," mean?.
Well, when we visualize comfort, sometimes we think of God putting His arm around us.
and patting us on the back or patting us on the shoulder and saying, "Hey, it's going.
to be okay.".
And it does mean that..
That's a part of what it means in this verse, but it means more than that..
It means so much more..
It means that God has a solution to our brokenness, a solution that is so beautiful and so perfect.
and so complete that when we turn from seeing our brokenness to beholding God's solution,.
we will literally go from mourning to dancing, as Psalms 30 verse 11 says..
The solution is redemption..
Through Jesus, we can be personally healed of our brokenness and have the ashes of our.
failure transformed into the beautiful threads of His robes of righteousness..
That is true and complete comfort..
As we are made aware of our brokenness, we are comforted by the resolution of His redemption..
And that's true for the world around us as well..
It's true for our neighborhoods, our families..
The hope for transformation in my neighborhood and in yours is the redemption of Jesus, the.
transformative power of the gospel, healing one heart, one family, one block at a time..
You see, Jesus wasn't a politician..
He didn't promise a chicken in every pot or more shovel-ready jobs..
Jesus was the Savior of humanity and He promised something infinitely better..

$^{321}$He promised love in every heart, peace in every home, healing for every wound..
And as we, the followers of Jesus, become aware of the brokenness around us, we can.
be comforted by the understanding that we get to be a part of the solution..
We can be a small part of the grand redemptive plan of God..
We can be a part of this kingdom of heaven that we've been talking about these past few.
weeks being present in the right here and right now..
How does that happen?.
Well, remember, the kingdom of heaven is the rule and reign of God..
As we obey God and as His presence is active in our lives and we bring the kingdom of heaven.
into contact with our friends, our families, our neighbors, as we serve them and care for.
them and share the love of Jesus with them, we have the opportunity to invite them to.
begin a relationship with God through grace and faith..
They'll have the opportunity to experience the redemptive work of God in their life and.
to experience the healing and restoration of Jesus for themselves..
Followers of Jesus get to be a part of God's solution because they're able to participate.
in the redemptive plan of God..
That is true and complete comfort..
You get to be a part of God's solution in your neighborhood, on your block, in your.
family and in your own heart and life..
You get to be a part of the work that God desires to do around you to rebuild hearts.
and lives and homes and through those rebuilt lives to rebuild the schools and the businesses.
and the hospitals and everything else in the community..
You get to be a part of the solution..
Now, I understand we don't want to mourn because it offends us..
It offends our sense of comfort and security..
It offends our pride..
But if we don't learn to mourn, we never get to be a part of the solution..
So let's do it..
Let's be poor in spirit..
Let's be aware of our need for God..
Let's allow ourselves to see the brokenness in our own lives and the lives of those around.
us..
Let's even allow ourselves to feel the sting of that brokenness..
Let's not blame others..
Let's acknowledge our guilt and our part and our sinfulness..
Let's mourn and then let's move to be a part of God's solution..
Let's rejoice in the beauty of God's comfort..
Rejoice in the beauty of God's solution..
Rejoice in the beauty of God's redemption..
Let's go out and love and serve and care for those in need..

$^{361}$Let's share the hope of Jesus..
That is true comfort..
Blessed are those who mourn, for they shall be comforted..
I want us to close out today by reading from Isaiah chapter 61..
This is another scripture passage that talks about mourning and comfort and I think it.
summarizes that relationship between those two things really well as we've been talking.
about it today..
Now Isaiah was a Hebrew prophet that lived hundreds of years before Jesus and he wrote.
this book and this chapter to God's people who had been in mourning and chapter 61 is.
a prophetic look forward towards a time when the comfort will come and it's actually about.
Jesus himself..
He's talking about the coming Savior, the coming Messiah..
This is what it says, "The spirit of the Lord God is upon me because the Lord has anointed.
me to bring good news to the poor..
He has sent me to bind up the brokenhearted, to proclaim liberty to the captives and the.
opening of the prison to those who are bound, to proclaim the year of the Lord's favor.
and the day of vengeance of our God, to comfort all who mourn, to grant to those who mourn.
in Zion, to give them a beautiful headdress instead of ashes, the oil of gladness instead.
of mourning, the garment of praise instead of a faint spirit, that they may be called.
oaks of righteousness, the planting of the Lord, that he may be glorified..
They shall build up the ancient ruins..
They shall raise up the former devastations..
They shall repair the ruined cities, the devastations of many generations.".
And that's just the first four verses..
I'd encourage you to take some time later today and read the rest of that chapter..
It's such an encouraging chapter..
But here's the cool thing..
The Bible records a story where Jesus himself read that chapter..
In Luke chapter 4, it talks about at the very beginning of Jesus' ministry, he goes to the.
synagogue in his hometown in Nazareth, and he reads those words that we just read..
He reads Isaiah 61, and for the first time, those words were read by the person they were.
written about..
Jesus was there to fulfill them..
This is what Luke chapter 4, verses 20 and 21 says..
After he read that passage, it says, "And he rolled up the scroll and gave it back to.
the attendant and sat down..
And the eyes of all in the synagogue were fixed on him..
And he began to say to them, 'Today this scripture has been fulfilled in your hearing.'".
That is true comfort..
That is something to rejoice in..

$^{401}$Jesus is the comfort..
He brought the victory..
He brought the kingdom..
And that's something worth rejoicing in, worth being excited about..
You can go ahead and shout about that, and it's not going to bother anybody but your.
neighbor, so go right ahead and shout..
That scripture, the comfort, has been brought..
It's been fulfilled in our hearing..
Jesus brought it with him..
Blessed are those who mourn, for they shall be comforted..
They'll experience the comfort of God's redemptive plan through Jesus Christ..
Let's pray..
God, thank you that you're honest with us..
Thank you that you tell us the truth, that you don't let us just look away from our brokenness.
or from the brokenness of the world around us, that you don't let us just ignore it and.
try to wish it away..
You're honest about our need for you..
And God, I pray that you'd help us to be honest about it too..
Help us to be poor in spirit..
Help us to see and acknowledge our need for you..
God, help us to see it so deeply that it even causes us to mourn and that it breaks our.
heart because of how broken we are and how broken the world around us is and how wrecked.
everything is because of sin and selfishness and all the terrible things that are in our.
world, God..
But as we mourn, God, I pray that we would not be hopeless, we would not be in despair,.
but we would look to you and find your hope, that we would find your solution, that we.
would rejoice that your solution is present in our lives to heal us, to restore us, to.
change us..
But through relationship with you, we get to be a part of your solution, a part of your.
redemptive plan, and a part of what you are doing to change the world, God..
And I pray that we, your people, would be a part of it even now, even in this moment.
of COVID-19 with all the fear and uncertainty and all the need..
I pray that even now, we, your people, would be a part of your redemptive plan to love.
and to serve and to care and to share hope with those around us..
God, I pray that we would be the pouring spirit, we would see our need for you, and we would.
be comforted as you transform us and make us a part of your redemptive plan for our.
families, our communities, our neighborhoods, and our world..
In Jesus' name, amen..
[MUSIC PLAYING].
♪ ♪.

\newpage



\section{}
\label{sec:VkJzbPZVpM0}
\textbf{2 十一少被賣記原生家庭(二)}
\newline
\newline
連結: \href{https://youtube.com/watch?v=VkJzbPZVpM0}{\texttt{https://youtube.com/watch?v=VkJzbPZVpM0}} ~~~~ 語音日期: 2020-05-05
\newline
\newline
\hyperref[sec:qk8r6sElCdc]{\small{< < < PREV SERMON < < <}}
~
\hyperref[sec:index_chronic]{\small{[返順時目]}}
~
\hyperref[sec:index_scriptual]{\small{[返順卷目]}}
~
\hyperref[sec:3aPT0T6_ZVI]{\small{> > > NEXT SERMON > > >}}
\newline
\newline
$^{1}$上次說到,人生就像一部戲.
而我們每一個都有自己要扮演的角色.
今次我是從《創世記》第37章2-11節帶給大家.
來看看主角約瑟悲埋的前因.
主角約瑟的父親和他四個妻妾一共生了12個兒子.
而主角約瑟是排行第11.
約瑟和他父母的弟弟都是父親最愛的妻子所生.
約瑟是親生父母年老時第一個所生的大兒子.
故此約瑟至少就深得父親的特別同愛.
父親親手為約瑟做了一件彩色的衣服.
在《聖經》中從未提及過父親親手為其他兒子做過衣服.
似乎好像暗示了父親做衣服給約瑟是有特別的意思.
有可能是想將家族的繼承權交給約瑟.
童年時的約瑟因為年少無知.
經常在父親面前告發兄長的罪行.
在《聖經》中沒有記錄過10個兄弟有什麼罪行.
但約瑟告發兄長這個舉動.
足以使兄長至少排斥約瑟.
不僅如此 約瑟先後還做了兩個夢.
這兩個夢的確加深了十兄長對約瑟的疾度.
在《創世記》第37章5到11節.
那裡提及約瑟第一個夢.
約瑟夢見10個哥哥向他來跪拜.
當約瑟發第二個夢的時候.
他夢見太陽,月亮和11顆星同時都向他來跪拜.
太陽,月亮分別代表了約瑟的親生父母.
而11顆星代表了約瑟11個兄弟.
約瑟因為年少就將兩個夢都告訴他家人.
結果父親雅各聽完後.
就當著兄長和母親面前即刻責備約瑟.
但另一方面 父親雅各亦將約瑟的話藏在心裡.
而約瑟10個哥哥聽完約瑟所發的這個夢後.
就更加疾度約瑟.
兄長疾度約瑟的程度所帶來的嚴重後果.
竟然成為日後約瑟有被殺的危機和被賣的結果.
故事說到這裡 不知道大家有沒有引起你們的聯想呢?.
為人父母的你 當你要處理到子女之間的衝突的時候.
你們會公平處理嗎?.
當你們要處理這些事情的時候.
究竟是對事還是對人呢?.

$^{41}$大家有沒有想過 父母親處理子女之間衝突的問題.
如果處理得不恰當的話.
是會造成親兄弟之間的不和呢?.
在故事裡的主角約瑟.
被兄長恨他而所帶來的嚴重後果.
竟然是向同父異母的親弟約瑟動殺機.
手術之間可以為了權.
為了名 為了味 反目 承受 甚至骨肉相殘.
還記不記得在聖經裡所記載.
人類的第一宗殺人事件嗎?.
都是發生於兩兄弟之間.
起因都是因為嫉妒.
在故事裡 主角 兄長來向約瑟動殺機.
主要原因是與原生家庭所產生的問題有關.
約瑟童年和青年時期.
都有他幸福的一面.
因為親生父母對他的痛愛.
但亦由於約瑟的父親對約瑟的偏愛.
而令約瑟受到兄長的嫉妒和排斥.
甚至導致約瑟日後有被殺的危機和閉脈的結果.
約瑟自出母胎至17歲.
17年來在他身上所發生的事情.
都是不由自主的.
例如 他的親生母親是父親最愛的妻子.
父親偏愛著自己 親手為自己做採衣.
兄長至少就對他 排擠他.
他做夢 夢見十個哥哥跪拜他.
連親生父母都要跪拜他.
父親吩咐他去找兄長.
他只有聽命.
結果就被兄長埋到埃及.
從此之後就與原生家庭分開.
流落異鄉 做埃及人的奴隸.
這些種種事情的發生.
都不是約瑟有自由可以選擇.
亦不是約瑟有能力可以改變.
下一次我會和大家再探討一下.
造成約瑟閉脈的原因和約瑟原生家庭的問題.
究竟兩者之間有什麼關係.
請不吝點贊訂閱轉發打賞支持明鏡與點點欄目.

\newpage



\section{}
\label{sec:3aPT0T6_ZVI}
\textbf{3 十一少被賣記:原生家庭(三)}
\newline
\newline
連結: \href{https://youtube.com/watch?v=3aPT0T6-ZVI}{\texttt{https://youtube.com/watch?v=3aPT0T6-ZVI}} ~~~~ 語音日期: 2020-05-05
\newline
\newline
\hyperref[sec:VkJzbPZVpM0]{\small{< < < PREV SERMON < < <}}
~
\hyperref[sec:index_chronic]{\small{[返順時目]}}
~
\hyperref[sec:index_scriptual]{\small{[返順卷目]}}
~
\hyperref[sec:I_xWBBrkSzQ]{\small{> > > NEXT SERMON > > >}}
\newline
\newline
$^{1}$上次提到約瑟原生家庭的問題.
是造成主角約瑟被兄長埋的原因.
究竟約瑟的原生家庭問題在哪裡呢?.
而家庭的問題又怎會導致約瑟被埋呢?.
約瑟的家庭裡面兄弟之間出現問題.
主要原因是兄長對約瑟的疾度.
而造成兄長對約瑟的疾度.
是因為他們親生父親雅各一夫多妻.
造成同父異母一共12個兒子.
但親生父親雅各最愛的妻子.
就只有約瑟的親生母親.
所以造成其餘10個兒子.
就是一個既關係複雜又不健康的家庭裡成長.
雅各的10個兒子.
自幼的自我形象已經很低.
又缺乏安全感.
10個兄長自我形象低加上沒有安全感.
都是造成日後約瑟被疾度和被排擠.
還有就是約瑟年少又無知.
經常來向父親告發兄長的罪行.
引致兄長更痛恨約瑟.
10兄長自我形象低,無安全感.
除了造成疾度,排擠,約瑟外.
加上兄長很擔心,亦很害怕.
父親終有一日將家族的繼承權傳給約瑟.
引致雅各10個兒子為了爭父愛.
又為了爭繼承權.
於是就疾度約瑟,排擠約瑟.
甚至有殺他的念頭,賣他的行動.
弟兄姊妹.
人出生之前有很多事的確不由自主.
我們不能選擇自己出生時的父母是誰.
自己出生排行第幾.
在什麼年代,什麼家庭背景出生.
不過人出生後並不是每樣東西都不能選擇或改.
當然我們無法指改變原生家庭.
但人可以有自由選擇.
用什麼方式,用什麼態度,用什麼想法和行動.
去面對和處理.
就好像故事中雅各這10個兒子.

$^{41}$雖然他們無法只改變他們出生的家庭背景.
例如人出生就在一個一夫多妻制的家庭裡長大.
一個父親的愛同時要和幾個不同母親的兄弟分享.
而自己的親母又不是父親摯愛的妻子.
10兄弟至少就在一個缺乏安全感的環境裡長大.
弟兄姊妹.
雖然現實有很多人和事.
是無法只隨著自己的心意來成就或改變.
不過我們可以改變我們自己.
改變我們自己對人,對事的想法和看法.
亦有自由來選擇自己的反應和行動.
例如約瑟的10個兄長.
雖然無法只改變自己原生家庭.
但他們可以有自由改變自己.
例如對親弟約瑟的質度,想法.
和對親弟殺和賣的反應行動.
不過人的確很軟弱.
就像聖經的拜子和墨子的比喻.
墨子被這個世界周圍的拜子影響.
墨子的成長受到身邊拜子的影響變得非常受虐.
人其實一樣.
在這個聚的世界生活.
不能單靠自己的力量就勝過聚的人有.
約瑟的兄長就像這個例子.
他們警隊親弟來動殺機.
大家試想一下.
約瑟當時只是一個17歲的少年人.
大家回想一下當年自己16,17歲的時候.
我們說了多少錯的說話得罪了人.
我們又做了多少錯的事來傷害了人.
我們說錯話,做錯事.
絕大多數原因是因為當時年少無知.
少年時期的約瑟當時的思想也非常簡單.
他自幼就深得父母的愛.
他從小到大都沒有需要主動去爭取.
因為至少他就得到父母親完全的愛.
但他10個兄長就很不同了.
這10個兄長至少知道.
自己的母親並不是父親的最愛.
當約瑟出生後.

$^{81}$就更讓這10個兄長感到非常不安.
他們的不安是因為缺乏了父親完全對他們的愛.
這10個兄長為了爭愛,爭權,爭利.
以至有殺親弟的念頭.
甚至埋親弟的行動.
約瑟和家人生活了一共17年.
可以想像一下.
年輕人自出生至17歲.
在身邊的家人有極之愛他們的親生父母.
但同時也有極之恨他的兄長.
兄長恨約瑟的程度.
竟然嚴重到要殺他.
故事說到這裡.
有沒有引起大家什麼聯想呢?.
大家會不會有一個疑問.
對於原生家庭所帶來的負面影響.
而這些負面影響又會帶來有多少嚴重的後果呢?.
下一次我會和大家再深入探討.
約瑟的原生家庭負面影響.
所帶來有什麼後果.
謝謝大家.
請不吝點贊訂閱轉發打賞支持明鏡與點點欄目.
\newpage



\section{雅各書}
\label{sec:I_xWBBrkSzQ}
\textbf{2020/5/3 智慧廚房 vlog 信息 \#3 (雅各書 3) - 陳崇基牧師 (updated)}
\newline
\newline
連結: \href{https://youtube.com/watch?v=I-xWBBrkSzQ}{\texttt{https://youtube.com/watch?v=I-xWBBrkSzQ}} ~~~~ 語音日期: 2020-05-09
\newline
\newline
\hyperref[sec:3aPT0T6_ZVI]{\small{< < < PREV SERMON < < <}}
~
\hyperref[sec:index_chronic]{\small{[返順時目]}}
~
\hyperref[sec:index_scriptual]{\small{[返順卷目]}}
~
\hyperref[sec:v2v5pAP5WjA]{\small{> > > NEXT SERMON > > >}}
\newline
\newline
$^{1}$歡迎各位又再回來我們智慧廚房.
我們邊煮邊講.
最重要的其實不是煮什麼.
而是我們講什麼.
所以今天呢.
在還沒講煮什麼之前.
我們先看看聖經.
因為最近那兩集我們在講《雅各書》第一章和第二章.
這次應該是第三章了,是不是?.
他說:我的弟兄們,你們不應該有太多人作教師.
因為知道我們作教師的.
將受更嚴厲的審判.
我們在許多事上都有過錯.
假如有人在言語上沒有過錯.
他就是完全的人.
也能夠控制全神.
這裡第三章就轉了另一個 topic(話題)了.
在講我們怎麼講話.
講到不要太多人做教師.
為什麼呢?.
因為在當時的社會裡面.
教師是很重要的.
教師只是少數的人.
能夠懂得寫,懂得讀,懂得教.
所以當時教會.
很尊重那些教師.
其實教師就是當時教會領袖,長老或者執事.
就是教會的領袖.
他在這裡講.
就說不要搶著做.
為什麼呢?.
因為到將來上帝要審判我們.
查問我們的時候.
就會更加嚴厲.
就會 strict(嚴格)很多的.
這裡為什麼呢?.
他就跳到第二節了.
最自然就是.
為什麼不要搶著做教師呢?.
因為我們呢.

$^{41}$在言語上我們每個人都會有過錯的.
有誰在言語上沒有說過錯的話呢?.
我們經常說病從口入.
禍從口出.
所以呢.
我們要想一下.
當時為什麼雅各會警告人們說這些話呢?.
或者警告當時的教會.
不要搶著做領袖.
不要搶著做教師.
這個我們邊主邊講的時候.
我們就會解釋一下.
想一想.
[音樂].
今天會煮什麼呢?.
就是跟舌頭有關的東西了.
會不會是牛舌呢?.
或者豬舌呢?.
不知道大家是不是喜歡吃內臟的人.
有很多人很討厭吃內臟的.
不好意思.
今集我們要跟內臟有關.
因為講到舌頭.
所以今次我們煮什麼呢?.
我們就煮鴨舌.
煮鴨舌.
鴨舌不知道大家喜不喜歡吃呢?.
如果喜歡的話.
你在 comment(留言)那裡說.
喂,我很喜歡吃.
如果你不喜歡吃的話.
你給它 thumbs down(下降).
哇,今次不行.
但是我們迫不得已了.
因為經文呢.
講到舌頭.
我們一定要煮舌頭有關的東西.
那麼買了一包新鮮的鴨舌.
10個2號3號.
我還沒買過這些東西.

$^{81}$覺得真的很貴.
不過無論如何.
今天我們要做滷水鴨舌.
講到鴨舌的時候.
我在 research(研究)裡面告訴我.
原來有很多很多的做法.
可以是椒鹽.
滷水.
可以是花雕酒.
可以是啤酒鴨舌.
可以是玫瑰露鴨舌.
可以是醬包.
還是簡單的鴨舌呢?.
原來有很多很多不同的做法.
那麼今次我們就做滷水鴨舌了.
首先要鴨舌.
我們大概不需要多.
300克就夠了.
大概半斤左右.
另外我們就需要一些老抽了.
我們準備了3湯匙的老抽.
然後呢.
就需要冰糖.
20克冰糖.
薑呢.
就需要4塊的.
我們切了4塊薑.
然後煮這個滷水的時候.
我們需要的水是125毫升.
最後我們煮滷水的時候.
後面要加紹酒.
2個湯匙.
聽說不只是煮滷水鴨舌.
滷水的食料呢.
就是一定要好的醬油.
還有醬油要到位的.
什麼意思呢?.
就是說不可以太多.
太多就會浪費嘛.
太少可能又不夠.

$^{121}$吸收不夠味道.
所以那個的份量就要調得對的.
好,那麼我們準備了這些材料的時候.
我們就可以開工了.
不過話說回來.
我覺得這個鴨舌真的挺貴的.
說到鴨舌這麼貴.
我就想起智慧文學裡面.
經常提到智慧的言語.
《雅各書》當然是其中有智慧文學色彩的書信.
智慧文學經常提到智慧言語.
我想起一句話.
《箴言二十章十五節》.
他說:雖有金紙許多紅寶石和寶器.
但充滿知識的嘴唇.
比這一切更加寶貴.
講到智慧的言語不是這麼簡單的.
原來物以罕為貴.
其實這個時代智慧的言語真的是物以罕為貴的.
你想一下今時今日很多人.
說一些話都是不負責任的.
或者是全無智慧的.
或者是完全顛倒是非的.
最近我聽到好像說.
香港有一個法官判了一個斬傷人的人.
判詞的時候居然說這個人蓄意傷人的時候.
他是高尚情操.
真的沒想到.
言語可以這樣顛倒是非的.
或者那些一國之君.
現在那些國家領袖.
特別是在疫情裡面說很多話.
完全都不負責任的.
我在搜查資料的時候.
原來鴨舌曾經是宮廷的食物.
原來在清朝的時候.
慈禧太后最喜歡吃的一碟菜.
就是鴨舌.
是清燉鴨舌.
就是用鴨湯來燉鴨舌.

$^{161}$我不知道你有沒有摸過鴨舌呢?.
其實鴨舌很小的.
因為鴨嘴裡面.
我們看到的鴨舌其實只有這麼多.
裡面的軟骨.
包括在買到的鴨舌裡面.
所以真的很小的.
《雅各書》也是說這個.
我們繼續說.
三章第三節開始.
他說我們將雀環加入馬嘴.
使馬順服.
就能夠駕馭整個馬的全身.
他說再想想.
雖然船隻是很大隻的.
又被狂風來催逼.
不過舵手只是用小小的舵.
就可以隨意來操作.
照樣他說人的舌頭.
雖然只是一個很小的肢體.
卻會說誇大的話.
他說想想.
星星之火可以燎原.
舌頭就是火.
在我們八體當中.
這麼多部位裡面.
是不二的世界.
能夠污穢全身.
將整個生命.
在運轉當中焚燒起來.
而且是被地獄之火去點燃的.
真的很厲害.
他說到各樣的飛禽,走獸.
自然界裡面的昆蟲.
水族裡面的東西.
都可以被人馴服.
不過人總是無法制服舌頭.
這裡翻譯得挺好的.
他說舌頭是什麼呢?.
舌頭是諜不休的惡物.

$^{201}$是充滿致命的毒素.
真的很厲害.
不過你想想.
在人生當中.
不僅智慧的言語是少見的.
甚至我們人生的經歷裡面.
我們也被很多沒有智慧的言語.
傷害過.
這裡講了三個例子.
他說馬的爪環.
就是扣在馬的嘴巴上.
騎馬師可以將馬.
他要說他去就去.
停一停.
另外一個比喻就是.
船的舵.
一個很小很小的後面的舵.
可以將整個船操縱.
第三個比喻就是火.
那種山火.
你想想那些山火.
多麼一發不可收拾.
你想想美國加州的山火.
不同地方的森林的大火.
亞馬遜的大火.
是多麼破壞力多大.
但是只不過是小小小小的火種開始.
所以智慧文學經常警告我們.
你講話要想清楚.
你打電話給別人.
去講那些話.
無論是閒言閒語等等.
都要想清楚.
因為那種破壞力是多大.
不單止破壞別人.
可能自己的生命也會隨即破壞.
按照我們的食譜.
首先我們要將鴨舌.
先洗乾淨.
洗乾淨之後.

$^{241}$我們將鴨舌放在熱水裡.
大概煮兩分鐘左右.
不可以煮太久.
因為聽說煮太久.
鴨舌就會縮得太厲害.
那些不喜歡內臟的朋友.
覺得這一集非常難受.
不過請你繼續下去.
因為亞國也覺得很難受.
亞國說到舌頭的時候.
他覺得.
為什麼我們永遠都明白不到.
我們需要有智慧的言語.
剛才比較大一點.
現在開始慢慢收起來.
好,鴨舌可以拿上來.
然後放在筲箕裡.
我們要將它過冷河.
然後等它乾一乾.
現在我們來做豉油黃汁.
鴨舌已經準備好了.
瀝乾水份.
我們將水放進鍋裡.
然後老抽三湯匙.
放在一起.
然後薑.
冰糖等等.
我這裡做了一隻八.
加了一顆八角.
希望香一點.
少許酒先不要放.
煮到冰糖融化的時候.
我們就可以準備將鴨舌放進去.
煮的時候我們再看看亞國書.
第九節.
他說我們用舌頭來稱頌我們的主和天父.
又用舌頭來咒詛.
照著神的形象被造的人.
然後第十節.
同一張嘴和舌頭.

$^{281}$竟然又稱頌主又去咒詛人.
我的弟兄們,這是不應該的.
同一個泉源.
怎麼會湧出甜水和苦水呢?.
又是另一個比喻.
他說我的弟兄們.
無花果樹怎麼會結橄欖呢?.
葡萄樹又怎麼會結到無花果呢?.
鹹水也不會有甜水出來.
第十三節.
他說你們中間誰是有智慧的有見識的.
就應該有美好的生活.
明智的溫柔將自己的行為表現出來.
又說到作為一個人.
我們需要在行為生活上.
表現到那種智慧.
不要只有講智慧.
你記得他開頭第一節說.
你千萬不要亂爭著做教師.
不要想著帶頭.
因為很多時候我們說的話.
神就會有更大的審判.
我們需要知道有一個智慧的生命.
是必須由我們的言語開始.
我們煮冰糖蓉.
那些薑的味道都出來了.
我們就可以將鴨舌加入去.
你看到那些豉汁.
不是很多.
不用這樣.
不斷地浸著.
說到智慧的言語.
我們可能自然會想.
聖經裡面一定是叫我們.
說話要有禮貌.
溫溫柔柔.
不要那麼容易得罪人.
如果得罪人.
我們的舌頭就不會那麼有智慧.
是不是這麼簡單呢?.

$^{321}$其實你想深一層.
其實也不是這麼簡單.
你只看雅各書就知道.
雅各書很多時候說的話都很有挑戰性.
絕對不是今天說的那些政治正確.
其實你想想.
整本的聖經.
由舊約到新約.
都有很多人說真理.
說出一些公義的話.
他們說的話.
很多時候都會得罪人.
特別是為真理為公義發聲的時候.
由舊約到新約.
很多先知一直排隊.
他們每個人都說出智慧的言語.
並且說出的話都是土上帝的喜悅.
但很多時候都忠言逆耳.
事實上有些人說雅各.
是 James the Just.
你記得我們第一集的時候說過.
為什麼叫 James the Just.
就是義者雅各.
因為他說的話都是充滿公義.
為公義發聲.
為真理發聲.
有人甚至說雅各是新約裡面的亞摩斯.
如果熟悉舊約聖經的話.
就知道亞摩斯是所謂小先知書裡面的其中一本書.
亞摩斯說很多話都是針對當時古代以色列人的權柄和君王.
雅各就好像亞摩斯一樣.
不過他是新約時代的.
其實想想耶穌也是這樣.
耶穌也說很多話.
是針對當時的宗教政治權柄的人.
法利賽人.
大祭司等等.
就想想為什麼雅各說這麼多智慧的言語呢.
這個我們就需要看聖經的時候.
要想想上下文.

$^{361}$如果你看第十四節看下去你就知道了.
他說如果你們中間心存黑白的質度和自私.
就不要誇口了.
也不要說謊來抵擋真理.
這種的智慧不是從天上來的.
是屬地的,屬血氣的,屬鬼魔的.
因為凡有質度和自私的地方.
就是他們當時的教會.
就必有擾亂和各樣的壞事.
他說這些是屬地的,屬鬼魔的事.
天上來的智慧.
首先就是純潔的,聖潔的.
其次就是和平,溫柔,謙遜.
並且是沒有恩慈和善果.
沒有偏袒,沒有虛偽.
這些就是締造和平的人.
用和平所栽培出來的異果.
你看到亞國的子女行間.
你就知道當時教會發生了什麼事.
有時候人們去到彼此的質度.
彼此用言語來攻擊.
甚至去陷害其他人.
或者去毀壞其他人的名聲.
亞國下一章甚至說他們殺人.
你猜他們是不是真的用刀去殺人呢?.
當然不是.
只不過他們是 commit murder.
用言語來打敗別人.
抬高自己然後踩低別人.
不過在亞國書的時候.
我又看看兩者之間有什麼文學.
是很像亞國書裡面所說的東西.
其中一本就是叫做.
Testament of the Twelve Patriots.
中文翻譯出來可以叫做.
十二族長的遺訓.
十二族長的遺言.
這個是兩者之間.
猶太人的創作文學.
Creative literature.

$^{401}$這個著作是講什麼呢?.
就是講亞國的十二個兒子.
他們是一家人.
這十二個兒子老了的時候.
他們有些遺言.
臨死前說了些什麼呢?.
回憶自己一生裡面.
其實就是兄弟之間有很多的嫉妒.
很多的爭鬥.
不喜歡這個不喜歡那個.
特別他們不喜歡的那個就是.
虐殺.
想害死虐殺.
想殺死虐殺.
當然上帝保護了他.
然後最後創世記.
你知道結局.
大團圓結局.
那個十二族長的遺訓.
就去教導下一代.
教導那些以色列人.
我們一家人.
我們不可以彼此去爭鬥.
不可以彼此去嫉妒.
這個就好像亞國書裡面所說的.
就是告訴教會.
你們是十二支派的人嗎?.
你們是以色列的民族.
你們就好像一個屬靈的家庭.
你們的言語要有屬天的智慧.
能夠造就大家.
能夠建立大家.
所以你看第三第四章的時候.
你會發現亞國其實不是把東西炒在一起亂來的.
第三章一到十二節.
他講到那種惡言.
abusive language.
我們要提防這件事情.
因為這是屬魔鬼的.
不屬天的.

$^{441}$然後那一段呢.
他就講到我們之間不可以有.
自高自大.
以及去嫉妒.
去攻擊對方.
然後第三部分呢.
他就講到和平.
然後他又回到講到.
弟兄姊妹之間的關係.
教會之間的關係.
不可以有嫉妒.
最後這一段的最後一部分.
他又回到講那種 abusive language.
我們需要留意.
就好像你把所有的醬料.
這些好像不相干的東西.
組合在一起的時候.
大家去配搭的時候.
就成為一個主題了.
好,我們煮到差不多的時候.
我們就加一點紹酒進去.
我們應該差不多了.
最後我們準備一些麻油.
然後呢.
開大一點火.
這叫做收汁.
其實就好像我們煮一頓飯.
我們把不同的醬料.
不同的材料放在一起.
有一個好的配搭.
其實這個也是亞國書的那個.
智慧人生.
言語要有智慧.
我們行為上要行動實踐出來.
然後要遇見考驗等等.
我們需要有智慧去分辨.
怎樣去成長.
好,應該差不多了.
我們加一點麻油.
我們可以上碟了.

$^{481}$你看到沒有多到汁.
不會整個汁都在.
浪費了.
不錯.
這些就叫做簡簡單單.
又一餐了.
真的很簡單.
不過說起智慧的言語.
就絕對不是簡單.
我們試一試.
好香.
很多人就這樣拿一碟鴨舌.
可能送啤酒.
或者喝杯茶.
這樣做一個小吃.
不錯.
不過我覺得可以再好一點.
可能調味料.
有些人.
我剛才加了八角.
或者有不同的陳皮等等.
這個就可以.
悉隨尊便了.
煮鴨舌好像很簡單.
不過呢.
你想想一個智慧的舌頭.
就需要很多操練.
這裡說教會裡面.
我們應該用我們的言語.
顯出智慧.
能夠彼此去建立.
有時候真的我們心裡面.
有很多的疾度.
有很多的憤怒.
或者對其他人不滿.
或者我們自己的高傲.
我們自己的Pride.
這次煮了這頓簡單的.
就想想一個不是那麼簡單的問題.
那個問題就是什麼呢?.

$^{521}$就是我們怎樣才可以.
培養到更有智慧的言語.
你記得第一章的時候.
雅各都已經說了.
你們要快快的聽.
慢慢的說.
再慢一點動腦.
是不是?.
雅各說這個智慧是什麼呢?.
就是時間.
操控那個時間.
聽到要快快的聽.
說話呢.
慢慢來.
沒什麼所謂.
更加不重要的.
就是我們裡面的不滿.
我們的怒氣.
有什麼智慧的言語呢?.
我可以參考.
新約裡面說到.
教會裡面需要有的智慧言語.
其實有很多很多的篇幅.
比如說彼此鼓勵.
彼此去安慰.
彼此去代討.
代討也是用我們的舌頭代討.
讚美神.
因為教會最終的目的.
我們這個家的目的.
就是要高舉上帝.
上帝的榮譽.
上帝的名聲才是最重要的.
不是我們個人的名聲.
所以這個星期.
不妨你想一想.
怎麼可以 practice(練習).
什麼叫做智慧的言語.
在這個星期當中.
不妨你寫一個 text(文字).

$^{561}$或者寫一個電郵.
去欣賞一個人.
另外我們可以.
將我們的智慧言語延伸出去.
不只是教會.
是外面的社區.
現在有很多有需要的人.
需要你去幫助.
需要你去安慰.
或者說一些公義真理的話.
瓦割當時的年代.
那些就叫做孤兒寡婦.
今天可能就是一些被打壓的人.
或者被欺壓的人.
說到智慧的言語.
這一集相信也是.
主得最簡單.
但是也是最難去應用.
最難去實踐的.
因為好像瓦割那樣說.
我們有誰可以立得自己的石頭.
如果立得到的話.
我們就是完全人.
瓦割又說回這件事.
我們需要一個完備的生命.
就是整全的生命.
我們需要在神面前.
來到謙卑.
因為我們每一個.
言語上總會有過失.
總會去得罪人.
總會說一些不能夠建立人的說話.
我們需要來到禱告.
我們同心得有祈禱.
上主我們多謝你.
因為你是那一位說話的神.
並且你的說話每一句.
都是滿有智慧滿有生命的.
主就是因為你的真道.
你所說的話.

$^{601}$你的福音將我們重生.
求你自己親自幫助我們.
我們得到重生的生命的時候.
我們自己的言語.
我們的說話.
都是能夠去建立自己的國度.
建立自己的教會.
建立自己的家庭.
建立屬靈的家.
也能夠建立到身邊其他有需要的人.
幫助我們的說話帶有恩典.
幫助我們的說話.
出自我們一個純潔的心懷.
也幫助我們在這個世界當中.
為你的真理.
為你的國度.
為你的公義去說話.
求你繼續幫助我們.
培養我們有智慧的言語.
我們這樣祈禱奉主耶穌基督的名字.
《城匯網台》是屬於城匯社區教會.
我們位於美國芝加哥.
我們網台的目的很簡單.
希望在網上做到更多好的節目.
好的資訊,祝福更多人.
使更多人認識我們主耶穌基督.
這個時光希望您能夠和我們同工同行.
下面的描述欄有不同的電子奉獻方法.
您看到我們的節目製作.
軟件硬件都是非常簡單的.
您的電子奉獻絕對能夠幫助我們繼續和發展這個時光.
最後,不要忘記.
讚好,分享,留言和訂閱.
請您按下訂閱和鈴鐺.
下一次您就會知道我們有什麼新的動向.
願上帝祝福您.
謝謝大家,謝謝各位!.
\newpage



\section{雅各書}
\label{sec:v2v5pAP5WjA}
\textbf{2020/5/10 智慧廚房 vlog 信息 \#4 (雅各書 4) - 陳崇基牧師}
\newline
\newline
連結: \href{https://youtube.com/watch?v=v2v5pAP5WjA}{\texttt{https://youtube.com/watch?v=v2v5pAP5WjA}} ~~~~ 語音日期: 2020-05-12
\newline
\newline
\hyperref[sec:I_xWBBrkSzQ]{\small{< < < PREV SERMON < < <}}
~
\hyperref[sec:index_chronic]{\small{[返順時目]}}
~
\hyperref[sec:index_scriptual]{\small{[返順卷目]}}
~
\hyperref[sec:teCFMiZwZp0]{\small{> > > NEXT SERMON > > >}}
\newline
\newline
$^{1}$各位歡迎來到智慧廚房.
這次我們是第四集了.
所以講到《雅各書》第四章.
你看到我們已經預備好.
今天要煮什麼菜了.
或者煮什麼.
但是我們先看《雅各書》第四章.
先從一到三節.
他說你們中間的爭執和打鬥是從哪裡來的呢?.
不是從你們肢體中好鬥的私慾來的嗎?.
你們放縱貪慾.
如果得不到就殺人.
你們嫉妒.
如果一無所有就打鬥爭執.
你們得不到因為你們不求.
你們求也得不到.
因為你們的動機不良.
要把所得的耗費在你們的私慾上.
《雅各書》講了三章了.
現在到第四章了.
他就將我們人性裡面的私慾.
從第一章開始.
一邊開口一邊開口給我們.
告訴我們我們人性裡面.
我們有很多的黑暗面.
我們需要去留意.
你可以想象到.
當時這個讀者.
或者因為教會被迫害.
散居到不同地方的教會.
那些屬靈的群體.
裡面發生了什麼事情.
特別《雅各》提到.
他們彼此去嫉妒.
彼此去爭鬥,打鬥等等.
裡面有很多的私慾出來了.
現在在疫情當中.
其實我們很多時候.
我們只是看正面.
我們就會覺得.

$^{41}$好吧,多留一些在家裡.
是不是?.
多一點時間跟家人一起.
但是我們也看到很多的新聞.
說家暴.
婚姻的問題.
家人的關係.
張力就越來越大.
每當在困難當中.
在一些艱難的裡面.
很多時候不只是表現出人的美善的一面.
反而是表現了人性醜惡的一面.
當時的教會.
我相信也是這樣的.
面對很多打壓.
被趕散到不同的地方.
那些群體自然受到很大的壓力.
所以那裡有一些有權有勢的人.
就會欺負一些無權無勢貧窮的人.
那些貧窮的人.
就會在那種情況裡面.
可能會靠攏一些有權有勢的人.
以致他們自己得到保障.
其實今時今日這個世界也是一樣.
你看看一些新聞.
一些頭條.
最近我看到一個頭條.
他說政界裡面一些的問題.
粗暴霸位奪權.
就是人性黑暗的一面.
不再講道理.
誰的權力大一點.
他就是真理.
亞國就說我們屬靈的群體.
我們真的不可以這樣.
不可以跟隨世界那一套.
那麼一個這麼苦的世界裡面.
怎麼樣呢?.
所以我今天就煮一些甜品.
因為這麼苦是不是?.

$^{81}$所以今天我就煮這個.
Pineapple Upside Down Cake.
不知道你有沒有聽過.
是一個反轉來的菠蘿蛋糕.
這是一些西人的甜品.
首先講一下蛋糕本身.
我們有什麼需要準備呢?.
第一.
一有三分一杯的 all-purpose flour.
然後就一個 teaspoon 的 baking powder.
半個 teaspoon 的鹽.
四分一 cup 牛油.
等一下會融化的.
half cup 的白糖.
四分一 cup 的 brown sugar.
然後就大蛋.
vanilla extract.
sour cream 就是四分一 cup.
四分一 cup 的菠蘿汁.
好,另外就講到蛋糕的 topping.
就是上面那層.
那麼我們需要什麼呢?.
四分一 cup 的牛油.
等一下也是會融化的.
半個 cup 的 brown sugar.
然後這個食譜就是說.
要有一罐切好的菠蘿.
那麼我就看到新鮮的菠蘿.
不如這一次就用新鮮的菠蘿.
你看這個新鮮的菠蘿.
已經熟得很厲害.
已經很香了.
然後我們有一些罐裝的 cherries.
這些很少單獨去吃的.
這些 cherries 櫻桃.
通常是雞尾酒加一顆進去.
然後就蘸著.
那麼這個 pineapple upside down cake.
剛剛加這些.
跟那些 pineapple 和那些 cake 一起.

$^{121}$就很好吃的.
唯一一個煮法我會吃的 cherry.
就是這個 pineapple upside down cake.
那麼因為這個叫 upside down cake.
所以我們是反過來做的.
先做那個面.
其實就是先做底.
然後就反轉它.
那麼這個很有意思.
那麼我們先將這些牛油融化.
然後我們就倒它進去這個 pan 那裡.
那麼將它塗得均勻一點.
牛油放進去了.
那麼我們將那些焦糖就鋪進去.
我們先切菠蘿.
這個是新鮮菠蘿.
很便宜的.
但是我平時是這樣切的.
然後再起一些釘出來.
那麼就想起我的雅各書.
開頭第四章那裡這樣說.
他說我們那些的絲肉.
全部在那個釘裡面.
是全部跑出來了.
這個是真的.
當我們是受著壓力的時候.
我們很多時候.
就不知道怎麼做.
或者呢.
我們將我們的黑暗面.
無限的放大出來.
這裡甚至說.
他說也求不到.
因為我們是網球.
和學本就譯為.
是網球.
新譯本就譯為動機不良.
或者動機是邪惡的.
主要是他講到就是我們.
很多時候因為被我們的絲肉.

$^{161}$去敗壞的時候.
連我們的祈禱.
都是為了我們自己的好處.
單單為了我們自己的慾望去禱告.
那麼在這裡很值得反省.
到底我們在疫情當中.
我們的祈禱.
到底是為了上帝的國.
還是純粹為了我們舒適.
當然.
我們在疫情裡面.
我們有很多的禱告都是好的.
譬如我們求上帝去供應.
求上帝來看顧我們.
保守我們.
醫治我們等等.
這些都沒有問題的.
沒有人會在祈禱的時候說.
不如我們為了大家.
中了肺炎病毒來禱告.
沒有人會這樣的.
是不是?.
但是沒錯.
是一個很好的提醒.
到底我們.
我們想離開這個疫境.
想上帝來念穩我們的時候.
我們有沒有想過.
最終的目的是為了什麼呢?.
是為了上帝的國.
上帝的真理.
上帝的公義.
祂的福音.
能夠得到彰顯.
還是純粹我們想生活得舒服一點.
我們自私一點.
我們只是不想這麼辛苦而已.
那麼在這裡講到我們門徒的價值觀.
應該是不同這個世界.
就是基督的門徒.

$^{201}$耶穌也是這樣教導我們去禱告的.
在《主禱文》裡面.
願人都尊你的名為聖.
願你的國降臨.
願你的旨意行在地上.
而不在天上.
就不純粹是我們自己個人的益處.
或者個人的好處來去求.
四章三節這裡的原文是挺妙的.
新譯本未必能夠譯出來.
四章第三節.
他說你求也得不到.
因為你們動機不良.
和合本就保留了兩個求字在這裡.
他說你們求也得不著.
因為你們妄求.
其實兩個求字.
在原文有一點點分別的.
很 subtle(細膩).
希臘文的 voice(聲音)是有分別的.
就是有 active voice(活動的聲音).
有主動的詞態.
有被動的詞態.
也有中間的詞態.
middle voice(中間的聲音).
那麼你們求也得不到.
那麼這個就是主動了.
就是說你們祈禱.
得不到.
因為什麼呢?.
因為你們妄求.
或者動機不良的去祈求.
那麼後面那個求.
就是中間的 voice(聲音).
就是餵著自己的意思.
在上文下理裡面就告訴我們.
雅各就是針對一些人.
他們去祈禱.
他們希望上帝給他們一些東西.
滿足他們的嫉妒.

$^{241}$滿足他們的私慾.
損毀其他人也好.
或者他去拆毀其他人的生命.
鞏固自己的勢力.
就是一種弱肉強食的一種勢態.
那麼 again(再一次).
在疫情當中.
在當時的困境裡面也一樣.
在困難裡面.
很多時候那些人就….
黑暗面就出來了.
就是….
會真的很弱肉強食.
因為不保住自己.
還保誰呢?.
是不是?.
忍不畏己.
雅各就告訴我們.
我們神的子民怎麼可以這樣呢?.
怎麼可以這樣呢?.
那麼想一想.
你最近的祈禱.
到底是怎麼樣的呢?.
你是為了什麼祈禱呢?.
或者你沒有什麼祈禱.
那麼為什麼沒有什麼祈禱呢?.
是不是因為上帝沒有….
按照你的心意而行呢?.
那麼這個也是一個私慾的表現.
是不是?.
好,那麼我們繼續去切這個菠蘿了.
因為我就沒有那個….
做菠蘿心的工具.
我才發現沒有.
但是不要緊的.
我們不一定要圓形的嘛.
是不是?.
切它薄薄的.
因為這個是給那個….
topping用的.

$^{281}$你預備的時候要調轉.
那麼其實呢.
最後吃的時候就….
調轉才是對的.
那麼 upside down.
其實為什麼會選 upside down 蛋糕呢?.
因為呢.
雅各書第四章呢.
這裡呢.
講到的價值.
天國裡面的價值呢.
真的跟世界的價值是 upside down.
是完全不同的.
世界叫我們.
弱肉強食.
總之有….
有權的就說了算.
沒有權的就被人欺負.
那麼雅各就說.
No No No.
因為耶穌不是這樣子的.
四字大節.
他說.
淫亂的人.
你們不知道與世俗為友.
就是與神為敵嗎?.
所以與世俗為友的.
就成了神的守敵了.
聖經說.
神愛他那安置在我們裡面的靈.
愛到嫉妒的地步.
你們想這話是徒然的嗎?.
但神所賜的恩更大.
所以聖經上說.
神抵擋驕傲的人.
賜恩給謙卑的人.
四字六節.
這裡可能你第一次讀的時候會覺得.
好像有一點蹺口.
特別那句是什麼呢?.

$^{321}$那句話.
神愛他那安置在我們的靈.
愛到嫉妒的地步.
很奇怪.
上面他講到.
教會裡面屬靈的群體.
散居不同的地方.
受壓抑的時候.
他們就會.
彼此的嫉妒.
彼此的傷害大家.
這裡他又用了同一個主題.
又講到嫉妒.
但是他講的嫉妒不是人性黑暗的嫉妒.
而是神的靈.
神將他的靈放在我們裡面.
聖靈.
是一個有嫉妒的靈.
這個很奇怪.
你想到嫉妒的時候.
一想起好像應該是負面的事情.
怎麼會神是嫉妒的呢?.
我覺得嫉妒一般來說都是負面的.
都是黑暗的.
愛是不嫉妒不自誇.
但是這裡講的嫉妒是有一點不同的.
這裡講的嫉妒是一種關係上.
盟約上的嫉妒.
舉個例子很簡單.
人世間.
甚至人間裡面.
都有一些嫉妒是應該.
或者是健康的.
是不是?.
你想像一下一個婚姻.
是不是?.
甚至可能是一個感情.
是不是?.
拍拖的.
給大家的一個承諾.

$^{361}$一個關係裡面.
那是不是一個健康的婚姻.
健康的一個 committed(承諾)的關係.
都是有一些嫉妒的呢?.
是不是?.
就是你的女朋友.
男朋友.
你的伴侶.
都不會到處去搞這個搞那個的.
是不是?.
如果他們這樣的時候.
你不會舒服的.
是不是?.
你會發怒.
那些是意怒來的.
那些是應該的.
所以他這裡講的就是這一種的 jealousy(嫉妒).
這一種的嫉妒.
在聖經裡面不同的地方.
十誡裡面都有講.
就是上帝是一個忌邪的神.
其實那個意思就是說嫉妒的神.
因為他跟我們之間有一個的盟約.
所以上文下理為什麼他去針對呢?.
我們基督徒不可以以世界以世俗為由.
我們要以上帝為我們的朋友.
你記得朋友這個 theme(主題).
這個主題是《雅各書》第二章.
是出現過的.
第二章講什麼呢?.
第二章講信心是有行為的.
他舉了一個例子就是亞伯拉罕.
亞伯拉罕信神.
不只是單單在頭腦上他的行為走出來.
同時亞伯拉罕被稱為上帝的朋友.
是一個 friendship(朋友).
你想清楚.
friendship(朋友).
今天我們講的就是很兒戲的事情.
是不是?.

$^{401}$隨時可以 unfriend(不和)一個人.
是不是?.
或者不跟他做朋友.
是不是?.
不跟他聯絡.
消失.
以前的社會.
講到 friendship(朋友)的時候.
不只是我們那些拍肩膀不簡單.
他們在講盟友.
alliance(聯盟).
就是政治上的盟友.
社會裡面.
關係裡面.
大家彼此去 protect(保護).
彼此去保護.
彼此有一種排他性的關係.
這些才叫做真正的朋友.
所以神跟我們之間有一個盟約.
一個 friendship(朋友).
或者應該這麼說.
alliance(聯盟).
一個同盟.
我們去保護上帝的明星.
上帝也在他的恩典.
在他的愛中來保護我們.
我們彼此之間有一個互相定立的.
毀身的一個盟約.
不可以兩樣.
一腳踏兩船.
所以不要做一個淫亂的人.
所以他為什麼講到.
說不可以世俗為友呢?.
因為如果我們以世俗為友的話.
我們就是出賣了上帝.
是不是?.
我們應該以上帝的國的價值.
上帝的真理.
上帝的公義.
上帝的慈愛.

$^{441}$他的恩典.
為我們的盟友.
是不是?.
不是去靠攏世界的價值.
世界叫我們弱肉強食.
是不是?.
我們去到霸位.
霸權.
用暴力.
我們作為神的子民.
我們不是這樣的.
所以以神為友.
就不是說.
我們跟上帝是好朋友的.
隨隨便便都可以.
不是這個意思.
與上帝為友.
是跟他結盟的意思.
有些以前華人古代的文化.
結義.
就是生與死都不會放棄對方的.
桃源結義就是這個意思.
給我們一個很好的提醒.
與世俗為友.
特別是一些.
我們在世上有權力的人.
有時候也看到一些基督徒高官.
去靠攏世界的權力.
靠攏一些打壓教會.
或者是說方.
去到弄權的一些人.
這是一個很大的引誘.
與世俗為友.
我們自己也是.
當我們安定.
我們算是比其他人的條件好一點的時候.
我們不要漠視世界裡面發生的事情.
而我們就助紂為虐.
在我們說的事情做的事情裡面助紂為虐.
或者在我們的無言沉默的裡面.

$^{481}$我們也助紂為虐.
這就是亞國.
二者亞國.
James the Just.
他告訴我們.
好像先知的語氣.
你不可以與世俗為友.
你要以上帝為你的盟友.
所作所言的.
一定要去推動上帝的國度.
上帝的價值.
好,這個就是.
我們 upside down cake 的下面.
先做下面,做底.
我們只需要將那些菠蘿.
我們就放進去.
好,現在先加上一些cherries.
會不會.
我想我買錯了.
這隻是jam來的.
不過應該也OK的.
如果那些菠蘿是圓形的.
我們放在中間.
那就更加漂亮.
我覺得我做的也很漂亮.
好,那這個topping就搞定了.
就是底先搞定了.
好,現在將cake mix.
我們放在一起.
首先是麵粉.
接著是baking powder.
鹽.
好,我們再加牛油.
第四樣.
第五樣就是糖.
糖.
我怕我的糖已經少了.
我們不用加那麼多.
因為西人的甜品甜到.
太甜了.

$^{521}$ground sugar.
將ground sugar放進去.
然後一隻大的蛋.
我將大的蛋放進去.
雲呢拿.
然後有sour cream.
最後就是菠蘿汁.
好,將這些東西.
加在一起的時候.
我們就可以將它攪勻了.
好像像樣子.
好,我們再看第七至第十節.
他說你們.
應當要馴服神.
去抵擋魔鬼.
魔鬼就逃避你們.
你們應當親近神.
神就親近你們.
罪人要去潔淨你們的手.
三心兩意的人.
要去清潔你們的心.
你們要受苦.
悲哀哭泣.
把歡笑變為上童.
把快樂變為優秀.
你不要在主的面前去謙卑.
祂就會使你們去高聲.
所以這裡講到.
高聲的方法.
不是去毀壞別人的聖源.
或者不是去傷害身邊的人.
以致抬舉自己的地位.
他說在上帝面前.
我們謙卑的時候.
我們就自然會高聲.
不是世界的價值差這麼簡單.
而是世界的價值.
弱肉強食.
那些嘔鬥.
嫉妒.

$^{561}$那種世界觀背後是誰呢?.
就是魔鬼.
我們在上帝的吻友的時候.
我們就以魔鬼為我們的敵人.
這裡第七至第十節.
他說你要馴服神.
要去抵擋魔鬼.
你要去親近神.
神就親近你們.
要去潔淨.
要去清潔你們的心.
要受苦等等.
這裡連續好像打拳一樣.
一拳一拳地打下來.
有個厲害的廚師.
一刀一刀一刀這樣.
快刀斬亂麻的.
但是很清楚地告訴我們.
喂!你醒醒定定.
清清楚楚.
講了整整十節.
講到原來我們人最大的問題.
質度.
就是我們的自高.
我們的ego.
我們的自我.
我們放不下.
別人比你厲害.
你又不喜歡.
你就去質度他.
你想去踩低他.
你都去毀壞他.
這個就出自我們的私慾.
上帝的世界是upside down.
是完全不同我們見到的世界.
好,我們把蛋糕混合.
我們就把它放到下面.
是要平均的.
其實弟兄姊妹.
教會這個群體.

$^{601}$大家走在一起.
其實就好像一個餅一樣.
煮菜都是這樣.
有些東西當然會多一點.
有些東西會少一點.
在教會裡面群體都是一樣.
有些ingredient.
就好像眼睛看上去就多一點.
有些就好像很細微的.
是不是真的不關重要呢?.
不是的.
如果沒有了你又不行.
所以你不去這樣看的時候.
我們很多時候就落入那種彼此質度.
彼此不高興.
你又做多過我.
你好像多一點人去重視.
我又沒有那麼多人重視.
很多時候我們的看法不知不覺地.
其實就是那種世俗的看法.
這裡說什麼?.
說得很對.
這裡說.
弟兄們你們不要互相毀謗.
人若毀謗弟兄.
判斷弟兄就是毀謗律法.
就是判斷律法.
如果你們判斷律法.
就不是實行律法的人.
而且是變成了審判官.
立法的審判的原來只有一位.
那個就是神.
就是上帝自己.
就是那個能夠拯救人.
也能夠毀滅人的神.
你們這樣判斷你的論社.
你是誰呢?.
原來我們這樣不明白上帝的心意.
這個不僅是以世俗為由.
不僅是被魔鬼留有地步.

$^{641}$甚至是我們 playing guard(裝作守護者).
甚至我們當自己是上帝.
當自己.
我就是立法的那個.
我就是判斷審判的那個.
Right?.
群體生活就是需要將我們自己放下.
Right?.
這裡最針對的始終都是我們人的自我.
我們在上帝面前是不是一個謙卑的.
是不是願意謙卑.
降服在上帝自己的安排.
他的角色的安排.
他的時間的安排上面.
如果是這樣的話.
我們就自然在上帝的國度裡面.
就是被升高.
好,我們在等焗爐熱的時候.
應該幾分鐘而已.
我們再看《雅各書》四章最後那段.
13節開始.
他說你們說今天或明天.
我們要到某城去.
在那裡住一年.
做生意賺錢.
其實明天怎樣你們並不知道.
你們的生命是什麼呢.
你們本來是過眼雲煙.
轉瞬之間就消逝了.
你們倒不如說.
主若願意.
我們就可以活著.
作者是或作那是.
但現在你們竟然張狂自誇.
這一切自誇都是邪惡的.
人若知道該行善事.
卻不去行.
這就是他的罪了.
又是說到說話.
但是他的說話就表達了.

$^{681}$他內心世界裡面的慾望.
你可能會想起一些traveling businessman.
但是在以前古代.
所謂這些traveling merchants.
這些到處去不同的地方.
做生意或維生的人.
都是一些上流社會的人.
即是上層裡面1\%或2\%.
有錢的人.
有權力.
有勢的人.
他們才會說出這樣的說話.
或者生命裡面有這樣的安排.
這些是一些上流社會的人.
亞國就像先知那樣警告他們.
針對他們.
說些什麼呢.
他們生命的本相就是過眼雲煙.
好像一口氣.
我們這一口氣.
我們現在經常戴口罩.
我們都感覺到我們一口氣.
這裡給我們一種角度.
即是怎樣去看生命.
生命真的是unpredictable.
是不能夠預測的.
沒有誰想過.
今年開年的時候.
會想到現在.
那個世界會是這樣的狀況.
所以他說應該說什麼呢.
他說主若願意.
即是有上帝的心意許可.
有他的帶領.
有他的開門安排.
你才可以做到那件事.
當然他不是說我們經常口頭禪.
我們主若願意.
我們明天就上學.
主若願意我們明天就結婚.

$^{721}$主若願意我們明天有工作.
而我們生命裡面一個基本的心態.
一個attitude.
好,我們的焗爐已經350度了.
我們可以將蛋糕放進去.
大概應該是30分鐘左右.
就應該成形了.
說回亞哥書第四章結尾.
他當然不是說我們懶懶行.
每件事都是等上帝安排.
人有自己的責任.
我看這裡是更加明確的.
提醒我們需要有一種心態.
不要等明天才行事.
因為如果我們這樣說的時候.
我們好像可以預測到明天會發生什麼事.
很多時候我們不是這樣.
我們覺得不如我們先延遲.
明天先吧.
先事奉上帝明天吧.
我們行善明天吧.
上帝感動我們做的事.
等等.
無論你說什麼都好.
其實下意識上我們就是覺得.
我掌管一切.
我的旨意是行頭的.
不是上帝的旨意行頭.
今天想一想.
神要你該做什麼呢.
請你謙卑的行出來.
不要遲延上帝在你心裡的工作.
好,看看行不行.
拿一支牙籤.
如果出來是清潔的.
是可以的.
好,什麼叫做Pineapple upside down cake呢.
就是要把它反過來.
這樣才像樣.
正如基督徒對世界的觀念.

$^{761}$我們對人生的觀念.
是反過來的.
因為上帝的國和地上的世俗是不同的.
好,四邊都切開了.
我們把這個盤子放上去.
然後就回去了.
最緊張的時刻就到了.
哎呀.
哇,不錯.
菠蘿黏上去了.
It's not bad.
就是這樣.
這個就是叫做Pineapple upside down cake.
好,可以上碟了.
這樣.
先試一試.
嗯,挺香的.
看看找些什麼可以送我這個Pineapple upside down cake.
原來是你們,這麼巧.
我想找些東西,像是wood cream去送我這個Pineapple upside down cake.
不過在這裡見到你.
心理學家最近在疫情裡面說.
如果人們跟冰箱說話.
是沒有問題的.
不過如果冰箱對著它說話.
是有些問題的.
要看心理醫生.
我又沒有去到這個地步.
不過在過去這幾個月.
從一月到現在.
我都很多時候.
跟自己自言自語.
我也想告訴你,我自己跟自己說過些什麼.
喂,你是誰.
不用左望右望了.
我是2020年5月的你.
我是你4個月後的將來.
即是你可以預知未來.
那你可不可以告訴我這幾個月會發生什麼事呢.
當然可以說一下.

$^{801}$不過不能透露太多細節.
籠統地告訴你會發生什麼事.
那你想要good news還是bad news.
good news.
好消息.
good news.
是芝加哥不塞車.
這麼奇怪.
是的,一點都不塞.
芝加哥沒有交通.
還有油價是2元一個gallon.
那你就可以去road trip.
經常去玩.
哈哈哈哈.
那就可以經常去多人多.
到時你就知道了.
到時你就知道了,油便宜不等於去到地方.
另外一個good news.
就是大英姐妹.
或者朋友.
很多人都會work from home.
這樣是多些時間在家.
可能會好些.
那可以是一件好事.
那可不可以說一下教會.
大概會怎樣.
教會都ok.
因為大英姐妹.
找到一些很有創意的方法.
來服務社區.
服務世界.
你一向都這樣祈禱.
那老婆子女會有什麼冬瓜豆腐.
家人.
沒有什麼冬瓜豆腐.
不過.
你都要準備一下.
你女兒去了大學.
她很快會回來.
她不是在大學drop out.

$^{841}$她不是drop out.
不過她很快會回來陪你.
可能陪你到一個地步.
可能你覺得.
嫌見她多得太多.
那麼奇怪?.
你兒子也是,他會經常留在家.
不知道好事還是壞事.
現在新年我有些plan.
港島有plan,我自己生活也有plan.
還有我這個basement也有很多plan.
那這些計劃都sound good.
港島計劃.
或者basement怎樣用.
那是好的.
下次請團契來這裡.
請小組來這裡,多開心.
我怕到時這個basement.
請人來,人們都不理.
那不是很浪費,我已經整好了.
那就不會浪費basement.
basement一定會.
非常實用.
你沒想過.
會怎樣實用.
你穿圍裙.
沒有,我剛剛煮完菜.
你變了.
沒什麼.
你以前不是經常說要試新事物.
要創作.
又不想不傳統.
我告訴你.
你很多機會試.
有些教會的弟兄姊妹跟我說.
今年五月想去以色列.
你想去以色列.
就算你去到也沒有人會跟你去.
你講得好像很嚴重.
無論發生什麼事.

$^{881}$我都希望你身體健康.
如果你身體不健康.
你死都像.
你講的東西這麼不吉利.
什麼事.
又很多事要給你做.
可能比以前還要多.
我不用計劃.
無論發生什麼事.
無論多嚴重.
神自己的時間一定是美善.
縱使上帝的時間.
不一定按照我們的時間表.
希望你喜歡一個比較輕鬆的方法.
告訴你過去幾個月.
確實生命裡面有很多的失控.
好像亞國書所說.
明天明年發生什麼事.
實在不是到我們掌握.
但是亞國書聖經告訴我們.
很清楚.
他說我們的人生.
不是被一些盲目的命運擺佈.
而是每一天.
每時每刻.
都受到這位愛我們的天父.
縱光之父.
所有美善的事物.
都是從他而來.
由這位天父帶領我們.
所以最近我想起.
這首詩歌.
一首比較舊的普通話詩歌.
叫做《神的道路》.
我希望將這首歌成為我們最後的回應.
神的道路 高過人的路.
神的意念 高過人的意念.
他的心裡有難度.
他的時間不錯.
一步一步帶領你前途.

$^{921}$他離大地的更近.
他讓雲在空中漂浮.
是他創造生命啟示.
是他延續教育之父.
身教萬物更啟示.
耐心等待必要跨過.
啊 啊 啊 啊.
耐心等待要跨過.
可以同心低流其道.
上主我們多謝你.
因為我們的人生是由你帶領的.
因為你的本性始終都是美善.
永恆地美善.
求主你繼續帶領我們人生.
縱使在逆勢當中 在逆境裡.
我們實在不知道明天會發生什麼事.
將來不是在我們手裡掌握.
但我們的將來.
每一步 每一時 每一刻.
都在你的安排裡.
祝我們感恩.
讓我們繼續的忠心來跟隨你.
祈禱奉主明求.
母親節快樂.
希望上帝賜給每一位.
特別是母親.
有祝福和平安.
在這個星期裡.
能夠經歷上帝賜給你的恩典.
希望能夠在祈禱會.
或者在我們的主日崇拜.
在網上能夠見到你.
上帝祝福.
\newpage



\section{雅各書}
\label{sec:teCFMiZwZp0}
\textbf{2020/5/17 智慧廚房 vlog 信息 \#5 (雅各書 5) - 陳崇基牧師}
\newline
\newline
連結: \href{https://youtube.com/watch?v=teCFMiZwZp0}{\texttt{https://youtube.com/watch?v=teCFMiZwZp0}} ~~~~ 語音日期: 2020-05-18
\newline
\newline
\hyperref[sec:v2v5pAP5WjA]{\small{< < < PREV SERMON < < <}}
~
\hyperref[sec:index_chronic]{\small{[返順時目]}}
~
\hyperref[sec:index_scriptual]{\small{[返順卷目]}}
~
\hyperref[sec:N2MGZTvagIA]{\small{> > > NEXT SERMON > > >}}
\newline
\newline
$^{1}$各位歡迎回來《智慧廚房》.
我們邊煮邊講.
這一次我們繼續講《雅各書》第五章了.
這一次《智慧廚房》已經是第五集了.
這一次我們先看第五章《雅各書》.
一到第六節.
這裡講你們富有的人.
應當為那將要臨到你們的災禍來哭泣.
哀號.
你們的財物扭壞了.
你們的衣服吸住了.
你們的金銀生鏽.
這鏽要成為空高你們的鐵證.
又要像火一樣吞吃你們的肉.
你們竟在這末後的日子積聚財寶.
漢那工人為你們收割莊稼.
你們竟然剋扣他們的工資.
那工資必要為他們敷冤.
收割者的呼聲已經達到萬軍之主的耳中了.
你們在世上窮奢極侈養肥了自己.
竟不知屠宰的日子到了.
你們把二人定罪殺害.
但他並沒有反抗.
這裡又回到亞國書一個很重要的主題了.
就是對 wealth(富有).
就是富有的人.
有一些審判或者一些很嚴密的警告.
這裡講到當時富有的人.
不僅窮奢極侈.
很有錢.
並且他們是剋扣那些工人的工資.
這個要了解當時的社會.
當時的人.
大部分在社會裡面都是很窮的.
正如我所講.
只有1\%或者少於1\%的人.
是比較有錢的.
而那些人跟平民百姓比較起來.
他們得到的財寶和產業.
是多很多很多倍的.

$^{41}$所以大部分的窮人.
他們當時是一天就出一次工資的.
你聽過耶穌講的一些比喻.
跟那些葡萄園工人.
那些在田地裡面打工的比喻.
你們知道了.
那些人是早上去開工.
幸好有工作.
然後在最後日落之前.
下班的時候就給他們一天的工錢.
所以當時是那些 day wage labor.
day wage labor 是什麼呢?.
就是做一天的工作就出一天的工資.
那麼少了一天的工資.
就代表一家大小就沒東西吃了.
這裡就是那個情況了.
就是說那些有錢的人.
他們連一天的工資.
都不肯給那些工人.
來欺壓他們.
所以這裡的雅各就好像先知的口氣.
你記得我說過的雅各.
二者雅各.
James the Just.
就好像舊約的先知.
叫做 Amos.
Amos.
Amos 就是當時古以色列時期的一個先知.
他是國家裡面的元首.
那些少數的貴族.
囤積著很多財富的.
那幫人去警告他們.
那麼雅各呢.
也是用同一個口氣.
用同一個警告來告訴他們.
說什麼呢?.
他說你們這些人.
你們應該從喜樂變成哭泣和哀號.
因為很快很快上帝要去審判你們了.
你們的金銀銅鐵就會生鏽了.

$^{81}$你們的衣服穿得這麼華麗的時候.
又會被蟲蛀等等.
這些就好像耶穌說的那些口氣.
在地上積聚財寶的人.
上帝就一定會來 keep them accountable(保護他們).
最後會有 judgment day(審判日).
他形容世上奢華極恥的人.
就好像一頭肥牛一樣.
就是你一尾養一尾養他.
那麼他其實呢.
最終的結局是怎麼樣呢?.
就是被屠宰了.
那麼這個是很大的 warning(警告).
一個很大的警告.
當 judgment day(審判日).
審判來到的時候.
上帝就要去問責.
上帝就會施行他的審判.
我們要了解聖經當時的時代.
特別講到經濟的時候.
就需要留意我們現代人的經濟.
跟古代人的經濟的分別.
我們現代人看經濟是很不同的.
因為我們自工業革命等等的發達.
或者 capitalism(資本主義)發達的時候.
那麼自由市場裡面.
我們永遠都覺得資源是無限的.
或者資源會繼續去開發.
會繼續去 produce(生產).
繼續去 multiply(多元化)的.
Right?.
永遠都覺得資源是不會有限的.
在古代的經濟裡面.
只有一個餅去分的.
那麼那些有錢有權有勢的人.
就先將他分了 99\% 先.
他們吃光了.
然後剩下的 1\%.
2\% 就給那些貧苦大眾.
99\% 的人口.

$^{121}$來分他一點點東西.
所以是一個很不均衡.
資源很不足.
很罕有的一個社會的狀況.
所以在聖經裡面.
經常聽到一些叫做仇富的話.
就是那些先知,雅各.
很不妥的那些有錢的人.
因為他們有權.
又有勢.
又有錢.
那麼就來欺壓當時的.
那麼這裡特別講到什麼呢?.
他說不僅他們這麼做是錯的.
他甚至說是末後的日子.
末後的日子是什麼呢?.
新約裡面這個 term(條件)是一個 technical term(技術性的條件).
是耶穌第一次來.
跟耶穌第二次來中間的一段時間.
就是說在聖經的概念裡面.
這一段時間.
現在我們身處的這一段時間.
就是世界歷史裡面最後一段時間.
這個就叫做末後的日子.
last days(最後一天).
之後上帝要做的事情是什麼呢?.
就是耶穌會回來.
然後會有最終的審判.
是不是?.
所以他在講這些富有的人.
不僅他們在倫理上是有邪惡的.
並且他們 doing this thing at the wrong time(做錯了事情).
就是現在末後的日子.
他竟然不知道.
上帝會來審判.
他們逃走的日子.
會突然間來到的.
那麼這個是聖經裡面的末世觀 eschatology(末世觀).
就是告訴我們.
因著現在這個末世.

$^{161}$我們現在的生活應該是怎麼樣的.
今天可能我們賺錢.
或者我們開發一些東西.
做生意等等.
就不會像以前那種看法了.
就是現在可能做生意開發.
是靠你自己的能力.
是一些正途可以做到的事情.
所以我們要做一點所謂 translation(翻譯).
就是這段聖經.
雅各去警告有錢的人的時候.
我們應該怎麼去形容出來呢?.
那就不可以照字搬紙.
因為那個時代文化不同了.
但是那個屬靈的功課其實也在的.
那個信息是什麼呢?.
就是那種物質主義.
再大間一點的房子.
我們再有多一點東西.
我們再有多一點 gear(裝備).
我們再有更加大的.
更加好的東西.
那麼我們心靈就會滿足了.
其實就是警告我們那種物質主義.
那種試探.
就是會侵蝕我們的心靈.
並且會將我們跟主的關係越來越拉遠.
使得我們好像當時富有的人.
我們不知道.
we are doing this at the wrong time(我們在做這件事在錯的時間).
就是上帝的審判來到的時候.
我們怎麼去交賬呢?.
那麼也要提醒我們.
在我們這個比較富足的時代.
我們有沒有去實踐慷慨.
generosity(慷慨).
我們有沒有實踐那種知足.
contentment(慷慨).
就是我們現在有的東西.
其實已經很好了.

$^{201}$那麼這個是一個很大的提醒.
就是聖經裡面由始至終.
從舊約到新約.
都有一個原則.
所有我們的一分一毫.
都是屬於上帝的.
最終的把握權都是上帝的.
我們只不過是一個 bank manager(銀行經理).
就是他放在我們的口袋裡面.
我們也去管理而已.
我們有一種很大的 condition(條件)就是.
我十一房獻了.
那麼其餘的十分之九.
就一定是我的了.
上帝是沒有 say(說)的.
上帝是不需要過問的.
我喜歡怎麼用就怎麼用.
那麼這個呢.
不是聖經裡面的財富的觀念.
這個不是上帝要我們做的管家.
其實百分之一百都是上帝的.
遲早我們會見主的.
我們怎麼交貨呢.
他其中一個問的問題就是.
我們的 worth(值).
所以你說到跟隨主的時候.
做一個門徒.
decipher(記載).
一定會 touch on(觸碰).
我們的錢包.
一定會 touch on(觸碰)我們的錢財.
所以我們真的要想一下.
到底我們跟隨主的時候.
我們有沒有知足.
有沒有實踐慷慨這個原則.
所以真的簡簡單單有一頓飯.
其實真的是我們應該有的.
生活的原則.
是不是有這樣的時候.
不過.

$^{241}$大部分的時間.
我們應該是簡簡單單.
我們用剩了的東西.
或者我們有的財富.
我們怎麼去祝福其他人呢.
一個 art center(藝術中心)的世界觀.
每一個我們跟隨主的人.
我們不斷地需要去學習這件事.
所以今天我們煮什麼呢.
煮一些很簡單的東西.
有多簡單呢.
就是煮蛋了.
這次煮一個早餐.
平時不知道大家早餐吃什麼.
我的早餐是很簡單的.
我喝一杯咖啡.
黑咖啡.
不加任何東西的.
另外呢.
就吃一碗.
這麼大一碗的.
屋苗.
那些的麥片.
就這麼加水.
然後叮它一分鐘.
就是這麼簡單.
我真的很簡單.
一碗屋苗.
然後一杯咖啡.
有時候呢.
我會煎一隻蛋.
弄一隻蛋字.
這麼簡單.
這次呢.
就做一個比較特別的蛋.
想告訴你.
簡單的其實都可以很有趣.
可以很好吃的.
大家可能試過了.
叫做Sauvéé Omelette.

$^{281}$就是非常厚的煎蛋.
那麼.
不過呢.
我就看了不同的食譜.
有些是兩隻蛋.
有些是五隻蛋.
有些是七隻.
有些是十隻的.
就煮給一位的人的.
因為我想簡單嘛.
還有我想將那個食譜呢.
就教成是一隻蛋.
所以需要什麼呢.
需要蛋.
一隻.
那麼這些多了的.
一隻蛋.
然後呢.
我們呢.
需要一點點鹽.
那麼也呢.
是需要一些糖.
那麼到煎的時候呢.
我們需要呢.
牛油或者呢.
橄欖油.
那麼我選了橄欖油.
為什麼呢.
等一下告訴你.
蛋黃蛋白分開.
那麼原來就是這樣子而已.
就是這麼簡單.
倒來倒去.
那麼將蛋黃打發的時候呢.
我們加一點點鹽就夠了.
那麼.
吃多鹹呢.
你自己啦.
你自己決定.
那麼將蛋黃呢.

$^{321}$先打發掉.
然後蛋白呢.
我們再加呢.
糖呢.
一點點.
那麼然後呢.
我們就要.
將蛋白.
去打發.
打發掉.
打發到它變泡.
當然如果你用電的mixer.
就會快一點啦.
是不是.
但是這次我們考一下自己耐性.
為什麼要考耐性呢.
因為《雅各書》呢.
五章第七至十一節.
講什麼呢.
就是講忍耐了.
你記得第一章也是講到忍耐.
試探.
試煉來到的時候我們需要忍耐.
七至十一節這裡講什麼呢.
他說.
所以弟兄們你們應當忍耐.
直到主來下.
這裡呢.
他頭六節.
我想常文下里大概是警告當時一些.
可能甚至不是教會的權貴.
先知經常都做這件事情的.
去挑戰那些.
就算是教會以外的權貴.
那麼來到這裡第七節呢.
他就講弟兄們了.
就是弟兄姊妹.
講到教會裡面.
那些大部分可能是受欺壓的那些人.
他說.

$^{361}$所以弟兄們你們應當要忍耐.
直到主來.
看啊.
農夫等候著地裡寶貴的出產.
為他忍耐.
直到獲得秋林春雨.
你們也應當忍耐.
好像農夫那樣忍耐.
堅定自己的心.
因為主在內只近了.
弟兄們不要彼此的抱怨.
免得你們受審判.
看啊.
審判的主已經站在門前了.
弟兄們你們應當效法奉主的名說話的先知.
以他們為受苦忍耐的榜樣.
看啊.
那些忍耐的人.
我們稱他們為有福的.
你們聽過約伯的忍耐.
也看見了主賜給他的結果.
知道主是滿有憐憫和仁慈的.
那麼我們可能會覺得.
雅各無緣無故轉了主題.
說忍耐.
之前就在說那些.
就是有權有勢的人.
他們是欺壓人等等.
其實這些都是猶太人的背景.
猶太人很多這些章節.
這些經文都是說到.
二人.
就是窮人受苦的時候.
他們可以做些什麼.
當時的窮人當然沒有勢力了.
他們就要去仰望主了.
那麼這裡主題的變化.
其實就是這些猶太人文學經常有的.
審判講完了.
然後就勸勉那些等待著的人.

$^{401}$他們要繼續堅定他們的心.
譬如詩篇37篇.
詩篇73篇.
如果你回去看一下的話.
你知道這兩篇詩篇都是講到.
不要心懷不平.
你要繼續去等候主.
因為有永恆的結局.
你才會有盼望.
那麼.
縱使你做不到什麼.
你還是可以等候主.
因為主永遠都是公義的.
那麼這裡又要講回.
我們不要斷章取義.
要留意當時的社會背景.
那些窮人沒有任何勢力的.
他們沒有任何的 opinion(意見)可以講出來.
所以聖經叫他們去忍耐.
不是叫他們去死忍.
因為當時他們確實真的沒有這樣的能力.
去反抗等等.
也不能夠照字搬字說.
哦!你今天就是總之就是死忍.
我們什麼都不要出聲.
我們就沉默不作聲.
不是這樣的意思.
而是當時的人確實真的沒有這樣的能力.
那麼這裡就讓我們想到.
就是聖經裡面有兩方面要我們留意的.
就是.
社會裡面不公不義的事情.
我們是需要努力的來改善.
努力來實行社會公義.
social justice(社會公義).
去幫助別人,是不是?.
同一時間也告訴我們.
on this side of heaven(這方向的天空).
在新天新地來到之前.
我們都不會有一個完美的世界的.

$^{441}$但是兩件事情都要做.
我們一定要 keep(保持)做兩件事情 intention(意圖).
一方面我們努力來改善.
另一方面我們等候主的真正公義.
祂新天新地來到.
如果你說一條直線.
然後第一次基督來的時候.
祂死而復活.
成就了救人.
那麼神國就在地上展開了.
到祂最後回來的時候.
就是 second coming(第二次來臨).
就是主會再來的時候.
between first and second coming(第一次和第二次來臨).
中間這一段就是所謂 last day(最後一天).
末世的日子.
上帝的國度已經來到展開了.
但是還沒有完全的實現.
所以在這個狀態之下.
就有所謂 already(已經).
but not yet(但還沒有).
已經展開了.
但是還沒有完全實現的一個張力.
我們一方面繼續將福音傳開.
一方面繼續為主作供.
讓這個世界成為一個更美的世界.
但同時我們有一個張力就是.
我們知道怎麼做都好.
永遠都不會完美的.
但不代表我們不做.
因為我們知道上帝的國度已經展開了.
並且將會一定實現的.
所以我們現在做的事情才會有意義.
他這裡就用了一個比喻.
就去比喻我們是農夫.
農夫每一天努力的去耕耘.
而且他知道.
氣候,溫度,土壤.
那些需要成長的成分.
都不是他自己能夠完全控制得到的.

$^{481}$但不代表他不做.
正因為他做.
但是他同時也需要上主來看顧.
他才能夠有修成的.
我們今天在地上的工作也是一樣.
我們儘管為上帝的真理,他的福音.
他的天國的價值.
我們去努力去做事情.
但同時我們也知道.
所有的土壤,外在的因素.
不是我們掌握的.
唯有上帝去恩待.
這樣才會成事,會有修成的一天.
但是因為上帝是滿有憐憫和滿有人志的.
他是那一位公義的主.
所以我們現在做的事情.
一切都是有價值的.
所以無論我們現在的情況是怎麼樣.
我們的 circumstances(情況)是怎麼樣.
我們都需要繼續去等候.
等候將來的結局.
就是主宰回來的那個結局.
或者上帝去成就他自己心意的那個結局.
可能是一兩年後.
可能是十年,二十年後.
可能是超過了你一生.
我們也不會看到那個成就的.
但是我們照樣默默地去耕耘.
這個就是基督徒應有的人生觀.
我們攀望於主.
同時我們在地約天.
將上帝的國利帶來.
另外他提到一個榜樣.
就是約伯阿扎.
聖經裡面的約伯記的約伯.
如果你回去看.
舊約裡面的約伯記.
你不是很明白這麼多的東西.
不是 get(得到)到很多約伯的忍耐等等.
不過正如我所說.

$^{521}$雅各的 context(角辯).
是猶太人的 context(角辯).
我們之前不是提過.
雅各很多時候會引用一些當時猶太人.
耳熟能詳的故事嗎?.
那麼在兩約中間.
有一個創作的文學就叫做 testament of Job.
上次我們提過 testament of twelve patriarchs.
就是十二族長的遺訓.
那麼也有一本叫做約伯的遺訓.
猶太人想像他臨死之前.
他跟子女說什麼.
就是說他自己一生怎麼忍耐.
來等候上帝的國度.
等候終局來到.
所以他就用當時他們讀者.
了解的故事.
告訴他們.
喂,約伯也是這樣而已.
我們這些家傳戶曉的故事.
我們就應該要效法約伯.
不斷地去忍耐的時候.
那麼你就是屬於上帝的國.
你就是一個有福的人.
這個就好像耶穌說的登山寶粉一樣.
事實上雅各經常都引用耶穌他自己的 half brother.
所說的話.
他最熟悉的.
跟他生活了三十年.
那麼雅各應該很熟悉他哥哥耶穌的價值觀.
或者他的教導.
或者他人生的智慧.
所以這裡說.
就像登山寶粉一樣.
那些忍耐的人.
甚至在考驗當中.
在迫害裡面受忍耐的人.
他們都是有福的.
都是屬於天國的.
那麼他說這裡忍耐.

$^{561}$不是死忍.
他說是看見什麼呢?.
主次級他們的結局.
看見上主給他們的 ending(結局).
或者英文有時候會譯作 climax(結局).
或者 ESV.
他會譯作 purpose(目的).
那麼.
從第一章到第五章最後.
都是說一樣的.
他說我們要成長的時候.
人生在世我們會有難處.
會有逆境.
但是我們要 keep the end in mind(留意結局).
就是我們要知道上帝有他自己的 purpose(目的).
縱使我們現在未必明白.
將來可能也未必明白.
可能到天家才會明白的.
但是我們要 keep the purpose in mind(留意結局).
[音樂].
那麼你看到我這裡.
將這些.
蛋白一邊漏一邊發的時候.
它就開始變泡了.
第五章十二節到十六節.
他說我們弟兄們最要緊的就是.
不可起誓.
不可指著天起誓也不可指著地起誓.
任何的誓都要禁絕.
你們的話是就說是.
不是就說不是.
免得你們落在審判當中.
你們中間有人受苦嗎?.
他就應該禱告.
有人心情愉快嗎?.
他就應該歌頌.
你們中間有人患病嗎?.
他就應該請教會的長老來.
讓他們奉主的名.
為他沒有祈禱.

$^{601}$出於信心的祈禱.
可以使病人康復.
主必叫他起來.
他若犯了罪也必蒙赦免.
所以你們應當彼此認罪.
互相代求.
這樣你們就可以痊癒.
二人祈禱所發的力量.
是大有功效的.
其實這裡預設了什麼呢?.
就是一個基督徒要成長的時候.
他一定要需要在社群.
一個關係的網絡當中來成長.
每一個基督徒都不是一個孤立的個體.
那麼這個呢.
給我們有很大的反省.
就是我們現在疫情裡面.
我們就知道了.
社會學家叫做 social recession(社會沮喪).
社交上的大蕭條.
我們想見人也見不到.
現在我們一直在家隔離幾個星期.
都已經不是很好過了.
人是被造的時候.
上帝給我們的形象.
是有那個群體性在裡面的.
你內向也好.
外向也好.
都是需要群聚的.
這裡就講到.
祈禱.
講到大家走在一起.
你們中間受苦的.
就是大家彼此來禱告.
就算有愉快的時候.
就是好景的時候.
我們都需要彼此歌頌.
你們中間有人患病的.
他就要請教會的長老.
來到.

$^{641}$那麼他就奉主的名來抹油來祈禱.
那麼是要出於信心的祈禱.
使他康復.
主叫他起來.
如果他犯罪的時候就要彼此.
去饒恕.
大家能夠去痊癒.
大家能夠得到平安.
得到康復.
那麼這裡或者可以講一講.
他說什麼是抹油呢?.
所以為什麼我選了橄欖油.
橄欖油是什麼呢?.
當時的古代的社會.
就跟我們現代社會很不同的.
那麼.
現代的社會.
我們有很多不同的藥物.
是不是?.
古代的社會.
很多時候藥物的選擇是很少的.
橄欖油是其中一個.
那麼.
所以很多時候他們受傷.
或者皮膚乾燥等等.
特別是巴勒斯坦地.
當然我們不能夠照字搬子.
有些人會覺得.
我們現在也照抹油.
做一個祝福.
那麼我又不會極力反對這件事情.
不過.
這是一個解經的原則.
我們不可以照字搬子.
正如新約書上經常叫我們彼此親嘴問安.
現在疫情當中你更加不會這麼做.
就算沒有疫情.
我們也不會這麼做.
因為他們那種中東的做法.
不能夠照字搬子.

$^{681}$下來我們今天現代的社會裡面做法的.
是不是?.
所以呢.
其實這裡講到什麼呢?.
就是互相的 support(支持).
互相的得到醫治.
這個對於一些散居的猶太人.
歷世歷代因為他們失去了自己的國家.
尤其重要的.
他們就彼此 support(支持).
寫信去 support(支持)大家.
就是這個群體.
寫信給這個群體.
support(支持)他們.
或者呢.
他們散居在不同的地方.
雖然他們在城市裡面是分開了住.
但是每逢到安息日的時候.
他們就會去河邊.
小人就去河邊.
他們就去聚集.
來一起禱告.
他們有會堂的話.
多一些人的.
他們就走在一起.
來去讀經.
來去禱告.
聽道等等.
鼓勵大家.
那麼這個就是那個群體性的緊要.
那麼為什麼講到散居呢?.
因為記得第一章第一節.
雅各怎麼形容這些猶太背景的基督徒呢?.
他說散居的十二支派.
他說你們才是真正的以色列的群體.
相信耶穌的理所當然的群體.
散居也好.
你也需要走在一起.
來彼此鼓勵彼此建立.
所以這裡.

$^{721}$十二至十六節經常講.
你們中間.
你們.
以及互相.
彼此.
所以十二節這裡開頭講什麼呢?.
他說不要起誓.
其實就是在講那個群體.
那個真理的緊要.
那麼又解釋一下.
他說憑信心的禱告.
是大有功效的.
那麼很多時候我們可能會執著這句.
我們說.
總之我們加多一點 faith(信心).
為什麼不照我們的心意去做呢?.
為什麼呢?.
因為可能我們沒有信心.
所以我們加多一點信心.
那麼這個也是大家要看那個 context(觀點).
你也知道雅各書是智慧文學來的.
看箴言等等.
智慧文學要留意.
很多時候他在講的是 general truth(總真相).
不是 guarantee(保證).
是.
普遍來說是這樣.
不是保證.
你一定這麼做.
就是一條 formula(方程式).
你可以控制上帝.
或者行你喜歡行的事情.
只要你拍著胸口夠信心就可以了.
不是這樣的意思.
所以有個上文下理你就明白了.
這裡是在鼓勵我們.
人生裡面的智慧是什麼呢?.
我們需要緊靠主.
要在上主的面前.
我們有信心來祈求.

$^{761}$是.
信心 again(再一次).
不是我們自信有多大.
信心就是我們將我們的焦點放在一位信實的上帝身上.
所以我們有網上祈禱會.
我真的鼓勵你.
你會準時來.
我們網上已經方便了.
是不是?.
而且也是同樣有效.
是不是?.
那麼.
我們在禱告裡面去大家紀念.
以及大家在禱告當中去牧羊.
每一個大英之母.
那麼現在我們可以將這些東西放在一起了.
我們發到現在已經起泡了.
那麼我們現在可以將它們.
是.
混在一起.
然後我們就將蛋黃放在一起.
將這些.
蛋黃和蛋白.
混在一起.
現在將蛋白蛋黃.
混在一起了.
那麼我們只需要.
把鍋子輕熱.
然後我們加一點點油.
橄欖油.
在古代社會.
巴勒斯坦社會.
很有醫學的價值.
好,那麼我們差不多熱到這裡了.
那麼我將它.
放進去.
小一點的鍋子應該會比較好看.
但是我只有一隻比較大的.
應該也可以的.
其實廚房有很多智慧的.

$^{801}$有些東西真的不能急.
你煮一些東西.
或者煮一些更加爛的東西.
其實真的要等的.
那麼反蛋要有些忍耐的.
這裡說二人祈禱.
雅各不是在說二人是完美的人.
其實是我們屬於神的人.
所以為什麼他說以利亞.
跟我們的性情是一樣的.
他很有趣的.
就是以利亞有很多祈禱的神蹟.
死人復活.
為什麼特別說這裡呢.
就是祈禱三年六個月不下雨.
祈禱又接著又下雨.
那麼我相信他的意思就是.
在以利亞一生裡面.
最考驗到他忍耐祈禱的就是這個神蹟.
就是一定要堅持下去.
同樣我們也是一樣.
我們在教會裡面禱告.
我們都需要堅持到底.
還可以的.
樣子可能要再 improve(改善)一下.
要去改善一下.
但是還是挺泡的.
好,這次第一次試驗.
很滑的.
早晨流流.
這個是一個非常好的早餐.
最後兩次他說.
你若有人是偏離了真道的.
如果有人使他能夠回頭.
你就應該知道.
使罪人從祈禱上面轉回的.
就是拯救了他的靈魂.
給他有永生.
脫離了死亡.
也都遮蓋了很多的罪惡.

$^{841}$這裡講到我們教會.
一個群體.
是一個 restoring community(復原社會).
就是人來到我們教會的時候.
我們盡量將福音帶給他.
將上帝的命志介紹給他.
教導他到底主是怎麼樣的.
我們最終的目的就是 restore(復原)罪人.
以致他能夠回來主的身邊.
這個就是我們一直努力要去建立教會的一個目的.
所以作為教會.
他就提醒我們要不斷地去到彼此的建立.
建立這個群體使本地的教會的 community(社會).
能夠繼續的 strong(強)下去.
雖然在疫情裡面我們都遇見不同的困難.
就是我們現在 worship(慈禱)又不可以實體.
我們的祈禱會又在網上.
但是在不同的情況.
不同的挑戰之下.
我們不要忘記.
這個就是我們最主要的目的.
是一個 restoring community(復原社會).
一個更新的群體.
所以這個星期給大家的功課.
你可以做些什麼.
能夠去建立我們這個群體.
我挺喜歡有一個團體最近做的.
就是每個人 assign(指定)另外一個人.
去鼓勵他.
然後說一些事情.
是實在的.
很實在的.
很 specific(具體)的.
來去鼓勵那位弟兄或者姐妹.
assign(指定)給他.
他就去寫一段 text(文字)給他.
然後鼓勵他繼續去中心的地方事奉.
這些話.
好像很老土.
也好像很黑板.

$^{881}$你收到當然得到依次得到鼓勵.
但同樣你寫給別人的時候.
你的言語出了去的時候.
你建立到人的生命.
這個就是雅各書所說的.
我們能夠將人帶到主的面前.
我們結束了這個雅各書和智慧廚房五個星期的系列.
我們希望能夠在主的話語當中.
我們得到鼓勵.
在考驗的裡面.
我們信心能夠成長.
並且在考驗當中.
在疫情裡面.
我們能夠繼續在這裡投靠神.
在這裡事奉.
我覺得最好就是在廚房裡面結束這個系列.
唱一首胡影詩歌.
這首胡影詩歌叫做《投靠》.
讓我們成為一個敬畏主.
有智慧的人.
投靠 你翅膀陰霞.
在你光中 不必驚怕.
唯獨敬畏你的方可得以靠.
別再恣意來轉向你.
投靠 你翅膀陰霞.
在你光中 不必驚怕.
唯獨敬畏你的方可得以靠.
別再恣意來轉向你.
洪水即將泛濫之時.
是我主 所著為王.
願天歡喜.
願地上快樂.
如英戰士上台.
等候我主的必不守軌.
願你甘霖與我.
只用一切.
好 我們重新低頭祈禱.
上主我們求你幫助我們.
在考驗當中 逆勢的裡面.
我們能夠繼續依靠你.

$^{921}$讓我們恆久的忍耐等候.
等候你自己的來到.
等候你的國度.
並且我們禱告.
在我們的所言所行當中.
讓我們能夠繼續.
在你光中 依靠你.
禱告在我們的所言所行當中.
來行你的旨意.
願你的教會.
我們每一個都能夠在.
每一個地方繼續榮耀神.
我們做一個投靠你的人.
做一個有智慧的人.
我們這樣祈禱.
奉主耶穌基督的名字.
全體委員會審議文件.
\newpage



\section{}
\label{sec:N2MGZTvagIA}
\textbf{4 十一少被賣記:原生家庭(四)}
\newline
\newline
連結: \href{https://youtube.com/watch?v=N2MGZTvagIA}{\texttt{https://youtube.com/watch?v=N2MGZTvagIA}} ~~~~ 語音日期: 2020-05-26
\newline
\newline
\hyperref[sec:teCFMiZwZp0]{\small{< < < PREV SERMON < < <}}
~
\hyperref[sec:index_chronic]{\small{[返順時目]}}
~
\hyperref[sec:index_scriptual]{\small{[返順卷目]}}
~
\hyperref[sec:12OrdZOBpSQ]{\small{> > > NEXT SERMON > > >}}
\newline
\newline
$^{1}$今次我再和大家從原生家庭的問題.
去了解負面影響所帶來的後果.
每一個人都有原生家庭.
原生家庭可以帶給我們正面.
但同時亦可以帶給我們負面的影響.
你父你子.
有這樣的父親,就會有這樣的兒子.
在故事當中主角約瑟原生家庭負面影響.
所帶來的家庭問題.
最主要找出根源原因.
就是從約瑟的祖宗三代開始起.
約瑟的祖宗三代.
是猶太祖父亞伯拉罕.
祖父以撒和父親雅各.
約瑟的三代祖宗.
同時都有一個共同點.
就是口齒薄皮.
他們都是偏愛某一個的子女.
我們先看看太祖父亞伯拉罕.
在《創世記》第十六章.
亞伯拉罕因為妻子撒萊不能生育.
於是接受妻子的建議.
娶侍女哈格為妾.
結果亞伯拉罕有了第一個兒子.
再看二十一章.
多年後撒萊雖然年老.
但因為得到神的賜福.
撒萊終於和亞伯拉罕生下了以撒.
以撒是亞伯拉罕和撒萊年老時所生.
亦是他們的第一個兒子.
所以他們對以撒寵愛猶加.
因為這個緣故.
他們興褻了早年時.
和妾哈格所生的長子.
不僅如此.
撒萊還故意撫慰哈格和其兒子.
最後還把哈格和其兒子趕出家門.
亞伯拉罕的家庭問題.
都出自於亞伯拉罕的厚齒薄皮.
而以撒呢?.

$^{41}$以撒自出於母胎後.
可說是得天獨厚.
因為以撒深得父母親的疼愛.
例化道理神.
有這樣的父親.
便會有這樣的兒子.
亞伯拉罕的厚齒薄皮帶來的家庭問題.
傳到下一代.
去到以撒的身上.
《創世記》的二十五章.
二十到二十八節.
那裡說到.
利伯加和以撒身著一對上輩胎.
大兒子爾素.
小兒子雅各.
父親以撒偏愛著大兒子爾素.
而母親利伯加偏愛著小兒子雅各.
我們看看《創世記》的二十七章.
母親利伯加因為偏愛著雅各的緣故.
便用手段幫雅各從父親的手中.
騙取了長子的名分和得到祝福.
結果雅各很成功奪取了長子的名分.
雅各用欺騙的手段從父親的手中.
騙取了長子的名分和得到祝福.
雖然雅各很成功騙取了長子的名分和得到祝福.
不過雅各卻要付出很大的代價.
雅各的一生再沒有機會和母親利伯加有團聚的機會.
由亞伯拉罕的厚齒薄皮家庭問題.
就像一種胃串病一樣.
傳到下一代以撒的身上.
再由以撒的身上傳到下一代雅各的身上.
雅各一夫娶四個妻妾.
一共生了十二個兒子.
但雅各只是偏愛弱室.
只為弱室親手做採衣.
父親的厚齒薄皮.
造成頭十個兒子至少缺乏安全感.
至於弱室.
因為年少無知.
又被父親寵愛.

$^{81}$所以至少意氣風發.
這樣就更加深了在家庭中兄弟的不和.
在雅各的家裡竟出現了.
撒底的念頭.
賣帝的行動.
弟兄姊妹.
弱室對上三代所發生的家庭問題.
就可以足以證明.
原生家庭的問題.
會帶來很大負面影響.
結果更加會造成家庭的悲劇.
在弱室的家庭裡造成弱室被賣的原因.
都是因為父親厚齒薄皮.
子女得不到父母親完全的愛.
以致兄弟不和.
甚至嚴重到引起殺忌.
雖然人的罪性的確造成很多家庭的問題.
連這些家庭的問題.
好像傳染病一樣.
一代一代地傳下去.
但是去到弱室的一代.
因為有神的介入.
使弱室能夠避過了兄長的殺害.
雖然弱室要被賣到埃及.
但是弱室在埃及大大經歷到神的賜福.
分享完弱室被賣的原因.
下一次我會和大家再一次探討一下.
弱室被賣之後所遇到的經歷.
多謝您收睇時局新聞,再會!.
\newpage



\section{}
\label{sec:12OrdZOBpSQ}
\textbf{5 十一少被賣記:有神同在(一)}
\newline
\newline
連結: \href{https://youtube.com/watch?v=12OrdZOBpSQ}{\texttt{https://youtube.com/watch?v=12OrdZOBpSQ}} ~~~~ 語音日期: 2020-05-26
\newline
\newline
\hyperref[sec:N2MGZTvagIA]{\small{< < < PREV SERMON < < <}}
~
\hyperref[sec:index_chronic]{\small{[返順時目]}}
~
\hyperref[sec:index_scriptual]{\small{[返順卷目]}}
~
\hyperref[sec:uQ6L96b0_II]{\small{> > > NEXT SERMON > > >}}
\newline
\newline
$^{1}$約瑟被賣後的幾次大經歷.
第一次經歷是約瑟被賣後.
去到法老的內臣.
波提弗的家裡做奴隸.
在《創世記》第39章.
一到第六節.
那裡說到約瑟因為有神的同在.
雖然約瑟被賣後.
他沒有被賣.
但他有一個祖先.
就是法老.
那裡說到約瑟因為有神的同在.
雖然約瑟被賣到埃及做奴隸.
但非常受到主人的器重.
當時約瑟的主人.
就是埃及人法老的內臣波提弗.
波提弗經常與法老接觸.
應該是一個舉足輕重.
有影響力的大官.
而約瑟當時.
只不過是一個被賣到埃及人的奴隸.
奴隸的身份非常卑微.
可以說是一無所有.
就算連他自己的命都不屬於他自己.
但神一直與約瑟同在.
神賜福也保守約瑟.
於是約瑟百事順利.
弟兄姊妹.
之前我們都有對約瑟被賣的前因探討過.
他被兄長恨而所帶來的嚴重後果.
竟然是做到手足之間可以畏道權,畏道名.
來反目成仇,甚至屈辱相殘.
人的一生有很多東西的確是無可選擇.
不單是出生之前,我們是無可選擇.
即使我們出生之後.
人的一生經歷都一樣是無可選擇.
就好像故事裡的主角約瑟.
他被賣到埃及.
雖然被賣到埃及了做奴隸.
但神卻依然與約瑟同在.

$^{41}$於是約瑟在主人的眼前蒙恩.
還得到主人的肯定.
約瑟起初是以一個奴隸的身份來被賣.
後來他就再被賣到去法老的奈神.
Bolv提佛的家裡做奴隸.
是埃及人的家.
與埃及人同工作,同生活.
由於神與約瑟的同在.
約瑟就在主人的眼前蒙恩.
Bolv提佛認為自己的家業能得到恩福.
全是因為約瑟有神的同在.
於是約瑟就深得住主人的賞識.
主人還升約瑟做總管.
還把家裡各樣大小的事都交由約瑟一個人來管理.
大家留意一下.
Bolv提佛與約瑟兩個人的背景.
地位身份都是非常懸殊.
若非約瑟有神的同在.
以Bolv提佛的背景.
根本不會賞識約瑟.
更加不會重用和信任他.
試想一下.
當初約瑟能夠免受到輕掌,殺害.
雖然仍然賣到埃及來做奴隸.
但他是由一個非埃及人的奴隸身份.
可以晉升為一個總管.
管理著埃及法老內臣的一切家業.
可想而知約瑟不單是大能不死.
神還用到約瑟在埃及裡.
大大經歷到與神同在所得到的福氣.
由約瑟被賣到埃及來做奴隸.
到升為總管這段期間.
約瑟由一個青年人.
變成一個壯年人.
在這段日子裡.
約瑟先後經歷過有兩次被陷害.
在約瑟的青年期時.
被掌興賣到埃及來做奴隸.
在約瑟的中年時.
他被主母陷害被主人鎖在監獄裡.

$^{81}$約瑟當年被害是因為原生家庭問題.
父親只是偏愛約瑟.
父親忽略了其他十個兒子的感受和需要.
加上約瑟年少無知.
不懂人情世故.
又經常向父親提高掌興的罪行.
結果造成被兄長賣到埃及來做奴隸.
但神一直與約瑟同在.
神賜福給約瑟.
約瑟因禍而得福.
最初是由一個奴隸的身份.
慢慢來到晉升成為總管.
約瑟當年要離開自己的家鄉.
去到另一個地方來生存.
是約瑟自由的選擇.
還是他不由自主呢?.
電影姐妹.
走入人生的路上的確有很多時候是不由自主的.
你有沒有像約瑟一樣.
不單要離開自己的原生家庭.
還要一無所有.
獨自來到一個陌生地.
沒有家人,親戚,朋友,工作.
甚至連自己的性命都不屬於自己.
下一次我會再和大家深入探討.
究竟約瑟流落異鄉後有什麼經歷.
\newpage



\section{}
\label{sec:uQ6L96b0_II}
\textbf{6 十一少被賣記:有神同在(二)}
\newline
\newline
連結: \href{https://youtube.com/watch?v=uQ6L96b0_II}{\texttt{https://youtube.com/watch?v=uQ6L96b0\_II}} ~~~~ 語音日期: 2020-06-16
\newline
\newline
\hyperref[sec:12OrdZOBpSQ]{\small{< < < PREV SERMON < < <}}
~
\hyperref[sec:index_chronic]{\small{[返順時目]}}
~
\hyperref[sec:index_scriptual]{\small{[返順卷目]}}
~
\hyperref[sec:sSMekSYd36k]{\small{> > > NEXT SERMON > > >}}
\newline
\newline
$^{1}$在上一次講到約瑟雖然是不由自主要留在異鄉.
但因為有神的同在.
究竟有神的同在帶來約瑟有什麼後果呢?.
在故事當中的主角約瑟一直以來都有神的同在.
所以當他被賣之後.
可以由一個埃及人的奴隸升為埃及人的總管.
在約瑟壯年的時期他做了一次避害.
不過這次的避害不是因為兄長的疾到.
而這次約瑟的避害是因為約瑟收到試探.
這個經歷是約瑟做埃及人主管之後所發生的.
約瑟當時是受到他的主母色友試探.
在《創世記》第39章7-20節.
那裡講到由於當時約瑟還未婚.
而且他長得非常英俊又瀟灑又能幹.
並得到主人的器重和信任.
以一個非埃及人的身份.
可以由奴隸晉升為總管.
管理著埃及人家裡所有的人和事.
於是就令到主人的妻子對約瑟產生了貪念.
於是這個主人的妻子就去色誘他.
弟兄姊妹.
人受試探都外乎幾樣.
權,名,利和色.
當時的約瑟深得埃及人的器重.
約瑟在家裡做主管是一人之下萬人之上.
魔鬼如果要成功試探約瑟.
是毋須用到權和利.
魔鬼要成功試探約瑟.
必須用到色誘這種方法.
加上約瑟年輕又未婚.
用色誘了引誘他試探他.
的確是最容易令約瑟跌倒.
不過,試探的結果.
竟然是約瑟拒絕了主母對他的色誘.
約瑟成功拒絕主母主要有兩個原因.
第一,約瑟得主人的器重和信任.
約瑟知道不可以做出任何不忠於主人的事.
第二,約瑟拒絕主母的原因.
亦是最主要的原因.
究竟是什麼呢?.

$^{41}$使約瑟能夠勝過試探.
主要原因是約瑟心裡有神.
約瑟知道神願你一直與祂同在.
約瑟知道神是耶和神,是以色列神.
神是聖潔的,所以約瑟不敢犯罪來得罪神.
還是怎會呢?.
神是無所不知,無所不在,亦是無所不能.
我們應該時時刻刻監察著自己的心思和意念.
不要以為,不要說出口,不要將它化為行動.
神就會不知,不理,不管,不介入.
神是創造的主,主人是有主權.
主人是按著自己的時間和方法.
讓我們來經歷時年.
同時我們也要知道.
魔鬼的試探是事先沒有需要來通知我們.
或得到我們的同意.
試探是隨時隨地都可以發生.
魔鬼不會輕易放棄對我們的攻擊和引誘.
只要我們是魔鬼試探的對象.
魔鬼就會不斷試探我們.
即使今天我們成功拒絕魔鬼的引誘.
但不代表魔鬼會放棄對我們下一次的引誘.
在故事中,這個主角約瑟.
就是天天受到矚目的色誘.
大家留意的是魔鬼的試探.
是天天都會發生.
而且是約瑟生活和工作的地方.
可以說是約瑟避無可避.
今天我們又有沒有在生活或工作的地方.
經常受到試探呢?.
故事看到這裡.
可能你會問約瑟為何不積極反抗.
例如向主人告發主母色誘他.
為何約瑟只是消極.
你都避開主母就算了?.
如果當初約瑟真的有積極反抗的話.
例如向主人投訴主母對他的色誘.
或許就不會讓主母有機會你無告約瑟.
最後還要被主人囚禁.
弟兄姊妹,很多時候,當有事發生.

$^{81}$如果當時沒有積極,正視,採取適當行動處理的話.
可能真的會造成日後更多一發不可收拾的後果.
尤其是人與人之間的誤會.
可能只是出自言語之間.
又或者是文化,家庭背景等因素.
所以要避免誤會,最好儘早釐清.
至於約瑟為何沒有南克正視和處理當時主母色誘他.
下一次,我們會繼續和大家再深入探討.
\newpage



\section{}
\label{sec:sSMekSYd36k}
\textbf{7 十一少被賣記:有神同在(三)}
\newline
\newline
連結: \href{https://youtube.com/watch?v=sSMekSYd36k}{\texttt{https://youtube.com/watch?v=sSMekSYd36k}} ~~~~ 語音日期: 2020-06-16
\newline
\newline
\hyperref[sec:uQ6L96b0_II]{\small{< < < PREV SERMON < < <}}
~
\hyperref[sec:index_chronic]{\small{[返順時目]}}
~
\hyperref[sec:index_scriptual]{\small{[返順卷目]}}
~
\hyperref[sec:y63dV5dKyc4]{\small{> > > NEXT SERMON > > >}}
\newline
\newline
$^{1}$今次我試同大家分析一下.
為何藥室沒有立刻正視和處理當時主母對她的色誘呢?.
我個人認為有以下幾個原因.
首先,當年的藥室幾乎被兄長殺死.
最後就被兄長賣到埃及做奴隸.
他的遭遇對一個當時年青人來說是一個相當大的打擊.
一直以來,藥室在親父母的愛和保護之下成長.
加上他年少無知,經常在父親面前告發兄長.
結果因為被兄長質盜,要經歷被殺的危機和被賣的結果.
從此之後,要離開自己的親生父母,留落異鄉成為奴隸.
雖然藥室一直有神的同在,使他在埃及的主人眼前蒙恩.
可以由一個奴隸升為總管.
不過,對當時藥室只是一個年青人時就遭遇這些事情.
對他來說的確造成心靈上很大的創傷.
創傷似乎未因為時間的過去而得到醫治.
所以造成藥室的心理上留下了很大的陰影.
例如當時藥室遇到主母色友他,這件事為例.
從他的處理方法我們可以知道.
藥室因為知道有神的同在,所以他不敢犯罪得罪神.
於是他就有勇氣每一次都成功地拒絕主母對他的色友.
但他卻沒有勇氣,向主人告發主母的罪行.
導致藥室沒有勇氣向主人告發主母的罪行原因?.
有可能是因為他害怕,一旦告發了主母的罪行.
結果就像當年一樣,自己因為年少無知.
因為心得父母的偏愛,在父親面前告發兄長.
帶來的結果就是被兄長殺和賣.
中年的藥室再不是一個無知的年青人.
他有被自己的親人殺和賣的經歷.
藥室在過去的經歷中,新心靈都被受傷害.
從此之後,亦成為他很大的心理上的障礙.
過去的日子,要離開父母.
被迫留在一個非自己的出生地,生活,工作.
在這段日子,他需要克服很多生活上對他的挑戰.
相信大家曾經做過新移民.
剛來到一個新的地方.
要落地生根的確需要面對很多生活上和工作上的挑戰.
有沒有試過不敢向警方舉報.
或者不敢出庭做證人,指證一些罪犯?.
甚至有時候,連發聲表達自己的不滿都不敢.
在故事中的藥室.

$^{41}$不能夠立即處理主母對他的色誘帶來的後果.
就是讓這個主母有腦收成奴的機會.
看19到20節,那裡說到.
主人信妻子的話,於是就將藥室囚禁.
藥室在一次的經歷到一無所有.
由一個總管,變成一個囚犯.
甚至有生命的威脅危機.
弟兄姊妹,人生會有不同的經歷.
大多數的經歷都不可以有預知.
就好像藥室今次所遇到的.
並不是他自己預料之事.
更加不是他自己所選擇.
同時亦沒有能力來改變.
故事說到這裡,大家會不會認同.
「戲如人生」?.
如果人的一生真是一部戲.
人的一生,由開始到結局.
都只可以跟著一本不能改寫的劇本.
無論你喜歡,還是不喜歡這部劇本.
你都沒有自由來再選擇.
你都沒有辦法來改變.
不過,信主重生後的基督徒.
生命就不一樣.
重生後的基督徒是有新的生命.
新的生命是屬於靈性的生命.
靈性的生命是由聖靈所掌管.
在聖靈掌管下,有神的同在.
亦有神的賜福,帶領和保守.
即使在人生路上再遇到有多少難過的關口.
因為有神的同在.
一切所遭遇的都會化險為夷,因禍得福.
\newpage



\section{}
\label{sec:y63dV5dKyc4}
\textbf{10 十一少被賣記:因禍得福－下(一)}
\newline
\newline
連結: \href{https://youtube.com/watch?v=y63dV5dKyc4}{\texttt{https://youtube.com/watch?v=y63dV5dKyc4}} ~~~~ 語音日期: 2020-06-18
\newline
\newline
\hyperref[sec:sSMekSYd36k]{\small{< < < PREV SERMON < < <}}
~
\hyperref[sec:index_chronic]{\small{[返順時目]}}
~
\hyperref[sec:index_scriptual]{\small{[返順卷目]}}
~
\hyperref[sec:LqhzJpn7Fyo]{\small{> > > NEXT SERMON > > >}}
\newline
\newline
$^{1}$上一節提到,酒精得到王的恩典.
不僅醉得黑到赤面.
總可以復職,能夠在王的身邊.
再慰到王來祭酒和試酒.
是王所信任的高官.
不過,酒精復職後.
她忘記了曾經在監獄裡一起生活.
既得到照顧和關心她的約瑟.
酒精亦忘記了約瑟曾經為她來解夢.
約瑟在監獄裡一直期望酒精.
能夠復職後可以紀念她.
但酒精卻忘記了約瑟.
雖然約瑟被酒精忘記.
但神卻沒有忘記約瑟.
神只是延遲了約瑟兩年後.
才被釋放,其實是有神的心意.
我們先看看41章,一到第七節.
埃及王所發的夢.
試酒兩年後,王就發夢見到自己站在河邊.
另外有七隻肥牛.
從河裡上來,同時站在河邊.
接著還有七隻瘦牛.
從河裡上來,同時站在河邊.
後來這七隻瘦牛.
就把這七隻肥牛吃了.
再看八到第十六節.
第二章,王的夢境.
在第十六節.
第二天早晨,法老因為對所發的夢.
心裡感到非常不安.
於是按照埃及文化.
叫博士和術士為他解夢.
結果呢?.
博士和術士都不能為王解夢.
在這時候,酒精記起當年.
為他解夢的年青人.
一個被拐,被埃及當奴隸.
又被主母誣告.
甚至最後被主人囚的年青人.
約瑟.

$^{41}$酒精同時記起.
自己忘記了對約瑟的紀念.
寫道約瑟要多留在監獄裡空等兩年.
在人的眼中,約瑟這兩年.
被囚在監獄裡可以說是沒有希望.
白白消磨心智,浪費時間浪費生命.
在神的眼中,他並不是這樣看.
神是要約瑟經歷兩年.
被酒精忘記的日子.
神的用意是想試煉約瑟.
希望能夠加強約瑟的信心.
不僅如此,約瑟也能等到酒精記得他.
酒精還將他推薦給威王來解夢.
我們看看16,25和32節.
當約瑟為了王解夢的時候.
約瑟強調.
「不在乎我,上帝,別像平安的話回答法老」.
約瑟對法老說.
「法老的夢乃是一個」.
「上帝已將所有作的事都指示法老」.
「王所發的夢是因上帝命定者事,而且速速成就」.
從約瑟這番話可以聽得出.
約瑟其實是一個不會奪取神榮耀的人.
他會將榮耀歸神.
同時他亦很有信心.
在埃及王的面前將解夢的能力傳歸他所信的以色列神.
約瑟亦向王坦言.
「將上帝要做的事都指示王」.
弟兄姊妹,我們有沒有像約瑟一樣呢?.
當從神的裡面我們得到領受.
或者得到人的稱讚的時候.
又或者我們有從神而來的恩賜,能力,恩典甚至福分.
我們有沒有心存感恩的心.
來將榮耀頌讚歸給父神呢?.
當你向未信的朋友或者.
有其他宗教信的人面前將榮耀歸神的時候.
其實就是為主作見證.
亦都是向未信的人傳福音呢?.
我們跳去看看25到39支.
那裡說到約瑟替王解夢.

$^{81}$還向王指出.
王所發的夢其實是有一個意思的.
就是埃及人將會有七年的豐收.
亦都有七年的饑荒.
約瑟強調這件事的發生.
是有神的意思在裡面.
約瑟同時亦提出需要選擇一個有智慧的人.
派他管理埃及.
結果約瑟就被王選擇選.
王委任約瑟為埃及首相.
大家可能會覺得很奇怪.
為何會選擇約瑟?.
選約瑟因為約瑟不單能解夢.
同時亦能指出解決埃及饑荒的方法.
最重要的是.
王他很相信神是與約瑟同在.
亦認同約瑟有神同在的不只是王一個.
而其餘的大神同時都認同.
可以說王和各大神都能夠印證.
印證約瑟能夠解夢的能力是從神而來的.
而且他們相信夢的預告.
就是神必定將要發生的事.
如果想知道發生甚麼事的話.
下一次我會繼續和大家探討.
\newpage



\section{}
\label{sec:LqhzJpn7Fyo}
\textbf{8 十一少被賣記:因禍得福－上(一)}
\newline
\newline
連結: \href{https://youtube.com/watch?v=LqhzJpn7Fyo}{\texttt{https://youtube.com/watch?v=LqhzJpn7Fyo}} ~~~~ 語音日期: 2020-06-18
\newline
\newline
\hyperref[sec:y63dV5dKyc4]{\small{< < < PREV SERMON < < <}}
~
\hyperref[sec:index_chronic]{\small{[返順時目]}}
~
\hyperref[sec:index_scriptual]{\small{[返順卷目]}}
~
\hyperref[sec:YbuM_vsj2Sc]{\small{> > > NEXT SERMON > > >}}
\newline
\newline
$^{1}$上一次和大家分享過.
如果人的一生真是一部戲.
人的一生由開始到結局.
都只可以是跟著一本不能改寫的劇本.
無論你喜歡還是不喜歡這本劇本.
你都沒有自由再選擇.
也沒有能力再改.
不過,信主重生後的基督徒.
生命就不再一樣了.
重生後的基督徒是有新的生命.
新的生命是屬於靈性的生命.
靈性的生命是由聖靈的掌管.
有神的同在.
有神同在的生命.
就有神的賜福,保守和看顧.
即使在人生的路上遇到有多少難過的關口.
因為有神的同在.
一切所遭遇的都會化險為夷.
因禍得福.
在《十一少被賣記》的故事裡.
我們可以從主角藥室.
他所遭遇的經歷.
雖然藥室要再次了受害.
但同樣神是依然與藥室的同在.
神賜福給藥室.
藥室雖然經歷了試探.
被誣告.
也被囚在監獄裡.
但到最後反而他可以因禍得福.
我們先看看《創世記》第40章.
一到第八節.
故事當中出現了兩個人物角色.
第一個是酒正.
另一個是善長.
酒正是負責製酒.
是一位製酒師.
是專門為埃及王製酒和試飲.
另一位是善長.
是負責做餅.
是一位做餅師.

$^{41}$專為埃及王製餅.
他們經常與埃及王的接觸.
據資料顯示.
他們是深得著王的信任高官.
不過.
經文並沒有交代.
為何他們會被埃及王拘留在監獄裡.
他們被困在監獄的這個地方.
就是護衛長府裡的監獄.
正正是當時藥室被囚禁的地方.
加上護衛長指定藥室.
要在監獄裡好好照顧他們.
於是藥室就有機會認識他們.
並且在監獄裡他們可以一起生活.
再看看第五到第八節.
在經文裡並沒有直接交代.
他們在監獄裡相識的日子究竟有多久.
只是交代到其中有一晚.
酒正和膳長同時做夢.
做夢後的第二天早上.
藥室看到他們悶悶不樂.
情緒低落.
於是藥室關心他們.
問他們究竟發生了什麼事.
大家留意一下.
藥室當時的反應和行動.
當藥室看到酒正和膳長悶悶不樂.
情緒低落時.
藥室關心他們.
問他們究竟發生了什麼事.
藥室之所以有這樣的反應和行動.
可以證實藥室是一個非常有責任心的人.
因為當時神要藥室在護衛長眼前蒙恩的緣故.
護衛長派藥室管理所有囚犯.
特別是對身邊的酒正和膳長.
藥室聽護衛長的吩咐.
很忠心地去辦事.
另外藥室其實是一個非常關心人的人.
所以當藥室看到酒正和膳長悶悶不樂.
情緒低落時.

$^{81}$藥室就問他們究竟發生了什麼事.
大家不要忘記.
當時藥室雖然在護衛長眼前蒙恩.
派他管理監獄裡所有囚犯.
但當時藥室依然是囚犯身份.
而且無期地來判刑.
在一般人眼中.
藥室將來沒有前途和希望.
藥室在監獄裡.
藥室的肉身固然失去了自由.
但在監獄裡.
藥室的心靈並沒有失去自由.
究竟為什麼呢?.
下一次我會帶大家在這方面再探討一下.
究竟藥室在監獄裡.
為什麼藥室的心靈並沒有失去自由呢?.
\newpage



\section{}
\label{sec:YbuM_vsj2Sc}
\textbf{9 十一少被賣記:因禍得福－上(二)}
\newline
\newline
連結: \href{https://youtube.com/watch?v=YbuM_vsj2Sc}{\texttt{https://youtube.com/watch?v=YbuM\_vsj2Sc}} ~~~~ 語音日期: 2020-06-18
\newline
\newline
\hyperref[sec:LqhzJpn7Fyo]{\small{< < < PREV SERMON < < <}}
~
\hyperref[sec:index_chronic]{\small{[返順時目]}}
~
\hyperref[sec:index_scriptual]{\small{[返順卷目]}}
~
\hyperref[sec:dNRqxYjeS0o]{\small{> > > NEXT SERMON > > >}}
\newline
\newline
$^{1}$[音樂].
今次我再繼續來和大家探討一下.
藥室如何關心酒精和腺脹.
藥室是非常細心留意到酒精和腺脹的表情.
於是就主動用行動來關心.
問候一下酒精和腺脹.
結果酒精和腺脹很放心.
將自己的難題都和藥室分享.
當時藥室和酒精和腺脹最大的差別.
除了文化和語言之外.
其實最大的差異就在信仰方面.
和藥室同在的是以色列的神.
而酒精和腺脹所信的是埃及的日神.
當時要酒精和腺脹完全接受藥室的想法和提議.
又或者接受藥室為他們解夢並不是一件容易的事.
加上大家要明白的是埃及人的文化.
埃及人認為夢是埃及人的神向著埃及人預告.
將來必定要發生的事.
所以酒精和腺脹極想有人為他們解夢.
因為他們很相信.
亦很想知道將要發生甚麼事.
尤其在監獄裡的日子.
對他們當時來說是毫無安全感.
對他們對將來來說亦毫無盼望.
在這種情況下他們更加想知道.
夢見究竟代表甚麼.
將要又是發生甚麼事.
在經文的九道第十一節.
那裡說到酒精夢見在自己面前有一棵葡萄樹.
在樹上有三根枝子.
好像發芽了,開了花.
而上頭的葡萄都已經成熟.
當時黃的杯在酒精的手裡.
於是酒精就將葡萄握在黃的杯子裡.
還將杯子遞在黃的手中.
這就是酒精所夢見的事.
究竟夢見的是預告甚麼呢?.
再看十二道十三節.
約瑟替酒精解夢.
約瑟解釋那三根枝子其實代表了三天的時間.

$^{41}$原來三天之內黃會讓酒精出監.
而且酒精還可以復職.
為黃再次製酒和試酒.
再看十四道十五節.
那裡說到約瑟解夢後.
向酒精求紀念和試因.
既然約瑟相信酒精三天之內可以復職.
約瑟就很希望酒精能夠為約瑟.
向埃及王求情.
讓約瑟可以早些防監.
不過為何約瑟還要將自己的身世.
和不幸的遭遇都向酒精分享呢?.
約瑟當時分享自己的身世和不幸的遭遇.
目的是希望酒精能紀念他.
這樣的話約瑟就可以透過酒精.
在黃的眼前蒙恩.
我估計約瑟當時將自己的遭遇.
向酒精分享還有一個原因.
他相信和很想加深酒精對自己的印象.
這樣酒精就可以紀念他.
於是約瑟在酒精面前強調自己是一個希伯來人.
因為被拐賣到埃及當奴隸.
又被主母誣告.
又被主人囚禁.
再看16-17節.
是善長所發的夢.
善長見到約瑟解酒精的夢解得這麼好.
於是就讓約瑟為他解夢.
善長當時夢見自己頭頂上有三籃白餅.
而籃子的裡面.
還有法老王敲各樣的食物.
還有飛鳥來吃.
在他頭上籃子裡面的食物.
大家有否留意.
善長肯放心讓約瑟解夢其實是有原因的.
善長認為約瑟解夢解得好.
因為酒精有好的結局.
可以在三天之內出獄.
並且可以復職.
於是就加強了善長的信心.

$^{81}$讓約瑟為他解夢.
但結果呢?.
看18-19節.
說到約瑟為善長解的夢.
原來三個籃子代表了三天的時間.
在三天之內.
王必斬善長的頭.
還將善長的頭掛在木頭上.
還有飛鳥來吃善長身上的肉.
看完最後20-23節.
那裡說到到了第三天.
因為是王的生辰的緣故.
王就設宴.
將酒精和善長從監獄中召出來.
結果酒精被王來復職.
為王再次祭酒和侍酒.
但善長卻是被王處死.
酒精和善長的結局.
完全應驗了約瑟當日所說的一切.
如果想知道.
當約瑟所解的夢.
應驗之後會有什麼事情發生.
下一次我會繼續和大家.
深入的來探討.
\newpage



\section{}
\label{sec:dNRqxYjeS0o}
\textbf{Kingdom Values 10 - Sam Johnson - A Conversation on Race}
\newline
\newline
連結: \href{https://youtube.com/watch?v=dNRqxYjeS0o}{\texttt{https://youtube.com/watch?v=dNRqxYjeS0o}} ~~~~ 語音日期: 2020-06-22
\newline
\newline
\hyperref[sec:YbuM_vsj2Sc]{\small{< < < PREV SERMON < < <}}
~
\hyperref[sec:index_chronic]{\small{[返順時目]}}
~
\hyperref[sec:index_scriptual]{\small{[返順卷目]}}
~
\hyperref[sec:qRHTClh4NmQ]{\small{> > > NEXT SERMON > > >}}
\newline
\newline
$^{1}$Family, it's good to be with you again..
And as we've been having this conversation over the past few weeks and talking about.
all the stuff that's going on in our world, I've been trying to point you to some voices.
in the church and some voices in the African American community that can speak to these.
issues really well..
And so today I wanted to sit down with a couple of my really good friends..
And this is Wade Norman and Eric Cummings..
And they're good friends of mine and they're men who love Jesus and men who are serving.
Jesus with their lives..
But they're also black men who've grown up in Chicago and been here their whole lives..
And so these things that are happening right now in the news that we're all thinking about.
and discussing and trying to understand how to respond to, these events are realities.
of their experience that they've dealt with for a long time..
And so I really want you to be able to hear their voice and hear their hearts and hear.
what God has done in their lives and how God has taught them and hear how God wants to.
teach us right now and the things that God would have to say to us..
So guys, to start, why don't you just share..
I know these moments feel different for everybody..
And for you guys, these moments are a part of your experience and have been..
And so if you could just share what that feels like a little bit, some of the emotions they.
have to navigate, some of the pain that's involved, some of what that is..
You..
Yeah, you go..
Yeah, for me, it's like, it was a constant reminder of, in the Christian sense, this.
place is not my home..
But then also like historically, knowing like, you know, the scars of racism in this country,.
it's like, you know, this country wasn't necessarily built for me or..
And so it's, you know, it's like stuff like this happens..
And it's like, there's a part of you that's like really shocked or like, you know, in.
tears about it, but then it's also like not surprised, like you almost get numb to it.
to some degree when it's like, you've been kind of living that way..
But yeah, a lot of it is a lot of it is really complicated now, I think, with so much like.
I was saying to somebody like, we're at a time now where like, everybody's voice can.
be elevated..
You know, social media has made this platform where now people can speak and share about.
injustices or things or issues that are happening..
But what is created is like, yeah, everybody has maybe the equal opportunity to share or.
say things about something..
But it's not necessarily equally valued or like equally worthy of listening to..
And so it's like this, yeah, this hard part where like social media has done really well.

$^{41}$at like, making us like, more aware of these things..
But it's like, at the same time, it's also made it kind of easier to almost ignore it.
if you want to..
So it's yeah, so it's hard..
It's a lot of different emotions..
But what I find also just being an African American man is expressing emotion is very.
difficult anyway..
Like it's one of those things that can, you know, that's kind of shunned within our own.
community..
We have to be strong for our families, we have to be strong for, you know, our wives,.
our kids..
But then strong for those who, you know, we don't know if they have ill intent toward.
us..
You have to, you know, you can't really show weakness..
And so that part is very hard because allowing, you know, stuff to build up, you know, in.
a lot of ways it can poison you if you don't have an outlet..
And thankfully, you know, I found Jesus..
And so like, I feel like even if I can't express an emotion around a bunch of people, I can.
always express it to God, you know, and know that He hears me, you know, that He finds.
my voice valuable..
But I can't imagine, you know, my brothers who don't know Christ, you know, I mean, it's.
like torture in a lot of ways..
And like just to piggyback off of what Wade was saying, the key word he said that really.
resonated with me was "no.".
So I mean, we've been through this so many times..
I know it's been the times where it's like, I say, you know what, that's another killing..
I'm like, it's like you said, we're so used to it..
It's like we see it on social media and you just go pass it..
But this one with George Floyd, I can say personally, this one hit me extra hard because.
a few months before, somebody, I just want to say I was, somebody called the police on.
me and they're like, once you experience it, it's like you can start putting your feet,.
putting yourself in their shoes..
And it's like, wow..
And you're going through, you're trying to fight off that numbness..
You're fighting off anger, you're fighting off fear, you're scared..
And you got all of these emotions you're trying to figure out, but still try to go to work..
Yeah..
And it's like, you're still trying to do a nine to five and still keep it professional..
And it's like, it's hard because it's like, you want to grieve..
Like you were saying before, like as black boys, we was told, like, you're not supposed.

$^{81}$to show your, your emotion..
You're not supposed to cry..
But it's like, what are you supposed to, like, how am I supposed to like share these emotions?.
How am I supposed to feel?.
So it's like, I'm learning how to share my emotions more, but at the same time, it's.
like, I don't want to show too much..
Like, it's like a double edged sword..
Like you show too much emotion, you're that angry black man..
Or you're an angry black person..
If you do show it on other sides, like, oh, you're too soft..
And so it's like, we're struggling to figure out what the mean, like, what should we do.
to like show how we feel about everything..
So it's the first part is always numb..
It's like, when I remember the first day when George Floyd died, I looked at the video for.
like 10 seconds and I said, I don't want to look at this today..
But then it kept on coming back up and I was like, okay, let me sit down and like watch.
this video..
And then all the emotions started coming in..
So it's like, you go through it every day, but it's like just to see it again and again,.
it's like it hurts..
So, because you think it could have been me..
Yeah..
So..
One of the things we've been talking about a lot of different things and responding..
And one of the things we talked about was in Romans 12 and the verse that says that.
we're to weep or those who weep..
And so that's one of the things we've been focusing on grieving because we know that.
that is this process..
It's grief for the African American community..
It's grief..
It's mourning, it's weeping..
And we've talked about how it's not just grief over this incident..
Obviously, it's the whole history of injustice, particularly in our country..
It's grief over that..
Yeah..
And one of the things that we've talked about in that is how believers and people in the.
church who are outside the African American community, how they can respond in a right.
and genuine way..
We see obviously a lot of different responses in the world right now, but how can we do.
that well and do it genuinely and do it in a way that's truly supportive and loving..

$^{121}$And so that's something I would love to hear you guys thoughts on how other believers can.
support you guys and your community and where you're going through..
The first step, I will say, is to become a listener, to actually come and listen to someone.
share how they can help..
Because we don't need people to come in to be a savior for us..
That's having a problem now..
But we need you to come in and listen and hear how people are feeling through this..
And over the last few weeks, I've actually had people come to me and do that..
And it's like, wow, like you actually want to sit down and listen and see, like, I know.
that we live in a broken world..
I had one person actually come out and tell me that he said that I've been racist and.
it's because the way I've been raised and what can I do or how can I help?.
So even I had to even encourage him at that moment, like you saying this and acknowledging.
that you have sinned or this sin is still in your life..
It's like you're ready to make a change..
But also when you make this change, you have to be able to listen to someone else..
So it's all you have to be able to humble yourself to listen to your brother and want.
to learn and actually, yeah, just shut up..
That's the thing, just shut up when you're talking..
Because sometimes what I've noticed, like people want to share and it's like you're.
still fighting on not listening..
And it's like one challenge I did, like an exercise with somebody, like when you are.
the exercise is you, somebody say one thing and you have to start a hand, you have to.
start your sentence with that last letter of that last sentence..
And they didn't understand why I was doing it..
And then after a while, we kept on doing it and they noticed like, okay, like what was.
that for?.
And I was like, you had to listen to everything I said to see how you're going to start your.
thing..
So you have to be able to sit and listen to everything the person is sharing, not just.
the one point you want to talk about, but they share a lot of points they want to talk.
about..
So it's a learning process and it's going to be hard..
That's the thing, it's going to be hard..
It's not going to change overnight..
Like you do it for one day, it's not going to be the same the second day..
It might change with whoever you talk with..
So it's the first two steps is acknowledging where you're wrong or where you're saying.
it is..
And then the second part is being humble enough to sit down and listen to your brother, to.

$^{161}$your African American brothers, actually share how they feel about the situation..
Yeah..
That's awesome..
Yeah, no, I think you guys are on the right track of what the Romans 12 of like grieving.
alongside your brothers..
I think, yeah, if I would add anything to listening to be lamenting alongside..
And really to go into what Eric's saying is like creating an environment of empathy and.
transparency where your African American friends do feel like they can express themselves without.
repercussions or feeling blindsided or like letting their guard down..
I mean, honestly, that's one of the biggest things we need, a place to be who we are..
And it's, yeah, we have to put on so many faces when we step out into the world..
I was watching this one pastor talk about even like because of the Ahmaud Arbery thing,.
he was like, "In my mind, when I go outside, I've already played the mental game of what.
can I wear when I go on this run so that people know that's what I'm doing and there's no.
issue at all or there's no mistake that he's out there running, guys..
He's not robbing anybody..
He's not trying to do anything.".
And it made me remember like something I did or something I still do actually..
And I never really associated with race necessarily..
I just associated with, oh, it's just a common courtesy to do this..
So when I go into a store, well, one, I never go into a store that I'm not willing to buy.
anything in..
And so the next thing is if I go into a store and if I'm going to be in there longer than.
10 or 15 minutes, I have to buy something in there..
So whether it's a soda, stick of gum, otherwise there's this attitude of like, "I don't want.
people to think I'm in here shoplifting..
I want people to know I'm in this store just browsing, looking for something, didn't find.
it, but I'm not trying to get over on anybody.".
And it's, yeah, I don't know..
Hearing that pastor made me think there's so many areas where the African-American has.
to do that and put on this face of like professionalism sometimes, or maybe you do have to be the.
angry black guy because you're in a situation that calls for that..
But it's almost always a mask of who you are..
It's never, you never, you honestly, unless you maybe go to counseling or have opportunities.
to talk with people that have been through the process and have given you opportunities.
and made that safe environment to express yourself..
I know there's some African-American men that never had that chance, and it plays out in.
different ways, whether that's domestic abuse or heavy alcoholism, or it's gambling, all.
these other sins are really deep..
But in a lot of ways, they're surfacy to what some of the issue is..

$^{201}$It's just like, "I can't be a man..
I can't be myself.".
And so, yeah, so lamenting is definitely one..
And then trying to create, yeah, creating a space..
And I think, yeah, the way you create that space is with empathy and transparency..
People that are transparent, that is one of the most contagious things, I think, especially.
for believers..
Every Bible study that I love going to is because men were honest and transparent about.
the real issues that we deal with as believers..
So I think, yeah, I think those are good ways to do it..
- That's good..
One of the things that hit me as you guys were speaking about that is how some of those.
things you're describing are like basic gospel values, basic values that should be in the.
church, that Koinonia type of fellowship..
And I think, I mean, that directs us to the hope of the gospel, obviously, but it also.
directs us to the failure of the church to nurture those things..
It's sad that we haven't done a better job of building those spaces and building them.
for everybody..
And that points an error of our own failing in the church, and that's really sad..
So I'm also interested to hear from you guys as we go forward how believers, and again,.
especially believers who are outside the African-American community, can really continue to be involved.
and engaged and supportive..
Because obviously there's a lot of attention right now, and there's a lot of response,.
and that's good..
There's a lot of positive response, which is really good, but it's easy for those things.
to just be a moment and to just be kind of an event, and then they fade away..
And in a few weeks or a few months, people forget about them and then go on..
And then the realities don't change..
And in the church, we want to continue, by God's grace, to be living out the gospel and.
seeking God's justice and doing those things..
So I'm interested to hear from you guys what are some ways that people can continue to.
be involved and to respond..
That's actually my biggest fear..
That people....
Oh, what?.
Like everybody's protesting..
Yeah, what's next?.
And it's like, okay, what are you going to do next?.
Everybody's on the protest, but when a protesting dies down, like you said, what's your next.
move?.

$^{241}$I was just on a Bible study when we was talking about this..
It starts with talking about community..
What was that guy named?.
Bon Hopper..
Oh, Dietrich Bon Hopper..
Dietrich Bon Hopper..
He was talking about....
It was a quote saying, "If you have a dream for your community, you'll die, but if everybody.
else is in it, if you have a community, you'll thrive.".
So we were talking about community..
You need to create that community of not just people that look like you, but people that.
doesn't look like you, with different races, different age groups, so you'll be able to.
grow and learn from that..
And so I think they were college students..
Best thing about college students, they're able to....
There's different events on your campuses..
Go to those different events and start able....
I won't say put yourself out there, but take a chance to start talking to other people.
that doesn't look like you..
You're going to fail at first..
It's basically like that, because you're stepping out your comfort zone, you're coming into.
somebody's comfort zone..
Everybody's uncomfortable, but after a while, they see that you're not doing it because.
of what's going on, and it's really coming from your heart..
People can tell the difference..
Once they see the difference, like, "Okay, I'm going to let down my guard just a little.
bit, but I'm going to still keep my distance..
I'm going to still keep it up.".
But it's not a....
Like I said before, it's not an overnight change..
If you're serious about being the change, nobody said it's going to start overnight..
It's long suffering..
That means suffering long..
But going through the changes, going through the hurt, and through the journey..
I was talking to somebody last night, like, we always want to get to the end of the journey,.
but the whole experience and the journey is what creates you, or what makes who you are..
So when you're working towards that, it's going to be hard..
But I mean, if you walk with God and you faith, He's going to provide for you..
But if it's just for hidden knowledge or pride, it's just going to die..
It's not going to work..

$^{281}$Like, it's a true heart change..
So I hope I answered that question..
I think you answered it well..
Yeah, that's good..
It made me think of something else that might be helpful is one of the big things, I think,.
in our country and why a lot of racism or just prejudice in general is able to continue.
like generally is because we love categories, right?.
We love grouping..
You think about the news, right?.
They had the looters, the rioters, the peaceful protesters..
It was all these different categories..
And honestly, the African-American experience is pretty used to that..
We are used to carrying the weight of our race all the time..
And so one of the most helpful things is to treat each black friend you have as an individual..
You know, like even me and Eric, we are very similar, but we're also really different,.
like distinct..
And it made me think of this one training I went to, and it was like, "Which black are.
you?".
And it's like, so even within black, there's categories..
And this may be helpful for you, but there's the ABC..
So the A black is assimilated..
On the streets, we will call you a sellout..
This is the person that is basically joined with the majority culture..
Most of their friends are probably white..
Most of the stuff they like is like pop culture type things..
Then you have the B black, which is like bicultural..
So this is the person that's pretty much able to walk in both worlds..
They recognize that they're black and are able to be black around blacks, but they also.
are able to kind of turn on, "What does it mean when I'm around people that aren't black?".
And then there's the C black, which is like contextualized or culturally black..
So if you were like, "Oh, how do I spot this person out?".
This is the person that has the Malcolm $\times$shirt on..
This is the person that is never going to probably come to a church that's mostly white..
This is the person that's at all of the events that Eric was mentioning on college campus..
If it's a black event, they're probably there..
They probably have peace signs and things like that..
And they are very aware of their race and are very active and vocal..
And so one thing that could be helpful is discovering who are the ones around you, because.
all of us interact with these things differently on an individual basis, but even in those.
groups we interact differently..

$^{321}$All this is general, so none of this obviously can stick on anybody..
Let's just say, for example, with all the stuff that's been happening, the assimilated.
might, let's say all of us say the exact same phrase..
We say, "Man, we watch the news and we say, 'This is messed up.'".
I, by culture, say, "This is messed up.".
And assimilated blacks say, "This is messed up.".
And then let's just act like Eric is a C black culture..
And he say, "This is messed up.".
Well, we all said the same phrase, but what did we mean when we said it?.
I say, "Man, this is messed up..
I can't believe the media is showing all these people looting and rioting and stuff.".
And then the assimilated black says, "Wait, what do you mean?".
"No, I'm talking about the way they're treating us police officers, blah, blah, blah.".
And then the culture black is like, "I just can't believe they didn't riot all these other.
places or do all these other places.".
Now, once again, that is general, but yeah, even within the black community, we're all.
coming at this from different angles, depending on how we were raised..
And then what our connections is to our own race..
Some of us, honestly, the bicultural, we have an identity crisis, to be honest..
We can walk in both worlds, but we feel like we belong to neither..
We're not black enough for the contextualized friends..
We're not white enough for our assimilated friends..
And each group probably feels some type of animosity toward others..
It's an everyday struggle..
It's a battle..
It is..
It's a battle..
It is..
I remember being at day camp and almost having to prove my blackness to somebody that came..
They thought I lived in Beverly..
It was something about that..
They thought I listened to....
I did listen to this Hillsong, but I was like....
I mean, there's a lot of people that listen to Hillsong..
I like Hillsong..
They make good music..
But it is funny..
It is like this, "Oh, you have to prove, 'Oh, no, I'm black..
I'm black.'".
Or whatever..

$^{361}$But one of the things Eric said that made me think when you said community that is missing.
in all of this is, yeah, we want to run to solutions..
And a pastor was talking about this..
He said, "Don't be so quick to run to a solution that you forget to run to relationships.".
He was like, "Because solutions can be a quick fix, and you forget about it.".
We get to 2030, and 2020 is 10 years past, and it's not on our mind anymore..
He said, "At the core of it all, do you care about people?".
And it just made me think people is the problem with all of this..
If we're honest, you remove people, and there's none of these things..
There's none of these issues..
But you put people who are sinful, have a sinful condition in a situation with each.
other, and sin is bound to happen..
And one of the things is not only people, but opportunities for power..
And so power is this thing where it only means something when people are involved..
You can't be powerful by yourself..
One of my favorite conversations is when he's having his disciples, he's like, "In the world,.
these people in high leadership, they lord it over each other..
They lord it over those that are lower than them.".
He said, "But that won't be the way with you.".
He's like, "Actually, you are going to go to the bottom.".
He's like, "Everybody is racing to this top, and lording it over each other.".
He's like, "You are going to race to the bottom.".
He says, "If you want to be great, you got to serve everybody.".
And what we see in this country, obviously the reversal of what Christ said, if you want.
to be great, this is who you have to step on, basically..
And it's people who were created and made in the image of God..
So yeah, I think the other thing too is to try to think what you would do if it was your.
people..
That's one of the biggest things that has always gotten me..
What if this was my brother that was killed?.
What if this happened to my sister?.
I think sometimes when you empathize in that way, you don't even need to ask a person what.
to do..
You just do it..
And that's when you learn that, yeah, we may be a different color, but we all want the.
same thing..
Another helpful thing is to understand just because things may be equal doesn't mean they're.
necessarily equitable..
That's another thing that I've been learning recently is things can be equal, but still.
inequitable..

$^{401}$Let's say we're all looking over this fence and I'm short..
I'm five, seven and a half..
Sam is pretty tall, six, two, six, three..
Eric might be a little shorter than me..
So maybe there's two crates here..
And the equal thing is if we all just had a crate, but I still can't see over the thing..
But it's equal..
We're all equal..
It's like the equitable thing would be maybe Sam doesn't get a crate because he can see.
over the fence..
Eric might get two crates and I get one crate..
So a lot of it is thinking about it in that way too..
Not just, "Oh, let's make it equal for everybody," because it's not necessarily the most helpful.
at times..
Sometimes it's like, "No, this needs to be equitable.".
But it's a lot of really awesome hopeful things that I've seen, like Eric just in the past.
few weeks, that is very different from this situation than situations that we've been.
used to..
It's a lot more people doing something..
I feel like this is a very action-packed time to be a Christian because all eyes are on.
you right now..
What will you do?.
What will you say?.
Because people are going to remember it..
That's beautiful, guys..
I appreciate you sharing your hearts..
I know that's not easy a lot of times..
Everyone can be honest about anything in our lives, but especially about things like this.
that are really part of the core of who you are..
So I just want to thank you guys for that..
I really appreciate you guys..
I'm honored by your friendship..
I appreciate what both of you had to say about relationship..
I know that's been a big part of, for me, understanding more of God's heart and the.
gospel and what it means in relationship to all the peoples of the world..
It's having relationships with different people, and that's been really important..
I'm just always amazed at Jesus' love for us and the fact that, like you were saying,.
it comes back to people..
That's what Jesus cares about..
He cares about hearts and lives, and that's what's important to Him..

$^{441}$He's not really in any way attached to our politics or to our philosophies or our ideas..
He cares about hearts and lives, and that's a part, something we can't and should not.
miss in the church..
I hope you enjoyed the conversation that we had today, and I just want to thank Wade and.
Eric again for coming and having that conversation with us and opening their hearts and being.
so honest and giving us the perspective of their life experience..
That's so good for us, and I hope it was encouraging to you..
I hope it was thought-provoking..
I hope it helped you to see some things that maybe you hadn't seen before..
I really am grateful to both of them for sitting down and doing that with us..
I really appreciate both of them..
As I said at the beginning, I'm honored by their friendship..
They're men who love Jesus and who are giving their lives to serve in Jesus and to serve.
in His church..
They're also great friends of mine..
That was actually the noises that you could hear in the video..
That was the Z's making a lot of noise because Wade and his wife have young children, and.
they're great friends with Jessmyn and I's kids..
Our kids were playing together and having a great time, and so the Z's were making a.
lot of noise..
You could hear that on the video..
But I want to point us to this heart of unity that Jesus has for His church as we close.
it out today..
I want us to consider that..
Unity is hard work in the church..
It's always been hard work, but it's something that Jesus wants for us, and it's something.
that is available for us if we're willing to pursue it..
Pursuing it means sitting down with people who have different experiences than us, who.
have different perspectives than us, and listening, and really listening and hearing, not just.
being in the same room while they're talking, but actually really listening and hearing.
their heart, and practicing some empathy and saying, "I want to try as much as possible.
to think about what life might look like from that perspective, and to understand what my.
brother or my sister has gone through and what that feels like for them.".
To allow that to happen, even if it's difficult, even if it's uncomfortable, even if it reveals.
something about my own heart that I don't really want to see, we need to do that..
We need to do that in order to pursue unity..
We need to allow ourselves to lay down some of our own preferences and some of the things.
that we enjoy and some of the things that we're familiar with in order to understand.
or to accommodate the feelings or the preferences or the familiarity of a brother or sister.
in Christ..

$^{481}$That's one of the things that we have to do to work towards unity in the church, to.
pursue unity..
It's not easy, but God can do that work in us and through us..
We talked last week about how God in these moments wants to do the hard work of heart.
change in our lives, and that's one of the hard things that God wants to do right now..
One of the works of heart change that He wants to do in our lives is helping us to understand.
one another, helping us to hold more loosely to some of our preferences and to some of.
the ways that we're used to doing things and to some of our cultural affinities and be.
willing to take on more of God's kingdom cultural affinities..
That's really important for us to remember..
If we are living according to the values of the kingdom, if we're living according to.
the values of heaven, then all of us are going to be a little bit uncomfortable, no matter.
what cultural background or what our life experience is..
God's going to ask all of us to lay some things down in order to love one another well, and.
that's what we need to be seeking and pursuing in this moment in the church as we continue.
to respond in all the ways we've been talking about over the past few weeks..
We want to pursue unity, and that comes through really caring for one another, really hearing.
one another, really loving each other well, and choosing to give up some of those things.
that maybe are important to us in order to love someone else well..
I'm going to close out just by reading a few verses from John 17..
This is a prayer that Jesus prayed for the church..
I'm going to start in verse 20..
This is John chapter 17, verse 20..
I do not ask for these only, but also for those who will believe in me through their.
word that they may all be one, just as you, Father, are in me and I in you, that they.
also may be in us, so that the world may believe that you have sent me..
The glory that you have given me I have given to them, that they may be one, even as we.
are one, I in them, and you in me, that they may become perfectly one, so that the world.
may know that you sent me and loved them, even as you loved me..
That's a beautiful, powerful passage..
It's a prayer that Jesus was praying for us, for the church..
And in that prayer, Jesus tied the manifestation of God's glory to our unity..
He tied those things together, that God's glory and God's power would be revealed in.
people from all different backgrounds being unified in the church..
And he tied our testimony in the church to our unity, that the world around us, people.
who are outside the church who don't know Jesus, would know the supernatural power of.
the gospel and the reality of Jesus because of the unity in the church..
So it's really important for us to walk in it and to live it out..
It's really important for us to seek it and pursue it in moments like this, even though.
it might be challenging, even though it might feel difficult or frightening to us, or we.

$^{521}$don't understand, or we say, "Hey, what does that mean for me?.
What do I have to lay down?".
We need to pursue and seek after and hunger for unity and for love for one another..
Let's pray..
God, I do pray that your church would be one..
I pray that we would listen to one another, that we would love one another well..
I pray that you would do the hard work of heart change in our hearts and in our lives.
to help us to really care for one another and to help us to be unified in pursuing the.
values of the kingdom, to be unified in pursuing the heart change in our own lives, but then.
the transformation of our families, our neighborhoods, our communities, our cities, and even our.
nations that comes as we respond to the work of Jesus in our heart, as we respond to the.
healing that only He can bring to our hearts..
God, I pray that it would happen..
I pray that we would be healed and that we would love one another and that we would be.
one..
In Jesus' name, amen..
[BLANK AUDIO].
\newpage



\section{}
\label{sec:qRHTClh4NmQ}
\textbf{12 十一少被賣記大團圓－上(一)}
\newline
\newline
連結: \href{https://youtube.com/watch?v=qRHTClh4NmQ}{\texttt{https://youtube.com/watch?v=qRHTClh4NmQ}} ~~~~ 語音日期: 2020-07-05
\newline
\newline
\hyperref[sec:dNRqxYjeS0o]{\small{< < < PREV SERMON < < <}}
~
\hyperref[sec:index_chronic]{\small{[返順時目]}}
~
\hyperref[sec:index_scriptual]{\small{[返順卷目]}}
~
\hyperref[sec:hwLDlGYTpto]{\small{> > > NEXT SERMON > > >}}
\newline
\newline
$^{1}$《十一少》被賣寄資因禍得福.
故事的發展可以說是峰迴路轉.
在上一集提到主角約瑟.
就在七年的豐收期.
將收到來的糧食儲在倉庫裡.
等待七年的寒年一到.
埃及全國不僅沒有因為大饑荒而缺糧.
埃及還反而成為各地前來買糧的地方.
昨年約瑟的親父雅各一聽聞消息後.
就立即吩咐十個兒子到埃及來買糧.
看看42章,一到第五節.
經文沒有明確指出為何雅各十個兒子沒有主動.
來向父親提出.
反而是父親主動吩咐十個兒子到埃及來買糧.
不過,雅各十個兒子沒有主動.
去埃及來買糧.
不過有一點清楚的來指出.
雅各當時只是吩咐他十個兒子去.
卻留下一個兒子在自己身邊.
自從約瑟出事後.
雅各最愛的兒子就是約瑟的親弟別雅文.
別雅文是約瑟同父同母的親弟弟.
也是雅各失去約瑟後.
唯一最通愛的兒子.
雅各當日之所以會失去約瑟.
雅各自己也要負責任.
因為是雅各吩咐約瑟一個人去找十個兄長.
我相信在過去的二十多年以來.
雅各對當日的決定一定很後悔.
雅各為免再一次失去自己所親愛的兒子.
雅各今次就只有吩咐十個兒子去埃及買糧.
從雅各今次的反應和行動.
可以表明在二十多年以來.
雅各依然對他的兒子非常厚此薄彼.
雅各家庭的根要問題似乎沒有改善過.
雅各由始至終關心的只是別雅文一個人的安危.
事實上從雅各的反應我們能夠體會到.
時間未必一定能使人走出傷痛.
在故事中雅各一家還是活在昔日的傷痛中.
弟兄姊妹我們現在有沒有活在昔日的傷痛中呢?.

$^{41}$我們每一個人未必是為人父母.
但我們每一個都是為人子女.
我們有沒有好像雅各十二個兒子一樣.
依然活在原生家庭的陰影之下呢?.
雅各吩咐十個兒子到埃及.
十兄弟只有聽父親的話.
即使當時的十兄弟已經成家立室.
也是為人父母的中年人.
可見當時雅各的家.
一家之主依然是父親雅各.
再看6-17節.
當十兄弟來到埃及朝見埃及首相.
當然要行埃及跪拜之禮.
這個跪拜之禮的舉動.
正正應驗了當年約瑟在年輕時所做的夢.
記得當年約瑟曾夢見自己的十兄弟.
向他來到跪拜.
在十兄弟面前的埃及首相約瑟.
不單是埃及一人之下萬人之上.
同時也是主持買賣糧的最高領導人.
當時十兄弟跪拜埃及首相.
根本他們無法了認.
當時的埃及首相就是他們的弟弟約瑟.
畢竟他們兄弟已分隔二十多年.
加上當時的約瑟已經不再是一個17歲的年輕人.
約瑟在結婚時自己十個兄長時.
他已經是一個中年人.
還有就是當時的約瑟已經被埃及文化同化了.
約瑟當時說的是埃及話.
穿的是埃及衣服.
吃的都是埃及食物.
約瑟一切都是按照著埃及的生活和文化.
當然約瑟的十個兄長根本無法認出.
埃及的首相就是分隔了二十多年的十一弟約瑟.
不過約瑟卻一眼就能認出十個兄長.
如果想知道約瑟也將十個兄長之後.
將會有甚麼事情發生.
請大家留意下一次我和大家的講解.
\newpage



\section{}
\label{sec:hwLDlGYTpto}
\textbf{13 十一少被賣記大團圓－上(二)}
\newline
\newline
連結: \href{https://youtube.com/watch?v=hwLDlGYTpto}{\texttt{https://youtube.com/watch?v=hwLDlGYTpto}} ~~~~ 語音日期: 2020-07-06
\newline
\newline
\hyperref[sec:qRHTClh4NmQ]{\small{< < < PREV SERMON < < <}}
~
\hyperref[sec:index_chronic]{\small{[返順時目]}}
~
\hyperref[sec:index_scriptual]{\small{[返順卷目]}}
~
\hyperref[sec:BmdnpSdGr5s]{\small{> > > NEXT SERMON > > >}}
\newline
\newline
$^{1}$上一節提到,因為大饑荒的緣故,雅各就馬上吩咐十個兒子到埃及來買糧.
當這十兄弟見到埃及的首相的時候,由於這十兄弟未能當時認出埃及的首相就是約瑟的緣故.
於是他們十兄弟都向當時埃及的首相行埃及要行的跪拜之禮.
看看第九節,雖然當時十兄弟未能認出約瑟,但約瑟卻是一眼就能認出他們就是自己的十個兄長.
不單止如此,約瑟還想起當年年輕時候所發的夢.
於是約瑟一口指定他們就是奸細.
似乎約瑟對當年所發的夢,約瑟不單止沒有忘記,反而一直都放在心裡.
當年的約瑟所夢見的是十兄長和父母親都要跪拜他.
同時亦因為約瑟所發的夢,而加深了當年十兄長對約瑟的疾度.
最後還引致約瑟被賣,成為埃及的奴隸.
事隔了那麼多年,十兄弟就因為父親雅各的吩咐而來到埃及買糧.
本應動機是很明確的,不過約瑟偏偏就要為難他們,說他們是奸細.
約瑟不單止不接受當時十兄弟對他的解釋.
還要求他們所有人都要留下,只准其中一個人回到迦南.
回到迦南後,還要把最小的弟弟帶來埃及,證明他們不是奸細.
否則他們所有人都不可以離開埃及.
於是約瑟就將他們十兄弟囚在監獄三天.
故事沒有交代過十兄弟在監獄三天內發生什麼事情.
亦沒有交代約瑟在三天內究竟在想什麼.
去到十八到二十六節,說到三天後.
約瑟向十兄弟為了表明自己是一個敬畏上帝的人.
於是就准他們帶著糧食回迦南.
讓他們可以夠活在迦南的家人.
不過有附帶條件.
條件是要他們十兄弟當中.
其中一個人留下並囚在監獄裡.
直至等到他們把最小的弟弟從迦南帶來埃及.
三天前後,約瑟所定的條件先後有所改變.
在三天前,約瑟的條件只可以其中一人回迦南.
還帶最小的弟弟從迦南來到埃及.
三天後的條件呢?.
約瑟只要他們其中一個人留下.
其餘所有人都可以帶著糧食回迦南.
不過依然要將最小的弟弟從迦南帶來埃及.
大家有沒有想過.
事隔三天,約瑟為何要開出這樣的條件呢?.
試想想,如果約瑟只是放一個兄弟回家.
試問一個人可以帶足夠的糧食回迦南嗎?.
另外,約瑟當時顧念的是自己的父親雅各.
一下子要父親失去九個兒子.

$^{41}$父親能否承受他的傷痛呢?.
還有,約瑟如果要試十兄弟,有沒有兄弟情.
留一個在這裡和留九個在這裡.
基本上不是太大的分別.
十兄弟當中,留下其中一個人囚在監獄裡.
其餘人回去後,就要將最小的弟弟從迦南帶來埃及.
換算,你就是那十兄弟,你又會怎樣來這裡做呢?.
下一次我會帶大家來看看.
故事裡的這十兄弟對約瑟所開出的條件.
究竟有什麼反應行動.
請不吝點贊訂閱轉發打賞支持明鏡與點點欄目.
\newpage



\section{}
\label{sec:BmdnpSdGr5s}
\textbf{15十一少被賣記大團圓－上(四)}
\newline
\newline
連結: \href{https://youtube.com/watch?v=BmdnpSdGr5s}{\texttt{https://youtube.com/watch?v=BmdnpSdGr5s}} ~~~~ 語音日期: 2020-07-06
\newline
\newline
\hyperref[sec:hwLDlGYTpto]{\small{< < < PREV SERMON < < <}}
~
\hyperref[sec:index_chronic]{\small{[返順時目]}}
~
\hyperref[sec:index_scriptual]{\small{[返順卷目]}}
~
\hyperref[sec:F_RgXPWQknI]{\small{> > > NEXT SERMON > > >}}
\newline
\newline
$^{1}$上一節提到,弱勢得到神秘的機會.
能夠面對昔日,傷害過他的十兄長.
今次就帶大家去了解一下.
弱勢的心靈又如何得到意志.
在一般情況下,人被困在傷痛的情緒當中.
人很難冷靜,客觀分析.
客觀分析人和事.
故事沒有交代過.
弱勢當年要流落異鄉.
做埃及人奴隸的時候.
有沒有受過埃及人的虐待.
不過,在埃及的日子裡.
弱勢的確無親無故.
試想想當年,弱勢只是一個十七歲的年輕人.
幾乎還要死在親兄長的手下.
在他身上所發生的事.
對一個十七歲的年輕人來說.
的確造成很大的心靈上的陰影.
當弱勢聽到大哥留便重提當年的事.
有殺自己的念頭,有埋自己的行動.
弱勢的痛苦經歷就像電影一樣.
一幕一幕在弱勢的腦海中浮現.
弱勢的傷痛記憶觸動到他的情緒而大哭.
是極之正常的事.
我估計當時弱勢會在想.
為何十兄長會不顧親情.
要這樣對待自己.
究竟自己做錯了甚麼事呢?.
弱勢退去大哭的反應行動.
估計當時的弱勢並沒有心理上的準備.
說要和十兄長相認.
於是弱勢便轉身退去大哭一場.
以免被他十個兄長發現弱勢的身份.
弱勢的反應行動一方面.
除了證明當時弱勢未想和十兄長相認.
可能是對十兄長未完全的信任.
畢竟十兄長當年對他曾經做過很大的傷害.
雖然事隔了二十多年.
而弱勢仍然未能釋懷.
於是他便計劃試一試十兄長的兄弟情.

$^{41}$另一方面,弱勢是一個非常重情重義的人.
他對原生家庭並沒有真正忘記.
畢竟十兄長都是自己同父異母的兄長.
加上親父母和親弟弟現在遇到饑荒.
弱勢豈可見自己的親人死而不救呢?.
可想而知,當時的弱勢和十兄長會面時的心情非常複雜.
雖然當時的弱勢在原生家庭生活後十七年的時間.
而在埃及,定居生活已經超過二十年.
而且當時的弱勢已經有自己的家庭.
不過弱勢沒有因為這個緣故而忘記了親情.
弟兄姊妹,當我們面對著自己的傷痛時.
我們會選擇逃避.
還是像故事中的弱勢一樣勇敢積極面對呢?.
弱勢大好一場後,他再回來.
按照先前的計劃,他選擇了西面.
要西面,一個人留在埃及.
雖然弱勢吩咐將糧食裝滿在他們的器皿裡.
又將各人的銀紙歸在各人的口袋.
還給他們一路上的食物.
於是他們就可以離開埃及,回到迦南.
帶大家猜想一下.
當時九兄弟帶著什麼心情離開埃及,回到迦南呢?.
而在迦南的途中,又有什麼事情發生呢?.
下一次,我會再帶大家一起繼續探討下去.
\newpage



\section{}
\label{sec:F_RgXPWQknI}
\textbf{16十一少被賣記大團圓－上(五)}
\newline
\newline
連結: \href{https://youtube.com/watch?v=F_RgXPWQknI}{\texttt{https://youtube.com/watch?v=F\_RgXPWQknI}} ~~~~ 語音日期: 2020-07-06
\newline
\newline
\hyperref[sec:BmdnpSdGr5s]{\small{< < < PREV SERMON < < <}}
~
\hyperref[sec:index_chronic]{\small{[返順時目]}}
~
\hyperref[sec:index_scriptual]{\small{[返順卷目]}}
~
\hyperref[sec:WV2ED6LTYMM]{\small{> > > NEXT SERMON > > >}}
\newline
\newline
$^{1}$上次說到日蝨留下西蓮一個.
要放其餘九個回到迦南.
今次就帶大家來看看.
九兄弟回迦南的途中究竟有什麼事情發生.
而回到迦南之後.
當九兄弟把埃及買糧的事.
和手上的要求都告訴父親的時候.
父親雅各有什麼反應呢?.
看看以上的資料.
看看27節.
在回程的途中有人打開口袋.
發現口袋裡依然有自己的銀紙.
就告訴其餘的兄弟.
他們膽戰心驚.
彼此說究竟上帝對我們做了什麼.
可想而知當時十兄弟非常害怕.
他們很害怕有一天.
神會為當日他們所犯的罪.
而向他們來到追討.
而且他們很認為.
很有可能就是這次到埃及買糧這件事.
所遭遇的一切.
畢竟他們離開迦南去到埃及買糧的時候.
是十兄弟完整一齊的去的.
但結果.
在埃及買糧之後回來.
就只有九兄弟.
加上埃及的首相.
並沒有收取到他們任何買糧的錢.
唯獨是要求他們留下西面一個兄弟.
還要求他們九兄弟回到迦南之後.
要將父親最愛的兒子別雅文離開父親.
將他們帶到埃及.
看29節.
當他們回到迦南.
就將埃及所發生的事.
和埃及首相的要求都告訴父親.
當其時父親聽完之後.
他有什麼反應?.
換算你就是父親雅各.

$^{41}$你又會有什麼反應?.
第36節.
當其時雅各對九個兒子說.
「你們使我喪失了我的兒子」.
似乎雅各一直對當年.
失去愛子約瑟的事仍然耿耿於懷.
認為十個兒子當年.
使他失去了最愛的兒子約瑟.
如今西面又被扣留.
還說要將他身邊唯一的愛子別雅文帶走.
雅各不得不懷疑.
十兄弟的兄弟情.
大兒子為了讓父親來到相信.
於是大兒子劉便就向父親來承諾.
如果不能將十二弟別雅文.
平安從埃及帶回來的話.
劉便就願意賠上他自己兩個親兒子的性命.
大家有沒有留意?.
雅各的大兒子劉便.
和他父親雅各的反應行動.
對比之下.
做父親的雅各的確只是偏愛著一個兒子.
相反大兒子劉便.
他是以自己兩個親兒子的性命作為擔保.
為的是要讓親弟西面.
可以離開埃及平安回到迦南.
大兒子劉便為了父親的信任.
可以放心讓別雅文去埃及.
向埃及的首相證明他們十兄弟不是奸細.
就用自己兩個兒子的性命作為擔保.
雖然大兒子劉便有這樣的承諾.
但父親雅各依然不肯.
也讓別雅文離開自己.
可想而知.
雅各當時是多偏愛著別雅文一個兒子.
他為了保護這個心愛的兒子.
他不惜也要放棄和其餘十個兒子的親情.
甚至犧牲他們的生命.
弟兄姊妹.
雅各的反應提醒了我們.

$^{81}$一家人可以團圓和一家人可以復和.
不是單方面的.
一家人的團圓和一家人的復和是雙方面的.
即使一方願意.
另一方也要同意.
當時大兒子劉便願意.
但父親雅各卻不同意.
故事發展到這裡.
對大家有甚麼提醒?.
有神許可之下.
我們需要勇敢面對昔日的傷痛.
這樣才可以和神和人和自己復和關係.
如果神現在向你發出挑戰的話.
你又願不願意回應神向你的挑戰呢?.
下一次我會和大家繼續來看看.
雅各家是如何可以大團圓結局.
多謝您收睇時局新聞,再會!.
\newpage



\section{}
\label{sec:WV2ED6LTYMM}
\textbf{17 十一少被賣記大團圓－中(一)}
\newline
\newline
連結: \href{https://youtube.com/watch?v=WV2ED6LTYMM}{\texttt{https://youtube.com/watch?v=WV2ED6LTYMM}} ~~~~ 語音日期: 2020-07-06
\newline
\newline
\hyperref[sec:F_RgXPWQknI]{\small{< < < PREV SERMON < < <}}
~
\hyperref[sec:index_chronic]{\small{[返順時目]}}
~
\hyperref[sec:index_scriptual]{\small{[返順卷目]}}
~
\hyperref[sec:egka2gEIyik]{\small{> > > NEXT SERMON > > >}}
\newline
\newline
$^{1}$上一節提到父親雅各不肯讓最愛的兒子別雅文離開自己.
即使世面被扣留在埃及.
但兒子留便又用自己兩個兒子的性命作為擔保也好.
父親還是不能改變自己的心意.
可想而知雅各當時是偏愛著別雅文.
不惜放棄和犧牲對其他十個兒子的親情甚至性命.
今次說到日蝨和完新家庭的團聚.
我們看看43章,一到第五節吧.
迦南饑荒的情況愈來愈嚴重.
即使之前在埃及所帶回來的糧食都不夠.
於是父親雅各見到情況如此緊迫.
就再一次吩咐九個兒子來到埃及買糧.
按照故事前後發展推算下.
上一次和今次的買糧距離時間不會多過兩年.
因為按前文,可年有七年長的時間.
而按後理,饑荒還有持續五年的時間.
由上次由迦南到埃及買糧.
到現在再打算第二次到埃及.
前後不會多過兩年的時間.
由於上一次帶回來的糧食.
因為路途遙遠,人手又不足.
可以帶回來的糧食是很有限.
估計當時帶回來的糧食只能供應一段短的日子.
當時雅各也真的沒想過.
饑荒的情況會持續不斷,而且愈來愈嚴重.
當雅各家再一次缺糧饑荒.
唯一的解決方法就只好再一次到埃及來買糧.
不過不知道雅各是忘記了.
還是一直沒有放在心裡.
在上一次他吩咐十個兒子.
到埃及買糧後,結果只有九個兒子平安回來.
西勉自從離開迦南,一直被扣留在埃及.
當時埃及首相只容許九兄弟離開埃及回迦南.
左相還開出條件,要帶別雅文到埃及.
當時父親雅各拒絕.
將自己唯一的愛子別雅文帶到埃及.
即使當時大兒子流變.
自願用自己兩個兒子的性命作為擔保.
父親也沒有改變心意,於是就不了了之了.
直至饑荒的情況愈來愈嚴重.

$^{41}$父親雅各不能不吩咐九個兒子再次到埃及買糧.
雅各其中有一個兒子叫尤大.
尤大也是同時當年提出要殺約瑟的那個.
尤大向父親重提埃及首相當年的要求.
若父親不用讓別雅文再到埃及買糧.
埃及首相不會和他們會面,更不會答應買糧.
在經文的八到第九節.
尤大強調除非父親肯讓別雅文同去.
否則不能夠救活一家.
尤大還將自己的性命作為擔保.
他承諾若不能平安.
將十二弟別雅文平安帶回家.
尤大願意將自己的性命作為補償.
父親雅各終於批准.
將最好的乳香,蜂蜜,香料,活藥當作禮物.
又要他們加倍的大銀子同去.
為何要加倍呢?.
難道雅各擔心埃及首相不賣糧?.
還是雅各擔心別雅文的安危呢?.
雅各要加倍的原因.
一方面是用來買糧用.
但另一方面是想歸還上次買糧的錢給埃及首相.
因為上次埃及首相並沒有收取他們買糧的錢.
而上次和今次最大的分別.
是雅各終於肯別雅文同去.
大家有否留意.
在上次,大仔劉便同樣向父親請求.
甚至劉便以自己兩個親生兒子的性命作為擔保.
但結果依然不獲父親的批准.
但今次呢?.
尤大向父親同樣有請求.
希望父親准別雅文同去.
結果竟然沒有到父親的批准.
大家認為,父親雅各為何肯改變心意呢?.
當時饑荒的情況的確很嚴重.
再沒有糧食供應的話.
雅各的家,從老至少都會盛名不保.
到時雅各同樣會失去他的愛子別雅文.
而且不單止失去一個人的生命.
而是整個家族的存亡.

$^{81}$雅各作為家族之首,他是有責任的.
來保護家族裡每一個大大小小的平安.
究竟今次雅各十個兒子再一次到埃及買糧.
又會有什麼事情發生呢?.
下一次我會繼續和大家繼續一來探討下去.
\newpage



\section{}
\label{sec:egka2gEIyik}
\textbf{20 十一少被賣記大團圓－中(四)}
\newline
\newline
連結: \href{https://youtube.com/watch?v=egka2gEIyik}{\texttt{https://youtube.com/watch?v=egka2gEIyik}} ~~~~ 語音日期: 2020-07-06
\newline
\newline
\hyperref[sec:WV2ED6LTYMM]{\small{< < < PREV SERMON < < <}}
~
\hyperref[sec:index_chronic]{\small{[返順時目]}}
~
\hyperref[sec:index_scriptual]{\small{[返順卷目]}}
~
\hyperref[sec:fHYowIViTao]{\small{> > > NEXT SERMON > > >}}
\newline
\newline
$^{1}$上一節提到.
約瑟親耳聽到當年想殺他.
又最沒有兄弟情的二哥猶大.
竟以自己的性命來擔保.
為的是救同父異母的弟弟.
別牙文.
案情的案例.
約瑟都是被害者.
約瑟是最有資格.
為當年他的遭遇來犯案.
約瑟也是最有資格.
為當年他遭害來取回公道的人.
不過約瑟不單沒有這樣做.
反而他還主動向傷害他的人提出.
不要再為當年的事憂心.
約瑟之所以有這樣的反應行動.
除了是受猶大感動之外.
還有就是他認為有更重要的事要處理.
究竟約瑟還有什麼更重要的事情要處理呢?.
因為當時迦南遇到饑荒已經有兩年的時間.
而饑荒會持續多五年.
約瑟很擔心.
擔心是迦南的親生父母.
於是約瑟就將神的旨意.
向兄弟來堂開.
說明自己是神預祖的計劃.
安排他先來到埃及.
為的是救活迦南地的一家.
約瑟對他過去的遭遇.
從他的反應我們可以看出.
約瑟心裡一直都相信神是信實的.
無論他過去所遭遇的一切.
他都相信神與他同在.
而且神一直保守他.
也帶領著他.
約瑟也要他的兄弟趕緊回老父那裡回報.
饑荒的問題是會持續多五年.
勸老父要帶同所有人和生畜.
離開迦南來到埃及定居.
還要老父知道的是.

$^{41}$約瑟不但沒有死.
反而還成為了埃及的首相.
不論在地位,勢力方面都足以保證.
整個的雅各家族可以安居於埃及.
約瑟和親兄弟比雅敏.
互相在對方的頸上大哭.
約瑟亦向其他兄長親嘴並相抱大哭.
訂姐妹.
有沒有留意到.
如果要一個已經分裂的大家庭.
可以有大團圓的結局.
是不可以只要求一方面願意.
大團圓的結局是需要雙方都願意.
承認,接受.
每一個人包括自己都不是完美.
亦要為自己的過犯而應罪和悔改.
更要為過去無論是自己或別人的過犯而釋懷.
訂姐妹.
聖經教導我們要彼此饒恕.
彼此不是單數.
彼此是雙數.
彼此饒恕不單只是饒恕別人.
重要的是要饒恕自己.
在故事中主角約瑟.
他確實做得到.
既饒恕別人亦饒恕自己.
所以約瑟一點也沒有苦讀.
約瑟全心要救活雅各家族.
為的就是要馴服神的心意.
可見約瑟既愛神,愛自己亦愛別人.
當十兄弟帶十二弟別雅文回到迦南後.
就把埃及的事全部報告給老父雅各知道.
於是雅各就帶同子女,孫和所有後裔.
包括生畜和財物.
去到別士巴這個地方.
雅各在別士巴這個地方.
是先向耶和華獻祭.
大家留意一下.
雅各向耶和華獻祭這個舉動.
其實是帶有相當高的屬靈意義.

$^{81}$雅各帶同子女,孫和所有後裔.
包括生畜和財物.
要離開迦南去到埃及這個決定.
並不是一件簡單的事.
雅各是先來到寇高耶和華神.
確認了得到神的准許,神的賜福和保守.
雅各才採取正式行動.
舉家的離開迦南去到埃及.
大家若想知道.
約瑟和父親雅各和一家團聚後.
會有什麼事情發生.
請不要錯過下一次和大家的講解.
\newpage



\section{}
\label{sec:fHYowIViTao}
\textbf{22 十一少被賣記大團圓－下(二)}
\newline
\newline
連結: \href{https://youtube.com/watch?v=fHYowIViTao}{\texttt{https://youtube.com/watch?v=fHYowIViTao}} ~~~~ 語音日期: 2020-07-06
\newline
\newline
\hyperref[sec:egka2gEIyik]{\small{< < < PREV SERMON < < <}}
~
\hyperref[sec:index_chronic]{\small{[返順時目]}}
~
\hyperref[sec:index_scriptual]{\small{[返順卷目]}}
~
\hyperref[sec:6gEpaWbkTWE]{\small{> > > NEXT SERMON > > >}}
\newline
\newline
$^{1}$(《約瑟》).
上一節提到當約瑟知道父親.
已經去到歌山這個地方.
約瑟就親自去到歌山迎接父親.
當約瑟和新家庭團聚之後.
接下來他就要安頓雅各家族.
當時雅各家一共有70人.
在經文47章11-12節說.
約瑟因為得到法老同意之下.
約瑟就安排雅各家留在歌山定居.
因為歌山的地形適合蓄木.
雅各和約瑟團聚之後.
一直在歌山過著蓄木生活有17年.
當雅各知道他自己的死期將到.
雅各就請求約瑟為他做一件事.
究竟是什麼事呢?.
原來雅各有一個心願.
就是死後的屍體.
能夠在家欄埋葬.
好讓他和祖父和父親能夠同葬.
大家有沒有留意.
猶太人價值觀念和中國人傳統文化有些相似.
告老還鄉.
死後葬在自己親人旁邊.
故事中的雅各希望將來可以死後.
能夠和祖先同葬.
於是他就留下自己的遺願.
給自己最愛又最信任的兒子約瑟.
在經文的48章說.
雅各臨死前.
他就祝福約瑟兩個兒子.
約瑟的大兒子叫嗎哪西.
小兒子叫爾法列.
雅各要將兩個孫歸入自己的名下.
究竟是什麼意思呢?.
雅各將約瑟兩個兒子歸入自己的名下.
就是要約瑟將兩個兒子過繼給他.
雅各這樣做是想約瑟可以分得更多的產業.
可見約瑟由始至終都是雅各最愛的兒子.
於是雅各在臨死前.

$^{41}$就為嗎哪西和爾法列祝福.
按照傳統.
嗎哪西是大兒子.
理應祝福的右手放在他的頭上.
而左手放在小兒子的頭上.
但雅各竟然用他的右手.
放在小兒子爾法列的頭上.
而左手就放在大兒子嗎哪西的頭上.
雅各這樣做不是因為他腦眼分花.
而是因為雅各知道.
將來小兒子爾法列的支派.
會比大兒子嗎哪西的支派強大.
看看49章.
那裡說到.
當雅各賜福給爾法列和嗎哪西後.
雅各就召集其餘所有子孫聽他的遺言.
雅各對他所有子孫宣佈.
長子劉弁因為之前曾放縱情慾.
所以從此失去了長子的名分.
而次子西冕和三子利美.
因為他們曾用刀劍秘密殺害人命.
結果要分散在各支派當中.
同時失去了長子的名分.
故此長子的名分就落入四子猶大的身上.
大家有沒有發覺.
雅各所立的遺囑和中國傳統文化很相似.
就是很看重長子名分.
因為父親的產業.
一般由長子一個人繼承.
又或者當分家產的時候.
長子是可以得到雙分的.
記得當年雅各年輕的時候.
曾經用繼騙長子的名分.
獲得父親以殺的祝福.
如今雅各年老.
在離世之前.
他是親手祝福猶大.
同時亦親手將大的福氣.
賜給弱勢的小兒子爾法列.
雅各的一生.

$^{81}$因為曾經經歷過.
爭長子的名分.
而使兄弟情破裂.
還要離開原生家庭流落到異鄉.
雅各傷痛的經歷.
由一代傳到一代.
雅各因他的偏愛弱勢緣故.
而帶來給他們兄弟情破裂的傷痛.
所以雅各臨死之前.
雅各信著上帝的心意和啟示.
不再用自己的心意來祝福他的子孫.
雅各死前向所有子孫交代.
免得他們和下一代再一次要經歷到.
因為名利和權.
而導致兄弟不和所帶來的傷痛.
頂子輩.
福和和大團圓.
其中要達至的目標.
就是要傷痛不代存.
或許我們可能在人際關係方面.
家庭,工作甚至教會方面.
受過傷害.
如果我們願意經歷福和和大團圓.
我們就要留心.
在生命上的見證.
有沒有將傷痛帶到自己的家庭.
工作甚至教會裡.
難道我們真的想將傷痛的經歷.
一代傳到一代?.
如果想知道故事發展.
請留意下一次和大家的講解.
\newpage



\section{}
\label{sec:6gEpaWbkTWE}
\textbf{23 十一少被賣記大團圓－下(三)}
\newline
\newline
連結: \href{https://youtube.com/watch?v=6gEpaWbkTWE}{\texttt{https://youtube.com/watch?v=6gEpaWbkTWE}} ~~~~ 語音日期: 2020-07-06
\newline
\newline
\hyperref[sec:fHYowIViTao]{\small{< < < PREV SERMON < < <}}
~
\hyperref[sec:index_chronic]{\small{[返順時目]}}
~
\hyperref[sec:index_scriptual]{\small{[返順卷目]}}
~
\hyperref[sec:oFXXJCLRQwo]{\small{> > > NEXT SERMON > > >}}
\newline
\newline
$^{1}$《亞國》.
上一節提到亞國死前.
向所有子孫交代清楚.
免他們和下一代再次要經歷到.
因為名利權而導致兄弟不和所帶來的傷害.
在經文的50章1到2節說.
亞國祝福和留下遺言之後.
亞國就安然離世.
日蝨為父親的離世而大哭.
今次日蝨的大哭是生離死別的哭.
日蝨為父親辦身後事.
就吩咐醫生用香料分割父親的屍體.
同時日蝨也遵行父親的遺願.
將父親的屍首送回迦南安葬.
當日蝨處理父親亞國的身後事後.
有一件事發生.
由於當時十兄弟良心不安.
怕父親去世後.
日蝨會因為他們以前的所作所為.
而向他們報復.
於是十兄弟就打發人.
來向日蝨求情.
從十兄長的行動可以知道.
原來過去十七年的所謂團聚和復和.
並不是真實和完全.
至少十兄長以為日蝨在十七年前肯原諒他們.
是全因為老父的緣故.
如今老父已逝世.
但認為日蝨再沒有任何復仇的故記.
而事實上日蝨在十七年前.
和原生家庭團聚時.
就已經原諒十兄弟.
日蝨之所以能夠自發,耐心,真誠原諒十兄弟.
主要的動力不是基於和父親的關係.
主要來源的動力是來自同實的關係.
日蝨向十兄長道說清楚.
他對十兄長說.
「不要怕,我豈能代替神呢?.
從前你們的意念是要害我.
但神的意念亂是好的.

$^{41}$是要傳許多人的成名.
成就今日的光景」.
日蝨向著他們的承諾.
必定養活他們和他們的孩子.
日蝨回望當年被兄長埋的經歷.
日蝨不是不知道當日.
十個兄弟對他不懷好意.
但日蝨更加知道的是.
神容許發生.
神讓日蝨出生在一個複雜環境的家庭.
日蝨知道.
也信服神更高的意念.
日蝨清楚知道.
神是想借用這個機會.
使雅各整個家族因為日蝨.
成為埃及的祖上.
而使全個家族得救.
神祂祈告一象.
而日蝨也信服神.
甘願成為神的棋子.
被神所使用.
日蝨能夠釋懷.
還真心原諒得罪和傷害過他的人.
今日可能我們都會因為昔日的事.
而不能釋懷.
似乎我們還在怨天尤人.
不僅覺得人欠了我們.
有時候甚至會覺得神欠了我們.
如果人有自我反省的能力.
重新可以回顧自己的人生.
起與跌.
或許大家就會發現.
在我們沒有留意的時候.
原來神一早就已經動弓.
看不見並不是神不知不理不回應.
只不過神並不是安然.
按照著你的心意.
你的方法.
你的時間.
所以你看不見.

$^{81}$又或者是你自己拒絕了神的介入.
和神處理的方法.
能夠將一段已經破壞的關係.
可以復和.
還可以有大團圓結局.
是必須要先來明白神的心意.
更重要的是信服神的心意.
日蝨死之前.
他依然向他的兄弟.
遵守當年的承諾.
日蝨死的時候是110歲.
而且還是五代同堂.
但是長子爾法年的後裔.
只是壯丁有四萬零五百名.
亦都應驗了當年.
雅各向爾法年的祝福.
雅各家自從離開迦南到埃及.
一共在埃及生活了四百年的時間.
四百年裡的日子.
並沒有被埃及的文化所同化.
反而雅各家族越來越強大.
由起初雅各家七十人到埃及定居.
到他們出埃及之前.
已經有數百萬人.
在《民數記二章》三十二節記載.
他們出埃及的時候.
只是南丁已經有六十萬人.
十一兆被賣寄的故事.
已經來到尾聲.
從約瑟原生家庭的問題.
到他經歷到神同在的過程.
去到最後.
約瑟能夠和父親和兄弟大團圓.
主角約瑟.
這位十一兆的日生.
都是由耶和神一直引領著他.
在神引領下.
讓他出生在一個關係複雜的大家庭.
在神的引領下.
而令他被賣到埃及.

$^{121}$遭奴隸 遭囚犯.
在神的引領下.
從埃及奴隸 囚犯的身份.
逆轉成為埃及的宰相.
在神的引領下.
可以和兄弟復和 一家團聚.
在神的引領下.
能夠幫助雅各家族成為最大的民族.
今天 我們有沒有容讓到神.
引領著我們的一生.
經歷神在我們生命當中的引領.
同時 有體會在關係當中.
復和 團圓的美好呢?.
請不吝點贊 訂閱 轉發 打賞支持明鏡與點點欄目.
\newpage



\section{}
\label{sec:oFXXJCLRQwo}
\textbf{14十一少被賣記大團圓－上(三)}
\newline
\newline
連結: \href{https://youtube.com/watch?v=oFXXJCLRQwo}{\texttt{https://youtube.com/watch?v=oFXXJCLRQwo}} ~~~~ 語音日期: 2020-07-07
\newline
\newline
\hyperref[sec:6gEpaWbkTWE]{\small{< < < PREV SERMON < < <}}
~
\hyperref[sec:index_chronic]{\small{[返順時目]}}
~
\hyperref[sec:index_scriptual]{\small{[返順卷目]}}
~
\hyperref[sec:2kNd6PvPzMU]{\small{> > > NEXT SERMON > > >}}
\newline
\newline
$^{1}$上一節提到,約瑟向十兄長開出的條件.
要他們兄弟當中留下一個.
其餘九個回到迦南.
將迦南最小的一個弟弟從迦南帶到埃及.
看21-22節.
那裡提到十兄弟的對話.
在對話當中,大哥劉變重提當年所發生的事.
當年十兄弟有殺約瑟的動機.
當年十兄弟有賣約瑟的行動.
重對話當中可以知道.
十兄弟對當年的所作所為.
在過去二十多年以來.
不單一直沒有忘記過.
反而還記在心裡.
而且還害怕有一天.
罪會向他們來到追討.
弟兄姊妹,你們猜想一下.
十兄弟在過去二十多年以來.
除了害怕有一天,罪會向他們追討之外.
你們會認為他們有沒有對當年的所作所為有悔意呢?.
或許他們對四日的所作所為真的有悔意.
又或者他們希望可以有大團圓的結局.
可惜的是,在過去二十多年以來.
他們似乎沒有實際行動過.
不過,真正能寫到事情.
能夠大團圓結局的人.
似乎只有故事裡的主角,約瑟.
看看二十三到二十六節.
當約瑟聽完十兄弟的對話後.
他就轉身退去大哭一場.
大家有沒有想過.
約瑟聽完之後,為什麼約瑟需要退去呢?.
約瑟為什麼要大哭一場呢?.
約瑟當時已經是一個中年男人.
也是埃及一國的首相.
一人之下,萬人之上.
還有他可以負責買賣糧的最高領導人.
李英,以約瑟當時的身份.
他既有權有勢,又有一個幸福的家庭.
根本是沒有需要為了二十多年前所發生的事.

$^{41}$而退去大哭一場.
除非約瑟當年遭遇的禍患.
所帶來給他的心靈傷痛.
一直都沒有機會正視和處理.
即使約瑟現在已經有身份,有權,有勢.
也有一個幸福的家庭.
不過約瑟所擁有的一切.
都未能補償他當日的心靈傷痛.
鄭英姐妹,有些時候上帝會容許某些事情的發生.
推駛我們必須要面對昔日的傷痛.
就好像故事裡的約瑟.
和十兄弟舉起的重聚.
是神給約瑟一個機會.
來面對和處理昔日的傷痛.
鄭英姐妹,如果神讓我們有機會.
要和曾經傷害過我們.
又或者得罪了我們的人復和的時候.
你又會如何選擇呢?.
逃避?還是勇敢而正面地面對呢?.
如果要復和的話.
就必先要面對和正視昔日自己所受的傷痛.
處理心靈傷痛的過程.
就好像處理傷口一樣.
在過程中因為觸動到感覺神經.
所以會有痛的感覺.
不過,痛是人的正常反應.
當有適當處理,傷口就會很快地痊癒.
同樣,處理心靈上的傷口也是一樣.
在過程中,要重提昔日的痛苦經歷.
始終也會觸動到人的情緒反應.
如果要有效地處理和復和人的關係.
痛的過程是必須的.
不過,處理的時候也必須要適當.
否則會帶來更大的反效果.
下一次我會帶大家一起探討一下.
謝謝收看!.
\newpage



\section{}
\label{sec:2kNd6PvPzMU}
\textbf{18 十一少被賣記大團圓－中(二)}
\newline
\newline
連結: \href{https://youtube.com/watch?v=2kNd6PvPzMU}{\texttt{https://youtube.com/watch?v=2kNd6PvPzMU}} ~~~~ 語音日期: 2020-07-07
\newline
\newline
\hyperref[sec:oFXXJCLRQwo]{\small{< < < PREV SERMON < < <}}
~
\hyperref[sec:index_chronic]{\small{[返順時目]}}
~
\hyperref[sec:index_scriptual]{\small{[返順卷目]}}
~
\hyperref[sec:QNtUzTDF7W4]{\small{> > > NEXT SERMON > > >}}
\newline
\newline
$^{1}$上次跟大家說到.
父親雅各因為迦南饑荒嚴重.
終於肯讓到愛子別雅敏.
與其餘九個兒子.
同去埃及買糧.
究竟十兄弟今次到迦南.
再見藥室又會有什麼事情發生呢?.
看看第十四節.
在出發去迦南之前.
雅各為了十個兒子.
祝福希望神使到十個兒子.
都能在埃及的手上.
眼前來到蒙恩.
還釋放他們的弟兄西勉.
使別雅敏可以平安回來.
如果結果雅各真的會失去.
十個兒子的話.
雅各都願意接受.
似乎雅各今次的決定.
是帶著一個先死而後生的心情.
雅各的父親雅各.
也願意接受.
先死而後生的心情.
雅各已經作出了心理準備.
和最壞的打算.
看看第十五節.
那裡說到.
今次他們出發去埃及.
讓人可以看見的是.
九兄弟所帶的禮物.
包括銀紙和濕衣弟別雅敏.
不過呢.
今次他們出發去埃及.
讓人看不到的.
究竟是甚麼心情去埃及.
而別雅敏被父親吩咐和九兄長.
同去埃及.
究竟別雅敏又知不知道.
所有的隱情呢?.
今次一去之後.

$^{41}$又說可以平安回來呢?.
當九兄弟帶同別雅敏來到埃及.
約瑟見到親弟別雅敏.
就吩咐人預備筵席.
要和他們一起來共進午膳.
在午膳期間.
九兄長就將迦蘭帶來的禮物.
送給約瑟.
又向約瑟來下拜.
當時約瑟有甚麼反應呢?.
約瑟就問候他們.
也問候他們的父親雅各還健在嗎?.
當約瑟知道父親還健在的時候.
又親眼見到自己和父母的親弟別雅敏.
第29節.
約瑟對別雅敏說.
小兒讓上帝賜恩淵你.
為何約瑟會稱別雅敏為小兒呢?.
有可能約瑟和別雅敏的年紀相差很遠.
無論如何.
約瑟稱別雅敏為小兒.
正好表現出約瑟對別雅敏的親切.
約瑟愛別雅敏.
還為別雅敏帶上祝福.
加上約瑟是一個重情又重義的人.
所以一時感觸下.
約瑟就轉身到自己的房間裡大哭一場.
大家還記不記得上一次.
為何約瑟要退去大哭一場呢?.
約瑟當時已經是一個兩個兒子的父親.
也是埃國的一國首相.
還有就是.
他也是買賣糧的最高領導人.
以約瑟當年的身份和地位.
而是一個幸福的家.
根本沒有需要為二十多年前所發生的事而大哭.
除非約瑟當年所遭遇的禍患.
所帶來給約瑟的心靈傷害.
一直都沒有機會正視和處理.
聖經沒有交代過.

$^{81}$約瑟當年要留在洛爾鄉.
究竟有沒有被埃及的人虐待過.
試想一下當時約瑟.
只不過是一個十七歲的青年人.
當年他身上所發生的事.
對一個十七歲的年輕人來說.
的確造成心靈上很大的傷害.
約瑟第一次退去大哭的一場.
是聽大哥流便.
重提當年的往事.
當年.
十兄長有殺自己的動機.
也有賣自己的行動.
約瑟的過去痛苦經歷.
好像電影一樣一幕一幕.
在他腦海中浮現.
上一次約瑟退去大哭.
是因為他傷痛的記憶.
觸動了他的情緒.
這次約瑟的大哭.
我認為是有分別的.
這次約瑟的心情和目光.
完全只是注視一個人的身上.
就是他親弟別雅敏.
在《聖經》中沒有交代過.
約瑟和別雅敏兩兄弟.
究竟年紀相差有多大.
也沒有交代過.
兩兄弟的感情究竟有多深.
不過很肯定的是.
他們是同父同母所生的親兄弟.
而最重要的是.
約瑟為何一直排斥.
質疑出賣約瑟.
只是十兄長而不是別雅敏.
所以我認為這次約瑟的大哭.
純粹是因為親眼見到自己的親弟而感觸.
約瑟上一次是傷心的感觸.
這次是開心的感觸.
約瑟哭過之後.

$^{121}$約瑟再回來.
和他的十兄弟共進午膳.
想知道故事會繼續如何發展下去?.
請留由下一次我和大家的講解.
多謝您收睇時局新聞,再會!.
\newpage



\section{}
\label{sec:QNtUzTDF7W4}
\textbf{11十一少被賣記因禍得福－下(二)}
\newline
\newline
連結: \href{https://youtube.com/watch?v=QNtUzTDF7W4}{\texttt{https://youtube.com/watch?v=QNtUzTDF7W4}} ~~~~ 語音日期: 2020-07-08
\newline
\newline
\hyperref[sec:2kNd6PvPzMU]{\small{< < < PREV SERMON < < <}}
~
\hyperref[sec:index_chronic]{\small{[返順時目]}}
~
\hyperref[sec:index_scriptual]{\small{[返順卷目]}}
~
\hyperref[sec:X74nEQWZ90I]{\small{> > > NEXT SERMON > > >}}
\newline
\newline
$^{1}$上一次講到.
王和各大臣都相信.
神與約瑟的同在.
於是就任命約瑟做埃及的宰相.
我們一起再看看40到46節.
那裡講到.
王就派約瑟管理住王的家.
除了王一個之外.
所有王的子民都要聽從約瑟.
王還派約瑟治理埃及全地.
如果沒有約瑟的命令.
任何人都不可以主作主張.
王又將自己手上所戴的戒指.
都脫下來戴在約瑟的手上.
王手上的戒指其實是一個印記.
是代表著身份和權力.
王又將細麻衣今年都給約瑟穿上.
細麻衣和金鏈都表示著尊貴.
王要將所有大臣和人民都知道.
約瑟做了宰相.
於是王就安排約瑟坐在自己的車上.
向所有人宣告.
王又為約瑟改了一個埃及的名字.
王亦將祭司的女兒下嫁給約瑟.
埃及祭司的身份和地位在埃及是非常高的.
約瑟能夠成為祭司的女婿.
同時都可以提高了約瑟的身份.
和鞏固了約瑟在埃及的地位.
大家移民來到外國.
多數都會替自己改一個英文名.
大家知道改英文名有什麼意思嗎?.
有些移民落地生根.
除了改名投入工作生活之外.
還會進入到本地的社群.
嘗試一下同文化.
例如和本地人做朋友.
甚至結婚生兒育女.
約瑟的情況就是要和埃及人來去同文化.
否則的話.
約瑟沒有辦法只能夠當上埃及人的首相.

$^{41}$管治埃及.
所以王就將祭司的女兒下嫁給約瑟.
埃及祭司的身份地位高.
能夠成為祭司的女婿的確是助長.
約瑟日後治理埃及.
約瑟日前.
有一個埃及囚犯被王提升成為埃及首相.
還得到王的信任.
還用治理埃及傳遞.
所有埃及人除了王一個之外.
都要聽命於約瑟.
王將財富甚至祭司的女兒都下嫁給約瑟.
又為約瑟改一個埃及名.
約瑟殉失之間.
他重獲曾經所失去的權,名,利和自由.
但這次他所得的.
比之前在埃及的主人波提弗的家中所得的.
更多更尊貴.
但從人的角度去看.
約瑟他因禍得福.
理應故事的發展來到此應該結束.
不過呢.
從神的角度.
約瑟的故事還未完結.
因為約瑟在埃及做首相.
是要經歷七年的風盛.
但同時亦要經歷七年的饑荒.
這樣的安排其實一直有神的旨意在內.
我們一起看看53到57節.
那裡說到七年的風修過去後.
隨即就有七年的饑荒.
當時除了埃及外.
各地都一樣有大饑荒.
不過約瑟在七年風修期間.
因為他預知有饑荒的緣故.
於是他就儲起了大量的糧在倉庫裡.
以防七年饑荒的來臨.
結果.
七年的饑荒年引起了埃及以外的人.
都來到埃及買糧.

$^{81}$就連約瑟的原生家庭.
都因為大饑荒的缺糧的緣故.
於是約瑟的十兄長都來到埃及買糧.
當然.
十個兄長來到埃及買糧.
必定會再遇上失去了聯絡的親弟約瑟.
兄弟再會面將會有什麼事情發生?.
想知道故事的發展.
就不要錯過下一次和大家的講解了.
\newpage



\section{}
\label{sec:X74nEQWZ90I}
\textbf{21 十一少被賣記大團圓－下(一)}
\newline
\newline
連結: \href{https://youtube.com/watch?v=X74nEQWZ90I}{\texttt{https://youtube.com/watch?v=X74nEQWZ90I}} ~~~~ 語音日期: 2020-08-02
\newline
\newline
\hyperref[sec:QNtUzTDF7W4]{\small{< < < PREV SERMON < < <}}
~
\hyperref[sec:index_chronic]{\small{[返順時目]}}
~
\hyperref[sec:index_scriptual]{\small{[返順卷目]}}
~
\hyperref[sec:xOb08azcrmI]{\small{> > > NEXT SERMON > > >}}
\newline
\newline
$^{1}$上次說到雅各舉家到埃及之前.
雅各帶同子女,孫子和所有後裔.
包括生畜和財物.
到遍士巴這個地方.
雅各在遍士巴這個地方.
他先來到鄉耶和華獻祭.
雅各求告耶和華神.
究竟雅各為何要先求告耶和華神呢?.
我們在經文46章3-4說.
神在異象當中和雅各的對話.
神說:我是神,就是你父親的神.
你下埃及去,不要害怕.
因為我必使你那女成為大族.
我要和你同下埃及去.
在雅各的異象當中,雅各很清楚知道.
離開迦南,去到埃及定居是神的心意.
而且神應許過一定與雅各同行到埃及.
也會與雅各同回到迦南.
神又在異象當中鼓勵雅各.
也提醒雅各.
神的應許將來雅各的家庭.
必成為大族.
於是雅各安心帶同所有人和生畜.
乘坐法老的車離開迦南.
在人的立場.
雅各當然重視與約瑟的大團圓.
但雅各更看重的是.
神是否喜悅整個家庭.
究竟喜不喜悅整個雅各家族.
由離開迦南,移居到埃及.
當雅各在比斯巴憲制後.
雅各肯定耶和華神的心意.
耶和華願意同在同行.
雅各才肯放心舉家前往埃及.
當年,大家舉家,移民.
又或者要到外國讀書.
甚或打算到其他城市,國家發展時.
不知道大家有沒有像雅各一樣.
先來求告神.
很清楚知道是神的心意.

$^{41}$才最後決定和行動呢?.
又或者,當我們要復和一段已經破壞的關係.
希望已經破壞的關係可以得到復和.
最後能夠大團圓結局.
不知道大家有沒有試過邀請神的介入呢?.
甚或我們有沒有為別人的復和而祈禱.
讓神介入在一段已經破壞的關係中呢?.
經文29-30.
當約瑟知道父親已經到達歌山這個地方.
約瑟就前往歌山,親自迎接父親.
可見當時約瑟極之心急.
畢竟約瑟和親父已經分隔二十多年.
一旦父子重逢.
大家可以想像一下當時的情景.
一個四十多歲的中年男人.
竟然像小孩子一樣.
伏在父親的頸上大哭.
不知道大家有沒有試過呢?.
像約瑟一樣伏在自己的父親頸上大哭.
我相信當時的場面都很感動.
而且是一幅相當美麗的圖畫.
彼此之間的情景流露.
又是否我們渴望經歷的場面呢?.
還是只是在戲劇中我們才可以看見.
當父子重逢,一家人可以團聚之後.
接下來,故事又會如何發展呢?.
下一次,我會再和大家繼續探討.
\newpage



\section{}
\label{sec:xOb08azcrmI}
\textbf{19 十一少被賣記大團圓－中(三)}
\newline
\newline
連結: \href{https://youtube.com/watch?v=xOb08azcrmI}{\texttt{https://youtube.com/watch?v=xOb08azcrmI}} ~~~~ 語音日期: 2020-08-05
\newline
\newline
\hyperref[sec:X74nEQWZ90I]{\small{< < < PREV SERMON < < <}}
~
\hyperref[sec:index_chronic]{\small{[返順時目]}}
~
\hyperref[sec:index_scriptual]{\small{[返順卷目]}}
~
\hyperref[sec:EEqL82eh5Po]{\small{> > > NEXT SERMON > > >}}
\newline
\newline
$^{1}$「上次講到約瑟因為高興九兄長能夠將別雅敏從埃及帶來.
於是就為他們設吾宴.
在用膳的時候約瑟因為開心而感觸.
就走回房間大哭一場.
當約瑟哭完之後.
再回來和十兄弟一起共午膳.
接下來又發生什麼事呢?.
看看44章,一到第三節.
約瑟他吩咐人將糧食裝滿在他們的口袋裡.
約瑟又吩咐人將各人的銀紙都放回在各人的口袋.
不過大家小心留意一下.
約瑟同時吩咐人將銀貝放在那個少年人的口袋.
也就是別雅敏的口袋裡.
到了第二天,他們起程要回迦南了.
在回程途中,約瑟吩咐人去追捕他們.
根據第九節說.
無論在任何人的口袋裡搜出銀貝.
這個人就要死.
結果人贓並獲.
在別雅敏的口袋裡就找到這個銀貝.
於是九兄弟只有斯里依福回到埃及.
斯里依福是什麼意思呢?.
原來斯里依福這個行動表示他們相當憂傷.
他們為了小兒子別雅敏所遇的事來到憂傷和難過.
看看44章,14到45節.
當他們回到城之後.
是猶大和埃及的首相約瑟的對話.
從猶大的對話中可以聽得出的是.
猶大不願意別雅敏一個人來受罪.
所以就向埃及的首相來求情.
希望首相能夠紀念別雅敏.
在家中裡還有一個老父.
以下就是猶大請求的內容.
猶大這樣說.
「別雅敏是老父在晚年時候所生的一個兒子.
別雅敏的同父同母的哥哥已經死了.
別雅敏的親母就只剩下她一個兒子.
父親非常痛愛別雅敏.
如果父親死了.
父親必定會傷心至死.

$^{41}$猶大也為了救家族和所有的人.
更不願意見到老父失去了最愛的兒子.
而傷心至死.
猶大就用自己的性命.
向老父來擔保.
他擔保可以安全將別雅敏帶回老父的身邊.
如果你就是約瑟.
聽到這個兄長猶大的請求.
你會有什麼反應呢?.
猶大就是當日提議要殺約瑟的人.
在約瑟的心目當中.
十兄長之中.
猶大是最沒有兄弟情的一個.
也是約瑟最不信任的一個.
如今 約瑟卻親耳聽到.
猶大衛了親情.
竟然將自己的性命作為擔保.
大家認為對約瑟一個這麼重情.
又這麼重義的人來說.
約瑟對猶大這番請求.
會有什麼反應和行動呢?.
看完45章.
一到第15節.
約瑟吩咐所有人退去.
只剩下他自己和十兄弟的時候.
約瑟就和十兄弟相迎.
約瑟和十兄弟相迎的時候.
他大聲地吐哭.
這次是約瑟第三次在哭.
聲音大得連外面的埃及人.
和法老家中的人都哭了.
可想而知 當時的約瑟.
已經有心理準備.
還有行動和兄長團圓和復和.
不過 當時兄弟都被約瑟的表現嚇呆了.
於是約瑟就向他們表明.
自己的真實身份.
就是當年他們賣他的那一張照片.
約瑟不但表明自己的身份.
看完第五節.

$^{81}$約瑟先向他們指出.
不要再為當年的事而自憂自恨.
當時約瑟這番話足以證明.
約瑟當年被兄長出賣的事.
其實約瑟已經釋懷了.
約瑟的身份和身份.
其實約瑟已經釋懷了.
弟姐妹 冤尊你就是約瑟.
你會像約瑟一樣.
肯放下二十多年來兄弟之間的怨恨嗎?.
如果想知道故事發展會是怎樣.
請不要錯過下一次的講解.
\newpage



\section{}
\label{sec:EEqL82eh5Po}
\textbf{聖地網上行 - 步主腳蹤 [第三集] 約旦河 - 曠野}
\newline
\newline
連結: \href{https://youtube.com/watch?v=EEqL82eh5Po}{\texttt{https://youtube.com/watch?v=EEqL82eh5Po}} ~~~~ 語音日期: 2020-08-22
\newline
\newline
\hyperref[sec:xOb08azcrmI]{\small{< < < PREV SERMON < < <}}
~
\hyperref[sec:index_chronic]{\small{[返順時目]}}
~
\hyperref[sec:index_scriptual]{\small{[返順卷目]}}
~
\hyperref[sec:QHoIKk7ZHEQ]{\small{> > > NEXT SERMON > > >}}
\newline
\newline
$^{1}$各位我們又來到聖地這個地方了.
今天是第三集.
第一集我們去了白尼恆耶穌出生的地方.
第二集我們去了那撒勒,就是耶穌成長的地方.
在那裡大概有三十年的時間.
然後才出來公開傳道.
其實在公開傳道之前.
他的表哥施洗約翰.
其實已經在猶太地的約旦河一帶開始施洗.
叫人悔改受洗.
然後預備天國預備彌賽亞的里林.
約翰福音一章二十八節.
這是在約旦河外約旦河東伯大利.
Bethany約翰施洗的地方所作的見證.
伯大利約旦河東的伯大利叫做Kasar al-Yahu.
意思是說The Tower of the Jews.
就是猶太人的亡流.
在耶利哥以東十公里或者五六英里左右.
因為耶穌受洗是耶穌事蹟裡面最重要的一個.
四福音都有記載.
所以自四世紀以來.
很多教會和宗派會來到Kasar al-Yahu去朝聖.
設立很多教堂和修道院.
來紀念耶穌受洗的事蹟.
所以這裡又叫修道院之地.
Land of Monasteries.
因為有很多修道院在這裡.
現在看到的這一帶至少有七間舊教堂.
不過1967年以色列六日戰的時候.
以色列佔據了這個地方.
因為戰局一直很動盪.
所以在這裡以軍埋設了三千多個地雷.
所以過去五十年都成為軍事重地.
朝聖者和遊人都來不了這個地方.
直到2011年才逐步開放.
所以你看到有不少地方仍然被攔住.
很多地方都有地雷的告示牌.
直到2019年政府才決心要清理所有地雷.
所以現在已經差不多全部開放了.
你看到很多教徒來自世界各地來這個地方受洗.

$^{41}$有些第一次 有些是重新受洗.
這裡我們看到一些埃塞俄比亞裔的信徒.
受洗是一個很隆重的儀式.
所以他們在這裡慶祝.
然後我們看到下面就是約旦河.
對面就是今天的約旦王國.
只是一條小河去分界.
對面就是聖經所說的約旦河東或約旦河外.
Beyond the Jordan.
也有古教去紀念聖經的事蹟.
施洗約翰的事蹟.
或舊約先知以利亞.
你記得在這裡是升天的.
我所說的所謂約旦河東和所謂約旦河外.
為什麼呢?.
因為大家知道河流在平原的地方不停地改道.
歷史上早期的文獻已經記錄了約旦河改道好幾次.
所以到底我們現在站的地方是約旦河東還是約旦河西.
要視乎考古的發現.
但因為這裡仍然是非常敏感的地方.
所以還沒有這樣的研究.
這裡我們和兩個以色列的女兵聊天.
河的兩邊都是打手的森嚴.
今天的約旦河有兩個可以公開受洗的地方.
第一個就是這裡,其實是阿爾也河.
另一個地方就是在北邊的,叫做雅典尼.
就是小約旦河的意思.
北邊加里里海出約旦河的地方就是雅典尼的地方.
1981年以色列政府特別開設.
因為這裡南面的地方是軍事重地.
所以開設了另一個地方.
那裡你可以看到雅典尼的水因為在北邊.
所以清潔很多.
也有較多教會去那個地方.
現在這一處是阿爾也河.
就是約旦河的下流.
所以水的泥比較多.
加上有很多分流.
有工業污染等等.
水就比較足.

$^{81}$相片可以看到我只是浸腳.
然後就洗腳.
不要下水,我沒有他那麼大膽.
小心一點會比較好.
來到這個地方是施洗約翰開始傳道的地方.
猶太的曠野的其中一部分.
信息的中心是來自以賽亞書第43章.
馬達福音第3章記載.
他說那時施洗的約翰出來.
在猶太的曠野傳道說.
天國近了,你們應當悔改.
先知以賽亞所說.
在曠野有呼喊者的聲音.
預備主的道收藉他的路.
所以先知以賽亞所預言來到的先鋒就是施洗約翰.
在約旦河東的曠野.
其實有很多聖經故事都是跟這裡有關.
譬如你想起舊約的先知.
以利亞在這裡升天.
神應許末後天國來到的時候.
以利亞會再來.
那時人們會迴轉成為新的神國的子民.
新的以色列民.
怪不得施洗約翰的裝束是很像以利亞.
除了以利亞的故事在這裡發生.
跟施洗約翰有關.
這裡約旦河的渡口也是昔日約書亞帶領以色列人.
過約旦河.
入主迦南去建立天國的時候.
這些都是聖經裡面舊約先知的預言.
上帝的應許.
神災來的時候.
神國要建立的時候的景象.
你會想起過河過海.
出埃及.
想起出埃及之後到曠野.
然後過約旦河.
然後想起先知來到以利亞等等.
這些景象都是當時很多猶太人所盼望的.
在附近有一個地方叫做昆蘭.

$^{121}$早早有一群猶太人搬到這個地方.
他們很可能就是一世紀裡面經外文獻提過的愛錫里人.
Asim.
這群人是脫離早已經腐敗的耶路撒冷聖殿的體制.
他們不滿當時的權貴.
祭司來靠攏羅馬帝國欺壓人等等.
所以他們就搬到曠野昆蘭這個地方.
另外設立一個新的群體.
他們的身份都是認為他們是新的以色列民.
昆蘭這個地方當然留下了二十世紀裡面最大的考古發現.
叫做死海孤卷.
這些昆蘭人是很自律和自潔的.
泛物公庸.
每一天他們都會在潔淨池受洗.
等候自己天國王者救主的來臨.
所以有人認為.
施洗約翰甚至耶穌曾經和這些昆蘭的社群有所接觸.
昆蘭的遺址在曠野裡面生存是很難的.
但他們建了很多引水道.
然後把這些水放在一些井.
在這個遺址上也留下了很多房屋.
他們抄寫聖經的地方.
在這個遺址上有一個很重要的特色就是潔淨池.
他們按照樓梯走下去.
然後經過水受洗.
然後再上來的時候就得到潔淨.
起碼禮儀上的潔淨.
他們就預備在社群裡面生活.
預備去抄寫聖經.
預備去親近神.
預備去看神的華語.
我再講施洗約翰.
施洗約翰在約旦河受洗.
是一次性的受洗.
這個猶太人的傳統很多時候是外邦人入教的時候就需要受洗.
猶太人是不需要.
但這裡特別的地方就是施洗約翰的洗禮是為所有人而設的.
連猶太人自命是神的子民的後裔.
這些人都要在施洗約翰的事奉下去受洗.
換言之對施洗約翰的事奉來說.

$^{161}$猶太人和外邦人是沒有分別的.
當上帝的天國 上帝的統治要來臨的時候.
每一個人都要悔改受洗.
當時施洗約翰面對的政局.
聖殿的權貴是貪腐.
羅馬軍隊是濫權.
本地的王 希律王等等是很腐敗.
施洗約翰不單止是叫人悔改受洗.
並且叫人成為神國的子民去行公二號聯運.
所以他講了很多先知的講道.
他也是救主的先鋒.
因為這個緣故施洗約翰吸引了很多猶太人成為他的門徒.
但他只不過是救主的先鋒.
當救主來到的時候.
耶穌來到的時候.
他就將人指向彌撞亞.
他就退出.
沒多久耶穌就出現了.
馬太方第三章十三節.
那時耶穌從加利利即北面.
來到約旦河約翰那裡即南面.
要受他的洗.
約翰想要阻止他.
說我應該受你的洗.
你怎麼會來我這裡呢.
這是一個很真實的記載.
因為按照猶太人的傳統.
甚至初期教會的傳統.
施洗者的地位是高過受洗者的.
所以如果是虛構的話.
一定會刪除這一段.
不過四福音都有記載.
耶穌是接受約翰的施洗.
所以為什麼約翰想要攔住他.
耶穌怎麼回答呢.
他說暫且這樣做.
我們理當這樣履行全部的義.
和學本就是諸般的義.
於是約翰就答應了他.
耶穌受了洗.

$^{201}$立刻從水中上來.
忽然天為他開了.
他就看見神的靈.
好像甲子降在他身上.
又有聲音從天上說.
這是我的愛子.
我所喜悅的.
這就像是一個公開的見證.
告訴人耶穌基督真正的身份是什麼.
新譯本就叫做理行全部的義.
和學本就是諸般的義.
其實是什麼意思呢.
義就是按照上帝的心意.
意思就是人子.
耶穌基督要完全認同我們每一個罪人.
我們每一個都需要悔改.
受洗進入天國.
雖然他沒有罪.
但因為他認同我們.
作為我們的代表.
他就要經歷我們所經歷的事.
正如以色列人所經歷的事.
也是我們每一個人所經歷的.
你記得上兩集說過.
耶穌的出生.
是在暴君大希律的追殺下出生.
正如以色列人曾經在埃及.
被暴君法奴所追殺一樣.
出生之後.
耶穌也是流亡在異地.
在埃及.
然後就離開埃及.
是上帝帶他出來的.
正如神曾經救以色列人.
離開埃及.
慰勞之地.
曾經被流亡.
但又回歸.
以色列人經過紅海或約旦河.
好像死裡重生一樣.

$^{241}$耶穌也是經歷約旦河的受洗.
以色列人經歷曠野四十年.
受上帝來試煉他們.
耶穌同樣在曠野四十天受試煉.
所以耶穌成為我們的代表.
能夠拯救我們一群罪人.
耶穌受洗後.
他就去到曠野受試探四十天.
《馬太方第四章》記載.
隨後耶穌被聖靈帶到曠野.
受魔鬼的試探.
耶穌禁食了四十晝夜.
就餓了.
試探者前來對他說.
你若是神的兒子.
就吩咐這些石頭變成食物吧.
耶穌回答.
經常記著.
人活著不是單靠食物.
是靠神口裡所出的一切.
你看到猶太的曠野.
滿地都是石頭.
所以耶穌肚餓的時候.
這個試探者.
撒旦魔鬼.
就來引誘他.
就要他濫用他的能力.
然後將這些石頭變成食物.
第五節.
隨後魔鬼帶耶穌進了聖城.
使他站在殿的最高處.
對他說.
你若是神的兒子就挑下來吧.
因為經常記著.
神為了你.
會吩咐自己的使者用手托著你.
免得你的腳碰到石頭.
耶穌對他說.
經常有記著.
不可試探主你的神.

$^{281}$這裡說殿的最高處.
大概是耶路撒冷聖殿的聖殿山.
聖殿平台的東南一角.
下面一望下去就是吸羊谷.
是當時最高的地方.
這裡又是魔鬼來試探耶穌.
叫他濫用自己的權柄.
來試探神.
第八節到第十一節.
最後魔鬼帶著耶穌.
上了一座極高的山.
將世界各國的榮華都賜給他.
並且對他說.
你只要跪下來拜我.
我就把這一切都給你.
但耶穌說.
撒旦走開.
經常記著.
當拜主你的神.
單要事奉他.
於是魔鬼就暫時離開耶穌.
有天使來服侍他.
這裡馬太記載耶穌的三個試探.
很有意思.
因為三個試探是越來越高的.
由開始石頭變成食物.
然後帶他到聖殿的最高一樹.
然後到最高的山.
更值得留意的是.
耶穌三次的回應.
都是出自《申命記》六至八章.
那一處說的是什麼呢.
就是摩西引導了以色列人教導他們.
在進入迦南之前.
他們需要去中心敬拜主.
就跟他們回顧以色列人的故事.
你記得那個時跡.
當以色列人從埃及.
上帝救他們出來之後.
他們在曠野.

$^{321}$以色列人失敗過三大次.
第一次是什麼呢.
就是他們不信神.
沒有食物的時候他們發怨言.
詩篇後期這樣記載.
他們忘記了神的作為.
不仰望神的意志.
反而在曠野大起貪慾.
在曠野沙漠裡試探神.
第二次失敗就是他們駐造偶像.
詩篇第106篇也是這樣說.
他們在賀列山做偶.
來到敬拜主神的偶像.
他們竟然忘記了拯救他們的神.
第三次失敗就是他們沒有水.
詩篇第106篇這樣說.
他們去抗拒神的靈.
以色列人在曠野三次失敗.
但是這裡講清楚.
耶穌是真正的以色列人.
真正的神的意志.
他在曠野三次得勝.
所以告訴我們所有在耶穌基督裡面的人.
我們都能夠分享到這個特性.
在耶利哥附近有一個山叫做.
試探山Mount of Temptation.
上面是一個歷史悠久的修道院.
紀念的當然是耶穌40日受試探.
說起這些修道院.
其實是教會歷史上.
信徒學到的功課的遺產.
這些修道院到底是怎樣來的呢?.
修道主義最高峰的時期.
就是公元後四至六世紀.
最盛行的時候.
單在巴勒斯坦地區.
就有73間大大小小的修道院.
這些修道士保留了很多.
教會歷史的信仰遺產.
譬如聖經的手抄本.

$^{361}$一些學術哲學的著作.
或者修道士做一些醫療接待的服事等等.
並且也培育了很多出色的學者.
好像有一個很出名的John of Damascus.
大馬士革的約翰.
他生於大馬士革.
但他餘下一生的歲月.
就是在伯利恆附近的一間修道院.
叫做Massaba.
專心著作.
事實上是去港島的.
並且出去跟別人辯論.
宗教信仰和社會政治等等的題目.
值得留意的是.
修道院盛行的時間是公元後四至六世紀.
當時是羅馬帝國拜占庭的時候.
那個國家已經定了基督教成為國教.
對信仰有什麼侵蝕呢?.
慢慢地信仰變成很官方.
變成很形式化.
甚至很政治化.
修道院其實是一種抗衡的生活.
是一種尋求更新的運動.
修道就是你暫時或一生.
你決心脫離帝國都市的建制.
你就引退到曠野.
過一些清靜簡單的生活.
禁食等等.
甚至有些修道士會在山洞裡生活.
他們只會去迎見上帝.
在曠野裡.
他們的生活只靠上帝.
耶穌也有類似的習慣.
祂很自律很恆常.
暫時避開複雜表面華麗的建制人群.
所以科林書這樣描述耶穌的靈性生活.
不是靠耶路撒冷華麗的建制聖殿.
祂時常獨自退出來.
然後單單去到親近神.
然後才回到城市裡面面對人群.

$^{401}$路加福音第五章.
但耶穌的名聲越發的傳揚出去.
成群的人就來到去聚集.
要聽祂的說道.
並且要洗他們的疾病存餘.
第十六節.
耶穌卻退到曠野去禱告.
有些人會因為宗教的建制很不滿或受傷.
一拼就跟上帝割席.
但你可能需要問.
會不會這樣是以偏概全.
會不會失去了很多東西.
修道主義修道院就是邀請我們退到曠野.
在那裡沒有任何令我們分心的東西.
沒有東西可以去怪責.
你單單去到曠野.
只有你和上帝單獨相遇.
這叫做搶不走靈魂的信仰生命.
在今天我們的世界很混亂很壓迫.
耶穌也是一樣.
所以跟隨他的教會.
跟隨耶穌的門徒.
都應該是這樣.
在你的生命裡面.
有沒有曠野的空間.
你能夠單獨和上帝相遇.
現在我想大家感受一下修道院是一個怎樣的氣氛.
我帶你去不同的地方.
看看不同的修道院.
還有背景音樂是一班東正教修士唱詩篇141篇.
修道院.
大小姐,您是來自哪里的?.
我出國,访问寺廟的..
你又回來,几年了?.
大小姐,您是來自哪里的?.
我出國,访问寺廟的..
你又回來,几年了?.
我出國,访问寺廟的..
你又回來,几年了?.
我出國,访问寺廟的..

$^{441}$你又回來,几年了?.
我出國,访问寺廟的..
你又回來,几年了?.
我出國,访问寺廟的..
你又回來,几年了?.
我出國,访问寺廟的..
你又回來,几年了?.
我出國,访问寺廟的..
你又回來,几年了?.
我出國,访问寺廟的..
你又回來,几年了?.
我出國,访问寺廟的..
你又回來,几年了?.
我出國,访问寺廟的..
你又回來,几年了?.
我出國,访问寺廟的..
你又回來,几年了?.
我出國,访问寺廟的..
你又回來,几年了?.
我出國,访问寺廟的..
你又回來,几年了?.
我出國,访问寺廟的..
你又回來,几年了?.
我出國,访问寺廟的..
你又回來,几年了?.
我出國,访问寺廟的..
你又回來,几年了?.
我出國,访问寺廟的..
你又回來,几年了?.
我出國,访问寺廟的..
你又回來,几年了?.
我出國,访问寺廟的..
你又回來,几年了?.
我出國,访问寺廟的..
你又回來,几年了?.
我出國,访问寺廟的..
你又回來,几年了?.
我出國,访问寺廟的..
你又回來,几年了?.
我出國,访问寺廟的..

$^{481}$你又回來,几年了?.
我出國,访问寺廟的..
你又回來,几年了?.
我出國,访问寺廟的..
你又回來,几年了?.
我出國,访问寺廟的..
你又回來,几年了?.
我出國,访问寺廟的..
你又回來,几年了?.
我出國,访问寺廟的..
你又回來,几年了?.
我出國,访问寺廟的..
你又回來,几年了?.
我出國,访问寺廟的..
你又回來,几年了?.
我出國,访问寺廟的..
你又回來,几年了?.
我出國,访问寺廟的..
你又回來,几年了?.
我出國,访问寺廟的..
你又回來,几年了?.
我出國,访问寺廟的..
你又回來,几年了?.
我出國,访问寺廟的..
你又回來,几年了?.
我出國,访问寺廟的..
你又回來,几年了?.
我出國,访问寺廟的..
你又回來,几年了?.
我出國,访问寺廟的..
你又回來,几年了?.
我出國,访问寺廟的..
你又回來,几年了?.
我出國,访问寺廟的..
你又回來,几年了?.
我出國,访问寺廟的..
你又回來,几年了?.
我出國,访问寺廟的..
你又回來,几年了?.
我出國,访问寺廟的..

$^{521}$你又回來,几年了?.
我出國,访问寺廟的..
你又回來,几年了?.
我出國,访问寺廟的..
你又回來,几年了?.
我出國,访问寺廟的..
你又回來,几年了?.
我出國,访问寺廟的..
你又回來,几年了?.
我出國,访问寺廟的..
你又回來,几年了?.
我出國,访问寺廟的..
你又回來,几年了?.
我出國,访问寺廟的..
你又回來,几年了?.
我出國,访问寺廟的..
你又回來,几年了?.
我出國,访问寺廟的..
你又回來,几年了?.
我出國,访问寺廟的..
你又回來,几年了?.
我出國,访问寺廟的..
你又回來,几年了?.
我出國,访问寺廟的..
你又回來,几年了?.
我出國,访问寺廟的..
你又回來,几年了?.
我出國,访问寺廟的..
你又回來,几年了?.
我出國,访问寺廟的..
你又回來,几年了?.
我出國,访问寺廟的..
你又回來,几年了?.
我出國,访问寺廟的..
你又回來,几年了?.
我出國,访问寺廟的..
你又回來,几年了?.
我出國,访问寺廟的..
你又回來,几年了?.
我出國,访问寺廟的..

$^{561}$你又回來,几年了?.
我出國,访问寺廟的..
你又回來,几年了?.
我出國,访问寺廟的..
你又回來,几年了?.
我出國,访问寺廟的..
你又回來,几年了?.
我出國,访问寺廟的..
你又回來,几年了?.
我出國,访问寺廟的..
你又回來,几年了?.
我出國,访问寺廟的..
你又回來,几年了?.
我出國,访问寺廟的..
你又回來,几年了?.
我出國,访问寺廟的..
你又回來,几年了?.
Just give me Jesus.
《與生俱來》 作词:李宗盛 作曲:李宗盛.
演唱:李宗盛.
Give me Jesus.
当我安静在主前.
心中唯一的思念.
只要耶穌.
拿走我的世界.
拿走我的一切.
只要耶穌.
是.
\newpage



\section{}
\label{sec:QHoIKk7ZHEQ}
\textbf{聖地網上行 - 步主腳蹤 [第二集] 拿撒勒}
\newline
\newline
連結: \href{https://youtube.com/watch?v=QHoIKk7ZHEQ}{\texttt{https://youtube.com/watch?v=QHoIKk7ZHEQ}} ~~~~ 語音日期: 2020-08-26
\newline
\newline
\hyperref[sec:EEqL82eh5Po]{\small{< < < PREV SERMON < < <}}
~
\hyperref[sec:index_chronic]{\small{[返順時目]}}
~
\hyperref[sec:index_scriptual]{\small{[返順卷目]}}
~
\hyperref[sec:KnqUXPbp8Kk]{\small{> > > NEXT SERMON > > >}}
\newline
\newline
$^{1}$各位團友,我們再一次來到聖地以色列巴勒斯坦地.
你記得上一集我們查訪過耶穌出生的事跡.
我們去過伯利恆.
你記得故事內容大概是這樣.
約瑟瑪利亞當時他們在北面加利利一個小鎮叫做拿撒勒.
因為羅馬帝國要迫人民回鄉到布水.
所以約瑟利亞帶著有生人的瑪利亞.
他就回到自己的家鄉猶太地的伯利恆.
在那裡耶穌剛好出生.
可能因為在伯利恆有其他家人親屬照顧.
所以他們暫時留在伯利恆.
誰知到耶穌大約兩歲的時候.
大希律王就想追殺耶穌.
然後他們有神的指引.
他們一家就逃難到埃及成為難民.
大希律不久就死了.
他的國土就分給他三個兒子.
第一個兒子叫做阿基努.
他管理猶太地,撒瑪利亞和以東地.
第二個兒子叫做安提柏.
叫希律安提柏.
他管北面的加利利省和約旦河東的一處比利亞.
然後東北面的巴山地.
現在我們叫做哥蘭高地.
就交給他定愛的第三個兒子.
叫做希律菲力菲勒.
約瑟聽聞大希律王死了之後.
當時耶穌還小.
本來他們想一家回到伯利恆.
可能他覺得至少那裡有本族的親友照顧.
可以照顧小孩.
所以想回到伯利恆.
但他知道阿基努來接管猶太省的時候.
心裡非常不安.
說起阿基努.
他和父王大希律其實都一樣.
生性非常凶殘.
結果阿基努只管了猶太省九年這麼短.
即是公元前四年至公元後六年.
羅馬帝國蓋薩將阿基努革職.

$^{41}$然後流亡到歐洲一個大的地方.
省份叫做高盧.
在那裡的一個城市就是今天的奧地利維也納.
那裡放逐王.
你放心和莫扎特是沒有關係的.
因為是不同年代的事.
無論如何.
之後猶太省由羅馬帝國直接委派一個總督.
來直接管轄.
所以為什麼有秦府或總督彼拉多等人.
他們直接管理猶太省這個地方.
言歸正傳.
神再次指示約瑟.
叫他全家回到北邊的加利利納薩勒.
馬太福音第二章.
這樣解釋耶穌在伯利恆出生.
然後在納薩勒成長.
馬太福音第二章.
希律死後.
這個說的是大希律.
埃及主的使者在夢中向約瑟顯言說.
納薩勒要殺害這孩子的人已經死了.
起來帶著孩子和母親回以色列地去吧.
約瑟就起來.
帶著孩子和母親回以色列地去.
只是聽見阿基努接觸他父親希律作了猶太王.
他就不敢到猶太地去.
猶在夢中得了指示.
於是往加利利境去.
來到納薩勒城住下.
這樣就應驗了先知所說的.
他必稱為納薩勒人.
如果你熟悉舊約希伯來文聖經的話.
你就知道先知沒有一處是明文指出.
彌賽亞救世主是由納薩勒城而出的.
不過在《爾賽亞書》第十一章是這樣說的.
他比喻彌賽亞作為大衛樹幹的枝條.
即是樹的其中一個branch.
大衛國雖然流亡.
但不會最終滅亡.

$^{81}$好像大樹被斬倒下之後.
它不會枯死.
然後大樹的樹幹的枝葉.
就會落地又再生根成為大樹.
這就是《爾賽亞書》第十一章所說的.
從耶西(大衛的父親).
樹幹必生一根聯亞.
從他的根而出的枝條.
Nasser必結果子.
所以大衛樹的枝條的分枝Nasser.
就成為了彌賽亞的稱銜.
其實這不單止《爾賽亞書》這樣說.
在昆蘭發現的死海古卷裡.
也提到彌賽亞的稱號是Nasser.
是大衛樹的枝條.
另外枝條Nasser和納薩勒Nasserine.
發音相似.
所以那些人會把這兩樣東西.
混在一起說.
這是當時猶太人或希羅文學.
常常慣用的辯證方法.
同音似音的字.
他們就放在一起.
所以你看看小小的主耶穌.
他的人生裡是迂迴曲折.
道成肉身的主宰.
他正在經歷難民的生涯.
是居無定所的.
耶穌的經歷就是每一個以色列人的經歷.
曾經走難到埃及.
然後上帝又帶他出埃及.
然後又流亡又回歸.
天國仍然是向前行的.
在天父自己安排裡.
逐步實現舊約聖經裡的應許.
所以這裡告訴我們.
神最終是信實不變.
其實你想想我們的人生.
有時我們會覺得生於亂世.
政局很不穩.

$^{121}$人生路上有順.
很多時也有逆.
我不知道你有沒有經歷過神的帶領.
但是這個故事.
耶穌我們的主道成肉身.
就是經歷了神的帶領.
其實我們自己回望.
在我們自己人生當中.
其實都可能會見到.
神自己的帶領,神的信實.
(阿拉伯首都).
今天第一個去的地方就是拿撒勒.
我們一起飛到以色列北面的城市.
今時今日的拿撒勒不是小城.
是一個很大的城市.
人口有七萬人.
被譽為以色列的阿拉伯首都.
因為七萬人裡面.
很多都是阿拉伯人.
百分之七十是伊斯蘭教徒.
百分之三十是基督徒.
是一個充滿活力的中東城市.
一世紀的拿撒勒.
只不過是一個小鎮.
很可靠的估計.
告訴我們當時最多只有五百人.
在聖地研究當中.
我們的資料是從哪裡來的.
包括聖經,初期教會.
教父寫的文獻.
考古的發現,地理等等.
也包括一些早期的初代基督教朝聖的紀錄.
頭幾百年有不少歐洲各地的基督徒.
他們都來到聖地朝聖.
尋訪聖經裡面的事蹟.
包括耶穌的事蹟.
有一個公元後六世紀的朝聖者.
是這樣記載.
他說拿撒勒的女子.
比其他地方的女子更加漂亮.

$^{161}$為什麼呢.
因為他們有瑪利亞所賜給他們的恩賜.
不過這方面我幫不了你.
因為我沒有這些影片.
不過最有吸引力的女人.
在拿撒勒她敢認第二.
沒有人敢認第一.
(聖經).
現在我進了一個教堂.
叫做報喜堂.
Church of Annunciation.
這個教堂是紀念天使加拔蓮.
向瑪利亞報喜的.
(聖經).
這個報喜堂.
一向都是天主教的聖方濟各.
Franciscans來管理的.
當他們設計動土興建的時候.
他們說有三個大的重點.
可以總結成為3M.
3M不是我們美國公司的3M.
3M是什麼呢.
第一個M就是Multinational.
即是這個教堂代表了很多國際的國家.
第二個M就是Modern.
要有現代感.
第三個M就是Mystery.
即是有神秘感.
你在教堂的內外就感受到這個國際性.
什麼叫Multinationality.
就是你會找到有不同國籍種族.
本土化的聖母瑪利亞像.
張貼在教會的內外.
(聖經).
我們進入教堂的時候.
你就感受到建築的風格是非常現代的.
非常Modern.
神秘感你進入教堂就感受得到.
這個教堂是分兩層建築的.
底層特別設計到是很幽暗的.

$^{201}$蓋住在第一世紀的民居考古遺跡上.
這裡可以看到一個山洞.
一個古跡.
傳統來說這裡是瑪利亞的居所.
是不是真的這個地方呢.
當然沒有可能去確認.
但是剛才所說.
納薩達只不過是一個很小的村莊.
就算不是這個屋.
位置其實都很相近.
祭壇上寫著「道成肉身」.
(聖經).
我們再上去上層.
經過一個很漂亮的樓梯.
現代的裝飾是一個很華麗的教堂.
上層這個教堂可以容納很多人聚會.
連接上下層.
就是一個八角形的天窗.
叫做Oculus.
意思是眼睛.
教堂的頂部可以看到.
是一個圓錐形的拱頂.
把太陽光從上而下射入教堂.
光是從天降下的.
上層代表上帝.
下層代表人間.
之間所連結的就是天窗.
接通了天地神人之間.
這個很有神秘感.
上帝的啟示就是從上而下帶下來.
向瑪利亞這個童貞女報喜.
在天主教的心目中.
當然瑪利亞是很偉大.
我們其實都可以猜測.
瑪利亞是否特別聖潔.
是否特別蒙大恩的女子.
但你回去聖經看的時候.
你就要謹記.
路加福音第一章四十五節.
「這相信的女子是有福的」.

$^{241}$換句話說上帝體重瑪利亞.
不是因為她本身無罪.
也不是因為她本身比其他女子.
更加美或者更加神聖.
而是她對主的信心.
對上帝話語的信服.
你想想這就是聖經信仰.
基督教信仰的本質.
天上的主向我們啟示.
感動我們的時候.
我們在地上的人.
我們謙卑地相信進行.
看完報喜堂在教堂出來之後.
其實你有時間可以去另一端.
有一個教堂叫聖約瑟堂.
那裡當然是紀念約瑟.
教堂的底層是一個洞穴.
上一次說過.
當時的人住在山區.
他們很多時候屋後都是山洞.
這個地方相傳是約瑟的工作室.
他的工作室.
也是約瑟曾經住過的家.
大家可以去猜一下.
等一下我們會告訴你.
約瑟的工作大概是怎樣.
怎樣也好.
那瑟是耶穌成長的地方.
體現道成肉身的道理.
在耶穌成長的經歷裡.
他過著百分之一百平凡人的生活.
你可以想像耶穌.
由出生到長大的經過.
他要學走學行.
學說話.
要學習.
要爭取經驗.
耶穌是百分之一百的人.
沒錯也是百分之一百的神.
但他的人性是百分之一百.

$^{281}$除了沒有罪之外.
他一切都和我們一樣.
所以耶穌有資格成為我們人類的代表.
來到那薩拉的時候除了看補氣堂之外.
有一處地方是一定要去的.
是一個基督教機構所創建出來的.
就叫做Nazareth Village.
是那薩拉村莊.
這是一個很有意思的地方.
這是一個露天的博物館.
他將所有第一世紀的生活狀況.
再重新創建出來.
你走進去的時候.
有不同的義工.
他們會穿上第一世紀的衣服.
做一些第一世紀人會做的事.
還有房屋等等.
很值得在那裡講解.
到底一個住在那薩拉的人.
他的生活是怎樣的.
再一次,這就是人性的傳言經歷和體會.
就是我們的主耶穌基督道成肉身.
他一切的經歷.
這裡有一個第一世紀的墳墓.
是這裡的工作人員親自挖出來的.
模仿第一世紀山洞的墓穴.
耶穌被埋葬的地方就是好像這樣.
附近也有果園.
不同的果樹種了出來.
在山坡上有梯田.
這些梯田仿似聖經時代第一世紀有的梯田.
這裡有一個打和牆的Flushing Floor.
收成了的大麥小麥等等.
就會放在這個和牆上.
然後將它打碎.
然後就拿到裡面的麥子.
有些人會去示範.
到底怎樣去打鞠.
怎樣去壓鞠.
再上一點就有一個猶太人的會堂.

$^{321}$仿效第一世紀會堂的格式所做出來的.
說到耶穌的成長.
《路加福音》就加插了一個獨有的情節.
就是《路加福音》第二章.
每年逾越節.
祂父母都上耶路撒冷去.
當祂十二歲的時候.
他們按照節期的慣例照常上去.
過完了節.
他們回去的時候.
孩童耶穌仍留在耶路撒冷.
祂父母卻不知道.
還以為祂和同行的人中間.
走了一天.
就在親戚和熟人中找祂.
沒有找到就轉回耶路撒冷找祂.
過了三天才發現祂在聖殿裡.
祂的父母見了非常驚奇.
祂母親說:孩子,為什麼只能待我們?.
你看,你父親和我都擔心地在找你.
祂說:為什麼找我?.
你們不知道我必須在我父的家裡嗎?.
《路加福音》二章記載.
耶穌十二歲的時候.
跟祂父母一起去耶路撒冷的聖殿.
去守禦節.
要了解這個情節.
可以幫助我們去耶路撒冷的博物館.
在耶路撒冷的博物館.
我們看到一個1:50的.
第二聖殿時期的耶路撒冷.
所謂第二聖殿.
即是第一聖殿.
所羅門的聖殿已經被毀.
第二聖殿就是猶太人.
在波斯帝國時期.
他們回去耶路撒冷重建的聖殿.
直到第一世紀.
當然大多數是執政的時候.
他們討好猶太人.

$^{361}$然後大興土木.
大大加建耶路撒冷的聖殿.
但我們都會稱為第二聖殿.
在這裡1:50的模型.
我們就能夠看到.
到底這個故事的情節是怎樣發生的.
現在我們看的地方.
大概是鑒欖山上.
鑒欖山是對著耶路撒冷的聖殿.
我們一直看過去.
前面就是聖殿.
耶路撒冷的聖殿外邊.
有一個很廣闊的外院.
是大氣律去加建的.
供朝聖者去敬拜所用.
所以不同的角落.
不同的地方.
就會有一些猶太人.
拉彼和學生來聚集.
一邊講聖經.
群眾就會有問有答.
這段聖經大概是發生在.
露天的平台上的一角.
當時耶穌是十二歲.
接著就十三歲.
按照猶太人的傳統.
男子進入十三歲的時候.
就會成為成人.
那些比較富有的猶太人家庭.
甚至會在這個時候.
舉辦大生日會.
所謂Baptist.
畫面上可以看到.
很多時候在西場.
就會看到一群人大肆慶祝.
這裡我們看到有一個猶太人的男子.
他是Baptist.
第一次去到西場.
這個地方.
其實也是一個公開的會堂.

$^{401}$他去供讀律法書.
猶太人很緊張成人禮.
約瑟和瑪利亞當然是比較清貧的家庭.
我想他們請不到那麼多人.
來大肆慶祝.
但你要留意.
耶穌出生之後.
他們照樣按照律法書所講.
想去耶路撒冷聖殿.
奉獻去到恆國來.
每年他們都按照節期的慣例.
這裡也說他們照常上聖殿.
有如今天一個基督教的家庭.
他們會敬錢的去敬拜.
會盡力去遵守主的懷願.
約瑟和瑪利亞就是一個敬錢的小家庭.
到了十二歲.
耶穌就在會堂或者家庭裡面.
在父母身上.
其實已經學了不少希伯來的聖經.
所以他去見拉比的時候.
有問有答.
其實他就是在繼續學習聖經.
在增長自己的智慧.
增長自己對律法的認識.
這套故事的情節.
大概可以想到.
好像舊約裡面的薩姆爾.
他在回眸裡面.
也是照樣跟著祭司去學習.
敬拜神.
在回眸裡面成長.
或者更加近的更刻意的.
路加其實是將耶穌的成長.
對比施洗約翰他Cousin的成長.
這裡告訴我們.
耶穌的成長跟施洗約翰有相似.
但也有很多不同的地方.
有什麼不同呢.
這裡是路加福音記載耶穌說的.

$^{441}$第一句話.
為什麼要找我呢.
你們不知道我必須要在我父的家裡.
這句話很特別.
就代表他知道彌賽亞的身份.
彌賽亞的使命.
他以神為他的天父.
這裡就是我父的家.
這些猶太人不會說的說話.
但從耶穌的口中出來.
他說他們不明白他所說的話.
耶穌就跟他們下去.
離開耶路撒冷回到那撒勒.
並且順從他們.
他母親瑪利亞把這一切事都存在心裡.
耶穌的智慧和新娘.
以及神和人對他的喜愛都不斷增長.
但我想問一個問題.
為什麼路加特別記載這個獨特的故事呢.
除了和施洗日的成長有對比之外.
因為路加記載耶穌的事蹟.
耶穌成就的救恩的時候.
他想將祂放置在耶路撒冷的聖殿裡.
所以你留意路加福音開始的地方.
是耶路撒冷的聖殿.
結尾也是耶路撒冷的聖殿.
這個情節特別的地方.
就是耶穌進入成人階段.
就留下了伏筆.
就是說這位耶穌比聖殿更大.
所以今天一個人敬拜神.
去親近神.
就不是去到地圖上的GPS一點.
耶路撒冷的聖殿.
而是在耶穌裡面.
所以你今天敬拜神.
主要不是在聖殿裡.
而是在基督裡面.
因為基督已經成就了聖殿所成就的.
基督的聖殿.

$^{481}$按照福音書的記錄.
耶穌之後花了大概十七年的時間.
默默地在拿撒勒過著平凡的生活.
拿撒勒這個地方不單是小.
也沒有人看過.
甚至是一個貶義詞.
一個一世紀的猶太歷史學家.
叫做Josephus約瑟夫.
他曾經在他的寫作裡提過.
加利利這個地帶.
有204個城鎮.
他居然一隻字都沒有提過.
有拿撒勒.
所以你想想.
在猶太人在世界眼中.
拿撒勒是一個多麼渺小的地方.
約翰福音第一章四十六節.
這裡是耶穌首次來召祂第一批門徒.
其中一個門徒叫做費力.
費力遇見耶穌之後.
他立刻找拿撒勒.
他告訴蘭大爺.
我們真的找到聖經裡預告會來的彌賽亞.
他就是拿撒勒人耶穌.
他說拿撒勒還能出什麼好的呢.
費力說不如你先來看.
你親自來經歷.
你就知道他有什麼特別.
你可以看到拿撒勒人耶穌.
是一個多麼貶義的詞語.
使耶穌上十字架的時候.
彼拉多寫的一個牌釘.
在十字架上寫什麼呢.
他說猶太人的王.
拿撒勒人耶穌.
他說你多醜.
居然有一個拿撒勒人.
一個這麼低微的人.
自稱是猶太人的王.
或者初代教會時期.

$^{521}$人們看不起教會.
看不起基督徒的時候.
就說他們是拿撒勒教黨的人.
都是貶義詞.
所以你看到那些第一代的基督徒.
他們是很寫實地將事情的真相導出.
如果他真的想去虛構一些故事.
去哄人信的時候.
他就不會介紹耶穌是拿撒勒人.
說回拿撒勒這個地方.
本身的泥土都不是肥沃.
所以農產一定是有限.
可能是這個原因.
迫使約瑟他主要的職業是木匠.
馬太夫的這裡怎麼說呢.
那裡的人質疑耶穌的身份的時候.
他們就問了幾條問題.
他們就說他不是木匠的兒子嗎.
他母親不是瑪利亞.
他弟弟不是雅各 約翰 西門 猶大嗎.
即是說我們認識他們全家.
他只不過是一個很平凡的木匠的兒子.
中文英文的翻譯通常就譯作木匠.
carpenter.
但原文的意義其實是廣很多.
原文是tecton.
即是說不單單是木匠.
而是石匠 包括工匠 建築工人.
有些像今天的handyman 或者建築師傅.
我們可以進一步問.
約瑟和年輕的耶穌當時的工作情況是怎樣的呢.
如果要回答這個問題.
我們可以宏觀一點看.
看看拿撒拉的地理.
我們可能會有多一點答案.
拿撒拉雖然有一個很小的村莊.
但它附近是一個大城市.
事實上是加利利省裡面最大的城市.
是加利利省的省府.
叫做Sephoris.

$^{561}$Sephoris 或者中文可以譯作齊波利.
齊波利是位於東南西北的交匯點.
往東北走.
就會經加利海 加伯隆等地.
然後一直往敘利亞的大馬士革.
然後就通往中東的米索不達米亞世界.
由齊波利向西南走.
就會經地中海.
然後南下經沿海平原.
直接到埃及.
所以它是一個匯點.
在第一世紀希律安提柏.
他在加利利作王的時候.
就定都於齊波利.
這裡不單止是交通的匯點.
你也看到附近一個很肥沃的山谷平原.
在猶太人的文獻裡面是很有名的.
叫做Batotva Valley.
它是位於夏加利利的山渠.
Lower Galilee.
夾住東邊加利利海的提比利亞.
西邊就是地中海的海法市中間.
一大片的平原.
是加利利山區最好的.
最肥沃的谷地.
所以從拿撒拿.
你只需要跨過一個山.
往北走大概90分鐘.
你就會到達齊波利.
約瑟瑪利亞一家.
從埃及回到拿撒拿的時候.
這個時候就是希律安提柏.
定都於齊波利.
然後大興土木.
他急需大量的工匠,木匠等等.
自然就是需要附近小鎮的工匠.
你這樣就解釋到為什麼.
約瑟的本鄉本身是南面的猶太地.
伯利恆人.
為什麼他一直在拿撒拿生活.

$^{601}$就是為了謀生.
齊波利的考古發現非常豐富.
你進入古城的時候.
首先你會發現一條大街.
很直的.
兩邊有大大小小的商鋪.
是一個商業中心.
山頂上有一個後期十字軍的網樓.
最主要吸引我們的是.
在公元後二至四世紀的時候.
齊波利變成一個非常富庶的城市.
有一個大型的猶太人會堂.
也有一些大戶的豪宅.
是發掘了出來的.
進去之後.
你會發現很多古老的建築物.
都是發掘了出來的.
進去這個豪宅飯廳.
上面鋪了很精細的馬賽克地板.
整個聖地都找不到這麼精細的馬賽克.
甚至土耳其和意大利.
都找不到這麼精細的馬賽克.
這裡最有名的就是上面的美女.
從任何角度看都好像看著你.
所以有人說這是加里里的蒙羅麗莎.
Mona Lisa of Galilee.
可見齊波利是一個很富庶的地方.
所以當時的人說.
它是Ornament of the Galilee.
這裡在山丘上有一個羅馬劇場.
在一世紀的時候開始動工興建.
你看到留下來的遺跡好像很小.
但其實它的直徑有45米.
按照這些遺跡來推斷.
這個羅馬劇場可以坐四千五百人.
很自然會想.
到底耶穌是不是當時其中一個的工人.
因為這些羅馬希臘劇場的演員.
每逢表演的時候都會戴上面具.
戴上面具的演員叫做Hypocrite.

$^{641}$聽起來這個字好像英文的hypocrite.
馬太方第六章第二節.
這裡耶穌去諷刺和指責法利賽人.
因此你們施捨的時候.
不可到處張揚.
好像偽君子.
這是新日本的翻譯.
和合本就是假冒為善.
在會堂和街上所作的一樣.
以博取眾人的稱讚.
我實在告訴你們.
他們已經得了他們的賞賜.
天父不會再賞賜他們.
所以智慧行間.
你會看到耶穌絕對不是一個.
什麼都不懂的人.
雖然他在一個小村落裡長大.
但他很緊貼當時人的文化和生活.
他說這些法利賽人是hypocrite.
是假冒為善.
他們就像那些羅馬劇場的演員.
有兩塊臉.
他們所取悅的是觀眾的讚賞.
不是天上唯一一位神的讚賞.
我們再看看這裡的地景.
約瑟和年輕的耶穌.
很可能每天來往那沙那.
和齊玻璃上班下班.
自從耶穌十二歲上耶路撒冷之後.
福音書的情節都沒有再提約瑟這個人.
暗示約瑟已經一早離開了世界.
作為長子.
耶穌自然就要承擔起當時的家計.
他要繼續養家.
直到其他的弟弟妹妹長大成人.
他才離開家裡.
然後公開傳道.
這可以解釋為何耶穌到三十歲左右.
他才開始公開事奉的職事.
中間有十七多年的時間.

$^{681}$我們叫他做Silent Years of Jesus.
沉默的歲月.
就是沒有驚天動地的事.
只是默默地付出.
默默地去生活.
這就是道成肉身的意思.
我們一起結束祈禱.
親愛的主.
因為你是道成肉身神的意志.
你是百分之一百經歷我們人所體會的.
主我們有的掙扎.
我們有的思念試探.
你自己都體會過.
主我們求你繼續幫助我們.
讓我們在我們生活當中.
我們需要默默地耕耘.
默默地事奉.
默默地作供.
求你給我們耐性和忍耐.
也讓我們知道每一天的生活.
縱使是黑板.
縱使不會有驚天動地的事發生.
這些都是你自己不斷推動的國度.
藉著我們的生命.
彰顯你自己天國的榮美.
主在平凡的當中.
用我們忠心的力量跟隨你.
祈禱奉主命救.
字幕志願者 李宗盛.
出品:桑海紀工作室.
本集完結.
\newpage



\section{}
\label{sec:KnqUXPbp8Kk}
\textbf{聖地網上行 - 步主腳蹤 [第八集] 登山變像}
\newline
\newline
連結: \href{https://youtube.com/watch?v=KnqUXPbp8Kk}{\texttt{https://youtube.com/watch?v=KnqUXPbp8Kk}} ~~~~ 語音日期: 2020-08-26
\newline
\newline
\hyperref[sec:QHoIKk7ZHEQ]{\small{< < < PREV SERMON < < <}}
~
\hyperref[sec:index_chronic]{\small{[返順時目]}}
~
\hyperref[sec:index_scriptual]{\small{[返順卷目]}}
~
\hyperref[sec:vLHlHGB1MwM]{\small{> > > NEXT SERMON > > >}}
\newline
\newline
$^{1}$歡迎各位來到聖地網上行第八集.
上一次耶穌和門徒去到凱撒利亞菲勒比.
在一個外邦人約旦河東的地方.
今天我們叫做哥蘭高地.
在這裡耶穌問門徒一個很重要的問題就是人子是誰.
這個問題是很重要的.
因為這是科幻書下半節要處理的問題.
就是耶穌真正的身份.
因為如果搞錯了耶穌真正的身份.
就會影響門徒對自己的身份的瞭解.
所以耶穌問他的門徒到底他自己是誰.
彼得回答說你是基督是永生上帝的兒子.
門徒跟著耶穌有幾年.
見過耶穌的神蹟騎士.
見過他趕鬼.
行很多的義能.
又應許建立他的教會.
陰間的權勢不能夠勝過.
對於門徒來說這個王者的國度已經來臨了.
很快就會到.
但是突然間耶穌轉口.
人子沒錯是基督是王者.
但是他要上耶路撒冷來受苦.
並且要被殺.
十六章第二十一節馬太福音說.
從此耶穌才指示門徒.
他必須上耶路撒冷去.
受祭司長文士許多的苦.
並且被殺第三日復活.
彼得聽到這個消息之後.
他不能夠接受.
我相信其他的門徒都不能夠接受.
彼得就拉著耶穌勸他.
萬萬不可.
因為這事必不臨到你身上.
耶穌轉過來對彼得說.
撒旦推我後面去吧.
這一句話撒旦推我後面去吧.
就是在曠野裡面.
耶穌三次受撒旦試探.

$^{41}$最後一次在馬太福音.
在高山上魔鬼將萬國展現在耶穌面前.
祂說只要你拜我.
不要上十字架.
我就會賜萬國給你.
耶穌說撒旦推我去吧.
居然這句話.
是耶穌跟彼得說的.
你想想有多大的震盪.
門徒很不能夠接受.
這位王者居然會受苦受死.
也很懷疑.
很震盪他們的信念.
到底這個主耶穌.
是不是上帝的王者.
天國的王者.
耶穌就說.
你推我去吧.
因為你不是體貼神的意思.
並且他之後就告訴門徒.
你們都是和我一樣.
你跟隨我的話.
你都要背上十字架.
於是到十六章最後結束的時候.
耶穌就說.
我實在告訴你們.
站在這裡的人在未嘗死未知前.
必看見人子降臨在他的國裡.
接著下來第十七章.
就是一個很出名的故事.
就叫做登山變將.
所以這個登山變將.
一方面來回應門徒的懷疑.
也繼續回答一個問題.
就是耶穌是誰.
他說初嘗天國的榮耀.
就好像你看一本小說.
或者看一部長篇劇.
中間的情節是很幽暗的.
你就忍不住.

$^{81}$去看一本小說.
看完結局.
原來主角是不會死的.
並且會有好的收場.
然後你又回到中間.
繼續看下去.
你就有力去完成那個小說.
完成那個情節.
其實這個登山變將.
也有一點像這樣的.
可能我們人生裡面.
經歷不同的苦痛.
或者經歷一些不好的際遇.
或者很多的不幸.
或者像門徒一樣.
他們在經歷一些.
震盪他們的事情.
這個登山變將.
就是將榮耀的結局.
給他們一個觀察.
讓他們看一看.
以致門徒能夠繼續跟著.
繼續下去.
登山變將.
登山變將這個故事.
第一節就講到.
在高山上發生.
當然自古以來.
有不同的建議到底這座山在那一處.
有一個建議.
就是在他柏山.
他柏山是一個.
很孤立的山峰.
在夏加利尼的一個平原的地方.
在一個很特別的山.
景色是非常漂亮的.
575米高.
1900呎高.
自四世紀開始.
有些古教堂.

$^{121}$建築在這個地方.
紀念耶穌登山變將.
但是一些文獻.
告訴我們一世紀.
在這處.
有羅馬軍隊在山上駐兵.
所以這個可能.
未必是那麼大.
另一個建議.
就是在麥倫山.
上到來.
上加利尼.
這些大山區的地方.
這裡有一個山峰.
叫做麥倫山.
1200米高.
4000呎高.
這是另一個建議.
當然第三個建議.
就是在.
Seserea Philippi.
在凱撒尼亞菲勒比.
上面的高山.
就是黑門山.
2800米高.
9200呎高的山.
有什麼地理線索.
可以告訴我們.
仍然是在凱撒尼亞.
菲勒比附近的高山.
發生的呢.
因為在故事的情節上.
其他方案書都是一樣.
沒有什麼位置上的移動.
在凱撒尼亞菲勒比附近.
就是有高山.
並且是以色列地.
裡面最大最高的山峰.
在耶穌的變相.
他將祂榮耀彰顯的時候.

$^{161}$就有摩西和以利亞出現.
猶太的文學認為.
這兩個人很特別.
他們兩個都沒有死.
以利亞當然很特別.
因為他們離世的情節.
記載在.
列王記下第二章.
這裡說到.
他們在約旦河東.
有一些戰車.
火馬車.
接以利亞升天.
大家記得這個地方.
靠近上一次我們提及過.
就是耶穌受洗的地方.
以利亞當然沒有死.
直接被提上天.
所以末後的時候.
他又會再降臨.
那摩西又如何呢.
大家記得摩西在.
申命記最後一章.
明明是說他死了.
並且埋葬在尼泊爾山.
尼泊爾山在約旦河東的平原.
其中一個山峰.
但是.
大家知道摩西是一些.
很重要的人物.
所以兩約中間的文學.
不斷去猜想.
摩西是怎樣的.
摩西其實沒有死.
有些文學說摩西其實又是像以利亞一樣.
被接上天.
所以當時猶太人的心裡面.
這兩個人不單是.
代表律法和先知.
代表整個的教育.

$^{201}$這兩個人也是很奇特的人.
也是末世裡面很重要的角色.
並且你留意.
地理上這兩個人.
都是在聖地以外在約旦河東的人.
摩西從來沒有踏足過應許之地.
他只是在約旦河東的尼泊爾山.
來遠望應許之地.
然後就離開世界.
以利亞當然是.
北國以色列.
阿哈耶西別時候做先知的.
但是以利亞本身是一個什麼人呢.
他是約旦河東的人.
他的鄉下.
就是基列.
聶王記上十七章第一節.
就是說基列寄居的提斯比人.
以利亞.
所以以利亞的鄉下.
就是基列地.
Gilet.
Gilet在約旦河東的山區.
或者基列這個地方.
是泛指約旦河東.
聖地以外東邊.
這個地區.
Gilet這個字.
大家是很熟悉的.
因為最近新冠病毒疫情的緣故.
我們聽到美國有一間公司.
叫做Gilet Sciences.
Gilet的科學公司.
這個是製藥的公司.
他們有一種特效藥.
叫做Rendezivir.
你問為什麼.
這間藥廠.
這間藥公司會叫做Gilet.
用了聖經這個字.

$^{241}$Gilet.
因為聖經裡面曾經提過.
Gilet這個地方.
因為約旦河東比較高的地方.
林木比較茂盛.
雨水比較多.
連聖經的時代.
Gilet都盛產了一種乳香.
一種膏藥.
Balm of Gilet.
耶利米斯第八章曾經提過.
他說Gilet那裡是有乳香的.
有些能夠治病.
治傷口的膏藥.
為什麼你不去那裡.
才罵那些意識的人.
為什麼你去拜別神.
去其他外國的地方.
去找治你的藥物.
無論如何.
摩西,以利亞是約旦河東的人.
他們是聖地以外的人.
和耶穌在這個故事裡面一起出現.
這個登山變象.
耶穌將它榮耀顯出的地方.
也是在約旦河東的地方.
碰到柱水角.
都像一道糖蒸蛋一樣.
一門張紗布.
左手按住勾子.
左手往前拉.
左右拉一下.
雙手合二.
看完這些山景色都很美.
但我們想進一步問一個問題.
如果馬爾羅曼山是黑門山.
約旦河東的高山又如何呢.
可能重點不是山在哪裡.
而是山上發生什麼事才是重點.
古人的概念是高山是天與地的聯繫.

$^{281}$天與地中間.
就是神在這裡降臨.
神在這裡對人說話.
尤其是在馬太福音.
以前都曾經提過.
馬太福音提過七座山.
比如八福山.
八福山就是耶穌上山的時候.
祂坐下.
然後門徒來到祂的根前.
祂就開口教訓他們.
就像耶穌摩西上山一樣.
將律法來到此下.
然後馬太福音最後.
二十八章第十六節.
所謂大聖名山.
就在加尼尼的山頭上.
耶穌約定他們.
然後將大聖名賜給他們.
所以馬太福音第十七章第一節.
也提到登山變場在高山上.
情節是什麼呢.
十七章第一節.
過了六日.
耶穌帶著彼得,雅各和雅各的兄弟約翰.
暗暗地上了高山.
就在他們面前變了形象.
臉面明亮如白日頭.
衣裳潔白如光.
忽然有摩西,以利亞向他們顯現.
與耶穌說話.
彼得對耶穌說:主啊.
我們在這裡爭好.
禰若願意.
我就在這裡搭三座棚.
一座為禰.
一座為摩西.
一座為以利亞.
說話之間.
忽然有一朵光明的雲彩.

$^{321}$遮蓋他們.
且有聲音從雲彩里出來說.
這是我的愛子.
我所喜悅的.
你們要聽他.
門徒聽見就俯伏在地.
極其害怕.
耶穌進前來摸他們說.
起來不要害怕.
他們舉目不見一人.
只見耶穌在那裡.
下山的時候.
耶穌吩咐他們說.
人子還沒有從死裡復活.
你們不要將所看見的告訴人.
這個記載發現摩西代表律法說.
以利亞代表末世的先知.
即是復興摩西律法的先知.
所以兩個人代表整個舊約.
來到耶穌面前.
即是耶穌是整本舊約所指向的那位承傳者.
摩西以利亞兩個人在山上.
和耶穌在榮耀裡面顯現.
提出一個很重要的重點.
就是這兩個人另外的特徵.
就是兩個人都曾經在山上.
在舊約和榮耀的神相遇.
並且是同一座山.
這座山叫西乃山.
或者叫赫烈山.
Mount Sinai或者Mount Horus.
這座西乃山在現在的西乃半島南面.
叫Jabut Musa.
即是Mount of Moses.
記得舊約裡面的情節.
就是說以色列人出埃及之後.
他們經過西乃的曠野.
然後來到西乃山或者赫烈山.
神就在山上和他們顯現.
摩西也上山在榮耀當中.

$^{361}$和神說話對話.
出埃及記二十四章十二節.
耶和華對摩西說.
你上到我這裡來.
上山來到這裡.
我要將石板.
並且我所寫的律法和誡明詞給你.
你就可以教訓百姓.
雲彩遮蓋山六日.
第七日神就在雲中.
去召摩西和他說話.
耶和華的榮耀停在山頂上.
在以色列人的眼前.
形狀如烈火.
在這裡發生什麼事呢.
就像耶穌登山變像的故事.
你看到有山 高山.
過了六日.
這裡也有說到雲彩停在山上.
遮蓋山六日.
第七日就說話.
就像馬太科一樣的十七章情節.
在這裡很長的幾章情節裡面.
說到摩西在山上發生什麼事.
他不僅見到榮耀的神.
不僅見到上帝的榮光.
並且他自己的面也發光.
這些情節.
就說到摩西變像.
到耶穌變像的時候.
其實都是一樣.
其實說到耶穌登山變像.
就好像帶著以色列民.
門徒來到西乃山再經歷一次.
說完摩西就說到以利亞.
以利亞有一段經歷.
也很像摩西.
就是在聶王紀上帝十九章.
那裡說到以利亞逃亡.
因為阿哈和耶細別要追殺他.

$^{401}$他從以色列走到南面的西乃山.
他起來到神的山.
上帝的山就是赫烈山.
他在那裡進了一個洞.
就住在洞中.
然後阿基華就想到他.
他說以利亞你在做什麼.
以利亞就回答他.
我為耶和華大發熱心.
以色列人背棄你的約.
毀壞你的祭壇.
用刀殺你的先知.
現在只剩下我一個.
我的命也可能沒有了.
耶和華就對他說.
你出來站在山上.
在我面前.
那時耶和華從那裡經過.
在他面前有裂封大作.
崩山碎石.
耶和華卻不在風中.
風後有地震.
耶和華卻不在其中.
地震之後有火.
耶和華也不在火中.
火後就有微小的聲音.
以利亞聽到這個聲音.
就用愛依蒙傳念.
出來動口.
有聲音對他說.
以利亞你在做什麼.
他的回答都是一樣.
最後耶和華對他說.
你回去.
從曠野的路去大馬騎士.
去到那裡.
你要高納不同的人.
來審判以色列人.
這些在舊約里叫做.
Theophanies.

$^{441}$即是上帝顯現的異象.
人們見到這些榮耀的時候.
他們初嘗將來天國的那種榮美.
但一個很重要的.
不是看見什麼.
當然有描述很多看見的情況.
有榮光 有裂封 有火等等.
有白光等等.
但似乎重點在最後是說話.
留意摩西.
見到榮耀之後.
經常出現一句話.
就是耶和華曉遇摩西.
去吩咐摩西.
告訴他你要去見會莫.
按照我的說話去見會莫.
你要行律法.
要教導你的子民.
要遵守約裡面的規條.
最重要的當然是上帝頒布十誡.
他們來到西乃山的時候.
上帝最重要頒布的.
就是十誡.
十誡本身的原文不是誡命.
其實是十句說話.
以利亞的經歷也是一樣.
以利亞的異象.
我們經常想起.
見到狂風大作.
一些很奇特的現象.
但最重要的重點是什麼.
就是火候有微小的聲音.
那聲音就是耶和華和以利亞說.
你回去.
做我要你做的事情.
所以這些所有的.
Theophany就是上帝榮耀顯然的異象.
重點不是看到什麼.
而是你要聽.
聽什麼才是重要.

$^{481}$你看看耶穌登山變象的經文.
重點不是視覺.
而是聽說.
第三節.
有摩西有以利亞出現.
和耶穌說話.
彼得在這情景裡面.
他也不知道自己在說什麼.
總之他說見三座棚.
他也未瞭解.
原來基督和摩西和以利亞.
其實是不同的人物.
比摩西大.
比以利亞大.
彼得就當作三個是一樣的.
彼得在榮耀裡面.
他很興奮見到這些榮耀.
他就想將這些榮耀留在山上.
但最重要的是說話.
第五節.
彼得說話之間.
忽然有一朵光明的雲彩.
去遮蓋他們.
並且有聲音.
有聲音在雲彩裡面說.
他說:這是我的愛子.
我所喜悅的.
你們要聽他.
這裡記得.
在《馬太福音書》第一次提到.
我所喜悅的.
我的愛子.
就是耶穌受洗的時候.
就是天上有聲音.
天父開始耶穌在地上.
他愛子在地上的工作.
有神跡.
有歧視.
有治病.
有趕鬼等等.

$^{521}$這裡第二次.
天上有聲音.
就開展耶穌第二段最重要的工作.
就是耶穌上耶路撒冷.
受死.
受羞辱.
被釘十字架的旅程.
第八節.
不見一人.
只見耶穌在那裡.
我們很多時候覺得.
就是視覺.
只見耶穌在那裡.
其實這樣結束榮耀的意象.
不見一人.
只見耶穌.
其實就是說.
你唯一要聽的.
就是耶穌所說的話.
所以下山的時候.
第九節.
耶穌吩咐他們說.
就是又說話.
我們讀到這些天上的意象.
天上的榮耀.
我們覺得很奇妙.
很吸引我們.
記得很多年前.
我幫一個老伯做一個安息禮拜.
因為他是外州來的.
我一直不認識他.
他的兒子到了.
我便和他談話.
我問老家父他何時信主.
他說很多年前.
他的父親患重病.
垂死的時候.
他突然有一個經歷.
他突然飛上天.
另一個境界.

$^{561}$他見到很多大光.
很多榮光.
聽到天上很動聽的聲音.
然後他突然下來.
病好了.
他便信了耶穌.
自此之後他沒有離開主.
直到他離開世界.
這種死過活回.
這種特別的生命經歷.
不是每個人都有.
但我又想起.
同時也有些人.
又是經過大能不死.
但他依然固我.
很湊巧地又死不去.
但他依然是沒有信耶穌.
分別在哪裡.
不是見到的東西.
就是他肯領受神的說話.
肯領受神的奉迎.
所以其實神一直和我們說話.
整本《馬拉夫人》都強調.
悶透是不可或缺的特徵.
不是見到耶穌.
或者藉著耶穌做很多神跡奇事.
不是經歷到很奇妙的事.
而是聽上帝的說話並且進行.
這個登山變相的故事.
讓我們看到將來最終.
基督的榮耀,天國的榮耀.
但其實焦點不是耶穌的威榮.
而是要聽清楚,想清楚.
到底基督的身份是什麼.
就是基督必須要受苦.
因為沒有受苦就不會有救贖.
上帝不會救贖我們這群罪人.
沒有救贖也不會有榮耀.
所以主題不是德聖的威榮.
而是異獄耶穌.

$^{601}$就是神所喜悅的,揀選的.
以塞瓦殊的受苦的獨人.
我們跳一跳去路加福音看九章.
同樣記載登山變相這個故事.
九章第三十節.
忽然有摩西爾利亞兩個人與耶穌說話.
他們在榮光里顯現.
談論耶穌去世的事.
就是祂在耶路撒冷要成就的事.
去世的事英文翻譯是departure.
其實原文路加特別用一個字.
是exodus.
出亞埃及.
就是exodus.
特別加上後面一句短的說話.
在耶路撒冷將要成的事.
耶穌的死.
就有如出埃及.
你記得第一次出埃及.
就是神救祂的子民.
從埃及捆綁中得到釋放.
從埃及一直經過曠野,山乃山等等.
然後進入了應許之地.
去到耶路撒冷.
這個就是完結.
第一次的出埃及.
耶穌的出埃及.
新的出埃及.
new exodus.
就是拯救祂的子民捆綁.
真正的釋放.
就是從耶路撒冷開始.
就是耶路撒冷十字架.
直到天上的榮耀.
所以是vertical的.
耶穌的死和復活.
才是真正的出埃及.
真正讓我們的生命得到釋放.
馬太福音就將這個故事放下了伏筆.
將登山變象和十字架的情景放在一起.

$^{641}$用同樣的字眼來描述.
如果我們跳到馬太福音27章.
那裡講到十字架上發生了什麼事.
有幾個共通點.
第一.
在登山變象看到榮耀.
他們害怕.
在十字架的面前.
他們看到耶穌死在十字架上.
他們也很害怕.
在登山變象.
天上的聲音就是.
這是我愛子我所減損的.
你要聽他.
27章54字.
也有聲音.
這是地上的聲音.
就是伯夫長和一切看守耶路撒冷的人.
看到地震所經歷的事.
他們就極其害怕.
他們說這是上帝的意志.
這是神的意志.
也提到以利亞.
只不過以利亞沒有來.
另外一個共通點.
在十字架面前.
也有三個門徒.
三個是他的門徒.
所以你看到十字架下面的情景.
其實是跟登山變象那種榮耀.
威榮的情景是一個平衡.
十字架上的耶穌才是榮耀的主.
遲些我們會去到耶路撒冷.
就是耶穌被釘的地方.
耶路撒冷有一個古教堂.
歷史歷代蓋在上面.
就是紀念耶穌被釘和埋葬復活的地方.
那裡叫聖母堂.
Church of Holy Sepulchre.
Sepulchre是拉丁文.

$^{681}$即是母月的意思.
我記得有一次.
在耶路撒冷.
我跟一個東方教會歷史學者的專家.
實地講解當中發生了什麼事.
他在教堂外面指著教堂上的十字架.
他說你有沒有想過.
十字架的Eastern Orthodox的設計.
就是金色代表榮耀.
但設計十字架的人中間放了黑色.
一些黑色的寶石.
設計十字架的人明白一個道理.
耶穌的榮耀是要經過苦難.
要經過十字架的痛苦,受苦.
才會彰顯出來.
不是整個都是金光閃閃的.
我們很多時候想起榮耀.
或者想起神的榮耀彰顯.
神的工作來到了.
人人很多時候都想起.
你是多威風.
我們教堂多大.
多麼華麗.
最近有一個新聞大家可能聽過.
在土耳其伊斯坦堡.
有一個很大的古教堂叫聖索菲爾.
是一個大教堂.
是一個爭議.
是一座在拜佔庭時期已經有的大教堂.
有千年歷史.
當時拜佔庭時期的王傑斯丁去興建.
直到現在是一個曾經在世界上最大最華麗的教堂.
享福成名.
到今天去土耳其遊覽的人.
都會說這個偉大的建築.
當時古時居然可以做到這麼大的教堂.
贊不住口.
最近的爭議是什麼呢?.
就是土耳其的總統.
將這個地方想正式變成清真寺.

$^{721}$我們想起神的榮耀.
很多時候就想起這件事.
建築多麼的輝煌.
沒有了就很損失.
但你知道教堂不是教會.
教堂只是建築物.
上帝所建立的教會是一群信徒.
是一群願意去行上帝旨意的一群人.
正如耶穌我們的主一樣.
耶穌的榮耀是如何表現出來的呢?.
不是他金剛不壞之身.
不是打不死.
而是在他的苦難裡面.
在他受苦.
在十字架上.
羞辱的裡面.
赤裸裸的在十字架上.
成為眾人的救贖.
這才是上帝榮耀的彰顯.
如何顯出上帝的榮耀?.
你問是否要多花錢.
我們就顯出上帝的榮耀.
我們有很多地.
有頭有面的.
社會裡面是成功地位的人.
我們就有上帝的榮耀.
教會很多時候就容易掉入這種危機裡面.
真的.
你有頭有面.
你是成功人士.
你就很可能很多時候.
在教會裡面聲大點.
甚至做領袖.
這是一個很容易掉入的危機.
因為我們用人的榮耀.
去衡量神的榮耀.
或者我們將它調轉.
真正認識耶穌.
知道耶穌的榮耀.
顯出就是祂默默地去遵行神的話.

$^{761}$甚至死在十字架上.
我們有多像耶穌?.
我們是否明白耶穌的身份?.
我們是否明白我們作為基督徒.
作為一個追求神的人.
我們真正追求的目標是什麼?.
上帝最近對你說了什麼話?.
我們是否有聽?.
讓我們聽了上帝的話.
並且進行.
我們低頭祈禱.
因為禰是那位說話的主.
禰對我們每時每刻.
即使是聖靈.
禰都會提醒.
禰都會說話.
主一求禰讓我們有個好的聽覺.
聽見神對我們的感動.
禰叫我們所作的.
我們願意遵從.
即使在人間.
不是什麼大事.
不是什麼很大的榮耀.
但讓我們默默地遵守禰的說話.
好像我們的主耶穌基督.
以福基督一樣.
我們要有一個堅定的信念.
\newpage



\section{}
\label{sec:vLHlHGB1MwM}
\textbf{聖地網上行 - 步主腳蹤 [第七集] 凱撒利亞腓立比}
\newline
\newline
連結: \href{https://youtube.com/watch?v=vLHlHGB1MwM}{\texttt{https://youtube.com/watch?v=vLHlHGB1MwM}} ~~~~ 語音日期: 2020-08-27
\newline
\newline
\hyperref[sec:KnqUXPbp8Kk]{\small{< < < PREV SERMON < < <}}
~
\hyperref[sec:index_chronic]{\small{[返順時目]}}
~
\hyperref[sec:index_scriptual]{\small{[返順卷目]}}
~
\hyperref[sec:jXnsElF6xnM]{\small{> > > NEXT SERMON > > >}}
\newline
\newline
$^{1}$我們來到第七集聖地網上行報主腳鐘.
今次我們來到耶穌事蹟裡面一個很重要的轉捩點.
上一次我們講到馬太福音第十四章.
分封王希律聽見耶穌的名聲開始對他有威脅.
於是福音書的記載就說他五餅而餘.
餵飽五千人之後在加利利湖的風暴上.
他就在水上行走.
然後他去到加利利湖的西北岸旁邊.
一個叫格尼撒勒的地方來治病.
到馬太福音第十五章.
耶路撒冷的宗教權貴政治權貴派人反對耶穌.
你會感覺到對耶穌反對的勢力越來越強大.
所以到十五章第21節.
這裡講到耶穌去到一個地方叫做泰爾西頓的境內.
你記得那個故事就是一個加拿大外邦的婦人.
對耶穌是很有信心的.
那裡有一個神蹟.
完了這個加拿大婦女的故事.
耶穌又回到加利利湖的海邊.
但她繞過猶太人的境地.
當然那裡的反對勢力因為越來越強.
所以耶穌大部分時間都是用外邦人這個地方來做事.
馬太福音第七章第31節說.
耶穌又離開了泰爾的境地經過西頓.
然後他就經過我們今天叫他哥蘭高地的地方.
就經過低加波利亞.
又是外邦人的地方.
就是加利利海.
應該是加利利海的東邊或東南邊.
在這裡有另一個餵飽四千人的神蹟.
這也是外邦人的地方.
十六章耶穌就渡船過去猶太人的地方.
去了這個莫大拉.
然後就叫門徒和他一起渡去另外一個地方.
過去對面.
應該就是白菜大或外邦人的地方.
你可以看到當這些反對勢力繼續威脅耶穌和他的門徒的時候.
耶穌大部分時間在這兩章裡面.
他都是留在外邦人的地方.
或者是加利利湖的東邊.

$^{41}$加利利湖很近的西北邊.
這些全部都是外邦人羅馬希臘人的地方.
馬太福音第十六章十三節.
在讀新漢語譯本.
耶穌來到凱撒利亞菲勒比境內.
他問門徒說人們說人子是誰.
他們說有人說是施洗的約翰.
有人說是以利亞.
也有人說是耶利米和先知中的一位.
耶穌對他們說那你們呢.
你們說我是誰.
西門彼得回答說.
你是基督是永活神的兒子.
經過幾集我們希望大家明白.
聖經經文記載一些地名.
一些地理的細節.
是有它特定的意思的.
也告訴我們要明白那段事跡.
我們一定要先了解背景.
那個地方才能夠掌握裡面的內容.
這裡耶穌和門徒來到一個地方.
叫做凱撒利亞菲勒比.
或者是和合本的翻譯叫做.
蓋撒利亞菲勒比.
Caesarea Philippi.
這裡是加利海北面.
25英里或者40公里那麼遠的.
在黑門山西南的山腳.
這個黑門山高達2700米.
或者9000呎那麼高.
是以色列聖地裡面最高的山.
它處於北邊.
也有高山.
所以雨量是很充足的.
特別是冬天到四五月之間.
經常有積雪.
其實是以加利式來說是一個好地方.
天朗氣清的時候.
你在加利利湖.
在南面的地方.

$^{81}$你遠遠就可以看到一個積了白雪的山峰.
有時候回到北峰的時候.
就將黑門山的寒氣,水氣帶到地上.
滋潤土地.
這個就是聖經裡面說的.
黑門山的禁路.
黑門的禁路.
我們來到凱撒利亞菲勒比.
這個地方的景色是非常漂亮的.
而且是一個天然的大水泉.
其實是約旦河的源頭之一.
在這個地方的盤石底下.
你可以看到一個很大的盤石.
一個很大的山.
在山腳當中.
就有一個很大的山洞.
這裡就是水泉源頭的地方.
以往的水位比現在更高.
這個山洞的水泉是深不見底的.
這裡很可能是舊約聖經裡面經常提到的.
巴力加德.
Belgard.
或者是巴力黑門.
Belhamon的地方.
就是敬拜雨神巴力.
Bel的地方.
這個就是巴力的神像.
在舊約的以色列裡面.
他的右手是拿著雷電的.
因為他是暴風雨神.
所以他是像雨水的.
也是像繁殖.
給人有生命的別神.
在以色列在舊約裡面是非常繁茫的.
然後到了希臘的時期.
大概公元前300年左右.
這個地方就改了名叫做.
潘尼亞斯.
Parnias.
因為這個地方就轉了.

$^{121}$變成了希臘的一個山神.
叫做潘.
所以這個地方就變成了Parnias.
今天的阿拉伯話的地名就保存了.
Parnias就叫做Barnias.
就Barnias.
不是Parnias.
因為阿拉伯文是沒有P音的.
所以就變成了Barnias.
在希臘的神話裡面.
這個潘神潘.
他是半人半獸的.
是一個很恐怖的神靈.
事實上英文.
慌張.
恐慌.
其實是從潘神這個名字而衍生出來的.
這個潘神就是半人半羊的.
他在大自然裡面活躍的神靈.
所以拜祭他的地方.
在希臘的世界裡面.
通常都是大自然.
洞穴,山洞.
這些洞穴是通往地下的陰間的.
拜祭潘神有一個特別的儀式.
就是獻祭的人會帶著祭物.
然後將祭物扔進去.
這個深淵裡面.
這個地下陰間裡面.
如果這個祭物不浮上來.
就代表潘神接納了他的獻祭.
但如果附近的水源有血.
或者這個祭物浮上來的話.
就代表神靈不要他的祭物.
可能要再做過.
大希律王得勢之後.
羅馬帝國就將這塊地賜給了大希律.
大希律王就在這裡建設了一個.
用白混石建造的很宏偉的羅馬皇神殿.
獻給他的老闆蓋薩或者凱薩.

$^{161}$大希律死了之後.
大家記得了.
他的國土就分給了他三個兒子.
猶太省就給阿基努.
加利利就給安提柏.
東北面的哥蘭高地一帶.
就屬於第三個兒子.
菲力.
由始至終都是愛邦人.
希臘人.
羅馬人的陣地.
菲力就在這裡擴大了羅馬皇帝王神廟.
然後就加建了城鎮.
就將這個地方命名為凱薩里亞.
凱薩里亞就是給他的老闆.
來慶祝他的老闆.
菲力就是慶祝他自己.
因為他叫做菲力.
所以就叫做錫瑟利亞菲力牌.
所以你看到一個很大的盆石.
很大的岩洞.
而且有水從地下水湧出來.
這個岩洞在希臘的神話裡面.
就是通往陰間死亡的地方.
在這棵岩壁上面.
就雕了一個很大的圖案.
叫做Nichus.
就像神位一樣.
有很多是給潘神和其他的神靈.
有一個當然是給羅馬帝王來供奉他.
羅馬帝國統治佔領了這個地方的時候.
就將羅馬帝王廟建立在這裡.
在這個山洞的地方.
供奉的主要就是羅馬帝王神.
就將帝王廟蓋在希臘的神廟潘神廟上.
這個就是帝國的政權.
超越了以往希臘的神話和宗教.
在一世紀裡面.
拜別神.
拜羅馬帝王的宗教的中心.

$^{201}$有了這個地方的背景地理背景.
我們就了解一下這個故事.
裡面有些深層的意思.
耶穌首先問他的門徒.
他說人們說人子是誰.
就是民間裡面他們在說耶穌到底是誰.
門徒就回答他.
有些說你是屍體若寒.
雖然被斬.
但是會復活的.
有些就說他會沒後來的先知.
耶利米,以利亞等等.
因為當時的耶穌.
他們正在預期著.
這位王者利塞亞會來到.
這個王者在《但以利書》經常有說.
在但以利的意象裡面.
這個人子.
Son of man.
就會來征服地上的國度.
並且他是一塊石.
這塊石會戰勝征服地上.
去迫害神子.
去迫害神子民的勢力.
同時這個人子也是先知的綽號.
所以耶穌傳道的一段時間.
他不斷地說自己是人子.
因為這個稱號.
能夠將兩方面的意思混在一起.
無論如何.
耶穌傳道一段時間.
民間不斷地猜想到底他是誰.
耶穌問完門徒.
問民間如何看待自己的時候.
他說你們要如何看待我.
跟隨耶穌的門徒.
不斷地要去明白.
耶穌真正的身份.
其實這個不是門徒第一次去確認.
耶穌的身份.

$^{241}$你記得上一次在加里尼湖的風暴裡.
所有的門徒在船上都承認.
你真是神的兒子.
The son of God.
為什麼彼得去到應信.
有一個confession這麼重要呢.
因為這是第一次門徒.
全面地去了解耶穌的身份.
其實這些門徒是代表教會.
代表基督徒.
代表我們在跟隨主的路上.
我們要不斷經歷.
不斷對主有進一步的認識.
更深的體會.
很多時候這些歷練有成有敗.
但我們對主的認識.
要進一步去認知.
對主的認識在信心上的程度.
不是一時一刻間會發生的.
要經歷一些時間.
要經歷一些失敗.
早前我們對主的認識.
肯定和現在是不同的.
以前一定是有所不足.
有所缺乏的.
現在我們對主的認識.
應該是越來越有把握.
越來越充分.
夏文也可以告訴我們.
這裡告訴我們.
門徒當時正在掙扎.
到底耶穌是誰.
甚至到耶穌復活之後.
中間也有人繼續懷疑.
所以福音書我們去跟隨主的腳蹤.
跟隨著這個故事的時候.
其實是邀請我們不斷對主有更深的認知.
好像一個三歲的小朋友.
他對父母的認識.
和一個二十三歲.

$^{281}$四十三歲的成年人.
對父母的認識.
之間一定是有差別的.
差別不一定是對與錯.
而是程度 深度上.
你未信主之前.
你可能對主的認識是.
耶穌是一個歷史人物.
很偉大 很有智慧等等.
有沒有錯呢?未必一定是錯.
但一定是有限制.
但當你領受了主對你的啟示.
每一個人信主 認信主的時候.
都是得到天上的啟示.
我們對這位主對我們的關係認識就更加深了.
你信了主 久了之後.
你明白這個關係.
越來越知道誰是主 誰是僕.
誰是王 誰是民.
這是一個永久的關係.
所以為什麼凱撒利亞 菲勒比 彼得的認信.
是有再跨前一大步.
但有沒有需要繼續學習.
需要繼續去了解的地方呢?當然有.
十六章第十七節.
耶穌對彼得說.
約拿的兒子西門 你是有福的.
因為啟示你的不是凡人.
而是我在天上的父.
這裡用了一個稱呼.
約拿的兒子西門.
這是阿難民的稱呼.
很強調彼得本身的人性.
你得到耶穌是王者 基督的啟示.
不是屬血氣的人 凡人能夠給你.
而是天上特別的啟示.
這不是你自己聰明一點才能夠了解.
而是天上特別的啟示.
我們信主的人.
我們就是有一個特別的啟示.

$^{321}$是一個恩典 是一個福分.
十八節 我還告訴你 你是彼得.
我要在這塊盆石上建立我的教會.
死亡或陰間的權勢勝不過他.
這段聖經無論如何說都是很富爭議性.
最大的爭議就是天主教一直強調.
指的是彼得.
教會建立在彼得之上.
因為彼得是第一任的教皇.
而之後的教皇不斷有階層 不斷有延續.
這才是真正的教會.
怎樣解釋都有爭議性.
不過我們值得留意的是希臘文的原文.
希臘文的原文兩個字的音很相似.
但字是有所不同.
彼得這個字是petros.
是一個陽性盆石的意思.
基法石的意思 陽性.
然後在這塊盆石上.
這個盆石是陰性的petra.
作者很肯定是有刻意的work play.
即是同音字 似的音字.
但那個字是有不同.
一方面要表達彼得是重要的.
確實彼得是十二門徒之首.
是初代教會的領袖.
所以值得留意.
彼得的重要性是時間上的重要.
不是地位上的重要.
是chronologically.
但hierarchically不是地位上的高超.
其實你留意一下耶穌所說的話.
你都略略知道.
耶穌說什麼呢.
我要將我的教會.
又說回希臘原文.
就是字的次序.
我的 my church.
是放在前面的.
不是放在後面.

$^{361}$刻意放在前面.
代表教會是屬於他的.
my church.
不是peter的church.
另外後面盆石這個字.
音相似但字不同.
這個不會是指彼得自己.
不會是在彼得上建立耶穌的教會.
如果又是指彼得的話.
耶穌最自然應該會說.
你是彼得.
我要把我的教會建造在你身上.
用第二身.
不會是on this rock.
所以如果耶穌說將我的教會建立在彼得之上.
這樣說就很奇怪.
所以為什麼特別用第二個字.
就是音相似的字.
不同的字.
指到盆石.
正如我開頭所說.
經文一定要按照地理去了解.
為什麼故事的地方背景.
是在凱撒尼亞菲勒比.
剛才解釋凱撒尼亞菲勒比.
是在菲律賓的時候.
就是說這裡有一個大的盆石.
是以色列最高的山.
黑門山.
最大的盆石.
而當中在盆石上.
有羅馬帝王神廟.
有潘神廟.
都是建立在大的山.
大的盆石rock上.
是代表帝國的勢力.
代表死神的威嚇.
耶穌說我要將我的教會建在這盆石上.
其實指著的是黑門山的盆石.
在這盆石上.

$^{401}$這個在字是一個前置詞.
有在上的意思.
表達是高過.
是超越的意思.
英文可以翻譯成above.
或者over.
就是說耶穌說我的教會.
要超越這個死亡的權力.
或者也可以.
這個前置詞.
有英文against的意思.
就是opposition的意思.
就是說我要建立我的教會.
是對抗.
是顛覆.
去抗衡一個假的宗教.
假的勢力.
而且這個帝國的勢力.
以死作威嚇的帝國的勢力.
是不會勝過這個教會.
對於當時的教會來說.
他們很多時候面對惡勢力.
面對一些有組織的.
禁聖打壓等等.
這一段的技術就很有意思.
所以為什麼呢.
我們今時今日.
二十一世紀的教會.
我們都需要知道.
特別初代教會的歷史.
我講到church history教會的歷史.
人人都覺得很悶.
但是你要用這個角度去看.
基督教的教會歷史.
其實是一些活生生的見證.
為什麼去聖地呢.
為什麼來這個地方呢.
因為你隨便看一看這個地方.
所有教會的歷史就呈現出來.
之前有一集我們講到曠野.

$^{441}$我們看到曠野裡面有修道院.
有教會的存在.
你看那些修道院.
每一間都有幾百年.
千多年的歷史.
甚至早前有分天主教.
通靖教 基督教.
這些修道院.
這些古教堂.
見證著信徒的存在.
你看這些歷史的時候.
你就知道.
很靠近東方的教會.
特別是東正教.
就是去到見證著.
有多少權勢.
有多少幫國的改變.
改朝換代.
但是教會依然留在這裡.
今日的教會在不同的地方.
其實面對著很多的威脅.
在西方的教會.
面對著很多權利上的引誘.
在東方的教會.
面對著中共去拆十教.
去拆教堂.
用金錢 用名利去收買領袖.
去將他們禁聖.
國安法去用恐嚇來打壓.
去禁聖他們.
但是在這裡耶穌所講的.
就是一個他的保證.
死亡陰間的權勢.
是不能夠去勝過的.
相反教會是可以勝過.
這些帝國的勢力.
羅馬帝國.
多強硬的打壓也好.
現在羅馬帝國已經消亡了.
但是教會仍然存在.

$^{481}$我們今天看回威脅教會的權勢.
其實都是一樣.
有朝一日他們會覆亡.
但是教會.
所有認信永生神的兒子.
彌賽亞王者基督的.
他們依然會存在.
第十九節.
耶穌說我要給你天國的約事.
凡你在地上所禁止的.
或者和合本所捆綁的.
也將會是天上所禁止的.
所捆綁的.
凡是你在地上所許可的.
或者和合本釋放.
也將會是天上所許可的.
或者釋放.
這裡說到天國的約事.
意思就是開門關門的權柄.
什麼意思呢.
其實這裡釋放捆綁.
或者新漢語的官.
許可禁止.
這些字眼.
是猶太人拉比文獻裡面的法律用語.
釋放的意思就是說.
在他的教導裡面.
拉比的教導裡面.
就是許可赦免的意思.
捆綁,禁止.
就是定位有序的意思.
就是不准做的意思.
也是一種權柄的象徵.
表面上我們可能會以為.
教會在地上說什麼.
天上就要聽話.
這裡又是說到原文的語法.
其實是刻意的.
做得很奇怪的.
因為這裡第二句.

$^{521}$是用將來完成式.
Future Perfect Tense.
是刻意這樣用的.
如果是簡單的將來式的話.
即是Simple Future的話.
即是我們剛才以為的意思.
就是彼得教會說了算.
但將來完成式就不同了.
就是說彼得所做的決定.
必會或將會在天上已經.
Perfect Tense.
已經決定了.
所以我們按照將來完成式.
Future Perfect Tense.
可以這樣譯譯的.
凡是你在地上所禁止的.
所捆綁的.
也將會是天上已經禁止的.
或者已經捆綁的.
凡是你在地上所許可的.
所釋放的.
也將會是天上已經許可.
或者釋放的.
換句話說.
主權就在天上.
而不是在地上.
天是先行的.
換句話說.
這裡是說彼得.
代表教會.
是一個忠心的管家.
就是天上的父神.
去託管我們的律法.
我們地上的執行.
是神已經有的決定.
要簡單理解.
舉個例子.
你知道科林書二號的記載.
甚至到司堂傳的記載.
主要想指什麼呢.

$^{561}$就是教會去到傳道.
傳福音.
傳一個叫人悔改歸主的福音.
對願意悔改的人來說.
就是開門給他.
許可.
對一些不願意認罪的人來說.
當然就是關門.
禁止或者捆綁.
另外一個實例就是.
初期的教會.
將福音傳到外邦人中間的時候.
發生什麼事呢.
就是耶路撒冷.
猶太人的教會.
他們就開會了.
他們想了解這件事.
他們討論什麼呢.
就是外邦人信耶穌.
信福音.
要不要他們又要守猶太人的律法.
他們才能夠成為神的子民.
教會開會之後.
就理解到.
原來信耶穌悔改的福音就可以了.
這個是教會去作一個決定.
但這個早已是神天父的意思.
就是要普世的人去認識他.
去成為他的子民.
這裡講到天國的約事.
講到在地上許可的權柄.
開門的權柄.
就代表教會在屬性上.
是一個很大的權柄.
這裡提醒我們的教導.
是否跟上帝的心意一致.
20節.
之後耶穌就祝福門徒.
不可對任何人說他是基督.
21節.

$^{601}$從此耶穌才指示門徒.
他必須上耶路撒冷去.
受長老祭司長問是許多的苦.
並且被殺.
第三日復活.
這裡提醒門徒.
他們對耶穌的認識.
雖然是越來越進步.
但有許多地方他們仍然未了解.
他們需要經歷去了解.
這裡講得很清楚.
耶穌說他是基督.
但暫時不准門徒傳開.
因為猶太人的民間對基督.
王者的了解實在太過片面.
不清楚.
所以如果這個時候傳出去.
就會惹起許多民間的誤會.
甚至連這些跟了耶穌幾年的門徒.
他們都未明白基督的意思.
所以這裡這樣說.
從此.
從這個時候開始.
耶穌就要指示門徒.
他必須上耶路撒冷受苦.
受死.
然後第三日才復活.
從此.
這裡就是科林書情節的一個最大的轉捩點.
科林書的作者說.
從此他就定義要去耶路撒冷.
所以這裡就是耶穌全路的工作.
在北面加里的結束.
之後有很多的細節.
有很多的內容.
就是回到南面耶路撒冷發生.
我們結束的時候.
看二十四節至二十六節.
接著耶穌對門徒說.
如果有人要來跟從我.

$^{641}$就應當寫記.
背起他的十字架來跟從我.
在盧家福音甚至加了一個形容詞.
天天背起他的十字架.
是一個習慣.
是每一天的經常.
我們要跟主一起受苦.
放下我們自己的意識.
天天這樣做.
二十五節.
凡是要救自己生命的.
必喪失生命.
凡是為我喪失生命的.
必找到生命.
人就算賺得全世界.
卻賠上自己生命有什麼好處.
還能拿什麼換取自己的生命.
這裡說到兩種生命.
你救地上的生命.
即是我們現在有血有肉的生命.
如果喪失了永生的話.
那種真正的生命.
永久的生命.
有什麼好處.
就是救地上的生命.
保存自己現在的安好.
而賠上我們所喪失的.
就是永恆的生命.
沒有什麼好處.
我們作為主的門徒.
要記起.
既然耶穌背十字架死.
才會有榮耀.
我們跟隨主的時候.
其實都是一樣.
拿起十字架.
即是捨己的意思.
去遵行神旨而為上的意思.
\newpage



\section{}
\label{sec:jXnsElF6xnM}
\textbf{聖地網上行 - 步主腳蹤 [第九集] 榮進耶路撒冷}
\newline
\newline
連結: \href{https://youtube.com/watch?v=jXnsElF6xnM}{\texttt{https://youtube.com/watch?v=jXnsElF6xnM}} ~~~~ 語音日期: 2020-08-27
\newline
\newline
\hyperref[sec:vLHlHGB1MwM]{\small{< < < PREV SERMON < < <}}
~
\hyperref[sec:index_chronic]{\small{[返順時目]}}
~
\hyperref[sec:index_scriptual]{\small{[返順卷目]}}
~
\hyperref[sec:PFNNLzN6LA8]{\small{> > > NEXT SERMON > > >}}
\newline
\newline
$^{1}$今天是聖地網上的第九集.
我們時間上.
在說猶太人的御節之前大概一個多星期.
御節是一個很重要的節期.
散居不同地方的猶太人.
都會視這個為大節.
因為法書裡面說到.
猶太人要回到耶路撒冷的聖殿守節.
加里尼的猶太人.
他們在猶太人的聖地上.
他們會用幾天的時間.
由北面下到南面的耶路撒冷.
因為猶太人和撒瑪利亞人之間.
有很多的隔膜和仇恨.
所以大多數的猶太人都會避開.
不會經過撒瑪利亞的山地.
在一條更近的路下去耶路撒冷.
而是沿著約旦河谷.
一直南下一直往北.
一直南下一直到死海的北岸.
那裡就看到耶利哥.
然後直著耶利哥上山路.
就上到耶路撒冷.
你記得在耶利哥附近.
曾經發生過什麼事呢.
就是撒瑪利亞約翰約翰.
他開始幫人撒瑪利亞的工作.
和耶穌受洗的地方都在這裡.
馬太福音第十九章第一節.
耶穌說完了這些話.
就離開加利利.
來到猶太的境界約旦河外.
有許多人跟著他.
他就在那裡.
給他們的病人治好了.
這裡講到耶穌最後的一段路程.
上耶路撒冷.
在約旦河谷一直南下.
兩旁可以看到很多的高山峭壁.
過往幾十年甚至百多年.

$^{41}$大希律王或者之前的瑪加比王朝.
就在這些山崖上興建了不同的碉堡.
Palace Fortress.
作用當然是守住東方線.
Palace Fortress的另一個用途.
就是給貴族和王室的人來享用.
因為約旦河峽谷是地球最低的峽谷.
所以在冬天的時候.
這裡的氣候就會比較溫和.
對於這麼多代的貴族.
猶太人的以前的王.
這些權貴人士.
他們通常會在冬天的時候.
暫時離開高山上的耶路撒冷.
來到這裡避寒.
所以這裡這麼多的碉堡.
都是一些冬宮.
是給人避寒的.
在耶穌的時代也是一樣.
猶太的貴族仍然在使用這些碉堡.
從北至南可以發現.
有很多不同的Palace Fortress.
這些堡壘和皇宮的碉堡.
我們可以看到有亞歷山大堡.
Alexandria.
在約旦河東岸的馬基路.
Makaris.
這裡就是施洗日漢被斬的地方.
下面就有耶利哥.
雖然耶利哥是地上的一個城鎮.
這個新約的耶利哥.
也有大希律.
甚至之前馬加比時代建的皇宮堡壘.
再下面就是塞浦斯.
Cyprus.
可以守護整個耶利哥的平原.
再下面就是一個堡壘.
叫做赫加利亞.
Hycania.
再下面就是一個很有名的堡壘.

$^{81}$叫做馬薩達.
這些都是比較大的堡壘.
我們也沒有提到.
耶路撒冷也是一個堡壘.
防守非常嚴密.
另外我們曾經提過的.
Herodium.
希律堡.
這些都是大希律.
或者之前馬加比王朝.
來鎮守東方線.
或者是給自己貴族享用的城堡.
適逢預節的時候.
民族情緒很高漲的時候.
猶太人在路上.
去耶路撒冷首節的時候.
他們一定會傾談.
講論.
溢述猶太人王國的政治問題.
也很可能會聽到路上的人來探論.
以前馬加比王朝的時候多厲害.
我們擊潰了希臘的軍隊.
然後自立成為一個國家.
雖然只有100年左右的時間.
你會聽到別人說.
以前大希律的家族.
甚至現在他的兒子.
來接管他的國土的時候.
他們也有半個猶太人的國土.
得到的國土幾乎像.
以前所羅門王的版圖一樣.
這些政治故事.
一定是猶太人很熟悉的.
也是門徒很熟悉的.
怪不得耶穌說上耶路撒冷的時候.
每個人都想著.
這個王者來到了.
到他建立這個國土的時候.
到底誰坐第一位.
誰坐第二位呢.

$^{121}$所以上耶路撒冷的那幾天.
門徒都在爭誰坐第一位.
誰坐第二位.
耶穌用很多的方法去教導他們.
天國不是這樣的.
是要去服事.
並且耶穌三次來預言.
自己上耶路撒冷.
並且死在十字架上.
要先受苦.
然後才得到榮耀.
死在裡面服務.
我們繼續南下的時候.
就經過耶利哥.
剛才說過.
耶利哥是大希律所建的皇宮寶殿.
這裡也是大希律最愛的東宮.
沿著耶利哥.
一直要上山上耶路撒冷.
路程大概有20公里.
或者15英里的路程.
要經過猶太曠野的山.
攀登的海拔大概需要900米.
或者3000尺左右.
要的時間大概6-8個小時.
所以應該一早上他們就出發.
然後花一整天的時間.
一直上山.
因為到了晚上.
要經過這條路的話.
就會有危險.
你記得耶利哥和耶路撒冷之間.
這一段山路.
經常有強盜出現.
好撒瑪利亞人的比喻.
就是用一個這樣的寫實情景.
來講述出來.
在這一段路.
我們一直上耶路撒冷的時候.
我們就看到一間博物館.

$^{161}$叫做好撒瑪利亞人博物館.
這裡當然不是真的故事發生在這裡.
而是四世紀開始.
有一些拜占庭的基督徒.
來這裡建了一間教堂.
現在這個博物館.
收藏了很多在聖地裡找到的馬賽克地板.
上耶路撒冷首節的時候.
猶太人會唱什麼詩歌呢?.
你記得有上行之詩.
詩篇120至134篇.
這15篇的詩篇.
就是時常提到.
石安,撒安,耶路撒冷.
上行的朝聖者.
一路途上會經歷不同的困難.
威脅或者危險.
這些上行之詩就鼓勵他們.
要去倚靠做天地的主.
但你要留意.
耶穌上耶路撒冷的時候.
祂不是去一個中性的地方.
因為在福音書一直交代.
耶路撒冷或者猶大.
這個地方是敵對耶穌的地帶.
他們曾經有文士,法理才人,宗教領袖.
被派到加尼尼質問和試探耶穌.
你可以想像一下.
耶穌上耶路撒冷的途中.
聽到或唱上行之詩的時候.
會有什麼感受.
我們一路向耶路撒冷上山進發的時候.
快要到達耶路撒冷的時候.
就可以看到一個山脊.
這個山脊叫做橄欖山.
有幾個山峰.
Mount of Olives.
橄欖山的後面或西面.
就是耶路撒冷石安山.
就是他們的目的地.

$^{201}$我們跟隨約翰福音的編排.
耶穌未騎到進入耶路撒冷之前.
祂先來到橄欖山的東面.
在我們前面有一個小村莊.
叫做伯特利.
Bethany這個字的意思是.
Bethany.
意思是House of the Poor.
即是窮人的村莊.
或者House of the Fix.
即是無花果之家.
無花果的村莊.
距離耶路撒冷大概三公里或兩英里.
事實上每一次耶穌來到猶大地.
或耶路撒冷的時候.
祂只會留在城外.
在伯特利這個小村莊.
你記得這裡有什麼主要的事跡發生過嗎?.
在伯特利這個地方.
在巴勒斯坦地區.
在西岸的一個教堂.
你會略略知道.
這三幅圖像是.
馬大,瑪利亞和拉薩路.
這裡是聖Nazareth Church.
即是聖拉薩路教堂.
是一個現代的天主教教堂.
建造在古代教堂的遺跡上.
這些遺跡可以追溯至公元後四世紀.
教堂內也有不同的馬賽克和圖像.
教我們想起耶穌在伯特利這個地方.
和一個小家庭有很密切的關係.
你記得故事的背景.
就是拉薩路的死和復活.
叫耶穌趕來這個地方.
希望能夠救回他的弟弟拉薩路.
亦有耶穌和馬大,瑪利亞在他家裡趕到.
這裡的教堂的設計.
裡面沒有任何窗口.
形成了一個比較黑暗的環境.

$^{241}$讓我們想起這些墳墓的設計.
僅僅在圓頂上有玻璃的光透下來.
讓我們想起拉薩路的復活.
就是耶穌的生命,復活,光.
帶給他的門徒新的生命.
拉薩路復活的神蹟發生在伯特利.
約翰福音記載了耶穌有七個神蹟.
而拉薩路復活約翰福音第十一章.
就是第七個最後一個.
亦是最大的神蹟.
第一個神蹟在約翰福音.
就是在迦那的婚禮裡面.
耶穌將水變成酒.
所以整個約翰福音.
按照他記載耶穌的神蹟奇事的時候.
首先就是一個婚禮開始.
而最後到第十一章.
就是在喪禮裡面行出他的神蹟.
所以是頭尾有呼應.
因為約翰福音的神蹟裡面.
叫人想起一些摩西出埃及的事蹟.
你記得摩西出埃及的故事.
第一個神蹟是什麼呢?.
第一災就是水的神蹟.
將水變成血.
耶穌在約翰福音第一個神蹟.
也是水的神蹟.
將水變成酒.
摩西最後一個出埃及的神蹟.
就是跟死跟生有關.
第十災就是殺長子之災.
審判埃及.
然後以色列人離開埃及有生命.
在約翰福音裡面.
耶穌最後一個神蹟.
就是他的朋友拉撒路從死裡面復活.
我們之前說過.
耶穌有好幾次來過耶路撒冷.
這次絕對不是第一次.
但是這次是最後一次.

$^{281}$在他最後的一個星期.
所以他進耶路撒冷的時候.
有特別的安排.
我們先看馬太福音第21章.
這裡提到一個地方叫做伯伐奇.
就是在伯特利再往前走一點.
伯伐奇英文叫Bethphage.
或者Bavarcki.
意思是House of Unripe Figs.
就是未熟無花果樹的村莊地方之家.
記得耶穌在這個星期裡面.
曾經做過什麼呢?.
就是耶穌就坐不結果的無花果樹.
無花果樹就是代表不結果的耶路撒冷聖殿.
那些反對敵對耶穌想殺害他的那群人.
猶太人的拉比文學也曾經提過伯伐奇這個地方.
因為猶太人拉比很緊張耶路撒冷聖城的邊界.
為什麼這麼重要呢?.
因為祭司要預備聖殿裡面的感恩祭.
奉獻的麵餅,陳設餅等等.
一定要在耶路撒冷城的以內去做.
而伯伐奇拉比的討論裡面.
說這裡是最東邊的邊界.
也是屬於耶路撒冷.
也是一個安石的路程.
大概半英里或者四分三英里的距離.
所以來到伯伐奇的時候.
就獻上合乎聖殿月立的祭物.
這裡當然有一點寓意.
就是耶穌是那個羔羊祂自己獻上.
經文也提到耶穌交代門徒.
說拿那隻驢過來.
密碼就是主要用他.
就是對驢的主人說主要用他.
他就會放那隻驢給你用.
這裡可見耶穌對於一些已經信祂的門徒.
早有預備了一隻驢.
來到伯伐奇的地方.
這裡有一座教堂.
紀念耶穌在附近騎驢進入耶路撒冷.

$^{321}$你看看教堂上的牆壁和畫像.
都是讓我們想起這道聖經的事蹟.
包括門徒拿著驢.
包括他經過耶利哥發生了什麼事.
這個人在樹上應該是撒該.
在耶利哥那裡順了主.
然後慷慨地還錢給人.
甚至奉獻出來.
不如我們看看耶穌進入耶路撒冷的情況.
約翰福音第十二章.
第二天有許多上來過節的人聽見耶穌將到耶路撒冷.
就拿著中暑之外去迎接他.
喊著說.
和撒拿奉主名來的以色列王是應當稱頌的.
耶穌得了一個驢驢就騎上.
如經上所記的說.
錫安的民也不要懼怕.
你的王騎著驢驢來了.
這些是門徒起先不明白.
等到耶穌得了榮耀以後.
才想起這話是指著他寫的.
並且眾人果然向他這樣行了.
這裡的事蹟我們會叫做榮進耶路撒冷.
Triumphant Entry of Jerusalem.
裡面提到的幾件事是值得我們留意的.
中暑之人群將他們的衣服.
綁一些中暑之.
然後散在路上.
然後揮動去歡迎耶穌皇者進城.
這種情景一定讓猶太人想起.
大概二百年前.
就是西門瑪加比去擊敗希臘的軍隊.
然後海船回歸進入耶路撒冷.
因為瑪加比一書這樣記載.
當時的戰勝者西門瑪加比.
擊退了軍隊之後.
群眾手拿著中暑之.
然後高唱海歌.
他們歡迎得勝者回到耶路撒冷.
然後結成聖殿.

$^{361}$成立瑪加比的王國.
另外民眾唱一首詩歌.
和薩那和薩那.
奉主明來的以色列王子應當稱頌.
這裡直接引用一篇詩篇叫做118篇.
猶太人有自己讀經的方法.
或者有一些經文他們會組成一起讀.
113到118.
六首的贊美詩.
就是他們預設 Passover 的時候.
他們一定會唱和讀.
詩篇118就是最後一篇.
這裡提到神的子民要在困難中重建聖殿.
求上帝幫助和作王.
和薩那的意思是Save Now.
就是拯救我們.
現在就來拯救我們.
所以用到這篇詩篇的時候.
就說到要重建耶路撒冷.
王者終於來了.
Save us now.
去拯救我們.
另外在事跡的記載裡面.
有一個意境就是騎著路.
這裡是唯一一次描述耶穌騎著路.
因為一向都是假設.
或者知道祂是走路的.
那麼耶穌從加利利下來.
和門徒一起來到守節.
祂已經走了一百多里路.
現在在白髮期來耶路撒冷.
只不過是十分鐘步行.
其實是不需要騎著路的.
另外猶太人進入聖城的慣常方法.
都是走路的.
以表尊重.
所以祂特意安排一條路.
其實是想作一個宣告.
以致門徒以後再想起.
這個就是我們所信的耶穌.

$^{401}$是耶利賽亞.
祂是王者.
騎著路進入耶路撒冷.
你會想到什麼呢?.
你想起舊約裡面.
首先是耶路撒冷.
大衛首先是建都於耶路撒冷.
大衛的王國流亡之後.
滅亡之後.
都沒有真正出現過大衛的君王.
大家在兩約中間.
都是在等候大衛的子孫.
大衛的兒子.
來重建消亡已久的大衛國.
你也想起大衛的兒子.
所羅門登基的時候.
他也是騎著路的.
在基分泉下面.
在耶路撒冷和加cum 山中間.
有一個峽谷.
下面就是基分泉.
這個地方就是.
所羅門被高納作王.
成繼大衛的寶座.
另外這裡提到經上所記.
就是幾百年前.
有一位先知叫做撒迦利亞.
他預言將來會有一位王者.
來統治以色列.
撒迦利亞書第九章第九節.
那裡提到耶路撒冷的民.
你們應當歡呼.
汗牙利的王終於來到了.
他是公義的.
並且私人拯救.
並且他是謙謙和和的騎著路.
下面提到這位王者.
來到的時候.
他就會帶來和平.
平息戰爭.

$^{441}$所以騎著路就是應驗.
幾百年前說過的利塞的認許.
並且耶穌刻意騎著路入城.
顯明他的天國.
他的王位.
不是靠武力帶來的.
而是騎著謙卑的路.
不是騎著戰馬.
說起戰馬.
有件趣事.
這個白伐奇的天主教堂.
建於十九世紀末.
其實它是蓋在一個.
十二世紀十字軍時期的遺跡上.
當時十字軍的古教堂.
裡面最重要的物件.
就是中間那塊大石.
這塊方形的很高的大石.
被譽為基督上馬之石.
即是Mountain Stone of Christ.
十字軍的信徒設想.
耶穌上路的時候.
就要用這塊石來做踏腳石.
我想當時的十字軍人.
他們看聖經就不是很小心.
因為他們忘記了.
耶穌上的路是一條很低微.
很矮的路.
不是他們那些高高的戰馬.
這麼大塊石.
上戰馬就對了.
但是上路就不對了.
耶穌是謙謙和和的.
騎著道路.
不是騎著戰馬.
誰會騎著戰馬進來呢?.
如果再看遠一點的話.
就是每逢預節的時候.
羅馬的官員.
總督本雕彼拉多.

$^{481}$就會帶著軍隊.
羅馬的戰馬進城.
他就在西面進城.
耶穌就在東面進城.
但是羅馬的進城方法.
就是鎮壓異己.
看著猶太人不要造反.
並且會去審案.
我們知道之後的幾天.
耶穌被捉拿的時候.
帶到彼拉多面前.
這不是湊巧的事.
平時彼拉多都會定期來耶路撒冷.
特別是節期之前.
去審案.
耶穌騎著道路.
永進耶路撒冷作王.
我們就想起一個基督教的節日.
就是復活節之前的主日.
叫做 Palm Sunday.
中樹節或中樹之節的主日.
我們想起 Palm Sunday.
就是耶穌當時永進耶路撒冷.
或者每一個年度.
去慶祝中樹節主日的時候.
都有很多歡呼聲.
可以看到那些人.
在耶路撒冷就拿著中樹枝.
他們會唱和剎那和剎那和聲.
但是否就是這麼簡單?.
是否就是神來拯救我.
現在你去得勝我的敵人.
在四福音裡面.
都有記載永進耶路撒冷的一件事.
但是約翰福音就有些特別的地方.
因為約翰福音的記載裡面.
沒有花什麼篇幅或細節去講.
耶穌有什麼預備.
叛徒有什麼預備.
耶穌對群眾有什麼對答.

$^{521}$全部約翰福音都沒有提及.
很簡短.
因為約翰福音想做一件事.
就是想集中處理周邊的人.
特別是耶路撒冷裡面的人.
如何去回應耶穌騎路進入耶路撒冷.
我們讀讀約翰福音第十一章.
這裡是剛剛耶穌叫拉撒路從死裡復活.
行完神蹟.
然後還沒有永進耶路撒冷.
還沒有騎路進入耶路撒冷.
中間發生什麼事呢?.
十一章四十五字開始.
那些來看馬尼亞的猶太人.
見了耶穌所做的事.
就多猶順他的.
但其中也有人去見法利賽人.
將耶穌所做的事告訴他們.
祭司長和法利賽人聚集開公會說.
這人(耶穌)行好些神蹟.
我們怎麼辦呢?.
若這樣猶著他.
人人都要順他.
馬人也要來奪我們的地土和我們的百姓.
內中有一個人名叫蓋亞法.
是本年做大祭司的.
對他們說.
你們不知道什麼.
都不想一個人替百姓死.
免得通國滅亡.
就是你們的逆處.
蓋亞法這話不是出於自己.
是因他本年做大祭司.
所以不知覺地.
他預言了耶穌將要替這一國國民死.
不但替這一國死.
也要將上帝四散的子民都聚集歸一.
從那天起.
他們公會的人.
包括大祭司.

$^{561}$就相意要去殺害耶穌.
他們有什麼反應呢?.
他們覺得耶穌很威脅他們的地位.
因為這裡說到.
他們是和羅馬帝國合作串謀.
耶路撒冷是祭司長握羅馬的政權中心.
這不是一個忠誠的地帶.
有如我們今天說的.
白宮.
北京.
這代表著人有很多利益.
有很多政治的權勢.
有很多的財富.
公會的人說什麼呢?.
他說如果猶著耶穌去做他現在做的事情.
人人都會相信他.
所有的百姓都會歸他.
那麼羅馬人會怎麼做呢?.
羅馬人會來搶回我們的土地.
以及我們的百姓.
你有沒有留意他們說什麼?.
他說這些是我們的東西.
還有這裡說到地土.
其實不只是指國土這麼簡單.
地方這麼簡單.
地土特別是指聖殿.
在《士徒行傳》.
在以前《兩約》中間.
說到地方,地土的時候.
很多時候都是指聖殿這個地方.
無論如何他們覺得耶穌對他們很大衛脅.
所以他們千方百計想保存什麼呢?.
就是他們的現況.
他們的Status Quo.
我不知道你想起Status Quo的時候.
會想起什麼呢?.
我們人人都會有所謂的.
「現況」,「現狀」Status Quo.
而Status Quo很多時候就是.
我們一己的平安.

$^{601}$我們一己的快樂.
或者我們可以隨便.
按照自己的心意.
去分配我們的資源,資金.
可能這些就是我們的Status Quo.
或者這麼說.
我們其實跟當時的耶路撒冷權貴.
可能都不是太大分別.
因為你想想他們當時工會商議的是什麼呢?.
未必是他們眼中.
或者在社會的眼中是彌天大罪.
我們最重要的是什麼呢?.
是想固存大局.
留意一下大祭司說什麼.
當時的宗教局局長或者政治局局長.
留意他子女行間想說什麼.
總而言之.
他想保存Status Quo.
覺得耶穌出現.
打破他們改變了他們的現況.
他們的利益或者他們的福利.
他們有的好處就會被威脅.
其實可能我們都會這樣想.
或者我們會這樣對待信仰.
我們可能會覺得.
信主.
認耶穌為我的主.
其實很多時候我們都會淪落到一個地步.
只不過是為我的福利.
為我天堂的福利.
可以這麼說.
我們是以一己的現況為至上.
我現在的Status Quo就最好了.
我是信耶穌.
不過我不想改變我的現況.
或者我不想信耶穌.
我想除滅耶穌.
原因就是因為我為了我現在的現況.
我想保住.
這種以一己為出發點.

$^{641}$所謂保存大局.
就是和上帝的國剛剛相反.
王來到就理應震動我自己一切的生活.
震動我自己本以自我為中心的生活角度.
所以他會影響和威脅我的Status Quo.
可能我以往生活的方式.
是不合乎神心意.
是不能彰顯到天國.
耶穌來到的時候.
我歡迎祂進入我的生命當中.
不只是歡呼聲.
而是真的有一個重整的架構.
所以和薩那和薩那.
主作王主來救我.
不只是主德勝我的受敵.
以至我可以安於現況.
我的Status Quo.
和薩那和薩那主作王.
原來就是主耶穌來到我的生命當中.
是祂勝過我.
以前初信主的時候.
讀過一本書叫做《主居我心》.
My Heart, Christ's Home.
這裡有一個比喻.
是說到在一間屋.
我本身是主人.
我邀請了耶穌進入我的生命當中.
進入這間屋的時候.
祂開了每一個屋的角落.
我的客廳.
即是我有的關係.
我的工作室.
我娛樂的地方.
我工作的地方等等.
每一個房間我都需要交出鑰匙.
讓主耶穌來作主.
這就是所謂權力的重整.
哪一個拿鑰匙.
哪一個作主.
哪一個的旨意先行.

$^{681}$當耶穌進入我們的生命當中.
我們有沒有一些地方.
仍然未被祂作主.
我們低頭祈禱.
主我們求你原諒我們.
赦免我們.
很多時候我們將福音和信仰.
簡約或斬簡度.
只為了我一己的好處.
為了救我 為了獨善其身.
主我求你親自來幫助我們.
讓我們認真思想.
當你進入我們的生命當中.
當我們稱你為我們的主.
我們的君王的時候.
是你親自來掌管.
我們生命的每一個角落.
包括我們所謂的宗教生活.
我們去敬拜.
我們去團結我們的教會.
包括我們的職業.
包括我們社會.
政治中.
我們需要彰顯自己國度的真理和公義.
主你做我們整個生命當中的主.
求你親自來幫助我們.
我們這樣祈禱.
奉主耶穌基督的名字.
\newpage



\section{}
\label{sec:PFNNLzN6LA8}
\textbf{聖地網上行 - 步主腳蹤 [第十一集] 最後晚餐}
\newline
\newline
連結: \href{https://youtube.com/watch?v=PFNNLzN6LA8}{\texttt{https://youtube.com/watch?v=PFNNLzN6LA8}} ~~~~ 語音日期: 2020-08-28
\newline
\newline
\hyperref[sec:jXnsElF6xnM]{\small{< < < PREV SERMON < < <}}
~
\hyperref[sec:index_chronic]{\small{[返順時目]}}
~
\hyperref[sec:index_scriptual]{\small{[返順卷目]}}
~
\hyperref[sec:c1irNOxf2YQ]{\small{> > > NEXT SERMON > > >}}
\newline
\newline
$^{1}$我們來到耶穌生平的最後一個星期.
我們講過耶穌榮進耶路撒冷.
然後他結成聖殿.
之後幾天按照福音書的記載.
耶穌與耶路撒冷的宗教權貴.
發生了很多爭議和辯論.
我們來到猶太人的預節的時候.
按照猶太人的規矩.
最好預節的晚餐.
就是在耶路撒冷城的.
最高的地方.
就是在耶路撒冷城裡面吃的.
所以馬可福音第十四章就講到.
耶穌如何去預備預節晚餐.
除夭節的第一天.
就是宰殺預節羊羔的日子.
門徒問耶穌.
你要我們到哪裡去.
為你預備預節的晚餐.
他就差派兩個門徒對他們說.
你們進城去.
必有一個人拿著一瓶水迎面而來.
你們跟著他.
無論他進哪一家.
你們就對那家的主人說.
老師說.
讓我和門徒吃預節晚餐的客房在哪裡.
他必給你們看到.
樓上一間擺設好又預備好的大房間.
你們就在哪裡為我們預備吧.
這裡很簡單地去記載.
他說密碼就是有一個人.
男人.
拿著一瓶水迎面而來.
打水通常是女人的工作.
所以見到一個男人拿著一瓶水.
是非常罕見.
可以見到耶穌一早已經預備了.
在耶穌的紀念傳道裡面.
他都吸引了一些少數的.

$^{41}$較為上層的人.
所以來到耶路撒冷城的時候.
他想找一個大的房間.
甚至是樓上的房間.
來到做最後晚餐的時候.
這個是可以理解的.
這裡說到樓上一間大房間.
這是新漢語譯本.
其他中文譯本沒有說得很清楚.
這是上樓房.
是upper room.
我們知道在當時耶路撒冷的地理.
上城或者聖殿西邊這個地方.
這裡的房子比較大.
而且是比較貴族的地方.
所以比較可能這個上樓房會在這一區.
今天去耶路撒冷遊覽的時候.
通常會去到現在的西山.
Western Hill.
就是現在舊城耶路撒冷的西邊.
西南角.
這裡有一個古舊的石樓房.
很多時候傳統都叫上樓房.
就是upper room.
這個傳統比較遲一點.
比較後期一點.
是主後五世紀才開始出現.
所以不是很可靠.
我們極其量只知道當時五世紀.
這裡有一所教堂.
然後十二世紀.
十字軍來到耶路撒冷.
他們就在這所教堂上建築另外一所教堂.
所以你看到這個所謂上樓房.
有一些十字軍的拱門.
之前我們有好幾集都見過.
十六世紀的時候.
伊斯蘭教徒入主聖地中東.
然後來到耶路撒冷.
就將這個地方改建為一個清真寺.

$^{81}$在牆上不同的地方不同的地點.
有一些清真寺的擺設和裝修.
另外我們也知道.
到第一世紀之後.
石安山這個地名有些混淆.
所以慢慢地就將這個所謂的石安山.
本來是耶路撒冷東山或聖殿山的地方.
就遷移到西邊.
所以今天我們叫這個西山.
或者叫石安山.
但其實石安山原本不是向這裡.
所以有這些混淆的緣故.
後期有很多傳統就引發出來.
所以這個上樓房下面居然是大衛的墓穴紀念的地方.
這是猶太教今時今日都認為是一個聖地之一.
當然這個不是可靠.
因為聖經沒有錯有說.
大衛是葬在石安山.
但因為石安山的地方.
因為後期的混淆已經轉移了.
所以樓上就是基督教的所謂聖地.
紀念這個上樓房.
因為伊斯蘭教徒也來過.
所以這裡就變成了一個清真寺.
所以我覺得雖然這裡不是最可靠的地方.
但卻是點出了這個地點最有趣的地方.
就是三大一神宗教交會點.
就正是在這個上樓房.
基督教,猶太教和伊斯蘭教都有份.
這就是聖地的特色.
想起耶穌和門徒的最後晚餐.
你可能會想起這幅名畫.
就是十五,十六世紀.
達文西所畫的最後晚餐.
當然耶穌的時代一定不會這樣一字排開來吃.
這幅名畫有當時的背景.
如果你想知道更多的話.
我們城匯網台下面有一條連結給你.
我們有一段短片是講到藝術和聖經的關係.
討論到這幅名畫.

$^{121}$無論如何聖經所形容的預節晚餐.
當時用的風格應該是羅馬帝國最流行的風格.
我們叫做Triclinium.
三邊臥餐座.
即是在飯廳的房間.
三邊的牆有較為低的臥餐座.
像梳化一樣形成一個U形.
中間的空位讓僕人進出服侍和上餐.
客人是斜臥在左側.
然後用右手拿食物.
腳向外伸.
所以僕人可以為客人洗腳.
主人家坐的不是中間的臥座.
很多時是坐在旁邊的臥座.
因為是對著門口.
所以是最矚目和最尊貴的地方.
如果你是主人的話.
你的左邊就是貴賓位.
右邊通常是另一個朋友.
或者是家中最年輕的一個.
在三臥餐座的擺位.
按照客人的尊卑.
由最大的到最小的.
最靠近門口的就是最卑微的位置.
我們在意大利的Pompeii.
有很多三邊臥餐座的飯廳.
牆上有些濕壁畫.
當時是最流行的顏色.
叫做Pompeii Red.
龐貝的紅色.
在意大利的地方.
很多時牆上有神像和人像的裝飾.
或者動物的繪圖.
羅馬是在儲錢63年.
佔領以色列巴勒斯坦地.
之後我們見到有不同的猶太貴族.
開始引入這種飯廳的格局.
例如我們之前看過很多大希臘的雕堡.
Palace Fortress.
全部的雕堡都有這些客廳.

$^{161}$現在我們看到馬薩大的情況.
在馬薩大的北邊有三層的大庭院.
我們來到北邊最北的客廳.
我們看到有些柱來裝飾.
上面本身有外頂.
在客廳後面的牆.
也裝飾了龐貝的裝飾.
是龐貝紅色.
當然在猶太人的眼中.
這些人像和神像是不可以放上去的.
所以很多時都是靜色.
在這裡一邊和貴族一起吃飯.
餓在這裡一起享受.
還可以欣賞到死海北邊的風景.
是多麼壯觀.
同時我們在耶路撒冷.
也找到這些三臥餐桌的格局.
在舊城的猶太人區.
考古人員在七零年代已經開始找到一世紀的豪宅.
我們來到一個叫Wall Archaeological Museum.
裡面是一個地庫.
上面是現在居民的地方.
下面因為發掘了這些東西.
成為了一個博物館.
我們來到這個地庫時.
發現這裡所發掘出來的.
本身是一個幾層高的住宅.
面積是5500呎平方呎.
600平方米.
在牆上也找到一些很熟悉的.
是Freshco 濕壁畫.
顏色是Pompeii Red.
當時猶太的貴族不斷引入羅馬式用餐的方法.
在這裡的豪宅也看到了主後70年.
耶路撒冷聖殿聖城被毀時的火燒跡象.
在這麼大的豪宅裡面.
有不同的飯廳.
這裡只是一個比較小的.
地上鋪了馬賽克.
你可以想像在這裡有三邊放的梳化.

$^{201}$這些長臥椅.
成為一個U字形.
在這個小房間中間有一個瀑人可以出入.
後面的桌子像水吧一樣.
石桌子.
也有兩個猶太人用的潔淨水缸洗手.
你這樣想像的時候.
我們就可以明白當時耶穌和門徒最後晚餐的情景.
特別是有些經文我們未必能夠看得明白.
我們有這樣的圖畫的時候就能夠明白.
我們先選幾節.
我們去約翰福音第十三章.
這裡耶穌已經透露了猶太會出賣他.
22節.
門徒彼此對視.
不曉得耶穌所說的是誰出賣他.
有一個門徒就是耶穌所愛的.
他執我在耶穌的胸懷之前.
即是旁邊的意思.
這應該是約翰.
然後24節.
西門彼得向約翰示意.
即是說他們比較遠在對面.
他問耶穌這句話到底是指誰.
是誰出賣他.
西靠近耶穌胸前的門徒就問他主要是誰.
耶穌說要和我爭一點食物.
遞給他的就是出賣我的.
於是他就佔了一些食物.
然後遞給加里蘭西門的兒子猶大.
猶大吃了之後撒旦進入了他的心.
然後他當然離開了這個產場.
所以你可以看到耶穌坐在主席位.
然後他在右邊或是在他胸前.
應該是約翰.
後面就是猶大.
可以看到耶穌在最後的晚餐.
似乎他都是說我想愛這個受敵.
希望猶大能夠回轉.
不要背叛他.

$^{241}$當他仍然是他的密友.
他的好朋友.
給他最後悔改的機會.
彼得當然就在對面.
他不能夠直接親自問耶穌.
所以要示意對面的約翰.
所以彼得就坐在三河桌最卑微的地方.
近門口的地方.
這些羅馬式的餓座吃法.
很強調地位的轉變.
階級的轉變.
當時甚至有些希臘羅馬哲學家.
居然去辯論這些問題.
寫哲學文章去辯論.
坐位安排的問題.
所以難怪耶穌在最後晚餐時.
門徒仍然在爭誰大誰小.
我們再問多一個問題.
約翰福音或福音書一般來說.
其實不是寫給當時耶穌的門徒.
而是寫給第二代初代的教會.
當時的教會都是家庭的教會.
會眾來到家庭聚會.
聚會,聽道,吃東西的時候.
遇上洗腳的問題.
亦遇上安排座位的問題.
在豪宅裡面.
到底主席位誰坐?.
或者誰坐地下?.
耶穌和門徒的晚餐.
每個星期都在當時的家庭教會發生.
這種情景寫出最後晚餐的細節.
其實就是教導當時的教會.
如何坐?如何安排?.
如何彼此服侍?.
來顯出神的光?.
耶穌和門徒洗腳.
當中的情節有沒有深一點的意思?.
我們要留意.
在十三章第六節.

$^{281}$耶穌來到西門彼得的腳踐.
彼得對他說.
主啊,我的腳你也要洗嗎?.
耶穌回答說.
我所做的你現在不理解.
以後就明白了.
彼得說.
你絕不可洗我的腳.
永遠不可.
耶穌說.
如果我不洗你的腳.
你跟我就沒有關係了.
西門彼得說.
主啊,那就不要只洗腳.
連手和頭也洗吧.
耶穌說.
洗過躁的人就不需要再洗.
只要把腳洗了.
全身就乾淨了.
你們就是全身乾淨的人.
當然不是每一個都是乾淨的.
十一節.
耶穌已經知道誰出賣他.
這就是猶大.
你越讀就越亂.
第一.
耶穌要洗彼得的腳.
梳理一下問題.
他起初是不想的.
是反對的.
這個反應是可以理解的.
第二.
我們想問.
洗過腳的先有關係.
到底是什麼意思.
我們可以看一些動字.
約翰福音來到這個地方.
他特意把鏡頭變成慢鏡.
去留意每一個動作.
十三章第四節.

$^{321}$耶穌起身離席.
脫下外衣.
又用束腰的手巾擦乾.
洗完他們的腳.
穿上外衣.
這個說離席.
脫下外衣.
Lay aside.
在約翰福音出現過好幾次.
特別是第十章.
第十章第十一等等的經文.
耶穌說.
我是好牧人.
好牧人為羊捨棄生命.
這個就是脫下的意思.
Lay aside.
十五節.
我為羊捨棄生命.
同一個字.
十七節.
父愛我因為我把生命捨棄.
好再次取回來.
捨棄是同一個字.
取回來就是穿上外衣的意思.
同一個字.
十八節.
沒有誰能奪去我的生命.
是我自己捨棄的.
我既有權捨棄脫下.
亦有權再取回來.
這個就是洗完腳穿上外衣.
取回來穿回外衣的意思.
十三章第一節.
耶穌既然知道要離開這個世界.
他就要愛門徒愛到底.
底這個字.
Telos.
就是The End的意思.
就是耶穌在十加七言裡面.
其中一句說話.

$^{361}$在約翰福音裡面十九章三十節.
他說成了.
It is finished.
就是The End.
很清楚.
這個洗腳有個所謂的預表.
就是說耶穌為門徒洗腳的時候.
其實就是想指向很快會發生的.
就是他為門徒為了人.
在十字架上去死.
十三章十五節.
這裡說我給你們作了榜樣.
這個榜樣也是一個特別的字.
在聖經裡面.
甚至在兩章中間的文獻.
譬如《瑪加比書》.
經常都用到榜樣的字.
用到殉道的情節裡面.
就是那個人願意去犧牲.
為了國家.
為了保護子民.
他願意走出來犧牲.
約翰福音是寫給後期初代的教會.
我們知道約翰福音成書的時候.
是一世紀九十年代.
當時教會開始面臨壓力.
面臨打壓.
起碼約翰福音的修信讀者.
應該是伊弗所的猶太人教會.
他們面對著不同的壓力.
撒旦就是趁機去打散教會.
所以寫約翰福音的時候.
特別留意這一點.
這裡說耶穌如果要得勝的話.
我們不要自己顧自己.
我們要去學習去犧牲.
去給錢洗腳就是這個意思.
使教會能夠得以建立.
縱使遇見黑暗.
遇見打壓的時候.

$^{401}$在疫情延長的時候.
我相信各地的教會.
都是人心很不安.
加上有很多其他的動盪.
教會是其中一個受到打擊的群體.
甚至香港教會都面臨很多陰影.
國安法的陰影等等.
如果你這樣看的時候.
約翰福音就是寫給當時.
在不安的教會裡面.
你都要去教導他們提醒他們.
不單叫我們今天照辦主管.
其實是說誰想做大.
就要甘願去犧牲.
同一章的下文十三章.
講完洗腳.
耶穌就說我給你們一個新的命令.
要彼此相愛.
我怎樣愛你們.
你們要彼此怎樣相愛.
眾人就見到你們這樣犧牲的話.
這樣彼此服侍的話.
就是我的門徒.
在疫情當中.
我們屬於信主的人.
到底我們可以為信仰的群體做什麼.
怎樣去捨己,怎樣去犧牲.
以致我們在疫情當中.
能夠讓人見到.
這群人所說的神,所說的信主.
不是風花雪月的東西.
而是我們真的要理直.
願意服侍.
繼續在艱難當中堅定教會.
第十三章第八節.
彼得說你絕不可洗我的腳.
永遠不可.
耶穌說如果我不洗你的腳.
你跟我就沒有關係了.
為什麼洗腳又說有關係.

$^{441}$中文和學本就譯作與我無份.
這個字是一個inheritance.
基業的字.
是當時地產物業的字眼.
舊約耳色人入主聖地.
入主應許之地的時候.
有份是什麼意思.
代表在神的國土上有share.
約翰用這個字來形容.
耶穌的洗腳.
耶穌給我們的洗禮.
是給我們能夠在天國裡面有份.
所以啟示錄幾次最後用到這個字.
有份這個字,有關係這個字.
形容將來新的聖殿降臨.
新的耶路撒冷.
新天新地來到的時候.
我們接受了耶穌的救恩.
才會在裡面有基業,有份.
因為講到基督的救恩是最核心的.
不是其他外在的禮儀.
用水要做的事情.
你記得我之前說過.
約翰福音是寫給伊弗所的教會.
當時的群體.
是猶太人背景.
是寫於一世紀末的時候.
九十年代.
當時什麼大事發生.
剛才提過.
除了有壓力,有打壓之外.
就是在主后七十年的時候.
耶路撒冷的聖殿被毀.
對於猶太人有一個很大的衝擊.
因為猶太教有史以來.
這麼多年有很多實體的儀式.
就是他們信仰的記號.
是他們的核心.
第一樣我們想起的就是聖殿.
上聖殿去敬拜,上聖殿去祈禱,去獻祭.

$^{481}$是人生裡面最高的追求.
第二樣就是節旗.
猶太人的節旗對於他們來說.
是等於他們民族的記號.
所以按照節旗就要上聖殿去敬拜.
預節,主辦節,午巡,收割節等等.
第三樣實體的事情.
就是猶太教特別在兩約中間.
衍生了一種觀念叫做潔淨禮.
用水來潔淨.
所以每逢敬拜神的時候.
他們會經過一些受浸池.
洗過身.
潔淨了才能夠親近神.
我們之前去過的地方.
很多猶太人的遺址,地方.
都有這些潔淨池.
有些石吸進去.
然後經過之後就得到禮儀上的潔淨.
並且當時猶太人的家庭用的器具.
很多時候都是石器具.
Stone vessels.
拉比是這樣教的.
Stone vessels.
石器是不會沾染不潔的.
所以最好就是用石器.
無論如何聖殿被毀之後.
對於猶太人的身份,民族有一個很大的震盪.
到底我是誰?.
我還是不是神的兒女?.
神的子民剩下來的實體的記認.
就是用水潔淨.
所以有些文獻甚至說.
他們每一家去到首禦節的時候.
無論他們散居哪個地方.
猶太人首禦節的時候.
總之他們做足功夫.
用水洗淨之後.
每一家首禦節的時候.
就好像一個聖殿一樣.

$^{521}$一個小的寺廟.
來通往上帝,遇見上帝.
所以在《約翰福音》的背景.
你就可以明白.
為什麼他提出一些東西.
是其他福音書沒有提.
特別強調針對幾方面.
就是真正的潔淨.
罪人得到潔淨.
能夠有上帝的同在.
能夠來到上帝的聖殿裡面.
有份得到耶穌的救恩.
所以由第一章開始.
就提到這些潔淨禮.
Purification Ritual.
幾乎水潔淨的問題.
是貫穿整本的福音書.
我們很快看看.
第二章水變酒.
大家很熟悉這個故事.
但特別提到第六節.
就是那裡放了六口石缸.
就是猶太人行潔淨禮用的.
就是初一來到的時候.
水要變成酒.
接著就是潔淨聖殿.
第二章尾.
真正的潔淨就是耶穌的死而復活.
他的身體.
建立一個新的聖殿.
他就是聖殿.
第三章就是耶穌和猶太人拉比.
彌哥底姆去辯論.
就是水和靈才能夠重生.
這裡指向耶穌真正能夠成就的潔淨.
使我們能夠從水.
他所賜的水.
到聖靈裡面得到重生.
第四章.
去到撒瑪利亞井旁.

$^{561}$耶穌和婦人來交談的時候.
特別說到什麼呢?.
就是說到活水.
他說活水從哪裡來的呢?.
當時的聖殿也是有活水的傳源.
他說不是在耶路撒冷聖殿.
也不是在撒瑪利亞人的基尼森山上.
在耶穌的救恩裡面.
去到敬拜才能夠遇見上帝.
親近上帝.
第七章.
就是駐彭哲.
猶太人一個很大的節期.
這個最大的節期.
末期的時候.
猶太人就會祈求上帝賜水.
耶穌就高聲說.
人若渴了可以渡我者裡來.
我就賜給他什麼呢?.
在福中流出的活水江河.
他就是那個聖殿流出的江河的主.
第九章.
出現一個生下來是盲眼的人.
耶穌醫治了他之後.
就叫他到西羅亞池裡面洗.
他一洗.
回頭他就看見了.
所以這裡說到.
耶穌給的潔淨是比超越了.
猶太人西羅亞池的水好.
到十三章.
就是為門徒洗腳.
就談論到潔淨的問題.
就是耶穌的潔淨才能夠洗門徒.
在聖殿當中有份.
到十九章.
就是十字架的情況.
就是耶穌已經死了.
然後三十四節.
其中一個兵丁就用槍.

$^{601}$來刺耶穌的臘膀.
有血和水流出來.
如果你問.
聖經的信仰最核心的是什麼?.
是否一些外在的實物?.
最核心的我們一定會說.
是耶穌基督自己.
是他十架信的導引.
在疫情當中.
全世界的教會.
都受到不同程度的衝擊.
以前的教會我們習慣了實體聚會.
崇拜小祖團契.
甚至我們傳福音.
都是一些實體信仰的表白.
但在疫情當中.
特別是現場的疫情.
我們就失去了這些習慣.
或者暫時中斷了.
對於我們來說.
大體上是一種失去.
因為我們都不想這樣繼續下去.
但在疫情當中.
我們可以問一個更加核心的問題.
在疫情當中.
是否可以鍛鍊.
你對主,對神的信心?.
更加能夠鍛鍊到.
當我們失去了這些實體的移民的時候.
是否我們突然對神的心.
迷失了,失去了?.
換句話說.
我們以往的信仰.
是否只是建基於這些.
摸得到,看得到實體的東西?.
如果是這樣的話.
在疫情當中.
這個時候我們就需要重整.
重整一下我們信仰最核心的東西.
是否十字架上的耶穌?.

$^{641}$沒有了基督的話.
你想想.
怎樣有實體外在的東西.
其實都是空殼.
既然如此.
現在我們沒有了這些外在的東西.
我們應該回到.
真正基督與我們的關係.
你看了我們這個聖地網上.
有很多集.
其實我們就是想認識這位主宰.
認識耶穌基督.
十字架上的基督.
如果你今天還是有些迷失.
我請你回到耶穌基督的裡面.
可以同心低頭祈禱.
神在我們求你來幫助我們.
無論我們現在在哪一個地方.
我們開始認識你.
更加看你為寶貴.
主要是復興我們的心.
讓我們的信仰不單單建立在.
人造出來的一些移民或制度.
主要我們的心要歸你.
我們要相信你.
求你復興我們的心.
去接受耶穌基督的救恩的時候.
我們的生命得到更新,改變.
你潔淨我們以至我們有新的生命.
你都會去服侍你.
主要是復興我們每一個.
無論我們世界各地哪一個地方的信徒.
讓我們真正能夠成為你的門徒.
詞:陳曦 曲:陳曦.
十字架 身體光芒.
溫柔又慈祥.
帶著主愛的力量.
向著我照亮.
我的心不再隱藏.
完全地白裳.

$^{681}$願主愛來教導我.
在愛中的自由釋放.
我願意相扶.
在你愛的懷抱中.
我願意相扶.
你是我的主.
永遠在你懷抱中.
你是我.
你是我的主.
我願意相扶.
在你愛的懷抱中.
我願意相扶.
在你愛的懷抱中.
我願意相扶.
在你愛的懷抱中.
我願意相扶.
你是我的主.
永遠在你懷抱中.
你是我.
你是我的主.
你是我.
你是我的主.
\newpage



\section{}
\label{sec:c1irNOxf2YQ}
\textbf{聖地網上行 - 步主腳蹤 [第四集] 加利利 - 迦百農}
\newline
\newline
連結: \href{https://youtube.com/watch?v=c1irNOxf2YQ}{\texttt{https://youtube.com/watch?v=c1irNOxf2YQ}} ~~~~ 語音日期: 2020-08-29
\newline
\newline
\hyperref[sec:PFNNLzN6LA8]{\small{< < < PREV SERMON < < <}}
~
\hyperref[sec:index_chronic]{\small{[返順時目]}}
~
\hyperref[sec:index_scriptual]{\small{[返順卷目]}}
~
\hyperref[sec:aXI0__JmVi0]{\small{> > > NEXT SERMON > > >}}
\newline
\newline
$^{1}$各位好,我們今天來到第四集,我們繼續聖地望上行.
上一次講到耶穌在施洗約翰那裡受洗.
然後四十天在猶大曠野受試探.
那時候他還沒有正式公開傳道.
馬可方第一章十四節怎麼說呢?.
他說施洗約翰被捕以後.
耶穌才去到加利利宣講神的福音.
他說時而後近,神的國近了,應當悔改來相信福音.
所以是講到施洗約翰被逮捕之後.
為什麼呢?我們要交代一下那個證字.
我們如果再重溝一下時間表的話.
我們知道耶穌一直在猶太地跟約翰同工.
不久之後施洗約翰就北上到薩馬利亞省附近.
一個叫做靠近撒冷的艾亂斯西.
約翰福音第三章二十三節這樣說.
然後約翰又將他的工作擴展到加利利省.
所以一步一步北上到加利利的時候就發生了一些事情.
因為那裡施洗約翰譴責希律王安提柏王.
發生什麼事情呢?.
本來安提柏王Helotantipus.
他的原配就是阿拉伯的一個民族叫做Neptunian的公主.
他居然去幽妻.
然後另娶同父異母的兄弟.
肥力的太太.
太太的名字叫希羅底.
娶了他為妻.
可想而知這件事情是多麼讓猶太人不滿.
一來猶太人很注重婚姻.
另外這個舉動是犯了猶太人認為不能夠犯的家族的亂倫.
所以施洗約翰就公開去指責.
指責的本質不單單說他是淫蕩或者不道德.
而是說希律安提柏王是不配做猶太人的王.
這個說法就很不得了.
因為施洗約翰作為一個民間的領袖.
他們想顛覆王權.
指責希律王或者去詆毀他.
當然不單是王很憎恨他.
皇后希羅底也視他為眼中釘.
因為約翰在民間的勢力越來越大.
所以就將他收監.

$^{41}$但怎知道到某一天.
就是希律王生日他舉辦了一個生日會.
福音書怎麼記載呢.
公主薩諾米在眾人面前跳舞.
可能她跳得很好或者艷麗非常風采.
大得王的歡心.
開心之際希律王竟然發誓.
說這個公主求什麼他就給什麼她.
皇后就去叫女兒施洗約翰的人頭.
最毒就是婦人心.
馬達福音第十四章有這樣記載.
希律王最後放棄.
打發人在監獄裡斬了約翰.
將約翰的頭放在盤子上拿給女兒.
女兒就拿給皇后母親.
施洗約翰的門徒就將屍首埋葬.
然後就告訴耶穌.
約翰被處斬這件事不單單是聖經有提到.
第一世紀的猶太歷史學家約瑟夫也有提過這件事.
他說約翰被處斬的地方是在約旦河東的希律王王宮堡壘.
那個地方叫做馬基魯.
這個王宮位於今天的約旦王國.
在死海東邊的一個山頭上.
今天你去的那個地方你也可以看到馬基魯的遺跡.
耶穌回到加里尼接續施洗約翰的運動.
既然他已經下監.
那群眾就需要有一位領袖出來.
所以耶穌就上了去.
路加方第四章十四節.
耶穌滿了聖靈的能力回到加里尼.
他的名聲就全遍四方.
他就在加里尼各個會堂教訓人.
眾人就稱讚他.
加里尼當然有很多猶太人的城鎮鄉村都有會堂.
當然其中一個有猶太人會堂的地方就是耶穌長大的拿撒勒城.
我們來過拿撒勒村.
我們說這裡有一個猶太人的會堂重建了出來.
仿照一世紀猶太人會堂格式重建.
你在門口進去的時候.
你就會發現這個會堂有羅馬式長方形的大廳.

$^{81}$這些會堂是猶太人社群的中心.
門口很多時候是面向耶路撒冷的.
大堂周圍三面有長石凳圍著.
人們坐下來聽道聽律法書等等.
中間有一張小的桌子或是一個講台或是一個座位.
是給主席用的.
當時的聖經是寫在被卷上的.
所以讀經的人就會捲開書卷然後供讀出來.
這個就很寫實.
路加福音四章這樣記載.
耶穌來到拿撒勒自己長大的地方.
照著習慣進入會堂.
在安石的時候他就起來讀經.
有人把以賽亞先知的書卷遞給他.
他剪開了書卷找到一個地方.
上面就記載了什麼呢.
他說主的靈在我身上.
因為祂告我去傳福音給貧窮人.
差遣我去宣告被擄的得食房.
黑眼的得漢奸.
受壓抑的得自由.
又宣告主月立的希年.
他就把書卷捲好.
然後交給侍從.
坐下來.
然後眾會堂的人就注視他.
他就跟他們說.
這一段聖經今天仍然在你們中間.
就是說他就是那個救主以賽亞.
來到成就以賽亞眾先知的應許.
救主到了.
耶穌打開的地方就是以賽亞書61章58章等等.
那裡應許受苦的猶太人神的子民.
終會有救恩來到.
四章十九節這裡也說.
主月立的希年.
這個jubilee.
這裡引用利未記二十五章的jubilee.
律法書說到什麼呢.
就是說到一個社會社群的願景.

$^{121}$一個vision.
就是每逢五十年的時候就要取消債項.
這些是上帝很喜歡看到的事情.
是上帝所月立所接納的事情.
所以我們看耶穌這一段拿撒勒的講道.
我們就不要這麼快將它私有化.
這裡其實是在講律法書裡面.
是去啟示出一個新的社群.
新的國度.
那個奴僕是得釋放受壓的能夠得到自由.
那個社會是有一個希年來到.
就是貧者不會越貧.
富者不會越富.
是一個均勻了的社會.
一個按照上帝的律法來行出來的一個社群.
其實這個就是希望教會能夠做到的.
這些拿撒勒人他們聽見耶穌所講的事情.
他們就不接納上帝所月立的.
這個很諷刺.
這個情節大家可能很熟悉.
路加福音第四章二十八節.
會堂裡面的人聽到就很生氣.
然後趕耶穌出城.
帶他在山上拉他到山崖.
想推他下去.
耶穌就在他們中間走了.
離開了.
當地的傳統.
就是將這個事蹟放在拿撒勒南面一個山崖上.
這個是唯一一個山崖.
這裡叫做跳崖山.
我們大概可以想到當時會堂裡面的人.
可能都是有權有勢.
他們覺得這個神國要來到地上.
又說要有欺憐等等.
其實是在威脅他們自己的地位.
或者不想便宜別人.
不想得到改變.
我們會不會是這樣呢?.
神國來到的時候.

$^{161}$不只是我們罪得赦免就罷了.
我們有永生就算了.
而是需要我們在生活上作改變.
我們作為基督的門徒.
我們不可以只求保持原狀.
但又希望天國來到地上.
我們又不想慷慨.
我們又不想祝福有需要的人.
又不想行公義.
又不想關心身邊的人.
我們只顧自己私有化的得到祝福.
這個不是天國的意思.
所以這裡告訴我們.
到底我們是哪一班人.
到底我們是接受主.
還是拒絕耶穌.
我們跟著耶穌的腳踏.
我們離開了拿撒勒來到加伯隆.
這個地方是屬於天主教聖方濟各所管的.
Franciscans.
這裡一世紀有很多遺跡.
來到門口的時候有個牌很清楚.
Capernaum, hometown of Jesus.
加伯隆,耶穌的城鎮.
我們想一想.
不是拿撒勒才是他的城鎮嗎?.
因為他長大花了那麼多年.
聖經有沒有說.
這個真是耶穌的家.
耶穌的城鎮.
馬太福音第九章第一節.
耶穌從另一個地方坐船.
回來到自己的城.
His own town.
His own city.
加伯隆這個字.
英文是Capernaum.
是兩個字組成的.
村和安慰的意思.
安慰之村.

$^{201}$安慰的意思.
你記得舊約有一本書叫拿紅書.
那個名字其實是安慰的意思.
我們從上面看加利利湖的時候.
我們發現其實有很多漁港.
在當時第一世紀.
至少有15個漁港碼頭.
圍著湖的每一個海岸.
其中一個很出名的就是莫大拉.
Madeira.
你記得莫大拉馬尼亞就是從那裡出來.
加伯隆就是那麼多漁港裡面.
最大的漁港.
考古的人員去到古代.
當時加伯隆的人口.
大概千多人.
雖然對我們今天現代的尺度來說比較小.
但是一個很繁盛的小漁村.
來到加伯隆的時候.
有兩樣東西值得看.
第一就是進到左邊.
有一個白色的加伯隆會堂.
是用白色很漂亮的石頭來造的.
我們要講一講石頭.
這裡附近的地質.
是以色列北面一帶.
本地的石頭全部都是黑色的.
這些黑色叫做Besot.
玄武岩.
以前很遠古時代留下來的火山岩.
都算是很硬的.
是不差的建築材料.
所以大部分的房屋.
都是由這些本地的黑色的Besot來造的.
你看到這個會堂是用白灰石.
White limestone去造的.
你知道是一項很大的投資.
因為要從很遠的地方.
搬這些石塊來建築這麼大的會堂.
這個會堂入口有三個大門.

$^{241}$進到大廳.
大門向南.
向著耶路撒冷的方向.
大堂三邊有石凳.
你看到有些很高的.
哥林多式的石柱.
頂著上面二樓.
兩邊都有另一個大的房間.
所以容納了很多人在這裡.
馬科科一章二十一節.
他們來到加伯隆.
耶穌隨即在安息日進入會堂教導人.
這個會堂會不會就是耶穌在加伯隆教導人的會堂呢?.
不過大家都知道這個會堂.
其實是建於公元後四五世紀.
所以一定不是耶穌當時建的會堂.
不過我們來到聖地做旅客的時候.
我們要學習一件事.
每個地點都是一個考古的土幽.
即是說是一種人造的山幽.
怎樣人造的山幽呢?.
因為人類隨著不同的世代.
隨著不同的年日.
不同的人在上面不斷來建築.
來拆毀.
然後再重建.
又再拆毀.
所以一層一層的土層的遺跡.
疊在一起.
好像一個蛋糕一層一層.
我們來到這個地方的時候.
我們需要知道其實要揭去下面.
到底下面發生了什麼事呢?.
所以考古人員來到這裡的時候.
沒錯這個白色的會堂是四世紀末或五世紀所建的.
但是他們發現下面有一個一世紀的石的地基.
而這個一世紀的地基是由黑色的玄武岩.
這些Bazooka來建的.
所以大概這個地方.
同樣都是一世紀的會堂地方所在.

$^{281}$雖然只剩下一個地基.
但這個就是加伯隆當時一世紀的猶太人會堂.
加伯隆第二樣值得去看的就是離這個會堂不遠的.
只是隔一條街.
就好像這個現代教堂.
好像太空船一樣.
他蓋著一個八角形的遺跡.
這個八角形的遺跡.
就屬於四五世紀的拜占庭時期的古教堂.
當聖方濟各的人來到這裡買了這個地物業的時候.
他們開始發掘的時候.
他們發現這個四世紀的古教堂之下.
就有更加早期的遺跡.
這裡特別的地方就是.
這個地點是經過好幾個階段的加建.
最初的階段就是一世紀的上半葉.
這裡只是一個很簡單的房屋.
有兩個天井.
以前的房屋就是這樣的.
一個家庭貼著另一個家庭.
每個家庭就圍著一個天井一個階層.
我們在地圖上看到.
我們看到有一個大房.
但到了一世紀的下半葉.
這個大房就成為了一個家庭教會聚會的房間.
然後進入第四世紀的時候.
這個大房兩邊就加了兩條柱.
頂著一個房頂.
一個拱門一個大門.
房屋四周就加了圍牆.
牆上又引上了白灰.
很罕有的.
因為這裡沒有一個建築物是這樣的.
牆上刻了很多的字.
就是一些基督教的字句字眼.
譬如基督等等.
很可能是早期的朝聖者.
來到這個地方.
觀察到一個教堂.
他說這是彼得之家.

$^{321}$變成了教堂.
到了第五世紀.
這個地方就蓋上了三角八角形的建築.
當時拜占庭時期的教堂特色.
就是用一個八角形來標籤.
某些耶穌的事蹟或聖經的事蹟.
所以從早期的文件或基督教的傳統下來.
這個很可能就是彼得的家.
馬科音第一章二十九節.
耶穌和門徒一出會堂.
就從雅各和約翰.
到了西門和安德烈的家裡面.
就是用會堂一出去.
他們就到了會堂.
這個值得留意的.
就是會堂和彼得的家是非常近的.
彼得和安德烈是兩兄弟.
兩個家庭.
你留意那個房屋.
最初的時候是有兩個天井.
三十節西門的惡母正在發燒躺著.
他們立刻告訴耶穌.
耶穌走到他們面前.
他拉著他的手扶起他們.
熱就退了.
他就服侍他們.
到了黃昏的時候.
不斷有人帶有病或被鬼附的人.
帶到耶穌面前.
全城的人聚集在門外.
耶穌就醫好他們各樣的病.
趕出很多的鬼.
並且不許鬼說話.
因為鬼認識他們.
如果進入天主教的教堂.
這是一個很特別的設計.
因為中間是玻璃.
讓你在上面可以看下去.
所謂彼得的家.
如果留意到兩邊是有拱門的柱.

$^{361}$加建了一個大的屋頂.
把這個地方本身是家庭房間.
成為教堂聚會的地方.
為什麼耶穌會選擇加白龍.
為祂的傳道基地呢?.
馬太福音第四章十二節.
耶穌聽見施主約翰下了監就退到加利利.
後又離開拿撒拿往加白龍去.
住在那裡?.
那地靠海在西布倫拿弗塔利的邊界上.
這是要應驗才知道以賽亞所說的.
西布倫地拿弗塔利.
就是沿海的路約旦河外.
外邦人的加利利地.
那在黑暗裡面的百姓看見了大光.
在死刃之地的有光發現照著他們.
從那時開始耶穌就傳到這裡來.
天國近了你們應當悔改.
如果馬太要依書直說的話.
直接交代就是他來到加白龍.
耶穌就開始傳道就行了.
為什麼要加中間那一節呢?.
十四至十六節.
其實就是引用以賽亞書第九章.
這裡提到西布倫地和拿弗塔利地.
是黑暗的地方.
以色列人在舊約分地的時候.
兩個支派西布倫和拿弗塔利.
分了以色列北邊的加利利附近的地方.
為什麼說他們是黑暗的地方呢?.
因為在舊約的以色列歷史裡面.
這個地方是第一次被外邦人.
亞述帝國入侵的地方.
就是西布倫和拿弗塔利.
當時是公元前732年.
有個亞述王叫做提拉帕利茲.
他入侵北國以色列.
以色列人第一次大敗.
很多人在那裡喪生.
並且很多人流亡到亞述.

$^{401}$這是舊約歷史裡面第一個黑夜.
亞述帝國不單止將以色列人流亡到第二個地方.
也將其他戰敗的國的外邦人.
搬到加利利這個地方.
為什麼呢?經文經常提到.
加利利是外邦人的加利利.
是外邦人來到或是外邦人踐踏的地方.
加利利就是這個地方.
現在天國終於來到了.
按照以賽亞的認許.
耶穌到了,天國近了.
他們得到安慰.
受苦的百姓終於見到光.
所以在加伯隆這個地方.
耶穌傳道之外.
他也做了很多安慰的工作.
經文也提到沿海的路.
這個字其實指當時一個很出名的商貿路線.
叫做Villamares.
其實是貫穿南北的一條大道.
北面是敘利亞和中東米索不達米亞.
南面是直接通往埃及的路線.
中間是經過以色列整塊的土地.
特別是沿海平原.
這是一世紀的交通命脈.
所以耶穌從拿撒勒搬到加伯隆.
其實是增加了他的能見度和策略.
這是一個國際性的地方.
所以耶穌開始傳道一兩天.
聖經說幾十里外的其他城鄉的人.
就會聽聞他的消息.
這個Villamares是一個國際的路線.
軍隊,商貿是必經之路.
我們知道古代的路.
由一個大城市連接另一個大城市.
東面.
我們看到的是白菜大這個地方.
Bazeta.
這是希律費力的省府.
左邊是提比利亞Tiberias.

$^{441}$這是加利利省希律安提柏的第二個省府.
之前是在齊波利的Sephiroth.
我們去過.
城市與城市之間如何連接呢?.
由白菜大第二個省府連接提比利亞加利利的省府.
中間就是加伯隆.
這是一個國際的交匯點.
我們了解了福音書的一些情節.
為什麼會有羅馬伯夫長來求耶穌醫好他的僕人呢?.
伯夫長就是管理100個羅馬士兵.
所以他在加伯隆這個地方.
就代表羅馬帝國的利益所在.
這裡是有軍隊駐守的.
是軍防重要的地方.
另外,你記得耶穌的12個門徒.
其中一個叫馬太或利美.
他是做哪一行的呢?.
他是做稅理,收稅.
即是說這裡是有關稅的辦公室.
是羅馬帝國的利益.
一個省去到另一個省需要收關稅.
並且在這裡捕捉到的魚或做生意商貿的收入.
就需要交稅.
這裡是國際的會點,是城市的會點.
最好做什麼呢?就是傳訊息.
所以你記起耶穌開始傳道的時候.
很多的教導都是講道.
你記起其中一個最出名的叫登山寶訓.
在這裡附近山上有一間教堂叫伯福堂.
紀念耶穌在這裡駐名的講章.
登山寶訓.
下次我們有機會去.
登山寶訓的內容很多都是家傳戶曉的句子.
例如愛你的仇敵,多走一里路,walk the extra mile等等.
你記得登山寶訓的主旨是講什麼呢?.
就是門徒的生活,我們基督徒的生活.
就是要勝過宗教的領袖.
民事法是錯人的那些假冒為善的人.
另外門徒的生活應該要勝過外邦人.
跟外邦人有分別.

$^{481}$五章47節.
你若單請你弟兄的安,比人有什麼長處?.
就是外邦人也不是者,是吧?.
所以你要完備這麼多,像你的天父完備一樣.
愛你的仇敵.
讀章第七節.
你們禱告的時候不可像外邦人.
用很多重複的話.
我們在這個地方的生活應該跟其他人有分別.
我們的價值跟其他人有分別.
我們所繼續的事.
首先要跟其他人有分別.
跟猶太人不同,跟外邦人也不同.
加利利湖分兩邊.
西邊就是耶穌傳道,門徒傳道.
很多時候集中的地方.
就是猶太人的地方,加帕隆等地方.
東邊就是外邦人住的地方.
其中一個城市叫格拉森.
你記得講群鬼的地方就是那個地方.
所以這裡就是夾著外邦人和猶太人的地方.
我們在這中間,門徒的生活應該要分別出來.
這裡我想起登山寶範其中一段.
六章二十五節.
所以我告訴你們.
耶穌說不要為生命憂慮吃什麼喝什麼.
為身體憂慮穿什麼.
生命不勝於飲食嗎?身體不勝於衣裳嗎?.
也不積蓄在倉內.
你們的天父尚且養活他.
你們不比飛鳥貴重得多嗎?.
你們哪一個能用思慮洗壽數增加一刻呢?.
何必為衣裳憂慮呢?.
你想野地裡的百合花長起來.
它又不勞苦又不枉善.
然而我告訴你們,所有門極榮華的時候.
它所穿戴的,還不如這花一對.
你們小信的人野地裡的草今日還在.
明天就丟到路裡.
神還給他有這樣的裝飾.

$^{521}$何妨你們呢?.
所以不要憂慮說吃什麼喝什麼穿什麼.
這都是外邦人所求的.
我們是有不同的.
因為你們知道你們的需要.
你們的一切,你們的天父是知道的.
但我們應該做些什麼呢?.
三十三節,你要先求他的國和他的義.
這一切東西都要加給你.
所以不要為明天憂慮.
因為明天才有明天的憂慮.
一天的難處,一天當頭夠鳥.
但在這裡為什麼會提到天上的飛鳥呢?.
野地的草我們上一集看過.
因為加利利這個地方不單止是人類的交會點.
外邦人,羅馬人,猶太人等等.
加利利也是一個大自然的交會點.
因為每逢每年那些季候鳥南北遷移的時候.
它就會經過以色列這個地方.
而集中的地方就是加利利這個地方.
按估計每年這裡有五億乘千上萬的飛鳥經過這個地方.
舊約仙子耶利米也指責過猶太人.
他說連空中的灌鳥也知道自己來去的季候.
班寇,燕子,伯鶴也遵守遷移的時令.
他知道那個時機.
他就遵守,很隨意地飛過這個地方.
他說我的子民卻不曉得.
耶和華的法則.
換句話說.
這些子民和上帝立約.
知道上帝是信實的時候.
居然他們不曉得神的律法.
不去忠於上帝的說話.
耶穌在登山寶訓裡面叫門徒.
看看天上的飛鳥.
你們很熟悉的.
每年季候鳥經過的時候.
你們就知道你們應該憂慮的第一件事.
不是跟外邦世界的人一樣.
我們先求的是什麼?.

$^{561}$Priority.
就是上帝的國.
上帝的義.
在我們生命當中.
在我們世界裡面.
我們可以行出上帝的心意.
聽上帝的說話.
這是我們門徒第一個的價值.
我們所憂的,所求的.
應該是上帝的國度.
能夠來到我們的生命當中.
來到我們社會.
來到這個世界.
上帝今天要求你們做什麼?.
你是不是將祂放在第一位?.
祂對我們的心意.
是不是放在你的優先次序上?.
又或者你想想.
我們花多少思緒,花多少心情.
花多少時間.
是在世人所關心的東西.
還是我們花的精神.
是在上帝的國度.
上帝要人聽見福音.
上帝要教會被建立.
上帝要公義,要真理.
在社會上彰顯出來.
所以今天最後有一個這樣的邀請.
就是我們先求上帝的國.
和神的義.
我們結束後一起祈禱.
上帝我們求你幫助我們.
讓我們先求你的國和你的義.
在這個世界當中.
特別是在疫情裡面.
很多時候我們所求,所想的.
與世人沒有分別.
主一求你幫助我們.
讓我們站在國際的地方.
讓我們在世人當中.

$^{601}$我們先求你的國,你的義.
我們先求你所關心的.
你的福音,你的公義,你的真理.
能夠彰顯出來.
求你幫助我們為你的國度發聲.
求你幫助我們.
我們的行為,我們所行,所想.
是去行出你的公義.
說出你的公義的說話.
我們這樣祈禱,奉主名求.
《情歌》.
情好飛了.
百修不蹤.
你田枯雙赤養活他.
野地青草.
田園百合.
露水半段站.
他都悉心裝成.
難道你不悲.
了易珍貴.
受累那可是你的歲月變多.
難道你一生.
就是這樣憂慮.
吃什麼喝什麼穿什麼.
你們要先祈求他的國與他的義.
你所需的都要加給你.
明日只有明日的憂慮放雨.
一天的難處一天都可救了.
你們要先祈求他的國與他的義.
你所需的都要加給你.
明日只有明日的憂慮放雨.
一天的難處一天都可救了.
難道你不悲.
了易珍貴.
受累那可是你的歲月變多.
難道你一生.
就是這樣憂慮.
吃什麼喝什麼穿什麼.
你們要先祈求他的國與他的義.
你所需的都要加給你.

$^{641}$明日只有明日的憂慮放雨.
一天的難處一天都可救了.
你們要先祈求他的國與他的義.
你所需的都要加給你.
明日只有明日的憂慮放雨.
一天的難處一天都可救了.
明日只有明日的憂慮放雨.
一天的難處一天都可救了.
「明日之憂慮」字幕.
\newpage



\section{}
\label{sec:aXI0__JmVi0}
\textbf{和你約定 [1] - 傾o下國安法}
\newline
\newline
連結: \href{https://youtube.com/watch?v=aXI0__JmVi0}{\texttt{https://youtube.com/watch?v=aXI0\_\_JmVi0}} ~~~~ 語音日期: 2020-09-03
\newline
\newline
\hyperref[sec:c1irNOxf2YQ]{\small{< < < PREV SERMON < < <}}
~
\hyperref[sec:index_chronic]{\small{[返順時目]}}
~
\hyperref[sec:index_scriptual]{\small{[返順卷目]}}
~
\hyperref[sec:u3NKvcUYcmI]{\small{> > > NEXT SERMON > > >}}
\newline
\newline
$^{1}$歡迎所有觀眾來與我們聊天和思考信仰和現實生活的問題.
大家知道2020年6月30日,香港國安法生效.
我們作為基督徒應該如何思考和看待這個問題.
聖經有沒有一個角度讓我們看到政治社會的問題.
神在聖經裡面應該有什麼原則讓我們使用.
所以今天我找了Anthony,他是在Vancouver的.
他在文革會議中,在溫哥華的短訊中心.
Helen和Selwyn在多倫多的,他們在盛會中在教會中計劃.
我姐姐在芝加哥.
我們幾個走在一起,我們想談這些議題.
其實這些議題是很黑的議題.
我們發現我們不能避開.
我覺得上帝對我們很好,用不同的衝擊讓我們思考.
我覺得在國安法的時候,在橫行教會的圈子.
香港當然非常敏感,北美其實也是很敏感的.
我想探討一下這個問題,為什麼.
為什麼你覺得在你的經驗和觀察裡面.
在北美的教會,這個比較遠一點,隔了一個太平洋.
為什麼你覺得這裡是這麼敏感呢?.
其實當我在牧會的時間,其實也接觸過很多.
香港的社會運動近年,都是由2014年比較開始熾熱.
雨傘運動開始.
其實當時的議題,當然當時都是一個很龐大的運動.
但是在北美來說,那個反而沒有那麼靠近.
但是當時去反應的時候,我聽到基本上是一個聲音.
從教會出來的,就是我們有國內的士工.
我們有宣教的士工.
所以我們不出聲,我們不想影響我們國內的士工.
所以這個理由是當時一個,也算是一個很可取.
也算是一個很合理的理由.
但是一直發展到去年,去到今年的時候.
你會發覺沒有人再用這個理由.
因為十字架又被拆,教會又被封.
所以在這個處境裡面,你不可以再說是因為我.
為了宣教,為了福音的緣故而不出聲.
甚至應該為了福音的緣故而出聲.
但是我想更加浮現到教會那個實際的問題.
就是我們其實不是很知道怎樣跟政權去抗衡.
而是怎樣去對付政權.
特別我們香港過來的,我們過去在殖民地的地方長大.

$^{41}$我們是受惠於英國一個基督教的國家.
所以我們去面對一個不是基督教的政權的時候.
我們是不知道怎樣面對的.
而我們有恐懼,最好我們不要出聲.
所以我想這個是浮現了移民的教會,香港的教會.
華人的教會的核心問題.
面對政權的時候.
所以很多時候,一扔扔出來的那一句就是.
政教要分離的.
但是很多時候他們所說的政教分離.
就是在說政治和宗教.
但是其實真正的政教分離應該是政權和教會.
或者是State and Church.
和Religion and Country.
那個是我一些觀察的內容.
說到政教分離的時候.
我想原意不是說教會不理政權.
而是政權沒有任何權力應該去牽涉教會.
去控制教會.
所以政教分離不是我們不講政治就叫政教分離.
政教分離是保護宗教自由和自由.
所以我想都是來自不同的理解.
我們每次說到政治的時候.
我們覺得和信仰是沒有關係的.
信仰應該是一些很熟悉的東西.
甚至一個移民教會.
我想更加一個本身.
就算你是Miguelito也好.
其實我想參與在主流社會的程度相對是低的.
很多時候,特別是第一代移民來的時候.
住在一個安穩的生活.
其實已經很難了.
或者說收成期.
要享受這裡的.
有的你那麼久以來得來的成果.
可能對社會的參與.
無論是因為語言或文化的不同.
其實那個參與性是相對低的.
變相在教會.
除了一兩件的事件之外.

$^{81}$在加拿大或美國.
我相信都是.
那個參與.
教會或信徒參與的政治.
普遍是少的.
我想這個也很相關.
剛才Andy 說的.
我們在香港移民過來的時候.
回想我小時候在香港成長.
除了1989年的時候.
上街之後.
其實我也想不到有很多大眾一起去參加一個政治運動.
在香港,我從小到大都聽到.
我對政治沒有興趣.
與我無關的.
總之我不搞了.
我想這件事也很深入.
我們這些香港人.
甚至移民到海外.
一樣都是這樣.
突然間我想起.
在這幾年,我們在2016年回到加拿大.
曾經我們對加拿大一些越來越有爭議的議題.
現在都流行一些petition.
我們都在教會的WhatsApp群組.
小組.
提出了一些.
當時傳統是.
首先是這樣說的.
多謝你去參與政治.
不過這些petition你不要再發出來了.
我現在要叫Elder Sport.
要看看我們可不可以再發出這些東西.
當然.
他說完之後就沒有再跟進.
但我想反映了.
其實橫教會對.
這個我剛才也說過.
不是一些敏感的.
好像國安法.

$^{121}$是在現場的一些政治議題.
我們的取態.
在教會的取態.
都是我們不要再說了.
不要搞了.
我就叫Elder Sport.
不過其實都是不了了之.
很多時候.
我們都要.
不同方面的取態和意見.
我想拿出來談.
其實是一個好的事.
但是我想.
缺乏了這樣的空間和平台的時候.
變了就會.
不要了,最後大家可能吵架收場.
面左左.
我想這個是.
中國人都是喜歡和諧.
大家不要互相和氣.
和諧就最好了.
我想這個都反映在裡面.
何諧.
何何諧諧.
但我從另一個角度.
就是其實對於政治.
這件事來說.
其實對於中國人來說.
我們是沒有政治的觀念.
我們在儒家的思想.
一直流傳下來的時候.
我們是沒有政治的觀念.
當我們接觸政治的時候.
其實我們已經.
被灌輸到政治是骯髒的.
是污穢的.
這個是我們.
從小到大想政治的時候.
就想到的,污穢,骯髒.
這個再加上.

$^{161}$我們在教會裡面.
現在我是一個訓練人.
做報道宣教的機構.
我都很多時候研究.
我們上一代.
其實傳福音的方法.
就是很多時候信耶穌得永生.
香港我們.
經常看到十字架都是這幾個字.
信耶穌得永生.
信耶穌最重要的是什麼?就是上天堂.
這種概念來說.
變成我們一個很出世.
而不入世的信仰.
信耶穌.
就是為了我將來死了.
上天堂.
我今生拿了這個護照.
That's it, OK.
我不會再思想.
我怎樣回應.
聖經裡面所說的作肝作炎.
是在時代裡面.
去發聲.
做福音的見證.
我們所想的福音是什麼呢?.
就是.
confession prayer.
你say confession prayer.
信耶穌,信主,決志.
很重要,因為你講了那個.
你就上到天堂.
但今次的生命是怎樣呢?.
是不會理的.
這種影響我們就是性俗義分.
我們就以為.
這就叫做聖潔.
在聖經裡面的聖潔.
耶穌所說的,你要跟從.
那種像Christ like的生命.

$^{201}$我們沒有什麼追求的.
我們不是要Christ like的生命.
我們是要.
In Christ,一個的eternal life.
將來的那個最重要.
所以這個是我們.
我明白.
上一代的時間.
傳福音的環境是很不同.
香港或其他.
東南亞的地方的貧窮.
或物質上的需要.
那種情況的底下.
重點我們就說.
我們信耶穌的一個永恆的盼望.
那個是對的.
但那個其實是某一部分.
我們一半講了福音.
但我們沒有講到今生的部分.
所以這個其實是影響了.
我們怎樣回應政治的議題.
還有我們很怕.
我們不懂得去碰這些東西.
某程度上.
我們有屬靈潔癖.
我們有潔癖.
這些全部都不要碰.
變成我們不像耶穌.
我自己的角度來說.
耶穌不是這樣的.
祂有潔癖就不會來拯救我們這些罪人.
祂最辛苦的就是.
祂聖潔的.
祂要來到我們這裡.
沾污了祂.
但其實我們的信仰就是要我們.
在祈禱的聖餵裡.
去改變這個世界.
我有一個相關的.
自我問自己的問題.

$^{241}$如果是這樣的.
如果信仰是這樣.
為什麼耶穌.
或者上帝.
在我們信耶穌的時候.
祂不直接.
這樣就搞定了.
做一步做少了一步.
我才是最完美的.
我小時候在教會長大.
我真的想過這個問題.
我信了主.
為什麼還不放我走?.
其實最想的就是不用考試.
不用讀書.
生老病癒.
哈哈哈.
考試.
去讀書這麼辛苦.
帶我走,我相信了你.
天堂有永生,天堂有多好.
天堂又不用讀書.
我們的信仰在某程度上.
是這樣的.
我聽到.
如果我們談一些社會政治的議題.
很多時候.
首先我們對舊恩.
對神的國.
的理解已經重新重新.
另外我聽到一個很實際的.
關鍵就是.
我們想和和諧諧.
不想吵架.
一談政治很多時候就吵架.
我就想下一半我們可以.
談一談.
我們用什麼角度.
回到聖經裡面.
以至聖經是提醒我們.

$^{281}$告訴我們我們應該怎樣處理.
這些問題,怎樣處理.
怎樣處理政權.
其實時間.
對談判.
其實都重要.
去到現在很多事情已經很白熱化,很兩極化的時候.
怎樣可以仍然製造.
其實你要去對話的.
是要有一個安全的空間.
要大家願意.
放下身段也好,放下成見也好.
放下真相也好.
坦誠地一起去.
seek for the truth.
但這個seeking for the truth.
的空間.
甚至在教會裡面.
對政治的議題.
究竟存不存在呢.
這個是.
我覺得是一個問題.
一般最實際來說.
有時候教會都很講實際.
實際就不要講.
我給一個例子.
不關國安法的事.
我想反映到其實.
教會怎樣去處理conflict.
或者處理controversial.
我不開名這間教會.
在加拿大有一間教會的一個會眾.
他就贏了.
香港叫六合彩,這裡叫六四九.
我不是在說30元.
我是在說以千萬計.
一注中.
在教會裡面.
當然會慢慢揚了一些東西.
但這個中了.

$^{321}$你們可以想像一下.
在會眾當中就會問.
為什麼他會買呢.
教會怎樣去面對呢.
教會是否深受他貢獻呢.
最重要是11貢獻給教會.
我也是這樣說.
不要理他.
你有你的政治取態.
之後再談.
先收起來.
教會其實是義淫紛紛.
我記得我星期日.
在專用私手上課.
差不多是.
他確認了這件事.
這件事是真的在發生.
剛剛就中了.
不過我們以後不要再說這個話題.
哈哈哈哈.
不要再說.
有聲音.
對不起.
整個我估計他已經解釋了一分鐘.
兩分鐘.
是完全沒有.
當時我聽完之後.
我們需要去解釋.
不是去判斷那個人的行為.
而是你怎樣去教我們.
面對這件事.
其實華人的圈子.
華人的文化就是.
有什麼就掩蓋.
不出聲就當沒事.
但我們一直沒有去解釋.
去到國安法的問題.
我們不懂得用什麼態度.
去處理這個.
這麼大的問題.

$^{361}$我們缺乏了一些工具.
或者其實也有.
有時可能拿羅馬書.
十三章.
經文回答.
有沒有機會可以重新再去看.
甚至其實很簡單.
例如這個.
我拿.
馮任坤的.
《食敬書》.
其實他有提到.
你是可以解釋到.
無論是違反也好.
起碼拿出來看.
有些金句.
回答了.
沒有什麼好說的.
討論就只能停在那裡.
拿羅馬書十三章.
其實不是想去釋經.
是想拿羅馬書出來.
我告訴大家不要說.
總結了一個的方法.
其實俗一點來說.
教會.
華人教會.
我們慣常用一個方法.
就是大石仔殺死蟹.
河蟹不動的時候.
我們就大石仔殺死蟹.
這樣就解決了問題.
蟹是最好的.
蟹不可以再動.
蟹不動.
石頭就壓住了牠.
這樣就沒事了.
但其實是傷害.
我看到.
我們自己做牧者.

$^{401}$我們也要去檢視.
我願不願意.
在電影展會面前.
有些大議題我真的不知道.
我承認我真的不知道.
我不需要壓下去.
我們不如一起討論.
或者牧者與牧者之間.
我們大家可以彼此多些交流.
彼此分享.
這個議題.
我們怎樣去處理呢?.
我們不需要.
統一地去處理.
我們不是天主教會.
天主教會也不是統一.
我們不需要統一地去處理.
但我真的不懂.
或者Mark.
我們談談吧.
但我們連談也不談.
我很喜歡.
我們的聖禮.
無論是甚麼宗派.
我們一定有聖餐.
你們怎樣稱呼也好.
我自己的宗派是Communion.
Communion的意思.
就是融合.
我們要Communicate.
我們就要融合.
這個在世界是不可能的.
是因為耶穌基督藉著祂的身體.
藉著我們的福音.
能夠將我們融合在一起.
而這個我覺得很奇妙.
主耶穌當臨上十字架之前.
祂很緊張地為大祭司祈禱.
為祂眾門祈禱.
祂不斷說我在你裡面.

$^{441}$你在我裡面.
我年紀小一點的時候.
我經常想這個問題.
我怎樣在你裡面,你在我裡面呢?.
可以達至的.
唯有的,我們將我們自己完全的破碎.
甚至將固體變成液體.
我才可以我在你裡面.
你在我裡面.
但如果沒有這個願意破碎.
願意放下自己的態度的時候.
我很難我在你裡面,你在我裡面.
我們就大家拼了.
拼不成就拼在石頭裡.
通常都是主任牧師.
醫心牧師.
就會摘下來就解決了問題.
但其實我們到今天.
我們知道問題是不會解決的.
我們可不可以從另一個角度去看看.
國安法的影響.
我想在香港你見到.
有很多牧者因為這樣緣故.
被移動辭職也好.
被辭職也好.
即是.
我剛剛看完這個《士途行傳》.
可能被提到去別的城市也好.
其實是帶來很多的恐懼.
也有時可能很介懷.
會不會變成要勾結外國勢力.
這一點對香港教會.
特別牧者來說.
其實影響可以很近身.
但其實北美呢.
北美這方面牧者會不會都有類同的感受.
或者恐懼呢.
我想絕對會.
例如在這裡.
我們認識的一些牧者.

$^{481}$他們說國安法不是我不想講的.
是不想回香港的.
你拿條口紅.
很實際的感覺.
因為你看現在通緝的人.
可能講了一些東西.
或者支持某一些東西的時候.
他就會被通緝.
這個國安法是伸延了全地球.
所以我們在這裡講的東西.
全宇宙.
全世界.
全宇宙.
都有這個恐懼.
有這個負擔.
有這個重擔.
是很真實的.
以往你不會覺得.
有什麼在香港那邊會影響到我們在加拿大那邊.
現在你是很近身.
自從我們從來沒有想過.
原來一條法律在中國那邊立了.
或者在香港那邊立了.
如果不立了我這邊原來會被人抓的.
還要是我們的政府不知道它可不可以保護到我們.
以前我們想著.
我回中國拿著加拿大護照回去.
回鄉證的話就可以保護到我.
現在是不行的.
這些的確.
如果從一個牧者的角度.
他們其實是擔心的.
不只是他們自己的安全.
事工怎樣去做.
其實是整個新常態.
教會是新常態.
差傳是新常態.
我們屬靈的生命都是新常態.
不過在這個新常態裡.
究竟上帝要我們怎樣在裡面去突破.

$^{521}$或許我們過去那些不是常態.
我們以往那些可能我們是走歪了的歪態.
現在上帝是將我們撥亂反正.
回到一些常態.
怎樣去做教會.
信徒怎樣去做.
怎樣做一個有血有肉.
造成肉身的信徒.
好像耶穌基督的.
不只是說的.
來到這個地上能夠去回應到各方面的需要.
其實不只是政治議題上的一個問題.
而是其實很多人在這段時間裡.
很多無論是情緒上.
經濟上.
或者是.
不只是說政治的.
其實全人都有很多的需要.
而這個holistic gospel.
我們怎樣能夠帶到這個generation的人.
上帝其實給了我們一個機會.
今天想說最後一個問題.
就是聖經神學給我們有甚麼工具.
或是角度.
因為我們有很多不同的.
不同宗派有不同的神學背景.
不同的持守.
但其實當我們在說.
我們的神學是apply.
在生活上.
我們不是像在象牙塔裡的神學思想.
或者是神學的鑽研時.
我們就要明白.
其實我們有甚麼神學.
就是theology.
就會出現甚麼教會論.
就是theology.
甚麼教會論.
就會出現甚麼宣教學.
就是theology.

$^{561}$所以其實我們今天怎樣去看.
福音究竟是甚麼.
我們怎樣去回應.
究竟是甚麼.
我們怎樣去活出福音.
是messiology的內容.
是影響我們怎樣成為一個教會.
如果我們今天看教會.
只是一個shopping mall.
一個consumer聚集的地方.
我們出來的messiology出了問題.
而同時我們的神學都出了問題.
究竟我們信的一個甚麼神.
所以我們今天怎樣活出.
我們無論是國安法.
無論是其他社會的議題.
我們怎樣去回應.
我們怎樣去tackle.
用甚麼的立場.
其實反映著我們的教會.
會成為一間怎樣的教會.
更加是我們信的一個甚麼神.
我們的上帝究竟是一個怎樣的上帝.
所以那個是我在神學思想上.
去一路回應這些話題的時候.
我一路很緊張.
不只是我們要做到甚麼.
因為那個doing.
其實是define我們的being.
究竟是甚麼.
教會今天的being究竟是甚麼呢?.
所以那個反而是我緊張的一件事.
我覺得你剛才開頭說得不錯.
很多時候教會的招牌就是.
信耶穌得永生.
是.
由那裡開始我們已經有一個概念.
就是上帝就是來救我離開這裡.
而不是來這裡作王.
這個是完全是兩回事.

$^{601}$上帝要我們來這裡作王.
我們是代表祂在地上.
作光,作電,作見證.
還有去發聲.
變成我們傳福音給別人的時候.
就會是entail.
就是你走.
信了耶穌就走.
自然我們會遇見一些社會和政治的問題.
我們就覺得不需要談.
避.
避而不談.
我想因為這件事是我們不懂得做的事.
如果是我們耳熟能詳的事.
或者我們可以engage.
這個conversation.
但是在教會.
因為剛才說到我們環社會.
整個背景我們是不懂得怎樣去碰政治.
不懂得怎樣去面對這些議題.
在教會裡面從來不教.
又有些弟兄姊妹去從政.
教會都未必是一個很支持的姿態.
你怎樣去教弟兄姊妹.
我們要在世上做炎紅光.
但其實你又不要去equip他們.
甚麼真真正正在世上做炎紅光.
又不是教會裡面.
不只是送飯.
不只是關心盧叔.
我們是good stuff.
但你真的走入社會的時候.
為甚麼基本上是完全失去了他的發聲功能呢.
如果一定要拿出抗議牌.
而是面對制度的不公義.
社會的腐敗.
我們有沒有engage過.
有沒有去教導他們.
怎樣去分辨甚麼是真理.
怎樣行出真理呢.

$^{641}$或者怎樣去看一個帝國.
重返聖經.
特別是很多時候是不看的舊約裡面.
怎樣去活出一個信徒應有的生命.
我想這是一個很好的決策.
我們在這裡結束.
我們希望能夠引起大家去問天國觀.
神學.
怎樣成為一個地上的教會.
怎樣做宣教.
我們希望這件事不是單單是國安法我們要經.
而是引起我們對世界的關注.
多謝您收睇時局新聞,再會!.
\newpage



\section{}
\label{sec:u3NKvcUYcmI}
\textbf{聖地網上行 - 步主腳蹤 [第五集] 加利利 - 撒種者湖灣/八福山/抹大拉/希布斯}
\newline
\newline
連結: \href{https://youtube.com/watch?v=u3NKvcUYcmI}{\texttt{https://youtube.com/watch?v=u3NKvcUYcmI}} ~~~~ 語音日期: 2020-09-06
\newline
\newline
\hyperref[sec:aXI0__JmVi0]{\small{< < < PREV SERMON < < <}}
~
\hyperref[sec:index_chronic]{\small{[返順時目]}}
~
\hyperref[sec:index_scriptual]{\small{[返順卷目]}}
~
\hyperref[sec:m3nKOXb4VNI]{\small{> > > NEXT SERMON > > >}}
\newline
\newline
$^{1}$歡迎來到聖地網上行的第五集.
我們在這麼多集中希望你們能夠了解到.
原來聖經的地理和背景的處境是非常重要的.
往往因為我們現代的信徒.
跟聖經的距離有差別.
所以我們很多時候看聖經.
我們會錯過了這些地理,文化,背景的細節.
但如果我們這樣略過這些細節.
我們就會看到聖經的背景是非常清楚的.
我們就會失去了聖經本身很豐富的意思.
也可能有這樣的危機.
就是將自己的想法.
現代的想法加進聖經裡面.
所以我們不是單單旅遊.
我們其實是去到賢經.
去看看聖地不同的地方.
然後我們就去連結聖經裡面的經文.
我們上次去到跟隨耶穌的腳蹤.
我們來到加利利.
在你下載的地圖.
北邊的地方就是加利利.
當然中間最重要的地方就是加利利湖.
耶穌就是用加利利湖的其中一個漁港.
叫做加帕隆.
成為祂的傳道基地.
所以耶穌的教導和講道.
都是用當時的地理,地景.
日常人能夠接觸到的事物來做材料.
今天我們去讀《馬可福音》第四章.
耶穌有到海邊教導人.
有一大群人聚集到他那裡.
因此他上船坐下來.
船在海中.
群眾都朝著海站在岸上.
他用比喻教導他們許多事.
在教訓中他說.
你們聽著.
有一個殺種的出去殺種.
殺的時候.
有的落在路旁.

$^{41}$小鳥飛來就吃掉了.
有的落在泥土不多的石地上.
因為泥土不深.
很快就長起來.
但太陽一出來就把它曬乾.
沒有根就枯萎了.
有的落在荊棘裡面.
荊棘長起來就把它擠住.
它就結不出果實.
有的落在好的土壤上.
生長繁茂結出果實.
三十倍,六十倍.
有的是一百倍.
耶穌又說.
有耳可聽的就應當聽.
今天的教師通常都會站起來去講書.
但第一世紀的背景就不是這樣.
猶太人的社會裡面.
特別是猶太人律法的拉比.
他們教學的姿勢是坐下來的.
海邊當然是指加利利湖邊.
有很多人來到他聽到.
你可以想像耶穌就把小艇撐出去.
他就面對著岸上的群眾.
這個很實體的形容加利利湖的地理.
如果我們畫一條橫切面的話.
由左至右,由西至東.
畫在加利利湖上.
你就會發現湖的四周都是斜的山坡.
這些加利利湖的湖岸.
自然天然的山坡.
就變成了天然的劇場.
耶穌活躍的地方就是加伯隆.
就是湖的西北岸.
就是西邊猶太人的地方.
只要我們在加伯隆步行南面.
大概15分鐘左右.
就會發現一個小湖灣.
你看看這裡的地景.
環境有什麼特別呢?.

$^{81}$你看到半圓形的湖灣.
然後背後當然是斜的山坡.
這個小的山坡就好像成為了一個.
天然的羅馬的劇場.
事實上有些人來過這個地方實體做研究.
做一些聲學上的研究.
Acoustic的研究.
顯示這個湖灣上的山坡.
其實馬路以下可以坐到5000至7000人.
加上馬路上的山坡.
其實坐了可以有1萬至1萬4千,1萬5千人.
比較一些在地中海世界裡面.
找到著名的希臘劇場.
那些劇場可以坐到1萬4千人.
其實相比之下.
加里里湖是一個天然的劇場.
絕對是不遜色的.
你可以想像一下這個地景.
耶穌在小船艇上撐出去一點.
所有山坡上的群眾都會看到他.
並且在山坡上距離耶穌大概100多米的山上.
都可以聽得清楚他所說的話.
你想一想聖經裡面重要的宗教聚集的地方.
例如西乃山羅十界,班布十界.
士劍的基利森山.
還有士羅,庭柏舊約的回望地方.
耶路撒亞等等的地方.
其實之所以成為宗教的中心.
其中的原因當然是神學上的原因.
歷史上的原因.
但最重要的原因其實是一些很實際的因素.
就是這些地方都有山與谷.
成為了一些天然的劇場.
人群可以聚集並且聽到.
所以有些人說這個加里里湖的小湖灣.
叫做殺種者的湖灣.
是耶穌說的一個比喻的處境.
當然是不是真的在這裡.
這個不是重點.
因為整個加里里湖的湖邊.

$^{121}$全部都是一些天然的劇場.
更加重要的是有一個處境.
你就會看到比喻裡面說什麼.
講道的第一句是耶穌說什麼.
他說你們聽.
他跟天然劇場上的聽眾說.
你們要聽.
比喻裡面說到的種子.
就是殺在路邊.
有些是土淺的石地.
有些是荊棘裡面.
有些是好的土壤.
所以在天然山坡上.
劇場上的群眾.
他們一看就看到.
在常生活中觀察到的東西.
路邊土淺的石地.
荊棘.
好的土壤等等.
耶穌是用當時的天時地利.
用當時的科技.
用當時人們見到.
捉到摸到的東西來講道.
殺種的比喻有什麼特別地方呢?.
特別是在馬後福音.
就是耶穌第一個比喻.
第一個parable.
這個比喻就是挑戰我們.
你願不願意成為一個願意聽到.
並且願意接受的人.
不是再挑戰一些外界的環境因素.
在劇場裡面沒理由看不到.
沒理由聽不到.
留意馬可方的上文.
馬可刻意將一些不聽到的人說出來.
他刻意將法利賽人.
就是當時宗教的權貴.
他們的耳朵是聽到耶穌所說的話.
眼睛是看到耶穌所做的.
但是他們不願意接受.

$^{161}$不願意去承認.
甚至想去除滅耶穌.
第一個比喻就是問我們.
到底我們心裡面是否願意去接受主的道.
我們的心是否一個好的土壤.
好的泥土.
願意讓這些道.
塞在我們的心裡面.
能夠去成長.
所以最後的時候.
有意可停的就應當停.
我們再次出發到湖灣的山上.
這個地方上次提過了.
就是叫做八福山.
Mount of Beatitudes.
這裡紀念耶穌最著名的港島.
登山部分Sermon on the Mount的地方.
這裡天主教的教堂.
環境是非常優靜.
下面就看到很漂亮的加利利湖.
這個現代教堂是1936年建造的.
它的設計是很簡樸.
一點都不浮誇.
我覺得與登山部分所講.
與世不同的價值.
那種教導的精神其實都很符合.
過去一千六百多年.
初期的信徒特別拜占庭時期.
他們來到這個地方.
他們就要去讀耶穌的八福.
讀耶穌的登山部分.
到底是不是這個地方.
耶穌就站在這裡講登山部分呢?.
當然我們無從稽考.
但是二千年前.
耶穌和門徒所見到的景色.
山與水其實與今天相似.
我們欣賞這個教堂之餘.
我們不如就去讀讀.
登山部分的頭一段八福.

$^{201}$這是新日本的翻譯.
馬太福音第五章.
耶穌看見群眾就上了山.
他坐下之後門徒來到他跟前.
他就開口教訓他們.
心靈疲乏的人有福了.
因為天國是他們的.
哀痛的人有福了.
因為他們必得安慰.
溫柔的人有福了.
因為他們必承受地土.
愛慕公義如基如學的人有福了.
因為他們必得飽足.
憐憫人的人有福了.
因為他們必蒙憐憫.
內心清潔的人有福了.
因為他們必看見神.
使人和平的人有福了.
因為他們必稱為神的兒子.
為義遭受迫害的人有福了.
因為天國是他們的.
人若因我的緣故辱罵你們.
迫害你們並且捏造各樣壞話.
誹謗你們你們就有福了.
你們應該歡喜快樂.
因為在天上的賞賜是大的.
在你們以前的先知.
他們也曾這樣迫害.
這裡提到福,Blessing.
不是我們華人想像的那種.
有錢現在可以發達.
而在聖經裡面說有福.
是最終得賞賜的意思.
雖然現在未必是順境.
但他們知道有盼望.
將來一定上帝會賞賜他們.
所以他們現在就能夠.
繼續忍耐和相信.
這裡提到的動作也很重要.
這裡提到耶穌上山去講登山部分.

$^{241}$猶太人一想起的時候.
就會想起舊約裡面.
很多時候都會上山.
摩西上西乃山.
他就拿上帝的律法.
頒佈律法的條文給神的子民.
所以馬太福音也特意將耶穌的講道.
看成是摩西上山拿律法的一樣.
馬太福音由開始到現在.
五章聖經都形容到耶穌.
好像另一個摩西.
新的摩西.
甚至比摩西更大的救主.
耶穌出生的時候.
摩西的出生也是一樣.
摩西出生被暴君法路追殺.
耶穌出生也是一樣.
摩西從埃及出來.
經過曠野才進入迦南地.
耶穌也是一樣.
下了埃及.
然後上帝帶他出來.
然後來到以色列這個地方.
並且四十天受他試探.
耶穌在這裡一樣上山.
他就去頒佈天國的律法.
就是說聽了又願意遵循神話語的人.
這些才是真正天國的子民.
真正他的門徒.
到底我們是不是真正耶穌的門徒呢?.
在這個地方.
我們可以看到.
耶穌的律法.
在這個地方.
我們可以看到.
耶穌的律法.
在這個地方.
我們可以看到.
耶穌的律法.
在這個地方.

$^{281}$我們可以看到.
耶穌的律法.
在這個地方.
我們可以看到.
耶穌的律法.
在這個地方.
我們可以看到.
耶穌的律法.
在這個地方.
我們可以看到.
耶穌的律法.
在這個地方.
我們可以看到.
耶穌的律法.
在這個地方.
我們可以看到.
耶穌的律法.
在這個地方.
我們可以看到.
耶穌的律法.
在這個地方.
我們可以看到.
耶穌的律法.
在這個地方.
我們可以看到.
耶穌的律法.
在這個地方.
我們可以看到.
耶穌的律法.
在這個地方.
我們可以看到.
耶穌的律法.
在這個地方.
我們可以看到.
耶穌的律法.
在這個地方.
我們可以看到.
耶穌的律法.
在這個地方.
我們可以看到.

$^{321}$耶穌的律法.
在這個地方.
我們可以看到.
耶穌的律法.
在這個地方.
我們可以看到.
耶穌的律法.
在這個地方.
我們可以看到.
耶穌的律法.
在這個地方.
我們可以看到.
耶穌的律法.
在這個地方.
我們可以看到.
耶穌的律法.
在這個地方.
我們可以看到.
耶穌的律法.
在這個地方.
所以麥當勞的莫大拉.
是盛產鹹魚.
並且釀製過的魚.
用鹽醃過的魚.
是馳名海外的.
甚至一些羅馬.
希臘的地理學家.
都說這個地方.
是很出名的.
譬如其中一個地理學家.
叫做Strabo.
他提過莫大拉這個地方.
是一個名叫Tarachia的地方.
那裡的湖供應上好的魚.
用鹽醃過的.
然後岸上有不同的果樹.
所以看到莫大拉的時候.
會想起鹽和用鹽醃過的魚.
這些鹹魚是馳名海外的.
莫大拉的麥當勞.

$^{361}$過去十年有很多新的發現.
其中一個最重要的.
就是一個一世紀的猶太人會堂.
這個對大家來說.
絕對不陌生.
你看過之前的幾集.
我們看過多少個會堂.
加伯隆有一個很大的.
那撒勒有一個重建的.
你知道這些猶太人的會堂.
有一個格式.
三面可以坐人.
中間就是港台的地方.
這個會堂特別的地方.
就是發現一個方塊的石.
有裝飾的.
學者叫做莫大拉石.
這個莫大拉石.
幾邊就有裝飾.
我們簡單說說.
有兩個大的水瓶.
封城的水.
而且刻上一個三腳的.
七燈台.
大家不知道能不能看到.
這個很特別.
這個七燈台的雕刻.
至今是最古舊的.
最早期的七燈台的石刻圖.
畫出這個燈台的人.
很可能都是親眼見過.
可能是祭司.
在聖殿裡面見過七燈台.
這個古漁港.
另一邊就是一個很蓬勃的漁市場.
一個市場.
裡面有很多的水池.
當然有些不是養魚的.
有些是猶太人的潔淨池.
但也有一些是塑膠的.

$^{401}$一些是魚池.
那些漁民捕捉了的魚.
就會放在魚池上面.
暫時不可以.
因為魚池是很大.
所以就放在漁池上面.
暫時養著牠們.
這個就考考大家.
看到這個圖的時候.
不知道大家能不能找到一個石碼頭.
和一些停船用的石.
當時的水當然比現在多.
所以已經浸到上來.
這個地方.
細心一點看看.
這裡一排的石.
其實就是一個碼頭.
碼頭的邊上凸出來的.
幾塊石.
石上面有洞.
那些洞就可以給繩綁在上面.
小的漁船就可以.
停泊在岸邊.
所以不同的漁港.
捕捉到的魚就會運到這裡.
然後養著牠們.
或者去到處理.
曬乾醃乾牠們.
然後運到羅馬世界.
不同的地方.
所以耶穌說地上的鹽.
那些門徒.
有很多漁夫.
在加里里的漁民或農民.
他們想到的是什麼呢.
鹽就想到莫大拉.
這個最著名的港口.
因為是用鹽最多的地方.
享負盛名.
他們走遠至羅馬.

$^{441}$整個地中海世界都會讚賞他們.
所以他說.
你們是地上的鹽.
其實就是問我們有沒有那種生命.
可以揚名四海.
還是基督徒.
我們的明星.
是臭名遠播.
---!--.
(翻譯有啟用).
(嘉日和宅不當處理!!).
(後面還會有捷徑,請大家小心).
(字幕由【阿波羅網】提供).
耶穌又繼續說,馬太第五章十四節.
你們是世上的光,見在山上的城市無法隱藏的.
人點了燈,不會放在良器底下.
而是放在燈台上,就照亮一家人.
照樣,你們的光應當照在人前.
讓他們看見你們的好行為.
又頌讚你們在天上的富.
在古代的近東世界,成見在山上是非常常見.
但是耶穌說到的地方,說登山部分的一帶是加里尼湖.
我們去探測一下那裡的地景.
湖的旁邊有不少城鎮,多數是漁港,十幾個.
但是沒有一個漁港,沒有一個城市是建在山上的.
因為這裡的地景是湖邊.
但是湖的另一邊是羅馬城市,建在山上.
你記不記得加里尼湖是分開兩邊的.
西邊是屬猶太人的地方,東邊是外邦人的地方.
我們去一下這個地方.
這個地方叫做Hippos,Susita.
中文翻譯可以叫做Hippos,或者Susita.
希臘話,或者希伯來語,意思是馬.
Hippos是希臘話,Susita就是希伯來語.
因為這座城是坐落在馬鞍上的山上.
當然需要一些想像.
重要的地方是耶穌的時代,這個城市是非常重要的.
重要的地方是Susita這個地方,Hippos.
它有自己的硬幣,可以鑄造自己的硬幣.
在高高在山上,你會看到這個城市很有羅馬味道.

$^{481}$有羅馬文化的豪華生活方式,是一種帝國生活的形式.
建造在山上,告訴山下一帶的人,這是羅馬盛世的好處.
如果你靠攏羅馬,你就可以得到這些享受.
留意一下,這裡加里尼湖的東邊地理,就是羅馬帝國的東方邊境.
當時的文件告訴我們,羅馬帝國在這裡設立了十座自治的城市.
Autonomous Cities,Hippos是其中一個的城市.
雖然聖經沒有直接提及希布斯這座城.
但福音書裡面有提及過所謂十座城,叫做低加玻利,英文是Decapolis.
馬太福音第四章,這是登山補份之前,登山補份第五章.
這裡是第四章,就是耶穌來到加里尼做什麼呢?.
走遍加里尼,耶穌在各會堂教訓人.
我們看到在莫大拉有會堂,在加伯隆有會堂.
祂就傳遍天國的福音,去醫治百姓各樣的病疾.
祂的名聲就傳遍敘利亞,敘利亞就是加里尼湖的東北邊.
有許多人從加里尼來本地.
還有對岸的耶路撒冷猶太人,約旦河外的人都來跟隨祂.
再留意一下希布斯的地景,對面就是耶穌講道的地方.
登山補份就在對面,加里尼湖的另一邊.
猶太人的地方,那裡有加伯隆,有莫大拉.
希布斯的這一邊就是外邦人,羅馬人的地方.
對於聽到的人來說,門徒來說.
一講到建在山上的城市,在加里尼這一帶只有希布斯.
而他們在對岸一看就看到.
加里尼湖.
用一下你的想像,你會看到羅馬城市.
有兩三層樓高的建築物,在對面遠遠就能看到.
日間也能看到,晚上也能看到.
晚上當這個城市點綴一些光的時候.
當羅馬劇場有一些表演的時候.
他們就看到一些光,對面就看到.
成建在山上,不能夠隱藏.
作為羅馬城市,特別是在防線邊境的羅馬城市.
它有什麼作用呢?就是政治作用.
它是代表羅馬京城來發揚羅馬帝國的好處.
來讚揚羅馬帝王,其實是羅馬帝國的大外宣.
是政治的宣傳,設立在這裡其實是告訴人.
羅馬有多好,其實當時一些文獻都說.
羅馬京城是世上的光,自稱為世上的光.
因為每一個羅馬的城市就代表世上的光.
耶穌就把這些帝國的語言反過來說.

$^{521}$祂說我們作為門徒,我們才是世上的光.
而世上的光不是理念上的信主那麼簡單.
世上的光是有好的行為.
好的行為不只是做一兩件好事.
好的行為是整個生活,我們發自內心.
發自一個已經有上帝的生命出來的生活狀態.
就是我們跟上帝有一個特別的關係.
所以經文這裡說你們在天上的父.
我們在發揚我們在天上的父.
這是馬太夫人第一次用天上的父,Heavenly Father.
就代表我們門徒跟上帝有一個特別的關係.
是以主的律法為我們中心.
因為猶太人的典籍裡,拉比都認為.
律法就是世上的光,原來我們就是上帝的律法.
我們的生活就是上帝的道.
就是讓人看到上帝的國度,天國.
有人說得很對,他說未信主的人是看不到耶穌.
我們就是他看到的耶穌.
未信主的人是讀不到聖經.
他們讀的聖經就是我們的生命.
我們就是他們的聖經,我們就是他們的上帝律法.
如果我們的生活是遵循上帝的律法.
遵循主的教導,我們身體力行出來的.
我們的目的就是宣揚上帝的好處.
我們每天所想的,所求的.
最終是要去榮耀上帝,去發揚上帝的好處.
去告訴全世界的人,上帝有多好.
上帝有多愛我們.
我們一起低頭祈禱.
上帝我們求你幫助我們.
讓我們的生命,我們的生活,我們的心態.
有著你和我們在恩典當中的關係.
你就是我們生命當中天上的父親.
主是讓我們的生命,我們的生活.
來榮耀上帝.
讓我們知道我們的新身份.
就是世上的炎.
能夠在世界各地不同的地方.
能夠發揚你的美名.
讓我們的生命成為光.

$^{561}$在世界當中黑暗的地方.
我們成為你的點光.
求你幫助我們,讓我們作為你的門徒.
不只是獨善其身.
而是我們為著世上的人.
願意去發揚上帝在天國的價值.
讓我們生命,我們的心態.
我們願意遵從主你的教導.
我們這樣祈禱奉主名.
\newpage



\section{}
\label{sec:m3nKOXb4VNI}
\textbf{聖地網上行 - 步主腳蹤 [第十三集] 十字架}
\newline
\newline
連結: \href{https://youtube.com/watch?v=m3nKOXb4VNI}{\texttt{https://youtube.com/watch?v=m3nKOXb4VNI}} ~~~~ 語音日期: 2020-09-06
\newline
\newline
\hyperref[sec:u3NKvcUYcmI]{\small{< < < PREV SERMON < < <}}
~
\hyperref[sec:index_chronic]{\small{[返順時目]}}
~
\hyperref[sec:index_scriptual]{\small{[返順卷目]}}
~
\hyperref[sec:o6x_Tz0d6Ww]{\small{> > > NEXT SERMON > > >}}
\newline
\newline
$^{1}$我們由第一集開始追溯耶穌的生平.
現在已經來到耶穌的最後一天.
或者祂最後的幾個小時.
現在耶路撒冷的面貌當然與第一世紀不一樣.
我們今天去到耶路撒冷舊城.
看到的城牆是公元後16世紀.
土耳其奧圖曼時代所興建的.
耶穌的時候是公元後30年代.
耶路撒冷的城牆是這樣的.
上一次我們說到.
彼拉多是在城西的大希律王宮審耶穌.
之後就交給羅馬兵丁去戲弄邊打.
然後耶穌就要背起自己的十字架橫樑.
上刑場割割他.
這是歷史裡面第一世紀.
耶穌所走的十字架的苦傷路.
via dolorosa.
拉丁文意思是way of sorrow.
即是傷痛受苦之路的意思.
當你聽到via dolorosa的時候.
可能會想起一首很出名的.
八十年代的基督教的英文金曲.
叫做via dolorosa.
或者你聽到苦傷路的時候.
你就會想起今天在耶路撒冷裡面的.
via dolorosa的苦路.
這條苦傷路是由聖殿的西北角.
一個叫Antonio Fortress羅馬的兵營.
叫做安東尼亞堡開始的.
之前的人以為彼拉多.
就是在Antonio Fortress審耶穌.
所以苦路歷時歷代就由這個地方開始.
但這條苦路的傳統.
由公元後十四至十八世紀.
幾百年不斷地演變出來.
當時被這些聖地的朝聖者.
幫助他們實地體會一下.
耶穌由受審到去迎場.
各個他中間路上的遭遇.
或者你可以這樣想.

$^{41}$這個傳統的苦路via dolorosa.
其實是一種受難節的禮儀.
苦傷路不是一個實體的地理.
不是geography of facts.
而是一個信心的地理.
geography of faith.
所以這條苦路有十四個站.
根據經文,教會的傳說.
或者實際的想像.
來重構耶穌十字架的經歷.
全長大概六百米即二千呎.
今天我們一起走這條苦傷路.
via dolorosa.
.
苦路的第一站是安東尼亞堡.
在第一世紀的時候.
大希律王在這裡修建聖殿的時候.
加了一個大的兵營.
以至羅馬軍隊能夠在這裡駐守.
並且監視著聖殿山上一舉一動.
馬科方第十五章.
彼拉多要眾人起悅.
釋放巴拉巴及他們.
將耶穌鞭打了.
交給人釘十字架.
羅馬兵丁就將耶穌帶到.
去衙門的大院裡面.
然後叫齊全營的兵.
將耶穌穿上紙袍.
又用荊棘編作冠冕給大尚.
又恭賀他.
說恭喜猶太人的王.
又拿一斤的毀子打他的頭.
套他抹在他的臉上.
屈膝拜他.
戲弄完之後就脫下他的紙袍.
叫他穿上自己的衣服.
帶他出去釘十字架.
這些劇本或形容.
全部都是典型的羅馬帝國.

$^{81}$釘十字架的程序.
來到第二站.
這裡紀念耶穌背起他的十字架.
當時羅馬或希臘作家會提到十字架的刑罰.
例如這個Protag.
他說每一個死刑犯被判十字架的人.
必須自己背起十字架.
指著十字架的橫樑.
自己走到刑場受死.
告訴被征服的人民.
羅馬帝國.
羅馬這個世界已經征服了你.
來到虎廬第三站.
這裡的想像是耶穌第一次跌倒.
以前有一個羅馬高官叫Cicero.
他將一個羅馬公民捆綁在刑事裡.
將羅馬公民釘十字架.
他說十字架是羅馬帝國最嚴厲的刑罰.
是為了非羅馬人民.
尤其是政治反意見分子而設.
羅馬公民是不需要受十字架的刑罰.
因為這樣代表羅馬帝國自我摧毀.
(羅馬公民被釘十字架).
虎廬第四站.
想像到耶穌母親瑪利亞.
在虎廬的途中遇上耶穌.
因為她在這裡跌倒.
這是一個母愛的人性的刻畫或想像.
因為瑪利亞確實跟著耶穌.
在十字架底下到達十字架的刑場.
第五站.
紀念古利諾西門被羅馬兵強迫為主背十字架.
古利諾是今天北非利比亞的地方.
他剛剛入城時就見到耶穌背著十字架.
教會以往多年的傳統都以古利諾西門為門徒的典範.
因為他為耶穌背十字架.
但經文裡說得很清楚他是被強迫的.
會不會有些背景幫到我們.
當時羅馬王巡行時不會一個人.
而是一定會有十二個護兵.

$^{121}$有時二十四個.
拿著一些木棍.
用木棍紮成代表軍隊團結.
上面加上一個斧頭.
代表帝國殺頭的武器.
這些木棍的作用.
陪著皇帝一同行行.
其實是想放大皇帝的尊貴和地位.
你這樣看的時候.
你就知道羅馬軍人迫西門背起十字架.
拿起這支木棍跟耶穌同行.
其實不是榮耀他.
相反是嘲諷和羞辱他.
所以加上西門強迫他背十字架的時候.
其實是放大耶穌釘十字架的恥辱.
所以鏡頭似乎不是在古利諾西門身上.
而是在耶穌身上.
他才是我們門徒的典範.
因為他興翰羅馬帝國十家的羞辱.
所以我們可以看到.
羅馬軍人的羞辱和羞辱.
是在西門的十字架上.
和羅馬軍人的羞辱和羞辱.
是在西門的十字架上.
和羅馬軍人的羞辱和羞辱.
是在西門的十字架上.
和羅馬軍人的羞辱和羞辱.
是在西門的十字架上.
和羅馬軍人的羞辱和羞辱.
是在西門的十字架上.
和羅馬軍人的羞辱和羞辱.
是在西門的十字架上.
和羅馬軍人的羞辱和羞辱.
是在西門的十字架上.
和羅馬軍人的羞辱和羞辱.
是在西門的十字架上.
和羅馬軍人的羞辱和羞辱.
是在西門的十字架上.
和羅馬軍人的羞辱和羞辱.
是在西門的十字架上.

$^{161}$和羅馬軍人的羞辱和羞辱.
是在西門的十字架上.
和羅馬軍人的羞辱和羞辱.
是在西門的十字架上.
和羅馬軍人的羞辱和羞辱.
是在西門的十字架上.
和羅馬軍人的羞辱和羞辱.
因為你想想.
這些傳記背後.
為何會是一個女孩.
你讀福音書的時候.
十字架的技術應該有一個印象.
剩下的是女門徒.
到廉文和耶穌的同行.
最多都是加多一個.
當時最年少的約翰.
為主所愛的門徒.
是他們才會跟主到底.
居然是一群女子婦孺.
是他們才會跟主到底.
居然是一群女子婦孺.
是他們才會跟主到底.
居然是一群女子婦孺.
來到第七站.
這裡就是想像耶穌第二次跌倒.
在意大利一個港口.
找到一個屬於二世紀.
很早期的一個圖畫.
上面畫著一個被釘十字架的人.
雙臂是剪開的.
你見到他皮膚全部沒有了.
只剩下骨骼.
另外一個屬於二世紀的.
一個圖畫叫Alexandros Graffiti.
圖畫上你見到一個人被釘在十字架上.
但這個人像的頭是小驢頭.
十字架下有一個人.
他的手是高舉的.
這是敬拜的姿勢.
圖畫上畫了一些粗略的文字.

$^{201}$Alexandros Sebeti Deon.
敬拜神.
羅馬人一直有一個傳說.
猶太教是敬拜驢子.
當然這不屬實.
因為羅馬人對猶太人的歧視.
第一代的基督教是出自猶太教.
羅馬人順理成章.
把基督教所拜的耶穌當成驢子.
或者一些侮辱嘲諷基督教的圖畫.
十字架的道理對外邦人來說.
就是愚蠢.
對猶太人來說是一個半腳石.
怎麼可能有人相信.
試一試當自己是一個百分百的羅馬公民.
從小到大你所接受的教育.
所接收的資訊.
每天看到的景象都是羅馬帝國的威榮.
最吸引的成功榮耀就是羅馬所定的.
羅馬給予你平安.
給予你好吃好住.
給予你安定.
所以你一定要跟著羅馬帝國的世界.
去生活.
這就等於美國的American Dream.
羅馬帝王說我才是神的兒子.
我才是救主.
我才是萬人的君王.
我使你安定繁榮.
所以你要信靠我.
但基督教的信仰.
宣告一位顛覆羅馬帝國.
世界規律的神.
來到第八章.
這裡講到耶穌對耶路撒冷的婦女說話.
路加福音第23章.
有很多百姓跟隨耶穌.
內中有很多婦女.
這些婦女為耶穌哭號痛哭.
耶穌轉身對她說.

$^{241}$耶路撒冷的女子不要為我哭.
要為你們自己.
和自己的女兒去哭.
這裡預言.
公元後七十年.
耶路撒冷被毀.
因為耶路撒冷不肯悔改.
不肯承認.
這個被釘十字架的耶穌.
才是真正的君王.
寧願靠著親愛的.
羅馬帝國.
以保現在的平安.
誰知在一代之內.
耶路撒冷被毀.
第九章.
想像耶穌第三次跌倒.
十字架的過程.
被審,被鞭打,被虐待.
到他上刑場中間.
有很多創傷.
耶穌在十字架.
大概幾個小時.
你便可以知道.
他之前受過多少折磨.
我們來到第十章.
觸頭地.
這裡紀念主被削去衣服.
我們現在在聖母堂.
Holy Sepulchre的外面.
這是羅馬十字架的程序.
紀念主.
是程序.
劇本寫好了.
被判釘十字架.
你便背起十字架.
自己上刑場.
最後軍人將你的衣服削去.
赤裸裸地將你釘上.
去示眾.

$^{281}$告訴別人.
我們羅馬帝國才是特性.
很多時候我們想起耶穌被釘十字架.
他也可能會這樣想.
是完全被動,完全消極.
但這裡似乎不是.
福音書所堅持的.
好像是相反的.
例如按照世人的常理.
世界的常理.
被釘十字架的人.
必然是咒詛其他人.
但他居然說.
父啊赦免他們.
十字架七言之一.
十字架每一句.
幾乎每一句.
都是反轉羅馬帝國.
釘十字架的程序.
這些反轉過來.
很詫異的情況.
尤其在約翰福音出現.
約翰福音第十二章早說了.
現在這世界受審判.
這世界的王.
要被趕出去.
撒旦要被趕出去.
我若在地上被舉起.
這裡其實是說釘十字架.
被舉起就會吸引萬民.
來歸我.
耶穌說.
原是指自己要如何在十字架上死.
這個世界.
耶穌說.
是羅馬帝國給我們的世界.
他的規律.
耶穌在十字架上.
被高舉.
相反的.

$^{321}$他好像被羞辱.
但他在神的眼中只有榮耀.
好像登基一樣.
被提升.
聖國這個世界反而會賜給人.
我們一班罪人.
得到榮受.
(十字架).
來到第十一站.
我們進入聖母堂.
這裡有一個富堂.
上面的圖畫表明.
十一站是主被釘在.
十字架的木頭上.
關於釘十字架.
最重要的考古發現.
是1968年.
在耶路撒冷的古墓中找到的.
考古人員在墳墓中.
找到一個石骨盒.
我們曾經說過.
盒子上刻有名文.
約哈蘭.
其中一個骨頭.
是腳筋骨.
是一支長的鐵釘插穿了.
你可以想像.
約哈蘭被釘十字架的形態.
左右腳筋都有一支鐵釘穿過.
其中一口的鐵釘.
釘在十字架的木頭上.
可能是彎曲了.
這支落葬的時候.
也拔不出來.
手骨不見有明顯的損傷.
所以推斷.
這個人被釘十字架.
雙臂被綁上.
腳上各自釘在木頭上.
當時不同的文獻都說.

$^{361}$釘十字架有不同的形狀.
耶穌的形態是怎樣的呢?.
初期的文獻或科學書中說到.
祂手上有釘痕.
應該是祂左右手的手腕.
和兩隻腳的腳筋.
都被釘上.
來到第十二章.
這裡是在剛才那一章的隔離.
這個赴堂的樹.
裝飾得很華麗.
並且紀念耶穌在十字架上死亡.
約翰福音第十九章這樣記載.
耶穌見母親和她所愛的門徒.
站在十字架旁邊對母親說.
母親看你的耳字.
又對那門徒說看你的母親.
從此那門徒就接她到自己家裡去了.
這時以後.
耶穌知道各樣的事已經成了.
為要使經常的話應驗.
就說我學了.
有一個器皿盛滿了醋.
放在那裡.
他們就拿海藻.
把醋綁在牛膝草上.
送到她口中.
耶穌嘗了那醋就說.
成了.
便低下頭將靈魂交付上帝.
在這裡看你的母親.
十字架上的人.
很多時候會交代最後的遺囑.
譬如物業或身後事等等.
耶穌照樣在十字架上做這件事.
他將母親交給所愛的門徒.
很應該是約翰.
約翰就將她接到自己的家裡.
自己的在哪裡出現過?.
在約翰第一章.

$^{401}$第一章說.
耶穌來到自己的地方.
自己的人都不接待他.
但是接待他的就是信他名的人.
他就像他們權柄作神的兒女.
所以這裡似乎不單單照顧瑪利亞.
而是照顧自己的家人.
所以這裡似乎不單單照顧瑪利亞.
這個寡婦.
是在十字架底下形成教會.
我們就是一群在十字架底下的人.
成了出現了多少次?.
三次.
耶穌知道是已成了.
已經成了.
並且經常說應驗.
其實原文是成了.
Telos.
他自己就說了.
十教七原其中一句.
成了.
是成就了上帝的救恩.
上帝的大計在十字架上成就.
買熟我們這一群人.
使我們能夠成為他們的子民.
來到第十三站.
這裡就紀念到耶穌的屍體.
從十字架被拿下來.
放在石上.
或者這個地方.
東正教等等的傳統.
信徒來到就喜歡摸這個地方.
或者親這個地方.
《約翰福音》第十九章.
上文下理說到.
羅馬兵丁.
羅馬軍人.
想去到肯定.
三個被釘上十字架的人.
耶穌和兩個強盜.

$^{441}$是真的死了.
三十三節.
只是來到耶穌那裡.
就見他已經死了.
不打斷他的腳.
這些事成了.
為要應驗經常的話說.
他的骨頭一斤也不可折斷.
骨頭一斤也不折斷.
是一句舊約的說話.
形容預節的羔羊.
出埃及記第一次預節的時候.
就說羔羊的骨頭一斤也不折斷.
就是因為這隻羔羊.
能夠得勝埃及帝國.
以至他們能夠得到自由.
在上帝的約下成為他的子民.
我們在十字架上見到的耶穌.
是一個真正的預節羔羊.
得勝這個世界.
苦路的最後一站第十四站.
這裡就是聖母堂最中心的地方.
這裡就是主耶穌的身體.
被安葬在一個新的墳墓裡.
當然今天的面貌和第一世紀的面貌.
已經完全不同.
這一帶當時第一世紀.
是一個墳墓.
看到的是一個圓形的大廳.
中間有一個小的教堂.
小的建築物.
蓋在耶穌的墓穴之上.
約翰福音第十九章.
「這時以後有亞利瑪泰約瑟」.
是耶穌的門徒之一.
怕猶太人.
就暗暗地作門徒.
他來求彼拉多.
要把耶穌的身體領去.
彼拉多允准.

$^{481}$他就把耶穌的身體領去.
又有尼哥德姆或尼哥底姆.
先前夜裡去見耶穌.
他帶著末曰沉香約一百斤來到.
他們兩個人就照著猶太人的殯葬規矩.
把耶穌的身體用西瑪布加上香料裹好.
在耶穌釘十字架的地方.
附近有一個園子.
園子裡面有一座新墳墓.
是從來沒有人葬過的.
只因是猶太人的預備日.
索伯就來了.
只因那一墳是相近.
他們就把耶穌安放在那裡.
這裡有幾處特別的地方.
第一是新墳墓.
一個未釘十字架的犯人.
很多時候都是葬在城外的亂葬崗.
但這裡說到耶穌是一個全新的墳墓.
在石山掘出來的.
在那裡被安葬.
約翰和其他科林書有少許不同.
他沒有特別看舊大石.
因為在約翰的神學當中.
十字架的死和埋葬.
絕對不是失敗.
復活也不是扭轉失敗.
約翰看十字架.
就是榮耀.
就是德勝世界的榮耀.
另外就是沉香沒藥約一百斤.
大家記得第一集.
我們曾經去過希律堡Herodium.
那裡是大希律王的墳墓地方.
約瑟夫告訴我們.
大希律死的時候是風光大葬.
除了金銀無數之外.
也強調有多少沒藥和沉香.
這是君王的葬禮.
這個墳墓的地方是源子Garden.

$^{521}$七十四億本.
用這個字來形容大衛君王的墳墓.
所以耶穌死的時候被安葬.
人的眼中.
羅馬世界的眼中.
看似是失敗.
但在上帝眼中.
祂是真正的君主.
門徒四散.
似乎一切都是失敗.
但這裡也給我們一線曙光.
約翰福音告訴我們.
有兩個人開始真正成為門徒.
阿里瑪太約瑟和尼哥德姆.
這兩個人有什麼特別.
他說他暗暗地作門徒.
因為怕猶太人.
尼哥底姆約翰福音第三章.
夜間來見耶穌.
因為怕猶太人.
十二章四十二節告訴我們.
為什麼有些人信了耶穌.
他們卻不敢認自己是門徒.
他說官長當中有些信耶穌.
只因法理錯人的緣故.
就不承認.
恐怕被趕出會堂.
失去了很多地位.
金錢或安定.
這是因他們愛人的榮耀.
過於愛上帝的榮耀.
這就是羅馬帝國所規定的世界.
之前我們說到出賣信仰.
其他人物的角色都表明.
受到這個世界的牽制.
在江湖裡面.
太祭司 哥瓦法 彼拉多 猶太等等.
因為金錢.
因為想安定.
他們就簡化他們的信仰.

$^{561}$或者出賣耶穌基督.
出賣了真理公理.
但耶穌說什麼?.
他說我勝了世界.
真正能夠改變世界.
能夠得勝世界.
耶穌說就是他上十字架.
在十字架上.
耶穌所求的.
只是父神的名聲.
父神的榮耀.
遵照父上帝在他生命中的心意.
他勇敢地走出來.
即使要輕看十字架.
你看了這麼多集.
到了第十三集.
你看到這一位.
甘願為我們犧牲的耶穌.
你應該看到神的愛.
神的恩典.
我們接受耶穌的一群人.
我們不單得到這個永生.
並且他告訴我們.
我們不需要怕.
在這個世界當中.
很多時候我們受不同的牽連.
但如果我們願意跟隨耶穌.
就應該活像他一樣.
堅持著愛神的榮耀.
多於世間上的榮耀.
無論世間的榮耀.
在你的生命當中.
翻譯出來的時候.
可能是一些安定.
一些好處.
或者群眾給你的肯定.
這些現實.
不一定是真實.
耶穌說.
我來到是開另一個國度.

$^{601}$是一個藐視羅馬帝國秩序的.
The Order of the Empire.
我們藐視他.
不會被他擺佈.
就算上十字架.
也不會威脅到他.
輕看羞辱.
拿神的榮耀為自保.
這才是真正的王者.
真正的自由.
真正門徒的生命.
我們願不願意這樣做?.
讓我們低頭祈禱.
因為你的獨生子耶穌基督.
輕看羞辱勝過世界.
在祂裡面.
我們能夠得到這個恩典.
得到這個特性.
主我們求你.
讓我們接受這個恩典.
為我們每一個聽見這些訊息.
我們求你今天成為我們真正的主.
讓耶穌成為我們生命當中的主.
並且你這位輕看羞辱.
甚至上十字架.
來到榮耀神的救主.
是我們生命的態度.
是我們生命門徒的典範.
求你幫助我們來到學校耶穌基督.
輕看世界的榮耀.
讓我們在我們生命當中.
只求你自己的名聲.
只求祂會順服你.
我們感激到,奉主名禱.
《輕看羞辱》 詞:陳家麗 曲:陳家麗.
人浪翻騰 士兵穿插.
千夫只向這個他.
受死前被鞭打 劇痛掙扎.
罪寬苦的群眾.
正在冷嘲戲弄 唾罵欺壓.

$^{641}$遍體傷痕 壓割痛血.
荊棘冠冕 靜似磚.
承受世人罪災的 沉重身影.
步步艱辛的路徑.
背負四周冷汗 肆虐苦食.
榮耀君王救主幾代.
成名以十家結組.
無罪高揚何故遭再殺另有.
全為愛不退歡服.
決定走上火流之路.
殘酷傷害無悔經過.
作犧牲 正為你我.
遍體傷痕 壓割痛血.
荊棘冠冕 靜似磚.
承受世人罪災的 沉重身影.
步步艱辛的路徑.
背負四周冷汗 肆虐苦食.
榮耀君王救主幾代.
成名以十家結組.
無罪高揚何故遭再殺另有.
全為愛不退歡服.
決定走上苦難之路.
殘酷傷害無悔經過.
作犧牲 正為你我.
顫怒 被血洗.
實戰 神大計.
無家贖罪救恩.
不知覺覺慘.
榮耀君王救主幾代.
成名以十家結組.
無罪高揚何故暗逃再某人.
為著愛不退歡服.
決定走上苦難之路.
曾在Fia Dolorosa 喝苦杯.
正為你我.
\newpage



\section{}
\label{sec:o6x_Tz0d6Ww}
\textbf{浩劫重生 【第1集:挪亞洪水】 陳崇基牧師 - 主日講道視頻系列}
\newline
\newline
連結: \href{https://youtube.com/watch?v=o6x-Tz0d6Ww}{\texttt{https://youtube.com/watch?v=o6x-Tz0d6Ww}} ~~~~ 語音日期: 2020-09-29
\newline
\newline
\hyperref[sec:m3nKOXb4VNI]{\small{< < < PREV SERMON < < <}}
~
\hyperref[sec:index_chronic]{\small{[返順時目]}}
~
\hyperref[sec:index_scriptual]{\small{[返順卷目]}}
~
\hyperref[sec:9fgDv4ygGRU]{\small{> > > NEXT SERMON > > >}}
\newline
\newline
$^{1}$我相信有一個最家傳戶曉最多人知道的.
就是《創世記》第六至第八章.
這裡記載一段的事蹟.
是羅亞方舟.
羅亞的大洪水.
我不知道你想起這個故事的時候.
你會想起什麼畫面.
不如我們看看聖經如何描述.
這個故事的背景.
我們先看看.
創世記第六章.
見事敗壞了.
凡有血氣的人在地上都敗壞了行為.
神就對羅亞說.
凡有血氣的人.
他的盡頭已經來到我面前.
因為地上滿了他們的強暴.
我要把他們和地一同毀滅.
你要用過飛木做一隻方舟.
方舟的做法乃是這樣.
要長三百爪 闊五十爪 高三十爪.
當羅亞六百歲.
二月十七日那一天.
大淵的泉源都裂開了.
天上的湘湖也淌開了.
四十晝夜江大雨在地上.
洪水泛濫在地上四十天.
水往上漲.
把方舟從地上飄起.
水勢浩大.
在地上大大地往上漲.
方舟在水面上飄來飄去.
水勢在地上極其浩大.
天下的高山都掩沒了.
雖然羅亞方舟大洪水這個故事.
是如此家傳戶曉的聖經故事.
但它是給現代人有最多誤解的一段記載.
可能你久不久會聽到人說.
說找到方舟的殘骸.
又或者有人這樣報導.

$^{41}$說找到一個普世全球洪水的證據.
在二十世紀有一個很出名的故事.
當時是1930年代.
一個英國考古學家叫做Wooley.
他去到南米索不達米亞一帶.
他去到那裡發掘聖經的古城吾爾.
即是阿巴拉罕的鄉下.
當他在吾爾城那裡向下發掘的時候.
他找到一層八尺厚的洪水積成的淤泥.
土層的下面是一個更遠古的人類的遺跡.
土層的上面是後期的遺跡.
換句話說.
這個八尺厚的純粹淤泥.
洪水積成的泥土.
就好像表示有一場大的洪水.
曾經在這個地方發生過.
當這個消息一公佈出來之後.
在當時的社會當時的世界就轟動一時.
有人認為終於找到聖經裡面說的.
羅亞大洪水的跡象.
不過繼續發掘下去的時候.
就發現在其他同一帶的城鎮.
就沒有在同一個時期找到洪水破壞的土層.
換句話說.
淤泥所找到的洪水的跡象.
只不過是一個地區性的洪水.
淤泥這個故事就告訴我們.
我們現代人去看羅亞大洪水的記載的時候.
我們經常都會去問一些現代人會問的問題.
證據歷史裡面的考證.
可能你會聽過.
有人說世界上很多的民族.
很多的傳說說曾經有過大洪水.
所以應該是有過全球大洪水.
但你細心想想.
這樣證明不了全球曾經有過大洪水.
特別是你講到羅亞的經文.
因為在經文引申出來.
這個大洪水有40週夜.
或者甚至一年這麼長.

$^{81}$所以如果你要證明曾經有過全球性的洪水的話.
你一定要找到所有的古代的城鎮.
他們這些古城的文明.
曾經一度同一時間.
在不同的城鎮裡面.
發生過一些洪水毀壞的斷層.
但是我們從來沒有找到這些.
第二你可能會想起.
有一個阿拉納山的傳說.
創世記第八章這麼說.
他說7月17日.
羅亞方舟停在阿拉納山上.
所以在土耳其的東邊.
有一個地方.
一個山叫阿拉納山.
傳說是羅亞方舟停泊的山頂.
但是這個傳說.
歷史的資料告訴我們.
只能夠追溯至公元後十世紀才出現.
所以這個阿拉納山.
並非一個最可靠的歷史傳統.
況且你細心看看聖經裡面的記載.
他說方舟停在阿拉納山.
阿拉納山的原文不是單數.
而是眾數.
意思並非在阿拉納山峰之上.
方舟停泊了.
真正地理的意思.
是在阿拉納群山.
或阿拉納山區.
方舟就停泊了.
可以是山頂.
也可以是半山.
也可以是山居.
(為什麼要追溯至公元前十世紀?).
很多時候.
我們缺乏歷史,地質,考古的證據.
我們就會訴諸於一個方法.
就是說這是一個神蹟.
如果是神蹟的話.

$^{121}$我們就說不需要想一些水力,地質等等.
不同的科學證據.
但如果你這樣想的話.
每樣東西都訴諸於神蹟.
有幾樣東西就不通了.
首先第一.
聖經裡面洪水,羅亞,方舟的技術.
確實是講很多大自然,自然的描述.
不是講到神蹟的歧視.
接著如果我們真的訴諸於神蹟的話.
你就退一步去想.
為什麼我們還要花這麼多時間.
去找這些洪水的遺跡.
或者是地質上的證據.
甚至考古的證據.
因為沒有人需要用這些證據.
去證明神蹟曾經發生過.
如果你這樣想的話.
你就似乎會比較近一點.
去問清楚到底經文想描述些什麼.
我們看回剛才大洪水,羅亞,方舟的經文.
你會留意到有很多普世的字眼.
這些普世字眼講到.
凡地上的活物無一不死.
毀滅天下等等.
現代人可能會想起.
一定是全球性的事件.
但是你留意聖經裡面.
很多時候講到普世的字眼.
只不過是形容一些地區性發生的事件.
譬如創世記第四十一章.
講到當時的埃及和中東.
發生很大的饑荒.
甚至是天下的饑荒.
各處的人都去到埃及問藥室取糧食.
了解歷史的就知道.
在這段時間根本沒有一個全球的饑荒發生.
只不過是在中東,埃及一帶地區性的饑荒.
所以用這些普世字眼不一定等於全球.
用一些普遍的字眼.

$^{161}$只不過是告訴我們那個intensity.
那個程度是非常嚴重的意思.
不如我帶你去一個地方.
你可能會更加了解.
當時近東的人想些什麼.
以及他用什麼字眼修辭.
來形容他們的世界.
這個地方你甚至沒有去過.
你都應該會聽過.
叫做大英博物館.
我們來到大英博物館.
我們進去其中一個展館五十五.
這裡是米索不達米亞近東的一個展館.
其中有一個手掌大的文物.
刻上了一個巴比倫古代的地圖.
代表著巴比倫人怎樣看待他們當時的世界.
這個地圖屬於舊約聖經的時期.
所以他們想的東西應該很像.
舊約時代的以色列人所想的東西.
你看到這個地圖有兩個大圈.
代表著他們整個的世界.
是一個圓形的平面.
地圖用不同的圖案來代表不同的地方.
加上標籤.
外圈之外有八個三角形.
代表著當時平面世界以外的靈魂.
右上方的三角形刻有日頭不見之地.
即是西面日落的地方.
兩個圈中間標籤了海.
這些海洋環繞著整個大地.
近東文獻中常出現的海洋叫做宇宙海.
Cosmic Ocean.
這些就是地的極處.
因為地的邊界之外就是人不能去到的大海.
內圈就是米索不達米亞這個地方.
有不同的城鎮.
世界的中心當然是中間的巴比倫城.
這個地圖當然不是按照科學的尺度畫出來.
但卻表達了古近東人眼中的看法.
他們所知的世界.

$^{201}$他們所知的known world.
是一個平面.
一個大的圓盤.
四周地極都是高崇入雲的山脈.
地極之處就是海洋.
這些海洋或宇宙山.
Cosmic Mountains.
就是地的根基.
因為他們的深處是立在陰間.
而山峰頂著的就是一個上面的鑊形的天空.
叫做穹瘡.
在他們的印象之中.
這個好像是一個倒轉了的鑊.
然後頂著上面的大水.
聖經時代的人.
他們看世界.
不會有地球的概念.
他們看世界是一個平面.
聖經當然是神.
用當時的人能夠明白的語言和概念.
來啟示他的真理.
所以你這樣想的時候.
羅亞方舟大洪水.
其實只不過是一個地區性的大洪水.
另外我們要了解.
羅亞方舟大洪水這個技術.
其實在聖經的頭十一章裡面找到.
創世記一至十一章.
其實帶領我們進入一個古代的世界.
叫做米索不達米亞.
米索不達米亞是囊括今天的東土二期山區.
敘利亞,伊拉克直到波士灣.
本身是一個很熱很乾的地方.
沒什麼可能成為人類定居的樂土.
但是米索不達米亞.
Mesopotamia是希臘文.
意思是說兩河流域.
在這裡有兩條大河.
第一條是底格里斯河.
第二條是幼發拉底河.

$^{241}$因為這兩條大河.
將整個地帶變成了人類定居的理想之地.
亦都將這個地方成為了.
世界人類最早的文明遙覽.
在這裡人類最初發展農業,捕魚,河水運輸,貿易.
甚至是契營民治,城邦和宗教等等.
河流當然帶來很多的好處.
但亦都帶來不同的危險.
每逢雨季或春天高山積雪融的時候.
就很容易有洪水.
所以米索不達米亞裡面的文獻,禮版.
包括聖經都經常講論這些大水災,大浩劫.
為什麼會有水災.
水災當中誰能夠幫助我們.
現在我們所知道的近東的神話,技術裡面.
提到有大洪水,水災,浩劫.
起碼有四,五個這些文獻.
所以如果這樣看.
聖經記載洪水.
就是跟近東的洪水神話做一個爭辯.
因為近東的洪水故事都是講什麼呢.
洪水發生的原因是因為人口太多.
噪音太大,觸犯了神靈.
以致神靈不能安睡.
它發怒的時候就降水災.
有一個近東神話叫做.
阿達哈斯史詩.
它提到洪水的原因是什麼呢.
其中有一句說.
大地變得太廣大,人民太多.
人的喧鬧聲太大了.
因為這個原因.
印尼這個大神靈感到很不安.
他就召集眾神來跟他們說.
人的喧鬧聲已經太大了.
他們的噪音使我心裡很不安.
因為他們的噪音,我無法入睡.
於是他就降下洪災.
來毀滅人類.
聖經就去爭辯.

$^{281}$說這是一個荒謬的原因.
一個真正的神.
怎會因為睡不著而發大怒.
無緣無故地殲滅人類.
聖經提出一個完全不同的詮釋.
是一個道德的解釋.
為什麼會有洪水.
《創世記大六章》.
凡有血氣的在地上都是敗壞了行為.
神就對羅亞說.
凡有血氣的盡頭已經來到我面前.
因為地上充滿了他們的強暴.
我要把他們和地一同毀滅.
其實你看《創世記》.
由第四章開始.
人已經背叛了上帝.
人去犯罪.
遠離上帝.
罪就進入世界.
然後強暴就繼續在神的世界裡蔓延.
去污衊上帝美好的創造.
開始有謀殺.
然後繼續有各項暴力.
甚至地上掌權的都滿有強暴.
強暴這個字是Hamas.
今天我們來形容恐怖分子.
帶來破壞.
這個字在聖經裡指的是傷害人的行為和態度.
是指社會裡的不公義.
就是說人類因為背叛上帝.
現在來到一個自我毀滅.
已經罪惡滔天.
上帝不得不施行公義的審判.
所以水災的原因.
有這樣的浩劫的原因.
是一個道德公義.
是個人社會不公義的問題.
羅亞如何得救呢?.
當然聖經說得很清楚.
羅亞蒙恩.

$^{321}$上帝給他恩典.
羅亞亦是這樣記載.
羅亞是個義人.
在當時世代是一個完全人.
與神同行.
這些字眼不是我們現代人說的.
是一個完美無罪的人.
一個義人.
所謂完全正直的人.
一個按照神的真理.
供你來做審判做決定.
也不欺壓身邊的弱者.
反而去行公義的一個人.
所以聖經對洪水的技術.
目的不是告訴你.
曾經有過一個水災.
或者給你一些細節.
以致你能夠重整重構.
當時歷史事件的水災.
因為洪水的記載.
並不是洪水事件的本身.
什麼意思呢?.
我舉個例子.
美國的911事件發生之後.
有不同的技術.
這個技術不是告訴你.
曾經有過911事件.
而在911事件裡面.
找出那個意義.
譬如這張很出名的911相片.
它不是告訴你.
曾經有911事件發生.
而是告訴你.
因為這件事的裡面.
我們怎樣去找回那種意義.
聖經的洪水記載.
也有相似的地方.
目的不是告訴你.
曾經有大洪水發生.
有,是人都知道.

$^{361}$最重要的就是告訴你.
在水災當中.
怎樣找到神學生命的意義.
在近東的神觀裡面.
神的權能是怎樣表達出來的呢?.
往往是神可以去掌管大水.
能夠控制混沌,混亂.
這樣表達出神的權能.
所以神創立世界的時候.
是將水分開.
然後使汗地露出來.
這就真正表達到神是沒有權能.
這是創世記第一章的記述.
不是科學報告.
而是神學的宣告.
所以為什麼洪水的故事.
在近東的文獻.
在聖經裡面時常都會出現.
不單單是因為那個地區.
經常會有洪水.
而是哪一個可以在洪水氾濫的時候.
仍然掌管一切來作王.
近東的洪水神話.
往往都會講到神靈.
在洪水氾濫的時候.
他們是不知所措的.
聖經對羅亞洪水故事的技術.
有一些取笑這些假神.
所以作者寫的是神學.
是生命的神學.
當你生命遇見混沌.
遇見災劫的時候.
遇見水災的時候.
到底你信靠的是哪一個?.
你看羅亞洪水的技術.
你就可以見到那種文學的修辭.
那種結構.
結構是怎樣的呢?.
首先七章十七節講到洪水氾濫四十天.
然後故事最後的時候.

$^{401}$第八章第六節.
就講到洪水過了四十天.
然後在七章十七節至十八節.
描述到洪水不斷上漲.
八章第五節.
相應呼應洪水繼續消減消退.
然後到第三個元素.
講到高山被水淹沒.
然後相對地第八章四至五節.
講到高山山頂露出來.
七章二十一節.
講到水勢浩大一百五十天.
相對應的另一個描述.
就是八章第三節.
就是水退去一百五十天.
而中間第八章第一節.
就講到神顧念羅亞.
聖經和近東文學都是相似.
經常有三文字的結構.
中間的元素是最重要的.
這裡八章第一節就是最重要的.
神紀念羅亞和羅亞方舟裡一切的走獸生蟲.
神紀念他的世界.
神去重整他的世界.
神紀念忠於他的約的人羅亞.
縱使在浩劫的當前.
他都不需要去怕.
(十二章 第十一節).
講起浩劫的時候.
可能人生想到很多不同的際遇.
可以是災病 可以是疫症.
可以是水災 可以是任何天災人禍.
這些都是突然其來.
沒有預示的.
我記起好幾年前.
我有一個這樣的經歷.
我一個人去到以色列.
我進入了巴勒斯坦的地區.
今天叫西岸西岸銀行.
我去探望一個巴勒斯坦的基督徒.

$^{441}$一個弟兄我和他相識.
所以我去到他山上的農場.
那天是春天 很大霧.
幾乎伸手不見五指.
我去到他的農場的時候.
我找不到他.
因為他剛剛離開了.
他太太在那個山上 在農場上.
和他傾談的時候.
我聽到很多的狗在狂叫.
我們望出去的時候.
發現有八九個蒙著臉的人.
他們拿著武器來搞事.
我馬上想到.
經常在報章和新聞上.
預見和看見的一些情況.
因為在一個戰爭區域.
一個戰時的區域裡面.
是沒有旺貫的.
很多時候有這些知識份子.
他們會做一些不法的事情.
他們拿著武器可以傷你.
可以殺你.
也不需要負任何的責任.
這個農場主人的太太也是我相識的.
一個女人就馬上衝出去.
對抗這些人.
我也跟著出去.
但在那一剎那.
我想到這是一個生死關頭.
我那時人生第一次這樣祈禱.
如果你要我走.
要我離開這裡.
我也可以離開.
因為我知道我會去哪裡.
我從來這麼大個人.
都沒有試過這樣去祈禱.
但因著耶穌給我的盼望.
祂給我的認許.
我很深信.

$^{481}$在那個關頭裡面.
即使有任何的不測.
即使我會喪生.
我仍然知道.
神是我生命的依靠.
我仍然是有平安.
我不知道你遇見浩劫的時候.
有沒有想起過.
神是跟你同在的.
神去紀念你.
還是你仍然是方方失失.
你完全不知道.
有什麼可以做來靠你.
讓今天這一段經文.
給你一個提醒.
投靠主的人.
上帝去紀念那些守約的人.
你是否一個守約的人.
是否忠於上帝對你的恩若.
是否忠於上帝給你的救恩.
以至你遇見浩劫的時候.
你都不需要懼怕.
因為你知道最終是神去紀念.
跟你同在.
我們同心得有祈禱.
在疫情裡面.
我們遇見不同的浩劫和威脅.
這些可能就是聖經裡面說的洪水.
這些天災人禍.
這些是我們人不能掌管的事.
但我們知道我們的生命在你的手中.
求你自己親自幫助我們.
忠於所約.
我們蒙受你的恩典的時候.
我們與你同行.
縱使洪水氾濫.
我們仍然知道你是掌管一切.
你也是掌管我們的生命.
縱使浩劫當前.
我們仍然最終不怕.

$^{521}$因為我們的生命在你的手中.
求你讓我們再一次來到你.
投靠你.
在我們生命遇見很多不測的事.
我們仍然與你同行.
我們這樣祈禱.
奉主耶穌基督的名字.
接下來一首回應詩歌.
我希望這首詩歌的幾分鐘時間.
能夠跟著歌詞.
洗上帝再一次成為你生命當中的頭銬.
洗上帝成為你的主宰.
Zither Harp.
投靠你至旁陰霞.
在你光中不必驚怕.
唯獨竟為你的方可得依靠.
別再自以來轉向你.
投靠你至旁陰霞.
在你光中不必驚怕.
唯獨竟為你的方可得依靠.
別再自以來轉向你.
洪水即將泛濫之時.
是我主佐著為我.
願天歡喜.
願地上快樂.
如影指指上騰.
等候我主的必不受愧.
是你今臨如路.
只讓也在.
Zither Harp.
投靠你至旁陰霞.
在你光中不必驚怕.
唯獨竟為你的方可得依靠.
別再自以來轉向你.
洪水即將泛濫之時.
是我主佐著為我.
願天歡喜.
願地上快樂.
如影指指上騰.
等候我主的必不受愧.

$^{561}$是你今臨如路.
只讓也在.
洪水即將泛濫之時.
是我主佐著為我.
願天歡喜.
願地上快樂.
如影指指上騰.
等候我主的必不受愧.
是你今臨如路.
只讓也在.
是你今臨如路.
只讓也在.
Zither Harp.
\newpage



\section{}
\label{sec:9fgDv4ygGRU}
\textbf{神的名字1\_Elohim (靈思牧禱)}
\newline
\newline
連結: \href{https://youtube.com/watch?v=9fgDv4ygGRU}{\texttt{https://youtube.com/watch?v=9fgDv4ygGRU}} ~~~~ 語音日期: 2020-09-30
\newline
\newline
\hyperref[sec:o6x_Tz0d6Ww]{\small{< < < PREV SERMON < < <}}
~
\hyperref[sec:index_chronic]{\small{[返順時目]}}
~
\hyperref[sec:index_scriptual]{\small{[返順卷目]}}
~
\hyperref[sec:_8xftiiwoRs]{\small{> > > NEXT SERMON > > >}}
\newline
\newline
$^{1}$當準備這個新的靈修系列的時候.
我被這個美麗的教堂彩色玻璃所吸引.
它使我想像到神的光照進它的電.
照亮電裡面每一位屬於它的子民.
包括你和我.
特別是充滿衝擊.
甚至好像很沒有希望的時候.
基督教和其它宗教一樣最大的不同.
就是我們所認識的神.
並不是高高在上,冷漠無情的.
相反它一次又一次屈尊降卑地來到人世間.
尤其是當它的子民走在最幽暗的日子的時候.
他們體會到它的同在.
而不禁驚嘆呼喊它是一位怎樣的神.
而這些驚嘆呼喊.
就成為聖經裡面所說的神的名字.
在這個靈修系列裡面.
我想和大家從舊約聖經去看一些神的名字.
沒錯 其中包括我們以上所說的驚嘆呼喊.
亦都包括了神自己曉喻祂自己是一個怎樣的神.
總之每一個名字都有祂獨特而深入的意思.
而這些名字我們不能輕忽.
而且當我們牢牢記住的時候.
更加會成為我們心中很大的力量.
今天要和大家說的第一個名字就是阿羅卿.
阿羅卿在舊約聖經裡面出現.
有大概2600次.
出現這麼多次.
究竟神想我們從這個名字認識祂甚麼呢?.
首先L 原文的字根是強而有力的意思.
聽起來很威猛.
但其實不是那麼獨一無二.
為甚麼這樣說呢?.
因為當時以色列周圍的國家的神都自稱是L.
都是強而有力的.
但特別的是以色列的神自稱為阿羅卿.
要強調祂自己有兩個獨特的地方.
第一 L是用大寫寫的.
指明祂是眾神的神.
是獨一的神.

$^{41}$是凌駕一切的神明的.
而在《創世記》第一章第一節裡說.
「起初神創造天地」.
這個神我們中文時時譯為上帝.
是告訴我們.
唯獨這個上帝是創造天地萬物.
其他神明沒有這種能力.
第二個獨特之處.
阿羅卿是一個複數字.
《聖經學制》指出.
原來這個複數是指明神是獨一的神.
更加是三位一體的神.
即是我們今天所認識的.
聖父,聖子,聖靈.
這個奧妙的道理是貫穿整本《聖經》.
總的來說.
阿羅卿這個名字強調上帝是創造天地萬物.
獨一的真神.
而祂的三位一體.
體現了在祂創造天地.
救贖因罪而破碎的這個世界.
和引導祂子民進入真理的裡面.
面對今天動盪不安的世界.
面對生命中的狂風巨浪.
能夠知道我們的神.
是一位這樣的神.
對你有甚麼提醒呢?.
歡迎你現在在Live chat.
或者在九一的群組裡.
表達你即時的感受.
認識神的名字.
會使我們屈膝跪下.
敬拜祂威嚴的聖潔.
認識神的名字.
會讓我們得到極大的安慰和確據.
認識神的名字.
更加使我們知道祂是一位怎樣的神.
而我可以向你肯定.
當你越認識祂的時候.
你會越感激祂.

$^{81}$敬畏祂.
你的生命會更加得到很大很大的激勵.
現在我們一起祈禱吧.
創造天地萬物獨一的真神.
我們齊集在此.
獻上我們的敬拜.
聖父,聖子,聖靈.
和在我們生命裡面賜愛的引導.
求你赦免我們.
雖然我們認你為主.
但很多時候因為本性的軟弱.
轉去敬拜其他的神.
求聖靈在我們生命裡面.
時時提醒,激勵.
使我們行事為人.
能夠尊你的名為聖.
將你當得的榮耀歸給你.
我們先求你的國和你的義.
願你的旨意成全.
奉耶穌基督的名禱告.
阿們.
\newpage



\section{}
\label{sec:_8xftiiwoRs}
\textbf{浩劫重生 【第2集:洪水之後】 陳崇基牧師 - 主日講道視頻系列}
\newline
\newline
連結: \href{https://youtube.com/watch?v=-8xftiiwoRs}{\texttt{https://youtube.com/watch?v=-8xftiiwoRs}} ~~~~ 語音日期: 2020-10-05
\newline
\newline
\hyperref[sec:9fgDv4ygGRU]{\small{< < < PREV SERMON < < <}}
~
\hyperref[sec:index_chronic]{\small{[返順時目]}}
~
\hyperref[sec:index_scriptual]{\small{[返順卷目]}}
~
\hyperref[sec:fE8jDxIIYZM]{\small{> > > NEXT SERMON > > >}}
\newline
\newline
$^{1}$上一次我們去看羅亞方舟洪水浩劫的時候.
我們說按照古近東的背景.
舊約的經文.
考古地質等等的線索.
似乎沒有一個普世全球的洪水曾經發生過.
但也有很多人繼續問.
這些記載的背後.
是否有一個真實的歷史事蹟.
或者有一個真實的.
很大的水災曾經發生過.
你可能聽聞過一個理論.
叫做黑海洪水假設.
黑海是位於土耳其北部.
這個理論這個假設認為.
以前的黑海其實是一個淡水湖.
或者我們可以叫做黑湖.
湖的淡水.
是經過一個很窄的峽谷.
叫做博斯普拉斯海峽.
流入去.
地中海.
但大約公元前5000年的時候.
當時全球的水位因為冰川融化.
海洋的水位.
就突然間升高.
就造成一個大的洪水.
突然間地中海的海水.
就倒轉流回.
現在的黑海.
估計當時.
那個流入湧入的水力.
是尼加拉瓜瀑布的400倍.
所以短短的一年裡面.
那個本身是淡水湖的黑湖.
就變成了今天的黑海.
當時的水位上漲了幾百呎.
湖的面積.
突然增加了30\%.
以致淹沒了很多湖岸四周的古城鎮.
現在考古學家也都在那裡做一些深海的探索工作.

$^{41}$在黑海的海底裡面找到很多古代的遺跡.
所以有學者認為.
這場黑海的大浩劫大水災.
就成為了中東一個很大的災難死傷無數.
漸漸成為古近東每個洪水故事的背景.
當然我們不能夠去肯定這個黑海的論說.
就是.
洪水背後的歷史事件.
正如上一集我們所說.
洪水的記載.
不是歷史事件的本身.
而是一個.
對於那個歷史事件.
的文學神學的解釋.
事實上你可以想像一下.
為什麼一個.
水災.
會對近東的人.
聖經時代的人.
成為一個這麼大的震撼.
除了水災會構成很大的破壞之外.
因為.
近東的世界觀是一個平面的世界觀.
他們在一個平面的地上生活.
上面一個窮倉.
他們的想像就是這個窮倉頂住外面的大水.
讓它不會去消滅這個世界.
但當黑雨暴雨不停.
或者河水氾濫的時候.
對於一個近東人來說.
就好像天塌下來.
世界好像沒東西一樣.
或者整個世界.
好像回到未創造之前那種.
只有大水空虛混沌的情境.
如果留意聖經如何描述.
羅亞洪水這個事蹟.
你會發現有很多字眼.
是形容創造天地改變天地的.
在羅亞洪水這個技術裡面.

$^{81}$我們看看一些字眼.
發現原來在創世記第一章出現過.
比如創世記第一章第一天的創造.
第一章第二節.
地是空虛混沌.
淵面黑暗.
神的靈.
原文是神的風.
混行在水面上.
講到羅亞洪水的故事.
洪水開始停了.
八章第一節.
神紀念羅亞和方舟裡面一切的走獸生俗.
神叫風.
原文是靈.
吹在地上.
水勢漸落.
淵源.
天上的窗戶閉塞.
天上的大雨.
也止住.
接著我們看第一章第八節.
神稱空氣為天.
有晚上有早晨.
這是第二天的創造.
講到羅亞洪水的技術.
八章第二節.
淵源和天上的窗戶都閉塞了.
天上的大雨也止住了.
創世記第三天的創造.
一章第九節.
神說天下的水要聚在一處.
使寒地露出來.
事就這樣成了.
講到羅亞洪水的技術.
都是一樣.
洪水漸退了.
八章第三節.
從地上漸退的時候過了一百五十天.
水就漸消.

$^{121}$七月十七日.
方舟亭在阿拉納山上.
水又漸退.
到十月初一日.
山頂就現出來.
是同一個字.
露出來.
現出來.
創世記第一章講到神創造.
第五天的時候發生什麼事呢.
一章第二十二節.
神賜福給這一切.
有雀鳥生在地上.
八章第六節.
講到羅亞洪水漸消的時候.
在方舟那裡.
他把雀鳥放出來.
首先是烏鴉.
然後在地上飛.
甲子.
然後也在地上飛.
然後第三次也是放甲子出去.
再看下去.
第六天的創造.
跟羅亞的洪水技術.
就更加相似了.
一章第二十六至二十八節.
大家很清楚了.
就是神造人的時候.
他是按照神的形象所造.
然後講到羅亞的洪水完結之後.
神跟羅亞和大地立約.
第九章第六節.
他說人是按照神的形象所造.
一章二十八節.
神賜福給人類.
叫他們生養眾多.
要治理者地.
要偏門者地.
到第九章講到羅亞洪水之後.

$^{161}$神也賜福給羅亞和他的兒子.
叫他們生養眾多.
要偏門全地.
創世記第一章講到創造.
羅亞洪水氾濫的時候.
就是一種反創造.
或者de-creation.
有人會說un-creation.
洪水過了之後.
就是一個再創造.
是recreation.
神用洪水這個號節.
來再造天地.
其實在舊約裡面.
經常有一種創造.
反創造.
再創造.
這樣的方式來描述.
號節災難.
譬如猶太人亡國.
戰敗的時候.
都是用這樣的方法.
來描述.
視乎你身處不同的地方.
一些大的號節.
不一定是水災.
可能是疫症.
可能是絕症.
可能是一些發生在我們身上的意外.
或者一些生命裡面的失去.
其實每一類的號節.
在某程度上.
其實都是一種de-creation.
然後再recreation.
剛疫情開始的時候.
有人說新冠病毒.
是要逼整個人類.
來面對一個很不想面對的問題.
就是死亡.
就是我們人.

$^{201}$是不堪一擊的.
不一定是新冠病毒.
以致我們驚慌.
以致被迫要面對死亡.
要面對病症.
在生命裡面.
其實我們有很多事.
有不同的事情會打擊我們.
可能是一些情緒的問題.
可能是一些健康的問題.
可能是事業等等.
我們面對著.
生命裡面一些我們不能夠去管理.
不能夠去控制的因素.
我們實在是非常脆弱.
在這個星期.
2020年10月初.
一直將新冠病毒.
去輕言的美國總統特朗普.
連他自己都感染了.
所以這句說話說得很對.
這個病毒的大流行.
是人人都受影響.
是不分階級.
不分季節.
我們每個人都需要面對.
人生的脆弱.
面對死亡.
疫症可能就帶給我們.
一個de-creation的機會.
就去拆毀我們以往.
以為我們是很堅強.
來到再重新建立.
重新re-create我們.
以致我們知道.
原來我們生命裡面.
是需要面對信仰的.
面對上帝創造主的.
所以當有浩劫.
在我們生命裡面發生的時候.

$^{241}$或者我們可能要問這個問題.
就是上帝在我們生命當中.
有沒有做一些un-creation.
de-creation.
去到重新再整頓.
重整我們的人生.
以致上帝能夠再一次創造.
再一次在我們人生的觀念上.
人生的方向目的上.
再一次重新去塑造.
今年2020年.
有一個學者.
拿了一個獎叫做登普頓獎.
Templeton這個獎.
是專門頒發給一些.
在宗教信仰和科學之間.
有很大貢獻的人.
今年2020年得獎者.
是美國國家衛生院主任.
Francis Collins.
Francis Collins可能你聽過這個名字.
他是現在全世界最頂尖的人類遺傳基因學者.
當他還是一個年輕的醫生的時候.
他是一個無神論者.
他就被委派去照顧一個.
患有嚴重心臟病的年老婦人.
雖然這個女士每一天心胸會非常痛.
是難以置信的.
但是她絕對以她的信仰.
耶穌基督為她的中心.
這位女士和Collins分享過好幾次福音.
有一次她終於問這個醫生.
她問Francis Collins.
我和你說了好幾次福音.
你還是避而不談.
我想問你醫生.
到底你信的是什麼?.
What do you believe?.
就是因為這四個字.
令Francis Collins突然改變他人生的方向.

$^{281}$這四個字喚醒他去想想.
到底他人生的觀念.
他人生的信仰.
他到底是信什麼?.
當時信無神論的他.
他突然有一個生命的倒轉重整.
之後他就成為了一個基督徒.
甚至他現在致力於結合信仰和科學理性.
有一個網站叫做Biologos.
它就是發起人.
其實Biologos想做的.
就試探討聖經宗教信仰和科學.
如何連接.
Francis Collins的得獎致詞裡.
他說科學在人生的終極生死上.
特別當我們遇見浩劫的時候.
其實沒有什麼用處.
因為在這個領域裡.
是宗教信仰的領域.
是宗教信仰給我們的智慧.
所以在他得獎的致詞裡.
With all your mind.
盡你的理性去尋求真理.
去尋求神.
如果你有時間的話.
我希望你去看一看他的致詞.
因為四個字.
英文是What do you believe?.
就是你所信仰的是什麼.
令到Francis Collins在他人生裡.
有一個大的逆轉.
我們在看浩劫.
我們當然要想一想.
神藉著不同的災難.
不同的浩劫.
為我們帶來一種什麼的de-creation.
然後有re-creation.
給我們一個翻天覆地的人生方向的概念.
你可能還沒有信主.
你看見這個訊息的時候.

$^{321}$我想問你.
What do you believe?.
你信的是什麼?.
如果你是已經信了耶穌的基督徒.
其實這個時候.
也是很認真的要思考.
去想清楚.
上帝給我們生命當中有浩劫.
有疫情.
到底我們所信的是什麼?.
What do you believe?.
你是否真正去嚴肅的.
認真的去對待你的信念?.
如果你再細心看看羅亞洪水這個故事.
你會發現有些細節.
我們很少留意.
有幾處地方提到日期.
以色列人的農曆.
是他們的月曆.
我們看看這五段聖經.
第一段七章第十一節.
羅亞六百歲的時候.
二月十七日那一天.
大淵的泉源裂開.
窗戶打開了.
四十晝夜有大雨降在地上.
正當那一天羅亞就入放晝.
這是洪水的開始.
第二段又提到二月.
八章第十四節.
二月二十七日地乾了.
洪水就完結了.
即是洪水氾濫在地上.
有一年零十日.
為什麼這麼特別?.
在摩西五經裡面提到二月發生的事.
只有另外兩處.
第一處就是出埃及記第二十六章.
那裡提到二月十五日.
以色列人傳會中.

$^{361}$他們出埃及在曠野的樹上.
他們發怨言.
然後神首次在那一天.
來到去恭應瓦剌.
二月還發生什麼事呢?.
在民數記第二年二月二十日.
雲彩從法季的帳幕收上來.
以色列人就按照前面的站.
來到前行.
離開曠野.
雲彩就停在前面巴蘭的曠野.
換句話說.
上帝在第二年二月的時候.
神願意和他們的子民.
帶著他們在曠野的樹上同行.
二月特別的地方就是.
原來在以色列巴勒斯坦地.
春天兩個修成的時候.
大麥和小麥中間修成的時候.
一月和三月中間就是二月.
沒有什麼特別的宗教節日.
但是卻是提醒神的子民.
神在一年裡面.
兩個修割節中間.
來到去恭迎他們的需要.
洪水故事的記載.
就是告訴方舟上的羅亞一家.
代表以色列的子民.
在浩劫的當中.
神仍然和他們同行.
去保守羅亞和他們一家.
也保守方舟上不同的動物.
代表整個大地.
說到春天修割節.
二月的時候.
另一個在巴勒斯坦地.
經常出現的現象.
就是雨後彩虹.
因為春天是雨季末後的時候.
如果你有去過以色列地的話.

$^{401}$下雨有彩虹是經常出現的.
請你不要搞錯.
這裡不是說人類歷史裡面.
第一次出現彩虹.
一向都有雨後彩虹.
但是洪水之後.
神就用這個經常見到的彩虹.
叫人去看見這個彩虹的時候.
他明白神和大地.
神和他們的子民重新立約.
這是神學的意義.
洪水災劫.
神仍然保守他們的子民.
保守他們的創造.
第三段聖經就是第八章第四節.
七月十七日.
方舟停在阿拉納山上.
按照以色列的律法.
七月就很特別.
因為七月十五日.
就是鑄棒節.
鑄棒節是慶祝什麼的呢?.
是有一個星期.
紀念以色列人曾經在曠野.
神和他們在曠野漂流四十年.
和他們同行.
所以現在到Sukkot的時候.
以色列人就會在外面搭橡棚.
紀念神的看顧和同在.
今天這個視頻訊息.
是2020年十月四日.
上星期五的晚上.
就是以色列Sukkot鑄棒節的第一天開始.
我們華人有中秋.
對於猶太人來說.
就是七月鑄棒節的時候.
七月十七日正是鑄棒節.
慶祝的一個星期的中間.
羅亞的故事說.
方舟就在這一天停下來.

$^{441}$停載這個字.
有安全休息的意思.
即是不用顫抖,不用顫抖.
不用那麼害怕.
就是在鑄棒節的中間.
去到慶祝上帝的保守.
紀念上帝以往去到同行的時候.
他們就想起原來在上帝裡面.
我們能夠有平安.
其實羅亞這個名字.
Noah 也是這個意思.
即是說有安息的意思.
第四段八章第五節.
這裡說水遊漸退.
十月初一日.
山頂都現出來.
羅亞就開了方舟的窗戶.
然後放甲子出去.
最後甲子就拿了一個新的.
剛剛生出來的橄欖葉子.
羅亞就知道地上的水退去了.
這裡說十月.
十月是以色列龍更的播種期.
三個月之前即是七月.
就是慶祝秋天的收割.
三個月之後.
農夫就盼望有春季的收割.
農夫就會在十月的時候.
盡他的能力來播種.
播了種之後.
他就很著緊.
希望看到農作物發芽.
所以如果你想像一下一個農夫.
播了種之後.
他就每一天會看著田地.
希望發芽能夠生長.
或者這樣問.
上帝會不會再次大地一個新的生命.
打開窗戶的時候.
看到什麼?.

$^{481}$山頂出來.
他打開窗戶然後將白甲放出去.
白甲就將新生出來的橄欖葉子帶回來.
換句話說是一個應許神去更新.
整頓過他的大地.
第五段是記載有日期的.
就是八章第十三節.
就是正月初一日.
地上的水乾了.
羅亞就撤去方舟的觀海.
看到地是乾的.
正月一日.
當然就是新年最大的大年初一.
以色列的農民和人民在正月初一.
這個時候就是以色列的雨季尾聲.
即是雨水漸退.
田地開始變成綠色.
看到有新的盼望.
在五經裡面看到正月初一的時候.
就會說到新的盼望.
新的生命.
我們人民的生命需要更新.
出埃及記十二章.
說到出埃及的第一日.
就是正月初一.
就是以色列人有新的開始.
出埃及記四十章.
也說到第二年的初一.
那裡就是會幕的建成.
上帝的榮耀降臨.
敬拜的生命開始.
民數記二十章也提到正月初一.
那裡說的是第一代漸漸要過去.
他們叛逆,不順.
漸漸死在曠野.
第二代就要興起.
他們要憑著信心做一個典範.
去邁向迦南地.
新年就是一個更新的時候.
很明顯聖經的作者不是死板的.

$^{521}$好像記者那樣.
放下這些日期.
他是特意將洪水浩劫的經過.
套入以色列人的龍耕.
龍曆的循環裡面.
告訴人們.
羅亞的水災時間.
一年五個日期.
要想起上帝看顧的報仇.
上帝的恩典.
上帝的約沒有消減過.
洪水的故事更加套入.
以色列人每年有的.
敬拜首節的月.
那個Liturgical Calendar.
整個羅亞的故事.
水災就是按照.
以色列人敬拜首節的月曆進行.
其實記載的故事就是教導後世的人.
縱使有一年之長的浩劫.
在我們敬拜首節的時間.
我們不要忘記上帝的約.
逾越節,新年,收割節,鑄棒節等等.
不要忘記上帝對我們的看顧.
不要忘記體會到.
上帝的實在和恩典.
同時你和我都要守約.
如果有人問你.
在你人生當中.
有什麼日期是最深印象.
可能你會說是喜慶的事.
但可能仔細想一想.
你一定會想起.
在你人生中的一些低谷.
一些日期.
你遇見一些挫折.
無論是感情,事業,身體健康,意外等等.
或者某年某月某日.
你很記得你收到的電話.
是醫生打給你的.

$^{561}$是一些不好的消息.
一些你個人的浩劫.
你一定會想起這些日期.
或者對全人類來說.
2020年是一個浩劫災難之年.
是我們一生都會記得這一年.
疫情到現在.
至少都有一年.
但你要知道.
這個羅亞的故事.
神學的意義.
就是告訴我們.
上帝往往是用浩劫.
來敲我們生命的門.
來拆毀我們.
去複製我們.
以致我們能夠再重新被創造.
以致我們的生命能夠被更新.
想清楚你所信的是什麼.
如果我們所信這位上帝.
是一位守約的上帝的話.
我們仍然會真正的跟隨他.
我們同心低頭祈禱.
因為你是那一位守約的神.
縱使在浩劫當中.
在我們人生的挫折當中.
你仍然願意與我們同行.
主我們求你.
你的恩典再一次在浩劫當中.
在我們人生的低谷當中.
拆毀重建.
拆毀又再重建.
幫助我們認認真真地面對你.
認真地追求你.
叫你自己的子民.
叫你自己的兒女.
能夠更加體會你自己的實在.
我們這樣祈禱.
奉主耶穌基督的名字.
孤單懼怕的心.

$^{601}$今天願教給主你.
是你的恩典與愛.
讓我可得到安慰.
憂傷可遂的心.
今天願教主手裡.
是你的變相與痛.
讓我可得已痊癒.
醫治我傷.
在你的痛災裡.
讓我痛與你為裡.
我心知道你是我神.
恢復我心.
在你的永遙裡.
讓我得見你完美.
我的身處.
一得已否認.
孤單懼怕的心.
今天願教給主你.
是你的恩典與愛.
讓我可得到安慰.
憂傷可遂的心.
今天願教主手裡.
是你的變相與痛.
讓我可得已痊癒.
醫治我傷.
在你的痛災裡.
讓我痛與你為裡.
我心知道你是我神.
恢復我心.
在你的永遙裡.
讓我得見你完美.
我的身處.
全然康復的你.
我靈數聲.
要轉你命為星.
我靈歌唱.
要頌讚你的名.
班容此生.
成就你的志.
我願為你再奔跑.

$^{641}$我靈數聲.
要轉你命為星.
我靈歌唱.
要頌讚你的名.
班容此生.
成就你的志.
我願為你再奔跑.
醫治我傷.
在你的痛災裡.
讓我痛與你為裡.
被你的愛.
捉摸我靈.
恢復我心.
讓我如影輾轉.
讓我得見你完美.
我的身處.
全然康復與你.
讓我痛與你為裡.
讓我得見你完美.
我想知道.
你是我的孫.
Zither Harp.
\newpage



\section{}
\label{sec:fE8jDxIIYZM}
\textbf{Kingdom Values XXIV: Matthew Chapters 5 - Pastor Sam Johnson}
\newline
\newline
連結: \href{https://youtube.com/watch?v=fE8jDxIIYZM}{\texttt{https://youtube.com/watch?v=fE8jDxIIYZM}} ~~~~ 語音日期: 2020-10-06
\newline
\newline
\hyperref[sec:_8xftiiwoRs]{\small{< < < PREV SERMON < < <}}
~
\hyperref[sec:index_chronic]{\small{[返順時目]}}
~
\hyperref[sec:index_scriptual]{\small{[返順卷目]}}
~
\hyperref[sec:x6dSk6R1byI]{\small{> > > NEXT SERMON > > >}}
\newline
\newline
$^{1}$Good morning family. Today, after a few months in Matthew chapter five,.
we are finishing out this chapter. Now it's not the end of the sermon on the.
mount. There's two more chapters of this sermon to go..
And so we're going to continue on with looking at the kingdom values that are.
found in this sermon. But we're closing out this chapter today..
And Jesus closes out this chapter with this amazing statement that he makes..
So we're going to look at in Matthew chapter five, verse 48..
It says this,.
you therefore must be perfect as your heavenly father is perfect..
Okay, so wait just a minute here. I want to make sure we read this right..
It sounds like this verse is saying that we're supposed to be perfect..
Like God is perfect. We're supposed to imitate God in holiness..
We can look at it again and that's pretty much exactly what it says..
You therefore must be perfect as your heavenly father is perfect..
Now what is going on here? This is an impossible command..
This is an impossible ask. This is an impossible request..
So what is Jesus saying in this verse?.
I want us to dig into this passage today and examine what Jesus is saying to us.
when he calls us to be perfect in the way that God is perfect..
First of all,.
this word for perfect in Greek is a word that means complete and lacking nothing.
or in need of nothing and it could also mean mature,.
but it also does mean perfect. So when we read it, we could say, well,.
maybe God's just calling us to be mature..
Maybe God's just calling us to be complete,.
totally equipped to serve him and follow him. Yeah, that's a part of it..
That's there, but we can't escape the straightforward reading of this verse..
It says we're supposed to be lighting holy lives and that is a really,.
really challenging statement for us..
And as we look at the passage that surrounds it,.
we see that it gets even more challenging..
Remember when we read verse 20 it says, for I tell you,.
unless your righteousness exceeds that of the scribes and Pharisees,.
you will never enter the kingdom of heaven..
That was the beginning of this section where Jesus talks about the old Testament.
law and our relationship to it. And in all these ways,.
when Jesus talks about our anger, when Jesus talks about our lust,.
when he talks about how we break our promises,.
when he talks about how we retaliate and take revenge on one another and all of.
these things that Jesus has listed for us in this section,.

$^{41}$he's demonstrated how we fall so short,.
how our righteousness is so short of God's standard..
It certainly doesn't exceed the righteousness of the scribes and Pharisees..
And those two statements,.
verse 20 and verse 48 that we read this morning kind of bracket this section..
Jesus starts the section by saying, listen,.
if you think you want to be good enough to please God,.
then your righteousness has to exceed the righteousness of the scribes and.
Pharisees. As we talked about before,.
he's basically saying your righteousness has to exceed an external conformance..
It has to be more than something superficial..
It has to be more than something that is just outward and just for other people.
to see. It has to be something real..
And then he closes out this section with verse 48,.
the last verse of chapter five where he says,.
you therefore must be perfect as your heavenly father is perfect..
So Jesus is setting the bar impossibly high for us..
He's setting the standard impossibly high for us,.
but he's doing that to help us see how much we need God's work in our life and.
to help us see that this is impossible without God's work in our heart and in.
our lives. Notice in this verse,.
Jesus doesn't call us to just be a little bit better than other people..
Look at verse 47 and verse 46. He says,.
for if you love those who love you, what reward do you have?.
Do not even the tax collectors do the same? And if you greet only your brothers,.
what more are you doing than others? Do not even the Gentiles do the same..
Jesus isn't saying, Hey guys,.
if you can just do a little better than the tax collectors,.
if you can just do a little better than the Gentiles,.
if you can just do a little better than the people who don't know God and the.
people who don't love God, if you could just do a little better than them,.
then you're doing great. No, he says,.
I want you to be perfect in the way that God is perfect..
That's the standard. Jesus also doesn't say, Hey,.
I just want you to be a little bit better than you were last week..
I want you to just improve a little bit and just become a little bit better than.
you were. If you can do a little bit better job than you were last week,.
then that's what I want from you as a, as a subject of the kingdom,.
as a disciple, as a, as a follower of Jesus,.
just be a little bit better than you were..

$^{81}$But Jesus doesn't give us that either. Jesus says, no,.
I want you to be like God. I want you to be like me..
I want you to live out the values of the kingdom and I want you to live them out.
in every area of your life. That's what Jesus calls us to..
Now that's something that we cannot do. It's impossible..
And the fact that it's impossible and the fact that Jesus calls us to it,.
even though it is impossible points to what Jesus wants for our lives..
See, Jesus doesn't just want to improve us..
He wants to live through us..
And that's a central truth of Christianity..
That's a central value of the kingdom..
It's a central reality of following Jesus is that Jesus doesn't just want to.
improve us. He wants to live through us..
A lot of times we get this wrong idea about our faith,.
about being a Christian and about the Bible..
And we think that the Bible is kind of like a self help manual and it's a self.
improvement book and it's there to help us be a little bit better version of.
ourselves. It's there to turn us into me 2.0..
That's a little bit better than I was last week or last year, whatever..
Well, there's a little bit of truth in that. Yes,.
God is working in our life over time. Yes, it's a gradual process. Yes,.
we grow in small steps. And so it's a week at a time. It's a day at a time..
It's a year at a time. Yes, that's true..
But the reality is Jesus teaches us clearly in the scripture..
He doesn't want us to be me 2.0..
He doesn't want us to be a slightly improved version of ourselves..
That's not what he's trying to grow us towards..
He is trying to grow us towards his life in us..
He wants us to go away..
He wants us to disappear and he wants to live his life through us..
In Matthew chapter 16 verses 24 and 25 it says this,.
then Jesus told his disciples, if anyone would come after me,.
let him deny himself and take up his cross and follow me..
For whoever would save his life will lose it..
But whoever loses his life for my sake will find it..
And when we hear the word cross, we think of Jesus's death on the cross..
And we think of something beautiful. We think of our redemption..
But when Jesus said this to the disciples, the cross had not yet happened..
So they didn't think of redemption. When they heard the word cross,.
they thought of what it meant to them..

$^{121}$And to them it was an instrument of torture and death..
They had seen people crucified by the Romans. They knew what the cross meant..
And so when Jesus said, take up your cross to them,.
they understood that he was saying, listen,.
I want for your life to disappear..
I want for those natural things, those sinful things,.
all those old pieces of you to die and to be gone so that I can live through you..
That's what Jesus is calling us to. He doesn't just want to improve us..
He wants to live through us. He wants all that old sinful stuff,.
everything from verse 20 down to verse 48 that Jesus called out in our lives,.
the anger, the hatred, the lying, the broken promises,.
the lust, the retribution and revenge, and the way we treat other people..
All of that. Jesus is calling all that stuff out and he says, Hey,.
I want that stuff to be gone. I want that stuff to die..
I want it to encounter the cross,.
to encounter the instrument of death so that I can begin to live through you in.
the place of all your selfishness and all your sin and all your brokenness and.
all the ugly stuff,.
all the ways that you hurt one another and all the ways that you reject God..
I want those things to be gone so that I can live through you..
And that's what Jesus is pointing us to in this passage..
He's helping us to see that our righteousness has to exceed this righteousness.
of the scribes and Pharisees and that he wants us to live and in a way that is.
imitating God, that is imitating Christ..
It's not enough just to try to do a little bit better..
It's not enough just to try to be better than someone else..
Jesus wants us to completely disappear and he wants to live through us..
God wants to see less and less of me every day and more and more of Jesus every.
day. And in your life, God wants to see less and less of you every day..
Less and less of your pride. Less and less of your selfishness,.
less and less of your anger, less and less of your lust..
And he wants to see more and more of Jesus in your life every day..
More of Jesus's love, more of Jesus's care for others,.
more of Jesus's selflessness in your life,.
more of Jesus's self sacrificing love..
That's what God wants in my life and in your life..
Now, how does that happen?.
If we say, okay, that's what Jesus is telling us in this passage..
How does that happen? When we see this standard,.
when Jesus sets the standard impossibly high,.

$^{161}$there's a few ways that we can react to it..
One is to go to a place of despair and say,.
I can't meet that standard. I'm a bad person..
I'm a failure and I'm just going to live in despair..
That's not what God has for us..
Jesus doesn't put the standard out there to make us feel bad and to discourage.
us and to put us in a place of despair..
The other thing that we often do,.
and this is what Jesus actually talks about in Matthew chapter six,.
this is the last verse of chapter five..
But when we start talking about chapter six, next week, we'll see this..
Jesus begins chapter six by talking about hypocrisy..
And that's the other response. When we see how high the standard is,.
sometimes we respond by saying, well,.
I'm going to go ahead and pretend that I meet the standard, even though I don't,.
I'm going to try to just do my best and act like I do and tell other people that.
I do. And when we do that, we become a hypocrite..
We start acting like we're something we're not..
We start acting like we're, we're perfect and we never could be perfect..
That's hypocrisy..
But the way Jesus wants us to respond to this is through surrender to him,.
through saying, God, I can't live this out..
I can't exceed the righteousness of the scribes and Pharisees. God,.
there's certainly no way I can be perfect as my heavenly father is perfect..
I cannot do this. And so I'm going to let myself disappear..
I'm going to take all those old, selfish,.
sinful parts of me to the cross and let them die there and surrender to Jesus so.
that he can live through me..
And surrender is a choice..
I said that last week that nothing changes until you make a choice and we have.
to choose to surrender. We have to choose to say, God,.
I'm going to stop trying. I'm going to stop saying,.
maybe if I just try a little bit harder,.
maybe if I just figure out a few more things about the Bible,.
maybe if I just wrap my head around it or maybe if I just try to be a better.
person or maybe if I just work harder at hiding my sin from other people..
No..
When we stop all of that stuff and say, I'm going to surrender to you..
God, I can't do this..
I'm asking you to come and to live through me and to do this in my life..

$^{201}$I really can't think of a better way to express what this surrender looks like.
than the words of Jesus in another sermon that he gave..
And this is in the book of John in chapter 15..
Jesus is talking about abiding in him,.
surrendering to him, giving up so that he lives through us..
And he's talking about the peace and the rest and the joy that we find in that..
So I'm going to read this. This is John chapter 15 verses one to 11..
Jesus says, I am the true vine and my father is the vine dresser..
Every branch in me that does not bear fruit, he takes away..
And every branch that does bear fruit, he prunes that it may bear more fruits..
Already you are clean because of the word that I have spoken..
You abide in me and I in you as the branch cannot bear fruit by itself unless it.
abides in the vine. Neither can you, unless you abide in me. I am the vine..
You are the branches. Whoever abides in me and I in him,.
he it is that bears much fruit for apart from me, you can do nothing..
If anyone does not abide in me,.
he is thrown away like a branch and withers and the branches are gathered,.
thrown into the fire and burned. If you abide in me and my words abide in you,.
ask whatever you wish and it will be done for you. By this,.
my father is glorified that you bear much fruit and so prove to be my disciples..
As the father has loved me, so have I loved you abide in my love..
If you keep my commandments,.
you will abide in my love just as I have kept my father's commandments and abide.
in his love. These things I have spoken to you that my joy may be in you and that.
your joy may be full..
What a beautiful description of what it looks like to surrender to Jesus..
And Jesus is clear here in verses three and four. He's talking to his followers..
He's talking to his disciples. He's talking to the people of God,.
the subjects of the kingdom, those who already have relationship with him..
He's talking to you and me. We see that in verse three. He says,.
already you are clean because of the word that I have spoken to you..
He's talking to those of us who already know Jesus. We've already been cleansed..
We've had our sins forgiven. And then in verse four, he says,.
abide in me and I in you as the branch cannot bear fruit by itself unless it.
abides in the vine. Neither can you, unless you abide in me. I am the vine..
You are the branches. Jesus is saying he's the main stock..
He's the roots and the stock of the tree. And we're just the branches..
We can't do it on our own. He says that in the passage, apart from him,.
we can do nothing. But if we will surrender to him,.
and let him live through us,.

$^{241}$then beautiful things will start to happen in our lives..
Beautiful fruit will start to come. And that's what Jesus is calling us to..
He says that if we're abiding in him, then we will keep his commandments..
We will do what he has for us. So maybe you say, okay, what does that look like?.
If I believe that I can't be perfect, that's my heavenly father is perfect,.
that I have to surrender to Jesus in order to do that,.
that I have to start to disappear, that he has to start to live through me..
But what does that look like? How do I even do that? Well,.
the first step is to pray, to talk to God and say, God,.
I want to turn it over to you. I want to surrender to you..
I know I can't do it. I need you in my life. Please take control..
Please begin to live through me. Please begin to bring about fruit in my life..
That's the first step is just to, to ask and to,.
to confess your surrender to God..
But then what does it look like in the day to day of our lives? Well,.
it begins to look like surrendering to these values of the kingdom that we've.
been talking about over the last few months. Instead of saying,.
I'm going to live out my life according to my own values..
I'm going to live out my life in the ways that I think are best..
I'm going to surrender to the values of God's word,.
to the values that Jesus has taught me, to the values of the kingdom..
And I'm going to seek to live those out..
A great starting point is the ones that we've covered in the beatitudes..
Let's think about those for a second..
Blessed are the poor in spirit for theirs is the kingdom of heaven..
Surrendering to that saying, God,.
I'm not going to live from a place of pride or from a place of seeking my own.
righteousness. I'm, I want to be poor in spirit..
I want to understand how much I need you. I want to live from that place..
Blessed are the meek for they shall inherit the earth. Surrendering to saying,.
God, I don't want to live from a place of thinking I'm better than someone else..
I want my own understanding of my poverty of spirit to cause me to treat.
everyone else with love and respect because the rate made in your image..
And I know I'm just as broken as anyone else..
And I'm going to surrender to living out that value instead of trying to live.
from a place of pride and trying to assert myself and trying to give myself.
advantage over others. I'm going to be meek..
Blessed are those who hunger and thirst for righteousness for they shall be.
satisfied. Saying, God, I'm, I'm going to surrender to that value. I'm,.
I'm going to stop hungering and thirsting after all the stuff of the world,.

$^{281}$all the stuff that I, that I crave because of my sin,.
all of the money or the power or the popularity or the lust or whatever those.
things are. God,.
I want to stop living for those things and I want to hunger and thirst after.
righteousness. I want to surrender to that..
Blessed are the merciful for they shall obtain mercy. God,.
I want to live in a way that I see the brokenness of others and that it breaks.
my heart. God, I want to surrender to that value..
Blessed are the pure in heart for they shall see God. God,.
I want to surrender to having a single minded desire towards knowing you and.
stop being distracted by everything else in the world and stop being.
contaminated by all these other desires and all these other things I want..
But God, I want to surrender to only desiring you..
That's hard because our heart wants, our flesh wants everything else..
We want this and that we want to achieve. We want to be comfortable..
We want to have things. We want to be wealthy..
We want all these different things. But can we surrender to saying, God,.
I want the purity of my heart to be found and only desiring you..
Blessed are the peacemakers for they shall be called the sons of God..
Can we surrender to that and say, God,.
I want to live a life that brings Shalom into the lives of others and the verses.
10 to 12 that talk about being persecuted for righteousness sake. God,.
I want to surrender to that..
I want to surrender to the fact that if I'm following you,.
there's going to be people who are uncomfortable with that,.
that it might bring about some persecution and I'm going to live it out anyway..
And I'm, I don't want to give up when things are difficult and following you..
That's a great starting point in surrendering to the values of the kingdom..
Start with the beatitudes..
And then as you continue to know God's word and study the scripture,.
you'll see more and more values of the kingdom that God wants you to surrender.
to as you continue to follow Jesus and learn more about him and deepen your.
relationship with him..
God will open up more and more of his truth and you'll see more and more things.
that you can surrender to more and more things about you and your,.
your selfishness and the old part of you that can pass away, that can die,.
that can disappear..
And more and more of Jesus that can come in and live through you..
And that's what Jesus is saying in this verse when he says,.
be perfect as your heavenly father is perfect..

$^{321}$He's not saying that we're going to try hard enough to do this..
He's saying there's no way you can do it. You've got to surrender..
You've got to abide. You've got to let Jesus do it through you..
You've got to let him live his life through your life..
He doesn't want to improve us and make us a little bit better..
He wants to radically transform us so that he is living in our place..
He is living through us. So let's live into this..
Let's be honest that we're not going to be good enough on our own,.
that we're not going to try hard enough to be all that God wants us to be or.
even to be all that we want to be ourselves..
But we have got to have God's help. We've got to have God's work in our life..
We've got to have the empowerment of the Holy Spirit in our life..
And the reality is that that's what God wants all along..
God never wanted us to do it. God never expected us to do it..
God wants us to disappear, to fade away..
He wants our selfishness, our pride,.
all the things about us that have driven our decisions and charted our course.
before we knew Jesus..
God wants those things to fade away and he wants Jesus to begin living through.
us..
God wants us to surrender to Jesus and let him live his life through us..
God wants us to surrender to the values of the kingdom and say, God,.
I want to live out the values of the kingdom instead of just trying to live out.
the things that are important to me, the things that I value,.
the things that the culture around me values..
I want to live out your values in my life. I want to prioritize those things..
God wants us to surrender to relationship with him and say, God,.
I'm going to prioritize knowing you over everything else in my life..
God wants us to surrender to abiding and saying, God,.
I can't do this on my own apart from you..
I really can't do anything..
So I'm going to rest and I'm going to let you live through me so that I can be.
who you want me to be..
That's what God calls us to and that's what Jesus is pointing us to in this.
chapter..
And we have to cry out to God and ask God to help us to do it..
And we have to ask God to give us the courage and the strength to surrender and.
the courage and the strength to step out of the way and let Jesus live through.
us. So let's do that. We'll do it right now,.
but let's continue to do that over this week..

$^{361}$Let's continue to ask God to be at work in our life,.
to ask God to give us the strength and courage to surrender over and over to.
continue to surrender,.
to make that journey of surrendering to Christ and being closer to.
him and becoming more and more like Jesus and more and more of who God wants us.
to be. Let's pray. God, we do just profess our need for you. God,.
we know that we cannot be who you want us to be or even who we want to be.
on our own without your help. And so God, we just cry out for your help..
We ask you to give us the strength to surrender to you and God,.
I pray that we would do it..
I pray that we would surrender to the values of your kingdom and that we would.
live out what's important to you,.
that we have prioritized the things that you prioritize. And God,.
I pray that we would surrender to orienting our lives around you and really.
focusing everything on you and focusing everything on our relationship to you,.
God. And as we do that, God, we ask for you to live through us..
May your love shine through us. May your peace shine through us..
May your unselfishness and your selflessness shine through us..
May your humility shine through us. God,.
may you live through us so that we ourselves, but more importantly,.
the people around us see less and less of us and more and more of Jesus in our.
lives. God, our world needs that so bad..
Our world so badly needs to see more and more of Jesus..
May we be your people and may we be the people who,.
who live that out, who show more and more of Jesus to our friends,.
our families, our community, our city, and to the world around us. God,.
In Jesus' name, Amen..
[BLANK AUDIO].
\newpage



\section{}
\label{sec:x6dSk6R1byI}
\textbf{聖地網上行 - 步主腳蹤 [第一集] 伯利恆}
\newline
\newline
連結: \href{https://youtube.com/watch?v=x6dSk6R1byI}{\texttt{https://youtube.com/watch?v=x6dSk6R1byI}} ~~~~ 語音日期: 2020-10-07
\newline
\newline
\hyperref[sec:fE8jDxIIYZM]{\small{< < < PREV SERMON < < <}}
~
\hyperref[sec:index_chronic]{\small{[返順時目]}}
~
\hyperref[sec:index_scriptual]{\small{[返順卷目]}}
~
\hyperref[sec:vkydkBsADbY]{\small{> > > NEXT SERMON > > >}}
\newline
\newline
$^{1}$各位團友歡迎來到聖地.
我們現在在聖地的上空.
可能你會問.
到底旅遊和信仰.
看聖經有什麼關係呢?.
正正經經看聖經就行了.
何必要飛到這個地方.
走到這麼遠呢?.
我不知道你有沒有這樣的感覺呢?.
每一次你打開聖經的時候.
你都會覺得勝經的世界和我們的世界有一點分別.
是不是?.
聖經本身的語言其實和我們很有分別.
雖然已經翻譯成了.
現代的語言.
你知道本身聖經舊約是希伯來文.
或者有一點點亞蘭文所寫.
新約聖經就是希臘文所寫.
語言都不是我們現在會用的語言.
是不是?.
所以有一個差距.
我們一定需要去連接.
聖經的世界和它的文化背景.
和我們今天的文化背景都很不同.
我們不是說將聖經.
框在以前的文化背景裡面.
而我想告訴你.
神是藉著當時的.
處境和文化.
和當時的人說話啟示他自己.
我們需要明白當時的處境.
我們就能夠明白.
上帝的遺言和真理是什麼.
我們才能夠應用在.
此時此地.
就像你平時去旅行.
你都要明白當地的語言.
風土文情等等.
今天我們看聖經其實都是一樣.
要了解我們的信仰其實有一點像去旅遊.

$^{41}$我們要明白那個處境.
那個文化.
歷史甚至地理.
你才會仔細發現.
聖經的作者.
神要藉著他們說什麼.
光說也沒用.
整個旅程裡面.
我不斷會有一些示範給你.
給你一些影片.
和連結聖經.
詩篇第125篇第一節.
以靠耶和華的人.
好像石案山總不動搖.
永遠屹立.
群山怎樣圍繞著耶路撒冷.
耶和華也照樣圍繞著他的子民.
從現在直到永遠.
接著他說.
惡人是不會得逞的.
義人才會得逞.
因為耶和華保護他.
看顧著他的道路.
所以在耶和華裡面我們就有平安.
他說眾山都圍繞耶路撒冷.
什麼意思呢.
首先我們飛到第一個景點.
就是耶路撒冷.
舊城O city的上空.
上面一個方形就是我們今天看到的.
舊城耶路撒冷.
下面是一個長的山尺.
叫做Eason Hill東山.
上面其實就是.
大衛城耶路撒冷.
石案山原本的地方.
在這個地方你也找到一些位置.
屬於大衛或是大衛之前的時代.
迦南人時期的一些城牆.
我們下去大衛城.

$^{81}$裡面.
就看到.
大衛城的一個廣場.
我們下去大衛城.
裡面.
我們可以在上面.
做一個360度的遙望.
你就知道C邊125邊.
到底想說什麼.
現在我們下來下面.
向北面望就是聖殿山上的青真寺.
這個深色的圓頂.
就是在我北面.
你看到高度比我高.
向右邊來看的時候.
你會另外有一個山脊.
也是比我高.
叫做迦南山.
迦南山和耶路撒冷.
是對面東面的山.
然後我們往右邊看.
就是一些很密集的屋.
這個叫做Silvan Village.
是一個巴勒斯坦人的村落.
另外我們看.
南面.
就是山頂上.
就是UN的總部.
然後下面我們看到.
一些大衛城的考古的遺跡.
然後向右邊.
我們向西邊看.
遠遠有一個教堂.
那裡就是西山.
耶路撒冷的西山.
也是比我站立的地方.
更高.
我們現在回到北面.
你看到.
這個大衛城.

$^{121}$站在那裡.
我們360度這樣看的時候.
就看到真的中山環繞著.
石安山在耶路撒冷.
所以詩人是從地理地景.
引伸出神學的信念.
就是說這些山.
真的在物理狀態下圍著他.
然後他就想起.
原來不單止是物理狀態.
是山圍著.
想起的是耶和華.
去保護著.
一起倚靠他的人.
這裡就領略到詩篇的訊息.
我們是否去倚靠上帝呢?.
首先我們在聖經的地圖入手.
很多時候你會留意在電子聖經.
或者印出來的聖經.
都會有一些聖經的地圖.
但是你可能未必.
能夠去了解到.
這些地圖給了你什麼資訊.
今天我們首先下載了.
這個地圖的時候.
我想你跟我一起做一些雜作.
在這個地圖的右手邊.
你去寫下.
聖地這個字.
很多時候人們會叫這個地方.
Holy Land.
聖地.
也有很多人叫它以色列地.
或者叫巴勒斯坦地.
這個字在公元後第二世紀.
猶太人第二次革命後失敗.
羅馬王就將整個土地改名.
叫它巴勒斯坦.
去驅逐大部分猶太人.
舊約聖經裡面.

$^{161}$也會叫這個地方.
迦南地.
也會叫它應許之地.
或者甚至留納漁物之地.
這些地名.
都是來自舊約希伯來聖經.
今天很多人都會叫這個地方.
做聖地.
也有它的意思.
因為這個地方.
確實是一個宗教色彩.
非常豐富的地方.
這塊長長窄窄的土地.
是三大一神信仰宗教的中心.
這三大宗教.
分別是猶太教,基督教.
和回教.
或者穆斯林的聖地.
所以從宗教來看.
這個地方是富有神聖色彩的地方.
現代人很多時候.
會叫它近東或者中東.
為什麼呢?.
我們回到200年前的時間.
當時十九世紀末的時候.
西方歐洲的國家.
當時是強大的帝國.
從歐洲人的眼中來看.
東方是奧特曼帝國.
或者奧特曼帝國.
土耳其的帝國.
當時十九世紀末.
中日有甲午戰爭.
土耳其的奧特曼帝國.
已經有四百年的歷史.
但在這個時候就開始色微.
開始瓦解.
所以在西方歐洲的大國.
的眼光來看的時候.
他們就關注到亞洲政權變改的問題.

$^{201}$於是乎有很多歐洲人.
歐洲的國家就覺醒到.
有兩個東方的問題.
一個是遠一點叫做遠東.
一個是近東.
慢慢二十世紀初的時候.
西方的政權就開始留意到.
原來在波斯灣和印度一帶.
都是一個國際關係的戰略地區.
譬如英國會在印度駐兵.
所以他們慢慢在語言上.
加插了一個字叫中東.
二次大戰之後.
英法兩國.
就開始對印度的政權.
英法兩國就瓜分了.
敘利亞,黎巴嫩等等的地方.
以及巴勒斯坦地.
因為巴勒斯坦地是接壤印度.
所以中東就包括了巴勒斯坦這個地方.
以上幾分鐘告訴你.
在過去的歷史.
政權的變遷裡面.
使得這塊土地有很多不同的名字.
但是無論有多少不同的名字.
都改變不了以色列巴勒斯坦地的地理特色.
因為它是處於亞洲,非洲和歐洲.
三個大陸板塊的交通樞紐,交通中心.
所以單從地理,地質來說.
的確是中間之地.
是古代世界的中心.
所以在這個地方.
有很多不同世界的商品.
商隊或者軍隊.
都會以它為一個地橋,一個踏腳石.
是人口來往穿梭的一個地方.
從這個地理角度去看.
你就會了解更加多.
為什麼聖經的故事會這樣發生.
譬如亞述王是東邊的帝國.

$^{241}$他們為什麼會來耶路撒冷.
然後去圍攻希西加王.
希西加王在世界的舞台上.
只是一個非常小的角色.
而耶路撒冷本身不是一個很矜貴的地方.
但是亞述王要攻打希西加.
因為這個小王阻礙亞述王來到南征埃及.
因為他不合作所以就要攻打他.
所以神簡選往往不是人眼中的霸權.
或者人眼中的帝國.
而是一些很卑微的地方.
很卑微的人物.
( 錄音: 音樂 ).
《新命紀》第十一章第八節這樣說.
以色列人要進入迦南地的時候.
他們說你們進去這個地方.
這個地方不像你們出來的埃及地.
在埃及地那裡你們殺了種.
然後用腳來灌溉.
好像菜園一樣.
因為埃及地利羅河三角洲是非常平的大田園.
用腳去挖泥土水就會進來.
但是他們這樣說.
你們要進去得遺孽的地方.
是有山有谷之地.
是天上與水之淵的地方.
就是耶和華神眷顧的地.
從歲首到年終.
如果你們留心聽我今天吩咐你們的誡命.
哀耶和華你們的上帝正在事奉的話.
祂必按時降雨.
使你們可以有豐壽.
所以你們要謹慎.
免得你們的心受到迷惑.
你們就離開了神.
去事奉別的神去敬拜他們.
叫耶和華的老氣向你們發作.
祂就將天封閉.
不再降下雨水.
地就不再生產出土產.

$^{281}$使你們從耶和華賜給你們的美地上.
你們的日子就會很快滅亡.
《申命記》的描述很配合迦南地本地的地理.
因為那裡不會有大河.
如果你要找大河的話.
你可以去埃及有尼羅河.
你可以在那裡安居樂業.
你可以去米索不達米艾.
都有兩條大的河.
那裡耕種是很方便的.
但是在迦南地你要生存的話.
你要去繼續安居的話.
你就需要靠雨從天上.
要靠天上下來的雨水.
這是本地的氣候.
每逢十一月開始雨季的時候.
他們就靠地中海的雨雲.
由西至東移動.
然後將雨灑在迦南地上.
以致他們能夠有農產品.
能夠有豐收.
如果沒有雨的話.
寒季就會來.
他們就會有饑荒.
所以要在這裡生存的話.
他一定要去事奉去敬拜耶和華.
因為耶和華是唯一一個來管理雨水的.
你不可以受別的神.
別的宗教來影響.
因為他們不是真正去控制雨水的.
所以迦南地是一個屬靈的教室.
是一個生命的教室.
是因為生命當中來操練我們.
去思念我們.
以致我們對耶和華.
對主忠心.
其實如果你這樣引申的話.
我們人生裡面.
我們經歷很多的事情.
其實都是一個生命的教室.

$^{321}$上帝給我們的鍛鍊.
給我們的考驗.
就是要我們的信心對祂忠心.
要我們不斷地緊緊地跟隨祂.
現在我們試試做地圖的雜作.
這個其實好像當一些桌面遊戲的畫板.
你將標籤放上去就可以了.
我們首先由上至下.
由北至南來講解一下.
上面看到一個湖叫加里里海.
在希伯來的原文沒有分淡水或是咸水.
都是叫做海.
其實就是地球裡面.
海拔以下最低的一個淡水湖.
這個叫加里里海或是加里里湖.
加里里的山區就在附近.
有兩個城鄉.
一個在加里里海的旁邊.
在湖邊的叫加伯隆.
另一個在加里里海的南邊山區.
就是拿撒勒.
兩個地方你記起.
就是跟主耶穌的成長和傳道很有關係.
由上至下我們看到.
由加里里湖出來有一條河.
一直由北至南流入下面的死海.
這條河叫約旦河.
river Jordan.
流到南邊就是死海.
死海是地球裡面.
海拔以下最低最低的地方.
因為水進去這個地方.
是不會出來的.
是死路一條.
所以一定會蒸發.
然後水裡面的礦物和鹽.
就會留在死海裡面.
所以鹽分是非常濃的.
不會淹死你.
但一定會打死你.

$^{361}$現在我們從左至右.
由西至東.
我們有什麼標籤呢.
最大的海當然是地中海.
然後在以色列巴勒斯坦地圓海的一片地方.
我們叫它圓海平原.
是一個很開放平地的地方.
其中一個很重要的港口叫蓋薩尼亞.
然後我們再往東移.
去到中央山區.
Central Hill Country.
這裡有兩個城鎮.
你需要留意.
北面是耶路撒冷.
南面是伯利恆.
耶路撒冷和伯利恆.
都在新約耶穌福音書的故事裡.
經常出現.
然後再往東移.
離開中央山區.
去到猶大曠野.
猶大曠野也是耶穌試探的地方.
然後再往東移.
我們越過約旦河河谷.
上來到約旦河東的高原.
約旦河東的高原再往東移.
就會遇上真正的沙漠.
就是阿拉伯沙漠.
不如我們不要只看地圖.
我們給你一些短片.
你猜一下這段片是從哪裡出來的.
看看你猜中了多少個.
第一個.
沿海平原對著地中海.
第二個.
這個是猶大曠野.
沒錯,這裡是約旦河.
有些教會的人在洗禮.
這裡看到一個很大的劇場.
是對著地中海的.

$^{401}$這個是蓋薩利亞港.
即是沿海平原.
這裡從高山跳望.
看到加里里湖或加里里海.
這裡在加里里湖乘船.
很漂亮.
這裡有些梯田.
是加里里山區.
拿撒勒這個地方.
然後有牧羊人帶著羊.
這個也是回到猶大曠野.
猶大曠野是牧羊的地方.
這個很簡單.
這個是四海.
這個呢.
這個就是約旦南部的沙漠地帶.
你看到迦南地確實是一個多元的地景.
這是一個很難生存的地方.
但你需要生存活著的話.
你就需要單靠控制雨水的上帝.
單單事奉他.
其實你這樣想就是我們生命的縮影.
是一個屬靈的教室.
去鍛鍊我們提醒我們.
只需要單靠這位上帝.
如果你用一個大一點的景觀來看.
你會發現這裡不是建立地上帝國的樂土.
如果你要建立地上帝國.
你可以找其他不同的角落.
米索不達米亞可以去埃及.
但這裡是活在地上帝國的陰影下.
這裡的神的子民.
他需要單靠上帝.
並且去抗衡地上帝國的價值.
我們對以色列巴勒斯坦地理有一個概括認識的時候.
現在我們可以追蹤一下耶穌的行蹤.
按照馬達科罕第二章所說.
耶穌出生的時候.
就是大希律王Hildegrave執政的時候.
大希律王是主前三十七至四年間.

$^{441}$做猶太人的王.
推斷耶穌出生的時候.
大概是主前公元前六年的時候.
大希律王本身不是純正的猶太人.
他是南面以東民族的一個人.
但是這個人很有政治的手腕.
以及野心勃勃.
他一生靠攏著羅馬帝國.
然後代替羅馬帝國來強硬地控制猶太人.
大希律王的生性是非常妒忌和忌才.
時時刻刻怕別人謀朝篡位.
所以他一生想盡方法去滅聲.
去殲滅一切的政敵.
包括自己的家人.
大希律後稱為古代最偉大的建築家.
因為單單在巴勒斯坦地.
他有十幾個建築了的堡壘和皇宮.
薩瑪利亞,蓋薩利亞,海港.
馬薩大,耶利哥,耶路撒冷等等.
這些堡壘和皇宮不單是自己享受所用.
最重要是鎮壓當地反羅馬的勢力.
也是鎮守羅馬帝國東邊的邊境.
單從這些建築我們就知道.
甚麼是活在帝國的陰影下.
我們首先來看大希律所建的其中一個皇宮.
我們由耶路撒冷出發.
我們去耶路撒冷以南大概七至八英里.
十二公里的地方.
這裡叫希律堡.
希律堡是大希律建的皇宮堡壘.
以自己命名是他心所愛的地方.
也是最後他埋葬的地方.
希律堡建立在一個天然的山峰之上.
很有大希律建築的本色.
因為他將一個天然的山峰削平.
然後在山峰上建立四個瞭望塔.
有七層樓那麼高.
在希律堡底下山腳的地方.
就有不同的宮殿和房間.
還有大型的羅馬花園和淡水游泳池.

$^{481}$游泳池的水是由附近的百里行城.
用人工的引水道帶進來的.
所以從一方面可以讚嘆這個人.
確實是一個很偉大的建築家.
但是從這些建築背後.
你就可以看到什麼叫做帝國狂魔霸權.
這是聖誕節的背景.
當時馬利亞已經有新人.
約瑟和馬利亞本身在北面加利利.
一個小鎮叫做拿撒勒.
因為當時羅馬帝國要強制的.
由上而下的.
要天下的人民都回到自己的鄉下報名上策.
報名上策是一回什麼事呢.
其實是羅馬帝國一種由上而下強制的法令.
保障國家保障帝國的安全.
也保障帝國的稅收.
所以約瑟和馬利亞就由北面的拿撒勒.
花了幾天的時間.
長途不捨地回到自己的家鄉巴黎行.
好 讓我先帶大家去巴黎行.
這裡就是耶穌出生的地方.
也是整個基督教信仰法院的地方.
巴黎行位於耶路撒冷以南.
大概五英里或者九公里.
剛才我們去過希律堡和.
巴黎行相隔三英里或者五公里.
我們來到主誕堂Church of Nativity.
這個古舊的教堂.
紀念耶穌出生的地點.
這個教堂的底層是主後四世紀的遺跡.
我們現在站在大門的面前.
可以看到這個教堂是歷史悠久.
你看到一個很大的方形的門口.
其實可以追溯至主後五世紀.
即是一千五百多年之前所建的堂皇門口.
但到了十二世紀的時候.
十字軍來到聖地這個地方.
他們把這個教堂的門口收緝了.
成為一個教室的拱門.

$^{521}$你看到這些尖頂的拱門.
就是十字軍的建築風格.
但到了十五世紀.
Ottoman Ottoman時期的時候.
因為防止強盜騎馬進入教堂搶劫.
把門再收緝成為一個更小的門口.
這個很有意思.
因為當遊客進入這個地方的時候.
就好像告訴每一個人.
你必須謙卑地進入這個教堂.
因為這裡是紀念主耶穌謙卑降生的聖地.
這個古教堂一層一層蓋著一個小山洞.
這個符合當時古巴斯坦地的房屋設計.
因為很多時候這些小屋.
都是由一些小山洞加建出來的.
裡面是馬棚.
外面是人住的地方.
雖然現在已經是面目全非.
但藉著不同最早期的文獻.
考古歷史等等不同的記載.
一直都沒有質疑過.
主誕堂是耶穌出生之地的可靠性.
現在我們在大堂正在收緝的時候.
現在2020年收緝的工程已經差不多完工了.
下面保存了主後四世紀的馬賽克地板.
馬賽克是甚麼呢.
是一些小石頭裝嵌出來的石地氈.
非常精細的手工.
現在我們進入教堂最中心的地方.
就是要走幾級下去地底.
這裡就是山洞的地方.
這裡看到一個祭壇.
下面是一個十四角的星.
為甚麼是十四角的星呢.
如果你記得勝經怎樣記載的話.
耶穌的家譜是由阿巴拉罕到大衛有十四代.
由大衛到被擄巴比倫的時候也有十四代.
被擄到巴比倫到耶穌降生的時候也有十四代.
十四也是大衛名字的數目符號.
祭壇的後面就是馬槽的地方.

$^{561}$我們讀一讀《馬太福音》第二章.
希律王執政的時候.
耶穌生在猶太的伯利恆.
那時有幾個占星家博士從東方來到耶路撒冷說.
那生下來作猶太人的王在那裡.
我們看見他的星出現.
特來朝拜他.
希律王聽見就心裡不安.
傳耶路撒冷的居民也是這樣.
東方大概是米索不達米亞的地方.
而當時羅馬帝國的外敵勢力就處於這個地方.
他們是一個外投來的勢力.
他們來到耶路撒冷來拜訪耶穌.
但希律聽見的時候當然非常不安.
疑心大作.
因為老虎頭上釘塞拿.
或者太歲頭上洞土.
可能他會想難道我臨終之前.
又有另一個想挑戰我的王位.
你可以看到馬太科罕第二章記載.
耶穌的出生這個描述是很滑稽和諷刺的.
或者可以這樣說是一種嬉笑怒罵的方法.
可能是當時在教會裡面傳閱的一個頭條新聞.
就是告訴教會只要他們在耶穌基督裡面.
他們就能夠得勝所有壓迫他們的權勢.
剛才我們去過一個地方叫希律堡.
在羅地安 距離伯利恆非常近.
我們提過大希律在東邊的邊界.
有一連串的堡壘來把守森嚴.
但這幾個博士竟然初到貴境.
不知頭不知路.
他們卻能夠逃離大希律的五致散.
原來背後有上帝指示這些博士.
叫他們不要回去.
從其他的路逃走了.
而小孩耶穌也是一樣.
一個手無寸鐵的小孩.
上帝指示他們離開這個地方.
走到埃及避難.
上帝勝過地上的君王和政權.

$^{601}$如果這樣看的話.
耶穌的作王天國的加林.
就是顛覆地上一切的勢力.
當耶穌來到的時候.
任何人都要讓位.
如果我們接受耶穌基督成為我們生命的主的時候.
我們同樣要讓位.
怪不得有史以來的獨裁者.
獨裁國家都不喜歡人們慶祝聖誕節.
都不喜歡人們想起聖誕節的故事.
小孩子耶穌.
伯利恆裡面的孩童.
他們是被打壓的教會.
他們代表黑暗帝國勢力所壓迫的人.
他們受盡威嚇 受盡引誘.
甚至是殉道.
不過這個故事最終告訴我們.
屬於主的教會.
屬於真理公義的天國子民.
縱使在帝國的陰影下.
他們仍然去敬拜.
仍然跟隨這位真正的王者耶穌基督.
他們會出天國門道的生命.
他們最終是會勝利的.
我們看完伯利恆.
我們又看完聖誕節的經文.
到底我們自己是站在哪一方?.
我們是否站在地上.
一定會失敗的權勢那一方.
還是我們站在這位真正的王者耶穌基督那一方?.
我們同心抵抖祈禱.
主我們謝謝你.
我們能夠去紀念你出生的地方.
也看見聖誕節的故事.
主你再一次提醒我們.
你是那位真正的王者.
你是那位真正的主.
求你繼續幫助我們.
我們歷世歷代不同的時候.
我們都活在不同的帝國的陰影下.

$^{641}$主我們求你親自幫助我們.
站在你那一方.
跟隨你.
作你天國的門徒.
以耶穌基督為我們真正的王.
真正的主宰.
神我們也為著.
我們現在活在這樣一個時候.
可能感到很黑暗.
感到很多的威逼.
很多的壓力.
主我們求你.
用你的恩典來保護.
主我們為了香港這個地方.
特別最近國安法的通過.
令到社會的氣氛非常之幽暗.
神我們求你親自保守這個城市.
保守特別當地的教會.
讓他們真的站在主耶穌王者的那一方.
在很多的威逼利誘下.
主讓他們堅守.
來敬拜你.
堅守著他們的信仰.
來為你天國發聖.
主我們在不同處境下.
我們求你繼續幫助我們.
忠心地來做你的門徒.
做你天國的子民.
站在你聖利那一方.
我們這樣祈禱.
奉主耶穌基督的名字.
我們去到不同聖地的地方.
亦看過聖經不同的故事.
我們希望你覺得今天的旅程.
是滿載而歸和非常之豐富.
我們希望下一集能夠和你準時的出發.
我們繼續去尋索耶穌的大大小小事跡.
我們下次再見.
多謝各位收看.
\newpage



\section{}
\label{sec:vkydkBsADbY}
\textbf{浩劫重生 【第3集:出埃及十災】 陳崇基牧師 - 主日講道視頻系列}
\newline
\newline
連結: \href{https://youtube.com/watch?v=vkydkBsADbY}{\texttt{https://youtube.com/watch?v=vkydkBsADbY}} ~~~~ 語音日期: 2020-10-12
\newline
\newline
\hyperref[sec:x6dSk6R1byI]{\small{< < < PREV SERMON < < <}}
~
\hyperref[sec:index_chronic]{\small{[返順時目]}}
~
\hyperref[sec:index_scriptual]{\small{[返順卷目]}}
~
\hyperref[sec:O6HbgHZahfI]{\small{> > > NEXT SERMON > > >}}
\newline
\newline
$^{1}$今年2020年見到天災人禍接二連三的發生.
可能你和我一樣都會想起聖經出埃及記十災故事.
十災是以屬能詳的聖經故事.
講到埃及帝國的法老不容許以色列人離開埃及.
於是上帝就降十個災禍懲罰審判埃及.
十災這個故事是記載在出埃及記七至十二章.
這裡講到水變血的災禍.
災禍是話災,實災,營災,蓄逆之災,滄災,薄災,黃災,黑暗之災.
最後第十災是殺長子.
最近有人在網上留言.
他說見到十災的情景就好像2020年我們所遭遇的一樣.
但我們如果退一步想一想.
出埃及記十災故事的歷史背景的話.
會不會十災只是當時亂世冰山的一個.
可能我們這樣想的時候.
我們會對十災的記載有另一個體會.
出埃及記的年代.
你如何估算聖經的日期,年日等等.
都是在公元前1500至1200年之間.
這個時期就是近東歷史最關鍵的時候.
我們叫它做Labyrinth Age.
青銅時代後期.
當時的近東世界進入一個最興盛的時候.
大約在公元前1340年.
有一艘商船在東地中海的海域航行.
這艘船載滿著很多名貴的貨物.
不知道是什麼原因.
可能是遇到風浪.
可能是船長大意.
又或者他們正在逃避海盜.
這艘船經過土耳其南邊一個小鎮叫烏魯布倫.
在這個海域突然間擱淺.
整艘船連同貨物沉在海底.
直至1982年.
有一位潛水員意外發現這艘沉船.
經過十年超過二萬多次的潛水考古工作.
才把這艘商船的貨品一一打撈上來.
船上的貨物是價值連城的.
至少有149個迦南風格的長身陶瓷大缸.
這些大缸裡面載滿著各種的食物.

$^{41}$有杏仁 無花果 橄欖 葡萄.
另外在船上亦找到上等的木材.
是埃及而來的.
找到象牙 荷馬牙十多條.
亦有象牙製的化妝用具 首飾 金銀寶石 馬路.
在船上亦找到不同的武器.
有箭頭 矛頭 匕首 軍用利劍.
最特別的是在船上找到未經提煉的銅和石.
銅的重量是十噸 石的重量是一噸.
銅是來自當時的希臘半島 塞浦路斯 克里特島等等.
石 天 主要的來源是來自今天的阿富汗.
青銅是當時最堅硬的材料.
是怎樣提煉出來的呢.
就是將銅和石以九比一的比例化合而成.
你可以看到在船上的銅和石足以生產十一噸的青銅.
是當時最先進武器的原料.
這艘沉船的貨物來自七 八 甚至更多的主要國家地區.
表示在青銅時期是一個盛世的時期.
由東面開始有以蘭國 巴比倫 亞述 赫人帝國 埃及.
塞浦路斯島國 米洛島國 希臘的邁西尼國.
是當時的G7和G8.
我們從一些文件 泥板等等.
也知道當時這些國家彼此通商 彼此有外交的關係.
在軍事上也有很多競爭.
是一個國際盛世的時候.
但很奇怪 在青銅時代中末時期 短短五十年.
我們估算大概是一二二五至一一七五年之間.
這些國家文明好像骨牌一樣突然逐個瓦解.
到底有什麼勢力使這些歷史悠久 強大的文明政權好像一夜之間覆亡.
會不會是地震 旱災 戰爭 經濟崩潰 科技轉型.
政變 人口遷徙 還是疫症大流行.
這裡是土耳其的赫人帝國Hittite Empire的京城Hetusa.
赫圖沙本身是一個堅固防守的大古城.
歷史非常悠久 但突然到了公元前1200年.
就被入侵 被大火所燒毀.
不單止赫圖沙是這樣.
其他在土耳其 敘利亞 以色列 巴哈斯坦等地的大城邦.
海港 貿易中心都遭受同樣好像一夜之間的崩潰.
在埃及帝國有不同的紀錄提到有一浪又接一浪的入侵外敵.
這些外敵叫做Sea Peoples 海人.

$^{81}$有很多是藉著陸路入侵中東不同城市.
也有很多是藉著海路入侵海港和埃及帝國.
他們是由不同的種族所造成的.
埃及提到這些名字的時候 有很多我們今天仍然很陌生.
我們不是太清楚這些族群的來頭.
不過其中一個叫做Pelaset.
這個所有歷史學家都認定.
就是聖經提到的非利士人Philistines.
埃及的法老蘭瑟三世Wemsis III.
曾經和這些海人Sea Peoples在尼羅河的三角洲大戰過好幾場.
他將他們擊退.
然後這些非利士人入主了迦南地的南面.
大家熟悉聖經的話就知道他們以後成為古以色列人的心腹大患.
曾經我們以為這些神秘的海人Sea Peoples.
就是使青銅文明突然崩潰的唯一原因.
但如果你看看這些壁畫.
你會看到這些Sea Peoples是帶著小朋友和女人.
似乎他們都是被迫逃避一些的豪傑和天災人禍.
被迫遷徙到不同的地方.
在同一個時期 在赫人帝國末期.
有一個赫人的皇帝叫做穆西尼斯二世.
穆西尼斯二世留下了一個這樣的祈禱.
他當時向赫人的神靈來祈求.
他說赫人的雨神 我的主 赫人的忠神.
到底你們做了什麼.
在赫人之地你們釋放了瘟疫.
在我父 我兄弟和我的當前.
我不斷看到有死亡.
在我父 我兄弟和我的當前.
我已經有20年之久.
根據這些文件 歷史資料 線索等等.
現在我們認為青銅文明突然崩潰.
原來是一場完美的風暴 是一場perfect storm.
有氣候的變化 climate change.
在希臘半島有一連串的地震.
國與國之間有爭戰 經濟崩潰 科技轉型.
也有疫病大流行 人口遷徙 甚至政變.
以至當時的超級大國 包括埃及帝國.
也一落千丈.
你看看當時的情況.

$^{121}$你問一問 到底和今天有多少相似的地方.
所以歷史學家稱這一段時候.
為地中海的末日浩劫 Mediterranean apocalypse.
剛才我們用了一個大歷史的圖畫.
我們看了之後 再看十災.
似乎意義會有一點不同.
似乎十災不只是上帝審判埃及.
也表明神在人類的大歷史舞台背後工作.
十災只不過是一個縮影.
上帝要廢亡立王 要改朝換代.
他要告訴埃及帝國 要告訴整個世界的人類.
到底哪一個才是歷史的主宰.
如果你細心看 聖經裡記載十災.
是用不同的字句 很多時候是重複的.
在結構裡表達這個時候.
在十災的結構裡 你可以將十個災禍分成三組.
即是三三三一這樣分組.
每三個災禍成為一個組合.
每一組的頭一災 都有同樣的字句開始.
首先我們看第一組 即是第一災水變血.
在七災第十四節 明天早晨.
法老出來水邊的時候.
你摩西要在河邊迎接他.
到第二組 即是第四災 八災二十節.
耶和華對摩西說.
你清早起來 法老來到水邊.
你就站在他面前.
到第三組 開頭是第七災.
九災十三節.
耶和華對摩西說.
你清早起來 站在法老的面前.
然後我們看每一組的第二災.
先看第一組 即是第二災八災第一節.
你進去見法老對他說.
到第二組 即是第五災.
同樣的字句 你進去見法老對他說.
到第三組 第八災.
同樣的字句 你進去見法老對他說.
然後每一組最後的災禍.
即是第三第六第九.

$^{161}$都沒有事先預告警告法老.
結構越來越清楚.
再看第一組 頭兩災.
有埃及的法術師.
他們都做到摩西和亞倫所行.
到第二組 頭兩個都說到.
埃及和以色列人有分別有分開.
營災和蓄逆之災.
沒有落在以色列人身上.
到第三組 頭兩災 即是第七第八災.
兩個災禍都強調在埃及是前所未有.
九章第十八二十四節.
類似都說到 埃及開國以來沒有這樣.
整體來看 作者記載十災的時候.
他不是單單寫一些事實的報道.
而是用故事 用敘記來講神學的信息.
十災的故事.
就是要表明上帝比埃及的神和法老更大.
但每一組的中心都有特定的信息.
第一組 頭三災.
就是神的能力超越埃及的法術師.
超越他們的邪術.
到了第三個災禍 失災.
埃及的法術師承認.
行法術的也用邪術來生出失子來.
卻是不能 連他們都承認.
這是神的指頭 神的工作.
第一組看完 我們看第二組.
第二組所強調的.
就是上帝的偉大 他的權能.
是在於他有能力去區分以色列人和埃及人.
上帝是完全控制著大自然的災禍.
大自然的勢力.
他可以降災在這裡.
保護另一個地方 是療癒之象.
然後看第三組的災禍.
又有一個特定的信息.
就是上帝的能力是無與倫比的.
是沒有人可以和上帝相比的.
到最後第十災.

$^{201}$當然是技術最強.
和最嚴重 打擊最大.
而且是上帝直接出手的.
是不經過 不藉著摩西和亞倫所做的.
神親自降滅投生的災.
去擊打法老.
因為當時埃及的法老自稱為神.
他是保證埃及帝國繁榮興盛.
上帝特別用這個最嚴重的災禍.
來打擊法老.
所以不是單單記述一些歷史的事件.
而是用故事來寫神話.
當以色列人想像一下.
他們看到整個世界開始瓦解.
有很多的豪傑災禍臨到的時候.
整個世界崩潰的時候.
這個故事告訴他們.
上帝依然掌權.
在人眼中好像很混亂.
但在上帝眼中.
祂完全操縱一切.
但你繼續看十災聖經的記載.
十災似乎不是單單來教訓埃及和法老.
真正的目的是什麼呢?.
在第十災降下之前.
耶和華的很長篇的交代.
一個猶太人的節日.
叫做預月節.
十二章第十一節.
這是耶和華的預月節.
因為那一夜我要巡行埃及地.
把埃及地一切投生的.
就是敗壞埃及一切的神.
我是耶和華.
這血要在你們所住的房屋上作記號.
我一見這血.
就越過pass over.
預月的意思.
就越過你們.
我擊殺埃及地投生的時候.

$^{241}$災殃必不臨到你們身上.
滅你們.
這是預節晚餐的記載.
第一次的預節他們殺了一隻羊.
大家吃羊羔的時候.
他們把羊羔的血.
抹在門尾的上面和兩旁.
以致滅命的天使.
越過以色列民.
越過神的子民.
大家留意.
這裡說的是耶和華的預月節.
預月節不只是保我們平安.
是與耶和華立約.
最重要是吃預節的時候.
慶祝這個節期的時候.
就是你成為神國的子民.
你再看看以色列人過了紅海之後.
埃及的軍隊是全軍覆沒.
十四章31節是怎樣記載的.
以色列人看見耶和華向埃及人所行的大事.
就敬畏耶和華.
信服他和他的僕人摩西.
所以十災預節pass over.
過了紅海.
不只是我們獨善其身.
我們沒事就好了.
其實上帝就是說.
你要成為我立約的子民.
你要忠於我的約.
你要信服我.
你要敬畏我.
你看見這些浩劫,十災,災禍的時候.
看見世界好像很混亂.
很多災禍臨到.
你要知道我仍然掌權.
不只是保你平安.
而是我救贖你.
你能夠成為我的子民.
你就在地上見證我.

$^{281}$剛才我們說.
經文很多時候用一個結構來表達神學.
剛才我們看看.
是三三三一的結構.
或者我們可以這樣想.
到底十災真正的教導目的訊息是什麼.
我們可以考慮一下.
另一個結構是七三的結構.
頭七災經常會出現一些字眼.
譬如七章第五節.
耶和華說我伸手攻擊埃及.
將以色列人從他們中間領出來的時候.
埃及人就知道我是耶和華.
好像十災是教訓埃及人.
但過了第七災到第八災.
到第八災開始的時候.
最後這三災的字眼就很不同.
十章第一節.
耶和華對摩西說.
你進去見法老.
我使他和他神僕的心剛硬.
為要在他們中間顯我這些神跡.
並要叫你將我向埃及人所做的事.
和在他們中間所行的神跡.
傳於你的兒子和你孫子的耳中.
好叫你們知道我是耶和華.
轉過來就不同了.
不是埃及人知道我是耶和華.
而是你們知道我是耶和華.
似乎真正的訊息.
最終要教導的.
就是我們.
神的子民.
留意一下十災的描述.
好像是叫我們想起.
希伯來人,以色列人.
舊約裡面一些很重要的經文.
我們回到第一災.
水變血.
七章十九節.

$^{321}$耶和華曉如摩西說.
你對亞倫說.
把你的象身在埃及所有的水以上.
就是他們的江河池塘上.
叫水都變成血.
在埃及遍地.
無論木器石器都變成血.
他說埃及所有的水.
Pools of water.
這個字是一個很特別的字.
之前只出現過一次.
就是創世記第一章.
一章第十節.
神稱海地為地.
稱水的聚處.
Pools of water.
Gathering of water.
為海.
上帝捍衛是好的.
這個字就叫我們回到創世記第一章.
另外叫水都變為血.
Let there be blood.
or become blood.
Let there be.
在十災裡面用了好幾次.
又是創世記第一章經常出現.
Let there be.
要有光.
要有空氣等等.
到第二至第四災.
即是嘩災和實災和凝災.
影響大自然的領域.
就是水,地和天.
又是叫我們想起創世記第一章.
第五災就是蓄逆之災.
影響地上的生存.
叫我們想起上帝在創世記裡面.
第一章第六天的創造.
第七和第八災.
即是薄災和黃災.

$^{361}$擊打的是田間蔬菜.
這句說話又叫我們想起.
創世記第一章第三天的創造.
到第九災.
黑暗之災.
光與黑暗分開.
這個起碼是創世記第一章.
創造的記載.
十災當然就是反創造.
去到uncreate整個埃及和當時整個世界.
但同時也是新的創造.
很久之前的猶太人拉比.
其實早已想到這方面.
當他們看見十災的記載的時候.
他們就問一個問題.
為什麼是十個的災禍.
為什麼不是七個的災禍.
當然十和七都是特別的數字.
但為什麼特別提到十呢.
提到十的原因.
其實又是回到創世記第一章.
因為在創世記第一章.
上帝創造天地的時候.
也是有十次他說話.
十次都是說上帝說.
下一次上帝再說十句說話是什麼時候呢.
就是在出埃及記二十章.
當以色列人離開了埃及.
他們去到西乃的曠野.
在西乃山.
上帝頒佈十誡.
其實舊約聖經說到十誡的時候.
不是ten commandments.
而是ten words.
是十句說話.
十個生命的原則.
要上帝新的子民.
按照這些原則.
這些道理來生活.
以致在萬國當中.

$^{401}$以色列民.
神的子民.
在地上能夠見證上帝.
所以最後你這樣看的時候.
你看完十災的故事.
你就要問一個現在我們應用的問題.
就是上帝在歷史的舞台當中.
他輝煌立和.
將世界不斷瓦解重整.
其實要說清楚一件事.
就是上帝要創造一班新的子民.
都是見到很多災禍.
很多翻天覆地的事發生.
整個世界被重整的時候.
其實上帝就是要創造自己的子民.
新的創造就是我們這班在基督裡面的人.
在疫情當中.
翻天覆地的裡面.
各地都有不同的新的屬靈的模式.
以往我們習慣了的實體的聚會.
是其中一個例子.
今天我們就需要新的適應.
我們或許都會想.
什麼時候才回到以前的正常.
但是疫情過了之後.
還是不會有正常.
就算有正常.
都是一個新的正常.
一個新的常規.
我們現在要做一些創新的事.
就是上帝迫我們去有新的經歷.
新的屬靈的經歷.
以往習慣了實體的模式當中.
我們有的經歷.
現在我們可能是混雜有實體.
更加有妄想.
這個時候弟兄姊妹.
所有神的子民.
我們就要學習.
如何去自己去更新自己.

$^{441}$在不同新的模式當中.
我們要學習去認識上帝的旨意.
並且以前我們有的心思意念.
可能我們只是在想.
上帝可不可以幫我防災.
讓我平安.
今天我們可能在想.
今天是一個新的時代.
上帝令到世界翻天覆地之後.
再重整過之後.
我自己應該如何學到去成長.
學懂如何自我學習.
學懂如何在我現在的位置裡.
去見證和事奉祂.
或許上帝要你做一些新的事情.
在這個時候.
其實是最佳的時候.
你去祈禱.
求主親自用你.
可能用你一些完全沒有想過的方法.
完全沒有想過的崗位.
求上帝幫助我們.
求主也幫助我們在疫情裡.
在這些翻天覆地的浩劫當中.
我們心意更新重新過來.
我們繼續盡心盡力愛主.
縱使沒有了以往我們所依靠的東西.
但今天我們真心歸向主耶穌基督.
在我們主耶穌基督的心意和使命下.
我們用盡不同的新方法.
去見證祂,事奉祂.
我們同心低頭祈禱.
不同的浩劫.
來幫助我們反省.
去迫使我們前進成長.
主,我們過往很多時候.
只是當你給我們安書,平安.
今天我們在疫情當中.
在人生不同的境況裡.
我們知道你仍然是掌權.

$^{481}$你才是那位歷史的主宰.
我們見到世界裡有很多不同的災患.
有很多不同的浩劫.
我們知道你仍然是掌權.
我們要去敬拜你,事奉你.
幫助我們在這個新的世界當中.
在新的創造裡.
你重新更新我們的生命.
以致我們繼續在我們的崗位中.
忠心來事奉你.
縱使是一些新的事奉,新的領域,新的工作.
求你教導自己的子民.
我們的弟兄姊妹.
每一個的教會.
我們在這個世界當中.
我們能夠繼續堅持.
去督順你,去信服你,去事奉你.
祈禱奉主耶穌基督的名字.
《祈禱》.
在困苦艱難.
仍欠卑禱告.
要安靜.
知道你是神.
無罪羔羊.
曾被殺願犧牲.
為救贖.
流血捨身.
在困苦艱難.
仍欠卑禱告.
要安靜.
知道你是神.
無罪羔羊.
曾被殺願犧牲.
為救贖.
流血捨身.
令我如月死亡.
越過風浪.
神恩手把我恩寵.
勝過危疾災禍.
潔淨一生過錯.

$^{521}$藉基督寶血保護我.
在困苦艱難.
仍欠卑禱告.
要安靜.
知道你是神.
無罪羔羊.
曾被殺願犧牲.
為救贖.
流血捨身.
為救贖.
流血捨身.
令我如月死亡.
越過風浪.
神恩手把我恩寵.
勝過危疾災禍.
潔淨一生過錯.
藉基督寶血保護我.
令我如月死亡.
越過風浪.
神恩手把我恩寵.
勝過危疾災禍.
潔淨一生過錯.
藉基督寶血保護我.
願救恩歸於.
坐在寶座上羔羊.
耶穌.
願救恩歸於.
坐在寶座上羔羊.
耶穌.
願救恩歸於.
坐在寶座上羔羊.
耶穌.
願救恩歸於.
坐在寶座上羔羊.
耶穌.
令我如月死亡.
越過風浪.
神恩手把我恩寵.
勝過危疾災禍.
潔淨一生過錯.

$^{561}$藉基督寶血保護我.
令我如月死亡.
越過風浪.
神恩手把我恩寵.
勝過危疾災禍.
潔淨一生過錯.
藉基督寶血託活我.
願救恩歸於.
坐在寶座上羔羊.
耶穌.
願救恩歸於.
坐在寶座上羔羊.
耶穌.
\newpage



\section{}
\label{sec:O6HbgHZahfI}
\textbf{神的名字2\_Yahweh (靈思牧禱)}
\newline
\newline
連結: \href{https://youtube.com/watch?v=O6HbgHZahfI}{\texttt{https://youtube.com/watch?v=O6HbgHZahfI}} ~~~~ 語音日期: 2020-10-14
\newline
\newline
\hyperref[sec:vkydkBsADbY]{\small{< < < PREV SERMON < < <}}
~
\hyperref[sec:index_chronic]{\small{[返順時目]}}
~
\hyperref[sec:index_scriptual]{\small{[返順卷目]}}
~
\hyperref[sec:Lh5hqFxfj84]{\small{> > > NEXT SERMON > > >}}
\newline
\newline
$^{1}$想像一下你爸爸是全世界最傑出的一個藝術家.
他做了一個很創意,很輝煌的玻璃球.
就千分萬負地跟你說.
你只可以看,不可以摸的.
但是你就好像明知故犯那樣.
結果就砰一聲.
整個玻璃球掉在地上,粉碎了.
你也受到很重的傷.
在這個時候,爸爸很難過.
長嘆一聲說.
我要將這個玻璃球修補還原.
也會醫好你的傷.
這個是一個虛構的故事.
但是我很想用它來做一個比喻.
比喻起初上帝創造完美的天地.
在伊甸園裡面.
作為人類代表的亞當·夏娃.
她說她吃了禁果.
整個創造都好像這個玻璃球一樣.
粉碎,毀壞.
人類早晚自己也受了重傷.
一發不可收拾.
但是也在這一刻.
上帝向人承諾.
他說會將這個毀壞的創造.
還原和救贖過來.
其中包括醫好人的重傷.
在這一集.
我就想和你講到上帝的另一個名字.
叫做「Yahweh」.
在中文聖經也叫做「耶和華」.
你看看上帝怎樣去履行祂以上的承諾.
「Yahweh」這個名字.
在聖經裡面出現了6800次有多.
強調祂是以色列的神.
和以色列有親密的關係.
不過我們不能停在這裡.
因為詳細地研讀.
「Yahweh」是和全人類.
包括你和我.

$^{41}$都可以有一個很親密.
而且是萬萬不用錯過的關係.
所以上帝承諾.
要將這個被人毀壞的創造.
還原和救贖.
而祂的計劃就是從一個民族開始.
首先祂找到一位.
在米索不達米亞的族長亞伯蘭.
呼召他離開本鄉本族.
去一個祂為他們預備的地方.
耶和華對亞伯蘭說.
你要離開本地,本族,本家.
往我所要指示你的地方去.
至此耶和華和他呼召出來的這一族.
後來我們稱祂為猶太民族.
建立了盟約的關係.
從這個關係.
讓人看清楚「Yahweh」這個名字.
三個深重的含義.
而第一個含義.
就是「Yahweh」不單只是上帝自身的名字.
所謂Personal Name.
更加是上帝與以色列人立約的名字.
就在《創世記》第17章所記.
耶和華向亞伯拉罕.
即是之前的亞伯蘭顯然對祂說.
「我就與你立約.
使你的後裔極其繁多.
我要與你並你世世代代的後裔.
建立我的約.
作永遠的約.
要作你和你後裔的神」.
這個約的精髓就是.
如果你忠於我.
服從我的命令.
我就必定祝福你.
而通過「Yahweh」和以色列民的關係.
我們看到這位立約的神.
果然真是世世代代都履行祂曾經答應過的.
就是永遠不離開和撇下忠於祂的指紋.

$^{81}$而第二個.
「Yahweh」這個名字的含義就是.
「我是自有永有」的意思.
當「Yahweh」在西那山上向摩西顯現的時候.
祂揭示自己奧妙的一面.
正如《出埃及》裡說.
上帝對摩西說.
「我是自有永有的」.
而英文亦叫做「I am who I am」.
「自有永有」的意思就是.
「Yahweh」超然的本質.
祂的存在和身份.
是不依賴任何事物.
或者是任何的幫助.
祂是無始無終.
永不改變的.
而第三個.
「Yahweh」這個名字的含義.
是上帝願意與所有忠於祂的人立約.
雖然「Yahweh」的救贖.
祂的創造計劃.
是由祂與以色列民立的恩約開始.
但卻不是停止於此.
正如在《創世記》第十七章.
祂對亞伯拉罕很大聲地說.
「並且地上萬國都因你的後裔得福.
因你聽從了我的話」.
「萬國的後裔」指的是.
以色列民族以外的世界上的人.
只要他們願意忠心的敬拜.
也就是說「Yahweh」.
「Yahweh」的恩約就與祂有份.
你留意到嗎?.
原來藉著「Yahweh」這個名字.
上帝強調「自有永有的祂」.
已經履行了祂在伊甸園裡的承諾.
祂要修補和救贖.
這個被罪所破壞的創造.
包括醫好人類的重傷.
救贖在罪裡面的人類.

$^{121}$而祂沒有開空頭支票.
祂承諾過的事.
祂已經實現.
藉著「Yahweh」這個名字.
上帝呼喚你和我和祂立約.
祂說「凡敬拜祂的.
和忠心於祂的.
祂必定與我們同在」.
無論在多艱難的日子裡面.
祂都不會離開或撇下我們.
最後我很想問一句.
認識到這樣的「Yahweh」.
令今日面對人生各樣挑戰和衝擊的你.
有什麼安慰和鼓舞呢?.
所以特別送給你一個Fun Fact.
因為「Yahweh」.
耶和華是上帝親自向以色列人.
啟示自己的名字.
而猶太人覺得是神聖不可侵犯的.
因此不敢直呼這個名字.
只是用「子音」,「字母」.
這個來代表.
亦因為他們避免用口.
說出這個聖經的名字.
所以往往採取一些希伯來名號.
例如.
Adonai.
意思是「My Lord」.
或者是.
Hashem.
意思是「The Name」.
因為這樣.
這個名字.
原本的發音已經失傳了.
現在學界認為.
這個名字較為可能的發音.
就是「Yahweh」.
即是中文譯為.
「耶和華」.
英文譯為「The Lord」.

$^{161}$現在我們一起來祈禱吧.
「耶和華」「Yahweh」.
因為你用你的名字向我們說明.
你是一位多麼顧念我們的神.
當人類落在似乎不能逆轉的困境裡.
非常迷失和痛苦的時候.
你主動地答應.
還願救贖.
而你答應答應過的.
已經實現.
雖然這個實現.
是要你付上沉重的代價.
就是在你之下.
你的愛子耶穌基督為我們.
寫在十字架上.
維持我們衷心的讚美.
你是走若斯馳愛的神.
我們又求你赦免我們.
因為我們往往因為軟弱不忠.
我們會離棄了和你立的約.
叫你傷心.
在這裡我們願意一再立志.
忠於你和我們所立的約.
即使落在人生各種困境裡面.
我們仍然堅守不變.
我們的禱告是奉救主耶穌的名.
阿們.
\newpage



\section{}
\label{sec:Lh5hqFxfj84}
\textbf{Kingdom Values XXVI: Matthew Chapter 6 Continued - Pastor Sam Johnson}
\newline
\newline
連結: \href{https://youtube.com/watch?v=Lh5hqFxfj84}{\texttt{https://youtube.com/watch?v=Lh5hqFxfj84}} ~~~~ 語音日期: 2020-10-19
\newline
\newline
\hyperref[sec:O6HbgHZahfI]{\small{< < < PREV SERMON < < <}}
~
\hyperref[sec:index_chronic]{\small{[返順時目]}}
~
\hyperref[sec:index_scriptual]{\small{[返順卷目]}}
~
\hyperref[sec:XYSdVdRz0wM]{\small{> > > NEXT SERMON > > >}}
\newline
\newline
$^{1}$Good morning family. It's good to be with you again virtually. We're going to be back in Matthew.
chapter 6 this morning and continuing to talk about what Jesus is focusing on at the beginning.
of this passage and that is this theme of hypocrisy and how we as his followers, as subjects of the.
kingdom, as believers in Jesus should seek to avoid hypocrisy in our lives. And Jesus gives us some.
very challenging and convicting illustrations to help us see how we fall into hypocrisy..
And he lays out in the first half of this chapter three different illustrations and he points out.
how we can fall into religious hypocrisy in our giving or our good works, in our prayer,.
and in our fasting. And so he really covers a wide range of what we would think of as our.
devotion to God and how it shows us that it has to be from a right heart, from a right motivation,.
or it can become hollow and it can become something that is empty and something that is even.
full of hypocrisy, something that's pretending to be something that it's not. Last week we kind of.
just did an introduction. We looked at Matthew chapter 6 verse 1 where Jesus says, "Beware of.
practicing your righteousness before other people in order to be seen by them, for then you will.
have no reward from your Father who is in heaven." And we just kind of introduced this concept of.
hypocrisy and what it can look like in our lives. And we began to define hypocrisy in a biblical.
sense as the act of pretending to be something that we're not or acting as if we're something.
that we're not. We're pretty familiar with that definition. Actually the word hypocrisy that we.
use in English has roots all the way back in a Greek word that was the word for an actor or someone.
on stage. And so it carries that connection, that meaning of pretending or acting to be something.
that you are not. And that's a part of the word. But the other way we defined it biblically is that.
hypocrisy is when we deceive ourselves into thinking that we are meeting God's standards.
when we aren't and we can't and we don't. And hypocrisy is when we begin to think that we are.
satisfying God's standards, that we are righteous on our own, that we are right with God on our own,.
that we are doing everything that we should be doing in life apart from God. And we cannot do.
that. We cannot live rightly with God apart from God's help. We've talked about that so much over.
the past few weeks that Jesus wants to live through us and that's the only way that we can.
live rightly before God is when Jesus is living through us. And when we forget that and begin to.
think it's us who are going to make ourselves right with God, that we're going to earn something or.
that we're going to be full of good works on our own, then we begin to walk in hypocrisy. We begin.
to deceive ourselves into thinking we can meet God's standard when we can't and we don't..
Today we're going to continue on and look at verses two to four that kind of fill out this.
first illustration that Jesus gives us of hypocrisy. We'll read verse one again. He says,.
"Beware of practicing your righteousness before other people in order to be seen by them, for.
then you will have no reward from your Father who is in heaven. Thus, when you give to the needy,.
sound no trumpet before you, as the hypocrites do in the synagogues and in the streets, that they.
may be praised by others. Truly, I say to you, they have received their reward. But when you.
give to the needy, do not let your left hand know what your right hand is doing, so that your giving.
may be in secret, and your Father who sees in secret will reward you." Jesus is calling out.
our hypocrisy here and calling out the fact that we do so many things to be seen of others and to.
impress others. We discussed that a little bit last week that Matthew 5, 16 talks about letting.

$^{41}$our light shine before men. But what is the focus in that verse? It's so that others may see our.
good works and glorify our Father which is in heaven. That verse is talking about that light.
that shines. That is Jesus living through us. It's saying that as Jesus lives through us,.
it will draw people's attention to us and it will cause them to see God and to glorify Him..
But these verses are describing something totally different where we do what we think are good works,.
but they're selfishly motivated and they're selfishly focused and they ultimately put the.
focus on us. When people look at our selfish good works, they don't see Jesus, they see us. God is.
not glorified, we are glorified. And that's what Jesus is pointing out here when He says.
that in verse 1 that if we do things just to be seen by others, that that's not going to honor.
God. There's not going to be any reward from God. He says that in verse 2 that the hypocrites do.
things so that they'll be praised by others, so that other people will recognize them and say.
good things about them. And He says that they've received their reward, that the only reward.
they're going to get is the hollow praise of other people. And He says when you're giving,.
just to remove that temptation, do it in secret. Don't even let others know what you're doing to.
give away or what you're doing to care for someone else or what you're doing to share with someone.
else, because it's not about what other people see and perceive and think about you, it's about.
your heart before God. That's what Jesus is calling us to. We're going to be talking more and.
more about hypocrisy over the next few weeks because Jesus gives us so much time in this passage, so.
much time in this Sermon on the Mount. So as we talk about it, we're going to kind of break it.
out into a few different sections. And the topic that I want us to think about this week is how.
hypocrisy, and specifically religious hypocrisy, how maybe that has affected us and how maybe that's.
hurt us and how we can deal with that in our lives. So I want us to consider that today..
And I think when we consider religious hypocrisy and hypocrisy that we've encountered,.
maybe if we grew up in church or even if we didn't, just in some connection that we had to.
church or to people of faith, maybe we saw something that didn't quite line up where.
someone's words and their actions didn't agree. Maybe we saw something where someone asked us to.
live up to a standard that they weren't living up to. Or maybe we saw an even worse form of hypocrisy..
But the point is, those things may have hurt us, may have affected us in some way. And probably.
all of us have dealt with that in our heart on some level or another. And when we think about.
those things, we have to see that it starts with a misplaced standard. That humans, we misplace.
our standards of righteousness. We misplace what the source of truth in our lives is. And that can.
lead to really dangerous and disastrous consequences. You know, some people do practice a self-righteous.
and dishonest religion. And that's what the Pharisees were doing. They were practicing a.
religion that was not honoring to God. And I'm using the word religion intentionally because.
the Pharisees were not living out a relationship with God. They were practicing a ritualistic.
religion. And it was not honoring to God and it was self-focused and it was ultimately.
self-righteous. It was a way for them to make a distinction between themselves and others and to.
try to show people that they were righteous and somehow better than other people. And to try to.
earn favor with God and to try to do good works that they felt like would get them into God's.
favor and also ultimately into heaven. And that's not what Jesus calls us to in the scripture..

$^{81}$And unfortunately, a lot of people and even people in the church still try to relate to God in that.
same way. They try to live out a self-righteous ritualistic religion whereby they try to do a.
bunch of good works that they think will please God and do good works that they think will make a.
distinction between themselves and someone else and make them better than someone else and give.
them a sense of pride and righteousness in their goodness. And that's not what Jesus calls us to..
Jesus actually tells a story about this and it's recorded in Luke chapter 18. In verse 9, it says.
this, "He also told this parable to some who trusted in themselves that they were righteous.
and treated others with content." That verse right there could be a great definition of religious.
hypocrisy. Trusting in ourselves that we're righteous and treating others with content..
Jesus goes on to say this, "Two men went up into the temple to pray, one a Pharisee and the other.
a tax collector. The Pharisee standing by himself prayed thus, 'God, I thank you that I'm not like.
other men, extortioners, unjust, adulterers, or even like this tax collector. I fast twice a week,.
I give tithes of all that I get.' But the tax collector standing far off would not even lift.
up his eyes to heaven but beat his breast saying, 'God, be merciful to me, a sinner.' I tell you,.
this man went down to his house justified rather than the other. For everyone who exalts himself.
will be humbled, but the one who humbles himself will be exalted." Jesus gives us a perfect picture.
of religious hypocrisy here. He gives us this picture of this man who is doing a lot of.
external good works. From the outside he seems like the perfect guy. He's praying, he's fasting,.
he's giving tithes, but it's all for his own pride. It's all for the image that he puts forward to.
other people. It's not because his heart is in a place of humble devotion to God. And then there's.
this other man, this tax collector. In the first century Jewish culture, the tax collectors were.
hated people because they were people who were cooperating with the Romans to oppress their own.
people. So they were some of the most hated people in the society. They were considered.
totally corrupt and totally bad people. Jesus is saying, 'This guy who everybody thinks of as one.
of the worst people in culture and in society, he is being honest before God and he's saying, 'God,.
I can't even lift my eyes to look at you. I know how guilty I am. I know how broken I am. I know.
how much I need you.' And he's crying out in humility and that God actually hears and responds.
to this prayer and this cry of humble brokenness instead of this show of spiritual pride and.
self-righteousness that the Pharisee gives. And that still happens today. There are still people.
in the church, and maybe we've been those people at different times, who think that we've done.
something good and we've done a few right things in our life. And whether it's our church attendance.
or whether it's something that we've been a part of, we feel like it makes us better than someone.
else and we begin to operate in a hypocritical, self-righteous way. And that's wrong. That is not.
what Jesus wants for us. That is a misplaced standard. God's standard of holiness is a standard.
that applies to our heart, not a standard that applies to our external good works. Yes, when our.
heart is in the right place, beautiful good works flow out of that. It's not that our works don't.
matter, but it doesn't start with our works. It starts with our heart. And as humans, we can so.
easily fall into a self-righteous religion where it's all about the works. It's all about keeping.
the law. It's all about ticking the boxes. It's all about the external things that we do and say.
and don't do and don't say. And that's not what Jesus is calling us to. Another area that we see.

$^{121}$our misplaced standards is sometimes there are people put forward hurtful standards, impossible.
standards that are not God's standards. And we see that in the Pharisees. Clearly, that was one of.
the things that they were very good at. And Jesus calls them out for that multiple times in the.
Gospels. They were telling people to keep certain laws and customs and rituals and traditions that.
God never asked anyone to keep. And yet they were getting onto people for not keeping those laws or.
those customs or those traditions. And that happens today sometimes in the church. Sometimes people in.
the church choose to put a really strong emphasis on things that God never put an emphasis on..
And they call people out for maybe the way they dress or some type of other cultural influence..
They say, "Oh, you can't be a Christian and do this." Or, "You can't be a Christian and do that..
You can't be a Christian and listen to this music or that or whatever." And people get really caught.
up in these external emphasis. And they ask people to adhere to standards that are not God's.
standards. And that's another misplaced standard. We need to go to Scripture to find out what God.
says are the standards of His kingdom and not take standards from the culture around us and not allow.
even good and well-intentioned things in our lives to become standards that we try to enforce on.
someone else. God may call us to a standard of holiness or righteousness or to a specific.
action or to avoiding a specific thing in our life in order to be right with Him. But if that.
standard is not in the Scripture and if He didn't call everybody else to it, then we don't need to.
try to enforce it in other people's lives. And that's another way that we sometimes have misplaced.
standards in the church is trying to make people keep standards that are not God's standards..
And sometimes in the church, maybe we've experienced hurt from trying to keep an impossible.
standard. We talked about that a few weeks ago, Matthew 5, verse 48, where Jesus says we're to.
be perfect as our Father is perfect. We cannot keep that standard. And if we try to live up to.
that in our own strength, if we try to just be perfect and good and do everything right through.
trying really hard and through working at it as hard as we can, we will ultimately be disappointed.
and frustrated and crushed because we will not be able to do it. And yet some of us have done that..
And that's a part of the hypocrisy that sometimes happens in the church. People misunderstand and.
they forget about God's grace. Or maybe they never understand God's grace in the first place..
And they think that they're actually supposed to attain holiness through effort and through.
trying hard and through being something that they think they can be. And that was never what Jesus.
called us to. It's always been about God's grace. And Jesus made that clear in the beginning,.
that we couldn't keep that impossible standard of holiness, that it was only through Him,.
through His sacrifice, we can be forgiven, we can be cleansed. And as He lives through us,.
as we've been talking about, we can begin to live out the good works of the kingdom. And that's how.
the gospel works in our lives. But sometimes we forget that or we misunderstand it, and we try to.
live up to a standard that we cannot meet. And we get so frustrated and so hurt. Sometimes that.
leads to us walking away from faith or from church or from God and saying, "It's impossible..
It's too hard. It can't be done. I'm done with it." Jesus never asked us to live up to an impossible.
standard. He asked us to repent and be forgiven and let Him live through us. Sometimes our misplaced.
standards are in our understanding of God. Sometimes we're hurt because we misunderstand.
God. Maybe we think God is just some angry deity that is just looking to judge and condemn and.

$^{161}$hurt us. That's a misunderstanding of who God is. And obviously, that will cause us to behave in a.
way towards other people that is hypocritical and out of balance. Or maybe we think God is just some.
type of Santa Claus God that He just rewards good behavior. And if we just do the right stuff, then.
we get what we want in life. And He gives us the relationship we want and the job we want and the.
house we want and the financial situation we want. That's a wrong perspective of God. And if we start.
to believe that, even in small ways in our mind, then we fall into hurt because we say, "Wait a.
second. I did what I thought I was supposed to do. I lived out the 'good behavior,' and then I didn't.
get the present that I wanted. I didn't get the present I wanted for Christmas. I didn't get the.
relationship I wanted. I didn't get the job I wanted. I didn't get into the school I wanted.
to get into. I didn't get the house I wanted." But that's not who God is. And if we thought that was.
who God was or someone else told us that was who God was, if someone told us that God was just.
rewarding good behavior, if someone told us that God was just all about health and wealth and.
prosperity, if someone told us that God was just trying to give us everything we want, then that.
hypocrisy can hurt us. Because we could say, "Well, I thought God said my life was going to be easy..
I thought God said nothing bad would ever happen to me if I followed Him." That's not what God said..
Jesus said that He came that we would have life and have it to the full. And having life to the.
full sometimes means experiencing joyful blessing from the hand of God. But having life to the full.
sometimes also means experiencing loss and adversity and difficulty and allowing that to.
shape us and to grow us and to form us and to mold us. And God uses those things in our life..
And so sometimes the misplaced standard is that we have gotten some weird idea of God. And maybe.
we got it from someone in the church, maybe from a preacher or from a leader in the church who.
said something or taught something that wasn't really true to the scripture, and it caused us.
to get this wrong idea about God. And that form of hypocrisy caused us to be hurt. There are so.
many different ways that we can be hurt by hypocrisy in the church, by religious hypocrisy..
Another one maybe is just spiritual abuse. Maybe there was a leader in the church that we observed.
or that we were a part of or someone that we read something that they had written or we heard some.
of their teachings online or whatever that they were just wrong and they were abusing people..
Maybe they were manipulating people or maybe they were using people for financial gain or.
something else. There are those terrible things that happen in the church. There's stories of.
sexual abuse even in the church. And those things are terrible and they're completely wrong and they.
totally dishonor Jesus. But they do happen. First of all, there are people in the church who don't.
actually know Jesus. They're not actually followers of Jesus. They're not subjects of.
the kingdom. The Bible talks about that, that there are wolves among the sheep. There are.
people in the church who aren't truly in a relationship with God. And obviously people.
like that are going to abuse people and they're going to take advantage of things and they're.
going to try to take money from people or manipulate people. That's going to happen..
But there are even people in the church who are followers of Jesus but they're.
misunderstanding God's word or they're misinterpreting scripture in some way and.
that's causing them to do things that are damaging to others. Or there are people in the church who,.
maybe they're a follower of Jesus but they are in bondage to sin and they haven't gotten free.

$^{201}$from sin the way they should have. And because of their captivity to sin, they are taking actions.
that are wrong. And that happens in the church. There is hypocrisy in the church and it doesn't.
help us as believers to try to hide it or to try to cover up or to try to ignore it or to try to.
pretend it doesn't exist. That actually makes it worse. We need to be honest about it and we need.
to do what Jesus did and say, "Hey, this is not right. This is not the heart of the gospel for.
people to do things to be seen by others, for people to hide who they really are, for people.
to pretend to be something they're not, for people to stand up one day and act righteous and tell.
other people to be righteous and then turn around and do sinful or hurtful things to people. Those.
things are wrong and we need to be clear about that in the church. We need to call those things.
out. We need to lament over those things and mourn over those things and repent and turn away and.
turn our eyes back towards Jesus and towards following Him. Hypocrisy is so damaging and.
that's why Jesus is calling it out. So I want to encourage you if you've encountered hypocrisy in.
the church and if it's been hurtful to you, I want to encourage you. I want to say sorry. I.
want to apologize for that. I want to say as I just did that those things are real. They do happen.
in the church and it's wrong and it shouldn't happen and I'm sorry if that's happened to you..
But I also want to point you towards Jesus. You know, one of the things that we can have misplaced.
when we encounter hypocrisy is misplaced expectations. Sometimes we have an expectation.
of a leader or another person in the church and we kind of put them on a pedestal because we think.
they're going to do everything right. We think they'll never let us down and we think they're.
really close to perfect. And when we do that, when we have misplaced expectations, we will at some.
point encounter disappointment. People will always fall off of those pedestals that we put them on.
and so we shouldn't do that. We should understand that even followers of Jesus who are genuine and.
who love God with all their heart and are seeking to follow Him, they are still human and they're.
still going to make mistakes too. We shouldn't put them on a pedestal. Only Jesus is perfect. Only.
Jesus is never going to disappoint us. Only Jesus is always going to be with us and never leave us.
or forsake us. Humans cannot take that place, cannot take that role that only Jesus can be in..
And so we shouldn't put other people on a pedestal. Another way that our misplaced expectations.
express ourselves is sometimes we get caught up in the hypocrisy, right? Sometimes we maybe.
drink some of the Kool-Aid that someone else is drinking about how righteous they are and we start.
to drink some of our own Kool-Aid and say, "Yeah, I'm pretty righteous too. I think I'm really.
getting this figured out and I think I'm really becoming a really great Christian and I'm doing a.
fantastic job. I know I'm doing way better than the people around me so I'm in pretty good shape.".
We can start to feel that way and we can misplace expectations in ourselves. And when we do that,.
we're going to experience deep hurt and frustration because sooner or later we are going to fall off.
that pedestal and we're going to see our own brokenness and our own sin. And we can't hide it.
from ourselves because it's in our heart. We may hide it from someone else. We may act and pretend.
in front of others that we don't have any struggles but we know our own struggles and.
we'll feel ourselves falling away from God and we'll say, "I'm not who I want to be." And we'll.
be crushed by that misplaced expectation in ourselves. We were never supposed to expect.
ourselves or anyone else to live the perfect life. Jesus is the only one who can do that..

$^{241}$What we should expect of ourselves and others, including leaders in the church, is to live a.
life that is surrendered to the rule and reign of God. To live a life that is surrendered to Jesus.
and to let Jesus live through them. That's what we can expect. And to be someone who is willing.
to repent when they make mistakes and who is willing to be forgiven by Christ and to be.
cleansed. That's what we can expect. But we should not expect ourselves to be on some pedestal of.
perfection because if we do, we will fall off and we'll be hurt by that fall. So how do we deal.
with the pain of hypocrisy? If we have been hurt in some way by self-righteous religion and we've.
experienced the pain of that, how do we deal with that? And how do we deal with that in a way that.
allows us to continue to have a soft heart towards God and a soft heart towards relationship with God.
and a soft heart towards God's church? I think one of the first things that we need to think about.
and understand is that hypocrisy often causes us to have a misplaced pain. What I mean by that.
is that sometimes we are hurt by things that happen in the church. We are hurt by things that.
other Christians or people who profess to be Christians may inflict on us. And we misplace.
that pain onto Jesus or onto God and assume that we've been hurt by God or we've been hurt by Jesus..
And that's simply not true. Jesus is not the source of self-righteousness and hypocrisy..
Jesus is not the source of our pain. Jesus loves us and desires good for us. But so often in our.
culture and in our own minds, we create these caricatures of Jesus where we think Jesus is.
something that he's not. We have an image of Jesus that's different from the reality of who Jesus is..
You know, in our culture, people blame Jesus and Christianity and God for basically every evil of.
our culture, right? People sometimes think Jesus is some type of agent of European imperialism,.
or they think Jesus is a source of discrimination, or that Jesus didn't care about women, or that.
Jesus promotes tribalism, or that Jesus promotes conflict between people groups. And people accuse.
Jesus of all these different things. Those are caricatures of Jesus. Those are images that are.
created that are not true to the reality of who Jesus is. And if we want to know the true Jesus,.
we need to look at the Word of God and see who Jesus actually was and what he actually said..
And when we do that, we see something very different from the caricatures of Jesus that.
we see in our culture and that sometimes exist in our mind. Jesus was not European. He was Jewish..
He was from the Middle East. Jesus was not pursuing political power in any way. He was.
full of spiritual power and he was pursuing the healing of human hearts. Jesus deeply cared about.
women and elevated them to the place that God originally gave them, which is equal in worth.
and essence and value to men. Jesus was clear about that. Jesus actually called people out for.
their prejudice and their racism and their discrimination. Jesus spoke against it. Jesus,.
as he is in this very passage we're looking at today, called people out for their hypocrisy.
and their pride and the distinctions that they made between themselves and other groups of people..
Jesus called people out for their ethnocentrism. Jesus did all of these things. Jesus was not who.
we often make him out to be in our mind. And Jesus was certainly not some type of angry or unkind God.
that wanted to inflict pain on people or condemn them or kick them out. That's not who Jesus was..
But sometimes when we've experienced that in the church or we've experienced that from.
other people who are Christians or profess to be Christians, we think that that came from God or.
it came from Jesus. And that's simply not true. If we go back to the Word and examine who Jesus is,.

$^{281}$what he said, what he did, we find that he loves us and his message was one of grace, forgiveness,.
and calling us to live in right relationship with God, with the one who created us..
Now, yes, God has standards and those standards are clear in the Bible, but it's also clear in.
the Bible throughout the Bible, but especially in the ministry of Jesus, that those standards.
are for our own good. Those standards are not to hurt us or to restrict us. They are to protect us.
and to care for us. Jesus is the good shepherd. He is caring for us and he loves us. And it's clear.
throughout the ministry of Jesus that the grace of Jesus is the only hope that we have for meeting.
God's standard. We are never supposed to try to be perfect and work it out and try hard enough and do.
it all on our own. That was never the emphasis. We are supposed to come to Jesus in humility and say,.
"God, I can't do this. I ask you to forgive me. I ask you to cleanse me. I ask you to live through.
me and make me into something better than what I am." That's the design from the very beginning..
That's the gospel. We were never supposed to turn into some self-righteous super person who's better.
than someone else. That's a caricature that we've made in our own minds. John 1 is clear that Jesus.
is full of grace and truth. And I love that about Jesus. Jesus never lied to us and told us we were.
any better off than we actually are, but he also never offered anything less than complete redemption.
and restoration to a perfect relationship with our Father in heaven. That's who Jesus is. He always.
loves. He always cares. He always heals. And so if you've experienced hurt and you felt like maybe.
that was coming from God, it wasn't. If you felt like that was coming from Jesus, it wasn't. Jesus.
cares for you. He wants to heal you and he loves you. So I would encourage you, if you have.
experienced that hurt, to understand that that wasn't coming from God. It was probably coming.
from an imperfect person who inflicted that hurt on you, probably because of some hurt in their own.
life. Or maybe because they completely misunderstood God's grace. Or maybe they completely.
misinterpreted God's truth in the scripture. And their own confusion, their own hurt, and their own.
pain caused them to communicate something to you that simply wasn't true. That did not come from.
God. And then I would encourage you also to seek to practice forgiveness. Forgiveness is a totally.
radical kingdom value that Jesus taught us. And if you've experienced hurt because of hypocrisy,.
it can be really hard to forgive someone who has been a hypocrite, who has pretended to be.
something that they're not, and pretended that they were better than someone else. That can be.
really hard to forgive. But that's what Jesus calls us to. Jesus calls us away from the old.
human way of vengeance and getting people back, and calls us to forgiveness and to love. And so.
I would encourage you to seek that. Even if you've experienced something difficult, like spiritual.
abuse in some way, or some other type of abuse at the hand of a hypocritical person, that's.
something that you can seek forgiveness for. You may need to do some counseling and deal with that,.
especially if it was something serious. And there may be consequences of that. Maybe that.
relationship with that person will never be repaired. Or maybe that person, if it was serious.
enough, may actually suffer other consequences in some way. Some people have to actually suffer.
legal and civic consequences because of the wrong things that they've done. But that doesn't mean.
that you can't forgive someone. And that's what Jesus calls us to. To understand that that pain.
was not inflicted by God. It was not inflicted by Jesus. And that even the person who inflicted that.
pain was full of their own brokenness and their own sin. And we need to forgive. And that can be a.

$^{321}$difficult thing. It can be hard for us to do, but it gives us freedom as we forgive hypocrisy in.
others. Let's pray. God, thank you that you are so honest with us about our brokenness as humans..
And God, I pray for anyone who has been hurt by hypocrisy in the church, God. God, I pray that.
you would help us not to ignore it or to act like it didn't happen or to try to explain it away. God,.
help us to be honest about how wrong it is and how devastating it can be. And God, if we've been.
hurt by it, God, I pray that you would help us to see that that did not come from you. It came from.
another human who was hurt and broken in their own life. And God, help us to not place that pain on.
you and understand that you love us and you want healing for us. God, help us to take that pain to.
you. It is real pain. Help us to take it to you and let you heal and restore in our heart. And God,.
help us to find forgiveness. Give us the power to forgive. Even if someone has hurt us deeply,.
God, even if there are consequences of that hurt that may not change, help us to forgive.
and be free of that load of bitterness or hurt or frustration or pain that we carry around. God,.
God, I pray that you would help us to more and more move away from pretending to be something.
we're not, pretending that we meet your standards when we don't, and move towards a place of.
honesty over our own brokenness, our own need for you. Move us towards a place of poverty of spirit.
where there is no hypocrisy in us, where we are transparent and open and honest, and where we let.
you live through us in a way that shares your love and your grace with the people around us..
In Jesus name, Amen..
\newpage



\section{}
\label{sec:XYSdVdRz0wM}
\textbf{浩劫重生 【第4集:出埃及路線】 陳崇基牧師 - 主日講道視頻系列}
\newline
\newline
連結: \href{https://youtube.com/watch?v=XYSdVdRz0wM}{\texttt{https://youtube.com/watch?v=XYSdVdRz0wM}} ~~~~ 語音日期: 2020-10-20
\newline
\newline
\hyperref[sec:Lh5hqFxfj84]{\small{< < < PREV SERMON < < <}}
~
\hyperref[sec:index_chronic]{\small{[返順時目]}}
~
\hyperref[sec:index_scriptual]{\small{[返順卷目]}}
~
\hyperref[sec:c2n9YZYOe80]{\small{> > > NEXT SERMON > > >}}
\newline
\newline
$^{1}$上一次我們提過出埃及這件事.
是在公元前1500至1200年間發生.
當時的世界都是面臨史無前例的動盪.
近東有很多人正在遷徙移居.
想找一個安居之所.
我們舉個例子就是海人Sea People.
其中一組人就是非利士人Philistines.
在大歷史圖畫中以色列人是其中一群在遷徙.
他們在埃及被欺壓多年.
離開時前無去路後有追兵.
但有特別的地方是他們發現上帝親自帶領他們.
對於今天你和我面對社會政局的不穩.
有一種被迫的感覺流散在世界不同的地方.
同樣面對很多改變身份改變和未知之數.
這樣看出埃及這個故事對我們來說是非常切身.
但當我們看出埃及這個技術時.
我們一定要梳理一下歷史的問題.
為什麼呢?因為在舊約聖經中.
經常提到出埃及這件歷史事蹟.
可以這樣說在舊約歷史中出埃及這件事.
就有如新約的耶穌基督從死裡復活.
歷史性都是這麼重要.
但在過去幾十年越來越多人質疑.
為什麼沒有證據證明出埃及這件事.
或者可能出埃及是很後期虛構的故事.
我記得好幾年前在香港電台上受過訪問.
主持人問了一個問題很有道理.
他說出埃及這麼大件事總會有少少線索.
居然一點線索都沒有.
有些道理我們必須承認.
暫時到現在都沒有直接的歷史證明.
但我們這樣說也好.
我們都需要留意幾點.
第一就是證據的不存在.
不一定等於不存在的證據.
這個你動高床板就要想一想.
The absence of evidence is not the evidence of absence.
因為很多歷史事跡.
包括聖經或聖經以外的歷史事跡.
都未必找到考古的證據.

$^{41}$我們不能夠踏時光機回去以前.
將歷史放在化驗室裡面去證明多一次.
何況當時的背景.
埃及帝國是獨裁的政權.
他們不會將失敗的事跡寫下.
第二件事我們需要留意.
當我們想起考古證據的時候.
我們對考古證據的要求一定要實際.
假設以色列人出埃及的話.
你想一下他們會帶著什麼.
他們好像難民一樣離開埃及.
我不知道你有沒有做過服務難民的工作.
我記得有一次我參與一個難民服務的工作.
去到機場接一家六口從剛果難民營來的難民.
六個人他們入境下機的時候.
每個人你想一下他們帶著什麼.
他們只是帶著一個小的膠袋.
就是這麼多.
所以你再想一想.
作為難民甚至作為平民百姓也好.
過了一百年之後.
你看我們這班平民百姓的人.
都一定是蕭聲匿跡的.
就是說在考古歷史裡面是看不到我們的.
所以為什麼有這麼多皇帝想建皇陵.
大墳墓有兵馬俑金字塔等等.
因為他們想留芳百世.
所以沒有考古證據不代表你不存在.
再有一點需要留意.
如果你真的按照聖經所說.
以色列人在埃及裡面住了好幾代.
如果他們出埃及的時候.
你想一下他們帶著什麼實物出來.
衣服,瓷器,日用品等等.
但這些全部的東西.
都必定是埃及製造的.
是在埃及製造的.
所以他們到曠野就算留下這些實物.
我們找到的.
我們只會當有一班埃及人.

$^{81}$他們經過留下過.
舉個很簡單的例子.
我的子女都是在美國出生長大的.
我看遍他們擁有的實物.
我找不到哪一樣可以直接表明他們是華裔第二代.
除了他們現在在生的樣子.
像亞裔的人.
但他們說話,吃東西,想的東西.
都是美國本土化的.
所以我想如果他們真的化了灰的話.
他們的實物跟其他白人,黑人.
這裡的美國人是沒有任何分別的.
所以我們對考古證據的要求.
是需要實際一點的.
我們不會找到一塊石頭.
上面寫著摩西到此一遊.
你說沒有直接的歷史考古證據.
是不是就這樣算了.
其實也不是的.
在歷史的探索裡面.
我們再問一個問題.
有沒有間接性的,情況性的,背景的旁證.
即是circumstantial background的evidence.
如果這樣問的話.
可能真的有.
因為如果出埃及記.
真的是後期虛構著作.
即是到波斯時代,公元前500年的時候才寫.
他的描述裡面一定會用一些後期的字眼.
但是如果出埃及記.
真的寫當時的歷史事跡.
他的描述一定是配合當時的語言.
當時的環境.
如果是這樣的話.
歷史可靠性就會高很多.
你問前前後後.
又說早期又後期的描述.
好像很難明白,很抽象.
其實一點都不難明白.
我播一下一個短片.

$^{121}$你就知道有什麼問題.
是,駱Sir.
失憶的大軍要二十天才到達白里浦.
但每日的軍餉二十萬兩.
我們的國庫已經所剩無幾.
浸暫時不能給太多.
為今之計.
只要交了十五,十六和十七這三天的軍餉.
其他的就慢慢想辦法.
駱Sir,讓盧平替你換一盆熱水.
周兄長,跟浸再檢查一下.
看各地各省還有沒有人欠稅.
浸要把水行一次過追回來.
實不相似.
明眼人一看就明白.
因為浸是不會用check這個字.
看起來好像很好笑.
但其實這裡點出了一樣東西.
就是如何鑒定一個歷史文件可靠性的原則.
古話說出埃及記.
如果歷史可靠的話.
它裡面的描述一定是吻合公元前1500至1000年之間.
那時是New Kingdom Period.
是古埃及帝國的新王國時期.
那時是埃及最後輝煌的時期.
看當時埃及的歷史文件.
用什麼字眼.
描述什麼地理.
然後再比對一下.
出埃及記我們就知道歷史的可靠性.
舉幾個例子.
出埃及記一章第十一節.
於是埃及人派督公管轄他們.
即是以色列人.
加重他們的負擔和害他們.
他們為法老建造兩座積火城.
就是彼東和蘭室.
使他們因做苦工覺得命苦.
無論和禮 造磚 做田間各樣的工.
在一切的工上都嚴嚴待他們.

$^{161}$這裡細緻的描述.
造磚 和禮和一切的工.
我帶你去一個埃及的墓穴.
叫做Tomb of Rekhmire.
雷克米爾墓穴.
這個人是公元前15世紀.
即是1400多年的時候.
在埃及做大官.
在他的墓穴有很多壁畫.
裡面描述很多奴隸.
包括戰俘.
他們在田間做很多事.
這個跟出埃及記的描述非常吻合.
你看到圖畫這一部分.
有人在池裡取水.
有些奴隸在和禮.
有些在造磚.
包括田間各樣的工作.
細緻地看的時候.
埃及的督工是拿著木棒來管轄他們.
另一方面我們可以看地理.
一章第十一節.
兩座積火城.
叫彼東和蘭室.
蘭室是法老的名字.
有一個很出名的法老王.
叫蘭室二世 或蘭室大帝.
Ramesses II.
蘭室城真的存在.
Ramesses 或 Ramesses.
在舊約聖經裡只提過五次.
都是在摩西五經裡提過.
而我們知道蘭室城.
Pyramuses 城市.
在埃及的文獻裡出現過.
也是蘭室二世的法老所建的京都.
說起蘭室二世的法老.
他就很有名.
因為很多人都認為.
這個很可能是出埃及的法老.

$^{201}$但說到考古.
蘭室城我們知道的.
只有二百年的歷史.
由公元前1270至1070年之間.
為什麼最後會被荒廢呢.
因為當時利羅河在三角洲平地經常改道.
利羅河對埃及人來說.
是很重要的.
等於今天的高速公路.
交通 貿易 軍事的輸流都是靠利羅河.
所以利羅河改道.
城市在兩百年之後就被荒廢.
埃及人很會復原.
他們就把本身在蘭室城的石碑.
有蘭室二世的名字的石碑.
搬到北面十二英里.
建另外一個城鎮.
那個城鎮叫索安.
Zoan 或希臘文是丹尼斯.
說起索安城或丹尼斯.
如果你記得或接觸過.
1981年有一套電影叫《奪寶奇兵》.
In the End of Jones' Waders of the Lost Ark.
裡面提及過丹尼斯的古城.
電影裡的橋段是二次大戰時.
德國納粹黨要法國考古學家.
去丹尼斯索安找藥櫃.
你看看以下短短兩三分鐘的片段.
你告訴我有什麼真有虛構.
Dr. Jones, I've heard a great deal about you..
Have you?.
Professor of Archaeology,.
Expert on the occult,.
and I just want to say it,.
obtainer of rare antiquities..
One way of saying it,.
why don't you sit down,.
you'll be more comfortable..
Thank you..
Yes, you're a man of many talents..

$^{241}$You studied under Professor Ravenwood.
at the University of Chicago..
Yes, I did..
Dr. Jones, you must understand.
that this is all strictly confidential..
I understand..
Yesterday afternoon,.
our European sections intercepted.
a German communique.
that was sent from Cairo to Berlin..
You see, over the last two years,.
the Nazis have had teams of archaeologists.
running around the world.
looking for all kinds of religious artifacts..
Hitler's a nut on the subject..
He's crazy..
He's obsessed with the occult..
And right now,.
apparently there's some kind of.
German archaeological dig.
going on in the desert.
outside of Cairo..
We've got some information here,.
but we can't make anything out of it,.
and maybe you can..
Tannis Development Proceeding.
Acquire Headpiece.
Staff of Ra.
Abner Ravenwood, U.S..
Nazis have discovered Tannis..
What does that mean to you, Tannis?.
Well, the city of Tannis.
is one of the possible resting places.
of the Lost Ark..
The Lost Ark?.
Yes, the Ark of the Covenant,.
the chest the Hebrews used to carry.
around the Ten Commandments..
What do you mean, Ten Commandments?.
You're talking about the Ten Commandments?.

$^{281}$Yes, the actual Ten Commandments,.
the original stone tablets.
that Moses brought down out of Mount Harab.
and smashed, if you believe in that sort of thing..
Did you guys ever go to Sunday school?.
Well, I....
No, look..
The Hebrews took the broken pieces.
and put them in the Ark..
When they settled in Canaan,.
they put the Ark in a place.
called the Temple of Solomon in Jerusalem,.
where it stayed for many years.
until all of a sudden, whoosh, it's gone..
Where?.
Well, nobody knows where or when..
However, an Egyptian pharaoh....
Sisak..
Yes, invaded the city of Jerusalem.
around about 980 B.C.,.
and he may have taken the Ark.
back to the city of Tannis.
and hidden it in a secret chamber.
called the Well of Souls..
Secret chamber?.
However, about a year after.
the pharaoh had returned to Egypt,.
the city of Tannis was consumed by the desert.
in a sandstorm which lasted a whole year,.
wiped clean by the wrath of God..
What does this Ark look like?.
I don't know..
There's a picture of it right here..
That's it..
Good God..
Yes, that's just what the Hebrews thought..
What's that supposed to be coming out of there?.
Lightning, fire, power of God or something..
You don't understand Hitler's interest in it..
Oh, yes..

$^{321}$The Bible speaks of the Ark leveling mountains.
and laying waste to entire regions..
An army which carries the Ark before it.
is invincible..
What is the truth?.
If you've been to Sunday School,.
you'll know that there's a Ark..
And you'll know that.
900 years before the Temple was built,.
there was a pharaoh named Shishak..
He really attacked Jerusalem.
and took away the treasures.
from the Temple and the palace..
This is what the 14th chapter of the Book of Kings.
has told us..
The Ark's whereabouts.
is another big topic..
We won't talk about it today..
But anyway,.
when it comes to Tennis,.
Denis, Soin, Zouan,.
these cities,.
it's true that since 1928,.
the French archaeologists.
have been excavating..
They've been doing it for years..
In the early years,.
the French archaeologists.
thought that Soin, Zouan, and Tennis.
were the epicenter of the Egyptian War..
They found the remains of the dead.
and the stone tablets..
But now they know that it's not true..
Soin and Tennis.
are an ancient city of 1,000 years..
But they were built in 1070..
So it's a post-historical city..
The later books of the Bible.
also mention Soin and Tennis..
In the 12th chapter of the 78th chapter of the Book of Psalms,.

$^{361}$it says that God was in Soin,.
the fields of Soin,.
the fields of Tennis,.
and was doing a miracle.
in front of the Pharaoh..
This is very important..
The 78th chapter of the Book of Psalms.
is about the plague in Egypt..
But it says that the place where the plague happened.
was Soin, not Nanset..
So when the poet wrote this poem,.
he recalled the plague..
He wrote it after 4,500 years.
after the plague..
At that time, Nanset had already ceased to exist..
So the poet only knew.
that the plague had happened.
in the city of Soin near Nanset..
But when you go back to the time of the plague,.
you see the details of the geography.
mentioning Nanset..
This town really existed at that time..
So when it was recorded,.
the reliability of history was shown..
You may think it is very complicated..
But it is not complicated at all..
It is also very simple..
For example,I have lived in Chicago for 20 years..
I only know the geographical location and geography.
of Chicago in the past 20 years..
You can ask some old Chinese here..
They can tell you that.
40 or 50 years ago,.
the name of the restaurant in the same place.
was "Meng"..
Today we have really determined.
the location of Nanset..
It is a small village in Alibaba.
called Cantia..
Here we find the stone wall of Nanset..

$^{401}$It also has its royal name..
There are Nanset's great soldiers on the stone wall..
They are attacking its enemies..
It is said that the Bible often mentions.
the new-made shield of the great man..
It is the political idiom.
that the French used at that time..
In Nanset,we also find.
the military horse-drawn carriage..
There are also many large stones for horse-drawn carriages..
It is estimated that the horse-drawn carriage.
can stop 500 horses.
and more than 200 chariots..
Nanset is estimated to be.
as large as 10 square kilometers..
Now it is a field..
But it was the largest city in the world at that time..
The Egyptian Bible also contains.
many special idioms and words.
that the Egyptians used at that time..
But there is no word for the Persian period..
On the contrary,the books written later.
such as the Book of Ezekiel and the Book of Nehemiah.
often contain Persian words..
We have just spent about 15 minutes.
to talk about the historical reality.
of the Egyptian Bible..
We don't need to doubt it..
Now we can see.
the Egyptian route of the Israelites..
How did God lead the Israelites to Egypt?.
They set out from Nanset..
In the 13th chapter of chapter 17,.
it is mentioned that when Pharaoh went there with his people,.
the roads to Philistine were narrow..
God did not lead them where to go..
Because God said,.
"I am afraid that when the people see the battle,.
they will regret it and turn back to Egypt.".
So God led the people to go astray.

$^{441}$and went the way of the Red Sea..
Here is a road called the road to Philistine..
This road was also mentioned.
in the Egyptian documents at that time..
The other name of Egypt is.
the Way of Horus..
It is the way of Horus, the Egyptian god..
The way to Philistine was the most convenient way..
It took a few days to reach the land of Philistine..
But the Bible says.
God did not lead the Israelites to the fastest way..
Why?.
Because God said,.
"I am afraid that when the people see the battle,.
they will regret it and turn back to Egypt.".
This is a good example..
On the road to the northwest of the Mediterranean,.
we found a lot of soldiers from the New Kingdom..
More than ten soldiers..
There were Egyptian war horses on this road.
to guard the land..
It was a military stronghold..
So the Bible tells us.
that if the Israelites took this shortest way,.
they would not have been well equipped.
and would not have fought well..
They would have regretted it and turned back to Egypt..
The way God led people to go.
was not necessarily the shortest and fastest..
It was not necessarily our plan..
Sometimes it was not even in line with the law..
The other way to Egypt.
was the way to the shore in the Bible..
This is what we saw in Genesis..
Do you remember the story?.
When Haggai went out,.
he walked on the way to the shore..
So the Israelites started from Lancester.
and walked to the southeast..
According to the Bible,.

$^{481}$they should have gone out of Egypt through the Shu'i Basin.
and then went to Galilee..
But what did the 14th chapter say?.
It said, "And suddenly there was a great storm..
And God commanded the Israelites to turn back.".
It means to go north..
They went to the sea of ​​red..
It was actually in front of the sea of ​​grass..
The Israelites were afraid of the storm..
So they turned back to Egypt..
So they walked on the way to the south..
God suddenly called them to turn back..
Why?.
When they were about to go out,.
God called them to turn back..
What was God's will?.
In the 14th chapter, verse 3,.
God said, "And Pharaoh said,.
'The Israelites were lost on the ground,.
and the wilderness had trapped them.'.
In the 14th chapter,.
God said, "I will make Pharaoh's heart strong,.
and he will chase them..
And then I will be glorified.
in Pharaoh and his whole army.".
The Egyptians knew that they were in the wilderness..
So the Israelites walked on..
God's command seems to be very complicated..
It's not humanly reasonable..
But what is the ultimate goal?.
It is to be glorified to God..
God's purpose is not to make us feel sad and empty..
It's not that people think.
they are doing well and that's it..
God's command is needed.
to make us learn to trust Him..
At the beginning, we said that.
there are many changes in life..
Especially in the pandemic,.
we have encountered many unideal situations..

$^{521}$We all want to go to an ideal place..
We always say that we don't like this situation,.
don't like this place,.
don't like this context..
So we want to walk away..
We think that's OK..
Of course, this is humanly reasonable..
But this story seems to tell us.
what is the most important thing in God's eyes..
The most important thing is not what we think..
Everything is fine and good..
In chapter 14, verse 30,.
it says, "On that day, Jehovah saved the Israelites.
and took the hands of the Egyptians.".
It says that they crossed the Red Sea..
What did the Israelites see?.
They saw Jehovah's victory..
What did the Israelites see?.
They saw Jehovah's great deeds to the Egyptians..
They knew how to glorify Him,.
how to trust God and His divine power..
When we encounter difficulties,.
we think that we are leaving this unideal situation..
We are looking for an ideal place..
If God really wants to lead us to our ultimate goal,.
without disasters, without worries, without sickness,.
we can retire and live happily,.
I think God is the most unsuccessful prophet..
The most successful one should be Jehovah..
Because in the Egyptian Empire,.
if you put your trust in Him, you can live in peace..
So, you see, when the Israelites left Egypt,.
they were not able to enjoy their eternal peace..
In fact, after the Israelites left Egypt,.
they crossed the Red Sea and the Sea of ​​Grass..
It was a trial..
Because God's purpose was to prepare us.
to become a group of people who were faithful to His will..
This is the greatest blessing..
The greatest blessing is not that we have no disasters, no sickness,.

$^{561}$no epidemic..
The greatest blessing is that we can become the children of God,.
and be with Him..
We can be so glorious.
that we can become people who serve God..
We are now facing different disasters..
Whether it is an epidemic or a downturn in life,.
we may need to ask.
whether we will still be loyal to God.
after these disasters and downturns..
We just need to be there to answer the call..
We have not learned to be more faithful to God..
Or when we have passed the command of God,.
we will be more faithful to God..
Personally, the same is true for the Church..
God said, "Do not waste the path He has led us on.".
No matter what the disaster is,.
let us, the Church,.
be able to experience new lessons..
Let us, the Church,.
become a Church that is more in line with His will..
Some scholars said it was very interesting..
They said, "The main thing that God wants to do in the Exodus.
is to take Israelite out of Egypt.".
God did not want to take the Israelites out of Egypt..
After they left Egypt,.
they still had a heart for Egypt..
They wanted to return to the Egyptian Empire.
because there was a stable and prosperous place..
God wanted to take Egypt out of the Israelites.
and take the Pharaoh out of Egypt.
and take away the idols.
like Moses,.
who was a prosperous man.
and a man of peace..
So that we can truly return to God.
and be faithful to Him.
and serve Him,.
and be faithful to Him.
Are we willing to do that?.

$^{601}$Let God take away these idols.
and let our hearts truly return to Him..
Let us pray with our hearts..
God, we know that Your way.
and Your great command.
is not what we humans think..
Your ultimate goal.
is to make our lives belong to You..
God, we pray that You take away.
the things we long for in our lives.
like stability and prosperity.
and the benefits that the land can give us..
God, let our lives truly belong to You..
Let You be our true Lord..
You continue to lead our lives.
and the path ahead of us..
No matter what we are going through,.
let us know that the ultimate goal.
is not only to be free and carefree.
but to be faithful to You..
Continue to train us,.
continue to help us,.
and nurture us.
so that we can become.
more fit for Your will..
Let us pray..
God bless you..
《神的道路》.
作詞:李宗盛 作曲:李宗盛.
神的道路.
高過人的路.
神的一念.
高過人的一念.
他的心裡有藍圖.
他的世界不作物.
一步一步帶領的前途.
他離大地的更近.
他讓雲在空中漂浮.
是他創造生命氣息.
是他隱世照顧自負.

$^{641}$神走萬物和愛棲息.
愛心等待你要歡呼.
啊 愛心等待要歡呼.
神的道路.
高過人的路.
神的一念.
高過人的一念.
他的心裡有藍圖.
他的世界不作物.
一步一步帶領的前途.
他離大地的更近.
他讓雲在空中漂浮.
是他創造生命氣息.
是他隱世照顧自負.
神走萬物和愛棲息.
愛心等待你要歡呼.
啊 愛心等待要歡呼.
啊 他離大地的更近.
他讓雲在空中漂浮.
是他創造生命氣息.
是他隱世照顧自負.
神走萬物和愛棲息.
愛心等待你要歡呼.
啊 愛心等待要歡呼.
啊 要歡呼.
「愛心等待要歡呼」.
\newpage



\section{6:5-6}
\label{sec:c2n9YZYOe80}
\textbf{Kingdom Values XXVII: Matthew Chapter 6\_5-6 Continued - Pastor Sam Johnson}
\newline
\newline
連結: \href{https://youtube.com/watch?v=c2n9YZYOe80}{\texttt{https://youtube.com/watch?v=c2n9YZYOe80}} ~~~~ 語音日期: 2020-10-28
\newline
\newline
\hyperref[sec:XYSdVdRz0wM]{\small{< < < PREV SERMON < < <}}
~
\hyperref[sec:index_chronic]{\small{[返順時目]}}
~
\hyperref[sec:index_scriptual]{\small{[返順卷目]}}
~
\hyperref[sec:y_7ZwDcfiBI]{\small{> > > NEXT SERMON > > >}}
\newline
\newline
$^{1}$Alright family, we are back and we are once again in Matthew chapter 6..
Continue to look at this kingdom values study that we've been going through for a while..
Looking at the values of the kingdom of heaven that Jesus wants us, if we are his disciples,.
his followers, if we are subjects of the kingdom, then he wants us to live out these values..
In the past few weeks we've been talking about hypocrisy and how Jesus in the first part.
of Matthew 6 is calling us to avoid hypocrisy..
He's pointing out how damaging and devastating it is in our lives..
He's also pointing out a little bit just how silly and foolish it can be for us to live.
as hypocrites and yet it's clearly something that all of us as humans do..
And as we talked about a few weeks ago, as humans we're all experts at it..
We all are hypocrites..
That hypocrisy can also easily come into our life as followers of Jesus..
It can easily come into our faith and when it does it can have absolutely devastating.
consequences..
And last week we talked some about that..
How if maybe we've been hurt by hypocrisy in the church, maybe we've been hurt by someone.
else who treated us with hypocrisy or someone who lived hypocrisy in front of us and maybe.
it damaged our faith in Jesus or maybe it just hurt us in some way..
We talked about ways that we can get past that..
We can redirect our focus back to Jesus and we can begin to experience forgiveness even.
for those who have hurt us..
We can begin to take those misplaced standards and misplaced expectations and misplaced pain..
We can see that Jesus is not the source of our pain..
Jesus is not the source of hypocrisy or self-righteousness..
That's never something that Jesus called his people to live out..
Jesus wants us to live out the opposite..
He wants us to live out authenticity and transparency and honesty and love for one another..
And we can learn to walk in forgiveness..
And this week we're going to talk about something that's probably going to be a little bit more.
convicting for us, a little bit more challenging..
It's going to hit a little bit closer to home for us and that is the hypocrisy that we exhibit.
in our own lives..
Last week we were more talking about hypocrisy that someone else has exhibited and how that.
might have affected us..
But this week we're going to focus on our own hypocrisy and how we live it out and how.
we are dishonest about the fact that we're living it out and how much it hurts us and.
hurts the people around us..
So let's jump into this here..
We're looking specifically at verses five and six of Matthew chapter six today, but.
we're going to go ahead and read the first four verses as well just to give us kind of.

$^{41}$a big picture of what Jesus is saying here..
So it says this, "Beware of practicing your righteousness before other people in order.
to be seen by them..
For then you will have no reward from your Father who is in heaven..
Thus, when you give to the needy, sound no trumpet before you as the hypocrites do, in.
the synagogues and in the streets, that they may be praised by others..
Truly I say to you, they have received their reward..
But when you give to the needy, do not let your left hand know what your right hand is.
doing, so that your giving may be in secret, and your Father who sees in secret will reward.
you..
And when you pray, you must not be like the hypocrites, for they love to stand and pray.
in the synagogues and at the street corners, that they may be seen by others..
Truly I say to you, they have received their reward..
But when you pray, go into your room and shut the door and pray to your Father who is in.
secret, and your Father who sees in secret will reward you..
Now as we've discussed hypocrisy the last few weeks, we've kind of defined it out of.
this passage..
And we've defined it in a few ways..
We've defined it as pretending to be something we're not, or acting as if we are something.
that we're not..
And we've defined it as deceiving ourselves or others into thinking that we can meet God's.
standard of holiness, that we can be who God wants us to be, when we can't meet that standard.
and we don't..
Not on our own, not without Christ..
With Christ living through us, we can, because Jesus lived the perfect life..
But on our own, we can't..
And when we deceive ourselves into thinking we can, that's hypocrisy..
And then another great definition of hypocrisy that we pull straight out of scripture is.
from Luke chapter 18, verse 9..
And this is Jesus telling a parable, and it says, "He also told this parable to some who.
trusted in themselves that they were righteous and treated others with content.".
And I think that's a great definition of hypocrisy as well..
Trusting in ourselves that we're righteous and treating others with content..
And we've also defined it as putting the spotlight on us, wanting to glorify ourselves instead.
of glorifying God..
We talked about the distinction between Matthew 5, 16, where Jesus tells us to let our light.
shine before men so that they will see our good works and glorify our Father in heaven..
And then here in Matthew 6, 1, where he says that don't do things to be seen by others..
God doesn't want us to do things to glorify ourselves, to put the focus on us..
He wants us to do things to glorify him and put the focus on him so that people will see.

$^{81}$him, so that Jesus will be seen through our lives and that God will be glorified and people.
will be drawn to Christ..
And today I want us to think about a few different aspects of what we're actually doing when.
we live dishonestly, when we pretend that we're something that we're not, when we pretend.
to be more righteous than we are, when we pretend that we're the ones who are going.
to fulfill God's standard..
When we do that, we're living out hypocrisy..
But sometimes we think it may not be that big of a deal..
It's okay if people think I'm a little bit better than I am..
And we don't see how serious that can be in our life..
We don't see what's actually at the root of it and what's actually taking place..
So we're going to talk about that today..
And the first thing I want us to see is that when we live out hypocrisy, there is a matter.
of stolen glory involved..
Stolen glory..
Because we know if we understand the scripture that all glory belongs to God..
God is the creator..
God is the redeemer..
God is the one who loves us and gave himself for us..
That's what the scripture tells us..
The glory belongs to him..
And yet we as humans have a desire to steal some of that glory for ourselves..
We want the focus to be on us..
We want the attention to be on us..
We want to glorify ourselves..
That's our pride..
We talked about that the first week that we spoke about hypocrisy..
How pride is the root of our hypocrisy..
And that pride causes us to want to steal some of God's glory and redirect it at ourselves..
And that was the sin of the devil in the beginning..
Remember he wanted to be like the most high and that's why he rebelled against God..
That's what he tempted Adam and Eve with in the garden..
He said, "Hey, don't you want to be like God and know good and evil?".
And that's a core temptation for us in our life every day..
And it's certainly the root of hypocrisy is that temptation to say, "I want to be like.
God..
I want to be the center of attention..
I want to steal some of His glory..
I want to be full of pride and focus everything on myself.".
And that's what we're doing when we begin to live in a way where we're pretending to.

$^{121}$be something that we're not..
We're stealing the glory that belongs to God and trying to direct it on ourselves..
We see this in this verse here where Jesus says that the hypocrites love to stand and.
pray in the synagogues and at the street corners that they may be seen by others..
And all of us have that in us..
We want to do our good works..
We want to do things to be seen by others..
We want others to see what we're doing..
We want to do it in front of others..
We want to do it before others so that they will see us and regard us as good and regard.
us as special, regard us as righteous, and in some way glorify us..
That's what we want..
That's at the root of our hypocrisy..
We want to do those things before others..
Now some of us may say, "Well, I don't do that.".
You know, the examples that Jesus gives us here are pretty over the top, right?.
When He talks about someone giving to the needy, He says that they're sounding a trumpet.
in the street to draw attention to themselves so everybody will see it and to amplify what.
they're doing..
When He talks about people praying, He says they're standing and praying on the street.
corner loudly so that everybody will look..
We may say, "Okay, wait a second..
I don't do that..
You know, when I help someone out, I don't sound a trumpet..
I don't go out and pray in the middle of the street so that people will see what I'm doing.".
But Jesus is pointing out not just the specific situation..
He's pointing out the heart behind it..
And we definitely share that heart behind it..
I mean, let's be honest..
Probably none of us have ever gone out and sounded a trumpet before we did something.
to help someone else..
But I can't think of very many ways at all to more powerfully and visibly focus attention.
on ourselves than to use some of our modern technology, right?.
And we do it..
We go out to serve..
We go out to help at the soup kitchen or the shelter or this place or that place..
And we use our social media or some other form of technology to let everybody know what.
we're doing..
That's really the modern equivalent of sounding the trumpet, right?.
And again, I'm not saying that it's wrong for us to let other people see our good works..

$^{161}$Matthew 5:16 says we should let others see our good works, but they should see it in.
such a way that the glory is focused on God..
So I'm not saying it's wrong for you to take a picture of something you're doing and share.
it with other people..
That's fine..
But if your heart motivation in doing that is to focus people on yourselves, to glorify.
yourself, then you're sounding a trumpet..
You're doing it in order to glorify yourself..
And that is wrong..
That's what Jesus calls out in this passage..
And we see this replete throughout our culture and our society..
Think about it..
The celebrities, all of the politicians, all of the people in our culture and society,.
that they won't even show up to help..
They won't even do something unless the cameras are going to be there..
If it's not a photo op, then to them it never even happened..
It's not worth doing if no one's going to know about it..
And sometimes we live that way ourselves..
We think, "Well, yeah, I kind of want to be a part of this, but I'm so busy..
Oh, but someone else is going to be there and they're going to see me doing it..
Yeah, that makes it a lot more appealing to me.".
Or, "I'll have an opportunity to share it with other people and they'll know what I'm.
doing.".
Yeah, that makes it a bigger deal..
And that shouldn't make it a bigger deal..
We should be committed to doing the right and to living in obedience to God and to living.
out what God wants us to do, whether or not it gets noticed by anybody..
And we should be committed to the fact that if it does get noticed by someone, we should.
direct that attention towards God, not towards ourself..
But that's a difficult thing to do because we all want to do our good works before others..
Now some of us might say, "Well, wait a second..
I'm not even sure that I even try to do good works..
I don't even know what that means or what that is or that I'm even interested in that.".
Now we've talked about this in the past few weeks, that as believers, as followers of.
Jesus, as subjects of the kingdom, that God is interested in the condition of our heart..
We don't do good works to earn favor with God..
We don't do good works to please God..
We don't do good works to get to heaven..
None of that..
We do good works because God is changing our heart..

$^{201}$If we really have relationship with God through Jesus, God is at work changing our heart..
And out of that change, out of that transformation, good works come and beautiful good works come..
But the good works are not the starting point..
God's change in our heart is the starting point..
And the good works flow from that..
They come from that..
But some people may say, "Well, I'm not interested in good works..
I'm not interested in trying to be a good person..
I'm not out there trying to please anybody..
I'm just doing my own thing and I don't even care what people say about it.".
And that sounds good for a moment..
It might sound like, "Okay, yeah, I mean, that kind of gets rid of some of the hypocrisy.
and pride, right?.
If you don't care, then you're not going to be a hypocrite.".
But that's not even really true..
The reality is all of us have some standard of morality that we live by..
We all have some things that we think are good as opposed to other things that we think.
are bad..
Even if you say, "Okay, I'm an anarchist, so I don't think there should be any rules.
about what's good in society and I don't think there should be anything that should be a.
good work in society.".
Well, the reality is you have in that position defined a good work, right?.
Because you said that it's good to have no rules and that that's what's best..
And that if you do try to have rules in society and enforce those rules on other people, then.
that's bad..
So if you've, from your position, defined a good work and a bad work, even if you say.
that you're an anarchist and you don't care and you don't want any rules and you want.
chaos, you still believe that some things are good and some things are bad..
All of us do..
So we can't hide behind the fact that we may say, "Well, I'm not as interested in religion.
or I'm not really caught up in Jesus and so I don't really care, so I'm not a hypocrite.".
That's actually not true..
All of us have some system of morality, some things that we think are good and some things.
that we think are not good..
And the reality is whatever that system is for whoever you're talking about, all of us.
as humans want to do the things that we define as good in front of other people so they will.
see them and praise us..
Even the anarchist wants to go out and create anarchy for other people to see so that other.
people can praise what they're doing..
All of us do that as humans..

$^{241}$All of us live in hypocrisy in that way..
We want to steal God's glory and direct it to ourselves..
Another way we really see our hypocrisy in stealing God's glory is in this human nature.
of always seeing the problem in someone else and not seeing our own issues..
We always excuse ourselves and our sin and our injustice and we see the problems of others..
Jesus talks about this later in the Sermon on the Mount when he talks about the fact.
that we as humans see the little specks of dust in someone else's eye and overlook the.
giant log that's in our own eye..
And he makes that illustration for us..
But we do that all the time..
We say, "Oh, not me..
I would never do that..
I would never be a hypocrite..
I would never do good works in front of others..
I would never sound a trumpet..
That's just not me at all.".
And yet we turn around and see what everyone else is doing and point out what they're doing.
and get frustrated about what they're doing..
We see this so much in our culture all the time..
You know, right now we're just weeks away from a big election and this is so heavy in.
our politics..
Hypocrisy is so heavy in our politics..
It's so heavy right now that you could literally cut it with a knife..
It's so thick in the air..
Every other commercial on TV is a political commercial..
And if we use our definition from Luke 18, 9 where Jesus says that hypocrites are those.
who trust in themselves that they're righteous and treat others with contempt, man, that.
would be a great definition of every single political commercial we see..
Every commercial you see is a candidate demonstrating by their words and the pictures of themselves.
why they are righteous and why the person they're running against is contemptible..
And that's exactly what our politics is full of..
And our politics is full of this "not me" attitude..
The reality is throughout our politics in both parties there are things that are deeply.
dishonoring to God..
And yet both parties are full of people who are convinced that they are completely righteous.
and that all the hypocrites and all the sinners are in the other party..
And that's how everybody acts..
But the reality is it's not just in our politics..
That's how we all act as human anyway..
We are all convinced that the problem is with everybody else..

$^{281}$It's not me..
I'm doing right over here, God..
It's not me..
Because we have stolen God's glory and we have said, "You know what?.
I'm doing it right..
I'm getting it right..
I'm being obedient..
I'm righteous.".
It's all those other people who are deserving of guilt or shame or punishment or something.
like that..
But I'm the good person over here..
We've stolen God's glory..
The other thing that we steal in our hypocrisy is we want to steal God's throne..
We want to steal God's rule..
We want to steal God's control..
We want to take control of our image..
We want to take control of how other people perceive us..
We don't want to admit that we could have a weakness or a failing in any way..
We don't want anyone to think or imagine that we could ever have a problem..
So we try to hide it..
We try to pretend we're something that we're not..
We practice hypocrisy because we're trying to control our image..
We're trying to steal God's rulership, God's throne..
We try to control our success and not by defining success in the way that God has defined it,.
but defining it according to our own definition, which is how other people view us, how much.
we can accomplish and achieve here on the earth in this life..
And we try to control our success by controlling what other people see of us..
We try to make sure that we're successful because we give this tightly controlled image.
and perception of ourselves to other people..
And we are taking away God's control over our life that rightly belongs to God..
We try to control our righteousness and we do that by trying to redefine what is actually.
right and wrong..
Like I said earlier, everybody has a standard of morality..
The only distinction is where that standard comes from..
Now Christians are unique because our standard comes from beyond ourselves..
It doesn't come from within us..
It comes from God's Word..
It comes from God Himself..
That's our standard of morality..
And Christians also are unique because we understand that we can't live up to God's.

$^{321}$standard, that we all fall short and only through Christ can we achieve rightness with.
God..
But for most people, their standard of morality comes from within themselves..
They want to be in control of what righteousness is..
They want to be in control of what right and wrong is..
And they try to do that by defining it for themselves..
And they're stealing that throne from God..
God is the one who says what is right and wrong..
God is the one who defines righteousness, not us..
We try to define those standards for ourselves and as we do that, we are stealing away the.
authority that only belongs to God..
The other thing that we try to steal is our destiny..
The reality is that Jesus is the only way to God..
The reality is that the grace of Jesus Christ is the only way to satisfy God's standards..
Those things are so clear in the Bible..
There is no other way that we can have relationship with God except through Jesus..
There's no other way that we can achieve holiness and that we can be right except through Jesus.
Christ..
And we don't achieve that holiness through our good works..
We achieve it when we surrender to Christ and His righteousness washes over us and covers.
us..
And there's no other way apart from that for us to be righteous and to be right with God..
And yet we don't want that to be true..
We don't want Jesus to be the only way to God..
We want to make a way on our own..
We don't want to believe that we might have to rely on the sacrifice and forgiveness of.
Jesus for righteousness..
We want to somehow prove ourselves righteous apart from God..
And as we do that, we try to take control of our eternal destiny in a way that we cannot.
because it only belongs to God..
Jesus is the only Redeemer..
And Jesus is the only way for us to experience health and shalom and healing in this world..
And yet we try every other way..
And we try to fix things in every other way..
And it doesn't work, but we keep trying and keep wanting to pull destiny away from God.
and the future of what it means to be human..
We want to pull that away from God and define it ourselves..
And we keep failing..
You know, Jesus in Matthew 5 and 6 and 7 throughout the Sermon on the Mount, Jesus is giving us.
this new way to be human, this way of love, this way to love God and love those around.

$^{361}$us..
And it's really a new way because it's been around since the beginning..
It's the old way that God called us to in the beginning that we never lived out..
And so it seemed like a new way when Jesus was expressing it in Matthew 5..
It seemed like a new way because what people had been living out was their own way..
And that way was the way of vengeance and retribution..
Ever since Cain, the first siblings, Cain and Abel, Cain decided, "You know what?.
The way to fix my problems is to take revenge on my brother, is to get him out of the way,.
is to make him hurt..
And that's what's going to fix things.".
And ever since then, that has been the cycle of humanity..
People have tried to fix problems with revenge, with retribution, with war, with fighting,.
with violence, with oppression, and with hurt..
And when someone hurts someone else, people say, "Well, the way to fix that is to hurt.
them right back.".
That's been the story of humanity since the beginning of time..
That's what we as humans have chosen again and again and again..
And Jesus has said clearly in the Sermon on the Mount, "No, there's another way..
I want you to choose the way of love, the way of peace, the way of relationship with.
God, the way of caring for one another.".
Jesus is calling us to that..
But we don't want that destiny for ourselves..
We say, "No, God, I want to go back to my way, which is revenge, retribution, getting.
even..
It's taking my revenge on someone else.".
We try to steal destiny, and we don't have that authority..
And when we try to do that, it ends in disaster every time..
And sadly, even though we've seen it end in disaster a thousand times before, we still.
think we're going to try it one more time, and this time it'll actually work, and it'll.
actually come out right, and it'll actually make us feel better..
And then it doesn't..
But we still try it again, and we turn back to it again and again..
We want to pull destiny away from God and try to pull it down to ourselves and think.
that we're going to fix our lives, we're going to fix our future, we're going to fix the.
world, we're going to fix everything..
And we can't..
Now, what Jesus is pointing out to us in this passage is that hypocrisy is so damaging,.
it's so devastating, it's so wrong..
It's ultimately, as we discussed today, a direct attack on God's authority and on who.
God is..

$^{401}$It's trying to steal from God, and that is completely wrong..
And what Jesus is pointing us towards is the fact that we have got to avoid it as subjects.
of the kingdom, as followers of Jesus, because it is a powerful and constant temptation to.
us as human beings, and we've got to strive to avoid it..
So I want us to think about how we can avoid it as we close out today, how we can avoid.
these things..
We've talked before about the best way that we can avoid pride is by living in a place.
of poverty of spirit, living in a place of seeing and admitting our own spiritual bankruptcy.
before God, seeing and admitting our need for God and saying, "God, I understand that.
I am broken, that I am full of sin in my heart, that I'm full of selfishness, that if I just.
had my own way, I would be hurting other people, I would be trying to glorify myself, I would.
be trying to do all this stuff.".
And when we're able to be honest and transparent before God and see our own brokenness, our.
own need, then that gets rid of that pride and it gets rid of that hypocrisy..
We've got to do what Jesus suggests in this passage, which is to hide our good works..
That's actually a really good thing to do, is as much as possible, and this isn't like.
a law, you don't need to try to make it into some hard and fast rule that you have to follow.
in every instance, but as much as possible, when you go out to obey God, just do it in.
a way that other people aren't going to see or aren't going to notice, so that that temptation.
to glorify yourself doesn't take hold..
When you go out to obey God, you should go out to obey God with the motivation of pleasing.
God and God alone..
You should go out to obey God because it's the right thing to do..
You should obey God in your life, whether it's in the privacy of your home or whether.
it's serving in your community in some way, you should do that because it's the right.
thing to do because you want to love God, you want to love others, you want to glorify.
God..
Don't do it because other people are going to see it and notice it..
As much as you can, do it in secret..
Do it so that other people don't know..
If you're going to go serve in your community, maybe you don't need to take all those pictures.
and post them online..
Again, that doesn't mean it's wrong to do that, but just think about that..
Just ask yourself, "Hey, can I do this in a way that the focus is not on me?".
If someone does notice what you're doing, if someone does notice you serving and obeying.
God, then direct that attention towards God..
Don't let people just say, "Oh, you're a great person..
You do such a great job.".
That feels good and there's a part of us that is like, "Yeah, it doesn't really hurt for.
someone to glorify me a little bit, for me to get a little praise, but hey, it's going.

$^{441}$to be damaging to you and to them..
It's damaging for others to glorify us and to think we're something we're not, and it's.
really damaging for us to think we're something we're not.".
So when you get a little bit of that praise, redirect it to God and be honest and say,.
"Hey, I'm just as broken as anybody else..
I'm only here doing this because of God's work in my life..
It's not me..
I'm not a great person..
Don't glorify me..
It's because of God.".
Redirect that praise to Him..
Don't steal God's glory..
Don't steal God's throne..
Don't try to steal God's control or God's rule in your life..
Don't try to steal the destiny that God has for you..
Don't try to redirect it into your own will..
Remember Jesus, here in just a few verses, He gives us the Lord's Prayer..
What does He pray at the end of the Lord's Prayer?.
"Father, your kingdom come, your will be done on earth as it is in heaven.".
Don't try to make that into, "My kingdom come and my will be done.".
That's dangerous and that has devastating consequences..
Don't steal the control of destiny from God..
That belongs to God and God alone..
The throne belongs to God and God alone..
Glory belongs to God and God alone..
If we will do these things, if we will seek to live in poverty of spirit, if we'll seek.
to take our hands off of God's glory and God's throne and the destiny that God has for us.
and let God be in control of those things, if we'll be honest and transparent about our.
own weaknesses and our own struggles and our own failings with the people around us, we.
will find ourselves living in a much more authentic way..
We will find ourselves with very little hypocrisy in our life..
And that's what Jesus is calling us to in this passage..
Let's pray..
God, I pray that you would give us poverty of spirit..
You would help us to be aware of our own need for you, to be aware of our spiritual bankruptcy,.
the fact that we have nothing before you, the fact that we can't do anything on our.
own apart from you..
God, I pray that you would help us to be honest about our own sins and our own struggles..
Help us to be honest with the people around us..
Help us not to cover things up or to hide things, but help us to be able to be honest.

$^{481}$and say, "Hey, I struggle as well..
I'm imperfect as well..
I've messed up too..
But I need to repent too..
I need God's help in my life.".
Help us to be honest and transparent with our struggles, God..
Help us not to operate from a place of pride, but from a place of authenticity, God..
God, help us not to steal any of your glory..
Help us not to redirect any of the glory that should go to you towards us, God..
Help us to always point to you and always live in a place of humility..
God, help us not to steal your throne, your control, God..
Help us not to grab that and try to act like we are going to be in charge when we cannot.
be in charge..
What a disaster that would be if we were in charge, God..
God, help us not to steal the destiny that you have for us and our world and try to think.
that we should direct that..
Help us not to think that our will should be done..
Help us to always desire for your will to be done, God..
Help us not to put ourselves in a place that we should not be in..
Help us not to pretend to be something we're not..
Help us not to think we are righteous and treat others with contempt..
God, help us to avoid hypocrisy and to live honestly and openly as your disciples in a.
way that glorifies you, God, and that causes people to see you and to desire you and to.
want to know more of who you are because of what they see in our lives, God..
May you be glorified by the testimony of our lives, and may people see more and more of.
you in us. In Jesus name, Amen..
\newpage



\section{約拿書}
\label{sec:y_7ZwDcfiBI}
\textbf{經緯線 [釋經視頻] EP1 約拿書 - 介紹,文學體裁}
\newline
\newline
連結: \href{https://youtube.com/watch?v=y-7ZwDcfiBI}{\texttt{https://youtube.com/watch?v=y-7ZwDcfiBI}} ~~~~ 語音日期: 2020-10-30
\newline
\newline
\hyperref[sec:c2n9YZYOe80]{\small{< < < PREV SERMON < < <}}
~
\hyperref[sec:index_chronic]{\small{[返順時目]}}
~
\hyperref[sec:index_scriptual]{\small{[返順卷目]}}
~
\hyperref[sec:WmRxFI5hs7M]{\small{> > > NEXT SERMON > > >}}
\newline
\newline
$^{1}$各位好,各位好,現在是經緯線的時候了.
究竟經緯線是什麼呢?.
不過在介紹之前,不如讓大家認識一下我們是誰.
我叫Tony,王家漢牧師.
我在卡加里,既沒有養人教會.
也是在三一神學院讀我的教育學的博士學位.
我把時間交給大家,其他的同工一起介紹自己吧.
我是Mark Sir,陳淑姬,我在Chicago的.
成匯網台就是我們成匯社區教會,芝加哥教會的一個事工.
我們歡迎每一位guest,我們錄了很多program,能夠去祝福更多人.
大家好,我是卡城的環浸信會的黃師母.
大家好,我是Benny,首先我就是一個愛主的好弟兄.
另外一樣,我現在在Trinity Western讀神學.
不是讀神學,是在一個program中入讀.
順道一提,希望大家教會捐多點錢給神學院.
那些神學課程就不用那麼貴.
完全明白你的感受.
說到這裡,多謝幾位在不同的位置,不同的崗位.
一起來做這個節目,經緯線究竟是一個什麼節目呢?.
Mark Sir,你叫我來做什麼呢?不如介紹一下吧.
經緯線其實就是一個網上識經的視頻.
我們希望能夠會集不同的讀經的人.
嚴肅的讀經,嚴謹的想去識經的人.
我們會嘗試每一集大概二十多二十五分鐘.
我們會講一段的聖經.
我們希望用一個不是靈修式的看聖經.
希望用不同的識經工具去看聖經文學.
大家看這個program的時候.
聽到我們如何討論那段聖經.
可以多一些的insight.
亦可以幫到你如何apply一些識經的工具去看聖經.
好,麻煩Mark Sir為我們在這個program定了一個調子.
其實就是我們嚴謹去讀聖經,嚴謹去看聖經.
我知道蔡卓其實我自己和其他兩位都是.
三位都是有心讀聖經的人.
我都想很快的,你可不可以講一兩句.
為什麼你會參與經緯線的工作呢?.
先讓我講吧,因為我知道聖經是一個文學.
但如果你在教會裡面講聖經是文學.
你就會被人罵,在這裡你就不會被人罵.

$^{41}$很希望給大家一個不同的角度去理解聖經的經文.
黃師母你自己有什麼原因推動你參加這個節目呢?.
首先多謝Mark Sir邀請,如果不是他邀請就沒有這個機會.
另一方面我自己很熱愛查經,我希望在查經裡和大家分享.
藉著分享可以刺激到大家對查經產生一個熱誠.
我自己心底裡有一個希望.
是幫助弟兄姊妹能夠明白聖經的內容.
不單止明白當時當聖經寫好了之後.
講給當時第一代的讀者聽.
對當時第一代讀者的意思.
同一時間對我們今天,我們過了很多年之後.
再看經文的時候對我們今天的意義.
用英文說我們既要知道What it meant.
同一時間也希望知道What it means today.
這個也是我自己的盼望,透過經衛線能夠達到的.
事不宜遲,我們進入正題,經衛線.
我們揀選了第一卷書就是約拿書.
為什麼會看約拿書,為什麼大家這麼有興趣研究約拿書呢?.
不如大家分享一下.
我想約拿書是最fun to read,最過癮.
而且是比較細,大家都很清楚.
所以我們第一卷就用約拿書.
我覺得在華人讀經的傳統中.
經常都在掙扎約拿書的歷史問題.
那條魚吃了三天,是不是真的會發生?.
上帝是不是可以做到這件事?等等.
我覺得如果這樣看,我們經常看聖經.
歷史裡面會不會發生這樣的事?.
我想我們錯過了聖經文學的重要意思.
就不在那裡了.
你講到這裡的時候,也令我想起一件事.
在傳統華人來說,我記得我小時候聽到的時候.
很多時候一講到約拿書,我們很快就想到耶穌自己引用約拿書.
就說以色列人要見神蹟,唯一的神蹟就是.
當若拿被魚吞了之後,三天在漁福的神蹟.
當然在科幻書中講到明,那個神蹟是指著耶穌基督的死和復活.
我這樣說不是否定耶穌基督所說的話.
只不過我們要明白一件事.
耶穌當時看到舊約的聖經.
而他借用舊約聖經這個概念引用在自己身上.

$^{81}$耶穌是這樣用這段經文.
但是我相信同一時間對於當時舊約的讀者來說.
回到第一生,當這本宣之書講的時候.
我相信當時的讀者還沒有看到耶穌這一面.
所以我自己希望先從去理解耶穌之前的舊約.
當這本宣之書向當時的以色列人講的時候.
當時的人聽到什麼,他有什麼反應,有什麼感受.
以致我們讀這段經文的時候能夠清晰地說明.
最想說的是我們想看看這本書對當時第一生的讀者有什麼意思.
我除了被你拉下水之外.
因為約拿書是很短的.
你慢慢讀的時候會發覺它裡面用了很多.
在其他的聖經書卷裡面都會運用一些文學手法去寫.
我覺得這是一個非常好的入門.
用一個比較批判的角度去看經文.
我自己為什麼會選約拿書.
因為之前港獨的時候,四張書我都查過.
我發覺我很喜歡這張書,最主要是開放式結尾.
讓大家去想,究竟結局會是怎樣.
所以我就選了這本書.
我們有四個不同的角度去看這本書,相信會很豐富.
開始之前,我們一起留意一下約拿這個人.
在這段經文的第一章第一字已經說到約拿是亞米泰的兒子.
當我看的時候,我自然就問.
究竟約拿身處在一個怎樣的時代呢?.
這是值得去想的,Mark Sir你有沒有什麼見解?.
我想我們要去看這本書的時候,就要想一下歷史的背景.
為什麼要特意提亞米泰的兒子約拿.
可以是其他名字都可以,可以是其他先知都可以.
為什麼一定要亞米泰的兒子先知約拿.
這本書是有關他的.
原來約拿這個人不單止是約拿書裡面提過的.
在聶王記下也有提過.
那時候就是耶羅邦二世的時候.
約拿就是做耶羅邦二世的時候的先知.
就是聶王記下十四章.
那裡說到耶羅邦二世就像一個所羅門的翻版.
所羅門的二世那樣.
因為他說到他修復了以色列的邊界.
由哪個地方到哪個地方.

$^{121}$以前屬猶大的地方,他現在拿回來了.
以色列的版圖大到像以前所羅門那麼厲害.
裡面其中一節,二十五節說什麼呢?.
正如耶和華說,上帝直誅他的僕人.
加特希弗人亞米泰的兒子先知約拿所說.
所以他是宮廷的先知.
他說了一些話,說耶羅邦二世多厲害.
真的應驗了,就是耶羅邦這麼厲害.
剛剛那段時間是790至760的BC.
那時候亞述國,就是以色列人很憎恨的敵國.
亞述帝國就是色微.
所以耶羅邦二世的時候,以色列很光榮很強.
然後他們的敵國就一直插水,搞不定.
所以是一種民族主義情緒,國家主義情緒很高漲的時候.
那時候他做先知.
所以我想約拿書就一定要在這個context裡面去看.
如果這樣說,其實我們看到約拿書裡面本身.
約拿自己已經是一個,我們可以說是一個愛國先知.
可以說在他的眼裡面,當日欺負他的亞述大帝國倒塌了.
那我這個國家終於站起來了,成為我們人民站起來了,多厲害.
那說到這件事的時候,我就想一下.
舊約聖經這樣說,其實還有沒有其他的文獻.
是支持類似說這些的話.
又或者是讓我們看到,好像是約拿在當時的舊約時代.
或者是當時以色列的時代,是受人尊敬.
因為畢竟他之前說了一個,以色列國的一個好信息的先知.
不知道大家怎麼看呢?.
這個我懂,而且一定要在這裡說.
因為在教會說就會被人說是異端.
因為有時候查經看聖經,其實不妨可以去參考聖經以外的文獻.
其中一種文獻就叫做《賜經》.
《賜經》裡面有一本書叫做《多比傳》.
在《多比傳》最後的章節,就是章節14.
他就說到Tobit快要死了.
他快要死之前,他跟他的兒子Tobias說.
我深信約拿先知,他對於亞述國和尼尼微城的預言一定會應驗.
之後他就死了.
到了章節14的最後,就輪到Tobias快要死了.
他臨死的時候就很開心地說.
現在終於死得瞑目了.

$^{161}$因為我終於看到約拿先知所說的預言應驗了.
所以從這本書裡面可以看到約拿先知所說的在以色列人心裡面有一定的份量.
你說得對.
其實有一件事順帶一提.
雖然對於我們基督新教的基督徒來說.
我們可以肯定我們的舊約聖經和我們所謂一般舊約的《賜經》.
在我們心目中的份量當然很不同.
不過如果從《賜經》的角度來說.
他的確提供了一個歷史的場景.
或者一個角度讓我們看到當時的猶太人和以色列人.
如何看待他們自己的時代.
從一個歷史的角度來說,這是一個很好的參考.
順便說一件事是很真實的.
作為一個先知,他的幾代都經歷著亞述被欺凌.
現在亞述倒下,自己的國家強大了.
當中的民族情緒應該是很高漲的.
我覺得這是一個歷史的台.
設了一個stage,英文叫做台戲的第一幕.
變成了幫我們去理解,稍後再看約拿書.
說到這裡,尤其是歷史背景.
我突然間很有興趣,很想知道約拿的名字是甚麼原因.
因為像中國人一樣,我們每個人改名的時候.
我相信我們的父母都有意思的.
不知道還有哪一位對約拿的名字有所研究.
可以分享你的成果給我們聽.
不如我講一下吧.
《聖經》裡面的人物的名字其實都很有意思.
那約拿當然也不例外.
很多歷史的人都會用一些有良好品格的鳥獸來改名.
而約拿這個名字其實取自於甲子.
大家都知道甲子代表著和平,善良和溫馴.
如果約拿真是人如其名的話.
他應該是很善良,而且是很追求和平的.
但不要用中文,因為中文放甲子就是放飛機的意思.
我相信當時他父母幫他改名就是這樣.
似乎看上去他是放飛機的人.
也是,也很適合,真的,的確是.
沒有改錯名字.
不過有一樣東西我想問你,黃師母.
亞米泰的兒子亞米泰又有什麼意思呢?.

$^{201}$其實當時亞米泰的兒子.
他父親的名字和他自己連在一起.
其實這樣連在一起相當有意思.
首先看看他父親亞米泰.
其實亞米泰人是一個稀白來文.
他的意思是那個字根和阿們同屬於一個字根.
所以亞米泰的兒子代表著一個真理之子的意思.
我們講完歷史,又講完名字.
我想講一下文體.
我們對約拿所處身的歷史和人名的認識.
不知道大家對這本書的文體有什麼理解呢?.
這就是我們希望大英姐妹看完聖經.
人類看完一本書的時候.
就要用文體入手,genre.
因為好像你看報紙一樣.
不同的genre,不同的題材.
看廣告有不同的預期.
看八卦新聞有不同的預期.
看頭條有不同的預期.
約拿書是一本諷刺文學,sceptre.
什麼叫諷刺文學呢?.
就是希少羅馬的政治卡通.
將這些東西誇大來說.
或者去marking取笑一些東西等等.
如果你這樣入手,說這是諷刺文學的問題.
你看約拿書可能有不同的味道.
就留意一下有沒有希少羅馬誇大的細節.
你剛才說到那是一個諷刺文學,sceptre.
其實你覺得他在諷刺哪一個呢?.
我們要一路看下去,好像一個圈.
會不會諷刺一件事.
然後你看看其他的detail.
會不會又諷刺其他東西呢?.
我們都會談到約拿書特別的地方.
就是和其他小先知書不同.
因為通常小先知書都是先知說的東西.
很少說先知自己.
唯獨是約拿書是說先知自己.
不說他說什麼.
我們先打一個問號,繼續看下去.

$^{241}$會不會就是諷刺先知自己呢?.
還有如果你將約拿跟其他11個小先知去比較的時候.
英文有一個詞叫asore.
他是一個特別的asore,特別不同的.
因為每一個先知都是他們要去救那些人.
他是救不了的.
唯有約拿他才可以救到.
這個就是很諷刺的東西.
What happens when a prophet becomes successful.
其實這個問得挺好的.
因為如果我們再留意一下.
在約拿書整個故事上.
我們要問一個問題.
自己覺得是不是真的成功呢?.
這是一個好問題.
我暫時先賣個關子.
不過剛才Mark Sir說的諷刺文學.
我想我會再加多一重.
就是我們如果就這樣看經文的時候.
我們發現這一卷小先知書.
剛才Benny說得很好.
是12卷小先知書裡面其中一卷很不同的.
Mark Sir和Benny都有說過.
這卷小先知書是一個敘事文學.
不是真的說的,是一些對話.
或者是一個詩歌體寫出來.
但這是一個敘事文學.
而且是有一個很明顯的起承轉合.
如果你看故事就知道.
有一個beginning起.
有一個成development.
有一個轉差不多去到一個climax.
最後合去到一個conclusion.
是一個很清晰的故事.
所以我覺得如果將它拍成電影或歌劇.
是一個很吸引的故事.
所以我想我除了從一個.
所謂諷刺的文學的手法.
我也鼓勵大家從一個故事文體當中.
一起去理解這段經文.

$^{281}$就這樣我們簡單地.
未入經文就說這麼多東西.
我們已經知道這裡的歷史.
是一個愛國先知在他的國家站起來.
在他的敵人倒下的時候.
所發生的一個事件.
他也說到這本書裡面的文體.
是一個諷刺文學.
也是同一時間.
是一個敘事文的版.
很清晰的一個例子.
知道這麼多東西了.
各位觀眾.
我們應該是時候進入經文.
不過這樣.
給大家休息一下.
我們下一集.
我們就會進入經文約拿書第一章.
嘗試去看這個敘事文學.
這個諷刺文學的開端.
所以請大家有機會的.
再看約拿書.
接著我們一起進入經緯線.
\newpage



\section{約拿書 1:1-6}
\label{sec:WmRxFI5hs7M}
\textbf{經緯線 [釋經視頻] EP2 約拿書 1\_1-6}
\newline
\newline
連結: \href{https://youtube.com/watch?v=WmRxFI5hs7M}{\texttt{https://youtube.com/watch?v=WmRxFI5hs7M}} ~~~~ 語音日期: 2020-10-30
\newline
\newline
\hyperref[sec:y_7ZwDcfiBI]{\small{< < < PREV SERMON < < <}}
~
\hyperref[sec:index_chronic]{\small{[返順時目]}}
~
\hyperref[sec:index_scriptual]{\small{[返順卷目]}}
~
\hyperref[sec:lgq1HtuVP6w]{\small{> > > NEXT SERMON > > >}}
\newline
\newline
約拿書 1:1-6
\newline
\begin{longtable}{cl}
\hline
\hline
章節 & 經文 (和合本修訂版)\\
\hline
1:1 & \begin{tabularx}{0.7\textwidth}{X} 耶和華的話臨到亞米太的兒子約拿，說： \end{tabularx} \\ \\ \relax
1:2 & \begin{tabularx}{0.7\textwidth}{X} 「起來，到尼尼微大城去，向其中的居民宣告，因為他們的惡已達到我面前。」 \end{tabularx} \\ \\ \relax
1:3 & \begin{tabularx}{0.7\textwidth}{X} 約拿卻起身，逃往他施去躲避耶和華。他下到約帕，遇見一條船要往他施去。約拿付了船費，就上船，與船上的人同往他施，為要躲避耶和華。 \end{tabularx} \\ \\ \relax
1:4 & \begin{tabularx}{0.7\textwidth}{X} 耶和華在海上颳起大風，海就狂風大作，船幾乎破裂。 \end{tabularx} \\ \\ \relax
1:5 & \begin{tabularx}{0.7\textwidth}{X} 水手都懼怕，各人哀求自己的神明。他們把船上的貨物拋進海裡，為要減輕載重。約拿卻下到艙底，躺臥沉睡。 \end{tabularx} \\ \\ \relax
1:6 & \begin{tabularx}{0.7\textwidth}{X} 船長到他那裡，對他說：「你怎麼還在沉睡呢？起來，求告你的神明，或者神明顧念我們，使我們不致滅亡。」 \end{tabularx} \\ \\
[1ex]
\hline
\hline
\end{longtable}
$^{1}$各位好,又回到經緯線了.
我們上一集已經介紹過.
關於約拿書的歷史.
約拿的名字的意思.
以及也說過一些在約拿書的文體.
需要留意的.
就是一個諷刺文學.
再加上一個敘事文學的結合來看這本書.
事不宜遲.
我們再進入看這本書.
我們一起看的時候.
不如我們先留意第一至第六節.
大家有留意第一至第六節嗎?.
很有趣.
好像一開始.
上帝已經叫約拿做些事情.
大家知不知道他叫約拿做些什麼?.
就是去尼尼美.
這個就是去敵國.
當時我們剛才說過的那個 context.
當時的歷史背景.
就是開始色微.
以色列過去開始漸大.
居然叫他去亞述的尼尼美.
叫他做什麼?.
又是惡人.
那幫惡人要去跟他們傳道.
叫他傳道.
他不知達到我面前了.
達到你面前.
聽完之後其實只是一個很普通的宣講.
但好像約拿做了一樣挺有趣的行動.
大家有沒有留意到?.
走佬.
走佬?.
是雀草.
我覺得很奇怪.
在經文我們特別留意的時候.
我們發現到.
人家叫他尼尼美.

$^{41}$尼尼美接著去了約帕.
接著去了他斯.
其實這三個地方是哪裡來的?.
亞述帝國是在以色列的東邊.
約帕港是以色列其中一個舊約裡面有的海港.
他就在那裡搭船.
然後去他斯.
他斯就是近現在的西班牙.
總之就是最天涯海角最西面的地方.
這個地理裡面已經給我們一些irony.
耶和華叫他去東.
他就走去西.
還是進西的地方.
是,很有趣.
還有大家看到他說.
遇見一隻船.
但是我看其他譯本好像不同.
有沒有人可以幫我確認這個想法?.
這個就是我們為什麼要查經要看不同的譯本.
希伯來文的意思就是說他找一隻船.
He found a ship.
所以和合本也有一點預見.
味道出不來.
是.
不過如果我們不認識希伯來文又怎麼辦呢?.
各位有沒有辦法救救我?.
救救小Tony?.
其實也可以參考一下新譯本.
新譯本的解釋就是找到一隻船.
新譯本就準確很多了.
他有心intentional.
刻意去找一隻船.
這個很significant的.
這個很重要的.
這個就是誤殺和謀殺的分別.
是.
你偶遇一件事和你去找一件事是完全不同的.
就好像你去小巴站.
你遇見一輛車去他司.
和你找一輛車去他司是完全不同的.

$^{81}$這麼多輛巴士我總是找一輛最遠的走過去.
這個就是讀文學的一個特別之處.
還有其他的clue.
譬如他去買船票.
他買船票這件事根本就是一個預謀的事.
他是想這樣做的.
他invite我們讀者去想.
由約拿的家鄉的地方.
走到南面去約泊.
他途中當然要錢.
拿著八百通.
原來還有船價.
他根本是非常intentional.
非常刻意.
齊了錢就走了.
所以不是intentional.
差在他還沒有book airbnb.
我book了也不book的.
我有一個很有趣的問題.
我在看的時候.
他真的很買船票.
還有我發現了一件事.
如果我們看一至六節裡面.
你留意.
有一個日本叫呂鎮忠的日本.
我也挺喜歡用的.
他講到約拿逃往他司被永恆主的面.
然後他說下到約泊.
然後就下到船.
然後下到倉底.
似乎聖經第一生寫的時候.
不是為了看.
是為了聽.
如果我聽的時候.
下 下 下.
出現了三次.
一直走著.
好像躲起來一樣.
挺有趣的.
這是一個修辭的問題.

$^{121}$有時候聽牧師講到.
他經常會數著.
這個名字出現了多少次.
或者這個字出現了多少次.
他就不是講到差.
拿來浪費時間的.
其實是本身有關聯的.
剛才Tony講到.
以前的聖經是在一個公眾場合讀出來的.
當你不斷去講這些字的時候.
那些聽的人就會明白到.
原來這件事是很重要的.
剛才講到他下到約泊.
下了船.
還有下到倉底.
你就會有一個很自然的反應.
如果你要找上帝的話.
你不是看地上的.
你應該是看上天的.
所以這件事又是另一個很有味道的.
文學的手法.
嗯.
如果就這樣看.
起碼現在手上這些資料.
起碼我們都肯定.
亦都經文自己有講.
是若拿要逃避上帝.
我就在想.
為什麼他選擇要避開上帝這個呼召.
很清晰的.
叫你去利利美.
去宣講道理.
不講.
逃跑.
大家覺得為什麼呢?.
這就是我們第一集講的那個context.
一定會inform到我們.
告訴我們我們需要拿什麼message.
因為我們上一集講過歷史的背景.
你就是需要了解這個故事.

$^{161}$在那個處境裡面.
現在是大以色列國.
大以色列民族主義抬頭的時候.
要走去利利美.
去我們敵國.
他們是在式微的.
你就想想.
就是這樣才make sense.
他逃避.
嗯.
很有趣.
好像他在這一刻.
很想逃避這個使命.
很有趣.
但不知道其他呢.
除了這件事之外.
還看到一些關於約拿的事情嗎?.
有一件事其實我想補充一下.
嗯.
就是剛才Mark Sir講到.
這個利利美城是東邊的.
就是inland的.
但他就要走一條水路.
猶太文學裡面.
那個水啊.
大海啊.
都是其實在講一些很chaotic的東西.
是一些不好的東西.
所以這個又是一個諷刺.
沒錯.
利利美城是一個很不好的地方.
但約拿逃跑.
他要去一個更加不好的地方.
是啊.
對.
因為海裡面.
Yarm.
經常都講到邪惡的源頭.
那些混沌的源頭.
嗯.

$^{201}$你不去那些惡人那裡.
你就走去那些混沌的源頭.
更加混沌.
他心裡都挺硬的.
如果這樣講.
寧願去一個更混亂的地方.
都不要去利利美.
其實都可以看到約拿.
有一方面的性格都看到.
他都挺敢作敢為的.
你知道他被呼召去做一個先知.
是吧.
是代言人.
理應是有他自己的使命.
但是他神吩咐他要去利利美城.
有他的差事的時候.
他卻是不聽上帝的話.
拒絕上帝還要逃跑.
其實都看到他很有膽識.
敢作敢為.
那倒是.
我看的時候會覺得.
如果我在.
我自己有一刻想.
他是先知.
又是知道上帝的心意.
他為什麼這麼敢做呢?.
敢作敢為.
他不是不知道神的性格的.
這個我想我暫時沒有一個答案.
不過我覺得是一個挺好的方向去想.
究竟最終約拿想了些什麼呢?.
我們留個關子在這裡.
我們再看第四至第六節.
在第四至第六節裡面.
我們又發現.
當約拿上了船.
突然之間起大風.
狂風大浪.
甚至船都幾乎破壞.

$^{241}$兩位剛才說得對.
去了一個很 chaotic 的地方.
現在真的吹到媽媽都不認得了.
當中有三類人.
各自有不同的反應.
有一批水手.
有船長.
還有約拿自己.
不知道大家在看這幾節的時候.
你留意到些什麼呢?.
如果我們看一個諷刺的文學.
我們就留意這些 character.
要問一個問題.
為什麼這些 character 和想諷刺些什麼呢?.
我覺得很真實的是.
做得水手.
你跑馬三分險.
你應該都很知水性的.
但是害怕.
第一,那個風浪真的是很大.
可能因為你知水,你知死.
這個我大家看到的.
如果你不熟水就不理會.
腳就睡著了.
如果你知道熟水性.
熟行船的.
就知道這是大風浪.
這是第一件事.
接著我又看到.
那些水手不斷求自己的神.
都很有趣.
除了他們之外.
不知道大家看不看到其他.
比如約拿.
或者船長的反應是怎麼樣呢?.
約拿已經下去船艙睡覺了.
當時雖然環境打風打得很大聲.
那些人都要求救.
求生.
但是約拿呢.

$^{281}$看到約拿竟然在艙底睡著了.
而且還睡得很熟.
我可能在想.
怎麼可以睡得這麼熟呢?.
外面應該一群這麼熟水的水手.
嚇得真的.
嚇得臉都青了的時候.
一定在很大聲求告上帝.
喊救命的時候.
你應該聽到的.
又是船上浪下浪下.
招昏.
很奇怪的.
按道理應該是不合情理的.
是的.
但是它是這樣發生的.
所以這就是重點.
他不是要你去重新建構.
到底當時有沒有可能發生這件事.
而是文學的手法就是有一個諷刺的地方.
這些外邦人.
他們各自急急的去求自己的神.
有懶的時候求神是很自然的.
誰知道這個很正統.
我想是超正統的先知.
非常有道德的.
非常神學正確的人.
他們在做什麼呢?.
在那裡睡著了.
那個諷刺的高峰.
我覺得就是那個船長.
一個不信耶和華的船長.
竟然下去問他.
全世界都在拜神.
你不在那裡睡覺.
為什麼不拜神呢?.
喂喂喂,對.
就好像有一個牧師在睡覺.
飛機就在失事.
旁邊的一個無神論主義者說.

$^{321}$喂,祈禱啊,救命啊.
快要死了.
剛才我想到的就是.
或者是文學手法的技巧.
就是用一些誇張的人物的行為表現.
使讀者更加看得到對比性的強烈.
達到那個效果.
我們作為一個讀者.
剛才也看到了.
無論是船長的反應.
一群水手的反應.
還有約拿自己本身的反應.
其實我們不要去追究.
實際情況怎麼可能不合理.
但是如果我們作者用誇張的手法.
用技巧的手法.
想帶出那個果效的時候.
這個就是那個原意.
我相信作者真的做得到.
喂喂喂,你留意.
有時候那個故事.
希伯來的故事.
不回答或者不說話.
都是一種回應.
你見到他到第六節.
甚至你再下幾節.
他都還不回答.
他的沉默就是說.
作為讀者.
一個這樣的人怎麼會不回答呢.
尤其是他是先知的角色.
他是一個孫郎的角色.
其實這個技巧手法非常好.
令我們作為讀者很有興趣想追下去.
其實是一個不錯的寫作手法.
所以他的結局都是開放的.
我欣賞他就是開放的.
我們讀者可以很開放自己去想.
所以如果這樣說.
就這樣看這裡的時候.

$^{361}$我已經發現了一個很有趣的圖畫.
約拿在這一刻.
他在船艙底.
船長叫他的時候.
有些人在祈求.
他們滿天神在求.
而在猶太人的傳統.
他們都相信.
讀這本書的人都是猶太人.
都知道這班人拜的是假神.
但為了拜真神.
他們竟然不求.
很有趣.
他們會問:搞什麼啊先知.
為什麼你不出聲.
我覺得這個問題開始浮出來了.
為什麼你的心這麼硬.
為什麼死都不求呢.
他被迫才會答.
第七節:聖上的人彼此說.
第八節:中人對他說.
搞什麼鬼啊你在哪裡.
第九節:他才說.
他才說出來.
或者遲一點.
那個情況要去到多壞.
你才願意去見到.
這個都是很有趣的經文.
我們已經看了六節經文.
已經發現很多問題出現了.
容許我做一個簡單的總結.
在這六節經文裡.
我們已經看到若拿知道上帝邀請他.
或者甚至嚴厲一點說.
呼召他去利尼尼大城.
去宣講上帝的道.
他不要做.
他當他去西部.
買船票.
特意要坐一艘船離開利尼尼.

$^{401}$有那麼遠就遠.
接著風浪大至.
一直到船艙底的時候.
到眾人在求救他們的神.
他都在睡覺.
但是竟然一個.
誤信耶和華神的船長走下來.
拍拍他.
喂,你上來一起.
祈禱求.
是很有趣的一個圖畫.
在這六節經文裡.
已經很豐富地表達出一個問題.
我們直接開始問.
為什麼.
或者開始諷刺.
為什麼這位先知.
李英不是同神最親.
為什麼他會去到這樣的光景.
第二就是為什麼他要.
那麼心硬逃避.
什麼原因呢?.
各位,這是值得我們去看的.
接著我們會.
在第七節.
去到第十六節的經文當中.
很有趣的就是.
我發覺去了船上.
差不多是冤有頭債有主.
不只是中國人的思想.
因為船員就一起.
究竟是誰搞出來的.
所以他們就用.
制籤,就是搖籤.
就搖了出來,求籤.
終於就.
知道是若拿.
那麼.
在這裡的時候.
我們就要問.

$^{441}$知道了之後.
大家留意一下.
這個若拿.
和船員的對話的時候.
不知道大家看到什麼呢?.
制籤那裡我想說.
你可以想像到.
他們拿些牙籤出來.
有多少人拿多少字位.
然後撕一條.
就短一點的,然後放在那裡.
喂,大家每人抽一次.
誰知道他抽了一支長的.
原來是你.
你不能走.
為什麼這麼像周星馳.
就是周星馳.
就是這麼搞笑.
好了.
說到這裡的時候.
起碼現在知道.
是若拿,是你了.
那麼你應該怎樣解釋呢.
這個很問題.
但是時間所限.
不如這樣吧.
我們就留到下一次.
下一次,精彩了.
下一次的我們經衛線.
看約拿書的時候.
就問抽了一支短籤的若拿.
怎樣面對這群信仰.
怎麼搞好呢?.
我們一起繼續研究這個約拿書.
順勢在這個經衛線當中.
探索這本書的真義.
\newpage



\section{6:16-18}
\label{sec:lgq1HtuVP6w}
\textbf{Kingdom Values XXVIII: Matthew Chapter 6\_16-18 - Pastor Sam Johnson}
\newline
\newline
連結: \href{https://youtube.com/watch?v=lgq1HtuVP6w}{\texttt{https://youtube.com/watch?v=lgq1HtuVP6w}} ~~~~ 語音日期: 2020-11-01
\newline
\newline
\hyperref[sec:WmRxFI5hs7M]{\small{< < < PREV SERMON < < <}}
~
\hyperref[sec:index_chronic]{\small{[返順時目]}}
~
\hyperref[sec:index_scriptual]{\small{[返順卷目]}}
~
\hyperref[sec:LPCHL652i_w]{\small{> > > NEXT SERMON > > >}}
\newline
\newline
$^{1}$Family, I hope that you are doing well this morning and we are going to continue on with.
Matthew chapter 6, continue on with this study of the kingdom values that Jesus points us.
to in the Sermon on the Mount..
Now the last few weeks we've been talking about hypocrisy and specifically religious.
hypocrisy, hypocrisy in our faith, in our relationship to God..
And we've been talking about how serious it is, how it can be such a strong temptation.
to us as people of faith, people who take our faith seriously..
It can be a really strong temptation to pretend to be something that we're not..
And that's one of the reasons that Jesus gives it so much space, this topic..
He gives this topic so much space in Matthew chapter 6..
But we have also talked about how devastating and serious the consequences of it can be.
for us..
And that's another reason that Jesus gives this topic so much space..
So we're going to finish up today, but we're going to get just a little bit out of order.
with the verses because as Jesus is talking about this topic of hypocrisy, in the middle.
of this, he pivots and he talks about the correct way to pray because the passage we.
were discussing last week, he tells us not to pray like the hypocrites..
And then he gives us a definition of how to pray, a template for prayer..
And then he goes back to talking about hypocrisy and kind of finishing that up..
So what we're going to do is we're going to skip ahead to verse 16, finish up this topic.
of hypocrisy, and then we're going to go back to the section on the Lord's Prayer and kind.
of consider that separately on its own..
So this morning, we're going to skip ahead to verse 16..
Last week, we finished up there at verse 6, but now we're going to be down into verse.
16..
And it says this, "And when you fast, do not look gloomy like the hypocrites, for they.
disfigure their faces that their fasting may be seen by others..
Truly, I say to you, they have received their reward..
But when you fast, anoint your head and wash your face, that your fasting may not be seen.
by others, but by your Father who is in secret, and your Father who sees in secret will reward.
you.".
Now, as we've discussed hypocrisy, we defined it biblically and we defined it in a few different.
ways..
We defined it as pretending to be something that we are not or acting as if we are something.
that we are not..
We've talked a lot about how we can do that in our relationship to God and our faith,.
that we can pretend to be righteous, pretend to be holy when we are not, and how that is.
damaging to us, but also to the people around us..
We talked about how hypocrisy is also deceiving ourselves or others into thinking we can meet.
God's standards when we can't..

$^{41}$Hypocrisy happens when we misunderstand God's grace and we decide that we can meet God's.
standards of holiness and righteousness in life through our own effort, through just.
trying hard and through trying to produce good works and through trying to be a good.
person..
And we can't meet God's standards..
That only happens as we repent and as we surrender to Jesus and let Him live through us..
We talked about that..
We talked about our other definition of hypocrisy, that it's really turning the spotlight on.
ourselves instead of God, trying to glorify ourselves instead of trying to glorify God..
It's not that God wants us to hide all of our actions, but God wants us to make sure.
that in our actions that He is glorified, that the attention is put on Him and not on.
us, that our actions should not be to glorify ourselves, and especially not in our relationship.
to God and our faith..
And we use this beautiful definition from the scripture that Jesus gives us from Luke.
18 verse 9 that says this, "He also told this parable to some who trusted in themselves.
that they were righteous and treated others with contempt.".
And that's a great definition of hypocrisy in our faith, when we begin to trust in ourselves.
that we're righteous and treat others with contempt..
And that's one of the devastating consequences of hypocrisy, is that we treat others with.
contempt..
We're not kind and caring and loving towards others in the way we should be..
As we've discussed it, we've gone through a few different things..
We talked the first week about how maybe we've been hurt by hypocrisy in the church..
We've talked about how maybe there were misplaced standards of righteousness and who defines.
righteousness that we had in our life..
Maybe we were trying to live up to some standard that wasn't actually God's standards..
It was just other humans that said, "You should be like this if you're a Christian..
This is what it means to be righteous or to be a follower of Jesus or to be a Christian.".
And that wasn't God's standard, it was someone else's..
And trying to live up to that damages us..
We talked about those misplaced standards..
We talked about misplaced expectations and how if we fall into hypocrisy or we follow.
someone who is in hypocrisy, that we can put expectations on them or on ourselves that.
God never meant for us to have..
We can be deeply hurt and disappointed when we or that person we're following falls short.
of those expectations..
We talked about misplaced pain..
The fact that Jesus is not the source of hypocrisy..
Jesus is not the source of the things that harm us..
But sometimes when we encounter hypocrisy in the church or in our faith journey, we.

$^{81}$think that that comes from God or we think that that comes from Jesus..
And so we think that that pain is the fault of God when it's not..
It's the fault of other broken humans who have misinterpreted scripture or maybe they've.
just been plain wrong or maybe they didn't even know Jesus in the first place..
And because of that, they've lived a life of hypocrisy..
And then last week we got really convicting as we talked about our own tendency towards.
hypocrisy and how easy it is to fall into that, but also how serious it is..
We talked about the fact that when we fall into hypocrisy in our life, that ultimately.
we're stealing glory from God..
We're directing glory towards us that belongs to God and God alone..
And we talked about the fact that ultimately we want to steal God's throne..
We want to steal control from God..
Control belongs to God and God alone in our life and in the lives of those around us..
But we want to steal that away and give the image, give the perception that we're in control,.
that we can control what happens in our life and we can control our own success and our.
own righteousness and those things..
And that's just not true..
And we talked about the fact that we also want to steal destiny, God's will for our.
life, God's purpose and plan for our future..
We want to steal that away and we want to insert in its place our own purpose and plan.
for our future..
And that happens when we live in a place of hypocrisy, acting like we can control our.
own destiny when we can't..
That's something that has to be surrendered to Jesus..
And we talked about some of the practical applications of how to avoid hypocrisy through.
seeking humility and through doing what Jesus talks about, through being willing to do good.
works in secret that no one sees, that are hidden, because our heart is to do them out.
of love for God and love for others, not out of the recognition that we will get from anyone.
else and how important that heart is, that we do it because we want to please God and.
we want to love others, not because we're trying to please any other human..
And today we're going to talk about one other aspect of hypocrisy that comes from the verses.
we're looking at today..
That is a misappropriated burden, taking on a burden we're not supposed to take on..
And we see some of that in verses 16 to 18 here..
I'm going to read them again..
It says, "When you fast, do not look gloomy like the hypocrites, for they disfigure their.
faces that their fasting may be seen by others..
Truly, I say to you, they have received their reward..
But when you fast, anoint your head and wash your face, that your fasting may not be seen.
by others, but by your Father who is in secret, and your Father who sees in secret will reward.

$^{121}$you.".
So what's going on in this passage?.
Jesus is talking about fasting..
And maybe if you haven't been in church a lot, you're not sure what fasting is..
Or maybe if you have been in church all your life, you still don't really know necessarily.
what it is..
Well, fasting is something that God called his people to a long time ago in the Old Testament..
And it continued in the New Testament..
People have continued to practice fasting as a way to draw near to God..
And fasting is when you voluntarily deny some of your physical appetites in order to seek.
after God..
And usually it's fasting from food..
It can be from other things as well, but a lot of times it's associated with fasting.
from food..
So people say, "I'm going to not eat for a certain period of time," or "I'm going to.
only eat a few certain things for a certain period of time.".
And it's done not just for health reasons, although that might be a part of it, but it's.
not just done because it's going to help us to be healthier..
It's done for spiritual reasons to say, "Okay, I'm going to deny some of these physical appetites.
so that I'm actually able to seek God more and to be focused on him and to really grow.
in my relationship with him and draw near to God.".
That's what fasting is..
And that's what Jesus is talking about in this passage..
So fasting involves intentionality..
It involves a desire to know God..
And it involves some sort of sacrifice and saying, "Okay, I'm going to sacrifice some.
of my desires in order to become closer to God and to draw near to him..
And I'm going to demonstrate that I'm serious about God by denying some of those desires.
and appetites in my life.".
So when Jesus says that the hypocrites, when they go to fast, they are actually making.
it very clear and plain that they are fasting..
When he says they disfigure their faces, what Jesus is saying is that the hypocrites were.
intentionally looking very exhausted and very ragged so that anybody who looked at them.
would say, "Oh, look at that guy over there..
I mean, he's been fasting..
You can see it in his face..
He looks so tired..
He looks like he hasn't even been able to clean up..
I mean, what a spiritual guy.".
And again, it was something that was done for outward display to impress other people..

$^{161}$And as we spoke about the last few weeks, to glorify themselves, to put the attention.
on them instead of putting the attention on God..
That's what Jesus is referring to here..
And so he says in verse 17, "When you fast, anoint your head and wash your face.".
What Jesus is saying is, "Don't be like the hypocrites that you try to put your spirituality.
on display for others and try to show people the depth of your devotion or how intense.
your sacrifice is so that other people will look at you and say, 'Wow, what a committed.
person..
Their devotion to God is so deep..
Their sacrifice is so great.'".
Jesus is saying, "No, don't do that..
When you fast, go ahead and wash your face and look normal and go about life normally.".
Because what is going on in a fast is supposed to be about your relationship to God..
It's supposed to be about drawing near to him..
It's supposed to be about denying your physical appetite so that you can connect more with.
God..
It's not supposed to be about other people and impressing them in what they think of.
you..
And so it's literally counterproductive to the fast itself for you to make it about other.
people's opinion of you and impressing other people with it..
That's what Jesus is pointing out..
And he's saying, "Your father knows that you're fasting..
If you're sacrificing in some way to connect with God, God knows that..
And he will reward you..
He will reward you with a deeper relationship with himself..
But you shouldn't do it to be seen by others.".
Now we've already discussed a lot this principle of not doing things to be seen by others and.
how easy it is to fall into hypocrisy there..
So what I want us to think about today is just this principle of sacrifice..
It's involved in fasting and it's involved in other parts of our walk with Jesus..
And sometimes we think that that sacrifice is something that we're supposed to take on.
in order to basically be some type of super spiritual, super human, super Christian..
And that's totally not true at all..
And that's part of what Jesus is pointing out here..
He's saying that was a problem with the hypocrites..
They were trying to take this sacrifice on and let everybody know how great they were.
sacrificing so that everybody would see that they were a superhero..
And that's not the point..
Jesus is the hero, not his followers..
And we do that too..

$^{201}$Sometimes we misappropriate this burden..
We try to take on some burden that we were never supposed to have..
And we try to shoulder up under something and show that we are doing this great sacrifice.
for God..
And we get a little bit of a Messiah complex..
We get a little bit of this idea in our head that we're actually supposed to go out and.
save the world..
Or we're actually supposed to really transform someone's life..
Or we're actually supposed to turn someone around and take them from a place of struggle.
and darkness and confusion to a place of beautiful light and flourishing and peace..
And we're not supposed to do that..
God is supposed to do that..
Jesus is supposed to save the world..
Jesus is supposed to transform lives..
Jesus is the one who takes people from hurt and brokenness and affliction to peace and.
joy and beauty..
Jesus is the Messiah, not us..
And we begin to operate from that place of having a little bit of that Messiah complex,.
thinking that we're supposed to save the world..
First of all, there's an obvious amount of arrogance that plays into that..
The fact that we would even think that in the first place..
The fact that we would even think that we're able to save the world or transform someone's.
life in the first place shows that there's some pride in our heart that is not pleasing.
to God..
We need to be in a place of saying, "God, I understand that apart from you I can do.
nothing.".
Like Jesus talked about when he talked about abiding in Christ..
And we need to understand that..
We need to say, "I'm not the one who's here to change things..
I can't save other people..
I can't even bring about transformation in my own heart, much less someone else's heart.".
So we need to get away from that arrogance of feeling like we're supposed to bear this.
burden of saving the world..
But we also have to get away from the hypocrisy of that Messiah complex, thinking that we're.
going to save the world..
That causes us to have to convince ourselves and the people around us that we don't make.
any mistakes, that we're perfect..
And that leads us straight into all the things we've been talking about the last few weeks..
This pretending to be something we're not..
This acting like we meet God's standards when we don't..

$^{241}$This treating other people with contempt because we think that we're more righteous than someone.
else..
And so we have to stay away from that type of thinking because it leads us straight into.
hypocrisy..
And we also have to stay away from that type of thinking because it is impossible..
Thinking that we are going to save the world, having that Messiah complex, that leads us.
to a place of impossibility..
We're basically trying to pick up a load that we cannot carry and we will ultimately be.
crushed by it..
If we start to convince ourselves in our own mind that we're supposed to save this person.
or that person or this situation or this neighborhood or whatever, we are trying to shoulder a load.
that will crush us..
As I said a moment ago, we cannot even transform our own heart, much less someone else's heart,.
much less transform a city or a neighborhood or a community or whatever..
Now Jesus can..
That's the beautiful thing..
That's the good news is that Jesus can and he wants to work through us..
We get to be instruments..
We get to cooperate with God in seeing transformation happen in our neighborhood and in our family.
and in our friends' lives..
It's not that we're not going to be a part of that, but we have to understand that we.
do not do that..
We do not start it..
We do not bring it to pass..
Jesus starts it..
He uses us to bring it to pass, but it's a work that he and only he can do and it's a.
work that he does as he lives through us..
If we feel like we're supposed to save the world or save this person or that person or.
this situation or that situation, we are taking on a burden that was never ours to carry..
We are misappropriating that burden and ultimately we're going to find ourselves crushed by that.
burden..
And so we may say, "Well, you know, look at my sacrifice..
I want other people to see that..
I want other people to know that I'm carrying this heavy load," but we'll fall into hypocrisy..
We'll fall into arrogance..
We'll fall into impossibility and we'll be crushed by the magnitude of that impossibility..
We cannot carry that burden..
I think that's good for us to think about as followers of Jesus..
We should not pick up burdens that we were never supposed to carry..
We need to pray for the people in our lives, our family, our friends, our community, our.

$^{281}$city..
We need to love and we need to serve with the love of Jesus..
We need to let Jesus live through us..
But we cannot feel like that we have to make change happen and if it doesn't happen, then.
it's our fault and that we have failed..
That is not the role that God has called us to..
God has called us to be an instrument, to be a vessel, to be someone that He can use..
God has not called us to be the Savior, the Fixer, the Healer..
We can't do that and we need to understand that..
We need to operate and live our lives from a place that says, "I want to cooperate with.
God, obey and be sensitive to Him and put myself in the right position so that He can.
live through me, but I'm not here to fix it..
I'm not the one with all the answers..
I'm not here being proud and arrogant and thinking that I'm going to transform things..
I understand that I am a broken human who God is at work in my life and as He works,.
He can use that work to help someone else, but it's not about me.".
Now speaking of misappropriated burdens and a Messiah complex and hypocrisy, I want to.
close out today with just a few thoughts about the upcoming election..
Don't worry, I'm not going to try to influence your voting in any way or try to berate your.
political views in any way, so this is not going to be one of those weird, awkward conversations..
But I just want to give us some encouragement and some things to think about..
Obviously we're just days away from November 3rd and this election and this is a really.
important election..
You don't need to hear that from me..
I'm sure you've heard it from many people already..
So I want to begin just with encouraging you to vote..
Hopefully you've already been able to vote through mail-in voting or early voting, but.
if you haven't, I would encourage you to go vote on Tuesday..
I understand, I know our political system in the US is really full of deception and.
hypocrisy and division right now and sometimes that makes it kind of distasteful to be a.
part of and sometimes we want to disengage, but it is something that's important for us.
to participate in as followers of Jesus..
It's something that we want to seek to encourage kingdom values in our world through our politics.
and in the United States, God has given us this opportunity to do that through voting.
and so that's a really important privilege and responsibility that we should participate.
in and we should take part in and take that seriously and make sure that we do it..
God has delegated authority to earthly governments for the purpose of protecting human lives,.
and promoting human flourishing and that's something that God wants us as his followers.
to seek to do, to seek to promote human flourishing through the way we participate and the way.
that we vote..

$^{321}$So we want to do that..
Now obviously this election is going to be really hotly contested and I have no idea.
what's going to happen or what the results will be, but I just want to talk about some.
of the things that we might feel after the results because this is an election that people.
feel very passionately about, very strongly about..
People have deeply held emotions about this election and our politics is so polarized.
and so divided right now that most people have a side, they have a team, they have a.
candidate or they have a group of candidates or a party that they are really emotionally.
invested in and so if you're one of those people, then depending on the results of the.
election you might feel a lot of things after Tuesday, on Wednesday or through the rest.
of the week..
And I just want to give us some perspective and some thoughts about those things, especially.
if we're experiencing some negative emotions afterwards..
I understand that we may not know on Tuesday night or Wednesday morning or Thursday who.
the winner of many of the races are, especially with the pandemic and such an unusual year.
in voting..
It may be days or weeks or even possibly months before some of these races are decided, but.
the reality is the way our political system in the US is set up, the elections fall into.
categories of winners and losers..
Once the final tally is in, there are winners and losers..
So first let's talk about the scenario of your side or your candidate or your group.
of candidates or your party or whatever being on the losing side, feeling like that you.
have lost in this election..
And I want to encourage you to not fall into despair..
Don't fall into frustration with other people..
Certainly don't fall into frustration that causes you to dehumanize the other side or.
the other party and feel like they're the source of all the problems in your life or.
in our country or in the world..
That's simply not true and it's dangerous for us to think that way..
God always wants us to regard others as made in His image and as worthy of our love and.
respect and our kindness and our service..
Even if we disagree with them, even if we think they're doing something wrong, God still.
wants us to love them, to pray for them, to care for them in any way that we can..
And if they don't know Jesus, to seek to connect them with God so that they can come into relationship.
with Jesus..
And that's how we should respond..
So if you feel like your side is losing, don't get upset or angry at the side that wins,.
even if they are completely unkind and arrogant in their victory..
And don't misappropriate the burden of saving souls or saving our country and think that.
the soul of our country is won or lost with elections..

$^{361}$The reality is Jesus is the only one who saves souls and He's the only one who can save the.
soul of our country..
And yeah, in the United States we have an executive branch in our government that has.
a president at the top of it..
But as a follower of Jesus, you have a foundational identity that is much deeper than your ethnic.
identity or your national identity..
And so you don't serve a president, you serve a king..
Jesus is the king that you serve and the kingdom of God is unshaken by elections..
And so the reality is that on the morning of November 3rd, you'll wake up as a subject.
of the kingdom of heaven..
You will wake up as someone who is loved by God and who's been given a mission to love.
God and to love and serve the people around you, to let Jesus live through you, to love.
those around you..
And on the morning of November 4th after the election, November 5th, every day going forward,.
you will wake up with that exact same identity..
You will wake up as a subject of the kingdom of God..
You will wake up as someone who is loved and forgiven by God..
You will wake up with this mission of loving God and loving others and allowing Jesus to.
live through you to love others..
That will not change..
So there's no reason for you to feel despair even if you feel like your side has lost in.
the election..
You haven't lost anything..
Your identity is unshakable and it lies at a level that is untouched by human decisions..
God is not changed..
God is still on the throne..
Jesus is still in control of the world and your identity in Christ is completely untouched.
by the election results or anything else..
Now let's consider the other scenario and this is perhaps an even more dangerous one..
What if you're on the winning side?.
What if your side wins and your candidate wins or your party wins or the candidates.
that you prefer win?.
One of the dangers there is that you might feel this sense of victory and kind of self-congratulation.
and slip into a sense of complacency or a false sense of security because you feel like,.
"Well, my side won..
The people that I want to be in office are now in office and everything's going to get.
better..
Everything's going to turn out great..
The world's going to be just fine.".
And that can cause us to get to a place of complacency in our obedience to God and that's.

$^{401}$really dangerous for us as followers of Jesus..
Our culture has attached a strong element of morality to our politics and our voting.
process and sometimes there's this underlying assumption that as long as you vote and vote.
in the correct way, which of course is defined differently by many different groups and many.
different people, but as long as you vote in the correct way, then you've done some.
type of moral good and you've done your moral duty and that's all there is really to do.
to fix things or to be a part of positive change in the world is just to go out and.
vote and to vote in the correct way..
And that simply isn't true for the followers of Jesus..
For the followers of Jesus, voting is very important, but it's only a very small piece.
of a much bigger picture of how God wants us to be involved in seeing change happen.
in our world and in seeing people cared for and in seeing the love of Jesus go out into.
our communities and our neighborhoods..
And the reality is that no matter how great you think your candidate is and your party.
is and even if every race is run by your party and every candidate that you want to get into.
office gets into office, there will still be poverty even after those people are installed.
in office..
There will still be racism and discrimination after those people are installed in office..
There will still be violence and people killed and children killed after those people are.
installed in office..
There will still be substance abuse and addiction..
There will still be people abusing one another and oppressing one another..
There will still be coronavirus even after your favorite candidate is installed in office..
Your candidate winning doesn't mean an end to evil or corruption or anything else..
So you need to understand that as a follower of Jesus, God still has so much that he wants.
to use you to do in our world..
Our world is still broken and Jesus still wants to transform your life and use you,.
live through you to love and serve and transform the lives of others..
Don't become complacent or comfortable because your side or your party or your candidate.
does well on Tuesday..
Understand that God still has so much that he wants you to cooperate with him to be a.
part of..
And remember that regardless of what happens on Tuesday, your identity is unchanged..
You're not a winner or a loser..
You're forgiven and loved..
You're covered by the grace of Jesus..
Your mission is unchanged..
You're called to love God and to love the people around you..
You're called to let Jesus transform your life and live through you to love others and.
transform the lives of people around you..

$^{441}$You're called to live out the values of the kingdom..
So go out and vote..
Vote the values of the kingdom and then go right back to living out the values of the.
kingdom..
Continue to love God..
Continue to love others with the love of Jesus..
Continue to serve and care for your family and your community and your neighbors and.
your church and your city..
Don't despair..
Don't get frustrated..
Don't get complacent and comfortable and apathetic..
Don't misappropriate the burden of saving the world..
Don't think that burden belongs to you..
Don't think it belongs to your political party..
Don't think it belongs to your favorite candidate..
It doesn't..
Most of the people involved in politics probably have some level of Messiah complex..
They think they can save the world, but they can't..
There's only one Savior and it's Jesus..
And the good news is that he is mighty to save..
Let him do his work in your heart and in your life and in the lives of people around you..
Let's pray..
God I pray that you would help us not to personally misappropriate that burden and think that.
we are supposed to save the world or transform lives or change everything..
We cannot..
And it is arrogance and pride and hypocrisy for us to think we can..
But God, also help us not to misappropriate that burden and place it on something else.
or someone else, on our politics or on our political party or our favorite candidate..
Help us to see that only you can transform lives..
Politics at best can maybe modify people's behavior, but that's really all it can do..
Only you can change a heart and change someone from the inside out..
And God, that's what our world desperately needs..
And so I pray that you would help us to cooperate with you, to put ourselves in a position so.
that you can live through us, so that you can truly be transforming our hearts and live.
through us to transform the lives of people around us..
God, may you be glorified in our lives and may we be used of you to love and serve others.
and to connect others with your kingdom so that more and more people are able to know.
you and to walk in your grace and to walk in your love and to experience your goodness.
and so that we see more and more of your shalom in our world..
God, in Jesus' name..

$^{481}$Amen..
\newpage



\section{}
\label{sec:LPCHL652i_w}
\textbf{浩劫重生 【第5集:亂局中的以利亞】 陳崇基牧師 - 主日講道視頻系列}
\newline
\newline
連結: \href{https://youtube.com/watch?v=LPCHL652i_w}{\texttt{https://youtube.com/watch?v=LPCHL652i\_w}} ~~~~ 語音日期: 2020-11-10
\newline
\newline
\hyperref[sec:lgq1HtuVP6w]{\small{< < < PREV SERMON < < <}}
~
\hyperref[sec:index_chronic]{\small{[返順時目]}}
~
\hyperref[sec:index_scriptual]{\small{[返順卷目]}}
~
\hyperref[sec:rfX8QtfxB9g]{\small{> > > NEXT SERMON > > >}}
\newline
\newline
$^{1}$我們想起2020年的時候.
我們相信不會單單想起天災人禍接二連三的發生.
我們也會想起世界大亂.
有很多未知之數.
不單止是有疫情.
我們也會面對政局,社會很多不明朗的時候.
聖經其實有很多這樣的記載.
講到政局很混亂.
很不明朗.
令我想起以利亞的時代.
當時的背景.
是阿哈作以色列王.
聶旺基上第十六章.
猶大王阿薩三十八年.
暗利的兒子阿哈登基作以色列王.
暗利的兒子阿哈在薩馬利亞京城作以色列王二十二年.
暗利的兒子阿哈行義和還眼中捍衛惡的事.
比他之前的列王更甚.
犯了尼伯的兒子耶路播恩所犯的罪.
他還義為兄.
又娶了西頓王潔巴力的女兒耶洗別為妻.
去到事奉敬拜巴力.
在薩馬利亞京城建造巴力的廟.
在廟裡面為巴力祝壇.
阿哈也做了阿舍拉.
是一種生育的女神.
他所行的是惹惹禍禍以色列神的怒氣.
比他之前的以色列諸王更甚.
這是聖經對阿哈王正職的評估.
阿哈王朝大概是公元前八百七十至八百五十年之間.
當時已經分了南北國.
北國就是以色列.
京城在薩馬利亞.
或者另一個王宮耶斯利.
南邊就是猶大國.
京城在耶路撒冷.
阿哈王作以色列王的時候.
他藉著很多的征戰和外交.
將他的勢力擴張到很大.
他的影響力是由地中海伸展到約旦河東.

$^{41}$由北邊的旦山區.
一直伸展到南邊的摩亞地.
是一個歷史最輝煌的時候.
幾乎可以跟所羅門聯合王國的時候相比.
因為他擴張了版圖.
他是雄霸著沿海的貿易路線.
以及約旦河東的國際貿易路線.
如果從一個經濟發展的角度來看.
你就知道為什麼阿哈王會娶西倫的公主.
夜使別為妻.
因為首先從地理來講.
以色列的海岸線是比較直的.
是沒有一些深的大的海灣.
全部都是淺灘小的海港.
在東地中海這個世界裡面.
大的海港.
深的海灣.
全部都集中在北邊.
推羅,西頓,比布魯斯,烏加利等等.
這些海岸線的海港.
我們統稱為.
Felicia.
在聖經的世界當中.
Felicia就是地中海遠洋貿易的霸王.
相等於今天一些大的海港城市.
比如香港,新加坡等等.
所以阿哈王娶夜使別為妻.
其實是一種外交的手段.
以色列國如果要繼續發展的話.
他就一定要跟一些海洋貿易的城邦建交.
西頓就是當時其中一個最強的海洋貿易的城邦.
以前的宗教,經濟,政治等等是融為一體的.
是分不開的.
Felicia也是拜巴力的大本營.
所以邦交就要混合耶和華的信仰.
和巴力的信仰.
巴力是風雨的神.
在烏加利這個海港中.
我們找到很多當時古代的禮版.
最出名的一段文獻就是巴力的神話.

$^{81}$這就是巴力.
巴力的右手是拿著雷電的.
他有時會站在牛犢的上面.
就是說巴力在神話當中是控制著風雨的.
他是賜雨給大地的.
所以巴力有一個綽號.
就是Rider of Cloud.
是駕運者.
所以當阿哈做王的時候.
以色列不但有巴力廟供奉巴力這個偶像.
也供養了很多巴力的假先知.
同樣也是殺害耶和華的真先知.
在上帝眼中他看到的不是一個輝煌稱霸的國度.
在上帝眼中看到的是一個最邪惡的暴君.
是一個離經背道的時期.
所以神就藉著以利亞宣告三年的寒災.
巴力這位先知到京城和阿哈王講.
十七章第一節.
基列所寄居的提比斯人以利亞.
他去到見阿哈講.
他說我指著所事奉永生的耶和華.
以色列的神起誓.
這幾年我若不禱告.
必不降奴不下禦.
其實這句話就是耶和華藉著以利亞.
來挑戰阿哈的王朝.
也是挑戰他所拜的巴力.
所以三年的寒季過了之後.
列王記上帝十八章.
這裡我們就讀到一個故事一個事蹟.
就是我們耳熟能詳.
就是以利亞和幾百個巴力的假先知來比試.
比試的地方是加密山Mt.Karmel.
為什麼在這裡.
因為看地圖的話你就明白.
加密山Mt.Karmel.
就正正在弗里基佛利時亞和以色列國Israel中間這個地方.
這個可以說是一個爭霸地區.
並且每年雨季來到的時候.
地中海的雨雲吹來以色列的時候.

$^{121}$第一個著陸的地方就是加密山.
在近東的文獻也告訴我們.
加密山是一個拜巴力的地方.
它可以稱為Bale of the Head of Karmel.
就是加密山頭上的巴力.
所以如果要比試哪一個真正控制風雨的話.
是巴力還是耶和華呢.
加密山就是最好的地方.
於是在加密山.
雙方耶和華的先知以利亞一個.
和巴力的假先知幾百個.
就各自捉起祭壇.
巴力的先知由早到晚就去求巴力顯靈.
叫巴力來降火.
證明巴力才是風雨淚電的神.
第十八章第二十八節.
說到巴力的假先知.
大聲地求告.
按照他們的規矩用刀和槍自割自刺.
直至身體流血.
到午後即到天晚的時候.
他們狂呼亂叫.
但是沒有聲音.
巴力也沒有營運.
之後到以利亞.
他重建了毀壞了的耶和華的祭壇.
最後他去祈禱.
十八章第三十八節.
於是耶和華就降下火來.
即是有雷電.
然後燒盡梵濟,木柴,石頭和塵土.
又燒乾溝裡面的水.
然後以色列的眾民看到.
耶和華真的這樣降火的時候.
他們就俯伏在地.
他說耶和華是神.
耶和華是神.
以利亞在加密山上.
大聖巴力的眾先知.
之後我們看到一些很生動的描述.

$^{161}$就是以利亞上到山頂.
然後祈禱耶和華降雨.
十八章四十五節.
他說:煞時之間.
天就因為風雲黑暗降下大雨.
阿哈王就坐車往他的王宮耶斯列.
耶和華的靈降在以利亞的身上.
他就像得勝一樣.
去到蓄氣邀奔在阿哈的前頭.
直到耶斯列王宮的城門.
那個意境你可以想像得到.
就是先知以利亞得勝了.
他就在前頭帶著阿哈君王的前尺.
回到他國家的王宮.
這個時候就是帶領著整個國度能夠復興的時候.
誰知就來了一個反高潮.
第十九章第一節.
阿哈將以利亞一切所行的.
以及他用刀殺了眾假先知的事.
都告訴了耶和華王.
耶和華王就差遣人去見以利亞.
告訴他明日約在這時候.
我若不使你的性命像那些人的性命一樣.
願神命重重的降伐於我.
以利亞見到這個光景就害怕.
然後就逃命了.
到了猶大的別士巴.
將他僕人留在那裡.
自己就在曠野走了一天的路程.
來到一棵樹叫做羅騰樹之後.
他就坐在那裡求死.
耶和華王說罷了.
求你取我的性命.
因為我不勝於我的列種.
他就躺在羅騰樹下睡著了.
有一個天使拍他說.
起來吃吧.
和炭火燒的餅.
他就吃了喝了.
仍然躺下.

$^{201}$耶和華的使者第二次.
輕輕拍他.
起來吃吧.
因為你當酒的路甚遠.
他就起來吃了喝了.
仗著這個飲食的力量.
走了四十晝夜.
到了神的山就是賀列山.
開始在加密山見到那個.
威風凜凜.
戰無不勝的以利亞.
一個這樣的英雄.
突然之間變成一個出走的.
抑鬱的耗盡的失敗的人.
掉進了低谷.
在他和上帝的互動之間.
又同樣見到上帝關心他的需要.
帶他走四十天的路.
帶他去到神的山.
在那裡最安全的時候.
上帝才和他說話.
向他顯現.
然後再恢復以利亞.
一個很多人能夠有同感的經歷.
就是當我們跌入低谷的時候.
上帝怎樣去眷顧.
怎樣去供應我們.
怪不得這個故事.
是最多用來做心理輔導.
做心理治療的材料.
我們看看聶旺記上帝十九章第九節.
以利亞來到神的山.
西乃賀列山的時候.
在那裡進了一個洞.
就住在洞中.
耶和華的話想到他.
以利亞 你在這裡做什麼.
以利亞就回答.
我為耶和華萬軍之神大發熱心.
因為以色列人背棄了你的約.

$^{241}$毀壞了你的壇.
用刀殺了你的先知.
只剩下我一個.
他們說要尋索我的命.
耶和華說.
你出來站在山上.
在我面前.
那時耶和華從那裡經過.
在他面前有烈風大作.
崩山碎石.
耶和華卻不在風中.
風後有地震.
耶和華卻不在其中.
地震過後有火.
耶和華也不在火中.
火後有微小的聲音.
以利亞聽見.
就用外衣蒙上褸.
出去站在洞口.
有聲音對他說.
以利亞 你在這裡做什麼.
他說.
我為耶和華萬軍之神大發熱心.
因為以色列人背棄了你的約.
毀壞了你的祭壇.
用刀殺了你的先知.
只剩下我一個.
他們還要尋索我的命.
耶和華對他說.
你回去.
從曠野的路往大馬士革去.
到了那裡就高納哈舍作亞蘭王.
高納令氏的孫子耶護作以色列王.
並告阿伯米荷拉人沙發的兒子.
以利沙作先知接續你.
將來躲避哈舍之刀的.
必被耶護所殺.
躲避耶護之刀的.
必被以利亞所殺.
但我在以色列人中.

$^{281}$為自己留下七千人.
是未曾向巴黎屈實的.
未曾向巴黎親隨的.
一個很長的記載.
我們當然看到.
上帝如何去眷顧他的先知.
他的僕人.
在這裡供應他一切的需要.
去跟他對話.
然後好像輔導他一樣.
去恢復他.
但是似乎那個描述.
不只是單單這樣.
你留意一下這裡出現了什麼字句.
或者什麼意境.
是令我們想起.
在舊約裡面另外一個人.
我們有曠野.
有神供應40天.
以利亞來到神的山.
以利亞來到這裡的時候.
他躲在盤石岩洞裡面.
然後神的榮耀經過.
也說到以色列人背若立若的問題.
我們想起誰呢.
就是摩西.
因為摩西.
大利以色列人出了埃及之後.
來到賀列山的時候.
他的經歷正正就是這樣.
出埃及記第十九章.
他說到了第三天的早上.
在山上神就顯現了.
有雷電.
有閃電.
有蜜雲.
並且有覺星神大.
在西大山樹.
有霧煙.
有火.

$^{321}$神在火裡面降臨.
然後地大大的震動.
號覺星漸漸高.
漸漸大聲的時候.
摩西就說話了.
上帝就有聲音來答應他.
上帝就和百姓來到立約.
跟他們說話.
這個其實就是摩西.
看到神顯現的情況.
現在以利亞好像另一個摩西.
在同樣神的山的樹.
看到神來到顯現.
幾乎是一模一樣.
但是他有很大的分別.
耶和華說.
你出來站在山上.
在我的面前.
那時耶和華從哪裡經過.
好像在摩西面前經過一樣.
在他面前.
有烈風大作崩山碎石.
耶和華卻不在風中.
風後有地震.
耶和華卻不在其中.
地震後有火.
耶和華也不在火中.
火後有微小的聲音.
耶和華卻不在其中.
說了三次.
以色列人預期上帝.
有權能降臨.
有權能彰顯他的大能.
但這個故事告訴了以利亞.
上帝不一定要按照我們的權能神學來顯現.
你想想當時以利亞在想什麼.
以利亞當然在想.
過了加密山的比試之後.
已經證明了耶和華就是神.
耶和華是有能力的.

$^{361}$在加密山上可以使暴君阿哈屈服.
照樣的神就應該去降火降雷電.
使邪惡的皇后耶洗別也屈服.
為什麼上帝沒有這樣做.
相反耶洗別可以威脅並殺害剩下的先知.
所以為什麼以利亞會立即逃亡.
因為他發現上帝不是按照他的方程式來做事.
似乎告訴我們以利亞為上帝大發熱心.
可能是一件好事.
但很多時候這種熱心往往會成為.
上帝一定要按照我的工程式來做事.
事實上如果你想想.
這些火,雷電,地震在巴力的神話中.
都是這樣形容巴力的.
每逢巴力出現的時候會有什麼現象.
就會有大風,火,雷電.
當你看清楚的時候.
你會發現耶和華所做的就是想告訴先知.
我耶和華是真神,不像巴力.
我耶和華是真的神.
也不是一個更厲害的巴力神.
我之所以是真神.
就是不會一定按照你的預期來做事.
可能這個故事就是告訴我們.
我們需要改變一下我們對上帝的觀念.
我們的theology.
可能這樣會幫助我們在谷底.
在低潮的時候,在我們掙扎的當中.
在亂局當中.
幫助我們從低谷上升.
讓我們認定上帝依然是上得過.
雖然上帝未必一定會按照我的預期.
按照我的工程式來做事.
在這個星期是我們美國總統的大選.
無論大選的結果如何.
總會有一大批人不高興.
也總會有一大批人開心.
如果你是因為大選的結果而感到非常失望.
我想在這裡告訴你.
請你小心.

$^{401}$如果大選的結果令你失望.
你覺得上帝不按照你的預期去做.
你對上帝失去了信心.
你就更加要小心.
但如果因為大選的結果而你非常開心.
你也可能更加要小心.
請你不要搞錯.
不要以地上的政權來代表上帝的國度.
因為沒有一個黨派.
沒有一個政權能夠代替上帝的國度.
其實我們可能也會犯了以利亞有的毛病.
想起上帝是否應該這樣去工作.
只要他崩山碎石.
做一些驚天動地的事.
他就把暴君的政權扭轉過來.
既然可以扭轉過來.
一切就會復甦.
上帝的國度一定會臨到.
這不是我們很多時候.
基督徒面對政權的時候.
我們有的事態.
我們覺得上帝一定要按照方程式來進行他的計劃.
上帝一定要這樣來彰顯他自己的權能.
這裡告訴我們.
上帝卻不在火中.
卻不在這些崩天動地的事情裡面.
上帝真正的工作.
他能夠繼續去延續他天國裡面的計劃.
是微小的聲音.
是他的說話.
你想想賀列山,西乃山.
最重要的都是這一件事.
在第一次上帝降臨於西乃山的時候.
他有驚天動地,地震,有風有火等等.
他向摩西和眾民顯現的時候.
最終要做什麼?.
就是他要頒布什戒.
他要跟上帝的子民納藥.
以至這些子民能夠成為萬國的光.
按照他的說話來生活.

$^{441}$去彰顯上帝的真理和公義.
上帝的子民第一件事要做的.
就是要聽上帝的說話.
聽上帝的聖音.
以利亞另一次像摩西那樣的經歷.
見到上帝親自降臨.
見到火,地震,風.
上帝卻不再記住.
要說明一件事.
如果我們繼續推進上帝的工作.
我們不一定要上帝.
去做一些驚天動地的事情.
去行神蹟.
才叫做彰顯上帝的榮耀.
真正去推進上帝的計劃.
就是一群願意單單事奉上帝的人.
單單願意聽上帝的話.
遵行上帝的話語的人.
留意一下.
賀烈山,以利亞所見到的.
有四樣.
烈風大作.
也就是說他不在風中.
地震也說他不在其中.
有火也說他不在火中.
火之後有微小的聲音.
這四樣東西你再平衡去看.
你就會知道微小的聲音.
就是上帝吩咐以利亞的四件事.
第一件就是高哈舍作亞蘭王.
亞蘭是外邦的政權.
高哈舍作亞蘭王.
就是說神仍然掌管外邦的政局.
第二是高耶護作以色列王.
就是上帝在掌管以色列的內政.
第三是要高納爾尼沙作下一代的先知.
這些人.
就是高哈舍,耶護和新的先知以利亞.
就是會代上帝審判阿哈王朝和耶世別.
第四就是留下七千人.

$^{481}$這些是未向巴黎屈膝和親聚的.
這些就是跟隨耶護話說話先知的群體.
如果你這樣平衡去看的話.
上帝真正能力的彰顯.
就是他的說話.
就是上帝所說的話語必定會兌現.
如果要看到上帝的話語兌現.
你也不一定立刻看到.
因為如果繼續追這個故事的話.
你要讀到聶望基下的時候.
你才會有揭曉.
阿哈王朝,耶世別王合.
受審判被殺.
要到聶望基下第十章才能揭曉.
第十章說什麼?.
由此可見.
耶護話指著阿哈加所說的話一句沒有落空.
因為耶護話藉著他僕人以利亞所說的話都成就了.
這句話在上下文出現了幾次.
就是上帝最後會兌現他所說的話.
而你要明白.
故事發展到這裡的時候.
以利亞早已不在人世.
問題就是這樣.
到底我們是否堅持聽神話語的人.
我們是否先知的群體.
是否真正上帝的代言人.
明白上帝所說的.
遵守上帝所言.
無論上帝今天是否打雷降火這樣工作.
還是他不會這樣工作.
上帝對我們的要求都是一樣.
就是你要聽我的說話.
行公義,維護我的真理.
作為一個先知的群體.
作為那七千人.
是不會向巴黎屈膝的.
最後在過去的一個星期.
在美國大選.
有很多不同很紛亂的訊息.

$^{521}$令我經常想起這篇詩篇.
這篇詩篇我想用它作一個結束.
詩篇146篇.
你們要讚美耶和華.
我的心你要讚美耶和華.
我的心要讚美耶和華.
我還活的時候要歌頌我的神.
你們不要依靠君王,不要依靠世人.
他一點不能幫助.
他的氣一斷就歸回塵土.
他所打算的當日就消滅了.
以雅各的神為幫助.
仰望耶和華他神的.
這人便為有福.
耶和華做天地海和其中的萬物.
他守成實直到永遠.
他為受我的神源.
似食物與饑餓的.
耶和華釋放被囚的.
耶和華開了核子的眼睛.
耶和華負起被壓下的人.
耶和華喜愛義人.
耶和華保護己居的.
父子孤兒和寡婦.
卻使惡人的道路彎曲.
耶和華要作王直到永遠.
錫安娜你的神要作王直到萬代.
你們要讚美耶和華.
求主幫助我們.
讓我們能夠成為一群.
不向巴黎屈膝.
只敬拜耶和華的群體.
去聽他的說話.
遵行他的心願.
我們同心大力祈禱.
上帝求你親自幫助我們.
讓你自己的同在能夠彰顯.
就是在你微小的聲音.
在你的話語當中.
向我們彰顯.

$^{561}$主要是幫助我們.
不要跌落一個試探裡面.
去到依靠地上的君王.
地上的政權.
真正的國度.
真正的榮耀是屬於你的.
只有你才是我們敬拜跟隨的主宰.
求你讓我們成為你喜悅的人.
讓我們相信你.
你自己的說話必定會彰顯.
你的國度必定會彰顯.
求你自己的旨意臨在地上.
我們這樣祈禱奉主命咁.
忘戀裡難以安靜.
靈內充斥了世俗理想.
但我確信你的話.
是我唯一的方向.
如若我願意打開心.
誓聽主聖音就夜禱告.
求被聖靈拆悔改造.
讓你來引領我路.
我心裡渴望你.
謙卑到你的面前.
只想聽到你的聲線.
耶穌我的心渴想你.
感恩獻上眾贈.
活生歸拜保佐賜.
唯願貼近你的心意.
耶穌我的心永屬你.
忘戀裡難以安靜.
靈內充斥了世俗理想.
但我確信你的話.
是我唯一的方向.
如若我願意打開心.
誓聽主聖音就夜禱告.
求被聖靈拆悔改造.
讓你來引領我路.
啊~.
我心裡渴望你.
謙卑到你的面前.

$^{601}$只想聽到你的聲線.
耶穌我的心渴想你.
感恩獻上眾贈.
活生歸拜保佐賜.
唯願貼近你的心意.
耶穌我的心永屬你.
其實你是想撫緩我.
從未會離棄不顧我.
但我心冰冷封閉.
假裝聽不到你聲音.
誰像你仍然等待我.
遺忘往來回應你的愛.
讓我將所有意念交上.
真正信服你的旨意.
耶穌我的心渴想你.
啊~.
我心裡渴望你.
謙卑到你的面前.
只想聽到你的聲線.
耶穌我的心渴想你.
感恩獻上眾贈.
活生歸拜保佐賜.
唯願貼近你的心意.
耶穌我的心永屬你.
耶穌我的心渴想你.
\newpage



\section{}
\label{sec:rfX8QtfxB9g}
\textbf{Kingdom Values XXX: Matthew Chapter 6 Prayer - Pastor Sam Johnson}
\newline
\newline
連結: \href{https://youtube.com/watch?v=rfX8QtfxB9g}{\texttt{https://youtube.com/watch?v=rfX8QtfxB9g}} ~~~~ 語音日期: 2020-11-14
\newline
\newline
\hyperref[sec:LPCHL652i_w]{\small{< < < PREV SERMON < < <}}
~
\hyperref[sec:index_chronic]{\small{[返順時目]}}
~
\hyperref[sec:index_scriptual]{\small{[返順卷目]}}
~
\hyperref[sec:MzNtNdJRGgQ]{\small{> > > NEXT SERMON > > >}}
\newline
\newline
$^{1}$Hey family, we are continuing on this morning with this study of the Sermon on the Mount,.
the Kingdom Values that Jesus teaches us..
And we're looking specifically again at this section in Matthew 6 on prayer, where Jesus.
is teaching us how to pray..
And this is what we commonly call in the church, the Lord's Prayer..
It's a model prayer for us, kind of a template that Jesus gives us in learning how to pray..
And that's so important..
Prayer is just so important for us..
Prayer has always been important for believers..
This is one of the fundamental pieces of our health, our spiritual health as Christians,.
as followers of Jesus, is to have a strong and a vibrant prayer life..
And last week I mentioned that I think this is something that a lot of times we're missing,.
especially in the modern world, especially as believers in North America and the United.
States..
We don't pray enough..
We a lot of times don't really know how to pray..
And especially as young people, we don't pray enough and we don't really know how to pray..
And so this is just huge..
This is so important..
This is so big for us to get..
And I'm excited that we're looking at it..
And I hope that we can really get a lot from what Jesus is teaching us here, because it.
will be such a blessing and such a benefit to our spiritual life and our spiritual health.
and our relationship to God..
Now last week we talked about the fact that Jesus was teaching us that prayer is a personal.
conversation with God and how important it is for us to understand that..
Sometimes we don't really know what prayer is and it's confusing..
Obviously it's also difficult sometimes because there's discipline involved..
It's not easy to set aside time to pray..
It's not easy to focus on praying..
Our mind can wander, all these things..
And so just those obstacles in and of themselves can really keep us from praying because we're.
like, "This is hard..
It's kind of difficult..
I don't know what I'm doing..
It seems confusing to me.".
And so we talked about how important it is just to understand that it's simple..
It's not necessarily easy, but it is simple..
Jesus teaches us that prayer is simply a personal conversation with God..
It's not some type of ritual..

$^{41}$It's not something that we have to do all these certain steps and get everything right..
It's not like something where we have to say the perfect words or say the right words or.
say them in the right order..
And if we don't say the words according to the formula or say them in the right way,.
then it's not going to work..
That's not prayer at all..
Prayer is simply having a conversation with God, sharing our hearts with God, pouring.
out the needs, the desires, the fears, the trouble, everything that we're dealing with..
We share that with God in conversation and that's what prayer is..
We talked last week about the personal nature of that conversation..
Jesus makes it clear at the beginning of this prayer that we're talking to our Father who.
loves us and who is close to us and who wants to be involved in our lives, who wants to.
hear us..
Our Father in heaven is someone that we have a close and personal and intimate relationship.
with..
Jesus made that really clear for us..
And then we spoke last week about how Jesus teaches us that our prayer should be focused.
on God..
And He begins this prayer by focusing our hearts on God, by talking about God first,.
by thinking about God first and how important that is to begin by focusing on God and who.
He is and His love for us as Father, but also His greatness and His goodness and His power.
and His holiness in our lives..
That's so important..
And then the other thing we discussed last week was just that as we begin to focus on.
God, it really informs our praying..
As we begin to think about God as our Father, as a loving God who cares about us, as we.
begin to think about God as powerful and great, but also as good..
And as we begin to think about God as holy and righteous and all these things, it really.
affects the way we pray..
It really changes the way we pray..
It really changes what we're going to ask God for..
And so it's really important for us to do that, to begin prayer from a place of focusing.
on God and who He is and letting the reality of who God is change the way we pray and change.
the requests that we have and the things that we're going to ask for..
So this week, we're going to continue on and look at the next verse here..
And we're going to continue to learn what Jesus has to teach us about prayer..
And I took advantage of some of the really nice weather this weekend..
I went on a little bike ride on the Lakefront Trail, so I've got a little footage of that.
for you..
But I just want us to really consider what Jesus is teaching us here today and try to.

$^{81}$really understand it and practice it..
I want us to think about how we can practice these things in our life and in our prayer.
life and how we can grow in our prayer life with God..
So we're going to begin today looking again in Matthew chapter six..
We're going to start at verse seven..
This is Jesus speaking..
He says, "And when you pray, do not heap up empty phrases as the Gentiles do, for they.
think that they will be heard for their many words..
Do not be like them, for your Father knows what you need before you ask Him..
Pray then like this, 'Our Father in heaven, hallowed be your name..
Your kingdom come, your will be done on earth as it is in heaven.'".
And I really want us to focus on verse 10 today and just one big truth that God is teaching.
us in this verse..
Jesus is really teaching us here that prayer is not about us notifying God..
It's more about God changing us..
And we sometimes think that prayer is kind of about us informing or notifying or letting.
God know what's going on in our lives..
Maybe you've prayed that way..
I've prayed that way a lot..
Maybe you've heard other people pray that way..
What we're doing in prayer is kind of telling God what's going on and telling God what we.
need and what we want and all those things..
And that is an important part of prayer..
I'm not discounting that in any way..
Like I said last week, God loves us..
He cares about our needs..
He wants us to express the needs of our heart to Him..
That's hugely important..
And in the model prayer here that we are talking about a few verses later, we'll get to that.
section where Jesus begins to teach us to express our needs to Him..
So that's important..
I'm not saying it's not about that..
But I'm saying that's not the primary purpose of prayer in our lives..
And sometimes we think it is..
Sometimes we think the primary purpose is that I can just tell God what's going on and.
what I want, what I need..
And that's not what prayer is about..
Jesus is showing us here at the beginning of this prayer, He gives a foundational and.
primary place to thinking about God's will and purpose in our life..
And that's something that God wants to do through prayer..

$^{121}$As we begin to practice it, as we begin to spend time with God in prayer, as we begin.
to use prayer in conjunction with God's Word to really think through our circumstances.
and think through what's going on in our life, God is going to use prayer to change us..
And that's really important for us to understand..
Prayer is not as much as we think about us informing God, and it's a lot more about God.
changing us..
And we see that in verse 8, Jesus says, "Hey, God already knows what you need before you.
ask Him.".
Well, if God already knows what we need before we ask Him, why does He even want us to ask.
Him?.
Why does He even want us to pray in the first place?.
God already knows everything..
He already knows what we're going through..
He knows the intricacies of our challenges..
He knows the complexity of our pain..
God knows all that..
So why are we supposed to tell Him?.
Why are we supposed to pour our heart out?.
Why are we supposed to let Him know what's going on?.
Why are we supposed to cry out to Him with our heart?.
Because as we do that, something will change..
And it won't be God..
It's not that we're going to change God with the things we say..
It is us..
God wants to use prayer to change us..
And that happens in so many ways..
First of all, prayer begins to change our relationship to God..
As we begin to practice it, as we begin to do it, as we begin to spend time in conversation.
with God and pouring our hearts out to Him, we get closer to Him..
It changes our relationship because we grow in that relationship..
We get closer..
That relationship gets stronger and deeper and tighter..
That's what happens..
Just like in a human relationship, when you have a friend and you begin to spend more.
and more time with that person and you guys share more and more conversation, and over.
time in that conversation you begin to talk about deeper issues of your heart and difficult.
things that you've gone through and past experiences and things that have troubled.
you or things that have been a challenge in your life, that changes that relationship..
It changes that friendship..
That friendship deepens and it gets stronger and you become closer to one another..

$^{161}$The same is true in our relationship to God..
As we begin to practice prayer and we begin to practice it more and more and we spend.
more time with God and we open our heart to God more and more and we begin to share more.
and more of what's going on in our life, it begins to change our relationship to God..
We get closer to Him..
That relationship gets stronger..
It deepens..
It gets more mature..
Prayer changes our relationship to God..
Prayer as our relationship begins to change, prayer changes our perspective on our circumstances..
See when we see a circumstance, especially a difficult circumstance in our life, when.
we see some difficult situation and we say, "This is a problem in my life..
It's hurting me..
It's a big deal and I don't know what to do about it.".
Maybe it's some type of issue at our job..
Maybe we lose a job or there's some relationship at our job with a boss or a co-worker that's.
really a problem..
Maybe it's something in school where school is not going how we want it to or we're afraid.
that we're going to experience some sort of failure or rejection in some way..
Maybe it's in a relationship with a friend or in a romantic relationship or maybe it's.
in some other area of our life..
But when we see a difficult circumstance, that can cause us to feel fear..
That can cause us to feel worry..
That can cause us to feel anxiety..
All different types of things that that can cause in our heart when we see that circumstance.
and focus on it..
But when we begin to go to God and focus more on Him, suddenly our perspective on that circumstance.
begins to change..
We begin to see God's love for us, His goodness, His care for us..
The fact that He is bigger and stronger than that circumstance, that situation, that obstacle..
The fact that He can get us through it..
The fact that no matter what happens, He will be with us even if things feel like they're.
not going well for us..
God's presence is in our life..
Suddenly our perspective on that situation begins to change..
Prayer changes us because it changes our perspective on our situation..
As we grow in that relationship with God, as God changes our relationship and we get.
closer to Him, then our perspective begins to change even more..
In that way, we can understand prayer as a relational conversation that changes our perspective..
And that's something that we can illustrate with a human relationship..

$^{201}$Maybe you've had a friend or a family member, someone that you knew well, and you had a.
good relationship with them..
And you guys had a different perspective on some issue..
You had a different idea about something..
And as you sat down and began to talk about that issue together, it changed your perspective.
or it changed their perspective..
Maybe you've experienced that..
Maybe you've experienced the fact that when you have a relationship of trust and empathy.
with someone and you begin to share your heart with them and they share their heart with.
you, that begins to change your perspective on the issues that you talk about or the issues.
that you're facing in life..
Maybe it causes you to say, "Okay, I see this in a different way that I didn't see before.
because I had this relational conversation with this other person.".
And that is something that happens in prayer..
As we begin to know God more closely and walk with Him in a deeper way, and as we begin.
to have these conversations with Him through prayer, God changes our perspective..
We start to feel differently about some of the issues in our life..
Some of the things that used to be really important to us are now less important..
And some of the things that used to be just kind of side issues that we didn't even really.
notice now become of primary importance in our life because we begin to value what God.
values and we begin to value less all those things that we in our selfishness value and.
that the world around us in its selfishness values..
That's what happens in prayer..
It is this relational conversation that changes our perspective..
And as that happens, prayer begins to change our desires..
And this is really, really important..
That prayer begins to change the things we want and the things that we desire in our.
life..
Let's look at Psalm 37..
This is a great psalm..
And it's actually got a lot of good truth for us right now in the moment that we're.
in with all the things going on and all the schemes and the plots and everything that.
happens that makes us frightened in our world..
This is really good for us to read..
This is the first, we'll read the first seven verses..
It says this, "Fret not yourself because of evildoers..
Be not envious of wrongdoers, for they will soon fade like the grass and wither like the.
green herb..
Trust in the Lord and do good..
Dwell in the land and befriend faithfulness..

$^{241}$Delight yourself in the Lord and He will give you the desires of your heart..
Commit your way to the Lord..
Trust in Him and He will act..
He will bring forth your righteousness as the light and your justice as the noonday..
Be still before the Lord and wait patiently for Him..
Fret not yourself over the one who prospers in His way, over the man who carries out evil.
devices.".
Like I said, it's a great passage for us in the world that we live in..
But I want us to think specifically about verse four, "Delight yourself in the Lord.
and He will give you the desires of your heart.".
There's something very interesting in this verse and something very interesting in our.
understanding of prayer as we consider this verse..
We all have desires, right?.
We all want stuff..
And for a lot of us, those human desires are we just want to be healthy and wealthy and.
comfortable and for everything to go great in our life..
We want it to be easy..
We want things to go well..
We want it to all work out and for our life to be comfortable and happy for us, right?.
That's what we want..
But this verse is telling us that we first need to delight ourselves in the Lord and.
we need to make what's important to Him important to us..
We need to find joy in God and in His truth..
And as we do that, He'll give us the desires of our heart..
Now what is this verse saying?.
Is it saying that if we delight ourselves in God and find joy in His truth, that then.
He'll just make us healthy, wealthy, and our life will be easy and comfortable?.
He'll give us everything we've ever wanted..
We'll be a millionaire and we'll get all the things we want from life..
I don't think that's what this verse is saying..
I think what this verse is saying is that as we delight ourselves in God and begin to.
rejoice in His truth, the desires of our heart will become aligned with God's desires for.
our life..
And we will end up seeing those things brought into our life..
We will see God's plan at work in our life..
And as our desires are now aligned with God's desires, we will get what we want from life..
We see this in the prayer that Jesus is praying here in Matthew 6 and verse 10..
Jesus says, "God, your kingdom come, your will be done on earth as it is in heaven.".
Jesus is saying, "God, I am not asking for you to give us what we want..
I'm asking for you to do what you want..

$^{281}$And I'm asking for us to become aligned with what you want.".
And we want to see that happen..
See as we begin to pray, prayer changes our desires..
Maybe we go into prayer thinking, "Okay, God, here's what I need in this situation..
I need my boss to change, or I need my teacher to change in this way, or I need this person,.
my friend, to change in this way, or I need you to give me this, or I need you to give.
me that, or I need you to fix this, or I need you to fix that in my life.".
Maybe we go into prayer that way..
But as we begin to spend time with God, focus on Him, and see more of who He is, and see.
more of how He works in our life, and as we think more about how He loves us and how He.
is working all things to good in our life, even the difficult things, and He uses even.
the challenges to grow us and to strengthen us and to bring beautiful things in our life,.
maybe those prayers begin to change because those desires begin to change..
And maybe we say, "Okay, God, you know what?.
Maybe you don't need to change my boss or my teacher right now..
Maybe I need to learn how to respond to them in a better way..
Even if they're wrong, even if they're doing something that's difficult for me and it's.
not my fault, maybe I can grow in this situation..
God, maybe what's going on in my life is not about me just needing to get something else.
and me needing you to give me something else..
Maybe it's me learning how to be content with what I have, or maybe it's me learning how.
to give to someone else, or maybe it's me learning how to serve and using what you've.
done in my life to serve others.".
See as we begin to pray, it begins to change the desires of our hearts, and as those desires.
of our heart change and begin to align with God's desires for our heart, suddenly we'll.
see more and more of those desires getting met..
That's what Psalms 37, verse 4 is saying..
If we will delight ourselves in God, if we will find joy in His truth for our life, then.
the desires of our heart will suddenly move to a place where we are wanting for our life.
what God wants for our life..
And as those desires change, then we'll start to see them come to pass..
Instead of just wanting to be rich and comfortable and have everything go great in our life,.
we'll say, "God, what do you want for my life?.
How do you want to grow me and change me and use even the difficult, hard circumstances.
in my life for your glory and for my good?".
As we begin to desire that, that desire will come to pass..
We'll say, "Yeah, I can see this is actually happening..
God is using these things in my life to glorify Him and to benefit me, and I can see it, and.
now the desires of my heart are coming to me.".
Prayer changes our desires, and that's so important..

$^{321}$Prayer changes what we ask for..
As our desires change, we start to ask for different things..
Instead of just saying, "God, I just want everything to be great in my life," we start.
to say, "God, I actually want to know you more, and whatever it takes for you to grow.
me and take me to a place of knowing you more, that's what I want..
Even if it's not easy, even if it doesn't feel great, that's what I'm asking for.".
Instead of saying, "God, can you just do this in my life because it would be great for me,".
we start to say, "God, can you do something in my life that would allow me to serve and.
love someone else, even if it's not great for me?".
Instead of saying, "God, can you make everything work out where I get what I want in life,".
we start to say, "God, can you help me to connect someone else with you, and can you.
help me to care for others?.
Can you help me to give more and to love more and to serve more?.
Can you help me to be more focused on what's important to you in my life instead of just.
what my own selfish desires are?".
We start to ask for completely different things, and that's good..
When I was a young man, I was a follower of Jesus as a teenager, but the things I asked.
for back then are not the things I ask for today..
The things I asked for in my life back then, some of them were good, but a lot of them.
were pretty selfishly motivated and selfishly oriented..
I'm glad that those things have changed, that I ask for different things today..
There's still a lot of things that I ask for right now in my life that are selfishly oriented.
and selfishly motivated, and I'm sure God will change more of those as I continue to.
grow and walk with Him..
But it's good that God has changed the desires of my heart..
It's good that God has changed what I ask for in prayer..
As you spend time with God and your relationship to Him changes, your perspective changes,.
your desires change, it will change what you ask for..
You'll start asking for some different things..
The way you pray when people are sick, the way you pray when things are difficult in.
your life, the way you pray when things surprise you because they don't go your way will change..
You will ask for different things..
And as you begin to change and ask for different things, your choices will change as well..
Prayer changes our choices..
We ask for different things and we begin to say, "God, I'm going to choose something.
different..
I'm not just going to choose the thing that maximizes my own pleasure and comfort and.
selfishness and wealth..
I'm going to choose the thing that maximizes your glory in my life..
I'm going to choose to do what would honor you in my life..

$^{361}$I'm going to choose to do what's best for the kingdom of God in my life.".
It changes our choices..
And as it changes our choices, prayer begins to change our outcome..
It begins to change the direction of our life..
It begins to change where we're headed..
It begins to change the choices that we're making and it begins to change who we are..
Who we are as a person begins to fundamentally change..
And what we're doing with our time and with our resources and with our energy, with our.
focus begins to fundamentally change..
Prayer changes who we are..
We get to a place where we can say genuinely what Jesus said in verse 10, "God, your kingdom.
come, your will be done on earth as it is in heaven.".
That's a hard thing to say..
It's hard to say, "God, I'm going to surrender my will to you and I'm going to say, 'God,.
this is what I want or this is what I desire,' but I'm saying, 'Your will be done, your kingdom.
come on earth as it is in heaven.'".
It's an easier thing to say in those moments when you don't care, when you say, "Okay,.
whatever happens in this situation, I don't really care.".
It's not that difficult..
But in the situations where you really want something to happen, where you're really praying.
for something, where someone's sick or where there's something going on in your life that's.
really important to you and you're saying, "God, can you please do this?.
This is what I want..
This is what I need..
This is what I want to see happen.".
It's difficult to get to that place where you say, "God, this is what I want..
This is what I want to see happen, but I'm going to pray that your will is done and that.
your kingdom comes on earth as it is in heaven.".
That's a difficult prayer to pray, but as we begin to spend time in prayer, as we begin.
to practice prayer, we begin to let God change our relationship to Him, change our perspective,.
change the way we look at our circumstances, change our desires, change what we ask for,.
change our choices..
That's what we'll start to see..
We'll find ourselves in that place where we can say, "God, this is the desire of my heart..
Yes, I'm going to pour it out to you..
I'm going to be honest with what's going on..
I'm going to be transparent about my feelings, and I'm going to let you know all of what's.
going on, but I'm also going to come to this place where I'm saying, 'God, may your will.
be done and not mine..
I'm going to come to a place of surrender and of allowing you to work in my life.'".

$^{401}$And that's what's huge about prayer..
Jesus points us to this right up front in the Mahalo prayer, that prayer is less about.
us telling God stuff and more about God changing our heart..
The telling God stuff is important, but what's even more important is that God wants to change.
who we are and how we relate to Him, how we relate to the world around us, and God does.
that through prayer and through this practice of spending time with Him..
So I want us to begin to practice this this week..
I want us to intentionally say, "You know what, we are going to, as we pray, ask God.
to change our relationship to Him, to change our perspective, to change our perspective.
on our circumstances, to change our desires, to change the things that we're asking for,.
to change some of the choices we make, to get us to that point where we can say, 'God,.
your kingdom come, your will be done on earth as it is in heaven.'".
And Psalm 37, verse 4 gives us a good idea of how that happens..
It happens when we start to say, "God, I'm going to delight myself in you..
I'm going to rejoice in your truth..
I'm going to focus on your truth..
I'm going to focus on knowing your truth and understanding what's important to you, and.
as I do, I'm going to shift my desires to align with your desires.".
So let's practice that this week..
Let's focus on knowing God's truth and on knowing what's important to God in our world.
and in our life and what He wants for us..
We've already talked so much about these values of the kingdom..
That's what God wants for us in our lives..
Do we have those values in our life?.
Are our desires in line with those values of the kingdom?.
Let's think about that, and as we do, let's begin to let God change our perspective and.
our desires so that we can become more and more of who Jesus wants us to be, and we can.
move more and more to that place of saying, "God, may your kingdom come and your will.
be done on earth as it is in heaven.".
Let's pray..
God, I do ask that you would change us through prayer, God..
I ask that you would help us to first of all just practice prayer..
Just do it..
And as we do it, that our focus would be on you, and as we begin to focus on you and on.
your truth and on who you are and your love for us, that we would begin to change in our.
perspective..
That we would begin to change in our desires, God..
I pray that you would change the desires of our heart to align with your desires..
And as that happens, God, I pray that we would change the things we ask for, that we would.
begin to ask for things that glorify you instead of things that just satisfy ourselves..

$^{441}$And God, I pray that as we do that, we would see your kingdom come and your will being.
done on earth as it is in heaven, God..
Your will is perfect..
Your will is good for all of us in a way that our own will is not, even when we imagine.
that we know what would be best in a situation..
We don't, but you do..
So God, we pray that your will would be done..
We pray that your kingdom would come on earth..
We know it's already that way in heaven, that your will and your power and purpose is already.
fully manifested in heaven..
And so we pray that through us, your people, it would be manifested on earth as well..
In Jesus name, amen..
Amen..
[BLANK AUDIO].
\newpage



\section{約拿書 1:7-16}
\label{sec:MzNtNdJRGgQ}
\textbf{經緯線 [釋經視頻] EP3 約拿書 1\_7-16}
\newline
\newline
連結: \href{https://youtube.com/watch?v=MzNtNdJRGgQ}{\texttt{https://youtube.com/watch?v=MzNtNdJRGgQ}} ~~~~ 語音日期: 2020-11-14
\newline
\newline
\hyperref[sec:rfX8QtfxB9g]{\small{< < < PREV SERMON < < <}}
~
\hyperref[sec:index_chronic]{\small{[返順時目]}}
~
\hyperref[sec:index_scriptual]{\small{[返順卷目]}}
~
\hyperref[sec:_ZlpJCFpNOg]{\small{> > > NEXT SERMON > > >}}
\newline
\newline
約拿書 1:7-16
\newline
\begin{longtable}{cl}
\hline
\hline
章節 & 經文 (和合本修訂版)\\
\hline
1:7 & \begin{tabularx}{0.7\textwidth}{X} 船上的人彼此說：「來吧，我們來抽籤，看看這災難臨到我們是因誰的緣故。」於是他們就抽籤，抽出約拿來。 \end{tabularx} \\ \\ \relax
1:8 & \begin{tabularx}{0.7\textwidth}{X} 他們對約拿說：「請你告訴我們，這災難臨到我們是因誰的緣故呢？你做甚麼行業？你從哪裡來？你是哪一國的人？屬哪一族？」 \end{tabularx} \\ \\ \relax
1:9 & \begin{tabularx}{0.7\textwidth}{X} 他說：「我是希伯來人，我敬畏耶和華，天上的神，他創造了滄海和陸地。」 \end{tabularx} \\ \\ \relax
1:10 & \begin{tabularx}{0.7\textwidth}{X} 那些人就大大懼怕，對他說：「你做的是甚麼事呢？」原來他們已經知道他在躲避耶和華，因為他告訴了他們。 \end{tabularx} \\ \\ \relax
1:11 & \begin{tabularx}{0.7\textwidth}{X} 海浪越來越洶湧，他們就問他說：「我們當向你做甚麼，才能使海浪平靜呢？」 \end{tabularx} \\ \\ \relax
1:12 & \begin{tabularx}{0.7\textwidth}{X} 他對他們說：「你們把我抬起來，拋進海裡，海就會平靜了；我知道你們遭遇這大風浪是因我的緣故。」 \end{tabularx} \\ \\ \relax
1:13 & \begin{tabularx}{0.7\textwidth}{X} 然而那些人竭力划槳，想要把船靠回陸地，卻是不能；因風浪越來越大，撲向他們。 \end{tabularx} \\ \\ \relax
1:14 & \begin{tabularx}{0.7\textwidth}{X} 於是他們求告耶和華說：「耶和華啊，求求你不要因這人的性命使我們滅亡，不要使流無辜人血的罪歸給我們；因為你－耶和華隨自己的旨意行事。」 \end{tabularx} \\ \\ \relax
1:15 & \begin{tabularx}{0.7\textwidth}{X} 他們把約拿抬起來，拋進海裡，海的狂浪就平息了。 \end{tabularx} \\ \\ \relax
1:16 & \begin{tabularx}{0.7\textwidth}{X} 那些人就大大懼怕耶和華，向耶和華獻祭許願。 \end{tabularx} \\ \\
[1ex]
\hline
\hline
\end{longtable}
$^{1}$又回到經緯線了.
上一集我們講完第一章的一至六節.
看到約拿仙子上帝呼召她來不來.
往北走 誰知道她一下子就屈了一尾十.
就從南邊去了西邊 去了他斯.
在大風大浪之下 最終上帝都不肯求.
最後水手門抽籤的時候 一抽就抽到她.
就說她是有問題那個.
但是上一次停在那裡.
我心裡面有些覺得很怪的就是.
抽籤搞什麼?不知道大家怎麼看.
我覺得你覺得怪.
因為你又好像我們那樣看周星馳大.
我覺得整個戲劇就是好像大家抽籤.
就是一幫人在抽籤.
其實支支籤都是一樣的.
他們抽了之後就撕了一塊.
唯獨約拿那個她是純良像甲子.
她又沒有撕.
OK OK 我這樣就不是了.
這個不知道周星馳再早一點的黃禎.
都是差不多的.
歷史是不是這樣?.
就是那些諷刺文學.
就是有些驚喜 有些很突然的差異的地方.
如果我回來就是於是他們就制籤.
居然制出了一個難題.
很有趣.
Mark Sir你覺得如果從考古和歷史那一套.
其實我這個純粹是開玩笑.
但是你從考古裡面你會看到些什麼呢?.
那個宗教習俗.
以前那些外邦人,近東人.
每逢有大能臨到的時候.
他們一定是覺得得罪了那個神.
所以要找出那個罪.
到底誰是犯錯.
所以他們很自然的就會是.
這個不差異的.
那些水手就在那裡抽籤.

$^{41}$到底誰做錯事.
誰沒有初一,十五,沒有上香等等.
他們就要抽籤.
找出誰是誰是壞人.
有趣 都挺有趣的.
這個故事就開到了.
去到這裡抽到的時候.
約拿就和這幫水手的對話.
在對話當中不知道大家發現了些什麼呢?.
那個第八節都挺生動的.
五個問題是吧?我沒有看錯.
他們一味問 一味問 一味問.
你會明白那個意思.
就是不斷地質問.
你從哪裡來的?在哪個族?等等.
就是那個意思.
我再讀一段經文.
那段經文是這樣說的.
請你告訴我們.
這災臨到我們是誰的緣故.
你以何事為業.
你從哪裡來.
你是哪一國哪一族的民.
你可以當四個問題.
不過四至五個都沒有所謂.
但是我同意好像被人.
好像被老鏈再審他一樣.
不知道大家看的時候.
再看下去的時候.
你從約拿的回應.
你發現了些什麼呢?.
我就發現了約拿.
無論是多少條問題都好.
他都只是回答.
他說他是希伯來人.
也是一個敬畏神的人.
還有他所敬畏的這位神.
是滿有能力.
能夠創造海洋和陸地的一個神.
我就留意到.

$^{81}$如果大家認識一些原文.
其實原文的命令就是.
他說希伯來人.
我是.
耶和華創造天地的神.
我敬畏的.
就是他的ethnic background.
他是走之的.
然後是神.
這個創造天地的神.
我敬畏的.
他有兩個問題是避開的.
他沒有回答那些人.
他說是因為我的緣故.
所以災難來到你那裡.
還有他沒有說過他做了什麼.
那就說回.
剛才Marx說.
原文就是說.
我是希伯來人.
希伯來人就是我.
其實他從頭到尾都沒有認自己是壞人.
他也沒有認自己.
去到雀草是一個不好的事情.
我想會不會當時約拿也很維護.
No matter what.
他跟神當時的關係怎樣也好.
但其實他對自己的身份.
他很尊重.
也對上帝很敬畏.
所以他無論怎樣也好.
他不想在任何人面前.
尤其是不同自己.
同種族的人.
在他們面前有辱自己上帝的名呢?.
這個可能性也有的.
還有剛才大家說好的.
在希伯來人.
剛才說他擺的第一個就是.
我是希伯來人.

$^{121}$他將民族身份擺得很高的.
沒錯.
希伯來人我也是.
如果用我們現在這樣說.
有點像某個國家的小粉紅.
他很維護那個身份.
很強調那件事.
或者白人上帝.
也是.
那時候就是希伯來人.
你也分得出其實他的身份.
一個人的身份是先說的.
例如你是基督教.
你是中國人.
還是加拿大人.
所以這三樣看他怎樣擺的時候.
你也會看到他自己的身份.
這是一個好點.
因為他會不會有一個伏筆.
在這裡就是他的ethnic.
就是他的種族的identity.
高過他的theological identity.
他的神學的identity.
這個也值得我們考究.
各位觀眾可以留意.
你可以記下這個問題.
這個是好問題.
究竟是不是那個種族的身份.
比他對上帝的身份高.
現在暫時沒有答案.
但起碼是一個好的observation.
還有我覺得他也挺大派的.
一出牌就說.
我的神就是滿有能力的.
創造海洋和陸地.
其他的神.
因為大家近東文化.
我記得所有的神都是.
一個神是管一件事.
或者管幾件事.

$^{161}$就是管生仔.
管農作.
管海.
跟我們中國人一樣.
做君管做頭.
每人管一範.
他現在就說我什麼都管.
其實是一張大牌來的.
那些水手也acknowledge的.
在上一節他直接說.
這麼厲害你搞定他吧.
這個就是他搞笑的地方.
突然間就流淚.
他們就大大驚怕.
說他死火了.
受死犬的氣.
他問了一個問題.
他繼續問你做了什麼事.
你搞什麼鬼.
但他又不回答.
還有一個問題.
就是你以何事為業.
他沒有回答過.
他沒有說過.
第十節B.
那個narrator龐白就說.
因為他已告訴他們了.
代看故事的就知道.
原來他們已經知道了.
他們說你搞了什麼事.
搞什麼鬼.
說到這裡的時候.
我們就開始發現.
其實如果約拿不肯說.
或者不肯回答.
silence不肯回答.
究竟我以何事為業的時候.
其實都有一個反應.
他都繼續在.
就像我們上一集所說.

$^{201}$可能他對這個呼召.
是可以說.
仍然是沒有辦法接受的.
他這樣說是吃不下的.
不知道大家認不認同.
這個吃不下的.
這份工作他認為打不過.
打不過.
繼續逃避你吧.
他要辭職.
似乎這個辭職不太行.
因為第十一到第十六節的時候.
很嚴重的.
因為大家在第十一到第十六節的時候.
如果大家再去留意的時候.
經文那裡.
船員就問他們.
我們當向你怎樣行使海浪平靜.
這話是因海浪越發翻騰.
就是說說完了.
判判完了.
越判越厲害.
似乎到了這個時候.
當水手這樣問的時候.
約南是怎樣回應的呢?.
他都沒有認衰仔.
他只不過是.
算了你扔了我.
應該沒事了.
雖然這件事是因我而起.
不單止他沒有認衰仔.
反而作為一個讀者來說.
我會看到他都挺有義氣.
挺良善.
他的名字是甲子.
順良如甲子.
他又做回他自己的名字的行為.
他寧願自己一個人犧牲.
可以救到整船人.
其實如果Benjamin這樣說.

$^{241}$他其實那個.
正正回答了之前水手問的.
就是這災難到.
我們是因誰而起.
最後他肯承認了.
扔我下去.
扔我下去就搞定.
如果他英雄一點的話.
結束這個故事.
他應該自己跳下去.
不用叫人扔他下去.
他因我而過.
其實也是.
他這樣的緣故.
濾動那班水手好像做了Cure.
逼那班水手.
逼那班水手做壞人.
後來他自殺.
寸於他不疑.
從這個角度看.
他真的不是很良善.
後來那些水手.
他們也大頭盔.
他們也說.
我扔不是我要扔.
他叫我扔他不關我的事.
又看到那班水手這樣.
你看13節.
就是那個對比.
現時那些水手.
是竭力去撐船.
是不可能.
他也不想扔下去.
是.
大家說得很有趣.
又是一個irony.
就是說.
那個水手就說.
你做我吧.
但他就說不可以.

$^{281}$不可以這樣.
你是和華先知.
你又說他有神經病.
扔下去更糟糕.
所以為什麼他不自己跳下去.
好問題.
這個問題.
值得反省.
不止他不立即扔下去.
14節說什麼.
他們變求高野和.
是.
看那班水手的反應.
大家求.
有兩件事我發現.
第一.
始終沒有回應.
水手的一條問題.
三條答了.
有一條始終沒有答.
你以何事為業.
他死口不肯認自己.
他不肯說.
但是.
還有一件事.
我發現.
他就是.
水手的反應都不同了.
就是由初初.
求告自己的神.
現在變成.
求告上帝.
那個落差很大.
是.
還有他也是.
我覺得.
我的領袖.
因為.
因為Jonah.
叫他們這樣做.

$^{321}$而當他們這樣做的時候.
風浪真的停了.
於是他們就.
有理由去拜野和.
這個也是.
一個.
差不多像.
悔改一樣的.
事.
還有這個挺諷刺的.
大家有沒有想過.
一個在理性上.
才知道認識上帝.
這麼深的.
一個才知道.
但是他的回應.
和一群.
知道的很少.
他所知道的可能只是說了幾句.
就信服做了.
而得到的那種熱門受上帝拯救.
兩個的落差很大.
還有.
我想那個經文的.
作者也有那個字眼.
在那裡.
大大敬畏.
那些水手大大敬畏耶和華.
剛才你記不記得他被迫.
約拿爾被迫回答的時候.
耶和華我敬畏.
你就問.
到底誰敬畏耶和華.
非常好.
大大敬畏.
剛才他說.
希伯來人我是也.
大大敬畏.
所以如果你這樣說.
我覺得會有一點.

$^{361}$預測.
去預測將來整個故事.
是怎樣發展的.
都挺有趣.
其實像不像.
一個伏筆在那裡.
我嘗試過代入.
我是一個猶太人.
去看的時候.
一個他們認為這麼好.
這麼信仰純正的先知.
但是.
他竟然逃避上帝的呼召.
他寧願被人.
扔下海.
他也不去.
但這班水手當他真的.
迫到沒辦法.
他真的做了.
這班水手竟然因為.
他下水的緣故.
他們悔改.
很搞笑.
好像是不是有一點伏筆.
究竟將來.
會發生的事情.
我們這樣看到尾.
我們發現有一點很有趣.
因為在我的聖經.
一章十七節.
黏在了二章一節.
成為一個大段落.
一章十七節至二章十節.
我不知道大家的.
聖經本是怎樣.
看到這裡的時候.
會不會有一個問題.
為什麼不是將一章十七節撥上去.
一章十六節黏在一起.
這個問題這麼聰明.

$^{401}$其實很多時候.
很多時候在聖經.
它的分段.
不是文學本身的分段.
所以你說.
十七節和十六節.
如果我去分的話.
我會將十七節.
放在第二章.
第二章的開始.
就不應該有十七節了.
明顯你的故事的完結.
就應該是撥上去.
好啊.
我們不如這樣.
我們就停在這裡.
我們下一次的研究.
我們才會再看.
由一章十七節.
我們去到二章的第十節.
這個就是我們下一次.
如果再釋經的時候.
我們會看這個部分.
不過到最後.
這段經文.
對於我們今天的應用.
有什麼意義呢?.
我想第一件事我想說的.
就是我們第一集開始.
我們說要從文體的genre.
體材出發.
我知道有很多傳統的信徒.
不是很習慣.
我們這幾集所說的.
就是有些稀少老馬.
這麼搞笑的東西.
因為是神的話語.
聖言.
如果不用這個方法去看.
你就會錯過了很多.

$^{441}$作者原意.
他希望刻畫.
或者想諷刺的對象.
不要覺得勝經.
就是一定很正經.
其實有時候.
就是用一些.
我們想不到的文學手法.
諷刺.
來表達一些.
很重要的點.
好啊.
如果這樣說.
大家用Marx的角度去想.
大家會想到.
這十六節經文.
有什麼點子.
值得留意呢?.
其實說到諷刺.
很多時候都是用一個對比.
剛才很明顯的對比.
就是若拿和瑞秀的對比.
若拿作為一個先知.
他是不是應該要行公義好憐憫.
全謙卑的心.
與神同行.
但他不是.
相反那些水手.
他更加有那個compassion.
他更加願意.
行公義和好憐憫.
其實另一個我想的東西.
就是若拿.
他已經是先知.
他知很多.
他經歷很多.
但他心裡面是不斷抗拒上帝.
甚至乎在船裡面.
他那種.
用了一個不好聽的發爛咋的情形.

$^{481}$就是在發爛咋.
就是完全發爛咋的.
但相反那些水手.
他知說的就像別人的屬靈四定律.
四句就這麼多.
但是他按他所知的.
去做出他所做的.
我講這件事.
對於那些代替論.
約拿書的經文來說.
我好深看到.
這個新聞還未講完.
知道他正在逃避這個使命.
即是說其實知多少.
和我們承擔使命.
不一定是正比.
因為承擔使命是一個信服.
其實講的是信服的問題.
我覺得這件事.
在經文可以看到.
一班外邦的水手.
我們可以說是屬靈的.
蒙恩的邊緣人.
他知不多.
但他願意信服去做.
他仍然抗拒.
這是我直接看到的圖畫.
當中的領袖.
我自己就給了一個很大的提醒.
就是一個對比.
亦都是一個對比的手法.
若拿和這班水手.
其實若拿和上帝.
很親密.
認識耶和華神.
他自出母胎.
開始已經是.
但相對這班水手.
是外邦人.
其實是未認識耶和華神是誰.

$^{521}$但當他們認識了之後.
他們的反應.
就好像現在.
我們在教會看到的.
一些誤導朋友.
變成了初信的朋友.
他們接受耶穌基督.
然後就很投入教會.
很追求.
那種態度.
相對之下.
教會有些.
像我們這些.
可能老油條.
信了很久很久.
但可能已經失去了初心.
或者那種服事神的熱誠.
很多時候教會.
要你去服事.
明知是神的呼召.
叫你去服事.
但可能就會逃避或許很多藉口.
但相反一些初信的.
他們的初心.
因為仍然是很初心.
很熱誠.
反而他們很投入教會的服事.
很熱心.
有這種的看法.
在裡面.
我做一個小小的總結.
在這裡我們看到.
約拿書的先知.
在第一章我們看到.
他一直逃避上帝.
叫他所做的事.
就算周遭的人給很多問題.
他都逃避.
這個正正是我們要問.
為何約拿逃避呢?.

$^{561}$可能我們暫時沒有答案.
這個問題可能值得我們去想.
甚至乎可以這樣說.
很奇怪.
或者真是諷刺的.
一個逃避的先知竟然帶了一班外邦人.
悔改.
不過是一個意想不到的方法.
想都沒有想到.
我想這個也是一個伏筆.
我覺得是一個伏筆.
我們看到這裡的時候.
我們可能都會問這個問題.
究竟我是否正在逃避一些東西呢?.
究竟我像不像約拿.
我有時會沒有留意到.
其實我們身邊有很多初信甚至無道的人.
其實他們是蒙恩的邊緣人.
其實從他們的行為.
從他們的簡單的信心.
我們看到一個真信.
那種的真是美呢?.
這個值得我們去反省.
各位.
我們下一次.
我們就開始進入.
一章十七節.
到二章十節.
又鼓勵大家了.
在我們下一次的經緯線出之前.
好好再看那段經文.
容許我們.
下一次在經緯線當中.
再尋索這段經文.
對昔日.
和對今天我們的意義.
下次再見.
拜拜.
\newpage



\section{約拿書 1:17-2:10}
\label{sec:_ZlpJCFpNOg}
\textbf{經緯線 [釋經視頻] EP4 約拿書 1\_17-2\_10}
\newline
\newline
連結: \href{https://youtube.com/watch?v=-ZlpJCFpNOg}{\texttt{https://youtube.com/watch?v=-ZlpJCFpNOg}} ~~~~ 語音日期: 2020-11-14
\newline
\newline
\hyperref[sec:MzNtNdJRGgQ]{\small{< < < PREV SERMON < < <}}
~
\hyperref[sec:index_chronic]{\small{[返順時目]}}
~
\hyperref[sec:index_scriptual]{\small{[返順卷目]}}
~
\hyperref[sec:ksA3W2OSofQ]{\small{> > > NEXT SERMON > > >}}
\newline
\newline
約拿書 1:17-2:10
\newline
\begin{longtable}{cl}
\hline
\hline
章節 & 經文 (和合本修訂版)\\
\hline
1:17 & \begin{tabularx}{0.7\textwidth}{X} 耶和華安排一條大魚吞下約拿，約拿在魚腹中三日三夜。 \end{tabularx} \\ \\ \relax
2:1 & \begin{tabularx}{0.7\textwidth}{X} 約拿在魚腹中向耶和華－他的神禱告， \end{tabularx} \\ \\ \relax
2:2 & \begin{tabularx}{0.7\textwidth}{X} 說：「我在患難中求告耶和華，他就應允我；我從陰間的深處呼求，你就俯聽我的聲音。 \end{tabularx} \\ \\ \relax
2:3 & \begin{tabularx}{0.7\textwidth}{X} 你將我投下深淵，直到海心；大水環繞我，你的波浪洪濤漫過我身。 \end{tabularx} \\ \\ \relax
2:4 & \begin{tabularx}{0.7\textwidth}{X} 我說：『我從你眼前被驅逐，然而我仍要仰望你的聖殿。』 \end{tabularx} \\ \\ \relax
2:5 & \begin{tabularx}{0.7\textwidth}{X} 眾水環繞我，幾乎淹沒我；深淵圍住我；海草纏繞我的頭。 \end{tabularx} \\ \\ \relax
2:6 & \begin{tabularx}{0.7\textwidth}{X} 我下沉到山的根基，地的門閂將我永遠關住。耶和華－我的神啊，你卻將我的性命從地府裡救出來。 \end{tabularx} \\ \\ \relax
2:7 & \begin{tabularx}{0.7\textwidth}{X} 我心靈發昏時，就想起耶和華。我的禱告進入你的聖殿，達到你面前。 \end{tabularx} \\ \\ \relax
2:8 & \begin{tabularx}{0.7\textwidth}{X} 那信奉虛無神明的人，丟棄自己的慈愛； \end{tabularx} \\ \\ \relax
2:9 & \begin{tabularx}{0.7\textwidth}{X} 但我要以感謝的聲音向你獻祭。我所許的願，我必償還。救恩出於耶和華。」 \end{tabularx} \\ \\ \relax
2:10 & \begin{tabularx}{0.7\textwidth}{X} 耶和華吩咐那魚，魚就把約拿吐在陸地上。 \end{tabularx} \\ \\
[1ex]
\hline
\hline
\end{longtable}
$^{1}$各位好,經衛線又回來了.
上一次我們看完約拿書的第一章.
現在我們會看約拿書的第一章十七節到二章十節.
不過說之前,Benny你後面那一章是什麼?.
這一章是感恩節火雞.
因為最近天氣很冷,臨近美國感恩節.
我想營造一下人間溫暖.
另外一件事就是我們今天查的經文.
當中的二章二至九節.
其實是一個感恩的詩歌.
感恩詩歌?你可不可以多說一點.
為什麼他會唱感恩詩歌?.
因為我有做功課,我去找過一些解經書.
他都在參考到感恩詩歌裡的一些詞彙或情節.
都是在當時耳熟能詳的詩篇裡找到的.
譬如詩篇第42篇,第69篇和第88篇.
都是說到大衛被追殺,然後耶和華拯救了他.
於是他就去感恩,之類的.
這些還是希伯來人很出名的感恩詩歌.
感恩節就是一定會唱的歌.
還有很多是義人受苦的時候.
他們會有這樣的祈禱.
義人受苦得救的感恩讚美.
好像我們現在這樣說.
如果我們幾個坐在一起.
突然悶悶地想唱感恩詩的時候.
我們會唱我們的年代的詩集的歌.
是的,差不多.
我不同年代.
你不同年代?.
差不遠而已.
我都有看這首詩.
這首詩很有趣.
以一個講故事來說.
它有很多文字的工具.
用了很多詞彙,字.
勾起了很多思想.
特別是山根,深淵.
很多重複類似的字出現了.
不知道大家覺得作者為什麼會用這麼多字.

$^{41}$他想說些什麼呢?.
我們看希伯來人的詩歌時.
經常要留意那些平衡句.
因為A句和B句是平衡的.
而下一句就是擴張第一句所說的東西.
譬如第二章一節.
他說「遇見患難」.
一節的B就不單單是患難.
說到陰間.
希伯來人說.
陰間的深淵,即是死人才會去的地方.
第三節A和B.
深淵,深處,不單是深的谷.
也是海的深處.
海,大水.
一直擴展下去.
很多混亂的東西.
大水,深水,混沌等等.
我也覺得是.
因為有一節是很深刻的.
就是第三節和第五節.
我嘗試用講故事的感覺.
說出來,我閉上眼.
我想到一個圖畫.
大家可以試試.
大水環繞我,波浪,熊濤都慢過我身.
珠水環繞我,幾乎已沒我.
深淵圍住我,海草纏住我的球.
我嘗試閉上眼.
當時的讀者是沒有讀的,是聽的.
聽出來的時候也是充滿著凶險.
很危難的情況.
很折墮的情況.
這首詩歌,既然大家說過是感恩的詩歌.
我想藥娜其實是感恩.
要獲神對她的拯救.
特別是在死亡陰間拯救她出來.
這可以說是感恩詩的特色.
當你看的時候.
尤其是在第六節.

$^{81}$她說得很清楚.
我下到山根,地的門將我永遠關住.
講得白一點,就是她快死的時候.
就是希伯來人的underworld.
陰間的情況.
她的轉接變得很快.
耶和我的神,你卻將我的生命從坑中救出來.
這首詩歌裡面,可以看到藥娜的信心.
特別留意到第七節,第九節.
我們再細心留意的時候.
某程度上都是藥娜自己在這個屬靈經歷下的感想.
不知道大家看到甚麼呢?.
她是一個很威風的人.
她在經文裡面說.
那些不信的人.
她是一個很敬虔的人.
她進到聖殿裡敬拜.
但那些不信的人,就不會進到聖殿裡敬拜.
言下之意,她就是想說.
我是一個這麼敬虔的人.
雖然我現在都這麼折墮.
但我就是Bet,耶和華一定會來救我.
她就是這個意思.
其實都有這個情況.
因為我留意到,藥娜肯定了.
她在發恩的時候,祈禱都去到上帝的面前.
很厲害的,她很穩定的祈禱,上帝都聽到.
還有,你也說得對.
信奉虛無之神的人.
我們這樣說,非以色列人,就是外邦人.
他們一定會離棄,連愛他們的主.
其實那個翻譯不是連愛,是離棄上帝.
而相反,藥娜就不會離棄.
而且她又要獻祭,又要還願.
她真的很敬虔.
你們每個人都會放棄,但我不會放棄.
最後一句,我很喜歡.
救恩出於耶和華,多爽啊.
一鎚定音,多棒.
我想熟悉希伯來斯歌的人.

$^{121}$就會聽到這首歌是一首串燒的感恩歌.
是一首很美的感恩歌.
其實真的很美.
感恩字,其實只說這段都很棒.
不過我聽著聽著覺得怪怪的.
二至九節是很美,但總覺得怪怪的情況.
不知道大家有沒有這樣的感受呢?.
是怪怪的.
因為你沒理由被人追斬追九條街的時候.
你去唱感恩詩.
就算你看回剛才的詩篇.
42篇,69篇,88篇.
唱這首歌的人,即是大衛.
他已經脫了險才會唱.
他不會在危難中唱這首歌.
危難時他會喊救命,但他不會感恩.
因為那不是感恩的時候.
奇怪,譬如那首詩是蓋式.
好像Benny所說,已經過了難.
這是其一,Past tense.
第二,非常格格不入.
作者是把這個禱告放在餘福裡面做的.
不是在餘福裡面被套出來,去到寒地才做.
我也同意,雖然有些人跟我們說.
他在餘福裡面被拯救了,但其實在餘福裡面還是危險的.
deliverance還沒完,還在危險當中.
應該還會有些求救的聲音.
所以是一個很奇怪的.
這是一個很不現實的寫照,一種描述.
怎會在餘福裡面做到這樣的祈禱.
回到第一集,我們經常想的問題.
華人讀經經常想,有沒有可能呢?.
而不現實呢?.
有什麼理由在三日三夜,在魚裡面呢?.
不止在魚裡面,還在魚裡面作出這樣的串燒.
這麼美的感恩詩歌.
沒理由在街市唱「我心尊主為大」.
在街市也不可以.
剛才你那個比喻多好,被人追著狗條街的時候.
沒理由你可以唱到這麼美的歌.

$^{161}$挺有趣的.
黃師傅,你覺得在這首詩,有沒有地方覺得很怪?.
我想這首詩固然是一個感恩詩.
如果純粹感恩就怪了.
但剛才也聽到,處境,場景,時間.
說著說著就有很多問題出來了.
不得不去質疑,究竟這首詩是否真的感恩的詩呢?.
還有留意第八節的時候,可能會找到答案.
因為在第八節的時候,我們看到一個很大的對比.
他好像若拿已經覺得自己是一個很淑齡的.
他淑齡到一個地步,覺得自己沒有問題.
但相對地那些外邦人覺得他們是不會去敬拜神.
他們會離開上帝.
但問題是,第一章誰是真正敬拜上帝?.
大家記得吧,是外邦人.
第一章是誰是逃跑先知?.
是若拿.
是他離棄神.
所以如果在那個情況下他下了餘福.
他真的要祈禱,應該是一個認罪的祈禱.
他沒有的.
他整首詩,我不是否定他的讚美.
但他最要做的事情是沒有了.
究竟他為甚麼沒有認罪呢?.
他是不是覺得沒有問題呢?.
這是我聽的時候覺得很大的一個詫異的問號.
還有你剛才說的那首感恩詩.
他用了很多希伯來人傳統感恩的歌詞.
認罪詩也有很多,出名的認罪詩有51篇.
偏偏他沒有用.
我們讀者從第一章讀到這裡.
應該預期若拿這個時候是要懺悔的.
上帝就要拯救他.
如果這樣說,一般的猶太讀者覺得.
李英正常情況就是他在餘福當中.
有一篇類似詩篇51篇.
然後套出來,然後才有這首詩歌.
這首這麼美的串燒詩出來.
但中間好像人家數學做素小時候.
不知道是否香港來,跳step,跳skip,跳skip了他.

$^{201}$沒有了的.
真的有一點out of place.
如果這樣看的時候.
很多傳統的想法有一點不同.
如果是這樣,大家覺得.
這首詩這樣理解,重點是想說什麼呢?.
我們以往接觸這篇詩篇的時候.
我們有以往的印象.
可能我們都跳過step.
我們看了第一節,我們就跳到第十節.
就是將魚吞了若拿.
然後神就吩咐魚將若拿套到上帝身上.
然後他才作這麼美的thanksgiving prayer.
上帝拯救了我.
但作者就不讓我們這樣做.
作者要我們知道.
這個很美的感恩詩在餘福當中去做.
這是其中.
我覺得我們是有意無意地.
我們篡改了經文.
將第十節搬到前面.
以為魚吞了他.
然後他就在沙灘上.
開感恩會.
但事實上不是這樣.
我們一向都說有諷刺文學.
到第二章有什麼諷刺的元素呢?.
第一就是第一節的禱告.
這個最諷刺.
因為第一章在海欖上.
他反而不禱告.
外邦人是很禱告的.
大大的驚威.
第二章他就不斷說禱告.
這個就變成了一種很大的鮮明的諷刺.
如果這樣看的時候.
不知道大家有沒有感受到.
若拿這個先知在第二章的表現.
引我們去問一個問題.
究竟他現在是否在文學作品.

$^{241}$或者聖經作品裡面.
在這一刻其實是針對若拿先知.
其中一個的問題.
熟齡的問題.
不知道大家有什麼看法呢?.
我想像Ben這樣說.
大家可能也有這個意思.
如果我們將這首感恩的詩歌.
拿走了出來.
不理會上下文.
我們只看一首感恩詩.
這首詩歌在感恩會裡面.
應該是可以用得著的.
但是如果將它放在敘述的前後.
上下文下裡.
應用的神學就完全不同.
如果不理會上下文.
你會覺得這是感恩的.
死裡感恩逃難的很偉大的詩歌.
但如果放在敘述的前後.
就是一個諷刺.
諷刺一個可以作出這樣的歌曲的人.
一個這麼敬虔的人.
滿口都是一些很熟悉的句語.
很熟悉的祈禱,禱詞.
反而是在諷刺他.
但一個這麼熟悉的人.
卻看不到自己的問題.
我們預期他應該是先認罪.
誰知他不是認罪.
感恩死不去.
如果這樣說.
會不會構成了這首約拿書第二章.
其實是針對自義這個問題呢?.
或者是諷刺一群自以為很敬虔,很愛神.
又相信神很愛他們的一群人.
不要搞錯.
所有聽眾不要搞錯.
而是有一個格格不入的地方.
會不會是一個東西.

$^{281}$我們這樣看.
在諷刺一個情況就是.
他們理應要認罪.
卻不去認罪.
卻用一堆感恩的說話.
覆蓋了自己.
來逃避.
真的要面對上帝那一刻的處境呢?.
其實我覺得他們都很值得同情.
因為相信他們都未必知道自己是自以為的.
所以其實神都很愛他們.
走失了他們.
是真的一些迷失的.
我覺得這是一個好點.
所以我覺得為什麼會有諷刺文學.
就是要用諷刺文學才能打醒這些人.
他用諷刺的手法.
其實就是想透過這個方法.
來提醒他們.
好像新約裡面所說.
有意聽的就應當聽.
但是法利賽人就永遠聽不到.
如果用新約的語言.
就會說這些是法利賽人.
在十字路口的禱告.
類似.
耶穌都有用這些諷刺手法.
是呀,經常用在福音書.
在意義上.
約拿書第二章的神學都是這樣.
我們可以用在一個應用當中.
在現代意義上.
很有可能這本書是給基督徒看的.
而這班基督徒可能.
他自己覺得是很敬虔的.
覺得是很愛上帝的.
但是他從來看不到.
他其實在他生命當中.
仍然有些位置.
其實是沒有去處理他自己的問題.

$^{321}$這也是值得這首詩給我們.
這首歌是當頭棒學的.
我們看這段聖經.
不需要一定要指罵人.
首先我看回我自己.
是不是我這樣.
我是不是去到教堂敬拜.
我就舉手很興奮.
我不知道我本身有很多問題.
我是不想的那些.
不肯的那些.
我正在接受上帝.
但我去到教會很興奮的時候.
你看看那些不去教堂的人.
多離棄你.
我就不是了.
其實都是一個生命的反省.
特別是一些老信徒.
我記得上一次結束的時候.
我都記得說過.
我們這班老油條.
特別是信得久的時候.
和一些初信.
仍然有初心的一班.
初信弟兄姊妹.
其實都是的.
我們往往容易迷失.
這個都可以引申出.
這些叫做ritual.
那些的危險.
認識了.
認識了一些屬靈的儀式.
我們已經很熟習了.
聖餐怎麼拿我們不會拿錯.
不會搞錯步驟.
祈禱我們一定要走哪一句.
其實有一種很大的危險在.
危險在.
我們用這些就像阿公一樣.
我們遮蓋了我們的民團.

$^{361}$這個遮蓋.
是真的.
自己不知道.
習慣了.
在教會20年或20年以上.
很多的.
一些術語.
我們經常說基督徒的術語.
是我們自己才聽得懂.
微信的親朋戚友.
他們都不知道我們在說什麼.
其實那些.
其實我舉一個很悲哀的例子.
就是美國.
大家都知道美國南方.
其實都是一個.
在傳統基督教來說.
是一個比較保守.
但也是一個很熱心的社區.
但是很奇怪的.
黑奴.
地方出現得最厲害在哪裡?.
就是南方.
其實那些基督徒.
我覺得可惜的是.
他們看不到.
遮蓋了一些東西.
他以為那不是一個問題.
這個是值得我們去反省.
還有一個很像歷史的例子.
在深圳的人.
他們祈禱可以用King James的英文去祈禱.
很神奇.
我聽過.
很純粹.
很漂亮.
你知道King James是Doubt the true law.
Doubt the true law.
但是.
有一種很漂亮的ritual.

$^{401}$是不是你裡面.
其實就是包裝著你.
遮蓋著你裡面有很多問題.
這個我稱之為ritualism.
就是一個形式主義.
當一個形式變成一個主義.
變成僵化的時候.
我們就很多東西看不到.
聖餐也好.
我們祈禱也好.
就算是靈修也好.
如果成為一個.
形式主義.
例行公事.
我們就看不到那種自省.
正正是今天.
這個茶經.
值得我們去反省.
從約拿書的第二章去看.
約拿這個奇怪的祈禱.
提醒我們.
會不會有時候我們走進了形式主義.
而忘記了.
最終在上帝面前的真誠.
不過黃牧師不是叫你不要去祈禱會.
不是叫你不要去茶經.
我繼續做的.
大家不要誤會.
大家都明白的.
所以說是形式主義.
我們某些形式在這裡.
但他主要是僵化.
忘記了他本身的初心.
今天我們看到第二章.
時間過得很快.
我們就要跳入第三章.
所以請大家.
好好準備.
我們下一次的經緯線.
我們進入約拿書的第三章.

$^{441}$謝謝大家.
\newpage



\section{}
\label{sec:ksA3W2OSofQ}
\textbf{浩劫重生 【第6集:急難中的希西家】 陳崇基牧師 - 主日講道視頻系列}
\newline
\newline
連結: \href{https://youtube.com/watch?v=ksA3W2OSofQ}{\texttt{https://youtube.com/watch?v=ksA3W2OSofQ}} ~~~~ 語音日期: 2020-11-15
\newline
\newline
\hyperref[sec:_ZlpJCFpNOg]{\small{< < < PREV SERMON < < <}}
~
\hyperref[sec:index_chronic]{\small{[返順時目]}}
~
\hyperref[sec:index_scriptual]{\small{[返順卷目]}}
~
\hyperref[sec:UwZdh9361sc]{\small{> > > NEXT SERMON > > >}}
\newline
\newline
$^{1}$面對政局不安兵荒馬亂的時候.
不知道你通常會怎樣回應.
或者想想上帝想你怎樣回應.
我們今天看一段聖經的記載.
講到一個君王叫做希西加.
在他的事蹟當中.
當時是動盪不安甚至兵臨城下.
到底他怎樣回應.
或者在這段的記述當中.
上帝想他怎樣回應.
這段記載第十八章第一節.
以色列王以拉的兒子.
荷西亞第三年.
猶大王阿哈斯的兒子希西加登記.
他登記的時候年二十五歲.
在耶路撒冷作王二十九年.
希西加行耶和華眼中看為正的事.
效法他祖大位一切所行的.
這裡稱讚希西加.
和以前的王截然不同.
之前的王包括他的父王.
阿哈斯是敬拜別神.
當希西加登記的時候.
他做了什麼呢.
就是宗教的改革.
第四節這麼說.
他廢去休壇毀壞署帳.
砍下木偶打碎摩西所造的銅蛇.
因為當時以色列人.
迎向銅蛇銷香.
希西加稱銅蛇為銅塊.
希西加倚靠耶和華以色列的神.
在他前後猶大的烈王中.
沒有一個及他敵.
因為他專靠耶和華總不離開.
謹守耶和華所吩咐摩西的誡命.
在希西加之前.
猶大地不同的地方居然有不同的神廟.
當時人們將耶和華獨一的信仰.
混合了其他宗教的元素.

$^{41}$所以不單單在耶路撒冷有聖殿去敬拜神.
在其他地方他們敬拜一些別神.
並且混合了耶和華的信仰.
例如猶大南面一個駐防城阿拉特阿拉.
在這裡我們找到一個廟.
廟內有一個庭園.
庭園有一個大的祭壇.
祭壇按照摩西律法的尺寸來造.
在祭壇內我們找到一個致勝所.
正方的致勝所外面有兩個香壇.
一大一小.
在這個致勝所內有兩塊直立的石碑.
我們相信這兩塊石碑分別代表耶和華.
和亞瑟拉女神.
大概在希西加的時期.
這個聖廟被封閉和荒廢.
另外在別士巴南面的駐防城.
找到祭壇的壇角石.
證明在耶路撒冷以外.
人們私自建造祭壇燒香和獻祭.
這些都是在希西加的時期停用的.
譬如在阿拉特的聖所.
這些香壇很小心地被用灰土埋住.
所以封閉了這些不同的廟宇.
這些都是希西加做的宗教改革的跡象.
事實上他們這樣做是希望挽救猶太國.
因為隨從別臣的風俗.
就是使得上面北國以色列滅亡的原因.
十八章第九節.
亞述王將以色列人擄去亞述.
都因他們不聽從耶和華.
他們神的話違背他的約.
就是耶和華藉著僕人摩西所吩咐他們.
在希西加執政的時候.
他就徹底地在全國做了宗教改革.
十八章第七節這樣繼續記載.
耶和華與希西加同在.
他無論往何處都進刀亨通.
他也背叛不肯事奉亞述王.
希西加攻擊非利人直到加薩.

$^{81}$不單止他實行宗教改革.
希西加也在猶太國改變以往一直有的外交政策.
就是斷絕與亞述帝國建交.
因為以往的朝政都是靠攏亞述帝國.
當時亞述帝國是超級大國.
正在向西面擴張.
所以公元前七二二年的時候.
北國以色列亡於亞述.
上面十個支派就被流亡.
只剩下南面的猶太國.
你可以想像一下當以色列亡國的時候.
對於南國猶太國希西加來說.
壓力是有多大.
因為亞述帝國是當時的超級大國.
不斷地步步逼近.
並且以色列亡國之後的十幾年.
對於猶太國也有很大的震撼.
第一方面就是難民湧向南面的猶太國.
我們知道耶路撒冷本身的面積.
只是在狹窄的東山上.
在希西加作王的時候.
就擴張至西山.
估計當時的人口增加到最多二萬人.
第二方面希西加也在城的北面加建城牆.
在1970年的時候.
猶太人在舊城耶路撒冷發現一道寬牆.
石牆的位置.
這一個寬牆是闊有七米.
按照闊度估計.
這幅石牆是高達八米.
這道寬牆是向著耶路撒冷的北面.
最脆弱的地方.
防守著外敵的入侵.
第三方面.
希西加也預計亞述帝國會再來犯境.
所以就去建設水道.
這條水道是由大衛城的山下.
基分泉基漢水道開始.
這條水道一直將這些泉水改道穿過石山.
流到大衛城的南面西羅亞茨.

$^{121}$這條水道可以看到它迂迴曲折.
因為兩方面一起加工建造的時候.
他們需要靠當時的聲音.
來斷定他們要向哪一個方向挖掘.
最後能夠在中間相遇.
當時是鐵理時代.
一個最偉大的人工水道工程.
烈旺記下第20章第20節記載.
希西加有其他勇力和其餘的事跡.
就是他如何挖瓷挖溝和引水入城.
記載在猶大的烈旺記上.
這條水道有五百多米長.
出口有一個碑文.
我們叫它西羅亞碑文.
這裡記載水道的故事.
翻譯出來時這裡寫說.
斧頭彼此激打.
還有三爪的時候.
要再激打再挖掘的時候.
有人聲聽見他呼喊對面的同門.
石的右邊有缺口.
就在那一天.
工人的斧頭彼此激打.
將水道通了.
水就從源頭流到水池.
長度達一千二百爪.
這個冥文特別之處就是.
通常在近代世界.
冥文工程的領工人.
通常都是君王.
但這個冥文所提的只是一班工人.
兩組工人如何努力將水道建成.
似乎暗示希西加是一個愛民之王.
暗示希西加不是一個.
很喜歡自我榮耀的皇帝.
這是什麼.
(水).
我們很多時候有種錯覺.
我們覺得人已經做了我們要做的事.
我們已經盡了責任.

$^{161}$我們祈禱上帝一定會讓我們平安.
但不一定是這樣.
希西加做好了所有防戰的策略.
結果亞述帝國依然犯境.
在公元前七零一年.
亞述國真的來到猶大的境內.
在列王記下第十八章第十三節.
希西加王十四年的時候.
亞述王西拿基納上來.
攻擊猶大的一切堅固性.
將城攻取.
猶大王希西加就差人去拉吉見亞述王.
他說我有罪求你離開我.
犯你所罰的我必定承擔.
於是希西加就將聖殿和皇宮的金銀.
拿給亞述王.
希望他能夠退軍.
為什麼特別提拉吉這個地方.
因為在地理上拉吉是耶路撒冷的第二把交椅.
如果你看地理的話.
你會知道沿海平原上去中央山區耶路撒冷的地方.
是直著幾個山脊.
其中一個最方便能夠上到耶路撒冷的.
就是直著拉吉.
拉吉就是把守著一個最重要的關口.
如果拉吉淪陷的話.
軍隊就可以直接上到耶路撒冷.
這個征戰的紀錄.
其實亞述王也有提過.
在芝加哥近東的博物館.
你進去亞述的展館.
你可以看到一個文物叫Synagogue Prism.
西那基納國柱.
那裡記載著一段這樣的說話.
因為猶大王希西加不服於我定壓.
亞述王我派大軍攻打.
靠著軍力和我的勇力.
連取屬於他的四十六個堅防成仁.
擄走二十萬零一百五十六人.
擄亂男女連同馬匹,盧駒,駱駝,牛羊.

$^{201}$數目無可估計.
我又把希西加困於他的京城耶路撒冷之內.
猶如龍中鳥.
拉吉淪陷不只是聖經提過.
在亞述的歷史裡也是大事宣傳.
因為亞述王回到尼尼美的時候.
他在皇宮的牆壁上.
就弄了一個很大的浮雕.
大事的慶祝他攻打拉吉這件事.
如果我們一起去大英博物館亞述展館.
就可以看到這些浮雕.
就是如何紀念亞述王攻打拉吉的圖像.
這些士兵不斷射箭上拉吉防守城.
用這些脫石機器攻打城牆.
建了攻城的斜坡.
然後想上這個斜坡攻打城市.
拉吉的猶大民眾就像以卵擊石一樣反抗.
在畫的另一邊.
你也看到一些平民百姓在拉吉淪陷之後.
被擄去被流亡.
在城外也有很多殘殺的行動.
亞述帝國亞述軍隊最著名的就是他們的殘暴行為.
你看到這裡有幾個將軍被釘十字架.
在拉吉的城牆門外示眾.
再過一些我們看到一些士兵或猶大人民被劍山削皮.
在浮雕的另一邊.
文字上寫著我亞述王批准這些殘殺行為.
這裡就是拉吉.
你看到這個土優的形狀和地理.
其實和大英博物館的亞述浮雕很相似.
這裡是亞述軍隊所建的斜坡.
是用來攻城的.
上面也找到很多箭頭和石頭.
是當時他們擊戰時留下的遺物.
拉吉淪陷後亞述王派大軍到耶路撒冷.
在耶路撒冷的京城城外罵戰.
不知道大家有沒有聽過人罵戰.
這是很長的經文.
在列王記下十八十九章都有看過.
經文太長了我們只抽幾句.

$^{241}$總括來說是一種trash talk.
十八章第十九節.
拉伯沙基其實是亞述問你的.
即是高級酒精或者是commander的意思.
你們去高蘇欺世嫁孫.
亞述大王如此說.
你所倚靠有什麼可仗賴的呢.
你說有打仗的計謀和能力.
我看不過是虛話.
你到底倚靠誰才叛變我.
十八章二十五節.
現在我上來攻擊毀滅者地.
豈沒有耶和華的意思呢.
耶和華吩咐我.
我上來攻擊毀滅者地.
拉伯沙基又和城上的欺世嫁朝廷的官長說.
我主亞述王差遣我來.
豈是單向你和你的主說這些話.
不也是對這些坐在城上的百姓士兵來說話.
就是與你們一同吃自己的糞.
喝自己尿的人說話.
三十二至三十五節.
欺世嫁勸你們.
說耶和華必拯救我們.
你們不要信欺世嫁的話.
列國的神有哪一個.
救他本國脫離亞述王的手呢.
他們曾救撒瑪利亞脫離我亞述王的手嗎.
難道耶和華能夠耶路撒冷脫離我的手嗎.
這些是馬戰的垃圾話.
亞述其實也不想廢一兵一卒.
可以攻城.
是一種心理戰.
或者按照聖經所說是褻瀆的話.
當兵臨城下的時候.
欺世嫁有什麼回應呢.
十九章第一節.
對他說.
欺世嫁如此說.
然後下文.

$^{281}$在上帝休息著以賽亞回答欺世嫁禍.
第一句說話.
不要怕.
然後說了上帝的認許.
上帝的承諾.
你留意到這段說話當中有很多次序.
那個chain of command.
亞述王軍隊.
他的chain of command.
就是由亞述王自己.
將說話傳給他的將軍.
然後將軍就將他的說話帶給欺世嫁.
同樣的.
上帝藉著他僕人先知以賽亞.
又將他自己的說話帶給欺世嫁.
用亞述帝國的chain of command.
來回應這個危機.
其實背後想說的是.
在整個說話的次序.
命令的次序.
chain of command裡面.
到底哪一個說話才會兌現.
哪一個命令才是有效.
到底是亞述王還是上帝的說話.
然後到第十五節.
欺世嫁向耶和華禱告說.
在在以基督帕上的耶和華以色列的神.
你是天下萬國的神.
你曾經創造天地.
耶和華求你執意而停.
要請西拿基勒打發使者來辱罵永生神的話.
耶和華說亞述諸王果然使列國和列國之地變為荒涼.
使列國的神像都永在火裡.
因為他們本不是神.
乃是人手所造的.
是木頭石頭所造的.
所以滅絕他.
耶和華說你不同.
你是神.
現在求你救我們脫離亞述王的手.

$^{321}$使天下萬國都知道唯獨你耶和華是神.
一個滿著信心的禱告.
去認定耶和華是獨一的真神.
獨一的救主.
只有他才能夠在急難當中成為我們最後的依靠.
神藉著以賽亞去回應這個禱告.
第三十二節.
所以耶和華論亞述王如此說.
亞述王必不得來到這城耶路撒冷.
也不得在這裡射箭.
不得拿盾牌到城前.
也不能夠捉女攻城.
他從哪條路來.
必從哪條路回去.
必不得來到這城耶路撒冷.
這是耶和華說的.
因我為自己的緣故.
又為我僕人大衛的緣故.
必保護拯救這城.
我們看這裡的戰況很緊張.
結果是怎樣呢.
十九章三五節很簡單地交代.
當夜耶和華的使者出去.
在亞述的營中殺了十八萬五千人.
清早有人起來一看.
全是死士.
亞述王西那基納就不認回去.
住在尼尼微.
有一天他在神廟裡叩拜的時候.
他兩個兒子就用刀刺殺他.
然後他另外一個兒子.
以撒哈頓就接續他做王.
當然在耶路撒冷打敗仗.
然後退軍.
亞述的文獻當然不會提.
在極權的社會裡面.
那些皇帝永遠都不會認輸.
不過在字裡行間.
其實就看得出.
西那基納其實是認了自己是打敗仗.

$^{361}$為什麼呢.
第一 拉吉只是一個小鎮.
怎麼會打完拉吉之後.
回到京城大肆宣揚.
然後在浮雕大大宣揚自己打勝仗.
而一隻字都沒有提過.
他打過耶路撒冷.
第二.
他說將希西加圍困如同龍中鳥.
這句話其實已經代表.
他沒有進過耶路撒冷.
並且亞述的學者告訴我們.
每逢亞述的文獻提到龍中鳥這句話.
其實就代表亞述皇帝打敗了.
他攻不了那個城市.
龍中鳥只不過是保存面子的說法.
這是一個很神奇的事跡.
其他的歷史學家都有提過.
有一個叫希羅多德的希臘歷史學家.
他後期說回這件事.
他解釋當中有鼠疫.
使亞述君敗亡.
當然這些全部都是逆賊.
因為經文裡面最重要的訊息就是.
這是上帝自己的作為.
更重要的是你看看那些字眼.
他說當夜 耶和華的使者出去.
這裡有晚上 有使者.
然後敵軍敗亡.
你會想起什麼技術.
就是出埃及既定 十災.
同一個道理.
一個冒犯上帝的君主.
一個暴君 欺壓神的子民.
當日埃及帝國的法老被剪除.
這裡亞述帝國的皇帝也被剪除.
我希望藉著這個技術給你和我一個鼓勵.
當時最大的鼓勵不一定是.
告訴你多祈禱.
不一定是告訴你不要害怕.

$^{401}$而是藉著一個技術 一個故事.
當你去細讀這個故事的時候.
上帝的話語帶領你進入這個技術當中.
你就好像身在其中一樣.
然後你出來的時候.
你的心境 你的角度會完全不同.
因為在技術裡面.
我們遇見一位行神跡的上帝.
一位在背後 執政掌權.
都是在祂手中的大能的神.
我們就能夠有信心在上帝面前禱告.
如果你去到一個地步.
你現在覺得很無助 很無能.
你什麼方法都試過.
我想給你一個邀請.
就好像這個故事所給我們邀請一樣.
來到上帝的面前禱告.
並且禱告的方向不是脫離苦難這麼簡單.
而是要讓萬國 讓身邊的人.
認定耶穌基督才是真正的神.
這就是希西加禱告的精髓的地方.
十九章第十九節.
那天耶和華說.
我們的神現在求你救我們脫離亞述王的手.
使天下萬國都知道.
唯獨你耶和華是神.
不是單單我們脫離苦難.
而是在經歷當中.
最後人們知道 天下萬國知道.
唯獨耶和華才是神.
另外一方面.
很多時候我們留意地上的政局很混亂的時候.
我們就很關注一個問題.
到底哪一個國度 哪一個君主是最有權勢.
最有權勢的我們就想靠攏他.
但很多時候我們看見世界的政局發展的時候.
我們很少留意.
原來上帝其實是國際關係中的一個玩家.
並且他是最大的玩家.
這其實就是聶王記的信息.

$^{441}$聶王記告訴我們.
上帝才是最高的權力.
神要去到立王廢王的時候.
哪一個強國如果逆他意的話.
就必定會被打敗.
這其實就是聶王記提醒神的子民.
告訴他們 特別是猶大國.
很多時候他們陷入一個試探的當中.
特別是猶大國的結尾的時候.
他們就左望右望.
到底是否亞述是最強國.
如果亞述是最強國的話我們就事奉他.
如果是巴比倫最強國的話我們就事奉他.
如果是埃及最強國的話我們就事奉他.
所以整個猶大的朝廷.
很多時候就是轉來轉去 擺來擺去.
哪一邊強我們就轉去那一邊.
聶王記是寫給已經流亡的猶太人.
主要是解答他們一個問題.
就是為什麼我們現在流亡被擄.
是不是因為上帝沒有能力.
聶王記就回答他們.
他說不是上帝沒有能力.
而是我們三心兩意.
上帝當然有能力.
他曾經救過希西加.
猶大國脫離亞述帝國.
但我們作為神的子民.
我們有沒有想過.
到底我們是否真的敬拜上帝.
還是我們仍然是三心兩意.
故事邀請我們專心地.
來敬拜事奉耶和華.
行他的話語所吩咐的.
在地上我們作為天國的子民.
專心地去依靠他.
我們同心低頭祈禱.
神我們求你親自幫助我們.
讓我們專心去信靠你自己.
因為唯獨你才是真正的神.

$^{481}$唯獨你才是那位最有權柄的人.
主要這就是我們需要的宗教改革.
幫助我們專心一意尋求跟隨你.
讓你的教會讓熟你的子民.
專一敬拜.
專一去稱頌你.
以至於在我們的生活.
我們的生命當中.
能夠讓人知道.
唯獨耶和華你才是真神.
祈禱奉主命救.
\newpage



\section{}
\label{sec:UwZdh9361sc}
\textbf{神的名字4\_Yahweh Yireh 耶和華預備 (靈思牧禱)}
\newline
\newline
連結: \href{https://youtube.com/watch?v=UwZdh9361sc}{\texttt{https://youtube.com/watch?v=UwZdh9361sc}} ~~~~ 語音日期: 2020-11-17
\newline
\newline
\hyperref[sec:ksA3W2OSofQ]{\small{< < < PREV SERMON < < <}}
~
\hyperref[sec:index_chronic]{\small{[返順時目]}}
~
\hyperref[sec:index_scriptual]{\small{[返順卷目]}}
~
\hyperref[sec:kB2Ko_XJaY8]{\small{> > > NEXT SERMON > > >}}
\newline
\newline
$^{1}$人生難免困境.
特別在2020年的疫情和社會動亂裡面.
很多人直接受到打擊.
其中有失業的,生意倒閉的.
因為封鎖義大利嚴重家庭問題的.
當然還有不幸染上病毒.
痛苦掙扎甚至死亡的.
不過老實說.
封鎖之前人生其實不是.
也有很多的挑戰和患難嗎?.
而封鎖之後.
我們可以重啟所謂正常生活的時候.
人生是不是就一帆風順呢?.
我們都知道不是.
正因為這樣.
神的名字「Yahweh Yaira」.
中文譯為「耶和華爾勒」.
亦是「耶和華預備」的意思.
對我們非常之寶貴.
值得反思.
說到「耶和華爾勒」.
我們一定要從《創世記》22章.
一個真實而經典的故事說起.
這個故事圍繞著一個爸爸.
面對人生極大的困境.
這個爸爸很多人都熟悉.
就是我們先前所提過的亞伯蘭.
後來稱為亞伯拉罕.
上帝呼召他離開本鄉本族.
去到引領的土地.
建立一個與上帝立約的民族.
這個民族會逐漸繁衍.
亦會駛到萬國得福.
亞伯拉罕就是以一顆堅定的信心.
跟隨上帝去了.
但是有一天.
亞伯拉罕落在一個不可思議的困境裡.
原來上帝忽然要試驗亞伯拉罕的信心.
他吩咐他去到摩利亞的一個山上.
獻上梵祭.

$^{41}$梵祭是用火燒的祭.
是被燒的聖祭.
要整隻獻上.
而上帝這個要求.
對亞伯拉罕的打擊是幾方面的.
第一 因為上帝這次要求獻上的祭神.
不是動物.
而是亞伯拉罕心愛的兒子 以撒.
第二 在古時的近東.
兒子就是晚年的寶藏.
上帝竟然叫亞伯拉罕放棄這個寶藏.
他自然就會想晚年怎麼辦.
第三 以撒是上帝曾經允許給亞伯拉罕的兒子.
使他民族成為大國的關鍵.
現在竟然成為憲制.
那麼神對他這個民族的應許會變成怎樣呢?.
我們可以想像.
亞伯拉罕的心情是很複雜的.
以前上帝叫亞伯拉罕去什麼地方他都去.
而上帝一直就這樣祝福他.
但是這次獻上以撒的這件事.
似乎將他什麼祝福都拿走.
亞伯拉罕還能不能夠聽上帝的命令這樣去呢?.
亞伯拉罕實在陷入很大的迷茫裡面.
他痛苦掙扎.
一方面他相信有必須被繁祭的陽光.
但是另一方面他實實在在地遵從上帝的命令.
在祭壇上他拿起刀來準備將以撒獻上.
就在這個清君一法的時候.
有天使從天上呼叫亞伯拉罕說.
「你不可在這童子身上下手」.
「現在我知道你是敬畏神的了」.
「因為你沒有將你的兒子留下不給我」.
在這個時候亞伯拉罕抬頭一看.
竟然見到有一隻公羊.
兩隻角扣在稠密的小樹叢裡面.
亞伯拉罕深深知道.
這是神所預備的新祭.
他很興奮地將這個公羊拿下來獻為繁祭.
代替他的兒子.

$^{81}$我們都很為亞伯拉罕而高興.
因為最後是大團圓結局.
兒子不需要喪命.
而他親身經歷了耶和華特別為他所預備.
所以聖經繼續有這麼多的記載.
它說:亞伯拉罕給那地方起名叫耶和華爾.
直到今日還有人說.
在耶和華的山上有雨被.
細心地解讀.
這是一個非常震撼的信心故事.
好像我以前提過.
上帝曾給亞伯拉罕的祝福.
是會隨著獻上兒子以殺.
一下子通通都失去.
而上帝正是以這一點.
來考驗亞伯拉罕的信心.
看看你的信心是有條件還是沒有條件.
你敬畏我是否因為有好處呢?.
如果這些好處都沒有了.
你又會怎樣呢?.
而亞伯拉罕最後顯示他的信心就是.
即使沒有好處.
我也敬畏遵從上帝.
結果就是在似乎最失去的日子.
也是他最體會上帝預備的日子.
請注意.
耶和華爾辣.
爾辣原文yara這個字的意思.
除了是我們時時說到的預備,供應.
更加有「看到」的意思.
即是耶和華看到亞伯拉罕對他全心的遵從.
他觀察所有的細節.
以至他預備,供應他所需要的.
弔詭的就是.
我們往往以為.
我們要得到上帝的預備.
才會對他全心地遵從.
而上帝卻告訴我們.
我們如果先全心地遵從.
必定會得到他的預備.

$^{121}$就好像亞伯拉罕一樣.
「也為耶華」.
「耶和華爾辣」.
意思是上帝看到.
信心只是掛在口頭上.
不過是空談.
但是如果你用行為表現出來.
就好像亞伯拉罕一樣.
上帝必定看到.
「也為耶華」.
「耶和華爾辣」.
意思是上帝供應.
即使在不可思議的困境裡面.
你仍然對上帝忠心忠誠.
他必定供應給你.
「耶和華」今天.
怎樣為你和我預備.
我不知道.
但是我知道上帝的這個名字.
帶來好心的提醒.
就是在人生困境裡面.
只要你堅定相信和遵從上帝.
就好像失去一切保障的時候.
也只是你體會上帝預備的時候.
讓我們一起禱告.
上帝啊.
今天我們都仍在疫情的威脅.
和政治的動亂裡面.
很多人對前途感到迷茫.
也都受到很多的打擊.
你一一都看到.
「耶和華爾辣」.
多謝你藉著亞伯拉罕的故事.
一再的提醒我們.
即使今天面對多大的困境.
我們仍然要敬畏你.
遵從你的話.
以示我們親嘗到你的預備.
以示到我們的眼睛一再明亮.
看到耶和華真的與我們同在.

$^{161}$「耶和華爾辣」.
因為當這個世界因罪而殘缺敗壞.
我們人人都承受其中的痛苦.
沒完沒了的時候.
你竟然將你獨身兒子耶穌基督獻上.
聖位贖罪祭.
讓願意悔罪投奔你的都得到拯救.
願我們贖你的子民.
我們不忘記數算這個救贖的恩典.
也不忘記數算你在我們過去生命上的恩典.
以致即使我們在無論當前日子多難.
我們仍然堅持到底.
經歷你的預備.
永要歸於你的名.
我們禱告是奉救主耶穌基督的名.
阿們.
多謝您收睇時局新聞,再會!.
\newpage



\section{約拿書 3:1-10}
\label{sec:kB2Ko_XJaY8}
\textbf{經緯線 [釋經視頻] EP5 約拿書 3\_1-10}
\newline
\newline
連結: \href{https://youtube.com/watch?v=kB2Ko_XJaY8}{\texttt{https://youtube.com/watch?v=kB2Ko\_XJaY8}} ~~~~ 語音日期: 2020-11-20
\newline
\newline
\hyperref[sec:UwZdh9361sc]{\small{< < < PREV SERMON < < <}}
~
\hyperref[sec:index_chronic]{\small{[返順時目]}}
~
\hyperref[sec:index_scriptual]{\small{[返順卷目]}}
~
\hyperref[sec:EWFiHkhUrkU]{\small{> > > NEXT SERMON > > >}}
\newline
\newline
約拿書 3:1-10
\newline
\begin{longtable}{cl}
\hline
\hline
章節 & 經文 (和合本修訂版)\\
\hline
3:1 & \begin{tabularx}{0.7\textwidth}{X} 耶和華的話第二次臨到約拿，說： \end{tabularx} \\ \\ \relax
3:2 & \begin{tabularx}{0.7\textwidth}{X} 「起來，到尼尼微大城去，把我告訴你的信息向其中的居民宣告。」 \end{tabularx} \\ \\ \relax
3:3 & \begin{tabularx}{0.7\textwidth}{X} 約拿就照耶和華的話起來，到尼尼微去。尼尼微是一座極大的城，約有三天的路程。 \end{tabularx} \\ \\ \relax
3:4 & \begin{tabularx}{0.7\textwidth}{X} 約拿進城，走了一天，宣告說：「再過四十天，尼尼微要傾覆了！」 \end{tabularx} \\ \\ \relax
3:5 & \begin{tabularx}{0.7\textwidth}{X} 尼尼微人就信服神，宣告禁食，從最大的到最小的都穿上麻衣。 \end{tabularx} \\ \\ \relax
3:6 & \begin{tabularx}{0.7\textwidth}{X} 這消息傳到尼尼微王那裡，他就從寶座起來，脫下朝服，披上麻布，坐在灰中。 \end{tabularx} \\ \\ \relax
3:7 & \begin{tabularx}{0.7\textwidth}{X} 他叫人通告尼尼微全城，說：「王和大臣有令，人、畜、牛、羊都不可嘗任何東西，不可吃，也不可喝水。 \end{tabularx} \\ \\ \relax
3:8 & \begin{tabularx}{0.7\textwidth}{X} 人與牲畜都要披上麻布，切切求告神。各人要回轉離開惡道，離棄自己掌中的殘暴。 \end{tabularx} \\ \\ \relax
3:9 & \begin{tabularx}{0.7\textwidth}{X} 誰知道神也許會回心轉意，不發烈怒，使我們不致滅亡。」 \end{tabularx} \\ \\ \relax
3:10 & \begin{tabularx}{0.7\textwidth}{X} 神察看他們的行為，見他們離開惡道，神就改變心意，原先所說要降與他們的災難，他不降了。 \end{tabularx} \\ \\
[1ex]
\hline
\hline
\end{longtable}
$^{1}$各位好,各位好,現在是經衛線的時候了.
上一次我們看完約拿書的第二章.
從約拿一個我們覺得相當之奇妮的祈禱.
我們看到一樣東西,會不會有時我們都好像約拿一樣.
在那個祈禱或者一個生活當中.
我們表面上有很多的敬虔,很多的讚美.
但是其實背後是隱藏著我們去逃避.
去承認自己需要認罪的一個情況.
將我們的屬靈生命變得形式化了.
而今天我們就轉到第三章了.
而在第三章裡面很有趣.
因為在第二章的第十節.
當約拿做完祈禱之後.
耶和華吩咐大禹一吐.
就吐了約拿出來.
吐了約拿出來之後,第三章就開始了.
究竟當約拿在一個旱地的時候又發生了什麼事呢?.
約拿在旱地上的時候所發生的事就是.
耶和華再一次的來呼召她.
在這次的呼召裡面,神的呼召的內容沒有變.
她的心意也沒有變.
變的就是約拿變了.
在第一章的時候,約拿是逃避耶和華神的呼召.
但是來到第三章的時候,約拿就是聽耶和華神的吩咐.
聽耶和華神的吩咐,去了利利眉.
但是到了利利眉之後.
究竟大家有沒有留意她那段跟著的記錄有什麼特別呢?.
Benny看到什麼呢?.
我覺得這段記錄是很過癮的.
首先經文說到利利眉是一個很大的城.
要三天才能走完.
約拿只能走一天.
然後她就跟別人宣告.
她說再等四十天.
四十天,利利眉就傾覆了.
這些跟平時先知所說的有一點不同.
因為她只是說一個consequence.
她沒有說那個解救的方法.
其實我想問,其實約拿是不是在敷衍呢?.
因為原文,希伯文宣告第四字的內容只有五個字.

$^{41}$真的很厲害.
是先知裡面最短的一篇道.
五個字就搞定了.
這些在教會裡面就能解僱了.
我們做牧師的,如果上台講那五個字.
十秒鐘就下台,大家就會瞪著你,瞪著你發生什麼事.
還有一件事,Mark Sir.
很多時候在經文裡說.
這個城要三天才能走完.
好像很大,要走一天.
其實在考古方面,我不知道你對這方面比較熟悉.
究竟你在現在考古學裡面,怎樣理解這件事情呢?.
完全是不會的.
因為利利眉城怎麼大,都是幾個小時走完.
不可能會走三天.
我想第一方面,大家要記得.
約拿書是一個諷刺文學.
是有些誇大的修飾.
所以就說它極大.
你就不需要很現實地去看.
另外三天,當時近冬.
那些王要發一個使者去另外一個地方宣告訊息.
他應該會留在那裡三天.
就是說利利眉城很尊重.
這個耶和華臣.
這個王,差派來的先知.
他就在那裡三天時間.
但他走了一天.
但他走了一天,進城走了一天.
宣告訊息,五個字就搞定了.
其實說起來,就像你老闆叫你.
你在那裡代表我三天.
但你講完一天就撇了,走了.
真的很敷衍.
講完之後,躲在自己的Airbnb裡面.
OK,一篇五個字的道.
叫你去三天,就只有一天.
但似乎這個所謂這麼敷衍的行動.
似乎那個反應又不是想像中這麼敷衍.
大家留意到當時那些人聽到這個訊息之後.

$^{81}$那些人的反應是怎樣呢?.
第五節第一個動詞就是信服.
They trust,利利眉人就信服了神.
除了信服之外.
他好像那個利利眉城的人.
好像全城人不理大小,年老年幼都在披麻.
這些又是一些keyword.
作者經常用的.
第四節就是若拿進城,走了一天.
宣告說,Call out.
他只講了五個字.
但那些利利眉人信服神.
宣告又是Call out.
宣告了很多東西.
很大contrast,一個五個字.
差不多是五節,五句.
很大段.
如果你留意這段利利眉城的人的反應.
很有趣的.
他先是從最大到最小.
然後訊息傳到利利眉王.
然後王又下了寶座,脫了朝服.
穿上披麻,坐在灰中.
他又要向全城人宣告.
好像是由下而上,到皇帝都聽都去這樣做.
是一個很大的復興.
作者用了很誇張的手法.
我想他想突顯,想講出那個重點.
有一點像你可以想像.
他只是說五個字.
誰知在大球場,郭培里那些.
一千人下來決戰.
Yes, yes, I am.
然後一群人湧下來.
他們自己唱,不用你唱.
你沒有預備師班,他們已經下來了.
很有趣,不知道大家有沒有留意.
在當中,尤其是王的那封信.
寫給所有的王與大臣有靈的.
寫給全城人.

$^{121}$他說人不可以吃東西.
固然啦,這個生畜,牛羊.
不可以吃草,又不可以喝水.
生畜和人都要披麻布.
其實那個形象這一刻都覺得很奇怪.
很誇張.
人不可以吃東西,人不可以喝水.
你如何控制動物不吃不喝呢?.
我也想知道.
是,是畜生.
這個又回到諷刺文學.
我們經常提到諷刺文學有什麼特色.
就是誇張.
還有很多動物會出現.
諷刺一些東西的時候.
譬如你想起那個Balamb.
Balamb的假仙子.
也有隻驢仔會說話.
這是一些諷刺的手法.
在約拿書一直看一章到這裡.
都是有的.
大魚就聽耶和華話.
叫牠吐出來就吐.
叫牠吞牠就吞.
然後這裡有牛,有羊,有生畜.
這些就好像很聽話.
他們好像人一樣去悔改.
生畜很聽話.
生畜很聽話,的確很聽話.
那個生畜是諸狗不如.
還有一件事是真的.
Benny在這個裡面.
他裡面說到城裡面未成的人.
都是那幫人.
你看到後面他講.
他說人要切切靠高神.
各人回頭離開所行的惡.
丟棄手中的強暴.
或者神轉意後悔.
不發烈怒使我們不自滅亡.

$^{161}$也未可知.
其實那幫人某程度上沒有把握.
他只說的是薄一薄.
試試看.
看神會不會網開一面.
是這樣的一個表達.
這個就約拿應該要跟他們說的.
因為先知應該這樣說.
你看耶利米永遠都是這樣說.
你後悔的話.
上帝就不會剿滅你們.
但搞笑的是.
約拿的影響力是比他們厲害.
說了五個字已經全國.
耶利米真慘.
說完人們都不聽.
下次上台講道理.
你給雷劈算了.
不用說了.
為什麼.
剛才約拿應該要說的.
他們要悔改的說話.
沒有說到.
但竟然是由外邦人的口說出來.
如果我是猶太讀者.
我會噴血的.
為什麼會這樣.
為什麼會知道.
我會立即噴血出來.
或者噴水.
因為有一個很諷刺的.
很大的對比.
先知不說.
要說的就不說.
不懂的.
外邦王去說.
是的.
這樣就真是有趣.
而且最後結果.
我們更加看到在第十節說.

$^{201}$他說.
於是神察看他們的行為.
見他們離開樂道.
他就後悔不把所說的哉.
講於他們了.
又說後悔.
轉念.
真是改變.
真是不出手.
大家說得對.
一個敷衍五個字的道.
就令到全城悔改.
就令到上帝最後回心轉意.
一個改變.
想到這裡的時候.
短短十節的經文.
這十節的經文對你來說.
在第三章來說.
給了你什麼提醒.
或者一些想法.
讓你去了解這段經文呢?.
我想很強烈對比就是.
城中的人的悔改.
包括那些的心畜.
雖然很誇張的手法.
但是其實也表明了.
他們的敬虔.
對耶和神的敬畏.
當他們和舊之後.
襯托出若拿的先知.
他的行動.
他所做的事.
反映出他根本不敬虔.
但他自以為敬虔.
有這個道理.
其他呢?.
其他你會看到什麼訊息.
你覺得這段經文給了你嗎?.
我會想像第二章.
如果你read out of context的話.

$^{241}$你不理會文體是一個諷刺的文體.
你可能去到第三章.
你就是一個宣教大會會用的經文.
你看看我們願意去傳的話.
嘩!很多人會悔改.
上帝得榮耀.
你就會是mysticological application.
但如果你看文體的genre.
是諷刺文學的時候.
其實就是一個照妖鏡.
target會不會是我們呢?.
我們就是一班不明白上帝.
是一個有compassion的人.
只要人去悔改的話.
他就會去憐憫.
但我又沒有去悔改.
因為target主要不是外邦人信主.
主要是在諷刺一個.
敬虔所謂先知的人.
很認識上帝的.
他不會悔改.
我覺得這個照妖鏡其實是很timeless的.
因為你看一些secular的organization.
他們做錯事的時候.
他們很快會有一個press conference去認錯.
反而一些教會的團體.
他們很多事情就是掩蓋著.
但最後掩蓋不了就爆出來.
這個我覺得是很sad的.
我看到一件事其實大家都說了.
我從另一個角度去想.
約拿棄港.
這裡沒有說為什麼棄港.
但是有一樣東西肯定的.
我想他沒有預期到.
他這次這樣去說五個字的道.
就變成一個萬人空巷的報道會.
可以這樣說.
其實反過來我們可以問一個好的問題.
其實誰都可以回應上帝的拯救.

$^{281}$就算那些人或者那些群體.
可能你未必會覺得他會聽這個救恩.
因為你看到你未必成.
大家都知道當時亞述曾經的強國.
雖然開始衰落但他始終是強.
而且手段都不是惹事的國家.
但是沒有想過他聽了一個訊息.
五個字他就悔改.
其實我覺得對於我們今天來說.
一個現代意義.
其實會不會有時.
作為我們一個信徒,老信徒.
我們已經定了一些框架.
有些人永遠都不會信耶穌.
有些人永遠都不會悔改.
將他們判了死刑.
但是你又怎麼知道.
我又怎麼知道呢.
會不會是我們的自意.
或者會不會是我們已經有一個固定的框架去看那些人.
所以我們就懶得做.
但是誰知道神是可以透過我們所做的事.
就算是懶得做都好.
都可以產生這麼大的變化.
當然這個重點不是他們信耶穌.
重點是問我們.
神是可以出手的.
我們怎麼看呢.
為什麼我這樣說呢.
因為我想起一個真人真事.
大家都知道過去這段時間.
有很多宗教機構的報告都說很多回教徒信耶穌.
而他們信耶穌有很多見證.
其實沒有人和他們傳統福音.
但心裡不知道為什麼他們聽到福音.
甚至是做夢見到耶穌.
最後追尋耶穌.
我曾經聽過一個真人真事的故事.
是在瑞典發生的.
就是一群在中東的難民.

$^{321}$信耶穌去到瑞典.
他們走進瑞典的教會.
他們見到瑞典的教會死氣沉沉.
他們就去問瑞典的牧師和教友信徒.
問他們.
你信耶穌是這樣的嗎.
你不是有種火熱的喜樂.
你不是要悔改的嗎.
令到那群牧師和那群在瑞典的信徒.
一下子當頭拍下鶴.
不知如何回應.
我想某程度上.
有少少相似.
可以這樣說.
可能這一段的經文.
是一個諷刺.
其實是諷刺了一群人.
你沒有想過原來上帝的悔改.
可以去得很徹底.
也可以去得很深.
哪管我們這群老油條.
其實可能從來都沒有想過.
神就令到這件事情發生了.
或者真的諷刺了我們一群.
信了主久的老油條.
我們究竟對神的救恩.
即是所謂的恩典.
我們究竟理解有多少.
真正正甚麼是恩典呢.
或者好像你剛才說悔改.
是不是就不在我們的字典裡面.
其實我們才需要悔改.
我記得有一次很多年前.
那時候我在另一個中派服事.
我是做中派的區的區委員.
我和另一個牧師去探他.
去鄉鎮教會探他們.
我見到一個很大的陣勢.
我們大家站在那裡看.
然後我突然間心血來潮.

$^{361}$我講了一句說話.
我說你有沒有留意.
我們和人決戰信耶穌的時候.
講很多愛很多恩典.
但一受浸一浸下去之後.
一扯上來.
那些愛和恩典都留在水缸裡.
那些恩典和愛都在水缸裡.
然後我們講甚麼呢.
沒有了.
其實有時候會不會都是一樣.
那種自然就在這裡.
這個是我想到的提醒.
還有另一件事我覺得.
反過來講就是.
原來上帝可以.
祂的慈愛和恩惠.
如果神要工作是阻隔不了的.
我記得有一句說話.
我不知道是誰人講的.
是中世紀有人說.
中世紀有人說.
有人說是馬丁路德說.
但那句說話是甚麼呢.
God can draw a straight line.
with a crooked stick.
翻譯出來就是.
上帝可以用一支很曲的棍.
來畫一條直線.
我相信這句.
都可以讓我們看到另一個點.
就是上帝的工作.
其實是阻不了祂的.
但其實當上帝工作時.
正是一個照妖鏡.
照我們怎樣看自己的生命.
我想這個也是一個重點.
很多人都會問.
上帝要尼尼微尾的人悔改.
可以不用若拿.

$^{401}$為甚麼一定要逼若拿去呢.
其實重點不是.
上帝要教的課程.
主要不是尼尼微尾的人悔改.
其實就是這個才是.
會不會代表我們呢.
其實若拿這次悔改的機會.
神給的機會.
本身是一個因點.
讓若拿真真正正去認識.
甚麼是舊因.
甚麼是耀和神的因點.
所以看回之前的對比.
他在難處的時候.
生死的關頭的時候.
耀和神是怎樣救他.
同樣耀和神.
也是怎樣救尼尼微尾的人.
其實去到這裡.
我們可以再想.
如果這個故事.
整個約拿書這段經文四章.
主角是上帝與若拿.
去了三章.
究竟若拿知道了.
上帝後悔不降和.
給尼尼微尾的人.
究竟若拿變成怎樣呢.
他的經歷.
下回分解.
當然下回.
還說不吃飯.
今天還未吃飯.
吃宵夜.
我們一起吃拉麵.
我們就去到這裡.
下一集就說到第四章.
現在在第三章我們就看到.
原來問我們.
是否看不看到原來上帝的工作.

$^{441}$和看不看到在這個過程裡面.
原來我們會否好像若拿一樣.
需要重新反省.
甚麼叫做愛.
明明真的明白上帝的愛和上帝的恩典.
是否我們自己已經畫了框架.
還是上帝的愛和框架可以很闊呢.
我們就留在這裡.
下一次我們一起準備好.
看第四章.
約拿書第四章.
然後進入經維善.
下集再見.
\newpage



\section{約拿書 4:1-11}
\label{sec:EWFiHkhUrkU}
\textbf{經緯線 [釋經視頻] EP6 約拿書 4\_1-11}
\newline
\newline
連結: \href{https://youtube.com/watch?v=EWFiHkhUrkU}{\texttt{https://youtube.com/watch?v=EWFiHkhUrkU}} ~~~~ 語音日期: 2020-11-25
\newline
\newline
\hyperref[sec:kB2Ko_XJaY8]{\small{< < < PREV SERMON < < <}}
~
\hyperref[sec:index_chronic]{\small{[返順時目]}}
~
\hyperref[sec:index_scriptual]{\small{[返順卷目]}}
~
\hyperref[sec:vU00DNbRAss]{\small{> > > NEXT SERMON > > >}}
\newline
\newline
約拿書 4:1-11
\newline
\begin{longtable}{cl}
\hline
\hline
章節 & 經文 (和合本修訂版)\\
\hline
4:1 & \begin{tabularx}{0.7\textwidth}{X} 這事令約拿大大不悅，甚至發怒。 \end{tabularx} \\ \\ \relax
4:2 & \begin{tabularx}{0.7\textwidth}{X} 他就向耶和華禱告，說：「耶和華啊，這不就是我仍在本國的時候所說的嗎？我知道你是有恩惠，有憐憫的神，不輕易發怒，有豐盛的慈愛，並且會改變心意，不降那災難。我就是因為這樣，才急速逃往他施去的呀！ \end{tabularx} \\ \\ \relax
4:3 & \begin{tabularx}{0.7\textwidth}{X} 耶和華啊，現在求你取走我的性命吧！因為我死了比活著更好。」 \end{tabularx} \\ \\ \relax
4:4 & \begin{tabularx}{0.7\textwidth}{X} 耶和華說：「你這樣發怒，對嗎？」 \end{tabularx} \\ \\ \relax
4:5 & \begin{tabularx}{0.7\textwidth}{X} 約拿出城，坐在城的東邊，在那裡為自己搭了一座棚。他坐在棚子的蔭下，要看看城裡會發生甚麼事。 \end{tabularx} \\ \\ \relax
4:6 & \begin{tabularx}{0.7\textwidth}{X} 耶和華神安排了一棵蓖麻，使它生長高過約拿，影子遮蓋他的頭，使他免受苦難；約拿因這棵蓖麻大大歡喜。 \end{tabularx} \\ \\ \relax
4:7 & \begin{tabularx}{0.7\textwidth}{X} 次日黎明，神卻安排一條蟲來咬這蓖麻，以致枯乾。 \end{tabularx} \\ \\ \relax
4:8 & \begin{tabularx}{0.7\textwidth}{X} 太陽出來的時候，神安排炎熱的東風，太陽曝曬約拿的頭，使他發昏，他就為自己求死，說：「我死了比活著更好！」 \end{tabularx} \\ \\ \relax
4:9 & \begin{tabularx}{0.7\textwidth}{X} 神對約拿說：「你因這棵蓖麻這樣發怒，對嗎？」他說：「我發怒以至於死，都是對的！」 \end{tabularx} \\ \\ \relax
4:10 & \begin{tabularx}{0.7\textwidth}{X} 耶和華說：「這棵蓖麻你沒有為它操勞，也不是你使它長大的；它一夜生長，一夜枯死，你尚且愛惜； \end{tabularx} \\ \\ \relax
4:11 & \begin{tabularx}{0.7\textwidth}{X} 何況這尼尼微大城，其中不能分辨左右手的就有十二萬多人，還有許多牲畜，我豈能不愛惜呢？」 \end{tabularx} \\ \\
[1ex]
\hline
\hline
\end{longtable}
$^{1}$各位好 各位好 《經衛線》又回來了.
我們上次看完約拿書的第三章.
在那裡看到約拿講了一堂五個字的道.
接著就全城悔改.
真是厲害的葛培莉 太厲害了.
Danny 我記得你之前在開頭問過一個問題.
如果有一個先知一事奉大有成果的時候.
究竟先知會怎麼想?.
不知道你現在還想起這個問題.
你是有什麼回應?.
我這個問題我不停的想.
情理上作為一個先知 你救了那麼多人應該燒炮仗.
開了派對.
但他不是 他在發脾氣.
經文說他聽到之後 他知道之後他大大不悅.
沒錯 我們現在去到經文的第四章.
大大不悅.
真的很有趣 我都不明白.
那麼多人下來決自信主 他居然發脾氣.
絕對厲害 絕對是一個好問題.
我再看這段經文第四章 約拿就祈禱.
大家有沒有覺得他開頭講的那段話很熟悉.
那段話是這樣說的.
他說 耶和華 我在本國的時候豈不是這樣說.
我知道你是有恩典 有憐憫的神 不輕易發怒.
有豐盛的慈愛 並且後悔不降所說的災.
這段話很熟悉.
大家想不想起誰講這番話?.
我知道 因為我對這個問題.
剛才所問的問題就深有研究.
原來這句話是在出埃及記第34章6-7時講過的.
那個情況是怎樣呢?.
就是以色列人拜完金牛道.
摩西就弄破了那兩塊石頭.
是法板.
我意思都不清楚.
那個故事繼續的時候.
上帝就叫他再作兩塊石頭上來.
來見我面.
但他一見面的時候.

$^{41}$上帝就跟他說這兩句話.
目的就是想告訴他.
其實上帝本身的位格是怎樣的.
有因點的.
其實若拿在神學上知道上帝是誰.
但你再看下去.
這就很諷刺了.
他就跟著說.
所以我急速逃往他施去.
耶和華現在求你取我的性命吧.
因為我死了比活著還好.
很怪的.
真的很諷刺.
他知這麼多聖經.
我們都不是第一次見.
第二章是這樣的.
感恩串燒歌.
誰知最後就弄了一個最正統的神學宣告.
在以色列人裡面.
其實很奇怪的.
你說他跟上帝不親近.
你又不是.
他又認識這麼多.
他知道上帝有這個位格.
所以才會走人.
但你說他跟上帝不親近.
這個人是非常奇怪.
Mark Sir,你覺得作為一個諷刺文學來說.
這裡算不算是一個大諷刺呢.
我想一定是.
他引用這段聖經或這句的諷刺文是摩西說的.
你又想想.
是先知.
摩西是以色列國裡面最大的先知.
沒有一個像摩西.
到底他像摩西還是不像摩西呢.
我覺得他像反摩西.
Anti-Moses.
Anti-Prophet.
真的很搞笑.

$^{81}$也很諷刺.
然後耶和華就這樣說.
我在經文第四節這樣說.
你這樣發怒合乎你嗎?.
你覺得.
為什麼耶和華會說這麼有趣的話.
耶和華上帝這句問題想帶出什麼呢?.
你聽下去這句話.
其實好像很斯文.
他就說那又沒有理由.
但其實不是這麼簡單.
因為這個字合乎你嗎?.
希伯文的字根叫Yatov.
Yatov這個字在哪裡出現過呢?.
就是在創世記四章七節.
就是哥仁和阿伯他們都去獻祭.
但上帝就不太願意.
因為他們都不太願意獻祭.
於是哥仁就拉長了臉.
上帝看到你這麼不開心.
他就說為什麼你的臉色也變了.
如果你若行得好.
喜不蒙喜悅.
你若行得好.
這個字就是Yatov的字根.
所以其實上帝在約翰書裡面.
他說你合不合理.
其實那個語調是更加嚴重.
不是說你有沒有理由這麼簡單.
其實是你覺得你做不做得好.
又是一個發怒的人.
又是問這個問題.
發怒合乎嗎?.
哥仁那時候就拉長了臉.
現在又起火了.
又說有問題.
剛才看的時候的確很生氣.
剛才我們看到的時候.
哇 起火了.
然後約拿出城了.

$^{121}$在第五節到第八節的時候.
起火了.
然後就坐在城的東邊那裡.
搭了一個棚.
要看情況是怎樣.
基本上就死心不息.
還要看有沒有彎轉.
然後就很有趣了.
上帝就安排了幾件事發生.
又安排了一棵披麻.
在約拿的棚上拉上來.
披麻在當時來說.
雖然不是很清楚是哪一類型植物.
但應該是屬於蘭藤類的植物.
而且葉子是很大的.
所以可以遮得住陽光.
遮住了頭.
然後就有了這個披麻之後.
約拿就大大快樂.
因為減開了它的苦楚.
但第二天.
上帝又安排了一條蟲來咬這條披麻.
以致它枯乾.
最後上帝再安排了炎熱的東風.
和日頭.
約拿的日頭曬到它快要發瘟.
很有趣的.
如果從一個靈性的經文當中.
上帝安排上帝安排上帝安排.
很清楚.
不過不知道怎樣理解這三個字言物呢?.
其實這三種.
剛才Tony所說的.
都是神的創造物.
但我們也留意到這三個創造物的出現.
它們的消失.
它們的生與死.
都是聽命於耶和神.
我們不會驚訝.
因為我們說的諷刺的故事.

$^{161}$之前也有很多創造物走出來.
牛羊會披麻蒙灰.
然後有一條大魚已經出來了.
所以最後又有小動物出來.
但很有趣的.
它自己搭棚子.
打算可以遮陰.
但是要那棵樹.
不是那棵樹.
是那些披麻才可以遮陰.
不知道它搭棚子是否廢了.
很諷刺.
還有我們說聖經是讀出來的.
剛才這三句說話.
神安排 神安排 神安排.
明顯聽眾聽到一隻神.
安排這些創造物的工作.
而這些創造物肯聽.
第八節說.
又在求死.
我死了比活著還好.
很有趣.
因為聖經是聽的.
這句說話和第四章第三節.
有些差不多的說法.
第三節說.
現在求你取我的命吧.
因為死了比活著還好.
如果是聽眾來說.
他會再一次提醒.
又求死了.
又死了.
很有趣.
都說它是先知.
先知求死的有沒有聽過.
以利亞.
最大的影子都在他身上.
誰知他就是一個最大的反對政府.
不要求死就求死.
這又是一個諷刺.

$^{201}$以利亞求死是因為他被追殺.
這個求死是因為他任務完成.
這個我們真的不太理解.
完全是諷刺.
這個求死其實是他不想救人.
是不敢救人.
真的先知完全反轉.
還有一件事.
如果聖經讀的方法是聽的.
你看第九節.
第九節耶和華的問題再出來.
差不多了.
你因者披麻.
法勞合乎你嗎?.
你這樣法勞是否對的?.
又出現了.
如果有聽眾說.
又提醒一件事.
究竟約拿做了這個決定.
這樣法勞合乎你嗎?.
是否對的?.
又出現了.
很有趣.
我們看到有幾個諷刺已經出現了.
又在發脾氣.
這個發脾氣的先知.
最後經文第十節到第十一節.
這個我覺得我們需要去留意.
這個就是耶和華.
去回應這個發脾氣的先知.
經文這樣說.
這披麻不是你栽種的.
也不是你培養的.
一夜發生 一夜乾死.
你尚且外色.
何況這理你未大成.
其中不能分辨左手右手的.
有十二萬多人.
並且許多生畜.
我豈能不外色呢?.

$^{241}$相當之有趣.
其實這個首先講明.
這披麻根本和約拿是沒有關係的.
你都為一棵這樣的披麻去外色.
你今天難道我耶和華上帝.
不可以為你們未成的人.
和生畜來外色他們嗎?.
不過很有趣.
有一個字我也解不通.
十二萬 如何理解呢?.
其實十二在猶太背景是否有意義呢?.
Mark Sir 你會不會知道?.
我想有幾個人.
十二當然是十二之派.
perfect number.
左右手也是.
分別左右 分別善惡.
經常有這句話.
猶太的傳統就是.
一個人去到成人.
knowing good and evil.
knowing right and left.
去到一個負責的年齡.
age of accountability.
變成成人.
他們會讀律法.
他們會在上帝面前負責.
如果你這樣看的話.
這十二萬不能分辨左右手的人.
他們spiritually是children.
他們不認識神.
我聽過這個解釋.
還有沒有其他解釋說十二萬?.
在我的理解當中.
這十二萬左右手都不懂分辨的人.
應該是一班小朋友.
在神的眼中是相當寶貴.
我想我們可以從一個角度去理解.
不論他是小朋友也好.
或者好像Mark Sir所說.

$^{281}$在屬靈上不能分辨左右手的一班成年人.
兩者都有共同特性.
在這裡來說.
他們是未認識耶和華上帝的人.
我想第三個諷刺是什麼呢?.
很有趣的.
不知道大家有沒有想過.
這個punch line是多重的.
耶和華上帝尚且能夠去愛.
這些外邦人.
一些不認識上帝的人.
但是這個若拿仙姿.
竟然不明白耶和華的心腸肺腑.
造成一個反仙姿.
在這個情況下.
去頂撞耶和華上帝.
這是另一個諷刺.
可以說反過來用另一個說法.
這位仙姿寧願看上一棵.
在當時來說.
非常平凡的一棵欄藤植物.
也不看重那12萬人.
和當中在利利尚裡面的生畜.
我想這是一個很大的諷刺.
另外一個諷刺我覺得是.
那12萬人本來不認識耶和華上帝.
但是他願意悔改.
反而這個說得出.
出埃及的三心四象.
他反而是有很多人.
他其實說明了.
他知道上帝會做些什麼.
所以他故意不去跟進.
其實很有趣.
如果我們看完整個故事的完結.
曾經有一個學者將這段約拿書分段.
他說在約拿書來說.
他有一個很有趣的分段.
他分段的.
他說在經文裡面有幾個細節.

$^{321}$斬下來很少人去想的.
第一他說.
他將第一章一節到第二章十節.
為一個大段落.
在這個大段落裡面.
第一章一至三節.
就是上帝和若拿去說話.
然後就是第一章四節到第十六節.
就是若拿和外邦人一起.
第三節就是.
二章一至十節就是.
耶和華再見若拿.
在漁福當中.
然後另外一段第二段.
也是差不多的話.
三章一至三節就是.
耶和華和若拿再呼召他.
然後第三章四節到十節.
就是若拿再和外邦人一起.
而這個今天所讀的經文.
四章一至十一節.
就是耶和華再和若拿說話.
在這個好像很平衡的鏡面裡面.
你卻看到一些很奇怪的.
就是這個若拿在整個故事裡面.
都是很明白耶和華的心想做些什麼.
他不明白.
而且可能我們剛才之前所說.
我記得第一集.
我們說他是小粉紅的先知.
我想耶和華上帝在挑戰他.
我是愛以色列人的仇敵.
尼里美成的人.
不可以嗎?.
我想這個是對當時.
如果是猶太人的讀者.
可能整個呆了.
不知道大家怎麼看呢?.
因為我剛才說點出一些東西.
我們說摩西這個神聖人物.

$^{361}$和喬納也是一個神聖人物.
是有一個平衡的.
我想他會看摩西那段.
就是以色列人犯罪.
但是耶和華就憐憫他們.
這個喬納就覺得.
這就對了.
因為我們是以色列人.
我們是選擇人.
當然我們是值得這個憐憫.
但是他們就會說.
那當然不會.
他們不值得.
他們覺得憐憫永遠都是我們要.
應付我們.
但不會應付我們的鄰居.
或者我們的敵人.
這個也是值得我們反省的一點.
我覺得這個約拿書.
是一個很大的諷刺.
其實這段經文我看的時候.
就給我們一些應用.
我們轉進去.
有一樣東西我看到.
約拿是我們今天的信徒.
尤其是教會領袖或牧師.
一定要留意的.
因為約拿的確在信仰上.
在神學上.
我想沒有人可以質疑他.
但是我們對他最上帝的心.
和對他的信仰行動.
我們有很多疑惑.
在研究一個信仰成長當中.
有三個orthodox是要很清楚的.
一個orthodox是正規的信仰.
一個orthopathy是正規的對神的心態和信服.
第三個orthopractice就是正規的行為.
跟著你的信的緣故.
約拿的orthodox是可以的.

$^{401}$但如果兩個都出問題.
我想這個在我們作為基督徒.
我們真的要留心.
會不會有時候我們信得久.
背聖經很容易的.
我們背了很多.
《流浪上口》.
《流浪上口》神學書.
如果你多讀幾本.
又可以說說卡巴.
又可以說說鮑納夫.
又可以喜歡說John Stork.
J.I.Pekka都可以.
但是不等於我們的生命.
和對神的情感.
就好像有以上.
剛才所說的神學大師那種情操.
這個是我們值得留意的.
或者讀完神學院.
做了教務那些.
他都會在這些上面出問題.
剛才你說的那三樣東西.
其實都會連繫到上帝對他創造的情.
例如在約拿書裡面.
他不只是救了利尼尼城裡面的人.
那些豬狗羊貓全部都救了.
其實要反映.
我們今天做人.
上帝放了我們在這裡.
是不是只是為了我們的快樂.
而罔顧了這個環境.
我覺得這個要反映.
這是一個好點.
因為整本書的結局是很過癮的.
並有很多生畜是這樣結束的.
因為上帝說他很在意利尼城裡面的小孩.
或者宗教的小孩.
分不到左右手的十二萬人.
還有很多那些生物.
他都很在意的.

$^{441}$豬狗羊貓.
因為所有都是屬於神的創造物.
所以剛才丹尼所說的情.
其實都是基於耶和神的愛.
也因為基於耶和神的愛.
才有這個恩典.
整個宗拿四張書都是在說神的恩典.
所有的創造物.
都在神的恩典底下被創造.
也都在神的恩典底下被拯救.
特別是人類.
如果我們這樣看的時候.
其實約拿書.
可以這樣說是耶和華上帝.
展現出他那份恩慈慈愛的一本書.
而在這本書裡面.
正正是反照出先知.
甚至我們剛才給了名字.
反先知的約拿.
他那份的自然.
他那份為了要堅信自己那種.
你可以說他對利利未人那種.
民族情感那種的憎恨.
而抗拒他們.
拒絕他們.
希望他們不進入上帝的恩典範疇裡面.
我想這個是值得我們今天.
去好好思想.
上帝的恩典和慈愛.
是給所有人.
但是會不會有時候我們信一主太久.
我們已經定死了一些人.
是不能夠.
也都不可以接受上帝的恩典.
同時剛才大家都說得很好.
會不會也都同一時間.
我們都沒有留意到原來上帝的愛.
不單止是愛我.
他愛整個他所創造的世界.
如果從一個末世論來講.

$^{481}$最後上帝再回來的時候.
他回來是重新再做.
將整個世界重做.
我想這件事是值得我們做基督徒.
怎樣對我們現在所處身的世界.
都要有一個責任.
這是上帝的國度.
好啊 我們四章看完了.
完沒呢?.
還沒.
我們下一集.
也是最後一集.
分享我們幾個人.
在看約拿書當中.
我們自己的心路歷程和感受.
同時也會介紹一些資源給大家.
希望幫助大家.
進深入去約拿書.
所以下一次和大家進入經緯線.
去看看我們看完約拿書之後.
我們四個人的感受.
和一些資源介紹.
不要錯過了.
拜拜.
\newpage



\section{約拿書}
\label{sec:vU00DNbRAss}
\textbf{經緯線 [釋經視頻] EP7 約拿書 - 研經資料推介}
\newline
\newline
連結: \href{https://youtube.com/watch?v=vU00DNbRAss}{\texttt{https://youtube.com/watch?v=vU00DNbRAss}} ~~~~ 語音日期: 2020-11-27
\newline
\newline
\hyperref[sec:EWFiHkhUrkU]{\small{< < < PREV SERMON < < <}}
~
\hyperref[sec:index_chronic]{\small{[返順時目]}}
~
\hyperref[sec:index_scriptual]{\small{[返順卷目]}}
~
\hyperref[sec:DlruT1gd8EY]{\small{> > > NEXT SERMON > > >}}
\newline
\newline
$^{1}$各位好 《經衛線》 約拿書的研讀.
今天是最後一個的錄影.
相信大家在之前已經聽了我們.
看過我們講過關於約拿書裡面四卷書裡面的分享.
也都很深深的看到原來上帝.
就算是舊約聖經裡面也說上帝無愛嗎?.
在約拿書裡面很突顯出上帝那份的恩典和慈愛.
同時反照出一個約拿.
也是我們的寫照 究竟我們的生命.
是不是離開了上帝的心意.
不明白上帝的慈愛呢?.
我想這個是整本書的重點.
但是為什麼我們還要錄這段視頻呢?.
因為我們都很想跟大家分享.
我們這四個的講員 你可以這麼說.
我們這個心路歷程在錄這段錄影的時候.
這個視頻當中我們的體會和經驗.
我想問一下大家 各位都這麼多次了.
錄影了這麼多次 不知道大家在整個過程裡面.
對於約拿書的內容也好 或者在這個過程裡面.
大家有什麼的領受和體會呢?.
我先說 因為這個經緯線 我開播的位置.
我找了Tony拍檔幫忙 找了兩位嘉賓.
我們四個一起聊.
其實我們經緯線的精神就是.
我們有一些讀經的工具 解經的工具.
其實我們很多時候都沒有學到用.
我覺得有一些可惜的就是.
在教會成長 屬靈的成長裡面.
基督徒很少被人教.
被人叫你讀聖經就有 叫你讀.
但是怎樣去讀呢?就很少說.
How to read.
我們想做一個經緯線.
無論是處理哪些經文也好.
我們覺得有一些工具是可以用的.
其實你看到我們這麼多集.
只是看約拿書一本這樣的書.
從一個文體genre出發.
你覺得是fun to read.

$^{41}$不單是fun to read.
其實都是一些很inspiration reading.
好像趙妖鏡我們說了幾次.
聖經其實是很活著的.
他跟我們說話的.
神藉著他來說話.
所以我覺得這個不單是每一位觀眾.
我們自己也是.
我們去apply這些工具的時候.
我們發現有很多新的東西.
以前我們沒有留意過.
我找Tony李的時候.
我們決定看約拿書開始.
我就碰到一篇文章.
是英文的.
他說這本書是去attack Jewish tribalism.
猶太人自以為是的那種部落主義.
我為什麼這麼吸引到我的注意呢?.
就是我覺得現在在世界裡面.
社會,我們美國國家.
哪裡都是.
有很tribalism的.
Facebook,互聯網,social media.
有一種你不喜歡我就不跟你談.
沒有對話可言的.
一味在罵你罵他.
這個Juner的時候.
我就希望invite每一個讀者.
這本書是否給我們反省這件事.
是很不對勁的.
上帝是否真的這樣去對世界呢?.
我們很多時候都說小粉紅.
是不是上帝真的會這樣看呢?.
耶穌會這樣看嗎?.
我覺得是一個很需要反省的題目.
另外就是在猶太人的傳統.
每逢雍潔泊的時候.
他們都會在會堂synagogue讀Juner的書.
我覺得這是猶太人的智慧.
我們作為神的子民.

$^{81}$我們要看這本書就是看.
我們有沒有repent.
我們有沒有改變.
這個對我來說是一個很大的反省.
很坎坷.
雍潔泊是不是整天都在縮短時間內?.
是,9月,10月左右.
其實你講得很好.
Mark Sir 我記得曾經有人引用過林肯總統.
就是當年打南北戰爭的時候.
有一個故事是有一個將軍走來.
因為南北戰爭是北面贏的.
那個將軍就說.
God is in my side.
在我們的side那邊.
但是林肯總統就說.
不是神在我們的side.
是我們在神那邊的side.
是很不同的.
有時候你說得對.
那個部落主義就是說.
只有我對別人是錯.
也都是有一個情況就是說.
去到一個地步就是沒有了一個外.
沒有了一個.
我們可以代替上帝說話.
將上帝那份慈愛抹走了.
這個也是一個很好的提醒.
這個又是你回到讀經那裡.
我覺得最恐怖的就是.
In the name of God.
We can do a lot of things.
我們可以intend很多的東西.
In the name of God.
In the name of religion.
這本書我們不要只當它是一個人物的處理.
因為主要真的不是講一個人物.
是講一群人.
是講一種主義或是一種ideology.
一些很powerful很dangerous的ideology.

$^{121}$因為我覺得這本書就變成了.
不會像平時那種感覺.
小孩子看的書很有趣很過癮.
其實是在address我們很多的issues.
political issues.
identity issues.
nationalism.
Whatever that is.
我們面對差不多所有的issue.
都可以relate to the Jonah inside of us.
非常認同Jonah是非常有深度的一本書.
是一本故事.
我自己的體驗就是.
人需要去珍惜神的創造.
特別是人類的生命.
在Jonah裡面我發現.
先知經常都求死.
同時他也拒絕耶和神的呼召.
也不肯去救其他的人.
特別是利美城裡面的一群.
左右手都未必分到的12萬的人.
其實看到這件事有一個很大的對比.
對比就是神的愛.
他反而很珍惜他的創造物.
特別是我們人類.
他可以將自己獨生的兒子耶穌降世為人.
釘死十格.
救我們這一群的罪人.
所以這也是我這次去讀Jonah的時候.
一個新的體驗.
我很echo Gareth的Ecological Concern.
這是我這次第六第四章的時候才發現.
有很多創造與神之間的互動.
神很憐憫他整個的創造.
一個ecological的compassion.
我其實跟Mark Sir的看法有點相似.
大家看我就知道我喜歡跳機.
因為你在教會看到很多時候.
有些是既定的idea.
你經常都密不密的.

$^{161}$但你很明顯是錯的.
你看到我們幾個好像在揶揄.
其實我們都花了很多時間去unlearn以前學到的東西.
等於張無忌被人廢了武功才可以學到新的東西.
差不多這樣的東西.
還有我們這次展示給大家的tools.
其實都不是一些很state of the art的東西.
我們會quote一下.
做節目都要扮一下東西.
會quote下希伯來文.
quote下原文.
但其實最主要的就是.
你懂不懂得分析那一段文字.
而裡面的tools不是說你要有大學程度才可以做到的.
其實你中學程度就可以做得到.
至於你說self reflection.
我都是同意Mark Sir的.
那個tribalism的問題.
我覺得在華人教會界.
他很盛行一種dispensationalism的東西.
就是凡是以色列你都會支持的.
如果你去看新聞的時候.
你會看到巴勒斯坦地區會有很多很亂的事情在發生.
在教會裡面他只會說這個是上帝的應許.
所以以色列人就要復國.
但我覺得在這句話背後其實忽略了巴勒斯坦人的辛酸.
就好像以前的17th century,18th century.
他有一個廢除黑奴的制度.
其實他有一個symbol,一個plaque.
寫著一個黑奴,一個黑人.
他被鎖住手鐐和腳鐐.
而下面的行字寫著.
Am I not a man and a brother?.
我覺得這個我們要反思.
因為如果我們忘記了He is also a man and He is also a brother.
你就忽略了上帝對人的慈愛.
和剛才也說過對他創造的慈愛.
我想這件事其實都是我們值得去反省的.
我剛剛記得早前我看一段新聞.
荷蘭的教會為了在第二次世界大戰的時候.

$^{201}$他們大部分的教會或者基督徒.
沒有為了當時在荷蘭的猶太人站起來.
以致到按統計有七成的有荷蘭籍的猶太人死在集中營裡.
當時來說只有極少數的荷蘭的基督徒.
去保護這些猶太人.
有很多都是按照法例去做.
告訴蓋世太會聽.
那些猶太人在哪裡,就把他們都抓去集中營.
我想這個其實也是我們基督徒要很小心去思考的.
剛才大家都提過.
會不會我們已經將一批人定死了.
他不會也不能領受上帝的恩典.
這件事我們要很小心的.
我們可以自限一個章節.
舉例有一首詩歌我經常聽完.
我經常有一個問號.
宣教在中國.
先講明一件事.
我是很希望中國的人信耶穌的.
我一定要講明這件事.
我自己也很支持這件事.
我會問為什麼只有一首宣教在中國.
為什麼沒有宣教在荷蘭.
宣教在日本.
宣教在世界的歌.
一首不重要.
為什麼只唱一首而沒有所有的關注呢?.
當然你說要集中正常.
OK的我理解.
但是那個不能夠成為了一個理由.
而去忽視了原來世界上有其他人.
其他的群體都需要耶穌基督.
這件事是我們華人有一種主義.
很容易就是我只顧著自己的族群.
不是不要你顧著自己的族群.
但是如果只是這樣就有問題.
甚至去到一個地步.
其他的他們信不信耶穌我不管.
或者他們有沒有領受上帝的救恩我不管.
又或者另一個想法令我想到.

$^{241}$當教會我們在.
你可以說我們用很多的預算建建築物.
建很多的運動場.
沒問題的.
真的沒問題.
但是我建完之後.
只是為了自己會友會員的使用.
而從來沒有想過.
這個建築物是怎樣可以成為.
那個社區裡面任何階層.
任何種族的人的祝福的時候.
這個就是我們值得去問的問題.
是不是他們不能夠去.
不值得受上帝的恩賜.
是不是他們不可以接受上帝.
這是我讀這本書想到的東西.
還有另一件事我覺得.
在整個讀書過程裡面.
其實我們是一個色經的群體.
在座當中的弟兄姊妹朋友.
你看我們視頻是十多二十分鐘.
但是我告訴你.
在我們背後不止十多二十分鐘.
可能是很多個小時.
台前一分鐘.
台前一分鐘台後很多時候.
但是這個是好的.
為什麼呢?.
其實剛才Ben和Benny說的有點相似.
就是當我們四個都是帶著不同背景的基督徒.
牧者或者評述一起去查經的時候.
當我們打開心懷.
用一些適當的工具的時候.
其實我們是對這段經文是越來越有一定的掌握.
但是同一時間很吊詭.
又越來越謙卑.
因為我們發現我們所知道的.
其實是有限的.
但是當我們這樣去在這個過程裡面.
我們越來越成長.

$^{281}$我們可以這樣說.
可以說我們今天之後.
不等於我們會停止研究約拿書.
可能更加是挑戰我們繼續研究約拿書.
不過我們是怎麼樣呢?.
是confidence with humility.
就是我們有一個信心.
就是我們知道我們花了時間去學習.
有神的帶領.
有聖經的帶領.
也是用那個適當的工具.
同一時間我們是有一個謙卑.
因為我們知道我們所知道的其實是很有限的.
我聽到同一件事.
我們不是在找對答案.
因為我覺得這個比較難.
但是我們是在找一個最合理的答案.
就是至少在我們這個時刻.
我們懂得的東西裡面.
我們怎樣去找一個比較合理的答案.
我自己可能是比較傳統的一個解經的.
或者是識經或者是查經的一個老信徒.
始終OIA都是離不開的.
還有就是永遠都是.
我們這麼有限.
我們都是一個被創造之物.
我們要去了解和認識我們的創造主.
一個無限的主的時候.
唯獨是我們是要謙卑在他的面前.
去求這個智慧去明白.
同時我也相信.
我個人相信.
上帝有喜可不會將他自己.
來躺開給我們去認識呢?.
或者用另一個說法就是.
我們謙卑下來的時候.
我們去更加理解上帝的真理.
我們有個信心.
但是同一時間我們要謙卑.
因為知道可能我們過了若干而再查.

$^{321}$若然有些新的體會.
可能修正我們過往某一些.
不是很清楚的看法都未定.
但是有些我們清楚了.
我們肯定的.
我想這件事我覺得在這裡.
我自己個人一個學到.
還有我相信這件事.
其實給我們個人在這個言經當中.
在若拿給我的很大的提醒就是.
我想你會看到我時很針對.
就是究竟一個人除了明白信仰的真理之外.
他對宗教的情感和所謂的信仰的行為.
三者的結合.
在若拿當中成為一個照妖鏡.
給我們嘗試去看.
究竟我站在神面前.
究竟我所信的道.
我對神的那份情感.
和我對神那種行動的回應.
是不是真的反映出上帝那份恩賜.
我想約拿書是給我們.
我同意你所說的.
我們不用翁潔篷我們常常都要看.
就是反省自己.
使我們看到我們所行所做的.
都是反映出上帝那份的慈愛.
我可不可以多說一件事.
最後一集.
我比較情緒化一點.
我想說回.
譬如我們在教會裡面傳福音.
我發覺我們很多時候都有一個既定的模式在.
但我們常常都在追求多人.
但其實我們很少去練習那個情緒.
如果你看回教會歷史.
譬如剛才說過的十七十八世紀.
那時候是有一段復興的時候.
為什麼復興呢.
就是因為他真的練習了那個情緒.

$^{361}$譬如解除了奴隸制的問題.
另外一樣就是我們常常放在眼裡.
我們十二年免費教育.
十二年免費教育就是那段時間出來的.
也出了很多很經典的.
John Newton is my hero.
譬如那首Amazing Grace.
就是那段時間出來的.
是一個罪人自己審查了自己.
然後寫出來的一個詩歌.
這樣都是好的.
其實讓我們今天去回望.
不知道大家的聽眾或者觀眾.
視頻觀眾你處身在什麼地域當中.
我相信你將它應用回剛才我們所說的.
可能你想到很多的思考.
但我相信就是.
不論你是處身在哪一個國度都好.
容許我們去問自己.
究竟我們的心思維.
是不是反映出上帝那份的恩慈.
好像若拿一樣.
口是先知.
但其實所做的一切是反先知.
我想這個是值得我們去思考的.
到最後我們都有一些資源去介紹給大家.
以致幫助大家在看約拿書.
能夠可以在家裡也好.
或者你和弟兄姊妹一起看的時候.
能夠更加盡心.
我不如Benny你先介紹好嗎?.
好,我想介紹一本很普遍的工具書.
就是Gordon Fee寫的.
叫做《How to read a Bible for all its worth》.
他也挑機的.
他在這本書裡面說.
釋經的竅門是如何去正確分析這段文字.
他說你未看這本書之前.
你應該看一本書叫做《How to read a book》.
怎樣看一本書.

$^{401}$因為他說.
例如在第一集我們說過文體是很重要的東西.
他整本書就是教你怎樣從一個文體,文學分析的角度去分析經文.
裡面很少提到希伯來文,希臘文.
只不過是純粹教你.
或者磨擦你.
令你有基本的閱讀能力.
Benny,我讀PhD的時候.
第一本書就是要讀《How to read a book》.
然後我在神學院讀的第一本釋經書就是《How to read a Bible for all its worth》.
Good choice.
黃師母你有沒有什麼介紹呢?.
一本翻譯書.
它的英文是《Every man's Bible Commentary》.
這本翻譯書.
不知道大家有沒有見過.
另一本書可能大家都會有的.
是《十二》的先知書主義.
是唐柚子牧師所著的.
這兩本書基本上.
我也有參考過一些舊約的精纜等等.
大家不妨可以去看一下背景.
從舊約裡面去看一下.
我自己也有好用幾本的.
我就介紹一本.
這本我覺得是因為它是很新的.
這本書叫做Jonah God Sandler Mercy.
寫的人就是Kevin Youngblood.
Kevin Youngblood其實現在是美國哈頓大學的聖經科教授.
他自己的PhD是在美南浸信會神學院的PhD.
如果你知道他學校的背景.
你也知道他所說的話也不會走到哪裡去.
但是我喜歡他一件事.
雖然他的信仰是保守的.
但是他的入手點是非常新.
他處理很多文體文學上的理論.
還有他處理整個書的結構.
這本書的英文不難看的.
尤其是對於一些死很害怕的.
英文很難.

$^{441}$其實英文不是很難.
還有很重要的.
有很多圖畫的.
有很多的chart幫大家.
我覺得對於一個入門的朋友.
尤其是我們的視頻是廣東話的視頻.
可能當中有些朋友.
尤其是在香港或者其它地方.
可能英語不是你最強的.
但是你都需要一些新的書.
又容易進入入手點的.
我覺得這本是一本挺好的入門書.
而且這本書其實出版了不久.
2013年左右.
大約都是7年左右的書.
所以他的學術方面.
學術方面都追得比較近一點.
更加重要.
你引用這本書.
你不會被人說是異端.
這個很重要.
我的經驗.
Mark Sir你呢.
我介紹兩本都不是方牌的書.
不代表是異端.
這個我們想挑戰大家.
我們成會教會沒有這樣的包袱.
不一定要看保守.
保守不代表你是對的.
若拿也保守.
大錯特錯.
我介紹一本是猶太人的識經書.
我經常預備聖經的研究.
或者講道都會用猶太人的識經書.
我怎樣都要看一看.
這本就叫JPS.
Jewish Publication Society出的一個系列的舊約.
希伯來聖經的識經書.
這本就是Jona.
所以這些猶太人的傳統.

$^{481}$我覺得我們應該去欣賞.
他們有些識經的智慧等等.
另外一本就是一個Families Theologian寫的.
叫做Philip Trouble.
Philip Trouble已經去世了.
他是一個表表者.
用的工具去識經就是用修辭學.
所以這本書叫做Rhetorical Criticism.
Context Method in the Book of Jona.
他就將修辭的批判學.
就應用在約拿書.
我們城匯網台.
特別我都經常送大包.
Philip Trouble就將整個的約拿書.
將對稱的結構.
做了一個表出來.
我就將它翻譯成中文.
如果你在下面的描述欄填那張表.
我會給你連結.
你就可以拿到我的表.
在PowerPoint的表格.
並且要訂閱我們的網台.
將這麼好的工具送給你.
非常好.
有一件事要補充.
剛才那本How to Read the Bible.
其實它有中文版.
大家可以買中文版.
如果英文不沙家的話.
我們都很感恩.
其實這個約拿書的完結.
但是還沒有完結.
為什麼還沒有完結呢?.
因為經緯線是繼續有的.
我們就等下一輯經緯線出來的時候.
可能有不同的主持.
有不同的組合.
有不同的書卷.
但是目標仍然沒有變.
希望幫助大家用合切的工具.

$^{521}$去閱讀神的話語.
以至這個神的話語.
透過聖靈幫助我們.
去更加明白聖經.
以及在生活上去反省,應用.
所以各位.
暫時暫別.
但是等下一輯經緯線的時候.
希望大家繼續參與.
在經緯線當中.
明白上帝的真理.
拜拜.
\newpage



\section{}
\label{sec:DlruT1gd8EY}
\textbf{浩劫重生 【第7集:流亡被擄的以西結】 陳崇基牧師 - 主日講道視頻系列 7}
\newline
\newline
連結: \href{https://youtube.com/watch?v=DlruT1gd8EY}{\texttt{https://youtube.com/watch?v=DlruT1gd8EY}} ~~~~ 語音日期: 2020-12-02
\newline
\newline
\hyperref[sec:vU00DNbRAss]{\small{< < < PREV SERMON < < <}}
~
\hyperref[sec:index_chronic]{\small{[返順時目]}}
~
\hyperref[sec:index_scriptual]{\small{[返順卷目]}}
~
\hyperref[sec:h_zZvugRB1g]{\small{> > > NEXT SERMON > > >}}
\newline
\newline
$^{1}$每當我們想起一些浩劫的時候.
我們自然就會想起.
在我們世界的政局裡面有不穩.
或者可能近身一些.
在我們生活當中的狀況.
很多的改變.
或是經濟或是關係.
或者甚至是我們自己命途的前途.
觸摸不到的很多未知之數.
聖經每每提到一種類似情況的時候.
其實是講到上帝會做一些新的事情.
今天我們要去看一個人.
他是一個先知.
Prophet.
但是我們首先要搞清楚什麼叫Prophet.
什麼叫先知.
先知並不是拿著水晶球.
然後去微卜先知.
知道明天預測什麼時候會下雨.
真正的先知.
其實就是神的代言人.
是代替上帝說話的.
每一個真的先知.
就是接受了上帝對他的啟示.
應用在他自己的生命當中.
並且在危機的時候.
在轉變的時候.
他代替上帝向其他人宣告神的心意.
今天我們看這個先知.
他叫做以西結.
Ezekiel.
以西結書第一章.
當三十年四月初五日.
以西結在加巴魯河邊避老的人中.
天就開了.
得見神的異象.
正是約瓦根王避老去.
第五年四月初五日.
在加勒底人之地.
加巴魯河邊.

$^{41}$耶和華的話特特臨到布西的兒子.
祭司以西結.
耶和華的靈降在他身上.
很多先知書開始的時候.
他首先會提到幾樣東西.
第一是地點.
這裡提到以西結.
他身處在加勒底人.
即是巴比倫人之地.
加巴魯河邊.
這個地方其實是位於今天的伊拉克.
尼泊爾尼泊爾的古城.
距離當時巴比倫帝國京城.
東南大概六十英里.
是一個外邦人的地方.
我們去看這個遭遇的時候.
我們要想起當時猶太人在想什麼.
對於以西結來說.
猶太人認為一個外邦人的地方.
當然就是一個不潔的地方.
一個遠離上帝的地方.
當時發生什麼事呢.
猶太人經歷了一個國家個人的浩劫.
因為當時巴比倫軍隊來到他們的國家.
然後將他們打敗.
猶太人流亡到異鄉.
很多時候他們在異鄉當中.
就會去到河邊聚集.
因為既然沒有聖殿.
他們就去有活水的地方.
在那裡敬拜.
在那裡祈禱.
提到地點.
也提到時間.
當時的時間就是約瓦根王被擄的第五年.
我們運算出來的時候.
我們知道這是公元前593年.
這王記下最後那兩章.
二十四二十五章.
就提到這個事蹟.

$^{81}$猶大王約瓦根和他母親.
神僕領袖太監一同出城.
投降巴比倫王.
巴比倫王就拿著他.
那時是巴比倫王第八年.
被擄之後.
巴比倫王就另立一個人來做猶大王.
那個人叫做西底加.
這裡提到約瓦根王被擄第五年.
意思就是說.
以西結是在那個被擄流亡的經歷裡面.
其中的一分子.
就是說他現在已經身處巴比倫有五年的時間.
當然也是非常近耶路撒冷被滅亡.
亡國的時候.
這裡也提到.
第三十年.
經文沒有詳細解釋.
到底這三十年是什麼意思.
或者這三十年.
倒數推翻前三十年的話.
大約是公元前632年.
當時是約西瓦王.
做西人.
去做宗教改革的時候.
是那段最輝煌的歷史.
也是最後一次去做宗教的改革.
希望能夠挽救上帝的國度.
挽救猶大國.
以至猶大國不會滅亡.
但當然我們知道最後是失敗的.
可能三十年是這樣的意思.
也可能三十年是指到.
當時近東的國際秩序.
因為如果運算的話.
公元前632年.
當時是亞述帝國.
最後一個大帝身亡的時候.
後世人稱艾什勃·貝爾尼波.
是亞述末代的大帝.

$^{121}$當他陣亡之後.
亞述的帝國就開始正式沒落.
然後後起來統治天下的.
就是巴比倫帝國.
可能三十年就是在說.
那種朝政上改朝換代.
天下大亂的之後三十年.
或者也有可能.
三十年是指到以西結本身的年齡.
當時他只是三十歲.
為什麼這樣說呢.
因為按照民數記四章.
祭司在三十歲的時候.
開始在耶路撒冷聖殿去事奉.
無論如何對以西結來說.
國際的政局交錯.
民族的命途.
個人的前途.
都是在經歷翻天覆地.
在這個時候.
上帝的話就臨到他.
臨到一位本身是做祭司的.
他現在流亡於異地.
失去了聖殿.
失去了他本身應該事奉的地方.
在這裡他就聽見上帝的話.
你可以看到四月初五日.
大概是今天西曆的七月八月.
我們知道在米索不達米亞巴比倫.
就是現在的伊拉克一帶.
夏天經常會有一些很大的沙漠風暴.
以下一個短片.
可以讓你感受一下.
這種沙漠風暴的情景.
無論怎樣也好.
以西結告訴我們.
他當時身處的時代.
就是一個國際政局交錯的時候.
民族以至個人的命途都翻天覆地.
但在這個時候.

$^{161}$上帝的話就臨到他.
一個本身是祭司.
流亡異地的以西結.
第一章四節.
見狂風從北邊刮來.
隨著一朵包括閃爍火的大雲.
從其中的火裡.
發出好像光耀的正感.
又從其中顯出四個活物的形狀.
他們的形狀是這樣.
有人的形象.
各有四個臉面.
四個翅膀.
他們的腿是直的.
腳掌好像牛獨肢體.
都燦爛如光明的銅.
這四個活物的臉和翅膀乃是這樣.
翅膀彼此相接.
行走並不轉身.
軀角直往前行.
至於臉的形象.
前面各有人的臉.
右面各有獅子的臉.
左面各有牛的臉.
後面各有鷹的臉.
他們軀角直往前行.
行走並不轉身.
至於四活物的形象.
就如燒著火炭的形狀.
火把的形狀.
火在四活物中間.
上去下來.
有火的光輝.
從火中發出閃電.
這活物往內奔走.
好像電光一閃.
我們看完這些描述的時候.
我們現代人會覺得.
是一種科幻電影.
是很難明白的.

$^{201}$其實對現代人.
甚至對古代人來說.
都是非常難明白的.
因為以前曾經有一些猶太人的拉比.
他說如果你未夠三十歲的話.
你就不可以讀以西結書.
因為以西結書裡面.
有記載了很多很難明白的異象.
所以我要解釋一下.
以西結書提到的異象.
提到有四活物.
這些Living Creature.
這些活物.
各有四個臉.
有人.
獅子.
牛.
和鷹的臉.
我們想起就會覺得是一些怪獸.
一些怪物.
但是當時近東世界.
那種圖像的觀念.
很多時候會將不同地上的活物.
混雜在一起.
在這裡我們看到.
在敘利亞巴勒斯坦地所找到的.
皇宮裡面的象牙雕塑.
你看到有的活物.
就是有人的臉.
牛的身.
獅子.
牛.
不同的組合.
也有翅膀.
在這裡我們看到.
亞述帝國的皇宮門口.
很多時候有一些守門口的活物.
叫做Lamasu.
也有人的頭.
牛的身.

$^{241}$獅子的腳.
鷹的翅膀.
很多時候這些Lamasu.
也會有五隻腳.
在側面看的時候.
你會以為牠在行走.
在正面看.
你看到牠是直立.
牠也看守著.
守衛著.
很多時候近東的圖像.
圖畫.
他們會將這些不同國方的活物的能力.
混雜在一起.
就是想表達霸權的權威.
這裡看到四個活物的形狀.
四塊的臉.
首先有人.
人在近東的世界裡面.
甚至我們現今的世界裡面.
都是屬於世界裡面最高的活物.
所以我們叫他做萬物之靈.
是不是?.
另外有獅子的臉.
就是野獸裡面最高最大的.
另外有牛的臉.
就是家畜類裡面.
domesticated animal.
裡面是最高最有能力的.
最後就是鷹.
就是在鳥獸當中作王.
所以將這些臉面.
將這些形狀配合在一起的時候.
雖然我們覺得會很怪.
但其實他最主要想表達的.
就是那種能力.
那種權力.
以色列看到這個異象的時候.
他也嚇了一跳.
然後我們知道.

$^{281}$在以色列書繼續看下去.
他再看到類似的異象的時候.
他就澄清了.
他就明白.
原來這些活物是什麼呢?.
原來是基魯柏.
以色列書第十章第二十節.
他說我在加巴魯河邊所見的.
以色列神榮耀以下的活物.
我就知道他們是基魯柏.
四活物就是天使基魯柏.
你記得基魯柏在聖經裡面.
是看守守衛聖殿之聖所的.
他是守住約桂的.
而約桂就是耶和華神的王座.
歷代之箱第十八章十八節.
這裡是約桂的裝飾.
有精金.
然後有金紙做的基魯柏.
基魯柏的翅膀是張開的.
然後他們是遮蓋著耶和華的約桂.
其實這裡的原文意思.
就是說基魯柏的金箭射.
Golden Chariot.
就是約桂其實是神的寶座.
也是上帝的戰車.
原來以色列所見到的.
就是他作為祭司.
應該在耶路撒冷聖殿之聖所見到的約桂.
見到上帝的座駕.
見到上帝的戰車.
因為上帝是乘著基魯柏天使.
來出戰的一位神.
講完四活物之後.
他又繼續描述在這個意象裡面.
見到的戰車.
一章十十五節.
有一段很難明白的意境.
我相信這裡其中一句最難翻譯的.
就是輪中套論.

$^{321}$我相信NIV的譯法是最貼切.
Wheel intersecting wheel.
就是耶和華的戰車.
你可以想像是一個車輪.
然後在九十度的垂直.
又有另外一個車輪.
在輪裡面再包著另外一個車輪.
換句話說.
想表達的就是.
這個戰車是隨意前後左右上下.
是不需要調頭的.
一個輪中套論.
代表著耶和華的戰車.
耶和華自己是可以無所不在.
四個輪滿有眼睛.
這裡是在描述耶和華的戰車.
他的車輛是無所不知的.
就是說上帝知道了所有事情.
並且他是無處不在.
你是不能夠逃避的.
你留意在整段的意象裡面.
四字是經常出現的.
地的四極.
神是可以去到地的周圍.
不同的地方.
上帝不受任何的空間地理所限制.
形容完戰車了.
我們越來越清楚看到的.
就是車上到底有哪一個.
第二十二字.
活物的頭以上有窮倉的形象.
好像可畏的水晶.
活物行走的時候.
我聽見翅膀的響聲.
像大水的聲音.
像全能者的聲音.
有一像軍隊的聲音.
窮倉之上有寶座.
寶座以上有彷彿人的形狀.
我看到他腰下面.

$^{361}$好像晶金一樣的耀眼.
二十八字.
下雨的時候.
雲中彩虹的形狀怎樣.
這個形狀就是這樣.
然後他說得很清楚.
他說這就是耶和華榮耀的形象.
我一見到就俯伏在地.
又聽見說話的聲音.
越來越看清楚.
原來他見到的就是榮耀的神.
在寶座之上.
在他的座駕裡面.
他御駕親征來到他的宗教.
一個現代的比喻.
就是今天美國總統.
所乘搭的飛機.
空軍一號Air Force One.
大家知道.
總統上了那架飛機.
那架飛機就可以叫做空軍一號Air Force One.
換句話說.
這裡說的是耶和華.
他榮耀的出現.
他坐著他的戰車出戰.
如果比較一下.
舊約裡面經常提到.
上帝榮耀出現的意象.
你就可能會更加清楚.
譬如以賽亞.
在大六章的時候.
以賽亞在聖殿敬拜的時候.
他也見到上帝榮耀的意象.
讀章第一節.
他說當烏西亞黃崩那年.
我見到主坐在高高的寶座上.
他的衣裳睡下遮滿聖殿.
我們見到這個意象裡面.
以賽亞想像.
上帝是一個很大很大的形象.

$^{401}$他坐在寶座上來掌權.
又或者你可以比較一下.
所羅門王.
來到獻第一聖殿的時候.
神的榮耀也都降臨.
《耀王記》上第八章.
祭司將耶和華的約鬼.
抬進去內殿.
就是致勝所.
Holy of Holies.
放在兩個基督派.
Cherubim.
的翅膀底下.
祭司從聖殿出來之後.
有雲彩充滿耶和華的殿.
甚至祭司不能夠站立公職.
因為耶和華的榮光.
是充滿著這個殿.
寶座上的上帝.
是很威嚴.
非常榮耀的景象.
甚至今天的猶太人.
他們最嚮往的情景.
其實都是舊約裡面.
第一聖殿獻殿的時候.
神的榮耀充滿著他的殿.
在耶路撒冷裡面.
充滿著他自己的威榮.
如果你這樣比較的時候.
你就更加能夠知道.
以色列所見的榮耀.
是有多大.
因為最大的分別.
就是我們剛才所說的.
那些異象.
看到神榮耀的異象.
所羅門聖殿獻殿的時候.
以賽亞在敬拜裡面.
看到高高在上的寶座.
這些神的榮耀的彰顯.

$^{441}$全部都是靜態.
是static.
以西結在加巴魯河邊.
在外邦人被擄的人中.
他所見到的榮耀異象.
就完全是暫時.
因為是動態.
輪中套輪.
並且行走不需要倒轉.
他的榮耀的出現.
最令被流亡於巴比倫的猶太人驚訝.
就是上帝居然.
在我們被擄的人中間.
以西結和其他的猶太人.
在巴比倫異鄉當中.
過著一個流亡被擄的生活.
什麼叫做exile.
什麼叫做流亡.
我們還沒有給他一個定義.
我相信最好的定義.
就是今天在洛杉磯.
教希伯來聖經的一個女猶太學家.
她叫做Tamara Eskenazi.
她為流亡做了一個這樣的描述.
她說什麼叫做exile.
流亡不單單是無家可歸.
相反你本身是有家.
但你的家園被敵人強佔.
exile流亡不僅僅是無根.
相反你本身根深蒂固在一個地方.
而現在卻被連根拔起.
被迫活在異地.
你是暴露無遺.
在痛苦的寒夜中.
被人恥笑的當中.
你很想回到原本原根的土壤.
但因為現在流亡.
你知道你所屬的故土.
但卻深知不能回到那裡.
這就是猶太人流亡的意思.

$^{481}$一種最難受的感覺.
想起自己的家.
想起自己的belonging.
自己所屬的地方.
卻回不去.
我相信我們想起猶太人被露的經歷.
我們也能夠認同.
因為今天很多人都有類似的感受.
那種感覺就像猶太人生活在異鄉一樣.
你問問你自己.
你現在身處的那個地方.
你住的那個地方.
你自己最關心的新聞.
消息,故事或政局.
是甚麼呢?.
很多時候不是你住的那個地方你最關心.
而是你本族,本家.
你自己本身根源的地方.
而你外在地身處一個流亡的地方.
又或者想想另一個情況.
就是你本身在這個地方土生土長.
但你見到這個城市.
土生土長的社會是不斷改變.
與以前不同.
突然醒過來的時候.
你覺得這個地方很陌生.
你經常覺得這個不是我熟悉的地方.
你很想走.
或者你很不熟悉這個.
我一向很認識的地方.
以前很嚮往這裡的生活.
但突然醒過來.
你覺得這個地方變成很陌生.
不知道如何自處.
又或者你可能想想.
一種更加難受的流亡經歷.
就是我們人生的內在.
我們受到一些情緒的打擊和困擾.
無論是一些工作的改變.
因為疫情,困難.

$^{521}$我們的生活大受打擊.
或者是一些關係上的失去.
突然的消逝.
無論是哪一種流亡的經歷.
外在的也好.
內在的也好.
或者內外一起來的.
我們都會問到底怎麼辦.
以西結他見到的異象.
就是能給出一個答案.
就是告訴以西結和猶太人.
他們在流亡,被擄當中.
原來上帝可以跟他們同在.
上帝的榮耀.
不會被困於耶路撒冷的聖殿.
不受地理,時間,空間的限制.
他居然來到加巴魯河邊.
一個不潔的地方.
被擄,流亡的地方.
跟他的子民同在.
流亡的時候我們怎麼辦.
很多時候我們會想起.
或者很夢想.
我們去了一個更好的地方.
這個沒錯.
人之常情.
很多時候我們想脫離現在的處境.
我們去了一個新的處境.
那就好了.
以西結書這個異象告訴我們.
最重要的.
就是上帝到底是否跟你同在.
如果上帝能夠跟你同在的話.
你縱使在流亡當中.
你更加會見到神自己的榮耀.
神自己的實在.
以西結的異象最令猶太人驚訝的地方.
不單是上帝的榮耀彰顯.
並且上帝的榮耀離開耶路撒冷.
自己的城鎖,自己的聖殿.

$^{561}$去到外邦人中間.
如果你看看以西結書的內容.
你就可以知道.
神為什麼會離開自己的城鎖.
自己命所居住的地方.
就是因為祂討厭猶太的國民.
討厭耶路撒冷的人.
他們知道聖殿在他們中間.
但是他們的心不是敬拜耶和華.
他們敬拜的是必神和很多的偶像.
他們把聖殿當成一個建築物.
成為他們的平安符.
他們心裡面想.
有聖殿的話.
我們就會有保證.
有聖殿的話.
耶路撒冷就不會有外敵來入侵.
有聖殿的話.
耶和華神的榮耀一定在我們中間.
這個所謂舊約裡面的錫安神河.
Zion Theology.
那種意念就是.
總之有聖殿在我們裡面.
我們就有上帝來保護.
就算現在流亡.
就算現在是被擄.
在以西結的人民當中.
也有些假先知不斷有說.
說我們的流亡只是暫時.
很快我們就會回去.
我們被擄的時間很快就會結束.
以西結領受了上帝的話.
他說.
你千萬不要自己騙自己.
如果你繼續.
你去敬拜偶像的話.
不以上帝為你的主.
不以神的道為你生活的中心.
上帝一定會離開你.
這些外在有的建築物.

$^{601}$或者有的移民.
是不會保證到神與你同在的.
以西結對人民的指責.
很多時候在書中找到一句話.
叫做可憎的事.
Abominations.
經常出現的.
當然他就是說到.
人民不看上帝在眼內.
不以耶和華為主.
這是一個很大的警告.
就是告訴我們.
到底我們是不是很多時候.
好像猶大國的國民.
很多時候我們可能都會這樣想.
如果我們回到以前的生活就會.
如果我們有以前.
能夠去碰到的東西.
無論是一個教堂的建築物.
無論是一些移民.
我們就可以好好地去敬拜神.
只要我們不是像現在這樣流亡.
這種難受的經歷當中.
我們回到以前那樣.
我們就好了.
就等於猶太人.
去到問一個問題.
他說只要我們能夠回到耶路撒冷.
有那些宗教的外在記號.
我們就可以有神在我們中間.
以西結書說.
不是這樣的.
神不會在我們中間.
如果我們的心一直不歸他.
我們仍然是三心兩意.
仍然不想歸向神的話.
不以他的道為我們生活的中心的話.
上帝一樣可以離開.
現代的例子.
比比皆是.

$^{641}$你看看歐洲有很多宏偉的教堂.
神在不在中間?.
當然不是.
你看看以往有很多宗教的符號.
宗教的建制.
上帝在不在中間?.
不在中間.
神的榮耀不會因為人的制度.
不會因人所做的這些institution而留下.
相反他看到的.
就是問我們的心是否真正歸他.
以西結書想說什麼?.
他說相反.
如果就算我們是被擄.
是被流亡的時候.
雖然我們不見到有君王在上.
不見到有我們平時的朝廷.
不見到有聖殿.
祭司都不能夠上班.
神在不在我們中間?.
以西結說.
只要我們是專處偉大.
我們是按上帝的道行事為.
相反我們會看到神的榮耀臨到.
神的榮耀同在.
縱使在流亡當中.
只要有上帝和我們在一起.
只要我們做一個不折不扣的神的子民.
我們說真話.
說上帝的道.
哀上帝的道.
恨上帝的道.
神就會來到我們中間.
一章二十八字.
這是二將最後的一節.
他說這就是耶和華榮耀的形象.
我一見到就俯伏在地敬拜.
聽見一位說話的聲音.
之後就是以西結被呼召的故事.
他接受神的書卷.

$^{681}$在被擄的人當中.
神的說話.
不單單是幻想回到以前.
現在縱使在流亡當中.
也接受上帝的說話.
按照上帝的道來生活.
為上帝做代言.
到底我們今天的信徒.
在不同的浩劫當中.
在疫情當中.
到底我們是否以上帝為我們的重心.
求主幫助我們.
無論我們是哪一個境況.
我們渴想有上帝的同在.
並且我們知道.
我們一定要行在上帝的道中.
上帝的榮耀才會臨到我們.
無論我們處於什麼境況.
教會處於什麼境況.
讓我們真真正正.
做一個不折不扣神的子民.
我們一起低頭祈禱.
神可以求你親自幫助我們.
在我們不同的地方.
我們所經歷的.
好像是在一種浩劫,疫情當中.
神可以求你幫助我們.
不只是單單去到幻想.
去到夢想.
以往有一種正常的生活.
而是我們此時此地.
縱使好像是流亡的當中.
我們仍然來到尊利偉大.
讓你的道能夠成為一種實在.
在我們生命當中被展露出來.
讓你的榮耀臨在.
讓我們成為一個.
來敬拜你,事奉你的人.
今天就是來事奉你的機會.
讓我們所處於這種狀態.

$^{721}$你所安排我們的崗位.
我們盡心竭力地去行動.
我們這樣祈禱奉主明救.
孤單擁抱的心.
今天願教給主你.
是你的恩典與愛.
讓我可得到安慰.
憂傷可遂的心.
今天願教主修女.
是你的變相與痛.
讓我可得已痊癒.
醫治我傷.
在你的痛災裡.
讓我痛與你為裡.
我心知道你是我神.
恢復我傷.
在你的榮耀裡.
讓我可得見你完美.
我的身處.
別得已復原.
孤單懼怕的心.
今天願教給主你.
是你的恩典與愛.
讓我可得到安慰.
憂傷可遂的心.
今天願教主修女.
是你的變相與痛.
讓我可得已痊癒.
醫治我傷.
在你的痛災裡.
讓我痛與你為裡.
我心知道你是我神.
恢復我傷.
在你的榮耀裡.
讓我可得見你完美.
我的心已全然康復於你.
我來掃聖.
要尊你名為聖.
我來歌唱.
要頌讚你的名.

$^{761}$假裝用此生.
成就你的志願.
我願為你再奔跑.
我來掃聖.
要尊你名為聖.
我來歌唱.
要頌讚你的名.
假裝用此生.
成就你的志願.
我願為你再奔跑.
醫治我傷.
在你的痛災裡.
讓我痛與你為裡.
被你的愛捉摸我裡.
恢復我傷.
讓我如影輾轉.
讓我得見你完美.
我的心已全然康復於你.
讓我痛與你為裡.
讓我得見你完美.
深知你是我的全.
\newpage



\section{}
\label{sec:h_zZvugRB1g}
\textbf{浩劫重生 【第8集:亂世中的耶利米】 陳崇基牧師 - 主日講道視頻系列 8}
\newline
\newline
連結: \href{https://youtube.com/watch?v=h_zZvugRB1g}{\texttt{https://youtube.com/watch?v=h\_zZvugRB1g}} ~~~~ 語音日期: 2020-12-07
\newline
\newline
\hyperref[sec:DlruT1gd8EY]{\small{< < < PREV SERMON < < <}}
~
\hyperref[sec:index_chronic]{\small{[返順時目]}}
~
\hyperref[sec:index_scriptual]{\small{[返順卷目]}}
~
\hyperref[sec:wsAyaNl6nbo]{\small{> > > NEXT SERMON > > >}}
\newline
\newline
$^{1}$我們來到浩劫重生的第八次最後一次.
今次我們仍然是看舊約的先知書.
我們看一本書叫做耶利米書.
這裡記載一個叫做耶利米的人.
他做先知的時候的經歷.
他的講道.
他的事奉的生涯.
耶利米有另一個綽號叫做淚眼先知.
哭泣的先知.
weeping prophet.
為什麼呢?因為他是眼見.
他自己的國民不肯悔改.
不肯聽上帝的話.
最後淪陷亡國.
耶利米身處的時代就是一個國破家亡的時候.
甚至當時整個近東的世界.
國際的關係.
都處於一個天下大亂的當中.
時勢有多艱難我們可以想象.
同樣我們今天面對的世界.
我們面對很多不安動盪.
或者可以從耶利米書得到一些啟迪.
耶利米書很長的.
我們主要可以分成兩大段.
第一段就是耶利米對自己的人民的說話.
第一至第25章.
就是警告他們要他們悔改.
如果他們不悔改.
上帝的審判就會淪落.
26至45章.
就是講一些故事.
就是耶利米事奉的時候.
猶大的國民.
特別是一些君王.
不肯聽上帝的話.
最終引致上帝的審判.
第二大段就是46章開始.
耶利米這裡的矛頭.
由自己的國民指向周圍的列國.
一個一個的城邦列國.

$^{41}$來發語言攻擊他們.
46章第一節.
又說論列國的話.
臨到先知耶利米.
於是由埃及開始講.
到非利士人.
摩亞.
阿們.
以東國.
大馬士革.
基達和夏索.
這些是當時沙漠裡面的遊牧民族.
以蘭.
這些外國的列國.
通通會被巴比倫軍所滅.
但最後50至51章.
連巴比倫帝國都會傾覆.
神也會審判.
最後52章有一個副錄.
就是講到上帝沒有忘記他自己的子民.
就是第二大段.
我們想多花一點時間.
耶利米這裡想講什麼呢?.
就是上帝也會審判這些列國.
我們自然就會引起一個問題.
我們就會問.
先知針對自己本國的內政.
這個很自然.
但為什麼會針對外邦的政權呢?.
神和外邦之間有什麼關係呢?.
其實你看看不同的先知書.
你也會發現.
這些先知.
很多時候都會發聲針對外邦的政權.
外邦的國家.
以賽亞書13至23章.
以西結書25至32章.
都是發預言來攻擊周圍的鄰邦.
周圍的列國.
連一些所謂的小先知書.

$^{81}$一些比較短的先知書.
亞摩斯.
沃巴底亞.
拿紅書.
都是針對這些列邦的政權.
我們不需要去想像.
這些先知一定要周遊列國.
去到不同外邦的國家做大事.
當然有這樣的可能.
但是主要這些論列國的預言.
其實最終是要告訴神的子民.
告訴他們一些很重要的信息.
但是有些什麼要點呢?.
第一.
就是這裡想點出一種.
舊約聖經裡面.
很重要的神號.
就是神是地上萬國的主.
在近東的世界裡面.
人想起神明的時候.
就會想起一些地域性的神明.
因為只會在他們自己的地方.
這些神明才有主權.
但是呢.
在舊約的聖經裡面.
那個神學.
就是說耶和華上帝.
不僅是一個地區裡面的神明.
不僅僅是猶大國耶路撒冷作王.
他是全地的君主.
所以耶和華當然和以色列民.
是有特別的關係.
這個沒錯.
但是因為神是全地的主.
所以他會向列國.
向列邦來問責.
我們看一些例子.
46章第二至第六節.
論到關乎埃及王.
法魯尼戈的軍隊.

$^{121}$這軍隊安營在.
幽法拉底河邊的加基米市.
是巴比倫王利布格尼撒在猶大王.
約西亞的兒子約瓦敬.
第四年所打敗的.
你們要預備大小盾牌.
往前上陣.
你們套上車騎上馬.
我為何看見他們.
埃及的軍隊驚慌轉身退後呢.
他們的勇士打敗了.
急忙逃跑.
驚嚇四圍都有.
這是耶和華說的.
他們在北方幽法拉底河邊.
半跌瀑島.
這是耶和華的預言來攻擊埃及帝國.
當時的歷史背景是怎樣的.
埃及帝國和亞述帝國已經在式微.
兩軍想聯盟阻止巴比倫軍隊往西擴張.
結果埃及的軍隊.
上到一個地方叫加基米市.
在公元前605年.
打了一場很大的戰役.
這場戰役在近東歷史中非常決定性.
戰果是巴比倫軍隊大勝埃及.
耶利米的預言說.
耶和華早已有說過.
不單止如此.
巴比倫帝國也會繼續南下入侵埃及.
46章第13節.
耶和華對先知耶利米所說的話.
論到巴比倫王利布格尼撒攻擊埃及帝國.
住在埃及的人民要預備老去時所用的物件.
因為拿忽這個大京城必成方牆.
且被銷毀.
無人居住.
埃及是肥美的母牛都.
但出於北方的毀壞.
內都了內都了.

$^{161}$其中僱傭的士兵.
好像圈內的肥牛都.
他們轉身逃跑.
站立不住.
因為遭難的日子.
追討的日子.
耶和華審判的日子.
已經臨到.
這裡說到上帝率領巴比倫軍隊.
來審判埃及.
在耶利米書中有一個情節.
就是猶大人擄掠了耶利米.
想帶他逃去埃及.
猶大人以為只要我們逃難到埃及.
我們就能逃離巴比倫帝國的五字山.
神的審判就不會臨到我們.
耶利米對他們說.
沒有用的.
因為上帝是無所不在.
上帝是全地的主.
繼續在第47章.
巴比倫軍隊繼續南下的時候.
耶利米又預言到.
另外一個城邦淪陷.
就是非利士地.
考古人員.
在非利士地其中一個城市.
叫做阿什基倫.
發掘了很多年.
在巴比倫時期的土層.
發現了當時滅城的跡象.
找到火燒的灰燼.
這些完整的瓦器被打爛.
埋沒在灰燼裡面.
亦發現一個很完整的人骨.
在瓦瓶的旁邊好像在避難.
但其實這些灰燼.
這些毀壞的跡象.
都是巴比倫軍隊進行毀滅的冰山的一個.
第47章第4至5節.

$^{201}$論到非利士地.
他說.
加薩要成為降屠.
平原裡所成的阿什基倫.
要歸於無尤.
但上帝是公義.
最後.
連巴比倫帝國.
是上帝審判的工具.
都會受審判.
你留意耶利米蘇整段.
講到列國的預言.
你放在地圖上再看的時候.
其實他是在說.
由西而向東.
所講的列國的預言.
是包括整個世界.
整個近東.
連最遠的伊蘭國.
是巴比倫東邊的.
在猶太國是以東500多英里.
是和猶太國沒有什麼瓜葛的.
都在上帝列國的預言當中.
整段論列國的預言裡面.
第一個針對的國家.
就是埃及.
最後針對的是巴比倫.
兩個是當時世界裡超級大國.
super powers.
中間是一些小國和城邦.
意思是說.
沒有一個強國.
可以強到超越上帝的王權.
也沒有一個小國.
是小到上帝視而不見.
耶利米說到天下政局大亂的時候.
猶太國就正正夾在中間.
在兩個大的super power的中間.
如果你身處.
可能在香港,台灣.

$^{241}$一些類似地理政治的環境當中.
你可能看經文的時候.
你會有相同感受.
可能會在一些大的super power.
國力的中間.
帶來很多不安,混亂.
如何算?.
上帝藉著耶利米所說的話.
對他自己的子民說.
神說我是萬國的主宰.
46章第18節.
他說君王.
名為萬君子耶和華的說話.
他發出這些烈國的預言.
最終是告訴神的子民.
告訴你和我.
你要去投靠耶和華.
你要去投靠全地的真正主宰.
上帝說我會保守你.
但你不可以三心兩意.
(上帝說我會保守你).
這些攻擊烈國的預言.
第一件事告訴我們.
當然上帝才是地上萬國的主宰.
是我們唯一的依靠.
第二方面.
也很清晰地告訴我們.
上帝會向烈國.
甚至是外邦的國家問責.
為什麼?.
其中一個最明顯的原因.
就是這些烈國曾經在.
以色列人身上.
在神的子民身上行惡.
巴比倫帝國就是其中一個最明顯的例子.
上帝雖然興起巴比倫.
作為他審判的工具.
但是上帝依然在這個.
邪惡的人民裡面.
他會報復 會去審判.

$^{281}$耶利米蘇51章第24節.
耶和華說.
我必在你們眼前報復.
巴比倫人 迦勒底居民.
在錫安所行的諸惡.
35節 錫安的居民說.
巴比倫以強暴待我.
願這罪歸給他.
所以耶和華如此說.
我會為你申冤.
我必定會為你報仇.
巴比倫必成為亂雞野狗的主主.
背後有一個預節.
就是上帝和他的子民的關係.
如果有人鬥上帝的子民.
上帝一定會按時審判出手.
撒迦利亞書2章8節.
表露無遺張這個關係.
萬君子耶和華說.
在顯出榮耀之後.
差遣我去懲罰那擄掠你們的列國.
我你們的就是我他眼中的同人.
即是說鬥上帝子民的.
就好像鬥上帝眼中的同人.
The apple of the eye.
好像有人鬥你眼珠的時候.
你就一定會出手.
你一定會還擊.
上帝說如果這些列國.
在我的子民中間做欺凌.
犯事,做惡事的話.
他一定會按時審判出手.
但你綜合這些先知攻擊列國的預言.
你就可能會引起另一個問題.
就是說這些外邦國家.
有一些沒有錯.
是直接得罪上帝的子民.
去到得罪猶大國.
所以上帝就審判他們.
但也有其他的列國.

$^{321}$是和猶大國沒有直接關係.
那為什麼上帝又要去審判他們呢?.
我們可能在《以賽亞書》第24章.
有些答案.
《以賽亞書》第24章第五節.
地被其上的居民污衊.
因為他們犯了律法.
廢了律例.
背了永約.
這裡上文也說了一大堆.
攻擊列國的預言.
上帝就擊打這些外邦國家.
因為這些列國.
什麼原因?.
他們背了永約.
你越看越奇怪.
神和這些外邦國家.
什麼時候有過約呢?.
上帝和以色列人有過約.
這是理所當然的.
我們叫它西乃.
摩西之約.
Mosaic Covenant.
或者之前的阿拉罕之約.
之後的大衛之約.
其實三約都是同樣指.
同一個神和祂子民之間的關係.
只不過在聖經歷史裡有不同的演進.
最後當然實現在耶穌基督身上.
以色列人和神之間有立約.
我們很清楚.
納了約也不是免死金牌.
不可以任意放蕩.
因為約永遠都是有兩者.
一起建立的關係.
近東的盟約也是這樣.
一定要兩者一起忠於盟約.
上帝告訴以色列人.
你忠於我的約就有祝福.
你背約的話一定會有審判有酒俗.

$^{361}$以色列人當然是離棄上帝.
拜偶像.
行不宜.
甚至殺真的才知道.
流無辜人的血.
泯滅上帝的地方.
上帝警告一次又一次.
不肯悔改.
結果當然就是有審判流亡.
但是以塞瓦在這裡.
告訴我們這些外邦的政權.
背了神的永約.
Eternal Covenant.
這個字很特別.
因為舊約另一處地方提到.
Eternal Covenant.
就是創世記第九章.
這裡我們曾經提過.
在洪水的浩劫.
洪水之後.
神跟新的人類.
諾亞就是代表.
跟他重新立約.
跟大地立約.
這個叫做彩虹之約.
我們經常想起彩虹之約的時候.
我們就想得很浪漫.
上帝有一個彩虹之約.
就永遠不會再審判.
到底事實是不是這樣.
既然提到約的時候.
就一定有兩方面.
你再看創世記第九章.
你會發現近東條約所用的字眼和條文.
九章第七節.
你們要生養眾多.
在地上昌盛繁茂.
As for you.
那個強調的地方是.
As for you.

$^{401}$是你們.
正如神跟第一個人類亞當是代表.
跟他立約.
你們要好好管大地.
現在跟第二個人類.
新的人類諾亞是代表.
跟他立約.
As for you.
你要做的責任就是要好好管理大地.
從九到十節.
上帝對諾亞和他的兒子說.
我與你們和你們的後裔立約.
並且跟一切的活物大地來立約.
你見不見到有兩個的parties.
有兩方面一起進入這個永約.
近東的條約是講到兩方面的責任.
如果你看看神跟阿巴拉罕立的約.
都是同樣有這樣的字眼.
創世記第十七章第四和第九節.
他說我與你阿巴蘭立約.
要你作多國之父.
As for me.
第九節.
你阿巴蘭.
你的後裔.
世世代代要遵守我的約.
As for you.
我和你進入這個關係.
神有他的責任.
人同樣有他的責任.
創世記第九章說到.
整體的人類的責任就是要管理大地.
好像第一次創造一樣.
要好好管理大地.
如果不好好管理大地的話.
上帝的審判就會倒台.
第九章第六節.
凡流人血的.
他的血也必被人所流.
因為上帝做的.

$^{441}$是按照自己的形象去做.
這裡提到人類的社會.
人類的政權.
不能夠流無辜人的血.
相反.
流人血的.
濫殺無辜的.
就必被人所流.
就是說人類的社會裡面.
有相當的克制的制度.
去到克制的強暴.
一種所謂的司法制度的起初.
克制人類的強暴謀殺等等.
為什麼特別要提這方面.
強調這一點.
強暴.
你記得為什麼上帝會降洪水嗎?.
是因為人類的強暴.
六章第十一節.
西街在神面前敗壞.
地上滿了強暴.
就是亂禁殺人.
亂禁流人的血.
流無辜人的血.
人的責任.
就是要維繫秩序.
神的責任.
當然就是在第九章之後說.
神說我納若.
凡有血氣的.
不會再被洪水所滅絕.
記號就是採洪.
所以說永若.
是一個免死金牌.
上帝會做他的部分.
人類.
世俗的社會.
世界的政權.
人類就要負上他的責任.
所以如果你看看這些先知書.

$^{481}$攻擊列國的預言.
你就會留意到.
上帝向他們問責.
我們看一些樣本.
哈巴谷書第二章十二字.
列邦指責巴比律.
以人血建城.
以罪孽納殷的猶和了.
即是他濫殺無辜.
建造他們的帝國.
建造他們的城市.
那洪書第三章.
禍災者流人血的城.
充滿著荒詐和強暴.
搶糧的事總不止.
又是指責到外邦.
神要審判列國審判外邦.
就是他們流人的血.
濫殺無辜.
不尊重人命.
因為人就是神的形象.
所以上帝最終要審判列國.
最終的原因.
就是這些列國自高狂妄.
一次又一次.
這些先知攻擊列國的預言都提過.
四十六章第八節耶利米書.
埃及將利羅河漲化.
將江河的水翻騰.
埃及說什麼呢?.
法老說我要將法滿入到一個地步.
去掩蓋遍地.
我要毀滅城邑和其上的人民.
即是埃及是自高自大到一個地步.
他說我好像一條河一樣.
可以江河一樣掩滅整個的地.
然後其他城邑和其他人民.
我就可以將他們殺戮.
阿摩斯書第二章.
另一個樣本.

$^{521}$耶利華語之說.
摩訝三番四次地犯罪.
我必不免去他的刑罰.
因為他將以東王的骸骨焚燒成灰.
我要降火在摩訝.
要燒毀他的宮殿.
摩訝必在喧鬧,吶喊,摧國之中死亡.
我必會剪除他的審判者.
他的君王.
其上一切的領袖,一切的殺戮.
這是耶和華說的.
這裡很明顯.
摩訝所犯的罪.
是跟以色列無關.
因為他得罪的是以東.
另一個外邦.
將以東王的骸骨挖出來.
焚燒成灰.
這是神人共憤的事.
非常殘暴的事.
近東的火葬已經少有.
還將骸骨拿出來.
然後火燒成灰.
成為lime.
指的意思可能是將灰燼.
那些骨灰.
變成了plaster.
變成了在房屋上.
想抹牆的那些plaster.
那些白灰.
這是很殘暴,很強暴,很泯滅人性的事.
跟以色列人無關.
但是上帝說.
你背了我的永約.
所以我要來審判你.
列國因為心高氣傲.
因為他們泯滅人性.
所以上帝的審判就這樣.
或者可能因為錢多.
或者因為他們富裕.

$^{561}$所以令他們心高氣傲.
上帝也會來問責審判.
講到摩訶耶利米蘇48章第七節.
他說摩訶.
你因倚靠自己所做的.
以及自己的財寶必被公取.
巴比倫帝國也是一樣.
富有到一個地步.
萬國與他行刃.
上帝就來審判問責.
51章第十一節.
講到巴比倫的末代晚年的時候.
耶和華定義攻擊巴比倫.
因此是耶和華報仇的日子.
就是為自己的奠報仇.
你們要樹立大旗.
攻擊巴比倫的城牆等等.
巴比倫以前的古城.
牆的厚度.
內牆有21尺厚.
外牆有20尺厚.
是兩層的牆.
幾十尺厚.
甚至跨厚度到一個地步.
是堅不可摧.
城牆甚至有250座城樓保守.
所以巴比倫王可以誇口.
說沒有東西可以入侵.
但是大家知道近東歷史.
公元前539年.
波斯的軍隊.
不費一兵一卒.
全城投降.
我們看了耶利米當時的政局.
那種混亂不安天下大變.
也在他的先知書中.
抽了一些經文.
去講審判列國的預言.
我們需要問.
到底這些經文和我們今天有什麼關係.

$^{601}$先知書請你留意.
不是預先寫好的報紙.
所以我們不能說.
巴比倫是今天的巴格達.
所以先知針對當時的政局.
就等於針對今天的政局.
是兩回事.
因為歷史背景不同.
當時上帝審判的工作不同.
先知書就記載下來.
但是先知書不是直接寫給我們21世紀的人.
所以我們不能夠回到先知書.
我們就說他正在預測今天中東的政局.
你說如何去形容.
先知書.
現代的意思.
就不是預測.
不是拿著早已寫好的現代報紙去看.
而是一個範例.
paradigm.
不是預測.
而是一個paradigm.
我們看回耶利米時代.
列國行強暴.
泯滅人城.
不尊重上帝的主權.
犯了上帝的榮譽.
上帝就要問責.
所以巴比倫.
一個這麼大的帝國.
短短七十年裡就敗亡.
今天我們見到這個範例.
我們就可能要去到景城.
我們就想起今天世界上.
有很多不同的列邦.
列國.
橫行無忌.
以為自己的權勢.
他就可以做很多不公義的事.
行很多強暴.

$^{641}$先知書說有一樣東西沒有改變.
就是上帝照樣會行審判.
上帝會照樣出手.
會去行公義.
我們所信的.
就是這位萬國私行公正的主宰.
值得我們去信教.
值得我們繼續去宣揚他.
我們同心,低頭祈禱.
神啊,你是無所不在的上帝.
我們今天見到很多的強權.
流無辜人的血.
見到很多講真理的人.
去捍衛公義的人.
他們被人收監.
被人視為無物.
被掌權者誘餌.
我們見到很多關心社會.
關心城市的人.
被禁聖,被抹殺.
高高在上的人.
視而不見,視若無睹.
甚至落井下石.
我們見到教會.
我們的弟兄姊妹.
在不同的地方.
在中國大陸.
在一些伊斯蘭國家.
受盡侮辱.
受盡欺壓.
在試煉當中.
主要我們見到這些不義的事.
上帝我們不單單心痛.
我們更加無助.
面對局勢的動盪.
我們在未知之數當中.
我們好像被歷史的巨輪所擺佈.
但我們深知.
你今天同樣的掌權.
你是在列國當中.

$^{681}$掌權施行審判的公義主宰.
主要是為了在監獄中.
被強暴,被流放的人.
黑眼的德漢賢.
這是你對我們的應許.
求你讓這些受欺壓的人.
見到你的救恩.
求你幫助他們.
以你為他們的依求.
主要你講過.
你的國度要來到.
被擄的會得釋放.
你成為他們的寶藏.
神求你成為他們的神.
主要是你的國度.
你的權力.
求你彰顯.
在黑暗中讓我們見到光明.
見到盼望,見到公義.
今天在不安的當中.
我們依然來到深信,倚靠你.
求你幫助我們.
成為一個忠於所有的人.
我們今天依然為你發光.
在我們的光位上我們要尊重.
祈禱奉主耶穌基督的名字.
\newpage



\section{路加福音}
\label{sec:wsAyaNl6nbo}
\textbf{疫年聖誕信息 (路加福音二章) - 陳崇基牧師}
\newline
\newline
連結: \href{https://youtube.com/watch?v=wsAyaNl6nbo}{\texttt{https://youtube.com/watch?v=wsAyaNl6nbo}} ~~~~ 語音日期: 2020-12-26
\newline
\newline
\hyperref[sec:h_zZvugRB1g]{\small{< < < PREV SERMON < < <}}
~
\hyperref[sec:index_chronic]{\small{[返順時目]}}
~
\hyperref[sec:index_scriptual]{\small{[返順卷目]}}
~
\hyperref[sec:eJV9zLgjq90]{\small{> > > NEXT SERMON > > >}}
\newline
\newline
$^{1}$各位聖誕平安.
今時今日講這句說話可能會有很不同的意思.
或者今年2020年的聖誕節正是提醒我們.
以往一直這麼多年西方的世界文化已經將聖誕節無限浪漫化.
我們一想起聖誕節就想起火樹銀花.
每個人都會很開心.
聖誕是否真的這樣.
今時今日疫情仍然橫掃全球的時候.
我們想起聖誕就非常不一樣.
或者這個可以提醒我們.
本身第一次聖誕應該的處境是怎樣.
路加科第二章這樣說.
在百里行之野地裡有牧羊的人夜間按照更次看守羊群.
有主的使者站在他們旁邊.
主的榮光四面照著他們.
牧羊的人就甚懼怕.
那天使對他們說不要懼怕.
我給你們大喜的訊息是關乎萬民的.
因為今日在大衛的城裡為你們生了救主就是主基督.
你們要看見一個嬰孩包著布臥在馬槽裡那就是記號了.
忽然有一大隊天兵和天使讚美上帝說.
在至高之處榮耀歸於神.
在地上平安歸於他所喜悅的人.
眾天使離開他們升天去了.
牧羊的人彼此說我們往百里行去.
看看所成的事就是主所指示我們的.
他們急忙去了就尋見馬里亞和約室.
又有那嬰孩臥在馬槽裡.
既然看見就說天使論者孩子的話傳開了.
凡天見的就察現牧羊之人對他們所說的話.
馬里亞卻把這一切事傳在心裡反覆思想.
牧羊人回去了因所看見的一切事.
正如天使向他們所說的就歸榮耀於上帝讚美他.
說起聖誕節不少教會和基督徒會在聖誕節有崇拜.
有特別的慶典.
這當然不代表我們說十二二十五號西曆就是耶穌出生那一天.
我們不會搞錯的.
初期的教會只是到了公元後四世紀開始.
才有聖誕節慶祝的傳統.
由第一世紀開始初期幾百年.

$^{41}$基督最重要的事件.
教會所紀念的全部都是集中在耶穌的死和復活裡面.
但我們會問耶穌誕生這麼大件事.
福音書的記載有什麼交代.
特別說到耶穌出生的時間.
到底第一個聖誕節是什麼時間.
如果不是十二月二十五號的話.
應該是什麼時候.
很多人自然會指出.
在路加福音第二章第八節.
伯利恆的野地有牧羊的人.
夜裡按照庚次看守羊群.
他們就會指出一段經文說.
不可能是冬天.
到底是不是呢?.
可能我們要看看地理和氣候.
在伯利恆的野地有牧羊的人按照庚次看守羊群.
其實一年四季都能夠這樣做.
就算是雨季或冬天的時候.
為什麼呢?.
我們只需要明白一點.
以色列巴勒斯坦地的地理和氣候就會知道.
這個又長又窄的以色列地.
由地中海的海皮.
沿海平原這個地方.
一直往東走的時候.
你會發現地勢會越來越高.
直到最高的地方.
就是一個中央山脈.
Central Hill Country.
上面有什麼城鎮呢?.
就是耶路撒冷和伯利恆.
伯利恆就在中央山脊的分水嶺之上.
為什麼叫分水嶺呢?.
因為每逢雨季來到的時候.
冬天雨水就會落在分水嶺以西的土地上.
但當我們過了分水嶺.
地勢就會一直下降.
會落到油帶礦業和死海低谷的地面.
所以礦業和死海峽谷.

$^{81}$即使是冬天和雨季都會比較亂和乾旱.
今時今日巴勒斯坦地的牧羊人.
都會在冬天的時候出外牧羊.
你查一查天氣報告就知道了.
在聖誕的星期.
氣溫是怎樣的呢?.
伯利恆最低氣溫是攝氏5至7度.
或是40至45度.
最高氣溫是攝氏14至17度.
或是57至63度.
但當你過了分水嶺.
你下去死海峽谷.
這裡的氣溫.
這個星期最低氣溫是攝氏13至15度.
或是53至57度.
但最高氣溫是攝氏21至23度.
或是68至72度.
在整個星期都是天晴.
只是星期三至四會有少許的驟雨.
如果你這樣看.
其實溫度和氣候都是非常溫暖.
對牧羊人甚至羊仔.
應該不會有很大問題.
所以我們不能單單說.
牧羊人在野地看守羊群.
就不可能在冬天.
但想清楚這句話.
是跟當地的地理.
農耕的運作很有關係.
因為巴勒斯坦地大麥小麥收割的時候.
大概是六月的時候結束.
收割結束之後.
當地的農民很多時候會跟牧羊人有一個合同.
允許牧羊人帶羊群進入自己的農田.
因為收割之後.
剩餘的穀物.
牛羊羊群就可以吃掉.
同時牛羊也可以在農田留下肥田鳥.
因為這樣的合作.
我們了解原來這段經文背後有熟悉背景的人.

$^{121}$一聽就知道.
野地有牧羊人.
就代表這個時間大概是七月八月甚至九月的時候.
這個時候正是夏季的時候.
全年最乾旱和最難捱的時候.
或者這裡說到耶穌誕生的時候.
神的加冕跟我們很有意思.
因為正是一個最難捱的時候.
耶穌來到我們中間.
天國就降臨.
可能在今年的聖誕裡.
大家都不是最歡喜.
相反可能是最孤獨.
最抑鬱的時候.
因為疫情或者因為政局的不利.
但是這個時候.
這個故事告訴我們.
主就是在最難捱.
最辛苦的時候就臨到我們.
或者另一方面.
這段經文裡面.
講到不是真的重點在時間那裡.
而是留意某一些字眼.
路加福音第二章開始記載聖誕的故事的時候.
他講到什麼呢?.
當那些日子.
凱撒奧古士都有旨意下來.
叫天下人民都報名上策.
就是居里樓作敘利亞巡撫的時候.
頭一次行報名上策的事.
所有人各歸各城報名上策.
留意到這裡的重點.
是講到天下所有的人.
全部的人.
眼前的現實.
就好像羅馬帝王凱撒說了算.
這是羅馬帝國的政治宣傳.
經常有的口號.
天下所有人都是受他嚇管的.
但是這個聖誕故事裡面.

$^{161}$告訴我們什麼呢?.
二章第十節.
那天使對他們說不要懼怕.
我報給你們大喜的訊息.
是關乎萬民的.
因為今天在大衛的城裡.
為你們生了救主就是主基督.
忽然有一大隊天兵和那天使讚美上帝說.
在至高之處榮耀歸於上帝.
在地上平安歸於他所喜悅的人.
留意到神在做什麼.
你怎麼解釋都好.
明事或暗事.
這裡也講到一種跟羅馬帝國相反的宣傳.
就是真正的救主.
真正的平安.
真正管治天下萬民的.
是耶穌.
祂才是全人類的救主.
這裡講到今天.
你說有什麼特別.
今天對第一代的讀者來說.
就是羅馬帝國的作主.
就是祂說了算.
一切都是在祂操控之下.
連耶路撒冷聖諜也是在祂操縱之下.
救主就是羅馬帝王.
就是靠祂才有平安.
能夠風山水起.
但是天使跟他說什麼.
真正的救主.
真正的王者.
是一個新的大衛君王.
真正祂出生的地方.
大衛城.
是原本的百里行.
大衛所出生的地方.
這個出生於靶槽的嬰孩.
才是王者.
今天不是明天.

$^{201}$老家說今天.
是神降臨.
神得醒的時候.
不是明天.
這是路加福音經常用的字眼.
就是告訴我們.
雖然眼前好像是帝國.
是黑暗掌權.
但是今天.
神在地上彰顯.
神在地上得勝作.
祂誕生.
因為整個歷史.
整個宗末.
都在上帝的手裡.
老家經常提到今天這個字眼.
有好幾次.
最後一次.
是在路加福音二十三章.
耶穌被釘在十字架上.
跟祂一起釘的強盜.
垂死的時候.
有一個這樣的對話.
耶穌對那個強盜說.
我實在告訴你.
今天你要和我在樂園裡.
眼前的現實.
好像是羅馬帝國玩弄了.
好像是它控制了所有的東西.
但是耶穌在十字架上宣告什麼.
祂說今天你要和我在天國裡.
或者這個就是聖誕的奧秘.
或者這個就是為什麼.
到了第四世紀的時候.
教會會選擇十二月夜五號.
來慶祝聖誕.
為什麼會有一個.
慶祝聖誕的傳統出現.
昔日去到定夏.
一個聖誕日期的時候.

$^{241}$不是因為他們查根究底.
明確地知道.
這個就是耶穌誕生的那一天.
而是當時的教會.
是取用了民間的宗教慶典.
當時是拜太陽神節.
是冬節的時候.
就是日光是最短.
重回日光最長的日子.
世人就在敬拜太陽神.
因為太陽神可以給他們有光明的盼望.
初期的教會.
就取用了這一天來慶祝.
耶穌可以得勝黑暗死亡的權勢.
能夠讓我們重見生命的光明.
黑暗好像是掌權.
眼前的現實好像是帝國主事.
但是耶穌的降臨.
就代表神的得勝.
代表上帝的勝利已經展開了.
今年的聖誕.
我們每一個人都會覺得失控.
原先的計劃.
可能要去哪一個地方.
我們要被打消.
心裡面可能記掛的人.
可能不在我們身邊.
或者可能是病危.
是ICU.
我們生命裡面家庭的問題.
經濟的問題.
很多時候都是揮之不去.
因為有什麼可以做.
在聖誕節裡面.
聖誕的故事.
可能要我們留意的.
耶穌誕生的時候.
牧羊人有什麼回應.
二章第十五節.
我們不妨再看.

$^{281}$眾天使離開他們.
升天去了.
牧羊人彼此說.
我們往伯利行去.
漢漢所成的事就是主所指示.
告訴我們的事.
他們急忙去了.
就尋見瑪利亞和約瑟.
又見到乃英孩臥在馬槽裡.
既然看見.
論者孩子的話全開了.
凡聽見的就察而牧羊之人.
對他們所說的話.
瑪利亞將這些事.
藏在心裡反覆思想.
牧羊人回去.
因所聽見所看見的一切事.
正如天使所說的.
就歸榮耀於上帝.
讚美他.
在路加.
和路加的第二集.
使徒陶傳.
經常會出現一個主題.
讚頌上帝.
這次不是第一次.
之前瑪利亞.
已經有做過.
也不是最後一次.
總共在路加和使徒陶傳.
出現了二十八次.
很多時候.
這些讚美的說話.
讚美的聲音.
在迫害教會.
受打壓的時候.
他們見到上帝的工作.
在重重困難裡.
將上帝的道傳開的時候.
他們見到上帝仍然在作旺.

$^{321}$所以最後一次出現.
並非偶然.
在使徒陶傳第二十八章十四節.
這裡說到.
保羅來到羅馬.
那裡有一些.
信主的弟兄來接待他們.
保羅見了他們.
就做了什麼呢?.
並且放膽傳講上帝的國.
將主耶穌基督的事.
教導人.
並沒有人禁止.
好像每樣東西.
都有很多艱難.
好像黑暗的權勢.
是贏了.
但是他說今天.
仍然可以哼出讚美的歌.
請你留意.
讚美敬拜的歌.
不是自嗨.
不是自我感覺溫馨.
而是海歌.
是勝利的歌.
宣告主耶穌是得勝.
宣告這位王者.
是主耶穌基督.
不是羅馬帝國.
不是疫病.
不是苦難.
不是死亡.
這才是聖誕.
是聖誕.
聖誕有另一個名字.
叫做Event.
將臨節.
我們等候主的駕臨.
我們既知.
二千年前.

$^{361}$耶穌王者.
曾經駕臨.
死而復活成就救恩.
他坐在上帝的右邊.
他會審判死人活人.
他會再回來.
這個Event.
就是我們仰望耶穌.
再駕臨.
他會審判死人活人.
這位主宰的王者.
今天如何提醒我們.
當我們仍然鬆懈.
仍然與上帝的關係.
搞不清楚的時候.
這個聖誕節就是讓我們.
回轉到警城.
如果我們眼前有很多困境.
很多失望.
很多烏雲籠罩我們的時候.
今天.
讓我們看得通.
眼前的現實.
不可以決定我們生命的.
方向與狀態.
應該是主的駕臨.
主作王.
耶穌基督這位在馬槽裡.
誕生的主宰.
才是歷史的王者.
是我們值得信求.
求主幫助我們.
在這個章林節.
聖誕節裡.
讓我們轉向真正的主.
耶穌基督.
在祂裡找到盼望,找到光明.
我們同心大道祈禱.
上帝在聖誕節當中.
我們看到.

$^{401}$全地全球的人.
都被疫情所擺佈.
甚至可能在.
世界不同的角落.
不單止有疫情.
也有政治動盪.
甚至看到政權濫用疫情.
不顧人命.
神我們也仰望你.
因為你提醒我們.
耶穌基督.
誕生在馬槽裡.
這位主.
才是我們生命的主宰.
曾經你駕臨.
曾經你將光明帶到這個世界.
並且你應許.
會再臨.
主讓我們來仰望你.
讓你會再回來.
會審判活人死人.
你是公義的主.
會降臨.
這個盼望,這個忠梅.
成為我們今天生命的方向.
我們將我們的信靠.
將我們的仰望.
放在你身上.
求你繼續幫助我們.
看到你天國的光明.
並且我們坐在地上.
不斷有這個盼望的時候.
我們的生命能夠彰顯.
上帝你自己的榮光.
奉耶穌基督的名求.
阿們.
\newpage



\section{}
\label{sec:eJV9zLgjq90}
\textbf{作神的子民 - 再讀歷代志 (直播重溫) - 陳崇基牧師}
\newline
\newline
連結: \href{https://youtube.com/watch?v=eJV9zLgjq90}{\texttt{https://youtube.com/watch?v=eJV9zLgjq90}} ~~~~ 語音日期: 2020-12-29
\newline
\newline
\hyperref[sec:wsAyaNl6nbo]{\small{< < < PREV SERMON < < <}}
~
\hyperref[sec:index_chronic]{\small{[返順時目]}}
~
\hyperref[sec:index_scriptual]{\small{[返順卷目]}}
~
\hyperref[sec:n3sfAz_OcWA]{\small{> > > NEXT SERMON > > >}}
\newline
\newline
$^{1}$各位好,希望大家聽得清楚.
首先跟大家say hello.
因為這是2020年最後一個主日.
希望今天能夠在live跟大家見面.
不單止是和大家分享訊息.
也可以在live和大家互動.
每一次做live的時候.
我們都會有少許互動的時間.
大家知道有一個live chat的box.
就是YouTube有一個live chat的box.
我想大家先熱身一下.
第一,這裡這麼多位朋友來到的時候.
首先把你代表的城市的名字打上去live chat.
今天由一開始到現在.
我看到有北美東岸和西岸的朋友.
我看到有佛羅里達州和西岸的朋友.
另外也有加拿大的朋友.
我看到有Vancouver,Amy,Celine.
一早就到了.
另外有多倫多,Sunny和其他朋友.
有香港,有夜貓.
香港的人很晚的.
另外也有新加坡,還有不同城市的朋友.
請大家把你代表的城市打上去.
這裡有很多,應該有人找.
也有芝加哥的朋友,我們城匯的本地的弟兄姊妹.
首先歡迎各位.
今天是2020年最後最後的一個主日.
我們希望能夠有一個互動的時間.
和和大家分享神話.
不知道你們想起2020年最後的時候.
有什麼經文.
當我們說訊息的時候.
我希望今天的時間比較緊迫一點.
訊息可能只有十多分鐘.
然後有什麼問題.
有什麼comment(留言).
大家可以在live chat和我們互動.
首先試一下快問快答的時間.
我們先熱身一下.

$^{41}$問你一條問題.
通常都是一些選擇題.
你把答案打在live chat.
可能是1,2,3.
把答案打在那裡.
然後快問快答.
看誰先答對.
答對的有獎品送.
我和後面的義工弟兄姊妹談過.
送什麼.
你知道我們的UV class.
我們的sunday school case(日曆學期).
就搬了上online.
然後剛剛麗仙有報告.
說我們和北美神學院.
會做一些certificate(證書)的program(程序)等等.
所以我們live class的page(專頁).
大家可以去register(登記).
可以去報名參加不同的class(課程).
裡面有一個check out(退款)的時候.
選了class check out的時候.
有一個叫做discount code(折扣碼).
我們會送這個discount code給大家.
有越南的朋友也有.
哇 厲害.
Hello 你好.
香港很多朋友.
我們先做個熱身.
譬如我出一條問題.
譬如這條.
就說如果你要找多一些.
承換網上課程的資料發展的話.
你可以在哪個網址找到呢.
那哪個網址找到呢.
我們繼續出這個選項的時候.
第一個就是slash ourpassers(我們的通行證).
第二個是slash class(我們的課程).
第三個是urbanvoicechurch.com(urbanvoicechurch.com).
那哪一個是正確的答案呢.
你們現在就打下去.

$^{81}$可以嗎.
先打開.
first question(第一條問題).
如果第一個答對的話.
我們會把獎品送給你.
不過這個是熱身.
答案當然是2.
如果答對第一個的話.
用我這裡的live chat box(直播箱)作準.
就是李小貓 香港夜貓.
但是正如我所說.
這個只是熱身.
所以大家留意.
我們等一下有一個短的訊息.
正如我剛才所問.
當你想起2020年的時候.
你會想起什麼經文.
想起什麼聖經裡面的說話.
2020年相信一定是我們人生當中.
在世界的歷史裡面.
一個很大的turning point(轉折點).
一個很大的轉折點.
天翻地覆.
搞到我們民不聊生.
搞到我們很混亂.
所以你想起聖經的歷史.
其實你想起不同的年代.
對於舊約的神的子民.
舊約的教會.
都有很大的衝擊.
其中一年就是.
主前586年.
這一年.
是對舊約的神的子民.
一個最大的轉折點.
因為巴比倫軍隊來到耶路撒冷.
將耶路撒冷毀滅.
以及燒了聖殿.
神的子民當時是猶大國.
被露到異鄉.

$^{121}$在異鄉流亡.
我們知道舊約的歷史.
流亡了大約70年.
然後又改朝換代.
整個天下.
巴比倫帝國瓦解.
然後波斯帝國.
統治了當時的天下.
當時的群島.
在舊約的歷史裡面.
公元前586年.
確實是一個最decisive(決定性).
最決定性的一年.
連今天猶太人.
也在紀念這一天.
第一聖殿.
耶路撒冷的聖殿被毀.
我想起了這一段經文.
歷代至下36章.
記載了舊約歷史最末段的時間.
第17節開始.
所以耶路撒冷.
使巴比倫人.
即是迦勒底人的王來攻擊他們.
在他們的聖殿裡面.
用刀殺了他們的壯士.
不連戳他們的少男處女老人白手.
耶和華將他們.
都交在迦勒底王的手裡.
迦勒底王就將上帝的殿裡面大小的器皿.
耶和華殿裡面的財寶.
王和眾首領的財寶.
都帶到巴比倫去.
然後巴比倫人.
迦勒底人就燒了上帝的聖殿.
拆毀了耶路撒冷的城牆.
用火燒了宮殿.
或者裡面的財寶.
裡面的器皿.
然後犯脫離刀劍的.

$^{161}$迦勒底王都擄到巴比倫.
作他和他子孫的補被.
直到波斯帝國興起.
這是應驗先知耶利米所說的話.
他說土地享受安息.
因為土地荒涼便受安息.
寂滿了七十年.
這裡記載主前586年.
猶大國神的子民終於亡國.
這麼大的事情.
我們覺得最重要的那一年就是2020.
對於舊約的歷史.
就是公元前586年.
歷代志是什麼時候寫的呢?.
是神的子民.
猶大國的國民.
被流放到巴比倫.
然後過了幾十年之後.
再回到猶大的地方.
回到耶路撒冷.
再重建聖殿.
再重建聖城.
你想起舊約裡面.
以斯拉,尼希米那段時間.
當時的人就在想.
到底我們是誰?.
上帝曾經將我們流放了.
現在我們又回來了.
但我們仍然在帝國的陰影下.
改朝換代.
很多事情我們還未知之數.
到底我們應該怎樣定位呢?.
到底我們是誰?.
是一種很大的混亂.
歷代志是分開幾段的.
整個歷代志上下.
由加普開始.
由亞當的加普開始.
一直追溯到族長.
到以往的以色列歷史.

$^{201}$然後建立他們聯合的王國.
就是薩姆爾,蘇魯和大衛.
所羅門的時間.
然後南北分裂.
然後到第三段.
記載猶大國的國史.
由南北分裂之後.
到猶大國的末年晚年.
特別是到三六章最後.
就是我們剛才所讀的經文.
就是講到巴比倫軍隊兵臨城下的時候.
耶路撒冷被燒毀.
聖殿被毀.
最後有兩節.
就講到在波斯帝國改朝換代之後.
這個塞魯士.
Syrus.
頒佈諭旨.
讓這些被流放的猶太人.
能夠回到自己的故土.
所以段落的脈絡大概就是這樣.
其實歷代志的作者.
在整本書裡面.
他會引用不同的紀錄.
不同留下的紀錄.
一些歷史的紀錄.
或者一些先知說過的話.
這裡只是四段.
如果你數一下歷代上下.
有很多地方.
作者都說.
我說了這些東西.
寫在哪裡呢.
比如九章第一節.
寫在以色列的諸王記上.
就是King of Israel的紀錄.
二十九章二十九節.
這些事情記載在哪裡呢.
先知撒母耳的書上.
拿丹先知等等他們的書上.

$^{241}$以賽亞所記的書上.
有二十幾三十處的地方.
不斷地在講這個東西.
所以歷代志的作者.
就是那些被迫回歸.
回到耶路撒冷.
他們要重整他們的生命.
重整他們的群體.
他就去搜集不同的史記.
不同的書卷.
你可以想像到.
可能是當時公元前586年.
他們逃難的時候.
那些文士就要拿著那些.
除了撿自己的東西之外.
就要拿著那些書卷.
就是宮廷留下了.
先知書裡面寫過的東西.
宮廷裡面的史記.
那些紀錄.
就全部拿走.
然後被人拿去巴比倫.
然後再帶回來耶路撒冷.
然後他們重新寫這段歷史.
回憶過往的歷史.
到底看到了什麼呢.
看不看得見到上帝有的角色.
或者上帝有的工作.
同樣我們今天去看2020年.
到底我們怎樣回憶我們這一年.
裡面發生過的事.
經歷過的事.
到底我們有沒有看到.
上帝的足跡.
或者上帝的工作.
說到歷代志.
歷代志是希伯來聖經.
Hebrew Bible裡面最後一本書.
這個表.
你先看右邊.

$^{281}$右邊就是我們現在用的舊約聖經.
我們按照那個內容來分.
這個是七十四億本.
開始就是這樣.
將目錄這樣安排.
分四個.
律法書Torah就是五經.
然後有歷史書.
就是我們所說的書士德薩王等等.
然後有智慧詩歌.
白詩森傳國.
然後就是我們說的大小先知.
塞埃潔丹.
然後有十二小先知.
這個右邊的方法就是我們一向看的舊約歷史.
那個canonical order.
就是正點的目錄.
但是左邊就是猶太人用的Hebrew Bible.
叫做Tanakh.
就是分開三個不是四個.
五經是沒有變.
但是他們就將我們所說的歷史書.
當作是先知書.
前先知書就是書士薩王.
後先知書就是我們今天所說的先知書.
大小先知.
塞埃潔丹.
然後有十二小先知.
第三部分就是Ketuvim.
叫做Writings.
意思就是寫作.
由詩篇開始一直排到歷代字.
放在最後.
所以剛才我們看的聖經歷代字.
First and Second Chronicles.
作為一本書放在Hebrew Bible最後.
所以為什麼猶太人用的Hebrew Bible叫做Tanakh.
因為是Torah.
Nephilim.
和Ketuvim.

$^{321}$所以TNK就是Tanakh的意思.
歷代字寫在最後.
其實就是讓神的子民再想一想.
回憶他們整個舊約的歷史.
由亞當開始 加普開始.
說到猶太國的興衰.
到最後 上帝在他們身上做了什麼.
所以是一種回憶和前瞻的神學著作.
讓我們回想過去發生過的事.
但是用上帝的角度去看.
一個theological perspective.
到底上帝曾經在我們身上做過什麼.
縱使有浩劫 有苦難.
縱使將來有不同的uncertainty.
有不同的未知之數.
到底上帝當中是怎樣呢.
我們作為神的子民.
我們應該怎樣去define我們自己.
怎樣去為我們自己作定義.
所以這就是歷代字.
我們等一下看看歷代字.
到底有什麼中心值得去反省.
所以我們現在又快問快答.
講著講著我們又問問大家.
因為這些問題幫助我們了解這些訊息和資訊.
第一條問題.
現在是禮真的.
我要將問題.
Question 1.
Will.
先打上去.
問第一個問題.
第一個問題就是.
希伯來文 舊約聖經.
將所有書卷分成哪三類.
第一是律法書 歷史書 智慧書.
第二是律法書 前先知書 後先知書.
第三是律法書 先知書 和寫作.
到底哪一個才是正確的答案呢.
剛才不是真的.

$^{361}$剛才是玩玩而已.
現在是真的.
到底希伯來聖經.
將所有書卷分成哪三類呢.
答案很顯然就是第三.
律法書 先知書 和寫作.
所以叫做Tanakh.
因為有Torah 先知書 和寫作.
你去過我live的朋友.
你會知道問題很容易.
但這次不是真的.
這次比較難.
第三就是正確答案.
分開Torah 先知書 和寫作.
最後的那本就是剛才所說的.
First and Second Chronicles.
歷代字.
第二條問題準備好了嗎.
第二條問題.
舊約歷史最大的轉捩點是耶路撒冷第一聖殿被毀.
這件事發生在哪一年.
第一是公元前931年.
第二是公元後70年.
第三是公元前586年.
這年是哪一年呢.
在舊約歷史最大的轉捩點.
答案很明顯是586年.
這一年就是舊約聖經經常想起.
流亡然後再回歸.
這是第一聖殿被毀.
第二個選項.
我們很容易搞錯.
公元後70年.
這是第二聖殿被毀的時候.
第一聖殿被毀是舊約歷史最決定的時刻.
很多朋友答對了.
繼續努力.
第三條就是舊約歷史書.
在希伯來文的聖經中.
希伯來文的書是哪一類的書.

$^{401}$第一是先知書.
第二是律法書.
第三是前先知書.
這個很難.
我們現在所看的歷史書.
書士德澤王等等.
在希伯來文中.
到底哪一類書是屬於哪一類的書卷呢.
答案是哪一類呢.
這個有點困難.
其實是前先知書.
先知書不是很正確.
分為前先知書.
前先知書就是希伯來的聖經.
去看歷史書.
為什麼呢.
要知道看歷史書的時候.
他不只是告訴你以前發生過什麼事.
而是用先知的角度告訴你為什麼如此發生.
所以有些說話是.
他們不聽先知的說話.
不聽阿哈的說話.
不聽以利亞的說話.
所以以色列國或猶大國.
一直在下降.
所以是一種先知書.
是宣告上帝的說話.
當他們不聽上帝的話.
不去敬拜上帝的時候.
不去跟隨上帝的約書的時候.
他們的生命裡面就有很多問題出現.
我們去到最後就是歷代志的中心.
在這裡最後想點出幾點.
有三點.
第一 歷代志想告訴回歸了的神的子民.
他們真正的神國身份.
其實是去redefine神的子民的identity.
什麼叫神的子民 什麼叫基督徒.
什麼叫教會.
今時今日在我們現在這個時代.

$^{441}$我們的身份是用地上的東西來決定.
有些人用他社會地位 用他賺多少錢來決定.
用他的身份來決定 用他自己的權力來決定.
甚至用他的政治來決定他的身份.
歷代志就是去挑戰我們.
我們真正的身份是在上帝的國度裡面.
我們真正的身份是屬於上帝的國.
所以一開始就說 由第一章至第九章.
就說加普 由亞當 亞當開始一直傳下去的時候.
他們的歷代志就說我們真正屬於的是誰.
不是任何帝國.
巴比倫或者波斯.
或者將來改朝換代的希臘或者羅馬.
這些會變 這些會expire.
但我們神國的身份是不會變.
我們真正效忠的是上帝 不是地上的東西.
第二就是一個敬拜的呼召.
一個call to worship.
如果你看歷代志的話.
你就知道他們經常說聖殿.
聖殿的敬拜.
怎樣回歸到上帝面前 真正去worship 真正去敬拜.
第三就是說到神仍然似盼望.
有一個問題剛才提過的.
有一位朋友.
就是我們面對很多苦難 面對很多改朝換代.
很多未知之數的時候.
怎樣去定位.
歷代志最後那兩節就很有意思.
故意這樣去結束整本書.
三學章二十二節.
就是波斯王 Cyrus 的元年.
當時就是539年.
耶和華要去應驗耶利米茲所說的話.
就激動波斯王的心.
然後他就下詔全國.
然後他就說什麼呢.
他說天上的神已經將天下的萬國賜給我了.
又祝福我在猶大的耶路撒冷為他建造殿院.
你們中間凡是他的子民就可以上去.

$^{481}$上回去.
Go up to Jerusalem.
願耶和華你們上帝與他同在.
被擄之後.
上帝的手仍然去帶領著.
仍然去引領著自己的子民.
沒有丟下過他們.
將他們帶回來.
是一個恩典的結束.
但也是一個開放的結束.
一個開啟的開放的結束.
看完《易代志》再去定義自己的時候.
到底我們真正屬於誰.
然後到最後這兩節.
就告訴我們經歷過公元前586年.
我們亡國 被擄 流放 再回歸.
仍然在帝國底下.
但是上帝仍然在助我們.
上帝仍然去帶領著我們.
這樣去到回憶歷史.
可能給我們一些角度去看2020年.
在2020年裡面.
發生了很多不同的事情.
可能你的個人生命.
發生了很多不同的事情.
可能你的教會.
有很多不同的衝擊.
世界有疫情影響.
工作裡面遇見很多不同的挑戰.
或者家庭遇見很多不同的危機.
不同的危機.
怎樣去回頭看.
我們仍然有一個theological perspective.
有一個神學的角度.
我們能夠去理解.
理解上帝在我們身上.
仍然沒有放過手.
是和我們同行.
我們先低頭祈禱.
上主我們多謝你.

$^{521}$當世界翻天覆地.
改朝換代的時候.
在我們面對很多未知之數.
我們很難做一些打算的時候.
但我們知道你仍然是引領著歷史.
特別是引領著屬於你的子民.
上帝求你再一次提醒我們.
你依然是歷史的主宰.
縱使疫情橫掃全球.
縱使我們每個人的生命.
好像被捲入一個漩渦.
但我們依然知道.
你是那位帶領著.
每一張歷史的寫作的神.
上帝求你請你親自與我們.
每一個屬於你的子民同在.
當我們和你還未搞清楚關係的時候.
主要求你成為我們的神.
讓我們認識你.
我們為每一位今天來敬拜的弟兄姊妹.
我們祈禱.
特別是在一些地方.
一些城市.
有很多不同的動盪.
有很多令人不安的事情發生.
主要求你親自與他們同行.
我們為香港這個地方.
面對很多政局的不穩.
很多令人很心酸.
不義的事情發生.
求主你親自與你的子民同在.
求你親自幫助他們.
將他們帶到保守的底下.
求你親自幫助他們.
在他們的城市.
縱使在黑暗當中他們能夠發光.
主要為美國在最近在 Tennessee.
也有炸彈的事情發生.
在政局裡面也有很多不安.
神讓我們再一次來到你面前.

$^{561}$我們能夠知道在上帝你的手中.
我們有平安.
在一些社會的不安裡面.
我們的身份定義於天國子民的身份.
不是地上的政局.
地上的東西.
而是真正屬於你的.
你都可以敬拜你.
求你親自與我們每一個弟兄姊妹同在.
無論我們身處哪個地方.
求你自己仍然在我們生命當中坐和.
幫助我們成為你自己的子民.
成為你的僕人.
去聽你的說話.
去遵你的說話義行.
我們這樣祈禱奉主名求.
好,我們完結了剛才的快問快答.
我知道有些人說不是我贏嗎?.
為什麼去了另一個?.
不要緊的.
如果剛才贏了我們宣佈的那些.
請你email我們info@urbanvoicechurch.com.
你說我贏了.
我們有紀錄的.
所以我們將優惠碼給你.
其實都是十元的優惠碼.
我們每一個報名去參與課程.
連成會的人都是.
他們都沒有什麼權利.
他們去報名.
他們就能夠一起在牧場教室學習上帝的話語.
在這裡有三分鐘的時間.
告訴大家有什麼新的動向.
我們網上的進修課程.
剛才麗仙都已經報告了.
你去到我們的slashclass的網頁.
你就可以找到.
你就可以去報名.
譬如你報名參加這個.
你就可以輸入特價碼.

$^{601}$沒有都不要緊.
其實都已經很便宜.
因為我們的意象是什麼呢?.
我們希望我們的主流學或者網上課程.
能夠是affordable(可接受).
和flexible(彈性).
正如我們和Northern Seminary(北東教會)合作的時候.
大家的意象其實都是一樣.
我們希望能夠使全球的人.
發展這個士工的時候.
讓每一個神國的子民都能夠得到助紂.
好,這個是我們新的一年.
2021年的時候.
我們希望能夠做到的.
今天是live(直播).
我們希望大家是ok(可以).
是喜歡的.
如果你覺得繼續可以做live(直播)的.
請你給我一個1好嗎?.
我知道有很多朋友已經register(登記)了class(課程).
Vicky已經register(登記)了.
也有一些朋友問一些class(課程)的問題.
2021年我們就會繼續有live(直播).
同樣我們有新的series(系列)叫做《舊約》網上行.
因為去年.
不是去年,是今年.
去年年中的時候就做了《聖經》的網上行.
《部主腳鐘》.
《耶穌的生平》.
我決定一月開始.
我們就會有一個《舊約》網上行.
由《舊約》開始.
我們一起直接Google Earth.
來看《聖經》的故事.
下個禮拜.
12月31日.
禮拜四.
我們就有一個countdown(倒數)的時間.
希望能夠在這裡.
給大家一些網上的資源.

$^{641}$特別是我們網上的課程.
能夠幫助大英姐妹.
在不同的地方.
我們都能夠得到主的祝福.
在這裡我想點一首歌給大家.
是Brenda為我們預備的這首歌.
《耶穌的生平》.
耶穌萬有主宰.
惟有你之榮 此生尊貴.
榮耀憐憫公義.
來讓我向你舉手讚美.
耶穌獨有的主.
從沒有哪個能與你相比.
傳遞跪拜呼福.
齊眾讚向你屈膝敬拜.
懷氣當耶穌寫生.
講悲事拯救.
願意犧牲一切救贖我.
並賜予聖靈光照與報受.
重新建築何須生命.
懷氣當耶穌掌管世間所有.
全命大能都歸我主.
惟有你背受 懂得的稱頌.
永美聖潔永貴與你.
懷氣當耶穌掌管世間所有.
全命大能都歸我主.
惟有你背受 懂得的稱頌.
永美聖潔永貴與你.
惟有你背受 懂得的稱頌.
永美聖潔永貴與你.
\newpage



\section{}
\label{sec:n3sfAz_OcWA}
\textbf{舊約網上行 - 上帝的足跡:EP1 起初}
\newline
\newline
連結: \href{https://youtube.com/watch?v=n3sfAz_OcWA}{\texttt{https://youtube.com/watch?v=n3sfAz\_OcWA}} ~~~~ 語音日期: 2021-01-12
\newline
\newline
\hyperref[sec:eJV9zLgjq90]{\small{< < < PREV SERMON < < <}}
~
\hyperref[sec:index_chronic]{\small{[返順時目]}}
~
\hyperref[sec:index_scriptual]{\small{[返順卷目]}}
~
\hyperref[sec:lVpUws1JssE]{\small{> > > NEXT SERMON > > >}}
\newline
\newline
$^{1}$歡迎各位參加《舊約》網上行.
繼上一次《聖地》網上行.
部主腳鐘.
我們一連十四集追蹤主耶穌的生平.
我們這一次再整裝待發.
就算疫情還沒退.
我們仍然被困在家中.
我們仍然可以在旅程裡面.
我們去查訪一下.
聖經舊約的世界.
以及讓上帝在舊約歷史當中留下的足印.
陪我們走過這個逆境.
網上行有一個很重要的 idea(概念).
不知道大家記不記得.
就是每一次我們看聖經的時候.
就有如去一個陌生的地方旅行一樣.
我們不可以用自己的觀念套到聖經上.
要先明白聖經當地的風土人情,背景等等.
因為聖經是神對人類的啟示.
神向人溝通的時候.
是用當代人能夠明白的語言,概念等等.
所以我們先要明白聖經古代的世界.
今天我們才能夠明白.
神藉著本聖經想跟我們說什麼.
這次是舊約網上行的第一集.
當然我們就由創世記開始.
說到創世記我們會想起歷史.
不單單是以色列的歷史,它的起源.
比如族長Patrius的歷史.
阿布拉罕爾薩雅國.
我們更加會想起萬有的歷史.
宇宙,人類的起源.
創世記所記載的就是神創造天地.
說到人類的開始.
人類如何喪失了Paradise,樂園.
以至離開了上帝.
人類文化的開始等等.
這些主題.
不單單是我們現代人切身感到興趣.
更加是古代世界文獻裡面經常會討論.

$^{41}$這就是創世記.
特別第一至十一章有的背景.
這個世界叫做什麼呢.
美索不達米亞.
Mesopotamia.
人類的文明,文字寫作,數學,天文地理.
耕種畜牧,商業法治,城邦的建立等等.
是由一個地方開始.
這個地方不是埃及,不是希臘,不是羅馬.
不是中國,不是印度.
而是美索不達米亞.
Mesopotamia這個字是兩河流域的意思.
這片土地位於底格里斯河和幼發拉底河之間的肥沃平原.
即是今天的伊拉克,波斯灣,科威特,敘利亞一帶.
人類定居在這裡有近於一萬二千年的歷史.
在古時這個地方氣候穩定.
土壤肥沃.
長年也有河水供應.
很適合農業,城鎮的發展.
所以有史以來這個地方就出現人類的文明.
也出現人類文明第一次的創新改革.
改變了以後人類的歷史.
大約在六千年前.
好像一夜之間這些小的農村.
發展成為世界歷史裡最早的城市.
公元前大概四千至三千年之間.
美索不達米亞這個地方冒起很多城市的幫國.
彼此來競爭.
曾經在阿卡德王薩爾岡一世.
薩岡第一世之下統一過.
然後分裂成為古巴比倫和古亞瑟帝國.
儘管戰爭不休.
這個地方仍然在文明上不斷發展.
這裡的人建立了宮殿.
一些大規模的建築.
包括神廟,神塔.
創世記第十一章曾經出現過巴別塔.
這些塔廟.
亦發展了算術.
創建了一分鐘有六十秒的時間單位.

$^{81}$一個小時有六十分鐘.
一個圓形有三百六十度.
亦有一些六十進制的計算方法.
巴比倫人用了這些複雜的數學系統去研究生長.
將地球分成十二個週期.
十二個月.
這是後期希臘人採用的方法.
造成今天我們稱為十二星座.
古巴比倫人亦將一週分成七日.
但最大的文明影響力.
就是美索不達米亞出現的文字.
最初可能是一些黏土上標記著一些很簡單的圖案.
做一些簿記.
一些bookkeeping.
但到了公元前三千年.
這些圖案就發展成了一種複雜和靈活的書寫系統.
我們後世人稱它為契形文字.
之後的三千年裡面.
這種字體就引申出很多不同的語言.
包括蘇美爾語言.
巴比倫 赫人 波斯文等等.
用途非常廣泛.
古巴比倫王漢莫艾比.
就是用契形文字形成了初代的司法系統.
這是漢莫艾比發典的石碑.
裡面有一些說話就像摩西五經的律法.
譬如是以眼還眼 以牙還牙.
作為文明的發源地.
美索不達米亞有不少文獻討論人類的歷史.
在某程度上是與創世記一至十一章所討論的相似.
都是講到世界人類起源的歷史.
但在講歷史的時候.
在近東的歷史寫作大概可以分為四種不同的種類.
或者我們了解一下的時候.
可以幫助我們解釋創世記頭十一章.
講到世界起源的經文.
哪四類的歷史寫作呢.
第一類我們可以叫當代歷史寫作.
Contemporary History.
即是講到的事情是近期發生的.

$^{121}$或者是仍然發展中的事態.
現代的例子就像一些新聞報道.
譬如上星期美國國會發生過的事情.
最近每個地方的疫情等等.
這些都是屬於當代歷史.
第二類的寫作是近代歷史.
就是Previous History.
即是比較遠一點的人事物.
很著重記載哪年哪月哪日發生過什麼事.
譬如聖經裡面的一些例子.
聶王記提到一些外邦的君王.
來攻打耶路撒冷.
埃及王,亞述王等等.
所以這些記載在經外的文獻是可以做參考的.
第三種的歷史寫作.
就是遠古歷史.
Legendary History.
就是講到一些時間性更加遠古的人事物.
很多時候記載的時候的描述.
都是添上了一種英雄傳記的色彩.
現代的例子也有很多.
譬如美國的歷史.
講到林肯總統.
開國的元老.
Founding Fathers的歷史人物.
都是有一種文化英雄的色彩.
那種Cultural Hero.
聖經裡面的例子也有.
譬如阿巴拉罕.
信心之父.
以色列的族長.
帶有一種傳記的色彩在書中.
第四種的歷史寫作.
就是原始歷史.
Proto History.
越來越遠的時候.
目的不是記載某年某月發生了什麼事.
這些不是寫作所有的焦點.
目的不是記載事件的詳情.
而是用歷史來教導世界觀.

$^{161}$歷史去教導宇宙的觀念.
教導你如何去看待這個世界.
譬如有什麼例子呢.
近東文獻裡面有巴比倫的神話.
講到世界的起初.
就是一種教導世界觀的歷史寫作.
帶有神話性的.
好像Mythic一樣.
所以有一種超現實的描述在書中.
講到歷史的時候.
想起創世記1至11章.
就是近似這種原始的歷史.
所以我們不需要太著緊.
這些源在哪裡和什麼時候出現過.
因為這不是寫作的目的.
目的在哪裡呢.
就是想給我們一個世界觀.
猶太人的拉比說得很對.
他說創世記是用故事來講哲學.
用故事和歷史的事蹟來講世界觀.
所以要讓你去了解.
上帝想你了解的世界觀.
要你明白你在世界上存在的目的.
你的角色.
講到宇宙觀世界觀.
絕對不只是古人才有興趣.
現代人也很有興趣.
我們講到世界觀或是宇宙觀的時候.
現代人也有這種看法.
可能你也聽過1977年的時候.
美國太空總署發射了兩部太空船.
叫做Voyager 1和2.
這兩部小的太空船一直在外太空進發.
越來越離地球遠.
時至今日仍然在服務.
已經有43年以上.
不斷傳送外太空資料和相片回到地球.
直到2020年12月11日.
這兩部太空船距離地球大約有22億公里.
剛剛離開太陽系的時候.

$^{201}$當時是1990年2月14日清人節.
Voyager 1是最後一次能夠看到地球最後一面.
拍了一張相片.
那張相片就是一個世界觀.
在一個浩瀚的太空裡面.
地球只不過是一個很渺小的藍色一點.
Pale blue dot.
這就是現代人的世界觀.
近東的創世故事一直都是教一個宇宙觀.
一個世界觀.
所以如果這樣看的話.
創世記頭一二章其實都是有同一個目的.
我們先讀幾段經文.
這裡講到伊甸園.
伊甸園是一個怎樣的地方.
第二章第八節開始.
耶和華神在東方的伊甸立了一個園子.
把所做的人安置在那裡.
耶和華神把各樣的樹從地裡長出來.
可以閱人眼目.
其上的果子可作食物.
園子裡面當中有生命樹和分別善惡的樹.
有河從伊甸流出來支援那園子.
從那裡分成四道.
第一道叫做彼順.
是環繞哈菲拉全地的.
那裡有金子.
在那裡金子是好的.
有很多寶石.
第二道叫做基墳.
是環繞古實全地的.
第三道叫做底格里斯.
Tigris.
是亞述東邊的.
第四道叫做幼法拉底河.
耶和華神把那人安置在伊甸園.
使他修理看守.
在這段描述.
伊甸園是什麼地方.
伊甸這個字.

$^{241}$就是喜悅的意思.
Delight的意思.
所以我們今天叫做Paradise.
留意這裡的描述.
伊甸園也是河流的源頭.
分成四道支援全地.
也有樹是一個盛產有生命的地方.
有寶石有黃金.
回到第二些經文.
以西結書第28章.
說到太二王一個很有錢的海港王.
他說曾經將伊甸.
上帝的原子裡面.
放他在那裡佩戴各種寶石.
又有藥櫃上的基魯柏.
將他安置在哪裡.
神的聖山之上.
所以說到伊甸的時候.
就是說伊甸是一個聖山.
一個Holy Mountain.
回到創世記第三章.
人類墮落反叛之後.
人類的代表亞當被趕出去.
又在伊甸的東面.
上帝安置了基魯柏.
四面轉動發火焰的劍.
把守著生命樹的道路.
所以伊甸園的門口是向東的.
是向日出之處.
第二章剛才看的經文.
人被安置在伊甸園做什麼.
修理看守.
這兩個字是在舊約聖經裡面.
說到祭司打理聖殿會幕的工作.
這些特點全部都是近東的聖所的描述.
說到神廟的時候都是這樣去形容.
所以伊甸園是什麼地方.
是聖所 是聖殿.
是Holy of Holies.
是至聖所有上帝同行同在的地方.

$^{281}$人類被造的時候.
原本的目的 原本的身份.
就是上帝的祭司.
我們被造的時候.
我們敬拜有君尊的祭司來事奉神.
我們再看其他經文.
看第一章.
第一章一至到第五節.
起初神創造天地.
地是空虛混沌 淵面黑暗.
神的靈蘊能在水面上.
神說要有光就要有光.
上帝看光是好的.
就將光暗分開.
神稱光為晝 稱暗為夜.
有晚上有早晨 這是頭一天.
我們再看創造日的第七天.
二章第一節.
天地萬物都造齊.
到第七天.
上帝造的功已經完畢.
就在第七天揭了他一切的功安息了.
上帝賜福給第七天.
定為聖日.
因為在這天上帝揭了他一切創造的功.
就安息了.
這些描述肯定不是科學的報告.
所以不需要用我們現代人科學的眼光去看.
還有這裡有很多文學的修辭.
你會留意到第一至第三日和第四至六日.
所做的事都是平衡的.
最後第七日.
就是最精彩的.
就是上帝安息.
安息的意思是進入聖殿.
登基作王的意思.
並且在創世記第一章的創世記創世記第一章所描述.
有很多七或七的倍數.
舉個例子.
一章第一節在希伯來文有七個字.

$^{321}$所以中文翻譯也很好.
起初神創造天地.
七個字.
第二章第二節有十四個字.
七個倍數.
或者天和地各自出現了二十一字.
是就者人成了這句話出現了七字.
身看著是好的.
七字.
是一些很精密很絕美的文學創作.
古代的近東只有一類文件.
經常會用七段來形容.
就是講一些建築物.
一些Architecture.
是甚麼建築物才會用七段來描述呢.
就是聖殿.
最簡單的例子就是所羅門第一個的聖殿.
在耶路撒冷所建造.
用了多少年.
用了七年.
再找舊約其他的經文.
想想哪裡會有經文用七段來形容一些工作.
我們看出埃及記.
二十五至四十章.
這裡說到甚麼呢.
上帝吩咐摩西和以色列人去建造回眸.
Tabernacles.
在以色列南部的曠野有一個一比一的回眸.
紀念以色列人在出埃及後.
在曠野經過的地帶.
由上帝在回眸裡帶領他們.
在這個一比一的模型裡.
可以在外面走進去外院.
然後在外院有祭壇.
再可以進去.
就是進入聖所.
Place of Holy.
進入聖所.
在聖所可以看到金燈台.
陳設餅等等.

$^{361}$再進入至聖所.
裡面有基督教所遮掩的藥櫃.
法板就在裡面.
這裡就是上帝的寶座.
出埃及的二十五至四十章.
描述回眸建造的時候.
耶和華神用了七段的說話.
在他的指示裡面.
教導摩西和以色列人如何建造回眸.
這七段的說話.
第一段是回眸裡面的設備.
譬如藥櫃,陳設餅,燈台.
如何點燈.
或者祭司所有的條例.
他們的裝飾等等.
第二段是回眸的捐獻.
然後是高柔,香等等.
最後第七段是安息日的律法.
七段的說話當中.
都是以一句開始.
就是耶和華曉遇摩西.
耶和華對摩西說話.
留意第一段.
講到回眸的設備的時候.
是和光有關.
譬如如何燃點七燈台.
留意第七段最後一段.
就是會莫見造指示的結束.
就是最高峰的時候.
頒布安息日的律法.
再提以色列人.
你要在安息日尊我為和.
來敬拜我.
你見不見到有一個這樣的平衡.
就是和創世記第一章.
創造天地的平衡.
又是七日又是七段.
第一天上帝做什麼.
上帝說要有光就有光.
創世記第一章的創造天地.

$^{401}$第七天就是安息日.
上帝是這樣作供.
所以以色列人都要似上帝的形象.
六天工作.
第七天要安息要去敬拜.
創世記第一章想說的.
就是整個宇宙就像聖殿.
好像回眸一樣.
是神聖的.
整個世界都是敬拜上帝.
上帝創造萬有美好的意思是什麼.
是Holy.
是Sacred.
再看近東的背景.
創世的神話裡面.
神明做天地之前.
一定是和大水的水神爭戰.
打敗了之後才會做世界出來.
例如巴比倫的神話裡面.
有一個很有名的叫做.
Illumina Elysia.
意思是起初世界未創造之前發生什麼事.
就是巴比倫的大神馬杜克.
他要擊敗水神.
那個邪惡的水神.
然後將水神的屍體分開.
將一半放上天.
一半放在地上.
創造天地將水分開.
這個就是創世記裡面的描述.
你記不記得.
做天地.
然後巴比倫的神.
就會在聖所安息.
去登基做.
巴比倫每逢新年的時候.
都有一個慶典叫做Akitu.
Akitu的意思是大麥播種收割的時候.
春天的時候.
他們就會背誦Illumina Elysia.

$^{441}$即是紀念到或者希望.
盼望這位神明.
每年再打一次仗.
去克服混沌大水.
然後他可以作王.
可以祝福我們接下來這一年.
創世記.
就是在一個這樣的背景裡寫成.
就是講出另外一種的世界觀.
你留意創世記一章第一字.
第一個字是什麼字.
希伯來文是Bericci.
中文翻譯做起初.
英文in the beginning.
不是我們現代人心目中.
時間的時間軸上的第一點.
Bericci是什麼意思.
in the beginning是什麼意思.
其實是兩個字加起來.
第一個是前置詞.
in是什麼意思.
即是再.
第二個字就是開頭的意思.
Bericci.
在舊約裡面出現過好幾次.
我們只看一字的經文.
耶利米蘇第26章第一字.
那裡講到猶大王.
約西亞的兒子.
約瓦敬登基的時候.
Bericci.
就是他登基開頭的時候.
中文翻譯登基就很硬.
君王開始作王.
要去整頓他的王國的時候.
英文翻譯應該是.
in the beginning period of the reign of the king.
即是說君王作王的時候.
他去管理他的國度.
kingdom的時候.

$^{481}$你看看創世記第一章第一字.
就知道他想講什麼.
上帝是君王.
起初神創造天地.
起初就是上帝是君王.
來整頓宇宙的國度.
地是空虛雲遁 冤冤黑黑.
神的靈運行在水面上.
神作王.
你留意那個描述.
有什麼是和文化不同的.
和巴比倫文化不同.
是沒有戰爭的.
宇宙裡面沒有任何的勢力.
可以去與上帝對抗.
大水 神的靈運行在水面上.
水是沒有神聖的.
然後記載太陽 月亮 星宿.
這些在巴比倫眼中.
都是一些神明要去膜拜的.
但在創世記第一章.
就完全沒有了神聖.
完全只不過是萬物的其中一個.
神說要怎樣就怎樣.
相比之下創世記的故事.
神創造萬有的時候.
一切都是平靜.
一切都是順上帝的旨意.
事就這樣成了.
有七次講到.
你要看明聖經的話.
你就先要過這一關.
先要看清楚創世記第一章.
想講的神學.
就是你要認識.
我藉著聖經啟示的神.
祂才是真正的主宰.
真正的君王.
唯一一個值得敬拜.
唯一一個值得去追求.

$^{521}$去景仰的.
不是地上的君王.
不是巴比倫帝國.
不是政權.
不是地上有勢力的東西.
或者吸引我們的東西.
搞清楚.
唯一一個要去敬拜.
要去進程去事奉的.
是一個創天造地的上帝.
除祂以外沒有別的神.
會不會去到動盪不安的時候.
可以試到我們.
到底真正敬拜的是誰.
今天是2021年1月10日.
上星期大家知道美國發生什麼事.
美國的政局.
以至影響整個世界的政局.
都有不安有所震盪.
我有一個牧師朋友.
是白人來的.
他住在美國南部.
他是很熟悉政治的.
美國的政治.
他甚至說自己是政治人物.
冷飯菜汁.
每樣東西都混在一起.
住在南部.
你想像到他的政治環境是什麼.
但是這位朋友.
他時常都強調.
我不屬於哪一個黨.
因為作為基督的門徒.
我是獨立的.
他前兩天在他的Facebook發帖.
他說早在20年前.
他已經不再追聽極端的政治論壇的電台.
或者今天的YouTube.
原因就是.
他以前發現.

$^{561}$他聽這些節目的時候.
他發現這些今天叫KOL.
是很擅長做什麼的.
就是煽動聽眾的情緒.
他也不會看黃金時段的政壇節目.
因為他發現幾分鐘的報導裡面.
很多時候都是誇大的.
將小小的事實發展成一些陰謀.
或者一些荒謬的論點.
目的是什麼.
目的就是去煽動情緒.
就是等你看完今晚.
你很急不及待.
再看明天.
明天再不及你再看下一天.
就有一點點像吸毒.
他在Facebook裡面.
他就說.
奉勸各位.
奉勸他的朋友.
不要被這些煽動情緒牽制.
其實如果這樣做的話.
我們只不過是向這些所謂KOL.
這些極端論壇的人.
他的金庫來進貢.
不要為了他們所偏愛的政黨或者政客.
來奉勸我們自己.
我有一位朋友說.
如果你是基督徒的話.
希望你多用一點時間和精力.
去讀透觀聖經.
去吸收上帝的話語的資訊.
使你成為一個更好的門徒.
去參與一些基督徒的群體教會.
去學習事奉.
總比花心情,花情緒,花這麼多體力.
為了他們喜歡的政客去護航好.
如果你不是基督徒的話.
我邀請你和我們一起同行.
認識這位創天造地的主.

$^{601}$你認識的當中你會發現.
我們教會真正的門徒.
並不是一些完美的形象.
因為我們都是有破碎,有罪性的人.
我們活在一個有罪的世界.
我們只不過是不斷跌跌碰碰.
我們學習希望明天比我們昨天更懂上帝恩典.
其實你想想這番說話很有意思.
特別是今時今日.
無論你住哪一個地方.
你想想地上的權力,勢力.
最終想做的都是吸引我們.
吸引我們去追隨它.
吸引我們去敬拜它.
然後慢慢成為我們心中的偶像.
就是要我們向它效忠.
向它獻上我們的敬拜.
當這些有權勢的人.
他們做對事的時候.
我們就要去稱讚他們.
當他們做錯的時候.
我們就自然要花心力去為他們護航.
你問心一層.
會不會是我們給他們太多的注意.
太多的榮耀.
根本他們就不配.
又或者.
你可能聽過一些故事.
我也聽過不少.
一些的牧者或者一些的傳代.
起初的時候可能是窮傳道.
一個窮學者.
是有赤子之心.
不計較名利.
但是做著做著.
多年之後得到賞識.
甚至得到國家的推捧.
慢慢慢慢就變了.
古今中外這些故事.
比比皆是.

$^{641}$地上的權勢.
到底有沒有消耗我們太多的心力.
我不是說不理世事.
而是有沒有消耗我們的心力.
而是令我們沒有盡心盡力去敬拜事奉神.
沒有去問清楚.
到底我們真正敬拜的那個.
是不是創天造地的主.
是唯一的.
可能不只是說到一些權勢的問題.
說到消耗我們.
可能是一些最受歡迎的男女歌手.
一些地上的偶像.
我們問一問.
到底我們真正敬拜的是誰.
真正決定我們今天是否平安.
真正決定我們今天有沒有喜樂.
有沒有生命.
是不是世上這些人.
人生為了什麼.
這就是創世記第一章.
想去挑戰我們.
問一問你自己.
在去年或最近.
消耗你最多的心力.
最影響你心情上落的.
到底是不是上帝的名聲.
是不是上帝有沒有被人敬拜.
被人尊崇.
還是其他東西.
今年上帝給我一個挑戰.
我也向上帝許願.
我神可以盡力來事奉你.
求你來用我.
一個謙卑的.
一個很謙卑的.
很不配的一個器皿.
求你來讓我事奉你.
讓我敬拜你.
讓我盡心盡力來傳揚你的名.

$^{681}$我希望這個許願.
這個立志.
也是我們每一位弟兄姊妹的立志.
讓上帝的名在我們心所想的.
我們生命裡面最終想追求.
因為祂值得我們敬拜.
我們一起低頭祈禱.
神啊你是那位創天造地的主.
雖然世界因罪而得到破壞.
但是你仍然救贖.
仍然讓我們能夠去敬拜你.
主啊我們求你來幫助我們.
做我們在生之年.
無論我們現在的境況是信還是疑.
主幫助我們盡心盡力.
來慰藉你的名聲.
慰藉你的聖名.
來奉上最好的.
盡心盡力來敬拜和事奉你.
以致我們生命能夠成為你的僕人.
能夠去服侍你.
能夠去尊崇你.
祈禱奉主名.
\newpage



\section{}
\label{sec:lVpUws1JssE}
\textbf{神的名字6\_耶和華賜平安 Yahweh Shalom (靈思牧禱)}
\newline
\newline
連結: \href{https://youtube.com/watch?v=lVpUws1JssE}{\texttt{https://youtube.com/watch?v=lVpUws1JssE}} ~~~~ 語音日期: 2021-01-13
\newline
\newline
\hyperref[sec:n3sfAz_OcWA]{\small{< < < PREV SERMON < < <}}
~
\hyperref[sec:index_chronic]{\small{[返順時目]}}
~
\hyperref[sec:index_scriptual]{\small{[返順卷目]}}
~
\hyperref[sec:u79BpJLNaZg]{\small{> > > NEXT SERMON > > >}}
\newline
\newline
$^{1}$當說到平安 我們的期盼是甚麼呢?.
或者農曆年的揮春可以解答這個問題.
就是「催促平安」,「四季平安」,「老少平安」.
再加上「平安大吉」.
簡單來說就是「無風無險」.
「人生一帆風順」,「興旺安泰」.
當然 這樣的平安 我們人人都嚮往.
說起來 在聖經裡面「平安」「Salom」這個字.
希伯來原文的意思就是「圓滿」,「健全」,「興旺」,「安康」.
比較我們中國人對平安的期盼.
其實是有過之而無不及.
只可惜 在我們人生的經驗裡面.
持久的平安似乎是不可能實現的.
圓滿的平安似乎是失落了.
為甚麼呢?.
基督教的信仰對於這個失落的平安.
有詳盡的交代.
事實上 聖經一開始就揭示這個問題.
聖經第一卷書《創世記》.
蝦頭就說 神所創造的一切本來就是圓滿的.
但是因為人類的始祖在伊甸園收射的斯坦.
選擇了為抗神的命令.
自此 整個人類都墮入罪的深淵.
整個世界都因罪而殘缺 敗壞.
人彼此傷害 忍受痛苦.
這些不單止是遙遠的歷史.
就看我們今天人際關係.
夫妻 父母 子女 朋友 同事 甚至是教會.
其中的傷害 甚至是痛苦.
都是真實體會得到的.
我們又看到 人因為無窮的慾念.
帶來自身的破壞 以及為害其他人.
都發生在家庭 工作場所 在城市 國家.
比比皆是 你和我都捲在裡面.
這些殘酷的現實 無情地告訴我們.
當人類始祖犯罪.
以至跟神的關係破壞的一刻.
亦都是人失落 平安的一刻.
聽起來好像很令人喪氣.
但是美妙的就是 在人失去平安的一刻.

$^{41}$上帝就立刻宣告.
祂說會賜平安之子 修補人類的破碎.
恢復失落的平安.
在這一集 我就是要和你分享.
上帝的另一個名字 Yahweh Shalom.
中文譯作耶和華沙龍.
亦即是耶和華賜平安的意思.
原來耶和華沙龍這個神的名字.
是《舊約聖經》裡的一個軍事人物.
基奠呼叫出來的.
而基奠是以色列人進入迦南地之後.
其中一位的士師.
即是當時的軍事領袖.
不過這個領袖的心情就不是那麼踏實.
因為他正面臨大敵 憂心忡忡.
士師記第六章記載.
基奠當時正在打谷.
打谷是很正常的衛生工夫.
只不過基奠打谷就在不正常的地方進行.
因為通常當時打谷是在室外.
而基奠卻在釀酒的屋裡打谷.
是因為他和他的百姓都很驚怕.
他的兇猛的敵國米奠人.
隨時攻佔他們的土地.
以及毀壞他們的農作物.
而奇妙的事就發生了.
就在這個時候.
上帝的使者來探訪滿心恐懼的基奠.
對他說.
大能的勇士要說和你同在.
他必幫助你戰勝米奠人.
拯救自己的同胞.
照常理.
看到上帝的使者聽到這句話.
基奠應該是很興奮激勵的.
奈何他心裡因為太多疑惑恐懼.
而使他堅持要見到神蹟才相信.
只見神蹟很快就出現在他眼前.
就是使者伸出手裡的匠.
用火從盤石出來燒盡了祭壇上的玉.

$^{81}$以及無烤餅.
這個時候使者也消失了.
基奠終於相信了.
但是他心裡更加驚惶.
因為他想到.
我這麼多質疑冒犯了神的使者.
更加冒犯了上帝.
我死定了.
不過也在這個時候.
上帝親自對他說.
你不用懼怕.
你必不自死.
就在這個和上帝的相見.
使得基奠豁然大悟.
原來自己是這麼信心軟弱.
自己是這麼不配.
聖潔的上帝仍然無條件赦免和接納他.
因此他就在各個地方築起一座壇.
紀念他親身體會到的耶和華慈平安.
而基奠所體會到的耶和華慈平安.
其實是回響了在伊甸園裡面的景象.
就是即使罪人是多麼污穢不配.
多麼軟弱.
上帝願意修補人類的殘酷.
恢復人失落的平安.
今天你和我其實也像基奠一樣.
我們可以體會到.
從上帝而來的這一份平安.
就像我開始的時候所說.
我們今天說平安.
往往是期盼人生一帆風順.
興旺安泰.
可惜的是.
在這個因罪而殘缺敗壞的世界裡.
我們得不到這樣的平安.
而可喜的是.
上帝主動來修補.
使我們重新得到這種平安.
如何達成呢?.
就是上帝賜下他愛子耶穌基督.

$^{121}$在十字架上為人承傳了舊恩.
拆毀了人和神因.
罪而隔斷的圍牆.
上帝和人的關係因而得到恢復.
人重獲平安.
難怪耶穌基督被稱為和平之子.
雖然聖經說的平安並不是無風無險.
就像今天.
你和我永遠的抵受疾病或壓迫.
甚至戰爭等等的災害.
但是因著和上帝的恢復關係.
我們就能夠嘗到伊甸園當初那份.
圓滿就是神所賜的平安.
而且藉著耶穌基督.
這個破碎墮落的世界.
有一天都會過去.
我們會看到新天新地.
就像啟示錄21章所說.
神要拆去他們一切的眼淚.
不再有死亡.
也不再有悲哀,哭號,疼痛.
因為以前的事都過去了.
在這一天.
我們將會完全恢復真正的平安.
就像Shalom希伯來愚民所說的.
圓滿健全興旺安康.
多謝您收睇時局新聞,再會!.
\newpage



\section{}
\label{sec:u79BpJLNaZg}
\textbf{經緯線 - 耶穌的比喻 (EP1 撒種的比喻)}
\newline
\newline
連結: \href{https://youtube.com/watch?v=u79BpJLNaZg}{\texttt{https://youtube.com/watch?v=u79BpJLNaZg}} ~~~~ 語音日期: 2021-01-13
\newline
\newline
\hyperref[sec:lVpUws1JssE]{\small{< < < PREV SERMON < < <}}
~
\hyperref[sec:index_chronic]{\small{[返順時目]}}
~
\hyperref[sec:index_scriptual]{\small{[返順卷目]}}
~
\hyperref[sec:xKw__Td71oI]{\small{> > > NEXT SERMON > > >}}
\newline
\newline
$^{1}$各位歡迎來到《經緯線》.
這是網上的《釋經時評》.
我們這個新的系列.
就是《新約聖經》.
大家可能看過我們一些約拿書的《釋經時評》.
我們開始一個新的系列.
所以有舊約也有新約.
所以今天我們去查一些新約的經文.
以致大家能夠認識更多.
我們先介紹一下自己.
我是Mark Sir 陳重基.
我是城會社區教會的牧師.
城會網台是我們教會的事奉.
侍工的其中一個.
我們希望能夠做到一些這樣的視頻.
讓更多人認識主.
其他三位嘉賓介紹一下自己.
我先說吧.
我是Pastor Jackson.
我是宣道會石安堂的主任牧師.
很開心有機會和大家一起來到賢道神話語.
我是Zelwyn.
我20年前回教會的時候.
經常問很多問題.
問著問著就發覺有時候問太多問題.
有時候不可以問太多問題.
到現在自己做牧師.
去席堂.
發覺都很需要創造空間.
大家一起討論聖經.
不要怕問問題.
很重要.
不怕問題少年少女.
大家好.
我是Joyce.
我之前在城會回教會.
現在搬來加州讀書.
這次很開心有機會和三位牧師一起看聖經.
平時都很少機會可以直接問問題.
所以.

$^{41}$不怕做問題少女.
我們這個色經視頻的系列.
是一個新的系列.
大家可能看過約拿書.
舊約的茶經色經的視頻.
這四個人會開一個新的系列.
這個系列是處理查考耶穌的比喻.
說到比喻的時候.
大家應該很熟悉.
我想兩位牧師.
或者Joyce.
都應該查過.
或者甚至講道講過.
沒錯.
我講過了.
《馬太福音》十三章.
二十五章的比喻.
我也有一個比喻的系列.
是的.
Zelwyn應該也講過了.
有講過有查過.
Joyce應該也查過了.
我們處理《馬太福音》十三章的一堆比喻.
這次我們講第一個比喻.
就是薩種的比喻.
簡直熟悉到爛.
十三章一.
非常熟悉.
糾結一下.
閉上眼睛都應該可以知道內容了.
有一個薩種的人出去.
他就薩種.
薩種在不同的地方.
這個其實有點像巴勒斯坦的地理.
因為你見到那些田園很多時候在山邊.
所以薩種的時候.
有些落在路旁.
有些在荊棘.
有些就在石邊.
然後有些在好的土地.

$^{81}$是的.
我也經常聽到.
一講就說有四種土.
如果薩種在路旁.
真的聽了.
但是聽不入耳.
馬上忘記了.
或者薩種在石地或者淺土.
就是聽了.
很開心.
但是很快有迫害.
不開心的東西.
很快就不可以繼續生長.
第三就是在荊棘上面.
聽了很想做的.
很想跟著神的話要做的.
但是就因為給錢.
或者事業.
各樣東西.
就不可以繼續結出果子.
最好的就是在好土上面.
可以結果實.
通常我覺得記得查這段經文.
通常就會問.
大家是哪一種土呢?.
通常我真的沒有聽過.
或者我自己也沒有說過.
我在好土上面.
就會說.
不敢說.
自己說自己就是好土.
你也好土.
我想大家是很正直.
承認自己一定有金錢的顧慮.
金身的憂慮.
其實是很真的.
Joyce你講得.
因為是很好的討論.
例如說突然之間.
因為耶穌也解釋了比喻的意思.

$^{121}$特別講到世上的思慮.
錢財的迷惑.
大家就會很自然地.
對,很有認同.
我上班遇到很大.
我們上不了教會.
所以人們就會很自然地談到.
還有一些兒女.
如果家裡有兒女.
又說我遲回教會.
就變成了這樣的情況.
我們剛剛留意到.
我們有四格.
四個人.
我們代表四種土.
哈哈哈哈.
但是我們這個《釋經》視頻.
我們想帶出什麼呢.
就是我們不是一種零修飾的讀經.
給大家一些《釋經》的工具.
或者《釋經》的approach.
一些進路.
其中一個就是看context.
就是看上下文.
這個很多時候我們講道.
甚至查經.
我們不能夠處理.
我就留意到.
這裡撒種的比喻.
講了比喻本身.
其實很有趣.
這裡其中一些特別的地方就是.
耶穌解釋了比喻的意思.
原來有些上下文.
可能我們很多時候會錯過.
大家會嗎?.
會的.
其實在我們看.
剛才Joyce開玩笑的時候.
也不是開玩笑.

$^{161}$我們查經的時候也是很認真.
可能一開始就會控制著別人.
先不要說應用.
先不要說你是什麼土壤.
查經早就還沒開始.
是啊.
都看一看《聖經》.
但是其實有一樣東西可以幫到大家的.
就是我想很簡單的一個觀察.
就是第十三章裡面.
大家都會留意到.
其實一連串的比喻.
耶穌一直在講比喻.
這個已經是一個很好的觀察.
然後我們再一路放大它.
另一個放大就是你發覺.
十三章一開始的時候.
就說到有一群人來到耶穌面前.
耶穌就開始做這個教導的工作.
其實這個開始很特別.
原來《聖經》的學者就將這些教導.
其實《馬太福音》裡面有很多耶穌的教導.
不過《馬太福音》特別的地方就是.
耶穌的教導.
就是把馬太放在一起.
所以就有五次一些很長的教導.
舉個例子.
就是一個我們大家很熟悉的登山寶訓.
是的.
從五章到第七章.
大家都知道是一個很長的一段教導.
登山寶訓是第一個.
第二個其實大家也很熟悉的.
就是第十章.
《馬太福音》的第十章就是耶穌差遣門徒.
一段很長的說話.
這個就是第二段.
第三段就在我們這裡.
《馬太福音》的第十三章.
就是一連串耶穌講到的那些比喻.

$^{201}$然後到第十八章.
第十八章給大家一些提示.
就是門徒差誰是最大的.
然後也講到瑤書的功課.
其實那裡講的是什麼呢?.
十八章就是講一個信仰的群體應該是怎樣相處的.
所以這個就是十八章.
然後到最後.
也是大家很熟悉的.
就是第二十四至二十五章.
就是鑒欖山上面.
耶穌講到一些教導.
譬如說來的時候有些人不見了.
那些是不是大家很熟悉呢?.
所以總共是五段的聖經.
而在這五段很長的教導裡面.
大家留意十三章的教導正正在正中間.
多謝Jackson給我們整本書的脈絡.
就是五大教導.
現在看比喻.
十三章原來就好像整個《馬太福音》的高峰在中間.
在中心.
有很多的比喻在這裡.
最出名的比喻就是行仙.
就是撒種的比喻.
不過我留意到第十節.
那裡說有耳可聽的就當聽了.
第九節.
然後第十節.
門徒進前來問耶穌.
對眾人說什麼用比喻呢?.
然後耶穌就給他解答.
我覺得這個設定之前是很公開的.
每個人都聽到的.
然後到第十節突然間.
有些門徒跑去問耶穌.
問問題.
問問是少年.
是不是也很有趣呢?.
好像很神秘.

$^{241}$耶穌又不跟每個人說.
只是跟門徒說.
只是跟他們解釋.
什麼天國奧秘.
我想這個很值得.
很多時候被奧秘的字吸引了.
就像剛才大家所說的.
這些只是圈內人才可以有的.
只是細小的群體.
他所揀選的.
或者甚至是用了一個predestination的角度去看.
就是上帝欽點了.
你要知道就知道.
不知不知就沒得選.
將整個球交給上帝.
但我想其實很值得去看.
剛才Mark Sir說的.
其實繼續下去.
第十五節.
其實那裡是在引用二創夏書.
或者是paraphrase.
或者是直接引用的.
他也是在說二創夏書一個很重要的經文.
而他在說什麼呢?.
其實就是二創夏書第六章.
我不知道大家有沒有聽過.
有時候你看新約聖經.
如果你每次看到有舊約聖經.
有人這樣說.
你就不應該繼續skip掉它.
最喜歡的.
不明白的.
最需要的反而就是在這個位置停一停.
回到舊約.
沒錯.
還要像剛才所說.
不單止看那一節.
這是直接引用.
第九節.
第十節.

$^{281}$看這兩節算了.
也不是.
你應該要去看那一段經文的上下文.
舊約的經文的上下文.
對,這才是很重要.
當然我這段時間不是足夠講得很詳細.
其實如果你很快看一看二創夏書第六章.
其實那裡是在說.
以賽亞領受上帝給他的calling.
作為一個先知進入到以色列民.
當中宣講上帝說話.
但是這群是什麼人呢?.
如果你看第九第十節.
其實那段話也很灰.
是說這群人是怎麼也不明白的.
你告訴他什麼他也是當聽不到.
耳邊風.
所以二創夏十一節繼續問.
要到什麼時候呢?.
我們也很熟悉的.
就是到成一方來無人居住.
是非常絕望的.
是被老之後才.
沒錯.
但是其實也有一點點的希望.
非常非常少.
但是其實是存在的.
因為到了十三節.
最後他也說是一個聖潔的種類.
或者新日本的聖潔的苗裔.
是說有一個holy seed.
一個sacred seed.
就是上帝仍然留下一群人來跟從他.
Selvin你說得很對.
因為馬太福音十三章就是在引用二創夏書第六章.
但是很多時候我們看二創夏書第六章.
我們看一到第八節.
我再請餐喜我.
宣教大會就結束了.
原來下面那一節才是戲肉.

$^{321}$沒錯.
就是在說以色列民不聽上帝的道.
這也是.
不僅新約.
在比喻上看夏文.
舊約回頭.
我們也要看上夏文.
你剛剛說的這件事.
第一如果誰在宣教大會.
還是培靈會.
要差遣我.
就像你所說的一定要聽下一句.
差遣你的時候.
如果沒有人聽你的話.
沒有結果的.
先救不活他.
但其實你看回十三章.
耶穌也是這樣的.
剛才你說門徒進前.
一開始說這群人來聽.
再之前十二章尾.
是說誰是我的母親弟兄姊妹.
真正的家庭成員是誰.
其實說了一個遵行我天父旨意的人.
十二章五十節.
但到了十三章尾.
又說了這個先知本成本沒有人尊重他.
其實耶穌也是面對這樣的情況.
我想這裡其實說的是.
耶穌就好像去舊約的時候.
一個二創先知一樣.
其實耶穌就是有一個先知的身份.
他來到其實是宣講上帝的審判.
他的審判的臨到.
但同一時間明知道這群人.
都是不聽的.
都是有這麼深的反應的時候.
他也選擇了用這個比喻.
其實這一段是說出整個原因.
原因不是不想他們聽.

$^{361}$剛才我們說的小軒子就不聽.
外國人就不聽.
其實不是的.
其實他是反而用這個方法.
一個prophetic tools.
來幫他們去明白.
或者這麼說.
未必一定每個人都聽.
但當你有心去問的時候.
你就可以知道多一點.
所以他這裡有一句這樣說.
是不是第十二節.
凡有的都要加給他.
沒有的連他有的都要拿走.
當他肯去seek for more.
去尋求問多一點的時候.
那個revelation就會臨到他那裡.
他可以更加明白.
這句話經常都亂抽出來.
拿來應用.
凡有的就要加給他.
有錢的就繼續給他.
沒有的什麼都要奪去.
資本主義的社會就是這樣.
其實是說你要明白天國的奧秘.
明白kingdom of god的時候.
不單止是說真正的明白.
他這裡道出.
其實你看回十三到十五節.
是一些重覆的字.
其實核心是指向.
什麼是真正的聽見看見.
其實不是水掛背.
聽完就算.
不是這樣.
而是真正進入你的心.
真正的明白過來.
這個才是整個引導我們去.
要有的回應.
我想換句話說.

$^{401}$因為剛才Joyce說了一句話.
我都在想.
她說說到這麼神秘.
那群門徒來到.
但是Selwyn就分享給我們一個方向.
其實耶穌說的話是那群人的情況.
很多人都不會聽.
但是那些肯聽的人就有福了.
其實就是這個意思.
所以當她說到那群門徒走到前面的時候.
其實我覺得有一個象徵式的意義.
當然他真的走到前面.
走到耶穌的前面.
但是也代表了他們真的接近耶穌.
他們的身體語言告訴你.
他們很樂意聽耶穌說話.
打開了心門.
另外我看到的就是.
剛才Selwyn說的.
他們這群人.
究竟講不講給他們聽呢.
耶穌明知道他們不聽的.
耶穌都講給他們聽.
其實這個比喻的因由就是.
好像說他們聽不明白.
看不見.
那我講你做什麼呢.
這個就是我們的態度.
父母跟子女說.
講了十次都不聽我不講了.
讓你碰一下板吧.
對耶穌明知道他們不聽.
但是耶穌都講給他們聽.
因為耶穌很渴望.
如果你們能夠聽到.
那就好了.
所以其實這是一個警告.
但是也是一個勸勉或者鼓勵.
來吧聽我講.
我現在就告訴你.

$^{441}$這是第十節.
門徒進前來問耶穌.
剛才Jackson說的大圖畫.
登山寶藩之前.
其實也有一句這樣的說話.
門徒到他跟前來.
他就開口教訓他們.
就像講到這些弟子.
是一種理想的弟子.
我們應該走到耶穌面前.
主動地去聽.
是不是說好徒.
其中一個表意.
就是會去問問題.
會去想知道多一點.
很好啊你這麼說.
其實我在這裡想的時候.
是想知道多一點.
回想起剛才Joyce你一開始說的時候.
就是說.
那些人查經.
說上班.
最壞是工作.
以致我上不了教會.
就賴工作.
然後就賴子女.
讓我遲到上不了教會.
總之.
是很真實的.
我們說這樣說那樣.
到頭來其實就不關我的事.
但是當我們回看這段經文的時候.
你會發現將那個責任.
其實放在我們那裡.
我們選擇了.
工作而不上教會.
就是這麼一句.
我們選擇了兒女而不上教會.
所以責任不是兒女,不是工作,不是金錢.
是我們自己.

$^{481}$所以在這個經文裡面.
我們看到將那個責任.
在泥土那裡.
我們肯不肯來到前面.
我們肯不肯聽多一點.
看多一點.
其實在那個人的責任那裡.
沒錯,我想這裡其實就是.
一個上帝的宣判.
已經在那裡,就是結局.
其實那四種的土壤.
不是你選不選.
都是你怎樣去回應.
但是其實是告訴你.
結局.
當你肯聽的時候就是這樣,當你不肯聽的時候就是這樣.
這個這樣的.
一個這樣的宣判.
其實已經臨到.
而在這個過程裡面.
其實上帝是告訴我們.
正如當年的.
以色列一樣.
他仍然將一個.
中文這樣翻譯.
Remnant,正如的人.
仍然留下在這裡.
繼續去傳講.
他這個天國.
但是怎樣能夠成為.
其中一份子呢.
這個是就好像Justin這樣說.
up to我們來做一個選擇.
而在耶穌這個setting就看到.
當你可以去打開你的心.
來去聽.
當你願意.
不是只是聽完就算.
而是放在心裡.
而真的走出來.

$^{521}$遵從我天父的旨意的那一班人.
他就變成了上帝所揀選正如的人.
那我想這個是一個.
很震撼.
很需要我們去反思的一個問題.
很震撼.
我覺得其中一個.
可以應用這段經文的地方.
就是.
其實一個很好的例子.
現在很多人都透過Zoom.
因為疫情.
我們用很多網上的聚會.
小組,全契.
我們那時候.
到現在已經有七,八個月的時間.
到現在.
人們就會有一種感受.
很累.
見得太多了.
我們就不去了.
不要叫我了.
不要問我什麼了.
其實這些真的有危險.
因為我們應該是經常的聚會.
我們自己也要.
來自我反省.
究竟我是因為.
我不想見了.
我不想再看人了.
我不想再講信仰了.
我們自己要做一個反省和檢視.
究竟這些是不是我們的excuse呢.
所以我有個建議.
大家應該要留意.
如果有太多聚會.
可以cut down(停止).
我不是一個殘忍的牧師.
但是有些事情是不能cut的.
崇拜是不能cut的.

$^{561}$如果有聚會和弟兄姊妹溝通分享.
那些最好不要cut.
其他的都cut了.
除了這些之外.
我們知道了更多上文下理.
以及以賽亞書的背景.
對於parable(反省).
有沒有其他的反省呢.
是否人才或學業.
阻礙我們去神的面前.
我想可以看這個比喻的時候.
好像我們剛才一路強調.
就是上下文.
其實馬太福音.
它特意鋪排到一個地步.
就是引我們進入上下文.
讓我們知道.
到底我們應該用什麼心態.
去接受天國的來道.
或神的道的來道.
或耶穌的來道.
因為我們知道.
天主的道的來道.
像剛才Jackson這樣說.
有幾個block(區域).
現在我們處理的是最中間的區域.
這個區域其實是.
從十一章開始的.
其實一路十一章到這裡.
撒種比喻.
馬太福音安排了很多不同的故事.
那些故事都是和.
耶穌說完道之後.
人們的反應有關.
法利素人聽完.
耶穌說的話.
有些人會反對.
有些人會挑戰.
有些人不接受.
一路在準備.

$^{601}$聽撒種比喻.
問到底這些實際的故事.
是否我們的肅認.
我們是否聽到.
上帝的道的時候.
我們就反對.
甚至可能想殺耶穌.
想除滅.
其實是一個挑戰.
到底天國的來道.
是怎樣的人.
有人這樣說.
The only conversions that count.
in the kingdom.
are those confirmed by the life of discipleship.
這個可以.
總結整個我們今天所說.
究竟什麼為之.
真的去聽到耶穌的道.
什麼為之.
成為一個門徒.
不是你信了耶穌那一刻.
而是你一生可以活出.
這個真正.
上帝道理在你心裡面.
結出果子的一個生命.
希望大家聽到這裡的時候.
和我們一樣.
我們都是願意.
來到耶穌的面前.
去聽天國的道理.
走出來做天國的子民.
希望下一集能夠再見到大家.
多謝各位.
請訂閱支持啊.
謝謝.
\newpage



\section{}
\label{sec:xKw__Td71oI}
\textbf{舊約網上行 - 上帝的足跡:EP2 叛變}
\newline
\newline
連結: \href{https://youtube.com/watch?v=xKw--Td71oI}{\texttt{https://youtube.com/watch?v=xKw--Td71oI}} ~~~~ 語音日期: 2021-01-19
\newline
\newline
\hyperref[sec:u79BpJLNaZg]{\small{< < < PREV SERMON < < <}}
~
\hyperref[sec:index_chronic]{\small{[返順時目]}}
~
\hyperref[sec:index_scriptual]{\small{[返順卷目]}}
~
\hyperref[sec:ke055dTE6Oc]{\small{> > > NEXT SERMON > > >}}
\newline
\newline
$^{1}$歡迎參加《舊日網上行》.
我們這次是第二集.
在《網上行》我們就是去不同的地方.
去到聖經的世界.
我們去聽一下別人的觀點.
我們不要硬是將我們自己現代人的觀點.
套到聖經身上.
其實你想深一層.
《網上行》是一個超時空的旅程.
我們回到以前.
地理,考古等等的資料幫助我們.
去到古代聖經的世界.
我們先聽一下以前的人他們在想什麼.
不要覺得古代的人還沒有開發.
沒有文明.
我們現代人就不同了.
這個想法.
就是希望在《網上行》裡面我們能夠改正.
聽一下古代聖經世界.
直到他們來啟示真理的時候.
我們就能夠了解.
上帝今時今日會跟我們說什麼話.
《網上行》 歌手:李宗盛 詞:李宗盛 曲:李宗盛.
上一次去比對過一些近東的.
講到起源的歷史.
近東很多時候有一些神話.
但是我們發現這些記載,文獻.
他們的技術手法.
講的主題.
往往都有跟聖經有很多相似的地方.
但是你再想深一層.
其實本質上是有很大的差別的.
所以聖經就是跟當時近東的神話.
去做對話.
或者甚至做一些爭辯.
我們叫做批判.
或者是一種辯證.
就是說要跟他們爭辯.
要去抵抗或者抗衡他們的想法.
我們讀《創世記》一至十一章的時候.

$^{41}$我們上一次看過.
說可以算是近東的起源史.
Protohistory 或者Primeval History.
主要是講一些世界觀.
或者去處理一些人生的問題.
人生的意義等等.
第一樣處理的主題就是宇宙的起源.
這裡處理到底神跟我們之間有什麼關係.
我們看過近東的創世史詩.
Elumar Elish.
我們知道裡面講到.
神靈在創造這個世界之前.
有一個神靈的大戰.
就是巴比倫的神贏了.
他殺死了邪惡的水神.
把水分成上下.
然後宇宙按照這個想法.
本身是邪惡的.
正邪不斷地去鬥爭.
所以每年都要再做一次.
每年都要再打一次仗.
人之間亦是彼此的鬥爭.
創世記提出批判.
他說神創造一切都是好的.
宇宙原本就是有神的同在.
甚至整個宇宙都是上帝的聖殿.
第二方面提到的主題就是人類的起源.
人的本質.
到底是問我們自己.
到底我們的價值,我們的自由,我們的尊嚴.
到底是哪裡來的.
近東的文學.
很多時候都會講到這些神靈.
因為是多神論.
五時花六時滅的.
又彼此去打鬥的.
彼此去鬥爭的.
你最好就是拜對那個神.
如果不是的話.
你就會走霉運的.

$^{81}$人類呢.
為什麼會被造呢?.
就是神靈想造一些奴隸出來.
你都要服侍他們.
因為神靈也有自己的需要的.
祂會餓.
初一十五你就要上香.
你去餵飽祂.
希望祂能夠祝福你.
創世記就提出增辯了.
一章第二十七節.
神就照著神的形象.
自己的形象來做人.
乃是照著祂的形象做男做女.
這是一個很革新的思想.
就是說每一個人.
整個人類.
都是有神的形象.
代表神來管理大地.
第二章我們提過.
修理,看守,原子.
是祭司軍尊的職事.
就是說.
不只是當時的軍王才可以說我們是天子.
我們才是神的形象.
整個人類都是神的形象.
男女平等的.
有尊嚴的.
是抗衡當時的文化.
第三個主題.
就是處理苦難的起源.
生老病死.
特別為什麼會有死亡.
這些自古以來.
我們人都是在掙扎這個問題.
就算我們現在二十一世紀.
我們有高科技有文明.
是不是更加幸福呢.
可能某程度上是.
但是本質上是不是呢.

$^{121}$本質上是不是現代的人.
更加有智慧.
生活得更加有尊嚴.
我們很清楚.
是不是呢.
現在面對疫情.
社會政治很多的動盪.
無論東方的世界還是西方的世界.
你說是不是生活更加幸福呢.
一年多的疫情.
我們看到天災.
甚至這個疫情.
大部分都是因為人禍.
疫情裡面.
全球的人就被迫面對一個事實.
就是死亡.
苦罪的問題.
死亡的問題.
近東的文化.
他們是怎麼去回應的呢.
在古代近東的世界裡面.
有一個神話故事.
叫做 Gilgamesh Epic.
基加美斯史詩.
當時是非常流行的.
因為我們找到一些.
Cuneiform tablet.
這些詩形文字.
刻在泥板上.
很多地方在中東近東找到.
譬如巴比倫.
亞述利尼尾城.
敘利亞爾瑪艾瑪的古城.
遠到西邊.
土耳其北邊.
赫人帝國的布赫斯奎.
那裡也找到.
Gilgamesh Epic.
那些泥板的碎片.
甚至到以色列地.

$^{161}$米吉多.
那裡也找到一些文字.
是當時古代的時候抄寫的故事.
Gilgamesh Epic 到底是講什麼呢.
就是講在南米索不達米亞.
這個烏魯的地方.
有一個王.
他叫基加米斯.
他多年去欺壓他的人民.
人民叫苦的時候.
神靈就回應了.
就做了一個野人.
叫做安基道.
Ankidu.
這個本身是一個赤身露體的活物.
是野人.
之後呢.
他就變成了文明的人.
住在城市裡面.
穿上了衣服.
不再赤身露體.
代表了人類文明的進步.
於是這個 Ankidu 就跟 Gilgamesh.
不打不相識.
成為了好朋友.
之後這兩個朋友.
就去周圍歷險.
這些古代的故事.
經常有一些歷險故事.
就表明了英雄的偉大.
但是在這個歷險的故事當中.
最後 Ankidu 就不幸死了.
於是死亡就很困擾了 Gilgamesh.
那麼這個王.
基加米斯.
他就決心要去四處去找.
長生不老的藥.
那麼他去到海極之處.
最後他找到一個.
死不去的先人.

$^{201}$這個先人叫做 Unapishtim.
意思就是說.
He who saw life.
就是他經歷了生命的.
因為連大洪水他都沒有死.
於是乎他成了先.
那麼.
這些近東的故事.
經常提到上古之年.
就是神靈.
就無端端.
就降大洪水.
去到淹沒大地.
就是視人命如無物的.
那麼誰知道有一個其中一個神靈就.
向這個 Unapishtim.
就去通風報信.
他就做了一艘船.
然後逃過這個大能.
他就成仙了.
那麼這個 Unapishtim.
就去指示基加米斯.
就說海底有一個仙草.
或者 magic plant.
magic tree.
好像生命樹一樣.
你吃了就會長生不老的.
那麼這個基加米斯.
結果就真的在海底的樹上.
他就找到了這一棵生命樹.
但是呢.
誰知道.
無意之中他就睏了.
睡著的時候就被蛇偷了.
這個 magic plant 這個生命樹.
結果就很傷心.
回到皇宮就唯有認命了.
人必定會死的.
於是人的生存目的是做什麼呢?.
就是盡情享受.

$^{241}$在這裡建造城市.
大興土木.
流芳百世等等.
你聽完這個故事.
你立刻就會想到.
跟聖經某一些主題.
都是接觸過的.
是不是?.
創世記很多時候.
都會跟這個 Gilgamesh epic.
做一些比對.
最明顯的當然就是洪水故事了.
為什麼會有天災呢?.
神無緣無故降災禍的.
是無法解釋的.
聖經的神學呢.
就非常不同.
羅亞洪水的故事.
是因為人犯罪.
暴力淹沒全地.
神就用公義來審判地上的人.
創世記三至十一章.
就是去處理另外兩個人生.
人類自古以來到今天.
都是在掙扎的問題.
就是死亡和永生的問題.
幾乎整個故事的主題.
這個 Gilgamesh epic.
都是環繞著死亡.
死亡的真相.
所以為什麼要去找呢?.
長生不老的.
永生的生命樹.
所以在近東的故事裡面.
就經常看到一些生命樹.
在近東的長華.
不同的雕刻.
你都會找到這些生命樹.
是一種魔法的 magic plant.
讓人能夠得到永生.

$^{281}$人沒有永生.
在這個 Gilgamesh epic 裡面.
就因為人不幸運.
在故事裡面.
很搞笑的.
就是一個很狡猾的蛇.
偷了這個 magic plant.
使得人要認命.
會有死亡.
我們看創世記第三章.
起碼我們會 get(明白)到.
聖經說.
人為什麼沒有永生.
以及為什麼要面對死亡.
三章第一節.
這裡說到人的墮落.
或者人的叛變.
又說神所做的.
唯有蛇比填冶一切的活物更加狡猾.
蛇對女人說.
上帝豈是真說.
不許你們吃圓中所有樹上的果子嗎?.
女人對蛇說.
圓中樹上的果子我們都可以吃.
唯有圓當中那顆樹上的果子.
神曾說.
你們不可吃也不可摸.
免得你們死.
蛇對女人說.
你們不一定死.
因為神知道.
你們吃的日子眼睛就明亮了.
你們便如神能知道善惡.
於是女人見那顆樹上的果子好作食物.
也越人眼目.
且是可喜愛的.
能使人有智慧就摘下果子來吃了.
又給他丈夫.
他丈夫也吃了.
他們以人的眼睛就明亮了.

$^{321}$知道自己赤身老體.
便拿無花果樹的葉子為自己編作裙子.
在近東的神話故事.
邪惡有時是蛇代表的.
是常傳的.
根本這個世界的原則本質.
由開始都是邪惡的.
聖經提出另一個角度.
蛇.
無論如何說都是symbol of evil.
他出現的時候.
聖經說得很清楚.
他只是創造物之一.
他不能夠與神相比.
是被創造主上帝所控制.
所容許的.
當然我們會繼續追問.
為何一切都是美好.
為何會有蛇出來.
為何會有所謂邪惡.
邪惡的源頭.
聖經沒有滿足我們人的好奇心去解答.
因為其實解答了也沒有用.
但提出的角度.
聖經說.
堅持邪惡不會高過神.
不能夠挑戰上帝.
另一方面.
蛇出現.
只不過是帶來試探.
帶來試探的機會.
試探這個故事.
創世記第三章.
表明人.
創造人的時候.
有道德的自由和責任.
人不是機械人.
在神美好的世界.
被破壞.
是因為人的罪.

$^{361}$沒有一個人能夠推諉.
我們在這個故事裡面.
現代人很多時候看多了.
我們以為女人多嘴的問題.
女人是先犯罪.
因為她連累了我們所有人.
掉進罪裡面.
你小心讀經.
三章第六節.
他說摘下果子來吃.
又給他丈夫.
他丈夫也吃了.
原文有兩個字.
就是說這個丈夫跟他在一起.
說得很清楚.
人的責任.
就帶來罪惡的世界.
周座就自然來.
創世記一至第三章.
你看清楚.
亞當夏娃所為.
其實是人類的代表.
作者在寫什麼呢.
在每個人的故事.
Every man's story.
或者Every woman's story.
亞當 亞當 希伯來文.
是人類的意思.
所以聖經譯作那人.
就是The human.
其實是說我們每一個人.
都在這個故事當中.
我們都掉進了這個試探裡面.
並且用了我們的自由.
用了我們的道德上的自由.
來去為抗上帝.
所以創世記一至第三章.
說到人的墮落的時候.
就很強調人的責任.
這些跟近東的故事是完全不同的.

$^{401}$如何說到人的責任呢.
就是用分辨善惡處的果子.
Tree of knowledge of good and evil.
近東的故事.
剛才看了Gilgamesh epic.
經常提到生命樹.
或者Magic plant.
給人有永生.
聖經就相反.
沒有提過生命樹.
提過反而強調的.
就是分辨善惡處.
創世記的故事裡面.
有很多象徵的記號.
我們不要單單直解.
所謂禁果不是一個蘋果.
蘋果是怎麼來的呢.
禁果是因為後來拉丁文.
蘋果這個字.
還有邪惡這個字.
Malum和Malus.
音相似.
所以西方的基督教.
就有一個這樣的傳說.
禁果是一個蘋果.
肯定不是說.
吃了之後才能夠分辨好壞.
分辨到是非.
神不想人分辨是非.
所以不能吃分辨善惡樹的果子.
肯定不是這樣.
因為創造人的時候.
神的形象.
有自由的.
有道德的抉擇.
是分辨到善惡的.
創世記二章已經有提到.
我們就問.
到底原文.
希伯來文.

$^{441}$這個knowledge of good and evil.
或者determining good and evil.
分辨善惡是什麼意思.
我們只是舉一節經文.
三號耳機下十四章十七節.
這裡形容大衛君王.
我們上文一下.
先不理會.
這裡形容君王是什麼呢.
他說我主我王能夠.
辨別是非.
分辨善惡如同神的使者.
所以是一些皇帝的用語.
形容一些君王.
能夠去斷定是非.
他就是那個尺度.
他可以斷定是非.
你看看創世記一至第三章.
由第一章開始.
誰去分辨萬物.
神稱光為宙.
稱暗為野.
神才是君王.
在創世記第三章.
講到試探的故事的時候.
就是人想去分辨善惡.
即是說人想推翻神的主權.
想自我作主.
所以分辨善惡.
不只是吃過水果那麼簡單.
其實是一種叛變.
有位色經家說得很對.
他說這個試探的故事.
就是The Guardian of God.
即是這個創造主.
我不會認他為神.
是我自己來斷定善惡.
斷定是非.
這是人類的叛變.
叛變之後的結果就是死亡.

$^{481}$死亡不只是醫學的經驗.
不只是肉身的死亡.
死亡在聖經裡面.
是一種屬靈的崩潰.
整個世界本身創造的時候.
你看到頭七天的規模.
是整整有條的.
是一個很和諧的世界.
神,人,大地都是和平的當中.
但是人決定叛變了.
很大的事情.
將整個宇宙,整個創造破壞.
第三章九至十九節.
這裡就交代人叛變之後.
神以公義來審判.
我們不詳細說.
但這裡的結構很清楚.
有五個段落.
第一個對話.
耶和華審問那人.
三章第十二節.
審問女人.
三章十四至十五節.
審問舍.
大家知道這不是對話.
是神耶和華宣判他的失敗.
不是一個對話.
三章十六節.
耶和華宣佈因為罪闖入了世界.
人叛變會有很多痛苦.
女人會有痛苦.
三章十七至十九節第五段.
耶和華對那人說.
大地會受咒詛.
神,人和大地三方面.
全部破裂.
我們很多時候想起罪.
我們都會想起個人做的一些動作.
一些個人行為.
當然這是確實的.

$^{521}$但聖經裡形容罪是更加大的.
形容Sin.
意思是罪惡的狀態.
我們每一個都是有罪的人.
不單止是我們一些個人的犯罪行為.
做錯的行為,不義的行為.
聖經甚至說我們活在罪惡的狀態下.
咒詛.
所以為什麼人需要神的救贖.
我們看幾節經文.
創世第三章第七節.
人叛變之後.
他們二人的眼睛就明亮了.
這裡一定有些幽默.
因為蛇說眼睛會明亮.
果然是明亮.
但結果是不好的.
他們才知道自己是赤身露體.
這個字又出現了.
在這裡是負面的意思.
是羞恥.
犯罪羞恥.
相反第二章尾出現赤身露體的時候.
正面的意思是坦誠.
沒有阻隔.
人的和諧.
關係上的和好.
但叛變之後.
人赤身露體.
知道自己羞恥.
然後人之間就有隔膜.
他們就拿起無花果樹的葉子.
為自己編作群字.
這是人為自己安排的救贖.
創世第三章第二十一節.
我們看看神為叛變的人做了什麼.
神為亞當和他妻子用被子做衣服及藤椿.
人用自己的方法遮蓋.
是無花果的樹葉.
能撐多久.

$^{561}$又是取笑的方式.
神的遮蓋就是被子.
無論如何.
都是暗示有一個生命的犧牲.
神為罪人都能夠預備一個更好的遮蓋.
當然我們知道這是預表.
將來耶穌基督為我們所做.
三章第二十三節.
神變打發他那人.
人的代表.
出去伊甸園.
耕作他自己自出自土.
於是把他趕出去了.
又在伊甸園的東邊安設基路柏.
和四面轉動發火焰的劍.
要把守生命樹的道路.
我們現代人的理解.
就好像他犯了罪.
神就將罪人踢出自己的家.
然後就關門了.
希伯來文是很妙的.
你知道以前的希伯來文.
寫的時候本身是沒有音標的.
沒有響音.
所以如果音標同樣是用另一種方法.
去加響音的話.
這一節可以這樣翻譯.
耶和華在伊甸園的東邊居住.
你看到Var-yes-ko.
不是之前的讀音.
但一樣解得通.
意思是神在伊甸園的東邊.
在伊甸園的外邊.
他捨棄了他自己本身的聖所.
去和人同坐.
神願意.
你看到人永遠和他隔絕.
死亡的狀態.
所以為什麼要禁止他.
罪人去拿生命樹呢.

$^{601}$這是一個因點.
因為在罪的狀態裡面.
你再拿生命樹.
繼續永活下去的話.
你就會在一個罪的狀態裡面永活下去.
神居然離開自己的聖所.
住在伊甸的東邊.
和罪人伴翼的人同坐.
伊甸的故事最後一節.
不是以神的公義來結束.
是神的因點.
第二個近東的文學和聖經有接觸點的地方.
就是死的事實.
人既然會死.
那麼我們餘下的生命.
是為了什麼呢?.
基加米斯的史詩最後一段.
他說當神靈創造人類的時候.
已經定了人會死.
神靈自己才會有永生.
所以基加米斯就勸他喝飽食醉.
建立城市.
米索不達米亞.
因為是人類文明的發源地.
其中一個最重要的文明標誌.
就是城市.
Civilization.
就是由赤身露體的野人.
變成一個穿上衣服的文明的人.
這些有點像古代的人民主義.
Humanism.
你看《創世記》四章開始.
其實就是討論這方面.
四章開始.
有兩個家譜.
我們不可以詳細講解.
但可以列出一個家譜.
第一個家譜.
就是亞當接待了該隱.
之後傳下來的家譜.

$^{641}$這個就是越來越和上帝.
遠離的一個家族.
一種人類.
你記得該隱和阿伯的故事.
該隱殺了自己的兄弟阿伯.
結果該隱離開耶和華的面.
建造自己的城市.
第七代是拉密.
追溯這個故事的時候.
你會看到越來越多暴力.
整個家族不單沒有上帝的同族.
跟著這條線.
那些暴力遮蓋大地.
最後帶來洪水的滲透.
第二個家譜.
就是亞當再生瑟特.
阿伯死了.
然後四章二十五節.
亞當與妻子同房.
他就生了一個兒子叫做Zed.
意思是上帝給我立一個兒子代替阿伯.
因為該人殺了他.
二十六節.
瑟特也生了另一個兒子叫做爾諾斯.
這個時候人去求告耶和華的名義.
爾諾斯再生蓋南.
然後繼續生下去.
第七代就是爾諾.
五章二十一節.
爾諾活到六十五歲.
就生了瑪圖撒拉.
爾諾生了瑪圖撒拉之後.
就與神同行三百年.
並且生兒養女.
爾諾共活了三百六十五歲.
爾諾與神同行.
神將他娶去.
他就不再世上.
然後繼續生下去到第十代.
就是羅亞.

$^{681}$我們看這些家譜的時候.
我們很多時候有現代人的curiosity.
我們覺得勝經說到以前的古人很長命.
很長壽.
幾百年那麼長壽.
其實這些數字就是象徵性的數字.
不是實數.
但我們最重要知道的是什麼呢.
就是近東的人.
聖經時代的人.
要寫家譜來做命.
這裡看到對比兩個家譜.
特別是第七代.
兩個家譜第七代都講多一點.
該隱的家譜.
是暴力以至洪水審判.
因為人所建造的.
是為自己的命.
失德之後的家譜.
第七代就是Eno.
就是與神同行.
代數是三百十五.
即是全心全意.
整個生命.
他反而經歷到所謂永生.
神將他娶去.
家譜是做什麼的.
家譜是身份的定義.
定義你的identity.
到底你是屬於創造主.
你是屬於他的名下.
你是不是宣揚他的名.
失德的後裔.
開始呼求要話的名.
呼求的意思是宣告要話的名.
去公開宣告.
去高舉上帝的名聲.
即是說這班人.
得到上帝的救贖之後.
他們堅持以神的名聲為第一.

$^{721}$有人說在洪水之前.
上古的信心偉人有三個.
第一個是阿伯.
阿伯獻上最好的祭.
相信神.
去敬拜神.
最後當然被那個人所殺.
生命很短.
為主殉道.
但他卻獻上最好的祭.
為上帝所遇難.
第二個.
剛才我們看的以羅.
與神同行.
你見到整個家譜裡面.
都是強調一件事.
死亡.
不斷有死亡.
就像我們今天見到.
外面很多死亡的象徵.
無論是COVID-19也好.
人的脆弱.
但是他與神同行.
猶如在伊甸.
與上帝同行一樣.
仍然經歷到永世.
經歷到生命.
第三個貴人.
就是羅亞.
二人.
信上帝.
他的結果是什麼呢.
雖然洪水審判了他的世界.
但他免於這個審判.
恩典裡得救.
亞伯信神.
他自己信到死.
以洛信神.
他不用死.
羅亞信神.

$^{761}$其他人死.
無論結果如何.
屬神的人.
定義是什麼呢.
就是求告耶和華的命.
宣告耶和華的命.
如果繼續追創世記的話.
你去到第十二章.
亞伯拉罕去見濟澹.
求告耶和華的命.
這個家族.
這一群人.
不斷地去求上帝的名聲.
能夠被彰顯.
不是我們人的名聲.
不是我們人的文明.
而是上帝的榮耀.
我們不單止可以反思.
社會政治上和我們的關係.
其實我們個人也可以反思.
求告耶和華的命是什麼意思.
你為什麼祈禱.
你不是單單為工作順利.
還是你反過來.
你去求的.
你的願景.
就是在工作上.
你可以怎樣將神的名聲高舉.
讓人看到你的工作.
看到你背後的神.
家庭 兒女其實都是一樣.
我們求的是什麼.
我們當然.
人之常情.
當然想他們每個人都乖.
有好的成績.
好的品行.
但求告耶和華.
宣告耶和華是什麼意思.
不是單單希望他們做個乖兒女.

$^{801}$而是希望我的家庭.
我的兒女.
在別人看到他們的時候.
神的名聲被宣告.
你為什麼去祈禱.
你的人生方向是怎樣.
聖經給我們邀請和挑戰.
求主幫助我們.
我們一起低頭祈禱.
神多謝你.
因為你的救贖.
在這個罪惡的世界裡.
有死亡.
我們看到人脆弱.
看到我們自己.
在生命當中其實很短暫.
求你繼續幫助我們.
我們所宣揚的是你的名.
幫助我們的家庭.
我們的關係.
我們的學校.
我們的工作.
每一個場所.
社會 政治.
宣揚耶和華的名聲.
求告你的名.
特別在一個政治不安的地方.
動盪的時代.
主求你讓我們.
繼續堅持高舉神的名聲.
高舉神的名.
你的名得永遠.
我們這樣祈禱.
\newpage



\section{}
\label{sec:ke055dTE6Oc}
\textbf{經緯線 - 耶穌的比喻 (EP2 麥田稗子的比喻)}
\newline
\newline
連結: \href{https://youtube.com/watch?v=ke055dTE6Oc}{\texttt{https://youtube.com/watch?v=ke055dTE6Oc}} ~~~~ 語音日期: 2021-01-29
\newline
\newline
\hyperref[sec:xKw__Td71oI]{\small{< < < PREV SERMON < < <}}
~
\hyperref[sec:index_chronic]{\small{[返順時目]}}
~
\hyperref[sec:index_scriptual]{\small{[返順卷目]}}
~
\hyperref[sec:1qsacx0Vy78]{\small{> > > NEXT SERMON > > >}}
\newline
\newline
$^{1}$Hello 大家好,歡迎回來《經位線》.
上次我們看了一個大家都很熟悉的比喻.
就是殺種的比喻.
看到原來提起上下文和舊約的經文.
可以更加有深度和貼切地理解經文.
我自己學了原來那個比喻是說.
我們接受天國應有的心態.
和接觸耶穌的責任.
這個星期我們會看接下來.
《馬太福音》13章24至30節.
拜祭的比喻.
這個比喻就是說.
耶穌就是跟人們說.
原來天國就好像有一個人.
把種子殺在田裡.
人們睡覺的時候.
仇敵就殺了不好的種子.
就是拜祭在田裡.
到他們發苗的時候.
那些拜祭都和麥子一起生出來.
那些僕人就問主人.
我們要不要去拔走那些幣子出來.
就像除雜草一樣.
但是主人就說不用了.
因為如果你拔那些幣子的時候.
可能那些麥子都會拔出來了.
所以就應該等到收割的時候.
才去除那些幣子出來.
然後把火燒掉.
然後麥子就留在倉庫裡面.
原來這個比喻還沒有完結.
就是等到後面.
耶穌說了其他比喻之後.
36到43節.
耶穌才再解釋這個比喻是為什麼.
我自己看第一眼.
我就馬上去Google.
找究竟拜祭和麥子有什麼分別.
因為我都沒有去耕田.
原來拜祭和麥子是很相似的.

$^{41}$但是要等到他們成熟一點的時候.
才分得出他們哪個是麥子.
哪個是拜祭.
問題就是他們的根會纏在一起.
就是拔走那個拜祭也要拔走麥子.
所以耶穌就這樣說.
媽媽我知道你去過以色列.
我想知道多一點究竟拜祭和麥子的分別.
或者你有沒有去耕田.
去體驗一下究竟是怎樣的回事.
我想我去過以色列的都應該有去過.
有些叫做Biblical Garden.
我去過一個最大的那個叫Neoketimi.
在以色列大概平原一點的地方.
他們的猶太人就將整個以色列.
裡面有的花草樹木都種在那裡.
我記得有一次2,3月的時候.
就是麥子吐水的時候.
就是春天的時候.
還沒有到 Harvest 收割的時候.
我們就在麥田旁邊.
他就撿一些麥水.
還沒有熟的.
就分給我們.
我們城市人就覺得.
哇 這些麥子很有趣.
又告訴我麥子拜祭這個比喻是怎樣的.
你看到照片就知道.
麥子是這樣的樣子.
然後拜祭是很相似的.
但是吐水的方法是有些不同的.
因為野麥拜祭就是.
你可以叫它醜陋一點.
所以農夫就一分就分出來.
還有很多study.
看回這個比喻.
都是很貼切當時農耕的社會.
現在有很多地方還是有傳統耕田的方法.
有些study說.
有些收割出來的麥.

$^{81}$其實有10\%都是壞的麥.
毒麥 拜祭這種.
所以我想.
不單止當時那些人在struggle(抗爭)這些東西.
有很多傳統社會還在struggle(抗爭)這些東西.
耶穌就地取材.
祂就說出一些.
當時農夫一聽就聽得明白了.
這樣就把這個比喻帶出天國.
我自己一看的時候.
就發現這次和上一次不同.
就是門徒沒有馬上上前來問耶穌.
究竟這個比喻是為什麼.
但是耶穌就是再說多一點.
門徒才來問耶穌究竟是為什麼.
是因為耶穌太忙.
祂沒有時間.
被門徒馬上問祂.
還是為什麼呢?.
我想我們讀經的時候.
很少留意當時古代的人.
文學的結構.
我們是linear(直循環)的.
我們說完A就要說B.
說完introduction就要說內容.
然後要說conclusion(結論).
很多時候大家讀過聖經.
或者查過聖經的.
都知道有些叫做giastic structure.
就是那些交叉的並聯.
就好像這裡展示出來.
整個《馬太福音》十三章.
是很對程的.
就是開頭和結尾是首尾呼應.
然後就進入第二部.
第二部就和下面的.
尾二那部分是對程的.
第三段.
今天我們這個默田拜子的比喻.
就對稱另外一個比喻.

$^{121}$就是十三章47.
就是那個殺亡的比喻等等.
然後最中間的那個就是那個戲肉.
比喻的背景.
和天國來臨的實現.
fulfillment(滿足)等等.
如果你有這樣的chart.
大家一起查經就知道.
為什麼會這樣.
好像很散亂.
其實一點都不散亂.
其實很緊接的.
真的不知道.
如果不說.
我想問一下.
其實墨子和拜子.
是代表什麼呢.
如果細心去看.
看到耶穌都有答案.
這個其實都是一個比較罕有.
因為在《馬太福音》只有三個比喻.
是耶穌.
起碼在記載上.
他是將那個解釋來去指明的.
所以我們這次看到這個比喻的時候.
沒理由不看耶穌所說的.
耶穌說什麼呢.
我們來看看.
在第36節開始.
他那裡就門去問他.
然後耶穌的回答就說.
殺好種就是人子.
田地就是世界.
好種就是天國之子.
拜子就是惡者之子.
殺拜子的仇敵就是魔鬼.
收割的時候就是世界的末了.
收割的人就是天使.
其實他將重要的元素解釋了出來.
當然聽完的時候.

$^{161}$都有很多引申出來.
譬如你看看一些一般的識經書.
你會發覺有不同的理解.
包括殺落去.
或者那個場景.
可以是整個世界.
正如耶穌所說.
有人會說在教會裡面.
有好人有壞人.
甚至乎在你自己裡面.
都有好人的部分和不好人的部分.
我比較傾向一個叫Snogras的人.
他所理解.
就是說那個.
所理解的其實是說世界.
其實是說什麼呢.
那個比喻是說.
就正正在這麼充滿邪惡的世界裡面.
竟然上帝的國度是臨在的.
是Presence的.
這個是整個比喻一個核心的信息.
而背後引申出來的.
其實你看到耶穌的解釋後半.
更加著重了其實是審判.
就算是這樣.
其實剛才Mark Sir也有提到.
其實你看回上下文.
就知道他環繞他的比喻上半部分.
就是說將有好有不好.
混在一起.
也都在一個不起眼的情況下.
天國就是這樣展現出來.
到後半在這個解釋之後.
其實是說要分.
有一天要去審判.
要將它再分開.
所以這個其實是.
那個耶穌去解釋背後兩個很重要的東西.
Sarwan你剛才這樣說的時候.
我突然間想起一件事.

$^{201}$我覺得很有意思.
我們看這個比喻的時候.
就覺得是一個好的田地.
竟然有不好的敗子在當中.
但是你剛才這樣說.
又給我們另一個角度去看.
就是雖然有不好的敗子.
但是竟然又有好的物質在當中.
這個兩個不同的角度是可以的.
這個是一個很好的點.
因為我想很多時候.
即使我自己有時看這些比喻.
很容易被那個次序.
如果看回他原先那個比喻.
第幾個節?24節.
是說天國好像一個殺好種的人.
在他的田地裡面.
之後才有一個仇敵.
但是有人都解釋其實那個時序.
有時在比喻那裡我們不可能太斟酌.
他講了另一個例子.
耶穌在末世的筵席裡面.
邀請了一班好人.
然後才邀請了一班什麼人都來.
是不是一定要按這個時序呢?.
其實不是的.
比喻就是我們要抓住那個重點.
有些重點要抓緊.
有些重點其實不要太過轉牛角尖.
另外一個讀這個比喻.
或者我們要去查經,教學的時候.
有一樣東西我們需要留意.
就是那個時間性.
在比喻裡面.
時間是佔了一個很重要的部分.
耶穌解釋比喻的時候.
在第39節就講得很清楚.
即是收割的時候.
是這個時代的終結.
即是不可以現在來到的.

$^{241}$馬上拔出那些敗子.
要等到收割的時候.
所以那個時間稍後.
也都在第41節的時候.
不止41節.
有幾節.
兩三節來解釋.
41節,42節,43節.
那個收割的時候.
對那班敗子,那班壞人.
他們的下場是怎麼樣的.
我想這裡我自己的看法.
在當中我覺得是這樣的.
正如剛才Mark Sir所講到.
在比喻當中的一連串13章.
有一個鋪排.
在第35節的時候.
耶穌說我要開口用比喻說話.
說出創世以來所隱藏的時刻.
在當中講到一些天國裡面的奧秘.
而這個時間就是其中一個奧秘.
什麼意思呢.
當時有些信徒.
他們很期待神能夠即時去做審判.
在這裡都寫得很清楚.
那些死人半島的事.
新漢語譯本在41節說.
死人半島的事和作惡的人.
第41節.
可能當時的信徒都很想有這個.
想對付,對抗這班人.
或者可能心裡很不明白.
神為什麼讓他們還在那裡.
這樣來搞事.
但是這個比喻就正正告訴.
這班信徒.
給他們另一個角度.
好像剛才Saron所說的.
神不是靜默的.
神不是不理的.

$^{281}$不過在時間上是會有遲的.
是等最後大審判那一天.
神才會去做最終的審判.
另外我自己又覺得.
回想起約翰福音.
耶穌在大祭司的祈禱裡面.
如果大家記得在祈禱裡面.
這樣的時間性再一次出現.
什麼時間性呢.
就是在耶穌大祭司的討問裡面的時候.
耶穌為一班門徒來祈禱.
祂說我不求他們離開這個世界.
大家不知道有沒有印象.
祂不求他們離開這個世界.
祂要在這個世界.
耶穌還說了一句話.
我差遣他們來到這個世界.
所以如果用這個比喻裡面的時候.
一班信徒的角色就是被差遣.
來到這個世界.
他們做了什麼呢.
是不是被人欺負呢.
我又覺得不是.
經文裡面第43節.
我覺得很有意思.
在新漢語的日本就這樣說.
二人在他們父的國裡.
要像太陽一樣發出光芒.
所以我覺得反映我們一班的教徒.
是一個好像很無奈.
充滿罪惡.
甚至被人欺騙.
欺壓的世界裡面.
信徒是不是很無助呢.
絕對不是.
信徒可以做的就好像太陽一樣.
來照耀我們身邊的人.
照耀這個世界.
其實你剛才說這個.
對上帝是不是在做事.

$^{321}$剛才和兩個女兒在車裡聊天.
上帝是不是在睡覺.
有時女兒都會這樣覺得.
我們都這樣覺得.
不是當代的信徒覺得.
是我們今時今日.
你看看今年.
都是一個黑暗的時候.
無論在北美.
因應疫情.
到香港這一年發生的事情.
其實很多.
好像剛才Justin這樣說.
很多的壓迫.
很多的無奈.
在這個時候.
我們都會問.
其實上帝是不是在做事.
你睡了覺.
或者上帝的國度是不是在這裡.
其實這個比喻就是答這個問題.
就算在這個這麼黑暗的時候.
原來上帝的國度都是存在.
不過同一時間.
未曾完全看見.
就是夏文那兩個比喻說的那件事.
不是去問.
我們在這樣的時候.
我們怎樣.
應該做什麼.
來去剔除.
根治所有邪惡.
在這裡再一次提醒.
其實最後的審判在上帝那裡.
但是.
剛才Justin所說的.
亦都很重要提醒.
又不是說我們完全.
Totally passive.
什麼都不做.

$^{361}$那無所謂了.
就是有人會說.
將它apply在教會.
教會裡面有好人有壞人.
不要緊.
任由它吧.
耶穌都是這樣.
這個又不是重點.
重點是.
當你看見上帝的國度.
耶穌說祂是存在的.
而最後審判一定會臨到的時候.
我們會選擇什麼.
其實和我們上次所提到的很類似.
你今天怎樣選擇.
所以為什麼耶穌最後.
除了剛才Justin所說.
二人要發出光來.
那個很大的祝福之外.
還有一句.
有耳可聽的就應當聽.
其實上次出現過.
其實正正就是說.
你怎樣選擇呢.
充滿邪惡發生.
逼不臨到你的時候.
你怎樣選擇呢.
我的input就是.
我們這次第二集.
我們在說比喻.
大家都有提過.
kingdom.
就是國度利道.
還有邪惡等等.
我看Anti-Right的解釋都不錯.
他說比喻好像那些舊約的.
就是天啟文學.
天啟文學.
譬如他特別說到.
但以理書.

$^{401}$但以理書有很多天啟文學的元素在裡面.
經常都是說.
有個故事.
有個dream.
有個人發夢.
然後沒有人解得明.
有些奧秘在那裡.
然後有個但以理出來解夢.
那些天啟文學.
經常都是寫給神的子民.
叫他們很黑暗的時候.
很難受的時候.
就告訴他們上帝現掌權.
他就說回.
其實耶穌的比喻.
都有這樣的結構在那裡.
就是很黑暗.
特別我們今天查的這個麥田敗子.
就是黑暗的時候.
撒旦掌權.
好像他做了很多事情的時候.
有很多奧秘.
好與壞在一起.
然後審判又會來到.
天國一定會彰顯.
就叫我們這些信徒.
好像Sharon你這樣說.
有耳開聽的就叮噹聽.
我想我自己都會好像.
好像28節那裡的僕人.
很想問耶穌.
其實我們要不要去拔掉那些敗子.
要不要把不好的立刻除去.
但是發現原來不是現在.
不是我們的工作.
正如她的比喻.
都很強調這件事.
主人回答29節.
不必.
恐怕現在你抽出來的話.

$^{441}$什麼都玩完了.
兩樣都玩完了.
這個時候一定要共存.
暫時共存.
下面都解釋的時候.
耶穌都有說.
世界沒了的時候.
上帝會猜派他的使者.
41節人子會猜顯他的使者.
要挑選出來.
天使做的這件事.
不是我們現在需要做.
我想起那些耕田的.
他們殺種的時候是12月.
然後他們收割的時候是5月.
等了大半年.
很多個月去等候.
我覺得這個比喻都有一個重點.
他經常重複.
就是好種.
一味都是說好種.
我們聽比喻的人是不是有耳可聽.
就當聽.
做回天國的子民.
做回我們的本份.
我想這裡說一個選擇.
其實就好像上一次我們說的這個.
殺種的比喻.
不過那個subtle point有一點點不同.
重點有一點點不同.
上一次說的就不是很多時候.
一般可能聽到的教導.
就說你去選擇你屬於哪種土.
有時候有人說笑我是叫荊棘.
荊棘組的.
整個小組都好像荊棘一樣.
只是說說去遊船河.
吃好東西.
那些荊棘組.
沒得選擇的.

$^{481}$上一次說的重點是.
你會不會應上帝的道呢.
這次其實有一點點不同.
這次都是一個選擇.
這次就是說.
在這個世界裡面.
那惡者 撒旦的工作是存在的.
不需要去推 抵賴.
但上帝的工作同樣存在.
當你在上帝的恩典下.
你接受了這個福音.
你成為天國的兒子 兒女.
天國的一份子的時候.
你的選擇又到了.
你的選擇就是.
你會怎樣去活出.
這個上帝給你的恩典.
你怎樣去活出一個.
不是去做那敗子的一個人.
而是你成為那個.
好像剛才Mark Sir說的那個好種子.
你選擇仍然堅信上帝的工作.
仍然堅信上帝的審判會來到.
所以我今天仍然選擇做一個二人.
其實你今天這個時間.
做一個二人.
第41節就說到.
使人叛道的事和作惡的人.
其實這句話.
甚至耶穌接著解釋那個將來.
是和馬太福音第24章.
是很相似的.
因為那裡如果大家記得的話.
耶穌在末世的時候說了一些先兆.
其中一些先兆就說.
有很多人是被叛奪到的.
互相出賣 彼此憎恨.
有很多假先知出現會迷惑眾人.
不法的事真多.
很多人的愛心都冷淡了.

$^{521}$如果對我自己來說.
就覺得Make sure我自己.
就像剛才西蒙說的.
我要做一個選擇的時候.
我自己就會提醒自己.
這個比喻就是Make sure我自己.
我是一個叛奪到別人的人.
和沒有做一個作惡的人.
而第24章是一個很好的支持這段經文.
因為最後耶穌在第24章13節就說.
堅忍到底的必定得救.
堅忍不是死忍 冷忍 別人欺負你.
而是那種對神的相信.
對神的信仰上的堅忍.
以致神會出手 不是我出手.
可以回應Mark Sir所說的話.
我覺得這樣就可以在這個比喻上.
有一個新的角度.
和一個更深的對神.
更深的相信在裡面.
好啊 希望大家都可以有盼望.
而這個盼望真是來自神對我們的應許.
這一集就到這裡 下次再見.
多謝您收睇時局新聞,再會!.
請不吝點贊訂閱轉發打賞支持明鏡與點點欄目.
\newpage



\section{}
\label{sec:1qsacx0Vy78}
\textbf{舊約網上行 - 上帝的足跡:EP3 文明}
\newline
\newline
連結: \href{https://youtube.com/watch?v=1qsacx0Vy78}{\texttt{https://youtube.com/watch?v=1qsacx0Vy78}} ~~~~ 語音日期: 2021-02-03
\newline
\newline
\hyperref[sec:ke055dTE6Oc]{\small{< < < PREV SERMON < < <}}
~
\hyperref[sec:index_chronic]{\small{[返順時目]}}
~
\hyperref[sec:index_scriptual]{\small{[返順卷目]}}
~
\hyperref[sec:tF7ESjGDOjk]{\small{> > > NEXT SERMON > > >}}
\newline
\newline
$^{1}$這是舊約網上的第三集.
我們開始這個系列的時候.
我們解釋過看創世記的頭幾章.
我們要去到另外一個地方.
就是今天的伊拉克日大.
叫做美索不達米亞.
去看看當地當時一些的文獻.
特別是創世的神話.
近東人如何述說他們遠古的歷史.
世界的起源人類.
起源人的問題等等.
你可能心裡面有一個疑問.
為什麼不可以單單看聖經呢?.
聖經不是已經很清楚了嗎?.
為什麼要參照這些虛幻的神話故事呢?.
參考他們如何看遠古的歷史.
好像降低了聖經的神聖地位.
但是你記得.
每一次我們翻開聖經的時候.
我們其實就好像去到另外一個地方旅行一樣.
不可以就這樣拍照打卡就算了.
我們要去看看當地的人在想什麼.
就將聖經比對當時的一些著作.
甚至一些神話.
如果這樣想的話.
反而會令到聖經更加真實.
因為聖經的啟示.
是在一個真實的時空裡面發生的.
神說不是空中流過.
一定有它的背景.
所以聖經很多時候會用當時古代.
近東人熟悉的概念.
來講述上帝的真理.
表面上好像相似.
但是想清楚.
其實本質上是完全不同的.
因為聖經是針對當時的思潮.
當時的文化.
去提出一些批判.
一些的爭辯.

$^{41}$這正是上帝啟示的焦點.
上兩次我們看完《創世記》一至到第三章.
這次我們看清楚一點.
第四章.
一至到第三章.
人類代表亞當夏娃伴便流亡於伊甸園之外.
這次看第四章.
我們去到聖經的世界.
繼續與神同行.
Zither Harp.
《創世記》第四章.
是一個很熟悉的故事.
該隱與阿伯.
King and Abel.
該隱是農夫.
阿伯是畜牧的.
兩兄弟去敬拜神去獻祭.
該隱的獻祭沒有被上帝閱立.
阿伯的獻祭.
上帝就接受了.
我們先問為什麼.
似乎經文有一些交代.
四章第一節.
有一日那人與他妻子夏娃同房.
夏娃懷孕生了該隱.
她說我得了一個男子.
又生了該隱的兄弟阿伯.
阿伯是牧羊的.
該隱是種地的.
有一日.
該隱拿地裡的出產為貢物獻給耶和華.
阿伯也將羊群裡投生的羊的滋油獻上.
耶和華看中了阿伯和他的貢物.
只是看不中該隱和他的貢物.
該隱大大發牢.
變了殮食.
這裡有一個字.
是貢物.
或者一個更好的翻譯.
是進貢的意思.

$^{81}$貢品.
Tribule.
Tribule是什麼呢?.
在近東世界裡面.
是獻給皇上的.
就是要將最好的東西獻給一個郡王.
原文是Minha.
Tribule是要將最好的獻上.
阿伯就將投生的最好的獻上.
該隱的對比很明顯.
他只是將地裡的出產.
其中一部分獻上.
故事繼續發展下去.
該隱就妒忌,發怒.
四章八至九節.
這裡說他們兩個人在田間的時候.
該隱謀殺了他的兄弟.
然後去審問該隱的時候.
該隱說.
我不知道我的兄弟在哪裡.
我豈是看守我的兄弟的話.
我是不是一個brother keeper(兄弟保護者).
這是聖經記載的第一宗謀殺.
創世記第三章是一個平衡.
第三章人類犯罪.
罪闖入世界.
夫婦之間的關係就破裂了.
人類最基本的關係的單位.
因此而破壞.
現在呢.
第四章.
就是兄弟的關係破壞.
整個community(社會)的破壞越來越嚴重.
該隱甚至比亞當更加心硬.
亞當起碼也肯承認.
但是該隱不僅不承認.
也去抗議神的審判.
神的判決.
故事被我們追溯.
該隱在敬拜裡面.

$^{121}$沒有一張最好的獻上.
神與人之間的關係破裂.
人與人之間.
兄弟之間破裂.
地也受咒俗.
創世記第四章繼續記載.
該隱的歷史.
這一段是比較少讀的經文.
我們花一點時間去看.
四章第十七節.
該隱與妻子同房.
他的妻子懷孕生了二郎.
該隱建造一座城.
按照他的兒子的名.
將那城叫做以洛.
以洛生二娜.
二娜生米胡雅利.
米胡雅利生馬土薩利.
馬土薩利生拉麥.
拉麥娶了兩個妻子.
一個叫阿大.
一個叫洗拉.
阿大生雅伯.
雅伯是鑄匠棚.
牧養畜牧之人的祖師.
雅伯的兄弟名叫尤伯.
他是一妻彈琴吹簫之人的祖師.
洗拉又生了土伯該隱.
他是打造各種銅鐵利器的.
土伯該隱的妹子是拿瑪.
一個加普.
作者就去追溯.
由阿當到該隱之後的後代.
這裡很簡單的列出.
這個加普.
提出他們的成就.
有說城市.
叫做以洛.
另外到了第七代.
是拉麥.

$^{161}$他就成為了第一個.
去到破壞上帝所創造的.
自立的婚姻.
你記得創世的第二章尾.
是有提到的.
另外他們是文明的始祖.
畜牧業.
彈琴吹簫.
與樂.
銅鐵利器等等.
為什麼會提到這些事情呢?.
因為近東的神話.
講到人類起源的歷史的時候.
總會提到.
人類文明的起源.
什麼叫做人類的文明 civilization(文明文明).
在近東的世界裡面.
就是建造城市.
city.
就是人類可以聚居不再分散.
是建立文化.
政體.
社會.
有結構.
第二.
是提到農業.
發明耕種.
就是不再四處採摘食物.
而是他們掌握了農業的科技.
畜牧業.
就是掌握了馴養動物.
不再可以單靠四處打獵.
這幾件事情.
都是人類文明的支柱.
可以定居有城市.
可以成為真正的.
所謂現代的人.
近東文學看到人類文明的建樹的時候.
有一個特點.
就是永遠都會將它.

$^{201}$deify(神聖化).
蘇美爾的《創世神話》.
有一首詩歌.
是去到讚美一件事情的.
讚美什麼呢?.
讚美鋤頭.
hoe(鋤頭).
或者 pickaxe(鋤頭).
為什麼會讚美鋤頭呢?.
因為.
那首詩歌.
去到說鋤頭使萬物興旺.
有鋤的話.
就有好的大麥.
又有泥磚.
人類就能夠好好地生活.
這首詩是說鋤頭的.
是開發好的田地.
然後也讚美.
賜給人類鋤頭的神靈.
現代人看這些古老的詩歌的時候.
我們會覺得很搞笑.
為什麼會讚美一個這麼簡單的鋤頭呢?.
因為.
對於古代的人來說.
這些工具的發明.
掌握到銅鐵.
然後帶來有農業灌溉.
可以收足運河.
去到建立原始的城市.
其實有一點像.
我們今天大讚.
什麼新的科技.
新的電話.
近東人說到文明的時候.
就一定會將它神聖化.
因為神靈是有功能的.
近東裡面神話甚至有一個特別的 title(頭銜).
是叫做 Apkalu.
Apkalu 意思就是說智者的意思.

$^{241}$在美索不達的神話裡面.
有七大個 Apkalu.
因為半神半人.
下半身是魚.
在大洪水之前.
就將人類建立文明的智慧.
這些神靈去教導人類.
譬如寫作,數學.
建築城市,寺廟,制定法律.
怎麼去做耕種,收割等等.
然後有一種說法.
近東的神話說.
因為這些神靈 Apkalu.
就去到建立了七大個以前最重要的古城.
甚至在一些神話當中.
當時的王.
就是這些 Apkalu.
就是說他們開始人類的文明.
所以近東的一些君王的家譜.
就說這些君王成為了一些文明的英雄.
他們又叫做 Apkalu.
所以你回去《創世記》第四章.
你就明白了.
為什麼會突然間提到人類的文明起源.
畜牧業,各種不同的工業.
但是你要留意.
聖經不同的地方在哪裡.
聖經沒有將這些東西神話化.
反而聖經是去神話化.
去 demythologize 它.
一直以來人類的學者.
都知道遠古的人類.
原本是靠採摘,打獵,衛生.
大概到了一萬年八千年之前.
突然間好像一夜之間.
當然實際上是過了很多年.
但是在大歷史裡面.
有如一夜之間.
就出現了農業革命.
這個我們叫做 Neolithic Revolution.

$^{281}$就是新石器時代的革命.
就是有耕種出現.
有農業,畜牧業等等.
人類在遠古以來.
可能有意無意地實驗.
他就發明了耕種.
怎麼去馴養動物.
人類就可以聚居.
有 surplus(超過).
有剩餘.
就建立了聖真.
建立了社會.
最近有一個研究.
是收集種子的數據.
就是說野生的種子.
經過人手種植的時候.
會越來越大.
那些粟米的種子會比野生的種子大十五倍.
黃豆的種子大七倍.
大麥,小麥都是一樣.
當然種子大不一定收成多.
但通常會是這樣.
因為在新石器時代的時候.
人類實驗發明了農耕.
甚至其他的工業.
畜牧,銅鐵等等.
近東的人就會將這些東西神話化.
如果你問創世記的話.
怎麼回應呢?.
聖經創世記第三天的創作.
地上的創造.
神話要有青草.
以及結種子的材塑.
結果子的樹木.
各從其類.
果子是包著核的.
事就這樣成了.
於是地就發生青草.
以及有種子的材塑.
各從其類.

$^{321}$結果子的樹木.
各從其類.
果子包著核.
神看著是好的.
有晚上有早晨.
第三天.
人在創造主的底下.
知道大自然其實有一個原則.
就是他會不斷地去 multiply(多層次).
不斷地去生長,生產.
人就要做國家.
好好管理榮耀神.
但最進入世界之後.
人類的文明有什麼不同了?.
你看這個故事.
是關於該人的後代文明科技的發展.
沒錯人類是管理大地.
可以有文明的演進是好事.
但故事告訴我們.
卻不能夠救人類脫離罪惡和死亡的咒詛.
如果你這樣想.
近東的古人會有一種錯誤.
就是將這些科技文明神聖化.
其實你想想.
現代人.
我們都有相似的地方.
我們總覺得這些科技.
好像給了我們一些救贖.
一些 redemption(復仇).
聖經說.
頂多只是普遍的恩典.
common grace(普遍恩典).
可以減低我們一些.
咒詛帶來的痛苦.
疾病帶來的痛苦.
但最終不是救贖的恩典.
科技多發達.
給我們方便.
但都不能夠除去罪惡.
生活多方便.

$^{361}$都不能夠根治我們人的內心.
相反.
科技文明進步.
可能會增加地上的痛苦.
你看看蓋因後代的故事發展就知道.
蓋因第二代.
謀殺.
人論破壞.
四章十六節.
於是蓋因離開耶和華的民.
住在伊甸東面羅德之地.
Land of Nod.
羅德.
這是希伯來文的字.
其實意思是流亡.
wandering(流亡).
就是說離開了上帝.
就沒有家.
沒有根.
伊甸的東面.
再往東遷移.
創世記裡面.
往東走通常都不是好事.
是遠離伊甸遠離神的面.
這不是 physical geography(物理地理).
這是屬靈的 geography(物理地理).
spiritual geography(靈性地理).
然後這個故事繼續發展下去.
第七代是拉麥.
拉麥和他兩個的妻子.
去到搶手的詩歌.
阿大.
洗拉.
聽我的聲音.
他說.
有一個年輕人壯年人傷我.
我就將他殺死.
虐殺該人是遭報七倍.
殺拉麥.
就遭報七十七倍.

$^{401}$這首詩.
是對比另一首詩.
由創世記開始到現在.
有兩首詩出現過.
第一首詩就是在第二章.
南唱給女兒聽.
你是我骨中的骨.
肉中的肉.
是去到歌頌家庭的美好.
上帝創造的原意和理想.
但是這首詩拉麥.
就是歌頌自己的強暴.
這是戰士的口吻.
所以有人說這首歌叫做 Song of the Sword.
是刀劍之歌.
原本的音樂是歌頌上帝的美好.
歌頌神的.
但是現在的詩詞.
是歌頌暴力.
歌頌暴政.
打銅打鐵.
就變成殺人的武器.
繼續看故事的發展.
你就知道了.
到了第六章的時候.
神的兒子.
罪惡滔天.
然後就招致上帝降洪水的審判.
一個墮落的世界.
文明.
不僅不能夠讓人類救贖.
更加會使邪惡欺壓的事情更加泛濫.
由戈癮到拉麥.
這是繼續走下坡.
罪惡到破壞世界的一個故事.
所以後期到了兩者中間的時候.
猶太人甚至是戈癮這個家族.
為撒旦的化身.
黑暗佔上風.
欺壓的事情繼續增多.

$^{441}$罪惡滔天.
怎麼辦?.
最近在 Christianity Today 去講一個新聞.
就是一個牧者.
叫做 Francis Chan.
他在去年去了香港宣教.
誰知不足一年.
他的簽證不被批准.
所以就叫他回美國.
無論你對這個宣教事件的看法是怎樣.
很明顯香港這個城市已經不會放穴.
不僅是宣教士.
基督教的人士受影響.
差不多同一個時間.
Francis Chan 被拒絕簽證回來的時候.
亦有很多社運民主運動的人士被捕.
有很多雖然已經被保釋.
但亦有很多被還押在監獄.
這些突變我們問.
到底上帝在哪裡?.
公理在哪裡?.
不要說其他地方.
甚至我們本地這個地方.
芝加哥這裡.
我幾日前和本地的學生的士工開會.
亦有其他教會在這裡.
我們都同一個感受.
在疫情的蔓延下.
這些士工的工作受到很大的挫折.
好像黑暗是佔上風.
怎樣辦?.
創世記的家譜.
第四章到第五章.
會否和我們說一些話.
我們讀家譜的時候.
很多時候我們都覺得這些是古老十八代的事.
聖經的家譜為何要說這些.
有這麼多名字.
首先你要知道家譜在聖經裡面.
不是叫我們去計算遠古歷史的年代.

$^{481}$家譜可以這樣形容.
其實是亂世裡面寫的神學宣告.
近東的神話裡面告訴我們.
這個宇宙是不會有公理的.
本身就是邪惡.
本身就是有暴力.
誰的拳頭大誰就說了算.
這些全部是近東的神話.
都是說這件事.
創世記說甚麼呢?.
創世記的家譜是要帶出另一個歷史觀.
創造主仍然是統領著人類歷史.
在創世記的起源史裡面.
有一個特別的用意.
就是去說出這些故事.
去見證上帝在人類歷史裡面曾經有過作為.
並且仍然繼續工作.
上帝仍然繼續推展他自己的計劃.
創世記首幾章的脈絡.
我們大概可以這樣分段.
你就知道他想說甚麼.
由天地的起源說到人類的起源.
第三章就說罪惡的開始.
第三章結尾的時候.
上帝恩典究屬的開始.
第四章就是惡者的後代.
但最後說到義者的後代.
失德等等.
然後到第六章就是罪惡滿盈.
但最後六章第八節.
神恩代羅馬是義者的結局.
這個結構是一個綜合性.
一個ascetology.
英文這個字是來自一個叫ascaton.
一個希臘文的字.
即是the end的意思.
人類歷史寫出來的時候.
是有一個end.
一個最後的目的地正在去.
是上帝在統領.

$^{521}$那個終末就是上帝的審判.
上帝的公義一定會彰顯.
兩個家譜做一個對比.
這個更加鮮明.
第四章說到該人之後的後代家譜.
第七代就是強暴欺壓的事.
越來越多.
但創世記的作者就加第五章.
第五章就是失德的家譜.
新的人類被救贖的人類.
和上帝要同行.
第七代就是以羅.
這裡起碼提出兩件事.
第一件就是創世記告訴我們.
邪惡是有限期的.
你看到該人的後代.
直到拉密.
然後到最後神的兒子罪惡滔天.
第六章.
最後上帝公義的審判洪水來到.
第二就是屬神的世代.
必會繼續延續下去.
失德以羅時一直帶到第十代.
就是和牙.
他是得救的一代.
一個屬神的世代不斷延續下去.
第五章我們不花很多時間.
裡面有很多的壽數.
是很大的數目.
這些不是實數.
這些是近東人.
去記載一個人壽數的時候.
很多時候會用一些完美的數字.
特別是一些君王.
去記載他的作王年期的時候.
都是用一些很大的.
甚至成千上萬的數字.
象徵的數字告訴人們.
他們是完美的.
第五章的家譜.

$^{561}$講到屬神的世代.
都是用一些象徵的完美的數目.
活在地上.
他說這一群屬神的人.
他們是信心的英雄.
上帝繼續去統領帶領他們.
繼續延續屬神的子民.
縱使好像黑暗佔上風.
但是上帝繼續帶領著他的公眾.
問題就是這樣.
到底我們屬於哪一類人.
有人說每一個人的墓碑.
都有兩個數字.
第一個是我們出生的年日.
第二個就是我們離開世界的那一年.
每一個人的墓碑都有兩個數字.
但最重要的可能不是這兩個數字.
最重要的是中間那一條.
中間那一環.
我們的生命.
得到上帝的恩典.
得到上帝的救贖的時候.
我們是為什麼去生活.
我們是去敬拜哪一個.
我們是跟隨哪一個.
我們忠於的是哪一個.
到底是上帝自己的王權.
還是地上的文明.
地上的君王.
我們全心去敬拜他.
敬拜我們的主耶穌基督.
因為祂是我們的君主.
我們同心帶領祈禱.
神我們求你幫助我們.
忠心地去敬拜跟隨你.
縱使我們活在一個很混亂的世代.
在一個亂世當中.
很多令我們不安的事.
在我們身處不同的城市.
我們見到很多的改變.

$^{601}$情況是每況愈下.
主要是因為這些事我們很不安.
我們見到人類有罪的人.
在這個世界有很多欺壓的事出現.
主我們求你讓我們繼續信靠你.
讓我們知道你繼續延續你的工作.
你自己的國度.
讓我們知道邪惡亦有限期.
我們繼續去等候你的公義降臨.
同樣是我們的生命.
我們活在世間的時候.
讓我們繼續去稱頌你.
去事奉你.
去跟隨你.
祈禱奉主名.
我們繼續去等候你的公義.
祈禱奉主名.
我們繼續去等候你的公義.
\newpage



\section{}
\label{sec:tF7ESjGDOjk}
\textbf{經緯線 - 耶穌的比喻 (EP3 葡萄園工人的比喻)}
\newline
\newline
連結: \href{https://youtube.com/watch?v=tF7ESjGDOjk}{\texttt{https://youtube.com/watch?v=tF7ESjGDOjk}} ~~~~ 語音日期: 2021-02-12
\newline
\newline
\hyperref[sec:1qsacx0Vy78]{\small{< < < PREV SERMON < < <}}
~
\hyperref[sec:index_chronic]{\small{[返順時目]}}
~
\hyperref[sec:index_scriptual]{\small{[返順卷目]}}
~
\hyperref[sec:3j_OvRyfEuc]{\small{> > > NEXT SERMON > > >}}
\newline
\newline
$^{1}$Hello 大家好.
歡迎來到第三集的經位線.
耶穌的比喻系列.
我們今天會看《馬太福音》第20章1至16節.
耶穌說天國有一個園主.
一早就請工人到葡萄園工作.
他和工人商量好.
一整天的工錢是一個硬幣.
然後派他們到葡萄園工作.
上午九時.
園主出去.
見到街上的人沒有工作.
就跟他們說.
你們也來我的葡萄園工作.
我會給你們合理的工資.
人們就去了.
在正午和下午三時.
園主又出去請工人到葡萄園工作.
到了下午五時.
他又再次出去.
見到有人站在街上.
問他們經常在這裡做什麼.
沒有工作.
人們就回答他.
因為沒有人請我們.
園主就說.
你們也來我的葡萄園工作.
到了傍晚.
園主就和管理工人的人說.
叫工人過來拿錢.
最遲來的人先拿錢.
最早來的人排到最後才拿.
下午五時的人就得到一個硬幣.
最早來的人以為可以拿到更多錢.
誰知只得一個硬幣.
他們就發牢騷.
最遲來的人工作一個小時.
我們曬了一整天.
為什麼還是拿同樣的錢.
園主就和其中一個人說.

$^{41}$朋友,我沒有欠你什麼.
我們不是說好是一個硬幣嗎.
拿著你的錢就走吧.
我給你之後來的人也是一樣的工錢.
是因為我願意.
難道我用我的錢不可以嗎.
因為我慷慨.
你是不是眼紅啊.
最後耶穌就說.
為守的將要殿後,殿後的將要為守.
我有一個問題.
為什麼這麼有趣.
要三個小時.
園主就要出去叫人來工作.
我想我們一直幾集都想告訴大家.
看這些比喻有一個文化背景.
就是很地道的.
我們今天未必會明白.
我們可能一個月出兩次糧或一次糧.
以前的工人每天早上都要出去工作.
晚上就有錢拿回家吃東西.
沒有就很慘了.
一天一天這樣去的.
所以第一節說的.
本地人當時一聽就明白了.
因為每個人都是這樣上班的.
去到市場等人請.
特別是葡萄園或者田地要收割的時候.
例如葡萄園在巴勒斯坦收割大概是9月10月左右.
所以他們就在公子站那裡.
他們就等人請他們.
還有你剛才說的時間.
其實都是當時的人去計算時間的方法.
如果去到古代.
你就知道有一些日規.
很多時候就會分12格.
或者三格三格這樣分.
因為以前說時辰的時候是三小時三小時.
所以新的翻譯本我覺得是挺好的.
就是早上6點.

$^{81}$6am,然後9點,12點,3pm.
最後6點就下班了.
所以剛才整個比喻是12個小時的工作.
我們之前一直都說.
看相聞下你很重要.
為什麼這次我們會從13章一跳就跳到20章.
前面後面究竟發生什麼事.
可以讓我們理解這個比喻會好一點呢?.
我想如果大家一直有跟著我們這個比喻的系列.
第一集我們其實做了馬太福音的一個大圖畫.
在馬太福音裡面有五個很長的耶穌的教導.
大家應該記得了.
其中一個一開始就是登山寶訓.
第五章至第七章.
第二個就是去猜險門圖.
第十章.
第三個大分段就是講天國.
第十三章.
然後就是講天國群體的相處.
第十八章.
最後就是橄欖山.
耶穌的回來.
那個比喻.
第24章到第25章.
這個比喻就結在第四和第五章中間.
沒錯.
我們很快就在十八章到二十章裡面.
都留意到一樣東西.
就是天國一群人的相處.
然後天國的價值是怎麼樣.
我們都看到這些東西.
一直由十八章一直流傳到我們今天的二十章.
我自己也有留意到一點東西.
除了這個博客之外.
太大了我分不清.
但是我看到在第十九章最後的一節.
寫著在前的章要在後.
在後的章要在前.
這個是二十章.
這個比喻最後也有這樣說的.

$^{121}$George你很厲害.
你開始留意一下.
如果我們看正經的.
看前面和後面.
就可以明白多一點.
我們看到這樣包著的時候.
其實這個很明顯是其中一個主題.
如果我們看長一點.
再看十九章.
再看遠一點的時候.
我們就更加明白.
耶穌說我們今天選擇二十章的比喻.
那個意思和重點.
因為如果大家很快看看十九章的時候.
你會發覺十九章有幾樣東西.
因為我們剛才說十八章是第四個分段.
在第四個分段的時候.
開始說天國的價值觀.
十九章裡面我們都看到這件事.
一路延續,延伸著天國的價值觀.
大家很快看看十九章的時候.
你會看到有人來問耶穌.
究竟可不可以離婚.
耶穌說不可以離婚.
就是天國的價值觀.
然後在第十九章,第十三節就開始說.
神是很愛兒童的.
很愛孩子的.
然後我們又看到另外一個.
也是很有名的一個故事.
有人來問耶穌.
那個人很有錢.
但是那個人為了錢.
不肯跟隨耶穌做門徒.
所以變成我們全部都是在看天國的價值.
天國價值就是.
就是小朋友.
神是愛小朋友的.
天國價值就是.
你們盡可能都不應該離婚的.

$^{161}$在天國的價值裡面.
因為看到有錢的人跟隨著是很難的.
因為他們看到地上的價值.
而不是天上的價值.
來到我們這裡的時候.
也是另外一個天國的價值觀.
值得留意一件事.
其實在二十章一開始的第一句.
和合本有這個「因為」這個字.
其實你如果看回.
其實它就代表原文的「加」.
不是加減乘除的「加」.
希臘文的那個「加」.
對,for.
通常譯作.
其實這個正正就是去引出.
其實十九章三十節說什麼.
其實是同一個對話.
耶穌沒有斷過.
起碼在馬太的表達上沒有停的.
是同一個conversation.
然後他似乎給了一個example.
什麼是在後會在前,在前會在後.
去到尾十六節又再重覆.
就好像剛才Justin這樣說.
是一個inclusive去重覆.
原來這個比喻其實和這個proverb.
這一句話其實很有相關.
所以在這個比喻裡面.
我想當我們去讀這個比喻.
我們有一條問題可以問自己.
這條問題也是我問自己的.
就是一邊讀的時候.
因為一邊都在說耶穌在和門徒.
教導天國的價值觀是什麼.
所以我們可以問自己.
幫我們去解釋這個比喻.
就是說究竟在這個比喻裡面.
究竟耶穌想和門徒說的價值觀.
是什麼價值觀呢?.

$^{201}$我們如果從這個角度看.
可能就會發現到一些新的事物.
這句話我們就知道好像很重要.
在前的要在後,在後的要在前.
很圓的,但是說完都不太明白.
究竟是怎樣.
耶穌給了這個例子.
有什麼在這個比喻裡面.
我們可以留意,可以明白一點呢?.
我想整個比喻其實再看.
我想一個很關鍵的,用這個起承轉合.
就是他設了一個局.
就是有原主去找工人來工作.
剛才Mark Sir說的很普遍.
當時就是這樣.
在這樣的情況下.
當然不是每樣東西都是現實的.
為什麼他會出去五次呢?.
是不是應該計劃得好一點呢?.
這些其實是超過現實生活.
但是計劃本身就是這樣.
有些東西你要看現實生活.
有些東西你就要知道它是假裝的.
但是就是這裡他想帶出一個情況.
就是他連最後請那班.
他都說你來吧,你照樣來吧.
你來完之後我不會去虧待你.
接著那個高潮就來了.
當他去到要給錢的時候.
他第一,他叫最後那個被請的來拿錢.
不單是這樣,先後次序不重要.
你都想著一開始那個一元雞.
這個一毫子,五毫子.
總之沒理由一元.
誰知道他就是給一元.
還要給開頭那批的人知道.
那班人就開始去發牢訴.
其實有一個很特別的字.
第十一節他們是他們得了之後就埋怨家主.
如果你熟悉舊約這個埋怨.

$^{241}$其實就是正正是以色列人在曠野埋怨上帝.
就是因為他們埋怨不能進加拿大的處境的反應.
這班人就是明明和家主約定了那麼多錢.
可能他忘記了.
但竟然因為旁邊那個拿著一元雞的時候.
他就不開心,他就去埋怨.
所以這樣看來整個高潮.
不單單是說家主主人有多好.
都不是最重要的一件事.
因為坦白說都是一元雞.
他不是給一千萬他們.
不是不做事給錢你.
而是你做完事給錢你.
不過做得少的,做了一個小時的.
他都照樣給那麼多錢,這是他好人.
但這裡所說的更加重要.
其實是說為什麼這個元主這麼好人的時候.
竟然都有份被請,有份得到這個好處.
竟然去埋怨,去嫉妒,去不開心.
旁邊那個都同樣得到這個好處.
這件事其實是整個高潮要帶出.
也符合剛才Justin所說.
你看看上下文,特別在說之前彼得去問.
他說撇下所有東西跟隨你,將來怎樣.
應該更扭轉.
這個是實現了人的計算.
你說得對,細雲.
眼紅.
人的計算,覺得不值得.
做了那麼多時間.
正如我們剛才所說的.
天國的價值觀原來不是這樣.
可以坦白說一句.
眼紅,妒忌,不值得,不值得.
在天國裡面是沒有份的.
你看看我們剛才看了上,我們看下文.
我們看到在二十章之後.
很好笑,耶穌教導完.
你們應該聰明了,這群門徒.
誰知道在馬太的編排下.

$^{281}$馬上弄一個現實例子給你看.
讓你知道他們這群人真的不明白.
現實例子就是二十章的第二十節.
西比太尼兒子的母親.
兩個兒子上前來.
希望跟耶穌說.
將來你的國來的時候.
我可以坐在你左邊和右邊做左右丞相.
又來那些世俗的眼光.
又來那些世俗的計算.
然後第二十四節又是Selwyn剛才說的.
比喻那些情緒.
十個門徒走漏了他們.
還有剛才Joyce說的那句在前在後.
二十七節又出現.
垂燕為首的,這個就很熟悉了.
人子離,不是要受人服侍.
而是要服侍人.
這個觀察也好,你下文看完.
你到馬太福音的脈絡.
已經是最後一段.
二十一章就已經榮盡耶路撒冷.
其實最後一個星期.
天國真的會到,藉著他們的死和復活.
帶來的.
我蠻喜歡剛才Selwyn說的細節.
十一節說埋怨.
在以色列曠野的人.
埋怨就不能進入迦南.
我覺得這裡上下這樣放下來.
我們說比喻經常有一個.
Eschatological urgency.
一個末世最後的一個hour.
這個比喻說.
最後一個hour發生了什麼事.
現在真的到耶穌.
事奉的最後一個hour.
滿途有跟著他.
天國就到了,到底你能不能進入天國.
你是否擁抱到.

$^{321}$天國的價值觀.
在前的在後的在前.
其實這裡你說得對.
其實就是你認不認識上帝.
認不認識上帝的角度.
這個比喻一開始是說天國.
好像是這樣.
其實也是說一個天國會是怎樣.
這裡特別一個字.
其實很有意思.
就是或者我們作為一個.
跟從耶穌的門徒.
與其去眼紅其他人.
與其去用這個.
公不公平來計較的時候.
或者這個更加值得學.
他第十五節.
因為我作好人.
你就紅了眼.
Evil eyes.
去直譯.
其實兩個都有一個.
你可以refer去.
箴言裡面,或者明多一點.
一個是在箴言的.
二十八章.
第二十二節.
他說人有惡眼.
Evil eyes.
想要急速發財,卻不知窮乏.
必臨到他身.
即是說你那個.
去度.
那個.
眼紅的眼.
就會落入這樣的結局.
另外一個在箴言二十二章.
第九節,他說.
眼目慈善的,就必蒙福.
因他將食物分給窮人.

$^{361}$就是說其實是.
good eye.
誰是有好的眼.
其實他就是會去.
看得見有壞的人.
將他有的去分享.
其實這個正正.
如果你又看回上文.
十九章,少年才主.
去問的時候.
他說耶穌是好的.
耶穌說.
只有一個是好的.
那個是天賦自己.
其實就是.
真正有good eyes.
那個就是上帝.
我們有沒有good eyes,還是有evil eyes.
我們看得見的.
是好像上帝那樣.
將有的.
去分享給其他人.
其實是一個self-centered.
只看見自己所謂掛著一個公平的原則.
其實都是掛著自己的益處.
這樣去看.
進入天國.
就好像傑森說,其實你不能進入天國.
這個挺好.
這個又帶回傑森說的.
larger context,大一點的下文.
下文,未進耶路撒冷之前.
就是醫治兩個核子.
為什麼最後天國來到之前.
要那個核子.
我的眼睛能看見.
看見什麼.
我們經常看到醫生醫好.
那個人可以看東西.
最主要看什麼,就是spiritual eyesight.

$^{401}$就是看到耶穌來是做什麼.
我們跟隨耶穌的應該做什麼.
和他看見之後.
是去跟從耶穌.
exactly,very good.
我又學到一樣新的東西.
原來除了我們見到舊約的經文.
就可以看那個章節.
是什麼意思.
來對比喻有深入的了解.
原來找到一些重點的字眼.
去看那些字在哪裡出現過.
都可以對於理解.
就是段經文.
有一個深入的.
understanding.
所以,總括我們今天說過的東西.
其實這個比喻的重點.
是否在說.
神,原來神和天國的價值.
對於我們來說.
是很不一樣的.
很多時候我們以為.
最頭,最先,最叻那個.
肯定有多些獎賞.
可以站第一.
反而不是.
是否最有錢.
第一個來工作的.
就有最多的獎賞呢?.
又不是.
原來在天國的價值來說.
就是最遲的.
反而可以.
最先的.
反而又要等到最後.
我想.
我自己看這個經文.
最後我自己的反省.
有點像一個.

$^{441}$dynamic,一個motion.
給我自己一個反省.
不是說究竟我是前還是後.
而是我們時時刻刻.
作為基督徒都可以問自己.
究竟這一刻,如果我是後的時候.
其實我可以是前的.
但不是靠.
自己的努力或爭取什麼.
而是看自己的心態.
我是否和神的角度是align的.
我會不會紅眼.
我會不會和別人搶大.
我會不會在事奉上.
我覺得我要有title.
變成我們經常.
可以有這個evaluation.
變成我們在後的時候.
就可以是前.
不是說你在後.
我就定了你永遠在後.
我覺得經文的flow.
是可以這樣.
其實很多時候我們的重點.
就是完全放錯了.
在神的葡萄園裡面.
工作不用叫我.
怎知道就放在.
是否公平.
但其實這個公平.
都不是我們理想的.
好,我們今集就說到這裡.
下次再見.
喜歡的話請訂閱支持啊~.
by bwd6.
\newpage



\section{6:30-34}
\label{sec:3j_OvRyfEuc}
\textbf{Kingdom Values XL Matthew 6\_30-34}
\newline
\newline
連結: \href{https://youtube.com/watch?v=3j-OvRyfEuc}{\texttt{https://youtube.com/watch?v=3j-OvRyfEuc}} ~~~~ 語音日期: 2021-02-15
\newline
\newline
\hyperref[sec:tF7ESjGDOjk]{\small{< < < PREV SERMON < < <}}
~
\hyperref[sec:index_chronic]{\small{[返順時目]}}
~
\hyperref[sec:index_scriptual]{\small{[返順卷目]}}
~
\hyperref[sec:OwmCWGNsDAQ]{\small{> > > NEXT SERMON > > >}}
\newline
\newline
$^{1}$Hey family, as always I am praying for you. I hope that you're warm and safe this week and.
enjoying the winter. If you're someone who likes to get out in the weather and enjoy the snow,.
then I hope you're doing that. If you just like to stay in and stay warm and sit by the fireplace and read or.
get on the couch and and watch a movie that you love, then I hope you're enjoying that. I hope that you're.
feeling well and being safe and warm even during this cold time. I also want to wish you a very happy.
Lunar New Year. I know all of the holidays have been so different this year.
because we're not able to gather together in the way that we normally would..
But I hope that you are enjoying this time and that you're able to connect with family and friends, even if it's only virtually..
Hope that you're able to connect with the people that you love and to share.
with them how much you love them and how grateful you are for their presence in your lives. And also hope that you're able to just.
enjoy celebrating and looking forward to the next year and what it could mean for us..
I hope that however you're celebrating the Spring Festival this year, I hope that it's joyful and.
encouraging and a blessed time for you and for the people you love..
Valentine's Day too, if you're celebrating that, I hope that you're enjoying that, celebrating that with someone that you love as well..
And I want to remind you as always to connect with us. You can send us an email at info@urbanvoicechurch.com..
We'd love to hear from you all..
But today we're going to be looking again in Matthew chapter 6, still looking at these kingdom values, and we're finishing up.
the topic of anxiety. We've been talking about this for a few weeks now. So let's look at Matthew chapter 6 verse 25 here..
We're going to read all the way down to verse 34. It says this,.
"Therefore I tell you do not be anxious about your life, what you will eat or what you will drink, nor about your body,.
what you will put on. Is not life more than food and the body more than clothing?.
Look at the birds of the air..
They neither sow nor reap nor gather into barns and yet your heavenly Father feeds them..
Are you not of more value than they? And which of you by being anxious can add a single hour to the span of.
life? And why are you anxious about clothing? Consider the lilies of the field, how they grow..
They neither toil nor spin, yet I tell you even Solomon in all his glory was not arrayed like one of these..
But if God so clothes the grass of the field, which today is alive and tomorrow is thrown into the oven,.
will he not much more clothe you, O you of little faith?.
Therefore do not be anxious saying, 'What shall we eat?' or 'What shall we drink?' or 'What shall we wear?'.
For the Gentiles seek after these things and your heavenly Father knows that you need them all..
But seek first the kingdom of God and his righteousness and all these things will be added to you..
Therefore do not be anxious about tomorrow, for tomorrow will be anxious for itself..
Sufficient for the day is its own trouble..
As we talked about.
anxiety over the past few weeks, we talked about the fact that.
Jesus is not telling us here to be naive or out of touch with reality..
He understands that our life is going to involve trouble and difficulty and challenges..
We live in a broken and fallen world and.
he understands that we are going to be aware of those things and even concerned about them. And that's okay. That's actually healthy..
We need to be aware and concerned of those challenges of life..

$^{41}$But Jesus doesn't want that awareness and concern to grow into something.
that damages us and that causes us to to make unhealthy decisions..
And that's one of the big things that we see a lot of times when we are.
trapped by worry and fear and anxiety..
I've been there in my life and maybe you've been there or seen someone who was there in their life. When we begin to.
get anxious and worried and fearful, we begin to make unhealthy decisions that we would not.
make otherwise. And those decisions can be damaging to us or to the people around us.
because they're motivated by fear and they're motivated by our worry. And that fear and worry can cause us to do some really.
even uncharacteristic things. It can cause us to do some strange things that just.
are not what God has for us. And that's one of the reasons that Jesus is.
calling us out here and saying, "I don't want you to live in anxiety.
because I don't want you to to be motivated by fear..
I don't want you to make decisions and do unhealthy things in your life that come about because of.
your fear and your worry." And in these verses, Jesus is lovingly acknowledging our need and our weakness..
And he's saying, "Hey, I get it. I get that you get afraid. I get that you worry.".
There's so much beauty and love in the fact that Jesus is talking about this with us..
And he's getting in the middle of our struggle and our mess and he's loving on us and I appreciate that so much..
But he's also showing in these verses that his desire for us is not to get stuck in anxiety and fear and worry..
His desire for us is not to be under bondage to it and not to be.
crushed by it. That's not what he has for us..
It may be a lifelong challenge for us,.
but he wants us to experience freedom and healing as we go through life and even victory in this area of our life..
And.
Jesus is pointing us to the fact that he.
cares for our physical needs and he cares for our emotional and spiritual needs and he cares about the things that are.
beyond our control even. We see this in these verses. Jesus is saying, "Look, I know you have stuff that worries you..
I know you worry about the stuff that you need and I care about that..
I know that you worry about things that happen that are hurtful or harmful or difficult and you have no control over them.
and I care about that." I think there's such a loving heart here.
that's displayed in these verses and we talked about how as we move away from anxiety, we need to change our focus..
We need to let that change in focus bring about a change in perspective and we need to grow in our trust of Jesus..
And what I want us to see this week is that Jesus is really telling us to pursue.
the heavenly kingdom and leave the rest to him. We can see this in verse 30..
It says this, "But if God so clothes the grass of the field which today is alive and tomorrow is thrown into the oven, will he not.
much more clothe you, O you of little faith?".
He's saying, "Look, God is caring for nature and God will much more care for you..
So have faith. Keep your focus in the right place. Keep your eyes on the heavenly kingdom.".
And in verse 32, he says this, "For the Gentiles seek all after all these things and your heavenly father knows that you need them all.".
Jesus is drawing this contrast when he uses that word Gentiles..
He's talking about the people who are outside the kingdom of God, the people who don't believe..

$^{81}$And he's saying, "Look, the people who are outside the kingdom,.
they're the ones who get so worried about all this stuff..
They're the ones who are really focused on their food and their clothing and those are huge issues to them.".
And for us in the modern world, we could add some things to that, right?.
We could say it's the people outside the kingdom who we expect to be focused on materialism and wealth.
and all of those physical things and status and prestige and what kind of house you live in.
and what kind of car you drive and your image on the internet and what kind of meals you're eating.
and do people think you're cool or do people think you're good or do people think you're concerned about the right issue socially.
or that your politics are in the right place..
All of those things that get so wrapped up in who we are as people,.
those are the things that people outside the kingdom of God care about..
Why do they care so much about those things? Because that's all they have..
They don't know the kingdom of God. They don't know the deeper reality of God's truth..
And so all they have is the physical material things that are right in front of them..
All they have is their status or how other people feel about them..
All they have is the opinion of someone else..
All they have is the attention and adoration of someone else and that's it..
But Jesus is saying, "Look, it shouldn't be that way for you because you have a different perspective..
You know the kingdom of heaven..
You should be pursuing the heavenly kingdom and leaving the rest to God.".
In verse 33, Jesus says this very clearly, very plainly, and very profoundly..
He says, "But seek first the kingdom of God and His righteousness, and all these things will be added to you.".
That is one of the most powerful and profound verses in the Sermon on the Mount..
Jesus is saying, "Look, if you will pursue the values of the kingdom,.
if you will pursue rightness with God, and His righteousness.".
That word righteousness means living a life that is right with God,.
living in a way that is pleasing to God, living in a way that is obedient to God..
If you will pursue the values of the kingdom, if you will pursue a life that is right with God,.
if you will seek to obey Him in your life, then God will add in all the rest..
All those things we just talked about that are important to us as humans, God will add them in..
God will take care of our physical needs..
He'll make sure we have enough to eat and enough to wear..
He'll make sure we have a place to live, and He'll make sure we have transportation..
He will even make sure that we have relationships with other people that are healthy..
God will add in all of those things if we will simply pursue His kingdom and the values of His kingdom..
God will take care of the rest, and that's what Jesus is pointing us to in this verse..
And Jesus is also reminding us here to take our hands off of what we can't control..
Verse 34 is another verse that is just so profound and so powerful..
It says this, "Therefore do not be anxious about tomorrow, for tomorrow will be anxious for itself..
Sufficient for the day is its own trouble.".

$^{121}$I love that verse..
Jesus is saying, "Why would you be worried about what's going to happen next week, or next year,.
or five years from now, because you have no control over next week, or next year, or five years from now..
You can't do anything about it..
And by the time you get to next week, and next year, and five years from now,.
you'll have a whole new set of circumstances, and a whole new set of problems probably,.
that you don't even know about right now..
So why try to figure all that out and worry about that now?.
That day is going to have its own trouble..
It's going to have its own issues to deal with..
For today, you should simply trust and walk with God, and let Him be at work in your life,.
and let Him bring you peace, and comfort, and joy..
And you can worry about next week, and next year, and five years from now, when you get there..
Now Jesus is not saying that we shouldn't plan, or we shouldn't prepare for things in our life..
There's other passages of Scripture that talk about that..
That's something that we're supposed to do..
That's a healthy thing for us to plan and prepare..
But even as we plan and prepare, if we are followers of Jesus,.
we should do that with an understanding that anything could change..
We should hold loosely to our plans and our preparations..
We should understand that, "Okay, this is what I might like to do..
This is what I might plan to do..
But this could change..
God could change the course of my life..
God could change where I'm headed..
God could change what's happening in my life.".
Or something else could happen, something external that's beyond my control could happen,.
that God might allow to come into my life, that might be difficult or challenging..
So it may end up being completely different than what I want it to be..
And that's okay..
As believers, we can plan and prepare, but we need to hold very loosely to those things..
Take our hands off of what we can't control,.
and stop being anxious about the future when we have no control over that..
When I was a young man, I had a lot of plans for my future in terms of career and education.
and all these things..
And God changed those as I went along..
When I was an undergrad, I was planning to go on to grad school.
and to pursue a career in science, and God changed that..
And I ended up going to grad school and pursuing a degree in divinity,.
and then going into the ministry..

$^{161}$And at the time when I was going into those things,.
when I was in high school and the first few years of college,.
if you would have described to me my life now,.
I would have said, "That doesn't make any sense..
I don't think that I would ever live that life.".
But now looking back, I can see how good it's been for me,.
and how God has been good every step of the way,.
and how all the things that God has brought into my life,.
the relationships and the work that He's given me,.
and the ministry that He's given me, has all been so good for me..
And I wouldn't want to change anything..
But when I was a young man looking forward, I never would have planned that pathway..
I would have planned a completely different pathway that, to be perfectly honest,.
was really focused on me and was a much more selfish pathway..
And to be perfectly honest, if I had lived out that pathway,.
I would probably be a very selfish and proud and arrogant and self-absorbed person..
I'm really glad that I didn't live out that pathway..
But I didn't know what was best for me..
God did. God changed a lot of stuff..
And so when we think about our future, we should say, "You know what?.
Instead of trying to focus on all these things I can't control,.
I'm going to allow God to be in control..
I may plan and prepare, but I'm going to hold loosely to those plans and preparations,.
and I'm going to let God do what's best in my life,.
because He knows what's best for me..
He sees things that are coming in a way that I do not see..
I'm going to take my hands off and let God drive..
I'm going to let Jesus take the control in my life.".
And that's really important for us, and it's a way that we can avoid a lot of anxiety..
If we'll stop trying to put our hands on it and force it to work.
and make it work out the way we want it to, life's not like that..
We can't control life..
We can't just bring it all down and mold life to our will..
We cannot do that..
And if we'll take our hands off and say, "You know what?.
I'm going to trust the Lord to control and to work, even in the difficult moments.".
That's something that can relieve us of so much anxiety and give us so much peace..
You know what Jesus is saying here in verses 33 and 34.
is so counterintuitive to how we feel about our circumstances..
He's saying that as we release control and let go of our desire for control,.

$^{201}$we actually find peace..
And that seems just backwards and upside down to our minds..
When we are struggling with anxiety, a lot of times we feel like if we had more control,.
we would find peace..
We feel like if only we had more control over that situation,.
only if only we had more control over our future,.
if only we had more control over that relationship,.
then we could stop worrying about it..
And we feel like more control would bring us peace..
And Jesus is saying here that actually what brings us peace is when we release.
and we stop trying to control everything..
We stop trying to be in charge of it and we release and trust God and His love and care in our lives..
And that also has such a powerful effect on our behavior..
Like we talked about at the beginning,.
when we're motivated by fear and worry and anxiety,.
we can make really unhealthy decisions because we're being motivated from an unhealthy place..
And we've experienced this in our lives..
Think about it..
Some of the worst financial decisions that you've made in your life.
probably came during a time of anxiety and fear..
And that anxiety and fear motivated you to do something that was rash or impulsive.
or just wasn't a good idea..
Something that maybe you never would have done otherwise..
And it ended up being unhealthy for you..
The same is true in our relationships..
When we get anxious about a relationship and feel like we need to control that relationship more.
and we're motivated by fear and anxiety in that relationship,.
suddenly we become manipulative and we become difficult to interact with..
And the relationship suffers because of our anxiety and our fear..
So many things in our life are like that..
We convince ourselves that it would be so much better if only we had more control..
The reality is we don't control anything in our life anyway..
And any control that we think we have is just an illusion..
But we still think we should control it..
And we still want to control it..
But the truth is that when we release that control and say,.
"God, I can't do this. I can't make it happen..
I don't want to act out of fear of what might go wrong or what could happen..
I want to trust you and trust your goodness and trust your love and your care for me.".
When we begin to relate to our life and our circumstances and our relationships and our future.

$^{241}$and all the needs that we have from that place, that's when we really do experience peace..
Peace comes through release, through letting go of those things,.
not through grasping for more control..
And in verse 33 here, I love this verse,.
"But seek first the kingdom of God and his righteousness,.
and all these things will be added to you.".
Jesus is calling us to loyalty to kingdom priorities..
You know, we talked about in the verses that precede this section on anxiety,.
where Jesus is talking about our focus..
We talked about how Jesus in those verses was calling us to loyalty to kingdom values..
And Jesus wanted us to be loyal to the values of the kingdom.
and to seek to live those out and to be obedient to God..
And in this section, Jesus is calling us to loyalty to kingdom priorities..
He's saying, "Okay, if these values are what you value in your life,.
if you value the things that are important to God,.
then your priorities, the things that you organize your life around,.
the things that you say, 'This is number one and this is number two and this is number three,'.
will reflect those values of the kingdom.".
And that's what Jesus is saying to us in Matthew 6, verse 33..
"Seek first the kingdom of God and his righteousness, and all these things will be added to you.".
That's a fair question to say, "Well, what is the kingdom of God and his righteousness?.
What are those kingdom priorities?".
Well, what do we know from Scripture?.
Jesus told us that the first commandment is to love God with all of our heart and soul and strength..
And the second is to love our neighbor as ourself..
And that those two commandments, if we would really live them out,.
that that would basically encompass all of the law and all of God's instructions to us in the Old Testament..
So that's a great starting place to say, "Okay, if I want to be loyal to kingdom priorities,.
then the number one thing in my life is God and my relationship to him..
And I want to love him with everything that I have..
I want to love him with my mind. I want to love him with my heart..
I want to love him with my will and my choices and my obedience..
I want to love him with all of the whole of who I am as a person..
I want all of that to holistically love God..
And then out of that love for God, I want to love others..
I want to love my family, my siblings, or my parents, or my spouse if you're married,.
or my children if you have children..
Or I want to love my friends, the relationships that God has brought into my life,.
these people that God has allowed me to have relationship with..
I want to love my co-workers. I want to love the people that I go to school with..

$^{281}$I want to love the people at my church..
I want to love the people in my community, in my neighborhood..
I want to love people who have needs in my community, in my neighborhood, that I can serve and care for..
Though that is a kingdom priority, and to say that is more important,.
loving others is more important than me getting rich, or than me being comfortable,.
or than me stacking up a lot of possessions, or than me having an image of being some amazing person.
through my social media accounts, or than me feeling good about myself.
because I'm trying to do things to impress other people or to perform to some standard..
It's more important for me to love others..
It's more important for me, as Jesus said, to lay down my life for someone else..
Remember, Jesus said that greater love has no one than this,.
than to lay down your life for your friends..
It's more important for me to sacrifice for someone else than to just seek my own good,.
and my own wants, and my own desires..
That's what it looks like to have loyalty to kingdom priorities..
And that's what Jesus is calling us to in this verse..
He's saying, "Look, if you will seek the kingdom of God first,.
if your loyalty will be towards the things that are important to God,.
and you will order your life in that way,.
if you will say, 'First and foremost, it's about loving God and my relationship to Him,.
and secondly, I'm going to love others, and other things are going to fall in behind that,'".
Jesus is saying, "Then God will add in everything else you will need,.
and you can finally stop being anxious,.
because you won't have to worry about the future..
You won't have to say, 'Well, is God going to take care of me?'".
Because you will be so in love with God, and so caught up in your relationship to Him,.
that you will know that His presence is what sustains you,.
and what encourages you, and what gives you joy and peace,.
whether the future is good,.
or whether the future is something that you might think of as bad..
You can say, "You know what? I'm going to enjoy God's presence in my relationship with Him.
if I'm healthy, or if I'm sick..
I'm going to enjoy my relationship with God and His presence in my life.
if I'm rich, or if I'm poor..
I'm going to enjoy my relationship with God and His presence in my life.
if I'm in a place of comfort, or a place of discomfort.".
And that is the testimony of Christianity throughout the ages..
There have been so many people in the history of Christianity.
who have lived and been imprisoned for their faith,.
or been persecuted, or even tortured or killed because they believed in Jesus,.

$^{321}$and yet in the midst of that, they were joyful,.
and they were full of peace and full of comfort..
That makes no sense, but it's because their priority was on loving God,.
and they were so caught up in their relationship to Him,.
that the earthly and physical realities didn't bother them anymore,.
and they had peace and comfort in the midst of something difficult,.
and uncomfortable, and painful..
And all of us can experience that..
We can experience that, whether things are good or bad..
We can find peace and comfort and joy in life..
And we can learn to love others, whether things are good or bad..
Most of us as humans, we love others conditionally,.
and when we feel good, and when someone else is making us feel good,.
then we love them back, and we treat them kindly..
If they're nice to us, then we're nice to them, and we all get along..
But if someone's unkind to us, then of course we're not going to be kind to them..
And if things are not going well in our life,.
then we're going to be grouchy and cranky, and we're going to snap at other people..
And Jesus is saying that if we will have loyalty to kingdom priorities,.
and say, "I'm going to love others, regardless of my circumstances,.
even if I'm mistreated, or even if I'm going through something difficult,.
I'm still going to love others,".
that we will experience God's goodness, and we will experience God's peace..
And we will not be so worried and anxious about what's going on..
And God will add those other things that we need into our life..
So I want to close out this study we've been doing on anxiety.
by just asking you to think about what it is that brings you anxiety..
And this is a question for me too. I grapple with this as well..
And I want to point you to the fact that whatever it is, God is bigger..
God's love is stronger..
God's grace is deeper..
God's power is greater..
And if you're having trouble believing that right now,.
if the problem, if the obstacle seems bigger than anything else,.
then I want to point you back to Scripture..
Go back to God's truth..
The Scripture tells us God's truth about reality as it really is,.
and about who we really are, and about the way the world really is..
Look at that and see what it says about you..
See what it says about your circumstances, and bring that into your situation..

$^{361}$So if you're anxious, and you're worried, and you're afraid,.
because of your past, maybe something that you've done,.
or something that's been done to you,.
and you feel like you can't get over it, you feel like you can't recover from it,.
you feel like you can't be forgiven, or you can't be clean,.
or you can't get free from it, or you'll never get over it or get past it,.
let me tell you what the Scripture says..
The Scripture says that God's grace is greater, and deeper,.
and stronger than anything that we've gone through..
Whether it was created by our own sin, or created by someone else's sin,.
someone who harmed us in some way, God's grace is greater..
And healing, and freedom, and victory is possible because of Jesus..
That's what the Scripture says..
If you're anxious, and worried, and fearful because of your future,.
and you think about things in your life that you want to happen,.
and you're afraid that they might not work out,.
and they might not go well for you,.
and they might not end up being what you need, or what you want in your life,.
look at the Scripture..
What does God say?.
What does Jesus say right in this passage?.
That God's love and care is complete and all-encompassing,.
even for the smallest animals and for the flowers of the field..
And if God is that good to them, you matter much more to God..
And He will be much better to you..
He will care for you..
He will love you..
He will protect you..
He will guide you..
He will be with you..
And even the difficult things that He allows to come into your life.
will be for your good to strengthen you, to grow you,.
to make you more like Jesus,.
and to bring about God's glory in your life..
Maybe you're anxious, or worried, or afraid.
because of relationships in your life,.
and things that you feel like are not going well,.
and they're beyond your control..
What does the Scripture say about that?.
The Scripture says that God is able to change any heart..

$^{401}$God is able to transform any situation..
God is able to defend people, even from lies and false accusations..
God is able to protect people,.
even from people who are out to do evil.
and who are intentionally trying to harm someone else..
Whatever you're dealing with in your relationships,.
God has the power to bring healing..
God has the power to protect you..
God has the power to restore..
God has the power to fix it..
And God wants to do that..
You can trust Him..
Maybe you're anxious, and afraid, and worried.
because you look at the world around us.
and you see all of the mess..
Maybe it's the COVID-19 pandemic,.
or maybe it's political unrest and turmoil,.
or maybe it's injustice,.
or maybe it's just things that go on.
in different parts of the world.
that you care about, but you can't fix them..
You can't change them..
You can't do anything about them..
What does the Scripture say about those things?.
The Scripture says that God holds the world in His hands,.
that God laughs at the nations.
because they're so weak before Him,.
and that men make their plans,.
and even evil men, and government leaders,.
and people in positions of power make their plans..
But ultimately, God rules over it all..
And God is a God of love, and care, and justice,.
and commitment to His people..
And so God takes His people safely.
through all of the turmoil, all of the circumstances,.
all of the plagues, and pestilence,.
and everything that happens..
That's what the Scripture says..
Whatever it is that causes you to be afraid,.
causes you to worry, causes you to be anxious,.

$^{441}$God is bigger..
You can take it to Him..
You can pour it out before Him..
You can trust Him..
And I would encourage you to do that in your life..
As we talked about that very first week,.
Jesus lays it out for us in this passage..
We've got to change our focus,.
put our eyes on God, and how great and good He is,.
and how great His love and care is for us..
Allow that change in focus.
to start to change our perspective,.
where suddenly the problems in our life.
are not such a big deal,.
and our relationship to God is what strengthens us,.
and anchors us, and sustains us,.
and gives us joy and hope,.
even in difficult times..
And allow that change in perspective.
to really start us growing in trust for Jesus,.
and saying, "God, I see who You are..
I see what You've done..
And because of that,.
I trust that You will take care of me..
I trust that You will love me..
I trust that You will bring good things into my life,.
and not evil..
And I trust that You will be in control of any situation,.
even the most difficult things in my life.".
And we can experience the peace,.
and the comfort, and the hope.
that God has for us.
when we begin to put that type of trust in Him..
I'm praying for you.
as you walk through anxiety, and fear,.
and worry in your life..
And you can pray for me..
And I know that God wants us to find peace in Jesus..
Let's pray..
God, thank You for Your truth,.

$^{481}$and for the comfort that it brings us, God..
I pray that You would help me,.
and help us to learn to release control,.
and to just trust You,.
and to say, "God, we are going to let go.
of these things in our life that cause us fear,.
these things we can't control,.
these things that knock us around,.
and get us sideways,.
and just get us confused,.
and cause us to be afraid and anxious..
God, we're going to release those things,.
give them up to You,.
let You have them,.
let You take them from our hands..
And we're going to trust You to do what You do..
What do You do?.
You do good things, God..
You work miracles..
You are a fixer..
You are a healer..
You are a repairer..
You are a builder..
You are the one who takes our brokenness.
and makes something beautiful, God..
And our world is full of brokenness,.
our lives are full of brokenness,.
but we trust You to fix, to heal,.
to restore, to repair..
And God, we thank You for that..
We praise You for that today..
And we ask that we would grow in our trust of You,.
and find peace, and comfort, and hope.
as we put our trust in You,.
to be at work in our hearts and in our lives..
In Jesus' name..
Amen..
\newpage



\section{}
\label{sec:OwmCWGNsDAQ}
\textbf{經緯線 - 耶穌的比喻 (EP4 十個童女的比喻)}
\newline
\newline
連結: \href{https://youtube.com/watch?v=OwmCWGNsDAQ}{\texttt{https://youtube.com/watch?v=OwmCWGNsDAQ}} ~~~~ 語音日期: 2021-02-16
\newline
\newline
\hyperref[sec:3j_OvRyfEuc]{\small{< < < PREV SERMON < < <}}
~
\hyperref[sec:index_chronic]{\small{[返順時目]}}
~
\hyperref[sec:index_scriptual]{\small{[返順卷目]}}
~
\hyperref[sec:UecWhHGM39A]{\small{> > > NEXT SERMON > > >}}
\newline
\newline
$^{1}$大家好,歡迎回來第四集的經緯線.
第一集我們看了回應神的話語應有的心態.
接著兩集看了兩個天國就好像的比喻.
這次來到第四和第五集.
我們會看兩個在《馬太福音》第25章的比喻.
今天我們會看25章1至13節《十個的同類》的比喻.
這個比喻說到天國好像有十個同類.
拿著燈去迎接新郎.
其中五個是愚蠢的,有五個是聰明的.
五個愚蠢的顧著拿燈,沒有預備油.
聰明的拿了燈,又拿了足夠的油.
新郎遲遲都沒有來,他等到很睏,還睡著了.
到半夜,突然有人叫新郎來,出來迎接他.
同類醒了,就點燈.
愚蠢的同類對聰明的同類說.
分些油給我們,因為我們的燈快熄了.
聰明的同類說.
不行,我的油不夠大家用,你們自己去買吧.
當愚蠢的同類去買油的時候,新郎就到了.
準備好的同類就和他們一起進入了婚宴.
然後門就關上了.
後來當其他的同類回來的時候.
就叫:主啊主啊,開門給我們吧.
但是新郎就說.
我實在告訴你們,我不認識你們.
最後耶穌就說.
所以你們要警醒.
因為你們不知道我回來的日子和時辰.
這個故事就是這樣.
我想和大家一樣.
我都有一個問題,究竟發生什麼事呢?.
為什麼當時的人會等新郎等到半夜這麼晚呢?.
這個就是我們以往一直查比喻的時候.
我們發現聖經耶穌的比喻.
就是和當時的生活.
他們的習俗是很有關的.
當時第一代讀者一聽就知道.
因為當時的結婚,婚嫁,飲宴的習俗.
不一定和我們現代人一樣.
但是現在很多很傳統的巴勒斯坦村莊.

$^{41}$他們都是這樣的.
他們一早就去接新娘到別的村.
伴郎跟著新郎就去接.
接完之後在那裡就慶祝一大輪.
整天都在跳舞.
跳完舞他們就回家了.
他們那種婚嫁的習俗.
接新娘回到新郎家的村.
那個接送是很重要的.
很慶祝,很歡樂,很快樂的.
然後最後晚上回來才很晚才會有婚宴.
所以婚宴通常都是很晚才會發生.
伴郎的同女需要在伴郎迎接新郎回來.
所以他們就要預備油,火把.
其實不只是古以色列是這樣.
當時的希羅文化其實很多時候都是這樣.
所以當時第一代讀者就知道.
那五個沒有預備油的,有沒有搞錯.
這些是很常見的東西.
明白,明白.
我們看了幾集都學到.
有時候原來這些蠢的,聰明的伴娘.
或者這些油的東西都未必是最重點.
幫助我們去理解那個比喻是什麼意思.
反而是看上文下理.
反而可以令我們認識,明白得透徹一點.
我們看了24章,好像比喻是在回應24章很多的問題.
Jackson可不可以跟我們解釋多一點這個比喻.
其實如果大家都很熟悉.
25章我們在開頭的時候都說過.
25章有三個比較詳細的比喻.
這是其中一個.
剛才Joyce也說到24章裡面都有.
不過在24章裡面看上去只有幾句.
而且內容也沒有25章的三個那麼豐富.
25章的比喻裡面再加上24章整章的教導.
其實耶穌也是在回應門徒問的一條問題.
就是在第24章的第三節.
這條問題不停地問.
就是他們問主耶穌究竟何時回來.

$^{81}$這條就是24章第三節問的問題.
然後我們看到在24章的36節.
到25章的第48節.
都是圍繞著門徒問的這條問題.
就是主回來的時候是何時.
不過我們在看的時候發現了一件很有趣的事情.
就是主的答案.
我們通常問主耶穌的事情.
主耶穌也不是給我們想要的答案.
這裡也不例外.
主的答案不是告訴門徒祂何時回來.
主的答案是在比喻裡面重覆重覆地說.
重覆重覆地說.
24章是36節.
42節是44節.
25章是13節.
都是說主耶穌回來的日子是不知道的.
是突然的.
所以主就說你們不要再問.
不要再糾纏我何時回來.
不過主耶穌說了一件很重要的事情.
就是在這個比喻裡面.
我們可以為那一天做好準備.
所以做好準備是很重要的.
其實有一點也很特別.
如果你看回24章的尾章的比喻.
是一個忠心還是不忠心的僕人.
那個是說他們以為主人會遲來.
先開派對玩玩.
反過來.
今天我們看這個.
他以為他會早來.
誰不知他就睡著了.
他還沒來.
他就道出一個重點.
無論他是遲來還是早來.
剛才說的其實沒有人知道.
耶穌自己也說.
他連他自己也不知道.
在這樣的情況下.

$^{121}$你是否準備好是很重要的.
其實還有一樣東西是值得提一提的.
大家不要忘記.
馬太福音是寫給第一世紀的一班信徒.
寫給他們的時候.
一定有原因的.
馬太用這麼長的篇幅.
來到24,25章也在說這件事的時候.
其實很明顯地告訴我們.
第一世紀的信徒當時是很渴望.
很迫切.
也希望耶穌回來.
我們也可能聽過.
在很多早期教會的書信裡面.
我們也看到.
在保羅的書信裡面.
在彼得後書當中.
也看到保羅和彼得也說.
主耶穌的回來你們是不知道的.
不過你們要警醒.
你們要警醒等候.
所以在當時來說.
真的好像剛才Selvin說的.
有些信徒真的很渴望著回來.
甚至可能有些已經渴望到想放棄的心態.
有些甚至已經離開了真道.
在這些比喻裡面.
正正提醒當時的信徒.
記不記得我們說過.
主耶穌是會回來的.
所以你們要警醒.
要警醒.
就是這樣.
所以對我們來說也很迫切.
其實我們跟第一世紀的信徒.
心態也可能是這樣.
有些很想他回來.
現在回來.
有些已經等到放棄了.
尤其是現在疫情.

$^{161}$放棄了.
我們看數據.
人們也說將來教會重開的時候.
一定少了人.
如果是這樣的話.
究竟怎樣為之聰明.
怎樣為之蠢.
我們現在可以開燈.
關燈按個鍵就行了.
不用準備游泳.
究竟他在說.
跟有一個準備好的心態.
有什麼關係呢.
其實這裡說的有沒有智慧.
其實跟準不準備好.
很有關連.
很多時候我們以為有沒有智慧.
純粹是地上的智慧.
或者.
他很聰明.
但是.
這裡說的智慧.
其實是一個.
很關乎到.
英文叫ascaton.
希臘文.
那個末日.
是不是你已經準備好.
你有沒有在你的.
心態上.
在你每天的.
行徑上.
其實有那個末日.
在你的心中.
這個很重要.
這裡說的這個比喻就是.
如果你知道耶穌會回來.
不單止會回來.
還要在你不知道的日子.
回來.

$^{201}$所以你要準備好.
那就是有智慧.
所以說的是一個心態.
一個的毀身.
甚至是說.
整個的生活.
不是某一刻.
我有準備好游泳.
我加油的時候.
有下車的時候有很多桶油.
不是關於這個.
而是關於你整個的生命.
以致到你.
get yourself prepared.
and fit for.
上帝的國度.
上帝的管治.
隨時臨到.
其實現在都在.
但是.
你準備好沒有,這件事很重要.
所以.
有一個quote.
有人說是馬丁路德說的.
他說.
只有兩個日子.
就是今天和那一天.
那一天就是剛才正正是Selvan你說.
煮回來的那一天.
他就說用那一天的心態.
去過每一天.
去過每一天的今天.
明白,那我們如果要學.
像聰明的童女.
是不是要準備夠很多油.
別人問我都說.
因為不是的話就不夠用.
這些都是我們查比喻.
會引申出來的問題.
或者會很容易跌落去.

$^{241}$就是我們覺得有些什麼moral.
有些什麼.
好像醫學語言那樣有些什麼可以學習.
這裡就好像教你.
去學那五個聰明的自私一點.
是這意思.
所以好像Justin一開始這樣說.
我想不是.
只是去拿那些.
bit and bits.
要看一個比較整全的圖畫.
特別是上下文.
所以他在這裡說.
你們要警醒.
那一天的時刻.
是你們不知道的.
這是一個很大的hints.
你需要去之前.
和之後說過這個相關.
說話的地方去看.
如果說到這裡的時候.
有一個keyword就是馬太科.
就是主呀主呀.
這個比喻裡面.
主人就和那些同類說.
我實在告訴你們我不認識你們.
這句話聽過馬太科.
也有一點印象.
你記得登山寶訓.
第一個block.
就是耶穌的教導五至七章.
七章最後的時候說什麼.
也是耶穌說這句話.
不是每個人叫我主呀主呀.
就能進入天國.
因為我是不認識你們從來都不認識你們.
你說有行乘職有路.
按我的名做很多事.
我根本不認識你們.
登山寶訓也出現.

$^{281}$最後的時候就是聰明.
和愚蠢的.
建房屋.
最後結束的時候.
聽了我的道有遵行的.
有kingdom lifestyle.
這些人才是屬於我的門徒.
回到這裡.
再提醒你.
如何預備呢.
這個首尾呼應很有趣.
因為我們也在去馬太福音的結束.
第七章開頭.
接著就到第25章.
其實這個問題也問得很好.
就是你發覺.
耶穌說我不認識你們.
當我再想想這個比喻的時候.
你發現一件事很有趣.
就是如果這些童女.
她們對主人.
對新郎有心的時候.
他們就真的會有準備.
他們也知道那個culture.
會知道各方面.
可能會遲.
玩到瘋了.
他們連準備都不去.
其實你看到他們根本就沒有心.
對主人的關係也不是一般.
我們發覺這個和主人的關係.
一般般.
也是貫穿著我們幾個的比喻.
我們會看到的.
接著下一集.
我們又是看到同一樣東西.
就是那個人和主人的關係.
也是一般般.
還要賴主人一會.
所以換句話說.

$^{321}$其實你和上帝.
或者在天國裡的關係.
其實是沒有一般的.
有時候.
剛才Joyce也說過.
是不是.
可以.
這樣的時候.
我都信了耶穌.
沒有問題,夠不夠油.
是不是最後.
連天堂也不能去.
很多時候我們會想.
問這些問題.
如果真的很認真進入.
這個問題的核心.
其實.
如果純粹想找一個shortcut.
you ask the wrong question.
剛才Mark Sir也說.
這是一個心態.
一個生命.
一個你和上主的.
一個長遠的關係.
不是一刻的.
沒有辦法.
我想這件事.
就是一線之差.
如果我們不是.
斟酌一些很細的細節.
來看應用的話.
整個比喻.
放回今天的處境.
我們又應該怎樣.
理解這段比喻來應用呢.
我想一個很重要的.
就是肯定.
主耶穌的回來.
剛才Justin也說.
特別是現在的pandemic.

$^{361}$有些可能就會.
很灰心.
遲些可能連教會也不想再回來.
有些可能.
很渴望耶穌.
回來,但是.
也好像.
在那裡待不住.
我想回到.
核心.
這個耶穌的.
回來,這個天國的.
臨在,其實是.
事實來的.
這個acknowledgement是很重要.
我想當有這個.
你就要知道.
這個天國,就是耶穌自己說.
已經在我們當中.
但又未完全來到.
already but not yet.
我們就要持守著一個.
心態去準備好.
我們既已知道.
有這個現實.
甚至有人說.
用現實可能不是最正確.
這是將來.
一個future,但是.
插入了.
我們今時的時刻.
在這樣的情況下.
我們就應該活出將來.
不會發生的事.
這樣的心態,這個就要準備好.
我想這件事是很重要.
是時刻.
在我們所做的事情.
在我們怎樣去活每一天.
都相關.

$^{401}$就不是只靠雙手.
等一天,其實是等死.
等天堂來.
而是你知道.
上帝的國度.
是已經臨在.
又要將來.
完成的.
我們今天怎樣活呢?.
當然是說可不可以more specific.
有時很難,因為其實.
這就相關到每個人.
在自己的處境裡面怎樣去活.
剛才我講到那個今天和.
那天,即是quote.
應該是馬丁路德講的.
那時我都講過這篇道.
在這篇道裡面,我最後.
自己有一個反省.
對一班信徒來說.
我們經常說將來要得到主的.
稱讚,得到主的賞賜.
我們現在這個比喻都是圍繞著這些.
在第二.
第三.
其實將來我們都會和主一齊.
歡樂慶祝,特別在這個比喻當中.
結論.
原因是因為我們今生的努力.
結論原因.
是因為我們今生現在的忠心.
所以將來就可以得到主的.
稱讚,所以將來就可以.
主一同來到他的歡樂.
希望看完這個比喻.
之後,大家都會可以.
重新審視自己和神的.
關係,警醒準備好.
因為知道耶穌是一定.
會再回來的,我們下次再見.

$^{441}$拜拜.
\newpage



\section{}
\label{sec:UecWhHGM39A}
\textbf{Kingdom Values XLI - Judgement}
\newline
\newline
連結: \href{https://youtube.com/watch?v=UecWhHGM39A}{\texttt{https://youtube.com/watch?v=UecWhHGM39A}} ~~~~ 語音日期: 2021-02-21
\newline
\newline
\hyperref[sec:OwmCWGNsDAQ]{\small{< < < PREV SERMON < < <}}
~
\hyperref[sec:index_chronic]{\small{[返順時目]}}
~
\hyperref[sec:index_scriptual]{\small{[返順卷目]}}
~
\hyperref[sec:bJj9ij22kAs]{\small{> > > NEXT SERMON > > >}}
\newline
\newline
$^{1}$Hey family, here we are again enjoying this beautiful Chicago winter. You know.
for a while I thought we were gonna have a really mild winter this year and it.
seemed that way all the way up until about the middle of January and wasn't.
that cold, wasn't much snow, seemed like things weren't gonna be that bad and.
then winter just came roaring back and it has not let up since. It has been very.
cold. We've had a crazy amount of snow that we're all trying to dig out from.
under and we're shoveling snow every day and I know it's probably been a little.
bit tiring and challenging for all of us and I'm praying for you as you go.
through it, not just the challenges of shoveling snow and encountering those.
things but also just kind of the emotional challenges of being isolated.
and stuck in the house and not able to do things compounded with already being.
isolated because of the pandemic. I know sometimes in the winter it just seems.
dark and dreary and so I'm praying for you all as you encounter just all this.
winter weather that you'll stay encouraged and that you'll be able to.
enjoy it at least a little bit. And we're once again gonna be looking in Matthew.
today at the Sermon on the Mount. Today we're looking at Matthew chapter 7 so I.
just want to start off by commending you and encouraging you to give yourself a.
little round of applause or later on today make yourself a little trophy.
something like that because you've come a long way. The Sermon on the Mount is.
three chapters in the book of Matthew, Matthew 5, 6, and 7. It's the longest.
continuous sermon that's recorded in the Gospels, the longest continuous sermon of.
Jesus. There's a lot of content there and we've been going slowly through it.
taking it a verse or two at a time and so I know it's taken us a while to get.
through it. So I commend you for sticking with me, I commend you for your patience.
and perseverance and Matthew 7 is the last chapter we're starting in this week.
so there's a lot at the end of the tunnel, we're getting close to the end.
but Matthew chapter 7 is full of a lot of really good truth for us. I want us to.
really dig into this chapter and enjoy learning what God has for us and I know.
for me studying through the Sermon on the Mount it's it's really kind of life.
changing. There's so much deep and powerful truth in it and it's just been.
so transformative and such a blessing for me. I pray it's been the same for you.
and this Matthew chapter 7 is no different. In Matthew chapter 7 the first.
section of this passage Jesus is really talking about judging others and so.
that's what we're going to be talking about today, judging others and this is a.
really important topic and it's really important in our culture and so we're.
going to take our time and look at it over the next few weeks, go through it.
kind of slowly and try to really discover the truth that God has for us.
on this topic and really get a healthy understanding of what Jesus is saying in.
these verses. Today we're just going to look at Matthew chapter 7 verse 1 and.
verse 2 so let's read these together. Jesus says, "Judge not that you be not.

$^{41}$judged for with the judgment you pronounce you will be judged and with.
the measure you use it will be measured to you." Now these verses are very well.
known Bible verses. Matthew 7 1 is often said to be one of the most popular Bible.
verses in the entire Bible. A lot of people know this verse and sometimes.
people say it's the most popular verse of people who are not Christians, people.
who are not a part of the family of God because they love to quote it to say, "Hey.
don't judge, you can't judge, Jesus said don't judge lest you be judged." So you.
may have heard this verse before, maybe you've heard it as someone has said it.
to you or maybe you've heard someone on TV say it or something like that. It's a.
verse that people say a lot but it's also misunderstood in a lot of ways and.
it's difficult to interpret. And what I believe and what I want us to examine.
today is that Jesus here is not saying that we should not judge at all but he's.
saying that we should exercise right judgment when we judge and what he's.
really calling out is having a judgmental attitude. He does not want us.
to have a judgmental attitude and what I think what most people think of when.
they say judging is really a judgmental attitude. Obviously a lot of people in.
our culture have a very negative perception of judgment and of judging.
and so a lot of people say, "Well you can't judge others." But I think a lot of.
times what people are talking about is actually a judgmental attitude that is.
proud or that is hypocritical or that is self-righteous and those are definitely.
things that Jesus wants us to avoid and we're going to talk about that over the.
next few weeks. But Jesus is not saying that we shouldn't judge or that we.
shouldn't make judgments or that we shouldn't exercise discernment. I really.
don't think Jesus is saying that in this verse. And if you just have such a.
negative connotation with with the idea of judgment that it's hard for you to say.
the word or think about it judging something or some circumstance, then you.
can just use the word discernment. Discernment is a biblical word and I.
think it's a really good word for us to think about. Discernment refers to being.
able to see or recognize or distinguish or apprehend things and so discernment.
is a synonym of judging and it's really an important quality for us.
Christians to have and to be able to exercise. So if you're uncomfortable with.
the concept of judging then you can talk about discernment. I'm going to use.
those two terms almost interchangeably as we discuss this. And I think to start.
it's helpful to think about what judgment is. We can say, "Well judgment is.
when you judge something, when you make a decision about something.".
And that's true, but I think on a deeper level it helps us to understand that.
judgment is ultimately comparison to a standard. It's when we compare something.
to a standard and say, "Does this line up with the standard? Is this right.
according to the standard?" And for Christians, for followers of Jesus, for.
the subjects of the kingdom that we've talked about in the Sermon on the Mount,.

$^{81}$that standard ultimately comes from God. It's God's truth. For people who are not.
followers of Jesus, that standard could be a lot of different things. It could be.
their own feelings, it could be what someone else has told them, it could be.
all types of different things. It could be cultural values. And for believers and.
people in the church, a lot of times that's true too. Instead of sticking to.
the standard of God's truth, we start to hold on to a lot of different cultural.
values or our own feelings or what someone else has taught us. And that's.
not always healthy. A lot of times that can lead to really unhealthy attitudes.
and unhealthy judgments. And that can lead to that judgmental attitude that we.
want to avoid and that Jesus is calling us to avoid in this passage. So we want.
to remind ourselves and to always be mindful of the fact that when we are.
thinking about discerning something or judging something, we're talking.
about comparing it to God's truth and to God's standard. And remember that God's.
judgment is always rooted in God's justice. God has a standard of justice.
that is above and beyond ours and God's justice transcends our idea of justice..
And we need to understand that. That God's judgment is always rooted in God's.
standard of justice, not ours. And we also need to understand that God's justice is.
rooted in His love. And so even when God exercises judgment, when He compares.
something to a standard, when He asks us to compare ourselves to a standard,.
that's always rooted in love for us. God does not do that to punish us. God does.
not do that to expose our weakness or our failing. God does that because He.
loves us and He wants to grow us and strengthen us and transform us into.
something more beautiful than what we have been in the past. God's judgment and.
justice are always rooted in His love. Now you might look at this verse on the.
face of it and say, "Hey, judge not that you be not judged." I don't see anything.
in there about having a judgmental attitude. So why am I saying that Jesus.
is calling out having a judgmental attitude, not just calling out exercising.
judgment in this passage? I want to remind you of how we interpret passages.
in the Bible. We talked about this with some of the other passages in the.
Sermon on the Mount that are a little bit more difficult to interpret. There's.
a lot of passages in the Bible we may look at and say, "I don't fully.
understand it," or "it seems confusing," or "is it saying this because it seems like.
that doesn't agree with other parts of the Bible?" And so there's a lot of.
principles of biblical interpretation. We went through some of those when we.
talked about some of those other passages, so we're not going to go.
through all of them again today. But I'm just going to point to a couple that we.
can really see in this passage. And that is it's always important to examine.
verses in their context, in the context of the verses around them. And that's.
really helpful to look at the passage and see what the passage is saying. And.
it's also helpful to examine verses in light of other scripture and other.

$^{121}$verses and passages throughout scripture and compare and say, "Okay,.
did these agree?" And that helps us to interpret a verse that may seem unusual.
to us when we first see it. And this is one of those verses. If we just look at.
Matthew 7, we say, "Judge not that you be not judged." Jesus is clearly.
saying that we should never exercise judgment. But when we look at the context.
of this passage, this section in Matthew 7, Jesus is clearly talking about judging.
others and is clearly not saying that we shouldn't do it. He's just calling us to.
do it in the right way. So in verse 6, he's talking about judging and knowing.
when it's appropriate to say something or to reveal something to someone else..
That's an issue of judgment. In verses 15 to 20, he's talking about identifying.
false teachers by the fruits and by the actions of their life. That's clearly.
something where he's telling us to exercise judgment and discernment. So the.
context helps us to see that, but the rest of scripture helps us to see that.
too. Throughout the scripture in the Old and New Testament, God calls on us to.
discern, to exercise judgment. He just wants us to do it in the right way..
There's even a verse in John 7, verse 24, where Jesus says, "Do not judge by.
appearances, but judge with right judgment." So again, I think the focus of.
this passage is not that we should not judge at all, but it's rather that we.
should not have a judgmental attitude. And having a judgmental attitude is.
really unhealthy. It's really terrible. It causes a lot of problems. And having a.
judgmental attitude is when we make judgments or just try to discern things.
with a spirit of pride and arrogance, when we don't exercise grace towards.
other people, when we don't see how much God has forgiven us, when we don't, we're.
not meek, when we're condemning of other people, when we try to sentence other.
people or hold other people accountable, when we judge based upon what we can see.
and think that we know someone's intent or heart based upon the outward.
appearance, all those things are wrong. And they lead to hypocrisy and.
self-righteousness and all of these hurtful things. And sadly, a lot of that.
type of judgmental attitude is found in the church and has been found in the.
church. And so Jesus is not saying that he doesn't want us to exercise discernment..
He's saying he wants us to avoid that judgmental attitude. And we're gonna.
examine all of those things over the next few weeks. But today specifically, I.
want us to start off by looking at some unhealthy attitudes about judgment that.
we probably have on some level. And these are attitudes in the church. They're also.
attitudes in culture at large. And sometimes they can kind of slip into our.
way of living and our way of thinking and we don't even realize it. But they.
can cause us to be unhealthy in different ways. So the first thing I want.
us to think about is just the attitude that, you know, no judgment at all is.
acceptable. And this is a pretty common cultural attitude that it's wrong to.
judge anything or anybody. It's unkind in some way. And it can seep into our.

$^{161}$life. Sometimes in the church as believers we can find ourselves saying,.
"Is it okay for me to even make a decision about certain things or to even.
try to discern what's right or what's wrong?" And I think it's clear that.
judgment is a necessary part of life. It's an essential part of our daily life..
You probably already made a lot of judgments this morning. When you woke up,.
you probably looked at the alarm clock and based upon what time it was, you made.
a judgment decision as to whether or not you could go back to sleep or whether.
you had to get up and get ready to do whatever it was you wanted to do today,.
including get on your computer or your phone or your device so that you could.
view the stream. You made a judgment decision. If you went outside or.
if you're going to go outside later, you'll make a judgment decision then..
You'll look at your phone or something else that tells you the temperature or.
maybe you'll stick your head out the door and just determine the temperature.
that way and you'll make a judgment decision based upon that about what you.
need to wear in order to stay safe today. And so all of us make judgment decisions.
all of the time and ultimately they're essential for our life. We couldn't.
function without them. We have to make them in order to be safe and that's.
important. And we have to realize that the reality of our world as God has.
created it is that not all actions have the same outcome. And that's something.
that we don't always want to admit. We want to try to make a world in our own.
minds where even if we do something that's not a good idea or it's.
unhealthy, it's still going to turn out great for us. But we can't make a world.
like that. We can't make our world be that way because that's not the world.
God created. God created a world where not all actions have the same outcome..
And to go back to the illustration I used just a moment ago, if you go outside.
on a day like today in shorts and a t-shirt, that is not going to have the.
same outcome as if you go outside in a warm coat. Those are two different.
actions and they're going to feel very different depending upon which one you.
choose. And so making judgments helps you choose the action that's going to be.
healthy and right for your life and that's a good thing. We need to do that.
as humans and God wants us to do that. And the reality is that all people make.
those types of judgments. So saying that, "Well, we can't make any type of.
judgments. We can't discern anything that to even try is unhealthy or unkind.
or something like that." That's not really a good argument because all of us have.
to make those types of judgments about life just in order to function and to be.
safe and to do all the things that we need to do. But maybe we say, "Okay, so.
yeah, we have to make judgments about our daily life so that we can function. We.
get that." But the problem is when we begin to judge other people and we begin.
to maybe say something that is not affirming of someone else. Well, again, I.
would say that that's an unhealthy attitude towards judgment because we.

$^{201}$have to discern things and make decisions about other people. That's a.
reality of life. Every relationship that you're in involves making some level of.
discernment and judgment in that relationship. Why? Because relationships.
involve trust. And to trust someone you have to have some knowledge about them..
And so any relationship, whether it's a friendship or someone at work or someone.
at school or a much deeper relationship like a marriage, is going to involve.
communication. It's going to involve response to that communication. And during.
that process you're discerning, you're saying, "Hey, is this person someone that.
I can trust? If I share the secret with them, are they going to share it with.
someone else? If I become transparent or vulnerable in this way, will I be safe?".
Those are judgment decisions and all of us do that in relationships. All of us.
have to make decisions about who we can trust and who we can share our life with.
and what level we're going to share things with. And there's nothing wrong.
with that. That's a necessary part of living life and forming relationships in.
the world that we live in. And ultimately relationships would be impossible without.
some form of discernment and judgment. All of us understand that in our.
relationships. And again, not all actions have the same outcome. So at some point.
we have to be able to discern what's going on in the life of someone else. A.
simple example of this is like driving. If you pull up to a four-way stop or a.
red light, one of the things that you do is you're looking at the other drivers.
as they come up to the stop. And even if you have the right-of-way and it's your.
turn to go, if you see a driver that isn't slowing down and looks like.
they're going to blow through the stop sign or that they're going to go through.
the light even though it's red, then you make a judgment call to protect yourself.
and anyone who's riding in the car that you're in. And you say, "I'm not going to.
go even though I have the right-of-way, even though it's my turn to go. I'm not.
going to go because I see the action of someone else and it could harm me or it.
could harm the people around me." That's a judgment call and we have to make.
those. They're a part of our life. And even in the last year with the pandemic,.
we've seen ourselves as humans making so many judgment calls where we look at.
other people and their actions and say, "Hey, are their actions going to protect my.
health? Are their actions going to protect the health of the people around me?" And.
we make decisions based upon that. We say, "Okay, well here's what I need to do.
because of what I see someone else doing." Or, "Here's a situation that I can't be.
in or something that I can't go to because it might endanger my health." And.
so the reality is we're all making those types of judgments and discernment-based.
decisions all the time, even when they involve the actions of other people. And.
that's not wrong. It's okay to make those types of decisions. Again, what we want.
to avoid is being judgmental in a wrong way towards other people. So maybe we say,.
"Well, okay, I get that we need to make judgments in life and even some.

$^{241}$judgments about the actions of others are acceptable, but the problem is moral.
judgments that when we judge things as good or bad, that's really a problem." And.
I think that's a criticism that has been leveled at Christians and at the church.
a lot that people have said the problem with Christians is that they make all.
these moral judgments and they're always trying to say what's good and bad and.
that's really the problem. And it just turns people off and it.
makes people upset. Well again, I would say that that argument's not really a.
great argument to say that no moral judgments are acceptable. The reality is.
that everyone has a moral standard in their life. And honestly, sometimes we can.
see our moral standard or someone else's moral standard much better by our.
reactions than by our actions. What do I mean by that? Well, I mean, we might.
say, "You know what? I get really irritated at someone else because they judge me and.
I took some of their food out of the fridge, my family member, and they're so.
mad at me about it. What's the big deal? It's just a little bit of food." We might.
say, "It's not a big deal. My actions aren't that much of a problem." But.
our standard is revealed by our reaction. If someone takes something from us, that's.
when we reveal our true standard. We might say, "Oh, my standards are easy." But.
when someone wrongs us, how do we react? And when someone steals from us, suddenly.
we're angry and they need to be held accountable and they need to.
pay everything back to the last penny. A lot of times we can see our moral.
standards more by our reactions than by our actions. And the reality is everybody.
has a moral standard. Everybody has things that they think are right and.
things that they think are wrong. And everybody makes moral judgments..
Everybody looks at things in life and says, "Hey, this is right or this is wrong.".
We again have seen that so much in the last year in our politics. You know, as I.
said, a lot of times this claim about moral judgment not being good or healthy.
is leveled against Christians and people in the church. But we've seen a lot in.
the last year, a lot of moral judgments. People saying what is good, what is bad,.
what is right, what is wrong. And the majority of those people are not.
Christians. They're not people in the church. Many of them are not even.
religious people, but they've made very strong statements about what is right.
and what is wrong. And that shows us that all of us have some standard of morality,.
some standard of what we think is right and is wrong. And we're willing to make.
judgments based on that standard. So the real question is, what standard are we.
going to compare to? It's not a question of whether or not we're going to compare.
things to a standard, whether or not we think anything is right or wrong. All of.
us have some things that we think are wrong and off-limits for humans. And all.
of us have other things that we think are good and healthy for humans. So.
that's not the question. The real question is, what standard are we.
comparing to? Whose standard are we going to use? And as I said at the beginning,.

$^{281}$for followers of Jesus, our standard is God's truth. For other people, there's a.
lot of different standards. But for followers of Jesus, we want to go back to.
God's truth. And we have to understand as believers that making moral judgments.
and discernments is actually an essential part of our faith. Again and.
again in Scripture, we're called to not be deceived. We're called to cling to the.
truth. We're called to be obedient to God. We're called to live righteously. All.
this stuff. Well, how do we do that if we are not able to make moral judgments and.
decisions that involve discernment, decisions that involve seeing and.
recognizing and distinguishing and apprehending right from wrong? And so.
ultimately, we have to make moral judgments in life. And Jesus is not.
banning us from doing that. He's just calling us to do it in the right way and.
in a healthy way and in a way that is not unkind to others. In a way that's.
respectful of others. And as we said at the beginning, in a way that is rooted.
and grounded in God's love and grace. God's love and grace for us and God's.
love and grace for others in the world. So like I said, I want us to close out.
here seeing that Jesus is calling us to exercise right judgment and to not have.
a judgmental attitude. Next week, we're gonna get more into the specifics of how.
to avoid having a judgmental attitude. But for this week, I just want us to see.
how important it is that we are able to discern things in life and that we're.
able to make healthy judgments and healthy decisions and that we are able.
to discern what is right and wrong. What's good for us and what's not good.
for us. And that we're able to look at the world around us and have some.
clarity about what's happening. That's so important and Jesus is calling us to.
that in this verse. And that's something that we should desire and that's.
something that we should practice and something that we should grow in as.
believers and as followers of Jesus. We should grow in our ability to discern..
It's something that we can practice and it's a skill that we can develop by.
God's grace. And so we want to do that. Let's pray. God, thank you that you give.
us the ability to discern things, to see, and to distinguish, and to recognize. And.
God, we are so weak and limited in ability, God. Our vision is so short and.
so God, we need your help. We need your help so much in discernment, in making.
judgments, in understanding what's right and wrong. And so God, I pray that you.
would help us. That you would help us to grow in our ability to do that. That you.
would help us to grow in our ability to look at your word and to see what your.
word is saying about the issues of our life. That you would help us to grow in our.
sensitivity to listening to the Holy Spirit and how you want to use the Holy.
Spirit to speak to us. God, please help us to grow in discernment and in making.
right judgments so that we can distinguish right from wrong in our.
lives and we can have loving relationships with others. And God,.
please help us to not be judgmental and unkind and disrespectful in any way to.

$^{321}$the people around us, God. Help us to exercise right judgment and right.
discernment as we go through life. In Jesus' name, Amen..
\newpage



\section{}
\label{sec:bJj9ij22kAs}
\textbf{經緯線 - 耶穌的比喻 (EP5 才幹的比喻)}
\newline
\newline
連結: \href{https://youtube.com/watch?v=bJj9ij22kAs}{\texttt{https://youtube.com/watch?v=bJj9ij22kAs}} ~~~~ 語音日期: 2021-02-24
\newline
\newline
\hyperref[sec:UecWhHGM39A]{\small{< < < PREV SERMON < < <}}
~
\hyperref[sec:index_chronic]{\small{[返順時目]}}
~
\hyperref[sec:index_scriptual]{\small{[返順卷目]}}
~
\hyperref[sec:jf7dAWwLNe0]{\small{> > > NEXT SERMON > > >}}
\newline
\newline
$^{1}$大家好,歡迎回來《經緯線》.
今天就到這個系列的最後一集,第五集.
我們會看《馬太福音》第25章14-30節.
《才幹的比喻》.
這個比喻說到天國好像有個人要出玄門.
他就叫了他的奴僕回來,把產業交給他們.
按照他們的才幹,分別給了一個人5000銀幣.
一個人就給了2000,一個人就給了1000.
接著就走了.
有5000銀幣和有2000銀幣的兩個人.
就馬上去做買賣,賺了一個雙倍.
但是拿1000銀幣的那個人就去挖了一個洞.
就把錢埋在地底.
過了很久之後,主人就回來跟他們計數.
5000和2000銀幣的那個人就跟主人說.
你看看你給我的錢,我又幫你賺了一個雙倍.
主人就跟他們說.
做得好,你真是良善有中心的奴僕.
既然你在小事上中心,我就會把更大的事交給你們.
過來分享你主人的快樂.
(歡呼聲).
然後,得到1000銀幣的那個奴僕又上前去跟主人說.
主人,我知道你很苛刻.
沒有播種的地方都要收割.
沒有殺種的地方都要有收穫.
我就很害怕,於是我就把你的錢埋在地底.
看看這裡就是你的錢.
主人就說,你這個又可惡又懶的奴僕.
既然你知道我在沒有播種的地方都要收割.
沒有殺種的地方都要收穫的時候.
你就應該把錢存入銀行.
等我回來的時候,至少可以連本帶利給我.
然後主人就說,還說收回那1000個銀幣.
給那賺了總共有一萬元的奴僕.
因為凡有的,還要給他更多.
使他豐富有餘.
凡沒有的,連他僅有的,也要奪去.
把這個沒用的奴僕扔出去.
(歡呼聲).
在他那裡就要哀哭切齒.

$^{41}$這個比喻就說到這裡.
都很精彩,好像在看電視劇一樣.
我想問問Mark Sir.
究竟那些有錢人是不是真的.
搭飛機去旅行,要出雨門的時候.
就會把東西分給那些人呢?.
因為以前古代,我們可以想像.
沒有我們現在那麼容易的.
上路,或者去很長的旅程.
因為古代的財主,地主.
他們很多時候都要去旅行.
或者好像大熙和王那些.
他們的子女,他們有時候要上京去羅馬.
所以有時候要去很長的旅程.
但是以前旅行就沒有我們現在那麼方便.
很多時候會有很多險阻,或者很多未知之數.
所以你可以想像,你會知道.
講這些比喻就是當時很熟悉.
上路,當然要把東西交下.
因為不知道會發生什麼事.
路上有很多不確定性.
所以主人理所當然地就交給三個僕人.
把家產,東西叫他們管理.
剛才你說的五千,二千,一千.
你看原文應該是五個他連得.
兩個他連得,一個他連得.
一個他連得其實是磅銀.
每個他連得大概七十幾磅,七十五磅左右.
如果你這樣算的話.
一個他連得大概是二十五萬美金.
是很多錢,是很大一筆錢.
主人其實這樣分法.
你今天算,分了二百多萬出去.
他自己的家業給三個僕人管理.
所以為什麼有些現代的翻譯本.
就把他,不是譯做五二一.
因為五二一我們感覺好像很少.
就把他譯做幾萬個銀幣等等.
總之你就知道,這個很多錢交給他.
交給一些僕人的時候.

$^{81}$就希望他能夠好好管理.
他很大筆的家業.
這個比喻放在這個位置.
又和之前的比喻或者跟之前的篇章有什麼關係呢?.
如果我們一直有跟開這個節目的朋友.
大家都記得在《馬太福音》裡面.
都是分了五個耶穌很主要的教導在我們當中.
在五個教導裡面.
我們現在來到第五個教導的時候.
就是集中在耶穌的回來.
但不是集中在什麼時候回來.
是集中在我們一班作為跟從主的人.
可以怎樣的來準備.
在這個比喻裡面.
就開始似乎讓我們看到多一點.
來到他怎樣的準備.
在這個比喻當中.
有一樣東西大家留意.
就是在二十五章裡面.
他是處於第二個位置.
跟著第三個比喻就是山羊和綿羊的比喻.
都是說到主將來審判的時候.
那一班人是得到獎賞的.
所以在我們這個比喻裡面.
我們看到有主人離開.
將他的錢,身家交給三個僕人.
我會從這個關係裡面入手.
因為我上次那一集都說到.
那些同類的五個愚蠢.
他們不是因為智商愚蠢.
是因為他們來到跟主人,新郎沒有關係.
如果他們跟他有關係的時候.
他們就應該有準備.
在這裡其實都讓我們看到.
第三個僕人其實都很不認識主人.
大家有沒有留意.
剛才Joyce讀聖經的時候.
有一句說話就是.
主人是按照著他們的才幹來比的.
換句話說.

$^{121}$主人是認識這些的僕人.
跟著在這個比喻裡面.
我們看到第三個僕人.
其實如果我們細心讀的時候.
其實那個僕人是非常懶惰.
他諸多藉口.
因為主人說.
如果我真的像你所說.
我是這麼橫蠻無理的.
沒有種的地方要收割.
沒有散的地方要罪孽.
既然你知道我是這樣的時候.
假如我真的像你這樣說的時候.
那為什麼你不將錢放在銀行裡面.
所以其實換句話說.
這個僕人根本.
他說每一句話都是假的.
是藉口.
是一個又懶又惡的僕人.
我想剛才Justin說的.
一件事很重要.
是說一個關係.
所以如果你看看那個結果.
其實都是反映一件事.
如果你是良善忠心的.
那個結果.
主人說.
不單止是分更加重要的東西給他做.
更加重要的是進入我的喜樂當中.
其實是在關係當中.
你可以繼續升華.
可以繼續去有份.
但是那個又懶又惡的僕人.
最後要怎樣.
是放他去黑暗當中哀哭切齒.
是不能和主人有份.
我想這件事就導出關係的重要性.
沒錯.
如果繼續從這個關係去想的時候.
正如開頭所說.

$^{161}$安著財綱.
主人是認識那三位的僕人.
換句話說.
就是說其實主人是對三位的僕人.
都是有期望的.
希望他們藉著給他這一筆錢.
是期待著他們有些動作.
希望可以搞一些東西出來增值.
去投資也好.
主人是期待著他們這樣做的.
所以那個比喻.
這個關係,期望.
不要輕看主人對你的期望.
都是充斥在故事背後的氣氛.
其實Jason你這樣說.
我們回想第一集.
我們說這個殺種的比喻.
其實也有相關很類似的.
也是說有增長.
有配數增長.
甚至乎其實是一句是重複了的.
就是凡有的就再給他.
沒有的就連他有的也拿走.
其實就是在說這個.
從天國的角度.
對一個屬於天國的人的期望.
我想這個比喻也說得很清楚.
我們真的要很清楚.
究竟神對我們的期望是什麼.
還有我們自己有沒有.
達到祂給我們的期望.
讓我們可以清楚知道.
將來天國也有我們的氛圍.
不過我之前也有聽過.
會不會在這個一千銀幣的羅浮身上.
也可以學到一些東西.
就是主人既然是這麼嚴苛.
又這麼壞.
他好像在剝削人類.
有這麼高的期望.

$^{201}$所以他好像在抗議.
不行.
我就把錢放在這裡.
給你就算了.
就什麼都不做.
不要給他.
越有錢就有錢.
這個有時候有一些.
如果你看過那些叫做社會科學的評論.
就是說古代的經濟.
社會的運作.
那種方法去看.
是真的好像Joyce你這麼說的.
就是那一千個才是抵抗地產霸權的分子.
那個才是英雄.
但是我想如果這樣看的時候.
沒錯有那個背景.
首先我們看那個比喻的時候.
不可以每樣東西都拿出來作比較.
因為那個比喻本身是有它的方向.
不是每樣細節都可以比較.
或者每樣細節都有一個教導在裡面.
如果你說耶穌是不是仇父呢.
可能不是這裡說的.
有很多地方都有說.
就是說有錢的人有禍等等.
潔淨聖殿等等.
這個比喻說出來不是說縱容剝削人的主人.
而是他想拿這個當時人能夠一聽就明白的故事.
他在說什麼呢.
就是在說不是經濟的問題.
而是說我們怎麼預備.
主人再回來我們主宰.
為什麼我們一定要看上下文呢.
以致我們比較或者想去應用的時候.
我們不會搞錯了.
拿了其他不需要有的資訊來應用.
Mark Sir你說得很對.
突然間你這樣說的時候.
令我想起一件事.

$^{241}$我想我們現代人面對這些比喻的時候.
我們有一個距離.
我也是從關係入手.
距離但是根據你剛才說的東西.
我們有一個很大的距離.
因為我們聽到僕人的故事.
我們會想起我們那些僕人.
我沒有僕人.
就是想起我們長大的飛龍,印蓉.
我們想起這些東西.
但是當時的僕人不是我們那些飛龍,印蓉.
當時的僕人是住在主人家.
根本是跟他一起生活.
根本是當他是一家人.
Exactly.
所以當主人和僕人.
這樣分類錢財的時候.
我們可能會說.
搞錯了,主人給他一袋金,五袋金.
好像老鄰居他或者什麼.
其實完全沒有這個意思.
主人是很清楚他們每一個人的性格.
每一個人的才剛.
這樣給他們讓他們用.
所以應該將這些印蓉,飛蓉的感覺.
全部拿走.
反而代入第一世紀.
這樣看是差不多.
我想其實也好像.
之前我們幾集也有說過.
我們很多時候將重點放在不標準的細節上.
但是錯過了大圖.
還有一些應該要看的細節.
例如是一些舊約我們不太看得懂的經文.
或者是重點重複的字眼.
但是其實更加要從上文下理.
看得到那個大方向.
從而理解得到究竟那個比喻真正的意思是什麼.
是的,我想用這個做一個例子.
我想聽得最多.

$^{281}$應該也未必是你剛才問的問題.
可能也有人問過這個問題.
但是我想可能聽得最多的是.
你有多少的屬靈恩賜.
因為說的是這個才華.
你有什麼恩賜你就要用盡它.
某程度上你可以有些這樣拉到這個位置.
但是更加重要的就好像你這樣說.
其實那個核心重點是在說什麼呢.
有時候都是資治到你拆開了就回不來了.
剛才我們看上文下理連上一集.
其實是在說24,25章.
這個耶穌末日末世的時候.
有一天祂要再回來.
究竟天國是怎樣呢.
其實如果你聽了這兩集你大約會明白.
第一就是要準備好.
第二就是在等候的時候你不知道.
所以你要準備好.
上一集我們說到有智慧.
智慧其實也不是純粹街邊的智慧.
而是真的對天國來臨的認識.
而去改變你的生活.
這個很重要.
不是沒有關係,不是到時去天堂才算.
是今天就已經有關.
因為天國今天已經臨在我們當中.
來到這個比喻其實是在說當我們知道上帝的主權的時候.
其實是同樣加了一個責任給我們.
每一個有份在天國裡面作為天國的子民有其應有的責任.
我想這件事其實是很多時候我們都.
如果只是像上一次那樣說.
只是想著天堂.
信耶穌我可以永生.
或者香港那時候說得水牛.
其實就是這樣,只是想著將來.
其實是說今天你的責任.
你被分配的位份.
你應該是怎樣去做好你的本份.
因為如果你記得剛才約略提過.

$^{321}$馬太福音十三章其實和這個有很大的呼應.
十三章說的就是一個種子.
就像上帝的說話.
他來到我們作為一個土壤.
我們來到選擇怎樣去聆聽.
當我們選擇了認真聆聽.
剛才也提過.
那個聆聽其實是聽完而行出來的.
當我們這樣的時候.
那個種子,因為我們的生命就好像變成一個種子一樣.
我們就是負起這樣的責任.
我們被信託了這個位份.
我們就帶著這個使命來去活.
這裡其實正正就是說.
你有這個位份.
你代表著天國.
你怎樣來去活出來呢.
收起來還是按著上帝所給你的能力來去活出來.
這更加重要.
所以其實這個天國.
不只是將來的事.
其實今天就在我們生命當中.
這很重要.
這個天國在我們今天的生命當中.
其實就是一個我們已經領受了上帝給我們的國度.
我們怎樣去按著看到將來要發生的事來去活.
而你今天的選擇就影響到將來你是得著獎賞.
還是要出去外面.
跟主人沒有關係.
這個是今天的選擇.
不是到時我選擇了.
你的生命其實是要重新再生.
Selwyn你說得很對.
這個天國其實是一個現在的現實.
應該是進入了我們今天.
影響我們今天以後怎樣生活.
所以有時聽到一些人說.
我退休了或者有空才去事奉.
我才去傳道.
做天國的事.

$^{361}$現在暫時沒有空.
如果正正是這樣想的時候.
就正正是這個比喻所針對.
你根本就是不知道什麼是天國.
你跟天國是沒有分別的.
如果是這樣的時候.
你就變成一個沒用的僕人.
上帝用不了你.
希望這個系列不單止幫到大家.
去接觸比喻的時候.
應該有什麼技巧或者有什麼工具可以用到之餘.
真的可以認清到天國的現實是怎樣.
神對我們的期望又是怎樣.
並且我們真的要有智慧來去警醒.
因為真的不知道那天會什麼時候來.
也不是用自己的智慧或者價值去計算.
由今天開始回應神.
請訂閱魔人SAVI的頻道及按旁邊那個鈴鐺.
\newpage



\section{7:1-5}
\label{sec:jf7dAWwLNe0}
\textbf{Kingdom Values XLII: Matthew 7\_1-5}
\newline
\newline
連結: \href{https://youtube.com/watch?v=jf7dAWwLNe0}{\texttt{https://youtube.com/watch?v=jf7dAWwLNe0}} ~~~~ 語音日期: 2021-03-01
\newline
\newline
\hyperref[sec:bJj9ij22kAs]{\small{< < < PREV SERMON < < <}}
~
\hyperref[sec:index_chronic]{\small{[返順時目]}}
~
\hyperref[sec:index_scriptual]{\small{[返順卷目]}}
~
\hyperref[sec:QRUeCLAvrYk]{\small{> > > NEXT SERMON > > >}}
\newline
\newline
$^{1}$Hey family, it's great to be with you again. I hope that you had a good week.
I hope that you're encouraged and experiencing God's blessing in your life this week.
I'm praying for you. And as always I want to encourage you to reach out to us.
You can send us an email at info urban voice church comm we'd love to hear from you.
Any comments or questions you might have or any way that we can pray for you and this week?.
We're back in Matthew chapter 7 looking at this section where Jesus is talking about judgment.
And we're talking about how important it is to not have a judgmental attitude.
That's what we're really going to focus on this week how to avoid being.
Judgmental and Jesus is giving us a lot of truth in these verses about how to avoid being judgmental.
So let's look at these verses together. We're gonna start in verse 1 and read down.
To verse 5 this is Matthew chapter 7.
It says this judge not that you be not judged for with the judgment you pronounce you will be judged and with the measure you use.
It will be measured to you. Why do you see the speck that is in your brother's eye?.
But do not notice the log that is in your own eye or how can you say to your brother?.
Let me take the speck out of your eye when there is a log in your own eye you hypocrite.
First take the log out of your own eye, and then you will see clearly to take the speck out of your brother's eye.
Now last week we kind of did an intro we started to talk about what?.
Judgment is and what biblical discernment is and we talked about the fact that judgment is really.
comparing something to a standard and for us as believers.
That standard is God's truth that we we know from God's Word and so for us to exercise discernment.
To make judgments about things we compare them to God's truth and to God's Word.
And we also talked about the fact that that's something that is important to our lives in a lot of ways.
We talked about the fact that judgment has a very negative connotation.
and.
That's okay. It has a negative connotation for good reason it has a negative connotation because.
People and often people in the church have been very judgmental of others.
And so that has caused people to have this reaction to it and to say that's something bad.
But even though it has a negative connotation.
We it's still something that we have to do we have to exercise discernment.
We have to exercise judgment in order to function in order to live life in order to make decisions in order to have.
Relationships with other people and in order to know what's right and wrong and what's healthy for us, and what's not healthy.
So we talked about that last week all of those things and this week what I want us to start to do is to start.
To see how we can avoid having a judgmental attitude. That's really I believe what Jesus is saying in this verse in verse 1.
He's not saying that we shouldn't judge.
He's saying that we should not have a judgmental attitude.
And he's laying out a lot of ways that we need to avoid having a judgmental attitude, so I want us to begin to.
Examine that today.
The first thing I want us to see about a judgmental attitude. That's unhealthy is that a judgmental attitude is.
Graceless it doesn't see God's grace in my life or God's grace in the life of others a judgmental attitude is.
Graceless you know verse 2 Jesus says for with the judgment you pronounce you will be judged and with the measure you use.

$^{41}$It will be measured to you.
Now what is Jesus saying in that verse is he's saying that you know however you judge someone else is how God is going to.
Judge you if you judge harshly God will judge you harshly.
I don't think that's really what Jesus is saying because we know biblically that God judges us all according to the same.
Standard his standard of justice and righteousness and that by the way none of us could ever measure up to it.
That's only because of Jesus that any of us can be accepted in God's sight God is not a biased judge.
He judges us all by the same standard.
So I don't think this this verse is saying that God's standard will change depending upon how we treat others.
I think this verse is actually really similar to Matthew 5 14.
Where Jesus is talking about forgiveness and in that verse Jesus is not saying that if we don't forgive others God will refuse to forgive us.
But what Jesus is saying is that in light of how much God has forgiven us.
That we should be able to forgive others, and I think Jesus is saying something very similar here in Matthew 7 - he's saying look.
Look at the measure that God has judged you with.
Look at the mercy and grace that God has poured out into your life.
Look at your own need and how much you fell short of God's standard and how much you messed up and how much.
you sinned and how broken you were and.
That's the type of judgment that you need to pronounce on someone else.
You need to be full of grace when you look at someone else's life.
And I think that's a huge way that we can begin to avoid being.
Judgmental in our life and avoid having a judgmental attitude is by practicing awareness of God's grace in our life.
We should not ignore the activity of God's grace in our heart and in our life.
And we should extend grace to others.
Because when we examine our life and say God has forgiven me of so much God has covered.
So much of my mistakes God has healed so much of my brokenness.
God has restored so much of the stuff that I did foolishly or I wasted or I did something completely.
I was headed in the wrong direction completely and God.
Turned me around and restored me and he fixed it or I did something unkind to someone else and I I harmed a relationship.
And God restored it and healed it and God brought forgiveness and all this stuff when we start to see.
God's grace and mercy and how deeply it was poured into our heart and our life.
Then we can begin to extend grace to others and I think that's what Jesus is saying in verse 2 here.
He's saying that we need to be aware of God's grace and extend that to someone else.
And so that's one of the big ways we can avoid being judgmental is by practicing awareness of God's grace.
The next thing I want us to see about a judgmental attitude is that a judgmental attitude is arrogant. It's proud and.
The way that we can really avoid that type of attitude is by practicing.
poverty of spirit we talked about poverty of spirits the first beatitude in Matthew chapter 5 and.
We talked about what it really is. It's an awareness of our spiritual bankruptcy before God.
It's an awareness of our need and when we are aware of our spiritual bankruptcy before God it gives us a deep.
humility and.
Humility is so important humility is probably one of the most defining.
characteristics of a true Christian a true believer a true subject of the kingdom of heaven is that.

$^{81}$That person should be humble and that person should live with a deep humility a deep awareness of their own need for God.
And if we are full of humility.
We will not be having a judgmental attitude towards others.
Even when we begin to see something in someone else's life that is unhealthy.
We will respond to it with a spirit of love and we will respond to it honestly with the spirit of brokenness.
when we see someone else who is who is.
Struggling or who has some need in their life. We should never have a spirit of saying. Oh, I'm better than that person.
Yeah, they're having a problem. I would never have that problem.
If our heart is the heart of Jesus our heart says I am so.
broken and sad over this need in their life and how hurting they are and how.
They've made this mistake and how it's damaged their life and and I'm I'm broken up and and I'm in tears over it.
Not an attitude that says yeah, I knew it. I knew they were gonna mess up or look at me compared to them.
I'm so much better. That is a judgmental and proud and arrogant attitude.
And Jesus is calling us to an attitude of humility that when we make discernments.
we should do that with a deep humility and.
even if we're in a position where we are able to discern something and judge something that someone else cannot see and God allows us to.
See some of his truth in a way that someone else doesn't we still are gonna communicate that.
From a place of humility and love and care for someone else, you know.
another thing about a judgmental attitude is that it's focused on the brokenness of others or the problems of others and.
one of the ways that we can avoid that is by practicing seeing our own brokenness first and.
Jesus makes it so clear for us here in this passage in verses 3 to 5.
He says this why do you see the speck that is in your brother's eye?.
But do not notice the log that is in your own eye or how can you say to your brother?.
Let me take the speck out of your eye when there is the log in your own eye.
You hypocrite first take the log out of your own eye and then you will see clearly to take the speck out of your brother's eye.
Jesus uses this metaphor for us. That's so engaging and beautiful, but it's it's humorous and.
He does that on purpose. He says look this is what you look like.
When you go out with a proud self-righteous attitude and you don't see your own brokenness.
You don't see your own need and you start trying to help someone else.
You look like somebody who has a giant.
Splinter a log size splinter in your eye and you're looking at someone else saying hey.
there's a little speck of dirt in your eye, let me brush it out for you and.
Everybody in the room is looking right at you and going. Hey, do you not see there's a log in your eye?.
Do you not get this and that's what Jesus is describing for us?.
He's saying stop looking at the brokenness of other people first stop seeing their problems first and stop thinking.
Oh, this is all them and I'm here to fix it because I'm such a great person.
Jesus is saying no, that is not a kingdom value. The attitude is wrong. It's proud and arrogant.
Jesus is saying see your own need first.
You know one of the things that I've seen in my life.
Maybe you've seen this in your life.

$^{121}$Is it sometimes the the things that irritate us the most about other people?.
Or the situations that irritate us the most or the personality types that irritate us the most.
One of the reasons they irritate us so much is.
That that's who we are.
sometimes the people that we clash with the most are the people that are most like us and.
There's something about seeing our own sin and our own.
weakness and our own.
Idiosyncrasies.
Reflected in someone else that just drives us crazy and we want to call it out and we want to say hey stop being arrogant.
Because we're arrogant and we want to say hey stop being unkind because we're unkind.
And we want to say hey stop being unjust because we're unjust we want to say hey stop favoring other people.
Because we are favoring other people.
You know a lot of times that is a big red flag to us.
When we find ourselves being really irritated at someone else.
It's a big red flag that that.
Same issue is probably something in our life and we should use it as a red flag to say wait.
I need to go back and examine my heart. Why is this bugging me so much?.
Why is it bugging me so much that I see this person at work who's?.
Unsuccessful and and I don't feel like they deserve it and it's just making me so mad.
Maybe it's because.
I've taken some shortcuts and I've done some things to try to.
Put myself in a position that I didn't deserve and so when I see someone else behaving in that way.
Then it just makes me angry. Why is it that when I see someone else being manipulative?.
It makes me so upset and I say they've got to stop being so manipulative.
Maybe it's because I'm manipulative in some of my relationships. And so I know how to recognize it in someone else.
And so when I see it in them, I know exactly what's going on.
And that's why I'm so frustrated and angry and I I lash out at them so much.
oftentimes the brokenness of someone else is reflecting our own brokenness and.
We want to try to go fix their brokenness and Jesus is saying no you examine your heart first.
You let God deal with your brokenness first. You let God deal with your sin first.
You take your need to God and then after God heals you and restores you and helps you.
You'll actually be able to help someone else because first of all now you can actually see clearly. I love the.
Illustration Jesus gives us here of stuff in your eye because obviously if something's in your eye, you can't see clearly.
so how are you gonna help someone else when you're half blind because there's stuff in your eye and.
That really helps us to apply what Jesus is saying here first.
We've got to let God deal with our brokenness and our sin and then after that maybe we can go help someone else.
Because we can see clearly.
But also now we have a testimony if God has worked on your brokenness and helped you to to be healed and to be restored.
Now you can see clearly.
What how you can help someone else, but you also have a testimony to share with someone else and you can say yeah.

$^{161}$You know what?.
I'm just as broken as you are and I used to struggle with that exact same thing.
And this is how Jesus brought healing and restoration of my life.
And this is how Jesus helped me and you have an attitude of grace and humility.
And true concern for that person and what's going on in their life and you actually have a testimony something.
That's helpful that you can share with them instead of just saying hey, you're wrong and I don't like it.
I don't like what you're doing, right?.
That's a judgmental attitude and Jesus is.
Calling us to examine our own hearts first and we have got to practice that in the church.
We've got to look at our own need first before we try to fix.
Everybody else around us and another thing that we see about a judgmental attitude is a judgmental attitude is.
Misjudging this is huge. This is so important for us to understand as humans.
We look on the outward appearance.
That's what we see and that's all we can see and we misjudge a lot of things based on the outward appearance.
And we've got to stop doing that, you know in first Samuel 16 7.
this is when the Prophet Samuel has gone to anoint the next king of Israel and.
God tells him to go that he's gonna be one of the sons of Jesse.
And so it's Jesse's son start coming out and the Prophet Samuel's looking at Jesse's sons.
And he sees this guy named Eliab who apparently was you know.
big and strong and handsome and everything that Samuel thought a king should be and.
Samuel says oh, this is probably the guy I should go ahead and anoint him because you know, this is King material right here.
and in verse 7.
God says this.
But the Lord said to Samuel do not look on his appearance or on the height of his stature.
Because I have rejected him for the Lord sees not as man sees.
Man looks on the outward appearance, but the Lord looks on the heart.
God is the only one who can see our hearts and see our thoughts and see our intentions and see our motivations.
God can see that the Bible is clear about that Genesis 6 5 says the Lord saw.
The wickedness of man was great in the earth and that every intention of the thoughts of his heart was only evil continually.
that verse is before God's judgment in the days of Noah and.
That verse tells us that God can see our thoughts and God can see our intention and God knows if it's evil.
Or if it's wicked God knows what's in our heart. He knows what's in our head.
He knows what's on the inside of us. The Bible is clear about that.
But here's the thing we as humans do not and that's one of our biggest.
Problems in judging other people. That's one of the ways that we get ourselves into the most trouble is.
We begin to think that we are able to judge someone else's intent.
Someone else's motivation someone else's thoughts why they did something and answer all these questions that we cannot answer.
When Jesus talks about exercising right judgment.
The only thing we can do as humans is compare someone else's actions to God's standard and say was their action.
Right, was it healthy we can do that.

$^{201}$Okay, but we do not have any grounds to.
Say why they did it or to judge their intents or their motivations or their thoughts.
I don't believe you can find any place in scripture to support.
The view that we as humans are able to judge the intents and thoughts and motivations of others.
Only God can do that and yet we try to do that all the time.
We judge intent all the time. You hear us as humans saying. Oh, well, he thought this. Oh, well, she thought that oh.
he did this because and.
Then we list a motivation or we list an emotion or we list a thought or we list an intention.
we have no idea if they did that because.
We say oh well this happened because this person felt this way. We don't know how they feel.
We say this happened because they were afraid.
Well, maybe so but we also it could have also been because they were excited. We have no idea.
What is going on inside the heart of another person?.
We can discern and make judgment decisions based on someone's actions. We can look at someone else's life and say.
We can compare it to the standard of God's truth and say okay that action was not healthy.
God has said that that's not a healthy action for us to undertake.
God has said that if we undertake that action it's gonna harm us. It's gonna harm the people around us.
We can look at someone's action and know if it was healthy or not.
But we do not know what's in their heart and we have got to stop doing that.
that's such an unhealthy thing for us to do and I would say a great way for us to.
Avoid that type of judgmental attitude is to to practice banning ourselves.
from judging intent and motivation and thoughts.
When you find yourself observing someone else's actions observing someone else's behavior.
Whether it's in life you you actually see someone do something that's wrong or are unhealthy or unkind to someone else.
Or maybe it's you know on social media or something whenever you.
Observe someone in some way and you start to think oh.
They were thinking this. Oh, they said that because oh they did that because oh, I know why that happened.
They intended to do this.
Oh.
They were being manipulative and they said this to make me think that when you start to go through those mental.
Gymnastics and you start to tell yourself that you know.
What people are thinking or what they intended or what their motivations were or what their thoughts were?.
You need to say wait, I I can't do that. I don't actually know. I don't know why they said that.
I don't know why they did that. I don't know why they reacted that way.
Sometimes we just assume well this person reacted negatively to me because they they don't like me.
maybe they reacted negatively to to me because.
They were having a terrible day and there was all these things that happened in their day that I didn't observe and had no idea.
About or maybe they totally misunderstood what I said.
Maybe they thought I said something that that I didn't even say I remember a long time ago when I was in college.
I had a bad ear infection and I happened to be.

$^{241}$leading this Bible study and there was quite a few young people there was like high school students and.
I was leading and we were in a room and I was having people go around and introduce.
Themselves and we were sitting in a big circle. There were some adults in the circle as well.
and so the young people were introducing themselves and.
answering some icebreaker questions and we got like halfway around the circle and.
One of the young people said kind of quietly.
Why don't you let the adults introduce themselves because there were some adults in the circle and I thought.
that this young person said.
there's some people outside who need your help and so I got up and.
Walked out of the room left the circle walked out the door and went outside.
I was outside looking around for a while couldn't find anybody then I came back into the room.
Everybody was staring at me like I was an alien and I said, yeah, there's no one out there.
I couldn't find anybody and.
Everybody started laughing and people are like, what are you talking about?.
And so I was like, I don't know and then they explained to me that this young person had said.
Why don't you let the adults introduce themselves? I.
undertook a really strange action.
Because I totally misunderstood what this person had said and that happens in life all the time.
things happen because of.
misunderstandings things happen because of.
fear or frustration or hurt things people do things because of something they've gone through in the past in a.
Situation our relationship reminds them of that past hurt and they respond in a way that might be uncharacteristic.
We have got to stop always thinking that we know.
Why people are doing things or we know what they intend or we know what their thoughts are.
We know what motivates them because we don't.
God does and Jesus calls us to right judgment.
It's not that we can't ever look at someone's action and say that that's that's wrong or that's unhealthier.
That's not good for them.
We can and we actually have to do that if we love people in our lives. There are moments where we have to say.
That action is unhealthy. It's harming them. It's harming the people around them and I care about them.
So I'm gonna be honest about that. We have to do that in life.
But even in those moments we do not know what their thoughts and intents are and we shouldn't act like we do.
We've got to learn.
To ban ourselves from doing that when we find ourselves doing that we have to say hey.
I need to stop because I don't know what they're thinking. I don't know why they did that.
I I don't have that type of wisdom only God knows that and so I'm gonna focus on helping them to.
Respond to God and to make right decisions. I'm gonna focus on loving them.
I'm gonna focus on caring for them, but I'm not gonna attribute thoughts and intents and motivations to them because I don't know what they are.
Unless those are expressed to me in some way.
I have no idea what is in someone's heart and we've got to live by that to avoid a judgmental attitude.

$^{281}$Now there's a lot more aspects of avoiding a judgmental attitude that we're going to talk about next week.
But those are some big ones for us to practice this week.
We need to practice awareness of God's grace and extend that to others.
We need to practice poverty of spirit and humility so that we're not arrogant and proud.
We need to practice seeing our own brokenness first like Jesus taught us and we need to practice banning ourselves.
from judging.
Intent and motivation and thoughts and things like that that we just don't know and that we're not able to know.
We don't have the ability to know those things.
So let's work on these things this week in our life as we encounter them as we exercise judgment and make discernments.
All of us have to do it all the time every day throughout the week. And so as we do that this week.
Let's really start.
allowing Jesus to work on our heart so that we will not have a judgmental attitude so that we'll be kind and loving and caring and.
Operate from a place of love as we interact with other people and even when we see some someone doing something unhealthy.
That's something that God has told us is not a behavior for us to engage in we can have a spirit of love and kindness.
as we as we interact with that and as we encourage our friends and our loved ones and the people around us towards.
Right behaviors. We need to have a spirit of love and care for others at all times and a spirit of humility a.
Spirit of seeing our own need for God a spirit of grace. Let's pray.
God, I just pray that you would give us an awareness of your grace and how much you've extended grace to us and.
That we would live from that place of knowing how much mercy and grace you poured into our lives and we would extend it to others.
God, I pray you'd give us poverty of spirit.
Help us to see our own need for you and help us to have meekness that comes out of that.
because we see our need for you or we're able to operate from a place of humility and.
And to see our own brokenness first God.
I pray you'd help us to see our brokenness before we look at the brokenness of others help us to examine our own heart first.
God and God help us not to misjudge others because we only see the outward appearance.
Help us not to think we know what's going on in the heart when we don't and help us to be honest about that.
God.
help us to love and care for the people around us and help us to even when we need to exercise discernment and judgment and.
even when we need to.
To help someone else to move towards you and to move towards right relationship to you God help us to do that.
In a way that comes.
From love and in something that is rooted in your love and your love for us and your love for others God help us to always.
Operate from a place of love God and even as we exercise discernment God help us to grow in this right now this week.
help us to grow in.
avoiding a judgmental attitude and having.
Discernment and right judgment that honors you God in Jesus name. Amen.
[BLANK AUDIO].
\newpage



\section{}
\label{sec:QRUeCLAvrYk}
\textbf{舊約網上行 - 上帝的足跡:EP4 方舟}
\newline
\newline
連結: \href{https://youtube.com/watch?v=QRUeCLAvrYk}{\texttt{https://youtube.com/watch?v=QRUeCLAvrYk}} ~~~~ 語音日期: 2021-03-02
\newline
\newline
\hyperref[sec:jf7dAWwLNe0]{\small{< < < PREV SERMON < < <}}
~
\hyperref[sec:index_chronic]{\small{[返順時目]}}
~
\hyperref[sec:index_scriptual]{\small{[返順卷目]}}
~
\hyperref[sec:rLZ_I4dgAlo]{\small{> > > NEXT SERMON > > >}}
\newline
\newline
$^{1}$如果你要找一個全世界最受歡迎的聖經故事.
相信一定會想起羅亞方舟這個故事.
當小朋友還小的時候.
曾經去買過一些兒童用品.
總會找到羅亞方舟這個水上的動物園可愛的設計.
並且這些設計裡面一定有一隻動物.
就是長頸鹿.
羅亞方舟Lois Ark.
相信是一個最Iconic的聖經故事.
今次舊約網上行第四集.
我們一起來研究創世記1至11章.
是整本聖經的導言.
中間6至8章就是交代方舟的故事.
羅亞洪水的技術比其他事跡更加長.
可見它有多重要.
我們提過創世記1至11章有一個很重要的目的.
就是要和當時美索不達米亞近東的神話思潮做爭辯.
到底當時世界的人如何去想紅水這個災難.
聖經提出什麼辯駁.
當中有什麼特別的神學信息.
歡迎你跟我們舊約網上行.
在美索不達米亞這個世界.
河水氾濫不是罕有的事.
今天的南美索不達米亞.
伊拉克南部和波斯灣是三河一片.
印象中是不會有河水氾濫的事.
很大程度是因為在過去幾十年.
伊拉克,土耳其,敘利亞這三個國家.
在上游經驗了大約有二十多個大大小小的水壩.
舉個例子.
2018年土耳其剛剛啟用了Elisu Dam大水壩.
引起了國際關注.
因為大大阻截了Tigris River下流的水流.
造成了很多生態上的破壞.
但是自古以來.
這個地方兩河流域.
是常有河水氾濫這件事.
紅水的季節是三月至八月之間.
三月春天會暖.
北面的高山會積雪溶解.

$^{41}$春雨連長的時候.
河水就會上漲.
五月的水位很多時候是最高.
一直到七八月才會減少.
而美索不達米亞這個盆地.
它的傾斜度是非常低的.
非常平坦.
紅水的季節的時候.
可以水浸的地方是整個地帶的三分之二.
舉個例子1954年的時候.
春天暴雨連長.
加上高山的積雪溶解.
就形成了很大規模的水浸.
歷世歷代.
其實有很多古代的文獻.
都提過很多次的大水災.
人類的文明.
創造之一.
就是這些君王.
說是如何去做一些防水的提拔.
例如有一個君王叫漢莫拉比.
Hammurabi.
他是一個舊巴比倫帝國的君王.
在他的黑文上這麼寫.
他說我漢莫拉比是大能的王.
是巴比倫的王.
是四方的王.
我建造這個土牆.
讓我的神.
諸神歡喜.
就用土堆.
來抬起在Sipir的城牆的頭.
這是一個防洪牆的建造.
所以在這一帶的地方.
經常都會找到一些紅水沉積的土層.
例如這個列表.
就讓大家看到.
例如在基斯城邦裡面.
大概在公元前2800年.
間間斷斷有一些紅水造成的土層.

$^{81}$在Swarupa.
在Rurik也有.
但是很多時候同一時期.
未必會在其他的城市找到.
無論如何.
這裡告訴我們.
這些紅水.
並非一些罕有的事.
但是如果你翻閱一些近東的起源史.
遠古的歷史記載.
講到紅水的時候.
包括創世記的羅亞故事.
似乎記憶著一個古時.
一場極大的水災.
或者可以說是一種千年一次.
甚至五千年一次的大紅水.
淹沒了當時所知的世界.
造成很大的破壞.
這些大的紅水的事跡.
似乎是流傳下來.
以至世世代代有一個很深刻的記憶.
慢慢就成為寫作.
在泥板出現.
在文件裡面出現.
所以為什麼近東的遠古起源史的時候.
往往就會提到一個大的紅水.
近東的紅水故事.
譬如我們所熟悉的.
Gilgamesh Epic.
基加米斯神話.
裡面也提到過一個紅水故事.
有什麼元素在裡面呢?.
第一就是有方舟.
方舟建造.
並且將眾生帶上方舟.
當然近東的世界也會帶上什麼呢?.
就是一些文明的創建.
譬如水手.
手工藝.
不同的技術帶上方舟.

$^{121}$也有水象.
水退的描述.
譬如水退的描述.
Gilgamesh Epic.
就提到船上的紅水英雄.
放著三次.
去看看水退了沒有.
並且會有下船的記載.
就是紅水生還的英雄會獻祭.
等等這些元素你會發現.
跟創世記紅水的故事.
有很多相似的地方.
當你看到這些比對的時候.
我們立刻引起一個現代人經常問的問題.
就是有沒有抄襲的可疑.
就是說聖經是不是抄近東的紅水故事.
因為這麼相似的緣故.
現代人一提到抄襲.
我們假設就是偷了別人的資料.
因為我們現代人有一個預設.
就是知識產權.
但是我們這樣問的時候.
其實是問錯問題了.
因為在古代近東的世界裡面.
沒有這種我們現代人的概念.
極少的著作.
是會有自己署名的.
不同的民族裡面的 scribes.
他們的文士.
都會抄寫並且保存著.
上古之年.
以往祖宗帶下來的文學著作.
所以他們抄寫.
並且有時候會刪改.
並且跟這些文學對話.
是一種保留神聖文學的傳統.
所以即使聖經.
是跟這些神話,禮版相似.
也不足為奇.
所以我們要去決定.

$^{161}$說誰先寫.
誰寫後.
後的抄前的.
其實這些問題.
對於古人來說.
是完全沒有意思的.
創世記怎麼去回應當時近東洪水的神話呢?.
怎麼在這些思潮當中.
提出另外一種嶄新的神學.
完全不同的一種世界觀.
我們看看近東的神話洪水的記載.
譬如Gilgamesh Epic.
第十一塊的禮版.
就描述大洪水的故事.
透露了為什麼會有洪水.
原因是什麼.
加上其他近東的洪水故事.
一致的原因都是說.
因為當時人口增加.
有很多的草蔭.
就擾眾神清夢.
睡不著的時候.
眾神就決定將人類殲滅.
想以根清靜.
其中一位神靈.
就告密於一位犯人.
就叫他造一隻方舟.
結果大洪水真的來了.
連眾諸神都怕了.
然後他們就後悔了.
因為沒有人來獻祭.
眾神就開始餓了.
誰知道呢.
那個生還了的.
所謂洪水的英雄.
就在方舟上.
他沒有死.
然後洪水過了.
他就下船.
去到獻祭.

$^{201}$眾神聞到香味的時候.
就決定.
不再降大洪水了.
你看到這些近東的神學.
怎麼去想神的時候.
就好像神靈跟犯人一樣.
睡不著就會發脾氣.
然後做錯事.
有我們犯人的需要.
看到洪水也會害怕的.
大水混沌來到的時候就搞不定了.
聖經的經文.
就是提出爭辯.
就是想清楚神.
一個值得我們敬拜的神.
是一個怎樣的神.
為什麼會有洪水.
在聖經裡面講到洪水的原因.
就不是噪音.
而是人的罪.
六章第一至第八節.
我們選這些經文來看.
他說人在地上開始增多的時候生兒養女.
神的眾子看見人的女子未來.
就隨意的跳船取作為妻.
耀號說人基屬血氣的.
我的靈就不永遠住在他裡面.
但是他的日子還有一百二十年.
在那時候的日子.
以後的日子.
有巨人在地上.
神的兒子跟人的女子結合.
就生出上古英武有名的人物.
耀號說看見人在地上罪惡很大.
終日所思想的盡都是邪惡.
於是耀號說就後悔做人在地上.
心裡面憂傷.
耀號說我要把我所創造的人.
地上來消滅.
無論人大地的一切都要消滅.

$^{241}$因為我後悔做了他們.
只是羅漢在耀號眼前蒙一.
這段聖經怎麼講都好.
都是其中一段最難解的經文.
有些人說神的兒子.
神的眾子與女子交合.
就好像天使跟凡人交合一樣.
這些沒錯.
後期.
到了特別兩約中間.
猶太教的時候.
他們在詮釋這段經文的時候.
就發展了一些魔鬼論.
或者天使論等等.
但是相信對第一代的近東讀者來說.
他們聽見神的兒子的時候.
就好像我們華人傳統的文化一樣.
就會想起當時的君王.
天子.
這些人類文明的創立者.
繼承者.
上一集我們講過.
這些人有多麼重要.
叫做Abkhazia.
這些君王.
他們建立人類的文明社會.
但是在創世記眼中.
這些文明.
這些君王.
就是敗壞了神美好的世界.
怎麼解釋都好.
這個創世記的作者.
似乎是用一些特別的字眼.
叫我們想起創世記第三章.
我們比對一下.
第三章第六節.
人的第一次叛變.
他說女人見那火樹上的果子好作食物.
就越人眼目.
使人有智慧.

$^{281}$就摘下果子來吃了.
六章第二節.
講到洪水發生的原因.
人的罪惡的時候.
都是用幾個同樣的字.
神的眾子看見.
人的女子美好了.
美麗了.
就隨意挑選.
如果你看到他們運用一些同樣的字眼.
就是想表述.
洪水氾濫之前.
就是另一個人類的大叛變.
不守本位任意作主.
去挑戰上帝的權威.
跟第一個的叛變都是一樣.
你再看看創世記第一章.
到第六章.
脈絡的時候.
就講到人類罪惡的氾濫.
你可以看到罪惡破壞的範圍越來越大.
由開始第一次.
阿當,夏娃是個人被破壞.
然後到第四章.
就是該人殺他的兄弟.
是家庭被破壞.
關係被破壞.
然後到第六章我們剛剛看的經文.
就是神的眾子的君王.
就講到社會也被破壞.
甚至是全地每個人都敗壞.
是一種 progression(進步).
就是罪惡的氾濫越來越大.
但是神是怎麼回應呢?.
你看看第六章.
五到八節.
這裡有一個很清晰的結構.
當時古代的作者.
沒有一些 highlight(亮點).
或者 boldface(複雜).

$^{321}$沒有 italic(詳細).
去突出重點.
他們用一些關鍵詞.
重複詞或者結構.
這幾節經文.
有什麼結構呢?.
你就看到有一種交叉.
並聯的結構.
又說看見人的罪惡.
然後又說後悔傷心.
然後呢.
心裡面憂傷.
然後他就說.
我要將他們審判.
然後就對稱.
又再講.
神後悔傷心做了他們.
然後最後的時候.
他就看見羅亞.
他就蒙恩.
這一段又是另一段很難解的經文.
但是對我們的神學是很重要的.
因為很多時候我們覺得難解.
是因為經文裡面一些的翻譯上.
造成一些誤會.
這裡講到後悔.
對我們現代人來講.
後悔就好像做錯事一樣.
英文有其他的翻譯是.
grieve(痛).
或者 sorry(抱歉).
其實希伯來的原文就是.
Lame(雷蒙).
這個字是怎麼用的呢?.
簡單地去陳述.
你可以找回這個字同樣的用法.
在其他的經文裡面是怎麼用的.
其實最清楚的一節經文.
就是《史詩》的二十一章第十五節.
那裡上下文講到.

$^{361}$以色列人內戰.
變雅敏支派的人死了很多.
於是乎人口就越來越少.
幾乎要滅族了.
他們就講了一句這樣的話.
眾人為變雅敏支派難過.
Lame(雷蒙).
因為耶和華在以色列眾支派裡面.
有了缺口.
是差了.
平衡不回來.
所以有一些的色經家認為.
這個Lame(雷蒙).
我們中文翻譯做後悔的時候.
或者英文翻譯做 sorry(抱歉)的時候.
不是講到做錯事.
想起然後後悔.
而是一種 accounting(資料搜集).
核數 auditing(核數評估)的字眼.
就是說.
這個字是一種用語.
來表述一種在創造界裡面.
一種的 check and balance(檢查和平衡).
一種的核數.
就是人類最大惡極的時候.
就平衡到世界的秩序.
就有所缺失.
怎麼平衡呢?.
在聖經的舊約的神學裡面.
就有兩個方法來平衡.
第一個方法就是因憐上帝去赦罪.
或者是公義的審判.
這裡講到洪水的時候.
就講到上帝用公義來審判有罪的世界.
創造主是要對創造界來問責.
這個在近東的世界裡面.
也是一種很特別的世界觀.
這個世界是由創造主獨一的上帝.
來去有 check and balance(檢查和平衡).
並且這裡提出了神的聖情.

$^{401}$講到憂傷.
傷心欲絕.
就是肝腸寸斷.
你想起洪水的警告的時候.
我們好像經常想起.
就是世界有危難會來.
有 crisis(危機).
但在聖經裡面.
你詳細去看的時候.
真正有 crisis(危機)的.
是上帝自己.
耶和華傷心欲絕.
因為所做的人罪惡滿盈.
上帝感到傷心.
但這裡就是聖經提出一個截然不同的神觀.
theology(道德).
另外一個世界觀.
就是人如何去面對這位上帝.
神是怎樣的呢?.
這裡講出神心裡面的感受.
祂的本性.
神對人的要求當然是很清楚.
是會向人類作道德上的問責.
但同時.
人的罪.
人破壞這個世界.
上帝的心情是怎樣的?.
祂是會傷心.
會肝腸寸斷.
快要到受難時.
這起不來想起神對人的愛.
神對人的恩典.
一個完美的彰顯.
耶穌在十字架上.
來為我們受苦.
人應該怎樣面對這位創造主.
這位上帝.
是否應該在祂的恩典裡面.
來迴轉.
來與祂建立關係.

$^{441}$有一首英文詩.
寫得很好.
不是很多人知道.
叫做《神的眼淚》.
Tears of God.
如果有興趣的話.
可以在下面看看.
講到人的敗壞.
這個世界的墮落.
上帝到底心裡面在想什麼?.
祂是流眼淚.
現代人看方舟的故事.
往往重點是搞錯了.
我們很多時候會問一些物理.
科學等等的問題.
方舟有多大?.
是不是真的可以容納全世界的動物.
每一個物種在方舟裡面.
怎麼可以一年在洪水裡面餵飽這麼多動物呢?.
怎麼處置很多物理的問題呢?.
我們就思想一大堆這些問題.
然後就想出一大堆護轎的答案.
譬如我們想起方舟.
它的長度.
差不多等於Titanic的一半有多.
現代人自然會有一些爭論.
這麼大的木船.
排水量四萬三千噸.
一定會在洪水裡面破損的.
其實呢.
會不會就是我們現代人問錯了問題?.
好像我們問創世記六日的創造一樣.
是不是真的六個二十四小時.
或者創世記的五張家譜.
那些很大的壽數一樣.
那些數字到底是不是按照我們現代人的字面意思.
還是我們應該回到古近東的時候的背景來瞭解.
講到方舟的尺碼.
這些近東的洪水故事.
譬如Gilgamesh Epic.

$^{481}$裡面的洪水英雄生還者叫做Unapishtim.
他做的方舟.
長闊高是一樣的.
是一百二十爪.
你算出來是等於今天的一百八十尺或者五十五米左右.
排水量其實比聖經裡面的羅亞方舟大三至四倍.
並且那個方舟是七層高的.
形狀可以是正立方體.
或者是當時近東的塔廟的sigra.
有這麼多層.
所以是比羅亞方舟更加不切實際.
1985年的時候.
又發生了一件事情.
在大英博物館發生.
就有一個人拿著一盒東西.
是一些古物.
就帶到大英博物館的館長叫做Irvin Finkel.
Irvin Finkel去檢驗這些古物的時候.
就找到一塊泥板.
他一看這塊泥板的時候.
他就嚇了一跳.
原來內容是講大洪水的故事.
他甚至運用了很多高科技.
將泥板所有的細節.
所有的details.
將它翻譯出來.
裡面這塊泥板提到做方舟的尺碼.
他就說長度直徑是222尺.
那艘船是圓形的.
好像大半個足球場一樣.
之後他甚至找了一些航海的專家.
特別是古代的航海研究的專家.
就嘗試去照這個尺碼來造船.
他們運算之下.
這些專家說在物理上.
根本這艘船是沒有可能的.
極其量只能夠造到這艘船的五分之一的尺碼.
最後甚至有電影商去資助.
去拍一個紀錄片.
就按照這個新的泥板.

$^{521}$去造這艘圓形的船.
但很可惜.
這個泥板不是Akira 現在的家司的手冊.
沒有詳細講清楚逐步怎樣去造.
但很幸運地.
原來伊拉克的南面.
仍然保留著上一代造圓形的船.
叫做Goovers.
這種手法.
就是傳統用一些草.
去摘掉.
曬乾之後摘成一些菌狀的成熟.
然後一路打圓圈.
就形成一個碗狀的船.
這些還沒有失傳.
然後上面放上一些櫻青.
這些原來就是.
今天在伊拉克南面.
一些民族叫做沼澤.
阿拉伯人Marsh Arabs.
仍然採用這樣去造屋造船的手法.
有一點像香港以前的水上人.
但因為這裡的政局不穩.
所以他們的紀錄片.
沒有在這裡做實驗.
結果就找了印度一個村落.
在那裡找了一些人.
按照阿拉伯人的古法.
造一個碗狀的船.
結果2017年.
這艘船正式下水.
為什麼要講這個故事呢.
我們就想知道.
到底以前古人的例版.
這些故事.
他們記載方舟的尺碼的時候.
很多時候都不是實際的尺碼.
不是我們現代人心想的.
而是他們想起一些神學的意義.
這些尺碼往往有宗教的意味在.

$^{561}$譬如我們剛才看到.
Gilgamesh Epic.
那裡講到洪水的故事.
方舟是一個浮水的Sagara.
是一些巴比倫的塔廟.
是一種聖所.
如果是這樣的話.
我們回頭看創世記.
可能就發現了不同的東西.
羅亞方舟.
是怎樣記載它的建造的呢.
六張十六折.
方舟上面要做天窗.
向上有一爪.
方舟的門要在旁邊開.
方舟是分上中下三層.
你想起上中下三層.
你就想起近東.
希伯來人的宇宙觀.
有天,地和地下.
近東的神廟的記載.
甚至會說這些神廟.
塔廟.
是會有天窗的.
你就會看到.
當時的人想起聖殿.
聖所,廟的時候.
就是連接天上.
上帝寶座的地方.
所以在大水裡面.
羅亞的方舟是什麼呢.
就是一個整個宇宙的.
縮影.
將一些動物.
潔淨與不潔淨.
羅亞一家其實是代表著萬物.
代表著一個新的人類.
進入一個新的世界.
上帝藉著洪水的災難.
重建這個世界.

$^{601}$就由方舟開始.
羅亞方舟.
也是長方形的.
想起長方形的時候.
似乎就指到.
以色列人所熟知的聖所.
會幕是長方形.
聖殿是長方形.
你看看那個描述.
都是與會幕.
與聖殿的建造相似.
比如羅亞建造方舟的時候.
最後一段經文.
讀章22節.
羅亞就這樣行了.
神吩咐他的.
他都照樣作了.
這句話在很多地方出現過.
出亞洲.
出東亞.
很多地方出現過.
出埃及記有好幾傳.
我們只是看一段經文.
40章第16節.
這裡講到會幕建成的時候.
就是摩西就這樣作了.
凡事耶和華吩咐的.
他都作了.
講到建造會幕.
讓上帝臨在.
上帝同行.
建造會幕.
是一個神聖的世界.
等於創世記第一章.
去到建造宇宙的時候.
創立宇宙的時候.
都是一個神聖的聖所.
所以創世記記載.
這一段羅亞方舟的故事的時候.
其實是講到神要藉著水災.

$^{641}$去到更新.
重整這個世界.
羅亞與上帝立約.
就是參與上帝重建的工作.
以前我們曾經講過.
洪水的故事.
神學的中心訊息是什麼.
大水是近東人.
很普遍的理念.
講到創造.
就是將大水分開.
所以聖經這樣記載.
神的靈魂在水面上.
這是創造.
到洪水氾濫的時候.
就是一種反創造的創造.
大水再淹沒.
當時所知的世界.
但講到重創造的時候.
就是神的風.
神的靈.
將水吹乾.
再一次露出來.
大水再一次被分開.
是一種重創造.
伸直著洪水的災難.
去到重整這個世界.
羅亞來參與.
上帝重整世界的工作.
疫情已經有年多的時間.
其實這些災難.
都是上帝重新組織這個世界.
組織我們屬靈的生命.
組織教會.
組織教會的使命.
羅亞可以參與上帝的工作.
我們又如何去參與上帝的工作.
你看羅亞洪水故事的記載.
洪水有一年長的時間.
你越看的時候.

$^{681}$就越覺得奇怪.
裡面沒有講羅亞很多的心情.
他做的事.
那些情節.
我們現代人最想知道的.
反而聖經沒有講.
或者當時近東的洪水故事.
都會有講洪水英雄做了什麼.
但聖經很多時候隻字不提.
所以如果拿羅亞洪水這個故事.
去拍電影的話.
是很難的.
2014年的時候.
Russell Crowe拍了一套電影.
叫做《羅亞》.
加了很多創世記洪水故事.
是沒有講到的情節.
那些想像全部都是取材質.
猶太人後期的一些傳說.
至經的內容.
這樣才可以拍到一齣電影出來.
所以如果你看羅亞洪水故事的時候.
所謂研究羅亞的角色.
幾乎是沒有可能的.
經文著眼所記載的是什麼.
很簡單地說.
羅亞是遵循上帝的說話.
被神去記錄.
不是他做了什麼大事.
有什麼威風.
而是單單對上帝遵從.
遵照上帝的吩咐去奉行.
創世記七章六節.
當洪水在地上氾濫的時候.
羅亞是六百歲.
六百我們曾經講過.
也是巴比倫完整的數字.
近東洪水的故事.
那些英雄甚至會變成神仙.
一個不是羅亞洪水聖經的神學.

$^{721}$羅亞最被上帝紀念的就是.
得到上帝的恩典.
然後他遵循上帝的說話去奉行.
最近我去了一個弟兄的安息禮拜.
我發現很多人去分享.
這位弟兄的安息禮拜.
這位弟兄的短短生命.
如何觸動他們的生活的時候.
異口同聲地說這位弟兄很安靜.
很謙卑.
從來不誇獎自己.
沒有做過什麼轟轟烈烈的事.
很細心地關心身邊的朋友.
關心教會.
叫人查經.
去影響身邊社區的人.
不是很大的事.
但它觸動了很多人的世界.
疫情,災情當中.
可能我們也會有這樣的感覺.
覺得我們有很多事都做不到.
很多重要的大事.
我們想做的,我們辦不到.
我們感覺到很悲動.
但很多時候.
在上帝面前.
最重要的,或唯一一樣最重要的.
就是我們安靜地去遵命.
遵循上帝的指示.
六章第二十二節.
我們用這一節聖經結束.
「墮阿就這樣作了.
神吩咐他的,他也照樣作了」.
求主幫助我們成為一個.
遵照祂說話而行的人.
縱使我們有一些小責任.
我們也遵照神的吩咐.
盡心地做一個忠於祂,敬拜祂的人.
縱使在疫情當中.
我們依然來做一個順命的子民.

$^{761}$參與上帝重建世界的工作.
繼續去事奉祂.
我們低頭祈禱.
因為你是那位對我們說話.
去吩咐我們的神.
你不單是創造主,也是再創造主.
在我們疫情當中.
我們有不同的難處.
有不同的受困.
但我們要知道.
在你面前,在你眼中.
最重要的就是我們默默地.
遵照你的話語而行.
在我們生活中大小的角色.
我們做一個遵從你心意的人.
讓我們參與你在世上重建的工作.
讓我們知道你自己對我們的心意.
聽到你對我們說話的時候.
凡你所吩咐的.
我們也照樣來做.
求主你親自幫助我們.
祈禱奉主名禁.
\newpage



\section{}
\label{sec:rLZ_I4dgAlo}
\textbf{經緯線 - 耶穌的比喻( EP6 研經資料推介)}
\newline
\newline
連結: \href{https://youtube.com/watch?v=rLZ-I4dgAlo}{\texttt{https://youtube.com/watch?v=rLZ-I4dgAlo}} ~~~~ 語音日期: 2021-03-06
\newline
\newline
\hyperref[sec:QRUeCLAvrYk]{\small{< < < PREV SERMON < < <}}
~
\hyperref[sec:index_chronic]{\small{[返順時目]}}
~
\hyperref[sec:index_scriptual]{\small{[返順卷目]}}
~
\hyperref[sec:w_y0M2H_ndU]{\small{> > > NEXT SERMON > > >}}
\newline
\newline
$^{1}$Hello.
經過我們五集《經位善耶穌的比喻》系列後.
希望大家接下來看聖經.
都可以從《上文下理》舊約的背景和文化背景入手.
我自己很喜歡用Cultural Background Study Bible來看.
至少可以很簡單的馬上有轉世告訴我.
當時的背景是怎樣的.
好像少了一些代溝.
和幾千年前在以色列人住的文化.
除了這些之外.
我們在哪裡可以找到更多的工具呢?.
Mark Sir有沒有什麼好介紹的?.
我知道《比喻》的歷史經書有很多.
兩位都會說.
我其中一本想介紹的就是《Short Stories of Jesus》.
是一位女猶太的色經家.
她不相信耶穌.
不過她對猶太人的典籍很清楚.
Amy Drew Levine.
她有一個很重要的點.
就是每次解釋《比喻》的時候.
她說耶穌說的比喻是很 provocative(挑釁性).
是一些很挑釁性的故事.
我覺得整本書的主題就是.
看到比喻是如何的 provocative(挑釁性).
如果不是的話.
你是看不到的.
她的意思就是這樣.
所以我覺得這個主題是挺好的.
逼我們去看.
不要當耶穌的比喻是一些衣索寓言.
教小朋友不要守株待兔.
教他們一些好故事.
好的 moral(道德).
教他們做一個好人.
其實就是用天國的價值來挑釁我們.
挑戰我們.
逼我們活出天國.
有些人當頭棒學很多時候.
我覺得這本書很幫助我們去想這方面.

$^{41}$Jason呢?.
我介紹兩本經典古舊的書給大家.
很多時候我們都會看一些比較新的學者.
他們有些新的研究.
因為不停都會有新的東西發生.
但是我們都不要忘記這些經典.
因為這些經典都給我們一些很好的基礎.
例如其中一本.
我想很多人都認識的.
就是楊木篤木施寫的《亂世風言》.
這本書.
他不是說了聖經所有的比喻.
他選了幾個.
但是我喜歡的就是.
我們可以學到的.
有時候我們都會很 admire(敬佩)一些講道的人.
為什麼他看完故事.
為什麼這麼快可以有一個生活的例子出來呢?.
為什麼可以把生活的例子和聖經拉上關係呢?.
楊木篤木施在這一方面.
我自己覺得他是很有心得的.
你看完之後.
都對由經文到現實生活.
如何連結.
有一個很好的學習.
另外一本經典.
就是Gregg Bromberg寫的《Interpreting the Parables》.
為什麼這也是很好的一本書.
可以幫我們打好基礎呢?.
因為很多時候我們看比喻.
有一個很大的問題.
就是.
究竟這個比喻.
我們其實在系列裡面也有講過.
究竟這個比喻的重點在哪裡呢?.
是一個重點?.
兩個重點?.
還是三個重點呢?.
然後我們就說.
不是啊,四五個都可以的.

$^{81}$我覺得這本書很好.
因為這本書.
他告訴你.
這個比喻只有一個重點而已.
然後這個比喻是有三個重點.
我就覺得非常好.
因為.
起碼他給了你一個根基和基礎.
他框住了你.
你就不會亂說話.
所以就送這兩本經典給大家.
Sulphan有什麼好介紹的?.
我介紹一本比較進階一點的.
不是因為我是進階的.
不過剛才Mark Sir和Justin講的.
其實很重要.
其實是因應大家不同的程度.
有一本比較理解一點的.
叫做Klein's Logos.
他是一個作者.
他這本書叫做Stories with Intent.
其實和剛才Justin講的很類似.
和Mark Sir也是.
其實是有一個故事.
耶穌講的故事是有一個intention的.
所以就找到那個重點.
他是一個comprehensive guide to the parables of Jesus.
幾年前應該是18年就出了一個second edition.
你介紹這本書很厚.
Sulphan.
哈哈哈.
你厲害.
你又有second edition.
是厚一點的.
但為什麼這麼厚?.
這個值得一提.
我有幾樣東西很喜歡.
第一,他有很多的background information.
剛才Mark Sir講的cultural background.
其實有一個我們比較忽略.

$^{121}$不過這個比較進深一點.
是一些相關經文.
ancient manuscript related的東西.
有時很重要,有時可能次要一點.
但其實是講到很多雷同的說法.
讓我們看到其實耶穌不是作一個故事出來的.
無論在他之前還是之後.
其實都有類似這樣的教導.
或者一些這樣的用法.
關於一些經文或者道理.
但更加重要的是他列了很多不同的可能性.
這個很重要,有時你去明白.
不是只是別人塞進去給你的.
而是你自己也要去理解.
為什麼呢?.
我也有想過,有沒有可能呢?.
或者這幾個不同的可能性有什麼分別.
為什麼我要這樣選擇呢?.
我想這個作者就做了一個很好的幫助.
去幫你去理解和選擇.
他會去回答這些問題.
所以我想他用了一個很特別的做法.
他不是逐節解,不是這樣的類型的commentary.
他是用問題.
通常一個parables有七至十個問題.
我覺得這些問題很特別.
因為通常我看完一個parables有問題.
當我打算去問的時候,他原來已經問了.
看完一個就引申到另一個問題,他原來也有問.
他可以read到你的mind.
知道什麼問題是你很關心.
而他真的會去很仔細回答.
最後一個session就是如何adopt在今時.
我想這件事,如果你看多幾個就發現.
其實很多時候我們看parables太貪心.
好像剛才Joyce一開始這樣說.
很多時候我們就這裡抽一些,那裡抽一些.
其實就忽略了那個中心.
其實原來看得多,發現耶穌說的比喻其實很盜慾.
都是很環繞某一個天國的特質.

$^{161}$不過很多時候我們就跳出去,找了其他的東西.
反而忽略了一個重點.
剛才的書籍我相信會幫到大家去找到最想.
剛才Mark Sir說的,provocative.
很盜慾的,應該是這樣的.
希望這些書籍會幫到大家.
這個比喻的series我們在這裡完結.
我們城會經常送大包的.
我們記得一,二集的時候說過.
馬達方案第十三章好像很亂.
說一下這個,然後不解釋,之後才解釋.
其實有沒有問你的呢?.
說到經常說我們要上下問.
我們有一個這樣的.
一個交叉並聯的結構.
如果你有興趣的話,在下面的連結.
你訂閱了就可以下載.
送給你,希望你能夠做一些教學的工具.
或者自己研究的工具.
希望這個系列真的幫到大家.
可以成為一個真正天國的子民.
不夠後,像一些門徒那樣上前來.
有多些問題就可以用你的工具書.
下次再見,拜拜.
\newpage



\section{}
\label{sec:w_y0M2H_ndU}
\textbf{舊約網上行 - 上帝的足跡:EP5 巴別}
\newline
\newline
連結: \href{https://youtube.com/watch?v=w_y0M2H_ndU}{\texttt{https://youtube.com/watch?v=w\_y0M2H\_ndU}} ~~~~ 語音日期: 2021-03-09
\newline
\newline
\hyperref[sec:rLZ_I4dgAlo]{\small{< < < PREV SERMON < < <}}
~
\hyperref[sec:index_chronic]{\small{[返順時目]}}
~
\hyperref[sec:index_scriptual]{\small{[返順卷目]}}
~
\hyperref[sec:hPoAv0AxRzM]{\small{> > > NEXT SERMON > > >}}
\newline
\newline
$^{1}$創世記一至十一章是最吸引人的情節.
其中一個除了羅亞方舟之外.
就是巴別塔的故事.
我們先讀讀這段經文只是九節長.
創世記十一章一節開始.
那時全世界只有一種語言.
大家說同樣的話語.
他們向東遷移的時候.
在士拿地發現一塊平原.
就住在那裡.
他們彼此說.
來我們造磚把磚燒透吧.
他們就把磚當作石頭.
然後用石漆當作灰泥.
他們又說.
來我們建一座城造一座塔.
塔頂要通天.
我們要為自己立名.
免得分散在全地上.
然後說.
來看看世人建造的城和塔.
然後說.
看他們同是一個民族.
有一樣語言.
他們一開始就作者事.
以後他們所做的一切.
就沒有可以攔阻他們的了.
來我們下去在那裡.
混亂他們的語言.
使他們聽不懂對方的話.
於是耶和華把他們在那裡.
分散在全地上.
他們就停止建造那城.
因此那城的名就叫巴別.
因為耶和華在那裡.
混亂了全地所有人的語言.
又從那裡把他們分散在全地上.
一個很吸引.
令人聯想的故事.
你可能看過一些類似的藝術作品.

$^{41}$是跟巴別塔有關的.
這些巴別塔的畫像.
都是表達不同的時代.
那個畫家或者背景的想像.
這些塔是圓形的.
帶著很多歐洲的風味.
但是聖經的巴別塔到底是怎麼樣的呢?.
聖經當然不是直接寫給我們的.
而是寫給當時神的子民.
但是聖經卻是為歷世歷代的神的子民.
包括你和我而寫的.
So the Bible is not written to us.
but it's written for us.
所以我們第一步要去明白.
當時古近東的世界.
去問問這個巴別塔.
在當時的世界是什麼意思.
歡迎你跟我們往上行.
我們進入聖經的世界.
《聖經》 詞曲:李宗盛.
我們看巴別塔的故事的時候.
首先梳理一些基本的問題.
就是.
會不會這個故事就是很硬性的.
去告訴我們.
人類的歷史百科全書呢?.
是不是提出.
告訴我們.
人類萬國的語言.
真正的歷史.
linguistic history(語言歷史)是什麼呢?.
如果你看一下上下文的話.
似乎不是這麼簡單的.
為什麼呢?.
因為上文第十章創世記.
講到羅亞的家譜的時候.
羅亞有三個兒子.
閃含雅弗.
講到他們後代的時候怎麼說呢?.
十章第五節.

$^{81}$他說雅弗的後代散佈各地.
各人隨著自己的謊言.
自己的語言.
住在自己的地方.
二十節.
講到含的子孫.
也是按照自己宗族謊言住在自己的地方.
三十一節.
閃的後裔.
也是按照自己的謊言.
所以到第十一章的時候.
突然間又說全世界都有同一個的語言.
其實時序是不一樣的.
因為聖經的作者.
不是真的硬性地寫一個 chronological order(時序).
他是在寫一些主題的東西.
我們講過.
創世記一至十一章.
是一種起源史.
在當時近東的文化裡面.
經常會有的.
他著眼的不是在寫一個人類的百科全書.
發展史.
而是教導一個世界觀.
一個 worldview(世界觀).
是告訴我們.
神,人,世界.
應該有什麼關係.
告訴我們怎麼看世情.
巴別塔這個故事.
出現的地方在哪裡呢?.
十一章第一節.
他說全世界有同一個語言.
有同樣的華語.
他們向東遷移.
這裡記得 spiritual geography(心理哲學).
凡是向東的.
在頭十一章.
也不是好東西.
到了一個地方叫做士那地.

$^{121}$在那裡平原居住.
士那.
Shina.
指的是當時近東文獻的蘇美爾.
Sumer 這個地方.
就在美索不達米亞南邊.
現在伊拉克南邊一帶的平原.
在平原上有不少這些廟塔的遺跡.
至少現在發掘出來的有二十五個.
在平原上最壯觀最顯耀的.
就是這些廟塔.
當時人稱為 Sikarut.
其實是 Arkadian 的 Sakaru.
這個意思.
就是說向上建造.
To build high.
他建造這些廟塔的時候.
那時候廟塔下面旁邊.
就會有聖殿的神廟.
也是一個城市的中心.
所以為什麼這個故事去交代的時候.
就是有這樣的背景.
城市.
廟塔.
是全部一起說的.
這些廟塔.
不是讓人上去聚集的.
人不會花很多時間在上面的.
因為塔不是為人而建造的.
是神逗留和神居住的地方.
所以每個城市都有自己的主堡的神明.
這個地方是他住的.
不是人去的地方.
有時候.
那些祭司.
可能會上去廟塔那裡.
來打理一下.
因為這些 Sakarut.
就是為神明而做的.
讓天上的神可以下來.

$^{161}$可以說是一個很大的樓梯.
或者是一個電梯.
象徵著一個聖山.
讓神明方便他下來.
然後他下來的時候.
他就可以在旁邊下面的廟那裡.
接受拜祭.
這些 Sakarut 的名字.
告訴我們到底那個建築物.
那個廟塔.
到底想做些什麼.
譬如在 Nippur.
當時近東很大的宗教中心.
那個廟塔的名字叫做聖山之居宇.
就是 House of the Mountain.
吾爾.
你記得.
阿布拉罕就是吾爾出來的.
這個地方.
那個廟塔叫什麼呢?.
就是 Temple whose foundation creates aura.
就是創造光線.
創造光的基座的廟宇.
這個是一個很大的廟塔.
有六十米.
乘以四十五米那麼大.
高有三十米那麼高.
所以有好幾層樓高.
當然在說到近東的時候.
這些廟塔最出名的就是巴比倫那個.
也是這個創世記十一章故事所針對的.
它的名字是什麼呢?.
就是天地基座的樓宇.
The House of the Foundation of Heaven and Earth.
因為美索不達米亞本身.
天然來講是沒有木材.
以及很缺乏石頭的.
跟以色列巴勒斯坦地就不同了.
所以在南美索不達米亞.
一些龐大的建築物.

$^{201}$就要用泥磚來砌出來.
因為這一帶是比較潮濕的緣故.
所以一些重要的大的建築.
好像這些廟塔那樣.
那麼高的時候.
這些泥磚就需要燒硬它.
然後有很多時候就加上這些 bitumen(磚磚).
這些瀝青.
以致可以防水.
這個故事怎麼去描述這些工程呢?.
第三節.
他說他們彼此說.
我們要做磚把磚燒頭.
他們有把磚當作石頭.
就把石漆當作灰泥.
他們又說我們要建一座城.
做一座塔.
塔頂要通天.
是為了自己來立名的.
免得分散全地之上.
當時的人一聽的時候.
就會發現這是一個建築物的故事.
是一個很偉大的建築.
很宏大的一個工程.
但是以色列的希伯來聖經.
這個故事說出來的時候.
好像一至十一章那樣.
就是跟這些美索不達米亞的思想.
這些思潮.
去做爭辯.
你說建一個廟塔.
建一個大的工程.
有什麼大問題呢?.
首先.
在聖經裡面.
神聖的地方.
sacred place.
不只是這個地方出現.
每次說到sacred place.
神聖的地方.

$^{241}$就是神同在.
神居住的地方.
是一個高舉神明星的地方.
聖經強調很多次.
我相信在摩世五經.
最重要的.
講到這一方面的.
就是.
申命記.
申命記那裡講到.
他要敬拜神.
要重陽神的話.
就要怎麼樣呢?.
去到.
耶和華裡面的神選擇.
作為立他名的地方.
就是說.
那個地方.
his name dwells there.
之後這個地方我們知道.
就是從會幕變成耶路撒冷.
在耶路撒冷是神揀選.
去到立他名的地方.
要敬拜的話.
唯一一個的sacred place.
就是.
神所揀選.
看看巴別這個故事.
你就知道.
是人自己選擇.
一個的sacred place.
去到建造出來.
眾人的動機是什麼呢?.
為自己立名.
去到表揚自己的名聲.
在這個故事當中.
連同.
名字和名字相似的音的字.
出現了好幾次.
十一章第二節.

$^{281}$他們來到東邊遷移的時候.
去到史達地發現了一個平原.
就住在那裡了.
Sharm.
希伯來文是Sharm.
你聽起來.
你知道.
咦? 跟名字是同音.
十一章第四節.
他們又說.
來我們建造一座城造一座塔.
塔頂要通天.
Shemayim.
Heaven.
Sky.
又是跟Shem.
Sharm.
是同樣的音.
他們便為自己立名.
Shem.
這個就是名字.
免得分散全地.
第七節.
來我們下去.
在那裡混亂他們的語言.
使他們聽不懂對方的話.
在那裡.
Sharm.
第八節.
於是耶和華說.
他們從那裡.
Sharm.
分散全地.
他們就停止了建造.
第九節.
因此那一城的名.
Shem.
他的名字.
或者他的 reputation.
就叫做巴別.

$^{321}$因為耶和華說.
在那裡.
Sharm.
去到混亂全地.
所有人的語言.
又從那裡.
Sharm.
把他們分散到地上.
重點在哪裡呢?.
就是名字.
到底這一座塔.
他們所建造的工程.
是為了哪一個名而去建造的呢?.
如果你看看近東其他的.
建廟或者建塔.
這些冥文的時候.
你就知道當時在位的.
君王.
到底他們是為了什麼.
去建這些廟塔呢?.
舉兩個例子.
第一個.
在 Lagash這個古城.
有一個王叫做 Gurdia.
有一個人像.
他的大腿上面.
其實是一塊板.
那裡挑了一個廟宇的藍圖.
他記錄自己的功勞.
自己的風高偉績的時候.
就誇耀自己的工作.
他自己為自己立大名.
於是乎他能夠與眾神同植.
他說什麼呢?.
他說.
The great name.
that he made for himself.
第二個例子.
就是巴比倫王.
利保加里斯.

$^{361}$當他重修了巴比倫的廟塔之後.
他在石碑上面寫什麼呢?.
他說.
我堅固了高聳入雲的聖廟.
和巴比倫的圍牆.
我為自己的王位.
我的王權立下永久之名.
Everlasting name for my reign.
但是這個故事告訴我們.
巴比倫帝國.
當他建立他的霸權的時候.
起到差不多.
或者甚至起到最高峰的時候.
以為真正可以統領天下.
世界人民都講他的話的時候.
寫了一張第五字.
要說.
下來.
要看一看世人建築的城和塔.
怎麼看都好.
下來這個動作.
必然是一個諷刺和批評.
是一個 political satire(政治詐騙).
因為這個塔.
這條天梯.
原來只不過是這麼低.
要勞煩天上的耶和華下來.
他才可以察看人這麼大的工程.
到底真的有多大.
十一章第九節.
結尾的時候.
他說.
因此.
那城的名.
他的名聲.
就叫做巴別.
因為耶和華在那裡混亂了.
傳遞所有人的語言.
從那裡分散了他們.
你留意這裡.

$^{401}$也有一個所謂 work play(工作).
就是同音的字.
似音的字放在一起.
那城名叫巴別.
Babel.
英文.
我們翻譯成 Babylon.
其實這個就是巴比倫城的名字.
意思是什麼呢?.
就是神明的大門.
Gate of God.
就是全世界的中心.
神明上帝的大門.
通往上帝的那一處.
就是 Babylon.
就是 Babel.
但是這個故事.
經文.
卻將它顛倒過來.
將它識破.
就是說.
對於耶和華來說.
這個所謂巴比倫.
Gate of God 的地方.
其實不是真的通往上帝.
而是 Bala.
是混亂.
是一種 chaos(混亂).
而不是帶領人去認識上帝.
還有什麼可以幫助我們.
去了解創世記十一章.
巴別塔的故事呢?.
就是我們將故事裡面的一些語言.
一些語句,用詞.
去比對一下.
其他近東的一些文獻.
特別是一些帝國所寫的文獻.
一些帝王留下的一些名文.
譬如這裡舉一個例子.
我們可以比較一下.

$^{441}$在亞述王.
薩干二世.
薩干 the second(薩干二世).
他曾經征討了很多不同的民族之後.
他就留下了一些這樣的說話.
在一個 inscription(記載).
一個名文上面講.
他說.
世界四域的眾民.
講外語的.
講不同語言的眾民.
住在高山的,住在低地的居民.
我都憑著我的主神亞述神.
Ashore(阿述神)的說話.
憑著我的權杖的力量.
我都統治了他們.
並且我使他們有同一個的嘴巴.
one mouth(一個嘴巴).
就是說講同一個的語言.
將他們這些的眾民安置在這裡.
就是Dusarukin(杜撒魯金).
其實是亞述其中一個的京都.
現在叫做Crossiped(克羅西佩).
你可以看到.
這些說到同一把嘴.
同一個的語言.
one language(一個語言).
其實就是一種帝國統戰的一種語言.
《創世記特首》有一章巴別塔的故事.
其實就是將這些語言.
將這些用語.
然後再顛覆反轉過來.
他說耶和華才是眾民的真正主宰.
不是地上的君王.
所以你們看到這個故事.
可以叫它做political sceptre(政治帝國的批判).
或者政治帝國的嘲諷.
一種諷刺.
你以為這是天上的大門.
誰知道呢.

$^{481}$其實是一種沒有意義的廢話.
你以為是Bible(聖經).
但其實是Ballot(聖經).
二十世紀納粹黨統治德國的時候.
當時有一個教會所出來的運動.
叫做Confessing Church Movement(認信教會運動).
是認信教會運動.
這些教會運動.
都是反對納粹黨去控制教會的.
其中有一個牧師叫做Hermann Goethe.
Hermann Goethe.
是在柏林的一家路德會裡面務會的.
他前任的牧師叫做Martin Niemöller.
Niemöller可能大家都聽過.
他也是一個反對納粹統治教會的牧者.
1937年的時候就被抓進了集中營.
你可能聽過Niemöller的名字.
或者聽過他寫的一首詩.
那首詩叫做First They Came.
起初他們來了.
這首詩是講什麼呢?.
他說當納粹來抓共產黨員的時候.
我保持沉默.
因為我不是共產黨員.
當他們關押社會民主黨的人的時候.
我保持沉默.
因為我不是社民黨員.
當他們來抓工會的會員的時候.
我也保持沉默.
因為我不是工會會員.
當他們來抓猶太人的時候.
我仍然保持沉默.
因為我不是猶太人.
但當他們來抓我的時候.
就再沒有人來為我講說話.
這首詩很出名.
因為在巴別的統治.
一個大統帝國當中.
很多時候就會逐漸地使人不講說話.
或者只講一種說話.

$^{521}$Goris是繼承了Nimrod的目的.
在這段時間.
他講了很多的講道.
這些講道編成了一本書.
那本講道的系列叫做.
Sermons in the Time of World Crisis.
就是世界危機中的正道.
其中一篇的講道.
就是講解巴別塔這個故事.
他說這個故事.
其實就是導出歷史歷代帝國.
以及這些帝王的歷史.
這些帝王是用大統的方法.
單一的語言.
單一的政制.
單一的文化.
來統戰各地的人民.
美其名就是合一.
美其名就是平安.
去祝福人類.
在歷史當中經常會出現.
譬如羅馬盛世.
Pas Roman.
譬如德國的盛世.
或者大英的盛世.
甚至美國帝國的盛世.
今天你可以加上可能中國.
其他不同的大小的國家.
講到這些帝國.
一直在歷史裡面出現的時候.
就好像建立巴別塔一樣.
就是到頂尖的時候.
它總會倒塌.
因為在基因裡面.
是有那種腐敗在.
你記得1945年之後.
納粹德國統一.
歐洲的夢破滅之後.
人人都在講什麼呢?.
就是世界的政府.

$^{561}$World Government.
經濟的大同等等.
World Order.
但是今天.
過了這麼多年.
仍然看到世界是四分五裂.
就好像有一個定律在中間.
就是這個巴別塔故事所講出來的.
其實聖經的作者.
是一些很入微的觀察.
因為他看到大大小小.
這些逃亡敗瓦的巴別塔留下了.
這些遺跡看到的時候.
就發現這些帝國.
由近東開始.
不同大小的帝國.
巴比倫.
亞述.
埃及.
波斯等等.
他們的命運都是一樣.
其實實際上.
這個巴別塔的故事.
就是在預言今天的世界.
也是在預言將來會發生的事情.
猶太人這個民族.
經歷了很多這些改朝換代.
看到歷史裡面.
有很多的帝國.
去到欺壓過他們.
然後他們曾經亡國過.
被擄.
然後又回歸.
在這些這麼多的世紀的經歷裡面.
有時候你看看他們.
特別是猶太拉比.
解經的時候.
你真的會看到他們有一種.
獨到的地方.
猶太的拉比文學裡面.

$^{601}$有一種叫做 Midrash.
Midrash 其實是一種猶太人的拉比.
解經的方法.
就是很多時候在這些經文當中.
他會 read between the lines.
然後很多時候會運用自己的想象.
甚至將那些故事.
套進去經文裡面.
可能你聽不慣.
未接觸過.
你會覺得很奇怪.
因為很多時候他們解經的時候.
或者他們所謂講道的時候.
就會很多時候生安白造.
但是呢.
如果你聽明白那些猶太拉比.
解經的目的的時候.
為什麼要加插這些想象進去經文裡面呢?.
其實他們很多時候的目的就是.
想將它放大鏡.
放在經文上面.
讓你不會錯過那個細節.
加很多的想象.
加很多其他想到的故事進去.
以至你不會錯過細節.
特別是一些我們經常耳熟能詳的故事.
有一個 Midrash 是解巴比塔故事的.
拉比將前前後後的家譜的年份加起來.
當然這個是完全不合乎解經的原則.
但是無論如何.
他們將年份加起來的時候.
他們就說.
巴比塔建造的時候.
阿巴拉罕原來是在生的.
當時阿巴拉罕大概是五十歲.
這些完全沒有根據.
是憑空想像.
但是 Midrash 的故事是講什麼呢?.
他說有一天.
這些工人很忙來建造巴比塔的時候.

$^{641}$阿巴拉罕就去到巴比塔的地盤.
阿巴拉罕四處看的時候.
他說.
這些是什麼呢?.
就指著上面有很多的泥磚.
那些巴比的工人就去搭阿巴拉罕.
就說.
我們建造了一座很高很高的廟塔.
其他人看到的時候.
他們就會說.
哇.
他們真的很厲害.
他們的名聲真的很厲害.
誰建造這座塔這麼高的呢?.
是我們.
所以不用石頭.
也不用木頭.
因為石頭.
上帝可以做出來.
木頭.
上帝也可以做出來.
但是泥磚呢.
每一塊的泥磚.
都有蓋印.
是我們.
人做出來的.
是帝王自己做出來的.
所以用泥磚做出來的.
塔頂高天.
我們這樣做是希望.
上一代.
上帝用洪水毀滅了.
我們建造了一個高高的塔的時候.
洪水就不會再毀滅我們.
阿巴拉罕就說.
Fukele.
算了.
因為什麼呢?.
因為耶和華的名才是堅固的保障.
The name of the lord is.

$^{681}$the mighty tower.
the fortified tower.
甚至阿巴拉罕在這裡.
是在引用了真言.
幾百年之後.
所羅門寫的真言.
阿巴拉罕居然厲害到一個地步.
他可以立刻引用.
當然這些解經.
是很勉強的.
但是你不要忘記.
猶太拉比的Mirage.
他想做些什麼.
就是希望你不要錯過.
經文裡面的一些細節.
因為貫穿整個故事.
上下文的.
都是跟名有關.
Name.
人的名聲.
還是上帝的名聲.
所以阿巴拉罕看到這些巴別塔.
在建造的時候.
他就發現.
為自己立名.
地上的巴比倫帝國.
他的大統.
他以為能夠稱霸天下.
其實是很短暫.
其實是很愚昧.
你一天看不透這件事.
你就會成為巴別塔的奴隸.
最後.
會被上帝瓦解.
分散.
人類文明的發源地.
就是美索不達米亞.
當時第一代.
第一次在人類歷史裡面出現的帝國.
如果這個巴別塔的故事.

$^{721}$是寫於當時古代.
特別是舊巴比倫最強盛的時候.
這個巴別塔的故事.
就是一個先知的預言.
一個Prophecy.
就是說上帝終必會瓦解.
這些霸權和帝國.
或者.
可能到了後期的時候.
猶太人流亡被擄回歸之後.
再看這個巴別塔的故事的時候.
就是看到巴別倫帝國.
在短短七十年之後就覆亡.
然後人們看到巴別倫以前留下的光彩.
留下的輝煌的東西.
就是這些廟壇.
剩下的只不過是泥磚.
廢墟.
是一些很深刻的歷史教訓.
巴別塔的故事.
說到人為自己的名去建築.
為自己的帝國.
為人的Reputation去建立.
跟這個對立相反的.
就是神的名.
The name of Yahweh.
耶和華的名.
你看看巴別這個故事之後.
就說到什麼呢?.
閃的加堡.
十一章第十字.
以下是閃的後代.
閃.
Shem.
就是明治的意思.
明治的後代繼續延續下去的時候.
你就會去到阿不蘭.
亞伯拉罕.
就是上帝呼召亞伯拉罕出來.
亞伯拉罕就回應上帝的呼召.

$^{761}$亞伯拉罕和他的子孫.
所代表的是什麼呢?.
就是一個高舉上帝.
明治的世代.
一個呼求上帝.
去敬拜創造主獨一上帝的世代.
求主幫助我們.
我們能夠成為這個世代的一份子.
縱使我們看到這個世代是多麼艱難.
這個年代是多麼黑暗.
願意我們只高舉一個名字.
就是Hashem.
就是上帝耶和華.
耶穌基督的名字.
我們同心低頭祈禱.
我們的主耶穌基督.
你曾經被定罪.
被無辜的囚禁.
被凋棄.
被無辜的審判.
被凌辱.
人類的痛苦.
你每一點都嘗過.
我們知道你仍然會與.
為易受迫害的人同在.
也不會忘記這些.
在這個時候我們在世界各地.
有很多陷於受壓迫的情況的人.
我們去禱告.
為了香港的47個人.
更多無辜被拘禁的人.
我們去禱告.
為了在緬甸抗爭的人.
為了他們的自由和尊貴而站出來.
他是捱著軍政府的子彈的人.
主要是為了他們去禱告.
求你自己親自的.
叫我們在這些不安的動盪當中.
賜給我們信心.
保守我們不會陷於失望的裡面.

$^{801}$也喚醒我們的良心.
使我們能夠去明辨是非.
去持守正義.
持守公義真理.
拯救我們離開這些怯懦.
懼怕.
主我們知道在你裡面.
我們永遠有盼望.
親愛上帝你是創造和救贖我們的.
你更加是萬國的主宰.
你是萬國的審判主.
救贖主.
主一是你見到.
也知道惡人所做的一切.
也痛恨世上所有不義.
自大.
邪惡.
主一是我們知道那些誠心尋求你的人.
你必定會和他們同在.
求主你親自私下.
要書.
憐憫.
親自來和他們同行.
為那些無辜受監禁的人.
為他們特別的親人.
他們的家人.
主我們求你親自連續他們.
主願你親自來釋放.
使他們恢復他們的自由.
主我們知道無論遇見什麼事情.
主你讓我們堅忍和盼望.
但願你自己.
你的聖名.
是被稱頌.
是被高舉.
我們奉主耶穌基督的名字.
祈求.
阿們.
\newpage



\section{}
\label{sec:hPoAv0AxRzM}
\textbf{舊約網上行 - 上帝的足跡:EP6 離別吾珥}
\newline
\newline
連結: \href{https://youtube.com/watch?v=hPoAv0AxRzM}{\texttt{https://youtube.com/watch?v=hPoAv0AxRzM}} ~~~~ 語音日期: 2021-03-16
\newline
\newline
\hyperref[sec:w_y0M2H_ndU]{\small{< < < PREV SERMON < < <}}
~
\hyperref[sec:index_chronic]{\small{[返順時目]}}
~
\hyperref[sec:index_scriptual]{\small{[返順卷目]}}
~
\hyperref[sec:uhU9sAnkAzo]{\small{> > > NEXT SERMON > > >}}
\newline
\newline
$^{1}$歡迎來到《九月網上》第六集.
過去五集我們研究過《創世記》一至十一章.
你可能會覺得好像沒有什麼行程沒有什麼去過的地方.
因為《創世記》一至十一章的背景.
都是集中在美索不達米亞的南面.
一至十一章是整本聖經的導演.
告訴我們世人多次叛變上帝的故事.
第一個故事我們看過.
就是人類代表亞當夏娃叛變上帝想自立為神.
被審判之後流亡.
神有什麼反應呢?.
神反而是應許去救贖.
並且願意離開伊甸與罪人同在.
有一節很重要的經文是三章第十五節.
「我要使你和女人彼此為首」.
你的後裔和女人的後裔亦要彼此為首.
「她要傷你的頭,你要傷她的腳筋」.
這裡講到一個爭戰.
一個人與蛇的爭戰.
人與邪惡的爭戰.
女人的後裔會傷蛇的頭.
蛇會傷她的腳筋.
有少許伏筆告訴我們.
上帝如何救贖人類.
應許有一位女人的後裔.
與邪惡的爭鬥.
結果女人的後裔會得勝.
然後展開整本聖經的故事.
第十一章是巴別塔的故事.
人為自己立名.
上帝來混亂和分散他們.
最後是阿拉罕的出現.
這一大段我們稱為族長的歷史.
Patriarchal History.
裡面主要講到幾個以色列人.
神子民的始祖事跡.
就是阿拉罕以撒亞國和十二派.
這段歷史大約在公元前2000至1500年.
背景當然是古近東的世界.
當中最重要的角色當然是阿拉罕.

$^{41}$想告訴你.
無論你看過聖經還是沒看過.
信聖經還是不信.
你都需要認識一下阿拉罕的事跡.
為什麼呢?.
我帶大家去耶路撒冷的舊城.
舊城大概是一個正方形.
分開四個區域.
右上角是穆斯林區.
右下角是猶太人區.
左邊的上下就是基督教和亞美尼亞區.
這四個大區域代表著世界三大宗教.
三大獨一神的信仰.
而三個區域裡面都有一個神聖的地標.
穆斯林區當然是金頂清真寺.
猶太人的區域Jewish Quarter.
當然最神聖的地方就是西牆Western Wall.
而基督教就是聖母堂.
就是耶穌基督死而復活的地標.
三大獨一神的信仰.
代表著今天世界的人口成千上萬甚至數億.
所以你要了解今天世界的狀況.
或者世界的宗教政治動向.
你不能夠不了解亞伯拉罕的事跡.
但我想最重要的是.
《約聖經》裡面告訴我們.
每一位信奉耶穌基督的人.
都是以亞伯拉罕為我們屬靈的父親.
因為我們就是亞伯拉罕的子孫.
亞伯拉罕是我們所謂的信心之父.
所以他的事跡就是一個信心的旅程.
教導我們如何跟隨神跟隨耶穌基督.
多謝.
亞伯拉罕的母親.
亞伯拉罕的父親.
亞伯拉罕的母親.
亞伯拉罕和亞伯拉罕其實是同一個人.
亞伯拉罕的意思就是尊貴之父.
然後上帝為他起了一個新的名號.
叫做亞伯拉罕.

$^{81}$意思就是萬國之父.
亞伯拉罕出現的時候.
就是創世記第十一章尾的時候.
沒有交代他童年或青年的時間.
一開始的起點就是加勒底人的吾爾.
這是南米索不達米亞一個歷史悠久的古城.
我們所知道的吾爾這個地方其實也有很多.
在公元前2500年的時候.
當時是吾爾城邦最高峰的時候.
我們叫他做第一皇朝.
第一皇朝.
在1922至34年之間.
有一個英國考古學家叫做Wooley.
他去到吾爾這個地方發掘的時候.
就發現了很多舉世聞名震驚世界的出土.
那裡他找到16個皇陵.
這些皇陵.
當時的皇的墳墓.
其中兩個的皇陵特別令到新聞界震撼.
因為裡面找到幾十具屍體.
最著名的一個墳墓.
就是屬於一個女皇.
裡面有74具屍體.
6具是男人.
其他都是女子.
裡面有很多名貴的陪葬品.
譬如金樹山羊.
女皇的頭飾.
戰車.
用金箱建造的樂器等等.
屍體很整齊.
擺放在墳墓裡面.
並且穿上了華麗的衣服.
有頭飾.
Wooley開始發掘這個地方的時候.
他起初去解釋.
這些年輕的男女.
被埋到女皇的墳墓裡面.
原因是.
是一種陪葬的人製.

$^{121}$即是軍王一聲令下.
他們全部進入墳墓裡面.
然後將他們殺死.
以致他們成為陪葬的人製.
所以Wooley稱這個為死亡大坑.
Great Death Pit.
之後.
Wooley再檢驗這些屍體的時候.
發現很整齊的擺放.
沒有掙扎的跡象.
於是他太太給他一個意見.
他太太認為.
這些宮女和侍從.
其實是一起服毒.
去到自殺.
然後去陪葬.
因為有一個新的解釋.
1928年的時候.
英國有一份報紙叫.
Illustrated London News.
那裡有一個畫家.
將這個情況.
用一個畫像來表達出來.
你可以看到.
這些黑白的油畫.
講到這些隊伍.
進入墳墓裡面.
然後他們服毒自殺.
然後每一個都倒閉.
以致這麼多年之後.
Wooley再找出來.
這些是很毛骨悚然的情景.
當時是驚動一時.
但近年.
在University of Pennsylvania.
用CT Scan去掃描幾個頭骨.
發現幾個頭骨都有一個.
30毫米的斧頭造成的創傷破口.
所以現在學者相信.
這些宮女和侍從.

$^{161}$不是服毒而死.
而是先被殺死.
然後整整齊齊的.
穿好衣服.
有頭飾.
然後放進墳墓裡面.
作為女王的陪葬.
怎樣也好.
這些都是告訴我們.
當時古米索不達人米亞.
吾爾的王朝.
甚至那裡的背景.
有一些神人祭的儀式.
怎樣也好.
是一個非常富庶的地方.
你看到那些金器.
這些名貴的珠寶.
全部都不是米索不達人米亞.
本有的天然資源.
全部都是從第二個地方.
藉著貿易運來的.
所以這個吾爾的地方.
確實是一個非常富庶.
當時是一個很繁榮的世界.
但最終在公元前2000年的時候.
瓦解了.
也宣告了蘇美爾文明的沒落.
阿巴拉罕的事跡.
是在這個時間之後.
在吾爾最輝煌的時代之後才出現.
雖然吾爾在阿巴拉罕的時候.
不再稱霸米索不達人米亞.
但仍然是一個非常重要的經濟宗教中心.
這些發現可以幫助我們.
想像一下阿巴拉罕當時的背景.
吾爾的文化.
創世記十一章這樣記載.
以下是他拉的後代.
他拉生阿巴蘭 拿學和哈蘭.
哈蘭生羅德.

$^{201}$哈蘭在他父親他拉之前死在他出生之地.
就是迦勒底人的吾爾.
阿巴蘭和拿學各自娶妻.
阿巴蘭的妻子叫做撒萊.
拿學的妻子叫做密格.
是哈蘭的女兒.
哈蘭是加密和易加的父親.
以往近東的家譜的作用就是追溯一個民族的歷史.
聖經裡面的家譜就是追溯神的子民.
由亞當 瑟特 羅亞直到閃的後代.
現在追溯到他拉 即阿巴拉罕的父親.
這些家譜的名字我們知道.
全部都是屬於西閃族.
即敘利亞一帶的民族名字.
所以他們不是南美索不達米亞本土的名字.
為什麼呢?.
似乎我們所知道的歷史背景.
在公元前大概二千年左右.
當時有一個遷徙的現象.
就是從美索不達米亞的人眼中.
有些西方人 敘利亞的人.
他們很大批逐漸遷徙到美索不達米亞兩河流域的地方.
這個現象叫做阿摩利人遷徙.
Amorite Migration.
所以他們很多時候因為不同的原因.
這些阿摩利人遷徙到南邊.
然後落地生根.
甚至進駐在南美索不達米亞一帶.
建立自己城邦.
所以他拉和阿巴拉罕.
似乎都是來自這些阿摩利人的民族.
我們看看家譜的名字.
都可能會猜到他們宗教信仰的背景.
他們不是敘利亞人.
他們是阿摩利人.
他們的名字是Tala.
Tala這個字Tala.
很可能是一個月神的名字.
或者是一些很典型的阿摩利人的名字.
我們也知道.

$^{241}$唔爾或者將來他們去到哈蘭的古城.
都是當時拜月神的宗教地方.
事實上在聖經之後也是這樣交代.
在約書亞記二十四章第二節.
耶和華以色列的神使這樣說.
古時你們的列祖.
說到阿巴拉罕那河的父親Tala.
他們住在大河的那邊.
在幽伐拉底河的美索不達米亞地方.
他們是事奉什麼?.
別的神.
就是其他的眾神.
所以屬靈上.
我們看到這個家譜.
似乎發現這個家族開始已經死亡.
Exspiritual death.
Tala生了三個兒子.
阿巴拉罕那河和哈蘭.
聖經的家譜通常都是很工整的.
好像我們看《創世記》第五章等等.
每逢有離題或者附插了一些細節的時候.
就有特別的意思.
這裡有兩處地方值得留意.
第一.
就是哈蘭在父親死之前去世.
這是非常不幸的.
這是Physical death.
另外一處就是撒萊.
就是阿巴蘭的妻子.
不能夠生育.
是沒有孩子的.
當時在古晉通的社會裡面.
絕子絕孫是非常不幸.
因為沒有人去延續你的名字.
所以我們看到這個家譜告訴我們.
阿巴拉罕的世代.
來到這個時間.
在屬靈上已經死亡.
在身體上也死亡.
好像我們看星球大戰.

$^{281}$最後的The Last Jedi.
沒有人再繼續這個神的世代.
好像一件事都玩完了.
31節.
塔拉帶著她的兒子阿巴蘭和她的孫子.
就是哈蘭的兒子羅德.
以及她的媳婦撒萊.
一同出了迦勒底的梧意.
想要去迦蘭這個地方.
他們到了哈蘭這個古城.
住在那裡.
32節.
塔拉死在哈蘭.
她的壽數是205歲.
從這個故事的情節來看.
告訴我們.
塔拉和阿巴蘭罕.
有一個神的呼召.
所以整個家離開了這個古城.
本身想去的目的地就是迦蘭.
他們應該按照當時古代商貿路線.
或者沿著幼發拉底河不斷北上.
經過一個大的城邦.
叫做馬里.
是一個大的阿摩利亞城邦國家.
那裡出土了很多泥板和文獻.
裡面提到一些雜族.
與創世記族長阿巴拉罕爾薩亞國.
他們的風俗非常近.
從馬里一直向左上角.
向西北的方向.
來到哈蘭這個大的古城.
這是美索不達米亞的北面.
不知什麼原因他們停留在這裡.
到底是叛逃而廢.
還是他們回到自己的阿摩利亞家鄉.
他們留戀在那個地方.
無論如何.
都好像是一個死路一條.
神的西多能否延續.

$^{321}$就是這個時候.
上帝說話.
創世記十二章第一節.
對阿巴蘭說.
你要離開本地,本族,父家.
到我指示你的地方去.
然後有七句說話.
就是上帝對阿巴拉罕的應許.
我必使你成為大國.
賜福給你.
使你的名為大.
你也要使別人得福.
給你祝福的.
我必賜福給他.
興翰你的或者周助你的.
我必周助他.
第七句.
地上的萬族都要因你得福.
神賜給阿巴拉罕一個大國.
你也應許告訴他.
會有一個大的名字.
別人會因為他而得福.
萬國會因為他而得福.
這個其實都是回應創世記一至十一章.
裡面不同的主題.
你記得人類在這頭十一章裡面.
不斷地叛變上帝.
並且是一種愚昧的行動.
帶來的只是審判和咒詛.
亞當夏娃我們說過.
他們叛變是想得到智慧.
想變成神.
誰知道以為是一個福氣.
但是帶來的就是咒詛.
第四章.
就是戈仁.
他的敬拜獻祭.
得不到上帝閱立.
他就說想為首.
誰知道他也是流亡.

$^{361}$第六章.
一個比較難解的經文.
就是說當時的日子.
有所謂巨人.
這些人是神的兒子.
是君王.
他和人的女子結合.
生了上古英武有名的人物.
這些君王也是罪大惡極.
因為以自己的名為大.
上帝說要洪水氾濫去審判這個世界.
十一章我們說過巴別塔的故事.
就是人為自己來立名.
一直來說都是人類.
想高超自己.
說自己要為自己來立名.
甚至超越神.
在這裡有一個調轉.
就是上帝叫阿不蘭的名為大.
以至他能夠使萬國得福.
即是說要使萬國認識創造主.
上帝對亞伯拉罕的呼召.
開展了上帝救贖人類的計劃.
但是有一樣很重要的東西.
就是人要回應.
這裡唯一一個命令式.
就是離開.
離開不易.
或者可以以這段經文去翻譯.
上帝對亞伯拉罕說.
你要離開.
以至我能夠在你身上實施實行.
我祝福萬國的計劃.
亞伯拉罕如何回應?.
十二章第四節.
阿不蘭照著耶和華吩咐他的去了.
羅德也跟他同去.
阿不蘭離開哈蘭古城的時候已經是七十五歲.
阿不蘭帶著他的妻子撒萊和侄兒羅德.
以及他們在哈蘭古城裡面.

$^{401}$所積蓄的一切財物所獲得的人口.
一同出去.
要到迦南地.
或者當時是跟著國際的路線.
一路在哈蘭古城南下.
在阿不蘭這個大古城城邦的東邊經過.
然後去到另一個大的站.
就是大馬士革Damascus.
後來他們終於進入迦南地.
他們可能去到第一個大的站.
就是夏索這個大城.
到第六節.
阿不蘭穿過那裡.
到了事件的地方.
摩利的橡樹那裡.
當時迦南人住在那裡.
耶和華向阿不蘭顯然說.
我要把這地賜給你的後裔.
阿不蘭就在那裡.
為向他顯然的耶和華築了一座祭壇.
後來他也在那裡遷到南邊.
就是伯特利東邊的山地.
在那裡搭起帳篷.
西邊是伯特利東邊是艾城.
他在那裡也為耶和華築了一座祭壇.
呼求耶和華的名.
阿不蘭又在那裡起程.
再南下漸漸移到南地Negev這個地方.
迦南地的地理.
有一個中央的山地.
這是一個天然的分水嶺.
這條山路是由士劍這一帶開始.
士劍本身地理是一個十字路口.
士劍這個城市到今天.
還是一個繁榮很大的巴勒斯坦城鎮.
它是一個十字路口東南西北的交匯之處.
然後再去看伯特利和艾這個地方.
其實就是中央山地中間的另一個十字路口.
它連絡東邊耶利哥平原死海一帶.
經過中央山地就可以往西.

$^{441}$去到迦南地沿海貿易的國際路線.
所以又是另一個交匯點.
去到南地也是一樣.
Negev這個地方也是迦南地另一個交匯點.
連絡阿拉伯地約旦河東邊的貿易.
和沿海國際平原的路線.
所以三個地方都是迦南地天然的交匯點.
士劍,伯特利和艾和南地.
在迦南地做什麼呢?.
在迦南地做什麼呢? 祝祭壇.
去呼求耶和華的名.
什麼意思呢?.
呼求可能是call upon的意思.
即是祈禱向上 去到上帝呼求.
求上帝給恩典或幫助.
但是呼求也可以是call out的意思.
calling out是什麼意思呢?.
即是在眾人面前宣揚神的名聲.
宣揚神的大名.
要眾人認識他所敬拜的上帝.
聖經講到信心之父阿巴拉罕的事跡的時候.
第一件事要交代的就是他的旅程.
journey.
很多時候我們會說什麼叫做旅程.
就是由一個地方去另一個地方.
這個就是行程.
但很多時候文學研究都會有更加仔細的定義.
什麼叫做旅程呢?.
一個定義就是adventure 一個歷險.
什麼叫做歷險呢?.
歷險就是一個行程.
由開始到完結的時候.
你去一個地方再回來.
你自己本身沒有什麼大的影響和改變.
譬如在阿巴拉罕那段時間.
當時埃及有一個很有名的.
耳熟能詳的並且找到很多抄本的旅程故事.
當時埃及是中王國的時候.
Middle Kingdom.
有一個文學著作很有名.

$^{481}$叫做西魯克故事.
Tale of Sinuhe.
簡單來說這個故事是講一個埃及的大臣叫西魯克.
因為政變的緣故他離開了自己的家園埃及.
他逃亡到迦南地那裡寄居.
跟那裡的城邦有很多族長去周旋化險為夷.
最後埃及的政局穩定下來.
他就回到國家.
然後回到自己的家園升官發財.
然後風光泰壯.
這個故事形容迦南地的情況.
就好像創世記阿巴拉罕等等故事的描述相似.
所以可以做一些歷史背景的研究.
從文學看來西魯克的故事.
可以算是一個自傳的故事.
一個冒險的故事.
因為以動作為主.
去了一個旅程增廣見聞.
然後就像我們今天在網上看到的旅行日誌.
去完之後人生多了很多內容.
拍了很多照片等等.
但是他的生命的角色根本上沒有什麼大的改變.
小孩子小的時候.
有時候會陪他們看一些兒童的卡通節目.
很多時候就是一種冒險.
是幻想出來的.
下午很悶.
他們幻想去到另一個美地.
一個世界那裡可能很有趣.
有很多不同的有趣的情節.
見識多了.
回到家就依然很窩.
這個叫做冒險.
亞巴肯的故事不是一個歷險.
不是一個冒險.
而是一個探索.
你說冒險和探索有什麼分別.
探索是一個探索的旅程.
就是以主角的演變為導向.
不是他所見所聞這麼簡單.

$^{521}$而是說這個角色在行程當中有什麼改變.
旅程當中甚至未結尾也好.
他不斷在改變.
或者完了行程的時候.
你完全不認得這個人.
因為這個人的改變實在太大.
其實你有沒有想過.
聖經的信仰.
去認識神.
基督教的信仰.
其實不是一個歷險.
不是一個冒險.
而是一個探索.
一個任務.
在這個旅程當中.
你會發現自己不斷有改變.
不斷對世界的事情.
對世情.
對生老病死.
對名氣.
對錢財.
對成功失敗.
對人生的事情.
都會不斷改變.
到最後可能是完全不同的角度.
因為在你身上.
應該是看到上帝的恩典.
上帝的明星.
看到上帝不斷在你身上工作.
所以在這個旅程這個任務當中.
你是不斷地演進的.
我們回到《創世記》十二章第一節.
這裡有什麼特別的地方.
你看看經文裡面.
神對阿布拉罕的呼召是怎樣的.
通常叫一個人去到起行.
進入一個旅程的時候.
最重要講述的.
應該是目的地.
他的destiny.

$^{561}$你參加一個旅行.
你最關心的是什麼.
當然是目的地.
那裡的氣候是怎樣的.
風土人情是怎樣的.
習慣是怎樣的.
目的地是我們最關注的.
但是在這個旅程的呼召裡面.
神呼召阿布拉罕.
怎樣講目的地呢.
他說到我指示你的地方去.
你看《上下文》當然我們知道.
就是迦南地.
但目的地其實不是講得那麼清楚.
沒有什麼地圖或藍圖給他看.
總之就是說.
你會成為泰國.
你現在不育.
誰是後代.
所以你看阿布拉罕的故事的時候.
你就不斷猜.
到底誰是繼承人.
到底是羅德還是其中一個僕人.
是不是以實馬利呢.
所以開頭的時候沒有講清楚.
之後故事才慢慢去到交代.
這個旅程的描述.
往往跟我們常理相反.
你看《十二章》第一節.
所關心的是什麼.
是起點.
是起始的地方.
或者是旅客本身.
你要離開本地本族 父家.
你看原文.
應該是這樣翻譯.
你要離開你的地.
你的族群.
你的父家.
唯一一個命令.

$^{601}$就是離開.
猶太人的拉比文學.
經常討論這句話.
就是離開.
整本聖經.
這個命令.
是出現過兩次.
除了這裡之外.
第二次是何時出現的呢.
就是創世記第二十二章.
一個很熟悉的故事.
也是令我們非常費解的故事.
就是阿布拉罕獻意殺.
二十二章第一節.
「這事事以後.
臣要試驗阿布拉罕對他說.
阿布拉罕回答.
我在這裡.
臣說帶著你的兒子.
就是你的獨一.
你所愛的.
意殺.
你離開.
到摩里亞地去.
獻意殺故事.
又是用到這個字眼.
離開.
並且三次說你的.
你的獨一 你所愛的.
你的意志.
兩段很清楚.
是頭尾呼應.
十二章第一節.
臣第一次對阿布拉罕說話.
整個敘述的歷史故事.
到二十二章.
這個獻意殺故事.
又是耶和華.
最後一次跟阿布拉罕說話.
之後就沒有直接跟他說話.

$^{641}$兩次都是說什麼.
Lack luck.
是離開.
或者中文可以這樣翻譯.
擺脫.
擺脫你的地方.
擺脫你令你安心的地方.
擺脫那些最綁著你的東西.
那些可能是你最愛的東西.
你留意他不斷強調.
你的.
你自己.
是綁著,枷鎖著你的.
以致你不能夠去宣告上帝的名.
以致你不能夠去經歷上帝在你身上的工作.
你的地.
你的民族.
你的富家.
很可能也是你最心愛的東西.
這些可能就是在我們旅程裡面最大的難阻.
讓我們體會不到上帝的工作.
所以可能這樣翻譯.
是什麼叫離開.
離開你本族,本地,富家.
Journey out of yourself.
從你自己的生命當中.
擺脫自我一個長征旅程.
信心之父.
第一個旅程.
就是告訴他你要看見上帝.
你要和上帝建立關係.
你要與上帝同行.
在一個旅程裡面.
最重要的不是目的地.
最重要的是起點.
是我們自己.
當我們決定跟隨耶穌基督的時候.
其實起不起也是一樣.
耶穌對我們的呼召.
Follow me.

$^{681}$跟隨我.
是一個旅程.
最重要的不是目的地.
而是我們能否起步.
我們能否放下一些綁著我們的東西.
以致我們看見上帝在我們生命當中.
不斷有的工作.
在這個quest裡面.
這個探索當中.
我們的生命.
我們的信心.
不斷地改變.
最近你和上帝有什麼祈禱.
或者你對上帝有什麼要求.
很多時候不是祈禱多少的問題.
有時候可能越祈禱得多.
反而不是一件好事.
因為視乎你怎樣祈禱.
視乎你祈禱的方向.
很多時候我們祈禱的.
很多時候都是一些安全的禱告.
Safe prayer.
上帝我有難你救我.
上帝你醫我.
上帝祝福我兒女.
祝福我的家庭.
保護我等等.
這些有沒有錯.
不一定錯.
但很多時候不知不覺.
反而會讓我們的旅程.
我們的探索停下來.
變成上帝是為我服務.
最終是為自己臨命.
不是為上帝建制壇.
不是去叫出上帝的名字.
不是宣揚上帝的名字.
宣揚上帝在我身上所做的事.
在這個信心之父的第一個旅程.
就是給我們的挑戰.

$^{721}$在旅程當中.
上帝要實行他在你身上的計劃.
使更多人認識上帝的偉大.
認識上帝的名義.
不如我們試試一些比較危險的祈禱.
危險的祈禱就是.
上帝你可否容我?.
使用我?.
叫人在我身上.
見到你的名偉大.
見到更多人藉著我.
他們可以得福.
可以認識創造主你自己.
一個危險的祈禱.
一個探索旅程的禱告.
邀請你和我們一齊來到去.
這樣禱告.
這樣經歷神.
我們同心得著祈禱.
因為你是那位帶領我們.
人生旅程探索的神.
主,我們來到你面前.
我們也求你幫助我們.
我們所求所想的.
是一些所謂危險的祈禱.
好像亞伯拉罕第一個旅程.
告訴我們一樣.
我們要見到你自己的名偉大.
你要使我們成為一個.
祝福其他人的生命.
叫其他人欣賞.
你自己在我們身上的工作.
你自己的名聲.
能夠被宣揚.
叫更多人來認識你.
叫更多人能夠得福.
祝福我們.
祝福你自己的子民.
你的教會.
讓我們成為一個.

$^{761}$好像亞伯拉罕探索旅程的生命一樣.
你在我們生命當中所做的新工作.
叫我們的生命能夠祝福更多的人.
祈禱奉主耶穌基督的名字.
\newpage



\section{}
\label{sec:uhU9sAnkAzo}
\textbf{考古探秘 (發現新死海古卷) -  陳崇基牧師 Mark Sir (2021.03.21)}
\newline
\newline
連結: \href{https://youtube.com/watch?v=uhU9sAnkAzo}{\texttt{https://youtube.com/watch?v=uhU9sAnkAzo}} ~~~~ 語音日期: 2021-03-22
\newline
\newline
\hyperref[sec:hPoAv0AxRzM]{\small{< < < PREV SERMON < < <}}
~
\hyperref[sec:index_chronic]{\small{[返順時目]}}
~
\hyperref[sec:index_scriptual]{\small{[返順卷目]}}
~
\hyperref[sec:neXNn__inUY]{\small{> > > NEXT SERMON > > >}}
\newline
\newline
$^{1}$各位好,很久沒有跟大家講過聖經考古探秘的節目了.
因為以前曾經做過直播,做過Live.
因為工作時逢太忙,所以已經沒有再做了.
過了一段時間,我想不如做一些短片.
一些短的節目,15分鐘左右,20分鐘.
跟大家分享一下這方面的資訊.
今次想跟大家分享一下一個廣傳的新聞.
在網上大家可能看過了.
在以色列猶大曠野的一個洞穴裡面.
所謂找到一些新的Dead Sea Scroll Fragments.
一些死海古卷的碎片.
這個地方到底是哪一處呢?.
在猶大曠野的其中一個山洞.
這個洞穴,考古學家叫做恐怖洞穴.
為什麼叫做恐怖洞穴呢?.
等一下就告訴你了.
在這個地帶,叫做猶大曠野.
說找到新的死海古卷碎片的Fragments.
為什麼會說是死海古卷呢?.
大家知道20世紀最大的考古發現.
一定是Dead Sea Scroll.
找到的地方就是在這裡.
就是昆蘭這個地方.
我們知道當時第一世紀前後.
一些猶太教派移居到這個地方.
他們不喜歡耶路撒冷.
所以他們就建立自己的社群.
在昆蘭這個地方.
並且在這裡可能抄寫了很多聖經.
在這些山洞那裡就藏起來了.
以致2000年之後.
20世紀中1947-48年的時候.
很巧合地就找到了.
一個最大的考古的發現.
幾十年之後.
一直都沒有找到聖經的抄本.
突然60年之後.
就在下面這一處.
南面這一處其中一個洞.
就找到一些聖經的抄本.

$^{41}$所以這些都是一些很重要的東西.
這裡就是昆蘭.
就是死海古卷發現的地方.
這個洞穴在溪谷旁邊.
在一個懸崖峭壁那裡.
是一個很難進入考古工作的地方.
據報導故事是這樣的.
其中一位考古工作人員.
他說一開始在洞裡面搜索的時候.
想去洗手間.
去旁邊洗手間的時候.
發現原來地上不是普通的泥土.
有一些跡象是羅馬的草鞋.
是鞋子.
撥開之後他就發現.
證明洞裡面仍然有很多文物.
可能還沒有發掘出來.
所以就決定更加徹底地去清理洞穴.
所以發現這個所謂的死海古卷的地方.
就是叫做恐怖洞穴.
你見到下去這個地方.
是需要像爬山一樣從上面掉下去.
因為是沒有石級和其他方法進去的.
為什麼叫做恐怖洞穴呢.
這個洞穴為什麼會這樣命名呢.
因為當時在1960年和61年.
其實考古學家已經去過這個洞穴.
做了考古工作.
當時一去的時候.
就發現這個洞穴裡面有40具人的屍骨.
有男有女有小孩.
所以就非常恐怖.
這些鞋骨屬於第二次猶太人反羅馬革命的時候.
或者稱為Bacobab Revolt.
公元後大約132至36年間.
這些革命分子應該想逃避羅馬的追兵.
所以就藏身在這個隱密的洞穴裡面.
結果羅馬士兵就知道他們躲在這裡.
然後在山頂上駐紮守住.
不讓他們走.

$^{81}$其實在山的上面.
山崖的頂上也找到一個大約150平方尺的羅馬士兵營.
結果大家猜到了.
那些革命分子和家人.
年老年少的都死在這個洞穴裡面.
2000年後考古學家就在這裡發現.
是非常恐怖的景象.
這次再回到這些洞穴重新發掘.
想徹底一點的時候.
就發現到一些聖經抄本的碎片.
大約20多份.
是非常殘缺的.
要組織出來的時候.
就好像要砌圖一樣.
放大鏡等等.
李英這些聖經抄本.
都是第二次猶太革命分子.
他們帶去和留下的文物.
是屬於這個時期.
因為Radio Carbon Dating.
都是鑑定這些殘卷.
是公元後二世紀.
就是Bakufa Revolt的時候.
這些新的抄本.
上面能夠讀到的文字.
就知道有十二小先知書裡面的一些經文.
現在我們知道的.
就是有撒迦利亞書和拿紅書裡面的一些經節.
其中一塊碎片.
寫著這句話.
就是撒迦利亞書第八章十六節.
大家看到這些字.
和翻譯出來的時候.
就知道是這一節的經節.
就是你們應當行者些事.
各人要向論社說真話.
在城門口要憑著誠實和公平.
施行審判.
值得留意這些抄本碎片.
經文都是希臘文.

$^{121}$但是每逢到神的名節的時候.
就會用Paleo Hebrew.
就是古希伯來文寫出來.
就好像我們今天寫中文.
然後寫著寫著寫著.
G-O-D或是The Lord L-O-R-D.
兩者的混合.
同時據報導.
這些經文和七十四億本的經文有些相差.
這裡就表示.
希伯來聖經的抄本傳統.
其實不是單一的.
是多元的.
可能有幾個不同的傳統.
平時我們用的就是馬索拉抄本.
就是Masoretic Text.
是屬於其中一個大的傳統.
但是其實我們在二十世紀.
找到死海古卷之後.
已經明白到多元抄本傳統的現象.
無論如何.
這個由大曠野洞穴的發掘工作.
是在1960年至61年的時候已經開始.
所以以色列的考古人員.
已經在這一帶逐個山洞搜索過.
而且這些洞穴找到的考古發現.
都是代表不同的考古歷史時期.
譬如最遠古的就是Calgolithic Age.
就是公元前4500至3500年的時間.
這個時間是剛剛人類開始用銅器.
然後另外一個時期就是Iron Age.
就是1200至560年左右.
然後死海古卷的時間就是羅馬時期.
公元前一世紀至公元後.
Late Roman Period.
就是到三百多年的時候.
在這些由大曠野的洞穴裡面.
其實都經常找到第二次猶太人革命的文物.
譬如一些硬幣 箭頭等等.
很多年來以色列都已經在這一帶山崖那裡.

$^{161}$做一些很有系統的工作.
每一個洞穴再回去.
然後徹底地去搜查去發掘.
是非常艱辛的工作.
因為要進入這些洞穴不單止危險.
要攀山繩才能進入.
亦都很細小和非常嚴熱.
所以因為這個緣故就非常艱辛.
但今次這個發現就是一個很大的成績.
代表著這一項工作是應該繼續進行的.
因為裡面的洞穴不知道還有什麼.
其實到現在來說.
由大曠野山崖上的洞穴.
只有一半已經進行了徹底發掘的工作.
即是還有50\%要做的.
所以仍然有很多工作要進行.
還有沒有其他的發現呢?.
是有的.
第二個大的發現就是在報章上經常說的.
就是一個小朋友的屍骨.
你見到那些骨頭都是比較短的.
所以肯定是一個小朋友.
未知他是男還是女.
但那個鞋骨是用布包裹的.
而且那個屍體是曾經經歷過某些處理.
仍然保留著一些頭髮和表皮.
可能是這個緣故.
有些傳媒就會說這是木乃伊.
無論如何由大曠野這個地方.
其實是一個非常熱和乾燥的地方.
正因為這個緣故.
所以是一個考古的天堂.
因為所有放進去的東西.
都是進入了一個很理想保存文物的環境.
說回這個兒童的鞋骨.
有6000年歷史這麼久.
即是你推算公元前4000年的時間.
那段時間就是考古歷史裡面的銅石器時代.
4500至3500年的時候.
如果你將這個時期對比整個聖經考古時期的話.

$^{201}$我這裡為你預備了.
這個時間就是銅石器時代的時候.
當時的人開始用銅器.
開始掌握如何煉銅.
所以這個不是太驚訝的事.
在猶大曠野的山洞裡找到.
Calaculic Age的Burial.
找到一些墳墓或安葬屍骨的地方.
為什麼呢.
因為1961年去探索這些山洞的工程.
在南邊的一個溪谷叫Nahal Mishmar.
它在哪裡呢.
你會看到昆蘭在這裡.
剛才我們說的恐怖的洞穴就在這裡.
然後隱居底在北邊一點點.
死海旁邊.
另外就是我現在所說的.
就是Nahal Mishmar這個地方.
這裡有幾個山洞.
幾個山洞裡面都找到不同的屍骸安葬的地方.
也是這個時期的.
就是4500至3500年.
Calaculic Age.
就是這個同石器同庸的時代.
所以在這裡所謂Cave of Horror.
這個恐怖的洞穴裡找到另一具屍骨.
其實都不是驚訝的事.
因為已經找到二十多具屍骨.
在不同的山洞裡.
這是其中一個例子.
也是一個小朋友.
葬在剛才我說的.
Nahal Mishmar這個溪谷.
其中一個山洞裡.
你看到這個圖表.
告訴你其實有骷髏骨頭的地方.
都是找到葬屍體的地方.
這裡下面有.
並且在上面.
平原上面.

$^{241}$礦業平原上面.
還有北邊一點點.
有點的地方就是一些Cave Site.
就是一些山洞的地方.
所以你看到有很多的山洞.
並且這些山洞可能有很多是用來.
藏寶或者在這個Calaculic Age的時候.
是用來安葬屍體的.
在這個Nahal Mishmar這個地方.
也找到很多的寶藏.
為什麼找到寶藏的山洞呢?.
考古學家叫做藏寶山洞.
Cave of Treasures.
找到什麼樣的寶藏呢?.
就是找到很多銅器.
剛才我不是說.
Calaculic Age就是開始.
去練銅.
懂得去掌握這個技術.
在這個Nahal Mishmar的一個山洞.
當時剛剛進去的時候.
就是拍到這張黑白的照片.
這些銅器堆在樹上.
藏在那裡.
洗乾淨 整理了.
就變成了這些這麼漂亮的銅器.
所以這些是挺驚人的.
因為在這個山洞裡面.
總共找到442件的銅器.
在伊斯利福博物館也可以看到.
好像銅像 銅冠.
或者一些珠鏈等等.
都是用來做裝飾的.
手工非常精緻.
好像絕無僅有.
因為在這個時期.
Calaculic Age正正是剛開始鑄造銅器.
都沒有找到這麼大批的鑄銅器皿.
你記得勝經裡面曾經說過.
打造銅器的開始.

$^{281}$有沒有記得.
這個創世記的第四章.
尾的時候就說.
土伯該人就是打造各種銅鐵器皿的匠人.
無論如何.
在猶太曠野.
這個藏寶山洞.
K of Treasures.
Mishmash這裡找到的.
似乎告訴我們.
背後可能有一個人.
或者有一群人.
他們真的非常懂得.
在當時掌握的技術.
才能夠做到這麼多這種銅器.
這裡也告訴我們.
在這一區猶太曠野.
很可能在銅器時代的時候.
有一些游牧民族.
或者一些族長.
他們就用這些山洞來藏寶.
或者用這些山洞.
都是用來做一些墳墓.
代表一個古代的游牧民族.
在這個地帶非常活躍.
其實我們也有其他的數據可以連接一起.
譬如在隱磯底.
大家記得在哪裡.
隱磯底就是在死海的西岸.
很多時候人們去到以色列的時候.
都會逛逛自然教育徑.
有很多很美的景色.
因為那裡有水泉.
但很多時候我們沒有留意.
在上面很少人會去的.
就是一個神廟.
就是在教科學時代的時候.
一個露天的神廟.
這個考過家多年來已經找到了.
已經找了幾十年.

$^{321}$就是一個露天的神廟.
有兩個大的建築物.
中間有一個類似石造的水盤.
建成用的地方.
這裡可能是小一點的.
可能是一個workshop.
或者一個supply room.
儲藏室.
但是那個大的長的.
就值得留意了.
因為中間有一個.
一進去.
圓形的一個地方.
可以放一個偶像的地方.
這個就是重構圖.
剛才看到一個大的長形的室.
所謂的century.
裡面有一個可能是放置偶像.
去拜祭的地方.
這個就是教科學時代.
當時這一帶的人.
他們宗教上的一些表達.
我們有一些學者.
就這樣去推斷.
當時教科學時代的時候.
或者一些遊牧民族.
他們有時候會聚在一起.
聚居在這個地方.
可能在死海的岸邊.
他們就會在那裡聚集.
或者在那裡暫時居住.
然後再上一些.
在Angeli Springs.
有水泉.
他們就建立了一個.
他們共用的神廟.
在那裡有宗教的儀式.
宗教的活動等等.
再上一些就是.
這些黑色的地方.

$^{361}$就是一些小的山洞.
我們說如果搜索這些山洞的時候.
很大可能也會繼續找到.
一些墳墓.
或者一些藏寶藏的地方.
不知道會找到甚麼.
怎樣也好.
我們想用一個地帶的方法.
很多時候看這些考古的新聞.
我們會留意.
在那一點.
在那個地方找到甚麼.
而我們沒有去到所謂connect the dots.
想想在那個地區裡.
到底有甚麼共通的地方.
剛才說了Calculific age.
所以說找到一個很舊的.
Calculific age的小朋友的屍骨.
如果你知道其他地方.
都有同樣的情況.
你就知道原來這一帶.
可能是指標著一些民族.
當時用這個地方.
可能在這個地方是非常活躍的.
第三個重要的發現是甚麼.
就是這個草籃子.
這個更加古舊.
計算出來的時候.
是10,500年這麼舊的大籃子.
你會問會不會有可能.
這麼久.
如果你運算出來的時候.
這個用草織成的籃子.
是公元前8,500年.
有可能的.
為甚麼呢.
正如我所說.
又大礦的地方又乾又熱的地方.
所以任何的古物.
文物放在這裡都很容易得到保存.

$^{401}$這個草籃子其實都是一樣.
你會看到在泥土慢慢揭開的時候.
就找到這個草籃子.
是一個新石器時代.
因為公元前8,500年.
是屬於這個新石器時代.
我們看回剛才我們看了銅石器時代.
之前那個時代.
8,500年至4,500年.
就是這個新石器時代.
新石器時代.
新石器時代到底發生了甚麼事.
最大的改變就是.
人類開始有耕種.
人類開始懂得如何播種.
以至能夠建立到以農耕為中心的城鎮.
所以這個大的草籃.
可能都是跟農業有關.
所以在這個區域.
其實由新石器時代開始.
我們已經看到有些人開始聚居.
以及有農業的活動.
這些都不是很詫異的事.
為甚麼呢.
因為你再回到這個區域.
又是用區域的方向.
你會看到有些甚麼.
就是當時有些很古舊的城市.
其中一個最古舊的當然就是.
耶利哥.
因為耶利哥是世界上最古舊的一個城市.
因為找到土丘考古的土層.
可以追溯至公元前一萬至八千年的時候.
這個草籃在哪裡找到.
就是在這個cave of horror找到.
所以同期的時候.
到底大家有甚麼聯絡呢.
就是講到這個新石器時代到底發生了甚麼事.
有草籃.
有城鎮的建立.

$^{441}$有農業的開始.
你看到整個地區這樣看的時候.
你會發現原來不是一些孤立的發現.
而是跟整個region.
一個歷史上不同時期的一些演變.
無論如何.
現今來說這個草籃還有蓋子.
是世界上最古老的一個草籃.
是非常重要的發現.
希望告訴大家.
用甚麼方法去聽這些新聞.
以及用甚麼角度去分析.
希望今次這個regional approach.
給大家能夠connect the dots.
讓大家聽這些考古新聞的時候.
能夠多一點資訊.
明白當時的世界發生了甚麼事.
說不定也可以跟聖經做一些連接.
希望下一次和大家再分享.
\newpage



\section{}
\label{sec:neXNn__inUY}
\textbf{舊約網上行 - 上帝的足跡:EP7 試煉}
\newline
\newline
連結: \href{https://youtube.com/watch?v=neXNn--inUY}{\texttt{https://youtube.com/watch?v=neXNn--inUY}} ~~~~ 語音日期: 2021-03-23
\newline
\newline
\hyperref[sec:uhU9sAnkAzo]{\small{< < < PREV SERMON < < <}}
~
\hyperref[sec:index_chronic]{\small{[返順時目]}}
~
\hyperref[sec:index_scriptual]{\small{[返順卷目]}}
~
\hyperref[sec:MgtgOAT6lXs]{\small{> > > NEXT SERMON > > >}}
\newline
\newline
$^{1}$上一集我們跟隨阿布拉罕的足跡.
我們看到上帝呼召他.
然後他順命一家離開自己本地本族本家.
未索不達米亞.
離開吾爾哈蘭.
終於來到迦林地.
到了迦林地的時候故事是這樣發展的.
他經過不同的地點.
士堅,伯特利艾中間和南地.
在不同的地點他也建立祭壇宣揚上帝的名.
《創世記》十二章繼續這樣記載.
第十節 當時那地發生饑荒.
阿布蘭就下到埃及去.
在那裡寄居.
因為那地的饑荒十分嚴重.
迦林地我們知道是靠雨水來生存.
所以一年兩年幾年缺乏雨水的話.
就會造成旱災農作物失收.
但這裡有一個問題.
古時的時候發生的旱季.
是不是可以有跡可尋.
到底靠什麼數據.
用什麼方法讓我們知道.
以前可能曾經有不同的時段.
雨水不夠所以造成旱季.
答案就是花粉.
我們知道雨水充足的時候.
那些木本植物.
即是大樹,高樹就會繁盛.
而且他們的花粉就會更加多.
但如果到旱季缺乏雨水的時候.
這些木本植物的花粉就會相繼減少.
並且矮灌木因為他們需要的水份比較少.
所以矮灌木的花粉自然會增加.
過去多年其實有些學者不斷做研究.
來收集這些花粉的數據.
譬如在北面加里海地區附近的.
一些古代花粉的數據.
還有南面死海一帶.
搜集了古時花粉的多少,數量等等.

$^{41}$在2014年的時候.
有一篇論文來總結了出來.
你看到這個圖表是由下至上這樣看的.
由下至上即是公元前2500年至1500年之間.
這1000年的花粉的數量是有多也有少.
那條灰色線代表木本花粉在不同時間的數量.
木本花粉即是大樹的花粉越少.
就代表雨水越少.
灰線就會向左.
如果花粉越多的話.
就代表雨水越來越多.
灰線也會向右.
你可以看到在公元前2000至1800年之間.
確實有下降的跡象.
即是在這200年裡面.
雨水不充足.
以致可能會造成農作物失收,旱櫃.
我們提過在創世記續漲的時期.
阿巴拉罕,以撒,瓦各.
大約在公元前2000至1500年之間.
在這段時間的初期.
你看這個列表的時候.
你就知道迦南地確實有旱櫃的出現.
所以可以看到聖經的記載是相當之寫實.
所以這次我們跟隨阿巴拉罕的足跡.
繼續看他的故事有什麼發展.
創世記12漲之後的故事.
我們說有饑荒,有旱災在迦南地出現.
於是阿巴拉罕發生什麼事呢?.
12漲第11節就交代了.
快要到達埃及的時候.
阿巴蘭對他的妻子撒萊說.
看啊,我知道你是個美貌的女人.
埃及人看見你的時候必會說.
這是他的妻子.
他們就會殺我.
卻叫你活著.
請你說你是我的妹妹.
使我因你的緣故可以平安無事.
我的性命也可以因你的緣故得以保存.

$^{81}$果然,阿巴蘭進入埃及的時候.
埃及人就注視那女人.
因為她十分美麗.
法老的一些大臣看見撒萊.
就在法老面前稱讚她.
於是那女人被帶進法老的宮內去.
阿巴蘭因這女人的緣故就被休待.
得了許多牛羊,公牢,母牢,瀑皮和駱駝.
又說因為阿巴蘭妻子撒萊的緣故.
就用嚴重的災病打擊法老和他的全家.
於是法老把阿巴蘭照料下來對她說.
你向我作了什麼?.
你為什麼不告訴我她是你的妻子?.
你為什麼說她是你的妹妹.
以至我娶了她作妻子?.
現在你的妻子在這裡.
帶她走吧.
法老為了阿巴蘭的事,吩咐神父.
他們就把阿巴蘭和他的妻子.
以及他所有的一切都送走了.
是一段很有趣的故事.
這裡是交代阿巴蘭一家.
到了埃及所謂大能不死.
從埃及回到迦南.
並且繼續走的時候就到了南地.
十三章第三節又說到他的行程.
他從南地一站一站往前走.
到北面的伯特利.
到了伯特利和艾中間.
就是他先前去支立帳棚的地方.
也是他竹頭一座祭壇的地方.
阿巴蘭在那裡求告耶和華的名.
行程就說完了.
這個故事說到迦南地有旱災.
阿巴蘭要逃難.
就到了埃及.
埃及的氣候地理當然是非常不同.
因為有利羅河不是靠雨水.
所以很多時候在迦南地有旱災的時候.
很多人都會去埃及.

$^{121}$阿巴蘭也不例外.
他到了埃及的時候就發生了一件事.
他心裡面有一個策略.
就是不如說他太太其實是妹妹.
讀這段經文的時候.
無論如何讀都覺得很費解.
因為我們自然的反應就是這樣也可以.
想保平安當然很明白.
但是要保自己的命.
可以有很多方法.
可以首先不去埃及.
或者留在迦南地.
應該也死不去.
又或者就算去埃及也好.
可不可以花點錢.
多給人一些禮物.
這樣別人就不會搞你的家.
最令人費解就是.
為什麼可以叫自己太太做妹妹.
這樣就可以過關呢.
太太當成妹妹.
其實類似的情節在創世記出現了三次.
這些情節很清楚表明.
聖經不是直接寫給我們的.
因為我們不知道當時的背景和文化.
經文首先去到交代.
就是說撒萊就是美貌的女人.
美貌也不一定是我們現代人心中所想的女神.
也可以是雍容華貴的意思.
撒萊當時應該有六十五歲.
所以就算看美貌.
我們也不一定用現代人去看美貌.
我們繼續問.
有什麼背景可以幫助我們理解一下.
阿巴拉罕的行徑和策略.
我覺得起碼有兩方面可以提出來.
第一就是一些古代的文件.
第二就是一些現代人類學的研究.
第一方面.
一些古代的文件.

$^{161}$告訴我們當時的城邦或者軍王.
互相有邦交.
而邦交往來.
很多時候是靠外交婚姻的習俗.
很多時候邦交結盟的方法.
就是一個國與另一個國去結親.
一個王或一個族長派一個女兒或公主.
和另一個邦國去結親.
結婚當中一定會有嫁妝禮金等等.
於是兩個王就聯盟.
舉一個例子.
大約在公元前1500年的時候.
巴比倫王與埃及王之間有一些外交的書信.
這些書信是一些埃及宮廷外交的書信.
在埃及一個首都叫做阿嗎哪找到的.
有很多這些泥板都是交代當時一些外交的關係.
其中一封書信就是阿嗎哪泥板.
我們簡稱EA.
發言的地方就是阿嗎哪這個地方.
第四封.
怎麼說呢?.
這封書信就是巴比倫王.
就是中東的另一邊.
軍王投訴埃及王.
說埃及王不願意將公主派去給巴比倫.
巴比倫已經送了公主去結親.
但是埃及王不禮尚往來.
所以影響了兩國的交情.
這封書信說什麼呢?.
巴比倫王說.
我的兄弟還沒有埃及的公主被送過來.
總之你給我一個美女.
就將她當成你的女兒也好.
為什麼我的兄弟遲遲都不送上一個女人.
就是一個公主呢?.
我是不是應該像你一樣.
強留公主不給你呢?.
我當然不會.
我的女兒就在我這裡.
我不會留下一個不給你的.

$^{201}$這裡你會看到一個很有趣的外交文件.
巴比倫王投訴埃及王.
說他不遵守國際的條例.
一些國際法.
他有他的做法.
就是不派公主來.
我們知道這個信件之後的結果是什麼呢?.
就是埃及王始終都沒有派公主去.
只不過是給他一些禮物.
你可以看到當時的邦交.
很多時候都是靠外交的訓人.
派一個公主去.
或者派一個女兒.
或者派一個妹妹去.
你就明白了.
亞伯蘭為什麼會將他太太.
說成是他妹妹呢?.
因為太太當然不能夠結親.
只有妹妹才能夠結親.
這是第一方面.
是一些古代的文件能夠給我們一些背景.
第二方面是如何看到當時古代的習俗呢?.
就是一些人類學的研究.
有一個人類學研究.
就是訪問一些中東的民族.
中東的傳統的一些習俗.
其中一個故事就是說.
去到百里行南邊一個小村叫做阿特斯.
訪問了一個女人.
那個女人嫁給了一個阿拉伯人.
當時是第一次世界大戰爆發的時候.
在中東這一帶.
政局非常不穩.
這個家庭就舉家離開自己的村莊.
搬到約旦河東一帶.
住在一些遊牧民族中間.
那個丈夫就叫這位太太.
不如妳是我的妹妹.
我們以兄妹相稱.
因為不同的原因.

$^{241}$這個丈夫在遊牧民族中間生活的時候.
得罪了當地人.
當地人對這個女人這樣說.
他就警告她.
他說如果這個男人是妳的丈夫的話.
我們早就將他殺死.
因為他是不懂世情.
這個女人也聽到當地的族長彼此商量的時候.
就說因為這個男人是她的哥哥.
所以我們就放過他.
如果他是她的丈夫的話.
我們早就將他殺死.
但人類學者想明白.
這些古代的習俗.
原來在阿拉伯的社會中.
在傳統的文化中.
一個女人的兄弟.
比她的丈夫更加重要.
你會問為什麼.
因為在這些傳統的遊牧民族的文化中.
丈夫死了.
女人可以找到另一個.
但是親生的兄弟.
是天賜的恩典.
是不能夠取締的.
就是brother是irreplaceable的.
死了就不會再有.
因為這種觀念.
世世代代在遊牧民族.
或者傳統的中東文化.
都會保存著.
或者保護著一個人的兄弟.
如果有這些背景的時候.
我們就可以了解.
阿巴罕為什麼到了埃及的時候.
他怕死.
他立刻想到這個策略.
無論如何.
這個故事的交代就是.
魏蘭到了.

$^{281}$饑荒有打擊來到的時候.
到底如何去自處呢?.
看到阿巴罕如何去應對.
留意阿巴罕自己在言談之間.
他顯露出自己心中的動機是什麼.
十二章第十一節.
快要到埃及的時候.
阿巴蘭對他的妻子撒內雪.
看啊.
我知道你是美貌的女人.
埃及人看見你的時候必定會說.
這是他的妻子.
他們就會殺我.
並叫你活著.
請你說你是我的妹妹.
使我因為你的緣故可以平安無事.
我的性命也可以因你的緣故得以保存.
聖經的敘事手法我們要留意.
就是很多時候會說到.
人物有直接的對話.
留意一下這裡的記載.
就是聖經第一次記載阿巴罕直接說話.
好像一套電影的主角.
第一次開聲有對白的時候.
希伯來的故事很多時候是靠著對白.
去透露出人物的動機和心態.
你看到這裡所強調的就是他自己.
一種自私一種自我的想法.
無論如何是一種軟弱脆弱的時候.
再看法路回答.
十二章第十八至十九節.
於是法路就叫阿巴罕來.
對他說你對我做了什麼.
為什麼不早早告訴我他是你的妻子.
為什麼你說他是你的妹妹.
以至我娶了他做妻子.
現在就有難臨頭了.
現在你的妻子在這裡.
帶她走.
阿巴罕有什麼回應.

$^{321}$這裡就是沉默無言.
所以無論如何.
都是阿巴罕軟弱.
阿巴罕錯在先.
創世記十二至五十章.
在族長的故事裡面.
最獨特的地方就是.
當記載以前以色列民族.
以色列國的長老.
他的始祖原老歷史的時候.
你會發現.
不只是說他光輝的地方.
而是切切實實地說出他的缺點.
他所做錯的事.
這一段太太當為妹妹的故事.
在創世記的族長故事裡面.
出現了三次.
第一次這個情節發生.
就是十二章第十至二十節.
就是阿巴罕來到埃及的時候.
然後第二次記載一個這樣的情節.
就是二十章.
就是阿巴罕沒有到埃及.
這次好像好一點.
去了基拉爾.
就是加拿大南地的其中一個城邦.
第三次就是講到以色列.
同樣發生的地方也是在基拉爾.
這次可能更加好.
因為他是順命留在加拿大.
二十六章一至三節.
有一個這樣的交代.
就是從前阿巴罕曾經有過一次饑荒.
阿弟又有饑荒.
以撒到了基拉爾的那裡.
到了非利士人阿比米納王的那裡.
耶和華向以撒顯現.
他說你不要下到埃及了.
要住在我所指示你的地方.
你所居住的地方.

$^{361}$我必定與你同在等等的說話.
我們看到三個故事.
都是記載著以色列國的國父.
阿巴罕 以撒.
這些應該是最威風的人.
誰知說出他們有血有肉.
會有脆弱的時候.
這就是聖經最獨特的地方.
你試試看一些國家.
無論是古代還是現代.
這些國家官方的歷史.
一定是有讚無嘆的.
以免傷害國情.
不單止是這樣.
這些開國功臣.
很多時候在歷史著作和記載裡面.
都是一些非凡的人物.
未必一定會去到神級.
但都是與我們平常人不同.
是一種英雄的人物.
並且很多時候將他們偶像化.
國家的元首 總統 社會的領袖 元老.
甚至很多時候我們牧者.
我們很容易將他們偶像化.
放在一個非凡的層面上.
結果這些領袖倒下的時候.
我們就不知道該怎麼辦.
可能我們對神的信仰也會受挫.
或者我們追求的方向也會搞錯.
甚至我們可能會怪責上帝.
聖經有這個獨特的地方.
記載這些國父.
所謂以色列的神子民的族長.
是很真真實實地去記錄.
不是說元老有多厲害.
不是說這些創始人有多厲害.
他們只不過是和我們一樣.
《篤刀創世記》記載了太太和妹妹的故事.
三個故事來看.
十二章 二十章 二十六章.

$^{401}$其中一個最特別的地方就是.
這些外邦的王發現了真相之後.
就理應將這個族長殺死.
但是他們沒有.
比如十二章 法老放阿布蘭走.
二十章十六節 另一個故事.
亞比米勒王甚至給了一些遮羞的錢.
給他一些禮金.
就放他走.
二十六章都是一樣.
外邦的王反而保障了二十.
因為什麼緣故呢.
不是因為人是英雄.
不是因為他厲害.
而是耶和華的恩典.
人是無可掛的.
《創世記》十二章第十七節.
剛才我們看過了.
再看一次.
是耶和華因為阿布蘭妻子殺害的緣故.
去保護她.
緣故這個字可以有不同的翻譯.
可以是說話.
如果翻譯成說話的話.
可能是殺害的哀求 她的禱告.
或者可以更籠統的翻譯.
就是事情.
就是耶和華緊張去保護著殺害.
所以她做了一些事情去拯救他們.
有一個恩典的拯救.
阿布蘭落埃及避難.
神又帶她出埃及.
這個絕對是聖經裡面第一次提到出埃及.
或者有種預表.
將這個故事寫成.
好像以後在出埃及發生的.
以色列人出埃及的情節.
有相似的地方.
譬如這個饑荒的緣故.
落埃及避難.

$^{441}$耶和華降災打擊法老傳教.
這個災字也在出埃及的十災裡面出現.
你想想.
其實兩個故事好像有相似的地方.
也有不同的.
因為後期出埃及所記載的.
那一次法老是一個暴君.
他心硬.
所以上帝用恩典去拯救以色列人出來.
但是阿布蘭這裡.
第一次出埃及就不同了.
這裡明明是阿布蘭犯錯.
是他先錯.
是他先在先.
雖然有不同的地方.
相通的地方你應該可以看到.
就是上帝恩典的拯救.
無論是人的軟弱也好.
是因著上帝的守弱也好.
都是上帝願意加入拯救.
我相信這個也是故事裡面的重點.
不是說到族長有多威風.
而是真正想說出.
我們神是那位信神和守弱的神.
我們人生的信仰道路上.
我們很多時候都會有跌倒的地方.
我們其實都好像族長一樣.
他其實和我們一樣.
都是一些凡人.
是有血有肉.
其實這樣去看這個故事.
反而對我們一些脆弱的人有幫助.
因為我們跟隨上帝的時候.
跟隨神跟隨耶穌的時候.
我們不只是拍胸口去強化我們自己.
我們可以這樣做.
但我們知道最後都會軟弱.
我們會跌倒.
聖經裡面告訴我們.
恩典或者福音.

$^{481}$是真正告訴我們.
我們本身就是有血有肉.
會有脆弱的時候.
而上帝的恩典.
上帝的福音是什麼?.
就是不如面對我們這個軟弱的時候.
我們的起點就不再是.
我們自己行不行.
或者是我們自己嘗試再努力一點.
而是我們的重心放在哪裡?.
我們的起點放在哪裡?.
就是耶和華的恩典.
原來只是上帝有能力來供應.
上帝有能力來保守.
當我們更加去認識上帝.
體會到上帝的恩典的時候.
我們才可以在我們的試煉當中.
我們明白如何自處.
如何去面對.
你看十二章這裡的結構.
其實那個中心都是說.
如何去面對試煉.
你看到十二章一至三節.
就是耶和華的呼召.
叫他離開 去.
叫亞伯拉罕離開本地.
半族半教.
第二段就是四至九章.
就是那種的movement.
就是亞伯拉罕離開他自己本地.
去到迦南地.
然後第十節就是那地有饑荒.
並且饑荒非常嚴重.
那個就是一個轉折.
然後十一節到十七節.
也是另一個的movement.
就是在迦南地到埃及.
十八至二十章又出現同一個字.
開頭出現的.
就是走或者離開.

$^{521}$去這個字.
這三文字的結構.
經常在《創世記》出現.
第一段我們看到.
就是耶和華吩咐亞伯拉罕離去.
你要走.
然後第二段我們看到.
又是一對的.
是並行的.
上面就說到要進入迦南地.
然後下一段就說到到埃及.
中間就是最重要.
轉折的地方.
就是有饑荒有試煉來到.
如果將每一段的內容.
總結一下出來的時候.
你就更加清楚.
十二章一至第三節.
就是亞伯拉罕需要成為萬國的祝福.
所以上帝吩咐他.
你要離開.
四至第九節.
亞伯拉罕就相信上帝.
只靠上帝報往迦南地.
然後中間的時候.
亞伯拉罕就面對試煉.
就問一個問題.
到底是否願意繼續信靠上帝.
結果一個情節發生.
就是亞伯拉罕不信上帝.
十一至第十七節.
他就落到埃及.
然後十八至二十節.
總結的時候.
亞伯拉罕就成為了萬國的咒詛.
中心轉折的地方.
是很明顯的.
就是危機試煉來到的時候.
到底你怎樣回應.
這一章的中心.

$^{561}$是告訴我們面對試煉的時候.
我們是否願意相信神的話.
相信上帝的應許.
還是只相信自己的謀略.
自己自編自導自演.
或者一意孤行.
上一集我們開始亞伯拉罕這個故事的時候.
我們就說聖經告訴我們.
這個信心的旅程.
是一個探索.
一個quest.
在這個探索當中.
我們經歷上帝的恩典.
特別是危機在我們軟弱脆弱的時候.
我們更加發現上主給我們的恩典.
安排和帶領.
以致我們會不斷地改變.
很多時候我們叫亞伯拉罕做信心之父.
我們很多時候將主角放在人身上.
說他有多厲害.
說他有多信靠上帝.
這個可能也是對的.
但只不過是一半.
或者我們應該叫亞伯拉罕.
做經歷上帝恩典之父.
因為在他自己碰碰撞撞的當中.
他經歷了神自己的恩典.
神自己的報仇.
他這樣看的時候.
我們更加明白.
原來在我們探索的旅程上.
真正的主角就是上帝自己.
求上帝在我們屬靈的歷程中.
在我們的quest當中.
這個探索當中.
讓我們更加豐富地體會上帝自己的恩典.
以致我們的生命真的宣揚.
不是我們自己的信心有多強.
而是上帝是多麼信實.
多麼有能力,多麼有恩典.

$^{601}$主在這個信心的歷程中.
我們去感恩.
因為從你自己呼召我們開始.
我們就有這樣的恩典來跟隨你.
求你親自幫助我們.
因為我們的生命當中.
總會遇見饑荒.
總會遇見危機和試煉.
而且我們知道.
本著我們有血有肉的人.
我們實在很多時候是不信的.
但我們知道你仍然是信神.
我們知道你的恩典仍然是足夠的.
求你繼續幫助我們.
在我們跟隨你的腳蹤裡.
在這個信心的歷程當中.
我們愈來愈見到上帝自己是值得我們信求.
見到上帝你自己的榮耀.
你自己的恩典.
你的能力保守.
叫我們的生命真的宣揚你自己的名.
不是因為我們有多厲害.
而是本身我們就是一個脆弱的人.
我們需要依靠你.
叫我們在跟隨你的探索路程當中.
我們更願意讓人見到.
上帝你自己的恩典.
上帝你自己的偉大.
奉主名求.
阿們.
\newpage



\section{}
\label{sec:MgtgOAT6lXs}
\textbf{考古探秘 (聖墓堂) -  陳崇基牧師 Mark Sir (2021.03.31)}
\newline
\newline
連結: \href{https://youtube.com/watch?v=MgtgOAT6lXs}{\texttt{https://youtube.com/watch?v=MgtgOAT6lXs}} ~~~~ 語音日期: 2021-04-01
\newline
\newline
\hyperref[sec:neXNn__inUY]{\small{< < < PREV SERMON < < <}}
~
\hyperref[sec:index_chronic]{\small{[返順時目]}}
~
\hyperref[sec:index_scriptual]{\small{[返順卷目]}}
~
\hyperref[sec:_gtM1xmdWw0]{\small{> > > NEXT SERMON > > >}}
\newline
\newline
$^{1}$我們到了受難節和復活節.
在基督教最重要的節日裡.
我們紀念耶穌基督在十字架上受刑受苦.
然後受死埋葬第三天復活.
想起耶穌基督在人類歷史裡.
為我們受死和埋葬復活的事蹟時.
我們就一定想起聖母堂Holy Sepulchre.
你去過耶路撒冷舊城這個地方.
你相信一定去過這個景點.
叫聖母堂.
我們進入聖母堂這個地方時.
在庭園的樹木.
我們從南面的入口.
因為當時十字軍的時候.
將入口改成向南.
大部分你走到Holy Sepulchre.
看到的都是十字軍留下的.
當然大家知道.
這個教堂所蓋著的考古遺跡.
不是十一二世紀開始的.
而是更早甚至是到第一世紀.
所以說到耶穌基督受難,受死,埋葬和復活.
就是指這個地方.
就是這麼大的古教堂所蓋著的地方.
你走進這個教堂.
好像有一點感覺很容易迷失.
因為好像一個迷宮.
我們從南面按照這個地方地圖走進去.
你見到進入這個入口.
你轉右的話.
就上到這個叫做Chapel of Calvary.
這裡就是加勒山的禮堂.
這裡是右邊是天主教.
左邊就是Greek Orthodox.
希臘東正教的禮堂.
這裡就紀念耶穌基督在十字架上受刑,受死.
然後在這個禮堂下來的時候.
我們就會見到一塊很大的長方形的石.
這一塊石就是對著南面的門口.
很多東正教和天主教的人進來朝聖的時候.

$^{41}$就會親手去碰這一塊石.
因為這塊石紀念耶穌在十字架被師生帶下來的時候.
然後在這塊石上預備他的屍體.
然後再過去左邊.
另外一個大的禮堂就是聖母堂最中心的禮堂.
一個圓頂的Rotunda.
一個圓頂紀念的教堂.
裡面當然大家知道有一個Elegy.
裡面就是標記了耶穌基督的埋葬和復活的墳墓的地方.
當然這個教堂裡面所紀念的,所標記的.
不單止是耶穌基督的死而復活那麼簡單.
也埋藏了很多教會歷史的遺跡.
如果你從聖母堂進去的時候.
你轉右經過一個通道.
一直走下去就會到Helena Chapel.
Helena就是君士坦丁·Constantine的媽媽.
故事說到主後四世紀三百多年的時候.
皇母Helena來到聖地這裡尋訪耶穌的事蹟.
在不同的地方建立一些教堂.
在這個Chapel裡面.
你會看到走下去左右這些石牆上面.
都刻了很多小的十字架.
以往一直都認為這些小的十字架.
就是中世紀的時候.
基督徒在不同的地方來到耶路撒冷朝聖.
有一個聖母堂的時候來到聖母堂.
他們就所謂到此一遊在這裡做一個紀念.
然後在石上畫上這些十字架.
以往我們一直都認為這些朝聖者來到聖母堂.
他們就自由地按照自己的心意.
在石牆上畫上十字架作一個紀念.
但是在2018年的時候.
有一個考古的研究.
這個研究是做什麼的呢?.
就是在這個地方將這些十字架.
用一些高解像的攝影技術來紀錄它.
就是拿了這些刻在石牆上的十字架的所有數據.
發現在聖母堂裡面找到的刻上的十字架.
原來不是很隨機的.
就是各人有自己雕塑的方法.

$^{81}$因為你將這麼多個十字架來做一些統計.
拿了它物質上的數據的時候.
就發現原來這些十字架都是頗一致的.
長,闊,高或深度的雕刻都是頗一致的.
所以這些數據去到組成之後.
就認為只有三至四個或幾個雕刻的人的手筆.
所以你看到這個圖有些十字架.
都是刻在石牆上.
紅色這條線你見到右邊的圖表.
就是告訴我們到底它的高度有多少.
還有深度等等不同的尺寸.
藍色這條線也是.
你見到紅色藍色都是頗相似的.
所以是頗一致的.
下面這個圖我們見到是一個綠色的圖表.
就是將這些十字架用一些高解像的攝影.
拿了它的數據之後就計算在一個圖表上.
你見到綠色這條線就是這麼多個十字架.
總共在研究裡面就拿了61個十字架.
拿了它的橫切面.
發現它們的尺寸和深度都是相似的.
所以似乎不是自己朝聖者刻這些十字架.
而是可能有些專人刻這些十字架.
甚至這個可能是告訴我們當時中世紀的信徒.
來到耶路撒冷朝聖的聖徒經濟運作是怎樣的.
可能有些基督徒在歐洲或不同地方.
他們不能來到耶路撒冷朝聖.
但他們很想來.
但可能在他們本地的教堂.
他們就可能托人要求在耶路撒冷的聖母堂.
這個基督教裡面最神聖的地方.
幫他們畫上一個十字架.
這些都是一些中世紀.
或教會歷史留下的遺跡.
讓我們知道當時的歷史和背景等等.
其實歷代以來甚至到今天.
不同的宗派都會很寒蟲去聖地.
去聖母堂朝聖.
正如剛才所說.
有些天主教或東正教徒.

$^{121}$進入聖母堂時會親手摸到這些東西.
這個當然不是根正教的傳統.
根正教的傳統很多時候沒有什麼實物.
都是比較屬靈化.
但在巴士坦地這個地方有一個特色.
無論講什麼信仰.
猶太教或基督教.
很多時候要拿一些很屬靈化.
即是親手摸到的東西.
一個地方是特別神聖.
在那裡祈禱是特別接近上帝.
例如去西牆 去聖母堂等等.
這些傳統都是一些所謂的聖母堂傳統.
一些朝聖的傳統.
當然對我們根正教來說.
雖然未必有這麼重的傳統.
我們去聖地的原因.
不是因為那裡特別靈驗 特別神聖.
而是我們最主要想做的.
就是去到聖地 巴士坦地.
了解聖經背景的歷史.
這個也是跟其他宗派有的朝聖傳統.
其實都有共通的地方.
因為我們相信我們的信仰.
是一個歷史的信仰.
Historic Christianity.
來到首蘭節和復活節的時候.
我們更加想起這件事.
但同樣去到新約聖經的時候.
我們更加找到新約聖經裡面.
一些更加革命性的概念.
就是不再將我們的信心放在地圖上某一點.
再不在地圖上某一個地方.
那裡是特別神聖一點.
特別是親近神一點.
因為聖經裡面告訴我們.
不再是在一個地方.
而是在基督裡面.
不再是在地球上某一個地方.
而是在聖經裡面.

$^{161}$我們就來到聖經裡面.
你記得約翰福音第四章的一句話.
是耶穌和撒瑪利亞婦人在井旁對話.
那裡說什麼呢.
那裡說事後將到.
那時你們敬拜天父不在這山.
不在Mount Gerizim.
撒瑪利亞人認為的聖山上敬拜神.
也不在耶路撒冷.
也不在猶太人以為最神聖的地方.
耶路撒冷聖殿的地方.
你們敬拜你們所不知道的.
我們所敬拜的我們所知道的.
因為救恩是從猶太人出來.
然而耶穌說什麼呢.
這一句就是一個更加革命性的話.
事後將到.
現在就是了.
那靠心靈按照真理敬拜父的.
才是真正敬拜的人.
不再是在哪一個聖山.
甚至是在耶路撒冷.
是在福音裡面.
是在基督的救恩裡面.
因為父在尋找這樣敬拜他的人.
神使靈敬拜他的.
必須靠聖靈按照福音的真理.
來敬拜他.
在受難和福祉裡面.
讓我們更加深明白.
耶穌基督是為我們所做的.
我們能夠在基督裡面敬拜神.
去到成為神的子民.
\newpage



\section{出埃及記}
\label{sec:_gtM1xmdWw0}
\textbf{考古探秘 (蘇伊士運河與出埃及記) -  陳崇基牧師 Mark Sir (2021.04.10)}
\newline
\newline
連結: \href{https://youtube.com/watch?v=_gtM1xmdWw0}{\texttt{https://youtube.com/watch?v=\_gtM1xmdWw0}} ~~~~ 語音日期: 2021-04-11
\newline
\newline
\hyperref[sec:MgtgOAT6lXs]{\small{< < < PREV SERMON < < <}}
~
\hyperref[sec:index_chronic]{\small{[返順時目]}}
~
\hyperref[sec:index_scriptual]{\small{[返順卷目]}}
~
\hyperref[sec:FvI2UPC3nfc]{\small{> > > NEXT SERMON > > >}}
\newline
\newline
$^{1}$又來到《聖經考古》的時間了.
接下來這幾集想跟大家分享一下.
在埃及的幾宗新聞.
這幾宗新聞跟《聖經考古》.
《聖經的背景世界》有些關係的.
第一宗新聞想跟大家分享的就是.
在2021年3月底的時候發生的.
話說有一艘大貨輪從紅海.
一直往北走到航行的時候.
進入了蘇伊士灣.
進入了蘇伊士運河.
據報導進入了運河的時候.
就遇上有沙漠的強風吹襲.
導致貨船打側.
擱淺了在蘇伊士灣南面的地方.
嚴重阻塞了世界海洋貿易的交通.
你知道這條運河多麼重要的時候.
你就知道什麼叫做阻礙地球轉了.
幸好擱淺了六天之後.
終於解困了.
多謝搶救的工人.
還有剛剛春天潮漲的水位比較高一點.
幫這艘大的貨輪能夠解困.
不會阻礙著整個世界的海洋貿易.
你問問這個蘇伊士運河.
跟考古《聖經的背景世界》有什麼關係呢?.
也有關係的.
先講一下這個蘇伊士運河.
其實是十九世紀中的時候才興建和啟用的.
全長是160公里.
是世界上最長沒有水閘的運河.
這條運河是博通上面北面地中海.
以及下面蘇伊士灣紅海的水域.
其實是歐洲和亞洲之間.
最重要的水運河.
沒有了蘇伊士運河的時候.
你就知道了.
這些貨船要繞過非洲的南面.
所以時間.
需要用的燃油等等都是大大的增加.

$^{41}$因為有蘇伊士運河的緣故.
這些時間都會大大的縮短.
第一點就是想講的.
就是蘇伊士運河.
還沒有興建之前.
其實歷世歷代都有人嘗試去挖一條運河.
去連接下面的紅海水域.
以及上面尼羅河三角洲地中海水域.
所以19世紀所興建.
所啟用的蘇伊士運河.
其實不是第一次有人想到的.
以往自古以來.
都有人想博通這兩個水域.
最早可以追溯至公元前18世紀.
當時古埃及是中王國的時候.
是Middle Kingdom.
有一位法老叫做St. Nusret.
他嘗試用一個尼羅河支流.
將它挖通.
這個尼羅河支流叫做Wadi Toumala.
你在地圖上看到很深色綠色的.
這個尼羅河支流.
就是叫做Wadi Toumala.
這個法老嘗試再將這個溪谷延伸.
以至到達大虎湖.
Great Bitter Lake.
然後下面有一個當時自然的水道.
將這些不同的水域.
不同的湖會博通.
以至到達下面.
紅海蘇伊士灣這個地方.
這個法老第一次嘗試這麼做.
當然當時的技術有限.
到了公元前的十世紀.
另一個法老又想做大工程.
這個法老就是Nicole.
是尼哥.
他想打通這個水道.
當然也不算成功.
所以因為這麼多的嘗試.

$^{81}$現在我們回頭看的時候.
我們就會形容這個所謂的工程.
這個運河.
就叫做縱法老的運河.
Canal of the Pharaohs.
因為當時古時的技術.
科技有限.
不能夠接通.
直到什麼時候才能夠接通呢?.
就是到波斯帝國.
Persian Empire的時候.
就是大利烏王Darius I.
他征服了埃及之後.
他下令將Wadi Toumala的運河.
再一次發掘.
然後就駁通尼羅河三角洲的水域.
以至地中海的水域.
以及南面紅海的水域.
大利烏王是波斯帝國的帝王.
挖通了運河之後.
他就立了幾塊石碑.
其中一塊石碑上面的碑文是什麼呢?.
這裡翻譯出來就是說.
大利烏王說.
我是波斯人.
從波斯出發.
我征服了埃及.
我下令在尼羅河挖出這運河.
接通波斯開始的海洋.
因此.
當這條運河按照我的命令被發掘開的時候.
船隻就按照我的意旨.
從埃及經過這條運河.
到達波斯.
如果你放大來看的時候.
你就知道這條運河對波斯帝國有多重要.
波斯帝國在歷史當中.
是第一個最大的帝國.
波斯帝國就是由印度的西北面.
一直延伸到今天的土耳其.

$^{121}$以前小亞細亞.
安納托尼亞這個地方.
包括埃及.
將這條運河挖通的時候.
因為波斯帝國的中心.
幾個的京都.
都是在現在的伊朗的地方.
舒山是其中一個.
佩瑟布拉斯是其中一個.
所以如果你看到他.
征服了埃及之後.
他能夠挖通這一個運河.
他的商貿海洋路線就會有利了.
因為他從波斯灣出發.
繞過沙地阿拉伯半島.
上去紅海北面.
他就可以進入尼羅河的三角洲.
以至能夠直通地中海.
所以這個運河第一次在古代裡面.
真的打通了.
這個人就是波斯王大尼烏王.
這個石碑上面.
你會發現.
有波斯太陽神的形象.
一個太陽.
以及下面有一些太陽光線射下來.
太陽也是有翼的.
下面有一對面對面的形象.
就是大尼烏王自己.
中間有一個有埃及風格的.
一個長橢圓形的皇命框.
叫做Kathulj.
上面當然是大尼烏王的皇銜了.
當然這條運河挖通之後.
在Wadi Thumala這裡.
當然一定會繼續要去維修.
因為河流下來的泥土.
不斷會有淤泥塞住運河.
所以不久之後也塞住了.
直到公元前三世紀.

$^{161}$就是公元前大概二百多年的時候.
當時就是另外一個王.
就是多利米王Philadelphus.
他再下令清理和挖通這個Canal.
話說這個Canal.
這個Wadi Thumala的運河.
有很多名人.
古代的貴族都曾經在這裡暢遊過.
其中一個就是埃及艷后.
或者埃及妖后Cleopatra.
第二點跟聖經考古有關.
聖經的世界有關的.
就是講到出埃及的路線.
如果我們放大去看的時候.
你會發現出埃及的路線.
裡面提到過紅海這個事跡.
如果舊的聖經地圖就會告訴你.
這些以色列人出埃及的時候.
他們從蘭室這個歌山地開始出發.
他們一直往南.
經過下面蘇伊士灣這個所謂的紅海.
我們今天中英文聖經.
很多時候現代翻譯是這樣的.
就是Red Sea Crossing.
就是過紅海.
我們現代人一聽到紅海的時候.
我們自然當然就會想起蘇伊士灣.
因為這裡真的叫Red Sea.
但其實這個傳統.
是來自公元前二三世紀.
七十士日本的希臘文的希伯來聖經翻譯出來的.
當時的七十士日本翻譯的猶太人.
他們大致上都是在埃及這個地方居住.
很可能他們翻譯的時候.
他們覺得摩西大以色列人過紅海.
那個紅海就是南面的Gulf of Suez.
這個Red Sea.
所以他們就將希伯來文的原文翻譯成Red Sea.
其實就不是真的Red Sea.
真正的原文其實這裡看到了.

$^{201}$就是Yam Suv.
希伯來文可以翻譯為草海或者是蘆葦海.
所以真正的所謂過紅海.
或者過Yam Suv.
以色列人過蘆葦海的時候.
那個水域.
那個水體應該是在這一處的.
為什麼呢?.
因為將聖經的地理資料.
將考古的資料結合的時候.
你會發現他們從蘭塞這個地方出發.
一直往西南走.
到了以上的時候.
就發生了一件事.
什麼事呢?.
就是出埃及記第十四章.
十四章第一節說什麼呢?.
去到這個地方的時候.
耶和華就曉言摩西了.
他說你要吩咐以色列人轉回.
Make a U-turn.
回頭.
所以如果你按照這樣的數據的話.
他們就會在以上這一處.
他們會回到北面.
然後在一個海或者湖.
一個水體面前.
來安營.
在海邊安營.
靠近海邊安營.
那個海邊.
在《上文夏利》我們知道.
只是提到Yam Suv.
當時新王國出埃及記的時候.
也有提過Yam Suv這個水體.
這個水體就在蘭塞的東面.
在這裡.
當你今天去到這個地方的時候.
你會發現只有一些很淺的小湖.
或者甚至是乾地.

$^{241}$為什麼呢?.
這就是蘇伊士運河所產生的問題.
就是十九世紀中的時候.
刮通這條運河的時候.
就將這條運河附近一帶的.
本身有的古代的一些湖.
一些水域.
全部沽乾了.
所以你今天去到那裡.
只不過是一個乾地.
以及一條運河.
但是結合了《聖經考古》.
最新的一些研究.
我們就發現原來有很多這些古代的湖.
就是在運河北至南部署在這裡的.
所以摩西大以色列人出埃及的時候.
過紅海Yam Suv.
草海路尾海的時候.
其實大概應該是這個地方.
所以他過了這個湖或者海的時候.
他們就到了西乃的曠野.
一直向南面走.
就到了南面的西乃山.
所以在這個所謂出埃及的故事的事跡裡面.
是一種 itinerary(故事).
一種神帶領著以色列人出埃及的.
一種路線的旅程的事跡故事.
裡面有很多更深的屬靈的信息.
就好像十四章第一節這裡所說的.
耶和華曉玄摩西.
他說吩咐以色列人.
這個時候突然間要轉會.
去到一個地方是前無去路.
後有追兵.
那麼你看到了.
這個帶領的故事.
旅程的故事.
就講到上帝帶領神的子民的道路.
往往不是我們人心想一定是這麼順利.
一定是這麼方便.

$^{281}$或者是最快.
神帶領人走的道路.
往往最重要的目的.
是彰顯神自己的奇事.
神自己的榮耀.
講到這裡.
下次繼續跟大家探討聖經考古.
\newpage



\section{}
\label{sec:FvI2UPC3nfc}
\textbf{Kingdom Values XLVII: Sermon on the Mount}
\newline
\newline
連結: \href{https://youtube.com/watch?v=FvI2UPC3nfc}{\texttt{https://youtube.com/watch?v=FvI2UPC3nfc}} ~~~~ 語音日期: 2021-04-17
\newline
\newline
\hyperref[sec:_gtM1xmdWw0]{\small{< < < PREV SERMON < < <}}
~
\hyperref[sec:index_chronic]{\small{[返順時目]}}
~
\hyperref[sec:index_scriptual]{\small{[返順卷目]}}
~
\hyperref[sec:U5GUh9ATRrs]{\small{> > > NEXT SERMON > > >}}
\newline
\newline
$^{1}$Welcome family, it's great to be with you again..
Today is part two of this message that we're looking at asking and receiving in Matthew.
chapter 7..
We're going to be in verses 7 to 11 again..
So let's dive into this and look at what God has to say to us..
Starting in verse 7 it says this, "Ask and it will be given to you..
Seek and you will find..
Knock and it will be opened to you..
For everyone who asks receives, and the one who seeks finds..
And to the one who knocks it will be opened..
For which one of you, if his son asks him for bread, will give him a stone?.
Or if he asks for a fish, will give him a serpent?.
If you then who are evil know how to give good gifts to your children, how much more.
will your Father who is in heaven give good things to those who ask him?".
You know, last week as we began to look at this passage and discuss it, we talked about.
kind of the obvious and direct reality of this passage, that Jesus is just challenging.
us to ask God, and that we don't do that enough..
And a lot of times it's because we don't have faith..
We're not actually confident and trusting God that he actually will be present in our.
life, meet our needs, bring the transformation that we need or the provision that we need,.
whatever it is that we might be asking for..
A lot of times we don't trust that God can do that..
And so we talked about some of those obstacles to faith that we need to deal with..
Maybe we don't believe that God is powerful enough..
Maybe we don't believe that God is good enough..
Or maybe we don't really believe that we're weak and we think that we can still find the.
solutions and the answers and fix our lives on our own without God..
Whatever those things are, we need to learn to put them out of our life so that we can.
begin to truly trust and have faith in God and begin to ask him for the things that we.
need in our life, the big and the small things..
And today I want us to consider, what are we supposed to ask for?.
Because if we're learning the lesson that we talked about last week and we're saying,.
"Okay, I want to learn to have more faith..
I want to learn to have more confidence and trust in God..
I want to learn to ask God..
I want to learn to rely on him for everything.".
Then we get to the question of, "But what do I ask for?.
What's important for me to ask for?".
And that's a really valid question for us to ask..
James chapter 4, in the first eight verses, it says this, "What causes quarrels and what.

$^{41}$causes fights among you?.
Is it not this?.
Your passions are at war within you..
You desire and do not have, so you murder..
You covet and cannot obtain, so you fight and quarrel..
You do not have because you do not ask..
You ask and do not receive because you ask wrongly to spend it on your passions..
You adulterous people, do you not know that friendship with the world is enmity with God?.
Therefore whoever wishes to be a friend of the world makes himself an enemy of God..
Or do you suppose it is to no purpose that the scripture says he yearns jealousy over.
the spirit that he has made to dwell in us?.
But he gives more grace, therefore it says God opposes the proud and gives grace to the.
humble..
Submit yourself therefore to God..
Resist the devil and he will flee from you..
Draw near to God and he will draw near to you..
Cleanse your hands, you sinners, and purify your hearts, you double-minded.".
This passage is giving us a picture of God's love for us and how God really does want us.
to be in relationship with Him and not distracted by everything else in our world..
How God even has a spirit of jealousy and says, "Hey, I want you for myself..
I love you and I want you to be in relationship with me.".
But in the beginning it echoes some of the things that we've been talking about..
In verse 2 it tells us that we don't have because we don't ask..
It's saying, "Listen, there's a lot of things in your life that you don't have just because.
you're not asking..
You're not even asking God for them.".
But then in verse 3 it says that there are things that we ask for and don't get because.
we're asking for the wrong things..
We are asking for things that we just want to spend on ourselves..
We're asking for things to gratify our own flesh and our own lusts and our own passions.
and our own pleasures and our own comfort..
And that's not necessarily what God wants to give to us..
And so there are things that we ask for and don't get because we're asking selfishly..
And that's what I want us to think about today as we learn to ask, "What are we going to.
ask for?".
And I think Jesus gives us some clues in this passage and also just in the context of the.
entire Sermon on the Mount..
We see that Jesus is pointing us towards what we should ask for..
So when we look at verse 9 of Matthew chapter 7, Jesus gives us this illustration of someone.
asking for a gift..

$^{81}$And Jesus says, "Which one of you, if his son asks him for bread, will give him a stone?.
Or if he asks for a fish, will give him a serpent?.
If you then who are evil know how to give good gifts to your children, how much more.
will your Father who is in heaven give good things to those who ask him?".
Now part of this illustration, Jesus is just pointing us to the fact that God is good and.
he wants our good in our lives..
And that God is going to act in that way that we can trust the goodness of God, just like.
we talked about last week..
And that's a powerful truth here..
But I think there's also a powerful truth here for us to consider about what is good.
in our lives and what is best..
God knows us in a way that we don't even know ourselves..
And God is going to work in our lives in a way that is truly for our best, not just for.
what feels good to us in the moment..
And that brings up some different questions here..
What does it mean to ask God for truly good things?.
What does it mean to ask God for truly what's best for us in our lives?.
I think all of us can think of plenty of things that we would ask God for..
And so if we just say, "Okay, God wants me to ask him for things and that's the end of.
it," we would probably all say, "Okay, I know what I want to ask for..
I want to ask for comfort..
I want to ask for stuff..
I want to ask for success..
God give me a lot of money..
Give me a big house..
Give me a few cars to drive..
Give me good health..
Give me time to sit on the beach with my family..
God make things easy..
Give me relationships where no one's unkind to me..
And make everything work out at my job..
Or better yet, make me not even have to go to work anymore..
Make me no longer have to do school.".
If we were just going to ask for whatever we wanted, it would probably be a pretty long.
list of things that were focused on our own selfishness and on our own comfort, right?.
That's just the honest reality of most of us as humans..
And that's borne out in the way we live our lives..
The things that we pursue and the things that we try to fill our lives with are often based.
and focused on our own selfishness or comfort..
They're based on trying to gain things and possessions and prestige and stuff, comfort,.

$^{121}$pleasure for ourselves..
And they're based on trying to insulate us from things that would hurt us or make us.
uncomfortable or make life difficult or make anything challenging in our life..
That's how most of us operate..
That's how most of us live our lives..
But is that really what's best for us?.
Well there's no question that that's something that feels good in the moment, right?.
It feels good to have an easy life..
It feels good to not deal with challenges and adversity..
It feels good to not have anyone in your life be sick..
It feels good to not have anyone be unkind to you..
There's no question that that's what feels good in the moment is pleasure and comfort.
and the pursuit of those types of things..
But is that what's truly best for us?.
And so when Jesus is pointing us to the fact that God is good and he's going to give us.
good things, he's going to give us bread, he's going to give us a fish, he's going to.
give us what we truly need in our lives to sustain us spiritually, that Jesus is actually.
teaching us to pursue holiness..
Jesus is teaching us to pursue what's truly important in our lives..
And I think there's a really powerful truth here as well..
That God is ultimately waiting for our response..
I've talked about this a lot, but God is not coercive..
He is not..
He's a God of love and he's not a God who forces us to do things..
God has all the power in the universe, but he doesn't use that power to force us to do.
things..
He can control things and he turns the course of history wherever he wants it, but he doesn't.
force us to love him..
He doesn't make us enter into relationship with him, even though he could, because God.
is a God of love..
And that's true in our lives as followers of Jesus as well..
We see it true in the world at large that there are many people who do not love God.
and who are not in relationship with him, who are not a part of the kingdom of God..
God is not going to force them to be in relationship with him..
But even those of us who have begun relationship with God and come into the kingdom and begin.
to live our lives in obedience to him, he's still not going to force us to change..
He's going to invite us to change..
He's going to invite us into transformation, but he's not going to force us to make it.
happen..
God is waiting for us to respond..

$^{161}$Just like when you send an invitation to someone and you say, "Hey, I want you to come to my.
birthday party," or "I want you to come to my graduation party," or "I want you to come.
to my wedding and my retirement party," whatever it is, you send the invitation, you say, "Please,.
RSVP," and then you wait for their response..
That is what God is doing to us..
God is inviting us..
He's saying, "Hey, I want you to come into relationship with me..
I want it to be deep and fulfilling and satisfying, and I want it to be transformative in your.
life..
And this is what I'm inviting you to..
I'm inviting you to live out the values of the kingdom and to experience all the goodness.
of what it means to be a child of God..
But I'm not going to force you to do it, and I'm going to wait on you to respond.".
And that's what God is doing..
Because God is waiting for us to respond, we have to begin to pursue holiness..
We have to begin to choose to draw near to God, like the book of James talks about..
If we will draw near to Him, then He will draw near to us..
But He's waiting for us..
I think that's really powerful..
A lot of times when we look at our life and we say, "Well, I've asked for something that's.
not coming to pass," or "I want to grow in my faith, and why doesn't God just make that.
happen in my life?".
Sometimes we assume that it's God's fault, that like, "Well, this hasn't happened yet.
because God just doesn't really love me or He doesn't really care about me, or He doesn't.
want to meet this need, or He doesn't want to help me grow in my faith, or He doesn't.
want to change this thing in my life.".
And actually the opposite is true..
God is right there the whole time saying, "I love you..
I want to go deeper in fellowship with you.".
We're the one who is not responding to God..
We're the one who's holding God at arm's length, and we're the one who's saying, "God, I don't.
want to go any further in my relationship with you..
I don't want to give you any more control over my life..
I don't want to give you any more access into my heart.".
God is waiting for our response, and Jesus is teaching us that we need to pursue holiness.
and draw near to God..
God wants to give us those things that are spiritually sustaining to us, the bread and.
the fish, but we don't have them because we're not responding..
We're not asking God for those things..
We're not saying, "God, instead of just making me comfortable and giving me everything I.

$^{201}$want in my life, can you give me the things I truly need?.
Can you give me holiness?.
Can you give me righteousness?.
Can you help me to live rightly before you?.
Can you help me to live out the values of the kingdom, even if it means walking through.
adversity?".
But we don't ask for that, and God is waiting for us to ask..
God is waiting for our response..
There was this pastor of a previous generation that spent a lot of time preaching and ministering.
right here in Chicago..
His name was A.W..
Tozier..
He wrote a lot of excellent books, so if you haven't ever read any of them, I encourage.
you to check some of them out..
I remember reading one of his books and coming across this interesting statement where he.
said, "Every man is as holy as he wants to be.".
On the face of it, that statement sounds kind of crazy and almost heretical, right?.
It almost sounds like, "What is this guy saying?.
Is he saying that we manufacture holiness, that we make holiness happen?".
We know that is not true in any way..
That it's only through Christ and what he's done on the cross that we could ever draw.
near to God..
It's only through Jesus that we could ever have relationship with God..
It's only through Jesus that we could ever experience holiness..
It's not that we manufacture holiness, and that's not what Tozier was saying..
What he was saying is that when there is a deficit in our holiness and the rightness.
of our relationship to God, it's because there's a deficit in our desire..
It's not because God doesn't want us to be holy..
It's not because God doesn't want to teach us and grow us..
It's not because God doesn't want to take us farther..
It's because we don't want it..
We don't want to go any further..
We don't want to grow any more..
We don't want to become more holy and any closer to God..
And that is a powerful truth..
I think Matthew 7 is really pointing us towards that here..
Jesus is saying, "Listen, I want you to learn to ask God for good things in your life, and.
you need to know what to ask for..
Ask for the good stuff..
Ask for holiness..

$^{241}$Ask for the spiritually transformative things in your life, because that will cause you.
to grow into who God wants you to be..
And if you're not growing into who God wants you to be, it's not because God doesn't want.
it for you..
It's because you don't want it, and you're not moving towards it..
You're afraid of it, or you're holding on to your sin, or you're letting something else.
hold you back..
And I think when we look at these verses in Matthew 7 in the context of the Sermon on.
the Mount as a whole, we see that Jesus is teaching us that God desires us to possess.
and live out Kingdom values..
If we say, "Okay, Jesus is telling me to ask," and if we just pulled this passage out, Matthew.
7, 7-11, and looked at it out of context, we could say, "Well, what is Jesus telling.
me to ask for?".
Anything, right?.
Like I said earlier, maybe it's just stuff and comfort and pleasure..
But when we look at it in the context of the entire Sermon on the Mount, it begins to open.
up what we're supposed to ask for..
Jesus tells us to ask, "What has He been talking about this whole time?".
He's been talking about the values of the Kingdom..
He's been talking about poverty of spirit..
He's been talking about meekness..
He's been talking about being a peacemaker..
He's been talking about storing up treasures in Heaven, seeking first the Kingdom of God..
He's been talking about being pure of heart..
He's been talking about having a clear spiritual vision..
All these things that we've discussed over the last year that Jesus has been talking.
about, that's the context of this statement where Jesus is asking, "It will be given you.".
He's saying, "Hey, here's the values of the Kingdom..
I'm laying them out in front of you.".
And as you listen to them, yes, they are things that you do not currently possess in your.
life..
So ask, and it will be given to you..
Ask God for them..
Seek them, and you will find them..
Look, and God will open that door and begin to give you those values of the Kingdom in.
your heart and in your life and begin to allow you to live them out..
And that's one of the things that Jesus is teaching us here..
Jesus is teaching us what to ask for..
We should ask for the values of the Kingdom..
We should seek the values of the Kingdom, and we should ask God to begin to give them.

$^{281}$to us so that we possess them and we live them out in our lives..
Another thing that Jesus is teaching us here is that God is more interested in our spiritual.
well-being than our material well-being..
That's really important for us to understand..
When we think about prayer and we think about asking God, I think all of us have had experiences..
If you've been a follower of Jesus for very long, I'm sure you've had an experience where.
you've asked God for something, and you didn't get what you wanted, or you didn't get it.
in the way that you wanted..
And that's challenging for us..
That challenges our faith..
That causes us to say, "Okay, God, can you answer prayer?.
Do you want to answer prayer?.
Why would you not answer this prayer?".
It gives us a lot of questions, right?.
And I think we need to understand that God is more interested in our spiritual well-being.
than our material well-being..
Sometimes we ask for something that may be what's best for us materially and physically,.
and in the moment we may ask God for some stuff or for some pleasure or for some comfort.
or to work something out in a way that's great for us..
And maybe that's not what God wants to do in our life at that moment..
Because maybe God sees that giving us that thing or that position or that situation or.
that job or that educational achievement or whatever it is in that moment would actually.
not be what's best for us spiritually..
Maybe that would cause us to be proud or cause us to look down on other people..
Or maybe it would just cause us to miss out on a lesson that we desperately need to learn,.
that we've got to go through some adversity to learn..
And so God in his love for us is more concerned with our spiritual well-being than our material.
well-being..
And sometimes God says no or sometimes God says wait when we ask because he has a lesson.
that's so valuable for us to learn spiritually that we would not learn if he just said yes.
and gave us the thing that we want..
And we have to remember that when we're asking God..
Jesus is pointing us to that in the Sermon on the Mount..
He's saying yeah, there's a lot of things that you could ask for, but here's what you.
really truly need, what's really truly going to be good for you and what's really truly.
going to bring peace and fulfillment and satisfaction and joy into your life is to possess and live.
out the values of the kingdom..
You need to ask God for that..
And God's more focused on you being healthy spiritually than he is on you just getting.
everything you want out of life..

$^{321}$And that's really valuable and powerful for us to know..
There are times when we just want things to be easy and simple and quick and exactly how.
we want them to be..
And we ask God for that..
And that's okay..
It's okay to ask God..
But God doesn't necessarily say yes to that because God wants what's best for us in a.
deeper way..
And that's what God is going to move towards..
You know, God is more interested in his kingdom than ours..
Sometimes we ask God things that really just have to do with our own kingdom and our own.
pursuit and what's going to benefit us..
And God loves us and he wants good for us, but he's more interested in the goals of his.
kingdom than he is in ours..
And that may sound selfish..
We might say, well, how could God be so selfish and not want what's best for me?.
But if we're honest with scripture and who God says he is, if God really is all powerful.
and all knowing, then how foolish would it be for him to allow our goals and priorities.
to outsize his goals and priorities..
His kingdom should come first..
And so it's only right for God to be more interested in the goals of his kingdom than.
in the goals of our kingdom..
And the reality is that God is more interested in his glory than ours..
And again, if we're honest with scripture, that's only right..
That's exactly how it should be..
God should be more interested in glorifying himself than in glorifying us..
And so as we ask for things in our life, we need to understand that God is going to respond.
and God is going to teach us and to grow us and to love us..
But it may not always move in the direction of our glory or the goals of our kingdom or.
even our comfort and pleasure..
Sometimes it will, but there will be many times when it does not because God is more.
interested in our spiritual well-being..
He's more interested in glorifying himself and in advancing his kingdom than in just.
giving us what we want..
So I want to give us a few things to think about as we close out here..
As we begin to think about asking God and what we should ask for, I would say, don't.
be afraid to ask God for anything big or small..
God cares about your life and he wants to hear the desires and needs of your heart..
So pour those out to him..
Even if those are small things, pour them out to God..

$^{361}$If they're emotional needs, pour that out to God..
If they're physical realities, pour it out to God..
Even if you're not sure and you say, "Well, I don't know if I'm just asking for my own.
comfort or pleasure.".
There's nothing wrong with asking..
If that's not what God has for you, then God will reveal that to you..
But it's just like little children are not afraid to ask a loving parent for a gift,.
even if they know the answer might be no..
And we are God's children..
He loves us..
So don't be afraid to ask for gifts, even if it's something that may not be right for.
you at the time..
God can tell you no..
Don't be afraid to ask for anything big or small..
Pour your requests out to God..
Consider that privilege that you have of going before the Almighty God of the universe to.
make your requests known..
And then I would encourage you to, if you're not sure whether or not you should ask for.
things, ask for things that you know are in alignment with kingdom values..
Ask for things that you know are in alignment with God's will..
You know, we can say, "Well, I don't know if I should ask for this in my work or for.
that in my work.".
Well, if one situation in your job is going to give you an opportunity to display poverty.
of spirit or to be a peacemaker, then you should probably ask for that because it's.
in alignment with the values of the kingdom..
There are certain things that we don't have to wonder about whether or not we should ask.
for or God wants us to have them if they're already in alignment with kingdom values..
We don't have to ask God and say, "God, do you really want me to help the poor and the.
needy?.
Is that something that you want me to do?".
The answer is already yes..
God's made it clear that that's a kingdom value that we should be helping those who.
are in need..
So you might need to ask some specific questions about how you do that and who specifically.
you help..
But you don't have to ask whether or not you're supposed to help because you already know.
the kingdom value there..
And so think about kingdom values in your life and some areas, some of those situations.
that you may be confused about and you may be asking God's questions..
Maybe those questions are already answered just by looking at and aligning yourself with.

$^{401}$the values of the kingdom in those areas..
And ask for things that match God's priorities, that align with God's priorities..
What's important to God?.
God wants the people around you to feel his love, to feel his forgiveness, to feel his.
grace, to feel his goodness..
God wants the people around you to come into relationship with him..
And so if you say, "Well, I don't really know what to pray for my family," or, "I don't.
really know what to pray for my friends or my co-workers or my classmates," that's a.
great place to start..
Start asking God for those things..
You know what God's priorities are for those people in your life..
God wants them to experience his love and his forgiveness..
God wants them to come into relationship with him..
God wants them to be served by you, by the members of the family of God..
Start asking for that stuff..
You don't have to wonder if God's interested in those things..
We already know that God is..
Those are priorities in God's heart and in God's kingdom..
And ask for God's glory, things that are in alignment with God's glory..
Again, we don't have to wonder if God is going to be behind things or in things that are.
glorifying him..
We know that God is going to be in things that are glorifying him..
And so when you face a situation in your life and you're like, "I don't know what to pray.
here..
I don't know which outcome to pray for," that's a great place to start..
Say, "You know what?.
There's this situation before me and this could happen or that could happen, and I'm.
not sure what God wants in my life.".
Begin to pray, "God, would you be glorified?.
I don't know what I'm supposed to choose..
I don't know what's going to happen..
I don't know the future..
I don't have control over it..
But would you be glorified whatever happens?".
Begin to ask for things that are in alignment with God's glory..
The reality is none of us have as much control over our lives as we think we do..
We are not able to control the situations that are presented to us in life, but we can.
control the way that we respond to them..
And God wants us to ask for his help in every situation that we face..
God wants us to ask for his help when we have a need..

$^{441}$God wants us to ask for his help when we are afraid, when we are confused, when we don't.
know where to go, when we need guidance, when we experience some type of deficit in our.
ability and we say, "God, I want to love this person, but I don't know how.".
Or "God, I want to live out the kingdom values in this situation, but I don't have the capacity.
to do that.".
God wants to help us in those situations..
God wants us to ask for his help in every situation of our life..
And we need to learn what to ask for..
The things that align with the values of the kingdom, the things that align with God's.
priorities, the things that align with God's glory..
We need to learn to seek holiness, to seek to draw near to God..
We need to pursue that because God is waiting for us to respond to him..
If there's not holiness and righteousness in our life, it's not God's fault..
It's on us because we're not moving towards him..
He's always moving towards us..
And we need to really seek and ask God to allow us to possess and live out these values.
of the kingdom..
Let's pray..
God, I thank you for the privilege that you've given us to ask you..
And I pray that we would learn to ask you for the right things..
I pray that we would learn to ask you for the things that are truly best for us, for.
the things that are spiritually valuable, not just for our material and physical well-being,.
but that we would ask for the things that really are going to be best for us..
I pray that we would learn to ask for things that align with the values of the kingdom,.
that align with your will and your priorities, your goals, and your glory in our life, God..
And God, I pray that you would help us to begin to ask and seek to pursue holiness in.
our life, to pursue deeper relationship with you..
And I pray that you would help us to ask that we would be able to possess the values of.
the kingdom, that you would give us these values of the kingdom that we've talked about.
in our hearts in a real way so that we can live them out and that we can experience you.
using us to bring healing and love and encouragement and salvation to the world around us, to the.
people around us, God..
In Jesus' name, amen..
[Music].
[MUSIC].
\newpage



\section{}
\label{sec:U5GUh9ATRrs}
\textbf{考古探秘 (法老黃金遊行) - 陳崇基牧師 Mark Sir (2021.04.18)}
\newline
\newline
連結: \href{https://youtube.com/watch?v=U5GUh9ATRrs}{\texttt{https://youtube.com/watch?v=U5GUh9ATRrs}} ~~~~ 語音日期: 2021-04-18
\newline
\newline
\hyperref[sec:FvI2UPC3nfc]{\small{< < < PREV SERMON < < <}}
~
\hyperref[sec:index_chronic]{\small{[返順時目]}}
~
\hyperref[sec:index_scriptual]{\small{[返順卷目]}}
~
\hyperref[sec:Ks1q0nCc4X4]{\small{> > > NEXT SERMON > > >}}
\newline
\newline
$^{1}$你好,又到了《聖經考古探秘》的時間.
今次繼續跟進一些埃及發生的新聞.
這些新聞與《聖經考古》,《聖經》的歷史背景有什麼關係.
4月3日的時候,大概兩個星期前.
埃及上下一個大的國家新聞.
就是埃及的開羅舉行了法老黃金遊行.
就是Golden Parade of Pharaohs.
這個遊行組織已經很久了.
就是為了在埃及的博物館裡面.
22具木乃伊.
將它搬到開羅市南面.
另外一個新建的博物館裡面.
為了這22具木乃伊.
來到搬家的時候.
埃及上下整個政府就做了一場很大的show.
每一具木乃伊就放在一個專車裡面.
車程大概有四十多分鐘.
因為怕這些木乃伊裡面的屍骨會腐化.
所以當局就有專人來護送.
並且那個棺木就放在氮氣的特製箱裡面.
以免氧化等等有損壞.
每一輛車上面就有特製的避震.
這些木乃伊運送的時候.
你可以看到有射燈,有表演者,有音響,交響樂.
盡顯生前古埃及帝國的威榮.
這22具木乃伊.
18個是法老王.
有4個是皇后.
所以總計有22個.
全部所代表的歷史.
就是古埃及歷史裡面的17,18,19,20王朝.
大概的時間就是公元前1600至公元前1000年的時間.
所以所代表的歷史有五,六百年之多.
這22個法老巡行的時候.
就按照他們當時統治的時代.
所以我就在這裡預備了一個埃及的列王表.
埃及歷史的列王法老列了出來.
我先處理18個法老.
先不要管皇后.
第一張是埃及的第二中介期.

$^{41}$Second Intermediate Period.
這個時間大概是公元前1650至1550年.
這一百年歷史學家叫做Second Intermediate Period.
這個就是中皇國Middle Kingdom之後的.
一百年的時間.
特別是埃及是不統一和四分五裂的.
所以你看到這裡有一些問號.
或者有一些並列的混亂.
左邊你看到就是15王朝的時候.
15th Dynasty.
這一百年就是在下埃及.
就是尼羅河三角洲的地方.
進入了Hysso Period.
希克索時期.
Hysso的意思就是Foreign King.
就是外族的君王.
外族的統治者.
在這一段時間前前後後.
有很多亞洲裔的閃族的人.
遷徙到埃及的三角洲.
他們有很多定居在一個地方叫做Afaris.
現在叫做Tel Adaba這個地方.
如果你看Google Map就知道.
在尼羅河三角洲的.
聖經叫做哥山地這個地方.
接著Premse這個附近.
旁邊就是Afaris.
就是Tel Adaba這個地方.
在這個時間.
大概公元前1500的時候.
有很多亞洲人.
Asiatics.
搬了到三角洲.
所以他們人多的時候.
甚至是管治埃及最肥沃的土地.
然後趕走了本土的埃及人.
本土的埃及人就南下.
他們南下去到.
Thieves這個地方.
被趕到Thieves這個地方.

$^{81}$他們就管了上埃及.
埃及人自己管的王朝就是十六.
十七的王朝.
就跟希克索斯這個時期是平衡的.
這些本土的埃及人就列出了這些名字.
去到最後這個黃色的地方.
我highlight了的.
這個就是Secondary Dao.
中文可以譯作薩奇尼克道王.
這個就是十七王朝末的時候.
他是May2那個.
Second to last的埃及王.
他在Thieves這個地方.
就是底比斯的作王.
由他開始就起義.
想攻打尼羅三角洲的外族人.
他就出擊過好幾次.
他是第一個帶領埃及人去反擊的.
去到Avaris這個地方.
Tel Adaba.
這個時間大約是公元前1560年的時候.
好幾次游擊嘗試去到反擊.
都是失敗的.
這位的Secondary Dao.
是一個帶領者.
他也有稱為勇者.
The brave的意思.
他的木乃伊是非常著名的.
為什麼呢.
因為他的頭顱有好幾處斧頭造成的傷口.
曾經有些學者懷疑.
他在反擊的戰場上.
死於沙場.
所以被人打或者被人攻擊.
以致頭顱破了.
有好多的傷口.
但最近醫學上用CT scan再一次掃描他的頭骨.
發現似乎不是在打仗的時候.
留下的傷口.
不是在搏擊留下的傷口.

$^{121}$而是可能是戰爭的時候失敗.
被捉拿了.
然後處決.
這些傷口在刑場上.
被處決的時候留下的傷口.
非常殘忍.
這些斧頭造成的傷口的形狀.
也很配合Hicksall.
這些Foreign King用的斧頭的形狀.
非常有可能是他反擊的時候最後一次失敗.
然後被希克索斯人來刑場.
在刑場上監獄處決.
然後才變成了穆拉爾.
才埋葬了.
Segundavis死了之後.
他的兒子Khamos.
繼續去做埃及人的起義反攻戰.
Khamos在十七王朝.
是最後最後的一個王.
他繼續和Hicksall作對.
想反攻.
結果到了第十八王朝.
就是Arkmos.
瓦克摩西一世.
他成功擊退了管治了一百年的外族希克索斯人.
把這些亞洲人趕走了.
把外族人都趕走了.
統一了上下埃及.
之後埃及進入了一個新的時期.
叫做New Kingdom Period.
新王國的時期.
進入了十八王朝.
18th Dynasty和1920王朝.
這些highlight就是在這次黃金遊行裡.
有份的法老的穆拉爾.
有幾個很出名的.
譬如Thutmose III.
圖特摩斯三世.
圖特摩斯三世不單是一個很英明的.
很威的一個法老.

$^{161}$也在亞洲其他地方.
做了很多的征途.
另外有一個很出名的就是Ramesses II.
蘭瑟二世.
這個也叫蘭瑟大帝.
因為他管治了埃及很長時間.
你算一下有多少長.
有六十多年的時間.
這個Thutmose III或Ramesses II.
很可能就是出埃及記裡面的法老.
也更加有可能.
現在新的學者的研究.
更加有可能就是Ramesses II.
因為出埃及記的文本.
有很多是講到的用詞.
或者地理的名稱等等.
都是很配合十九王朝.
At-Night-Tin Dynasty.
無論如何.
這個跟聖經的背景更加有關係.
除了是講到出埃及的法老之外.
你記得在創世記最後幾章.
和出埃及的第一章.
那個歷史背景.
或者你看經文的感覺.
就是你覺得創世記最後那幾章.
那時候的埃及法老.
特別是對約瑟和雅各加.
他們在加拿大亞洲Asiatic.
他們到埃及避難的時候.
是特別友善的.
如果你拼在大歷史上.
你就知道了.
因為當時是Hexo Period.
Hexo Period是什麼呢?.
是外族人在管理埃及.
他們在做法老的時候.
所以都是在講一些同族.
相似的非埃及人.
你看到創世記幾章之後.

$^{201}$到出埃及的第一章.
你發現那個情況.
突然間因為八字路轉變.
出埃及的第一章第八節.
就講到那時.
有一位不認識約瑟的新王興起.
統治埃及.
這個新王興起雖然沒有講名字.
但這裡似乎不是純粹一個王出現.
有些埃及的學者.
兼聖經學者都指出.
這個不只是單單有另外一個法老出現.
而是有一個新的王朝出現.
有剛才那個歷史背景的時候.
你就發現這個很可能是十八王朝的出現.
因為到十八王朝出現的時候.
他們攻打成功.
將這些亞洲人.
希特索斯人.
外族人趕走了之後.
當時新王國New Kingdom的埃及帝國.
他們的外交政策就完全改變了.
對外族非常敵視.
甚至去奴役他們等等.
這個就很配合出埃及的第一章所講.
因為一章第八節說有新的王.
或者新的王朝.
新的Dynasty出現之後.
希伯來人就被奴役.
然後造尊等等.
這些都是新王國當時的背景.
環境是非常之吻合的.
希望這個新聞能夠讓大家想起.
原來創世記和出埃及的背後.
有一個古埃及的歷史背景.
講到這裡.
下一次繼續.
\newpage



\section{}
\label{sec:Ks1q0nCc4X4}
\textbf{舊約網上行 - 上帝的足跡:EP8 亞伯拉罕與羅得}
\newline
\newline
連結: \href{https://youtube.com/watch?v=Ks1q0nCc4X4}{\texttt{https://youtube.com/watch?v=Ks1q0nCc4X4}} ~~~~ 語音日期: 2021-04-20
\newline
\newline
\hyperref[sec:U5GUh9ATRrs]{\small{< < < PREV SERMON < < <}}
~
\hyperref[sec:index_chronic]{\small{[返順時目]}}
~
\hyperref[sec:index_scriptual]{\small{[返順卷目]}}
~
\hyperref[sec:QTJ66JH0gPw]{\small{> > > NEXT SERMON > > >}}
\newline
\newline
$^{1}$我們這次繼續追溯阿巴拉罕的故事.
上兩次我們提到阿巴拉罕離開吾爾哈蘭米索不達米亞來到迦南地.
然後這裡有饑荒.
他們暫時到了埃及避難.
在那裡發生了第一次太太當是妹妹的失敗故事.
在危機裡面上帝用恩典來保守.
然後他們一家又回到迦南地.
阿巴拉罕的故事.
可能我們現代人看下去的時候.
我們會覺得很鬆散.
有些事情說了.
然後又暫時不說.
然後之後再看.
看了這麼多集你應該明白了.
聖經的編排.
很多時候不一定按照我們現代人.
直線.
按時序硬性編排.
有很多時候整個故事的情節會不斷向前.
但是故事的中心.
很多時候不一定在尾段.
而是在中間.
特別在創世記.
特別很多時候.
作者會取材一些故事.
是一對一對記載的.
構成一種三文字或者交叉.
Chaostic的結構.
阿巴拉罕的故事主要的內容.
可以在創世記十二至二十二章找到.
如果這樣列出來的話.
大家可能會清楚很多.
一直按著ABCDE.
來進發.
然後在下面的一半.
又是影照上面ABCDE.
不過是調轉次序.
譬如你看A.
就是哈塔拉的後代.
有一個是加普.

$^{41}$然後最後就說到阿巴拉罕的兄弟.
拿學的後代.
然後第二段.
就是說到阿巴蘭信心的旅程.
他的啟示.
十二章開始.
這裡就影照.
就是他最後一次信心的旅程.
就是憲爾薩的旅程.
就是他信心的高峰.
二十二章一到十九節.
第三段.
就是說到撒萊.
或者撒拉.
阿巴拉罕的太太.
在外邦的宮廷.
然後危機得到化解.
然後這個相對於下面的.
另外一個也是撒拉的故事.
在外邦的宮廷裡面.
這次在基拉爾.
危機也得到化解.
所以你這樣一段一段看的時候.
你就會發現.
原來一直邁向一個.
最重要的中間位置.
中心.
就是神.
和阿巴拉罕來到立約.
講完十二到二十二章了.
主要的內容.
交代了.
然後到二十三章.
就講到阿巴拉罕買了一個墳地.
去撞他太太.
撒拉.
然後到二十四章.
就是為以撒來娶妻.
因為要找一個繼承人.
然後第三段.

$^{81}$最後一段.
就交代阿巴拉罕的逝世.
你可以看到阿巴拉罕的事蹟裡面.
一段一段.
一對一對的記載.
所以今天我們看到.
阿巴拉罕和羅德的故事.
都是一樣的.
我們看十三章開始.
然後看到十四章.
影照到下面.
就是十八十九章.
所以歡迎大家來到.
舊約網上行.
阿巴拉罕和羅德的故事.
我們先看創世記第十三章.
十三章第一節.
阿巴蘭帶著妻子和所有的一切.
以及羅德.
一同從埃及上到南地.
阿巴蘭有許多生畜和金銀.
他從南地一站一站的往前走到巴特里去.
去到巴特里和艾城中間.
就是他先前支搭帳篷的地方.
也是他竹頭一座祭壇的地方.
阿巴蘭就在那裡叩告耶和華的名.
和阿巴蘭同行的是羅德.
也有羊群,牛群和帳篷.
因為上帝的恩典.
他們一家就有許多的財物.
很多的牛群,羊群,兩叔侄.
有很大的上帝恩典的祝福.
來到巴特里這個地方.
就出現了另一個危機.
危機是什麼呢?.
十三章第六節.
他們說:我們容不下他們.
因為他們的財物很多.
所以他們再也不能住在一起了.
牧羊阿巴蘭和牧羊羅德的兩群牧羊人就起了紛爭.

$^{121}$當時迦南人和比利時人都在那裡.
這裡說到兩方面的牧人有爭執.
因為資源可能不夠.
我記得有一次.
我去到以色列的時候.
我就一個人跟一個阿拉伯人遊了一天.
他帶我去到巴特里南邊.
在希伯倫附近的山地.
來探訪他的牧羊場.
到了之後他就在牧羊場跟我說.
就說近日他的兄弟.
到了放牧的時候.
就跟另外一個家族的牧羊人有爭執.
為什麼呢?.
因為不夠草食.
就是大家去同一個地方.
所以這種情況.
其實在以色列巴斯坦地是非常寫實的.
這裡說到阿巴蘭和摩爾羅德之間.
他們的牧羊人又起了紛爭.
那麼阿巴蘭就提出了一個解決方案.
十三章第八節.
阿巴蘭對羅德說.
你我之間不可以有紛爭.
你的牧人和我的牧人也不可以有紛爭.
因為我們是知七.
全地不都是在你面前嗎?.
請你離開我吧.
你若向左我就向右.
你若向右我就向左.
這個解決的方案.
我們一定要在聖經的地理來了解.
現代人我們看地圖.
以北面為我們定方向的方位.
但是在近東在聖經的地理裡面.
定方向不是北面.
原來是在東面.
就是日出的地方.
在他們的前面.
所以將地圖扭過來.

$^{161}$我們面向前面在東面的時候.
你就明白阿巴蘭提出的方案有什麼意思.
他說前面是東面.
他說如果你向左的話我就向右.
你向右我就向左.
他的左邊自然是我們今天想的北面.
右邊自然是南面.
所以大家面向前面面向東的時候.
阿巴蘭的提議.
就是想告訴羅德.
我們仍然留在神的應許之地迦南地裡面.
我們一起分享.
你往北的話我就向南.
你向南我就向北.
所以阿巴蘭是一種慷慨地去解決這件事情.
任由他的質疑比他小.
比他地位低一點.
先去選擇他認為北面是最好的.
他就去那邊.
然後阿巴蘭就會去另外一邊.
但是羅德看到的是什麼呢?.
創世記第十三章第十節.
羅德舉目.
看見約旦河整個平原直到所以有水灌溉.
在耶和華毀滅索多瑪和俄摩拉之前.
這地好像耶和華的原子也像埃及地一樣.
他看到的是什麼呢?.
他看到前面.
他不是看左也不是看右.
不是向北也不是向南.
他向東面繼續看.
羅德看到的是約旦河整個平原.
好像埃及一樣.
剛剛他們從埃及上來.
埃及有什麼特別呢?.
就是埃及的農田收成是靠利落河.
利落河每年安時安候就會氾濫.
於是土地就會有水灌溉.
收入是可以預測的.
所以相比之下在迦南地就不同了.

$^{201}$在迦南地就要靠天雨.
是很難預測的.
所以一定要靠雨水的神.
之前伯拉罕和羅德來到迦南地的時候.
發生什麼事你記得嗎?.
就是旱災有饑荒.
不夠雨水.
如果你看地圖的時候.
從伯特利這個方位一直向東.
約旦河的平原死海附近一帶.
一直看過去就是羅德所見到的.
雖然我們不是很能夠確定索多瑪和俄摩拉.
索爾真實的地點在哪裡.
但其實地理的資料已經足夠了.
索多瑪和俄摩拉兩個的雙城.
放在死海的東邊一帶.
現在是一片礦業.
大概在阿布蘭時代.
公元前大概二千年的時候.
水源是充足的.
並且這個地方是富庶的.
羅德見到的.
想要的是什麼?.
就是剛剛從埃及回來的那種安定.
那種能夠預測的經濟.
約旦河整個平原也像埃及一樣.
聖經這樣說.
十三章第十一節.
羅德選擇了約旦河整個平原.
他就向東移動.
他們彼此分開.
再看回這個地圖.
在伯特利.
你見到這個平原直到遠遠的索爾.
羅德說他要往東走.
所以他就一直靠近對面的索多瑪.
經文絕對表示.
羅德的選擇打了一個很大的問號.
十三章十二節.
阿布蘭住在迦林地.

$^{241}$羅德就在平原的城市中間.
漸漸的羅爾象棚直到索多瑪.
第十三節特別提到.
在耶和華看來.
索多瑪的人是最大惡極.
參考一下這句話.
後面出現過.
之後第十五章.
神對阿布蘭說.
何時他的後裔可以得到迦林地.
入主迦林地.
十五章十六節.
神對阿布蘭給予一個應許.
之後到第四代的時候.
他說你的後裔必定會回到迦林地.
因為阿摩利人.
迦南人.
他的罪孽還未滿盈.
所以說到羅德的選擇的時候.
聖經在十三章十三節.
提到在耶和華的眼中看來.
索多瑪的人是最大惡的.
就是說.
索多瑪的人.
這個城市已經罪惡滿盈.
上帝的審判很快就會到.
快過審判迦林地.
索多瑪平原的城市.
正在步向滅亡.
羅德有什麼心態呢?.
他就說.
不想錯過這個良機.
這是一個發達致富賺快錢的.
一個窗戶.
上帝的應許暫時放下.
是一個機會主義者的抉擇.
最大的問題.
在羅德的眼中.
就是上帝的應許.
屬靈的事物是次等的.

$^{281}$相對於眼前的富貴.
財富.
能夠發達的.
這個才是最重要的.
這就是羅德的心境.
於是兩叔侄就分家了.
但是亞伯蘭又怎樣呢?.
亞伯蘭就留在迦林地.
這是信心的舉動.
十三章十四節.
羅德離開亞伯蘭之後.
耶和華對亞伯蘭說.
你要舉目.
又是同一個動詞.
之前說到羅德舉目.
他看到的是約旦河平原.
瑣多嗎哪種豐盛和致富.
但是亞伯蘭舉目做什麼呢?.
神說.
你所在的地方往東南西北貫寒.
你所看見的.
我都必定賜給你和你的後裔.
直到永遠.
我要使你的後裔.
多於地上的塵土.
地上的塵土能夠數算的話.
才能夠數算你的後裔.
所以你起來.
中環走遍這地就是迦林地.
因為我必把這地賜給你.
於是亞伯蘭就漸漸的樂而張鵬.
他由伯特利的地方搬到南邊一些.
去到希伯倫的萬里象樹附近.
在那裡居住.
又在那裡建一座祭壇.
創世記很多時候用兩個角色去對比.
互相比較.
這裡我們看到亞伯蘭.
或者亞伯拉罕和羅德的比較.
我們再問多一條問題.

$^{321}$為什麼要提到羅德的事跡.
我們看亞伯拉罕的故事的時候.
其實一開始羅德已經出現.
十一章最後的時候.
說到他拉的後代.
我們記不記得曾經提過羅德.
他拉生亞伯蘭,那賀和哈蘭.
哈蘭生羅德.
所以亞伯蘭和羅德是兩叔侄.
當然羅德的父親哈蘭早已死了.
然後十一章三十一節.
說到他們起行離開梧椅.
然後帶著的就是哈蘭的兒子羅德.
十二章四至五節.
亞伯拉罕蒙夫趙離開哈蘭來到迦南地.
十二章第四節說什麼呢?.
亞伯蘭照著耶和華的吩咐就去.
羅德也跟他一起去.
然後他帶著妻兒等財物來到迦南地.
兩叔侄我們可以假設.
叔父要照顧侄兒.
所以羅德就跟著.
十三章第一節.
就是我們剛才看到的經文.
亞伯蘭和妻子所有的一切.
以及羅德從埃及回到迦南地.
一直提到羅德的原因.
因為在亞伯拉罕的故事裡面.
由頭到尾有一個很重要的中心.
是神與亞伯蘭立約的繼承人.
那個Heir of Inheritance.
因為亞伯拉罕撒拉的太太沒有兒子.
我們自然會問.
那個繼承人相信應該是第一個.
是羅德.
這個也是當時古近東的文化習俗.
之後你記得一路的故事.
亞伯蘭就不斷去問.
到底誰是他的繼承人.
有說到是不是我的僕人.

$^{361}$以利以謝.
是不是夏甲所生的以實馬利.
然後到以撒才知道.
那個才是真正的繼承人.
然後故事一直繼續下去續長的故事.
說到雅各有十二種子.
到底誰是真正的繼承人.
是不是約瑟還是長子.
原來是第四個就是猶大.
每一個故事是在交代屬靈的祝福.
屬靈的約.
神的約的繼承者.
不單止繼承者是上帝親自用恩典諫算.
也很多時候說到子孫的繼承者有什麼特質.
以至他能夠繼承那個承約.
羅德是第一個最有潛質.
成為神揀選.
神恩典的繼承者.
他是最好的選擇.
但是他有沒有活出阿布拉罕後代.
阿布拉罕繼承者的生命.
有沒有那種眼光.
不斷看羅德的故事記載的時候.
你會發現聖經花了不少的筆墨.
來形容羅德的眼光.
他所見到的東西或動作.
我們看少少經文.
十三章第十一節.
羅德選擇了約旦河整個平原.
往東遷移.
你記得從創世記開始到現在.
往東遷移也不是好事.
但在原文特別說到.
他選擇了約旦河整個平原.
其實是選擇了自己.
為自己來選擇這個平原.
然後下一節.
尖尖內耳張膀.
直到索多瑪.
越來越接近罪惡之城索多瑪.

$^{401}$十四章十二節.
羅德就住在索多瑪了.
所以他不單單接近索多瑪.
他住在裡面.
索多瑪我們說過了.
是一個富庶的地方.
當時是一個很有收入.
田產豐盛的地方.
自然就惹起很多競爭.
於是第十四章就有打仗發生.
我們暫時不說很多詳情.
十四章說到.
在米索不達米亞有四個王入侵.
攻打索多瑪平原附近的五個王.
四個王打五個王.
他們應該是按照當時的王道.
那個叫Kings Highway.
聖經經常提到的.
就是一個國際南北的.
在約旦河東平原的商貿路線.
在那裡一直下來的時候.
就發生打仗.
然後十四章這樣交代.
四王與五王交戰的時候.
四王就將索多瑪.
俄摩拉所有的財物.
和一切的糧食都掠走了.
連擇義.
羅德.
羅德的財物也都拿走了.
當時他住在索多瑪.
住在索多瑪.
沒錯可能發財.
但是損失也很多.
差一點連命都沒有.
誰知有人來報信.
告訴阿卜蘭.
你的擇義有難了.
於是十四章這樣交代.
有一個人逃跑.

$^{441}$告訴希伯來人阿卜蘭.
阿卜蘭一聽到擇義被人擄去的時候.
他就調動他精煉的壯丁三百十八軍.
一直往北追.
追到他彈.
趁夜的時候就攻擊四王.
然後將他們潰敗.
追趕到甚至大馬士革不變的河壩.
將一切的財物都奪回來.
然後帶回.
最重要就是擇義羅德和他的財物.
和他的家人.
阿卜蘭就從希伯來人出兵拯救羅德.
羅德才可以拯救自己.
然後講到另外一節的故事.
十九章.
這裡又是關於羅德的故事.
今次上帝的審判.
真的臨到索多瑪和阿卜蘭.
派了兩個天使來滅城.
十九章第一節開始的時候.
講到晚上黃昏.
那兩位審判的天使到了索多瑪.
當時羅德在哪裡呢?.
羅德正坐在索多瑪的城門口.
坐在城門口.
是當時的長老的動作.
就是說不僅是羅德住在索多瑪.
成為他們的一分子.
甚至成為他們的長老之一.
中間一些緊張的情節我們不講了.
總之就是神的審判要臨到.
這兩個天使就去催羅德.
告訴他快點走.
如果你不走的話.
你就會沒命了.
誰知道在十六節.
羅德猶豫不決.
於是兩個天使因為要寬容羅德.
就拉著他的手和他的家人.

$^{481}$然後就帶他們出城外面.
然後到了十八節.
這裡有一段對話.
很有趣.
但是講到羅德自己的心裡面在想什麼.
羅德對兩個天使說.
我主你不要這樣.
你看見我在你眼前是蒙恩的.
我又讓你去到大師前外.
你救活我的性命.
可惜我還逃不到山上.
恐怕這個災禍會追趕到我死.
看了這座城很近.
他看著另一座城.
一座很小的城.
他問天使可不可以逃到那個城.
以至他的性命能夠得到存活.
那座城就是索爾.
二十二節.
至於漢奸.
你看見他的看法.
他的生命在想什麼.
由始至終.
羅德從來沒有意識過他能夠生還.
是因為其他人去因待他.
第一次生還.
羅德在索多瑪被四王擄掠的時候.
死不去.
是因為他叔父亞伯蘭.
黑夜來搶救他.
第二次生還.
上帝真的要滅絕索多瑪.
俄摩拉的時候.
要傾覆這個城市.
是因為什麼緣故羅德能夠得救.
十九章二十九節.
神毀滅那平原上的重城.
傾覆羅德所住過的重城的時候.
什麼原因.
神紀念亞伯拉罕.

$^{521}$把羅德從傾覆中救出.
是因為亞伯拉罕的原因.
羅德才能夠在傾覆中撿回性命.
這裡告訴我們這個角色.
他像不像一個立約群體.
神的子民的一部分.
這裡告訴我們.
羅德完全沒有一個屬靈的眼光.
亞伯拉罕和他相比.
是對立的.
我們是不是這樣.
看這個故事的時候.
我們有沒有發現我們自己像誰.
我們是不是好像羅德一樣.
不如我們看看羅德的結局.
這個結局在創世記十九章最後.
羅德因為怕住在索爾.
他就與兩個女兒一同離開索爾.
上了山住在那裡.
就在洞裡面住.
本身是想逃難到索爾一個小城.
於是他又怕.
上帝的審判又會臨到.
連那個城也滅了.
所以就躲在洞裡面.
這裡發生了什麼事.
就是兩個女兒有一個方案.
他們以為整個世界的男人都死了.
於是按照當時近東的文化.
女子正在想的是傳宗接代.
他們就想起沒有男人了.
於是剩下一個是父親.
於是就灌他喝酒.
使他喝醉了之後.
就跟他同睡.
以致他說什麼呢.
我們可以藉著父親保存後代.
所以兩晚都令父親喝醉了.
然後兩個女兒就跟父親發生戀戀.
三號節.

$^{561}$羅德有兩個女兒.
都從他們的父親懷了孕.
大的女兒就生了摩亞.
成為摩亞人的始祖.
小女兒就生了另一個兒子.
叫做辮那米.
就是阿滿人的始祖.
無論如何.
這是最後講到羅德的結局.
是近東男性最大的羞辱.
還有一件事.
你記得有些以其人之道.
還治其人之身的情節.
如果熟悉索多瑪俄摩拉.
被滅亡的故事第十九章的時候.
你記得有一個情節.
就是那些狂徒圍困羅德一家的時候.
在方亂裡面.
羅德出了一個什麼方案呢.
他說看啊.
我有兩個女兒.
都是處女.
我將她帶出來.
你們高興怎樣待她.
就怎樣行罷.
只是兩個天使.
他做一個所謂的家主.
想保護著他所接待的兩個客人.
他說你不可以在他身上做什麼.
這些是不是一個很精明.
想清楚的方案呢.
可能不是.
但在方亂裡面.
他就推他兩個女兒出來.
任人主治.
當然這個故事情節.
沒有發生這些事.
但到了這裡山洞裡面.
就要以其人之道還治其人之身.
就是反過來羅德的結果.

$^{601}$就是被兩個女兒任意對待.
連他自己都不知道.
另外羅德想套到一個小城.
叫做索爾蘇.
希伯來原文的意思.
其實是索瑞的意思.
是小城.
是一個很小的城.
所以中文翻譯得很貼切.
索就是很瑣碎.
怎樣說得好聖經的經文呢.
是用一種諷刺.
或者一種笑柄.
說羅德本身開始的時候.
就是一個聖若.
上帝祝福.
一個屬靈的繼承者.
到最後居然變成索爾.
是微不足道.
無關痛癢.
在拉比的文學裡.
曾經這樣說過.
說羅德和阿布拉罕.
相貌相似.
因為是兩叔侄.
但不過是表面上相似.
內心完全是判若兩人.
羅德以為把握住機會.
做一個機會主義者.
就出賣了自己的寧靜.
出賣了自己的生命.
自己的良心.
最後得到的結局就是羞恥.
你看羅德.
絕對不是大奸大惡的人.
但這可能是一個最大的提醒.
我們很多時候.
只不過好像羅德一樣.
我們不做大奸大惡的事.
那就算了.

$^{641}$但深入在我們的骨子裡.
很多時候被我們的慾望.
被我們眼前所見到的利益.
看為優先.
被那些利益來牽制.
阿布拉就完全判若兩人.
起碼在這些情節裡.
阿布拉首先想做的.
就是他願意慷慨地.
放棄迦林地裡最好的土地.
任由他的擇兒羅德來選擇.
一種好像放棄.
不去爭.
但最終他有沒有失去呢?.
反而他是得罪了.
十三章十四節.
羅德離開了阿布拉之後.
耶和華對阿布拉說.
你要舉目.
從你所看見的地方.
看見的地方.
我都要賜給你和你的後裔.
直到永遠.
我要使你的後裔.
多於地上的灰塵.
若能數算地上的灰塵.
才能夠數算你的後裔.
你起來中環走遍這地.
就是我賜給你這個地.
不是現在能夠兌現.
但阿布拉罕以信心來等候.
阿布拉罕代表著真正神的子民.
代表著我們.
如果要成為祝福.
成為萬國的祝福.
就要先學懂什麼?.
放下那種短視.
不要只想著.
我好怕得不償失.
我想把握住這個機會.

$^{681}$眼前的利益很重要.
這就是諾德.
阿布拉罕說.
跟隨耶和華的方式不是這樣.
其實你想想.
聖經的故事裡面.
這一次是第一次.
亦都繼續記載下去.
一次又一次告訴我們什麼?.
就是上帝的子民.
神的子民.
基督的門徒.
應該有什麼樣的標準?.
舉個很簡單的例子.
九曰大衛放棄能夠立事.
可以奪取王位的機會.
不殺掃弩.
最後得到什麼?.
不是失去王位.
反而是神賜給他王位.
神去應許他.
他得到榮耀.
到新約.
主耶穌自己.
興翰十字架的羞辱.
最後得到什麼?.
不是羞辱.
反而是得到上帝給他的榮耀.
其實耶穌對門徒這樣說.
都是一樣.
我們結束的時候.
看馬達科罕第十六章.
馬達科罕第十六章說到.
耶穌對門徒的挑戰.
如果有人願意跟隨我.
就當寫記.
背起他的十字架來跟隨我.
凡要救自己生命的.
就必喪掉生命.
但為我捨棄生命的.

$^{721}$犧牲的就必會得著生命.
人若賺得全世界.
卻賠上自己的生命.
有什麼好處?.
人魂能拿什麼?.
換取自己的生命.
求主幫助我們.
成為一個真正神的子民.
有上帝屬靈的眼光.
有信心的眼光.
讓我們所換取的.
是永恆的生命.
是永恆價值的東西.
我們同心低頭祈禱.
主我們求你來恩待我們.
因為我們的眼光很多時候都是短視.
我們所見到的很多時候都是前面.
試探給我們的.
或者機會給我們的.
一些短暫的利益.
主幫助我們.
我們的眼睛是望見你的影響.
望見你所賜的祝福.
求主你幫助我們.
我們能夠作為真正的門徒.
去捨棄.
去見到你自己的眼光.
見到你的事物.
見到你所珍貴,有價值,有永恆價值的事物.
主讓我們繼續去追求.
在我們信心的旅程上.
我們不斷地去更新我們的心意.
更新我們的眼光.
以致我們追溯的是你自己的祝福.
追求的是你天國的一切.
祈禱奉主耶穌的名字.
\newpage



\section{}
\label{sec:QTJ66JH0gPw}
\textbf{舊約網上行 - 上帝的足跡: EP9 凱旋回歸}
\newline
\newline
連結: \href{https://youtube.com/watch?v=QTJ66JH0gPw}{\texttt{https://youtube.com/watch?v=QTJ66JH0gPw}} ~~~~ 語音日期: 2021-05-03
\newline
\newline
\hyperref[sec:Ks1q0nCc4X4]{\small{< < < PREV SERMON < < <}}
~
\hyperref[sec:index_chronic]{\small{[返順時目]}}
~
\hyperref[sec:index_scriptual]{\small{[返順卷目]}}
~
\hyperref[sec:afSlbD_al4Y]{\small{> > > NEXT SERMON > > >}}
\newline
\newline
$^{1}$歡迎來到舊日網上行的第九集.
我們繼續追蹤神在阿布拉罕身上有的事跡.
上一次我們看到阿布拉罕和羅德的故事.
今天我們繼續看創世記第十四章.
這裡說到四個近東的王.
來攻打迦南的五個王.
之後就將羅德,即是阿布拉罕的侄子,擄去.
阿布拉罕出兵相救.
然後他海船回歸回來的時候.
有撒冷王來迎接阿布拉罕.
這個撒冷王叫做麥基西德.
而阿布拉罕也獻上十分之一.
怎樣看這段經文也好.
我們也會感覺到有些離奇.
並且這段經文引起很多不同的爭議.
第一方面的爭議就是歷史性的爭議.
我們先看一下創世記第十四章.
第十四章一至第四節.
當暗拉非王作士拿王.
阿洛作以拉薩王.
基達魯瑪作以蘭王.
提達作戈印王的時候.
他們興兵攻打.
所多馬王比拉.
阿摩拉王比薩.
阿馬王是拿.
洗遍王善爾別和比拉王.
比拉就是所爾.
這五個王在西靈谷會合.
西靈谷就是嚴海.
他們服侍了基達魯瑪已經十二年了.
到第十三年就背叛了.
這是一段非常獨特的經文.
因為創世記的經文裡面.
很少會提及這麼多的君王.
好像給我們.
一個.
當時一些歷史的人物.
你這樣想.
就說既然有這麼多歷史的人物.

$^{41}$我們只要找到經外的文獻.
提到這些王的名字.
我們就知道阿巴拉罕的歷史性.
和他的事跡和時間等等.
因為.
我們大概知道.
阿巴拉罕是屬於公元前.
二千至一千八百年之間的人物.
所以考古的歷史.
如果在這段時間.
所謂青銅時期.
中期.
找到一些經外的文獻.
提到這些名字的時候.
我們就似乎可以再一步.
縮窄.
到底知道阿巴拉罕是屬於誰的時期.
的確有很多學者.
曾經花過不少功夫.
譬如士拿這個地方.
士拿王.
巴比倫.
以拉薩這個地方.
應該是阿述.
因為古代文獻提到阿述這個地方.
是叫做阿拉薩.
即是阿述這個字.
以蘭應該是今天的伊朗.
戈印就比較濃濃.
戈印.
大概應該是土耳其一帶.
因為那裡的王名.
或者客人帝國的王名.
很普遍.
都是提達.
或者用提達來開頭.
這裡這麼多的王.
其中一個最有機會.
用經外文獻來比較.
就是阿姆拉菲.

$^{81}$阿姆拉比.
如果你翻譯出來.
不拿響音.
就好像巴比倫王.
這個阿摩利人.
漢莫拉比.
阿姆拉比的名字相似.
這個是最有可能.
但學界仍然有很多的爭議.
所以未能夠.
改觀定論.
講完近東四個王.
又講一下五個王.
他們是屬於死海.
嚴海一帶.
索多瑪,阿姆拉等等的城邦.
到今天我們都很難去確定.
不過在經文裡面.
正如上一次講過.
從地理.
從一些描述.
我們大概知道是哪一個地方.
無論如何.
都引起很多歷史性的問題.
不過我們一定要承認.
我們對這段時期.
所謂Middle Bronze Age.
我們仍然是非常模糊的.
有很多的缺口我們仍然未認知到.
所以雖然聖經.
裡面一些的細節.
不能夠和一些經外文獻.
去配合.
吻合.
是不足為奇的.
但是經文的資料.
地名地理.
描述等等.
其實都是非常寫實的.
意思就是說大體上是.

$^{121}$吻合這整段時間.
或者這樣說.
看完經文的時候.
我們知道應該是屬於當時.
熟悉當時政局.
環境的人才能夠寫得出.
第二方面.
我們可以問.
在聖經的經文.
脈絡記載裡面.
為什麼在這裡會特別記載.
創世記第十四章.
打仗這回事.
經文的讀特.
給我們讀者的印象.
就是有一點格格不入.
為什麼無端端.
要提到這件事.
如果我們仔細.
留意這個故事的.
用詞的時候.
上文下理我們會發現.
作者是刻意.
有心的.
將阿布拉罕.
擁有索多瑪和.
羅德的故事.
配合在這裡.
加上麥基駛德.
迎接阿布拉罕的故事.
跟其他的.
經文的結構.
是很有關的.
我們看過整個阿布拉罕的.
故事的敘述.
那個結構是一種三文字.
的結構.
ABCDE然後再反過來.
去記載.
第一對.

$^{161}$前後互相呼應.
第二和尾二.
互相呼應.
所以今天我們看到創世記第十四章的時候.
發現是屬於.
D部分.
和十八十九章.
有一些回應.
所以如果我們一對來看的時候.
創世記十四章和.
十八十九章.
D的部分.
我們發現有些用詞.
是共通的.
古代的作者.
沒有我們今天的作者.
那麼方便.
我們可以強調一些內容.
用螢光筆highlight.
或者在電腦的文件.
用boldface等等.
去強調某一些的重點.
古代的作者.
用什麼方法去強調呢.
就是用一些重複.
相似發音的字眼.
等等的方法.
去將一些共通的地方.
重要的地方.
強調出來.
我們列出某一些兩段經文.
共通的詞彙.
第一方面.
所以這個地名.
Sor.
十四章看過.
十九章第二段落也看過.
鹽這個字.
就更加清楚了.
剛才讀經的時候已經看過了.

$^{201}$十四章第三節.
西定谷就是鹽海.
Sea of Salt.
十九章第二段.
跟他一起看的.
就是羅德的妻子.
在他後面向後一望.
離開所多瑪的時候.
不捨得.
財物於是就變成鹽柱.
第三樣.
就是逃到山上.
這個動詞.
形容所多瑪王戰敗之後.
十四章第十節.
向山上.
逃跑.
十九章也是一樣.
這裡說到羅德離開.
逃離所多瑪的時候.
他說.
我還逃不到山上.
怕著災會追上.
所以同一個動詞.
兩段經文段落也出現過.
另外.
摩利象樹.
這個地方.
十四章十三節.
摩利象樹那裡.
十八章也說.
向阿巴拉罕顯現的時候.
摩利象樹那裡.
第五方面.
公義這個字.
這個很清楚.
因為撒冷王十四章十八節.
說給我們聽.
就是公義.
然後.

$^{241}$神的三個使者.
要去查實.
所多瑪的罪行的時候.
要快去施行審判.
又說十八章十九節.
他說我揀選了阿巴拉罕.
也吩咐他的子孫.
他的家屬要去遵行我的道.
要去秉行公義.
好叫我.
應許阿巴拉罕的話可以實現.
秉行公義.
這個字出現了.
也接著出現了.
七次.
因為這裡說到阿巴拉罕.
為所多瑪祈求的時候.
他說如果有義人的話.
神豈會.
將惡人和義人一起除滅呢.
義字.
義人這個字就出現了.
七次.
第六方面.
十這個數目也在兩段經文裡面.
各自會出現.
十四章第二十節.
阿巴拉罕將所得的一切.
這些戰略的略物.
拿出十分之一.
來吸了.
默記誓得.
十八章第二段也出現了.
阿巴拉罕說請我主不要動怒.
用我.
問者最後一次.
假如在這個城.
所多瑪找到十個義人.
又如何呢.
你會不會立刻除滅他呢.

$^{281}$神回答我恩者十個義人.
的緣故.
也不會毀滅.
所以十這個數字.
也在兩段裡面刻意出現.
最後第七樣可以提到的.
就是接待的情節.
有hospitality的.
內容.
十四章第十七至二十四節.
這裡最清楚就是阿巴拉罕海船回歸之後.
回到迦南地.
默記誓得王.
應該是耶路撒冷.
出來迎接阿巴拉罕.
稱頌他.
阿巴拉罕就獻上十分之一.
十八十九章就出現了.
兩次的接待的情節.
十八章開始.
阿巴拉罕接待三行客.
所謂三個神的使者.
十九章當使者.
去到所多瑪的時候.
羅德就去接待神的使者.
所以換句話說.
我們可以將十四章.
阿巴拉罕.
用來救羅德的海船回歸的戰事.
是和十八十九章.
可以一起.
合起來去參考.
我們看完.
兩段的經文.
十四章和十八十九章.
是有共通的地方.
也似乎在結構裡面.
作者要我們將兩段.
合起來去看.
這樣的話我們就可以.

$^{321}$問多一個問題.
這段經文有什麼意義呢.
十四章的內容是很直接的.
那個戰事.
是發生在應許之地.
以外的.
上一次我們講過了.
羅德自己向東走.
向前走.
離開神的應許之地.
離開阿巴拉罕的範圍.
地圖可以看到.
四王是從近東來到的.
攻打五個王.
包括所多瑪和摩賴.
將五個王的財物.
拿走.
連同所多瑪的羅德.
也被俘虜.
之後發生了什麼事呢.
創世記第十四章十四節.
阿巴蘭一聽見.
擇兒被人擄去.
就抽掉他家裡生養.
精煉的壯丁.
三百一十八人.
一直追到達.
阿巴蘭和他的僕人.
趁夜分隊攻擊他們.
結果擊敗了他們.
又追趕他們直到大馬士革.
不變的河把.
將一切財物都奪回來.
也把他的擇兒.
羅德和羅德的財物.
以及婦女和族人.
都奪了回來.
這個故事很多時候.
在聖經的大標題.
都會叫做阿巴蘭罕永救羅德.

$^{361}$但其實這個.
你想想只是故事的.
一半.
另一半其實阿巴蘭罕也救了.
所多瑪的眾王.
將他所有的掠物財物.
都拿回來了.
之前我們不是碰過.
見過所多瑪.
是一個什麼地方.
創世記13章13字.
在耶和華看來.
所多瑪人是罪大惡極的.
但是在這書.
我們見到這個故事.
神居然藉著.
阿巴蘭罕來拯救.
裡面的人.
拯救裡面的人的生活.
將掠物拿回來.
這是拯救裡面.
告訴我們神給了.
所多瑪多一段時間.
誰知道羅德沒有離開.
罪惡之城.
所多瑪王眾王更加不用說.
不單止沒有離開.
更加沒有.
稱頌耶和華.
完全意識不到.
原來是藉著阿巴蘭罕.
他們才拿回他們有的東西.
但是神的審判.
也沒有立即來到.
看第二篇的故事.
第十八章.
阿巴蘭罕去禱告.
跟耶和華說.
因為知道耶和華要去審判.
所多瑪.

$^{401}$他就說到最後.
十個義人在城裡面.
你會不會滅城.
神當然回答.
不會貿然將義人.
和惡人一起去除滅.
但是.
我們看到整個故事的發展.
所多瑪的結束.
就是到最後.
限期的時候.
神的審判.
必然會來到.
因為阿巴蘭罕的神.
這位上帝就是審判全地的主.
他不單單是迦南地.
審判迦南的神.
也是審判每一個地方.
每一個民眾的神.
這位公義的主宰.
就是審判全地.
今天.
我們在世界不同的地方.
找到很多.
看到很多不義的事.
義人.
好像越來越少.
惡人越來越囂張.
有時候.
我們看到.
有些地方.
是方華連篇.
是非顛倒.
見不到公義.
我們有一種的試探.
我們就會說.
何必要去到.
持守真理公義.
或者良心.
誰作王.

$^{441}$我們就靠誰.
哪裡有錢.
我們就靠近那一處.
像羅德一樣就是機會主義者.
哪裡有錢.
我們就去所多瑪.
故事的脈絡講得很清楚.
阿巴蘭罕在.
祈禱和神對話的時候.
十八章第二十五節.
審判全地的主.
豈可不行公義嗎?.
我們活在一個.
難見公義的世代.
我們不要忘記.
上帝依然是審判全地的主.
他仍然會施行公義.
問題就是這樣.
我們願不願意去信教他.
去等候他.
繼續去堅持.
我們去行公義.
去相信主.
第二方面的意義.
似乎是指向.
特別顯出阿巴蘭罕的財力.
他的勢力.
十四章第十七節開始.
阿巴蘭攻擊了基大魯馬和他聯盟的王.
回來的時候.
所多馬王出來.
在沙美谷迎接他.
沙美谷就是帝王谷.
撒冷王麥基斯德也出來.
帶著餅和酒.
他是至高神的祭祀.
他給阿巴蘭祝福說.
願創造天地的主.
至高的神.
賜福給阿巴蘭.

$^{481}$請你把人交在你手裡的.
至高的神是應當稱頌的.
阿巴蘭就把所得的一切.
拿出十分之一來.
給了麥基斯德.
二十一節.
所多馬王對阿巴蘭說.
請你把人交給我.
至於財物你拿去吧.
阿巴蘭卻對所多馬王說.
我已經向創造天地的主.
至高的神耶和華.
舉手起誓.
凡是你的東西就是一條線.
一條鞋帶我也不拿.
免得你說.
我使阿巴蘭發了財.
我什麼都不要.
除了僕人吃掉的以外.
但與我同行的阿賴.
以實國萬利.
所應得的份.
讓他們拿去吧.
這裡說到阿巴蘭永久.
海船回歸.
或者.
近東的文獻.
可以幫助你解答一下.
為什麼聖經的作者.
這麼刻意.
寄述阿巴蘭海船回歸的戰績.
記得我們曾經說過.
一個近東的.
大的古城.
叫做馬里馬里.
位於今天的.
敘利亞北面.
記得我們說過.
在創世記十二章的時候.
阿巴蘭離開吾爾.

$^{521}$去哈蘭的時候.
經過馬里的地方.
馬里的大城.
是一個很重要.
很有策略性的地方.
是一個很富裕的地方.
阿巴蘭海.
由吾爾去到哈蘭的時候.
應該是經過.
馬里這個地方.
當時的時候.
大概是主前一千八百年左右.
馬里是一個.
阿摩利人的城邦.
你也記得我們說過.
阿巴蘭海.
也可能是阿摩利人之一.
這個故事.
有關馬里是這樣的.
1933年的時候.
八月的時候.
有一群敘利亞的農夫.
想去埋葬他剛剛去世的親人.
在這裡他們想.
去發掘一個墓地.
誰知道.
掘一掘的時候.
發現了一些古代的文物.
一些神像等等.
就通知當時.
法國的軍部.
因為當時是法國佔領著敘利亞.
消息就傳回巴黎.
在羅浮宮.
博物館.
就立刻.
派考古的學者.
去到當地.
做一些探測的工作.
之後很快他們就開始正式考古的發掘.

$^{561}$然後發現了什麼呢.
就是發現了馬里的.
大皇宮.
有三百多間的房間.
也有保存著.
一些濕漆的壁畫.
最重要的.
發現之一.
就是有二萬二千塊的.
褻瀆文字的禮版.
學者就將這些禮版.
叫做馬里皇室文獻.
Royal Archives of Mali.
時至今天只有幾千塊的禮版.
被翻譯成為.
現代的語文.
還有很多工作要做.
馬里文件要提的.
為什麼呢.
因為這些是阿摩利人族長和軍王的寫作.
是一些外交的書簡.
或者一些.
政治社會裡面的文件.
很多時候就講到一個真正的族長.
當有的質素.
他是一個英雄.
一個帶領著國家的人民的英雄.
是一個怎樣的人.
就是一個領兵出征的戰士.
一個warrior.
一個願意出去.
帶領著他的子民.
去打仗的一個人.
所以在阿摩利人的文化裡面.
就講到一個真正的英雄.
就是會去打仗出征.
並且將掠物拿回來.
所以在創世記十四章的時候.
就提到阿巴拉罕的時候.
就特別將.

$^{601}$阿巴拉罕這場戰事放進去.
就說明.
對當時的文化來說.
阿巴拉罕就是當時的.
真正的族長.
真正的軍王.
真正的王者.
所以當阿巴拉罕回來的時候.
眾王都承認他.
是才雄細口.
是很有實力的.
經文裡面.
就記載了兩個王的反應.
我們先講一下.
所多瑪王.
當他見到阿巴拉罕.
帶著所有的財產回來的時候.
他就說.
我是想得回人口.
當時的人口是多麼重要.
但所有的掠物.
他知道是阿巴拉罕應得的.
他就說不如你拿去吧.
你留著它.
但是阿巴拉罕做了什麼呢.
十分豪氣.
二十二節.
阿巴拉罕對所多瑪王說.
你們的主 至高的神.
耶和華舉手起誓.
凡是你的東西 一條線 一根鞋帶.
我都不會拿的.
免得你說我是阿巴拉罕.
發了財.
我什麼都不要 除了我僕人所吃過的東西.
或者我同行的盟友.
他們應得的份.
讓他們拿走.
這裡很顯出阿巴拉罕.
不是靠自己的能力.

$^{641}$不是靠自己的財力.
而是他對耶和華的信靠.
以致他能夠得勝.
所以在所多瑪王的面前.
他認定.
使他封城的.
只有一位.
不是地上的君王.
而是創造天地的主 耶和華.
另一個王有什麼反應呢.
這個王叫做默基西特.
剛才我們說過.
這是一個很奇特.
很突然出現.
有突然消失的一個人.
特別是在舊約的經文.
或者猶太人的思想裡面.
特別是後期的經文.
就做了很多神學的擴充.
expansion.
又加了很多想像.
imagination 特別到了新約的時候.
死海古卷等等的時候.
加上新約的希伯來書.
就會認為.
默基西特是一個真正的王.
真正的祭司 也是基督的儒表.
但是如果我們從.
創世記第十四章.
來看的話.
這個人不是很特別.
默基西特就是.
公義的王.
這是他的名字 或者是他的綽號.
其實他本身只是一個迦南人.
在耶路撒冷.
來做.
他的勢力應該.
很強大.
在2000年的時候.

$^{681}$有一位以色列的.
考古學家叫做.
他在耶路撒冷發掘了幾十年.
2000年的時候.
他在耶路撒冷.
發現了一個大石的.
城樓 有16米.
乘16米.
這些石牆最後有.
7米 如果你.
去過耶路撒冷的話.
你去到.
基奮泉 中間一段經過.
看到這些大石陣.
這兩處.
大的石牆.
是平衡從山頂向東.
一直伸延到山腳.
作用是.
保護著耶路撒冷.
唯一的水泉 就是基奮泉.
基漢泉.
兩處大的石牆中間.
有一個窄的通道.
使城裡面的人.
不需要出去城牆.
以外 去到山腳.
的基奮泉來打水.
他們可以經過.
這個窄的通道.
安全地被城牆保護.
特別到戰亂的時候.
更加有用.
這個城牆.
是屬於麥基斯德.
阿布拉罕時期的耶路撒冷.
所以告訴我們.
這個Malchizedek.
其實是一個非常強大.
很有勢力的.

$^{721}$本地的郡王.
本地的迦南王.
當他看到.
阿布蘭海船.
回歸的時候.
迦南的王看到了什麼.
十八字.
撒冷王麥基斯德也帶著餅和酒出來.
他是至高神的濟士.
他給阿布蘭祝福.
願創造天地的主.
至高的神賜福給阿布蘭.
把敵人交在你手裡.
至高的神是應當稱重的.
阿布蘭就將所得的一切拿出十分之一.
給了麥基斯德.
在麥基斯德.
這個迦南本地的王眼中.
他就把所有的東西.
拿出十分之一.
在迦南本地的王眼中.
他就不單單看到.
阿布蘭的勢力有多大.
而是他看到阿布蘭海船.
背後的神.
是更大.
至高神.
在迦南的宗教裡面.
經常出現的一個稱號.
這個稱號.
是一個至高神的稱號.
換句話說.
這裡看到麥基斯德.
這個迦南王.
他就在這裡表明.
阿布蘭海船的神.
是一個獨一的主.
一個萬有萬聞的主宰.
他看到阿布蘭的德性的時候.
他不是看到一個人有多威風.

$^{761}$而是看到.
阿布蘭海船背後的神.
是有多威風.
賽夫.
可不可以讓人.
讓我們自己.
去親近神.
去敬拜神.
還是會讓我們.
更加遠離神.
你有沒有留意.
在這段故事裡面.
經常出現一個字.
我讀幾段聖經.
看看你能不能找到重複的字.
有一個強調的地方.
重點在哪裡.
十四章第十一節.
四王就把索多瑪昂波拉所有的財物.
以及一切糧食都拿走.
第二節十四章十二節.
連阿布蘭的侄子.
羅德和羅德的財物.
也帶走了當時.
羅德主宰索多瑪.
第三段.
十四章十六節.
將一切財物都奪回來.
也把他侄子羅德和羅德的財物.
連同.
父女族人.
都奪了回來.
十四章二十一節.
應該越來越清楚了.
索多瑪昂對阿布蘭說.
請你把索多瑪昂的財物.
交給我.
至於財物你拿去吧.
你看到.
經常出現一個字.

$^{801}$就是財物.
財產.
其實之前兩章已經出現過.
十二章.
當阿布蘭來到迦南的時候.
那裡怎麼交代呢.
阿布蘭帶著他的妻子.
侄兒羅德.
以及他在哈蘭所積蓄的一切.
財物.
獲得的人口.
一同帶來到迦南地區.
十三章第六節.
之前.
保守住阿布蘭的一家.
保守住薩萊.
十三章第六節說什麼.
那地容不下他們.
阿布蘭和羅德住在一起.
因為他們的財物很多.
又是說到財物.
又是說到財產與財物.
又是說到財產.
越來越多.
因為上帝的恩典.
因為上帝的帶領.
以至阿布蘭漢的財富.
是越來越多.
在疫情的一年.
我知道有很多人.
都是受損失.
教會受到損失.
不只是我們教會.
也不只是各地的教會.
因為疫情困難的緣故.
受到損失.
最近我們也聽到.
城市其他的基督教機構.
一些的士工.
都是入不敷出.

$^{841}$因為奉獻.
不穩定等等.
疫情帶來的因素.
有很多人在疫情裡面.
有損失.
但也可能在疫情裡面.
你是得益的那一班人.
可能股市升.
你賺幾個開.
甚至我認識的一些人.
跟我說.
不知為何縱使疫情困難.
但卻沒有損失.
甚至可能股票.
一直升上去.
多了財力.
多了勢力.
我們到底.
怎樣去自處.
可能像.
所多馬王.
保存自己的財產.
暫時保存.
最後的結果是什麼.
也是滅亡.
還是.
我們面對著.
財富多了的時候.
特別在疫情困難裡面.
我們發現原來我們未必是.
那一班最受損失的人.
反而我們會.
多了財產.
但這些財產.
是否可以令我們敬拜神.
還是讓我們遠離神.
或者這麼說.
我們這些財產.
我們越來越富強的時候.
我們是否真的.

$^{881}$可以讓人看到.
出名的那個.
不是我們自己.
是上帝自己.
主去幫助我們.
讓我們有上帝的保守.
有上帝的看顧.
去供應的時候.
沒有損失.
反而可能財物多了.
讓我們去到.
貢獻在神的國.
使神的名.
在我們的生命,財富當中.
更加能夠.
彰顯出來.
我們同心低頭祈禱.
上帝.
你就是那位創天造地.
你都祝福我們的神.
讓我們知道我們今天所享有的.
恩典,所享有的祝福.
來去敬拜.
來去告訴.
人們知道你是那位.
使我們富足的神.
你是那位唯一一位.
創天造地的主.
你是那位唯一一位.
使我們生命能夠富足的上帝.
主叫我們來敬拜你.
叫人看到.
我們的生命,我們的.
勢力越來越大.
越來越穩定的時候.
越來越富強的時候.
讓人看到我們背後.
是你自己.
是你的名.
是應當稱頌的.

$^{921}$求主你親自來幫助.
你自己的子民.
在各地不同的弟兄姊妹.
主我們有著你的恩典.
大靈當中.
讓我們來去.
真正敬拜,奉獻給你.
祈禱奉主耶穌的名字.
\newpage



\section{7:13-14}
\label{sec:afSlbD_al4Y}
\textbf{Kingdom Values: Sermon on the Mount L Matthew 7\_13-14}
\newline
\newline
連結: \href{https://youtube.com/watch?v=afSlbD_al4Y}{\texttt{https://youtube.com/watch?v=afSlbD\_al4Y}} ~~~~ 語音日期: 2021-05-18
\newline
\newline
\hyperref[sec:QTJ66JH0gPw]{\small{< < < PREV SERMON < < <}}
~
\hyperref[sec:index_chronic]{\small{[返順時目]}}
~
\hyperref[sec:index_scriptual]{\small{[返順卷目]}}
~
\hyperref[sec:_taiMeC2wK0]{\small{> > > NEXT SERMON > > >}}
\newline
\newline
$^{1}$Good morning family and welcome. I'm glad and excited that we get to open God's.
Word together again this morning. I want to remind you as always that we'd love.
to hear from you. We'd love for you to connect with us. You can reach out, send.
us an email at info@urbanvoicechurch.com. If there's any way that we can pray for.
you and encourage you, we'd love to know that. Or if you just have comments or.
questions, we'd love to hear those as well. So please reach out to us. We'd love.
to hear from you. This morning we're going to be looking at Matthew chapter.
7 verses 13 and 14. These are really significant verses in the Sermon on the.
Mount and really significant verses in our lives and in our faith and our.
journey as followers of Jesus. And they're important in understanding the.
application of the values of the kingdom to our lives. So I want us to examine.
these closely today and to really see what God has to say to us this morning..
Let's read these together. Matthew chapter 7 verses 13 and 14. "Enter by the.
narrow gate, for the gate is wide and the way is easy that leads to destruction..
And those who enter by it are many. For the gate is narrow and the way is hard.
that leads to life. And those who find it are few." Now what is Jesus saying here?.
What's going on in this passage? The metaphor is pretty obvious here. He's.
talking about two paths that are very different. One that is narrow and.
challenging. Another one that is broad and wide and easy. And Jesus is pointing.
us to an important truth about what it means to be a follower of Jesus, to be a.
subject of the kingdom, and to live out the values of the kingdom. At this point.
in the Sermon on the Mount, this is towards the end, Jesus is beginning to.
kind of wrap up the sermon and he's giving us some application. So this.
section, these verses, are really speaking to what it looks like for us to begin to.
live out the values of the kingdom that he's already been speaking about for the.
last almost three chapters of the book of Matthew here. And Jesus is giving us a.
description in this metaphor and it's almost a warning of what it looks like.
to live out the values of the kingdom in our world. That it's not going to be easy,.
it's not going to be simple, it's not even always going to feel great. And Jesus is.
telling us that. And I think it's important to understand that Jesus is.
telling us that up front because Jesus was never trying to talk us into doing.
something. Jesus was never trying to pull a bait-and-switch. The gospel has always.
been up front about the commitment that is involved. And sometimes us, Christians,.
believers, followers of Jesus in the church, when we share the gospel, we.
kind of want to hide that. We kind of want to obscure that and make it seem.
really easy to people and say, "Hey, it's going to be great. You'll get everything.
you ever want." And that is true. The rewards of the gospel are indescribable..
They truly are. But that doesn't mean that it's easy. There is very real.
commitment involved and there's often suffering involved. And that's what Jesus.
is telling us here. Jesus is being very honest, very direct, very up front with us.

$^{41}$about what it means to live out the values of the kingdom in our broken.
world. And that it's not going to be easy. It's not going to be simple. And Jesus is.
telling us some hard truths here. And I love that Jesus does that. There's so.
many beautiful things that Jesus tells us. So many beautiful things in the.
Gospels. So many beautiful things right here in the Sermon on the Mount that.
Jesus tells us and describes for us. But Jesus doesn't just tell us the good.
stuff. He also tells us the challenging things. Jesus also says to us what we.
need to hear but don't want to hear. Jesus tells us the hard truths and I.
love that. And that's what he's doing in these verses. You know God's Word is.
oftentimes encouraging and comforting to us. But it's also oftentimes challenging.
to us. Sometimes it's both at once. And this is one of those passages that is.
challenging. There is comfort and encouragement in it as well. We'll see.
that today. But it's challenging. It is really challenging. So in these verses I.
think there's some important things for us to understand that Jesus is.
communicating. And first of all I want us to consider the fact that Jesus is.
telling us the narrow way is not easy. That's what Jesus is saying. The narrow.
way is not an easy way. You know part of the reason the narrow way is not.
easy is that it's in opposition to our natural desires and selfishness. Following.
Jesus, living out the values of the kingdom, is really in opposition to the.
things that we naturally want to do. We've been talking about this for so.
long now as we study these values of the kingdom. But living a life of humility, of.
poverty of spirit, that's in opposition to our natural desires which are.
towards pride and self inflation and self promotion. Things like that. Being.
humble and living out poverty of spirit is in opposition to our selfish desires..
So many of the values of the kingdom really oppose our natural desires and.
for that reason it is not easy to walk in the narrow way. It is not easy to live.
out the values of the kingdom. And the reality is that this verse communicates.
that there are times when when the narrow way feels closed in. It feels.
tight. It feels difficult to maneuver through. In verse 13 this word for narrow.
the Greek word gives us this this idea of narrow or straight. And narrow or.
straight because there are obstacles on the sides. It's kind of you can think of.
like a narrow corridor through a stone passageway in a castle or something like.
that. Or you can think of maybe walking on a trail. I don't know if you've ever.
hiked a trail like this but I've hiked on some trails where you're walking.
through cliffs kind of and you're on a trail where it's just straight wall on.
this side and straight wall up on that side. It's very narrow and very closed in..
And that's part of what the Greek word in verse 13 is communicating. That there.
are times when it's gonna feel tight. It's gonna feel difficult to navigate..
It's gonna feel difficult to go on through because it's not necessarily.
easy. And Jesus is telling us that. He's being up front with us about that. And.

$^{81}$this concept of the narrow way also communicates the fact that there is.
suffering involved. In verse 14 where it says for the gate is narrow. That's.
actually a completely different Greek word from the word in verse 13. And the.
word in verse 14 is a word that communicates suffering. It's a word that.
is translated in other parts of the New Testament as affliction or as trouble or.
as suffering tribulation. And so it's saying that this gate is not just.
narrow. It's difficult. That there's suffering involved. This path is.
challenging to travel. It's hard to travel at times. Just like verse 14 says.
the way is hard. It's not easy. And Jesus is being up front with us and telling us.
that. That if we are going to commit to live out the values of the kingdom. If.
we're gonna commit to really being a follower of Jesus. If we're gonna commit.
to really being serious about obeying him in our life. Then it's not always.
gonna be easy. It's gonna be in opposition to our own natural desires..
And there are times when it's gonna feel tight and closed in. And we're gonna.
feel like it's hard to push on through. But there's also times when it's gonna.
involve suffering. And when it's not gonna feel good in the moment. And when.
it's going to hurt. And it's going to involve affliction. That word in.
verse 14. The Greek word as also carries this idea of being pressed. Just like.
grapes were pressed to extract the grape juice. It's a word that gives us.
that word picture of being pressed. So that something is gonna be extracted..
And sometimes that's what following Jesus means. God lets us go through.
difficulty and challenges and affliction. So that he can extract impurities and.
things that are not healthy from our life. And fill us with goodness. So.
that he can refine us. And there is this concept of suffering that's carried in.
these words and in these verses. Another thing we see in these verses is that.
this is a narrow gate and a narrow way. And that's important for us to.
understand. Because sometimes we might think of following Jesus as like a.
funnel kind of. Or maybe as an uphill climb. And we might think of it as it's.
wide open and easy in the beginning. And everything's great. And then as we get to.
more advanced levels it gets a little bit tougher. Or that's when it gets.
challenging. And it is an uphill climb. The journey of faith. Of following Jesus..
So that's not necessarily a bad metaphor. But we need to understand that it's.
narrow and straight from beginning to end. And Jesus is letting us know that in.
this verse. That there's only one way to heaven. Jesus himself said that he is the.
only way to the Father. He's the way the truth and the life. It's narrow in the.
beginning. It's narrow to get in. It's only by God's grace and by faith in.
Jesus that we can come into salvation. It's a narrow way. It's pretty.
exclusive to get in that narrow gate. But then as we travel on it, it's still.
narrow. It's not something that as we grow more advanced then maybe we'll.
have an experience of challenge or affliction or trouble. No. God wants to.

$^{121}$grow us from the very beginning. And he's gonna allow adversity in our lives to.
stretch us and to grow us from the very beginning. And so we may say, "Well I.
wasn't expecting this. I wasn't ready for this. Sure I thought maybe someday in my.
journey of following Jesus that I'd have to deal with something challenging or.
something difficult. Or God would let me face something that was was difficult..
But I wasn't expecting this right now because I'm young or I haven't been a.
believer very long." But it's not a funnel. It's something that from the beginning.
God wants to grow us. And God wants to challenge us. And God wants to stretch us..
And we need to understand that. And it doesn't have moments in the middle where.
we just say, "Okay I've been on the narrow way for a while so now I get to just.
take some time off and float." And that's not there. It's always going to be a.
narrow way beginning to end. Now that doesn't mean that it's not full of.
blessing and reward and joy because it is. We're gonna discuss that a little bit.
later on. So don't get the wrong idea that following Jesus is just some type.
of difficult drudgery and it's always just sad and hard. That's not the.
reality of what it means to be a follower of Jesus. But understand that.
the way is narrow and it is from beginning to end. You know it requires.
commitment. And that's why Jesus is telling us this up front. Why he's.
telling us this in the beginning is that it requires commitment for us to step.
onto the narrow way, to choose to follow him, and to stay on the narrow way. To.
continue to follow him. Jesus is speaking to people that were very familiar with.
faith, very familiar with the Old Testament, very familiar with God. But.
they weren't necessarily familiar with the values of the kingdom that he was.
giving them. And he was being honest with them up front and saying, "If you want to.
step into this life, if you want to begin to live out these values and.
follow God in the way that I'm calling you to, it requires commitment. And you've.
got to understand that." And that's just as true for us. We need to understand.
that there's a commitment required to enter the narrow gate and to walk on.
the narrow way. And that brings us to the next thing I want us to consider. And.
that is that the narrow way requires intentional pursuit. The narrow way.
requires intentional pursuit. Now this is an obvious point and it's not something.
that I need to belabor or explain a lot because it's pretty clear. But I just.
want to bring it out for us and remind us that if you are not intentionally.
staying on a narrow path, then you're off the path. I think a good illustration is.
if you think of maybe being on a small path in the forest that's not real well.
marked. I've done a lot of hiking. I love to hike with my family and we do that a.
lot. And we've hiked a lot of different types of trails. And some trails you hike.
on that are pretty wide and maybe they're even a gravel trail or they're.
partially paved. It's really easy to see where the trail is and stay on it. But.
some trails you hike that are not well traveled and they're not well marked,.

$^{161}$they may be a very small path. It's not very wide. Just a little footpath through.
the woods. And if there's been leaves blowing around and branches that have.
fallen and things like that, it's easy to lose the path. Sometimes trails like that.
are marked on, they have little marks of paint on the tree or something else on a.
tree to let you know where the path is. But the point is you have to be.
intentionally pursuing that path to stay on it. You have to be looking at the path.
to know where it's going and you have to even be looking at the trees and.
other things around you for markers of where that path is. And if you're not paying.
attention, if you're not being intentional, you will get off the path..
And that's just as true in our journey of following Jesus. There are so many.
distractions from the narrow way. There are so many temptations that the devil.
puts in front of us to try to pull us off of the narrow way. There are so many.
off-ramps that the devil puts in front of us and says, "Hey, here's an easy way to.
just get off the narrow way and start traveling on the broad way." And the.
narrow way requires intentional pursuit. We have to be intentional about staying.
on it. Be paying attention. We have to be intentional about our relationship with.
God and keeping it healthy and our connection to Him and staying in His.
Word and praying and staying connected with other believers in the church and.
checking all of our beliefs against Scripture and against God's truth so.
that we don't start to believe erroneous things or things that are going to hurt.
us or hurt someone else. We've got to be intentional about all that stuff because.
if we're not intentional about staying on the path, then we're going to be off.
the path. It's just as simple as that. Another thing that's really important.
for us to understand, especially for us as younger people that live in our.
modern culture, is that the narrow way is deeply countercultural. The narrow way is.
deeply countercultural and that's reflected in these verses. Jesus says in.
verse 13, "The gate is wide and the way is easy that leads to destruction and those.
who enter by it are many." Those words, the word for wide means that it's broad,.
it's flat, it's easy to travel on, and the word for easy means that it's.
spacious, there's lots of room, it looks really nice. And there's a lot of.
people on this path. Verse 13 is very clear about that. That means the narrow.
way is not popular. It's not a path that is going to be crowded and it's probably.
not the path that many of your friends and acquaintances and co-workers and.
classmates and the people around you are going to be on. And it's probably not the.
path that your friends and classmates and co-workers and acquaintances and.
maybe even family members, the people around you, it's probably not the path.
that they're going to encourage you to go towards. It's probably not the path.
that they're going to support you in walking in. If anything, they're probably.
going to try to pull you away from the narrow way and encourage you to walk on.
the Broadway because it seems so much easier, right? They say, "Hey, this road is.

$^{201}$flat and easy to travel. It doesn't have the challenges. It doesn't oppose your.
selfish desires. On the Broadway, you do whatever you want. If you feel like.
doing something, you do it. If you feel like gratifying yourself in some way,.
then you engage in that activity. It's easy. And look, there's so many people on.
the Broadway and we're having a great time and this is where you want to be..
It's great over here. All the cool kids are here. So come here. Come to.
this way. Don't try to walk on that narrow, difficult path." That's the.
metaphor that Jesus is putting in front of us. And if we're honest, that's the.
world that we live in. Most of the people around us are not on the narrow way..
They're on the broad path and they are calling to us, sometimes implicitly, but.
sometimes explicitly in things that they say to us. They're calling to us to leave.
the narrow way and to come walk on the broad path. The narrow way is not.
supported by our culture and by the values of our culture. The values of the.
kingdom are upside down and backwards to the values of our culture. And so our.
culture does not encourage us to live out and walk in the values of the.
kingdom. That's just the reality. If you don't believe me, let's just take a.
second and go through some of the kingdom values that we've been talking.
about. Let's just look at the Beatitudes themselves. Think about these..
The very first one, "Blessed are the poor in spirit, for theirs is the kingdom of.
heaven." Think about that. If you say, "Okay, I'm ready to start a humility club at my.
school or at my job and I'm going to say, 'Hey, this club is all about poverty of.
spirit and humility. Let's sign up. Let's do this.'" How many people.
do you think are going to be excited about joining the humility club? No one..
That's not a value of our culture. If you want to start a club to promote.
yourself, then probably a lot of people would be interested in learning about.
that. But humility is not a value of our culture. What about purity of heart? Is.
that a value of our culture? Is that a value that your friends are going to.
encourage you to walk in and to live out? No, it's not. These things that Jesus.
talks about in the Beatitudes are really upside down and backwards to what we.
live out and practice in our culture. What about being meek and gentle? Is that.
a cultural value? It's not. What about hungering and thirsting after.
righteousness? Is that a cultural value? What about being a peacemaker? In our.
modern moment of turbulence and anger and frustration and division, everywhere.
you look, it is not a cultural value to be a peacemaker. And if you try to be a.
peacemaker, some people are going to be very upset at you and very angry at you..
Look at what's going on in our politics and in our society here in the.
United States. People make money. People fundraise. People get votes and get.
elected off of division and fear. And so if you're spreading peace, people don't.
want to see that. People in power want others to experience division and fear.
because that keeps them in power. And so if you're a peacemaker, that's not going.

$^{241}$to be a popular value. What about these values later on in Matthew chapter 5.
where Jesus talks about loving our enemies and turning the other cheek to.
someone who has wronged us? Are those popular values? Do you think you can get.
those off the ground with anybody in our culture? No. No one wants to hear that..
Certainly no one wants to practice it. But people don't even want to hear that.
expressed. People oftentimes don't even want to hear forgiveness expressed in.
our culture because we are a culture that wants to take revenge and seek.
retribution anytime that we've been wronged. No one wants to hear someone say,.
"Hey, Jesus calls us to forgive." No one wants to hear that value in our.
culture. What about in Matthew chapter 6 later on where Jesus talks about.
storing up treasures in heaven? Do you think that's a value in our culture? Do you.
think that's something that people are going to be excited about and encourage.
you to live out? Is making decisions on earth with your finances that have.
eternity in mind? Is anyone going to think that's a good idea? No. People are.
going to say, "Hey, this is for you. You use this to take care of you. You.
don't use this to store up treasures in heaven. That doesn't make any sense.".
What about in Matthew chapter 5 when Jesus talks about anger and lust and he.
says, "Hey, you should not even entertain these thoughts in your mind because.
that's sinning." What if you try to live those out? Think about that. Young people.
think about that. If you have a conversation with your friend and you.
explain to them, "Hey, I'm trying to even avoid sexual sin in my mind, to not even.
have thoughts of lust in my mind because I feel like that's sinful. And so because.
of that I want to avoid pornography or because of that I'm really careful about.
the situations I put myself in." That's something that people will make fun of.
you for. I've had those conversations. I've been made fun of that. If you say,.
"Hey, I want my sexuality to honor God and because of that I'm only going to have.
certain types of relationships or I'm going to wait until I'm married to.
engage in sexual activity." That's something that people will mock you for..
That is not a cultural value and we need to understand that. If we are serious.
about living out the values of the kingdom, the culture around us does not.
share those values and is not going to support them and it's not going to.
encourage them and a lot of times will oppose them and mock them. And that's a.
part of walking on the narrow way and we need to be ready for that. The narrow way.
often leads to persecution. Jesus told us that in Matthew 5 verses 11 and 12.
when he said that we were going to experience persecution. And here he's.
expanding on that and letting us know why. Because we're walking on a way that.
doesn't make sense to people who don't know God. We're walking on a way that is.
narrow and difficult and everyone else is on the broad and easy way. Another.
thing I want us to see here in these verses is that Jesus only describes two.
paths. He says there's the narrow way and the broad way and he doesn't describe.

$^{281}$anything else. And that's important for us to understand as well. There is no.
middle way. Sometimes we think, "Well I can kind of keep a foot in both worlds. I.
can kind of do both. I can kind of be involved in following Jesus and.
in living out kingdom values a little bit. But I can also be involved in.
living out the values of the world around me." But there is no middle way and.
Jesus has already talked to us about that in the Beatitudes when he told us.
no one can serve two masters. So that's not really an option even though we.
want it to be an option. It's not. Jesus doesn't give us that option. He only.
talks about two paths and he's clear that only one path leads to life. One.
path leads to destruction and there's no ifs, ands, or buts about it. Jesus makes it.
clear the broad way, even though it's easy and popular and seems great, it.
leads to destruction both in this life and in eternity. And the only way that.
leads to life is the narrow way. And that's what I want us to see as we close.
this out. That there is a reward to the narrow way. Now so far the things we've.
talked about can seem discouraging and burdensome like, "Wow this sounds so tough.
and I don't know if I want to be on the narrow way because it just sounds really.
grim and it's something that I really want to live out." But we need to.
understand that there is a reward to the narrow way and it's an amazing and.
indescribable reward. Now first of all there's a reward of everlasting life..
Jesus says that the narrow way leads to life and that word has multiple.
connotations here and so it means everlasting life. Scripture is clear.
about that. If you think about John 3 16, it says Jesus is the way to eternal life..
John 5 24, where Jesus says that eternal life comes through believing in.
him. There's many other verses that talk about Jesus being the way to eternal.
life. But that is a part of the reward. That when our earthly life is over and.
our body dies, our soul will not die but will go into the presence of God in.
heaven and live eternally in a perfectly beautiful place with no sin, no.
injustice, no disease, no tears. All the terrible things that we experience in.
this life will be gone and we'll have eternal life in the presence of God. That.
in and of itself would be reward enough even if life was full of suffering and.
challenges to achieve that. It would be reward enough. But that's not the only.
reward. Jesus is clear that there is life in this life that is a reward here. Like.
I said, this word that the narrow way leads to life has multiple.
connotations and it refers to our life in the here and now as well. In John 10.
10, Jesus talks about that. How he came so that we could have abundant life. That.
word abundant means have life to the full and we could experience what we.
were truly created to experience. And that is a part of the reward of the.
narrow way that we get to experience and live out if we follow God and live out.
kingdom values. We experience that in this life. Life to the full. The life that.
we were created for. A life of perfect and healthy fellowship with the God who.

$^{321}$made us and a life of healthy fellowship with the people around us. And that's.
something significant because what truly gives meaning to life? Is it the food we.
eat? Is it the money we make? Is it the things that we go do? I think all of us.
if we examine for very long we can see that what truly gives meaning to our.
life is the relationships that we have. That's why we enjoy the food if we're.
enjoying it with with people that are important to us. That's why we enjoy the.
experiences and the concerts and the events that we go to is if we're doing.
them with people that we love. That's what makes life meaningful is those.
relationships and the narrow way living out the values of the kingdom gives us.
the opportunity to experience truly healthy relationships. Like we talked.
about last week a truly healthy relationship with God the one who made.
us as we love him first and out of that we can start to live the golden rule and.
truly love the people around us and experience truly healthy relationships.
in our marriages, in our families, in our communities, in our cities, in our world..
So this promise of life is for the here and now as well and it also gives us.
something significant in the suffering and adversity that we face and that is a.
purpose. The reality is everyone experiences suffering and adversity in.
this world because our world is broken by sin. Even if you're not a follower of.
Jesus you're gonna go through some difficult times and if you are not a.
follower of Jesus if you're not living out the values of the kingdom then it's.
so difficult to find the meaning and the purpose in that suffering. It feels like.
there is no meaning and purpose in it. Why does this even happen? And that's.
why so many people can end up angry at God and say, "God why would you allow evil.
to be in our world? Why does this stuff happen? Why do you let any of these bad.
things take place?" But for a follower of Jesus, for someone who's living out the.
values of the kingdom, we can see that God is using that adversity and that.
suffering in our lives to grow us and to stretch us and to make us more like.
Jesus and to make us more able to love God and to love others. And so that gives.
purpose even in the challenging difficult moments of our life and that's.
a huge part of the blessing of walking in the narrow way. It's not that the.
narrow way is easy, it's challenging, but we see the meaning in the challenges. We.
see the purpose in the challenges that God brings into our life and it gives us.
that sense of shalom that we talked about last week. It's so powerful that.
God desires for us to have that sense of wholeness and that sense of.
health and prosperity and peace and justice and of permanence in our.
life because that only comes through right relationship with God and through.
seeking to live out the values of the kingdom. And I want us to see this, I.
don't want us to miss this because this is really big. We are not alone on the.
narrow way. Verses 13 and 14 tell us that there are many people on the broad.
way and sometimes when we read that and then read the words that follow we just.

$^{361}$assume that well I guess the narrow way means you're by yourself. It doesn't say.
that. It says those who find it are few. It doesn't say that those who find it.
are alone. And so even though there are times when the narrow way can feel.
lonely because it is countercultural and we're not encouraged to.
walk on that way by many of the people around us, we are not alone. First of all.
we have that relationship with God that strengthens us and sustains us and keeps.
us going. But secondly God through his body, through the church, through the.
the family of God as the New Testament calls us, puts all these people around us.
who love us and care for us and support us that we can fellowship with, that we.
can be encouraged by, that we can be cared for by, and that we can in turn.
care for them and encourage them and serve them and help them. We are not.
alone on the narrow way and that's encouraging to us in the times of.
challenge and it also is a part of that that abundant life. Life to the full is.
experiencing those truly healthy relationships that God gives us in in.
the family of God and the church. And then we have the opportunity to truly.
love the world around us. We talked about that last week when we were discussing.
the golden rule that it's so hard for us as humans to love and care for anyone.
else in a truly unselfish way. We always have some type of selfish motivation at.
the heart of what we do but through the work of Jesus in our heart and the power.
of the Holy Spirit in our lives we can actually be a conduit of God's love to.
others and we really can love the people around us, the people who don't know.
Jesus. We really can love and serve them in a truly unselfish and loving way and.
have life-giving relationships with people who are outside the church and.
that's a part of that life to the full. So there is great reward to walking in.
the narrow way. It is totally worth it and God calls us to it. We need to.
understand that it's not easy. We need to understand that there's commitment.
required. We need to understand that there's intentionality required that we.
have to pay attention and stay focused to stay on that path. We need to.
understand that it's countercultural and that we are not going to be encouraged.
in the narrow way by most of the people around us but we need to understand that.
it is full of reward and beauty and goodness because God loves us and he is.
full of grace and he's extending that grace to us even as he calls us to walk.
on the narrow way. Let's pray. God thank you for this truth and just for how.
clearly and directly you lay it out in this passage that you call us to walk on.
the narrow way. It's not an easy way. It requires commitment. It requires.
intentional pursuit. We have to follow after you and God it is full of.
challenges. There's a lot of difficulties in life but there is so much reward and.
so God I pray that we would not be surprised by the challenges or the.
adversity of our life. We would see what you're saying here and understand that.
that it's going to be challenging at times and we would understand the.

$^{401}$commitment and the intentionality that's required and we would.
dedicate ourselves to having that commitment and living that out.
intentionally and God as we do I pray that we would experience your reward in.
our lives. That we would really begin to be so saturated with the values of the.
kingdom that they're flowing out of our life in a life-giving way to others. That.
we would live out these values of the kingdom that when it's difficult, when.
we're made fun of, or when people don't understand or oppose us that it would.
not stop us from living out the values of the kingdom and God that we would.
experience your life to the full. We know we'll experience that one day in heaven but.
even now in this life that we would experience it to the full. That we would.
live in healthy relationship to you. That you would just pour your love and your.
life into us and that we would be a conduit of your love and your life to.
others. That we would be able to to have healthy and life-giving relationships.
with those around us. God in Jesus name, Amen..
[Music].
(music fades).
[BLANK AUDIO].
\newpage



\section{}
\label{sec:_taiMeC2wK0}
\textbf{舊約網上行 - 上帝的足跡: EP10 等待之約}
\newline
\newline
連結: \href{https://youtube.com/watch?v=_taiMeC2wK0}{\texttt{https://youtube.com/watch?v=\_taiMeC2wK0}} ~~~~ 語音日期: 2021-05-30
\newline
\newline
\hyperref[sec:afSlbD_al4Y]{\small{< < < PREV SERMON < < <}}
~
\hyperref[sec:index_chronic]{\small{[返順時目]}}
~
\hyperref[sec:index_scriptual]{\small{[返順卷目]}}
~
\hyperref[sec:17rNRRMvNT8]{\small{> > > NEXT SERMON > > >}}
\newline
\newline
$^{1}$我們來到《舊約》網上的第四集.
上一次我們提到《創世記》第十四章.
在這章的聖經裡面提到亞伯拉罕出兵打敗近東五王.
擁救羅德然後海船回歸.
就證明了亞伯拉罕來到迦南地的時候.
他的勢力其實都是很雄厚.
並且是一個軍事的領袖.
今天來到《創世記》第十五章.
我們留意十四十五章兩章的經文其實有一些共通的字詞.
這裡告訴我們兩章的聖經是連接著的.
首先最明顯的就是十五章第一節.
就是「這時以後」打完仗之後發生什麼事呢?.
神對亞伯拉罕說十五章第一節.
他說「我是你的盾牌」.
這個「盾牌」這個字在聖經裡面經常出現.
就是耶和華對軍王作的保護是他們的盾牌.
特別是跟打仗有關.
但是十四章十九節.
「盾牌」這個字在希伯來文也出現過.
當時打完仗回來.
亞伯拉罕遇見默基西德王.
在耶路撒冷默基西德王向亞伯蘭祝福.
他說願創造天地的主至高的神.
賜福給亞伯蘭.
將受敵交在.
在希伯來文同一個字「盾牌」.
你的手裡面至高的神是應當稱頌的.
第三方面我們留意到.
神對亞伯拉罕說.
我必大大的賞賜你.
「賞賜」這個字通常是用在軍兵的報酬上.
其他的舊約聖經也有這樣用.
這裡是因為亞伯拉罕海船回歸.
是一個軍事領袖.
所以神必定會大大的賞賜他.
第四項就是318這個數字.
我們記得在十四章十四節的時候.
亞伯拉罕出兵去救他擇兒羅德的時候.
就抽調了他家女生養精蓮壯丁.
318人就一直追趕到達.

$^{41}$318這個字剛好是15章這裡.
「以利以謝」這個名字.
希伯來字母的總數.
「以利以謝」就是亞伯拉罕其中一個的僕人.
所以你看到「以利以謝」.
這個字每一個希伯來字母加在一起的時候.
就是總數318.
這些都不是巧合的事.
就是告訴我們十四章到十五章是緊接著的.
十五章想交代什麼呢?.
十四章就說亞伯拉罕是一個英雄人物.
他是軍事的領袖.
十五章就想交代兒子的認許.
和土地的認許.
當時近東的文化要建立一個邦國.
就需要一定有後代.
也需要有國土.
不可以單單有軍事的勢力.
神認許亞伯拉罕成為大國.
藉著這個大國去使萬族得福.
如何能夠應驗出來呢?.
就是十五章要交代的情節.
亞伯拉罕的故事.
由十一章開始到十二章.
一直發展到現在.
都是有一個危機.
這個危機一直以來都未能解決.
首先神對亞伯蘭最初的承諾.
就是呼召他離開本地本族本家.
但是他去的時候.
他是沒有遺旨的.
一方面是殺來不辱.
當然羅德是跟著亞伯拉罕一齊來到迦南地.
但是上幾章我們說過.
羅德是自己來DQ自己的.
所以到這個地方.
仍然沒有一個繼承人.
也沒有土地.
如何能夠成為大國呢?.
上帝會通過什麼方法.

$^{81}$使亞伯拉罕成為大國.
使萬國得福.
換句話說這裡問一個問題.
就是神的話是否可信?.
第十五章我們先讀透半段.
你會發現在這張聖經裡面.
特別的地方就是.
經常有對話.
之前沒有這麼多.
但第十五章十多次提到說話.
話語的字眼.
我們可以將頭一段這樣分段.
首先耶和華說話.
對亞伯拉罕有一個賞賜的認許.
然後就是亞伯蘭的回答.
他有說話.
他對他沒有遺旨.
擔心後裔會來的訴求.
然後耶和華又回答.
他在說話裡告訴亞伯拉罕.
這個後裔會為他預備.
是由他所生出來的.
之後上帝帶亞伯拉罕去看一個景象.
這個景象可以叫做眾星的意象.
最後就稱亞伯蘭為義.
這裡就是一個立約.
我們先看經文.
十五章第一節.
這時而後.
耶和華在意象中臨到亞伯蘭說.
亞伯蘭你不要懼怕.
我是你的盾牌.
你的賞賜是很大的.
亞伯蘭說主耶和華.
我一向都沒有孩子.
你還能賜給我什麼呢?.
這樣承受我家業的.
就是大馬士革的義理者了.
亞伯拉罕又說.
你既然沒有給我後裔.

$^{121}$那生在我家中的人.
就是我的繼承人了.
按照當時近東的社會.
兒子當然是相當重要.
有兒子的話.
就代表兒子有社會家庭的責任.
來照顧父母.
顧父母.
並且死後.
有兒子.
你也能買下.
並且紀念到父家的名字.
幾乎是我們現代人說的.
一些人壽保險.
或者醫療保險.
另一方面.
有後裔就是去繼承家族的財產.
繼承家族的名字.
特別對於一些君王皇室.
或者一些大的家族來說.
兒子就相當重要.
十五章第二節.
亞伯拉罕向上帝申訴.
他說我既然沒有後裔.
沒有孩子的時候.
有什麼賞賜都沒有用.
按照當時近東處理家業.
繼承權的法律文件.
一個人如果沒有自己的兒子的時候.
他可以在法律上.
收養另一個成人.
作為他的兒子.
這個人就成為他家業的繼承人.
在這裡.
亞伯拉罕就向上帝提出.
說這個大馬士革人.
以利爾謝.
作為他的繼承人.
你記得在吾爾到哈蘭.
來到迦南的旅途當中.

$^{161}$未進入迦南地.
應該會經過大馬士革.
或者大馬色這個地方.
相信可能在那裡.
他收養了這個僕人.
亞伯拉罕就說.
不如就立這個以利爾謝.
做我的繼承人.
有什麼反應呢.
第四節.
耶和華的話又臨到亞伯拉罕.
說這人必不會做你的繼承人.
你親生的才會成為你的繼承人.
神的意思是.
不要用你自己想的方法.
來成就我在你身上的工作.
我有我的安排.
於是耶和華就給亞伯拉罕一局警長.
第五節.
於是亞伯拉罕出到外面.
數點眾星.
你能不能把他們數得清楚.
又對他說你的後裔將要這樣眾多.
之前上帝已經用類似的說話.
來確保兒子的認許.
十三章第十六節.
如果你記得的話.
神曾經跟亞伯拉罕說過.
他說我也要使你的後裔如同地上的沙塵.
塵沙那樣多.
人若能夠數算地上的塵沙.
我才能數算你的後裔.
今次的意境就有少少不同.
就是叫他看天上的眾星.
問他能不能夠數得準.
在中世紀有一個很有名的猶太拉比叫Rashi.
他解釋這段經文.
亞伯拉罕因為自己沒有後裔.
他等得不耐煩.
他想用自己的方法.

$^{201}$或者用以前他自己曾經用過的方法.
來解決這個問題.
以前的方法就是占星去計算未來的方法.
因為亞伯拉罕在加勒底人的吾爾.
即是巴比倫人的地方.
那裡最有名的是什麼呢.
就是占星術.
就是用星塵的軌跡來預知未來.
我們在巴比倫等地方.
找到很多例子是跟占星學有關.
Rashi就說上帝在責備亞伯拉罕.
有如對他說請你放棄你自己的占星方法.
因為上帝要用他自己的方法.
去實現他對你說的話.
可能會是這個意思.
但我相信這個故事.
也不是說在黑夜裡看星星這麼簡單.
或者是數星塵這麼簡單.
你看看故事在十五章的時間敘述.
他說他在外面向天觀寒.
數算眾星.
到底是何時發生的.
我們想起數算眾星滿天星的景象.
自然我們想起黑夜.
特別是在一個中東的地方.
曠野的地方.
好像在一個沙漠裡.
一望上天就一定有很多很多的星塵.
但如果你看完十五章.
特別是看下文.
單單按照故事敘述的時間.
你就可能再深一層想一想.
因為看下文的時候.
神在獻祭的時候.
就和亞伯拉罕來立約.
是什麼時候發生呢.
十五章第十二節.
在日落的時候.
伯蘭尼沉睡.
然後有很大的黑暗落在他身上.

$^{241}$然後下文就去交代獻祭的事.
所以下面獻祭的時候就是日落的時候.
換句話說.
至少在故事敘述的鋪排上.
是日間的時候.
即是說上帝在白晝的時候.
可能是日當頭的時候.
就帶亞伯蘭出去外面.
叫他看看天上的星.
看看他能不能夠數到.
這些星星就是你的後裔.
有多少.
換句話說.
這裡就很幽默.
叫他觀星.
是日當頭的時候.
日間白晝的時候來觀星.
根本看上去就不會看到有任何星星.
怎樣數呢.
換句話說.
這個異象就是要去挑戰亞伯拉罕.
他要憑眼見還是憑信心.
去見到上帝說話承諾的兌現.
不是能夠憑眼見的.
因為你看見天上的景象.
你是不會見到任何星星的.
你只能夠憑信心去相信.
上帝說了的話一定會兌現.
上帝說了的應許一定會實行出來.
未必一定會按照我的方法.
但一定會按照上帝的時間表.
一定會實現出來.
要挑戰亞伯拉罕的信心.
十五章第六節.
他就回應.
亞伯蘭信耶和華.
耶和華就以此為他的意料.
很多時候我們看這節聖經的時候.
我們很快就會跳到新約.
就是說一定是因信稱義.

$^{281}$因為亞伯蘭信耶和華.
所以他就成為了一個義人.
其實不需要這麼快.
義這個字righteousness.
什麼意思?.
正確.
正當的回應.
正當的反應.
回應什麼?.
就是回應神的應許.
回應神的承諾.
他的話語.
有一個正確的回應.
這就是亞伯蘭的信心.
在亞伯蘭的故事當中.
信神在這一處.
意思就是將你的將來交給上帝.
神說我必定賜給你後裔.
也會賜給你迦南地這個土地.
並且會使你成為萬國得福.
能夠見證到我的.
他就是這樣相信.
上帝將來必定會實現他的說話.
我們再看十五章的下一段.
就是七至二十一節.
你會看到這個列表告訴我們.
其實脈絡都是差不多的.
就是有對話.
然後有一個景象出現.
最後是立約.
不過十五章下一半所說的.
不是兒子的應許.
而是土地的應許.
我們看看經文如何發展.
第七節.
耶和華有對亞伯蘭說.
耶和華曾經將你從迦勒底人的吾爾領出來.
為要把這地賜給你作產業.
跟上一段一樣.
第一段說到兒子的應許.

$^{321}$亞伯蘭就申訴.
聽到之後他就說.
我怎麼知道我會有後裔呢.
神就用說話來回應他.
然後又展現出一個天上眾星的異象.
現在第二段就說到土地的產業.
這個應許.
第八節.
亞伯蘭就申訴.
他說我聽到你說將會賜這塊地給我.
但是主耶和華.
我憑什麼能夠知道你真的將這塊地賜給我.
神又用說話來回應.
第九節.
你給我拿一頭三歲的母牛.
一隻三歲的母山羊.
一隻三歲的公綿羊.
一隻班溝和一隻初鴿.
亞伯蘭就照做.
將這一切拿來.
每樣都當中劈開一半一半的相對擺列.
只是鳥沒有劈開.
比較小的祭物就不劈開.
有資料下來.
落在這屍骸上.
亞伯蘭就將他們嚇走.
一落的時候.
亞伯蘭沉沉的睡著.
忽然有可怕的大黑暗落在他身上.
耶和華就對亞伯蘭說.
你要確實地知道你的後裔.
必在外地居住.
也必服侍那地的人.
那地的人輔待他們四百年.
他們所要服侍的那國.
我要親自懲罰.
後來他們必帶著許多財物從那裡出來.
你必在長壽被人埋葬.
平平安安回到你列祖那裡.
到了第四代.

$^{361}$他們你的子孫必定會回到這個地方.
因為阿摩利人的罪孽還沒有滿盈.
簡單說一些重點.
這裡有許多的祭神.
是一個很隆重的祭禮.
用了全日的時間.
去等待.
在這種景象.
你看到一些祭物.
一半一半擺開的時候.
其實是近東的立約情景.
兩個王來立約.
很多時候會將這些祭神.
一半一半放開.
然後在中間經過.
然後就起誓.
我們有一些例子.
譬如一個古城在今天的敘利亞.
叫做阿拉拉克.
是公元前大概1800年的時間.
阿巴拉罕就是這個時代.
找到很多的泥板.
其中一些泥板.
有這樣的字句.
去形容兩個王立約的時候.
看到什麼呢?.
大王就說.
我將阿拉拉克這個城市.
賜給小王.
將這塊地賜給他.
然後就向他起誓.
又切割羊頸.
有祭神出現.
也有起誓的說話.
如果我收回這塊地.
我是會走的.
但是將來這個小的王.
叫做耶利念.
反叛的話.
他也會放棄他的土地.

$^{401}$所以他要去守約.
亞述一些泥板.
也告訴我們.
同樣一些的祭禮.
就是一個亞述王.
和一個小城邦的王.
來到立約.
這個小城邦的王.
叫做馬提爾路.
立約裡面的一些條文這麼說.
在祭禮上指著.
祭禮的頭.
他說這個頭不是屬於羊的頭.
是屬於小王馬提爾路的頭.
是他兒子的頭.
代表著是他土地人民的頭.
如果他們背約的話.
他的命運就好像.
這隻羊一樣.
他的頭被切割出來.
耶利米蘇三十四章.
也同樣有一個這樣的描述.
耶和華指控背約的以色列民眾.
他說他們沒有去到里行.
在我面前所立的誓言.
立約的誓言.
就是他們將牛獨劈開一兩半.
並且從兩半中間走過去的時候.
立下的誓約.
這裡說到.
祭物分開兩半.
立約的雙方在中間走過.
然後起誓.
以性命自己來擔保.
不會背約.
所以十五章十八節.
耶和華與阿不蘭立約.
原文其實是.
Cut a covenant.
切割立約.

$^{441}$就是斬開這些祭神.
來起誓立約.
看回這個故事.
到底發生什麼事.
十五章第十七節.
日落天黑的時候.
忽然有無煙的路.
燒著的火把.
從祭禮的玉塊中經過.
就在那時候.
耶和華與阿不蘭立約.
跟祂說.
這塊地會賜給你的後裔.
就是埃及河直到丘法拉底河之間.
這塊迦南地.
火把煙.
經常是上帝出現的景象.
在祭物中間經過.
代表上帝.
用自己的生命來保證.
保證祂的說話.
祂的應許.
一定會成就.
即是對阿不蘭立約說.
我是掌管你的將來.
一定會在你身上.
成就我的工作.
我是掌管神的子民.
每一個的將來.
即使你現在看不見.
你現在摸不清.
在亞伯拉罕的故事中.
一直都是纏繞著信心的問題.
信心是什麼?.
有人用信心來下一個定義.
可以說我拍胸口就是信心.
或者我只熟悉一些基督教的術語.
就叫信耶穌.
或者有些人覺得.
信仰是反理性的.

$^{481}$不講理由.
在亞伯拉罕的故事中.
信是什麼?.
信的元素就是一個效忠.
一個偉神.
等候或者相信.
上帝的應許一定會實現.
在等候當中.
去認定神的說話.
神的應許.
必定會兌現.
即使現在眼睛不能夠見到.
再看大一點的圖畫.
可以想想神對亞伯拉罕說過什麼.
就是給他後裔.
給他土地.
但記得第十二章的時候.
上帝真正的應許是什麼?.
要使亞伯拉罕成為大國.
所以要有兒子的應許.
要使他成為大國.
所以要有土地的應許.
但最終要做什麼?.
就是要藉著亞伯拉罕.
使他成為祝福.
使他能夠祝福萬國.
在他身上能夠見證到.
上帝自己的榮耀.
這才是上帝對亞伯拉罕的應許.
再看遠一點.
神的呼召.
叫亞伯拉罕成為萬國的祝福.
有應許之地.
有君王會出現.
最後會有救主出現.
其實這些應許.
你看看整個故事.
都沒有在亞伯拉罕的生命裡面兌現.
甚至整個救約到最後真正成就的時候.
是在耶穌基督的十字架上.

$^{521}$但在亞伯拉罕身上.
要讓他的信心.
雖然要等候.
雖然未必現在能夠見到他兌現.
但他相信的.
就是上帝最終一定會成就他每一句說話.
每一句說話都不會落空.
這就是亞伯拉罕的信心.
一個等候的信心.
一個願意堅定盼望的信心.
疫情影響全球.
在疫情裡.
等可能成為一個常態.
我們每個人都在等.
等著什麼時候疫情會過.
等著什麼時候某些事情會重開.
隔離又要等.
等多久.
我們在等候裡.
實在不容易過.
但亞伯拉罕的信心說.
認定的是神必定會帶領著未來.
遠勝過我們自己能夠把握障礙.
現在人的眼睛看不清楚.
但我們深信.
神會完成自己的工作.
會繼續去見證自己.
現在能否見到?.
可能能見到的東西.
剛剛相反.
有新聞報導.
甚至說不少教會.
因為過去一年多的疫情解散.
我認識有一兩間在附近的教會.
因為捱不過疫情.
人開始散.
以至教會也解散.
有很多工作受到很多的攔阻.
在等候裡.
有很多未知之數.

$^{561}$有時真的可以是一種黑色的夜晚.
一些很難過的黑夜.
我不知道你現在面對的黑夜是什麼.
面對人生一些挑戰的關頭.
或者你在等待一些東西.
在未知之數裡掙扎.
可能面對一些很難決定的關頭.
摸不清將來.
可能在等待一些結果.
無論是工作,去留,醫療報告等等.
可能是一些你很關心的人.
你等候他迴轉.
希望神在他身上工作.
或者可能你事奉了一段時間.
見不到那個效果.
在等待當中.
你覺得好像很浪費.
但是亞伯拉罕的故事說.
對人來說.
等候是浪費.
但對上帝來說.
等候是他的工作.
這樣看的時候.
未必能夠縮短我們等候的時間.
也未必能夠減輕等候的苦楚.
但起碼可以給我們一個大的角度.
一個屬神的角度.
去面對現在心靈裡的黑夜.
猶太人拉比去讀亞伯拉罕的故事時.
說神在整個故事裡十二至二十二章.
給亞伯拉罕十個不同的試煉.
這些試煉往往是試煉他信不信神的應許.
但一定要等候.
在等候當中.
神往往在他身上繼續工作.
去鍛鍊他對上帝的信求.
最後亞伯拉罕才能夠成為.
以色列民的信心之父.
成為祝福.
以至他的後代.

$^{601}$他所代表的.
能夠成為萬國的祝福.
能夠在萬國當中見證神.
等候.
信心.
這些是最基本的功課.
上帝要我們每一個都經歷.
宣道會有一個已故的牧師叫做托薩.
他曾經說過一段這樣的說話.
他說神從來都不趕時間.
Not in a hurry.
因為神沒有一個期限.
迫使他一定要今天工作.
但我們如果深知這個道理的時候.
我們的心靈可以得到安撫.
以至我們不會慌亂.
在等候當中.
我們焉有盼望.
希伯來書十一章第一節.
我們用這一章的聖經來結束.
這裡說一下信.
就是對所盼望的事的把握.
是看不見之時的明證.
然後十一章列出一個又一個信心的偉人.
信心的榜樣.
請你不要搞錯.
這些不是完美的人.
他們都是在生活中的掙扎.
在聽到上帝的話.
知道上帝一定會藉著他們去見證神.
在這些掙扎軟弱的當中.
不是完美的人.
但有一個共通點.
每一個都是繼續相信.
神必定會承受他的諾言.
十一章第八節.
就說到亞伯拉罕.
因著信亞伯拉罕在夢召的時候.
聽命往他將要承受遺孽的地方去.
他去的時候還不知道.

$^{641}$往那裡去.
一步一步的憑信心走.
因著信亞伯拉罕在應許之地寄居.
好像異鄉.
承受同樣應許的以撒雅各一樣.
是住在帳棚裡面的.
因為他們等候的那座有根基的城.
就是神所設計所建造的.
一直繼續數不同的信心偉人.
十三節.
他說這些人是存著信心死的.
在他心靈裡面還沒有得著所應許.
只不過是從遠處用信心見到.
就表示歡迎.
承認自己在世上只不過是客許.
這樣說話的人.
是表明他們尋求另一個的家鄉.
一個更美將來天上的家鄉.
所以神就說.
我不以他們為羞恥.
因為神已經為他們預備將來的一座城.
在未知的裡面.
在黑夜的當中.
我們相信神必定會按照他的說話.
按照他的計劃來行事.
可能我們現在未必見到上帝的國.
但我們很肯定.
上帝會按照他的說話.
來成就他的工作.
去榮耀他自己的名.
求主幫助我們.
也去用盡我們.
讓我們現在的處境.
無論你覺得很倒楣.
或者很不順利.
我們向著一個方向.
祈求上帝仍然可以使用我們.
使我們成為祝福.
使我們在身邊的人當中.
見證神自己.

$^{681}$我們同心 佳德 祈禱.
上主我們求你.
在我們信心被試煉的時候.
特別是我們心靈裡面.
可能過著不同的黑夜.
遇見不同的挑戰.
有很多未知之數的當中.
主讓我們認定.
你自己的說話必定會兌現.
你自己的認許必定會成就.
主你要我們像阿爸一樣.
能夠被你所用.
成為祝福.
被你所訴訟.
能夠成為一個信心的人.
我們的信心.
就是你自己的話語.
必定會兌現.
求你親自來幫助我們.
在我們信心的旅程當中.
在我們現在這個角色裡面.
無論是白晝.
還是我們正在經歷一些黑夜.
主你去用我們.
使我們能夠見證你.
我們感激到方鑽名求.
\newpage



\section{}
\label{sec:17rNRRMvNT8}
\textbf{舊約網上行 - 上帝的足跡: EP11 聖約的神}
\newline
\newline
連結: \href{https://youtube.com/watch?v=17rNRRMvNT8}{\texttt{https://youtube.com/watch?v=17rNRRMvNT8}} ~~~~ 語音日期: 2021-06-21
\newline
\newline
\hyperref[sec:_taiMeC2wK0]{\small{< < < PREV SERMON < < <}}
~
\hyperref[sec:index_chronic]{\small{[返順時目]}}
~
\hyperref[sec:index_scriptual]{\small{[返順卷目]}}
~
\hyperref[sec:hPfz5WlNiPQ]{\small{> > > NEXT SERMON > > >}}
\newline
\newline
$^{1}$今日夠約網上行,我們繼續看阿拉罕的故事.
我們看創世記第十六至十七章.
曾經我提過阿拉罕這個故事可以用一個三文字的結構來了解.
所以我們今天看十六十七章正正在這個三文字結構的中間.
十六章第一節.
阿布蘭的妻子撒萊沒有為他生孩子.
撒萊卻有一個婢女,是埃及人名叫夏格.
撒萊對阿布蘭說.
請看,若我使我不能生育,求你親近我的婢女.
或許我可以從她得孩子.
阿布蘭就聽從撒萊的話.
阿布蘭的妻子撒萊是不能生孩子.
你記得在阿布蘭的故事開頭已經曾經提過這個事實.
十一章,撒萊是不育的.
當時的社會太太不育當然是一個咒詛,一個羞恥.
因為在一個農耕的社會當然需要人手.
甚至在古代近東社會裡面.
一個父母希望有兒子.
是希望當他年老的時候.
能夠跟他傳宗接代,辦理身後事.
最主要是紀念和延續父家的名.
所以在近東社會裡面,一個妻子最主要的責任就是傳宗接代.
生孩子,生得越多越好.
有一點像我們古代傳統的華人文化.
當時的社會如何解決沒有後代的問題.
男人當然很簡單,很多時候直接去休妻離婚.
另外娶更多的女人.
這些都是當時的風俗習慣.
從我們看阿布拉罕的故事開始.
我們已經知道阿布拉罕的家族是一個不小的家族.
人手和財力都是非常豐富的.
所以他也可以按照當時的風俗去做.
男人說完了,女人又如何呢?.
女人按照當時的風俗.
一個很普遍的做法就是為丈夫找一個妾侍.
而這個妾侍所生的孩子就會歸給這個證實.
法律上他就成為這個家族的繼承人.
在今天的伊拉克有一個城鎮叫做Kirk.
在西南面大概十三公里左右.
在Tigris River的上游,底格拉斯河的上游.

$^{41}$有一個古城叫做Nuzi,老支.
在1920年代曾經發掘過有七千多塊的.
蝕形文字的泥板.
這些泥板屬於的時代大概是公元前1500年左右.
這個就是族長時期.
阿布拉罕以撒雅各這段時間.
這些泥板有很多種類.
有些法律的條文,或者商務的記錄.
讓我可以在視窗看到當時的風俗習.
其中一塊泥板講到什麼呢?.
他說到一個妻子叫做Kinimnu.
她生了孩子的話,這個夫君Shinimma就不可以另娶妻子.
因為她生不了孩子.
但如果她生不了孩子的話.
這個妻子就要為夫君找一個僕婢作她的妾.
而這個僕婢的後代是不可以送走的.
因為是歸她名下.
所以任何妻子為了夫君所生的孩子.
甚至是為了代人母的妾事所生的.
都可以承受土地,房屋,產業等等.
另外一個例子就是舊亞述帝國.
大概公元前1800年,更加近阿布拉罕這個時間.
泥板記錄了什麼呢?.
就是一個丈夫娶了一個妻子.
假如兩年之內這個妻子未能為他生子.
他可以買一個僕婢為他生子.
這些內容跟創世記族長的家庭故事,不育的故事非常相近.
所以在這個時候薩奈.
或者只不過是動用他當時婚約的條款.
一些風俗習慣來解決母後代的問題.
我們聽聽薩奈的話說什麼.
十六章第二節.
薩奈對她丈夫阿布蘭說.
請看,耶和華使我不能生育.
求你親近我的婢女,或者我要從她得孩子.
但在薩奈的心目中,誰錯了?.
他認為是耶和華,使他不能生育.
之前是阿布蘭的信心受考驗.
到這一章十六章,就是薩奈的信心受到考驗.
可能年老的薩奈以為上帝允許他能夠生子.

$^{81}$到這個時候已經等不及了.
他可能以為很快,但這個時候有一點心急.
十六章第三節,這個故事繼續發展下去.
阿布蘭在迦南地悶了十年.
他的妻子薩奈把自己的婢女埃及人夏甲.
給了她的丈夫阿布蘭為妾.
這個夏甲可能是之前十二章.
基方走難到埃及的時候.
上帝由恩典救他們出來.
出了埃及之後,帶著很多財產.
也有很多的僕人,僕婢.
這個夏甲可能是其中一個.
第四節,阿布蘭與夏甲親近.
夏甲就是懷了人.
夏甲見自己有了新人,就興翰他的主母.
薩奈對阿布蘭投訴.
他說我因你受屈,把我的婢女送到你懷中.
他見到自己有新人,就興翰我.
冤孽在你我中間主持公道.
阿布蘭就跟他說,這個婢女在你手下.
你想怎樣待她就怎樣待她.
於是就有家暴出現.
這位夏甲婢女很不幸地逃走.
我們剛才看的例版.
就是當時處理這些問題的法律條文.
但聖經就講得深入很多.
不單是講表面上的風俗習慣.
而是講到人性的醜陋.
夏甲見到自己有新人,就去看低主母.
主母又去報復.
如果你從來沒有看過聖經或仔細看過的話.
你應該會發現,你細讀的時候.
就是那種人性真實的一面.
很多時候就是在這些故事裡面交代出來.
是沒有分析過的.
創世記每一個家庭都是不健全,都是破碎.
因為我們每一個都是有血有肉的人.
如果你的家庭是不健全.
或者你這樣才會更加了解聖經的訊息.
聖經的訊息告訴我們.

$^{121}$不是因為我們完美,所以神的恩典才會到.
正是因為我們破碎,神的恩典又臨到.
這就是上帝能力所在.
故事繼續發展下去.
夏甲這個時候應該是粗生大細,行動很不便.
但他要逃亡,他本身是埃及人.
所以他應該在南地這個地方.
一路沿路回到埃及,這條路叫做樞義之路.
經過西乃半島北面的一個礦業叫樞義.
一直到埃及一個很出名的溪谷叫做瓦迪圖馬納.
他就會回到埃及.
創世記十六章第七節.
發生什麼事?.
耶和華的使者在礦業的水泉旁邊.
就是樞義路上的水旁邊,遇見夏甲.
他就問他,到底你在哪裡來?.
你要去哪裡?.
夏甲回答,我是從我的主母撒萊面前逃跑出來的.
耶和華的使者就對他說.
你回去,回去主母的那裡,伏在她的手下.
可能他的意思是,回去就不要再輕看你的主母.
儘管在她的權下繼續生活.
神也繼續應許和祝福夏甲和他生的孩子.
十一節,耶和華的使者對夏甲說什麼?.
他說,漢,你已經壞了人,你要生一個兒子.
你要為他起名叫做以實瑪利,Ishmael.
以實瑪利的意思是神聽見的意思.
所以他解釋,耶和華聽見你的苦情.
這個兒子將來會怎樣?.
將來他為人必像野狼,他的手要攻打人人的手.
也要攻打他,他必住在眾弟兄的東邊.
這些描述都是說近東的遊牧民族.
說到以實瑪利將來會成為遊牧民族的支派.
因此這個井叫做比爾拉海內.
比爾是什麼意思?希伯來文是井的意思.
拉海就是活著者的意思.
內大概是看見,看顧的意思.
所以這個井是紀念他和神的相遇.
知道上帝是看顧他的.
這個地方就在迦底斯和巴列的中間.

$^{161}$之後夏甲當然回去了.
十五節,如何交代?.
夏甲就給亞伯蘭生了一個孩子.
亞伯蘭就給夏甲所生的兒子起名叫做以實瑪利.
這個起名的動作就代表亞伯蘭在這個時候.
就接納以實瑪利作為繼承人.
當時神仍然未啟示.
未說要從撒萊所生的以實才是真名天子.
十六節,亞伯蘭八十六歲的時候.
夏甲就為他生了以實瑪利.
我們要問一個問題.
為什麼要記載一個這樣的故事?.
為什麼要記載以實瑪利這個故事?.
首先我們要搞清楚.
我們作為現代人不要亂用以實瑪利這個故事.
以為以實瑪利就是現在的巴勒斯坦人.
或者阿拉伯人.
他的後裔就是現在我們看到的那些住在中東的人.
為什麼呢?因為這個是不合歷史的.
到了大概公元前五世紀四世紀的時候.
巴比倫帝國,波斯帝國.
佔領了日南河東西半島等等的礦業地方的時候.
以實瑪利包括以東,摩亞,阿們等等.
這些世紀民族.
包括這些遊牧民族.
其實已經被消失了或者被同化了.
在舊約裡面.
以實瑪利從來都不是指今天的阿拉伯.
阿拉伯這個名字.
這個標籤.
泛指散居在一個很大的領域.
西拉半島,阿拉伯半島.
很多遊牧民族都是泛指阿拉伯.
這個不是一個單一的一個種族的標籤.
因為到了主後七世紀伊斯蘭帝國興起的時候.
就夾雜了很多不同的地方.
很多不同的語言,宗教,文化.
才醞釀了阿拉伯的身份.
所以現代伊斯蘭教的一些說法.
說穆罕默德是以實瑪利的後人.

$^{201}$這是不合乎歷史的.
所以這個以實瑪利的故事.
就跟今天中東發生什麼事.
沒有什麼直接的歷史的關係.
我們要記得看聖經.
聖經不是直接寫給我們.
聖經要明白的時候.
首先我們要問.
第一代的讀者聽見這個故事的時候.
他們會有什麼反應.
創世記成書的時候.
最早大概是聯合王國的時候.
當時是大衛,所羅門已經在以色列國建立了國家.
是太平盛世的時候.
四周就有很多不同的民族,不同的國家.
創世記去講這些列祖的故事的時候.
其實很多時候就是想回答一個問題.
他就說看這些鄰國,民族.
這些附近周圍的國家.
其實都是我們的遠房親戚.
譬如我們看過羅德.
羅德的後裔我們知道之後交代的就是摩訶和阿們.
以素以東之後所交代的故事.
這裡講到以實,瑪利.
這些都是神普遍的恩典.
就是阿布拉罕的子孫的遠房親戚.
以色列國在這個時候已經建了國.
在這些鄰國的中間.
就要問一個問題.
到底以色列神的子民應該有什麼身份.
有什麼使命.
你記得神呼召阿布拉罕的時候有什麼目的.
叫萬國得福.
所以神的子民聽見這些列祖的故事的時候.
他就要明白怎樣使萬國得福.
能夠認識耶和華.
你有沒有看這個故事的時候想起.
好像有點奇怪.
我們明明知道這個故事的結局.
結局就是最後以實,瑪利.

$^{241}$不是真命天子,不是真正的繼承人.
以色才是.
但是神都多方的去眷顧.
因為這裡講的是一種神學.
就是以色列的神.
這位阿布拉罕的神.
就是萬國的神.
萬國的主宰.
最明顯創世記最結尾的時候.
就是神藉著約瑟.
不單止救了亞國家.
也救了所謂天下的人.
神的子民.
可不可以去祝福萬國.
這就是創世記想答的問題.
這些列國能不能藉著以色列民.
能夠去認識耶和華.
能夠得到真正的祝福.
或者我們今天套用這個情境.
我們作為基督徒.
你有沒有聽過很多人信主的故事.
往往是因著什麼原因.
是因著他們在信主的人身上.
得到一些祝福.
得到一些服侍.
得到一些擺賞.
神的子民.
由古至今.
都是活在萬國之中.
可不可以我們意識到這方面.
以致我們能夠明白.
我們今時今日的使命.
有一點像耶穌所說的.
在登山寶訓.
《馬達》第五章第十六節.
祂說照樣你們的光要照在人們的面前.
在眾民萬國的面前.
讓他們看到你們的好行為.
還稱讚我們天上的父.
來到《創世記》第十七章.

$^{281}$亞伯蘭九十九歲的時候.
又說向祂顯然對祂說.
我是全能的神.
你要在我面前行事為人.
你要作完全人.
我要與你立約.
要使你的後裔人丁慶和.
亞伯拉罕的故事由開始到現在.
神是一步一步地去揭示祂的計劃.
去啟示祂的計劃.
現在來到第十七章.
就是整個亞伯拉罕的故事裡面.
啟示的最詳細的一章.
裡面有什麼元素呢?.
我們列出幾樣.
第一樣就是一個新的名字.
亞伯蘭就變成亞伯拉罕.
或者音可以變成解作.
多國之父.
Father of all nations.
Many nations.
撒萊變成撒拉勒.
音也像多國的母或者皇后的意思.
近東的文化裡面.
名字不單單是一個標誌.
名字往往是表達一個人存在的目的.
他的destiny.
他將來要成為怎樣的一個身份.
講出亞伯拉罕和他的子孫神國的子民.
他的destiny是為了什麼?.
是為了萬國.
第二個元素是什麼呢?.
第二個元素就是君王會從亞伯拉罕而出.
換句話說.
之前所說過的.
我要祝福你.
要使你成為大國.
這裡有進一步告訴亞伯拉罕.
因為將會有皇室從你身上出來.
比之前神所認許的更進一步啟示.

$^{321}$第三就是以撒.
終於看到他的名字.
十七章第十九節.
神說你的妻子撒拉.
真的要為你生一個兒子.
你要給他起名叫以撒.
我要與他建立我的約.
作他後裔的榮約.
至於以撒瑪尼.
我也應找你.
看來我已經賜福給他了.
使他昌盛子孫.
會繁多眾多.
他會有十二個族長.
成為一個大國.
但是我和你的約.
和以撒所建立的約.
這個以撒是明年這個時候.
撒拉要為你而生.
更加清楚.
說這個真命天子.
真正會延續這個聖約的.
就是以撒.
沒有其他人.
只有艾瑟.
第四個元素就是約.
約.
covenant這個字.
在十七章出現很多字.
代表雙方來建立.
在近東的約裡.
永遠有兩方面一起建立.
十七章也有這些字眼.
十七章第一節.
我是.
我要與你立約.
我要等等.
as for me.
神是這樣說.
但同樣也要求阿布拉罕.

$^{361}$as for you.
你要在我面前行事為人.
去做一個完全信靠我的人.
第五個元素就是約的記號.
就是割禮的出現.
十七章第十三節.
對阿布拉罕說什麼呢?.
他說你在家裡生的.
你用銀紙買的.
都一定要受割禮.
要受circumcision.
我的約要刻在你的身上.
作為永遠的約.
我們知道這個割禮.
不是希伯來人發明的.
當時近東很多不同的民族.
都會行割禮.
譬如在埃及的壁畫.
有集體行割禮的儀式.
可能是一些成人禮.
是一種這樣的儀式.
神就是用當時阿布拉罕.
所熟悉的禮儀.
然後加上一個新的意義.
這個新的意義.
可能對我們現代人來說.
不是那麼能夠接受.
但這些在近東的社會裡.
就能夠接受.
在一個男人的生殖器.
那裡割了羊皮.
這個意思有一點點詞.
想去表達.
就是你能夠有後裔.
你能夠有後代.
是全靠我和你立的約.
是耶和華的能力.
使你能夠繁殖.
能夠成為大國.
今天我們當然是阿布拉罕的後裔.

$^{401}$因為加爾泰書告訴我們.
每一個我們的信心.
是和阿布拉罕一樣.
我們就成為了阿布拉罕真正的後裔.
他屬靈的後裔.
你看看今天這個世界.
有多少人在歷史歷代裡.
信奉耶穌基督.
去信奉神.
真是阿布拉罕的應許.
已經在我們眼前實現.
並且會繼續延續下去.
一天到耶穌基督還沒有回來之前.
神的工作也會延續下去.
會有更多人被上帝拯救.
更多人歸於上帝的名下.
到底你是否這樣呢?.
如果我們看了阿布拉罕.
這段信心旅程的一部分.
我們可以有些take home.
有些應用.
我們作為阿布拉罕屬靈的子孫.
到底有什麼提醒呢?.
阿布拉罕的故事.
當然是一個story of faith.
講述我們如何跟隨主.
如何跟隨神.
我們可能會覺得.
如果我們像阿布拉罕.
神可以直接跟他說話.
這樣就好了.
我們不需要在基督徒的旅程裡.
在我們信心的道路上.
我們經常要猜測.
但你不要弄錯.
當你看阿布拉罕的故事的時候.
你數一下神直接跟他說話.
有多少次?.
十隻手指都能數到.
並且每一次神直接跟阿布拉罕說話的時候.

$^{441}$跟他啟示之後.
往往不是回答很多問題.
反而會產生更多問題.
所以我們今天不需要羨慕.
阿布拉罕的時候.
有神能夠直接跟他說話.
太好了.
今天我們有聖經.
今天我們有聖靈來帶領我們.
今天我們有身邊的弟兄姊妹.
和聖靈的群體來帶領我們.
其實這個更加清楚.
我們再看《創世記》第十六章第十六節.
十六章最後一節.
阿布拉罕八十六歲的時候.
下格維他生了以實瑪利.
這個時候阿布拉罕認為.
起碼在他眼中.
以實瑪利就是後繼之人.
但到了十七章第一節.
我們再下一節就是十七章.
阿布拉罕九十九歲的時候.
又說了向他顯現.
你要作完全人.
我要為你納若.
之後就跟他說什麼.
其實以實才是真正的意志.
八十六歲.
以實瑪利生了.
九十九歲.
接著到一百歲.
以實就生了.
試格十三年.
人生有多少是十三年.
這十三年裡面.
你再想想.
阿布拉罕一直以為.
以實瑪利由他出生到他十三歲的時候.
他就認為你是這個.
沒有其他人.

$^{481}$你說是不是神在玩弄他呢?.
是不是在耍他呢?.
當然不是.
這個旅程是想說什麼?.
就是想道出信心是什麼.
信神.
不是因為我已經掌握了所有細節.
所有將來的事.
阿布拉罕很明顯沒有一個水晶球.
以實瑪利出生.
他當時以為這個就是真正的繼承人.
以實瑪利出生不是一件邪惡的事.
不是一件壞事.
這個過程當然有人的軟弱.
有人的破碎表達出來.
但是神仍然祝福.
神仍然祝福以實瑪利.
使他成為大國.
所以極其量.
以實瑪利的出現.
他的出世.
都是一個擾導.
一個detour.
但是上帝不斷推進他的計劃.
不斷在經驗裡面.
教導信心的功夫.
信心和將來很多事.
都是手牽手.
兩件事都是掛鉤.
因為要信神.
就要信得過神.
對你將來的帶領.
這個必定我們要面對.
今時今日現代的人.
我們現代的信徒.
可能我們面對的信心問題.
不是生一個孩子.
不是傳宗接代.
可能我們面對的.
是不同的東西.

$^{521}$去留問題.
我們工作的前景.
我們的兒女.
如果你有的話.
他們將來會怎樣呢?.
這個都是一個將來信心的問題.
家庭.
或者甚至社會.
政局,世界都一樣.
將來會怎樣?.
疫情好像快要過去.
那又如何呢?.
我們都是面對將來.
面對很多不確定性.
一些未知之數.
神往往都和我們說.
正如他和亞伯拉罕的工作一樣.
將來會怎樣?.
到時我才會揭曉.
但今天.
每一天你要做什麼呢?.
十七章第一節.
神和亞伯拉罕說什麼?.
祂說我是全能的神.
你要在我面前行事為人.
要做完全人.
或者我們現代人聽見完全人的時候.
可能有些誤解.
以為是完美.
或者一個更加好的翻譯.
是entirely.
即是全心全意.
整個人.
你都要信靠,忠於上帝.
是不猶豫的.
整體我們去交託給上帝.
去相信祂.
你和我去信神.
你和我去信主.
去信靠我們的上帝.

$^{561}$不是因為我們有水晶球.
我們知道所有將來的細節.
但上帝一定要說完所有細節.
我們才去跟隨他.
信神是因為我們知道.
上帝是那位守約者.
上帝是那位值得我們信靠.
是那位faithful的上帝.
疫情期間.
我相信我們全球的人.
都會面對生活大亂.
我們教會.
或者我的事奉,我的工作.
都有很多混亂.
我都不知道明天將來會怎樣.
但我知道.
上帝對我今天的要求.
就是叫我盡力去做.
疫情帶來很多未知之數.
有很多未知之數.
是我們不能預見到的.
是要一步一步上帝才能夠揭示.
一步一步上帝才能夠明顯地.
讓我們知道祂的帶領.
我要做的事.
就是要去相信這位主.
無論怎樣都好.
我都要盡力去做.
盡心盡力地.
去跟隨祂.
或者你現在面對的.
是另外一些方面的挑戰.
另外一些生活的混亂.
上帝對你說.
你要信得過我.
你要相信上帝.
因為我是把握著你的將來.
而你今天要做什麼.
就是要盡心盡力.
來被我所用.

$^{601}$疫情到現在已經有年多.
我都想分享一下.
我自己的一些心聲.
一開始有疫情的時候.
世界大亂.
教會工作等等.
不能夠實體照常舉行.
我們搬到網上.
真的.
當時我只是憑著一個很小的信心.
儘管做.
盡量去做.
能夠祝福更多的人.
居然上帝去.
少許地使用我們.
求主給我們一顆小小的信心.
無論你現在的處境是怎樣.
無論你能否見到將來.
或者你在混亂裡面.
有很多的不確定性.
你儘管跟上帝說.
就是上帝.
你會一步一步去啟開你對我的計劃.
但你對我此時此地要求.
就是我在你面前.
做一個守約.
行事為人.
你都可以見證和事奉.
被你所用.
這就成為我們的祈禱.
可不可以?.
我們同心低頭.
主,我們多謝你.
因為你是那位帶領我們將來.
並且一步一步去揭開你對我們美好計劃的神.
你是那位大能的主宰.
在我們面前遇見很多的思念.
未知之數.
有很多的混亂.
我們知道你必定能夠克服.

$^{641}$你必定能夠帶我們經過.
主,此時此地.
求你親自來幫助我們.
我們成為立約的子民.
來跟隨你的時候.
我們願意讓你盡力去使用.
盡心盡意的.
我們就在今天來事奉你.
去見證你.
願你藉著我們.
能夠使更多人來得到你的福氣.
我們感恩祈禱.
奉主耶穌基督的名字.
\newpage



\section{}
\label{sec:hPfz5WlNiPQ}
\textbf{Sam Johnson Kingdom Values LV}
\newline
\newline
連結: \href{https://youtube.com/watch?v=hPfz5WlNiPQ}{\texttt{https://youtube.com/watch?v=hPfz5WlNiPQ}} ~~~~ 語音日期: 2021-06-26
\newline
\newline
\hyperref[sec:17rNRRMvNT8]{\small{< < < PREV SERMON < < <}}
~
\hyperref[sec:index_chronic]{\small{[返順時目]}}
~
\hyperref[sec:index_scriptual]{\small{[返順卷目]}}
~
\hyperref[sec:p69k7rhhWPQ]{\small{> > > NEXT SERMON > > >}}
\newline
\newline
$^{1}$Good morning family and welcome. I'm excited that we're here again together to open God's Word..
And once again we're looking at the Kingdom Values found in the Sermon on the Mount, Matthew chapter 7..
And today is kind of a special day and an exciting day because we are finishing up this study of Kingdom Values..
We're finishing up the study of the Sermon on the Mount..
And we have been looking at it for a very long time now, for over a year..
There is so much truth in this passage and so we've just been kind of going really slow verse by verse.
and mining those truths, digging those things out..
And I'm excited that we've gotten this far. I appreciate you all for sticking with it..
I know this has been a long study and if you've been watching this from the beginning.
then you've been in this passage of scripture for a long time..
But it's really so good for us and has so much truth and so much that we've learned..
I hope that on this journey of walking through this passage that you've really learned a lot of God's truth..
It's been really transformative to my heart and to my life..
And I pray that it's been transformative to your heart and to your life as well..
But today we're going to finish it up..
We're in the last couple verses of Matthew chapter 7 verses 28 and 29..
So let's read these verses together and consider what they're saying to us this morning..
This is Matthew 7 verses 28 and 29. It says this,.
"And when Jesus finished these things, the crowds were astonished at his teaching..
For he was teaching them as one who had authority and not as their scribes.".
Now this is just kind of a little coda tacked onto the end of the Sermon on the Mount..
It doesn't seem like it's really that significant, these two verses..
But there actually is a lot of significant truth here that I want us to consider today..
What does it mean when it says that Jesus was teaching them as one who had authority and not as their scribes?.
In verse 28 we see that it's pretty clear that whatever way Jesus was teaching,.
it was really revolutionary to his listeners, to the people who were listening..
Verse 28 says they were astonished at his teaching..
They couldn't believe what they were hearing..
And verse 29 tells us that what was astonishing to them was the authority that he taught with..
The authority that he taught with..
And what was that referring to?.
Well, in Jesus' day and in his time, the scribes were people who knew the Scripture very well..
The Pharisees, many of them would have been scribes as well..
They were scholars of the Scripture and of the Old Testament Scripture..
And not just of the Scripture, but of all of the rabbinical teachings that they practiced in Judaism along with the Scripture..
And so the scribes were these people who were scholars that really knew the Bible,.
knew all these other rabbinical teachings, knew all this minutiae and these obscure facts about the meaning and all these things like that..
And so these scribes were people who taught and who expressed that scholarship to others in the synagogues..
And people were used to hearing this type of teaching..
But Jesus was teaching in a way that was very, very different..

$^{41}$I think what we see in this passage is that Jesus was teaching different because Jesus was teaching on his own authority,.
while the scribes taught based on scholarship..
The scribes taught in a way that's very similar to the way that we often teach today the Bible, the New Testament and the Old Testament..
By looking at the context, by looking at the history, by looking at what other people have said about the passage and how they've interpreted it,.
and trying to determine what the meaning of the passage is..
That's how the scribes taught..
They taught based upon their scholarship, based upon their knowledge, based upon the things they had read,.
based upon the things that they had been taught by others, based upon the knowledge that they had gained in their lifetime of study..
And there's nothing wrong with that..
But Jesus was teaching in a very different way..
Jesus was not teaching based upon his scholarship..
Jesus was teaching based upon his own authority..
We actually see this in Matthew chapter 5, in verses 21 and 22 and 27, where Jesus says, "You have heard this..
You've been taught this..
You've been told that it was this way..
But I say to you...".
That's an example of what we're talking about here..
Jesus was saying, "You've been taught by the scribes, by the scholars, by the people who say they understand the scripture,.
and they've taught you that things are like this, or that obeying God is like this, a relationship to God is like this..
But I am telling you that it's different..
I am telling you, based upon my own authority, that relationship to God is different than you think,.
or relationship to God is different than what you've heard, or relationship to God is different than what you've been told.".
Jesus was teaching on his own authority..
Now, how is it that Jesus could speak with this authority?.
And that was what was so astonishing to the people that were listening..
They were like, "I can't believe this guy, a carpenter from Nazareth,.
in Galilee, in this rural, unheard of place, is teaching with this kind of authority.".
Obviously, there was clear, deep truth and wisdom in what Jesus was saying..
And Jesus was able to really incisively and perceptively really cut through a lot of the cultural trappings.
and all the other things that were so confusing that the scholars got caught up in..
They argued about all the minutiae, just like scholars today argue about all the minutiae of different things..
And Jesus was able to cut through that and to speak clearly and succinctly to the issues..
So I think that was a part of what people were astonished at..
But they were also astonished at the fact that Jesus was not only teaching with amazing wisdom and amazing perception.
and amazing understanding of God's truth, but he was saying, "Hey, believe this because I'm telling you.".
He was calling them to believe it based on his own authority..
And that was just mind-blowing. It was just astonishing..
How could he do that? How could he claim to know that much about God and about relationship to God?.
How could he claim to know that much about the truth? How could he claim to know that much?.
And that's what was astonishing to the people..

$^{81}$And that's a question that I want us to examine today for us as well..
How do we know that Jesus was able to speak with such authority?.
How do we know that all these things he's saying in the Sermon on the Mountain, throughout the Gospels, really are right and true?.
And first, I want us to consider this, that Jesus speaks with authority because he is the Word of God,.
or the truth of God, or the knowledge of God. That's who Jesus is..
Let's look at the book of John, the very beginning of the book of John..
In chapter 1, verse 1, we'll read the first five verses. It says this,.
"In the beginning was the Word, and the Word was with God, and the Word was God..
He was in the beginning with God. All things were made through him,.
and without him was not anything made that was made. In him was life,.
and the life was the light of men. The light shines in the darkness, and the darkness has not overcome it.".
Now that first verse says, "In the beginning was the Word.".
And in the original language, that word is "logos," and it's this word that is talking about the knowledge, the truth of God..
And we know from Scripture that God's knowledge encompasses everything. God knows everything in the universe..
And then as this passage goes on, it says, "The Word was God," that that knowledge, that the truth of God,.
is a part of God's character, it's a part of who he is. And then it says, "And the Word was God.".
And then it says, "He was in the beginning with God." So we start to see that the Word is embodied in the Godhead..
It's embodied in a person. And as we read on through this passage, we start to see that the Word,.
the knowledge of God, the truth of God, is embodied in the person of Jesus Christ..
"All things were made through him, and without him was not anything made that was made.".
That tells us about the deity of Christ. "In him was life, and that life was the light of men.".
And as you go on and read in John 1, it becomes clear that it is talking about Jesus,.
the one who came to earth to be our Savior. Jesus is the embodiment of truth..
Jesus is the embodiment of God's wisdom, of God's knowledge, of God's understanding..
Jesus is the embodiment of those things. All of the knowledge and wisdom and truth of God,.
which is all of the knowledge and wisdom and truth in the universe, because God is omniscient. He knows all..
It's all wrapped up and embodied and personified in Jesus Christ. The scripture is clear about that..
And so when Jesus began to speak about truth and wisdom and knowledge,.
he knew exactly what he was talking about because it was him..
That's a very difficult thing for us to understand because we think of truth and wisdom and knowledge.
and those types of things as being outside of ourselves, as being some type of standard that.
that we may compare something to, but not as being a part of us. And that is true for us as human beings.
because we are imperfect and our knowledge is limited. And there are many things that we do not.
know and many things that we do not understand in many areas where we do not have wisdom..
But that is not true of God. God is different than us. God is completely beyond us in every way..
God is in a class by himself. And God ultimately is the standard of righteousness and holiness and goodness..
God himself is the standard. How do we know what goodness is? We know that because God has revealed.
to us who he is. He has revealed his character to us. And we know goodness because we know God..
And as we know God more, we begin to understand more of what goodness truly is. How do we know.
what righteousness is or what holiness is? How do we know what pleases God and what's the right.

$^{121}$way for us to live and what's the right way for us to treat other people? How do we know what.
justice is? How do we know those things? Do we define them ourselves? Do we try to imagine what.
we think would be best in the world or what we think would be a standard of right or what we.
think would be holy or sacred? No, those things are defined by God himself and by his character..
Who he is defines to us what is right, what is just, what is holy, what is sacred, what is sinful and.
evil and wrong and off limits. All of those things are defined by God. And God personifies those good.
things, those good attributes of his character in perfection. God is perfectly holy. God is.
perfectly just. God is perfectly righteous. God is perfectly good. He personifies those things.
completely in perfection. And in that way, God is the standard of what those things are in our lives.
and in the universe. God's the standard of goodness, righteousness, holiness, justice, everything else..
Jesus being fully God was the standard of righteousness and of holiness and of truth..
And so when Jesus began to talk about these things, he knew what he was talking about.
because this was him. This was who he was. This was his character. To begin to speak about the.
values of the kingdom was to begin to speak about his own character. Jesus was the earthly.
personification of these things. And so when he said in verse 21 and verse 27, "Listen, I know.
you've heard from other people that this is what it looks like to be right with God, or this is.
what God wants, or this is what righteousness looks like, but I'm telling you something different.".
He knew exactly what he was talking about and he meant exactly what he said because he was.
the personification of those values of that goodness, of that righteousness, of that holiness..
Jesus spoke with this authority and Jesus still in our day speaks with this authority because he.
is the word of God. He is the truth of God. He is the wisdom and the knowledge of God,.
personified and embodied in human form..
Jesus was also able to speak with authority because he is perfect relationship with God..
That's one of the huge things in the Sermon on the Mount that I hope you did not miss. I talked.
about it over and over again, is that Jesus is really pointing us towards relationship with God.
the Father. Throughout the Sermon on the Mount, Jesus repeatedly contrasts the emptiness of just.
external religious conformance with the vibrancy of living in relationship with the God who made.
us and loves us. And Jesus makes it clear that what God desires is relationship with us. That's.
what God desires between us and him. God created us for relationship with him, for a healthy and.
right relationship. And that's what he wants. And God is not trying to get us to just simply fall.
into a system of religious conformance and where we do external things. We do outward observances.
to try to make it look like we're a good person. God wants us to know him from our hearts and to.
love him from our hearts and to live rightly because we are full of love for him. And Jesus.
just repeatedly emphasizes that in the Sermon on the Mount. He repeatedly draws that contrast. He.
repeatedly talks about how dangerous hypocrisy is and how dangerous empty religion is and how.
dangerous it is to try to please God in other ways. And he repeatedly calls us back to relationship.
with God and the importance of relationship with God. And again, Jesus could speak with authority.
about this because who better to teach us about relationship with God the Father than the one who.
has a perfect relationship with God the Father. In the book of John, Jesus speaks about this in.
different places. In John chapter 10, verse 30, he says, "I and the Father are one." He's saying.

$^{161}$that he and the Father agree in every way. Their will is totally aligned. And we know the doctrine.
of the Trinity, it's an amazing and beautiful doctrine. It's difficult for us to comprehend.
sometimes, but God is one God in three persons. And so even though Jesus is a separate person.
from God the Father, he is one with God the Father in every way. They are aligned in every way and.
their desires are the same in every way and their activity is the same in every way and they are.
always in agreement. John 14, 9, it says, "Jesus said to him, 'Have I been with you so long and.
you still do not know me, Philip? Whoever has seen me has seen the Father.'" Again, Jesus is showing.
his disciples, his followers that he is one with the Father, that they are always in perfect.
agreement. Their relationship has no distance. It has no conflict. It has no tension. There is.
never a time where Jesus is headed in one direction and God the Father is headed in the other direction..
They are always in agreement. Jesus lives constantly in perfect relationship with God the.
Father. And so when Jesus begins to speak about how we can have right relationship with God in.
the Sermon on the Mount, he speaks with authority because he knows exactly what that is and he knows.
exactly how to achieve it. And Jesus spends so much time in the Sermon on the Mount speaking about.
that, how we can have right relationship with God. It is one of the major themes of the Sermon..
And as I said, I hope you have not missed that as we have gone through it. The Sermon on the Mount,.
Jesus is calling you into a loving relationship with God that manifests itself in the values of.
the kingdom and you living your life. But that is what it is about. It is about a relationship.
that Jesus is inviting you into. And that brings us to the last thing I want us to consider,.
which is that Jesus is speaking with authority because he is the way to God. If one of the major.
themes of the Sermon on the Mount is this relationship with God and how to walk in.
relationship with God instead of just trying to merely conform to external religious observances,.
then one of the big questions that that leaves us with is how does that relationship.
take place? How do we know the way? Jesus himself is the way. In the Sermon on the Mount,.
Jesus has laid out this striking and amazing vision of the kingdom of God. And Jesus has.
described the good life for us. He has described the good life for us as humans by articulating.
the values of the kingdom. He said, look, if you want to live a life that is full of fulfillment.
and joy and peace and all of these beautiful things that you desire, the shalom that you desire,.
then you need to live out the values of the kingdom. And if you do so, you'll be living in.
right relationship with the creator God who made you and in right relationship with the rest of the.
human beings around you. Jesus is painting this beautiful picture for us throughout the Sermon on.
the Mount. He's described it for us. But the truth that he's also been very honest with us about is.
that there is no way to the kingdom of God through our own effort or through our own merit. The.
kingdom of God is not accessible to any of us humans by our own effort or our own merit. And.
living out these beautiful values of the kingdom that Jesus has described for us, living out the.
good life and experiencing shalom in the way that we want is not accessible to any of us humans by.
our own effort or merit. Jesus is the only way to right relationship with God and right relationship.
with other humans in our world. And Jesus himself is the only way to the kingdom of God. We talked.
about what the kingdom of God is at the beginning of this study, that the kingdom of God is both the.
future kingdom of heaven, where it will be fully manifested, where there will be no crime, where.

$^{201}$there will be no hatred, where there will be no sin, no lying, no disease, no COVID. All of those.
things will be completely eradicated in heaven because the kingdom of God will be fully manifested..
And that is the future kingdom of God. And there is no way to the future kingdom of God in heaven.
apart from Jesus Christ. There's no way to get there based upon our own efforts or our own merit..
But there's also the present kingdom of God, whereby believers in Jesus, the subjects of the.
kingdom, those who have put their trust in Christ can manifest the kingdom here on earth as they.
submit to the rule and reign of God. And as they submit to the rule and reign of God in their own.
lives, they start to live out the values of the kingdom. And that is the present manifestation.
of the kingdom of God, even though it's not complete and full on earth, because there will.
always be brokenness, there will always be sin. We can be a part of the present manifestation of.
the kingdom of God, but the present manifestation of the kingdom of God, again, is only accessible.
through Jesus. There is no way for us to experience that through our own effort or through our own.
merit. We can't try hard enough to live out the values of the kingdom. We can't just try harder.
and learn to live out the beatitudes through our own effort or through trying to achieve something.
or through our attainments or through our accomplishments or through our education or.
through any other source. We can't buy it. Remember the book of Acts, the story of the man who tried.
to buy the gift of the Holy Spirit. We can't purchase it with money. We can't trade for it.
with anything we have. It's only through Christ. And the future kingdom of God is the same way..
We can't purchase it. We can't earn it. We can't get there by trying hard enough. It's only through.
Christ. And Jesus is so clear about that. Jesus is the only way to healing and shalom in our broken.
world. We live in such a broken world. We live in such a hurting world and we desire healing. And.
that's good. You know, if you don't desire healing in our world after the past couple of years that.
we've been through, then you need to examine your heart and see if you're still alive and healthy..
Because we have seen so much brokenness over the past few years. We've seen disease. We've.
seen COVID that has just taken so many lives and made so many people sick. And we've seen economic.
issues and troubles and struggles. We've seen all types of racial tensions and racism on display and.
evil things on display. We've seen so many people harmed. We've seen political turmoil. We've seen.
political unrest and oppression around the world and people being silenced and all these terrible.
things going on. Our world is deeply broken. And it's only right that as human beings, much less.
as Christians, that we should desire healing and shalom. And we should want to see things fixed and.
want to see peace and want to see wholeness and want to see health and want to see fullness in.
our world. Our heart should cry out for that. But here's the thing, as our hearts cry out for that,.
we have to see that the only way to experience it is through Christ. Every other avenue that we try.
to follow in order to experience those things on our own will end up falling short. Every road that.
we go down trying to experience healing and shalom on our own apart from Christ will end up being a.
dead end. We will not experience these things apart from Jesus. Jesus is the only way to healing and.
shalom. Jesus is the only way to the future kingdom of heaven. Jesus is the only way to the present.
manifestation of the kingdom of heaven here on earth in our lives. And Jesus is the only way.
to the beatitudes. If we want to live out the things that he talked about, being poor in spirit,.
being meek, being pure of heart, hungry and thirsty after righteousness, being a peacemaker,.

$^{241}$all of those beautiful things, being able to be strong and persevere during times of persecution.
and difficulty, all that stuff, that's only going to happen through Jesus. Any other pathway we'll.
try to find will not lead us there. And ultimately those other pathways will lead us to pain and to.
harm and to self-destruction. Jesus is the only way. Remember John 14, 6, Jesus said this explicitly.
to us. He said, "I am the way and the truth and the life. No one comes to the Father except.
through me." And that's why Jesus could speak with authority because he knew he was the way to God..
He didn't have to say, "Well, I think if you really dig into some of the history of this,.
or if you really examine the words here, I think it's saying that this is the way to heaven." No,.
Jesus already knew the way to heaven because he was the way to heaven. Jesus, one with the Father,.
incarnated, came down and took on human flesh in order to show us the way. And how did he do that?.
He did that by living a perfect sinless life. He never sinned. He was tempted in every way,.
just like we are, to sin. He was hurt. He was rejected. He was spit upon. He was mistreated..
He went through everything that we have experienced, but Jesus never sinned in any of it..
And at the end of his life, he chose to give his life for us, to lay his life down as a sacrifice.
in our place, taking on himself the punishment that we deserve because of our sins, because of.
all the stuff that we've done, all the way that we've mistreated everybody else, and all the.
self-destructive stuff we've done, and all the hateful things that we've done, and all the lies.
we've told, and all that. Jesus took that on himself, made himself our substitute. He took.
the penalty of sin and bore the wrath of God in our place so that because of what he did,.
we could be forgiven and we could have access to relationship with God. We could have access to the.
love of the creator God who made us. We could have access to restoration, to right relationship that.
had been broken in the beginning because of our own sin. It was our fault that we broke up in.
the first place. But God loved us and came to us in love and sent Jesus in order to open the door.
to heaven, in order to open the way for us to know God and have right relationship with him..
And that's the beauty of the gospel, but it's also the beauty that we see here in the Sermon.
on the Mount. The Bible is clear that this happens by grace and through faith, that we don't have to.
do anything, we can't do anything to try to earn it or to try to achieve it. We simply have to trust.
in Jesus that he has done that for us and that he is great enough and good enough to forgive our.
sins. And we ask him to come into our heart and to become the Lord and leader of our life and to.
begin to follow him in our life and to begin to seek to submit our lives to the rule and reign of.
God to the values of the kingdom. And as we do that, by God's grace, we begin to live out these.
values that we've been talking about over the last year. We begin to become so saturated with.
the work that Jesus is doing in our heart and life that the values of the kingdom are what come out.
of our lives. Whether it's in good times or bad times, whether things are easy or difficult,.
whether we're being challenged, whether we're rejoicing or whether we're mourning,.
the values of the kingdom are going to come out. God's love, God's peace, God's joy, humility,.
care for others, purity of heart, hungering and thirsting after righteousness. All of these.
beautiful things that Jesus articulated for us in this amazing sermon can come out of our lives..
And that's what we want and that's what we desire. Jesus spoke with authority because he is the word.
of God. He is perfect relationship with God and he is the way to God. And my prayer is that in my.

$^{281}$life and in your life, we would enter into relationship with God by grace and by faith in.
Jesus so that we can begin to truly live out these values of the kingdom. We can respond to the.
authority that Jesus spoke with, to the truth of his words, and we can begin to obey God and live.
in right relationship with him so that the values of the kingdom come out of our lives and flood.
every life around us and our family and our community and our cities, because that's what.
we desperately need in the world we live in. We desperately need the love of Christ. We desperately.
need the values of the kingdom. Let's pray. God, thank you for all that we have learned as we have.
walked through the Sermon on the Mount. God, thank you for your truth. God, thank you that we can.
trust your authority that Jesus spoke with authority because he knew what he was talking.
about because all of these things were who he was to the core of his being. And God, I pray that you.
would help us to trust. If we don't have relationship with you, God, I pray that we.
would start there by beginning relationship by faith and the grace that you've extended,.
the unmerited favor that you've extended to us. But for those of us who do have relationship with.
you and we've been walking with you, God, help us to trust you more and to trust that if we will.
really orient our lives around the values of the kingdom, if we will really surrender to living out.
those things that even though it may seem countercultural, it may seem difficult, it may seem.
even confusing at times, that you will work in us and through us and you will use what you are doing.
in our hearts to work in our world and to bring about healing and peace and love in the lives of.
others, God. We thank you for who you are. We thank you for this beautiful Sermon on the Mount.
that you gave us, God. Help us to believe you. Help us to trust you. Help us to live it out,.
God. Give us the grace to live out the values of the kingdom day by day. In Jesus' name, amen..
\newpage



\section{}
\label{sec:p69k7rhhWPQ}
\textbf{God's Goodness Episode 1 Gratefulness Towards God}
\newline
\newline
連結: \href{https://youtube.com/watch?v=p69k7rhhWPQ}{\texttt{https://youtube.com/watch?v=p69k7rhhWPQ}} ~~~~ 語音日期: 2021-07-04
\newline
\newline
\hyperref[sec:hPfz5WlNiPQ]{\small{< < < PREV SERMON < < <}}
~
\hyperref[sec:index_chronic]{\small{[返順時目]}}
~
\hyperref[sec:index_scriptual]{\small{[返順卷目]}}
~
\hyperref[sec:hd_nz1Oh6bM]{\small{> > > NEXT SERMON > > >}}
\newline
\newline
$^{1}$Good morning family and welcome..
I'm really excited that we're opening God's Word again together this morning..
And this morning we're going to start looking at something new, a new study here for the.
next few weeks, and that is really looking at biblical examples of our response to God's.
goodness, our response as the people of God to God's goodness and God's work in our lives..
And specifically this morning we're going to be focusing on gratefulness and on seeking.
to begin to have an attitude of gratefulness towards God..
So I know this is going to be really good for me..
Gratefulness is something that I always struggle with and something I need to focus on more..
So I'm excited..
I think it's going to be good for me and I pray that it'll be really good and encouraging.
for you as well..
We're going to be looking in Exodus chapter 15 this morning..
And Exodus chapter 15 is this song that's often called the Song of Moses and it's the.
song that the children of Israel composed basically extemporaneously on the spot, spontaneously,.
this song of praise and gratefulness to God just after they came through the Red Sea..
So if you're not real familiar with this story, I'll kind of give a quick narrative of it,.
but you could look at these chapters in Exodus for yourself, especially chapter 14, which.
tells the story of the Red Sea..
But the preceding chapters in the book of Exodus are the story of the children of Israel.
being delivered by God from the land of Egypt..
And if you remember that story, the children of Israel, the Hebrew people, the Jewish nation.
are in Egypt..
They've been there for generations and over time, what was originally a friendly relationship.
with the Egyptians has become an oppressive one..
And so Pharaoh has enslaved the Hebrew people..
And so they're slaves in Egypt and they're forced to do really hard labor to build for.
the Egyptians..
And also Pharaoh decides that he doesn't want their strength to grow, doesn't want their.
nation to grow..
And so he begins to kill their sons..
And so it's a really, really terrible, oppressive situation, the children of Israel, the Hebrew.
people are in..
They're God's people, but they're in a place of deep oppression and deep pain and deep.
hurt because they're enslaved and they're being treated unjustly in so many ways..
And eventually God raises up Moses to be a deliverer to them..
Moses goes and speaks to Pharaoh and tells Pharaoh that God wants Pharaoh to let his.
people go..
And Pharaoh resists and says no..
And so then Moses brings the plagues into Egypt, really God brings them into Egypt,.

$^{41}$but Moses is the messenger that continues to speak to Pharaoh and to the Egyptian leaders.
to tell them that God is bringing these plagues as a result of their rebellion against God.
and the fact that they won't let his people go..
And so all these plagues come against Egypt..
And again, if you're familiar with that story, you know what those are..
If not, you can look at it and read it..
But God does all these supernatural things basically..
Some of them are very fearful and terrifying, but he does them to show the Egyptians his.
own power and to force Pharaoh to let the Hebrew nation go free..
And so finally after the tenth plague, Pharaoh decides to let the Hebrews go and the entire.
Hebrew nation leaves Egypt and they go out into the wilderness..
Well Pharaoh quickly changes his mind and says let's go get them..
And at this time in world history, the Egyptian nation was really an empire, really a dynasty,.
and it was the most powerful empire on the face of the earth and they had the most powerful.
and fearsome army..
And so Pharaoh says we are going to go and get the Hebrews and bring them back and put.
them back to work and enslave them once again for us..
And he takes his enormous army to go get them..
And as I said, this was the most fearsome army in the world and they had the best technology..
They had chariots and other weapons of warfare that were the most advanced technological.
weapons that existed in the world at that time..
And so this was a very scary army and by now the Israelites had journeyed down close to.
the Red Sea and that's where they were when the army of Pharaoh caught up with them..
And God was leading them in a pillar of fire by night and a pillar of cloud by the day..
So they are literally relying on God day by day for guidance and for leadership and for.
provision..
Moses is leading them as a representative between the people and God, but God himself.
is leading them and protecting them..
But when Pharaoh's giant army comes up behind them, they are terrified because they find.
themselves closed in with the sea on one side and this massive fearsome army with all of.
these weapons of warfare on the other..
And they are afraid that they are just going to be destroyed and eradicated and killed.
and it seems like there is no way out..
So they cry out to God for deliverance..
God puts the thick cloud between his people and Pharaoh's army to keep Pharaoh's army.
away and then God begins to roll back the waters of the Red Sea..
And as I said, this story is in Exodus 14..
If you'd like to look at it, it's a really good story..
God moves back the water so much that there are basically walls of water on each side.
with dry land in the middle and the Hebrew nation are able to walk across on dry land..

$^{81}$And this is a large group of people by now..
The Hebrew nation has grown very large and very strong during this time in Egypt and.
so it's a huge group of people..
God holds back the waters long enough for them to get all the way across..
It probably took all night for them to cross..
And then God removes the cloud and the army of Pharaoh sees that the people are gone and.
sees the dry land in the middle of the sea and sees that they've been able to get out.
and decides, "Hey, we're going to pursue them.".
And begins to pursue and God lets back the water and the waters come rushing back in.
and destroy the army of Pharaoh..
And so God in one miraculous event, he delivers his people and he destroys the enemy that.
are seeking to harm them and to kill them..
So this is an absolutely amazing and miraculous event in the history of God's people..
And immediately after this, when they're safe on the other side of the Red Sea, they begin.
to sing this song of praise..
It's found in Exodus chapter 15..
So let's read some of it..
We're going to read the first 13 verses..
It says this, "Then Moses and the people of Israel sang this song to the Lord, saying,.
'I will sing to the Lord, for he has triumphed gloriously..
The horse and his rider he has thrown into the sea..
The Lord is my strength and my song, and he has become my salvation..
This is my God, and I will praise him, my Father's God, and I will exalt him..
The Lord is a man of war..
The Lord is his name..
Pharaoh's chariots and his host he cast into the sea, and his chosen officers were sunk.
in the Red Sea..
The floods covered them..
They went down into the depths like a stone..
Your right hand, O Lord, glorious in power, your right hand, O Lord, shatters the enemy..
In the greatness of your majesty you overthrow your adversaries..
You send out your fury..
It consumes them like stubble..
At the blast of your nostrils the waters piled up..
The floods stood up in a heap..
The deeps congealed in the heart of the sea..
The enemy said, 'I will pursue, I will overtake, I will divide the spoil..
My desire shall have its fill of them..
I will draw my sword..
My hand shall destroy them.'.

$^{121}$You blew with your wind..
The sea covered them..
They sank like lead in the mighty waters..
Who is like you, O Lord, among the gods?.
Who is like you, majestic in holiness, awesome in glorious deeds, doing wonders?.
You stretched out your right hand..
The earth swallowed them..
You have led in your steadfast love the people whom you have redeemed..
You have guided them by your strength to your holy abode.".
And I think there's just so much powerful truth in this passage for us about how to.
be grateful and how to learn to celebrate and reflect on and appreciate the goodness.
of God in our lives..
The first thing I want us to see here is that when we experience God's goodness in our life,.
when we go through something that causes us to reflect on God's goodness, we first of.
all need to really place our focus on God..
I love that about this song..
This is the song, like I said, it's often called the song of Moses, but this song doesn't.
talk about Moses at all..
It doesn't talk about the Hebrew nation at all, really..
I mean, it does some as their relationship to God, as the people of God, but it doesn't.
talk about them and their strength in any way..
It just talks about God and how powerful He is..
The focus in this is clearly on God because of the things that He has done..
And I think that's a really instructive thought for us when we go through things in our life..
A lot of times when something good happens and we experience God's goodness in some way,.
we begin to credit ourselves and we think, "Yeah, this happened because I did a great.
job..
I figured something out..
I tried harder..
I worked harder..
I am smarter than other people.".
We begin to credit ourselves or maybe we credit other humans around us and say, "Well, it.
happened because this person that's my friend or my family member or this person that I'm.
associated with in some way, they're really good at this.".
And that is true..
God gifts us and He gifts the people around us..
And it's certainly good for us to lift up the people around us and encourage them and.
cite their roles in any success that we might have..
But we have to be honest and begin to credit God..
In these verses you see again and again it talks about who is the strength in the song..

$^{161}$That's the Lord..
Who is the one who is worthy of praise?.
That's God..
Who's the one who has triumphed?.
It's the Lord..
Who is glorious in power and defeats the enemy?.
That's God and His right hand..
The focus is clearly on God..
And so many times when something good happens, when we get that promotion that we've been.
praying for, we do well in school and we finally graduate, or we get that degree that we've.
been working for for so long, or we graduate high school, or we get a new place to live,.
a new apartment or a new house, or we get a new car, you know, all those things that.
we pray and ask God for and we desire..
And then God brings those things into our life..
And then suddenly we just think that we did it ourselves and we forget to credit God and.
to begin to praise Him and to begin to focus on Him..
And you know, what's even more sad than that is oftentimes good things happen and we simply.
don't consider God at all..
Sometimes we just take things for granted..
Like it's just a coincidence that something happened in our life..
Like it's just a coincidence that that relationship in our life was restored..
It's just a coincidence that our health was preserved in some way..
We just act like all those things are coincidences..
I think it's pretty clear when you look at Scripture and the way God cares for His people.
that that's not the case..
God is at work and He is doing good things for His people..
And so we need to recognize that..
And we shouldn't just assume that everything's a coincidence or that it just worked out that.
way because, you know, we're a little bit better, a little bit smarter than someone.
else..
We should say, "Hey, this good thing that happened in my life, this is because of God..
And I believe that.".
Especially if it's something that we prayed and asked God for..
How self-absorbed are we if we pray and ask God for something and then He provides it.
in our life and we turn around and never say thank you and we just act like, "Yeah, this.
was supposed to happen anyway," or, "I expected this because I'm so great.".
And we never give the credit that is due to God..
We've got to learn to begin to do that more and more in our lives..
And we've got to learn to just start to pay attention to God's work..
Sometimes we're just ignorant of the fact that God is at work..

$^{201}$We don't even see it because we're just not paying attention..
Or sometimes even worse, we have an attitude of entitlement and maybe God does something.
really good in our lives and we just say, "Well, yeah, of course that's supposed to.
happen..
You know, it's supposed to work out that way for me because I'm a good guy..
I worked hard for this or this is what I deserve..
And so of course I was going to get that.".
And again, in all those different ways and circumstances, we forget to focus on God..
We forget to give Him credit for what He's doing..
We forget to praise Him and to thank Him for all of the good things that He does in our.
life..
And you know, another important thing that's really essential for us to do in order to.
be grateful is to begin to place our focus on God's work and on what God is doing in.
our lives..
And this is huge, especially in the midst of times of life where we may not readily.
see the goodness of God..
I understand that sometimes we say, "All right, I have no problem praising God and.
thanking God when God does something good in my life and when I can readily see God's.
goodness.".
And some of those examples I just gave and the last point are more examples of times.
that we can readily see God's goodness..
When we get a promotion or we get a new car or when we graduate, it's a lot easier for.
all of us to say, "Wow, this is cool..
This is good..
I'm so grateful to God," when something is going well in our life..
But the reality is even when things don't feel like they're going well in our life,.
God is at work and God is always at work for His glory and for our good in so many different.
ways..
And so we've got to learn in all circumstances of life, especially those difficult circumstances,.
those trials, those challenges, those difficulties, we've got to learn to put our focus on God's.
work and to say, "What is God doing here?".
And there are things going on in my life right now, even though it may not feel great, there.
are things going on that God is using and God is working and God is doing something.
and I can be grateful for that, even if the circumstance that I'm in the midst of doesn't.
feel great to me right now..
We've got to put our focus on God's work..
You know, we can do that by beginning to seek out areas of God's goodness in our lives..
I think we see this in the passage that we read here in Exodus 15..
Look at these verses where it talks about what God is doing..
It's talking about how God is, verse 7, God is overthrowing the adversaries and verse.

$^{241}$8, it talks about God standing the waters up and in a heap, piling the waters up..
And verse 9, it says, "The enemy said, 'I will pursue, I will overtake.'".
This is not just some overly rosy picture of the goodness of God that just says, "Yeah,.
we were walking through the park and everything was great and everything was fantastic, so.
we're going to praise God.".
This passage is saying, "No, we thought we were going to die..
The enemy was pursuing us and about to overtake us, and the enemy was saying that they were.
going to kill us with their sword and take the spoil, and that we were going to be destroyed.".
This passage makes it clear that this was a dark, challenging time..
This was a time of real difficulty in the life of God's people..
But even in the midst of that, God was at work..
And so we can find those times and those places in our life where God is at work, even when.
it's challenging..
We may have to search for it a little bit..
We may have to seek it out and say, "Okay, the circumstances are tough right now, but.
I believe God is at work, and I'm going to look for the areas where God is at work in.
my life.".
And if we'll begin to seek those out, we will begin to find areas that God is working and.
we can begin to praise God for..
Even when we're experiencing those trials and difficulties, we can be grateful..
We can be grateful for God's presence..
I've been through some difficult times in my life..
My wife and I have been through some very difficult circumstances, and the loss of family.
members and different things that we've walked through, different challenges that we've walked.
through..
And I can tell you, honestly, there's been a lot of times in the midst of those experiences.
where it didn't feel great to me, or it didn't feel great to my wife..
And where we've gone through those times of saying, "Wow, this does not feel good right.
now.".
There have been times when we've said, "You know, pretty much none of what we're going.
through right now feels good..
Nothing feels good right now.".
But even in those moments, we can be grateful for the presence of God and for the nearness.
of God..
And we can say, "You know what?.
Even though nothing feels good right now, and I'm discouraged, and I'm experiencing.
loss, and I'm experiencing grief, and I'm experiencing all of these challenging things,.
I'm so grateful that God is with me in this grief..
I'm so grateful that God is with me in this loss..
I'm so grateful that God is with me in this discouragement, in this setback, in this challenge.".

$^{281}$Because if God was not here, I would be completely alone, and I would be completely in the dark.
and lost on my own..
So even if we're in the toughest moments of our life, and nothing feels good, and everything.
seems to be going wrong, I get it..
That happens..
That's real..
And even in those moments, we can be grateful for the nearness of God, the presence of God.
that is with us, that sustains us..
We can be grateful for God's people, even in those challenging times..
God put us in this family called the church, and one of the things that we're supposed.
to do is love one another and care for one another, support each other, encourage one.
another, and hopefully we're doing that for each other..
Hopefully there are people in your life who love you and who care for you and support.
you when you go through challenges, and you can be grateful for that..
You can say, "God, I'm grateful for the people you've put in my life..
Even during a time of illness or during a time of grief, I'm grateful that these people.
are here..
I'm grateful that you're using them to encourage me and to build me up.".
We can be grateful for God's sanctifying work in our life..
That word "sanctification" is just a church word that means that God is purifying us..
He's changing us and making us more like Jesus..
He's taking away some of those things, those impurities, those things that are not healthy.
for us, those bad habits, that selfishness..
He's slowly taking that stuff out of our heart and out of our life and slowly turning us.
more and more like Jesus so that we can begin to truly love others and truly be unselfish.
and truly be focused on the right things in life..
The honest reality is in my life and I think in the lives of most believers, that sanctifying.
process works best when we're going through some suffering..
It's hard for us on a bright sunny day when everything is going our way to be unselfish.
and to be humble and to really care for others and to really focus on God because everything's.
great..
But when we're going through something difficult, suddenly it's so much easier for us to see.
our need for God and to focus on Him and to be humble and to love and care for others..
That sanctifying process a lot of times works best in times of struggle and in times of.
suffering and in times of trial..
We can be thankful for that..
When we're going through something difficult, we can say, "You know what?.
This doesn't feel good and I don't love being in this moment, but God, I'm so grateful that.
You're using it to work in me to make me more like Jesus, to strip away some of that unhealthy.
stuff, some of that stuff that causes me to be focused on myself, some of that stuff that.

$^{321}$causes me to be self-destructive in the choices I make or causes me to make choices that harm.
the people around me..
God, I'm so grateful that You're changing me and taking that away and making me more.
like Christ.".
Because if you're going through a challenge, if you're going through a trial, if you're.
going through a difficulty and you're a follower of Jesus, then I guarantee you that God is.
using it in your life to try to get you closer to Him because that's what God does..
We can be grateful for that sanctification that happens in difficult moments..
We can be grateful for God's strength and grace..
I don't know if you've ever gone through something challenging that you look back on and say,.
"I don't know how I got through that..
I've had experiences like that.".
And the truth is that God gives supernatural strength and grace to His people when we go.
through really difficult times of challenge and of grief and when we're stretched to the.
limit and we just feel like, "I don't even know if I can make it.".
God is there with us in those moments and He carries us and He gives us strength and.
grace..
And we can be grateful for that..
We can say, "You know what?.
The challenge I'm encountering right now doesn't feel fantastic, but God is getting me through.
and I'm thankful for the strength and the grace that He's giving me, He's sustaining.
me with.".
And we can be grateful for that growth..
God uses challenges and trials and afflictions to change us and to grow us..
Just like you won't strengthen your muscles, your physical muscles, without kind of stretching.
them and straining them a little bit..
That's literally what we do when we go work out at the gym or we work out hard in any.
other way, is we stretch and strain our muscles a little bit so that they'll become stronger..
And God does the same thing with us spiritually..
A lot of times He has to take us through something that stretches us and strains us a little.
bit in order to strengthen us so that we can grow and be stronger than we were before..
And we can praise God for that..
Even in a tough time, we can say, "God, I thank You and praise You that You stretched.
me, that You strained me, and that You're growing me.".
The next thing that I think is really important to help us begin to practice gratefulness.
in our response to God's work in our lives is to intentionally celebrate God's goodness..
When God does something good and we begin to notice that, we begin to see God's work.
in our life, and we begin to see God's goodness in that work when we say, "Okay, I see God.
at work, even in my difficulties..
And I'm paying attention and I've searched out for God's goodness..

$^{361}$And I do see it, even in this difficult time, I can still see how God's working.".
Or maybe in this really good time, I can see how God's working..
God blessed me and He brought this wonderful thing into my life, and I see God's goodness.
there..
Once we see it, and once we see how God is working, we should begin to intentionally.
celebrate it and to say, "You know what?.
I'm not just going to notice it, I'm going to celebrate it..
I'm going to begin to thank and praise God for His goodness.".
And you might say, "Well, what does that look like?.
What does it look like to thank and praise God for His goodness?".
Exactly what it sounds like..
You begin to thank God..
You can thank God in your prayers, silently..
You can thank God verbally and say, "God, I thank you for what you're doing in my life..
God, I thank you for the way you've saved me, for the way you're changing me, for the.
way that you're growing me into a new creature in Christ..
God, I thank you for the blessings that you've given me..
I thank you for my family, for my community, for my church, for all these different things.".
That's what it sounds like, thanking God..
And obviously we know that the church has a long and beautiful tradition of praising.
God in song, just like here in Exodus 15, where we lift our voices in song to thank.
God, to praise Him, to exalt Him, to talk about how good He is, how great He is, and.
how grateful we are for what He does in our lives..
And that's a beautiful way to praise God, is to praise God in song..
You can do that by yourself..
You can do that in a church service with other believers..
You can do it when you're driving in the car or when you're in the shower..
You can sing anywhere and thank and praise God..
And that's a good way to intentionally begin to celebrate God's goodness in your life..
And so I hope that you're doing that..
But you can also take a break from the busyness of life to celebrate and reflect and praise..
And I think this is really important..
Sometimes we do see God's goodness and we say, "Okay, that's something I know I need.
to praise God for.".
But our life is so busy that we just kind of go on by, and then a week and a month later,.
we've kind of forgotten about the good thing that God did..
And we see this example throughout Scripture, and we see it clearly in the Old Testament..
We see it clearly here in Exodus 15..
They got a cross and they were safe, and the people who were trying to kill them were destroyed..
They could no longer chase after them..

$^{401}$And they didn't just say, "All right, on to the promised land, on to the next thing..
Let's get to the next spot.".
They said, "No, we're going to stop..
We're going to take a break, and we're going to sing and dance here on the shoreline and.
praise God for what He's done and for how He delivered us and for His goodness that's.
been at work in our lives.".
And like I said, we see that throughout the Old Testament..
God clearly tells His people so many times to do that..
I think we need to practice that..
As modern believers, especially in the U.S., that is not something we practice well..
Our lives are so busy, and it's all we can do just to find some time to get to church.
and to sing two or three songs..
But we don't take breaks from the busyness of life to just praise God, to just reflect.
on how good He's been and to just thank Him for it..
And we need to do that..
And that's something we should practice..
If you're a single person, you should practice that and say, "Hey, I'm going to take a break.
from my work schedule, from my friend relationships, so that I can just reflect on God's goodness.
and begin to praise Him for it.".
And you can do that with other people..
You can do it with friends..
You can do it in the company of others..
But take a break..
If you're a married person, especially if you have children, say, "Hey, we as a family,.
we're going to take a break because God did something big in our lives, and we want to.
stop and we want to celebrate it..
Maybe it's just for a few hours..
It's a special dinner..
Maybe it's a half day or a day off.".
But take a break from the busyness of life and say, "You know what?.
I'm going to celebrate what God has done..
I'm going to reflect on it..
I'm going to think about it..
I'm going to consider how good God has been in my life..
And I'm going to talk about it..
I'm going to praise Him..
I'm going to thank Him for what He's done in my life.".
And you know, God has given us so many things in our life that sadly a lot of times, like.
I said, either we take for granted or we're just not grateful for..
And all of us have things to thank God for..

$^{441}$You may say, "Well, you know, life's been tough, and there's some things going on in.
my life that are challenging right now..
And so it's hard for me to feel like praising God.".
But the reality is all of us have things that we can praise God for..
So I want to think about a few of those things as we close out and just give us some reasons.
that we can praise God..
You know, God has given all of us as humans special talents and gifts and abilities..
And we can thank and praise God for those things..
You have special talents and gifts and abilities..
Maybe they're artistic..
Maybe they're creative..
Maybe they're music..
Maybe they're analytical..
Maybe they're scientific..
Maybe they're kind of interpersonal..
Maybe you're really good at relating to other people or you have a great sense of humor..
Whatever those things are, you can thank and praise God for those..
And it's easy for all of us to look at the talents and gifts that other people have and.
feel like we wish we were someone else and we wish we had their talents and forget about.
the beautiful talents and gifts that God has given us..
So you can and should praise God for who He made you to be and all the talents and gifts.
that He's poured into your life..
Another thing that all of us can praise God for right now is some of the freedoms that.
we enjoy in our country..
Today is July 4th..
It's a day that we celebrate the beginning of our country..
And we can all thank and praise God for the freedoms that we enjoy in the United States..
Obviously, our country is not perfect by any stretch of the imagination..
There's still a lot of problems and a long way to go..
But we do have a lot of freedoms here and people have a lot of good blessings in our.
nation..
And I think that's something that we can praise God for and be grateful for..
And we should praise God and be grateful for..
We can be grateful for COVID-19 being on the decline here..
I think that's something huge..
As I talked about earlier, there's so many things we take for granted that God does..
And I think it's easy for us to forget that a year ago and a year and a half ago, we were.
really worried about this illness and that we were praying and crying out to God for.
deliverance that God would deliver us and deliver our families and deliver our communities..
And God has answered those prayers in many ways..

$^{481}$And I don't say that to diminish the fact that there's still sick people..
I understand that..
And we need to continue to pray for those people..
And especially in other parts of the world where there's still higher levels of disease,.
we need to be praying for those places desperately..
But we can also celebrate and thank God that He's brought about so much deliverance from.
this pandemic and that a lot of things are reopening and people are healthy and so many.
people have been able to take the vaccine and there's much less disease and much less.
death..
Hey, we prayed and we asked for that..
And God has answered that prayer..
And we need to give credit to God..
Obviously, God used a lot of humans to do that..
And that's fine..
We should thank them too..
But we need to give credit to God and we need to praise Him for deliverance from this pandemic..
Just like the children of Israel in Exodus here, we're praising God for deliverance from.
a fearful moment..
We can praise God for deliverance from a fearful moment..
We can praise God for the beautiful summer weather..
I hope you've been enjoying the weather lately and getting outside..
And that's something we can praise God for..
We can praise God for the simple things like the fact that He gives us food and shelter.
and clothing and meets those basic daily needs in our life..
We can thank God for a lot of other beautiful things that He pours into our life that are.
so enriching..
We can thank God for life itself and for the beauty of being able to be alive and for the.
health that He gives us..
We shouldn't take that for granted..
The fact that God wakes us up every day and that we have a measure of health..
Some of us may feel healthier than others, but all of us have some measure of health.
that God gives us..
As long as we're here alive on this earth, we have some measure of health..
And we should be thankful for that..
We should be thankful for the beauty of family and of love and of relationships with other.
human beings that God has given us..
That is something that enriches our life in an unspeakable way and something that is just.
so amazing to experience..
And it's hard for us to describe..
I can't describe the beauty of that as I speak about it right now..

$^{521}$What the beauty of the relationships in my life, my marriage, my relationship to my wife,.
our relationship to our children and the love that we have for them and so many other friends.
and family members that we enjoy a relationship with..
That's a beautiful thing..
We can thank God for that..
We can thank God for the community of the church that we spoke about earlier, that community.
of His people that He's placed us in where we love each other and care for each other..
That's a beautiful thing..
And we can thank God for that..
We can thank God for the beauty of nature..
I am amazed by the beauty of nature and by the diversity and the creativity that is displayed.
and just all of the beautiful things God has made..
I love to see it..
I love the mountains..
I love the ocean and kind of everything in between..
And it just fills me with a sense of awe and gratefulness to God..
I love to look up in the night sky and see the stars and just be filled with that sense.
of awe and gratefulness to God..
God has put so much beauty in our world..
And He didn't have to..
He made it all gray and beige and really boring..
But He didn't..
He made it beautiful and amazing..
And I'm so grateful for that..
And we see so many beautiful things, even that God allows us as humans to make..
We see artistry and creativity that humans express themselves..
We can be grateful for that..
We can be grateful for the beauty that we see..
We can be grateful for culture and food and all the simple pleasures of life, sitting.
down with our family to enjoy a meal or to do something that is fun or listening to someone.
tell a really good story or something that just causes us to laugh watching a child play.
or something like that..
Those are all beautiful things that we should be grateful to God for..
He's poured so many blessings into our lives, even in the small moments of life..
And we often miss them because we're focused on what we don't have..
We're focused on what we want..
And we don't see all of the goodness that God has already put in our lives..
And we should certainly be grateful for God's love for us and the way He loves every one.
of us made in His image and the way He sent Jesus to die for us..
The salvation, the grace and mercy and forgiveness that God has extended to us, that because.

$^{561}$of Jesus we can be forgiven, we can be made clean..
All the mistakes we've made can be forgotten and put behind us..
Wow, that is something to praise God for..
That's something to be thankful for..
We can praise God for the rest of our lives and into eternity just for that alone..
There is so much to be thankful for in God's love for us and in the salvation that He's.
extended to us..
We can be thankful for the nearness of God, as we spoke about earlier, in those trials.
and those difficulties that God is always present with us, even when life is challenging..
And one of the things we can be thankful for is God's truth..
You know, we just spent a year studying through the Sermon on the Mount and the values of.
the kingdom..
I am so grateful for the values of the kingdom..
I'm so grateful that God revealed His truth to us..
Without it we would be lost and we'd be fumbling around in the dark and hurting ourselves and.
hurting everyone else..
I'm so grateful for God's truth and that He's expressed it to us and He's told us how we.
can live and please Him and be in right relationship with Him and right relationship with others..
I am so grateful for that..
The point is, there are so many things that all of us can be grateful for..
Even if in the moment things don't feel good and we think, "I don't know how much to be.
grateful for," we do..
And we've got to learn to begin to live out gratefulness and to practice it in our lives..
This is a great start to begin by placing our focus on God, placing our focus on the.
work that He's doing in our lives and all the good things that He's doing, and being.
intentional about celebrating God's work and God's goodness in our lives, stopping our.
lives, taking a break from the business just to reflect and just to praise and to thank.
God for who He is..
Let's pray..
God, I thank You right now for Your goodness and Your greatness in our lives, God..
You are so great and You're so powerful and so amazing..
And we see that so clearly when we look at the world and we see the amazing beauty of.
nature, but we see it so clearly when we look at each other and see the amazing intricacy.
of the human life, not just of our human body, but of the human soul and our emotions and.
all of the depth and the layers of that, God..
It's mind-blowing..
That's such a beautiful thing, God..
You are so great and so powerful to make all of that..
And yet You are also so good and You love us and You've extended us forgiveness and.
salvation and You've sent Your Son incarnated into this world so that we could have relationship.

$^{601}$with You and You've taken on our flesh and gotten on our level and come into our world.
so that we can know You because You are so good, God..
God, we praise You for that..
We lift You up..
We thank You, God..
We haven't done anything..
It's all You, God..
And there is no one like You..
There is no one like You in all of the universe, God..
And we just want to profess and proclaim that today..
And God, I pray that You'd help us to begin to reflect on and understand that and to think.
about it every day, to not take Your goodness for granted, but every day to be grateful.
and to be thankful and to intentionally celebrate how good You are, God..
In Jesus' name, amen..
[BLANK AUDIO].
\newpage



\section{}
\label{sec:hd_nz1Oh6bM}
\textbf{考古探秘 (耶路巴力銘文) - 陳崇基牧師 (2021.07.14)}
\newline
\newline
連結: \href{https://youtube.com/watch?v=hd-nz1Oh6bM}{\texttt{https://youtube.com/watch?v=hd-nz1Oh6bM}} ~~~~ 語音日期: 2021-07-17
\newline
\newline
\hyperref[sec:p69k7rhhWPQ]{\small{< < < PREV SERMON < < <}}
~
\hyperref[sec:index_chronic]{\small{[返順時目]}}
~
\hyperref[sec:index_scriptual]{\small{[返順卷目]}}
~
\hyperref[sec:yCNlGPEoTik]{\small{> > > NEXT SERMON > > >}}
\newline
\newline
$^{1}$又到聖經考古探秘的時間了.
今次跟大家談這個名文叫做耶路巴力名文.
Jerubbaal inscription.
找到的地方叫做Kubit al-Rai.
叫做阿拉圖有.
這個inscription.
你看到是寫在兩塊小的圖片上.
有一些字母.
這些字母經過鑑定之後.
我們知道是屬於公元前1200至1000年之間.
1200至1000年的BCE.
這段時間200年.
就是聖經考古時期的鐵器時代1.
Iron I的時候.
同時也是舊約聖經歷史的士司紀的時代.
你記得以色列人初初進入迦南的時候.
好像一盤散沙一樣.
慢慢他們在社會結構上邁向一個新的時代.
就是一個王朝的時代.
所以你看士司紀就是說他們很鬆散.
沒有王在他們中間.
各人行自己看為正的事.
然後到了薩姆爾紀.
你就找到開始有王朝出現了.
所以有薩姆爾,蘇羅,大衛等等的故事.
這個名文上面的字母.
今天我們叫它做原始的迦南文.
Protocanonite.
希伯來文.
今天我們看舊約希伯來聖經的時候.
那些Hebrew.
我們叫它做classical Hebrew.
公元前大概一千年左右開始演變出來的.
之前就是迦南的文字.
因為其實希伯來文這些字母.
本身就是一些迦南文裡面的字母.
這個名文在過去一個星期.
在不同的報章.
不同的社交媒體都有大肆報導.
說這個是新的發現.

$^{41}$其實在2019年已經出土了.
只不過真正的同行評審的論文.
研究的結果在這個星期才出來.
發現的地方當然就是Cubic Arai.
這個地方你看到圖片上看到.
就是阿拉圖丘的上空.
下面看到Area A和B.
你看到這個圖片的方向.
照片的方向.
其實是看著西面的.
如果天氣好的話會看到地中海.
所以比較近海平.
另外反過來看.
向東看的時候.
你看到遠遠是一個山脊.
是一個中央山區.
就是猶大山區向東面.
而他所身處的地方.
這些山都是比較矮小一點的.
這個地方在聖經裡面曾經提過.
叫做低山地.
或者是西山腳.
我們看看地圖.
如何給我們一些方向線.
首先我們去到這個迦林地.
然後我們放大到這個阿拉圖丘.
阿拉圖丘就在這裡.
你看到在這一區.
就是所謂西山腳.
低山地.
我這裡畫了綠色的地方.
就是沿海平原下非利士地的地方.
在非利士地這裡有五大最重要的城鎮.
就是從這裡左邊的加薩.
然後上面有艾克隆.
有艾克隆和加德.
據印哥利亞就是在加德出來的.
中間這裡就是西山地.
西山腳或者是低山地.
在東邊剛才那張相片看到的.

$^{81}$就是中央山脊.
比低山地更高.
上面你看到耶路撒冷,希伯來等等.
這個地方其實就是在聖經的歷史裡面.
西山地是一個緩衝區.
兩個主要的地帶.
沿海平原.
就是非利士人聚居的地方.
以及山地以色列人初入迦林的時候.
聚居的地方.
中間很多時候就是一個 transition zone.
一個過渡的地帶.
或者一個 buffer zone.
一個緩衝區.
所以這個 cubic array 這個地方.
就是在西山地那裡.
其實這個土優.
在十九世紀的時候.
英國來的聖地的探索家.
其實都已經知道了.
這裡是有考古的遺跡.
只不過到了2015年的時候.
才開始正式發掘.
為什麼考古團隊會進行發掘呢?.
因為他們在這個土優的面上.
面頭那一層.
他們發現到一些圖像.
是跟附近已經發掘了的土優.
是非常相似的.
是同一個時期.
就是 Iron I 的時期.
再看回 Google Earth 你就知道了.
阿拉土優 Array 就在這裡.
其他的土優在 Iron I 找到冥文.
或者找到一些遺跡.
就是譬如薩埃這個地方.
上面的 Kiyava.
一個基亞法的地方.
阿塞卡這個地方.
都是同樣在這個時期.

$^{121}$是有一些遺跡的.
這個基亞法的地方.
這裡就是伊拉谷.
Eli Valley.
大衛和哥利亞對壘的地方.
打仗的地方.
就是在這個地方.
所以在緩衝區裡面.
很多時候會有衝突.
但是無論如何.
在地理上.
Array 就是在緩衝區的西面.
在這個 Array 土優上面.
找到這幾年發掘出來的文物.
我們看到是一個迦南人的成真.
因為這裡找到一些迦南的實物文化.
這些文化都告訴我們.
那裡聚居的人.
是一些迦南人.
譬如這裡有一個例子.
就是找到一些迦南風俗裡面有的.
所謂 Bolambo 碗燈碗的裝置.
這些裝置很多時候都是在房屋的地基那裡.
賣在那裡的.
什麼叫碗燈碗呢?.
很簡單就是.
上面有一個碗蓋著.
中間有一盞燈.
右燈.
然後下面有另外一盞碗.
你看到左邊這幅圖就是.
看到上面蓋著的碗.
然後將碗打開.
下面就有一個油燈.
然後最底就有另外一個碗.
這些可能是保佑.
或者是辟邪用的.
一些迦南宗教的實物文化.
因為這個阿拉土優.
Cubic Array.

$^{161}$它是靠近非利士地.
所以也找到非利士人的雙色彩圖.
這些就是一些很著名很明顯的.
受到非利士人文化所影響的遺跡.
這些 Bicrom.
有兩種顏色的雙色彩圖.
那麼這個阿拉土優到底是怎麼找到那個命運的呢?.
原來就是考古人員去到發掘這個土優的時候.
他們就開始去到發掘一個地下坑.
叫做Silo.
下面就放著不同的瓷器.
不同的瓷缸.
來裝飾物或者飲品等等.
這個Silo.
我們看到四周圍有石塊來鞏固.
坑裡面有好幾個這些瓷缸.
這些瓷缸已經打爛了.
只不過是考古學家再拼起來.
在這些碎片當中就發現了這些冥文.
這些冥文是用一些古代的迦南字體.
字母來寫的.
或者是希伯來文.
前身.
很相似的.
這些幾個冥文.
如果你把它們比較成同時期的字母的話.
你就可以拼第一個.
可能就是肉.
第二個是鼻.
第三個是rage.
第四個是iron.
第五個是lamen.
所以用英文寫出來的時候.
拼起來就是ybr一撇和lamen是L.
你會看到有五個字母可以看得很清楚.
或者是重構出來的.
這個同行評審的論文.
就是在這個星期2021年7月12日星期一刊登出來的.
刊登的期刊叫做Jerusalem Journal of Archaeology.
是一個學術上的同行評審的論文的期刊.

$^{201}$這份論文有好幾個考古學家的名字.
其中有法國的主要的director是Gavrilo Gavrilo Gavrilo Gavrilo.
另外第一個主要的作者主要研究這個名文的就是Christopher Wollstone.
Christopher Wollstone不單是在美國.
在世界都是首位一指.
專門研究聖經考古時期的名文的專家.
所以如果你想看這個同行評審的學刊的文章的話.
在我們的description下面有個link.
你點進去就可以看到一個公開的學刊.
你就知道這個名文中間的細節是怎樣的.
怎樣去推敲到這些字體是從哪個時代出來的呢.
其實是將這些在考古發現到的名文.
將它擺列出來的時候就可以做一些比對.
我這裡有一個圖畫一個列表.
就告訴你第一行由幼稚組體.
是現代希伯來文的字母寫法.
今天我們拿出一本列印出來的希伯來文聖經.
我們學習希伯來文的時候就是學習這些字體.
Aleph, Be, Gimel, Tala, Hiv, Av, San, Heter, Yod.
然後有22個這些希伯來文字.
這個是現代的字體.
以前的字體是怎樣的呢.
不同的時代就有不同的寫法.
譬如將Gazette Calendar.
基式閱歷大概是公元前十世紀.
即是主前大概九百多至一千年的石板.
上面有一首詩我們叫它基式閱歷.
裡面有的字母我們擺出來的時候.
我們就知道原來Aleph.
以前公元前一千年九百年的時候是這樣寫的.
然後再過一百年就到米沙石碑.
這個是在約旦河東摩押旺留下的石碑.
裡面有的字母幾乎齊全.
22個對稱的希伯來文或者摩押文.
你看到Aleph的寫法又不同.
這個是穿透.
或者Mem的寫法就不同一百年前的寫法.
它的撇是彎彎的.
一直推敲下去.
有一個地方叫做Kentile Ajru.

$^{241}$在南邊的西乃半島東北邊的一個Way Station.
一個客旅所停的站那裡.
就找到一些石缸.
裡面有一些名文.
我們看到那些字體.
就不同之前米沙石碑或者基式閱歷的字體.
繼續下去我們看到Ceylon Inscription.
就是西羅亞名文.
這個是希西加的時候.
是七百年的BCE.
701 BCE.
希西加水道留下來的名文.
你看到Aleph和Beth的彎位都有一些不同.
所以你這樣推敲的時候.
你就發現一些Epigraph.
就是古名文的學家.
如何將一個名文找出它的時代.
譬如耶路巴力的名文.
你看到有Beth.
你要去比對一下.
到底Beth在這段時間.
還在基式閱歷之前的百多二百年.
到底是怎樣寫的呢?.
這些不同的Rage.
到底是怎樣寫的呢?.
在這個更加複雜的圖表.
你就更加能夠看到.
你可以自己去拼一下.
如何拼出耶路巴力的名文.
這個圖表是如何看的呢?.
由左邊開始看.
左邊這一欄就是今天希伯來文的字體.
Aleph,Beth,Gabriel,Dalek,Hepha,Sinai,Heth,Jerusalem.
然後繼續下去有22個希伯來文.
然後你會看到有BC15 BC13 BC13 13.
直到去什麼時候?.
直到去BC2.
BC2 Isaiah Scroll.
即是死海古卷的時間.
你看到只是看Aleph.

$^{281}$都有不同的演進和寫法.
如果我們說耶路巴力的名文.
是屬於公元前1200至1000年的時候.
我們所比對的就是.
Tel-Said裡面有的名文.
他的Aleph.
和Isbeth-Sattar的名文.
因為這裡是公元前13世紀.
1200年左右.
Tel-Said裡面是怎樣寫Aleph的呢?.
原來是好像A字一樣.
打橫的.
有時左有時右.
我圈上了.
框上了.
Beth是怎樣寫的?.
一個圓圈這樣卷下來.
Yod有不同的寫法.
可以這樣插下去.
打橫插下去.
或者這樣打直.
然後Lamed就是一個.
好像Beth一樣的.
但方向有時是不同的.
下面的Rage就像一支棒.
後期就像一支棋.
如果這樣看的話.
就像一個古命文學家這樣看.
這裡大致重構出來的就是Yod.
這個就是Rage.
第三個是一個凹下去.
中間有一個斜劃的Beth.
就是B.
第四個就是一個圓圈.
一點.
這個就是Iron.
第五個就是一個這樣的圈.
剛才你記得.
看Lamed的時候.
就是一個這樣的圈.

$^{321}$這個就是L.
所以你這樣拼出來的時候.
YRB一撇.
L的話.
Yod Rage Beth Iron Lamed.
這樣拼出來加上響音.
就是Jerubal.
耶路巴力.
為什麼耶路巴力的名字.
這麼受人著重呢.
或者這麼多人留意呢.
因為在聖經裡面.
確實有出現過這個名號.
或者這個名字.
耶路巴力什麼時候出現的呢.
在《事司紀》六章三十二節等等的經文.
其實在整個基典的故事裡面.
六至八章九章.
經常出現耶路巴力這個字.
這個名是什麼意思呢.
六章三十二節說.
當基典受到神的呼召.
他拆毀了他父家的巴力祭壇之後.
人們就很生氣.
然後就將他改名.
或者給他一個綽號.
就叫做耶路巴力.
當日基典被稱為耶路巴力.
意思說他拆毀了巴力的祭壇.
讓巴力與他爭辯吧.
所以他就叫做耶路巴力.
他跟巴力作對.
巴力這個風雨神.
或者假神會不會真的跟他作對呢等等.
所以你看到這個名字.
真的出現在聖經.
並且還出現在聖經的《事司紀》.
考古找到的時代.
是Iron 1 1200-1000 B.C..
在《事司紀》記載到耶路巴力.

$^{361}$這種名字配上巴力這個字.
又是在這一段歷史的時間.
在《事司紀》的基典故事.
出現了很多次耶路巴力的名字.
你看到名字的形式.
在這個時候是特別受歡迎的.
比如什麼巴力.
什麼Isht-Bael.
這種名字在撒母耳基.
蘇羅時代.
和《事司紀》時代出現比較多.
後期以色列人入主了迦南.
在皇朝時期開始了.
這些名字就沒有了這種特點.
另外一種特點就出現了.
就不再是用Bael巴力這個字.
當然我們想起巴力.
我們不一定要想起.
迦南宗教的風雨神Bael巴力神.
巴力的意思其實只是主人Master的意思.
所以耶路巴力這個字.
可以說是爭辯主人等等的意思.
不一定和雨神巴力有關.
但無論如何.
後期以色列人的名字.
你會發現不再有Bael的字加上去.
加上去就是Yah了.
Yahweh.
Yah-Hu.
就是耶和華.
或者短寫叫做Yah.
比如以利亞.
El-Yah-Hu.
My God is Yahweh.
我的神El-Yah-Hu.
就是耶和華了.
或者希西加.
等等這些字就開始有不同的形式出現.
所以單單看名字.
你都可以猜到.

$^{401}$知道它是屬於哪個時代.
所以找到耶路巴力名文.
就和聖經的歷史有一些吻合的地方.
當然耶路巴力肯定不會是基典.
但是如果你這樣去比對的時候.
考古在這段時間找到耶路巴力這個名字.
也在士師記初入迦南的時候.
這段歷史的時期.
找到耶路巴力的記載.
基典改了這個名字.
你就會看到聖經的經文.
很多時候其實是保留了一些很真實的歷史記憶.
\newpage



\section{}
\label{sec:yCNlGPEoTik}
\textbf{舊約網上行 - 上帝的足跡: EP13 所多瑪 - 審判與恩典}
\newline
\newline
連結: \href{https://youtube.com/watch?v=yCNlGPEoTik}{\texttt{https://youtube.com/watch?v=yCNlGPEoTik}} ~~~~ 語音日期: 2021-07-31
\newline
\newline
\hyperref[sec:hd_nz1Oh6bM]{\small{< < < PREV SERMON < < <}}
~
\hyperref[sec:index_chronic]{\small{[返順時目]}}
~
\hyperref[sec:index_scriptual]{\small{[返順卷目]}}
~
\hyperref[sec:rJBi5QVSDGc]{\small{> > > NEXT SERMON > > >}}
\newline
\newline
$^{1}$我們來到舊日網上的第十三講.
之前我們講到創世的故事是發展成怎樣呢.
亞伯拉罕在希伯倫的萬里長殊身的三位使者來到.
亞伯拉罕款待他們.
神就向撒拉報棄.
一年之後.
以撒就會出生.
之後發生了什麼事呢.
我們用少許時間解釋這段很長的經文.
然後看看如何去形容.
亞伯拉罕款待了使者之後.
三人就從那裡起身.
亞伯拉罕與他們同行.
要送他們一程.
所多瑪,俄摩拉等城市.
有人認為是在死海的北面.
但大部分的學者都認為是死海南面的東岸.
雖然現在還不能確認.
或者這些古城滅亡之後.
就是埋在死海之下.
但按照聖經的地理.
是在這一帶.
無論如何.
聖經的地理細節是合理的.
由希伯倫的山地萬里象樹往東.
一路下去死海的低谷.
過了死海.
對面就是所多瑪和俄摩拉.
為什麼神要去所多瑪,俄摩拉呢.
十八章二十節.
又話說.
控告所多瑪,俄摩拉的聲音.
甚大他們的罪惡極重.
我現在要下去.
看看他們所行的.
是否象那聲聞於我的控告之聲.
如果不是的話我自己也會知道.
神去所多瑪查證.
查證一下是不是罪惡滿盈面臨審判.
之後亞伯拉罕聽見這個消息的時候.

$^{41}$他就和神有一個很長的對話.
第二十二節.
其中兩個人往那裡動身向所多瑪走去.
亞伯拉罕站在耶和華面前.
亞伯拉罕就上前來說.
神啊!你真的要將義人和惡人一起除滅嗎?.
假如城中有五十個義人.
你還真要除滅那地方.
不為城裡面五十個義人而赦免那地方嗎?.
你絕不會作這樣的事.
把義人與惡人一同殺死.
把義人和惡人一同看待.
你絕不會這樣.
審判全地的主豈不會行公義嗎?.
神就回答他.
因為那五十個義人的話.
如果真的有五十個義人.
他就暫且不會審判所多瑪和摩拉.
然後亞伯拉罕繼續對話當中.
就好像講價一樣.
由五十位義人到四十位,三十位,二十位.
最後就十八章三十二節.
亞伯拉罕就說.
請主不要動腦.
用我問你最後一次.
假如在那裡找到十個義人.
那又如何?.
神就回答他.
因為這十個的緣故.
我不會毀滅那城.
然後十九章就很詳細地講述.
兩位神的使者到達所多瑪.
發生了什麼事情.
十九章第一節.
黃昏的時候.
那兩位使者到了所多瑪.
當時羅德坐在所多瑪的城門口.
羅德一看見他們.
就起來迎接他們.
俯伏在地.

$^{81}$羅德再三地來拯救他們.
這些使者才肯到他的家.
羅德就為他擺上筵席.
燒了母烤餅.
然後他們就吃了.
記得這些是款待客女.
Hospitality.
近東十分重要的社會價值.
這些是道德的指標.
這裡已經有一點暗示.
整座的所多瑪城.
只有羅德一個願意款待客女.
其他人到底去了哪裡.
十九章第四節就出現了.
他們沒有躺下睡覺.
那城裡面的人.
所多瑪的男人.
無論年輕的還是老的.
從各處來到去.
圍住羅德的房子.
他們就呼喚羅德.
對他們說.
今晚到你這裡的人.
將他們帶出來.
交出來.
我們要與他們同房.
你想像當時的情況.
可能非常緊急.
以及非常暴力.
第六節.
羅德就出來.
隨手把門關上.
到門口眾人那裡說.
我的弟兄們.
你不要這樣做我.
看啊.
我有兩個女兒.
沒有與男人同房.
讓我把他們帶出來.
你們高興怎樣待他就怎樣待他.

$^{121}$至於兩個客人.
來到我四下的時候.
你們不可向他們作甚麼.
我們現代人也好.
古代人也好.
我相信看到這個情節的時候.
都很難解釋.
羅德說.
你們想污辱我的客人.
你們千萬不可以這樣做.
不如我把我兩個女兒交出來.
這些是令人很費解的情節.
或者很費解的對答.
這裡當然預設了.
中東的款待文化.
人的天職.
主家要做些甚麼呢.
就是保障客人的安危.
或者.
可能這個情勢真的很急.
很暴躁的一個時間.
連羅德也不知道怎樣去解圍.
或者.
經文又再次想強調一件事.
這位愛城市.
愛索多瑪.
愛世界的羅德.
既然住在罪惡之都.
現在就面臨一些道德的困境.
無論如何.
這個故事繼續.
第九節.
眾人卻說滾開.
又說有一個這類寄居的人.
就指羅德.
竟然敢做我們的審判官.
現在我們要害你.
比害他們更加厲害.
眾人就向羅德猛烈的擠湧.
向前想打破他的房門.

$^{161}$這兩個使者伸手出來.
將羅德拉回到屋內.
並且將門關上.
然後就擊打門外的人.
使他們無論怒笑.
眼目昏花.
以至找不到門口.
這裡有些很緊急的時候.
帶天使插手進去.
以致有種拖延.
最後天使真的要去審判罪惡之都.
第十二節.
那二人對羅德說.
你這類還活有什麼人.
還有沒有你的女婿女兒.
和城裡面所屬你的人.
你們要帶他走.
因為我們快要毀滅這地方了.
控告他們的聖人在耶和華面前實在太大.
所以耶和華派我們去毀滅這地方.
之後整個晚上有很多拖延.
十九章十五節.
到天快亮的時候.
天空降落.
搞了一整晚.
兩位使者就催促羅德.
起來帶著你的妻子.
和這裡兩個女兒出去.
免得你因為這個城的罪孽.
你們一同滅亡.
羅德還猶豫不決.
兩個人因為耶和華的恩典.
寬容羅德.
就拉著羅德的手.
他的妻子的手.
和兩個女兒的手.
將他們帶出去城.
羅德到達.
所以的時候.
太陽已經在地面上升上來了.

$^{201}$二十四節.
當時耶和華就使得流黃於火.
從天上耶和華拿了江宇.
所多瑪和俄姆拉.
把那城全個平原.
城中所有的居民.
連地上生長的都毀滅了.
二十六節.
羅德的妻子在羅德的後面.
向後一望就變成炎柱.
向後一望.
可能是一些很生動的妙術.
意思是說.
羅德的妻子是凝望.
是很眷戀想留在罪惡之中.
之後她自己也滅亡.
然後故事的發展.
大家可能記得.
就是羅德和他兩個的女兒.
逃到所謂的山洞上.
然後在那裡避難.
然後亂倫就發生了.
就成為了摩訶和阿們人的祖先.
剛才所交代的二十四小時多的故事.
是非常緊湊.
但是也是一個非常令人吃驚.
最後的結局也是非常羞恥的.
罪惡之都所多瑪索當.
這個故事是家傳戶曉的.
因為英文sodomy.
中文可以譯作群間.
等等的譯法.
這個字.
sodomy.
就是由這個地名.
Sodom所多瑪而來.
第一方面.
在我們今時今日的文化裡面.
看這段經文.
一定是令這個世代的人很難接受的.

$^{241}$因為經文似乎很反同.
特別是對這個罪行.
作出很嚴厲的批判譴責.
甚至是神的審判.
在這個時代.
我們可以了解.
為什麼看這段經文會有不同的感覺.
因為同性戀.
可能已經被大部分的人所接受了.
大叔的愛見怪不怪.
所以也有同志識經出來.
所謂同志識經就是對這些經文.
再作一些新的詮釋.
或者會指出.
這段經文所針對的.
不是我們今時今日所說的同性行為.
而是一些欺壓的情況.
又或者十九章第五節.
那裡有一個同房動詞.
希伯來文是 yada.
有些同志識經會認為.
這裡所說的不是性行為.
而是 investigation.
而是調查的意思.
就是說 bring them out.
we may know them.
we may investigate them.
就是把他們交出來.
我們盤問一下他們.
調查一下是不是間諜來到我們城裡面.
想做些什麼.
跟性行為是無關的.
這個就是從希伯來文的文法入手.
但是如果你看上下文的話.
其實也有提示的.
這件事情是跟性一定有關係的.
因為羅德想解圍的時候.
他跟這些所多馬人說.
你們辱火攻心的話.
我就把我的女兒交出來.

$^{281}$還有最大的問題.
也是文法的問題.
因為希伯來文.
yada.
know.
之後每每有一個 personal object.
就是帶著一個人的時候.
就是 know someone 的時候.
就不是純粹調查某人.
或者知道某人.
希伯來文由始至終.
說到 know someone.
知道某人的時候.
其實也是跟性有關的.
所以這種色經我們不能夠說得嚴謹.
不如直接說.
我們希望去屈就現在的大的文化.
大的環境.
我們想退一步來看看.
聖經裡面.
教一些東西.
很多時候是不是人不能夠接受的.
或者再退一步再想一想.
如果聖經.
神藉著聖經來對我們有教導.
如果每一樣東西都合我們口味的話.
這個神還是不是主呢?.
神是聖傑是公義.
祂會私行審判.
人是需要悔改.
要認罪的.
這些這麼簡單的福音道理.
也是人不能夠自然接受的.
我就情願.
聖經裡面所說的.
神想教導的.
很多時候是不合我的口味的.
總比聖經每一樣東西都要切合我的口味好.
我才相信.
如果是這樣的話.

$^{321}$神就不再是主.
只不過是偶像.
是一種偽基督教.
剛剛這個星期.
2021年7月的時候.
報章上曾經有一些文章.
討論現代的宗教現象.
特別是美國.
美國的New Religion.
是一種Fake Christianity.
這個文章其實是想引用一個基督教的統計學家.
叫做George Barnard.
George Barnard經過調查之後.
他就形容今時今日的人對基督教信仰的理解.
是用三個英文字來總結出來.
Moralistic, Therapeutic, Deism.
就是道德治療的自然神論.
道德的意思是Moralistic.
就是說我相信耶穌.
其實就是我做一個好人就夠了.
治療性的Therapeutic.
其實意思是Feel good.
就是合我的心意.
神最重要的.
祂要做的就是要令我開心.
Deism, 自然神論.
就是神創造了這個世界.
祂不再參與在當中.
這些就囊括了今時今日很多時的趨勢.
就是人看基督教聖經的信仰.
只不過是一種治療性.
教你做一個很好的人.
一個Nice的人就足夠了.
這種現象或趨勢.
其實是有它的危機在.
按照George Barnard所說.
就是你將基督教信仰的生活.
慢慢地將它每一個部分打掉.
不再需要在教會為神.
不再需要去崇拜.

$^{361}$不再需要去宣教.
不再需要奉獻事奉.
作光作艷.
不需要再教養下一代去敬畏神.
讀經聽道.
只不過是繼續深化我們本身對信仰的了解.
一切心照.
一樣東西都不用信仰生活.
最後結果會怎樣.
就會甚麼都不打緊.
最後連基督也不需要.
基督在歷史十家空墳墓所做的事.
全部都不重要.
因為最後只不過是做一個Nice person.
去Feel good就足夠了.
這種是偽基督教.
是基督教外學.
但不是基督教本質的.
這種趨勢其實跟以前的.
可能上一代基要派的問題其實都一樣.
基要派說信耶穌就等上天堂.
現在就不用做任何事.
你想想那個動態.
那個原理其實都一樣.
就是將神信仰Fit入我自己的框架.
Fit入我的Box.
到最後我們出賣的可能就是福音.
就是基督.
引伸出另一個問題.
就是疫情繼續延長的時候.
就更加深化這種偽基.
甚麼我們都不需要.
最重要的是我的心照就行.
在這個報道裡面.
有一個畫家畫了一個illustration.
一條直立的彩虹的羽.
上面有一個少年.
然後其實是一座塔.
他這個作品叫Sun Barbie塔.
New Tower of Babel.

$^{401}$到底我們很多時候.
是否就是滑落了這種偽基.
不知不覺地.
其實我們就是在信偽基督教.
而失去了的就是耶穌基督.
不過當我們看這段經文的時候.
可能我們又要留意.
會不會有另一個偽基.
如果我們只是將這段經文當作反同的金句.
我們的感覺就是.
好了,上帝就是去剿滅這些人.
他們是死有餘辜的.
這些罪是最大最大的罪.
那麼這個偽基是甚麼呢.
就會令到我們失去了平衡.
就會自以為然.
其實自以為然又是很大的罪.
經文有甚麼是值得我們留意的.
當我們看過索多瑪的故事的時候.
我們當然不可以忘記一件事.
就是一個款待hospitality的故事.
就是經文裡面刻畫.
索多瑪是一個怎樣的社會.
一個怎樣的城市.
上文下理我們應該看得很清楚.
特別可能看第九節.
那個中人跟羅德說甚麼呢.
他說你滾開.
又說羅德是在這裡寄居的人.
竟然敢作我們的審判官.
現在我們要害你.
比害他們兩個客女更加厲害.
中人就極力往前推.
然後想打破他的房門.
索多瑪的罪不是單單淫蕩淫惡.
而是誤款待欺凌弱小.
特別是客女.
其他舊約的經文.
去到講到索多瑪的時候.
其實都是針對一些更加廣泛的論理問題.

$^{441}$譬如以賽亞書第一章.
這裡講到索多瑪的官長.
其實就是比喻作耶路撒冷的官長.
他說要聽人的話.
你們要使作自竭.
在我面前取消你們的惡行.
甚麼惡行.
就是你們要學習行善.
尋求公平公義.
去到解救受欺壓的.
被孤兒寡婦去辯屈去生.
以西結書.
另一本先知書十六章.
講到索多瑪也是講到耶路撒冷的罪惡.
就是他和他的眾女是心驕氣傲.
糧食飽足大享安逸.
沒有輔助困苦.
噴佛人的手.
到了兩約中間的一本書.
叫做希臘書Book of Jubilee.
第十六章講到索多瑪的罪.
就是不潔的罪.
就是逆性的淫亂.
然後猶大書新約.
講到索多瑪和摩拉周圍盛事的人.
他們是淫亂和逆性的情慾.
以至遭受永火的刑罰.
是後世人的尷尬.
聖經的傳統如果全部拿出來講.
索多瑪的罪就不只是淫亂.
而是心驕氣傲.
不肯幫助脆弱的人.
不肯為他們伸冤.
不肯私行公然.
其實對索多瑪的罪有一個更廣的理解.
在猶太人的拉比文學.
早就出現了.
猶太人的解經書.
有一套叫做Tosefta.
在Tosefta Sotah.

$^{481}$第三章講什麼.
那裡說索多瑪的官長.
因為糧食,銀子,黃金和寶石.
他們說都是出自我們的土地和城市.
所以我們不會容許其他人.
去分薄我們的貴寶.
他們說起來.
我們要去阻擋這些客人進城.
然後欺壓一些經過他們的人.
神就說什麼呢.
神就說我既然賜了你好處.
你居然忘記了我們城裡面的寄居者.
我必將你們從世上除滅.
講到索多瑪不單止是淫亂.
也是欺凌弱小.
去阻礙一些外來者.
去分薄他們的利潤.
分薄他們的財富.
人的傲慢.
神去審判.
這些教導其實什麼時候都不會受人歡迎.
但是我們要想清楚.
神不單止是慈愛.
神不單止是有恩典.
神也是聽見被屈者的呼喊.
如果人不悔改的話.
就必定會遭受上帝公義的審判.
你想想阿巴拉罕.
知不知道索多瑪的惡行嗎.
一定知道.
因為就在對岸.
可能以往有不少的客女.
都已經受盡索多瑪的凌辱.
所以上帝要下來查看.
並且下來施行審判.
神有沒有給索多瑪的人民機會呢.
我們不能夠只看這段經文.
就好像覺得神很兇.
一犯了什麼事.
沒有警告.

$^{521}$然後就滅了這個城.
整個的故事.
索多瑪什麼時候出現.
十三章已經出現.
就是阿巴蘭住在迦南地的時候.
羅德搬到了平原的城市索多瑪.
然後以他的帳篷直到索多瑪城.
第十三節.
在耶和華汗內.
索多瑪人是最大惡極.
當時審判還未到.
到第十四章你記得了.
就是北方的四王來攻打.
南方的五王.
索多瑪是其中一個.
就忽略了羅德和索多瑪的財物.
阿巴拉罕做什麼.
他就出兵來營救.
救了羅德和將所有的財寶帶回來.
上帝有給悔改的機會.
但是阿巴拉罕說什麼.
他說審判權力的主.
豈可不幸公義.
這段經文不單止告訴我們.
上帝是全然聖潔全然公義.
必定會私人審判.
也提醒我們作為神的子民.
所謂阿巴拉罕的子孫.
我們有什麼責任.
可能這才是故事最重要的地方.
我們回到故事的開始.
十八章第十六節.
三個使者起身.
他們向索多瑪觀寒.
阿巴拉罕和他們同行要送他們一程.
然後耶和華說什麼.
耶和華說我要作一件事.
這件事當然是審判索多瑪所有的城市.
豈可以瞞住阿巴拉罕.
阿巴拉罕必成為一個大國.

$^{561}$他要使萬國能夠得福.
我揀選了阿巴拉罕.
是吩咐他的子孫他的家屬.
要遵照我的道.
秉行公義.
好叫耶和華應許阿巴拉罕的話.
可以兌現.
能夠去應許阿巴拉罕的話是什麼意思.
當然就是要萬國得福.
要使萬國的人能夠悔改歸向主.
在一個罪惡滔天的世代裡面.
神的子民不是無助的.
我們有很多事情要做.
就是我們世世代代要秉行公義.
要遵照主話.
使話主的福音能夠實現出來.
因著神的子民.
因著你和我.
萬國能夠有機會去相信福音.
有恩典去悔改.
你留意索多瑪被滅亡被審判的經文.
前後有什麼交代.
很仔細交代什麼.
就是神和阿巴拉罕的對話.
所謂阿巴拉罕的祈禱.
阿巴拉罕對神說什麼.
他說你不會將義人和惡人一起去看待.
一起去剿滅吧.
因為你是行公義的神.
然後他就好像講價一樣.
他說如果這個城裡面有50個義人的話.
你會不會滅這個城呢.
又說當然不會.
然後繼續說45呢.
30呢.
20呢.
最後發生什麼事.
32節.
阿巴拉罕說.
請主不要動怒.

$^{601}$我再問你一次.
最後一次.
假如這個罪惡的城市裡面找到10個義人.
你會不會剿滅這個城市.
又說不會.
因著這10個的緣故.
不會毀滅那城.
阿巴拉罕是為了惡人.
為了索多瑪來求恩典悔改的機會.
相信他當時在紀念羅德.
在猶太人的傳統裡面.
這是最偉大的祈禱.
因為這個祈禱是一個典範 paradigm(平行).
以後的先知都是按照這個 paradigm(平行)來祈禱.
最明顯的就是摩西.
在抗耶.
以色列人硬著頸寒.
要去拜金牛.
去叛變上帝.
摩西說什麼.
你不要剿滅你的百姓.
在這裡求恩典.
求悔改的機會.
其實阿巴拉罕的這個祈禱.
硬性來說不算一個祈禱.
其實是想在這個對話裡面.
表達一個原理.
就是神的子民在世.
in the world(世界).
我們有的 impact(影響).
我們有的影響.
在城中的公義小眾.
righteous minority(正義少數).
是因著這些公義小眾的存在.
神會延遲他的審判.
給人悔改的機會.
疫情一年多.
改變了世界.
也引發了很多深層的個人.
社會裡面一些結構的罪惡.

$^{641}$一些邪惡的欺壓的事.
我們有時候真的覺得很無助.
有時候覺得我們可以做些什麼.
但這段經文告訴我們.
在上帝的眼中.
我們有很重要的角色.
原來在一個道德淪亡崩壞的城市裡面.
因著一群義者.
因著一群小眾.
righteous minority(正義少數).
就能夠有恩典.
有時機.
讓罪人悔改.
使萬國得福.
這就是我們作為阿拉罕子孫.
有的責任.
結束的時候.
我們看彼得後書第三章.
他說親愛的弟兄姊妹.
有一件事你們不可忘記.
就是主看一日如千年.
千年如一日.
主所應許的尚未成就.
有人以為他擔言.
其實不是擔言.
但是寬容你們.
不願一人沉淪.
乃願人人每一個罪人能夠悔改.
求主幫助我們成為.
義者小眾.
righteous minority(正義少數).
在地上.
in the world.
in our city.
我們為著上帝的福音而禱告.
為著上帝的國度而禱告.
求主藉著我們的禱告.
藉著我們的工作.
我們的使命.
我們的宣教.

$^{681}$使人能夠悔改.
得著救恩的福.
我們同心協力地祈禱.
因為你將我們安置在不同的城市.
在這些罪惡滔天的地方.
很多時候我們會覺得很無助.
但你卻在自己的話語裡面.
提醒我們.
我們有一個很重要的責任.
就是為了罪人的悔改而存在.
讓他們有時機來存在.
讓他們能夠聽到福音.
主我們為了他們而禱告.
為了未認識主的.
正在步向沉淪的不義的人.
或沒有盼望的人.
主我們為他們而禱告.
無論他們是我們家人.
在我們社會裡面的人.
是在上掌權的.
平民百姓的.
主你聽我們禱告.
願你給他們機會能夠悔改歸向你.
願萬人來悔改歸向你.
以致你的名得著榮耀.
我們這樣祈禱奉主耶穌的名字.
\newpage



\section{}
\label{sec:rJBi5QVSDGc}
\textbf{舊約網上行 - 上帝的足跡: EP14 守約者}
\newline
\newline
連結: \href{https://youtube.com/watch?v=rJBi5QVSDGc}{\texttt{https://youtube.com/watch?v=rJBi5QVSDGc}} ~~~~ 語音日期: 2021-08-09
\newline
\newline
\hyperref[sec:yCNlGPEoTik]{\small{< < < PREV SERMON < < <}}
~
\hyperref[sec:index_chronic]{\small{[返順時目]}}
~
\hyperref[sec:index_scriptual]{\small{[返順卷目]}}
~
\hyperref[sec:PyoQq9QIPJA]{\small{> > > NEXT SERMON > > >}}
\newline
\newline
$^{1}$我們來到《九一網》上的第十四講.
上一次我們講到《創世記》第十九章.
上帝去審判罪惡之聖所多瑪俄摩拉.
表明上帝是全地的主人.
現在來到《創世記》第二十至二十一章.
這個差不多去到阿巴拉罕故事的尾聲.
通常一個作者會在故事或者劇集的尾聲.
會交代不同的情節.
其實這兩章就是講到三個不同的事蹟.
第一個故事.
就是講到阿巴拉罕在外邦人的地方.
叫做基拉爾發生一些事情.
然後二十一章上半章.
就講到以撒出生.
以實瑪利離開.
然後二十一章下半章.
就講到阿巴拉罕在另一個地方.
叫別士巴.
那裡也發生了不同的事情.
我們先處理中間這段線.
就是處理以撒和以實瑪利的問題.
其實是一個繼承者的問題.
誰能夠繼承阿巴拉罕所立的大國呢?.
這個是故事的尾段.
之前你發現.
有很多故事的發展.
有很多的拖延.
有很多的等待.
有很多的障礙.
有很多的問題出現.
最後上帝就克服所有的困難.
所有的路障.
以至最後以撒就出生.
第八天受了割禮之後.
創世記二十一章第八節.
他說以撒漸漸長大了.
斷了奶之後.
斷奶過一天.
阿巴拉罕就大排筵席了.
撒拉看到埃及人下格.

$^{41}$被阿巴拉罕所生的兒子.
就是以實瑪利在欺欺.
就是去欺負他的時候.
就對阿巴拉罕說.
他說將這個婢女和她的兒子趕走.
因為這個婢女的兒子.
不能夠跟我的兒子以撒一同承受家業的.
阿巴拉罕就因為這個兒子以實瑪利.
非常煩惱.
神就對阿巴拉罕說什麼.
他說不要因為這個兒子和你的婢女而煩惱.
撒拉說什麼你儘管聽從.
因為以撒生的.
才可以稱為你的後裔.
至於這個婢女所生的兒子以實瑪利.
我也會使他成為一國.
因為他也是你的後裔.
我們看這個故事的時候.
很難禁會有一種感覺.
就是覺得很無情很無辜.
下格以實瑪利.
如果你先算誰錯的話.
一定不會是他們.
如果看回故事的發展的開頭.
是阿巴拉罕和撒拉去到獻祭.
才會使下格生了以實瑪利.
現在到故事的尾聲了.
根本就不關他們母子倆的事.
為什麼又要將他們趕走呢?.
好像很慘很慘.
第一我們需要留意的.
看這個故事的時候.
就是我們發現這個故事是寫給當時古近東的人.
他們明白當時的文化家族的產業是怎麼運作的.
可能對第一代的讀者.
古代的讀者.
這個是還可以接受的.
無論如何.
如果你看這個故事.
以實和以實瑪利.

$^{81}$如果單從一個分家產的角度來看.
沒錯的.
其實可以兩個兒子都一起分享的.
最多以實分兩份.
以實瑪利分少一點.
況且阿巴拉罕也不是窮的人.
他雖然不至於是一個帝國帝王.
但是家產也是不少的.
但是如果你留意.
我們在看上帝的國度.
上帝的預約一直向前邁進的時候.
這個就不單單是說分土地.
不單單是說分財產的問題.
是說到一個國度裡面.
有一個王位.
這個王能夠從萬國之父.
以色列阿巴拉罕的家族出現.
如果是這樣的話.
一國之君就只有一個.
所以只有一個繼承人.
所以為什麼創世記一直在說.
上帝所創立的天國.
要使萬國能夠得福.
這個天國的繼承人只有一位.
況且上帝的選擇原則.
往往是出人意外的.
人們以為第一個出生的長子.
就一定會有這個位分.
但是創世記在族長的故事裡面.
一直在說的所揀選的.
都不是人的價值.
不是按照人的原則.
譬如這裡神就立以撒為繼承人.
長子以撒瑪利就沒有份了.
之後的故事我們會記得.
雅各和以素.
以素是長子但是上帝沒有揀選.
他是揀選雅各.
然後雅各的十二種子.
揀選了第四個兒子.

$^{121}$是猶大.
他才能夠有這個位分.
君王的繼承人.
我們怎麼看這個故事的時候.
你就說.
都是覺得很慘.
為什麼要將他們踢出家門.
我們來看看故事的發展.
作者怎麼交代.
二十一章第十四節.
亞伯拉罕就動身了.
清早起來的時候他就拿餅拿水.
然後給了夏甲.
然後把孩子交給他.
就打發他走.
夏甲走的時候.
就去到別士巴的曠野飄蕩.
就是在那裡迷失了.
皮袋裡面沒有水.
夏甲就將這個孩子撇在小樹底下.
這個時候以撒瑪利.
已經是一個teenager.
一個年輕人.
但是我們可以想像.
在曠野迷路的時候.
可能他會中暑或者會生病等等.
所以去到這個時候.
母子倆就走投無路.
然後到十七節.
神居然又出現了.
神就聽見童子的哭聲.
他的使者就從天上呼叫夏甲.
對他說夏甲你不需要怕.
因為孩子在那裡所發出的聲音.
神已經聽見了.
扶起你的孩子盡力去保護他.
因為我要洗他.
也成為大國.
神就開了夏甲的眼睛.
就看到一口水井.

$^{161}$他就裝滿水然後給孩子.
神就與孩子同在.
他漸漸長大的時候.
住在曠野就成為弓箭手.
住在巴林的曠野.
他的母親在埃及.
就為他奪了一個妻子.
夏甲和以撒瑪利在曠野迷路.
卻是得救.
上帝有一個恩典.
普遍的恩典.
使得這個母子能夠生存.
並且將來能夠成為大國.
在他們山窮水盡.
死路一條的時候.
上帝就出現.
來給他們一個拯救.
人生很多時候不是這樣.
就是看到上帝的恩典.
將我們救活.
其實這個也是創世記本身有的神學.
一直有一個所謂國邦的觀念.
一個政治神學可以叫做.
就是神是全地的主.
所以上帝掌管的.
上帝賜福的.
不是單單他所揀選的國度以色列民.
而是萬國都在他的主權底下.
全地的各國國民國方.
得以生存.
是因為上帝普遍的恩典.
你記得創世記第十章.
有70個列國.
都是上帝普遍的恩典.
以至他們能夠擴散在不同的地方.
然後我們看見過.
羅德的後裔.
摩訶 阿們.
他們得以存活.
也是上帝普遍的恩典.

$^{201}$以素 以東國.
都是上帝給他們一個共同的恩典.
而以色列國就是在中間.
成為神的國度.
使萬國能夠得福.
所以這個繼承人的問題.
就在這個故事做一個交代.
真正的繼承者就是以實.
不是以實瑪利.
在這裡你就看到.
上帝用他的恩典.
用他的能力.
他不斷地克服不同的問題.
以至這個國度能夠繼續向前.
他的計劃能夠繼續延伸下去.
繼承人這個問題.
上帝在他的恩典.
在他的安排裡面就解決了.
真正繼承家業的.
繼承這個國度的.
就是以實一個人.
現在我們看故事第一段和第三段.
兩段的故事.
分別在兩個不同的地方發生.
兩個故事都是說到阿巴拉罕的家.
去到不同的地方.
去到基拉爾和別士巴.
發生一些困境.
一些的威脅.
我們先看第一段.
二十章第一節.
阿巴拉罕從那裡.
應該是希伯倫山地.
就遷往南地尼加夫.
就在加底斯和舒爾中間.
阿巴拉罕在基拉爾的時候.
曾經提到他的妻子撒拉說.
她是我的妹妹.
於是基拉爾王阿比米勒.
就把撒拉搶了.

$^{241}$基拉爾應該是今天南地的.
Tel Haro.
這個地方是在加薩.
和別士巴中間.
聽起這個故事.
我們就發現有些相似.
之前好像耳熟能詳.
聽過的.
是太太當妹妹的情節.
那個危機.
兩個故事都是相似的.
因為撒拉又陷在外邦的宮廷裡面.
在那裡遇險.
第一次在埃及.
這次是在基拉爾.
菲利斯王阿比米勒的宮廷裡面.
還有什麼相似的地方.
就是那個原因.
就是阿巴拉罕軟弱.
脆弱的時候他怕死.
生命受到威脅.
想用這個策略來求存.
真正原因為什麼會這樣做.
你聽回第七講你就知道了.
但兩個故事.
雖然有相似.
也有不同.
最不同的地方.
就是時間.
那個時間性.
在阿巴拉罕的旅程.
開頭的時候.
埃及發生外邦宮廷的危機.
那個真正的威脅.
就是阿巴拉罕會失去他的妻子撒拉.
如果沒有了撒拉的話.
就不成家也不成國了.
上帝的計劃就會告退.
這次在旅程的末端的時候.
撒拉這個時候應該已經有新人了.

$^{281}$只不過還沒有見到.
危機就不單單沒有了妻子.
而是如果阿比米勒這個外邦王.
占污了撒拉的話.
撒拉所生的兒子還會不會是阿巴拉罕的後裔呢.
阿巴拉罕期待已久的繼承人.
是不是真的能夠成就呢.
神的計劃會不會因為困境.
人的軟弱.
能夠成為一個很大的問號呢.
結果是怎麼樣呢.
第一故事的情節.
就是耶和華神去解困.
耶和華首先在第三節.
就跟阿比米勒王說什麼呢.
他說你是該死的.
因為你所接回來的這個女人是一個有夫之婦.
阿比米勒王當然是不服氣了.
因為他根本不知道內情.
所以第二次神就在夢中跟他說.
他說我也知道你作者是問心無愧的.
所以我才會攔阻你免得你得罪我.
所以你不可以碰這個女人.
現在你要將這個妻子還給阿巴拉罕.
因為他才會為你禱告.
你可以存活.
你若不將她歸還的話.
你當知道你跟你所屬的一切人.
都死定了.
兩次上帝去介入.
就是用死來作警告.
警告這個愛邦王不要亂來.
所以看回故事的之後發展.
你就知道了.
撒拉就被解困.
然後救了撒拉出來.
不僅如此.
阿巴拉罕也可以在基拉爾附近去寄居.
第一段危機解困之後.
因為上帝因憐上帝的能力.

$^{321}$第三段就發生在別氏壩.
同樣是講到阿巴拉罕跟愛邦王阿比米勒之間的關係.
二十一章第二二節.
那時候阿比米勒王和他的軍長菲葛.
就跟阿巴拉罕說.
他說我看到你所做的一切.
有神跟你同在.
現在不如你去起誓.
我們也去起誓.
然後你以真誠.
不要用詭詐來代我和我的子子孫孫.
我怎麼恩待你.
你也要怎麼恩待我.
阿巴拉罕就回答.
他說我願意起誓.
這些都是雙方起誓.
世世代代彼此恩待近東的立約的慣例.
之後發生了一些什麼危機呢.
二十五節.
阿巴拉罕因為阿比米勒的僕人.
強佔了他所挖過的水井.
就指責阿比米勒王.
水井的危機出現.
阿比米勒王當然不認賬.
他說我根本不知道這件事.
現在才聽見而已.
他不認賬.
於是乎兩個的權勢.
兩個的家族就有一些張力.
有一些威脅.
這個阿比米勒王.
肯定比阿巴拉罕勢力更加大.
跟之前的危機困境有些不同.
因為之前是撒拉在外邦的宮廷裡面.
有威脅.
這一次的威脅就是阿巴拉罕的生存.
阿巴拉罕挖了口井.
對於遊牧民族來說.
口井水源是非常重要的.
故事就是問到底阿巴拉罕在這些威脅下.

$^{361}$無助之中.
他可不可以真的在迦南地立足呢.
會不會被阿比米勒王來威脅呢.
最後這個故事發展.
就是有立約的祭典出現.
於是阿巴拉罕就獻來一些牛羊.
送給阿比米勒.
然後雙方說了一些話.
都是當時近東立約的時候的一些對話.
到三十節了.
意思就是說.
你接受了這些羊羔的時候.
就好作阿巴拉罕挖口井的證據.
第三十節.
因此那地方就名叫別士巴.
Bisheva.
因為兩個人在那裡起了誓.
大家就和睦共處了.
於是阿比米勒和他的軍長.
就不再威脅他回到非利士頓.
我們可以用一個人物性格的角度來看這個故事.
我們說阿巴拉罕族長為人比較平和的.
不想搞事的.
大家用和平的手法來處理就好了.
但如果我們這樣看的時候.
我們還未解答到這個問題.
就是這個故事為什麼要放在這裡.
為什麼要記載一個這樣的故事.
在阿巴拉罕故事的末端.
在結尾的時候.
為什麼要記載這樣的事跡呢.
說到在外邦的宮廷.
撒拉受到威脅.
說到在外邦人的地方.
阿巴拉罕的生存受到威脅.
當然其中一個原因.
就是向第一代的讀者.
當時古以色列人.
他們已經住在迦南地.
去到解釋到底別氏巴比赦瓦這個地方.

$^{401}$它的由來是怎樣.
比赦瓦的意思就是立約之井.
比赦瓦.
所以對於古以色列人.
這個是他們感興趣的故事.
原來一個這麼出名.
南面邊疆的一個城鎮.
比赦瓦.
原來有列祖阿巴拉罕的事跡發生過.
但可能也有另外一個原因.
不單止解釋地名的由來.
因為這裡說的是族長阿巴拉罕.
所謂信心之父.
他在水井的爭執裡面.
問題就是他能不能夠生存下去.
一個遊牧民族.
能不能夠在這裡暫居.
在這裡生存立足.
之前撒拉可以生存.
因為上帝的恩典.
這次阿巴拉罕能夠生存得到.
能夠在這裡寄居.
危機得到化解.
原來是因為上帝的恩典.
是因為上帝的祝福.
大家記得阿巴拉罕的整個故事敘述.
好像三文治的結構.
現在我們在看末端.
在處理的環節是C這個分段.
所以為什麼在阿巴拉罕的旅程的末端.
會提到太太當成妹妹的故事.
以及看到阿巴拉罕受到威脅.
能否生存這個故事.
看到他脆弱.
看到他面對外敵的時候.
他是多麼的無助.
如果你想想.
你寫一個英雄列祖的傳記.
會不會是這樣寫呢.
開頭的時候.

$^{441}$你可能會看他無能和軟弱.
但在傳記的結束的時候.
我們大概也不會這樣看.
但聖經有一個獨特的地方.
就是將他無能的地方.
將他軟弱表露無遺.
目的就是一個屬靈的異議.
如果不是上帝的恩典.
去保守的話.
他一切都不能成事.
甚至不能立足.
不能寄居.
這個可能才是重點.
記載一個人的故事.
記載信心之父的事蹟.
居然講到他不是豐功偉績.
而是他失敗的地方.
或者是他軟弱脆弱的地方.
可能這才是我們今天要聽的.
屬靈的故事.
今年2021年年頭的時候.
有一個天主教神學家叫做.
Hans Kuhn 去世.
Hans Kuhn 最出名的.
不單止是他的著作《等臣》.
有很多的著作和很多的書.
可能他最出名的地方就是.
他經常都不服天主教的教理.
經常都挑戰梵諦光.
甚至觸動教廷來開除他的教席.
Hans Kuhn 其中一本小的書.
叫做《The Church Maintained in Truth》.
就是教會在方的真理裡面被保守.
他講到不是教會去捍衛保守真理.
相反是真理是福音保守教會.
你看看教會歷史.
其實不是這樣.
教會經常犯錯.
會偏離的.
但是教會到今天仍然存在.

$^{481}$為什麼呢?.
不是因為教會厲害.
而是因為上帝的恩典和上帝的保守.
舉例中世紀.
教會腐敗贖罪券等等問題.
不是興起.
就是馬丁路德等等的改教者.
他們來改革.
然後叫人來歸回真道.
近年大家都知道.
有很多大型的教會.
大型的機構.
人人以往都是夢寐以求的.
覺得他們的成功就是赤道.
會眾是去追捧一些大的人物.
發生什麼事.
大家心裡有數.
榮譽多還是羞辱多?.
聖理往往充分透露.
神的工作原來最大的難阻.
往往不是魔鬼撒旦.
當然我們會說.
魔鬼撒旦在背後做很多事情.
但神工作最大的難阻往往是.
神的子民.
教會.
為什麼?.
因為教會不單專注為主.
不看重上帝的恩典.
我們很容易成為上帝的難阻.
成為了他工作不能夠前行的阻礙.
神將真理託付給教會.
門徒應該是忠心宣講.
捍衛天經地義.
但如果同時我們忘記了.
要天天悔改.
天天去認罪.
唯獨是恩典.
Great alone.
我們不能成為神的子民的話.

$^{521}$恐怕我們很容易成為福音的絆腳石.
成為神工作的難阻.
換言之.
你和我一生要戰戰兢兢地.
作成我們得救的功夫.
戰戰兢兢地.
怕我們自己因為軟弱.
因為我們的無能脆弱.
而我們偏離主導.
我們失去這種態度的時候.
我們就不能夠繼續向前.
如果不是每一天.
我們在上帝的恩典裡面.
謙卑.
悔改更新.
不斷調教我們自己的話.
問一問我們自己.
到底我們今天是否完全全心.
降伏在主.
降伏在上帝的計劃裡面.
降伏在祂的道的底下.
我們往往可能就成為阻攔阻.
特別信主久了的人.
越信得久.
就越容易僵化.
以為我們在屬靈上.
我們在旅程上.
我們比別人走得早很多.
我們食鹽多過別人吃米.
經歷多了.
豈不是我們很容易成為僵化.
教會很多時候就是上一代.
是我們.
我們以為自己走得久了.
反而可能是我們阻礙地球轉.
救主恩代我們.
幫助我們.
調教一下我們.
不要以為我們食鹽多過食米.
我們永遠都是站在高地.

$^{561}$穩如泰山.
想一想.
到底誰才是真正的.
信仰成長的旅程.
讓你選擇.
第一個選擇.
就是信仰的旅程.
是不斷去兼顧我自己舊有的信念.
聽慣聽熟.
以往我起初所認識的.
現在我聽回的時候.
都是安安樂樂的.
這是第一個選擇.
第二個選擇.
就是信心的旅程.
是不斷去拆毀舊日的信念.
拆毀今天的我.
拆毀昨日的我.
不斷去調教更新和挑戰我自己.
這是第二個選擇.
你覺得選擇一還是選擇二.
才是聖經所說的.
Spiritual Journey.
真正去到跟隨主的旅程.
肯定是後者.
因為後者才是活著的信仰.
不是一個僵化的信仰.
肯定是第二個.
我就寧願選第二個.
因為在第二個選擇裡面.
在這種旅程裡面.
所要證明的不是我有多厲害.
而是上帝的恩典有多大.
神的明星是值得去稱讚.
你的選擇會是怎樣.
謙卑我們自己.
在上帝的恩典裡拆毀舊日.
重新再更新.
去跟隨主.
結束的時候.

$^{601}$看這個故事尾段.
二十一章第三十三節.
這些危機都完結了.
上帝化解了這些所有問題.
三十三節.
阿巴拉罕在別氏巴栽了一棵垂絲柳樹.
這些是曠野常有的大樹.
他在那裡叩高耶和華永生神的名.
阿巴拉罕的旅程.
現在來到差不多末端了.
在跌跌碰碰裡面找到的不是自己有多厲害.
而是上帝的恩典.
往往很多時候都是三步向前.
然後兩步倒退.
慢慢慢慢在跌跌碰碰當中.
在危機,困境裡面.
見到上帝的恩典.
以至我們所說的.
是永生神的名字.
救主叫我們每一個.
特別是我們這些靈情比較長的人.
不要為自己而留名.
不是為自己去建立自己的legacy.
而是上帝在我身上.
讓後人看到是永生神的名字.
神值得稱頌.
不是我自己.
我們同心低流祈禱.
上主你是那位滿有恩典的神.
你是那位由始至終都是守約的主.
我們感恩因為耶穌基督.
祂的十家的救恩.
我們這一班的罪人.
才能夠成為你的神國的子民.
主謙卑我們每一個.
調教我們每一個.
讓我們的焦點放在神你自己的恩典.
你自己的偉大,你自己的名聲.
不是我們自己的豐功偉績.
求主你幫助我們.

$^{641}$拆毀我們今天.
讓我們在你的唯獨恩典裡面.
重新悔改.
重新上路.
來跟隨你.
願你在我們生命當中.
你的名字被宣告.
你自己的榮耀被彰顯.
祈禱奉主耶穌的名字.
\newpage



\section{}
\label{sec:PyoQq9QIPJA}
\textbf{Mapping the Bible  - 1 Geography of Faith (Pastor Mark Chan)}
\newline
\newline
連結: \href{https://youtube.com/watch?v=PyoQq9QIPJA}{\texttt{https://youtube.com/watch?v=PyoQq9QIPJA}} ~~~~ 語音日期: 2021-08-13
\newline
\newline
\hyperref[sec:rJBi5QVSDGc]{\small{< < < PREV SERMON < < <}}
~
\hyperref[sec:index_chronic]{\small{[返順時目]}}
~
\hyperref[sec:index_scriptual]{\small{[返順卷目]}}
~
\hyperref[sec:x_0xBEFerwg]{\small{> > > NEXT SERMON > > >}}
\newline
\newline
$^{1}$Welcome to the first of this series called Mapping the Bible..
Now this is a special series I've always wanted to do,.
to use Google Earth, videos, pictures, and different visuals.
to help us travel to the land of the Bible.
so we can appreciate, we can hear the messages of the Bible..
You can see in my background,.
we are right above in the air of the land of Israel,.
the Bible land..
Now of course the best ways to go about studying the land.
and how it relates to the messages of the Bible.
is to go to Israel yourself, right?.
To study the land, to walk on the ground,.
to experience yourself..
If you do that, I can almost guarantee you,.
you will see the Bible in a totally different new way..
But we are still in the pandemic.
and international travel is such a hassle..
So we cannot do that easily..
However, we can use technology like Google Earth.
to help us out..
Now you may think immediately,.
why do we need to study the land of the Bible?.
Can we just get to know God and read scripture directly?.
Why do we need to travel to the land of Israel, Palestine.
in order to get a better understanding of the word, right?.
Now I believe understanding the Bible,.
our faith has so much in common with traveling..
Now think about it..
Every time you open the pages of scripture,.
you turn to the word of God..
You should recognize there is huge difference.
between our culture and the culture.
of the Old Testament and New Testament, right?.
In fact, the original languages of the Bible,.
even though translated into modern languages.
were in classical Hebrew, not modern Hebrew,.
classical Hebrew, some Aramaic in the Old Testament.
and the New Testament was written in Greek..
So on top of linguistic language differences,.
our world and the biblical world are quite different.

$^{41}$in terms of cultural values, right?.
Customs and all that..
We're not saying the Bible was simply a product.
of ancient culture, trapped in the ancient time..
We're not trying to suggest that,.
but we need to understand when God reveals himself.
through scripture, God communicates in ways.
that are intelligent to the ancient audience.
using their languages, using their metaphors,.
using their images or their cultures..
So in order for us to get a better understanding.
what scripture says or what scripture may mean to us today,.
we need to take a step back and kind of time travel.
to understand, to learn about how the first audience,.
the ancient Israelites received the messages, right?.
Now, somewhat like traveling to a foreign place,.
a foreign culture, a foreign country,.
different from our own..
We first need to get some basic understanding.
of that culture, right?.
The history, the geography of that place,.
how those people live in that region and all that..
We may get those information from a guidebook.
like Lonely Planet or find information off the internet..
So now think about it..
Every Bible reading, every time we think about the scripture.
is actually a cross-cultural activity..
Now let me begin with Psalm 125, verses one and two..
Those who trust in the Lord are like Mount Zion,.
which cannot be shaken but endures forever..
As the mountains surround Jerusalem,.
so the Lord surround his people now and forevermore..
Then the Psalm talks about God granting peace and security.
to those who are committed to the way of God,.
contrasting the wicked whose way will perish..
But how do we understand the words,.
as the mountains surround Jerusalem?.
Now I am zooming into the Old City of Jerusalem..
The yellow line you see on the map.
outlines the current Old City wall of Jerusalem..

$^{81}$The city wall dates back to the Ottoman period..
South of the Old City is the long, narrow Eastern Hill..
This is where the Bible refers to as City of David,.
Mount Zion, ancient Jerusalem..
Now if we are standing on the Eastern Hill,.
on an elevated point on the City of David,.
and take a 360 degree roundabout look,.
you will be able to appreciate Psalm 125..
So we're looking north now..
We see the Al-Aqsa Mosque, Temple Mount in front of us..
Notice the elevation is higher than where we are,.
the City of David..
Now turn a little bit to the right,.
northeast is Mount of Olives..
Elevation again is higher..
Continue to the east, we see the Palestinian village,.
Silwan, on the foothill..
Now looking south, we see a mountain range..
This is the southern approach to Jerusalem..
On top of the hill is the United Nations Headquarters..
Again, the elevation is higher.
than where we are standing now..
Below us, there are some archeological remains..
They date back to the pre-Israelite,.
Canaanite, Jebusite period,.
called the Step Stone Structure..
Now southwesterly, we see a high hill called Western Hill..
On top of it is the Catholic Church Domitian Abbey..
Now continue turning back to where we started,.
we see the Temple Mount again..
Now you should be able to appreciate.
the mountains surrounding us is hugging Jerusalem..
It's the City of David, right?.
Being surrounded by higher, elevated mountain range..
So this psalm attributed to King David,.
but whoever wrote this psalm originally.
has to be someone who is familiar.
with the local physical settings of Jerusalem, right?.
Otherwise, he would not be able.
to write something like that..

$^{121}$The psalmist is using local geographical terrain features.
to talk about something beyond geography,.
talks about theology, talks about faith, right?.
If we put our trust in God, look around you, right?.
Just like the Lord Yahweh will surround Jerusalem.
and protect Jerusalem,.
just like the mountains are protecting it, right?.
If we put our trust in God, the Lord God will protect us,.
will grant us security, will grant us peace..
(gentle music).
Now continue with our orientation of the land..
This land of the Bible has many different names and labels..
Sometimes it is called Israel, sometimes called Palestine,.
and Palestine was a name given by the Roman emperor.
after the defeat of second Jewish revolt.
in the second century AD..
The Old Testament also has different names and labels.
to describe this land..
Land of Canaan, the promised land,.
land flowing with milk and honey,.
and many call this land the Holy Land today..
It is a religious symbol, a religious label..
And yes, this is a sacred place.
because it is a sacred home of three different.
major monotheistic faiths, Judaism, Islam,.
and Christianity, of course..
And some still call it Near East or Middle East..
Why?.
Well, about 200 years ago in the 19th century,.
from the perspective of the European powers,.
this region is, of course, the East, right?.
Now, back then, it was under the control.
of the Ottoman Empire..
At the end of 19th century,.
the first Sino-Japanese War broke out far away from Europe..
And then the Ottoman Empire,.
which lasts for about 400 plus years, started to crumble..
So to the Westerners, to the Western European powers,.
there were two different Eastern problems, right?.
There were two regions in the East.

$^{161}$that have some political struggles..
One was the Far East..
The Near One is naturally called Near East..
Now, then in the 20th century,.
the Western countries recognized that the region.
of Persian Gulf and India was strategic.
to their political interests..
So people started using terms like Middle East.
to describe Persian Gulf..
And then as times go by,.
because Palestine is connected to Persian Gulf and beyond,.
the definition of Middle East enlarges.
to include the regions between Persian Gulf.
all the way to the Mediterranean Sea..
So yes, the land of the Bible has many different names,.
but one thing remains constant..
It is the geographical dynamics, the geographical feature..
This is a land connecting three different continents..
As you can see, Africa, Asia, and Asia Minor in Europe..
It is the center, the land bridge of the ancient world..
It's a springboard, right, to go from one region to another..
It's a place where people are connected,.
a place where people come and go all the time..
Now, with this geographical perspective,.
I would like to suggest with some biblical geographers.
would say this land is appropriately described.
as the land between..
It's always the land between,.
people come and go all the time..
And the best way, the best value of this land.
is not to build an earthly empire, but to communicate,.
to tell the world a message, to broadcast a message..
A Bible geographer rightly says,.
this land, the best value is to build a podium..
It is to communicate..
So if you think about it,.
isn't that why God would choose this land.
to begin his redemptive program?.
Isn't that why God would place his people.
to start his redemptive purpose in this place, right?.

$^{201}$Now let's look at a second passage today, Deuteronomy 11,.
Deuteronomy chapter 11..
Now the book of Deuteronomy was the sermons.
of the servant of the Lord, Moses..
When Moses was standing in the east side of River Jordan,.
he was not allowed to enter Canaan,.
but he is called by God to tell the Israelites,.
the next generation, as they go into Canaan,.
they have to pay attention to a few different things..
Chapter 11, verses eight and following,.
"Observe, obey all the commands I am giving you today,".
Moses said to the Israelites,.
"so that you may have the strength to go in.
"and take over the land that you are crossing.
"the Jordan to possess.".
Verse 10, "The land you are entering to take over.
"is not like the land of Egypt from which you have come,.
"where you plant your seed in Egypt.
"and irrigated it by foot as in a garden,.
"a vegetable garden.".
Verse 11, "But the land, the land of Canaan,.
"you are crossing the Jordan to take possession of.
"is the land of mountains and valleys.
"that drinks rain from heaven..
"It is the land the Lord your God cares for..
"The eyes of the Lord your God are continually on it.
"from the beginning of the year to its end.".
13, "So if you are faithfully obeying the commands.
"I'm giving you today to love the Lord your God,.
"to serve him wholeheartedly,.
"then I will send rain," God says, "on your land..
"I will send rain on your land in this season,.
"in this right timing..
"Autumn and spring rains will be sent to you.
"and you may gather in your grain,.
"your harvest, your new wine, and your oil..
"I will provide grass in the fields.
"for your cattle, for your flock,.
"and you will eat and be satisfied.".
Now think about this passage..

$^{241}$This passage is full of geographical details..
Look at the land of Canaan..
It is totally different from Egypt, right?.
Egypt, the delta, is flat with all the branches of Nile,.
because the Nile will flood annually,.
constantly supplying water to farmer..
It is not like Egypt..
Canaan is not like Egypt..
Nor is it like Mesopotamia,.
where you got two big rivers, Tigris and Euphrates..
The land of Canaan has no big river..
The biggest river you can find is the Jordan River, right?.
But it is very small..
It flows through the dry and thirsty land.
of the wilderness of the Jordan River Valley..
So there's not much use of it in terms of farming..
So in order to survive in this land,.
Deuteronomy says you need to understand.
this is different geographically from Egypt,.
where you come from..
In Canaan, you have to depend on rain from heaven..
Around the rainy season of October and November,.
rain clouds will move from west to east, all right?.
Rain clouds will bring in water to water the land,.
so you can farm, you can harvest,.
you will have food to eat..
But when rain is not enough,.
is when rain is not sufficient,.
you get drought and famine..
So the book of Deuteronomy says.
you have nothing to depend on in the land of Canaan,.
but God who controls the rain,.
the God who sends rains at the right time..
You got to depend on God, you got to trust in him..
You do not turn away from the God of heaven,.
to the idols, to other false gods in the land of Canaan,.
be it Baal, the storm god, or any other false god..
You got to trust in God, you got to focus on him..
Only Yahweh the Lord you can trust, and he will send rain..
So you think about it, to survive in this land is difficult..

$^{281}$It is not an easy land..
It is never a paradise..
And God settles his people in this land.
to teach them, to train them..
The land must wait for God to open the heaven and water..
The land is a spiritual classroom,.
a spiritual mentor, to train the faith of God's people,.
to teach them, to discipline them,.
to fix their eyes on God the creator alone..
(gentle music).
- Now, isn't that true for most of our life journey.
as we follow Christ?.
What COVID has taught us is that life is full of crises,.
challenges, temptations, and trials,.
and all those tough situations can happen to us,.
every single one of us, unexpectedly, suddenly..
We will be hit by a storm, right?.
And faith in the Bible has always been faith.
in the midst of difficulties..
Never easy..
But following Christ is like that, right?.
God put us in a spiritual classroom for a good purpose..
In all these difficulties and all these challenges,.
pandemic included, God has a good purpose..
It's a life-giving purpose..
The purpose is to draw us closer and closer to him..
So when we face difficulty, when we face challenges.
or uncertainties in our futures,.
whether we are going to school, determining our future,.
whether we hit rock bottom in our career,.
remember that God is mentoring us..
Like living in the land of Canaan,.
we have to depend on him..
We have to trust in him..
And knowing that God is always with us,.
God is always trustworthy..
Let me close by reading Psalm 125 again..
Psalm 125..
"Those who trust in the Lord are like Mount Zion,.
"which cannot be shaken, but endures forever..

$^{321}$"As the mountains surround Jerusalem,.
"so the Lord surround his people,.
"both now and forevermore.".
Let's pray together..
Father, we give you thanks.
for all the experiences that we have..
You put us in a spiritual classroom to train our faith,.
to pray to you, to depend on you,.
knowing that you are the faithful God.
who always cares for us..
Thank you for the lands of Canaan that we experience now..
No matter what challenges we are facing,.
trials and temptations,.
we pray that we will continue to fix our eyes on you,.
to follow you, to commit our lives to trust in you..
And pray that you will strengthen our faith..
And as we do so, we draw closer and closer to you..
In Christ, we pray..
Amen..
\newpage



\section{}
\label{sec:x_0xBEFerwg}
\textbf{舊約網上行 - 上帝的足跡: EP15 最後試煉}
\newline
\newline
連結: \href{https://youtube.com/watch?v=x_0xBEFerwg}{\texttt{https://youtube.com/watch?v=x\_0xBEFerwg}} ~~~~ 語音日期: 2021-09-01
\newline
\newline
\hyperref[sec:PyoQq9QIPJA]{\small{< < < PREV SERMON < < <}}
~
\hyperref[sec:index_chronic]{\small{[返順時目]}}
~
\hyperref[sec:index_scriptual]{\small{[返順卷目]}}
~
\hyperref[sec:bgfQNbxQms4]{\small{> > > NEXT SERMON > > >}}
\newline
\newline
$^{1}$我們來到九月望翔第十五講.
今次我們去到阿布拉罕信心旅程.
Quest 探索最後一段.
就是創世記第22章.
我們先來看經文.
.
摩尼亞地區,在我所指示你的一座山上,把他獻為梵濟..
阿布拉罕清早起來,預備好了爐,帶著兩個銅簸和自己的兒子伊薩,.
也配合好了獻梵濟用的柴,就起程到神指示他的地方去..
阿布拉罕對他的銅簸說:.
你們和爐留在這裡,我與孩子要到那邊去敬拜..
然後回到你們這裡來..
阿布拉罕把獻梵濟的柴放在他兒子伊薩的身上,.
自己手裡拿著火與刀,他們二人一同向前行的時候,.
伊薩問他父親阿布拉罕說:.
爸爸!.
阿布拉罕回答:.
我兒有什麼事?.
伊薩說:.
你看,火和柴都有了,可是獻梵濟用的羊羔在那裡呢?.
阿布拉罕回答說:.
我兒,神自己必親自預備獻梵濟用的羊羔..
於是二人繼續一起前行..
他們到了神指示他的地方,阿布拉罕就在那裡捉了一座祭壇,.
擺好了柴,捆綁了自己的兒子伊薩,.
又把它放在祭壇的柴上..
阿布拉罕伸手拿刀,要宰殺自己的兒子的時候,.
耶和華的使者從天上呼叫阿布拉罕說:.
阿布拉罕!阿布拉罕!.
阿布拉罕回答說:.
我在這裡..
天使說:.
不可在這孩子身上下手,一點也不可害他..
現在我知道你是敬畏神的了,.
因為你沒有留下你的兒子,你的獨生子不及我..
看見在後面有一隻公綿羊,.
兩角扣在稠密的小樹叢中..
阿布拉罕就去取了那隻公綿羊,.
把牠獻為梵濟,代替自己的兒子..
阿布拉罕就給那地方起名叫耶和華爾納,.

$^{41}$直到今天仍還說在耶和華的山上必有預備..
耶和華的使者第二次從天上呼叫阿布拉罕說:.
耶和華說:.
我只著自己起誓,.
你既然作了這事,沒有留下你的兒子,你的獨生子,.
我必定賜福給你,.
必使你的後裔繁多,.
將天上的星海邊的散,.
你的後裔必佔領受敵的城門,.
地上萬國都要因你的後裔得福,.
因為你聽從了我的話..
阿布拉罕回到他的同步那裡,.
他們就起程,.
一起到別士巴去了,.
阿布拉罕就住在別士巴..
這段聖經創世記二十二章獻以撒的故事,.
相信是聖經裡面最棘手,最難解的一段經文..
他在讀大學的時候,.
曾經讀過一科哲學,.
老師對基督教印象十分差,.
是非常反教的..
有一次他就用創世記二十二章,.
去駁斥聖經,駁斥整個基督教信仰..
他說上帝如何會命令阿布拉罕來獻以撒,.
又殘酷又荒謬,.
萬不講理,.
神明明將以撒賜給阿布拉罕,.
之後又說要收回,.
好像在玩弄他一樣..
說真的,聖經裡面有不少的內容,.
確實對我們現代人來說,.
是有些無稽,.
或者一聽到的時候,.
我們覺得非常難解..
原因其中一個,.
是我們用現代人的角度來看經文..
聖經裡面的一些內容,.
我們認為是無稽,.
或者我們認為,.
為什麼會做這些事,.

$^{81}$寫這些東西呢?.
我們會覺得很困惑,.
除非我們嘗試退一步,.
來放下我們現代人的眼光,.
嘗試用古代聖經世界的人,.
為什麼會寫這些故事,.
而他們聽這些故事的時候,.
未必有我們現代人的反感,.
相反,可能他們一聽這個故事,.
一明白裡面所做的事,.
他們就會馬上導出這個聖經章節,.
這個故事,.
到底跟我們的信仰,.
跟神的關係,.
有什麼道理在裡面..
故事裡面,.
有些奇怪的地方,.
第一,因為故事,.
我們讀者最想知道的東西,.
都沒有交代,.
譬如,.
撒拉,媽媽撒拉,.
到底在哪裡呢?.
阿巴拉罕做父親,.
她居然可以帶著兒子,.
離開幾天,.
然後做一些這樣的事情,.
然後撒拉母親,.
一句話都沒有說過,.
其實你看看那個描述,.
其實是很surreal的,.
是一種超現實的描述,.
怎麼會在現實生活裡面,.
是這樣發生的呢?.
其實古代的猶太拉比,.
去解釋這段經文的時候,.
他們都有這樣的看法,.
有一個拉比這樣說,.
他說原來阿巴拉罕是偷偷地,.
帶著他的獨生子離開的,.

$^{121}$結果幾天之後,.
回到家裡,.
大家安全了,.
阿巴拉罕才向撒拉交代清楚,.
到底發生了什麼事,.
撒拉一聽之後,.
她心血少,.
可能就嚇得昏迷,.
然後她就一直病,.
然後就死了,.
所以這些拉比就說,.
為什麼22章之後所交代的情節,.
就是撒拉被埋葬,.
因為他說那個刺激太大了,.
對於母親來說嚇死了她,.
第二方面我們覺得奇怪的,.
就是以撒,.
以撒在整個故事裡面,.
她說的話做的東西,.
好像都是很超現實的,.
當時來說,.
這個故事的刻畫裡面,.
以撒應該是一個年輕人,.
已經可以走三天路,.
還有擔柴上山,.
但是對一個年輕人來說,.
可不可以說聲都不吭,.
會自甘願意的,.
躺在柴上,.
然後任由宰割呢?.
之後當然也有人不同的猜測,.
在兩約之間,.
一些猶太人的傳統裡面,.
他們會說,.
以撒其實是甘願犧牲自己的,.
第三方面,.
我們覺得這個故事奇怪的地方就是,.
到最後故事的結束,.
當然我們作為讀者,.
知道整個故事的來龍去脈,.

$^{161}$最後那個結果的時候,.
我們知道上帝根本是對憲以撒,.
以撒要獻為梵祭這回事,.
是完全不感興趣的,.
其實早就預備了代替的羔羊,.
至少在這個故事裡面,.
不是明事的話,.
肯定是暗事,.
暗事就是耶和華根本不批准,.
以色列的宗教裡面有獻人祭這回事,.
因為那個故事本身的結束,.
就是說有代替的機制,.
上帝去預備羔羊,.
人敬拜神的時候,.
是不需要去跟隨當時迦南人宗教的一些風俗,.
所以不單止這個故事是暗事,.
舊約的聖經是不容許這些所謂獻人祭的,.
一些這麼邪惡的事情發生,.
在舊約的其他律法,.
不同的經文,.
包括先知書,.
都是有一律的去譴責這些獻人祭的事情,.
譬如申命記第十八章,.
那裡說你們以色列人,.
在你們中間不可以有人把自己的兒女獻作火祭,.
也不可以有占卜算命,.
法術,行邪術,念咒,問鬼,.
過陰,行巫術等等這些迦南地宗教的風俗,.
先知書後期其實都是這麼說的,.
但是一個故事說到獻長子,獻兒子,.
阿巴拉罕獻爾撒這個故事,.
到底有什麼背景呢?.
舊約裡面有一個這樣的故事,.
可能很少我們會留意的,.
就是聶旺基下第三章,.
這裡記載了一場戰事,.
就是以色列人,猶大人,.
以及以東人,.
三個國家,三王,.
來夾攻攻擊摩訶人的戰事,.

$^{201}$這些軍隊,聯軍去打摩訶王的時候,.
摩訶王快要被打敗的時候,.
聶旺基下第三章發生什麼事呢?.
26節,摩訶王見陣勢甚大,難以對敵,.
他就率領七百個拿刀的兵,.
要衝過陣去到以東王那裡,.
卻是不能,便將那應當接觸他作王的長子,.
在城上獻慰犯制,.
以色列人遭遇耶和華的大怒,.
於是三個王就離開摩訶王,.
各自回到自己的地方,.
耶和華的大怒到底是指什麼呢?.
有很多不同的可能性,.
第一個可能性就是摩訶的軍隊,.
可能看到王將接觸他作王的長子都殺了,.
都去到獻慰犯制,.
即是說他們緊急到死到臨頭的時候,.
摩訶的軍隊就不願被打敗,.
不顧一切奮戰反擊,反敗為勝,.
可能耶和華的大怒就是這個意思,.
對以色列人來說,.
這就是耶和華對他們的審判,.
第二個可能性就是以色列軍隊,.
可能看到摩訶王居然做這種事,.
殺死自己的長子,.
這些是馬目不人的事,.
令人非常嘔心,.
所以軍隊看到的時候,.
他們都覺得我們一定要退兵,.
居然天下上有這麼殘忍的事真的會發生,.
無論如何,這件戰事就這樣完結了,.
但這裡提到就是獻長子這件事,.
說出這個故事的時候,.
又沒有什麼交代,.
原因就在近東的社會裡面,.
確實是有這樣的事情,.
並且古代的讀者,古代的聽眾一聽到的時候,.
就知道這就是當地迦南地,.
有一些這麼邪惡的宗教風俗,.
所以當時的第一代讀者,.

$^{241}$他們一看這個故事,.
聽到這些情節的時候獻長子,.
他們一看就明白了,.
當時的習俗似乎是一些臨急,.
到最後一種戰術,.
就是將自己的長子獻上,.
長子代表什麼呢?.
就是代表最大的弟弟,獻給神靈,.
獻上之後就迫使神靈去回應,.
迫使神靈去照你的心意去走,.
所以長子有一個這麼大的價值,.
在埃及也有一些浮雕,.
在這裡我們看到,.
是埃及一個在Karnak的神廟,.
那個紀念法老蘭塞爾世,.
他去攻打一個迦南地的城市,.
叫做阿什基倫的時候,.
戰況是非常激烈的,.
你看到有城牆,.
城牆以外就是埃及的很多軍兵,.
這些就是埃及王戰勝之後,.
然後回到埃及,.
就有一些這樣的浮雕,.
是當時埃及帝國的中央電視台,.
再講到埃及王是多麼的威猛,.
埃及的戰是特別大的,.
好像導彈一樣,.
是很長,很厲害的,.
然後埃及王在圖上是看不到的,.
因為他在右邊,.
但是你看到的,.
就是埃及王騎著戰馬的戰車,.
馬的前腿,.
你看到是很有動感的,.
一隻很大的戰馬,.
去到遊亂阿什基倫的城市,.
你再看看這個圖畫,.
戰果是仍然未揭曉的,.
為什麼呢?.
因為城門還未開,.

$^{281}$你看到城牆上,.
迦南人的長老,.
他們就舉手向天祈禱,.
他們不是向埃及王求情,.
是向天來呼求他們的神明,.
希望神明可以立刻介入,.
拯救他們,.
右上角你可以看到,.
可能是城裡面的王或王子,.
他手裡拿著香道,.
代表著他的哀求,.
來祈禱,.
你留意城門的左右,.
你會留意到有個小孩被吊下來,.
一個在左邊,一個在右邊,.
你看到這個圖畫的時候,.
你會發現這些小孩是沒有掙扎的,.
即是說他們已經死了,.
即是說這裡在刻畫一些東西,.
就是在緊急的時候,.
特別是在戰爭當中,.
有獻人祭這些事情發生,.
迫使神明出手,.
這是什麼宗教邏輯呢?.
什麼叫宗教邏輯?.
當然神明賜那個獎給你,.
然後他也有權要得到,.
但對於人來說,.
這個獎字是天下之間最高價值的祭禮,.
是最高價值,最寶貴的東西,.
獻獎字就是要迫使神明扭轉乾坤,.
聽你講,是一種交易,.
神與人之間的一種交易,.
這當然就是迦南宗教的邏輯,.
這正正不是聖經信仰的邏輯,.
所以我們了解這些獻嬰孩的風俗的時候,.
或者可以幫助我們舒緩一下,.
如何去看待這段經文,.
現代人會感到抗拒,.
但可能第一代的讀者,.

$^{321}$他們聽見這些故事情節的時候,.
就會想起這些背景,.
亞伯拉罕獻耳札的目的,.
當然我們知道最後是有代替的羔羊,.
即是上帝在聖經信仰當中,.
在以色列的宗教當中,.
應該不會有這樣的事情發生,.
而是耶和華已經提供了一個代替的機制,.
用羊羔來贖回投生的兒子,.
所以用這種背景來寫成,.
亞伯拉罕獻耳札這個故事,.
當然應用的時候,.
不是給我們,.
或者給古代的聽眾聽,.
要他們照辦主管,.
因為看清楚這個故事的時候,.
根本故事的目的就不在這裡,.
而是用這樣的背景,.
來導出一個信心的旅程,.
信心的本質是甚麼,.
所以我們看完這些背景的時候,.
我們就回到二十二章第一節,.
經文講得很清楚,.
他說:.
「這些時候,神要試驗亞伯拉罕.」.
試驗這個詞,當然是講到一些礦物,.
它的雜質要用火燒來除去,.
以致將真正的純正的元素留下,.
好像一些精金,精銀一樣,.
或者今天我們會想起考試測驗,.
本身的目的不是壞的,.
相反的,是必須要的,.
因為目的,考試測驗,.
就是要去確認,.
提煉出一些最好的元素,.
亞伯拉罕信心的旅程也是一樣,.
由他開始的時候,.
他離開吾爾,離開哈蘭,.
當時是七十五歲,.
到了以撒出世,是一百歲,.

$^{361}$現在來到創世記第二十二章,.
以撒已經長大,是一個年輕人,.
亞伯拉罕在這個時候,.
在近東的社會當中,.
是有了所有最好的東西,.
不單止有了財物,.
但是在近東的社會裡面,.
財物也不是最重要的,.
更重要的就是兒子,.
有個兒子,有個長子,.
這個才是最要緊的,.
這個有一點像華人的傳統思想,.
生一個兒子,就是為了養兒防老,.
所以有個兒子,.
就好像有今天一堆養老金,.
是我們最寶貝的東西,.
是我們最倚靠的東西,.
上帝祝福亞伯拉罕,.
你會不會想到,.
往往成為亞伯拉罕,.
去信靠神,去跟隨神的難處,.
其實想想我們自己也是一樣,.
往往上帝給我們很多的祝福,.
你想想,問心那句,.
是不是很多時候這些祝福,.
難了我們去跟隨神,.
年輕的可能就會去求,.
說有一份很好的工作,.
有個很好的career,.
很快上帝真的給我們一份很好的career,.
我們風生水起,.
佔據了我們生命的,.
往往就是這些物質的東西,.
而賜給我們工作的上帝,.
我們往往就將他撥走,.
或者未結婚的時候,.
我希望有個伴侶,.
上帝真的賜給你伴侶,.
到你有伴侶之後,.
你就有很多的牽掛,.

$^{401}$很多的掣肘,很多的難處,.
變成了一個情況,.
就是神給你的祝福,.
其實很多時候,.
你騎劫了我們的生命,.
舉個例子,.
好像你送給你的小朋友,.
兒女一些禮物的時候,.
誰知道他反過來,.
他看重的是你所送給他的禮物,.
而不是送禮物的你,.
這就成為一個最大,.
關係上的難處,.
整段的耳塞的故事,.
耳趾這個字,.
出現了十字,.
很強調那種情感的依附,.
那種的依戀,.
你所愛的,.
你最鍾情,朝思暮想的,.
在你生命當中不能取締的,.
獨一無價之寶,.
不能夠取締,.
神要亞伯拉罕,.
就是獻上這一種,.
當他不肯獻上的時候,.
這就會成為他一生裡面,.
身隨身的難處,.
我們自己信心的道路,.
有沒有一些我們覺得是無價之寶,.
是不能被取締,.
好像是我們的心肝寶貝,.
朝思暮想不斷去.
鉗制著我們,難阻著我們,.
這個故事就是說,.
如果你肯將這些獻上的時候,.
你不計這些東西,.
你的信心的焦點,.
是對準著賜福給你的神,.
不是神所賜給你的福,.

$^{441}$你的信心才能夠去到一個真實的地步,.
22章第16節,.
那裡說什麼,耶和華說,.
我指著自己起誓,.
你亞伯拉罕既然肯走這一步,.
來到金元璋你最深愛的,.
你以為不能取締你最鍾情的,.
來到獻上,.
沒有留下你的意志,你的獨生子,.
這個才是一個真實的信心,.
沒有留下不給,.
沒有留下不給回賜福的主,.
這種信心才是真實的信心,.
這個故事,其實最主要不是難解,.
不是要舒緩我們現代人,.
甚至古代人覺得很難啃的地方,.
而是最困難的地方,.
就是對我們的挑戰,.
你自問一下,.
神賜給你的祝福,.
是不是佔據著你自己的生命,.
佔據你的生命多過神自己,.
神所賜的禮物,好處,.
是否已經騎劫了你生命的大部分,.
還是這個故事要去挑戰我們,.
獻上繁祭,.
絲毫不漏的,.
我們知道上帝只是託管給我一些祝福,.
要我管理一段時間,.
可以是兒女,可以是財產,.
可以是你等了很久很久有的祝福,.
終於來到了,.
到底神是否讓我們看清楚,.
給我們一種試煉,.
讓我們知道真正主權是賜福的神,.
我們的信心照樣不是神的祝福,.
而是祝福的神..
第二方面,.
這裡似乎也告訴我們,.
挑戰我們,.

$^{481}$想清楚我們跟隨神的時候,.
這個旅程當中,.
我們有的回報,.
阿巴罕已經走到這個故事最後的一段,.
也是故事最後的高峰,.
神給他的試煉也是最後的一段,.
你記得在阿巴罕的故事,.
由十二章開始,.
第一次他去到踏上這個探索旅程的時候,.
神跟他說什麼呢?.
十二章第一節,.
對阿巴罕第一次講話的時候,.
他說你要離開你本地本族本家,.
到我所指示你的地方,.
我必使你成為大國,賜福給你,.
使你的名為大,.
你也必使別人得福,.
給你祝福的,我必祝福他,.
救助你的,我必救助他,.
地上萬族都要因你得福,.
第一次的呼召,.
第一次踏上這個step of faith,.
已經是困難的,.
但是你聽著聽著就明白,.
再難也好,.
起碼頭一次對他的認許,.
也是有incentive,.
有獎勵的,.
再難也好,.
起碼都是覺得值得的,.
或者多多少少,.
這些祝福,這些認許,.
都會迴盪,.
有incentive當然可以踏上信心的旅程,.
但是到最後這裡,.
第22章,.
就沒有了incentive,.
沒有了獎勵,.
相反上帝給他最好的祝福,.
他要拿回,.

$^{521}$拿回這些所謂的獎勵,.
到底他是不是願意繼續去,.
在舊約聖經裡面,.
去這個字,.
是一個很特別的字,.
在希伯來的原文,.
出現在這個地方,.
第22章,.
和第12章,.
沒有其他地方出現過這個去字,.
lekkaleka,.
我們說過這個字,.
可能一個翻譯比較好一點,.
就是journey of yourself,.
就是擺脫你枷鎖的一個探索,.
你肯走這個quest的時候,.
就是從頭,第12章,.
到尾,第22章,.
都是lekkaleka,.
都是你journey out of yourself,.
就算你走了很多年,.
你都仍然要求上帝,.
你都去練正我們的信心,.
要我們去信服上帝,.
不是為了有回報,.
不是為了有著數,.
你想想我們基督徒的生命裡面,.
很多時候我們和神的關係,.
往往搞亂了,.
往往有問題的時候,.
就是我們在說回報,.
和神建立關係,.
我花時間,花心機,.
是不是有回報,是不是有亮光,.
如果沒有亮光,我就不做了,.
結果我們很多時候就不會去親近上帝,.
我們去敬拜主,.
很多時候我們慢慢慢慢,.
就會心裡面有一種要回報,.
所謂要著數,要benefit,.

$^{561}$今天編導好不好聽,.
是不是動我心,.
今天的音樂敬拜的詩歌,.
是不是動人,.
走時候,願結聚會,.
我們是不是覺得有得著,.
這些未必是壞事,.
但是往往我們用了這樣的心態的時候,.
我們就是在計算值與不值,.
值或者不值,.
這種是一種要回報才會跟隨的心態,.
我們不知不覺就會說著數,.
你想想很多時候,.
初順的弟兄姊妹很熱心,.
他們很多時候都不會去想會不會有回報,.
相反,是一些事奉得久了,.
有沒有人稱讚,.
我有沒有title,.
或者我自己是否覺得過人,.
有沒有fulfillment,.
在我們這個亂世當中,.
我們會計較,會計算,.
如果要說一句公道的話,.
說一些公義的話,.
恆公義好憐憫,.
我們說不好,.
或者對我來說未必有好處,.
我何必要做呢,.
或者說到金錢,.
財物的奉獻,.
去support或者支持某一些你覺得有意義的工作,.
或者聖靈感動你去支持的工作,.
你可能會搞一搞,.
我們就會說有沒有好處,.
或者我們可能真的很有餘的時候,.
我們才會拿出一點點,.
又是一種說回報,.
說有沒有benefit,.
我們才會走出來,.
是conditional,.

$^{601}$是說條件的,.
如果你是有疑慮的,.
你就應該很清楚,.
特別是疑慮小的事情,.
是因著什麼緣故,.
他去做你所吩咐的事情,.
你就明白他是不是真心,.
到底是不是為了獎品,.
好像一個交易,.
一個exchange,.
一種這樣的宗教邏輯,.
就是迦南宗教的邏輯,.
而耶和華在舊約聖經,.
甚至新約聖經,.
我們耶穌基督的信仰當中,.
都是不同的,.
我們說的是上帝先安我們,.
然後我們才甘心情願,.
要去回報,.
無論代價有多高,.
我們都願意去回報,.
願意去跟從,.
聖經的信仰,.
由始至終都是這樣,.
這樣才是一種真實的信心,.
我們現在說的是舊約的聖經,.
你知道在舊約的聖經信仰,.
只有Old Testament,.
在舊約的聖經當中,.
是沒有永生這回事,.
因為上帝進程的啟示當中,.
在舊約是沒有清楚啟示有永生,.
在新約的時候,.
耶穌基督死了裡面復活,.
為我們成就了永生,.
eternal life,.
我們就知道,.
我們最終會去哪裡,.
我們的destiny,.
就是去到天家,.

$^{641}$這是一個很大的回報,.
因為這樣的回報,.
很多時候我們才願意去事奉主,.
去跟隨主,.
但在舊約的以色列人來說,.
一個真信耶和華,.
歸向耶和華的人,.
他是因著什麼,.
去做一個這樣的決定,.
是沒有永生的incentive,.
沒有所謂永恆生命,.
去天家的這些著數,.
人歸向耶和華,.
主要的原因,.
就是因為神是真神,.
是因為這個跟我們立約的神,.
是值得我們相信,.
值得我們去毀身的,.
主是我真正生命的主,.
我才commit我自己,.
這是舊約裡面的信仰,.
只是這種信仰,.
就真的給我們一個很大的挑戰,.
到底是為什麼我們去跟隨神,.
二十二章,.
這裡是獻耳札的故事,.
給我們另一層的挑戰,.
就是我們的信心,.
是否一個真正去跟隨神的信心,.
不計回報祝福與著數,.
二十二章第十六節,.
耶和華說,.
我指著自己的起誓,.
你既作了這事,.
亞伯拉罕你沒有留下你的兒子,.
你的獨生子,.
我必定賜福給你,.
必定使你的後裔繁多,.
像天上的星,.
海邊的沙,.

$^{681}$你的後裔必定會佔領受敵的城門,.
地上的蠻族,.
都會因你的後裔來得福,.
因為你聽從了我的話,.
不計回報,.
這裡是最後一次,.
神對亞伯拉罕說話,.
對他有的應許,.
你留意這些經文,.
中間一些字眼,.
都是在亞伯拉罕探索旅程開頭的時候,.
神對他說話,.
不同的應許一些重要的字眼,.
譬如你的後裔,.
必定會是天上的星,.
和海邊的沙,.
你有兒子,.
你的後裔會變成君王,.
地上的蠻族會因你得福等等,.
都有重複的,.
但在這22章,.
最後一次,.
神對亞伯拉罕說話,.
最特別的地方就是,.
分別就是次序,.
第一次,.
神應許祝福亞伯拉罕,.
《創世記》第十二章,.
是亞伯拉罕踏上信心第一步,.
之前說的,.
十二章先跟他說,.
蠻族要因你得福,.
你要成為大國,.
你要有大名,.
說完之後,.
聽了,.
亞伯拉罕就用信心回應,.
這次最後一次說話,.
次序就跟上次相反,.
就是亞伯拉罕先獻以殺,.

$^{721}$不計回報,.
神之後才跟他說這番說話,.
給他最大的祝福,.
不計回報的信心,.
不計著數,.
不計有沒有得益的信心,.
可能這種信心,.
才讓我們能夠經歷到.
神最大的祝福,.
信耶穌,.
當然有著數,.
但如果只是講回報,.
才去跟隨神,.
只是講有著數,.
我們才去到事奉主的話,.
才肯為神去到.
行我們的旅程,.
這只不過是創世記十二章,.
幾步的信心,.
到了現在,.
你要繼續向前的話,.
你就看到真實的信心,.
到底本質是怎樣,.
就是不能夠講回報,.
事奉上帝,.
跟隨主,.
去到為主去拜上,.
不計回報,.
就是這種paradox,.
這種張力,.
你放下獻上當飯盞,.
你才可以更加體會到.
上帝的偉大,.
更加使到神的名,.
能夠被彰顯出來,.
求主幫助你和我,.
讓我們跟隨耶穌基督的時候,.
我們是一種不計回報,.
甘願獻上的一種信心,.
這樣我們才能夠體會到.

$^{761}$真正的上帝的能力,.
經歷到神自己的工作,.
我們同心低頭祈禱,.
上主我們求你.
來赦免我們,.
很多時候我們跟隨你的時候,.
我們的怨虐,.
我們很多時候只是講回報,.
有什麼incentive,.
有什麼對我們的好處,.
我們才會肯走出來,.
主要求你來試驗我們的信心,.
因為我們弟兄姊妹,.
我們身處在任何一個處境當中,.
主要是你的試煉來,.
你幫助我們見到我們的試煉本身,.
就是要我們經歷你的行動,.
主要是我們心靈裡面,.
有什麼在佔據我們,.
無論是我們一些在你身上,.
接受到的祝福,.
主要是如果這些東西.
是佔據了,.
騎劫了我們的生命的時候,.
求你親自幫助我們,.
給我們恩典,.
給我們獻上,.
而獻上之後,.
我們會見到更加是你自己,.
更加是你自己能力的彰顯,.
你自己工作的彰顯,.
願你在我們生命當中,.
你的名為大,.
願意我們跟隨耶穌基督的時候,.
我們是不後悔,.
我們願意獻上最好的,.
給我們的事奉,.
對我們生命裡面,.
我們不斷去毀身於利,.
叫我們的生命,.

$^{801}$能夠使別人能夠得福,.
使你的名在我們生命當中,.
為大..
祈禱奉主耶穌的名字..
\newpage



\section{}
\label{sec:bgfQNbxQms4}
\textbf{2021.08.29  God's Goodness - Pastor Sam Johnson}
\newline
\newline
連結: \href{https://youtube.com/watch?v=bgfQNbxQms4}{\texttt{https://youtube.com/watch?v=bgfQNbxQms4}} ~~~~ 語音日期: 2021-09-06
\newline
\newline
\hyperref[sec:x_0xBEFerwg]{\small{< < < PREV SERMON < < <}}
~
\hyperref[sec:index_chronic]{\small{[返順時目]}}
~
\hyperref[sec:index_scriptual]{\small{[返順卷目]}}
~
\hyperref[sec:XIq2yvbK8oo]{\small{> > > NEXT SERMON > > >}}
\newline
\newline
$^{1}$I'm really blessed to be here and excited to see friendly faces, and it's.
it's a blessing, it's encouragement to my heart. It's really cool to be in this.
room with you all again, and just during the time of worship, just thinking about.
all of the time that we've had in this room in the past of worshipping the Lord.
together and fellowshipping and sharing meals and and even playing games in.
times of just fun and joy, and it was just bringing to my mind the beauty.
of the family of God and of the community and the fellowship that we.
share, and so it really is good to be here again in this room with you all and.
and to fellowship and to worship together again. That's a beautiful thing..
I'm blessed to be here this morning, and it's like being back with family, which.
is really a really a good thing. This morning we're gonna be looking in.
Matthew chapter 10, and we're gonna be talking about generosity, and I had.
started this series a month or so ago about our response to God's goodness in.
our lives, and so this is a part of that series, and we have, we're kind of.
picking up where we left off. It's been a few weeks since we talked about this,.
but it's fine because each one of these topics really stands well on their own,.
and there is so much of God's goodness in our lives, so much more even than we.
see or identify a lot of times, and God wants us to respond to Him in response.
to that goodness that He pours into our lives, and so the first few weeks.
that we talked about this, we talked about the response of gratefulness, how.
important it is for us to be grateful for God's goodness in our life, how.
that's a response. Then we talked about repentance and how important it is for.
us to respond to God's goodness with repentance and aligning our lives with.
God's righteousness and with what God is doing in our hearts and our.
lives and in our world, and then we talked about faith. We talked about the.
fact that our response to God's goodness should be one of faith and that God's.
goodness in our lives gives us an archive of testimony that we can go to.
again and again in times of uncertainty or times of challenge and have faith.
because we see how good God has been in our lives, and today we're going to talk.
about generosity, that one of the responses that we should have to God's.
goodness in our life is generosity and beginning to give to others. So we're.
going to look at verse 8 of chapter 10 here, and I'll read this.
verse and we'll talk briefly about the context. The verse says this, "Heal the.
sick, raise the dead, cleanse those who have leprosy, drive out demons, freely.
you have received, freely give," and this is Jesus speaking here, and the context.
of this passage, Jesus is speaking to his disciples and he's sending them out in.
this passage, he's sending them out on kind of the first church mission trip.
that ever took place. It wasn't the first mission trip, undoubtedly. God's.
mission has been in the world since the beginning, but this was the first.
instance of a mission trip of Jesus' followers. So Jesus is sending them out.

$^{41}$in Matthew chapter 10. He's giving them all these different instructions and.
that's the context of these words. And in verse 8 here, he reminds them that they.
should freely give to others as a response to what God has given to them..
And that's the principle I want us to think about today, that God has been so.
good to us in so many different ways and out of that we should have a response of.
giving to others. We should have a response of generosity. Another way that.
this is phrased in some different translations is it says, "You didn't.
pay for this. You did not receive the grace of God through payment, so don't.
ask anyone else for that. Be willing to give freely because you have received.
this freely." And as we talk about generosity today, I want us to think.
about it in a few different ways. Obviously there's generosity in giving.
with our material resources or our finances. That's probably the one that we.
think of the most when we hear the word generosity. We think of giving gifts to.
someone else, some material item, wrapping a present for someone or we think of.
giving money or something like that. And that's certainly an application of.
generosity and one of the applications that we're talking about.
today. But more than that, there are also applications of just our time and energy.
giving our time and energy, our presence to other people, but also giving of our.
skills and talents and abilities. God has given each one of us unique skills and.
talents and abilities that we can use to glorify him and we're to be generous.
with those as well. We're to be generous in giving those things away to others as.
well. So when we talk about generosity today, I want us to think about it in all.
of those ways, not just in terms of our money or our finances, but in terms of.
our time, our energy, our presence, and also the talents, the skills, the.
abilities that God has given us, that God has uniquely gifted us with. And we could.
also think of it in terms of giving grace and forgiveness to others. That's.
one of the most precious gifts that God has given to us. And that's something.
that it's tough sometimes for us to be generous with. It's easier for us to.
give away some money to someone than it is for us to forgive someone who's.
wronged us. And so we have to think about generosity in those terms as well. That.
when Jesus says, "Freely you have received and I want you to freely give," one of the.
things that we have received is forgiveness and grace and mercy from God..
And that's one of the things that he's challenging us to give to others. So.
we're gonna go through just some points about what generosity should look.
like in our lives as followers of Jesus, and I'm going to illustrate most of.
these from Scripture, from other Scripture passages that illustrate these.
well to us and show us what it looks like. And we'll move through some of.
these. We probably won't get through all of them today, but these are just to give us,.
begin to give us a picture of what generosity should look like in our lives.
as a follower of Jesus, and also to help maybe remove some of the misconceptions.

$^{81}$or perhaps excuses that we might have when it comes to being generous and when.
it comes to giving away. So the first thing I want us to see is that our.
generosity is really a response to God's overwhelming goodness in our lives, and.
that's what we see right here in Matthew 10:8, that Jesus is saying, "Your.
generosity needs to be motivated by what God has done for you." And that's really.
important for us to understand. We'll see that as we look at some of the other.
passages of Scripture today, but if we're trying to give for other reasons, if.
we're trying to give in order to look good in front of other people or in.
order to somehow make ourselves feel better or in order to give ourselves a.
reputation or any of those other reasons, that's ultimately going to be something.
that's self-defeating. It's not going to accomplish the purpose of glorifying God.
and of loving others in our lives very well, because it's selfish. It's.
focused on ourselves, and so our giving really needs to be motivated by God's.
goodness in our lives, and as we see what God has done for us, we respond to that..
We're able to give with the right heart, a heart that gives freely and.
gives with a motivation to glorify God. The next thing I want us to see is that.
our generosity is not dependent on our own abundance or lack, and this is.
important for us to think about. I know in my own life, that's one of the most.
simple excuses that I often find to not give to someone else, and as I said,.
whether I'm talking about giving my money or some material possession or if.
I'm talking about giving some of my time or some of my energy into a.
relationship or some of my presence in someone else's life or giving.
forgiveness or giving some of my skills and talents, one of the easiest.
excuses for me is to say, "Well, I don't really have that much. I can't.
really give into this because I don't have that much. I don't have that much.
material resources to give or I'm kind of tired right now and so I don't feel.
like I have the energy to give or I don't want to give of my skill because.
I'm not that skilled in that area." Obviously, there are limitations that we.
all have in terms of time and our resources, and so I understand that, but.
at the same time, we have to understand that God is more interested in our heart.
than He is in our resources. We see this illustrated really well in the.
story that's in Mark chapter 12. This is verses 41 to 44, and I'll read these to.
you. It says, "And he sat down opposite the treasury and watched the people putting.
money into the offering box. Many rich people put in large sums, and a poor.
widow came and put in two small copper coins which make a penny. And he called.
his disciples to him and said to them, 'Truly I say to you, this poor widow has.
put in more than all those who are contributing to the offering box, for.
they all contributed out of their abundance, but she out of her poverty has.
put in everything she had, all she had to live on.'" And this is Jesus who's telling.
His disciples that she's given more than anyone else because she's given out of.

$^{121}$her poverty and she's given with a genuine and sincere heart. And that.
story is really a powerful reminder to us that we may look at our life and say,.
"Well, I don't have a lot to give right now. I don't have a lot of material.
resources to give, or I don't have a lot of extra time, or I'm not really.
skilled in that area." You know, someone has a need in their life that's kind of.
in a specific area, and that's not something I'm good at. It's not so much.
about what you have or don't have. It's not about your abundance or your lack..
It's really about your heart and are you gonna respond to God with a heart that.
is genuine and that says, "You know what? I want to love someone else. I.
want to help. I want to serve. I want to be a part of what God is doing, and I'm.
doing that because of God's goodness in my own life." And if you respond in that.
way, then whatever you have to give will, it'll count. It'll be a blessing..
It'll help someone else. It will be something that honors and glorifies God..
And that's really clear in this passage. So don't let that stop you..
Don't say, "Well, I don't really have a lot right now." Or another version of that.
that I've used a lot in my life is sometimes I've said, "Well, I'd love to.
give in this area, but I'm gonna wait until I get more in this area of my life.
or in some other part of my life, and then I'll feel free to give." And again, a.
lot of times that's just an excuse. I've been there myself where then I've.
gotten to that point that I wanted to be at in terms of resources or in.
terms of time, and then I still didn't give. So don't let your feelings of.
saying, "I don't have enough," keep you from giving. Begin to give, and if you.
don't have much to give, that's fine. This lady in the story Jesus told,.
this poor widow, gave very little. What she gave wasn't going to be noticed by.
anyone in terms of the amount, but it wasn't the amount that was significant..
It was her heart and the fact that she was giving with the right heart. Another.
thing that's very simple but it's really clear in Scripture is that our.
generosity as followers of Jesus, if we have generosity, if we freely give with a.
right heart in response to God's goodness, that generosity leads to reward..
And this principle is throughout Scripture. There's a couple verses in the.
Old Testament that are real clear about this. Psalms 112 5 says, "It is well with.
the man who deals generously and lends who conducts his affairs with justice.".
Proverbs 22 9 says, "The generous will themselves be blessed, for they share.
their food with the poor." Now we don't give in order to get a reward..
If that's our motivation, then as I said at the beginning, it's gonna end up.
being self-defeating. We're not gonna experience God's goodness and.
glorifying Him in the way we want to. But Scripture is clear. There is a reward if.
we will be generous with the right heart, if we'll care for others, if we will.
choose to give of our time and our energy and our presence and our.
giftedness and our resources with a right heart responding to God's work in.

$^{161}$our lives, then there will be a reward from that. God will bless us and give us.
peace and oftentimes give us provision. And I think that's significant. Sometimes.
in that first point that we discussed, we say, "I'm not sure I can do this right.
now because I don't have a whole lot." And sometimes we find ourselves.
finding out that as we begin to obey God and be generous in the way that He is.
putting on our heart, that God provides for us in another way and we have more.
than we thought. Whether it's resources or energy or even those talents and.
skills that God's gifted us with, we find out that God fills in what we.
perceive as lack in a lot of other ways when we just begin to be obedient in.
practicing generosity and in giving. And I want us to see as well that our.
generosity helps to break the hold of material possessions on our souls. And.
this is big. We live in a world that's a very material world. The things that are.
important to us that can really grab a hold of our focus and our.
attention and our time and our relationships, they're material things..
And we're not supposed to be enslaved to material things. We're not. As followers.
of Jesus, we're supposed to have a heavenly perspective and see that the.
eternal things are more important than the material things. We're supposed to be.
able to see that people's lives and their relationship to God and our.
relationship to God is much more important than the things that we.
possess or the jobs that we have or the titles that go after our name. All that.
different material stuff. We're supposed to have a different perspective but it's.
easy to get caught up in everything going on around us in the world that we.
live in and to begin to have material things grab a hold of us and to begin to.
commit all of our time and our energy and those gifts that God has given us..
Instead of using those gifts to glorify God and to give to others, we can end.
up just using those gifts to pursue material possessions and to pursue our.
own reputation, to pursue our own ambition, our own pride. And one of the.
beautiful things about obeying God and generosity is that it helps break that.
hold. We find ourselves saying, "I know I'm supposed to give but it's kind of hard.
because I want to keep this for myself." That's that hold the material possessions.
have on our soul or that earthly thinking and pride and ambition and.
things like that have on our soul. But as we begin to obey and to give, we find.
ourselves saying, "You know what? I guess that really wasn't as important as I.
thought it was and now that I'm giving it away, it's not holding on to me.
as much and I'm able to see that I didn't really need that as much as I.
thought in the first place or that thing I was pursuing in my life. It isn't as.
big of a deal as I thought it was." And it's important for us to give for.
that reason. In 1 Timothy 6 verses 17 and 19 it says this, "As for the rich in this.
present age, charge them not to be haughty," or that word means proud, "nor to.
set their hopes on the uncertainty of riches, but on God who richly provides.

$^{201}$us with everything to enjoy. They are to do good, to be rich in good works, to be.
generous and ready to share, thus storing up treasures for themselves as a good.
foundation for the future so that they may take hold of that which is truly.
life." And in that passage, there's instructions being given to the church,.
to the body of Christ, and there's a whole lot of really good stuff in here..
He's saying if you're a rich person in this present age, he means if you're rich.
in the world's goods and the common definition of what we would think of as.
wealth and rich. First of all, he doesn't say like if you're rich, you're evil and.
you have to burn all your stuff. He doesn't say that. He doesn't say it's.
sinful to have material possessions, but he says that it can be unhealthy and it.
can be sinful. So he says if you are one of these people that's wealthy in this.
present age, don't be proud or arrogant about it. Don't set your hope or your.
security or your sense of stability on those material possessions because.
obviously, as we know from other parts of Scripture, they can rot, they can be.
stolen, they can disappear. But he says to set your hope on God because God is the.
one who richly provides you with everything that you enjoy. If you have.
material possessions that you enjoy and you enjoy your games and your toys,.
whatever they are, if you're a young person, maybe they're Legos. If you're an.
adult, maybe it's a boat or a car or something like that. But if you have.
those things that you enjoy, if you have experiences that you enjoy, if God has.
blessed you in some way, you need to understand that God is the one who has.
provided you with everything that you enjoy anyway. And then in verse 18, he.
says you should do good and be rich in good works. You should have an abundance.
of good works, not just of material possessions, but have an abundance of.
works and actions that honor God. And it says be generous and ready to share. And.
as you do that, then you're storing up treasures for yourself as a good.
foundation for the future in verse 19. And then at the end of the verse 19, he.
says, so that they may take hold of that which is truly life. And I love that part.
of the verse. He's saying if you're just caught up in the material possessions.
and what they mean to you, you're not really living. And if you begin to obey.
God and being generous and in freely giving out of God's goodness, then you're.
gonna begin to truly experience life and experience joy and experience peace and.
actual real contentment. And all of these things that are gonna be really healthy.
and good for you come as you give. And it begins to break that hold of material.
possessions and material aspirations on your soul, and you begin to be free from.
those things. Another thing for us to see is that our generosity should not have.
strings attached. In 2nd Corinthians 9, it talks about giving. We'll look at these.
verses more a little bit later on, but in verse 7, it says, each one must give as.
he has decided in his heart, not reluctantly or under compulsion, for God.
loves a cheerful giver. And we need to understand that when we decide to.

$^{241}$respond to God's goodness by giving, whether it's giving of our possessions.
or our time or whatever, we shouldn't do that with strings attached. If we're.
being obedient to God and giving to someone else, then we just need to give.
and not be worried about what follows. And a lot of times, even as followers of.
Jesus, as believers in Christ, we can give with strings attached. We can say, "Well,.
yeah, I feel like God's putting on my heart to give financially to someone.
else, but they don't use it in a way that I think is right. You know, I'm never.
giving to them again." Or we can say, "Yeah, I feel like God really wants me to give.
some of my time to this person in this relationship, and they're hurting and.
they're going through something, and I really know that God wants me to.
give and be generous with my time and my presence and my energy. But once I.
started doing that, you know, they didn't really respond how I wanted them to, and.
so I think I'm done with it. I don't want to give any more time to them." Or, "Yeah, I.
felt like God wanted me to give some of my talents, and I began to share.
those talents with people in my church or in my community, but no one really.
appreciated it the way I thought they were going to, and so I think I'm done.
giving." There are so many ways that we give with strings attached, and we don't.
always realize that, and that's not how we should give. If we're giving as a.
response to God's goodness, our only motivation is to honor God and to.
rejoice in what he's done for us, and it shouldn't matter how the gift is.
received or how it's used or whether we get praised or recognized for it. And I.
understand that all of us want that. I do too. That's normal. It's human nature. We.
want the gift to be used well. We want to be recognized for it. We want for it to.
be received well. I get that, but we have to learn to let go of that and say, "I'm.
just gonna give in response to God's work in my life. I'm gonna give in.
obedience to him. I'm gonna give because Jesus has freely given to me, and there.
are no strings attached, and if it goes great, then I'm gonna praise God and.
rejoice for that, and if people appreciate what I've done, then I'm gonna.
be grateful for that. But if no one seems to appreciate my gift or recognize it.
and it's not even received well, that's okay. I didn't give it to get.
something in return. I gave it out of obedience to God, and there are no.
strings attached. Another thing we should consider is that our generosity should.
align with kingdom values. Malachi 3.10 is a verse that talks a lot about giving,.
and this one that is quoted a lot. It says, "Bring the whole tithe into the.
storehouse that there may be food in my house. Test me in this, says the Lord.
Almighty, and see if I will not throw open the floodgates of heaven and pour.
out so much blessing that there will not be room enough to store it." That verse is.
a lot of times talked about in the context of giving so that God.
will give back to you, and that is a part of what is portrayed in this.
verse. But I just want us to think about the beginning of this verse. It says,.

$^{281}$"Bring the whole tithe into the storehouse." God is calling his people to.
give, but he's calling his people to give into his kingdom and into the work that.
he is doing, and I think that's important for us to consider as well. That if we.
say, "Okay, I want to give some of my material possessions or my time or my.
energy or I want to give some of my skills," we should ask ourselves, "How can.
I give those in a way that aligns with the values of the kingdom? How can I give.
those in a way that honors God and glorifies him and aligns with the work.
that he is doing in the world around me?" And sometimes you may say, "Well, I feel.
called to a certain place or a certain type of service. I feel called to be.
involved in a specific way," and that's great. And as you explore that, you may.
say, "Okay, there's some different- there's an organization that's a part of this.
that I could partner with or there's some other group that's doing.
this work that I could be a part of." And I think this is a great question for us.
to ask as we look at those opportunities to say, "Okay, well does this.
organization or this group or this movement, does it align with some of the.
values of the kingdom?" Because if not, then maybe it's not the place for us to.
invest and to practice obedience to God through generosity. We know we're.
supposed to be doing this. We know we're supposed to be giving. That's clear. But.
it is important where we give and what we're giving to and what we're.
partnering with and participating in. And one of the beautiful things is in most.
any kind of domain of culture, if there's something you want to be a part of that.
God puts on your heart, you'll probably be able to find a Christian group that's.
engaged in that area. A lot of the biggest relief organizations and things.
like that in the world, especially in the western part of the world, started out as.
Christian groups. So there's a whole lot of groups that you can find and say, "Hey,.
these people are doing good work and it aligns with the values of the kingdom.
and I can be a part of this." But sometimes, maybe not. Sometimes you say,.
"Okay, well this is the issue that God's put on my heart that I really want to.
give to financially or I want to give my time and energy to or my skills to." But.
there's not really anything in this space that aligns with kingdom values.
and some of the things going on in this space actually are a little bit in.
opposition to kingdom values. Well don't be afraid to let God draw you into that.
space and start something new and to say, "Hey, it's time for something to happen in.
this space that allows people to use their gifts and their generosity to.
glorify God and to point to Jesus at the same time and to align with the values.
of the kingdom." That might be a part of what God is doing in your heart and why.
God is stirring your generosity in the first place is maybe to meet a need that.
no one else is meeting or maybe to meet a need that people are responding.
to but no one is responding to it with the love of Christ and with the values.
of the kingdom. So I think that's a great question for us to consider as well as.

$^{321}$we think about generosity is how can we give our time and our energy and our.
resources in a way that really aligns with the value of the kingdom. And a.
similar one that's real clear in Scripture is that our generosity should.
be in Jesus's name. Mark 9 41 says this, "Truly I tell you, anyone who gives you a.
cup of water in my name because you belong to the Messiah will certainly not.
lose their reward." And I think that's powerful. Our generosity should be in the.
name of Jesus. That doesn't mean that everything we do we have to tell someone.
about Jesus or tell them, "Hey, this is in the name of Jesus," or obviously there's.
a lot of acts of generosity that God's gonna put before us and call us to and.
allow us to be a part of that we won't be able to share that and that's fine..
But any time that we can, we should let people know that this is done in Jesus's.
name on one level or another. First of all because that glorifies God, but.
secondly because it gives people an opportunity to see that it's not just.
you out there being a good person and trying to do something great for someone.
else, but it's actually the love of Christ that's in you and that's at work.
causing you to care for someone else. And that gives the people who are a part of.
that generosity an opportunity to respond not just to you but to respond.
to the love of Christ and to say, "Oh wow, this impacted my life and it meant.
something to me and I see now that it's not just because this person gave to me.
because they're a cool person but it was because of Jesus." And what does.
that mean? What does it mean that Jesus would love me and that he would motivate.
someone I don't even know to give to me or to get involved in my life or to care.
for me in some way? It opens a door for other people to be connected with God.
and to hear the message of love and forgiveness and grace that Jesus has for.
them. And so I think anytime we can, we should practice our generosity in the.
name of Jesus. And then one more, we'll just do one more today and we'll finish.
up the rest in a future message because there's a lot of good stuff here about.
generosity, but our generosity should always be to glorify God and that's.
really important. In Peter, in the first chapter of 2nd Peter verses 2 and 3, it.
says, "This grace and peace be yours in abundance through the knowledge of God.
and of Jesus our Lord. His divine power has given us everything we need for a.
godly life through our knowledge of him who called us by his own glory and.
goodness." Peter's saying, "Hey, anything that you are able to do in response to.
God that is good for others, that's coming through the divine power of.
Christ anyway. So don't go out there and act like it's you. Don't go out there and.
do good works in your family or your community or in the lives of your.
friends or with other people and give them your time and energy and resources.
and then act like, 'Yeah, I'm just a great person. That's why I'm doing this.' It's.
not. It's not you. And if it is, if it's centered on you, then it's not gonna be.
healthy anyway. It's gonna end up hurting someone else or it's gonna end up.

$^{361}$hurting you. Do those things because God is the one who enables you to do it and.
give the glory to him. You know in Matthew 6, we talked about this when we.
did the Kingdom Values study during the pandemic, but the first four verses here.
Jesus says, "Be careful not to practice your righteousness in front of others to.
be seen by them. If you do, you will have no reward from your Father in heaven. So.
when you give to the needy, do not announce it with trumpets as the.
hypocrites do in the synagogues and on the streets to be honored by others..
Truly I tell you, they have received their reward in full. But when you give.
to the needy, do not let your left hand know what your right hand is doing so.
that your giving may be in secret. Then your Father who sees what is done in.
secret will reward you." That's a Kingdom Value for generosity right.
there. Do it to honor God and to truly love your neighbor and don't let.
anybody see it. Don't let people know what you're doing. That's one of the best.
ways to love and serve and care for others and to give your time and to give.
your energy and to give your resources and to give of your skills and talents.
and abilities is in a way that's not gonna lead to you being recognized or.
you being put on some type of pedestal. Give in a way that God is glorified and.
that other people are truly empowered and enabled and that you don't end up.
being the hero of the story because you're not the hero of the story. Jesus.
is the hero of the story and we want to always practice our generosity in that.
way. There's a lot more. We'll talk about some of this in the future. I'll finish.
up this message talking about generosity. But generosity is such a powerful.
Kingdom Value. It's so clear in Scripture that our lives as followers of Jesus.
should be defined by it. That we should really be generous and overwhelmingly.
radically generous. And as we began discussing this morning, it's because God.
has been overwhelming and radically good to us. And we should respond to that.
goodness by freely giving to others. And we should ask God to really enable us.
and empower us to do it in the right way so that it honors him and it truly loves.
the people around us. Let's pray. God, we thank you first of all for how you have.
freely given to us. God, your grace and your forgiveness are things that we.
could never have earned and we could never have deserved. And yet you freely.
gave them to us. All of us have fallen short, have made mistakes, have sinned,.
have messed up. And yet through Christ, you wiped those things away. You cleaned.
up our past. You erased our sins and you opened the pathway to relationship with.
the God who created us and loved us. And that's the deepest wonder of the.
universe, but it's also the most amazing gift of the universe. And God, we have, if.
we are followers of Jesus, if we've entered into that relationship with you,.
then we have already received a gift that is really of unspeakable worth. We.
can't even begin to comprehend, much less communicate, how much worth is contained.
in the gift of grace and forgiveness. And so God, may we, out of response to that.

$^{401}$overwhelming goodness in our lives, may we begin to give to others freely. May we.
begin to lead lives that are defined by generosity and by a spirit that freely.
gives of our money, our material possessions, but also our time, our energy,.
our presence, our skills, our....
\newpage



\section{25:31-36}
\label{sec:XIq2yvbK8oo}
\textbf{Welcoming Jesus in Disguise (Matthew 25\_31-36) - Keith Draper}
\newline
\newline
連結: \href{https://youtube.com/watch?v=XIq2yvbK8oo}{\texttt{https://youtube.com/watch?v=XIq2yvbK8oo}} ~~~~ 語音日期: 2021-09-13
\newline
\newline
\hyperref[sec:bgfQNbxQms4]{\small{< < < PREV SERMON < < <}}
~
\hyperref[sec:index_chronic]{\small{[返順時目]}}
~
\hyperref[sec:index_scriptual]{\small{[返順卷目]}}
~
\hyperref[sec:9svuWXYcy2M]{\small{> > > NEXT SERMON > > >}}
\newline
\newline
$^{1}$to bring the word to us..
Does love have an open door?.
Welcome stranger.
Give us your tired and poor.
All those in need.
Come every weary soul.
Find your haven, liberty is.
For all who can breathe free.
Oh what a beautiful world, what a beautiful world.
Welcome to America!.
Does love have an open door?.
Welcome stranger.
Does love have an open door?.
All those in need.
Oh what a beautiful world.
We're here!.
Oh what a beautiful world.
Oh what a beautiful world, what a beautiful world.
Welcome stranger.
Well it's great to be with you here today..
I feel like it's been a while though.
when I realized pulling into the parking lot that.
there was another church that had parking and there were signs about where I could park and couldn't park.
and you know even being in this time post-pandemic and yet new waves of the pandemic.
ministry has just changed tremendously during this time..
The video that you saw really represents what my work is these days with World Relief..
World Relief is a local global organization so we are based here in Chicagoland.
but we have offices in 16 locations around the United States..
I'm off to the Sacramento office to speak here coming up..
Mark I know has some root there..
And then we have 20 offices around the globe.
so we're involved in quite a few settings..
But the subject matter that I'm speaking to you today about is really something that is dear to my heart.
not just something that I do for my job..
And yet I also respect the fact that we're in a pulpit here.
that this is a place where God's word is conveyed..
So you don't want to, I don't want you to hear me doing an infomercial..
I really want you to hear the scriptures today as we speak together..
And yet I think it will weave well with the conversations that we have around.
what's happening in our world and the mass displacement of people..

$^{41}$You may have heard this famous quote before..
I'm wondering if it's one you're familiar with..
The quote is, "Each of them is Jesus in disguise.".
Each of them is Jesus in disguise..
This was Mother Teresa's explanation of her ministry serving the poorest of the poor.
in Calcutta, India..
This is what she said of the vulnerable people that she encountered.
amidst her own struggles and doubts about her faith..
Her life's work, no matter how poor or destitute those she worked with were,.
was to see them as they truly are, as blessed children bearing God's image.
no matter where they've come from or who they are..
Brandon Heath, a contemporary Christian artist known for at least two hits,.
"Give Me Your Eyes" and "The Light in Me,".
was moved by Mother Teresa's words to write a song in 2008 called "Jesus in Disguise.".
Yet today, when I hear the words "disguise" and "vulnerable people,".
it's not about ministry these days..
It's about fear..
Fear of people who have been displaced..
They are labeled as a haven for terrorists who are intent on harming us.
but disguising themselves to sneak into our country..
In light of 9/11 this weekend, I can certainly see our reasons for caution here..
But a couple things I think that we should know,.
and that is that it's not likely that our country would have that happen through those we serve..
Our country's track record on refugee resettlement is a strong and stable work..
Three million people have actually been brought into this country over the past 30, 40 years now,.
since 1980 with the Refugee Act, and not one citizen has died at the hands of a refugee..
So that vetting process has been pretty thorough and extensive through the years..
Inciting fear of certain people in blanket ways to cast blame on them is bigotry and racism,.
and it should not be tolerated and should not happen, especially among those of us that are Christians..
Our country and its leaders have not always gotten that right either, and for that I do apologize..
This is not an easy topic, though, I want to say to you to discuss,.
so I pray that you hear what I say in the right spirit..
Just as it is the government's responsibility to properly vet people that are entering this country,.
it is the Christian's responsibility to welcome the stranger..
If we expect the government to do their job, then we should do ours..
Most of us who identify as Christians have never really thought about migration or forced migration in particular,.
but it's in the Bible. From its beginning, in the book of Genesis,.
we find a man and a woman violating the laws of the land and placed in removal from the Garden of Eden,.
which is forced displacement. At the end of the Bible, we find John writing the book of Revelation in Patmos,.
where he is exiled as an expatriate for his beliefs..

$^{81}$The Bible is full of story after story of forcibly displaced people..
Joseph was forcibly displaced as a victim of human trafficking..
Moses was forcibly displaced for taking justice in his own hands..
Moses certainly carried that out. David was forcibly displaced as a political refugee on the run from King Saul for 12 years of his life..
And Jesus was forcibly displaced to Egypt with his family soon after his birth to find safety from profiling and ethnic cleansing as a refugee..
And the first missionary in the Bible, if we think of it this way, Philip ran from Samaria or ran to Samaria.
because of religious persecution that had broken out in the city of Jerusalem..
Priscilla and Aquila, who were Paul's business partners and ministry partners together,.
were forcibly displaced to leave Rome and go to Corinth because of their beliefs..
That's how Paul met them, is because Caesar had thrown out all of the Jews from Rome..
Forced displacement is in the Bible, but very few Christians have recognized it or understood what it means for us and how we treat them..
World Relief for Us has a 40-day Bible reading and prayer challenge just to help churches understand more of our responsibilities in this area.
before God regarding the strangers that are in our midst..
Most Christians just have not had their eyes open to this reality..
I think one of the main reasons that we see that this day and these attitudes is because of a lack of biblical education around this area..
But the other thing that I see that we're living in today is a time of great media influence..
And media is extremely powerful in these days..
Polling data really tells us today that Christians are more influenced by media and by politics than they are by the Bible and by their faith.
when it comes particularly to this issue..
I can tell you though story after story of lives that have been changed by welcoming our newer neighbors..
The Bible is clear. There is no question with the new society that Jesus brought in that we are to welcome strangers..
Matthew 25 is probably one of those signature texts, one of those greatest texts to point this out to us..
I believe the church needs to really wake up to this reality and our responsibility before God to do what we can to help..
So this morning I'd like us to look more closely at Jesus' words in Matthew chapter 25..
We're going to be looking at the end of that chapter beginning in verse 31 when God brings true justice to all of the world..
I'd like to read to you the words of Jesus beginning in verse 31 of Matthew 25..
It says there, "But when the Son of Man comes in His glory and all the angels with Him, then He will sit on His glorious throne.
and all the nations will be gathered in His presence..
And He will separate the people as a shepherd separates the sheep from the goats..
And He will place the sheep on His right hand and He will place the goats on His left..
And then the King will say to those who are on His right, 'Come, you who are blessed by the Father,.
inherit the kingdom prepared for you from the creation of the world..
For I was hungry and you fed me. I was thirsty and you gave me drink..
I was a stranger and you invited me into your home. I was naked and you gave me clothing..
I was sick and you cared for me. I was in prison and you visited me.'.
Then these righteous ones will reply, 'Lord, when did we ever see you hungry and feed you,.
or thirsty and give you something to drink, or even a stranger and show you hospitality,.
or naked and give you clothing? When did we ever see you sick in prison and visit you?'.
And the King will say, 'I tell you the truth, when you did it for one of the least of these, my brothers and sisters,.
you were doing it to me.'.

$^{121}$Then the King will turn to those that are on His left and He'll say, 'Away with you, you cursed ones,.
into the eternal fire prepared for the devil and the demons..
For I was hungry and you didn't feed me. I was thirsty and you didn't give me a drink..
I was a stranger and you didn't invite me into your home..
I was naked and you didn't even give me clothing..
I was sick and in prison and you didn't visit me there.'.
Then they will reply, 'Lord, when did we ever see you sick or thirsty or hungry or stranger or naked and not help you?'.
And He will answer, 'I tell you the truth, when you refused to help the least of these, my brothers and sisters,.
you were refusing to help me.'.
And they will go away into internal punishment, but the righteous will go into eternal life.".
So let me just break this apart a bit for us to think about what the teaching here that Jesus is giving..
Really the immediate setting that we see here is the passage is the throne room of God for the entire world..
It's the place where all of ruling of society will happen and that throne room is God's throne room..
There's a broader setting though here in this particular text and the broader setting is the end of time..
All things have been finished and time has come to an end..
But there's a greater challenge even being extended just beyond this section of Scripture in this whole chapter..
Because we have here the parable of the ten virgins, we have the parable of the three servants, all in this chapter..
And the broader challenge is that we need to be well prepared and using what we've been given wisely..
The characters introduced to us here, other than Jesus, are the angels who were with Him.
and all the nations that have been gathered together in Jesus' presence..
Imagine that experience..
All of the nations gathered, I can't even fathom what it would be like to see all of the nations gathered together before Jesus..
And He is center stage in this..
He's center stage in the scene that we're hearing here in all of His glory preparing for the final judgment..
The opening scene and the only scene in this picture is Jesus sorting things out..
In particular, Jesus as the chief shepherd is separating the sheep from the goats, the blessed from the cursed,.
those who will go into eternal punishment and those who will go into eternal life..
What's at risk?.
Well, both futures and a vast inheritance that's been prepared before creation began..
But note this, very interesting..
An inheritance is not something that is earned, is it?.
But it's something that's given..
So Jesus is sharing His Father's inheritance by His choice,.
not as some earned income for what these people have done to deserve that inheritance that He's given them..
Yet the fundamental character of those who receive this inheritance and those who won't.
are those who are blessed and those who are cursed, their character couldn't be more different..
The difference seems to center on those who are willing to help and those who are not..
Ones whose hearts are tender toward those who are vulnerable around them.
and those whose hearts have become calloused, whose love has grown cold toward those hurting around them..
Today, we might call it compassion fatigue..

$^{161}$With all the suffering we see and are exposed to in our globally collected community,.
it certainly requires of us great discernment, supernatural power,.
and even seeing God as our source to respond appropriately to all of these situations that we see,.
asking God to help us so that we might help others..
In what ways are we supposed to help?.
Well, we might think about the grand things that we should be doing,.
but really this passage of scripture is super practical..
Pretty much it means just meet basic needs, nothing very extravagant..
Give someone some food or give them some drink if they're hungry and thirsty..
If someone's a stranger, find a way to make them feel welcome..
If they don't have clothing, give them some clothes to wear..
If they're sick, care for them..
If they're in prison, go to see them..
These all sound like really simple actions, but they are essential for life..
But why should we do this?.
Why should we be engaged in this kind of good behavior that we think is here?.
And it's so interesting how Jesus really challenges us to think about it..
He says we should do this because he identifies so closely with these vulnerable people.
that the very act of helping them is helping him..
What are the consequences?.
Well, the consequences are pretty stark here..
The cursed are sent away empty-handed into eternal punishment,.
and the blessed are invited to come to receive the inheritance of eternal life..
Pretty stark differences there..
So let me ask this really as we close here..
I want us to think about making sure that what I'm sharing with you is on good theological footing..
So I want to really draw my conclusions from Baker's Evangelical Dictionary of Biblical Theology..
It sounds a pretty credible source for finding this kind of conclusion..
And it's specifically the words of a Chicago theologian whose name is Scott McKnight..
And he says this about these Bible verses that we're looking at here today..
"Jesus demanded that his followers obey him,.
and that those who were obedient to his will would be called the righteous..
It's also the case that he called his followers to live a life of love,.
love for God and love for others..
As Jesus compared--showed his compassion for others by healing them and helping them,.
so Jesus called on us to do the very same thing..
As Jesus was willing to reach out to help people of all nations,.
so the disciples are charged to make their love transnational in nature..
Thus the foundation of the tension of the church that it had with Israel.
over the symbolic value of the nation of Israel.

$^{201}$may have really rooted Jesus' teaching to love and minister to all nations,.
not just to Israel..
This universal love would be a fundamental factor we see here in the final judgment of God.".
So what's he saying there?.
He's saying, according to Jesus, loving God by loving others is the greatest commandment..
And doing it among the nations may be one of the greatest challenges for the church today..
Jesus said in John's Gospel that the world will know us by our love,.
and our love for God is expressed in how we love our neighbor..
But who's our neighbor?.
Well, fortunately, Jesus answers that question too, doesn't he, in the Scriptures..
In Luke's Gospel, Jesus says that a loving neighbor is best illustrated by the Good Samaritan,.
a person who shows mercy to an injured person he doesn't even know,.
maybe not even respected,.
and doesn't even share his faith as the priest and the Levite did who passed him by..
May I say to you today, you and I are that injured person on the side of the road,.
left for dead in our sins until Jesus showed up and showed us mercy..
He extended a helping hand to us and gave generously so that you might be well cared for.
and enjoy the riches of eternal life..
The big question is this, though..
Should you and I do anything less for someone who's in a similar situation?.
And how do we show our love?.
Very simply, we meet real needs we see in very tangible ways..
Expecting nothing in return until the final judgment when Jesus will make all things right..
So may I suggest to you one way you might respond to this eminent need,.
and an eminent need that's in our own backyard right now as people begin arriving this fall..
So please watch this video and consider what God would have you do..
By now we have all seen videos and photos of the evolving crisis in Afghanistan..
We have heard about the many thousands of people who have been forced to leave behind everything they have or have known..
People who bear the image of God..
If you are anything like me, it breaks your heart and makes you ask, "What can I do? How can I help?".
At Water Libra, we understand, we believe that you are meant to be a force for good right here in Chicagoland..
Over the past 20 years, we have guided churches and individuals to welcome, befriend,.
and help support over 9,000 Afghan immigrants in vulnerable situations in communities all across our nation..
Over 500 here in Chicagoland..
Right now, Water Libra is actively preparing to accept 750 people from all over the world to Chicagoland in the coming year..
There's over 150 coming from Afghanistan..
We need you..
This is how you help..
Pray. Pray for people fleeing Afghanistan and other conflict zones around the world..
Pray for those who have recently arrived here in Chicagoland..

$^{241}$And pray that God will reveal how you should respond..
Volunteer. Gather others from your church to serve in high-need volunteer roles..
The most urgent needs are on our website..
Donate needed items..
Together, we can furnish the apartments of every refugee family who arrives to Chicagoland through Water Libra..
Give. Financial donations provide many of the vital services and ongoing supports needed by those arriving from Afghanistan and other parts of the world..
All of these ways to partner are outlined on our website, along with how to register for Here to Welcome, an ongoing way to learn more and get involved..
We believe in the church..
We believe that we are called to support in moments of crisis like the situation in Afghanistan and the earthquake in Haiti..
We look forward to continuing to share with you how you can help change the world..
There are displaced people all over the world..
In fact, the current numbers that I see are about one person every two seconds being displaced in the world,.
which brings us to almost 80 million people..
People all across Asia, people across Europe..
We see famines that have created movements of people, earthquakes and hurricanes that move people around..
Maybe what we're talking about here locally is not where your heart is..
Maybe your heart is for another group of people who've been displaced, whatever part of the world you're in..
What's God calling you to do as a result?.
How is he calling you to act to be his hands and his feet and really to care for him as you care for others?.
Let me pray for us..
Father in heaven, we thank you so much for the challenge that we have from your son Jesus and how we care and love others..
Lord, we know that the world will know us by our love,.
and so I pray that we would not be characterized as people who make stereotypes and maybe lay aside or turn the other way.
when we see situations that are difficult, but we step into them and we act in ways that you've called us to act..
I love the way in the scriptures at the Church of Macedonia that it says that they helped with what they were able.
and beyond what they were able because of the generosity of their hearts..
I pray, Lord, that our hearts would not be cold,.
but that our hearts would be warm toward those who are vulnerable around us and find ways to help them.
and bring glory to you and serve you as we serve them..
Thanks so much for this day, God..
Speak to us and instruct us, even as we saw at Pentecost, challenge our hearts to the things that you would have us do..
And I pray all of this now in Jesus' name..
Amen..
[BLANK AUDIO].
\newpage



\section{}
\label{sec:9svuWXYcy2M}
\textbf{舊約網上行 - 上帝的足跡: EP16 撒拉的埋葬}
\newline
\newline
連結: \href{https://youtube.com/watch?v=9svuWXYcy2M}{\texttt{https://youtube.com/watch?v=9svuWXYcy2M}} ~~~~ 語音日期: 2021-09-13
\newline
\newline
\hyperref[sec:XIq2yvbK8oo]{\small{< < < PREV SERMON < < <}}
~
\hyperref[sec:index_chronic]{\small{[返順時目]}}
~
\hyperref[sec:index_scriptual]{\small{[返順卷目]}}
~
\hyperref[sec:pfpxXxJmbDw]{\small{> > > NEXT SERMON > > >}}
\newline
\newline
$^{1}$今次我們來到《九曰網》上的第十六講.
又是講到《創世記》第23章.
上一次我們看到《創世記》22章.
獻意殺這個故事.
我們講過阿拉罕殉心的旅程.
幾乎已經去到完結.
之後的兩章.
23章,24章.
其實是記載一些 transition.
一些過渡的內容.
但不代表是不重要的.
今次我們看23章.
就是記載阿拉罕的妻子.
撒拉她的死和埋葬.
這裡也是聖經第一次記載到一個很詳細的埋葬葬禮的故事.
23章第一節.
撒拉享壽一百二十七歲.
撒拉一生的壽數.
她死在迦南地的基列.
阿爸就是希伯倫赫本.
阿拉罕為她哀動哭號.
後來阿拉罕從死人面前起來對黑人說.
我在你們中間是外人是寄居的.
求你們在這裡給我一塊地.
好好埋葬我的死人.
使他不在我眼前.
這裡第一個我們需要問的問題.
就是埋葬到底在當時近東的文化裡面是什麼意思呢?.
或者對於我們現代人來說.
可能埋葬只是搞一些身後的事.
在近東的文化裡面.
在那種社會的氣氛裡面.
埋葬不單單是一些身後的辦理手續.
而是在一個榮與辱的文化.
honor and shame 裡面.
在舊約的世界當中.
葬禮可以分開五類.
由一些最尊貴的.
到一些最羞恥的葬禮方法.
第一種當然就是.

$^{41}$在舊約的世界裡面.
最好的葬禮方法就是自己家族.
一個family tomb 一個自己的墓穴.
然後你就歸到你的先人那裡.
這是最榮譽的一種葬禮.
第二種的葬禮就是一種親屬的墓穴.
或者你未必有錢.
能夠有自己家庭自己家族的墓穴.
你就可以投靠你的親戚.
在親屬家族的墓穴裡面安葬.
或者第三個層面.
就是在別人的家族的墓穴裡面安葬.
由最榮譽的.
一直向著比較沒有那麼榮譽.
比較沒有那麼honor的一種葬禮.
第四種就算是一種.
感覺上是咒詛.
就是一種坑墳或者坑洞.
將你的屍身埋葬.
不能夠歸回家屬或者家人.
然後最底層就是最羞恥的.
真的算是一種咒詛.
很多時候在戰亂或者被擄的時候.
沒有正式的葬禮.
有一點像橫屍街頭.
在一些動亂的情況.
事實上都會發生.
是捍衛一種羞恥的一種咒詛.
我們就想問.
撒拉的葬禮或者阿布拉罕希望撒拉有的葬禮.
是哪一類呢?.
很簡單.
我們應該想到是第一類.
因為二十三章七節開始.
不斷去交代阿布拉罕.
希望能夠買到自己家族的墳地.
在希伯倫這個地方能夠安葬撒拉.
當世紀二十三章第七節.
阿布拉罕向那地的客人下拜說.
你們若有意叫我埋葬我的死人.

$^{81}$使他不在我眼前.
請聽我的話.
為我求所覺的兒子爾弗倫.
把田地上那麥比拉洞給我.
他可以按著足價埋給我.
作我在你們中間的墳地.
當時爾弗倫正坐在客人中間.
對阿布拉罕說.
不然我主請聽.
我送給你這塊田.
連田間的洞也送給你.
在我同族的人面前都送給你.
可以埋葬你的死人.
阿布拉罕卻對爾弗倫說.
你若應允請聽我的話.
我要把田價給你.
求你收下.
我就在那裡埋葬我的死人.
整個故事的對話.
是跟這塊墓地的買賣有關.
就如爾弗倫怎樣形容.
他說給阿布拉罕那塊田地.
不單給墓穴.
還給你整塊田地.
你會覺得很慷慨.
字面上看是很慷慨.
但其實裡面隱藏著很多風險.
因為如果送給阿布拉罕的話.
業主仍然保留權利.
可以隨時收回.
這樣就會損失墓穴和地的權利.
或者這裡的字裡行間.
我們應該這樣翻譯.
送字可以翻譯成賣的意思.
所以整個對話裡面.
不代表爾弗倫這個客人是慷慨.
他說不如我賣給你.
明明阿布拉罕只要求一個洞穴.
一個墓穴.
他說要整塊田地賣給你.

$^{121}$其實是想逼阿布拉罕.
去買整塊田地.
為何會這樣做呢.
如果參考當時同期的客人的地產.
墓穴,田地的買賣法律條文.
我們就特別看到.
原來關於這些房地產買賣.
哪一個做那塊地的地主.
就要幫那塊地去納稅.
向城邦的王納稅.
或者向當地的政府納稅.
所以這個爾弗倫.
如果以這樣的背景來看.
他就是希望逼阿布拉罕.
不單止買一個洞穴.
雖然阿布拉罕需要的.
要求的只不過是一個洞穴.
去埋葬撒拉.
但爾弗倫就覺得這樣很吃虧.
不如把整塊地都賣給阿布拉罕.
以至整塊地的產業.
都是屬於他的.
但稅也是歸於阿布拉罕.
來到第十四節.
就是爾弗倫回答阿布拉罕.
他說我主請停.
即四百赦克勒的銀紙的一塊地.
在你我中間還算什麼.
只管埋葬你的死人吧.
阿布拉罕就聽從了爾弗倫.
照著他在客人面前所說的話.
把買賣的通融銀紙.
平了四百赦克勒.
給了爾弗倫.
所以作為一個寄居者.
他比較處於弱勢當中.
他就被迫不單止買一個墓穴.
也要買一整塊的田地.
第二個問題.
我們就問四百赦克勒銀紙.

$^{161}$到底是貴還是便宜.
這些很多時候我們現代人讀經的時候.
我們非常關注的.
價錢是四百赦克勒.
我們知道在近東的社會裡面.
到了波斯時期.
才會有通行的貨幣.
之前古人是用大大小小的銀錠.
一粒一粒一塊一塊的.
他們做買賣的時候就要秤重.
重量重量的赦克勒.
這個赦克勒其實就是.
一些度量衡的重量單位.
每一個赦克勒.
大概是11.33克.
所以四百赦克勒的話.
這裡就說的是4.5千克.
大概是十磅的純銀.
按照一些計算.
現在的話.
這個務月和加上整塊田.
就是62萬4千美元.
所以以弗林對阿布拉罕說的是什麼呢.
60多萬美元的田地.
在你我中間還算什麼.
我們聽起來可能覺得很昂貴.
但是按照當時的經濟條文.
或者一些文件.
這一塊地不算很便宜.
也不算很貴.
但是算是可以接受的.
第三方面我們可以問的.
就是經文的重點在哪裡.
我們起初說的23章.
就是撒拉的死亡和埋葬.
但是如果留意智慧行間.
整段的經文.
雖然撒拉的去世.
是一個很哀傷的一段時間.
但是經文是輕描淡寫的.

$^{201}$真的說的其實不多說.
埋葬的葬儀的葬禮的條文.
智慧行間的引言或者結論.
給我們看到了什麼呢.
重點在哪裡.
我們看第17節.
於是麥比拉萬里前.
以弗林的那塊地.
和其中的洞穴.
並田間四處的樹木.
都定準了歸給阿布拉罕.
乃是他在客人面前.
並城門出入的人面前.
來賣妥的.
等等的經文.
告訴我們重點.
真的不是說如何哀悼撒拉.
也不是怎麼說埋葬他.
而是整段的經文.
都是在說地產的買賣.
一塊墓地.
田地的廣價協議等等.
在城門口發生的.
城門口就等於今天的律師樓.
如果你去過.
買過一些房地產.
買過樓的話.
你closing的時候就要去律師樓.
就要簽名.
簽到你手軟.
有很多很繁複的手續.
其實這段聖經.
大部分的細節.
都是在說這件事.
在城門口.
如何去買妥.
交易清楚.
所以有很多章節.
都是在說阿布拉罕.
如何去買這塊地.

$^{241}$為什麼這麼著重這件事呢.
你記得土地的認許.
是在阿布拉罕整個故事.
其中一個最重要的主題.
第一次出現土地的認許.
就是十二章第七節.
當時阿布拉罕回應耶和華的呼召.
離開本地本族本家.
去到迦林地的時候.
第七節.
要向阿布拉罕顯然說.
我要把這地迦林地賜給你的後裔.
然後在不同的經文也提過.
十三章十七章.
在阿布拉罕的故事最後的時候.
到現在這一章.
再說清楚一點.
因為說到神的認許.
真的進一步的實現.
因為這裡說的是一個商業的交易.
買到一塊地.
是族長阿布拉罕.
以色列人的所謂國父.
得到第一塊地.
擁有一個家族的墳墓.
當然對當時來說.
是非常重要的.
因為一個家族的墳墓.
就是為了將來的下一代.
下一代而設的.
舊約世界裡面的墓穴.
是一些石山挖出來的.
屍首存在墓穴當中的時候.
當屍身腐壞之後.
一年之後只剩下一些骸骨.
家人往往會回到墓穴那裡.
然後將這些骨頭收藏.
在舊約的世界.
第一聖殿的時期.
這些骨頭會儲存在一個角落.

$^{281}$或者儲存在下面的坑裡面.
一個儲存坑.
所以舊約聖經裡面.
經常有一句俗語.
叫做歸回你的列祖.
世世代代延續.
可能這就是墨比拉的意思.
這個洞叫做墨比拉.
因為墨比拉的意思是雙倍.
不是只為了你這一代.
而是為了下一代.
甚至再下一代.
所以之後土地的認許.
不斷地在創世記出現.
族長過了阿拉罕之後.
以撒雅各也有一個這樣的認許.
想說歸回列祖.
重新安葬在這個洞.
代表上帝沒有忘記他的認許.
沒有忘記他對你說過的話.
在創世記最結尾的時候.
約瑟到了埃及.
大家記得.
他期盼的是什麼.
是回去重新安葬在迦南地.
一直這個土地的認許.
不斷地一步一步地.
上帝守著這個認許.
直到約書亞記.
以色列人已經進入迦南地.
21章第45節約書亞記.
就是耶和華的認許.
致福給以色列家的話.
一句也沒有落空.
都應驗了.
他們已經進入迦南地.
所以在這裡.
剛剛開始有土地認許的第一步實驗.
就是告訴我們.
阿拉罕信心之父.

$^{321}$雖然他快要離開這個墓台.
但是他買了一塊地.
剎下了一粒種子.
就是一粒信心的種子.
你說阿拉罕在迦南地擁有第一塊地.
到底怎樣去形容.
關我們什麼事呢.
是不是說希伯倫這個墨比拉洞.
是最神聖的地方.
因為第一塊地就在那裡.
所以今天如果我們霸佔這個地方.
我們擁有這個地方的時候.
就好像上帝一早寫了地契.
現在我們拿回了.
這個地方永遠都是屬於我們的.
這個地點是特別輕鄉一點的.
希伯倫的墨比拉.
只要有土地的霸權的話.
我們就心安理得.
這個的確是今天一些比較激進的分子.
他們怎樣看希伯倫墨比拉洞這個地點.
有些很激進的猶太人.
可能會覺得這個地方.
真的很異常的神聖.
早在二千年之前.
我們今天見到這個墨比拉洞的地方.
大希倫當時統治的時候.
已經在這裡建立了一個很大的石建築.
就是包著裡面一個山洞.
紀念這個族長墓地的地方.
當然我們不能夠去證明.
還有你找什麼標準去證明.
根本沒有什麼可能.
但是裡面今天你進去的時候.
就有不同的座標.
紀念到族長每一位先祖墳地葬的地方.
在近代的歷史.
這個地方以巴兩方.
發生了很多很嚴重的衝突.
死傷不少.

$^{361}$所以為什麼今天這個建築物.
就分成了兩半.
一半就成為猶太人的會堂.
另一半就成為伊斯蘭教徒的清真寺.
那你說是不是經文就是這樣的意思呢.
誰霸佔這個地方就特別輕香.
就好像上帝在創世記二十三章.
寫了一封地契.
給我們今天拿著.
我們就可以說擁有這個地方的權利.
或者我們應該這樣想.
到底這個故事.
創世記二十三章.
買到第一塊地.
對於舊約既斯利人來說.
是有什麼關係呢.
我們按照整個的文理.
上文下理和脈絡.
我們可能會看得比較清楚一點.
我們說過阿巴罕的故事.
是十二章開始.
到二十二章.
他的信心最高峰.
在猶太人的拉比文學裡面.
其實也有提及過.
就是說在阿巴罕整個信心的旅程裡面.
裡面的故事告訴我們.
上帝去試驗阿巴罕的信心有十次.
他試驗了這些信心之後.
然後才到第二十三章.
就是我們今天看的這段經文.
就是阿巴罕買入第一塊地.
承受地土.
第一次一個initial的fulfillment.
土地的認許.
初步的實現了.
意思就是說.
從脈絡我們知道.
怎樣才能夠留在這個地方.
能夠承受地土.

$^{401}$能夠拿到這塊地.
就是你需要有一個阿巴罕的信心.
就是十二章到二十二章所交代出來的.
就是你需要留在這個地方.
你可以享受這個地方.
能夠與神同在在迦南地.
向萬國能夠見證耶和華.
使萬國能夠得福.
先決條件是什麼呢.
就是純.
有阿巴罕那種純.
一種沒有條件.
願意去跟隨到底.
堅信到底的純.
才能夠使以色列人在迦南地.
能夠日子長久.
你甚至整個聖經來看.
你就發現舊約裡面.
重要的人事預物.
包括祭司.
聖殿.
祭禮.
甚至所謂基業土地.
迦南地得到的地權.
其實都是後事的影移.
因為這些人事預物在舊約.
都是指向將來.
一種好像路標一樣.
好像你還沒有去到將來的目的地.
有的相片.
就是指向一個真相.
真相什麼時候出現呢.
就是耶穌基督來的時候.
利塞亞救世主來到.
耶穌基督就承受了.
實現了舊約裡面.
所有這些的應許.
祭司最偉大最完美的.
就是耶穌基督.
祭禮最完美的.

$^{441}$就是耶穌十字架上所獻上的.
以前聖殿.
就是神與人去到連繫.
在那裡敬拜神的地方.
今天敬拜神就不是說地方.
是在基督裡面敬拜神.
以前是要有信心才能夠享受到.
Life in the land.
現在要有信心才能夠享受到.
Life in Christ.
在整本聖經一脈相承的.
都是告訴我們.
從亞伯拉罕的故事開始.
就是有一種這樣的信求.
去堅信,爭取到底.
才能夠承受到.
上帝給我們的同在.
能夠住在上帝的生命裡面.
今天我們在基督裡面.
不是這樣.
我們一定要有信心.
我們才能夠經歷到.
上帝給我們的盼望.
才能夠經歷到上帝給我們的平安,喜樂.
各道裡面的生命.
《墨落上文》.
看了亞伯拉罕整個故事的《墨落上文》.
和看了新舊約整個聖經的大圖畫.
我們才能夠更加明白.
這段舊約的聖經.
提到亞伯拉罕買墨比拉洞.
第一塊土地的時候.
對於我們今天的信徒.
我們每一個世世代代的教會.
有什麼應用的地方.
提到亞伯拉罕要憑信.
他才能夠得到迦南地的第一塊土地.
對於舊約的記者來說.
對於我們今天的人來說.
一墨相承.

$^{481}$都是在說一種無條件的忠誠到底.
跟隨耶和華.
我們才能夠住在神的生命當中.
《穆乖府希伯來書第十一章第八節》.
那裡這樣說.
他說恩爵信.
亞伯拉罕在蒙召的時候.
聽命往他所要承受的地方去的時候.
他出去的時候.
都不知道自己要去哪裡.
恩爵信.
他在認他之地寄居.
保障在異鄉.
與承受同樣認他的其他族長.
以撒雅各一樣.
都好像寄居一樣住在象棚.
他們一路整個生命都是寄居.
整個生命都是期盼著.
以信來盼望著.
將來上帝要承受的東西.
就是要等候著一個有根有機的城.
這個城就是神所帶來.
神所設立的.
不是地上的一個地方.
而是將來新天新地.
神的應許完全降臨的神.
這個就成為我們歷世歷代的信徒.
無論經過我們每一個的時代.
是順利也好.
是很艱難也好.
我們有一個不會改變的期盼.
這個期盼.
這個盼望.
能夠推動我們不斷來恩爵信.
繼續去跟隨.
繼續去相信上帝.
好像希伯來書第十一章.
那裡講到信心的英雄.
歷世歷代的舊約的英雄.
他們為什麼可以被上帝捍衛.

$^{521}$是榮譽呢?.
因為他們有那種信心.
你都可以期盼著上帝.
他們的國度.
他們的應許將會.
將來一定會實現.
我們今天身處在這個世界.
無論是因為有疫情.
或者你生命裡遭遇不同各樣的挑戰.
不同各樣的試煉.
我們仍然可以有期盼.
我記得在我讀神學的時候.
記得有一次組會.
一個神學教授.
他一篇講道.
他的題目是講什麼呢?.
他說How can we sing a happy song in a sad world?.
是一個充滿苦難.
充滿悲哀的一個世界.
我們為什麼可以敬拜上帝的時候.
我們是sing a happy song.
不是因為我們視而不見.
我們對世界的苦難.
或者對世界的不義.
我們掩面不寒.
我們不理會.
總之我們躲在一起.
我們很開心的.
我們自嗨.
不是因為這樣的緣故.
而是因為我們有恐.
有盼望.
我們的期盼是按照上帝曾經所講所應許的.
在基督裡面已經實現了.
並且將來新天新地.
基督再回來的時候.
會完全實現.
因著這個期盼.
因著這個盼望.
我們深知上帝仍然在推動他的計劃.

$^{561}$仍然在應驗他所講過的應許.
他必定會帶來公義真理的國度.
因為上帝所講的.
我們不需要去擔任.
我們有這個期盼.
就算是今天.
我們預見的.
是多麼不順境的時間.
我們依然可以像這些信心的偉人.
繼續的向前.
再看看阿巴拉罕根隨神整個的旅程.
由他75歲.
到回應上帝的呼召.
到他死的時候175歲.
創世記25章第7節.
有100年一個整數的時間.
整個的時段當中.
我們見到他以信心來根隨神.
也被上帝多方來試煉.
在他生平當中.
讓我們見到的.
是一個真正信心的本質.
在他整個的故事當中.
他最後只是實現了.
上帝整個計劃裡面的一小部分.
只是初步.
整段的經文創世記23章.
阿巴拉罕都說自己.
只不過是一個外國人.
一個外人.
如果你這樣看的時候.
看回整幅上帝整個計劃的圖畫.
阿巴拉罕只是應驗了一小部分.
實在有時我們在神的計劃當中.
在整體的工作裡面.
我們真的只是做一小部分.
真的.
有時我們的角色.
我們的時間.
我們的崗位.

$^{601}$在上帝不同的處境.
安排我們不同的時段裡面.
我們只是上帝整個藍圖裡面.
佔一小部分.
但其實這個不是最重要.
最重要的是我們是否忠於.
今時今日這個照明.
也不一定在我們今生.
可以看到上帝的計劃成就了.
或者我們只看到上帝的計劃.
成就的一小部分.
阿巴拉罕並不是這樣.
他只是看到整個救約的應許.
的一小部分.
最後的時候只得到一塊墓地.
但在上帝眼中.
就是一個信心的見證.
巴伐希伯來書十一章十三節.
結束的時候看看.
神如何看待阿巴拉罕.
如何看待我們只要忠於所信.
忠於我們今時今日的呼召.
在上帝的計劃當中.
我們成就一小部分.
神如何看待.
他說這些人都是存著信心死了的.
還沒有得著所應許的.
只不過是從遠處看見就表示歡迎.
又承認他們在世上是異鄉是客去.
表明他們所尋求的是另一個教養.
一個更美在天上的教養.
所以神不以他們傾他為神.
而覺得羞恥.
神覺得很honor.
很honor能夠有這些人.
來稱他為神.
但願我們今天都是一樣.
見到阿巴拉罕的信心.
我們就是學校這種信心.
忠於所同.

$^{641}$今時今日.
上帝所安排我們的崗位.
我們忠於所同.
做好我們自己的一小部分.
讓上帝自己得榮耀.
我們同心低頭祈禱.
上帝求你親切幫助我們.
在我們今時今日這個時代當中.
你所給我們的呼召.
你所安排我們有的時間.
甚至我們的受數.
讓我們忠於所同.
忠於你對我們的呼召.
求你帶領我們.
在我們世界各地不同的崗位當中.
不同的角色裡面.
我們以信心去仰望你.
\newpage



\section{}
\label{sec:pfpxXxJmbDw}
\textbf{Anointed by God - Part 1 (1 Samuel 17) - Tommy Lee}
\newline
\newline
連結: \href{https://youtube.com/watch?v=pfpxXxJmbDw}{\texttt{https://youtube.com/watch?v=pfpxXxJmbDw}} ~~~~ 語音日期: 2021-09-21
\newline
\newline
\hyperref[sec:9svuWXYcy2M]{\small{< < < PREV SERMON < < <}}
~
\hyperref[sec:index_chronic]{\small{[返順時目]}}
~
\hyperref[sec:index_scriptual]{\small{[返順卷目]}}
~
\hyperref[sec:JJV1ojc4_5E]{\small{> > > NEXT SERMON > > >}}
\newline
\newline
$^{1}$[Music].
All right. Thank you very much..
Well, good morning from Chicago, everyone..
For those who are listening from wherever you are,.
I hear some of you are in Hong Kong..
I hear from some of you guys are in Canada, different places..
And those who are in the audience, good morning and good evening..
My name is Tommy Lee, and I get a chance..
I run multiple organizations, but one of the things I am very proud of is.
I am the president and founder of Resource Global..
And one of the things I get a chance to do at Resource Global is.
we work with young Christian marketplace leaders in different global cities..
And so as a result of that, I get a chance to go to Jakarta, Indonesia,.
when pre-pandemic, in Kuala Lumpur, in Singapore, as well as Austin, Texas, Chicago..
And I was recently just in Nairobi, Kenya..
So I literally just came back from there last week,.
and I had a chance to really be able to see what God was doing in Nairobi, Kenya..
And it is phenomenal how different global cities are continuously grown..
I'd love to share a little bit more in the future in terms of what Nairobi is,.
but the one thing I will probably see, there is a lot of growth.
in countries like Africa as well as Asia..
But the one thing that it is, is what does it mean that as the countries continue to grow,.
how do marketplace leaders, how does the local church continue.
to really be able to proclaim what God is doing in those cities.
and interact with those cities..
So, man, it has been just a joy and journey..
And with my role, I get a chance to travel a lot..
And it's not always easy..
A couple of weeks ago, it's my fourth straight week traveling, right?.
It's my fourth straight week traveling. I'm tired..
A couple of weeks ago, I was in Grand Rapids..
Tomorrow, I jump on a plane, I go to Los Angeles..
Just got back from Nairobi, and it has not been always very easy.
when you're jumping on planes all the time..
But in the midst of all these things, I get a chance to meet wonderful people..
Let me just show you one person, all right?.
Here's this one person I want to talk about..
And they are going to put that right there..
His name is Ray Bakke..
I'm not sure if you guys have ever heard of the name Ray Bakke..

$^{41}$Ray is someone--it ties in with our sermon today--.
what I would call as God's anointed, God's professor, global professor..
A couple of weeks ago, maybe about four or five weeks ago,.
I get an email from Ray's assistant saying that,.
"Tommy, Ray has made a turn for the worse..
Ray has pancreatic cancer. Ray is dying.".
As a result of that, he asked if you had time to go see him.
in Bellingham, Washington for one final meeting..
Well, here's the thing..
If someone tells you that they would like to see you for the final time,.
you do everything you can to go out there..
I have gotten to know Ray for many, many years..
He is a global professor that has ministered and taught many students.
all over the years, traveled to Hong Kong..
Some of you are sitting there watching from Hong Kong..
You may have been familiar with Ray and his travels and his teachings..
But the one thing Ray has always taught me as a professor is,.
"Tommy, you have to learn to what?.
Respect and honor the culture that you're in..
You have to learn to 'exegete to community.'".
What does he mean by that?.
Too often--and I say this as an American standing here--.
too often as an American, when I travel to your city in Hong Kong.
or in Canada, all of that stuff, as Americans, I feel like I have all the answers..
As Americans, I feel like you have to learn from me..
As Americans, I'm your primary funder. I'm the leader..
You learn from me..
Whereas some of you who are sitting in Hong Kong or Vancouver.
or wherever you are in the world are sitting there,.
"Wait a minute. That's not right..
You fail to understand our culture..
You fail to understand--you bring your American ideas.
and you expect all of us to learn.".
That was the one thing Ray has always taught me..
He says, "Tommy, when you go to different global cities,.
you are a learner of that city..
You are a student of that city..
God has been at work in that city way before you ever got there,.
way before any Americans ever got there..
He has a plan for that city..

$^{81}$And so when you go to that city, listen and hear..
Ask questions. Honor the people there. Honor the culture there..
Always be respectful of the culture..
Hear what God is doing and always be asking yourself,.
"How can God be working in your life?.
How is God working in the city and how can you be part of it?".
That's what Ray taught me, and I will always remember that from Ray..
Ray is what I call God's anointed,.
and it's a position that all of us are in today..
And this is a topic we're going to talk about actually.
for the next three times I'm with you,.
and I'm going to be with you every single month,.
this September, October, November,.
and we're going to talk about a three-part series on the life of David..
Today we're going to talk about a story that all of you are familiar with,.
the story of the battle of David and Goliath,.
which I would say if you're young or old,.
you're probably familiar with this story..
It is a story, I would say, probably aside from the crucifixion.
to the birth of Jesus, it's a story that we're all familiar with..
The story, the battle of David and Goliath,.
good versus evil, underdog versus the giant..
But it's more than just a battle between little man versus giant..
But it talks about how God's anointed, God uses his anointed.
to continue to call him out and use him,.
begin to prepare him for the biggest stage of being the king of Israel..
And so we're going to talk about part one of that.
and lessons that emerge, principles that we're going to talk about.
that emerges of being God's anointed..
We're going to talk about part one..
We're going to talk about the next time I show up,.
we're going to talk about part two as we finish up the rest of the story..
And then part three is one of my favorite ones is what happens.
in the aftermath of a giant's killing..
Because you think that after the giant's killing, story ended, happy ending..
But it's far more than that..
And God is in the midst of preparing a king..
God's anointed with David..
It's a story of you as well, too, because God has a plan for you..
God chooses and wants you..

$^{121}$We all have the potential to be anointed by God to do something great..
We all have the potential to be anointed by God to make a difference.
and impact this world because he is in the midst of writing your story..
He is in the midst of writing a chapter in your story,.
a story that is so unique that my story is so different from your story,.
that all of our stories are very, very unique..
And so as a result, within that journey, what is that story?.
How is God calling you?.
How is God making a difference in your life?.
And that's what we're going to look at today..
As always, I will probably say is as you're getting to know me,.
you will find that I will need a lot of water in the midst of it..
And so I always will tell people, please be patient with me during that period of time..
I think I told you in the last time I had a sermon, I spoke to you guys..
Cancer really affected my body..
Three years ago, I was diagnosed with cancer..
And in the midst of that cancer, right in my sinus,.
it pretty much destroyed all my tears, most of my hearing..
That's why a lot of times if you're here, I speak very loudly.
because a lot of times I can't really hear..
But more so, it really took away most of my saliva and most of my taste buds..
And so as a result, it is a part of that journey that God has written..
And so a lot of times I will ask you for a little grace as we continue on as well too..
Hey, let's take our story over to 1 Samuel 9 as we begin..
1 Samuel 9, and we're going to begin our story with that..
And we're going to go over to 17 eventually, but I'm going to start with 1 Samuel 9..
And the story begins in verses 1 to 2..
"That there was a man of Benjamin, whose name was Kish, the son of Abel, son of Zerah, son of Bekroth, son of Ephiah, a Benjamite.".
Very, very important. Remember who that is..
He is from the tribe of Benjamin. He is a Benjamite..
Always remember that. We're going to get back to this..
"A man of wealth. And he had a son whose name was Saul, a handsome young man..
There was not a man among the people of Israel more handsome than he..
From his shoulders upward, he was taller than any of the people.".
So the story begins with the selection of Saul as the leader..
He is a Benjamite. Very key. Comes from a very influential, very, very wealthy family..
He comes from a family of means..
He is the tallest. He is the most handsome. He is the most athletic and natural because of his outward appearance..
Coming from a very, very lucrative family, very well-to-do family, he is chosen king of Israel..
And the story goes, and we don't have time to go through all this, you know what happens with Saul..

$^{161}$Through a series of mistakes, that ego starts getting in the way..
The character flaws begin to show within Saul..
And God deems him as no longer the right fit as a leader..
As a result, he is looking for a replacement..
There is a prophet Samuel. He says, "Samuel, I have disqualified Saul..
Even though he is elected by the people, even though all of his outward appearance, it checks the boxes..
I have disqualified him as king of Israel..
And so I need you to go find the next king..
And I want you to go to Bethlehem to go do it.".
And so Samuel goes over to Bethlehem, and he is now called to the son of Jesse, the Bethlemite..
And there Jesse comes in..
And he arrives and he says to Jesse, "Bring all your sons in. I want to talk with them. I want to see them.".
And the story goes, he brings in three of his sons..
And the first one to show up is Eliab..
Well, Eliab is the oldest of all of the sons..
Eliab is the strongest, the tallest of all the sons..
And here's what 1 Samuel 16 tells us..
"But the Lord said to Samuel, 'Do not look on his outer appearance, on the height of his stature,.
because I have rejected him..
For the Lord sees not what a man sees..
Man looks on the outward appearance, but the Lord looks on the heart.'".
So he goes through all the sons, not him, not him, not him..
And he says to Jesse, "Jesse, do you have any more sons?".
Jesse remembers, "Oh, there's young David.".
We all know the story. There's young David..
He's out in the sheep. He's tending the sheep..
"Hey, call David in.".
David walks in..
"There Samuel anoints him with who watching is, his brother, including Eliab.".
That's also very important. Note that as well, too..
We don't know exactly how he anoints him, but he anoints him..
Maybe he whispers to David, "You're going to be the next king.".
And there the story ends..
First thing, first principle that emerges out of God's anointing..
God always looks at the condition of the heart..
God does not look at outer appearances..
God does not care how wealthy you are..
God does not care how good-looking you are..
God does not care how rich you are..
God does not care how successful you are..

$^{201}$God does not care how well you are, speaking-wise, your talents, whatever it is..
God does not look at the physical appearance..
But he always looks at the condition of your heart..
God is looking to use men and women wherever you are, whatever part of the world that you are..
God is looking for men and women with hearts of humility..
God is willing to look at individuals with hearts of generosity, others first, integrity..
Men and women who are willing to stand up, who are willing to give up their lives to serve the Lord..
God is looking always, always with the heart, and the condition of the heart versus the outer appearance..
Why? Why not outer appearance?.
Because, simply, he's created you..
Number one, he's created you..
He's given you all your strengths..
He's given you and placed you and given you the riches of your life..
He's given you the talents that you have..
He's given you the position you are, how wealthy you are, whatever it is..
God has given you all that, so that's number one..
Number two, if you base your skills, and God looks at your skills, there's something you can brag about..
There's something you can say, "Look at what I can do.".
And that's not it..
Ephesians tells us, "For it is by grace you have been saved, not by yourselves.".
Why? So you can't boast..
God looks at the heart of the individual..
And never about your skills, never about your strengths, never about the things that you have, all of these different things..
Let me give you a couple of stories and illustrations of what I mean by this..
There's a pastor that I get a chance, so if you don't know me, I mean, instead of running five different organizations, I get a chance to run five different organizations, right?.
So instead of one, I run five different organizations..
But one of the jobs I get a chance to do is I get to work with different ministry leaders on conferences that are held throughout the world..
And in the midst of all that, I get a chance to interact with pastors, speakers that have books, radio, that have great podcasts, that have a pretty well-known, and that's either good and sometimes bad..
I've met different people, I say this, who have written books that I said, "Wow, I can't believe I met this person..
I want to meet with them and realize, oh, wow, they're not exactly what I thought they was.".
So one situation comes in..
I had this conference I was overseeing..
And this pastor comes in, and he's very well-known..
He has a large social media following..
He has book deals. He's always on the speaking circuit..
He had to do, for this particular conference, he had to do one workshop, and he had to do one 10-minute devotion..
That's all he had to do..
And so as a result of that, he then charges \$5,000 in expenses for his time..
He traveled first class..
He won a large honorarium..

$^{241}$He didn't like the hotel that we put him in because he thought it was not nice..
He checked himself to a \$500 hotel..
He started eating expensive meals, and suddenly his assistant sends me his receipt..
And for one day and a half, I'm looking at the expenses, \$5,000..
And so naturally, I don't have the budget for this because I wasn't authorized..
So I go back to his assistant and says, "I can't do this..
We're going to have to reject it.".
Again, very talented..
Huge book writer. Huge following. Everything like that..
Calls me up one day, and it's a Sunday now..
I'm in Houston. I'm at an event..
Calls me up and says, "Oh, why is this guy calling me? Coming from church?".
He says, "Tommy, let me just say this to you..
You have disrespected me..
Highly disrespected me because you're questioning my integrity and questioning my finances.".
I said, "Yeah, you're right. I'm questioning your finances.".
He says, "Here's why you're doing this..
Because your life stinks..
You probably live a miserable life..
Your kids probably hate you..
Your wife probably hates you..
You probably have no friends..
And that's why you're disrespecting me..
So I want you to do something like this..
I want you not only to pay my reimbursement,.
but I need you to go on social media and apologize for disrespecting me,.
or I will make your miserable life more miserable than it is right now.".
God honors the heart..
It doesn't matter what skills you have..
It doesn't matter what you do or not able to do..
God honors the heart..
Second story that I have..
I know an elder of a church..
Very strong, very, very involved..
Involved in finance, all of that stuff..
Very dedicated to all of this stuff..
But inwardly, he and his wife have problems all the time..
Why?.
He can't control his anger..
Because behind closed doors--hey, he knows what to do on a Sunday, right?.

$^{281}$He knows what to do at churches on a Sunday..
He knows what to say..
He's there on Sundays..
He's there helping out with finances..
He's on the elder board..
But his wife says to me, "Behind closed doors.".
He yells and screams to the point sometimes the neighbors are ready to call the police..
He cannot control his anger..
He belittles..
He's never happy..
God sees this..
We think no one has seen all these things..
But God sees it..
Because God is always after men and women whose heart belong to him,.
whose heart are full of integrity,.
whose heart is devoted to him..
God always looks at the heart..
Never, never on the outer appearances..
One last thing I just said to you..
I just came back from Nairobi..
When you go to Nairobi, you see roads galore..
You see overpasses..
You see new buildings being built, everything..
It's a brand-new city..
But then when you talk with the people there, the city is poor.
because it's ruined by corruption..
Because you have people who are what they say are Sunday Christians.
who all go to church..
Eighty-five percent of the population in Nairobi says they're Christians..
They all go to church on Sundays..
The president, the politicians all go to church on Sundays..
Monday through Friday, they do whatever they want..
And so as much as the money, which by the way, a lot of the money.
and all the changes in Nairobi is funded by who?.
The Chinese..
And as a result of the Chinese, they poured in all of this money..
And the politicians in Nairobi who are supposedly Christians.
are benefiting from all this, but that money is not going to all the city..
And so you have people who are poor..
You have people in slums..

$^{321}$You have people who are struggling..
But only a few people are getting wealthy..
God's anointed for all of us..
God chooses to anoint you..
You have the potential, wherever you are, to sit in there,.
to be anointed by God to do something great..
You really do..
What is the state of your heart?.
Where is your heart?.
Is God working your heart?.
I don't care how talented you are..
I don't care how wealthy you are..
I don't care how educated you are..
I don't care how well you speak..
God cares about that heart..
Let's go on..
The story goes..
And the story goes to verse 7, chapter 17..
And it goes to verse 1..
And we're reading from the ESV..
And now the Philistines gathered their armies for battle..
Go into 17..
And they were gathered, or Soca..
Now the story goes that he's anointed David..
And now there's a battle which belongs to Judah..
And a camp between Soca and Ezekiah in Ephesus, the Mim..
And Saul and the men of Israel..
Now Saul was the king now of Israel..
We're gathered..
And they camped in the valley of Elah and drew up a line of battle.
against the Philistines..
And the Philistines stood on the mountain on one side..
And Israel stood on the mountain on the other side with a valley between them..
And there came out from the camp of the Philistines a champion named Goliath the Gath,.
whose height was six cubits, or over nine feet tall in his span..
Let me give you a little background on this in order to understand fully this story..
Take a look at this, right?.
Let me give you a background on this..
Here's-- and if you guys get a chance, over there, there is a PowerPoint over there..
Well, the PowerPoints over here that I have is-- thank you very much..

$^{361}$Look at the two descriptions..
I gave-- the left-hand column is the Philistines..
The Philistines were coastal people..
They were sophisticated cultural people..
That was the Philistines..
The Philistines had learned how to use iron to make weapons and tools..
They were a master of iron..
They had a monopoly of iron..
But contrast that to the Israelites..
The Israelites were primitive people..
They were composed of shepherds and farmers..
They had no experience in working with iron at that time..
The Israelites were so, in some sense, appeared so inferior to the Philistines.
that they lived in the mountains because the fertile valleys were dominated by the more powerful Philistines..
And, at the same time, here's something important to remember..
Remember what the Philistines are key at..
They were great blacksmiths..
They were great at iron..
Let me give you a point..
1 Samuel 13, 19, all right?.
Here's the way it reads..
Not a blacksmith could be found in the entire whole land of Israel.
because the Philistines had said, "Otherwise, the Hebrews will make swords or spears.".
So all Israel went down to the Philistines to have their plow shares,.
maddox, axes, and sickles sharpened..
The price were two-thirds of a shekel for sharpening plow horses and maddox.
and a third of a shekel for sharpening forks and axes and repointing goats..
So on the day of battle, not a soldier with Saul and Jonathan had a spear or sword in their hands..
Only Saul and his son Jonathan had them..
What is the picture that's being painted?.
Understand the context is very key..
Philistines came from Gath, had all of the supply of swords and iron..
They were masters at that..
Israelites were the primitive people who lived in the non-fertile valleys in the mountains..
They were farmers. They were shepherds..
The only one who had weapons was Saul, and Jonathan had the swords..
Because why? The Philistines had a monopoly on everything..
So now you're asking, as we read into it,.
so why in the world are the Philistines invading the Israelites?.
Are they just doing it for fun, or what's wrong with it?.

$^{401}$Let me show you a series of three pictures, alright?.
Here's a series of three pictures right here..
If you guys could see all these things, here's the first picture..
You see the Israelite camp, and you see the Valley of Elah..
If you go to Israel right now, the Valley of Elah is still there right now..
The Israelite camp is where the mountains are, and you see the Philistine camp..
Let me actually give you a closer picture as well, too..
You see right there, Gath is to the west of the Philistines, right?.
You see the Israelite camp, again the Valley of Elah, and also the Philistine camp..
Let me actually then now go to a third map, and now you'll understand why this battle was taking place..
Take a look at this..
Right in the middle, you see where Ezekiel and Soca was?.
That's where your Philistines and the Valley of Elah and all of that are..
You see Gath right there..
What do you see?.
The Israelite is right there in the mountain, and where does it lead?.
It becomes a gateway to what?.
Look up north..
You see Bethlehem, you see Jerusalem, you see all of Bethel, all of these different places, Shiloh..
It is the gateway to the Jordanian Valley..
It is a key strategic point..
The Assyrian Empire at that point in time, many of the empires went through this gateway to make their trades..
What does the Philistines know?.
That as much as they have the fertile valley, if they can control that gateway, it opens up the trade routes for them to control the gate routes..
It's like the tax collector Matthew, right?.
So they need to control that gateway..
And who's standing in the way?.
The Israelites..
So as a result, let's bring in everyone and get rid of the Israelites..
And you know what?.
We're such a powerful army..
We're not even going to waste our time..
Let's not waste our soldiers..
Let's not have anyone get hurt..
Let's send in a big guy, Goliath..
Let's have him battle their biggest one because it probably doesn't stand a chance..
And whoever wins, if we win, we'll leave..
Which we probably are guaranteed that we won't win..
If we lose, we leave..
If we win, you're out..

$^{441}$And so Goliath shows up, and boy does he show up..
For 40 days, he comes in and out with the same challenge..
Bring your best soldier..
Bring your best hero and let him battle me and let's see who wins that out..
All for this particular valley..
And now we go to 14 to 18..
David, in verse 14, was the youngest..
The three eldest followed Saul, but David went back and forth from Saul to feed his father's sheep at Bethlehem..
And for 40 days, the Philistines came forward and took his stand morning and evening for 40 days in verse 14 and now 16..
And Jesse said to David his son, "Take for your brothers an epath of his parched grain in these ten loaves.
and carry them quickly to the camp to your brothers..
Also take these ten cheeses to the commander of their thousands..
See if your brothers are well and bring what?.
Some token from them.".
Which, you're going to actually, that plays an important role for them..
"Bring some token for them.".
If point number one is that God chooses to look at the heart of the individual,.
when God anoints an individual, he looks at the heart of the individual and never the outer appearance,.
then number two, you never know when God comes calling..
You never know when God comes calling..
You never know when God is ready to call you..
You never know when that chapter in your book turns..
You never know what adventure he has..
You never know what plans he has for your life..
And so we approach our life with open hands..
We approach our life with great anticipation..
We approach our life sitting there each and every single day saying,.
"Lord, what is it that you want me to do?.
How can I continue to serve you each and every single day?".
Because you never know when God comes calling..
You never know when he has a plan for you..
So you live with an anticipation because why?.
And we always talk about this because you are a steward of what God has given you.
because your life is not yours, but you are a steward of what God has given you..
And when we are stewards, we are always anticipating to serve him in every aspect of our life..
It is a story of Nicodemus when he is called by God and he interacts with Jesus..
He never knew that through that interaction in John chapter 3,.
when God says, "For God so loved the world," it will open his life..
Or it becomes the whole idea of Nehemiah..
We read about Nehemiah at 60-something years old as the cupbearer of the king..

$^{481}$Nehemiah is called to rebuild the walls..
He is 60-something years old..
God calls him to rebuild the walls..
He has had a lucrative job as cupbearer of the king..
Now he is called to rebuild the walls..
But what is the key is he hears one day that Jesus is coming to Jericho..
He wants to see Jesus in the midst of it..
Jesus invites himself to his house..
And Jesus says, "I want you to do something. I want you to stay as a tax collector..
But I want you to make a difference in your city.".
In contrast to Matthew where he says, "Matthew, pick up your everything as a tax collector..
Follow after me..
Peter, James, and John, move away from being fishermen..
Follow after me.".
You never know when God comes calling..
You never know when he is calling you to do something big..
So you live your life with anticipation..
You live your life with open hands..
You live your life constantly praying, "Lord, what do you call me to do?".
But more so, very important too,.
you are faithful to where he has placed you at..
That wherever he placed you at, whether it's a stay-at-home mom,.
whether it's in a job that you cannot stand,.
whether it's a dead-end job,.
whether it's in whatever situation that you're in,.
you're faithful to where he has placed you at..
But you're always anticipating..
You're always looking..
You're always praying..
You're always looking, asking that God may make sense and close and open doors for your life..
Here's the one thing for those who are sitting there,.
especially those who are in their mid-20s to mid-30s,.
because I get a chance to talk with you guys quite a lot..
If you're sitting there in your mid-20s and 30s,.
a lot of times you're always wondering what God is wanting you to do..
You're always anticipating..
You're always saying, "God, what do you want me to do?.
I want to know what God wants me to do.".
But you're not content..
Be content with where you are..

$^{521}$Learn why God is--and we'll talk a little bit about that later..
Be faithful to where God has placed you,.
but always be looking..
If number one is God works in those who are anointed.
by finding those individuals.
and looking at identifying where the heart is,.
then number two, you never know when God is going to work..
You never know when God's calling, so you're anticipating..
Number three--and we're going to work at the final thing,.
and then we're going to move on..
Story goes, in '28, Iliab--we remember Iliab, right?.
Iliab was the oldest son of Jesse..
He was the oldest brother of David..
Now, Iliab, his eldest brother, heard when he spoke to the men,.
and Iliab's anger was kindled against David,.
and he says, "Why have you come down?.
And with whom have you left those sheep in the wilderness?.
I know your presumption and the evils of your heart,.
for you have come down to see the battle.".
And David said, "What have I done now?.
Was it not but a word?".
Point number three--the anointed leader will be criticized by others..
The anointed leader will not be free of criticism..
The anointed leader needs to learn to have thick skin.
and learn that he or she can never make people happy, everyone happy,.
and you are sometimes prone to criticism,.
and you are not free of criticism..
You're going to have to learn that..
Too often, a lot of times, we like to be liked, right?.
Especially for those who are Asian sitting here, we like to be liked..
We don't like conflict..
I don't like, for me personally, when people criticize me..
My defenses get up..
I sit there questioning all of these different things..
I like peace..
I don't like to deal with conflict..
We don't like when people don't like us..
But sometimes, as God's anointed,.
you're going to have to make decisions that will never make people happy..
You may make decisions that people would disagree with,.

$^{561}$that are critical with..
People may not like you because they're jealous of you..
People may not like you because they just don't like you..
God's anointed will not be free of criticism,.
and you're going to have to learn to have thick skin..
You're going to have to learn not to take it personally..
But as God's anointed, you are making decisions.
that glorify and honor the Lord in all that you do,.
and that's a very fact..
You're going to have to deal with criticism..
Just like, for instance, Iliab over here doesn't like David..
Why? We don't know..
Does he have an intent against David?.
Does he not like the fact that David was anointed by Samuel and he wasn't?.
We don't know. I would guess probably so..
And so he's critical..
Look at what he says to him..
This little flock that you have,.
these few sheep in the wilderness,.
he's very critical of him..
You're not going to make people happy..
And so as a result, how do we learn to have thick skin?.
How do we learn not to make things personal.
but to continue to go forth with what God has called us to do?.
If number one is the fact that God's anointed,.
God looks at the heart of the individual.
and never the outer appearance,.
then number two is that you never know when God comes calling..
And you live with open arms, you live with open hands,.
you live with anticipation, asking that God is working..
But number three, you've got to learn to be thick-skinned..
You've got to learn not to be taking things personally..
You've got to learn that sometimes people will be critical..
Out of jealousy, out of fact they disagree with your decisions,.
whatever it is, but you're going to have to learn through all that stuff..
God's anointed means that you are a steward of what God has given you..
And for David, he had to learn these lessons before he became king..
And so as a result of that,.
now next time I come back we're going to talk about part three..
There's a second half of the battle against Goliath..

$^{601}$And then we're going to talk about the aftermath of the days..
But for you, look, here's the one thing I would say to you for those who are listening..
Wherever you are in the world as you sit there,.
always know that God has a plan for you..
Always know that he chooses to want to anoint you to do something bigger.
than you can't even imagine for the gospel..
Now you're going to have to figure out as we wrap up what that calling is..
But that means you've got to be asking and living in anticipation.
of what that practically means in your life..
But you're going to have to do business with your heart..
Because no matter what your skills are outwardly,.
it's the heart that really matters..
So as we sit here, do we want to be used by God? Yeah, we do..
Do we want to be doing great for God in everything that we do? Yeah..
I think all of us will say yes..
But what does it say to your heart?.
Are you anticipating for God to work in your life?.
Are you learning to not take things personally?.
Because God's anointing will not be free of criticism..
How well are you doing in those three areas?.
Let's come back next month and talk about the rest..
Let's pray..
Father, thank you for how you work with us each and every single day..
Thank you for the words that you give us as we look at the life of David.
and how you anoint God, how you anointed David.
and you bring him on the journey..
And we pray that as we face our own calling,.
the journey that you have given to us,.
that we learn from the lessons of David..
May you use us in ways that we can't even desire.
to make a difference for you and for the gospel and the glory of God..
Thank you for the story of David..
We lift these things up to you and give off our heart..
We pray for anticipation, and we pray for thick skin.
to continue to honor you..
\newpage



\section{}
\label{sec:JJV1ojc4_5E}
\textbf{舊約網上行 - 上帝的足跡: EP17 以撒與利百加}
\newline
\newline
連結: \href{https://youtube.com/watch?v=JJV1ojc4_5E}{\texttt{https://youtube.com/watch?v=JJV1ojc4\_5E}} ~~~~ 語音日期: 2021-09-25
\newline
\newline
\hyperref[sec:pfpxXxJmbDw]{\small{< < < PREV SERMON < < <}}
~
\hyperref[sec:index_chronic]{\small{[返順時目]}}
~
\hyperref[sec:index_scriptual]{\small{[返順卷目]}}
~
\hyperref[sec:YtsGA4E1iho]{\small{> > > NEXT SERMON > > >}}
\newline
\newline
$^{1}$我們來到《九一摩》第十七講.
今次我們看《創世記》第十四章.
我們先看第一節.
亞伯拉罕老邁了,年事已高.
耶和華在一切事上都似乎及他.
亞伯拉罕對他家中最老的.
管理他所有產業的老婆人說.
我要你指著耶和華天地的神起誓.
不可從我現在居住的迦南人中.
為我的兒子娶他們的女子為妻.
你可以回到我的故鄉我的親族那裡去.
為我的兒子以撒娶妻.
《創世記》二十四章.
可以說是亞伯拉罕為他兒子以撒尋妻的故事.
最後當然就是找到利伯加Rebecca.
二十四章雖然很長.
但其實只有三個段落.
第一個段落很明顯就是第一節至第九節.
就是亞伯拉罕吩咐他的老婆人.
叫他回到他的故鄉找一個妻子給以撒.
然後故事就發展到第十節至第六十一節.
這個地方就是在美索不達米亞.
這個老婆人聽命就到了哈蘭附近.
然後在水井邊祈禱.
為什麼呢?.
因為水井通常是一些談婚論嫁的典型場景.
type線.
而且打水通常是日落天晚的時候.
打水的人當然全部都是女人.
所以這個老婆人到了水井邊的時候.
她就不知道到底找誰了.
她就求上帝給她一個引證.
她就求問神.
她說如果我問哪個女孩子給我水喝.
這個女孩子不僅會樂意打水給他.
甚至打水給她的樂徒的話.
這個就是上帝的引證了.
之後的故事居然真的發生了.
所以這個老婆人立刻拿了這些嫁妝,金銀珠寶.
然後就去問她的身世.

$^{41}$然後就到了李百嘉的家.
談婚論嫁完了之後.
最後就過了大禮.
就好像我們華人的文化一樣.
就是嫁娶的時候.
以往傳統的都要過大禮.
要有生雞,檳榔等等.
也出現了很多的樂陀.
你留意一下.
在這張裡面出現樂陀.
因為近東的社會裡面.
樂陀就好像今天名貴的房車一樣.
所以是一個很隆重的婚嫁.
然後這個故事就到了最後一段.
就是李百嘉願意離開.
然後到了迦南地.
62-67節就說到李百嘉和以撒.
好像是新人邂逅一樣.
兩個人就拉上天窗.
就是天作之合.
到底如何去看這個故事呢.
《創世記》24章.
是不是一個愛情故事那麼簡單呢.
聖經當然很少描述男女之間的浪漫情節.
24章有一點這樣的味道.
譬如在62節.
那時以撒住在南地.
剛從比爾拉海內回來.
天漸漸漸漸漸晚了.
以撒就出來田間默想.
舉目一看的時候.
就看到一些樂陀正在來.
李百嘉也舉目來看見以撒.
就急忙下了樂陀.
就問那個老僕人.
這個田間來迎接我們到底是誰呢.
老僕人就回答他.
這個就是我們的主人.
就是你將來的丈夫.
李百嘉就急忙的拿柏子夢想念.

$^{81}$就是夢想一句的沙.
然後僕人就將所辦的事.
所遭遇的一切.
就去稟告主人以撒.
以撒就將李百嘉領進去他母親撒拉的帳篷.
娶了他為妻.
並且愛他.
還說以撒自從他母親不在了.
這才得了安慰.
很多時候我們看《創世記》二十四章.
我們就直接了當的.
我們就說這裡一定是在說.
找一個配偶.
如何尋找配偶的典範.
因為明明就是為以撒找一個好的妻子.
我們就直接應用了.
我們就說不要找那些本地的家男的女子.
因為他們拜別神.
不要找他們為妻.
要回到本鄉同親族的人.
另一方面.
魯僕人去尋找這個女子的時候.
就不斷去祈禱.
很倚靠神的.
所以我們就說.
這段聖經還不是告訴我們.
如何去找一個合適的配偶.
一個典範.
這樣應用也是無大礙的.
但是你想清楚的時候.
應用的原則很多時候就很選擇性.
因為方便我們的.
我們就應用.
不方便的.
我們就各自一旁.
方便我們的當然就是.
找李百嘉一個很好的女子.
很有德性的.
很有個性的.
應該找一個這樣的人做你的配偶.

$^{121}$但是不方便我們應用的是什麼呢?.
你留意一下.
這裡是說內族通婚.
這裡是說家父全然安排整個婚事.
新人是不認識大家的.
如果你這樣.
單單看這段聖經.
以為是一個找配偶的模範.
我們照樣煮一碗做.
似乎是錯過了.
可能未必留意到.
聖經這段經文真正的重點在哪裡.
我們知道.
聖經的文化.
聖經的世界.
是跟我們今天現代的世界截然不同.
內族通婚是當時古近東文化.
一些理所當然的事情.
因為婚嫁不是只有兩個人的事.
而是保障家族的資源.
肥水不流別人田.
留下人力.
留下資源在自己的家族裡面.
另外家長安排更加不用說.
如果我們真的想應用一些聖經的原則.
如何今天去摘偶的話.
其實聖經除了這一章.
還有更加識切的經文.
譬如智慧文學等等.
既然不是一個尋找配偶的典範故事.
我們就說.
一定是一個尋求上帝旨意的典範榜樣.
因為魯伯人去到水井邊的時候.
他求一個引證.
希望上帝給他一些不尋常的遭遇.
於是他就能夠確定.
能夠決定上帝要他選擇哪一個路向.
好像他說.
這個女子如果肯打水給我.
甚至他自願地為樂徒打水.

$^{161}$那個就是真命天子.
結果故事真的這樣發生.
類似這樣的情節.
其實聖經有其他地方都出現過.
最著名的當然就是機電的故事.
你記得在士師記裡面.
說到機電.
神父召他去打仗.
他很害怕.
於是他就需要上帝給他引證.
於是他就求問上帝.
當中有一個草地上放羊毛的一個際遇.
於是他就說.
我確定了上帝給我的心意.
有時候我們以為.
這些尋找上帝問引證的.
就是我們今天尋求上帝的典範.
我們就照辦主觀這樣做.
我們問上帝給我們一些引證.
聖經對事情的描述.
不一定等於對我們的命令.
prescription.
所以description不等於prescription.
聖經形容的不等於.
要我們照辦主觀.
命令我們今天這樣做.
正如一個報章.
如實報道一宗案件.
不等於容許或命令我們照辦主觀這樣做.
你記得這兩個例子.
老僕人創世記十四章和機電的故事.
都是舊約歷史比較初期的時候.
到了王國的時候.
上帝似乎不是用這種模式去顯明他的心意.
神猜派他的先知.
很清楚地教導他的子民.
然後繼續追溯舊約歷史的發展.
到第一世紀新約的時候.
聖經起碼舊約的聖經.
上帝最權威的啟示都已經寫了下來.

$^{201}$所以聖經成為神最終的啟示.
我們明白神心意的標準.
當然我們說神一定可以有能力.
用不同的方法去顯明他一些特有的心意.
例如一些異夢,異象,奇妙的遭遇,人證等等.
神當然有這樣的能力去做.
我們也不能排除.
但如果你想清楚原則的時候.
當我們問一些人證.
求一些非比尋常的事情發生的時候.
其實有一個危機.
危機就是要迫上帝按照我們的選擇.
按照我們的尺度來行事.
是我對上帝開一些很特定的要求來控制上帝.
危機就在那裡.
但如果你回到聖經的話.
聖經告訴我們神啟示的權威在哪裡.
就是他自己的道.
是聖經本身.
是唯獨聖經.
所以無論我們有什麼人證.
多麼不尋常的經歷也好.
最後我們還是回到神在聖經裡面很清楚的啟示.
其實你看看聖經.
聖經的心意叫我們作綱作嚴.
要我們見證真理公義.
要我們行事為人.
我們的心態動機要純全是聖潔的.
以基督為我們的主.
這些全部生活裡面的原則.
其實聖經已經啟示得很清楚.
尋求上帝的心意.
其實最好的方法就是我們明白神的心意.
就是藉著聖經來明白.
我多讀一些聖經.
明白上帝的喜好.
明白上帝對我們的要求.
其實已經足夠使我們有一個合乎聖經的人生觀.
一個價值觀.
當然我們一定要祈禱.

$^{241}$但祈禱不應該是問上帝.
到底我要選A還是選B.
反而我們應該問上帝.
求上帝給我們一個純全的動機.
給我們做一個負責任.
有理性的決定.
給我們一個受教肯聽的心.
起碼聽一些信任得過的.
屬靈長者.
其他弟兄姊妹.
甚至在各行各業.
不一定是信主和不信主.
有些更有經驗的人.
客觀地去分析環境.
做一個負責任.
合乎理性的一個有智慧有智慧的決定.
其實這個就是上帝對我們已經啟示了的心意.
既然不是一個尋找配偶的典範故事.
亦不是一個尋求上帝旨意的典範故事.
那麼這個故事的焦點到底在哪裡呢?.
我們問一問這個故事裡面有什麼角色.
第一個就是魯浦人.
他沒有名字.
最多他只是一個敬虔的人.
但故事到最後的時候.
焦點就不是放在他身上.
會不會是第二個角色呢?.
阿巴拉罕.
沒錯阿巴拉罕就是命令魯浦人去韋爾薩找一個女子.
但是你發現阿巴拉罕命令魯浦人之後.
他就從台上或者從故事消失了.
魯浦人帶一個女子回來的時候.
他不是稟告阿巴拉罕.
他居然是向新的主人爾薩去匯報.
好了,那你問.
會不會是第三個角色才是故事的焦點呢?.
就是爾薩.
爾薩在整個故事裡面都是比較被動的.
當然他不能夠親自去到迦林地.
這個我們需要問的問題.

$^{281}$為什麼不能夠離開迦林地?.
似乎就是阿巴拉罕之後的一代爾薩.
他就是那個認許的承受者.
他不能夠離開迦林地.
強調上帝是次列祖這一塊地方的認許.
我們看了三個角色.
最後一個角色當然就是利伯加.
會不會焦點就在他身上呢?.
當然我們曾經提過.
利伯加是一個表表者.
他是一個有德行的女子.
除了這一點.
你有沒有發現有些動詞.
經常出現去描述利伯加.
二十四章五或者三十九節.
這裡老僕人說了什麼呢?.
他說如果那女子不願跟我去的話.
怎麼辦呢?.
二十四章五十一節.
說到拉班等等他的家人.
他說看見利伯加就在你老僕人面前.
你可以帶她去照著耶和華所說的.
作你主人兒子的妻子.
再看多一節五十五節.
拉班母親的家人.
他跟老僕人說.
讓這女孩子和我們同住多幾天.
或是十天.
然後她可以離去.
五十八節.
家人就把利伯加叫來問他.
你願意與這老僕人同去嗎?.
他回答.
我願意去.
這裡越來越看得清楚.
就是去這個動詞.
經文多次環繞一個問題.
問題就是這個女子利伯加肯不肯離開.
到最後他自己回答.
他說我願意去.

$^{321}$這個動詞.
你記起在阿爸有凱旋針的旅程裡.
開頭和結尾都是強調.
第一次創世記十二章.
最後一次創世記二十二章獻以殺.
就是離開或擺脫他的愛子.
神最大的祝福.
以殺這個最大的束縛.
離開.
肯去.
去哪裡?.
去神指示他獻以殺的地方.
誰知到最後.
他得到的束縛更大.
除了動詞之外.
你留意一下這段經文.
有一些字眼.
是和地方有關.
二十四章第三節.
阿爸去到吩咐老僕人.
你要指著耶和華天地的主氣勢.
你要往什麼?.
我的本地本族去.
為我的兒子而殺.
娶一個妻子.
阿爸亦對老僕人說.
你要謹慎.
不要帶我的兒子回到我故鄉的地方.
耶和華天地的主曾經帶我離開.
附加本族之地.
對我說.
對我氣勢.
將這塊地賜給你的後裔.
所以有這些平衡的動詞和形容詞.
將阿巴拉罕信心的旅程.
和利伯家連在一起.
正如阿巴拉罕離開本地本族本家.
願意去.
現在新的一代.
利伯家願不願意去.

$^{361}$結果他也願意離開本地本族本家.
如果你這樣看.
創世記十四章.
並不是純粹分家的故事.
而是一個呼召的故事.
呼召新的一代.
呼召利伯家.
他離開本地本族本家.
讓上帝在他生命當中.
能夠成為萬國的祝福.
舊的一代.
上帝曾經呼召了.
亦以信心他們回應了.
今次新的一代.
到底願不願意.
繼續承傳上帝的工作.
繼續回應上帝的呼召.
疫情已經一年多.
似乎使整個世界都停頓.
使很多工作都慢了下來.
甚至我們可能有時覺得.
上帝的工作都會慢下來.
但如果你留意.
在這一代.
上帝仍然去到呼召人.
新的一批人.
你都要受裝備.
你都要服侍.
在上個星期.
有一個調查出現.
在北美的神學院.
有二百多間.
按照統計.
自二零一六年起.
到現在.
只有七間有所增長.
其中一間就是.
我們城會教會.
有所合作的北美神學院.
其中一間有所增長.

$^{401}$幾年之前.
北美神學院決定.
賣掉他們附近的.
一個很大的校舍.
將它縮小.
然後他們開展了.
網上的教學.
疫情就更加加速了進度.
有很多上一代.
傳統的神學院的思維.
仍然觀望態度.
希望如果疫情過去了.
我們一起回復正常.
對於一些向前看的人.
對於一些不斷發展.
不斷尋求新的改變的人.
北美神學院就決定.
不單止做實體.
更加主攻網上的教學.
我們現在知道.
疫情到現在.
不可以只單單是實體.
實體和網上都需要兼顧.
所以今年北美神學院.
居然在他們的學位.
研究生的層次.
人數比過去十年更多.
是破紀錄的.
你說是不是上帝.
繼續呼召人.
呼召他的工人.
裝備他們.
以至他們能夠成為上帝對萬國的祝福.
李百嘉就是代表新的一代.
他願意回應這個呼召.
開設本地本族附加.
這也是他最後能夠接受的祝福.
事實上祝福這個字.
貫穿創世記二十四章.
而祝福這個字.

$^{441}$亦與李百嘉這個名字很有關係.
因為李百嘉這個字.
就是Ribca R-B-K-H組成.
而祝福在希伯來文.
就是Barak B-R-K.
有三個字母都在李百嘉這個名字出現.
而且到二十四章最後.
李百嘉得到的祝福是什麼祝福.
他說:於是他們把自己妹妹李百嘉.
她的乳母阿布拉罕的僕人和同來的人都打發走了.
他們給李百嘉祝福.
對他說:我們的妹妹願你作千萬人之舞.
願你的後裔佔領受敵的城門.
兩點.
第一點千萬人能夠祝福千萬人.
第二點就是佔領受敵的城門.
如果你記得的話.
創世記二十二章獻以殺故事.
到最後神對阿布拉罕的祝福是什麼.
亦都是很多人能夠得到祝福.
上帝說我必定賜福給你.
使你的後裔繁多.
千千萬萬天上的星海邊的沙.
你的後裔必佔領受敵的城門.
地上萬國要因你得福.
因為你聽從我的話.
你肯去.
故事說到阿布拉罕的祝福.
又臨到李百嘉新的一代.
因為他願意回應.
而且故事的地方都有所改變.
故事開始的時候就在迦南地.
然後就去到美索不達米亞.
然後又回到迦南地.
不止這樣.
也進入了撒拉的帳篷.
上一章我們說過.
是撒拉第一代的女族長死和埋葬.
過去了.
現在李百嘉就取代撒拉.

$^{481}$成為新的女族長.
是新的一代接棒的故事.
創世記二十四章.
是創世記裡面最長的一章聖經.
總共有六十七節.
我們就問.
為什麼要這麼詳細記載一個婚嫁故事.
原來是一個呼召接棒的故事.
等於你看奧運會四成一百接力賽.
最關鍵的焦點在哪裡.
就是慢鏡來看清楚接棒的過程.
接棒的一刻才是最重要.
就像李百嘉接上棒.
用一個很長的故事.
慢鏡來記載.
他肯離開 肯去.
回應上帝的呼召.
今天上帝仍然呼召很多人.
新的一代 新的工作 新的崗位.
甚至是一些舊的崗位.
仍然用一個新的態度來承接.
上帝仍然呼召你和我來服侍他.
李百嘉 Rebecca.
其實就是女性的阿博艾漢.
Female Abraham.
上帝呼召的不單是男人.
也是呼召女人.
我們經常有一種感覺.
聖經是男人當道.
英雄通常都是男人.
這個無可厚非.
因為在近東的社會裡.
男人是重心.
但如果你仔細留意一下.
只看摩西五經.
其實一半的章節 一半的章節.
57\%都有提及女性.
創世記主旨是記載.
神呼召阿博艾漢這個家庭.
成為神的選民.

$^{521}$總共有50章.
創世記有39章.
你留意一下 都是有提及女性的.
佔78\% 或多或少.
女性 神藉著她們來工作.
神呼召男人.
也呼召女人.
弟兄或姊妹.
只要我們願意.
求主幫助我們.
無論我們今天遇到的境況.
有疫情或沒有疫情.
讓我們今天有一個新的心態.
新的工作.
對我們一個再一次的呼召.
我們跟上帝說.
我願意離開.
我願意去 去跟隨你.
我們同心得道祈禱.
因為你每時每刻都呼召你自己的工人.
呼召我們.
我們能夠去跟隨你.
求你繼續幫助我們.
讓我們有信心 遵命.
聽見你呼召的時候.
我們說我們願意去.
離開我們的本地.
離開我們本有的舒服的地方.
願意接上棒.
願意去繼續傳承你自己今天將來的工作.
求上帝你藉著我們.
使更多人能夠得到福氣.
使更多人能夠認識你.
得到你的恩典.
祈禱奉主明求.
\newpage



\section{3:1-21}
\label{sec:YtsGA4E1iho}
\textbf{Rejoice in Knowing Christ (Philippians 3\_1-21, 4\_4) - Jonathan Barney}
\newline
\newline
連結: \href{https://youtube.com/watch?v=YtsGA4E1iho}{\texttt{https://youtube.com/watch?v=YtsGA4E1iho}} ~~~~ 語音日期: 2021-11-01
\newline
\newline
\hyperref[sec:JJV1ojc4_5E]{\small{< < < PREV SERMON < < <}}
~
\hyperref[sec:index_chronic]{\small{[返順時目]}}
~
\hyperref[sec:index_scriptual]{\small{[返順卷目]}}
~
\hyperref[sec:MkL9xiQNZJM]{\small{> > > NEXT SERMON > > >}}
\newline
\newline
$^{1}$Awesome. Well thank you very much for having me..
It's great to be here and be with your pastor, Pastor Mark..
Yeah, it's really neat for me because I was in a, so I'm in the process of.
planning a church, all people's church in the West suburbs, and Pastor Mark was.
one of the trainers, the facilitators, so it's really cool to come and be with.
you all now. So well let me pray for us before we get in..
Father in heaven, I pray that you would be with us now..
Bless the words of my mouth and the meditations of all our hearts. In Jesus.
name we pray, amen. Amen. What we believe about God is the most important thing.
about us. What we think about God, it's the most important thing about us. Do we.
see God as distant? Do we see God as a some kind of cosmic killjoy? It matters..
It matters our view of God. I got four points for us today as we look at what.
Paul is saying, what Paul is saying, and just a little background here. We're.
jumping just right into the verse, these verses in Philippians chapter 3, and this.
probably comes out of my church had been going through the book of Acts for like.
a year plus, and then we jumped into Philippians because in the book of Acts,.
kind of the last half of it, we get a picture of into Paul's life as he is.
going around much of the known world at the time, 2,000 years ago, sharing about.
Christ, and the Holy Spirit is making disciples, and it's this amazing thing.
that God is doing. And then he finally gets to Rome where he's under house.
arrest, and he's under house arrest for two years as he's looking to stand.
before the Emperor Nero. And the Emperor Nero, many of you have heard of him, he was a.
pretty bad dude as we know through history. And so he's waiting, he's waiting.
for two years, and we get a few letters. We get these amazing letters of.
scripture that Paul writes during this time as he's just waiting in prison. We.
get this letter of the Philippians, we get Ephesians, we get Colossians, we get.
some of these letters. And so he's writing from under house arrest, and I.
saw that's important, that's important. And so he writes the.
Philippians, he hears how they're doing, and he has to write to them. And so.
that's where we get, starting off right here in verse 1, to write the same.
things to you is no trouble. The same thing. So he's previously instructed them,.
he knew them, he was with their right, Paul and Silas, maybe you remember that.
story growing up. And he wants to write them the same things, especially he wants.
to pinpoint, laser focus on the gospel, on the gospel. See the gospels is as Tim.
Keller says really well, not the ABCs of Christianity, but the A through Z of.
Christianity. He doesn't want them to be distracted by some other doctrine or.
teaching or something else because there's been false apostles, there's.
false teachers, and people are spreading lies about that. And so he wants us to.
look right at the gospel. And what he's talking about, I'm reminded today is.
you guessed it, no not Halloween, it's Reformation Day. It's Reformation Day..

$^{41}$Yes, yes, of course. Yeah, so I'm a nerd and I got a really bad joke for you. So.
it's Halloween and some kids come to the door, they knock on Martin Luther's door..
All right, they knock knock. "Who is it?" "Ah, trick or treat!" "No, no kids, I said 95.
feces, not 95 Reese's, okay? I don't have any candy for you. I got the gospel though.".
Okay, yeah, bad joke, bad joke. I just had to though, I couldn't resist. All right, so.
at least Pastor Mark is laughing with me, he gets the joke. All right,.
Reformation Day reminds us of how Martin Luther, John Calvin, these reformers, these.
others, they point us back to the truth because the church had gone way off and.
they reminded us of two huge things. That our trust is in the Word of.
God alone. That's our authority, God's Word. We got to find truth there, nowhere.
else. That's the ultimate place. And then secondly, that we are saved by grace. We.
are saved by grace through faith. And that's what really he is wanting to.
pinpoint and remind the Philippians of today. And I think we need that.
desperately as well. So that's what's coming up. That's what's going on. That's.
why the stupid joke. And Martin Luther, I love he said this, he says justification..
Justification is the biblical Christian word that just means that we are saved.
by grace, we're accepted by God, we've declared and made righteous by his work,.
not our. He says justification is the article by which the church stands or.
falls. But he wants them to remember the gospel. He tells them this, "Therefore my.
brothers and sisters, you whom I love and long for, my joy in my crown, stand firm.
in the Lord this way, dear friends." He wants them to stand firm on the gospel..
Stand firm in what I taught you. What Silas and I, what we taught you when we.
shared the gospel. And that's the same thing for us, brothers and sisters. Stand firm on.
the gospel. So that's our first point that I want us to get. We want to.
rejoice in the Lord. We rejoice in the Lord. We need to know the gospel. We need.
to stand firm on the gospel and nothing else. Friends, what are you standing firm.
on? What are you standing firm on today? Is it anything else other than the gospel?.
Then that will really be shaken. I can tell you from experience that will be.
shaken. All right. And we're just gonna continue to explore this even more..
Second point is this, Paul has for us, is to reject works righteousness. If we.
really want to rejoice in the Lord, to take joy in knowing Christ, because.
that's what this passage is all about, then we have to reject works.
righteousness. We see that in verses 2 through 6. Look with me there. He says,.
"Watch out for those dogs, those evildoers, those mutilators of the flesh.".
Wow. He's got some strong language there. Strong language. He's not, he's not holding.
back. Dogs. All right. So he calls them dogs. So what's going on here, all right?.
What's going on here? They talk about circumcision before and maybe some.
of you get a little squeamish now. So there are these Jewish.
Christians and people who are coming in. There was this huge tension between Jews.
and non-Jews. Jews and Gentiles at the beginning of the church. And a lot of.

$^{81}$these Jewish Christians, some Jewish Christians were saying, "All right, you.
have to be circumcised in order to be a follower of Christ." And in fact, they so.
tied it to it that it was part of their justification, it was part of their salvation..
That, all right, in order to know that you're really saved, you've got to be.
circumcised. All right? Maybe some Christians today, you think of, you think.
of, all right, you have to be baptized in order to be saved. It's like, all right,.
you know, you get saved because you, or you become saved, you become a follower.
of Christ, and then you get baptized. All right? But that doesn't mean it's not.
works-based. And so that's what he's getting at with this. That's what he's.
getting at. And he calls them dogs because back in the ancient world, 2,000.
years ago, it wasn't like your sweet little puppy dog, you know, it wasn't your.
nice pet dog. 2,000 years ago, all right, dogs were gross. They were filthy. No one.
had them as pets. They ran wild in the streets, okay? They were disgusting. All.
right? So he's saying, "The dogs." And you know what? It was a ethnic slur,.
actually, that Jews would use for non-Jews. They'd say, "Man, those dirty,.
uncircumcised Gentiles, those non-Jews, they're dogs." So Paul turns it around, he.
flips it around, he says, "You know what? Who really are the dogs? These people who.
are saying, 'Oh, you have to be circumcised in order to be a Christian.'".
Because what did you have to do to be a Christian? To believe, to be saved by.
grace through faith, right? So he's saying, "All right, these evildoers, these.
mutilators of the flesh." He goes on, he says, "We are the circumcision, who.
worship by the Spirit of God in glory in Christ Jesus and put no confidence in.
the flesh." So he's saying, "You need to reject works righteousness. You're not.
saved by this ritual, all right? No, you're saved by faith." Circumcision in the Old.
Testament was never meant to be a thing that, "Oh, that's a work that.
saved you." No, it was a sign. It pointed to the reality that you had been.
saved, that sin has been cut away in your life and now you're a follower of God..
I love this. He goes on, verse 4, let's keep going, "Though I myself have reasons.
for such confidence." So now Paul is, all right, he's worked up here. He's like,.
"All right, if you want to say some confidence in the flesh of what our.
works are, he doesn't want them to be that way." So he says, "All right, I have.
some confidence in my flesh that I could talk about. If someone else thinks he has.
reasons to put confidence in flesh, I have more." All right? So then Paul goes.
into this and he has seven things that he can have confidence in, all right, that.
he was this amazing Jewish person, all right, circumcised on the eighth day..
All right, you were supposed to do that according to Leviticus, according to the.
command of God, of the people of Israel, of the tribe of Benjamin. All right,.
Benjamin, that tribe was one that sometimes stayed more faithful to God.
when other tribes fell away from God throughout Old Testament history. He was.
a Hebrew of Hebrews. All right, what does that mean? All right, we know, we already.

$^{121}$know you're Jewish, Paul, okay? No, so what that meant was that he was not only.
Jewish, but he actually spoke Aramaic and his parents spoke Aramaic. That was the.
Hebrew language of the day. So he's saying, "I don't just speak Greek, I speak,.
man, I'm..." So he's building his credential, all right? His resume is building. He's like,.
"Man, wow." All right, he keeps going though. In regard to the law of Pharisee,.
Pharisees were regarded as the strictest kind of Jewish sect at the time. There are.
a few kind of branches, if you will, of Judaism, and Pharisee, that was strict, and.
he was climbing that ladder. As for zeal, persecuted the church. I persecuted the.
church. Man, I was so zealous for my faith in God, I was such a good Jewish boy that,.
man, I held the cloaks for those, I held the coats, I was the coat rack for those,.
who stoned Stephen, the first Christian martyr. You can go back and look at.
that, it's like Acts 7, I think. I persecuted the church. As for.
righteousness based on the law, faultless. As for righteousness based on the law,.
faultless. All right, so this doesn't actually mean that Paul was without.
sin. You can read Romans 7, you can read a lot of the rest of his letters, he is.
very... He says, "I'm the chief of all sinners. I'm the worst." Right? But he said, "On the.
outside, I kept all the laws. On the outside, it looked like I was perfect.".
That's what he's talking about here. So he's telling, all right, this is....
Paul's supposed to sound just crazy. All these things, though, they sound very,.
very good. The Jews would be wowed by his resume, you know? Maybe for some of you,.
maybe some of you young people who are looking to go to college, or went to.
college, or going to get your master's, or whatever. Man, you did so many.
extracurriculars, right? And you're Asian too, I know how it is. No, I'm just kidding..
Hey, hey, yeah, exactly. I shout out because I grew up in San Francisco Bay Area. I was.
talking to Pastor Mark before on this. I was just back after 11 years, San.
Francisco Bay Area. And I got a lot of my Asian friends there, man. All the Asian.
love, man. I'm telling you. All my friends were Chinese, and all my friends,.
most of my friends were Asian because I was in the honors AP classes, right? They.
were in the honors AP. So they were all Asian. I'm just saying, I'm just saying. I.
know how it is. I know your parents and how your parents are. Yeah, there you go..
But yeah, maybe you were building that resume to try to impress your parents..
But maybe you keep doing that to try to impress God. Man, God, I've done this for.
you. I've sacrificed. I went on that mission trip. Man, I'm praying hard. I'm.
reading my Bible. It's just the same thing like circumcision. All right, I'm.
all right. That'll make me earn your approval. I know for me, I was just.
the same way, a huge way. I remember when I was in college, I became part of this.
church. And it was a very passionate church. There were a lot of good.
things there. I remember I was coming out of a, I've shared with many people in my.
church, I deal with situational depression. And I had the worst, my.
first bout when I was 17. And I would have been depressed out of my mind. I had.

$^{161}$considered myself before that a mature Christian. But God really humbled.
me. And so I have all that baggage and I'm going, I come out here to the Midwest.
to go to school. I go to Wheaton College. And right, it's just this, supposed to be.
this evangelical flagship that just, and there's, I'm on a dorm floor and there's.
40 other Christian guys. And you know, we were praying in the classroom. I have.
professors who are Christian. I go to chapel three times a week, three times..
Chapel. I get involved in this passionate church, right? I go on mission trips. Man,.
God totally changed my life through this dorm discipleship group, these.
guys. He started to bring me out of this depression and really restore my.
relationship with him again. This group was so intense that it meant three.
nights a week as a very rigorous academic school. Three nights a week. This.
meant, and they totally changed my life. I got real with these other guys and.
they got real with me about their lives and we encouraged each other. And so I.
just wanted to share my faith with others as well. And I would go to some.
local campuses, college campuses, and share my faith with them. And I would go.
one to three times a week to share the gospel with people at campuses. I started.
leading small groups with my church. I was passionate and I would fast.
one or two times a day because I want to set apart time to meet with God. I would.
spend two or three hours in the morning meeting with God. You know, after a couple.
years of all that, I remember at one point I was trying to spend just a few.
minutes with God. This was after I was married and my wife and I, our picture of.
the world and our worldview of Christianity had been totally turned upside.
down just a few months before because we had realized some very dysfunctional.
things about that church. It was very passionate but it was so works.
oriented. And I had this huge burden on my shoulders of self-righteousness.
because I was a Pharisee of Pharisees, guys. See, I had an impressive resume too..
On the outside I looked really good, but man on the inside I felt like dead man's.
bones. And that's Paul's story too. And so he's passionate about telling the.
Philippians here, "Man, trust in the gospel of grace." Preaching the gospel to.
yourself. It requires us, you see here, "reject works righteousness." He says, "All.
right, reject it. Don't put it, base it on anything else, your own works, your own.
goodness, whatever, other than Christ's goodness, His grace." All right, let's keep.
going. Third point, "Remain in faith righteousness." Remain in faith.
righteousness. We see this in verses 7 through 11. "But whatever were gains to me.
I now consider loss for the sake of Christ." So if you will, he's got.
these two columns, right? Your gains and losses column. He says, "All right,.
just as I went through all these gains, he went through all these gains. I'm a.
Hebrew of Hebrews. I was circumcised on the eighth day. I was a Pharisee. I.
persecuted the church. Man, I was there. I was so passionate for God." He's got.
all these gains. He says, "You know what? They're in the losses column now..

$^{201}$Every single one of them. Every single one of them. Where gains to me now I.
consider as a loss for the sake of Christ. What is more?" He goes further,.
verse 8, "I consider everything a loss because of the surpassing worth of.
knowing Christ Jesus my Lord, for whose sake I have lost all things." Remember,.
he's in house arrest in Rome. Man, he's been shipwrecked a couple.
times. He's been beat with rods. He's suffered the lashes. He's been stoned. One.
time they thought he was dead. He was just been stoned. They were about to.
make him, they were about to sacrifice to him. They were calling him Zeus and.
Barnabas, him Hermes and Barnabas Zeus. They were about to make sacrifices to.
him. And then some Jews come in and they start stoning him. Man, what a.
contrast there. They thought he was dead, but he ends up getting up. This guy had.
been through so much, right? He says, "All right, because of this grace, because this.
free gift of grace, not works righteousness, lost all things because I.
want to know Jesus. I want to know Jesus. I consider them garbage." So not only.
are they in the loss column, but I consider them garbage. That's this word.
that means refuse. Refuse. So I want to say then, just think about what this word.
means. They didn't have running water or a sanitation system like we do now..
Nuff said. Refuse. He says, "All these good works, I consider them refuse and be found.
in him." Verse 9, "Not having a righteousness of my own that comes from.
the law, but that which comes through faith in Christ." The righteousness from.
God on the basis of faith..
It's amazing. It's good news. It is good news. We're learning to preach the.
gospel to ourselves. We need to learn to preach the gospel to ourselves to know.
this. So what is it for you? Maybe is it your own purity? Maybe that your.
parents are Christian? That you're Christian? That you go to church? Maybe.
you've done some significantly good things in your life? Your knowledge of.
the Bible? Your giving? Your overall just you're a good moral person? Is.
that what you're basing your acceptance with God in? Or is it in Christ's work?.
Verses 9, he talks about the life of righteousness. That we're saved by faith.
alone, Luther said, but that faith that saves is not alone. So how do we remain.
in faith righteousness? Yeah, we keep on rejecting it. We find ourselves in that..
We encourage ourselves in that. That it's by grace. We keep on preaching that to.
ourselves, but we also want to continue. We want to obey. We want to keep.
following Christ. And this is amazing what he says here. Let's keep going on..
Verse 9, he says, "And be found in him, not having a righteousness of my own that.
comes from the law, but that which comes through faith in Christ." I want to know.
Christ, yes. I want to know the power of the resurrection and participation in.
sufferings, becoming like him in his death, and so somehow attaining to the.
resurrection of the dead. That I may know him. Man, we make life about so many.
things, guys, right? But do you, every day, do you get up and say, "All right, I want to.

$^{241}$know Jesus today. At the end of today, I want to know Jesus." He says, "I want the.
power of his resurrection." What's the power of his resurrection? The power to.
save sinners. That sinners are saved. That, all right, we get the righteousness of.
Christ because he got all our sin on the cross, and God punished him for that. But.
then that means we get the righteousness of Christ. That's amazing. The power of.
the resurrection is also God saves us through Jesus on the cross, but then he.
also helps us to live a life of love, of loving others, of loving God. It's the.
power to live holy lives. That I may share in his sufferings. Man, did you see.
that this last week? Did you say like, "Man, you know what? I want to share in the.
sufferings of Jesus." I doubt most of you said that. Man, we might not be persecuted.
in a way that Paul is, but we will in other ways. If we're sharing the gospel,.
we can do that. If we're living as Christians, as we just go through.
hardship in our lives because of the choices that we make to follow Jesus,.
it's gonna be difficult. It is hard to be a Christian. The way is narrow, and it's.
hard, but it leads to life. He says, "Becoming like him in his death." Becoming.
like him in his death. That we also can die to ourselves. We can die to our sin.
to follow him. That he would attain the resurrection of the dead. So don't get.
Paul wrong. He's saying, "Once saved, always saved," but we have to endure to the.
end. He's encouraging that you have to continue to follow Christ. It's the cross.
before the crown. That we die before we can live, and we do that each day in our.
lives. So how do we, again, rejoice in knowing Christ? We remain in faith.
righteousness. In our righteousness, we are righteous because of our faith in him,.
not our own works. Fourthly, fourth point, he says to reside in his work of.
righteousness. We rest in his work of righteousness. We see this in verses 12.
and also through the end of the passage in 4.1. How do we do this? How do we do.
this? I love there's this, there's this line in the book of Hebrews, and he's.
talking about the rest of God. The rest of God. The Sabbath rest for God's.
people. And he says, "You know, man, you guys need to rest." He says, "How do you do that.
though? You strive." What? "You work really hard to not work." So what he's saying is,.
you got to work, you got to fight to remain in his righteousness. That it's.
not your work to be accepted by God. You got to work at that. You got to work to.
also live a life of holiness as you continue, because he wants to strain.
toward the goal. We see this in verse 12. He says, "Not that I have already obtained.
this." Alright, so I'm not with Christ fully now. I'm not in heaven. I'm.
not arrived at that goal. I'm not perfect, but I press on to take hold of which.
for which Christ Jesus has taken hold of me. Alright, so I can I can strive for.
that because, man, I'm already Christ. He took hold of me. Can you get out of.
Christ's hand? No. How absurd would that be? How can you go out of God's hand? So I.
want to keep on going to make Christ my own. I want to press on, take hold of it..
It's like this picture, right? I love the great artist, sculptor, Michelangelo,.

$^{281}$right? You guys probably know that famous statue of David, picture of David,.
and Michelangelo was asked, "Alright, how do you do this?.
You produce these amazing works, sculptures and works of art. How do you.
do this?" He says, "I just..." See, back then what they'd get is this big.
rectangular or block of marble, just huge massive block of marble, alright? And he.
says, "I would look at the block of marble and I would see the man in there. I would.
see the man in there and I just have to try to call it out. I have to like bring.
it out of there." Isn't that beautiful? But how he does that, there's this block of.
marble and he would have to go over there and have this huge block and slowly chip.
away at this block of marble all around it to get this amazing, just very.
intricate, detailed body of David that we see. That's what God does to us. Holiness,.
striving toward God, is as we continue to reside, to rest in his work for us,.
we have to press on toward the goal. It's painful at times. He's chipping it away.
at our sinful nature, but that prize, he says, "The prize is the fullness of.
rewards and joys in heaven." And so he sets his eyes on that, the fullness of.
the rewards and joys in heaven. It's this perfect life with Christ forever..
Verse 17, he says this, he says, "Join together in following my example,.
brothers and sisters, just as you have us as a model. Keep your eyes on those who.
live as we do." Alright, so look at me, look at your Christian leaders, look.
at your pastor here, look at the leaders and mentors that you guys have. To imitate.
them, imitate them, to have godly discipleship. Maybe some of you need to.
say, call up or text someone in the church and say, "You know what? I need to.
talk to you. I love, I need your help in following Christ. Let's do this together..
Let's follow Christ together." That's huge. That's massive. That has totally changed.
my life, walking with other believers, walking with other men of God who I want.
to be like, and we can sharpen one another..
Verse 18, "For as I have often told you before and now tell you again, even with.
tears, many live as enemies of the cross of Christ. Their destiny is destruction..
Their God is their stomach. Their glory and their shame, their mind is set on.
earthly things." So what's he talking about here? So these are these Judaizers,.
again, who, they're on a self-salvation project. They're saying, you know, you need.
to add these dietary laws, you need to add this circumcision, whatever, so it's.
Jesus plus something equals your acceptance with God, rather than just.
Jesus equals your acceptance with God. He's wanting to add these things to them..
And Paul says their destiny is final judgment. They worship themselves. That's.
what we're talking about. They worship their belly. They just worship.
themselves. Man, when I was so burdened by my own self-righteousness and just.
living, I was just always navel-gazing, as I like to say. I just.
always like, "Man, how am I doing? Am I doing good? Am I doing enough for God? Man,.
that's so exhausting, right? That's so exhausting." He says the things that we.

$^{321}$can do, the earthly things, the works of the flesh, they're consumed with earthly.
things, rather than Christ's external work, His work. Verses 20 through 21, here's.
another thing that the Philippians could have boasted. He says, "But our citizenship.
is in heaven." You see, Philippi, they thought they were all that, because they.
were a Roman colony. So they're like, "You know what? I get free taxes. How cool.
would that be, right? Free taxes." I mean, what am I saying? Free tax. That's what we.
get all right. No. They had freedom from taxes. They wouldn't get taxed, because.
they were a Roman colony. And so everyone there who was born there, you were a.
citizen automatically. You had this special thing. It was the status at.
the time. And so he's saying, "All right, your citizenship is in heaven. It's not.
with Caesar. It's not with Rome, right? So our identity, right? Our identity is not in.
our performance. It's not in other people's approval of us. It's our identity is in.
Christ, in heaven. And from there, we eagerly await a Savior, the Lord Jesus.
Christ, who by the power that enables Him to bring everything under His control.
will transform our lowly bodies so that they will be like His glorious body." So he.
says, "All right, there's the prize. There's the prize." He's using that imagery of the.
ancient world. They were really big into sports, just like we are today. And.
there was that image of a wreath that at the finish line, that they would be.
given this crown, this wreath. And he's saying, "All right, there's this prize you can go for.
of reaching Christ, of running the race." And he's saying, "Yeah, I know it's hard to.
transform your lowly bodies because right now we are dying every day, right?.
We're dying every day. We're fighting against sin. We're slugging it out.
and different hardships in our lives. But he says, "One day, we'll have a glorious body..
We're gonna have a glorious body. Perfection is coming." He says, "Therefore,.
brothers and sisters, with you whom I love and long for," he wants to see them.
again, "my joy and my crown." There's that crown, that imagery, right? "My joy and my.
crown. Stand firm in the Lord in this way, dear friends." He says, "Keep on going..
Remain in the gospel. Remain faithful." And maybe some of you need to hear that.
now. You know, the uncertainty of maybe things with COVID, uncertainty maybe with.
things with your job, things with how your future is gonna be maybe after.
graduation, with a relationship you're in, with your family dynamics. I don't know.
what's going on with you. Just stay firm in the gospel. Just stay firm in following.
Christ because we can rejoice in Christ because His work for us, not our work,.
that saved us. That we are totally accepted by God. Just as Christ is.
accepted, He treats us the same way. We are accepted before God. That is a good.
news. I'll close with this. You know, it reminds me, we can rejoice in Christ..
We can rejoice in knowing Him, that we know Him. It's like this. It can be.
like marriage. It can be like falling in love. Have any of you fallen in love?.
She's saying that, you know, as her husband, I think, is next door. You know,.
you'll do, you'll do. People in love are crazy. They'll do anything. They're nuts,.

$^{361}$right? You know what I'm saying? They're nuts. They're crazy. You'll do anything..
They start acting weird ways, right? I just, you know, you're like, you're like,.
"Oh, I thought I knew him. Like, I thought he would never, you know, we were still.
gonna be friends. Now it's like, she came in and it's like, who are you?" You know,.
they start doing crazy things. You know, public displays of affection..
I thought, I never thought that they would do that. People in love are crazy. You'll do.
anything for that person. You didn't care about anything else. But imagine this..
Imagine, each of you, you're married to Mr. Law. Imagine with me here. You're.
married to Mr. Law. And Mr. Law, he likes everything very particular. He likes.
everything very particular. So he would come home every day at 5. He eats the.
same breakfast every morning. He comes out at 5. He demands a kiss and dinner on.
the table. You always had to look just right for him. Otherwise, you'd hear from.
it. You could never really be real in your authentic self with Mr. Law. No, you.
definitely couldn't be. All the time, you just felt like you were striving to be.
good enough. You never felt good enough for him. Like, sure, yeah, he'd married you,.
but no matter how desperately hard you try, you always felt like you were just.
letting him down. And in fact, he would tell you that. He would tell you that. You.
were never good enough for him. You're always working on yourself. You're always.
trying to gain his approval. But then one day, Mr. Law died. And Mr. Faith.
Righteousness came. I sound like Pilgrim's Progress, John Bunyan or.
something right here. Jesus comes. Mr. Righteousness by faith. Mr. Grace. He.
comes and he sweeps you off your feet. Indeed, it seemed he always loved you, no.
matter all your deep shortcomings. He always put your needs above his..
Every day, every moment, he was dying to himself for you. He always took out the.
trash. He cleaned the bathroom. With him, you were still so very in touch with.
your massive shortcomings. You just felt even more like, "Man, I'm so messed up." But.
you knew how much he loved and accepted you. You always found yourself actually.
just loving him more. Loving life more. Loving other people more. And you were.
slowly even changing and becoming a better person. That's what Paul talks.
about in the letters. And I think what he's getting at here, he says in.
Romans, he writes this. He says, "Likewise, my brothers and sisters, you have.
died to the law." You were married to Mr. Law. You died through the body of Christ.
so that you may belong to another. That's marriage language. To him who has been.
raised from the dead that you may bear fruit for God. For a while, you are still.
living in your flesh. Your sinful passions aroused by the law were at work.
in your members to bear fruit for death. But now you're released from the law..
You're no longer under this. You have to earn your way. But Christ earned.
your way for you, having died to that which held us captive. Friends, I want.
you, I want myself to know the joy rejoicing to knowing Christ. Here's how.
you do this. Paul says right here, "Rejoice in the Lord.".

$^{401}$Reject works righteousness. Remain in faith righteousness. Rest in his work of.
righteousness. Let's pray. Father in heaven, we thank you that you sent Jesus.
and we're reminded of that today. Lord, that what you do, you did what the law.
couldn't do. Lord, I pray that we would rejoice. We'd have so much joy flowing.
out of our lives because we know Jesus. That we would make our life about that..
That we would revel in what you've done. That we would reject any other kind of.
righteousness, any kind of way of making ourselves acceptable to you God, whatever.
that may look like. In Jesus' name I pray, amen..
\newpage



\section{}
\label{sec:MkL9xiQNZJM}
\textbf{We are Jonah (Jonah 4) - Brock Lockenour}
\newline
\newline
連結: \href{https://youtube.com/watch?v=MkL9xiQNZJM}{\texttt{https://youtube.com/watch?v=MkL9xiQNZJM}} ~~~~ 語音日期: 2021-11-08
\newline
\newline
\hyperref[sec:YtsGA4E1iho]{\small{< < < PREV SERMON < < <}}
~
\hyperref[sec:index_chronic]{\small{[返順時目]}}
~
\hyperref[sec:index_scriptual]{\small{[返順卷目]}}
~
\hyperref[sec:GChzvOXTH_M]{\small{> > > NEXT SERMON > > >}}
\newline
\newline
$^{1}$Let's welcome him..
(audience applauding).
- Hello, hi everybody..
Thank you, Pastor Mark, for praying,.
and thank you, James, for leading worship..
Fun fact, I texted James this morning.
to see what he was wearing.
so that we could make sure we were matching..
(laughing).
That's not true, we just have a similar fashion sense..
(laughing).
Thank you both..
So today, I wanna talk about the Book of Jonah,.
but before we do that, I wanna tell you guys a story..
This is a story that I haven't really told before..
It's been something that happened.
and was decently significant in my life for a while,.
but I didn't really see it as something.
that was timely or necessary at the time..
But I think the themes of this story fit really well.
with the themes of the Book of Jonah,.
so I'm bringing it out today..
I want you to think back to about a year and a half ago..
This would be right as Chicago is getting out.
of our state-mandated lockdown, right?.
We're all kind of trying to figure out.
how to do this post-pandemic world,.
even though it was still mostly in the pandemic,.
and we're trying to figure out how to continue to function.
even though this virus is spreading..
And we've been locked up for two and a half months..
Things are kind of tense, right?.
And then, right as quarantines are lifting,.
the killing of George Floyd is basically televised.
and thrown everywhere online, and protests erupt..
Peaceful marches are everywhere over our city..
It was kind of a hectic time..
Not only were we living through a massive worldwide pandemic.
like once in every hundred years,.
but we're also experiencing crazy amounts of social unrest.

$^{41}$and calls for new things to happen.
to bring justice about in our country..
It was a bit of a crazy time..
Now, for those of you who don't know,.
I am a photographer, kind of..
I get to do part of that for my job,.
and I also get to do it on the side..
So I love taking photos of cool places in Chicago..
I love taking photos of cool places when I travel,.
and I also really enjoy taking photos of worldwide events..
So I loved it when the Cubs won the World Series.
because I got to go, and I got to take photos.
of entire crowds at Wrigleyville.
just celebrating and getting excited..
I loved it when the pandemic happened.
because I got to take really stark photos of empty shelves,.
no toilet paper in sight..
It was great..
It's like a photographer's dream.
to live through these kinds of world-changing events, right?.
So when all these protests started happening,.
when these marches began happening,.
I'm like, "I gotta be there..
"I'm a photographer..
"I can't one day teach my kid how to use a camera.
"and then not have photos.
"of some of the most significant historical events.
"that I've lived through.".
So I went..
I went to different protests.
that were happening around the city..
I documented it, and I took photos,.
and I partook in those..
I even accidentally came across.
some of the looting that was happening,.
and I documented that a little bit,.
although a little more scared..
And then what ended up happening next.
was quite possibly, quite possibly.
the worst experience I've ever had.

$^{81}$trying to photograph anything.
because in addition to protests that were peaceful.
and marches put on by churches.
and different organizations and groups,.
there was also a lot of panic and fear in our neighborhoods..
And that's what I stumbled on one night..
Sammy and I, my wife, were coming home from Carbone,.
which is the best hidden gem in Chicago..
It's right by Riccobini's,.
and it is so much better than Riccobini's,.
so if you're ever in town, go there..
It is the best..
So we were coming back from that,.
and we were driving on 31st and Princeton,.
and we noticed at this intersection.
was a massive crowd of people, like massive..
I'm talking 40 to 50 people,.
and remember, we've just come out of lockdown,.
so this is not normal, right?.
What is this massive crowd doing?.
And we finally drive through, and we get through,.
and we notice along our street.
there's a bunch of people sitting outside.
of their doorsteps, people sitting on lawn chairs.
right in front of their buildings..
People are kind of alert and waiting..
So we find out that apparently.
one of the very peacefully planned protests.
was going to walk through just a corner of our neighborhood,.
just a tiny spot of our neighborhood..
And the place where we live just erupted..
People were not okay..
That crowd, they were walking around with baseball bats,.
crowbars, table legs that they had unscrewed.
from their table..
I even saw a machete..
There were people sitting outside their buildings.
with pistols..
I even saw somebody with an AK..
I am not kidding, a semi-automatic weapon..

$^{121}$So being the genius that I am,.
and against the wishes of my much smarter wife,.
I decide I'm gonna go take photos of this..
How often do you see somebody walking along a street.
right next to a police officer with a pistol.
or a freaking machete?.
Not very often..
So I take my camera and I go out, I go down our street,.
and it's tense..
It's, well, not yet..
It's not actually yet tense..
Right now it's pretty light..
I get up there and it's by Franco's Restaurant.
and everybody is having a good time..
They're joking with each other..
They're showing off the table legs or the crowbars.
or whatever random assortment of weapons.
they decided to bring to this little gathering..
Oh, it's a little insane..
And they're all drinking too..
They're in this group and they have,.
like some of them have crowbars in one hand.
and beer in the other..
And it's getting to them..
So I'm a little anxious about this..
So I make sure to do all the best practices..
I say, "Hey, I'm just a photographer..
"I've been photographing everything that's happening.
"from the pandemic to the protests..
"Just so you know, I'll be taking photos..
"If you don't want me to take your photo,.
"just let me know.".
And I get a lot of people saying, "Yeah, you're fine.".
And then some people saying,.
"I'd prefer not to be in photos.".
I'm like, "Okay.".
And then we hear a noise..
And two young kids, I'm talking 10 and 13, tops,.
ride their bikes past our crowd..
And two of the guys in the crowd.

$^{161}$throw their beer bottles at these kids.
and yell, "We don't want you in our neighborhood.".
Everything just shifted in a moment..
The tone just took a 180..
And suddenly, things were not jovial anymore..
People weren't messing around..
There was this weird macho and this weird,.
almost like desire for violence.
or desire for letting out their aggression in this crowd..
It shifted massively..
And in that moment, some of the people.
who were a little more under the influence turned around.
and they pointed their weapons at me and they said,.
"You can't have that here..
"You need to get out of here..
"Some of these people have city jobs, man..
"They don't want me taking photos of them.
"doing illegal and terrible things.".
So I sense it and I try to push back,.
like, "Hey, it's public space," but they're not having it..
So I start to turn around.
and I don't wanna walk directly home.
because I'm only one and a half blocks away.
from where I live..
So I walk the other direction.
and two of the guys start following me..
My genius mind was like,.
"I should get the memory card out of my camera.".
So I do that and that's what I'm thinking of..
And I'm thinking, "I need to find a safe spot.".
I see a streetlight off to the corner.
and I head there thinking,.
"Surely no one's gonna mess with me.
"in broad streetlight 'cause it's dark at this point.".
So I get there and the two guys show up.
and there's a group of people.
kind of hanging around that streetlight.
and they all form a circle..
And out of the home by that corner comes this older guy..
And if you've ever seen a mafia movie.

$^{201}$and you think of like that old mob boss,.
that's the image of this guy..
About yay high, balding head, combed over..
Like I don't think he had a cigar,.
but if he did, it would have made perfect sense..
So he comes out, he's in his white shirt with his overalls.
and he decides to be the one.
to take the charge questioning me..
And on his shoulder is a sawed off shotgun..
So this whole crowd surrounds me.
and starts asking me questions about where am I from?.
Am I actually in this neighborhood?.
Do I live here?.
What, let me see your driver's license..
Prove that you're not from somewhere else in the city.
and you're not trying to like screw us over..
These guys start getting intense..
They start yelling at me..
I'm anxious, I'm worried..
I'm trying everything I can to calm them down..
And eventually I just say,.
"Look, I am just gonna leave, all right?.
"Clearly you don't want your photo taken..
"I will not take your photo..
"I haven't taken any photo since it got intense..
"I'm just gonna leave.".
And they're like, "Yeah, you do that.".
I leave and people trail me for just a little bit.
but then I zigzag around.
so that nobody can see where I'm going.
and I come home and I am shook up..
I am like anxious about this..
Even now thinking about the encounter.
it gets me really, really tense..
I haven't told this story.
because I think there are a lot of other stories.
at that time that were more worthy of telling.
and more timely and I don't think the world needed me.
to tell this story or anything like that..
It wasn't that groundbreaking or anything..

$^{241}$There were other narratives that needed.
to be lifted up during that time..
So I kept it to myself but the reason I'm bringing it out now.
is because the themes of that story.
of hatred and division and aggression.
tie in really well to the themes of the book of Jonah..
I won't talk about what that is now.
but I'll bring it back in just a minute..
So let's pray..
Oh God, thank you so much that you are our God,.
that you love us, that you protect us,.
that you have grace for us,.
that you have given us your word.
so that we may know who you are.
and that you give us wonderful gems of Bible books.
like the book of Jonah..
I pray that you'd be with us as we talk about it..
Be with me as I speak about it..
Don't let me say anything that is against.
or contrary to what you stand for, God..
Let each of us in this room examine our own hearts.
in the light of you and what you stand for, God..
May you have your way with each of us..
We love you, God, and we pray this in your name, amen..
Now, if I were to ask any of you.
what the book of Jonah is about,.
most of us think about VeggieTales, right?.
Most of us think about the pirates who don't do anything.
and we think about the VeggieTales movie,.
the bankrupt their studio, right?.
And in that movie, you have a bunch of vegetables.
singing fun songs about how awesome God is.
and how he gives second chances..
And it's so much fun,.
but it's also a very watered down version.
of the actual book, right?.
Trying to describe the book of Jonah as,.
oh, Jonah's scared of Nineveh, he runs away,.
God gives him a second chance.
and everything turns out great..

$^{281}$Describing it that way is like trying to describe.
Avengers Endgame as somebody steals jewelry.
from an old guy..
It's not really gonna cover all the bases..
It might technically be true,.
but it doesn't give you the full scope of what's there..
And that's a shame because the book of Jonah.
is crazy cool..
It is funny, it is extreme,.
it is almost satirical in nature..
And it is kind of one of a kind in the Bible..
And when we look at it and we kind of view it.
through the lens of a children's book,.
we do it a massive disservice..
I gotta give credit to Tim Mackey from The Bible Project.
because his work on unraveling and decoding this book.
and helping me see it in the way.
that maybe the person reading it at the time would have.
has changed the way that I look at the book forever..
So props to him..
Now, the part we wanna talk about today.
is at the very end of the book..
And we can't talk about the very end of the book.
without discussing what goes on for the rest of it..
So in order to talk about the rest of the book.
and to sum it up,.
I'm gonna use the most sophisticated.
storytelling tool I have, puppets..
(audience laughing).
Now, my wife was very much against this,.
but for me, that was very personal.
because for those of you who don't know, I was homeschooled..
And we don't have sports when you're homeschooled..
When you're homeschooled,.
you don't join the volleyball team,.
you join the puppet team..
I'm not kidding..
I was in a traveling puppet troop..
We had neon puppets and we were basically the varsity team..
I was pretty good..

$^{321}$In fact, the puppets I had,.
you actually had arm rods to deal with as well..
And they were way more advanced than these level one puppets..
But we're gonna work with what the Moody Library has.
and we're just gonna go with that..
Sound good?.
So bear with me as I give a bit of a comical recap.
of the book of Jonah up to this point..
Sound good?.
Cool..
Let's recap Jonah..
This is a bit of an overview that Tim Mackey did..
It is wonderful..
We're gonna go through the spark notes of this.
as we tell this story..
This is Jonah..
You can tell because he has a necklace of a whale.
and it says Jonah..
I don't think that's canonical,.
but the puppet says it, so we're gonna go with it..
This is Jonah..
And Jonah's a prophet..
For those of us who don't know,.
prophets of the time were basically.
God's mouthpiece to Israel..
So God would need to say something to people,.
but because God prefers to work through people,.
through his representatives on the earth,.
that's what it means that we're the image of God.
or images of God,.
he uses humans to work with other humans..
So he had prophets..
It was this class of people that would.
speak on behalf of God..
And Jonah was one of those prophets..
Now if you're in ancient Hebrew.
and you're reading the book of Jonah,.
when you say, oh, this is a prophetical book,.
you're expecting to open it up.
and hear all the proclamations of judgments.

$^{361}$that God has for Israel,.
or all the warnings or all the encouragements that he has..
If you look at the books ahead or behind of Jonah,.
they're just chapter upon chapter of the prophets.
saying to Israel what they need to know..
But then the book of Jonah is different..
It's a narrative..
And not only that,.
it's kind of intense and insane and extreme.
in the best way possible..
So at the beginning of the book,.
God sends Jonah a message..
It says, "And the word of the Lord came to Jonah..
"Go to Nineveh and that great city.
"and tell them the judgment I have for them..
"In 40 days, they're going to be destroyed.".
Now, Jonah doesn't like this..
Jonah dislikes this so much.
that rather than going to Nineveh over here,.
he turns the other direction and heads towards Tarshish..
He gets to a port of Joppa,.
he gets in the boat,.
and he sails away from where God has asked him to go..
We were going over this with our students.
at Urban Voice recently,.
and we were like, what do you think is Jonah's motivation?.
Why is he running?.
Why does he not want to do what God has asked him to?.
And the general consensus from the students.
and from just about anybody who reads it.
is that Nineveh's a bad place and Jonah's terrified..
What Nineveh used to do is when they conquered a city.
and they took the walls down,.
or they conquered a city.
and they would take the leader of that city.
and they would string,.
I don't know how graphic I should be,.
they would cut them up in creative ways.
and string them along the wall of that city..
So the leaders would not only be killed and bleeding out,.

$^{401}$but they'd also be publicly humiliated in the process,.
and everybody in that city would have a very clear reminder.
of what has happened to them..
So our students are like, well, Jonah's terrified..
They don't want that to happen to him..
I wouldn't want to go tell Nineveh,.
hey, you're gonna be destroyed in like 40 days..
So Jonah's running away..
He gets in the boat, he falls asleep,.
and then the Bible says God sent a massive wind,.
a great wind, a storm to hit this boat,.
and things are terrifying, things are freaky..
It's so bad, in fact,.
that the sailors, who are apparently.
on a different diet than Jonah, are terrified,.
and they begin praying to their God,.
and they're like, ah, what is happening?.
And they're like, is everybody praying to their God?.
Do we have all of our bases covered?.
What's that one guy down deck that's sleeping?.
Go wake him up..
And the captain comes and says, Jonah, pray to your God..
What are you doing?.
He's like, what, what, oh, okay..
So they're all praying to their gods..
Notice, Jonah ain't saying a word..
They're terrified, and Jonah knows what's going on,.
but he's just keeping his mouth shut..
He's willing to let them all figure it out themselves..
So they cast lots, they play a game of chance,.
and Jonah draws the short straw..
So in that day and age, oh, crap, our God's mad at you..
And Jonah says, yeah, I am a Hebrew,.
and I worship the God Jehovah, Yahweh,.
and basically, he's really angry at me.
for running away from him..
That's why all this is happening..
So they're like, oh my gosh, this is terrifying..
And they're such great people.
that they actually try to row the other way..

$^{441}$They're like, maybe we can row out of it..
And Jonah's like, no, what you have to do is,.
now at this point, you'd think Jonah maybe understands.
that God's mad at him..
Jonah maybe knows what he has to do..
Maybe Jonah's gonna tell them,.
hey, we should turn this boat around,.
go the other way, and I'll go to Nineveh..
But what's Jonah say?.
You should throw me overboard..
You should throw me into the sea and leave me to drown..
Jonah would rather die than go to Nineveh for some reason..
They don't wanna do it..
They try to row the other way..
It won't work..
And with a prayer to God, like God,.
forgive us if this results in this man dying,.
they throw him overboard..
The storm clears..
And the first subversive thing that happens,.
the thing we're not expecting,.
is that all the sailors are like, crap,.
that must be the real God..
And they make covenants towards him,.
and they make promises..
Those sailors become believers in Jehovah..
Even though they would have been of other nations,.
even though they would have been of other tongues.
and different religions,.
they come to know God through it..
So Jonah's drowning..
And this is the part of the story everyone loves..
God sends a fish,.
who is apparently a very happy whale,.
to swallow Jonah up, right?.
God is not going to let Jonah die,.
'cause honestly, Jonah dying would kinda be like him.
getting one over on God, like,.
ha ha, you can't make me go to Nineveh..
So God swallows, God has the fish swallow Jonah up,.

$^{481}$and while in the belly of the whale,.
Jonah recites different psalms..
He prays to God..
Now, on first read, it looks like he's praying,.
looks like he's repenting,.
but the more you read it, he never actually apologizes..
He never actually commits to be better, really..
All he's really doing is quoting psalms.
about how amazing God is, how forgiving God is,.
how salvific and loving God is..
And after three days, God tells the fish.
to vomit Jonah out onto dry land..
That's what the text says..
There's vomit in the Bible..
Vomit is a key element of this story..
So the fish vomits Jonah onto dry land,.
and Jonah decides to go to Nineveh..
So he's going into Nineveh,.
and we're told Nineveh is a great city..
It would take three days to walk all the way through Nineveh..
But Jonah, he walks one day into it..
Just one day..
It's a three-day trip city, but he goes in one day,.
and he says, oh, sorry,.
and he says, "40 more days,.
"and Nineveh shall be overturned.".
Hey, in 40 days, you're all going to be destroyed..
Not a very convincing sermon..
Jonah didn't put his best work towards this..
He's not exactly attempting to convince anybody..
He's fulfilling an obligation..
He's doing the barest minimal possible..
So he goes part of the way into the city,.
and he says, hey, you guys are all going to be destroyed.
in 40 days, and then that's it..
But the next part is also surprising..
We wouldn't have expected the sailors to turn to God,.
and we also wouldn't have expected Nineveh to,.
but what happens is that word gets all the way,.
oh boy, word gets all the way to the Ninevite king,.

$^{521}$and he's like, oh, crap, guys,.
we're going to be destroyed in 40 days,.
and he puts out a decree to all of the people.
in that nation..
He says, guys, we need to dress ourselves.
in sackcloth and ash, which is the thing you would do.
at that time if you were grieving and in despair,.
and he says, and we need to pray to this God.
and ask him, because maybe he'll relent..
And so the king has every man, even every cattle,.
every freaking cow, be in sackcloth and ashes and repent..
And at the end of chapter three, it says,.
and God saw what they had done and relented..
So now we're at the point of the story where we want to be..
Now we're at the very last chapter,.
and this chapter is kind of like a key..
It unlocks some of Jonah's motivations..
It tells you what he was thinking.
throughout this entire ordeal..
So that's where we're going to pick up..
But it displeased Jonah exceedingly, and he was angry,.
and he prayed to the Lord and said, oh, Lord,.
is this, is not this what I said.
when I was yet in my country?.
That is why I made haste to flee to Tarshish..
Pause for a second..
Jonah's telling us why he didn't go to Nineveh..
He's explaining it himself..
He wasn't afraid of Nineveh..
He knew that God was going to forgive Nineveh if he went..
Jonah continues, for I knew that you are a gracious God.
and merciful, slow to anger and abounding in steadfast love.
and relenting from disaster..
Therefore now, oh Lord, please take my life from me,.
for it is better for me to die than to live..
Jonah's throwing a temper tantrum..
Jonah is like a child who doesn't like a decision.
his parents has made, and the Lord said,.
you do well to be angry..
Jonah went out of the city and sat to the east of the city.

$^{561}$and made a booth for himself there..
He sets up tent..
He sat under it in the shade till he should see.
what would become of the city..
So remember he said in 40 days,.
you guys are gonna be destroyed..
He's camping out there for 40 days just out of spite,.
just because he wants to watch them be destroyed..
Now the Lord God appointed a plant.
and made it come up over Jonah,.
that it might be a shade over his head.
to save him from his discomfort..
So Jonah was exceedingly glad because of this plant..
Notice when God is giving Jonah some favor,.
Jonah's happy, he's a happy boy..
He's great with this..
But when dawn came up the next day.
and Nineveh was not destroyed,.
God appointed a worm that attacked the plant.
so that it withered..
When the sun arose, God pointed a scorching east wind.
and the sun beat down on the head of Jonah.
so that he was faint..
And he asked that he might die and said,.
it is better for me to die than to live..
Apparently Jonah keeps going back to this.
whenever he's in an argument..
But God said to Jonah,.
do you do well to be angry for the plant?.
He said, yes, I do well to be angry, angry enough to die..
And the Lord said, you pity the plant.
for which you did not labor,.
nor did you make it grow which came into being in a night.
and perished in a night..
And should I not pity Nineveh,.
that great city in which there are more than 120,000 persons.
who do not even know their right hand from their left.
and also much cattle?.
And that's the end of the book..
That's how it ends..

$^{601}$In some translations, which I love,.
it's and even the cows?.
Then it just ends..
You end with God being like,.
and what about the cows, Jonah?.
So this is just the weirdest, strangest dud.
of an ending to a book, but it is so fascinating..
It is so fascinating..
Because what this final chapter reveals about Jonah.
is that Jonah was not afraid that Nineveh would kill him..
Jonah was afraid Nineveh would be forgiven..
Jonah was so sure that if he went to Nineveh,.
like God asked him to, that God would forgive Nineveh..
And Jonah didn't like that..
And he was willing to run away for it,.
leave his entire life for it..
He was even willing to die for it..
What is interesting here is that Jonah refuses.
to accept a God who loves people that Jonah hates..
Say that again..
Jonah refuses to accept a God who loves people Jonah hates..
I wanna go back to that story, kind of..
I wanna go to a few days after that happened..
Samy and I were going out to eat.
once restaurants were allowed to open up..
And we were at a pizza place next to our restaurant,.
next to where we live..
And it's in the same neighborhood that this event happened..
And we're eating our pizza,.
and in walks one of the guys from that night..
In fact, this was one of the guys.
who even threw his bottle at one of the teenagers..
And one of the guys who followed me.
midway through when I tried to leave..
I noticed myself tense up..
I noticed my breathing get shallow..
I was getting anxious..
I asked Samy, "We finished our food?".
And then I asked her if we could leave..
And I remember, because he was the one.

$^{641}$who had been saying, "We don't want you here..
"These people have city jobs.".
And from that point, for the next few months,.
something in me had shifted..
Something in me had changed..
And something about the way that I viewed my neighbors.
and the people on my street had altered..
I didn't want to pass them while I was walking my dog..
I didn't want to talk to them when I was parking my car..
I didn't want to run into them at any of the businesses..
In fact, when I'm doing all of those things,.
there is a very real part of me that despises them..
There's a very real part of me that sees.
some of those faces that I remember from that night.
and feels very real anger and very real vitriol.
and almost hatred, almost like I despise their very being..
I hate the fact that they're in my neighborhood..
I hate the fact that I live among them..
And all these thoughts are going through my head.
as I walk my dog past them,.
as I shop next to them at the gas station.
or fill up my vehicle..
In that moment, guys, I had also become Jonah..
We read this book and we think it's a little insane.
or a little extreme,.
but we ourselves are prone to the exact same thing..
I could not look at my neighbors.
with the eyes that God looks at them.
because of my anger towards them..
I could not accept that God also loves them.
because of what I perceived as the kind of people they were..
And I lived like that for a while..
In this moment, Jonah believes he knows better than God..
That's the root of it..
God, I know better than you..
And if you were to make something bad happen to these people.
that would be better, that would be justice..
If you were to make it so that they couldn't do that.
to people anymore, that would be justice..
It's this belief that we know better than God..

$^{681}$And it's also the same exact motivation.
as Adam and Eve in the garden..
What do Adam and Eve in the garden do?.
God tells them something bad is going to happen.
if they do something, but they decide,.
I think something good is gonna happen if I do that..
So often in Genesis we're told, and God said it was good,.
and God said it was good, and God said it was good..
And then with the apple or the fruit.
that he's told them not to eat,.
it says they looked and saw that it was good..
That's a verbal code for us,.
that they looked and they're making the moral judgment call..
The disobedience didn't happen at the bite of the fruit..
The disobedience happened when they decided.
to pick it from the tree..
To pull this even closer to our story,.
at the very beginning it says,.
but it displeased Jonah exceedingly..
See that little A?.
That means there's a different way.
that it's said in certain passages..
Some passages don't say that it displeased Jonah exceedingly..
Some passages say Jonah found it exceedingly evil..
And what happens right before this?.
God relents, God does not destroy Nineveh..
So God doesn't destroy Nineveh,.
and in some of the passages we have,.
it was seen as exceedingly evil to Jonah..
Jonah is now the one making these moral judgment calls..
But who are we to look down on Jonah for that?.
'Cause we do that all the time..
I know that this has happened to me,.
and I know it's probably happened to you..
People who have hurt us, people who have maybe betrayed us,.
or people who have inconvenienced us..
We have this desire in ourselves to not see them be treated.
with the same love and care that we ourselves have known..
Here's the other weird part about this passage..
Belief in God was not Jonah's problem..

$^{721}$Jonah was very aware of God's existence..
Jonah knows God, Jonah hears from God,.
yet he doesn't understand God's values.
and isn't on board with his mission..
A person can fully believe in God,.
yet not understand his values.
or be on board with his mission..
That should terrify us..
There are people walking around, sometimes even us,.
with the Christian word on their shirt.
who are hurting people, who are causing pain or division,.
and all the while they have the name Jesus on their shirt..
That should scare us..
And it should remind us that belief in God,.
belief that he's real and belief that the things.
in the Bible are how they happened,.
isn't the thing that saves us..
It's us saying we submit ourselves to your mission,.
to your values, and to your way..
You are our God, not me..
When we believe we know better than God,.
we are choosing Adam, not Jesus..
Interesting side note, our students, some of them,.
I've had interesting conversations with them.
where we'll talk about what it means to be a Christian,.
and almost none of our students are hung up.
on the idea of believing in Jesus..
Some of them wanna talk apologetics..
Some of them wanna talk about the archeology and all that,.
but the majority of them, the majority of,.
I would say people, the thing you can't get behind.
isn't the belief about whether or not it's real,.
but it's the submission to another being's way.
of doing things that's hard..
And we struggle with that,.
even though we may believe in him..
Here's another thing we learn from this passage.
is that we cannot fully love others like God loves them.
until we learn to see ourselves the way God sees us..
We cannot fully love others like God loves them.

$^{761}$until we learn to see ourselves the way God sees us..
When we view ourselves as higher than or better than others.
because of whatever metric.
is in front of our noses at that time,.
we are doing the same thing the Pharisees did..
We're doing the same thing Jesus rebuked the Pharisees for..
We are saying we are better.
because of something we've earned,.
because of something we've done,.
because of something I am,.
and you are worse because of something you are.
or something you've done..
Paul in the New Testament,.
an apostle who would write.
to many of the different churches,.
he understood this well..
He understood that he can't love other people.
until he sees himself the way God sees him..
Paul has this really interesting message at one point.
where he says, "I am the chief among sinners.".
And that's crazy and weird to the people of the time.
because this is the guy who's leading the charge.
in bringing the gospel to Gentiles..
This is like earliest Christian celebrity, right?.
That guy, the one who's writing most of the New Testament?.
And he says, "I am the chief among sinners..
I'm the greatest sinner I know.
because I know my thoughts..
I can't read your thoughts..
And I know who I am and I know how dark I can be..
I know how evil I could become.".
And that sort of knowledge and understanding of oneself.
leads to a humility that says,.
"I can love other people and love them better.
and find them more important than I find myself.".
Paul had that, Jesus had that..
Sometimes we struggle to have that..
Finally, we don't get to decide who is worthy of God's love..
We're not the gatekeepers..
We are not the one who makes the rules..

$^{801}$When we accept that God has loved us.
and done something for us.
to bring us into a better state of being,.
to bring us from a state of death to a state of life,.
we agree with God that that is a possibility for everybody..
Even the people who we despise,.
the people who have hurt us, who have betrayed us,.
the people who are on the opposite end.
of our political spectrum,.
the people who do evil things in the world,.
we have to believe with God.
that they are as eligible for his grace as we are..
We can't be the ones who decide.
who is and who isn't worthy of that..
And the way we go about our lives needs to reflect that..
I know that's sobering for me..
I say all this,.
I say all this having told you about this struggle I've had,.
especially in the past year with my neighbors..
And I wanna tell you how this sort of ended.
and wrapped up for me..
One of the things I love about our neighborhood.
and about the churches in our neighborhood.
is that we do this thing called Pray Bridgeport..
And it's amazing..
It's all the churches in the area.
who are willing to come together.
and we pray for our neighborhood..
And it's freaking beautiful..
We got Anglicans over here, we got Baptists over here,.
we have people of all different like cultures,.
nationalities, all different denominations,.
and we're all coming together to pray.
that God would love the people in our neighborhood..
And it's amazing..
And a few months after that incident in my neighborhood,.
I was at a Pray Bridgeport at our park at Armor Square,.
which is right where I live,.
which is right where this incident had happened, basically..
And they had set up prayer at different stations..

$^{841}$So you would go to one station.
and pray for a specific thing,.
go to another station and pray for a specific thing..
And I was going about these stations.
and I get to a station.
that is praying for the political.
and social divide in our neighborhood..
And it's run by some of the pastors.
at Progressive Baptist Church,.
who are some of the people who are organizing.
the peaceful marches and the protests of that time..
They're organizing these things.
and they lead us in prayer for the very people.
who I had just been hating in my own heart for so long..
The very people I didn't wanna think about,.
the very people I was honestly terrified.
that God would ask me to talk to at some point.
because I didn't want to.
and I didn't want them to know about God..
I didn't want to have to be forced to talk about him.
to someone who I find so despicable..
And these people who were even organizing.
these protests and marches,.
lead us in prayer for that, for them asking God,.
would you work in their hearts?.
Would you help them to see the error of their ways?.
And would you grip them for a greater purpose?.
Make them people who bring about good in your world..
And that was the prayer that I couldn't pray at that time..
And it was me coming face to face with the idea that.
Jesus loves them as terrible.
as they might be at their worst..
I'm also terrible at my worst..
And let me clarify, this is not to say in any way.
that Jesus permits some of the kinds of things.
they were doing or the violence or the animosity,.
but that if we have a biblical example.
of God using one of his prophets.
to reach such an evil city as Nineveh,.
we need to know that God can reach.

$^{881}$even those who we see are evil..
And if we're not getting on board with that,.
we are being against God and trying to gatekeep.
and hold things that we want for ourselves.
and not letting other people have it..
Here's the final point..
Jesus is kind of like the new Jonah..
Do you guys remember when Jonah.
was in the belly of the fish?.
He was in the belly of the fish.
for three days and three nights..
Our alarm bells should be ringing..
That is like a biblical hyperlink.
to another theme in the Bible..
Who else was in the dark for three days and three nights?.
That's Jesus..
We're supposed to see parallels between Jonah and Jesus..
Both are, in a sense, messengers to a people.
delivering a message from God.
of repentance for their own good..
And Jonah goes into the belly of the fish.
for three days and three nights,.
and Jesus goes into the grave.
for three days and three nights..
But the difference between the two.
is that where Jesus decides.
that he is going to take up this cross.
to bear this mantle and to love the people he's sent to,.
Jonah despises them..
Jonah fights tooth and nail against God,.
even as he's given the job.
of letting these people know that they need to repent..
But Jesus walks to the cross..
Jesus is consistent..
Jesus is steady..
Jesus continues to love the people.
who are doing him active harm..
And when Jesus is strung up on the cross,.
in much the same way the Ninevites.
would have strung up their enemies on their walls,.

$^{921}$Jesus utters, "God, forgive them,.
"for they know not what they do.".
May we have that kind of a heart..
May we have that kind of a desire for people,.
that when we see evil.
and when we see bad behavior in the world,.
we don't despise them,.
we don't stoop to their level of making our own enemies,.
but we understand that God's transformational love.
and his work in our hearts.
is something that also could be done in theirs..
And may we hope for that..
It might not always happen,.
but might we pray for that and desire that,.
not just for the people who we see as worthy,.
but for the people who God loves, which is everybody..
If you're a Christian in this room,.
at all times, as an ambassador of God's kingdom,.
you have a choice to be a Jonah or to be a Jesus..
Not that you can be Jesus, but you know what I mean..
Are you going to love the people you have been sent to,.
sacrificially trusting that God is good,.
or are you going to fight him tooth and nail.
for your own way?.
May we choose to be more like Jesus and not Jonah..
Let's pray..
God, thank you for this convicting passage..
Thank you for this strange novelty of a book..
Thank you that it's so different.
from many of the other parts of the Bible..
And thank you that it is such a mirror to us..
Thank you that when we look to it,.
it is so subversive and it points out.
the flaws in our own heart so well..
Thank you that it allows us to see.
the areas where we need to grow..
And God, I pray for each of us,.
would you give us a heart that is similar.
to the heart of Jesus,.
that when we look at people who are lost,.

$^{961}$when we look at people who are in the dark,.
we don't despise them for what they might do.
to ourselves or to other people,.
but would we have pity and compassion and empathy.
and seek to show them the love that you have first shown us?.
God, we pray that you would use us mightily..
We pray for each of those people.
who we all maybe have a difficult time loving..
Would you be with them?.
Would you work in their hearts?.
And also would you work in ours.
to know how to love them well and to represent you better?.
We love you, God, and we pray this in your name..
\newpage



\section{}
\label{sec:GChzvOXTH_M}
\textbf{Anointed by God - Part 7 (1 Samuel 17) - Tommy Lee}
\newline
\newline
連結: \href{https://youtube.com/watch?v=GChzvOXTH-M}{\texttt{https://youtube.com/watch?v=GChzvOXTH-M}} ~~~~ 語音日期: 2021-11-16
\newline
\newline
\hyperref[sec:MkL9xiQNZJM]{\small{< < < PREV SERMON < < <}}
~
\hyperref[sec:index_chronic]{\small{[返順時目]}}
~
\hyperref[sec:index_scriptual]{\small{[返順卷目]}}
~
\hyperref[sec:eQLpWVai5F0]{\small{> > > NEXT SERMON > > >}}
\newline
\newline
$^{1}$- You know, one of the things, good morning everyone..
My name is Tommy Lee and I get a chance to lead.
this wonderful organization, Resource Global.
and the Growth Center and we get a chance to partner.
with Pastor Mark and his church..
I love the humility of Pastor Mark.
and I love how even simple things of just coordinating.
some of these setups here on Sunday morning.
is something that I've learned from Mark as well too..
Thank you James and you know,.
it's my first time meeting James in person..
I got a chance to meet Phoebe also in person..
Our whole life these days is online..
Our whole life is on Zoom and everything like that.
so when you get a chance to meet each other in person,.
it is a joy..
I will probably say James is, when we chat with him,.
he had graduated from Moody Bible Institute in Chicago.
and now he works in the marketplace..
And I think all of us in this day and age,.
this is the world of the modern day pastor,.
don't you think?.
And a lot of times there will still be a need.
of the full-time pastors, there will still be a good chunk.
of those who are full-time pastoral ministry.
who went to seminary and are doing that full-time.
as a vocation..
But more and more as I encounter leaders.
from different parts of the world,.
they're wearing different hats..
They're your bi-vocational pastor,.
whether or not their church won't be able to afford them.
in the full-time salary, so they're doing that.
and then from that point on, they're doing a job.
or there are different business leaders out there.
who are working in the area of finance or accounting.
or tech, but they love the church..
And so they're doing that full-time and then on the weekends.
or on the time, they're working..
There's this one friend I have, Shundren Thomas..

$^{41}$Shundren is the president.
of Northern Trust Asset Management..
Shundren is so successful in the area of private equity,.
but yet on the weekends, he's the associate pastor.
of his dad's church..
There's another guy, Dr. Nicholas Pierce..
He is a MBA, executive MBA professor.
at Northwestern University..
He does a great job in the area of education,.
but yet on the weekends, he's the associate pastor.
of his church..
We're seeing that more and more as well too.
and that's why a lot of times with the Growth Center,.
I get a chance to partner with Pastor Mark.
and do all of these wonderful things.
that you guys are doing..
Hey, let me just, before we dive into our sermon today,.
two, couple of things I did wanna share with you.
as well too..
Let's bring up the first couple slides here..
Next slide here..
We are starting a certificate..
I went to, for me, I spent 10 years building cell towers..
That was my job, those big telecommunications tower..
That was what I did for 10 years..
But during the evening, I went to seminary.
and it was three grueling years of evening classes..
I walked away from seminary realizing.
that sometimes a seminary is not something.
that people can afford..
It's not even something people can have time for..
So what are some of the core classes that you need.
from some of the key teachers, whether it's Bible,.
theology, but leadership, culture, justice,.
integration, faith, and work that you need.
that allows you to really be able to gain the knowledge.
and education that you need..
So we've actually created a nine-month cohort program..
If some of you guys who are listening.
or sitting here who are interested,.

$^{81}$I'll send you more information that Pastor Mark.
to send over to you once a week for nine months.
and you could decide, hey, I wanna do month one,.
but I don't wanna do month two..
I could do month four taught by some practitioners,.
marketplace leaders..
You have people like Ashundra Toms who I was talking with,.
Jillian Anderson who works at Gallup..
Her job is working with Southwest Airlines.
in the area of Strength Finders..
You have Pastor Albert Tate who does a great job.
of expository preaching in LA..
Andy Crouch who wrote the wonderful book,.
"Culture, Making, and Power," and he does a great job.
talking about the role of Christians integrating in culture..
Or Dr. Crawford Luritz who teaches a wonderful class.
in leadership that he taught over at Trinity,.
"Leadership as an Identity.".
Another one, next slide as well too,.
that I wanna point out is Mr. Dave Evans..
If you have ever never talked to Mr. Dave Evans,.
he worked for Steve Jobs many years ago,.
created the mouse for Apple..
Now he teaches at Stanford University.
called "Designing Your Life," a wonderful class.
that is always filled with people as well too..
All talking about design thinking,.
how God uses design thinking to really be able.
to help you process the calling in your life..
So that is it..
I will send you guys more information..
Let's go, actually, let's hold off on all the slides..
Two other things as well too..
Is number one, and again, ministry has changed, doesn't it?.
And so I think a lot of times one of the things.
I've also been wondering, and I love what Urban Voice.
is doing with your church, I've been wondering, Mark,.
how do you plant churches in the meta universe?.
We've now done it over in Zoom, and we're getting used to it.
but around the world with government restrictions,.

$^{121}$access, whatever it is, Roblox has become.
one of the most popular platforms in all of the world..
I sat down with a private equity analyst.
who studies all of the Asia markets,.
and she says, Tommy, you need to look at how ministry.
interacts with the world of Roblox in the meta world..
Because a lot of times Katy Perry just did a live concert.
there and had over three million people attend..
Another individual on Fortnite, Travis Scott,.
had a live concert there and had 30 million people attend..
Another individual over in Bellevue in Seattle,.
or in the Washington area, created a youth ministry..
He draws 40,000 youth from around the world.
for this youth ministry..
And so as a result, as we begin to look at ministry,.
what are some of the creative ways,.
because the traditional church model is decreasing..
A lot of the American churches,.
I speak mostly for American church,.
their decline are coming from the younger people.
who are disengaged or frustrated with where the church is..
So a couple of thoughts as well too,.
but look, it is always in ties in with our sermon.
as well too..
All of us are anointed by God to do something unique..
All of us are called by God to contribute to this world..
How that looks like for each one of us.
is gonna be a little bit different..
Today is a part two of our sermon called Anointed by God..
We're looking at a popular story that all of us,.
whether you're sitting here or you're sitting at home.
watching this, is a story of David and Goliath..
A story that all of us have always been familiar with..
Besides the crucifixion, besides the birth of Jesus,.
David and Goliath is a story that we've grown up hearing..
The story of the little boy with the slingshot.
who takes on the nine foot giant in Goliath and wins..
And later on becomes the king of Israel..
We love that story..
It has a great happy ending..

$^{161}$But within that story, there are some unique things.
that God teaches us and how God anoints a young boy..
And through that journey, he allows him to have.
one turning point in his life.
that changes the course of his life..
He is anointed by God..
And we talked about part one two months ago..
It may seem like eternity, part one two months ago..
And today we're gonna talk about part two.
of anointed by God..
Some of the things, just to refresh your memory.
of what we talked about part one.
and we'll go on to part two..
I talked about three things,.
six insights on being anointed by God..
Two months ago we talked about the first three..
So let me just review it..
Number one, when God looks at the heart,.
when God anoints a leader,.
he looks at the heart of the individual.
versus the outer appearance..
When God anoints a leader,.
he looks at the heart of the individual.
versus the outer appearance..
And we talked about the story of Saul..
If you guys all remember Saul,.
Saul was a man from Benjamin..
Benjamin had a wonderful reputation..
Benjamin was the elite of all of the tribes..
Benjamin was a very wealthy tribe..
That's where Saul came from..
And Saul came from an extremely wealthy family,.
an extremely powerful family,.
and now he was brought in and selected as king..
But sometimes having the outer appearance, right,.
doesn't qualify you to be king..
Having the talents, having the smarts,.
having the wealth by appearance alone,.
outer appearance, does not qualify you for leadership..
And as a result, through a series of mistakes.

$^{201}$that Saul makes, disqualifies him for leadership..
And so then God sends Samuel the prophet.
over to Bethlehem to meet Jesse and his family..
And he says, "Jesse, there when you meet this guy,.
"a Bethlehemite by the name of Jesse,.
"you will find your first king.".
And there Jesse goes, and the story goes,.
some of you guys, he brings all of his son,.
Jesse calls all of his son, the first three oldest boys,.
and the first one is Eliab, and he looks at Eliab,.
he is the oldest, he is the strongest of all of the boys,.
and Samuel says to him, "Surely, surely, that is the king.".
And God reminds him, says to Samuel,.
"Do not look on his appearance,.
"or on the height of his stature,.
"because I have rejected him..
"For the Lord sees not as man sees..
"Man looks on the outward appearance,.
"but the Lord looks on the heart.".
And so Eliab was in it..
And remember, later on, Eliab becomes very key in the story..
And he looks at his second son, that was in it as well too..
And he goes for the third son, that was in it..
And Samuel then says, "Jesse, is there anyone else?.
"Oh, I forgot, there's David, the young boy.
"who's out in the shepherd field, call him.".
And there God, Samuel anoints him as the king..
Story goes on, sometime later on, a battle takes place..
And this battle takes place in the Valley of Eliab,.
between the Philistines and the Israelites..
One of the things, if you guys remember,.
that I talked about, the differences.
between the Philistines and the Israelites,.
and we're gonna talk a little bit about.
why Eliab was so important in today..
But look, the Philistines were coastal people..
The Philistines were sophisticated cultural people..
The Philistines learned how to use iron for weapons.
and tools and maintain monopoly on the markets..
They were masters at working with iron..

$^{241}$As a result, they had the technology and tools.
to make them one of the strongest,.
one of the most elite tribes at that time..
They could make shields, they could make swords,.
they could make spears, they could make armor, everything..
Compare that to the Israelites..
The Israelites were primitive in comparison..
They were not a professional army..
They were shepherds and farmers..
They had outdated weapons..
Probably the only two that had even any armors.
or shields or swords was Saul, their leader, their king,.
and Jonathan, their son..
That's it..
So you're having a ragtag bunch of guys.
who were farmers and shepherds against this group,.
and the odds were great..
Why Eliab, as you said?.
Guys, if you get a chance,.
take a look at this particular map,.
and you'll see, next slide, right there..
Look at where Gath is..
Look at where Azekah is in Soko..
That right there is the Valley of Elah..
If you look at north, what do you see there?.
It leads to a lot of things..
The richness of the Jordanian Valley..
Thank you, guys..
And the reason why is where Israelite was positioned.
led to a passageway into the Jordanian Valley.
into much riches..
And so the Philistines had to defeat the Israelites.
and get rid of them in order to gain access.
into this passageway..
They were standing in their way..
And so they brought all of their armies.
from Gath over to the Valley Elah..
And this is why they were so weak..
They were so weak..
The Israelites were so weak..

$^{281}$The Philistines are looking at it and says,.
you know what, they're not worth our time..
We're not gonna get our guys into battle.
and try to fight them..
Let's just get it easy..
So they called in their big giant Goliath.
and say, Goliath, we're gonna call, bring you in..
And Israelites, bring in your best champion,.
which was supposed to be Saul..
And we're gonna have a duel..
Whoever wins gets access..
If we win, Israelites, we'll go away..
But if you lose, you go away and we take over..
Easy..
So for 40 days, this happens,.
which leads us to the second point..
That you never know when God is calling..
Because if the first point that we had.
is that when God anoints a leader,.
he looks at the heart of the individual.
versus the outer appearance..
The second point is this one..
You never know when God will come calling..
Here is David who just got anointed by God.
as the next king..
Samuel whispers to him, you're gonna be the next king..
And what happens?.
Does David then become the next king?.
No, he goes back to the field..
He becomes a shepherd boy..
And then Jesse comes calling him and says,.
hey, come here, come here..
I need you to bring food to your brothers..
And from that point on, bring these 10 cheeses.
to the commander and their thousands..
Bring some food to your brothers, make sure they're well..
And bring some token from them..
You never know when God is calling you.
to do something amazing in your life..
So God looks at the outer appearances..

$^{321}$But yet God calls you..
We look at the story of Nehemiah at age 60 something, right?.
He's a cupbearer to the king..
And now he's called because he has a burden.
to rebuild the walls of Israel..
Peter, James, and John are fishermen..
God calls them to take up and leave their job as fishermen.
and to come and follow me..
But then we also see Zacchaeus, don't we?.
Zacchaeus is a tax collector,.
not just any tax collector, but the chief tax collector..
And somehow or another he encounters Jesus.
because there's something about his life.
that he doesn't like to see..
And what happens with Zacchaeus?.
God calls and says, Zacchaeus,.
I need to spend time in your home..
But yet he doesn't call him into full-time ministry..
He doesn't call him just like he calls James.
and John and Peter to pick up everything and follow..
Matthew calls him as a tax collector,.
pick up everything, follow him..
But Zacchaeus is different..
Zacchaeus, I need you to go back.
and become the tax collector again..
But I want you to be a better version of the tax collector..
You never know when God is calling..
You never know what he's calling you to do..
And that calling may not sometimes be vocational ministry..
Maybe you're calling like I was talking with James..
Graduate from seminary.
and now work in the corporate marketplace,.
but still using his skills as a worship leader,.
a wonderful worship leader, and in ministry as well too..
It all looks differently as well too..
For me, I grew up, I worked 10 years.
in the corporate world building cell towers..
I always thought I was supposed to be a minister..
I thought I was supposed to be a pastor..
After three years, I realized,.

$^{361}$God, you're not calling me to be a pastor..
That is not what I'm called to do..
Worked at Moody Bible Institute as seminary.
for a number of years under the president..
I thought that was it..
After three years, that ended, and God closed that door..
And now for the last 12 years,.
I get a chance to run multiple organizations,.
not one, but two, four, five organizations,.
because I realized God had wired me differently..
That the traditional way of doing church.
was not what my calling was..
But I get a chance to work with marketplace leaders..
I get a chance to launch new organizations..
That's what God was calling to do..
You never know what God is calling..
You're always sitting there looking and seeing.
and asking God to really work in your life..
And the third thing, the anointed leader.
will be criticized by others..
That you won't make everyone happy,.
and your goal is not to strive to make everyone happy,.
because we saw the story of Eliab,.
where when he sees David, Eliab, the oldest brother,.
he sees David and says, why have you come down?.
With whom have you left those few sheep in the wilderness?.
That God's anointed can be criticized.
because we're never gonna meet everyone's expectations..
There are things that we do or don't do.
that's not gonna make everyone happy..
We're gonna not meet the expectations.
that people have of us..
And so God's anointed needs to learn how to,.
as we talked about, needs to be thick-skinned..
You can't take things personally..
You have to learn to hear what people have to say,.
learn to listen to feedback,.
but you have to be thick-skinned.
and continue to strive to do what is best.
to honor the Lord in everything that you do..

$^{401}$And there it goes to verse 34..
But David said to Saul in verse 34.
in 1 Samuel chapter 17, in verse 34,.
in the story of David and Goliath,.
David said to Saul, as soon as he volunteers.
to kill Goliath, Saul says, what in the world are you?.
Who are you again?.
You're this tiny little shepherd boy.
coming to bring cheese to his brothers..
But David said to Saul,.
your servant, in chapter 17, verse 34,.
your servant used to keep sheep for his father..
And when there came a lion or a bear.
and took a lamb from the flock,.
I went after him and struck him.
and delivered it out of his mouth..
And if he rose against me, I caught him by his beard.
and struck him and killed him..
Your servant has struck down both lions and bears,.
and this uncircumcised Philistine.
shall be like one of them,.
for he has defied the armies of the living God..
And David said, the Lord who delivered me.
from the pawl of the lion and from the pawl of the bear.
will deliver me from the hand of this Philistine..
And Saul said to David, go and the Lord be with you..
Which goes to my fourth point..
God uses your past to prepare you for the future..
Why do I say this?.
Some of us are always in the midst of looking.
to see what God is doing, what God is calling us to do..
We're anxious..
Across the board, when I travel,.
I get a chance to travel to Southeast Asia, Nairobi,.
I hear the same thing from young leaders,.
especially those in their mid-20s and their mid-30s..
They always say to me, Tommy,.
I wanna know what God is calling me to do..
I just wanna know what God is calling me to do..
And one of the things I've always had to tell them,.

$^{441}$I said, you gotta learn to be patient in this sense..
That sometimes God doesn't tell you what to do..
God, sometimes God leaves you in limbo.
because remember, you are on a journey.
that he wants you to be..
And so as a result, he wants you to learn to be faithful.
to where he has placed you right now..
And he uses that experience to prepare you.
for something that he wants you to do..
I say this, we look at scripture, Nicodemus, right?.
Nicodemus was in his mid-60s before what?.
Before God says, I need you to build this wall for Jerusalem.
but what was Nicodemus doing this whole time?.
He was a faithful what?.
Cupbearer to the king..
He served the king faithfully..
He did his job well..
He was a great employee..
And when he finally got that burden to rebuild the wall,.
what happens?.
His boss, the king, who Nicodemus served faithfully well,.
was his first and only donor, funded his entire project,.
gave him the vacation days and the time off.
to go and do his job because what?.
He was a great employee..
When I graduated from Washington University in St. Louis,.
I graduated with a business degree..
And there I was looking for a job.
and I got hired in my job.
and it was a company by the name of American Tower.
and they built cell towers for a living..
And I thought, great, I come from a Ivy League school,.
Washington University in St. Louis..
What is my job?.
Am I gonna oversee a team?.
Am I gonna oversee projects?.
Tell me what it is..
They said, yes, your first job..
You see all those files that we just got?.
Verizon Wireless, at that point in time,.

$^{481}$we acquired 200 towers and all of their documents,.
all of their files in one room,.
unorganized construction drawings,.
all of this random drawings..
You see that?.
Organize that room..
What?.
I just graduated from university..
I graduated with a management degree..
You want me to do what?.
Lord, what are you calling me to do?.
What can I learn from this?.
So I said, Lord, if you would put me in here,.
I'm gonna do the best job I could possibly do..
And I organized that file room pretty well..
But more so, I learned how to be organized..
I learned the discipline of spotting.
and creating structures that I never learned before,.
day in and day out..
For over six months, I did that..
And so after I did that, my boss said to me,.
you did great, now I need you to do something different..
I need you to create a system to scan everything..
Really?.
I gotta scan everything?.
Back then we had scanners and we had to take card copies.
and sit there and scan it..
And so I did that for another four months..
Oh, God, I look back now, it's 12 years later..
God used that time in my life to create me a discipline,.
to be organized..
He helped me to spot how to create patterns,.
how to organize different things,.
how to organize different projects,.
how to really be able to use Excel, all of that..
To this day, I still use that..
One of the things that people tell me all the time,.
they say, Tommy, you organize your team so well.
and you are able to keep in line everything..
How do you manage so many different projects?.

$^{521}$And I think back to all those months.
of being in a file room, learning grudgingly.
and saying, Lord, if you've placed me,.
or how can I be faithful to where you have placed me.
in my life?.
You never know..
We sit there and we grumble..
This is, we're above this..
No, where has God placed you?.
How can we be faithful?.
How do we continue to be a good witness,.
to honor and faithful to where he has placed us.
in whatever job?.
Because you may be sitting here right now,.
watching this, you're saying, I can't stand my job..
I can't stand my boss..
I hate everything that I'm doing..
What is God teaching you?.
Why does he have you right there?.
Why is that so important during this first time?.
What is he showing you?.
God uses your past to prepare you in the future..
One little thing, I get a chance to work.
with young leaders around the world,.
especially Indonesia, I love Indonesia..
I get a chance to work with these young leaders.
who are now starting businesses.
and they're growing in massive, not only wealth,.
but power that will influence their city..
They're running family businesses.
and what they call unicorn companies that are growing big..
And I get a chance in the honor of serving.
and working alongside with them to integrate their faith.
in everything that they do..
But you know why I'm doing that so well,.
or I've gotten a chance to work?.
Because I sat in a church in Chinatown, Chicago.
for 18 years, working with youth in high school and college..
And I learned how to work with them.
because during those 18 years in that private.

$^{561}$with those 100, 200 students that came in.
and out every single year,.
I learned what it meant to be a leader..
I learned what it meant to work with developing leaders,.
working with younger leaders,.
what it meant to work with volunteers,.
what it meant to plan games,.
what it meant to plan events,.
to sit there and plan, all of those things.
led me to what I am doing now..
You never know, you never know..
God has powerful plans,.
and here's my encouragement for those who are watching,.
God has powerful plans in your life..
But he may not have revealed it to you right now..
So are you faithful to where he has placed you.
at that moment because God uses your past.
to prepare you for the future..
Second thing is in 1 Samuel 17, 38 through 40,.
then Saul called David with his armor..
He put a helmet of bronze on his head.
and clothed them with a coat of mail..
And David strapped his sword over his armor,.
he tried it, and he tried in vain to go.
for he had not tested them..
Then David said to Saul, "I cannot go with these,.
"for I have not tested them.".
So he put them off, then he took his hand,.
took his staff in his hand,.
and chose five smooth stones from the brook.
and put them in the shepherd's pouch..
His sling was his hand and he approached the Philistine..
Let me show you another picture..
Let's show the last slide over there.
of what that sling looked like, that slingshot..
Right there, that's what the slingshot did look at..
Couple things..
Ironically, Saul chooses to give David his armor.
and sword and everything like that..
Here's why..

$^{601}$Let me read to you, Judges 20, verses 15 through 17..
Judges 20, 15 through 17..
The people of Benjamin, remember which tribe Saul was in,.
the people of Benjamin mustered out of their cities.
on that day 26,000 men who drew the sword.
besides the inhabitants of Gibeah.
who mustered 700 chosen men..
Among all those were 700 chosen men from Benjamin.
who were left-handed..
Everyone could sling a stone at a hair and not miss..
The tribe of Benjamin were experts at the slingshot..
Ironic, isn't it?.
That Saul, where is he from?.
The tribe of what?.
Benjamin..
The very tribe who knew how to use a slingshot,.
they used it in battle..
They were skilled marksmen with the slingshot..
He gives David a sword and a shield, armor, all of that..
David puts on him..
Does it work?.
Lesson here is this..
What works for others may not necessarily work for you..
What advice people give you, what they tell you to do.
may not necessarily work for you..
And does that mean as leaders we sit there.
and we don't listen because hey,.
we're gonna do it our way?.
No, you listen..
One of the best things, qualities about leadership.
is you're listening to people..
You're always trying to improve yourself..
You're always trying to be the best version.
that God calls you to be..
But in the midst of that, always knowing that sometimes.
the answers that they give you, it may not work for you..
That's how it works..
Let me give you an example..
My mother is trained in Hong Kong as a teacher..
And she has a very, very stern way of teaching..

$^{641}$Growing up, we always had to get straight A's..
And she had a chopstick on her finger.
that if we got a math problem wrong,.
that chopstick went on her finger..
If we got a spelling word wrong,.
we had to sit there and write it out 100 times..
Well, and she's always reminding us,.
you have to get straight A's, education, education,.
all that you do..
Well, now I have a daughter and she's seven years old..
Well, she goes to my mom's school over here in Bridgeport,.
here in Chicago..
My mom takes her way because that's her way.
of doing education, that she believes.
is work on straight A's..
And she tries to use this on my daughter..
Not to say it doesn't work,.
but I keep telling my mom all the time,.
you have to understand that getting straight A's.
is not always the end all..
Just driving, studying without interaction.
is not always the end all..
But in her mind, no, no, no, this is the best..
You have to study all the time..
To the point, my daughter doesn't like going to her house.
anymore because my mom is at the same school..
She picks her up and she sits there.
and she makes her study from three o'clock.
all the way to five and then when she wants her.
to go to her house on Sundays to do homework,.
what do you think is gonna happen with my daughter?.
She cries every single time she has to go.
to her grandmother's house because she doesn't wanna go..
And I keep telling my mom, maybe you need to lighten up..
No, no, you need to make them study all the time..
It's gonna look different, folks..
What worked for them may not necessarily work for you..
And so as a leader, hear from people..
Listen to what people have to say..
But you gotta figure out, is it gonna work for me?.

$^{681}$Is it gonna work for me?.
In all of the leadership that you have..
As a leader, those who are anointed by God,.
you're gonna hear advice from problems,.
you're gonna hear solutions to the problems.
that you deal with all of the time..
And some of them, listen to it..
Some of them, will it work?.
Probably not..
You gotta figure out that decision..
But as leaders, on the other flip chart,.
if you have a team, don't insist that your way works..
Don't insist that your way works..
What works for others may not work for you..
Your job as a leader, for God's anointing,.
God has given you people to shepherd.
to help them find their voice,.
to help them learn how to think..
I work for one ministry leader,.
and Caucasian guy, older guy, his way was always the best..
His style was always the best..
You couldn't make the decision unless you do it,.
run it by him for him to give you the okay..
That way worked for him, probably for a period of time..
Doesn't work for everyone..
What works for others may not necessarily work for you..
And the last one I want to look at,.
is the battle goes on, David takes a slingshot,.
from afar, he defeats Goliath..
And David said to the Philistine, right before that,.
you come to me with a sword and with a spear,.
and with a javelin, nine foot tall Goliath..
But I come to you in the name of the Lord of hosts,.
the God of the armies of Israel,.
whom you have defied..
God uses your past to prepare you for the future..
What works for others may not work for you..
We have to understand that..
And thirdly, the Lord's anointed is a concern.
about upholding the glory of God,.

$^{721}$and honoring God and he or she goes out before the Lord,.
in everything they do..
You're a steward of who God is..
You're a steward of what God is calling you to be..
Your life is no longer yours..
I pick on James because I see him right in front of me..
James who's at Moody Bible Institute,.
that life isn't his..
His calling to work in the marketplace,.
that life and that calling isn't his..
It's a calling that God says,.
that God, what do you want me to do?.
I graduate from a seminary,.
but Lord, what is it that you want me to do?.
His life is coming before the Lord with an open hand.
to say, Lord, my life is yours..
I give it to you..
But that decision has to come early on, doesn't it?.
Again, going back to Indonesia,.
one of the, I get a chance,.
one of this 24 year old individuals,.
I had a chance to sit down and talk with him,.
I had a mentor who sat with him..
He had started and he had invested \$10,000.
and he got an investment of \$10,000 to work on a startup..
Well, that startup grew,.
that he sold his business for \$132 million..
But one of the things that we talked to him about.
before he did this startup,.
and I said, so and so,.
whose life does it belong to?.
Whose life do you belong to?.
Because if that life is the Lord's,.
then you gotta make a decision right now,.
what you're gonna do with that money.
before your company goes really successful.
and you make tons of money..
Because when you make that decision,.
once you become rich,.
you won't make that decision..

$^{761}$You're gonna have to make that decision early on,.
before you become wealthy,.
before you get that job,.
before, whose life do you belong to?.
Because when we're right in that mist,.
it's not gonna be very good..
And we're gonna make decisions that benefit us..
I share this with you because for many of us,.
money is a huge issue, isn't it?.
We worry about money..
We like that life..
My wife and I, early on,.
when we first started Resource Global,.
we had no money..
I still remember the early days of Resource Global.
when we first started 12 years ago..
I sat there, I had maxed out my credit cards,.
not very good for a startup..
I had my last amount of \$200 left to get a hotel,.
to meet with the donor, all of those different things..
But my wife and I said,.
can we live on one salary?.
And can we always use that rest of the salary that we have.
to glorify the Lord in everything that we do?.
To continue to support it for ministry initiatives.
because Lord, that money is not ours..
We're a steward of the money and the time and gifts.
that He has given us..
Fast forward 12 years later,.
I've gotten doing pretty well financially..
This organization has grown to a point where it was zero.
and now it's about \$1,000,001.5 million..
Things are going well..
My wife is doing well at work as a teacher..
Her salary continues to increase..
Does our giving change?.
The answer is no..
We're still committed, doing the same thing..
And I'm not standing here bragging about it..
It's a commitment that we have to make..

$^{801}$It's a commitment that all of us have to make in our lives..
Whose life do we belong to?.
And that our life belongs to the Lord..
And if the Lord calls us that it's about others first.
and never ours, then it makes it easy..
That the Lord's anointed, as David talks about,.
is concerned about upholding the glory of God.
and honoring the God, that I come to you.
in the name of the Lord, the host,.
the God of the armies of Israel,.
whom you have defied, and Lord,.
you are calling me to do this..
So the answer makes it easy..
That the money that I have is not mine,.
but rather it's yours..
And so when Mark and church is so gracious.
to give me an honorarium every single time I speak,.
the answer is simple..
Mark, you use it back for your church.
because this money is not mine..
I do it out of humbleness because of what you have provided.
for me in the cross..
This person needs help, great..
This money isn't mine, Lord..
And the Lord blesses me, probably not in financial ways,.
but in ways that I can't even imagine..
Whose life do you belong to?.
\newpage



\section{詩篇 46:1-11}
\label{sec:eQLpWVai5F0}
\textbf{What's up (詩篇46\_1-11) - 陳德晶傳道 [HeArea心靈手作室同工]}
\newline
\newline
連結: \href{https://youtube.com/watch?v=eQLpWVai5F0}{\texttt{https://youtube.com/watch?v=eQLpWVai5F0}} ~~~~ 語音日期: 2021-11-17
\newline
\newline
\hyperref[sec:GChzvOXTH_M]{\small{< < < PREV SERMON < < <}}
~
\hyperref[sec:index_chronic]{\small{[返順時目]}}
~
\hyperref[sec:index_scriptual]{\small{[返順卷目]}}
~
\hyperref[sec:I4XJBn_dGH4]{\small{> > > NEXT SERMON > > >}}
\newline
\newline
詩篇 46:1-11
\newline
\begin{longtable}{cl}
\hline
\hline
章節 & 經文 (和合本修訂版)\\
\hline
46:1 & \begin{tabularx}{0.7\textwidth}{X} 神是我們的避難所，是我們的力量，是我們在患難中隨時的幫助。 \end{tabularx} \\ \\ \relax
46:2 & \begin{tabularx}{0.7\textwidth}{X} 所以，地雖改變，山雖搖動到海心， \end{tabularx} \\ \\ \relax
46:3 & \begin{tabularx}{0.7\textwidth}{X} 其中的水雖澎湃翻騰，山雖因海漲而戰抖，我們也不害怕。（細拉） \end{tabularx} \\ \\ \relax
46:4 & \begin{tabularx}{0.7\textwidth}{X} 有一道河，這河的分汊使神的城歡喜，這城就是至高者居住的聖所。 \end{tabularx} \\ \\ \relax
46:5 & \begin{tabularx}{0.7\textwidth}{X} 神在其中，城必不動搖；到天一亮，神必幫助這城。 \end{tabularx} \\ \\ \relax
46:6 & \begin{tabularx}{0.7\textwidth}{X} 萬邦喧嚷，國度動搖；神出聲，地就熔化。 \end{tabularx} \\ \\ \relax
46:7 & \begin{tabularx}{0.7\textwidth}{X} 萬軍之耶和華與我們同在，雅各的神是我們的避難所！（細拉） \end{tabularx} \\ \\ \relax
46:8 & \begin{tabularx}{0.7\textwidth}{X} 你們來看耶和華的作為，看他使地怎樣荒涼。 \end{tabularx} \\ \\ \relax
46:9 & \begin{tabularx}{0.7\textwidth}{X} 他止息戰爭，直到地極；他折弓、斷槍，把戰車焚燒在火中。 \end{tabularx} \\ \\ \relax
46:10 & \begin{tabularx}{0.7\textwidth}{X} 你們要休息，要知道我是神！我必在列國中受尊崇，在全地也受尊崇。 \end{tabularx} \\ \\ \relax
46:11 & \begin{tabularx}{0.7\textwidth}{X} 萬軍之耶和華與我們同在，雅各的神是我們的避難所！ \end{tabularx} \\ \\
[1ex]
\hline
\hline
\end{longtable}
$^{1}$大家好, 大家平安!.
很高興能夠跟大家分享訊息.
不過當日我收到WhatsApp邀請時.
我當場大叫了.
WhatsApp? 到底發生什麼事?.
我還以為自己在睡覺.
什麼事呢?.
因為我不太認識大家成會社區教會.
我上網查看才知道我的中學老師吳君才牧師.
在當中服侍.
到底WhatsApp發生什麼事?.
不單在我身上發生.
也很經常在我們的生活中出現.
有時當我們面對一些病痛.
或者我們面對一些人際關係或工作的阻滯.
我們都會發問這個問題.
這兩年更加多.
我們會問WhatsApp.
細細的細菌竟然大大影響我們的生活和工作.
令我們要封城, 要去插墨線.
WhatsApp只發生十天.
竟然令世界變了.
將二十年所打造的和平推倒重來.
WhatsApp一個巨峰竟然可以將整個州搞到停電.
什麼事發生?.
WhatsApp不單在我們生活中.
我想起今天跟大家分享的一段經文.
詩人都經歷了要問到底什麼事發生.
突然來的時期.
詩篇四十六篇.
我們其實都熟悉這段詩.
不過只記得其中一節.
就是第十節在當中.
但我們看詩不要只斷章取義.
只要看完整篇.
才可以明白整個意境.
整個詩人想表達的想法是什麼.
而這篇詩有兩個「細拉」.
就當作一個標點符號的分號.
或者在一首歌中的優子符號.

$^{41}$有一樣東西的時候.
很明顯地將一段詩斬開了.
將詩分成三個段落.
第一段是二至七節.
都好像一首詩或一首歌的正歌.
內容是相似的.
是有一種平衡在當中.
不過意境是不同的.
而第三段就好像是副歌.
八至十一節就看到.
連說話的主動權是哪一個都不同了.
似乎好像在回應第一二節當中所說的.
而我們看回第一段二至三節.
就看到整個環境是怎樣的.
二至三節說.
所以地需改變.
山需搖動到海深.
其中的水需攀均分層.
山需因海漲而戰鬥.
你見到這裡是一幅海嘯地震的圖畫.
本來最不會變,最穩陣,令人傲嬌.
有安全感的地,山.
竟然突然變了.
本來很寧靜的海.
突然在攀均分層.
搞到亂七八糟.
海在古代是象徵著一些邪惡的勢力.
例如魔鬼,敵人,天災,人禍.
剛剛說完天災.
人禍是更加靠近身體並刺入你的心.
在第二段四至七節.
這裡提到人禍.
說到兵臨城下.
是突然之間的.
打到你.
第六節說.
愛邦軒陽,列國動搖.
和之前的天災一模一樣.
都是曹苑巴閉.
都是兵荒馬亂.

$^{81}$很亂的戲當中.
說真的.
人身處在一個亂七八糟的天災人禍當中.
除了不知道將會發生什麼事.
心裡的感受一定是一種恐懼,不安.
很緊張,憂心忡忡.
很失控的狀況.
很想找一些方法去幫自己.
而那些方法甚至可以是非理性的方法.
我們看到新聞的報導.
在阿富汗要撤走的時候.
兵臨城下.
人們不惜一切.
要坐飛機.
但怎麼坐?.
坐在飛機外殼外面.
想著可以逃離這個地方.
但你也知道.
飛機一起.
他們就會跌倒.
很非理性.
但他們都寧願要去坐.
鄧小妹.
說真的.
不用等到我們要打仗的時候.
才會覺得亂七八糟.
每一天我們看著WhatsApp.
Signal.
Facebook.
IG.
還是24Hours的新聞頻道.
你都會看著一切.
你都會有一份恐懼,不安,焦慮在當中.
所以我都很想你這一刻靜一靜.
問一問自己.
WhatsApp.
如果你身處在一個突然而來失控的事情當中.
或者一個訊息當中.
你會發生什麼事?.
不過我們再看一看.

$^{121}$我們看到46篇.
私人和同伴.
整個群體.
在這個天窗和淫窩裡.
心態卻不是像我剛才所形容.
那種比較負面的感受.
它是很平靜.
很有一份安息在當中.
所以我們都要問這個問題.
WhatsApp.
他們發生什麼事?.
為什麼可以這樣?.
這麼特別的呢?.
可以幫我們在現在.
或者我們將會面對的種種事情的處境當中.
我們來看第一段.
這一段是三節.
這裡再重複我所說.
這是一個地震.
這個地震往往都會令人有一份恐懼.
但是私人這裡很特別.
直接告訴你.
我不害怕.
我們不害怕.
我們不害怕有一份安全感.
它的原因.
第二節已經說了.
所以我們看回第一節就知道.
那個原因原來是因為神是我們的避難所.
是我們的力量.
是我們在患難中隨時的幫助.
避難所這個字.
在原文的字根.
其實是一個倚靠和一個信賴.
是說比地更加可靠.
可以安居的保障.
力量在原文裡面.
也可以解一個避難所.
但是是帶著一種能力在當中.
可能是好像山.

$^{161}$尹回泰山在當中.
可以強而有力地去保障.
而幫助.
這裡是說一種主動拯救的動作.
隨時是一個很足夠,徹底,合時的形容.
在波濤洶湧的海裡面.
神就好像可以讓我們抱緊那一個救生圈.
說真的.
神是可靠,堅固,幫助.
誰不知道的.
但是我猜你和我也經歷過.
面對一些困難.
面對一些突然來的事.
我們仍然有份恐懼在當中.
斯人在這裡很特別提及我們一件事.
他說:神是我們.
原文更加說:神是屬於我們.
是belong to us.
說的是甚麼?.
是說人和神的關係.
是我們和神的關係.
不是我們在頭腦裡知道.
而是因為這個關係.
令我們不用害怕.
有安全感.
所以我第一件事很想提及.
為何斯人有一個秘訣.
能夠安靜,有一個安息.
這裡就是我們抱緊屬於我們的神.
抱緊屬於我們的神.
我相信大家都見過或聽過.
不少做女朋友的.
都很喜歡問她的男朋友.
「你是否愛我的?」.
我猜你也知道.
那個男生其實很疼愛那個女生.
連那個女生也很疼愛.
但不知為何總是要問這個問題.
「你是否愛我的?」.
我有時去觀察.

$^{201}$特別在婚前輔導時.
我會觀察到.
可能是關於我們的成長背景.
我們以往的一些經歷.
令我們總是缺乏信心.
沒有安全感.
我們有很多事在頭腦裡知道.
但在關係的連結.
能否緊緊地抓緊.
總是差一點.
所以我也很想提及.
我們不用害怕.
我們不是要抱著其他東西.
也不是要確認甚麼.
我們要抱緊神實一實.
要知道祂是和我們在一起.
那個關係很結連,很緊密.
不過我也想你記住一件事.
如果要抱緊神實一實.
你要知道神未必每個人都要給.
這個要有關係才行.
所以不是在我們的名義上.
或是在我們的頭腦裡知道.
我們一定要建立那個關係.
抱緊到上帝實一實才行.
所以這一刻我也要問你們.
失去你和神的關係.
又會發生甚麼事?.
你是觸碰它?.
你是知一知?.
還是你真的抱緊它實一實?.
第二段,四至七節.
他說這裡的狀況.
是由大自然轉到打仗.
圍成圖畫.
我們留意阿富汗事態的發展.
你就知道守城比進攻更難.
當中很多考慮.
亦令很多人憂慮.
但第四節你竟然會發現一件事.

$^{241}$有兩隻字是很不相襯的.
整個意境.
他說「歡喜」.
他們一群人很歡喜.
這個歡喜令私人和他的群體.
一群人都能夠安心在當中.
而那個原因.
我們看到這裡很清楚.
神不是在哪裡.
這次他不是遠房親戚的關係.
很遠,隔山打牛.
他在這裡親自和我們同在.
第七節他這樣說.
萬君之夜和我與我們同在.
也該神是我們的避難所.
是那一陣子.
他已經在當中保護他們.
而第四至第六節.
這裡更加詳細地描述.
神如何和他們一起.
第四節好像是一個寓意.
是神的恩典和祝福臨在當中.
好像一條河圍繞整個耶路撒冷城.
第五至第六節.
神在其中同在.
讓城堡比山更堅固.
不動搖的在當中.
古時天亮才開始打仗.
因為晚上看不到.
隨時殺錯人.
神在早上.
在要打仗的時候.
神就在那裡.
很及時.
給予他們一份幫助.
在當中請求.
中賢萬邦都像海一樣.
吵得亂七八糟.
在當中喊戰.
但是神一發聲.

$^{281}$你應該聽到的.
砰!.
其他東西都鴉雀無聲.
城堡比他們大.
比他們厲害.
所有東西都很平靜.
所以如果這樣說的話.
你知道.
第二方面.
叫我們在無論什麼處境當中.
不用憂心.
關鍵就是我們要容許神與我們同在.
可能剛才你都很想問一件事.
我們一直都抱緊著神.
我與神的關係很密切.
但是為什麼?.
面對很多事情.
總是會有焦慮.
有憂慮在當中.
其實憂慮,失眠.
甚至這些骨病,焦慮病.
都在現代社會很普遍.
香港四分之一的人.
基本上都有這方面的狀況.
我不太清楚美國是怎樣.
但我所知應該都是差不多的狀況.
我不知道你有沒有試過.
幫一些抑鬱傾向的人.
我們都不知道他是否抑鬱.
但看到他有這個狀況.
我相信你一定會這樣.
你都很想提議去見一見醫生.
或者找一找輔導員.
但相信你都有這樣的經驗.
不是每個人都會馬上說好.
他們都會覺得這是好事.
不過現在不用.
讓我努力一下.
或者他會再說很多原因.
我現在沒有時間.

$^{321}$我不想吃完藥之後很睏.
我不想讓人知道.
有很多原因在當中.
但他們怎樣才能解決問題?.
不是他們自己嘗試.
他們親自找醫生和輔導的時候.
讓他們同行.
那些問題才慢慢解決.
他們肯吃藥.
病情才開始好轉.
我亦見過很多人一好轉.
就馬上停藥不肯見.
反而令病情加深.
我覺得很多時候.
我們要去接納醫生.
其實和接納全能醫治我們的主很相似.
因為我們的性格,經驗,想法,迷思.
我們會不自覺地去抗拒神進入我們的生活.
所以我都很想邀請你.
你去檢視一下你的心靈.
雖然你知道你和神之間的關係.
但是現在神在你那裡.
什麼事發生了?.
你不容許神上和你一起嗎?.
你問一問.
第三個段落.
八至十一節.
我形容這一段是那個副歌.
天災都好像大軍壓境.
人禍兵臨城下.
我剛才都說過.
每一個人都會有很多情緒反應.
很想做一些事.
但是這裡你明顯看到的.
是甚麼?.
神叫我們只做一件事.
那個主動權不是形容我們.
這裡是神親自出聲,去做事.
叫我們只做一件事.
是說你來去看神的作為.

$^{361}$第八節這裡是這樣說.
你們來看.
這個萬軍的耶和華的作為.
神已經親臨在當中.
在這裡負責任.
在這裡做事.
所以你在當中不用再做任何事.
反而要做的就是去看.
看祂怎麼做.
第八節,第九節就說.
祂做了什麼.
祂使地怎麼去荒涼.
祂止息刀兵直到地極.
去折弓,去斷槍,去把戰車焚燒在火中.
是神才能做到的.
祂確實是我們軍隊的元帥.
甚至是親自去打仗的.
不只是指這一點.
祂親自去打仗.
可以搞成這樣.
如果神是我們的元帥.
祂不單只能夠做所有事.
在當中再給我們一個命令.
除了看之外.
給我們一個命令.
在第十節.
大家很熟悉的.
所以我猜你們也很想讀一次.
所以我請你們一起跟著我去讀.
你們要休息.
要知道我是神.
我必在外邦中被尊崇.
在遍地上也被尊崇.
你們要休息.
在新日本所譯的.
你們要住手.
不過我更加喜歡廣東話.
或者加上英文.
就是stop 收手.
但為什麼在這裡看東西的時候.

$^{401}$就叫我們收手.
我們正在看的.
為什麼要收手.
而且在當中.
想想收手這個字.
其實通常是對惡人說的.
停!停!.
不要再做壞事.
為什麼叫我做.
我想起的.
相信不少弟兄姊妹也會有子女.
當你的子女小的時候.
很想幫你煮飯.
這個動機很好的.
是好事來的.
但他一進廚房.
就弄得廚房亂七八糟.
糖和鹽都混淆了.
然後就弄了一個地獄的料理出來給你.
你看到他差點下了.
你會怎麼做呢?.
你一定會跟我說.
你多疼你的子女.
都會跟他說.
停!收手!不要亂來!.
你們坐在一邊.
看東西就行了.
這個世界.
我們經常都會好心做壞事.
不是你不好.
但有時可能我們還不夠霸氣.
如果是這樣的話.
我們不如坐下來.
憑信心去看萬軍的耶和華.
怎樣令到亂七八糟的事情.
弄回原狀.
學收手對我們城市人來說.
是一個相當困難的事情.
我也在學.
剛才介紹我的時候.

$^{441}$我叫一個社區傳道.
跟你們教會都很吻合.
為什麼會這樣呢?.
因為我剛剛離開了服侍多年的教會.
很想創新一下.
想透過教授一些手作.
去接觸或關懷新世代人的心靈需要.
但有一天.
我自己在想.
如果這樣做.
我就少了很多機會趕到.
我其實真的想了想.
信完這一句.
就沒有再想.
但過了幾天.
我收到陳牧師的WhatsApp.
我估你明白.
為什麼我當初真的很驚訝.
真的要問.
What's up? 什麼事發生?.
我什麼都沒有做過.
但我相信自己幫我做.
幫我提醒我.
幫我看去的錯誤.
第二支門我再提.
我們可能在亂七八糟的社會裡.
我們對神的要求和命令.
寫了一句.
我們要坐定.
Let go and let go.
放手.
停止我們正在做的一切.
當我們停止.
神才可以去做所有的事情.
當我們停止.
我們才可以看到神的作為在當中.
不過要停.
或者我們要收手.
真的不容易.
我的經歷.

$^{481}$是要我們有很多的時間和空間.
去察現神的作為.
所以在最後的時間.
我想大家再用多一點時間.
去試一試.
我們有安靜的時間.
停一停.
放下你手頭所有的東西.
問一問自己.
What's up?.
上帝在我們身上又發生了什麼事.
願主賜給幫助我們.
\newpage



\section{路得記}
\label{sec:I4XJBn_dGH4}
\textbf{亂世生存手冊 (路得記) - 朱莉傳道 [基督中心堂佐敦堂 ]}
\newline
\newline
連結: \href{https://youtube.com/watch?v=I4XJBn_dGH4}{\texttt{https://youtube.com/watch?v=I4XJBn\_dGH4}} ~~~~ 語音日期: 2021-11-17
\newline
\newline
\hyperref[sec:eQLpWVai5F0]{\small{< < < PREV SERMON < < <}}
~
\hyperref[sec:index_chronic]{\small{[返順時目]}}
~
\hyperref[sec:index_scriptual]{\small{[返順卷目]}}
~
\hyperref[sec:oiXSurvMOKc]{\small{> > > NEXT SERMON > > >}}
\newline
\newline
$^{1}$也多謝教會的邀請.
讓我有這個機會來和大家分享神的話.
我的名字叫朱莉.
你們一些的頂梓梅應該會記得我的.
因為我曾經在Trinity讀過書.
認識了你們一些的頂梓梅.
所以這次能夠藉著這個機會.
用這樣的方法去連繫,是一件很開心的事情.
我大概在五年前左右就搬回香港.
現在就在香港一邊募會一邊進修.
如果你們有留意香港的消息.
可能都會知道香港正經歷著新的一波移民潮.
而這一波的移民潮.
很多人的評論說會比過往任何一波的移民潮都更加厲害.
有更多的人選擇離開香港.
我相信你們中間很多都是這樣移民過來美國.
或者是移民家庭長大的.
今天我剛剛想和大家分享一個關於移民的故事.
是舊約聖經,路得記.
如果你們已經回到教會一段時間.
相信你們對路得記一點都不會感到陌生.
今天我想和大家一起思考的一個問題是.
為什麼路得記會在聖經裡面?.
為什麼路得記會有什麼價值.
在聖經這麼重要的基督教經典裡面.
有一本這樣的書.
或者另一個問題就是.
路得記到底有什麼神學價值呢?.
首先我們要知道路得記是記載一個.
在大時代裡面一個普通家庭掙扎求全的故事.
的確有峰迴路轉,高低起伏的際遇.
但其實我們每一個人.
我們平常的人所經歷的都差不多.
再加上其實路得記的主人翁.
即是路德.
在聖經裡面不是有很重要的地位.
因為在其他的經典裡面.
是非常非常少再提到路德.
難道這本書記載了神一些很重要的作為.
以致要留在聖經裡面?.

$^{41}$有趣的地方是什麼呢?.
神在路得記似乎是缺席的.
當這個小家庭面對人生很多打擊的時候.
都會期望神親自介入扭轉命運.
但很可惜在路得記裡面.
我們看不到神有這些作為.
我們看不到有任何神跡騎士.
最要命的是什麼呢?.
是神異常地沉默.
連差遣先知向小家庭宣告安慰都沒有.
而在整本路得記.
上帝直接的行動.
上帝Divine Direct Action.
只有兩件事.
即是說在整本路得記.
只有兩處經文的動詞.
是以上帝為主題.
第一次就是一章第六節.
「他」即是拿訛米.
與兩個兒婦起身.
想當初拿訛米是小家庭的媽媽.
因為饑荒.
她與丈夫和兩個兒子.
從猶大地的伯利恆遷移到摩押地.
結果十年過去.
丈夫和兩個兒子都相繼離世.
剩下拿訛米和兩個兒婦.
就是這三個.
即是一屋都是寡婦.
然後第一章第六節就說.
他們起身要從摩押地歸回.
因為她在摩押地聽見耶和華.
眷顧自己的百姓.
賜糧食與他們.
諷刺的是甚麼呢?.
經過一連串的打擊.
即是人生的際遇.
現在聽到鄉下又有東西吃了.
結果拿訛米就打算回流.
這是第一次在路得記裡面記載的.

$^{81}$Divine Action.
主動的介入.
就是神眷顧耶和華.
眷顧自己的百姓賜糧食給他們.
第二次神的介入在哪裡呢?.
就要跳到第四章.
第四章十三節.
「於是波亞斯娶了路德為妻」.
波亞斯是我們的男主角.
路德是誰呢?.
路德當然是我們的女主角.
就是跟著拿訛米回猶大鄉下的摩亞義父.
他跟著婆婆回鄉.
在鄉下遇到婆婆的親戚波亞斯.
結果他們相遇而且相愛.
結果再嫁.
與他們與他同房.
耶和華使他懷孕生了一個兒子.
所以最後是一個大團圓的結局.
神使路德懷孕.
這是第二次在路得記的Divine Action.
第一次和第二次是一頭一尾的.
路得記只有四章.
第一次是第一章.
第二次是最後一章.
但中間我們完全看不到.
神有怎樣介入人類的歷史和命運裡面.
在人的困境裡面.
我們期望神來做更多的事.
但除了一開始神眷顧百姓.
和最後神使沒有生育過的女子懷孕.
當然這些都很重要.
但中間我們看不到神做了什麼.
相反在路得記.
我們看不到神的Divine Action.
但我們看到的是人與人之間彼此的恩代.
不是神恩代人.
是人與人之間彼此的恩代.
我們繼續看第一章第八節.
「那奴婢對兩個兒婦說:.

$^{121}$「你們各人回娘家去吧」.
當時那奴婢奶奶打算回鄉下.
她非但沒有要求這兩個媳婦服侍她中老.
反過來她勸她兩個媳婦回他們的娘家.
甚至祝福他們有重新開始.
所以她說:.
「你們各人回娘家去吧.
願耶和華恩代你們.
將你們恩代已死的人與我一樣」.
首先你要留意「恩代」這個字.
英文譯作Loving Kindness.
其實在聖經裡面是一個非常重要的字.
例如形容神自己的慈愛.
「因為慈愛」就是這個字.
「慈愛」這個字就是「恩代」這個字.
但我想留意的是甚麼呢?.
在這八節的路得記.
那個表達方法很奇怪.
我再讀這一節給大家聽.
「願耶和華恩代你們.
將你們恩代已死的人與我一樣」.
這個以誰為恩代的標準呢?.
是「將你們恩代我一樣」.
即是希望神賜福你們.
好像你們賜福我一樣.
即是用這兩個兒婦.
甚至是摩訶的兒婦.
為那個Role Model.
為那個慈愛,恩代的標準.
如果我們平時這樣溝通.
我們會怎樣呢?.
「求神賜福給你.
好像神賜福給我一樣」.
即是神賜福給你.
好像神賜福給我一樣.
但這裡不是.
是說「願神恩代你們.
好像你們恩代我一樣」.
原來這正正就是路得記的神學信息.
藉著人與人之間彼此的恩代.

$^{161}$成就了神自己的恩賜.
我再說,路得記神直接的參與.
是非常少的.
甚至有人這樣形容.
在路得記神是passively present.
祂在這裡,但祂是passive.
祂在背後策劃.
正如你讀下去.
你會留意到很多字.
「恰巧」,「剛剛」,「正好」等等.
但在路得記裡.
神是passive的.
祂是present,但祂是passive.
但更加active的是甚麼呢?.
是我們看到主角之間彼此的恩代.
不單是口說祝福.
更加是從對方的角度出發.
付代價地去恩代對方.
成就了神的恩賜.
甚至最後推動了救恩的計劃.
所以某程度上可以這樣說.
路得記是一個甚麼呢?.
是一個教科書.
教你如何做一個好人.
如何做一個恩代人的標準.
這些是可以在路得記裡看到的.
但時間關係,我今天只會說.
其中一個彼此恩代的原則.
在路得記裡,我們可以看到.
彼此恩代有其中一個原則是甚麼呢?.
口裡的祝福要化為行動.
我們繼續看一些經文.
因為時間關係,我們想看完四章.
但我們會跳來看.
例如第一章第九節.
「願耶和華使你們各在新夫家中得平安」.
第八節我們讀了.
「那訛米」就勸新婦離開.
祝福他們.
如何祝福他們呢?.

$^{201}$祝福他們各在新夫家中得平安.
即是甚麼呢?.
即是祝福他們再嫁.
這個很難得.
這個婆婆老實說.
她不去考慮自己的前途.
相反她考慮的是新婦們的前途.
竟然祝福這兩個寡婦的儀式有新婦.
留意得平安.
祝福他們在新夫家中找到安寧.
然後我們跳到第二章.
第二章結果.
兩個媳婦一個聽勸離開了.
另外路德,即是我們的主角.
緊緊跟著那訛米.
他不肯走.
他一定要跟著那訛米.
陪著那訛米一起回到鄉下.
去到第二章他們就回到鄉下.
男主角波亞斯在第二章出場.
原來當時路德為了謀生.
剛剛去到波亞斯的田地.
去拾取收割人掉在地上的墨水來謀生.
波亞斯作為一個田主.
即是一個大財主.
他明顯聽過路德的事蹟.
甚至欣賞他.
以至吩咐他的工人恩待路德.
我們跳到二章十到十二節.
二章十到十二節.
路德就苦服在地叩拜.
對他(波亞斯)說.
「我既是愛邦人,怎麼蒙你的恩?.
怎麼憐恤我?」.
因為田裡的工人受了波亞斯的叮囑.
要好好待路德.
路德就來感激波亞斯.
但波亞斯怎麼回答呢?.
十一節.
波亞斯回答說.

$^{241}$「自從你丈夫死後.
凡你向婆婆所行的.
並你離開父母和本地.
到數不認識的民眾.
這些事人全都告訴我了」.
留意,原來路德做的那種友善.
是會有回報價值的.
是因為路德怎樣恩待他的婆婆或婆婆.
現在波亞斯就回報給他.
我們留意第十二節.
是波亞斯的願望.
波亞斯怎樣祝福路德.
十二節.
「願耶和華照你所行的賞賜你.
你來投靠耶和華以色列神的賜旁下.
願你滿得祂的賞賜」.
波亞斯是一個很榮幸的人.
他也同樣祝福路德.
「願神祝福你」.
為什麼呢?.
因為「你來投靠耶和華以色列神的賜旁下」.
留意「賜旁」這個字.
剛才我們說了「Rest」.
即是「得平安」.
但這裡有另一個字我們要留意.
「賜旁」這個字.
在神的賜旁下.
我們跳到第三章.
「當拿訛米知道路德和他的親戚波亞斯.
在田裡相遇.
而且波亞斯厚待路德」.
到第三章一節.
「路德的婆婆拿訛米對他說.
女兒,我不當為你找個安身之處.
使你享福嗎?」.
原來這裡拿訛米就開始想了.
我要幫你找個安身之處.
一個「Rest」.
我要「Seek rest for you」.
使你享福.

$^{281}$這個「Rest」是甚麼呢?.
正正是第一章第九節.
我們剛剛讀的.
「願耶和華使你們各在辛苦家中得平安」.
起初這是拿訛米的願望.
是希望神為你們預備平安.
結果到三章一節.
「我豈不應該為你找平安或安身之處?」.
口中的祝福要化作行動.
接下來我們看到拿訛米.
在第三章後面開始為路德籌劃.
叫路德有一連串的預備.
去到收割的田裡.
去到三章第七到第九節.
收割完之後第七節.
「幫阿司吃喝完了.
心裡歡暢就去睡在麥堆的旁邊.
路德便悄悄地來掀開他腳上的被.
躺我在那裡.
到了夜半那人忽然驚醒.
翻過身來.
不料有女子躺在他的腳下.
他就說:你是誰?.
回答說:我是你的被裡.
路德,求你用你的衣襟遮蓋我.
因為你是我一個至近的親屬」.
我們可能有很多懷疑.
在這裡發生什麼事呢?.
好像很曖昧.
為什麼路德會這樣說呢?.
簡單來說,這是另類的求婚.
其實就是路德向邦亞斯求婚.
是求邦亞斯用他的衣襟.
代表保護來遮蓋路德.
其實是求你求婚.
然後求你保護我.
這樣的意思.
但這裡特別的地方是什麼呢?.
求你用你的衣襟這個字.
其實正正就是我們剛才讀的.

$^{321}$二章十二節的邦亞斯口中的祝福.
你來投靠以色列神的翅膀.
原來翅膀這個字.
其實在這裡就是衣襟.
即是Wings這個字.
即是翅膀,Wings這個字.
所以邦亞斯的Wish.
邦亞斯的祝福.
你來投靠以神的翅膀.
願神厚待你.
但路德說,你來兌現你的Wish.
不單是神的翅膀.
是你的衣襟來遮蓋路德.
所以原來Kindness,恩慈.
在路得記裡不只是停留在口中.
其實和新約很多教導是一致的.
是要在口中的祝福化作行動.
替對方著想,甚至付代價的.
這樣的一種恩慈,彼此來相待.
我講到這裡,可能大家就會想.
這個路得記的神學價值.
最後是否說Be a good man.
做個好人就這麼簡單呢?.
是否說,我們總之彼此恩待.
最後就會大團圓結局呢?.
我去到第四章.
結果的確是大團圓結局.
我們看第四章十一到二十二節.
四章十一到最後的二十二節.
「在城門坐著的眾民和長老都說.
我們作見證」.
其實這個簡單去畫圖.
是一個婚禮.
就是幫助使徒路德最後結成婚盟.
所以有在城門坐著的眾民和長老.
就作見證了.
然後他們怎樣呢?.
他們都有一個Blessing.
他們都有一個Well-wishes.
「願耶和華使進你家的這女子.

$^{361}$像建立以色列家的拉潔尼亞二人一樣.
又願你再以法他德亨通.
在伯利恆德名聲.
願耶和華從這少年女子賜你後裔.
使你的家長.
哈瑪從猶大所生法勒斯的家一半」.
所以這是一個婚禮Blessing.
在婚禮裡祝福這對新人.
當然焦點是甚麼呢?.
就是後裔.
十二節「願耶和華使你得後裔」.
到十三節.
真的成就了.
這次是神親自來成就Blessing.
「於是幫助使徒取了路德為妻.
與他同房.
耶和華使他懷孕生了一個兒子」.
所以神在路得記做了兩件事.
第一章第六節和第四章十三節.
「神使路德懷孕」.
跳到十六節.
「我米就把孩子抱在懷中.
作他的養母.
論者的婦人說:.
我米得孩子了.
就給孩子起名叫俄彼得.
正俄彼得是耶西的父.
耶西是大衛的父」.
如果最終路得記的神學思想是甚麼呢?.
「好人有好報」.
就是這本書的Take-home lesson.
就是這本書的中心思想.
其實停在十七節就已經完美了.
大團圓結局.
「那我米都有兒孫奉養他中老」.
而且還說了甚麼呢?.
說原來這個俄彼得成為大衛的祖先.
超級完美.
最後提醒我們.
記住了好人的句號.

$^{401}$是不是這樣呢?.
當然你會發現.
原來路得記不是停在十七節.
後面還有十八到二十二節.
十八到二十二節是甚麼呢?.
是一個家譜.
「法律師的後代記載下.
法律師生希斯倫.
希斯倫生蘭.
蘭生阿米拿達.
阿米拿達生拿遜.
拿遜生撒滿.
撒滿生波亞斯.
波亞斯生俄彼得.
俄彼得生耶西.
耶西生大衛」.
很多人會覺得這個家譜.
似乎是一個appendix.
是一個很多餘的東西.
只不過可能是正式的.
去介紹大衛的譜系.
其實不需要的.
因為十七節已經介紹了.
俄彼得和大衛的關係.
即是代表甚麼呢?.
路德和大衛的關係.
但為何最後要加這個家譜呢?.
原來最後十八到二十二節的家譜.
才是路得記裡面最重要的一部分.
他將路得記.
他把所有事物放在角度.
他將路得記.
告訴我們路得記要怎樣去解讀呢?.
原來要留意這個家譜.
家譜是怎樣開始的?.
十八節.
四章十八節.
「法勒斯的後代記載下面」.
「法勒斯生希斯律」.
為何這個家譜是從法勒斯開始的呢?.

$^{441}$如果我問.
就這樣問大家.
到底你們有多少人.
可以馬上說出法勒斯是誰呢?.
我懷疑不是很多.
因為法勒斯其實不是一個很出名.
在聖經裡面很重要的人物.
但他父親其實就是一個很重要的人物.
他父親是誰呢?.
猶大.
如果我要很荒謬地去介紹大衛的家譜.
想將一些重要的人物連繫在一起.
家譜的作用某程度上就是這樣.
我會不會從法勒斯開始?.
我會從猶大開始.
去強調大衛和猶大的連繫.
但不是的.
這裡提到法勒斯是怎樣解釋呢?.
還要怎樣呢?.
法勒斯不是在家譜裡面第一次出現.
法勒斯在路得記的前文.
在第四章已經出現了.
我們剛才也讀了.
不知道大家有沒有留意到.
在第四章十二節.
「冤爺和華從這少年女子賜你後裔.
使你的家將他媽從猶大所生」.
法勒斯的家一半.
所以前面已經提到法勒斯.
而且是婚禮的祝福.
所以法勒斯也說了是猶大所生.
但是有一個額外的資訊.
一個不需要的額外資訊是什麼呢?.
使你的家將他媽從猶大所生.
無端端提到他媽.
當然的確法勒斯是猶大和他媽所生.
但他媽如果大家有印象.
記得勝經在《創世記》三十八章.
今天不會有時間去看.
他媽這段故事不是很光彩.

$^{481}$為什麼呢?.
因為他媽其實是猶大的媳婦.
某個程度上他們的結合.
媳婦和公公的結合.
生下法勒斯.
這樣的結合是有一些爭議的.
為什麼這群長老竟然在婚禮裡.
在城門口這樣祝福.
幫助斯和勒斯呢?.
你試試在婚禮上寫一個婚禮卡.
恭喜你的朋友.
希望你好像他媽一樣.
我懷疑你朋友也做不到.
所以很奇怪.
無端端提到他媽.
換言之他講了什麼呢?.
其實第十一節也有提到兩個女性.
就是拉潔和莉雅.
這個更容易理解.
因為拉潔和莉雅的確是十二支派的國母.
即以色列的國母.
十二支派是從他們兩個人而來的.
所以提到兩個女性我們明白.
但為什麼要提他媽呢?.
而且特別提到法勒斯是他媽所生.
然後家譜就從法勒斯開始.
其實就證明了張大衛.
證明了原來是他媽和路德的子孫.
將他媽和路德放在一起.
而聖經裡面.
只有兩處的聖經.
整本聖經不只是舊約.
整本聖經只有兩處的經文.
是同時提到他媽和路德.
另一處又是一個家譜.
不知道大家有沒有印象.
《馬太福音》第一章.
耶穌基督的家譜.
又再提到他媽.
又提到路德.

$^{521}$又提到大衛.
更重要的是耶穌基督.
就是這些人的後代.
所以原來這樣看的話.
我們就明白路得記不只是指向大衛.
其實路得記更加要.
那條線要一直畫.
一直畫到新約.
是指向耶穌基督的.
路得記發生在士司的時代.
按照基督教的聖經編排.
路得記在士司紀和撒母耳紀上下的中間.
好像一個橋樑連接這兩段歷史.
當時路得記發生的背景.
是士司的年代.
如何是當時任意而行的國人.
沒有王管的一個很混亂的時代.
然後路得記記載了這些小人物.
路德,拿俄米,波瓦斯.
如何彼此恩代的一個事情.
原來就在不知不覺裡面.
他們就推動了歷史的軌跡.
當他們堅持為善彼此.
用實際行動彼此恩代的時候.
他們服侍了他們那一代的人.
亦為下一代人預備.
最後如何呢?.
在撒母耳紀我們看到大衛的出現.
我們以為大衛的出現就是答案.
既然士司記說沒有王管.
大衛的出現.
而且撒母耳記提到大衛是合神心意的王.
我們會以為終於解決方案出現了.
在大衛那裡.
但如果大家記得勝經.
撒母耳紀裡面是充滿了大衛的缺失.
甚至最後到撒母耳紀下二十四章.
是以大衛的缺失來終結.
即是說原來一個沒有王管的.
連一個大時代裡面.

$^{561}$一個混亂的時代.
我們以為一個屬靈的領袖.
一個合神心意的王出現.
就可以解決問題.
最後都是失望終結.
路德紀的答案.
或者這個時代的答案.
不是在於一個王的出現.
一個所謂合神心意.
但其實在他人生裡面仍然有很多缺失的王的出現.
更加重要的是甚麼呢?.
路德紀是指向新約的耶穌基督.
原來要解決那個時代.
混亂,各人任意而行.
沒有王管的年代.
不是任何一個我們覺得可以解決的領袖出現.
就可以解決這麼簡單.
是要等到耶穌基督.
去到新約.
耶穌基督的出現.
他的降生,受死,復活.
成就了神最大的恩慈.
就是袋屬的救恩.
成為了我們最近的親屬.
原來路德紀不是只連接兩段歷史這麼簡單.
他告訴我們.
有些事是可以盼望的.
不是盼望一個王的出現.
是盼望那個萬王之王.
耶穌基督的出現.
才可以解決人類的問題.
不知道大家有多久沒有回香港.
我相信起碼兩年.
因為疫情的期間.
自從我搬回香港之後.
我就住在長洲.
不知道你們有沒有多少人.
記得長洲或者去過長洲呢?.
長洲有兩樣東西很出名的.
兩樣東西很多的.

$^{601}$第一是單車很多.
第二是野狗或者狗很多.
而這兩樣東西.
我都害怕的.
狗當然.
特別是野狗.
牠會無視.
牠會突然在你旁邊出現.
我本身害怕狗.
就會很驚恐.
另一樣東西多的是單車.
因為長洲沒有車.
所以很多人是以單車代步.
在橫街閘巷裡.
單車通過的時候.
牠就會提醒你.
所以牠就會響鈴.
提醒你.
我在後面.
我要通過你.
但每一次聽到這些聲音的時候.
我就會害怕.
我就想避開牠.
我在想應該怎樣避開牠.
讓牠不要撞我呢?.
後來我發現.
秘訣是什麼呢?.
無論是對付單車.
特別是突然出現的響聲.
或者是突然出現的野狗.
其實秘訣都只有一樣.
只有一樣東西.
可以讓我避開這些煩惱.
不要理牠.
繼續走.
你試一下.
如果狗突然在你旁邊出現.
然後你反應得很快.
你彈開.
甚至想怎樣養的話.

$^{641}$想拿把傘出來等等.
保護自己等等.
你猜狗會怎樣呢?.
同樣道理.
你猜如果有單車在後面.
你又是想避開牠.
你就會動來動去.
移來移去.
這樣會更危險.
所以原來.
面對很多時代.
會擾亂我們的事情.
很多驚恐.
很多不安.
很多不確定.
我們需要做什麼呢?.
繼續走我們前面的路.
是.
身處在一個大時代裡面.
無論在香港.
無論在北美.
我們都不知道疫情沒完沒了.
我們看到整個世界的局勢.
無論是政治,經濟.
都有很多動盪.
都足以令到我們不安.
但我們都只是小人物.
我們都平凡如路德.
如波亞斯.
如納俄米.
但我們作為跟隨主的人.
透過路得記.
我們可以提醒這一點.
堅持為善.
堅持彼此恩待.
服侍這一代.
守護下一代.
有些事我們做了.
的確我們看不到後果.
就好像家抱.

$^{681}$路德可能都不知道.
大衛會成為他的兒孫.
或者大衛作王的時候.
路德可能都已經不在了.
更何況耶穌基督的出現.
所以我們不只是期望好人有好報.
因為的確.
好人我們未必看到好報.
但神知道.
我們堅信神仍然在工作.
堅持的成果交託給神.
盼望神的計劃.
終必成就.
在這個世上.
我們一起祈禱.
讓我們有這個機會.
一起來打開你的話語.
來看路得記這本有趣的書.
可能的確我們會覺得.
這本書只是一個有趣的故事來看.
但神透過這本書再一次提醒我們.
活在大時代裡.
主你交託我們的責任.
就是堅持走在你的路上.
彼此來到恩代.
默默的來到守候.
因為我們相信你仍然在工作.
你對這個世界的計劃仍然未完成.
主你求你幫助我們.
讓我們忠心的成為你計劃的一部份.
以致有一天我們能夠見到.
你自己公義的彰顯.
你計劃的成就.
祝福我們無論身處何方.
讓我們見到你自己的恩惠.
見到彼此之間恩慈的結果.
禱告奉耶穌基督的聖名.
阿們.
\newpage



\section{}
\label{sec:oiXSurvMOKc}
\textbf{舊約網上行 - 上帝的足跡: EP19 以掃與雅各}
\newline
\newline
連結: \href{https://youtube.com/watch?v=oiXSurvMOKc}{\texttt{https://youtube.com/watch?v=oiXSurvMOKc}} ~~~~ 語音日期: 2021-11-18
\newline
\newline
\hyperref[sec:I4XJBn_dGH4]{\small{< < < PREV SERMON < < <}}
~
\hyperref[sec:index_chronic]{\small{[返順時目]}}
~
\hyperref[sec:index_scriptual]{\small{[返順卷目]}}
~
\hyperref[sec:cKHN4UF2tsA]{\small{> > > NEXT SERMON > > >}}
\newline
\newline
$^{1}$開到《九日妄想》第十九集,今次來到創世記另一個新的段落.
二十五章第十九節開始.
以下是阿布拉罕的兒子以撒的後代.
阿布拉罕生以撒,以撒娶利伯加為妻時正四十歲.
利伯加是巴旦亞蘭地亞蘭人比土里的女兒.
是亞蘭人拉班的妹妹.
以撒因為自己的妻子不生育,就為她懇求耶和華.
耶和華應允她,她的妻子利伯加就懷了人.
不能夠生育,然後神即地使他能夠生育.
能夠生出下一代,能夠繼承下去.
這是創世記特有的主題.
主要是講到人的無能,不能夠生育,不能夠延續下去的時候.
神就藉著神職,藉著他的工作,能夠將神的子民傳承下去.
這裡講到以撒的禱告,神應允了他.
表示這個家族能夠延續下去.
能夠將來成為以色列國,能夠以色列民.
之所以能夠存在,神的子民能夠繼續發展下去.
是全部因為上帝的恩典,不是靠人意.
這個故事開始的時候,講到上帝賜給以撒和利伯加後代.
然後就交代這個後代的出生和成長的經過.
原來他們是有雙胞胎,二十五章第二十二節.
雙胞胎在利伯加的福中彼此碰撞.
他就說,若是這樣我為什麼活著呢?.
好像是一個很苦很辛苦的壞人.
他就去求問耶和華,耶和華就回答他.
他說兩國在你的土裡,兩族從你的福中要分出來.
將來這族必強過那族,大的要服侍小的.
二十五章第二十四節.
到生產的時候,在利伯加的肚腹裡面果然是生出雙生子.
先出來的全身赤紅有毛,像毛衣一樣.
父母就為他起名叫做爾素.
在這段經文裡面有很多諧音字出現,很相近的字.
譬如紅色這個字,就是和爾素將來有瓜葛的.
因為將來他的destiny就是會住在以東地.
以東就是有紅色的意思.
大家可能去過現在約旦王國的下面.
就是以往的以東地.
在以東地有一個最出名的古城就是佩特拉城.
那裡最著名的當然就是公元後一世紀.
這個勒伯提人的民族,君王的墓地.

$^{41}$佩特拉城有另一個稱號叫做玫瑰城.
因為在以東地全部找到的岩石都是紅岩石.
紅色是一個諧音字,講到以東爾素將來居住的地方.
毛這個字,希伯來的音,也是像西爾.
西爾是一個地方,也是將來的以東地.
所以講到紅色,講到毛的時候.
不單止講到爾素出生的情況.
也好像一個預言一樣.
講到人和民族將來會有什麼結果.
然後二十五章二十六節.
隨後他的弟弟也出來.
他的手抓住爾素的腳跟.
因此父母為他起名叫做雅各.
有一個諧音字,像腳跟這個字.
利伯加生了兩個兒子的時候.
以撒年正六十歲.
二十五章二十七節.
講到兩個孩子的長大.
他們的經過,他們的成長.
兩個孩子長大的時候.
爾素善於打獵,喜歡生活在田野.
雅各為人比較安靜,常常在帳棚裡面.
以撒愛爾素.
因為很喜歡吃爾素所打來的野味.
利伯加媽媽偏愛雅各.
講到爾素和雅各.
一個在田野做獵人,喜歡打獵.
然後雅各就常常躲在帳棚裡比較安靜.
其實也沒有什麼正負面去做一些評語.
只不過是交代一個哥哥喜歡在外面.
然後雅各喜歡在帳棚.
當然二十八節就來了一個隱藏的原因.
將來有的一些張力.
家庭裡面的張力,來自父母偏愛,視單一個的兒子.
無論如何,這裡交代.
就是以撒之後的兩個後代.
創世記二十五章,在這個新的段落開始.
一直環繞著一個家族的問題.
這個主題就是長子的名分.
提到長子的名分,現代人可能比較陌生.

$^{81}$因為我們比較西方化,民主化了.
大體上我們都不會按照以前的傳統社會的習俗.
但是在近東的社會.
甚至今天在近東一些比較傳統的村莊的民族裡面.
都是非常盛行的一個習俗.
長子,男丁.
特別是一些大家族,在雨傘區內.
是一些無與倫比的特權.
幾乎是一種間選.
如果我們想像這些背景的時候.
可能我們可以用一些現代人.
我們比較習慣聽過的一些概念去了解.
我在香港長大的.
香港有一些所謂丁屋.
一些爭丁權的問題.
以往曾經有一套劇集.
Sitting Forward.
就是講到這些丁權的爭奪.
所以如果你想到丁權丁屋的時候.
或者有些相似的就是近東社會那種長子的文化.
因為丁屋的地權.
財產的繼承.
是按照那個長子.
按照那個男丁去傳承下去的.
所以是一些很寶貴的東西.
是一種特權.
在近東的文獻.
其實也有時候會提到這些長子的名份.
不單單聖經有提過.
譬如這裡舉一個例子.
在創世記的族長時代.
大概公元前2000至1500年間.
在一個地方叫做馬里.
Mauri.
在現在的伊拉克的北面.
聖經沒有提過.
但是有很多出土的文獻.
是很像族長時代的那種習俗.
其中一份文獻提到這些所謂丁權.
講到這些習俗.

$^{121}$他說如果希臘倫.
希臘倫.
那個丈夫的名字.
還有亞利林的妻子.
阿利圖.
夫妻之間收養了很多的兒子.
因為當時的大家族.
他們繼承他們的財產的時候.
他們都是靠著南丁的.
所以很多時候會接收.
會收養很多其他人做他們的合法兒子.
那些可能是一些親友.
可能是一些非親屬.
但是這些種子當中.
只有那個長子.
叫做艾哈提尼.
艾哈提尼這個長子.
才是主要的承繼人.
可以從他父親的遺產當中.
得到雙份.
Double portion.
他的弟弟就將剩餘的財產均分了.
就是兄弟如兄弟的意思.
所以這裡讓我們看到.
長子的特權.
長子的名份.
在近東的社會裡面.
是一個很常見的習俗.
一種文化.
說到長子的名份.
可以分開有兩方面.
作為一個長子.
他的名份帶來的.
當然第一方面就是財產.
就是那種祝福.
物質上的祝福.
長子通常會得到家富.
雙倍給他的家產.
所以一個家庭.
如果他有十個小孩.

$^{161}$這樣才算.
因為有長子.
長子是要得到雙倍的遺產.
所以分家產的時候.
家富就不會將家產除十.
會將他除十一.
雙倍雙份的.
Double portion就會給長子的第一半.
長子就會承繼家富更多的財產.
所以你想起Double portion的時候.
你會想起繼承權.
譬如你記得之後舊約的故事.
以利沙和以利亞.
一個徒弟一個師傅.
以利沙有什麼祈禱有什麼請求.
就是林師傅未離開之前.
他說什麼.
求你使你的零.
雙倍的Double portion的林道岸.
所以這裡以利沙的禱告.
在做什麼呢.
他希望能夠成為Firstborn.
能夠繼承以利亞先知的職份.
是一個很榮耀的職份.
第二方面長子的名份.
我們就會想起領導權.
所以不單單是物質上的福氣.
也是一種Spiritual的Status.
一種Social Status.
一種社會最高的地位.
長子就是繼承家庭最高的領導.
在他的轄下.
他的帶領之下.
母親弟妹工人.
都會在他的領導之下.
所以在社會當中.
在物質裡面.
長子是一個很尊貴神聖的地位.
有了這些背景資料的時候.
我們就看看以下的經文.

$^{201}$29節就開始講一個以素能像的故事.
他說有一次.
雅各正在煮豆湯的時候.
以素從田地打獵回來.
他非常之累.
以素就跟雅各說.
求你把紅豆湯給我喝.
因為我已經疲乏得很.
因此以素的名字又叫做以東.
因為紅豆湯紅色.
雅各就跟他說.
你先要把你長子的名份賣給我.
以素就說.
我快要死了.
我的長子名份對我有什麼益處呢?.
雅各就說.
你先向我起誓.
以素就立刻向他起誓.
然後把自己的長子名份賣給了雅各.
於是雅各就把餅和紅豆湯給了以素.
以素吃了喝了.
起來走了.
以素就輕看了他長子的名份.
首先我們覺得很奇怪的就是.
為什麼以素很累回到家裡.
他很餓.
為什麼不吩咐一下.
大家族嘛.
有很多工人煮東西給他吃.
為什麼好像受到雅各一個人去威脅呢?.
可能實際情況未必我們表面上這麼想的.
可能這裡的情況有一點像.
創世記之後的故事.
你記得約瑟的故事.
哥哥也去了幾天的路程遠的地方放養.
四周四景都沒有人.
雅各就派約瑟去探查一下.
到底他的哥哥是不是平安.
走了幾天的路程.
約瑟終於找到他的哥哥.

$^{241}$在那裡方圓十里都沒有人.
然後就發生了一些家庭的慘劇.
所以這裡說到以素和雅各.
可能也是同樣有這樣的情況.
就是可能以素去了很遠的地方打獵.
然後雅各隨後可能帶著其他的工人.
去找以素.
然後在那裡搭了一個帳篷.
以素打完獵很久之後.
他就回來了.
很累.
然後又用紅豆湯.
雅各就去威脅著他.
還在帳篷裡面威脅著他.
因為以素就在野地作主.
雅各就躲在帳篷裡.
所以當以素跟隨力進回到帳篷裡面的時候.
在方圓幾里都沒有人的時候.
雅各就在帳篷裡面佔上風.
煮紅豆湯.
就是用一些野豆或者野菜蔬去煮成的湯.
是當時最基本最簡陋的一些食物.
雅各為人比較內斂.
比較思想型一點.
心思熟慮一點.
他可能從小到大.
大家兄弟都知道大家的個性.
他就抓到以素哥哥的致命傷.
在以素最軟弱的時候.
他就欺騙了.
拿了.
搶了這個長子的名分.
經文只是告訴我們.
叫他起誓.
以素就立刻起誓.
但是沒有詳細的來交代這個交易.
是怎麼進行的.
可能有其他的僕人在.
他就起誓了.
或者有一些條文的簽署等等.

$^{281}$經文就假設了是一個正式的交易.
只是告訴我們交易的事實.
長子的特權.
我們講過有兩方面.
一方面是財產.
另一方面就是社會的領導權.
家族的領導權.
這裡似乎告訴我們雅各是第一次騙取長子的名分.
似乎牽涉財產.
而沒有的就是家族的繼承權.
那種領導權.
所以你記得故事繼續發展下去的時候.
創世記27章的時候.
才講到雅各第二次騙取以薩加夫的祝福.
所以你想像一下.
這種背景這種情況.
近東的人.
當時古代的人.
熟悉所謂丁權.
熟悉所謂長子名分.
聽到這個故事的時候會有什麼反應.
你會不會為了一些小小的湯水.
你會出賣所謂丁權.
如果我們都是丁屋的家族.
從小到大我們知道財產領導上有一種特權.
那種很有價值的地位.
很尊貴的財產.
我們會不會為了出錢一丁.
然後放棄丁權.
所以如果你這樣看.
這個故事對當時的人來說.
是不可思議的.
怎麼會將一些與生俱來.
有的價值有的地位.
為了一碗紅豆湯這麼基本的東西而出賣.
聽到這個故事.
相信古人一定會有這樣的反應.
長子的名分怎麼可能會當為等閒事.
第34節.
這裡怎麼形容耳掃.

$^{321}$他說於是雅各將餅和紅豆湯給了耳掃.
耳掃做了什麼.
吃了喝了就起來走了.
耳掃就這樣興旱了他長子的名分.
這裡用了多少個動詞.
五個連續短的動詞.
吃喝起來走了興旱.
或者應該說藐視.
藐視了長子的名分.
為了解決和舒緩前面短暫的痛苦.
肚餓口渴.
就將一個最有價值的東西.
就興旱了藐視了.
藐視不是不知道價值.
不小心放過了.
這種不是藐視.
這種不是興旱.
興旱藐視的意思是明知道某一個地位.
那種價值就將它當成等閒事.
這裡我們只看耳掃的刻畫.
你看償文夏利.
其實都是強調耳掃是否配授長子的名分.
你看文理的脈絡.
二十五章這裡說一個故事.
三十四節結束的時候.
就說耳掃興旱長子的名分.
然後二十六章一到三十三節.
整個故事說到耳剎.
是傳言繼承亞伯拉罕的承諾.
繼承神的財產.
然後二十六章三十四至三十五節.
又再提到耳掃.
說他娶了外邦的女子為妻.
並且帶來這些客人的女子.
常常會使家父和家母傷心難過.
夾在中間頭和尾都是刻畫耳掃.
到底是否真的配授長子的名分.
我們知道《創世記》舊約.
直接寫給當時古代以色列人.
屬神的子民.

$^{361}$《創世記》要寫給這些人的時候.
就揀選不同的情節.
這裡揀選了俗長一些故事.
說到耳掃興旱長子名分的故事.
到底想直接跟以色列人說些什麼.
或者延伸下來.
由舊約到新約.
我們作為教會會有什麼教導.
對於以色列人來說.
直接寫給他們一個長子的故事.
講到長子名分的問題的時候.
故事當然是想以色列人了解自己的身份.
原來其他的經文告訴我們.
以色列民 以色列國是神的長子.
我們看《出埃及記》四章二十二節.
這裡說到你要對法老說.
摩西地神差派要跟法老對抗.
對法老說什麼呢.
耶和華者人說.
以色列是我的長子.
原來以色列人整個民族的子民.
就是上帝的第一端.
是上帝的恩典.
白白地將他們從遺老之地.
埃及救出來作為他的長子.
他的第一端.
在上帝的財產當中.
在上帝的國度裡面.
他們是多麼寶貴.
《出埃及記》十九章.
繼續有一種主題.
十九章他們快要來到西大山的時候.
摩西上到神的那裡.
耶和華在山上呼召他.
你要跟亞國家和以色列人說什麼呢.
第四節.
我向埃及人所行的.
你們都看見.
看見我怎樣將鷹一樣.
背在我翅膀上.

$^{401}$帶領你來到這個地方西大山.
現在你們若是實在聽我的話.
遵守我的約.
你們就會在萬民當中.
作為我屬我的產業.
這是一個特別的字.
Treasure.
或者有些英文翻譯會叫做.
Crown Jewel.
就是那種名價之寶.
就是掌上名珠的意思.
以色列人在神的眼中.
是一種很寶貴.
一種長子的名分.
說到Crown Jewel.
Treasure.
想想你自己.
一些很寶貴的東西.
藏在櫃底.
或者在銀行的保險箱.
總之一些金銀珠寶.
或者對你來說很有價值的東西.
上帝就是這樣看著以色列人.
看著他們自己的神的子民.
因為全地都是我的.
你們要歸我作為君尊的子子.
聖潔的國民.
在上帝眼中.
以色列人就是他們的長子.
所以寫一個這樣的故事.
告訴他們.
不要輕看.
不要去藐視這個尊貴的身份.
以色列人本身就是神的長子.
神的掌上名珠.
以素輕看長子名份這個故事.
就是一個很大的屬靈的警惕.
就是告訴他們.
不要以為你與生俱來被上帝揀選.
你就不會失去這個長子的名份.

$^{441}$這個屬靈的產業.
你要守約.
如果你是背約的話.
你就會失去更多.
失去更大.
繼續延伸下去到新約.
說到神的子民.
教會.
基督徒.
你和我.
也用了舊約這個.
以素輕看長子名份的故事.
其中一處就是希伯來書.
這裡用了以素的故事.
來作一個警告.
第十二章.
你們要小心.
免得有人失去了神的恩典.
免得有苦根長起來.
纏擾你們因而污染了許多人.
又免得有人成為淫亂和貪戀世俗.
好像以素一樣.
為了一丁點的食物.
竟把自己長子的名份出賣了.
你們知道.
後來以素想要承受祝福卻被拒絕.
雖然帶著眼淚去尋求.
但也沒有反悔的餘地.
以素作為長子.
講條件就必定先會得到神的記業.
但因為眼前一丁點的福祿.
想舒緩一下.
就放棄了永恆的產業.
你說這是多大的損失.
希伯來書本身就是寫給猶太基督徒.
在第一世紀的時候.
猶太人信了主.
他們得到神的救恩.
但卻因為在當時第一世紀的社會裡面.
受著很多以往猶太教給予的壓力.

$^{481}$挑戰和引誘.
所以有很多猶太基督徒.
因而為著眼前想舒緩一下.
好受一點.
就回到猶太教.
就放棄了基督.
這是一個很大的損失.
So great is salvation.
得到基督.
認識到基督.
卻因為眼前一丁點的好處.
而放棄永恆.
這是一個很大的錯失.
你想想.
在世上有哪一個民族.
是更有條件去認識彌賽亞.
認識基督.
豈不是猶太人.
他們有律法書.
有先知書.
有祭禮.
有聖殿等等.
這些全部都是指向基督.
指向彌賽亞的將來.
為什麼會有這麼多條件.
能夠成為神的長子.
能夠得到這個屬靈永恆的救恩.
就會為了世上一點點的舒服.
而失去永恆.
忽略了這麼大的一個恩典.
希伯來書.
一個很重要的警惕.
其實也是寫給我們.
我們有上帝的恩典.
能夠去認識神自己.
認識耶穌基督的救恩.
傳言是一個恩典.
長子的名分.
恩典臨到我們.
求主給我們這個警告.

$^{521}$讓我們知道在難處裡面.
我們不要輕易放手.
把握著上帝給我們的恩典.
去珍惜它.
求上主給我們重整我們的價值觀.
不要因為前年的警況.
不同的困境.
我們就捉不住神的恩典.
求主幫助我們.
繼續去維持我們成為神的子民.
在恩典當中.
我們繼續往前.
在恩典裡面我們靠主.
去信靠祂.
讓我們不要失去這麼大的救恩.
永美的救恩.
在耶穌基督十字架裡面為我們預備.
我們同心低流祈禱.
因為你的恩典.
是我們不配受的恩典.
主,我們求你幫助我們.
在我們遇見困難,難處的時候.
遇見隱憂的時候.
我們不會輕易放過.
這個長子的名分.
這個屬靈的,這麼寶貴的救恩.
一個永恆的救恩.
求主你親自來堅定我們的價值.
堅定我們的觀念.
以致我們真的能夠抓緊你的恩典.
繼續在你的恩典當中.
來跟隨你,事奉你.
祈禱奉主並救.
\newpage



\section{10:19-25}
\label{sec:cKHN4UF2tsA}
\textbf{The Church is the gathered (Hebrews 10\_19-25) - Rev. John Yi}
\newline
\newline
連結: \href{https://youtube.com/watch?v=cKHN4UF2tsA}{\texttt{https://youtube.com/watch?v=cKHN4UF2tsA}} ~~~~ 語音日期: 2021-11-23
\newline
\newline
\hyperref[sec:oiXSurvMOKc]{\small{< < < PREV SERMON < < <}}
~
\hyperref[sec:index_chronic]{\small{[返順時目]}}
~
\hyperref[sec:index_scriptual]{\small{[返順卷目]}}
~
\hyperref[sec:BxAa7Ok6gjk]{\small{> > > NEXT SERMON > > >}}
\newline
\newline
$^{1}$our time to John..
Let's welcome John..
(audience applauding).
- Everybody again at Urban Voice..
And in fact, this is, you know,.
since I've been coming out monthly since August,.
this is the most of Urban Voice I have seen since August..
So thank you for choosing to come on out today.
and join with your brothers and sisters in worship.
and in the hearing of God's word..
With that, let's begin with just a quick word of prayer.
as we go into God's word, okay?.
Heavenly Father, thank you so much..
We do thank you because not just 'cause it's Thanksgiving,.
but because this is an occasion to remind us.
to pause for a moment..
Perhaps it's in the midst of a season of grumbling,.
of discontent..
In a season where we're longing for more or better,.
help us now to stop and just recall.
all the great things you have done..
And then to think about the great things.
you have done for us..
And make those recollections cause us to draw near to you,.
to take a step forward, inch forward..
And knowing as we're reminded today in the word.
that we are going to be looking into,.
that God, you have made a way through your son.
for us to draw near to you and to do it boldly..
And we wanna do that today..
And it's in Jesus' name we pray, amen..
So brothers and sisters, you know,.
as already been mentioned a few times,.
this is a cause for great celebration today..
A wonderful thing is happening today..
I will call, you don't have to call it this,.
I will call it the new Urban Voice, okay?.
Don't change your signs, don't do anything like that,.
don't change your website, but the new Urban Voice,.
a new chapter for Urban Voice is starting today..

$^{41}$You know, the Peniel Multiethnic Church.
is joining forces with Urban Voice Community Church..
Two, what I believe genuinely,.
two great works that God has been doing..
And you know, the world might call this a merge.
or a merger, but you know, let's face it,.
it's really more like two families that are coming together..
It's like a marriage..
It's two becoming one..
And so praise the Lord for doing this great thing, amen?.
So I don't have a whole lot of time,.
'cause I know everybody is like, there's a luncheon today?.
So probably like, please, keep it short..
And I will do my very best..
In fact, I am gonna do more than my best..
I'm actually gonna keep it short today..
So there's a lot to say,.
but I'm gonna let the Holy Spirit speak to us.
through His word..
So if you'll turn in your Bibles,.
or turn your Bibles on on your phone apps,.
to Hebrews chapter 10, verses 19 through 25..
Hebrews 10, 19 through 25, and I'll read that for you..
And it says, "Therefore, brothers and sisters,.
"since we have confidence to enter the most holy place.
"by the blood of Jesus, by a new and living way.
"open for us through the curtain that is His body,.
"and since we have a great priest over the house of God,.
"let us draw near to God with a sincere heart.
"and with the full assurance that faith brings..
"Having our hearts sprinkled to cleanse us.
"from a guilty conscience,.
"and having our bodies washed with pure water,.
"let us hold unswervingly to the hope we profess,.
"for He who promised is faithful.".
Then verse 24, "And let us consider.
"how we may spur one another on toward love and good deeds,.
"not giving up meeting together.
"as some are in the habit of doing,.
"but encouraging one another,.

$^{81}$"and all the more as you see the day approaching.".
So there are three let us statements,.
three invitations, three exhortations.
that the author of Hebrews gives us..
First, we see in verse 22, "Let us draw near to God.".
Then in verse 23, "Let us hold unswervingly.
"to the hope we profess.".
And then third, "Let us consider.
"how we may spur one another on.
"towards love and good deeds.".
But before we get to these invitations,.
let's stop for a moment,.
looking starting at verse 19 at the foundational conviction..
You know, we're pointed to this in verse 23, right?.
What is this hope we profess?.
What is the foundation of this invitation?.
And it is this, the gospel of Jesus Christ.
is the foundation of the church's.
intimate relationship with holy God,.
and with one another..
The gospel of Jesus Christ is the shared foundation.
of the church's intimate relationship.
with God and with one another..
The church, we're told in verse 19, just by the address,.
"Who is the author talking to?".
And he says, "Brothers and sisters.".
If you don't know who your brothers and sisters are,.
just look to the left, look to the right,.
look behind you, go ahead, you can do that,.
it's all right, look around..
These are your brothers and sisters.
who share this foundational conviction.
in the gospel of Jesus Christ..
Is the writer of Hebrews writing to a bunch of strangers.
who don't know each other, or some,.
just kind of like group of people.
that happen to be together one day?.
No, of course not, he's specifically addressing.
a known and exclusive, and I'll say it,.
an exclusive group of people..

$^{121}$We hear a lot of people say about like,.
"Our church is inclusive, we want everybody to come.".
But the reality is, is the church is exclusive..
You know why?.
Because you are not really a member of the church,.
the body of Christ, unless you are in Christ..
And you have to have the gospel to be in Christ..
You have to believe who Jesus is,.
believe in him and what he has done on the cross,.
to lay down his life for sinners like us..
And then to rise again in victory on the third day..
To give us new life, eternal life,.
to guarantee the forgiveness of our sins,.
and a relationship with holy God..
I mean, who do we think we are?.
We know how messed up we are..
We know how weak and flawed we are..
Who do we think we are that we think.
we should have a relationship?.
And yet you know what the gospel writer John says?.
That we have the right, the right..
We have the audacity because of Jesus.
to claim to be children of God..
Let that sink in for just a second..
That ought to just blow your mind..
Who do we dare to be to say we are children of God?.
That we would have the right to be called children of God?.
Formally we might call one another members of a church.
or the church..
Says therefore brothers and sisters,.
we have confidence to enter the most holy place..
Here's the foundation, the gospel..
It's probably maybe explained in a way.
that you haven't thought of it..
But it says we have confidence to enter.
the most holy place..
It's a picture of the temple..
In the temple is the smallest room..
It's called the Holy of Holies or the most holy place..
And it's a place where one person once a year,.

$^{161}$the high priest gets to go in..
Once a year..
And if he's not worthy, if he has done,.
if he has not prepared himself to go into the presence.
of God, then he is to be killed..
Killed..
He would die..
But it says we have confidence, not just the one person,.
but we have confidence to enter the most holy place.
by the blood of Jesus, by a new and living way.
open for us through the curtain that is his body..
Whose body?.
Jesus Christ..
And since we have a great priest over the house of God,.
so Jesus is the high priest with a capital T,.
the capital H, capital P..
He is a superior high priest..
He's not like one of those of old..
He is even better..
You could say he is the super priest..
He says because Jesus is our high priest,.
he says let us draw near to God with a sincere heart.
and with the full assurance that faith brings..
And that's where the gospel comes in again..
We don't claim this because we're good enough..
We don't claim the right, this boldness,.
this confidence to enter into God's presence.
because we did enough, we measured up..
No, it's because the high priest, Jesus,.
did everything necessary so that by faith.
in what he has accomplished for us,.
by faith in who he is, we ride on his coattails.
and we have confidence because the high priest.
has done what's necessary..
So we boldly, confidently enter into this holy of holies,.
this most sacred place where even when the high priest.
doesn't do it right, he's at risk of losing his life..
We go in confidently..
And why?.
Because we're told our hearts have been sprinkled.

$^{201}$to cleanse us from a guilty conscience.
and having our bodies washed with pure water..
Because the church has this great high priest.
in Jesus Christ, let us draw near..
Here's the first invitation..
Let us draw near to God..
How many in this room have been away at college.
and you're back for Thanksgiving break?.
There's about a half dozen people in the room..
My daughter's back from college too..
I know what college is like..
I was at college..
And if you have gotten real busy with studying,.
playing, sleeping, whatever it is,.
and it has drawn you away from God,.
let me call you back..
Let me invite you..
Draw near to God with a sincere heart.
and with the full assurance of faith,.
absolute trust in Jesus..
On that day of atonement when the high priest.
goes into the temple, the people of God together.
experienced a national revival..
They had confidence to say, now we can face another year.
and we can expect God to bless us..
We can expect God to be present..
We can expect God to show us favor and prosperity..
We can expect, on the other hand, God will not kill us..
For our sin, he's forgiven us..
They could collectively breathe a sigh of relief.
like, oh, it's all good with God..
And what does this passage tell us?.
If you read this whole section of Hebrews,.
it's just this whole explanation of,.
look at all that the Old Testament prophets.
and priests have done..
All the animals that were sacrificed,.
and it had to be done every year over and over again.
because sin became a problem over and over again..
But now we have this high priest who does it once and for all.

$^{241}$and his sacrifice is so superior,.
he never has to do it again..
Once and for all, Christ has died for his church..
Once and for all..
That's as long as there's a high priest.
who has done it right..
A high priest that was acceptable to God,.
who followed every law,.
obeyed every command to the utmost..
But the church, as I said,.
has an even greater high priest,.
the Son of God, the Son of Man..
The superior high priest, the super priest..
For sinners, we cannot move forward..
We cannot move forward in a relationship with God..
We cannot go boldly into God's presence.
without great danger..
We just cannot unless we have a great high priest..
And for those who don't know, I'm here to tell you,.
there is a great high priest and his name is Jesus..
If you have not trusted in him,.
don't you dare even think about approaching God..
It's nothing but danger..
But if you have Jesus as your high priest,.
your Savior, who didn't slaughter a lamb,.
he himself was slaughtered..
He shed his own blood, his own body broken for us..
Trust that high priest 'cause he has made us a way..
He's made us guilt-free and he has purified us from sin.
so that the righteous judge might declare us righteous..
Righteous..
And that's what he does when we have put our faith.
in the righteous one, Jesus..
He has made us new, or he has made a new and living way.
for sinners to have intimate relationship with holy God..
This is the gospel of Jesus Christ,.
that he has secured God's forgiveness for us,.
that he has secured for us on the cross a birthright..
As I said before, the birthright that we would dare.
to claim to be sons and daughters of God..

$^{281}$Sons and daughters of God..
This is the gospel of Jesus and this is that foundation.
that we share together..
So let us draw near to God with a sincere heart.
and the full assurance of faith..
When you have trusted Jesus,.
you don't have to walk carefully,.
you don't have to walk sheepishly shy, afraid..
No, we can approach him as our God and father..
Second, in verse 23, let us hold unswervingly.
to the hope we profess..
That same gospel, for he who promised is faithful..
The church is holding unwaveringly.
to the gospel of Jesus Christ,.
her hope and her conviction and her claim..
The church is unshakable trust in the faithful one..
We trust the Lord Jesus Christ, that he has done enough,.
that he has done it right..
We trust his finished work on the cross..
There's no one else that we trust, is there?.
There's nothing we can add to perfectly finished work,.
is there?.
Whose gospel is the power of salvation.
for those who believe?.
It's the gospel of Jesus..
It's the gospel of Jesus, amen?.
Do we really believe that?.
Do we really believe that without this gospel,.
we are doomed, we are hopeless?.
So then we are able to spur each other on..
We are called to push and to press,.
to encourage and to inspire, to provoke each other,.
to move forward because we share this fundamental commitment.
to Jesus Christ and his gospel..
How many times in recent days have you heard somebody.
call people to unity?.
We have to stop being so divided..
I just heard it recently..
You know, when the Kyle Rittenhouse verdict was read,.
there was a call for peace and unity..

$^{321}$Let's not fight about this..
Let's be united..
Let's hold together..
And you know what?.
No matter who makes that call, whether it's the president,.
whether it's somebody influential in your life,.
a celebrity, no matter who makes that call, it's hopeless..
We can't be united..
There's too much that divides us..
Being a Chicagoan doesn't unite us..
Being an American doesn't unite us..
Being Chinese doesn't unite us..
A person, a citizen of the world doesn't unite us..
It all divides us..
But you know what unites us?.
Our hope, the gospel of Jesus Christ..
It's the only thing that unites the church..
So here's a little warning..
I'm not a prophet, so later on don't say,.
"Man, that John Yee was a prophet..
"He knew what was coming.".
I just know, and you know human nature..
There is going to be things that will come up.
that will divide us..
There are gonna be things that will happen.
that will cause us to say,.
"I don't wanna be with that person anymore..
"I don't wanna have anything to do.
"with those people anymore.".
But the question we need to ask ourselves is,.
but do we agree?.
Are we united in the gospel?.
Are we all one in Christ?.
Are we?.
Only the gospel of Jesus Christ can unite people..
Only Jesus is worthy of such commitment,.
loyalty, and devotion..
Only Jesus..
So when division appears again and it's imminent,.
let's hold together..

$^{361}$Hold on to the hope we profess, the hope we claim..
I know the gospel is preached here at least once a week..
Hold to that hope we profess..
Hold on to the promises of God in Jesus Christ.
to give us that new life, eternal life,.
to forgive us for our sins, to declare us righteous,.
and to say, "You are my children..
"You are my son, you are my daughter.".
Let us hold together to the one faith,.
Jesus, our superior high priest..
Let's hold on to him..
Third, let us consider in verse 24.
how we may spur each other on towards love and good deeds,.
not giving up meeting together.
as some are in the habit of doing,.
but encouraging one another,.
and all the more as you see the day approaching..
Mark, I'm about to say something.
that is probably gonna,.
that could potentially lead to saying,.
we wanna cancel your December visit..
Folks, what excuses are we making?.
People at home, what excuses are you making.
that seem like a good enough excuse.
to forsake meeting together?.
There are people all around the world,.
COVID or not, who have imminent threats to their life,.
and they are still meeting together..
I think it's time that a lot of people,.
the church needs to repent for all the excuses.
that they make for not meeting together..
I'm gonna say it right here, enough is enough..
Mark, people need to just give up their membership.
and say, you know what, I'm not coming back..
Stop making excuses..
Stop saying it's gonna be next month,.
or when the next wave of COVID is passed,.
or whatever it is, when my toe doesn't hurt anymore,.
or when I get a new pair of pants.
that I can wear to church..

$^{401}$Let's stop making excuses..
Gather together..
Let's stop forsaking the meeting together..
How are we gonna spur each other on.
when we are not together?.
How?.
How are we gonna encourage each other.
when we're not meeting together?.
The church faced all kinds of dangers..
And you know, I always find it interesting.
when I speak to my friends who come from places.
where the church was severely persecuted,.
and you're one of them, people who come from China,.
people who come from Vietnam, Romania,.
and so many other countries around the world..
You know when it's the most dangerous to be a Christian?.
When you're meeting together..
It is..
And yet they still do everything possible.
to meet together..
Because they know that there is a day coming.
when the Lord will return,.
and he will set all things right..
The earth will be his footstool..
He will sit on his throne as the righteous judge,.
and all will be well..
And those who are members of his body,.
those whose names are written in the Lamb's Book of Life,.
have nothing to worry about..
We encourage each other towards that day..
I know it's hard right now, I know you're suffering right now.
but the day is coming when Jesus will make all things right..
Let's hold firm to our hope together..
Let's hang on to his promises together..
Let's spur each other on to good works..
Let's spur each other on towards love,.
not giving up meeting together..
You know in Acts chapter two, verse 42,.
there's this description of the church..
They devoted themselves to the apostles' teachings.

$^{441}$and to fellowship, to the breaking of bread and to prayer..
Everyone was filled with awe at the many wonders.
and signs performed by the apostles..
All the believers were together and had everything in common..
They sold property and possessions.
to give to anyone who had need..
Every day they continued to meet together.
in the temple courts..
They broke bread in their homes.
and ate together with glad and sincere hearts,.
praising God and enjoying the favor of all the people.
and the Lord added to their number daily.
those who were being saved..
Until the Lord returns,.
the church is encouraged by being united..
They encourage each other when they are united.
by being together..
Don't forsake gathering as some are in the habit of doing..
Yes, worship together on Sundays.
and yes, study the word together on Mondays.
and then yes, break bread and pray together on Tuesdays.
and let's make sure that every day of the week.
we are sharing and caring, then breaking bread again..
Let's make sure that we keep worshiping together,.
encouraging each other, spurring each other on.
towards love and good works..
Do this until the day of the Lord..
The gathering of believers,.
which is what the word church literally means,.
ecclesia means gathering,.
it should give us a taste of heaven..
I'm not overstating that,.
it should give us a foretaste of heaven..
Gloria and I, my wife and I,.
we experienced that when we were in seminary..
Wow, how amazing it is to be in a place.
where everybody's kind of running in the same direction..
And then we've experienced that in churches.
we've been members of..
We've experienced that when we're surrounded.

$^{481}$by like college students,.
when we were going to college together.
and all the different things that college students go through.
and we're all running together because of Jesus..
An urban voice, 2.0,.
can and does experience this foretaste of heaven too..
Experience Jesus and all his benefits..
And so many of those are right here in the body of Christ..
Right here in the body of Christ..
I'm gonna ask you to bow your heads with me and pray..
I wanna give you a chance to just tell God.
how thankful you are for the gospel of Jesus Christ..
And if you don't know Jesus,.
then you don't yet have that to be thankful for..
But if you wanna experience the power of the gospel.
to change your life, to give meaning to your life,.
have you been wondering about your own identity?.
Who am I and why am I here?.
What's the meaning of all this?.
Jesus has that answer..
He is the answer..
And I'm inviting you to come,.
come, taste and see how good Jesus is..
And you can pray something like this..
God, I don't need that preacher to tell me I'm a sinner..
I know there's so much wickedness in my heart..
There is rebellion in my heart..
There is hate in my heart..
There's jealousy in my heart..
There are all these things that I wish were not there,.
but it's true..
There's so much sin in my heart..
And I cannot carry it anymore..
I cannot carry that guilt anymore..
And I'm asking you, Jesus, I'm trusting you by faith.
that you will relieve me, you will save me from that guilt..
And that's not all..
That you will unite me with the creator of the world,.
the one who knit me together in my mother's womb,.
the one who has fearfully and wonderfully made me,.

$^{521}$that I can know my creator, the God of the universe..
Lord Jesus, I want to know you..
I want to be a child of God..
Forgive me for my sin..
Forgive me for my offenses..
And declare me by faith, declare me righteous,.
pure, and holy..
And God, would you show me the new life that you promise?.
The new life today and tomorrow,.
but also the new life that comes.
when the day of the Lord comes,.
when all things are made right,.
that on that day I will be counted.
among the children of God too..
Thank you for Jesus, the perfect high priest.
who laid his own life down on the cross for me,.
paying the penalty for my sin,.
although he was without sin..
And thank you that he proved.
that he even has power over death.
when he rose again on the third day..
And I join the body of Christ, the church,.
the people who are around me, my brothers and sisters,.
in longing and waiting for that day,.
the day of the Lord, when you will return..
And you will claim me as your own..
God, as Urban Voice Community Church continues forward,.
I pray that you will bring them.
into a new season of fruitfulness,.
a season where many, many hear the good news.
and come to faith..
People who speak English and Mandarin and Cantonese.
speak all manners of languages.
that are around us in Chicagoland,.
so many, too many to count,.
that they will hear the good news.
and they too will be saved..
Lord, protect this church.
because there is an evil one who roars.
and he is like a lion looking for whom he might devour..

$^{561}$But we know that we have Jesus..
The Bible calls him the Lion of Judah, a greater lion..
And we ask you, Lord, hold your church together..
Unite us under Jesus..
Lord, as we go to lunch today,.
I'm just gonna pray it right now..
Let us have a great meal together and great fellowship..
Thank you for the bounty that will be on the table..
Thank you for the bounty that is around the table,.
the family of God..
Lord, we praise you, we celebrate you..
We thank you for the cross..
We thank you that because of Jesus,.
we are a family, the Urban Voice family..
And it's in Jesus' name we pray, amen..
\newpage



\section{7:1-23}
\label{sec:BxAa7Ok6gjk}
\textbf{Keep your heart from the evil more than from COVID (Mark 7\_1-23) - Rev. Jaeuk Jeong}
\newline
\newline
連結: \href{https://youtube.com/watch?v=BxAa7Ok6gjk}{\texttt{https://youtube.com/watch?v=BxAa7Ok6gjk}} ~~~~ 語音日期: 2021-11-29
\newline
\newline
\hyperref[sec:cKHN4UF2tsA]{\small{< < < PREV SERMON < < <}}
~
\hyperref[sec:index_chronic]{\small{[返順時目]}}
~
\hyperref[sec:index_scriptual]{\small{[返順卷目]}}
~
\hyperref[sec:CNrSX78PCg8]{\small{> > > NEXT SERMON > > >}}
\newline
\newline
$^{1}$My dear brothers and sisters, good morning..
Okay, you know what, last week actually was the first Sunday worship we had together after the merger..
So it seems like, you know what, it seems like transplanting our organ..
So it'll take some time..
So when you look at each other, actually they are you..
Okay, we are one body..
It means that they are part of you, you are part of they, part of them..
So you need to accept each other as one body of Christ..
Okay, so this is my homework for all of you guys..
So before you go home after this worship, as one body, practice oneness..
Like what?.
Meet three strangers around you..
So urban voices, urban voice, you guys, you meet your new brothers and sisters..
Learn them by their name..
Look at their eyes, ask their job, and ask their lives, okay?.
Do the same thing..
Okay, so Puneon, former Puneon, now you are urban voice, don't forget that..
So ask the other brothers and sisters how they are doing now, okay?.
This is practicing oneness in Christ..
So I don't know, can we have the title on the screen, today's message?.
The message of -- the title of today's message is like this..
Keep your heart more from the evil than from COVID..
Did you get it?.
What is the title?.
Keep your heart more from the evil than you are trying to keep from -- keep yourself from COVID..
So our COVID situation is going on and on..
And this COVID situation is the most inconvenient and fearful and dreadful period of our lives..
So what did you experience?.
Do you remember last year, March, COVID situation started in Chicago and then a few months, for the following a few months, everything was closed down..
Do you remember lockdown?.
There was no people on the streets..
And all the public places were closed..
We couldn't do anything..
There was COVID-19..
And how about these days?.
I have three children..
They go to the public -- they go to the school..
Whenever they have any symptom, like a very light symptom or a mere little fever or just if they -- even if just they sneeze, they need to take the COVID-19 test..
Okay?.
So COVID-19 is now -- has been so serious for all of us..

$^{41}$Even after one and a half years later now, more than that, what did you learn?.
What is the greatest lesson through this COVID-19?.
What did you learn?.
I learned about the importance of cleanliness..
So you guys place some -- what is that?.
Sanitizers in front of the doors everywhere..
What do you think of?.
Cleanliness..
Through this COVID-19, what we learned most is cleanliness, the importance of cleanliness..
And in today's story, Jesus and the Pharisees and teachers of the law, they were arguing about the true cleanliness..
Why?.
When Jesus' disciples, they joined eating the food, Jesus' disciples, they were maybe -- they were too hungry..
So they didn't wash their hands..
So what did they do?.
Jesus' disciples didn't wash their hands, and they were beginning to eat the food..
When the Pharisees and teachers of the law, when they looked at them, they were shocked..
What are you guys -- what are you doing?.
We Jewish people, if we do not wash our hands, we do not eat any food..
That is our cleanliness..
We are holy people..
That's what the Pharisees were thinking..
Right?.
They were super holy people, super clean people..
So they were talking about some cleanliness..
But when Jesus heard it, Jesus began to talk about another kind of cleanliness..
What is that?.
Jesus said, when you can read your Bible, please keep your Bible open, and you can go back to the verses again and again together..
When you look at verse 15, Jesus says like this, "Nothing outside a man can make him unclean by going into him..
It is rather, it is what comes out of a man that makes him unclean.".
Did you hear that?.
So here, when Jesus heard the Pharisees and teachers of the law, Jesus is giving a totally different interpretation about the true cleanliness..
So, what is that?.
The Pharisees and teachers of the law, they were talking about physical cleanliness, visible, which is visible..
Visible, physical cleanliness..
But Jesus was talking about what?.
Spiritual, invisible cleanliness..
Which is more important?.
Pharisees and teachers of the law, they were keeping the law so well..
So they thought we are holy people..
But when Jesus was looking at them, Jesus knew that they were actually unclean..

$^{81}$Who is unclean?.
The Pharisees and teachers of the law..
For an example, Jesus was pointing out their fault..
You are washing your hands before eating the food, but actually you are not holy people..
You are unclean, dirty people..
Why?.
Because look at what you are doing now..
What is that Jesus is pointing out?.
Jesus talked about Korban..
Everybody, Korban..
Korban..
C-O-R-B-A-N..
Korban is some amount of money, very precious gift set aside only for God..
So the people, they set aside quite a lot of money, and they called it Korban, which is a gift to God..
Then what happens?.
If you set aside a lot of money for God, it means that you don't need to support your parents..
So many Jewish people, they were avoiding the duty of supporting their parents, making their money as a kind of Korban..
But actually in their tradition, they were not supporting their parents..
At the same time, they were not worshipping God with their true genuine heart..
That is their heart..
So when Jesus was talking about this, what did Jesus do?.
This is not the only thing..
When you go down to verse 13, Jesus says that this is not the only thing..
Actually, Jesus says, "You are doing many other things like this Korban.".
So look at you guys..
You are trying to keep the law, but actually you are ignoring the law..
You are not keeping the law..
So are you really clean?.
Jesus was asking the Pharisees and teachers of the law..
Then my brothers and sisters in this place, which is more important?.
Physical cleanliness, visual cleanliness..
When you guys are coming to the church, if you do not wash, people will look at you..
"Oh, look at the guy. Oh, smell. Did you brush your teeth? We don't like it, right?".
But how about the cleanness?.
What is the true cleanness?.
Physical, visual cleanliness, one side..
The other side, spiritual, invisible cleanliness, which is more important..
Actually, you cannot compare..
If you are not clean in your body, it's okay..
People will frown, make a frowning face at you..

$^{121}$That's it..
But if your heart is not clean, what happens?.
Because of this dirtiness, people go to hell..
If our conscience, if our mind and heart is not cleaned by Jesus' blood,.
if our heart is not clean, then what happens to the people?.
They will go to hell..
Spiritual uncleanness gives us that penalty, which is more important..
What else?.
If our spirit is not clean, what happens?.
It governs all the areas of our lives..
Actually, our spiritual condition shapes and decides the whole quality of our lives..
Are you now painful because your spirit is not clean?.
Are you now anxious or fearful because your spirit has a problem?.
Your spiritual condition decides everything in your life..
I had a couple of days of short trip this week, and then I couldn't brush my teeth every time, three times per day..
When I didn't do it, when I came back from the trip, I realized my teeth..
I felt a little pain in my gum because I'm an old man..
But listen, if our spirit is not clean, day by day, moment by moment,.
we harvest what we sow when we do not take care of our spirit..
Then let's think about this..
When you read the scripture together, what is the source of the spiritual uncleanness?.
If you do not brush your teeth, what is the source?.
Yes, a little tiny piece of the food inside and between the teeth..
That is the source, right?.
Of the stinky smell..
But how about the source of the spiritual uncleanness?.
What is that?.
This is very important..
My brothers and sisters, listen to me very carefully..
This knowledge is really, really important..
You young people, you need to listen very carefully today..
What is the source of the spiritual uncleanness?.
Verse 15..
Can we go back to verse 15?.
It's very important..
Do you see any verse?.
Okay, verse 15, listen to me..
It is what comes out of a man that makes him unclean..
Nothing outside of man makes him unclean..
Rather, it is what comes out of a man that makes him unclean..

$^{161}$What did that mean?.
When disciples heard this verse, they didn't understand, just like you..
They didn't understand..
How come?.
Something out of us, it makes unclean..
It is not from outside, but from our inside, that is the thing which makes us unclean..
The disciples didn't understand the first explanation..
So when Jesus was alone with the disciples, Jesus was explaining one more time..
From verse 21 to 23..
Listen very carefully, you young people..
This is very important knowledge..
Verse 21..
For from within, out of man's heart, come evil thoughts, sexual immorality, theft, murder, adultery, greed, malice, deceit, lewdness, envy, slander, arrogance, and folly..
All these evils, verse 23, all these evils come from inside..
Where?.
Not from outside, but from inside and make a man unclean..
What does that mean?.
All these evil thoughts, what are the evil thoughts?.
When they enumerate, when Jesus was enumerating all these unclean evil thoughts, Jesus summarized the first one, evil thoughts..
What are the evil thoughts?.
Coming out of your mind, inside..
Number one, sexual immorality..
All kinds of dirty desires, sexual desires, dirty desires..
Where does it come from?.
And then theft, what is that theft?.
You want to steal..
Are you paying the tax honestly?.
Even so many Christians, they are deceived by the Satan, by the evil spirit..
And even paying the taxes, they are not honest..
So when you hear the word theft, you are thinking about really huge rich people?.
No..
So many of us, we are choosing the wrong things..
Theft, that is stealing some money in whatever ways..
Or murder, what is murder?.
Jesus said, if you have a hatred inside you, if you have a small hatred inside you, it's just stealing..
It's killing people..
Adultery, adultery..
Greed, what is greed?.
You desire more than enough..
You have already enough..

$^{201}$Who is now hungry?.
Now in our society, you are hungry?.
Pastor Mark, Rachel is now hungry..
She raised, okay, not that kind of meaning..
You have enough money to buy food, but what I'm saying is, like, nobody is hungry in our society now..
But people are desiring more..
People are desiring more and more..
There is no satisfaction..
There is no thankfulness..
Greed, how about malice?.
What is malice?.
Young people, what is malice?.
Okay, you did the wrong thing to me..
Okay, I'm going to hate you until I die..
Malice..
How about deceit?.
Are you honest in dealing with your job?.
Dealing with any responsibility?.
Deceit, lewdness, what is lewdness?.
Sexually dirty in words and thoughts and behavior..
Envy, what is envy?.
You have already enough, but you think that the person, what the person has now is what I now want..
So you continuously desire what others have..
Envy..
Until when?.
Until you die..
You are never fully satisfied..
You never understand that God is the wonderful provider..
And then slander, what is slander?.
You talk back, you're biting persons, any character, anything behind their back..
Okay, how about arrogance?.
Pride and folly..
All kinds of foolish things..
You choose foolish words, foolish behavior, but you think that it's okay..
When you look at other people, when they are doing it, you are shocked..
Why they are doing such a foolish thing?.
But actually, that's what we are doing..
Then brothers and sisters, where do all these evil thoughts come from?.
Today's passage is very important..
It is not from outside, but it is from inside..

$^{241}$Inside..
Who is inside?.
So when you go to the scripture, 1 John 3, verse 8-9..
When you go to the scripture, 1 John 3, verse 8-9, it goes like this..
He who does what is sinful is of the devil..
Did you hear me?.
He who does what is sinful is of the devil..
Because the devil has been sinning from the beginning..
What does this verse talk about?.
Any sin, every sin is from the Satan..
It's from the enemy..
Because when God made us, we were made in God's image..
What does that mean?.
We love, we desire the true love..
We desire the holiness..
We desire the cleanness..
That's how God made us..
But all these kinds of evil and dirty thoughts come from the Satan..
So where is Satan?.
In your heart..
Inside you..
In your thoughts..
In your emotions..
They have built their house in there..
So this is a very important understanding..
Jesus said it is not, it is what comes out of a man that makes him unclean..
So who is inside us?.
The Satan, the evil spirit..
You may wonder, how can it happen?.
Yeah, that's what the Bible reveals..
That's the truth the Bible reveals about..
Then, from today, when you have that kind of dirty and evil thoughts, emotions, imaginations, and feelings, and all those things,.
what should you think?.
Oh, you Satan..
Are you trying to enter my heart?.
Again, you should understand..
All those evil desires is not your desire..
You need to understand that those desires are from the Satan..
Okay?.
That is the teaching..

$^{281}$That's what Jesus teaches today in today's passage..
Then, let's think more about this..
How does Satan's uncleanness enter us?.
Okay, Satan can enter through our thoughts, through our emotions, especially through our hearts..
Then how does it happen?.
What is our heart?.
Our heart is like a cockpit..
Do you know cockpit?.
Cockpit is the small place in front of the airplane where the pilot flies the whole airplane..
It's a very small place, cockpit..
Our heart is like a cockpit..
Our heart is the center of our being..
In our heart, what happens?.
In our heart, we decide and we discern which is right and which is wrong..
At the same time, we know that, "Oh, that's good. That's not bad. That's good. That's bad.".
Our heart is the place we discern and we decide everything..
So our heart is the center of us..
Right?.
Then this heart can be stolen by the evil spirit..
So actually, when we were talking, Pastor Mark read some parts, but actually read a little bit more than I expected, but it was good..
Why?.
Because when Pastor Mark read the following passage, it shows two people..
One is a little daughter of a woman..
The little daughter was possessed by an evil spirit..
Possessed..
The Bible says possessed by the evil spirit..
Second case was like this..
What is the second case?.
The last story. Do you remember?.
What is that?.
Death man..
Jesus says the death came from the evil spirit..
So when you look at these people, when you look at the Bible, you need to understand that much, much more than we can think, ever think..
There are a lot of evil possession cases..
We see like a physical condition..
Like some psychological issue, mental issue, or a physical issue..
But behind that, a lot of spiritual issues too..
Okay?.
Then, in another story, when you go to Mark's Gospel, chapter 7, when Jesus was expelling the evil spirit, Jesus says, when Jesus called the spirit, Jesus used evil spirit..
But in Greek language, it says unclean spirit..

$^{321}$Do you get it?.
So evil spirit, what they are doing is all unclean things..
Dirty things..
We love cleanliness, but actually, so many dirty things are coming from our hearts because of the evil spirits..
Now you are getting it a little bit?.
Then, how does the evil spirit enter our hearts?.
How does the evil spirit enter our hearts?.
How?.
When you think..
When you think, the evil spirit enters your thinking as if it were your thinking..
Actually, it was not part of your thinking, but the evil spirit comes into your thinking slightly, very secretly..
Okay?.
Or your emotion..
When you are really angry, evil spirit enters your heart..
Gives more anger..
More anger..
Breaking your things..
Breaking the relationship..
Breaking the most precious thing in your life..
Or some strong desire..
Very strong sexual desires..
Now young people, you need to listen..
In your days, if you do not keep your heart pure, you are going to lose everything in your life..
If you do not keep your heart pure, you are going to lose everything in your life..
Because it is from the Satan..
I'm telling you..
Whatever you cannot do in front of other people, it is from the Satan..
Listen. What is sin?.
Sin is everything you cannot do in front of other people..
Because it's dirty..
It's not right..
That's why you cannot do it in front of other people..
So through your thoughts, through your desires, or some words..
Some words are trying to enter your heart..
I experienced a lot of times..
Now I know how to fight this spiritual battle..
So many years..
And God gave me a victory..
I couldn't do it..
But God gave me the victory over a long period of time..

$^{361}$But still, the Satan tries to enter my heart..
I can feel it..
Or sometimes, very grotesque, evil pictures..
Some dirty thoughts and pictures in your heart..
You know what?.
This is not yours..
But it is the Satan's..
It is the evil spirit in disguise..
And then all these things, dirty, evil things, they are all lies..
You need to understand..
They are not true actually..
But they are all lies..
About you..
Lies about what?.
Lies about you..
You guys, at least as far as I know, all the brothers and sisters in here..
You love Jesus..
You are saved..
You want to be holy..
But that kind of all dirty, evil desires and thoughts, they are not yours..
They are from the Satan..
So what is that?.
It's a lie..
So John's Gospel chapter 8 verse 44 says like this..
John's Gospel chapter 8 verse 44..
You belong to your father, the devil..
Here you is Pharisees and the teachers of the law..
Okay?.
Just religious people..
Religious people..
You belong to your father, the devil, and you want to carry out your father's desire..
He was a murderer from the beginning, not holding to the truth..
For there is no truth in him..
Listen here..
When he lies, when the Satan lies, he speaks his native language..
For he's a liar and the father of lies..
Who is a Satan?.
Satan, whenever he opens his mouth, Satan is a liar..
He never says any truth about you, about God..
Okay?.

$^{401}$So whatever the Satan says, it's a lie..
It's a lie..
And he's the father of lies..
His native language is not Cantonese..
His native language is not Mandarin or Korean..
His native language is lie..
Do you understand?.
How many times did he lie to you?.
Do you understand?.
And then when you listen to the Satan, what will happen in your life?.
What is the time now?.
I'm sorry..
When does the worship end?.
Ten more minutes?.
Okay, ten more minutes..
So listen..
When you listen to the Satan or the unclean thoughts and desires, this is the things that will happen in your life..
You're going to break your relationship, broken relationship..
And there are all kinds of shame, shame and guilt in your life..
We had a great man in our country..
His name was Park, just a family name..
He was a mayor, mayor of Seoul city, metropolitan city, 10 million people..
He was a mayor, very powerful man..
He was an activist when he was young, so people respected him..
He had everything now, mayor of Seoul metropolitan city..
But you know what?.
He yielded into the wrong desire and what happened to him..
He had a sexually wrong thing with his secretary..
And the secretary was going to expose the sin..
And this guy, he went on to mountain, he committed suicide..
Just one single wrong desire killed him..
My brothers and sisters, this is what the Satan does..
And then, listen, how can we keep our hearts clean?.
My brothers and sisters, how can we keep our hearts clean?.
Number one, guard your heart..
Everybody, guard your heart..
Guard your heart..
Proverbs chapter 4, verse 23 goes like this..
"Above all else, guard your heart, for it is the wellspring of life.".
You guys, I have Chinese brothers and sisters here,.

$^{441}$but I was shocked when I saw you guys..
Because many brothers and sisters, they have a camera,.
the security camera in their houses..
I don't think they need it because they are not rich people..
But many brothers and sisters, they have a security camera in their houses..
They lock the door and then it is not enough..
So they put the security camera in their houses..
Why? To guard, to guard..
How about your heart?.
Are you setting a guard in your heart?.
Or does every dirty thing come in, goes easily?.
How is your heart? You need to check..
How can you keep your heart?.
What is the guard?.
If you do not know the word of God,.
then you don't know how to keep your heart..
When you check God's word, you need to know that,.
"Oh, that is right, or that is wrong.".
You can set a guard according to God's word..
That's why you need to study God's word day and night..
Young people, read God's word..
Don't go into any blog or Facebook..
Facebook after Facebook, it will kill you..
Listen. Read God's word. Read God's word..
Second one, how can we be clean in our hearts?.
Number two, throw away the sources of uncleanness..
Throw away all uncleanness..
The source of uncleanness. What is that?.
The Bible, Mark's gospel, when you go down more,.
in the same chapter, chapter 9, Jesus said,.
"If your hands do something sinful, cut it off..
If your legs do something sinful, cut it off..
If your eyes are doing something wrong, pluck it out.".
There is a seriousness, seriousness of uncleanness, sins..
1 John 3, verse 9 goes like this..
"No one who is born of God will continue to sin.".
If you are born of God, you should stop doing the same sin again and again..
If you are repeating the same sin again and again,.
it means you are listening to the enemy..
This is the enemy's lies..

$^{481}$You cannot overcome..
That is who you are..
You've been enjoying it for more than a decade..
You cannot get out of that..
My husband cannot change..
My wife cannot change..
That's a lie..
When Jesus' blood touches us, we should change..
Like what?.
What is throwing away all the sources of uncleanness?.
When you have a pile of dung..
Do you like dung?.
You don't like it?.
If you have a good pile of dung, flies will come, surely..
But if you clean it up, no more flies..
In the same way, if you throw away all the evil things from your heart,.
from your lives, from your habits, then the evil spirit will leave you..
Third one, come near to the clean gathering and people..
Come near to the clean gathering and people, just like this..
Like what?.
When Jesus was walked about, people wanted to touch Jesus' cloak..
Just Jesus' clothes..
Why?.
They knew that even if I touch Jesus' body or Jesus' clothes,.
I'm going to be healed..
Unclean things are contagious..
At the same time, cleanliness is contagious, too..
So if you come to the clean people, our heart becomes cleaner and cleaner..
Number four, wash your uncleanness with Jesus' blood..
Our guilt and shame of sin is cleansed only by Jesus' blood..
If you have guilt, if you have any kind of shame, I had a lot..
In my 30s, I was a sinner..
I had been going to the church for a long time, but I had a lot of sins inside me..
Guilt, guilt and shame..
I had to be covered with Jesus' blood fully, and I got free..
I was freed..
1 John 1, verse 7 says, "The blood of Jesus, the Son, purifies us from all sin.".
What does that mean?.
Cleanses us from all sins..
Okay, last one..

$^{521}$So you know what?.
What do you do?.
If you are repeating the same sin again and again, you want to get out of it, but you cannot do it,.
you need to cry out, "Jesus, save me from my sin and cover me with your sin!".
"Remove all the shame and guilt! Make me whole! Make me clean!".
Okay?.
Cry out in your prayer..
Last one..
How can we keep ourselves clean all the time?.
Be filled with the Spirit..
Ephesians 5, verse 18 says, "Do not get drunk on wine, which leads to debauchery.".
Instead, be filled with the Spirit..
If you are drinking Facebook after Facebook, all kinds of Instagram and those things,.
you are drinking the unholy, unclean things..
How can you be filled with the Holy Spirit?.
It means that when you sit at home, when you are eating, when you are sleeping, when you are studying,.
you are filled with God's thinking..
You are filled with God's emotion, God's Word, God's idea, God's feeling..
Then you are Spirit-filled..
How can that happen?.
If you give the center of your heart to God..
If not, you cannot be filled with the Spirit..
We are easily filled with the Spirit of the world..
We need to be very careful..
Okay?.
So, let me close, my dear brothers and sisters..
We don't like unclean vessels..
If we have an unclean container, we cannot take the food from it..
We don't like it..
If you look at an unclean person, we know that the person has a problem..
But as I told you, as I challenged you, how is your heart?.
How clean is your heart?.
You need to ask today..
And you need to receive the cleanness only coming from the blood of Jesus Christ,.
which is our hope..
Let's pray..
Our Father God, you made us after your image,.
holy and righteous..
Fearfully, wonderfully, you made us..
But Lord, we have lived in this world, and we were stained..

$^{561}$We got all kinds of uncleanness, and we didn't know that it was a sin..
We desired it. We accepted it. We received it. We welcomed it..
So our hearts became so defiled..
But today, Jesus, Holy Spirit, come upon your people..
Remove all the uncleanness coming from the Satan..
Let us understand that all these evil desires are not from us..
They are not ours, but they are from the Satan..
So give us the right understanding..
At the same time, Holy Spirit, be our God in our hearts..
Give us the desire to love you more, more than anything else on this earth,.
that we can love you with our clean hearts,.
and make us clean and healthy in our body and spirit when we love you, Lord,.
when we pursue you..
In Jesus' name we pray. Amen..
\newpage



\section{以賽亞書 7:13-17}
\label{sec:CNrSX78PCg8}
\textbf{以馬內利的兆頭 (以賽亞書 7\_13-17) - 陳崇基牧師}
\newline
\newline
連結: \href{https://youtube.com/watch?v=CNrSX78PCg8}{\texttt{https://youtube.com/watch?v=CNrSX78PCg8}} ~~~~ 語音日期: 2021-12-07
\newline
\newline
\hyperref[sec:BxAa7Ok6gjk]{\small{< < < PREV SERMON < < <}}
~
\hyperref[sec:index_chronic]{\small{[返順時目]}}
~
\hyperref[sec:index_scriptual]{\small{[返順卷目]}}
~
\hyperref[sec:e3EJDUy_J3k]{\small{> > > NEXT SERMON > > >}}
\newline
\newline
以賽亞書 7:13-17
\newline
\begin{longtable}{cl}
\hline
\hline
章節 & 經文 (和合本修訂版)\\
\hline
7:13 & \begin{tabularx}{0.7\textwidth}{X} 以賽亞說：「聽啊，大衛家！你們使人厭煩豈算小事，還要使我的神厭煩嗎？ \end{tabularx} \\ \\ \relax
7:14 & \begin{tabularx}{0.7\textwidth}{X} 因此，主自己要給你們一個預兆，看哪，必有童女懷孕生子，給他起名叫以馬內利。 \end{tabularx} \\ \\ \relax
7:15 & \begin{tabularx}{0.7\textwidth}{X} 到他曉得棄惡擇善的時候，他必吃乳酪與蜂蜜。 \end{tabularx} \\ \\ \relax
7:16 & \begin{tabularx}{0.7\textwidth}{X} 因為在這孩子還不曉得棄惡擇善之先，你所憎惡的那兩個王的土地必被撇棄。 \end{tabularx} \\ \\ \relax
7:17 & \begin{tabularx}{0.7\textwidth}{X} 耶和華必使亞述王臨到你和你的百姓，並你的父家，自從以法蓮脫離猶大的時候，未曾有過這樣的日子。 \end{tabularx} \\ \\
[1ex]
\hline
\hline
\end{longtable}
$^{1}$各位電視節目好.
在這兩個星期我們會看看聖誕節有關的經文.
這些我們不能夠避免的.
因為已經到章林節了.
還有幾個星期就到聖誕節.
剛才報告有提過.
我們有什麼措施.
或者有什麼再要考慮的事情.
我們很快會在群組跟大家說的.
我們知道在疫情當中.
經常都是反覆.
經常都是來來往往.
我相信也不會很快會有結束的.
我們有一些心理準備.
不過結不結束也好.
在聖誕節期間.
我們知道有一樣東西是肯定的.
有一個事實就是基督曾經降世.
在我們中間.
所以在這幾個星期我們會看看聖誕節的一些經文.
想起聖誕節的經文.
想起這一段.
就是《以賽亞書》第七章十四節.
因此主自己必給你們一個少頭.
看哪必有同女懷孕生子.
他要給他起名叫做以馬來尼.
一個以馬來尼名字的英孩.
我們經常在聖誕節的時候.
提到這一段聖經.
我們說這段聖經是講到基督的降臨.
是吧.
就中了一半的.
但是有另外一半.
我們今天可以看看.
就是我曾經跟大家講過.
在聖經裡面.
每逢有一些經文出現.
都是不是直接寫給我們的.
而是為我們而寫.
直接寫給當時的人.

$^{41}$所以《以賽亞書》第七章.
是寫給當時以賽亞的群眾.
以賽亞一定不是21世紀的人.
是主前大概七百多年的時候.
大概二千七百多年之前的一個先知.
他所講的話 所講的講道.
所以我們需要了解一下.
到底這裡為什麼講以馬來尼.
以馬來尼呢.
所以有一些背景.
我們一定需要去了解.
在時間表上很簡單地告訴大家.
以賽亞書一至三十九章的背景.
我們知道在大概三千年之前.
聖經裡面記載.
蘇羅 大衛 所羅門.
這段時間大概有一百年左右.
也不夠一百年.
那個所謂聯合王國的時候.
就是當時以色列國十二支派聯合的時候.
大概不夠一百年.
所羅門死了之後.
聯合王國就分裂成為南國和北國.
北國叫什麼名字.
叫做以色列 是不是.
或者叫做以法列.
很多時候講到以色列的時候.
很多時候會叫以色列做以法列.
因為上面十個支派.
北面十個支派.
以色列十個支派.
最大的那個就是以法列.
最有影響力的就是以法列.
所以通常我們叫做以法列.
那就不會那麼混亂了.
是不是.
北國叫做以色列.
南國叫什麼.
猶大國 是不是.
Judah.

$^{81}$京都就是在耶路撒冷.
北國的京都就是聖馬利亞.
在聖馬利亞.
在這段時間裡面.
就是一段很混亂的時候.
同時之前聯合王國.
或者剛剛分裂的時候.
世界都是沒有一些超級大國.
super power出現.
直到大概八百多BC.
到七百BC的時候.
那段時間就很不同了.
就是黑色的那個box.
就是亞述帝國出現.
由於亞述帝國出現.
漸漸慢慢強大的時候.
它就去到不斷影響著.
西面的一些小國.
包括以色列國和猶大國.
所以我們看聖經.
看到舊約的歷史的時候.
特別是聶王記.
聯合王國之後.
歐羅馬納斯達克羅門死了之後.
我們就發現.
有一些超級大國出現的時候.
它就不斷的想要稱霸世界.
這段時間就正是.
以賽亞亞書.
以賽亞亞事奉的時間.
特別是講出.
以賽亞亞一至三十九章裡面的.
講章 預言等等.
所以你看到這個圖表.
你就發現在這段時間.
猶大國和以色列國.
在一種危機 危機中.
講起危機.
我們大家很清楚.
因為今天就算生命裡面.

$^{121}$自己沒有危機.
無論一些親友.
他在一些危難裡面.
或者我們自己身體.
發生一些不同的毛病.
我們處於一些不同的困境也好.
我們都會叫它做危機.
在世界裡面.
更加就是我們.
纏繞著我們年多.
繼續會纏繞我們的疫情.
在這樣的危機裡面.
所以就在這段時間.
以賽亞亞去到事奉的時候.
就講出這些話.
一至三十九章.
一段很長很長的.
一些講道或者一些言論.
剛才我們看那段聖經.
二馬內地.
童女懷孕生子.
就是在這個時候出現的.
在這個危機當中.
以賽亞亞就將這段話.
帶給當時的南國猶大國.
當時的王就叫做阿哈斯.
阿哈斯王可能大家沒有什麼印象.
你記得的只需要是這個.
就是他一個很壞的王.
所以阿哈斯王就接受了.
聽到以賽亞亞所講的話.
這個圖表告訴大家.
整個以賽亞亞書頭一半.
一至三十九章.
講了不同的預言.
但是在裡面記載了.
兩個亞述帝國所造成的危機.
第一個就是阿哈斯王.
那個很長的字.
提勒拉的那個.

$^{161}$Tiklapalisher.
你先不要理他.
那個是亞述王的名字.
就是他造成南國猶大國很大的危機.
然後就十三至三十五章.
又講了不同的預言.
到最後三十六至三十九章.
就另外一個的亞述帝國的危機.
一個亞述帝國所造成的危機.
當時造王的就是希瑟加王.
而且當時亞述帝國的帝王.
就是塞納基納.
無論如何都記載了兩個不同的危機.
歷史裡面的危機.
就是猶大國要面對這些危機的時候.
到底他怎樣去到部署.
怎樣去到聽見上帝對他說的話.
亞述王我們講過了.
就是一個很壞的王.
記得他是一個很壞的王.
是個壞小子.
講到什麼呢.
講到他不行大衛所行的.
慎看為正的事.
反倒去到跟隨以色列諸王的道路.
就是跟了北國.
北國也是很壞的.
南國這個猶大王亞述亞述.
也是很壞.
更加去到焚燒自己的兒子為祭.
好像那些迦南人.
好像以前那些外族人所行的可憎惡的事.
在丘壇這些樹底和山崗上面燒香.
來獻祭給別神.
這個就是亞述王.
正是在他造王的時候.
就發生了亞述的危機.
亞述步步逼近的時候.
猶大國就處於一個危機.
當時發生什麼事呢.

$^{201}$就是因為你看到右上角亞述帝國.
不斷向西擴展.
向地中海東面的一帶以色列巴勒斯坦地進發.
擴展的時候就影響了當時的小的國家.
那些小的國家要做什麼呢.
在亞克斯作王的時候.
北國以色列和亞蘭國這些小的國家.
就想有一個聯盟.
現在說的是當時的國際形勢.
這些小的國家看到一個很大的超級力量來到的時候.
影響他們和威脅他們.
他們很多時候會聯盟一起聯軍.
來阻擋亞述帝國向西擴展.
所以亞蘭和以色列國是其中兩個小的國家.
沿海也有 譬如說推羅等等.
聖經裡面經常看到這些小的城邦的名字.
這些小的國家就會組織在一起.
聯軍想擊退或者阻擋亞述帝國超級力量來到.
在亞克斯作王的時候發生什麼事呢.
就是以色列和亞蘭已經組成了聯軍.
他們想要南國和猶大國也跟他們一起玩.
要他們一起和他們聯盟來阻擋亞述帝國.
亞克斯就左左右右十五十六.
他不知道怎麼做.
在這個時候在他國家的朝廷裡面.
有很多不同的意見.
到底是和這兩個小國聯盟.
去擊退亞述帝國.
還是有其他外交的政策可以搞定呢.
結果他說不肯.
於是上面的以色列國和亞蘭國.
聯軍先想攻打猶大耶路撒冷.
拿了他們然後廢了亞克斯王.
不過他們喜歡的王.
然後他們三國或者更多的小國.
就能夠一起聯盟來阻擋亞述.
所以就在這個冰林城下的時間.
就是到以色列書第七章.
如果你看看經文七章第一節.
那裡說烏西亞的孫子約坦的兒子猶大王阿哈斯.

$^{241}$在位的時候.
亞蘭不要管他什麼名字叫李迅.
還有利瑪利的兒子.
是以色列國的國王比加.
就是說亞蘭國和以色列國.
就上來攻打耶路撒冷.
卻不能把他攻陷.
所以就上去圍城.
七章第二節有人告訴大衛加雪.
亞蘭已經在以法連帝.
就是以色列地安營.
於是發生什麼事.
王的心和人民的心都震動.
好像樹林裡面的樹木一樣.
被風搖動一樣.
怎麼辦呢.
在這樣的大勢裡面.
當冰林城下要去到圍剿耶路撒冷的時候.
阿哈斯王不知道怎麼辦.
所以在這段時間就是一個不知所措的時候.
這個時候先知就出現.
就是以賽亞.
以賽亞就被上帝去猜牌.
就去到和阿哈斯講話.
大家要想清楚.
搞清楚什麼叫先知.
什麼叫Prophet.
先知不是說今天下雨了.
幾點鐘會下雨.
所以是準的.
先知不是Weather Forecast.
不是天氣報告.
先知是什麼呢.
先知有一個說法叫Crisis Manager.
就是說當國家或者人民.
生命裡面有crisis來的時候.
先知就去到和他講說話.
告訴他應該怎麼管理.
他現在面對的困境.
所以那個crisis就成為了teaching moment.

$^{281}$就是教導的時候.
上帝要面對危機的你.
你面對困境不同問題的時候.
上帝要和你講話.
而和你講話的媒介就是先知就是Prophet.
Prophet就是做什麼呢.
就是做Crisis Management.
帶你去了解到底現在的困局是怎麼樣.
於是以賽亞就去見阿哈斯王.
記得這個壞人的王.
Ahes.
以賽亞就去找Ahes和他說一些話.
七章第三節.
耶和華對以賽亞說.
你帶著你的兒子.
原來以賽亞有個兒子.
叫做斯亞瓦述.
意思就是說.
被擄的人一定會回來等等.
總之帶著他的兒子去到哪裡呢.
去到上池的水溝盡頭.
就是漂泊在田間的大路上去迎接阿哈斯王.
就是說先知知道阿哈斯王.
在水溝的盡頭.
在大路上.
他在這裡做什麼呢.
應該就是當時耶路撒冷.
如果你說快要有人攻陷你的時候.
最重要的是做什麼.
就是Defense.
你的Defense在耶路撒冷最重要的就是你的水.
你的Water System.
這裡說到上池的水溝.
可能是耶路撒冷北面的一個水池.
就是當有人圍城.
你的City in siege的時候.
被人圍困的時候.
最重要的是什麼呢.
就是要有水.
所以皇帝和他的內閣.

$^{321}$就去到耶路撒冷不同的地方視察一下.
他們的防衛系統.
先搞國防.
這個時候就是National Defense的時候.
搞清楚水源夠不夠.
檢視一下的時候.
以賽亞就去那裡找他.
然後說你去找他.
對他說什麼呢.
他說你要謹慎.
因為這個皇帝來到的時候.
他的計謀是不能夠暫時的.
就是說耶和華一定會為你爭戰.
在這個危機裡面你不需要怕.
因為這兩個皇帝.
絕對不能夠成就他們的計謀.
亞蘭的首都就是大馬士革.
以法尼就是以色列.
必定會被毀滅.
你們若不相信.
就必定不得建立.
先知做什麼呢.
不是跟他說明天天氣的預測.
而是幫他去管理他的crisis.
在這個危機裡面.
以賽亞就是跟阿克斯王說.
搞清楚.
你要用上帝的眼光.
去檢視一下你現在面對的危機.
不要用你自己的眼光.
用上帝的眼光來看一下上帝.
在你身上在做什麼.
在世界上在做什麼.
阿克斯在這個危機裡面.
就是怕到臉青青.
他怎麼回應.
第七章第十節.
耶和華有告訴阿克斯.
你向耶和華的神去求一個兆頭.
或顯在深處或顯在高處.

$^{361}$即是sky is the limit.
即是你儘管求一個兆頭.
儘管求一個sign.
以致能夠confirm.
我在你中間在做什麼.
我在你中間做的救贖.
阿克斯王你儘管求.
阿克斯王說什麼呢.
阿克斯王說我不求.
我不試探耶和華.
大家記得很奇怪.
這個是壞人的王.
他說了一句這樣的說話.
他說我不求任何兆頭.
因為我不會去test the lord.
我不會去試探耶和華.
好像是一句很敬畏神的話.
但是如果你看看其他的經文記載.
列王記或者歷代寺.
你就知道為什麼他不求兆頭.
為什麼他不想耶和華插手.
原因就是.
例如列王記下第十六章.
那裡說什麼呢.
一段很長的經文.
那裡說阿克斯已經猜派了使者.
去見亞述王.
你記得地圖上面兩個小的國家.
亞蘭和以色列圍攻他.
阿克斯王找到一條什麼辦法.
當時朝廷就決定.
我們找我們的敵人的敵人.
找亞述帝國來幫我們.
於是他就在聖殿和皇宮那裡.
挖了很多的金和銀.
就送了去亞述王那裡做禮物.
第九節.
於是亞述王就答應他.
我來幫你.
於是敵人的敵人.

$^{401}$亞述帝國就去攻打亞蘭國.
以致他們就退兵.
耶路撒冷就沒事.
你想想這個阿克斯王.
他寧願去求敵國.
都不願意去求耶和華.
事情的下文也都有交代.
這個以色列看到.
這個阿克斯王不願意去投靠耶和華.
十三節.
以賽亞就去警告大衛家.
警告亞述.
大衛家你們要聽.
你們使人厭煩還算小事.
還使我們的神耶和華厭煩嗎.
有沒有搞錯.
你求敵國都不去求問耶和華.
七章第十四節.
所以主都會給你一個兆頭.
給你一個sign.
這個sign是甚麼呢.
他說阿罕就是我們聖誕節讀的經文.
必有同女懷人生子.
他要給他起名叫以瑪萊尼.
到他曉得棄惡擇善的時候.
choose good and evil的時候.
他這個小孩一定會吃乳酪和蜂蜜.
因為這個孩子曉得棄惡擇善之前.
他未懂得choose good and evil之前.
你所畏懼的兩個王.
就是亞蘭王和以色列王.
他們的地方就會被棄住.
就是說他們一定會失敗.
這句話很多時候我們只cut了十四節出來.
放在聖誕卡上面.
但上文下理你見到原來是有其他的內容.
但這裡幾乎每一句都需要解釋.
第一同女懷人生子.
同女在希伯來的原文.
是年輕女的意思.

$^{441}$當時十三,十四歲就會結婚.
十三,十四歲就是young woman.
marriable age.
一個年輕剛出嫁或者適合出嫁的女子.
這個是同女.
所以不是我們想同證女瑪麗亞.
首先當時的時候.
是說一個年輕女.
必定會懷孕生子.
或者這麼說.
原文其實不是將來懷孕生子.
其實是現在已經懷了人.
有一個女孩子已經懷了人.
有一個女子已經有了新人.
這個女子是誰呢.
這裡沒有交代清楚.
大概應該是阿哈斯知道是誰.
可能是她後宮裡面的其中一些妃嬪.
總之有一個女子懷了人.
已經可能是見肚了.
阿哈斯是知道的.
然後她生出來的這個小孩.
她會叫她的名字叫做「爾瑪萊蒂」.
「爾瑪萊蒂」的意思就是說.
God with us.
三個字組在一起.
Emmanuel.
三個字組在一起.
你說好像很神聖.
因為是「爾瑪萊蒂」.
但其實聖經裡面的舊約.
那些人的名字.
很多時候都有神的名字.
譬如說「伊西喬」.
就是God strength.
就是上帝加力給他.
「伊西喬」最後是EL.
所以好像Emmanuel那樣有個EL.
或者「丹尼爾」.
就是God judges.

$^{481}$上帝會以公義來審判的.
這個就叫做「丹尼爾」.
最後有EL.
所以無端端有EL.
或者有神的名字.
不代表他一定是神聖的.
這個小孩其實可能是一個普通的人.
「爾瑪萊蒂」.
第四樣東西.
就是這個小孩懂得去選擇好與壞.
棄惡損善.
在猶太人傳統.
大概十二,三歲的時候.
Barmitzvah的時候.
就會叫他做已經去到一個.
Age of Accountability.
就是說這個人已經成為一個成人.
十三歲就算是一個Adult.
Adult就可以.
就要按照上帝的律法.
在上帝的底下.
來去棄惡損善.
去選擇好與壞.
換句話說這裡說什麼呢.
就是說這個小孩還沒懂得選擇好與壞之前.
這兩個國家都不會再威脅到你.
換句話說很快這兩個國家就會敗亡.
你記不記得那個意思.
有個女子已經有了新人.
而且她生出一個孩子.
這個孩子將來會叫他名叫做「爾瑪萊蒂」.
因為上帝與他同在.
就是這個意思.
但是這個孩子還沒到十二,三歲之前.
很快就.
這兩個最後的字.
這兩個王的地就會被棄住.
必有乳酪和蜂蜜.
這兩種的食物.
就是指當時國泰民安的時候吃的東西.

$^{521}$因為如果打仗走難的話.
你就不會有乳酪吃.
也都不會有蜂蜜吃.
簡單來說就是說.
已經有一個兆頭給了阿哈斯.
就是有一個baby.
在媽媽的肚子裡面.
阿哈斯你們知道的.
告訴你這個baby.
還沒到十二,三歲之前.
這兩個國家都一定會敗亡.
軍事的威脅就會完結了.
所以這個孩子.
真的可以叫他做「爾瑪萊蒂」.
就是上帝居然與我們同在.
去保住我們.
整段的兆頭.
所謂「同女懷孕生子」.
意思就是說不出幾年的話.
不需要幾年.
很快就出.
上帝就會為你解困.
那個神蹟在哪裡.
你看清楚這裡.
上下文看完的時候.
背景看完.
這個神蹟就不是有「同女懷孕生子」.
平常的事.
神蹟在哪裡.
神蹟就是上帝願意去拯救不信的大衛家.
即是大衛家不願意去投靠耶和華.
就算在危機當中.
都不願意去投靠.
但是神蹟出來.
上帝仍然去到保守著友大國.
和大衛家.
上帝恩典大過人的不淑.
即是阿哈斯多壞都好.
上帝恩典居然願意與他同在.
對於面對著危機的人.

$^{561}$有甚麼提醒.
對阿哈斯來說.
就是你面對困境的時侯.
你就有兩個選擇.
第一 就是好好地轉向神.
你多不事多不順都好.
上帝給你一個機會.
你要回到上帝的恩典.
第二個選擇就是你繼續去背若.
這個就是阿哈斯的選擇.
選擇繼續遠離神.
選擇繼續在不信裡.
你問為何耶穌出生的時候.
馬太科罕第一章.
明明說到這段聖經應驗.
不是一早寫了.
好像天理報告那樣.
幾百年前就說.
將來會有個同女懷孕生子.
然後叫他做二馬來尼.
這個是下一個星期要解決的問題.
就是馬太科罕怎樣去看這段聖經.
現在我們看回這段聖經.
當時的背景時.
我們發現原來說的東西.
有另一個意思.
另一個意思就是.
這個二馬來尼的小孩出生.
原來是一個兆頭.
告訴阿哈斯.
告訴不信處於困境.
不肯轉向耶和華的人.
告訴他們.
不出幾年上帝的恩典.
就會來拯救你.
請你轉向神.
可能我們看這段聖經看得太多.
我們看得太熟了.
我們就很被耶穌出生.
聖誕卡上面的情況太過吸引.

$^{601}$我們就經常覺得.
這裡一定是說同女懷孕生子等等的事.
其實如果你看看經文.
從頭到尾想說些什麼.
其實重點就不是同女懷孕生子.
重點就是你願不願意轉向神.
在你困境裡面.
在你面對危機的時候.
你願不願意趁這個機會.
來轉向神.
你看到這裡的脈絡.
其實都很清楚.
頭兩節 七章一至二節.
就是猶大國民受威脅.
亞蘭和以色列國.
來圍剿耶路撒冷.
然後對下面結束的時候.
他既然不肯相信的時候.
猶大國民將會受威脅.
就是亞述真的會.
不單止吞滅了以色列國.
猶大國都會受到威脅.
一對的 A和A.
然後B就是有個兒子出現.
你有沒有留意.
七章裡面不只是伊瑪利亞的小孩.
就是以賽亞的兒子.
都帶上去的.
為什麼要帶上去.
其實就是想配合B下面的.
上面以賽亞帶著他的兒子.
斯亞瓦述.
就去講出這個敵國.
北國以色列 亞蘭國.
他們的聯合會一定會失敗的.
他們的計謀.
然後告訴你下面七章第十六節.
就是這兩個敵國因為耶和華恩典的拯救.
就一定會敗亡.
然後再去到中間.

$^{641}$七章第十九節.
就是上帝拯救的應許.
然後給他一個笑頭.
七章第十三至十五節.
那個就是以瑪利亞的經文.
重點不是那個同女懷孕三節.
重點是中間阿斯不願意相信.
七章第十節到十二節.
這個才是重點.
在危機裡面.
能不能夠拿到一個教訓的時刻.
明不明白上帝要我們去轉向他.
去投向他.
還是只顧著自己.
我有我自己的答案.
有我自己面對困境的方法.
你不知道這個困境來到.
其實上帝最關心的是你的信念.
是你跟他的關係.
疫情到現在.
多年時間.
將會踏入第二年.
大家都不想.
在這個困境好像一個骨頭老一樣.
但我想過聖經裡面給我們一個角度.
一個perspective.
就是在這個crisis裡面.
上帝所concern的是什麼.
不是危機退去.
而是危機裡面的你.
有沒有轉向他.
還是我們這個困境在疫情裡面.
我們只不過是等疫情過去就好了.
我們依然顧我.
我們沒有改變過.
我們跟上帝的關係沒有深過.
沒有任何進展.
好像亞克斯一樣.
疫情的你 我想問.
你跟上帝的關係是怎樣.

$^{681}$現在 將來.
這是關鍵.
上帝在crisis裡面有的teaching moment.
我剛才開始的時候都講過.
疫情不單止快要有兩年.
更加你要知道.
亦都會繼續反復.
當我們 特別在美國北美.
我們有疫苗.
好像有些曙光.
但是都還沒有見到出路.
現在有Omicron.
希臘文的字母只有24個.
Omicron第15個.
跳了幾個 跳了C 跳了什麼原因都好.
還可以有多幾個 不知道.
但是如果你看這段所謂聖誕經文.
在疫情裡面.
上帝容許這件危機發生.
我們怎樣去manage.
Manage不是去解困這麼簡單.
而是上帝看見危機中的你.
你可以轉向它.
疫情之前 你跟上帝的關係.
無論是怎樣都好.
在疫情當中.
有沒有令到你得到試煉.
有沒有那種teaching moment.
還是我們只是自己過自己的生活.
不願意上帝去參與.
我們中文經常說.
Crisis這個字叫危機.
有危亦有機.
有威脅 亦都有生機.
阿克斯王就是在這個危機裡面.
沒有轉向上帝.
沒有好好manage他的crisis.
我不知道你是怎樣.
他沒有好好manage他的crisis.
結果又大棍.

$^{721}$最終都是滅亡.
重點在哪裡呢.
七章第九節.
我想這個都算是一個重點.
七章第九節.
他說 你們若不相信.
就必定不得建立.
如果你是faithless的話.
你就不會捱得很久.
我們的faith是怎樣.
我們同心低頭 祈禱 預備我們自己.
上帝 我們多謝你.
因為你容許我們心靈當中.
遇見不同的困境 不同的問題.
求你親自去挑戰 去提醒 幫助我們.
你是那位以馬來尼.
你是那位與我們同在的神.
最重要的.
不是我們危機 疫情得到解決.
而是在危機當中.
我們會否轉向你.
全心全意地去投靠你.
求主你憐憫.
憐憫我們很多時不順.
我們對你的信心是很軟弱.
求你幫助你自己的子民.
轉向你.
讓我們以信心去投靠.
知道你是那位比我們任何面對的困境.
更大的神.
讓我們更加深來認識你.
幫助我們 經歷這些疫情.
經歷我們心靈裡不同的困境的時候.
我們經歷了之後.
我們與你之間的關係.
更加清楚 更加清晰.
更加深來認識主 你自己的大能.
祈禱 奉主命救.
今天是我們聖餐的主日.
用這個恩典 拿起這個餅的杯.

$^{761}$都可以提醒我們.
上帝給我們的恩典.
即使我們很多時不順.
但是上帝的恩典.
祂的請求沒有改變過.
在座我們有一個聖餐的套餐給你.
你拿了之後.
我們等下一起同心領受.
網上的大英姐妹.
跟我一起參加這個崇拜.
如果你的生命.
是與上帝搞得很清楚.
耶穌基督是你的救主的時候.
請你和我們一起在家中.
同心領受.
領受的時候讓我讀出一段聖經.
預備我們的心.
哥林多前書第十一章.
二十三節開始.
我當我傳給你們的.
原是從主領受的.
就是主耶穌被賣的那一年.
他拿起餅來祝謝之後就張開.
說這是我的身體為你們張開的.
你們應當這樣行.
為的是紀念我.
飯後照樣拿起杯來.
這杯是用我血所納的新約.
你們每逢喝的時候.
應當這樣行.
為的是紀念我.
你們每逢吃這餅喝這杯.
就是宣揚主的死情.
等到他來.
你們應當如此行.
為的是紀念基督.
不是單單拿起餅和杯.
當如此行.
就好像基督所做的一樣.
為了教會.

$^{801}$為了他愛的我們每一個.
去寫記.
應當如此行.
是向耶穌基督的榜樣.
這個餅和杯就是紀念到.
主對我們的恩典.
讓這個恩典幫助我們.
讓我們的心態.
我們的生活.
好像基督一樣.
我們同心靈受主的餅.
我們同領主的杯.
我們就在這節帶來祈禱.
主讓我們今天.
在困境 在不同的危機當中.
來投靠你.
讓我們更深認識你.
叫我們去愛你的話語.
愛去聽你的說話.
主幫助我們去管理我們的危機.
管理我們的困境的時候.
讓我們知道你和我們同在.
主讓我們轉向你.
讓我們不再在那種不順.
去拒絕你的那種淑女的狀態裡.
搞好我們的生命.
求主引導我們.
全心全意來愛你.
去體會你.
去認識你.
祈禱奉傳明教.
最後我們一首回應的詩歌.
叫做《投靠》.
希望在場的弟兄姊妹.
藉著這首詩歌我們有少許時間.
去到蜜桃.
然後我們才去散會.
家中的弟兄姊妹同樣.
求上帝賜給大家平安.
\newpage



\section{1:18-25}
\label{sec:e3EJDUy_J3k}
\textbf{The birth of Jesus Christ (Matthew 1\_18-25) - Rev. Mark Chan}
\newline
\newline
連結: \href{https://youtube.com/watch?v=e3EJDUy-J3k}{\texttt{https://youtube.com/watch?v=e3EJDUy-J3k}} ~~~~ 語音日期: 2021-12-13
\newline
\newline
\hyperref[sec:CNrSX78PCg8]{\small{< < < PREV SERMON < < <}}
~
\hyperref[sec:index_chronic]{\small{[返順時目]}}
~
\hyperref[sec:index_scriptual]{\small{[返順卷目]}}
~
\hyperref[sec:lTIqtT1KU5c]{\small{> > > NEXT SERMON > > >}}
\newline
\newline
$^{1}$And let's turn to the word..
One question I wanna ask you is that.
where do you find in the Bible.
the story of the birth of Jesus Christ?.
All right, if you can talk to your neighbor,.
where do you find it, the first Christmas story,.
all right, so-called first Christmas story,.
where do you find in the Bible?.
Where?.
All right, so Old Testament or New Testament?.
New Testament, right, which book?.
Luke, okay, right, the third book of the New Testament,.
and then the first book, right, the Gospel of Matthew..
Basically, you find the stories of the record,.
the historical record of the first Christmas,.
the birth of Jesus Christ, the details over there..
So let's turn to the first book of the New Testament,.
the Gospel of Matthew, chapter one, verse 18..
Here we go..
This is how the birth of Jesus Christ came about..
His mother, Mary, was pledged to be married to Joseph,.
but before they came together,.
she was found to be with child through the Holy Spirit..
19, because Joseph, her husband, was a righteous man.
and did not want to expose her to public disgrace,.
he had in mind to divorce her quietly..
Now, everybody loves the Christmas story, right?.
Everybody loves the Christmas story..
A young Jewish couple, about 14 years old, 15 years old,.
back then, people marry very early in their lives,.
and in Jewish marriage, there were two stages of marriage..
The first stage, the first period, is about a year long,.
is called betrothal or engagement..
It's about one year long,.
and in the beginning of that one year,.
when a young couple get engaged,.
usually the parents arrange them..
So they arrange them, and then they talk about.
the bride price and all that details,.
and then the young couple will be,.

$^{41}$technically, legally, they will be married,.
and we call them betrothed or engaged,.
and take about a year to prepare for the real wedding,.
which is the second stage,.
and the couple will consummate their marriage relationship..
So it is at this point in the first stage.
of betrothal and engagement.
that we find the first Christmas story..
So Mary was found to be with a child..
She was pregnant..
And then for Joseph, and you think about it,.
it's quite a tension, it's quite a tension,.
because all you plan for, all you arrange ahead of time,.
didn't work out, right?.
Mary was found to be with a child..
And for Joseph, for the husband,.
in this stage, in the first stage, the betrothal,.
Mary and Joseph, even though they are so-called engaged,.
but they are officially and legally,.
they are husband and wife..
So at this point, Joseph had two options, all right?.
One option is to divorce Mary properly,.
grab her to the city gate, to the court of the day,.
and divorce her and just shame her,.
because it's presumed that Mary is with child.
before they come together, before the second stage..
So it's presumed that Mary would have committed adultery..
And in the traditional society,.
and actually in some areas in the land of Palestine today,.
it's so conservative that whoever found.
to be committing adultery would be stoned to death..
So that's the first option for Joseph..
The second option for Joseph is to divorce Mary privately..
Probably would ask two or three witnesses.
and come together in a private setting to divorce her,.
give her a divorce certificate..
I think Joseph still had compassion toward Mary..
At this point, I think he was trying to keep it down.
to divorce Mary quietly, to save Mary from public disgrace..
So here's what happened..

$^{81}$And Joseph, if you think about from his perspective,.
would he believe that Mary had a child from the Holy Spirit?.
Would he believe that it was really a miracle from God,.
because we got another version,.
Mary's version from the Gospel of Luke, right?.
At this point, Joseph knew nothing about the annunciation,.
about the news that the angel brought to Mary..
So Joseph knew nothing about it,.
and it was pretty difficult for him to accept the reality..
But this is where God came with his divine intervention..
He sent an angel in chapter one, verse 20 and following..
But after Joseph had conceded all of this,.
an angel of the Lord appeared to him in a dream and said,.
Joseph, son of David, do not be afraid.
to take Mary home as a wife,.
because what is conceived in her is from the Holy Spirit..
It's a virgin birth, it's a miracle, only God could do..
She will give birth to a son,.
and you are to give him the name Jesus..
And the name Jesus means he will save,.
he will redeem his people from their sins..
Now with that, Joseph obeyed God.
in kind of a difficult situation..
Even though in the midst of uncertainty,.
you just don't know what happened next, right?.
So Joseph was kind of an example here,.
example of obedience and compassion..
So he took Mary as wife,.
and essentially he was adopting Jesus to be his legal son..
Now remember, Joseph is from the house of David..
He's one of the descendants of the Davidic dynasty,.
the house of David..
So in adopting Jesus to be his legal son,.
Jesus would be son of David as well..
Now why did all this happen?.
Matthew continued, 1:22..
All this took place is to fulfill the will of God..
What the Lord has said through the prophet..
Here Isaiah 7..
The virgin will be with child,.

$^{121}$and will give birth to a son,.
and they will call his name Emmanuel..
Well, we sang the song,.
which means in Hebrew, God with us..
Now why did this happen?.
Matthew is quoting a prophecy.
spoken 700 years before the time of Christ..
It was back in the time of the Davidic dynasty,.
and it was about 700 BC, before Christ..
It was a wonderful prediction, right?.
700 years before,.
already prophesied the birth of Jesus Christ..
Now that is only half of the story,.
because as we know, the Bible is not written to us, right?.
The Bible is not written directly to us,.
but the Bible is written for us..
Now with that said,.
the Bible always has a historical background..
There are tons of history,.
tons of background behind the Bible..
So we have to understand why Matthew.
is quoting this prophecy from Isaiah 7..
So Isaiah 7, let's turn to that..
What happened at that time was that in about 700 BC,.
there was,.
the dynasty of David was already divided.
into North and South..
The Northern Kingdom is called Israel..
The Southern Kingdom is called Judah..
Now what happened in 700 BC was that.
the Northern Kingdom, Israel,.
and other pagan nation called Aram,.
they decided to form a sort of a military coalition.
against Judah, against Judah in the South..
Now Judah was a much smaller nation,.
and didn't stand a chance..
So back then, the king of Judah called Ahaz,.
he was freaking out..
He was freaking out because these two pagan nations.
are invading Jerusalem, invading the nation of Judah..

$^{161}$So it was about 700 BCE..
This little country was under military attack.
by two coalition, Israel and Aram..
So what happened was that when the king was freaking out,.
did not believe in God, did not trust in God,.
then the Lord told Isaiah to prophesize to him,.
to speak to King Ahaz..
Verse 10, the Lord spoke to King Ahaz,.
and he said, "Ask the Lord your God for a sign,.
"whether it's the deepest depth or the highest heights.".
So the sky is the limit, you can ask anything..
You can ask a sign from God that God will confirm.
that he is with you..
God is protecting you, and you can trust him..
But King Ahaz said, "What?.
"I will not ask, I will not put the Lord to the test.".
Now that may sound really godly, right?.
I don't wanna test God, I don't wanna tempt God, right?.
But in fact, we know King Ahaz was an evil king..
He did not trust in Yahweh, he did not worship God..
He has actually did something to save himself..
So he scraped the gold and silver.
from the temple in Jerusalem,.
and gave it to another bigger king,.
which is Assyria up on the upper right..
Asked Assyria to help him and save himself..
So he has already taken money out of the temple treasury.
to bribe another enemy nation, Assyria, to help him out..
So he did not trust God,.
and hiring the Assyrians to help him out..
Now hiring the Assyrians to help him out.
is like taking money out of the offering box, right?.
Taking money out of the offering box.
in order to rob the church..
So Ahaz was a really bad king,.
he did not want to tell the prophet..
Now, but we're at this point, Isaiah continue his word..
Essentially he's saying, even though you don't trust me,.
you don't trust God, how so, David?.
Is it not enough to try the patient of man?.

$^{201}$Will you try the patient of my God also?.
Therefore, the Lord will still give you a sign..
The Lord will still give you a sign,.
some sort of evidence to tell you that he is with you,.
even though you are faithless, but God is faithful..
And that sign is the virgin will be with child.
and will give you a son..
And you will call his name, Emmanuel, God of us..
Verse 16, and this child will be born,.
but before the boy, this newborn child,.
knows enough to reject the wrong and choose the right,.
the land of the two kings, remember Israel and Aram, right?.
It's trying to attack Israel, trying to attack Judah..
These two kings will be laid waste..
Now, I highlight here, verse 16,.
before the boy knows enough to reject the wrong.
and choose the right..
It means when a child is born, continue to grow,.
and at some point, he will know what is right and wrong,.
right, what is black and white..
Now, how old is the boy, how old is he,.
going to tell the difference between right and wrong?.
Not very old, right?.
Probably in a couple of years, right, two years old..
Know enough to tell what is right and what is wrong..
So essentially, the prophecy of Emmanuel,.
in the context, in the background of Isaiah 7, 700 BCE,.
is saying in a couple of years,.
God will save you from trouble..
God will save you from the attacks.
of the two northern kingdoms..
God will save you, and you will,.
you have to put your trust in him..
So God is saying, this boy will be born..
This prophecy will come true in your day..
In a couple of years, God will save you..
So you have to put your trust in him..
This is a test of faith..
But King Ahaz, as we know, he failed the test, right?.
In the midst of difficulties,.

$^{241}$in the midst of difficult circumstances,.
he chose to trust himself, right?.
Hire the Assyrians to help him out..
He failed the test..
But God giving this sign, this newborn baby,.
is to tell the house of David, he will be with them,.
even though the circumstances look very difficult,.
very challenging..
And then we fast forward 700 years later,.
in the first century, when Christ was born,.
Matthew was telling us, this prophecy fulfilled again.
in the time of Jesus' birth..
Now you're saying, okay, let me ask the question..
So were there two virgin births?.
One 700 before Christ, and then one in the first century.
when Christ was born?.
No, no, there were not two virgin births..
All right, because when Isaiah talks about this prophecy,.
he used the word for virgin, English is translated virgin,.
is actually young woman only..
So virginity is not required, it's only implied, all right?.
So the word is amar in Hebrew..
Actually in the Hebrew dictionary,.
there's another word, technical word for virgin,.
which is betulah..
Isaiah 7 did not use that..
So it's vague enough, it is broad enough,.
to say that in times of Ahaz, in times of 700 BCE,.
the young woman, not virgin, will bear a child..
And then fast forward 700 years,.
in the time of the birth of Jesus,.
the final fulfillment of this prophecy came true..
You got my point?.
700 years before, there was an initial fulfillment.
in times of King Ahaz..
Now in the first century, finally we see.
the ultimate fulfillment of the word of God..
So the promise in the Old Testament,.
always have, or often has, multiple fulfillments..
When God says something, sometimes it was initially fulfilled.

$^{281}$in his own time, but then there are multiple fulfillments.
until the time of Christ fulfilled all of them..
For example, the first temple built.
by the son of David, Solomon, right?.
The first temple was prophesied to be built,.
and it was indeed built by King Solomon, the first temple..
And then several hundred years later,.
the temple was destroyed by the Babylonians..
The people were in exile..
The Jewish people were in exile in Babylon..
And then 70 years later, they returned.
to the land of Palestine, they rebuilt their own home,.
and they rebuilt a second temple, another fulfillment..
And then in the time of Christ, Christ is saying,.
he's responding to the Jerusalem temple,.
built by stone, built by hands..
He said, "If you destroy this temple,.
"I will rebuild it in three days.".
Now what Jesus was saying was that,.
"I will rebuild my own body..
"I will resurrect from the dead,.
"and my body will be the new temple.".
That is the ultimate fulfillment..
Now what I'm saying is, in the Old Testament,.
when you find some prophecies,.
it's not necessarily it has to be fulfilled.
many, many years later..
There are many different multiple fulfillments along the way..
And this is also true for Isaiah 7, the Emmanuel prophecy..
So when Matthew penned the gospel,.
he was checking many different accounts,.
many different stories, some he had heard before,.
because he followed Jesus for a few years..
He's one of the 12 disciples..
He knows something about the story of Christ..
He also gathered and collected different stories.
from different disciples of different traditions..
He put together the gospel of Matthew..
Now the question for you is,.
the very beginning I asked you,.

$^{321}$where in the New Testament you can find.
the story of first Christmas, right?.
One of them is gospel of Matthew, of course, right?.
Now, why was gospel of Matthew written?.
Why did Matthew decide to write a gospel of Matthew?.
Where you find the life of Christ,.
different story stitched together,.
come together as one book, gospel of Matthew..
Again, it's the historical background..
Because in the first century,.
when Matthew wrote the first gospel, the gospel of Matthew,.
the young church was in crisis..
It was under the Roman Empire..
It was being persecuted..
And the young church being persecuted,.
the people in the church, many of them are Jewish believers,.
they need to reconsolidate their faith..
They need to understand what they really believe in..
That's why Matthew decided to write the gospel of Matthew..
To tell the church what you should believe..
And to strengthen the faith of the church..
So Matthew is collecting different stories of gospel..
And then he traced the story and find a certain pattern..
For Jesus, when he looks at the birth of Jesus,.
the first Christmas, he found that Jesus.
is indeed the son of David..
Because Joseph is the son of David..
And that's why Matthew, in chapter one,.
if you turn to the very first chapter of the New Testament,.
it's not very appealing..
It's the genealogy of Jesus, right?.
So he traces from Abraham, Isaac, Jacob,.
generation after generation to King David,.
the establishment of the kingdom of God on earth, King David..
And then he traced after David, Solomon,.
and all the kings in the divided kingdom..
And then finally it comes to the Babylonian exile..
When everything seems to be over,.
when the nation of Judah was attacked.
and destroyed by the Babylonians,.

$^{361}$everything seems game over,.
he continued to trace the genealogy..
No, God is still faithful..
And generation after generation from exile,.
he traced another generation of Sababel,.
Jewish people who come back to Jerusalem.
and rebuild the temple..
And generation after generation doesn't end..
And finally it comes to Joseph and Mary..
And then Jesus..
That's why in the Gospel of Matthew, chapter one,.
after tracing the genealogy of Jesus Christ from Abraham,.
David to exile to Jesus Christ,.
verse 17, it says,.
there were 14 generations in all from Abraham to David..
And then from David to the exile to Babylon,.
there were 14 generations..
And then from exile to the coming of Christ,.
to Christ, Jesus, there were also 14 generations..
So he divided the entire Old Testament,.
all together about 39 books together..
Matthew divides the entire Old Testament history.
into three periods, 14 generations for each..
Now the question for you, why 14?.
Why 14?.
All right, why 14?.
Now in the ancient Hebrew culture,.
they didn't have the Arabic numerals,.
like what we do have, one, two, three, four, and so on..
Every letter of the 22 letters of Hebrew,.
each alphabet represents a numerical value..
Like the first one, aleph, is one..
The second one, bay, is two..
Gimel is three..
Dalet is four..
And then hey is five, and so on..
Val is the sixth letter, is number six..
So they used different letters to represent different numbers..
Now David, King David's name,.
it was originally written without the val..

$^{401}$So you take out the a and i, you've got dvd..
So David was originally written as dvd or dwd..
D is the fourth letter, and v, or w, is the sixth letter..
So you add them all up, and you get 14..
So 14 is the royal number for the King David's dynasty..
So when Matthew was tracing the Christmas story,.
he find that the long genealogy of Jesus Christ.
actually is 14 generations for three periods..
And of course, as the son of David,.
the real son of David, the final son of David,.
his number is also 14, right?.
14..
The birth of Jesus Christ is so much more.
than the miracle birth of virgin birth..
Okay, so note this..
Of course, Jesus was born as a virgin,.
the divinity of Christ..
He is the deity, he's 100\% God and 100\% human..
What a miracle..
It was indeed a miracle..
But the prophecy in Isaiah 7, the Emmanuel sign,.
is so much more than just the miraculous nature.
of the birth..
It is the timing of the birth..
700 years before Christ, in time of King Ahaz..
It was just at the right time, God save his people..
Here in the first century, God did the same thing..
Right, so Matthew is trying to rethink,.
trace the pattern in the first Christmas story..
And the story doesn't stop there,.
because in chapter two of the Gospel of Matthew,.
you remember the so-called three wise men, right?.
The three wise men from the east,.
they came to Jerusalem to worship the newborn king..
They came to the town Bethlehem..
Now Bethlehem is the hometown of David, right?.
Remember that?.
The hometown of David was Bethlehem,.
and Christ Jesus was born in Bethlehem..
So this, Herod the Great, at the time,.

$^{441}$it was Herod the Great, the evil king,.
he was an other pharaoh-like oppressor,.
trying to persecute the baby Jesus, right?.
The boy Jesus, trying to kill him..
So Joseph took his entire family, Jesus and Mary,.
became a refugee, and went down to Egypt by night..
So Matthew, thinking about this story,.
in chapter two, verse 15, why did this happen?.
Jesus became a refugee, went down to Egypt..
He said, "So it is, to fulfill what God has said.
"through the prophet.".
And he is quoting from Hosea, chapter 11..
It says, "Out of Egypt I call my son.".
Out of Egypt I call my son..
Now you check, you go back if you have time,.
go back to the book of Hosea in the Old Testament..
It was not about prediction..
It was not about predicting 700 years later,.
or so many years later,.
then Jesus will come out of Egypt..
The entire chapter is talking about.
God's commitment to his people..
When they were suffering in the land of slavery,.
in the land of bondage, to be slaves for Pharaoh..
In Hosea 11, God is saying, "I will not abandon you..
"I will be with you..
"I will rescue you out of Egypt..
"There will be an exodus.".
So Matthew is chasing the Christmas story..
He is seeing a pattern..
Jesus went down to Egypt, and he will come out of Egypt.
because God will not forget his people..
In Jesus, there will be the real exodus..
In the story of Jesus, he also went through exodus,.
just like the Old Testament people of Israel..
They went down to Egypt,.
and they came out of Egypt in the exodus..
And in Jesus, finally, we see the real final exodus..
And the story continues..
King Herod tried to kill the boys.

$^{481}$under two years old in Bethlehem..
As you can see, it's horrible..
It's a horrible story..
Under two years old, all the boys, he slaughtered them..
2.17, Matthew sees another pattern..
He said, "This happened is to fulfill.
"what was said by Prophet Jeremiah..
"A voice is heard in Ramah, weeping and mourning..
"Rachel, weeping for her children.
"and refusing to be comforted.
"because the children are no more.".
You see the pattern here..
He's quoting Jeremiah 31..
And Jeremiah was about 500 years before Christ..
What Jeremiah was talking about in chapter 31.
is that when the Babylonian slaughtered the Jewish people,.
killed them one by one,.
and captured them at a place called Ramah,.
and a place called Ramah,.
and Ramah happened to be just close to,.
right next to Bethlehem..
It was tragic for the nation.
because the nation was finally, seems like game over..
And all the men of Judah was being captured at Ramah,.
and they marched from Ramah all the way out to Babylon..
It was in the time of exile..
It was in the time of weeping and loss..
Lost the men, the children of Judah..
Jeremiah 31..
You keep on seeing the context there..
If you check what Jeremiah 31 has to say next,.
it's not just about the exile..
Because 31, Jeremiah 31,.
is about the return, the hope of return..
Even though the people of God is in exile,.
God will wipe away their tears..
God will comfort them again,.
will bring them back to the land,.
will give them back to homeland..
You see the pattern here as Matthew sees it..

$^{521}$In Jesus, we see the real return,.
the true return to homeland..
The true home is found in Jesus Christ..
He will wipe away our tears of exile..
Early on, we were praying for our community..
Tragic news that we receive..
So in the middle of last week,.
one of my friends called me..
He's a very close friend of the victim's family..
He called me and said, "Hey, Mark,.
"so-and-so just died from a gunshot.".
I could not believe that..
What a tragic news..
And other one hitting our community..
So in times like this,.
we always ask ourselves, where's our home?.
Where do we go that we have real security,.
that we'll have real hope?.
We'll find this place to be really like our home..
We fear our insecurity..
We fear about the craziness going on in our world..
We don't feel like home..
That is what happened to the Jewish people..
Yes, they returned from Babylonian exile,.
but empire after empire, they're persecuting them..
They were living, we were in a very, very hard life..
They were still living in exile..
Matthew was saying in the Christmas story,.
if you hope for a real return,.
it's not a political return..
It's not a historical return..
It's a spiritual return..
You return to Jesus Christ,.
you will find that your tears, your mourning be over.
because the eternal hope that we have in Jesus Christ..
Christmas story continues..
After King Herod died, Jesus and his family,.
Joseph took his whole family out of Egypt.
and back to the land of Israel..
And eventually they settled down.

$^{561}$in a place called Nazareth in Galilee..
It's a very small town..
And Matthew concludes the entire Christmas story.
with this saying..
Having been warned in dream by the angel,.
Joseph withdrew to the district of Galilee.
and he went and lived in a town,.
settled down there called Nazareth..
So was fulfilled what was said through the prophets..
He will be called a Nazarene..
Now, the prophets talk about Nazarene..
Actually, no..
If you check the entire Old Testament, 39 books,.
not one place mentioned the town of Nazareth..
But Matthew is saying, hey, the prophets are actually saying.
that the Messiah will be a Nazarene..
Now, here we need to understand how the Jewish people,.
especially the rabbis, how they read scripture..
Sometimes they will pick different words.
and connect different dots..
Some of the words, they sound alike, sound very similar..
They will connect them together..
And what we find in Isaiah 11,.
in fact, there is a prophecy,.
there is a promise to the house of David.
that sounded like Nazareth, 11.1..
It says a soot will come up from the stump of Jesse..
Okay, the stump of Jesse..
The tree is being chopped down..
It's like the nation of Judah being destroyed, right?.
It's like a dead tree..
The stump of Jesse..
Why Jesse?.
Jesse is the father of David, okay?.
It's the dynasty, the forefather of the dynasty, the nation..
The stump of Jesse is being chopped down..
However, there is a branch grow from the stump..
And the branch in Hebrew is Nasi..
Nasi will grow and the spirit of the Lord.
will be with his servant.

$^{601}$and the servant will be the redeemer, the Messiah..
Nasi will grow out of death, out of the ashes..
God is still there with them, with the people of God..
And the dynasty of David, the house of David,.
will not be forgotten..
So actually, this is the prophecy Matthew is referring to..
A Nasi will grow up..
The Nasi will grow up out of the ashes, out of darkness..
The nation of Judah, the people of God,.
the kingdom of God will come..
This is actually how the Jewish rabbi.
would like to read scripture..
This is one of their techniques..
It's called mitrash..
Mitrash means seeking..
In Hebrew, it means to stop and to ponder..
That's what mitrash means..
It's a technique of reading and interpreting scripture..
Connecting the dots..
You see the pattern here?.
This whole story of Israel in the Old Testament.
found in the story of the first Christmas..
Find in the life of Christ..
When you look at the life of Jesus,.
you look at the entire Old Testament story,.
and you see the faithfulness of God..
Mitrash, to ponder, to stop, and to ask the question,.
what child is this?.
What is this?.
Who is this Christmas child that we are celebrating,.
that we are worshiping?.
What does this child bring you, bring to your life?.
You see the parent, Matthew C., the son of David,.
continues until ultimately fulfilled.
in the life of Jesus Christ..
You see Exodus, Moses come out of Egypt,.
having freedom, true freedom..
You find the ultimate fulfillment in Christ..
Exile, the tears of exile will be wiped out in Jesus Christ..
We turn from exile..

$^{641}$The true home that we find is in Jesus and no other..
Our true home is found in Christ..
Look around the world, as I said..
All this unrest and all this crisis and pandemic continues..
Sometimes in situation like this,.
we always say to ourselves or say to each other,.
I wanna move somewhere else, right?.
Have you ever thought about it?.
Move out of Chicago, then you will be safe, right?.
Move out of Chicago, move somewhere else,.
somewhere beautiful, right?.
Somewhere secure, and then you will be safe..
The Bible is telling us the true home is found.
no other than in Jesus..
When you have Jesus Christ in your life,.
then you will have eternity..
That's what Mitrash is talking about..
That's what Mitrash is saying..
To stop and to ponder, who is this baby Jesus.
that we worship every Christmas?.
If you choose to trust in Jesus Christ,.
you will find the true security..
Don't look anywhere else, look up..
The child born in Christmas, he gives us the true freedom,.
the true Exodus, the true return, the true home..
Lastly, as we close, I wanna go back.
to the very first sentence of the New Testament..
We often miss this first verse..
A record of genealogy of Jesus Christ,.
the son of David, the son of Abraham..
In Greek, the original language of the New Testament,.
the record of genealogy is Bibliogenesis..
Now, if you know a little bit about English,.
Bibliogenesis, Bibliogenesis,.
you find words that are sounding so similar..
The book of Genesis, the first book of the New Testament,.
the first book of the entire Bible is the book of Genesis,.
the book of new beginning,.
because God created heaven and the earth..
He created the entire cosmos..

$^{681}$That is the original creation..
But then the first sentence of the New Testament,.
it says Bibliogenesis, the book of Genesis,.
the real Genesis, the new creation,.
found in Jesus Christ..
When we accept Jesus Christ as our Lord,.
when we follow Jesus Christ, we have new beginning,.
and this is the only true new beginning we need..
Now, who is Jesus to you?.
Who is Christ to you?.
Stop and ponder..
Put your trust in him..
Put our trust in Christ that we'll find new beginning.
to live as the people of God, right?.
To be his sons and daughters,.
to have hope and security wherever we are,.
no matter how dark the world out there,.
we still have light because Jesus Christ is the true light..
Let's all pray together..
Oh God, we give you thanks.
for the fact that you are our true light..
And in this first Christmas story,.
we see the hope that you offer to all of us.
in the name of Jesus Christ..
Thank you for sending your son to be born in this world,.
to bear our suffering, to bear our struggles,.
to wipe out our tears, and to grant us true freedom,.
the true exodus, true return, true home..
And I pray that all of us will put our trust anew in you..
We'll again choose to trust in you,.
again choose to believe in you,.
and help each one of us to be your true son and daughter,.
to live by your word, to live by your scripture,.
to follow and obey your will..
In this Christmas, we pray that you will continue.
to renew our hearts, renew our trust.
in our Lord Jesus Christ,.
and he will continue to lead our lives..
In Christ's name we pray, amen..
Well, God bless..

\newpage



\section{1:18-25}
\label{sec:lTIqtT1KU5c}
\textbf{耶穌的誕生 (太 1\_18-25) - 陳崇基牧師}
\newline
\newline
連結: \href{https://youtube.com/watch?v=lTIqtT1KU5c}{\texttt{https://youtube.com/watch?v=lTIqtT1KU5c}} ~~~~ 語音日期: 2021-12-13
\newline
\newline
\hyperref[sec:e3EJDUy_J3k]{\small{< < < PREV SERMON < < <}}
~
\hyperref[sec:index_chronic]{\small{[返順時目]}}
~
\hyperref[sec:index_scriptual]{\small{[返順卷目]}}
~
\hyperref[sec:PUn4RKNOgHU]{\small{> > > NEXT SERMON > > >}}
\newline
\newline
$^{1}$下星期已經是聖誕的主日.
雖然這兩年聖誕節快樂的時候比較難出口.
因為都很難快樂.
雖然我們在中間 特別是實體的.
我們都是沒有穿沒有褥.
可以坐在這裡.
這個都是上帝給我們的恩典.
上星期五晚.
在我們唐人街的社區那裡.
有一個燭光的儀式.
是一個紀念.
不知道大家有沒有去 有誰有去.
如果沒有去 請你去為家庭禱告.
為我們城市裡面的不安.
還有去搜望一下我們身邊的鄰舍.
特別是一些比較軟弱的朋友.
特別留意他們.
不要躲在我們家裡.
同時我們出入都要小心.
另外一些在芝加哥以外.
我不知道你的情況是怎樣.
但願上帝賜給大家平安.
馬達科罕第一章第十八節.
我們今天看到這段經文.
上帝的話語裡面記載.
所謂第一個的聖誕.
第一章十八節.
耶穌基督的降生.
這個故事是這樣的.
耶穌的母親瑪利亞許配了約瑟.
他們還沒有成親.
瑪利亞就從聖靈懷了恩.
她丈夫約瑟視過異人.
不願以張揚視他羞辱.
就打算暗中與他解除婚約.
當時猶太人的婚姻有兩個階段.
第一個階段就是訂婚.
所謂engagement.
第二個階段就是成婚.
就是婚禮的時間.

$^{41}$中間大概一年多時間.
首先雙方家長會談婚論嫁.
談完之後就會去到訂婚的階段.
當時的社會.
那些人是非常早婚的.
十四十五歲.
我們今天teenager的時候.
他們就會談婚論嫁.
就會建立他們的家室.
因為當時的人.
大概他們的壽數.
只不過都是四十歲左右.
所以十幾歲的時候.
就會談婚論嫁.
所以這裡說到瑪麗亞和約瑟的故事.
大概我們可以想像.
起碼瑪麗亞.
雖然後期有些傳統說.
約瑟是比較老的.
但起碼瑪麗亞是當時.
一個年輕的少女.
十四五歲.
開始要談婚論嫁的時候.
進入一個訂婚的階段.
當時的訂婚就不同.
我們今天的訂婚.
今天我們的訂婚可能說一下.
不一定要承諾.
當時已經不是.
已經在法律上是合法的.
所以要解除一個訂婚的話.
就等於離婚.
等於divorce.
所以在這個情況裡面.
你就可以想像.
瑪麗亞突然間發現自己有新人.
路加芳告訴我們.
天使米加勒向他報喜.
就是同貞女會懷孕生子.
這個是瑪麗亞那方面.

$^{81}$但約瑟那方面是怎樣呢.
對於約瑟來說.
就在兩難之間.
因為他未知情的時候.
他只能夠有兩個選擇.
兩個選擇都是離婚.
但選擇可以是公開的離婚.
這是第一個選擇.
公開的離婚.
就是說會給瑪麗亞帶來很多羞辱.
甚至當時一些很傳統的社會.
一些角落.
會將他處罰.
因為這樣瑪麗亞.
起碼會懷疑他有通姦.
然後被石頭打死.
所以這個時間是一個很大張力的時候.
第二個選擇是什麼.
都是離婚.
只不過是可以私下的離婚.
私下的意思就是說.
可能找兩三個見證人.
按照摩西律法.
去解除這個婚約.
就不會在城門口公開的離婚.
我相信約瑟仍然對瑪麗亞有多少信任.
多少熱情.
希望能夠不會令到他太過難受.
所以對於約瑟來說.
都是一個很大的張力.
因為約瑟未知道另外一半的故事.
天使只是先向瑪麗亞報喜.
瑪麗亞還用信心答上.
約瑟未知道的時候.
對人來說是很難接受的一個現實.
當一些人生裡面有的計劃.
訂婚 有婚嫁.
你已經計劃了一些東西.
突然間發生這樣的事.
是很難去接受一個這樣的情況.

$^{121}$所以為什麼在一章二十節的時候.
天使就向約瑟報告.
他一直在煩惱這件事的時候.
想著這些事.
講到約瑟.
主的使者就在夢中向他顯現.
大衛的子孫約瑟只管放膽將你的妻子.
記住已經是立法上的丈夫和妻子的關係.
將你的妻子瑪麗亞迎娶過來.
因為他懷的人是從聖靈而來的.
他必生一個兒子叫他喜明.
做耶穌 因為他會拯救自己的子民.
從罪惡裡面出來.
約瑟在這個艱難的時刻.
接到一個更加艱難的命令.
就是上帝跟他說.
你只管放膽迎娶瑪麗亞過來.
她的新人原來是童貞女.
聖靈一個神跡來去化身.
所以你將她迎娶過來的時候.
本質上就是說.
將這個嬰孩 耶穌.
將他修養過來.
legally將他adopt.
當他為你法律上合法的一個兒子.
所以看見約瑟之前掙扎的.
就是想暗中解除婚約.
有情 也有理.
他也踏上這一步.
將瑪麗亞迎娶過來.
於是乎 如果你這樣計算的話.
約瑟因為他是大衛家族裡面的後人.
耶穌被領養過來的時候.
耶穌就成為大衛之子.
為什麼這樣發生呢.
馬太就告訴我們.
原因就是整件事的發生.
上帝在背後策劃去錄orchestra.
就是應嚴主席先知.
這裡說到以賽亞所說的.

$^{161}$必有童女懷孕生子.
並且他的名叫以瑪麗亞.
以瑪麗亞翻譯過來就是神與我們同在.
上個星期我們已經解釋了.
當時以賽亞書第七章.
這個背景是什麼呢.
大概在主前730年左右的時候.
當時猶大國受到兩個另外一個大國的威脅.
於是以賽亞就將這個預言帶給當時的阿哈斯王.
一個壞人的王.
七章第十三節.
以賽亞書說什麼呢.
他說大衛的家因為他們不信.
所以令到上帝很煩擾.
但是上帝依然給你一個笑頭.
就是必有所謂童女懷孕生子.
他的名叫以瑪麗亞.
在這個孩子還未懂得去選擇善與惡之前.
所畏懼的兩個王兩個國家.
以色列和亞蘭國.
來影響著威脅著猶大國的時候.
早就被棄絕.
所以這裡說到一個嬰孩會出生.
叫做以瑪萊尼.
在一個很短的時間.
就是這個孩子未能夠選擇棄惡選善.
就是一個小孩由他出生到他長大.
大概多少歲他能夠分辨到.
到底什麼叫做好什麼叫做惡.
什麼叫做對和錯.
不是很久.
什麼叫做對和錯.
所以在一段很短的時間.
那段預言的意思就是說.
一段很短的時間這兩個的威脅.
就會被棄絕.
這兩個國家都不會再有.
歷史確是這樣發生.
在七八年BC的時候.
當時亞述帝國就去到侵佔以色列國和亞蘭國.

$^{201}$剩下的就是猶大國.
所以就是說.
上帝說可不可以信一信我.
讓我在困境裡面為你去解決.
然後他就給了你一個兆頭.
這個兆頭就是這個童女會懷孕生子.
所以這個是頭一半.
這個二創亞書第七章的背景.
到了馬達科那講到耶穌出生的時候.
就是快速地七八年.
然後又有另一步的所謂的實現.
又是一個年輕女孩懷孕生子.
那你就問.
那不就是有兩個童貞女懷孕生子嗎.
又不是.
為什麼呢.
因為希伯來文在二創亞書第七章.
他用了一個字.
其實是比較含糊一點的.
童女這個字其實不應該翻譯成童貞女.
因為那個字媽媽.
其實是年輕女孩的意思.
適婚年齡的年輕女孩.
在希伯來文裡面有一個字.
是真的講到童貞女的.
是Betula.
但是二創亞書卻沒有這樣用.
所以在阿哈斯的時候.
主前七八年的時候.
講出這個語言的時候.
可能是一個普通的女孩.
在宮後去生一個兒子.
有不同的解釋.
有些人就認為可能是二創亞另外一個兒子.
就快出生了.
這個就作為一個意馬利的兆頭.
有人會認為是另外一個王.
可能是希西加.
也蠻合理的.
可能是另外一個兒子出生的時候.

$^{241}$這兩個北面威脅他們的國家就會淪陷.
可能是Sakai.
但無論如何那個原文其實不是說一定是童貞女.
所以最初的時候這個意馬利的兆頭.
這個語言.
就在阿哈斯的時候已經Initially Fulfilled.
就是說初步已經應驗了.
上帝告訴阿哈斯.
你看看這個小孩很快就出生了.
已經懷孕了.
當他出生的時候.
還不懂得去選擇Good and Evil的時候.
這兩個威脅你的國家就會玩完.
困難的時候.
你要全心來信我.
阿哈斯當然在信心上.
信心的考驗上.
他就失敗了.
往往我們就是這樣.
但上帝仍然是憐憫大衛家.
仍然憐憫當時大衛的皇朝有大國.
仍然去拯救.
仍然去維持著大衛家.
其實舊約裡面有不少的應許.
都是有Multiple Fulfillment.
舉個很簡單的例子.
上帝預言說會有第一個聖殿會興建.
是誰做的呢.
就是大衛之子所羅門.
所以有First Temple.
然後First Temple到586BC的時候.
又被拆毀了.
巴比倫帝國來到.
去殲滅耶路撒冷的時候.
然後他們又被迴歸回來的時候.
猶太人又去重建那個聖殿.
又另外一次的Fulfillment.
到耶穌來的時候.
你記得約翰福音第二章說什麼.
他說你們拆毀這個石造的殿.

$^{281}$70年你們才能建起來.
我第三天之後.
我就會將它重建.
其實他指著自己的身體.
所以一個Ultimate Fulfillment.
就是在耶穌的身體上.
就是他死而復活之後所誕生的.
他的國度.
他自己的子民.
他會再重建起來.
所以只是聖殿這個應許.
都有Multiple Fulfillment.
所以這個以馬來的兆頭.
有Multiple Fulfillment.
絕對不是一些很罕有的事.
其實每一個大衛的皇駕崩.
然後有另外一個大衛的子孫.
再上去登基的時候.
都是另外一個Fulfillment.
所以每一次大衛的皇朝.
繼續有後繼.
來登基的時候.
其實都是另外一個Fulfillment.
直到最後一個Ultimate Fulfillment.
就是我們的主耶穌基督.
他是一個永遠大衛的君王.
所以看這些預言的時候.
我們明白這些原則的時候.
你就明白馬太在做什麼.
馬太福音.
為什麼會有馬太福音寫成.
你記得馬太福音寫成的時候.
就是主後一世紀大概六十年代的時候.
就是AD六十幾年的時候.
當時發生什麼事呢.
就是當時的背景.
就是當時初代的教會.
受到很多的迫害.
有很多的引誘.
就是一個初生的教會.

$^{321}$岌岌可危的時候.
馬太就是其中一個福音書的作者.
就要整理一下我們的信念.
我們要堅定我們.
鞏固我們到底信的是什麼.
所以馬太就將當時.
每一個耶穌的故事.
來到去搜集起來.
當然他自己也有很多.
因為他跟過耶穌有幾年時間.
當他搜集起來的時候.
就將這些故事.
來到去編排妥當.
以致當時的教會.
看了馬太福音的時候.
他就能夠去重整他們的信念.
到底信耶穌是一回什麼事.
耶穌對於我們來說是一回什麼事.
所以馬太去整頓耶穌生平的時候.
就發現不是亂七八糟的.
就是耶穌的故事.
可能這個知道一點點.
那個知道一點點.
當他搜集起來的時候.
他發現原來不是亂七八糟.
而是有一個的模式在那裡.
有什麼模式呢.
首先第一樣.
約瑟 耶穌地上的.
合法的父親.
adopt他那個.
就是大衛的子孫.
約瑟是大衛的子孫.
所以耶穌絕對是大衛的子孫.
所以大家留意.
馬太福音開始的時候說什麼.
加普.
我們平時不會看的一些經文.
或者我們覺得很悶的經文.
但是說到加普 為什麼呢.

$^{361}$他就從亞伯拉罕.
以撒雅各.
一直數算下去.
跟著那些加普.
跟著那些故事.
那些加普的記載.
然後就傳到大衛.
由大衛一直跟著大衛的子孫.
所羅門.
然後繼續數列王.
阿哈斯 希西加等等.
到最後586BC.
就是被擄巴比倫.
被遷到巴比倫的時候.
王國的時候.
好像一切都玩完了.
一切都沒有了.
上帝的認許都不會兌現的時候.
但是好像死國翻山一樣.
在灰燼裡面.
大衛的國度又要繼續.
神又帶被擄的猶太人.
回到耶路撒冷.
這條線就沒有斷過.
幾乎斷 但沒有斷過.
繼續傳下去.
有所羅巴巴.
然後他們再重建耶路撒冷的聖殿.
Second Temple.
最後一直傳下去.
每一代的時候.
就去到約瑟.
就是耶穌地上的父親.
所以第17節.
馬太福音的家譜最後那一節.
他說這樣從阿巴罕到大衛.
有14代.
大衛到被擄巴比倫.
王國的時候14代.
王國的時候已經玩完了.

$^{401}$但誰知被擄巴比倫的時候.
一直到基督.
都是有14代.
中間不是真的有14代.
其實是跳了幾代.
但為什麼要選14代呢.
大家知道.
為什麼他將舊約整個歷史.
整本舊約39卷書.
你不用看完.
這裡一節已經幫你看完.
就是三個階段馬太福音.
14代由阿巴罕到大衛.
大衛到被擄巴比倫.
被擄巴比倫到基督.
全部都是有14代.
三個的14代.
為什麼會有三個的14代.
因為當時沒有阿拉伯的數字.
01234我們今天用的數字.
每一個希伯來的字母有22個.
每一個字母就代表一個數字.
由1至9.
然後10 20 30這樣繼續上去.
你如果找回David這個字.
當時的寫法是沒有響音的.
即是DVD或DWD.
D就是第四個字母.
W或V就是第六個字母.
計算大衛的數目.
這個名字的數目.
你就計算是14.
所以計算耶穌基督的家譜的時候.
就發現他就是皇室的後代.
耶穌就是Son of David.
並且是正式的Son of David.
Son of David.
猶太人被擄回歸.
回到以色列巴勒斯坦地.
第一世紀的時候.

$^{441}$他們好像回到自己的家.
但是他們是沒有了自己的家園.
因為世世代代都是被一個帝國.
然後再另外一個帝國來統治.
生活都是很苦.
有如今天在不同地方的教會.
鎮的子民.
在中國大陸.
在緬甸.
在一些穆斯林的國家.
他們都受到很多的迫害.
當時的猶太人.
就在羅馬帝國的軍事的佔領下.
生活當然很苦.
回到故鄉.
但卻好像沒有家園.
即是You are home.
But not yet home.
回到家裡.
好像沒有了家裡.
這些是我們今天很多人可以心同感受.
剛才我們提過.
在我們這個城市裡面.
上個星期發生的事.
一個七十多歲的.
辛辛苦苦奉公守法.
在建立家園的一個老伯伯.
走去買一份報紙.
就被人槍殺.
你眼見到的今天就是暴力.
這些殘暴或者疫情.
很多的混亂等等.
你覺得這個是不是家.
這個不是home.
你不感覺到這裡是home.
雖然你好像回到家裡.
但真正的home在哪裡.
剛好星期五我看了一套紀錄片.
講到香港.
那個紀錄片裡面.

$^{481}$告訴我們.
也讓我知道.
近年到底香港發生什麼事.
為什麼香港人會走.
為什麼香港有大的移民潮.
為什麼我認識很多香港的朋友.
包括教務.
不少都是流散到不同的地方.
家在哪裡.
在這樣的時代.
在這樣的時間.
我們經常聽到人說.
我想搬.
但我想問你想搬去哪裡.
哪裡才是家.
哪裡才是你真正可以belong的地方.
猶太人在第一世紀.
雖然回到自己的故土.
但仍然不覺得是一個家.
因為仍然是盼望著.
上帝要駕臨.
上帝要來到我們中間.
因為那個大衛的皇朝.
那個天國.
好像還沒有實現.
上帝明明說過會兌現.
為什麼到底今天會不會兌現.
會不會恢復上帝的國度.
大衛的國度.
馬太去到搜集耶穌的故事的時候.
他就見到那個家譜.
十四 十四 十四的時候.
他就想寫出來.
告訴教會 告訴我們.
這個皇者.
是沒有其他人.
只會是基督.
所以基督的駕臨.
基督的降生.
就代表皇者真正到了我們中間.

$^{521}$所以看這個馬萊利的兆頭.
童貞女懷孕生子的時候.
你不要只看.
這個是神跡.
因為我們經常關心.
那個nature of the birth.
就是virgin birth.
就是一個神跡.
其實出生的耶穌是有神性.
當然這個是.
但最主要的以塞俠書.
由他本身有的context.
到他第一世紀.
馬太再去引用的時候.
其實是在說上帝對他自己的國度.
對他的子民的信實.
他不會放下大衛家.
馬太都沒有停在這裡.
看完家譜.
看完以塞俠書第七章的時候.
到第二章.
你記得有什麼故事嗎.
有一個東方尼的博士.
朝拜新生王.
來到哪個地方.
伯利恆.
伯利恆是什麼地方.
Bethlehem.
是大衛的老家.
耶穌基督出生的地方.
是大衛的老家.
是大衛的hometown.
Bethlehem.
大希律就好像.
你看馬太第二章.
就是另外一個的pharaoh.
是一個埃及的暴君.
就是試圖去殺耶穌.
追殺他.
然後約瑟在漏夜的時候.

$^{561}$就帶著一家.
當時耶穌大概兩歲.
應該懂得走.
帶著他就逃難到埃及.
二章十四節.
約瑟就起來.
在漏夜的時候.
就帶著孩子和母親.
做難民.
下了去埃及.
住在那裡.
直到大希律死的時候.
馬太說為的是應嚴處.
直著先知所說的話.
我從埃及召我的兒子出來.
這裡就很奇怪.
因為這句話.
在舊約裡面只出現過一次.
就是荷西亞書第十一章第一節.
如果你check那裡上文下理.
不是講語言的.
不是講prediction.
上文下理講什麼.
上文下理說.
上帝會帶他的子民出埃及.
上帝會對他們信實.
上帝會去看顧他們的憐憫.
不會放棄.
並且會守他的承諾.
應嚴處是什麼呢.
整個故事就是說.
出埃及.
上帝一定會帶領他的子民.
離開遺老之地.
apply到基督.
看到基督走難到埃及.
在埃及一定會出來.
所以他就quote荷西亞書.
就是上帝一定會faithful.
上帝會守他的承諾.

$^{601}$然後帶耶穌基督出來.
也帶所有跟隨耶穌基督的人.
所以這裡又講到.
神對他自己百姓的信實.
他會faithful.
當你覺得沒有人理你.
丟下了你的手.
上帝說我的兒子我會從埃及召他們出來.
故事又繼續了.
這個所謂聖誕的故事.
大希律死了之後.
然後他就殺了.
那裡還沒死.
大希律抓不到耶穌的時候.
他就去殺死伯里衡.
裡面兩歲或以下的男孩.
當然是一個很殘酷的故事.
二章第十七節.
馬太說什麼呢.
這就應驗了耶利米先知所說的話.
在拉瑪聽見有聲音.
有痛哭有極大哀號的聲音.
拉傑為他的兒女哀哭.
不肯受安慰.
因為這些兒女都不在了.
被殺了.
但這裡就引用耶利米書三十一章十五節.
發生什麼事.
上文夏利發生的事就是在公元前586年的時候.
巴比倫來攻陷耶路撒冷.
耶路撒冷淪陷的時候.
戰敗的時候.
在拉瑪這個地方.
或馬克這個地方.
就是拉住這些戰俘.
這些猶太人的戰俘.
就將他們帶到外邦巴比倫.
是流亡.
是一個哀悼的時候.
馬太去引用這段聖經.

$^{641}$就說正如現在發生的事.
在今天在伯利恆.
伯利恆以色列的母親.
他們為了誰哀號.
為了他們失去的兒子去哀號.
正如當日耶路撒冷淪陷一樣.
那些英賢有什麼滿足感.
如果看耶利米書三十一章上文夏利.
你就知道其實他在說什麼.
他說雖然在拉瑪這個地方.
其實拉瑪就在伯利恆附近.
雖然在拉瑪這個地方我們會流眼淚.
但是上文夏利說什麼.
上帝會帶他被擄的子民.
重新再返回來.
這個滿足感就在基督這裡找到.
仍然.
就是說被擄的眼淚.
上帝一定會插進去.
上帝會紀念他子民的眼淚.
然後大希律終於死了.
約瑟就帶了耶穌和瑪利亞整家.
從埃及出來.
返回以色列地.
最後他們就去了一個地方.
加利利的北面一個小鎮.
叫做拿撒勒.
瑪太結論聖誕的故事.
指示聽見阿基努.
就是大耶希律的另外一個兒子.
又是很殘暴.
他會接續他的父親.
大希律作猶太王.
所以他就不敢回猶太地.
在夢中得到天使的指示.
就往加利利境內.
來到一個小鎮叫做拿撒勒.
由最後加句.
這樣怎樣.
就應驗了先知所說的話.

$^{681}$他必稱為拿撒勒人.
這裡最有趣.
因為拿撒勒在整本舊約聖經裡面.
沒有提過.
在一個這麼小這麼卑微的小鎮.
沒有提過說什麼應驗.
有什麼應驗可言.
這個就是猶太人讀經竅妙的地方.
猶太人的拉比.
特別是兩個中間.
他們很多時候找一些相似發音的字.
就將它串連起來.
差不多一樣發音的字.
就將它串連起來.
拿撒勒的音像什麼.
納西德既然舊約裡面沒有出現過.
就像二賽亞書第十一章.
那裡說什麼.
他說從耶西的樹幹.
必生一根嫩芽.
Nazir.
Nazirut.
拿撒勒.
Nazir.
嫩芽.
從這個根而出的枝條.
必定會去到結果字.
耶西的樹幹是什麼意思.
大家知道耶西是誰嗎.
大衛的父親就是耶西.
從大衛的父家必會生什麼.
一根的嫩芽.
發生什麼事呢.
耶西那棵樹就像大衛的家族一樣.
在亡國的時候被人砍下來.
剩下的樹幹.
好像什麼都玩完.
什麼都完蛋了.
樹都被人砍下來.
死了.

$^{721}$但是他說從耶西的樹幹.
必定會再生嫩芽.
Nazir.
然後大衛的家又再重生.
死灰復燃.
幾乎就是死而復燃.
說到耶穌在拿撒勒長大的時候.
他是Nazirut.
拿撒勒人像什麼音呢.
就是嫩芽這個字.
Nazir.
告訴神的子民.
好像什麼都玩完的時候.
原來上帝仍然忠於他的承諾.
忠於他的說話.
這些都是猶太人很常見的一些色經技巧.
叫做Midrash.
Midrash是什麼意思.
Midrash就是think over的意思.
就是seeking.
就是去尋求的意思.
connecting the dots.
就是看見耶穌的故事.
看見舊約裡面發生的事.
就將這些dots來connect.
連串在一起.
馬太就是在做這件事.
告訴教會.
耶穌出生的故事原來是這樣的.
將這些dots來connect.
connecting the dots的時候.
他看到的是什麼呢.
看到耶穌基督.
看到上帝自己的信心.
Midrash就是這樣的意思.
尋求.
其實一個可能更加好的意思就是Ponder.
就是停下來.
來思考一下.
放慢生活的腳步.

$^{761}$在上帝的話語裡面.
我們去思考一下.
到底上帝在做什麼.
思考什麼.
聖誕的孩童.
聖誕的嬰孩.
到底是誰.
到底是你的誰.
在你生命裡面會帶來什麼.
或者你看見聖誕的故事.
有不同的dots.
耶穌是大衛所adopt的兒子.
大衛之子.
Son of David.
由大衛的角度來的時候.
上帝在我們心裡面開始作王.
去管理我們生命的時候.
我們才會有愛有平安有喜樂.
出埃求.
在耶穌誕生的故事裡面.
看到摩西.
看到法老.
看到出埃求在痛苦裡面得到釋放.
真正得到自由.
在基督裡面才真正有自由.
那個freedom.
流亡.
一切都玩完了.
原來可以再回來回歸的.
真正有家.
真正的家在基督裡面才找到.
真正要拆去我們無家的眼淚的.
是得耶穌基督.
因為我們只有祂是我們的伴.
哪裡是我們的家.
搬去哪裡.
馬太華.
基督.
以馬來黎.
祂才是我們的家.

$^{801}$有上帝的同在.
你才會有家.
最後我們結束的時候.
我們回到一處經文.
很少我們會留意的.
就是新約的第一句話.
馬太福音第一章第一節.
大衛的子孫.
阿布拉罕的後裔.
耶穌基督的加譜.
有什麼好留意.
加譜的原文.
Biblios Genesis.
或者Biblios Genesis.
你一聽到這個字.
你就想起舊約裡面第一本書.
Book of Genesis.
全新的創作.
新約第一句話.
在基督裡面.
Biblios Genesis.
New Creation.
真正在基督裡面.
才會有New Beginning.
有真正的開始.
不是你搬去哪裡.
才會有真正的開始.
不是你有什麼改變.
才會有真正的開始.
不是你轉工才會有真正的開始.
在基督裡面.
In Christ.
你才會有New Beginning.
在我們生命裡面.
上帝開始作王的時候.
我們才會有一個真正的New Beginning.
我們邀請耶穌基督.
加lum 的時候.
讓祂在我們生命裡面.
有一個新的開始.

$^{841}$我們同心協力地祈禱.
主愛多謝你.
因為你是以馬來里.
是我們真正的家.
無論我們身處在哪一個地方.
不安或平安.
是順或是逆.
主讓我們在聖誕的.
你的加lum .
讓你加lum 在我們生命裡面.
在我們生命當中.
給我們一個新的開始.
主我們可以面對不同的困境.
可能仍然繼續.
但因著以馬來里.
因著基督與我們同在.
你會擦乾我們的眼淚.
你會給我們有信心繼續向前.
求你親自幫助我們.
每一個屬你的子民.
讓我們能夠在你裡面.
在神的同在裡面.
我們能夠得到安穩.
繼續的宣揚你.
跟隨你.
祈禱奉主名禁.
\newpage



\section{}
\label{sec:PUn4RKNOgHU}
\textbf{2021 年十大聖經考古發現 (陳崇基牧師)}
\newline
\newline
連結: \href{https://youtube.com/watch?v=PUn4RKNOgHU}{\texttt{https://youtube.com/watch?v=PUn4RKNOgHU}} ~~~~ 語音日期: 2022-01-02
\newline
\newline
\hyperref[sec:lTIqtT1KU5c]{\small{< < < PREV SERMON < < <}}
~
\hyperref[sec:index_chronic]{\small{[返順時目]}}
~
\hyperref[sec:index_scriptual]{\small{[返順卷目]}}
~
\hyperref[sec:6HyiM6YA71c]{\small{> > > NEXT SERMON > > >}}
\newline
\newline
$^{1}$2021年剛剛過去了.
到底在這一年裡有什麼重大的聖經考古發現.
可以幫助我們認識聖經的世界呢?.
雖然在疫情當中.
過去這一兩年.
對考古的工作有很多的限制和阻滯.
但仍然有很多重要的聖經考古發現出現了.
所以2021年有什麼十大聖經考古發現.
在此和大家分享一下.
當然有不同的人會說出最大的十大考古發現是什麼.
在此我也是按照自己的想法.
將這十大考古發現拿出來.
首先第一個我想分享的是一個叫做Selucid Fortress.
這是希臘帝國在巴勒斯坦地南邊所建立的堡壘.
這個堡壘在拉吉附近.
拉吉在哪一處呢?.
大家可以看到拉吉有耶路撒冷.
左下角就是拉吉拉基什.
對於耶路撒冷和猶太人來說.
拉吉是一個非常重要的.
很有策略性的城鎮.
這裡本身就是希臘軍隊進駐的地方.
這個堡壘的年份是公元前112年.
這個時間到底發生了什麼事呢?.
大家記得在兩約中間大概167至63BC的時候.
猶太人擊退了統治他們的希臘帝國.
然後他們就去到智立.
成為了他們自己瑪加比的王朝.
或者叫做Hasmonean Dynasty.
總共大概有160至63BC一百年左右的時間.
他們是獨立的.
他們獨立之後.
世世代代仍然在東南西北擴展.
這個堡壘按照策略是公元前112年.
就是瑪加比王朝的中間的時候.
他們就向南邊以東地的這處.
包括拉吉.
去攻打的時候.
就找到這個堡壘.
然後就把它摧毀了.

$^{41}$所以在這裡找到的不單止是一些石頭.
丟出去的石頭.
還有一些武器.
也找到瑪加比軍隊.
在攻佔摧毀這個要塞的時候.
留下了一些遺物.
這個是第十大.
第九是什麼呢?.
第九就是在埃及發現的.
在埃及發現了一個法老叫做合弗拉.
他留下了一個箭臂.
這個箭臂完完整整的埋藏在地上.
聽說是一個農夫在耕田的時候.
他一掟就掟到一個這樣的箭臂.
一個石頭的箭臂.
發現的地方是在尼羅河三角洲的東邊.
這個Ismailia這個地方.
當時就是法老帝國.
埃及的尼羅河三角洲的東邊防線.
留下這個箭臂.
這些箭臂很多時候都是一些Boundary Marker.
就是一些邊界的石頭.
留下了法老的腳印.
告訴人家這裡就是他的東防線.
這個法老叫做合弗拉夏.
他作王的時候是589至570年的時間.
這段時間跟聖經的背景很有關.
因為你記得在586 B.C.的時候.
就是巴比倫帝國在殲滅耶路撒冷的時候.
這個法老作王的時候正是在這一段時間.
耶利米書44章30節都提及過這個法老.
耶和華說我必把這個法老叫做合弗拉夏夫拉.
或者有其他的名字和其他的串法.
交在他的仇敵手中.
和他尋索他性命的人手中.
那是誰尋索這個法老的性命呢.
其實就是巴比倫王利布加尼沙.
這裡耶利米說什麼呢.
耶和華說會把法老這個帝國交給巴比倫帝國.
也會把猶大王西底加.

$^{81}$這個末代的猶大王.
交給他的仇敵.
就是利布加尼沙的手中.
這裡講清楚.
日神的子民不聽神的話的時候.
上帝就會按照他自己的方法.
興起巴比倫帝國來管教他.
不僅管教他.
也管教其他的國度.
法老和弗拉也逃不掉了.
最後也被巴比倫軍隊吞併了.
以至埃及帝國就完結了.
這是第九大.
第八大是什麼呢.
就是發現一些紫色的布隧.
發現的地方叫做庭拿谷.
這裡有一個銅礦場.
有很多年這個銅礦場.
都是在考古學家的研究當中.
你看到這個庭拿谷叫做Slays Hill.
就是這個奴隸谷.
這個地方在Google Map上.
都是烏黑色的.
因為這庭拿谷是非常多的銅礦.
今天去這個地方的時候.
發現除了沙石外.
在礦業南邊.
以色列的南部.
非常乾燥的地方.
當中也有很多留下的.
一些銅礦剩下的殘骸.
在這裡找到什麼呢.
在這裡找到一些布隧.
有紫色染上去.
這些紫色的布隧.
都是非常難找到的.
因為在古代裡面.
染料最名貴和最難找到的就是紫色.
其實今天也是一樣.
我認識一些做印刷.

$^{121}$做油墨或者染料的朋友.
跟我說過.
今時今日最重要.
最難得到的染料.
其實都是紫色.
正如我們中文說.
紫氣東來一樣.
所以在《實驗經》的世界裡面.
每逢紫色出現的時候.
就是黃色才會有的顏色.
大家記得耶穌基督在上十字架之前.
羅馬兵兵來戲弄祂的時候.
就把祂戴上紙袍.
穿上紙袍戲弄祂.
說黃色是釘十字架的.
紫色就代表尊貴.
在《聖經》裡面.
一次又一次出現.
有紫色的時候.
都是一些很名貴的產品.
所以我們叫它Royal Purple.
或者有一個特名叫做Tyrion Purple.
推羅的紫色.
或者泰爾的紫色.
為什麼呢?.
因為原來在巴勒斯坦地近東.
有一個地方出產紫色染料.
就是泰爾或者推羅.
泰爾推羅在哪裡呢?.
就是在Fleisher的地方.
就是菲利基海港的地帶.
這裡有很多古代的大的海港.
有深水有深灣的海港.
其中一個就是泰爾推羅.
另外一個你可能不記得的就是西頓.
或者上面有其他的海港.
這些海港在海岸的邊界.
就會找到這些骨骼.
這些Murray's Shell.
提煉這些紫布料.

$^{161}$或者染料.
是怎麼去做的呢?.
是非常花功夫的.
因為你要去找這些裸殼.
然後拿裡面的肉出來.
在裡面的分泌線.
才可以有一點點這些紫色的染料.
所以你要找很多這些背殼.
才能夠集成.
然後才可以染到一件衣服.
成為紫色或者紅.
紫紅等等.
在這裡找到這些紫色的布料.
是非常珍貴的.
不單是珍貴.
也告訴我們當時的社會結構是怎樣.
當我們看這些考古的發現的時候.
我們不單要知道地區.
也要知道這個考古的發現.
是代表著當時古代的社會有什麼的延長呢?.
發現的地方正如我所說的就是天那谷.
你看到這裡是很乾旱和非常冷的地方.
是當時的異東地.
發現的時間.
這個紫布料的時間就是公元前十世紀的時候.
當時是大衛所羅門建立聯合王國的時候.
你看到在南面找到紫布料.
就是說北面的這些布料或者染料.
是經過通商或者不同的方法.
去到長距離來到南面.
所以有這些長距離或者中距離的交易.
我們可以看到當時的世界是非常活躍的.
還有在這個地方找到的紫布料.
不單是這些東西是珍貴的.
代表著當時的社會已經有階層.
Stratified.
就是說有高層貴族出現.
因為他們才會有這些紫布料.
並且有些礦工工人.
所以這個社會的結構.

$^{201}$這些紫布料就告訴我們.
已經有所謂成熟到可以有幫國.
有政治的國家的結構出現.
所以你看到下面的這些Radio Carbon的Dating.
給我們的時間就是這些紫布料.
是屬於公元前1114至933.
這個就是大衛所羅門的時間.
有什麼這麼重要呢.
這些紫布料告訴我們.
社會結構已經有幫國出現.
可以有幫國的形成.
State Formation.
這個就回答到很多時候.
人們質疑大衛所羅門.
在這段時間就是他們聯合王國的時候.
公元前十世紀的時候.
很多時候都會質疑到底會不會有.
大衛王朝所羅門的聯合王國出現.
這些發現都會告訴我們.
當時一定是有可能.
並且做得到一個幫國的形成.
如果在以東地都已經有.
何況在大衛所羅門的以色列巴勒斯坦地.
另外一個發現就是第七個.
就是拉吉的初期字母的名文.
每當在考古發現中找到一些名文有文字的話.
好像這一塊拉吉初期的.
Early Alphabet的Inscription.
你都會找到這些字母.
都會覺得非常重要.
因為早期的時候.
有人寫字是代表著當時有文士出現.
Scribes.
也有這些所謂政府裡面需要的人員.
因為他們是做記錄.
他們能夠寫故事.
或者做一些交易等等.
代表當時的社會.
是有一個複雜的程度在那裡.
不單止是這樣.

$^{241}$這些Alphabet的Language.
這些字母的語言.
是繼續在發展的.
繼續到普遍的民眾使用.
這些Alphabet的演進過程.
在近東.
特別巴勒斯坦地.
我們首先找到的是公元前1800年的時候.
在埃及的西乃半島.
裡面有一個寶石的礦場叫做Selbitel Cardam.
在這裡找到一些由象形文字.
轉成一些字母的痕跡.
一些名文.
這是最早期的公元前1800年的時候.
然後我們知道在公元前十世紀.
比較早的時候.
真的有古希伯來文出現.
中間好像有些Missing link.
到底中間是怎樣演進的.
這次在Lakish拉吉.
找到這些初期的字母名文.
就被我們聯繫到.
因為原來中間那幾百年.
不斷有字母文字的演進.
字母文字有什麼重要呢.
因為這些比其他的C型文字.
埃及的象形文字.
Khayelkhiv.
Khuneiform.
都比較簡單一點.
並且傳播力比較容易一點.
因為用字母來寫文字.
不需要記那麼多的character.
只是有限的character.
然後組織成英文.
或者當時的希伯來文.
或者迦南文等等.
他們用一些Alphabetic的Script.
然後做成一些語言.
在這塊拉吉找到的.

$^{281}$小小的字母名文裡面.
你就可以找到起碼有六個字母.
在這裡看到一個圓圈.
這個是Iron.
然後有一個正方形.
然後有一個缺口.
這個就是Bate.
第二個字母.
然後這裡有一個長方形.
上面有一格下面有一格.
是凍穿了的.
這個就是Dalit.
如果你這樣Iron Bate Dalit.
拼出來就是Avat.
就是僕人奴隸的意思.
可能是某個人的名字的其中一部分.
這些是西閃族的姓名.
才會有的Avat這個字.
或者下面這裡找到.
Noon,Pate和Tar.
這裡拼出來就是Noppet.
就是Honey.
蜜糖.
可能也是其中一個名字的一部分.
怎樣也好.
看到有這樣的文字記錄.
就告訴我們不單止文字演進.
繼續到Iron Age的時候.
十世紀公元前.
有希伯來文出現.
這些都是非常重要的發現.
並且告訴我們當時的社會.
是有讀寫的能力.
雖然只是少數的人士.
文士才會有.
但也可以有這樣的技術.
而這個技術.
能夠去建立幫國.
也是非常重要的.
我們看另外一個命運.

$^{321}$再快速一點.
公元前大約1100年的時候.
這個就是.
時詩的時候.
時詩記錄裡面記載.
一個人叫基典.
他的綽號叫耶路巴力.
在今年.
也找到一個.
耶路巴力的命運.
之前的節目也提過.
這個Jeroboar.
就是在這塊.
小小的.
瓦片上.
找到Jeroboar.
這個字.
翻譯出來就是耶路巴力.
耶路巴力正如我所說.
是基典其中一個綽號.
但找到這個地方的.
土優.
就是.
Qubit Aray.
這個地方.
和剛剛提到的拉吉.
所以大部分學者.
雖然知道這是時詩記的一個名文.
也很湊巧地.
提到基典的綽號.
耶路巴力.
但是應該.
不會直接指.
那位.
基典時詩.
因為基典這位時詩.
是在北面耶穌平原.
夏加利尼這個地方.
和米地人爭戰的.
所以地理上.

$^{361}$相隔很遠.
但時間上是一樣.
這裡告訴我們.
在時詩記的年代.
又出現了.
另一個耶路巴力.
可能是另一個人的名字.
可能也是.
另一個人的名字的一部分.
但是很湊巧地.
基典在這一段時間.
公元前1100年.
都是叫做耶路巴力.
接著是第五個發現.
第五個發現就是在.
耶路撒冷大衛城的東面.
找到屬於.
希西加和約西亞時代的城牆.
這一段的城牆.
大概有幾十呎長.
還有.
都是一個人的身高.
留下來的.
這個是屬於Iron Age.
希西加和約西亞的時候.
左下角.
這裡就看到.
City of David.
大衛城.
在東面就是所謂的.
Citadel of the Ascension Wall.
即是東面防守的石堆.
堆砌出來的石牆.
我們知道在.
希西加和約西亞的年代.
耶路撒冷是由.
本身聯合王國的時候.
大衛所羅門的時候.
只不過是由聖殿山.
直到下面.

$^{401}$的大衛城.
這一段長長窄窄的.
東山上.
耶路撒冷只有這麼少.
但到希西加的時候.
公元前大概八世紀末的時候.
公元前七百多年的時候.
耶路撒冷就由本身.
Temple Mount和City of David.
延伸到西面.
包括西山和.
現在O-City南面的.
一段的.
景地.
這個耶路撒冷.
在歷史裡面.
是最大最大的.
耶路撒冷.
所以在東面.
找到這一段的City Wall.
代表在耶路撒冷.
這一段時間.
希西加和約西亞的時候.
是歷史裡面.
最頂峰.
最受防護的.
耶路撒冷.
第四個.
大的發現.
我認為是水晶石的.
印記.
印章.
或者是印章.
在上面找到.
一個黑印.
刻畫了一隻小的.
甲子和一棵.
香枝樹.
在這裡大概可以看到.
一隻甲子.

$^{441}$畫在這裡.
然後這裡有一個樹幹.
繞了上來.
有五個果實.
這個樹.
是什麼種類的樹.
專家花了一點時間.
去找.
尋訪之下他們認為.
這個就是聖經裡面.
所說的乳香香枝樹.
這個印章.
是屬於公元前第一世紀.
在聖殿山.
西南角的.
退水道.
找到這麼小的.
印章.
大概長度只有十公分.
米米書.
曾經提過.
這個香枝樹.
Balsam tree.
原來在基列這個地方.
在約旦河東.
一個山地.
盛產一種香枝樹.
出了一種乳香.
一種叫.
Balm of Gillette.
基列所出的.
Balm.
是醫藥的成份.
當時的名藥.
有什麼頭暈身患.
可以找這些乳香.
來塗.
或使用.
就能夠治病.
耶利米書八章二十二節說什麼.

$^{481}$他說.
基列沒有乳香嗎.
為什麼你要走去別的地方.
去找到痊癒呢.
基列已經有了.
耶和華說.
難道那裡沒有醫生嗎.
我的同胞.
這些猶大人為什麼不會得到痊癒.
因為他們所敬拜的.
不是耶和華.
敬拜別神.
走去拜其他人.
找其他人來救他.
所以這裡有一個第一次.
在文物上出現.
乳香.
香枝樹的.
一個刻畫.
這個印章.
可能屬於某一個名人.
或是富商.
他自己可能有一個香枝樹的.
園林.
或是出產乳香的工業.
利都耶路撒冷的時候.
是第一世紀.
他可能就.
不經意地掉了.
這個重要的.
名貴的印章.
二千年後被考古學家找回.
第三個發現.
就是.
猶大曠野的.
洞穴裡的.
發掘結果.
這個曾經在以前的節目也有提過.
就是找到一些所謂.
死海古卷.

$^{521}$這個死海古卷不是昆蘭社團的死海古卷.
那些是公元前.
一二世紀的文物.
這個所謂死海古卷在猶大曠野.
其中一個山洞裡.
找到.
是希臘文所寫的.
1832至1836.
當時.
是猶太人第二次.
起義.
反對羅馬政權.
叫做Bacuba Rebou.
在這裡.
義士.
走難的時候.
他們就帶著這些經卷.
這些經卷.
收藏在山洞裡.
所以在猶大曠野.
一個這麼乾燥的地方.
收藏得非常之好.
現在在山洞.
的法規工作裡.
找到.
這裡有撒迦利亞書八章十六節.
先知.
告訴人民.
要說真話.
講真理.
在城門口要按照誠實公平來行審判.
另外.
在其他的山洞裡.
同一區在猶大曠野.
也找到一個更加遠古的.
公元前四千年的時候.
一個小孩子的骸骨.
更加重要的一個發現.
就是公元前八千五百年的時候.
當時是.

$^{561}$新世代的時候.
剛剛人們掌握到.
耕種.
所以在這些山洞裡.
找到一個草織成的.
籃子.
裡面是用來裝一些.
穀類或收成品.
等等.
因為在猶大曠野.
這麼乾旱的地方.
所以這個草籃.
一直被保留.
你看到好像原封不動.
考古學家拿出來.
這裡告訴我們.
原來猶大曠野.
當時在公元前八千五百年.
的時候的氣候.
和農產品非常豐富.
有一個這麼大的籃子.
來收藏一些.
種植物.
收穫等等.
第二大發現.
就是釘十字架的釘子.
留在一個人的腳骨.
這個腳骨.
你看到是一個羅馬釘.
十字架的鐵釘.
穿過了.
可能埋葬的時候.
因為拿不出來或很難拿出來.
所以直接把它留在.
屍骸那裡.
然後考古學家.
在墳墓裡找出來.
時間是什麼呢.
找到這個.
人的腳骨.

$^{601}$有一個釘穿過.
就是公元後二世紀.
二世紀AD或CE.
大概.
百多二百年.
這個僑崩找到.
還有地方在哪裡找到呢.
就是在英國找到.
你看到巴勒斯坦地在下面.
然後英國.
就在上面.
Fensterton這個地方.
找到這個.
人的腳骨.
有一個鐵釘穿過.
當時是什麼時代呢.
是公元後二世紀.
羅馬帝國.
不斷向四邊.
來擴張.
直到英國.
都有羅馬帝國的遺跡.
這些是非常難找到.
雖然文獻.
經常提到羅馬帝國.
用酷刑十字架.
去懲罰異見分子.
但很少找到.
考古的文物.
直接和這個有關.
1960年的時候.
在耶路撒冷一個地方.
一個墓穴.
都找到一個僑崩.
有一個腳骨.
是一個猶太人.
他的命運.
當然都是被釘十字架.
處死.
他的僑崩.

$^{641}$是用鐵釘穿過.
如果按照.
這樣的證據.
就知道當時.
十字架被釘.
的形狀是怎樣.
兩隻腳夾在.
梁木上.
鐵釘會打過.
僑崩.
左邊一口右邊也有一口.
所以當時.
十字架的.
做法大概應該是.
這樣的形狀.
所以想起.
主耶穌為我們釘十字架的時候.
大概的重構.
形狀會是這樣.
最後最重要的.
發現是甚麼呢.
我就選擇了這個.
莫大拉的猶太人會堂.
Medlar 莫大拉.
這個地方.
以前已經找到.
第一個猶太人會堂.
誰知在一個小小的小鎮.
這個小小的漁港.
找到另一個.
第二個.
猶太人會堂.
因為政府要鋪路.
加闊公路的時候.
所以就做了一些.
考古的發現.
這些發掘的工作.
就找到第二個.
猶太人會堂.
在上面.

$^{681}$見到第一個猶太人會堂.
你今天去.
Medlar 這個地方.
莫大拉這個地方.
Medlar 這個地方.
你就找到第一個已經整理好.
可以遊覽.
可以見到的猶太人會堂.
亦找到一個.
叫做 Medlar Stone.
有工匠就將.
這個七燈台.
和兩個.
大的水罐.
代表封城.
就挑在這塊石上.
這些都是.
第一個會堂的發現.
現在甚至找到第二個會堂.
就代表這個.
Medlar 的地方.
是多麼重要.
或者我們.
看看其他的文獻.
去到講回Medlar 這個地方.
原來都是不少東西.
約瑟夫.
看到Medlar 這個地方.
就叫它做.
Tarachia.
其實翻譯出來.
希臘文就是.
鴛鵡的鎮.
Place of Fish Sorting.
以後.
猶太人的經典Talmud.
都提過這個地方.
原來叫做Medlar Nunia.
即是.
Fish's Tower 的意思.

$^{721}$Tower of Fish.
即是魚流.
將魚釣上來.
捕上來之後.
在漏放上面.
然後吹乾.
就做了一些鹹魚或乾魚.
這個Medlar 的地方.
所出產的魚.
或者它的醃鹹魚.
其實很出名.
連一些羅馬的歷史學家.
地理學家都提過.
這個叫做Strabo.
在他的大典講到自然的歷史.
就提到Tarachia.
這個地方.
就說嘉里尼湖.
出產一些很好的鹽鹹魚.
Josephus 再提到約瑟夫.
他說這個Tarachia.
Medlar 的地方.
就滿有很多造船的材料.
和很多很厲害的工匠.
你看到這些McDonald's Stone.
你就會看到.
真的有工匠在這裡打工.
還有在這個.
Medlar 的港口工作.
所以當時的猶太人.
是一個很重要的地方.
一個城鎮.
甚至在海港可以停到.
很多的.
這些小艇.
漁船等等.
Medlar 在哪裡呢.
Medlar 就在嘉里尼海的.
西邊的岸邊.
一個港口.

$^{761}$在這個地方.
原來你看回剛才的文獻.
和當時.
第一世紀怎樣提到.
Medlar 這個地方.
你就會看到他們是享負盛名.
他們出產的.
鹽鹹魚.
鹽魚.
去到羅馬京城.
人們都聽聞.
出產的魚都很好.
醃的鹽都很好.
這個背景.
你就可以看到.
耶穌說到登山寶訓.
我們要作光作鹽的時候.
原來對當時的聽眾來說.
是一些很熟悉的東西.
作光作鹽.
就是享負盛名.
就是鍾愛馳名.
人們聽到你的名字的時候.
聽到教會.
聽到信徒的時候.
聽到天國的指紋的時候.
他們會想起什麼.
是否享負盛名 馳名鍾愛.
讚好 很大的提醒.
十大的聖經考古發現.
希望能夠幫你反省.
聖經的世界.
也讓你能夠明白.
聖經的華語.
\newpage



\section{1:6-12}
\label{sec:6HyiM6YA71c}
\textbf{Following God in Tough Times (Job 1\_6-12) - Rev. Mark Chan}
\newline
\newline
連結: \href{https://youtube.com/watch?v=6HyiM6YA71c}{\texttt{https://youtube.com/watch?v=6HyiM6YA71c}} ~~~~ 語音日期: 2022-01-07
\newline
\newline
\hyperref[sec:PUn4RKNOgHU]{\small{< < < PREV SERMON < < <}}
~
\hyperref[sec:index_chronic]{\small{[返順時目]}}
~
\hyperref[sec:index_scriptual]{\small{[返順卷目]}}
~
\hyperref[sec:deow9FxV6oI]{\small{> > > NEXT SERMON > > >}}
\newline
\newline
$^{1}$of God. And today, as we get together to worship together, I would like to share a few words.
with you. And the title of the sermon is "Following God in Tough Times." We're going to turn.
to the book of Job, chapter 1. If you have your Bible with you, please turn to your Bible,.
to the book of Job, which is found at the very middle of the Bible. Now, we are living.
in a tough time. The pandemic has entered the third year with no end in sight. We are.
having bad times whenever we are suffering, whenever we are not feeling good about what.
is happening in our lives. We sometimes think of the poor Job. We turn to the book of Job..
Sometimes we want to find some wisdom. Job talks about something that is universal for.
the humanity, the universal question of human suffering. Suffering is real. I don't know.
what you are facing, but there are many facets of suffering. At least there are three areas.
we can think of about human suffering. The first one is personal suffering. Of course,.
we are all familiar with it. It comes with different forms and different kinds. Emotional.
suffering, mental challenges, loss of jobs, or you're facing the future, you find a lot.
of uncertainties, you just don't know how to cope with it. Sickness and disability of.
all kinds. These are personal suffering, and it is universal. Everyone is familiar with.
it or gets some taste of it. The second category is disasters, natural disasters. COVID is.
one of them, although it was partly caused by human failures, but it is partly also a.
natural disaster. Think of what the Americans had experienced some weeks ago because of.
tornadoes. Their families, their houses got destroyed, and they became homeless and displaced..
These are some of the universal experiences of natural disasters. The third category is.
systemic evil. Sometimes we don't think of that, but it is so real. Think of people living.
under oppressive regimes. Churches, minorities being persecuted, human rights being taken.
away from them. This is systemic oppression. We think of all these kinds of sufferings.
as we turn to the book of Job, because poor Job suffers a lot. Chapter 1 opens with this.
verse to introduce this main character called Job. Chapter 1. In the land of Uz, there lived.
a man whose name was Job. This man was blameless and upright. He feared God and shunned evil..
He had seven sons and three daughters. He owned a lot of things, thousands and thousands.
of oxen and donkeys, a very rich man, and a large number of servants. He was the greatest.
man among all the people of the East. That's how the Bible introduced the character Job..
Then the book talks about how Job was a blameless man morally or ritually or in terms of worship..
He does everything he can to make sure he feared God. He is blameless in every single.
way imaginable. Then the book of Job turns the attention to heaven. All of a sudden,.
the focus turned to heaven. Verse 6 and following. One day the angels came to present themselves.
before the Lord, and Satan also came with them. The Lord said to Satan, "Where have.
you come from?" Satan answered the Lord, "From roaming through the earth and going back and.
forth in it." Then the Lord said, "Have you considered my servant Job? There is no one.
on earth like him. He is blameless and white. He is a man who fears God and shuns evil.".
Then the conversation continues with Satan asking a question. "Does Job fear God for.
nothing?" Satan replied. "Have you not put a hat around him and his household and everything.
he has? You have blessed the work of Job's hands so that his flocks, his herds, are spread.

$^{41}$throughout the land. But stretch out your hand and strike everything he has. He will.
surely curse you to the face." The Lord said to Satan, "Farewell. Then everything he has.
is in your hands, but on the man himself do not lay a finger." Then Satan went out from.
the presence of the Lord..
This story continues in the later context. It talks about disasters and suffering one.
after another fell upon the life of Job. His children died one after the other in tragedies,.
and finally he fell really, really sick. This creates many tensions for us because whenever.
we suffer or we see someone suffer, we always ask the question, "Why?" "Why did Job suffer?".
From this very short portion of scripture from chapter 1, it seems like we can blame.
it all on Satan, blame it all on the devil because of the evil one. He is doing something.
to destroy our lives. Does that explain suffering? Does it really answer the question, "Why.
we suffer? Why did innocent suffer?" It seems not because it seems like God allowed Satan.
to bring suffering to Job. So it really doesn't address the question of why Job suffered or.
why the innocent suffered..
Now about this character Satan, there is also something that we need to pay attention to..
The English word translated "Satan", if you check the Hebrew behind the word, it is not.
"Satan", it is "Ha-Satan". Now in Hebrew, you don't have to be an expert to understand.
this. In Hebrew, "Ha" is a definite article, in English equivalent to "the". And "Satan",.
of course we get it from Satan, right? "Satan" originally was Hebrew. Now it simply means.
challenger or adversary, someone who opposes, someone who raises a question. Now it can.
be a good character, it can also be an evil character, it can be neutral. So remember,.
Satan as we know it from the New Testament, the evil one, the capital evil one, from the.
book of Revelation for example, Satan does not become a personal name or a personal entity.
until the time between the end of the Old Testament and the beginning of the New Testament,.
so-called intertestamental period. Satan does not exist in the mind of the ancient Israelites,.
the ancient Jews, until the intertestamental period before the New Testament. So in this.
time in the Old Testament, there is no understanding of Satan as a personal character. So "Ha-Satan".
is actually not a personal name, that's why it's "the Satan", right? So a better translation.
could be "the challenger". So a better translation would go, "Does Job fear God for nothing,.
the challenger in the heavenly court, ask the question?" So the challenger asking a.
neutral question, perhaps not an evil one, is a question which is legitimate, because.
the challenger is asking a question, saying, "God, if you continue to give good things.
to Job, of course he will fear you, of course he will worship you. If you keep on blessing.
Job, you will ruin his motive of being good, and eventually Job will fear you only because.
of the benefits he will receive from you." So the challenge is legitimate, right? The.
challenge is really legitimate. If you think about young children, you should understand.
that. Children will behave because they want rewards from you, they want benefits from.
you. And eventually, a child will behave only for the benefits he gets from you. So the.
challenge is a valid question, is a valid concern. And what is driving the righteousness,.
the worship of Job, right? That's the essence of the question. So the challenger is questioning.

$^{81}$how God operates the cosmos. Some sort of a challenge to God's way or God's policy.
of managing the universe. And then, from the later context, we know God readily accepts.
the challenge. God said, "Okay, Job is in your hand. His things is in your hand. Take.
his things out. Take his things away. See if he continues to worship you." So this.
short story kicks off the rest of the entire book, a very long book, the book of Job..
Now let's take a look at this really long book of Job. What is it all about? The first.
three chapters' introduction is the suffering of Job. They have a dialogue between Hasatan.
the opposer, the challenger, which is a legitimate challenge that he presents, and then God himself..
First three chapters is a short introduction. And then there's cycles of cycles of cycles.
of dialogue between three friends of Job. They're debating why Job suffered. And then.
chapter 28, there's a wisdom song. And then chapter 29 to 41, another cycle of speeches..
Job and then there's a guy who popped up called Elihu. Probably he represents the ancient.
Israel, how they see suffering, why the innocent suffer. And then lastly Yahweh, God himself,.
show up and then conclusion. Now as you can see, a big chunk of the book is conversations.
or actually debates between Job and his friends. Why does the righteous suffer? The three friends.
are saying, and they insist, that Job suffered because he has done something wrong. He has.
sinned against God. If Job is willing, the friends argue, willing to confess his wrongdoing,.
whatever that is, something you can think of or you can just make it up, just confess.
your sin, Job, and then you will be fine. Now if Job confesses his sin, the friends.
are saying, your suffering will end and you will be okay and you'll get your stuff back..
It's very simple, it's very easy. It should be the end of conversation. But of course.
not, right? Of course not that simple because Job, in the entire chapters, in the entire.
debate, Job is saying, I am innocent. I insist that I am innocent. I do not deserve this.
kind of suffering. I do not deserve this kind of pain. So the debate prolongs and continues.
and continues. Now let's pause and ask why this prolonged.
conversation is there and what it really means. As I have told you before, the Bible is not.
written to us. It was written to the ancient audience, right, in the Near East, in the.
ancient Middle East. Now in the ancient Middle East world, people believe that their relationship.
with God or gods, the false gods, follow some sort of principle and we can summarize saying.
the principle is retribution principle. Some sort of, you know, I relate to God in this.
kind of retribution dynamic. You scratch my back and I will scratch your back. You know,.
this kind of, I will do good things and then the gods will give me good things. If I do.
something wrong, I offend the gods, then I will suffer because I did some bad things..
So good things lead to good stuff that I will receive. Good things lead to benefits I will.
receive from the gods. If I do something bad or unintentionally I step the toes of the.
gods, then I will suffer. If that's the case, then I have to do something better,.
do something good so I can appease the gods. So many times, as you can see, even in today's.
world, in the ancient world or in the contemporary world, this kind of principle is what we assume..
Sometimes this principle works, right. We do good things, you know, there's a good.
return. But many times, as you can tell, it doesn't work, right. We are in the midst.

$^{121}$of pandemic, right. We are suffering because of other people's wrongdoing, right. Because.
of other people's mistakes. That's why in the entire book, Job continues to maintain.
his sinlessness, his integrity. He even demands to interview God because he, Job himself,.
he wants to understand why he suffers. He has done nothing wrong, he said. So after.
30 some chapters, God finally shows up. He did show up and answer to Job. But the way.
he answered is very interesting. It's not really answering the question of suffering..
It's not really trying to explain the reasons why the innocent suffer, why Job suffered..
Let's take a look. 38 verse 1 and 4. The Lord Yahweh answered Job out of the storm. So God.
showed up. Verse 4. He asked a bunch of questions. One of the first questions is, "Where were.
you when I laid the earth's foundation? Tell me if you understand. If you're wise enough,.
tell me where were you when I created the world? Since you question my policy of running.
the universe, running the cosmos, where were you when I laid the foundation of the earth?".
And a bunch of questions about creation, all the awesome things that God has done since.
the beginning of time. And then 38, 25. This is really interesting. It's a rhetorical question..
It means the answer is always assumed. 25. God presents the questions. "Who cuts a channel.
for the torrents of rain and a path for the thunderstorm to water a land where no man.
lives, a desert with no one in it, to satisfy a wilderness, a wasteland, a desolate wasteland,.
and make it sprout with grass?" You see what is going on here? God is saying, "Do you.
know the principle of rain in this earth?" God is saying, "Do you even notice there are.
times when rain falls on no man's land? No one will get the benefit. No one will get.
to drink the water." And God's way of operating the universe is beyond the so-called retribution.
principle. Good things lead to good return. Good things lead to good stuff and benefits..
And God is saying, "Take a look at the rain and the snow. Sometimes it falls on no man's.
land. No one will get the benefit." Why? Do you understand why? It's way beyond the principle.
of retribution. And then 39, 13. This is one of my favorites. Talk about the big birds..
Talk about the ostrich. 38, 39, 13. "The wings of the ostrich flap joyfully, but they.
cannot compare with the pinions and feathers of the stork. She lays her eggs on the ground.
and lets them warm in the sand." Understandable, right? Very wise. "But unmindful that a foot.
may crush them." Sometimes the mother ostrich crushes the eggs and kills the birds, kills.
the babies. What is going on here? Sometimes a wild animal will come by and pick up the.
eggs and destroy the babies. If God creates everything out of his retribution principle,.
such things would not happen, should not happen. If God creates everything in his wisdom, we.
should be able to understand why. And this question is saying, we just don't understand.
why God would do such things. God would create stupid ostrich, right? They are so stupid..
They are not of wisdom. And why God created them in such a way? Now all these questions.
God presents to Job is saying rhetorically, you know the answers? We just don't know..
It's way beyond our finite wisdom. We just cannot comprehend, cannot fully understand.
the wisdom of God by which he managed the universe. And then at the end in chapter 40,.
verse 3, Job, after hearing all the questions being challenged by all these rhetorical questions,.
Job answered the Lord, "I am unworthy. How can I reply to you? I put my hand over my.

$^{161}$mouth. I spoke once, but I have no answer. Twice, but I will say no more." So why suffering?.
Why Job suffering? At the very end of the book, there's no answer. The book is not.
intent to help us how to think about suffering. It doesn't give reasons for suffering. It.
doesn't really, it doesn't even defend the actions of God, right? Called theodicy. Defending.
the deeds of God, why there is suffering in the world if there is a good and all powerful.
God. Why evil exists? Why suffering exists? The book of Job doesn't really give a philosophical.
analysis of the question of suffering, the question of evil. So the book is not about.
how to think about suffering. The book is about how to think about God in our suffering..
It's not to help us think about suffering. It's to help us how to think about God in.
our suffering. The book expands our theology, expands our understanding of God. Presenting.
to us God is wise, God's wisdom is infinite. How he ordered his cosmos, how God managed.
the universe is beyond our finite mind, is beyond our comprehension. We cannot fully.
understand it. God's wisdom is higher than our wisdom. It's ultimately incomprehensible..
When I was in seminary, one of my teachers, one of the best teachers was John Frame. Dr..
John Frame is one of the best theologians out there. He said one thing that changed.
my way of worship since then. One of his theology classes, he said, "God is incomprehensible.".
We were studying theology. We were studying the attributes of God, the character of God..
We always focus on God is love, God is righteous, God is truth. All those things are good and.
they are all correct and biblical. But here, Frame says, "So often we forgot that one.
of the attributes of God is incomprehensibility." God is incomprehensible. And Frame continues.
to say, "If God is incomprehensible, his incomprehensibility is incomprehensible.".
How incomprehensible God is, is way beyond our comprehension. It's incomprehensible..
We are finite and God is infinite. God is to be trusted and his wisdom is to be trusted..
Even though we don't know all the answers, especially the answers for the question of.
suffering, the question of evil. Now, ultimately, the book of Job doesn't really offer much.
encouragement or comfort to suffering. Now, many Psalms in the Bible do, especially Laments.
or Psalms of Confidence and all those Psalms. You go to the book of Psalms, there are many.
Psalms that help us to get encouragement and comfort in our suffering. But the book of.
Job is not designed to do that. It just doesn't cut it. The book of Job doesn't really offer.
a defense of God in the midst of suffering. There's no philosophical analysis to human.
suffering. If you approach the book with such expectations, why I suffer, why the righteous.
suffer, why are we in the pandemic? We will be disappointed because the book is not about.
how we think about suffering. The book is about how we think about God in the midst.
of suffering. How do we think about God in tough times? How do we think about him, our.
theology? That determines our life, our ethics. Theology determines ethics. How we understand.
God determines how we behave, how we live our lives..
Now, let's go back to chapter one, the valid question presented by Hasatan, the challenger..
Verse 9, chapter 1. Does Job fear God for nothing? It's a valid concern, it's a valid.
question. The question is how we think about ourselves, how we relate to God in our suffering..
Do you worship God because you get benefits from him? Does Job worship God, fear God for.

$^{201}$the benefits that he has? Or he really worships God and fears God for nothing, not for his.
own benefit? Is there something called disinterest righteousness, disinterest faith goodness.
that is not concerned by our own benefits, the good things that we can get? Does Job.
worship God because he gets the benefits from doing so? Does Job worship God only when he.
experiences prosperity, blessing from God? Is Job's righteousness really disinterest?.
Does Job worship God no matter what, good times or bad times? That is the most important.
question we should ask. Do we follow Christ all because of the benefits that we have?.
Now, don't get me wrong, we got a lot of benefits from being in Christ, right? We believe.
in Jesus, we believe in Christ, we receive the redemption. I think there are a lot of.
benefits that we have. Eternal life, spiritual blessings, the honour and dignity for being.
children of God, the privileges to approach God directly in Christ in prayers, the will.
of God that we know from the will of God. These are all the benefits that we have from.
being in Christ. But remember this, all of these benefits are not earned by us. We don't.
deserve them. And they are freely given to us in Christ, in the cross. The book of Job,.
essentially, is not a book on suffering. Again, there are no answers to the question, why.
we suffer. The book of Job is not a book on suffering. The book of Job is a book on discipleship..
Discipleship in the midst of tough times. The question, do you worship God, do you fear.
God for nothing? Or do you fear God because you want to get good stuff from him? When.
good stuff is taken away, we leave God alone. It's a book on discipleship. Discipleship.
in tough times. In calling us to follow Christ even in tough times, not just in good times..
Perhaps this is the more important question we need to begin asking as we start a new.
year. Perhaps this is the most important question. Why do you worship God? Is it because of the.
benefits that you receive? That's why you follow Christ? Or if God is God, he is worthy.
of our worship no matter what. He is worthy of our attention no matter what. He is still.
the centre of our lives no matter what. Do we follow Christ even without the benefits?.
This is a question of purifying our faith. Centring our faith on Christ alone. Not on.
the good stuff that we have we get from Christ. Again there are many benefits but remember.
we don't earn them. Freely given to us in grace. But do we really worship God for nothing?.
That should be the essence of faith. God is awesome. He is worthy of our worship. Christ.
is worthy of our following, of our faith and obedience no matter what. Thick or thin. That.
is the essence of faith. This is the most important question we should ask as we celebrate.
New Year. As we start a new year. We can set a lot of goals whenever we come to New Year..
You want to lose weight. You want to live a healthier lifestyle. You want to hit this.
goal or that goal in your career or in your school. Those are really good and worthy goals..
But the most important question we need to ask is presented by the book of John. Do you.
fear God for nothing? Is God really worthy of your worship no matter what? I read a New.
Year's prayer written by someone anonymous. I find this prayer really meaningful and I.
would like to end today's worship by reading this prayer together. A New Year's prayer..
May God make your year a happy one. Not by shielding you from all sorrows and pain, but.
by strengthening you to bear it as it comes. Not by making your path easy, but by making.

$^{241}$you sturdy to travel any path. Not by taking hardships from you, but by taking fear out.
of your heart. Not by granting you unbroken sunshine, but by keeping your face bright.
even in the shadows. Not by making your life always pleasant, but by showing you when people.
and their causes need you most and by making you anxious to be there to help them. God's.
love, peace, hope and joy to you for the year to come. Let's all pray together..
Dear God, you are the awesome and greatest God. You are the creator. You are everything.
and you are infinite and ultimately incomprehensible. You are worthy of our worship, worthy of our.
attention, worthy of our obedience. And God help us as we start a new year. As this prayer.
is not by praying that you take all these bad things and sickness and pandemic out of.
our lives. Of course that's a worthy prayer, but that is not the main thing. It's that.
no matter what circumstances we find ourselves in, with sorrow or with pain, with darkness.
and struggles, we pray that you will grant us faith. Grant us obedience. Help us to follow.
you no matter what. In good times or in bad times. In good times or in tough times. Help.
us to follow you faithfully and worship you. In Christ's name we pray. Amen..
[BLANK AUDIO].
\newpage



\section{傳道書}
\label{sec:deow9FxV6oI}
\textbf{逆要作門徒 (2) — 再思傳道書 (陳崇基牧師) 2022.01.09}
\newline
\newline
連結: \href{https://youtube.com/watch?v=deow9FxV6oI}{\texttt{https://youtube.com/watch?v=deow9FxV6oI}} ~~~~ 語音日期: 2022-01-11
\newline
\newline
\hyperref[sec:6HyiM6YA71c]{\small{< < < PREV SERMON < < <}}
~
\hyperref[sec:index_chronic]{\small{[返順時目]}}
~
\hyperref[sec:index_scriptual]{\small{[返順卷目]}}
~
\hyperref[sec:qJYnzOjf_Ec]{\small{> > > NEXT SERMON > > >}}
\newline
\newline
$^{1}$各位好,我們今天繼續看看.
逆境當中如何去做一個門徒.
今天我們看聖經裡面.
智慧文學的其中一本書叫做.
《傳道書》.
《傳道書》在聖經裡面這麼多的書裡面.
算是最有哲學味道的.
因為是在探索人生的意義.
最出名的一句話就是虛空的虛空.
好像我們這裡第一章開始這麼說.
耶路撒冷作王大位的兒子傳道者的言語.
傳道者說什麼呢?.
虛空的虛空.
凡事都是虛空的.
一切所有的勞碌.
在他日光之下.
Under the sun 的勞碌.
有什麼益處呢?.
因為一切都是虛空.
我們看看這本這麼有哲學味道的書的時候.
我們需要知道整本的傳道書.
《Ecclesiastes》.
其實是有兩把聲音在裡面的.
什麼兩把聲音呢?.
就是這個傳書的結構告訴我們.
一到第十一節.
這裡其實就是一個父親.
真正的智者.
跟他的孩子來交婚.
來叮嚀他.
然後一到十二節開始.
他就去引用一個傳道者的講話.
這個傳道者的希伯來文.
叫做 Kohelet.
這個字是什麼意思呢?.
有很多的討論.
很可能這個字是來自希伯來一個文字.
叫做召集 assembly 的意思.
所以 Kohelet 這個人呢.
我們譯做傳道書.

$^{41}$可能是一個召集人.
就是有一群門生在他面前.
他就好像一個哲學家.
一個老師來教導他們.
召集他們向他們講學.
講的就是他領悟人生的意義.
所以引用他的說話.
由十二節開始.
直到十二章第七節.
都是真正的智者父親.
去引用 Kohelet.
這個傳道者所講的言論.
然後最後才有一個結束.
到底傳道書.
十幾章裡面講什麼呢?.
或者引用傳道者.
為什麼這麼長篇大論.
引用他說的話呢?.
因為這個 Kohelet.
他在做的事情就是尋找人生的意義.
所以他講了很多方法.
在他的人生裡面.
嘗試去找人生終極有什麼意思.
譬如即是行樂.
或者勤奮去做工.
儲蓄很多的東西.
這些可能在人生裡面帶來不同的喜樂.
不同的滿足.
所以二章一節我們給一點樣本.
我們看看這個 Kohelet 想講什麼.
二章第一節.
他說心裡面他曾經這樣籌算.
他說我們來用享樂來試一試.
試一試享樂裡面有什麼益處.
有什麼意義呢?.
想不到最後的時候.
他就說一切都是虛空.
他說心裡面他籌算.
不如用酒來舒暢一下.
就是喝酒.

$^{81}$喝世界上最好的一些紅酒,白酒.
所有最好的酒都喝過了.
但是他發現.
全部一切都是虛空.
或者第四節他說.
他嘗試去到大事的發展.
就是在事業上.
去到做很多的成就.
栽種葡萄園.
然後積蓄金銀.
以及積蓄很多的財寶.
各省列國的財寶都在他面前.
眼睛所見的心裡面所喜愛的.
全部都不制止.
他一味的來勞碌.
然後賺取生命裡面當中有的所謂的擁有物.
然後他就說.
在一切的勞碌裡面.
想不到都是虛空的.
察看我所做的一切.
勞碌所成就的.
想不到都是虛空.
好像如飛而去.
就是一切都是暴風.
日光之下.
是沒有任何益處.
為什麼這麼消極呢?.
哲學家傳道者Cohen.
他提出兩個原因.
為什麼在日光之下.
一切都是沒有價值的.
一切最終都是沒有意義的.
第一.
他就說.
因為日光之下沒有任何益處.
原因就是第七章第十五節.
他說因為在世間上有很多不義.
人生裡面有很多.
做過好人也未必有好報.
傳道者就說他的經歷告訴他.

$^{121}$世間往往就是如此.
七章十五節他說有義人行義.
反智滅亡.
有惡人行惡.
誰知道就會有長壽.
就是生命當中虛度之日.
所見到的.
這些是多麼沒有價值.
除了在世間上有不義.
使人生最終沒有意義.
一切都是虛無.
更重要的一個因素就是死亡.
九章第五節.
他說.
活著的人知道必死.
就是我們死定了.
死的人毫無所知.
也不再得到上智.
你看到死就使人毫無所知.
無人去紀念.
所以一切傳道者說.
在生命當中就是沒有價值的.
就像我們英文說.
Life sucks and then you die.
中文可以這麼翻譯.
做人很濕滯.
然後你就要死了.
做人很苦.
然後就要死.
傳道者所說的.
就是說他經過.
甚至追求學識.
有最好的學位.
快樂.
人生當中工作有成就.
有財寶.
有地位等等.
一切最終都是虛空.
因為最後什麼呢?.
最後我們都是要死亡.

$^{161}$所以你回到第一節的時候.
你就發現了.
為什麼他說傳道者.
好像是一個所羅門的化身.
所羅門就是做什麼的?.
在耶路撒冷作王.
大衛的兒子.
好像是.
這位就是傳道者.
就是Cole Helen所說的.
為什麼會提到所羅門呢?.
因為在舊約裡面.
最風光最榮華的.
沒有誰.
唯有是所羅門.
就像今天的Bill Gates.
今天的花花公子的老闆.
總統等等.
這些人都好像是所羅門的化身.
所羅門甚至是聖經裡面最有智慧.
最有財富的.
最多 party(派對)的.
你記得所羅門有多少個妻子.
享盡榮華富貴.
但是到最後呢.
離經背道.
一切都是虛無的.
甚至連自己人生的方向.
都不知道是為了什麼.
小孩子的時候.
我想我們大家都去過沙灘邊那裡.
砌一些沙的女寶.
砌一些 sand castle(沙灘).
我們都出汗了.
很努力的.
然後我們砌了一些很好的堡壘出來.
但是呢.
一有大浪來到的時候.
一切都化為無有.
一切變成原狀.

$^{201}$之前所做的.
就沒有任何的痕跡.
這就是傳道者的領悟.
人生有多麼風光.
哪管你今天有多少錢.
哪管你今天學識有多高.
哪管你在世上.
你享盡了所有的快樂.
是不是就是不枉此生呢?.
到最後有一個浪打過來的時候.
什麼都等於.
你最後存了多少個 number(數字).
有多少萬億也好.
都是乘零的.
time zero(時間零).
一切都是虛空.
所以整本書呢.
傳道書.
經常有一個字眼出現的.
30次至多.
就是希伯來的翻譯成為.
虛空這個字.
heaven(天空).
就是你試一下想.
你用手去掩著你的嘴巴.
然後你呼氣的時候.
或者你說話的時候.
你感覺到.
你感覺到那些氣出來.
或者我們戴著口罩.
我們會感覺到那些氣.
那口氣其實就是heaven(天空).
就是這麼短暫.
傳道者就說人生就是heaven(天空).
這個字heaven(天空).
其實呢.
也讓你想起.
聖經裡面有一個人物.
他的名字.
創世記第四章.

$^{241}$有一個人叫做阿伯.
Abel.
其實就是這個字heaven(天空).
你記得他的故事.
你也不能夠說他.
是有什麼人生可取.
因為他出現了.
就不見了.
出現遍時.
如飛的人生.
就好像我們廣東話那樣說.
好像只有一個空隙.
那麼這個就是傳道者.
就像一個哲學家那樣.
來到告訴我們.
他領悟的人生.
因為人會死.
會老去.
所以呢.
無論你做過什麼都好.
無論你賺取了多少東西都好.
你以為在日光之下可以找到意義.
找到真正目標.
其實一切都是沒有的.
到了最後十二章第一節.
傳道者.
這個 Kohelet.
還是他在說話的.
他在說什麼呢?.
他就幫我們去領悟一下.
人生最終的結局.
如果想不起來這個結局的時候.
我們就不會認為.
人生在日光之下是沒有意義的.
我們就以為.
我們可以靠著享樂.
任何的努力.
我們就會得到所謂 meaning(意義).
所謂意義.
十二章第一節他說什麼呢?.

$^{281}$他說你要趁著年幼.
衰敗的日子尚未來到.
就是你所說.
毫無喜樂的那些年日.
未曾臨近之先.
你就要當去紀念.
做你的創造主.
不要等到日頭.
光明.
月亮.
星宿.
變成黑暗.
雨後雲彩返回的時候.
你才去領悟.
這裡在說什麼呢?.
這裡明明是天亮的.
日頭有光明.
但是你的眼睛看到的是黑暗的.
其實他在說我們每個人.
都要經歷一個結果.
大部分的.
都會到最後老去的時候.
明明是天亮的.
但是我們看上去就是黑暗的.
就好像我們要戴老花鏡.
我要戴bifocal一樣.
眼睛都是蒙了蒙了.
開始老花.
電視要越來越大了.
傳道者說這個時候一定會到的.
就是我們老去的時候.
第三節.
看守房屋的發顛.
就是震動.
這裡房屋這個字.
就不只是石屋,磚屋.
我們住的木屋.
而是.
比喻作我們人的身體.
我們這個身體就好像一個房屋一樣.

$^{321}$房屋會發抖.
會震動.
然後有力的屈伸.
推磨的稀少.
就止息.
有力的屈伸.
就是以前很大力的.
但是老了之後.
我們就彎腰.
我們就寒背了.
原因就是因為我們繼續老去.
然後最後我們就會死.
推磨的稀少就止息.
推磨是以前農業社會裡面.
社會用的東西.
他們收成的時候.
將穀物放在兩塊大石的中間.
然後把它轉.
把它壓碎.
那些穀物.
那些大小麥就成為了麵粉.
那麼.
他講到老去的時候.
推磨的稀少.
就會止息.
其實就是用另一個的寓意.
來講出什麼呢?.
不僅是那些起落的聲音越來越少.
推磨的石頭.
所磨的聲音.
其實大部分的學者.
都認為.
這裡是比喻我們的牙齒.
就是我們的牙齒.
會好像推磨一樣.
去壓碎一些東西.
但是當我們老去的時候.
牙齒也會稀少.
牙齒也會止息.
好像我的牙醫經常跟我說.

$^{361}$牧師.
你晚上當然是經常磨牙了.
你經常這麼精神緊張.
你的牙齒已經看得到.
你是睡得不好.
還有.
你是很 stressful(壓力)的.
這裡就是傳道者所說的.
推磨.
就越來越少.
甚至是中止.
往窗戶外看.
就都昏暗.
就是.
在街上看出去的時候.
你發現.
全部都是昏暗的.
看東西不清楚.
街門關閉.
不敢出門.
我們不同的地方.
可能有不同的不安.
有不同的動盪.
有疫情.
就好像我們.
街道不能夠出去.
關上門.
我們自我關閉.
因為疫情的緣故.
不同的原因.
其實就是在說.
在不同的境況.
甚至人老了的時候.
社交的圈子.
都自然會縮小.
以前那些老友.
大家都看藥了.
沒有時間在一起.
推磨的響聲微小.
是不是?.

$^{401}$就是會越來越少.
得到修成.
或者牙齒.
漸漸減少.
雀鳥一叫.
人就起來了.
為什麼呢?.
因為睡不著.
天朦剛的雀鳥.
小小聲唱歌的時候.
我們一大早就要起來了.
所以練腦的時候.
就需要去晨運.
或者睡得不好.
唱歌的女子也是.
衰微.
人怕高處.
以前去到哪裡.
路上不會有驚慌.
現在就不是了.
下樓都看著每一步.
怕跌倒.
因為人老的緣故.
杏樹開花.
就是杏樹開花是白白的.
白色的.
這裡另一個比喻.
就是人老了的時候.
頭髮就會自然白.
是不是?.
渣滿成為重擔.
渣滿本身是很輕盈的東西.
現在老了.
就會變成重擔.
人所願的都廢去了.
其實這裡是在講情慾.
講我們的慾望.
老了的時候.
全部都會廢掉.
會減退.

$^{441}$因為人歸到他永遠的家.
雕喪的在街上往來.
是不是?.
越來越多去喪禮.
婚禮,婚嫁,喜事越來越少.
更多的就是在殯儀館見老朋友.
是不是?.
第七節.
塵土仍歸於地.
靈乃是歸於次靈的上帝.
如果你仔細讀一讀的話.
你就會發現.
這裡有很多的字眼.
在講什麼呢?.
講到白天.
光明.
絕亮.
星星.
光明黑暗.
人的靈.
神賜人心靈的靈.
這些字眼讓你想起什麼呢?.
想起希伯來聖經裡面.
一些很基本.
很基礎的經文.
foundational(基礎的)的 passages(詞語).
是什麼呢?.
創世記一二章.
這裡講到創世記一二章的創造秩序.
這些在智慧文學經常出現的.
提到 creation order(創造命令).
上帝所創立的這個 order(命令).
但是傳道者 Kohalut.
他就將它反過來講.
就是說人老或者面臨要死亡的時候.
就是一種反創造.
Anti-creation(反創造).
明明是光明的變成黑暗.
是不是?.
跟創造不同.

$^{481}$創造明明是黑暗變成光明.
但是因為老去.
因為人生的事實.
我們必定會死的一定會是虛無的.
所以一切會反過來.
不像原本的創造.
是一種反創造.
上一年.
我有機會去探望一些老朋友.
一個老弟兄.
現在已經七八十歲.
意猶狼.
然後因為中風.
活動非常不便.
但是因為他仍然要照顧.
一個九十多歲的媽媽.
還有她的爸爸.
他們身體更加不適.
是不是?.
有很多的病痛.
所以每一天.
全日幾乎八個小時.
十個小時.
都是要照顧兩老.
一個七十多歲已經病重的老人家.
還要照顧一個更老的老人家.
你們可以想象.
生命是多麼的辛酸.
所以他有很多的怨氣.
他就告訴我.
他說有一次有這樣的遭遇.
他在街上.
看到一些教會的朋友.
因為很久沒有見面了.
閒談之間.
他的朋友就知道.
原來他的兩老父母仍然在世.
就是九十歲長壽.
因為華人出於禮貌.
所以那位朋友就跟這位我的老朋友.

$^{521}$禮貌地說.
你父母還在世.
你真的長壽.
真的有福了.
我這位老朋友的反應就是.
他很生氣.
然後他立刻就破口大罵.
去罵對方.
就是說你怎麼知道這是有福呢?.
因為他每一天都要去照顧兩老.
很辛勞.
看到他們受苦.
很痛苦.
他說你怎麼知道這是有福呢?.
因為他人之常情.
這樣的反應.
所以幾天之後.
他就很不好意思打電話給這位朋友.
說對不起.
那天我因為太過 frustrated(煩惱).
我太過沮喪.
所以我就有一些脾氣.
但是你可以看到.
這是一個很天然.
很自然的一種經歷.
人長壽.
如果只有今世的話.
到底老是福還是禍呢?.
或者我們不要說年老了.
就是我們自己的 physical age(年齡).
是長壽或者年輕也好.
甚至我們今天.
處於一個這樣的時代.
有疫情.
政治社會很多不安.
過了幾年而已.
世界上發生了很多事情.
譬如香港發生了很多事情.
使得我們都是無所適從.
實際上我們可能過了只是五年而已.

$^{561}$但是因為有這麼多事情.
世代的崩潰.
世代的扭壞.
等等等等.
不安使得我們.
雖然只是好像過了五年.
但是人生裡面好像過了很多年一樣.
是不是?.
就算你是很年輕的.
經歷過一個這樣的時代.
這麼多的變遷.
這麼多的動盪的時候.
你也發現.
這個人生突然間非常之漫長.
好像老了很多年一樣.
很辛苦.
這就是傳道者 Coleridge 所說的.
因為一切都會扭壞.
不久之前.
我看了一個這樣的插圖.
這個插圖就是說.
哪一年的同學會.
他這裡說什麼呢?.
他說.
哪一年的同學會.
大家同學敘舊的時候.
五年之後.
成婚的一台.
未婚的有一台.
十年之後.
有孩子的一台.
沒有孩子的一台.
十五年之後.
原配的一台.
再婚的另一台.
二十年之後.
酒量好的一台.
酒量差的一台.
二十五年之後.
國內的一台.

$^{601}$搬到國外的又另一台.
三十年之後.
這個我不記得怎麼讀這個字了.
應該是吃肉的意思.
吃肉的一台.
然後吃素的又一台.
三十五年之後.
退休的一台.
未退休的又另一台.
四十年之後.
有牙的一台.
沒有牙的就坐另外一台.
四十五年之後.
自己來的一台.
輔助來的又另一台.
五十年之後.
說了要來的一台.
說了沒來的就空一台.
五十五年之後.
能夠來的.
就一台.
不能夠來的.
就是放一張照片在這裡.
紀念他.
六十年之後.
同學會已經湊不夠一台.
人生就是很短暫的.
所以.
要好好的去深入.
這個就是.
Helen 所想說的事情.
或者有很多時候我們面對死亡.
好像面對老去一樣.
好像這裡這樣的心境.
哪一年的同學會.
是不是?.
既然老去或者死亡.
是一定會來到的話.
很多時候我們的結論就是說.
since the day.

$^{641}$是不是?.
好好的過我們今生.
那就好了.
傳道者其實也是有這樣的看法.
五章第十八節.
我認為.
既善又美的.
就是人在神所賜他一生有限的年日裡面吃喝.
勞碌裡面自得其樂.
就是神的恩賜.
就是說最終一切都是虛空的.
不如現在我們自得其樂.
我們吃喝快樂.
那就算了.
是不是?.
或者我們將我們的生命裡面.
排排滿滿我們的時間表.
充斥它.
然後我們就會覺得.
很好的意義了.
現在我們珍惜我們的生命.
since the day.
是不是?.
我曾經有一個朋友.
剛剛回到香港.
他說他探很多的老朋友.
大多都已經退休了.
這位老弟兄就跟我說.
他說他們雖然退休了.
但是他們每天的時間都是排滿滿的.
這個說學跳舞.
那個又說學一些什麼新的東西.
然後又學烹飪等等.
整天都非常忙.
然後忙到一個地步.
大家朋友走在一起.
說到底在學什麼呢?.
大家都不記得.
大家都不知道.
since the day.

$^{681}$是不是?.
這個就是傳道者.
去到上.
告訴我們.
如果在日光之下的話.
life under the sun.
你儘管since the day.
但是最後.
你還是輸的.
最後面對死亡.
人生必定的結局的時候.
我們一定會死去.
英文我們說.
He who dies with the most choice wins.
是不是?.
傳道者.
不是不贊成這句話.
傳道者再幫我們想清楚一點.
再推後一點.
再按照邏輯來說.
離場的時候.
是不是擁有最多玩具的就是贏家呢?.
傳道者說人生的贏家.
就算有最多的玩具也好.
都必定死.
死定了.
所以.
你可以暫且忘記.
或者 postpone不想這個殘酷的現實.
傳道者的結論就是.
你是跑不掉的.
因為人生始終都是.
日光之下.
是去無.
你說這樣就好像整本書.
很消極.
是不是?.
但是正如我一開始所說.
傳道書整本是有它的結構的.
最後結案陳詞.

$^{721}$就不是十二章七節.
因為真正的智者.
所謂的父親.
他是在引用傳道者的講話.
引用完這一把聲音.
他十二章第八節的時候.
他就作一個結論.
其實這個父親想說什麼呢?.
他說.
這個傳道者所說的話.
其實.
是 100\% 是對的.
所以第八節這個父親再說一番話.
傳道者說完了.
十二章第八節.
他說傳道者說什麼呢?.
虛空的虛空.
一切都是虛空.
這個是傳道者的 conclusion(結論).
他的結論.
他說傳道者是有智慧的.
將知識教訓眾人的.
揣摩過人生.
查考過人生.
寫了很多的真言.
這個 Coherent(真誠).
所說的話是對的.
這是父親對他的二子所說的話.
真正的智者說什麼.
他說 Coherent(真誠).
所說的話一定是對的.
因為人生當中.
Life under the sun(生命在太陽下).
日光之下真的全部都虛無.
除非.
我們有一個.
日光之上.
Life above the sun(生命在太陽上).
一個永恆的國度.
一個神的 perspective(觀點).

$^{761}$你要看我們人生.
這樣人生才不會.
最終是虛空.
最終是如飛而去.
然後好像暴風一樣.
所以真正的結案陳詞.
就不是傳道者.
真正的結案陳詞.
是一個真正的智者所說.
說什麼呢.
十三節.
他說你們一切都聽見了.
聽完傳道者 Coherent(真誠)所說的話.
總而言之應當要敬畏神.
謹守祂的誡命.
因為這是每一個人.
有的本分.
祂的 purpose(目的).
人所做的每一件事.
連同一切隱藏的事.
無論是善是惡.
神都必審問.
結案陳詞的.
原來是真正的智者.
告訴我們不要 live under the sun(生命在太陽下).
要 live above the sun(生命在太陽上).
你看到他這裡說了三件事情.
人生的三大方面.
敬畏神.
第二方面.
是謹守祂的誡命.
第三方面.
要知道神必審問.
三方面.
有什麼這麼有趣呢.
或者有什麼值得我們留意呢.
敬畏神.
在希伯來的聖經.
舊約聖經.
經常出現這句話.

$^{801}$敬畏神.
就是智慧文學.
就是講到人生的哲學.
聖經裡面人生的神學.
就是敬畏就是跟上帝有一個正確的關係.
去敬拜祂.
經而生畏.
跟這位創造主有一個正確的關係.
神是放在我們手上的.
因為祂是創造主.
這個 wisdom literature(智慧文學).
然後呢.
keep the commandments(遵守誡命).
遵守神的誡命.
這就是希伯來聖經裡面頭一部分.
就是五經.
摩西五經.
去守上帝的誡命.
然後第三部分.
這就是希伯來聖經裡面的另一部分.
先知書.
就是說用一個永恆的角度.
去審視.
evaluate(審視).
我們現在生活裡面所做的事情.
去看清楚我們現在生命裡面的事情.
我們生活是不是 live in light of eternity(永恆的光明).
是不是有一個永恆的方向.
來過我們的生活.
希伯來聖經.
不像我們今天呢.
中文或者英文的聖經是分類的.
我們會分為體材.
歷史書一堆.
然後智慧文學一堆.
大小先知一堆.
律法一堆.
希伯來聖經.
是分開三個部分.
第一個部分.

$^{841}$就是律法書.
第二個部分.
就是我們今天叫做歷史書.
加上先知書.
叫做前後先知.
第三個部分.
就是其中一些書卷.
就是智慧文學.
傳道書就是其中一個.
這個就叫做 writing(寫作).
就是 holy writing 寫作.
傳道書最後三重的結論.
其實就是總結了整本的希伯來聖經.
經威神.
智慧文學.
有智慧的.
in light of eternity(永遠的光).
第二方面.
謹守上帝在約裡面.
給我們的誡命.
第三.
要知道將來上帝.
會有審判.
會有 judgement(判決).
所以.
按照上帝的心意.
來過你的生命.
智者就是說.
將上帝放回在你的生命裡面.
應當的位置.
就是第一位.
他是唯一的君王.
唯一的主宰.
然後其他的東西.
就會找到適當的位置.
這樣人生.
才不會到最後.
只是暴風如飛而去.
我想大家都聽過一個比喻.
叫做生命罐的比喻.

$^{881}$Jar of life(生命罐).
什麼意思呢?.
這個空罐 玻璃罐.
就好像代表你整個人生.
因為人生在世.
裡面有一些大的石頭.
也有一些小石.
一些沙塵.
這些沙粒.
是屬於一些很普通.
很瑣碎的事.
左邊這個.
就是說如果我們將一些瑣碎的事.
我們放為第一位的話.
我們嘗試把這些東西放進去.
無論是生命裡面的享樂等等.
一些我們很害怕的事情.
譬如今天疫情.
我們大家都很害怕.
我們想不要染病等等.
這些不是不重要.
這些都重要.
這些小石.
但是到最後.
如果你發現你的生命裡面.
充斥著的只不過是一些恐懼.
只不過是一些.
生命裡面的鎖匙的話.
大的石頭.
你跟上帝的關係.
就塞不進去.
但是這個Jar of life(生命罐)的比喻.
就是說如果你先.
將這些大的石頭放進去.
你的生命的罐裡面.
然後慢慢將這些小瑣碎的事.
再放進去.
你會發現.
你的生命裡面.
是能夠滿載這些東西.

$^{921}$將這些永恆的事物.
先放在.
投資在我們生命裡面.
其他所謂沙石的小事.
才會找到正確的位置.
傳道者.
或者傳道書的智者.
告訴我們人生的哲學.
在我們現在這麼不安的時候.
我們更加需要去搞清楚.
我們的優先次序.
敬畏神.
遵守祂的誡命.
要知道將來.
大小事情.
都會來到上帝面前問責.
我們這個生命.
是一個life above the sun(天上的生命).
日光之上的生命.
是一個有方向.
有真正意義的.
求主幫助我們.
我們跟上帝在這裡祈禱.
主我們求你幫助我們.
讓我們的生命.
不只是活在一個日光之下.
去讓其他很多瑣碎的事情.
來去箝制.
讓很多瑣碎的事情.
這些小事.
來去highjack(擺脫)我們.
來去騎劫我們的人生.
主願你成為我們的守衛.
幫助我們去敬畏你.
去認識你的話語.
遵行你的道.
也要知道.
我們將來要見你的面.
求你親自的幫助我們.
以敬畏的心.

$^{961}$到我們每一天.
以至我們的生命.
不是如飛而去.
以至我們的生命.
不是暴風逐形.
讓我們能夠把持到.
真正有永恆價值的東西.
你都敬畏敬拜跟出來.
祈禱奉主名求.
\newpage



\section{1:12-14}
\label{sec:qJYnzOjf_Ec}
\textbf{Pray extraordinarily to be filled with the Spirit (Acts 1\_12-14, 2\_1-4) - Rev. Jaeuk Jeong}
\newline
\newline
連結: \href{https://youtube.com/watch?v=qJYnzOjf_Ec}{\texttt{https://youtube.com/watch?v=qJYnzOjf\_Ec}} ~~~~ 語音日期: 2022-01-12
\newline
\newline
\hyperref[sec:deow9FxV6oI]{\small{< < < PREV SERMON < < <}}
~
\hyperref[sec:index_chronic]{\small{[返順時目]}}
~
\hyperref[sec:index_scriptual]{\small{[返順卷目]}}
~
\hyperref[sec:SlU_7mFzyiw]{\small{> > > NEXT SERMON > > >}}
\newline
\newline
$^{1}$Good morning brothers and sisters.
Shalom to all of you who are joining this worship through online urban voice members and all the other brothers and sisters in the world.
Yeah, peace to you. God cares for us. He knows us and he is with us today.
I'm gonna share God's Word from the book of Acts chapter 1 verse 12 to 14 and verse chapter 2 verse 1 to 4.
two short passages.
So please prepare your Bible again the book of Acts chapter 1 verse 12 to 14 and chapter 2 verse 1 to 4.
If you can read the passage with me together.
Then they returned to Jerusalem.
From from the hill called the Mount of Olives a Sabbath days walk from the city.
When they arrived they went upstairs to the room where they were staying those present were Peter John James.
Andrew Philip and Thomas Bethlehem you and Matthew James son of Alpheus and.
Simon the Zillah and Judas son of James.
They all joined together.
Constantly in prayer along with the women and Mary the mother of Jesus and with it with his brothers.
second passage.
chapter 2 verse 1 to 4.
When the day of Pentecost came they were all together in one place.
Suddenly a sound like the blowing of a violent wind came from heaven and filled the whole house where they were sitting.
They saw what seemed to be tongues of the tongues of fire.
That separated and came to rest upon each of them.
All of them were filled with the Holy Spirit and began to speak in other tongues as the Spirit enabled them.
Amen, this is the Word of God.
Thanks be to God.
Let me pray briefly.
Father God speak to us Lord. We need you Holy Spirit.
More than anything else what we now need is you Lord.
Convict our hearts do your mighty work upon our hearts. We are waiting for you.
Speak to us in Jesus name we pray.
Amen.
My dear brothers and sisters the title of today's message is.
pray.
Extraordinarily to be filled with a Holy Spirit one more time. The title is the title is pray.
Extraordinarily to be filled with the Holy Spirit.
Let me ask you this question.
What is the most extraordinary experience for you in your whole lifetime?.
As some of you may know my story the most extraordinary experience for me is that.
God gave me my wedding date when I was praying and also my wife's name when I was sleeping even.
Before I met her I didn't meet her yet.
So it was very extraordinary experience for me.
today's story in the book of Acts chapter 1 and 2 also talks about a.

$^{41}$Super extraordinary story that we cannot expect in our daily life easily.
It can happen but it's very extra ordinary.
today's story happened on the day called Pentecost or.
Woosungje in Mandarin meaning 50th day. It was the 50th day.
after the first day of the Passover festival.
the Jews in Jesus with Jesus days were.
Celebrating the day Pentecost is the day of celebrating their harvest and bringing God their first fruits.
So on Pentecost one of the top three most important Jewish festivals all the Jews who were.
Scattered around the whole world whole nations for centuries. They came to Jerusalem and they celebrated this festival together.
At the same time before this day of Pentecost as you can see on the screen you'll see a map.
Yes.
Jesus died on the cross and rose on the Sunday.
where when counting 50 days began.
Then Jesus for the following 40 days.
Jesus spent with his disciples and his where he went back to the heaven on the Mount of Olives.
Which lies just about a one kilometer away from Jerusalem.
Then Jesus 11 disciples.
Because Judas is carried out who had had a sold Jesus killed himself those 11 disciples and others.
Around 120 people in total they gathered at a place and they had even praying together.
Constantly for 10 days as in verse 14.
Okay.
So that was the 50th day after Jesus rose again.
Then the most extraordinary day for the disciples began in today's passage.
Actually these three things happen to them.
What is that number one a sound like the blowing of a violent wind?.
Came from heaven and filled the whole house.
Second thing they saw what seems to be the tongues of fire that separated and came to rest upon.
It's one of them Wow and the third thing all of them were filled with the Holy Spirit and began to speak in.
Other tongues actually in other languages as the Spirit enabled them.
But think about this.
Before the Holy Spirit came upon them who were the yes look at the screen that that's what happened to them.
Before the Holy Spirit came upon them who were the disciples?.
How were the disciples number one when Jesus went back to the heaven?.
They were very much afraid the Jewish religious leaders. They had killed Jesus at the same time.
They were still hating Jesus disciples. So they were afraid of losing their lives.
Number two, who are they? They had been following Jesus for more than three years and their teacher Jesus died.
Rose from death and went back to the heaven and that they didn't know what to do. They were at a loss.
They didn't know what to do number three.
They didn't have any power or wisdom to do a new thing in this situation.
In short, they were fearful not knowing what to do and powerless powerless.

$^{81}$Then.
Today after receiving the Holy Spirit into this passage, how did the disciples change?.
What does the change number one?.
They became so courageous and bold and they began to preach the goodness of Jesus even to the Jewish religious leaders.
Whom they were fearing.
number two the Holy Spirit.
Started the Christian Church on us on earth and the disciples were joining the most important work in human history.
In other words, they knew what they were doing.
Wow in the third change the disciples became also very wise and powerful.
Later when they opened their mouth and preached all the religious leaders and Jewish Jewish people.
They were surprised to see their wisdom and ability.
When the Holy Spirit comes upon a person what happens.
Yes a new heart driving away our fear and sadness and a new mind filled with wisdom and.
Good insight and a new life working with confidence.
Meaning direction and power. This is why Jesus says in Luke's gospel chapter 11 verses 13.
When you look at the screen, you can read it together.
If you then though you are evil know how to give good gifts to your children.
How much more will your father in heaven give the Holy Spirit to those who ask him?.
You know brothers and sisters the best gift you can ever get from God.
The best gift your God can ever give you is the Holy Spirit.
Please remember this.
then.
How do we receive the Holy Spirit if he received this greatest and most powerful gift in our lives?.
You know what when we receive it most of the problems we have will be truly solved.
then.
How do we receive the gift of the Holy Spirit?.
the answer is.
Surprisingly simple why in today's space Caesar? What did they do?.
prayer prayer.
Prayer when Jesus talked about the Holy Spirit as the best gift God can ever give us.
Jesus was explaining how to pray in Luke's gospel chapter 11.
Okay, he was explaining about prayer life in this chapter.
Jesus tells a parable a story in the story.
There was a guy who had a friend visiting him from a far place.
But he didn't have anything to feed his friend at midnight.
so this guy he went over to one of his neighbor's house at.
Midnight and then bang bang bang. He pounded on the door of his neighbor's house.
I'm sorry. I have a friend visiting me now, but I don't have anything to give him.
please give me three loaves of bread and.
Jesus says.

$^{121}$The neighbor cannot give the bread because this guy is his friend.
but he gives the bread to this guy because of his.
Boldness, did you hear me because of his boldness?.
How did the man receive the bread? Yes, remember these two words?.
absolute.
desperateness absolute desperateness and.
complete wholeheartedness.
Complete whole heartedness. He went over to his neighbor's house and midnight and banged on the door and cried out to his neighbor.
Give me some bread.
My dear brothers and sisters actually this was the same attitude when the 120 people they were praying.
Constantly together for 10 days.
They were praying constantly for their receiving the Holy Spirit because Jesus promised it.
They were praying for it.
120 people praying together for 10 days day and night. That is the key.
How can you receive the Holy Spirit?.
You should show.
absolute desperateness or.
Complete wholeheartedness in your prayer.
Yes.
You need to make an extra ordinary effort in prayer. Have you ever prayed like that?.
Yes, you need to make an extra ordinary.
effort in prayer.
When you hear this part some of you may feel quite confused or just want to reject this teaching.
Because you have heard over and over again.
Our faith is all about grace.
Yes, I know. Yes, we were saved by grace. We live our daily life only by grace.
I know but why should we make any extra ordinary effort in?.
prayer apart from grace and.
This is the place you need to make a clear distinction between the two different notions.
one is in.
Dwelling of the Holy Spirit as you see on the screen in dwelling of the Holy Spirit.
The other one is in filling of the Holy Spirit, right?.
As you know very well the in dwelling of the Holy Spirit started.
When we repented of all of all our sins and accepted Jesus as our Savior and Lord.
The Holy Spirit began to indwell us.
So if you are a true Christian the Holy Spirit actually lives in you now. Yes very well.
However, the in filling of the Holy Spirit is totally different dimension.
Think about this.
Think about this even before the day of Pentecost the disciples and all the other people.
120 people they were already saved and technically Christians.

$^{161}$Because they saw Jesus death and resurrection.
They believed in Jesus and they had been following Jesus.
So I believe that when they confessed their faith in Jesus with their mouth.
The Holy Spirit began to dwell in them.
They already had the indwelling of the Holy Spirit. Well, what's the matter?.
Even the Holy Spirit they began to dwell in them.
They were lacking in many ways as I shared in the above like what they were fearful.
They didn't know what to do. They were powerless.
So what they needed was what they needed was at the moment was more of the Holy Spirit.
Not just the indwelling of the Holy Spirit, but in filling of the Holy Spirit and as in today's passage.
After 10 days of labor of prayer.
God literally filled them up with the Holy Spirit.
They were transformed by is this in filling of the Holy Spirit and they became powerful and super effective.
So let me summarize it here.
Indwelling of the Holy Spirit happens when we accept Jesus as our Savior and Lord and the Spirit begins to.
dwell in us.
then.
When does the infilling of the Holy Spirit happen if you are already a Christian what we now need is.
Yes, not the indwelling of the Holy Spirit, but the infilling of the Holy Spirit when we ask him.
When does the Holy Spirit empower us fill fill us up when we ask when we pray?.
Extraordinarily like those people.
My brothers and sisters.
Why do we need to have?.
Extraordinary prayer to receive the Holy Spirit the infilling of the Holy Spirit.
It is because the Holy Spirit is the most.
Extraordinary person. Did you hear me? Why should we pray like that?.
Because he is the most extraordinary person you can ever invite into your heart.
He's already in you, but you are inviting the empowerment.
Like what you cannot invite the president of the United States into your messy dirty and smelly room.
Right just like that.
If you want to invite the US president into a room you need to make proper preparation.
before that.
The Holy Spirit is triune God the Holy Spirit is God.
God is absolutely holy.
Holy meaning totally different from all other creatures.
So inviting the Holy Spirit into our hearts needs our.
proper preparation and the extraordinary prayer is the best way of preparing ourselves for the gift of the.
infilling of the Holy Spirit.
Because of the following three reasons I'm going to talk about three reasons.
Number one. Why should we pray?.

$^{201}$Extraordinarily to receive the Holy Spirit to be empowered to be filled with the Holy Spirit.
First the Holy Spirit will fill your heart when you treat him as.
He is let me repeat the Holy Spirit will fill your heart when you treat him as he is.
Again, the Holy Spirit is God not your servant.
The Holy Spirit is the owner and the Lord of your life not your guest.
If you neglect the Holy Spirit in your heart and life.
The Holy Spirit already knows that if you neglect him, he knows that.
So if you are not ready to accept him as King and Savior and Lord.
He does not want to fill your heart fully and he cannot do any and it's something great.
Showing his presence in you.
In other words you should desire and treat the Holy Spirit as a.
Number one person in your heart in your life. If not, the Holy Spirit cannot do anything.
Okay, do you know this? Have you thought about this?.
Have you ever used to the search engine Google?.
Alright, when you when you look at the screen Google search engine, how many times do you use a day?.
Here now I have Kelly. How many times do you look at look into Google a day?.
Uncountably.
Uncountably because Google is a number one search engine in America.
But my brothers and sisters have you ever used to talk to go?.
What is that doctor go doctor go is one of the most used?.
Search engine in America number one Google number two being number three. Yeah.
Who number four doctor go but you never heard about the name doctor go because it is number four.
Number one and number four it has a whole whole a lot difference.
You never heard about it and you do not use it.
The same thing is happening in your spiritual life.
If you do not regard if you treat God as the number one person in your life.
You never heard about it. You didn't you never know him. You have nothing to do with him.
That's what is happening in your life.
so.
Until you make the Holy Spirit the number one person in your life. God cannot do anything.
So through your extra ordinary prayer.
you need to remove all the other idols and you love that you love more than God and.
Prepare your heart for the coming for the infilling of the Holy Spirit.
Second reason why should he pray extra ordinarily to receive the Holy Spirit?.
The Holy Spirit cannot fill your heart when it is too dirty or just a dirty when it is dirty.
The Holy Spirit cannot fill you up.
because God is.
incomparably holier than us.
Do you like the dirty smell?.
Till like the dirty odor we cannot endure it the dirty smell just like that.

$^{241}$God cannot endure our dirty sins. That is why in check at in the book of Acts.
Chapter 2 when people asked Peter, what shall we do Peter replied?.
Repent and be baptized and you will receive the gift of the Holy Spirit.
Receiving the Holy Spirit and the infilling of the Holy Spirit is closely related to.
repentance.
Repentance does not mean mean just the easy your lips or words on your lips.
True true repentance should bring the true.
changes in your life.
If you have repented you should try not to commit the same sins again.
If you commit the same sin easily.
There is the reason God cannot enter your life.
Powerfully the Holy Spirit is not working in your life powerfully because you are repeating the same sins.
Again, and again so easily not honoring God and not fearing God.
My story.
before 2006.
Even though I love the Lord.
I was doing a lot of ministry in the church as a lay person and I was but I was easily.
Repeating the same sins. I didn't fear God and I was enjoying some sins in my life.
I was enjoying some sins, but to but from 2007.
2007 in my desperate prayer in my extraordinary prayer.
I began to hate my sinning and I began to fear God as if he were right beside me.
When I began to deal with my sins seriously like this.
Through prayer, you know what God gave me a new level of freedom and joy.
Through the extraordinary prayer God gave me a new level of freedom and joy.
That extraordinary prayer freed me from my habitual sins and I was literally I was literally.
filled with a Holy Spirit.
Last the third reason why should we pray?.
Extraordinarily to receive the Holy Spirit.
because.
You can expect to be filled with the Spirit when you are literally again filled with him in your 24/7.
Do you get it?.
You can expect to be filled with the Holy Spirit when you are literally filled with him.
24/7 my brothers and sisters.
Becoming a Christian is done when we confess our faith with our mouth believing in our hearts.
Believing living as a Christian is another dimension.
It requires our ongoing true commitment in our hearts in all of our actions.
In all of our actions in other words to live as true Christians.
You need to be trained in your life a lot of training like what how to think how to speak to others.
how to work how to spend your time how to spend your money and how to how to.
Fight against your old self. Our old self is the one which is really hard to kill and then how to spend how to.

$^{281}$resist the evil spirits.
temptations.
Actually, there are so many areas for which you need to train yourself and then it does not happen.
Automatically you need to spend your best quality time.
How can you grow you spend your best quality time to grow yourself in faith and in godly?.
Lifestyle it requires your new heart and your new lifestyle.
So Timothy first Timothy chapter 4 second half of verse 7 says it says.
Train yourself to be godly.
Train yourself to be godly training.
However at the same time we all know that training ourselves in godliness is very hard.
This is not just a problem of spiritual matters.
Even think about this even taking of your body taking care of your body or.
Building a new set of a good habit is very difficult.
Just like now at the beginning of every year.
Sometimes you make very good plans very good goals.
But you know just a few days later or at most a few weeks later.
Many of us find ourselves not doing the good plan anymore.
So building a spiritually healthy habit is actually much harder than this physical habits.
Why is it so difficult.
in the same verse first Timothy chapter 4 first half of verse 7 says.
Have nothing to do with godless myth and old wives tales. Let me repeat one more time.
Have nothing to do with godless myth and old wives tales.
What is that?.
Apostle Paul already two thousand two thousand years ago.
He clearly realized the true main reason why people were not growing spiritually.
Why?.
Because they were spending all their time in listening to godless myth and old wives tales.
Greek myth the godless myth Greek myth were great stories and many even many.
Greek myths were made into films even in our days.
So like what so Greek myths were like today's blockbuster film. Do you like blockbuster film?.
I did or like great movies and how about old wives tales old wives tales were.
superstition.
superstition or.
Woman-centered stories which old woman told young women or children.
Repeatedly again and again. So like what like today's soap opera or.
dramas.
So my brothers and sisters, it's not a problem of our own own generation.
2,000 years ago when there is no TV or internet or cell phone.
People were already obsessed with the stories and they couldn't grow in their spiritual life.
Two forces are fighting right the story of the world and the story of God and.

$^{321}$How much more than how much more difficult is it?.
Is it to live a truly spirit filled life in our own generation? Why?.
Today people are constantly.
Literally 24/7 connected to all kinds of media.
So God knows this.
The Bible commends us to be spirit filled.
Being filled with God's words God's thoughts God's minds and God's desire for our life.
However, most of us are actually filled with the spirit of the world, especially through media.
Many of us are quite seriously addicted to all kinds of media and the small gods small gods of the world are.
Controlling you through media most of the time you need to know this.
The small gods of the world are controlling you through media most of the time.
That's why you feel quite empty in your heart empty and your life is disorderly.
You are busy with those media fidgeting all the time, but you feel that your life does not have any focus or.
direction.
Because you are filled with the spirit of the world day and night.
my.
Dear brothers and sisters.
This is why we need to pray.
Extraordinarily to get out of this addiction addiction to the world and addiction to the spirit of this world.
If you pursue them like now if you are not freed from the spirit of the world.
They will control you and drive you and they will drive you to meaningless and meaningless.
Moments and moments and your spirit and life will never grow. The problem is that you cannot grow.
Your relationship and your family life will not grow mature.
At the same time, you know this God has a special plan for each one of us.
But if you are addicted to the spirit of the world God cannot use you for his purpose.
Without being filled with the Holy Spirit God cannot use.
individuals or even a church.
We are all God's people whether it's an individual or a church. We need to be filled with a Holy Spirit.
My brothers and sisters, let me close.
In the beginning of every year.
When you are wondering how to bring true meaningful good changes in your life or how to be empowered.
To continuously practice those important necessary changes in your life.
Growing in prayer is the greatest solution.
When you pray extraordinarily as those.
120 people in today's passage.
God will fill you up with his spirit.
You know what and your life and your heart will surely experience the greatest change and the power.
That you cannot get through any other way.
Being filled with the Holy Spirit is the true cure for all our inner.
Brokenness, we are all broken and sick people and then all our sins and all our broken relationship.

$^{361}$Being filled with the Holy Spirit. That is the true cure.
So brothers and sisters make the Holy Spirit the number one person in your heart and in your desire.
the number one person.
Just like a Google and repent of all your sins thoroughly.
Until you truly until you feel truly clean in your conscience and free from your habitual sins.
You need to pray until you are freed from your habitual sins. And lastly.
Set aside good chunk of time and invest in prayer until you feel until you feel you are.
Obsessed you're obsessed with spending time with God. God. I want to spend time with you. I want to pray.
I want to read the Bible until you are obsessed with a holy desire.
Then with any other media in your daily life.
God bless you.
so brothers and sisters.
Take the proper posture.
Now we are gonna sing the response song.
God spoke to you through my message and this is the time you are speaking to God.
God this is my heart.
God this is my desire.
Please empower me. So we are gonna sing the song.
Lord, I need you as the song sings.
Please have the desire in your heart Lord. I need you more than anything else.
Grow my prayer life. Let's sing together.
Lord I come I.
Confess.
Bowing here I find my rest.
Without you I fall apart.
You're the one.
Take guys my one more time. Lord, I come Lord I come.
I confess.
Bowing here I find the red without you.
Without you I fall apart.
You're the one take guys my heart.
Lord I need you. Oh I need.
Every hour I.
My one defense my righteousness.
There's.
Your grace is more.
Where this is fun.
Is where you are?.
And where you are.
Lord I am free.

$^{401}$The holiness.
Is Christ in me.
Lord I need you.
Every hour I.
Do.
My one.
defense my.
Righteousness Oh God.
Hawaii me one more time Lord. I need you.
Lord I need you.
Every hour I.
Sing it out my one.
Defense my righteousness. Oh God, how I need.
So teach my songs to you rise to you.
Temptation comes my way.
When I cannot stand I fall on you.
Jesus you're my hope.
They so teach my songs to rise to you.
When temptation comes my way.
When I cannot stand I fall on you.
Jesus you're my hope and stay.
Lord I need you.
Oh.
Every hour I.
My one.
Defense my.
Righteousness Oh God, how I need one more time with all our hearts.
Lord I need you.
I need.
Every hour I.
My one.
Defense my.
Righteousness Oh God, how I need.
All or that is you would I need.
I need.
a.
Three hour I need.
My one defense I will.
Defense my.
Righteousness Oh God.

$^{441}$My one defense my.
Defense my righteousness Oh God.
Oh God, how I need you.
Brothers and sisters, let's pray together.
When you pray take a right posture, you can kneel down on your floor.
Please lower your lower your posture.
We are standing in front of the greatest King.
Who knows our hearts and desires and our life?.
Who understands the darkest problems in our hearts who understand the most serious problems in our life?.
So my brothers and sisters, let's pray together.
The prayer music is their prayer music.
Okay, let's pray in your heart first prayer as I shared this word brothers and sisters.
The reason God cannot empower us the reason we cannot be filled with the Holy Spirit is that we are.
Neglecting we do not honor him in many ways. There are so many other idols.
We are addicted by all kinds of media. You cannot spend the true quality time with God.
And then our heart is filled with all kinds of dirty things.
If we are not holy enough.
God cannot intervene in our lives. If you do not welcome him the Holy Spirit cannot.
Intervene in your life. So my dear brothers and sisters now, let's pray with the right attitude and posture.
God come into my heart.
Remove all the dark and dirty sins and remove all the idols inside my heart.
Let's have this prayer of repentance together about our idols.
About our addictions to media and our dirty sins. Let's have this time of repentance.
Our heavenly Father have mercy on us Lord. We are broken people.
We are broken people Lord. Even we cannot keep our the best promises to you. Father God.
We are all broken people Lord and we do not have any vision in our eyes Lord.
Our life and our time is so disorderly Lord. We do not know how to honor you Lord.
We do not know how to love you Lord. Your eyes are all the time watching over us Lord.
But we are forgetting you in our lives Lord. So God would you do a new thing upon us Lord?.
Have mercy on us Lord. In this year Lord.
Oh Father God do a new thing in all of us Lord. All of the brothers and sisters in our life.
What we now need is your spirit Lord. Holy Spirit.
We need you Lord. We need you. Come to your people Lord..
Come to your people and do a new thing in their lives Lord.
And my brothers and sisters.
Even for the most important thing.
Even for the really great things if you do not plan it.
If you do not make a good.
very clean plan about it.
You cannot do it.

$^{481}$So after this worship.
The Holy Spirit wants to bring the change into our lives. In other words.
He wants to lead you to the extraordinary prayer life.
So my brothers and sisters now when you pray make a decision in your heart.
God I want to pray more than ever before.
\newpage



\section{創世記}
\label{sec:SlU_7mFzyiw}
\textbf{創世記第2章 (2.1) - 背景釋經 (陳崇基牧師)}
\newline
\newline
連結: \href{https://youtube.com/watch?v=SlU-7mFzyiw}{\texttt{https://youtube.com/watch?v=SlU-7mFzyiw}} ~~~~ 語音日期: 2022-01-16
\newline
\newline
\hyperref[sec:qJYnzOjf_Ec]{\small{< < < PREV SERMON < < <}}
~
\hyperref[sec:index_chronic]{\small{[返順時目]}}
~
\hyperref[sec:index_scriptual]{\small{[返順卷目]}}
~
\hyperref[sec:l4PeWdLBiWQ]{\small{> > > NEXT SERMON > > >}}
\newline
\newline
$^{1}$這次講到聖經的創世記第二章.
第二章講到伊甸園.
亞當夏娃等等的故事.
上一次我們講到第一章.
就是講到這個釋經的問題.
就是我們釋經上需要有一些工具.
就是有一些概念.
有些概念可以幫助我們理解一段聖經.
無論是看創世記的頭幾章.
當然頭幾章已經是有它複雜的程度.
有它的難度.
但是其實新約也好舊約也好.
這些都可以幫助我們.
我上一次講過一些釋經的事業.
每次我們讀經.
看聖經.
教聖經.
我們都要去到有一些事業.
你當作是一些工具.
第一個工具就是.
聖經的世界不是我們現代人的世界.
這個應該大家都很清楚很明白了.
因為聖經的語言.
它的語境.
它的文化背景.
都跟我們是不同的.
縱使聖經已經翻譯成我們現代的語言.
中文或者看英文.
現在國際的語言無論如何.
但是我們都知道要有一些原文.
有一些古代近東的作品.
古代近東的一些文獻.
讓我們可以參照一下.
到底聖經跟當時的文獻怎麼去對話.
聖經跟當時的文化.
怎麼去做一些抗衡也好.
或者跟它有一些爭辯也好.
甚至聖經考古.
我就是讀聖經考古的.
所以有一些.

$^{41}$過去200年.
在中東在以色列巴勒斯坦地找到的東西.
怎麼去幫助我們.
能夠明白.
以前聖經的世界.
怎麼去表達一些信息.
所以做了一些這樣的比對研究.
讓我們能夠明白它的背景.
另外就是一些的文理.
上下文.
我們要知道題材.
知道一些結構.
修辭的技巧.
上次說過.
《創世記》第一章.
有很多作者.
希伯來作者.
他自己的修辭技巧.
他放在了經文裡面.
是要讀者.
要聽眾去留意的.
第三就是去了解一下解經的歷史.
過去教會由新約.
《史圖行傳》.
《史圖時代》.
到現在有2000年.
我們累積了很多的經驗.
沒有一個人可以看聖經是自己一個人看.
一張白紙去看的.
我們總會有所謂道統的共識.
過去教會.
對於神學的教義.
聖經的解釋.
有一些某程度上的認知.
這些認知的歷史可以幫助我們.
有時候教會可能會搞錯了.
未必解得這麼通.
但是很多時候在這些碟碰碰裡面.
上帝給我們一點點能夠明白.
聖經的一些信息.

$^{81}$這個大家可能會覺得.
聖經有時候解錯怎麼辦?.
很麻煩的.
大家知道.
聖經裡面.
我們讀神學的時候.
有一個教義叫做.
Clarity of Scripture.
就是聖經的清晰度.
聖經清晰度這個教義是什麼意思呢?.
Clarity of Scripture.
意思就是說.
聖經的啟示.
其實三歲小孩五歲小孩都懂得.
都明白的.
所以有一些很清楚的.
很重要的一些信息.
譬如耶穌基督是唯一的道路真理生命.
耶穌的救恩.
十字架的救恩.
這些是幾歲大.
剛剛信主.
還沒信主.
甚至信了主很多年.
大家是同宮.
大家是傳道的.
都會明白的.
這個就叫做聖經的清晰度.
但是聖經裡面.
我們大家都知道.
有一些東西是比較困難的.
有一些東西是要經過一些傳息.
我們才能夠掌握得到的.
那些比較有深奧的東西.
我們是需要去學習的.
所以這個沒有問題的.
在聖經裡面.
我們有一些很基本的教義.
三位一體的神論等等.
基督的救恩等等.

$^{121}$這些一聽就明白了.
但是在歷史裡面.
我們是會有一些學習.
會明白得更多.
特別是我們到了現在.
現代的時候.
我們更加掌握到更多的工具等等.
所以這幾個方面.
給大家一個概念.
每一次打開聖經的時候.
我都會分開這幾部分.
聖經背後的世界.
文化背景跟我有些不同.
所以有時候我看聖經看不明白的時候.
可能他的文化是不同我的.
現代人的文化.
或者他一些文學的結構等等.
聖經很多時候.
某些經文.
重點不是在最後的.
在中間.
好像三文治一樣.
這些不是我們現代人文學修辭的技巧.
所以我們去明白這些事件的時候.
我們會更容易明白.
聖經所啟示的信息.
有幾件事情先重複.
就是解經裡面.
解讀聖經的時候有一些要點.
第一就是我們相信聖經沒有錯.
因為上帝是真理的上帝.
他不會說錯話.
也不會說一些誤導的信息.
聖經沒有錯是我們所堅信的.
但是聖經沒有錯.
不代表我對聖經的詮釋是沒有錯.
對不對?.
想清楚.
我相信聖經沒有錯.
不代表我詮釋聖經.

$^{161}$我對聖經的某種詮釋.
解讀.
一定是沒有錯的.
所以我對聖經的詮釋.
不一定是絕對的.
我可以繼續修改.
繼續進步.
繼續認識更多的時候.
我就更加深入.
更加完備.
去明白聖經.
這些都是一些不是最淺白.
最容易明白的.
聖經清晰度的教義.
基本的真理等等.
因為在聖經裡面.
有一些比較難去掌握的經文.
為什麼我們需要去學習這方面呢?.
去有更好的詮釋.
因為我們去教導別人的時候.
或者我們自己.
我們吃飽了.
我們去餵養自己的時候.
我們需要了解聖經裡面的信息.
所以在聖經詮釋裡面.
首先聖經不是直接寫給我們的.
你說是不是將本聖經拿走了?.
不是的.
聖經是為我們每一代的信徒而寫的.
這兩者有一個分別.
這個分別也很重要.
聖經不是直接寫給我們的.
換句話說聖經.
要去明白它的信息的時候.
特別是一些深奧的經文.
我們就需要知道.
他們當時的背景.
他們當時的語言等等.
這樣我們才會了解.
上帝藉著經文.

$^{201}$跟以前古代的作者想說什麼.
第一代的讀者想說什麼.
而哪一件事是當時的文化背景.
哪一件事是啟示的時候.
是超越文化的背景的.
以至我今天能夠形容.
所以聖經雖然不是直接寫給我們的.
但是卻是為我們每一代的信徒教會而寫的.
第二就是我們避免將我們現代人的意識.
我們的意義.
套入聖經裡面.
我們很容易的.
因為特別是一些現代的字面識經.
字面識經很多時候是我們自己的文化背景.
我們的概念.
譬如一些現代科學的概念.
我們問的問題.
我們套入聖經.
我們要聖經去解答.
其實聖經想解答的問題.
往往不是我們今天在問的問題.
往往我們最關注的一些事物.
不是聖經直接要去解答的.
往往聖經就是要引導我們去解答一些更重要更基本的問題.
譬如我們人生的方向.
我們追求的意義.
我們跟上帝的關係.
這些更加基本更加重要的問題.
往往不是我們現代人所關注的.
但是聖經要引導我們去明白.
所以我們要去認識古代聖經的世界.
所以我們有考古.
有背景.
等等玄文的東西.
在這裡先講清楚.
就是聖經裡面.
用語言來.
人的語言.
上帝藉著人的語言.
向我們啟示他的真理.

$^{241}$是不是?.
上帝想介紹他自己.
用什麼方法呢?.
藉著人的語言.
很多時候我們沒有想清楚語言的本質.
我們需要去想一想.
或者當作是一些概念工具.
去幫助大家去了解聖經的語言.
聖經裡面的語文.
有一位20世紀很有名的英國學者.
一個基督徒的護教學家.
他叫做C.S. Lewis.
C.S. Lewis在他一本或者一份文章.
是沒有公開出過書的.
沒有出版過.
這個文章是講什麼的呢?.
這個文章是很簡淺的去剖析一下.
到底語言.
有多少種類的語言.
以及不同的語言其實都是在表達真理的.
C.S. Lewis認為有三種語言.
大致上.
第一種就是我們平常用的語言.
Ordinary language.
我們平時溝通的時候.
已經是非常有效的了.
我們就是用一些平常的語言.
第二種語言就是科學的語言.
Scientific language.
這些就是我們現代人很喜歡的一種語言.
我們覺得有時候.
科學的語言才是表達真理的.
或者表達真理更加好的.
C.S. Lewis認為這只不過是其中一種表述真理的語言而已.
第三種就是.
詩歌的語言.
有詩意的語言.
就是 poetic language.
他認為有三種方式.
三種表述的方法.

$^{281}$來表達真理.
所以他在討論不同的語言的時候.
他其實是在討論解經的問題.
什麼意思呢?.
舉個例子.
C.S. Lewis就說.
想著想著.
我們怎麼去描述寒冬.
Winter.
當然他是用英文的例子.
我這裡嘗試去應用中文.
他說有幾種方法來描述寒冬.
第一.
就是今天十分寒冷.
是不是?.
你說很簡單而已.
這種是什麼語言呢?.
是科學語言.
是平常日常用的語言.
還是詩歌語言.
C.S. Lewis就說.
是平常凡人用的語言.
是不是我們平常說的.
就是今天非常寒冷.
是不是可以表達真理呢?.
是可以表達真理的.
是準確地表達真理的.
但是也有另外一種語言.
就是去表達寒冬的.
那個語言是什麼呢?.
室外的溫度是攝氏-10.6度.
濕度是55\%.
那你說.
這是不是表達真理的語言呢?.
沒錯是不是?.
但是這種語言是什麼呢?.
這是一種科學語言.
科學語言.
跟剛剛的語言不同.
是不是?.

$^{321}$因為平常的語言.
是很直接地告訴你.
去表述一件事情.
平日用的.
我們不是這麼分析性的.
但是科學的語言就是.
一種分析度很高的.
較高的一種語言.
科學的報告.
生物,物理,化學等等.
這些語言都好像是另一種.
今天的溫度攝氏是多少.
是不是?.
準確度是10.6度.
是不是?.
這種是科學的語言.
第三種語言是什麼呢?.
就是詩歌類的語言.
這裡就用了黃安石的《梅花》.
一首詩.
是五言詩.
他就去表達寒冬.
講到雪景是怎麼樣的.
"長閣有幾枝梅花".
"令寒".
就是很嚴寒的冬天.
"獨自開".
"遙知不是雪".
"惟有暗香來".
他用詩歌來表達.
寒冬.
你就會發現詩歌表達的東西.
是不是表達真理呢?.
當然是表達真理.
但是.
不是分析性的.
不像科學性的語言.
也不像平日的語言.
因為詩歌想做什麼呢?.
有詩歌的修辭.

$^{361}$就是想給我們多一點感受.
多一點觸覺.
我們會看到什麼呢?.
他說看到遠遠的白色的其實不是雪.
其實是梅花.
為什麼呢?.
因為有暗香.
聞到梅花.
是不是?.
所以呢.
他就表達了一種意境.
就是冬天你想到的圖畫會是怎麼樣的.
正在下雪.
或者下了雪.
梅花上面有一些雪花.
但是呢.
看到白色的.
其實是看到梅花的.
所以他對於感受來說所說的東西就深很多了.
所以有幾種語言呢?.
三種語言.
我想問一下.
三種語言裡面哪一種語言.
是更能夠表達真理的?.
其實每個語言都可以表達真理.
是不是?.
視乎你用什麼角度去看.
C.S. Lewis 就說.
三種語言.
平日用的.
詩歌的.
以及科學的語言.
聖經是哪一種語言比較多呢?.
聖經就是平日用的語言.
太陽從東邊升起.
是不是?.
西斜.
日落.
這些都是我們平日用的語言.
也是聖經世界會用的語言.

$^{401}$詩歌的語言更加不需要說了.
詩篇一百五十篇.
以及有很多其他詩歌.
都是想表達真理的.
所以不一定是科學的語言才可以表達真理的.
對嗎?.
這些只不過是我們現代人所謂的偏見.
上一次我說了.
看到《創世記》第一章的時候.
J.I. Packer.
這個系統神學家.
已經去了英國的系統神學家.
他就說我們沉迷於什麼呢?.
現代人只是沉迷於.
理解創造界.
理解宇宙怎麼運作.
理解宇宙怎麼來.
而我們永遠都不會問.
宇宙背後誰是創造主呢?.
誰是創造主呢?.
Who?.
對嗎?.
WHO.
不是How.
不是宇宙怎麼來.
這些是我們現代人喜歡想的問題.
但是聖經不是要解答這個問題.
聖經是告訴我們.
這個浩瀚的宇宙背後.
到底誰是創造主.
值得我們敬拜的.
這是《創世記》第一章想表達的.
絕對不是科學的語言.
是一些平日用的語言.
加上有一點詩歌的修辭.
正如我上次講過.
所以整個創造.
主要關注的不是物質創造的過程.
這是現代人的議題.
科學的議題.

$^{441}$你說很慘啊.
聖經就有缺陷啊.
沒有講科學.
沒有解答現代人科學的問題.
你記得上一次講過.
不是缺陷.
其實反而聖經是引導我們問一些更正確.
更需要問的問題.
就是敬拜的是誰.
所以《創世記》第一章.
所講到創造的時候.
就是在講一個統治.
一個國度.
一個創造主的 kingdom(國度).
他是君王.
他去到建立規律.
使大地和天地之間.
井然有序.
井井有條.
這是聖經要去講的訊息.
告訴我們什麼是神學.
就是上帝統治一切.
這也是古晉東所有其他民族.
以色列以外的其他民族.
亞述帝國.
巴比倫帝國.
其他的民族.
甚至迦南人,埃及.
他們寫的創造的一些技術.
要解答的問題.
不是現代科學的問題.
所以聖經是想講什麼呢?.
是想講一種神學.
不是想告訴我們科學.
所以第一和第三天.
第四和第六天.
很明顯就是有平衡了.
第七天才是最高峰的時候.
最重要要講什麼呢?.
你記得第七天是沒有講晚上有早晨.

$^{481}$第七天.
進入了第七天的時候.
就是上帝坐在他的寶座上.
他登基.
他是正式大地宇宙的君王.
是不是?.
所以為什麼我們會敬拜呢?.
所謂安息日的時候我們會敬拜.
我們六天要作供.
但是我們安息日第七天.
我們就停下來.
因為我們知道誰掌管宇宙.
世界上發生什麼事情都好.
我們要去敬拜他.
為什麼呢?.
因為我們知道上帝仍然掌管一切.
是不是?.
他在聖殿.
他在致聖所裡面.
他的寶座.
兩個基督柏之上.
約櫃裡面.
上面.
他是坐著為王的神.
所以我們需要去敬拜他.
所以我們可以放下我們的工作.
知道一切都是上帝掌管的.
所以這是創造神學.
不是解答科學問題.
而是在講主權的問題.
在講權能的問題.
第二方面.
就是天地.
整個天地.
告訴我們.
是用很多七字.
或者七個倍數.
來描述的.
好像修辭詩歌的一些方法.
但是其實在當時聖經的世界裡面.

$^{521}$是經常有這一類的文獻的.
這一類的表述叫做建殿的文獻.
Temple Building Text.
無論是以色列.
講到第一個聖殿建造的時候.
所羅門的聖殿.
七年建造.
講到會幕七段.
講到天與地.
整個宇宙都是有七或者七的倍數.
其他的民族講到這些廟.
神殿的建立的時候.
都是用七段的.
七日等等.
經常都出現這些修辭的方法.
所以告訴我們整個宇宙.
在聖經的世界觀來講.
整個宇宙是一個神聖的地方.
整個宇宙就是上帝的居所.
祂有祂的同在.
以及彰顯祂榮耀的地方.
整個世界本身原創的時候.
就是一個神聖的地方.
Sacred Place.
一個神居住的地方.
是要敬拜祂的地方.
所以上帝安息就是祂登基.
統治萬人.
你比對一下其他的創世的神話.
你就可能更加清楚.
創世的第一章.
想怎麼表達上帝的權能.
其中一件在考古裡面發現的.
一些近東的文獻.
就是以往近東的文獻.
就是主前千八二千年的時候.
一些留下來的文獻.
這個文獻叫做 Eluma Elish.
中文的翻譯我只找到一個而已.
就是以前中文大學.

$^{561}$一個近東的學者.
日文學的學者.
叫做饒宗賢.
他翻譯了這個 Eluma Elish.
叫做近東開闢史詩.
其實就是巴比倫帝國.
巴比倫的民族.
有的史詩.
這個史詩是講到.
巴比倫人怎麼去看這個世界的開始.
天地開闢的時候.
有一點像創世記第一章.
第一行.
就是文獻第一行.
講到什麼呢?.
就是講到天上發生什麼事情.
天之高系.
饒教授是這麼翻譯的.
When on high.
講到上面什麼呢?.
有水神.
水神然後他們一男一女.
然後交替.
生成每一代的其他的眾神.
其他的眾神一路生一路生的時候.
就生到一個巴比倫最大的神明.
就叫做馬布克.
然後發現.
當時有一些問題.
就是那個水神.
他去威脅其他的神明.
這個水神叫做切勿.
Tiamat.
他其實就是大水混沌的化身.
女的水神.
然後馬布克.
就在眾神混亂當中.
他就興起了.
他就是做一個最有權能的神.
他就打敗了這個水神.

$^{601}$然後打敗了這個水神.
就將這些水分成兩半.
上面一半下面一半.
他就創造了天地.
然後最後.
就創立了人類.
喏.
你聽聽這些故事.
有一點像《創世記》第一章.
還有馬布克.
是在《耶利米蘇五十章》第二節提過的.
巴比倫的主神.
就叫做米羅達.
其實就是馬布克.
或者另一個稱呼.
一個綽號.
就叫做Bell.
比勒.
如果你看看《耶利米蘇五十章》第二節.
你就知道那裡出現了.
《耶利米蘇五十章》說什麼呢?.
就是說去殘害神的子民.
使得神的子民被擄於外邦.
毀滅了耶路撒冷的巴比倫的大神.
巴比倫帝國.
一定會覆亡的.
他就說了他的預言.
所以那裡提過馬布克就會覆亡.
巴比倫的大神會覆亡.
無論如何.
近東的開闢史詩是在說什麼呢?.
就是你看到有大水開始.
創造.
創立世界之前.
是有大水.
其中一個水神就是切勿.
他就是混沌.
然後有馬杜克出來.
巴比倫的大神.
戰勝了水神.

$^{641}$然後就將他砍開.
分成兩半.
將水分成兩半.
上面一半下面一半.
然後就有土地出來.
乾地出來.
然後世界就被創立了.
人類就出現了.
這裡很明顯就是跟《創世記》第一章.
有一些語句.
一些描述.
是相似的.
但是有很大分別.
有什麼大分別呢?.
在《創世記》第一章描述.
神創立世界的時候.
有沒有打鬥呢?.
沒有打鬥.
全部事情都很安靜.
開始的時候.
你記得.
起初神創造天地.
地是空虛混沌.
冤冤黑暗.
神的靈運行在水面上.
是不是?.
我們科學的頭腦就說.
為什麼宇宙開始的時候首先會有水.
會有混沌.
冤冤黑暗.
如果你看回當時的背景就知道.
他是在用當時背景的概念.
去講創造.
告訴人們.
真正創造世界.
真正掌管世界的.
那個才是真正我們要敬拜的.
那個就是以色列的神.
耶和華.
所以沒有什麼打仗.

$^{681}$是不是?.
神的靈運行在水面上.
大水之上.
然後到第二天的時候.
你記得第一天有光暗.
然後第二天發生什麼事情?.
就是將水分開.
但是水有沒有掙扎?.
就是有沒有東西可以抗衡上帝創造主的勢力.
是沒有的.
就跟近東的巴比倫的神話.
完全不同.
所以你看這些神話是很刺激的.
因為講到打鬥.
是不是?.
但是看創世的第一章.
一件事都很安靜.
因為萬物都在上帝的掌管下.
沒有一件事是跟祂有抗衡.
沒有一件事是可以反對上帝的勢力.
上帝說要有什麼就要有什麼.
是不是?.
現在上帝第一句話.
就是要有光.
是不是?.
然後發生什麼事情?.
就有了光.
然後事就這樣成了.
七次講到.
是不是?.
慎看著是好的.
沒有什麼水神打鬥等等.
你這樣的背景看創世的第一章.
你就越來越發現祂想講什麼.
萬物.
原來沒有神聖.
是不是?.
只有一位神.
然後萬物沒有跟這位創造主有任何鬥爭.
一切都在祂的權能下.

$^{721}$祂的掌管下.
這些壁畫.
經常在巴比倫帝國.
亞述帝國.
近東的外邦.
出現的.
就是這些神靈.
跟水怪和怪獸打鬥.
打完之後.
它打贏了.
它才創立世界.
創世的第一章就不是講這件事情.
所以我們回到創世記的神學的時候.
主要不是講科學的問題.
是建立世界.
誰統管萬有.
天地之間你比對古近東的神話的時候.
你就知道原來是上帝創造宇宙的時候.
是美好.
以及井然有罪.
並且全部事情都是聽上帝的話.
第三方面是講到人類的時候.
剛才我們記得.
講到近東的開闢史詩.
也是講到最後有人類出現.
創世記第一二章.
都是一樣.
是在講什麼呢?.
創立世界之後.
最後就有人類出現.
但是你要知道它所解答的問題.
不是人類來源的歷史.
不要以為聖經裡面就給我們家譜加起來了.
我們就知道聖經教我們.
說地球是很年輕的.
聖經不回答這些問題.
是不是?.
年老也好.
年輕也好.
這不是聖經要去處理的問題.

$^{761}$因為不是在講現代科學.
人類什麼時候出現.
多少千年出現.
多少萬年.
多少十萬年出現也好.
都不是聖經要去處理的問題.
最主要處理什麼呢?.
他要交代一個世界觀.
怎麼去看待這個世界.
怎麼看待人生.
主旨是做什麼呢?.
就是去解答人類存在的身份特性和角色.
人在這個宇宙裡面存在.
到底我們的角色是什麼呢?.
上帝原創我們的時候.
我們的身份是什麼呢?.
這就是聖經要解答的問題.
不是現代科學的問題.
如果你喜歡我們的影片.
就記得要點讚,留言和分享.
\newpage



\section{創世記}
\label{sec:l4PeWdLBiWQ}
\textbf{創世記第1章 - 背景釋經 (陳崇基牧師)}
\newline
\newline
連結: \href{https://youtube.com/watch?v=l4PeWdLBiWQ}{\texttt{https://youtube.com/watch?v=l4PeWdLBiWQ}} ~~~~ 語音日期: 2022-01-18
\newline
\newline
\hyperref[sec:SlU_7mFzyiw]{\small{< < < PREV SERMON < < <}}
~
\hyperref[sec:index_chronic]{\small{[返順時目]}}
~
\hyperref[sec:index_scriptual]{\small{[返順卷目]}}
~
\hyperref[sec:rm7wwoQR_Zw]{\small{> > > NEXT SERMON > > >}}
\newline
\newline
$^{1}$我想和各位分享一下.
因為我們要有上帝的華語.
我們的生命教會.
神的子民才可以被建立.
我們作為牧者 作為領袖.
我們希望能夠更加明白.
更加容易掌握到聖經裡面的說話.
我們想由《創世記》第一章試一試.
如果有什麼問題.
稍後可以再討論.
這個分享大概只有20-25分鐘.
我們不會耽誤大家太多時間.
因為知道大家事奉都很忙.
主要我們想分享一些工具.
一些進路 一些方法.
以致大家看聖經的時候.
能夠幫助自己去明白.
並且幫助到我們的會眾.
我們的組員去明白.
就用《創世記》第一章試一試.
我經常叫弟兄姊妹去想像一下.
其實我們看聖經是好像去旅遊一樣.
其實這是一個進路.
為什麼看聖經好像去旅遊一樣.
如果你想一下去旅遊的話.
特別是去一個陌生的地方旅遊.
無論你是去另一個國家.
另一個城市.
那裡的風土人情.
那裡的文化是不同你本地的文化.
語言也可能是不同的.
所以你去準備一個旅遊的時候.
你去那個地方.
你能夠有所謂的得著.
希望不會浪費你的時間.
你需要了解當地的文化背景.
當地的語言.
當地人的習俗等等.
以致你去到那個陌生的地方.
你能夠經過這個旅遊的經驗.

$^{41}$你就明白當地的人說什麼.
其實看聖經也是這樣.
每次翻開聖經.
其實都好像去一個陌生的地方.
一個你可能完全沒有去過的地方.
因為那個文化背景.
是跟我們現代人的文化背景不同的.
正因為這個緣故.
所以我們每一次翻開聖經的時候.
其實我們在說讀經,解經.
都是有一點像一個旅遊的經歷.
是不是?.
所以我們做一個好的門徒.
或者做一個紮實的基督徒.
有堅固信仰根基的基督徒.
我們很需要去讀明本聖經.
所以讀經,識經和活經.
正是跟我們的生命有關係的.
所以我們很多時候看聖經的時候.
大家可能在網上查經.
跟小組查經的時候.
很多時候我們會去到一個地步.
只是上帝直接對我說話.
很感動我.
這種很主觀的去了解.
但是我們要做一個紮實的門徒.
我們就要懂得去到一些工具掌握了.
如何去看明本聖經.
能夠解明本聖經.
以至去應用在我們生命裡面.
能夠活出本聖經.
所以讀經,識經,活經.
三樣東西.
什麼叫識經呢?.
識經大家可能在神學院去過.
你知道就是exegesis.
就是去解釋聖經的意思.
解釋聖經這個字本身.
識經.
其實就是從經文引出作者的意思.

$^{81}$跟詩意解經或者靈意解經是不同的.
靈意解經就是將我們現代人的意思.
讀入了聖經裡面.
很多時候我們就會撞板.
我們就會不是很明白.
或者掌握不了聖經裡面的意思.
所以這些絕對不是一些所謂象牙塔.
裡面的學術.
其實我們每一天做基督徒.
我們都需要明白.
如何去讀本聖經.
所以是需要掌握某些工具.
某些進路.
所以就好像去另一個地方旅遊一樣.
我們明白語言.
明白背景等等.
所以我們需要了解聖經裡面一些本質.
當然聖經是神的話語.
所以在這裡我列出了五件事.
我經常跟弟兄姊妹分享.
就是五件事我們了解聖經到底是一回什麼事.
第一.
我們相信當然是聖經無誤.
因為上帝的話語是無誤的.
但是相信聖經無誤就不等於.
聖經是論盡天文地理的.
聖經不會說現代科學的東西.
不會說我們現代人.
很多時候要問的問題.
反而聖經是要搞盡我們應該問的問題.
很多時候我們翻一本聖經.
譬如《創世記》第一章.
我們問科學性的問題.
但是那些問題不是經文想去解答的.
所以聖經想去幫我們.
或者引導我們問一些更加正確的問題.
所以第一件事我們知道.
就算聖經無誤.
就算我們相信聖經是神的話語也好.
不等於聖經什麼都說.

$^{121}$所以為什麼我們通過了解聖經.
讀經,識經這些步驟.
以致我們能夠聽清楚作者.
神藉著作者.
想跟我們說什麼.
神藉著經文想跟我們說什麼.
另外第二件事我們了解.
就算我們相信聖經無誤的話.
神的話語.
不代表不保證我們比其他人.
我們分析真理或者了解真理更加準確.
對嗎?.
我相信真理.
耶穌是道路真理心明.
是不是?.
這是恩典.
不代表因為我相信了福音的真理.
所以我比其他人更加懂得去分析.
所以我們說回如何讀本聖經的時候.
我們需要了解一些基本的進路.
基本的方法.
可以這麼說.
第三.
聖經雖然已經翻譯成為一些現代的語言.
中文就是我們今天能夠一看就看得懂的語言.
現代的語言.
但是大家知道.
聖經是古代的文獻.
古代中東,古代近東的文獻.
所以正如我們去旅遊一樣.
到了另一個地方.
他們對世界的觀念.
他們對文化的觀念.
他們的習俗,價值觀.
跟我們自然有差距和差異.
所以我們需要了解聖經的背景.
第四.
就是聖經不是直接寫給我們的.
但是聖經是為我們每一代的信徒而寫的.
這個分別是很重要的.

$^{161}$因為我們翻開聖經的時候.
我們知道聖經不是直接寫給我們的.
因為我們經常在講道教道的時候.
我們都說聖經本身的原文是什麼.
希伯來文,希臘文.
是其他的語言.
是古代的語言.
今天不會再說的語言.
所以不是直接寫給我們的.
但是卻是為我們每一代的信徒而寫的.
是不是?.
所以我們掌握了聖經本身的原意是什麼的時候.
我們才能夠明白如何去應用.
以及上帝的懷語才能夠在我們的生命裡面.
有它的功效.
第五件事.
就是我們看聖經的時候.
不能夠抽離背景和上下文.
如果不是的話.
我們往往就會任意曲解.
所以在這裡.
我想用一個這樣的視野.
來去了解聖經.
我們去了解聖經的時候.
我們需要想三個視野.
第一個視野就是聖經背後有它的世界.
聖經的世界就不同我們今天的世界.
所以我們每次打開聖經的時候.
我們就進入了另外一個世界.
那個世界就好像我們去旅遊.
去一個陌生的地方.
那個風土人情習俗語言.
觀念是跟我們有所分別的.
所以我們先了解背後的世界.
我們才能夠了解聖經到底想說什麼.
另外一個視野就是聖經裡面的世界.
聖經裡面的世界是什麼意思?.
就是經文裡面.
上文下理.
給我們一個故事.

$^{201}$說出整個故事從頭到尾.
裡面的人物.
有什麼發展.
有什麼改變.
舉個例子很簡單.
福音書就是.
如果你看福音書.
彼得剛出現的時候.
他是一知半解的.
對耶穌是誰.
但是到約翰福音最後.
他就對耶穌復活了的耶穌.
死而復活的耶穌.
他就更加了解.
到底他來這裡做了什麼.
所以你看到彼得這個人的角色.
是有好像看電視劇一樣.
好像看電影一樣.
角色是不斷有改變.
他慢慢明白了.
從一知半解到最後他明白.
耶穌就是基督.
就是救世主.
所以中間跌跌碰碰.
這就是聖經裡面的世界.
亞伯拉罕也是.
亞伯拉罕開頭撰寫的十二章.
夢上帝呼召的時候.
他當然是信心去回應.
去到神所指示的地方.
迦林地.
但是.
他慢慢地每一章過去的時候.
他就有跌又碰.
有時候會懷疑.
有時候會跌倒.
然後最後他就發現.
原來真正的信心是什麼呢?.
就是獻耳殺的故事.
所以從第十二章到二十二章.

$^{241}$你就看到亞伯拉罕這個角色.
這個人物.
他的信心有一個歷程.
有一個改變.
就好像我們自己剛剛順著.
到我們現在應該有一個歷程.
中間有跌又碰.
上帝的恩典一直帶領著.
這就是聖經裡面的世界.
所以要看上文下理就是這個意思.
然後我們第三個視野就是聖經前面的世界.
就是有一本聖經放在你前面.
放在你桌面.
然後前面的世界就是我.
我自己.
我每次翻開聖經的時候.
我都不是白紙.
因為我會有我自己僅有的觀念.
或者是一些錯誤的理解.
現在我看了聖經了.
我前面的世界.
我就會改變了.
是不是?.
等於雅各書說什麼.
看到看聖經就好像照鏡子一樣.
照到自己的樣子.
然後我就有改變.
有一個更新.
我們希望我們自己看聖經.
我們帶領別人去看聖經.
我們帶領別人去認識神的位置的時候.
都是希望有一個更新.
生命裡面的改變.
但是是有步驟的.
所以有三個視野是需要留意的.
我舉個例子.
創世記第一章第一節.
他說:起初神創造天地.
地是空虛混沌冤冤黑暗.
神的靈混藏在水面上.

$^{281}$這個大家都知道.
起初中文.
背後就是希伯來文.
希伯來文就是這個字Beresheet.
我們不需要一定要懂得希伯來文.
而我想給大家一個例子.
我們想起起初的時候.
我們想起什麼呢?.
我們一想起起初.
聖經裡面第一個字.
Beresheet 起初.
翻譯做.
我們想起.
現代人想起時間上的其中一點.
就是第一點.
起初神創造天地一定對.
太初上帝在第一點突然間創造天地.
首先我們明白聖經背後的世界.
當他們講起初的時候.
希伯來文講起初的時候.
是不是我們現代人想時間的第一點.
如果你找起初這個字.
在舊約聖經裡面用過的.
有好幾處地方可以給個例子.
就是耶利米蘇26章27章28章49章.
這裡說起初這個字是同一個字.
Beresheet.
但是中文的翻譯.
因為要翻得中文的文法通.
字眼要通的時候.
我們就不翻做起初.
在耶利米蘇26章第一字.
那裡說猶大王約瓦敬.
這個人名.
這個王名.
登基的時候.
登基這個字就是起初這個字.
Beresheet.
突然間你退後兩步去想的時候.
原來我們起初想的是現代人的.

$^{321}$所謂起初就是時間的第一點.
原來在希伯來文的使用裡面.
起初不是時間的第一點.
而是時間的開頭那段.
最初的那段時間.
初期.
所以不是時間的一點.
而是時間的一段時間.
就是時間表上面最初的那段時間.
每個王登基的時候.
開頭他還沒真正坐上王位.
作王的時候.
那段時間.
父王死了.
然後新的王登基.
那段時間.
作王到他真正在新年的時候.
他真的坐上王位慶祝登基了.
好像美國總統那樣.
選了他.
然後隔了一個多兩個月之後.
才在白宮正式登基.
那段時間.
就是叫做起初.
就是西底加作王的初期.
起初是要做什麼呢?.
他就是成立他的內閣.
他就是成立他政府的部門.
他縮整他整個政府的權力.
整個架構.
他重新再整理過.
然後再安排過.
以致他新年的時候.
正式登基了.
正式作王了.
元年開始的時候.
第一年第二年開始.
數元年的時候.
他就能夠去管理統治這個國度.
但是如果你這麼看的時候.

$^{361}$創世記第一章.
說出神創造天地的時候.
就是在說一個國度.
我們現代人想.
就是想到宇宙好像一個機器一樣.
是不是?.
特別是唯物論的世界.
好像一個機器一樣.
萬物都是有自己的原則.
一直這樣運作的.
但是創世記第一章.
給我們一個神觀.
或者給我們一個世界觀.
是完全不同的.
原來創造主宰好像一個君王一樣.
起初的時候他做什麼?.
他就去整頓世界的秩序.
這個世界是屬於他的.
而且這個神創造主.
就是在這裡統治的.
所以如果這樣退後一步想.
我們不要用科學的頭腦去想.
不要用一個現代人的科學頭腦去想.
用聖經背後的世界去想.
他們想起初是什麼意思.
原來起初就是君王開始做事.
直到第七天.
他去到安息的時候.
安息也不是說他累了休息.
創世記第一章.
一到第七天.
第七天.
上帝安息.
原來在古代中東的世界.
說到那些廟.
做那些神廟的時候.
那個神靈進入聖殿的致聖所.
他去登基作王的時候就叫做安息.
所以你看到創世記第一天.
第二天,第三天.

$^{401}$一直數過去的時候.
到第七天他安息.
那個 Shabbat.
安息日.
他作王.
就真的完全是.
他的內閣.
他整個國度都整頓過了.
有條有理.
整整有條.
第七天他就坐在皇位裡面.
就是他的致聖所裡面.
他登基掌管這個世界.
所以你看創世記第一章.
其實是在說一個國度的成立.
原來整個宇宙就是上帝的國度.
就是上帝要去統治.
並且去接受敬拜.
所以我們不用一個現代人去看的時候.
我們發現.
這樣才看得通.
看得通聖經.
特別第一章.
創世記想說什麼.
如果你看回.
一二三四五六七日的話.
原文的字句真的要直接翻的話.
其實也不是跟你數.
第一天做完的事情到第二天.
然後再過二十四小時又做另外一天的事情.
因為這些原文.
你真的翻譯的時候.
你去到第六第七天的時候.
他才有英文叫做the.
就是中文的貫詞.
就是很特別.
很特定的那一天.
第六天他才加一個貫詞.
然後第七天就加一個貫詞.
所以嚴格來講.

$^{441}$一二三四五六七天.
就不是時序順序告訴我們.
上帝在做什麼.
OK.
所以那個強調的地方就是第六和第七天.
第六天做了什麼.
就是地上的萬物和人類.
人是按照神的形象所做.
然後第七天就正如我剛才所說.
上帝安息.
不是因為他累了做事很辛苦.
而是坐進去他的皇位.
他正式登基.
在他的聖殿的致聖所.
他的寶座上.
他坐著去為王.
去做王.
所以我們這樣想的時候.
我們就不會用一個現代人去看.
哦一二三四五六天.
我們整天就去辯論.
到底跟科學所說的話是不是對的.
是不是吻合的.
這個不是聖經要回答的問題.
正如如果你今天出去買東西.
你如果按時序的話.
你可能第一兩件事.
你去電腦中心買完一兩件.
然後第三件事你去超市買一些食物.
然後第四件事你去書局買一本書.
然後最後你經過一個地方買一個電池.
這個就是按照時間去編序.
去到說出來的.
一二三四五六七.
是不是.
但是如果你是按照主題.
去編寫的話.
希望你看得清楚.
按照主題編寫的話.
你回到家的時候.

$^{481}$你媽媽問你買了什麼.
你就說我在超市買了一些食物.
然後我在電腦中心買了什麼.
然後買了一本書和電子.
你就不是按照時序來講出來的.
你就按照什麼.
你就按照主題講出來.
《創世記》第一章.
其實就是想做這件事.
如果看《創世記》第一章到第二章三節.
就是頭一節.
就是頭七天.
你就明白.
它有一些文學的結構.
它是想表達一件事.
而且它想講一件事.
不是我們現代人科學要講的東西.
第一天就是做光與暗.
然後跟第四天做光體.
大的光體和小的光體.
和星星.
天上的星星.
第一天跟第四天跟光有關.
第二天就是天和海.
兩個空間.
第五天相對於天和海裡面的活物.
第三天所做的就是地和菜蔬.
跟第六天是平衡的.
因為講回第六天做的動物和走獸.
以及最終創造的高潮.
就是創造人類按照神的形象所做.
到第七天.
就是神安息.
神真的作惑.
所以整個的創造.
《創世記》第一章告訴我.
這個宇宙這個世界.
神所創造出來.
是上帝自己的角度.
按照祂的編排.

$^{521}$按照祂的安排.
祂到最後作惑.
人就在創造主之下.
代表祂來管理萬物.
所以你看看《創世記》第一章.
有很多文學的修辭在裡面.
特別是希伯來文化.
是有它的意思的.
比如說第一日至第三日.
和第四日至第六日.
是平衡的.
剛剛說了.
然後第一章第一節.
是有七個字的.
我講的是希伯來文.
雖然中文也是七個字.
起初神創造天地.
是不是?.
但是希伯來文也是七個字的.
如果你有機會你回去數一下.
你不需要很懂希伯來文.
但是你知道那七個字是分開的.
中間有一個空格分開讓我們看到七個字.
一章第二節有14個字.
一章第一節有28個字母.
第七日就有35個字.
然後過來這邊.
我們知道神這個字.
希伯來文L.
在第一章七日的創造出現了35字.
然後天與地各自出現了20字.
事就這樣成了.
就出現了七次.
慎看著是好的.
也出現了七次.
你看到這些數字.
絕對不是隨機放在一起的.
是很有意思地放在一起的.
就是有一個七的或者是倍數.
是不是?.

$^{561}$對於希伯來人來說.
七就是完美.
並且.
你想想聖經裡面.
有一些的做事.
創造的一些技術.
建造的技術.
都是七段的.
或者有七字出現.
想不想到?.
有一個最明顯的例子.
就是所羅門建造第一個聖殿.
列王記上六章第三十八節.
用了七年的時間.
是不是?.
所以說到有七的時候.
說到建造建立的時候.
就是建造聖殿.
所以.
創世的第一章.
似乎告訴我們.
整個宇宙就是上帝的聖殿.
有祂充滿的地方.
有祂作王的地方.
有上帝同在的地方.
就是上帝原本創造這個世界.
就是想跟我們同在.
是不是?.
有祂自己的榮耀.
有祂自己同在在裡面.
這個世界.
你看這個世界的時候.
不是純粹一個機器.
是一個神聖的地方.
就是上帝住在裡面的地方.
因為聖殿就是讓上帝住在裡面.
是不是?.
上帝同在的地方.
所以不僅是創世的第一章.
創造宇宙是七段.

$^{601}$還有很多七字出現.
所羅門建造聖殿.
也用了七年.
是不是?.
還有什麼是有七段的呢?.
就是回眸.
這是以色列一個考古公園.
叫做在庭拿的考古公園.
在曠野.
他們在公園裡面放了一個一比一的摩西回眸.
你看到有外園.
然後內園.
裡面有聖鎖是長方形的.
然後進去裡面.
就有致聖鎖.
致聖鎖就是正方形的.
正方形.
就是在近東的藝術裡面.
建築裡面.
正方形就是完美的意思.
但是如果你留意.
聖經裡面記載回眸的建造.
也是有七段的.
埃及的二十五至四十章.
記載回眸不是亂記載的.
是有計劃地敘述的.
二十五章第一節開始.
耶和華曉如摩西.
告訴他回眸裡面的建設是什麼.
陳設餅啊 燈台啊等等.
三十章十一節.
又出現一句話.
就是耶和華曉如摩西.
講到回眸裡面的捐獻.
然後十七節.
講到銅盆的時候.
十七節又有一句話.
就是耶和華曉如摩西.
然後講到高遊二十二節.
耶和華曉如摩西.

$^{641}$講到裡面的香.
所點的香.
三十四節.
又是耶和華曉如摩西.
三十一章第一節.
講到回眸的技工.
耶和華曉如摩西講什麼.
然後最後又是講到索伯.
講到安息人.
三十一章十二節.
但是開頭那句.
就是耶和華曉如摩西.
所以上帝.
叫摩西做回眸.
又是用七段的說話來跟他說.
七段的段落來跟他說.
然後最後.
做了回眸.
就到安息日了.
安息日就是敬拜那一天.
為什麼要敬拜那一天.
因為那一天.
上帝安息.
就是他作王的時候.
他做我們的君王.
他接受我們的敬拜.
是不是.
所以如果這樣想.
創世記第一章.
或者甚至第二章.
講到宇宙的創造的時候.
就不是講到一個唯物的機器.
就是講到一個神聖的地方.
整個宇宙.
是一個.
英文叫做Sacred Place.
整個宇宙.
是一個.
聖殿來的.
聖殿.

$^{681}$以前在耶路撒冷的聖殿.
就是上帝榮耀充滿的地方.
就是上帝同在的地方.
上帝去登基.
去作王的地方.
所以創世記第一章.
絕對不是回答我們.
一些科學的問題.
反而是要我們回答一個更加重要的問題.
在1977年的時候.
77年了.
以前我的一個神學院的老師.
叫做J.I.Peck.
他已經過世好幾年了.
是一個很出名的英國裔的.
一個神學家.
他在一本很舊的書.
I Want to Be a Christian.
就是我想做一個基督徒.
那裡說什麼呢?.
他說創世記一至二章.
是沒有提供關於.
這個宇宙怎麼形成的答案.
不是正在回答我們現代科學的問題.
今天.
可能有人.
會覺得.
這些是聖經的缺陷.
聖經為什麼不說我們現代科學的東西呢?.
說一些光學啊.
黑洞啊.
Big Bang啊.
大爆炸啊這些.
但是.
J.I.Peck 說什麼呢?.
他說其實如果.
以大圖畫來說.
以長遠來說.
這些我們今天對科學的沉迷.
其實是錯的.

$^{721}$因為科學的沉迷.
只是專注什麼呢?.
這個宇宙是如何運作的.
而我們忽略了.
就是誰是最重要的.
最重要的問題.
看到宇宙萬物這麼精密創造的時候.
這麼漂亮的時候.
我們問的問題不是How.
不是問它是怎麼來的.
什麼時候來.
多少千億年之前.
爆炸來也好.
怎麼來也好.
這些不是最重要的問題.
問的就是.
看到宇宙萬物的時候.
問Who.
誰啊.
那麼.
所以.
他說與其批評.
創世記這些經文.
沒有滿足我們今天現代人那種世俗科學的口味.
倒不如讓經文來責備我們.
很多時候我們希望聖經答的問題.
提供的答案.
都是已經墮落了的慾望.
因為我們只是想了解大自然.
而我們不想明白更加重要的是什麼.
是Who.
真正的是哪一個.
是認識萬有的主宰.
所以創世記一字二章的信息是什麼呢?.
他說.
他說你是不是看到大海.
就是在創世的敘述裡面告訴你.
窮滄天空.
太陽,月亮,星星.
你看到飛鳥.

$^{761}$魚兒.
你已經觀察到這些自然景象.
當中有很多植物.
大大小小的活物.
你只是因為這些奇妙.
複雜性.
而去讚嘆.
人類呢.
又擁有所有能力和技巧.
兩性之間產生的魅力等等.
有很多感情等等.
這些是太棒了.
但是.
最重要的是什麼?.
創世記第一章就是要介紹.
萬物背後的那一位.
值得我們敬拜的那一位.
就是創天造地.
作王的主宰.
所以兩章的經文.
J.I.Pekka最後結束說.
他說既然你已經享受了這些創造的藝術品.
現在呢.
你必須要跟那一位藝術家.
去握握手.
既然你已經聽過動人的音樂.
那才算.
聽了很偉大的音樂.
但是現在我要介紹作曲家給你聽.
所以創世記的經文.
是向我們彰顯.
Who.
就是創造主.
而不是創造界.
是要教導我們認識和敬拜這位神.
不是教導我們物理科學這些事情.
所以創世記第一,二章的目的.
就好像詩篇.
《約伯記》等等.
是一些讚嘆的詩篇.

$^{801}$所以創世記看完第一章的時候.
我也好.
早完也好.
他們一定要get the point.
就是一定要明白重點.
重點就不是說.
到底是不是按照那六天去做的.
跟科學的結論是不是一樣.
這些全部都不是重要的東西.
因為所解答的不是科學和物理的問題.
不是宇宙如何來如何形成.
而是宇宙的背後.
那位值得我們敬拜的創造主.
所以在這裡和你分享一下.
就是我們看聖經的時候.
用創世記第一章做一個個案.
我們要分三個地方來看.
哪三個地方.
就是我們需要了解.
聖經背後的世界.
聖經背後的世界就包括文化的.
文化的打錯字了.
文化的背景.
還有他們的認知環境.
是跟我們不同的.
第二就是要去明白.
聖經裡面的世界.
文學的題材.
結構.
剛才那些修辭技巧.
七字等等.
然後聖經前面的世界最重要就是我們.
我和我的會眾.
我和我的組員.
要去理解.
原來我們以往所謂的認受歷史.
我們了解創世記第一章.
可能就是現代人的看法.
這個就要改變了.
聖經前面的世界.

$^{841}$就是我這個世界.
就要改變.
我要來到上帝面前.
就是要去敬拜祂.
創世記第一章.
是一個敬拜的神學開始.
不是說科學.
而是說我們作為人.
就是需要去敬拜.
去相信.
去信靠.
這個創天造地的主宰.
在祂面前做佛人.
在祂面前敬拜和事奉祂.
這個就是上帝創造我們的意思.
所以最後有一些反思的問題.
可以給我和其他的組員.
或者我之後的學生去想.
讀經解經這一方面.
與旅遊有什麼關係.
有沒有相似的地方.
剛才我開頭的時候已經說了.
每次我們翻開聖經.
靈修也好,查經也好.
與別人解經講道也好.
其實真的要帶別人.
去到另一個陌生的地方.
去旅遊.
我就好像導遊一樣.
我作為一個組員.
作為一個講道者.
作為一個教師.
就是一個導遊.
告訴別人.
這裡你要看看.
為什麼呢?.
因為這裡你看清楚.
跟你以前看的東西.
你老家看的東西是不同的.
這裡你聽清楚別人說什麼.

$^{881}$他的語言是不同的.
我就好像一個導遊一樣.
告訴你.
在聖經裡面有什麼你需要看得到.
第二.
看《創世記》.
如果我們.
不去參考聖經當代的一些寫作.
其實創造這件事情.
創造世界.
不僅是希伯來人.
以色列人才會說的.
每一個當時的民族.
都有他創造的故事.
就是這個世界的觀念.
是不是?.
其實就算現代人也是這樣.
我們也有不同的世界觀念.
我們有唯物論.
我們有唯心論.
我們有很多其他的對世界的觀念.
這些世界觀.
宇宙觀.
但是看《創世記》的時候.
《創世記》就是要挑戰我們.
本身有的世界觀.
我們以為這個世界只是一個唯物的機器.
不是的.
《創世記》說的.
是有一個君王來作王.
這個地方.
我們看到這個世界.
是一個神聖同在的地方.
應該有他榮耀的地方.
我們敬拜就是開始.
將上帝的榮耀彰顯出來.
第三.
就是有什麼文學修辭.
剛剛我說的.
七.

$^{921}$十四.
二十八.
三十五.
這些倍數.
這些都是文學修辭.
一到第三天.
四到第六天.
大家做完應該要去試試.
觀察到這些特別的地方.
是不是?.
以致他們能夠知道.
到底聖經想做些什麼.
第四就是第一章.
有什麼神學重點.
今天跟我們敬拜生活有什麼關係.
有什麼應用.
剛剛我已經說了.
今天也有什麼地方影響我們.
不能夠去敬拜神.
包括我們.
以為這個世界只不過是一個很大的機器.
這個宇宙只不過是一個很大的機器在自己運作.
這個是不會幫助我們敬拜的.
會讓我們.
不認識這個世界背後的創造主.
我們想想有什麼方法.
去幫助我們.
真正去敬拜神.
OK 我就暫停了.
我不想我職業病一直說就停不了.
所以在這裡.
跟大家分享了.
希望今天可以提醒一些鼓勵.
大家在這裡彼此學習.
或者大家有什麼問題.
可以在這裡.
我們可以交換意見.
分享一下.
如果你喜歡我們的YouTube頻道.
就記得.

$^{961}$Like.
Comment.
和Subscribe.
\newpage



\section{創世記}
\label{sec:rm7wwoQR_Zw}
\textbf{創世記第2章 (2.2) - 背景釋經 (陳崇基牧師)}
\newline
\newline
連結: \href{https://youtube.com/watch?v=rm7wwoQR-Zw}{\texttt{https://youtube.com/watch?v=rm7wwoQR-Zw}} ~~~~ 語音日期: 2022-01-25
\newline
\newline
\hyperref[sec:l4PeWdLBiWQ]{\small{< < < PREV SERMON < < <}}
~
\hyperref[sec:index_chronic]{\small{[返順時目]}}
~
\hyperref[sec:index_scriptual]{\small{[返順卷目]}}
~
\hyperref[sec:yR3tIvy_FRE]{\small{> > > NEXT SERMON > > >}}
\newline
\newline
$^{1}$所以我們看聖經.
創世的第一章.
26-27節.
他說我們要照著我們的形象.
看著我們的樣式做人.
使他們管理萬物.
全地.
地上所有的東西.
27節.
神就照著自己的形象做人.
但是照著他的形象做男做女.
這裡先解決了幾個問題.
第一就是一個宗數的問題.
有時候我們會說為什麼神會說我們呢?.
是不是暗示他三位一體呢?.
是不是暗示說有其他神在這裡呢?.
絕對不是.
因為希伯來的原文.
有宗數的時候不等於一定有幾位.
有宗數的時候.
很多時候是一種尊貴的宗數.
就是.
主這個字.
有時候在希伯來的聖經裡面出現.
是宗數的形式出現的.
但是文法是形式而已.
其實也是在說一個主.
所以這裡說我們的時候.
不是說真的有很多位神.
而是說一種尊貴的字術.
好像美國的總統.
說我們.
我們跟另一個國家開戰才算.
他說為什麼我們呢?.
明明白宮只有他一個人說話.
原因就是他代表他的內閣.
是不是?.
正如上帝.
是君王的話.
就有所謂他的天庭.

$^{41}$就是有他的內閣.
就是他的天使.
他的使者等等.
他就說我們.
所以這種是尊貴的宗數.
始終都是一位神.
是不是?.
因為.
二十六節開始的時候說什麼呢?.
神說.
說字的動詞是單一的動詞.
不是宗數的動詞.
第二.
更加需要處理的.
就是人這個字.
人這個字.
其實這個字.
本身是有貫詞的.
中文翻譯是翻譯不出來的.
英文可能也翻譯得到.
就是The Man.
加上The的時候.
就代表是一群人.
人類.
不是一個人.
也不只是男女兩個人.
而是做人類出來的.
是不是?.
一個人口.
一個宗數.
有一個The Humanity.
人類出來的.
如果你查一下.
聖經特別頭幾章.
用到人這個字的時候.
你就知道.
這個人字.
我們很多時候.
有時候會翻譯成阿當.
阿當其實是希伯來文.

$^{81}$所以大家知道.
我們其實是懂一點點希伯來文的.
阿當這個字本身就是.
希伯來文.
Adam.
意思就是人.
的意思.
這裡有一個希伯來聖經的作者.
做了一些同音字.
似音的字.
發音相似的.
就將它們放在一起.
人這個字.
Adam.
其實像什麼呢?.
大地.
土地.
Adam.
是從那裡出來的.
所以人和地.
他有暗示.
表述了人和地之間的關係.
人和地很有關係.
塵土歸於塵土等等.
如果你看完整個特別頭三章.
提到阿當這個字.
Adam這個字的時候.
很多時候.
不是提到一個人.
而是提到整個人類.
所以阿當這個字.
有好幾個意思.
首先.
地名可以是阿當.
你記得約書亞進入迦南的時候.
他們拿著一個藥櫃.
然後祭司的腳踏進入約旦河的時候.
水就在遠處北邊停下來.
因為在阿當城那裡.
約旦河的水就阻截了.

$^{121}$那個就叫做阿當.
地名而已.
第二個就是人名.
到第四章尾的時候.
才有一個人名.
叫做阿當.
但是很多時候是加一個貫詞的.
加一個貫詞的時候.
就不是指一個人和一個人名了.
是不是?.
譬如你說我的英文名是 Mark.
你不會說 The Mark.
你不會加一個 The 給我.
因為你在說一個人名的時候.
你就不需要加貫詞.
但是.
在聖經的用詞裡面.
往往用到 Adam 這個字.
就是會加上一個貫詞.
換句話說是說一群人.
或者一個代表.
而這個代表不只是一個人這麼簡單.
它是代表了很多的人.
所以我們中文翻譯得好一點的話.
其實應該翻譯做人類的.
人類.
就是上帝創造人類.
按照他的形象.
做人類.
有做男做女.
人類裡面有男也有女.
所以大部分好幾次出現的時候.
都是有那個貫詞.
那個貫詞就是 Ha.
H-A.
有其他兩次是沒有貫詞的.
但是上文下你告訴我們.
其實它的意思是有貫詞的.
譬如說二十六節.
和二十七節.

$^{161}$等等.
我們試一試看看.
我知道不是每個人都讀過希伯來文.
但是其實不是這麼難的.
我們試一試.
這裡我們試一試跟.
特別是二十六節和二十七節.
頭一節.
不是整段的二十六節.
也不是整段的二十七節.
二十六節的開頭和二十七節的開頭.
那麼這裡怎麼看呢?.
希伯來文是從右至左來看的.
所以頭這一行.
二十六節.
第一行就是希伯來文.
第二行就是英文.
用英文來將它發音拼出來.
第三行就是逐字逐字的.
平衡的.
用中文來翻譯出來.
OK?.
所以不是這麼難明白的.
每一節我們這裡列出.
就是有三行.
第一行是什麼呢?.
希伯來文.
第二行是英文的譯音.
第三行就是中文的.
就是逐字的翻譯.
OK?.
.
是什麼意思呢?.
就是他說的意思.
誰說的呢?.
希伯來文.
很多時候是動詞先的.
然後主格才出來.
第二個字.
就告訴你誰說的.

$^{201}$這個就是什麼呢?.
Elohim.
我們英文也叫God.
中文也叫神.
或者上帝.
OK?.
所以Fireman Elohim.
神說什麼呢?.
.
讓我們做.
讓我們創立創造.
.
創造什麼呢?.
Adam.
人類.
.
意思就是說.
照著我們的形象.
最後一個.
.
就是我們的.
就是這個意思.
照著形象.
我們的形象.
然後呢.
.
照著我們的樣式.
你也可以做什麼呢?.
人.
Right?.
OK?.
所以26節我們知道了.
原文就是這樣.
27節了.
就是那個作者.
那個敘述者.
就說回到底發生什麼事情.
上面說什麼呢?.
Fireman Elohim.
Nasir Adam.

$^{241}$就是說.
神引述創造主說了什麼.
讓我們做人類.
然後到27節.
他就去到報告.
告訴我們.
神做了什麼.
然後這裡第一個字就是動詞.
好像26節開始的時候也是動詞.
希伯來文很多時候都是動詞開始.
Fireman.
就是說他做了.
他創造.
主格什麼時候出來呢?.
通常第二個字就是主格.
就是Elohim.
就是神.
上面你看到有Elohim.
這裡也是Elohim.
然後就有一個這樣的符號.
英文.
翻譯出來的讀法就是at.
其實是翻譯不出來的.
是一個客體的標記.
就是說.
神做了什麼.
就給了一個客體標記給你.
然後告訴你做了什麼.
客體做了什麼.
就是Hadam.
這裡就是我所說的.
有貫詞.
Ha.
這就是希伯來文的原文文法.
裡面有一個叫The.
Ha.
所以你真的直譯的話.
就是The Man.
但不是說一個人.
是說一群人.

$^{281}$人類.
Fire ground Elohim at Hadam.
上帝做人類出來.
是照著他的形象和樣式.
是不是?.
好像上面那樣.
所以兩節.
頭一半.
26節頭一半.
27節頭一半.
想說什麼呢?.
就是上帝做人類.
我們經常以為.
或者小童兒童的聖經.
告訴我們.
創世記開始的時候.
只有一個男人和一個女人.
不穿衣服.
然後躲在樹林裡.
這是沒錯的.
但是.
聖經你看回原文的時候.
其實是做了一群人出來.
第一章26-27節.
上帝做了什麼?.
Hadam.
The humanity.
人類.
上帝創立了人類出來.
所以我們看回亞當的時候.
我們就要了解了.
聖經說亞當的時候.
不是說一個人.
一個角色.
他的故事.
不是在描述一個人.
當然歷史上有亞當這個人.
但是其實亞當的故事是在說什麼呢?.
是在說人的代表.
整個人類的代表.

$^{321}$到了第二章的時候.
我們看到那人.
和合本譯做那人.
也算是算對的.
因為那人是什麼意思呢?.
Hadam.
就是The man.
而這個人是代表.
是一個典型.
代表了整個人類.
就是說在亞當身上就看到每一個人.
在夏娃身上看到每一個人.
所以亞當夏娃這個故事.
不只是他們獨一的故事.
聖經想說.
主子要說什麼呢?.
亞當夏娃.
說的是每個人的故事.
英文叫做Archetype.
就是一個典型.
一個representative.
一個代表.
所以我們想起亞當的時候.
不是想起一個人這麼簡單.
我們想起一個人.
他有一個典型.
是代表很多人.
代表整個人類.
這些東西在文學上經常出現的.
這種手法.
就是用典型來代表.
來描述一些事情.
而他描述的時候.
雖然他只是說一個人.
但是其實他代表很多其他人.
舉個很簡單的例子.
母親.
你去問問小朋友.
或者你在班房裡面問問小朋友.
怎麼形容母親呢?.

$^{361}$他們會怎麼形容?.
或者很多其他人會怎麼形容母親呢?.
說她雙手是龜裂的.
說她能進廚房又能出廳堂.
小朋友可能覺得母親是白雪公主.
甚至有些人會形容母親是風箏的扇.
是不是?.
扯著她.
讓她不會飛掉.
或者形容母親是囉囉嗦嗦.
說個不停的.
他所描述的母親.
不是一個獨特的母親.
他所描述的是一個典型的母親.
全世界世上只有媽媽好.
媽媽是怎麼樣的?.
母親.
就是雙手龜裂.
因為工作很辛苦.
洗碗,燙衣服,煮飯.
手會裂.
能進廚房,能出廳堂.
是不是?.
很可愛的.
囉囉嗦嗦的.
就是在說一個母親的所謂 Archetype.
代表.
或者.
我跟你說一個故事.
我說從前有一個人.
他的品行沒有什麼修養.
說話的時候很大聲.
沒有禮貌.
沒什麼風度.
有時候會擺架子.
你說那個人叫什麼名字?.
那個人的名字叫做中國人.
你就知道我說的不是一個人.
我是在說一群人.
一個典型.

$^{401}$是不是?.
所以你看亞當的故事也是一樣的.
沒錯有歷史的亞當.
但是經文的主旨想說什麼呢?.
就是人類的代表.
人類的代表.
所以我們看亞當的描述的時候.
《創世記二三章》.
其實是在看我們自己.
每一個人.
都有一個這樣的描述.
就是人的墮落.
不只是亞當做了一些事情.
搞到整個人類都砸了.
讓整個人類全部都墮落了.
不只是他.
每個人都是這樣.
所以你想想.
《創世記》第一章.
第二章.
說的是哈阿達.
The man.
人類.
一個代表.
《創世記》第一章開始已經有一群人了.
但是第二章鏡頭就放大了.
焦點就放在一個人身上了.
那個我們叫他亞當.
那人.
其實就是全人類的代表.
你這樣的話.
你解決了不少我們以往.
自勉識經.
引申出來的一些問題.
《創世記》第四章十四和十七節.
這裡說什麼呢?.
該人.
你記得嗎?.
該人.
他的祭物上帝.

$^{441}$不去接納.
然後他就發怒.
殺了他的弟弟和父親.
十四節.
他要離開耶和華.
他就投訴什麼呢?.
他說如果有人看到我的時候.
就會殺我.
我會害怕的.
我們平日的自勉識經.
我們就說哪裡會害怕呢?.
有什麼人在呢?.
因為.
內裡只有三個人而已.
有一個人已經死了.
第四個人已經死了.
阿伯死了.
只有亞當,夏娃和該人.
何來會害怕呢?.
然後第四章十七節.
突然間說什麼呢?.
就說該人.
有一個娶了妻子.
妻子是誰呢?.
我們平日的自勉識經.
我們就說一定是兄妹.
兄妹.
所以我們就會說.
以前兄妹是可以的.
接著摩西律法來的時候就不行了.
亂倫就不行了.
其實是不是真的這樣呢?.
剛剛我們一直看原文的時候.
很仔細地讀的時候.
我們就知道這樣解釋是不需要的.
因為一開始已經有什麼呢?.
一章二十六二十七節.
上帝創造人類.
有男有女.
按照他的形象所做.

$^{481}$已經是一群人了.
然後第二章說什麼呢?.
鏡頭就焦點放在一個人身上.
而且這個人不是一個獨特的人.
而是一個人類的代表.
他是Everyman.
他是每個人的故事.
所以說到二章裡面.
哪一個人呢?.
他的描述等等.
他的角色等等.
是在說每個人.
原創的時候.
上帝給我們的心意是什麼.
Again.
再說一次.
有沒有歷史的亞當?.
當然有.
是不是?.
但是經文的主旨不是要處理這件事情.
所以你不可以做一個亞當的人物的.
角色的查經.
這樣是不行的.
因為亞當.
他原本的意思.
就是想告訴我們.
這是每一個人.
你看亞當夏娃的時候.
我們每一個人都在那裡.
亞當所謂墮落.
吃禁果.
姿善樂事的果子.
每個人都是這麼做的.
是不是?.
我們就在亞當裡面找到我們自己.
所以如果你這麼看的時候.
不僅沒有引申很多無謂的.
這些護教的問題.
那個人的妻子在哪裡來的.
為什麼他之後又見了一個城.

$^{521}$沒有人見什麼城.
等等.
這些問題.
因為我們平時看得不仔細.
仔細看的時候你就知道.
亞當亞當.
是夏亞當.
是人類.
The humanity.
發生什麼事情.
人類.
就是上帝原創人類的角色是什麼.
然後第三章發生什麼事情.
每個人的故事.
你再比對一下古近東的文獻.
其實都試探討同一方面的問題.
所有這些所謂創世的外邦神話.
近東的記載.
全部都試探討不是科學的問題.
也不是什麼時候有第一對人出現.
不是這些問題.
不是人類學歷史的問題.
而試探討人類身份.
角色.
他的本位是什麼.
所以這些近東的文學.
是給我們一種人觀.
你看看以色列以外的外邦.
巴比倫等等.
他們記載人類的始源的時候.
他們想解決什麼問題.
或者想回答什麼問題.
就是人類本身.
出現.
在這個世界宇宙裡面.
他的身份是什麼.
這種是近東的人觀.
很多時候這些近東的神話.
這些文獻.
講到人.

$^{561}$被創造的時候.
都是神靈的奴隸.
就等於我們中國民間的信仰.
神明.
土地也好.
誰都好.
初一十五的時候.
你就需要給他吃一隻白切雞.
因為那些神靈會肚餓.
那些神靈肚餓.
我們人就做什麼?.
做他的奴隸.
我們就要去肚餓他.
所以我們安事安候.
我們不要激怒他.
是不是?.
我們就要奉上去.
是不是?.
就是給他吃一些東西等等.
人在近東的觀念裡面.
都是一樣.
是做神靈的奴隸.
接著呢.
很多時候.
都是那些神靈打完之後.
馬道克打完了.
贏了其他水神.
他做人的時候怎麼做呢?.
就殺死一個邪神.
用他的血.
來做出來.
這些不是一些科學生物的表述.
記得C.S. Lewis說什麼?.
有很多其他方面語言的表述.
絕對在古代聖經世界裡面.
人們不是用邪神的血.
是不是紅色流出來的血.
主要不是說這件事.
不是說科學生物的問題.
而是說人的本位.

$^{601}$人的身份問題.
所以用邪神的血所做的.
所以人是邪惡的.
是不是?.
人本身就是邪惡的.
跟聖經的人倫很不同.
或者你可以這麼說.
《創世記》一二章.
說到人倫的時候.
論到人.
人性.
跟近東的其他觀念.
這些人生觀.
這些世界觀.
就是作爭辯.
就是告訴你不是這樣的.
真正人的本位是怎麼樣的呢?.
所以是在說一種什麼呢?.
叫做世界觀的問題.
World View.
我們每個人.
古代也好.
現代也好.
我們都需要面對.
解答這些問題.
這就叫做世界觀的問題.
我們一出生.
在找人生意義.
為什麼我會在這裡.
這些全部都是世界觀的問題.
World View的問題.
無論今天科技有多發達也好.
都跑不掉.
我們都要問這些問題.
要解答這些問題.
就是我在哪裡.
我是誰.
為什麼我會在這裡.
是不是?.
為什麼我會存在呢?.

$^{641}$這些是人生觀.
世界觀的問題.
這是最基本的問題.
是不是?.
世界有什麼問題.
然後呢.
這些問題.
發生什麼事.
有什麼答案.
這些都是世界觀的問題.
而且解答的方法在哪裡.
有什麼方向可以解答.
現在是屬於什麼時間.
是不是?.
等等.
這些都是世界觀的問題.
就是我們平日.
我們需要去.
最基本要去到.
解答的問題.
我們最基本要解答的問題.
不是世界的問題.
不是科學的問題.
而是自古到今.
我們都要解答這些世界觀的問題.
是不是?.
我們為什麼會在這裡呢?.
聖經其實就是想解答這個問題.
就是告訴我們一個正確的世界觀.
是不是?.
從創世開始跟我們說.
創造人類開始說.
上帝創造人類到底本身有什麼心意呢?.
這些就是聖經要解答的問題.
也是最重要最基本的問題.
在這裡舉一個例子.
美國.
在1977年9月5日的時候.
發射了一個叫做航行者一號.
Voyage One.

$^{681}$射了上太空.
然後它一直向著深的太空進發.
77年9月5日.
我們中間可能有些還沒出生.
77年9月5日就發射了.
到現在還在飛.
一路離開地球.
一路進發.
已經超越了我們的太陽系.
現在還繼續向著深的太空進發.
在1990年的時候.
2月14日.
就是剛剛航行者一號離開太陽系的時候.
還可以看到地球.
還可以看到我們現在生活的地球.
航行者一號.
美國太空總署就叫他拍一張照片.
Voyage Number One.
航行者一號就回頭拍了一張照片.
在他後面的鏡頭拍了一張照片.
拍了一張照片看到地球.
看到地球你看到嗎?.
這張就是那張照片.
這張底相.
就是這一點.
這張照片很出名.
叫做Pale Blue Dot.
就是淡藍點.
在一個浩瀚的宇宙裡面.
我們是什麼?.
Who are we?.
就是我們為什麼存在?.
我們為什麼在哪裡?.
原來只是一個浩瀚的宇宙裡面一丁點而已.
之後再航行者繼續飛的時候.
已經再看不到地球了.
我為什麼說這張照片?.
就是一個現代人的世界觀.
我們為什麼在這裡?.
不知道.

$^{721}$只是意外而已.
我們存在沒有目的.
只不過是一些物質的東西.
唯物論就是說這件事.
我們為什麼在這裡?.
我們沒有形而上的目的.
我們人生沒有本有的目的.
我們無緣無故就存在了.
廣東話叫做石頭爆出來.
為什麼我們在這裡?.
這個Pale Blue Dot.
其實就是在表達現代人.
特別唯物主義的世界觀.
為什麼我在這裡?.
沒有意思的.
為什麼我會存在?.
沒有意義的.
我只不過是一些物質撞擊.
隨機的結果而已.
是不是?.
我就在這一點裡面.
這是世界觀.
是不是?.
唯物論是世界觀.
聖經要表達的世界觀.
是截然不同的.
就是說創造主.
有心思的創造我們.
以愛創造我們.
以他的形象來創造我們.
所以我們在世生存.
是有尊貴的遺憤的.
是可以與上帝建立關係的.
這個世界觀是完全不同的.
這就是我們福音裡面的信息.
是不是?.
我們跟別人說耶穌基督的福音.
其實就是告訴他們.
我們要轉變我們對世界的觀念.
對我們人生的觀念.

$^{761}$與這位創造主.
建立應該有的關係.
而且這個關係的方法已經給了我們.
途徑就是耶穌基督的恩典.
他的救恩.
我們就拿回聖經裡面告訴我們.
這個正確的世界觀.
所以漢世紀第一二章.
說的是什麼?.
World View.
也不是說一些抽離的東西.
雖然不是解答現代科學的問題.
而是解答一些世界觀的問題.
現代.
行者也好.
唯物論.
科學.
給我們的人生觀.
也是解答的這個問題.
就是說不是這樣的.
應該是跟這位創造主.
再重新結立關係.
這個才叫做生命.
是不是?.
真正有的生命.
所以在創世記第一章.
我們繼續看下去.
二十六節.
剛剛說的 Adam 或者.
赫 Adam.
說到創造人類的時候.
上帝創造人類.
是做什麼的?.
是按照上帝的形象.
上帝的樣式來創造.
以往在我們二千年的教會歷史裡面.
有很多這樣的討論.
就是形象這個字.
以及樣式這個字.
當然是指不同的東西.

$^{801}$可能形象是更加具體一點.
樣式就不是.
會不會這兩個字.
分別指可能是.
神創造人的時候.
我們就像神.
像神什麼呢?.
可能形象就是指.
他的知情意等等.
就是神的一些本性.
我們是像的.
或者可能這樣的解法.
也不是最好的.
我們最近.
或者近代我們找到的是什麼呢?.
也不算近代.
1979年的時候.
在敘利亞北邊.
有一個古城那裡.
就作了這個東西出來.
這是一個石像.
是一個君王的石像.
是屬於主前八百多九百年的時間.
這個石像你看到下面.
是說這個君王.
他有一些黑文在這裡寫著.
有一些冥文.
那些冥文當然很小.
但是那些冥文其實是什麼語文呢?.
其實就是亞述文.
以及亞蘭文的翻譯.
這個石像說到.
他指到這個石像.
這個人像.
亞述文就是Sawu.
但是在另一邊.
有亞蘭文的翻譯.
因為當時主前十世紀九世紀的時候.
當時的國際語言.
其中一樣就是亞蘭文.

$^{841}$Alame.
聖經裡面其中有一些章節都是亞蘭文.
有冥文就是雙語的翻譯.
亞蘭文怎麼翻譯亞述文呢?.
就用了兩個字.
Sawu.
以及Dumut.
這兩個字正正是跟希伯來文的兩個字.
形象.
Sawu.
以及樣式.
Dumut.
就是同一個字.
換句話說這兩個字其實是同義詞.
是不是?.
這裡的人像.
當時他翻譯同一個字.
就用兩個字.
隨便是同用的.
意思是一樣的.
所以形象,樣式其實是同樣的東西.
是同義詞.
這就是近東的背景幫助我們.
不需要去掙扎很多細微的問題.
原來在說什麼呢?.
形象,樣式是同義詞.
形象,樣式是同義詞的話.
我們就問.
按照神的形象所做.
是什麼意思呢?.
形象就好像石像一樣.
就是.
很多時候軍王.
他到了.
統治了一個地方的時候.
或者他佔領了一個地方的時候.
他就會將他的石像放在那裡.
統治了另外一個地方的時候.
他就會將另外一個石像放在那裡.
就代表他的同在.

$^{881}$代表他的統治.
他的權.
他佔領了一個地方.
他就放一個石像在那裡.
代表這個就是我了.
我就在那個地方.
去管治那個地域.
所以人類是神的形象.
就是說.
神創造人.
人在地上.
就有一個君尊的地位.
代表上帝.
在上帝的主權之下去管理萬民.
所以我們人雖然微小.
但是我們卻是有上帝給我們的尊貴的身份.
詩篇第八篇你記得嗎.
人不算什麼.
你竟顧念他.
是不是.
比他做的時候.
比天使還微小一點.
多有尊貴的身份.
這個就是世界觀的問題.
不是科學的問題.
不是科學的表述.
是說人有尊貴的身份.
好像軍王一樣.
多謝大家收看.
如果你喜歡我們的YouTube頻道.
就記得Like,Comment和Subscribe吧.
\newpage



\section{創世記}
\label{sec:yR3tIvy_FRE}
\textbf{創世記第2章 (2.3) - 背景釋經 (陳崇基牧師)}
\newline
\newline
連結: \href{https://youtube.com/watch?v=yR3tIvy-FRE}{\texttt{https://youtube.com/watch?v=yR3tIvy-FRE}} ~~~~ 語音日期: 2022-01-25
\newline
\newline
\hyperref[sec:rm7wwoQR_Zw]{\small{< < < PREV SERMON < < <}}
~
\hyperref[sec:index_chronic]{\small{[返順時目]}}
~
\hyperref[sec:index_scriptual]{\small{[返順卷目]}}
~
\hyperref[sec:oVRXAcV7BoQ]{\small{> > > NEXT SERMON > > >}}
\newline
\newline
$^{1}$還有,你講到神的形象的時候.
近東的文獻經常都會提到神的形象.
這句話是什麼人說自己是神的形象.
我們看看亞述的文獻.
我這裡列出了三句.
三段亞述的文獻.
第一句.
我主的父王是比勒.
Bell.
你記得Bell是誰嗎?.
馬獨克.
巴比倫的大神.
我主的父王是馬獨克真正的形象.
但是我主的皇帝,皇上.
也是比勒真正的形象.
就是說誰是神的形象.
就是皇帝王才是神的形象.
第二句.
這句話是其他人稱讚過皇上說的.
皇上是萬國的主.
是太陽神沙瑪斯真正的形象.
誰是神的形象?.
天子就是神的形象.
皇上才是神的形象.
第三句.
更有趣.
人是神的影子.
而僕人就更加不是神的影子.
只不過是人在底下的影子.
草民.
低等的人.
但是皇上是誰?.
以撒哈頓是亞述王.
其中一個王.
以撒哈頓皇上就是神的形象.
Image of God.
原來是說皇帝才是神的形象.
埃及也是一樣.
埃及帝國的一些文獻.
萬有的主是閃耀的形象.

$^{41}$太陽城.
埃及有一個宗教的中心叫太陽城.
主神所創造的.
他生了他的王.
就是法老.
在地上做出他閃耀的榮耀的後裔.
為了人類的拯救.
法老是誰?.
就是活著神活著的形象.
Image of living God.
當然是那些假神.
然後另外一句.
埃及其中一個神明.
當時的愛邦是多神論.
阿蒙形容法老是什麼?.
你就是我活著的形象.
我看的創造.
這幾句話已經足夠讓我們知道.
到底聖經說出神的形象這句話的時候.
當時的人在想什麼?.
當時的人在想.
神的形象一定只是在上掌權.
上面那些有尊貴身份.
有神聖身份.
天子嘛我們中國人都這樣說.
祂才是Son of God.
祂才是神的兒子.
創世的第一章.
就說了一些很嶄新的理念.
一些很革命性的理念.
就是說全人類都是神的形象.
神的形象不是在上掌權者.
尊貴身份不只是在上掌權者的專利品.
所有的人都是神的形象.
所以你看到第一章.
原來你看清楚原文等等背景.
你會知道原來他在宣告一些東西.
不只是古代的時候才能夠應用.
今天也能夠應用.
就是世界觀的問題.

$^{81}$就是人是什麼.
人算什麼等於尖帝八鞭.
是有尊貴的身份.
君尊的地位.
好像王一樣.
在這位創造主君王底下.
我們就是創造界的君王.
有一個很尊貴的身份.
跟當時近東的人觀去作一些爭辯.
跟現代今天所想的人觀.
很多時候假設的人觀都是作一些爭辯.
今天我們不知道石頭爆出來.
人算什麼不知道.
創世記就是想表述一個尊貴的人觀.
人很有價值.
人很有尊貴的地位.
代表上帝是管理萬有的.
管理我們生活裡面所接觸的所有東西.
這是一個很有價值的觀念.
並且不是在上掌權者才有這樣的神聖地位.
所有的人.
所以單單那句話.
已經想挑戰當時的文化.
而去創立一種新的概念.
告訴我們一個正確的概念.
每個人都有權.
每個人都有尊貴.
有價值.
在神面前有價值.
這是第一章.
到第二章了.
再看清楚第二章想說什麼.
首先第二章指到神的時候.
就不單單是阿路歇.
就說到耶和華神.
Yahweh 阿路歇.
經常說Yahweh 阿路歇第二第三章.
為什麼呢.
因為聖經的作者.
你當他是摩西.

$^{121}$他寫的時候.
其實就是告訴你創造主是誰.
就是等於以色列和以色列文立約的神.
Yahweh.
耶和華這個字.
就是一個個人的名字.
一個個人的名字.
不是一個稱銜.
神就是一個稱銜.
有一個稱銜.
但是說到耶和華的時候.
就是上帝自己的名字.
正如我們今天說的.
稱銜就是牧師.
人們不太熟悉就叫我牧師.
但是如果那個人和我熟悉.
就會用我個人的名字來叫我.
叫我Mark Chan.
陳崇基.
叫我個Personal Name.
就是和我特別有一個個人的關係.
立約的關係.
耶和華就是一個Personal Name.
就是和神的子民有一個立約的關係.
Yahweh.
第二就是伊甸.
我們有很多對伊甸園的想像.
但是你要知道.
裡面有什麼意象.
是想表達什麼.
似乎不只是表達一些實物這麼簡單.
所以我們經常問.
伊甸園本身在哪裡.
在地球的哪處.
其實聖經不是要回答這個問題.
聖經是用一些當時很熟悉的一些意象.
一些意境.
一些術語.
來表達到底伊甸是什麼.
人在伊甸是什麼意思.

$^{161}$第一方面.
在近東的背景裡面.
所有的廟.
所有的神殿.
都是有河水的.
剛才我們說的人像.
甚至用了一個字.
是Eden這個字.
是動詞Eden.
伊甸這個動詞.
一撇DN.
其實就是Eden.
動詞是在描述什麼呢.
看上面下的翻譯.
應該是灌溉的意思.
使之豐盛的意思.
所以Eden這個字就是豐盛的意思.
生命應該是豐盛的生命等等.
上帝原本創造人.
放人類.
不只是亞當.
因為亞當是代表全人類.
每個人都應該在Eden裡面.
就是上帝的殿裡面.
所以所有的神殿.
都是有河水的.
所有的皇宮.
神殿.
神殿和皇宮很多時候都是放在一起.
都是有河水出現的地方.
神殿皇宮.
都是水泉的地方.
另外神殿也有金銀寶石.
在創世的第二章當然有.
什麼馬路石啊.
寶石啊金啊銀啊等等.
說到Eden的時候有金銀寶石.
說到Eden的時候也有一些活物.
所謂天使.
基路柏.

$^{201}$你記得第三章.
創世記第三章.
把守生命樹道路的就是基路柏.
有發火焰的劍.
把守著生命的道路.
所以有這些天使活物.
基路柏.
是當時一些我們叫composite creature.
你記得勝經裡面有不同的章節嗎?.
我一想起一些活物.
譬如四活物.
四活物我們今天現代人看.
會覺得很奇怪.
你記得以西結書第一章.
說到四活物的時候.
四活物其實就是.
承載著上帝的座駕.
這個叫四活物.
四活物有什麼樣式.
有什麼形象.
以西結書第一章.
就是有人的臉.
獅子的臉.
牛的臉.
和鷹的臉.
為什麼會將這四個大地的活物放在一起呢?.
其實就是表達權能.
表達什麼叫power.
因為人一定是萬物之靈.
最有權能的.
獅子就是野獸裡面最大的.
鷹在天空.
天空範圍裡面.
領域裡面最大的.
牛就是加速裡面.
最有力的.
這些全部都是近東社會裡面.
一想起這件事的時候.
牛啊鷹啊獅子.
就會想起權能.

$^{241}$想起power.
所以說到基路伯的時候.
形容基路伯.
形容天使的時候.
承載著上帝的座駕的時候.
就是滿有權能的.
所以你回去看看.
以西結書第一章.
他在被擄的人中.
看到這四個活物的形象的時候.
看到這四個活物有四塊臉.
人鷹牛人.
看到很大的權能.
原來.
以西結突然間.
在榮耀的意象面前.
發現原來上帝的榮耀.
不是困在耶路撒冷的聖殿.
我們被擄的人中.
上帝的榮耀仍然在我們這裡.
我們被流散到周圍的猶太人.
我們都可以看到神的榮耀.
因為神的榮耀是無處不在的.
所以神學在那裡.
是在說一種權能權勢.
所以在伊甸園.
也是在說上帝的住處.
致勝所.
有基路牌.
有金銀寶石.
你記得第一聖殿.
所羅門聖殿的廟室裡面.
聖殿有什麼裝飾.
也有這些裝飾.
樹啊.
風城啊.
有銅盤啊.
金銀寶石啊等等.
在說神特別同在的地方.
在說致勝所.

$^{281}$在會幕裡面.
最深的地方.
最中間的地方.
在聖殿裡面最中間的地方.
是致勝所這個地方.
這個就是上帝同在的地方.
伊甸就是這個地方.
所以我們看到伊甸.
是有四道河出來的.
有四條河.
你記得有.
一個Tigris.
有Euphrates.
有Valadi.
有基訓河.
等等.
比順.
四條河.
我們很多時候去找.
到底這四條河在哪裡.
如果你看地圖的話.
近東的世界.
這裡就是以色列巴勒斯坦地.
那邊就是埃及.
這裡就是沙漠.
今天的沙特阿拉伯北面沙漠.
這裡.
就是叫做米索不達米亞.
上面.
你看到這條河嗎.
就是Tigris.
下面這條河就是Euphrates河.
兩條河到這裡的時候.
就會匯聚.
然後流入波斯灣.
這個就是聖經的世界.
在創世的第二章.
你就找到Tigris.
和Euphrates.
另外那兩條河在哪裡.

$^{321}$一條叫基訓一條叫比順.
基訓.
整本聖經裡面出現了.
提到基訓的聖經裡面.
永遠都是提到耶路撒冷的水泉.
就是基訓泉.
所以在這裡.
比順.
大家就在猜測.
會不會在這裡.
以前古代的時候可能有一個乾了的河床.
你就會發現奇怪.
如果伊甸園.
是四條河.
的源頭.
根本你都放不到伊甸園在哪裡.
因為不可能.
如果放伊甸園在這裡.
有很多人猜估伊甸園應該在這裡.
這裡都是三條河的匯聚的地方.
還有.
這裡是下流.
這兩條大的河.
Euphrates和Tigris.
的源頭在哪裡.
在今天的土耳其.
會不會伊甸園在這裡.
也不是.
因為如果你說.
這裡流出兩條河.
但是比順和基訓.
兩邊根本.
距離很遠.
如果我們這樣想.
可能會想錯.
因為在近東的世界裡面.
去到表述.
去到描述聖殿.
是河水的源頭的時候.
往往不是在說.

$^{361}$一些實物的地理.
不是實際的地理.
而是在說.
一些宗教的地理.
或者我們聖經可以這樣說.
Spiritual Geography.
就是屬靈的地理.
不是在說地圖上找到的東西.
而是在說一種意象.
一種概念.
一種神學.
你看看.
在近東裡面找到的壁畫.
這個壁畫.
是阿巴拉罕時代的.
有一個王.
阿摩利人王.
聖經裡面經常提到迦南人是阿摩利人的地方.
是嗎?.
這個.
先尼林王.
這個阿摩利人的王.
大概在1700年左右的時候.
在一個地方叫馬里.
馬里是什麼地方呢?.
如果你找.
耶路撒冷.
地圖上找耶路撒冷就是這裡.
然後你找.
梧爾.
是阿巴拉罕的.
出處.
梧爾在哪裡呢?.
梧爾是在.
這個米索不達米亞南邊.
馬里.
剛才我說的.
那個阿摩利人的王.
他的首都.
他的地方.

$^{401}$在哪裡呢?.
馬里就在.
這個阿巴拉罕的上游的地方.
這個地方.
就找到一些壁畫.
就是.
講到這個先尼林王.
他去到.
登基的時候.
那些盛大的場面.
當然有很多.
這些濕漆的壁畫已經流傳不了.
但是你看到.
你看看下面.
下面就講到.
那個王.
或者那個神像.
拿著一個水樽.
一個水瓶.
然後就有些水飄出來.
有四條河.
在水瓶那裡走出來.
你說這個這麼厲害.
可以反地心吸力.
但是這些是異象的圖畫.
不是在講實物的地理.
科學.
所以你看看.
這些重構圖.
就看到水樽.
然後有四條河出來.
換句話說.
王宮這個聖殿.
是封城的源頭.
四條河.
在那裡出來.
當然那裡沒有四條河.
但是他用這個異象.
是在講一種.
所謂神學.

$^{441}$亞摩利亞的神學.
說這個王的地方.
他的聖殿.
他所敬拜的神明的聖殿.
就是.
四條河的源頭.
你這樣就明白了.
不要按照現代人.
想想.
在GPS或者在地圖上.
哪裡找到這天呢.
真正他想表達的是什麼.
封城的源頭.
就是四條河.
封城的源頭.
四條河的源頭.
就是伊甸.
有神同在的地方.
本身他的致勝所.
就是這個河水的源頭.
並且伊甸還有什麼呢.
除了金銀寶石.
河水的源頭.
就是他向東邊的.
他的出入口是東邊.
向東.
在近東的世界裡面.
每逢出入口向東.
都是勝所.
三章二十四節.
基督拍把手的地方.
就是東邊.
所以伊甸的源.
的門口.
入口是向東邊.
等於會幕的.
出入口.
都是向東邊.
聖殿的出入口.
也是向東邊.

$^{481}$其他近東世界裡面.
其他所謂神明的廟宇.
都是向日出之地.
向東邊.
再看看其他.
經文如何描述伊甸園.
講到伊甸園.
是聖山.
Holy Mountain.
向東.
是勝所.
在山上.
建造的通常都是勝所.
耶路撒冷在山上.
你記得上行之師.
120至134.
向山舉目.
然後去到耶路撒冷.
上行去敬拜.
去朝聖去耶路撒冷.
怎樣都好.
怎樣去哪個方向都好.
都是上山.
所以在山上.
以西結書28章.
形容伊甸是勝山.
所以伊甸園是至勝所.
Holy of Holies.
誰在勝所.
我們說是亞當.
其實不是.
是人類.
你記得我說.
那人安置在伊甸園.
二章十五節.
其實就是將.
每個人.
每個人有個地位.
有個身份.
那裡不是很擠迫嗎?.

$^{521}$不是.
是在講神學.
在講一些神學的宣告.
在講人.
被創造的時候.
願創是做什麼.
安置在伊甸園裡面.
即是神的至勝所.
誰才能進入至勝所.
祭師.
大祭師.
動詞怎樣形容.
叫人.
人類.
去到耕種和看守那原子.
耕種和看守.
我們一想起.
我在北美的時候.
有很多人讀聖經的時候.
就會想起我們的後院.
那裡種瓜種菜.
什麼叫耕種和看守.
和合本.
易造就是.
修理和看守.
Work and keep.
其實這兩個動詞.
是形容什麼.
你看摩西五經其他地方.
民數記三章第八節其中一個例子.
講到利未人.
看守和辦理帳目.
的會幕的事.
所以是形容.
祭師的職事.
祭師的工作.
這個才叫.
什麼.
看守和耕種.
所以人被造的時候.

$^{561}$第二章就講什麼.
第一章你記得管理萬物.
軍專的地位.
第二章是講祭師的地位.
你說祭師不就是獻祭.
不只是這麼簡單.
祭師要做什麼.
祭師是要把守聖所.
把守神同在的地方.
去遵上帝的道.
去把守神同在的地方.
這是祭師要做的事.
所以你記得.
以色列人出埃及之後去到西大山.
十九章第五節出埃及記.
上帝說我揀選你們做什麼.
我揀選你們做軍專的祭師.
你們就在萬國的面前.
去遵我的道.
遵我的約.
把守著我的同在.
我就在會幕當中.
和你們同在.
在聖殿之後聖殿.
在聖殿有生.
和你們同在.
要和你們同行.
你們要去見證我的.
這個就是祭師的身份.
所以.
原來我們第二章.
告訴我們人生觀是什麼.
做人.
上帝做人.
原來就是在伊甸.
去敬拜他.
事奉他.
一個是神聖.
祭師的身份.
二章第十六節.

$^{601}$耶和華吩咐他.
阿丹.
人類.
阿丹代表全人類.
代表我們每一個.
吩咐他說什麼.
願中國人的果子你可以隨便吃.
只是分別善惡的果子.
你不可吃.
因為你吃的日子就必定死.
所以如果你.
去聽我的道的話.
聽我的吩咐的話.
耶和華說什麼.
你就可以得到.
不老的生命.
吃個生命樹果子.
但如果你不聽的話.
你去吃分別善惡的果子.
那個知善惡樹.
你就是遠離我.
這個就是死.
死不是肉身.
死就是和上帝斷絕關係.
離開伊甸園.
沒有上帝的同在.
這句話.
我們經常說.
是不是就像在市場吃果.
禁果就是蘋果.
其實.
為什麼禁果就是蘋果.
因為.
教會中世紀的時候.
拉丁文邪惡和蘋果.
兩個字.
很像音.
所以才有這個傳說.
其實不是吃果子那麼簡單.
你要知道知善惡樹.

$^{641}$其實是代表一件事.
代表.
不是那個.
生果.
而是代表了一種生活的方式.
代表了一種抉擇.
那個抉擇是什麼呢.
如果你找一下.
知善惡.
或者分別善惡.
這句話.
在希伯來聖經裡面出現的.
是指到王的能力.
所羅門王.
在列王記上第三章九節.
那裡說.
求上帝讓他自慰.
能夠辨別是非.
辨別是非.
辨別善惡.
Good and Evil.
是王擁有的.
權能.
所以.
你看整個創世故事的時候.
你就發現如果拿這個.
禁果.
吃這個果子的時候.
就不是咬蘋果那麼簡單.
不是說這個.
吃這個蘋果.
吃這個知善惡樹的果子.
其實就是自立為王的意思.
就是說創造主的君王我不管他.
我就自我作主.
以己為主.
就是說.
我叛變.
我不想要這個王去管理我.
我想自立我自己為王.

$^{681}$這個叫做人的墮落.
更加好的翻譯.
更加好的描述.
應該是人的叛變.
人被創造不甘心是這樣.
不甘心在神權底下.
我要顛覆神的王權.
我要自己做神.
這個才叫做罪.
這個才叫做.
人的墮落.
我們說.
每個人都是犯罪.
世人都犯了罪.
我們想起罪的時候就想起做錯事.
說謊不應該說.
或者不應該想那件事.
不應該去偷東西的就去偷東西.
當然這些是罪.
但聖經裡面怎樣去定義罪呢.
是.
不要神的王權.
我要自立為神.
現代的例子很簡單.
如果我的兒子.
我的兒子已經十九歲.
他在外面.
每個人都稱讚他.
說他讀書好,說他人品好等等.
但唯獨是一件事.
我兒子回到家的時候不認他爸爸是爸爸.
那你說大不大罪.
很大罪.
不認主.
這個才是罪.
人的墮落.
人的叛變.
就是不承認.
神為神.
以偶像為神.

$^{721}$甚至以我自己為神.
以我自己為主.
這個才是罪.
所以人被造的時候.
是多麼的尊貴.
有一個祭司的身份.
但是.
你看到第三章的時候.
下一次我們講第三章.
就是講到人的叛變.
其實就是顛覆神的王權.
神怎樣反應呢.
神怎樣做呢.
這是第三章講的.
第二章主要講的是人倫.
上帝啟示聖經獨有的.
一種很有尊貴.
很有價值的.
人的存在.
人的身份.
解決人的身份.
特別去到.
表述女人的身份.
女人的身份.
我們現代人.
就經常去想.
一定是女人錯了.
因為她聽了蛇說.
然後她就吃了禁果.
搞到世界都砸了.
但是女人是什麼呢.
女人是肋骨所做.
亞當的肋骨.
是不是也是在講.
生物的.
做手術等等.
男人就是少了肋骨.
當然不是.
原文是什麼.
是一種建築的用詞.

$^{761}$就是那些.
會幕的家具等等.
有旁邊的.
有側邊的.
有另外一邊.
一對的一邊.
英文叫做the other half.
那人見到那個女人.
她說這是骨中的骨肉中的肉.
不是一條骨這麼簡單.
而是我的同等.
我的本體的同等.
我的equal.
跟我同等地位.
這就是女人.
女人是幫手.
你記得她去到.
為動物去改名管理萬物的時候.
找不到幫手.
神就為她做一個女人.
來幫手.
幫手是不是就是.
高等一點的地位.
當然不是.
因為你看原文.
幫忙這個字.
經常都是形容誰.
形容神自己.
神是高等一點還是低等一點.
神是高等一點.
撕邊撒邊.
我的心裡等候耶和華.
因為神是我們的幫助.
我的幫手.
我的盾牌.
女人不是低等一點.
這是很革命性的語言.
因為.
在近東的社會.
女人是低等的.

$^{801}$甚至不是人.
是男人的物業.
擁有品.
但是第二章.
所揭示的是什麼.
神的形象.
男與女.
是同等的.
更加好翻譯應該是聯盟.
Ally.
亞當.
被造是一個祭司.
他去敬拜上帝.
守著上帝榮耀上帝的同在.
他自己一個做不到.
要下蛙.
要女人來幫他.
是一個神聖的任務.
一個祭司.
神聖服侍上帝的任務.
所以你看到.
第二章和第一章.
其實揭示.
很偉大的宣告.
就是要我們每一個讀者.
每一個聽眾.
去聽清楚我們在上帝面前.
我們有的地位.
我們的身份.
所以有一些反思的問題.
到最後.
我們等一下.
如果可以的話我們就有些討論.
聖經其實.
是在比對一些.
神話的故事.
會不會.
令到大家很不舒服.
影響了聖經神聖的權威.
如果沒有.

$^{841}$是為什麼.
剛才我所說.
聖經寫的時候.
是針對當時的世界觀.
所以為什麼要.
去比對這些假的神話.
我們叫做.
為什麼要比對巴比倫神的神話.
他們根本不是神.
是虛無的東西.
為什麼要去比對.
為什麼要去比較.
去跟他們爭辯.
跟他們作對話.
告訴他們你們這樣了解人.
是不對的.
聖經說這樣了解人才對.
因為神創造人的時候.
原創是什麼?是君尊的祭祀.
是去敬拜上帝.
是跟上帝建立關係.
比對一下.
聖經和近東的一些故事文獻.
你就會發現.
要有一些相似的地方.
大水開始.
分開.
做天做地.
但本質上很大分別.
沒有神鬥論.
每樣東西都是聽上帝的話.
創世的一二章.
去比較我們今天.
世俗的一些世界觀.
剛才我說的.
唯物論.
有什麼不同.
會不會跟我們今天.
我們接觸的.
唯物論等等.

$^{881}$有很大分別.
上帝也要挑戰.
這種世界觀.
聖經創世的一二章.
挑戰這些觀念.
要我們知道什麼?.
人類在上帝面前.
有的本位.
君尊的祭祀.
那要怎麼應用呢?.
很多地方都可以應用.
我們有沒有追求聖潔.
追求敬拜神.
在上帝的致勝所.
我們跟上帝的關係.
應該怎麼建立?.
怎麼維繫?.
我們跟人的關係.
說到第一個家庭.
亞當夏娃.
所謂的.
男與女.
一齊有一個使命.
做祭司的使命.
教導其他人.
神的律法.
遵守上帝的.
聖約.
活在上帝的道裡面.
跟上帝同行.
甚至我們跟大地之間的關係.
我們跟大自然的關係.
因為我們說過了.
創造界.
就是神聖的.
好像上帝的聖所.
一樣.
我們不應該糟蹋這個世界.
我們不應該去剝削.
大自然的資源.

$^{921}$是不是?.
為了我們賺錢.
然後污染這個環境.
我們出到.
世界裡面.
出到大自然裡面.
走每一步都是上帝的物業.
都是屬於上帝的地方.
是神聖的地方.
我們怎麼去愛護.
這個.
護理.
這個世界.
跟我們自己.
上帝 其他人.
跟大地之間.
有什麼關係?.
這些都是創世的第一一二章.
要解答最重要的問題.
最簡單可以這麼說.
中文說.
天倫.
人倫.
和地倫.
天倫就是我們跟上帝之間的關係.
人倫就是我們跟我們的倫社.
我們家人.
我們身邊的人.
弟兄姊妹.
就是我們跟大地之間的關係.
創世的第一二章.
就是要揭示一個這樣的.
基本的.
一個世界觀.
以致我們活為.
神的子民 有神的同在.
如果你喜歡我們的影片 記得要讚好 留言和分享.
多謝您收睇時局新聞,再會!.
\newpage



\section{詩篇}
\label{sec:oVRXAcV7BoQ}
\textbf{逆要作門徒 (3) — 再思詩篇一二 (陳崇基牧師) 2022.01.23}
\newline
\newline
連結: \href{https://youtube.com/watch?v=oVRXAcV7BoQ}{\texttt{https://youtube.com/watch?v=oVRXAcV7BoQ}} ~~~~ 語音日期: 2022-01-25
\newline
\newline
\hyperref[sec:yR3tIvy_FRE]{\small{< < < PREV SERMON < < <}}
~
\hyperref[sec:index_chronic]{\small{[返順時目]}}
~
\hyperref[sec:index_scriptual]{\small{[返順卷目]}}
~
\hyperref[sec:gQBhlqNH6hM]{\small{> > > NEXT SERMON > > >}}
\newline
\newline
$^{1}$我們這次繼續講《逆妖作門徒》這個系列.
這次我們講到詩篇第一和第二篇.
我想這次我們先從一個新約的經文開始.
這個經文就是《士徒行傳》第十三章三十三節.
這裡上下文講什麼呢?.
就是保羅第一次的宣教旅程.
去到土耳其一個地方.
他就跟人們講到耶穌就是基督.
就是神所立的君王.
所以這裡三十三節就講到.
藉著耶穌的復活.
向我們這些作子孫的應驗了.
就如詩篇第二篇所記的.
就是你是我的兒子.
我今天生你.
為什麼要用這一節聖經開始呢?.
因為有一些古代的抄本.
另一個傳統的抄本.
不是講到詩篇第二篇的.
是講到詩篇第一篇.
所以這裡似乎告訴我們.
在新約的時候.
詩篇第一篇和第二篇.
可能是一起當成是第一篇的.
或者甚至第一,二篇.
本身是沒有數目的.
沒有編碼的.
事實上一些中世紀的抄本.
詩篇第一篇.
是用紅色字寫的.
是沒有數目的.
不像我們今天講的.
這個就是詩篇第一篇.
所以如果你這樣看的時候.
在新約的教會.
看到詩篇一,二篇的時候.
似乎是告訴我們.
這兩篇其實是一個preface.
一個序言.
或者告訴我們.

$^{41}$在這兩個中間.
收集了這些詩篇.
成為詩集.
可能有100多首.
148首.
最後文士或者編輯.
才加上一個introduction的序言.
就用了現在的詩篇一,二篇.
所以很明顯.
一,二篇之間有很多共通的地方.
以及整本的詩集.
整本敬拜的詩歌集.
作一個introduction.
作一個序言.
如果你看看.
在詩篇第一篇和第二篇.
確實有很多相同的字眼.
是有用的.
譬如第一章第一節.
這裡說有福的人.
這是新譯本的翻譯.
因為有福的人這個字.
share.
真的是詩篇第一個字開始的.
胡合本將它推到第二節最後.
因為這樣文法讀起來比較通.
但是確實是有福的人開始的.
不從愛人的計謀.
不暫聚人的道路.
不坐好羈少人的座位.
這裡有福的人.
又在第二篇的最後一節.
十二節又出現了.
梵透教耶和華的.
就是有福的.
share.
也是同一個字.
第二個相同的字就是.
第一章第二節.
晝夜勿縱.

$^{81}$這是新譯本.
或者思想神的律法.
二章第一節.
又提到這些外邦或者反對神的皇權的列國.
萬民為什麼胸謀妄想.
是同一個字.
勿縱或者思考.
和胸謀妄想.
就是計劃反對上帝的皇權.
是同一個字.
一章六節和二章十二節都是一樣.
第一章第六節提到惡人的道路.
必定會滅亡.
滅亡這個字在二章十二節又出現過.
所以一和二篇.
其實是整個的詩歌本.
好像一篇大門一樣.
一個序言.
就是你進入一間大屋的時候.
你進入一個聖殿的時候.
這個entrance.
這個門口.
進來的時候就給我們一些提醒.
一些警告.
或者一些挑戰.
所以這次我們看詩篇第一和二篇一起.
首先看第一篇.
第一篇很多時候我們叫做.
Wisdom Song.
就是智慧詩歌.
智慧詩歌是什麼意思呢.
當你看到一些詩篇.
講到神的律法.
或者義者惡者兩個生命給你做一個選擇.
大概就是在看智慧詩篇.
就是Wisdom Song.
就是教導我們挑戰我們.
詩篇第一篇最為人熟悉了.
所以在第一篇第二節經常說.
我們經常背的.

$^{121}$就是維曉為爺和律法的晝夜思想者.
吾便為有福.
他好像一棵樹坐在溪水旁.
按時候結果指.
葉子也不枯乾.
做的事情也是順利的.
坐在溪水旁的樹.
是想起什麼景象呢.
如果你Google search的話.
好像我去Google search.
我就找到這一類的相片.
這一類很美的.
給我們感覺很安息.
很安恕.
很安靜的.
但是呢.
在以色列巴士眼地想起.
坐在溪水旁的樹.
是不是就是這個意思呢.
其實呢.
應該不像我們剛才那些相片.
那麼浪漫的.
那麼美麗的.
在曠野當中.
譬如這兩張相.
都是在猶大曠野的隱基地.
那裡拍的.
這裡就是.
剛剛在曠野裡面有一個水泉.
在盆石裡面湧出一些水.
然後那些樹植物.
就在旁邊那裡.
要下很多苦工.
在逆境裡面.
它才可以生存.
要向下扎根.
很深的根.
才能夠生存.
能夠去生活.
因為在曠野裡面.

$^{161}$寸草不生.
這些盆石出水的時候.
就是生命的源頭.
在隱基底這個地方.
看到這樣的景色.
似乎就是告訴我們.
詩篇第一篇那種景象.
其實隱基底這個地方.
就是大衛被蘇羅追殺的時候.
高立為王之後.
有很多敵人.
要去追斬他.
他就去到隱基底這個地方.
然後在那裡去敬拜神.
詩篇第一篇真正的意境.
其實也不是說這麼安殊.
只不過我們浪漫化了.
其實詩篇第一篇.
你想起面對惡人的時候.
行鬼詐的人.
周圍環繞著他.
或者一些抵擋公義的仇敵.
扭曲真相.
無論是今時今日經常出現.
還是在以前的那個世代.
所以詩人就像一個王一樣.
就像大衛王一樣.
他就說不從不斬不助.
這些惡人的遺志.
所以你比較一下詩篇.
用這些意境的時候.
我們似乎就可以去到想像.
到底那篇詩篇給我們什麼挑戰.
在這種逆境當中.
我們應該怎麼做一個有決心.
來遵循上帝律法的人.
幾年前我去大英博物館的時候.
搜集資料.
我就看到這一幅圖畫.
這一幅圖畫讓我想起.

$^{201}$就是我們今天看的詩篇第一篇.
為什麼呢.
這幅圖畫上面有一個山.
然後有一個通道上去山上.
這個亞述王就站在那裡.
他在那裡敬拜.
上面有一個皇宮.
就像聖殿一樣.
有一個廟.
然後在山的附近周圍.
就有很多這樣的樹木.
是非常茂盛的.
這個是皇室.
一個皇室的花園.
聖殿這個地方.
這個就.
你看到王在山上.
這些樹四周環繞他.
就在王的同在底下.
這個起碼就是詩篇.
用相似的畫像.
相似的語言來跟我們說.
通常這些畫本身是有顏色的.
只不過是剝落了.
博物館就將這些顏色畫上去.
看到有河水.
溪水旁邊的.
就是這些很茂盛的樹.
然後這些樹就來到皇室裡面.
在廟宇附近.
就跟這個王能夠親近.
其實你看詩篇.
都是用類似這樣的語句來形容的.
譬如詩篇52篇第八節.
詩人說到什麼呢?.
至於我就像神殿中茂盛的橄欖樹一樣.
我倚靠神的慈愛直到永遠.
92章13節.
他們指的是義人.
義人必定將這些香柏樹.

$^{241}$很高很大的樹.
栽在哪裡呢?.
耶和華的殿中.
耶和華的院子裡面.
在我們神的院子裡面.
繁茂的生長.
講到一棵樹.
詩篇裡面的世界.
就是告訴我們.
在聖殿當中.
在神的同在裡面.
所以詩人就挑戰我們.
是一個這樣的生命.
一個跟神的同在.
在上帝的聖所裡面.
跟祂連結起來.
愛祂的律法.
以至祂能夠去面對.
外面那種很邪惡的世代.
所以在第一章第四節.
這裡就告訴我們.
惡人卻不是這樣.
詩人看回這些惡人.
看我們周圍的惡人的時候.
他們像什麼呢?.
就不像栽在耶和華殿中的樹.
不是在耶和華的面前.
不是在祂同在裡面.
慷僻被風吹散.
在義人的會中.
這些罪人都是一樣.
他是不能夠站立的.
因為將來去到審判的時候.
耶和華必定會去看顧義人的道路.
惡人的道路就必定會滅亡.
在巴勒斯坦地.
去到收割大小麥的時候.
農夫就會去到山上.
或者是空曠的地方.
禾場 風比較大的地方.

$^{281}$他們將這些麥子養起.
然後風吹過的時候.
有實質生命的就會掉下來.
沒有實質生命的慷僻.
就會被吹散.
這裡告訴我們將來的盼望.
一個圖畫.
用一些當時非常熟悉的景象.
告訴我們一個義人.
縱使現在未必一定是順利.
但是在將來.
上帝審判的日子來到的時候.
我們望著那個成為我們的盼望.
我們就能夠去面對我們現在的境況.
就是去到僅靠著耶和華.
去到喜愛他的律法.
去追隨他的律法.
去研讀他的律法.
以致我們的生命就像麥子一樣.
是不會怕被風吹散.
在1989年的時候.
當時的達賴喇嘛.
會見幾個猶太教的不同支派.
不同的所謂denomination的幾位學者.
在會見當中.
好幾個小時他就問了很多問題.
他最有趣問的問題是什麼呢?.
達賴喇嘛當時最有興趣想知道的.
就是猶太人生存的秘訣.
Secret Technique of Survival.
因為他看到猶太人的歷史.
就好像西藏民族有一點點相似.
但是他卻看到猶太人這麼遙遠的歷史.
這麼悠長的歷史.
面對著很多打壓.
聲電被毀.
失去家園.
失去政治的權力.
甚至被打壓等等.
遇見很多挑戰.

$^{321}$應該是早就消失了.
這個民族應該不在世上.
應該在歷史書裡面才找到.
但是為什麼反而可以建國呢?.
反而他們的身份能夠繼續延續下去呢?.
這些拉比很多時候都會有爭執的.
每個人的看法都會不同.
但是最後一個最重要的因素.
大家是同意的.
就是Torah.
就是律法書.
在個人和群體的生命當中.
成為一個中心.
當群體讓上帝的道成為他們的中心的時候.
他們就有一個所謂的.
Secret Technique of Survival.
在逆境當中能夠延續下去.
其實你想一下.
在教會的歷史也是一樣.
我們只看新約教會的歷史.
2000年的時間.
每逢有所謂的復興的時候.
都不是一些情緒高漲就叫做復興的.
千萬不要覺得那些Great Revival.
我們看回的時候.
那些人是很高漲的.
哭到幾天幾夜.
然後情緒非常多的表露.
這些不是一件錯的事.
但是每逢教會裡面有復興.
都是回歸聖經.
回歸到上帝的道的裡面.
所以這個詩篇第一篇就告訴我們.
面對著惡者很猖狂的時候.
特別是我們這個時代.
到底有什麼讓我們Survive.
可能不單止Survive.
讓我們Fly.
讓我們不單止生存下去.
甚至讓我們去到茂盛的去生長延續下去.

$^{361}$當然詩篇第一篇.
就不是純粹告訴我們一些ethics.
應做什麼不好做什麼.
其實也都揭示給我們一個.
一個中末的盼望一個的Eschatology.
詩篇第一篇有提過.
將來審判的時候.
惡人像康丕.
義人像一些實質的大小物.
第二篇都是一樣.
第一篇第二篇是一齊要看的.
第二篇很清楚.
就不是哲偉詩.
就是君王的詩篇.
Royal Song.
大衛君王登基的時候.
或者大衛後代在猶大國.
登基的時候經常都會用一些Royal Song.
登基之日的時候就會唱出.
第二篇就是很典型的Royal Song.
我們看幾節抽出來看一看.
看二章的第一節.
這裡說列國為什麼騷動萬民.
為什麼兇謀妄想.
世上的君王起來.
領袖聚在一起.
敵對是耶和華.
和敵對他所高立的君主.
外面有很多這樣的惡人.
其實這些惡人就是詩篇第一篇裡面所說的惡人.
他們就是去敵黨上帝的主權.
上帝的王權.
當這些事情發生的時候.
令到我們地上覺得很不安的時候.
二章四節就告訴我們.
坐在天上的居然就是發笑.
他說我已經在石安我的聖山耶路撒冷.
立了我的君王.
所以你們不需要怕.
然後二章第七節等等.

$^{401}$又說到受高者說話.
受高者說我是神所生的兒子.
這裡當然不是說本質上的生出來.
而是說什麼呢.
就是神立他為君主.
他就好像神的兒子一樣.
二章第十節.
這裡說現在君王.
下來地上跟這些君王說.
叫他們謹慎.
應當要去聽勸告.
什麼勸告.
就是要以敬畏的心來事奉耶和華.
因為耶和華必定會去制服列國.
去制服所有敵對他地上的勢力.
這些邪惡的勢力.
上帝必定會得勝.
所以十二節就有一句這樣的說話.
你們要用嘴親子.
親吻王子.
神所立的王者.
當然不是面貼面的親吻.
在中東裡面.
什麼叫親吻.
Kiss the son.
其實就是五體投地下拜.
要親吻他的腳.
就是說他敬拜上帝.
認定他的王權.
如果不是.
上帝的審判臨到的時候.
當這樣發生的時候.
列國的君王就會被毀滅.
這裡給我們一個所謂的中末論.
一個中末的角度.
告訴我們為什麼上帝值得依靠和敬拜.
另外一個圖片在大英博物館.
看到的時候.
我就想起這個CPN第二篇.
你看到這幅圖.

$^{441}$是亞述的皇帝.
他在自己的花園.
Royal Garden那裡.
坐在沙發.
然後他的皇后就坐在他的旁邊.
宮廷有很多人撥線.
去奏音樂.
來服侍他.
基本上就是這樣的圖像.
去表達帝王的豐盛.
他的生命力.
其實不單止這樣.
也展示皇帝能夠維持著國度的能力.
這個皇帝是很大能的.
坐在那裡.
他就可以維持著整個國度的太平.
如果仔細去看的話.
在左邊有一棵樹.
樹當然在皇帝的面前.
樹上掛著一個物件.
原來那個物件是敵人的頭顱.
我們覺得很安靜.
在探望世界的一個畫面.
原來這些情景告訴我們.
亞述帝王能夠戰勝國邦國族.
這些敵人能夠將他們打敗.
看到這些景象的時候.
那些帝國裡面的人.
就覺得亞述王真的很厲害.
可以平定四方.
所以信他的時候.
信靠他就可以了.
詩篇第一篇.
詩篇第二篇.
其實就是去到顛覆.
用這些帝國的圖像和語言.
來講出耶和華才是作王.
耶和華所立的君主.
他的王者.
才是作王的那個.

$^{481}$詩篇第一篇講到神的源子裡面.
詩篇第二篇就講到.
神所立的君王.
是最終會得勝的.
耶和華才是王者.
但是必須要用一個信心的眼睛.
才能夠看到.
因為你不用信心的眼睛看的時候.
你只要看到地上發生的事.
詩篇第二篇.
第二篇你會看到格式很工整.
分成四段.
用兩個角度.
兩個的角度來看我們今天的現實.
我們的現實.
地上看到的現實.
就是地上的騷動.
很多的不安.
很多反對上帝的勢力.
很多不合真理的事.
但是用信心的眼睛.
你就會看到天上的角度.
邀請我們用天上的角度來看世情.
我們的盼望.
不是地上我們肉眼看到的.
盼望是另一個真正的現實.
一個永恆的現實.
就是神必定會來施行公義.
施行真理.
並且去重整他的世界.
其實如果你看詩篇第二篇.
甚至看完整個詩集的話.
都是在說一個盼望.
在說我們在上帝裡面.
我們有的好.
詩篇共有五卷.
我們大家都知道.
這五卷是有不同數目的詩歌.
來去搜集起來的.
每一卷有五卷.

$^{521}$每一卷最後的時候.
都會用一個短的讚美.
最後的時候用最長的讚美.
來結束.
所以你看到有五卷的詩歌.
為什麼有五卷呢.
其中一個解釋就是摩西律法.
有五部書.
創出黎明新.
無論如何.
頭三卷是一個單元.
一個unit.
這三卷大致上你看的時候.
其實是說出大衛國度的興瑞.
由他開始被立作君王.
遇見很多的苦難.
很多的追殺.
挑戰.
但是他信靠上帝.
最後王國建立.
到第二卷的時候.
以七十二篇為結束.
七十二篇就是另一篇君王的詩篇.
黃色的.
君王的詩篇就是所羅門的詩歌.
就是一個禱告.
希望大衛的皇室.
是按照公義來治國.
然後第三卷就是七十三至八十九.
用一篇智慧詩又開始.
逢這些夾縫的地方.
卷和卷之間.
很多時候就是智慧詩篇和君王的詩篇.
是有心這樣串連在一起.
卷三最後的時候.
就八十九篇.
你就發現.
是大衛的國度結束.
是亡國被擄的時候.
所以你看八十九篇.

$^{561}$四十九篇.
那裡有一個禱告.
一個很慘的禱告.
就是說你上帝從前指出你的信實.
向大衛起誓要施的慈愛.
現在在哪裡呢?.
很明顯就是在描述一個亡國.
黑暗好像去到很猖狂.
在德勝的時候.
這裡說的是公元前.
主前五百八十六年的時候.
就是大衛的王國.
猶大國.
被擄與巴比倫的時候.
好像一切都玩完了.
但是到卷四開始.
九十篇開始.
就開始有一個.
好像一個盼望.
再一次重現一樣.
你記得九十篇是說什麼嗎?.
九十篇的標題.
是神人摩西的禱告.
你記得在舊約裡面.
五經出埃及記.
說到什麼?.
說以色列人背若.
幾乎要滅亡了.
因為耶和華的震怒要來到.
但是摩西作中保.
他去禱告.
為以色列人祈禱.
他是一個媒體.
將上帝的恩典.
將上帝的信實帶到叛逆的民眾.
上帝去赦免他們.
重新去建立他們.
然後到卷四最後的那一篇.
就是106.
就是在整本的詩集裡面.

$^{601}$摩西最後一次出現.
就是106篇.
如果你看106篇第23節的話.
這裡說.
所以他說要滅絕他們.
即是上帝要滅絕以色列民.
不是有他揀選的摩西.
在他面前站在破口上.
使神的怒氣轉燒.
恐怕他早已毀滅了神的子民.
一切都是上帝的恩典.
卷四其實就是回答.
卷一二三裡面有的神學問題.
為什麼我們見到好像.
黑暗張狂 黑暗德勝.
上帝的恩典.
臨到的時候.
有恢復國度的盼望.
大衛的王國可以被恢復過來.
雖然現在是在流亡.
但是深信神所立的君王.
一定會再來的.
所以有幾篇的讚美詩篇.
去到信靠上帝的詩篇.
Song of Confidence.
Trusting.
最後110篇.
有另外一篇的君王詩篇.
就是深信上帝的信實.
他給我們的盼望.
一定會實施出來.
一定會兌現.
現在我們在新約的信徒.
我們就知道.
這個當然是大衛的真正君主.
永恆的君主耶穌基督.
但是在舊約裡面.
尼賽亞還沒有來到的時候.
第卷四和第卷五.
仍然去到.

$^{641}$憑著這個盼望將來上帝的說話.
一定會兌現.
仍然去到發出讚美.
信靠的詩篇.
所以整本的詩集.
最後最後的時候.
仍然是以讚美.
以worship.
Hallelujah.
來結束.
因著這個盼望.
我們雖然看到的.
未必是上帝所有的話語都兌現.
但是因著這個盼望.
我們能夠繼續去到敬拜.
讚美這個主.
什麼叫盼望.
最近我看一個新聞.
這個新聞說到一個香港家庭移民的故事.
說到這些小朋友.
很興奮.
因為離開香港.
去到一個他們可能很喜歡的地方.
所以他們就因為知道要走.
所以他們就不理香港學校的功課.
學校的功課的壓力.
對他們來說完全是等於零的.
立即開始在家裡講英文.
然後上去YouTube看很多英文的資訊.
在家裡和父母都講英文.
當然這樣去看一個移民的故事.
可能是過分簡單.
每一個家庭去到移民一定是非常艱難.
裡面有很多苦頭.
特別對大人來說.
但是我想去做一個比喻.
去表達的是什麼.
就是那個future.
那個將來.
成為了這些小朋友.

$^{681}$改變他今天的動力.
知道將來會發生什麼事.
有那個盼望.
以至今天他生活的態度.
是完全不同了.
詩篇第一和第二篇.
就是整本詩集的序曲.
就是邀請我們見到這個盼望.
見到上帝自己仍然作王的時候.
邀請我們去到神的面前.
在聖殿裡面來敬拜.
我不知道你每一次敬拜.
你是否有這樣的信念.
每一次敬拜的時候.
我們知道上帝必定會再駕臨.
祂的國度必定會再駕臨.
我們就好像詩篇第一篇.
要接受這個挑戰和邀請.
我們是否以上帝的律法.
作為我們生命的中神.
詩篇第二篇.
君王的詩篇.
地上雖然眾矢不安.
但是敬拜我們天上作為為王的神.
祂在上面發笑.
祂仍然坐著為王.
已經立了祂的王者.
耶穌基督為君王.
祂必定是管理 統治萬有.
也統治我們生命.
我們也可以依靠祂.
2022年 這一年的年頭.
在科恩派的教會.
有一個這樣的調查.
這個調查是有些驚訝的.
這個調查說什麼呢.
就是說一些年輕的.
Evangelicals 福音教會的信徒.
比上一代更加願意.
想聽一些深入的In-Depth.

$^{721}$深入的講讀.
我不知道這個調查有多真.
不知道這種現象.
是否真的可以繼續維持下去.
我記得教會的復興.
上帝要來工作的時候.
就是上帝的律法.
會到聖經.
然後我們在律法底下.
遵從祂的旨意來生活.
上帝的王權來到我們中間.
詩篇第一篇 第二篇給我們的挑戰.
詩篇第一篇 第二篇.
開始的時候 打開門.
給我們一個選擇.
選擇我們是否願意去依靠神.
在惡人當道.
可能你見到很多不利的情況.
可能是一些個人的挑戰.
我們不只是在這裡生存.
不是去survive.
而是我們茁壯地去生長.
Fly 決心地去學校.
學習神的律法.
去敬拜這位唯一值得我們敬拜的君王.
最後我想用Peter Kuzmich.
這個東歐的神學家.
他曾經這樣定義什麼叫做盼望.
他說盼望是什麼?.
盼望是聽到將來天國樂曲的能力.
信心就是今天隨著樂曲起舞的勇氣.
希望是聽到將來的樂曲.
信心是勇氣去跳舞樂曲在現在.
可以同心協力去祈禱.
因為你對我們這一班不配的人.
你去到承諾.
承諾你是那位王者.
你是統管萬有的.
耶穌基督已經得勝.
主我們求你幫助我們.

$^{761}$讓你來帶領我們.
主叫我們喜愛你的律法.
聽見你的律法的時候.
聽見你的心意我們尊重.
求主你給我們這一代.
我們能夠回歸你自己的話語.
讓你的聖言引導我們的生活.
讓你的盼望在你的國度裡面.
將來耶穌基督再回來.
這個榮耀的迎去.
引導我們今天能夠有勇氣.
去事奉你 去敬拜你.
忠心的跟隨你.
祈禱奉主耶穌基督的名字.
謝謝大家收看 再見.
\newpage



\section{2:1-3}
\label{sec:gQBhlqNH6hM}
\textbf{Bring the Christ in the World (1 Kings 2\_1-3) - Tommy Lee}
\newline
\newline
連結: \href{https://youtube.com/watch?v=gQBhlqNH6hM}{\texttt{https://youtube.com/watch?v=gQBhlqNH6hM}} ~~~~ 語音日期: 2022-01-28
\newline
\newline
\hyperref[sec:oVRXAcV7BoQ]{\small{< < < PREV SERMON < < <}}
~
\hyperref[sec:index_chronic]{\small{[返順時目]}}
~
\hyperref[sec:index_scriptual]{\small{[返順卷目]}}
~
\hyperref[sec:wRAUIVuDJgw]{\small{> > > NEXT SERMON > > >}}
\newline
\newline
$^{1}$- Pastor Mark, thank you very much, welcome..
My name is Tommy Lee,.
and I'm born and raised here in the city of Chicago..
So 44 years old, been here,.
living here on the South side of Chicago,.
graduate from Washington University in St. Louis.
over at business school..
And then for about 10 years, I built cell towers..
But during that period of time, I said,.
Lord, what is it that you want me to do?.
Is it to be a pastor?.
Because I've always wanted to do full-time ministry.
and serve you..
And so at night, I went to seminary..
I went to Moody Theological Seminary for about four years..
And then after that, I graduated.
and I was a pastor for about three years..
And through that, I realized,.
God, you did not call me to be a pastor..
That just was not my skill..
That was not how you've gifted me to do it..
And so now in the last 12 years,.
I get a chance to utilize some of my strengths,.
utilizing some of my experience in the business world,.
in the ministry world..
So I get a chance to lead four different organizations,.
whether it's at Resource Global and at the Growth Center.
and different other organizations.
to continue to integrate how God has gifted me..
And that really is the whole idea,.
the theme today that we're gonna talk about.
for those who are in-person,.
those who are watching online..
What does it mean for us to be Christ in the world?.
God is giving you a story..
All of us have a unique story..
Our stories are different..
But within those stories,.
there's chapters of the whole idea.
of chapters of a role that God leads all of us..

$^{41}$And there is a journey that he uses.
as he continues to develop our passions.
and our strengths to ultimately follow him..
And today we're gonna look at a specific story today..
The story of David and Solomon..
That as David dies,.
he leaves him with three simple words,.
three simple phrases that Solomon needs to understand.
in order to be a successful king..
And that's what we're gonna look at today..
And so before we turn on,.
you know, the last,.
I've been with you all for three times already..
And the last three times I've talked about the history.
and life of David..
And the whole story started,.
which leads us to this point..
David is a young shepherd boy..
Samuel is told by God that Saul is no longer the king.
and Samuel, I need you to go to Bethlehem..
And I need you to choose the next king..
Samuel goes to Bethlehem and there he meets Jesse..
And Jesse parades his son.
and Samuel looks at his oldest son comes in.
and the Lord tells him,.
do not look at the outward appearances,.
but the Lord looks at the heart..
None of the brothers,.
the older brothers are suitable to become king..
Samuel then asked Jesse,.
is there someone else?.
He ushers in the youngest boy..
David is in the shepherd field..
He's in the sheep,.
tending the sheep in the field..
So he brings them in..
And there he is anointed as king..
If you are on the side,.
now, if you ever get a chance,.
there's this lady by the name of Dr. Laverne Tolbert..

$^{81}$And she wrote this nice little pamphlet that says,.
the title, I serve you just as my mother did..
And it looks into the Jewish history.
of why David is so faithful to God..
And he makes one phrase in the Psalms..
He says, I serve you just as my mother did..
And we talk and we look at the faithfulness of David.
all throughout scripture..
And that's what God is using him..
That's what God has seen in his life..
Story goes on..
Another sermon that I did with you guys,.
David and Goliath..
All of you guys are familiar.
with the story of David and Goliath..
David battles the giant because he is sent,.
he wasn't even chosen to be a soldier..
He is sent by his father to go deliver food.
over to his brothers..
And there he sees Goliath..
And he sees the Philistines coming against the Israelites..
And as we talked about,.
the reason why the Philistines are attacking the Israelites.
is the Israelites control the passageway.
to the Jordanian Valley..
In order for them to get past the Jordanian.
and control the passage,.
they have to get rid of the Israelites..
David is the only one brave enough.
to fight the Goliath the champion..
And ironically, Saul, who is the King of Israel,.
is a Benjamite..
And one of the things that we talked about is.
what do Benjamites do?.
They're experts in the slingshot..
But Saul did not think of the slingshot..
Saul is afraid of Goliath..
And so David kills him with the slingshot..
Now, you all think that the story after that is.
happy ever ending, right?.

$^{121}$And as a result, David ascends to become King.
and story is over..
No, it actually doesn't..
The story of 1 Samuel tells us then.
David is now chased by Saul because everyone loves David..
Songs are written about David..
Messages, everything..
Everyone is a fan of David, except Saul..
And Saul ends up whining to kill David..
And so as a result, from 1 Samuel all the way to 22,.
about 17 to 22, it tells us the story of David running..
David runs to his wife or actually more so before that,.
David loses his job..
Because Saul no longer wants him to work for him..
He tries to kill him..
He runs to his wife, Michal..
Well, Michal is married to Saul's,.
is married, Michal is a daughter of Saul..
So he can't stay there..
He runs to his mentor, Samuel..
Well, eventually Saul finds him there and he has to leave..
He runs to his best friend, Jonathan, who is the son of Saul..
Well, eventually he can't go there..
Everything he could depend on from his job,.
from his wife to his mentor, to his best friend,.
things that he could depend on are all gone..
All these crutches are taken away from him..
You know where he ends up?.
He ends up in the city of Gath,.
the home of the Philistines, the home of Goliath,.
who he just killed..
And he ends up pretending to be a crazy man..
After that, he can't stay there..
He ends up in 1 Samuel 22, in the cave of a dulem by himself..
Everything is gone..
All the stories, all the songs,.
all of his crutches are all gone..
That's where David is..
And he writes, if you guys have a slide,.
he writes Psalm 57 and he writes Psalm 142..

$^{161}$And there in 142, if you get a chance,.
let's take a look at 142, extra 142, is this..
Next slide..
Psalm 142, I cry out to you, O Lord..
Yeah, I cry to you, O Lord..
I say you are my refuge,.
my portion in the land of the living..
Attend to my cry, for I am brought very low..
Deliver me from my persecution,.
for they are too strong for me..
Bring me out of the prison,.
that I may give thanks to your name..
The righteous will surround me,.
for you will deal bountifully with me..
All his crutches are gone,.
and he has no one else to depend on,.
but the Lord, as he wrestles by himself.
in the cave of Doolam..
The Lord needed him to learn that journey.
of what it means to depend on him..
That David, when all the crutches are gone,.
when all the people that you could depend on are gone,.
can you still have faith and trust in me?.
David had to learn that lesson..
Once he learns that lesson,.
God brings a community around him..
He now then becomes king..
And not only does he become any king,.
he becomes one of the greatest kings in history..
The greatest king Israel has ever seen, second to Jesus..
But yet, everything doesn't go well, does it?.
He makes a serious mistake..
The man after God's own heart makes a serious mistake..
The biggest mistake that he makes is Bathsheba..
He has an affair with Bathsheba..
A son is born through Solomon..
He ends up killing Bathsheba's husband, Uriah the Hittite..
He doesn't make a lot of good decisions..
And he has to learn to deal with the consequences.
and the repercussions of his sin..

$^{201}$And that will affect his legacy..
That will continue to affect his kingdom.
and what history writes about him..
And so now the story ends..
And now we are over in 1 Kings 2..
David is now an old man..
Solomon, his youngest son,.
who is the son of Bathsheba and David,.
is gonna become king..
And we all know about Solomon..
One of the wisest kings in history..
David sits down with him..
And he gives him these three phrases..
When David's time drew near,.
he commanded Solomon, his son, saying,.
"I am about to go the way of all of the earth..
"Be strong and show yourself a man..
"And keep the charge of the Lord.".
Be strong, show yourself a man,.
and keep the charge of the Lord,.
walking in his ways and keeping his statues,.
his commandments, his rules, and his testimonies,.
as it is written in the Law of Moses,.
"That you may prosper in all that you do.
"and wherever you turn.".
And that's what we're gonna talk about,.
is three simple phrases..
Not gonna spend too much time breaking out the verse..
I wanna talk about what those verses mean.
in our culture today,.
not only for those who are sitting here.
in Chicago, Illinois,.
but those who are watching it.
from different parts of the world..
What do the phrases, be strong,.
show yourself a man,.
and walk in his ways,.
mean for us as Christ followers,.
each and every single day?.
What does it mean in our work?.

$^{241}$What does it mean in our family life?.
What does it mean in our personal life?.
And that's what we're gonna talk about today..
Be strong, show yourself a man,.
walk in his ways..
First one, be strong..
One of the things, a couple months ago,.
or actually probably about a month ago,.
my wife is a first grade teacher.
here in Chinatown in Chicago Public Schools..
About, man, about one o'clock,.
I'm on the Zoom call doing interview..
And I get a text from my wife saying,.
"Our school is on lockdown..
"There is a shooter somewhere inside the school.
"or outside the school, no one knows.".
That's it..
And so I'm on this Zoom call,.
and I call her back right away..
Nothing, get nothing..
I go back to my work at three o'clock,.
she calls me up,.
and she says to us,.
"There was somebody killed right outside our school..
"And Tommy, I think the person who was killed.
"is the father of our closest friend..
"The Lamb family who are pastors.
"right next door to this church,.
"their father and father-in-law were killed,.
"was killed and shot 22 times.
"by an individual for no reason at all,.
"but just driving down Chinatown,.
"sees them, shoots them,.
"gets out and shoots them in their 20-something times..
"And now that pastor and his family,.
"the pastor and his family who had to counsel people.
"all the time about how to deal with grief,.
"talk to them about how to deal with trials and tough times,.
"they are now the ones dealing with that grief..
"They are now dealing with that tragedy.".

$^{281}$In the weeks ahead, in the weeks now,.
I've gotten the chance to spend a lot of time.
with the family..
I sit with the daughter a lot,.
and we sit there on Friday nights,.
I see them quite a lot on Friday nights..
And she tells me about how life is now,.
that her father has now been murdered..
Her mother is now dealing,.
because literally her husband went to get a newspaper,.
walked three blocks to get a newspaper..
On his way home, he never returned home.
because he was shot..
And she's dealing with the grief..
And the daughter tells me, "Tommy, every day,.
"I remember the last conversation I had with my dad..
"I just replay the conversation over and over.
"and over again..
"And you know what makes it hard?.
"I'm trying to figure out and deal with not being angry,.
"not being bitter against certain cultures.
"or certain things that are happening..
"I'm trying to choose not to be angry at God.".
One of the things of being strong in our faith.
is that we continue to embrace compassion and love,.
and that we never, ever allow anger.
and resentment overcome us..
Whether it's the problems that we see in our world,.
or whether it's in our personal life,.
we continue to make a choice to embrace compassion and love..
And you never allow anger and resentment to overcome you..
And in today's society, that's a hard time, isn't it?.
For some of us who are sitting here as Asians,.
I can understand that you see some of the Asian violence.
going on in today's society as well, too..
Recently, I was driving down 35th Street..
Many of you guys come down 35th Street,.
you get off the highway..
I live about five minutes away from here..
I'm driving down 35th Street,.

$^{321}$and there's a car weaving in and out, in and out,.
and it gets right behind me and wants to turn..
But well, I have a no turn on red, so I am not turning..
A guy comes in, again, it's an African American guy..
And he comes right back to me, he rolls down his window,.
and he says to me, "This is why we effing shoot you people.".
And I was shocked, and I said to the Lord,.
and he drives off, and I left wondering,.
Lord, what does it mean to love.
the people who are hard to love?.
What does it mean for us to learn to forgive.
as you forgave, because you had to deal.
with the ultimate sacrifice..
Learning what it means to have compassion and love.
does not mean that we're a stepping stool for everyone..
It does not mean that people get to do whatever they want..
As Christians, we say, oh yeah, that's fine,.
we forgive them..
Being strong means that we love and have a compassion.
for everyone, but we do fight the injustices in our life..
That we do stand up and call out those injustices.
in our life..
I love, we just celebrate Dr. King, right?.
I love the words of Dr. King, and if you guys.
could put that quote on the screen as well, too..
Dr. King says this, I'm sorry about this,.
I should've made the font a little bit bigger,.
but I'll read it over to you..
He says this, "Wherever the early Christians.
"entered a town, the power structure of God disturbed.
"and immediately sought to convict them.
"for being disturbers of the peace and outside agitators..
"But they went on with the conviction that they were,.
"it was a colony of heaven, and they had to obey God.
"rather than men..
"They were small in number, but big in commitment..
"They were too God intoxicated to be astronomically.
"intimidated..
"They brought an end to such ancient evils.
"as infanticide and gladiatorial contest.".

$^{361}$Being a Christ follower means that we continue to love.
every single death, people with despite their skin color,.
despite their sexual orientation, despite all of that..
It means that we love people as Christ would love them..
But what it does mean is that we do stand up.
to those injustices that we see, that we continue.
to make a difference in the cities and structures.
that we make, that as Christ followers,.
that we make it a better place..
That our workplaces, our homes, our cities,.
our communities, whether it's in Chicago,.
whether those who are watching from other parts.
of the country, that you are integrating your faith.
into all that you do..
But you choose to do it with the spirit of love.
and compassion, and that you never make decisions.
out of anger and resentment..
And that's hard to do, isn't it?.
That is hard to do..
I recently had the board meeting.
for my organization Resource Global..
We just sat here..
And the board sat down with me and says,.
"Tommy, how's your family?".
Great, great, great..
But they said, "How are you doing spiritually?".
And I said to them, "I'm struggling spiritually.".
You know why I'm struggling spiritually?.
It's been hard to love people..
Because the more I sit in ministry,.
it's been now 15 years now,.
and the more I get involved in ministry,.
the more I deal with pastors,.
pastors that I've read in books,.
pastors that I've listened to on podcasts,.
the more I interact with Christians,.
they're hard to love..
They take advantage of you every single time..
They find ways to cheat you every single time..
And I've been so surprised by that..

$^{401}$And I said to the board,.
"I've grown disenchanted..
I have to choose to see beyond that.
and see Christ in everything I do..
And it's a choice that I have to make..
And it has not been easy.".
So for us, one of the things is be strong..
To really be able to continue,.
to really be able to love the Lord,.
to have compassion in all that we do,.
and never lead out of anger or resentment..
Secondly, show yourself a man..
Two things I'll say about that..
Show yourself a man is number one example,.
is be willing to take ownership.
of your mistakes and sins in a very humble way..
One of the ways of showing yourself a man.
is that we are all grown up..
It's not like my two-year-old daughter.
who doesn't like to share,.
and now I have to sit here and talk to her.
about what it means to share..
She's two years old..
She's learning..
All of us are adults..
It's being willing to take ownership.
of our mistakes and sins in a very humble way,.
knowing that we're not perfect,.
knowing that we're gonna make mistakes..
But when we make those mistakes,.
are we willing to take ownership of it?.
We look at David all throughout scripture..
A man after God's own heart,.
and he made a lot of mistakes..
And he had to take ownership of those mistakes..
Peter was always putting his foot in his mouth all the time..
Paul had to deal with the fact.
that he had to sit there and deal with Christians.
when he was a Gentile,.
and he messed with them all the time,.

$^{441}$and he persecuted all the time..
He had to live with that guilt.
after he became a Christ follower..
I have two friends, two very successful individuals,.
and these two individuals,.
married couple, husband, wife, they were Chinese..
They were extremely successful in the sense.
that they graduated from a prestigious university.
over in California..
Then they worked for a tech company..
Many of us used their computers..
They worked for a very extremely influential tech company.
and was very successful..
They then went to get their MBAs over here at university..
Again, and then from that point, after their MBAs,.
they got hired for an athletic shoe company.
that many of us wear their shoes,.
and also was in a very prestigious position..
The wife specifically oversaw all of the marketing.
for the World Cup,.
and eventually she was gonna be promoted.
for vice president of marketing in Asia..
Right, great story..
It all came crashing down..
Why?.
The husband got caught selling shoes on eBay..
A guy who earned over \$300,000 a year.
thought it would be exciting.
to make a little money selling shoes on eBay,.
and after a while, you get caught,.
and he got fired right away..
The wife who worked with them, HR called her in,.
and he said, "Did you know any of this.
"was going on with your husband?".
"Yes, I did.".
"At any point, did you actually report this to HR?".
"No, I didn't.".
"Hey, she had just gotten promotion..
"She just got all her papers get promoted,.
"vice president of marketing for Asia..

$^{481}$"They were gonna move to Hong Kong.".
He said, "You didn't report it..
"You're fired as well, too.".
Instantly overnight, everything that they had worked for,.
gone..
They struggled through their marriage..
They struggled with identity..
They struggled with the decisions they had to make,.
but they had to rebuild their life..
One of the things I've appreciated.
about this particular couple is they admitted.
with the husband, "I made some mistakes, some bad mistakes.".
How do we then rebuild our life?.
How do we humbly make amends for our mistake,.
and those were dumb mistakes,.
but rebuild our life and continue to go on?.
All of us are gonna make mistakes..
Are we humble enough to take ownership of our mistakes.
and sins in a very humble way?.
The other thing to be strong is this one..
Be willing to listen to voices around you.
because you're not perfect,.
and we all need accountability..
For Solomon, if he was to be strong,.
to show yourself a man,.
Solomon was the brightest of all the kings..
Solomon knew, and he called the shots..
Was he willing to listen to the voices around him.
because you're not perfect,.
and we all need accountability?.
I travel a lot for my work,.
and during the fall, I was gone a lot.
for about six weeks straight..
When I got home, I noticed a distance.
between I have a two-year-old daughter.
and a seven-year-old daughter,.
and I noticed that they were always going to their mom.
all the time, and they kinda ignored me..
My wife says to me, "Hey, Tommy,.
"you've been gone quite a lot,.

$^{521}$"and when you're at home, you also have meetings..
"You're out early in the morning..
"You have evenings, all that stuff..
"Have you noticed that the girls.
"aren't coming to you anymore?.
"You need to spend more time with your girls.".
You know what the easy, justifiable thing is?.
I've been to Nairobi..
I had to come back and do this and this and this..
When I come home, I have a late-night meeting.
with Jakarta and all that stuff..
Do you not know how hard I work?.
You have to deal with it..
Was there truth in what they had to say?.
I come back, and I meet with my team,.
and now we've gotten to a point with Resource Global.
where over 30 people, big, they sat down with me.
and says, "Tommy, as this organization continues to grow,.
"we need more structure, and you don't provide.
"that structure that we need.".
Again, you know what the temptation is, right?.
I'm the boss..
Guys, do you not realize I'm trying to make money?.
I'm out there meeting with donors,.
and I'm trying to make this thing survive..
Is there truth to what they have to say?.
Are they speaking any truth?.
Can I be humble enough to listen to them?.
And so as a result, I had to make changes..
I had to bring in different people.
and change their position and admit.
that there are certain areas I wasn't good at..
And be okay with that..
I say this all because sometimes it's hard.
for us to listen to voices..
And I will say this, for those who are older.
who are sitting there, sometimes the Chinese culture,.
the Asian culture, especially if you're older,.
says you've gotta respect your older elders,.
you gotta listen to your elders..

$^{561}$I encourage you, if you're older sitting there,.
learn to listen because when you work.
with that younger generation, they will actually respect.
and honor you more when you take time to listen..
I sit there and talk to my mom all the time,.
and she says, "Well, I want everyone to respect me.".
But her idea of respect is everyone just listens.
and does whatever she says..
But I tell her all the time,.
respect is that you communicate..
Respect means that you dialogue.
and be willing to admit you're wrong,.
that you're not always right..
And so for all of us sitting here,.
one, be willing and humble to make mistakes,.
but also be willing to listen to the hard things.
in your life, the times when you don't want to hear.
what people have to say, but is there truth.
to what they're saying about you and your character,.
even though it's hard to hear?.
Solomon had a hard time hearing that..
You know, eventually what was the downfall of Solomon?.
He got caught up too much in the love of women..
That became ultimately, his ego became so big.
that he stopped listening..
The one of the wisest kings who actually had an opportunity,.
I believe because of his talent, surpasses his father, David,.
but he refused to listen..
He refused to admit mistakes,.
and soon eventually gave in to sin as well too..
Last thing, be strong, show yourself a man,.
walk in his ways..
We have a couple minutes left..
I wanna show you a couple of slides here..
I show one slide here..
My friend Andy Crouch talks about this..
He says that as individuals, as human beings,.
we approach life with the exploitative attitude..
Our method in life is we wanna take all you can get..
Our attitude is I win, you lose..

$^{601}$We've been taught about that all our life, right?.
You gotta take care of yourself..
I hear that quite a lot..
You gotta take care of yourself..
And then once you take care of yourself,.
take care of your family..
Then whatever you have left over,.
then take care of everyone else..
And so we approach life with this method of,.
we're gonna take all you can get..
We gotta think about ourselves..
I win, you lose..
But Andy says this, and let's show the next slide..
He approaches life with this..
He says, as Christ followers, here's our attitude..
All right, oh, yes..
With the ethical attitude,.
the approach is do things right..
That our attitude is you win, I win..
Guys, let's keep that up there..
That we approach things with the whole idea.
is that you win, I win..
It is not a matter is I take everything I can..
It is the attitude of, it's always others first..
It is always about Christ loving other people..
It is always about you and I being a steward.
of what God gives me..
It's a choice that we make every single day.
to make sure that the other person wins,.
the other person succeeds,.
the other person is successful.
in everything they do at our expense..
But it's a choice that we have to make..
Solomon had to make a choice..
David had to make a choice..
David made a choice when he chose Bathsheba.
when he should have been out with the army battling,.
but he chooses himself.
and has a fatal mistake with Bathsheba..
Solomon chooses pleasure.

$^{641}$as a result destroys his kingdom as well too..
As Christ followers, being Christ in the world.
means you approach life with the ethical attitude..
That it is never about other, I mean yourself,.
but it is always about others..
It is a decision that you make, not only your heart,.
but with your mind to choose that.
in every single decision that you make,.
it honors God and glorifies God in every single way..
I would say this, and I say this with the young people..
If you're in your mid-20s.
and you're just graduating and working in a job,.
the decision that you're gonna have to make.
for those who are in their young mid-20s to early 30s,.
whose life do you belong to?.
Before you become rich.
and before you start ascending into the corporate world.
and all that stuff, whose life do you belong to?.
Are you gonna make a decision.
that in every aspect of your life,.
from your money to your career, to your dreams,.
your family life, is all given to the Lord?.
Because if you wait until you're successful,.
if you wait until things get stabilized,.
you'll always give in to the temptation.
of choosing yourself..
Being a Christ follower means that we make decisions.
that never benefit us..
We make decisions that glorify God in all that we do,.
that we are a steward of what God has given us.
at our expense, and it's a choice that you have to make..
One last illustration, and I'll wrap up..
I get a chance to travel a lot, all the time..
And one of the countries that I travel to is Nairobi..
I think I've shared a little bit of this illustration,.
but let me just build out this illustration at all..
When you go over to Nairobi, there are a lot of streets,.
everything is all remodeled, airports, streets, roads,.
highways, overpasses, everything..
It's all being funded by the Chinese..

$^{681}$The Chinese are giving so much money to build roads.
and infrastructures and technology..
But in the process of that,.
they are bribing other people, officials,.
and paying them so much money that those who are from,.
the Chinese are getting rich,.
but the African politicians are also getting rich.
at the expense of the people,.
because there are still a lot of people who are poor..
What you need is Christ followers in the workspace,.
Christ followers in the government,.
Christ followers integrating their faith,.
choosing others first, never about how much it can make me..
To make a decisions that benefits the country.
and our societies and our communities as a whole..
And that's what we need for all of you who are sitting here,.
Christ followers that makes a difference in a community..
That when you walk away, wherever you are,.
that in your workplace or in your home or in your cities.
or in your community,.
that community is left in a better position.
because you made decisions that didn't benefit you,.
but benefit them..
That's what it means to be Christ in the world..
That's what it means to make a difference in the world..
The cross of Christ brought us to where we are..
Christ died for our sins..
Christ is our example in all that we do..
Our job now is to continue to glorify,.
to be a steward of what he has given us..
Let's pray..
Father, thank you for all that you do..
Thank you for your words..
Thank you for your leadership..
Thank you for the cross..
Thank you for your grace..
As we remember the words that David gave to Solomon,.
be strong, show yourself a man, walk in his ways..
May we take those words and apply that.
in our life as well too..

$^{721}$Christ we pray, amen..
\newpage



\section{創世記}
\label{sec:wRAUIVuDJgw}
\textbf{創世記第3章 (3.1) - 背景釋經 (陳崇基牧師)}
\newline
\newline
連結: \href{https://youtube.com/watch?v=wRAUIVuDJgw}{\texttt{https://youtube.com/watch?v=wRAUIVuDJgw}} ~~~~ 語音日期: 2022-02-01
\newline
\newline
\hyperref[sec:gQBhlqNH6hM]{\small{< < < PREV SERMON < < <}}
~
\hyperref[sec:index_chronic]{\small{[返順時目]}}
~
\hyperref[sec:index_scriptual]{\small{[返順卷目]}}
~
\hyperref[sec:rx44uhcmvMM]{\small{> > > NEXT SERMON > > >}}
\newline
\newline
$^{1}$好,各位我們繼續看《創世記》第三章.
首先我們看第三章.
去分析一段經文的時候.
我們曾經提過.
再一次提醒大家.
有不同的視野可以幫助我們分析.
去看深入一點.
去看那段經文.
第一件事就是.
《聖經》背後的世界.
文化背景.
認知環境.
當時的人如何認知宇宙.
其他不同的等等的關係.
我們要去看看古近東的背景.
無論是看古近東的文學.
考古或者一些原文等等.
第二個視野就是《聖經》裡面的世界.
看看裡面的脈絡.
經文如何鋪排.
一些的題材和修辭的技巧.
往往修辭技巧都不是跟我們今天現代.
較為西化西方的方法去表達.
我們需要留意《希伯來聖經》.
如何去表達一些事情.
另外就是教會歷史.
整個群體去看一段《聖經》的時候.
看第三章.
有什麼值得我們參考.
一起去對話呢?.
首先提一提大家一些要點.
《聖經》無誤不等於我們對《聖經》的詮釋是無誤的.
從我開始認識《聖經》.
到我現在一直對《聖經》的理解.
對《聖經》的詮釋.
我都會一路前進.
希望能夠繼續改進.
所以我對《聖經》的詮釋不一定是絕對的.
一個牧者對一個經文的理解.
不一定是絕對不會錯的.

$^{41}$因為《聖經》才是無誤的.
不是我對《聖經》的詮釋無誤.
所以我對《聖經》的解釋.
詮釋繼續豐富.
甚至有修改.
這是對《聖經》的權威不會造成任何影響的.
第二就是我們說到《聖經》詮釋的時候.
《聖經》不是直接寫給我們的.
所以永遠都要回到《聖經》古舊的世界.
藉著不同的工具我們了解.
當時的原文的意思.
或者當時的背景等等.
避免用一些現代的.
特別是我們一些現代的字面的意思.
去解釋那段《聖經》.
就要認識當時的世界了.
原文等等.
甚至我們要去跟一些古近東的文獻做一些比對.
如果我們發現.
《聖經》跟一些古近東的文獻相似.
其實我們不需要驚訝.
其實我們預期應該是這樣的.
因為《聖經》當時的世界.
就是跟這些文獻.
當時的思維.
思潮.
去做爭辯,對話等等.
所以有相似的地方.
是我們應該預期的.
也是我們去比較.
去比對的時候.
發現有差異的地方.
往往上帝所啟示的超時空的真理.
就是永恆的真理.
就是在這種差異當中.
讓我們能夠看到.
所以第三方面我們需要留意的就是.
上一次講過了.
就是有不同的表述真理的語言.
Different kinds of languages.

$^{81}$所以我們現代人.
我們很著迷.
很著重,很強調.
就是一種科學性的語言.
甚至是一些很準確的.
歷史的一些語言.
以為這種語言才是表達真相的唯一一種方法.
沒錯,這種是能夠表達真相的.
但不是唯一一種.
而是有平常的語言.
太陽在東邊升起.
然後日落等等.
地上有四角.
這些都是平常的語言.
都是很準確.
以及很確實的.
可以有效地轉達真理.
包括詩歌的語言.
在這裡我要特別提醒.
特別強調一件事情.
就是在聖經是一個文學的時候.
我們不一定說因為文學.
所以跟歷史是對立的.
很多時候.
特別是上一代.
就以為你講到文學結構,修辭那些.
就一定是講虛構的東西.
聖經不是虛構.
當然這種了解對文學.
Literacy.
Literature.
這樣了解的時候.
是一種很膚淺的理解.
Again,又是一種以為科學的語言.
才能夠表達真理.
所以講到文學,修辭那些.
是跟歷史性沒有衝突和對立的.
正如上一節講過.
不同的語言可能表達的東西.
是更加豐富的.

$^{121}$比這些科學現代的語言.
表達得更加豐富.
好,我們看第三章.
第三章一到七節.
當然就是講到人的叛變.
蛇的引誘等等.
你看看蛇的說話.
是去質疑和顛覆神自己的美善.
去否定神的誡命.
上帝說你吃了知善惡恕的果子.
就必定死.
蛇就去否定.
然後蛇對女人的說話裡面.
對答裡面.
他是故意刪去了耶和華的.
我們講過二章開始.
神的名字就是 Yahweh Elohim.
耶和華神.
但是在蛇的口中.
就刪去了耶和華.
為什麼呢?.
耶和華是約的名字.
是神的子民跟神立約的時候.
我是耶和華神話.
耶和華就是一個 personal name.
一個名字.
Elohim 只不過是一個.
神只不過是一個 title.
一個名號而已.
真正的名字就是耶和華.
所以將它刪除的時候.
就是要人類去 set aside.
就是忘記.
遺忘.
就是他跟神的約的關係.
這就是蛇對人類引誘的本質.
分辨善惡.
再強調一次.
之前曾經提過.
不是吃果子這麼簡單.

$^{161}$吃果子只不過是代表一個更加真實的道理.
就是吃分辨善惡樹.
其實就是自我作主.
顛覆上帝的王權.
上帝不要再做神了.
我自己做神好了.
我去決定善惡好了.
所以是一個顛覆.
上帝本身有的 order(規律).
你知道在第一章已經有這個規律.
上帝很美善的.
很美好的來建立了.
就是神創造主.
君王.
然後人是被神的形象所造的.
就是創造界的君王.
君尊的祭司.
然後最底就是所管理的萬物大地.
但是在蛇的引誘故事裡面.
就是蛇將本身有美好的創造的規律.
就將它顛覆翻轉過來.
蛇在地上創造的活物.
他就說了算.
然後人受到試探.
叛變上帝的時候.
人其實就不是做神了.
其實就是做了蛇的僕人.
然後最底就變成了創造主.
誤視創造主為神.
另外我們看第三章的時候.
我們發現有很多的 wordplay(演繹).
一些的押韻字.
就是始音的字.
希伯來聖經經常用一些同音的字.
相似的字.
來連結著經文與經文之間.
第二章最後的時候.
我們看到.
阿當夏娃.
始祖.

$^{201}$赤身路體.
這個就是正面的赤身路體.
阿吾喑.
二章二十五字.
然後到第三章第一節.
又出現同音的字.
Aroom.
狡猾.
就是蛇是最狡猾的.
三章第一節.
所以兩個同音的字.
就將第二章和第三章連結在一起.
另外第三章第六節.
試探的本質裡面用了幾個簡單的詞語.
看.
好作食物.
畫.
好作食物.
就是good for food.
tove.
用了這兩個字.
然後這兩個字.
同音的字.
相同的字.
在第一章事實上出現的.
之前曾經提過.
神看著什麼是好的.
由他去決定.
但是試探就是.
人看見.
good.
好作食物.
人去自己決定.
就是自我作主.
自我為神.
三章第六節.
Tava.
月木.
看到果子是月木.
然後Tava這個字.

$^{241}$相似的音就是無花果.
無花果樹葉.
無花果.
Tena.
所以就是整個故事.
是串聯在一起的.
三章六節.
然後三章第七節.
有押韻字出現.
然後可喜愛的.
或者應該這麼說.
應該也叫做使人貪愛的.
三章第六節.
看到果子.
Hamad.
這個字出現在哪裡呢?.
就是在十誡.
最後的第十誡.
就是說不可貪戀別人的東西.
不可貪戀別人的房屋等等.
《心靈紀五章》.
再一次複述十誡.
第十誡也是用到這個字.
Hamad.
然後犯罪的時候.
就是看見是好的.
就拿來.
就吃了.
這是三章第六節.
這些關鍵字.
這些重要的字要記起來.
因為創世記經常用這些字.
六章開頭的時候.
講到神的兒子.
另外一個叛變的故事.
九章二十至二十八.
就是羅亞之後.
那個小的故事.
都是講到看好拿著.
這些關鍵字都會用.

$^{281}$所以特別在創世記一到十一章.
一個文學單元裡面.
經常會重複這些字.
在這裡我引用了兩句話.
在 Henry Poche 這本書.
In the Beginning.
他分析創世記開頭幾章.
一本很好的《色經》書.
雖然比較舊一點.
但是也講到.
第三章的故事是講什麼呢?.
這裡表明因為神和人類之間.
是有一個恩典之約.
人應該在恩典之約裡面.
上帝是主.
然後我們是僕.
有如創造主的兒子.
上主應有的特權是決定.
什麼是 good(好).
什麼是 evil(惡).
什麼是善.
什麼是好.
什麼是惡的.
只有神才有絕對的主權.
在上帝之下.
如果人類要得到真正的 Salon(滿足).
真正的喜樂滿足等等.
人必須要認定自己是一個悲造者.
僕人的身份.
這是我們的 identity(身份).
拒絕叛逆君主.
以及依附神的主權.
另外這個故事.
試探墮落的故事.
創世記第三章.
就是向我們展示了一個典型的情況.
一個典型的故事.
一個 Archetypal narrative(建築風格).
第二章已經提過這個概念.
就是叫 Archetype(建築風格).

$^{321}$亞當不是一個人.
而是代表人類的一個人.
是一個代表.
一個 Archetype.
所以所有亞當做的事情.
就是每一個人的故事.
夏娃也是一樣.
所以是一個 Archetype.
一個 model(模型).
一個典型.
在這個試探墮落的故事裡面.
就是可以讓我們從中得到一個每個人都適合的.
普及的教訓.
就是試探做什麼呢?.
試探就是嘗試在人類.
或男或女.
當中來引導出一種.
Henry Poseidon 說的一種女性的態度.
就是試探是想聽到蛇和女人的對話.
好像是被歡迎.
試探是被開放的一種心態來面對.
同情,同意等等.
試探呢.
他就說.
應該普及的教訓就是告訴我們.
不要用一個所謂女性的態度.
來面對.
而是抱著一個男性的態度.
來面對.
集中我們的火力.
集中我們的力量.
去跟這個試探正面交鋒.
刻意的.
堅決的.
來拒絕這個試探.
認定是上帝的誡明.
他的客觀性.
所以我們有一些這樣的 insight(觀察).
當然這個 Henry Poseidon 說的不是說.
重男輕女.

$^{361}$不是說男人比較厲害.
所以女人比較軟弱.
他只是比較 assume(假設).
就是 patriarchal culture(宗教文化).
聖經裡面背景的世界就是這樣.
一個父權的社會.
一個 patriarchal culture(宗教文化).
patriarchal society(宗教社會).
所以男人通常就是力量的表達.
女人就比較柔弱一點.
但是.
他這裡不是說男人好一點.
女人就差一點.
意思就是說用回當時的 culture(文化).
來表述出來的時候.
你會發現.
要去戰勝這個試探.
要去面對這個試探的時候.
就需要用一個當時所謂認定的男人.
masculine(男性化)的力量.
有其他的經文.
《耶利米哀歌》.
《耶利米書》.
甚至林前都會說到什麼.
我們去到得勝試探.
得勝序.
力量.
都是用一個男性化的一種刻畫.
所以千萬不要誤會他的意思.
以為男人比較厲害.
女人就差一點.
女人受誘惑.
所以他比較軟弱.
絕對沒有這種事情.
其實你看三章第六節.
我們經常誤解的.
現代人看一點就不看一點.
我們就以為是女人多嘴.
說女人的天性就是這樣.
所以搞到整個世界都瘋了.

$^{401}$你看清楚三章第六節最後.
一個很重要的字.
就是.
with her(她).
in ma(在我).
就是說整個試探的故事裡面.
你看到第六章最後.
所謂的adon.
就是和他在一起.
in ma(在我).
他也吃了.
所以所謂的adon.
就是那人.
本身應該做什麼呢?.
應該是保存.
守衛著聖所.
safeguard.
sacred place(聖地).
這就是祭司要做的事情.
去保存聖地.
保存上帝的同在.
所以應該是將邪惡的化身蛇.
就是將牠踢走.
是不是?.
應該將牠趕走.
但是誰知道牠是沉默地在身邊.
然後.
和他在一起的丈夫.
給了他果子.
他也吃了.
所以.
男人要負責.
女人也要負責.
所以每個人也要負責.
所以看到這個叛變的故事之後.
第八節開始.
就講到神的審判.
這裡提到.
說神在伊甸園裡面行走.
Mikhalak.

$^{441}$這個字.
是一個很特別的希伯來文文法的形式.
Hippie.
如果你懂得希伯來文的話.
在三章第八節出現.
也在五經其他地方出現.
這裡形容的是什麼呢?.
就是形容上帝在他的子民中間行走.
特別是在 Sacred Place.
在會幕聖所裡面行走.
你美記二十六章.
十一,十二節就最清楚了.
我要在你們中間立我的將王.
神說什麼呢?.
我心不厭氣你.
因為我救贖你.
我要在你們中間行走.
同一個字.
同一個形式.
我要作你們的神.
你們作我的子民.
因為我救贖了你們.
從埃及裡面領你們出來.
所以這裡講到上帝行走.
不只是在花園裡面走走.
而是行走就是.
他跟人應該有的關係就是.
在他們中間行走.
在人的中間.
神的子民中間.
神的同在.
所以這一句話.
或者動詞.
Hippeel.
Halak 這個字.
也在創世記其他地方出現過.
形容以洛.
跟神同行.
365年.
然後被神接去.

$^{481}$羅亞六章第九節.
講到阿巴拉罕等等.
都是講到.
跟神之間的約裡面.
我們的關係.
就是一種跟神同行.
聽他的律法.
忠於他的約等等.
我們看看有什麼細節.
我們需要留意的.
三章第八節.
這裡的翻譯.
通常.
英文也好.
中文也好.
都是天起了兩風的時候.
那人和他的妻子聽見耶和華在園中行走的聲音.
園民是講什麼呢?.
三章第八節.
Fires blew.
就是他們聽見.
就是They heard.
他們聽見.
聽見什麼呢?.
就是客體標記了.
At.
聽見什麼呢?.
就是Call Yahweh Elohim.
就是聲音.
哪個聲音呢?.
耶和華神的聲音.
這個聲音怎麼形容呢?.
Miphalek.
就是行走的聲音.
在園中行走.
Bagan.
Lawak.
Hayom.
就是我們譯作在風或者在零的那一天.
所以我們一般的翻譯.

$^{521}$現代的翻譯就說.
天起了兩風.
In the wind of the day.
這樣的意思.
吹起了兩風.
這一句話很特別.
就是Wuhak.
Hayom.
在希伯來的聖經裡面.
這兩個字一起用.
只有這一句.
三章第八節.
John世傑.
當然.
Wuhak.
是可以風或者零.
Hayom 就是日子的意思.
但是這兩個字加起來的時候.
是不是天起了兩風.
或者我們現代人聽起來的時候.
是不是比較浪漫一點.
好像在很漂亮的花園裡面行走一樣.
我們可以做一些比對.
雖然Wuhak Hayom.
沒有在希伯來聖經其他地方出現.
但是Wuhak Khor.
就有出現過.
Khor 就是聲音這個字.
和Wuhak.
就是在其他的經文出現過的.
耶利米蘇十章或者五十一章十六節.
就有提過.
他說他一發聲.
Khor.
這個字.
空中便多水.
去到激動.
神使到雲霧從地極上騰.
他造雷電.
隨雨而閃.

$^{561}$在他的虎腹那裡帶出風來.
Wuhak.
Right?.
所以你看到Wuhak Khor.
也有出現過.
Yom.
當然我們希伯來文.
一聽到Yom的時候.
就發現那一天.
就是日子.
什麼意思?.
Day.
但是在阿卡德文.
Akkadian.
就是Wuum這個字.
這個字可以有日子的意思.
也有暴風的意思.
Storm的意思.
就是說神明出現的時候.
他私人審判的時候.
就會有Storm出現.
這個字典可以讓我們知道.
有不同的意義.
這個CAD就是 Chicago Assyrian Dictionary.
在第二十冊.
Chicago Assyrian Dictionary很出名.
因為在芝加哥大學.
University of Chicago.
有一個近東考古的部門.
在中東做了百多年的考古工作.
所以他們出版了一個叫做CAD.
Chicago Assyrian Dictionary.
可以上網下載.
在第二十冊就講到Wuum這個字.
所以是有暴風的意思.
審判來到的意思.
以賽亞書二十七章第八節.
西方的二十七章二節.
都有這樣的用法.
所以可能一個更加好的翻譯.

$^{601}$貼切的翻譯就是.
第八節這裡講到.
就是上帝的審判終於臨到.
他們聽見.
耶和華的咆哮.
可以這麼翻譯.
就是他的雷聲.
他的聲音.
在弦中是吹動.
如同暴風來到一樣.
Wuak.
Haium.
所以第八節是一個審判的場景.
去到刻畫出來.
這裡講到人的叛變.
人的墮落之後.
就由本身上帝的原創.
神創造人類.
是不是?.
人與人之間應該有order.
是和諧的harmony.
神與人之間都是一樣.
人類與大地.
都是一個和諧.
一個shalom.
有神的祝福.
有神的同在.
亞當這個字.
人類這個字.
Adam.
其實就是地這個字.
也是同音字.
Adamah.
Adamah 就是地的意思.
所以人與地.
天論人論地論.
都應該有一種shalom.
這種shalom不是說沒有戰爭.
沒有張力.
就是shalom.

$^{641}$shalom是一切都在上帝的原創.
他的order裡面.
他的規律裡面.
但是人犯罪之後.
人叛變之後.
這一切.
所有的關係.
都受到破壞了.
甚至被完全的顛覆了.
所以在三章第八到第十九節.
講到審判的宣告的時候.
有痛苦.
有戀慕管轄.
就是兩性之間.
人與人之間.
就是講到之間的張力.
之間的糾紛.
有爭戰.
甚至之後我們看到有謀殺.
有暴力出現.
地受咒詛.
以往地就是為人效力的.
現在不再效力了.
人本身是塵土.
這個就是他的 identity(身份).
again 講到塵土.
不是講到一些生物學的問題.
講到我們的 identity(身份).
我們是會扭圍的.
那麼就注定會是終身路苦.
並且歸於塵土.
所以這段經文的目的.
John Wharton 在他的《壯世的識經》中說.
就是去解釋.
告訴我們人類.
個體也好集體也好.
都是流亡於伊甸園之外.
失去了神的同在.
以色列人聽到這個故事的時候.
就知道.

$^{681}$it's not supposed to be this way(這不應該是這樣的).
原本不應該是這樣的.
原本應該是這樣的.
但是現在是這樣的.
所以第三章就是去解釋這個現實.
但是在第三章裡面.
頒布這個審判當中.
不僅僅是一些負面的東西.
好像給我們一個盼望.
好像答應一個的應許.
因為耶和華對神所說的話.
裡面帶有一些跟終末有關的東西.
或者跟將來有關的東西.
他跟蛇說.
他說你做了這事.
就必受咒助.
比一切的生畜.
也受更甚.
你必用土之行走.
終身吃土.
當然這些也是.
就是.
羞恥的表達.
就不是真的一條蛇.
正如果果不是真的.
市場裡面的水果.
而是說到.
邪惡的化身.
就是最終會是終身吃土.
最終會失敗.
我又叫你和女人彼此為仇.
你的後裔和女人的後裔.
都要彼此為仇.
女人的後裔要傷你的頭.
你要傷他的腳筋.
這個後裔這個字.
就有很多的討論.
因為女人的後裔就是單數.
我們很多時候就會.
立刻跳到一個結論.

$^{721}$一個色經結論.
就是.
為什麼是單數呢?.
因為只有一位.
那一位是誰呢?.
最終就是耶穌基督了.
這是從文法引申出神學.
那麼.
最終可能這個也不是錯.
但是.
單單從文法引申出神學.
就不是最嚴謹了.
為什麼呢?.
因為後裔這個字.
單單在創世記裡面.
都是有單數和複數的意思的.
28章14節.
這裡跟雅各說.
他說你的後裔必定在地上的灰塵那樣多.
你.
單數.
必定向東南西北擴展.
地上的萬俗會因你和你的後裔.
是不是?.
複數.
單數複數一起用.
但是原文是單數.
所以很難單從文法來說.
這個 singular(單一).
所以將來只有一位.
往往很多時候是一種 collective singular(集體單數).
一種集體的單數.
就是指整個神的子民.
所以無論如何.
這個女人的後裔.
可以指她或者他們.
但是在基督教傳統的解釋裡面.
我們很多時候最後得到的結論.
就是一種尼賽亞的 interpretation(解釋).
Messianic interpretation(妙妙的解釋).

$^{761}$就是說最後是耶穌基督了.
是不是?.
傳統的教會就說這段聖經.
女人的後裔.
這段聖經.
就是元初的福音.
Proto-Evangelium.
就是 proto 就是元初.
第一個.
Engelion 就是福音.
福音派就是 Engelion 這個字出來.
福音.
所以是第一個的福音.
這個初期教會的領袖.
我們叫 Church Fathers.
其中一個就是 Justin Martin.
在朱後一六零年的時候.
他就是這麼去解釋的.
但是正如我所說.
單從希伯來文化去決定是不能夠的.
不如看看夏文.
怎麼去延續這個故事.
女人的後裔.
你就發現.
《創世記》怎麼延續.
就是加譜.
有兩個加譜.
兩條線出來.
就是一個是失特.
就是京拜耶和華的後裔.
失特.
然後一路繼續跟著五張的加譜.
就到了羅亞.
然後閃.
到十二章.
亞伯拉罕.
就是阿不蘭.
十一二章那裡.
然後一路繼續下去就是立約的子民.
亞伯拉罕.

$^{801}$以撒.
雅各.
以至以色列加.
出來了.
神的子民出來了.
所以繼續 trace(追溯)這個 story(故事)的時候.
後裔發現原來就是以色列國.
就是上帝的選民.
另一個加譜.
就是離開了上帝的.
就是 King.
這個該人.
到拉麥.
所以你看完整本聖經.
《創世記》以至超越《創世記》.
整本的聖經下用了一個所謂聖經神學.
Biblical theology.
意思就是說到了追溯.
在聖經由《創世記》到《啟示錄》.
整個大的敘事裡面.
這個 meta narrative(內容)裡面.
某一些的主題.
到底是怎麼發展的.
這個就是聖經神學.
Biblical theology 要做的事情.
如果你 trace(追溯)這個所謂後裔的主題的話.
你一直 trace(追溯)下去.
你發現原來就是在說神的子民.
然後到最後當然耶穌基督來了.
教會就是神的子民.
繼續延伸下去.
所以這是一個很特別的所謂的認許.
很奇怪.
舊約裡面從來沒有直接引用過這一段聖經.
新約也很少看過.
最明顯的一處.
只不過是羅馬書最後結束的一句話.
十六章第二十節.
賜平安的神快要把撒旦踐踏在你們腳下.
這個應該就是引用.

$^{841}$就是跟《創世記》第三章女人的後裔的認許.
有所互動.
引用它.
但是很明顯.
沒有說到一個很清楚的.
Messianic interpretation(妙妙的解釋).
就是這個女人的後裔就是基督.
當然.
教會的歷史.
就是這麼去認定.
這個也沒錯.
但是這裡.
至少第三章告訴我們一件事情.
就是雖然人叛變.
罪闖入世界.
周作有死亡.
但是最終.
是神結語的.
神的認許.
神的說話.
他自己救贖的工作.
是那個結語.
我們看看第三章這裡形容神的審判的時候.
這樣描述的時候.
你就發現有一些這樣的.
《希伯來聖經》特別在《創世記》經常出現的.
就是這些所謂三文字的結構.
三章九至十二節.
就是耶和華神跟男人的對話.
然後就是跟女人的對話.
三章十三節.
中間夾住的就是三章十四十五節.
就是女人後裔的認許這段經文.
就是耶和華神對蛇說的.
你發現.
他跟蛇說話的時候.
沒有提問的機會.
跟男人有.
跟女人有.
但是跟蛇呢.

$^{881}$所謂邪惡的化身.
是沒有任何提問.
沒有任何對話的.
他只是宣判.
他的結局.
最終會是失敗.
然後呢.
三章十六節.
又回到跟女人說話了.
說到她生產的痛楚等等.
然後最後呢.
十七至十九節.
又說回那人.
他們不再為她效力等等.
為什麼要提這件事情呢?.
就是一些所謂工具.
一些概念.
讓我們知道.
就是聖經記載一些事情的時候.
他的手法.
聖經裡面的世界.
很多時候不是我們比較西化.
現代人的方法.
因為我們最重要的重點結論.
我們通常放在最後的.
但是希伯來聖經.
他的敘事手法.
很多時候是放在中間的.
一種三文字的結構.
這裡就是一些概念.
讓我們能夠理解.
或者讓我們看聖經的時候.
縱使我們未必看原文.
但是我們會想一想.
可能那個重點不一定在最後.
當我們看到聖經記載的脈絡.
脈絡.
有時候不像我們現代人那樣想.
直線的.
很多時候又回頭.

$^{921}$再看其他的東西.
我們就會想.
可能那個重點不是在尾段.
而是在中間.
有一本的聖經書就是 Adel Berlin.
一位很出名的希伯來聖經學者.
他在修辭與傳識.
傳識聖經的 narrative.
Biblical narrative.
這本書裡面講到一個很重要的概念.
他說 You don't know what the text means.
until you know how it means.
就是說除非你明白經文表達意義的手法.
How.
這個就是How.
ABCBA.
12321.
這種三文字的結構.
How.
這種修辭的手法.
除非你明白這件事.
否則你是不會理解聖經的內容的.
你要知道聖經.
What it says.
首先很多時候要掙扎的就是.
How it says.
如何去將這些鋪排.
以至我們能夠更加掌握到聖經的意思.
那麼.
我們暫時講到這裡.
記住.
最後這裡.
惹禍華臣.
他的宣判.
他的說話.
是作一個結語.
不是邪惡的化身.
也不是罪.
不是人的破碎.
而是上帝自己的工作.

$^{961}$上帝自己的話語.
做一個結語.
如果你喜歡我們的影片.
就記得要Like,Comment和Share.
\newpage



\section{創世記}
\label{sec:rx44uhcmvMM}
\textbf{創世記第3章 (3.2) - 背景釋經 (陳崇基牧師)}
\newline
\newline
連結: \href{https://youtube.com/watch?v=rx44uhcmvMM}{\texttt{https://youtube.com/watch?v=rx44uhcmvMM}} ~~~~ 語音日期: 2022-02-01
\newline
\newline
\hyperref[sec:wRAUIVuDJgw]{\small{< < < PREV SERMON < < <}}
~
\hyperref[sec:index_chronic]{\small{[返順時目]}}
~
\hyperref[sec:index_scriptual]{\small{[返順卷目]}}
~
\hyperref[sec:Q90kQe4p0_k]{\small{> > > NEXT SERMON > > >}}
\newline
\newline
$^{1}$我們繼續看下去.
就是第三章的後半.
20-24節.
三章20節.
這裡說什麼呢?.
亞當及他的妻子起名叫夏娃.
夏娃其實是希伯來文.
夏娃.
我們翻譯為夏娃.
是希伯來文.
正如亞當是希伯來文.
亞當.
夏娃夏娃.
因為她是眾生之母.
大家發現第八至十九節.
就在說死亡,審判,奏坐等等.
神的審判.
怎知道到二十節.
突然間有一個口吻上.
就是語氣上突然間改變.
仔細看我們就明白了.
之前說奏坐.
死亡審判.
現在突然間.
很突然的.
轉了.
亞當為他的妻子起名.
這是一個信心的形訊.
Confession of Faith.
我講講布魯斯沃奇.
布魯斯沃奇是舊約的太斗.
他已經很老了.
一輩子都是研究舊約的.
布魯斯沃奇說什麼呢?.
他說亞當為夏娃命名.
是盼望的開始.
Beginning of Hope.
雖然故事裡面充滿著死亡,審判.
痛苦的生產.
人間的衝突等等.

$^{41}$但卻存一絲盼望.
就是女人的後裔與蛇敵對.
必定會戰勝邪惡的化身.
這就是一絲盼望的出現.
是不是?.
剛剛我們講到一些字眼.
甚至整個的ABCBA.
就是中間那段.
耶和華神對蛇的審判.
不提問的.
只是給他一個審判咒助.
這裡有一個戰勝邪惡化身的一絲盼望.
在這裡也看到上帝的恩典.
上帝的慈愛.
剛剛三章八至十九節.
全部都是講一些所謂負面的事情.
很黑暗的事情.
死亡等負面的事情.
但是二十節開始.
亞當的認順.
他的信心.
Faith.
他的信靠.
就是說這個女人.
Again是所有的女人.
就是說人類繼續下去.
Hawa就是生命.
繼續到下一代.
繼續在鬥爭裡面.
與蛇的邪惡.
爭戰裡面.
就會展開.
並且有一絲盼望.
我們看看過去一些的釋經.
一些的觀察.
三章二十一節.
這裡一直講耶和華為人做了什麼.
為叛變的人.
救贖他們.
怎麼施展他的恩典.

$^{81}$三章二十一節.
他做了就是耶和華神.
Yahweh Adonai.
Fire Ash.
Yahweh Elohim.
Ladam.
Uishto.
Ladam Uishto.
就是為亞當與他的妻子.
為那人與他的妻子.
做了什麼呢?.
就是後面這裡.
Isa.
Kadno.
Or.
被甲做的衣裳.
Fire Bishim.
為他們穿裳.
這是一個很仔細.
很有恩慈的動作.
所以為他們做了.
用被子做衣服.
給他們穿.
相對於他們自己為自己做.
犯罪的人.
叛變了,離開了神的人.
他們為自己做什麼呢?.
就是無花果的樹葉.
這裡我在以色列.
其中拍到的無花果樹.
到處都有.
只有這麼大而已.
一隻手掌這麼大.
他們自己為自己做無花果樹的樹葉.
作裙子.
遮自己的醜陋.
遮自己的羞恥.
神這裡是為人做.
他出手為人做.
並且是一些耐用很多的東西.

$^{121}$無花果樹葉.
以及被子.
Bruckman.
另外一位很有名的舊約的色經家.
他在《創世記色經書》裡面說.
判決之後.
神就為夫婦.
來做一件他們不能為自己所做的事情.
他們不能處理自己的羞辱.
但神能夠.
神亦願意.
神並且實行出來.
人不能夠為自己做什麼?.
就是遮自己的罪孽.
遮蓋自己的羞恥.
無花果樹葉.
一下子就沒有了.
很快就扭歪了.
但是神就用被子來做.
這當然是有一個意味的.
大家都很清楚.
就是 sacrifice(犧牲).
就是殺了一隻動物.
然後做被子.
當然不是很清楚在這裡說.
也不需要這麼清楚說.
因為最主要在這裡說什麼呢?.
就是主體.
主體是誰呢?.
是神.
神為叛變了他的人.
是不是?.
他認信了.
他就為他安排所謂的救贖.
用被子來遮蓋他們的罪惡.
那麼.
如果你看原文的話.
這個被子這個字.
就是 all.
怎麼讀呢?.

$^{161}$all.
o-r.
all.
有一個拉比.
拉比是很值得我們看的.
看看拉比怎麼識經.
其實很有意思的.
因為他們真的比較真實.
每一字每一句.
每一個音都比較真實.
有一個拉比叫做 Meyer.
他就是主侯.
大概二世紀左右.
他就在加利利附近.
現在去加利利都會有一個紀念的墳墓.
是為 Rabbi Meyer 來到去.
在那個地點.
有些人過時過節就會去.
猶太人就會去那裡紀念他等等.
他就將這個經文裡面說.
就是.
Gatno.
all.
他就解成另一個 all.
下面這個是什麼呢?.
你看到下面這個 all.
這個寫法.
跟上面那個差一點點.
差第一個字母.
本身這個馬蘇拉抄本.
就是舊約聖經.
就是用這個 iron 開始.
讀音是什麼呢?.
都是 all.
因為第一個字母是沒有音讀的.
是不發音的.
那麼他呢.
Meyer 就將它寫成 all.
讀音是完全一樣的.
但是他就轉了它做 aleph.

$^{201}$他說將這個皮子第一個字母 iron.
就轉成 aleph.
讀成 all.
意思就是說光明的以上.
為什麼拉比會做這件事呢?.
我們先不要管他到底怎麼做.
是不是對的.
就是這樣.
去改一下聖經的字眼.
然後有其他意思出來.
這樣是不是合法呢?.
另一件事.
他為什麼這麼做呢?.
就是因為他覺得很突然.
看完聖經的時候他覺得很突然.
為什麼耶和華會為叛變他的人.
去做一件這樣的事情.
並且呢.
做一件事是去幫他們遮蓋.
他們的羞恥.
所以他一聽到這個字.
Canon.
all.
的時候他就想起.
不是皮子.
他想起光明的以上.
光明的以上.
那麼可不可以成立呢?.
或者也有它的道理.
因為當時聖經不是每個人都有一本去看的.
是口傳的.
當時是 all culture(全文化).
就是聖經只有幾本.
是有一個讀經的人在會堂讀出來的.
聽的人呢.
只是用他的發音來聽經文.
不能夠眼睛去看.
所以他聽到.
Canon.
all.

$^{241}$的時候就是光明的以上.
所以他覺得這裡不僅是說皮革這麼簡單.
是說到 garments of light(光明的衣服).
就是給人的聖潔.
就是救贖.
帶給人的尊榮.
他本身只是羞恥.
現在上帝給他們的救贖.
是一些光明的以上.
是不是?.
當然這個未必是合法的做法.
但是我想告訴你.
看經文的時候.
我們不要錯過了.
那種 surprise(驚喜).
就是那些突然.
語氣上突然.
是不是?.
就是阿當應信.
之前就死了.
周祖.
但是到了這裡.
耶和華為人去到預備的.
居然是一些很尊榮的事情.
是不是?.
很有恩典的事情.
很有慈愛的事情.
三章二十四節.
這裡說什麼呢?.
就是.
耶和華神.
快要趕出.
將阿當趕出伊甸了.
他安設.
Phaistion.
他安設在東邊.
Makedim.
在伊甸園的東邊.
安設了什麼呢?.
Again 我們又看到這個at這個字.

$^{281}$就是一個客體的標記.
這個客體的標記標記了什麼呢?.
就是基魯柏.
發火焰的火箭等等.
劍等等.
他就把守著那個生命樹的道路.
是不是?.
他就安設了這件事情.
這個動詞.
就是Phaistion.
Phaistion這個動詞.
就是一個史譯文法的形式.
什麼意思呢?.
用英文翻譯的話.
這裡就是cause to dwell.
就是使一些東西.
史譯嘛.
使其他東西.
就放在那裡.
cause to place.
就是將一些東西放在那裡.
這個解釋就是說.
將基魯柏.
發火焰的劍.
轉動的劍.
就放在伊甸園的東邊.
東邊就是聖所的入口.
上帝就安設了.
所以我們中文也叫做安設.
是對的.
所以整句意思就是在伊甸園的東邊安設.
史譯.
安設什麼呢?.
史譯其他東西.
就是基魯柏.
He placed on the east side of the Garden of Eden.
Tribute.
那麼又要回頭想想.
希伯來聖經本身寫的時候.
希伯來文是很妙的.

$^{321}$本身是沒有響音的.
只不過到了中世紀的時候.
這些拉比.
猶太的拉比.
文士.
就是想保存希伯來文的讀音.
它的原文.
所以就將這些一橫一點等等.
就將它點上去.
本身希伯來文寫的時候是沒有音標的.
你是要從上下文.
來揣摩出來的.
那麼在拉比的文學裡面.
很多時候.
他們會有一些傳統.
就會有一個音標的標誌方法是不同的.
而那個音標不同標音的時候.
就會有不同的意思出來.
其中一個拉比的文學叫做.
Tarigan.
Targum.
Targum裡面有四段.
就是對第一個動詞的音標.
是有不同的方法的.
什麼叫做Targum呢?.
Targum就是到了主後.
一至七世紀左右.
在猶太人的會堂裡面.
因為當時的猶太人.
未必懂得聽希伯來文.
所以就需要有一個人去做翻譯的.
那個人就將希伯來文翻譯成亞蘭文.
以至當時操亞蘭文的猶太人聽得懂.
就是希伯來聖經.
這個就是Targum.
在Targum裡面有四段.
就將這個音標.
不是標做Faiz Khan.
就標做Faiz Khan.
OK?.

$^{361}$不是Faiz Khan.
是Faiz Khan.
那麼那個意思呢?.
就有分別了.
因為不是再是一個使役的動作.
而是一個簡單的動作.
簡單的動作就是說.
不是說耶和華放一些東西在伊甸園門口.
這裡所形容的就是一個簡單的動作.
耶和華自己做的動作.
耶和華做了什麼呢?.
Faiz Khan.
就是他居住.
他居住.
那麼這裡的意思就會變成.
耶和華在伊甸園的東邊居住.
在基魯柏和四面轉動發火焰的劍意上.
把守生命樹的道路.
那麼下面這一段學術文章我給你.
他就追溯這四段古老的希伯來翻做亞蘭文的Targum.
他就說這裡.
有一個可能忘記了的意義.
可能這個希伯來文有雙重的意義.
另一個意義就是在Targum重新找出來.
就變成了意思是什麼呢?.
耶和華離開自己的伊甸園.
離開自己的聖堂.
與罪人同住.
這個就很震撼了.
就是說不是把他丟出去.
然後關了門.
不讓他回到裡面.
耶和華寧願離開自己的聖堂.
伊甸園.
住在哪裡呢?.
住在伊甸園的東邊.
He settled east of the garden.
離開了自己的Garden.
就與人類居住.
這個大概看上下文.

$^{401}$你都會明白這個意思.
因為第四章發生什麼事呢?.
哥恩阿爸.
耶和華是不是就這樣關了門.
不理他.
原來耶和華撇棄了自己.
就是聖所.
榮耀的聖所.
聖潔的聖所.
與人類居住.
與他們重新建立關係.
這是上帝一向想做的事情.
是不是?.
住在人當中.
住在他的子民當中.
所以.
完了第三章你發現.
耶和華去了哪裡呢?.
他不是躲在自己的聖所.
他是出來了.
然後他繼續去建立.
救贖他的子民.
建立他自己的聖所.
是不是?.
由伊甸園.
我們追溯到與阿爸顯現.
伯特利等等.
上帝向人顯現.
到會望.
然後到聖殿.
然後到耶穌基督.
這個神永遠都想在人的中間.
神的子民中間.
來同在.
所以這裡給我們一種很震撼的訊息.
臨尾的時候.
原來本身我們說.
咦?.
他就是把守著生命道路.
生命是道路.

$^{441}$將他拋棄了.
不是的.
原來是神離開了自己的聖所.
好了,這裡可以最後做一些比對.
正如我開頭所說.
可以比對一些近東的文獻.
到底是怎麼去形容神與人之間的關係.
如何凸顯聖經的那些特別的訊息.
其中一個近東的文獻叫做.
Epic of Gilgamesh.
Gilgamesh Epic.
基加美修史詩.
這個史詩是一個在古近東的世界.
很出名的史詩.
很出名的小說.
最述的就是一個叫做烏克魯王.
基加美修.
Gilgamesh.
他是一個歷史人物.
大概是公元前2500年左右.
在一個古城Uruk的地方.
Uruk在哪裡呢?.
Uruk就在近著梧椅.
我先點出Uruk.
如果你看看Uruk在哪裡.
Uruk就在這裡.
為什麼我們要看地圖呢?.
因為.
《創世記》1-11章的背景就是美索不達米亞.
美索不達米亞的意思就是.
Mesopotamia.
就是兩河流域的意思.
兩河流域就是.
Tigris.
底格里斯河.
和Euphrates.
幼法拉底河.
這個地方.
因為全世界.
在人類歷史裡面.

$^{481}$第一個有文明的地方.
不是在美國.
不是在歐洲.
不是在中國.
而是在這裡.
在美索不達米亞.
那些叫做蘇美人等等.
我們有時候會看到歷史書.
但是你看到.
《創世記》1-11章.
裡面所講的東西.
或者回應的東西.
往往就是這些美索不達米亞裡面.
文獻和文化的背景.
所以我們講到這個史詩.
就是講到King of Uruk.
在大概主前2500年左右.
這個小說講什麼呢?.
內容大概就是講王.
Gilgamesh.
很威猛的王.
他去尋找永生的歷程.
他很怕死.
知道人的命運一定會死的.
所以他就去到攀山越嶺.
跟他最好的摯友.
來克服很多事情.
希望找到永生.
最後失望而回.
認定人一定會死的.
死就是人的命運.
結論是什麼呢?.
這個Gilgamesh epic.
就是說總之享受現在的人生.
盡情享受.
因為我們最後都會死的.
並且神明造人的時候.
在講這些巴比倫的神等等.
中神創造人類的時候.
中神就將死亡命定給人類.

$^{521}$卻將永生.
生命留在自己手中.
你看到這些字句的時候.
你就看到那些所謂的外邦.
他們的文獻的神學.
他們也有神學.
他們的宗教神學.
就是.
人.
被創造.
不僅無緣無故.
就是做奴隸.
是沒有尊榮的.
沒有神的形象.
甚至呢.
神明造人的時候.
其實就是不想人好.
死亡就是命定給他們.
永生就留在自己手中.
這就是近東其他經外文獻的人論.
如何去論人.
人性.
人的地位等等.
你這樣去比對的時候.
你就會發現.
《創世記》第三章.
一二三章裡面所講的.
是完全不同的.
人被造有神的形象.
有尊貴.
並且在神的同在裡面.
還做祭司.
在伊甸裡面.
有生命樹.
可以選擇.
所以人是有.
神是給永生.
是給人的.
只不過人因為叛變.
是自立為神.

$^{561}$是不是?.
分辨善惡樹的果子.
所以就變成了.
正如耶和華吩咐亞當一樣.
就是說什麼呢?.
你摘那一隻果子的時候.
你必定死.
你的收場就是一定會死.
但是呢.
其實神創造人的時候.
祂的心意是什麼呢?.
生命樹.
是不是?.
之前可能大家提過了.
啟示錄最後.
Again 又是 Biblical theology.
就是聖經神學.
你去追溯某一個主題的時候.
你發現最後會發展成什麼樣子.
最後新天新地又是什麼呢?.
就是生命樹.
是不是?.
跟近東的文學.
是自然不同的.
所以無論你怎麼去欽標祂也好.
這裡說的是什麼呢?.
上帝是願意與人同在.
願意在恩典.
在慈愛裡面.
去救贖人.
好了.
最後就是布什維奇所說的.
他在墮落的狀態裡面.
就是一個人墮落了.
有罪了.
人絕對不能夠永活下去的.
所以不可以再摘生命樹的果子.
因為這樣會變成一種咒詛.
永遠與神離不開.
死亡呢.

$^{601}$所以為什麼會有死亡.
不讓他有生命樹的果子.
死亡既然是一個審判.
也是一種釋放.
不會永恆地與上帝交絕.
死亡.
布什維奇在他的《舊約神學》裡面說.
死亡將人類從永恆靈性的死亡.
就是肉身的死亡.
將人類永恆的靈性的死亡解脫出來.
對亞當和夏娃來說.
他們吃了哪個果子.
就是分別性樹的果子之後.
靈性變成死亡.
死亡就是擺脫咒詛的一種方法.
肉身的死亡.
就是開闢了一個超越墳墓超越死亡.
肉身死亡的一個永恆的救贖之路.
因此.
肉身的死亡可以是我們的救贖.
以至我們不會被判於這個墮落世界裡面.
繼續延長我們的生命.
上帝最終想做什麼呢?.
與人同在.
賜給人.
是永生.
所以看完整個三章的時候.
撞死的第三章.
你會發現.
說了很多人的墮落.
人的背叛.
叛逆上帝.
然後上帝去審判.
花了很多的折數.
去說出上帝對罪的審判.
咒詛.
死亡等等.
但是到最後三章.
最後結案陳詞.
居然就是上帝的恩典.

$^{641}$就是上帝仍然願意與有罪的人同在.
他自己救贖的計劃.
仍然繼續在這段以外開始.
展開之後.
我們發現一章又一章.
《創世記》.
甚至整本《聖經》.
一章又一章.
到最後的時候.
神的恩典.
神的同在.
神的得逞.
是結案陳詞的.
你看《創世記》開始到最後.
整本《聖經》的結束.
新天新地啟示錄.
就是神在人間.
神的罪惡莫在人間.
神的同在.
神的生命殊等等.
最後結案陳詞的是上帝自己.
得勝的是上帝自己.
如果你喜歡我們的影片.
就記得要讚好,留言和分享吧.
\newpage



\section{箴言}
\label{sec:Q90kQe4p0_k}
\textbf{逆要作門徒 (4) — 再思箴言 (陳崇基牧師) 2022.01.30}
\newline
\newline
連結: \href{https://youtube.com/watch?v=Q90kQe4p0_k}{\texttt{https://youtube.com/watch?v=Q90kQe4p0\_k}} ~~~~ 語音日期: 2022-02-01
\newline
\newline
\hyperref[sec:rx44uhcmvMM]{\small{< < < PREV SERMON < < <}}
~
\hyperref[sec:index_chronic]{\small{[返順時目]}}
~
\hyperref[sec:index_scriptual]{\small{[返順卷目]}}
~
\hyperref[sec:q8ZgN14Imcg]{\small{> > > NEXT SERMON > > >}}
\newline
\newline
$^{1}$各位,我們這次繼續思考智慧文學.
《聖經》裡面有不同的書卷.
都是叫智慧文學.
我們看過《約伯記》,詩篇,傳道書.
這次我們看《箴言》.
《箴言》真的是一個智慧文學.
並且是一本選集.
什麼意思呢?.
就是說將不同的智慧言語收集在一起.
變成一本書.
總共是三十一章.
想起《箴言》智慧的時候.
就一定想起《所羅門王》.
因為寫了很多的《箴言》.
但是如果你看《箴言》書的話.
裡面不僅有《所羅門》的《箴言》.
也有其他外邦的君王的一些《箴言》.
特別是最後那幾章.
什麼叫做《箴言》呢?.
通常我們拿起《箴言》這本書的時候.
就會看到有一些很特別的格式.
一兩句,幾句.
加起來一個單元.
就是智者所說的話.
有一點點像我們華人文化裡面的格言.
或者一些短句.
我們今天講這篇讀到看《箴言》的時候.
正是快要苦年的時候了.
那麼.
我們新年的時候就很應景了.
因為我們新年就專門講一些這樣的祝福的話.
我們看看聖經裡面怎麼說智慧.
什麼叫做智慧的生命呢?.
Wisdom在聖經裡面的定義.
有一點點像我們今天講的EQ.
情緒的雙數.
就是我們怎麼去管理我們自己的情緒.
我們的內心.
我們的關係等等.
你知道這件事情也很重要.

$^{41}$EQ.
特別是在職場上.
大家知道.
不是IQ高就夠了.
最重要是有EQ嘛.
所以《箴言》有一點點像這件事情.
所以我們首先將《箴言》裡面的一些話.
拿出來思想一下.
並且看看整本《箴言》.
怎麼去理解.
以至我們能夠聽到《箴言》對我們的挑戰.
第一段我想跟大家去思考的就是二十六章十三至十六節.
這裡說什麼呢?.
他說冷惰人說.
路上有獅子.
不是老虎.
廣場上也有獅子.
冷惰人在床上翻動.
就像繞著門勾轉動一樣.
冷惰人將手放在盤子裡面.
連拿食物回口邊也覺得勞累.
冷惰人看自己比七個善於應對的人更加有智慧.
這裡說什麼呢?.
這裡就說到一個愚昧的人.
冷惰的人.
那麼.
中間那兩句.
就去描述他一些行動.
懶散.
躺在床上.
就像門勾一樣自轉.
轉來轉去也去不了哪裡.
他也是維持在原狀.
冷惰的人不想出力.
甚至前面有食物也不想拿.
可能要等人餵他.
這些都是嘲諷的手法.
頭尾這裡說到冷惰的人.
內心的想法.
他說外面有獅子.

$^{81}$外面有猛獸.
所以我就不想出去了.
意思就是說.
不僅是不想出門這麼簡單.
就是說我不想出去工作.
因為在浸燃這個世界裡面.
義者.
義人.
Righteous.
就是願意去工作並且是公認人的.
那些不是為了自己賺錢.
而是為了其他人去祝福群體的.
那些叫做義人.
反過來.
那些愚昧的人.
或者冷惰的人.
只是為了自己.
從自己的角度去想而已.
這些就是浸燃這一處去諷刺的.
並且最後那一句.
他還以為自己很聰明.
比七個智者更加聰明.
他覺得自己不用生活.
貪小便宜.
有人餵他.
自己不需要去勞力工作.
不用去生產等等.
你看我多聰明.
其實就是這樣諷刺他.
因為等於我們中文的格言就是.
濫牛屎尿多.
濫人明天多.
經常都想著明天.
這些就是很典型的一些浸燃.
想當頭棒的也可以提醒我們.
二十四章了.
這一段我就是很喜歡的.
因為這裡可以告訴我們.
到底浸燃裡面的智者.
跟其他人才知道等等.

$^{121}$有什麼不同.
以及他不同的地方在哪裡.
什麼叫做有智慧呢?.
我們看看.
二十四章三十節.
智者說什麼呢?.
他說我經過濫暴人的田地.
無知人的葡萄園.
只見到到處長滿窒礙.
地面長滿刺草.
很多野草在這裡.
石牆也倒塌了.
智者說我看見.
我就用心思想.
我看見的時候.
就領受了教訓.
教訓是什麼呢?.
在睡片時打頓片時.
抱手躺臥在片時.
你的貧窮就如強盜來.
你的缺乏就如拿兵器的人來一樣.
這裡有什麼特別的地方.
大家記得先知是怎麼領受上帝的話嗎?.
先知很多時候被聖靈感動.
然後帶他到天上.
他看到天上的景象.
上帝給他異象和訊息.
然後他下來.
直接跟神的子民說話.
耶和華如此說.
這就是先知.
智者呢.
The wise.
有些不同.
他沒有什麼看異象的.
但是他觀察世界.
他說他看見.
看見懶惰人的田地.
他觀察到一些東西.
就是看見人生百態.

$^{161}$他就領略到一些生命裡面的教導.
所以智者在這裡.
他觀察到什麼呢?.
他觀察到自然界有一個定律.
就是石頭永遠都是向地下.
向下跌.
就像水往低流一樣.
是不會反過來的.
石頭不會從地上無緣無故地上來的.
是不是?.
他看見這些田地.
這些城牆,石牆.
是日久失修.
看見只會向著一個方向破壞.
就是說他領略到,觀察到.
這個世界裡面有一股力量.
一些邪惡的力量.
是將我們生命裡面的秩序.
變成混亂的.
永遠你的 order(命令).
都是會自動變成 chaos(混亂)的.
是永遠不會回頭.
不會說由 chaos(混亂).
混亂.
變成有秩序的.
所以你看看你家就知道了.
是不是?.
你看看你家.
你沒有收拾多久.
你沒有收拾東西.
你懶散的時候.
就自然.
不需要你做什麼.
秩序就會變成混亂.
智者就領略到這個教訓.
從他觀察裡面的自然律.
他就看出一些更深層的東西.
所以這個智者不是一個 natural scientist(自然科學家)這麼簡單.
就是我們今天叫做 observation(觀察).
觀察到世界裡面的一些定律.

$^{201}$不是停在那裡.
他想出教訓.
他就領受了一些教導.
智者是一個屬靈的教士.
一個生命的教士.
就是他觀察了之後.
他領略到生命都是一樣.
領略到就是我們生命.
他在說我們的屬靈生命.
是永遠從秩序走向敗壞的.
從 order(命令)走向 chaos(混亂).
你不需要做任何事情.
孜孜然就會是這樣的.
你想想我們的屬靈生命.
豈不是這樣.
貧窮會自然來.
並且搶著來.
混亂會自然來.
死亡自動會來.
但是.
除非你相反.
相反你去到下功夫.
你去到修理.
不像這個懶惰人一樣.
你去到發展他.
你去到成長他.
你去到栽培他.
那你就會阻止這個邪惡的勢力.
來敗壞你.
我們要下功夫.
才能夠維持秩序.
如果打頓騙事的話.
我們不要理他了.
我們屬靈生命先放在一邊.
沒事的.
沒問題的.
我們就一定會走向敗壞.
你只能夠加能量下去.
有下苦功.
然後我們才能夠生命裡面有平安.

$^{241}$想想我們的個人.
我們的工作.
我們的家庭.
我們的事奉.
神和我們的關係.
其實都是一樣.
一定要加上我們的功夫.
我們要肯勤於鍛鍊.
我們才能夠有規律.
能夠領受基督裡面有的祝福.
是不是?.
所以智者就有這個分別.
觀察完之後.
他看到世上的百態.
我們應用在我們自己的生命裡面.
所以有時候.
有很多人看《箴言》的時候.
就會覺得很 secular(宗教).
很世俗.
聖經裡面有一本這樣的書.
我以為聖經裡面是很神聖的.
其實好像我們中文說.
勤有功棄無益.
有這種感覺是很自然的.
為什麼呢?.
因為在《箴言》裡面.
確實沒有很明顯提到以色列的歷史裡面.
一些重要的事件.
或者重要的體制.
比如說出埃及.
聖殿耶路撒冷敬拜.
節旗.
聖約的律法等等.
這些很多時候都不提.
所以給我們的感覺是有點 secular(宗教).
有些世俗.
就是跟我們這些華人文化的格言有什麼分別呢?.
叫我們做好人.
但是要留意.
有很多地方是有提耶和華的名字.

$^{281}$耶和華這個名字就是一個聖約的名字.
就是神自己的 personal name(個人名字).
就是說.
智者去看世界的 order(命令)的時候.
他不僅僅是一個唯物的觀念去看世界.
有這樣的規律等等.
是在耶和華的約底下.
去看見生命裡面有的規律.
所以有很多地方會提到耶和華的.
譬如頭九章.
頭九章是一個序言.
第一章第七節.
這裡就是我們很多時候會記.
或者會背的一段經文.
經威爺會說就是智慧的開端.
知識的開端.
慾望人就是藐視神給他的智慧和教訓.
鑒然.
有一個 setting(設定).
有一個處境.
就好像一個父親教兒女一樣.
所以就用一些智慧的說話.
去告訴他.
搞清楚你跟神的關係.
在你成長的當中.
特別是你成為聖人的時候.
你需要首要的第一件事.
你需要擺第一的.
就是你跟耶和華的關係.
為什麼呢?.
因為父親教他的兒子的時候.
兒子快要成家立室了.
以前十幾歲的時候.
就會成家立室了.
就會成為聖人了.
生命裡面需要做很多的選擇.
最重要的選擇是做什麼呢?.
所以就智者的父親.
在第一章到第九章.
就盡力地勸兒女.

$^{321}$去聽智慧的說話.
九章.
一個所謂的序言.
我們看看第九章說什麼.
到最後第九章.
就是序言的結束的時候.
第九章第一節.
智慧建造自己的房屋.
作成七根柱子.
他宰殺牲口.
調和美酒.
然後擺設筵席.
他差不多派幾個侍女出去.
自己又在城裡面高處呼喊.
叫山下的人做什麼呢?.
瑞士愚昧的人.
就是年輕的人.
或者屬靈生命不夠穩定的人.
可以到我這裡來.
請你接受這個邀請.
然後到第十節.
他說驚威耶和華是智慧的開端.
認識智慧者便是聰明.
你藉著我的智慧.
你就有生命了.
日子必定會增多.
並且你的年歲會加添.
這些都是舊約裡面經常有的.
就是 lengthen your day.
有生命了.
就是你要去聽智慧.
就是跟智慧去建立關係.
這裡就是第九章.
這個序言.
那裡講到介紹智慧.
有一個以人法的方法去介紹.
智慧這個字.
在希伯來文.
hochma.
其實是一個音性的名詞.

$^{361}$大家知道.
希伯來文 Hebrew.
是跟很多其他歐洲的語言.
譬如法文.
他們的 noun 名詞.
是有分陽和陰的.
男和女的.
智慧這個字的剛巧.
就是音性.
就是 she.
就是女士.
女性的她.
所以這裡.
我們的譯法.
就是把女性的她放在這裡.
我不太明白為什麼《和學本》沒有譯.
或者《新人本》譯成了.
她是一個死物的她.
但是原文是 she.
就是說這個智慧.
就像一個女士.
她就預備現席.
在房屋那裡.
邀請人進來.
然後.
她在城術裡面.
高處的地方去喊叫.
所以聖經在《箴言》.
智慧形容成為一個女士.
英文我們叫她 woman wisdom 或者 lady wisdom.
因為希伯來的原文 chokma.
是音性.
所以這裡的手法.
就是說這個女子.
擺好現席.
要請這些年輕人.
沒有智慧的人.
有一個 hospitality.
這是文學的手法.
為什麼會說成這樣呢?.

$^{401}$用智慧.
說到智慧.
形容成為一個女子.
這個很自然.
剛剛我所說.
一個父親教育一個兒子.
因為他要成家立室.
他就是想什麼呢?.
要娶老婆.
他想他的配偶.
所以兒子想起要配偶的時候.
就想起女子.
這裡就用了一個文學手法.
就是勸沒有智慧的人.
你要娶智慧.
就是說你要跟智慧.
有一個緊密的關係.
這樣你才會有生命.
所以你這麼看的時候就明白了.
為什麼在最後的箴言.
又說回一個女子叫做.
蔡德的夫人.
或者賢德的婦女.
其實就是一個首尾呼應.
說的不是真的說男和女的問題.
說的是生命裡面.
我們要做一個很重要的 decision(決定).
一個決擇.
我們要選取智慧.
要棄絕愚昧.
第九節這裡介紹這位智慧的女子.
Lady Wisdom.
你再看下去的時候.
又有另外一個女子出現了.
九章第十三節.
愚昧的婦人.
另一個了.
不是Lady Wisdom.
我們可以叫她Lady Folly.
愚昧的化身.

$^{441}$也是一個婦人.
她就在那裡宣揚.
她是愚昧的.
一無所知的.
坐在自己家門口.
坐在城的高處.
坐位.
呼叫過路的人.
走直路的人.
告訴她不要走直路.
誰是愚昧的人可以轉到我這裡來.
又對那些無知的人說.
偷來的水.
我也預備了現值.
偷來的.
偷來的水是甜的.
暗吃的餅.
偷來的餅.
多好啊.
十八節.
人若卻不知道.
有陰雲在他那裡.
他的客裡.
他的客人.
是陰間的心存.
就是說Lady Folly.
跟上面那個Lady Wisdom.
是完全相反的.
上面那個就是比心明.
下面這個.
似乎比心明有大排現值.
呼叫我們進去她家裡.
但卻是死亡和陰間.
有兩條路讓你走.
兩個人.
你可以去做一個回應.
一個是智慧.
一個是Folly.
就是愚昧.
留意這兩個女子.

$^{481}$有什麼比對.
九章四字.
講到智慧.
Lady Wisdom.
她所說的話的呼喚也是一樣.
誰是愚昧的人.
可以轉到我這裡來.
Lady Folly.
這個愚昧的女子.
同樣也是說同一句話.
叫那些人.
去到她那裡.
不要走你的路.
轉過來我這邊.
然後我們再去比較.
兩者都是有房屋的.
九章第一節.
智慧就建造房屋.
House.
Her house.
可以叫人去到她家裡.
九章十四節.
Lady Folly.
她坐在自己的家門口.
都是房屋.
是House.
坐在城的高處.
又去到呼喊.
更加重要的.
兩個女子的位置在哪裡.
智慧的女子.
九章第三節.
她派幾個侍女出去.
自己在城的高處呼喊.
愚昧的Lady Folly.
也是坐在自己家裡面.
坐在城中的高處.
同一個字.
同一個位置.
去到呼喊.

$^{521}$在什麼人.
或者高處是什麼意思.
如果你熟悉舊約的世界.
近東的世界.
你就知道城最高處是什麼地方.
每個城的高處.
都是聖殿廟所在的地方.
神的住處就是高處.
如果看耶路撒冷.
耶路撒冷城.
就是建在兩個山.
一個東一個西.
東山的最高處是什麼.
就是聖殿.
就是耶路撒冷的聖殿.
就是敬拜上帝的聖殿.
就是神的住處.
那個就是城裡面的高處.
如果你再想想.
有很多其他地方.
高處也是一些假神的廟宇.
一些偶像所建立的廟宇.
都是建在城最高處的地方.
無論是巴黎也好.
什麼假神也好.
所以這裡所說的.
智慧.
和愚昧.
其實就是在說神和偶像.
智慧的女子.
其實就是神自己.
就是耶和華.
是一個比喻.
比喻這位就是耶和華.
耶和華在聖殿裡面向我們呼喊.
叫我們去敬拜祂.
這就是智慧的開端.
人生智慧的開端.
是從敬拜認識神開始.
但是假神.

$^{561}$你們的巴黎.
偶像.
都在祂的高處向人去呼喊.
兩個的呼喊.
其實是兩個生命裡面的呼召.
一個是真的神.
主向我們的呼召.
一個是假的偶像.
對我們的呼召.
時時刻刻在我們生命裡面.
就是要我們做一個決定.
智慧還是愚昧.
所以說是一種屬靈的決定.
屬靈上去敬拜耶和華.
去事奉耶和華.
這就是智慧人生的開端.
如果我們不聽這個呼喊.
聽了愚昧女子的呼喊.
聽了假神的呼喊.
我們只會生命走向死亡.
所以第九章第十節.
這一段.
Lady Farley和Lady Wisdom.
中間第十節.
就找到我們平時經常看的一段經文.
就是敬畏耶和華是智慧的開端.
認識智性者便是聰明.
這是屬靈的挑戰.
挑戰我們到底是敬拜耶和華.
還是敬拜偶像.
到底我們是擁抱智慧.
還是我們擁抱愚昧.
到底我們是走向上帝給我們的生命.
這個豐盛又平安的生命.
還是我們走向死亡和陰間.
這是一個很重要的.
我們時刻都要去警醒.
做一個抉擇.
沒有選擇的.
沒有其他的選擇.

$^{601}$只能選擇A或B.
只能選擇Wisdom或Farley.
我記得我年輕的時候.
在大學的時候.
當時大學團體有一個導師.
經常和我們栽培和查經.
開始一兩次的時候.
我們幾位弟兄姊妹都覺得.
很新鮮.
但是久了之後.
就覺得有點悶.
因為習以為常.
偶爾有一次.
導師不舒服.
他就說取消.
我們就很開心.
我們很高興.
好像中了彩票一樣.
真是好了.
我們不用去看聖經了.
不用去做栽培了.
我們可以做我們自己的事情.
今天回頭想.
其實是很愚蠢.
幸好上帝的恩典.
憐憫.
我記得有一次.
有一個神學生剛畢業.
好幾年前和他聊天.
他說很辛苦.
終於讀完碩士學位了.
可以不用再讀書了.
我已經把神學的教科書全部丟掉了.
求主憐憫.
為什麼會認識神.
是一件苦差.
為什麼我們會將認識神.
放在一些可有可無的位置.
這些就是愚昧女子的呼喊.
如果我們聽了這些呼喊的時候.

$^{641}$生命一定會走向敗壞.
智慧的女子.
就向我們呼喊.
就是說你要去到謙虛虛心.
你也要下功夫.
你也要追求認識自成者.
這個動作.
認識真神.
去到敬畏祂.
去到事奉祂.
才會走向生命的道路.
最近好幾個星期.
我和一些創啟地區.
Creative Access的地方.
一些的牧者.
網上聚會.
我就講了一些色經的課程.
錄了.
甚至在城會網台那裡.
我錄了那一段.
之後我們討論的問答那一段.
我們就沒有錄了.
所以我錄了那些.
我們就在城會網台來發佈.
我發覺這些牧者.
真是很有心.
雖然他們在一個地方很不安.
有很多打壓.
有很多 censorship(刑事).
有很多禁聲.
但是很虛心學習.
甚至他們說.
可不可以再多一點時間.
加時.
這些牧者全部都比我更年輕.
二十幾三十歲.
每一個我看到都很有前途.
肯去認識上帝的說話.
肯去下功夫.
雖然他們沒有這樣的資源.

$^{681}$但是很願意來學.
有很多人說.
我們現在這個世界這麼亂.
有這麼多不安.
不如我們只說一些實際的東西好不好.
跟我們貼身一點的教導.
這些不是不好.
當然這些是應該的.
我們城會其中一個的核心價值.
勇毅真實和緊貼時代.
是需要 relevant(相關)的.
但是你想想.
看箴言的時候.
你細心思想一下.
可能最 practical(實際)的東西.
就是認識神.
認識神才可能是最 practical(實際).
最實際 最緊貼時代.
不是說其他 topic(題目).
我們不需要去鬥.
怎樣去到.
過一個好的情緒生活.
好的關係等等.
怎樣教養兒女.
這些當然是很重要.
但是箴言告訴我們.
其實是不是最 practical(實際).
對我們的生命.
能夠回應這個世代.
最重要的就是認識上帝自己.
認識自性者.
留意箴言怎樣結束.
三十一章.
三十一章講到.
才德的妻子.
第十節開始.
一直到最後.
有二十二節.
其實是字母詩.
由 Aleph 到 Tav.

$^{721}$每一個希伯來的字母開始一句.
我們很多時候是母親節的時候.
我們就會說.
仿效吧.
這裡是一個才德的婦女.
其實你看看這二十二節.
才德婦女的詩.
其實不是人做的事情.
人本身不能夠直接照字面去應用.
不可能的.
其實是什麼呢?.
是一個智慧生命的刻畫.
也不僅僅是女性.
正如我說的.
是一個擬人化.
是一個智慧生命的表現.
才德的女子.
誰能得到她的價值.
是高於諸葛.
高於雄葛石.
你留意才德的妻子.
這句話.
在希伯來聖經.
舊約聖經.
只出現過兩次.
這是其中一次.
另外一次在哪裡出現呢?.
就是在路得記.
路得記也是最後三章十一節.
那裡說路德是賢德的女子.
翻譯不同.
但是原文.
其實是一樣的.
說到才德的婦人.
說到路德.
你記得路得記的歷史背景是什麼嗎?.
第一章第一節路得記.
就是事先秉正的時候.
國度最亂的時候.
國人行自己看為正的事.

$^{761}$最不安的時候.
就說到路得記這個故事.
說到一個智慧的化身.
其實不是一些大的英雄人物.
是一些小的國色.
波瓦斯.
路德.
劉亞米.
小的家庭.
怎麼盡忠.
不是為了自己而生活.
而是為了其他的人.
是不是?.
一個智慧的化身.
路德就是一個才德的女子.
我們記得.
我們的現代中英文聖經.
我們的目錄.
路得記放在哪本書後面.
就是士師記.
所以第一章第一節就是.
史詩秉正的時候的路得記.
但是希伯來聖經.
猶太人用的聖經.
我們的舊的聖經.
路得記放在哪本書後面.
路得記放在《箴言》之後.
《箴言》之後.
看完《箴言》.
你就問.
什麼叫做一個智慧的生命.
猶太人的目錄就說.
你看看路得記就知道了.
所以《箴言》之後.
就說到.
才德的婦人.
你就要看看一個智慧的生命.
就是怎麼樣.
就是這一幫小角色.
他們不安在亂世當中.

$^{801}$仍然有勇氣.
去忠於大家.
忠於耶和華的約.
為著其他人.
是不是?.
肯思念不暈.
肯到幫助其他人.
去建立那個群體.
為了那些已死的人.
留名.
將那個家系繼續延續下去.
這就是路得記的故事.
是不是?.
然後最後的時候.
路德.
波亞斯.
居然就變成了.
大衛的祖宗.
神的國就因為他們這些小角色.
而去到成就在地上.
認識自成者.
辨識聰明.
智慧的生命.
在不安亂世當中.
才是最實際.
在基督教的神學著作裡面.
一本最經典的著作.
如果你還沒看過的話.
請你有時間去看看.
翻譯成中文的叫《認識神》.
《Knowing God》這本書是J.I.Pekka所寫的.
這位神學家.
以前在Region College教系統神學.
J.I.Pekka在《Knowing God》的第一章.
開頭的時候說什麼?.
他說認識神.
他說只有認識神的人.
真正認識主的人.
不是knowing about God.
是knowing God.

$^{841}$認識神的人.
才可以在不安在亂世.
甚至受迫害的時候.
有勇氣為主而活.
他就提到.
但以利.
但以利三個朋友.
就是在最不安.
在亂世.
在被擄的當中.
他們有勇氣來敬拜神.
為什麼?.
因為他們真認識神.
真是認識至聖者.
就是智慧的生命.
亂世不安亂世.
認識至聖者.
才是最有.
最重要的.
使我們生命扎實.
想到智慧.
我們就想起所羅門.
你記得.
所羅門.
作了很多真言.
所以為什麼.
真言說到智慧的時候.
很多時候都會提到所羅門.
但是你記得《律子書》裡面.
說到所羅門.
榮華富貴.
然後最後晚年是失敗.
國度分裂.
為什麼會失敗?.
為什麼會一敗塗地?.
智者.
給我們什麼提醒?.
我翻過真言很多地方.
我覺得這一句是最到肉的.
第十九章第十七節.

$^{881}$我們以這一節來結束.
他說:我而.
你若是停止.
聽受管教.
就必會偏離.
知識的言語.
一個基督徒.
怎麼可以走得遠?.
豈不是我們每一個人.
豈不是我們每一個階段裡面.
我們拒絕.
去到說.
Been there, done that.
我聽過了.
為什麼要再學?.
還有什麼我可以去學習的?.
什麼都懂.
一個基督徒走得遠.
豈不是拒絕這種.
Foolishness.
這種原會的想法.
自以為是的想法.
虛心的去學習.
就是說我今天仍然要去到.
認識自成者.
所以我的生命.
能夠遇見不同挑戰.
不同不安的時候.
我都能夠有勇氣.
有力量.
為主而活.
這句其實不像耶穌.
在登山寶訓.
最後結束的時候說的話.
智慧的人是怎樣的?.
你記得.
將他的房屋建在.
盤石之上.
就是那些聽了主的話.
肯去到遵行的.

$^{921}$那些人.
就雨林水沖來到的時候.
都能夠抵擋得住.
愚昧的人就將自己的房屋.
建在沙土上面.
雨林水沖.
一切一敗塗地.
就是因為停止聽受管教.
停止去虛心學習.
去認識主.
兩個的選擇.
一個宅.
一個闊綠.
一個智慧.
一個愚昧.
求主幫助我們和我.
讓我們在一個.
縱使不安的時代裡面.
我們真真正正去認識主.
認識自成者.
以致我們心靈有智慧.
有力量.
我們得求祈禱.
主求你帶領我們.
我們每一位.
我們在智慧裡面成長.
讓我們認識你.
為我們自成者.
認識你自己的律法.
認識你自己的言語.
讓我們真真正正的.
知道你自己.
我們和你的關係.
我們能夠緊密.
並且我們的生命.
有著智慧.
有著力量.
你都可以為你而活.
祈禱奉主名禱.
阿們.

\newpage



\section{創世記}
\label{sec:q8ZgN14Imcg}
\textbf{創世記第4章 (4.1) - 背景釋經 [比對研究的圖解,近東宇宙觀] (陳崇基牧師)}
\newline
\newline
連結: \href{https://youtube.com/watch?v=q8ZgN14Imcg}{\texttt{https://youtube.com/watch?v=q8ZgN14Imcg}} ~~~~ 語音日期: 2022-02-02
\newline
\newline
\hyperref[sec:Q90kQe4p0_k]{\small{< < < PREV SERMON < < <}}
~
\hyperref[sec:index_chronic]{\small{[返順時目]}}
~
\hyperref[sec:index_scriptual]{\small{[返順卷目]}}
~
\hyperref[sec:xeNKKwLqznQ]{\small{> > > NEXT SERMON > > >}}
\newline
\newline
$^{1}$我們今天繼續看創世記第四章.
我們會看一些原文和一些工具.
和大家分享一下或者示範一下如何使用.
其中一個工具我曾經說過是一個比對研究的工具.
Comparative Studies.
是甚麼來的呢.
Comparative Studies即是說.
我們明白聖經有自己的context(背景).
有它的背景.
Comparative Studies 比對研究其實就是拿回聖經當時古代的背景.
一些資料.
以致我們能夠去到參考的時候.
明白聖經對當時的讀者當時的作者.
當時的世界首先有甚麼意思.
以致我們不會那麼容易用我們現代人的觀念.
讀進去聖經.
第一課的時候都已經說了.
有很多這些所謂的地方.
不敢說是陷阱.
一些地方很多時候我們會掉進去.
以為我們讀了一些東西是聖經的意思.
但其實我們將我們現代人的意念.
概念讀進去了.
所以比對研究其實是一個工具.
那是怎樣去做的呢.
我這裡就引用兩個學者.
兩個學者都是美國的學者.
兩個學者都是美國一些很重要的福音派的神學院的學者.
這兩位的學者我這裡引用的.
都是一些可信的學者.
他們對聖經的權威不會質疑.
對聖經的可靠性.
對聖經是神的話語是上帝的默示.
這些不會質疑的.
為甚麼我要引用他們呢.
有時候人們未接觸過一些古近東的背景.
一些以前有的考古找到的資料.
他們會覺得為甚麼要找一些這樣的資料來跟聖經相提並論.
會不會貶低了聖經的權威呢.
聖經在我們心目中是很神聖的.

$^{41}$因為他是最高的權威.
我相信我們大部分的人聽到的時候.
我們作為基督徒或傳道同工.
我們絕對尊重聖經.
聖經對我們來說是神聖的.
所以我們先解釋一下這個工具的使用.
我先引用這兩位的學者.
Dwayne Garrett.
Dwayne Garrett是美國南方浸信會的舊約學者.
舊約的教授.
他在這本書中提到Rethinking Genesis.
即是重新再思創世記.
他在開頭提到方法論的時候.
考古學將古近東的文明恢復.
找到一些泥板或文字.
是當時人寫的.
是當時以色列聖經的世界裡的人寫的.
對聖經的研究有很大的意義.
在當代的美索不達米亞,迦南地,埃及的一些研究.
拿來研讀聖經的時候.
聖經的歷史,詩歌,智慧和語言等等.
就會產生一些新的清晰度.
我們更加有Clarity.
更加清楚他們想說些什麼.
稍後我會引用他多一點.
另外一位就是我其中一位老師.
叫做John Wharton.
他現在在Whitman College.
維敦神學院教授.
他這本書是講到古近東的思想.
和舊約的研究.
這本書已經是第二版.
當時我剛讀考古學的時候.
就讀John Wharton的一些課程.
他剛剛才開始出這本書.
那時候是2009,2008年的時候.
他說什麼呢?.
我認為忠誠的聖經讀者.
可以找到很多不同的理由.
我們不需要毫不關注,比對,研究.

$^{81}$不需要用這東西.
但是一個忠誠的聖經讀者.
是不能夠忽視比對研究對解經的重要性.
如果我們不先考慮古代認知環境的資料.
他叫做認知環境.
就是說當時的聖經世界裡面.
人對認知宇宙,認知世界的環境.
他們怎麼想事情.
這些資料.
如果我們不考慮的話.
我們現代的讀者.
就會自動將現代世界觀的參數.
強加於聖經的經文之上.
從而有一個很嚴重的曲解聖經的意義風險.
我先引用這兩個學者.
因為這兩個學者都是很紮實的科幻派的信徒.
我們說古近東的研究的時候.
沒錯有很多不信聖經的人.
非基督徒.
他們會從事這些研究.
但是他們的前設.
很多時候我們會打一個問號.
但是有很多福音派的聖經學者.
其實已經在用這些資料.
看看這本書什麼時候出版.
這本書第二版John Wotten.
2018年.
沒多久之前出版.
正如我所說第一版的時候.
2009年的時候出版.
上面這本Rethinking Genesis.
2001年的時候已經有了.
所以是20年前的.
已經有這些資料.
所以這些絕對不是新的東西.
希望大家不要以為.
這些很新的很奇怪很新奇的方法.
絕對不是新的.
這個比對研究已經存在了很久.
並且由開始.

$^{121}$這些考古學家去到古近東世界.
去找到這些文物.
那時候是什麼時候呢?.
18世紀末.
19世紀的時候.
那時候可能不是那麼成熟.
現在已經非常成熟了.
100多年的辯論.
100多年的研究.
所以已經很清楚.
到底這些古近東的文獻.
到底是怎樣幫助聖經的解釋.
這裡有一個很簡單的圖.
我將什麼叫做comparative studies.
嘗試去表達出來.
最大的圈圈.
橢圓形的就是近東的世界.
古近東的世界.
為什麼叫古近東呢?.
在這裡大家先溫習一下.
古近東就是.
其實今天我們指的中東.
由波斯灣.
至以色列巴勒斯坦地道.
埃及這個地方.
我們叫做今天的中東.
或者以前叫做古近東.
或者學名叫近東.
為什麼會叫近東呢?.
你知道如果你在中國的話.
你就叫做遠東.
什麼叫遠東?.
什麼叫近東?.
從哪個角度看?.
由歐洲的國家.
歐洲人的角度來看東方.
當時發生什麼事呢?.
大概18,19世紀的時候.
當時管這個地方的帝國.
叫做Ottoman Empire.

$^{161}$奧多曼帝國.
由土耳其.
管到巴勒斯坦地中東.
以及下面的埃及.
19世紀的時候.
奧多曼帝國.
土耳其的伊斯蘭帝國.
他們開始色彌.
開始瓦解.
所以對於那時候的強國.
這些歐洲的強國.
他們看到東邊有兩個政治的問題.
兩個世界政局的改變.
第一個當然是近著他們的東方.
就是奧多曼帝國.
開始瓦解.
西方的勢力就開始滲入這個地方.
所以在西方人.
那些歐洲人的眼中.
這個地方就叫近東.
Near East.
另外一個就是.
大家熟悉近代史.
就是甲午戰爭.
他們也意識到.
這些歐洲人也意識到.
這個地方也有一些政治的變遷.
所以在歐洲人眼中.
就有所謂遠東的問題.
和近東的問題.
所以我們學名字.
我們會說.
這個地方就是聖經的世界.
就叫近東.
所以這裡先溫一溫.
這個就是近東的世界.
近東的世界.
近東的世界有它自己的文化.
有它自己的思維的方式.
或者看宇宙的方式.

$^{201}$裡面有些什麼人呢.
譬如有巴比倫.
亞述.
埃及.
等等.
不同的迦南等等的人.
他們有他們自己的一套的意識形態.
去看聖經.
去了解這個世界的方法.
另外就有聖經.
就是古代的以色列.
或者舊約.
舊約聖經.
他們也在近東的世界裡面.
所以.
也是近東世界裡面.
其中一個的民族.
以色列人.
以色列人和他們的鄰國.
譬如巴比倫.
亞述.
埃及等等.
所以他們有很多地方.
是會相似的.
譬如有什麼相似呢.
語言上有相似的地方.
譬如我們之前幾堂講過.
講到神的形象.
樣式.
這些字.
雖然是希伯來語.
希伯來文.
但是和亞述的文字.
是有相似的地方.
Shamu.
上次講過了.
或者有些術語.
表達的方式.
都是一樣的.
譬如一些交叉的結構.

$^{241}$一些修辭.
寫作方法等等.
甚至對這個宇宙的觀念.
怎麼看這個宇宙.
都是非常相似的.
他們會用一些用語.
四個角.
他們看這個世界是一個扁平的世界.
所以是和他們有相似的地方.
但是一定是有差異的地方.
這就是為什麼聖經是上帝的啟示.
有什麼差異的地方呢.
看神.
宗教.
看人的地位.
看世界觀.
了解人的位置.
人的身份等等.
之前那幾堂我們都曾經講過.
這些就是差異的地方.
所以在這裡讓我們看到.
聖經和當時的世界有差異的地方.
也有一些是很相似的.
相似的地方.
是我們預期的.
這個我們絕對不應該驚訝.
是不是.
一個以色列人.
古以色列人.
和一個巴比倫人.
和他們的鄰居巴比倫人.
埃及人.
有很多東西是相似的.
相似的地方多於我們和一個古以色列人.
現代人和一個古以色列人.
相似的地方.
他們更加多.
所以他們走在一起就能夠溝通.
因為有很多觀念上語言上的東西是一樣.
另外當然有差異的地方.

$^{281}$但是上帝去到啟示.
啟示是什麼呢.
啟示其實就是溝通.
溝通是什麼呢.
溝通就是兩者之間有共同語言才能溝通.
正如我今天和你溝通.
你聽得懂我的語言.
我講得懂你的語言.
我們用的例子.
方法.
語言.
概念.
都是差不多的.
我們才能溝通.
所以神向以色列人溝通的時候.
藉著聖經的作者來溝通的時候.
他都是用當時的人.
他有的觀念來溝通.
所以聖經作為上帝對人的啟示.
對人溝通的時候.
往往不會跟他說一個完全.
跟當時世界完全不同的世界觀套上去.
反而是用當時人熟悉的觀念來溝通.
所以我們明白什麼叫比對研究.
比對研究就是將我們現在找到.
近東世界裡面一些文字.
一些作品.
一些文學作品.
都是一些視窗.
讓我們知道當時的人如何了解自己的世界.
而且那種了解的方法.
很多時候和古以色列人一樣.
所以就讓我們知道.
聖經是如何用當時的世界觀念.
來溝通.
啟示.
神學的訊息.
介紹上帝是怎麼樣的.
當我們明白這個原則的時候.
我們就會知道為什麼要做比對研究.

$^{321}$和比對研究為什麼有必要性.
正如剛才所說.
我們不去比對的話.
不去看當時近東世界的背景的時候.
我們就會用我們現代人的科學頭腦.
我們的概念套上古代的聖經.
舉個例子.
我們說近東的世界.
他們是怎麼看宇宙的呢.
所謂的Cosmology.
Cosmos就是宇宙的意思.
Cosmology就是宇宙觀.
宇宙觀就是說一個人.
如何了解這個宇宙的構造.
這個宇宙的結構.
這個宇宙的Make up.
就是如何建構出來的等等.
在埃及.
你看到他們在一些圖畫上.
形容天就像女神一樣.
蓋著地.
地就是這個躺在地上.
然後有一個空氣的神叫做Shu.
Shu就是將水分開.
然後有一個寒地在下面.
上面就是窮倉.
半圓形的地方.
當然他們了解這個世界的時候.
每一個部份.
就是有神聖.
天有天的神聖.
地是地神.
空氣是空氣的神.
所以他們就會這樣去構思他們的宇宙觀.
所以他們就會認為平地.
這個平地就有地代表的神靈.
天就是那個推著上面的大水.
外面全部都是水.
這個地方就是我們現在生活的地方.
古晉東人就會看.

$^{361}$我們這個世界就像一個氣泡.
一個bubble.
我們為什麼會在這裡呢.
就是因為這些神去頂著這些大水.
讓混沌不會進來.
我們就有寒地.
我們可以有耕種.
我們有四季等等.
我們有我們的生活.
所以為什麼要去拜那些神.
為什麼要去拜天神太陽神等等.
因為這些是靠他們去維繫著整個宇宙的規律.
使它井井有條.
所以他們了解的宇宙.
他們的大地就是一個平的大地.
左下角這裡.
你會看到一個圓形.
中間有一條混沌.
其實這個是巴比倫的地圖.
古巴比倫的地圖.
大概是600BC.
即是主前600年的時候.
他們怎樣了解他們的世界呢.
就是一個平的碟.
然後外面有三角.
就是不同的方向.
其實這裡有一些字標籤了大水的地方.
所以圓碟以外全部都是大水.
是一個平面的世界觀.
平面的去理解宇宙.
這些在古近東的世界裡面.
就算問剛才那兩位的學者.
整個古近東的歷史裡面.
所有的民族都會這樣去了解他們的大地.
就好像這個插圖.
大概是這個意思.
就是一個平的地.
然後四周圍就是一些窮倉的山柱.
Cosmic Mountain.
窮倉的山柱.

$^{401}$就是頂著上面的一個蓋子.
那個叫做Rakia.
在希伯來文叫做Rakia.
就是說窮倉.
在第二天.
上帝創造世界的時候.
就是窮倉.
其實窮倉在希伯來文的原文.
是一隻硬的蓋子.
是硬物.
硬物用來做什麼呢?.
就是頂著上面那些窮倉以上的水.
就是那些Cosmic Water.
就是四周圍都是水.
這些是一些很現象性的語言.
不是一些科學的語言.
現象性的語言就是說.
我們看到天是藍色的.
古代的人想起.
看到藍色的天.
就代表上面一定是水.
但是水不會湧下來.
混沌不會湧下來.
為什麼呢?.
因為上帝分開了.
用窮倉來頂著上面的水.
然後下面的水聚在一起.
就使寒地露出來.
是不是?.
然後有太陽.
星獸和月亮等等.
窮倉那裡有窗戶.
你記得勝經裡面有很多這樣的語言.
窗戶打開.
上帝就傾福祝福等等.
降雨等等.
這些窗戶就在這個鍋子裡.
有些窗戶打開.
他們是這樣形容的.
這些不是科學的語言.

$^{441}$之前我們也說過.
不一定要科學的語言才能表達真理.
這些是現象性的語言.
等於我們今天說.
太陽從東邊出來西邊下山.
他們都是這樣去形容的.
在大地之下就是陰間.
陰間之下當然就是水.
氣泡以外的東西.
全部都是水.
這個其實有點像.
很多古代的民族都是這樣.
中國古代的民族.
其實都是這樣了解.
他的世界.
不是一些科學性的語言.
是什麼?.
是一些現象性的語言.
這樣看的時候.
我們知道神去啟示他的說話的時候.
其實.
他就是用當時人熟悉的世界觀.
來表達真理.
所以聖經絕對不是說.
將這個世界觀.
古近東古人的世界觀.
就說要照抄.
照版轉換.
要我們現在做.
不會的.
如果上帝現在啟示他自己的話.
他就不會用這個世界觀.
用我們現在的現代人.
我們所熟悉的觀念.
來啟示.
但上帝啟示給.
古以色人的時候.
就是用他們當時知道的世界觀.
所熟悉的觀念.
這個就給了研究.

$^{481}$給我們一些資料.
能夠建構出來.
所以我們就不會這麼淪盡.
將我們.
現代的觀念.
套上聖經.
譬如.
這裡舉個例子.
以賽亞書40章22節.
神造大地的.
在在玄穹之上.
地上的居民.
好像乍猛等等.
在大地的玄穹之上.
我們很多時候.
就以為.
就是什麼.
就是我們想起太空.
是不是.
想起那個.
圓球.
這個地球是圓球.
這是我們的世界觀.
我們的宇宙觀.
我們就想起.
這個一定是聖經裡面說得很清楚.
就是說.
這個世界.
證明了科學.
聖經早就說了一些科學的真理.
但其實你看一下圓雲.
這個圓拱.
其實是圓形的意思.
圓碟的意思.
不是一個圓球的意思.
Yosef.
他坐在.
Yosef.
他坐著的.
He who sits.

$^{521}$就是那位坐著的.
阿訶.
就是在.
一個圓圈.
圓碟.
什麼圓碟.
大地的圓碟.
阿訶.
大地的圓碟之上.
他坐著.
他就在說那個.
鑊.
倒轉的鑊.
那個窮倉異常.
就好像上面右上角這一幅圖畫.
這個是John Walton.
其中一本解經書.
裡面他畫出來的.
找一個藝術家畫出來.
你看到大地嗎.
一個旱地.
一個平面的旱地.
有些宇宙山.
頂著上面一個.
硬的蓋.
頂著上面的.
大水.
下面就是陰間.
和地下.
下面的水.
好像蛇或者龍.
海怪.
以往的人想到下面.
很多東西很深不可知.
裡面有些大魚.
他們會敬拜這些東西.
創世的第一章也有說.
神造大魚的那個字.
其實是近東人.
他們會想起大魚是神聖的東西.

$^{561}$但是創世的第一章就說.
不是啊 只不過是上帝造的.
百物裡面其中一種.
萬物裡面其中一種.
沒有神聖的 不需要怕他.
不需要敬拜他.
但是大地的上面就是窮倉.
有些是窗.
有些是窗口.
倒些水下來.
上面是一個聖殿.
聖殿在地上的聖殿.
其實就是抄天上.
上帝的寶座.
所以神在大水之上.
就是這樣的.
意念.
這樣的畫面.
所以.
不是說我們.
現代人的.
宇宙觀.
所以不需要用一些.
科學的意念.
來套上去.
用古近東的.
宇宙觀就能夠了解.
到底真理想表達什麼.
就是上帝.
統管大地.
統管萬有.
這一個.
就是上帝所啟示的真理.
就是每一代的信徒.
我們都能夠形容得到.
我們都能夠明白.
神學的意義.
雖然今天的宇宙觀是不同的.
但是以賽亞雅.
四十章二十二節.

$^{601}$都是在表達同樣的真理.
就是上帝仍然是掌管統管一切.
是不是.
我們只不過是很渺小的人.
所以用不同的宇宙觀.
都是表達同一個的真理.
但是以賽亞書.
很明顯就是用一種.
古希伯人的.
近東的宇宙觀.
那麼你說.
創世記第一章.
不是在說從無到有的.
創世的教義.
是不是.
因為一開始的時候.
起初神創造天地.
地上空虛雲遁,淤面黑暗有水.
那水從哪裡來.
是不是.
我們就想是不是第一節和第二節中間.
有些東西.
我們智慧行間可以讀進去.
其實是不需要的.
一開始的時候就有水.
創世記第一章是這樣說.
那你說不是從無到有.
因為我們基督教的.
系統神學裡面.
我們有一個Creation out of nothing.
就是從無到有的教義.
是不是.
那教義去了哪裡.
其實很簡單的.
創世記第一章不是想表達這個教義.
所以你不能夠.
用創世記第一章.
去引申出從無到有.
創造的教義.
Creation is the hero.

$^{641}$就是out of nothing.
因為創世記第一章不是在回答這個問題.
創世記第一章要表達什麼.
是一個kingdom.
一個國度,上帝安排一切.
這個叫做創造.
真正.
所謂從無到有的.
創造的.
神學意念到底在哪裡.
才發展出來呢.
就是兩約中間才有.
瑪加比書.
七章二十八節.
兩經之間,兩約中間.
幾百年沉默時期.
一些猶太人的著作.
那時候他們才去到.
理解.
或者去建構.
這個從無到有的創造的教義.
這個就是瑪加比書.
二章.
瑪加比二書七章.
就是父親和子女說.
因為那時候他們受迫害.
他說.
孩子們我求你要仰望天地其中的一切.
要知道.
這是神所創造的.
不是來自已經存在的東西.
人類又何嘗不是這樣呢.
其實他在說.
上帝從無.
創造萬有.
既然如此.
我們相信上帝.
我們何必怕這些迫害的勢力呢.
因為上帝.
是比這些萬有的大.

$^{681}$所以我們可以信靠他.
這裡.
一字聖經就是表達.
一個最清楚.
在歷史的文獻裡面.
猶太人的文獻裡面.
最清楚有無到有的.
創世的觀念.
之前.
是不會這樣想的.
因為我們現代人.
一想到創造的時候.
我們就想到什麼.
就是想到一個物質的創造的過程.
一個物質的創造的過程.
的過程.
由無變成有.
原子物質分子應該怎樣組成.
等等.
這是我們現代人有的觀念.
但是古代人想的創造.
就是一種功能性的.
一種功能性的.
一種安排.
一種功能性的.
一種的管理.
統治.
其實我們今天.
用創造這個字.
或者英文用create這個字.
都有這個意思.
你老闆跟你說.
你去創造一個部門出來.
你去管理那個工作.
或者做那個project.
那個老闆叫你創造一個team出來.
一個隊工.
一個部門.
去管理某一個工作.
那個不是物質創造的過程.

$^{721}$那個是什麼.
是管理的過程.
所以創世記第一章.
都是在說management.
是統治管理的層面.
是functional.
當然你問一個古希伯來人.
問問先知.
如果我們坐時光機回去.
你問他那些物質從哪裡來.
他們肯定會回答你.
當然是上帝創造的.
他們一定會這樣說.
但是呢.
但是你要知道創世記第一章第二章.
想說的就不是那一封.
明白嗎.
OK.
所以你搞清楚這個理念的時候.
有什麼分別呢.
其實為什麼呢.
就是因為有比對研究出現.
比對研究看看巴比倫人.
怎麼看創世.
埃及人怎麼看世界的開始.
全部都是一種functional的方法去看.
功能.
即是創造主.
去將這些東西.
安排得井井有條.
所以就要敬拜他.
所以創世記第一章.
不是物質的.
創造的過程.
是誰統領萬有.
所以我們應該要敬拜他.
如果你喜歡我們的影片.
就記得要讚好,留言和分享.
Like, Comment 和 Share 啦!.
\newpage



\section{出埃及記 12:1-14}
\label{sec:xeNKKwLqznQ}
\textbf{逾越年月 (出埃及記12\_1-14) -  朱莉傳道}
\newline
\newline
連結: \href{https://youtube.com/watch?v=xeNKKwLqznQ}{\texttt{https://youtube.com/watch?v=xeNKKwLqznQ}} ~~~~ 語音日期: 2022-02-07
\newline
\newline
\hyperref[sec:q8ZgN14Imcg]{\small{< < < PREV SERMON < < <}}
~
\hyperref[sec:index_chronic]{\small{[返順時目]}}
~
\hyperref[sec:index_scriptual]{\small{[返順卷目]}}
~
\hyperref[sec:INsxhCoyb9U]{\small{> > > NEXT SERMON > > >}}
\newline
\newline
出埃及記 12:1-14
\newline
\begin{longtable}{cl}
\hline
\hline
章節 & 經文 (和合本修訂版)\\
\hline
12:1 & \begin{tabularx}{0.7\textwidth}{X} 耶和華在埃及地對摩西和亞倫說： \end{tabularx} \\ \\ \relax
12:2 & \begin{tabularx}{0.7\textwidth}{X} 「你們要以本月為正月，為一年之首。 \end{tabularx} \\ \\ \relax
12:3 & \begin{tabularx}{0.7\textwidth}{X} 你們要吩咐以色列全會眾說：本月初十，各人要按著家庭取羔羊，一家一隻羔羊。 \end{tabularx} \\ \\ \relax
12:4 & \begin{tabularx}{0.7\textwidth}{X} 若一家的人太少，吃不了一隻羔羊，就要按照人數和隔壁的鄰舍共取一隻；你們要按每人的食量來估算羔羊。 \end{tabularx} \\ \\ \relax
12:5 & \begin{tabularx}{0.7\textwidth}{X} 你們要從綿羊或山羊中取一隻無殘疾、一歲的公羔羊， \end{tabularx} \\ \\ \relax
12:6 & \begin{tabularx}{0.7\textwidth}{X} 要把牠留到本月十四日；那日黃昏的時候，以色列全會眾要把羔羊宰了。 \end{tabularx} \\ \\ \relax
12:7 & \begin{tabularx}{0.7\textwidth}{X} 他們要取一些血，塗在他們吃羔羊的房屋兩邊的門框上和門楣上。 \end{tabularx} \\ \\ \relax
12:8 & \begin{tabularx}{0.7\textwidth}{X} 當晚要吃羔羊的肉；要用火烤了，與無酵餅和苦菜一起吃。 \end{tabularx} \\ \\ \relax
12:9 & \begin{tabularx}{0.7\textwidth}{X} 不可吃生的，或用水煮的，要把羔羊連頭帶腿和內臟用火烤了吃。 \end{tabularx} \\ \\ \relax
12:10 & \begin{tabularx}{0.7\textwidth}{X} 一點也不可留到早晨；若有留到早晨的，要用火燒了。 \end{tabularx} \\ \\ \relax
12:11 & \begin{tabularx}{0.7\textwidth}{X} 你們要這樣吃羔羊：腰間束帶，腳上穿鞋，手中拿杖，快快地吃。這是耶和華的逾越。 \end{tabularx} \\ \\ \relax
12:12 & \begin{tabularx}{0.7\textwidth}{X} 因為那夜我要走遍埃及地，把埃及地一切頭生的，無論是人是牲畜，都擊殺了；我要對埃及所有的神明施行審判。我是耶和華。 \end{tabularx} \\ \\ \relax
12:13 & \begin{tabularx}{0.7\textwidth}{X} 這血要在你們所住的房屋上作記號；我一見這血，就逾越你們。我擊打埃及地的時候，災殃必不臨到你們身上施行毀滅。」 \end{tabularx} \\ \\ \relax
12:14 & \begin{tabularx}{0.7\textwidth}{X} 「你們要記念這日，世世代代守這日為耶和華的節日，作為你們永遠的定例。 \end{tabularx} \\ \\
[1ex]
\hline
\hline
\end{longtable}
$^{1}$等姐妹平安.
臨近新年.
我首先在主裡面祝福大家.
新年蒙新恩.
剛好今天想和大家看的這段經文.
也是和新年有關的.
在《出埃及記》12章1至14節.
我們一起讀這段經文.
《出埃及記》12章1節.
「耶和華在埃及地曉遇摩西亞倫說.
你們要以本月為正月.
為一年之首.
你們吩咐以色列傳會眾說.
本月初十日.
各人要按著夫家取羊羔.
一家一隻.
若是一家的人太少.
吃不了一隻羊羔.
本人就要和他甲碧的鄰舍共取一隻.
你們預備羊羔.
要按著人數和飯量計算.
要沒有殘疾.
一歲的公羊羔.
你們或從綿羊裡取.
或從山羊裡取都可以.
要留到本月十四日.
在黃昏的時候.
以色列傳會眾把羊羔宰了.
各家要取點血.
塗在吃羊羔的房屋左右的門框上.
和門尾上.
當夜要吃羊羔的肉.
用火烤了.
與無烤餅和苦菜同吃.
不可吃生的.
斷不可吃水煮的.
要戴著頭腿五臟.
用火烤了吃.
不可剩下一點留到早晨.
若留到早晨要用火燒了」.

$^{41}$十一節:「你們吃羊羔.
當腰間束帶.
腳上穿鞋.
手中拿杖.
趕緊地吃.
這是耶和華的逾越節.
因為那夜我要巡行埃及地.
把埃及地一切投生的.
無論是人是牲畜都擊殺了.
又要敗壞埃及一切的神.
我是耶和華.
這血要在你們所住的房屋上作記號.
我一見這血就越過你們去.
我擊殺埃及地投生的時候.
災殃必不臨到你們身上滅你們」.
我們讀完第十四節.
「你們要紀念這一天.
守為耶和華的節.
作為你們世世代代永遠的鼎勵」.
沒錯.
今天想跟大家一起來討論的.
就是逾越節.
我相信我們對逾越節都不是太陌生的.
我們都知道逾越節其實對猶太人來說.
是非常非常重要的一個節日.
我們未必會記得.
原來逾越節就是猶太人的新年.
其實你這樣讀下去.
跟我們中國人過農曆新年的習俗.
是有些類似的.
很注重親吻等等.
要全家人在一起等等等等的.
但是逾越節對於猶太人來說.
有什麼重要性呢?.
第一.
第二節.
「耶和華在埃及地曉喻摩西亞倫說.
你們要以本月為正月.
為一年之守」.
就是說.

$^{81}$這是神的吩咐.
是神曉喻摩西亞倫.
要以這個逾越節為一年的開始.
猶太人的月曆.
他們的年曆.
他們的曆法.
就要從此不再一樣.
或者這麼說.
可能猶太人以前.
因為他們在埃及地遺老四百年.
可能跟隨了埃及人的曆法很久.
但是經過《猜及記》十二章之後.
猶太人不再一樣.
他們要有新的時間觀.
是用一個逾越節來界定的時間觀.
當然也要說一點.
如節是什麼呢?.
「如節」英文是past over.
可能會更加清晰.
就是「越過」的意思.
剛才我們也讀到.
在十二節提到.
那一夜.
耶和華要巡行埃及地.
將埃及地一切投生的都擊殺.
但是因為有血.
在猶太人.
在以色列人的門框上.
以及門尾上.
在十三節就說.
「這血要在你們所住的房屋上作記號.
我一見這血就越過你們去.
災殃就不會臨到」.
以色列人.
所以這是紀念猶太人.
用血作為記號.
以致避過了滅長子災殃的一晚.
當然如節的背景.
《猜及記》十二章上文說什麼呢?.
就是十災.

$^{121}$可能大家都會記得.
十災裡面提到.
神如何藉著摩西和亞倫的手.
施行一切的一連串的神蹟騎士.
以致拯救以色列人.
來脫離法老.
即是老邑他們的額.
可以離開埃及遺牢之間.
然後去到《猜及記》十一章.
就已經預告了.
即將會有大戰.
就是最後的災.
前面是九災.
去到十一章就預言.
十二章將會出現第十災.
亦都是高潮.
因為是滅長子的災.
然後最終會迫使法老屈服.
放以色列人離開.
所以那個掙扎.
一直在增加.
去到十二章.
問題來了.
如果好像拍戲一樣.
去到十一章,十二章的時候.
就已經很緊張.
很等待看什麼?.
大戰.
那個法老和上帝.
或者是上帝藉著摩西亞倫的手.
那個終極的大戰.
然後最後第十災會怎樣展開.
奇怪的是什麼呢?.
我們就這樣讀第十二章.
只是讀了頭部份.
但如果你有時間再看下去.
你會發覺.
《猜及記》十二章是出乎意料地.
不合乎比例地記錄什麼.
如二節怎樣過.

$^{161}$怎樣吃東西.
怎樣預備.
各樣的事情.
我們在等的終極大戰.
卻只是短短幾節.
在十二章的二十九至三十六節.
如果你有機會回去仔細看.
只是短短幾節記載.
那個第十災的經過是怎樣發生.
就好像什麼呢?.
就好像你去看電影.
追一套劇.
或者看一套電影.
很緊張.
一直在等最後的男主角和奸角.
怎樣有最後一戰.
然後怎樣男主角海船而歸.
將奸角打到無法還手.
但是導演.
他將鏡頭轉移到哪裡呢?.
女主角在家不斷煮飯.
切菜.
當然是新年.
煎蘿蔔糕.
煎年糕.
預備連夜飯.
吃什麼等等.
如果你是觀眾.
你看到這裡的話.
你都會很生氣.
你都會說「導演搞什麼呀?」.
我們想看的是那個大戰.
但不是.
明顯地作者的重點是什麼呢?.
是在寓月節本身.
甚至不是第一次寓月節怎樣過.
當晚發生什麼事.
再回去看.
非常詳盡去說寓月節怎樣預備.
各樣的飯菜.

$^{201}$有什麼規定.
那個經過不是最重要.
最重要是什麼呢?.
寓月節不是只有一晚的事.
是以後要怎樣過.
這個要成為什麼呢?.
定例十二章十四節.
你們要紀念這一天.
首為耶和華的節.
作為你們世世代代.
永遠的定例.
如果我讀新漢語譯本.
另一個譯本.
會怎樣呢?.
大家聽一下重複.
有什麼是重複的字眼.
這一天要成為你們的紀念.
這個也是《差甲記》十二章十四節.
但是新漢語的翻譯.
你們要首為節期.
是屬於上主的節期.
且要給你們作世世代代.
永遠的律例.
是你們當首的節期.
然後你繼續看十二章下去.
你會發現重複又再重複.
這個要成為定例.
要永永遠遠的.
甚至已經離開埃及了.
去到應許地.
已經在流乃與墨之地.
仍然要繼續守節.
可見逾節本身.
或怎樣過逾節.
對神的選民來說.
是非常重要的.
不只是那一晚發生了什麼事.
是要重複重複.
往後的日子中.
還要重過重演當晚.

$^{241}$當然是用一些象徵的意義.
象徵的一些象徵.
去重演當晚發生什麼事.
十災其實已經成為歷史.
甚至以色列已經成功出埃及.
但為什麼還要不斷重演.
逾節那一晚的事.
甚至猶太人要用逾節.
成為他們的月曆.
他們的年曆的開始呢?.
因為節期和身份認同.
是有非常密切的關係.
因為節期和身份認同.
是有非常密切的關係.
為什麼以色列人可以避過.
十長節的災呢?.
是否真的那些血.
有那種避禍的能力.
好像我們中國人的傳統.
以前古代相信.
年輝春將福倒轉.
等等.
就是可以招財進寶.
就是可以有福氣臨到等等.
是否真的因為那些血.
在門裡面.
就能夠有那個避禍?.
為什麼我們現在.
基督徒或者現在的猶太人.
不這樣做呢?.
因為經文這樣說.
剛才我們讀的經文這樣說.
血只是一個記號.
那個記號是表明什麼呢?.
表明這一班是以色列人.
是神的選民.
是和他周遭的鄰居埃及人不同的.
所以這是一個特殊的身份.
而是因為這個特殊的身份.
讓他們避開了.

$^{281}$這個殺長子的災禍.
為什麼以色列人.
又要以這個禹節.
成為他們新的歷法的開始呢?.
又是表明他們和埃及人不一樣.
甚至他們去到迦南地.
去到英帝已經安定下來.
去享受當地的各樣出產.
他們仍然要以這個禹節.
成為那個歷法的開始.
和他周遭的迦南人不一樣.
因為這是神選民特別的時間觀.
其實我不需要再講太多的例子.
我們本身都很容易明白這件事.
例如我們在美國.
我們接觸很多不同的華人.
他們可能從不同地方來的.
可能從香港來的.
從台灣來的.
當然也有從大陸來的.
馬來西亞,新加坡.
我們都過農曆新年.
但是我們過農曆新年的方法.
是有微妙的分別的.
我經常都說笑.
我們吃年糕.
我們廣東人或者香港人.
是習慣了煎的.
但是原來台灣人有不同的.
馬來西亞,新加坡人也有不同的.
甚至其實日本,韓國.
受中華文化影響的那些.
他們都有他們那個.
跟我們又類近又不像的.
一些過新年的習俗.
憑著這些很微妙的分別.
其實我們就已經可以知道.
他們是從哪裡來的.
所以節日和身份認同.
是有非常密切的關係.

$^{321}$神的拯救.
不只是那一刻發生的事.
不只是那一晚.
到底神怎樣偉大.
怎樣厲害地施行了第十災.
以至很奇妙地將以色列人.
從埃及帶領出來.
更加重要的是什麼呢?.
是這件事構成了他們的身份認同.
而這個身份認同.
是繼續影響著以色列人的人生.
當他們過雨節的時候.
他們回想.
他們回想當時發生什麼事的時候.
他們再一次被提醒.
雨節不只是一個古代.
一個歷史上發生的事.
是仍然影響著他們怎樣認識神.
怎樣去認識自己.
We are who we are.
如果一個猶太人這樣去想的時候.
其實就是從雨節開始.
來界定他們的身份.
有一個這樣的名言.
是捷克文學大師米蘭昆德拉.
這樣說過一句話.
人與強權的鬥爭.
其實是一場記憶和遺忘的鬥爭.
其實對基督徒來說.
對我們的信仰歷程來說.
都是有異曲同工之妙的.
對我們來說.
都是一種記憶和遺忘的抗爭.
神知道我們很善忘.
所以就透過這些節日.
透過每個星期按時的崇拜.
透過每一天我們讀的聖經.
就是這樣每時每刻提醒我們.
其實我們是誰.
我們的身份是甚麼.

$^{361}$然後提醒我們.
神是capable去做甚麼.
我們日常生活可能很平淡.
但其實我們每一天都需要被提醒.
其實我們的身份是甚麼.
最後我再有兩點提醒.
從這段經文.
我不打算很詳細地講解.
到底如節的陽光等等有甚麼意義.
但在這段聖經裡面.
有少少提醒.
讓我們在新的一年.
如何可以繼續經歷神新的恩典呢?.
在《初一及記》十二章一到十四節.
當然是關於如節如何過.
不單是第一次的如節如何過.
是以後要成為定例.
而我們讀下去的時候.
就會發現它充滿了很多命令.
是指示你們要怎樣做.
我們簡單讀幾句.
十二章的第一節.
「在埃及地曉喻摩西亞倫」.
所以神吩咐摩西亞倫.
第二節.
「你們要以本月為正月為一年之壽」.
「你們吩咐以色列的傳會眾」.
由摩西亞倫到以色列傳會眾.
「本月初十日各人要按著附加取羊羔一家一隻」.
「若是一家人太少吃不了一隻羊羔」.
「本人就要和他隔壁的鄰舍共取一隻」.
「預備羊羔」.
「要按著人數和飯量計算」.
「要沒有殘疾」.
「一歲的供羊羔」.
「你們或從綿羊裡取或從山羊裡取都可以」.
「要留到本月十四日」.
「在黃昏的時候以色列傳會眾把羊羔載了」.
我繼續讀下去.
但我相信大家感覺到.

$^{401}$你可以感受到經文的指示意味很重.
後面還有不可等等.
所以怎樣過如意節.
不是要斬羊羔等等.
那些很嘈吵的事.
最重要是甚麼?.
就是要照主的吩咐去做.
怎樣過如意節呢?.
其實是一種測試.
of obedience.
我們能否照著神的吩咐去做呢?.
它真的很嘈吵.
真的很細微.
甚麼時候吃.
吃多少隻.
吃多大隻.
要怎樣分.
要怎樣煮都說了.
是細微到.
我們有時會問.
「神呀!你真是太高要求」.
另外呢.
牽涉很多人.
剛才說了.
從神吩咐摩西亞倫.
還有以色列的傳會眾.
其實中間還有.
後面十二章提到的長老.
這些官長.
他們都要負責傳遞神的話語.
一層一層地說.
以至全會眾都一起.
accord(同意).
去進行神的主義.
最要命.
聽神的話.
最難的是甚麼呢?.
最難的一點是甚麼呢?.
是和我們的經驗往往有偏差.
是和我們的認知有出入.

$^{441}$這是第十災.
而前面的九災.
其實以色列人甚麼都不用做.
就自動免疫.
你要看前面的九災.
以色列人的參與度是非常低.
可能在旁邊.
用俗語說.
他們在旁邊吃花生.
但今次第十災.
不單止他們要去參與.
他們還要做甚麼呢?.
他們還要獻羽節羊羔.
某程度上羊羔是一個祭.
要獻.
要付代價.
要跟著.
會有很多問題要問.
既然要穿上準備遠行的衣服.
但又不讓你帶剩下的食物.
我們去遠行吃甚麼呢?.
可以有很多問題要問.
但十二章是一段很長的經文.
上半部是指指示.
是一個命令.
下半部是指指示.
是一個執行.
為何要這麼長進呢?.
不知道大家看聖經會否覺得很沉氣.
指指示已經用了很多篇幅.
很耗時的去介紹.
為何去到後半部.
還有一個指指示的版本.
差不多重覆前面指示所說的話.
聖經要告訴你.
那些人緊跟著神的命令去進行.
所以你看十二章的下半部.
有幾句說話是重覆的.
就是耶和華怎樣吩咐摩西亞倫.
以色列就怎樣行.

$^{481}$這是不可或缺的記錄.
指示,執行是一樣的.
是一致的.
我最近發現了一件很有趣的事.
有一次逛街.
我收到有人發了一封訊息給我.
是一幅很有趣的圖.
我想回覆他.
我想說:笑死我了.
我用手機打了一封訊息.
想說:笑死我了.
我打也打不到.
為甚麼呢?.
是手寫的.
我想寫一個「笑」字.
但他不斷出「哭」字.
後來才發現原來這兩個字很相似.
我以前都沒有留意.
我想寫「笑」.
但他就出「哭」.
後來我放棄了.
我就索性錄一個音響訊息.
因為我正在逛街.
不想不知道怎樣打.
有時候真的哭笑不得.
因為我們想這樣.
神的吩咐是這樣.
但我們可能做完全相反的事.
有時候電話,Siri都不太可靠.
更何況我們人.
聽神的話不容易.
但這表明我們屬於神.
而且這是經歷神恩典的路徑.
聽神的話是經歷神恩典的不易路徑.
在新的一年.
但願我們去學習.
緊緊跟隨神的話.
認真去學習.
去明白神的話.
到底對我們有甚麼要求.

$^{521}$第二.
我們怎樣可以在新的一年.
去經歷神的恩典呢?.
12章42節到後半部.
真正十災發生的時候.
12章42節這樣說.
「這夜是耶和華的夜.
因耶和華領他們出了埃及地.
所以當向耶和華緊守.
是以色列眾人世世代代該緊守的」.
如果我們這樣看和合伴.
我們可能有時不太明白.
「這夜是耶和華的夜」.
然後又要向耶和華緊守.
是以色列世世代代該緊守的.
我們未必太清楚.
所以有時看英文.
或者看其他翻譯.
我們會更加明白.
如果用ESV的版本.
這句怎樣翻譯呢?.
「這夜是耶和華的夜」.
原來是翻譯成.
「A night of watching」.
A night of watching.
然後用NIV.
我們可能比較熟悉的版本.
NIV是甚麼呢?.
A night to keep vigil.
即是警醒.
所以原來這句說話是甚麼意思呢?.
它說我們要去警醒.
而且有趣的地方.
不只是這句.
「A night of watching」.
或者「keep vigil」.
即是警醒.
即是「wake up」.
即是「stay awake」.
不只是這句出現一次.

$^{561}$它說這是神.
或「keep vigil」.
但原來世世代代我們都要怎樣呢?.
我們都要謹守.
都要警醒.
最visible.
最visual的例子是甚麼呢?.
就是香港.
如果你有印象.
香港多數都住大廈.
大廈通常會有保安.
即是伯伯或叔叔.
他就是要「keep watching」.
而且要「keep vigil」.
要警醒.
的一個這樣的模式.
某程度上原來那一晚發生的事.
是以色列人有很多事要準備.
要跟著神的指示.
一樣一樣事去做.
但真正的行動.
這是神的事.
是祂「keep vigil」.
是祂要我們去「watch」.
的一個這樣的晚上.
但很有趣的這一句.
其實是雙向的.
一方面是我們要「watch」.
但另一方面我們「watch」甚麼呢?.
我們「watch」原來神在那一晚.
實現了祂應許亞伯拉罕的那件事.
在創世記的時候.
即是距離真正以色列人出埃及.
已經幾百年前的事.
當時要說.
應許了亞伯拉罕會帶領他的子民.
他的子孫.
亞伯拉罕的子孫.
從奴役他們的埃及離開.
可以去承受這個應許地.

$^{601}$所以儘管是經歷了幾百年.
但神仍然實現了祂的應許.
拯救了亞伯拉罕的子孫.
從埃及地出來.
所以這是甚麼呢?.
是以色列人都要「watch」.
神怎樣承受這個應許.
「stay awake」「keep vigil」.
然後另一方面.
然後浩浩蕩蕩離開埃及.
在最近的這段日子裡.
都不是很容易.
特別是股市的起伏.
是令人心驚膽跳的.
甚至有戰爭的風聲.
我相信大家如果有留意新聞.
都會知道的.
不單好像.
有時我們覺得好像很遙遠.
但的確這些戰爭風聲.
都會影響我們的經濟.
影響我們的生活.
甚至是疫症.
到底有些預測.
Omicron這個變種已經是.
pandemic的終極.
到底是否真的盡頭呢?.
對於身處亞洲的.
例如我.
都會有很多的疑問.
到底2022年將會是甚麼年呢?.
到底這些事情發生.
這些風聲有沒有奪去我們的平安呢?.
還是我們可以在一個變幻的世代.
一個很多捉不住.
很多變化的環境裡.
仍然不斷去經歷神生的恩典呢?.
我最近有一個很特別的經歷.
就是.
因為在聖誕之前.

$^{641}$我打算去剪頭髮.
我一向都不講究.
在髮型屋找甚麼理髮師.
專門去剪頭髮.
所以通常就剛好時間.
去邊頭隨便進一間髮型店.
去剪髮.
如常地.
我去到一個髮型店.
我未去過的.
我坐下來就留意到一些很特別的東西.
我發現那張桌子.
那張鏡子和桌面.
就貼滿了或者放滿了.
最近香港某個非常紅的偶像團體.
一個男團.
不知道大家有沒有聽過.
叫MIRROR.
如果你未聽過.
我鼓勵你去Google一下.
因為在香港爆紅.
在去年.
突然間有一個這樣的男團.
當然中間有一些特別紅的成員.
在香港爆紅.
是一個很有趣的社會現象.
我知道有很多社會學家.
都研究這個現象.
反映香港社會發生甚麼事.
無論如何.
我發現放滿了這些相片.
某些成員的相片.
應援物.
香港跟著韓國的追星文化.
影響了應援物.
做了很多海報,卡等等.
叫應援物.
放滿了在那張髮型師的桌子上.
因為我也有留意MIRROR等等.
所以我一看就看到.

$^{681}$我就不發聲.
繼續剪輯.
因為我怕我亂說.
我怕我說錯了話.
不知道他會有甚麼反應.
他在剪我的頭髮.
所以我就一直觀察.
到差不多剪完了.
我忍不住就問他.
「你是否鏡粉」.
「是否MIRROR的粉絲」.
幫我剪頭髮的那個.
大概三,四十歲的男性理髮師.
他一聽到我問這句.
他突然雙眼發光.
而且那個緊張.
那個熱情度是由零.
因為我一直剪頭髮.
那十多二十分鐘.
我們沒有交流.
除了問你要怎樣剪.
就一直剪,默默地剪.
所以是很平淡的過程.
但我一問這句.
他的緊張是從零.
在短短幾秒內.
可能一,兩秒內.
就爆到一百.
很開心,很激動.
雙眼發光地看著我.
「是呀!你又是呀!」.
就聊天.
甚至怎樣呢?.
在他的抽屜裡.
拿了一疊很厚很厚的應援物.
即是那些為某些成員做的海報.
要送給我.
他就告訴我.
他們是怎樣自資.
自己出錢做這些.

$^{721}$自己設計.
甚至自己拿去印.
自己再租車去拾取.
甚至自己再組織.
在某個特定的時間.
在哪個地福音派這些應援物.
為了希望更多人認識他的偶像.
亦都希望更多人和更多人分享.
他追星的快樂.
聊了幾分鐘.
後來我離開了.
離開的時候.
我拿著他送給我的.
兩張成員的應援物.
那些海報都頗大張.
我哭笑不得.
我站在街上.
我不知道拿著這兩張做甚麼好呢?.
當然我有些朋友都是他的粉絲.
所以都可以送給他.
當然我拿著這兩張東西.
站在街上的時候.
我突然間很感嘆.
很感慨.
亦都反省.
有多久我.
我說的是我自己.
很久沒有試過.
為我領受的這個偉大的救恩而興奮.
我有多久沒有試過.
忍不住向人分享神的奇妙呢?.
我有多久沒有試過.
是因著我自己的生活行為.
而別人認出我是基督徒呢?.
是一個很有趣的經歷.
讓我有很多的反省.
沒錯,這個救恩.
甚至這本聖經.
當是一本古老的聖經.
是一個已經傳頌千年的救恩.

$^{761}$但這個救恩對我們現在的生活仍有影響.
當我們年復年.
週復週.
日復日.
去評論.
去看神在過往歷史裡所作的作為.
祂怎樣拯救我們.
怎樣在我們人生過往的一些經歷裡.
施恩給我們.
祂今日仍然可以.
祂不是在歷史.
神不是活在歷史的神.
祂今日裡面仍然與我們同行.
仍然在做新事.
有沒有警醒.
有沒有緊緊跟隨祂的腳步.
有沒有再一次慰悼我們有這位神.
我們有這個基督徒尊貴的身份.
而感到興奮呢?.
但願在新的一年裡.
神再次來激勵我們.
讓我們存著信心.
帶著盼望.
進入2022年.
我們一起祈禱.
你是一個值得我們傳頌.
值得我們高舉的這位上帝.
但願我們不只是停在口頭上.
讓我們在我們的生活.
在我們的行事為人.
在我們每一天要做的決定裡.
在我們讀的聖經裡.
我們仍然被你激動.
仍然為我們所領受的救恩而歡喜快樂.
當我們在一個不知.
在一個掌握不到的世代裡.
讓我們緊緊跟隨你行.
因為我們知道上帝你仍然作事.
仍然作新事.
不單在這個世界裡.

$^{801}$更加在我們的生活裡.
你是那個沒有離開我們.
繼續與我們同行的上帝.
在今日寶貴的時間.
禱告奉耶穌基督的聖名.
阿們.
\newpage



\section{21:1-14}
\label{sec:INsxhCoyb9U}
\textbf{我想和你無所事事 (約21\_1-14) -  陳德晶傳道}
\newline
\newline
連結: \href{https://youtube.com/watch?v=INsxhCoyb9U}{\texttt{https://youtube.com/watch?v=INsxhCoyb9U}} ~~~~ 語音日期: 2022-02-15
\newline
\newline
\hyperref[sec:xeNKKwLqznQ]{\small{< < < PREV SERMON < < <}}
~
\hyperref[sec:index_chronic]{\small{[返順時目]}}
~
\hyperref[sec:index_scriptual]{\small{[返順卷目]}}
~
\hyperref[sec:b3sjDQwWdlg]{\small{> > > NEXT SERMON > > >}}
\newline
\newline
$^{1}$弟兄姊妹平安!我們還在農曆新年裡面.
所以先祝大家出入平安,身體健康.
知道參與崇拜的不單止是芝加哥的弟兄姊妹.
也有來自世界各地的弟兄姊妹.
我相信為何能夠一起崇拜呢?.
是因為這個疫情帶來的正面影響.
說起疫情 其實我們全世界都被它影響.
不知道你在當中有什麼感受呢?.
會不會害怕?會不會恐懼?會不會覺得很煩?.
會不會很憂愁?或者甚至有些憤怒在當中?.
我最近看過一篇文章來自紐約時報.
我覺得很有意思.
因為當中有一個組織心理學家叫做Adam Grant.
他指出一樣東西.
剛才所說的一些感受都是2020年.
全世界集體的一種感受.
但到了2021年就換來另一種感受.
這個感受叫做Language.
中文譯為「韋美」.
看完中文的不明白他在說什麼.
意思是什麼呢?就是覺得很無趣.
很沒人 做什麼都沒什麼人.
萬無目的 無所事事 在當中.
Adam Grant 觀察到整個疫情都一段時間.
有很多人都發現其實都已經在改善疫情的狀況.
但亦都很多人在一個這樣的時間.
他經常都在家裡想著看劇集.
看完《魷魚遊戲》就看《Don't Look Up》.
或者一直在看舊劇集《Friends》.
什麼都混在一起.
有些人又會選擇整天睡覺.
一天十幾個小時在床上.
要去洗手間都不會去.
廁所都不會起床.
作者說這些不是一個燃燒的現象.
因為燃燒我們會很沒精神.
但在當中我們還有精神去做事.
亦都不是一些抑鬱的狀態.
因為我們還沒有去到很絕望的感受.
但總是有一種無人.

$^{41}$萬無目的 無所事事 在當中.
不知道大家有沒有這種感受呢?.
其實我很坦白 我有的.
所以我都會嘗試去問一下.
如何去解決.
我解決問題通常會去看聖經.
找一些出路.
在當中我見到一群無所事事的人.
就是在《約翰福音》二十一章一至十四節.
在第一節裡面你會見到.
「這些時以後 耶穌在提比利亞海又向門徒顯現」.
這裡提及「這些時」.
是指二十章.
耶穌向瑪利亞,彼得,約翰.
又向一個門徒甚至多瑪顯現的事.
當時甚麼門徒.
就好像我們在2020年裡面的那種感受.
感受到恐懼,懷疑,憂愁.
所以你見到他們不斷在背後哭.
也怕得要關上門 不讓人進來.
但來到這裡.
來到提比利亞海邊.
門徒顯然不同了.
敢走出去.
但他們卻在當中不同.
他們這班人有的是西門彼得.
和稱為底土馬的多瑪.
並有加利利,迦南人拿達葉.
還有西米泰的兩個兒子.
還有兩個門徒.
都在一起做甚麼?.
無所事事 頹.
你見到彼得說「去打魚了」.
其他人竟然說「跟」.
情況像甚麼呢?.
今天無緣無故.
可能大家在打麻將.
突然說「很無聊 不過去喝早茶」.
「等茶樓開門」.
這班人又說「都幾好」.

$^{81}$就跟了.
他們去了打魚.
但到頭來只有一個頹.
你見到他們真的很無所事事.
到一個真是.
都不知道怎樣說的狀態.
但耶穌起耳後.
祂就出現.
祂就出手.
祂很刻意想幫這班無所事事的人.
祂怎樣幫?.
在四至五節可以說.
去暫住.
去叫他們出去打魚.
但你發現一件事.
我們無所事事.
但耶穌出現的同在.
祂是想跟我們一起無所事事.
我自己說到這一句.
自己都想偷笑.
其實覺得很有安慰.
耶穌不是罵我們.
祂是想跟我們一起.
感恩的是.
耶穌透過這種無所事事.
就在做門徒訓練.
很多人說一段門徒訓練.
其實是去嘗試.
回想當初.
去呼召他們做門徒的經驗.
但我不是很認同這個說法.
為甚麼這樣說呢?.
你想想作為當事人的作者約翰.
如果他真的想你每一個讀者.
去聯想這兩件事.
一定會有些提示在當中.
但整卷約翰福音.
連提當時呼召的經歷都沒有.
提甚麼?.
更加另一方面.

$^{121}$你會看到整個呼召裡面.
當初的呼召有彼得,安德烈.
雅各和約翰在當中.
但這裡的七個人.
耶穌沒有理由想做門徒訓練.
只對著那四個.
而其他三個就當做教會人士.
沒有理由的.
還有一個原因.
如果耶穌要說呼召工作的話.
第一句應該會問他.
你在做甚麼?.
但這裡耶穌說甚麼?.
你吃飯了沒有?.
問他在吃飯了沒有.
是完全不同的.
耶穌不是我不斷說廢話的人.
他叫他去打魚似乎另有一個目的.
如果你想知道是甚麼目的.
就看看有甚麼效果.
最快你看到的一個效果.
就是去想那班門徒認得是主.
想那班門徒認得是主.
按照亞當·格蘭特所說的語言.
其實要怎樣面對語言呢?.
首先做一件事.
就是要去認知它.
你能夠提到那些情緒.
去認知自己有甚麼情緒.
發生了甚麼事.
是第一個步驟.
但耶穌不僅幫助那班門徒.
去認知他們那一刻無所事事.
耶穌更加幫助他們認知自己就是主.
去讓那班門徒認得主現在實實在在.
存在.
正在工作.
七罩八積.
你看到那班門徒拿到耶穌是主.
認得是主.

$^{161}$他們的行動是怎樣的呢?.
快了很多.
彼得立即跳進水裡.
噗噗噗.
你看到人生是會有改變的.
當你拿到.
你知道發生了甚麼事.
是有改變的.
當我們很多人遇上一些很頹.
或者一些困局裡面.
上上都會問一件事.
耶穌你死了嗎?.
我不知道你有沒有問過.
我很認真跟你說一件事.
問這句話.
很值得.
因為當你問的時候.
你會知道.
耶穌真的為了我們死了.
死了在十字架之上.
降在陰間.
第三天.
滿有能力地復活.
祂是主.
祂是復活得性的主.
現在正在工作.
當我們知道之後.
我們在當中.
耶穌也帶了另一個層次.
另一個效果在當中.
一個的另一個效果.
祂很想讓那班門徒.
有得吃.
我們看下去.
我們看到.
有得吃.
就是耶穌特意走出來.
要他們打魚的目的.
所以在當中.
我想改一下題目或者更正.

$^{201}$耶穌不單止想跟我們無所事事.
耶穌很想跟你無所事事.
吃吃喝喝.
才是耶穌想做的目的.
我在這裡不是無緣無故地說.
或者隨便說的.
你看看耶穌四句話.
在當中.
在第五,第六,第十,第十二節.
耶穌說的四句話.
其實大部分都跟吃有關.
但是很好奇一件事.
明明耶穌整晚沒有發過屎.
對著那班門徒.
就叫他們去找東西吃.
叫這班沒有發過屎的門徒去找東西吃.
他們辛辛苦苦.
最後帶了153條魚回來.
那些魚.
你想一下.
你一天吃這麼多的魚.
我也覺得應該夠一個月吃的.
他們帶了這麼多外賣回來的時候.
但是耶穌怎樣?.
他早就預備了早餐.
這麼多東西吃不吃.
竟然都預備了早餐.
還要當初幫他們煮好.
等他們回來吃.
那為什麼叫外賣呢?.
最多就是叫他們「我加sung 」.
「我們坐在一起圍爐」.
當然門徒就不像我這樣.
馬上問「圍甚麼」.
這班門徒連問耶穌你是誰都不敢問.
他們怎麼敢問得出這句話?.
你說什麼?.
為什麼這麼無所事事,這麼無聊.
做這些事.
又要我去打魚.

$^{241}$但是自己又幫我預備了早餐.
他們不敢問.
他們很安靜地.
去容讓耶穌拿起餅來.
拿起魚來給他們吃.
我猜那一刻他們不單很安靜地等著.
更加想起一件事.
我再重申.
他們不是想起當初的枯燥.
但是想起另一件事.
大家知不知道是什麼事?.
這裡提到的是提比利亞海.
耶穌好像掌管一切.
但是這班門徒卻已經預備食物.
當中有魚,有餅.
並且分級探望.
約翰在這裡就刻意用了一些字眼提示我們.
你能否想到是什麼?.
我們想到魚.
我們不想到是很正常的.
但是那班門徒一想就想到了.
約翰福音六章一至十五節.
五名二元.
他們回想起那一天.
同樣在提比利亞海.
耶穌掌管一切.
並要門徒去預備食物.
約翰福音六章五節.
耶穌舉目看見許多的人.
就對費力說.
我們從哪裡買餅叫這些人吃呢?.
這裡耶穌很刻意地問他的本地人.
「費力,你在哪裡?從哪裡買餅呢?」.
這裡強調的字眼起源是.
「從哪裡來?」.
似乎想啟發費力去想一想.
去鼓勵他實踐,去找.
不過在六章六節.
費力的回答卻是.
「我只有二十兩銀子,這麼多而已」.

$^{281}$「要買一個餅叫各人吃飽是不夠的」.
你看到費力不是很知道.
耶穌問他.
他真的看見他的資源很緊絀.
但卻沒有回答耶穌問.
「從哪裡來?」.
另一個門徒安德烈就厲害一點.
在六章九節裡.
他回答耶穌.
他說這裡有一個孩子.
帶著五個大麥餅,兩條魚.
只是分級在許多人裡.
還算什麼呢?.
似乎他知道耶穌是資源的來源.
他會安排一切的.
但他的回答裡.
他有一些擔心.
擔心現實不夠需求大.
但最後耶穌怎樣?.
他就出手.
跟這段經文一樣.
他出手.
當他出手的時候.
這幫門徒就知道自己再一次衰老.
他再一次被這幫門徒提醒一件事.
不要再去看現在的現狀.
不要再去看那些資源是否緊絀.
他仍然要相信.
是主已經安排一切去供應.
我們要做的.
是讓耶穌去鼓勵我們去做應該要做的事.
五餅二魚不單止讓我們知道.
主是供應的主.
是主供應.
另一方面.
其實也讓我們回望一件事.
主是恩典的來源.
因為在六章四節.
五餅二魚整個片段裡.
你會看到約翰會刻意記載.

$^{321}$那時猶太人的禹節近了.
耶穌在幫大家一起預備這個禹節的一餐.
他行這個神蹟.
是幫大家去想起禹節.
大家都知道.
猶太人想起禹節就是恩典的記號.
每年猶太人都會回想事主.
去拯救以色列人出埃及.
不是人說事.
不是自己可以完美地處理事主.
所以當中.
我們猶太人每一年都會用禹節的晚餐.
去紀念那份恩典.
那份舊恩.
這份回想.
這份回望.
令到門徒記得起.
所有事都是恩典來的.
去知道自己頂多只是那個途徑.
那個器皿.
或者只是那個見證人.
去見證神那份恩典.
在我們當中.
而我們要做什麼?.
作為那份器皿的途徑.
我們只可以跟著.
以愛還愛.
以恩報恩.
可以想像這一餐圍爐.
應該是一份很靜態的.
一邊吃一邊都沒有人出聲.
不過每一個人.
內心卻不是無所事事.
而是不斷去回想,回望.
我想回望是一樣很重要的事.
我會想像.
當我們去看《聖經》的時候.
我們會想像.
當我們去看《聖經》的時候.
我們會想像.

$^{361}$當我們去看《聖經》的時候.
我們會想像.
當我們去看《聖經》的時候.
我們會想像.
當我們去看《聖經》的時候.
我們會想像.
當我們去看《聖經》的時候.
我們會想像.
當我們去看《聖經》的時候.
我們會想像.
當我們去看《聖經》的時候.
我們會想像.
當我們去看《聖經》的時候.
我們會想像.
當我們去看《聖經》的時候.
我們會想像.
當我們去看《聖經》的時候.
我們會想像.
當我們去看《聖經》的時候.
我們會想像.
當我們去看《聖經》的時候.
我們會想像.
當我們去看《聖經》的時候.
我們會想像.
當我們去看《聖經》的時候.
我們會想像.
當我們去看《聖經》的時候.
我們會想像.
當我們去看《聖經》的時候.
我們會想像.
當我們去看《聖經》的時候.
我們會想像.
當我們去看《聖經》的時候.
我們會想像.
當我們去看《聖經》的時候.
我們會想像.
當我們去看《聖經》的時候.
我們會想像.
當我們去看《聖經》的時候.
我們會想像.

$^{401}$當我們去看《聖經》的時候.
我們會想像.
當我們去看《聖經》的時候.
我們會想像.
當我們去看《聖經》的時候.
我們會想像.
當我們去看《聖經》的時候.
我們會想像.
當我們去看《聖經》的時候.
我們會想像.
當我們去看《聖經》的時候.
我們會想像.
當我們去看《聖經》的時候.
我們會想像.
當我們去看《聖經》的時候.
我們會想像.
當我們去看《聖經》的時候.
我們會想像.
當我們去看《聖經》的時候.
我們會想像.
當我們去看《聖經》的時候.
我們會想像.
當我們去看《聖經》的時候.
我們會想像.
當我們去看《聖經》的時候.
我們會想像.
當我們去看《聖經》的時候.
我們會想像.
當我們去看《聖經》的時候.
我們會想像.
當我們去看《聖經》的時候.
我們會想像.
當我們去看《聖經》的時候.
我們會想像.
當我們去看《聖經》的時候.
我們會想像.
當我們去看《聖經》的時候.
我們會想像.
當我們去看《聖經》的時候.
我們會想像.

$^{441}$當我們去看《聖經》的時候.
我們會想像.
當我們去看《聖經》的時候.
我們會想像.
當我們去看《聖經》的時候.
我們會想像.
當我們去看《聖經》的時候.
我們會想像.
當我們去看《聖經》的時候.
我們會想像.
當我們去看《聖經》的時候.
我們會想像.
當我們去看《聖經》的時候.
我們會想像.
當我們去看《聖經》的時候.
我們會想像.
當我們去看《聖經》的時候.
我們會想像.
當我們去看《聖經》的時候.
我們會想像.
當我們去看《聖經》的時候.
我們會想像.
當我們去看《聖經》的時候.
我們會想像.
當我們去看《聖經》的時候.
我們會想像.
當我們去看《聖經》的時候.
我們會想像.
當我們去看《聖經》的時候.
我們會想像.
當我們去看《聖經》的時候.
我們會想像.
當我們去看《聖經》的時候.
我們會想像.
當我們去看《聖經》的時候.
我們會想像.
當我們去看《聖經》的時候.
我們會想像.
當我們去看《聖經》的時候.
我們會想像.

$^{481}$當我們去看《聖經》的時候.
我們會想像.
當我們去看《聖經》的時候.
我們會想像.
當我們去看《聖經》的時候.
我們會想像.
當我們去看《聖經》的時候.
我們會想像.
當我們去看《聖經》的時候.
我們會想像.
當我們去看《聖經》的時候.
我們會想像.
當我們去看《聖經》的時候.
我們會想像.
當我們去看《聖經》的時候.
我們會想像.
當我們去看《聖經》的時候.
我們會想像.
當我們去看《聖經》的時候.
我們會想像.
當我們去看《聖經》的時候.
我們會想像.
當我們去看《聖經》的時候.
我們會想像.
當我們去看《聖經》的時候.
我們會想像.
當我們去看《聖經》的時候.
我們會想像.
當我們去看《聖經》的時候.
我們會想像.
當我們去看《聖經》的時候.
我們會想像.
當我們去看《聖經》的時候.
我們會想像.
當我們去看《聖經》的時候.
我們會想像.
當我們去看《聖經》的時候.
我們會想像.
當我們去看《聖經》的時候.
我們會想像.

$^{521}$當我們去看《聖經》的時候.
我們會想像.
當我們去看《聖經》的時候.
我們會想像.
當我們去看《聖經》的時候.
我們會想像.
當我們去看《聖經》的時候.
我們會想像.
當我們去看《聖經》的時候.
我們會想像.
當我們去看《聖經》的時候.
我們會想像.
當我們去看《聖經》的時候.
我們會想像.
當我們去看《聖經》的時候.
我們會想像.
當我們去看《聖經》的時候.
我們會想像.
當我們去看《聖經》的時候.
我們會想像.
當我們去看《聖經》的時候.
我們會想像.
當我們去看《聖經》的時候.
我們會想像.
當我們去看《聖經》的時候.
我們會想像.
當我們去看《聖經》的時候.
我們會想像.
當我們去看《聖經》的時候.
我們會想像.
當我們去看《聖經》的時候.
我們會想像.
當我們去看《聖經》的時候.
我們會想像.
當我們去看《聖經》的時候.
我們會想像.
當我們去看《聖經》的時候.
我們會想像.
當我們去看《聖經》的時候.
我們會想像.

$^{561}$當我們去看《聖經》的時候.
我們會想像.
當我們去看《聖經》的時候.
我們會想像.
當我們去看《聖經》的時候.
我們會想像.
當我們去看《聖經》的時候.
我們會想像.
當我們去看《聖經》的時候.
我們會想像.
當我們去看《聖經》的時候.
我們會想像.
當我們去看《聖經》的時候.
我們會想像.
當我們去看《聖經》的時候.
我們會想像.
當我們去看《聖經》的時候.
\newpage



\section{}
\label{sec:b3sjDQwWdlg}
\textbf{Family Sunday School - EP 1 (Between Christmas and Easter)}
\newline
\newline
連結: \href{https://youtube.com/watch?v=b3sjDQwWdlg}{\texttt{https://youtube.com/watch?v=b3sjDQwWdlg}} ~~~~ 語音日期: 2022-02-22
\newline
\newline
\hyperref[sec:INsxhCoyb9U]{\small{< < < PREV SERMON < < <}}
~
\hyperref[sec:index_chronic]{\small{[返順時目]}}
~
\hyperref[sec:index_scriptual]{\small{[返順卷目]}}
~
\hyperref[sec:bkEucgq83Sg]{\small{> > > NEXT SERMON > > >}}
\newline
\newline
$^{1}$(upbeat music).
(speaking in foreign language).
Lego?.
Play-Doh?.
Games?.
(speaking in foreign language).
(laughing).
(speaking in foreign language).
Jesus's birthday..
(speaking in foreign language).
A big section of our family Sunday school.
(speaking in foreign language).
Jesus's ministry..
(speaking in foreign language).
Okay, here we go..
(speaking in foreign language).
Happy birthday, Jesus..
(speaking in foreign language).
Friends are forever, friendships..
(speaking in foreign language).
Miracles..
(speaking in foreign language).
Stories and parables..
(speaking in foreign language).
Huge mind shit..
(speaking in foreign language).
Upside down, right side up..
(speaking in foreign language).
Birthday..
(speaking in foreign language).
(clapping).
(speaking in foreign language).
One, two, three..
(speaking in foreign language).
One, two, three..
(speaking in foreign language).
Friends..
(speaking in foreign language).
Disciples..
(speaking in foreign language).

$^{41}$12 disciples..
(speaking in foreign language).
, pause..
(speaking in foreign language).
, pause..
(speaking in foreign language).
Birthday..
(speaking in foreign language).
James..
(speaking in foreign language).
John..
Matthew..
Mark..
Simon..
Thomas..
Peter..
Andrew..
Judas..
Thaddeus..
Philip..
(speaking in foreign language).
Bartholomew..
(speaking in foreign language).
Lazarus..
Lazarus, Mary, Martha..
(speaking in foreign language).
Miracles..
(speaking in foreign language).
Miracles..
(speaking in foreign language).
(sighing).
(speaking in foreign language).
Right, water turns to wine..
(speaking in foreign language).
Ready?.
(sighing).
(speaking in foreign language).
(sighing).
(speaking in foreign language).
Jesus calmed the storm..

$^{81}$(speaking in foreign language).
Ready?.
(speaking in foreign language).
Okay..
(sighing).
(coughing).
(sighing).
(speaking in foreign language).
Lazarus, get up..
Lazarus, (speaking in foreign language).
Wow..
(speaking in foreign language).
Stories or parables..
(speaking in foreign language).
Parables..
(speaking in foreign language).
Sort of like Pictionary..
(speaking in foreign language).
(paper rustling).
(speaking in foreign language).
(paper rustling).
(speaking in foreign language).
The story of the lost lamb..
(speaking in foreign language).
One out of a hundred..
(speaking in foreign language).
Of the story of the parable..
(speaking in foreign language).
Okay?.
All right..
(speaking in foreign language).
Ready?.
(speaking in foreign language).
Okay..
First one, ready?.
(speaking in foreign language).
Discuss, okay?.
Go..
(speaking in foreign language).
Mind shift..

$^{121}$(speaking in foreign language).
Mind shift..
(speaking in foreign language).
Like a thinking bubble..
(speaking in foreign language).
It's total opposite..
(speaking in foreign language).
I'll tell you..
(speaking in foreign language).
A substitute, a lamb..
The ultimate lamb, that's Jesus..
(speaking in foreign language).
I can't do it alone..
(speaking in foreign language).
I need help..
(speaking in foreign language).
I need help from God..
(speaking in foreign language).
I need Jesus..
(speaking in foreign language).
Because he's perfect..
(speaking in foreign language).
He can substitute me for my sin and things I've done wrong..
Whoa..
(speaking in foreign language).
What do you mean?.
(speaking in foreign language).
Friendship, stories..
(speaking in foreign language).
When our attitude and our hearts change..
(speaking in foreign language).
God knows..
So I could send Jesus..
(speaking in foreign language).
We are touched by love..
And that's hard sometimes..
(speaking in foreign language).
Mind shift..
(speaking in foreign language).
Extension time, extension time..

$^{161}$(speaking in foreign language).
Family Sunday school..
(speaking in foreign language).
Okay, here we go..
Extension time..
(speaking in foreign language).
Friendship..
(speaking in foreign language).
Miracles..
(speaking in foreign language).
Stories..
(speaking in foreign language).
Mind shift..
(speaking in foreign language).
But through it all..
(speaking in foreign language).
We don't, adults don't,.
but to go through the journey together..
Okay?.
(speaking in foreign language).
Easter..
(speaking in foreign language).
All right, good..
Okay, after you go,.
(speaking in foreign language).
Okay, coming up next, Chinese time..
(speaking in foreign language).
American-born Chinese..
(speaking in foreign language).
So, (speaking in foreign language).
it's okay..
(speaking in foreign language).
Counting words..
(speaking in foreign language).
What is the difference?.
(speaking in foreign language).
If you're ever confused with (speaking in foreign language).
All those, I'm here for you..
I hope, I hope..
Okay, here we go..

$^{201}$Why is (speaking in foreign language).
as in pen, pencil, marker, (speaking in foreign language).
You can't say (speaking in foreign language).
but you can't..
That doesn't make sense..
(speaking in foreign language).
And then, (speaking in foreign language).
Google accounting word is about,.
you'll see, the shape of (speaking in foreign language).
Right?.
(speaking in foreign language).
This is how I remember..
(speaking in foreign language).
It's like a (speaking in foreign language).
(speaking in foreign language).
is like a branch, like a stick..
Right?.
Stick..
(speaking in foreign language).
Branch, a stick of wood..
(speaking in foreign language).
Lipstick..
Oh, big clue..
Stick..
(speaking in foreign language).
Chapstick..
(speaking in foreign language).
That's when you use (speaking in foreign language).
Okay?.
Yeah..
That's my clue for you..
Another thing..
(speaking in foreign language).
Sheets of paper..
50 sheets..
(speaking in foreign language).
100..
(speaking in foreign language).
150..
(speaking in foreign language).

$^{241}$How many sheets?.
(speaking in foreign language).
So sheets..
Like this..
Right?.
(speaking in foreign language).
Rectangular..
(speaking in foreign language).
It comes in sheets..
Those are called (speaking in foreign language).
Okay?.
(speaking in foreign language).
One sheet of paper..
(speaking in foreign language).
That's when you use (speaking in foreign language).
I know..
Not my fault, this is Chinese..
At least..
(speaking in foreign language).
Okay..
Ready?.
(speaking in foreign language).
Just a shape, right?.
Stick or thin rectangle..
(speaking in foreign language).
Also, (speaking in foreign language).
The size, right?.
(speaking in foreign language).
Like drops..
Drops..
(speaking in foreign language).
Little bits, little pieces..
(speaking in foreign language).
Like, (speaking in foreign language).
Skittles..
(speaking in foreign language).
Rock..
(speaking in foreign language).
Pebbles..
(speaking in foreign language).

$^{281}$Right?.
Little pebble..
Little rock..
Stone..
Maybe a little bigger..
(speaking in foreign language).
Stone, right?.
(speaking in foreign language).
Okay?.
Boulder..
Big rock..
Huge rock..
Boulder..
(speaking in foreign language).
Whoa..
Okay, difference..
(speaking in foreign language).
Right?.
(speaking in foreign language).
Got it?.
(speaking in foreign language).
But you probably say, (speaking in foreign language).
'cause you add the (speaking in foreign language).
it makes it smaller too..
I know..
This is why..
Cantonese is not easy..
Okay, last one..
Are you ready?.
Last one for today..
(speaking in foreign language).
Okay..
(speaking in foreign language).
No..
(speaking in foreign language).
Is it like a stick?.
No..
(speaking in foreign language).
Is it like a sheet?.
Rectangular sheet?.

$^{321}$(speaking in foreign language).
Here we go..
Listen up..
(speaking in foreign language).
What does that mean?.
(speaking in foreign language).
There's like a handle..
(speaking in foreign language).
There's a handle..
Okay?.
Just like (speaking in foreign language).
(speaking in foreign language).
Sword..
(speaking in foreign language).
As a handle..
(speaking in foreign language).
Pull out the sword..
Just like (speaking in foreign language).
There's a handle..
That's all I got, guys..
Hope that helps your Chinese, folks..
(speaking in foreign language).
But it's okay..
(speaking in foreign language).
See you next time..
Bye..
(upbeat music).
\newpage



\section{創世記}
\label{sec:bkEucgq83Sg}
\textbf{創世記第4章 (4.2) - 背景釋經 [亞達哈斯史詩 vs 創世記] (陳崇基牧師)}
\newline
\newline
連結: \href{https://youtube.com/watch?v=bkEucgq83Sg}{\texttt{https://youtube.com/watch?v=bkEucgq83Sg}} ~~~~ 語音日期: 2022-02-22
\newline
\newline
\hyperref[sec:b3sjDQwWdlg]{\small{< < < PREV SERMON < < <}}
~
\hyperref[sec:index_chronic]{\small{[返順時目]}}
~
\hyperref[sec:index_scriptual]{\small{[返順卷目]}}
~
\hyperref[sec:1enAW_IL3Aw]{\small{> > > NEXT SERMON > > >}}
\newline
\newline
$^{1}$今天我們看到創世記的時候.
我們想試試做一些比對研究的工作.
我和你分享一個史詩.
這個史詩是一個巴比倫的史詩.
這個巴比倫史詩叫做《阿達哈斯史詩》.
《Adrachasis Epic》.
這個是屬於公元前1800年的時間.
大概是阿布拉罕的時間.
大家記得阿布拉罕是從舊巴比倫的地帶出來的.
當時是舊巴比倫的帝國.
講到一個王叫做Adrachasis.
他的名字的意思.
阿達哈斯.
就是大有智慧的意思.
Exceedingly wise.
這個人.
有一個藝術家將這個故事畫了畫.
這樣就比較容易讓我們了解.
到底裡面發生了什麼事.
我們讀這個故事的時候.
想一想有什麼地方和創世記相似.
有什麼地方和創世記相似.
但是有很大分別.
這個《阿達哈斯史詩》.
開頭的時候是這樣說的.
他說起初地上有男神女神.
男神女神不是那些帥哥帥妹.
創世之前的眾神官.
所以有男神有女神.
有一些低等的神.
他們要耕作土地才有東西吃.
他們想神明會餓.
是需要吃東西的.
因為有很多苦工要做.
所以非常艱辛.
生活很困難.
然後每一位的神.
就有苦工要做.
有些神明要去挖土去耕種.
有些就要去田裡打水.

$^{41}$然後做一些灌溉的溪.
又要去拔野草.
好像他們要耕田一樣.
因為這麼辛苦的緣故.
這一班低級的眾神.
就很不開心.
他們就很不愉快.
他們聚集在一起.
就去討論.
怎樣可以減輕他們的工作.
所以這一班低等的神.
他們就去找那位上面的大神.
這個大神是一個很名字的神.
叫做恩基.
E-N-K-I.
恩基.
就問他到底怎麼搞.
恩基就是陰間大水的神.
他在家裡睡覺.
就吵醒了他.
然後在跟他商談之下.
他就建議了.
恩基建議.
他會做一些活物出來.
他們就會代替這些眾神去工作.
他們會做苦工.
然後他們會耕種.
耕作這些田地.
於是.
這些眾神.
都認為這個計劃非常好.
這個恩基.
就在他的大水之家裡.
拿起泥土.
就去造出人類.
然後他就將這些氣息.
吹進這些泥漿裡.
就造出了七對人類.
使他們成為活物.
但是又限制了他們的壽命.

$^{81}$只是這些眾神.
有所謂永生.
人就限制了他們的壽命.
於是.
這些人類就被當奴隸.
做了出來之後.
他們就在田間裡工作.
上面的眾神.
男女.
這些眾神.
就有東西吃了.
就等於一些中國民間信仰裡面.
我們去供奉一些神靈.
他們就有飽肚.
供應他們上面的飲食.
人類是做什麼的呢?.
人類就是眾神.
女神的奴僕.
就是做一頓夢的.
就是為了他們能夠得食.
討好他們.
於是人類就要在河流裡打水.
然後灌溉一些毫無生氣的土地.
在這裡挖土來耕種.
人類就很艱辛.
來到這裡工作.
使地能夠長出一些東西.
眾神就很高興.
這個故事繼續發展下去.
是很長的.
就說人類生養眾多.
然後噪音就很大.
因為人多了.
噪音就很大.
就擾神清夢.
有一個大神就很不高興.
經常睡不著覺.
因為人類越來越多.
就吵著他睡覺.
然後他就降一些災禍給你.

$^{121}$就殺了一些人.
然後又再降災禍.
又殺了一些人.
最後就有洪水出現.
洪水出現的時候.
之後我們會說一些洪水的故事.
近東的洪水故事.
誰知道有個人就被其中一個神通知了.
他就做了一艘船.
做了一艘船的時候.
他就在洪水裡面生還了.
就留下了.
這個故事我們暫時說到這裡.
你就會看到有很多地方.
是跟聖經的記載有相似的地方.
我們不需要嚇到.
有沒有搞錯.
聖經是不是抄襲呢.
或者有些人會覺得.
因為這個緣故所以聖經抄錄了其他古近東的故事.
其實不需要這樣想.
正如我剛才所說的.
我剛才的圖你記得嗎.
其實所有這些民族包括以色列人.
都是在同一個近東的世界裡面.
他們的文化是大家分享的東西.
所以沒有抄不抄襲的問題所在.
正如我今天所說.
譬如我說今天有酒今朝醉.
我不知道你知不知道那個典故是什麼.
但我們這樣說的時候.
我們不是抄襲原本說今天有酒今朝醉的詩人.
不是在引述他.
而是因為文化裡面.
世界裡面已經預設了.
一種這樣的人生觀.
或者一種這樣的觀念.
今天有酒今朝醉的觀念.
所以我們說的時候.
其實就是在同樣的文化裡面.

$^{161}$我們想分享的文化裡面.
我們說出一些共同有的觀念.
所以不是說我抄襲或是誰抄襲.
大家明白我的意思了.
因為文化是不會這樣運作的.
有很多東西在文化裡面已經根深蒂固的時候.
大家在同一個世界裡面去說一些觀念的時候.
就不是引述或者抄襲某人.
如果你來美國的話.
我就會告訴你美國有日光節省的時間.
在美國這個社會裡面.
因為我們有一個共同的文化.
我們一說Daylight Saving.
日光節省的時間的時候.
大家就明白了.
但沒有人說日光節省時間.
其實沒有人說這件事.
就是去引用第一次有日光節省時間的時候的條文.
只不過我們在說我們有一個同樣的文化.
所以文化是這樣運作的.
不是今天有些人說.
哦 這些聖經的故事.
好像巴比倫的故事.
一定是聖經抄襲之前的.
不是抄襲雜文.
而是在同一個傳統和文化的背景裡面.
有同樣相似的觀念.
我們回到這個學者.
最後的時候.
他就將亞特哈斯史詩.
分段.
你發現分段的時候.
這個學者是誰呢?.
浸會神學院的神學家.
他就將他分段分成五段.
首先你記得那個故事開始的時候說什麼?.
就是神創造人.
然後有危機出現了.
就是妖神清蒙.
人生養眾多.

$^{201}$於是神明就降溫營.
然後有一個恩基獻祭.
阻止了.
讓人們不會全死.
然後到第二個環節.
就是另外一個危機.
又是生養眾多.
然後神明就很不高興.
因為妖居清蒙.
然後又降一個旱災.
然後又再加劇這個旱災.
然後這個恩基又獻祭.
然後到第三段.
就是人類又是生養眾多.
又是妖神清蒙.
於是那個大神就降大洪水.
正如我所說.
亞特哈斯就讓恩基獻祭.
叫他造一隻.
所謂方舟.
他就得救.
得救之後亞特哈斯生還.
他獻祭之後.
人類又繼續生養眾多.
然後眾神就有一個妥協.
就說不要再降災.
只不過限制住人口.
不要太多.
不會吵著我們.
不會阻礙我們睡覺.
他就將.
《創世記》一至十一章.
其實是同一個脈絡.
同一個脈絡來講什麼.
這些.
頭十一章.
和亞特哈斯史詩.
他的寫作.
的.
體材.

$^{241}$其實都是一樣.
有一些起源史.
起源史.
近東的人想起起源史的時候.
就講到神靈創造人類.
然後.
人類遇見不同的災難.
最後就是洪水.
洪水之後.
生還.
這個就叫起源史.
他想起起源史的時候.
就是這樣寫作.
所以為什麼《創世記》一至十一章.
都是按照這樣的脈絡.
除了因為歷史性.
這樣發生之外.
其實就是用近東人.
怎樣想起源史.
你說很奇怪.
當然奇怪.
因為不是我們的觀念.
我們的觀念是想什麼.
大爆炸.
想起.
以前大地怎樣移動.
等等.
這些科學的問題.
他們想起的就是.
這一種的故事.
所以《創世記》一至十一章.
絕對不是驚訝的.
如果我們有比對研究.
明白當時的人.
怎樣寫.
所謂的Protohistory.
即是由人類起初.
那些故事的時候.
《創世記》一至十一章.
他的脈絡.

$^{281}$其實就是很清楚.
結束的時候.
我們就看看比對研究.
給我們怎樣看得更加清楚.
《創世記》一至十一章.
到底想說什麼.
比對研究.
可以給我們.
首先我們看.
阿達哈斯.
巴比倫的起源史.
和《創世記》一至十一章.
所有的故事.
經文敘述.
你就可以做幾樣.
不同方面的比較.
有相似的地方.
也有差異的地方.
相似的地方.
或差異的地方.
是什麼.
我們首先看神學.
怎樣看神.
神觀.
人怎樣看神觀.
在巴比倫的文學裡面.
剛才我們看的其中一個例子.
就是多神.
並且這些神.
互相作對的時候.
你都不知道信哪一個比較好.
只不過阿達哈斯.
信了那個叫做恩基.
那就死不去.
神呢.
他們想的神明.
就是按照人的形象.
所做的.
所以神會餓.
你就要初一十五.

$^{321}$給他東西吃.
否則他發怒就不行了.
最好就哄他開心.
他就會祝福你.
交易.
神呢.
這些神明是需要睡覺的.
所以吵著他就不好.
所以神靈.
神觀是按照人的形象.
我們需要睡覺.
我有小孩.
小孩多了他就吵著我睡覺.
我會發脾氣.
跟我一樣.
我會肚餓.
所以他們想神的時候.
就是用我們人的角度去想.
並且.
這些神明.
就不想人生用這麼多.
因為多的時候.
增加了.
人口就會吵著我睡覺.
所以神就很不喜歡.
這個就是.
巴比倫人的神觀.
你比較一下創世記.
說的話截然不同.
因為是一神.
並且神的旨意很清楚.
很一致.
不是出而反而.
神的本性.
是信實.
是真理.
神也是超越的.
所以聖經.
創世記第一章.
人是按照神的形象去做的.

$^{361}$不是相反的.
人似神.
但是神不似人.
神供應百物給人.
你見到第一章.
裡面有些形容.
神創造大地.
給人吃東西.
供應人的需要.
並且祝福他們.
我們要生養眾多.
範顯 天本傳遞.
與巴比倫人完全是.
截然相反.
所以用同一個方式.
去表達一些很不同的.
神學觀念.
人類.
他的角色 身份.
你記不記得.
有泥土造 造七對人出來.
那七對人.
不是科學的.
製造過程.
而是.
身份.
人是被造.
有什麼身份.
人被造是低端的人口.
低端的奴隸.
他們是做什麼的.
他們是服侍神明.
獻祭給他.
到時候.
帶些東西去廟.
然後獻祭給他.
很清楚有階級高低之分.
人類.
被造.
是沒有什麼計劃的.

$^{401}$因為有些眾神.
不高興.
工作太辛苦.
不如做些奴隸出來.
讓你不用那麼辛苦.
是後來的想法.
不是原本想的.
一種尊貴的過程.
但在聖經裡面.
第一二章.
有相似的地方.
人被造如塵土.
但是有君尊的地位.
神的形象.
人人都是平等的.
不是君王在上.
才是高位置.
高位置.
每個人都是平等.
說到起源史.
有什麼相似的地方.
格式脈絡.
結構相似.
主旨也相似.
巴比倫的.
這些史前史.
起源史.
都是想表達什麼人類的身份.
價值.
地位等等.
是世界觀.
聖經就去爭辯這種世界觀.
就是說聖經說的才是對的.
聖經的世界觀才是對的.
就是人人是平等的.
人被造.
有一個尊貴的身份.
尊貴的身份.
並且.
這些起源史很多時候會提到.

$^{441}$人類文明的開始發展.
建造城市.
有農業.
畜牧業.
或者.
一些煉銅.
煉鐵.
這些工業出現.
發明出現.
但是每逢.
說到這些文明的.
發展的時候.
都是眾神給人間的東西.
就是將這些文明的發展.
人類文明的發展.
Civilization.
捍衛神聖化.
大家記得.
在《創世記》第四章之後.
說到戈隱的家譜.
我們就看到戈隱的後代是什麼.
是畜牧業的先祖.
是音樂的先祖.
是不同工業的先祖.
等等.
這些如果我們明白.
比對研究的時候.
明白當時近東的人.
如何想起源史的時候.
如何說到人的文明.
我們就完全沒有任何驚訝的東西.
因為.
說起源史的時候.
一定會說到人類的文明.
人類文明的發展.
有城市的建設.
有畜牧業等等.
所以《創世記》第四章.
但是有一個分別.
這些人類文明就是去神華化.

$^{481}$可以這麼說.
就是說人類所建造出來的東西.
有畜牧業等等.
完全都是人類自己.
管理萬物的時候.
得到的.
結果.
並且《創世記》.
想焦點放在哪裡.
就不是去推崇.
這些文明.
人類很厲害.
能夠有科技.
能夠有城市.
有不同的發展等等.
《創世記》一至十一章.
是說罪惡.
越來越擴散.
破壞這個世界.
所以最後有洪水出現.
等等.
這些下一次我們會繼續.
去解釋多一點.
但是這裡是一個工具.
給大家去了解.
為什麼我們需要了解.
近東的文獻.
絕對不是因為.
我們當聖經.
和這些所謂.
異邦異教的文獻.
去相提並論.
而是我們知道當時的世界.
給我們一個視窗.
我們知道他們是怎麼想.
這個宇宙.
怎麼想人類.
而上帝是藉著.
這個觀念.
去表達他自己.

$^{521}$所以這裡是單單.
一種工具.
給大家一些參考.
暫時這次的課.
就說到這裡.
如果喜歡我們的影片.
就記得要讚好,留言和分享.
\newpage



\section{6:1-2}
\label{sec:1enAW_IL3Aw}
\textbf{First of the Seven Seals (Rev 6\_1-2) - Rev. Mark Chan}
\newline
\newline
連結: \href{https://youtube.com/watch?v=1enAW-IL3Aw}{\texttt{https://youtube.com/watch?v=1enAW-IL3Aw}} ~~~~ 語音日期: 2022-02-28
\newline
\newline
\hyperref[sec:bkEucgq83Sg]{\small{< < < PREV SERMON < < <}}
~
\hyperref[sec:index_chronic]{\small{[返順時目]}}
~
\hyperref[sec:index_scriptual]{\small{[返順卷目]}}
~
\hyperref[sec:X5Jubq7tdoc]{\small{> > > NEXT SERMON > > >}}
\newline
\newline
$^{1}$- Right..
Well, this is a time of craziness..
It's a time of chaos..
We still in the midst of a pandemic,.
two years long, right?.
Probably something at the end of a tunnel..
I don't know yet, right?.
We don't know yet..
So we just put on masks for now, right?.
Except those who are singing or speaking,.
just for the sake of clarity..
So you look around, there's so much uncertainty..
And last week, the horrible news.
of Russians invading Ukraine, that is horrible..
You read all the stories and all the news reports..
It's just time of chaos..
Sometimes in times of crisis,.
we will think of what the Bible says..
We try to find answers from the word of God..
And sometimes we think of words from the Bible.
saying in the end times,.
there will be war and conflicts on the earth..
Or sometimes we think of the book of Revelation, right?.
I'm sure you have read this book..
It's very easy to find as the last book of the Bible,.
the last book in the New Testament..
But it's a very strange book..
It's a very difficult book to tackle..
So today, we would like to start a kind of a mini series..
We'll see how it goes..
And we will take some passages from the book of Revelation.
and learn how to read it, learn how to understand it..
Especially we know the Bible,.
even though it's not directly written to us,.
but it's definitely written for us..
So how do we read that?.
I chose this book, book of Revelation,.
because at times of crisis,.
sometimes we go to this book and say,.
you see, this is what the Bible predicted..

$^{41}$And this is what the book of Revelation is all about, right?.
We look at the conflicts, happened last week,.
and we look at all these tensions,.
international tensions or pandemic,.
or just unrest, just look around..
And we go back to the Bible and say,.
hey, this is what the book says, right?.
This is what the book predicted..
So how true is that, right?.
So I want you to think with me how to understand this book,.
how to understand the book of Revelation..
If you remember last week,.
I told you to do something, right?.
Every day, turn to the Bible,.
10 to 15 minutes in prayer, right?.
Reading the word of God, you know,.
absorbing the word of God..
Did you do that?.
All right, I'll make sure you do it, okay?.
Make sure you do it..
But sometimes there are passages in the Bible,.
some portions in the Bible,.
it's very difficult to do devotion those out of, right?.
Including the book of Revelation..
So we wanna turn to this book, book of Revelation..
If you have your Bible in your phone app, turn to it now..
In your languages, there will also be passages.
we project on the screen..
We'll start with chapter one..
The revelation of Jesus Christ,.
which God gave him to show his servants.
what must soon take place..
He made known by sending his angel to servant John,.
the apostle John..
He testifies to everything he saw..
That is the word of God, the testimony of Jesus Christ..
Blessed is the one who reads the words of the prophecies..
The book of Revelation, this book..
Blessed is the one who reads the words of this book..
And blessed are those who hear it and take to heart.

$^{81}$what is written in it..
Because the time is near, the time of Jesus' second coming..
He came 2,000 years ago into human history.
and accomplished salvation for us on the cross..
But he will come back..
The book of Revelation is all about the second coming..
What will happen, what will soon take place?.
The first word, apocalypses..
This is the word very familiar to us, right?.
Because we get the English word apocalypse out of this word..
Remember two years ago when we first locked down,.
I went out with my two children to pick up some takeouts.
from a restaurant..
It was March, right?.
2020, right?.
2020..
So it was at night..
So we drove out there to a restaurant nearby..
We're trying to pick up some takeouts for dinner..
It was the first day of lockdown and we drove on the street..
And it was quite a sight..
Because as you know, in Chicago,.
you know, the traffic is crazy, right?.
All the time..
But it was the first day of lockdown.
and the street is totally empty..
It was like a ghost town..
It was kind of crazy, right?.
This is how we describe situations like this..
Apocalypse, right?.
It's the modern meaning of the word apocalypse..
Means doomsday, right?.
Destruction of all kinds..
But originally in the first century,.
when this word is spoken,.
the basic meanings means uncover, disclose..
Something which was unknown before..
Now it's being revealed..
So that is what apocalypse means..
The revelation of Jesus Christ..

$^{121}$Reveal something to us that is not known before..
What is it?.
What must soon take place.
before the second coming of Christ?.
And then the passage goes on to explain to us.
what is the situation..
Verse four, John the apostle,.
writing a letter to seven churches in the province of Asia,.
which is today's Turkey..
John, your brother, I am with you guys..
Company in suffering and kingdom and patient..
And patience..
And doings..
There are hours in Jesus..
Who, John is saying that I'm writing it.
from the island of Patmos,.
which is a small island.
off the coast of Ephesus from Turkey..
So he was exiled, why?.
Because the Roman Empire was oppressing the churches,.
the early churches in those days..
It was kind of the 1890s,.
so '90s after Christ..
So the Roman Emperor was instituting.
some sort of systematic persecution,.
systematic pressure on the early church..
So it was time of crisis for the church..
And then John was exiled.
because he was one of the leaders, one of the pastors..
He was exiled on Patmos, on the island..
And he was suffering with other brothers and sisters.
who were in the mainland under the Roman persecution..
So this is what the book of Revelation is all about..
Sometimes we describe the book of Revelation.
as some special type of literature..
We call this type of literature apocalyptic literature..
And what it says is it's a group.
or a collection of many different writings,.
different kinds of writings..
Some of them end up in the New Testament..

$^{161}$Some of them is outside the New Testament..
Some are written by Jewish writers,.
some are by early Christian writers..
And they have some commonalities..
You know, think of Old Testament, the book of Daniel,.
or some chapters in Ezekiel..
Some of the passages in the Gospels or Paul's letters..
Some common languages all these writings would use,.
including the book of Revelation..
So as we turn to the book of Revelation,.
it's a strange book, right?.
It talks about horses, it talks about locusts,.
it talks about all kinds of strange phenomenon.
that is foreign to us..
But you need to remember this..
This sort of writings, apocalyptic literature,.
is written to the people of God who are in crisis.
in order to encourage them,.
in order to let them see what the real reality is..
It's not what you see here now..
It's something happening in the future,.
and something happening above,.
in another realm, in the heavenly realm..
What is really happening?.
You see conflicts on Earth,.
but there's spiritual warfare,.
which is so real up there and in soon future..
So the book of Revelation or apocalyptic literature.
is trying to write to,.
to say to the people of God who are under pressure,.
their faith is being tested..
And the purpose is to help them see a different reality.
so that they will find hope in God..
They will not find hope on Earth..
They will be encouraged in times of chaos..
Don't we all need that?.
We need this today, right?.
We see conflicts and we see chaos.
and uncertainties,.
all sorts of crazy things happening in our world..

$^{201}$What is going on?.
And the apocalyptic literature,.
especially the book of Revelation,.
is trying to talk to us..
It's a message written for us.
so that we will understand what God is doing..
So the first three chapters,.
we're not going to go into details,.
but first three chapters is actually Apostle John.
writing seven letters to seven churches..
And they will circulate those letters.
and to encourage the churches.
who are under trials and temptations..
And then in chapter four,.
the Apostle John is being lifted up,.
being taken away from the Earth..
And he entered into the heavenly realm..
And he saw something like a Roman Empire royal court,.
like the king's court, right?.
Going to the White House up there..
So he looked around and see who is seated on the floor..
It's not the Roman emperor..
It's not Caesar..
He saw who is seated on the floor is God himself..
Is the God of history..
The God who's still in control..
And then it launched many, many different chapters.
talking about the visions that he saw..
And he reported back to the churches..
Let's take a look..
Chapter five..
Let me get this ready..
It's so dark here..
Hold on one second..
All right..
Chapter five, verse one..
See that?.
Chapter five, verse one..
Then I saw in the right hand of him,.
who sat on the floor,.

$^{241}$a scroll with writing on both sides.
and sealed with seven seals..
I saw a mighty angel proclaiming in a loud voice,.
who is worthy to break the seals and open the scroll..
But no one in heaven or on earth,.
under the earth could open the scroll.
or even look inside it..
So I wept and wept.
because no one was found who has authority,.
who was worthy to open the scroll.
to see what is written inside..
Now in the ancient world,.
important information is always written on the scroll..
Made of animal skin or paper..
And then they roll it up.
and then they put some lumps of clay.
on top of the roll up document..
And then someone who has authority.
with a ring or with a stamp,.
will put the stamp on the lump of clay.
and seal it, close it up,.
and archive it in the royal bank.
or royal building or temple, anywhere important..
In the, say in today's language,.
the White House office, right?.
In the government archive..
So whoever wants to check what was written in the document,.
he has to bring his stamp or the ring.
and to match with the impression on the clay.
to break the seal and open it up..
Now what John saw was this..
He saw a scroll with seven seals..
The document is closed up..
But he found no one who could open the scroll..
What is written inside the scroll.
is definitely words of judgment..
If you have some Old Testament background,.
you understand a scroll written,.
words with front and back all over the place..
Book of Ezekiel, talking about the word of judgment from God..

$^{281}$He will vindicate the oppressed,.
vindicate the people of God..
He will help the people of God.
and will bring judgment on the oppressor,.
on the evil forces that is pressing down on the church..
John found no one having the authority to open the scroll..
So what will happen?.
Nobody knows..
What is the fate?.
What is happening to the church?.
What will look like tomorrow?.
Nobody can tell..
That's why he was weeping..
Nobody can be found to open the scroll,.
to tell us who is in control tomorrow, right?.
Especially in the church.
when they were under so much pressure..
Then in verse five,.
then one of the elders who was worshiping in the throne room.
said to me, "Do not weep..
See the lion from the tribe of Judah, the root of David..
He has triumph..
He is able, he has the authority to open the scroll.
and his seven seals..
Finally, the lamb of God, who was slain and rose again..
He is also the lion from the tribe of Judah,.
from the tribe of David..
He is the real conquering king..
He has conquered death and evil,.
and he has the authority to open and break the seal,.
to help us see what will soon happen..
Who is in charge tomorrow?".
So chapter six, six verse one..
"I saw as the lamb opened the first of the seven seals,.
that I heard one of the four living creatures, the angels,.
say in a voice like thunder, 'Come!'".
Verse two..
"I look, and there before me was a white horse..
Its rider holding a bow..
He was given a crown, and he rode out as a conqueror,.

$^{321}$bent on conquest.".
Now definitely this is some strange images for us..
For sure, this is negative..
It's a world of judgment, right?.
So definitely it's judgment of God.
against the church oppressor, right?.
God will do something to stop the oppression..
God will do something to the evil empire..
And definitely in this image, we see war, white horse,.
rider holding a bow, conquest, military image..
Now here's what we struggle with,.
because we modern readers, we modern people,.
we often misunderstand the images.
from the book of Revelation..
There are all sorts of, we're trying to read our own meaning.
into the text, right?.
We try to impose our own understanding onto the text..
Let me demonstrate how,.
you probably heard something like this..
Many people would misinterpret the images here..
Let me demonstrate how to do it wrongly, okay?.
All right?.
Let me try to demonstrate how to do it wrongly..
And this is so wrong, okay?.
So you see in this image, there's color white..
There's white horse, right?.
And you see, no, military conflict, right?.
Conquest..
Let me demonstrate..
You see the flag of Russia has white color, right?.
White, blue, and red, right?.
White, bingo..
You got Revelation chapter six, white horse, right?.
And then if you Google white horse Putin,.
you find pictures like this, right?.
Bingo..
Revelation six, right?.
That's the first seal being broken, right?.
And then you say, "Ha ha, you know, prophecy fulfilled..
"2,000 years ago it was written there, right?.

$^{361}$"Chapter six of Revelation, the first seal, right?.
"Last week it happened, it broke the seal.".
(laughs).
You see a lot of misinterpretation,.
just so wrong these days, you know?.
Because we just don't know how to interpret.
this sort of, you know, imageries..
They're strange to us..
Now I want to tell you how to read the Bible,.
especially the book of Revelation..
The Bible, the book of Revelation,.
is not written to us, right?.
It's not directly written to us..
It was originally in Greek, not in English,.
not in Chinese, not in any contemporary languages..
It was written in Greek..
It's not written to us, but it's written for us..
So when we try to understand what is going on here,.
we should stop from reading our own meaning onto the text..
We should ask, how did the first audience,.
the early churches in the first century,.
how would they understand these images, right?.
So this is totally wrong, okay?.
Don't take it home, all right?.
Don't take it home..
This is the wrong interpretation..
But you see this kind of problem.
whenever we try to understand.
what the book of Revelation is all about?.
Don't do that..
Instead, we ask, how would the first century Christians.
understand the imageries?.
White horse rider with a bow..
Let me show you some..
Ancient coins..
Parthian coins, all right?.
Do you see a soldier holding a bow?.
Yeah, the triangular thing, right?.
It's not very clear, but it definitely is a bow, okay?.
And then underneath, horses, right?.

$^{401}$White horse..
And the king was riding on the horse..
These are Parthian coins..
And it just happened for the ancient Parthian empire,.
the sacred color happens to be white..
In their army, there's always white horses..
Let's take a look at the map..
The Roman Empire on the left,.
the Parthian Empire on the right..
For 400 years, this is the ancient, powerful empire..
It always posed a military threat.
to the eastern border of Roman Empire, the Parthian Empire..
In fact, the Roman historians,.
when they talk about the threat against the empire,.
the old Roman Empire,.
they would talk about the Parthian Empire..
In those days, even the Jewish people.
living on the east coast, the east border,.
the land of Palestine,.
from time to time when they talk about the last day,.
the final battle, right?.
Some big battle would happen..
They would say the Parthians would be involved..
In fact, in first century,.
there were two major battles.
fought between the Romans and the Parthians..
Two times, the Parthians successfully challenged Rome..
So most probably,.
the first seal, white horse, rider with a bow,.
will remind the first audience, the first century audience,.
the Roman audience, their greatest enemy,.
most likely, are the Parthians..
Now, you will continue and ask and say,.
hey, I know a little bit about ancient history..
The fall of Rome didn't happen until AD 400 something..
Right, 476..
And the fall of Rome was not caused by the eastern Parthians..
It's actually caused by the northern barbarians..
It was the barbarians' job, right?.
Brought down the Roman Empire..

$^{441}$Does it mean Revelation is wrong?.
No, because the imagery is very flexible..
The imagery is, don't even talk about taking over Rome, right?.
The imagery, the first seal,.
the white horse rider with a bow,.
it simply means military pressure or military threat..
It's a common image that John used in the first century..
They are scared by the Parthians,.
the Romans are scared by the Parthians..
To use this image is to communicate a message,.
a message of judgment,.
a message to remind the Roman Empire,.
remind the first audience who are under persecution.
by the Roman Empire, that Roman,.
the Romans are not that powerful, right?.
That's the first seal, that's the message..
It's the same thing that we use the image of.
Osama Bin Laden to challenge the American power, right?.
Now Osama Bin Laden is already gone, right?.
But the message is there, right?.
To remind us the weakness, the vulnerability of the empire..
It's the same thing happening here in chapter six,.
first seal..
So the message is, in God's eyes,.
the empire oppressing the church will fail..
It will fall..
It looks very powerful, but it's not..
The first seal is talking about this..
So it talks about the evil empire,.
the evil force will fall one day..
Not because of the Parthians, but because of God himself..
So as you read the book of Revelation,.
you will find out who is really the one.
who brings down the evil empire, 1115..
Now it talks about, it already talk about the seven seals..
Now it talks about the seven trumpets,.
and later on talk about seven bows, okay?.
Now we're right in the middle, seven trumpets..
The seventh angel sang his trumpet,.
and there were loud voices in heaven which said,.

$^{481}$the kingdom of God has become the kingdom of our Lord.
and his Christ, and Christ will reign forever and ever..
Christ is the one who finally, ultimately,.
will bring down all this evil forces,.
and behind all this forces is Satan himself..
So as you can tell, the book of Revelation.
is not news report written ahead of time..
Now don't go to the book of Revelation and say,.
this is newspaper, today's newspaper,.
written 2,000 years ago..
No, that is wrong interpretation..
That is a misunderstanding..
The book of Revelation also is not a weather forecast..
It's not telling us, okay, tomorrow three o'clock p.m..
it will rain, it will snow, or something like that..
The book of Revelation does not intend to tell you.
when things happen, and how things happen exactly..
It doesn't do that..
Apocalyptic literature doesn't do this..
It doesn't answer our questions..
Our most interesting question is when and how, right?.
When does it happen?.
When will Christ come back?.
What will it look like, right?.
These are our questions, the when and how questions..
The book of Revelation is answering some other question,.
some more important question, which is the who question..
Who is in charge, and who is in control.
in the human history?.
When you see chaos out there,.
when you see everything is not right,.
who is the real Lord, who is in control?.
So it's not weather forecast,.
it's not newspaper within the head of time..
No..
Perhaps a better analogy is to think.
Revelation as a syllabus..
You all go to school, right?.
You receive a syllabus on the first day, right?.
It tells you the professor's name,.

$^{521}$the course name, classroom, and then the schedule, right?.
Okay, it's Monday, Wednesday, Friday,.
which assignment will be due,.
and what lessons will be taught,.
and how you prepare for that, right?.
And then when's the midterm, when's the final exam, right?.
And then, let's say, for example,.
you walk into your classroom one day, right?.
It was the day of midterm..
So you sit in the classroom,.
and then the professor came in,.
the teacher came in and gave you the midterm exam..
Okay, you got your midterm exam,.
and say, "Wow, prophecy came true..
"It's really the midterm.".
Would you say that?.
Of course not..
Of course not, because as you understand,.
this is a syllabus, right?.
It told you ahead of time,.
not in the sense of prediction,.
but in the sense that it's a plan, right?.
The plan the professor has already given you.
in the first day, right?.
And it also tells you who is the agency.
behind the syllabus, who is making things happen,.
who is fulfilling or accomplishing.
all those things printed on the syllabus, right?.
Who is the planner, the teacher?.
Who is the executor, the teacher himself,.
the professor himself?.
So you walk in on the day of midterm,.
and say, "That's right..
"That's how he planned..
"We are working out and going along with his plan.".
And the teacher has the authority.
to execute exactly everything according to the syllabus..
Human history is divine syllabus..
What we see in human history,.
what we see happening in our world,.

$^{561}$is the syllabus of God himself,.
the heavenly Lord, the King of kings..
He has written the syllabus..
That is the message, that is the theology.
of the book of Revelation..
God is in control of the syllabus..
God is in control of human history..
So no matter what happens in history,.
no matter what happens in your life and in my life,.
even though we see churches being persecuted in China,.
in some of the Muslim countries,.
churches, brothers and sisters being persecuted,.
you go back to history, you go back to syllabus,.
you know that behind it, there is someone in control,.
and God is still in control..
Everything that we see in your life, in my life,.
is according to the entire syllabus.
set out by God himself..
God will do all things according to his word,.
according to his syllabus..
That's why we can trust him..
That's why we should find hope in him..
Now that's what Revelation.
or Apocalypse literature is trying to do,.
is to tell us we can find hope only in him..
He is the one who is guarding and accomplishing,.
planning, executing the syllabus, right on his timing..
Even though we have questions we cannot answer,.
even though we see crisis,.
we try to find out the reasons behind it,.
sometimes we don't..
But look, everything is according to the syllabus of God..
Everything is according to his plan..
So we can trust him..
We find hope in God so that we can inform.
our own present living..
We will not be hijacked by fear, right, and unbelief..
Because we know everything will be according to his word..
He is the faithful God who will execute everything..
Even though we don't have all the detailed answers,.

$^{601}$but yes, we can trust in him..
I'm sure in the past week or so,.
you got a lot of news stories.
about what is going on in our world..
You probably heard a lot of news stories.
or some really sad reports from the Ukraine crisis..
I received a lot too..
I want to share one of those with you..
The president of Ukrainian Baptist Theological Seminary..
Crisis came and he said,.
"Now the church will go underground..
"We had that under the Soviet Union.
"and the church did not forget,.
"should never forget what it means to be persecuted..
"We, the church, will rearrange,.
"reorganize and do what we always do.".
Preach the gospel..
Isn't that true?.
Isn't that true if we find our hope,.
our ultimate anchor,.
nothing in the world but Christ himself.
who is in control of human history..
If he is in control, then this hope,.
this promise will inform how we react in times of crisis..
We will reorganize..
We will do what we always do.
or what we always called to do,.
to preach the gospel,.
to live a life of Christ follower..
Nothing will change..
Yes, there are times in darkness.
and sometimes we feel darkness is so overwhelming.
and the evil forces seem invincible, very powerful..
But the book of Revelation is reminding us.
all these evil forces are not omnipotent..
They are not all powerful..
They are limited..
God himself is omnipotent..
God himself is all powerful..
We still preach the gospel..

$^{641}$We still live for the kingdom..
Don't be hijacked by all this fear and uncertainties..
Your history is the syllabus of God..
Yes, we have a lot of questions without answers,.
but that's okay..
Midterm will come..
The final exam will come..
God will do according to his word..
Let me close with chapter five, verse five..
Then one of the elders said to me,.
"Do not weep..
"See the lion of the tribe of Judah,.
"the root of David, has triumphed..
"He is able to open the scroll and the seven seals.".
Let us all pray..
Father God, our Lord Jesus Christ,.
you are the one in control.
and you are the only king that is worthy of our praises..
(pensive music).
We face times of uncertainty and crisis.
and we're living in a time of chaos..
No matter what we face,.
remind us once again with the heavenly vision.
that you are the one who is in control..
That our own personal life, world history even,.
is in your syllabus, is in your plan..
Yes, we see evil forces all around us..
We see chaos and crisis..
We even see some of the people.
who are living less fortunate than us..
They are facing persecutions..
Our dear brothers and sisters in other places,.
they are less fortunate than us..
Yes, we see a lot of things that are not going right..
And it seems to us the evil is having the upper hand..
But remind us once again that Christ our Lord,.
the Lamb of God, who was slain on the cross.
and rose again from the tomb,.
you are the one who is in charge..
And we pray that our own personal life.

$^{681}$will dedicate ourselves again to you..
I pray that you will increase our trust and faith in you..
You will increase our love for you..
We still do what we are always called to do,.
is to live according to your will,.
to preach the gospel, to build up your church,.
to serve you with our gifts..
And pray that you will continue to renew us,.
renew our perspective..
To live faithfully in your sight,.
to live faithfully according to your word,.
to hear your word and to apply it into our lives..
May our lives be pleasing to you..
In Christ name we pray, amen..
\newpage



\section{}
\label{sec:X5Jubq7tdoc}
\textbf{Family Sunday School - EP 2 (The Last Supper)}
\newline
\newline
連結: \href{https://youtube.com/watch?v=X5Jubq7tdoc}{\texttt{https://youtube.com/watch?v=X5Jubq7tdoc}} ~~~~ 語音日期: 2022-03-03
\newline
\newline
\hyperref[sec:1enAW_IL3Aw]{\small{< < < PREV SERMON < < <}}
~
\hyperref[sec:index_chronic]{\small{[返順時目]}}
~
\hyperref[sec:index_scriptual]{\small{[返順卷目]}}
~
\hyperref[sec:FcE9snFLM6s]{\small{> > > NEXT SERMON > > >}}
\newline
\newline
$^{1}$(upbeat music).
(speaking in foreign language).
Pharisees..
(speaking in foreign language).
Religious leaders..
(speaking in foreign language).
Hello..
(speaking in foreign language).
Okay?.
(speaking in foreign language).
, right?.
(speaking in foreign language).
(laughing).
(speaking in foreign language).
Yay, got it!.
(speaking in foreign language).
Okay..
(speaking in foreign language).
'Cause it's bubbles on it..
(speaking in foreign language).
This should be a naan..
Is it naan?.
Yay!.
(speaking in foreign language).
I got it..
(speaking in foreign language).
Nice job, assistant..
Next one..
(speaking in foreign language).
Seven grain, 21 grain..
(speaking in foreign language).
It is..
The whole wheat grain bread..
(speaking in foreign language).
Yes, I got it..
(laughing).
Next one..
(speaking in foreign language).
What is this?.
It feels slimy..

$^{41}$Oh no, it's slime bread..
(speaking in foreign language).
It must be Chinese..
(speaking in foreign language).
But it feels like slime..
Hold on..
(speaking in foreign language).
Tastes like a ball of bao..
But it doesn't have the ball of bao crust..
But it's slimy in the middle..
Don't worry, I'm the only one eating this bao..
It's slime in the middle..
What kind of bao is it?.
Okay, I have to take a look, okay?.
I'm gonna take a look..
I forfeit this round, I take a look, hold on..
Oh, I got it!.
It's a (speaking in foreign language).
So usually, when I go to the bakery,.
(speaking in foreign language).
It's a little bit oblong..
(speaking in foreign language).
What threw me off, it's circular..
(speaking in foreign language).
Tricky, tricky assistant..
(speaking in foreign language).
Or cocktail bun..
(speaking in foreign language).
Woo-hoo!.
This is good..
(speaking in foreign language).
Okay..
(speaking in foreign language).
Is actually midnight snack or supper..
(speaking in foreign language).
It's happening real time..
(speaking in foreign language).
Last supper..
(speaking in foreign language).
So start with the last supper..

$^{81}$(speaking in foreign language).
Now imagine..
(speaking in foreign language).
Is that it?.
That's it..
Bread..
Just bread..
(speaking in foreign language).
Bread and wine..
That's how we're starting today's story..
(speaking in foreign language).
The last supper..
(speaking in foreign language).
Passover..
(speaking in foreign language).
Last supper..
(speaking in foreign language).
The last supper..
(speaking in foreign language).
Okay, time for supper..
(speaking in foreign language).
Right?.
(speaking in foreign language).
the Bible..
(speaking in foreign language).
No way!.
(speaking in foreign language).
Really good friend, right?.
Really good friend..
Let's just see what happens, right?.
(speaking in foreign language).
He's God's son..
(speaking in foreign language).
He probably knows..
Actually, not probably..
He definitely knows..
(speaking in foreign language).
Right?.
(speaking in foreign language).
Basically, it's a group of CFO, basically..

$^{121}$(speaking in foreign language).
(sighing).
(speaking in foreign language).
(sighing).
(speaking in foreign language).
By Judas law..
(speaking in foreign language).
Tricky, tricky..
(speaking in foreign language).
Midnight snack..
(speaking in foreign language).
Next time, we got Peter and Judas..
(speaking in foreign language).
Before..
(speaking in foreign language).
His purpose..
(speaking in foreign language).
Prophet..
(speaking in foreign language).
None of us..
(speaking in foreign language).
So that's why we're starting here..
(speaking in foreign language).
Little side fact..
(speaking in foreign language).
Bread of life..
(speaking in foreign language).
Deep, right?.
So, we'll unpack it as we continue..
(speaking in foreign language).
Time at home..
(speaking in foreign language).
Name that bao..
Okay?.
(speaking in foreign language).
Who can name the most baos, right?.
(speaking in foreign language).
Right?.
(speaking in foreign language).
Me against you..

$^{161}$Miss Rebecca..
(speaking in foreign language).
From all the baos you can think of..
(speaking in foreign language).
Okay, ready?.
(speaking in foreign language).
Okay?.
(speaking in foreign language).
Bao..
(speaking in foreign language).
Okay..
(speaking in foreign language).
Bao..
(speaking in foreign language).
It's really a bao, just doesn't have the bao in it..
(speaking in foreign language).
(laughing).
(speaking in foreign language).
Mediterranean food, you know, pita..
(speaking in foreign language).
Is that bao?.
That's why (speaking in foreign language).
Hey kids, Chinese time..
Chinese time..
(speaking in foreign language).
I know..
(speaking in foreign language).
But it's not easy..
(speaking in foreign language).
But bao, like (speaking in foreign language).
like the hamburger, it's kind of like a sandwich..
(speaking in foreign language).
As you can hear, sandwich, (speaking in foreign language).
phonetically, (speaking in foreign language).
Sandy witch..
So, (speaking in foreign language).
That's why it's called a bao, right?.
Like (speaking in foreign language).
isn't actually a sandwich bao.
or like a (speaking in foreign language).

$^{201}$It's like a dumpling..
But (speaking in foreign language).
Wrapped around..
(speaking in foreign language).
Wrapped present..
(speaking in foreign language).
Right?.
Hello (speaking in foreign language).
Thank you so much..
Bye..
(upbeat music).
\newpage



\section{啟示錄 6:1-2}
\label{sec:FcE9snFLM6s}
\textbf{亂世啟示 (1) 第一印 白馬  (啟示錄6\_1-2) - 陳崇基牧師}
\newline
\newline
連結: \href{https://youtube.com/watch?v=FcE9snFLM6s}{\texttt{https://youtube.com/watch?v=FcE9snFLM6s}} ~~~~ 語音日期: 2022-03-07
\newline
\newline
\hyperref[sec:X5Jubq7tdoc]{\small{< < < PREV SERMON < < <}}
~
\hyperref[sec:index_chronic]{\small{[返順時目]}}
~
\hyperref[sec:index_scriptual]{\small{[返順卷目]}}
~
\hyperref[sec:jFHDnibTFuo]{\small{> > > NEXT SERMON > > >}}
\newline
\newline
啟示錄 6:1-2
\newline
\begin{longtable}{cl}
\hline
\hline
章節 & 經文 (和合本修訂版)\\
\hline
6:1 & \begin{tabularx}{0.7\textwidth}{X} 我看見羔羊揭開七個印中第一個印的時候，聽見四活物中的一個活物，聲音如雷，說：「你來！」 \end{tabularx} \\ \\ \relax
6:2 & \begin{tabularx}{0.7\textwidth}{X} 我就觀看，看見一匹白馬，騎在馬上的拿著弓，並有冠冕賜給他。他出來征服，勝而又勝。 \end{tabularx} \\ \\
[1ex]
\hline
\hline
\end{longtable}
$^{1}$在這個主日開始.
我們都要講一講另外一些的主題.
我們這些做傳道牧師的.
有時候都要按照世界發生什麼事.
然後我們要調整一下.
我們想講些什麼.
因為疫情已經有兩年多了.
現在在各地的地方.
剛才我們也提過.
香港疫情也是非常差.
我們這裡稍微放緩一下.
但是也有很多的未知之數.
上個星期又是戰火連天的時候.
俄羅斯入侵烏克蘭.
到底在這些不安的時候.
我們會想起些什麼呢.
很多時候我們會想起聖經裡面的一些說話.
民要攻打民.
國要攻打國.
或者有時候一些經文我們會想起的.
可能就是啟示錄.
因為在這些兵荒馬亂的時候.
我們就想一下.
到底耶穌什麼時候回來會收拾殘局呢.
希望祂快點回來就好了.
所以我們就在這裡加插這個系列.
我們想去思考一下.
到底啟示錄應該是怎麼講的.
特別是在一些不安的時候.
讓我們去看看到底怎麼去解釋.
特別很多時候我們看到啟示錄的時候.
我們環教會有個傳統.
就是我們很喜歡亂解釋的.
所以我們想思考一下.
到底怎麼去理解裡面所講的東西.
既然聖經講的東西應該是造就我們.
應該是有message for我們的.
到底我們應該怎麼去理解裡面的一些經文呢.
所以我們從今天開始.
我們看看啟示錄裡面的一些經文.

$^{41}$今天我會看第六章.
第一節到第二節.
這裡所謂打開七印的時候.
第一印就有一隻白馬出來.
到底這些類似的經文應該怎麼去解釋呢.
首先第一開始的時候.
我們先不看第六章 先看第一章.
這裡講到耶穌基督的啟示.
就是神賜給他的.
就是要把快要發生的事.
指示他的眾僕人.
耶穌就差派他的天使指示他的僕人約翰.
約翰就將上帝的道 耶穌的見證.
凡是自己所看見的都見證出來.
讀者書上預言的那些聽見者預言.
並且遵守書上所記載的.
都是有福的 因為事後近了.
這裡第一個字 啟示錄.
第一個字就是apocalypse.
這個字很熟悉.
因為我英文apocalypse這個字就是這裡來的.
我記得兩年之前.
我們剛剛疫情開始爆發.
我們有lockdown 有宵禁的政策.
我就和兩個小朋友去了餐館.
我們買了一些take out 外賣回家.
我們開車出去的時候.
發現整個城市好像一個ghost town.
鬼影都沒有.
芝加哥很少的 經常都塞車.
我發現 你記得疫情剛剛爆發之後.
不塞車的 你開多快都可以.
這個就是apocalypse.
因為今天我們說到apocalypse這個字的時候.
我們就是說doomsday 或者是世界末日的災難性的事.
但是這個字本身在聖經的時代.
apocalypse apocalypsis.
這個字其實是uncover的意思.
就是說啟示或者揭開的意思.
裡面有些什麼 它將蓋打開.

$^{81}$然後揭曉一些東西.
那些什麼呢 它說快要發生的事.
要指示他的眾僕人 於是就派天使去.
指示他一些將來一定會發生.
快要隨時會發生的事.
指示了他的僕人 並且去到吩咐.
在西亞西 或者今天的土耳其.
七個教會 那些教會名我們不說了.
你記得我們說啟示錄的時候.
經常都會說這七封書信 因為很重要.
因為啟示錄不是說將來會發生什麼事.
我們有什麼預測等等.
主要是要警告教會 因為當時教會.
是落在一些試煉 試探的當中.
當時羅馬帝國來迫害教會.
或者使教會有很多的壓力等等.
以至教會在水深火熱的當中.
所以他提到約翰 即是使徒約翰.
他和教會一起去患難.
在國度的忍耐當中.
是為了上帝的道 這位約翰.
曾經被困於或者被逼到.
拔摩的海島 土耳其海岸的外面一個很小的島.
當時被帝國的異見分子監禁在那裡.
約翰就是其中一個受害的.
你會看到這個啟示錄.
到底寫來是為了什麼的緣故.
就是為了當時的教會.
他們得到勸勉和得到提醒.
所以首先我們搞清楚.
我們打開啟示錄的時候.
我們不要將啟示錄當為.
一些預先寫好的新聞報道.
不是預先寫好的報紙.
啟示錄是一個題材.
屬於一種文學的題材.
叫做Apocalyptic Literature.
叫做天啟文學.
我們一講到天啟文學.
我們會覺得很神秘的東西.

$^{121}$但其實聖經裡面有很多這樣的文學.
譬如大義理書 舊約裡面.
以西結書等等的章節.
大義理書 舊約裡面的福音書.
某一些的段落.
保羅書信某一些的段落.
啟示錄當然就是全本都是.
所謂的Apocalyptic Literature.
到底什麼是Apocalyptic Literature.
它是用一些符號.
用一些天上的符號.
或者不是地上的一些語言.
來啟示出一些天上.
另一個境界 將來會發生的現實.
來揭曉出來.
以致讀者聽見這些訊息的時候.
他就知道怎樣去詮釋現在的處境.
現在怎樣兵荒馬亂 水深火熱都好.
我們怎樣可以在啟示文學裡面.
能夠見到上高發生什麼事.
天上發生什麼事.
另一個境界 屬於一個境界.
將來發生什麼事.
上帝的子民 教會.
受到迫害的時候 信心受到挑戰.
到底怎樣拿到這些訊息.
以致今天我們知道怎樣去詮釋.
所以這些天啟的文學.
不是叫我們抽離 離開這個世界.
而是叫我們聽完這些訊息的時候.
我們能夠參與在這個世界裡面.
知道上帝的工作仍然繼續.
所以首先搞清楚.
天啟文學不是什麼.
天啟文學是寫給當時的教會.
去鼓勵他們.
所以我們看第一至第三章的時候.
我們講過了 就是七教會的書信.
然後到第四章.
剛才我們崇拜開始的時候.

$^{161}$引用過一些經文.
聖哉 聖哉 聖哉等等.
就是上帝被提到天上.
約翰被提到天上.
看到天上的景象.
好像一個羅馬帝國的宮廷.
但他看著看著的時候.
他發現坐在宮廷裡面.
寶座上居然不是羅馬皇帝.
而是上帝.
昔在今在 以後永在的神.
我們大致上就知道.
到底啟示錄是一本怎樣的書.
在第一世紀末期的時候.
90年代的時候.
AD90年左右.
就是羅馬帝國在不同的情況下.
不同的程度裡面.
軟硬兼施.
來打壓當時的教會.
啟示錄就是寫給他們.
來堅定他們的信心.
所以頭四章我們不講太多.
因為太多人講頭四章.
並且很少會解錯.
但很多時候解錯就是後面那段.
所以我們繼續看下去第五章.
他在這裡看到不同的景象.
這些就是天啟文學很有特色的一些語言.
他說我看見那坐在寶座上的右手拿著書卷.
這書卷兩邊都是滿有字的.
有七個印封上.
我又看到一個大力的天使大聲的宣告.
誰配展開這書卷.
拆開它的封印.
在天上地上地底下.
因為沒有人配得去展開這書卷.
所以約翰就大哭.
以前的文件很多時候寫一些羊皮或者紙皮.
寫了之後他們就捲起它.

$^{201}$然後就會加上一塊泥.
然後就用一個戒指或者一個風刃.
在泥上面去按一下.
將這些文件去到封存之後.
就會放在聖殿或者一些皇宮的宮殿的儲藏室裡面.
就是archive了它.
放在當時的辦公室.
如果你要看回這本書裡面的東西.
看回書卷就要拆開這些印.
誰可以拆開印呢.
就是那些印很多時候有個名或者有個符號.
是代表那個人的.
等於我們上一代有圖章.
就要拿回那個圖章去鑒定.
然後才可以拆開這些印.
但是這本書書卷裡面所寫的東西.
如果你看回舊約的背景.
就是講到一些words of judgement.
就是對地上的惡霸或者一些邪惡的勢力.
去做一些審判.
這些是一些話語來的.
一些先知的說話.
所以為甚麼約翰很想看這本書裡面所寫的東西.
因為當時他們在水深火熱的當中.
被羅馬帝國打擊教會.
到底明天會怎樣.
有很多未知之數.
很多uncertainty.
看不到這本書裡面所講的東西.
那是不是我們要認命呢.
就是說這個帝國真的太大了.
他的勢力太厲害了.
我們的命運是不是就是放在帝國的施壓者的手裡.
所以為甚麼他大哭.
因為到底將來會怎樣我們不知道.
或者可能我們現在面對不同的情況.
不同的困境.
我們都可能會想一下.
到底將來會怎樣.
我們是不是就繼續受苦.

$^{241}$或者在不安的威脅當中.
所以約翰大哭.
第五章第五節.
其中有一位長老.
他就對約翰說.
他說不要哭.
看了猶大之派的師子.
大衛的根.
他已經得勝了.
他可以展開那一書卷.
揭開那七印.
所以不需要懼怕.
他說因為這個高揚.
耶穌基督.
被殺的而且復活的.
猶大之派大衛的君主.
他有權柄有authority.
能夠打開這七印.
所以我們今天就來到第六章.
將這些七印一路一路打開的時候.
你就可以想像.
打開第一印.
然後書卷打開一點.
見到裡面的內容.
然後再打開第二印.
再看裡面的內容.
然後逐個逐個打開的時候.
這個就是所謂我們經常聽到啟示錄的七印.
就是上帝要帶來這個打壓教會.
這個羅馬帝國.
將來會有什麼審判呢.
第一印打開.
今天我們的主題.
第六章第一節.
楊高打開揭開第一印的時候.
就見到四個回物當中其中一個.
發出很大的雷轟的聲音.
他說你來.
我見有一匹白馬.
騎在馬上的是拿著弓.

$^{281}$有官面賜給他.
他就出去.
得勝並且會再得勝.
這些symbolic language.
我們很多時候現代人就會很容易誤解.
因為我們聽了之後.
我們就會說今天到底是怎樣.
我們就將現代的意思.
我們就很容易塞進去.
因為你知道這些symbol.
這些語言.
其實很有彈性.
幾乎你塞什麼進去都可以.
但這裡說到白馬.
馬上的Rider.
騎馬的有弓.
得勝又得勝.
這裡很明顯.
就是在說有戰爭.
一定是有戰火.
所以就有這樣的描述.
但如果我們將現代人的意思.
我們塞進去.
就很容易.
首先因為有白馬.
有弓.
我示範一下如何亂解.
示範一下如何亂解.
很容易.
我找一個現代的位置.
俄羅斯的國旗有白色.
如果你Google一下.
Putin.
White Horse.
你就會找到這些相片.
因為他真的騎過白馬.
我們就會說.
嘩 啟示得多厲害.
2000年之前已經預測了.
今天俄羅斯會入侵Ukraine.

$^{321}$所以聖經是很準的.
多厲害.
但你留意到.
那個方法是完全錯的.
為什麼呢.
因為我們經常說.
聖經不是直接寫給我們.
對嗎.
聖經是直接寫給誰.
是當時水深火熱的教會.
第一世紀末的時候.
他們受打壓的時候.
到底他們會想起什麼.
他們見到白馬.
和聽到這些符號的時候.
他們第一樣會想到什麼.
就不是我們想到什麼.
然後我們塞進去.
所以我們解釋聖經.
特別解釋啟示錄的時候.
我們就先不要用我們現在.
我們怎樣了解.
然後把意思塞進去.
如果不是的話.
我們就很容易錯.
你不一定要放Russia.
你放NATO進去也可以.
因為你搜索一下NATO的標誌.
這個都是另一個帝國的勢力.
NATO北約.
都是白色的.
都是白色.
有一個Compass.
有一個指南針.
用現代人塞進我們想到的東西.
進去聖經.
應該想一下.
當代的人.
他們會想起什麼.
在這裡給你看看.

$^{361}$一些錢幣.
當時第一世紀的時候.
前後.
一些錢幣.
是一個Parthian的Empire.
一個羅馬帝國東面.
一個帝國.
這個源自摩斯地的一個帝國.
帕提亞帝國.
這些錢幣全部都是Document.
都是Text.
你去解釋這些錢幣.
看這些錢幣的時候.
你就知道他們想說什麼的訊息.
這些士兵.
這些軍王是拿著弓的.
三角形的一個射箭的弓.
軍王騎著馬.
Parthian.
在當時羅馬帝國的時代.
最出名的就是弓箭手.
Archers.
而且Parthian的Empire.
他們當時最神聖.
最Sacred的顏色.
就是白色.
每一隊的軍隊.
都是有騎弓箭手的.
所以如果你這樣看的時候.
你就發現.
當時他們想起的.
就不會是我們現在新聞裡面看到的東西.
想起白馬.
想起弓箭.
想起軍事的威脅.
他們就會想起羅馬帝國東邊.
這個Parthian的Empire.
大概有四百年的歷史.
事實上當時的羅馬作家.
猶太人的寫作.

$^{401}$他們說到末日大戰的時候.
他們都會說到.
或者會比喻到.
好像這些Parthians會來入侵羅馬帝國.
你知道第一世紀末.
如果你知道一點點歷史的話.
羅馬帝國和Parthians.
是經常有一些軍事上的衝突.
事實上有兩次的大戰.
兩次的大戰.
羅馬帝國都是大敗.
所以我們想起.
咦 他們說起Parthian的時候.
到底是什麼意思呢.
他們想起當時東邊的帝國.
因為經常會犯警.
會威脅羅馬政府.
如果你這樣看的時候.
你就不會這麼容易.
將我們現代人想到的東西.
幻想到的東西.
就扼落聖經.
當代的人會怎麼想.
還有一件事.
就是啟示錄.
天啟文雄.
是不是一種預測.
是不是就是天氣報告預測.
就不是.
如果你熟悉一點點.
羅馬帝國的歷史的話.
你知道羅馬帝國淪陷的時候.
是主要大概四百多年.
並且不是Parthians.
來攻佔羅馬帝國.
而是其他Barbarians.
那你說.
那不就是聖經錯了.
啟示錄不是這樣說的嗎.
那隻白馬是東邊來的.

$^{441}$但你留意一下.
其實不是在做預測.
只不過是在用一些當時.
羅馬帝國感覺到.
是一些威脅.
一些這樣的Image.
來去揭示一件事.
就是揭示羅馬帝國.
雖然外表很強.
但其實他是很vulnerable.
是很脆弱的.
等於今天你去說一些美國.
美國的帝國.
最怕的是什麼.
是什麼Image.
你拿Osama Bin Laden出來.
雖然他已經過去了.
但你一提到這些東西.
恐襲的時候.
就會對美國的人心.
你就會揭示到.
可能大家都不是那麼強.
可能裡面是很脆弱.
所以這個第一印第二印第三印.
這些這樣的language.
不是一些預測的報告.
其實目的想做什麼呢.
就是告訴我們.
迫害教會的羅馬帝國.
始終是會覆亡的.
雖然現在教會好像是在下風.
但是啟示要揭示的.
就是這些帝國的勢力.
黑暗的勢力.
背後是撒旦的勢力.
不會是德性教會.
這些這樣的勢力.
一定會傾覆.
就不是因為Parthians.
如果你繼續看下去的時候.

$^{481}$你就知道十一章第十五節.
是什麼勢力會令到帝國倒退.
會失敗.
十一章第十五節.
印講完的.
就講到七號.
你知道接著就有七碗.
七號的時候.
第七位天使吹過.
天上有大的聲音說.
世上的每一個國度.
成為了我們的主.
和我們所立的基督的國.
基督要作王.
直到永遠.
主呀全能的神.
今在昔在.
因為你執掌了大權.
作了王.
作了萬國的王.
因為當上帝去審判列國的時候.
基督就成為萬國的王.
所以這些天啟的文學.
主要不是天啟報告.
不是預先寫好的報紙.
因為所有答的問題.
不是when and how.
我們最concern的就是.
when and how.
這場仗什麼時候完.
疫症什麼時候完.
耶穌什麼時候.
什麼月什麼日回來.
這就是when的question.
怎樣回來.
有什麼徵兆.
how will it happen.
這些都是我們感到興趣的問題.
但天啟文學要答的問題.
不是when and how.

$^{521}$是who.
到底誰執掌王權.
誰統領你的明天.
誰統領你的將來.
誰掌管幕後的命運.
是基督.
是who.
或者我覺得有一個更適合的比喻.
每次我們打開聖經.
看天啟文學.
先知書等等.
我們不要當它是天啟報告.
應該當它是什麼.
我想應該當它是syllabus.
大家有讀過書.
上學的時候第一天.
老師就會給我們一個課程的大綱.
syllabus.
講得很清楚.
1月2月3月.
什麼時候什麼月什麼日.
有什麼deal.
中期考試是什麼時候.
老師就會說.
會跟著這個大綱去行事.
所以大家知道什麼時候要交功課.
中期的時候.
你就是一天.
你回到班房.
就是中期考試那天.
你回到班房的時候.
你就發現.
真的有考試.
你會不會跟你老師說.
你真是靈.
原來你講到今天.
什麼時候什麼月什麼日.
你預測的完全是準確.
你真是神人.
受我一拜.

$^{561}$你會不會這樣做.
你當然不會這樣做.
因為教授的大綱.
一方面是他的計劃.
但同樣是由他全權去執行的.
是他去行事的.
你中期考試那天回到班房.
大考那天回到班房.
真的有大考.
是應該的.
因為你知道誰在負責.
你看回那個大綱的時候.
你就知道原來沒有走過.
走歪過.
完全是教授和老師去到成就一起.
啟示錄就是想講這件事.
先知的文學.
天啟的文學就是想講這件事.
人類歷史的大綱.
那個大綱.
就是上帝所計劃的大綱.
是上帝的教學大綱.
就是誰有權力.
誰有權柄去到計劃.
並且去承傳這個詩詞所講的東西.
承傳歷史的就是歷史的主宰.
就是上帝.
就是被殺的羔羊.
復活的羔羊.
這個才是啟示錄.
讓我們世世代代.
基督徒 教會.
縱使在水深火熱的當中.
我們都能夠明白.
上帝仍然是掌權.
Who is in charge.
不是How and when.
Who is in charge.
你看到很多Uncertainties.
有很多不安的時候.

$^{601}$到底我們的Attention.
我們被騎劫的一些問題.
是什麼問題.
往往就是When and How的問題.
什麼時候我可以解困.
怎樣可以解困.
上帝你可不可以出手.
做這件事 做那件事.
這些都是When and How.
但是我們就Miss the whole point.
我們就沒有問Who.
到底誰掌管明天.
誰掌管你一切.
如果他是掌管你一切的話.
你就信教.
信得過他.
上個星期俄羅斯入侵烏克蘭.
我相信大家收到很多不同的消息.
很多不同的報道.
其實烏克蘭70\%都是基督教徒.
不一定是Prophets.
不一定是新教.
東正教等等.
侵略之後.
我相信大家收到很多不同的消息.
我自己都收到很多.
我們自己網絡裡面的消息.
其中一個浸會神學院的院長.
他說了什麼.
他說在烏克蘭的浸會神學院的院長.
他說教會現在我們會Go underground.
我們在蘇共的時候.
我們曾經都是這樣.
但教會從來沒有忘記過.
受迫害的時候.
到底應該怎樣生活.
他說我們會重新regroup reorganize.
重新再重組.
這個是我們教會一直會做的事.
我們並且會繼續傳福音.

$^{641}$這個就是啟示錄的神學.
Who is in charge.
時局更加惡劣的時候.
惡勢力好像很猖狂的時候.
好像很紛亂的時候.
或者我們很多時候會感覺到很無助.
帝國和帝國背後的撒旦.
啟示錄要跟我們不斷講清楚.
就是祂不是全能者.
祂不是all powerful.
因為耶穌基督.
死而復活的主.
祂才是無所不能.
祂才是全能者.
耶穌基督這位羔羊.
才是得勝者.
已經將撒旦將軍了.
祂已經沒有棋子.
有的只是掙扎的空間.
主仍然是作王.
因為上帝是編寫歷史大綱的神.
你的歷史大綱.
世界的歷史大綱.
並且成全者.
都是上帝自己.
啟示錄就是想告訴我們.
這件事.
聽清楚.
裡面對我們今時今日.
有什麼意義.
叫我們不要被恐慌去騎劫.
叫我們知道.
能夠打開七印的.
是羔羊.
是耶穌基督.
五章第五節.
我們用這一節經文結束.
長老中有一位跟我說.
他說不要哭.
看猶太之派的師子大衛的根.

$^{681}$他而得勝.
能以展開那書卷.
揭開那七印.
求主幫助我們.
讓我們的生命再一次.
去到神的面前.
知道哪一個是在任.
我們低頭祈禱.
上帝我們多謝你.
因為我們的主耶穌基督.
被殺的羔羊.
為我們犧牲的羔羊.
而在死裡復活的得勝的主.
祂是色在金在以後永在的神.
主我們來到敬拜祂讚美你.
因為你是掌管一切的神.
主在我們生命當中.
有很多不安有很多混亂.
在世界政局裡面.
有很多我們看不透.
不明白發生什麼事.
也問很多為什麼.
問很多為什麼.
主縱使這些答案我們未能夠有.
但我們知道你仍然是掌管一切.
求上帝你自己親自去幫助我們.
我們面前可能遇見有很多.
的未知不察.
或者遇見有很多生命裡面的很大的危機.
甚至我們紀念到世界各地.
包括中國大陸.
教會我們主內的弟兄姊妹是受打壓的.
自由受到威脅的.
來到烏克蘭.
在戰火裡面.
受困的這些平民百姓.
主我們有很多的為什麼.
但我們知道你仍然是那位.
編寫歷史大綱.
承傳歷史大綱的神.

$^{721}$求主你讓我們知道在你裡面我們能夠有安穩.
無論我們的生命會去到哪裡.
我們知道你仍然是作王的神.
求主你幫助我們來跟隨敬拜你.
在地上我們仍有我們的使命.
就是去依靠你掌管一切的主.
我們在地上我們有我們福音的使命.
求你幫助我們讓我們的生命.
從我們的崗位開始.
我們出發.
能夠成為你的使者.
我們這樣祈禱奉主名禱.
我們播一首回應的詩歌.
我們聽完之後.
我們會做一個很短的禱告.
當我們聽回應詩歌的時候.
請你都想一下.
到底在你面前遇見什麼.
上帝在啟示錄.
在經文怎麼去提醒你.
\newpage



\section{啟示錄 6:3-4}
\label{sec:jFHDnibTFuo}
\textbf{亂世啟示 (2) 第二印 紅馬 (啟示錄6\_3-4) - 陳崇基牧師}
\newline
\newline
連結: \href{https://youtube.com/watch?v=jFHDnibTFuo}{\texttt{https://youtube.com/watch?v=jFHDnibTFuo}} ~~~~ 語音日期: 2022-03-07
\newline
\newline
\hyperref[sec:FcE9snFLM6s]{\small{< < < PREV SERMON < < <}}
~
\hyperref[sec:index_chronic]{\small{[返順時目]}}
~
\hyperref[sec:index_scriptual]{\small{[返順卷目]}}
~
\hyperref[sec:shYH8FbN1Gk]{\small{> > > NEXT SERMON > > >}}
\newline
\newline
啟示錄 6:3-4
\newline
\begin{longtable}{cl}
\hline
\hline
章節 & 經文 (和合本修訂版)\\
\hline
6:3 & \begin{tabularx}{0.7\textwidth}{X} 羔羊揭開第二個印的時候，我聽見第二個活物說：「你來！」 \end{tabularx} \\ \\ \relax
6:4 & \begin{tabularx}{0.7\textwidth}{X} 就另有一匹馬出來，是紅色的；有權柄賜給了那騎馬的，要從地上奪去太平，使人彼此相殺；他又接受了一把大刀。 \end{tabularx} \\ \\
[1ex]
\hline
\hline
\end{longtable}
$^{1}$我們繼續看啟示錄.
上一次我們開始了第一次.
如果你有聖經的話.
翻開這本書.
很容易找到.
聖經裡面最後最後一本書.
啟示錄.
這次我們看第二個人.
因為我們在看七人的啟示.
我們想這樣做.
啟示錄不只是說裡面的內容.
而是我希望分享一下.
到底應該怎樣去接觸.
怎樣去看這本書.
因為很多時候我們覺得啟示錄.
是一些很神怪的東西.
或者一些我們看不慣的語言.
所以我都會跟大家解釋一下.
到底我們現代人應該怎樣去了解這本書.
以致我們不會.
上個星期我都示範了.
怎樣亂解啟示錄.
很容易解錯的.
因為任何東西都可以加進去.
所以我們需要有一些方法去了解.
到底上帝藉著這本書.
怎樣跟我們說話.
首先我們第一件事.
我們上個星期提過的.
就是啟示錄是屬於一種文學的種類.
叫做Apocalyptic Literature.
什麼叫Apocalyptic Literature.
過去幾十年.
其實有很多人嘗試去總結.
到底什麼叫做天啟文學.
有幾樣東西.
大概有五樣東西.
我們可以總結出來.
因為啟示錄只是其中一個.
天啟文學裡面其中一卷書.

$^{41}$有很多經內的.
或者經外的文獻.
都可以叫做Apocalyptic Literature.
首先什麼叫Apocalyptic Literature.
什麼叫天啟文學.
當我們看見啟示錄裡面.
一些很神怪的異象的時候.
第一樣東西當然就是Revelatory.
就是說是啟示性的.
就是說是上帝由上至下.
來啟示自己.
啟示末後會發生的事.
將來的事情來打開.
這個就叫做Revelation.
上帝親自啟示我們才會知道.
另外很多時候天啟文學.
都會有一個Narrative.
一個Story.
上一次我提過一個很熟悉的樣本.
一個Example.
就是《彈耳曆書》.
舊約裡面這本天啟文學.
《彈耳曆書》.
就是有一個Narrative.
我們從小到大在教會.
聽見一些故事的時候.
我們都聽見有一個框架.
有一個旅程.
啟示錄都是一樣.
史羅若汗有一個旅程.
到最後的時候告訴我們.
我們要敬拜的是誰.
然後第三樣東西就是.
有一個Angelic Being.
去到Media.
很多時候天啟文學.
就是講完一些異象.
不太明白.
然後有一個天使.
或者有一個使者.

$^{81}$來作一些講解.
無論是先知還是使徒.
第四樣就是Transcendent Reality.
什麼意思呢.
就是說將將來和將天上真正的現實.
那個現實不是我們現在.
Here and Now.
我們捉到摸到的.
這個會過去的.
而是天上From Heaven.
和將來From the Future.
那個Reality.
那個真正的Reality.
上帝是藉著啟示錄.
或者一些天啟文學.
來彰顯出來.
叫我們知道.
原來我們所追求的現實.
不是只是Here and Now.
不是只是在這裡.
他表述給我們看.
真正我們要去追求的東西.
不是地上的事物.
而是天上的事物.
將來的事物.
這些很多時候都是用一些.
一些符號的方法.
一些象徵性的方法.
來向我們展現出來.
最後一樣當然就是處境.
就是這些天啟文學.
不是預先寫好的報紙.
是有它當時的處境的.
啟示錄就是當時第一世紀的時候.
水深火熱的教會.
他們面對著不同的挑戰.
不同的試探的時候.
啟示錄就是要來去鼓勵他們.
還有很多時候是去挑戰.
甚至去指責他們.

$^{121}$譴責當時的教會.
他們因為對主的心已經忘記了.
以致他們能夠悔改.
回到主的面前.
當時第一世紀的教會.
如果你熟悉一點的背景的話.
就應該用這樣的背景來看啟示錄.
你看的時候就會明白.
為什麼啟示錄要講這些東西.
當時第一世紀當然有不同的羅馬帝王.
最出名的兩個.
是特別和打壓教會有關的.
第一個就是六十年代的尼祿.
第二個就是九十年代的時候.
一個叫做多米田.
無論你怎麼翻譯他也好.
還有這兩個的王.
對初期教會來說.
是給當時信仰的群體一個很大的衝擊.
很嚴重的壓力.
所以這些啟示文學.
其實就是想去到鼓勵教會.
安慰教會.
還有要挑戰教會.
如果這樣看的時候.
這本書或者會對我們今時今日的信息.
才能夠消化出來.
第二樣東西就是啟示錄的結構.
我們看這些故事的時候.
或者啟示錄裡面一些經文的時候.
很多時候我們就是.
take it out of context.
或者將上文下理.
我們拿出來.
然後我們就猜誰是666.
誰是那個獸.
或者那個獸的印記是什麼.
等等等等.
這些我們將來希望能夠去處理.
但是我們如果將它take it out of context的時候.

$^{161}$我們就不明白裡面發生什麼事.
譬如666這個經文.
到底在哪裡出現的.
就是在中間那個七記號那裡出現的.
十二至十四章.
如果你這樣分結構的時候.
就發現很多的七字.
七個的段落.
七乘七的段落.
開頭頭三章就是寫信給七個教會.
然後就是四至七章就打開七印.
將來發生的事.
末後將會發生的事.
上帝藉著耶穌基督.
祂的六生子再一次回來的時候.
second coming的時候會發生什麼事.
這個就是seven seals.
就是我們這兩個星期在講的這段經文.
然後就是七號.
seven trumpets.
八至十一章.
然後到中間就是十二至十四章.
那裡有七個記號.
或者七個的意象.
七個的visions.
那裡經常說.
and I saw.
有六字.
加上第一次就有七段.
就是看到一些什麼天上的景象.
原來上面有一個.
天上的一個cosmic battle.
是在打的.
原來下面看到.
我們here and now我們看到的struggle.
原來有一個天上的現實.
一個reality打開了.
這個就是十二至十四章所講.
然後講到上帝對帝國的審判.
十五至十六章.

$^{201}$就是七碗的意象.
然後十七至十八章.
就是講到巴比倫的沒落.
這個巴比倫就是代表羅馬.
羅馬帝國.
因為當時羅馬帝國是給教會很多壓力.
到最後十九至二十二章.
就是講到new creation.
就是七件的末事.
當然不是每個人都贊成這樣去分.
當然這個結構上有相銳的地方.
不過最重要的就是給我們一些背景.
講到context的時候就很重要.
每次我們讀經的時候.
我們都需要有context.
我們當然贊成大家背誦聖經.
但是你背誦聖經的時候.
請你不要只背那句很溫馨的那句.
你回想一下上下文.
為什麼會講那些東西.
上一次我們講到第一印的時候.
我解釋過了.
這些天體的意象.
不是我們現代人.
可以用我們的想法去壓下去的.
我們應該第一件事就是問.
當時的context.
當時的背景.
當時的人想起的是什麼.
第一印你記得是什麼.
白馬.
White Horse.
不是我們想起White Horse是什麼.
而是當時第一世紀的教會.
當時第一世紀的人.
他們會想起什麼呢.
很可能就會想起羅馬帝國東邊.
這個Parthian Empire.
對帝國很多時候有入侵.
或者有一些威脅.

$^{241}$你可能會問.
我看到有很多其他人解法不是這樣解的.
有些人會覺得白馬是.
敵基督的第一個.
因為在啟示錄裡面.
經常有一些Antichrist出現的.
不同的Antichrist.
最後那個當然就是最大的那個.
The Antichrist.
會不會是第一個Antichrist才是白馬呢.
所以我們發現好像有些眾說紛紜.
大家你講這樣解.
我就講另外一個解法.
我覺得大體上沒有問題.
為什麼呢.
因為天啟的意象很多時候是很有彈性.
和很象徵性很symbolic.
對於當時的人來說.
都可能有不同的方面去看法.
但我覺得起碼方法是對的.
起碼方法不是用我們現在在報章上看到的東西.
現在在世界上發生的事情.
然後加上去.
起碼方法methodology來說.
起碼首先用那個context出發.
首先我們知道當時的人要想些什麼.
我們就會更加明白那個經文的意思.
天啟文學的目的很多時候.
不是要我們讀者去做一些matching.
不是想我們去做一些配對.
最主要是要我們get到一個message.
正如我所講的.
不是一些預先寫好的天啟報告.
告訴我們將來順序會發生什麼事情.
因為很多時候這些vision.
就不是順著時間去寫出來.
譬如六章第一節.
這裡是我們的翻譯.
我看見羔羊揭開七印中的一印的時候.
上次我們看過這段經文.

$^{281}$六章第一節.
你有沒有留意那裡的keyword是什麼.
七印中的一印.
One of the seven seals.
My first seal.
行不行.
和學本 新學本都是第一印的.
當然他看下文就覺得是第一印.
但其實是沒有first這個字.
沒有首先這個字.
大家get到.
所謂第一印不是說時間上第一件事發生.
很多時候不是chronologically.
按照時間的順序來寫給我們.
因為他要我們做的是什麼呢.
不是要我們去配對.
不是要我們去做matching.
我們現代人去做matching已經不對了.
因為你加什麼進去都可以.
第一世紀的信徒.
聽見提斯祿的這些意象的時候.
他都不需要做太多matching.
可能這個方面這樣看就是Parthian.
可能另一個方面看就是基督等等.
但最主要的point重點不是matching.
而是message.
不是配對.
而是本身神學的訊息是什麼.
如果你這樣看的時候.
你就不用經常掙扎.
我要match這個或者match那個.
大家小學的時候我們經常做一些雜作.
都是做一些matching.
一條線去另一條線.
提斯祿的point不是做這件事.
而是在這些vision的背後.
他要你掌握的就是那個message.
今天你能夠在你生命裡面有提醒的一些message.
這樣的introduction之後.
我們就可以今天看第三節.

$^{321}$就是第二印.
Second seal.
楊高揭開第二個印的時候.
我聽見第二個活物說.
你來.
就另有一匹紅馬出來.
騎在馬上的得了權柄.
可以從地上奪去和平.
使人互相殘殺.
又有一把大刀切給他.
馬是什麼顏色.
紅色.
鮮紅色或者bright red.
我們可以示範一下怎麼壞鬼色經.
亂來.
什麼叫紅馬.
紅色你就找紅色的東西.
有權柄賜給他.
然後他又去殺很多人等等.
以我所知.
現在是2022年.
過去100年.
解啟斯祿的時候.
你看看一些坊間.
特別最popular最多人買的.
那些賣到成千上萬的書.
那些可以賺很多錢的書.
以我所知過去100年.
說真的紅馬的時候.
都是找一些紅色的政權代入去的.
對不對.
那你想一下過去100年有什麼紅色政權.
蘇共.
Soviet Union.
中共.
是不是.
我們就將這些共產國家.
共產的政權.
然後我們就說.
這個第二個印.

$^{361}$就一定是說到這些紅色的政權出現了.
特別在美國.
1950年代開始.
直到1991年.
蘇共解體.
在那段幾十年的時間.
一直在美國的解經.
美國看啟斯祿的時候.
很多時候我們就用了.
現在發生的事.
蘇共當時冷戰.
美蘇是冷戰.
所以每個人都說.
這隻紅馬就一定是啟斯祿的第二個燒.
有什麼問題.
就是現代人讀了我們的意思進去.
還有另一個更大的問題就是沒有control.
到底control在哪裡才行.
你要知道解經不單止是一個藝術.
解經也是科學.
解經也是science.
你不可以亂插進去的.
我們最喜歡平時查經的時候.
小組查經 團結查經.
我們就說你的領袖是什麼.
我的領袖就是蘇共.
你的領袖是什麼.
我就是中共.
你不能夠這樣.
因為什麼是control.
可以亂來的嗎.
這個就是問題了.
到底經文想說什麼.
啟斯祿我們再想想.
紅色 為什麼是紅色.
如果你看這四個人.
或者七個人.
很肯定紅色就是負面的.
如果你再看下去啟斯祿.
到十二章的時候.

$^{401}$你記得七重結構中間的時候.
最重要的意象.
那裡出現了.
又是有紅色.
就是大紅龍.
紅色的龍.
如果沒有control的話.
我們圍在一起查經的時候.
大紅龍是什麼.
我的領袖就是蘇共.
你的領袖是什麼.
我就是中共.
或者可能中國人.
因為中國人是龍的傳人.
這些又是沒有control.
完全不知道在做什麼.
完全亂來.
但如果看啟斯祿的上下文的話.
去到十二章那個大紅龍出現.
那裡講清楚是什麼.
是撒旦.
是魔鬼.
Satan.
所以這裡講到紅色的時候.
第二印.
負面的.
是迫害教會的背景的時候.
你就知道這個紅馬到底.
可能會是指到什麼勢力.
再看舊約的背景.
這裡我們這次不可以詳細的講.
這裡四個印四隻馬.
很明顯一個熟悉舊約的人.
他就會想起舊約裡面其中一本書.
因為只有一本書是提到有四個顏色的馬.
又是天啟文學.
那本書就是撒迦利亞書.
撒迦利亞書在一些意象裡面.
講到上帝要來審判列國的時候.
就有四隻馬.

$^{441}$四隻不同顏色的馬出現.
所以啟斯祿是在引用撒迦利亞的.
那個的message.
去carry forward.
就是說以前.
講到撒迦利亞時代的時候.
重建聖殿的時候.
遇見很多的迫害.
很多的難處.
上帝來掌權.
上帝去審判列國.
告訴先知.
告訴猶太人.
你們可以繼續去建立這個聖殿.
因為上帝仍然是掌權.
就將這個message carry forward.
去到第一世紀.
第一世紀教會.
又是遇見很多難處.
很多的艱難的時候.
甚至當時的猶太的拉比.
他們看到紅馬的時候.
講到撒迦利亞的紅馬.
都是講到什麼呢.
就是流無辜人的血.
所以為什麼就是紅色.
如果你有這些舊的背景.
和當時的背景的時候.
你就發現第一世紀的信徒.
他們聽見這個第二印的時候.
他們會想起什麼.
和羅馬帝國有什麼關係.
一個打壓他們的政權.
就是羅馬帝國當時.
是有很多的內戰.
例如在第一世紀的60年代.
當時的英國.
因為羅馬擴展到英國.
當時有些反羅馬的起義的革命出現.
就釀成有15萬人死了.

$^{481}$按照統計.
或者在處後60年代.
67 68 69這幾年.
在巴勒斯坦地區.
又是很混亂的.
因為有第一次猶太人起義.
去到反對羅馬的統治.
在那一年裡面.
在羅馬的京城.
又是有很多的內戰.
很多的黨爭.
單單在那一年.
就有四個的帝王出現了.
你見到由左上角開始.
加布.
然後到阿佛.
然後到維特利亞.
然後到最後的威斯比亞.
他就打壓了所有的異見分子.
將這些的政敵全部殲滅了.
然後他就可以作王比較長一點的時間.
10年.
在這段時間叫做.
四個帝王的一年.
就是在京都裡面.
在羅馬裡面.
就是非常之混亂的.
非常之不和平的.
因為有很多的文明戰爭.
很多的人被殺戮而死的.
第二天.
你都講到從地上奪去和平了.
你一聽到和平這個字.
Peace.
羅馬人就很明白.
因為當時的羅馬帝國的大外宣.
他們的政治宣傳的機器.
就是講你要靠羅馬帝國.
靠我們的刀.
靠我們的軍隊.

$^{521}$才可以有和平的.
當時經常講一種叫做.
他們的大外宣.
當時的羅馬帝國.
就是有一個Pax Romana.
這件事.
就要靠著羅馬的刀.
靠著羅馬的軍隊.
他的刀他的劍.
才可以在地上有盛世.
羅馬盛世.
這裡講到第二印的時候.
就說這個印出現的時候.
原來不是帶給地上平安.
而是奪去世上的和平.
所以就跟羅馬講的東西.
羅馬電視台講的東西.
或者羅馬帝國中央電視台講的東西.
是完全不同的.
就是說羅馬帝國的太平盛世.
其實就是吹噓.
上帝就是要去啟示一個現實出來.
去揭示一些東西出來.
我們講過啟示第一個字是什麼.
Apocalypse 揭開的意思.
以前我們不知道沒有留意的.
現在揭開了.
這裡揭開什麼.
就是揭開羅馬帝國醜陋的一面.
或者它的真相真實的一面.
以致我們不會被它欺騙.
為什麼特別要這樣做.
因為第一世紀的教會.
就正是在羅馬帝國的打壓的底下.
還有很多的引誘.
很多的試探.
你記得啟示錄開頭.
就是寫給教會的七封書信.
是不是.
這七間的教會.

$^{561}$或者這七個的群體.
在哪裡.
在今天土耳其.
就是當時的羅馬的東部.
我們叫做Roman East.
在那裡發生什麼事.
第一世紀就是羅馬帝國的勢力.
要向東邊不斷擴張.
不斷擴張的時候.
他們就去打壓那些異見分子.
因為擴張的時候.
就要去宣傳一件事.
去推廣一種叫做.
Emperor Worship.
帝王的崇拜.
或者Imperial Cult.
因為教會不聽話.
因為教會只是拜耶穌基督.
只是拜耶穌基督.
稱祂為主為神.
就不稱帝王為神.
所以就有很多的打壓.
有硬有軟的.
硬的當然是將那些領袖.
包括史陶約翰.
將他們收監.
將他們exile等等.
但是軟的打壓.
就更加可怕.
軟的打壓是什麼.
就是給你經濟上的好處.
給你經濟上的利益等等.
就將你的信仰沖淡.
你說硬的打壓厲害.
還是軟的打壓厲害.
很多時候可能就是軟的打壓更加厲害.
七教會的書信.
如果你回去看一下.
很多時候就是面對著這樣的情況.
就是帝國給教會的信徒.

$^{601}$有很多的人有很多的試探.
不要緊的.
你親羅馬.
不要緊的.
你拜羅馬的皇帝.
將耶穌一起拜都可以的.
就將教會的信仰沖淡了.
以致那些人離開信仰.
以致那些人不知道自己敬拜的是什麼.
這個試探起碼是一個更大的試探.
我們看一段七教會書信.
以弗所的書信.
第一封書信是寫給以弗所教會的.
我們讀一段.
二章第三節.
這裡是基督寫信給教會.
他說:你又忍耐.
曾經為我的名義的緣故忍受一切.
並不困厥.
然而有一件事我要責備你.
就是你已經離棄了你起初的愛.
所以你應當回想你是從哪裡墜落的.
並且要悔改作起初所作的事.
你若不悔改.
我就要來到你那裡.
把你的燈台從原處移去.
一般來說.
我們聽見這裡離棄了起初的愛.
我們就想起.
就是你對主的愛.
起初的愛.
有沒有沖淡到.
有沒有少到.
對主的熱誠.
因為我們聽起來.
起初的愛.
first love.
我們就想起.
spiritual temperature.
你對上帝的熱誠.

$^{641}$對人的關愛.
上帝兩個最大的誡明.
愛上帝.
和愛身邊的人.
我們通常就會想起.
想起這一方面.
我相信那個意義一定是包括這一方面.
但我這幾天去整理一下我的資料.
以往很多年收集了這些資料.
在已忽所成.
就找到很多這些文字.
這些文字有很多的石碑.
在當時已忽所成的一些貴族.
一些社會裡面比較地位重要的人.
他們就丟自己的名字在那裡.
或者自己的稱銜.
我是哪個officer.
我是哪個政府官員等等.
然後就加一個字叫做.
Philosophastus.
或者Philosophastoi.
就是眾數和單數的分別.
就這樣自己自稱的.
這些石碑很多時候就放在.
這些市中心.
大會堂或者城門口等等.
就是那個人是一個社會裡面比較高位的.
他就會這樣形容自己.
就會說自己是Philosophastus.
什麼意思.
Philo就是愛的意思.
你聽過philosophy 是不是.
就是愛那樣東西.
Philosophastus就是帝王的意思.
Caesar Philosophastus.
就是說這些人去到說自己是什麼.
Emperor lover.
我是向帝王去效忠的.
向這個權勢來去效忠.
有100多200個這樣的石碑.

$^{681}$這些高位的人都說自己是什麼.
Philosophastus.
我是愛羅馬帝王的人.
你這樣看的時候.
你就會發現.
起初的愛就不是單單我們心裡面熱誠這麼簡單.
而是有另外一個權勢.
去到霸佔了我們對耶穌基督的效忠.
有另外一個權勢.
有另外一個王去到爭奪我們對主的歡心.
Emperor lover.
如果你這樣看爾弗所這封書信的時候.
就不是單單溫馨的問題.
不是單單有沒有激情問題.
對主有沒有熱誠的問題.
就是問我們的心.
到底Allegiance在哪裡.
到底是放在Emperor那裡.
還是放在基督那裡.
到底是Caesar還是Christ.
這句話在約翰的口中出來.
就是起初的愛加上這個背景.
你再找一下約翰寫過的東西.
約翰福音.
約翰福音其中一節.
這裡說什麼呢.
快到受苦節了.
又是說到耶穌被賣,被釘的故事.
約翰福音第十九章第十二節.
當時猶太人在討論.
那些權貴在討論.
就是和彼拉多.
就是審耶穌的羅馬官長.
就和他們爭辯中.
到底是不是要釘死耶穌.
十九章十二節.
從那時起.
彼拉多是想釋放耶穌的.
因為找不到他有什麼要被釘死的罪.
那時猶太人卻喊了一句話.

$^{721}$如果你釋放這個人.
耶穌.
就不是凱撒的中神了.
Fellow Caesar.
如果你釋放這個人的話.
你就不再是Emperor lover.
你就不再是羅馬帝王的中神.
凡致命為王的.
就與凱撒為敵.
他們就這樣去要脅彼拉多.
如果你釋放耶穌的話.
你就不再是Fellow Caesar.
不再是愛帝國者.
不再是帝王的中神.
和彼拉多一樣.
彼拉多當時的處境.
猶太人逼他.
你不可以在羅馬帝王和耶穌之間.
作牆頭草.
你不可以十五二十.
十五十六.
你不可以去到.
到底我是愛耶穌多些.
還是愛凱撒多些.
不能作牆頭草.
Fellow Caesar.
起初的我.
啟示了想做些什麼.
就是想告訴我們.
這個帝國世間裡面有一些權勢.
有一些主人.
有一些勢力.
是想拿走你的寬心.
當時羅馬帝國.
就是用權與利.
第二人.
包括第一人.
就是想揭示這個真相.
這個帝國本身就是流無孤人血.
甚至你繼續看下去.

$^{761}$巴比倫的淫婦.
是撒旦大雄龍的工具.
是迫害我們教會.
殺害我們的弟兄姊妹.
你豈可對帝王效忠.
怎麼可以這樣選擇.
又選擇基督.
又選擇一個撒旦的工具.
去殺害教會的.
殺害我們主來弟兄姊妹的.
怎麼可以向他效忠.
怎麼可以變成他的愛人.
根本沒可能.
你應當回想.
回想什麼.
你從哪裡聚落.
你應該要悔改.
作起初的事.
大家坐在這裡.
我想問你.
有沒有一些東西.
是佔據著你對基督的效忠.
那些可以是人.
是或物.
到底什麼去搶著我們.
對基督的心.
受苦節復節就來了.
我希望我們所謂大齋期.
我們去看一些靈修等等.
我們去預苦期.
去等著預備一個受難節.
來到的時候.
不是一些單單儀式上的東西.
而是一個生命主權的挑戰.
到底你是為誰效忠.
第一世紀的信徒.
就是被引誘去效忠這個權勢.
得到不同的利益.
得到不同的方便.
你是怎樣.

$^{801}$思想一下你從哪裡聚落.
你對主的心.
起初的愛.
我們很多時候都想起.
可能我剛剛信主的時候.
但我想你想的.
可能不只是一樣.
對起初的愛.
我們離棄了.
要回想起初的愛是什麼.
是什麼時候.
我昨晚想到的就是.
會不會是疫情之前.
疫情之前的你是怎樣.
你想一下你對上帝的關係.
你和上帝的關係.
你可能還是有那種熱誠.
疫情之前.
疫情已經兩年多了.
還繼續持續下去.
繼續不退.
你起初的愛去了哪裡.
這裡說什麼.
你要回想.
從哪裡聚落.
悔改.
不是去等.
不是等我什麼時候再溫馨一點.
我就愛主多一點.
這不是重點.
重點是尊重.
是信仰.
做起初所做的事.
悔改做起初所做的事.
我們等會回應詩歌有一首.
雖然舊一點.
叫做永恆的你.
但我很喜歡這首歌.
因為裡面有一句說話.
權勢終變改.

$^{841}$唯一我們可以去效忠的.
應該是耶穌基督.
我們誠心低頭的祈禱.
主我們多謝你.
因為你是那一位.
願意愛我們到底的神.
主你是我們唯一.
應當去效忠的.
求你幫助我們來到會上.
我們是不是離棄了起初的愛.
起初對你那種忠誠.
對你那種唯一.
神我們求你幫助我們去悔改.
讓我們做起初所做的事.
主叫我們不要忘記.
叫我們不要去善忘.
叫我們知道我們是屬於你的.
讓我們生命.
是以你為我們的唯一.
去事奉你 去愛你.
我們這樣祈禱 奉主命.
我們這個時候是....
\newpage



\section{創世記}
\label{sec:shYH8FbN1Gk}
\textbf{創世記第4章 (4.3) - 背景釋經 [近東爭辯詩,創四章,該隱的後代] (陳崇基牧師)}
\newline
\newline
連結: \href{https://youtube.com/watch?v=shYH8FbN1Gk}{\texttt{https://youtube.com/watch?v=shYH8FbN1Gk}} ~~~~ 語音日期: 2022-03-08
\newline
\newline
\hyperref[sec:jFHDnibTFuo]{\small{< < < PREV SERMON < < <}}
~
\hyperref[sec:index_chronic]{\small{[返順時目]}}
~
\hyperref[sec:index_scriptual]{\small{[返順卷目]}}
~
\hyperref[sec:hkH9QTCOPtw]{\small{> > > NEXT SERMON > > >}}
\newline
\newline
$^{1}$好,我們繼續講第四章創世記.
我們很快講一講脈絡.
第四章創世記.
由第一章開始.
就是天地的創造,人類的創造.
人類的叛變,這是第三章.
然後我們進入第四章.
大家記得第三章之後.
人流亡於這些之外.
因為我們叛逆和犯罪.
但是上帝仍然出來.
和我們同行,在我們中間.
所以第四章預設了這件事.
上帝和神仍然與人建立關係.
所以第四章開始就講到該隱阿伯.
他們是一個耕田的,一個是牧人.
然後第四章,大家記得.
很多篇幅講到該隱離開上帝.
他也是叛變,然後他離開了上帝.
然後他就建立自己的城鎮.
所謂文明,有其他的文化出現等等.
然後到第四章尾就講到.
亞當的另一個兒子,就是瑟特.
他的後代繼續延續下去.
就是到第五章了.
我們上一次講過.
我們做了一些比對的研究.
我們要明白當時的人.
他們的作品,一些近東的文獻.
以至他們怎樣去講述一些.
遠古人類發生的事,和一些起源的事等等.
和聖經有沒有相似的地方呢?.
用的語言,或者一些主題.
甚至是討論一些手法.
有沒有相似的地方,然後有沒有差異的地方.
差異的地方很可能是上帝特別啟示.
他的真理的那個方面.
所以我們首先要知道.
聖經首先是寫給當時近東的人.
當時的世界.

$^{41}$所以他們去討論一些議題.
討論一些問題的時候.
不一定像我們現代人那樣去講.
好像我們創世的第一章那裡講到.
我們的興趣,現代人.
就是科學.
大爆炸等等.
但是近東的人,包括聖經當時的世界.
他們所討論起源的時候.
就不是我們今天.
未必是我們今天所著重的問題.
所以我們看看怎樣去比對的時候.
我們就明白,為什麼聖經會記載這些東西.
以及用這些手法來講.
甚至為什麼他的脈絡是有他的特色的.
創世的第四章裡面我們剛才提過.
有耕田,有畜牧.
然後兩個兄弟發生了什麼事等等.
這裡有沒有一些比對可以幫助我們.
為什麼會第四章.
講到人類的起源.
都是在講一些史前史,很久之前的事情.
就會講到畜牧和農業的.
就是該隱,阿伯,兩個行業.
我們可以比對一下一些在近東.
又是討論到一些所謂遠古歷史.
遠古人類做過什麼的.
這些所謂的史詩.
或者一些的爭辯士.
這裡有一首爭辯詩.
就是那個標題叫做Dumuzi和Enkintu.
就是兩個的神明來做一個爭辯.
有一大堆的這些詩歌.
我們可以叫它做爭辯詩歌.
通常就是兩者對答.
裡面所討論的問題.
或者所講的故事.
都是有兩方面的.
就是兩方面很多時候就是講到人類.
開頭那些的工業.

$^{81}$人類的開頭的行業.
這個詩句或者這個爭辯詩.
這個是當時近東一個很早期的.
一個蘇美人的一首詩歌.
講到什麼呢?.
就是講到兩個的神明.
一個就是畜牧的神明叫做Dumuzi.
另一個就是農耕的神明Enkintu.
這兩個的神明之間有一些的對話.
和有一些的爭辯.
那個爭辯是什麼呢?.
就是這個畜牧的神.
或者他代表著這些畜牧的人.
還有這個農耕的神明.
他當然是代表耕田的.
在城市裡面或者建立了聚居地方的.
他們自己有耕種.
他們之間有一個這樣的爭辯.
那個爭辯就是講到什麼呢?.
就是他們大家都是和這個Inanna.
就是依娜娜這個女神.
有一個這樣的三角的關係.
兩者爭辯就是我是不是比對方好一點.
然後就得到這個女神子的歡心.
所以他們的詩句就好像一個話劇一樣.
開頭就說有一個大一點的神明.
大家知道當時巴比倫或者蘇美爾.
他們都是多神論.
這個Inanna就是一個女神.
有一個哥哥就是太陽神.
就想這個Inanna去到許佩比.
這個畜牧業的神.
但是這個Inanna就不喜歡這個畜牧.
他就喜歡耕田的農業的那個.
所以那個畜牧業的杜牧之.
就去到說我想我所出產的東西.
我出產的東西比如一些羊毛等等的東西.
是比這個農神所出產的東西好.
這個農神就和他大家有一個這樣的辯論.
就是說我出產的東西比如稻米,豆,田裡面出產的東西.

$^{121}$是比你好的.
整首詩就是在說這兩個到底誰比較好.
誰比較厲害.
最後這個Inanna就選了這個農神.
許佩了給他.
但是這個詩最後的結束.
就說到這個畜牧神和這個農神.
大家都變成朋友了.
其實大家都需要大家的.
所以這個農神就給這個畜牧的神.
去用他的田地.
比如收成之後.
這些畜牧的牛羊就可以去農田那裡.
去吃剩下的東西.
所以兩個的行業.
兩個的工業.
都是大家變成朋友的.
大家都是好的.
最後大家都是向上面的神.
這個Inanna代表豐沃的神.
向他獻上一些頌讚.
那你就想.
為什麼要做這樣的東西.
為什麼要做這樣的比對.
這些神怪的東西.
這裡起碼這首爭辯詩.
可以讓我們看到.
當時的古近東的人.
他們想起人類的文明的時候.
人類之所以能夠繼續發展.
就需要這兩個行業.
一個就是畜牧.
一個就是農田.
人類文明的支柱.
就是靠這兩個行業.
所以說到一的史言史的時候.
很多時候就會提到這兩個行業.
就是畜牧和耕田.
並且這些人類文明.
在蘇美爾巴比倫.

$^{161}$這些所謂外邦的虛神的多神論的底下.
很多時候討論的時候.
就是將他們神聖化.
為什麼會有畜牧業.
為什麼會有農田.
因為不知道是誰的神明.
他就將這些東西賜給了人.
將這些技術.
將這些方法.
賜給了人類.
你記得上一次我們提過.
一些近東的創世的神話.
這些的史詩.
他們說到人類是怎樣的.
人類被造出來是做這些神明的奴隸.
所以他們去耕種.
去吃苦.
然後他們就獻上一些東西.
初一十五獻上一些東西.
去供應這些神明.
因為這些神明都會肚餓的.
所以這些人類的地上的文明.
就是靠這些神祇來賜給他們.
最後歌頌的是他們.
所以你就想一想.
這個就正是創世記第四章.
開始講到人類起源.
都是一些起初的故事.
就用回這些主題.
當然他所要表達的神學信息.
就完全是不同的.
但是在近東的世界.
聖經要去跟這些神話.
或者這些主題.
去講述.
去表達的時候.
就用回他們那種方法.
或者用他們所用的議題.
所以人類文明之所以有.
就是這些神祇.

$^{201}$來賜給人類.
再看創世記第四章的時候.
你就知道為什麼第四章要討論.
人類的文明.
有什麼個人的後裔.
他們是哪一個行業的畜牧業的父.
或者音樂.
或者畜牧或者耕田的父等等.
但是我們要留意.
雖然用同樣的議題來討論.
好像同樣的手法.
有相似的地方.
但是也有很多不同差異的地方.
再看多一篇近東的史詩.
這些也是跟人類的文明很有關係的.
就是講到農業.
農作業.
這個詩句.
我們現代人看就覺得很奇怪.
這個是稱讚你的噱頭.
大家知道噱頭.
在耕田是多麼重要.
一個很重要的科技或者工具.
人類發展到某個階段的時候.
他們懂得怎樣去煉鐵.
以致煉了鐵之後的噱頭.
能夠被發明出來.
以致他們耕種的時候.
就會方便很多.
以及是有效很多.
所以這裡有一個古代的近東的神話.
一些蘇美人的神話.
非常古代.
這個是在2000年代之前的.
這裡說什麼呢.
他說噱頭使人們繁榮.
一切的繁盛都是從噱頭而來.
因為有噱頭就有大麥.
然後可以製造獵網.
然後可以使人存活.

$^{241}$是年輕人的力量.
他描述到這些科技.
這些起初的農作工具.
就是一些很大的力量.
可以建造美好的房屋.
建造城市.
最重要是耕田.
噱頭可以擴展到美好的農田.
也可以征服到一些頑固的用地.
可以拔野草.
制服這些野草.
噱頭最後.
正如我剛才所說.
就是因利而下的另外一個神明.
就是命定的工具.
就是賜給人類的.
就是耕田這個行業.
就是神聖化了.
所以噱頭就有美名.
而且讚美這些神明.
這些是神明的名字等等.
所以又是說到人類的文明.
最後所頌讚的.
或者所高舉的.
當然是人自己的建樹.
但是也是神子.
這些神明.
你也可以賜給人類.
我們看了近東的史詩.
例如《鋤頭詩》等等.
我們就可以幫助我們了解.
為什麼《創世記》第四章的脈絡.
或者主題是跟類似的東西有關.
四章第一節.
這裡說阿當跟他的妻子同房之後.
就生了哥哥.
藉著耶和華幫助得了一個男兒.
然後生了他的弟弟.
阿伯就是牧羊.
哥哥就是種地.

$^{281}$這兩個主題.
很明顯就是跟我們之前看的神話故事.
兩個最大的人類經濟的支柱.
種地的耕種.
還有牧羊業.
然後說到兩個之間的關係.
當然有相似的地方.
但也有很多差異的地方.
有一天.
該人在地裡的出產.
拿來當作祭物.
獻給耶和華.
阿伯有把自己羊群中一些投生的羊的滋油.
拿來獻上.
耶和華看中阿伯和他的禮物.
只是沒有看中該人和他的禮物.
該人非常憤怒.
垂頭喪氣.
為甚麼神看中了阿伯的祭物.
但沒有看中該人呢.
以新約後期的解釋就是.
by faith 信心的緣故.
或者以往很流行的解法就是.
因為血祭.
因為他是畜牧的.
所以牧羊人獻上的羊的滋油等等.
是一種生命的犧牲.
耶和華看中了.
但該人只是農產的獻祭.
所以沒有看中.
原文可以有甚麼線索呢.
其實也有的.
原文已經解釋了一些.
例如說地裡面的出產.
拿來.
該人拿了這些東西.
ma priha dama mika la donai.
即是說ma priha.
即是從果實.
即是地裡面的出產.

$^{321}$田產.
地.
dama這個字在第一章第二章第三章都有出現.
地裡面的出產.
mika.
即是供物可以譯作.
或者tribute.
其實很多時候指的是一個低層的人.
例如人.
獻上給一個高級的王.
很多時候會用mika這個字.
獻給耶和華當然是非常適合.
因為是mika la donai.
獻給上級的.
應該獻上是最好的.
誰知該人只不過是從地裡面的出產獻上.
阿伯將自己的羊群中投生.
羊的自由.
mikurot suno mikheveim.
這個意思是.
mikurot.
即是從投生的.
mikurot這個字.
後面這個字就是第一章.
suno.
his sheep.
他的羊群.
以及他的自由.
所以是將當時近東最好的東西獻上.
所以很明顯在原文裡面.
有一個清楚的線索.
讓我們知道阿伯的兒子是一個合神心意的兒子.
該人只不過是隨便在地裡面的出產而已.
所以引申至聖經裡面所記載的第一個murder.
不代表是世界歷史裡面第一個murder.
只是聖經裡面記載的第一個.
亦要留意第三章和第四章有很多的聯繫.
作者不只是記載兩兄弟發生的事.
畜牧和耕種的事.
他連繫了三章.

$^{361}$三章就講到fall of mankind.
人類的叛逆人類的墮落.
所以有些字都是大家彼此呼應的.
並且是有心的呼應.
講的是該人的下場.
該人的故事.
其實是另一個人類的墮落.
另一個人的罪惡.
他的邪惡 他的叛逆.
所以三章十六節中了一個字.
很少在希伯來聖經裡面用的.
你要戀慕你的丈夫.
這個對女人所說的事.
審判的說話.
他卻要管轄你.
Tis you critique.
即是戀慕 中文譯為戀慕.
是一種很強烈的控制.
三章十七節.
地必因你的緣故受咒坐.
Hadamah這個字又出現了.
受咒坐.
Aruah這個字又出現了.
這三個字在四章.
講到該人的故事.
又說跟他講什麼.
又是對人說話.
三章這裡都是對女人和男人說話.
然後這裡就是對該人說話.
你若行得好 罪就伏在門前.
罪必定會戀慕你.
Tis you kato.
後面這個to這個o這個字.
其實是一個字尾.
所以這兩個字非常相似.
一個是講到兩者兩性之間的張力.
一個是講到罪.
罪去到戀慕你.
是一種很強烈的控制 制伏.
幾乎形容到罪好像一隻猛獸伏在你門前.

$^{401}$隨時要吞滅你.
你如果還跌落這個試探.
你就會失敗.
你一定會死.
所以你又要跟他爭鬥 制伏等等.
四章十一節.
又是講到該人的審判.
你要從這地受咒助.
所以這些字都有呼應.
除了這些很清楚的呼應外.
那裡這個字有呼應.
戀慕 地受咒助.
講出向東變.
在創世記的章節裡面.
逢向東都是不好的.
向東就是離開耶和華.
因為這段是向東變.
繼續向東走就是離開耶和華.
推石責任這些描述.
這些都讓我們想起三四章.
講到該人不是單單一個單獨的故事.
而是一個罪的擴散.
罪被描述越來越多.
所以到了第四章.
我們發現後面繼續講到該人的後代.
他的兒子爾諾.
然後他又建造了一座城.
按照他的兒子.
吸那城叫做爾諾.
爾諾生爾那 爾那生米胡雅利.
米胡雅利生馬圖薩利.
馬圖薩利生拉麥.
這裡記載一二三四五六七 七代.
這是第四章.
一直追溯離開了耶和華.
叛變的宗族.
由該人開始直到拉麥.
這裡我們不需要糾纏很多.
我們用現代人去讀經.
然後引申出一些護教學的問題.

$^{441}$該人的妻子是誰.
是不是他的妹妹.
是不是在律法之前可以跟妹妹結婚.
兄妹可以結婚.
這些全部不需要引申出來.
這些是我們自己找麻煩.
第一章已經講了.
人類被創造.
所以這裡很清楚.
他跟妻子.
然後有座城.
明明沒有人的 為什麼有城呢.
我們沒有說.
是那個人派人殺他.
其實經文已經是.
assume了一個populated world.
一個已經有人口.
有人居住的世界.
所以派人去尋仇.
殺他等等.
所以第四章.
追溯他的家譜.
然後講到拉麥.
第七代.
娶了兩個妻子.
一個叫阿達 一個叫西拉.
這個很明顯.
是一個對比.
跟二章第一章的時候.
一夫一妻 一男一女.
神所設立的婚姻.
有一個對比.
阿達生了阿伯.
阿伯是居住在象棚.
畜牧業的鼻祖.
阿伯的兄弟.
就是尤伯.
尤伯是音樂.
吹簫的鼻祖.
然後西拉的另一個妻子.

$^{481}$生了土伯.
戈隱.
土伯戈隱是打造銅鐵器的將人.
土伯戈隱的妹妹.
就是娜瑪.
這裡就講到人類的文明.
一些文明的建樹.
等等.
跟近東的士師記載.
人類的發展.
其實都是非常相似.
相似得來.
有很多差異的地方.
因為所追溯的.
就是戈隱的宗族.
離開耶和華之後.
他們半邊繼續.
聚去擴散.
然後越來越多.
這些罪惡.
越來越泛濫的時候.
怎麼看到呢?四章二十三節.
這裡另外一首詩歌.
這首詩歌.
又給我們去一個對立.
就是二章最後的時候.
亞人阿丹.
上帝為他預備一個幫手.
女子.
然後他說.
這是我骨中的骨 這是一條鋒利的肉.
是一首很美的詩歌.
四章二十三節這個詩歌就不同了.
這是一個半邊族的人.
拉麥出口的一首詩歌.
阿大.
和洗拉他這樣說.
這兩個太太.
聽我的聲音 拉麥的妻子聆聽我的話.
壯年人傷我 我必殺他.

$^{521}$可以譯作蓋式.
我殺了他.
青年人.
去損我 我必害他.
人若殺該人.
就遭報七倍.
要殺拉麥的話就必遭報.
七十七倍.
這裡你看到在詩歌裡面.
想表達的是什麼呢?.
是一種殺害.
殺氣騰騰.
罪惡的泛濫.
暴力的泛濫.
第六章洪水之前.
你會發現這件事會越來越多.
其實聖經的作者.
想我們.
跟著這個故事的脈絡.
以致到最後.
人好像自我毀滅.
越來越離開神.
叛逆.
墮落.
罪惡越來越擴散.
越來越滿盈的時候 洪水就來到.
神的審判就來到.
這裡又用七字.
第一章我們提過.
有很多這樣的七字.
出現了.
七倍等等.
這裡就是殺氣騰騰的.
說到他去到.
殺死那些傷害他的人.
所以那個暴力的文化.
就越來越多.
暴力好像很橫行的.
一個描述 一個世代.
不單止這樣.

$^{561}$如果你有時間可以數一下.
其實有77個.
字母的.
拉麥的詩歌.
你可以數一下 第一行有14 第二行有25.
很想去到.
完整的去表達.
這個暴力.
非常之厲害.
所以這裡有一個對比.
創世記.
神所創造的世界.
一切都是好的.
很多的.
七個字.
perfect.
美好的.
sacred 神聖的.
那人對這個.
女人的詩歌.
一個很美的詩歌.
一個合一的詩歌.
然後這個拉麥的詩歌.
就是一個.
死亡.
步向死亡.
和罪惡滿盈.
然後下一次.
我們看到失特的加普.
和這個.
在邪惡.
罪惡氾濫的一個世代.
一直延續下去的時候.
失特的加普.
就是一個宣告.
耶和華的一個宗族.
如果你喜歡我們的影片.
就記得要like comment和share.
Like, Comment 和 Share 啦!.
謝謝大家收看 再見.

\newpage



\section{}
\label{sec:hkH9QTCOPtw}
\textbf{Family Sunday School - EP 3 (Jesus Gets Arrested)}
\newline
\newline
連結: \href{https://youtube.com/watch?v=hkH9QTCOPtw}{\texttt{https://youtube.com/watch?v=hkH9QTCOPtw}} ~~~~ 語音日期: 2022-03-11
\newline
\newline
\hyperref[sec:shYH8FbN1Gk]{\small{< < < PREV SERMON < < <}}
~
\hyperref[sec:index_chronic]{\small{[返順時目]}}
~
\hyperref[sec:index_scriptual]{\small{[返順卷目]}}
~
\hyperref[sec:yncumdAd8iU]{\small{> > > NEXT SERMON > > >}}
\newline
\newline
$^{1}$[Music].
Hello, everyone! I'm Miss Rebecca. Welcome back to Family Sunday School..
This time we're going to review what we talked about last time..
Last time we talked about Last Supper..
And one of the points, one of the themes was "Jesus is the Bread of Life.".
Last time we talked about a lot of buns, right?.
Have you played the bun game?.
Have you learned that there are many different cultures and people all over the world that need buns?.
I'll give you a book to read. It's a cover and there's a little bit of content in it..
I'm using this book for my kids to read..
But you can also use other books in the library about buns or how to make buns..
This book actually talks about a lot of different kinds of buns..
If I give you a quick summary, you'll know how people need buns..
Why is the bun in Last Supper so important?.
The symbolism of the bun is that it represents Jesus' body..
Jesus' body represents the bun..
Why?.
We need to look at Old Testament..
In Old Testament, the God/Goddess was tied up by the Egyptians..
When they were slaves, they were God/Goddess..
They were taken out of Egypt..
They were taken out of Egypt and they could leave Egypt..
Leave slavery, flee Egypt..
They crossed the Red Sea. Moses crossed the Red Sea..
God/Goddess walked across the Red Sea..
They had to wait..
Wait for when to make a promise..
What did they eat?.
God/Goddess said,.
"It's better to be a slave. At least I have food. I'm starving.".
God/Goddess heard them complain..
If God/Goddess could bring you out of Egypt, what else couldn't God/Goddess do?.
At that time, manna would fall from the sky..
It's like a food bun falling from the sky..
In Old Testament, we see that the bun, the manna,.
if you eat it, you will be full that day..
But it doesn't mean you won't die..
Fast forward to the New Testament..
The Savior has come..
Jesus is born..

$^{41}$If you eat the bun, you will live forever..
Very different, right?.
What kind of bun can live forever?.
Pineapple bun? No..
Hamburger bun? No..
The rabbit bun? No..
The cow bun? No..
Only Jesus..
So why did Jesus tell the disciples,.
"My body is like a bun..
If you eat this bun, you will live forever.".
Because he knew the culture and the history of the bun..
He remembered to slowly connect with it..
Okay, this is the bun..
Next one..
The second point I mentioned last time is.
Do you remember Judas left the group to betray Jesus?.
Remember? Jesus' family..
I asked you last time,.
Why did Judas do that?.
What will happen next?.
Jesus knew..
Why did he do that?.
Why did he allow that?.
Remember?.
This is the great saving plan..
God's great plan to save people..
God's plan must be fulfilled..
God said he would do it..
Jesus said it would happen..
It will happen..
But does it have to be so bad?.
We will continue the story today..
Answer a lot of different questions..
And refer to different points..
I will ask you questions too..
But before we start, let's play a game..
Come with me..
Hello, let's play a game..
What am I?.

$^{81}$I will show you different photos..
Ready?.
Hmm, a wheel..
Hmm, a screw..
Hmm, what is this?.
Can you guess?.
Let's start..
Art rack..
We need to dry our paintings in art class..
There are many art class teachers..
You can dry many paintings in this rack..
You can open and close it..
Okay?.
Next one..
What am I?.
Focus..
Hmm..
Can you guess?.
Hmm, what is this?.
Hmm..
Hmm, what is this angle?.
What angle is this?.
Hmm..
Can you guess?.
Hmm, a big circle..
What is this?.
Let's start..
It's a stool..
Ha ha..
Can you see?.
This is the seat..
This is the angle..
This stool can be turned..
This is the screw..
This is the plate..
This is a foot..
Ha ha..
Okay, next one..
I guess you can guess..
Let's start..

$^{121}$Hmm, what is this?.
It looks old..
Hmm..
This is the screw..
This is not the screw..
This is a nut..
Hmm..
Hmm?.
What is this?.
What is the hint in this picture?.
This angle is a bit different..
What is this?.
Can you guess?.
Let's start..
Maybe many children have not seen this..
Or adults, parents have seen it..
Okay, let's start..
Radiator..
It's a radiator..
It's not a central air radiator..
There are many gadgets on the side..
This is the angle..
This is the coil..
This is the switch..
This is the big screw..
Wow..
We have so many "What am I?".
I want to remind you..
Hmm..
If you look at it from different angles,.
you can't see the whole picture from every angle..
Right?.
We think with this mentality..
We will talk about the Bible story..
Perspective..
Point of view..
Angle..
Alright?.
Bible story time..
Hi, it's story time..

$^{161}$Where did we leave off last time?.
Where did Jesus tell Judas,.
"Do what you have to do.".
What did Judas do?.
He betrayed Jesus..
Telling others where Jesus is..
Because many leaders don't want Jesus there..
They want Jesus to leave..
But they can't find Jesus..
They are afraid..
Because the public loves Jesus..
They like him..
Because he did a lot of good things..
But they can't find out where he is..
Only the people inside will know..
Right?.
In this video,.
Judas went to find Jesus..
Jesus went outside..
He called other leaders..
"The time is here.".
"I will be caught.".
"You have to protect yourselves.".
"Don't sleep.".
"Pray to the Lord to help you.".
"You can resist all the attacks.".
"All the temptations.".
But they fell asleep..
Jesus started praying..
He asked God,.
"I know this plan must be fulfilled.".
"But if it is possible, can we not do it?".
"But I know your will is my command.".
"God, whatever you want.".
He finished praying..
What happened next?.
Judas told the soldiers,.
"Hey, the one I love is Jesus.".
Wow!.
Let's see..

$^{201}$This is Judas..
The people next to him..
"I know this is Jesus.".
Not everyone has seen Jesus..
Not everyone knows who to catch..
Now he knows..
What happened next?.
Peter..
Jesus said to Peter,.
"Peter, you will deny me.".
"Pretend not to know.".
"Deny.".
"You know me.".
"You pretend not to know me.".
"Of course it's not me.".
"Never.".
This is Peter..
He picked up a sword..
Jesus said, "Protect yourselves.".
"Bring the swords.".
"Because there will be an attack.".
Before, they went around with Jesus..
They talked to people..
They didn't need swords..
Now they need them..
Peter took out his sword..
What did he say?.
He didn't say anything..
He just cut open the soldier's ear..
First, was his aim good?.
Wow, he cut it open..
I don't think the soldier expected this..
But we know from Peter's perspective,.
from his perspective,.
he loves Jesus very much..
He loves Jesus very much..
He is a good friend..
He must defend Jesus..
No matter what, he will protect Jesus..
"Hiyah!".

$^{241}$What do you think of Jesus?.
"Put the sword down.".
"Put the sword down, Peter.".
The soldier's ear was cut off..
Remember, he was bleeding..
He was bleeding..
Ouch, it hurts..
What does Jesus do?.
He didn't take him to the hospital..
He didn't need stitches..
He didn't need to do much..
He just put his hand on the soldier's ear..
Healed..
Right?.
If Jesus can be blind,.
he can also be able to see..
If he can cut off the soldier's ear,.
he can also be able to see..
Wow..
Jesus said,.
"I have to go.".
"This is part of the plan.".
"This has to happen.".
"I just had last supper with you.".
"This has to happen.".
Look, Jesus followed the soldier..
"Arrest.".
He arrested him..
We are in the modern world..
In the modern world,.
what is "arrest"?.
Right?.
Is it tying his hands?.
Or using a rope?.
Or carrying him out?.
No, no, no..
The Bible says,.
Jesus followed them..
He allowed them to arrest him..
Because if Jesus didn't allow it,.

$^{281}$how could he do it?.
He could cut off his ear.
and touch it..
He could do it..
But if he tied his hands,.
tied his hands,.
he couldn't tie him up..
Big point here..
Right?.
Hmm..
I wonder..
What are these two soldiers.
thinking?.
From their perspective,.
what's happening?.
What are they thinking of?.
From their perspective,.
they are thinking,.
"Hey, we are here to arrest you.".
"By accident,.
someone cut off your ear.".
"The blood is flowing.".
"But I want to arrest that person.".
"Wow,.
he is a fake.".
"He is not a real son of God.".
"Why did he touch my ear.
and it healed?".
Hmm..
We said.
we should think from perspective..
Right?.
I will give you some time..
Let's talk about it..
This soldier.
was cut off his ear..
He was touched by Jesus..
His ear healed..
But what about this soldier?.
Remember,.

$^{321}$this story is about.
perspective..
Angle..
Point of view..
From different angles,.
we see the same thing..
Just like the game..
We will understand.
the whole story.
more completely..
The Bible did not say.
how this soldier.
saw the story..
But we can think about it..
We can guess..
We can imagine..
Right?.
I will pause the video.
with this question..
If you are.
this soldier,.
you see so many things..
What is your perspective?.
What is your story?.
What will you think?.
What will you say to your wife.
when you get home?.
What will you say to your children.
when you get home?.
What will you say to yourself?.
Pause. Discuss..
I told you to pause..
Why did you pause?.
Pause. Discuss..
What will this soldier see?.
What angle?.
Perspective..
He has different perspective from Peter..
Different perspective from Judas..
Right?.

$^{361}$I expect you paused..
I will discuss with you again..
Here we go. Let's continue the story..
Do you remember who said,.
"Jesus, where are you going?.
I will follow you.".
Do you remember who said that?.
Yes. Peter..
And do you remember.
who said that?.
When we first started talking about the story,.
the story of Last Supper,.
Jesus asked this soldier to set up a table.
for the people to eat..
Who was that?.
John..
John is often.
the one.
who loves the most.
in the Bible..
The one who Jesus loved..
That was John..
What will they do?.
After all these things happened,.
what will they do?.
What will they do?.
Can you see where they are?.
Look!.
There is John..
He is the youngest.
among the people..
Youngest..
Who is this?.
Peter is holding a sword..
They are following him..
Right?.
His beloved Jesus was taken away..
But what are they forgetting?.
What is he forgetting?.
Who is this?.

$^{401}$Jesus!.
Jesus was taken away.
by them..
Do you think Jesus will be in trouble?.
Do you think.
what will John and Peter.
do next?.
What are they thinking now?.
Let's talk about this.
when we debrief..
This is the end of the story..
We will continue next time..
We just finished.
the story of the Bible..
Today's topic is.
perspective..
Point of view..
Now, let's think about.
the story just now..
If we look at it from different perspectives,.
different perspectives,.
different point of view,.
from different perspectives,.
there are many different angles..
Do you remember our game?.
We have many different perspectives.
to build a comprehensive story..
If we don't have a comprehensive story.
or we can't see it in detail,.
we are not comprehensive enough..
Right?.
In the Bible,.
when Jesus was taken away,.
when Jesus prayed.
in the garden,.
there are many different accounts.
in the Gospel..
In the Gospel,.
we see.
there is Luke,.

$^{441}$John, Matthew, Mark..
But if we.
can use our brain.
to think from different angles,.
what will we think?.
Although.
it is not necessarily mentioned in the Bible.
that John has a perspective,.
because John is.
an author in the Gospel..
He is a writer..
We can also see.
a little bit of Jesus' perspective.
because he is in the garden..
Right?.
We can imagine Peter's perspective..
We can also imagine.
Judas' perspective..
Let's take some time.
to think about it together..
First, we have.
Jesus..
So, pretend..
假設你是Jesus.
Hello, my name is Jesus..
現在發生這件事,你會想什麼?.
If you were Jesus,.
what would your story be like?.
如果你是耶穌的時候,.
你的故事會怎樣變?.
你的角度會怎樣?.
從Jesus的角度.
會怎樣看這個故事?.
Pause.
想一想.
再回來,我告訴你.
One version of.
耶穌可能會有.
這樣的角度來看.
第一個.

$^{481}$If I were Jesus.
記得我們看到他在.
花園裡面嗎?.
他還在祈禱.
他的角度就是.
A bit of.
什麼?.
Exhausted.
是不是? 很累喔.
神啊.
如果你這個計劃.
可以沒有這一部分.
可不可以呢?.
這麼辛苦,可不可以呢?.
其實,耶穌.
會不會真的很累?.
他做了這麼多神蹟,救了這麼多人.
但也有人.
要追殺他.
還有門徒好像.
怎麼說都不明白.
叫他不要睡覺.
祈禱.
因為那些人會來抓我.
要保護自己.
祈禱完,發生什麼事?.
睡不著.
他的角度,耶穌的角度.
會怎麼說這個故事呢?.
我們想起.
Judas,即是出賣耶穌的那位.
用錢來交換.
告訴別人.
耶穌是誰的那位Judas.
我們未必很清楚.
當時的那一刻.
出賣耶穌.
出賣耶穌之前,他在想什麼.
但我們可以幻想.
我們可以想想.

$^{521}$如果假設我們是Judas.
我們會在想什麼呢?.
我就帶他們去找耶穌.
我親那口,那個男人.
那個就是耶穌.
就是這麼多.
哇!.
剛剛親完他.
彼得就去.
把那個男人的耳朵切下來.
不知道是不是很壞呢?.
耶穌還救回那個soldier.
還弄好他的耳朵.
覺得很矛盾.
為了這麼少錢.
出賣耶穌.
好像很壞.
我們再想想.
如果從另一個角度來看.
另一個角度.
第二個觀點.
彼得.
彼得就是.
把耳朵切下來的那個.
記得他嗎?.
就是那個.
耶穌,你去到邊,我跟你去邊.
耶穌,你死,我跟你去死.
我永遠是你的朋友.
我不會否認你的.
I won't deny you.
我不會否認你的.
到這個地步.
如果彼得要說這個故事.
他會怎麼說呢?.
你們休息一下.
你們大家談吧.
如果我是彼得.
我怎麼會說這個故事呢?.
我怎麼重複給人聽呢?.

$^{561}$If I were Peter.
How would I tell the story?.
是啊.
耶穌就祈禱完.
找我們滿堂.
但是很失望,我們又睡著了.
他明明叫我們不要睡.
叫我們祈禱.
但是,期期下,真的很累.
我都不以為我會睡覺.
由這麼多人之中.
我當然要為了耶穌祈禱.
當然要為了大眾祈禱.
但是,又不是睡著了.
之後,哇,大件事了.
朱迪斯來.
親了耶穌一口.
帶了一些不知道什麼軍人來.
抓耶穌.
哇,真是大件事.
但是呢.
我呢,就不會.
站著不做事的.
生氣得我.
拿著劍.
斬了那個軍人.
蘇州的議席下來.
怎麼樣?生氣吧.
這樣對待耶穌.
耶穌做了多少好事.
耶穌愛多少人.
這樣對待他,不得了.
不行,我是不會容許的.
對了.
又是衰衝動.
又是衝動.
耶穌就說,叫我放下劍.
這樣,很厲害.
他一碰那個人的耳朵.
哇,明明在那裡流血.

$^{601}$這樣碰一碰,好像.
魔術.
但是,耶穌不是魔術師.
耶穌是神的兒子,就是神.
當然可以治好,對吧?.
哇,治好了,叫我放下劍.
那我當然.
聽聽話,放下劍.
但是,我還是不服.
我也要跟蹤他們.
看他帶耶穌去哪裡.
我會永遠跟著他.
永遠跟著耶穌.
看他帶他去哪裡.
OK,現在有很多.
不同的角度看,對吧?.
剩下誰?.
John,是不是剩下約翰?.
對了,John就是.
剛開始的故事.
他幫人.
設置桌子,對吧?.
耶穌叫他擺好位置.
我們要吃 Passover.
乖乖地,你知道John.
是全部.
門徒之中最小的.
經常John寫.
福音的時候.
他說.
Jesus is most loved, Jesus is beloved.
就是耶穌最疼的那位.
其實我們看到他的故事.
玩到.
Love, and follow,對吧?.
那如果他.
故事怎麼說呢?我們其實知道.
因為聖經裡面.
有John的帳號.
我們可以讀,但是假設.

$^{641}$如果你是John.
你會怎麼看呢?.
關一鈊,你們聊聊.
回來.
我給你我的版本.
我的版本給你聽.
但是我們有聖經可以看.
可以看John.
是有哪些細節有.
哪些細節沒有.
如果我是John.
If I were John.
我會怎麼說這個故事呢?.
OK.
耶穌祈禱完就出來.
我就.
起了身.
我已經睡了.
然後很多事發生.
明明.
要跟隨耶穌.
去很多不同地方.
為什麼有這麼多事發生?.
吃.
吃宵夜之前.
祂就每個人幫我們洗腳.
哇.
祂真的很疼我們.
之後又說了很多事.
說祂要走.
說我們要吃祂的身體.
要吃麵包.
喝酒.
喝酒要吃祂的血.
說著說著.
都好像重複很多不同的東西.
但是又不是太清楚.
祂想說什麼.
祂出來祈禱.
不知道祂想說什麼.

$^{681}$貿貿然.
來了一個Judas.
Judas吃了他一口.
彼得又.
割了人家耳朵.
耶穌又救了他.
又治好他.
然後耶穌就.
跟著那兩個軍人走.
只是這樣走.
跟著他們走.
我就很愛耶穌.
耶穌那麼疼我.
我當然要跟蹤他們.
看他們帶我和耶穌去哪裡.
我一定會在他們旁邊.
他們那麼疼我.
我一定會聽他們的話.
他們一定有東西要告訴我們.
知道是什麼來的.
因為.
我們那麼愛他們.
他們那麼愛我們.
一定要跟蹤他們.
慢慢小心地跟蹤他們.
OK.
我們回想起來.
有那麼多不同的角度的時候.
你覺不覺得這個故事.
這樣看是比較全面一點呢.
我們不會知道.
為什麼有某些部分.
聖經裡面沒有說.
為什麼某些角度聖經沒有說.
但是我們作為人.
我們作為想到東西的人.
感覺到東西的人.
其實可以幻想的.
我們讀其他故事書也是.
是不是.

$^{721}$如果讀Three Little Pigs.
那我說.
如果一個狼說一個故事.
會怎麼說一個故事呢.
那變了是你可以.
不同的角度去看一個故事.
好像我們玩的遊戲一樣.
你反轉一個故事.
來參考,來了解.
反轉看一個比較深入的細節.
你會看到些什麼呢.
看到很多不同的東西.
我們可以求神.
祈禱說神啊.
幫我們看到這個故事.
雖然很多次聽過這個故事.
但是如果用不同的角度來看.
可以.
多認識一些東西.
可以了解你多少呢.
求你幫我們.
謝謝耶穌.
故事怎麼繼續.
耶穌跟了那兩個軍人去哪.
之後會發生什麼事.
哇.
下一集再說.
OK 大家準備好了嗎.
中文時間.
我整個episode.
就講角度.
Perspective.
Angle.
其實中文角度.
即是甚麼意思呢.
其實就是.
角度.
中文角度即是甚麼意思呢.
角.
角度即是角.

$^{761}$Corner.
度.
角度.
你可以說.
一個triangle.
三角形.
三.
一二三.
三個角.
Three corners.
三個angle.
角度.
角度和.
角.
同樣的.
周圍看看.
你周圍的東西.
哪些有角呢.
甚麼東西有角呢.
我就叫我兒子去找.
你們可以.
拍住他.
OK.
我叫我兒子去找東西有角.
甚麼東西有角呢.
你告訴別人.
你找到這本書的角在哪.
這裡.
即是corner.
這裡有角.
這本書的形狀.
是長方形.
長方形有多少個角呢.
How many.
四個.
你數吧.
一二三.
四.
四個角.
角又可以想corner.

$^{801}$角又可以找角度.
angle.
或者perspective.
在我的可以.
例子就是在我的角度.
來看.
我覺得.
應該日日吃雪糕.
So from my point of view.
From my angle.
I think everyday we should eat ice cream.
Who thumbs up.
You agree.
但或者媽咪會講.
唔係喎.
從媽咪的角度來看.
日日食雪糕.
唔健康喎.
媽咪says wait.
From a mom's point of view.
Eating ice cream everyday.
Isn't healthy.
Or.
Agree.
Disagree.
Or so so.
我們可以有唔同個角度.
Let's see.
What else is a 角.
Corner.
Street corner.
街角.
So when you say I'll meet you at the corner.
我係街角等你啦.
Corner.
角.
Right.
And then when you go to the corner.
Stand in the corner.
Or put it at the corner.

$^{841}$角一頭都係有個角字.
Challenge.
If.
如果三角形.
Triangle.
係叫三角形.
Hexagon.
係乜ye .
中文hexagon.
係點講.
下次你答我.
See you later.
Bye.
大家而家到返.
At home extension.
係屋企點樣可以深入D .
講呢個故事俾小朋友聽呢.
我就呼籲大家.
lor 返大型ge 聖經.
Your Bible.
去返John 18.
約翰第十八章.
讀返gwo 度係講返Jesus is Arrest.
之前都可以睇返.
Last Supper gwo D .
我就想你讀俾小朋友聽.
或者一齊讀都好.
讀完就想問你地.
睇下.
Miss Rebecca 做gwo 個.
edition, gwo 個 version.
Of John ge 角度.
同其實神想我地讀John.
ge 角度有D 咩出入.
有D 咩唔同呢.
點解.
One, Teacher Rebecca 係幻想.
但係係.
聖經裡面gwo 個 perspective.
gwo 個角度 from John.

$^{881}$咁就我地知.
亦都可以講俾小朋友聽.
其實福音, Gospel 有四個 account.
係咪?.
都係可以講緊 Last Supper.
都係可以講緊 Jesus is Arrest.
都係可以講緊 Easter Story.
點解要四個呢?.
有.
Matthew, Mark, Luke 同埋 John.
一個咪夠囉.
點解要咁多個?.
la , The Bible 唔會錯架嘛.
係咪?.
連 Easter Story ge  Gospel 都有四個 account.
證明.
神都想我地睇唔同個角度ge ye .
係咪?.
抽出唔同ye 出黎.
係一個 story 都可以抽出唔同ge ye .
變jor 點呀?.
gwo 個 story 會有多D 意思, more meaning.
gwo 個 story 會豐富D , richer.
[Music].
\newpage



\section{10:38-42}
\label{sec:yncumdAd8iU}
\textbf{Choosing What is Better (Luke 10\_38-42) - Rev. Sam Johnson}
\newline
\newline
連結: \href{https://youtube.com/watch?v=yncumdAd8iU}{\texttt{https://youtube.com/watch?v=yncumdAd8iU}} ~~~~ 語音日期: 2022-03-14
\newline
\newline
\hyperref[sec:hkH9QTCOPtw]{\small{< < < PREV SERMON < < <}}
~
\hyperref[sec:index_chronic]{\small{[返順時目]}}
~
\hyperref[sec:index_scriptual]{\small{[返順卷目]}}
~
\hyperref[sec:NyL1j_DJIlI]{\small{> > > NEXT SERMON > > >}}
\newline
\newline
$^{1}$Thank you..
Well, it's great to be here today..
I'm very excited to be here..
And love being at UV..
It's been a little bit since I've been here in person, but it's great to see you all and.
to see family..
It's good..
Spent a lot of time with a lot of you that are here, and it's great just to be around.
family again and appreciate you being here today and enjoy the worship with you all and.
just excited to be here..
We're going to be looking in Luke chapter 10 today at this story of Mary and Martha..
We're going to talk about choosing what is better..
I'm going to pray real quickly as we begin, and then we'll start to look at this passage..
God, thank you so much for the opportunity to be here today..
Thank you for your word..
God, thank you for how it changes us, how it forms us and shapes us..
And I pray that would take place today, that as we consider your word together, as we read.
it, as we talk about it, that it would change who we are, that when we leave, we would be.
different than when we came in and that ultimately we'd be closer to you, God, in Jesus name..
Amen..
So this passage, just a few verses we're going to look at in verses 38 to 42 of Luke chapter.
10, a familiar story, but I think a lot of powerful truth for us in our lives about what.
we choose to do with our time and our resources that God has given us..
So let's read this together..
Luke 10, 38 to 42, it says this, "As Jesus and his disciples were on their way, he came.
to a village where a woman named Martha opened her home to him..
She had a sister called Mary who sat at the Lord's feet listening to what he said, but.
Martha was distracted by all the preparations that had to be made..
She came to him and asked, 'Lord, don't you care that my sister has left me to do the.
work by myself?.
Tell her to help me.'.
Martha, Martha, the Lord answered, 'You're worried and upset about many things, but few.
things are needed or indeed only one..
Mary has chosen what is better and it will not be taken away from her.'".
And that's what we're going to consider today, choosing what is better..
This is a really powerful command that Jesus lays on us in this passage, but it's a little.
bit perplexing to us too, and we're drawn to the power of what Jesus is saying here..
We want this..
We want to choose what is better, yet sometimes we're confused about how to do so..
What does it mean to choose better?.

$^{41}$What does better even mean?.
And I think that question is even more significant for us in the modern world because we live.
in a world that there are so many choices and we live in a society where there's a lot.
of affluence and there's so many options in every place that we look in life..
And whatever we're trying to choose in that moment, there are just kind of a myriad of.
choices for us to choose from..
If you think about something like when you get a new car..
Now I've never bought a new car..
I've only bought used cars, so there's less choices when you buy used..
But sometimes when I'm dreaming or I don't have anything else to do, I have gotten online.
and looked at, you look at Tesla or whatever and you go through, build your own and customize.
and all that..
And if you've ever done that or whatever your favorite car is, if you ever looked at a Ford.
Bronco online or a Corvette or whatever, there's so many choices that it takes forever just.
to get through all the choices because obviously you have to pick the car that you want, the.
body style, do you want an SUV, a van, a sedan or whatever?.
And then after that, you've got to pick the engine, the drive train, then you're going.
to pick the seats and the fabric and the color and what kind of radio you want and all the.
different technology options..
It's just kind of endless amount of choices to choose something like that..
And that is just kind of a silly illustration, but in every era of our life, we're faced.
with just kind of an overwhelming amount of choice..
When it comes time to eat, there's so many things to choose from when we're going to.
eat, right?.
Are we going to go to a restaurant?.
Are we going to get takeout somewhere?.
And if we are going to go to a restaurant, which one are we going to choose?.
Because the whole city is full of great restaurants..
If we're going to cook at home, what are we going to cook?.
There's so many different ingredients and recipes and options..
There's all these different choices..
There's ready-made meals, there's things at the grocery store, what are we going to do?.
When we think about how we're going to spend our time and entertainment and things like.
that, there's sporting events, you could go to the movies, you could go to the park, you.
could go to a concert..
There's so many choices when you have a free night that you're going to go do something..
Or even if you're just going to sit down and veg out on the couch and watch something on.
TV, there's limitless choices, right?.
If you have cable, you could channel surf for hours..
I don't have cable, but we have some streaming and streaming is the same way, right?.

$^{81}$That's the hardest part of watching a movie or a show on streaming is actually getting.
started because there's so many things you're like, "Ah, maybe we should watch this..
I want to check this out..
I heard this was good.".
When I was a teenager, we had this thing..
Some of you know what this is, but some of you, this is totally foreign to you..
When I was a teenager, we had this thing called Video Store..
It was a store with shelves like other stores, and on these shelves were a whole lot of DVDs..
Some of you might know what a DVD is..
It was just literally shelves and shelves of DVDs..
My friends and I, when we were in high school, we would do the same thing that everybody.
does with streaming now..
We'd go in there, spend like 45 minutes walking around the Video Store, picking up movies,.
being like, "Hey, do you want to watch this?".
"No, I've already seen that.".
Or, "No, I heard my friend said it wasn't any good.".
It took so long just to choose..
We still have that problem now..
Like I said, we get on Netflix or Disney+ or whatever, streaming, and there's just so.
many choices..
And so, what Jesus is saying here is actually really significant to us..
And the reality is, there's so many choices in our lives that at some points, it almost.
becomes like an existential threat to us..
We feel overwhelmed because we have so many choices to make..
We feel like if we make the wrong choice, we're going to miss out or we're going to.
wreck our whole life..
I remember when I was trying to choose a college feeling that way..
I was like, "There's so many colleges in the United States..
How could I possibly know which one is going to be best for me?".
And it almost seemed like, you know, if I don't make the right choice, the world might.
come to an end..
And we feel that way with a lot of our choices..
So what Jesus is saying here, that we need to choose what is better, is really a significant.
and powerful truth to us..
But it's difficult for us to know what He means..
So let's look at what's going on in this passage and begin to discover what Jesus is communicating.
to us..
First of all, we see in this passage, there's this contrast between the two sisters..
They each have a different response to Jesus in His presence..
He shows up at their house and they respond differently..

$^{121}$But the response, the contrast here is actually more than just their behavior..
It's really a contrast in their attitudes and in the way that they're treating Jesus..
Sometimes I think we could look at this and see it simply as a contrast between, well,.
it's a contrast between a contemplative life and a life of work..
And it's not really that..
It's not just a story about, "Don't be too busy and forget to spend time with people.
you love.".
It's not that..
I mean, those things are important, but the story is more than that..
It's really about the hard attitudes of Mary and Martha in this story..
It's about the posture that they take as they respond to Jesus..
And what's happening here, Jesus and His disciples are traveling and they come to this village.
where Mary and Martha live..
And the passage tells us that Martha invites Jesus and the disciples into her home..
Now that's good..
That's a really beautiful thing..
See, Jesus and His disciples, remember, they were itinerant ministers..
They didn't have a bunch of money..
They didn't have, Jesus didn't have a home of His own..
So Martha was being generous and extending hospitality, inviting them into her home..
And we can assume from the context here, Martha's probably a pretty well-off woman..
Since she has the space and the resources to welcome visitors into her home, she probably.
is fairly affluent..
And her name in Aramaic actually means ruling lady..
So she was probably a significant person, a person of influence..
And she's probably most likely, we see this in the story as well, that she's probably.
the one who kind of makes the decisions in the home and makes sure things get done..
She's doing something praiseworthy and good here by welcoming Jesus and His disciples..
She's showing hospitality, generosity..
Those are biblical values..
That's really good stuff..
But the word, the Greek word in verse 40 that is used to describe what she's doing is the.
word diakonia, which is a word that we get the word deacon from..
It's a word that gives this connotation of service or ministry..
So the point that I'm making here is that Martha's actually doing something good..
She's ministering, she's serving, she's doing something good..
The problem is not with what she's doing..
Because she's doing something that she should be doing..
She's a respectable woman, she's ministering..
But the problem is that she's missing the best thing..

$^{161}$She's doing something good, but even in doing something good, she's missing the best thing..
And that's the point that Jesus is making..
She's missing the one thing that is truly important, or as Jesus says in this passage,.
that is truly needed..
Now Mary, on the other hand, here's the contrast..
Mary is sitting at the feet of Jesus, and she's taking the posture of a disciple, a.
learner, or a follower..
She's responding to God's presence with faith..
She has joyful anticipation..
If you think about what it was like for her to sit there and listen to Jesus, I'm sure.
she's leaning in and desiring to hear every word that he's saying..
And she responds with a personal affection..
It's very close..
And there's trust involved, there's love..
She's able to forget about the other tasks and pressures that she has in life, and just.
focus on being present with Jesus..
She's valuing, this is big for us in 21st century America, she's valuing presence over.
productivity..
That's what she's doing in this moment..
And she's valuing a relationship over simply good works..
Ultimately, she's valuing Jesus..
And her response is very different..
And this is the response that we need in our lives, that Jesus is pointing out in this.
passage and calling us to..
Like Martha, we have a lot of good stuff in our lives, but sometimes we miss out on what.
is best because there's so much good..
And Jesus is calling us to discover the ways that we, like Mary, can choose what is better..
So there's three things I want us to consider in this that will help us to choose what is.
better in life..
And the first is that we have to learn to choose identity in Christ that's grounded.
in God's grace..
We have to choose identity in Christ that's grounded in God's grace..
This is probably, I think, the most critical piece of choosing better..
It's considering where we place our identity as human beings..
If our identity is not in Christ but in anything else, our choices about our priorities and.
our time, our resources, even our happiness, will ultimately be driven by the source of.
our identity..
And this passage, why is Martha so focused on all the stuff she has to do?.
As Jesus says, she's worried about so many things..
Why is she worried about all that?.

$^{201}$We don't know exactly..
The passage doesn't tell us..
But I think we can imagine that it could be, one of the reasons could be, that she, like.
so many of us, is finding identity in the work that she has to do..
You know, perhaps she really likes being the one who makes the decisions and keeps control.
of the household..
Maybe that's an important part of her identity and who she is..
Perhaps she feels a sense of pride in how competent and comprehensive she is in serving.
guests..
You know, her house is well kept..
Everything's clean in its place..
The meal is perfectly delicious that she's putting on the table..
And there's treats for people who come in the door..
And there's cut flowers..
And everything's decorated..
And all the servants in the house are working..
And everything's running like clockwork..
Maybe that's really important to her in terms of her identity..
Maybe she's kind of a Martha Stewart of her time, right?.
And that's important to her, that she knows everybody's looking at her being like, "Wow..
Her house is amazing..
And everything smells great..
And it's all clean..
And it runs perfectly.".
And maybe that's really a big part to her of who she is..
And she feels the pressure of maintaining that facade of perfection in front of other.
people..
Because she's like, "Everybody's expecting me to do everything and for it to be great.
and perfect..
And I have to keep up with that appearance..
I have to keep up that facade..
I have to keep up that image.".
And she's feeling that pressure..
Because ultimately, her identity is in the work that she's doing, not in Christ..
Maybe she's feeling a sense of significance because she's doing all this important stuff..
And if she wasn't doing those things, maybe she would feel kind of empty and wonder who.
she was, where her worth was found..
That's so true for us as well..
And I think that's so significant for us to consider..
Maybe she felt like all the good stuff that she was doing somehow was earning her favor.

$^{241}$with God..
Maybe she thought that what she was doing helped demonstrate to other people and to.
God that she was a good person because she was doing all this good stuff..
And maybe Martha isn't choosing what is better..
She's not choosing to sit at the feet of Jesus to be present with Him..
Because she feels like she wouldn't know who she was or other people wouldn't know who.
she was without all this activity that she's doing..
I think we have to examine our hearts because it's so easy for us to do this as well..
How often do we place our identity as human beings in stuff other than Christ?.
In our work, we certainly, it's easy to do, right?.
A lot of us as Americans, when people ask us, when we meet someone for the first time.
and they ask us about ourselves, we start describing our work like that's who we are..
That's not who we are as human beings..
The job that we do..
That's not our identity..
But we place our identity in that..
Sometimes it's in our education..
We feel like our education defines who we are..
The success or sometimes we, if we don't experience the success we want in education, we feel.
defined by failure..
That's not who we are..
That's not our identity..
Our accomplishments, our skill in different things and different parts of life, our competency,.
our success, our relationships..
Sometimes we feel like, oh, I have this identity because of who I know or who my friends are,.
who my family is..
Our life experiences, our travel, our money, our prestige, our social media feeds, our.
Instagram, those things, if we're not careful, we start to define ourselves and think that.
that is who we are..
That's our identity..
And that's not true..
Even good things in our life, like our family heritage or our religious faithfulness and.
church attendance, things like that, that are good things that God wants us to do, we.
can start to put our identity there and say, yeah, I know I'm a good person because I always.
show up at church and I try to help some other people and every once in a while I give some.
money or whatever..
But that's not who we are..
Our identity has to be found in Christ..
You know, if we try to pull meaning and significance and worth from these things, then we start.
to feel good when they're going well, when our work is going well or when our something.

$^{281}$good is happening in our life, when we get to take the trip that we want to take or our.
school is going well..
We feel good..
We feel like, yeah, everything's great..
But then when they're not going well, we feel worthless..
We feel like who we are as a human being has been destroyed or crushed in some way..
We feel miserable..
And it's because we're placing our identity in something that it shouldn't be in..
And all the while, those things that we're placing our identity in begin to take more.
and more of our time and focus and more and more of our priority in our life..
They start to consume us and they become our priorities..
Even if we don't plan for them to, they start to take up all of our time..
And I want to just point out again that a lot of those things I'm talking about are.
not bad things..
Not saying that..
Obviously, we know if we're putting our identity in destructive, unhealthy, sinful things,.
then that's wrong..
We shouldn't do that..
But that's not what I'm talking about..
I'm talking about good stuff that God actually wants us to have in our life..
It's just we get it out of proportion, out of focus, and it becomes bigger than it should.
be..
And that's what was happening with Martha..
As I said earlier, she was doing good stuff..
She just wasn't doing the best thing..
She wasn't spending time in the presence of Jesus..
And that happens to us too..
It's good to work hard..
It's good to be successful and competent..
It's good to be faithful in church..
Those are all biblical values..
There's nothing wrong with seeking those things..
The problem is when we begin to put those things in the wrong place..
And to help us keep those things in perspective, we really have to understand the grace of.
God..
We live righteously in life because Jesus has forgiven us and saved us and is at work.
in our lives..
We do not live righteously so that Jesus will forgive us and save us and be at work in our.
lives..
We have to keep that front and center..

$^{321}$Like when good works come out of us, it's because God is at work..
It's not because good works are coming out of us so that we can earn anything or so that.
we can achieve something with God or so that we can impress other people..
And when we put our identity in something other than Christ, whether we realize it or.
not, we're seeking to derive something from those things that God never intended for us.
to get out of those things..
And trust me, nothing in the world can support the weight of human identity..
It's a crushing weight for us to try to find our worth in something..
And Christ is the only thing that can support that weight..
If we try to pull worth and identity out of our job, we'll ultimately be disappointed,.
right?.
Because our jobs can be lost..
If we try to pull that out of a relationship, that relationship could change..
That friend or that person in our life may not be there anymore..
Maybe they move away..
That can evaporate..
Our success, our education, our prestige, those things can go away..
People can forget about our accomplishments..
Maybe our degree or our GPA won't mean that much to anybody..
Those things were never meant to hold the weight of our identity..
But Christ can..
And here's the truth of God's Word, and this is the important part, I think, for us to.
remember in terms of identity, that we are made in the image of God, that each one of.
you has unspeakable value and worth because you're made in God's image..
And by God's grace, if you have relationship with Jesus Christ, you've been redeemed by.
the sacrifice of Jesus on the cross..
All of your sin and brokenness and all the mistakes and all the mess ups and all of the.
stuff in your life that has hurt you can be forgiven and healed and restored..
That's the reality of your identity..
But even in that, you don't bring anything to offer God..
It's a gift of grace that Jesus does..
We don't have to earn anything..
We don't have to work for it..
We're forgiven and reconciled to God because of His love for us..
We're restored to relationship with our Creator, the God who made us and loves us because of.
His grace..
He does all the work, and we receive all the reward..
And that's the good news of the Gospel, that we don't have to do anything, but we're healed.
by His love and grace and mercy..
And that's your identity..

$^{361}$If you're in Christ today, that is your identity..
You are forgiven and healed and restored and made new and covered and protected and.
loved..
That's who you are..
Not because of anything you do..
Not because of anything that you have accomplished or will accomplish..
Because of God's grace and work in your life..
And so you can stop striving..
You can stop trying to deserve it..
You can stop trying to impress God or other people, and you are free to go sit in the.
presence of Jesus and enjoy His presence..
That's where Martha was at..
There was something that was, she felt like she couldn't go in there and sit at His feet.
because she had to do all this stuff..
She wasn't free to say, "Yeah, I can just sit in the presence of Jesus.".
And we can feel that tension as well..
You are free in God's grace..
If you find your identity in Christ, to say, "You know what?.
I can go sit in the presence of Jesus, and I don't have to keep up any appearance or.
any image..
I don't have to earn anything..
I don't have to try to be good enough..
Because He's loved me..
He's forgiven me.".
That freedom allows us to choose what is better..
The next thing I want us to see here is that we have to choose provision and peace in Christ.
that is grounded in God's goodness..
We have to choose provision and peace in Christ..
That's grounded in God's goodness..
Maybe another reason that Martha is working so hard is because she's preoccupied with.
getting things done because she feels the pressure of living in our broken world where.
there's a lot of problems and there's stuff that has to be taken care of..
Maybe that's why she can't go and sit at the feet of Jesus because she's like, "Hey, there's.
bills to pay, and there's stuff that's got to get done in this house..
And if I'm not doing it, who's going to do it?".
Maybe there were debts..
We don't know..
Or maybe she's just saying, "Hey, all the comforts of this home that Jesus is enjoying.
and the fact that Mary can sit in there and relax at His feet, that only happens because.
somebody's working, and it's me..

$^{401}$So I got to keep working.".
Maybe she's feeling that pressure of the fact that we live in a broken world and nothing.
comes easy, and so she just has to keep laboring..
And that's true..
Those are very real concerns..
I get that..
All of us have those types of concerns..
We face those issues..
There are financial pressures..
There are responsibilities to care for people in our families, in our churches, in our communities..
The list of things we have to do and the bills we have to pay kind of seems to go on forever..
And at times we can feel like we simply don't have the time to sit down and spend in God's.
presence because there's so much that has to be done..
But what do we see about Mary in this passage?.
She does make time..
Now is Mary sitting at the feet of Jesus because she's lazy and unconcerned?.
I don't think so..
I don't think there's anything in this passage or anywhere else in the Gospels where it talks.
about Mary that would cause us to believe that Mary's just lazy or she doesn't care..
That's not true at all..
It's not that she has no responsibilities..
I'm sure she had a lot of stuff she needed to do too..
It's not because her life was just full of leisure and she could just sit around..
It's not that..
What's going on with Mary is that she has all the tasks and the responsibilities..
She has stuff that she needs to do..
She has pressures, but she's just saying, "No, I'm going to set that stuff aside and.
go sit at the feet of Jesus..
And I'm not going to worry about them because I'm going to trust that if I'll make time.
for Jesus and make that important in my life, then he'll help take care of all the other.
stuff.".
She's finding the peace to sit in his presence by trusting that he's going to provide for.
the rest of the pressures of her life..
And that's really important for us as well..
Martha doesn't really understand this..
We see when she comes in, she says, "Lord, don't you care that my sister has left me.
to do the work by myself?.
Tell her to help me.".
It's kind of like Martha saying, "How can she just sit here and listen to you when there's.
so much that has to be done?".

$^{441}$But Mary's not sitting there because she doesn't care..
She's not sitting there because she's being spiteful towards Martha..
She's sitting there because she's trusting Jesus and she's saying, "This is really important.
in my life..
This is the most important thing..
And I'm believing that if I make time for this and if I make sure my relationship to.
Jesus is healthy and I'm spiritually healthy, then God's going to take care of the rest.
of it..
I'm going to be able to get through the rest of it a whole lot better.".
Doesn't mean her life was easier, that she didn't have problems..
That's true for us..
Just because we have a relationship with God doesn't mean God's going to make all our problems.
go away or we don't have to pay bills..
We have the same challenges everybody else has..
But if we will choose to find provision and peace as we trust God's goodness in our life,.
then we'll find ourselves saying, "You know what?.
I am busy right now..
I have a lot to do..
But I can take a minute to sit in God's presence and be encouraged and be blessed by that.".
And that's what Mary is doing..
And one of the cool things about that is as we do that, as we spend time in Jesus' presence,.
we actually start to get perspective on the problems in our life that they're not as big.
as we thought anyway..
Remember there's this old hymn that says, "Turn your eyes on Jesus..
Look full in His wonderful face, and the things of earth will grow strangely dim in the light.
of His glory and grace.".
That's really true..
When we're focused on the problems and the challenges in our life, they seem huge and.
insurmountable..
But when we start doing what Mary was doing and sitting at the feet of Jesus and focusing.
on Him, spending time in His presence, suddenly the problems are not that big of a deal..
They start to fade into the background..
And the third thing we see here is that Mary is choosing satisfaction in God's presence..
You know, sometimes we work and we do all this stuff because it just feels good, right?.
Maybe some of us have painful things in our life, trauma from our past that it's tough.
for us to deal with, and so it's just easier to work and to stay busy and to put our head.
down and to try to distract ourselves..
Maybe that was going on with Martha a little bit too..
Maybe she was discouraged from things in her life, but it was just easier to push ahead.
into her work than to actually deal with any of that..

$^{481}$And that happens to us..
The reality is that we all ache for better things in our world, because we live in a.
broken world where there's all types of pain around us, and it comes into our lives as.
well..
And God created us for a perfect world..
God created us for perfect relationship with Him, for perfect relationship with one another,.
for Heaven..
And so there's a part of us that looks at our broken world and is just, we're hurt by.
it, we're frustrated, we're angry at times, we don't understand it..
Our hearts break over it..
And because of that heartbreak, we start to seek satisfaction in different things..
There's a lot of unhealthy ways that we seek that, but sometimes we seek it just in work.
and business and trying to order our lives and our families and our future..
And we feel like, if I can just kind of get everything in order the way I want, then it'll.
feel good..
I'll feel satisfied if I can just order my own life..
And the sad reality is we're always let down in that, because none of us, no matter how.
OCD we are, none of us can control everything and order our own life and make it all work.
out the way we want it to..
God's always going to disappoint us or get sideways or let us down or break..
And we have to understand that we should be finding satisfaction in God's presence..
There's a reason that God, He did create us for good works, but as I said earlier, those.
good works are a result of what He's doing in us..
God created us not as human doings, but human beings..
God wants us to be with Him more than He wants us to work for Him..
And sometimes we feel like, okay, I want to follow Jesus..
It's all about the stuff I do..
And God says, "Hey, yeah, I've got some stuff for you to do, but first I want you to be.
with me..
I want you to spend time in my presence.".
And that's what Mary was choosing to do..
Mary is finding satisfaction in the presence of Jesus, sitting at His feet..
She is ultimately choosing, as I said earlier, His presence over productivity, over trying.
to do everything..
She's saying, "I'm going to be most satisfied just spending time with Christ.".
Mary is choosing relationship with Jesus as a disciple, a follower, and a friend..
She's choosing not to be distracted..
She's not putting her identity in her activity..
She's trusting Jesus rather than worrying about the pressures around her..
And look at how she's relating to Him..

$^{521}$She's taking the posture of a worshiper, sitting at His feet, welcoming His presence, listening,.
learning from what He's saying, being humble and attentive to His words, opening her heart.
to His truth..
She's giving the response of a follower..
I'm sure she had questions and fears and misgivings, but even in the midst of that, she's choosing.
to obey..
She's being a disciple..
She's leaving other priorities to be devoted to the time in God's presence..
She's choosing the response of a friend..
She's sitting, talking, laughing, just like a friend would, sharing time together..
She's experiencing the love of Christ and responding with love..
And she's choosing what is better..
So I think the question for us then is how do we make that choice?.
We want to choose what is better..
How do we make that choice?.
It seems kind of simple for Mary because Jesus was there, like, in the room..
So you know, if Jesus was right in the other room, we could all say, "All right, let's.
go in there and have a seat.".
That's a simple choice, right?.
But in our lives, Jesus is present, but He's not physically present in the other room..
So how do we do what Mary did and choose what is better?.
Well, it's really not as complex as it might seem to us..
Ultimately, we choose what is better by responding to the presence of Christ, just like Mary.
did, by responding to His presence in our life..
Jesus showed up at Mary and Martha's house..
Mary responded by going into His presence..
And first of all, if you don't have relationship with Christ, you respond to His presence by.
beginning relationship with Him, as we talked about earlier..
That comes by grace, through faith, through believing what He's done on our behalf..
It's not something we have to earn or do anything for..
It comes by grace..
But for those of us who already have relationship with Christ, who say, "Okay, I've already.
made a decision to know Jesus..
I already have relationship with Christ.".
Then Jesus is already present in our lives through the presence of the Holy Spirit..
He's with us every moment..
So how do we respond to His presence?.
I think the answer for us to consider today is that we choose to take time with Him..
We choose to take time sitting at His feet..
Now like I said, His physical presence is not here, but how do we listen to His words?.

$^{561}$Mary sat at His feet and listened to His words..
Well, God's given us the Bible..
It's His words spoken to us..
So if we want to sit at His feet, hear what He has to say to us, we spend time reading.
His word..
And we do it in the way that we see modeled in this story..
We take time..
We do it slowly..
We do it attentively..
We do it with a heart to hear and listen and understand, not just a heart that says, "Yeah,.
I know this is kind of important..
Like I go to church..
I'm supposed to read my Bible..
I'll check it off the list.".
But we say, "No, that's not the point of what I'm doing..
The point is I want to actually hear what God has to say..
I want to be in His presence.".
We ask questions..
We don't rush through it..
We pray..
And when I say pray, I don't just mean like, "God, thank you for this food..
Help me not to break my leg today..
Amen.".
But taking time to really pour out our heart before God..
If you think about a relationship, if there's any closeness in that relationship, there.
has to be communication, right?.
And the closer the relationship is, the more honest and open that communication should.
be..
So if you're going to have a relationship with Jesus, you should be able to talk honestly.
and openly about what's in your life..
And that's what I mean when I say pray..
I mean sitting down with God and saying, "God, here's what's going on..
I'm confused or I'm angry or I'm hurt or I'm frustrated..
I don't understand..
I'm afraid because of what's happening in my life.".
That's what I'm talking about, honestly and fully pouring out your heart before God and.
sharing your struggles and your doubts and your fears..
Opening up your heart and saying, "God, I don't understand this..
I don't know what's going on right now, but I'm asking you to work and I'm going to let.
you work..

$^{601}$I'm going to take my hands off of it a little bit and let you do something in my life.".
Letting God speak to us, listening to what he's going to have to say as we get in his.
word..
We choose to stay and continue in relationship and in prayer, right?.
That's one of the challenges we have as believers..
It's easy to be like, "Okay, God, I prayed once, read a little of the Bible..
Now I want to see the result.".
And it doesn't happen tomorrow and it doesn't happen next week..
And then we're like, "All right, I'm done..
I'm going to go try to do this myself or find another solution.".
But continuing to stay with God and saying, "God, I'm going to continue to let you work.
in my life.".
That's how he speaks to us..
That's how we experience his presence..
Choosing to find our identity in him like we talked about at the beginning, not in all.
the other stuff, but to find our identity in Christ..
Guarding our hearts against putting our identity in something else that's not going to be able.
to support the weight of our identity..
Believing his truth that we're loved by him, we're forgiven, that we can be healed..
Finding our worth in who we are, created in God's image and redeemed by Christ..
And we choose to prioritize our time in that way, right?.
If you have a relationship with someone that you love, that's important to you, you prioritize.
time with that person..
You make sure that it happens..
You put it on the calendar and you don't miss it..
And you ignore other stuff to make it happen, right?.
And that's how we spend time at the feet of Jesus..
That's how we choose what is better, like Mary was saying..
We say, "I'm going to make sure I get time with God in his presence..
I'm going to make sure I get time in the Word, praying, worshiping, spending time with God..
And I'm going to prioritize it and make sure it happens consistently.".
And we do that even when, and I would say especially when, the responsibilities and.
the stress and the pressures of life are piling up..
Because that's when it's easiest to be like, "I just don't have time right now to sit in.
God's presence.".
But that's when we need it the most..
Maybe Mary knew that in this story..
Maybe she was saying, "Hey, I actually have so much to do and that's why even more than.
ever I need to go sit at the feet of Jesus so that I can experience some peace that will.
get me through what's about to happen in my life.".

$^{641}$We've got to value that..
We've got to choose it..
And I think if we do, our relationship to Him will start to flourish in really life-giving.
ways..
And we'll really discover the reward of knowing Him..
And we'll experience His blessing..
And we'll experience a lot of things that seem so threatening to us not being as threatening..
That God will remove those things or they just won't be as big of a deal for us..
What we'll ultimately experience is what Jesus says in this passage, that it will not be.
taken away..
Everything else is going to come and go..
But the time that we spend with God, loving Him, the time that we spend allowing Him to.
work in our lives to love other people, that ripples through eternity, right?.
The bills we pay, even our jobs, our degrees, our cars, our houses, everything else, it'll.
be gone in a few years..
No one will remember it..
But our relationship to God is eternal..
The way that God uses us to love other people, that's eternal..
The souls of other people will be in heaven for eternity or they'll be separated from.
God..
But it's going to be an eternal impact..
And that's what I want us to see as we close..
If you think about Mary and Martha, I want to just kind of humanize this story for a.
second..
It's easy for us to read Bible stories in a kind of two-dimensional way, like, oh yeah,.
they're Bible characters..
They're like people on a felt board that move around..
But these were real people..
They really lived just like we do..
And I'm sure one day when they were older, they probably told people about Jesus, right?.
Because everybody in their culture would have known about Him and wanted to hear stories..
I'm sure they sat around with kids someday and told them, "Hey, this one time Jesus came.
to our house and He was passing through with the disciples and we invited Him in.".
And think about that..
When they would tell those stories some days, what do you think they told about?.
Did Martha say, "Yeah, Jesus and the disciples came over..
It was awesome..
The pitas that I made that day were just, they were perfectly crisp.".
Is that what she was talking about?.
I really doubt it, right?.

$^{681}$Or every candle, I mean the wicks were trimmed at the same height..
That's not the story that she would have told..
She would have talked about Jesus was present in her house, in the room..
Like, this is what He said..
This is what it felt like to be there with Him..
This is how loved I felt in His presence..
This is what He said to me..
This is how He healed something in my life..
This is how I was forgiven by His grace..
That's the stuff that doesn't go away..
That's what Jesus said here..
It will not be taken away..
And that's true for us..
We get so caught up in all this other stuff in life and a lot of it won't be important.
in a few years..
But if we'll choose, like Mary, to sit in the feet of Jesus, to take time in His presence,.
to value presence over productivity and everything else..
That's something that will not be taken away..
We won't lose it..
And we will remember that and we'll share it with others, what it feels like to be in.
His presence, to be forgiven, to be loved, to be accepted, and to be comforted, to be.
healed, to know Him and to grow in Him..
That's what Jesus is calling us to, to choose that in our lives and to sit in His presence,.
to choose relationships, to spend time with Him..
Let's pray..
God, we thank You, again, for Your truth, for how powerful and real it is..
And God, I just pray over all of us in this room, I know it is so easy to do everything.
else except spend time with You..
We have so many pressures, so many things we feel like we have to do, so many things.
we feel like we want to do..
It's hard to choose what is better..
And God, I just ask for Your help in my life and I ask for Your help in each one of our.
lives in this room, that You would help us to choose what is better..
Help us not to find identity in anything other than You..
Help us not to feel like we have to work all the time because there's so many needs..
Help us to trust Your provision, God..
Help us not to find satisfaction in anything else..
Help us to choose to take time to get in Your presence, to stay there, to linger there,.
to experience Your love, Your forgiveness, Your grace, to learn from You, to grow in.
You, to enjoy that time and to continue to consistently practice it throughout our lives..

$^{721}$And God, please grow us in ways that glorify You, that bring us closer to You and also.
make us more able to love the people around us and to love our families and our communities,.
God..
In Jesus' name, amen..
[BLANK AUDIO].
\newpage



\section{創世記}
\label{sec:NyL1j_DJIlI}
\textbf{創世記第4章 (4.4) - 背景釋經 (陳崇基牧師)}
\newline
\newline
連結: \href{https://youtube.com/watch?v=NyL1j_DJIlI}{\texttt{https://youtube.com/watch?v=NyL1j\_DJIlI}} ~~~~ 語音日期: 2022-03-14
\newline
\newline
\hyperref[sec:yncumdAd8iU]{\small{< < < PREV SERMON < < <}}
~
\hyperref[sec:index_chronic]{\small{[返順時目]}}
~
\hyperref[sec:index_scriptual]{\small{[返順卷目]}}
~
\hyperref[sec:5qb8PRS0GKE]{\small{> > > NEXT SERMON > > >}}
\newline
\newline
$^{1}$到第六章 我們看少少.
六章也有很多字眼.
是跟創世記一二三章有出現過.
六章 人在地上開始增多.
Adam 人類在地上開始增多.
Rath 或者是Ruth.
就是multiply的意思.
生兒養女.
神的眾子見到人的女子美麗.
隨意的挑選作為妻子.
神見到人既然是屬肉體.
我的靈就不會永遠住在他裡面.
他們的日子只剩下120年.
這個細節到第六章的時候我們才解釋.
這裡說到那些日子和以後的日子.
那些日子 洪水快要來臨的日子.
以後的日子有什麼.
有巨人在地上.
什麼叫巨人呢 這個是後話.
神的兒子和人的女子的結合.
可能就生了上古英武的有名人物.
然後有一些動詞和形容詞.
又是在創世記第一章出現.
又說漢建.
你記得第一章說到什麼.
又說漢建創造裡面的東西.
漢為好的.
但這裡就相反了.
漢建是什麼.
是罪惡 邪惡.
Ra.
特意用這個字的 Ra.
見到人在中日所思想的思念的盡都是邪惡.
Ra.
所以上帝見到的就不是美好.
見到中日的都是Ra.
所以有一個這樣的回音.
呼應.
第四章你去看一看的時候.
我們就發現.

$^{41}$相似的地方就是.
都是用一些這樣的主題.
畜牧業 農作.
然後有城鎮.
然後最後有歌頌什麼的.
近東的詩歌 爭辯詩等等.
就是說畜牧業農業是神子所賜的.
所以要去拜畜牧的神.
你要拜農田的神.
等於今天的民間信仰都是這樣.
你做那一行的就要拜那一行的神子.
說到城鎮也是神子所賜的.
最後歌頌的就是人類多偉大.
還有神子多偉大.
眾神多偉大.
創世的第四章就說到另一個主題.
用這些所謂的ideas 這些意念.
去討論什麼呢.
就是討論人類的文明.
但是你見到創世的第四章.
如果你用近東人的眼光去看的時候.
你就會發現創世的第四章說的話.
是很奇怪的.
因為你記得近東聽到這些東西.
誰會給我們農業.
誰會給我們畜牧業等等.
這些不是上天給我們的神子嗎.
等於什麼神農氏等等.
華人文化背後的神話.
但是創世的第四章.
完全沒有這些神話的味道.
為什麼會有農業和畜牧業.
是因為人類管理大地.
創世的第一章的後果.
為什麼會有城市.
該人會建立城市.
完全沒有說誰的神子次來.
而是他自己建築.
人類最高管理的後果.
但是第四章.

$^{81}$用這些主題的時候.
他在說什麼.
人類離開了神.
半變上帝的時候.
離開了神沒錯.
因為有地上的文化使命.
在世界上神給人有那種能力.
有那種智慧.
可以去管理.
可以去組織.
然後可以發展不同的文明出來.
但是越來越遠離神的時候.
這些文明會走向哪個地方.
就是暴力的增添.
罪惡最後是不滿意.
所以創世記對近東的述說的方法.
是做一個意識形態的批判.
Genesis 1-11.
As Ideology Critique.
Critique什麼呢.
就是批判當時的.
米索不達米亞帝國的思想.
帝國的意識形態.
為什麼會有城市.
為什麼會有聚居.
米索不達米亞的文明.
為什麼會有這樣的工業.
為什麼會有這樣的文明.
因為神子.
因為我們人類的偉大.
因為君王的偉大.
你記得第一章我們討論過.
神按照他的形象所做人的時候.
那句話其實都是一個意識形態的批判.
因為當時的米索不達米亞.
蘇美爾 阿卡利亞 以色列 伯布隆尼亞.
他們都說誰是神的形象.
君王在上者.
都是一個暗示也好明示也好.
一種意識形態的批判.

$^{121}$這裡有一本書有一個學者叫Milton.
他專門去看創世記第一章.
用近東的背景去看創世記第一章.
這樣看就比較通了.
為什麼創世記第四章無端端要說這些.
哥人 阿伯是這兩個的工業.
然後有城市 文明的建設.
第四章所說出來的就是.
遠離了神的時候.
罪惡滿盈是一種批判.
人離開了上帝.
文明雖然是有.
但是最終都會怎樣.
暴力會增添然後罪惡滿盈.
是暴往死亡的.
神是需要救贖人類的.
所以在第四章說到加普的時候.
最後這兩個的slide就是說這件事.
說到阿當和哥人離開了神的文明發展.
追溯到第七代拉麥.
暴力的擴張和罪惡的滿盈.
布魯斯沃奇說的是接下來.
說到第四章.
哥人的加普和後裔的成就.
加普是直線的共七代.
由哥人開始到最後是拉麥的四個子孫.
他的加普出了第一個至今術.
第一首詩歌.
離開了伊利蘭之後.
第一首詩歌拉麥的詩歌.
第一個城市.
他的加普都是人類文化的.
文明的象徵.
擁有偉大的文明卻沒有上帝.
同時罪術者很沉默地去爭辯.
那些異教的思想.
異教的神話.
這些異教的思想.
文化使命.
Cultural mandate.

$^{161}$神要給人去管理萬物.
第一章創世記所說的.
這些進步歸因於神靈.
或者所謂半神靈.
作者就是在和這些東西爭辯.
所以是一種意識形態的批判.
同時第四章尾的時候.
作者不只是說這個.
追溯到另外一個加普.
這個就是神的後代.
神國的後代.
就一直追溯下去.
去到最後羅亞.
然後分三個意旨.
這個有什麼特別呢.
因為第四章25節開始.
就說到神去到.
保存著天國的子民.
就是另一個的宗系出來.
由瑟特開始.
第四章25節.
亞當又和妻子同房.
這句話是第四章第一節開始.
現在又呼應.
又開始一個新的家系.
一個新的子民出來.
他生了一個兒子Ben.
兒子就是Ben.
吸他起名叫瑟特.
因為他說神給我立了另一個後裔.
Sara.
這個Sid.
Sara.
你記得在哪裡出現過.
就是三章十五節.
女人的後裔.
生了一個兒子Ben.
代替了阿伯.
因為那個人殺了他.
第26節.

$^{201}$瑟特也生了一個兒子.
Ben.
繼續下去.
Sara Ben.
一個子孫.
就吸他起名叫爾羅士.
那時人開始去呼求耶和華的名.
開始去敬拜神.
去應信上帝.
去宣告上帝的同在.
所以去到這裡的時候.
似乎作者去到邀請和挑戰.
挑戰讀者到底我們是哪一班人.
到底我們是這邊.
遠離上帝.
縱使有文明.
但仍然步向滅亡.
還是這邊.
有上帝的名在中間.
我們去敬拜他.
宣告他等等.
所以我們結束的時候.
看看創世記的家譜.
你就知道家譜都想講一些神話.
第一章開始.
天地的起源.
到第二章四節.
第一個分段.
人類的起源.
這個伊甸園.
然後第三章.
罪惡的起源.
人的叛變.
或者我們說人的墮落.
The Fall.
然後第四節.
就是講到恩典的起源.
神仍然與人同在.
給他有皮夾等等.
然後第四章又是另一個cycle.

$^{241}$就是惡者的後代.
這個蓋人.
然後最後的時候.
講完拉麥.
罪惡滿盈.
然後四章二十五節.
到第五章整個家譜.
講到義者.
失特的這個後代.
直到五章最後.
然後第六章第一節.
又是另一個故事出現.
就是罪惡的滿盈.
剛才我們看了一點.
然後義者蒙恩.
就是這個羅亞.
然後六章第九節.
另一個大的分段.
就是神的審判洪水.
就彌琳.
所以我們可以有些討論的問題.
我們不多講這些問題.
如果你可以看回這些問題.
可能跟我們的組員.
或者我們自己靈修也好.
或者我們教導也好.
我們可以試一下.
有一些這樣的反思.
如果你喜歡我們的YouTube頻道.
就記得Like, Comment和Subscribe吧.
拜拜!.
\newpage



\section{馬太福音 24:1-34}
\label{sec:5qb8PRS0GKE}
\textbf{亂世啟示: 橄欖山講道 - 終末論 (馬太福音 24\_1-34) - 陳崇基牧師}
\newline
\newline
連結: \href{https://youtube.com/watch?v=5qb8PRS0GKE}{\texttt{https://youtube.com/watch?v=5qb8PRS0GKE}} ~~~~ 語音日期: 2022-03-15
\newline
\newline
\hyperref[sec:NyL1j_DJIlI]{\small{< < < PREV SERMON < < <}}
~
\hyperref[sec:index_chronic]{\small{[返順時目]}}
~
\hyperref[sec:index_scriptual]{\small{[返順卷目]}}
~
\hyperref[sec:RE3zuT2jZmI]{\small{> > > NEXT SERMON > > >}}
\newline
\newline
馬太福音 24:1-34
\newline
\begin{longtable}{cl}
\hline
\hline
章節 & 經文 (和合本修訂版)\\
\hline
24:1 & \begin{tabularx}{0.7\textwidth}{X} 耶穌出了聖殿，正離開的時候，門徒前來，把聖殿的建築指給他看。 \end{tabularx} \\ \\ \relax
24:2 & \begin{tabularx}{0.7\textwidth}{X} 耶穌回應他們說：「你們不是看見這一切嗎？我實在告訴你們，這裡將沒有一塊石頭會留在另一塊石頭上，而不被拆毀的。」 \end{tabularx} \\ \\ \relax
24:3 & \begin{tabularx}{0.7\textwidth}{X} 耶穌在橄欖山上坐著，門徒私下進前來問他：「請告訴我們，甚麼時候有這些事呢？你來臨和世代的終結有甚麼預兆呢？」 \end{tabularx} \\ \\ \relax
24:4 & \begin{tabularx}{0.7\textwidth}{X} 耶穌回答他們：「你們要謹慎，免得有人迷惑你們。 \end{tabularx} \\ \\ \relax
24:5 & \begin{tabularx}{0.7\textwidth}{X} 因為將有好些人冒我的名來，說『我是基督』，並且要迷惑許多人。 \end{tabularx} \\ \\ \relax
24:6 & \begin{tabularx}{0.7\textwidth}{X} 你們也將聽見打仗和打仗的風聲。注意，不要驚慌！因為這些事必須發生，但這還不是終結。 \end{tabularx} \\ \\ \relax
24:7 & \begin{tabularx}{0.7\textwidth}{X} 民要攻打民，國要攻打國，多處必有饑荒、地震。 \end{tabularx} \\ \\ \relax
24:8 & \begin{tabularx}{0.7\textwidth}{X} 這都是災難的起頭。 \end{tabularx} \\ \\ \relax
24:9 & \begin{tabularx}{0.7\textwidth}{X} 那時，人要使你們陷在患難裡，也要殺害你們；你們又要為我的名被萬民憎恨。 \end{tabularx} \\ \\ \relax
24:10 & \begin{tabularx}{0.7\textwidth}{X} 那時，會有許多人跌倒，也會彼此陷害，彼此憎恨； \end{tabularx} \\ \\ \relax
24:11 & \begin{tabularx}{0.7\textwidth}{X} 且有好些假先知起來，迷惑許多人。 \end{tabularx} \\ \\ \relax
24:12 & \begin{tabularx}{0.7\textwidth}{X} 因為不法的事增多，許多人的愛心漸漸冷淡了。 \end{tabularx} \\ \\ \relax
24:13 & \begin{tabularx}{0.7\textwidth}{X} 但堅忍到底的終必得救。 \end{tabularx} \\ \\ \relax
24:14 & \begin{tabularx}{0.7\textwidth}{X} 這天國的福音要傳遍天下，對萬民作見證，然後終結才來到。」 \end{tabularx} \\ \\ \relax
24:15 & \begin{tabularx}{0.7\textwidth}{X} 「當你們看見先知但以理所說的那『施行毀滅的褻瀆者』站在聖地（讀這經的人要會意）， \end{tabularx} \\ \\ \relax
24:16 & \begin{tabularx}{0.7\textwidth}{X} 那時，在猶太的，應當逃到山上； \end{tabularx} \\ \\ \relax
24:17 & \begin{tabularx}{0.7\textwidth}{X} 在屋頂上的，不要下來拿家裡的東西； \end{tabularx} \\ \\ \relax
24:18 & \begin{tabularx}{0.7\textwidth}{X} 在田裡的，不要回去取衣裳。 \end{tabularx} \\ \\ \relax
24:19 & \begin{tabularx}{0.7\textwidth}{X} 在那些日子，懷孕的和奶孩子的就苦了。 \end{tabularx} \\ \\ \relax
24:20 & \begin{tabularx}{0.7\textwidth}{X} 你們要祈求，好讓你們逃走的時候，不遇見冬天或安息日。 \end{tabularx} \\ \\ \relax
24:21 & \begin{tabularx}{0.7\textwidth}{X} 因為那時必有大災難，自從世界的起頭直到如今，從沒有這樣的災難，將來也不會有。 \end{tabularx} \\ \\ \relax
24:22 & \begin{tabularx}{0.7\textwidth}{X} 若不減少那些日子，凡血肉之軀的，就沒有一個能得救；可是為了選民，那些日子將減少。 \end{tabularx} \\ \\ \relax
24:23 & \begin{tabularx}{0.7\textwidth}{X} 那時，若有人對你們說：『看哪，基督在這裡！』或『在那裡！』你們不要信。 \end{tabularx} \\ \\ \relax
24:24 & \begin{tabularx}{0.7\textwidth}{X} 因為假基督和假先知將要起來，顯大神蹟、大奇事，如果可能，要把選民也迷惑了。 \end{tabularx} \\ \\ \relax
24:25 & \begin{tabularx}{0.7\textwidth}{X} 看哪，我已經預先告訴你們了。 \end{tabularx} \\ \\ \relax
24:26 & \begin{tabularx}{0.7\textwidth}{X} 若有人對你們說：『看哪，基督在曠野裡！』你們不要出去；或說：『看哪，基督在內室中！』你們不要信。 \end{tabularx} \\ \\ \relax
24:27 & \begin{tabularx}{0.7\textwidth}{X} 好像閃電從東邊發出，直照到西邊，人子來臨也要這樣。 \end{tabularx} \\ \\ \relax
24:28 & \begin{tabularx}{0.7\textwidth}{X} 屍首在哪裡，鷹也會聚在哪裡。」 \end{tabularx} \\ \\ \relax
24:29 & \begin{tabularx}{0.7\textwidth}{X} 「那些日子的災難一過去，太陽要變黑，月亮也不放光，眾星要從天上墜落，天上的萬象都要震動。 \end{tabularx} \\ \\ \relax
24:30 & \begin{tabularx}{0.7\textwidth}{X} 那時，人子的預兆要顯在天上，地上的萬族都要哀哭。他們要看見人子帶著能力和大榮耀，駕著天上的雲來臨。 \end{tabularx} \\ \\ \relax
24:31 & \begin{tabularx}{0.7\textwidth}{X} 他要差遣天使，用大聲的號筒，從四方，從天這邊直到天那邊，召集他的選民。」 \end{tabularx} \\ \\ \relax
24:32 & \begin{tabularx}{0.7\textwidth}{X} 「你們要從無花果樹學習功課：當樹枝發芽長葉的時候，你們就知道夏天近了。 \end{tabularx} \\ \\ \relax
24:33 & \begin{tabularx}{0.7\textwidth}{X} 同樣，當你們看見這一切，就知道那時候近了，就在門口了。 \end{tabularx} \\ \\ \relax
24:34 & \begin{tabularx}{0.7\textwidth}{X} 我實在告訴你們，這世代還沒有過去，這一切都要發生。 \end{tabularx} \\ \\
[1ex]
\hline
\hline
\end{longtable}
$^{1}$我們一向都在講啟示錄的.
我已經講了兩次了.
本身我們繼續講下去.
但是我就想一下.
我們需要去明白啟示錄裡面所講的東西的時候.
我們需要有一個整體一點的.
我們聖經裡面叫做《宗末論》.
Eschatology.
講起這個字的時候我們覺得很難.
其實一點都不難.
什麼叫做Eschatology呢.
是Eschaton這個字出來.
Greek的希臘文這個字.
意思就是說《宗末》的意思.
或者終點.
就是最後The End的意思.
聖經裡面有很多的不同的經文.
是跟這個《宗末論》是有關的.
Eschatology有關.
其中一段除了在啟示錄之外.
就是《馬太福音》第24章.
所以我們今天花一點時間去看《馬太福音》第24章.
一個很出名的港獨叫做Olivet Discourse.
意思就是說在艦艦山上面的港獨.
一個整段的聖經.
整段的經文.
就記載耶穌基督在艦艦山上面面對著耶路撒冷.
來講了一篇的說話.
這篇的說話是講到.
耶路撒冷聖城.
聖殿的將來.
還有世界終末的將來.
所以因為在艦艦山上面講.
我們就叫它做Olivet Discourse.
一篇的港獨.
但是我們每次看一篇的經文的時候.
要看一下上下文.
上下文發生什麼事呢.
如果看一下《馬太福音》第21章的時候.
就是耶穌基督.

$^{41}$榮進耶路撒冷.
騎著路進去耶路撒冷的時候.
然後他進去聖城.
就結成聖殿.
你記得.
然後跟當地的耶路撒冷的權貴.
宗教的領袖.
就發生了很多的衝突.
你記得很多說有禍有禍有禍.
那些文士,法利賽人等等.
然後就出到城裡面.
出到城外面.
就講到這個艦艦山的港獨.
如果你記得我剛才所講的東西.
這個就是耶穌在地上.
最後一個星期發生的事.
榮進耶路撒冷就是Palm Sunday.
就快要來到了Palm Sunday.
然後就進到城裡.
有結成聖殿.
有衝突等等.
然後就是到第24章.
講到耶路撒冷的將來.
就是Olivet Discourse.
所以我們試一試看這一段聖經.
就是耶穌基督未上十字架之前的幾天.
來講的一段聖經.
第24章第一節.
耶穌出了聖殿.
往前走的時候.
門徒就將聖殿的建築.
去到紙壁去看.
我們現在看的方向.
這張想拍的地方.
就是在艦艦山上拍的.
是對著現在的Temple Mount.
上面有金頂.
稱為Tomb of the Rock.
以前他們看見的.
就是當時耶穌時代的聖殿.

$^{81}$差不多這個.
你看到四四方方的一個平台.
是當時希律的.
非常宏偉的聖殿.
你今天去到這個西牆的地方.
或者下去西牆下面.
有個叫西牆隧道.
很多時候旅客都會去看的.
那裡留下了很多.
城牆的大的石頭.
那些石頭是幾十呎長.
幾米深的.
最大的那塊是五百多六百多噸.
是一塊石頭.
所以那些門徒.
指著聖殿的建築的宏偉.
告訴耶穌.
他們好像鄉下人出城一樣.
哇 這個是多麼偉大的.
聖殿多麼偉大的建築.
其實是當時古代七大奇蹟.
Seven Wonders.
可以媲美的一個建築物.
大希律所建的.
所以在後期的拉比的文學.
甚至說如果你沒有見過希律的聖殿的話.
你是從來沒有見過.
有生之年都沒有見過.
一個宏偉華美的建築.
所以這個可以明白.
為什麼門徒出聖殿的時候.
就跟耶穌說.
你看看這個聖殿是多麼的宏偉.
耶穌的答案當然是跟門徒不同.
他說什麼呢.
他說你們不是見到這一切嗎.
我實在告訴你們.
將來在這裡在這個聖殿.
必沒有一塊石頭留在另外一塊石頭上面.
每一塊都要被拆下來.

$^{121}$這個是一個很驚人的宣告.
門徒就問耶穌.
暗啞底來問他.
當他坐在橄欖山上的時候.
門徒就暗中來問耶穌.
請告訴我們什麼時候會有這些事.
就是說聖殿被毀的時候.
另外李降臨和這世代的終結.
有些什麼的預兆.
有些什麼的兆頭.
這個就是很多時候我們所謂在末世.
兵荒馬亂的時候.
我們就會看這段聖經.
就是Oliver Discourse.
之後耶穌說.
聖殿被毀到底會有什麼情況發生.
當時會是怎麼樣.
還有世代的Asceton的終結會是怎麼樣.
所以我們首先要想一想.
正如我們之前說啟示錄都是一樣.
我們說看這段聖經的時候.
先不要當它是預先寫好的報紙.
想清楚當時第一世紀的門徒教會.
會怎麼理解耶穌以下說的話.
這段聖經比較長.
我也不想把它cut出來.
所以我們看多一點.
第四節.
耶穌回答他們.
回答門徒.
他說你們要小心.
不要被人迷惑.
因為許多人要假冒我的名而來.
說我就是基督.
要迷惑許多人.
你要聽見打仗和打仗的風聲.
這個是和盤.
你們要小心不要驚慌.
因為這是免不了的.
不過結局還沒有到.

$^{161}$民要攻打民國要攻打國.
到處有饑荒和地震.
這一切不過是苦難痛苦的開始.
那時人要把你們送去受苦.
殺害你們.
你要因我的名被萬民去恨惡.
那時許多人會失去信仰.
彼此出賣.
互相恨惡.
也有許多假先知出來迷惑許多人.
因為不法的事.
增多許多人的愛心就冷淡了.
唯有堅忍到底必然得救.
這福音.
天國的福音.
要全遍天下向萬民作見證.
然後結局.
The End.
才會來到.
這裡說了所謂末世的徵兆.
我想一些比較年長的.
或者甚至是和我這一代.
或者比我更前一代的.
可能聽過很多的講道.
講到末世末代的徵兆.
都是出自這個Oliver Discourse.
這裡提到什麼呢.
有假先知 有假基督出現.
有戰亂 有地震 有饑荒.
有迫害教會等等.
這些所謂的signs of times.
signs of end times.
我們方法一定要搞清楚.
第一世紀的教會.
大概AD30.
那一代有40年左右.
AD30至AD70.
他們怎樣去理解這段話.
假基督 民要攻打民 國要攻打國.
有地震 有饑荒.

$^{201}$教會被迫害 被打壓.
不法的事增多.
許多人失去信心.
離棄主.
愛心冷淡.
天國的福音 傳遍天下.
他們會怎樣去理解呢.
很多時候我很不理解.
現在我就很不理解.
以前讀神學的時候.
讀新約聖經的survey的時候.
我不明白為什麼我們沒有讀.
新約以外的一些文獻.
我自己畢業之後 傳道的時候.
我看回其中一個很重要的.
新約以外的歷史學家寫的東西.
就是Josephus 約瑟夫.
約瑟夫寫了很多東西.
都是新約背景不能夠沒有的.
我們要明白新約聖經的時候.
就需要明白他所寫的東西.
在約瑟夫的記載裡面.
在耶路撒冷被毀.
就是AD70前後發生了什麼事.
他很詳盡的去記載.
因為他首先是猶太人.
他也是祭司.
他也是去到參與.
第一次反羅馬革命的.
其中一個將軍.
他就帶領他的子弟.
去對抗這些羅馬軍隊.
然後他就明白.
當這些羅馬軍隊打過來的時候.
不夠他打.
所以他就變節.
親羅馬.
然後被羅馬將軍去縱容他.
所以他是有份參與.
在公元後主後70年.

$^{241}$耶路撒冷被毀.
羅馬軍隊去攻陷耶路撒冷這場仗.
他是做他們通風報信的.
所以有很多猶太人到今天.
都覺得這是一個出賣他們的人.
但是他完了這場仗之後.
他就寫了很多很多.
有關AD70裡面發生的事.
並且非常詳盡的講到.
當時有很多的假先知出現.
當時的戰亂.
羅馬帝國的外憂外患.
或者當時第一次反羅馬革命裡面.
出現的不同的黨派.
猶太人的黨派.
他們之間彼此去到狗咬狗骨等等.
地震.
也都記載不少.
除了他之外當然有其他羅馬的歷史學家.
饑荒.
嚴重的發生.
這個不單止做CIVIS有這樣看.
《士徒行傳》十一章也都有這樣看過.
迫害教會我們更加不需要講了.
所以可以見到.
在耶穌基督升天之後.
那一代40年.
由30至到70.
在耶路撒冷被毀的時候.
這40年到底發生了什麼事.
在做CIVIS的歷史著作裡面.
有詳細交代.
在羅馬的歷史學家.
他們留下的著作.
都有詳細交代.
當時是一個兵荒馬亂的時代.
所以可以見到所謂的徵兆.
其實在當代的教會已經看到.
你說不是啊.
第十四節說什麼.

$^{281}$天國的福音要傳遍天下.
這個肯定當時不會.
或者你看看其他的新約經文的話.
你就見到他們當時.
in some sense.
連保羅都認為.
這個福音已經是傳遍天下.
歌羅西書一章六節二十三節.
福音已經傳給天下的萬民.
普天之下.
所以in some sense.
對當時的第一世紀的教會來說.
他們有理據.
絕對有理據.
相信人子是隨時會再來.
特別是AD70年.
耶路撒冷聖殿被毀之後.
他們有理由.
絕對有理由相信.
在他們有生之年.
人子是隨時會回來的.
所以啟示錄所記載的.
就是這一段時間.
特別是AD70之後.
他們已經看到所謂的signs of times.
看到這些末代.
eschaton.
應該會發生的這些所謂的徵兆.
然後二十四章十五節我們繼續看.
耶穌說什麼.
當你們見到但以你先知所說的.
造成荒涼的可憎者.
站在聖地的時候.
讀者就必須要領悟了.
Let the reader understand.
當時就會有走難的時候.
那些日子如果是懷孕的.
是與養孩子的有和.
你應該祈求.
祈求什麼.

$^{321}$就是你們走難的時候.
不要是冬天.
以及不要是安息.
因為那時必定會有大災難.
即使從世界的開始到現在.
未曾有過的.
以後也必不再有的.
如果那些日子不減少的話.
沒有一個人可以存活.
但是為了選民的緣故.
那些日子必會減少.
那造成荒涼可憎者的.
這個是大義利書七章九章等等.
天啟文學裡面出現過的.
就是那一位去到在聖殿.
去到建立或者設立一個的idol.
去污衊聖地的.
那個王會出現.
如果你看一下舊約的歷史.
問一個猶太的讀者.
猶太人.
他們就會告訴你.
有好幾次在歷史裡面已經出現過.
586BC.
巴比倫軍打到來的時候.
毀滅聖殿.
污衊聖殿.
主前167年.
當時馬加比時代的之前.
希臘的一個.
大家聽到.
臭名遠播的一個王.
來到去污衊耶路撒冷的聖殿.
到AD70.
主後70年.
聖殿被毀.
其實就是歷史重現.
這一位造成荒涼的可憎者出現.
當時耶穌說到.
如果你見到這件事發生的時候.

$^{361}$就是AD70.
你就會見到什麼.
就是末期已到.
Last days已經到.
你們就要走難了.
希望那天走難的時候.
不要是冬天.
不要是安息日.
因為如果你這樣就走不了.
初期教會歷史.
聖經以外所記載的.
就是AD70之後.
耶路撒冷的一些教會的弟兄姊妹.
一些信徒.
就是走難去了.
耶路撒冷北面約旦河東以外的一個地方.
叫做Palva.
一個古城.
在那裡避難.
所以你又見到.
對當代的人來說.
是一個不是寫給21世紀將來2000年之後發生的事.
是當代眼前就見到.
24章29節.
那些日子的災難.
光過去.
AD70.
光過去.
太陽就變黑了.
月亮也都不會發光.
眾星會從天墜落.
天上的萬象震動.
那時人子的徵兆要顯在天上.
地上的萬族都要哀號.
你要見到人子帶著能力.
滿有榮耀.
駕駛著天雲降臨.
當豪統發出響聲.
他就猜顯他的使者.
將四方的選民.

$^{401}$從天的這邊到天的那邊.
都召集來到.
這些不會吧.
未發生吧.
如果你沒有看過新約聖經以外的記載.
你就不會知道.
對當時的人來說.
曾經有這樣的事發生.
剛才我提過約瑟夫.
約瑟夫的猶太戰士.
第五至六頭像.
如果你有機會就回去看.
他就說到當時發生什麼事.
特別就是六十年代末的時候.
耶路撒冷快要淪陷的時候.
約瑟夫曾經提過.
他說當日日落的時候.
就見到戰車和軍隊.
在雲上面出現.
是包圍耶路撒冷.
這是一個淪陷的空笑.
你儘管回去看.
為什麼會提到約瑟夫.
剛才我提過的一些細節之外.
就是他一個所謂.
不是新約教會裡面的人.
他是一個外國人.
因為他不是信基督.
但他見到的東西.
就是當時耶路撒冷被毀之前.
發生的天象和奇事.
所以我們看到現在的時候.
我們就發現.
這個Oliver Discourse裡面.
所提過的所謂Signs of the Times.
其實是第一代的信徒.
已經見到目前眼前的事實.
The last days has begun.
末世已經展開 已經開始.
如果繼續看下去的話.

$^{441}$你就會發現.
聖經裡面說的終末論.
到底是一回什麼事.
第六節 第八節 第十四 第二十一節.
這裡又好像在說一些.
不是AD70年的時候會發生的事.
為什麼呢.
因為耶穌說到有打仗的風聲.
這是少不免的.
不過結局The end 還沒有到.
不過是痛苦的開始.
其實原文是Birth pains.
生了一個兒子的媽媽就知道.
當你陣痛的時候.
孕婦生產的陣痛.
就會越來越劇烈 會加劇.
就是說所謂的Signs of the Times.
其實就是陣痛的開始.
會繼續加劇 會繼續Intensify.
第十四節.
福音要傳遍天下 向萬民作見證.
然後The end 結局才會來到.
二十一節.
那時必有大的災難.
是世界開始和現在未曾有過的.
以後也必不會再有的.
這個應該很清楚.
到我們今時今日.
主後七十年 耶路撒冷聖城淪陷.
當然是一件很不簡單 很艱難的事.
但絕對不是我們歷史裡面最艱難的時候.
在我們中間有些走過難的 你們就知道.
那些的戰爭 那些的迫害.
肯定是比主後七十年更加厲害.
所以耶穌在說什麼呢.
耶穌在整個的馬太福音二十四章.
Oliver Discourse.
就說到一方面 預言聖殿是被毀的.
但另一方面又夾雜了很多.
是世界終末的時候會發生的事.

$^{481}$所以一方面在說AD70.
另一方面在說The End會發生什麼事.
就兩樣混在一起說的.
這個就是在聖經裡面末世論.
或者終末論.
一個比較棘手的地方.
誰打誰 還是混為一談.
為什麼要混為一談.
其實混為一談當然有它的道理.
因為在猶太教裡面.
特別是在第二教的猶太教.
他們認為聖殿被毀 聖城被毀.
是一個世代的終結.
另一個世代的展開.
就是那個Old Age要離開了 結終.
然後那個New Age會到.
那個marker就是聖殿被毀.
另一個很簡單的原因.
就是你記得開始的時候.
門徒問什麼問題.
門徒就問兩個問題.
第一個問題就是這件事發生會怎樣.
就是說這些石頭被推倒.
聖殿被毀 華麗的聖殿被毀的時候.
什麼時候會有這些事.
這是第一個問題.
然後他們又問多一個問題.
你的降臨和這世代的終結有什麼預兆呢.
這些就是Sunday School的那些不好的學生.
無端端問 為什麼問兩個問題.
問一個 問完這個他回答完再問第二個.
我們今天就會很清楚.
但是他們就問多一個.
所以耶穌一回答就回答兩個問題.
所以在這個Oliver Discourse裡面.
就有夾雜了聖殿被毀當時發生的事.
也都有世界終末的時候發生的事.
但其實這樣看也有它的好處.
因為看到聖殿被毀的時候.
這個是鐵一般在歷史裡面的事實.

$^{521}$AD 70其實是一個預演.
好像你看一套戲.
未看之前你有看預告片.
到你看完整套戲的時候.
預告片裡面所發生的畫面就會重現.
或者你可以做一個比喻.
你去聽一個很偉大的音樂.
就是Beethoven的Symphony 5.
第五交響樂.
你有沒有印象這首交響樂是怎樣.
這個Pattern.
三個音再加一個.
這種音的Pattern.
就在整個交響樂曲.
由開始的時候的序曲.
去到奏出來.
然後你繼續聽每個Movement的時候.
這個Pattern又繼續不斷重現.
不斷去Intensify.
不斷去到更加強烈的演奏出來.
聖殿被毀AD 70.
直到今天就會有很多這樣的預兆.
這些空兆.
活現在我們眼前.
之後每一個世代.
在過去的二千年的教會歷史.
都會看見這些所謂的Signs of Times.
或者用一個圖表來表示的時候.
就可能更加清楚.
一個說話如果今天任何東西都未吸收到.
這句說話你記住了.
一定會幫到你.
就是叫做Already but not yet.
Already是什麼意思.
就是說這個末世.
Last stage.
Already已經展開了.
Already begun.
Not yet.
就是說已經繼續去逼近這個完全的實現.

$^{561}$但是Not there yet.
但是一定是每一天過去的時候.
就繼續逼近這個最終的基督的在.
所以所謂講到末世的時候.
我們說現在就是末世了.
末世就是你了.
其實真是在聖經裡面所講的.
The last days.
一早已經在這裡.
我們現在這個時間就是Last days.
聖經裡面所講的.
終末或者The last period.
其實就是基督的第一次來.
和基督的第二次來.
中間這一段時間.
Already就是告訴我們.
今天這個世代.
就是人類歷史最後的一段時期.
之後不會再有另外一段.
這一段時間就是最後一段的時間.
Already意思就是說.
我們已經看到有改變.
不像以前那樣.
世界的歷史也不像以往一樣.
因為這個是The last period of time.
Not yet.
也告訴我們天國是不斷地我們在逼近.
到某一天很快到.
趨時到.
我們就需要去交著.
所以我們經常說.
世代的終結.
末日什麼時候來.
那個ascendant什麼時候來.
到底有什麼徵兆呢.
剛才我已經解釋了.
這些徵兆.
第一代的信徒.
All of the discourse.
裡面已經看得很清楚.

$^{601}$戰亂,饑荒.
所以我們不要再當.
那個終末什麼時候到.
到底有什麼徵兆讓我們看到.
Signs呢.
到底有沒有Signs呢.
好像我們有個checklist.
我們說Signs number 1 and 2已經到了.
我們還在等3,4,5,6,7等等.
其實早就已經到了.
耶穌說什麼.
第30節.
那時.
The Ascaton的時候.
真是他再來的時候.
人子的徵兆要顯在天上.
地上的萬俗都要愛好.
看見人子滿有能力的駕魂而到.
不需要再等徵兆.
也不需要再look for.
這些所謂的signs of times.
因為耶穌說什麼.
耶穌說如果你再要等下一個徵兆.
到時已經太遲.
因為下一個徵兆.
就是人子的徵兆.
Signs of the Son of Man.
他就會到.
到你下一次見到徵兆的時候.
已經太遲.
這樣說.
第一代的教會.
是身處在末世的開頭.
The beginning of the end times.
如果你這樣想的時候.
我們現在的世代.
我們就是在末世的尾聲.
The end of the end times.
再沒有時間讓我們浪費.
每過一天的時候.

$^{641}$福音就更加傳到萬邦.
如果你說第一世紀.
In some sense.
對保羅來說.
已經傳到當時所知的世界.
The known world的萬邦.
我們今天更加知道福音傳到萬邦.
戰爭不安激放.
比以前更加intensify.
教會被打壓.
也會越演越烈.
中國大陸.
今年三月.
三月頭開始.
就實施一個互聯網禁令.
這些一早已經見到.
是The signs of end times.
很多時候我們說主座台.
我們覺得好像靠客.
好像令我們很驚慌.
我們已經有很多事情要怕.
我們又要怕疫情.
又要怕這個又要怕那個.
又要怕死.
為什麼還要說這些末世的事情呢.
其實聖經裡面提到終末的時候.
其實不是要我們怕.
只不過是要我們警醒.
是完全是為了你和我的好處.
我們小時候在家.
媽媽很多時候會說一句話.
等你爸爸回來.
等你爸爸回來這句話很有意思.
因為有兩個可能的反應.
第一個就是你會很害怕.
為什麼呢.
因為你做得不好有規矩.
等你爸爸回來其實想你做什麼.
就是想你悔改.
這是第一個反應.

$^{681}$第二個反應.
等你爸爸回來.
就是你十分期待.
為什麼呢.
因為拿了好成績.
等爸爸回來的時候.
就一定有獎賞.
做得好.
不怕他回來.
並且想他快點回來.
到底我們的取態是什麼.
是對我們生命裡面的一些警醒.
耶穌沒有停在一個橄欖山講道.
在講所謂的signs of times.
因為你記得夏文.
我們再有時間會繼續處理.
我們講起中末的時候.
就想起那個被提怎樣.
之後我們會繼續處理.
被提到底是否合聖經.
為什麼不合聖經等等.
但是他沒有停在橄欖山的講道.
之後你記得24章之後.
是講什麼呢.
好幾個的比喻.
才幹的比喻.
僕人的比喻.
就是說主在來會隨時發生.
你和我都要交仗.
十個同類的比喻中心.
僕人的五千,二千,一千的比喻.
就是叫我們.
我們要交仗.
我們要見主面.
我們可不可以見主面的時候.
上帝說你真的很善用.
我所給你的資源,金錢.
你有樂力奉獻.
時間你有樂力去善用.
你的才幹你用盡它.

$^{721}$就是說叫我們作為基督的門徒.
我們不是一些宗教消費者.
你知道什麼叫religious consumer.
就是看完之後覺得很好.
我很enjoy.
然後依然顧我.
這叫religious consumerism.
我們的呼召是做天國的門徒.
是用我們的金錢,我們的才幹,我們的時間.
來推進上帝的國度.
因為this is the last day.
不再有時間給我們浪費.
we are at the end of the end times.
在最後的終末.
不會再有時間來為天國浪費.
要把握現在這個最終末的時候.
最終末的一刻.
在烏克蘭受到俄羅斯入侵的時候.
大家見到很多新聞.
其中一段新聞大家應該見過.
就是一個爸爸.
這個爸爸做什麼.
這個爸爸說.
他要放下他的太太和女兒.
自己留下去保護國家.
是一個很sad的新聞.
但對於他來說.
this is the last day.
不再有時間浪費.
沒辦法了.
這個時候不打.
就不會再有時間打.
所以他願意留下.
我們都是一樣.
如果我們已經來到終末的終末.
請問你.
你還有多少時間浪費.
問一問自己.
剛才我們講到.
起初的愛心.

$^{761}$起初的時候.
比較一下.
起初的時候.
可能你對主.
對教會.
對福音.
你真的義不容辭.
甘心樂意.
多吃力多辛苦.
你都肯定.
不用逼你.
現在.
會不會在一些很混亂.
很不安的.
我們起初的愛.
我們就不知道掉到哪裡去了.
我們悔改.
以免主再來的時候.
我們是和地上的萬族.
一同哀哭.
或是我們主再來的時候.
我們可以和他交仇.
準備好我們自己.
和我們生命.
不再浪費.
今日主來.
我們說我們準備好了.
我們準備好了.
我們同心低頭.
上帝求你來警醒我們.
讓我們在主的恩典裡面.
我們長進.
讓我們把握我們現在最後的時間.
你都可以事奉.
去推進你自己的國度.
去建立你自己的工作.
以致你再回來的時候.
我們就是忠心良善的僕人.
我們這樣祈禱.
奉主耶穌基督的名字.

$^{801}$我們拜拜.
\newpage



\section{創世記}
\label{sec:RE3zuT2jZmI}
\textbf{創世記第5章 (5.1) - 背景釋經 [比對研究,亞當家譜] (陳崇基牧師)}
\newline
\newline
連結: \href{https://youtube.com/watch?v=RE3zuT2jZmI}{\texttt{https://youtube.com/watch?v=RE3zuT2jZmI}} ~~~~ 語音日期: 2022-03-15
\newline
\newline
\hyperref[sec:5qb8PRS0GKE]{\small{< < < PREV SERMON < < <}}
~
\hyperref[sec:index_chronic]{\small{[返順時目]}}
~
\hyperref[sec:index_scriptual]{\small{[返順卷目]}}
~
\hyperref[sec:mEHavpOJC9o]{\small{> > > NEXT SERMON > > >}}
\newline
\newline
$^{1}$今天我們繼續看《創世記》第五章.
我們看《聖經》的時候,不是單單看經文的內容.
我們希望能夠做到的是將一些.
《釋經》的工具或者方法.
能夠試用一下,試用一些不同的方法來思考經文.
以致我們在教導上能夠多一些不同方面的見解.
所以我們回來這個比對的研究就是.
他說《聖經》不是一個真空的狀態出現的.
《聖經》一向都不是由天降下來.
然後用一些非人話來溝通的.
總是會有它那個歷史的處境.
文化的處境,特地的狀況.
所以《聖經》的啟示我們不要想.
一定是非常特別.
當然特別是因為上帝的話語上帝的啟示.
有上帝的權威.
但是特別的地方我們不要想到.
跟當時的世界,歷史的處境是完全沒有關係的.
所以我們看回這個圖的時候.
我們就發現原來《聖經》是在近東的世界.
近東的文化裡面出現.
上帝是藉著近東的世界.
來對他自己的子民說話.
以及啟示他自己的信息.
有相似的地方也有差異的地方.
一個古代的以色列人.
大衛摩西等等.
《聖經》舊約裡面的以色列人.
他跟他的鄰舍.
埃及人,亞述人,迦南人,巴比倫人等等.
他們之間有很多.
在他們的文化裡面相似的地方.
相似的地方多於.
我一個現代人.
跟一個古以色列人之間相近的地方.
相似的地方是什麼?語言上.
用的術語和觀念.
如何看待這個世界.
寫作裡面修辭的手法.
這些全部都是古近東裡面.

$^{41}$所謂他們的common ground.
就是大家相似和共通的地方.
在這個處境裡面.
上帝就藉著當時的近東文化.
不同的方面來啟示他自己的信息.
所以有相似的地方也有很多差異的地方.
所以對於上帝人世界之間的關係.
《聖經》就啟示一個完全新的信息.
所以為什麼我們過去看一些所謂近東神話.
我知道有很多時候我們還是在一些很傳統的.
舊有的觀念裡面.
我們覺得為什麼要鬥這些東西.
因為一些巴比倫的神話.
甚至將來我們會看一些近東洪水的技術.
是很像羅亞洪水技術.
像到一個地步可能會覺得我們很驚訝.
為什麼我們驚訝呢?.
因為我們以為.
我們覺得《聖經》只有它去說話.
只有它去記載某些事件.
其實不是的.
其實在近東世界裡面有很多共通的東西.
他們都會大家去討論.
或者大家會有不同的見解.
《聖經》是其中一個.
當然是按照上帝的啟示.
按照上帝的權威來表達出來的.
一些不同的信息.
所以在這個比對的研究裡面.
我們看到《聖經》有相似的地方.
但是也有很多差異的地方.
差異的地方我們可以像這裡六項.
當然六項不是彼此之間對立的.
很多時候是overlap.
很多時候是大家有相近的地方.
大家可以看到《聖經》裡面有不同的經文.
以及有不同的信息.
是跟近東的背景.
近東的文化有分別.
分別可以是所謂的degree of difference.

$^{81}$就是一個分差的程度.
最少的程度當然就是一些.
《聖經》說一些另類的.
不同的一些見解.
對人生觀有些什麼見解等等.
這些就是alternative.
也有時候我們發現有一些經文.
《聖經》很著意要跟當時的文化.
宗教.
社會政治不同的意識形態.
來做對話的dialogue.
也有時候是做一些所謂的counter text.
就是那個經文是做一些反對的言論.
就是反對當時社會.
當時的文化.
他們有的意識形態.
也有parody.
就是去諷刺.
譬如是什麼呢?.
譬如在《先知書》裡面有時候很多.
有一些諷刺的說話.
是說到外邦人的偶像.
說這些偶像是用木頭做的.
但是他沒有生命的.
不能夠說話的等等.
所以有一些這樣的諷刺.
同樣也有一些更加強烈的批判.
所以一種critique.
上幾次我們也有說過.
《聖經》裡面的一些用語.
一些的概念.
如果我們有這個背景的時候.
你就會發現.
原來《聖經》是有去到challenge.
或者去criticize.
就是批判當時的.
巴比倫帝國他們做的.
或者說的一些意識形態.
譬如人人不是平等的.
《聖經》第一章.

$^{121}$《創世記》已經說了.
人人都是神的形象.
整體的人類都是神的形象.
甚至最強烈的就是一種的爭辯.
爭辯就是說.
去到.
近乎用諺語來攻擊他.
譬如《以賽亞書》.
有很多不同的章節.
說到一些的異教.
一些偶像的宗教等等.
也就是說獨一的.
這些的神都是虛假的.
所以做一些很強烈的所謂的爭辯.
一些的爭辯.
所以雖然這六樣東西.
或者甚至可以更多的方法來描述這些的差異.
有人用不同的字眼來去形容.
但是都是一個程度上的分別.
比較弱一些.
以及比較強烈一些.
所以我們如何知道聖經在做什麼呢?.
就是我們需要去了解.
這些其他所謂異邦.
他們留下來的一些文獻.
那麼.
Again.
我們是需要這麼做的.
因為我們明白了這個背景越來越多的時候.
我們就越清楚聖經在做什麼.
而且我們不需要怕.
去做一些這樣的近東背景的比對.
曾經有人問過.
那個人在讀神學的.
他就開始讀《創世記》.
他在神學院的課程裡面就很驚訝.
他發現原來講創世的故事.
不是希伯來聖經才講的.
原來巴比倫,亞述,雅及等等.
這些所謂異邦,倫邦.

$^{161}$都會講這些創世的故事.
並且有一部分是相似的.
當然也有很多差異.
但是他很掙扎的一個問題就是.
他就問為什麼聖經不是第一個講.
而是後這麼多的.
比如之前我們講過幾次.
我們講了一些巴比倫的近東的創世史詩.
創世的神話.
那些應該是比希伯來聖經.
創世記第一二章所記載的創世故事早.
他就有很大的掙扎.
為什麼這麼遲聖經才講.
為什麼聖經這麼遲才講.
是不是神的啟示呢.
神的啟示不一定是第一個講的.
神的啟示不一定是第一個用那個形式來講出來.
我就舉一個很簡單的例子.
你看新約的聖經.
保羅不是經常寫書信.
書信的格式是當時希羅社會.
希羅世界用的書信格式.
首先有一個寫信人.
然後收款者.
然後才下面再講書信的內容.
如果你看希羅文化的作者.
到寫信的時候.
其實格式是幾乎一樣的.
甚至用語也是非常相似的.
那你說保羅是不是全世界第一個用這個格式寫東西的人.
當然不是.
只不過他用了當時流行的格式.
來講神學.
和教會寫信.
甚至將它questionize.
將它基督教化.
即是開始的文安那些句語.
他加了原因為平安等等.
基督教的字眼.
基督信仰的字眼.

$^{201}$加了上去.
所以那個形式是很相似.
但是內容本質.
想講的東西就完全不同.
但是保羅絕對不是第一個寫信的人.
再看看新約的一些傳記形式.
好像福音書一樣.
他們是不是第一個寫耶穌的傳記.
耶穌的福音書的述史呢.
當然不是.
路加寫《少林寺傳》的時候.
他是不是講過.
我怎樣去收集這些不同的故事.
不同的基述.
然後我就和阿非羅大人.
按次序的再寫出來.
換句話說.
起碼寫《少林寺傳》歷史的時候.
他不是第一個寫.
有很多其他的account.
有很多其他的記載.
他就將它收集.
然後再寫出來.
所以不是先後的問題.
我們不要覺得.
因為我們經常講聖經是神的話語.
他就一定是第一個.
不一定.
往往就是在一個已經有的處境.
有這樣的世界文化裡面.
聖經是啟示.
上帝和人溝通的時候.
就是用當時人明白的語言和觀念等等.
我們看《創世記》第一章開始.
我們講到創世.
我們第一次的時候已經講了.
不是要解答一些現代人的問題.
而是和當時古代的以色列人.
在近東的世界裡面.
他們所接觸的很多創世的神話.

$^{241}$一些論述等等.
去到做一些.
你可以叫他另類.
也可以是一些反對.
甚至諷刺.
甚至去到polemic.
爭辯的情況.
所以不是講科學的東西.
而是講當時世界人所有的觀念.
他們所問的問題.
去回答他們誰才是創造主.
所以講到這裡的時候.
我們今天會看創世記第五章.
但是我們先看第一章.
我剛剛告訴你.
在聖經和當時近東的文化.
宗教等等.
去做一些爭辯也好.
或者是做一些對話.
或者是做一些對立也好.
你就可以明白.
你就可以想一想.
聖經第一章的一些記載.
到底是想做些什麼.
和以色列人想做些什麼.
這裡是一章第十六節.
我們慢慢讀.
.
這句話是什麼意思呢.
.
我們記得希伯來文.
開始的時候.
有動詞的.
他造了.
He made.
He create.
動詞之後就是主格.
主格就是動詞的主者.
就是Elohim.
然後就是這個at.

$^{281}$如果你不太記得.
很多我們看過的希伯來文.
這個一定不記得.
這個aleph加tath.
下面三點.
e的音短的e音.
就是at.
就是一個客體的標記.
所以從頭開始.
這裡首先有一個動詞.
He made.
via us.
Elohim就是Elohim made.
神造了.
造了什麼呢.
這個沒有翻譯出來的.
就是at.
就是告訴你之後的東西.
造了什麼呢.
.
是兩個的意思.
.
就是一.
.
就是two.
二的意思.
兩個.
.
這個字是光的字.
演變出來的.
其實就是光體的意思.
就是發光的東西.
可以叫做Lamp.
可以叫做燈.
發光的東西.
光體.
.
兩個的光.
然後後面這個.
Hagadonim.

$^{321}$就是形容詞.
形容這個名詞.
.
這個形容詞就是.
Hagadonim.
就是大的意思.
大的眾數.
因為這裡說的是兩個.
是眾數.
神造了兩個.
.
兩個的大光.
Great light.
然後他就繼續說下去.
這兩個的Great light.
這兩個的大光是做什麼呢.
.
第一個就是.
那個的大光.
這個又是一個客體的標記.
不會翻譯的.
.
Light of greatness.
就是光是大的.
.
這個動詞演變出來的.
.
就是管理的意思.
統治的意思.
.
就是說這個大光.
.
這個大光就是管時間.
就是白天的時間.
.
大光就是管這個晝.
白天.
然後就另外一個光.
因為有兩個光.
第二個光就是小的那個.

$^{361}$.
這裡又是同樣.
好像上面那樣.
多了一個N.
就是一個VAR.
一個N.
然後有一個客體的標記.
作者就說.
我現在說那個小的光.
.
這個字和上面那個一樣.
平衡.
就是那個光.
那個光怎麼形容呢.
就是小的.
上面就是客多.
就是大的.
下面就是客多.
就是小的光.
管理這個字又出現了.
但是管理的意思就不同了.
因為上面是客多.
白天的時間.
下面是客多.
所以這個小的光.
就是管理夜晚的.
就是晚上的時間.
然後下面就加兩個字.
.
就是眾星.
那你說.
在近東的文化裡面.
如果我們認識多了的時候.
怎麼看到它在做什麼呢.
很簡單.
就做兩個光.
一個大的管晝.
一個小的管夜.
然後就眾星.
但如果你明白近東的文化.

$^{401}$所有天上的萬象.
都是有神聖的.
太陽是太陽神.
月亮是月亮的神.
所以在《希伯來聖經》裡面.
太陽是有名詞的.
我們今天英文都會說.
Sun.
我們不會說Great light.
我們說月亮就是Moon.
月亮.
我們不會說Lesser light.
不會說Small light.
但這裡就是《希伯來聖經》作者在做的.
就不用Shamash.
Shamash就是太陽的意思.
也就是在近東裡面.
Shamash是一個神.
一個神子.
一個神明.
就是太陽神.
也都不會用月亮這個字.
他就用Lesser light這個字.
就是.
.
就是Small light的意思.
小的光.
就不用.
.
月亮這個字.
因為這個字.
月亮.
等於華人的農曆.
Lunar calendar.
又是一個神的名字.
或者指向一個神聖的entity.
所以沒有用Shamash.
也沒有用月亮這個字.
.
所以為什麼不用呢?.

$^{441}$明明有這些字在字典裡面.
他偏偏用了什麼呢?.
.
就是大光.
.
小的光.
意思就是說.
他故意不用這兩個字.
就是讓讀者知道.
這些全部不是神聖的東西.
避免不用這些.
當時近東的人.
就用了一些神教的名字.
因為說到太陽的時候.
近東人就想起.
是不是那個太陽神.
是不是那個迦南人的太陽神.
或者那個宗教裡面的月亮神.
他就不用.
他就用大光和小光.
然後最後呢.
.
就是眾星.
這裡幾乎是不說也可以.
但是他好像加了一個豬腳.
做為眾星.
為什麼呢?.
因為所有的眾星.
木星,金星,什麼星星都好.
都是一些敬拜的偶像.
都是在近東的世界裡面.
是有神聖的.
所以你看到這裡.
有很多學者都說是Demythologize.
就是將它去神化化.
那些當時的世界所預設的.
那些神聖的東西.
全部都去除掉了.
沒有了那些宗教的味道.
大光,小光,眾星.

$^{481}$只不過是上帝所做的.
事就這樣成了.
所以你看到他在做什麼.
就是我剛才的圖表.
就是這些東西.
他去到跟當時近東的世界.
太陽,月亮,星星.
全部都可以拜.
但是他寫創世的第一章的時候.
就說這些東西不是敬拜的對象.
連名詞都不說出來.
暗示地說這些東西.
值得敬拜的只是創造主.
不是創造物.
所以有一個很大的差異.
怎樣看到呢?.
你只能夠看到.
你只能夠多認識近東的背景.
不用我們現代人的眼光.
去進入進度.
先用古近東的世界的人.
在想什麼.
你就知道這個聖經.
創世的第一章.
剛才的例子.
就是在做什麼.
可能是一些另類.
或者反對.
一些暗地裡的諷刺也好.
甚至說是一種prolific.
怎樣也好.
程度上的問題.
有多強烈的問題.
你就看到.
他用的用語.
有很多時候都是很刻意的.
想過的.
去講清楚到底神學的信息是什麼.
我們很快去到第四章.
最後的時候.

$^{521}$我們說記載完這個該人的.
離開了上帝.
因為他們的disobedience.
他們的背叛.
所以他們離開了耶和華的棉.
然後一代一代地生下來.
生到第七代.
去到拉瑪.
第四章快要到尾的時候.
最後的時候.
我們提過的二十五節.
這個作者又回頭了.
回頭說什麼呢.
說亞當的另一個兒子.
就是瑟特.
後裔的一個兒子.
繼續生下去.
到以羅氏等等.
然後二十六節最後.
那時人開始來到呼求耶和華的命.
你記住.
不是拉瑪之後才說這些人的歷史.
是回頭再倒敘.
說回瑟特.
然後一直追隨這個宗系.
這個宗系開始.
來到去將耶和華的命.
宣告去敬拜.
還有去認信.
開始來到擴散.
開始有另外一個人類出現.
這個就是跟隨耶和華的.
去敬拜耶和華.
所以有一個倒敘.
回到瑟特.
四章完結之後.
五章第一節.
看看你聽不聽到.
作者又呼應之前四章.
說過的一些字眼.

$^{561}$五章第一節.
以下就是亞當後代的記錄.
如果你看那個原文的話.
其實是寫字.
寫字帳戶.
有人以前猜.
可能會不會是一塊泥板呢.
等等.
那些家譜的字.
好像近東的獵王表.
寄出來.
特別這裡有一個廣告.
說到亞當後代的記錄.
這個字.
這裡說上帝創造人的時候.
按照自己的樣式來做.
他創造一男一女.
創造他們的時候.
神賜福給他們.
稱他們為.
人類.
亞當一百三十歲的時候.
生了一個兒子.
形象樣式都與自己相似.
就給他起名叫做瑟特.
亞當生瑟特的時候.
之後還活了八百年.
然後生了其他的兒女.
亞當共活了九百三十年.
就死了.
這裡就是新的人類開始.
家譜開始了.
我們說過.
聖經特別在創世記.
經常用家譜來講神學.
所以由亞當到瑟特開始.
這個新人類.
這裡有一些字眼.
讓我想起創世記第一章.
創世記第一章.

$^{601}$神願創人的時候.
那個心意是做什麼呢.
是在管理大地等等.
敬拜他.
這個由亞當到瑟特開始.
就是敬拜上帝的新人類開始.
與這個該人的作一個對比.
很多時候家譜和家譜.
都會做一些對比和比較.
這裡最清楚了.
所以這裡講回形象樣式.
這是第一章有提及過的字眼.
人類等等.
亞當就是一個正確的名字.
開始一個所謂的歷史家譜.
記載他的家譜.
然後我們去到家譜的第1.
2,3,4,5,6,7.
第7.
就去到以洛.
以洛65歲的時候生了瑪圖撒拉.
生了瑪圖撒拉之後.
以洛和神同行300年.
並且生了其他的兒女.
共活了365年.
他與神同行.
所以神就把他娶去.
他就不在了.
這些字眼.
又叫我們回想.
與神同行.
伊甸園.
記得嗎?.
神去到walk.
在他們中間去行.
所以這裡又講回.
原本神和人之間.
應該有的關係.
這一句話.
walking with God.

$^{641}$創世記.
都會出現.
羅亞.
阿布拉罕等等.
甚至到先知書都有講這些說話.
譬如恆公二號.
憐憫於謙卑的心與神同行.
一個二者的生命.
有神的生命.
李家書六章第八節等等.
這裡就講到.
這個new humanity.
new community.
講到他第七代的時候.
加了一些字眼.
雖然.
阿普五章.
一直講.
他生了兒女然後他死了.
然後繼續講到.
下一代生了兒女又死了.
有個死字.
在一個死亡的.
真實的狀態裡面.
但是在第七代.
居然有一個不死.
有生命出現.
到底是.
飛了上去還是怎樣.
沒有很多仔細去講.
但最重要的.
就是講到生命的.
強調.
這個家譜裡面第七代.
拿了以諾出來.
講出.
這個是生命的.
經歷.
經歷到上帝的生命.
被神取去.

$^{681}$就不在了.
很明顯.
跟該隱的第七代.
就成一個對比.
是不是.
因為該隱的第七代.
就是什麼.
最後就是死亡.
暴力.
你記得.
撒埃默的遭報77倍.
他講話.
他那首歌.
跟他兩個的妻子說.
阿大和阿細.
那首詩歌.
77個希伯來的字母.
全部都是講暴力的.
強暴.
這邊就完全不同.
是講生命.
第七代和第七代作一個對比.
兩班人.
做一個對比.
一個是跟隨耶和華.
一個是離棄了耶和華.
然後我們繼續下去.
五章第28節.
出現了一個字.
在該隱的家譜裡面.
也出現了一個名字.
就是拉麥.
第九代.
拉麥1882歲的時候生了一個兒子.
叫他起名做羅亞.
然後他就說.
這兒子必使我們從地上的操作.
和手中的老虎得著安慰.
因為耶和華曾經就坐這地.
拉麥就生羅亞.

$^{721}$還活595年.
並且生其他的兒女.
拉麥共活了777歲.
就死了.
羅亞500歲的時候就生了閃含羽雅弗.
這裡就拉麥和拉麥.
來做一個比較.
剛剛就是第七代和第七代.
做一個比較.
然後第五章.
失特的家譜.
敬拜上帝的這班人.
又用一個同名.
拉麥.
來做一個比較.
為什麼會用拉麥呢.
我想提出兩份文件.
一份文件是91年寫的.
就是一個.
福音派的近東學者叫做Richard Hayes.
Richard Hayes.
去分析這個家譜的時候.
創世的家譜.
就說到拉麥這個人.
他就引用一篇.
再舊一點的一篇論文.
那篇論文就是一篇考古的論文.
那篇考古論文.
就是一個近東學者.
猶太的近東學者叫做Michael Kugin.
Michael Kugin去到.
研究.
那些初期有字母的時候.
有時候會找到一些.
學生寫的.
寫ABCD.
去到.
就好像我們剛剛學希伯來文.
我們就要學寫那些字母.
Consonant.

$^{761}$我們要學寫.
那22個.
其實在以色列.
找到很多這樣的圖片.
就是那些打碎了的牙缸.
那些大的圖片.
有時候就是一些.
我去俗語叫.
寫得不是那麼好.
好像可能小朋友學生.
寫的那些.
他去到.
Michael Kugin去到看一些.
圖片的時候.
他就發現.
因為Eventually.
這個希伯來的字母.
就有22個.
剛剛一半的時候.
就是Lambert開始.
Lambert開始.
譬如這裡右上角.
在四海古卷.
發現的地方困難.
這裡屬於.
主前一世紀的時候.
在山洞那裡找到.
很多這樣的圖片.
上面你看到這裡爛了.
但下面有一個長線畫下來.
這個就是Lambert.
Lambert在這些.
迦南文.
希伯來文.
永遠都是第二半.
開始的.
因為你算一下.
在希伯來的字母裡面.
有22個.
Alphabet.

$^{801}$Cov.
就是第一個.
然後12開始.
就是下一半的字母.
就是Lambert.
Cov Lambert.
Mamlun Summit.
所以你將22個字母.
斬一半.
第二半.
是Lambert.
所謂L開始.
所以在這些字母的寫法裡面.
都是這樣.
所以Lambert.
在中間.
其實結束上一半.
然後開始下一半.
所以在家譜裡面.
都是有刻意這樣做.
因為在.
兩個家譜裡面.
一半的時候.
就是誰呢.
Lambert.
然後下面.
第五章敬拜神的家譜.
也是由Lambert.
去做一個所謂分水嶺.
你可以看得很清楚.
第四章.
拉密.
是有說話的.
整個家譜裡面只有他有說話.
還說了77個字母.
說Violence死亡.
到這個家譜.
第28節.
作者刻意去.
讓他說話.

$^{841}$他說生在.
羅馬的時候.
希望這個.
就是我們的盼望.
Noah這個字.
是好像.
得安慰這個字.
的音.
所以在Noah.
出生的時候.
就是羅馬出生的時候.
就希望在這個人身上得到安慰.
就是救恩.
可以從他而來等等.
就是一些祝福.
不是咒詛.
是一個新的盼望的開始.
所以在這個家譜裡面.
你會看到.
這個拉密.
就是上一半和下一半.
的.
開始.
刻意將這個名字.
放在這兩個家譜的.
然後就說出.
整個家譜.
有很多不同的號碼等等.
我們都可以.
做一些比對的研究.
就是這些.
碼的碎數等等.
有沒有一些類似家譜.
或者一些獵王表.
我們能夠去做一些比較.
然後我們知道.
第五章特別.
想做些什麼.
和當時近東的世界有什麼分別.
我暫時說到這裡.

$^{881}$如果你喜歡我們的影片.
就記得要讚好,留言和分享啦.
如果你喜歡我們的影片.
就記得要讚好,留言和分享啦.
\newpage



\section{創世記}
\label{sec:mEHavpOJC9o}
\textbf{創世記第5章 (5.2) - 背景釋經 (陳崇基牧師)}
\newline
\newline
連結: \href{https://youtube.com/watch?v=mEHavpOJC9o}{\texttt{https://youtube.com/watch?v=mEHavpOJC9o}} ~~~~ 語音日期: 2022-03-15
\newline
\newline
\hyperref[sec:RE3zuT2jZmI]{\small{< < < PREV SERMON < < <}}
~
\hyperref[sec:index_chronic]{\small{[返順時目]}}
~
\hyperref[sec:index_scriptual]{\small{[返順卷目]}}
~
\hyperref[sec:mFPkQ1HAFUA]{\small{> > > NEXT SERMON > > >}}
\newline
\newline
$^{1}$好 今天我們繼續看第五章 創世記第五章.
上一次我們講到創世記第五章.
這個家譜.
這個家譜就是一個十代的家譜.
由亞當 失德開始.
一直講到第九代到第十代.
拉麥和羅亞.
這個家譜很有趣味.
我從小到大 可能年輕的時候在教會.
一直教的.
就是說這些家譜的人.
有很長壽的年代.
每一個 你看到他們壽終.
是很長的 雖然還不夠一千.
到底這個家譜想表達什麼呢.
家譜神學是做什麼呢.
在最近幾十年.
因為近東的文獻.
開始有更多的了解.
也做了一些比對研究.
你記得我們從開始到現在.
我們一直在做比對研究.
我們說比對研究這個工具.
能夠幫助我們找到一些相似的近東作品.
相似近東的故事 或者寫作.
能夠跟聖經做一些對比.
所以我們不需要怕.
去找這些文獻.
翻譯出來 大部分當然是英文.
我們可以很粗略地把它翻譯成中文.
我們試一下去做一些比對.
比對要做什麼呢.
比對研究就是將一些類似的文獻.
篩選一邊.
至少有兩方面.
第一方面就是相似的地方.
另一方面就是差異的地方.
所以我們就試一下.
繼續用這個比對的方法.
Comparative Studies.

$^{41}$我們看看創世記第五章.
有什麼神學的意義.
以致比對之後能夠更加顯明出來.
加普十代.
講到他幾多歲的時候生兒子.
生了之後活了多少年.
最後他就死亡了.
好像亞當一樣 九百三十年等等.
有沒有一些類似的記載.
在近東我們能夠做一些比對 比較呢.
如果熟悉近東的文獻的人.
可能都聽過.
這個我們叫做蘇美爾的列王表.
Sumerian King List.
這個Sumerian King List是做什麼的呢.
就是在一個prism.
一個四方形 有四面的.
有一些Cuneiform的褻形文字.
雕刻在上面.
其實是一些列表.
列出在不同的城市.
遠古之年.
不同的城邦.
比如這個埃利都.
有兩個王.
兩個王的名字是什麼呢.
然後他們各自統治了多少年.
然後這個所謂的城邦.
再次有霸權的時候.
就不再是埃利都了.
就去到另外一個城邦了.
就是巴德提比拉.
這個城邦.
裡面有三個王.
就是說在這一段時間.
在巴德提比拉這個地方.
當時的所謂霸權.
或者一個城邦.
可以有很大的影響力.
對整個南美索不達米亞.

$^{81}$就是蘇美爾地區.
有很大的影響力.
這個城邦有三個王.
然後就去到另外一個城市.
再去另外一個城市.
然後這段時間去到斯魯柏.
這個地方.
就有另外一個王.
烏巴拉圖圖.
然後就說洪水來到.
洪水之後就繼續記載.
有什麼城市和有什麼王.
首先就說基什這個地方.
然後就數下去.
數下去有很長很長的.
到後面.
其中一個王我們曾經提過.
叫做基加美斯.
他就在烏魯克.
就是烏魯克一個古城.
當時大概是3500BC的時候.
就是主前3500至4000的時候.
在那裡作王.
你見到他們列出.
不單止是城邦和城邦裡面的王.
也列出什麼.
列出他作王多少年的.
你見到這些年數非常之大.
28800年.
另外一個王就36000年.
然後就轉了去這個.
這個地方.
巴德提比拉這個地方.
然後這裡有四萬多年.
一直都是數以萬計.
數以千計的年.
然後洪水就來到了.
你見到有一個分水嶺.
前面遠古的一些列王歷史.
然後就到洪水.

$^{121}$洪水到了.
洪水之後.
那個王權又去到另外一個地方.
英文就是Kingship Descent from Heaven.
開頭這裡也有一句這樣的東西.
Kingship Descent from Heaven.
我們第一章的時候說到神的形象.
這些君王都是自稱為神的代表.
神的形象等等.
到洪水之後.
也有回這句東西.
Kingship Descent.
王權降下來.
從天降下來.
去到這個基什的地方.
這裡你見到有些Paladin.
有些格式.
洪水之前.
這些年數是很長的.
數以千萬計.
然後洪水之後.
比較實數一些.
比較現實一些的數字.
譬如基加美斯.
126.
之後的王更加少.
3年,5年,20年,40年等等.
比較現實一些的.
統治的一些時間.
我們見到這樣的記載的時候.
我們首先不要說這些是天方夜譚.
我們首先用一個文學的手法.
來做一些比對.
文學手法的意思是.
他們用的格式和手法.
我們試一下將它比對.
聖經特別是創世記第五章.
它的格式和手法是怎樣.
到底想表達什麼意義.
首先我們看這些近東的文獻.

$^{161}$我們很多時候第一個.
特別是傳統教會的環境.
我們聽見那些神話.
創世神話.
或者以往我們曾經分享過.
另外一些的文獻.
我們就說我們怎可以看這些.
這些是異教虛無的神所做的.
一定是假的.
我們不是真假的問題.
而是那個記載的手法.
希望你能夠掌握到那個.
comparative studies.
那個比對研究.
到底想做什麼.
如果你將這個蘇美爾獵王表.
Shemarion King list.
和創世記第五章的家譜.
去做一些比較.
有什麼差異呢.
第一樣的差異就是.
米索不達米亞.
很明顯是用60進制的.
我們曾經我記得說過.
這個米索不達米亞.
米索不達米亞.
他們發明了60進制.
或者一些量度的單位.
都是60.
一分鐘有60秒等等.
現在我們都在用的.
這些是第一個.
人類的文明歷史裡面.
蘇美爾人Shemarion所發明的.
所以他們的數字的單位.
往往都是60進制.
剛才你記得那些獵王表的數字.
很多都是60進制的.
看到嗎.
這些都是36000等等.

$^{201}$6萬呀 6000呀等等.
600呀等等.
所以他們是60進制.
《聖經》希伯來文.
可能就是10進制.
當然我們現在不是說.
絕對清楚到.
但是似乎都可能是10進制.
所以在數字的計算上.
有分別.
另外有什麼的差異呢.
有什麼differences.
就是蘇美爾獵王表.
是一個獵王城邦的獵表.
King list.
等於舊約裡面其他的經文.
一些的獵王表.
獵出一個王.
統治了多久等等.
然後呢.
這個創世記第五章.
就是一個宗族.
一個家譜的記載.
沒錯 有這些差異.
但是很明顯.
有些相似的地方.
如果你第一次看.
這個蘇美爾獵王表的時候.
你應該立刻可以留意到.
是有洪水的歷史性的.
總之有個大洪水.
有多大.
這個講到第六章的時候.
我們才再討論.
但是那個洪水.
這個災難性的事件.
好像是一個歷史的轉捩點.
一個turning point.
蘇美爾獵王表.
很清楚講到之前有一堆王.

$^{241}$然後呢.
就有洪水來到.
洪水來到之後就有一堆王.
還有什麼相似的地方呢.
就是洪水之前.
數字很大.
Sumerian king list.
你看到有這麼大的時候.
你都不會驚訝.
創世記五章的數字.
一點都不大.
最大的是969.
馬兔薩拉.
但是你去到比較蘇美爾獵王表.
king list.
Sumerian king list.
裡面的創世記五章的數字.
根本就是非常皮毛.
因為數以千萬計.
洪水之後就降下了.
那些年數就降下了.
這個蘇美爾獵王表就是這樣.
創世記五章都是一樣.
之前是很大的數字.
然後之後就比較小.
一直開始小下去.
所以毫無疑問.
有相似的地方.
手法上有不同的地方.
獵王和加普的體材.
似乎可以作一些比較.
我想在這裡提醒大家.
我們做一些比較的時候.
不是做一些所謂.
很多時候我們想起這些東西.
我們覺得.
或者一些傳統的教導的分途底下.
我們會覺得這些很自由派的東西.
不信聖經的東西.
不是這個意思.

$^{281}$大家信聖經.
大家都可以對聖經的權威很執著.
都有一個所謂high view of scripture.
我們去比較的時候.
我們發現可以用一個genre.
用一個所謂文學的體材.
來做一些比較.
我這裡想提一提.
我這裡給你的資料.
是2021年的時候.
在美國一個很出名的.
福音派的神學會議.
Evangelical Theological Society.
裡面發表的文章.
其中一個學者.
一個近東學者.
一個福音派近東學者.
現在是三個神學院教近東歷史的.
近東背景的.
舊有的聖經.
就是Richard Affeberg.
他在2021年的ETS.
Evangelical Theological Society.
去發表一份文章.
那份文章還沒有刊登出來.
因為太近了.
去年發生的 2021年.
那份文章就是在說比對研究.
所以你看到有很多近東學者.
甚至福音派學者.
我們都在做一些了解.
嘗試去明白那個所謂的logic.
那些古代的文獻.
聖經都是古代文獻.
聖經不是現代文獻.
是近東的文獻之一.
去做一些手法上的比較.
到底有什麼我們可以發現到.
或者可以留意到.
第一 他說的是什麼.

$^{321}$他這樣去做比對的時候.
似乎用意不是給後代讀者.
給他們一個列王表.
或者一個家譜.
而是去推算一些遠古時間表.
用意不是這樣.
用意是做什麼呢.
似乎那些洪水之前.
那些很大的年數.
很高的年數.
就是告訴讀者.
這些人是很遠古的人.
用這樣的手法去表達.
然後洪水之後.
就是想跟讀者告訴他.
這些實際的年數.
小一點的年數.
就是告訴讀者.
這些這樣的王.
如果你說的列王表.
蘇美爾列王表.
就是告訴他.
這些王就是近代的一些王.
越小的number就越近我們.
所以他們就用數字.
用數字去表達出來.
所以你看到有相似的地方.
家譜 創世記第五章.
似乎都是用同一個手法.
所以我們這樣比較的時候.
我就發現.
我們會不會有時候.
譬如我在傳統教會長大.
經常聽到一些護教學的回答.
就是說洪水之前.
因為天氣氣候是不同了.
跟現在不同.
洪水之後就不同了.
所以以前的人壽命長一點.
現在的人壽命短一點.

$^{361}$因為紫外線問題等等.
這些我想告訴你.
聖經裡面沒有說過.
只不過我們想像出來的.
更好的方法我認為.
就是Alfred Beck這裡所說.
就是先了解文件創世記第五章.
和Sumerian King's List.
不是現代科學的文件.
現代人不會這樣寫家譜.
也不會這樣寫歷史.
但近東人會這樣.
古近東人會這樣做.
這個Sumerian King's List.
你一比較創世記第五章.
你就更加清楚.
真的用一些這樣的手法.
再看一下這個第五章的續譜.
有十代.
格式很清楚.
每一代講到後代出生的年齡.
然後後代出生之後有多少壽數.
總壽數有多少.
如果你將這些數字.
將它的餘數除以60.
將這些數字除以60.
再剩下餘數.
譬如130除以60.
就剩下10.
800除以60就剩下20等等.
如果你將這些除以60的餘數.
再加起來的時候.
這裡有三個數字.
頭兩個數字.
你將60的餘數加起來.
然後將60的餘數再加起來.
另外一個數字.
兩個數字加起來.
156加209就是365.
似乎告訴我們.

$^{401}$不是偶然的事.
似乎聖經古代的作者.
用續譜.
用這些數字.
想去表達一些神學的意義.
這個似乎是經文的intention.
他的用意是想做這件事.
除了這些比較.
好像不是意外的數字之外.
另外一位學者叫做Candance Sparck.
她在這本書.
叫做Ancient Texts for the Study of Hebrew Bible.
一本五百多頁的書.
她搜集了所有近東的文獻.
將這些文獻列出不同的組別.
譬如有些是家譜.
有些是星象的文獻.
有些是歷史的文獻等等.
來幫我們做希伯來聖經的研究.
她在這本書的introduction.
開始的時候.
道論裡面提到什麼呢.
她說提到.
她就用這個創世的第五章.
來做一個比對的研究比較.
她說除了一個數字之外.
這裡出現過這麼多的數字.
29個裡面的28個.
所有家譜的數字.
最後的一個位的數目.
就是0257.
如果是這樣的話.
隨機能夠產生一個這樣的數字的列表的話.
機會率是非常小的.
4.5乘以10的to the power of minus 11.
是0.0011的4.5.
是吧.
就是意味著這些數字.
似乎不是一些一般的.
實際的.

$^{441}$我們今天現代人想的那種數字.
所以這個值得思考.
這個值得思考.
我們就面對一個問題.
就是如何去解讀.
如何去解讀.
如何去詮釋這些數字.
現代人一打開聖經.
因為已經翻譯了給我們.
我們往往就用我們現代人的pre-supposition.
用我們現代人的預期.
預設.
我們套上去.
我們就說馬多薩拉就是969年.
他有這麼長的壽命.
這個就是我們現代人讀經.
很多時候會有的過程.
用我們現代人的concept.
我們對數字.
他說多少就是多少.
報實數.
聖經無五.
但是如果對一個近東人來說.
你看看剛才的Sumerian King list.
他們會不會是這樣理解呢.
可能未必會.
是吧.
再看下去.
除了機會率之外.
還有剛才那些.
隨機會產生的.
九代除60的餘數.
加起來是365.
如果你喜歡我們的YouTube頻道.
就記得Like, Comment和Subscribe.
\newpage



\section{創世記}
\label{sec:mFPkQ1HAFUA}
\textbf{創世記第5章 (5.3) - 背景釋經 (陳崇基牧師)}
\newline
\newline
連結: \href{https://youtube.com/watch?v=mFPkQ1HAFUA}{\texttt{https://youtube.com/watch?v=mFPkQ1HAFUA}} ~~~~ 語音日期: 2022-03-15
\newline
\newline
\hyperref[sec:mEHavpOJC9o]{\small{< < < PREV SERMON < < <}}
~
\hyperref[sec:index_chronic]{\small{[返順時目]}}
~
\hyperref[sec:index_scriptual]{\small{[返順卷目]}}
~
\hyperref[sec:ng1CRwPs7JE]{\small{> > > NEXT SERMON > > >}}
\newline
\newline
$^{1}$很明顯 二落似乎告訴我們一些象徵性的數字 365 solar year 365.
五章三十一節 拉麥777 這個更加不需要說 為什麼會是777呢.
我們由第一章講到現在 之前另一個拉麥 該人的後代 第七代拉麥.
他也有很多七 租報七倍 租報77倍 他那首詩歌是77個希伯來字母.
所以用七字對希伯來的作者來說 希伯的作者來說 絕對不會陌生.
因為他們就是用七來表達一些信息 並且我們知道.
如果你將一些星象的所謂索望週期.
這個synaptic period 什麼叫索望週期呢 即是說那些觀星 星象學那些.
就是看到一顆星 天上的一顆星 然後用多少年回到天空上同一個地方.
這個叫索望週期 在巴比倫的文獻裡面 有很多這樣的日記.
很多的astronomical text 現在已經翻譯出來.
就是這些泥板 記載巴比倫人去觀察星象.
而且觀察星象的時候 有很多時候他們會將這些年數寫下來和計算出來.
那麼你會看到拉麥777 不單止是希伯來聖經的文化裡面 七是完美.
也是木星和土星的索望週期加起來的年數.
所以不單止是希伯來聖經的修辭手法777是完美.
可能在巴比倫的星象文獻裡面 777也是象徵著某些東西.
瓦列965就是金星和木星加起來 兩個數字加起來就是962.
所以我們在這裡思想到一件事 就是到底他們想用數字做些什麼.
我們先不去爭辯到底是不是真的這麼長命 這個所謂歷史性的問題.
往往很多時候我們方牌的信徒 方牌的牧者 我們看聖經的時候.
我們常常卡住一個地方 就是歷史性的問題.
這裡說的東西有沒有發生過 但是很多時候經文想表達的.
不是單單發生了一件事那麼簡單 而是告訴我們神學意義.
這個才是聖經的用意 你明白我的意思嗎.
譬如你看創世的第一章 六日七日創造.
我們第一個現代人會問的問題 是不是宇宙這樣發生.
宇宙大爆炸之後就7個24小時 然後就造成整個世界.
我們就搞錯了 因為很明顯 創世的第一章想表達的.
不是說創世真的有這個事實那麼簡單 是不是.
他們告訴我們創世 上帝創立這個世界有什麼意義在裡面.
這個才是聖經的用意.
所以我們可以這樣比對一下 這裡這塊泥板就很出名.
為什麼呢 這是一個巴比倫星象的記錄.
就是屬於主前164至163年.
記載了這個我們現代人很熟悉的哈雷彗星.
他們計算到 觀察到164BC 9月22至28日.
就在巴比倫的上空見到這個哈雷彗星.
所以他們記錄下來 你就可以看到.
這些巴比倫他們對星象很熟悉.

$^{41}$其實你看聖經也有這樣的感覺 是不是.
去到馬太福音 所謂聖誕的故事 第二章.
那些幾個博士從東方而來的 巴比倫地.
米索不達米亞 他們來的 他們看到什麼.
就看到星 大衛的星 然後就來朝拜新生王.
所以用一個這樣的星象的文獻來比對一下.
到底他們想用的數字做些什麼.
那些列王他們的壽命 他們的作王的時間.
是表達一些宗教的意義 神學的意義.
這個王是完美的 是不是.
36000 或者和60有關.
雖然我們未必找到一個方程式.
能夠拆解所有的這些數字.
但是很明顯有一些這樣的意義在這裡.
這裡 哈雷彗星出現的時候.
有些學者認為 當時正是馬加比起義.
馬加比猶太人起義 去打退希臘的軍隊的時候.
馬加比的起義就成功了.
可能他們當時就觀察到這個彗星的出現.
無論如何 不同的文化對數字和壽數的理解都不同.
現代人正如我所說.
我們很多時候都會預設那個數字就是實數.
但近東的數字是帶有其他的意義的.
起碼這個比對裡面 我們未必可以拆解所有的數字.
但起碼可以讓我們看到.
不同的文化對數字是有不同的理解.
我們這裡曾經引用過不同的學者.
其中一位就是John Wotton.
我記得十多年前上他課的時候.
他教古近東的課程的時候.
他提到他認識一個弟兄.
那個弟兄在美國東南亞一些部落的地方.
所謂蠻荒 未開化的民族.
去那裡宣教.
這個宣教士就去到東南亞的部落.
山區的部落.
他說有一年去到的時候.
他去到被人介紹.
他就上來講道 分享.
然後他分享完之後.

$^{81}$那個主席就介紹另外一個女士.
上台去到講道.
就說這個是什麼名字.
這個姐妹.
她是四十歲.
然後他就上來講話.
首先第一件事對一個西方文化的人來說.
介紹一個女士要講她幾歲.
這個已經很奇怪了.
非常奇怪的事.
這個宣教士就回到美國.
兩年之後他回到同一個宣教工場.
回到這個部落.
又見到同一個姐妹.
在聚會裡面.
主席又介紹這個姐妹.
介紹這個女人出來上台.
就說這個姐妹叫什麼名字.
她是五十歲.
他就上台講完.
然後這個宣教士就說.
喂喂喂 等等.
有沒有搞錯.
我兩年前來到的時候.
你說這個女士是四十歲.
現在兩年後.
我有日記的 記了.
她現在就變了五十歲.
到底發生什麼事.
完了聚會之後他就去詢問這些事.
那裡的負責人就跟他說.
你不明白嗎 我們這裡部落.
用數字介紹一個人幾歲的時候.
不是講她的實數.
而是一個專榮.
一個專榮的數字.
就是介紹一個人.
告訴人家她幾歲的時候.
不是講她真正的壽數.
而是一個榮譽.

$^{121}$等於一個很出名的講員.
我們介紹她上來.
所以這個宣教士就明白.
我們不可以用一個現代西方化的.
一個看法來看壽數.
很簡單 是吧.
所以他被人介紹上去.
這個宣教士 美國的宣教士.
被人介紹上去的時候.
他都會覺得真是好.
因為她三十多歲.
有人介紹她是四十歲.
他都有安.
所以數字我們不要立刻跳下去.
然後我們就說.
我們用現代人想問的問題.
就想用聖經去給答案.
我們用近東的角度.
來看這些數字.
其實這樣有好處.
我們不需要想一大堆.
我們搞出來的護教的問題.
我們不需要解答那些問題.
我們記得講第四章的時候.
我們都提過了.
該人的妻子是誰.
以往傳統的一代.
從小到大都是這樣教.
就說該人的妻子一定是他妹妹.
不過聖經沒有講.
所以以前兄妹可以結婚.
接著就不行.
律法來到就不行.
我覺得不是不需要的.
首先第一.
聖經根本就沒有提示.
沒有經文的數據.
第二.
我們已經提過了.
第一章.

$^{161}$上帝創造人類.
humanities.
有男有女.
所以你這樣看的時候.
我們不需要去引申出一些.
很難搞的護教學問題.
那些根本是不應該有的問題.
Bruce Walkie講到創世記第五章.
他從亞當·杜拉默.
如果把附生子的歲數.
和總壽數總和.
再分別除以60.
剛才我做了給你看.
根據巴比倫的60進制系統.
預數總和就是365.
這或許是代表完美的壽命.
即是說這一班人.
亞當·濕特開始直到羅亞.
這十代的人.
他們是敬拜耶和華的.
呼求耶和華.
call upon the name of Yahweh.
全部都是有完美的壽命.
一個人的年歲的週期.
可以和天星天體的運行週期.
來相配.
表明他們的壽命.
是跟著一個有意義的方向.
一個有意義的模式.
並且是以一個完整的週期來作結束.
The symbolism is significant.
不要miss this.
這個反而是創世記家譜.
宗族的家譜想表達的東西.
就是這一班人.
跟隨耶和華的.
就是完美的壽命.
有長有短.
看回這個家譜.
你也可以看到.

$^{201}$以洛65歲的時候.
生了瑪圖薩拉.
生了瑪圖薩拉之後.
以洛和神同行.
300年.
又是一些完美的數字.
除60.
生其他的兒女.
又活了365歲.
遠神同行.
他被神取去.
不再了.
他是第七代.
所以比較特別.
去到描述他.
在死亡當中.
一代一代會過去.
但是第七代.
居然得到所謂永生.
被神取去.
不需要經歷死亡.
這個第七代.
很明顯就是跟該人.
尼希耶和華的世代.
第七代有一個比較.
其他的名次.
很明顯這兩個家譜.
作者是有心去作一些對比.
跟隨耶和華.
呼求耶和華.
和不呼求耶和華.
世界墮落了.
最進入世界的是敗壞.
並且最惡甚至是.
繼續散播.
黑暗繼續散播.
繼續有敗壞.
但是這一班人.
繼續見證神自己的名.
這樣看.

$^{241}$就是聖經跟我們說的話.
到底這個家譜的讀者.
所謂完美的讀者.
作者寫《創世記》第五章.
他都會寫給某些讀者.
這個經文希望給一個讀者.
我們叫他做ideal reader.
即是誰應該看這段聖經.
看了之後他應該做一些什麼決定.
去跟隨神的敬拜神.
那個ideal reader.
就見到歷世歷代.
這班人在敗壞.
在黑暗.
在罪惡滔天的世界.
他們繼續去事奉耶和華.
繼續去呼求耶和華.
宣講耶和華的名.
這個神學的意義.
就是對我們說的.
第五章的家譜.
我們做一點結束.
就是敬拜耶和華群體.
你見到第五章.
經常都提一句話.
就是誰誰生了誰誰.
就死了.
死亡是真實的.
每個人都會死的.
你記得第二第三章《創世記》.
講到人叛變之後.
就是Doom to death.
注定會死亡.
並且離開了伊甸園.
exile.
被趕出伊甸園.
沒有上帝的同在.
是注定會死的.
但是這些家譜.
最特別的地方.

$^{281}$第五章就是不斷說死了死了死了.
但裡面他們是有一個完美的生命.
出現了多少次.
就死了.
有十代.
出現了九次.
因為第七代是以洛.
他沒有經歷死亡.
換句話說.
整體來說.
這些耶和華群體.
他們都是與神同行.
以洛就是一個表表者.
最後一代就是拉默.
又是一個對比.
五章二十八字.
拉默182歲的時候.
生了一個兒子.
就是給他起名.
這個就是第十代.
羅亞.
說者而此.
必使我們從地上的作為.
工作.
手中的勞碌.
得了安慰.
亦那含憫怒.
就是給了我們安慰.
羅亞這個字.
有一個這樣的音.
有一個這樣的意義.
就是給了我們什麼安慰呢.
就是從這個就坐當中.
死亡的裡面得到安慰.
這個很明顯就是呼應.
創世記第三章.
所以羅亞最終是得救的.
蒙恩六章第八章.
吾敬拜神的群體.
第七代表表者.

$^{321}$就是拉默去說話.
並且在他說話當中表達到.
強暴欺壓越來越多.
最終上帝以洪水來審判.
到底我們是哪一班人.
哪一個群體.
是否屬於耶和華這個群體.
以至我們見到死亡也好.
罪惡繼續擴散也好.
我們仍然是去見證上帝.
這個就給了我們一些現代的意義.
在我們今天這個崩壞的年代裡面.
崩壞的世代裡面.
我們在黑暗當中.
可能在不同的掙扎裡面.
我們可不可以去同神同行.
以至經歷上帝給我們的生命.
我們的壽數多長多短也好.
最重要我們可不可以去同神同行.
最後我想告訴你.
有一個人提醒我.
有一個比喻挺好的.
他說每個人的墓碑都有兩個號碼.
你的出生年和你死去的年.
有兩個號碼.
我們每個人都會有.
最重要的不是第一個號碼.
也不是第二個號碼.
最重要是中間那一劃.
就是你有生之年.
你怎樣去跟隨神.
去敬拜神.
怎樣用我們的生命.
這個才是最重要.
不是哪個號碼開始哪個號碼結束.
而是中間.
怎樣過我們現在的生活.
這個似乎就是第五章.
加譜神學.
給我們一個挑戰.

$^{361}$一個膚淺.
如果你喜歡我們的影片.
就記得要讚好,留言和分享.
\newpage



\section{}
\label{sec:ng1CRwPs7JE}
\textbf{從「信仰百川」到「港風之谷」(上):網絡的言論自由與不自由}
\newline
\newline
連結: \href{https://youtube.com/watch?v=ng1CRwPs7JE}{\texttt{https://youtube.com/watch?v=ng1CRwPs7JE}} ~~~~ 語音日期: 2022-03-16
\newline
\newline
\hyperref[sec:mFPkQ1HAFUA]{\small{< < < PREV SERMON < < <}}
~
\hyperref[sec:index_chronic]{\small{[返順時目]}}
~
\hyperref[sec:index_scriptual]{\small{[返順卷目]}}
~
\hyperref[sec:B4MkD7Z7rb0]{\small{> > > NEXT SERMON > > >}}
\newline
\newline
$^{1}$各位大家好,我們今天很開心有四位加上我分享的人.
第一位就是Mark Sir,陳松基牧師.
第二位是Gordon,Dr. Gordon.
第三位就是Tony牧師.
我想請幾位來介紹一下自己.
今天的題目就是從信仰百川到港風之谷.
很明顯是一個介紹一些新的主題.
不如逐個介紹一下吧,我先請Mark Sir.
其實每個人都叫我Mark Sir,不用命牧師.
我在美國的芝加哥牧警會.
一直以來二十多年都是牧警會.
一直都是一個本地教會的牧者.
有時候會發表一下自己的讀經,研經的心得等等.
所以我們有一個成員網台,這個就是很簡單的我的角色.
Tony呢?.
大家好,我叫Tony,黃家漢.
其實我一直都是在加拿大牧會的十幾年.
特別是我在加拿大的中部和西部牧會.
我現在在加拿大卡加里的牧羊教會.
本身除了牧會之外,我還會讀書和教書.
所以都忙到飛起.
有時候都會寫文章,寫詩.
有空間的時候我都會發表.
這是我的小小故事,是我自己.
多謝Tony,輪到Dr. Gordon了.
大家好,今天有三位牧師在.
我是一個很普通的小人物.
我叫自己做一個全身都沒有張力的人.
正職是做醫生的,但過去兩年就半退休.
在做一些幫人顧問的項目.
我也讀過神學.
我自己比較專注在教會的文化.
也有空閒時寫文章之類.
多謝你的介紹,我也自我介紹一下.
我還忘記說,這次應該是成會網台和我自己的頻道.
礦牧頻道的一個聯播.
我叫礦偉志,我已經牧了牧很多年.
做過不同的教會機構,最多就是做報道訓練.
我現在是雪梨牧會的.
我也是一個很普通的人.

$^{41}$大家都很謙虛.
所以為什麼會走在一起呢?.
我們這四位,就像我們主題所說.
由信仰百川說起.
不如我們逐個說一下.
大家在信仰百川的參與.
信仰百川應該有十多年了.
很久了,十多年了.
前陣子就結束了.
突然shut down了.
很多人都覺得很可惜.
當然有些人可能很開心,我不知道.
但這幾位一直都有參與.
不如說一下.
信仰百川你曾經有什麼參與.
其實你是怎麼看呢?.
為什麼會參與?.
或者在網上的世界.
其實你覺得現在網路的風氣是怎麼樣呢?.
還是請Mark Sir先說.
我在百川,像礦牧你所說.
都十多年了.
我都是寫我的文章.
最近當然就多了.
除了分享文章.
有時分享一些視頻.
一些短的影片等等.
我一直都很喜歡百川的氣氛.
因為我覺得是一個很好的空間.
給大兄姐妹,甚至我們這些牧者.
能夠發表一下,或者討論一下.
一些topic.
有時候在一些傳統教會的框架裡.
未必有空間.
或者未必有時間給我們去說這些東西.
所以我一直都很喜歡信仰百川.
不單止我自己可以發表.
我們都看到別人發表.
Tony可以寫一些詩.
Tony又可以發表一下他的文章.

$^{81}$我們又可以學到其他東西.
我覺得是一個互相學習的平台.
特別在這個世代.
是非常需要.
不是一言堂,是大家一起來去談.
我同意,其實.
你講就講吧.
其實就像Mark Sir所說.
當初參與信仰百川.
差不多是一開始我已經參與.
當中就有分享一些文章.
也有分享一些詩歌.
一詩,我自己喜歡寫詩.
我想在這個過程裡.
其實除了是一個空間.
給我自己去整理一些思緒.
譬如我自己寫通常是寫詩.
寫牧養心得.
寫教育,寫神學教育各方面的文章.
有時不是太多空間.
可以給我去發表.
就算我有東西寫的話.
我想百川有空間給我去.
整理我自己的思緒去分享.
也都看看別人對這些分享一些意見.
一些回饋,我覺得都很好的.
同時我也同意在百川裡面.
百家爭鳴.
事實上就算我們在百川裡面.
我不是所有文章我都同意.
但我就覺得起碼知道.
世界上有人想的東西和我很不同.
我覺得這個也是構成了一個學習的機會.
去看看別人對某些信仰的事情的看法.
的不同,以至去擴闊自己的空間.
那就不會有一點坐井觀天.
以為自己那套就是.
一知獨秀就看得到答案.
我想這個都是一個很好的場地.
也給我去反省.

$^{121}$如果對方敢問我.
對方有這樣的看法,我自己怎麼看.
我參加了信仰百川的時間.
不過短一點.
應該是六,七年左右.
之前主要只是看,沒有機會寫.
到六,七年之後就開始寫東西.
和大家說的都是一樣.
信仰百川有一個很特別的地方.
就是容許我們說一些.
不是在普通教會官方的渠道裡面.
看到的東西.
我覺得很多時候在華人教會裡面.
太過透視.
只有一方向去看東西.
信仰百川給大家一個機會.
可能接觸一些另外的思想.
其實我和一些大英姐妹聊天.
大家都發覺你講這些東西.
是他們在想的.
不過從來沒有機會說.
從來沒有機會拿出來討論.
信仰百川就是一個非常好的渠道.
給大家能夠去想這些問題.
好像Tony所說.
不是每樣東西都同意.
有些我說的東西.
也不期待別人同意.
但起碼可以拿出來說.
而不是掃在椅底上.
我自己其實沒有甚麼貢獻.
但我也是最初信仰百川的人.
沒有甚麼貢獻.
是因為我最初期寫了幾篇文章.
然後就下班了.
其實我最大的貢獻就是.
第一代信仰百川的Logo是我做的.
很多人都不知道.
因為我本身是做設計的.
我也要補充.

$^{161}$其實信仰百川創辦.
就是Daniel.
今天就不在我們當中.
他想低調一點.
張廣棟先生博士.
其實他當時也很好.
我也是被邀約.
我想大家也是被邀約.
成為百川的一個作者.
我記得初期好像是五十多個作者.
到後期可能加加減減.
也是差不多.
我就發覺.
同上百川的作用.
我也不重複.
到後來.
尤其是中期.
開始百川放在Facebook.
因為他本身的設定就是.
在文章上其實也可以留言.
但比較少人留言.
但在Facebook.
我發覺Facebook的發展很快.
最後好像是二萬多.
變成文章發了出來.
就會引起很多討論.
還有很多人引用.
很多話題也會碰到.
就像Tony所說.
不是每個話題所有東西的文章.
我們都會同意.
但是我想那個空間.
我想百川是.
可以這麼說.
真的在我們橫教會界.
在這幾年裡面.
是一個很重要的貢獻.
我也覺得是.
所以要肯定的.
這個努力.

$^{201}$我後來就搞頻道.
其實我是想走的.
有一段時間.
百川也跟作者說.
其實有轉型.
想進一步做.
那段時間我跟Daniel說我想走.
他又叫我回來.
他說你post視頻吧.
就是這樣.
所以後來就看到我很多視頻.
所以我覺得.
百川其實給了我們一個很好的平台.
那我就想.
看看大家有沒有一些體會.
除了說百川是.
剛才說有很多.
機會給我們去交流.
或者看到其他人的觀點等等.
我就觀察到一些網路的文化.
一直在轉變.
我就會覺得有些不是很好的文化.
其實我覺得.
現在我形容為.
一個網上的言論.
自由和不自由的時代.
因為自由就好像很多人都可以說很多話.
但是呢.
大家都知道三月.
不知道是幾號.
31號.
大陸都已經禁了網.
所謂.
在宗教上面.
但是我覺得其實要建立一個健康的文化.
就很重要.
在網路上.
其實大家在過往.
或者到這個為止.
你有什麼看法呢.

$^{241}$在網上的一種討論.
無論是文化風氣.
或者你的體會.
是怎樣呢.
大家又說說吧.
我講一點點.
拋磚引玉.
阿礦牧你和我在網上都有一些內容.
包括我想Gordon Wong和Tony都有.
有時候你都會收到一些留言.
你都會收到一些評論.
我的習慣就是.
我是open的.
他罵我都可以.
你可以給不同的意見.
我是會聽.
但是有時候當然我覺得.
如果大家是弟兄姊妹.
我覺得應該有一個基督教的倫理.
不要人身攻擊等等.
分享意見是絕對沒有問題的.
你們讀書的都知道.
有book review.
有preview.
大家可以給意見.
我覺得這個網路上的文化.
好像礦牧你這樣說.
都是需要去重視的.
不只是freedom of speech這麼簡單.
到底我們怎樣發言的時候.
是能夠造就到人.
另外我也留意到在網路當中的討論.
當然正如大家所說.
每一個人都可以有意見表達出來.
我覺得沒有問題.
就算是兩個人是相反的意見.
而去表達出他們的看法.
都沒有問題.
但我就留意到可能因為網路上的設計.
各方面變成和我們一般聊天很不同.

$^{281}$譬如我和你聊天的時候.
我就可以詳細分析告訴你.
我怎樣看.
或者我不是怎樣看.
但變成有時候因為網路的載具問題.
有時候那些回應一來變短.
二來可能就令人產生一種誤解.
變成說了一句就走了.
其實就不是很清楚你想說什麼.
又或者因為他想說一句話的時候.
你就會很不清楚他想說什麼.
這個就好像缺少了比較系統的去說自己的觀點.
這是一件事.
另外一件事我也觀察到.
我想剛才Mark Sir也引到出來.
我會在一種英文叫Courteous.
就是在網路上的一種禮貌.
以及在網路上怎樣去表達.
我想這件事我也留意到在網路文化當中.
怎樣說呢.
我會覺得大家不知道是不是因為網路.
你可以用一個秘密身份.
變成你可以說很多東西.
但因為你不需要負責任.
也可以說很多東西的時候.
其實一些說話.
你自己看完就是.
如果我真的掛著真面目出來.
你會不會說呢.
有時候我自己看完這些.
未必是別人留言給我的.
我看別人留言給別人的時候.
這個在我心底裡有時候會覺得.
你會不會願意為你所說的話負上一個文責呢.
這個我自己看到.
我想Tony一針見血指出了這個問題.
在網路上因為不用出名.
所以看到有很多近乎是盧文式的行為.
即使在基督教的某些網站裡也是.
有些人說話就是沒有理由的.

$^{321}$他要說就說.
說了之後又不負責任又不解釋.
因為他不用出面.
還有這種文化.
我不想說名字.
我覺得在基督教裡面其實是很早出現的.
在近二十年香港某幾個機構.
所謂基督教機構.
搞一些東西的時候.
當他們被人懷疑他們做事的誠信那些.
就有一班不知道什麼的打手出來護航.
這個是一個非常危險的事情.
可能這個背後也是受到大環境影響.
為什麼有基督教流氓呢.
因為有其他網路流氓.
這種方式滲透了.
教會界其實就是一個小社會.
我猜也是這樣.
我自己也收到很多不同的留言.
自從做YouTube之後.
我就剷了它.
也算是有個自由.
你可以剷掉它.
或者封鎖了那個人.
其實我的準則就是.
你說的話有沒有道理.
不是說你不能說.
我看到兩位幾位都是很有胸襟.
罵我都沒有問題.
不過他有沒有道理.
我就說道理.
有道理我就改.
或者檢視.
那種禮貌尤其是差的.
因為匿名或者說話不負責任.
甚至有些異端.
還有一件事我想說.
有些異端留言.
有些hidden agenda.
那些我發現了.

$^{361}$我都會封鎖的.
這個我跟你也是一樣.
我也發現有些異端是想做其他事的.
當然那些一定要剷.
是啊.
有時候我發現了一件事.
看別人留言給我也好.
或者這些留言當中.
我也留意到一個現象.
就是很多時候有些留言.
我也說了.
討論大家觀點不同是沒有問題.
但是有時候你會發現.
其實你會問一個問題.
究竟是持什麼事實.
什麼事實去做討論.
有些討論.
甚至有時候看某些的.
怎麼說.
一些的post.
你就會突然間有些.
退一步去問究竟.
究竟這個是不是fake news.
或者是說的話.
有時候看的時候.
你報導一件事.
你有沒有去很清楚地知道這件事情.
又或者是.
有一種我留意到的.
就是一種叫爆.
好像在爆別人的私隱的.
某一些的報導式的言論.
有時候在網絡我也看到.
有時候我也會退一步.
除非你自己能夠肯定.
你想說的話.
就好像水門事件.
這麼真實又這麼震撼.
否則的話.
你自己要留意小心一點說.

$^{401}$因為這些有時候.
我在網絡上也會看到.
有時候我也會問.
可能很多人都很想做水門事件的記者.
但是你要知道.
你拿著那份東西.
是不是真的這麼厲害.
這個我有時候也會問的.
我可以說.
這件事就是近來激動我.
我彈起的事情.
說疫苗那些.
如果你有看我facebook的.
你會看到我很忙的.
要駁那些.
拿那份報告出來亂寫那些.
你是醫生.
你是醫生你可以說得到的.
有些什麼可以多說一點.
讓大家知道.
找醫生說說而已.
找醫生說說.
有一份報告.
英國出了一份報告.
說疫苗怎麼樣影響免疫系統.
他錯誤地說了一句.
段宗翠說了一句.
剛剛這兩天又有一份說.
拿了美國的輝瑞.
有一個報告又說.
有些什麼厲害的副作用.
其實那份報告也不是這麼說的.
我覺得你有什麼立場都不要緊.
醫生其實很著重的.
就是病人自決.
所以你喜不喜歡打針.
那些沒問題的.
因為起碼在北美.
說病人自決是很重要的.
醫生不可以監病人做事.

$^{441}$但是你不可以拿那些.
那份報告明明不是這麼說的.
你不可以硬要冤枉的.
我經常都覺得.
我尊重每一個人的立場.
無論你相信什麼都好.
或者你接受哪一個理論都好.
只要你不要硬要亂來.
你講開一點.
Dr. Gordon.
我那時候想過.
疫苗有一個很大的爭論的時候.
我想過叫你來我教會那裡.
講一下講座.
評析一下一些混亂的意見.
那就主意現在就沒事了.
還有另外一件事.
我剛才留意網路文化.
我不知道大家有沒有留意.
我留意到一件事.
其實在這段日子的時候.
我們都有些網路.
突然間靜了.
我想尤其是在廣東話的網站.
我不講其他語文.
廣東話我們比較熟一點.
我也明白.
因為可能對某些網站.
或者有些情況.
其實都很大壓力.
所以變成我留意到.
以前有些網站是很活躍的.
或者一些言論是平時以前見到的.
但近兩年就好像消聲匿跡.
或者是沒有了聲音.
我會看到.
這個也是很明顯.
是有一個政治的影響.
我也會覺得.
雖然我自己搞這個頻道.

$^{481}$不傾向講太多這些.
也很容易被人標籤了.
我也要小心.
但是有些是事實.
很明顯這兩年是因為某些原因.
很多人會自我審查.
或者暫時不發聲.
所以百川一完結.
很多人都會覺得.
連百川都沒有了.
這是一個很可惜的事.
其實在網上文化.
空間可能在某些地區.
已經少了.
而且人們又沒有禮貌.
一種風氣就是亂說話.
而且很容易相信假新聞.
尤其是令人氣憤的基督徒.
我想你們北美也很多.
我們澳洲也是.
基督徒領袖.
是的.
領袖才會令人氣憤.
又相信假新聞又傳假新聞.
沒有分析任何東西.
令我們都覺得很困難.
所以我想醞釀.
回到主題.
百川結束了.
但是我們醞釀到現在.
去做港通之谷.
可能真的有一個需要.
或者我們有些.
可不可以說抱負.
哈哈哈哈 這樣啦.
\newpage



\section{}
\label{sec:B4MkD7Z7rb0}
\textbf{從「信仰百川」到「港風之谷」(下):建立健康網絡信仰文化的努力}
\newline
\newline
連結: \href{https://youtube.com/watch?v=B4MkD7Z7rb0}{\texttt{https://youtube.com/watch?v=B4MkD7Z7rb0}} ~~~~ 語音日期: 2022-03-16
\newline
\newline
\hyperref[sec:ng1CRwPs7JE]{\small{< < < PREV SERMON < < <}}
~
\hyperref[sec:index_chronic]{\small{[返順時目]}}
~
\hyperref[sec:index_scriptual]{\small{[返順卷目]}}
~
\hyperref[sec:IQh3_NuFQqI]{\small{> > > NEXT SERMON > > >}}
\newline
\newline
$^{1}$其實港福之谷是怎麼開始的呢?.
Gordon可不可以講一下?.
是Daniel找我去港北日資有一個這樣的構思.
問我願不願意參加.
因為我現在是半退休.
所以我的時間上就比較彈性多一些.
我就很樂意去參加.
我個人的構想也是覺得.
上百川結束了就非常之可惜.
教會真的需要一些這樣的平台.
還有可以作為一個.
在我自己的構思裡面.
就是一個資源中心.
就是idea的資源.
可以分析一些不同的idea出去給人看.
所以我就答應了去參加這個工作.
其實其他兩位你們的角色是什麼呢?.
其實Gordon你的角色是什麼呢?.
我們都忘記介紹一下.
我們在港福之谷的角色是什麼.
我現在的名字就很厲害了.
是Executive Officer.
其實就是做打雜.
一腳踢.
他還要那些東西.
我就是會看著所有的.
我想我自己比較熟悉做多一些的.
就是做一些Policies and Procedures.
Operation那邊的東西.
Technology方面就一邊學一邊做.
比如怎樣設置一個平台.
怎樣設置一個網站那些.
在Content方面.
我覺得我可以有一點.
從一個評訊圖的角度.
帶一些idea進來.
因為大家都是牧師.
我也有做牧師的女婿.
所以我也知道做傳道人.
是有一定的方法去想事情.

$^{41}$還有他們做事情也是有些限制.
我作為一個讀過一下神學的評訊圖.
我就希望可以帶一些.
從評訊圖角度的看法進來.
如果是這樣的話.
我也是會在Gordon的角度工作的.
不是的.
是我會報告給Board的.
我會報告給Board的.
Marsha你是什麼角色.
其實角色很少的.
如果可以分享大家的工作.
我就很樂意的.
真的.
在這個角度我是在Gordon的角度工作的.
Tony可能做了很多事情.
其實我自己也是在信任白宣文.
因為結束之後.
我心裡也覺得.
其實真的需要一個平台.
起碼一方面能夠將不同的.
關於基督教的資訊.
廣東話的基督教資訊能夠帶出來.
二來我更加想也希望.
這個平台能夠成為一個.
一個保存一些香港的基督教信仰文化優越的地方.
這個也是我當初的思想.
所以當Daniel和我談的時候.
最後原本是想做顧問的.
顧問在談的時候就知道有需要.
最後就臨危受面的組成了一個董事會.
我就是董事會的主席.
其實我隨時都願意交出來的.
你做的事情最多了.
Marsha也是我們其中一個董事.
我現在做的事情就是幫港豐支局和總幹事一起合作.
其實日常的工作是有限的.
而我自己其實是設好一個平台.
以及設好一個比較清晰的方向.
以至港豐支局前面的路怎麼走.

$^{81}$然後行政總裁給我們意見.
甚至是接著.
某些時候是比較大層面的.
也會盡量幫忙推動這個工作.
我很欣賞大家做很多背後的事情.
大家要知道開一個平台.
其實背後有很多事情要做.
只是聯絡一些人已經是一個全職工作.
我們現在聯絡的還不只是自己教會的人.
不只是自己城市.
是全球的.
當然我們聯絡過一些人.
那些人可能不便去到公開身份.
我們也知道那個氣氛.
無論是政治社會等等.
可能他們不是很方便.
我們絕對了解.
當然大家要知道.
設置一件事情是需要很多很多時間.
不是一件事最理想做到的就一定做到.
但我最欣賞大家先把這件事踢掉.
同意Mark Sir所說.
當我們找一些董事或各方面需要的人的時候.
當然我們想各個範疇都有.
什麼人都有.
我們都很希望.
但如果我們量度好才開發.
我們都怕那個就已經.
可能是安息主懷.
我不敢說是安息主懷.
不過可能氛圍又不同了.
所以永遠開始.
我經常說永遠開始.
先開始然後再內部修改或調節.
做得好一點.
所以我想其實我們很清楚.
我們知道裡面有很多限制.
或者有很多地方可以改善.
我們都在聽意見.
但我經常有這個想法.

$^{121}$我情願有一個不斷修改的論壇.
可以讓大家繼續發表.
總比我們一直都要自我完善.
但大家沒有了論壇.
到我們做好的時候已經沒有了空間.
我覺得就走失了時機.
我也說一下.
我也是小角色.
每個人都是小角色.
我的名字叫顧問.
說起來好像很誇張.
其實我們剛才都說了.
角色有什麼職責和做什麼.
其實很簡單.
我沒什麼時間有什麼貢獻.
所以變成顧問.
為什麼變成顧問.
可以叫做僅有的角色.
可以願意公開身份.
其實都很重要.
所以我希望可以貢獻多一點.
因為其中一個.
其實大家都忙的.
我也忙的.
但是因為港風之谷是用討論區來做.
我就不玩討論區的.
老實說.
所以我對這件事是非常不熟悉.
所以我在參與方面.
最初我真的不知道發生什麼事.
我只知道可以有些延續.
我想討論到我們搞一個討論區.
其實有些人覺得有點老派.
怎樣可以幫到這件事.
可能保留信用百川的精神.
但是也可能是改善我們剛才所說的.
不太好的網路風氣.
大家可不可以分享一下這方面.
剛才Tony提過一點.
以往百川可以做到的就是發表文章.

$^{161}$發表一些資源等等.
但是平台的性質不是互動性.
也不鼓勵任何互動性.
所以剛才我們也說過.
以前百川的文章.
多多點擊也好.
多多人看也好.
變成互動性減少了.
意見未必是作者能夠找到.
很多時候討論的平台就變成了社交媒體.
Facebook或其他地方.
他們發佈在自己的群組.
他們就可以討論等等.
所以我想這個廣峰的討論區.
不單止可以搜集到這些資源.
拿回這些文化.
文章等等.
那些互動的意見也可以放在裡面.
希望能夠成為一個更加寶貴的寶庫.
其實現在我們只是利用一個討論區的平台.
來做以前百川發表文章的工具.
但是在這方面就更加加強了討論的功能.
因為我們現在的創辦人可以一樣發表文章.
在公眾的地方.
我覺得暫時的樣子就沒有上百川那麼帥.
不過看看遲些能不能弄得更好.
但是剛才Mark Sir和Tony也說過.
發表了文章之後上百川的回應是比較弱的.
功能比較弱.
所以在Facebook也做了.
而Facebook是做得到的.
不過問題就是Facebook有很多限制是我們控制不了的.
它的algorithm或者將來會怎樣刪人.
那些是比較難控制的.
我們希望可以將討論搬回自己的平台裡面.
就可以給我們多一點互動.
以及不受Facebook的限制.
還有對一個組的對比.
相信我們當中也有不少人參與在其中.
但是也可能會有人有一個感覺.

$^{201}$是一種感覺而不是真實的.
有時候是憑感覺行事.
有一種感覺就是你們是不是在玩泥沙.
在亂吹一輪就算了.
自己為之為為.
但我想其實設計了這個論壇區.
第一你不能為之為為.
因為你進來就可以給意見.
發表你的一些想法.
但某程度上是要控制的.
你可以有意見.
當然有些規矩要守.
這就是另一方面要幫助.
我們也要學習在過程中.
可以有禮貌地或者是可以很有 Courteous 底下去作討論.
另外為什麼會有一堆顧問.
有些人有時候會想.
你們是不是又玩小龍基因.
或者是玩一些大腦主義.
其實我絕對不是.
因為有一些顧問的原因.
正是想讓進來看到網站的朋友們.
尤其是基督徒.
不論你是基督新教,天主教或者是東正教.
基督宗教的大圈子.
起碼知道這個網站其實有一班有信仰的牧者.
是 endorse 我用了一個字.
變成某程度上.
其實這個 endorsement 不是來控制論壇的討論.
其實相反你告訴進來的論壇.
這是一班有牧者支持的論壇.
大家在網上看到.
其實網上有很多異端論壇.
掛著基督教的名字.
一上去看.
其實你真的不知道它是什麼料子.
某異端大家都很熟悉.
不說它的名字.
不幫它做廣告.
我有一次是上網找一些見證.

$^{241}$打基督教見證.
一看裡面起碼有三,四成是某一個異端.
對於某些人來說其實都是一個擔心.
我想這班顧問其實都有一個作用.
就是他們是.
你可以說是起碼這個網站的設計裡面.
我們是一班有信仰的人.
當然討論的內容我們是開放.
給不同人有一定他們自己的看法.
我們是有一個開放的心.
起碼那個核心都是一班.
起碼不是一定是牧者.
其實是一班有肯定的信仰.
大家出來做這件事.
同一時間因為知道各種氛圍的緣故.
我們有一些服事的人.
各種原因是比較需要低調服事.
我們都有這樣的空間.
其實正正可能跟某些人說.
是不是大佬文化就剛剛相反.
其他有沒有補充?.
關於我們廣風之谷的討論區.
比如怎樣加入.
怎樣讓弟兄姊妹知道.
加入就很簡單.
你去我們的網站.
可以在Facebook裡面按連結.
或者直接去網站.
就是hkchristians.world.
這樣就會去到網站.
再點進去就可以進入討論區.
你要註冊.
但這些是差不多所有網絡上的.
都很標準的做法.
加入之後.
你就可以發表你的文章.
或者發表你的評論.
暫時來說.
我們有一個是公開原地的.
公開原地因為人人都可以看.

$^{281}$我們就比較上是限制了.
不是每個人都可以在公開原地發表文章.
你需要在討論區裡面發表文章.
如果你覺得那篇文章是想公開發表的.
你可以跟我們任何一個版主.
或者任何一個管理員聯絡.
我們就可以將文章搬過去.
我們信仰百川以前的作者.
我們很多時候已經設定了.
一個自動發表的權利.
他就可以在公開原地繼續發表文章.
就不需要再經過其他人幫他搬過去了.
好處就是.
每一個想發表的人都有一個渠道.
以前信仰百川如果你想加入的話.
他不邀請你比較難找到誰去聯絡.
除非你已經認識幾個人.
他就跟你接觸這樣你可以加入.
但是現在反而加入的渠道是簡單一點的.
如果你是一個新手的話.
如果你現在是信仰百川的作者.
應該沒有問題的.
應該你可以馬上繼續發表的了.
我們都好像會將一些舊的文章發表出來.
因為那時候突然.
完了信仰百川很多人很可惜.
找不回那些文章.
其實我們都是有這樣的想法的.
當然就要看那個作者.
自己決定.
但是就算澳洲這邊也有人馬上說.
很可惜啊.
很想找回那些文章.
我想會有相當的數量的作者.
和他們的文章會按次發表出來.
我相信這個也是可以服務到我們以前的目的.
剛才說了很多.
最後還有沒有一些補充.
或者在這個題目裡面分享的呢.
我想我一樣.

$^{321}$作為董事會主席.
其實就是一個呼籲.
如果你是一直支持開信仰百川的朋友們.
我也希望你嘗試給自己一個機會.
走進這個新的論壇看一看.
看看自己會不會繼續在這裡參與.
如果你覺得這個挺好的.
也值得我們支持的話.
我也很鼓勵你.
希望將剛才Gordon所說的網址.
和你身邊的朋友弟妹姐妹再去推介.
也很想和其他的廣東話.
或者中文的教育圖的網站的領袖們說說.
這是一個新的時空.
希望大家彼此鼓勵.
也歡迎你們來看.
給一些意見.
還有這個新的論壇.
希望大家進來看的時候.
已經註冊了.
可以發表.
還有大家知道.
我們希望功能和作用.
比以前更加enrich.
不會take away了以前的作用.
更加會加上一個互動性.
我們希望這個平台.
能夠儲存一些好的資料.
並且大家可以有更加健康的.
建設性的一些對話.
我鼓勵我們現在在看的.
Tube的所有的弟妹姐妹朋友.
多點去visit.
剛才我說我對討論區這件事很不熟悉.
我後來也很忙.
他們做了很多很努力的事之後.
我就進去看了.
我發覺也很user-friendly.
也不難用.
真的好像剛才所說.

$^{361}$可能現在版面還不是最好.
但是他真的可以serve到.
以前的purpose之外.
更加是一個加強版.
我們更加不單想加強版.
是一個相同的網路文化.
透過我們這個討論的平台.
健康一點去發展.
所以大家為這件事去祈禱和參與.
不如這樣吧.
我們也談到剛剛八個字.
到這裡為止.
多謝Mark Sir 城會網台.
多謝Tony 多謝Gordon.
有機會大家也可以給一些意見.
或者在我們YouTube留言.
或者上去討論區說說話也可以.
為了這條片.
好,謝謝大家,到此為止了,拜拜.
\newpage



\section{}
\label{sec:IQh3_NuFQqI}
\textbf{Family Sunday School - EP 5 (Jesus Goes to Trial)}
\newline
\newline
連結: \href{https://youtube.com/watch?v=IQh3_NuFQqI}{\texttt{https://youtube.com/watch?v=IQh3\_NuFQqI}} ~~~~ 語音日期: 2022-03-21
\newline
\newline
\hyperref[sec:B4MkD7Z7rb0]{\small{< < < PREV SERMON < < <}}
~
\hyperref[sec:index_chronic]{\small{[返順時目]}}
~
\hyperref[sec:index_scriptual]{\small{[返順卷目]}}
~
\hyperref[sec:XtJhkYaLq8U]{\small{> > > NEXT SERMON > > >}}
\newline
\newline
$^{1}$[music].
Hello everyone! I'm Miss Rebecca..
Welcome back to Family Sunday School..
This week is the fifth episode..
Do you remember what we talked about in the fourth episode?.
Yes, we get to know a lot of different animals.
and make a lot of different sounds..
But what's the most important thing?.
The dog!.
What is Peter doing?.
Yes, he is rejecting Jesus..
Did you notice that I'm wearing different glasses?.
Do you remember?.
In the last episode, I asked my parents to read other Bible verses..
I asked them to see if they could relate to the story I told..
Why do we need to see so many different verses?.
So that we can see them from different angles..
For example, I'm wearing these glasses..
I'm a little blurred when I'm looking at things..
Why?.
Because of the eyebrows, the nose, and the nose..
But if I wear glasses or wear clothes,.
will it be more clear?.
I'm wearing these glasses, but they are a little dirty..
I can't see very clearly..
I change my clothes, but I can see different things..
This pair of glasses is so clear..
When we turn our glasses, we can see different things..
Sometimes it's a little blurry, sometimes it's clearer..
Or something is magnified..
Because there are many different glasses for different purposes..
For example, when parents or elders are reading,.
they are a little blurry..
They can't read when they are close to reading..
But with glasses, they can read..
So, we magnify..
Right?.
So, today, let's see if you really read other Bible verses..
Did you find anything?.
Any new discoveries?.

$^{41}$Tell me..
No?.
It's okay..
I read it too..
But, Miss Rebecca has a lot of different Bible verses..
Right?.
When Miss Rebecca didn't talk about it in the last episode,.
she didn't mention one thing..
She didn't mention the cock-a-doodle-doo..
Peter looked at Jesus..
Jesus looked at Peter..
Oh, he really did it..
Miss Rebecca didn't mention it..
How?.
In that book?.
This book also mentions it..
What does it say?.
Peter cried and ran away..
He wept and fled the courtyard..
He cried and ran away..
So, when he left the courtyard,.
he didn't know what happened..
So, this is one thing..
He left..
Right?.
Who did it?.
John!.
Right..
There's one more thing..
Miss Rebecca found it in other Bible verses..
Do you remember?.
There were three times when people asked Peter,.
"Do you know Jesus?".
"No, I don't know him.".
"I don't know what you're talking about.".
Right?.
There were two ladies..
One lady asked,.
"You seem to be with him.".
"No, not me.".

$^{81}$"No, you really seem to be with him.".
The second lady said,.
"No, I don't know what you're talking about.".
Then, a little further away,.
there was a man..
He said, "No, I heard your accent.".
"Accent" means "accent.".
Sometimes, when we come from different places,.
we have different accents..
Even though we're speaking English,.
when we use English with different accents,.
we'll know where we're from..
Right?.
So, he heard her say,.
"No, you're from Galilee.".
Because in Galilee,.
the way they speak and their accents are a little different..
They can hear..
Right?.
Ms. Birkin didn't mention that..
So, there are many hints.
to let them know that.
Peter is with Jesus,.
but Peter is a pharaoh..
What will we talk about in this episode?.
Jesus went in..
He went into his house.
with the elders..
These elders are the highest leaders of the Jews..
Do you know that?.
When they went in,.
they had to judge him..
Judge him..
Trial him..
See what crimes he committed..
Those who arrested him.
thought he committed all kinds of crimes..
But is it true?.
Today, we'll see.
what happens when Jesus is judged..

$^{121}$But before we start,.
what do we do?.
Of course, we play a game..
Okay..
Let's play a game..
Get ready..
I have a musical instrument..
I will play this musical instrument..
This is called a xylophone..
In Chinese, it's called a wood instrument..
In the past, it was made of wood..
It made different sounds..
This one is made of iron..
When I play,.
you have to throw the ball..
You can find a ball or a toy..
You can find a leaf..
I play and you throw..
You don't have to follow the order..
Throw to the person you want to throw..
Throw again..
The object is to not hit the ball..
I play the wood instrument..
After the music,.
whoever holds the ball is out..
Okay?.
We'll do four rounds..
I'll do four rounds..
I don't know how many people are there..
Let's start..
Here we go..
[xylophone playing].
Done!.
Whoever holds the hot potato.
is out of the round..
Okay, next round..
Ready?.
You can switch places..
You can take it away..
You can do whatever you want..

$^{161}$Okay, second round..
Here we go..
[xylophone playing].
Done!.
Okay..
Who got it?.
What do you do if you get the hot potato?.
Yes, you're out..
Sorry..
You can continue next time..
Okay?.
It's not a game..
Sometimes we win, sometimes we lose..
[speaking in Chinese].
Okay..
[speaking in Chinese].
Round number three..
[xylophone playing].
All done!.
[speaking in Chinese].
Okay..
[speaking in Chinese].
Ready?.
Let's go..
[xylophone playing].
All done!.
Okay..
[speaking in Chinese].
Okay..
Last one standing..
Winner..
[speaking in Chinese].
Okay..
[speaking in Chinese].
All right..
Have fun!.
[speaking in Chinese].
Bible story time!.
[speaking in Chinese].
Bible story time..

$^{201}$[speaking in Chinese].
Hello Chinese time!.
[speaking in Chinese].
Let me tell you, break it down, break it down..
[speaking in Chinese].
Okay, ready?.
[speaking in Chinese].
At home extension time..
[speaking in Chinese].
Act it out..
[speaking in Chinese].
Okay?.
[speaking in Chinese].
That tells you, they got it..
[music playing].
\newpage



\section{6:1-17}
\label{sec:XtJhkYaLq8U}
\textbf{Ordinary Me (Judges 6\_1-17, 25-26) - Tommy Lee}
\newline
\newline
連結: \href{https://youtube.com/watch?v=XtJhkYaLq8U}{\texttt{https://youtube.com/watch?v=XtJhkYaLq8U}} ~~~~ 語音日期: 2022-03-28
\newline
\newline
\hyperref[sec:IQh3_NuFQqI]{\small{< < < PREV SERMON < < <}}
~
\hyperref[sec:index_chronic]{\small{[返順時目]}}
~
\hyperref[sec:index_scriptual]{\small{[返順卷目]}}
~
\hyperref[sec:pinIq8rUt5E]{\small{> > > NEXT SERMON > > >}}
\newline
\newline
$^{1}$(audience applauding).
- Thank you, good morning to all of you guys..
And for those who are sitting here,.
those who are watching, my name is Tommy Lee..
It is always an honor to really be here.
with all of you guys..
I graduated from business school,.
and then from that point, I spent about 10 years.
building cell towers, and that's what I thought.
I was gonna do all my life..
Then from that point, I was a pastor for about three years,.
actually working in a Korean church..
And I was there, I realized,.
and Pastor Mark does a great job as a pastor,.
I realized being called to be a pastor.
just wasn't something that God was calling me to do..
I love the Lord, I wasn't called to be a pastor..
And then after that, I got hired at this big Bible institute.
called Moody Bible Institute here in the city..
So I worked for the president's office, Dr. Michael Easley,.
I said, "Okay, this is what you're calling me to do, Lord.".
And I realized, after spending about four years there.
working with the bigger ministry, it was very top-down,.
it probably wasn't for me as well, too..
And so when Dr. Easley transitioned out,.
I transitioned out with him..
So now these days, I get a chance, I realize,.
I fit in this thing where I could do multiple things..
We started Resource Global that works.
with young Christian marketplace leaders.
in Singapore, in Malaysia, in Nairobi, in South Africa,.
and Austin, Texas, and also in Indonesia..
Now I also then got hired to work with Northern Seminary,.
where I get a chance to partner with Pastor Mark.
under the Growth Center for Church and Missions,.
their executive director..
Then I got hired by Family Foundation.
to run Together LA initiative over in Los Angeles,.
and most recently as the managing director.
of the Wheaton Center for Faith and Innovation.

$^{41}$at Wheaton College..
And I realized I actually do better.
when I could run four different things versus one thing..
And all of us are wired differently..
All of us are in a journey..
I think from what I heard from James last week,.
the speaker talked about what does it mean.
to give of your life, to give of your treasures,.
both financially as well as your time..
Today, we're gonna talk about a story.
that I have loved and I've come to love,.
and it's the story of Gideon in Judges chapter six..
I hope that you guys have not heard.
about this passage recently, and I'm repeating again,.
but man, I love diving in this story.
because it's a reminder of how God uses ordinary people.
like you and me to do extraordinary things..
So let's dive in the story of Gideon..
In Judges chapter six, verse one,.
it begins in verse one, "The people of Israel.
"did what was evil in the sight of the Lord.".
And you're gonna find in the book of Judges,.
that's a theme over and over again..
The Israelites are always doing what was evil.
in the eyes of the Lord..
They're always getting in trouble..
They're always doing something that forces,.
no matter what God has done to them,.
from the days of Moses to the days of Joshua,.
and judge after judge, you see the same theme.
over and over again, that the people of Israel.
did what was evil in the sight of the Lord..
"And the Lord gave them into the hand of Midian.
"for seven years," I'm reading out the ESV..
"And the hand of Midian overpowered Israel..
"And because of Midian, the people of Israel.
"made for themselves the dens that are in the mountains.
"and the caves and the strongholds.".
They had nowhere to turn, they just hid there..
"For whenever the Israelite planted crops,".

$^{81}$whenever it was harvest time,.
whenever they knew it was harvest time,.
"the Midianites and the Amalekites,".
not just the Midianites,.
but they brought along the Amalekites,.
"and the people of the east will come up against them.".
So already they're scared, already they're hiding in caves,.
already they're hiding in all of these different things..
Harvest time, they're invaded by the Midianites,.
the Amalekites, and all the people in the east..
"They would encamp against them.
"and devour the produce of the land as far as Gaza,.
"and leave no sustenance in Israel,.
"and no sheep or ox or donkey..
"For they would come up with their livestock in their tents..
"They would," I love this description,.
"come like locusts in number,.
"both they and their camels could not be counted,.
"so that they laid waste the land as they came in..
"And Israel was brought very low because of Midian,.
"and the people of Israel cried out for help to the Lord.".
In today's culture, we see a little bit of that, don't we?.
We see recently a lot of times the terror.
or the violence that's happening,.
and I am, here we go, in Russia and the Ukraine..
You see the huge city or huge country that we see in Russia..
You see tiny Ukraine that was once a part.
of the Soviet Union, the USSR back in the '80s..
You see Russia trying to put together.
the former Soviet Union and trying to exert their power..
They are already overrun and captured Georgia quietly..
And now here, Ukraine, right before they became,.
they applied to be members of NATO..
Here they are, and for the last month,.
they're bombing, they're bombarding all of these things..
Take a look at this map here and there..
Look at what they're trying to do..
And for them, it's a show of power..
And it reminds me of the Midianites..
Here the Israelites are sitting there,.

$^{121}$hiding in their strongholds, hiding in their caves,.
and they are overrun by an enemy that's far larger to them..
It is a show of power..
We see this clearly here in Russia.
with these pictures here and there..
It is Putin, their president,.
and that big table symbolizes how he perceives power,.
that there is a distance to how he sees power..
And it shows that I am in charge.
and you have to submit to me..
And when we look at this particular thing,.
God gave the Israelites into the hands of the Midianites..
And why is that?.
Because God uses events in our lives to get our attention..
There's sometimes in the midst of God calling us.
in our lives, we go in a misdirection, we get distracted..
We start giving into our own personal agenda..
We start doing the things we wanna do because it's my life..
I've worked hard, it's my decision..
No one could tell me what to do..
I dictate, it's all about me..
But when we become Christ followers,.
that rule is out the window..
Our life is no longer ours..
Our decisions are no longer ours..
But we submitted to the gospel, we submitted to the cross..
And along the way, God uses events in our life.
in order to get our attention, to bring us back to him..
And we see this clearly when he gives them.
to the hands of the Midianites.
because they keep turning their life to the Lord..
And sometimes in the midst of our life,.
God loves you so much..
God loves and wants the best for you.
that he will use hard times and trials in your life.
to turn you over to him..
All of us will deal with trials in our life..
It is not a matter of if, but rather a matter of when.
that God uses those trials to turn our life..
And for the Israelites, it had gotten to the point.

$^{161}$that they were so low because of Midian.
and the people of Israel cried out for help in the Lord..
When the people of Israel cried out to the Lord.
on the count of the Midianites,.
the Lord sent a prophet to the people of Israel..
And he said to them, "Thus says the Lord, the God of Israel..
"I led you up from Egypt.".
This is what I did from you..
"I brought you out of the house of bondage..
"And I delivered you from the hands of the Egyptians.
"and from the hands who had oppressed you.
"and drove them out before you and gave you their land..
"And I said to you, I am the Lord, your God..
"You shall not fear the gods of the Amorites.
"in whose land you dwell,.
"but you have not obeyed my voice.".
Here's a key one, verse 11..
Pay attention to this one.
because I really want you to focus on this..
"Now the angel of the Lord came and sat.
"under the terebinth at Opera,.
"which belonged to Joash the Abysserite,.
"while his son Gideon was beating out by impress.
"to hide it from the Midianites..
"And the angel of the Lord appeared to him.
"and said to him, the Lord is with you..
"Oh, mighty man of valor.".
If you have the NIV, it says mighty warrior..
If you have a pen, if you highlight things in your Bible,.
you take this right away and you highlight it..
And you remember this verse because when you look at this.
and God calls him a mighty man of valor,.
already you're sitting there,.
wait, what in the world is this?.
That doesn't make any sense..
Why is that?.
Let me show you something..
Pastor Mark, let me just say this..
You teach this class, right?.
All right, mighty warrior, look at this..

$^{201}$When the angel of the Lord, NIV says mighty warrior..
Take a look at this..
Pastor Mark, you teach a class on archeology..
He probably knows this a lot more..
I've been to Israel once in my life,.
but it was through that experience,.
I learned so much..
So let me explain to you the situation.
that is happening, right?.
It is harvest time..
And now when the wheat and grain.
and everything's ready to be harvested,.
you bring all of that, you use your sickle.
and you cut out all of the wheat.
and you bring it to what they call as a threshing floor..
As you could see, it is out in open air..
It is usually on top of a hill..
It's usually in an open area,.
in like a little hill area where the wind can blow.
and you lay everything out..
If you look at the picture, you see some of the oxen.
or some people will use a little hammer called flail..
And you're sitting there knocking.
or you have your oxen, your horses step all over it..
At that point in time, it separates the husk.
as well as the grain and it just separates it..
Then from that point on,.
you use what they call a windowing fork..
It's a big rake..
And you've seen pictures of this..
You take it and you lift it up as high as you can..
It's hard work..
The chaff is blown from the wind and everything like that.
and it blows in the air..
The husk, which is heavier, fall in one pile..
The grain falls in another pile..
And so as a result, you have to understand.
the threshing floor has to be out in the open air..
But where is Gideon?.
Gideon is in a wine press..

$^{241}$You're wondering, so what's wrong with that?.
Let me show you a picture of what the wine press looks like..
This is the wine press..
I get to the Garden of Gethsemane, right?.
And it's a garden, no, it is a tomb..
The tomb where Jesus is buried..
And there right in front is this wine press..
And I didn't understand what a wine press is.
until I saw this picture..
And the story of Gideon made so much sense..
Here's what happens in the wine press..
Right there on the top, they lay out all of the grapes..
From that point on, you see that hole.
that leads all the way to the bottom..
And that path then leads to another hole,.
which there's a bunch of jars or what they called cisterns..
Once you lay out all the grapes, it's about waist high..
Mark, I hope you're gonna have to correct me.
and let me know if I'm wrong later on..
And from that point, you lay out all the grapes.
and you step on it..
I've seen pictures where there's a contraption on top.
and you hold onto it,.
and you just keep stepping on all of the grapes..
And then eventually all of the juice flows.
from down that corridor and into what they call a cistern..
And so what's the picture of this?.
It is this..
The angel of the Lord calls him a mighty warrior..
He is not a mighty warrior..
Because when you read this passage, you're sitting there,.
Lord, are you sure you have the right picture correct?.
Do you not know where Gideon is?.
He's threshing wheat, not on the threshing floor,.
but in the wine press..
And you know why he's in the wine press?.
Because he's afraid..
Because he's a coward..
Because he's running away from the Midianites.
so that no one sees him..

$^{281}$Are you sure you have the right person, Lord?.
Mighty warrior?.
Mighty man of valor?.
I'm seeing a coward here..
And why is that?.
God's picture of me is different than the picture.
I have of myself or that others have of me..
God's picture of me is different than the picture.
that sometimes I have of myself.
or the picture that I have, that others have of me..
Let me give you an example of what I mean by this..
All of you guys have probably watched this movie, right?.
All of us are familiar, Avengers Endgame,.
and we're sitting there and they go meet Thor.
and we're expecting the God of Thunder,.
the one that we were familiar with all of the movies.
coming out and what do we see?.
Fat Thor playing video games..
And you're sitting there, what in the world happened?.
Because he had given up hope..
He was done..
But something saw Thor as the adventure that he was,.
the hero that he was..
Because deep inside, and you saw that encounter.
with his mother, there was a hero,.
despite a person who gave up hope..
Sometimes we do feel like a fat Thor..
Sometimes we are scared..
Sometimes how we perceive ourselves.
is different than how others perceive us..
Sometimes how others perceive us is different.
than how God perceives us..
For us, we see fat Thor..
God sees a mighty man of valor..
God sees in Gideon, not someone who's scared,.
not someone who's a coward,.
not someone who's trying to escape..
What he sees in Gideon is a mighty warrior,.
the one who's gonna lead his people there,.
the one who's gonna battle against him..

$^{321}$Let's keep going on..
And so he calls them, and here's the next part that I love..
And the angel of the Lord appeared to him and says to him,.
Lord is with you, almighty man of valor..
And Gideon said to him, please serve..
If the Lord is with us,.
then why then has all of this happened to us?.
And where are all his wonderful deeds.
that our fathers recounted to us saying,.
did not the Lord bring us up out of Egypt?.
But now the Lord has forsaken us.
and given us into the hands of Midian..
Underline this one as well too, verse 14..
And the Lord turned to him and said,.
go in this might of yours,.
and Gideon, I want you to save Israel.
from the hand of Midian..
Do not I send you..
And he said to him, please, Lord, how can I save Israel?.
But behold, my clan is the weakest in man's sight..
So out of all the clans, Lord, my clan is the weakest..
And I am the least in my father's house..
So pretty much what you're telling me, Gideon,.
is out of all the clans, your clan is the weakest,.
and you're the weakest one in your clan..
That's what you're telling me, right?.
So you know where that makes Gideon?.
Dead last..
He's the weakest person in the weakest clan..
Not very good..
And you're using him to lead the charge.
to battle the Midianites..
Powerful statement, next thing..
And the Lord said to him, but I will be with you..
And you shall strike the Midianites as one man..
And he said to him, if now I have found favor in your eyes,.
oh wait, and I am the least in my father's house,.
and the Lord said to him, but I will be with you,.
and you shall strike the Midianites as one man..
So what the Lord is saying is,.

$^{361}$he comes to him and he says, okay, I am the weakest clan,.
and I am the weakest in the weakest clan..
And you know what the Lord says to him?.
He doesn't disagree with him..
He doesn't tell him, you know,.
no, you're not the weakest clan..
So-and-so is weaker than you..
Here's what you have..
Well, your clan's not the weakest clan..
Your clan's pretty good..
No, he doesn't disagree with him..
In fact, he actually kind of agrees with him..
It's like you and I are saying, right?.
A lot of times is we talk, you guys talked last week.
about what it means to give God our time and our treasures..
It's like when you and I are sitting there,.
Lord, I'm not sure if I could serve you.
because I'm not smart enough..
That sometimes is our excuse..
I'm not talented enough, right?.
I'm not as strong as that particular person..
I can't be as charismatic..
I'm not very good in public like that..
I'm very shy..
And it's, God is saying, I know..
I'm not disagreeing with you, I know..
I know you're probably not the smartest person in the room..
I know probably you're not the most talented person.
in the room..
I know so-and-so probably is better up front than you..
I know you're probably shy..
I know..
I know all these things..
You know why?.
Because I created you..
And when I created you, I gave you those skills..
And when I created you, I made you different..
I made you unique..
I know..
Get in, I know you're the weakest clan..

$^{401}$I know you're the weakest of the weakest clan..
I know all these things because I put you there..
Because when you were born, before you were born,.
before any of us were born, I designed you.
and I mapped out your life..
I know exactly what you're capable of..
But here's the thing, and underline this,.
but I will be with you..
And that's the difference..
I want you to do that..
I know what your deficiencies are..
But when I call you, I will be with you..
And that becomes the key..
It won't be easy..
It's not gonna be always very fun to obey my will..
It's not always gonna be smooth sailing all the time..
But I see you differently than how you see yourself.
and how others see you..
I see you as a mighty man of valor, a mighty warrior..
And I'm calling you to fight the Midianites..
And I realize you're the weakest clan..
I realize you're the weakest in your weakest clan,.
but I will be with you..
And that's where all of us needs to stop and think,.
what does God want you to do in your life?.
And what is holding you back?.
Because all of us have insecurities..
All of us have struggles..
All of us have doubts..
But what is God saying to us?.
And how do we begin to walk confidently?.
Knowing it may not be easy, but walk confidently.
because that's what God has called him..
Which leads us to the next portion, all right?.
Next portion is this..
And so later that night,.
he calls him and let's go over to verse 25..
That night he prepares,.
and here's the, I think the most important part of it..
That night the Lord said to him,.

$^{441}$"I want you Gideon to take your father's bull.
and the second bull seven years old..
And I want you to pull down the altar of bow.
that your father has.
and cut down the Asherah pole that is beside it..
And I want you to build an altar to the Lord your God.
on the top of the stronghold here.".
And here's what it looks like with the Asherah pole..
Asherah pole..
And so what God is saying to this,.
"Lord Gideon, I see you threshing wheat in the wine press..
I want you to fight the Midianites..
I want you to lead the charge..
But before you do, before I send you to battle,.
I need you to tear down that altar..
I need you to put the proper altar in its place..
I need you to tear down that altar..
I need to build a proper altar.
that honors the Lord in all that you do.".
That's what I want you to do before you go out to battle..
If you were to apply it in our context today,.
I have a calling for you..
I want you to do something..
But before you do,.
we're gonna have to do some business in your life..
I'm gonna need you to tear down some of the altars.
in our life, specifically in our hearts..
And I'm gonna need you to put the proper altar in its place..
I'm gonna need you to make sure that things are right..
A lot of times within the altar is within our hearts..
There is no big altar that we have.
that we all go to like we have with the temple..
That altar is in the life..
God establishes that altar in our life..
And I have a plan for you..
I need you to establish the proper altar in its place..
What do I mean about that?.
I love the fact that God has a plan for us..
But frankly speaking, all of us struggles with heart issues..
And I probably was shared a little bit.

$^{481}$about this in the past..
All of us have heart issues in our life..
All of us struggle with skeletons in our closet..
The majority of the group here are Chinese..
I see our two Caucasian friends in here..
And I find that Caucasians are very vulnerable.
in their life..
Asians tend to be very private for the majority,.
especially if you're older, you're very private..
You don't let people hear, you don't show vulnerability..
I find our Caucasian friends are more so vulnerable..
And we like things and we wanna have this appearance.
that we have..
And so we live with what we call a shadow mission..
That our shadow and how we live our life.
are two separate things..
Our shadow reflects who we truly are..
But on the outside, we try to pretend.
like we're someone different..
We try to show people the persona or personality.
we want them to see, which is so different.
than what we struggle on the inside..
Some of the four or five basic heart issues.
that some of us have..
Before God could use you,.
he wants you to deal with some of these issues in your life..
That some of you are sitting here right now.
and you're dealing with the issue of guilt..
That there are things in your life.
that nobody knows that you've done..
And only you know..
You've hidden it pretty well from different people..
Maybe a thought, maybe something you've done,.
something it is and it bothers you..
God wants you to learn to deal with that guilt..
He wants you to learn the power of confessing your sin..
The power of allowing him to take away that guilt.
that is eating up your life all the time..
For some, it's the power of anger..
There is anger that you have because of a parent,.

$^{521}$a sibling, a friend, something that happened.
at work, you got bypassed for a job, whatever it is..
You feel like you got a short end of the stick,.
whatever it is, you're dealing with that anger..
Maybe greed..
Maybe it's greed that you're dealing with..
It's all about me..
I need my money, I want to be successful..
I want to have lots of money, I want that Tesla..
I want that home, I want that, maybe it's greed..
I want my crypto to grow huge..
And then once I get so much money,.
God, I'll come and satisfy you because I need to work life,.
I need to provide for my family..
It's that greed that over, and when you're sitting there.
listening to sermons about giving of God.
or your story, you're sitting there, yeah, I'll do it,.
but I got to make sure I take care of myself first..
Guilt, anger, greed, maybe it's jealousy..
Why does so-and-so always get the benefit of doubt?.
(thuds).
Why does so-and-so always seem to get.
the great opportunities and work?.
Why does so-and-so get this,.
or why does my sibling get this and that?.
Maybe it's jealousy..
And you're struggling with that because you're sitting there,.
why do I get the short end of the stick?.
Or maybe it's judgment..
Judgment of how you perceive people of different color,.
how you perceive different people that you struggle with,.
judgment toward other people that you see in your family,.
whatever it is..
Which one is it?.
All of us struggle at one point with each one of these..
For me in my life, I've watched this destroy my family.
in a great deal..
I have a twin brother and I have a younger sister..
I see my twin brother a lot, he lives in Phoenix, Arizona.
with his three kids and Sarah, his wife..

$^{561}$My younger sister is two years younger than me..
I haven't seen her in 15 years..
She stopped calling our home 15 years ago..
She got into it with my mom 15 years ago..
And that anger between her and my mom never subsided..
She walked away from home, never to come back..
I haven't seen her in any Christmas gatherings,.
any holiday gatherings..
When my two daughters were born, never got a note..
When I was sick in bed with cancer, never heard from her..
Because she can't let that anger go away..
Because why?.
My mom deals with that anger..
My mom growing up as a hardworking Chinese immigrant.
in 1972 came from China and when she came in here,.
she had nothing and she had to start all over..
And there was so much anger and resentment that she had.
that she never can let go of that anger..
And that anger, she took it out on us as kids growing up..
She was always hard on us..
It was never enough..
That A that we got, you better make it A plus..
That B you got, you're grounded..
That belt or that hanger, pick one..
We always had to do that..
The anger and my sister was never able to let it go..
I deal with that anger as well too..
I struggle with that anger..
I could fool you people easily..
I'm 44 years old..
I look decently well..
I dress nice, I have a successful career..
But inwardly, I struggle with that anger..
I grew up and we are in church here in the Bridgeport area.
where a lot of times there's a lot of diversity..
I came here, I was born in '77,.
came to Chinatown and Bridgeport..
And in my school at Robert Healy,.
I was probably one of only five Chinese kids in that school..
For years, I was chased after,.

$^{601}$I remember there was some Caucasian kids.
that would chase us down the block all the time..
They would just beat us up all the time..
Everywhere I go and every street I turn,.
I still remember when I was in fourth grade,.
there was this eighth grader.
who just wanted to pick on us all the time..
I went to one street, he found the stair..
Went to another street, found the stair..
It was rough when you're the only Chinese kid.
in the community..
I still remember there was just one guy.
who would ride on his bike.
and just pick on all the Chinese kids in the community..
I grew up a little chubby,.
I grew up not always being very smart..
And I heard the conversations, my sister was brilliant..
Perfect SAT, perfect ACT, perfect grades,.
straight A's, never got a B,.
went to Yale, got perfect score there..
I struggled with school..
I had to work hard for my school..
And I always heard people say,.
"Hey, you're not as smart as your older brother..
"You're not as smart as your younger brother.".
And I was made fun of at school.
because I was kind of chubby at that point in time..
And that made me live with self-esteem issues,.
which eventually became anger..
And I always lived with that anger.
where I wanted to prove people wrong..
I wanted to stick it to them.
and show them and do things..
As I grew in my faith,.
I had to come before the Lord and say,.
"Lord, where's my heart?.
"Lord, what do I need to do to lay it before the cross?.
"Lord, what does it mean to have my heart come before you?".
I had to learn slowly and slowly to lay it down..
Every day in my life, I still remember all the voices,.

$^{641}$all the people who said I couldn't do it..
I still remember their faces..
That's still there..
But it's a choice I have to make every single day to say,.
"Lord, give me the strength to do what is right..
"Allow my heart to continue to honor you.
"in everything I do.".
Because God wants you to tear down that altar.
and put the proper altar in its place..
God's picture of you is different.
than the picture you have of yourself.
and that others have of you..
God's picture of you is different..
He sees in Gideon a mighty man of valor, a mighty warrior..
What does he see in you compared to what you see in yourself?.
Go in the strength that you have,.
whatever you are,.
because I designed you with your strengths.
and your weaknesses..
Go in the strength that you have,.
and I will be with you..
I'm gonna be with you..
May not be easy, may not go smoothly,.
but I'm gonna be with you..
But before you do,.
put the proper altar in its place..
Tear down the altar.
and put the proper altar in its place..
God uses ordinary people to do his extraordinary work..
Let's pray..
Father, thank you for this church..
Thank you for the people who are sitting here..
Thank you for the people who are watching online,.
wherever they are in the world..
Allow us to hear you, Father..
Allow us to walk with you..
Allow us to continue to understand the vision.
that you have for us,.
to listen, to hear the vision,.
and continue to follow you,.

$^{681}$to understand that you have a calling..
Despite our inefficiencies, despite our baggage,.
you are there with us..
But more so, Lord, I pray for our heart issues.
that with people here in this room..
We pray with some of them who are struggling.
with the guilt or the anger or the jealousy or the greed,.
I pray that you would do surgery in their lives..
And allow us as we wrestle with those.
to have the courage, the strong strength from you.
to follow after you..
\newpage



\section{}
\label{sec:pinIq8rUt5E}
\textbf{放下自己的接待(路14:1–14)- 黃嘉漢牧師}
\newline
\newline
連結: \href{https://youtube.com/watch?v=pinIq8rUt5E}{\texttt{https://youtube.com/watch?v=pinIq8rUt5E}} ~~~~ 語音日期: 2022-03-28
\newline
\newline
\hyperref[sec:XtJhkYaLq8U]{\small{< < < PREV SERMON < < <}}
~
\hyperref[sec:index_chronic]{\small{[返順時目]}}
~
\hyperref[sec:index_scriptual]{\small{[返順卷目]}}
~
\hyperref[sec:CgJV5or0kgk]{\small{> > > NEXT SERMON > > >}}
\newline
\newline
$^{1}$各位電子妹你好.
很久沒有在城會教會和大家有訊息的分享了.
疫情關係我都沒辦法來到你資格講當中.
我們就可以在網路上將上帝的訊息和大家分享.
今天想和大家分享一個聖經訊息之前.
想問問大家有沒有聽過一個四字詞語.
叫做「其貨可居」.
什麼是「其貨可居」呢?.
原來這個四字詞語的典故.
出自於秦國的宰相呂不韋.
講故事吧.
當呂不韋是一個商人的時候.
他去到當時戰國時代的趙國做生意.
在那裡他認識到一位公子.
就是秦國的公子.
不過這個秦國公子很慘的.
他被派到趙國去做人質.
當呂不韋見到這個秦國公子的時候.
他就覺得「其貨可居」.
因為他認為可以在這個人的身上帶來大的作為.
所以他就花盡了精神,金錢,時間,人脈.
去幫這位很窮困的公子上位.
幫他可以爬回到秦國當中的位置.
最終這位秦國的王子.
都跟著呂不韋的幫助緣故.
他就成為了秦國的莊襄王.
可能你不認識他是誰.
但你一定聽過他的兒子的名字.
他的兒子就是秦始皇.
呂不韋因為幫了這個秦國公子成為莊襄王之後.
自然得到很多好處.
他不但成為了秦國的宰相.
風雲信候.
而且在河南一帶有很多戶籍.
你可以說又有財又有權.
的確很坦白來說.
這個人生的故事我們看到的就是.
對很多人來說.
去幫人其實夾雜著很多動機.
其中一個動機就是.

$^{41}$希望能夠透過幫助人.
自己得利益.
自己得好處.
那些今世的利益.
去做這些事情.
若果你只是要求他.
不如你默默付出.
不問代價.
不出聲去做的話.
可能我們都未必會去做這些事情.
事實上都是真實的.
坦白說.
有些人說.
牧師如只是為了今世的好處去做好事.
不如不要做了.
幫助人 接待人.
尤其是少數有困難的人.
為什麼要幫呢.
又沒有即時的利益好處.
教會為什麼要幫一些弱勢的人.
有困難的人.
又沒有即時的回報.
為什麼要幫呢.
其實聖經裡面.
有沒有給我們一些提醒呢.
如果可以的話.
我們一起去看.
路加福音第十四章第一節.
去到第十四節.
我今天講到是會特別針對.
第十二節和第十四節.
不過看完第一至十四節.
是會讓我們知道整個故事的脈絡.
容許我們透過這十四節經文.
來幫助我們的思想.
究竟幫助人 接待人.
和利益之間有什麼關係呢.
換句話說.
是否真的會做一件事.
都要這樣計計較較呢.

$^{81}$如果對這段經文有深入的理解.
我們必須從經文著手留意.
尤其要留意到耶穌當時.
處身在什麼情況下.
如果留意到.
路加福音第十四章一至十一節.
我們會發現一件事.
其實耶穌在那段經文裡面.
講到耶穌被法利賽人邀請去.
一個法利賽人家的首領當中.
去吃飯.
在耶穌出席宴會的時候.
他被人監視著.
當時我們就問.
為什麼耶穌被監視著呢.
好像是想在耶穌身上.
拿一些把柄.
或者對耶穌平頭品足.
當耶穌去到現場的時候.
在現場前面.
就發現一個患有水股的人.
患病的人.
耶穌當時對著這個患病的人.
問回律法師和法利賽人.
安息日治病可疑不可疑.
耶穌就叫這班法利賽人.
問了他們.
這班法利賽人和律法師.
一言不發.
一句話都不說.
耶穌就醫好了這個可憐的人.
就叫他離開.
叫他離開之後.
耶穌就再問這班法利賽人和律法師.
他問他們.
你們中間誰有老或有牛.
在安息日要在井裡.
不立時拉他上來.
這些又是不答的.
不答的問題很有趣.

$^{121}$耶穌問了兩問.
兩問都不答.
其實想說什麼呢.
其實這個醫人.
如果當時的宗教領袖.
在他們眼中.
是不在意這些在他們眼前出現的貧窮人.
患病的人.
可憐的人.
在他們眼中這些人.
即是不是人.
可以說比較突兀一點.
在他們世界裡.
這些可憐的人不值得他們去關注.
他們只想關注.
他們自己以為對的人的宗教權民.
去死守住就行了.
好了.
在耶穌在世的時候.
發現這些現職的客人.
大家在一起坐首位.
坐頭位.
最尊貴的位置.
其實在當時.
講的是自我身感膨脹的表現.
即是抬高身價.
即是覺得我夠class.
我才可以坐在這些高位上.
耶穌這樣提醒他們.
不要自抬身價.
要謙卑下來.
最後耶穌就向.
邀請他們去現職的回去坐首位.
講第十二節和第十四節.
為什麼頭十一節要這麼快講一講.
因為他給了一個概括的輪廓.
就是說.
原來耶穌講這番話的時候.
耶穌當然是一眼看著.
在現職上的人.

$^{161}$這一班人基本上.
對其他有需要的人.
是漠不關心.
而且很在乎自己的身份和利益.
他們很想做很多事.
在抬高自己的身份.
在這樣的處境下.
耶穌就講第十二節和第十四節.
當我們想起第十二節的時候.
我們會發現耶穌這樣說.
祂說.
你擺設午飯或晚飯.
不要請你的朋友.
兄弟.
親屬.
和付出的靈者.
恐怕他們也請你.
你就得了報答.
繼續找字眼.
大家聽上去是有些有趣的.
其實祂想說什麼給我們聽呢.
為什麼耶穌會這樣說呢.
原來在當時的宴會文化當中.
其實每一個國家都有宴會文化.
但是在這個近東文化的宴會文化.
宴會文化當中.
原來當時如果有人請我們去吃飯的話.
請吃飯就等於接待我們.
接納我們.
就表明我們和現職主人之間的思想和信念.
都是相等於被邀請的.
就是說你被請.
你就成為現職主人的自己人.
就是自家人.
就是說他們被請去吃飯.
唯有取悅.
希望能夠得到彼此的肯定和認同.
所以吃飯當時是一個很重要的文化.
互相撐著大家的活動.
所以在現職當中.

$^{201}$如果在宴會要互相撐的時候.
當然請朋友.
又會請兄弟姐妹.
也會請親戚.
大家要很明白.
與此同時.
在這些現職的人當中.
他們更會請一些志趣相投的富足的鄰舍出席.
為什麼呢.
為什麼要這樣做呢.
原來其實請這些人的目的就是.
希望鞏固自己在社區當中的身份地位.
和維持其他在社群當中的聯絡.
就是說請這些有錢的鄰居一起吃飯.
身份被肯定.
同一時間就有很多網絡出現.
有一個特關係網出現.
其實這個和我們現在有的東西.
都不是太大的分別.
如果你身邊有些富貴的朋友.
想一起吃飯的話.
你身邊有其他朋友.
就去吃飯.
你是一個名人.
你的身份自然地地位會高了一點.
還有.
這些有錢人或有名的人.
自然會認識其他有名的人.
關係就會怎樣呢.
建立更加大的網絡.
在這樣的時候.
固然是這樣做.
但現職的主人也有時候請的時候.
身份低一點的人.
不是低很多.
是低一點的人來吃飯.
為什麼呢.
因為當時的文化就是.
如果一個地位高的人請一個低一點的人來吃飯的時候.
身份低一點的人來吃飯的時候.

$^{241}$就真的有一個韜光.
所以去到這些現職的時候.
他們就要去尊崇那個現職的主人.
你叫他去奉承也可以.
你說用一個俗套的說話.
就在這裡拍他的鞋也可以.
在這樣的文化下.
但現職主人是永遠不會請一些身份低.
極之相當於很大的人.
尤其是一些地位很卑微.
很寒美的人.
為什麼會這樣呢.
因為身份低的人.
是絕對不能夠按當時擺現職的習慣.
向現職主人在經濟上作出彙報.
例如我們中環有一次很清晰.
我不知道大家還記不記得.
我們小時候跟父母.
或者自己去飲酒的時候.
或者我們結婚的時候.
去擺酒的時候.
人家來的時候.
總會封封利是給你的.
或者給你禮物.
我記得我小時候看電視節目.
有一點諷刺時代.
那個節目就是有些人進去喝酒.
大家記得那個圖片.
中式的酒樓.
就是一張長桌子.
在一號的快餐店.
在桌子面前.
人家登記.
就看看你給多少禮金.
一開.
十元就有.
你坐那邊吃個橙.
然後那個人就拿著.
有些人拿著橙.
然後拿著炒粉.

$^{281}$拿著炒麵.
諷刺時代.
就是你做的人情做得不夠.
就給少一點東西.
在這件事上.
如果你去到這些宴會當中.
如果你身分比較低一點的話.
你要去的時候.
你都要在經濟上.
帶一點回報給這個宴席自願主人.
不然就不可以去這個宴席.
所以如果你是平行人.
沒有錢的時候.
怎樣去這個宴席呢?.
去不到的.
就是說.
擺宴席是會虧錢的.
是計算著的.
你來就有錢入袋.
一定收錢便宜.
分業還多一點.
還有很多時候.
如果我去一個宴席當中.
我請一個身分很低的人去的話.
是等於將自己的身分壓低的情況.
是自我否定的表現.
就是說如果我請一個很平行的人來吃飯.
如果我是一個.
現在的上等人地位上.
我自然把自己的地位拉低.
就和這些平行人一樣.
所以.
在你明白了昔日的宴會文化的時候.
你就更加清楚.
為什麼當時身處宴會的人.
是最多的平行人.
剛才我說了.
冷漠得不得了.
而且很著意坐高位.
因為這個宴席.

$^{321}$本身的目的是怎樣呢?.
是要使用自己的時間和經濟的資源.
來幫助自己去極速上位的.
所以宴席主人亦都.
說真的一次.
這個宴席的目的.
就是希望去拉幫.
結派拉關係.
得到最高的利益.
然後上位.
而不是來理會其他人的需要.
這個對比我們教的.
所謂的外宴.
例如一起吃飯.
大家分享.
是完全兩回事.
或者可以去研發一些.
大家可以看到圖案.
就是一班商家出來吃飯.
大家看看.
我是哪個老闆.
我是哪個會食主.
大家在拉幫結派.
吃飯做生意.
誰是大老闆.
大家思考一下.
是怎樣的.
希望和這個大老闆吃飯的時候.
拉到關係.
做到生意.
差不多這樣的理解.
去謀求商機.
所以會出盡法寶.
出盡各種方法.
要怎樣得到利益.
所以爬上高位.
拉大佬.
又或者見到朋友.
千萬不要理我.
不要搞我.

$^{361}$我不想拉.
拉低身份.
一切都想說耶穌的份上.
耶穌指出.
如果不這樣做的話.
你可能已經得到地上.
暫時的利益.
鞏固了你的勢力.
賺到你今世可能要有的財富.
但耶穌說的是什麼.
耶穌和這個賢者主人說.
你不要再請那些朋友.
兄弟親屬.
富足的鄰居.
這裡不是叫他以後都不要請.
因為那個不要請.
不要請那個字的文化.
是有不叫單單請的意思.
就是說耶穌在這裡.
是想鼓勵人.
不要因為要自己得到地上.
那個的獎賞.
所以就用盡自己的利益.
來招待人.
從而想自己其實是透過.
做這些事而自己取利.
耶穌在鼓勵賢者主人.
一個更好的方法.
善用他手頭上的資源.
耶穌邀請賢者主人.
邀請什麼人呢.
貧窮的.
殘廢的.
缺腿的.
黑眼的人.
剛才所說的四類人.
全部都是地位身份很低.
經濟能力是沒有的.
邀請這些身份寒微.
和在產質的人.

$^{401}$出來去到議會.
是基本上一個反文化的行動.
應該叫Counter Culture.
就是你不會這樣做的.
一定不會的.
因為你當初是不單止利益得不到.
還要被人看到你自己.
是降低自己的身份.
放低自己的身段.
在當時一個很啟示個人尊榮.
和啟示恥辱的文化當中.
這樣做法.
既沒有利益可言.
而且還要是.
可以這樣說是褻瀆的.
被人侮辱的.
其實和.
有一點點像華人文化.
華人文化對於這塊面.
是很重視的.
所以你這樣做.
基本上和當時文化是相反的.
但是.
耶穌說給賢直的主人.
你這個接待關心的行動.
不單止不會帶給你有任何羞恥.
而且這個行動其實會帶來真正的福氣給他.
因為耶穌說出.
這樣做出來的原因.
經文這樣說.
因為他們沒有什麼可報答你.
到二人復活的時候.
你要得著報答.
神因著人願意.
放下自己的榮辱.
尊嚴.
或者自己的利益好處.
溫柔他們的資材.
來幫助一些.
永遠都不可以在他們身上.

$^{441}$報答他們的人的時候.
神體在眼內.
就在二人復活的時候.
就會得到神親自的報答.
要知道.
作為賢直主人的法利賽人的領袖.
在他們信仰的圈裡面.
他們是深信神會回來審判世界的.
審判人.
死人要復活.
要得著神審判的.
耶穌是指出當日.
當他願意放下身段.
而用他的資材.
能力去幫助一些最軟弱.
最弱勢的人的時候.
上帝就會按著.
他所行的美事來獎賞他.
要記住一件事.
這不是透過接待這些軟弱的人.
來換取任何利益.
這不是這個意思.
可能你初聽的時候有個感覺.
因為上帝在這裡是不用回禮的.
就不是像我剛才說的賢直那些人.
我貧窮一點.
我來回禮給主人.
沒有的.
在上帝眼中.
上帝是不需要回禮的.
他只是按著人有沒有私行憐憫.
而得到神的賞賜和應分.
上帝是有絕對的權力去做事的.
同一時間作為花曆初日.
一個在生命當中追求.
走在神的律法.
和在人面前對神的律法.
有一個正確解釋的人.
他的思想是要做神要做的.
覺得是對的事.

$^{481}$他覺得要做一些.
他解釋的舊約聖經.
是按著他認為最對的方法去解釋.
這是他的思維.
這種賢直主人.
必定看重他們自己所做的事情.
是不是得神喜悅.
耶穌指出.
若果是要得到上帝的喜悅的時候.
他的生命必須要離開罪惡.
去約聖的人.
施行幫助和接待人.
或許新約有另一段經文.
是特別將這個訊息凸顯出來.
就是《雅各書》第一章第27節.
說得很清楚.
在神我們的父面前.
那稱見沒有尖烏的地的賢成.
就是看顧在患眼中的孤兒寡婦.
並且保守自己不沾染世俗.
所以耶穌在鼓勵賢直主人.
要善用資源和放下身段.
來幫助在弱勢當中的人.
說完之後對我們有什麼意義呢.
其實耶穌對賢直主人所說的.
這番說話和教導對我們來說.
是很重要.
因為對我們來說.
作為教會.
很有趣的.
教會其實有著金錢,人力,物資.
甚至建築物的資源.
我們沒有想過.
我們有沒有想過.
每一個星期都會有人奉獻金錢.
當你在現在的視像上.
和我一起去敬拜的時候.
有一班人在這裡.
做完這些視頻.
你才可以看到我來去敬拜.

$^{521}$或者做很多工作.
就有人來參加.
物資.
其實教會有很多物資.
你試一下.
有哪個人生病了.
可以幫忙煮飯呢.
我相信教會總有幾個太太.
可能馬上說.
等我來.
我煮.
又或者說.
我車子有事.
有沒有人幫忙.
看看.
可能有一兩個人說.
我知道誰可以幫你.
你想想這些常生活中的小事情.
其實對很多人來說.
那麼多人可以幫忙.
是的.
因為我有群體力量.
還有.
你說我要租建築物.
是你租建築物.
但是你有一個建築物.
可以相聚.
可以聚集.
可以用來做很多工作.
我自己教會有一個建築物.
現在其實我是教會圖書館.
就是在當中.
我教會有建築物.
可以做很多事.
一樣.
那些都是什麼.
物資都是資源.
作為建築圖.
就算是貨物也好.
金錢上.

$^{561}$時間上.
人力上.
資訊上.
都有各種資源.
有的.
我們每一個人都有資源.
剛才我提到的資源很少.
就算簡單到.
我病了.
你幫忙煮粥給我吃.
可能在底層.
煮東西給你吃.
這個都是技能.
又或者說.
我不懂報稅.
你教我.
教會當中.
我懂報稅.
我教你.
技巧.
金錢上.
時間上.
可能有些時間收入.
都可以有些時間.
轉到出來幫忙.
所以你這樣看的時候.
其實每個人都有一些資源.
我們當然可以.
按不同的情況之下.
運用這些資源.
每一個人都可以用.
各用在自己身上.
用在教會身上.
沒問題.
不是錯事.
是正常事.
但是我們要十分留意.
我們怎樣運用這些資源.
這裡不是指我們.
對自己的資源不作出安排.

$^{601}$我們每個人都會安排.
我們的時間.
金錢.
教會都會計算.
都會有財政報告.
預算.
教會的事工.
我們都會做安排.
安排不是問題.
需要安排.
問題就是.
我們帶出什麼心態和動機.
去安排.
去運用這些資源.
我們是不是只用資源.
來幫助得到自己的要求.
又或者我們做一些工作.
是自己得到要求.
或者是透過做一些工作.
在當中循序取利.
利益不一定是錢.
可以各方面.
還是我們願意放下自己的利益.
和身段.
來幫助一些有需要.
弱勢的社群.
知道他們實際上除了說一句.
Thank you 多謝之外.
可能他們什麼回報都給不了我們.
尤其是當我們服務一些弱勢的群體.
服務一些有困難的社區的人.
可能真的最多.
他可以說一句Thank you.
那我們幫不幫他們.
耶穌在這裡提醒我們.
當我們和人建立關係.
接待幫助人的時候.
我們不要整天想著.
在腦裡有算盤和計算機.
整天計算著.

$^{641}$我付出多少.
我就可以得到多少.
但我們也有時會犯同樣的毛病.
我們會想.
我要做多少.
我就要得到多少.
我們要反過來留心.
我們是不是不求個人或教會的榮辱.
只是甘心付出來幫助有需要的人.
知道在上帝的眼裡.
在將來十幾日子回來的時候.
祂看到我們對祂的敬潔.
不是一把口說的.
而是在行動上表現的.
將來要得到祂的尊敬.
其實各位.
當我們接待別人.
幫助別人的時候.
不論是別人是從世界不同地方來的新移民.
甚至是難民.
又或者是社區當中.
特別病病的軟弱的朋友.
這些朋友我們幫了他.
他也沒辦法幫我們.
真的只能再說一句.
多謝.
Thank you God.
You help me.
幫我們說這幾句.
多謝.
但我們願不願意放下.
要得著自己利益的心態去幫助他們.
各位.
我相信耶穌希望我們的教會和我們的生活.
就是見識出.
神就是愛而出於我們的生命.
神就是愛的生命不是一把口說的.
神就是愛的生命.
本身就是一個需要付出.
但是默默的付出.

$^{681}$不是為了要求讚賞.
在社區當中得到一個名分和名號.
而是真心的.
不是為了那些利益.
而是為了我.
為了上帝的恩典.
而幫助他人.
耶誓幫助我們.
讓耶誓教導我們.
教導我們.
教訓我們的心.
帶著一顆願意的心去服侍.
令社區當中.
貧寒困難有需要的人.
耶誓賜福給大家.
\newpage



\section{}
\label{sec:CgJV5or0kgk}
\textbf{Family Sunday School - EP 6 (The King of Jews)}
\newline
\newline
連結: \href{https://youtube.com/watch?v=CgJV5or0kgk}{\texttt{https://youtube.com/watch?v=CgJV5or0kgk}} ~~~~ 語音日期: 2022-03-30
\newline
\newline
\hyperref[sec:pinIq8rUt5E]{\small{< < < PREV SERMON < < <}}
~
\hyperref[sec:index_chronic]{\small{[返順時目]}}
~
\hyperref[sec:index_scriptual]{\small{[返順卷目]}}
~
\hyperref[sec:deMiI9YWwUk]{\small{> > > NEXT SERMON > > >}}
\newline
\newline
$^{1}$(upbeat music).
(speaking in foreign language).
- We're back with a family Sunday school time..
(speaking in foreign language).
Hot potato..
(speaking in foreign language).
Pontius Pilate..
(speaking in foreign language).
(sighing).
(speaking in foreign language).
Especially..
(speaking in foreign language).
God's salvation plan?.
(speaking in foreign language).
The rest of the story..
(speaking in foreign language).
Drawing..
Let's go..
(speaking in foreign language).
What's the theme that connects them?.
Okay..
(speaking in foreign language).
(writing on chalkboard).
(speaking in foreign language).
Sahara Desert..
(speaking in foreign language).
Cactus..
Okay, we got it..
(writing on chalkboard).
Cactus..
Okay..
Next one..
(speaking in foreign language).
Ready?.
Here we go..
(writing on chalkboard).
(speaking in foreign language).
(writing on chalkboard).
(speaking in foreign language).
Porcupine fish..

$^{41}$(speaking in foreign language).
(writing on chalkboard).
Puffer..
(speaking in foreign language).
Puffer fish..
Okay..
Next one..
(speaking in foreign language).
Let's see..
Hmm..
(speaking in foreign language).
Pencil..
(writing on chalkboard).
Pencil..
Okay, here I go..
Next one..
(writing on chalkboard).
(speaking in foreign language).
(writing on chalkboard).
(speaking in foreign language).
You got it..
Rose..
(writing on chalkboard).
Okay..
(speaking in foreign language).
Okay, let's see..
(speaking in foreign language).
Okay..
Hmm..
(writing on chalkboard).
(speaking in foreign language).
Sword..
Okay..
(speaking in foreign language).
What is similar about these five drawings?.
(speaking in foreign language).
Cactus..
Yee..
Pufferfish..
(speaking in foreign language).

$^{81}$Sword..
Hmm..
(writing on chalkboard).
(speaking in foreign language).
Spiky and sharp..
Sharp and spiky..
Spiky and sharp..
Sharp and spiky..
(speaking in foreign language).
All right..
(speaking in foreign language).
Bible story time..
(speaking in foreign language).
Pontius Pilate..
(speaking in foreign language).
Pontius Pilate..
(speaking in foreign language).
Jesus..
(speaking in foreign language).
King of Jews..
(speaking in foreign language).
Or leader, government, right?.
(speaking in foreign language).
Bright, red and blue..
(speaking in foreign language).
I am the king of Jews..
(speaking in foreign language).
This kind of king..
Sure..
(speaking in foreign language).
Okay, continue the story..
(speaking in foreign language).
Watch..
(speaking in foreign language).
A whip..
(speaking in foreign language).
Ouch..
(sighing).
(speaking in foreign language).
Prepare..

$^{121}$(speaking in foreign language).
Prepare, prepare..
(speaking in foreign language).
Make fun..
Mock..
(speaking in foreign language).
No way..
(speaking in foreign language).
Like a nail..
Nail..
(speaking in foreign language).
But he is, anyway..
(speaking in foreign language).
Enjoy your debrief..
(speaking in foreign language).
Pictionary..
(speaking in foreign language).
(writing on chalkboard).
(speaking in foreign language).
A crown..
(speaking in foreign language).
Thorns..
(speaking in foreign language).
Thorn..
(speaking in foreign language).
No way..
(speaking in foreign language).
Why?.
(speaking in foreign language).
Tough story..
(speaking in foreign language).
A sign..
(sighing).
Like this..
King of Jews..
(speaking in foreign language).
Make fun of him, mock him..
(speaking in foreign language).
How does a king enter a town?.
(speaking in foreign language).

$^{161}$Symbol..
(speaking in foreign language).
Announcement?.
Over and out, the king of the Jews is coming..
10-4, 10-4..
All right, all security, let's go, let's go..
(feet shuffling).
(speaking in foreign language).
Carried?.
(speaking in foreign language).
Donkey..
(speaking in foreign language).
A donkey..
(speaking in foreign language).
Very special, very different..
(speaking in foreign language).
Bye..
Okay, everyone, Chinese time..
(speaking in foreign language).
Royalty, king, queen, prince, palace..
Prince, palace, crown..
(speaking in foreign language).
King, okay..
Queen, okay..
(speaking in foreign language).
So connected, okay?.
(speaking in foreign language).
Ready?.
Okay, king..
(speaking in foreign language).
Queen..
(speaking in foreign language).
Still has the wong, right?.
You know how I remember the difference?.
(speaking in foreign language).
Prince..
(speaking in foreign language).
Zi is what?.
Child..
King's child..

$^{201}$Prince..
(speaking in foreign language).
Now, where does a king sit?.
On a throne..
(speaking in foreign language).
What seat does the king sit in?.
The throne..
(speaking in foreign language).
Right?.
Where does the king live?.
(speaking in foreign language).
This is how I remember..
(speaking in foreign language).
As big as a park..
(speaking in foreign language).
And then last..
(speaking in foreign language).
A crown..
But you don't say like (speaking in foreign language).
like a king's hat, no..
(speaking in foreign language).
Like (speaking in foreign language).
That's how I remember..
It's not the same word, but that's how I remember..
Like champion..
(speaking in foreign language).
Everything has a wong,.
which connects it all to royalty..
(speaking in foreign language).
Okay, another reason to learn it..
Bye..
Chinese time..
(speaking in foreign language).
Home extension time..
(speaking in foreign language).
Make a crown..
(speaking in foreign language).
Construction paper..
(speaking in foreign language).
Decorate..

$^{241}$(speaking in foreign language).
Parade around, okay?.
(speaking in foreign language).
Who?.
Yeah, so..
Why?.
'Cause he loves us..
(speaking in foreign language).
Wire..
(speaking in foreign language).
But literally, (speaking in foreign language).
Okay?.
(speaking in foreign language).
(upbeat music).
\newpage



\section{}
\label{sec:deMiI9YWwUk}
\textbf{Family Sunday School - EP 7 (Good Friday)}
\newline
\newline
連結: \href{https://youtube.com/watch?v=deMiI9YWwUk}{\texttt{https://youtube.com/watch?v=deMiI9YWwUk}} ~~~~ 語音日期: 2022-04-04
\newline
\newline
\hyperref[sec:CgJV5or0kgk]{\small{< < < PREV SERMON < < <}}
~
\hyperref[sec:index_chronic]{\small{[返順時目]}}
~
\hyperref[sec:index_scriptual]{\small{[返順卷目]}}
~
\hyperref[sec:i4txNEY8_xE]{\small{> > > NEXT SERMON > > >}}
\newline
\newline
$^{1}$[Music].
Hello, everyone!.
Welcome back to Family Sunday School..
I'm Ms. Rebecca..
Today is the 7th episode..
Have you noticed anything different about Ms. Rebecca today?.
Her hair is different, right?.
She's wearing a pony tail..
Do you know what a pony tail is?.
"Ma" is horse. "Mei" is tail..
So what is a pony tail?.
Pony tail!.
Ma is pony tail..
Ma Mei!.
Pony tail..
And what color is Ms. Rebecca wearing?.
You can't see it now, but her top is black..
Her bottom is also black..
All black..
Why?.
Why is Ms. Rebecca wearing all black?.
We'll find out in today's story..
But first things first..
Let's talk about the last episode..
Let's repeat..
We talked a lot about crowns in the last episode, right?.
Did you make a crown?.
Is it full of gold and jewelry?.
It's shining..
But our Savior Jesus decided to wear what crown?.
Yes, it's very painful and tight..
It's like a new crown..
It's all thorns..
Thorns poking at his forehead..
It's bleeding..
He's even holding his own cross..
What does the sign say?.
King of Jews..
He's being mocked..
He's being nailed to the cross..

$^{41}$That's what we're going to talk about in today's episode..
What happens after being nailed to the cross?.
Can you guess why Ms. Rebecca is wearing all black?.
You can't or you don't know?.
It's okay..
Ms. Rebecca will smile after the Bible story..
Or you'll know in the middle of the story..
But before we talk about Bible study,.
we're going to play a game..
What's the game?.
Ms. Rebecca will pretend to be a game..
You have to guess..
Get ready..
Game time..
Activity time..
It's game time..
Ms. Rebecca will focus on feelings today..
You have to guess what Ms. Rebecca is feeling..
I'm going to cover my face..
You can guess..
We're used to wearing masks..
You can tell by the mask..
If you cover your face,.
you can tell by the eyes..
Let's get started..
Ready?.
What are you feeling?.
Yes..
Angry..
What are you feeling?.
Yes..
Ready?.
Number two..
That was a practice..
Ready?.
Number two..
What is it?.
Confused..
Next one..
Just with my eyes?.

$^{81}$Ready?.
Next one..
Confused..
Ready?.
Next one..
What is it?.
Sneaky..
Feeling sneaky..
Next one..
Can you guess?.
Yes..
Sad..
What is it?.
Disgusting..
That's gross..
Did you get it?.
Happy..
Mr. Becker is happy..
I hope you can learn facial communication..
You can express yourself by talking..
Your face, body language, and attitude can tell a lot..
Today we will focus on feelings..
Let's see what Mr. Becker is wearing..
What do you want to express?.
Happy..
Sad..
Angry..
Surprised..
I want to wear black..
Story time..
Get ready..
Story time..
Hello everyone..
It's Bible story time..
Where did we leave off last time?.
We left off in the middle of two statues..
I want you to notice where the sun is..
Yes, it's here..
Where did we start?.
Last supper..

$^{121}$We prayed in the garden..
Jesus was caught..
Peter denied him three times..
The rooster crowed..
It was sunrise..
After doing so many things, he was nailed to the cross..
Mourning..
Do you know when Jesus died?.
When did he die?.
Between 12 and 3pm..
The sky was dark..
12pm is usually lunch time..
Lunch time is high noon..
Why?.
Because the sun was so high..
It was very hot..
Why was the sky like this?.
It was dark..
What does it mean?.
Jesus died in the middle of nowhere..
Huge deal..
Jesus will die..
The time of waiting was also in the middle of nowhere..
The whole world knew..
Of course, this is God's plan..
Of course they knew..
But there were many other things..
Who was there when Jesus was nailed to the cross?.
John, John, Mary, and other friends..
There were also soldiers..
Those who helped the government..
Those who watched him and made sure he died..
There were two criminals..
They knew Jesus was innocent..
They knew Jesus was God's son..
He was innocent..
Remember me in heaven..
Jesus said to John,.
"Because you have accepted me,.
I will give you a place in heaven.".

$^{161}$This means,.
"Hey, this criminal knows..
Can't other people see it?.
Don't other people know?".
I don't know..
And one other thing happened too..
Jesus asked God,.
"Have you forsaken me?".
"Do you really want me to die like this?.
Do you really want me to suffer like this?".
"Yes, I do.".
But do you remember how Jesus prayed in heaven?.
"If you don't need this plan,.
if you don't need me to die,.
if you don't need me to suffer,.
please take this cup away.".
But this plan must be fulfilled..
Jesus has to die..
Before he died, he had a throat..
Under the cross,.
there was a soldier who brought vinegar..
Who drank the vinegar?.
It was very sour..
He gave it to him..
After he drank it, he said,.
"God, forgive them..
They don't know what they are doing.".
He didn't know what he was doing..
Okay..
Jesus suffered so much..
Before he died, he called God,.
and put a cross on himself..
Wow, what a big heart, right?.
Super big..
And then later,.
Jesus breathed his last breath..
And he died..
And then more happened..
You know what happened?.
The whole earth was shaking..

$^{201}$A huge earthquake..
Do you know what happened next?.
A Jewish man had a huge temple..
He had a curtain there..
He tore it apart..
Wow, not only the earthquake,.
the curtain was also torn..
And those who died for God.
started to come back to life..
They started to come out of their graves..
Wow..
What do you think people who see this will do?.
What is good?.
Surprise, right?.
That's a feeling..
Wow..
They were surprised..
They would say, "This is really the Son of God.".
"He is not an ordinary person.".
"This Jesus is definitely very special.".
"He must be the Son of God.".
"Look what's happening.".
But there are other feelings..
Are there other feelings?.
Marley, John, so many friends are down there..
Jesus died..
The cross on the cross is dead..
The sky is dark..
They feel so sad..
You see these two people..
These two men next to him..
They are wearing black..
One is called Joseph..
A very rich Jew..
He is also a council member..
He is also a leader..
The other is Nicodemus..
He is a leader..
In the previous story,.
Nicodemus sneaked out to find Jesus..

$^{241}$Because he believes Jesus is the Son of God..
He believes Jesus is the Son of God..
Nicodemus even called Jesus a teacher..
They asked,.
"Can you give us Jesus' body?".
"We will bury him.".
They buried him in the Jewish way..
But the Jewish way requires a lot of oil and money..
Without this rich man, Joseph,.
they would not be able to do it..
Many of the Marley are there..
Lots of Marys..
Mary of Magdalene..
Mary, Jesus' mother..
And other women who love Jesus..
They helped to find Jesus..
They found a tomb back here..
Back here..
In the past, rich people dug a hole..
They said, "After I die, I will be buried there.".
They found a person who had not been buried yet..
They said, "We will put Jesus in there.".
Look at their faces..
Do they want to cry?.
They put Jesus in there..
They put a stone in front of him..
They said, "Don't use the paper to steal his body.".
"We will show you that he will not come out.".
"He will die.".
This stone is very heavy..
It covers the entrance and the exit..
I can go in..
Because this tomb is like a cave..
There is no door behind..
There is only one door..
This door is covered by a stone..
Think about it..
How did they feel?.
How did the 12 disciples feel?.
John, Peter, and all the others..

$^{281}$They were with Jesus..
But now they are waiting..
But do you think the disciples remember.
Jesus once told them,.
"I will die..
It doesn't matter..
I will be resurrected in a few days.".
This is God's plan..
God has said this many times..
But I don't know if the disciples and Jesus' friends.
remember that moment..
Because a lot of things happened..
How did they feel?.
Confused..
Everything was very confusing..
There was a lot of mixture of emotions..
Right?.
This story is about Good Friday..
This is the day Jesus died..
It was a dark and gloomy Friday..
It was sad..
It was gloomy..
I don't know what happened..
I don't know what I could do..
But I don't know what I could do..
In three days, what will happen?.
He rises..
That's what makes it good..
It's good for you..
And it's good for me..
It's good for everybody..
Jesus' death is good for everyone..
Because his death is good for everyone..
Do you know why Mr. Buckethead.
is wearing a full body suit?.
Why is he wearing a full body suit?.
Because we usually celebrate.
Good Friday..
We celebrate Jesus' death..
In many cultures,.

$^{321}$people wear dark or dark suits.
to funerals..
People wear all black or all dark..
Because dark or dark suits.
often represent what?.
Which feeling?.
This feeling?.
Or.
this feeling?.
Yes..
Sadness..
We cry and be sad.
for the lost person,.
the dead person,.
the lost Jesus,.
the dead Jesus..
There is one point.
I want to mention to everyone..
Jesus' death.
is a huge deal..
Not just because.
other people are watching..
Mary,.
Jesus' earthly mother,.
watching her son.
being crucified..
Jesus called John,.
John, remember to take care of my mother..
Those feelings..
What's the good of having a son?.
Maybe confused?.
What happened?.
Jesus is God's son..
Why did this happen?.
Why did he die?.
Mary was very sad..
So sad..
Or some people were happy..
Yes!.
Getting rid of Jesus, crucify him..

$^{361}$Yes!.
Some people.
were angry..
What?.
Crucify Jesus?.
Or the two sinners.
next to him.
would think,.
I know I made a mistake..
Jesus is not wrong..
I know I made a mistake..
Jesus is not wrong..
What?.
There are many different feelings..
Many different feelings..
Right?.
But Heaven.
also has feelings..
Remember?.
Boom! Boom! Boom!.
Heaven is dark..
There is one more thing..
In the Great Hall,.
the Great Holy Temple,.
there is a bowl,.
a curtain,.
like this..
A thick curtain..
It is removed..
Do you know what this curtain.
was used for?.
In the past,.
the Great Holy Temple.
had a curtain.
to block out.
who could approach God.
and who could not..
Can you go behind the curtain?.
It means you can go..
It means you are.

$^{401}$being excluded from the Great Hall..
Not everyone.
can go..
So,.
in the story,.
it is said that the curtain was removed..
It means.
it no longer blocks you.
from going in..
Let's try. How hard is it to remove the curtain?.
Here is a piece of paper..
Remove..
Super easy, right?.
You can keep removing it..
You can remove.
both layers..
Here is a piece of paper..
Remove..
You can't remove it..
I can't remove it with this much strength..
Try to remove it..
You can't cut it..
The story does not say.
you have to cut the paper..
It is the curtain itself that is torn..
What does it mean?.
The curtain separates.
God from others..
Not the select few..
Not the chosen few..
Yes..
That's right..
What does it mean?.
Jesus' death.
is for everyone..
Now, everyone.
can accept Jesus and approach Him..
No longer a barrier..
Barrier no more..
No more blocking..

$^{441}$Wow..
Huge deal..
But also sad..
Right?.
Sigh..
Especially.
Jesus' good friend..
When a good friend.
is gone,.
when a good friend dies,.
of course you will be sad..
So Good Friday.
is the cost.
of salvation..
Wow..
The cost of.
washing away my sins..
Tear this cloth away..
That cloth..
It is worth it..
But the question is.
does Jesus die worth it?.
This is God's plan..
Of course.
He thinks it is worth it..
That,.
my friends,.
is a lot of love..
It is worth it..
It is worth the sacrifice..
But you know what?.
Jesus has said it will happen..
He also said.
He will be resurrected..
But sometimes.
we may not remember.
or be clear..
Many disciples.
may be very confused..
What is going on?.

$^{481}$But you and I.
because it is this side of the story.
we can know.
what will happen..
Next time.
continue the story..
Chinese time!.
OK..
When I was in the game.
I said.
many different sounds.
or feelings..
Phew..
In English.
we may not know many feelings..
Happy.
and sad..
But happy and sad.
there are many.
feelings in the middle..
I will teach you a few..
OK..
So next time.
how are you doing?.
I am fine. I am happy. I am sad..
But there are many.
feelings in the middle..
Next time.
when someone asks you in Chinese.
how are you?.
You can answer many different things..
Happy.
and sad..
You know what?.
Often.
my kids will say.
not happy or sad.
but a little annoyed..
Bothered..
Irritated..

$^{521}$That's one..
Annoyed..
Another one.
which sounds like.
a.
food..
Like vinegar..
Have you ever had vinegar?.
Sour..
A little sour..
Black vinegar is sour. White vinegar is sour..
Anyway, it's sour..
If your feeling.
is sour.
you feel sour..
You would make that face..
Very sour..
Why?.
Irritable..
Very irritable..
Usually.
if I feel.
very annoyed.
and then maybe.
I am irritable.
people would say.
Mr. Becker, you are very annoying..
Annoyed..
You know what that is put together?.
Annoyed..
The stresses.
of life. Stress..
Wow..
Chinese, so efficient..
Just a syllable..
Irritated..
Bothered..
Irritated..
Bothered..
Irritable..

$^{561}$Sour..
Stressed..
Stress..
So much meaning..
So rich..
Learn your Chinese..
Bye..
Hello, everyone..
At home extension time..
We talked about a lot of different feelings in this story..
Right?.
Why am I wearing black?.
For basically Good Friday..
Basically, it's Jesus' funeral..
Right?.
Or maybe you went to a funeral.
and wore black or.
white..
Depends on the culture..
Depends on the funeral..
How can you extend this learning at home?.
Why does color matter?.
Color can express.
the feeling..
Right?.
We talked about a lot of different feelings today..
Let's think about.
different paintings..
Art..
Painting..
What does red.
mean?.
Maybe danger..
You can ask your child,.
"What is red?".
Look around..
"What is red?".
Stop sign..
Red light. Right?.
Caution..

$^{601}$Feeling..
See the color?.
Next color..
Maybe something is yellow..
That's pretty.
uplifting. Happy..
Why?.
Let's see what is.
naturally yellow..
Sunshine..
Yellow sunshine..
Right?.
We can talk about these things..
Blue is a.
cooler color..
A cooler color..
Not a warm color..
What feeling?.
Maybe peaceful..
Maybe calm..
Or a calm feeling..
You can ask your child..
Ask your child.
and connect.
why Mr. Rebecca wore black today..
Why did he wear.
this color at the funeral?.
What is the connection.
between color and feeling?.
Rainbow color..
There are many colors..
But look again..
What does grey.
mean?.
What does grey mean?.
It's more somber..
Right?.
What does black.
mean?.
What does black mean?.

$^{641}$It's more mysterious..
Or.
it's more scary..
Or sad..
So,.
when you can use these colors.
to help your child.
to express.
their feelings,.
the attachment.
to the Bible story.
will be deeper..
But the ability to express.
will be stronger..
We have another.
at-home extension..
Draw what you hear..
We talked about different.
feelings today..
Different facial expressions..
Different looks..
Different colors..
How to connect colors and feelings..
Like I wore black..
It means sadness.
or sadness..
Using colors and feelings.
to tell a story..
This Bible story.
is for children to draw..
Draw what you hear..
There are many ways to do it..
Children can use any colored pen.
to draw..
But when they are drawing,.
they are drawing to listen..
Either to the story told by parents.
or to the story read in the Bible..
Jesus dies and was buried..
Or they can rewind.

$^{681}$to the story I told..
They can listen,.
not watch..
Listen, listen, draw..
I will show you..
I don't have to do the whole story..
But you will get the picture..
You can give them.
many different words..
There are many different scenes..
You will have a story.
maybe 3, 4, 5 versions..
You can write.
the writing.
and text together..
You will have your own story..
Just one way..
I will use this crayon effect..
I will say.
Oh, you need to use yellow..
Because.
when Jesus was nailed to the cross,.
what is it?.
It is mourning..
The cross of Jesus..
There is also.
a time..
But when Jesus was nailed to the cross,.
the time passed..
It was 12 o'clock..
High noon..
The time passed..
It became.
3 o'clock..
But what happened.
at 12 o'clock to 3 o'clock?.
Wow, it was a big deal..
The whole sky.
was black..
It was black..

$^{721}$Black..
Why?.
Because Jesus was nailed to the cross.
to die..
It was a big deal..
That's why..
Not only was it black..
It was bright and black..
But.
Wow, it was.
an earthquake..
During this time,.
children can draw.
and make a paper earthquake..
They can also.
shake the paper..
Or crumble the paper..
If there is a next scene,.
they can use the second paper..
Right?.
Let them draw.
until.
they can't draw anymore..
The story is over..
You have your story..
You can continue..
You can also repeat.
the previous episodes..
Or the next episodes..
You can have.
your own story..
You can write and draw..
Your own story..
Self-written, self-illustrated..
Okay? Give it a try..
[music].
\newpage



\section{24:36-44}
\label{sec:i4txNEY8_xE}
\textbf{The Last Days According to Jesus 2 (Matthew 24\_36-44) - Rev. Mark Chan}
\newline
\newline
連結: \href{https://youtube.com/watch?v=i4txNEY8_xE}{\texttt{https://youtube.com/watch?v=i4txNEY8\_xE}} ~~~~ 語音日期: 2022-04-04
\newline
\newline
\hyperref[sec:deMiI9YWwUk]{\small{< < < PREV SERMON < < <}}
~
\hyperref[sec:index_chronic]{\small{[返順時目]}}
~
\hyperref[sec:index_scriptual]{\small{[返順卷目]}}
~
\hyperref[sec:sGka_a0cCzQ]{\small{> > > NEXT SERMON > > >}}
\newline
\newline
$^{1}$So that's the announcements..
Again, we come before the word of God..
We gonna turn to the book of Matthew,.
the gospel of Matthew..
So if you have your phone app or you have your Bible,.
turn to Matthew 24, Matthew 24..
Last days, according to Jesus, okay?.
Last time I was talking about this,.
the Olivet Discourse..
It's a long sermon that Jesus gave..
What signs will happen when the Temple of Jerusalem.
was destroyed?.
So Matthew 24, the Olivet Discourse,.
is related to the broad subject of eschatology,.
the end time, right?.
The end time..
It's not complicated, but this word comes from.
the Greek term eschaton, right, as you remember..
It's about the end..
So here Jesus is answering questions.
posted by his disciples..
What would happen when the Temple of Jerusalem was destroyed?.
What are the signs?.
And then what are the signs of his second coming?.
So two questions are lumped together,.
and Jesus is giving answer to this twofold questions..
He is talking about a list of signs..
So conveniently, I list it here..
You will see false prophets..
You will see wars and rumors of wars,.
earthquakes, persecutions,.
and many people will turn away from faith..
Their love will grow cold,.
and the gospel will reach to the end of the earth,.
and there will be unusual astronomical phenomena..
The sun and the moon and stars,.
they will behave differently or unusually..
As I told you before, if we understand the background,.
the context of first century,.
especially 8030 to 8070, what happened there,.

$^{41}$all the signs happened before 8070,.
or around that time..
So the first generation of Christians,.
they have every good reason to expect.
the second coming of Christ in their own lifetime..
Jesus Christ is coming back any time..
So no more signs is needed..
Sometimes you may have heard some prophecy teachers.
telling you things like, okay, we are watching for this sign..
You know, we are getting closer and closer..
But look at what Jesus said, 2430..
At that time, the sign of the Son of Man.
will appear in the sky,.
and all the nations of the earth will mourn..
They will see the Son of Man coming on the clouds of the sky.
with power and great glory..
Now Jesus gave a list of signs..
He said, this is a list of signs that you guys,.
the first disciples, will see before 8070,.
or around that time..
So all the signs are already visible..
You don't need to wait for another sign..
As I said, some of the prophecy teachers today.
will be telling you, there was a sign we are watching,.
there was a sign we are waiting for, the end is coming..
Jesus said, no..
Some of the prophecy teachers will tell you,.
especially these days, you know,.
there were a lot of unrest around the world..
Some will tell you, okay, Russia invading Ukraine.
is another sign fulfilled before the second coming..
No, this is bogus, this is bogus..
No need to wait for another sign..
Jesus said, the next sign you will see.
is the sign of the Son of Man's second coming..
By the time you see the next sign, it will be too late..
It will be too late..
So that's why he is trying to convince,.
trying to teach us how to be prepared for his second coming..
Preparing for the second coming of Christ,.

$^{81}$not by looking for another sign..
All the signs we have seen, right?.
Especially even, you look around the world,.
there are many things similar.
to the list of signs happening right before our eyes..
Yes, of course, because we have seen it.
in the first century already..
So how do we prepare for the second coming?.
Not by looking for another sign,.
but by preparing in our own lives..
Okay, let's continue with this chapter 24, 36..
Jesus said, no one knows about that hour, that day,.
that day of his second coming,.
not even the angels in heaven, nor the Son himself,.
but only the Father, only God the Father..
He determines, he plans, he knows the time of the eschaton..
Only he knows..
So you would not know, even the prophecy teachers these days..
You hear here now from time to time,.
some prophecy teachers will tell you,.
okay, they can calculate the exact day.
and time of second coming of Christ..
Don't believe that, they are false prophets..
No one knows about the day and the hour,.
not even the angels, nor the Son, only the Father..
37, as it was in the days of Noah..
So it will be the coming of the Son of Man..
For in the day before the flood,.
people were eating and drinking,.
marrying and giving in marriage,.
up to the day Noah entered the ark..
They knew nothing about what would happen.
until the flood came and took them all away..
This is how it will be at the coming of the Son of Man..
What it will be like when Jesus comes back is unexpected..
The day of the second coming of Christ is sudden..
You would not know, I would not know..
And Jesus likens his second coming to the days of Noah..
Remember this story?.
Okay, our Bible study will go through that story later..

$^{121}$Genesis 6 to 8..
Chapter 6 to 8 of the book of Genesis..
Talking about the great massive flood.
coming in the days of Noah..
And when the flood comes,.
when the water swifts away some people,.
nobody knew, nobody expected it..
And Jesus has described what the people were doing.
in the days of Noah..
People were going through their daily routine as normal..
Right?.
We will be working, business as usual..
But when the flood comes unexpectedly,.
it swipes them away..
Two men will be in the field..
One will be taken away, the other left behind..
Two women will be grinding with a hand mill..
One will be taken and the other left..
Therefore, keep watch because you do not know.
on what day your Lord will come..
Now some will use this passage.
to support a theory called rapture..
Right?.
You probably have heard of that, right?.
If you have never heard of that, that's good..
All right?.
Because it's confusing..
Some of you may have heard that some people.
use this passage to support so-called rapture..
Now simply put, what rapture theory talks about.
is that when Christ comes again,.
he will come first in secret..
So nobody knows all the earth,.
all the people would not know..
He will come in secrecy..
And when he comes in secret,.
he will snatch away some people.
and some people will suddenly disappear off the earth.
and leave it behind those who do not believe in Jesus..
And that seven years of disaster,.

$^{161}$intense disasters will befall the earth..
It will kick in the seven years,.
the last seven years of human history..
Now if you read this passage carefully,.
I don't think the passage is teaching that..
I don't think the passage is teaching.
that Jesus is coming in secrecy.
and that will snatch away some people.
and the left behind are those who are doomed..
You read this passage carefully,.
two men are working in the field,.
one will be taken and the other left..
Now in the ancient society,.
men are working outside on farmlands, right?.
Sometimes father and son working together..
One will be taken away, perishing..
Because Jesus is talking about the flood..
The flood will take some people away..
And then those who are not being taken away,.
those are the people who are left behind..
So in that sense, it's good to be left behind, right?.
'Cause you don't want to be taken away,.
taken away by the flood..
Two women will be grinding with a hand mill..
Now sometimes, most of the time in the ancient society,.
the women will work inside the house..
So the hand mill is two big stone,.
you know, on top of one another,.
and then you kind of turn it, right?.
So two women will be working together,.
you know, opposite, sitting opposite against one another..
Probably a mother and a daughter..
One will turn 180 degree,.
and the other person will turn the other 180 degree..
And then come back on and on and on..
Hand mill, moving the hand mill,.
turning the stone again and again,.
it's daily routine..
One will be taken away,.
and the other will be left behind..

$^{201}$So if you read this passage carefully,.
it is good to be left behind, right?.
Because the flood will take away.
those who are unprepared for judgment..
Only Noah and his family in the ark.
were left behind, and they were safe..
All the rest of the people were being taken away..
That would be like the future,.
day of the second coming of Christ..
Nobody knows, business as usual, right?.
Look at this chart..
I wanna ask you the question,.
where are we right now?.
Now this is the New Testament.
or the biblical concept of time..
This is the biblical concept of history..
It is, as I said before,.
we are living in between the first coming.
and the second coming of Christ..
We are living right in between,.
between the first and second coming..
We are in this last period of human history..
We can call it already, but not yet..
Already the kingdom of God,.
the future kingdom of God.
has already in breaking into our history..
And we are getting closer and closer to the end,.
getting closer to the eschaton..
Now one important takeaway from this chart.
is that the biblical concept of time is always linear..
If you flip the Bible, if you study the Bible,.
you understand the Bible has its own concept,.
perspective on time, our time, human history..
And it is linear..
It is not cyclic..
It is not recurring again and again and again.
as some Eastern philosophies.
or Eastern religions have taught us..
It is not coming back and coming back and coming back again..
Not reincarnation, no..

$^{241}$Not recurrence of eternity, no..
The biblical concept of time is linear..
It is going toward an end, the eschaton..
So in the Bible, human history,.
time has the ultimate purpose..
There's a destination in human history..
We're going forward, coming to the end,.
reaching the eschaton..
And when the eschaton comes,.
the second coming of Christ, the yellow arrow,.
the second coming of Christ,.
the new heaven and the new earth.
as revealed in the book of Revelation will come,.
will be reality..
So we are all moving toward the end..
Practically speaking, there are two ways..
There are two ways you and I will see Jesus..
Okay, there are only two ways..
The first way, of course,.
is the second coming of Christ, right?.
That's what the New Testament, the olive discourse,.
what Jesus taught about on Mount Olive,.
are telling you..
The end of the age is coming to the end,.
will be conclude at the second coming of Christ..
And we will all see God face to face..
If he comes today, before this sermon ends, okay,.
we will see him face to face..
That's one practical way we will see Jesus..
And we will all, and each and every one of us,.
will need to give an account of our whole life..
That's one way..
The second way we will see Jesus,.
is that when your life ends,.
often without warning, often without warning..
Friday night, last week, I received a text.
from one of my pastoral staff,.
a brother I know from another church up north..
I've known him for years..
No signs of warning..

$^{281}$He isn't that old, okay?.
He's pretty healthy..
He was at home last Friday,.
but all of a sudden he fell unconscious..
All of a sudden..
His wife is trying to wake him up, right?.
But she was not able to wake him up..
He's now in critical condition..
So we were praying for him..
Since Friday night, and the latest news I know.
is that he's not hospitalized..
He got out of the hospital, thank God..
So he's heading home..
But it's a wake up call..
That our life will end, right?.
You guys all younger than me, you would say,.
"Heh, you go first," right?.
No, not necessarily..
If you go to cemetery,.
you look at the years on the tombstones, right?.
Many of them are younger than you..
Have you ever taken a tour to cemetery?.
I remember I was in Vancouver..
I was leading a youth group, right?.
There was one Friday night I got a creative idea.
because there was cemetery by the church, you know?.
So we took a tour..
We took a tour at night to the cemetery..
So I assigned the students..
They were about 20, 25..
So I gave them assignment..
I said, "Go around and survey the cemetery..
"Write down the years of their birth and their death..
"Come back, report to me.".
You know, no surprise..
They found quite a few younger than them..
Young people, right, high schoolers..
Take a tour to cemetery..
So we got caught, right?.
We all need to give an account to God himself, right?.

$^{321}$Two ways you will see Jesus face to face..
The yellow arrow, he comes..
Everybody will see him face to face..
Give an account..
Second practical way is that you go see him, right?.
You go see him first..
Ultimately, it's the same.
because we all need to give an account..
2444, Jesus said, "So you must be ready.
"because the Son of Man will come at an hour.
"when you do not expect it.".
When you do not expect, he will come..
Again, eschatology in the New Testament.
is not teaching us to escape..
It's not teaching us to get away from this world, right?.
Oh, this world is so much trouble..
I would rather Jesus come and then I would escape, right?.
I don't need to go to work tomorrow, right?.
I don't need to deal with this or that trouble tomorrow..
That's not the spirit..
That's not the essence of eschatology..
The teaching of end time is not to help us to escape, right?.
It's to help us to engage..
It's to help us to participate.
into what God is doing in this world..
To fulfill his will, to use our gifts..
So the Bible never tells us, right,.
to sell everything of yours,.
do nothing and wait for the second coming..
The Bible never does that..
The first generation of disciples, they never did that..
Even though they saw all the signs.
and they rightly expect the second coming of Christ.
happening in their lifetime..
They never just quit their job.
and just go to the mountain and wait for Jesus to come..
They engage in their lives..
They preach the gospel..
They use their gifts to build up the church..
They go everywhere with the gospel, right?.

$^{361}$They have good works out of their faith..
It's to engage, it's not to escape..
That is why after all this teaching,.
in the end of chapter 24 to 25,.
Jesus has a series of parables..
Now these are the key points..
Sometimes our focus is on the signs, right?.
Oh, there's another war coming..
There's rumors of war, right?.
There are earthquakes and all that, right?.
Sometimes our focus are being driven.
to those sensational stuff, right?.
All this end time talk are usually sensational,.
trying to get your attention,.
trying to shift your focus away from this world, right?.
Not looking for another sign, but Jesus focus on what?.
That is the conclusion of his entire two chapter long.
teaching on the last days..
He taught us a series of four parables here..
The first parable, 24, 45 to 51,.
parable of faithful and wise servant..
Always ready to give an account to his master..
When his master comes back, he is able to give an account..
The second parable, the 10 virgins,.
or more appropriately, is more like bridesmaids..
As according to the Palestinian wedding custom.
of first century, they have tons of bridesmaids.
waiting for the groom to come back.
from another village, right?.
When the groom leaves the village to the bride's village.
to pick up the bride, to come back often at night,.
often because of the celebration.
and because of the good time that they had,.
the groom will arrive late..
And the bridesmaid of the groom's village.
will have to be prepared with the oil ready.
because often he will come back at night..
And you have your torch and your oil ready at all times..
Stay watchful for the groom's coming..
The five wise bridesmaid, they got oil prepared..

$^{401}$Even the groom comes back delayed..
They are ready to go out and welcome.
and to be part of the banquet..
The five foolish one, they did not have the oil ready..
So they could not enter the banquet..
The third parable, 25, 14 to 30, parable of talents..
The master's giving his money and trust them.
to three of his servants..
One give 5,000, the other one 2,000, the other one 1,000..
The two with 5,000 and 2,000,.
they invest the money to increase the asset of the master..
The third one took the 1,000 and bury it under the ground..
And then the parable says, after a long time,.
long delayed period of time, the master comes back..
The 5,000, 2,000, they gave double of the asset..
The 1,000 are being cast away..
Not prepared to give an account..
Wasted the talent that the master has given him..
Now again, it's using the first century context.
to tell the parable..
But the essence, the teaching is easy to grasp..
Is that God has given us and entrusted talents.
and treasures into our lives..
Are we, in light of his coming,.
are we investing and doing everything we can.
to increase the asset of the master's property?.
As we build the church together,.
as we preach the gospel,.
as we build up our spiritual life, right,.
become more mature disciples,.
we're building up the asset of the master..
We're building up the kingdom of God, right?.
And trusted with the talent..
You and I need to work to increase the value of the kingdom..
The last parable, the sheep and the goats..
Yes, when Christ comes again,.
he is the shepherd..
He can separate..
He can tell who is real and who is not real question..
No, it's easy to look like a question, right?.

$^{441}$I mean, don't you know if you've been to church.
for some time, you can act like a question..
And maybe you act like a question to the point.
that you thought you were really a question..
Because it's easy to just follow a script, right?.
Sing the song, right?.
It's very easy..
Most of the songs that the leaders are leading us.
are usually easy to follow..
We can all sing, right?.
We can all pray the same prayers, right?.
We can all read the same Bible..
We can all do many things, right?.
But externally, we look the same..
However, at the core, we need to confront ourselves.
whether we are acting like a question.
or are we being a question..
It's not about doing..
It's our internal being,.
whether we really have a personal saving relationship.
with our Lord Jesus Christ, right?.
And that's why Good Friday is so important..
Because one day each year,.
it's not a meaningless routine or ritual.
that we just go with the motion.
and thought that we're okay..
Good Friday, why do we call Good Friday?.
But what's so good because a man die on the cross, right?.
There's nothing good about it..
But if you think about it,.
what the Son of God has done on the cross,.
it is good for you who receive..
It's good for me..
It's grace for me that I can be called.
the son or daughter of God..
It's a reflection time,.
it's a time of reminder to confront ourselves.
whether our call is real..
Right?.
You know, in many other parables in the book of Matthew,.

$^{481}$Jesus said in the last day, when Jesus comes back,.
there were many who approach Jesus and say,.
"I know you, I serve you,.
"I even cast out a demon in your name.".
Right?.
You know what did Jesus say?.
He said, "I don't know you.".
Because not everyone call me Lord, Lord,.
will enter the kingdom..
Only those who listen to my word and do my will,.
and they will be wise..
They will be wise and they will be ready.
to come into the kingdom..
So four parables, one after another,.
and trying to push us, trying to confront us.
with the ultimate question of life..
Right?.
Whether we are truly related to God,.
whether we are ready and prepared for the kingdom..
Last time I asked you the three Ts,.
your time, your treasure, and your talent..
Where is your heart?.
You know, Jesus said where your treasure is,.
where your heart is, right?.
You can lie to yourself..
Where's your time?.
Do you manage your time well?.
So that you can use your time to extend,.
to increase, to extend the kingdom of God..
Your talents, right?.
Every one of us has talent, right?.
Some of us have $5,000, some of us have $2,000,.
some have \$1,000, but there's no comparison here..
Because remember what the parable says, right?.
The one who received $5,000 gave $5,000, double..
The other one who received \$2,000 gave doubles,.
that's \$4,000, right?.
God is calling us not to be a fanfare..
God is calling us to be faithful, right?.
Faithful to what he entrusts to us..

$^{521}$So these questions I give it to you..
Let's be prayerful and be reminded,.
are we truly related to Jesus Christ?.
Are we listening?.
Are we doing his will?.
Are we prepared for his return?.
Let's all pray together..
Jesus, you are the son of God..
The glorious son of God..
And you will return someday, unexpectedly..
You will come..
Because we know that you came..
You came the first time..
To die on the cross, to raise from the dead,.
to accomplish salvation, the glorious salvation.
for all of us who believe..
Yes, you came..
That's why we know you will come..
According to your word and promises,.
your second coming will happen anytime..
Or we will see you anytime..
And Jesus remind us to be your faithful and wise servant..
Help us to examine what we have,.
our time, our talent, our treasure..
Help us to ask the question,.
are we using all that you have entrusted to us.
to increase the asset of your kingdom?.
To help extend, advance your kingdom..
Jesus, help us to be alert and to be watchful.
at all time..
Live every single one of our days.
responsibly and faithfully..
To do all things that will be pleasing to you..
May you continue to grow our faith and grow our church.
so that your name will be exalted in Christ's name we pray..
All right, guys, hope to see you again soon..
\newpage



\section{以西結書}
\label{sec:sGka_a0cCzQ}
\textbf{亂世啟示 : 歌革的預言 (以西結書 38-39 章) - 陳崇基牧師}
\newline
\newline
連結: \href{https://youtube.com/watch?v=sGka-a0cCzQ}{\texttt{https://youtube.com/watch?v=sGka-a0cCzQ}} ~~~~ 語音日期: 2022-04-04
\newline
\newline
\hyperref[sec:i4txNEY8_xE]{\small{< < < PREV SERMON < < <}}
~
\hyperref[sec:index_chronic]{\small{[返順時目]}}
~
\hyperref[sec:index_scriptual]{\small{[返順卷目]}}
~
\hyperref[sec:jwup3nNhxBw]{\small{> > > NEXT SERMON > > >}}
\newline
\newline
$^{1}$各位好 我們現在講啟示錄和天啟文學.
雖然我們不是看啟示錄.
但這兩次我們是做一個延遲.
做一點擾導 因為我們需要照顧一些概念.
或者講到末世的時候.
一些中末論的經文.
所以我們上一次和這一次.
我們都會處理一些在啟示錄以外的經文.
但和啟示錄是非常有關.
今天我們要處理一段比較難的經文.
叫做《哥格的預言》.
在以西結書38和39章.
講起這一段聖經.
相信大家比較陌生.
所以我們今天試一下用腦.
想一下如何去理解一些這樣的經文.
有時候聖經有很多很淺白的東西.
三歲小孩都看得明白.
John Fusick先過都會明白.
但也有一些比較深一點.
比較困難一點 棘手一點.
這一段的聖經是其中一段比較棘手一點.
不如我們先看看聖經.
38章第一節.
耶和華的話臨到我.
即是以西結.
他說 人子啊.
你要面向馬各地的哥格.
就是米切土巴的大王.
向他說預言.
主耶和華如此說.
米切土巴的大王哥格.
看啊 我要與你為敵.
我要把你調轉過來.
用鉤子鉤住你的腮甲.
把你和你的軍兵馬匹騎兵都帶來.
他們當中有波斯人.
古實人 忽人.
還有哥滅人和他們的軍隊.
北方極處的陀拉馬族和他的軍隊.

$^{41}$這許多民族都要跟著你.
你要聚集到你那裡的軍隊都要預備.
過了多日你必被差派.
到末後之年.
你要來到以色列上九荒涼的山上.
你和你的全軍.
並跟隨你的許多民族.
都要上來雨暴風刮來.
勿魂遮蓋地面.
一段經文看了之後不太明白.
但這裡很簡單.
兩章的聖經38,39.
以西結說.
其實說到將來.
所謂的艾斯吉時.
末中的時候.
將來以色列國會發生什麼事呢.
在古以色列國.
以西結在公元前大概五百六百年左右.
一位先知.
他是首批去到被擄巴比倫的.
是一個被擄的先知.
被擄巴比倫的時候.
上帝就去呼召他.
來講說話.
講了很多章到38,39章.
就講到將來會有一個所謂大戰.
這個王叫做哥格.
他氣勢磅礡.
就會出兵來攻打以色列盟.
然後這兩章就講到他都還沒出手.
神就將他打敗.
為什麼會講這段聖經呢.
因為這段聖經.
在最近因為.
俄羅斯入侵烏克蘭.
所以有很多時候在坊間上.
流行的解法就說.
這個哥格王其實就是俄羅斯王.
所謂的普頓.

$^{81}$所以我們需要去了解一下.
到底這段聖經是講什麼呢.
首先講一下翻譯.
翻譯很重要.
翻譯準不準確.
或者怎樣去.
連繫到我們今天聽到的一些字眼等等.
會不會有一些誤導的成分呢.
你看到和合本翻譯38章第三節.
講到哥格王的時候.
他翻譯了羅斯.
原文是Rosh.
米切土巴之王.
Nisi就是王子.
或者王的意思.
領袖的意思.
哥格就是他的名字.
然後你看到和合本的修訂版.
之後他改變了.
就將他翻譯成為哥格地的哥格.
就是米切土巴的大王.
那個key在哪裡呢.
就是Rosh這個字.
到底是翻譯成為一個地名.
好像和合本本身的翻譯.
還是翻譯成為一個形容詞呢.
大王.
因為Rosh就是head的意思.
在希伯來文.
就是領袖的意思.
如果翻譯成為一個地名的時候.
就好像有很多流行的坊間的解法.
因為是Rosh.
你轉了響音.
做了u才算.
就是Rush.
你就聽到是Rush.
因為這樣的緣故.
所以有些人就認為.
這個21世紀裡面有的新聞.

$^{121}$就是這裡說.
這個哥格王就是Rush.
或者Rush的統領.
我們怎樣去理解這些事呢.
首先我們先講一個方法.
我已經講了很多次了.
我們看一段聖經的時候.
我們需要由聖經當時的時局來入手.
就不可以用我們今天21世紀的時局來入手.
即使這些的預言.
或者天啟文學.
是有它的將來性.
可能它有一些的principle.
是在講將來.
也都會應用得到的.
但我們必須要用當時人的政局.
他們的背景入手.
Rush這個字.
在希伯來的聖經.
在近東裡面所有的文獻.
都不是指一個地方的.
沒有一個地方叫做Rush.
所以這個和修版.
就譯做大王.
Rush, the sea rush.
大的王.
統領的王或者領袖的王.
這個就比較貼切一點.
米切和土巴.
用一些語言linguistic的connection.
也都connect不到.
今天我們所講的所謂俄羅斯.
因為當時米切和土巴.
我們知道的地方.
大概就是今天的小亞細亞.
或者土耳其的地區.
沒有上到那麼北邊.
並且我們知道Rush這個字.
在主後11世紀才出現.
所以在近東.

$^{161}$以西結書寫成的時代.
是不會想到Rush.
俄羅斯這回事.
所以如果我們將21世紀.
我們現在想到的Rush.
我們就讀進去這一段聖經.
在方法上就本末倒置.
我們繼續看下去的時候.
我們就問一下.
到底這個所謂Gorg.
到底有沒有一些線索給我們.
讓我們知道.
他的身份是什麼呢.
不同的學者就有不同的見解.
你看到這裡列出.
其中的五個.
可能是Gugu.
可能是Gagu.
可能是Gaga.
在阿瑪拉德一些埃及的外交文獻.
就有Gaga這個字.
不是Lady Gaga.
或者是一些蘇美文.
Gugu意思就是黑暗.
Darkness的意思.
或者是一個暗號.
是指巴比倫.
所以不同學者有不同的理解.
你看到這些不同的提案.
不同的意見.
無論我們有沒有一個.
實在的答案都好.
都不是指到以色列最北面的俄羅斯.
你說不是.
因為地理上俄羅斯真的在北面.
有人說你把耶路撒冷畫一條直線.
一直向北的話.
就會差不多去到莫斯科這個地方.
你說不是吧.
應該是這樣吧.

$^{201}$很清楚.
因為聖經裡面說的這個Gorg.
是從北面極遠的處.
北面極處而來的.
所以我們就見到.
如果真的這樣的話.
一定是現在21世紀的新聞重現眼前.
北面極處到底是什麼意思.
不如看看經文.
到底先知書舊約裡面說到.
以色列的北面極處是什麼人.
或者什麼民族.
我們選了幾段聖經.
五段.
耶利米斯一章十三節.
耶路撒冷畫第二次臨到我.
看見一個燒開的鍋.
即是有大禍臨頭.
從北面而傾.
必有禍患從北方發出.
我要照北方的列國的蠻族來懲罰耶路撒冷.
因為他們拜偶像利希.
這時耶路撒冷說的.
二十五章第八節.
所以耶路撒冷說什麼呢.
他說你們不聽我的話.
看我必照北方的種族.
和我僕人巴比倫王.
來到攻擊這地.
繼續看下去.
耶利米斯五十章.
看下有一個民族從北方而來.
是從地極而來的.
北面的地極.
即是最遠的地方.
這個王帶著很多的人.
就會來到耶路撒冷.
會攻打耶路撒冷.
想問一下你說的是誰呢.
很明顯就是在說利布格尼撒.

$^{241}$巴比倫王.
然後一章第四節.
說回以西結書開頭.
有狂風從北面來划起.
二十六章第七節.
主人說如此說.
看下我必駛諸王的王.
就是巴比倫利布格尼撒.
領著他的兵馬軍隊.
許多人從北方而來.
你會奇怪.
因為巴比倫很明顯.
就是在耶路撒冷的東面.
直接在東面.
為什麼說到巴比倫.
來到攻打的時候.
就會說到在北面.
甚至在北面的極遠之處.
如果你知道地理的話.
你就知道.
從以色列的方向來看.
任何入侵的.
都會從北面而來.
因為他們會按著綠色的地方.
從埃及左邊.
以色列到上面的敘利亞.
然後到右邊的不達米亞.
這個地方.
這個叫做Fertile Crescent.
就好像一個月亮一樣.
彎彎的.
這個就是剛剛足夠雨水.
可以定居的一個地方.
所以那時候的文明.
所有的國家.
軍隊商隊.
都會按照這個彎月的方向.
來行軍或者行商.
所以對以色列人來說.
通常入侵的.

$^{281}$就一定是像黃線一樣.
那個黃的箭嘴.
就一定是在北面而來.
所以說到巴比倫王.
來到攻打的時候.
他就不是從東面來.
他就一定是從北面下來.
所以要了解聖經地理.
是很重要的.
最近我看了一篇文章.
是一個歐洲的歷史教授.
他就寫了一些文章.
為什麼呢.
因為他說.
來大學的一二年級的學生.
在歐洲的學生.
他們去到學歷史的時候.
原來他們已經對歐洲的歷史.
有相當的認識.
但是對歐洲古代的歷史.
認識的源頭.
是從哪裡來的呢.
那些學生原來是打機那裡來的.
有一個在歐洲.
很出名的.
或者很多人玩的.
一些的Strategy Game.
就講到一些歷史.
有些對有些錯.
所以他就寫文章.
來看看到底中間有什麼差異.
不一定坊間所講的東西.
就是對的.
也不是一定最多人點擊的影片.
或者最多Like的東西.
或者最多人聽的KOL.
就一定是對的.
當然我也不一定是對的.
但我想告訴你.
一定是有方法的.

$^{321}$在我們講到啟示錄.
和這些預言的時候.
中末論的時候.
我們一定先搞清楚那個方法.
是不是.
不是隨便選其中一個的講法.
那我們就行.
是應該要想想那個方法.
看聖經或者解聖經.
我們需要從當時的背景來開始.
講了這麼久.
一個歌格王到底是誰.
很肯定我們知道不是誰.
這個歌格王是誰都好.
就是和今天或者之前的蘇共.
俄羅斯Russia.
是沒有任何語言歷史和地理的關係.
如果你看清楚聖經的話.
你就發現這種坊間流行的.
說這個歌格王就是俄羅斯王.
或者以前的蘇共的話.
這種的詮釋.
你就要留意.
這種詮釋.
在二次大戰之後.
到現在都是非常流行的.
特別是二次大戰到90年代初的時候.
當時是美蘇之間有冷戰的時候.
這是我們近代史.
發生的事情.
有很多時候我們就將這些東西讀進去.
這種詮釋在二次大戰World War II之後.
是非常流行的.
從美國甚至在環教會.
都是非常流行的.
特別很多時候就配合了.
一些我們已經預設了的時間表.
一些的Timeline.
就是什麼時候我們會被提.
教會突然在地上消失.

$^{361}$然後七年大災難就會發生.
史無前例的最後一個七年的禍災.
就會來到地上.
是一個已經既定了.
首先這個已經Debatable.
到底是不是有一個這樣的Timeline.
或者聖經有沒有清楚講這個Timeline.
這些是豪華.
但是有了一個Timeline.
然後我們就將所有的寓言.
或者一些天啟文學的經文.
我們就塞進去.
所以我們由我們剛才所講的.
從地理從歷史從一些語言來講.
是這樣講不通的.
如果你將Rosh.
中文以前翻譯不準確的時候.
羅斯翻譯做.
就當成俄羅斯的話.
就等於將古代的近東文獻裡面.
提過Gaga.
我們就說是Lady Gaga.
是一樣這麼荒謬一樣這麼Nonsense.
這些我想問.
或者問一下各位.
你覺得這些是不是Biblical Interpretation.
是不是聖經的詮釋.
還是這些是Cold War Interpretation.
是用了我們現在的冷戰.
新冷戰.
之前World War II之後.
直到90年代初的時候.
就是以前的所謂Cold War.
那幾十年.
我們用了這樣的模式.
來讀入了聖經.
是一種Cold War Interpretation.
到底是屬於哪一種.
想清楚那個方法.
到底為什麼會搞成這樣.

$^{401}$或者解成這樣.
1992年的時候.
有一個聖經學者叫做Paul Tanner.
Paul Tanner就是一份福音派.
首先搞清楚.
剛才我說的那些解釋.
你不要說Mark Sir你是新派.
所以你喜歡新的東西.
這個Journal.
叫做Journal of Evangelical Theological Society.
是一個保守的福音派.
神學的期刊.
JETS.
Paul Tanner.
在1992年的時候.
他就寫了這篇文.
是一篇解經的文章.
解經文章的Title說什麼呢.
他說Daniels King of the North.
Do we owe Russia an apology?.
先說到.
《但義理書》.
他也有論到《二世傑書》.
38,39章.
他說到《但義理書》最後那兩章.
說到有一個King of the North.
這個King of the North是什麼呢.
冷戰的interpretation.
冷戰的詮釋.
就在.
大行其道.
在那幾十年.
甚至今天也有人說.
就說這個King of the North.
在《但義理書》裡面出現的.
就是Russia或者之前的.
蘇共.
但其實如果你按照.
當時的背景來看.
《但義理書》十一,十二章.

$^{441}$是說到的.
是希臘時代.
當時是Solusid Empire.
當時是希臘帝國.
去欺凌猶太人.
所以那段聖經.
背景不是現在冷戰.
不是蘇共.
而是當時的希臘帝國.
他就說.
解了這段聖經的時候.
他就說.
Do we owe Russia an apology?.
留意這篇文是.
哪一年寫的.
1992年寫的.
蘇共什麼時候解體.
91年.
沒錯.
92年他就寫這篇文.
很有趣的一個.
的title.
如果我們真的用.
一種co-op interpretation.
解了幾十年.
然後蘇共又解體的話.
我們就真是.
需要和Russia道歉.
事實上.
如果你看地圖的話.
你將哥格和他的七個聯盟國.
放在地圖上.
以我們現在知道的.
已經知道不少了.
一些古代的文獻.
當時以西結的文件.
他時代的文件.
你將這七個聯盟國.
放在地圖上.
其實指的是當時近東的世界.

$^{481}$當時他們的領域裡面.
比較邊緣的.
所謂國家的民族.
哥格到底是誰.
正如我所說.
學界是沒有什麼共識的.
如果寫博士論文寫哥格是最好的.
因為還沒有共識.
所以很多人說話.
說什麼都可以.
所以說到哥格的時候.
我們還不知道有什麼共識.
但是我們可以很肯定.
很肯定.
以西結書寫38,39章的時候.
就是說到當時.
被擄的以色列人.
他們需要將來回歸的時候.
他們感到最大的.
使他們恐慌的那種勢力.
哥格和他的聯盟.
是在回答當時以色列時代.
被擄要回歸的猶太人.
他們一個很切身的問題.
就是回歸的時候.
到底我們有沒有安全.
是不是.
我們是不是還在一個恐慌的底下.
我們看看38,39章.
到底脈絡是怎麼樣.
37章說什麼呢.
這個大家看正經的時候.
應該比較熟悉.
37章就是一個儀仗.
是一個枯骨.
那裡以西結看到.
四處都是骸骨.
有很多的骸骨死了在那裡.
然後呢.
是死無葬身之地.

$^{521}$在舊約和近東的世界裡面.
這個就是最極端的咒詛.
所以是一種比喻.
他看到的這些骸骨.
他說耶和華的手按在我的身上.
耶和華藉著他的靈帶我出去.
我在平原裡面.
就看到遍滿骸骨.
其實他就是在代表什麼呢.
是在比喻以色列人.
神的子民.
被擄到巴比倫好像死了一樣.
是受咒了.
然後到第26節.
26節他就說我要與他們.
立平安的約.
我要與這一班神的子民.
雖然被擄好像死了一樣.
與他們立一個永約.
要將他們帶回他們的原地.
安頓他們.
然後使他們人數增多.
在他們中間我再設立我的城鎖.
這裡是在說什麼呢.
以西結向這些骸骨.
去發預言的時候.
講先知的預言.
突然間他的異象裡面.
他看到這些骸骨.
就有神的氣息.
能夠復活起來.
就指出這些神的百姓.
被擄如死一般.
上帝使他們再活過來.
就應許這些流亡的神的子民.
必定有回歸的盼望.
以色列民能夠再活起來.
被擄可以回到.
Exile可以return.
然後才說.

$^{561}$3839章.
我們在看這個.
戈格大軍被神擊潰.
你留意一下當中的描述.
因為時間關係我們不能夠講很多細節.
裡面說到戈格大軍.
被擊潰的時候.
到他們死無葬身之地.
到他們被打敗.
然後屍首是騙人的.
所以有一個reversal在這裡.
37章就是本身神的子民.
好像死的一樣.
然後上帝因照他們.
使他們能夠再活過來.
被擄的人能夠得到回歸.
3839.
那些向以色列民敵對的.
就相反了.
這些大軍被擊潰的時候.
他們就是死無葬身之地.
然後40-48章.
就是講到神的同在.
講到新的聖殿.
然後用一些天啟的文學的方法.
來講上帝在他們的子民中間.
來同住.
如果你這樣看的時候.
上帝問一下你.
你就明白了.
38-39章.
雖然這麼難解的一段經文.
講到戈格大軍.
這個所謂戈格的魔頭.
代表的是什麼呢.
代表就是這些.
到底回歸的猶大人.
他們會不會有安全.
因為他們為什麼被擄呢.
就是被巴比倫帝國所殲滅.

$^{601}$被他們殘殺.
被他們踐踏.
然後他們才被擄到異鄉.
現在上帝說.
我帶你回去.
你能夠回歸了.
對於當時的猶大人來說.
他們最切身的問題就是.
到底我們這個恐慌裡面.
怎麼辦.
我們能不能夠面對.
38-39章就告訴他們.
上帝為你增賤.
上帝會照顧你.
在恐慌當中不重要.
因為上帝會幫你增賤.
這些切身的問題.
不單單是以色列時代的猶大人.
是對的.
就是我們今天都對的.
作為神的子民.
我們有恐慌的時候.
誰會照顧我們.
我們怎麼辦.
上個月我去學院上課的時候.
上了一個星期課.
認識了一個姐妹.
她剛剛在同一間學院.
她就在另一個程式.
因為吃飯的時候有人介紹.
就認識了.
是中國大陸來的一位.
教會的姐妹.
剛剛在我浸信會的學院那裡.
快要畢業了.
商談之後.
她就非常惶恐.
她說很不想回去.
因為她知道回去會被監視.
因為她知道回去.

$^{641}$是不能夠再做.
她的家庭是一個傳道人的家庭.
連傳道人的家庭都叫她.
父母都叫她.
如果你可以的話.
請你不要回國.
惶恐與不安到底怎麼辦.
我不知道你呢.
你有什麼恐懼.
你有什麼令你擔驚受怕的.
在英國有一個社會學家.
叫做Fruity.
Frank Fruity.
他寫了一本.
很多人知道的一本書.
叫做Culture of Fear in 21st Century.
這本書是疫情之前寫的.
疫情之前很多年.
他用一個社會學家去分析.
分析一下我們今天這個世界.
現代的世界.
有一種所謂Culture of Fear.
一種恐慌的文化.
我們現在現代的人.
我們因為戰禍而死.
因為大災難而死.
因為饑荒而死.
或者甚至因為疫症而死.
我不敢說香港或者那些疫症.
還是很掙扎的地方.
我敢說起碼這裡.
起碼我們.
因為這些原因而死的機率.
其實比以往的人已經低了很多.
對嗎.
有沒有誰會因為饑荒而死的.
你有沒有認識一個人是因為饑荒而死的.
我們所謂的饑荒就是.
明明一天吃五餐.
我們變成吃四餐.

$^{681}$明明可以吃宵夜.
但不能吃宵夜.
餓著肚子睡覺.
這些對我們來說是饑荒.
但真正的饑荒.
遠我們非常之.
想都沒有想過.
我們比任何以往一個時代的人.
因為這些原因而死的.
真命受威脅的.
機會率是低之猶大.
但是Ferrari說.
我們就是比任何一個時代的人.
恐懼感更加強烈.
由星期日到星期六.
我們聽見的新聞.
或者接觸到的一些資訊.
都是叫我們擔驚受怕.
恐懼恐慌的一些警暴.
或者KOL跟我們說的話.
傳媒接收到的一些訊息.
我們變成不單止接受了這些訊息.
甚至將它常規化.
有意無意我們就會將它加劇.
或者應該這樣說.
我們會用一個恐慌的角度.
來看任何一件事.
小朋友吃的東西我們怕.
我們吃的東西我們又怕.
出到街我們又怕.
然後有疫症就更加怕.
打仗四角亂我們又會怕.
問題就是我們會不會是一種恐懼文化.
不單止是這樣我們變成一種恐懼文化.
不單止是恐懼文化.
而是恐懼教.
不是恐慌的文化.
而是拜恐慌教.
Cult of Fear.
Frank Fluenti 有一次.

$^{721}$他說他去到華盛頓.
在一個玩具店.
他想買一份給小朋友的禮物.
因為去探一個朋友的.
剛剛添丁.
給錢的時候職員就跟他說.
Sir, will you be interested in our new product?.
即是我們有一個新的產品.
你會不會有興趣呢?想買呢?.
是什麼來的呢?.
原來是一種室內的頭盔.
是給小朋友戴的.
不是因為他需要戴.
而是怕他在家跌倒.
撞到這樣.
所以就給他一個indoor的helmet.
做得很漂亮.
他以為是開玩笑.
誰知道原來整個架都是那些產品.
Culture of Fear.
他另外一本書.
我想介紹給你.
叫做 Paranoid Parenting.
那個subtitle很有趣.
他說.
Why ignoring the experts.
may be best for your child?.
當然我不是說.
你一定要做這個 做那個.
Parenting 你不要做這個 要做那個.
我想像不到在這個時候.
我的小朋友剛出生.
我真的想像不到那種氣氛.
我的小朋友已經大了.
他現在大學了.
不過我想給你另外一個角度.
或者另外一些值得我們討論的東西.
Ferrity 說過一件事.
他說長大後一個小朋友.
到四五六歲的時候.

$^{761}$你去跟他或者多說一些.
或者太過強調.
一些你不要試那件事.
那件事太大壓力.
你這樣搞不定.
如果你不自覺地灌輸他的話.
對於社會學家來說.
這叫做有一個名詞.
叫做Cultural Script.
文化的劇本.
如果你四五歲開始去灌輸他.
他有一種這樣的Vulnerability.
即是很脆弱 這個世界很危險.
到九歲的時候.
這個小朋友就會自然地形容自己.
小朋友會說.
Man, I am so depressed.
I am overwhelmed.
I am so stressed out.
是一種Self-fulfilling prophecy.
因為跟了那個Script.
他自己慢慢跟著那個Script去反應.
令到他更加Vulnerable.
我不是說我們一定要做這件事.
不做那件事 說這件事不說那件事.
I am not talking about that.
只是給一些 可能另一些角度給我們看到.
或者Ferrity給我們看到.
我們是不是在一種Culture of Fear裡面.
一種Cult of Fear.
我們拜的沒有其他的.
我們只是拜恐慌.
Fear.
不久之前.
我兒子跟我說了一句很簡單的說話.
因為你用你的方式來養育我.
當他這樣說的時候.
他未必真的意思.
他可能有其他的原因.
他不一定是說我很有家教.

$^{801}$他可能是說.
可能他心裡的意思是.
你一直都不理我.
怎樣都好.
他跟我說什麼呢.
因為他在Take a break的時候.
他回到Chicago.
他就幫一些補習社.
跟那些小學生補習.
去幫他們.
他們跟小學生溝通的時候.
他就經常跟他們說一些Stories.
Stories of my life.
他就告訴他們.
他在哪裡做過一些瘋狂的事.
然後他就跌倒.
然後有疊疤.
然後可以展示給小朋友.
小朋友個個都很光滑.
他就說.
現在的小朋友很有趣.
好像沒有什麼故事.
So sad.
小朋友為什麼沒有瘋癲的事做呢.
好像我這樣.
他就說好像我這樣.
可以做一些瘋癲的事.
當然他做過一些所謂瘋癲的事.
跟我那一代.
甚至我上一代.
那些瘋癲的事根本就不能比較.
但他的意思就是.
為什麼無端端的Take risk.
冒險Adventure.
Exploration.
那種的經驗.
回到Ferrity的說法.
他說人生其實就是要這樣.
不可以消除Risk.
不可以消除冒險.

$^{841}$反而應該要問.
我們有沒有Courage.
有沒有勇氣.
去Take good risk.
或者教我們下一代.
是不是應該要Take good risk.
因為到他長大的時候.
他總要面對一些挫折.
他用了一個例子.
Ferrity用了一個例子.
就是一個人長大之後.
其中一個最大的Risk.
就是談情說愛.
你想想談戀愛Fall in love.
是一個很大的冒險.
你有沒有想過.
如果他一直沒有學習.
跌倒再爬起來.
那種經歷的時候.
到他長大了.
他一面對感情失敗的時候.
那就更加嚴重.
是不是.
當然我不是說我們要亂來.
我不是說我們不需要用智慧.
當然我們需要.
但是我們需要用一個角度去看.
我想我們需要問.
除了恐慌什麼之外.
我們有沒有質問過恐慌的本身.
Fear itself.
或者恐慌本身.
有沒有侵蝕我們的信仰.
侵蝕我們對上帝的信仰.
我們作為神的子民.
是不是仍然亡亡不安.
我們就只是被恐慌擺佈.
我講這些是因為我想告訴你.
38至39章以西結書.
不是講預先寫好的新聞報導.

$^{881}$其實是在解釋一個很切身的問題.
就是以西結時代流亡巴比倫.
必須回歸的人.
他們要面對很多恐慌.
巴比倫曾經擊敗過我們.
是一個黑暗的勢力.
殘暴過我們 踐踏過我們.
到底這個魔王怎麼辦.
我們回去的時候會不會有平安.
所以很多時候我們學者在猜.
那個哥格王是什麼王呢.
我就認為.
以西結在他的經文裡面是有提示的.
那個哥格王到底是誰都好.
起碼形容的方法.
其實就是形容巴比倫帝國暴君的方法.
我在這裡的表很簡單 列出六樣.
由北邊來的 尼泊爾加勒薩是由北邊來的.
哥格王也是從北邊來的.
有很多民眾 有很多軍隊 兩者都是.
來搶掠 兩者都是.
但是最後這個 帶領這個字.
耶和華帶領 是他引巴比倫王來.
也是他引哥格王來.
既然他引得他來的話.
上帝說 這些恐慌是我已經得勝了.
我得勝過巴比倫王.
我也可以得勝哥格王.
這個哥格代表的是一些死亡的魔鬼 即是魔頭.
他是不會贏的.
因為最終是耶和華作王 是上帝作王.
死亡 疫症 恐慌.
不應該是我們生命裡面 結案陳詞的那個.
結案陳詞的那個應該是.
那位掌管世界歷史的主宰 耶和華自己.
講得很清楚 這是上帝的保守.
不會拿走你要信心為主冒險的機會.
有恐慌 沒錯 有哥格的威脅 沒錯.
有黑暗的權勢 仍在 YES.
但是你要知道 有上帝的保守 有上帝的大力.

$^{921}$但你也要知道 上帝的大力是不會奪去.
你要憑信心來冒險.
Take good risk.
可能會跌倒的這種機會.
可能會在冒險當中跌跌碰碰的那種經歷.
上帝是不會因為你信教他 而拿走這件事的.
因為上帝不會在你身上 拿走你在危險裡面.
在危機裡面 見到上帝得勝的那個寶貴的經歷.
你看不看到 信心的道路.
被擄要回歸 或者上帝要帶你走下一個chapter.
走下一步 是要你生命裡面有上帝寫的故事.
最終你能夠見到上帝的榮耀.
這個豈不是以色列書最後結束的主題.
結案陳詞的 不是哥格 魔頭的恐慌.
不是fear 而是上帝自己.
最後一節 48章35節 以色列書最後一節.
從此以後這城的名字 必稱為耶和華的同在.
上帝住在他的子民中間.
可能有恐慌 可能有很多不安.
但是在不安裡面 我們知道上帝保守和大力.
而且上帝已經得勝.
來到受苦節 復活節.
提醒一下我們自己.
我們信心道路裡面 是否只是避 不敢 害怕.
怕這個 怕那個.
還是我們的信仰 是可以給我們能量 給我們資源.
我們可以學習 行信心的道路.
可以同心協力.
主我們多謝你 你的話語來提醒我們.
你的話語告訴我們 你與我們同在.
沒錯 我們生命的道路裡面.
好像往日的以色列人.
被迫回歸 回到故土的時候.
仍然有很多的恐慌 很多的不安.
但是我們知道 你自己的同在.
我們更加知道 你對我們帶領 你的保守.
不會拿去我們生命當中 一些不安的事情.
反而在不安的事情當中 讓我們能夠經歷到上帝的愛.
和你自己的恩典.
主我們感恩.

$^{961}$你就在復活節 受苦節 你見到我們十字架的恩典.
上主我們求你 讓我們謹記.
你就是那位與我們生命同行的神.
祈禱奉主耶穌的名字.
今天是我們聖餐的時候.
讓我們在主的禮物裡面.
\newpage



\section{創世記}
\label{sec:jwup3nNhxBw}
\textbf{創世記第6章 (6.1) [神的兒子的叛變] - 背景釋經 (陳崇基牧師)}
\newline
\newline
連結: \href{https://youtube.com/watch?v=jwup3nNhxBw}{\texttt{https://youtube.com/watch?v=jwup3nNhxBw}} ~~~~ 語音日期: 2022-04-12
\newline
\newline
\hyperref[sec:sGka_a0cCzQ]{\small{< < < PREV SERMON < < <}}
~
\hyperref[sec:index_chronic]{\small{[返順時目]}}
~
\hyperref[sec:index_scriptual]{\small{[返順卷目]}}
~
\hyperref[sec:Gv0nluM2dc8]{\small{> > > NEXT SERMON > > >}}
\newline
\newline
$^{1}$好 今次我們就看到第六章了.
我們由第一章現在看到創世記第六章.
六章我們首先在餘下的這幾個部分.
我們就嘗試去看第六章 一至到第八節.
講起這段經文 我們先讀一下.
是一段很難的經文.
解釋上有很多的爭論.
先讀一下.
這裡說什麼呢.
這裡就是說.
人在地上開始增多.
由生養女兒的時候.
神的種子看見人的女子美麗.
就隨意挑選娶作妻子.
而我話說人既然是屬肉體的.
我的靈就不永遠住在她裡面.
但她的日子還有一百二十年.
在那些日子和以後的日子.
有巨人在地上.
神的兒子和人的女子結合.
就生了上古英武有名的人物.
然後看見人在地上罪惡很大.
終日心裡思念的盡都是邪惡的.
於是然後說後悔.
做人在地上.
心中憂傷.
然後說我要把我創造的人.
從地上消滅.
無論是人還是牲畜.
是爬行的動物和天空的飛鳥.
我都要消滅.
因為我後悔造了他們.
只有羅亞在耶和華面前蒙恩.
這一段聖經.
起碼每一句.
其實你看清楚的話.
每一節每一個字.
都會有很多爭論.
所以我這裡就highlight了幾樣.
第一 神的種子.

$^{41}$到底神的兒子.
所謂神的兒子們.
到底是什麼人.
另外就是120年.
他的人類的日子還有120年.
這是他的壽命還是指到什麼時期.
第四節就是在那些日子和以後的日子.
換句話說.
你看文理上下文.
洪水的日子和洪水以後的日子.
有什麼呢?.
有巨人在地上.
然後神的兒子和人的女子結合.
生出上古英武有名的.
這個巨人是什麼身份呢?.
是什麼人呢?.
第六節.
就是翻譯的問題.
後悔做人在地上.
到底這個翻譯是不是需要更改呢?.
原文應該是說什麼呢?.
所以這裡有很多很多東西.
我們要去討論.
還有很多東西引起不同的爭議.
我們希望能夠去解答一下.
首先就要解答神的種子的問題.
到底他們是誰呢?.
我們就試試看原文.
我們只選一兩節來看.
神的種子.
第二節.
這裡說快逆轟.
他們看見.
你記得希伯來文是動詞開始.
然後就是主題.
Bene Elohim.
Bene就是二子們.
然後呢?.
你記得在創世的第一章.
我們已經預見了.

$^{81}$就是神的意思.
Bene Elohim.
Ebono Adam.
就是說他們看見.
神的二子.
門見到.
女子.
人類的女子.
那麼.
就看見他們什麼呢?.
Ketuvo Hina.
就是說他們是很好.
他們非常好.
那麼.
這個就是神的二子.
這個名詞.
那麼到底是.
什麼的身份呢?.
那麼在解經的歷史裡面.
你記得我們說到解經的事業.
有經文背後的世界.
我們說過是近東的世界.
經文上下文.
就是經文裡面.
聖經裡面的世界.
不是聖經前面的世界.
就是我們.
我們去看聖經解經的時候.
就不是一個個體那麼簡單.
我們是和整個基督教的教會的歷史.
教會一起來去閱讀聖經.
去詮釋聖經.
所以以往的解經歷史.
都留下不同的見解.
不同的解經家等等.
那麼誰是神的二子呢?.
神的眾子呢?.
那麼就是.
第一個解釋.
就是失特的後代.

$^{121}$你記得第四章尾開始.
說到他們去呼求耶和華的名.
所以這個早期教會.
是有這樣看的.
就是他們是神的眾子.
就是失特的後代.
然後敗壞了.
變成這裡所說的.
神的二子.
那麼所以這些人.
就是這些敗壞的一代.
然後他們就和人的女子.
這個就可能.
在這個詮釋的方法裡面.
就認為這個人的女子.
就是什麼呢?.
就是戈癮的後代.
所以有一個這樣的看法.
當然這個看法不是最流行.
現在比較流行的就不會這樣說.
其實你這樣想.
如果是失特的後代.
那麼為什麼他在第六章的時候.
突然間又變成了.
Bnei Elohim呢?.
何解呢?.
所以很多時候都未必會這樣解釋的.
加普沒有這樣連接下去.
所以很多時候.
這個是非常初期教會的時候.
是會這樣解釋的.
然後呢.
另一個解法呢.
就是非常流行的.
也是初期教會到現在.
都是非常受不同的解經家.
來得到他們的支持的.
那麼就是說到這些神的兒子.
其實是天使.
天上的使者.

$^{161}$那麼就是說.
這些的使者.
是敗壞的使者.
他們見到地上的人的女子.
美麗.
那麼他們就下來挑選.
拿了她來做妻子.
隨意挑選.
取作妻子.
那麼這個解法呢.
天使呢.
理據在哪裡呢?.
就是Bnei Elohim這個字.
這一句話.
或者這個名銜.
Sons of God.
那麼在哪裡出現過呢?.
在聖經.
舊約聖經幾處的地方出現過.
那麼這些出現的地方呢.
上下文呢.
很肯定.
都是說到天使.
就是神的宗旨.
譬如約伯記有幾節的聖經提過.
舊約其他聖經有提過.
Bnei Elohim.
Sons of God.
那麼通常都是說到天使.
那麼譬如這個約伯記一章第六節.
那麼你記得這裡是在天庭.
在上帝的Court.
上帝的宮廷裡面.
然後有神的眾子來到耶和華面前.
其中一個叫做Hazatan.
就是挑戰者.
其實不是撒旦.
中文英文都譯作撒旦.
其實是Hazatan.
就是The Adversary.

$^{201}$那個opposite那個人.
Challenger.
那麼他就問約伯.
為什麼會敬畏神.
起是無緣無故的嗎.
所以第一章就有一個這樣的故事.
那麼就帶出整本約伯記.
那麼在第一章第六節說什麼呢.
有一天.
神的眾子都來.
是立在耶和華的面前.
撒旦Hazatan.
這個其實更加好的譯法.
不應該譯作撒旦.
因為Hazatan.
The Satan.
如果真是一個正確的名字.
一個個人的名字.
是Satan的話.
就不應該加一個The.
不應該加叫做Hazatan.
怎樣也好.
他來到他們中間.
就問Hazatan.
你從哪裡來.
所以這裡上面.
我們這次所關注的.
就是神的眾子.
Bnei Haelachim.
你看到兩個字是一模一樣.
所以一章六節等等的經文.
都讓我們看到.
這個Bnei Haelachim.
似乎是指向天使.
也有很多的解經家是這樣看.
這是第二個的傳息方法.
所以這個傳息方法.
用的理據就是Bnei Haelachim.
你搜查一下希伯來聖經裡面.
Bnei Haelachim.

$^{241}$有沒有出現過.
哪裡出現過.
裡面指的是什麼.
天使.
既然其他地方指天使的話.
這裡照理都是指天使.
第三個的看法就是.
說一些的君王.
一些的獨裁者.
當時叫做Tyrants.
或者一些很霸道的.
一些專政的貴族.
一些君王的人.
Royal Beings.
Royal People.
王族.
你記得我們第一章.
說到神的形象的時候.
在當時的亞述文獻.
近東的文獻.
都有提到這些君王.
是神的形象.
Image of God.
既然他們是Image of God.
是神的形象的時候.
只有這些君王才是神的形象.
換句話說他們是肉身的.
所以在近東的意識形態.
那種近東外邦的神學.
他們就認為既然王是神的形象.
所以王的肉身.
就必定是神靈居住的地方.
在地上.
所以這些的君王.
很多時候起碼在他們的文獻.
亞述帝王留下來的文獻.
巴比倫帝王留下來的文獻.
近東的王會說自己是神的形象.
並且他的肉身是有神性的.
所以這個理據.

$^{281}$不只是說Bernie Hall.
在舊約聖經.
如何用.
就擴大範圍.
看看經文背後的背景.
理據就是這些字句.
希伯來聖經裡面的字句.
譬如這個禮案.
就是Bernie Hall.
不可以單單翻譯文化背景.
大家知道無論是講現代的語文.
我們翻譯的時候.
很多時候翻譯到字.
所謂的linguistic science.
我們翻譯到英文轉中文等等.
但很多時候文化的背景.
背後的意識形態.
是很難去翻譯出來的.
所以這個理據.
就是Bernie Hall.
看看背景.
會不會有些東西幫到我們呢.
在這裡我們可以看看一些史詩.
譬如這個Gilgamesh Epic.
基加美修史詩.
我抽了好幾段.
就是指到這些王.
是有神性的.
或者半神半人.
在這裡可以看到.
Gilgamesh是近東一個很有名的王.
他的名字真的有奇人.
歷史上真的有一個叫Gilgamesh.
並且留下了一些文獻.
是講到有關Gilgamesh.
基加美修君王.
怎樣形容他呢.
他說.
Who is there can rival his king standing.
講英文.

$^{321}$翻譯出來的英文.
And say like Gilgamesh.
It is I am the king.
Gilgamesh was his name.
for the day he was born.
Two third of him god.
and one third human.
這個在第一塊的禮版.
45行裡面這樣講.
他說誰能比美Gilgamesh.
王者的地位.
像基加美修說.
我就是王.
基加美修.
就是他出生的名字.
他的三分之二是神.
三分之一是人.
所以他是Two third god.
然後One third is man.
所以講到Gilgamesh的時候.
就講到他是有神性的.
很厲害的.
並且在其他地方.
講到他是一個很巨大的人.
他不單有神性.
他自己的身量也很大.
他的腳長三爪.
每爪是尺半.
他的腿是半個標桿之長.
他的步幅.
一步有多闊呢.
就是六爪.
大概九尺.
他前半部分.
那裡掉了.
不知道前半部分是什麼.
有多少爪.
你可以想像.
他是一個很巨大的人.
Triple cubit.

$^{361}$What is foot.
How far is leg.
Six cubit.
就是他的steps.
或者他的步幅等等.
所以就形容到他是一個君王.
有神性.
並且是一個偉人.
很大的人.
你見到這個法國羅浮宮.
在巴黎.
這個浮雕.
相信是Gilgamesh的石雕.
拿著一隻獅子.
這隻獅子是Lifesize.
真的是一隻獅子的大小.
你可以看到他.
按照這些神話來說.
他是多大個.
是一個巨人.
另外呢.
這裡提到.
又說到Gilgamesh.
Gilgamesh的肉體.
他的肉身.
是眾神的肉體.
又說回這句.
整首史詩.
經常出現的.
基加美修王.
就是三分之二神.
三分之一是人.
基加美修王.
就是Flesh of the Gods.
然後在另外一些version.
譬如Gilgamesh Epic.
不斷被文字抄寫.
去到不同的地方.
又很流行.
說到他的史詩.

$^{401}$他的威水史等等.
其中一塊的例版.
就是Hittite的version.
即是客人.
去到現在的土耳其.
說他有多高呢.
有十一爪高.
所以是一個很巨型的.
很大的一個人.
並且說到.
基加美修.
他作王的時候.
他是King of Uruk.
就是烏魯克古城的.
一個城邦的王.
他當時是一個獨裁者.
一個Tyrant.
他怎樣獨裁呢.
在其中一個例版.
屬於Old Babylonian Period.
即是舊巴比倫帝國的時期.
大概是早前1800至2000年.
這段時間.
當時是阿拉罕的時間.
有一塊的例版.
我們叫它Pennsylvania.
即是說因為現在在.
賓夕法尼亞州的大學裡面.
那個例版.
在第160項.
就說到他是一位獨裁者.
去到殘害殘暴的人民.
並且去到將那些準的妻子.
和他們去交合.
即是是他先.
然後才給這個新郎.
所以說到他是一些很暴君.
說到這個Gilgamesh.
所以如果你這樣去比較之下.
似乎這個方法.

$^{441}$又是一個近東比對的方法.
讓我們可以比較一下.
創世記第六章所描述的.
和Gilgamesh Epic.
有沒有一些平衡相似的地方.
中間有什麼不同呢.
你記得我們在說比對研究.
Comparative Studies的時候.
我們就一定要看一下.
Similarities和Differences.
有什麼是相似的地方呢.
就是神的兒子.
Sons of God.
Gilgamesh這個史詩.
就說到他又是神的育身.
即是他不是單單一個人.
三分之二是神.
三分之一是人.
然後創世記第六章.
去到描述神的兒子.
描述他們是怎樣.
他們隨意去挑選.
挑選這些女子作為妻子.
然後我們就見到獨裁者.
獨裁者就強奪這些女子.
並且這裡就提到.
創世記就說這些是會死的人.
會扭壞的人.
Mortals.
神就說他們的日子只是120年.
另外在基加美修的史詩.
同樣提到.
這些叫做冥命神.
即是死神.
Nunaki.
這些神就去命定了人的生死.
命定了Gilgamesh的生死.
最後六章創世記.
就去描述這些神的兒子是偉人.
或者有些翻譯做巨人.

$^{481}$原文就是Lebanese.
有時候翻譯做Giants.
巨人.
這個巨人就是七十四億本.
將希伯來文翻譯成希臘文.
就翻譯成Gegentides.
希臘文就是Giants的意思.
即是巨人的意思.
和Gilgamesh的經典裡面.
我們剛才所看的一些文字.
就有相似的地方.
致命英雄.
身高雄偉.
大而可畏.
步幅很闊等等.
就好像說他們是一個巨人.
所以就有三個不同的見解.
正如我所說.
第二個和第三個是最流行的.
Bernard Elohim到底是什麼人呢.
我們從經文裡面再看一看.
有什麼細節.
是我們需要留意的.
還有為什麼我會傾向於第三個的說法.
不是天使下凡.
和人的女子結合.
然後生下後代.
那些後代就叫偉人.
為什麼我會傾向於第三個呢.
當然這些解經的爭論.
二三也好.
甚至一也好.
其實都不會影響我們基本的神學.
所以在這裡我們有很多討論空間.
爭辯的空間.
但我想指出的是.
到第六章的時候.
我們留意的話.
在講到洪水審判之前.
上帝公義彰顯之前.

$^{521}$到底發生了什麼事.
招致上帝要將洪水帶來.
我們再看回第六章一至八節的時候.
你就留意到.
我這裡將一些字眼.
去highlight出來.
第一就是人在地上增多.
下丹.
是在講人.
第二 神的眾子看見人的女子.
下丹.
隨意挑選作為妻子.
又說什麼呢.
對於這件事.
這件事一定是在創世記的脈絡裡面.
是最大惡極的事.
這是一個幾乎末世審判來之前.
一個人類的危機.
招致上帝洪水的審判.
又說人屬肉體.
所以他在講人.
下丹.
不是在講天使.
第四.
這個就是我們下一部分.
之後我們可以解釋一下.
在那些日子.
以後的日子.
即是洪水之後的日子.
有巨人在地上.
神的兒子和人的女子.
又是在講下丹.
第五節.
耶和華又再評估人類地上發生什麼事.
看見不是天使.
是下丹.
在地上.
罪惡很大.
中人思想都是邪惡.
於是又說後悔做人.

$^{561}$下丹.
心裡面憂傷.
然後他自己又說.
我創造的人.
等等.
似乎在講指向的是誰最大惡極.
是下丹.
人類.
似乎不是天使.
所以單從這些經文去看細節.
似乎是指向這一方面.
還有.
在第六章.
有什麼字眼.
是我們從第一章到現在.
我們經常碰見的.
這個就是上下文.
怎樣解經都好.
我們都要配合上下文的細節.
第六章第一至四節.
我們抽出來看.
人在地上.
是下丹.
增多的時候生養女兒.
神的宗子.
Baneha Elohim.
看見Ra.
人的女子.
Mele.
Mele是Toth.
是Toth.
就隨意挑選.
取作妻子Lakak.
接著他就說其他的細節.
這裡有幾個字.
在三章第五和第六節出現.
第三章五至六節.
就是人類第一次的叛變.
你記得.
阿丹 阿蛙 那人與女人.

$^{601}$其實就是每個人.
代表著每個人.
蛇跟他們說什麼呢.
蛇對女人說.
你們要像神Kielohim.
Kielohim.
能知道善惡.
引誘他.
第六節.
於是女人見Ra.
樹上的果子是好Toth.
就摘下Lakak.
果子來吃.
又給他在一起的丈夫.
他就吃你.
你看到這裡.
每一個字.
的keyword都有重複的.
Elohim.
你們如神.
Kielohim.
神的種子.
Bnei HaElohim.
於是女人見到Ra.
就好像神的兒子見到Ra.
那樹上的果子是好的.
Toth.
希伯來文是Toth.
就是good的意思.
神的兒子.
漢見女子.
中文翻譯成美麗.
但原文就是Toth.
其實.
奇怪的.
如果形容女子美麗.
有很多好一點的字眼可以形容.
譬如你記得阿布拉罕.
阿布蘭到埃及的時候.
12章創世記.

$^{641}$就說他太太很美.
殺萊很美.
就告訴他.
因為你這麼美.
我就怕了.
別人拿了你.
然後將我殺了.
你很美.
那是Yafa.
是原文.
說他外面的外貌很美麗.
漂亮.
就用Yafa.
Yafo等等.
同組的字.
美麗.
Beautiful.
但這裡.
有一點.
你看到作者.
特意選擇Toth這個字.
為什麼他會這樣做呢.
似乎是回應.
三章五至六節.
人類的第一次叛變.
然後神的兒子又再一次叛變.
這是叛變的高峰.
然後洪水就來到.
如果你這樣看的時候.
似乎這裡記載了一個敘述.
如果你接受第三個說法.
這些是王族.
神的兒子是一些君王.
你會看到作者抽這個故事出來.
洪水之前發生.
可以說是一個樣本.
末世的那一代的人.
好像這些神的兒子.
他們就好像第一次.
人類叛變的時候所做的事.

$^{681}$看見.
好.
然後摘下來.
或者取編.
這些關鍵字全部都出現了.
所以告訴我們.
上帝公義.
末世臨到的時候.
末日臨到的時候.
發生了什麼事.
就是最大的叛變出現.
上帝不得不按照他的公義.
來去施行審判.
唯獨是羅亞蒙印得救.
所以就跟這個審判一起來.
所以為什麼在新約裡面.
很多時候回應或者呼應.
這一段聖經的時候.
洪水到來.
當時的人就吃喝嫁娶.
或者當時有很多的叛變.
所以如果你接受第三個說法.
第一代的讀者應該收到的訊息.
就是這個.
就是當時神的眾子.
這些皇室墮落得非常厲害.
叛變得非常厲害.
然後招致上帝的審判.
那些巨人啊.
天使啊那些是怎麼來的呢.
有很多都是在兩約中間.
特別是有一本書叫做《以樂一書》.
那個是在正典以外的.
兩約中間的猶太人的文獻.
那些是一些天啟文獻.
好像啟示錄那種的寫法.
就說到他們在傳適這一段聖經.
後期再傳適這段聖經的時候.
就是這些墮落的分子.
就是天使.

$^{721}$在天上墮落.
然後掉下去繁家人等等.
他們就加了這些傳適上去.
所以你看到解經的時候.
如果你看第六章.
這是第一層.
然後到兩個中間.
就到猶太人.
天啟文獻是怎麼解釋的呢.
又再加另一層.
新約去到說.
或者呼應這一段的聖經.
洪水到來或者審判.
又會說到天使墮落等等.
比如彼得後書或者猶太書等等.
說到天使墮落.
末世來到的時候.
這些天使會墮落等等.
會搞亂會腐敗整個人類等等.
他們就建基於在兩個中間.
猶太人的一些傳適.
所以在這裡我們可以算是一個.
exercise 作業.
讓我們知道怎麼去解經.
有不同的方面.
可以看近東文獻 上下文.
甚至看兩個中間有什麼發展.
然後新約怎麼去建基於這些發展.
然後再去到傳適這些經文.
我們處理了這些keyword.
和八節的經文內容.
我們做一點點的總結.
三至六章的墮落是什麼呢.
作者很明顯是揀選了三個故事.
三個敘述.
然後這三個敘述就代表著.
人類罪惡的泛濫.
罪惡的擴散.
所以第一個故事.
你記得是第三章.

$^{761}$就是亞當夏娃.
那人與女人.
這個是archetypal story.
就是告訴讀者.
每一個人都是在個人層面上.
我們去到叛逆上帝.
將罪作一個定義.
就是以神為神.
我們想做神.
然後到第四章.
就是講到一個家庭.
罪惡開始擴散.
由個人去到一個家庭.
就是該人殺了阿伯.
然後他的後代去到第七代的拉密.
殺氣騰騰.
暴力繼續增加.
然後第三個故事.
就是我們今天看的.
就是社會國家幫國的層面.
神的意志.
按照第三個的解釋.
這些皇族.
作者拿出來.
是做一個代表性.
就好像告訴我們.
當時的人類.
冰山的一國.
就好像這些神的意志.
任意挑選.
看為好的.
不以上帝為上帝.
自己去看自己.
有如神一樣.
神的意志.
所以你見到.
由開始上帝所創造的.
order的規律.
就變成了越來越多.
這些disorder.

$^{801}$這些反規律.
所以去到最後的時候.
就到六章第九節.
就是下一段.
大段洪水的故事.
上帝就把大水混沌.
拿回來.
就把世界變成.
好像第一章第二節.
大水的混沌.
water的chaos.
就把non-order.
帶回來.
去到審判.
當時的不敬虔的一代.
所以你見到一到第六章.
之前我們看過的分段.
我的分段就是.
第三章第一個罪的故事.
然後第四章第二個罪的.
再擴散反濫的故事.
然後到第六章.
罪惡滿盈.
神的意志.
到最後就是洪水瀝淋的時候.
你見到這裡有一個進度.
一個progression.
罪惡越來越大.
越來越擴散的時候.
上帝會施行審判.
告訴我們.
我們作為義者的後代.
如果我們是蒙恩的一班人.
相信上帝.
我們在一個這麼黑暗的時代.
一個這麼混亂的時代.
我們堅信上帝.
公義的審判一定會臨到.
並且他的救贖會臨到.
在六章裡面.

$^{841}$特別來到這個地方.
你見到有一個pattern.
去到所謂末日的時候.
洪水公義的審判臨到之前.
就有這些crisis出現.
世代會越來越昏暗.
罪惡會越來越滿盈.
但是蒙恩得救的人.
相信上帝的人.
就會得到上帝最後的救贖.
這個是一個eschatology.
一個末世論的pattern.
其實已經set了這個pattern.
for其他聖經以後發生的事.
特別到啟示錄也是一樣.
有一個這樣的終末論.
到終末來到之前.
世界一定會是罪惡滿盈.
但同時神是去到sustain住.
去到維持住他自己的子民.
但願我們每一個.
都是與上帝.
建立與上帝之間的約的一班人.
我們那班蒙恩得救的人.
繼續去信靠上帝最後的救贖.
一定會來到.
如果你喜歡我們的youtube頻道.
就記得like comment和subscribe.
拜拜.
\newpage



\section{}
\label{sec:Gv0nluM2dc8}
\textbf{「流散/流放」神學初探 (下):華人教會實況(2022-4-16)}
\newline
\newline
連結: \href{https://youtube.com/watch?v=Gv0nluM2dc8}{\texttt{https://youtube.com/watch?v=Gv0nluM2dc8}} ~~~~ 語音日期: 2022-04-16
\newline
\newline
\hyperref[sec:jwup3nNhxBw]{\small{< < < PREV SERMON < < <}}
~
\hyperref[sec:index_chronic]{\small{[返順時目]}}
~
\hyperref[sec:index_scriptual]{\small{[返順卷目]}}
~
\hyperref[sec:H5eDGDUsh5o]{\small{> > > NEXT SERMON > > >}}
\newline
\newline
$^{1}$各位大家好,我們有第二集.
上次就說到流放神學.
流散神學.
有很多精彩的討論.
我們不再詳細介紹.
我有幾位嘉賓.
我們有胡志偉牧師在英國.
我們有王啟泰牧師在澳洲.
他現在在墨爾本.
雪梨兩邊走.
芝加哥的Mark Sir 陳崇經牧師.
我們繼續討論.
今天的焦點想補充一點.
和落地一點.
上次的結論.
如果我說錯了大家可以correct我.
就是我們探討到流散神學.
其實真的很闊.
也有很多角度可以看.
也有很多的explorer.
所以我想大家會不會有些補充.
上次我們從聖經神學那方面去討論.
會不會在歷史神學方面.
或者在教會歷史上有些參照.
可能有某些族群或基督徒.
他們都經歷這些流散流放.
以致我們再豐富一點.
從這個角度去看這個題目.
或者先由胡牧師開始.
我還想先讓黃牧師講.
沒所謂不用這麼禮讓.
教會歷史上我相信有很多不同的例子.
在昆士坦丁還沒承認基督教作為國教之前.
當時已經有不少的主教.
因為政治的迫迫.
他們要逃亡.
這個我們在歷史上可以看到.
甚至他們有些選擇走.
有些選擇不走.
都有些書.

$^{41}$特別是在教父著作裡面.
我們可以看到.
後來當然到了中世紀的時候.
這種流放.
你也可以說修道院.
整個修道運動.
都是代表一群追求靈性的基督徒.
他們不滿當時的所謂建制或者教會的腐化.
他們就選擇將自己放逐.
組成群體在歐洲不同的地方裡面.
來建立修道院.
過一種群居的生活.
到了宗教改革的時候更加明顯.
其中一個最典型.
大家都熟悉的例子就是加爾文.
加爾文他逃難跑到日內瓦.
在那裡開始宗教改革.
甚至當時如果看回英國的時候.
因為我在英國.
我就會多一點了解.
當時有不少的所謂不服從國教的領袖.
當時避難就跑去荷蘭.
因為荷蘭很近英國.
所以不少人跑去那裡.
甚至翻譯聖經很著名的就是William Tyndale.
他將聖經翻譯成英文.
後來就殉道.
他一段時間都在荷蘭這個地方裡面.
過著一種放逐或者流放的生活.
後來再發展到5月3日.
一群清教徒坐船去到美洲新大陸.
去建立他們理想的教會生活.
所以在歷史裡面這種的逝世.
無論是被迫也好.
自己選擇也好.
是有不少的例子.
下面那些我就交給黃牧師補充了.
其實剛才黃牧師提到教父.
其實有一個著名的人物居普良.
其實他在大壁包裡.

$^{81}$第一次的時候他是逃跑的.
第二次他就留下.
最後就殉道.
其實他逃跑的原因很簡單.
第一他目標太大.
容易引起更大的對加泰基教會的迫害.
第二他覺得他離開了.
是更能夠幫助教會.
當然這與否不同意另一回事.
這也可能是教育離開某些地方的理由.
但是我覺得我們不要猜測動機.
確實在歷史上有這樣的情況.
剛才黃牧提到改革運動的時候.
其實有一個很大的群體.
就是信洗派.
安娜·貝布提斯.
他們在某個程度上從自己原居地被迫離開.
無論是日內瓦也好.
或者是慈雲里的地方也好.
他們都離開.
後來有些人去到美國.
北美洲.
甚至乎在近代的十八,十九世紀的時間.
在北歐的國家教會裡面.
所謂的non-conformist的離開.
又去到北美洲.
開始了有他們自己的教會.
差不多可以說是今天所說的移民教會.
而這些移民教會後來.
譬如其中一個就是我們叫做中國的基督教播道會.
就是其實是由美丹瑞過來的.
就是一個這樣的群體.
其實就像剛才在講的.
在人類歷史裡面.
其實遷徙可能是一個常態.
多於是一個孤立.
從歷史裡面我們可以看到.
唯一一個我比較有思想.
不是跟教會有關.
但是有另一個族群.

$^{121}$就是客家.
客家就是因為不斷地遷徙.
但很奇怪的.
跟我們剛才所描述的.
不是很相似的.
客家就是去到那裡就是那裡.
他就停留在那裡.
但他仍然維持著客家族群的特色.
如果你說我們剛才講.
提到回歸的問題.
有exile就有回歸.
這是一個很自然的思考.
究竟是一個怎樣的情況呢?.
我覺得流散神學也好.
流放神學也好.
這個始終都是我很關注的.
一個很值得思考的問題.
就算我們說基督徒作為一個所謂流散者.
其實我們有一個所謂的final destination.
就是回歸到上帝那裡.
我不用一個傳統的說法.
回到天教或者回到天堂.
回到上帝那裡.
本身也有一個回歸.
這個會不會也是我們可以探索的.
將來有機會.
我停留在這裡,不要插得太開.
也不會的.
Mark Sir有沒有什麼補充或者你的看法?.
我想我跟大家分享過.
我認為聖經裡面給我們一種theology of exile.
有不同的書卷,申命記等等.
哀歌等等.
都是對這個流放,流亡,流散的現象.
對當時神的子民有很多不同的衝擊.
以及有consolidation.
剛才黃木說的那個return那裡.
我也認為有這個theology of return.
因為你看到第二聖殿時期的時候.
不同的猶太族群.

$^{161}$其實他們都對這個return有不同的理解.
如何處理這個exile有一個不同的理解.
譬如四海孤軍的群體有不同的.
有些散去另外一個地方.
又未必回到耶路撒冷.
所以對於大家在第二聖殿時期的時候.
都已經看到到新約.
都有這種傳統.
什麼叫回歸上帝.
我很贊成黃木所說.
回歸上帝.
他們不一定return to the land.
但是到了第一世紀的時候醞釀了.
耶穌出來的時候.
他就represent另外一種.
不是傳統的聖殿式的回歸.
嗯.
大家都補充了多一些關於歷史上.
其實我想由新約的引申.
如果新約根本保羅都到處走.
約翰都是被放逐的.
或者再說回舊約.
其實一路延伸下去.
教會歷史.
很多族群有這樣的體會.
我們說回華人.
華人其實都是.
我想除了猶太人之外.
都是一個比較散居.
或者流散的民族.
我們本身已經是.
只不過可能.
我不知道我這樣read對不對.
華人到那裡就落地生根.
希望安居樂業.
有一點安定的感覺.
我也看到會不會.
從這個起點.
我們在海外的華人.
有時候在華人的教會.

$^{201}$其實都挺安定的.
傾向到了那個地方就安定了.
所以我不敢說以前沒有人提出.
去想這個流散.
那些流放的問題.
但是華人教會.
我也知道有些神學.
這個導向.
但是似乎不是很講.
因為我們追求安定多於流散.
可不可以這樣說呢.
你們在華人來說的背景.
會不會有些insight呢.
我先說吧.
因為我算老華僑.
我是1990年已經到了美國.
然後家裡移民.
那時候就沒得選擇.
我自己觀察在過去.
我在華人教會30多年.
自己在不同的教會目過會.
只是說美國這個地方.
每個處境都不同.
美國很多時候都是70-80年代.
開始有留學生.
來北美留學的時候.
就組織茶經班.
然後茶經班慢慢多了人.
成熟了.
他就建立教會.
很多美國的教會.
都是從這種傳統走出來.
就像礦牧你說的.
有很多的移民.
來到北美.
就是美國.
都是那種.
我們上一集說的.
上流式的移居.
來到這裡.

$^{241}$我們讀到書.
然後有條件留下.
比其他人好一些.
很多時候.
這種取態.
這種心態就帶進了教會.
所以在北美的華人教會.
所以每個人都說.
居然在北美這個地方.
華人教會是最保守的.
因為有一個這樣的現象.
就是他們將以往.
譬如香港有很多.
他們以往教會.
所熟悉的東西.
他們就帶進來了.
因為北美的移民心態就是這樣.
他們有很多同鄉會.
譬如我們芝加哥的中國城.
有很多同鄉會.
姓什麼都有同鄉會.
同鄉會其實做什麼呢?.
打牌,自己賭錢.
然後就維持自己的.
同聲同氣的感覺.
那種文化的身份.
所以教會.
都變成.
經常都是這樣的.
所以為什麼教會會很容易保守.
不會創新.
也不會重新思考.
因為來到的時候.
我們已經安定了.
我們就想有個泡沫.
我覺得我一直以來在華人教會觀察.
我覺得天國公民的身份.
是非常弱的.
特別是我這一代和上一代.
因為來到只是.

$^{281}$想安定.
連解經都想安定.
解經的方法都是想像以前.
牧師那樣教那一套.
譬如說肥辣比輸.
就要上上起落.
但他沒有留意保羅說的是天國的.
Citizenship.
我們是天國的公民.
所以變成北美華人教會特別在近代.
遇見很多世界上.
或者本地上的一些衝擊.
我們啞口無言.
我們完全不知道怎麼回應.
唯一的方法就是我們躲起來.
我們處理我們泡沫裡的事.
我們安定就好了.
我就說到這裡.
我覺得這是美國華人教會.
一個很大的弊病.
其實澳洲.
也不是好得多.
通常那些人說.
我還沒來澳洲.
我來了五年多.
那些人說澳洲.
差20年左右.
神學和教會的情況.
公平一點說.
我來到就覺得.
友好的地方.
他保留了一些以前.
我比較香港.
香港教會的一些.
很熱心.
有時間.
其中一個.
又願意去事奉.
但你說保守那方面.
都挺保守的.

$^{321}$我想回應一下.
很簡單說一句.
都有那些留學生.
以往留下來的.
北美華人教會有好處.
但我想說.
時代穩定的時候就沒問題.
時代穩定沒問題.
每樣東西都覺得很need.
但到最近.
好幾年.
衝擊大的時候.
就完全沒負.
完全沒回應.
明白,另外兩位呢?.
我去澳洲一段時間.
其實我93年來的.
不過我又回香港事奉.
所以.
澳洲有經驗.
但不可能代表澳洲的說法.
不過我看.
都有一點像Mark Sir.
所說的情況.
當然.
澳洲橫教會的.
發展比較興旺.
開始在80年代90年代.
的時候.
香港97問題引發.
所以這一批.
背景的移民.
都是應該是professional.
有錢一點的人.
在那時候過來.
真的來到就是想安定.
這是其中一個因素.
令澳洲的.
橫教會在某方面來說.
他不是像我們期望的.

$^{361}$那麼dynamic.
但我覺得有其他因素.
contribute to這個結果.
包括.
其實是.
第一就是.
始終語言是一個文化問題.
語言未必是最大問題.
反而文化的大問題.
是一個文化的問題.
所以我認為.
如果要找回自己的族群.
找回自己的族群.
他一定會找回自己的族群.
因為政治問題.
人們都是有秩序地移民.
但是最近的流放.
簡單來說就是逃難.
英國是重災區.
聽聞也有很大的影響.
反而澳洲這邊.
芝加哥那邊.
沒有這麼快.
始終疫情也在進來.
但是沒有英國那麼明顯.
這一群逃難的.
甚至還沒準備就要離開.
其實那種心態和以往.
似乎從未見過這樣的情況.
胡牧師你那邊.
雖然你是最新的移民.
我也是其中一份子.
我會這樣看.
我在6月3月中的時候.
流放到英國Bristol.
我會看到越來越多人.
特別是.
因為BNO visa的緣故.
英國放寬了.
如果不是的話.

$^{401}$我相信很多人都是選擇加拿大.
選擇澳洲.
也不會選擇英國.
所以因為這個政治因素.
令到英國門戶大開.
大量香港人就會走進來英國.
而我想.
這幾年來的.
肯定和我剛才所論述的.
那些所謂上流的經濟移民.
是有所不同.
他們不是來自內地.
或者來自香港.
他們只是為了有一個經濟上.
能夠有一個更加好的保障.
為了子女的讀書.
為了他們的生活等等.
所以他們這種是向上流.
基本上我們政治一定是保守的.
而早期那些苦工.
他們因為要謀生.
被賣豬仔.
去到美國.
去到一些外地裡面.
然後匯錢回到內地.
來養妻活兒.
是一種經濟的因素.
但現在出去的那些.
當然有部分都是向上.
為了子女.
更加享受英國的教育.
為了其他等等.
當然有些所謂的.
意識形態仍然是很建制的.
或者稱為.
都是藍色的那方面.
都會移民的.
但我會看大部分.
他是因為這種.
對現在的特區政府的不信任.

$^{441}$特別是對中國那種.
以往的中國夢.
已經破碎了.
他們就選擇.
在英國裡面生活.
所以這裡就帶來.
一個很大的衝擊.
所以現在我和一班.
就我所認識.
都有不少特別是在Bistro.
或者附近一些小的城市裡面.
比較多現在的選擇.
他們是回當地的西人教會的.
不是回華人教會的.
大家要明白.
我有些文章裡面都有論述.
因為華人教會的文化.
都仍然是停留.
舉一個最簡單的例子.
都會認為這是烏克蘭的衝突事件.
雙方都有責任.
所以我們祈禱的時候.
都是祈求兩方面盡快.
放下武器.
大家能夠和平過日子.
但我是覺得.
他們的統文是有問題的.
他們甚至不可以說.
這是普京.
這是惡人.
這是不義戰爭.
他們連最基本的.
都不敢說.
因為怕得罪.
有不少來自中國背景的.
信徒在裡面.
所以我看到.
接下來這種衝突.
是必然會白熱化的.
以往就像Mark Sir所說.

$^{481}$太平盛世的時候.
就不會有衝突.
但在亂世的時候.
對疫苗.
對防疫的政策.
很多問題.
已經可以說是政治化的.
所以接下來.
我也會擔心.
隨著這種.
所謂對中國的恐懼.
甚至那種情緒的反彈.
以前就是.
對華人非常友善.
但未來日子裡面.
可能是對.
包括我們這些流放在外的.
他都不懂分你是香港人.
還是大陸人.
還是甚麼人.
我們都一定要承受.
可能將來各地政府會出現的排華潮.
我們都會要承受.
那個後果.
這個一定的.
就因為在澳洲的華人.
都會覺得習近平.
是更加可信過澳洲總理.
在調查裡面反映.
大家可以看到.
這個就因為有不少.
去到海外的.
無論是中國人也好,香港人也好.
他吸收的資訊都仍然是.
大陸官方.
或者香港的TVB.
或者大陸的中央電視台.
所以你可以理解.
他們就已經.
被洗腦,長期接受.

$^{521}$所以普京.
來到統一烏克蘭.
是名正言順的.
甚至將來.
中國統一台灣.
都是合理的.
他們就是接受大一統的政治論述.
所以這個就要.
面對那個張力和矛盾.
我想張力一定有的.
一直都有.
但胡牧師說的.
是嚴重了.
我聽聞.
美國那邊.
是嚴重了排華.
特別是.
最近黑命貴的問題.
種族主義的問題.
當然亞裔.
受不同的壓力.
不同的衝擊.
我也參與過不同的示威.
或者遊行.
不同種族的問題.
你也看到.
好像胡牧剛才所說.
華人教會.
是有少許綁手綁腳.
不知道怎樣做.
至少語言上.
已經沒有這種語言.
怎樣去面對這些社會問題.
或者社會政治上的問題.
很多時候.
都會走回去.
我們希望.
大家不要吵架就算了.
我們繼續安定下去.
所以他講到烏克蘭.

$^{561}$俄羅斯的問題.
的時候都是一樣.
他最終的目的都是希望.
不要打.
不要升.
我們每件事都安定.
將移民的心態.
上流移民的心態.
帶進來.
去詮釋我們的信仰.
怎樣去回應.
黃牧師.
你的觀察.
在澳洲都是這樣嗎.
澳洲華人教會.
當然剛才.
兩位胡牧和Mark.
的分析.
是很合理的.
我有.
另一個角度去看.
其實.
怎樣說呢.
澳洲華人教會.
我以前都喜歡用一個說法.
叫華僑教會.
那就是.
你在華僑.
的社群裡面.
因為他覺得自己不是這裡的人.
所以他的關懷是有限的.
他們的少數弱勢群體.
所以他覺得自己天賦地.
是沒有足夠的聲音.
讓他們可以.
表達出來的.
所以.
大家可能覺得我比較同情.
我看到他們的困境.
當然.

$^{601}$我也同意.
信仰內涵.
是千年如一日.
沒有太大變化.
也許.
現在黃藍.
所謂的衝擊.
可能是一個喚醒.
澳洲華人教會對這方面的看法.
譬如我在教學裡面.
其中有一個機會.
和學生聊天的時間.
我都.
大家.
悉尼的教會.
他們就說.
等候著一批香港移民.
所以我們要預備好.
其實問題是.
你不要想他是一個數字.
也不要想他是一個勝利軍.
你應該想的是.
他們會帶著傷痕.
如果教會要去侮辱他們.
應該要考慮這件事.
首先.
當然我們剛才說的.
很高調的流放神學.
可能要再等一等.
我個人覺得.
對那班.
新到的人來說.
這是一個.
他們的傷痕要等一等.
等得到治療之後.
他們才有機會.
面對一些新的挑戰.
當然和他們一起去談流放的神學.
本身都可能是一個治療的過程.
我暫時的.

$^{641}$即時的思考.
都在這裡.
不過我覺得黃牧師的觀察.
都是對的.
如果再sharp一點.
我不知道這樣演繹對不對.
其實可能.
久一點的移民.
他們的身份認同.
或者他們.
可能上流的那一類.
多的時候.
和現在新來的那一班.
無論是英國重災區.
美加.
澳洲這邊.
是有分別的.
主要我們現在新來的那一班.
譬如我們教會都在想.
怎樣去接待這班.
我都叫新一點.
這幾年才移民.
和一直有緊密的接觸.
和香港的弟兄姊妹.
教會.
我們都會提醒.
就像黃牧師這樣說.
那班很多傷痕.
那班經歷很多不同的想法.
甚至他們還未預備好.
因為帶著小朋友.
快點來了.
我就發覺是有差異的.
我們內移民那班弟兄姊妹.
就不是.
馬上會知道.
可能理性上.
澳洲經常會回去香港.
因為近.
但他們都會知道.

$^{681}$這些情況.
但他們是否真的準備好.
我就發覺.
相信大家可能有這樣的經驗.
真的.
一起聚會.
小組.
差異就來了.
我可能做一個中間人.
我不是說太久.
也不是說太新.
那就可以有些經驗.
也可能獨立地.
有些組去幫助一些較為新的.
我又新.
他們更新.
我經歷過的這些事.
又可以幫他們.
我想要面對流散的一些.
或者流放的一些信徒.
有特別的需要.
我想可能我們教會.
要多想一些.
是給這班人的.
可能當宣教工場.
這班是特別的群體.
這樣做.
其實會不會英國.
是多了這些.
思考呢.
你們的衝擊來得很快.
我相信一定.
剛才阿貢所說.
有不少是帶著傷痕.
仍然有很多未處理的情緒.
在他裡面.
在英國這個地方裡面.
所以我想這個就變成.
而這裡的華人教會又未必.
因為.

$^{721}$所以現在也有一些.
群體成立了.
比方在Bistro裡面也有一個叫流氓聚會.
就是有一位傳道人.
他就帶著在weekdays裡面舉行.
所以我想.
未來也會建構.
多一些的平台.
就變成不是用一個所謂.
以往就是華人教會的形式.
現在也是一個.
西人的教會裡面幫他們有一個.
廣東話的崇拜.
混合的模式來做.
所以我會看.
會有很多不同時空的需要.
去具體幫.
這些香港人.
去處理他們內心裡面.
的一些.
要處理的一些情緒.
時間也差不多了.
Mark Sir.
有沒有補充呢.
我想.
同意剛才所說.
特別在一個.
比較.
絕對是動盪過以往的.
一種這樣的時代.
教會有沒有那種的resilience.
有沒有那種的彈性.
還是我們只是box in了.
以往那種的方法.
剛才我扣木你講的.
情況也是.
當然我們美國是比較少.
美國很難進來.
但是其他地方.
我相信也需要有很多.

$^{761}$不同的彈性的處理.
特別有更多.
同理心.
我們移民了以後的人.
我的勸告就是.
不要用我們以往那種的.
心態去理解現在來的人.
是完全不一樣的一種人.
同意.
我同意Mark Sir所說.
其實我也觀察到.
的問題就是.
他們可能會想回香港.
無論探親.
購物旅遊.
但是我懷疑現在出來的人.
是斷了路的.
他自己斷了路.
不一定是別人斷了路.
應該是完全不同的.
反而這裡有一個契機.
他們真的會原居地.
作為安身立命的地方.
以往的屬命官.
教會官,使命官.
可能是更加到辱和更加落地.
所以這個也是.
我現在在英國做的.
就是鼓勵他們融入在當地的.
社會裡面.
其中一個途徑.
就是我鼓勵他們.
我兒子,媳婦,孫女.
都是回當地的西人教會.
我自己也是.
我身邊有不少以前回我教會的弟兄姊妹.
她們全部都落腳在不同的.
西人教會裡面.
這個就是我現在在這裡.
不少香港人.

$^{801}$在不同的當地的教會裡面.
成為他們的一份子.
以致他們可以將這種精神.
帶入他們的教會裡面.
將來不成為當地的教會領袖.
甚至可以在教會裡面.
有一把香港人的聲音.
這個就是我嘗試要做的.
多謝大家.
各位的觀察.
分享.
我想小總結就是.
不可以鐵板一塊.
用以前的方法.
或者教會增長.
移民潮來了,我們教會又會多些人.
其實是一個.
不是兩變.
兩和質都變.
我想可以像胡博士這樣說.
應該是有不同的方式.
無論是西人教會.
橫教會,橫教會自己也.
在心態或者幫助那些人的需要.
也要更加開放.
去改變自己的看法.
去聆聽.
像Mark Sir這樣說.
甚至很多.
很hybrid.
他過來,然後繼續在香港的聚會.
在Zoom,我想有很多可能性.
去想想.
怎樣去牧養一班.
我們所謂流散.
流放的一班弟兄姊妹.
其實我想這個是有效的.
繼續我們還有很多東西.
可以探索和去體驗.
期望將來.

$^{841}$我們跨區.
繼續有多些.
討論經驗和分享.
時間差不多了.
有機會再做節目.
多謝大家.
拜拜.
\newpage



\section{}
\label{sec:H5eDGDUsh5o}
\textbf{「流散/流放」神學初探 (上):從聖經神學定義(2022-4-15)}
\newline
\newline
連結: \href{https://youtube.com/watch?v=H5eDGDUsh5o}{\texttt{https://youtube.com/watch?v=H5eDGDUsh5o}} ~~~~ 語音日期: 2022-04-17
\newline
\newline
\hyperref[sec:Gv0nluM2dc8]{\small{< < < PREV SERMON < < <}}
~
\hyperref[sec:index_chronic]{\small{[返順時目]}}
~
\hyperref[sec:index_scriptual]{\small{[返順卷目]}}
~
\hyperref[sec:17_oWOqNjfw]{\small{> > > NEXT SERMON > > >}}
\newline
\newline
$^{1}$各位好,歡迎各位來到我們廣風之谷的另一個論壇.
今天有幾位嘉賓加入我們.
除了我在芝加哥之外,還有三位嘉賓參與今天的討論.
我們先讓大家介紹一下自己,先讓康牧先介紹一下自己.
大家好,我是康惠志牧師.
我是廣風之谷的一個閒人來的.
我在雪梨牧會,我比較喜歡負擔做多一點傳呼音的工作.
好,那胡牧先來介紹一下自己.
各位你好,我是胡志偉牧師,我是廣風之谷的顧問.
因為下水的事情對一個退休牧師來說太多了.
所以我自由身,做我喜歡做的事情,寫文章,討論等等.
現在在英國的Bistro.
英國的Bistro,歡迎胡牧師加入我們.
另外,Pastor Tony,黃牧師.
大家好,我是黃啟泰牧師.
還沒有到胡牧師的退休的時候,快要了.
所以有時候做事都是有些責任.
也要顧全服務機構的一些關聯.
這些是身不由己的情況.
我叫做廣風之谷的顧問.
也想參與多一點,不過這段時間的教學和教務行政上都比較忙碌.
只能夠暫時做一個類似閒角的角色.
比礦牧更加閒.
神樂縣在悉尼,但是我的家在Melbourne.
不要問我如何處理這個問題.
暫時是這樣.
(黃) 安身.
黃牧師你在Melbourne多一點嗎?.
其實一半一半.
今天很高興有澳洲兩位牧師,Richard和Pastor Tony.
還有UK胡牧師.
我在芝加哥美國,真的夠分散了.
差不多整個地球.
我們今天想講的題目就是「散流放」這方面.
我相信在香港背景的信徒教會來說是一個很切身的問題.
大家知道香港的問題.
香港也有很多的信徒要掙扎流於去的問題.
也有很多已經離開了的.
我相信我們四位當中有很多的人.
包括我們自己,都是流散的一部分.

$^{41}$我們看看可不可以在這方面認真討論一下.
討論一下流散或者exile這個問題.
所以我們想在這裡討論一下這個問題.
最近有一間香港的神童院.
因為這個題目就開了一個碩士課程.
我相信這個舉動或者做法.
其實反映了這個題目是非常熱門的.
我們大家也需要很深思去思考一下.
我不知道三位有什麼的參考.
在這方面流散或者流放等等.
胡牧師先來,好嗎?.
我還沒有離開香港之前.
大概在2020年9月寫了一篇文章叫做流放神學的探討.
大家可以在網上找到這篇文章.
那時候我還沒有離開香港.
我那時候也已經思想.
因為隨著香港已經變得越來越陌生.
有點像聖殿被毀的那種情況.
所以很多人就四散.
在四散的時候究竟是怎麼樣?.
如果有留心看我當時還沒有離開香港所寫的文章.
其實我那個論述基本上就是看不離開香港.
其實那個教會群體已經是落在一種流放的狀況裡面.
有些人錯誤的理解.
就是以為走了,不在香港.
這些才是流散的群體.
所以我會嘗試.
因為我經常思考這個問題.
就會看流散和流放是不同的.
如果你稍微找一本中文最簡單的來看.
流放就是中國古代對一些官.
大家很熟悉的其中一個就是蘇東坡.
他當時和王安石在朝廷上有一些大家在政治上的爭論.
所以蘇東坡被貶,他就流放到中國的南方一些地方.
所以流放就是以前對一些人的一種刑罰.
就放你,放逐你去一個偏遠的地方.
在聖經裡面也有這種描述.
特別是你看一些中文的譯本,比方說馮仗.
馮仗在聖經裡面也有這麼說.
特別是《小林傳》第七章43節的翻譯.

$^{81}$他就說漢王流放,趕去巴比倫.
就是當時這群以色列人.
他們拜偶像的時候.
上帝對他們的審判就是將他們流放,趕到巴比倫.
所以我去看的時候.
然後在另一個大家很熟悉的經文.
就是在《遺德前書》第一章第一節.
他就說彼得耶穌基督的使徒.
自流散在本都加拉泰,卡伯多奇亞,亞細亞.
和比聚尼的過客.
他這裡就用兩個字,流散.
過客就是我們今天所說的那個所謂流放者.
這個流放者其實.
當然在中文和學本的翻譯裡面就是寄居者.
所以我就會看這個exile這個字在英文裡面.
他的那個內容是豐富過流散.
流散者很多時候就是指向一個人口的遷移.
就是people movement.
你由一個族群或者一班人.
由一個地方去到另一個地方.
當然你可以這麼說.
其實在80年代的時候.
王永信羅曼華年代.
說的這個華福運動.
其實就是在說海外華人那種流散的現象.
怎樣將這些流散的海外華人基督徒能夠凝聚.
以至他們仍然繼續有這種所謂宣教的使命.
但今天我們如果看回現在的情景的時候.
我會嘗試去看有三種類別的.
我會看第一個類別.
就是那些叫做流汗的苦工.
這個就是第一類的類別.
如果我們看回華人的移民歷史.
第二類的我會稱為上流的移民.
這些上流的移民會包括內地和香港的.
他是尋求經濟或者子女的教育.
於是就去了加拿大,澳洲或者英國其他地方.
第三類的我會稱為流放的群體或者流放的市民.
就是因為政治的緣故.
因為對他所處身的主流文化有一種強烈的不滿.

$^{121}$於是他們就要不情願地離開他原來的原居地.
我就會看2014,2019年社會運動之後離開香港的那一群.
就和前面那兩批的流散者是完全不同的.
是和那些流汗的苦工和那些上流的經濟移民.
是完全不同的心境.
所以當我們如果用流散來聚焦的時候.
有時候會放錯了焦點.
所以當然你可以說流散是神學的.
但是我會看流放就會有一個更加豐富的想像力.
是讓我看到的.
所以當我寫這篇文章的時候.
其實我就說在香港的飆途群體.
他都仍然是和我們流放的心態.
這種流放的心態就是基本上我是對主流建制的.
無論政治也好,教會的體制.
我是有一種強烈的不滿.
我是不想繼續按你們那種方式來生活.
所以它就是一種自我放逐的那種群體生活.
所以其實他開學位課程是沒有問題的.
但是其實重點是錯了.
不是給那些海外流散的群體.
而是給在香港那些的飆途群體.
他們繼續怎樣實踐一種流放式的.
那種門徒生活.
其實我兩年前的文章就已經在說這個重點.
你可以看回我整篇文章就不是說走了會是怎樣.
而是說留下在這裡.
怎樣就不是再成為一種聖殿式的那種勁敗群體.
而是所謂聖賢式的那種學習,聖道的那班作業滿途的市民.
這個就是我基本的論述.
我就停在這裡.
這個讓我想起其實舊約裡面.
胡牧你也看過.
去聖殿化.
聖殿地位的時候對猶太人的那種衝擊.
有很多不同的群體出現了.
然後也有很多不同的神學的回應.
有些可能離開了.
有些可能留下.
我覺得你說那個留在香港.

$^{161}$但是仍然在自身上是一種流放的舉牌.
是.
黃牧師呢?.
我同意胡牧特別提出這個區分.
我覺得很有意思.
因為就算在中文本身.
譬如我們說流放.
又或者剛才胡牧也用了一個自我放逐.
放逐是一個意思.
就是在漢語裡面.
就很明顯是一個帶有.
無論是財在家諸在他身上的那種所謂懲罰.
又或者自己因為某些理念.
和那個主流有差異.
他離開這樣一個含義.
但是這個問題就接下來牽涉到一個.
我覺得比較有趣的就是.
是不是流放就一定是有一種自我意識開始呢?.
還是他也是一個可以是一種被動的呢?.
又或者流散我們剛才說胡牧提到最初那兩層的移民.
其實他們在某個意義上.
當然他們也是有所謂的經濟的考慮.
有些也很可能純粹是因為他的意識形態上的選擇.
我相信特別是講到80年代90年代那一批.
多多少少有這樣的背景.
其實他們怎麼classify(分類)呢?.
在胡牧這個區別裡面來考慮的時候.
是不是一定是有自我意識或者帶著一個使命.
又或者是一種被迫這樣才叫做流放呢?.
而那些早年因為經濟的因素.
這些可能都是一個被迫的原因去賣豬仔.
他們都是一種流放還是流散呢?.
我覺得這個是一個引起我思考的.
不是等於胡牧所講的不合理.
只不過我會覺得是這件事.
譬如聖經那方面其實有一個area(範疇).
譬如我們講到比特前書或者甚至是法國書.
但其實有另外一本書挺值得討論的.
就是《仕途行傳》.
《仕途行傳》那幫被打散的那幫.

$^{201}$到隨傳隨走隨存的.
我們又怎麼classify(分類)呢?.
如果用剛才所講的定義.
我覺得是可以繼續討論下去的.
因為我們面對的現象是太複雜.
可能是一個比較單一的定義.
未必可以全部涵蓋在一起.
至於剛才提到的神樂院的課程.
我覺得胡牧提到.
其實會不會比那些留在香港的更好呢?.
但其實海外走的人可能未必一定像剛才所描述的.
他們是帶著一定的目標或者使命感離開的.
可能他們都需要這種education(教育).
如果撇開這些內涵來講.
從課程的設計來說.
如果你說他成不成一個課程呢?.
如果我覺得有四分之一的科目是和這些有關的.
直接有關的.
其實他是在某個意義上可以成立的.
除了他的動機.
我不懷疑動機.
因為我們很難懷疑別人的動機.
但是在這個角度來講.
是不是完全不合理呢?.
不好意思.
和胡牧好像有一點不同的看法.
不過我覺得這個正好是一個很好的討論機會.
多謝黃牧師的補充.
我想不同看法是可以的.
這個會幫我們在思考的過程裡面.
對問題會更加豐富.
我也有思想.
其實有些時候是重疊的.
我會看史林卷路加的敘事.
是多一些從一個流散的角度.
所以你看到後來.
比起早期的時候.
剛剛去世的廉澤傑牧師.
那個論文就是在講紫赤宣教.
其實他就是在講一種流散式的宣教學.

$^{241}$所以如果用流散的角度.
就是在講神的子民.
他們是帶著使命去做一些上帝的工作.
但是流放其實是沒有的.
就好像剛剛黃牧師提出來的討論.
這個是對的.
根本他是因為形勢而離開的.
根本他不是有一種像亞伯拉罕那種.
現在叫他要離開他原來的地方.
離開巴比倫.
現在神要指示他去的那個地方.
他不是帶著那種使命感.
所以我會看流放.
但是流放和流散的不同就是流放者.
他是有意識.
他有一種意識就是我的身份.
是和其他人有所不同的.
於是當他有這種意識的時候.
他有他的論述.
比方說我看早期賣豬仔的苦工.
如果他講他的故事.
或者他的故事有這個論述的時候.
他建構一些黑人的奴隸.
他建構黑人的神學一樣.
是可以的.
這個就是第二代了.
那些上流的經濟移民.
如果他下一代又建構他們這種中產的.
然後去批判.
那是可以的.
我就會.
這些不同的向度.
只不過是讓我們對整個現象.
對那個context有一個更加豐富的了解.
我看不到有任何的衝突.
我想回應一下.
我們的起點是在講某神學院.
剛才黃牧師說.
會不會有四分一的內容.
是講到這個的內容流散呢?.

$^{281}$我想這件事出來的時候.
很多主內同道.
包括各位可能都有不同的看法.
很多人都會問一個問題.
就是可不可以搞到一個碩士課程.
我想我們不會說不好的.
不可以這樣說的.
但是要一個碩士課程.
而大部分的都是像普通讀神學的科.
只不過多一兩科是在講流散的.
或者是這些題目.
但是我想我們這次這樣談.
就是想搞清楚究竟流散.
或者是很好.
剛才黃牧師說的.
那個都是一個很好的introduction和summary.
就是在講什麼.
涉及這麼多東西.
其實如果我們要講這個題目.
我覺得不要留於.
又是一個留字.
就像古老一代的色經.
一個留字.
第一點流散.
第二點流浪.
第三點流放.
這樣的形式.
我覺得可能是太快.
好像有一個conclusion.
我們有東西要講.
流散神學.
其實是很闊的.
亦都層面很深.
很簡單.
譬如說你現在看烏克蘭.
俄羅斯.
其實那些人都是在流動的.
流散流動.
因為他可能要打仗.
可能要避難.

$^{321}$又數回在歷史這麼多的情況.
那些人流來流去.
我是問一個問題.
就是這樣是否值得去探討多一點.
還是我們因為香港人燒到近身.
所以我們才aware到.
其實根本已經有很多人去探討過這些題目.
我問究竟for what.
我們可不可以在歷史上拿一些體會.
很多人比我們的經歷更加深.
其實回答這個問題.
我會這樣說.
其實開科目不是一個問題.
問題就是有沒有這樣的人可以教.
這個才是最關鍵的問題.
因為其實比方說現在在歐美不同的國家裡面.
比方說migration studies是一個master.
現在已經有很多這方面的研究.
因為他不單單是講地理學.
或者文化學.
或者人類學.
而是一個跨科目的項目.
在英國.
我有時間當然要講錢.
你可以讀一個migration studies.
這裡其實是牽涉到很多不同的discipline.
是一個跨科目的項目.
有政治學,經濟學,人類學.
再加上我們所謂的神學在裡面的時候.
問題就是.
是否需要經過很多的study.
然後才可以教出來.
這個就是關鍵的問題.
比方說我思考這個題目也有幾年的時間.
我也會陸續寫一些文章.
我也不可以說我已經是這方面的expert.
我仍然是一個學習的階段.
可能我的看法會錯的.
可以修正的.
我就會這樣提出來.

$^{361}$你一定是要一些很具體,很solid的.
才可以去教那些學生.
給他們一些.
如果不是的話.
我想這個也是一個考量名目.
為了吸收多一些學生的市場技巧.
還有網上教.
我想阿胡說得很對.
特別是在現在這個migration上.
全球的migration.
不只是香港一部分.
那種複雜性和自身性的考慮.
特別是你想說的interdisciplinary.
不只是單單聖經,神學那麼簡單.
還有很多其他的科目我們需要去接觸.
如果你要對這方面有深入的討論和思考.
不可以單單單向.
只有一件事.
或者像阿鄺木那樣.
只搜尋一個exile或者流放.
我們就把聖經混在一起就算了.
不可以這樣做.
我聽聞有一個這樣的課程.
仔細的內容是怎樣也不好知道.
剛才胡木所提的也是值得留意的地方.
當然總是有一個開始.
究竟他怎樣做是另一個問題.
反而我覺得在現在的階段.
說真的流放或流散的問題是存在已久.
如果嚴謹來說.
你追上去當被抓出獄的時候.
你都已經是流放中.
如果岔開一點容許我.
我在這個問題想的那一下.
其實流散,流放是人類常態多於特殊的狀況.
我們不去區分流放或流散.
從遷徙的角度來說.
在中國的歷史以中國為例.
長期是這樣.
百多二百年就遷徙一次.

$^{401}$所以這是一個很值得我們探索的問題.
避開了香港的問題.
香港的基督徒的問題.
作為一個後設的思考.
既然我們遇到這樣的人生的處境.
其實我們從這裡如何建構一個神學論述.
其實本身已經是很值得做的一件事.
以至有神學院開始.
No matter 它的動機是怎樣.
作為一個開始點.
其實是引發一個討論.
我們至少現在這一次的討論.
都是因為這件事開始.
其實是一個正面的事情.
在某種意義上來說.
同意黃木之的分享.
所以我會看.
特別在現在全世界的形勢.
最近有一個字出現了.
你可以上網去看看.
叫做Keynote Politics.
就是說一種流動的政治.
以往我們看政治是國族之間的衝突.
但現在全世界有差不多一億的人口.
都是在一種所謂的移動的狀況裡面.
特別是以往那種區分地域的觀念.
在現在這個世代可能已經不合時宜.
所以每一個人都會有不同的觀念.
比方說我現在覺得英國生活得不好.
明年我會去加拿大.
或者我下一年會去澳洲.
都可以的.
當然要視乎你的經濟條件各方面.
所以這種的移動是帶來很大的社會的轉變.
所以當我們去看的時候.
特別是從聖經神學裡面去看的時候.
剛剛我所寫的那篇文章.
就是看創世記1-11章.
我就已經以亞當作為例子.
所以我就會看這種永久的移動.

$^{441}$就是我們那種本相.
因為人犯罪之後.
他就已經是長期落在一種放逐的狀況裡面.
這種放逐使他和人之間.
和自己和大地.
就是以往最傳統的那種說法.
就是那種疏離感和敵視.
那種不信任.
這是我們墮落的人性要面對的世界.
所以我會看今天在外面的.
我們身體是流放.
但其實有些人在精神裡面.
或者心理裡面他是流放.
所以你可以這樣說.
我們是雙重流放.
特別現在不是他原居地的.
中國人或者香港人.
我們是雙重流放的一種處境.
所以你從這個角度去看的時候.
是可以有很多很豐富的討論.
由第一個人亞當被驅逐.
離開伊甸園開始來作思考.
所以接下來我也是會想著.
會按照聖經不同的經文.
在新約裡面也有很多經文.
就是在說我們的身份本身就是一個寄居者.
一個流放者.
一個exile的identity.
所以彼得的論述.
亞伯的論述.
希伯來斯的論述.
史特克文的論述其實是相當豐富的.
這個就是我們說的身份.
我們其實反而.
如果從另一個角度去說.
教會的身份其實.
從另一個角度去說.
不是那麼強烈.
但我們也建構了教會的神學.
所以我最近也思考這個問題.

$^{481}$我會看.
一想到教會的時候.
就去了一個建制.
所以有一本講論教會論的書.
它的標題是什麼呢?.
就是在說我們是sojourners and exiles.
我們就是寄居者.
和這個.
異鄉人.
就是那種描述.
同意.
有兩個是胡木剛才說的.
引發我的思考.
一個就是.
其實.
說到流散流放.
這個如果純粹是一種文字遊戲.
容許我這個字眼.
其實我們過去在說.
所謂的寄居.
所謂的.
客女等等.
其實我們已經有類似的論述.
現在多了流散流放.
其實意義在哪裡.
我們可能要問的.
是否我們更清楚.
還是我們要換一個角度來看問題.
這是第一個引發我的思考.
第二個就是.
其實當我們說流放或流散.
自然會想到回歸的問題.
回歸是什麼呢?.
是一個地理性的考慮.
或者好像猶太人.
他們想要回歸到巴勒斯坦的地方.
還是有另外一個層次的意義.
就是說回歸不一定是地域性的.
也不一定是環境的.
是心態的.

$^{521}$甚至乎是所謂的universal value的回歸.
這些等等的問題.
我覺得是引發我繼續思考下去.
多謝黃牧提出來.
我會看就是.
剛才所說的.
特別我們去看的時候.
那個流散,流放.
從聖經的角度來說.
是有很豐富的.
我們可以在當中不斷消化.
當然在這麼短的時間裡.
我們很難.
我想英文就會容易一點.
只不過中文就將它弱化了.
比方說寄居.
其實如果你看回.
希臘文的英文裡面.
它的描述就會比較清楚很多.
它是一個所謂的exile的狀況.
還是一個寄居的狀況.
中文翻譯裡面有時候也會有不同.
它是外鄉人,過客.
你翻譯做過客.
翻譯做異鄉人.
翻譯做寄居.
意思已經不同了.
我會覺得用exile.
英文這個字.
反而是更加準確.
exile中文翻譯.
我會將它翻譯做流放.
這個字.
它所表達的意思.
會更加清晰.
所以我現在就嘗試.
用exile的眼光.
如何去理解聖經裡面的敘事.
我就會嘗試這樣去做.
其實我可以這樣說嗎?.

$^{561}$多謝兩位牧師.
一個很豐富的分析和分享.
或者引發.
我會覺得會不會是.
其實現在說的也是一個基督徒身份問題.
究竟.
為什麼要搞一個定義呢?.
無論是流散,流放.
或者回到剛才胡牧師說的.
教會群體的問題.
定義我們自己.
究竟我們是什麼狀態.
清晰一點.
不是文字遊戲.
這些的名稱.
或者描述.
幫我們去理解我們的身份.
以致我們的方向.
個人是如何定位.
無論我們在海外.
你怎樣才算是海外.
看你定位在哪裡.
然後群體.
是一群信徒走在一起.
我們一起.
都要embrace.
究竟我們是什麼身份.
是很複雜的.
那些人的心態.
我們在外國.
看見那些信徒.
其實他們的想法.
他們的身份.
認同.
簡單來說.
用寄居者.
他沒有這個心態.
他可能是很安逸.
我在這裡.
他沒有使命.

$^{601}$他也不在乎.
是否流放.
我想這個我們要搞清楚.
是會幫到我們.
作為一個信徒.
或者群體.
究竟怎樣走.
好像黃牧師說的.
我們很高興.
今天有兩位牧師.
和我們不同的城市.
能夠去.
開始討論.
這麼重要的.
這麼切身的問題.
我覺得看聖經.
聖殿被毀之後.
都是一個神學醞釀.
consolidate.
一段很重要的時間.
希伯來聖經就是在被露的時候.
我相信在這個時候.
有這樣的時勢.
上帝都要我們去.
做一些神學聖經上的.
一些影離.
讓我們能夠更加深的在我們這個信仰裡.
打壓這個時機.
讓我們的神學.
反省上能夠更加成熟一些.
首先多謝.
再一次多謝每一位.
我們希望能夠下一次和大家一起.
去討論這麼重要的一個課題.
\newpage



\section{創世記}
\label{sec:17_oWOqNjfw}
\textbf{創世記第6章 (6.2) - 背景釋經 [一百二十年] (陳崇基牧師)}
\newline
\newline
連結: \href{https://youtube.com/watch?v=17_oWOqNjfw}{\texttt{https://youtube.com/watch?v=17\_oWOqNjfw}} ~~~~ 語音日期: 2022-04-25
\newline
\newline
\hyperref[sec:H5eDGDUsh5o]{\small{< < < PREV SERMON < < <}}
~
\hyperref[sec:index_chronic]{\small{[返順時目]}}
~
\hyperref[sec:index_scriptual]{\small{[返順卷目]}}
~
\hyperref[sec:ezSMl3tcxvs]{\small{> > > NEXT SERMON > > >}}
\newline
\newline
$^{1}$我們繼續處理第六章.
第六章我們說是一個很難解的經文.
我們先有一點重溫.
第六章第一節.
人在地上開始增多.
由生養而女兒的時候.
神的眾子看見.
人的女子是美麗的.
Tove.
不是Yavah.
不是Beautiful這個字.
因為希伯文有Beautiful這個字.
就是美麗漂亮的意思.
但是用回Tove這個字.
隨意來挑選.
強取作妻子.
可以這樣解釋.
強取作妻子Lakak.
然後在那些日子和以後的日子.
有巨人在地上.
神的兒子Benahalim和人的女子結合.
生了上古英武有名的人物.
你看到這個經文.
一開始說到洪水之前.
危機的時候.
上帝快要審判了.
一個罪的高峰的時候.
用一些keyword.
是像三章.
人的第一次叛變用的.
這裡記得說.
Kilohim.
像神一樣.
Elohim這個字出現了.
然後女人看到原本的字根.
Rahar出現了.
果子是好的.
Tove.
為什麼呢?.
這裡不用美麗.

$^{41}$因為第六章的point不是女人漂亮不漂亮.
而是人叛變的高峰.
所以用回第一次叛變的keyword.
漢經的果子是好的.
就是叛變是好的.
作食物就是宅.
所以這裡用回了取女子的動詞.
然後有一個呼應.
所以由第三章開始到第六章.
繼續罪惡蔓延.
所以可以這麼看.
一開始上帝創造宇宙萬有.
人類代表上帝在地上做君尊的祭司.
是一個規律.
但是人開始叛變.
第三章.
講到所謂亞當夏娃.
人類的代表.
其實是代表全人類.
個人的叛變.
然後到了第四章.
就講到該人和….
直到trace到第四章.
近尾的時候就是拉麥.
第七代.
就是家庭上的叛變.
就是罪的擴散.
就是兄弟….
哥哥殺了弟弟.
然後家庭是暴力的蔓延.
拉麥的第七代.
然後到第六章.
就是社會或者甚至是國邦.
上一次我們解釋神的兒子.
以我的看法.
背景和物理來說.
Bnei Halel.
似乎是指皇室統治者.
用一個這樣的故事.
這幾個故事來說出.

$^{81}$就是罪闖入世界之後.
人類的罪惡.
使得上帝的秩序.
就變成了反秩序.
由order變成disorder.
招致上帝的審判就是洪水的審判.
上帝就用大水混沌.
再一次去淹沒.
好像我們電腦按一下reset button.
重新再來過.
這個non-order又再出現了.
你看到大水混沌的時候.
就好像一章第二節創世記那樣.
那不是evil的東西.
就是non-order.
回到non-order.
然後洪水過後.
又再重新創立的order.
所以一章到六章.
就有一個這樣的脈絡.
你看到so far.
我們暫時看到是三個半變的故事.
三章四章和第六章.
第六章結束的時候.
就是六章第八節.
羅亞蒙恩.
然後六章第九節開始.
就有另一個Toledo.
另一個大的分段.
就講到洪水的故事.
所以你看到第三章至第六章.
就是罪的擴散.
邪惡的入侵.
以致世界變成了.
本身是order.
現在變成了越來越多disorder.
chaos.
上帝就用水來審判當時的世界.
我們看回這段經文.
六章第一節開始.

$^{121}$我們講過這一段經文.
雖然只有八節.
但幾乎每一句說話.
每一個字都需要去探討一下.
因為引起很多不同的爭論.
不同的解釋.
上一次我們講了Bernard Alguer.
就是神的眾子Sons of God.
如何去解釋.
有幾個解釋的方法.
另外會引起一些問題.
就是120年.
巨人是一回什麼事.
然後.
另外一個翻譯上很重要的問題.
就是於是耶和華說後悔做人在地上.
心中憂傷.
然後他又說我後悔做了他們.
什麼叫後悔呢?.
什麼意思呢?.
為什麼會是後悔呢?.
好像做錯了事情一樣.
我們今次就試試解答另外一些難解的地方.
看看近東的背景.
以及一些上下文.
一些解經的原則如何幫助我們.
來拆解這些爭論.
首先就是120年.
我們先看看第三節.
這個就是動詞.
Filemare.
就是He had and he said.
然後是誰說的.
主體出現了.
就是阿東拉.
然後阿東拉說什麼呢?.
He said no, you're done.
這樣不會住在.
You're done.
可能是abide.

$^{161}$或者remain.
住在裡面的意思.
什麼不會住在什麼呢?.
主體出現了.
Ruhi.
就是Ruha的第一身.
加了這個一撇和一點.
所以就變成了My spirit.
Ruhi.
Vardhan.
就是在人類裡面.
所以他說我的靈.
My spirit.
不會繼續住在什麼呢?.
人類的裡面.
Vardhan.
Nadhan.
意思就是說繼續下去.
或者indefinite的意思.
不一定是我們今天想的eternity的意思.
所以Olam的意思就是一代一代.
不會繼續一代一代.
無止境地.
住在人類裡面.
這個字.
You're done.
可以翻譯為住.
中文通常都是這麼翻譯的.
那個字根可能就是Dalat.
Nunun.
DNN.
You're done.
可能是這樣的字根.
我們參考一下Akkadian.
也可能有一個字跟這個意思是相似的.
DNN.
To become strong.
就是使之強盛的意思.
如果這樣.
到了參考的時候.

$^{201}$希伯來文參考Akkadian.
卡德文.
你就可以看到.
可能意思就是說我的靈會無限期地維持在人的心裡面.
不會繼續持續無止境地.
所以開始倒數了.
開始countdown的意思.
人的生命開始countdown.
或者人當時的世代開始countdown的意思.
所以那個120年就是跟這個有關的.
剛才上面第三節的開頭.
我們已經讀了.
然後他就解釋了.
Bashukam.
這個可能是because的意思.
或者indeed.
因為的意思.
因為什麼呢.
Huvashah.
因為他是flesh.
他是肉體肉血.
繼續說.
Vahayu.
It will be.
Yaman.
他的日子會變成什麼樣子呢.
就是120年.
Meat.
Extreme Shana.
Meat就是100.
然後加了一個Vah.
就是100 and 什麼呢.
20.
Extreme就是兩個10.
然後最後就加一個年.
120年.
120年是什麼意思呢.
可以有兩個的解釋.
他的日子還有120年.
兩個可能會有的解釋就是.

$^{241}$洪水降臨前的時段.
就是有120年.
或者當時的世代人類的壽數.
會定為120年.
兩個我想都通的.
但是第一個通常解釋的時候.
就說羅亞用了120年來建方舟.
也不一定.
因為聖經沒有這樣的指向.
方舟就用了120年.
用了多少年都好.
如果按照第一個解釋的話.
總之就是洪水之前建的.
但是這裡說的是120年的.
所謂Grace Period.
就是上帝被人悔改的恩典時期.
到最後洪水來到.
人類沒有悔改.
為什麼120年呢.
一直以來我們都在看近東的背景.
你記得我們看過一個阿達哈斯史詩.
Traharsis Epic.
裡面也用了一些60位數.
特別是說到什麼呢.
洪水來之前或者災難來之前.
譬如第一塊泥板第352行之後.
說到什麼呢.
人類已經被創造了.
之前我們曾經說過了.
我們不再說了.
就是一個Creation of Mankind的Story.
然後說到1200年還沒有過去.
Land extends.
就是農地或者土地繼續extend.
人多了.
人的噪音也多了.
阻礙著神.
妖神清夢.
神很生氣.
大神Elite.

$^{281}$降一個瘟疫下來.
就是有疫情.
有部分的人死了.
所以安睡可以睡覺.
人的噪音就沒有了.
然後第1塊泥板第416行.
過了之後又說了一些故事.
然後又說1200年又過去了.
然後到第2塊泥板.
故事繼續記載的時候.
第1,2行.
1200年又提到了.
還沒有過去.
又說人多了.
噪音又多了.
妖神清夢.
然後神又很生氣.
又降災.
降旱災.
然後有一個地方是break了.
泥板下面那一部分可能沒有了等等.
但是故事大家都猜到了.
上下文.
就是人類的噪音又多了.
然後人multiply.
噪音多了.
妖神清夢.
大神決定降洪水.
有一個洪水出來.
所以如果你這樣去參考一些近東的背景.
就是什麼事情在這裡發生.
就是神靈因為人口多了.
噪音增加.
妖神清夢然後發怒.
無緣無故地降災.
去毀滅.
裡面提到過這樣的時期.
降災之前.
有1200年.
如果你這樣看.

$^{321}$創世記是在呼應.
或者是大家在分享同一個傳統.
傳統如何述說一些史前史.
或者遠古的歷史等等.
之前我也提過了.
第五章的時候.
洪水大家都知道的.
但是洪水之前的記載之後的記載.
有些地方是相似的.
當然也有很多不同的地方.
當然用這些連數.
有些相似.
但是神學是非常不同的.
因為創世記說到上帝是心哀傷.
不願意毀滅.
願意給恩典.
但人不悔改.
最後上帝的公義審判就會來到了.
跟近東的神論很不同.
跟密蘇波特米亞的神學很不同.
多神論首先是第一個不同.
第二個不同是神的本性.
因為那些大神都不知道在做什麼.
總之就是睡不著覺.
然後就發怒.
然後就將這些災禍降給人.
也不喜歡人繼續多數.
因為人多數的話.
人繼續增加的時候.
就會增加噪音.
所以就會睡不著覺.
所以你看到那些神論是有很大的分別.
雖然記載一些地方的守法.
1200年,120年等等.
可能有些是相似的地方.
洪水提到.
災難提到.
神降災等等.
但是在裡面的神論是完全不同.
為什麼我們要這樣說呢.

$^{361}$因為我們說解經的時候.
除了要留意上下文之外.
一定要留意聖經背後的世界.
因為上帝所啟示的真理.
是藉著當時的世代.
當時的文化背景來啟示出來的.
所以他所溝通的一些事蹟.
或者做的事情,工作.
所以都是跟當時的文化環境.
是相似的.
但是所要表達,啟示的東西.
就是完全不一樣.
就是告訴人們.
神只有一個.
並且祂自己的本性是怎樣.
所以120年我們可以這樣去理解.
哪一個解釋都好.
我們知道為什麼120年不是亂來.
因為可能是跟一些.
米索不達米亞去記載一些史前史的守法.
或者一些的脈絡是有相似的地方.
如果你喜歡我們的影片.
就記得要讚好,留言和分享吧.
\newpage



\section{創世記}
\label{sec:ezSMl3tcxvs}
\textbf{創世記第6章 (6.3) - 背景釋經 [耶和華後悔?] (陳崇基牧師)}
\newline
\newline
連結: \href{https://youtube.com/watch?v=ezSMl3tcxvs}{\texttt{https://youtube.com/watch?v=ezSMl3tcxvs}} ~~~~ 語音日期: 2022-04-25
\newline
\newline
\hyperref[sec:17_oWOqNjfw]{\small{< < < PREV SERMON < < <}}
~
\hyperref[sec:index_chronic]{\small{[返順時目]}}
~
\hyperref[sec:index_scriptual]{\small{[返順卷目]}}
~
\hyperref[sec:fBoYPwedTkU]{\small{> > > NEXT SERMON > > >}}
\newline
\newline
$^{1}$第六章有多少棘手問題?.
起碼有幾個.
第一個是神的義字.
Bnei Ha-el-chem.
第二個就是120年.
第三個就是後悔.
這個字.
耶和華看見人在地上.
罪惡很大.
中日心裡面思想的盡都是邪惡.
於是耶和華就後悔了.
做人在地上.
罪惡很大.
做人在地上.
心中憂傷.
耶和華說.
我要把我所創造的人.
從地上消滅.
無論人或生畜.
爬行的動物.
天上的飛鳥.
我都要消滅.
因為我後悔.
做了他們.
後悔這個翻譯.
就很困擾我們.
或者困擾我們的弟兄姊妹.
在《聖經》的時候.
由第一章讀到這裡.
為什麼神會後悔呢?.
第六章已經有這麼大的難題.
到底應該怎樣去解釋.
後悔這個字呢?.
如果你看看不同的翻譯本.
都有很多不同的翻譯方法.
譬如英文.
NIV就會譯作.
What's grief.
即是他心裡面很哀傷.
Grieving.

$^{41}$或者.
Sorry.
NAS, RSV, NLT.
一個paraphrase.
一個異譯本.
都是這樣翻譯.
以前的King James就是.
It repented the Lord.
所以你看到.
這些新的或近代的英文譯本.
都不用這個字.
Repent.
Repent就是做錯事.
然後他們後悔.
悔改.
NET就用regret.
Regret就未必是做錯事.
Sorry其實不一定是做錯事.
英文我們說.
I'm sorry.
I was sorry.
for what had happened.
我看到有些事情.
發生在你身上.
譬如你的親友去世.
我們就會說.
I'm sorry for your loss.
不代表我做錯事.
只不過是表達一個情緒.
哀傷的情緒.
Grieve當然不一定是做錯事.
所以有不同的英文翻譯.
中文大致上.
都是用後悔這個字.
和學本,新譯本,當代聖經,異譯本.
都是用後悔這個字.
有個意思就是.
其實這個Nahem這個字.
有個意思就是.
轉移,回心轉意的意思.

$^{81}$Change of mind的意思.
和修本.
和學本的修訂版.
就用了不再是後悔.
就是遺憾.
可能這個.
會比較適切一點.
譯做遺憾.
有些事情發生你感到遺憾.
不等於你做錯事.
然後你改變.
遺憾.
可能失去了一些東西.
或者有些很不幸的事情發生.
你不想看到的事情發生.
但是卻發生了.
你感到遺憾.
所以和學本的修訂版.
就是這樣譯.
問題就是這樣.
是不是上帝後悔做了人呢.
就是人犯罪.
罪惡滔天.
是不是他現在後悔.
Nahem.
或者神後悔當初做人.
我們可以去到.
追查一下.
到底這個字.
Nahem.
到底是什麼意思.
首先.
有一些原則.
當我們查這些原文的時候.
查到Nahem.
是一個Nephal的Stem.
就是Hebrew.
希伯來文有不同的Stem.
有Core,有Hebrew,有Hebrew.
等等.

$^{121}$其中一個就是Nephal.
Nephal很多時候.
我們在文法的書上面.
Hebrew grammar的書上面的參考.
就會說Nephal.
就是Core.
簡單的主動詞的形式.
的被動.
Nephal就是被動.
但是.
這個不是很幫到我們.
因為主動被動.
Nahem是被動又怎樣呢.
幫不到我們怎樣翻譯.
怎樣理解那個字詞的意思.
有一個釋經的原則.
一個字有沒有意思.
或者什麼意思.
其實很多時候.
不只是打開字典.
我們就可以找到.
很多時候.
字詞的意義.
是取決於字詞的用途.
它的用處.
和上下文.
所以Again 又是在說一個東西.
Context.
永遠都是在說Context.
因為字的意思.
是取決於Context.
在上下文的語境裡面.
Context裡面是怎樣使用的.
我們才會知道.
那個的意思.
所以我們就需要.
追查一下這個Nahem.
Nephal Stem.
有其他的Stem.
有其他的形式出現.

$^{161}$我們要知道Nephal.
是怎樣用.
那個的處境 語境 Context.
是怎樣用.
似乎都是表達一些情緒.
軟式 安慰.
嘆息等等.
但是這個Nephal.
的Nahem.
很多時候出現於.
一些經文.
上下文文理是在說.
上帝的審判和上帝的恩典.
所以我們就要嘗試.
去做一些Investigation.
嘗試去做一些.
字詞的研究.
首先我們找一下.
當然出現了很多次.
Nahem這個字.
Nephal出現的.
Stem的地方.
這些經文是怎樣幫到我們呢.
首先是耶利米書第八章.
第六節 我留心聽.
但他們都不說真實的話.
沒有人為自己的惡行.
Nahem.
中文可以翻譯為後悔.
即是為以色列人去犯罪.
沒有一個以色列人.
為自己的惡行.
而Nahem.
Nahem是甚麼呢.
我們先不要理會他們是否後悔.
總之他們沒有為自己的惡行.
做錯的事.
Nahem.
竟然說自己.
做了甚麼然後自己偏行幾路.

$^{201}$十八章第八節.
又有這個字出現.
經文說甚麼呢.
如果我輪到的那一幫.
轉來他的惡行.
說的是犯罪.
惡行.
如果上帝說.
我就必Nahem.
不將我原定的.
災禍降給他們.
如果他們行我眼中的.
捍衛惡的事.
不聽從我的話的話.
我就必Nahem.
不把我應許的好處.
賜給他們.
說的是以色列人.
他們犯罪 如果他們轉回的時候.
如果他們不肯轉回.
不肯去Nahem.
我就會將災禍.
降下.
這是約裡面的.
其中一部分.
將Curses的咒詛.
降下來.
這些是以色列人.
第八章.
第十八章我們看完了.
然後到二十六章.
又是耶利米斯.
或者他們肯聽從.
各人轉來自己的惡道.
上帝說.
我就Nahem.
不把我因為他們惡行.
原想要私人的災禍.
降給他們.
繼續看下去.

$^{241}$二十六章 又是同一章.
十三節.
現在先知對他們說.
你們要改善你們所行所作.
聽從耶護華你們神的說.
神的說.
耶護華就必Nahem.
會心轉意.
引師譯道.
耶護華就必Nahem.
如果你回轉的話.
我就會Nahem.
五章所說的災禍.
給你們.
再繼續.
又是先知書.
阿摩斯書.
七章第二節和第三節.
先知見到有災禍.
黃蟲的災禍.
他就禱告.
主耶護華.
太弱小了.
怎能站立得住呢.
耶護華對這事.
Nahem.
耶護華說這事不會發生.
停止了.
因為太弱小了.
亞國家太弱小了.
繼續有災禍就滅族了.
五至到第六節.
因此我對主耶護華說.
求你停止吧.
因為亞國太弱小了.
怎能站立得住呢.
耶護華對這事Nahem.
主耶護華說.
這事也不會發生.
所以在這裡說到.

$^{281}$耶護華Nahem.
所以這些字.
都在這些地方出現了.
繼續看下去.
看多幾節.
約拿書第三章第八節.
這裡說什麼呢.
約拿終於聽話了.
來到尼尼微城.
亞述人的首都.
如果不迴轉的話.
就會有禍了.
亞述王說.
尼尼微城的王.
不論人畜都要披上麻衣.
各人要.
懇切的呼求神.
悔改.
離開惡行.
摒棄手中的強暴.
或者神會Nahem.
不發烈怒.
使我們不至於滅亡.
神看見.
他們所作的.
他們真的這麼做.
這些亞述人.
就悔改離開惡行.
神就Nahem.
不把所說的.
災禍降在他們身上.
所以這裡兩次.
用了Nahem這個字.
中文呢.
就譯做轉意或者回心轉意.
change of mind.
另外橙色這些地方.
也告訴你.
和創世記第六章.
有些.

$^{321}$同樣的字眼出現.
創世記第六章五節.
一路看下去的時候.
你會發現說看見.
上面又有神看見.
地上的罪惡.
惡行.
很大.
中日心裡面思想盡多.
是邪惡的.
又說後悔.
中文譯做後悔.
Nahem這個字.
做人在地上等等.
又說自己說.
我要滅絕他們.
因為我後悔.
Nahem 做了他們.
那時發生什麼事.
世界在神面前敗壞了.
地上滿有.
強暴.
Hamas.
拼棄手中的.
強暴.
說到這些亞述人.
拼棄他們手中的強暴.
神就會.
下面第六章.
講到洪水之前的世代.
地上滿有Hamas.
看到什麼.
惡行 罪惡滔天.
罪惡很大 都是敗壞了.
所以有很多.
同樣的字.
Intertextuality.
互相的.
經文與經文之間.
互相交涉.

$^{361}$互相呼應.
所以有些主要的字.
在亞述第三章出現.
亦在.
創世記第六章.
出現.
所以說做了一些事.
人繼續犯罪.
不聽神的話.
神就不會Nahem.
如果人轉回的話.
神就會.
Nahem 就不會降災.
Nahem 應該怎樣翻譯呢.
所以到現在.
我們看了很多.
與先知書有關.
或與審判救恩.
審判恩典.
這些上下文.
是有關的.
所以我們研究.
Nahem 這個字的時候.
我們就要留意.
在哪個context 如何使用.
另外有一些.
經文.
又是Nahem.
這個字出現的.
雖然不是先知書.
我們先看看.
世紀二十一章.
這裡發生什麼事呢.
簡而言之 這裡發生的事.
就是當時以色列.
那麼多支派有內亂.
有內戰.
因為變亞文人.
做了一些很陰險的事.
很錯的事.

$^{401}$其他的支派就攻打他們.
自己有內戰.
很多變亞文人死了.
這裡就發現.
變亞文的族人.
快要滅絕了.
所以要怎樣做呢.
他們要加一些人給他們.
讓支派不會失去了.
所以這裡.
又用了Nahem 這個字.
Nephal.
以色列人為他們的兄弟.
變亞文族.
Nahem.
說什麼呢.
今天以色列中有一個支派.
被剪除了.
我們曾經指著耶和華起過誓.
必不把我們的女兒.
給變亞文人作妻子.
那麼現在我們怎麼辦理呢.
使得他們城下的.
很少的人.
都有妻子留下他們的族在這裡呢.
就為了這個情況.
他們有規缺.
數字上有規缺.
人數上有規缺 幾乎滅族.
但我們之前又起了假誓.
就是起了不對的誓.
就是不把女兒.
讓他們有生兒養女.
所以他們沒有後代.
所以他們就為了這個.
Nahem.
21章第15節又有這個字出現.
Nahem這個字.
就是眾人為變亞文人.
Nahem.

$^{441}$因為耶和華使.
以色列眾支派中有缺口.
有這個規缺.
所以你看到.
Nahem這個字.
就是和這個規缺有關的.
起碼在這個文理上.
似乎.
這裡告訴我們.
如果你逆做後悔的話.
可能未必表達到.
那個意思.
因為首先現代人聽起後悔的時候.
就會聽到什麼.
做錯了事 所以這次要補鑊.
要改變等等.
似乎逆做後悔.
就不太對.
為什麼呢 如果你看剛才那段.
聖經時詩記.
他們想補鑊給.
變亞文支派.
這個字.
Nahem這個字.
Nahem.
似乎是一種.
核數的.
用詞.
一種核數的用詞.
一種accounting auditing的用詞.
審核一個帳目.
的時候會用到.
大家知道什麼叫做報記.
報記就是bookkeeping.
就是說我們有兩個column.
一個是負債.
一個是credit收入.
就需要最後有一個.
平衡的 做會計的就知道.
或者做核數的就知道.

$^{481}$如果debit是存在.
並且太多的時候.
你就要做adjustment.
做一些平衡.
所以有些制衡.
credit和debit.
要去到.
齊數才行.
所以這裡似乎是說到.
神去看人的時候.
就是.
他衡量或者.
核查人類的時候.
就是看見人.
有一個道德的帳目.
一個moral accounting.
一個moral auditing.
上帝來到去審問人的時候.
就是看他平衡.
如果不平衡的時候.
就有些事情要做.
人轉回的時候.
就會將這個所謂的帳目.
平衡.
check balance.
所以你看見那些先知書.
或者一些用到.
用到Lockhem.
False stamp的時候.
很多時候說到人要做什麼.
你記得了.
二十六章第三節.
就是要做什麼.
二十六章第三節.
除了用了Lockhem這個字.
就是去呼籲.
以色列民做什麼.
或者他們肯聽從.
各自轉來.
自己的惡道.

$^{521}$神說我就會.
Lockhem 我就會平衡.
或者二十六章.
十三節.
現在你們要改你們所行所作.
聽從耶和華.
你們神的話.
就必Lockhem.
不將警告你們的災禍.
降給你們.
約拿書三章第八節.
就更加清楚.
除了用了神會Lockhem.
為什麼呢.
因為人有什麼反應.
人聽見審判的宣判的時候.
先知說話的時候.
他們就說我們要披上麻布.
悔改.
然後我們呼求神.
去惡行披上我們手上的.
Hamas.
強暴.
神就會Lockhem.
就會轉移.
所以如果是這樣看的時候.
這裡似乎是在說.
一個.
或者在假設.
Lockhem這個字.
假設一個Moral accounting.
一個道德上的賬目.
上帝要去.
和人算賬.
就像這樣說.
既然人罪惡滔天.
這個平衡失去了.
越來越多Disorder.
越來越多罪.
越來越多Evil.

$^{561}$越來越多惡行的時候.
神就希望人做什麼呢.
就是要他聽從.
離開惡行然後行善.
悔悟和禱告.
阿摩斯書你記得禱告.
然後上帝去赦免.
這個是人要做的事情.
神在做什麼呢.
就是去審核人的道德.
審核人的心靈.
他可以私恩.
憐憫.
也可以.
將公義的審判帶來.
所以提到這個Lockhem的時候.
似乎這個字用在.
一個這樣的Context(境界)下.
先知書剛才我們看了很多.
加上事司紀.
講到有規缺.
不Balance(平衡).
變雅民族就快沒有了.
所以這裡似乎有個要點.
一個Theological Context(神學境界).
一個神學的紋理.
Lockhem出現的時候.
講到公義的上帝.
是看見人類的道德像木.
就是有規缺.
因為他們的惡行他們的邪惡.
上帝因為是公義的.
他絕不會使這些邪惡.
逍遙法外.
不問不問.
必定上帝會出手.
平衡道德的像木.
問題是怎樣平衡.
上帝當然可以自己開恩.
憐憫.

$^{601}$但是也可以審判.
所以這裡是假設.
或者預設了.
一些道德的像木.
來有一個規缺.
你記得《士司紀》裡面所講.
Lockhem.
人類必要去迴轉悔改.
求上主開恩.
來赦免.
似乎這裡就是.
講到Lockhem的時候.
很肯定不是在講後悔.
或者悔改.
不是上帝悔改.
而是上帝去審核人類的時候.
Evaluate.
Moral Evaluation(道德審核).
來看人的道德的狀態.
來看人的道德的像木.
所以才會用.
Lockhem這個字.
我這裡就引用John Wharton.
的Genesis Commentary.
莊世傑的《釋經書》.
如果你有時間.
你可以看一下.
第310頁和311頁.
他找不到一個.
很悉切的英文字.
來翻譯Lockhem.
中文也很難找到.
他的意思是.
很多時候Lockhem用的時候.
就是講到上帝審查人類的道德.
法國有規矩的時候.
這種的文理裡面.
就出現了Lockhem.
所以他就唯有.
這樣來解釋.

$^{641}$或者好像我們那些.
《聖經》的Study Bible.
加個Note在下面.
就跟讀者解釋清楚.
到底這裡發生了什麼.
翻譯很多時候未必可以.
全部完備地.
來掌握到.
或者溝通到整個文理.
加上這些字.
所表達的意義.
他就怎樣去解釋呢.
John Wharton就說.
要說是審核人類的道德像木.
因為趴在地上造了人類.
他這樣將他的意義出來.
將這段《聖經》的意義.
但他心裡面卻被折磨.
Torment.
處於很痛苦.
所以話說我要把我所做的人.
從地上消滅.
因為我造了他們.
我已經審核了他們的像木.
就是用Auditing的字.
來表達Lockhem的意義.
所以.
他就說.
要說是審核人類的道德.
所以.
這裡說的不是.
上帝做錯了事.
或者上帝造了人.
現在又後悔了.
起初造人早知道不造了.
其實是在說Auditing.
那些做核數師的.
看到一間很好的公司.
他們最後每年核數的時候.
不合數.

$^{681}$就是那種情況.
神也是這樣看人類.
到時候.
要去到Evaluation的時候.
就發現人是有規矩的.
所以神就.
Lockhem.
希望人會改.
所以我們做一點點的總結.
就是這裡說到.
Lockhem這個字.
後悔我相信不應該再用這個字去翻譯了.
正如King James以前是用.
It repented the Lord.
這個英文很明顯就是一個Miscommunication.
因為以前可能會少一點.
這樣的意思.
做錯事改變.
做錯事之後會心轉意.
現在用Repent這個字.
英文絕對有這個意思.
就是犯了罪.
要去到悔改.
因為這樣翻譯就不太對.
所以其他的譯本.
嘗試去到更正或者更新.
以致我們今天看到的時候.
不會有太多的誤會.
Was grieved NIV.
Was sorry 其他的英文譯本.
Regret it 這個可能是更好的翻譯.
剛才我說過了.
後悔就是大部分中文譯本都是這樣.
但是和修版.
將它譯成遺憾.
我相信遺憾是比較適切一點.
當然不可以掌握所有裡面的意義.
正如剛才John Watton所說.
但是可能有些修法的Bible會加上一個闊綠.
告訴大家.

$^{721}$遺憾是什麼意思.
就是神去到審核人類的時候.
心裡面憂傷.
他會以公義.
恩典.
來制衡人類的敗位.
所以希望能夠解答到一點.
這些棘手的問題.
遺憾是一個很好的譯本.
棘手的問題.
上帝見到人類的時候.
審核他們的時候.
道德上有規矩.
神是憂傷.
他會以公義.
或者恩典.
來平衡這個障礙.
在近東背景裡面.
就更加差異.
令我們覺得驚訝.
為什麼呢.
因為我們剛才看了一些.
在之前的session上.
我們看過.
為什麼會有洪水.
就是因為.
神靈睡不著.
所以他就亂來.
然後就降災.
降災之後 原來這麼嚴重.
死了這麼多人.
怎麼辦呢 怎麼收場呢.
連神靈都害怕.
所以這裡表達的神學是完全不同.
或者表達的世界觀是完全不同.
因為在聖經的世界觀裡面.
這個世界的order.
是一個moral的order.
就是一個道德的規律.
是上帝所設立的.

$^{761}$如果人是偏行幾路的話.
就會和上帝的關係會斷絕.
會影響世界.
會敗壞整個世界.
就會使得這個世界.
罪惡更加擴散.
以致去到一個地步.
可能難以收拾.
所以是表達另一種的world view.
世界觀.
表達另一種的神學.
就是神是一個moral being.
可以說是一個有道德的存在者.
他是會問責於人類的.
會keep我們accountable.
我們來到神面前.
我們就是.
如果有規矩的時候就會悔改.
好像先知警告以色列民一樣.
和我們一樣.
我們就是一個.
有道德的存在者.
悔改和警告以色列民一樣.
同樣我們知道這個世界.
如果因為罪惡被越來越破壞的時候.
好像我們今天見到的.
上帝會有公義.
來做平衡.
因為他不單是創造主.
也是審判主.
如果你喜歡我們的youtube channel.
就記得like, comment和subscribe.
\newpage



\section{馬太福音 28:18-20}
\label{sec:fBoYPwedTkU}
\textbf{艱險我奮進 (馬太福音28\_18-20) - 朱莉傳道}
\newline
\newline
連結: \href{https://youtube.com/watch?v=fBoYPwedTkU}{\texttt{https://youtube.com/watch?v=fBoYPwedTkU}} ~~~~ 語音日期: 2022-04-25
\newline
\newline
\hyperref[sec:ezSMl3tcxvs]{\small{< < < PREV SERMON < < <}}
~
\hyperref[sec:index_chronic]{\small{[返順時目]}}
~
\hyperref[sec:index_scriptual]{\small{[返順卷目]}}
~
\hyperref[sec:RE1JRWJg7GI]{\small{> > > NEXT SERMON > > >}}
\newline
\newline
馬太福音 28:18-20
\newline
\begin{longtable}{cl}
\hline
\hline
章節 & 經文 (和合本修訂版)\\
\hline
28:18 & \begin{tabularx}{0.7\textwidth}{X} 耶穌進前來，對他們說：「天上地下所有的權柄都賜給我了。 \end{tabularx} \\ \\ \relax
28:19 & \begin{tabularx}{0.7\textwidth}{X} 所以，你們要去，使萬民作我的門徒，奉父、子、聖靈的名給他們施洗， \end{tabularx} \\ \\ \relax
28:20 & \begin{tabularx}{0.7\textwidth}{X} 凡我所吩咐你們的，都教導他們遵守。看哪，我天天與你們同在，直到世代的終結。」 \end{tabularx} \\ \\
[1ex]
\hline
\hline
\end{longtable}
$^{1}$頂尖妹平安.
今天想一齊思考的這段經文.
相信是大家滾瓜爛熟的一段經文.
也是猜拳聯會或者粉輕會.
經常出現的主題.
《馬太福音》28章18至20節.
沒錯 就是《大侍命》.
雖然《大侍命》這個名稱.
其實在經文本身是沒有用到的.
不會說耶穌說.
我現在將《大侍命》頒給你們.
但是我們也很容易接受.
這段經文會被稱為《大侍命》.
為什麼呢?.
因為的確我們也認同.
這段經文所帶出來的訊息.
傳福音這個訊息.
是我們凡跟隨主的人.
一個很重要很重要的任務.
最大最大的使命.
我今天想嘗試從《馬太福音》28章.
即這段經文本身的上下文.
去討論到底《大侍命》的內容.
和重要性在哪裡.
《馬太福音》28章18至20節.
18節這樣說.
「耶穌進前來.
對他們說.
天上地下所有的權柄都伺給我了.
所以你們要去.
使萬民作我的門徒.
奉父子聖靈的名.
給他們施洗.
凡我所吩咐你們的.
都教訓他們遵守.
我就常與你們同在.
直到世界的末了」.
我相信大家都很熟悉這段內容.
而且很明顯.
《大侍命》的主要的.

$^{41}$關鍵的核心在哪裡呢?.
就是19節.
「所以你們要去.
使萬民作我的門徒.
奉父子聖靈的名.
給他們施洗」.
或者20節的上半節.
「奉父子聖靈的名.
給他們施洗.
凡我所吩咐你們的.
都教訓他們遵守」.
甚至你可能聽過.
這句句子在文法上.
唯一一句的命令句.
imperative在哪裡呢?.
就是「使萬民作我的門徒」.
而另外包括去.
施洗.
或者教導.
其實都是participle.
是分詞.
再直接一點.
再簡單一點.
這樣來說.
《大侍命》最核心的內容是什麼呢?.
就是「使人作門徒」.
而其他另外三個的分詞.
是怎樣去做到.
或者是manner.
或者是means.
怎樣去.
然後替人施洗.
教導別人.
而去達到使命的目的.
就是「使人作門徒」.
所以今天.
其實我只想專注在一件事.
怎樣「作門徒」呢?.
既然「使人作門徒」.
是《大侍命》最核心的關鍵.

$^{81}$我們必須要理解.
怎樣「使人作門徒」呢?.
當然.
門徒訓練這個課題.
在這幾十年裡面.
在基督教圈子是熱門話題.
你去坊間.
無論看書.
或者參與很多課程.
都有很多討論.
今天我只想看《馬太福音》28章.
本身這段經文.
特別是整章28章.
去討論原來.
門徒有一個必要的條件.
我們回到28章的開頭.
28章的開頭.
一至七節.
我們先讀一到七節.
一到七節是這樣說的.
「安息日將盡,七日的頭一日.
天快亮的時候.
活大拉的瑪利亞和那個瑪利亞.
來看墳墓」.
這裡提到.
整件事在哪裡發生呢?.
在墳墓旁邊發生的.
「忽然地帶震動.
因為有主的使者.
從天上下來把石頭滾開.
坐在上面.
他的相貌如同閃電.
衣服潔白如雪.
開手的人因他而嚇得.
渾身亂顫.
甚至死刃一樣.
天使對婦女說.
不要害怕.
我知道你們是尋找那釘十字架的耶穌」.
所以整個發生.

$^{121}$是在墳墓旁邊發生的.
而那個墳墓是誰呢?.
是釘十字架耶穌的墳墓.
第六節.
「他不在這裡」.
所以是一個空墳墓.
當然我們剛剛慶祝伊斯塔.
「照他所說的.
已經復活了」.
我們停一停在這裡.
第一節一直到第六節的上半節.
都是描述的.
即是在說設定.
發生了甚麼事.
開始有些對話.
真正天使的吩咐.
或者跟著前面的描述.
而怎樣有些相應的行動.
在下面開始.
第六節的下半節.
「天使繼續這樣吩咐兩個婦女.
你們來看安放主的地方.
快去告訴祂的門徒.
說祂從死裡復活了.
並且在你們以先往加里尼去.
在那你們要見祂.
看啊 我已經告訴你們了」.
所以有個描述的設定.
而真正的行動.
或者真正的指令.
真正要有一個相應的行動.
跟前面的描述有一個相應的行動是甚麼.
就是天使吩咐那兩個婦女.
去看安放主的地方.
然後去告訴祂的門徒.
有個模式的.
自己去看.
去看自己的.
然後要告訴其他人.
發生甚麼事.

$^{161}$我們看看婦女的結果是怎樣的.
第八節到第十節.
「婦女們急忙離開墳墓」.
其實她們很聽話.
又害怕.
又大大地歡喜.
混合了很多感覺.
跑去要報給祂的門徒.
第九節.
「忽然耶穌遇見她們說.
願你們平安.
她們就上前抱著祂的腳.
拜祂」.
所以這次不單是看到.
空墳墓.
真的看到耶穌了.
第十節.
「耶穌對她們說.
不要害怕你們去高掃我的弟兄.
叫她們往加里里去.
在那裡必見我」.
所以這個模式.
這裡有一個模式出來的.
首先天使吩咐.
「Go see yourself」.
你自己去看.
她們真的看了空墳墓.
甚至後來耶穌的確顯現給她們看.
然後看完相應的行動.
自己看完之後.
下一步是甚麼?.
是「Go tell others」.
告訴別人.
「Go 掃別人」.
然後那些其他人看完又怎樣?.
那些人告訴其他人.
要其他人做甚麼?.
去看.
第十節的下半節.
「叫她們往加里里去.

$^{201}$在那裡必見我」.
所以這個模式.
就是自己看.
告訴別人.
別人也去看.
就像一個循環一樣.
所以門徒必須要有一個條件.
在這裡告訴我們是甚麼呢?.
門徒必須要是見證人.
即是「witness」.
是自己去看.
然後告訴別人.
又讓其他人都成為「witness」.
是一個像是一個「chain action」.
即是一環扣一環.
但是所有的「chain action」.
這些一環扣一環的連鎖行動前.
第一個步驟是甚麼呢?.
是你必須要去看.
所以你必須要成為一個見證.
你看見證是看.
看完然後怎樣呢?.
就要做見證.
就要告訴別人.
以致這個連鎖行動.
是會繼續發展的.
有趣的地方是甚麼呢?.
如果我們這樣看.
原來門徒的首要條件是要成為見證人.
而這一段經文.
在28章的開頭.
這段短短的幾節經文裡.
已經讓我們看到有兩種見證人.
兩種截然不同的見證人.
當然首先我們留意到的.
就是在空墳墓旁邊的那兩個婦女.
末大拉的瑪利亞和那個瑪利亞.
但剛才我們一直讀的時候.
不知道你有沒有留意到.
其實還有另一個群體的人.

$^{241}$另一堆人在那裡.
其實他們都看見.
末大拉,瑪利亞和那個瑪利亞所見到的東西.
但他們有截然不同的反應.
原來在第四節這樣說.
看守的人就因他嚇得渾身顫抖.
所以原來在天使顯現給瑪利阿們看到的時候.
其實有另一群人在那裡.
那些是甚麼?.
看守的人.
甚麼是看守的人呢?.
是否香港人叫做Secure呢?.
其實在27章的尾.
即是在28章前面.
27章的尾有提及到.
我略略這樣說.
當時原來耶穌基督已經被釘十字架.
甚至被埋葬了.
祭司長,法利賽人等等.
就去跟彼拉多說.
要求他派人去看守墳墓.
因為他們害怕甚麼呢?.
他們害怕門徒靜悄悄地.
在晚上的時候偷走了耶穌的屍體.
就會造成很大的困擾.
和會迷惑眾人.
所以當時他們就決定.
派一些看守的人.
所以看守的人來做甚麼呢?.
就是要看守.
應該是耶穌的屍體.
這個就很ironic.
很諷刺.
為甚麼呢?.
因為那個看守的人.
其實他是on duty的.
他是一個職責所在.
他應該是要看守耶穌的屍體.
但第四節形容他們甚麼呢?.
他們反過來被嚇得像死人一樣.

$^{281}$而耶穌的屍體.
當然我們知道.
並不在墳墓裡.
甚至是耶穌復活了.
所以看到原來看守的人.
其實都在場景裡.
但就衍生出兩種不同截然的反應.
有兩種見證.
一個是誠實的見證人.
另一個是不誠實的見證人.
剛才我們說到誠實的見證.
這兩個婦女的確是按照耶穌的吩咐.
告訴門徒.
但看守的人原來他們也有告訴別人.
他們告訴了甚麼人呢?.
我們繼續看11至15節.
他們去的時候.
即是那兩個婦女.
去了告訴門徒的時候.
同時進行.
看守的兵有幾個進城去.
將所經歷的事都報給祭司長.
他們都go tell others.
他們告訴誰呢?.
就是祭司長.
12節.
「祭司長和長老聚集相見.
就拿很多許多的銀錢給冰冰說.
你們要這樣說.
夜間我們睡覺的時候.
他的門徒來說他偷去了.
倘若這話被巡邏聽見.
由我們勸他保你們無事.
冰冰受了銀錢.
就照所祝福他們的去行.
這話就傳說在猶太人中間直到今日」.
原來除了那兩個做見證.
去告訴其他的門徒.
這幾個冰冰都做見證.
告訴了祭司長和文士.

$^{321}$長老等等.
但他們的見證是甚麼?.
他們後來come up with一個甚麼?.
假見證.
他們都有go tell others.
他們都有將這話傳開在猶太人中間.
但他們的見證是甚麼?.
是配合宗教領袖的劇本.
而進行的一個假見證.
為了扭曲真相.
所以有兩種見證人.
一種是那兩個婦女.
而在當時社會裡.
可能大家都會有這個想法.
當時婦女是社會地位非常低微的階級.
所以甚至在court cases的時候.
婦女的見證份量是最低的.
但聖經特別的地方.
就是揀選了這兩個婦女.
成為第一位見證.
至少在《馬太福音》有一章這樣告訴我們.
而他們是甚麼?.
反而是誠實的見證人.
另一方面是那些官兵.
那些看守的人.
他們是搜集權力的一方.
他們甚至有官方的頻道.
去流傳他們的帳戶.
他們的故事.
當然是假帳戶.
是假的見證.
誰敢去質疑官兵的見證呢?.
聖經清楚的告訴我們.
他們是不誠實的見證人.
誠實的見證人和不誠實的見證人.
他們同樣見過兇分母.
他們同樣被警察報告.
將新聞,將訊息.
他們所看見的東西去傳開.
甚至他們同樣都驚嚇.

$^{361}$不單是兵兵驚嚇.
剛才我們有留心的話.
婦女其實是害怕的.
甚至我們沒有時間看下去.
那些門徒都疑惑的.
但為甚麼.
終於有一群人成為誠實的見證人.
另一群人卻是不誠實的見證人.
關鍵在哪裡呢?.
這兩群人的分別在哪裡呢?.
我們回到《大使命教》那段經文.
第十八節,28章,18節.
耶穌進前來對他們.
他們主要就是十一個門徒.
對他們說:天上,地下.
所有的權柄都次給我了.
所以你們要去.
所以,既然有所以.
即是有甚麼?有因為.
原來大使命的來源.
或者大使命的基礎是甚麼呢?.
就是建基於這句說話.
天上,地下所有的權柄都次給我了.
是甚麼權柄呢?.
這種甚麼權柄呢?.
我是老闆,我就有權柄洗我的下屬.
去差遣他們做些甚麼呢?.
聖經是形容.
如何形容這種權柄呢?.
是所有的.
是天上,地下的.
不留餘地的形容詞.
無所不包的.
是最高最高的權柄.
而耶穌基督說.
這個權柄已經次給了祂自己.
如果我們看另一段經文.
一段舊約的經文.
很多解經家這樣認為.
其實《馬太福音》28章這段大使命經文.

$^{401}$是借鏡於《但以理書》第七章.
《但以理書》第七章13至14節.
《但以理書》第七章13至14節是這樣說的.
「我是但以理.
見有一位仁子.
駕著天雲而來.
被領到僅僅上載者的面前」.
這位僅僅上載者當然就是上帝.
第14節「得了權柄」.
這個「仁子得了權柄」.
Has been given 權柄.
榮耀國度.
使各方或萬民.
各國各族的人都事奉他.
他的權柄是永遠的.
不能廢去.
他的國必不敗壞.
這裡有幾個形容詞.
其實跟28章很相似.
所以解經家相信.
可以很平穩地去看這兩段經文.
《但以理書》第七章和《馬太福音》第28章.
而《但以理書》第七章其實是說甚麼呢?.
是一連串的異象.
然後提到「仁子被賦予權柄」.
作為審判.
如果我們真的這樣一連串地看.
《但以理書》第七章和大使命所說的.
天上地下所有的權柄一連串地看的話.
耶穌所說的權柄.
不只是差遣門徒外出傳福音那麼簡單的權柄.
是一種甚麼權柄呢?.
是一種管理,審判.
是一個「to rule」.
去管理,去統治,去審判,to judge.
甚至是甚麼呢?.
是一種國度的權柄.
是一種永遠的.
是超過了天地之間一切權勢的權柄.
如果我們對耶穌所說的權柄.

$^{441}$這個大使命的基石.
有一個更確實,更闊的理解的話.
我們再回想剛才所問的問題.
誠實和不誠實的見證人.
分別在哪裡呢?.
這兩個人為何會做出相應的見證呢?.
分別在哪裡呢?.
就是他們到底在馴服哪個的權柄呢?.
他們在馴服耶穌的權柄嗎?.
還是在馴服祭司長,文士,長老,彼拉多的權柄呢?.
當然,說到這裡很容易說.
我們信了耶穌很久.
讀了聖經很多次.
或聽了很多道.
我們都很容易回答.
我當然是馴服耶穌的權柄.
但我們很現實地看.
其實這從來不是一個容易的選擇.
這往往牽涉到甚麼呢?.
是生命還是死亡呢?.
為甚麼呢?.
因為我剛才提到的祭司長,文士,長老.
法利錯人,彼拉多等等.
其實是當時當地擁有極大權勢的人.
他們甚至可以判人死刑.
耶穌基督的判刑就是這樣.
在這個背景裡進行的.
到底耶穌基督的權柄大嗎?.
還是這班祭司,文士,長老等等的權勢大呢?.
所以我們要將大使命放在馬太的背景.
特別是28章上下文的背景.
這樣看更加有意思.
28章一開始.
場景是墳墓.
是耶穌基督的墳墓.
我們剛才讀過的第五節是怎樣稱呼耶穌呢?.
稱呼耶穌是釘十字架的耶穌.
到了28章的結尾.
耶穌基督再一次與他的門徒團聚.
那裡有另一個形容詞形容耶穌.

$^{481}$是耶穌自己說的.
說他是有國度,有權柄的耶穌.
為甚麼可以在同一章裡有這麼大的反差呢?.
即是有一個被釘的耶穌.
與一個擁有天上地下所有權柄的耶穌.
怎麼可能有這麼大的分別呢?.
或者為甚麼那些門徒婦女驚恐他們而淚.
但他們卻有膽量去抗衡當時的權貴階層.
甚至去告訴他們一個不同的故事呢?.
因為中間發生了甚麼事呢?.
因為耶穌基督已經復活了.
從被釘十字架的耶穌.
到擁有至高無上權柄的耶穌.
中間是甚麼呢?.
是基督已經復活.
墳墓是空的.
而且再顯現讓門徒來看見.
所以基督耶穌復活的意義.
當然對我們生命有很大的影響.
或者對未來的盼望等等.
但我今天想提出的一件事.
可能我們會少留意的.
基督耶穌的復活有甚麼意義呢?.
是祂將世界的秩序重新洗牌.
祂升級了這個世界的權力架構.
當耶穌基督連死的權勢都勝過的時候.
其實世界上有甚麼權柄可以在祂面前誇口呢?.
當世上的權勢爪牙都無法將耶穌困在墳墓裡.
甚至天使說這是一個空的墳墓.
當世上的權勢嘲諷耶穌.
如果大家有印象.
耶穌是以猶太人的「王」這個罪名被釘的.
但是神親自為祂平反.
另一處大家都可能很熟悉的經文.
在《費納比書》第二章第八至十一節說.
「主耶穌全心順服而致於死.
且死在十字架上.
所以神將祂升為至高.
有次級那超乎萬名之上的名.
叫一切在天上的,地上的和地底下的.

$^{521}$因耶穌的名無不屈膝.
無不口稱耶穌為主使榮耀歸於父神」.
祂不只是猶太人的「王」.
所以福音是甚麼呢?.
其實福音就是認定耶穌是主.
祂不只是救主.
也不只是你和我生命的主.
祂是全地的主.
祂是被賦予天上,地下所有權柄的耶穌.
大使明必須要建基於這樣的認順.
想和大家看一張照片.
不知道大家是否認得這張照片在哪裡.
是!這張有耶穌是主的燈牌.
當然是在香港.
其實是在銅鑼灣的一個教會.
在2003年開始樹立的一個廣告牌.
但是照亮了銅鑼灣的夜空.
如果大家熟悉香港的地理位置.
可能會有印象.
其實不單是照亮了銅鑼灣的夜空.
更重要的是照亮了維園的夜空.
在一個很特別的地方.
維園維多利亞公園.
在銅鑼灣的裡面.
這個教會是2003年開始樹立的.
2003年在香港歷史中是一個很重要的一章節.
在2003年7月1日有50萬人上街.
有很多訴求.
反對很多事.
向政府表達一些意見.
自此之後維園這個地方.
成為很多民意表達的地方.
在過去十多二十年的日子裡.
香港風起雲湧的地方.
都是在這個燈牌下面發生的.
所以到底耶穌是主?.
發生了這麼多的事.
耶穌真是主?.
還是我們能夠說.
耶穌是主?.

$^{561}$認同耶穌是主.
祂是天上地下所有的權柄都賜給祂的那位主.
今天想跟大家分享的題目.
我講章的題目.
叫「奸險我奮進」.
我不知道有多少人認識這句話.
這句話是中文大學的前身.
新亞書院的創辦人錢穆教授所寫的.
錢穆教授在1949年.
中國大陸政權易守之後來到香港.
那是一個顛沛流離的時代.
他創立了新亞書院.
始自香港這個彈丸之地.
當時仍然是英國的殖民地.
去傳承中華文化.
是很艱難的一件事情.
所以他為新亞書院寫的校歌.
其中有這句話.
「奸險我奮進」.
其實就是勉勵學生要堅毅刻苦.
勇於承擔的這種精神.
亦都成為後來一代又一代中大人.
自勉的格言.
我不是中大人.
但每一次聽到這句話.
「奸險我奮進」都很有感覺.
其實這句話對我們這班基督徒.
對我們有傳福音使命的基督徒.
又何嘗不是這樣呢?.
為何傳福音會受人恨惡.
為何要面對患難,面對迫迫呢?.
因為傳福音在一句話之下.
最重要的是甚麼呢?.
其實傳福音就是向這個世界宣告.
其實有另一個更高,更大的權勢.
傳福音是甚麼呢?.
就是向這個世界所流行的價值觀.
去挑戰,宣告有一個更高的準則.
有一個更需要我們遵守的標準.
所以當我們想誠誠實實.

$^{601}$為真理做見證的時候.
怎會不面對迫迫呢?.
作門徒的確是要付上代價.
大使命不是用命令來做完結的.
大使命是以應許來做結束的.
神應許有同在的恩典.
不是我們自己頂.
不是將項目交給我們.
然後祂就翹起手.
祂和我們一同同在.
然後這位曾經為我們的罪.
釘在十字架有死而復活的主.
祂已經敗壞了掌死權的魔鬼.
和世上一切權柄的這位耶穌.
祂和我們一起.
所以我們有盼望能夠完成大使命.
在馬太福音28章結束的時候.
耶穌將大使命交給十一個門徒.
主要十一個門徒.
這是不可能的任務.
十一個門徒怎能夠使萬民作門徒呢?.
因為從來不只是這十一個信徒承擔.
這個門徒承擔.
不只是一代去完成.
馬太福音的開始第一章是家譜.
其實馬太福音的結尾28章.
是另一個延續的家譜.
由十一個門徒.
十一個使徒.
一代一代一代傳下去.
如果我們再回看經文的話.
經文怎樣說呢?.
耶穌基督教導我們的東西.
你們要教訓他們.
有我,有你們,有他們.
傳福音的使命.
是一個馬拉松的接力賽.
其實我們只需要做什麼呢?.
誠誠實實地走好我們這一棒.
然後將我們這一棒.

$^{641}$在我們不同的崗位.
在我們日常的生活裡.
在我們每一天的日常行動裡.
去認定祂是主.
唯有祂是主.
捉緊這一棒.
然後將這一棒交給下一代.
最終末期必會來臨.
有一天我們都能夠見證.
耶穌基督的應許成就.
我們一起來祈禱.
讓我們經過復活節.
讓我們不單止去慶祝.
你為我們所成就的救恩.
求你更加幫助我們.
去反思我們的生命.
到底在世上,在地上.
我們有什麼任務,有什麼使命要去完成.
以致我們能夠有一天.
去見到你的時候.
我們能夠歡歡喜喜.
坦誠地去面對你.
救主幫助我們.
在地上儘管有很多的艱難.
讓我們能夠成為你誠實的見證人.
將這個福音的棒.
一代一代地傳下去.
求你恩代我們.
求你祝福我們的教會.
求你祝福我們的社區.
我們這樣禱告.
是奉耶穌基督的聖名.
\newpage



\section{創世記}
\label{sec:RE1JRWJg7GI}
\textbf{聖經考古學 - 創世記,洪水,方舟 (2)  -  陳崇基牧師}
\newline
\newline
連結: \href{https://youtube.com/watch?v=RE1JRWJg7GI}{\texttt{https://youtube.com/watch?v=RE1JRWJg7GI}} ~~~~ 語音日期: 2022-04-29
\newline
\newline
\hyperref[sec:fBoYPwedTkU]{\small{< < < PREV SERMON < < <}}
~
\hyperref[sec:index_chronic]{\small{[返順時目]}}
~
\hyperref[sec:index_scriptual]{\small{[返順卷目]}}
~
\hyperref[sec:2oEQOesAhNs]{\small{> > > NEXT SERMON > > >}}
\newline
\newline
$^{1}$今日我想講第15段落.
第15段落說地質學不能夠證明.
不能夠支持全球的洪水.
這個就是Steven Mosher.
以前我讀考古學的時候.
我都是收他的課.
他創立了一個新的部門.
一個新的學術名稱叫做Geoarchaeology.
地質的考古學.
他做了很多工作 現在也繼續做.
就是去跟蹤出埃及的路線.
講述當時出埃及的時代.
在所謂的紅海.
其實不是紅海 是草海.
草海在哪裡呢.
過紅海的地方在哪裡.
怎樣去用考古學 用地質學去還原.
當時波西的時候.
那個古代的環境.
很興奮 用一些以前美國有的間諜衛星的相片.
用GIS等等.
用Google Earth等等.
去追查當時的情況是怎樣.
原來是跟出埃及的路線是吻合的.
做了很多這樣的工作.
他本身是一個專業的地質學家.
他說紅水地質說是什麼意思呢.
就是說在基督教裡面.
美國也有的.
一些很保守很右派的人士.
認為這個地球很年輕.
因為聖經告訴他們很年輕.
因為他們對聖經的傳息.
這個是Young Earth.
並且他們的傳息.
就是說現在世界各地看到很多的山脈.
大峽谷等等.
這些是什麼時候形成的呢.
不是科學家說很長時間.
水侵蝕才會形成的.

$^{41}$他們認為是曾經有一個全球的大洪水.
來整頓整個地球.
然後有這樣的大峽谷等等.
這樣的山峽出現等等.
這個所謂的紅水地質說.
Flood Geology.
他說他反對.
他反對這個紅水地質說.
Flood Geology.
完全是基於相信聖經必定要要求.
有一場全球性的洪水.
就是說他們的解說.
認為聖經如果要解的話.
只有一個可能性.
就是曾經有一個Global Flood.
有一個全球的洪水.
如果不是的話.
你就是新派了.
幾乎是說成這樣.
這些支持者訴諸於什麼理由呢.
就是一些洪水而來的地質的過程.
去到推測的假設.
而這些假設是沒有在聖經裡面.
在經文裡面找到的.
也都見不到.
在大自然裡面見不到.
我們不認為.
就是斯蒂夫·摩西說什麼呢.
我們不認為聖經敘述.
必定要求我們全息的時候.
結論是有一場全球性的一年的洪水.
也都找不到任何古代的岩層.
或者地景的特質.
只向曾經有一場一年性的.
一年之長的全球的大洪水.
這個是一個專業的基督徒的地質學家.
他的結論.
如果你想看的話.
你儘管買一本書.
電子書也有.

$^{81}$你買來看看.
他的立論就是這樣.
那你說.
不會說他是信聖經的.
怎麼會跟聖經唱反調.
我就想告訴你.
福音派的學者.
不一定認為聖經是要求我們.
洛亞洪水的經文.
要求我們在立論的時候.
有一場全球性的洪水.
他就說了哲學.
我這裡想加插的是建築學.
就是考古學.
考古學怎麼看呢.
首先我們想一想一些方法.
不能夠說每樣東西都可以證明.
隨便找到一些東西.
曾經被水破壞過的.
又說是全球的洪水.
又說在山上找到海洋生物的花石.
證明了全球的洪水.
我們首先想清楚.
共和黨的主旨.
方法論是什麼.
每樣東西都有方法.
方法論是什麼呢.
考古的方法論.
如果你說有一個全球的洪水.
Global Water.
水位上漲到像飛士峰一樣.
你要找考古的證據.
怎麼找呢.
首先第一樣.
方法來說.
首先我們是否需要決定.
斷定到洛亞洪水的時期.
時間是什麼.
創世記六節九章.
洪水泛濫一年多.

$^{121}$然後水就退去了.
寒地又出來了.
這個時間到底是什麼.
先決定了時間.
然後你才可以在近東找一些東西出來.
考古學就是做這件事.
我剛才說大衛的歷史性.
為什麼可以證明呢.
因為我們知道大衛大概是什麼時間.
他是主前大概一千年的人物.
你把聖經裡面的時間的資料.
去到組出來.
大概是1000BC.
主前一千年的時候.
所以大衛的王朝出現.
就一定是1000BC之後.
剛才說的彈城石碑.
Talden的石碑.
亞蘭人的石碑.
是屬於900BC.
所以是合理的.
所以在時間上一定要合理.
考古學一定要問.
就是什麼時候.
時間性是什麼.
所以你先要決定.
洛亞到底是什麼時期的人.
洛亞的洪水是什麼時候發生.
When did it happen.
如果你先不決定這件事.
那你就要龍門.
龍門就要你搬.
你說找到什麼就能證明.
所以你首先要先決定這個時期.
第二就是洪水的時間.
如果你決定到的話.
就必須配合近東的歷史.
我們現在對近東的歷史.
已經有很多掌握.
很清楚.

$^{161}$什麼時候 什麼年 什麼月 什麼日子.
什麼時候發生什麼事.
大概有什麼找到等等.
好了 說回第一.
洛亞洪水的時期.
可不可以計算出來.
可不可以用創世記第十章的家譜.
回頭計算.
可以計算到洛亞的洪水時間.
答案是不可以的.
因為家譜一定會跳袋.
並且我們現在知道.
我們研究近東民族的家譜的學者.
都說家譜最主要寫作的目的.
不是去講時間和歷史性.
而是講身份.
所以家譜的目的.
不是讓我們去計算時間.
沒有這樣的東西.
古代甚至近代.
這些很保守的近東的中東民族等等.
你去研究一下他們.
他們講家譜的時候.
不是讓你寫出來.
讓你計算時間.
而是家譜的目的.
是社會的目的.
即是說是一個社會性的.
民族性的.
一個名份 一個身份的宣稱.
所以你看聖經裡面提到家譜的時候.
通常是在講這個人有沒有資格.
qualification.
以色列回到耶路撒冷.
你看一下以色列記.
怎樣去介紹以色列.
或者摩西大小人出埃及之前.
又提到他有家譜.
這個人有沒有資格.
做我們的所謂的英雄.

$^{201}$能夠做我們的領袖.
所以就有家譜出來了.
所以這些家譜.
很多時候都不是歷史上很準確的.
也不需要準確.
因為最主要講的是身份.
既然你找不到.
或者決定不到羅亞洪水的時期.
那你找什麼都不可以配合到.
第三.
如果你真的要找到一些大洪水毀滅的證據.
其實就不在山上找.
所以就不要寄望在阿拉納山上找方舟.
找到一些木頭等等.
可能你聽過這樣的報告.
或者一些這樣的故事等等.
我可以告訴你很直接的.
證據不在山上.
不在山上找船出來.
而在哪裡找呢.
如果你真的想證明有一個全球洪水.
一年性發生的.
氾濫一年多.
毀滅整個全地.
證據就在哪裡呢.
不是在山上.
在地下.
為什麼呢.
因為你去近東的世界找找一些考古的土幽.
那些古城被毀滅.
有沒有一年裡面.
全部都被毀滅.
所有的土幽同一年被大洪水所淹沒破壞.
所以就造成一年性.
一年裡面出現的破壞的土層.
Destruction layer.
並且沒有一個土幽.
有很多土幽是幾千年幾萬年的歷史.
一層一層的上去.
是沒有斷過的.

$^{241}$不斷有不同的文化.
不同的人在上面建造.
有一個連續性.
如果你說有一個一年的大洪水.
什麼都毀滅了的話.
那不可能有土幽可以安然無恙.
跨越那一年.
對不對.
所以一定要找到destruction layer.
並且今天很多人想在阿拉納山.
Alberta 上面找到方舟的殘骸.
找一些木頭下來.
我就想說.
我們都要講方法.
意思就是說.
這個所謂現在在土耳其東邊.
阿拉納山的傳統.
在什麼時候才出現.
以我們所知.
所有的古代的文獻.
現在這個所謂Mount Everest.
阿拉納山.
是主要在十一二世紀才出現.
所以距離聖經所說的羅亞時期.
真的太遠了.
所以是不可求的.
OK.
以後期的想像出來的.
這些傳統是不可求的.
這個都是歷史的方法.
OK.
創世記第八章第四節.
七月十七日.
方舟停在阿拉納山區.
Ahrere Arare.
Ar 是一個前置詞.
On 意思是在上面.
Ahrere 就是Mountains的意思.
Arare 就是阿拉納.
這個在亞述的文獻裡面出現過.

$^{281}$在阿拉納 亞美尼地區.
即是今天的土耳其東邊.
大家想留意的就是Ahrere.
這裡看到一撇肉和兩點.
其實是一個中數.
方舟不是停在阿拉納山頂.
不是說一座山.
而是一個山區.
一個山區是一個很大的山區.
所以可以停在山頂.
可以停在山的中間半山.
也可以停在地下的山區.
Ahrere Arare.
所以留意原文.
不是我們今天聽到很多.
想去找方舟的人.
正名方舟的人所說的.
這個阿拉納山.
Again 是主後.
十一二世紀非常後期.
AD 1110-1200.
才出現一個這樣的傳統.
一個傳說 上面有隻方舟.
所以我們在說的不是方舟.
是說方法.
你才掌握到方法.
這個方法才是合理的.
你才可以去找證據.
你不可以貿然上山.
去找一些東西出來.
很多這些所謂的考古學.
其實沒有一個很嚴謹的方法.
我先說一些概念.
一些考古學的概念.
正如我所說.
我是讀考古學的.
考古學裡面.
有一樣東西很重要.
叫做Tull.
或者阿拉伯文Tull.

$^{321}$都是說一些所謂的土優.
山優.
但這些山優不是自然的山崗.
不是自然的山崗.
是人造的山崗.
是Artificial Mount.
Artificial Hill.
為什麼是Artificial呢.
為什麼是人工造成的呢.
這裡有一個很簡單的構圖.
可以告訴你.
一個考古學家去看這些資料的時候.
到底他如何斟酌和解釋這些東西.
你見到這裡有一個.
我叫它做考古的土優.
有不同顏色.
一個考古學家到場的時候.
他就會在這些土優的頂部.
向下挖掘.
一直向下挖掘的時候.
挖掘了第一層.
發現這一層的實物文化.
圖詞等等用品.
都是差不多一樣的.
這個就是一層.
很多時候用一個羅馬的數字.
去表達出來.
One First.
然後挖掘到再深一點的時候.
發現原來這裡有另外一種圖詞的用品.
另外一種風格出現了.
他就會說這是另一個文化.
越挖掘到底的時候.
就越遠古了.
因為以往的人.
一直從古到今.
一直在上面建造的時候.
往往未必會把東西清理好.
然後才在上面搭建一些東西.
他一直被那些塵土封住.

$^{361}$然後再在上面再建造.
所以那些土優就不斷升高.
你去以色列巴勒斯坦地區.
去中東.
經常見到這些東西.
你開車經過.
這個山不是一個自然的山.
經過長年累月.
很多不同的文化在上面累積.
越來越高.
所以考古學家一直向下挖掘的時候.
就會找到不同顏色的文化.
不同的文化.
舉個例子.
今天去某個地方.
挖掘的時候.
發現有iphone X.
每間屋都有iphone X.
iphone 9, iphone X.
在那裡找到iphone X.
你就知道那個地方.
就是屬於什麼時間.
但你繼續挖掘的時候.
你會發現沒有iphone.
沒有手提電話.
只有有線的電話.
以前我們舊一代的有線電話.
有圓形的電話.
你會發現這個就是另一個文化.
再繼續挖掘的時候.
發現電話都沒有了.
就更加古代了.
這個就是另一個文化.
所以就由這樣的方法.
你會看到一些考古的土優的報告.
會告訴你有不同的土層.
不同的土層就用不同的數字去表達出來.
最上面的土層.
就是一個比較近代的.
然後繼續挖下去的時候.

$^{401}$就會越挖越久.
越挖越遠古.
所以在這裡.
直到你去到藏岩.
即是自然的山的地方.
你就不會再下面有東西.
所以第一代在這裡住的.
就是5.
5就是Paleolithic.
Old Stone Age.
就是古石器時代.
即是二萬年前的.
然後就是後舊石器時代.
到Level 3.
這裡有一個例子.
就是到這裡我們找到什麼呢.
就是藍色.
先不要理會這裡的顏色.
最主要是淺藍色.
這裡你可以假設.
就是一個洪水經過.
然後就毀滅了這個Number 4的Level.
所以洪水經過之後.
就有很多淤泥積聚在這裡.
那些淤泥就找不到任何物件.
只是一些淨的淤泥.
你就發現了.
這個時期.
曾幾何時就有一些洪水經過.
留下一些很清楚的沙土.
然後沿著洪水.
這些沙土淹沒了所有東西.
在上面又有.
另外一些文化建造在上面.
所以這個土優的結構.
給大家一個概念.
就是如果有一個洪水毀滅的跡象.
你就在這個土優裡.
你會找到什麼呢.
你會找到這些破壞層.

$^{441}$Destruction Layer.
在洪水裡面留下的.
我們看看考古學的結論是什麼呢.
我引用一個學界的論文.
叫做North Flood Estate Extent and Divine Accommodation.
是一個很傳統很保守的福音派的.
在一個神學期刊裡面找到的.
叫做Paul Seely.
是一個神學期刊.
叫做Westminster Theological Journal.
他的結論很簡單.
他說地質學家.
甚至是考古學家.
有很多基督教的地質學家.
都不相信有一個證據.
可以表明有全球的洪水.
因此沒有了這些侵蝕的效應.
或者沉積的泥土.
是證明不了有全球的洪水.
他引用一些地質學家說.
如果在5000年7000年之前.
曾經出現了一場全球性的洪水的話.
這些洪水的侵蝕的沉積的果效.
是沒有可能消失的.
因為一定會在很多土壤留下破壞症.
另外他說地質的證據.
是證明不了全球的洪水的存在.
testify against a global flood.
這些房屋.
找到一些土壤例如耶利哥等等.
在整個近東.
沒有表明到沒有任何證據.
告訴我們同時經歷過一場洪水.
由主前9000年之前開始.
到阿巴拉罕時期.
這些土壤一直都有連續性.
有很多土壤都是一直連續下去.
例如耶利哥.
主前9000年.
直到主前幾世紀.

$^{481}$都是一直一層一層地疊上去.
所以既然沒有破壞就沒有被破壞過.
那就沒有可能有全球的洪水.
地質學,考古學等等的數據.
冰川學.
幾乎所有領域的合資格科學家研究.
當中大部分都是基督徒.
都是同意.
在過去十萬年.
100,000 years.
甚至更遠遠的時期.
都沒有發生過全球的洪水.
所以科學告訴我們.
我們聖經的傳息應該怎樣去面對這些問題呢.
還是我們掩著耳朵當作沒有聽過.
我們照舊.
好像以往一樣自免色竟.
這樣是不行的.
我們試一下猜測一下洪水的時間.
1960年的時候.
美國有一本很有影響力的書就出了.
叫做Genesis Flood.
即是創世記的洪水.
是保守神學院的教授Wilk.
和一個基督徒科學家Henry Morris所寫的.
他們認為是有全球的洪水的.
因為他們解說聖經.
傳適聖經的時候.
傳適羅亞洪水經文的時候.
他們認為只能夠有全球的洪水.
所以他們認為那些家譜大概可以讓我們猜到.
到底羅亞大概是什麼時間呢.
我們當他們可以這樣做.
這本書幾乎是今天的天書.
那些中右份子.
那些很傳統的高福夫的支持者.
都會回到這本書.
他們認為洪水的時間大概是主前7000至5000年的時間.
他說不可能再推到再前一點.
7000至5000.

$^{521}$這很合理.
為什麼呢.
因為我們看了這麼久創世記.
你記得第四章提過什麼嗎.
阿伯和該隱.
該隱是種地的.
阿伯是畜牧的.
這個耕種的工業.
人類發明耕種.
我們大概知道是什麼時候出現.
在近代歷史裡面.
畜牧.
Domestication of sheep, flock, cattle等等.
Domestication.
就是馴養這些動物.
有羊,牛等等.
我們大概都知道是什麼時候出現.
那個時間就是在新石器時代.
Neolithic period.
那個時代大概是主前8500至4500年的時間.
所以第四章已經有.
甚至有說打銅鐵等等.
該隱的後裔.
有這樣的工業已經出現的時候.
肯定在近東的歷史裡面.
一定是8500年之後.
新石器時代發生什麼事呢.
新石器時代告訴我們.
人類好像在一夜之間.
在一個大的時間表裡面.
一夜之間.
發明了農業.
知道怎樣耕種.
然後知道怎樣去馴養這些動物.
所以有羊有牛等等.
這個有時候我們叫做新石器的革命.
Neolithic Revolution.
因為突然間.
好像有一段時間.
人們可能做實驗.

$^{561}$扔種子去種植.
然後又懂得怎樣去做農業.
怎樣有種植物出現.
這個時間大概是8000至4000年.
主前的時候.
所以創世記第四章所描述的.
其實是fit進去這段時間.
fit進去近東的大歷史裡面.
如果是這樣的話.
羅亞就一定是這段時間.
或者之後.
不能夠再推它前面.
所以阿巴拉罕的時間.
大概主前2000年.
如果我們這樣去推算的話.
Henry Morris就認為.
大概洪水的時候.
應該是7000至5000年BC的時候.
這個也make sense.
make sense意思就是說.
sort of 順序.
有了第四章.
這些工業啊.
耕種啊等等.
然後才有羅亞.
你不能夠農門任你擺.
所以大概ok的.
7000至5000年BC這段時間.
如果是的話.
我們就問一下.
在近東的大歷史裡面.
考古的土優.
很多這種土優.
多到一個地步.
只是計以色列.
現在有考古學家.
發掘過的這些土優.
只佔5\%.
我們知道的考古的土優.
遠古歷史的土優.

$^{601}$很多個.
現在考古學家正正式式.
發掘過的.
只不過是5\%.
所以我們有很多人還沒有開的.
還沒有打開的.
所以我們知道這段時間.
有很多不同的土優.
新世紀之後的土優出現.
如果你說有全球的洪水.
就一定是這些土優.
這段時間的土優.
就出現一次性的.
同時的.
同一年的被洪水沖沒的.
這些實證.
這個考古學的方法.
你不能夠說.
牧師啊.
這些不會出現的.
因為那是神蹟.
耶穌在加里里湖.
在那裡行水.
我們不會說.
我們想找一些歷史證據.
證明這件事的事實.
當然不會的.
所以如果你真的在說歷史的話.
在這2000年裡面.
按照Henry Morris所說.
考古學裡面應該可以.
拿到這些證據.
所有這些土優.
都同一時間同一年.
有這些破壞層出現.
只不過是沒有的.
所以這位學者說什麼呢.
如果有一次的大洪水.
將一個大的方舟沖上山.
必定在全球各地.

$^{641}$同時.
記住一定是同時.
在那一年出現了一些.
很實在的考古證據.
可視是沒有的.
我們將這件事.
放在這個.
這個.
時間線上.
這些時間線.
就是我們現在知道.
近東歷史的時間線.
由之前.
舊石器時代.
到後舊石器時代.
到大概8500年.
之後.
就去到新石器時代.
就是我說的有龍耕啊.
有畜牧啊.
出現了.
這段時間.
然後就到.
同氣石器.
共用的時代.
接著就到.
這段時間.
大概已經去到2000年.
2000年你記得就是阿巴拉罕.
摩西之後.
這個就是中期的時候.
這個整個就是.
很簡單將近東的.
大家有共識的.
時間表.
如果你說.
估計這個.
裝飾的六至九章.
諾亞的時代.
是這個7000至5000BC的時候.

$^{681}$OK.
那你就要查什麼.
你就要查這段時間.
近東的.
這些考古的土有.
有沒有.
如果沒有的.
那就是說.
不可能是全球.
剛才我們看的神學奇行.
斯里就認為是2900.
那他呢.
都是make sense(合理)的.
因為他說什麼.
因為他說這些畜牧業啊等等不同的工業.
以我們現在所知的近東裡面.
出現的這些.
發展啊.
人類的不同行業啊或者技術的發展啊.
等等.
有這些同氣啊等等.
他們整船啊.
他認為是2900.
那又是斯里.
是一個.
都是保守的學者.
Westminster Theological Journal.
很保守的.
一份神學學界的奇行.
他說大概是2900.
3000BC.
那你就要查一下.
這些時間這些的土有.
有沒有這些跡象.
在1930年代的時候.
有一個英國學者.
叫做Wooley.
他就被派到.
去到.
這個伊拉克.

$^{721}$去到這個吾爾的古城.
就是阿巴罕的家鄉.
吾爾的古城.
去那裡做法規工作.
他其中一個呢.
當時1930年代算是早期一點.
考過的方法.
就釘了一個大的坑出來.
他就挖了一個大的坑.
那個大坑裡面呢.
他就看到一層一層一層的.
文化不同的土層.
其中他找到的就是.
一個八尺深的.
一個淨淤泥的土層.
是洪水流下的.
這個故事是很出名的.
當時30年代.
回到英國的時候.
世界各地的.
大的報章就去報導這件事.
就說他確認了.
聖經的水流故事.
大概在4000BC的時候.
因為按照那些土層.
大概在3500-4000BC左右.
所以就說.
考古學家終於找到了.
這個洛亞洪水的跡象.
考古的記錄.
在烏尼瓦的圖就是這樣.
就是剛才我們那個簡圖.
所表達的東西.
就是你見到這個坑位.
然後有很多一層一層的.
不同的所謂.
實物的文化.
不同的土質等等.
你見到最底下.
已經到背景了.

$^{761}$即是沒有人性證據.
即是下面是床岩.
是石地.
下面不會再有.
其他文化.
在上面一層一層的.
再計算的時候.
就發現第一層.
你見到有泥磚.
灰磚.
或者一些礦物質.
一些瓦器等等.
這裡算是第一層.
這個層大概在6500至3800BC.
就是烏尼瓦文化.
考古學家叫烏尼瓦文化.
無論如何.
總之有同一個風格的瓦器出現.
然後就是剛才所說的.
有8呎多的洪水造成的淤泥.
按照時代的時間的考證.
就是屬於主前3500BC.
這些叫flood deposit.
水沖過河水帶這些淤泥.
經過毀滅了這個烏比的土層.
然後留下了8呎的淤泥.
過了之後.
就是3300至2000BC.
這裡一大段是連續性的.
一層一層加上去的.
這些人類留下的.
實物的文化.
你見到有pottery.
或者有些房屋等等.
所以這裡你見到.
中間突然間有一個8呎.
但是很可惜地.
有很多人以為可以證明到.
全球的洪水.
誰知大家在其他土壤挖掘的時候.

$^{801}$同一個時期.
其他的土壤沒有這些所謂的flood deposit.
這裡曾經有一個大洪水是本地的.
但是毀滅了.
只不過是吳爾這個地方.
或者部分吳爾的土壤.
其他在吳爾附近的.
或者更遠的.
可能找到不同的洪水的淤泥.
可能是不同時期.
但不是這個時期.
或者找不到有任何洪水毀壞的跡象.
所以在考古上.
沒有一個證明.
去告訴我們有全球的洪水.
有局部的洪水.
肯定.
這裡證明了已經有.
或者在其他地方找到.
同樣的現象.
或者其他考古國家.
在伊拉克等不同的地方找到的.
這些現象.
都可以告訴我們.
有一個局部性的洪水.
所以如果看近東的背景資料.
你會發現.
歷史上曾經有一個很大的.
或者破壞性很大的洪水.
但那個洪水絕對不是全球的.
是局部的.
我們就引申很多問題.
有很多人問.
如果是沒有全球的.
如果只是局部性的洪水.
何必弄方舟.
叫它走上山就可以了.
這個很有趣.
這個可以繼續解釋.
但我這裡很簡單地說.

$^{841}$誤會了方舟是什麼意思.
到底什麼是方舟.
方舟是一個小宇宙.
是上帝重整的世界.
方舟是代表新的創造.
所以為什麼要有個方舟.
我們想方舟就是想避難.
聖經裡面所說的方舟故事.
是在說創造.
反創造 毀滅當時的強暴人類.
當時強暴的社會.
然後再重造一個新的出來.
就是用方舟重造.
方舟是有神學性的.
不只是避水這麼簡單.
所以這個很簡單.
這個後話我們以後有機會繼續說.
如果你喜歡我們的YouTube頻道.
就記得Like,Comment和Subscribe.
\newpage



\section{4:13-18}
\label{sec:2oEQOesAhNs}
\textbf{亂世啟示: 空中會見主 (帖撒羅尼加前書4\_13-18) - 陳崇基牧師}
\newline
\newline
連結: \href{https://youtube.com/watch?v=2oEQOesAhNs}{\texttt{https://youtube.com/watch?v=2oEQOesAhNs}} ~~~~ 語音日期: 2022-05-05
\newline
\newline
\hyperref[sec:RE1JRWJg7GI]{\small{< < < PREV SERMON < < <}}
~
\hyperref[sec:index_chronic]{\small{[返順時目]}}
~
\hyperref[sec:index_scriptual]{\small{[返順卷目]}}
~
\hyperref[sec:eLoBtjnY0Kg]{\small{> > > NEXT SERMON > > >}}
\newline
\newline
$^{1}$繼續這個《亂世啟示錄》的系列.
我們之前的Good Friday和Easter暫停了一點.
但是我們回到這裡.
我們看看和這些所謂的《宗末論》有關的.
Eschatology有關.
當我們想起這個主題的時候.
什麼叫做宗末的意思.
請你不要覺得這件事只不過是一個生命裡面的註腳.
或者是我們信仰裡面一些很悠閒的事情.
不等閒的事情.
其實整本聖經如果你看.
去揭示給我們看人類歷史的方向.
是一個直線來的.
是一條線不斷向前面一個目的.
這個就正是宗末的意思.
就是Eschatology的意思.
宗末就是說向著一個目的地來進化.
這個目的地就是上帝所計劃的一個目的地.
所以你不要覺得這個只不過是一些註腳的事情.
就是聖經裡面最後那本書才叫做《啟示錄》.
講到末世末日發生的事情.
到底世界裡面會發生什麼等等.
以為這些只不過是一個註腳.
其實整本聖經是向著這個方向來進發的.
最近有一個美國的電腦教授.
他就將聖經裡面63779節的.
這些互相呼應的經文.
比如某一些的經文會引用其他的經文.
他就將六萬多句的經文.
就將它plot出來.
由左至右.
左邊就是《創世記》開始.
然後到最後右邊就是《啟示錄》.
你就發現他這樣去到.
每逢有呼應的經文.
他就畫一個弧線.
呼應得越遠的就畫一個最大的弧線.
你就發現《創世記》頭兩章講到伊甸園的時候.
上帝與人同在這個描述的時候.
到最後的時候.

$^{41}$最後那兩章《啟示錄》.
也都講到上帝最終要與人同在.
最終要去到住在人的中間.
是有新天新地利道.
你就見到是有互相的呼應.
和互相的去到引用的.
所以你見到其實在講.
由始到終.
上帝是向著一個方向.
人類歷史是向著一個計劃.
就是將來這個新天新地會臨到.
中間就發生什麼事呢.
大概中間就是耶穌基督在十字架上死.
而墳墓裡面復活.
是舊約聖經的實現.
所以我們在看這個《宗末論》的問題的時候.
就不是一些等限事.
今天我們處理一段經文是新約裡面.
不是《啟示錄》.
因為我們還在處理《啟示錄》以外的.
一些講到《宗末》的經文.
就是《帖撒羅尼迦前書》第四章.
這一段經文有很多的討論.
特別在過去幾十年.
就講到這個被提的問題.
Rapture的問題.
就是會不會說這個世界的《宗末》.
來到之前.
就會有一個史無前例的七年的浩劫.
會臨到這個世界.
然後七年的浩劫之前.
這個主會先回來.
但是那個回來不是明顯回來的.
是秘密的回來.
Secret的Return.
然後他回來的時候就將教會.
來到所謂被提了.
將教會提走.
普世教會將他們提走.
然後這個浩劫才會展開.

$^{81}$到時候就會很嚴重.
這個就是《啟示錄》裡面.
七因七患.
裡面所講的浩劫.
七年浩劫完結了.
有很多的戰爭.
或者有很多的劫難等等.
基督才真正會再來.
有Second Coming.
那些細節我們不講了.
大概這個講法.
在華人的教會裡面.
是非常普遍的.
叫做災前被提論.
Pretributional Rapture.
你見到教會現在快要終末的時候.
基督會回來被提.
教會會被提.
但是那個被提是秘密的.
不是普世人會見到.
是基督徒會突然間消失.
你們坐在那裡的時候.
你突然間會消失.
只剩下一套衣服在那裡.
然後大災難七年就會臨到.
最後的時候.
基督才會第二次回來.
我記得年輕的時候.
大學的時候.
成長的教會.
不止一間.
都是講這個所謂時代論.
這個災前被提論.
尤其去到80年代90年代的時候.
我們比較成熟的那些.
我們知道80年代90年代的時候.
所謂世界大亂的時候.
你記得90年代.
有波斯灣的戰爭.
石油危機.

$^{121}$然後就有蘇共解體等等.
很多世界的大事就發生.
這種講法就是大行其道.
如果你記得的話.
有一些這樣的書籍.
叫做Left Behind.
在95年的時候開始出版.
用小說的形式來講這些災前被提.
然後有七年的大災難臨到.
你看到有些車的標誌.
In case of rupture happen的話.
這輛車就沒有人開.
不是說這輛車是自動車.
是會撞車.
因為你突然間消失.
In case of rupture.
你的Facebook就沒有人更新.
或者現在有些網站.
這個網站都進行了很久.
叫做Rupture Index.
去看一下.
什麼時候那個Index會高呢.
好像一些指數.
恆生指數或者股票指數.
主要的秘密的反應.
被提就會很快出現.
這裡值得我們去想一想.
會不會有時候.
我們有一些是既定了的時間表.
或者既定了的一些詮釋.
然後我們就將其他的經文.
Cut and paste 然後插進去.
剛才我們的表.
就給大家看到了.
有一些既定的中末的時間表.
我們先預設了.
然後我們看任何的經文.
有關中末的時候啟示錄等等的經文.
貼在前書第四章.
我們就將這些東西插進去.

$^{161}$這個講法或者這種論述.
有沒有搞清楚聖經裡面想講什麼呢.
我想問你.
這些時間表.
你覺得是Pro War 還是Anti War.
是好戰還是反戰.
你有沒有聽過最近香港的一首新歌.
我今早剛聽完.
作品的說話.
姜濤的那首歌.
是反戰的歌.
你覺得這些的備提論.
是好戰還是反戰.
是好戰.
我還記得八九十年代的時候.
有很多這樣的論述出來的時候.
波斯灣發生什麼事.
或者中東發生什麼事.
有很多這樣的論述的人.
就會說是要打的了.
沒辦法了.
一定要打的了.
直接應該要打了.
這樣的既定的時間表.
然後我們有一些既定的想法.
會不會影響我們怎樣看世事呢.
或者最重要的一樣東西.
是不是我們沒有搞清楚.
聖經想講什麼.
我不管你哪一排.
最重要是看清楚聖經背後的背景.
到底背景他想要提出一個什麼問題.
然後我們從那個方法來入手.
才能夠看得清楚.
啟示錄是這樣看.
貼心學家的前書也是這樣看.
這段所謂的備提經文.
我們一起看一下.
大兄姐們.
至於已睡了的人.

$^{201}$我們不願意你們不知道.
恐怕你們憂傷.
像那些沒有指望的人一樣.
你知道保羅寫那些新約的書信.
寫給教會的.
很多時候是一個回覆的電郵.
就是說你看到保羅回覆寫了什麼.
但是原本的電郵是沒有的.
所以你就用回覆的電郵來猜測.
或者去重構一下.
到底當時教會在問什麼問題.
當時貼心學家的教會.
可能在問這個問題就是.
在他們中間.
他們受了很多苦.
這個初生的教會.
有很多的基督徒.
因為有很多方面的打壓.
在當時的羅馬帝國的時候.
特別是一些有權勢的猶太人.
就去到貼心學家的城裡面.
對這些教會.
對這些基督徒.
就造成很多的影響.
很多的壓迫.
甚至有些人.
有些信徒是死了的.
是殉道死的.
因為福音的緣故.
因為信耶穌的緣故.
他們是殉道.
所以這群人你可以自然想到.
他們所有的教會.
所有的信徒都相信主會再來.
耶穌基督曾經從死裡面復活.
隨時祂會再來再回來.
那些貼心學家的信徒.
就可能會問.
到底基督榮耀再來的時候.
到底這些已死的信徒.

$^{241}$已經殉道了的信徒.
會不會錯過了.
基督的榮耀的國度呢.
因為他們有生之年的時候.
未必能夠等到.
看到耶穌基督的榮耀的再來.
所以他們會不會錯過了呢.
保羅的答覆就是.
他說一定是不會的.
十四節.
他說既然我們信了耶穌死了.
復活了.
那些已經在耶穌裡睡了的人.
上帝也必將他們與耶穌一同帶來.
我們照著主的話告訴你們一件事.
請你們記得.
一件事.
你們者活著還存留在主.
到主來臨的時候.
你們這班人.
絕不會在那已經睡了的人之先.
睡了的人當然是一個類比.
我們在基督裡面.
我們離開這個世界.
我們就算叫做睡了的人.
不代表我們沒有知覺.
離開這個世界.
我們會與主同在.
然後我們就在天上等待主的再來.
所以這裡不是說他們沒有知覺.
但是保羅主要想說的是什麼呢.
他們不單止這班在基督裡面要以睡了的信徒.
不單止他們不會錯過上帝基督榮耀的再來.
不單止不會.
反而會是主再來的時候.
主裡面先死的這班人.
他們會是第一批來去會見基督的.
不單止不會被收.
還會是第一批來去迎接基督的.
大家看下去.

$^{281}$因為召集另一法天使長的父聖一叫.
上帝的豪國一吹.
主必親自從天降臨.
那在基督裡面死了的人.
必先復活.
然後我們這些還活著還存留的人.
必與他們一同被提到雲裡.
在空中與主相會.
這樣我們就要和主永遠同在.
這裡很明顯見到.
就是說到天使長的聲音.
上帝豪國的聲音.
就是主會親自會來的.
但這個絕對不是秘密的.
沒可能解作秘密的.
因為是開放的.
是公開的.
不是秘密的回來.
神的豪國等等.
這裡有個字「被提」就是從這裡出來的.
所謂被提.
被提到底是什麼意思呢.
在空中與主相會.
如果你留意上下文的話.
就不是說被提到半空.
然後有些人消失了.
而是主裡面已死的人.
與我們還在世的人.
被提到半空.
去到會見這個正在下來的主.
就是中途我們上到去.
他降臨的時候.
中途我們上去會見他.
然後我們再下來.
上帝的國就會到了.
有個特別的字叫做相會.
是什麼字呢.
是一個專用的名詞.
專用的動詞.
形容當時羅馬帝國的軍王去打仗.

$^{321}$打仗完了他們海船回歸.
回到城的時候.
還沒進城之前.
那個城裡面的鄉親父老.
或者那些官長.
他們就會組織一些歡迎隊伍.
然後他們就會走到城外面.
將軍還沒回來.
還沒進城之前.
在中途與他相會.
相會之後.
歡迎了他.
然後就甲度歡呼地護送軍王.
或者那個將軍.
回到城裡面.
海船回歸.
就是一個這樣的景象.
那些達官貴人.
在古時在現代都是這樣的.
有些什麼大的人物駕到的時候.
總是會有一些很隆重的預備.
我們在芝加哥.
我們經常都會有這樣的情況.
是吧.
鎖齋突然間會封了.
突然間全部封了.
為什麼.
因為總統來.
是吧.
對嗎.
我記得剛來芝加哥的時候.
有些朋友很熟悉這裡的本地歷史.
就帶我到處去看.
他去到一個地方.
就叫做Rosevelle Road.
你知道Rosevelle Road.
有一段特別漂亮.
對嗎.
為什麼那段特別漂亮.
原來相傳在90年代的時候.

$^{361}$克林頓總統來芝加哥.
他是搭著軍用直升機.
來Grand Park降落的.
然後有歡迎隊伍去等他.
甲道歡呼他到了.
然後到了之前.
他還沒駕臨之前.
Rosevelle Road就是大裝修.
就是為了歡迎他.
是吧.
你住的區域之下是有錢人區.
是特別漂亮的.
當然這些達官貴人才有這樣的待遇.
我記得好幾年前我去完以色列.
回來的時候我跟老婆說.
我回來了.
我下機在奧希亞.
她就跟我說.
你自己坐車回來.
達官貴人才有這樣的待遇.
但羅馬帝國當時.
就是用一個這樣的方法.
來歡迎一些大官.
總統級的人來到城裡面.
In fact 這個字.
所謂會見這個字.
在《使徒行傳》第28章出現過.
當時保羅.
《使徒行傳》28章已經是最後一章.
說到什麼呢.
保羅要上告該殺.
要去到京城羅馬.
他就帶著福音去.
帶著耶穌基督的福音.
所以你看《使徒行傳》28章有很多地名.
不同的碼頭.
不同的港口.
他就去到一個地方.
叫做布丁尼.
在那裡上岸.

$^{401}$然後他就跟著黃色的那條線.
那條線就叫做 Via Appia.
是以前羅馬帝國最古老的那條.
所謂大道.
所謂彌敦道.
所謂的黃道.
就是一直會通向羅馬.
如果你看28章的話.
你就會看到它一步一步.
是接近羅馬.
28章第14節.
就這樣我們.
保羅說我們.
就跟路加一樣.
來到羅馬.
那裡的弟兄.
就是羅馬的弟兄.
聽見我們來的消息.
就去到阿比烏市和三館.
來迎接阿班提西斯.
我們.
保羅見了他們.
越發壯膽.
就進了羅馬城.
你看到右邊.
這個地圖右邊.
就是龐佩.
黃色那條線一直上去.
左邊左上角的就是羅馬.
三館和阿比烏市.
Form of Appia.
就是在左邊那裡.
你看到有個橙色的線.
就是那些羅馬的信徒.
聽見保羅來的時候.
他們就有一個迎接的隊伍.
歡迎的隊伍.
在中途跟保羅相會.
然後就escort他去到羅馬.
當然這個寫法.

$^{441}$這種描述.
就很有意思.
因為28章最後的時候.
講到福音.
福音就是宣告耶穌才是皇.
不是羅馬帝國的皇上.
他不是皇.
真正的皇.
是復活的基督.
所以就用這樣的方法去形容.
不是說保羅是大官.
而是保羅帶著的福音.
是一個大官.
因為福音講的是基督.
耶穌基督.
才是真正的君皇.
所以你看到.
在這裡講到.
空中遇見相會的時候.
其實相會這個詞.
應該可以更加好地翻譯.
應該叫做拜見的意思.
迎見.
緊見這樣的意思.
緊見哪一個.
就是基督.
他說已死的人必先復活.
以前的墳墓.
已死的人就葬在墳墓.
以前的墳墓永遠都是在城以外.
城牆以外.
所以你想知道.
某時某月某日.
那段時間.
那個城有多大.
你就去找那些墳墓.
找到墳墓你就知道.
大概那個城就在墳墓以內.
墳墓就在城牆以外.
所以你看到.

$^{481}$那個影像.
就是那個大官.
平面.
一直向城走近的時候.
他第一班會見到的人是什麼人.
首先復活的人.
當然就是那些已死的人.
在城外面.
然後和我們在城裡面的人.
就會組成一個.
歡迎的隊伍.
出去迎接這個君王.
最終目的不是走.
不是消失.
而是載他回來.
回到城裡面.
代表他的德成.
所以你用一個平面這樣看的時候.
你就將他90度反上來.
直線.
中面是怎樣看.
就是基督從天上降臨的時候.
下面我們那些已死的人.
在基督裡面已死的人.
和我們這班活著的人.
基督到來.
我們就會組成歡迎隊伍.
我們會上到空中.
去緊見這位皇上.
去拜見他.
與主相遇.
就是這樣的意思.
這些字眼.
都是當時帝國裡面.
有的外交的字眼.
其實說的是一件事.
本來就是借用當時帝國的語言.
或者他們的宣傳.
應用在基督身上.
就是告訴我們真正的帝王.

$^{521}$是主耶穌.
是一種顛覆的做法.
真正的帝王不是該殺.
是我們來迎接基督.
所以這裡不是說secret rapture.
這裡是說榮耀的基督降臨的時候.
我們去迎接他來到我們中間.
為什麼要說這些.
我剛才已經說了.
我說這些鐘末輪不是一些.
生命的注腳.
是告訴我們上帝在人類的計劃裡面.
是向著一個終點.
就是基督第二次再來榮耀的降臨.
當他降臨的時候.
我們準備好.
我們就去歡迎他.
在空中與之相遇.
然後escort基督來到.
然後新天新地就回到.
說鐘末.
說將來一定會發生的事情.
就是要我們live in light of that future.
為了將來這個榮耀的國度降臨.
我們而生活的.
任何今天我們所做的選擇.
任何今天我們所有的行動.
都應該要想一想將來的帝國.
將來新天新地來到的時候.
是一個怎樣的國度.
我們就是那班的國民.
我們的生活.
我們的選擇.
我們的價值.
就是in light of that future.
所以很多時候說到這些鐘末的時候.
聖經總是會提醒我們.
不只是一個天氣報告.
先寫好的報紙.
No.

$^{561}$他說與你今天生活有什麼關係.
剛剛過了父節.
甚至是32節.
過了哥林多前書.
又是說到復活的問題.
保羅說什麼.
他說從人的觀點看來.
我當日為了福音的緣故.
在異佛所和野獸去搏鬥.
受了很多苦.
對我來說沒有益處.
在人看來當然沒有益處.
為了什麼.
為什麼要為福音而捱苦.
付出這麼多.
如果死人沒有復活.
讓我們吃吃喝喝.
明天要死了.
因為我們是沒有攀往的.
我們只有here and now.
我們不如就吃吃喝喝.
快快樂樂的過完今天就算了.
這個起碼就是.
很多時候我們做人的方針.
珍惜今天.
我們吃吃喝喝.
然後我們明天已經沒有.
Beyond death的攀往.
沒有一些超越死亡的攀往.
算數.
他說不要被欺騙.
濫交朋友敗位品德.
你們要醒悟為善.
不要再犯罪.
因為有人不認識上帝.
我說者說.
是要你們羞愧.
寫給教會.
你想想居然要寫一些這樣的東西.
信了主這麼久的日子.

$^{601}$居然他不記得.
他還做了一個什麼.
讓我們吃吃喝喝.
明天要死了.
我們沒有live in light of future kingdom.
是要叫我們羞愧.
為什麼我們今天所做的決定.
我們所想的東西.
很多時候都沒有去到.
預備將來榮耀的國要降臨.
We are too earthbound.
只顧著我們現在吃吃喝喝.
不是說吃東西不要緊.
我不是這個意思.
而是挑戰我們.
到底我們認識這個終末的時候.
到底我們今天應該怎樣.
去做我們的決定.
怎樣影響我們生活的方向.
我們再處理多一段經文.
都是和這個終末要預備有關.
上一次我們說到馬太福音第24章.
24章你記得是一個list of signs.
一個列表.
耶穌說到說.
終末來到的時候會有什麼事情發生.
民要攻打民.
國要攻打國等等.
上一次說了背景.
已經是有什麼.
已經在第一世紀那裡.
文獻告訴我們.
已經是fulfill了.
不需要再等任何的sign.
不需要說現在有什麼大戰.
有什麼動盪.
一定是有些徵兆要實現.
其實那些徵兆早就實現了.
你再等下一個徵兆.
下一個徵兆就是人子要來的徵兆.

$^{641}$你再等的時候就太遲了.
24章37字.
所以耶穌說什麼呢.
他說.
羅亞的日子怎樣.
人子來臨也要怎樣.
在洪水以前的那些日子.
人朝上吃喝嫁娶.
直到羅亞進方周那日.
不知不覺洪水來了.
把他們都全衝去.
人子來臨也是這樣.
那時兩個人在田裡.
一個被接去一個被撇下.
兩個女人推磨.
一個被接去一個被撇下.
這裡不是在debate洪水到底是全球還是不是全球.
這個不是重點.
你看清楚上下文是在說什麼.
是人的預備.
人有沒有很好的預備.
不知不覺洪水就會來到.
吃喝嫁娶的時候.
洪水就會來到.
這裡又是另外一段所謂被提的經文.
因為你說兩個人在田裡.
一個被接去.
一個就留下.
兩個女人推磨.
當時的農耕社會.
有兩塊大石.
然後倒那些穀進去.
然後轉那個大石磨成為麵粉.
所以很多時候都是媽媽和女兒一起.
你轉180度她就轉180度.
然後你又轉180度她又轉180度.
是平時日常的生活.
但這裡你看清楚.
被接去是好事還是不好事.
被帶走.

$^{681}$洪水來到的時候被帶走是好事還是不好事.
不好事.
被留下是好事.
因為羅亞就是被留下.
所以想清楚這裡上下文說什麼.
不是那些secret rapture那些事情.
而是審判要來到的時候.
你有沒有準備好.
說了這些主在來.
也是告訴我們我們今天要去prepare.
prepare.
接受大家前書第四章第十八節.
這裡說什麼.
就是說完所謂被提.
就說基督榮耀降臨的時候.
我們準備好我們成為那個歡迎的隊伍.
讓基督來到的時候.
能夠和他一起進入一個國度裡面.
最後保羅四章十八節.
所以你們當用這些話彼此勸勉.
你記得他寫這段話是為什麼緣故.
因為他們中間有些在基督裡面已經過世的人.
他們很傷心很難過.
保羅說我們當然會傷心難過.
但是我們不應該沒有盼望沒有指望的傷心.
因為我們知道將來會發生什麼事.
Eschaton 綜合的時候會發生什麼事.
在主裡面我們認識的弟兄姊妹.
無論他是什麼原因離開了我們.
我們當然是傷心難過.
但在基督裡面有什麼是不同的.
我們可以彼此勸勉.
因為這班的信徒雖然比我們早走了.
但他們會是第一批看到榮耀的國度.
所以你想一下福音有多重要.
福音給我們一個盼望一個指望.
是超越今生今世苦難死亡的.
我們不是說我們有親人過去.
我們就不需要哭不需要憂愁.
不是這樣的意思.

$^{721}$要But not without hope.
保羅說不可以像沒有指望的人一樣.
你去使一個人認識主.
接受福音是一件多大的事情.
是一個永恆的事情.
另外他繼續說第五章第一節.
他說關於那日子那日期.
不用人寫信給你們.
因為你們自己明明知道.
主的日子來到會像賊在夜間突然來到一樣.
Again 又是一個比喻.
是突然來到.
不是說主耶穌是賊.
而是他來的時辰是你不會知道的.
突然來到一樣.
人正說平安穩定的時候.
禍患災禍必然臨到他們.
如同陣痛臨到懷胎的婦人一樣.
他們是絕逃脫不了的.
見主面.
要見上帝要去交帳.
這個是逃脫不了的.
不用人寫信給你們.
信了主的人.
特別是信了內的.
你應該知道.
為什麼還是為了今天我們慶祝.
明天就死的這種方式去生活.
Doesn't make sense.
如果你知道那個ascotine的時候.
你就會明白.
我們逃脫不了.
一定要見主面.
有兩個方法見主面.
第一就是基督再來.
我們大家一起見.
第二個方法就是你去見他.
你也可以突然去見他.
最近我們聽到不同的朋友的消息.
都是有類似這樣的情況.

$^{761}$就是你上不到的時候.
你就發現原來生命很脆弱.
你就是時候要結帳.
有人說得很對.
死亡見主面.
就和交稅一樣.
是逃不了的.
你可以拖.
但你永遠都不能夠逃避見主面.
所以每次講到主再來的時候.
不是一些豬腳的事.
而是講到你自己生命.
有些很重要的事.
要實行出來.
不是再拖延的.
你記得耶穌基督.
講完橄欖山.
講到之後.
他講什麼.
在馬太福音第24,25章.
就是講到比喻.
為什麼花這麼多時間講比喻.
為什麼不只是告訴我們天文台預測.
什麼時候主會回來就罷了.
一個又一個的比喻.
有三,四個比喻.
都是講到我們要去到預備好.
認見主.
到最後一刻.
到底我們有沒有準備好.
很多時候我們會聽到.
有些人說不用做這麼多的.
用不著為天國做這麼多的.
反正最後都要進入天堂.
反正我們都是恩典才能進入天堂.
信了就行了.
你口說信.
你就進入天國.
你做這麼多.
和我不做.

$^{801}$做有36,不做36.
如果你這樣看福音.
這樣看恩典的時候.
你就應該為你自己的生命.
打上一個很大的問號.
到底是不是真的明白.
什麼是福音 什麼是天國.
什麼是恩典.
因為剛才那種講法.
是一種廉價恩典.
cheap grace.
廉價恩典是救不了你的.
換句話說沒有行為的信心是救不了你的.
你不要自己欺騙自己.
本來這裡經常講這件事.
講到主在來的時候.
就是Don't deceive yourself.
不要自欺.
因為你的信心可能不是真正的信心.
那個是救不了你的.
所謂真正的信心是會走出來.
不是一種廉價的信心.
一定會有代價.
一定會肯付出的.
這種才是真正的信心.
不要自欺.
不要自欺不要以為我一定會見主面.
我一定會進入大明的天堂.
你現在做到像劇一樣.
你進去也是進入天堂.
我就是享受世界.
我今天也是進入天國.
跟你一樣.
所以你多吃虧.
我告訴你.
如果是一種廉價的信心.
起碼我本聖經裡面找不到這種廉價信心.
才可以救到你.
如果你的聖經找到的話.
我想你要查一下你的聖經版本.

$^{841}$可能你搞錯了.
你不是真的在看聖經.
或者你是斷章取義.
聖經裡面從頭到尾告訴我們.
信耶穌的.
進入天國的.
是要努力進去.
特別到末時來到的時候.
有很多人會失去信心.
因為生活艱難.
會失去信心.
會忘記了這件事.
就以為自己帶著廉價的恩典就可以進去.
不.
耶穌一次又一次的比喻.
說什麼.
十個童女的比喻.
你記得嗎.
十個都以為自己會進入賢直.
誰知家主來到的時候.
五個發現沒有預備.
他們不是那種廉價的.
他們能不能進去.
家主說什麼.
我從來不認識你.
僕人的比喻.
才剛的比喻.
五千,二千,一千.
從來沒有一個的說法.
說OK啦 都是靠恩典的.
有什麼所謂.
全部都是拜祭.
如果是拜祭的話.
就搞定了.
廉價的就搞定了.
就不用說這些比喻.
三個僕人的比喻.
拿到上帝給我們的資源的時候.
我們的寶物,我們的時間,我們的才華.
我們怎樣去使用的時候.

$^{881}$我們有沒有去到冒險.
為了去增加我們家主的資源.
去投資我們的時間.
投資我們的資源.
以致天國的事業.
天國的價值能夠越來越高.
童女有沒有拿到油.
有沒有聽了上帝的話.
又進行的.
肯付代價的來進行.
去走閘路的.
這些才是預備好.
可以進去天國.
才剛的比喻.
付出才能夠去見主面.
端不是一些cheap grace.
今天你聽到這個訊息.
是時候醒一醒.
特別是你信了久的.
是時候當頭棒.
想清楚.
到底你的信心是否真的救到你.
到底今天是否為了將來.
這個榮耀的天國.
榮耀的主降臨.
我們預備好.
可以花時間.
同心投地去祈禱.
如果你不是很清楚.
你是否真的信了主.
我希望在這個祈禱裡面.
我們能夠和你一起來接受主.
上主我們去感恩.
因為你是我們歷史當中.
世界當中真正的主宰.
主我們求你再一次來到.
讓我們看見我們對你的信心.
是否一個真的信心.
是否一個真正能夠去救我們.
真正能夠去進入天國.

$^{921}$真正為了將來的天國.
將來的價值.
將來的公平仁義.
真理而活.
主叫我們不要自欺.
求聖靈來提醒我們.
引進我們.
能夠進入真理.
真正認識你.
並且在我們生活當中.
主我們求你.
讓我們能夠預備好.
預備好你自己的在來.
無論是你自己榮耀的光臨.
我們能夠在空中與你相遇.
緊見我們真正的君主.
迎接你來到地上.
新天生的你到.
還是我們自己.
生命走到盡頭.
我們要去見你面.
主給我們你的恩典.
真真實實的來相信你.
為了你的天國.
為了你的福音.
我們願意去走閘路.
願意去真心去歸順你.
叫我們的信心.
在聽了你的道進行出來.
叫我們不要自欺.
讓基督成為我們的主.
我們感激到鳳主名禱.
我們以下一首詩歌給大家默契.
\newpage



\section{}
\label{sec:eLoBtjnY0Kg}
\textbf{比對研究 2: 必要, 發展, 方法   -  陳崇基牧師}
\newline
\newline
連結: \href{https://youtube.com/watch?v=eLoBtjnY0Kg}{\texttt{https://youtube.com/watch?v=eLoBtjnY0Kg}} ~~~~ 語音日期: 2022-05-05
\newline
\newline
\hyperref[sec:2oEQOesAhNs]{\small{< < < PREV SERMON < < <}}
~
\hyperref[sec:index_chronic]{\small{[返順時目]}}
~
\hyperref[sec:index_scriptual]{\small{[返順卷目]}}
~
\hyperref[sec:Wp2QtRm4GIw]{\small{> > > NEXT SERMON > > >}}
\newline
\newline
$^{1}$第二節也是講到比對研究.
但不是講上一次,上一次是有必要的.
因為帶著我們的文化背景.
發展是怎樣的呢?先講18世紀.
剛才我們講過18世紀之前.
對聖經背後的世界幾乎沒有認知.
都是帶著歐洲人的眼光去看.
18世紀之前能夠知道聖經舊約背後的世界.
有的資料其實不是很多.
有些甚麼資料我們可以在18世紀之前.
能夠知道希伯來聖經背後的世界.
例如希羅文化的作家埃及或巴比倫的古代文獻.
留給我們的例如Herodotus.
希羅多德,所謂的現代歷史之父.
Father of History.
他寫了很多東西,叫做History.
一大堆的寫作.
講述當時波斯帝國發生甚麼事,希臘發生甚麼事等等.
他有時會講述巴比倫的世界.
講述巴比倫的建築等等.
我們看到他所活的時代.
就是主前大概四五世紀.
這個Herodotus我們就可以有少少的資料.
他所寫的東西,告訴我們舊約的世界是怎樣.
或者Manifold,埃及的一個祭司.
他寫了很多東西.
講述他古埃及的種族,文化,文明.
幾遠古的文明,一路追溯一路寫出來.
他留下了很多東西.
今天有很多講到古代埃及的歷史.
都是用Manifold這個人留下的文獻重構出來.
這些是希羅或近東的作家.
猶太人的文獻都有.
有一個很出名的,第一世紀和耶穌時代.
在他時代的一個歷史學家.
猶太人叫做Josephus.
Josephus就是37至100AD.
主後37至100年的時候活著的人.
他寫了很多著作.
講述他自己猶太人文明的傳統.

$^{41}$他的文化.
所以他有一套書叫Antiquities of the Jews.
就是猶太古事.
中文翻譯出來.
是非常好的參考.
因為他由創世記開始講.
他講述創世記,摩西的背景等等.
另外一個猶太的歷史學家或哲學家.
就是Phylo.
Phylo也是公元前主前20年至主後50年的人.
也是耶穌時代的人.
他住在亞歷山大港.
就是Alexandria.
Alexandria在哪裡呢?.
就是在Rosetto附近.
Alexandria是當時最大的海港.
在東地中海裡.
他寫了很多猶太人的宗教信仰.
和希臘的哲學有甚麼的互動.
很多時候他們寫的就是和外邦人的世界.
來告訴他們猶太教的獨特.
猶太人的高尚.
所以在這些猶太文獻裡面.
有一些資料可以追溯希伯來聖經舊約聖經的背景.
但是很大的問題就是.
希羅的作家巴比倫的文獻.
和猶太人的文獻.
都是比較後期的.
後期寫的.
400BC至第一世紀.
並且我們知道他們寫的東西不是很準確.
往往帶著他們的有色眼鏡.
因為他們當時知道的東西不是太多.
所以用他們來做一些可靠的史料.
就會有些問題.
正因為這個緣故.
在18世紀之前.
最可靠的.
或者最多能夠知道近東的世界.
就只有舊約聖經.

$^{81}$舊約希伯來的聖經.
但是大家知道.
打開舊約聖經的時候.
很多時候舊約聖經不是去交代近東的歷史背景.
或者文化背景.
有很多東西是預設了.
作者和當時第一代讀者.
是預設了某一些文化.
或者某一些的習俗.
很多時候沒有跟我們解釋.
就算解釋也是解釋給當時初代的讀者.
所以說回巴比塔.
你看那個故事.
他沒有解釋為什麼會有巴比塔.
無端端就有巴比塔.
無理由下就會有巴比塔.
為什麼呢?因為他是預設了.
他第一代的讀者.
就是古代的以色列人.
他們知道近東的宗教禮儀.
或者宗教的建築物.
所以這樣的緣故.
十八世紀之前.
真的沒有什麼可以讓我們知道.
到底那個近東的世界是怎樣的.
直到十八世紀末.
就是1798年.
當時發生什麼事呢?.
當時最大的西方的大國.
這些大國的政治.
就是英國和法國.
英法很多時候大家去爭霸.
法國的拿破崙大家都聽過.
他在1798年的時候.
就佔領埃及和敘利亞等地.
他帶著的海兵.
就是軍兵.
有三萬六千個.
也帶著一大堆當時法國的.
那些所謂的科學家.

$^{121}$有167個.
他就由法國起行.
去所謂東征.
東征埃及地.
特別是利羅河三角洲.
所以他由法國出征.
就去到東地中海這個地方.
原因是什麼呢?.
原因就是當時英法.
大家之間就去到.
說誰是當時天下的霸主.
所以拿破崙就被派到.
這個地方去征服埃及.
征服敘利亞巴拉斯等地.
原因就是想.
阻截英國去印度.
當時經過非洲繞過去.
就是去到印度.
所以有很多商貿都是.
藉著這個方向去搭通.
在這裡上船.
經過陸路.
去到中東的東邊.
甚至是印度的西北邊.
所以因為他們在英法之間.
大家去到爭鬥霸權的時候.
拿破崙就去到這裡.
就希望阻截到英國.
阻截到英方去到稱霸.
在地中海的世界上稱霸.
所以拿破崙帶著他的軍兵和科學家.
就去到阿里山大廣.
然後他們攻佔到凱羅.
當時埃及是屬於哪個帝國呢?.
就是屬於土耳其的奧特曼帝國.
這個是1798年.
大概是1800年的時候.
當時他擊退當時守住埃及的阿拉伯人.
他擊退之後就拿了埃及的尼羅山國舟.
然後第二年他就嘗試佔據巴基斯坦地.

$^{161}$所以今天去巴基斯坦地.
去以色列.
如果去一個地方叫做約伯港.
那裡就有一些拿破崙的事跡可以說一下.
因為他由埃及一直上去.
拿了這個所謂世界的中心.
我們曾經說過巴基斯坦地.
其實就是歐亞和非洲中間的地橋.
上到這裡的時候.
霸佔了一個地方一會兒.
但最後英方就在另外一邊去擊退了法軍.
然後他們就佔領埃及.
這裡發生什麼事呢?.
在1799年的時候.
就是Bouchard.
這個Francois Bouchard.
其中一個的將軍.
拿破崙的將軍.
他就在Rosetta這個地方.
找到一塊石頭.
Rosetta我點了給你.
在亞歷山大港附近的一個港口.
他在這裡去拆一個舊的建築物的時候.
一個遠古的建築物的時候.
就發現一塊這樣的石頭.
這塊就叫做Rosetta Stone.
因為在Rosetta找到所以就叫Rosetta.
所以就找到這塊石頭的時候.
找到上面的一些命運.
有三段的命運.
正因為有三段命運.
所以這個將軍Bouchard.
他就明白這是一件很貴重的事情.
因為大家都知道.
尼羅河和埃及是一些很遠古的文明.
所以他就將這塊石頭運去Cairo.
因為法英之間有這樣的衝突的緣故.
所以最後英國就拿了這塊這麼寶貴的石頭.
叫做Rosetta Stone.
今天去大英博物館.

$^{201}$展館第一塊在你面前的.
就是這個Rosetta Stone.
說回這個阿波倫.
去到亞歷山大港和Cairo的時候.
他帶著167個科學家.
他們就做了很多這樣的研究.
和很多的整理.
將風景畫了.
然後將那些文物找到的東西.
這幅相就很出名.
就是羅阿波倫騎著馬去到這個獅身人面像.
這個Cairo這個地方.
附近有大金字塔.
這個港口就是Rosetta的港口.
你見到他們開了一些.
可能是木乃伊等等.
就有些所謂的記者.
就去記載這個阿波倫.
Description of Egypt.
有好幾冊.
這些所謂當時的科學的研究.
和survey等等.
找到這個Rosetta Stone.
為什麼這麼重要呢.
因為在18世紀之前.
沒有人知道那些古阿波倫的.
象形文字Hierarchy.
是怎麼翻譯的.
沒有人知道他講什麼.
那些象形文字我們看到很古怪.
有時候就好像畫一隻小鳥.
畫一條蛇等等.
這個Rosetta Stone.
其實連接到.
因為裡面講到一個王年.
就是一個Ptolemy V.
這個Thor.
勒密五世.
當時他在主前200至180年.
來去管理埃及地.

$^{241}$這個石頭就在1896年.
主前196年來去雕.
其實就是當時他的詔書.
就是這個Ptolemy V.
從父王Thor四世那裡.
就世襲了王位.
就雕了這個石頭.
這個石頭就是他們政府官方的詔書.
就去到說五世新的王上任.
他有多少的貢獻.
減稅等等.
上面最重要的地方.
就是這個Rosetta Stone.
有三個的語文來寫成.
因為當時196BC的時候.
當時的國際語言就是希臘文.
所以這個Rosetta Stone.
是當時的一些埃及的祭司.
他們去雕出來的.
講這個Ptolemy V的威水史.
上面那一部分已經破爛了.
這些只是重構圖.
可能會是這個樣子.
上面這裡就是一些埃及的象形文字.
Hierarchs.
中間就是當時平民百姓看的.
Demotic的埃及文.
草書來的.
埃及的草書.
下面54行比較完整的.
就是當時的希臘文.
留意這個時間是196BC.
七十四本已經產生了.
七十四本大概二百多BC的時候.
所以七十四本的Greek.
就跟這些Greek有相似的地方.
為什麼會三種語文寫成呢?.
其實三種語文都是講同一個的詔書.
就是這個Ptolemy V的威水史.
只不過用三種語言翻譯了出來.

$^{281}$正因為這個緣故.
當時的人懂Greek的時候.
他就漸漸可以拆解到這個所謂Hierarchs.
就是在金字塔找到這些象形文字.
在其他再遠古一些的石碑.
找到一些Hierarchs的象形文字.
就知道怎樣去解開埃及的論史.
正因為解開了這些象形文字.
所以我們知道埃及世界的背景.
譬如他們怎樣看創造.
他們怎樣去記載某一些歷史事件等等.
所以這個Rosetta Stone為什麼這麼出名呢?.
就是今天有一個翻譯的軟件.
也是叫Rosetta Stone.
這個現代的軟件.
你買這個軟件.
什麼語言你都懂.
譬如你要學西班牙文.
你就買Spanish的Rosetta Stone.
英文的Rosetta Stone.
或者法文的Rosetta Stone.
所以為什麼Rosetta這個字這麼出名呢?.
就是因為在當時找到這個東西.
然後拆解了第一個鑰匙.
去理解近東,特別是埃及的世界.
過了幾十年之後.
1840年的時候.
就有另一個人.
變得很出名.
這個人叫George Smith.
George Smith為什麼這麼出名呢?.
他本身是在銀行工作的.
他是做印銀紙的人.
所以他是很細心的.
但是他對近東的歷史背景.
是很有興趣的.
這是1840年,50年,60年代的時候.
當時西方的國家,英國特別.
在中東已經挖掘了很多東西.
拿回英國倫敦.

$^{321}$George Smith雖然在銀行工作.
但是他對近東的背景很有心.
所以經常去大英博物館.
下班的時候就去那裡做義工幫忙.
被那裡的大英博物館的主管賞識到他.
他覺得這個年輕人很聰明.
就請了他.
這個人原來是很聰明的.
他就是一個自己教自己.
自己學的近東學者.
當時找到這些泥板.
很多時候是從亞述古城找到的.
所以有一個新的學科叫做古亞述學.
所以找到很多泥板的時候.
就請了George Smith.
把這些泥板拼起來.
因為很多很多.
有一天,George Smith就在大英博物館.
找到一塊這樣的泥板.
他因為自學的緣故,他很聰明.
他就懂得看這些私營文字.
他就自己去拆解.
當他拿起這塊板的時候.
他就非常驚訝.
為什麼這麼驚訝呢?.
因為他一讀的時候就發現.
原來裡面就是Gilgamesh Epic裡面的洪水故事.
而且這個洪水故事有很多地方.
是跟創世記羅亞洪水故事非常相似.
所以當他發現的時候.
他就很興奮,他就跳到桌子上.
甚至有些報告說他立刻就脫衣服.
因為他很興奮.
興奮到一個地步.
二千多年,三千多年.
他是第一個去看近東洪水的故事.
所以就找到這塊泥板,叫做第十一塊泥板.
第十一塊泥板就是屬於Gilgamesh Epic裡面.
十二塊泥板裡面的第十一塊.
故事裡面就講到一個洪水的故事.

$^{361}$這個洪水故事跟創世記的洪水故事有多相似.
在一些學者的研究.
由70年代到現在,甚至幾年前.
一個猶太學者叫做Joshua Berman.
他這本書裡面都有引用這個平衡.
就將這個Gilgamesh Epic.
十一塊泥板裡面所看到的故事的脈絡.
和創世記的Genesis Flood Again的脈絡.
來做對比.
有17項的主題.
按照次序是完全一樣的.
幾乎一樣.
你看到第十一的line 14-19.
20-23,然後24-31.
所記載的都是這些.
我把它翻出來的時候.
你就會看到.
創世記洪水的故事和Gilgamesh的十一塊泥板.
11,神決定去降洪水.
6章第六節,還有第十四至十九行.
警告洪水故事的主角.
在創世記的故事就是羅亞.
但Gilgamesh的主角的名字不是羅亞.
而是Yunapesh.
第三,你繼續看下去.
整方舟,然後尺寸,兩者都有記.
然後繼續下去,主角聽命去整方舟.
命令他進入方舟.
然後進入方舟的時候帶同家人等等.
在兩個故事都是平衡的.
方舟的門關上也是平衡的.
有洪水的描述,大水的描述.
第七章和第九十六至十八章.
生態的破壞,雨水停止.
然後方舟停在山上.
第八章第四節就是阿拉納山群山.
然後Gilgamesh Epic就是另外一個山.
都是在北邊的.
繼續下去,主角開窗.
然後放料出去.

$^{401}$試測水到底是不是乾了.
然後主角出方舟獻祭.
神就聞到這些輕香的祭.
然後祝福主角.
你看到這些脈絡完完全全都是拼在一起的.
這個就非常令人驚訝的事.
因為我們平時只知道的就是羅亞洪水這個故事.
但我們卻不知道原來基加美修.
Gilgamesh Epic,Tablet 11.
講到這個人.
他所經過的洪水,他survive了.
他能夠生還.
這樣對比的時候你會發現.
原來那個故事的脈絡.
或者那個記載的方法形式.
都是非常相似的.
所以為什麼會有比對研究.
會造成一些爭議.
你就想到到底有什麼爭議.
我們下一節就繼續講這個爭議.
如果你喜歡我們的影片.
就記得要讚好,留言和分享吧.
\newpage



\section{1:1}
\label{sec:Wp2QtRm4GIw}
\textbf{Apostles' Creed \#1 (Genesis 1\_1) - Rev. Mark Chan}
\newline
\newline
連結: \href{https://youtube.com/watch?v=Wp2QtRm4GIw}{\texttt{https://youtube.com/watch?v=Wp2QtRm4GIw}} ~~~~ 語音日期: 2022-05-09
\newline
\newline
\hyperref[sec:eLoBtjnY0Kg]{\small{< < < PREV SERMON < < <}}
~
\hyperref[sec:index_chronic]{\small{[返順時目]}}
~
\hyperref[sec:index_scriptual]{\small{[返順卷目]}}
~
\hyperref[sec:io1AlESJGhY]{\small{> > > NEXT SERMON > > >}}
\newline
\newline
$^{1}$to a happy Mother's Day..
Let me begin by a word of prayer..
I don't know what kind of spiritual condition.
or what kind of emotional condition.
you're coming into the presence of God,.
but I want you to focus on God himself..
I want you to come in, not to see anyone else,.
but God himself, all right?.
And know that God is for you..
Know that God is present with you, all right?.
Let's pray..
Father, we want to praise you..
We want to come before you and bow our knees..
And I pray that you will be with us.
and help us to sense that you are with us,.
you are present with us..
I pray for each and every one of us..
Perhaps we have been fighting our battles.
in the past week, our past few weeks..
We are fighting our worries..
We are dealing with our challenges..
And sometimes our problems in our lives.
seem insurmountable..
It seems so big..
But we come before you, we know that you are bigger.
and you're biggest..
You're bigger than anything else..
And you got the whole wide world in your hand indeed..
And we pray that you will help us to put our trust in you.
and know that you are still in control..
You are still our king and our leader..
Pray that you will forgive our sins,.
our shortcomings, our weaknesses..
And we come before you anew,.
knowing that your grace is always new every single morning..
So we pray that you will renew our commitment to you,.
our trust in you..
No matter what barriers we are facing,.
we know that you are bigger than all those challenges..
And I pray that you're all powerful,.

$^{41}$you're all knowing and your attributes will encourage us.
to rely on you because you are worthy of our following..
So we pray that you will grant us grace.
to deal with our own depressed, unhappy,.
discomfort or discouragement..
And I pray that you will encourage each and every one of us.
in the spirit to continue to believe and trust in you..
Knowing that as we follow you,.
we need to deny ourselves, bear our crosses,.
just like what your son Jesus has done for all of us..
And I pray that we will be a community.
to encourage one another..
As we have a Bible study together today,.
we pray that in the word we'll grow together..
In the word we'll encourage one another.
to build up one another..
Build up your church for your namesake, for your glory..
Also wanna pray for the mothers in our midst..
I pray that you will bless and grant grace.
to every one of them..
Thank you for our mothers..
Thank you for your divine appointment.
that we have only one mother.
and the perfect mother for all of us..
And I pray that you will bless every one of them..
And thank you for using them, our mothers,.
to extend your love, to show your love..
Yes, we are all imperfect,.
but you have shown us the perfect love on the cross..
And thank you, Jesus, in Christ, we pray..
All right, today we're going to begin.
our new mini-series on the Apostles' Creek..
Now, James has told us we have our daily devotion.
on the Apostles' Creek..
So let's read this together..
All right, let's read this together.
and go through this together..
We have the whole Apostles' Creek on the screen..
Now I want you to read this verse with the mentality,.
with the motive to evaluate what we really believe, okay?.

$^{81}$Sometimes when we go to church for a long time.
and we've been to church for some time,.
we have gotten what we truly believe.
or we have truncated our own belief..
We might have thought or we have truncated our belief.
into just one or two things, and we make up our own idol..
And in doing that, we are not believing the gospel anymore..
And sometimes it is always good.
to reveal what we truly believe..
And the Apostles' Creek is written for such a purpose.
so that we'll do a daily or regular check.
of what you really believe, all right?.
Just like you go to the doctor..
Now tomorrow, I'll go to the doctor for a regular checkup,.
okay, nothing wrong with me..
You just have to do it once a year..
But the doctor will tell me, okay,.
your blood sugar level or your blood pressure.
or anything wrong with me,.
anything I need to pay attention to, right?.
So the Apostles' Creek is doing exactly.
this sort of purpose, to check your belief, okay?.
The Apostles' Creek, let me read the left column..
I believe in God, the Father Almighty,.
creator of heaven and earth..
I believe in Jesus Christ, His only Son, our Lord..
He was conceived by the power of the Holy Spirit.
and born of the Virgin Mary..
He suffered under Pontius Pilate,.
was crucified, died, and was buried..
He descended to the dead..
On the third day, He rose again..
He ascended into heaven.
and is seated at the right hand of the Father..
He will come again to judge the living and the dead..
I believe in the Holy Spirit, the Holy Universal Church,.
the chameleon of saints, the forgiveness of sins,.
the resurrection of the body, and life everlasting, amen..
This is the Apostles' Creek..
Now, it's a quick summary.

$^{121}$to tell you what our basic belief is..
Of course, it's not just that,.
but it's good enough summary.
to tell you what you truly believe,.
what are the central core beliefs, okay?.
Now, how many of you have been to Chicago for some time,.
for at least one year?.
How many of you?.
All right, no one is new to Chicago, right?.
Now, if someone bump into you.
and say maybe someone who's new to Chicago,.
first time in Chicago,.
and if that someone or your friend is a tourist.
bump into you and ask you,.
hey, how do you get to downtown?.
Can you tell him right away?.
Can you?.
(audience laughing).
Who cannot?.
Who can tell the someone.
that is the way from here to downtown?.
Who can say that?.
None of you?.
Jesus. (laughs).
How do you get to downtown from here, right?.
You can probably, immediately,.
you can answer the question, right?.
Now, if you cannot,.
you probably are deceiving yourself, right?.
You haven't been to Chicago, right?.
Perhaps you are being isolated for so long in COVID.
so you don't know the way, all right?.
Now, it's kind of like that, a puzzle squeak..
It's to tell us what the Bible is all about,.
what the gospel is all about..
No one can actually read the Bible.
from cover to cover in one sitting..
No one can understand every single detail in the Bible, okay?.
I have been teaching the Bible,.
I've been studying the Bible,.

$^{161}$I've been reading the Bible almost every single day, all right?.
There are still a lot of things that I don't understand,.
I don't comprehend, all right?.
Amen?.
Amen, right?.
So there are a lot of things.
that I still trying to figure out,.
and I'm struggling with the meaning,.
the significance of some of the Bible passages, right?.
We're gonna have a Bible study later..
We are going into the book of Genesis,.
the story of Isaac and Jacob, right?.
Right?.
Isaac and Jacob..
Some of the familiar stories,.
but there were a lot of details probably you don't know,.
or you have no idea that's in the Bible..
So no one can really understand the Bible.
or comprehend every single thing from cover to cover..
But the Apostles' Creed is to give us a quick summary.
to tell you and tell me what believing in Christ truly is..
All right?.
Now I want you to think about it..
When did you come to faith?.
When did you come to, you know,.
accept Jesus Christ as your Lord and Savior?.
When?.
All right, think about that..
How many years ago?.
What day?.
What year, what month, and what day?.
And what circumstances?.
That you come to know Christ..
All right?.
Now, how long have you been a Christian?.
It's a very crucial thing you need to ask yourself..
Okay?.
Now, think about how long you've been a Christian,.
how long you have known the gospel..
Have you really known the gospel?.

$^{201}$In the past two years, because of COVID,.
many of the young families in our church, okay,.
including some of you,.
and including some of those in the previous worship service,.
I haven't been able to, you know, see them face to face..
And I know maybe two years ago, all right,.
two years ago when the lockdown came,.
some of the young families actually had babies..
They had baby..
Or they had infant before, right before the lockdown..
All right?.
But when the lockdown came,.
I was not able to see them, right?.
Now, as the pandemic is easing out, hopefully,.
fading out a little bit,.
and then I see them again, right?.
Now, I did not expect the baby remain baby, right?.
It would be terrible..
After two years, if the newborn that I've seen.
before the lockdown two years ago,.
and then two years later, they remain baby, right?.
That would be terrible..
I need to call 911..
That would be so sad..
So after two years, when they show up again,.
when they come back to the church, right?.
No longer just seeing each other on Zoom or online..
They show up and they bring their young ones..
They are, what do you think?.
Should be walking, right?.
They should be taller, right?.
They should be tallest..
And indeed, they are now tallest, right?.
And for those who were two years old before the lockdown,.
now they are four years old, right?.
Now, it's a very natural thing to see growth, right?.
At least in the physical sense, right?.
At least in the physical sense, they got taller..
They got bigger, right?.
Before the lockdown, they could not talk..

$^{241}$Now, they at least, right?.
At least they can have a conversation, right?.
Now, it is the same thing, isn't it?.
For Christians, for us..
Rewind the tape, two years ago..
Where were you spiritually?.
Where were you in your relationship with God?.
How was it?.
Now, two years later, where are you?.
What do you truly believe?.
And have you been growing, right?.
Now, let me give you a very broad brush,.
a very just big picture of our faith journey, right?.
There are some aspects I'm going to tell you.
that if you know Christ,.
if you come to believe in the gospel,.
these aspects, you should see some growth, okay?.
The first one is your habit, all right?.
Your time management..
Before you came to Christ, you were your own boss, right?.
Every single day, 24/7, you were your own boss..
How do I use my time, right?.
How do I use every single minute that I have?.
After you know Christ, Christ is on the seat.
of ownership in our lives, right?.
He died on the cross to redeem us,.
and our life belongs to him..
Habit..
So two years ago, we went back, all right?.
Or how many years ago you come to know faith,.
come to know Christ?.
Your habits should have changed, right?.
In the, at least in the fundamental way, right?.
Before you know Christ, you used your time,.
you were your own boss, but after you know Christ,.
you need to relate to God..
You need to build your relationship with God..
You need to relate to him..
Have you been reading the word of God.
or getting a better understanding?.

$^{281}$Better, remember, the emphasis is better..
Have you been getting a better understanding.
of the word of God?.
Are you getting closer to God himself?.
Are you getting closer to the Lord and Savior?.
Week by week, we come to worship..
It's not because we, there's a checkbox.
in our weekly activities..
There's checkbox, okay, all right, I got it done, all right?.
And then I just put it aside..
We come to worship because we bow before the true Lord,.
the true King..
As believers, we are moving,.
should be moving in a progress..
From seeker, before you know Christ,.
you may be interested in the gospel,.
may be interested in the church stuff,.
and you come and you try to listen,.
try to understand what Christian faith is all about..
From seeker, you move on to a spiritual child,.
and then you move on to a spiritual adult,.
young adult or older adult..
And then you should move on to become a leader, right?.
That is your faith journey..
That is how your faith journey should be..
So at least there are fundamental change in your habit,.
in your time management, how do you use your life.
and your hours and your minute, right?.
Second, it's your relationship..
Relationship with other sisters and brothers in Christ,.
right?.
Interestingly, the Bible always tells us.
there were tons of verses that is about one and other..
It's about community..
When we come to know faith,.
we call each other brother and sisters.
or family members in Christ, right?.
And this family, this spiritual family,.
is above all the most important thing..
God will continue His redemptive work on earth.

$^{321}$through the church, not through the angels,.
not through the government, not through other agencies,.
but through the church, through the people of God..
That is us, right?.
It's not about numbers..
Remember in the first century,.
when you opened the pages of the New Testament,.
the churches that we encounter in Corinth, in Ephesus,.
in every single city in the Roman Empire,.
they were the small group church..
They're very, very small, tiny cell churches..
But God's plan is to use the church to extend His kingdom,.
extend His rule..
How do we tell the gospel to the rest of the world?.
God did not use angels..
God could have,.
but God choose to use His churches in every generation..
So relationship-wise,.
we should have some fundamental change..
There was some basic reconfiguration in our relationship.
when we come to know Christ..
Now, before we came to know Christ,.
church is nothing to us..
Maybe there's a cool thing, maybe there's a good thing,.
some sort of institute in the society.
that would help people to do better things or good things..
But when we come to know faith, we truly know faith,.
there should be some fundamental change.
in how we see the church..
Church becomes our priority..
Third, your emotional life..
There should be emotional growth in your own life as well..
We know Christ through the gospel..
We know that we are truly loved by God Himself..
We're truly loved by the unconditional grace of God..
Then we should learn, you know, in the progress,.
I am fully accepted by God Himself,.
and that is how I anchor my identity..
My value, my dignity, is not defined by my popularity..
It's not by how many friends I have..

$^{361}$It's not by how much money I earn..
It's not by my social status..
My emotional growth should go from.
before I came to know Christ, it was unsecure..
I'm trying to find things to anchor my identity..
The figures on your paycheck, how big is your house, right?.
How much things do you own?.
How many friends do you have?.
How popular are you on IG or on Facebook.
or anything like that?.
That was before Christ..
Now, after Christ, you came to know Christ,.
your emotional life should be always growing.
in light of the gospel, in light of how Christ loves us..
Because Christ loves us unconditionally,.
then I am a secure individual..
No matter how popular I am, no matter how unpopular I am,.
I know that I'm fully accepted by Christ, right?.
Now, since two years ago, because of the lockdown,.
because of the COVID, we have gone online, right?.
Now, you should know, right?.
Once you go online, your videos.
or whatever sermon you preach, right?.
Now it becomes something permanent, right?.
Now, I always get some comments on YouTube, right?.
Many, I would say, are positive..
But quite a few are negative..
Quite a few are negative, right?.
I've come to know this reality, right?.
This is the way it is, right?.
That's the world we live in, right?.
Now, in the beginning, you may feel very uncomfortable.
with some negative comments, right?.
Made by some users, made by some customers, right?.
Religious customers, that they have no commitment to you,.
right?.
They are somewhere else, I don't know where they are, right?.
They may be in Africa, they may be in Asia,.
they may be in Chicago, but they make some comments, right?.
They make some critical comments..

$^{401}$And then you kind of, you know,.
you don't feel good about it, all right?.
But if you are totally secure in Christ,.
all you do in your worship, in your service,.
in everything that you do,.
you are securely anchor your value,.
your identity, your dignity in Christ alone..
Then all these comments, will he, doesn't matter..
You know what I'm saying?.
People can criticize you,.
or people can give you some negative feedback..
But that should give you, you know,.
motives to improve even further, right?.
But you don't doubt yourself..
You don't doubt your own identity,.
because I am truly anchor in Christ himself, right?.
Emotional growth..
And then there's a purpose, fundamental change,.
in your purpose of your life..
Your general, overall orientation of your life..
Before you know Christ, you're all by yourself..
You're living your own life,.
according to your own agendas..
Now you belong to Christ..
And the New Testament continues to emphasize one fact,.
that Christ has redeemed you on the cross,.
with a great price, and you belong to him..
And your life and your mission is to glorify him,.
is to extend his kingdom, to share the gospel..
God is using you in your own way,.
that nobody else can replace you..
And you are executing God's own mission,.
to glorify God in your career, in your family,.
in your community, in your neighborhood..
That is a fundamental change in your purpose of your life..
Finally, your understanding..
Your understanding of the gospel..
You know, what Christianity is all about..
What does Jesus think it's all about?.
Before you know Christ, you probably were a seeker..

$^{441}$You were beginners..
After you came to know the gospel,.
you should have a better understanding,.
until you get a very solid understanding.
of what the gospel is..
Now remember what I asked you very beginning,.
the first question, the way in Chicago, right?.
The way in Chicago..
You have been in Chicago for some time,.
and you're familiar with the city, right?.
You can show other people the way..
Now, I wanna apply this question to the Christian faith..
If someone you bump into, know nothing about the gospel,.
and he asked you or she asked you,.
what this Christian thing is all about, right?.
How do I get there, right?.
How do I get to know Jesus?.
What is the way?.
Can you tell the person?.
Immediately, right away..
Now what I'm saying is,.
not only you should get a better understanding yourself,.
you should be able to explain it,.
to tell other people, right?.
If you cannot tell other people,.
hey, what is the way to God, right?.
Then you probably have some work to do, right?.
And I encourage you to do that, right?.
So, these are some of the aspects.
that we need to ask ourselves..
Our habit, how we use our time..
Relationship, we configure in such a way.
that we'll enhance the kingdom work of God..
The emotional growth for yourself,.
for your understanding of the gospel..
Are you making progress?.
Are you making any progress?.
Are you making any step forward.
in your faith, in your faith journey?.
The book of Hebrews, one of the New Testament letters.

$^{481}$written by a pastor to the church..
Most of them were Jewish Christians..
512, he says, "For though by this time,.
"you ought to be teachers.".
You need someone to teach you again and again.
the basic principles of the oracles of God..
"You need milk, not solid food..
"For everyone who lives on milk is unskilled.
"in the word of righteousness, for he is a child..
"But solid food is for the mature,.
"for those who have the powers of discernment,.
"trained by constant practice.
"to distinguish good from evil..
"Therefore, let us leave the elementary doctrine,.
"the basic doctrine of Christ, and go on to maturity,.
"not laying again a foundation of repentance.
"from that work and all that.".
Foundation of faith toward God..
We don't need the foundation anymore.
because you got the foundation already..
You see the progress here?.
Back then, you should be having milk, right?.
It's like a baby..
Well, you need a milk bottle to feed yourself..
You cannot have solid food..
But now, two years later at least, right?.
Two years later, you came to know Christ.
at least two years later before lockdown,.
before the pandemics..
Are you making progress?.
Or you're still having spiritual milk?.
You should be taking solid food, right?.
Otherwise, that would be strange, right?.
It would be strange..
Year after year, you're still staying.
at the elementary level..
We cannot do that..
The pastor says, "You need to make progress..
"Stop making excuse in your life.".
You know, we always have our own reasons..

$^{521}$Hey, I wanna postpone this or postpone that, right?.
God is our priority, and you should be making progress..
Now, as I told you, no one can really comprehend the Bible.
cover to cover in one sentence..
Nobody can..
Now, of course, to be honest,.
the simple message of the gospel is easy to understand,.
right?.
A three-year-old, four-year-old can understand..
God sold off the world, John 3:16..
Everyone can understand the elementary doctrine,.
but we are not called to stay the same, right?.
That is why we have the Apostles' Creek..
The Apostles' Creek is a quick summary,.
a useful summary of Christianity in 12 sentences,.
12 statements, and the creek itself,.
Apostles' Creek, has been around for 1,800 years..
Can you believe that?.
1,800 years, 1,800 years..
Even the Roman Empire did not last that long..
1,800 years, since AD 200-something,.
we had the most basic form of the Apostles' Creek..
Of course, it was not written by the first apostles..
It was the rule of faith laid down by the early church,.
the first century church,.
or the second generation of the Christians..
And because it goes back to a generation or two.
after the first apostles of Jesus,.
we call it Apostles' Creek..
Now, when we go through the Apostles' Creek,.
there were a little, here and there,.
I want you to understand what it is all about,.
Apostles' Creek..
Now, you can imagine in the first century,.
when the Christian church just started,.
they need some standard of faith.
for the Christians to grasp,.
so that they understand what the true gospel is all about..
If you don't have the right content, you cannot grow..
As I said, we truncate, or we reduce the faith.

$^{561}$into one or two things,.
you cannot grow strong as a Christian..
Couple weeks ago, Steve, remember that guy?.
Steve, he talk about the letter to the Ephesians..
Right at the beginning, he said,.
we need to love our Lord, our God,.
with all our heart and all our soul,.
and all our mind, right?.
Your thinking..
If your thinking is wrong, is truncated,.
you cannot relate to God, all right?.
Now, let's read the first sentence of the Apostles' Creek..
Okay?.
It says, "I believe in God, the Father Almighty,.
"Creator of heavens and earth.".
Did you hear that?.
I believe in God, the Father Almighty..
Let me unpack that..
First, faith, in the biblical sense,.
faith is not making something true..
Faith is not making something true..
Faith is affirming something that is true..
Okay?.
Okay, let's do some intellectual exercise..
Faith, believe in Christ, believe in the gospel..
Faith is not making something true..
Faith is affirming something that is already true..
How many of you believe in ghost?.
How many of you are afraid of ghost?.
All right, how many of you believe in ghost?.
You believe in ghost, right?.
I believe in ghost because the Bible tells me so, right?.
You believe in ghost,.
have you ever hunt one of those?.
You go out to the haunting house and hunt one of the ghost?.
No, right?.
It stays in the intellectual level..
You can believe in something, something that exists,.
the existence of something or someone,.
but you don't have the relationship..

$^{601}$It's all a matter of intellectual only..
That is how the apostles' quick begins..
I believe, I believe means not just intellectual checkbox..
I believe means I am believing into God..
I'm committing on myself to the relationship of God..
Now, how many of you know my wife?.
How many of you know my wife exists?.
She really does, she really does, okay?.
Even though sometimes I ignore it, all right?.
She really does exist, she really does exist..
But the fact that you believe my wife exists.
doesn't mean you are committed to her, right?.
For me, it's a totally different thing..
I believe she exists and I commit to her, right?.
You see the difference?.
Now, the Bible has two parts,.
the Old Testament and the New Testament, right?.
Have you ever wondered why there's a word testament there?.
The word testament comes from Greek word.
and that means covenant..
It is a lifelong commitment, right?.
It's the same thing as intellectually,.
I know my wife exists,.
but in terms of life, in terms of my life commitment,.
I am committing myself to her,.
not just affirming she exists only,.
but I'm committing to her, right?.
I believe in God, the Father Almighty,.
means that I'm not just affirming the existence of God..
There's no use, right?.
It doesn't impact your life..
I believe in God, the Almighty,.
means I'm committing myself to Him..
So the Christian faith is not about.
affirming the existence of God, right?.
I would be totally comfortable.
if someone check in our church,.
what do you believe in?.
What you don't believe in..
If he asks what I don't believe in,.

$^{641}$what the church does not affirm,.
I would be totally comfortable to answer and say,.
we do not affirm the existence of God..
Yes, we don't,.
but we affirm the God who exists and relate to us..
You see, there's a total different thing, right?.
Many people come to faith thinking,.
oh, this is an intellectual checkbox..
I believe in the existence of God,.
but it doesn't really impact your life..
You're not believing in the biblical sense..
You believe the God who commit himself to you..
His son died for you on the cross,.
and that's why you believe in God who exists..
It's a totally different thing.
than affirming the existence of God only..
You see the difference?.
Genesis 1.1..
In the beginning, God created the heavens and the earth..
You open the pages of the Bible..
Do you find it interesting?.
There's no intellectual argument in the Bible.
that you can find that would prove the existence of God..
The Bible doesn't waste time doing that..
The Bible simply assume,.
simply presuppose that only with God everything exists,.
only because of God himself, everything works..
On this presupposition, you have to choose.
whether you come to faith in this God, creator, or not..
That is what the Apostle Creeks is all about..
So you open the pages of the Bible..
You're always confronted with the God who created you,.
the God who exists,.
and he rules over the entire cosmos,.
entire life of yours..
And what is your relationship with God who created you?.
You owe your life to him..
How do you relate to this God?.
You're always confronted with this cosmic question..
Who is God to you?.

$^{681}$Now what does it mean to believe in God?.
I believe..
It means for us today that I will worship God and God alone..
I show that in my use of time,.
in my emotional status,.
in my configuration with the church, with other people..
You cannot be double-minded..
You gotta put your entire trust in him..
And secondly, what does it mean for us.
when we say I believe?.
I will continue to seek God's will.
in all the decisions in my life..
You ask God whether I should do this or do that..
You always have that in your mind..
Whether I should wake up today to worship him or not..
I'm always confronted with this question..
I set my orientation of my life to please God himself..
And thirdly, it means I'm making a commitment.
to know God and know him better every day..
I'm making a commitment to him..
That is what it means I believe, all right?.
Let's end with doing the Apostles' Creed..
Let's recite it..
Let's all stand together..
And then we'll close by prayer..
Let's all stand..
We're gonna read the English only,.
but you can do the other language if you want..
Apostles' Creed, let's do it together..
I believe in God, the Father Almighty,.
creator of heaven and earth..
I believe in Jesus Christ, his only son, our Lord..
He was conceived by the power of the Holy Spirit.
and born of the Virgin Mary..
He suffered under Pontius Pilate,.
was crucified, died, and was buried..
He descended to the dead..
On the third day, he rose again..
He ascended into heaven and is seated.
at the right hand of the Father..

$^{721}$He will come again to judge the living and the dead..
I believe in the Holy Spirit, the holy universal church,.
the chameleon of saints, the forgiveness of sins,.
the resurrection of the body, and the life everlasting..
Amen..
Let's all pray together..
Father, thank you for the opportunity to get to know you..
And I pray that every single one of us.
will make progress in our faith,.
to get to know you, to get to know your power..
I pray that you will help us to commit ourselves into you,.
as you have committed yourself into us..
You show us on the cross, you show us on the empty tomb.
that you are the Lord and Savior..
And I pray that our belief will continue to grow,.
how we see ourselves, how we see the world,.
how we see you as our king and our Savior..
So I pray that you will exercise your lordship in our lives..
Yes, we face a lot of different problems in our lives..
Our own personal growth, our own personal career,.
our relationship, our family, our future..
But we know that you above all, you are God alone..
And I pray that we will make our commitment into you.
and we'll grow into your word..
And I pray that as we gather to study your word,.
your spirit will continue to speak to us,.
to build us up, that we will be a solid Christians,.
we will be solid church that will please you.
and worship you..
In Christ's name we pray, amen..
Well, you may be seated..
\newpage



\section{}
\label{sec:io1AlESJGhY}
\textbf{Apostles' Creed \#2 (Psalm 19) - Rev. Mark Chan}
\newline
\newline
連結: \href{https://youtube.com/watch?v=io1AlESJGhY}{\texttt{https://youtube.com/watch?v=io1AlESJGhY}} ~~~~ 語音日期: 2022-05-16
\newline
\newline
\hyperref[sec:Wp2QtRm4GIw]{\small{< < < PREV SERMON < < <}}
~
\hyperref[sec:index_chronic]{\small{[返順時目]}}
~
\hyperref[sec:index_scriptual]{\small{[返順卷目]}}
~
\hyperref[sec:2oUtiGUEP5Q]{\small{> > > NEXT SERMON > > >}}
\newline
\newline
$^{1}$Hey, thanks Kelly..
Let me get my PowerPoint ready, alright..
Why don't we pray together?.
Father, we want to pray to you because you are the King, you are the Lord..
Even though sometimes we cannot see your hands, we know that you are still in control..
You are still the Lord of history, Lord of our lives..
That's why we come before you to worship you..
That's why no matter if it is sunshine or it is raining, we come to acknowledge your.
Lordship..
And we pray that you will help us to put our trust in you, to believe in you..
I pray for each and every one of us..
Perhaps we are facing some life challenges or crisis..
I pray that you will grant us the grace to see your authority, your power..
You are bigger than all the problems and challenges that we have..
So pray that we will put our faith in you, wholeheartedly trust you with our full heart,.
no reserve, because you are the Lord, you are God..
And pray that you will speak to us, pray that you will awaken our soul, our spirit, we will.
draw closer to you in your presence, in Christ and pray..
Well we are talking about Apostles Creed and this is a very, very ancient Creed, uh, lasts.
for over 1,800 years..
I mean, can you think of any movement in human history that will last so long, right?.
That will last so long, that would last for all the rise and fall of empires on the earth..
But Christianity, the Jesus movement, lasts for so long, right?.
And this Apostles Creed we're going to go through together, uh, we started last week,.
uh, lasts for almost 2,000 years, right?.
The Roman Empire, who used to persecute the early churches, they fell in the 5th century..
But the churches are continuing their movement and they still exist until today..
And that's why this is so valuable as we look at, uh, the Apostles Creed, uh, going back.
to almost 2,000 years..
And we said that, uh, the Apostles Creed believes in statements, uh, it begins with statements.
with "I believe this" and "I believe that.".
And we said that "I believe" does not mean intellectually agreeing upon some proposition..
In the Bible, "I believe" means "I am believing into God.".
That is to say, I commit myself to God..
I'm not just putting checks on check boxes and say, "Okay, I believe in this and I believe.
in that.".
It's not an intellectual proposition only..
It is "I am committing myself to the relationship with God.".
And last week, well, we said, um, "I believe in, uh, God, the Father Almighty, Creator.
of heaven and earth.".

$^{41}$And we clarify one thing in the Bible, faith is not making something true..
Faith is affirming something that is already true..
Okay, can you repeat after me?.
Faith is not making something true..
Faith is affirming something that is already true..
I believe..
I affirm something that is already true..
So that's why you don't find any passages in the Bible trying to prove to you the existence.
of God..
The Bible don't do that..
Instead, in the Bible, you find a God who exists..
You never find the Bible, never find the Bible or in Christian faith that "I believe in the.
existence of God.".
No, we don't affirm that..
Instead, we affirm, "I believe the God who exists.".
Can you repeat after me?.
We don't affirm the existence of God..
You say, "What?.
That's a church.".
Right?.
No, we don't..
We don't affirm the existence of God..
We believe in the God who exists..
That's a totally different thing..
The first statement, "I believe in the existence of God," is just an intellectual thing..
But if I believe or I'm believing into the God who exists, who relates to you, who loves.
you, who redeems you, that's a totally different ballgame, right?.
So that is what faith is all about..
And today, I would like to unpack the first statement here and using Psalm 119..
Sorry, Psalm 19..
Psalm 19..
Let me read the first six verses..
The heavens declare the glory of God..
The skies, the heavens, proclaim the work of His hands..
Day after day, they pour forth speech..
Night after night, they display knowledge..
There's no speech or language where their voice is not heard..
Their voice goes out into all the earth, worse to the ends of the earth..
In the heavens, He has pitched a tent for the sun..
God has done that..

$^{81}$God has pitched a tent for the sun, which is like a bridegroom coming forth from his.
pavilion, like a champion rejoicing to run his course..
It rises, talking about the sun, at one end of heavens and make a circuit to the other..
Nothing is hidden from the heat of the sun..
Now Psalm 19..
This is only the first third of it..
Psalm 19 is a creative rendition of Genesis 1, the seven days of creation..
And here the Psalm says, the creation of God, the entire cosmos, the entire universe is.
like a choir, like a symphony, like an orchestra..
Have you ever been to a symphony theater?.
Have you ever done that?.
It's like a symphony, singing day after day and night after night..
And this is some spiritual music playing, mysterious music playing, and you have to.
listen to it with years of faith..
This is not normal speech, not normal verbal language..
But there it is..
The creation is proclaiming something..
It's declaring something..
It's singing about God who exists..
God who is behind everything..
Some years ago, I was listening to a debate..
It happened in Great Britain, in UK..
So it was some years ago..
I was listening to a debate online between a Jewish Rabbi called Jonathan Sarick..
And then he was debating with the very famous, very well-known atheist, Richard Dawkins..
Both of them are UK people, so they did kind of a radio debate, online debate..
So I was listening to this, and it was a really, really heated debate..
And of course, the Jewish Rabbi, Jonathan Sarick, he was very famous..
He passed away not too long ago..
He was the chief Rabbi of Judaism in Great Britain..
And he was debating Dawkins..
And Dawkins was also a British scientist, and he is known for his atheist standpoint,.
right?.
That there's no God and everything is just matters, right?.
And everything is just coming together randomly, and that's why we have life..
And then, Sax was debating with him, and it was kind of a very heated, engaging conversation..
I remember one thing Jonathan said with Dawkins..
The Jewish Rabbi said, said to Richard, "Richard, religion is music, and you are tone deaf.".
And Dawkins said, "Yes, it is true..
I am tone deaf, but there's no music.".

$^{121}$I tried to sink in what the Rabbi is trying to say..
Faith is music, but if you're tone deaf, you cannot hear the music..
That's what Psalm 19 is trying to say..
The heaven is proclaiming, proclaiming the God who exists, the God who loves you, the.
God who made all of this, who is behind everything..
Faith is the ability to hear the music beneath the noise..
There are a lot of noises we hear today from our culture, from our social media..
It's a great thing we have great technology, but also create a lot of noises in our life,.
right?.
In order to hear God, you need to understand faith is the ability to hear, to penetrate,.
to hear the music beneath all this noises..
The Psalm says there is spiritual music of creation..
The sun rises in the east every morning and travels to the, to the west..
Okay, this is a normal, typical, ordinary language..
The Psalm is not teaching cosmology..
The Psalm is teaching us theology, talking about God..
In 1905, looking at the sun traveling from the east to the west, it's like a bridegroom.
coming forth from his tent or pavilion, like a champion rejoicing, running his course..
The pavilion behind it, it was the Hebrew word, hopa..
All right, can you repeat after me?.
Hopa..
Now, what is that?.
I always find it funny the Chinese would translate like that, all right?.
If you know Chinese, you know what I'm saying..
I don't know why, yeah..
But behind it, it was the Hebrew word, hopa..
And what is it?.
Now, if you Google that, if you know a little bit about Jewish wedding custom, even today,.
the Jewish couple would come under a tent with four corners, with the rabbi, with the.
officiant of the wedding..
So in the entire wedding ceremony, okay, the Jewish couple will come under that tent with.
the rabbi, with the officiant, and they will go through the entire ceremony under the tent..
The tent is called hopa..
All right, so they will officiate the wedding, the Jewish wedding, under that hopa..
And then when the ceremony concludes, the bridegroom with the bride will go out of the.
hopa..
And in Jewish tradition, it means a new creation, a new family is being created, right?.
A new relationship, new marriage has been launched..
So when they come out of the chamber or the pavilion, the hopa is a new creation..
So with that picture in mind, Psalm 19 is trying to say, every single morning when you.

$^{161}$wake up today, when you wake up this morning, the sun rises and will follow this course.
across the sky to the west..
And every day is a new creation of God..
Every day the sun comes out of the east..
It's like the bridegroom coming out of the hopa..
It is a new song for God..
It is a new song for his love, for his power..
So every day is a new day..
And then Psalm 19, not only use that kind of imagery, you know, in the ancient days,.
in the ancient Jewish society, or even today if we go to Israel, going to some traditional.
Palestinian villages, they will have weddings..
And like in the ancient days, all weddings happening in the village is a public event..
It's an open event..
Everybody in the village is invited..
Now it's totally different from today, right?.
Today if you are getting married or your friends are getting married, usually the people are.
being invited are very selective, right?.
Right?.
You have your list of guests..
You're trying to invite, you know, not as many as possible, right?.
It's always a private, selective event..
But in the ancient days, wedding is always a public event..
Everybody knows who is getting married next week..
Everyone is invited to the banquet..
And everyone is a witness..
Because nobody would be so stupid and say, "Who's getting married today?".
It's a totally public event, right?.
Psalm 19 is saying, "The music the creation is playing is a public event.".
There's no way you can miss it..
There's no way you cannot listen to the music..
If you put on hearing ears of faith, then you will pick up the spiritual music..
There's a guy in United States, very famous, very important person, Francis Collins..
He is the head of the NIH, National Institutes of Health, for some years..
And then he was also the head of the genetic research in the United States, the human general.
research..
In 2020, when the pandemic just started, in 2020, Francis Collins was awarded with the.
Templeton Prize..
Now the Templeton Prize was an award given to intellectuals who have made significant.
contributions to science and faith..
In that sense, these intellectuals who have received the prize every single year, in 2020,.

$^{201}$it was awarded to Collins..
He has made a lot of contribution to integrate science and faith, technology and faith..
And by the way, Jonathan Sachs, the UK Jewish Rabbi, he received the same prize in 2016..
So when Collins received that prize, he talked about how he came to faith..
Some years ago, he was a medical doctor and he visited one of his patients who was an.
old Christian lady on sickbed..
The lady was suffering enormously because of her conditions and she was dying..
And one day when Collins, as a medical doctor, he went into the room to talk to this lady..
This old lady asked Dr. Collins..
He asked him, "Doctor, I know I'm dying..
I'm going somewhere else..
I know where I'm going..
I know what I believe..
But doctor, what do you believe in?".
And it just dawned on him, as a medical doctor, this cosmic question confronted him and he.
didn't know how to answer..
And after that conversation, he took some time to really dig into what this faith is.
all about and why this lady who's dying should have no hope at all, right?.
Should have no hope at all..
But she is full of hope, right?.
And he, as a medical doctor, he couldn't do anything about it..
And he see this lady with hope beyond death..
And he asked himself, "What do I believe in?.
What is my belief?.
How do I answer that cosmic question, right?.
The life and death question..
What do you believe in?".
I will ask you, "What do you believe in?".
Right?.
You've been to church..
You probably grew up in a Christian family..
Good Christian family or bad Christian family..
It really doesn't matter..
You've probably been to church, good church or a bad church..
It really doesn't matter..
What do you believe in?.
And then Dr. Collins spent some, just a little bit, maybe a few years, and then he read C.S..
Lewis' Mere Christianity and find the faith..
So he came to Christ..
He came to believe in Jesus..

$^{241}$And he's still very vocal about what he believes..
He's an evangelical..
He's a believer..
Very good guy..
In one of his books, The Language of God, you see the DNA there, double helix..
He wrote this book..
And there's one paragraph I want to share with you..
He said, "While the long history of religious oppression and hypocrisy is profoundly sobering,.
the earnest seeker of truth, if you're earnest seeker of truth, must look beyond the behavior.
of flawed humans in order to find the truth..
Would you condemn an oak tree because its timber had been used to build battle ramps,.
weapons to kill other people?.
Would you condemn the tree?.
Would you blame the air for allowing lies to be transmitted through it?.
Would you judge Mozart's Magic Flute on the basis of poorly rehearsed performance by fifth.
graders?".
Have you ever been one of those?.
The fifth graders?.
"If you had never seen a real sunset over the Pacific, would you allow a tourist brochure.
as a substitute?.
Would you evaluate the power of romantic love solely in light of an abusive marriage next.
door?".
No..
No, you don't..
A real evaluation of truth of faith depends upon looking at clean, pure water, not at.
the rusty containers..
What do you believe in?.
Now I picked this paragraph because I really echo with some of the things that he said..
It's indeed really sad to see churches, some human institutes, corrupt by power, corrupt.
by sin..
They rise and they fall..
Sometimes when you hear news like that, you're just shaking your head, right?.
SMH, shaking your head..
You just cannot live with that..
You just cannot live with that, right?.
How could people do that, right?.
How could people do that?.
But it really strikes me..
He said about oppression, about hypocrisy, but he also invite us, if you're really earnest.
seekers of truth, you need to look beyond that, right?.

$^{281}$Sometimes we put our faith in human institutions..
Sometimes we put our faith in flawed humans, including your pastor, including your leader,.
including anyone that you look up to..
These people are as broken as we are..
Leaders and pastors, including myself, we all needed grace..
We all needed the love of God..
Every one of us..
Don't you want to put your faith in flawed humans?.
You put your faith, you put your trust in God and God alone..
Psalm 19..
We read the first six verses..
The first third is talking about the word of creation, the speech of creation, general.
revelation..
Everyone knows it..
It's universal..
And then seven to ten, it changed to words from the creator, words of creators, the special.
revelation, especially for you..
The words in the Bible, the words in scripture, those are the words from the Lord, the creator..
Then the last third is our response..
We as creature, how do we respond to the word of God?.
Let's read through it and we're going to go through the Apostles' Creed together and then.
we're going to sing a response song..
Okay?.
So Psalm 19, seven to ten..
The second half, second third..
The Lord of the Lord is perfect, reviving the soul..
The statue of the Lord are trustworthy, making wise the simple..
The precept of the Lord are right, giving joy to your heart..
The commands of the Lord are radiant, giving light to the eyes..
The fear of the Lord is pure, enduring forever..
The ordinances of the Lord are sure and altogether righteous..
They are more precious than gold, more precious than pure gold..
They're sweeter than honey, than honey from the comb..
The words from the Lord..
And then our response..
By your word, by them, is your servant wanted..
And keeping them there is great reward..
Who can discern his own errors, forgive my hidden faults..
Keep your servant also from willful sins..
May they not rule over me..

$^{321}$Then will I be blameless, innocent of great transgression..
What a humble prayer..
And he answered prayer with verse fourteen..
May the words of my mouth, the meditation of my heart, be pleasing in your sight, O.
Lord..
My rock and my redeemer..
Well, we're going to recite the Apostles' Creed together..
Let's do that..
Why don't we stand and do that, alright?.
If you will believe, we recite it together as if we will believe I'm believing into this..
Let's recite the Creed together..
I believe in God, the Father Almighty, creator of heaven and earth..
I believe in Jesus Christ, his only Son, our Lord..
He was conceived by the power of the Holy Spirit and born of the Virgin Mary..
He suffered under Pontius Pilate, was crucified, died, and was buried..
He descended to the dead..
On the third day, he rose again..
He ascended into heaven and is seated at the right hand of the Father..
He will come again to judge the living and the dead..
I believe in the Holy Spirit, the holy universal church, the communion of saints, the forgiveness.
of sins, the resurrection of the body, and the life everlasting..
Amen..
Rachel is going to sing, and we're going to join her together..
I believe in God, our Father..
I believe in Christ, the Son..
I believe in the Holy Spirit..
God, the truth, three in one..
I believe in the resurrection..
Heaven will rise again..
For I believe in the name of Jesus..
Let's do that together, I believe, from the beginning..
I believe in God, our Father..
I believe in Christ, the Son..
I believe in the Holy Spirit..
Our God is three in one..
I believe in the resurrection, that we will rise again..
I believe in the name of Jesus..
Let's all pray together..
Father, we praise you, we believe in you..
We are believing into you and putting our faith, our whole trust in you..

$^{361}$For you alone is trustworthy, not any man on this earth..
You are the God that we can depend upon..
You are the God we can rely..
For you are the almighty, you are the creator..
Let us praise you, let us hear the spiritual music of the universe..
Proclaiming your glory, proclaiming your awesomeness, your power..
Thank you for loving us..
And for the rest of the week, we are sent by your Holy Spirit to live for you..
And put our faith in you. And do your will to please you and worship you alone..
In Christ's name we pray, Amen..
\newpage



\section{詩篇}
\label{sec:2oUtiGUEP5Q}
\textbf{防禦性悲觀性樂觀 (詩篇 27 篇) - 朱莉傳道}
\newline
\newline
連結: \href{https://youtube.com/watch?v=2oUtiGUEP5Q}{\texttt{https://youtube.com/watch?v=2oUtiGUEP5Q}} ~~~~ 語音日期: 2022-05-16
\newline
\newline
\hyperref[sec:io1AlESJGhY]{\small{< < < PREV SERMON < < <}}
~
\hyperref[sec:index_chronic]{\small{[返順時目]}}
~
\hyperref[sec:index_scriptual]{\small{[返順卷目]}}
~
\hyperref[sec:V2NOXpx01xg]{\small{> > > NEXT SERMON > > >}}
\newline
\newline
$^{1}$大家好,大英姐妹平安!.
今天我想跟大家分享一個題目,如果不在香港的大英姐妹,不知道有沒有聽過,叫做「防禦性,悲觀性,樂觀」..
當然不是我創作的,是一個近年在香港的社交平台非常活躍的傳媒人,也是網路媒體100毛的創辦人林日曦先生所說的..
他在2019年12月在一個talk show叫做「悲但真」,sad but true的一個talk show裡面所說的其中一個金句,.
他說他本身是一個很悲觀的人,所以所有事情都是從最壞,最差來打算的,.
所以這種叫防禦性,悲觀,某種程度上就是期望管理,去保護自己..
但他再進一步,他說這不只是防禦性悲觀,還要防禦性,悲觀性,樂觀,他有個手勢會比較容易明白..
他所有事情想到最大鑊,最差,但又覺得這樣很正,所以就叫做防禦性,悲觀性,樂觀..
他用了一個例子,可以地道香港的例子,香港有一間茶餐廳叫澳洲牛奶公司,應該都很出名的了,.
當然他的食物是有一定的質素的,但比起他食物更加出名,是那裡的侍應哥哥招呼是出名差的,.
在很短的時間內,就會要求你下單,然後很短的時間內就送餐,然後你吃完很短的時間內,.
他就會清理桌子,等於逼你走,他就用了澳洲牛奶公司這個餐廳來做一個例子,.
是,食物是好吃的,但其實你知道他的侍應的招呼是惡劣的,但你相反覺得這樣是正的,就叫做防禦性,悲觀性,樂觀..
事實上,如果經過那個餐廳,長期是有一條人龍排隊等著去吃的,是一個很特別的例子..
今天想跟大家一起來看的是詩篇,27篇,我們也是一起來讀完這篇詩..
第一節:「也說是我的亮光,是我的拯救,我還怕誰呢?也說是我性命的保障,我還懼誰呢?.
那作惡的就是我的仇敵,前來吃我肉的時候,就半跌撲倒.雖有軍兵安然攻擊我,我的心也不害怕,雖然興起刀兵攻擊我,我必仍夠安穩..
有一件事,我尋求耶和華,我仍要尋求,就是一生一世住在耶和華的殿中,沾養他的榮美,在他的殿內求問..
因為我遭遇患難,他必暗暗地保守我,在他亭子裡,把我藏在他帳幕的隱密處,將我高舉在盤石上..
我現在得以昂首高過四面的仇敵,我要在他的帳幕裡歡言獻祭,我要唱詩歌頌耶和華..
耶和華,我用聲音呼籲的時候,求你垂聽,並求你連率我,應運我..
你說,你們當尋求我的面,那時我心向你說,耶和華,你的面我正要尋求..
不要向我掩面,不要發怒趕逐僕人,你向來是幫助我的,救我的神,不要丟掉我,也不要離棄我..
我父母離棄我,耶和華必收留我.耶和華,求你將你的道指教我,因我仇敵的緣故引導我走平坦的路..
求你不要把我交給敵人遂其所願,因為妄作見證和口吐空言的起來攻擊我..
我若不信在活人之地得見耶和華的恩惠,就早已喪膽了..
要等候耶和華,當壯膽,堅固你的心.我再說,要等候耶和華..
基督徒都很喜歡詩篇,但是我們讀詩篇往往會面對兩個問題..
第一,內容的方面,儘管我們都很愛詩篇,但是可能因為它是詩歌的題材,不及敘事文有故事性,或者一些議論文有邏輯性..
所以讀詩篇我們往往會怎樣呢?斬儉食,有些金句很得我的心的,我就記住它,.
但是就把它忽略了上下文的關係..
但這樣閱讀的話,有時我們會錯過了作者的重點..
第二方面,我們讀詩篇的時候會面對的問題是應用的問題..
怎樣將詩篇應用在生活上呢?特別是詩篇經常提到敵人這個元素..
誰作為現代人呢?誰是我的敵人呢?.
所以我們今天嘗試從這兩方面一起去解讀,嘗試去理解詩篇27篇對今時今日我們生活的處境有什麼提醒..
我們先講應用.這首詩一開始在第一節開宗明義已經告訴你什麼呢?.
第一節,他在講恐懼.第一節,我還怕誰呢?我還懼誰呢?.
當然有讀經的經驗的都知道這是什麼呢?是反問句,Rhetorical Question,.
他自問自答,其實代表什麼呢?他在面對恐懼..

$^{41}$第三節又提到,我的心也不害怕,我必營救安穩..
他從正面和反面在講同一樣東西,就是什麼呢?恐懼..
到最後14節又說,當壯膽,當堅固..
所以似乎整篇詩有一個主題,就是怎樣去面對恐懼..
恐懼什麼呢?我們看第二節..
如果直譯,將第二節直譯,會怎樣譯呢?可以這樣譯..
那作惡的Evil Doers攻擊我,就是我的對手和我的仇敵,前來吃我肉的時候,半跌,撲倒..
第三節,雖有軍兵,Army,安營攻擊我,雖然興起刀兵,War,攻擊我..
看不看到,詩人用詩篇,詩歌體裁是其中一個特色,他用很豐富的詞彙來形容他在面對什麼..
Evil Doers, Army, War,留意到什麼?很多軍事上的用詞來形容敵人..
而且,其實攻擊我這個字是重覆了幾次,所以是很針對性的..
再加上什麼?攻擊的對象是什麼?是我呀!.
我恐懼,詩人去恐懼,是當然的,面對這些事情..
所以我們讀詩篇的時候,我們很多時候會嘗試去想一個context,或者一個occasion,即是一個場景..
這裡一開始有提到作者是大衛,而且又用這麼多的軍事用詞,我們很容易就會聯想到是什麼呢?.
會不會是在某一個occasion裡面,大衛面對這些的真戰的時候,以至他寫的這樣的一篇詩呢?.
問題是,怎樣去應用呢?真戰似乎離我們很遙遠,至少在現在,戰爭對我們大部分來說仍然是很遙遠的..
而且是不是一定要好像大衛經歷這些的場面,我們才能夠認同他所說的感受呢?.
所以當我們應用詩篇的時候,我們會發覺有一個困難,我們identify不到敵人是誰..
作為現代人,我們怎樣relay到呢?.
難道說同事攻擊我,或者在工作的場景,同事來攻擊我,或者給我們什麼?.
老闆壓迫我們,朋友說我們壞話,或者股市跌穿了多少點,驗身報告不理想等等..
但我們都不會形容這些為敵人的,是不是?.
但是這幾年在疫情的期間,我們發覺有一個特殊的現象,特別是在香港..
因為香港大部分都是高樓大廈,我們出出入入都要搭電梯,是不是?.
以前在香港,你搭電梯,通常最喜歡站的位置是什麼呢?就是靠近電梯門,或者是按電梯的位置..
因為你通常不想一進去,到最進去,到時你要出來的時候,就要「麻煩借借,麻煩借借」這樣出來..
但是因為疫情的時候,我們發覺有個有趣的現象,就是那些人反過來,不是喜歡站近電梯門,不喜歡按電梯,喜歡站在裡面一點..
為什麼呢?避免按電梯,避免有機會接觸到一些病毒..
所以這件這樣的小事,其實已經證明了恐懼其實離我們不是很遠的..
而的確,這種恐懼,這種害怕是很實在的,因為敵人,病毒是無處不在的..
所以在這幾年,我們再讀這些詩篇的時候,我們發覺可能當恐懼變成我們日常生活的一部分的時候,.
大家生活在美國的時候,你現在開車經過油站,會不會看一看那個價錢?.
其實你未需要入油的,但是今天那個油價沒有那麼高,快點去入先這樣..
或者因為疫情,我們多留在家裡,大家會不會怕要磅重啊?等等等等..
原來最大的敵人不只是無形的病毒,很多時候在我們日常生活裡面,我們面對那個壓迫,.
無論是社會,或者我們的衣食住行,甚至我們的家人,或者我們身邊的朋友,我們身處的community,都會帶給我們很多的壓迫..
原來我們突然間發現,原來這類的詩篇不是離我們很遠..
問題來了,詩人有個modal answer,他說我不害怕,面對那麼多的恐懼..
詩人說我不害怕,為什麼呢?我們來嘗試整篇詩這樣去看,去理解這篇詩的內容..

$^{81}$我們剛才大概講了一至三節,我們去到第四節,你會留意到詩人的筆鋒是轉了的..
無端端第四節提一些無關痛癢的事,第四節,有一件事,我尋求耶和華,我仍要尋求,.
就是一生一世住在耶和華的殿中,瞻仰他的榮美,在他的殿裡求民..
無端端提聖殿,好像一個平行時空一樣,前面講的是打仗,講的是一些兵凶戰危,.
是要和敵人對峙的危險場面,第四節突然間非常靜態,彷彿有些聖詩在後面響起..
所以有些解經家嘗試這樣解釋,會不會第一到第三節是講一些過去的經歷,.
詩人已經面對戰勝,已經經過了這些戰役,所以現在是得勝,在聖殿裡敬拜,去感恩..
但奇怪的是,如果我們再讀下去,我們會發現敵人又再出現,所以到底是怎樣去理解呢?.
亦都有人理解為在危機裡,在一個很大的危險裡,詩人怎樣呢?.
躲藏在聖殿,去到第四節,躲藏在聖殿尋求保護,有些人這樣說,聖殿很躲藏的,.
是不是純粹是躲藏在聖殿,所以第四節這樣寫呢?.
我們再仔細去看第四節,第四節開宗明義是這樣說,.
「有一件事」,或者可以這樣譯,「唯有這一件事」,將焦點聚焦在一件事身上,.
「一件」是什麼事呢?「尋求」,仍要尋求,從過去到現在都沒有間斷地渴求尋求一件事,.
是不是兵凶戰危,所以要尋求保護,是不是希望祭司將我收起來,.
以致我不需要面對敵人的威嚇呢?是不是這麼簡單呢?.
不是,這裡說有一件事,那件事是什麼呢?.
是一生一世住在耶和華的殿裡面,有兩個目的,去瞻仰,去保護神的榮美,.
雖然我們不明白,其實怎樣看神的榮美在聖殿裡面,.
如果你熟悉聖經的話,你知道在聖殿裡面是沒有聖像的,.
不像我們在教堂裡面,我們會看到很多的石瓶,很多的雕像,很多的像藝術品一樣的畫等等,.
但是在聖殿裡面是沒有的,所以要看什麼呢?.
但是要保護,要看神的榮美,第二樣東西,在聖殿裡面的目的,.
在他的殿裡面求問,而且那種形容是怎樣呢?.
形容得好像祭司檢查祭物是否合乎標準,那麼小心求證地去尋求神,.
所以簡單來說,這裡說的是什麼呢?就是他極度,斯人極度渴求得到神,.
要尋求神,要尋求神的面,他的榮美可以是他的說話,可以是他的作為,.
總之最重要是什麼?神是我需要去尋求的對象,.
所以有一件事,面對恐懼,面對甚至無孔不入的一種恐懼,.
最需要的是什麼呢?不是求神去消滅敵人,不是求神來保護,.
不是說這些不重要,但是有一件事更加重要的是什麼呢?.
是懼怕應該懼怕的,和尋求應該尋求的,.
我們要fear the right fear,我們要找出什麼是最重要呢?.
這裡告訴我們,原來尋求神,懼怕神是應該最重要的,.
所以懼怕應該懼怕的,尋求應該尋求的,.
是斯人告訴我們面對恐懼的不義法門..
我們繼續看下去第五節:.
「因為我遭遇患難,他必暗暗地保守我,.
在他亭子裡把我藏在他帳幕的隱密處,將我高舉在盤石上.」.
我再說,當遭遇患難,其實我們期望什麼呢?.

$^{121}$我們期望神橫掃千軍,宣告我們勝利,.
我們希望神清除敵人,那些惡人遭報,.
但起碼在詩篇27篇不是這樣說,.
他反而是求你暗暗地保守我,將我藏在帳幕的隱密處,.
這裡「亭子」是很有意思的一個字,.
這個「亭子」是比較特別的字,.
一方面它平衡帳幕的隱密處,即是強調那種安全,.
但更加有趣的是什麼呢?.
「亭子」的相關字是住棚節的棚,.
住棚節有印象的話,是猶太人非常非常重要的節日,.
而住棚節的特色是什麼呢?.
儘管可能已經進入迦南地,或者已經安定在應許地,.
在住棚節的時候,他們要住帳幕,住棚,.
在住棚節住七天,紀念出埃及的日子,.
一方面當然是紀念神怎樣帶領他們出埃及,.
怎樣成功擊敗當時的霸權法老,.
讓他們能夠成功出埃及,.
但這個模式是紀念他們曾經遭遇的困苦,.
來讓他們提醒神的勝利,.
這是一個很特別的模式,.
所以不是完全撇除了患難,.
反而我們要記得那個患難,.
以致我們記得神的特性..
第六節:.
「現在我得以昂首高過四面的仇敵,.
我要在他的帳幕裡歡言獻祭,.
我要唱詩歌頌耶和華.」.
不是說仇敵全軍覆滅,.
還怎樣形容仇敵?.
那些仇敵是四面的仇敵,.
是被包圍的,.
但重點不是那些四面包圍我的仇敵,.
重點是我得以昂首,.
是很典型在古代裡的Honor and Shame,.
無論敵人的數字,形勢是多麼險峻,多麼多,.
但他們不能夠包圍我,.
因為我仍然能夠昂首,.
反過來我昂首即是他們降臂,.
我是被圍困,.
但我的頭部是最高的,.

$^{161}$我仍然可以昂首..
我覺得這是一個很重要的,.
基督徒需要知道,.
基督徒不是鴕鳥,.
我們不是阿Q,.
我們不是一定要正能量,.
或者我們一定要盲目地樂觀,.
總之我躲在聖殿裡一切安好,.
其他東西什麼都看不到,.
不是這樣的..
最近這幾年香港很多的一些研究,.
或者是世界性的一些民調,.
都形容香港是一個很悲觀的城市,.
香港人都活在很多的失望,.
很多的絕望裡面,.
甚至其中一個民調,.
在2021年前是這樣做的,.
就說香港是全球第二最悲觀的城市,.
我生活在一個這樣的環境,.
一個這樣的城市裡面,.
我不去悲觀,.
或者不去對前景有疑慮,.
是很奇怪的事情,.
基督徒都不能夠倖免去面對這些事情,.
但是基督徒需要知道的是什麼呢?.
我可以昂首,.
還未完的,.
我可以在他的帳幕裡面,.
寬然獻祭,.
那個寬然是什麼呢?.
是摧國那種的寬然,.
是慶典,.
是打聖像才會有Trumpet,.
那種的寬然來獻祭,.
所以我再說,.
受敵仍然在這裡,.
但是我的態度,.
我怎樣去面對,.
已經不同了,.
為什麼要提聖殿呢?.

$^{201}$除了你說聖殿是象徵神的同在等等,.
為什麼要提聖殿在一個對壘的形式裡面,.
我和敵人在一個對壘的場景裡面,.
為什麼無端端要提聖殿呢?.
我相信Urban Voice的UV的大家,.
應該都不會陌生,.
因為Mark Sir這方面的專家,.
可能都會有提過,.
在舊約的場景,.
在Asia Near East,.
在東亞的歷史文化背景下,.
雙方對戰的勝利取決於什麼呢?.
不是取決於他們有多大的兵力,.
或者武器有多厲害,.
或者他們的策略有多厲害,.
取決於他們背後的神明,.
哪一個更加有能力,.
所以當對戰的時候,.
他們贏的時候,.
是代表什麼呢?.
是代表他們背後的神明,.
比對方的神明更加有能力,.
更加可以賜給他所守護的對象,.
有勝利..
其實有一個例子,.
大家可能會更加熟悉的,.
在《撒慕爾記》上提過,.
當時以色列人和肥里士人爭戰,.
結果肥里士人贏了是怎樣呢?.
是將約桂魯去,.
然後擠在哪裡呢?.
擠在肥里士人所敬拜的大滾廟裡面,.
去宣告勝利..
所以其實是一種神明的對決..
為什麼這裡要提聖殿呢?.
沒錯,.
受敵固然是值得恐懼的,.
我無意去降低,.
我們面對現在各種狀況,.
帶給我們的恐懼,.

$^{241}$不是不斷說不要怕,.
不要怕,.
其實是有用的,.
但是我們需要知道的是什麼呢?.
我們需要知道的是,.
是懼怕,.
應該懼怕的,.
不要因為四面的仇敵,.
我們被四面的仇敵所圍繞,.
被環視,.
然後我們就忘記了,.
其實神有祂的保護和祂的特性,.
更大的恐懼是什麼呢?.
更值得我們害怕的是什麼呢?.
當我們面對每一天的各種壓迫的時候,.
我們忘記了我們背後的神,.
其實比一切都大,.
我們忘記了我們最需要尋求的是什麼呢?.
是耶穌,.
或者我們所信奉的上帝,.
祂的特性,.
因為這場仗,.
關祂事的,.
是屬於祂自己的征戰來的,.
所以為什麼要提聖殿呢?.
因為征戰屬乎耶和華,.
而這種征戰,.
也是在我們每一天生活裡面,.
當我們每一天面對各種各樣.
無以名狀的恐懼的時候,.
到底誰比較大?.
我們需要恐懼的到底是誰?.
如果是這樣的話,.
我們讀下去七至十二節,.
我們就會明白私人的重點,.
我們就會明白為什麼前面好像已經有特性,.
後面又再出現祈求的片段呢?.
因為私人是用另一個角度,.
再一次強調我們要懼怕,.
應該懼怕的,.

$^{281}$我們要尋求,.
應該尋求的,.
第七節,.
「耶和華,我用聲音呼籲的時候,.
求你垂聽,.
並求你連率我,.
應允我,.
你說你們當尋求我的面,」.
重複前面所說的話,.
「那時我心向你說,耶和華,.
你的面我正要尋求.」.
這裡轉了一個角度,.
前面是用第三個人的,.
是用「他」,.
現在是第二個人的,.
是「你」,.
所以我要尋求你的面,.
直接向神說,.
「我求你讓我尋求你的面.」.
第九到第十節,.
「不要向我掩面,.
不要發怒趕逐僕人,.
你向來是幫助我的,.
救我的神,.
不要刁棄我,.
不要離棄我.」.
四個「不要」,.
但其實說同一樣東西的,.
就是「不要讓我找不到你,.
不要離開我,.
我要找到你,.
我要黏著你.」.
甚至第十節說什麼?.
「我父母離棄我,.
耶和華必收留我.」.
即是人世間最重要的關係,.
父母和子女都比不上.
神和人之間的關係..
在中國近代的歷史裡面,.
曾經有一個時期,.

$^{321}$當父母和子女都可以成為敵人,.
一個瘋狂,.
一個人性最黑暗的年代,.
但儘管是這樣,.
耶和華仍然是最可靠的那個..
第十一,十二節,.
「耶和華,求你將你的道指教我,.
因我仇敵的緣故,.
引導我走平坦的路,.
求你不要把我交給敵人,.
誰其所願,.
因為妄作見證的,.
以及口吐空言的,.
起來攻擊我.」.
所以,其實很諷刺的,.
敵人從頭到尾都沒有消失不得已,.
還似乎是越來越猖狂,.
那些妄作見證,.
口吐空言的,.
再詳細去說明,.
這些敵人可以怎樣來攻擊我,.
正正因為有敵人的存在,.
正正因為這些敵人無休止的攻擊,.
反而驅使私人緊緊抓住神,.
「求你將你的道指教我,.
因為我敵人的緣故,.
反而是因為他,.
on behalf of敵人的緣故,.
引導我走易路,.
因為那些妄作見證的,.
來攻擊我,.
所以求你不要將我交給敵人.」.
最後,13,14節,.
「我若不信在活人之地得見耶和華的恩惠,.
就早已喪膽了,.
要等候耶和華,.
當壯膽,.
堅固你的心.」.
我再說,要等候耶和華..
第13節,其實希伯來文的意思是不完整的,.

$^{361}$「就早已喪膽了」,.
其實是應該沒有這句的,.
或者你可以看到是有「…」的,.
你看英文,日本就更加清楚,.
沒有這句「就早已喪膽了」..
所以可以怎樣理解呢?.
可以理解為一個反問句,.
一個Rhetorical Question,.
即是前面一開始詩篇是用一個反問句開始,.
去到後面用一個反問句來結束,.
可以怎樣翻譯呢?.
「我怎會不信在活人之地會得見耶和華的恩惠?」.
或者再正面一點來說,.
「我確實相信…」.
很多英文,日本就用這個方法來翻譯,.
「我確實相信…」.
即是「I believe我會見到耶和華的恩惠」..
所以最後的總結是甚麼?.
第十四節,「要等候耶和華」,.
說兩次「要等候耶和華」,.
為甚麼要等候呢?.
因為勝利還未到,.
我只是相信我將會在活人之地見到耶和華的恩惠,.
「我相信我將會見到」,.
而不是勝利還未到,.
不是我們已經可以得勝,.
我們可以舉起升旗,.
我們可以向敵人跨勝,.
因為敵人仍然在,.
但重點是甚麼呢?.
既然征戰是屬乎耶和華的,.
所以我現在相信,.
我將會看見,.
所以我現在已經可以聲稱勝利,.
雖然勝利還未到,.
所以要怎樣呢?.
壯膽,等候,堅固地等候最終得勝的來臨..
我剛才一開始提到,.
這個防禦性,悲觀性樂觀,.
是林日曦先生在2019年12月,.

$^{401}$2019年是香港經歷很風起雲湧的一年,.
在那一年12月,.
他舉辦了一個講座,.
說了這句話,.
在2020年7月,.
他又舉辦了另一個講座,.
那個講座叫「半厭世」,.
在2020年7月,.
他舉辦一個叫「半厭世」的講座,.
這次他來個再進階版,.
他不是說防禦性,悲觀性樂觀這麼簡單,.
他再加一件事,.
他說防禦性,悲觀性樂觀性宏觀,.
什麼意思呢?.
是,悲觀性樂觀已經很厲害了,.
很多東西很差,.
去到最盡,去到最差的時候,.
但我們仍然覺得「嘩!正喎!」.
為什麼我們會覺得「正」呢?.
因為我們要宏觀地看,.
因為事情要拉闊地去看,.
即使事情去到最壞,.
但當我們拉闊地看的時候,.
我們可以想像下,.
敵人,對方最終會有什麼下場和後果呢?.
而我們將來又會有什麼情況呢?.
當宏觀地看的時候,.
可能我們都能夠釋懷,.
可能我們今天面對很多令我們值得,.
令我們悲觀的事情,.
我們不單可以樂觀地看,.
我們還可以宏觀地看,.
而這些恐懼在這些因素底下,.
可能我們發覺不值得再恐懼,.
我再說,恐懼是真實的,.
特別是我們每一天要面對的事情,.
恐懼也不是我們要去壓制這種情緒,.
但當恐懼變成常態,.
我們不要讓它過度,.
不要讓它所剩,.

$^{441}$要和它共存,.
甚至最重要的是什麼呢?.
是有一件事要不斷提醒我們自己,.
或者我們都可以好像詩人那樣,.
宣告這件事,.
「耶和華是我的亮光,是我的拯救,.
我還怕誰呢?耶和華是我的性命的保障,.
我還懼誰呢?」.
但願這篇詩篇能夠成為我們今天,.
以致我們每一天面對各種各樣不同狀況的生活裡面,.
一種提醒,也是一種力量,.
我們一起祈禱,.
儘管敵人四面環繞我們,.
儘管我們對未來有很多未知,.
有很多說不出來的恐懼,.
因為你沒有離開過我們,.
而且你最終是那個德性的上帝,.
主求你幫助我們,.
讓我們相信我們必然會見到你自己的恩惠,.
以致我們願意壯膽地等候,.
甚至在我們日常生活裡面,.
我們堅持我們去學習,.
懼怕應該懼怕的,.
尋求應該尋求的,.
以致我們可以昂首來面對每一天的挑戰,.
求你祝福我們,.
奉耶穌基督的聖名,.
阿們..
阿們.
\newpage



\section{11:1-44}
\label{sec:V2NOXpx01xg}
\textbf{An Emotive Lord (John 11\_1-44) - Brian Ng}
\newline
\newline
連結: \href{https://youtube.com/watch?v=V2NOXpx01xg}{\texttt{https://youtube.com/watch?v=V2NOXpx01xg}} ~~~~ 語音日期: 2022-05-24
\newline
\newline
\hyperref[sec:2oUtiGUEP5Q]{\small{< < < PREV SERMON < < <}}
~
\hyperref[sec:index_chronic]{\small{[返順時目]}}
~
\hyperref[sec:index_scriptual]{\small{[返順卷目]}}
~
\hyperref[sec:ytVA464zDA8]{\small{> > > NEXT SERMON > > >}}
\newline
\newline
$^{1}$Well, it's really great to be here and as Pastor Mark said, I've known him for a while.
but some of you have not seen me before..
I actually live in the northwest part of the city, but this is where I grew up, so it's.
like coming back to home..
I also found out that I've reconnected with Mir, who I've met before, so it's nice to.
see her and realize that actually some of you are part of a group that's come to UV.
from another congregation and Jayuk is actually a former classmate of mine, so it's nice to.
have that multiple overlap that I realize..
And I actually saw Jayuk before after moving back into the city from Trinity, where we.
were at seminary, and realized that his kids were going to the same school as our kids,.
and so it's nice to have just these multiple connections..
When Pastor Mark asked me to preach today, I asked him what scenes or what books he might.
be covering so that I can maybe stay consistent with the message and maybe how I should cater.
what I was going to preach, and he wittingly replied that he had taught some things in.
Genesis before and then he's taught some things in Revelation, so pretty much anything in.
those books or in between is anything that I can choose from..
So that left me with a lot of freedom and a lot of no direction, but I hope that what.
I chose today will be helpful and relevant..
And the topic that I chose for today is actually something that we as Asians do not like to.
address very much, and it's the topic of, the softer topic of emotions..
And some of you are already probably thinking, emotions, feelings, what good are those?.
Why talk about it?.
And I think in my younger years, I would probably be on your side as well..
I would think how emotions and feelings were not really convenient and weren't really useful,.
and if anything, might be a liability in making decisions and just journeying through life..
But I think that this is really colored by our Asian perspective, and we tend to focus.
on things that are just more practical and monetarily profitable..
What I've come to discover over the years is this is actually rather unbiblical..
This is not how we're meant to be..
And emotions are very much part of how God designed us and how he wants us to be..
He wants us to feel..
And so the title of my message today is actually An Emotive Lord, because our God is emotive.
and expresses a range of emotions..
We who are also created in his image are also meant to display the range of emotions that.
we see him express in the Bible..
But because of our brokenness, we fall into elevating and prioritizing what is our cultural.
upbringing and what is our cultural lens, and that just affects how we live..
We fall short of fully living out how God designed us to be and how we're meant to emotionally.
also express our faith..
And I know that over the past several years, it's been a really hard time..

$^{41}$The pandemic has taken a toll physically, mentally, spiritually, and emotionally on.
us..
And people are really feeling very anxious, and they're feeling the fragility of life..
And I actually observe this often through the teammates that I have, that I interact.
with in my work..
I actually work at a startup as my full-time job, and my team is globally distributed across.
the world..
And in my role, I actually do a lot of these one-on-one checkups with every person that.
I can on a, if not weekly basis, then a biweekly basis..
And there's a particular teammate of mine that I had a conversation with recently..
And these meetings are meant to be a pulse check to just check the general mood of people.
across the company to see how everyone's doing..
And the conversation that I had is with my teammate, who is also a fellow believer..
It was right after Easter, and I spent some time just kind of asking him how he is..
And there's always some personal chit-chat that happens at the beginning of these conversations..
And he told me that his Easter was relatively quiet and relatively somber..
And the reason was because his parents had to abruptly actually leave and go see his.
maternal grandmother, because they had found out that she was on her deathbed..
Fortunately, his mom and dad were able to actually make it in time to see his maternal.
grandmother before she passed..
I offered my condolences, and I just really tried to care for him and ask how he was doing..
And he said that he was close to his grandmother..
He said that he has really fond memories of her and how he's helped to shape his own faith..
And he started to shift into talking about how despite her passing, because of our faith,.
because of the hope that Jesus Christ provides, that he was hopeful about the situation..
And this was all true, or this is all true..
And especially with his grandmother passing over Easter, there is this idea that this.
wasn't the ultimate end, that there's a hope, there's an alternative for his grandmother.
beyond just the passing in this earthly world..
And it certainly made what was more of a sad and morphal situation a little bit more hopeful.
with this idea in mind..
However, despite not speaking about his feelings, I could really sense this greater level of.
grief and sadness that was emanating from him..
And it was masked by this cover of hopefulness..
I felt that he was unnaturally overlooking the grief and just trying to maintain a positive.
outlook and then trying to shift toward work..
I think many of us would actually appreciate this approach, and some of us have probably.
adopted this approach in handling an experience of loss..
But this conversation actually sat with me for a while, and I think it really disturbed.
me and it caused me to reflect upon my own experiences in the past and my own dealings.

$^{81}$with sadness..
I recalled my memory of my grandfather, who was 93 when he passed..
And a lot of people, when they found out, expressed their sympathies, which I was very.
thankful for..
And my reply was usually, "Oh, you know, it's okay..
He lived a long life.".
And this was the response that I provided people that came to me with their condolences..
It helped me accept the situation, but it didn't really acknowledge my true feelings..
And over time, I've come to believe that this isn't really a biblical approach..
It's not really the way that God had intended for us to deal with situations like this..
And so today's passage is actually from John 11, 1 to 44..
You can't, I guess you can't see it..
And I was planning on reading it, but then I decided that actually what's better is for.
me to kind of just paraphrase and summarize the passage and then kind of hone in on what.
I want to aim in on..
And so John 11 tells a story of Lazarus's death and resurrection..
In the beginning part, Jesus is in the wilderness, and Lazarus is sick in a different part of.
Israel, you know, in Bethany..
And this is the village where Mary and Martha are from, who are sisters of Lazarus..
The sisters sent word to Jesus and told him that the one he loves is sick..
These three siblings are actually described as ones that Jesus loves..
And so Jesus had a very special relationship with them..
When he heard this, Jesus said that this sickness will not end in death, but it is for God's.
glory so that God's son may be glorified through it..
Jesus doesn't leave right away, but he actually does the opposite..
He actually waits even longer to make a move and to leave..
He says to his disciples that Lazarus has fallen asleep and he will wake him up..
And so his disciples thought that he meant that he fell into a natural sleep and that.
he'll get better..
But then Jesus then has to correct him and say plainly that Lazarus is dead, but he intended.
for this to happen so that they may believe..
And Thomas actually replies, "Oh, then let us all go so that we might die together.".
So I think they really missed the point, even though Jesus was correcting them..
And so by the time that Jesus arrives in Bethany, Lazarus had already been dead and entombed.
for four days..
And many Jews had come to Martha and Mary to mourn with him and to comfort them during.
this period of loss..
When Jesus saw all these people weeping, it's described that he was deeply moved in spirit.
and troubled..
And in verse 35 out of the 44 verses that we have here, Jesus wept..

$^{121}$He's moved by the situation and then he goes on to ask for the stone to be rolled away.
from the tomb and cause Lazarus out, which he does, fully resurrected, looking like a.
mummy..
Those around come to believe and Jesus and God are glorified in this situation..
There are actually very many preachable lessons here..
It's a very jam-packed passage..
Jesus knowing Lazarus' condition deliberately delays his visit to Lazarus until way past.
his death to magnify the impact of the miracle he knew he was going to perform..
Jesus shows his power to resurrect the dead in dramatic fashion as he calls folks to roll.
away this stone and signals his own resurrection through this story..
So much can be impacted and applied to our lives about suffering, perseverance, and eventual.
deliverance, how God's timing is perfect for a reason even in suffering..
But my focus is actually none of this..
What I want to hone in on is verse 35..
And this small part is in fact the easiestly memorized Bible verse in the whole of the.
Bible because it is the shortest Bible verse..
It is verse 35, which is Jesus wept..
And so now you can all tell Pastor Mark that you've memorized today's verse because it.
is so easy to remember..
But there is something very profound in this two-word verse..
It openly describes Jesus, who is fully divine and fully human, weeping..
It isn't shown in a negative way..
It shows a very human response of grief as completely acceptable..
It doesn't dismiss it..
It doesn't regard it as weakness..
It doesn't try to cover it with something nice..
It's shown as natural and completely embraced..
Now remember that Jesus knows the greater plan..
He knew Lazarus was going to die, but he also knew that Lazarus was going to be raised from.
the dead by him..
With that in mind, many might wonder, "Well, why be sad at all?".
Yet Jesus openly grieves..
Maybe he grieved at the individual death of Lazarus, or maybe his grief comes from the.
idea that there will be so much death still until his resurrection and his second coming..
But whatever the case is, we see here that he wept, that he was moved and troubled, and.
that he openly grieved..
And there was nothing wrong with that..
This act of Jesus affirms grief, sadness, and despair as appropriate responses and feelings.
that we should more readily embrace..
Now, God really expresses a range of emotions as we see from Jesus..

$^{161}$But sociologists say that it's common for people to perceive that God is like the fatherly.
figures in their lives..
If God is caring, patient, and concerned, then children will believe God has those same characteristics..
And the opposite holds true when a father, when an earthly father is harsh, judgmental,.
or absent..
Fortunately for us, or unfortunately for us, a typical Asian dad is one characterized as.
a dutiful, hardworking provider, sometimes a bit emotionally distant, and one whom we.
might, who might be critical and hard to please..
Emotionally available and accessible is not a usual description of Asian parents..
This is important to recognize because this lens can shade our perception of how God the.
father is to us and how he's captured..
We run the danger of projecting an idea of how God the father is through a worldly lens.
versus a biblical one..
This has risks of stunting our spiritual maturity and can lend itself to unhealthy perceptions,.
especially of how God, of how we think God sees us, how we see ourselves, and how we.
think we think God wants us to act..
When we perceive God the father through a traditional Asian parental profile, we can.
project him to be a figure that is harsh, unforgiving, demanding, and not easily satisfied..
We may make him out to be emotionally distant and unconcerned and dismissive of our feelings.
and cause ourselves to undermine our own emotional health and situations..
We need to remember that God expresses the emotions of righteous anger, judgment, and.
disappointment, but he also shows a positive range of emotions of concern, love, pride,.
and joy..
These are expressions that Asians, Asian parents, have a hard time struggling to display..
Our earthly relationships can really cloud how we see our relationship with God the father.
is..
And we can dangerously misinterpret God being dismissive of our feelings and just focusing.
on our performance like our earthly parents do..
We are a broken people and we need God..
And we need to be mindful of how not to be trapped in this lie that he doesn't care about.
our emotions and our feelings..
It's important to observe and know that God displays a range of emotions with very positive.
ones as well..
There's delight, there's satisfaction, there's joy..
And we need to be able to embrace these as well as display these ourselves..
We can establish that God openly expresses a range of emotions..
And knowing that, and knowing that we are made in the image of God, we must not restrict.
how we model a potential range of emotions..
It's not enough that we acknowledge the negative emotions that we tend to avoid..
We also have to embrace and express the positive ones thoroughly and toward ourselves..

$^{201}$You can ask yourselves, do you receive compliments well?.
Do you feel successful?.
Do you feel valued?.
Or do you often feel that you're not good enough?.
And you feel like an imposter about to be found out by someone and revealing your true.
nature..
Now, this is also very small..
It's an emotional wheel..
And it shows a large range of emotions from happy, sad, disgusted, angry, fearful, bad,.
surprised, from the center spanning out to just more descriptive ones of feeling isolated,.
victimized, horrified, furious, excluded, but also positively aroused, free, joyful,.
curious, successful, respected, valued..
Just think about some of these for a while and see if you can actually articulate some.
of these or can actually say that you've expressed a range of these emotions..
Coming from a background that is grounded in honor and shame, most of us probably inherently.
understand that we tend to be emotionally reserved, focused on gaining face, concentrated.
on our duties and responsibilities, are hard to please, and very critical of others as.
well as ourselves..
And this whole sermon has actually been a difficult thing for me to prepare because.
I can't honestly say that I'm good at this..
Emoting is something that doesn't come easily to me..
And I remember when I was younger, I'm very expressive emotionally..
But then over the years growing up, I've been conditioned to focus on what's proper and.
what I need to do in terms of my duty and focus on the family and focus on our collective.
group and to preserve face..
It's been an even greater challenge for me as I've become a parent and how I've unwillingly.
transferred this behavior to my own children..
And as much as I try to maintain a positive mindset over the years and be more expressive.
emotionally with my feelings, with my children, I find that I'm challenged by my past and.
I tend to revert back to just how I was raised..
And some of you are parents here, are preparing to be parents..
Before you become a parent, you think that you're going to prepare to be the most amazing.
parent the world has ever seen..
And we study, we seek out advice, and we plan to be the greatest in the world to our kids..
But the truth is that when the children actually come, it's chaos..
And there's nothing that can prepare you..
And you can study all you want, but it's just a different experience..
And what happens is that what we plan on doing in terms of best practices just turns into.
just trying to survive one more day..
And when that happens, all the stuff that you thought about, planned out, goes out the.

$^{241}$window and you revert back to your natural instincts..
You parent how you were parented..
And emotions for me are not given their proper attention..
And what I see is my own children beginning to avoid expressing their feelings for fear.
of being dismissed..
I don't sit with them and their emotions adequately enough..
I don't hold space for them in their moments of hurt and sadness..
And I don't offer praise as much as I should..
I fall into this performance-oriented person that our culture has trained us to be..
And even when we do something well, or even if we're the best at doing what we're doing,.
there's very little praise for us, right?.
One because that's what we were supposed to do, so why get any praise?.
Or the fear is that, well, if there's too much praise, we breed mediocrity..
People get too used to it and that it doesn't really have the same level of strength that.
it's meant to have..
I know someone who shared with her parents not too a while ago that she took some standardized.
tests and she rated in the 99th percentile..
Her parents' immediate response was, "Well, why weren't you in the 100th percentile?".
But she already, if you understand percentiles, being in the 99th percentile means that you're.
in the top 1\%..
But the immediate response was, "Why weren't you in the top 100?".
Which you can't be..
And the funny thing is that even though your parents don't praise you, what you see is.
that they actually brag about you in front of their friends, right?.
So they use your accomplishments and your performance as ammunition for those around.
them, but you get none of the credit..
You get even more burdened because then you have to continuously deliver these expectations..
As aspiring as Jeremy Lin was a number of years ago, I felt sorry for a whole generation.
that came after him..
And I can imagine the countless conversations that were had of, "Well, not only did Jeremy.
Lin get into the NBA and is six feet plus tall, he also went to Harvard..
What about you?".
You know, I can already see the pain of all these children that had to face this reality..
But more recently we had the Winter Olympics..
And of course we know that the focal point for Asians was not the athletics, it was who?.
It was Aileen Gu..
And I thought it was hard enough when Jeremy existed, but Aileen Gu just blows them out.
of the water, right?.
Here we have someone that is a multi-gold medalist in the Olympics, scored near perfect.
in her SATs, is going to Stanford, is a millionaire, is a model, reads, writes, and, you know,.

$^{281}$reads, writes, and speaks Chinese fluently, and then just add insult to injury, she can.
recite Chinese poetry from the Tang Dynasty, right?.
Like we are all condemned..
This just makes it incredibly hard for everybody that comes now from the Chinese household..
But this is why God is great..
He doesn't care about this stuff..
He loves us for not what we can do, but who we are, right?.
We need to embrace the idea that God expresses delight toward us versus persistently projecting.
this idea that God is a critical parent and just expects us to deliver and perform all.
the time..
We need to feel loved ourselves..
You know, in Matthew 3, 16 to 17, it says that as soon as Jesus was baptized, he went.
up out of the water..
At that moment, heaven was open and he saw the spirit of God descending like a dove and.
rest on him..
And a voice from heaven said, "This is my son, whom I love..
With him, I am well pleased.".
God delights in his son for being faithful..
Who here has heard their parents say to them that they delight in you?.
Who here as parents have shared that same expression to their children that you delight.
in your children?.
God wants us to be able to feel the good, like adoration, loyalty, love, desire, peace,.
success and joy, along with the negative that we feel..
The Bible clearly captures the full range of emotions..
The clearest comprehensive collection is in the 150 Psalms that we have in the Bible..
It speaks of lament, fear, pain, hate, but peace, joy, contentment, love, and so on..
The Song of Songs is a full love letter..
And Revelation starts with a grim reality but turns into expected hope..
We need to be able to cry when we feel despair and celebrate jubilantly when there are things.
worth being happy about..
King David is seen trapped in some pretty dark moments in the Old Testament..
But he's also shown to be depicted as someone that can shout for joy and dance like a fool.
in the streets..
It's not enough for us to start expressing the feelings now..
And we need to express the full range, but God also wants us to, sorry, to work deeply.
through our emotions..
So we are meant to see God as one who expresses a range of emotions..
And sorry that I kind of blasted through that..
But God also encourages us to express the range of emotions ourselves..
And then thirdly, God wants us to work deeply through our emotions..

$^{321}$He encourages us to work through some of our older painful emotions..
We're called to dig deep into our past at times to examine some of the harder things.
that we've been through..
Studies have shown that there are real physical and mental effects resulting from undealt.
with suppressed and repressed emotions..
Studies show that individuals who repress their emotions and suppress their body's.
immunity, making them more vulnerable to a variety of illnesses leading from the common.
cold to cancer..
There are real effects from undealt with emotions, not just on the emotional level, but that.
physically manifest themselves in physical form in our bodies..
And it's not that dealing with emotions is a nice to have, but it's closely interconnected.
with all that we are, all that we are, with our physical, mental, and spiritual well-being.
by design..
And we have a responsibility to address them..
Before Jesus was crucified, he predicted Peter would deny him three times before the rooster.
cried in the morning..
But Peter boldly claimed that he would follow Jesus both to prison and to death..
Now we know that this wasn't true because Peter indeed denied Jesus three times, and.
he was bitterly ashamed of himself for it..
It's described in the Bible that he wept bitterly after the rooster crowed when he realized.
that Jesus' prediction came true..
And Peter could have wrestled with this for the rest of his life..
But what we see captured at the end of John is that Jesus works through these shameful.
and guilty feelings that traumatize Peter..
After Jesus is resurrected and reappears to his disciples, he's captured at the end of.
the book of John, sitting Peter down over a meal of fish that they just caught..
And this is the passage..
They had finished eating, Jesus said to Simon Peter, "Simon, son of John, do you love me.
more than these?".
"Yes, Lord," he said, "you know that I love you.".
Jesus said, "Feed my lambs.".
Again, Jesus said, "Simon, son of John, do you love me?".
He answered, "Yes, Lord, you know that I love you.".
Jesus said, "Take care of my sheep.".
The third time he said to him, "Simon, son of John, do you love me?".
Peter was hurt because Jesus asked him the third time, "Do you love me?".
He said, "Lord, you know all things, you know that I love you.".
Jesus said, "Feed my sheep.".
We see Jesus here going deeper and deeper with Peter, probably because Peter was avoiding.
this..

$^{361}$He doesn't explicitly bring up the denial that Peter did, but he asked him three times,.
almost in a way, redeeming the three denials that he had..
Jesus actively works with Peter through his emotions to bring restoration, and over a.
meal that actually physically restores him as well..
Those who serve in ministry are not exempt from this work..
One of the requirements to graduate seminary for me was actually go through a number of.
counseling sessions..
This was very important because we have to make sure that whatever potential painful.
memories, experiences, don't surface and put our future ministry at risk..
And there are actually some churches that are really good about this that have a requirement.
for their pastors to just get counseling every so often to make sure that they're cared for.
and that they're addressing their emotional and mental health..
When we go through periods of dismissed emotions that get suppressed and repressed, there are.
triggers that could derail us, that can surface in ways and at times that we don't expect..
We are seeing a lot of the effects of isolation and depression amongst children, especially.
during these COVID times, and we have to be very mindful of the longer-term effects of.
this behavior..
So I have a task for you..
I want you to, and I'll leave this on here, I want you to look at this emotional wheel.
and pick an emotion here that is hard for you to express..
And I want you to just kind of process it over this week, and if you can, by the end.
of next week, to find someone that you can trust to actually just talk about it and try.
to talk about it in a way that you haven't before..
I challenge you to share it with someone that you can trust so that you can really, you.
know, exercise, take this exercise just really seriously..
Well, Pastor Mark said that I could preach about anything from Genesis to Revelation,.
and the entire Bible captures the full range of emotions..
So I think that I fulfilled my duty in spanning the whole of the Bible with emotions that.
span the whole of the Bible..
And I just want to say that it's okay to not be okay..
It's okay to not feel okay..
And there's nothing wrong with that..
And there's nothing wrong with feeling great as well..
I think that sometimes being Asian, when we feel too positively, someone just wants to,.
you know, squash that..
So it's okay to feel great when there's cause to celebrate and to be festive..
God actually expects it..
And God actually designed us to be that way..
So we deny ourselves the fullness of how we're meant to live and how God created us when.
we restrict ourselves in this area of life..

$^{401}$Embrace the way that the Lord has made you, understanding that it might be uncomfortable.
at times to work through some of the things that we like to avoid, especially because.
of our cultural conditioning..
But don't allow our cultural lens, our upbringing, to shade the kingdom-centered perspective.
that God wants us to have..
Be renewed through His Spirit..
Have this fresh perspective..
And don't let the worldly way that we've been brought up color the way that we're meant.
to understand the world from just a believer's perspective..
Well, that's it for me..
Let me just close in prayer..
Heavenly Father, we thank you for the way that you've created us and the ways that we.
sometimes don't like..
But we understand that that is the way that we were designed..
Help us here to really feel the way that you want us to feel..
And where there are maybe painful memories, help us to also work through them, through.
your strength and through your support..
Let us embrace the way that you designed us and feel love when we should feel loved, but.
also feel the yucky things when we are meant to feel them, and to not repress them or suppress.
them, but to openly grieve and to openly support those that are going through these periods.
of negative feelings..
We praise you and we glorify your name..
In your name we pray..
\newpage



\section{馬太福音 1:1-17}
\label{sec:ytVA464zDA8}
\textbf{耶穌的騎呢家譜 (馬太福音 1\_1-17) - 何建邦傳道}
\newline
\newline
連結: \href{https://youtube.com/watch?v=ytVA464zDA8}{\texttt{https://youtube.com/watch?v=ytVA464zDA8}} ~~~~ 語音日期: 2022-05-24
\newline
\newline
\hyperref[sec:V2NOXpx01xg]{\small{< < < PREV SERMON < < <}}
~
\hyperref[sec:index_chronic]{\small{[返順時目]}}
~
\hyperref[sec:index_scriptual]{\small{[返順卷目]}}
~
\hyperref[sec:lNtIgECCCGM]{\small{> > > NEXT SERMON > > >}}
\newline
\newline
馬太福音 1:1-17
\newline
\begin{longtable}{cl}
\hline
\hline
章節 & 經文 (和合本修訂版)\\
\hline
1:1 & \begin{tabularx}{0.7\textwidth}{X} 亞伯拉罕的後裔、大衛的子孫耶穌基督的家譜： \end{tabularx} \\ \\ \relax
1:2 & \begin{tabularx}{0.7\textwidth}{X} 亞伯拉罕生以撒，以撒生雅各，雅各生猶大和他的兄弟， \end{tabularx} \\ \\ \relax
1:3 & \begin{tabularx}{0.7\textwidth}{X} 猶大從她瑪氏生法勒斯和謝拉，法勒斯生希斯崙，希斯崙生亞蘭， \end{tabularx} \\ \\ \relax
1:4 & \begin{tabularx}{0.7\textwidth}{X} 亞蘭生亞米拿達，亞米拿達生拿順，拿順生撒門， \end{tabularx} \\ \\ \relax
1:5 & \begin{tabularx}{0.7\textwidth}{X} 撒門從喇合氏生波阿斯，波阿斯從路得氏生俄備得，俄備得生耶西， \end{tabularx} \\ \\ \relax
1:6 & \begin{tabularx}{0.7\textwidth}{X} 耶西生大衛王。大衛從烏利亞的妻子生所羅門， \end{tabularx} \\ \\ \relax
1:7 & \begin{tabularx}{0.7\textwidth}{X} 所羅門生羅波安，羅波安生亞比雅，亞比雅生亞撒， \end{tabularx} \\ \\ \relax
1:8 & \begin{tabularx}{0.7\textwidth}{X} 亞撒生約沙法，約沙法生約蘭，約蘭生烏西雅， \end{tabularx} \\ \\ \relax
1:9 & \begin{tabularx}{0.7\textwidth}{X} 烏西雅生約坦，約坦生亞哈斯，亞哈斯生希西家， \end{tabularx} \\ \\ \relax
1:10 & \begin{tabularx}{0.7\textwidth}{X} 希西家生瑪拿西，瑪拿西生亞們，亞們生約西亞， \end{tabularx} \\ \\ \relax
1:11 & \begin{tabularx}{0.7\textwidth}{X} 百姓被遷到巴比倫的時候，約西亞生耶哥尼雅和他的兄弟。 \end{tabularx} \\ \\ \relax
1:12 & \begin{tabularx}{0.7\textwidth}{X} 遷到巴比倫之後，耶哥尼雅生撒拉鐵，撒拉鐵生所羅巴伯， \end{tabularx} \\ \\ \relax
1:13 & \begin{tabularx}{0.7\textwidth}{X} 所羅巴伯生亞比玉，亞比玉生以利亞敬，以利亞敬生亞所， \end{tabularx} \\ \\ \relax
1:14 & \begin{tabularx}{0.7\textwidth}{X} 亞所生撒督，撒督生亞金，亞金生以律， \end{tabularx} \\ \\ \relax
1:15 & \begin{tabularx}{0.7\textwidth}{X} 以律生以利亞撒，以利亞撒生馬但，馬但生雅各， \end{tabularx} \\ \\ \relax
1:16 & \begin{tabularx}{0.7\textwidth}{X} 雅各生約瑟，就是馬利亞的丈夫；那稱為基督的耶穌是從馬利亞生的。 \end{tabularx} \\ \\ \relax
1:17 & \begin{tabularx}{0.7\textwidth}{X} 這樣，從亞伯拉罕到大衛共有十四代，從大衛到遷至巴比倫的時候也有十四代，從遷至巴比倫的時候到基督又有十四代。 \end{tabularx} \\ \\
[1ex]
\hline
\hline
\end{longtable}
$^{1}$各位大型姐妹.
今天很開心來到城會教會跟大家分享訊息.
這次是我第一次親身來到城會教會.
所以我很榮幸Mark Sir邀請我來到城會教會趕導.
我來到這裡為城會教會感恩.
原來這裡比我想像中更加Cyber.
更加Modern.
敬拜隊也很棒.
我覺得我來到芝加哥之後.
很少聽到廣東話的崇拜有這麼現代的敬拜隊.
通常是英文的Church.
所以我很榮幸很開心可以在這裡參與大家的崇拜.
和一起分享訊息.
今天我要分享的主題是耶穌的奇尼加普.
我要分享的經文是馬太福音的第一章1至17節.
我猜這段經文應該是或許是整本馬太福音裡.
大家最少看的經文.
是不是?.
我不知道是不是.
但我就是.
不知道大家是不是.
在我小時候通常零修看馬太福音的時候.
最好的方法是什麼呢?.
就是快點跳過這一段的加普.
我就可以看少17節了.
是非常非常化算的方法.
很快就可以看完整本馬太福音了.
不過馬太這位福音書的作者.
將耶穌基督的加普放在整本福音書的最開頭.
一定有它一個很重要的目的.
所以今天我們就留心去看這個加普.
究竟想說些什麼.
今天我的講題剛才說過.
是耶穌的奇尼加普.
所以我們會集中去看一下.
究竟這個加普有什麼奇尼的地方.
我們先去看一下這個加普的標題.
這個加普的標題是.
耶穌基督大衛的兒子.
亞伯拉罕的兒子的加普.

$^{41}$透過這個標題我們可以看到.
這個加普有兩個重要的人物.
第一個是大衛.
而第二個就是亞伯拉罕.
我們首先會去看一下.
這個加普究竟怎樣去圍繞著大衛.
不知道大家有沒有想過.
為什麼這個加普由亞伯拉罕去到耶穌.
會剛好有三個十四代呢.
不知道大家有沒有聽過.
為什麼是十四的.
為什麼是十四個數字.
大家有沒有聽過呢.
如果我們比較一下在歷代至上的加普的話.
我們會發現.
馬太刪減了不少的人物.
是特意砌三個十四代出來.
為什麼馬太要這樣做呢.
其實這一點是和希伯來文有關的.
希伯來文和中文是很不同的.
希伯來文每一個字母都是對應著一個數字的.
而在希伯來文裡面.
十四剛剛好就是對應著大V.
這三個字母.
也就是大衛這個名字.
所以你就明白.
為什麼馬太要特意砌呢.
這個加普是要分成三個十四代.
因為十四這個數字是一個很重要的意思.
就是代表著大衛的這個名字.
在希伯來文裡面的這個名字.
如果大V等於十四.
所以要砌十四代的話.
為什麼要砌三個十四代呢.
如果我們去看這個加普的話.
這裡有三個十四代.
我用了一個表格去列出來.
如果我們留意.
每一個十四代的最後一個名字.
第一個十四代的最後一個名字是大衛.

$^{81}$第二個十四代的最後一個名字是耶哥尼亞.
第三個十四代的最後一個名字是耶穌.
這三個名字都是和大衛王朝有關的.
而且是關係著大衛王朝最重要的事件.
第一個人物大衛.
就是代表著大衛王朝的建立者.
第二個人物耶哥尼亞.
是大衛王朝的最後一任的君王.
而第三個人名耶穌.
就是猶太人所盼望的大衛王朝的復興者.
就是大衛的君王.
那位彌賽亞.
所以我們就明白.
為什麼馬太要特意將一個加普編成三個十四代.
為什麼要特意刪減這麼多人物.
一定要砌三個十四代出來.
就是因為十四代代表著大衛.
而這三個十四代的最後一個人物.
都是和大衛王朝的關鍵事件有關.
說完大衛.
我們就看看這個加普.
怎樣圍繞著亞伯拉罕.
這個加普第一個奇怪的地方是.
就是它是很簡單直接的.
它是說一個生一個.
例如亞伯拉罕生什麼什麼.
什麼什麼生什麼什麼.
很少提及兄弟.
唯獨是有兩個地方.
這個加普提及了兄弟.
大家知不知道在哪裡.
是特意提及兄弟的.
是很奇怪的.
就是在第一個十四代的第四個人物.
和第二個十四代的最後的人物.
那裡分別提及尤大和他的兄弟.
和耶哥尼亞和他的兄弟.
如果這個加普是一個很簡單直接的加普.
不斷說A生B.
B生C.

$^{121}$C生D的話.
為什麼要特意在這兩個地方.
提及什麼什麼和他的兄弟呢.
不知道大家有沒有想過.
我們平時看聖經就會隨著阿伯這樣看.
我們很少很留心地去看經文的細節.
不過今天我就很想和大家看經文的細節.
其實馬太特意要在這兩個地方說.
什麼什麼和他的兄弟.
是一個很重要的原因.
尤大和他的兄弟.
令當時的猶太讀者想起.
想起以色列的十二個支派.
尤大是一個支派.
和他的兄弟就是其餘的十一個支派.
所以尤大和他的兄弟代表的是什麼.
尤大和他的兄弟代表的是.
以色列十二個支派的建立.
以色列的歷史的開始.
以色列十二個支派的開始.
而耶哥尼亞和他的兄弟代表的是什麼呢.
剛才我說過耶哥尼亞是大魏王朝最後一個君王.
那就是說什麼.
以色列滅亡.
所以那裡就是說.
以色列歷史作為一個國家的終結.
說的是以色列十二支派的終結.
所以在這裡說.
尤大和他的兄弟和耶哥尼亞和他的兄弟.
是為了令當時的猶太讀者回想起什麼.
回想起以色列以亞伯拉罕為始祖的國家的歷史.
所以你就會明白到.
什麼什麼和他的兄弟.
其實是呼應著亞伯拉罕.
在這個家譜裡面第二個的關鍵人物.
不過我們還有一個地方還沒有解通.
就是如果大家看我這個表格的話.
你就會發現.
在第一個十四代和第二個十四代.
都有什麼什麼和他的兄弟.

$^{161}$特別的地方是什麼呢.
為什麼第三個十四代沒有什麼什麼和他的兄弟呢.
還有第三個十四代有一些奇怪的地方.
雖然馬太說每一個section都有十四代.
但如果你留心一點數字我已經列出來了.
第三個十四代其實是沒有十四代的.
第三個十四代其實只有十三代.
那少了的一代是什麼呢.
是不是馬太算錯了呢.
就是他加減乘除不是很行.
就是加少了一個.
算漏了一個.
所以就少了一代.
最後一個section.
是不是這樣呢.
所以這裡我就很想跟大家說.
其實馬太特意在第三個十四代.
沒有提及什麼什麼和他的兄弟.
還有少了一代.
是一個很重要的原因.
如果大家能夠翻譯聖經的話可以翻譯.
不過不用翻譯也可以的.
因為我現在就可以直接說.
就是在馬太福音第十二章的五十節.
耶穌重新去定義.
哪一個才是他的兄弟.
耶穌說哪一個遵行天父的旨意.
哪一個就是他的兄弟.
在馬太福音二十八章十節那裡.
復活之後的耶穌.
和那些去墳墓找他的婦女說.
他和我的兄弟說.
叫他們去加利利.
他們在那裡會見到我.
在這段經文裡面所說的兄弟.
很明顯就是指耶穌的門徒.
如果耶穌的門徒.
就是那些遵行天父旨意的人.
就是耶穌的兄弟的話.
今天同樣地我們這一群遵行耶穌的教導.

$^{201}$跟隨他的人.
我們都是耶穌的兄弟.
所以我們這樣看的話.
我們就會明白.
馬太特意不在第三個十四代裡面去提.
「XX和他的兄弟」.
是因為這個家譜還沒完.
馬太特意不提「XX和他的兄弟」.
是想我們能夠進入這個家譜當中.
是想這個家譜能夠應用.
在我們這個時代當中.
因為這個家譜還沒完.
以色列的歷史.
並不是終結在耶哥尼亞和他的兄弟那裡.
以色列的歷史是繼續延續著.
不是只是停在耶穌的時代那裡.
今天我們作為耶穌的兄弟.
馬太科音重新定義了.
耶穌的兄弟就是我們.
就是我們這一群追隨耶穌.
跟隨他教導的人的門徒.
所以為什麼馬太特在這裡沒有提.
「XX和他的兄弟」.
就是因為這個家譜.
最後的那一代.
就是由我們這一代去繼續延續.
這個就是馬太特很想告訴讀者.
為什麼在第三十四代.
不再提「XX和他的兄弟」.
也沒有了最後一代.
是因為沒有了的那一代.
就是讀者.
沒有了的那一代.
就是耶穌的兄弟.
就是復活主的兄弟.
馬太透過這個家譜.
去說以色列的歷史.
是怎樣開始.
怎樣終結.
和怎樣延續.

$^{241}$在當時是一個很瘋狂的想法.
因為在那個時代.
在古代.
從來不會有人去想.
一個民族.
一個國家的歷史.
是可以由另一個不同的民族去延續的.
是不是.
我們在座每一個應該都不是猶太人.
不過馬太是在說.
這班不是猶太人的人.
是在延續猶太人的歷史.
以亞伯拉罕為始祖的以色列的歷史.
這是一個很瘋狂的想法.
在我們這個時代都覺得很瘋狂.
現在在現代的中國.
很講愛國主義.
沒有任何一個中國人會覺得.
中國人的歷史是可以由美國人去延續的.
沒有任何一個中國人會覺得.
中國人的歷史是由黑人去延續的.
是沒有人會這樣想的.
不過馬太就是要這樣說.
今天.
這個以色列民族的歷史.
是由我們這班不是屬於以色列的人去延續.
這個訊息對於我們今天的時代來說.
都是很重要的.
如果我們看今時今日.
我們這個時代其實和耶穌的時代一樣.
是充斥著各種的民族主義.
各種的愛國主義.
如果我們看看中國.
看看美國就會看到.
很多這樣的思想.
這篇經文不單挑戰當時的猶太讀者.
也挑戰今天我們這班人.
基督徒所延續的以色列.
並不是一個基於民族的國家.
我們不是要追求愛國主義.

$^{281}$我們不是要追求民族主義.
我們延續的.
我們效忠的.
並不是一個特別的.
不是一個.
任何一個特定的民族.
特定的國家.
而是一個由不同民族.
不同國家所組成.
效忠於耶穌基督.
這位大衛君王.
這位彌賽亞的上帝的國.
這一點對於我們這個時代.
這個被愛國主義.
被民族主義.
被種族主義充斥的年代.
同樣是很色切.
同樣是一個很大的挑戰.
接著我說完「XX和他德兄弟」之後.
我就想再去看這個家譜.
第二個奇怪的地方.
這個家譜第二個奇怪的地方就是.
這個家譜所提及的人物.
大部分都是男人.
主要是男人.
90\%多都是男人的名字.
唯獨是有五個女人.
出現在這個家譜裡面.
就是他馬拉合路德.
和烏尼亞的妻子.
就是拔斯巴.
和瑪麗亞.
有這五個女人.
如果這個家譜主要是提及男人.
為什麼無緣無故又要提及五個女人呢.
這五個女人的出現.
一直都困擾著歷代的解經家.
所以有不同的解釋.
去解釋為什麼這五個女人要在這個家譜裡面出現.
第一個解釋就是.

$^{321}$在第四世紀有一個基督教神學家.
我們叫教父.
第四世紀有一個教父叫Jerome.
耶佑米.
就提出一個解釋.
他說馬太提及這五個女人.
是因為她們都是罪人.
她們的存在是為了凸顯.
耶穌是為罪人而生的.
不過我自己就覺得這個解釋是很偏頗的.
因為在猶太的歷史上.
從來不會有人覺得路德是一個罪人.
路德是一個敬虔的典範.
路德是一個聖女.
在猶太的傳統裡面覺得.
所以如果你說路德也是一個罪人的話.
是一個比較難去.
比較牽強的一個解釋.
有人就說她們都是義人.
所以為什麼馬太要特意提及這五個人的名字.
不過這一點同樣都是很偏頗.
因為雖然在撒母耳記下.
他強調的是大衛主動和拔士班通姦.
但是都未至於去描述拔士班是一個義人.
一個敬虔的人.
所以如果我們說馬太記載這五個女人.
因為她們是義人.
這個解釋都是非常偏頗的.
最流行的解釋.
最多人支持的解釋就說.
這四個女人不包括瑪麗亞.
這四個女人都是外邦人.
不過我想說的是.
多人支持的解釋不一定是對的解釋.
為什麼呢?.
因為我們在舊約聖經有說.
拉合和路德都是外邦人.
同樣地在猶太傳統裡面.
雖然在舊約聖經裡面沒有提及他馬是不是外邦人.
不過在猶太傳統裡面.

$^{361}$我們都可以看到猶太傳統會認為他馬是迦南人.
所以頭三個女人說她是外邦人都是合理的.
他馬,拉合和路德都可以說她們是外邦人.
不過問題是什麼呢?.
問題是從來都沒有文獻去說拔士班是外邦人.
在舊約聖經沒有說拔士班是外邦人.
同樣地在猶太傳統裡面.
都沒有說過拔士班是外邦人.
所以如果我們說馬太記載這四個女人純粹是外邦人.
這個解釋其實都不是很合理.
因為最少拔士班不是外邦人.
所以我覺得這三個解釋都不太能解釋.
為什麼這五個婦女要被記載在馬太家譜裡面.
其實按照馬太的句法和編排.
我認為這五個女人應該分成兩組.
分別是第一組是他馬,拉合和路德.
第二組就是拔士班和瑪利亞.
這樣的分組有兩個原因.
第一我們剛才說過大衛是這個家譜的關鍵人物.
他馬,拉合和路德是排在大衛之前.
而拔士班和瑪利亞是排在大衛之後.
第二馬太在從XXXXXX的句式裡面.
有直接提及他馬,拉合和路德的名字.
但從來沒有提及拔士班和瑪利亞的名字.
馬太在這個家譜裡面如果大家有留意.
馬太在這個家譜裡面從來沒有提及拔士班的名字.
馬太稱呼拔士班做烏利亞的妻子.
至於瑪利亞雖然馬太有提及瑪利亞的名字.
但如果我們夠留心去看經文的話.
我們就會發現第十六節的希臘文.
是將瑪利亞的名字連到約瑟裡面.
馬太稱呼約瑟為瑪利亞的丈夫.
但在十六節的那句從XXXXX的句式裡面.
在原文裡面瑪利亞的名字是沒有出現的.
馬太的寫法是甚麼呢.
馬太的寫法是耶穌從他而生.
他即是瑪利亞.
換句話說.
馬太沒有在耶穌的出生裡面提及瑪利亞的名字.
而是用他這個關係代名詞去取代.

$^{401}$所以我們可以根據馬太的編排和句法.
將這五個女人分成兩組.
分別是他馬,拉合,魯德.
第二組是拔士巴.
即是烏尼亞的妻子和瑪利亞.
除了句法和編排之外.
這兩組女性各自有甚麼共通點呢.
我們一定要明白這些共通點.
我們才會明白.
為甚麼馬太要將她們放在家譜裡面.
我們首先看看他馬,拉合和魯德的共通點.
第一,我剛才已經說了.
她們都是外邦人.
第二,她們不單只是外邦人.
她們是一些主動行動的人.
在《創世記》的38章裡面.
她說到她馬主動扮成妓女.
去毆打她的老爺尤大.
和尤大交合.
是為了為她死去的老公.
即是尤大的兒子去留後代.
而拉合和魯德的故事我們比較熟悉.
在《約書亞記》裡面.
拉合主動拯救了以色列派去耶利哥打探的探子.
而在《魯德記》裡面.
魯德主動陪伴他的奶奶娜娜.
並且主動接觸波亞斯.
以致後來成為波亞斯的老婆.
得以為她的老公留後代.
第三,這三個女人不單只是主動行動.
而且是甚麼.
而且展現出對上帝的信心和德行.
尤大知道了她馬是為了自己的兒子留後代.
才扮成妓女和她交合之後.
就讚揚她馬比自己更有義.
拉合表現出對以色列上帝的信心.
所以她才會選擇去救耶利哥的探子.
而魯德就展現了對他奶奶娜娜的不離不棄.
展現了對他奶奶娜娜的忠心.
和堅持地去照顧娜娜.

$^{441}$以致被讚揚有這個媳婦比有七個兒子更好.
說完塔瑪拉和魯德之後.
拔士巴和瑪利亞又有甚麼共通點呢.
最明顯的共通點就是.
他們的老公都不是他們兒子的親生父親.
瑪利亞的兒子就不用說了.
瑪利亞即是耶穌並不是約瑟的親生兒子.
大家都明白.
至於所羅門.
由於加普強調的是.
拔士巴是屬於烏尼亞的.
拔士巴是烏尼亞的妻子.
所以根據加普的寫法來說.
所羅門並不是拔士巴的老公的兒子.
即不是烏尼亞的兒子.
他是大衛的兒子.
所以根據加普無論耶穌還是所羅門.
都不是他們母親的老公的親生兒子.
明白我所說的吧.
說得有點蹺蹊.
第二個共通點是甚麼呢.
第二個共通點就是拔士巴和瑪利亞.
都是一個很被動的角色.
在聖經裡面提到.
在撒慕爾記下的十一章.
強調的是大衛和烏尼亞之間的互動.
和大衛主動和拔士巴通姦.
拔士巴被描述到只不過是.
在遵行大衛的命令.
所以拔士巴是一個很被動的角色.
在整個撒慕爾記下十一章.
拔士巴只有一句對白.
這句對白還要打發人和大衛說話.
所以如果你想像一下.
我們在看劇集.
這是一個非常沒有氣氛的角色.
至於瑪利亞.
在《馬太福音》裡面.
在耶穌出生的故事裡面.
如果大家有留心.

$^{481}$《馬太福音》和《路加福音》是非常不同的.
《路加福音》瑪利亞是很主動的.
《路加福音》裡面瑪利亞是一個很有氣氛的角色.
但在《馬太福音》裡面.
瑪利亞是一個沒有氣氛的角色.
為甚麼我這樣說呢.
就是因為如果你留意耶穌出生的記載的話.
天使說話的對象並不是瑪利亞.
而是弱失.
OK.
之後決定是否休瑪利亞的是弱失.
發夢的.
即是發異夢的是弱失.
但耶穌和瑪利亞去逃避希律追殺的那個人.
是弱失.
所以有氣氛的是弱失.
而不是瑪利亞.
瑪利亞和拔士巴一樣.
都是一個毫無氣氛的角色.
甚至比拔士巴更慘.
一句對白都沒有.
在耶穌出生的記載裡面.
除了這兩個共通點.
最後一個共通點是甚麼呢.
就是拔士巴和瑪利亞.
第三個共通點是.
上帝親自.
Sorry 我忘記了按PowerPoint.
上帝親自為他們的兒子去改名.
這一點我們很少留意.
如果我們去看.
撒慕爾記下的十二章二十五節.
我們很少知道所羅門其實有第二個名字.
大家是否不知道.
如果知道的話可以告訴我.
知道是甚麼名字嗎.
所羅門第二個名字.
我們看撒慕爾記下的十二章二十五節.
那裡就說到.
上帝藉著先知拿丹賜給所羅門一個名字.

$^{521}$叫做耶底底亞.
這個名字說出來都不知道是誰.
這個就是所羅門的第二個名字.
同樣地.
耶穌這個名字其實就是天使.
即是上帝差遣天使叫弱失去改的.
所以耶底底亞和耶穌這個名字.
都是上帝去改的.
第四個共通點就是.
上帝都很愛拔茲巴和瑪利亞的兒子.
在撒慕爾記下的十二章二十五節說到.
上帝喜愛所羅門.
同樣地在馬太福音的三章說到.
當耶穌受施洗的汗洗禮的時候.
天就開了.
就有聲音傳出來說.
這是我的愛子我喜愛他.
所以同樣地.
這兩個女人的兒子都是上帝所喜愛的.
好了我們說了這麼多經文的東西之後.
我們是時候去問我們一個很重要的問題.
說完這些共通點之後.
我們是時候問.
那馬太想說些什麼呢.
馬太想說些什麼呢.
我們先說拔茲巴和瑪利亞.
我們剛才說過.
大衛王朝是這個家譜不斷圍繞的重點.
拔茲巴生的是大衛王朝以色列歷史當中.
一個最偉大的君王所羅門.
他是起聖殿的.
所以他是一個最偉大的君王.
在他的在位的時間裡面.
以色列的疆土,疆界是最大的.
而瑪利亞生的是復興大衛王朝的彌賽亞耶穌.
為什麼上帝要揀選這兩個女人.
去生這兩個以色列歷史當中最偉大的君王呢.
是不是因為她們漂亮呢.
是不是因為她們聰明呢.
是不是因為她們是二人呢.

$^{561}$撒慕爾記下沒有說原因.
同樣地馬太福音也沒有說原因.
拔茲巴和瑪利亞在這兩本書中.
只不過是沒有什麼氣氛.
非常被動的角色.
古代的信徒就很好奇.
為什麼上帝要揀選瑪利亞生耶穌呢.
於是在公元二世紀.
就出現了一本典內文獻.
叫做《雅各原始福音》.
這本典內文獻就嘗試去解釋.
上帝去揀選瑪利亞.
是因為瑪利亞小時候就已經很聖潔.
這本書強調瑪利亞的聖潔.
衍生出後來天主教的一個.
叫做聖母無原罪的教義.
和聖母終身同征等等的教義.
我在這裡不是想批評天主教.
因為我覺得天主教和我們一樣.
都是基督信仰裡面的等子們.
不過我在這裡想說的是.
這個天主教的解釋.
是抓不到瑪太福音的重點.
瑪太並不是想說瑪利亞特別聖潔.
他不是想說瑪利亞特別的公義.
不是想說瑪利亞是一個好人.
在這裡上帝揀選拔士巴和瑪利亞.
去生所羅門和耶穌.
並不是因為他們漂亮.
並不是因為他們聰明.
並不是因為他們是好人.
上帝的揀選其實是不需要原因的.
沒有原因才是真正的重點.
就好像在撒慕爾記上十六章七節裡記載.
上帝是按著他自己的心意去揀選大位.
並不是因為大位的心有多純潔.
瑪太將拔士巴和瑪利亞放在這個家譜裡面.
是為了凸顯上帝對大位王朝的揀選.
所羅門和耶穌的誕生.
並不是靠著敬虔或者顯赫的母親.

$^{601}$而是上帝無條件的揀選.
今日上帝揀選你和我.
同樣都是無條件的揀選.
同樣都不是因為我們的背景有多好.
被揀選並不是靠父幹.
亦都不是靠母幹.
上帝揀選我們是沒有原因的.
其實我想說的是.
沒有原因的揀選才是恩典.
不過做基督徒做得久.
我們很容易就會忘記了什麼是揀選.
什麼是恩典.
今日拔士巴和瑪利亞在這個家譜裡面的存在.
提醒我們.
我們只不過是被上帝揀選和拯救的平凡人.
就好像上帝揀選兩個平平無奇的女子.
去生所羅門和耶穌一樣.
我們是平凡人.
我們沒有什麼可以誇口.
這種平凡才是重點.
馬太特登要凸顯這兩個女子的平凡.
就是為了襯托出上帝的揀選才是最重要.
沒有原因無條件的揀選.
才是真正的恩典.
說完拔士巴和瑪利亞.
我們最後就要問.
為什麼他馬拉合和路德要出現在這個家譜裡面呢.
我們剛才說到.
猶大和他的兄弟代表的是以色列的建立.
而他馬拉合和路德.
正正就是出現在猶大和他的兄弟和大衛之間.
換句話說.
他馬拉合和路德這三個外邦女子的存在.
在說什麼.
在說以色列從建立開始.
就已經有外邦人的血統.
這個家譜呼應著馬太福音最後的大使命.
外邦人要加入成為耶穌重新定義的以色列的一份子.
這個英雄或者這個大使命並不是憑空出現.
而是基於這個家譜.

$^{641}$這個家譜已經埋下了伏筆.
外邦人要加入成為以色列的一份子.
並不是憑空出現.
而是基於以色列本身的歷史.
已經有這三個外邦女子在當中.
我想說的是.
馬太這個家譜對當時的猶太人來說是非常「撩交打」的.
是因為當時的猶太人普遍會覺得.
猶太這個民族.
才是上帝揀選的民族.
而上帝的揀選是基於他們的血統.
馬太竟然說以色列從一開始建立.
就已經有外邦人的血統.
你能想像到這個是多麼「撩交打」的信息.
而這三個女子另一個存在的目的.
是為了為馬太福音之後的.
外邦人的角色設定.
埋下了伏筆.
剛才我說過.
他馬拉合盧德都是主動表現出信心.
主動參與上帝救贖計劃的外邦人.
在馬太福音之後的序詞當中.
外邦人通常都是主動展現對耶穌和上帝的信心.
例如在十五章裡面說到的迦蘭婦人.
即使面對著耶穌的拒絕.
她仍然都展現出對耶穌的一個很堅持.
很堅決的信心.
相反在馬太福音當中.
猶太的領袖和群眾通常都是反對耶穌的.
今天我們作為外邦人.
我們很容易就會覺得猶太人拒絕耶穌.
活該他們死.
歷史上很多基督徒都是這樣想.
所以迫害猶太人.
而反猶主義就是這樣來的.
所以馬太福音在很長的一段時間裡面.
都被基督徒認為是一本反猶的.
針對猶太人的福音書.
但其實我想說的是.
馬太這樣寫不是想針對猶太人.

$^{681}$因為他自己也是猶太人.
他怎會針對自己呢.
馬太這樣寫其實是想針對那些自以為.
靠著種族 靠著血統.
就已經可以被上帝揀選的人.
其實今天可能我們才是馬太所針對的猶太人.
有時候做基督徒做得久.
可能我們都會好像那些猶太人一樣.
或者當時的一些馬太的猶太讀者一樣.
我們慢慢會覺得其他人是騎尾巴.
就好像猶太人覺得外邦人是騎尾巴一樣.
外邦人沒有好東西.
外邦人是不聖潔的.
外邦人是沒有好人的.
所以馬太在這裡特意強調的是什麼.
他馬拉合路德是好人.
他馬拉合路德是展現出敬虔.
展現出信心 展現出義恨的外邦人.
這是很顛覆當時猶太傳統.
對外邦人的一種黑板的印象.
一種Stereotype.
今天同樣地.
我們這些基督徒.
我們很容易成為了猶太人.
那種猶太人不是說血統上的猶太人.
我們都會好像那時的猶太人一樣.
對外面的人有一種Stereotype.
一種黑板的印象.
他們是騎尾巴.
我們是好人.
我們是被揀選.
他們是騎尾巴的.
回教會回得久.
我們慢慢都開始.
我們不會說出口.
但可能開始會有一種基督徒的優越感.
覺得其他人是騎尾巴.
我們才是被揀選.
就好像我們有時會覺得.
咦 你這些人.

$^{721}$原來你也是基督徒.
你這些人你都回教會.
會不會影衰基督徒.
我們內心很多時候都會有這些想法.
最後我想說一個例子.
去結束我今天的訊息.
就是我在香港事奉的時候.
聽過一位牧者分享.
他說他會堅決地.
拒絕為一位變性人去施洗.
拒絕的原因是因為他做了變性手術.
這位牧者就很堅持.
除非這位變性人做一個逆向的變性手術.
否則教會是不會接納他成為教會的會友.
是不會為他施洗.
這位牧者覺得這樣做是在堅持真理.
他覺得這樣做是在捍衛福音.
可能我們在北美的橫教會.
我們都有這些想法.
覺得這樣就是堅持捍衛福音.
但我想說的是.
我不是想在這裡支持LGBT.
我不是要在這裡說一個絕對的答案.
不過我想帶出一個反思.
這個例子很值得我們反思.
就是什麼才是上帝的救恩呢?.
是不是變性人就不可以是異人呢?.
是不是變性人就等於騎呢?.
今天如果變性是罪.
貪心,說謊都是罪.
我們每一個都是罪人.
究竟哪個比較騎呢?.
是變性人比較騎呢?.
還是我們這群又貪心,又說謊的人比較騎呢?.
今天馬太這個騎呢加譜在提醒我們.
我們這群基督徒.
並不是一群以一個民族為中心的人.
並不是一個以血統來連繫的人.
我們這群基督徒.
是一群追隨耶穌教導的耶穌的兄弟.

$^{761}$我們所效忠的是耶穌基督.
這位大衛君王彌賽亞.
上帝的揀選不是因為我們的血統.
不是因為我們厲害.
不是因為我們的背景好.
不是因為我們富幹或者武幹.
不是因為我們有顯赫的背景.
上帝的揀選是沒有原因,沒有條件.
我們是一群平凡人.
但上帝正正是揀選我們這群平凡人.
他在告訴我們.
我們只不過是一個騎呢人.
我們和其他騎呢人一樣.
都需要上帝的救恩.
以至我們能夠繼續延續.
耶穌所重新定義的以色列的歷史.
讓我們一起去禱告.
親愛的天父多謝你.
讓我們今天有機會去.
細讀馬太福音的加譜.
主今日你提醒我們.
我們正在延續以色列的歷史.
我們延續的並不是由血統組成的民族.
我們延續的是.
由耶穌你所重新定義的以色列.
主你讓我們在這個時代的洪流當中.
仍然可以選擇去效忠基督.
去效忠你.
讓我們知道我們所效忠的.
並不是任何的民族主義的思潮.
不是任何愛國主義的思潮.
不是任何種族主義的思潮.
因為教會是一個跨民族.
單單以基督為中心的一個群體.
一個新的以色列.
主你特別去提醒我們.
當我們講到揀選的時候.
讓我們了解.
我們只不過是一個平凡的人.
今天上帝的恩典覆蓋我們.

$^{801}$讓我們重新去認識什麼是恩典.
讓我們可以明白.
我們只不過是一個奇異的人.
我們需要的是你的拯救.
以致到我們走出這個社會的時候.
我們不是在看這個社會.
全部都是奇異的東西.
我們不再只是覺得北美.
是一個奇異的地方.
是一個敵對主的地方.
而我們看到這些奇異的東西.
同樣都可以彰顯你的恩典.
因為上帝你都要救贖這些地方.
你都要救贖這些人.
因為我們同樣都是這麼奇異.
我們同樣都是這麼一顏一乏.
需要你的拯救.
主多謝你求你更新我們的教會.
讓我們有一個新的目光.
去看我們教會自己的使命.
和去看這個社會.
今天馬太福音的家譜就提醒我們.
讓我們在這個時代.
在這個社會當中.
繼續延續耶穌的教導.
\newpage



\section{創世記}
\label{sec:lNtIgECCCGM}
\textbf{創世記第6章 (6.4) - 背景釋經  [洪水問題 \#1] (陳崇基牧師)}
\newline
\newline
連結: \href{https://youtube.com/watch?v=lNtIgECCCGM}{\texttt{https://youtube.com/watch?v=lNtIgECCCGM}} ~~~~ 語音日期: 2022-05-28
\newline
\newline
\hyperref[sec:ytVA464zDA8]{\small{< < < PREV SERMON < < <}}
~
\hyperref[sec:index_chronic]{\small{[返順時目]}}
~
\hyperref[sec:index_scriptual]{\small{[返順卷目]}}
~
\hyperref[sec:0vMb7YIM6HI]{\small{> > > NEXT SERMON > > >}}
\newline
\newline
$^{1}$我們繼續講洪水有關的背景.
上一次我們有少少的基礎.
我們講到一些聖經考古學的東西.
為什麼需要講這些東西呢?.
為什麼不直接講經文呢?.
因為背後有些背景.
或者有些工具我們需要了解的.
所以特別講到經文和歷史之間的關係.
我們上次提過考古學.
可以幫助我們最重要的是背景.
有時候確認東西未必可以確定很多.
確定到有些歷史性.
但是裡面有很多屬靈的真理.
神學真理都不能夠用歷史方法.
科學方法,考古方法來證明出來.
但是有時候.
一些考古學的發現或者一些證據.
等我們羅列了之後.
其實可以修正到我們.
幫助我們更加清楚去看聖經.
所以我們一直在講.
我基於這本書開始.
這本書就是《Lost World of the Flood》.
我相信現在只有英文.
但是它的分段很簡單.
每一個的提案,每一個命題.
講到創世記.
特別是講到洪水的故事.
我們可能從小到大聽了這個故事的時候.
或者一些小孩的聖經.
而我們覺得我們有一個印象.
但是我們要知道.
由我們第一章開始創世記.
我們就說聖經是古代近東的文獻.
當然有人說.
我不單單是古代文獻.
不單單是古代近東的文獻.
當然啦,這個是上帝啟示.
但是你要知道上帝啟示是藉著什麼去啟示出來的.
就是藉著當時的人寫作的手法.

$^{41}$當時的背景,當時的文化,他們的世界觀.
他們所屬的觀念.
來表述出來的.
所以我們要了解創世記.
縱使我們說就是神的話語永遠都通.
放諸於死海而皆準是永遠的真理.
但我們都要明白到底我們怎樣去看古代文獻.
所以我們要知道他們運用的手法.
他們怎樣去用不同的方法來表述真理.
以及表述真理的哪方面.
所以我們知道這樣的時候.
我們就需要一些工具.
我們就說比對研究和一些考古可以讓我們發現的.
例如米索不達米亞.
近東的洪水故事.
不單單是羅亞才有的.
不單單是創世記才有的.
有很多近東的故事.
我們繼續看下去的時候就會非常相似.
是跟羅亞的洪水故事.
因為很多時候我們沒有接觸過的時候.
我們就會覺得很嚇怕.
為什麼會這麼相似呢?.
是不是聖經抄襲他們呢?.
或者有很多坊間的說法就是.
一定是聖經抄襲這些遠古的例版裡面的內容.
這個就是我們慢慢慢慢才能夠去處理到的.
我們不能夠一時間就處理的.
所以我慢慢一步一步去解釋的時候.
之後我們會在不同的session.
我們就會說一下比對研究的歷史.
就是以前碰過什麼釘.
做錯過什麼.
學者不太明白的文化是怎麼運作的.
到現在我們明白了.
在做比較的研究的時候.
我們就不會怕.
原來這些聖經.
創世記等等的經文.
為什麼會跟這些近東的所謂.

$^{81}$虛無的神話.
好像很相似.
但是其實也有很多差別的地方.
所以我們慢慢一步一步才能夠去解答到的.
所以我們想看的主要是第九個命題.
就是這裡的作者.
就是特朗普和約翰·沃頓.
這兩位在美國裡面最著名的舊約學者.
他們說洪水是一個很大的災難.
但是不會是全球性的.
但是被描述成為一個全球性的.
以至達到一些神學的目的.
創世記神學的目的.
第四部分就是15至17.
我們只看第九個的命題.
我們繼續看下去.
這裡就說到.
被描述成為一個全地.
全球的洪水.
想達到什麼目的呢.
上次已經說了.
在科學,地質,考古上.
甚至這本書也有基督徒的地質學家寫了.
還是很有資格的基督徒的地質學家.
寫了沒有這樣的證據.
去證明一個全球的洪水.
我們很怕因為傳統是這樣的教導.
傳統是這樣的印象.
但是我想先告訴你.
我們由第一章已經開始說了.
聖經在上面.
聖經無誤.
但是聖經的傳釋是有很多的.
聖經的傳釋是有1,2,3,4不同的方法去傳釋.
聖經的傳釋.
某一個對聖經的傳釋.
不等於聖經的絕對權威.
聖經無誤.
不等於我們的傳釋就是無誤.
所以我們的傳釋是可以更改的.

$^{121}$正如我所說的.
聖經考古學.
或者我們以往一直看的經文等等.
其實是可以修改的.
可以更加完備的.
好像一個小朋友在這裡長大.
我今天對創世記的看法.
特別可能只是說第一章.
或者說到這裡.
洪水的故事.
我以往十幾歲的時候.
甚至我大學的時候.
我所解的或者了解的這段聖經.
這些的傳釋.
和我今天的傳釋.
我希望一定是不同的.
還有繼續進步.
繼續完備一點.
當我們越來越多工具.
越來越多資訊的時候.
特別我們生於這一代.
21世紀.
我們有很多近東的資料.
我們不能夠忽視的.
所以為什麼會說洪水的技術.
的描述.
其實不應該是全球的呢?.
其實在經文裡面.
也有些蛛絲馬跡.
讓我們看到的.
只不過我們以往看聖經的時候.
我們略過了.
為什麼呢?.
因為我們以往有一種的傳釋.
來到聖經面前.
我們就認為.
這些不會構成問題.
或者我們不會留意.
現在我們慢下來.
所以為什麼要看原文呢?.

$^{161}$原文的好處就是.
我們不需要留意的.
希伯來文或者希臘文.
舊約新約.
我們不需要留意的.
最重要是有一個working knowledge.
還有我們看原文的時候.
最好的一個好處就是.
我們可以慢下來.
不要以為我們以為.
以往的看法就對了.
因為我們正在看的.
一向都是翻譯.
所以我們有原文的時候.
就慢慢地.
好像停了.
煞了一個break.
車子停了下來.
我們慢慢地.
一個一個字.
一個一個字母.
我們看清楚.
到底在做什麼.
今天我們看六章四節.
有一本好的聖經書.
都是講這一類.
很少人會留意的東西.
都很新出.
2016年出的.
Ian Proven.
他還是在region college教的.
我以前在region college.
在加拿大Vancouver.
溫哥華那裡.
他是教舊約的.
Ian Proven是一個很出色的英裔學者.
英國裔的學者.
他寫了這本書.
Discovering Genesis.
在那裡點出.

$^{201}$有很多我們經常看不到的東西.
所以這本書.
不失為一本很好的.
一本的聖經解經書.
六章四節在這裡說什麼.
在那些日子和以後的日子.
講到洪水.
洪水還沒來之前.
在那些日子和以後的日子.
有巨人在地上.
神的兒子和人的女子結合了.
就生了上古英武有名的人物.
有巨人這個中文引用的翻譯.
Hannifinim.
在LX$\times$七十事易本.
將他翻譯.
即是主前大概二百五年.
將希伯來文翻譯成希臘文的時候.
就用Gegentis這個字.
這個Gegentis這個字就引申有巨人.
但是原文.
希伯來文就是Hannifinim.
Hannifinim是什麼意思呢.
Hannifal意思就是.
跌低的意思.
或者墮落的意思.
可以這樣解釋.
但是無論如何.
這裡講到這一種的人.
叫做巨人.
或者Hannifinim.
提到一樣東西很特別的.
就是說那些日子和以後的日子.
就有他們在.
那就是說他們.
沒有被洪水消滅.
對不對.
這個是最自然的看法.
對不對.
我們不能夠說.

$^{241}$洪水之後有些人就給他這個稱銜.
叫做Hannifinim.
有其他人都是這樣叫的.
但是他們已經滅亡了.
不是的.
六章四節這裡的作者.
寫得很清楚.
沒有什麼自相矛盾的東西.
起碼他自己不覺得.
在那些日子和以後的日子.
很肯定在講洪水之前.
和洪水之後.
有這些叫做Hannifinim.
翻譯做巨人.
這些是什麼人呢.
如果你查一下Hannifinim這個字.
在舊約裡面出現過的地方.
就在摩西五經裡面曾經有提過.
民數記十三章你記得那個故事嗎.
那個十二個探子被差派進去迦南地去窺探.
是不是神給我們這個地方.
我們可以打進去等等.
回來報信的時候.
這十二個探子.
十個是不信的探子.
他說在那裡迦南地.
我們看見了巨人Hannifinim.
他們是什麼呢.
就是阿納人.
或者阿納族人的子孫.
Bnei Anak.
有一個Anak這個名字.
他們是什麼呢.
他們就是巨人.
他們就是Hannifinim的後裔.
就是說他們一代一代傳下來.
然後這十個不信的探子.
對上帝沒有信心.
他們說我們在他們面前就像炸蔓一樣.
這個民數記.

$^{281}$洪水之後.
就講回這些人一直都存在的.
現在我們進去窺探的時候.
就看到Hannifinim.
《申命記》第二章.
以前有以米人住在那裡.
這族人人數眾多.
身體高大.
像什麼呢.
像阿納族人.
Anakim.
其實都是同一個字.
只不過是son of Anak.
這裡就是Anakim.
這以米人就是阿納族人.
就是相等的.
名字可能不同.
也算什麼呢.
Harafim.
就是利弗欽人.
Hannifinim也是這種人.
也是這種人算.
起碼我們這樣叫他.
但是摩訶人呢.
就用另一個名字叫他.
所以是同一族人.
同一種人.
但是有不同的名字.
但是無論如何我們現在追蹤的是什麼.
在跟蹤Hannifinim.
到底是什麼身份呢.
然後呢.
我們去到《創世記》第十四章.
第五節.
這裡就是四王打五王的時候.
亞伯拉罕剛剛到迦林地.
在永久羅德的一場仗.
他就說基大魯瑪和同盟的王都萊鳥.
擊敗了什麼呢.
利弗欽人.

$^{321}$Harafim.
Rafaim.
第十五章又提到他了.
就是在迦林地裡面.
有這些人.
耶和華就跟亞伯拉罕立約.
我已經將這塊地.
迦林地將來一定會賜給你的後裔.
就是什麼人什麼人.
迦林人 赫人等等.
利弗欽人 Harafim.
就是同一群人.
然後呢.
三章十一節新領記.
又提到了.
迦林地裡面有的族人之一.
就是利弗欽人 Harafim.
又是這些和Harafim有關係的人.
漁民裡面.
留下了誰呢.
留下了巴山王河.
所以呢.
對於創世記六章四節的那位作者.
我們當他摸西了.
總之經文呢.
是寫了那些日子和以後的日子.
很肯定是暗示.
明示暗示.
他們是生存了.
就是生還於洪水.
那很奇怪.
為什麼我們耳目沒有提過呢.
耳目沒有留意過.
其實猶太人拉比.
早就留意了這個細節.
他們就去掙扎到底發生了什麼事.
你知道猶太拉比很有趣.
有時候他們會.
在解經上.
嘗試去.

$^{361}$smooth out(清除)這些差別.
將它解釋.
怎樣解釋呢.
其中一個猶太人的拉比.
就是這樣解釋.
他說原來不只是.
羅亞和他一家在方舟那裡生還.
還有誰呢.
就是巴山王河.
西行.
Orc.
King of Bashan.
巴山王河.
他就坐在方舟的排水溝底下.
一塊木頭.
他拉著船邊.
然後他對羅亞和他的兒子發誓.
永遠做他們的僕人.
只要他生還就行.
只要他不會在洪水死就行.
誰知道羅亞就開恩給他.
然後在方舟轉了一個洞.
每一天就給他吃.
他才會避留.
你記得這個巴山王河.
就是摩西大以色列人.
快要進入迦南的時候.
他還沒有到耶利哥平原.
在這個.
我們叫約旦河東的高地.
就是今天的約旦王國.
就是在這裡.
在這裡.
那麼.
我先看看.
OK.
在這裡.
看到嗎.
他們遇見巴山王河.
然後摩西將他們擊打.

$^{401}$然後就贏了他們.
然後他們就佔領了.
這個約旦河東的地方.
他們就在這裡.
這裡就是.
耶路撒冷在這裡.
然後耶利哥就在這裡.
他們就過了約旦河.
然後在耶利哥的平原去扎營.
然後攻進去.
大致上那個故事就是這樣.
無論如何這裡說的是.
他們遇見了巴山王河.
為什麼巴山王河又是.
又是那些所謂巨人呢.
但是摩西之後.
為什麼又遇見他們呢.
猶太拉比就這樣解釋.
就是說.
他們其實是生還在洪水那裡.
這是一個線索.
讓我們知道原來猶太人拉比.
在解經的時候.
他們也留意到這個細節.
只不過我們現代人.
我們沒有留意.
另外如果你翻查一下.
有一個第一世紀.
耶穌時代的.
猶太的歷史學家.
他叫約瑟夫·祖西弗斯.
他寫了很多東西.
有關他的民族.
猶太人的歷史.
和當時.
第一次反羅馬革命的時候.
那些的技術等等.
在他的Antiquities.
就是猶太古史.
其中一冊.

$^{441}$他就是從創世記一直說的.
說到洪水故事的時候.
他說什麼呢.
羅亞和三個兒子.
在大洪水之前.
在八年出生.
他們首先從高山.
下到平原.
他們就在Mountains.
他們就下來.
平原那裡寄居.
然後他們又說服.
其他平原的人.
原本住平原的人.
他們因洪水.
而非常害怕.
非常不願意從高地下來.
好像仿效他們的榜樣.
他們最初居住的平原就叫做士拿.
無論如何對於祖西弗斯來說.
他解這段經的時候.
你就發現.
很有趣,為什麼不理他呢.
為什麼約瑟夫沒有覺得有問題呢.
因為他這裡清楚說了什麼.
羅亞生還之後.
和他的兒子.
從山上.
方舟亭北山上.
下來平原.
然後說服其他平原上的人.
就是說有人生還了.
他們為什麼生還呢.
因為他們上了高山.
因為他們很害怕可能又再有洪水.
所以他們不想學羅亞他們.
就下來平原.
這個非常有趣.
這個就引申一個問題.
到底希伯來的聖經.

$^{481}$希伯來的經文.
裡面.
是不是真的要求.
非要.
一個全球的洪水.
不可以呢.
所以為什麼我們要看原文.
看原文裡面到底是怎樣.
描述.
和看看古近東的世界.
他們的世界觀是怎樣.
我們才可以了解.
這個經文背後.
的事件.
到底是怎樣.
是不是一定要一個全球的洪水呢.
所以在.
這個洪水的.
技術裡面.
有幾節.
很容易我們就會.
現代人.
我們一讀就讀了.
就會認為一定是全球的洪水.
我們就要.
慢慢慢慢.
逐節逐節解.
並且有時候我們需要看看原文.
到底這些翻譯是不是真的.
可以這樣翻譯呢.
還是有另外一個翻法呢.
所以我們需要在.
接下來的幾個session我們看清楚.
Again.
我們最重要想做的事.
就是我們想慢下來.
看聖經有時候我們看得太快了.
特別我們.
做牧師和傳道人.
我們看得太快因為我們以為.

$^{521}$應該是這樣解.
我們回頭再停下來.
我們慢慢來看.
先看.
七章第21節.
犯在地上有血有肉的犯在所有的人都死了.
所有的東西.
所有的東西經常出現的.
我們就會覺得是.
All.
All the people.
All the land.
All the earth.
只留下羅亞和他同在的.
方舟裡面.
這裡所謂All.
即是一些普世性的字眼.
都出現在.
起碼最清楚的這四段.
第一段我們要解決的.
就是犯在地上.
All.
這些普世字眼.
所謂universal的language.
所謂Lost world of flood.
都是這樣去.
這兩位舊學者都是這樣停下來.
慢下來來看.
首先我們問一件事.
就是有些普世的字眼在聖經裡面出現的時候.
是不是就一定代表.
Global.
整個地球.
在這裡舉出幾個例子.
第一個例子就是.
創世記41章56節.
所以我們快速地看.
到最後的藥室的故事.
藥室做了宰相的時候.
埃及裡面做了一個很高的官.

$^{561}$然後他就有了糧食.
發現當時發生了什麼事.
他說當時.
當時.
當時發生了什麼事.
他說當時.
饑荒遍滿天下.
All the faces of the earth.
Al cor.
pane haos.
這個是.
前置詞.
on有什麼意思.
在.
Al cor pane.
All the faces.
就是那個面.
地面.
地下.
中文翻譯就是天下.
就是說遍滿天下.
然後藥室就開了各處的糧倉.
讓人們來拿糧食.
埃及的饑荒甚大.
他說.
各地的人.
All the earth.
Cor haos.
都往埃及去.
到藥室那裡就扔糧.
因為天下.
地面.
就扔糧.
因為天下的.
饑荒甚大.
又出現這個字.
Cor haos.
All the earth.
這個事件.
藥室的時候.

$^{601}$有七年的旱災.
饑荒.
藥室是什麼時間.
這個我們可以計算.
藥室的時間是出埃及之前.
我們現在知道.
藥室的時間.
大概是1500至1650年.
為什麼可以計算到.
150年.
因為大概有100年.
在這段時間.
埃及的古代歷史.
曾經在這100年.
被外族.
迦南人進駐了.
然後去管了那個地方.
霸佔了三角洲.
的地方.
所以如果看藥室的故事.
有很多的細節.
你就知道.
很可能在當時.
那100年左右.
希克索斯的時期.
Hixel意思就是foreign king.
是外族管埃及的時候.
這個不只是聖經有這樣的暗示.
埃及自己的古史.
都有這樣記載.
是一個很令埃及人.
一段100年多的時間.
因為是給了迦南人進去.
然後他們.
接管了利羅河三角洲.
將他們推了上山.
利羅河的上游.
所以這裡就是一個.
Hixel的時間.
1650至1500左右.

$^{641}$主前.
這100多年的時間.
在Hixel的時候.
所以就算看創世記.
最後那幾章.
藥室的故事.
你發現藥室被賣到埃及的時候.
那時候的法老.
是對他們所謂的迦南人.
藥室本身是迦南人.
他們本身在亞國家.
到達的時候特別對他們友善一點.
特別對他們友善一點.
親切一點.
因為是同族裔的人.
真正的埃及人在哪裡呢?.
在這裡.
他們趕了他們到這裡.
1550BC的時候.
這群本身是埃及人.
就擊退了這些希特勒索斯人.
將他們趕走.
然後就統一了埃及.
所以是有一個.
大的歷史在背後.
藥室的故事.
所以上帝在人類的歷史舞台.
背後在做事.
所以就配合.
藥室的時間.
無論如何.
藥室的故事裡面有一個情節.
就是饑荒.
在這個時候.
全球不可能有饑荒.
在古代歷史裡面.
不會這樣發生.
如果你懂古代歷史的話.
你就知道.
當時的饑荒大概在.

$^{681}$我們不多解釋.
總之在這個時候.
在這段時間.
整個東地中海.
這個地方.
是有一些氣候的轉變.
有饑荒,有缺雨.
或者旱災出現.
這些考古可以考證出來.
所以.
我們知道.
這個.
饑荒.
真的可以在.
歷史的視窗.
這個窗戶裡面發生.
所以你如果將這些.
蛛絲馬跡去配合.
其實藥室的故事.
是配合到.
大歷史.
聖經背後的世界.
背後的歷史.
所以這個不足為奇.
但是很肯定.
這段時間不是整個地球.
都有饑荒.
因為有其他的古歷史.
告訴我們.
比如西歐,亞洲,印度.
這些地方.
有他們自己的歷史.
我們都有記載.
所以當時的饑荒.
是全球的饑荒.
但是聖經怎樣.
用什麼字句來形容呢.
Koharis.
All the earth.
所以你知道.

$^{721}$有時候聖經.
很多時候用一些.
所謂universal的語言.
一些普世的語言.
都不一定是全球性的東西.
如果你喜歡我們的影片.
就記得要讚好,留言和分享.
拜拜!.
\newpage



\section{創世記}
\label{sec:0vMb7YIM6HI}
\textbf{創世記第6章 (6.5) - 背景釋經  [洪水問題 \#2] (陳崇基牧師)}
\newline
\newline
連結: \href{https://youtube.com/watch?v=0vMb7YIM6HI}{\texttt{https://youtube.com/watch?v=0vMb7YIM6HI}} ~~~~ 語音日期: 2022-05-28
\newline
\newline
\hyperref[sec:lNtIgECCCGM]{\small{< < < PREV SERMON < < <}}
~
\hyperref[sec:index_chronic]{\small{[返順時目]}}
~
\hyperref[sec:index_scriptual]{\small{[返順卷目]}}
~
\hyperref[sec:sMa4o_PUcpQ]{\small{> > > NEXT SERMON > > >}}
\newline
\newline
$^{1}$耶利米埃國主持586年耶路撒冷聖殿被毀.
說到四處有很多敵人.
耶和華發怒的日子.
無人逃脫 無人生還.
全部都殺盡.
是不是真的耶路撒冷所有人都死了.
是不是猶大地的人所有人都死了.
不可能 事實上聖經其他的章節都告訴我們.
還有很多人留下了.
雖然被擄 有些人被攻陷.
但是有很多人在那裡.
所以又是所謂universal language.
一些普世性的語言.
西班牙一章第二節.
說到神要去審判以色列民.
我必從地上徹底地去除滅萬物.
所有的東西 core.
除滅所有的活物.
我必從地面裡剪除它們.
然後殺死所有的居民 core.
在那裡除滅它們.
這裡又是說586年.
上帝的審判來到的時候.
很肯定不是每個人都死清光.
為什麼作者要指出這裡呢.
Norman和Walton.
因為我們的解經一定需要一貫 一致.
我們不能夠說那裡有普世性的語言字眼.
然後我們就說全球.
這裡就不是了.
這裡我們就把它解釋掉.
這樣解釋就很生硬.
所以我們要知道聖經的原則.
是需要consistent 需要一致.
如何處理這些core auras.
in all the earth.
即是普世universal的字眼呢.
我們需要知道.
聖經很多時候記載一些大的災難.
一些很厲害的災劫發生.

$^{41}$浩劫發生的時候.
很多時候都會用一些這樣的語句去形容.
比如耶路撒冷被毀.
天下有饑荒.
沒有可能.
信了也好 還沒信的歷史學家都會告訴你.
根本沒有可能.
1500-1600BC約瑟的時間.
居然是全球的饑荒.
如果你告訴別人真的會笑死人.
所以我們在解經上需要留意這些.
所以有些人就說.
沒有可能.
如果你要證明洪水是全球的.
你一翻開耶穌的說話就知道是全球的.
因為在巴塞科罕第十四章說到.
人子來的時候.
耶穌的second coming.
沒有人知道那個時辰.
羅亞的日子怎樣 人子來臨也要怎樣.
他就說.
因為人子來的時候一定是全球的.
所有的人都會知道.
所以羅亞的日子是這樣的.
人子的日子也是這樣.
人子是全球的 羅亞不是全球.
但是你看清楚夏文說什麼.
夏文說的不是.
the extent of the flood.
不是重點.
不是洪水有多大.
是不是大到全球.
他說洪水來的時候.
那些人有沒有準備好.
那個焦點是人.
是people.
不是說地理的東西.
你看清楚38節.
洪水以前那日子.
人照樣喝駕取.

$^{81}$直到方舟的日子.
就不知不覺洪水來到.
將他們沖走.
人子來臨的時候也是這樣.
為什麼呢.
因為兩個人在田裡面.
一個被接去一個被撇起.
兩個女人推磨.
一個被接去一個被撇下.
所以你們要警醒.
因為不知到的時候 主作利落.
想不到的時候 人子在那裡.
所以說的是我們有沒有準備好.
我們是不是ready.
不是說羅亞的日子是這樣.
你知道說到怎樣的時候.
比較的時候.
在比較什麼呢.
就是在說洪水來的時候.
那些人沒有準備好.
喝駕取.
by the way.
這裡是華人教會裡面.
很多時候用來講備題的經文.
不知不覺洪水來到.
就將他們全部沖走.
那時兩個人在田裡面.
一個被接去一個被撇下.
兩個女人推磨一個被接去.
或者被沖去.
一個被撇下.
我們很多時候就搞了備題神學出來.
什麼意思呢.
就是說七年大災難之前.
災前備題.
所以就會被接去.
你看清楚這個原文.
是和備題神學講的東西完全相反.
因為這裡被留下是好的.
因為被接去是什麼.

$^{121}$就是被洪水沖走.
你不要想著.
哦 備題.
一個消失了.
消失的是被審判.
撇下留下的是好的.
為什麼.
因為在說洪水來.
所以想清楚.
我們以為這裡是備題神學的證據.
誰知道原來不是.
剛剛相反.
所以備題是後話.
就是末世論的問題.
所以這個.
我想指出.
這段經文不要拿來做.
以為是支持備題論.
一點都沒有.
相反.
被提去的.
被接去的.
被沖走的.
那些是接受審判的.
所以這裡說的是人要準備.
所以我們也要知道.
那些普世字眼要怎麼處理.
有沒有研究給我們一些資料.
讓我們知道近東的歷史記載.
都會這樣.
當然有.
這裡我引用一本近東的學術書.
叫做Mesopotamian Cosmic Geography.
這個學者.
Helouris.
他專門做近東的研究.
他就拿近東的文學.
特別說到地理.
就將它收集了.
在這本學術的論文裡.

$^{161}$其中一個文獻.
就叫做薩爾根地理的文獻.
薩爾根地理文獻.
這是一批的文獻.
什麼時候發現等等.
這個屬於主前八世紀.
當時有一個Sagan II.
以塞爾書也有提過這個人.
這個亞述王.
還有他的後人.
他留下了一些記載.
他講到他自己帝國的時候.
講到亞述帝國.
他說亞述帝國是什麼.
薩爾根.
統治全世界的地方和人民.
薩爾根是宇宙之王.
征服了天下全地.
從日出到日落.
全地的總和就是薩爾根.
宇宙之王三次征服的地方.
所以又是用一些普世的字眼.
你說他懂不懂地理.
怎麼會講出這樣的字眼.
但起碼在他當時所知的世界裡.
在當時亞述近東的世界.
他征服了所有的版圖.
所以他也用了一些天下全地的描述.
我們再看看當時近東的人.
包括聖經世界裡的人.
他們怎麼看他們的宇宙.
他們的宇宙結構是怎樣.
他們怎樣去形容.
我們講過當時他們想的宇宙.
就是有大水包圍著.
中間有一個空間.
你記得第二天創造的時候.
創造什麼呢.
將水分開.
然後有一個泡泡出來.

$^{201}$這個就是創造主當時近東人.
用來描述創造.
為什麼有這個世界.
那種方法將水分開.
上面有一個窮倉.
窮倉在希伯來文.
意思就是硬的東西.
一些鐵片等等.
他就頂住上面窮倉的水.
讓它不會倒下來.
不會回落.
所以如果大水沖進來的時候.
就毀滅這個世界.
他們是這樣去概念化他們的宇宙.
當然不是說我們今天要照辦主管.
但當時的人 包括聖經世界的人.
是這樣去想像他們的宇宙.
所以聖經裡面很多時候.
用這些字眼 用這些概念.
來跟他們溝通.
埃及人都是一樣.
他們去描述他們的宇宙的時候.
上面有一個女的天神.
下面有一個地神.
中間有一個空氣神.
叫Shu.
他就去頂住 將水分開.
所以他們在主前 即是摩西的時候.
主前1500年左右.
即是大汗摩西的時候.
當時新王國 New Kingdom Egypt的時候.
他們就是這樣去形容他們的世界觀.
跟這個畫出來的簡圖都是一樣.
再看下面.
到主前700年.
巴比倫找到的宇宙地圖.
他們怎樣概念他們的宇宙.
他們的Cosmos.
他們的Cosmology是甚麼呢.
都是一個平的碟.

$^{241}$然後有水包圍著四面.
你將這個巴比倫宇宙觀的地圖.
畫了出來.
然後我們簡化了它.
裡面有很多標籤.
其實這裡有很多文字.
解釋這些標籤.
你見到這個其實是巴比倫的地方.
他們的世界.
中間巴比倫首都就在這裡.
中間有條河.
就是有法拉底河.
所以巴比倫的西面.
就有一條河剛剛穿過.
所以這個巴比倫.
然後有其他亞述在上面等等.
有其他的Geographical Feature.
有這些特色在這裡.
四周圍最重要的.
我們今天想看的.
就是海.
我們今天可以形容它.
或者這本書.
就會這樣形容近東的地理.
就是這些Cosmic Ocean.
就是再看剛才那個圖.
就是這些周圍包圍著的水.
他們都是這樣去將它形象化.
去概念化.
這塊地.
外面的就全部都是海.
其他的箭咀就是這些方向.
日出不見 日頭不見之處.
即是日落之處等等.
所以你見到是一個平面的世界.
他們所知的世界就是這樣.
所以洪水.
都是在說一個他們當時.
所謂known world 所知的世界.
所以剛才我們看了一些細節.

$^{281}$原文的細節.
近東的地理觀.
或者Cosmology 宇宙觀.
你見到的時候.
創世記不一定是要求一個全球的洪水.
相反 它所指的.
如果你有背景.
有原文的經文.
當然有很多其他的經文.
我們還要處理.
這是後話.
但是指向的是一個known world flood.
所知的世界.
有一個大的洪水.
一個大的災難發生.
如果你想有一本好的識經書.
又看得懂英文.
這本我覺得是最好的.
這本NIV Application Commentary.
John Watkinson.
他特別要做的好事.
就是分析這些經文的細節.
所以近東的背景下.
加上我們一些原文的細節.
經文要表達的是哪方面的真理.
好像我們看創世記一章.
六日創造 七日創造.
他想表達什麼真理.
千萬不要搞錯.
不是想表達科學真理.
也不是想表達地質學.
剛才提到.
Rakia 窮倉.
是一個硬的窮倉.
字本身就是一個.
叫做凸出來的東西.
凸出來的鑊.
這是古人用的.
一些很簡單的現象性的語言.
不是一些科學的語言.

$^{321}$窮倉的窗戶打開了.
有水掉下來.
有水降下來.
大淵的泉源崩裂了.
這些都是一些很現象性的語言.
不是說Geology.
不是說Science.
也不是說我們.
聖經也不是說.
近東的世界觀就是唯一的世界觀.
我們今天知道窮倉是通的.
上到太空.
也知道地球是圓球.
以前他們的宇宙觀.
他們的世界觀不是這樣看的.
所以聖經不是教世界觀的.
聖經教什麼呢?.
是神學.
是藉著當時近東的概念.
語言等等.
在講最重要的是什麼?.
神學.
神學才是最重要.
我不是說歷史背後沒有什麼.
而是聖經想表達的.
不只是單單告訴我.
告訴你.
某年某月某日發生了某件事.
是重要的.
但他想告訴你.
那個詮釋是什麼?.
怎樣解釋這個歷史裡面.
曾經有的大災難.
什麼原因有發生.
為什麼會發生.
神在做什麼.
最主要想啟示的真理.
在哪方面.
是神學的真理.
不是在講地質.

$^{361}$不是叫我們建構科學出來.
沒有這樣的事.
是宗教.
如果你看洪水的結構.
洪水敘述的結構.
你大概會明白.
大概會明白.
他不是在講現代人歷史的寫作方法.
或者不是在講現代人的地質.
歷史等等.
你看到這些數字.
7天,7天,40天,150天.
到下一節.
又提到這幾個數字.
就好像一個線形的結構.
或者我們叫它Chaostic Structure.
為什麼叫Chaostic Structure.
就是交叉的結構.
就是很多時候近東人寫的東西.
包括這裡.
洪水的故事.
很清楚.
ABCCBA.
如果你這樣寫出來.
C對C B對B A對A.
就是交叉.
所以我們就用一個Chaostic來形容.
三文治的結構.
或者線形的結構.
你就知道他寫的是什麼.
就是上帝,紀念,羅馬和大地.
因為方舟在世大地.
所以他要表達的是什麼.
是Theology.
有很多人問.
既然是局部洪水.
就叫人走上山.
為什麼要弄一個方舟.
為什麼要搞120年弄一個方舟.
我上次解釋了.

$^{401}$120年不是用120年弄方舟.
不知道用了多久.
為什麼要弄一個方舟出來.
我們就用現代人的Logistic.
事情能不能安排到.
這種角度.
你也可以看這個故事.
方舟到底是什麼.
為什麼方舟要在大地上.
一些代表性的活物.
要放在裡面.
方舟到底是什麼.
方舟不只是避難.
我們在想.
如果是局部的.
為什麼不走上山.
就弄一個方舟這麼麻煩.
這不是麻煩的問題.
是在寫神學.
在表達神學.
在事件裡面表達神學.
特別是跟近東其他的洪水故事.
都有其他方舟.
他們都有說到有隻船.
要跟他們爭辯什麼.
去解釋那個災難.
等於我們美國.
十幾年前曾經有911.
911發生了之後.
有很多不同的傳息.
為什麼會發生.
發生了什麼事.
上帝在哪裡等等.
有很多不同的傳息.
聖經就說.
這個才是真正的傳息.
就要寫這個大洪水.
背後發生什麼事.
這個方舟其實就是宇宙的縮影.
就是上帝在救贖.

$^{441}$更新大地.
不只是走難.
我們現代人一味想走難.
方舟做的就是re-creation.
這個是後話.
我們繼續解釋下去的時候.
我們就更加可以清楚一點.
但是我們需要忍耐一下.
無論如何.
這個是70年代的時候.
已經有人.
有一位很好的創世記的解經家.
叫做Gordon Bramham.
是英國裔的解經家.
今天還在生.
他寫了另一本創世記的識經書.
比較學術一點.
比較難一點.
但是無論如何.
他寫得很好.
70年代的時候.
他和其他學者.
其實都已經留意到.
這個洪水的敘述.
其實是像一把扇一樣.
一個很大的.
史詩式的.
很偉大的.
很高尚風格的.
一個表述的寫作.
看到A提到羅亞.
然後最後又提羅亞.
提到三個兒子.
然後又對下面有三個兒子.
剛才我說過了.
七天七天.
四十一百五十.
又有.
然後再反過來.
就是想讀者.

$^{481}$或者聽故事的人.
去留意到.
關鍵的地方在哪裡.
八章第一節.
Son remember.
這裡我引用一本.
很舊的書.
但是也很好的書.
就是講到古代近東的.
希伯來的寫作技巧.
等等.
修字等等.
他說這個Kaistic opposition.
就是這些扇形的形狀.
其實是一種high style.
就是用回當時近東的那種epic.
很厲害很宏偉的史詩.
好像我們今天聽的.
一些很偉大的交響樂曲.
那種表述的方法.
其實創世的第一二章也是這樣.
那種偉大是要這樣去表達出來的.
好像一個交響樂曲.
去表達出來.
用一種epic的style去表達出來.
當你聽到的時候.
嘩 這是一個很宏偉的故事.
裡面講的神和介紹的神.
也是一個很偉大的神.
在洪水氾濫當中.
他作為了為王的神.
這次這一節我們先講到這裡.
如果你喜歡我們的YouTube頻道.
就記得Like,Comment和Subscribe吧.
\newpage



\section{}
\label{sec:sMa4o_PUcpQ}
\textbf{比對研究 3: 必要, 發展, 方法   -  陳崇基牧師}
\newline
\newline
連結: \href{https://youtube.com/watch?v=sMa4o_PUcpQ}{\texttt{https://youtube.com/watch?v=sMa4o\_PUcpQ}} ~~~~ 語音日期: 2022-05-28
\newline
\newline
\hyperref[sec:0vMb7YIM6HI]{\small{< < < PREV SERMON < < <}}
~
\hyperref[sec:index_chronic]{\small{[返順時目]}}
~
\hyperref[sec:index_scriptual]{\small{[返順卷目]}}
~
\hyperref[sec:JkOQgNW0fh8]{\small{> > > NEXT SERMON > > >}}
\newline
\newline
$^{1}$好 我們繼續講背景 或者比對研究的發展.
上兩次講了必要性和發展的過程.
這些我們需要知道.
因為不是貿然有人想起.
為什麼要找近東的背景 文學來做比較.
上次我們講到找到這個Rosetta Stone.
先是法國人找到 然後去到英國人的手中.
所以今天去到英國倫敦大英博物館.
第一件展館裡面的物件就是這個Rosetta Stone.
我們說這個發現讓學者拆解古埃及的文字.
以至古代埃及文獻能夠被翻譯出來.
然後我們又提過在十一世紀末的時候.
這個George Smith拆解了紅水的泥板.
Gilgamesh Epic第十一塊.
說和聖經創世記紅水的內容非常相似.
我們可能有些困擾 但我們必須忍耐一會.
我們要了解背後的比對研究.
特別是將一些泥板 或者近東找到的文字.
和聖經對比的時候 有相似有不同的地方.
這個比對研究的進路 有它的歷史背景.
就因為這個緣故 我們需要了解發展的來龍去脈.
現在我們學界才知道 如何做正確的對比研究.
所以我們這次就說到1902年的時候.
1902年大概百多年之前.
在1月13日在德國的柏林.
有一個很出名的學術講座.
這個講座正是當時德國創立了近東學會.
Oriental Society.
在1902年1月13日 有一個這樣的講座.
一個發表的演講 就是Frederick Delitzsch.
德國的學者.
他是一個很出色的近東和聖經的學者.
他在1902年的時候.
在這個學會發表了一個講義.
這個演講叫做Bible and Bible.
就是巴別語聖經.
巴別代表巴比倫等等.
他講這個演講的時候.
正是有德國的帝王William II出席.
當時是一個很大的事.

$^{41}$他在這一年開始就去講幾場的講義.
有關聖經 有關近東背景的學術演講.
在他的講義開始的時候.
他就說 為什麼我們.
他在說我們德國人 或者歐洲的人.
千辛萬苦跑到老遠荒涼危險的地方.
到底是為了什麼呢.
要花這麼多的資金 在幾千年的歷史裡面.
在這些廢堆 考古的土壤裡發掘.
又不是為了賺錢 或者是賺錢和保障.
到底是為了什麼呢.
為什麼當時的各個國家.
歐洲 英國 法國 德國 美國等等.
去到這些早已被遺忘的土壤裡.
去爭取這些發掘權.
另外大西洋的兩岸.
為什麼有這樣的興趣和捐獻呢.
有增無減的.
當時正是近東學 或者亞述學.
開始掘起的時候.
就要在巴比倫 亞述的地方來發掘.
他回答這麼多的問題.
第一個說話 解釋背後的動機和目的.
只有一個答案.
縱使不是全部的答案.
那個答案就是聖經.
因為當時剛剛開始.
打開了古近東聖經的背景 這個世界.
正因這個原始 這樣的原因.
所以當時的歐洲學者.
古代近東聖經研究的學者.
就去到這些地方 或者支持這些發掘的工作.
和將這些找到的文獻翻譯出來.
希望能夠找回這個失落了的世界.
就是聖經原本 特別是舊約聖經原本的世界.
聽到他這樣說的時候.
你就發現好像對聖經很有善.
好像幫助到聖經去了解.
但是聽這個講座 或者現在不錄.
我們看回當時發生的事.

$^{81}$或者他演講的內容的時候.
你會知道Frederick Delos.
其實在說什麼呢.
找到這些近東的文獻的時候.
他就在講學上說.
他說舊約的以色列文.
就是希伯來聖經和他的文獻.
只不過是這麼多古舊的國家.
這些遠古的文明.
一個最年輕的一員.
就是說以色列是最後最後才出現的.
這些巴比倫 所謂巴別 比布 阿述.
其實比聖經遠古得很久.
他只不過是最年輕的一員.
所以舊約 以色列的道基督教導 宗教等等.
對巴比倫並非高尚 並非超越他們.
所以他就認為 特別是舊約聖經.
他們所講的東西.
全部都是比近東的文明 巴比倫等等.
是比他們更加後期的原故.
所以他們是後期的產品.
並且他去到說.
聖經的作者寫作的時候.
舊約聖經的作者.
譬如講到洪水的故事等等.
他們其實有一個巴比倫的文獻在他們面前.
並且聖經的作者是在抄襲巴比倫的文獻.
所以你一聽到的時候.
你會發現 開始的時候好像是很想明白聖經.
但其實不是的.
其實Delish想講什麼呢.
他說巴比倫的文明 或者近東的文明.
在思想上和歷史上是為首的.
他們是先行的.
他們是先創建的.
然後以色列的思想才出現.
希伯來聖經 先知等等.
他們才會講他們的經文 舊約聖經所講的內容.
所以他們是後期的 是落後的.
所以對於Delish來說.

$^{121}$他就認為 現在我們要再一次.
因為有這樣的例版 有這樣的背景.
我們要再一次衡量希伯來宗教信仰的高尚性.
他就認為一點都不高尚.
其實舊約的宗教只不過是後期的抄襲的產品.
換句話說 巴比倫 近東才是superior.
才是高尚一點 因為他們是開頭來的.
後期來的以色列的宗教 是低劣一點的.
是差一點的 是後期一點的.
因為這樣的緣故.
他們幾場的講學.
就引起當時德國和歐洲的學界.
特別是聖經學界 聖經研究學者.
就有很大的爭論.
有很多人就非常之反感.
更加反感的就是.
Frederick Delish.
其實是當時剛好已故 剛好已過世的.
一個很出名的路德宗的色經家.
他爸爸就是Franz Delish.
Frederick Delish講的Bible of Babel.
他是他的兒子.
如果你對聖經的解經的資料 解經書.
有些認識的話.
你會知道Franz Delish.
一個德國的解經學者.
是一個很保守的基督教路德宗的解經學者.
如果你知道這些色經書的時候.
可能你會遇到一套是這樣的.
是靠Kyle和Franz Delish.
一起寫一套舊約的注釋書.
這一套注釋書 今天還有人用.
百多年後還有人用.
並且是一個很保守的一套色經書.
也是很值得參考.
百多年前 比較舊很多.
當然他們當時沒有很多新的數據.
沒有新的方法等等.
但是在Bible of Babel的講座之前.
Frederick Delish的父親.

$^{161}$就是Franz Delish.
為什麼呢?.
Frederick Delish講完他的《巴別與聖經》演講之後.
很多人有很大的反應.
令很多人感到很驚訝.
並且感到非常之反感.
因為他的立論.
就是說巴比倫在思想和歷史上.
都是為首以色列為後.
他認為有一個現象.
就是Pan-Babylonianism.
即是說是泛巴比倫主義.
換句話說.
你將巴比倫近東的文獻.
包括埃及等等.
近東文明留下的文獻.
和以色列去比較的時候.
你會發現很多相似的地方.
相似對於Frederick Delish來說.
肯定是抄襲.
既然是抄襲的話.
就一定是掠貨.
他們是以次創作也好.
或者是跟人家瘋也好.
所以以往一直以來.
我們說18世紀或之前.
一直以來我們對聖經背後的世界.
是所知甚少.
只有舊約聖經才知道.
巴比倫的宗教.
或者近東的世界.
但是只有舊約聖經的時候.
我們就知道.
舊約聖經很多時候.
先知都會去批判巴比倫.
或者別臣.
外邦的宗教文化等等.
所以只有舊約聖經的時候.
18世紀或之前.
你會覺得以色列的宗教.

$^{201}$是一種創新的宗教.
去批判其他人.
並且因為只有舊約聖經.
才知道聖經背後的世界.
所以我們這些學者.
在以前沒有近東的發現之前.
就會覺得希伯來的信仰.
是一種很高尚的啟示信仰.
獨臣的信仰.
發現了這些東西之後.
這個Bible的lecture.
就將它扭轉了過來.
就是說希伯來聖經原來是劣貨.
原來是低劣的宗教.
真正高尚的.
要回去尋求的.
就是這些巴比倫的思想.
我們今天怎樣看回這個演講呢.
我們就回到1902年的時候.
為什麼這個演講這麼重要呢.
回看當時的背景.
原來當時是19世紀.
歐洲國家民族主義崛起的時候.
你知道Friedrich Dillard.
他是德國人.
在19世紀的時候.
德國的國勢開始起飛.
歐洲很多其他國家都開始爭霸.
因為有工業革命的緣故.
所以當時世界的super power.
那些超大的強國.
就是歐洲的大國.
比如大英帝國等等.
正是這個時候.
德國的民族主義就非常之高漲.
你可以叫他們當時.
因為工業革命他們是比較後一點.
英國先的.
然後才到法國.
然後美國等等.

$^{241}$德國是比較遲一點.
所以當時Friedrich Dillard.
去講這個演講的時候.
就是所謂的超英趕法.
或者超英趕美.
在民族主義這麼高漲的時候.
講到這些近東的研究的時候.
不單止在工業革命上.
德國是比較遲一點.
比美國,法國,英國遲一點.
他們在近東研究上.
都是比英法落後一點.
剛才我們講過.
Rosetta Stone是誰發明的.
是法國.
然後在中東開始發掘的時候.
哪個國家開始.
英國不是法國.
所以他們的涉嫌文字.
全部都在他們手裡.
德國的近東學會才剛剛成立.
1902年的時候.
成立沒多久.
所以當時德國和.
當時德國的帝王.
都想大力推動去中東發掘.
拿這些遺產.
當時的一個爭霸的形勢.
不單止工業革命.
需要去趕緊的.
去做霸主.
去爭霸.
連這些近東的研究.
都希望能夠去爭霸.
你記得.
他剛剛開始演講的時候.
講什麼.
就是為什麼要去發掘.
為什麼要去發掘.
他就是想德國的聽眾.

$^{281}$聽到的時候.
為了聖經的緣故.
去發掘.
或者去得到這些發掘權.
佔領權.
以致德國都可以冒起.
正是這個時候.
我們發現.
Frederick Deluge.
講了這個演講.
成為了一個比對研究的分水嶺.
當他講了之後.
就形成了很大的爭議.
為什麼呢.
因為在他演講裡面.
你可以很清楚看到.
他不單止認為.
希伯來聖經.
或者猶太教.
是後期的產品.
其實他認為是很不幸的.
為什麼會有舊約聖經.
或者有希伯來聖經.
因為當時的德國學者.
認為舊約聖經.
Old Testament.
是一些Old Testament的宗教.
是一些發展得後期一點.
和是律法主義的一種宗教.
所以他的演講裡面.
都有這樣的味道.
你也看到Frederick Deluge.
開始講學的內容.
都是一些反猶太主義的.
很不喜歡猶太人.
所以很不喜歡猶太人的經卷.
所以就選擇了舊約聖經.
或者希伯來聖經.
他呼籲大家去回歸.
既然發現了這麼多的泥板.

$^{321}$巴比倫的文明等等.
是先者.
是源頭.
所以就應該回歸.
米索不達米亞的根源.
不要跟舊約聖經.
因為舊約聖經是後期的產品.
因為當時也是德國民族主義崛起.
所以更加令到他們不喜歡猶太教.
或者不喜歡舊約聖經.
或者將這些東西全部放在一起.
如果你知道20世紀初的歷史.
就發現這種德國民族主義崛起的時候.
其實就是為了希特拉普魯.
1930年4年5代納粹黨執政.
所以這些聖經學者.
包括Deluge其中一個.
其實在神學上就是為了納粹德國黨派.
去冒起然後取得權力.
希特拉再出現.
為了這些勢力去做神學.
在聖經研究上的鋪路.
所以在這篇article.
Journal of Biblical Literature.
講到Centennial Review.
Deluge講Bible-Babel的lecture之後.
100年之後再回頭去批判的時候.
就發現原來當時的Deluge.
都是他自己時代的產品.
因為當時的德國民族主義上漲.
並且德國人認為他們的種族是最厲害的.
當德國人開始在工業革命趕到英國的時候.
很快地他們的經濟開始有發展.
變成了一些很大的鋼鐵的生產商等等.
他們就會認為德國人是很厲害.
德國人是最好的民族.
所以就有一種民族的歧視出現.
特別是他們歧視的是猶太人.
所以就知道之後有希特拉.
他拿了權之後.

$^{361}$就令到德國很憎恨猶太人.
然後就有反猶太的運動出現.
當然反猶太的運動之前也有.
但是到了德國希特拉的時候.
就更加去intensify.
更加濃厚.
並且有Holocaust大屠殺的事件出現.
所以我們發現.
我們的解經家 所謂的解經學派.
和他當時的時代也很有關係.
有時候會受到他時代的影響.
你記得這個圖.
我們說過解經 釋經的事業有三方面.
第一方面就是The world behind the bible.
聖經背後的世界.
第二方面就是聖經裡面的世界.
The world in the bible.
第三方面就是聖經前面的世界.
就是我們作為釋經者 作為讀者.
我們怎樣去認識聖經呢.
其實我們都會受制於.
或者受到我們這個時代的洪流所影響.
如果我們不小心的時候.
不要自我批判.
很容易就去到以往Delage所走的路.
受到他德國至上的大國文化的氣氛影響.
我們是最厲害的.
其他人都是低劣一點.
包括聖經的文明等等.
所以我們傳統有我們的背景.
這些都是需要我們自己去檢視的.
在Delage說完巴別與聖經的演講之後.
1902年之後.
就造成了一個很大的分野.
很大的爭論是什麼呢.
就是將巴別論的近東文獻.
他們怎樣說宗教 怎樣說創造等等.
與以色列的文獻 舊約聖經.
希伯來聖經Hebrew Bible的文獻作比較的時候.
就造成了兩個越來越分裂的情況.

$^{401}$越來越對立的情況.
第一個情況我們可以叫他做.
為了找相似地方的.
很著迷於找這件事.
Parallelomania.
Parallelomania的意思是.
找到巴別論的文獻.
與舊約聖經比較的時候.
有相似的地方.
就一味找相似的地方.
這些就成為了批判學派.
我們叫他們不信聖經.
他們就經常有一種這樣的趨勢.
那種趨勢或者傾向.
就是經常去找相似的地方.
找到相似的地方就是.
就是犯了巴別論主義.
巴別論是先行的.
所以以色列就抄襲他們.
另一個就是一班.
當時可以叫他們做confessional scholars.
即是那些本身已經是信聖經.
或者本身是研究聖經.
或者好像Frederick Dillard的爸爸.
Franz Dillard.
他信聖經.
他發現有些相似的地方.
其實都不行.
他應該有些不想去看相似的地方.
所以就形成了另一個趨勢.
就是parallelophobia.
會見到相似的地方.
我們就很怕.
因為這樣好像以色列的文獻.
是去拿走了它的獨特性.
即是見到聖經裡面的描述.
創世也好 洪水也好.
或者講到其他文明等等.
就會覺得很敏感.
就不想碰這件事.

$^{441}$所以這個比對研究.
由1902年開始就造成了兩個.
不需要有的對立.
那些critical scholarship.
就是一味去找相似的地方.
那些confessional scholarship.
或者信聖經的.
以聖經前設為先行.
他們就很怕這些東西.
所以就形成了兩種對立.
其實在今天我們看聖經的傳統.
講到教會.
很多時候都是偏向這種phobia.
這種畏懼狂 畏懼症.
很怕看到有這麼相似.
希伯來聖經就很不獨特.
上帝的啟示怎麼會是這樣.
並且上帝的啟示怎麼會遲一點.
為什麼先過他的外邦.
一些虛神的宗教.
或者文明留下來的文獻呢.
你記得我們講到洪水的故事.
就是Gilgamesh epic裡面第十一塊的禮版.
有17個地方是相似的.
並且故事的脈絡是差不多一樣的.
所以為什麼呢.
講到比對研究的時候.
很多時候我們不敢看這件事.
因為一比較的時候.
你會發現為什麼會這麼相似呢.
會不會聖經抄襲別人等等.
我們如果這樣看的時候.
你會發現好像Frederick Delage.
他的lecture裡面所講.
他說如果把兩件事比較有相似的地方.
就代表什麼呢.
代表以色列一定是抄襲巴比倫.
以色列一定是抄襲別神.
別邦 外邦的文獻.
就沒有了他們的獨特性.

$^{481}$變成了泛巴比倫主義.
通常都是這樣看法.
而且今天你可能會遇到一些人.
他們會說因為聖經洪水的故事.
跟近東的洪水故事很相似.
因為這個緣故.
所以一定是聖經抄襲他們.
今天坊間都經常聽到有些人這樣講.
是不是這麼簡單呢.
背後到底這些文學作品.
這些近東的資料.
跟聖經相比和比對的時候.
是不是背後的處境.
背後的關係.
是不是就是這麼簡單呢.
這個就下一次我們會繼續解答.
但是我們需要一點耐性.
是不是一定要這樣看呢.
找到A然後聖經代為B.
就一定是B抄A.
一定是A高過B.
文化是不是這樣的動態呢.
所以我們希望能夠繼續去解答這方面.
以致我們繼續看《創世記》1至11章的時候.
我們看到有很多相似的地方.
也有差別的地方.
應該怎樣去衡量呢.
如果你喜歡我們的YouTube頻道.
就記得Like, Comment和Subscribe.
拜拜!.
\newpage



\section{19:7-14}
\label{sec:JkOQgNW0fh8}
\textbf{Apostles' Creed \#3  (Psalm 19\_7-14) - Rev. Mark Chan}
\newline
\newline
連結: \href{https://youtube.com/watch?v=JkOQgNW0fh8}{\texttt{https://youtube.com/watch?v=JkOQgNW0fh8}} ~~~~ 語音日期: 2022-06-01
\newline
\newline
\hyperref[sec:sMa4o_PUcpQ]{\small{< < < PREV SERMON < < <}}
~
\hyperref[sec:index_chronic]{\small{[返順時目]}}
~
\hyperref[sec:index_scriptual]{\small{[返順卷目]}}
~
\hyperref[sec:s7F8M3fuVJY]{\small{> > > NEXT SERMON > > >}}
\newline
\newline
$^{1}$Hey guys, thank you very much for being here.
and we'll continue our series on Apostles' Creed..
And my quiz for you is what is the Apostles' Creed?.
What is it?.
It's an ancient document,.
has been around for almost 1800 years,.
and it's a summary of our faith..
It's a map of the Christian faith..
And if people ask you, what do you believe in,.
or sometimes you may be questioning,.
what do I really believe in, and what should I believe in,.
you go to a summary or a quick document.
like the Apostles' Creed,.
and then you'll find out the basics, the core,.
the must do, the must believe,.
the core beliefs of our faith..
Yesterday, I was blessed to attend a wedding.
in the West Suburb..
I've never been to that neighborhood,.
never been to that part of the Western Suburb,.
so I turned to Google Maps..
So I put in the destination, I put in the place,.
and I wanna find out where that is..
And of course, in our days and age,.
and it is very convenient..
You just key in the address.
and then you'll find out how to get there..
So it shows me about an hour driving,.
passing through different neighborhoods,.
those are the details, right?.
I don't really care..
And Google Maps doesn't show me that..
It really show me is how to take a turn.
and what road to take in order to get me there, right?.
That is the summary, that is the quick,.
you know, core knowledge that I need.
to get to my destination, right?.
So I didn't get lost, I came back one piece, right?.
So that is the good thing about summary.
and Apostles Creek is doing something like this..

$^{41}$Now, as the pandemic fades away, hopefully,.
I know we're all getting tired of it,.
but it is kind of being sort of under control, right?.
Those who are infected with COVID suffer,.
it depends on depending on your condition,.
suffer quite a bit..
Or some people suffer slightly..
But as the pandemic fades away,.
some of the people in this church are asking me,.
when are you going to Israel, right?.
When you're leading a group to Israel..
And then I talk to some people, I talk to some of you,.
and say, yeah, I think it's doable, right?.
If you get enough people, we can go next week, you know?.
The border is open and it's very convenient..
If you have a US passport, it's really convenient..
I still remember talking to people.
about leading a group to Israel..
I've been to Israel for more than 10 years,.
almost 20 times..
I still remember in the early days, 2008, 2009,.
when we go to Israel, it was not that convenient..
I still had that atlas..
Do you know what atlas is?.
Atlas of maps, right?.
No?.
Is that yeah?.
Atlas, you know, to print maps..
It's like a book, okay?.
It's like a book..
So it's different kinds of maps..
And then, of course, you have the big map, right?.
So when you travel, when you're driving,.
especially in places that you are not familiar,.
you need those maps, right?.
And then you have to plan ahead of time.
in order to know where to turn and what road to take.
and how to get to your destination..
But now, of course, it's a totally different ballgame, right?.
You go to the airport, you rent a SIM card,.

$^{81}$and then you plug into your smartphone.
and everything is there, right?.
You key in your address and you will get there..
So it helps you with your navigation..
Apostle's Creek is also a tool of navigation..
Sometimes you may wonder,.
I know many of you, or maybe most of you,.
have not really read the whole Bible,.
cover to cover, right?.
Now, when people ask you, what do you believe in,.
sometimes you may wonder, you have a question mark,.
hey, what do I really believe in?.
What is the heart of the gospel?.
What is the heart of the Bible,.
of the heart of Christianity?.
And it is the Apostle's Creek..
So last week, we were working through Psalm 19,.
the first part of Psalm 19..
Now we come to the second part of Psalm 19..
I told you the flow of Psalm 19 is very easy.
and it's very meaningful..
The first part is the words of creation..
The creation is singing and expressing their praises.
and worship to the creator..
That is what we call general revelation..
Everyone who walk in this world,.
every single square inch of this universe.
is the property of God, has a signature in it, right?.
Whether you look at the world in details,.
the micro perspective, or you look at the whole.
entire universe, the galaxies and all that stuff,.
the macro perspective..
Every square inch of this creation marks the property.
of the creator, God himself..
So the first six verses talking about general revelation,.
the universal revelation, everyone should know that..
Everyone should be able to listen to the music.
of the creation, singing about the creator.
behind all of this, the intelligent designer.
of the entire universe..

$^{121}$And then the second part is words from the creator..
And that is what we call the special revelation..
God is not just a distant creator..
He created the world, boom, and it's there,.
and have his own mechanism and laws and principles,.
and then the universe run on itself..
No, the Bible presents to us the God who is super involved.
in every single thing since the first creation.
until now, his hand is there,.
even though it's invisible to our own eyes,.
but God is always involved in the creation..
So we come to verses seven to 10, seven to 10..
Okay, let's read that..
The law of the Lord is perfect, reviving the soul..
The statue of the Lord are trustworthy,.
making wise the simple..
The precepts of the Lord are right,.
giving joy to the heart..
The commands of the Lord are radiant,.
giving light to the eyes..
The fear of the Lord is pure, enduring forever..
The ordinances of the Lord are sure.
and altogether righteous..
They are more precious than gold,.
much, than much pure gold..
They are sweeter than honey, than honey from the comb..
Now you can find three pairs of terms.
that describe the word of God..
Can you find them?.
Three pairs, so a total of six terms..
Six different technical terms from the Old Testament.
describing the word of God..
Can you find them?.
The first one being the law, right?.
Law..
Now behind this English translation, law, L-A-W,.
is the Hebrew word Torah, Torah..
Torah means more than, so much than the law, right?.
When we think of law today,.
well 55 miles per hour speed limit, right?.

$^{161}$Don't park there and all that stuff, right?.
You park there, your car get towed..
Those are the connotation that we have today.
when we hear the word law..
However, Torah means so much more than that..
Torah means instructions, direction..
You know, like your Google map,.
like a compass, compass on the map..
You know, showing you the instructions.
to teach us to go the right way..
The second term is testimony, right?.
Or a statue, right?.
Testimony, some of the English translation will do that..
Statue or testimony..
And the meaning of the word is reminder..
Are you forgetful?.
You are, right?.
I am very forgetful..
I'm very, very forgetful..
You know, my family members can testify to that, right?.
I'm very forgetful because I'm so focused, you know?.
I'm so focused..
I'm so focused on a few things, right?.
And then sometimes I need to forget other things, right?.
But testimony or statue is reminder..
Some sort of reminder in our daily life.
to bring us back to the relationship with God, right?.
Number three is precepts, precepts..
Now, all these six terms are pretty synonymous..
They could point to the same thing,.
but their connotation, their meaning.
is a little bit, you know, distinguishable..
You can just, by meditating upon different words,.
and you see how the word of God.
can be at work in your lives..
So number three is precepts..
Precept has the meaning of leading and overseeing..
Overseeing..
Are you a supervisor in your workplace?.
Or you work under a supervisor, right?.

$^{201}$Usually they suck, but..
Usually, right?.
Usually they suck..
But, you know, the word of God is the precept..
It's to give us the leadership or give us the guidance,.
especially we get lost or confused..
You go back to the word of God,.
he will give you the precepts..
Number four is commands, right?.
Or commandments..
Command is authoritative and clear..
And this, especially so, I remind you many times,.
the Bible is not written to us..
It's written to the ancient people in the Near East..
What they used to believe in are those false gods, right?.
Babylonian gods and Assyrian gods, Egyptian gods,.
many gods, but none of them,.
none of them is clear about what they want..
Sometimes they are like us, right?.
Because in the ancient Near East,.
in the pagan environment where ancient Israel was located,.
in their environment,.
the pagan gods are just like people..
Because the gods, the false gods,.
are made in the image of men..
All right?.
So you remember Genesis one..
We are made in the image of God, right?.
But in the pagan environment,.
the gods are imagined by people, right?.
That's why they are like people..
So they are unpredictable..
So commandment has the meaning of God being clear.
about his moral view..
His moral view is to help you live in holiness,.
righteousness, justice..
Everything is spelled out in the Bible, right?.
Everything is spelled out in this word..
God doesn't play hide and seek with us..
So sometimes when we say,.

$^{241}$oh, what is the will of God in my life?.
Sometimes we got the question wrong..
Of course, the will of God is clear..
He wants you and me to live in his image,.
to live holy, right?.
To be righteous, right?.
Your relationship has to be holy..
Your work ethics has to be holy..
Somehow you can display the wisdom of God, right?.
The image of God..
So commandment is number four..
Number five is fear, right?.
The fear of the Lord..
Now, if you're familiar with the wisdom literature.
in the Bible, which is right in the center,.
some Psalms and the book of Proverbs,.
Ecclesiastes, the book of Job,.
all of this are called wisdom literature..
They often have a phrase, the fear of the Lord,.
the fear of the Lord..
It means you have the right relationship.
with this holy God, right?.
Right relationship with him..
Then you are living according to his will..
The fear of the Lord, meaning wisdom..
You will live in your practical situation having insights..
And where that insight comes from?.
It's come from the Bible, the will of God..
Number six is judgment..
Now that has to do with God's judicial discernment..
If you read the first five book of Moses,.
sometimes the five books of the law, right?.
You find some case laws..
If this happen, then what is the judgment?.
If this happens, what should be the decision?.
And all those stuff..
Judgment has a sense of God's wise discernment.
of different situations..
So God's general will is always clear..
To have the people of Israel in the Old Testament.

$^{281}$and the New Testament church like us.
to be walking in his way..
And then this Psalm talks about the benefits,.
the effects of the will of God..
Reviving the soul, making wise the symbol,.
giving you joy, giving you fulfillment,.
enduring forever, altogether righteous..
So the Psalm is encouraging you to go into the will of God,.
to go into scripture, to find that treasure..
It takes some work..
It doesn't come to you automatically..
You need to be trained, you need to learn, right?.
This is what every disciple needs to learn, right?.
From the first day you believe in Christ,.
from the first day you acknowledge,.
I come to believe and faith in Jesus Christ..
I need to learn how to interpret,.
how to understand and apply the will of God, right?.
So that is what Psalm 19 trying to encourage us to do..
And that's why one of the core in the Apostles' Creek,.
I wonder if you can notice that..
If you read through the Apostles' Creek,.
it gives you the very core of our belief..
There are many things that can get you distracted,.
but these are the things that you always return back to..
These are the things that you build your foundation upon..
These are the most important things..
That is what the Apostles' Creek is all about..
I wonder if you read this, how many statements are there?.
12 statements, okay?.
12 statements..
Do you notice what stand out to you?.
Now, let me tell you,.
what is not being said in Apostles' Creek.
is how Christians behave..
You don't find this statement,.
I believe in the good Christians..
I believe in how good the church is..
You don't find this statement in our Apostles' Creek, right?.
There's not a single mention of the Christian testimony.

$^{321}$or the examples Christians have..
You don't build your faith upon people..
You don't..
I understand in recent years,.
churches, leaders, rise and fall..
I understand that..
Pastors, Christian leaders, they fail..
And sometimes it shatters our faith, right?.
Because so often, we build our faith.
on broken people just like us, right?.
We do..
There's a temptation to idolize institution,.
idolize the big institution, the big names, right?.
Here, the Apostles' Creek makes no mention.
how good the Christian behave, right?.
And what is being emphasized is what Jesus has done for us..
Out of this 12 statements, there are five statements.
on Jesus himself, on Christ alone..
In fact, most of the statements,.
if you take it by proportion, are about Jesus..
It really centers on the historicity of Jesus Christ,.
that he did something in human history,.
that Christ did something in our lives..
Five out of 12 statements are on Jesus himself,.
the historical Jesus..
He came to die on the cross..
He came to rise from the dead..
He came to accomplish our salvation..
He is our focus and he is our only foundation..
You see what I'm saying?.
You don't believe in the person of a Christian leader..
You don't believe in any institution,.
any institution you can see today,.
or any institution or organization, churches,.
parachurch organizations that have come to pass..
You don't believe in that..
They all rise and be gone..
What is eternal?.
What is most important is Jesus Christ himself..
God speaks and he reveals himself..

$^{361}$The Old and the New Testament are the historical record.
of what God has done in human history..
And God is still active, doing his things..
Maybe mysteriously, our eyes may not be able to pick it up..
But God is and God will be..
The focus is always on God himself..
So let me remind you, if anyone or anything.
that has disappointed you in your Christian walk,.
seriously, who cares?.
What's most important is Jesus Christ and his cross.
and his empty tomb..
Psalm 19 is telling us God is revealing himself.
in his word and the word ultimately.
is Jesus Christ himself, right?.
The word of God, ultimately, the revelation of God.
is his son, Jesus Christ..
That's why the creek is on Jesus, centered on Jesus..
So Psalm 19 is encouraging you,.
hey, put your faith in God and his word alone..
It ends with a prayer..
And I hope this is your prayer today..
This is your response..
The psalmist says, by them, the word of God,.
the revelation of God, right, ultimately is Jesus Christ..
Is your servant wanted?.
In keeping your word, there is great reward..
Who can discern his errors?.
Forgive my hidden sins, my hidden faults..
Keep your servant also from willful sins..
May they not rule over me..
That's interesting..
Then will I be blameless, innocent of great transgression..
May the word of my mouth and the meditation of my heart.
be pleasing in your sight..
Oh Lord, my rock and my redeemer..
You can pick up the attitude of the psalmist.
as he ends the psalm..
He realized his brokenness as well, right?.
He's tempted, a prone to great transgressions..
Just like us, right?.

$^{401}$Look deeply into your inner life, right?.
We are all ultimately broken..
We are all depraved..
We are all sinners..
I think about 400 years ago, there was a church reformer,.
what we call reformer..
His name is Martin Luther..
And he has something to say, really meaningful..
He said, ask question..
Because we're fully justified by the blood of Christ,.
by the cross, so we're 100\% same..
Because of the grace of God, I'm 100\% same..
Yet, I'm also 100\% sinner..
Even you come to saving faith in Jesus Christ,.
you are not from 100\% sinner become 75\% sinner, right?.
Doesn't take this account in our sinfulness..
We're still 100\% sinner..
And yet, we are fully justified in Christ..
The psalmist, how many times he described himself as servant?.
Servant..
Two times..
Two times in this short prayer,.
praying as a servant to the Lord, right?.
Center on the authority, right, of God..
Not himself, humble himself before God..
He realizes fallenness, he also realized he needed.
the instructions, the word of God..
Now, in a sense, whenever we think of the word of God,.
whenever we think of God, we should not stand,.
and we must not stand above the word..
We don't stand alongside with the word either..
We stand under the word..
When we get together to do Bible study,.
when you read the word, when you think and meditate.
on the word, you stand under, and always under,.
under scripture..
You come before the word of God as a servant,.
ready to obey, ready to listen, ready to even wrestle.
with the meaning of what the word means for us today,.
always being teachable, always being teachable..

$^{441}$Now, that does not mean we don't need to think..
That does not mean we don't need to use our brain..
We need to, right?.
Sometimes we have to read the word critically,.
not in the sense that we're against it..
Critically means we're thinking it through, you know?.
We're using our intellect, God's gift to us, right,.
to reason, to struggle, to wrestle with the word, right?.
However, the ultimate position that we take.
in front of the Bible is to stand under the word of God,.
to be teachable, ready to obey..
Let's all pray together..
I pray that we'll come to God as servants..
Dear God, you are the king..
You are the leader..
You are the Lord of all..
And we stand under you..
We sit under you..
And we position ourselves as your servants..
And we pray that as we meditate on your word,.
as we try to understand the contemporary relevance,.
how it applies to our lives today,.
we'll always stand under your authority..
And we praise you because you speak to us..
No mortal can speak about you.
because we are creatures..
We are humble..
We are actually nothing..
Our existence depends on you..
So we realize that we cannot speak about you..
It's just because of your grace, you speak to us..
You speak to us..
You introduce yourself to us..
And thank you for your revelation..
Thank you for your word..
From the time of ancient to our time modern,.
you are still speaking..
And as you speak,.
we pray that you will speak in our lives as king,.
as the ultimate creator..

$^{481}$And I pray that you will grant each and every one of us.
heart and spirit of obedience,.
the spirit of being your servant,.
ready to obey, ready to be taught..
We pray that you will grow our church in your word,.
that we'll worship you, we'll follow and obey your will..
And pray that in the rest of the week,.
you will send us into different situations,.
different positions, different ministries to serve you,.
to bring the impact of your gospel into our community..
In Christ, we pray..
\newpage



\section{啟示錄 6:5-11}
\label{sec:s7F8M3fuVJY}
\textbf{亂世啟示:第三至第五印 (啟示錄 6\_5-11) - 陳崇基牧師}
\newline
\newline
連結: \href{https://youtube.com/watch?v=s7F8M3fuVJY}{\texttt{https://youtube.com/watch?v=s7F8M3fuVJY}} ~~~~ 語音日期: 2022-06-01
\newline
\newline
\hyperref[sec:JkOQgNW0fh8]{\small{< < < PREV SERMON < < <}}
~
\hyperref[sec:index_chronic]{\small{[返順時目]}}
~
\hyperref[sec:index_scriptual]{\small{[返順卷目]}}
~
\hyperref[sec:vmXTNOZZzYc]{\small{> > > NEXT SERMON > > >}}
\newline
\newline
啟示錄 6:5-11
\newline
\begin{longtable}{cl}
\hline
\hline
章節 & 經文 (和合本修訂版)\\
\hline
6:5 & \begin{tabularx}{0.7\textwidth}{X} 羔羊揭開第三個印的時候，我聽見第三個活物說：「你來！」我就觀看，看見一匹黑馬；騎在馬上的，手裡拿著天平。 \end{tabularx} \\ \\ \relax
6:6 & \begin{tabularx}{0.7\textwidth}{X} 我聽見在四個活物中似乎有聲音說：「一個銀幣買一升麥子，一個銀幣買三升大麥；油和酒不可糟蹋。」 \end{tabularx} \\ \\ \relax
6:7 & \begin{tabularx}{0.7\textwidth}{X} 羔羊揭開第四個印的時候，我聽見第四個活物說：「你來！」 \end{tabularx} \\ \\ \relax
6:8 & \begin{tabularx}{0.7\textwidth}{X} 我就觀看，看見一匹灰色馬；騎在馬上的，名字叫作「死」，陰間也隨著他；有權柄賜給他們，可以用刀劍、饑荒、瘟疫、野獸，殺害地上四分之一的人。 \end{tabularx} \\ \\ \relax
6:9 & \begin{tabularx}{0.7\textwidth}{X} 羔羊揭開第五個印的時候，我看見在祭壇底下有曾為神的道，並為作見證而被殺的人的靈魂， \end{tabularx} \\ \\ \relax
6:10 & \begin{tabularx}{0.7\textwidth}{X} 大聲喊著說：「神聖真實的主宰啊，你不審判住在地上的人，為我們所流的血伸冤，要到幾時呢？」 \end{tabularx} \\ \\ \relax
6:11 & \begin{tabularx}{0.7\textwidth}{X} 於是有白袍賜給他們各人；又有話吩咐他們還要歇息片刻，等到與他們同作僕人的，和他們的弟兄，像他們一樣被殺的人的數目湊足的時候。 \end{tabularx} \\ \\
[1ex]
\hline
\hline
\end{longtable}
$^{1}$我們回到啟示錄的系列.
我們曾經說過很多跟啟示錄這本書沒有關係.
但是幫助我們去明白到底這個終末的事.
所以我們繞了一個圈.
我們回到啟示錄.
我們今天終於去到.
之前說了兩印.
我們這次是說第三,第四,第五的印記.
我們說過啟示錄是歷史的主宰.
耶穌基督去打開這些印記.
去告訴我們世界的終末到底會發生什麼事.
以及我們作為他的子民.
我們如何做好預備.
所以我們打開啟示錄第六章.
我們今天先看第三印.
之前我們看過第一印是白馬.
我們說到俄馬帝國看到白馬的時候.
有外環.
然後有紅馬就是內戰.
然後我們第三印就是我們今天看的黑馬.
顏色可能不是很對.
那裡有一點錯誤.
黑馬就是六章第五節.
我們先看.
揭開第三印的時候.
我聽見第三印第三個的活物這樣說.
騎在馬上的首類拿著天秤.
我看見,聽見四活物似乎有聲音說.
一錢銀子買一星墨紙.
一錢銀子買三星大墨.
又和酒不可以糟談.
說了什麼呢?說了三句話.
這個印記第三印.
說到一錢的銀子買一星的小墨.
一錢銀子是當時大概一個勞工一天的薪金.
通常一家需要的食糧每一天.
就是所謂一星的小墨.
一星就好像我們美國大概一億.
如果你懂得汽車維修的話.
那些一億的Motor Oil.

$^{41}$大概就是中文譯作一星.
所以這裡說到一個人大約每天要吃一星的小墨.
第二句話就說.
一錢銀子就要買三星的大墨.
是什麼意思呢?.
在第一世紀的農耕的社會.
農業為主的社會.
如果小墨是短缺的話.
就會用大墨.
就是Wheat不夠的話就會用Barley.
每一天大概一個家庭.
可能以不同的計算.
就是需要三星的所謂三個Quart的大墨.
然後才能夠生存.
這裡似乎告訴我們一些物資短缺.
和一些通脹的一些說話.
一些當時羅馬帝國第一世紀.
啟示錄寫成的時候.
那段時間那些人會很明白.
當你說到一錢銀子買多少東西.
然後小墨不夠.
然後你換大墨.
就知道糧倉不夠.
或者物價開始高漲.
等於我們今天說.
在美國的人.
一個Gallon的車油.
就是六塊錢.
是不是嚇怕你.
現在五塊錢已經牙痛.
類似一些這樣的說話.
告訴我們.
帝國到了這個時間.
上帝如果要去到私人的審判.
就會有一些這樣的情況發生.
之前兩個的印.
第一第二印.
就是有講到戰亂不安的事情.
然後第三印引致的情況.
就是當時的帝國社會裡面.

$^{81}$有一些不安的事情.
我們講過了.
啟示錄的時候.
不要這麼快當它是預先寫好的新聞報紙.
不是和我們今時今日來講說話的.
是當時的人很熟悉的一種環境.
如果翻查一下.
當時第一世紀一些的記載的話.
你會知道在不同的羅馬帝王.
不同的時間.
也都有羅馬帝國所謂天下.
有的糧倉短缺.
這個物價上漲的問題.
在不同的皇帝.
譬如尼祿的時候.
大概是主後60年代的時候.
也都有通貨膨脹等等的問題.
啟示錄寫的時候大概是第一世紀末90年代.
當時竇米田.
也有一種這樣的情況出現.
就是有通脹和有饑荒的出現.
甚至因為當時的社會.
很多時候有農田的人.
他們因為為了賺錢的緣故.
就會將他的田地不種大麥不種小麥.
其實大麥小麥才是最重要的.
但是因為賺錢的緣故.
將他的田由麥田變成葡萄園.
或者變成橄欖樹的園.
為什麼呢.
因為那些是比較貴重的一些物資.
還有賺錢會比較多.
所以有很多的農田都會轉變了.
然後變成葡萄園或者變成橄欖樹園.
特別是在羅馬東的地方.
Roman East.
什麼叫羅馬東呢.
就是今天的土耳其.
就是以弗所那個地方.
就算是羅馬東的地方.

$^{121}$其實就是啟示錄第一代的讀者.
七教會.
以弗所至奴隸家七個教會.
那裡有很多的農田.
都因為物價通脹等等.
然後繼續轉變了大麥小麥的最基本的食糧.
變成葡萄園等等.
所以到了公元後90多年的時候.
竇米田.
就是啟示錄寫的時候的帝王.
曾經下令過.
將那些葡萄園或者橄欖樹園.
就將他轉變成大麥小麥.
就造成社會裡面有很多的抱怨.
很多的抗議等等.
然後帝王要收回成命.
在一些經濟上很亂的情況.
可以看到這個印打開的時候.
就是想揭示什麼呢.
就是揭示羅馬帝國本身是當時世界天下.
最大的經濟體.
最大的一個經濟的系統.
好像是無堅不摧.
一些當時羅馬帝國的記載.
告訴我們羅馬已經是控制住很多不同的麥田.
還有很多食糧的出產的地方.
就是他們的供應鏈已經全部管理好.
還有控制好.
到了第一世紀末的時候.
就有一些的亂局出現.
啟示錄其實就是想一直一章又一章.
一個人又一個人.
就是告訴我們這個帝國是不是真的無堅不摧.
你記得第一代的讀者.
去看啟示錄的時候.
就是第一代的教會.
那個所謂七教會.
是一些很小很小的教會.
一些在家裡躲起來聚會的小組的教會.
在異佛所不同的家庭的角落聚會.

$^{161}$在老底家不同家庭的角落聚會.
封信就是寫給他們的.
一個小的教會.
一些很小的勢力.
面對著羅馬帝國.
很多時候軟硬兼施的.
去到克制他們.
或者去打壓他們.
對著這個哥斯拉的時候.
這個教會需要有什麼信息.
啟示錄就是講到中門.
就是告訴教會.
雖然他們面對著很多很黑暗的勢力.
來威脅著他們的存在.
但是這個啟示錄就是去到揭示那個中門.
就是最終你要看通這個帝國.
其實本身他們的能力不是那麼大.
其實帝國本身.
羅馬帝國.
都是不堪一擊.
到了上帝的審判.
上帝的工作來到的時候.
上帝的公義審判來到的時候.
每一個人打開就告訴我們.
這個帝國其實是很脆弱.
到底你靠的是什麼.
到底你怕的是什麼.
其實今天都是一樣.
你看我們身邊發生的事情.
過去幾年.
是不是.
一個疫情.
令到我們整個國際的秩序.
世界的秩序.
全部東西.
打斷了.
全部東西都中斷了.
那個supply chain等等.
又加上一場好像.
離我們很遠的一場入侵的戰爭.

$^{201}$就令到我們很多地方都混亂了.
當然我們需要做好.
我們的管家的責任.
錢財和一些物資.
我們需要用我們的智慧去做好.
但是我們需要有一個對的perspective.
就是我們投放的資源.
我們的精神到底是給什麼.
有時候我聽到朋友說.
很想提早退休.
是不是.
可能你也這麼想.
沒有問題提早退休.
但是我想多問一個問題.
或者啟示錄會多問我們一個問題.
提早退休為了什麼.
如果為了你可以好好享受的話.
我想這個不是啟示錄給我們的答案.
如果你是為了事奉神的話.
如果你是為了你將你的精力.
可以投放在一個真正不會滅絕的.
真正無堅不摧的天國.
Kingdom of Heaven的時候.
這個就值得做.
為誰辛苦為誰亡.
為了什麼.
給我們一個大一點的perspective.
一個綜末論.
就是給我們一個大的角度.
來審視一下到底我們現在做的決定.
我們想做的事情是什麼.
上一個學期.
我有機會去到學院那裡.
上一個研討會一個seminar.
我記得休息的時候.
我就和同學去到大家聊聊天.
其中有一個同學就去到投訴.
他生活裡面有一些不順利的事情.
原來是一些所謂電腦的事情.
另外旁邊有一個同學他坐輪椅.

$^{241}$他就大聲說.
First world problem.
什麼意思.
First world problem.
其實他就是開玩笑說.
我們在聊那些所謂生活裡面.
令我們很煩的事情.
令我們很擔憂的事情.
其實都是那些很安恕的人.
所謂first world裡面的人.
我們才會掙扎的問題.
有很多世界各地不同的地方.
他們所掙扎的不是這些.
是不是.
有很多世界各地的居民.
他們連家園都沒有.
水都沒有.
每一天連想倒杯水都沒有.
他們所面對的問題和我們面對的問題.
我們就需要有一個對的角度.
對不對.
我家最近有一些裝修完之後.
有好幾個小時沒有水.
因為小朋友大學他們休息回來.
因為裝修開始一早就沒有水.
我們一早就起床.
他們那些很難睡.
九點十點才起床.
已經早了九點十點.
一定會暈的.
他們起床.
我女兒起床就說.
為什麼沒有水.
我就跟她說.
這個是first world problem.
我們想一下.
在中東有很多的人.
連起床刷牙.
倒杯水刷牙都沒有.
first world problem.

$^{281}$我們有沒有一個對的角度.
來看我們現在面對的問題.
然後啟示錄就是逼我們.
問我們所有所樂的.
到底是為了什麼.
有多少是為了天國.
有多少是為了地上脆弱的帝國.
第三人.
我想更加重要的是第三人這裡說的.
除了那幾句話之外.
就是有天秤出現.
這個人的騎士是拿著天秤的.
天秤是代表著羅馬帝國的經濟體.
說到天秤的時候.
特別是以前古代的時候.
買東西是要秤的.
我們上一代也是這樣.
買財譜是要秤東西的.
所以有天秤出現的時候.
你就自然想起經濟的問題.
是第一世紀羅馬帝國經濟的問題.
但我想對一個熟知舊約聖經的人.
他聽見天秤的時候.
可能會有不同的想法.
會有不同的信息出現.
因為舊約聖經裡面.
天秤和帝國一起出現的一個故事.
就是在但義理書第五章.
如果熟知舊約的話.
他就會想起第五章上下這段故事.
這裡是什麼記載呢.
這個故事.
但義理書第五章.
就是講到白沙薩王.
他有豪門夜宴.
由第一節開始一直講.
他有多少的妃嬪.
有多少的官員.
有多少的人一齊來到豪門夜宴.
一個很奢侈的.

$^{321}$一個很豪華的生活.
甚至在耶路撒冷所劫回來的器皿.
在聖殿裡面用的器皿.
他都拿出來.
然後在那裡斟酒.
在那裡喝.
那些金銀器皿都拿出來.
他好像是一個隻手遮天的一個時間.
這個白沙薩王好像是沒有任何東西.
能夠去搖動他巴比倫帝國的權柄.
誰知道在豪門夜宴裡面.
你記得.
有一隻手在牆上寫字.
寫了一些字.
大家看不明白.
那個王就很害怕.
所以就傳召了但以理.
來到去豪門夜宴裡面.
就叫但以理來解通.
到底這句話是什麼意思.
二十五字.
所寫的文字.
就是尼尼微尼提克勒烏法爾辛.
其實本身是亞蘭文.
我們只不過是中文用了音譯出來.
尼尼微尼提克勒烏法爾辛.
解釋是什麼呢.
這文字的意思是這樣.
尼尼.
其實就是counting的意思.
今天你讀accounting的時候.
count 1 2 3.
尼尼就是亞蘭文尼尼counting的意思.
尼尼就是上帝已經數算了你國度的年日.
始國終結.
提克勒其實就是天平的意思.
就是waiting的意思.
就是稱東西的意思.
顯出你國度和你的缺乏.
然後比勒斯.

$^{361}$比勒斯其實就是法爾辛的單數.
這樣的解釋就是說.
你的國度就會被分裂.
being divided.
歸給馬代人和波斯人.
但以你這個預言發生了之後.
結果聖經怎麼交代.
第30節.
當夜迦勒底王白沙薩就被殺.
然後馬代人大劉氏.
62歲的時候就娶了迦勒底國.
就是娶了巴比倫帝國.
世界的歷史確實是這樣發生的.
因為這個巴沙薩王.
他就是末代的皇帝.
誰知波斯馬代人就去到巴比倫京城.
就不需要一刀一槍就拿了京城.
一個外表好像很堅強很穩定的一個帝國.
本身就是很脆弱.
當上帝的審判臨到的時候.
一切都被摧毀.
說回我們美國其實都是一樣.
在我們美國.
你看看我們本身好像很好.
特別是這幾年.
完全是揭示了這個國度.
是多麼的脆弱.
你說這些槍械管理.
我覺得這些是很正常的.
是常識的東西.
不是不讓你有槍.
可以 我的看法就是.
No problem.
我們有這樣的權利去武裝.
但是不是需要有對手.
I just don't get it.
背後是有多少政治角力.
多少黑暗的東西.
鍾活倫就要揭示一件事.
就是我們沒有一個宏觀來看我們生命.

$^{401}$如果我們沒有一個宏觀的時候.
我們很多時候為誰辛苦為誰忙.
我們就是為了一個會終結的國度.
一個很脆弱的國度.
我們辛辛苦苦拼命.
一輩子就是這樣浪費.
鍾活倫就是告訴我們.
天國才是我們應該要關心的事.
這個宏觀.
給了我們一個真正的目的.
第四印.
高揚揭開第四印的時候.
我聽見四個活墨說.
騎在馬上的名字叫做死.
陰間也隨著它.
有權兵切割它.
可以用刀劍機放瘟疫野獸.
殺害地上四分之一的人.
不要當它是預先寫好的新聞報導.
是講第一世紀.
或者用一個天啟文學的方法.
一些意象的方法.
去告訴我們這個帝國.
上帝的審判要臨到的時候.
就會發生什麼事.
有刀劍第一第二人.
有機方第三人.
然後隨之而來的.
肯定就是瘟疫.
或者殺害地上很多很多的人.
這匹馬是什麼顏色.
灰色馬.
可能有不同的翻譯.
可以翻譯得更加好一點.
其實本身的字可能是青色的意思.
即是青灰色的意思.
面青嘴青這樣.
羅馬帝國會帶來的是什麼.
就是死.
第五人.

$^{441}$高粱鐵開第五人的時候.
並為作見證而被殺的人的靈魂.
大聲喊著說.
神聖真實的主宰.
你不審判住在地上的人.
為我們所流的血 腎源要到何時呢.
於是有白袍刺給他們各人.
又有話吩咐他們還要歇息片刻.
等到與他們同作僕人的和他們的弟兄.
像他們一樣被殺的人的數目湊足的時候.
這個第五人就很詫異.
因為第一二三四人.
都是告訴我們帝國的終末的時候會發生什麼事.
上帝的審判臨到的時候.
這個帝國到底能不能站穩.
是站不穩的.
居然第五人的時候是跟什麼有關呢.
就是跟那些為見證 為福音見證而被殺的這些靈魂.
在上帝的祭壇底下他們來禱告.
說到這些殉道者.
我們曾經說過第一世紀.
軟硬兼施底下有很多殉道者.
在當時的教會已經差不多站不穩了.
要散了.
這些殉道者因為為了見證上帝.
而自己的生命被威脅.
這裡就揭示了一件事.
第五人就是上帝怎麼看他們.
這些好像在人的眼中.
當然也是兩種不同的角度.
一個地上的角度.
一個天上的角度.
地上的角度覺得這些是很失敗的.
你看看這個教會快要散了.
在帝國的淫亂底下.
他還能夠站多久呢.
在地上看就是看到有殉道者.
看到他們流去 看到他們死亡.
看到他們被殺 被監禁.
但是天上給我們一個另外的角度.

$^{481}$是說什麼呢.
百抱刺急覺.
一個異抱.
是受到上帝的接納的.
並且你留意他們呼喊的地方.
就是神的祭壇之下.
等於祭物在祭壇之上被宰割.
獻給上帝.
但是他們流的血在祭壇底下.
他們所呼求的.
是得到上帝的接納.
各人有百抱刺給他們.
神也是聽見這些殉道者的聲音.
去到回應.
你有沒有留意這裡是啟示了唯一一個禱告.
唯一記載的 引述的一個禱告.
一個prayer.
就是在這些祭壇底下的殉道者的聲音.
他們的祈禱.
祈禱什麼呢.
無辜者申冤的祈禱.
唯一一個的禱告.
就是願意上帝快點出手.
他的公義能夠彰顯.
公義不是純粹去到take care的帝國.
公義的是為了上帝的道義.
殉道的這些靈魂.
他們需要公義 需要有justice.
他們的祈禱是什麼呢.
他們等到什麼時候.
如果你熟悉教育的話.
你就知道詩篇裡面經常有這樣的禱告.
就是義人受苦的時候.
他們的禱告就是how long.
到底要等到什麼時候.
向上帝呼求 向上帝公義的彰顯.
一種的禱告.
上帝說 他聽了.
他聽了.
但是還要等.

$^{521}$要等到什麼時候.
還沒有湊夠數目.
什麼數目.
殉道者的數目.
還要等.
原來還有一些要為主殉道的人.
他們還沒有走上這條路.
為什麼會說這些話.
對於一個猶太.
或者有猶太背景的一些讀者.
在史詩第一代的讀者.
這些話不是陌生的說話.
因為在舊約和新約之間.
有大概幾百年時間.
這幾百年時間我們覺得是一些.
所謂的silent period.
我們以往曾經.
我們就說 以前我們就說兩約中間.
四百年沉默時期.
其實一點都不沉默.
這幾百年其實猶太人.
他們掙扎著很多神學的問題.
特別他們在不同的帝國.
改朝換代之下.
他們是受盡迫害的.
他們有很多的著作 小說.
或者一些文獻他們寫出來.
很多時候寫到他們這些義人受苦.
猶太人因為去抗衡.
去抗爭這個帝國的時候.
他們就有很多被殺.
很多是殉道的.
但是猶太人的殉道者.
這些作品裡面經常都說一件事.
就是仰望上帝公義的彰顯.
上帝你什麼時候會出手.
但是這些作品很多時候去反思的時候.
他們就會說 因為數目還沒有足夠.
Yes 我們還是受威逼利誘之下.
我們是受盡折磨的.

$^{561}$數目還沒有足夠.
到足夠的時候.
上帝就必定會去到申冤.
必定會有公義的彰顯.
其實這句話 這種說法.
就是說上帝聽見殉道者的申冤.
但是上帝會有打算.
上帝不會忘記那些忠心為他們事奉的人.
他們所拜上的.
現在已經是五月了 2022年五月.
一月開始的時候.
幾乎每一個星期.
我都會上網.
用我自己的時間.
和一些我們叫他做創啟地區.
Creative Access的地方.
一些教務 一些傳道人.
我們在網上有一個聚會.
我們和他分享一些.
怎樣看聖經等等.
這些課程.
我真的很佩服他們.
在我的人生裡面.
我從來沒有經歷過這樣的情況.
他們要上網去一個講座.
但是他們不想 不可以開鏡頭.
我們不開鏡頭.
就是因為我們在吃飯.
我們不開鏡頭就是因為不方便.
或者我們還沒有化妝等等.
他們不開鏡頭是有政治原因.
是不開鏡頭.
因為他們怕不同的打壓.
有很多創啟地區的人.
特別是監獄.
如果你去做一些這樣的事.
可能你做完.
你還有公安來找你的.
是不是.
大家很清楚 是不是.

$^{601}$但是我真的很尊敬這班人.
因為他們受的裝備或者受的教育.
對我來說.
我當然比他們更幸福 是不是.
我讀書比他們多.
我有的資源比他們多 是不是.
我可以得到的資訊比他們更加多.
但是他們來到.
我們幾乎每一個星期.
除了有一兩個星期我們休息.
但是他們很想來到.
在道上能夠得到裝備.
使他們去帶其他小組 教會.
他們能夠更加有能力.
這些人有這麼堅定的信心.
可以繼續向前走.
這簡直是我們不能想像.
我們這些所謂first world裡面的人.
我們幸福到一個極點 是不是.
很多時候我們這班信徒.
我們有少少風雨日.
我們就什麼都不記得了.
我們都不知道搞什麼鬼.
如果你是這樣的話.
啟示錄就是要挑戰你.
去問清楚一些要問的問題.
和一些你經常避免不想問的問題.
一些終極的問題.
Asketon 終末的問題.
到底你是向誰交仗.
你忙的是為了什麼.
你辛苦的是為了什麼.
我不是說我們生活裡面.
我們就不需要管家的問題.
我們需要好好管理.
我曾經說過 剛剛說完.
我們需要管理 我們有stewardship.
但ultimately for what.
ultimately 到底是為了什麼.
是為了一些一定會消逝的東西.

$^{641}$本身好像無堅不摧.
但他問清楚的時候.
本身就是很脆弱的東西.
無論是一些地上的替果.
是一些地上的人有的事.
到底是為了什麼.
啟示錄就是問我們這件事.
所以為什麼第五人就去說這件事.
不要只看著這些東西可能會倒塌.
這些只不過是house of cards.
第五人就告訴我們.
到底是為了什麼.
有些什麼是上帝會閱立.
上帝是會使用的.
居然就是在壇底下的.
曾為上帝的道作見證.
被殺的人的靈魂.
這些就是上帝會使用的.
上帝說心裡有數.
你不需要怕.
因為上帝一定會為每一個.
這些被殺的靈魂而去到新淵.
並且結局會快樂.
問題就是這樣.
到底你是誰.
你是和帝國一起滅亡的人.
還是你是為了上帝的道.
為了上帝的福音.
你寧願被上帝接納.
作為壇底下被殺的靈魂.
為了終極這個國度.
一個永恆不會變改的國度.
而去努力的.
到底你是誰.
精神,你的時間.
每一晚睡不著的時候.
到底你是憂什麼,受什麼.
有什麼的觀點可以正確地對我們.
就是忠愛.
天上的國度才是永恆.

$^{681}$這個才是上帝所預立會使用的.
求主幫助我們.
讓我們不要忘記.
給我們一個當頭棒.
特別在我們中間.
芝加哥以外網上的弟兄姊妹.
當然我不能夠知道你在哪裡.
不是很確定你的情況.
但是我在說我們.
第一世界.
我們這班人.
到底是為了什麼.
為了什麼而受.
為了什麼而去努力.
我們先低頭祈禱.
主讓我們來到你面前.
我們去悔改.
我們去改變,去更新.
知道你自己的終末的角度.
你的角度.
怎樣去看世事,怎樣去看我們的人生.
神啊,我們求你.
你的國度向我們彰顯.
看到我們天上坐著為王的.
是上帝.
天上的國度.
將來永恆的國度.
才是永恆不變.
主叫我們的精力,我們的精神.
投放的.
是在你的福音.
你的道上.
我們這樣祈禱,奉主命救.
我們播我們的回應詩歌.
藉著回應詩歌.
有一節的時間.
我們在上帝面前.
你可以去默禱.
願意你的祈禱是上帝.
讓我看見你終末.

$^{721}$天上的國度.
有佢使用我地.
\newpage



\section{}
\label{sec:vmXTNOZZzYc}
\textbf{Apostles' Creed \#4 (Jesus Christ) - Rev. Mark Chan}
\newline
\newline
連結: \href{https://youtube.com/watch?v=vmXTNOZZzYc}{\texttt{https://youtube.com/watch?v=vmXTNOZZzYc}} ~~~~ 語音日期: 2022-06-06
\newline
\newline
\hyperref[sec:s7F8M3fuVJY]{\small{< < < PREV SERMON < < <}}
~
\hyperref[sec:index_chronic]{\small{[返順時目]}}
~
\hyperref[sec:index_scriptual]{\small{[返順卷目]}}
~
\hyperref[sec:Bxy4FSvUGwA]{\small{> > > NEXT SERMON > > >}}
\newline
\newline
$^{1}$All right guys, let's turn to our Bible and we're going to work on the Apostles Creek and we've been.
Examining the basic of our faith the core beliefs of our faith and when people ask you what is.
Christianity all about and you should be able to give a ready answer, right?.
this is what the Creek of Apostles is all about right an ancient Creek almost 2,000 years old and.
Going back to almost the Apostolic Age.
So we are working through.
Possibly line by line.
Something like that and we started with the first statement and we're going to move on to the next section today.
But let's reveal the first statement. I believe God the Father Almighty.
Creator of heaven and earth. All right.
So I want you to really think about it because the Apostles Creek even though there's there are only 12 statements.
It's not a long document is pretty much a pocket size, right? You can fill into one screen on your phone.
But these 12 statements are the core beliefs the foundations of Christianity.
I believe in God the Father creator of heaven and earth and that statement the first one.
alone sets apart.
Christianity from all other religions.
Think about it. It sets apart.
Christianity from all other religions.
We live in the time of.
What we call pluralism, right? Everything should be considered equally, right?.
That's what pluralism means.
Or we value diversity.
We live in a time when we always hear people say.
All religions are the same right all religions are helping people to do good.
Right to help you improve.
And on all faiths the same, right?.
That is the cultural script. That is what the culture is telling us all.
Beliefs all religions are taught are teaching us how to be a better person.
That is the assumption right but this statement actually sets apart.
Christianity from all others think about Christianity and.
Buddhism.
Or Hinduism or Islam.
Are they all the same this statement says no because these.
Apostles creeds.
Begin with God being the creator of the universe and the world has a beginning.
the world has a.
beginning of a personal agent.
Someone has started the world God the creator it sets the Christian faith apart from many other.
Eastern religions.
for example Buddhism.

$^{41}$Is it also sets apart the biblical faith?.
from all other secular religions or.
ideologies right religions or ideologies.
atheism or.
materialism.
that belief.
essentially.
Matter is all there is right from eternal to eternal.
all.
matters.
But this statement says no.
The universe that we live in we are part of the universe.
exists all because of a.
personal God who created it and.
It has a reason to be.
This universe and you and me we exist. There was a reason right?.
But as I work through.
Some of the passages in the Bible that talks about the beginning of the universe you understand. It's not a scientific report, right?.
So get this. Okay. This is your take home.
Go out from this church.
The rest of the week don't even think that the Bible is this scientific report.
No, it is not is this an ancient document?.
It is a document from many many years ago.
however, God used those concepts to.
To to teach us.
That who is in charge and who started all right. So Genesis 1 and 2 are not.
Scientific reports. All right.
They are theological report the theological statements telling us who.
Is behind all of this?.
I was reading a very very old book our print book and.
Written by J. I Pecker. I think he said it's so right. So I want to quote it.
He in this book I want to be a Christian.
Written long time ago our print but he said it's so right.
He said Genesis wanted to tell us who?.
with so many answers about how.
Some today think this is a defect.
but in a long huge perspective of history our present-day scientific.
preoccupation with how rather than who looks very odd in itself.
rather than criticizing these chapters.
For not feeding our secular modern contemporary interest.

$^{81}$Scientific interest we should take from them Genesis 1  and  2.
Needed we built of our perverse passion.
For knowing the nature to creation without regard to what matter most mainly nature's creator.
We are obsessed with the how right?.
Technology science right they have their place, but we are sometimes.
totally obsessed by this.
We need to be rebuilt because we are interested in the how not the who the Bible tells us we should be concerned about the who.
Who do we worship and who do we follow? It's God himself the Creator.
He continues the message of these two chapters is this you have seen the sea the sky.
the Sun moon the stars and.
Wonderful complexity of the creation be it human beings and all that fantastic, isn't it? Well now meet the one.
Who is behind it all?.
Genesis 1  and  2 are telling us. Hey, you have enjoyed these works of art.
And now you must shake hands with the artist and you are thrilled by the music.
We now introduce you to the composer.
It is to show us the composer the artist the creator.
Rather than the creation. That's what Genesis is all about.
That's what Genesis 1  and  2 is all about and the Apostles Creek rightly.
Says that we believe in this Creator God himself.
Creator of heaven earth creator of our existence our world and then it continues to a huge section about.
Jesus Christ the Son of God.
Now you can tell from the highlight. This is the longest section of.
The whole Apostles Creek the longest section on Jesus Christ the central section the core of everything.
because there are short statements before the highlight and then there are short statements after it but.
these short statements.
sandwiched.
the middle and.
It is all about.
Christ himself Jesus Christ himself.
now Jesus.
It's a name. All right, that's the name.
in Greek.
or.
Greek is for Joshua and Joshua.
Was was originally a Hebrew name. So Joshua Jesus.
Mean the same thing.
It's a proper name, right just like you can call me mark, right?.
Now it tells us from the Bible that Jesus is.
the only Son of God and.
was Mary's son.

$^{121}$from virgin birth.
Now that means it is only possible.
with God's miracle.
the birth of Jesus a miracle child.
The accomplishment of the plan of God the redemption of God is only possible because God did.
Also tells us the humaneness of Jesus Christ his humanity.
He is from Nazareth.
in Galilee.
He was a carpenter a Jewish carpenter.
He worked for three years as a rabbi and was put to death as the statement says was crucified and died.
in a specific time.
in the time of Pontius Pilate.
right.
so all the statement tells us that.
Jesus was a.
historical person, right a.
You a historical person.
You are right someday you die. Hopefully you live some footprint in the book or so.
All right on a tombstone.
You are historical.
Jesus was historical but he's more than a historical figure.
Because what Jesus has said or what Jesus has done.
It's beyond what a human could do or could say.
So Jesus is the centerpiece of the Apostle's Creed.
And then I said Jesus is the name but Christ is a title.
All right, Christ is not a proper name. Just like pastor is not a name, right?.
You can call me mark, right? This is my proper name.
But my proper name is not pastor, right? The title is the pastor Christ is a title, right?.
Um, don't mistake it saying Jesus Christ because his last name is Christ.
Christ is not Chen. It's not Lee. It's not Smith or any last name, right? It's an office.
title.
It's from the Hebrew Bible from the Old Testament.
Messiah.
So Messiah is the Hebrew and then translated to Greek is Christ.
meaning anointed one.
In the Old Testament culture in the Hebrew culture whenever a king is being set apart.
To be to reign on earth to reign on behalf of God.
He is being set apart by pouring oil on his head on his hair.
Over his head that is anointment right to.
To set him apart as king.

$^{161}$Now so that's why Jesus Christ.
Means Jesus the Messiah, right Jesus the Christ the Messiah the king. It's a title.
only much later in the Christian.
Generations and generations after after Christ.
Christ becomes a proper name, right?.
from a proper.
Title become becoming a proper name. So that's why.
Even in the New Testament Jesus and Christ are synonymous.
You can call Christ Jesus or Jesus Christ, but you do know that the.
Word Christ means he is the king. He is being anointed by God. He died but rose again.
So he rightly reigns over all of us and even the universe when he comes back.
So whether we believe in Jesus Christ as the Bible says or not now, we still need to understand who Jesus really is. I.
Once had a conversation with a Buddhist monk in Chinatown.
Interesting, right?.
So we met in a community meeting.
Something like that and then after the meeting we kind of you know, shake hand and instills one another and then.
You know because they a Buddhist monk has a different dress, right?.
They got the rope right? I don't right and then I approach him and I.
Introduce myself. I'm I am.
My name is Mark, right and I'm Pastor Mark.
So he introduced himself and then we try to talk and then try to strike conversation and after the monk.
Learned that I'm a Christian church pastor and our conversation quickly.
Turn into as you can tell.
About our religions, right? And then this monk continues. He said oh.
Great your church pastor and we are the same.
And I was kind of astonished.
Yeah, really, you know, I didn't shake my head. Um.
He said we are the same.
Both your religion and my religion are essentially the same because we are teaching people how to do good, right?.
Oh, that's nice. You know, you're doing your part and I'm doing my part and I quickly answer, you know, respectively, of course.
Right, respectively. I said, um, you know what? Mr..
monk, uh.
I.
Forgot his name. I forgot. He has his title. Yes. It's a very formal title. I said, hmm.
Not really.
And then he was kind of stunned and you know, I disagree with him, you know.
And I started to talk about you know, who Jesus is I.
Started to talk about who Jesus Christ is.
Right. He died on the cross to save us.
And he rose again. I.

$^{201}$Talked to him. We are being saved not because we do good.
We have been saved because someone else.
Did something on our behalf so that we are justified.
Right, so we become righteous at least positionally we become righteous.
so I so I can't kind of counter and then.
finally, you know after.
my, you know impromptu preaching.
to him.
He said okay, maybe the two religions are with a different I.
Said of course, we're different right what we expect you you try to teach people to do good.
But you know is apple and orange right you can't say okay we are.
You know, you don't say you don't say dogs and cats are the same because both of them have four legs, right?.
Simply because of that right the comments the lowest common denominator, right? I think that's a logical fallacy, right?.
you got to.
Compare and contrast right in order to know the essence of the two, right?.
As I said in the beginning our culture is telling us everything is the same all are the same.
But you know, I bet to be the differ.
And the Bible does that too?.
Romans 5.
Romans 5 Paul is writing to the church. He said you see at just the right time.
When we were still powerless.
Christ died for the ungodly.
Very rarely will anyone die for a righteous man. Though for good men someone might possibly dare to die.
But God demonstrates his own love for us in this look at this.
While we were still sinners Christ died for us.
This is there's not a hint of.
Improvement. There's not a hint of.
Being better so I can come to Christ. No, there's not a clue here.
We begin with the premise that we are sinners.
It says we are sinners. We are powerless. We are condemned and.
The pastor says we are ungodly.
Let me focus on this.
Emphasize this I emphasize this to many many people even those who have come to church for many many years.
Somehow so we don't get it. You know, let me emphasize this the Bible.
the Bible.
The Bible is not about teaching us to do good.
Scripture.
It's not about teaching you and me to do good, all right, so people are pretty stunned when I said this, right.
I remember that one time I was in the church. I was talking to.
Quite a few people from from China.

$^{241}$Actually, I was in Tel Aviv. I was in Israel. They have a little Chinese church there.
So I was invited to to speak so after speaking I was just discussing with some some of the people.
There and some of them are Christian some of them are not and I said that in order to stun them, you know in order to.
Shock, you know, I said, you know, we're looking at the Bible right? The Bible is not about teaching us to do good and.
They were like what?.
Because they have some sort of you know, notion of sacredness, right? This scripture, of course is about it's something good, right?.
Something to teach us good. No, I said no, I never found a verse in the Bible that teach us to do good.
They said Willie. I said yes, I'm serious.
Look at me. I'm pastor. Okay. I'm telling you the Bible is not about teaching us to do good to be good.
Let me say this again. The Bible the biblical faith is not about self-improvement.
So that we can be a better moral persons.
No.
Never find any word or verse in the Bible. I.
Really don't I study Bible for.
quite some years.
I really don't find any self-help.
moral lessons in the Bible I.
don't.
The Bible is not telling us to do good. No instead.
the Bible.
Is showing us we cannot do good.
We are not good and never can we be good enough that's what the Bible says.
The Bible says you cannot be good and you never.
Never can you be good enough?.
We are sinners.
We are powerless.
We're powerless to do good.
We're condemned in our sin in death.
So, what can we do that's what the Bible says Jesus came.
Jesus came to save us. That is the good news.
That is the gospel. That is what the Bible is trying to do.
So when we come to the Apostles Creek trying to understand who Jesus is.
Consider.
Carefully.
Who Jesus is?.
If the Bible is the gospel is all about improving ourselves, you know being better so that we can come to church.
We can come before God.
Then Jesus didn't have to die, right?.
Jesus have to do whatever he has done or say whatever he has said, right?.
As we call in the Bible Jesus has said something and did something and many things that is beyond.

$^{281}$typical human.
Can do or can say I?.
Still think the way CS Lewis word this so well in his classic.
mere Christianity.
Now he was talking about who Jesus is and let me quote this kind of a long quote.
I was CS Lewis.
Once passed away many many years ago. He used to be a Oxford.
English.
literature guy.
Scholar he wrote this book mere Christianity.
He's talking about who Jesus is.
He said I'm trying here to prevent anyone saying the worldly foolish thing.
That people often say about him being Jesus Christ.
I'm ready to the people are saying I'm ready to accept Jesus as a great moral teacher.
But I don't accept his claim to be God.
That is one thing we must not say said Lewis.
He said a man.
who was merely a man and.
Said the sort of things Jesus said would not be a great moral teacher. Think about it.
He said right if.
Jesus said what he said.
He would not be a great moral teacher. No way.
He is.
either a.
Lunatic or.
He would be the devil help.
Think about it. Whatever Jesus said in the full gospel or in the New Testament recorded.
He.
Would not be a great moral teacher no way and.
You must make your choice Lewis continues.
Either this man Jesus was and is the Son of God or else a madman or something worse.
You can shut him up.
for a fool you can spill at him and.
Kill him crucify him as a demon.
Or you can fall as feet and call him Lord and God but let us not come with a nonsense.
About him being a great moral teacher.
He has not left that choice open to us and he did not intend to.
Now that is the well-known.
Trilemma of says Lewis.
Trilemma means there's only three choices that is before us only three choices either Jesus is a lunatic.

$^{321}$crazy man or.
Jesus is a liar.
someone who has deceived many many people even the sensible people or.
He is indeed what he came to be.
He's a lunatic or liar or the Lord of all.
He's got himself so think about it. I think this is pretty well word.
What the Bible is all about what gospel is all about? I?.
Often hear people say, you know when I invite them to come to church, right?.
They often respond. Oh, I'm not good enough to go to church, right? I'm not as holy as you are.
Right, you guys are holy you can come to the church. All right, but you know, I kind of you know.
I'm not gonna come right because this is too sacred.
No.
And they always respond oh, let me do better so I can come to church.
We think.
That's not what the Bible says.
And that is the very reason that we are not holy. We sinners that we come to church.
We come to the cross and Jesus died for us. I.
Think one of the best passages that talks about the person and work of Jesus Christ is.
Philippians 2.
Let me read that and expand on that.
Your attitude should be the same as the attitude of Christ Jesus.
Who.
Being in fairy nature of God did not consider equality with God something to be grabs.
But made himself nothing.
But this that this is trickier.
the I.
Think in English as well in Greek.
There was some something called reflexive pronoun, you know reflexive.
meaning of a word.
Reflexive meaning meaning there's an action that you do not to others.
But you do it to yourself.
That is called reflexiveness, right?.
Self-reflexiveness.
He made himself nothing.
He made himself nothing he chose to do that to himself.
No one forced him.
But he did it all because of his willingness.
To make himself nothing even though he was very nature of God.
He made himself nothing taking the very nature of a servant being made in human likeness.
Being found in the appearance of men.

$^{361}$He humbled himself and became obedient to death even death on the cross.
therefore.
God exalted him to the highest place and gave him the name that is above every name.
that at the name of Jesus every knee shall bow in heaven on earth and under the earth and.
Every tongue confessed that Jesus is the Lord Jesus Christ is the Lord to the glory of God the Father.
That is the good news.
That is what the church is preaching.
That is what the Bible is all about.
Jesus did it.
To us.
Accomplish the salvation for all of us and die for our sins and he did he did it willingly as a slave.
Without his own self-will.
But only will of the Father to save you and me who are powerless in our sin.
So that we can be called children of God now, how how would this good news?.
affect or change our way and our attitude I.
Think it will change and it should change all of us radically.
Think about all the relationships that we have with our friends.
Without other church members brother and sister in Christ.
In fact what this passage is all about is about church life verse 5 your attitude.
Your attitude as a church to one another.
Your attitude should be the same as the attitude of Christ Jesus. He did it to himself.
He laid down his own right to die on the cross as a slave. We ought to use this kind of attitude.
Toward one another to care for one another.
That will affect or that would change our relationship with one another.
To love and to build up each other and forgive one another right?.
Now to extend this application it will affect.
Our marriages as well.
Right two sinners living together broken.
Not easy.
But in light of what Christ has done.
Yes, we're broken but we are totally forgiven right.
Because of the grace of God that should change how we relate to one another right how spouses relate to one another.
How spouses should forgive one another to the love one another totally.
And how it should change our way to live.
We are the instrument of God the great price of salvation.
Jesus has paid for is to redeem us so that we do not belong to anyone.
But God himself alone.
We are his instrument.
We are to be used by God himself.
We are the instrument of the gospel.

$^{401}$You go back to your campus you go back to your family you go back to your workplace.
You take a totally different role because God is using us.
Using us as instrument to invite other people to experience the grace of God.
Let me end with the Apostles Creek and if you have the same conviction I invite you to recite.
The Apostles Creek with me.
The Apostles Creek.
Let's do it together. I believe in God the Father Almighty.
creator of heaven and earth I.
Believe in Jesus Christ his only Son our Lord.
He was conceived by the power of the Holy Spirit and born of the Virgin Mary.
He suffered under Pontius Pilate was crucified died and was buried. He descended to the dead on the third day.
He rose again. He ascended into heaven and is seated at the right hand of the Father.
He will come again to judge the living and the dead.
I believe in the Holy Spirit the Holy Universal Church the chameleon of saints the forgiveness of sins.
the resurrection of the body and life everlasting.
Amen, let's all pray together.
Father we kneel our belief and understanding of your love of your grace and thank you.
For what Christ Jesus our Lord has done.
He is indeed the Lord of all.
He is indeed our King.
And we pray that you renew.
our sense of your grace.
So great of a salvation that while we were sinners while we were powerless while we were your enemies.
You came to save us.
Forgive all of our sins and our guilt and to heal us and to bring us back to you.
So that you.
our Holy God.
we can stand before you as.
justified.
Righteous all because of what your Son has done for us.
And we will give you thanks and praise you because you deserve all praises.
And.
Lord indeed.
We are redeemed.
by your great love and we pray that.
With your great love with your grace you renew and use us.
In the rest of the week we can be your instruments.
To bring peace and love and your grace.
To our family to our marriage to our.
Neighborhood to our campuses and pray that you will send us.

$^{441}$send us with this mission and.
Depending on your grace. We will be used by you for your glory in Christ name. Amen.
\newpage



\section{}
\label{sec:Bxy4FSvUGwA}
\textbf{Apostles' Creed \#5 (Virgin Birth) - Rev. Mark Chan}
\newline
\newline
連結: \href{https://youtube.com/watch?v=Bxy4FSvUGwA}{\texttt{https://youtube.com/watch?v=Bxy4FSvUGwA}} ~~~~ 語音日期: 2022-06-21
\newline
\newline
\hyperref[sec:vmXTNOZZzYc]{\small{< < < PREV SERMON < < <}}
~
\hyperref[sec:index_chronic]{\small{[返順時目]}}
~
\hyperref[sec:index_scriptual]{\small{[返順卷目]}}
~
\hyperref[sec:1_MORjvbY98]{\small{> > > NEXT SERMON > > >}}
\newline
\newline
$^{1}$Hey guys, happy Father's Day to all of you..
And I'm sure today's Father's Day,.
you probably want to hear something about some Father's.
Day story, right?.
Anyone?.
Yeah..
OK, all right..
A couple of Father's Day stories..
Again, it's not about-- as preachers,.
sometimes we use our children as our [INAUDIBLE].
and without their permission..
But here, I give you a couple of stories..
It's not about them, OK?.
So it's not about them, it's about me..
It's about me..
As a father, you probably would realize.
that we are broken as well, right?.
We are broken as much as anyone else..
And we need the grace of God..
So I've got a couple of stories of brokenness.
to share to you on this Father's Day..
I shared a couple of stories in the 10 o'clock service,.
so maybe I should just do this right here, right?.
The first story, when my daughter was about one year old,.
I remember there was one afternoon I was working really,.
really hard..
It probably was a Saturday..
And as you know, as a pastor, the most busiest day.
of the week is Saturday, right?.
It's Saturday..
You're getting ready for today, right?.
So I was preparing sermon..
And for some reason, the little girl.
was making some weird noise..
And for 10 minutes straight, I don't know,.
she was not crying, was not in pain or anything like that..
But she was just in the crib, and she was just.
making some noise..
I don't know why..
So I was kind of--.

$^{41}$if you know me enough, I can be an angry man..
I can be very, very angry in moments of weaknesses..
I can unleash my anger..
So that's the broken part of me, right?.
So I went into the room, and I started pushing to her..
So I had a 10-minute impromptu sermon, right?.
I forgot what I preached, but I think the central idea was.
that you cannot be selfish, right?.
You need to care for other people..
We have needs, too, right?.
I need to get things done, right?.
So I unleashed that angry sermon for 10 minutes straight..
So I will remember that..
And then I came back to my room and kind of calmed down..
And then I said to God, dude, that was my broken moment..
I shouldn't do that, right?.
So I confessed my sin and went back to the room..
I said, I'm sorry..
I was angry, right?.
My anger was not checked, right?.
And of course, I don't think she remembered that..
But that's OK..
I think fatherhood is vocation, right?.
You know what vocation means?.
Vocation means calling, right?.
Calling..
It's a calling that God has on us.
as fathers for a certain period of time in our lives,.
sometimes shorter, sometimes longer..
But we are being called to be a father, right?.
But I want to tell you, if you're a father,.
or you have a father, we're not perfect..
And we don't need to be perfect..
We all need the grace of God, right?.
Amen?.
We all need the grace of God..
And believe it or not, fathers, daddies, God loves you, right?.
God accepted you..
Entirely..
Not because you're perfect..

$^{81}$Not because you're better..
Just because who you are, right?.
We're sinners..
First story of Father's Day..
Second story is about my son..
This I also vividly remember..
I think he was about five to six years old..
And as Asian families, we want our kids to learn instruments,.
right?.
A lot of Chinese and Asian families,.
they don't really appreciate music..
But they just want their children to learn music, right?.
That's very ironic..
But I'm not one of those..
But I have my daughter and son go through piano lessons..
I remember one afternoon before the piano teacher arrived,.
and he was practicing some very simple song, right?.
Very simple song..
But for some reason, he didn't get it..
He didn't get it..
He just have that note wrong, right?.
Well, if you know me, I'm quite a self-disciplined person,.
right?.
I know if I put myself into doing something,.
I will get it done, right?.
A lot of fathers are like that, right?.
And that's a good thing..
That's also a bad thing, right?.
Because if you know me as a friend,.
you probably would not think I'm too imposing, right?.
But in my family, everyone else knows that I'm kind of--.
I want to impose..
I demand, right?.
And to live with me is quite an experience..
It's quite--.
For real?.
For real, yeah..
People suffer because I have high demand on other people,.
especially family members..
So he was practicing piano, and he was doing this chord,.

$^{121}$and something like this, right?.
[PIANO PLAYING].
Yeah, and he kept doing it wrong..
And to me, this is unacceptable, right?.
And at that moment, it was a broken moment..
It was a moment of weakness..
I unleashed my anger again..
So I'm yelling at him, right?.
Yelling at him..
I said, there's no way you cannot do this..
If you repeat that phrase, that measure, for 50 times,.
then you will do it, right?.
And then I forced him to repeat that once and twice and third..
After quite a few times, I realized.
that there was something wrong with me, right?.
And it just dawned on me..
What I want him to do, does it really.
matter compared to what God wants him to do?.
Do you know what I'm saying?.
It's not that important what I want..
It's what God wants, right?.
So after that heated moment, it was really heated..
Again, he doesn't remember that..
God is good..
[LAUGHTER].
After that piano lesson, I walked again..
I prayed, and I confessed my sin..
And I said, God, give me courage to apologize..
You know how that goes?.
A 37-year-old man, you need to apologize to a five-year-old?.
Like a real apology?.
It's not like, I'm sorry..
No, it's a real apology, right?.
So I walked up to him and I said, I'm sorry..
I was angry..
Can you forgive me?.
It was a humbling experience..
Humbling experience in the sense that we are all broken..
I'm broken..
I need the grace of God, right?.

$^{161}$We all need the grace of God..
Now I'll share these stories to you on Father's Day,.
because of course, some of us are fathers..
And some of you have fathers, right?.
They're still alive, right?.
Living with your father, living with your daddy..
We are broken..
We need to share the grace of God..
So that's the first I want to share with you before we.
go to the Apostles' Creek..
And no matter where you are in your life stage,.
our fathers are broken..
We, as fathers, we are broken..
And we need the love of God..
And the love of God is there..
1 John 3, verse 1..
"See what great love the Father has lavished on us,.
that we should be called children of God.".
And that is what we are..
The great love of our heavenly Father has lavished on us..
So that is a happy Father's Day to you..
But I want you to know that God loves you..
God's grace is always sufficient, right?.
Great love, lavishing on us..
Well, let's come to the loving presence of God..
Let's pray..
Father, we give you thanks because you love us so much..
You love us so much that your Son, your only Son,.
died on the cross for our sins, for our brokenness..
And we thank you..
Thank you for your loving presence..
Thank you for this relationship that we have,.
that we are called your children, children of God..
In our brokenness, in our humbleness, in our sinfulness,.
we meet your grace..
We experience your love..
And thank you for the cross..
Thank you for what you have done for all of us..
I especially want to pray for the daddies and fathers.
in our congregation..

$^{201}$I pray that your grace will be lavishing on them,.
that it is not easy to be a father, especially.
in this generation..
I pray that your grace will be sufficient for them.
to meet those challenges..
And even though they are not perfect,.
even though we are broken, we know that you love us fully..
You accept us at the foot of the cross..
And we thank you for this grace..
In Christ's name we pray..
Amen..
All right, so the end of Father's Day sermon..
OK, all right..
So we go back to the Apostles' Creek..
All right, I'll cut it short..
And I'll make sure we get out after 1225..
OK, we come to this section, the virgin birth of Jesus Christ..
The Apostles' Creek, again, is 12 statements..
It's like a road map of what Christianity is all about..
And here it says, "I believe in Jesus Christ as only Son,.
our Lord..
He was conceived by the power of the Holy Spirit.
and born of Virgin Mary.".
Jesus' birth is a miracle birth..
Now, as modern people, we hear miracle birth, all right,.
virgin birth..
We automatically think in scientific terms, right?.
Have you ever asked the question,.
what's the DNA of Jesus?.
Right?.
Now, definitely not from Joseph, right?.
But where's it from, right?.
What does it look like?.
Now, we come to the virgin birth or miracle birth.
with the scientific thinking..
But I think for the biblical readers of this first audience,.
they would not think in terms of scientific DNA.
and all that stuff..
They would be thinking along the line of the Old Testament..
The Old Testament, if you remember,.

$^{241}$there were tons of miracle birth stories, right?.
Can you name some of them?.
Isaac, right?.
Jacob and Esau..
And there are many other miracle births..
If you track the miracle birth stories.
of the entire Old Testament, every time there's.
a child miraculously birth, it is.
because God is advancing another chapter in his redemption.
plan for mankind..
Isaac was born, opened another new chapter in Abraham's story,.
in the story of Genesis..
And then Esau and Jacob was born as the twin brothers,.
miracle birth again..
And another chapter was open..
So they would be thinking about God is doing something new..
So when the Bible talks about the virgin birth of Mary,.
when Jesus was birthed, the Bible.
is talking about finally God is fulfilling his promises..
Whatever he has said in the entire Old Testament,.
the covenant with Abraham, the covenant with David,.
and all this stuff, God is faithful..
And only God can do that..
That is why the birth of a child is a miraculous birth..
Virgin birth, of course, is the miracle birth, right?.
So again, we don't come to the Bible always.
with our modern mindset, right?.
I think I have talked about it in the past few parts.
of the sermon series..
I think it's sad..
It is really sad that even today, in 21st century,.
we always think that faith and science are in disharmony..
They don't go together..
And they must be being polarized..
Again, when we come to the miracle birth of Jesus,.
virgin birth of Mary, we don't need.
to think in scientific terms..
There is no contradiction between miracle.
and how the natural world works..
They can be in harmony..

$^{281}$I was born-- I grew up in a very conservative church..
And I remember back then in the Bible studies.
and in Sunday schools, we always have.
to feel the tension, the huge tension,.
between faith and reason, faith and intellect..
We opened the Bible, the first few pages, Genesis..
There's always this harmony, kind.
of a tension between what Genesis 1.
has to say about the beginning of the universe.
and then what natural science has to say..
Now, I want to tell you, I want to emphasize this..
I think this is really, really sad that we have.
to pitch one against the other..
There's no such thing, because the Bible is not.
a scientific textbook..
It doesn't tell us what the early history.
of the physical universe..
The ancient people, the biblical world,.
does not concern with scientific questions.
that we concern about..
So there has to be some way we can bring the two together.
in harmony..
Again, as I said before, Genesis 1.
is answering the who question, not the how question..
The how question is answered by scientific enterprise..
The who question, the theological enterprise,.
is what scripture is trying to do,.
trying to tell us who behind the universe.
and who we should worship..
In 2020, I mentioned this before,.
Dr. Francis Collins accepted Templeton Prize..
It was during COVID..
It was a huge concern, so it was a private kind of event..
It was live streamed on YouTube..
And Collins said, as a medical doctor, in his whole life,.
he has been atheistic..
He doesn't believe in God..
And then one day, he was challenged by his patient..
And then it just dawned on him that science did not help much.
with the profound question of meaning..

$^{321}$What do you live for?.
The meaning of life, the value, right?.
Science cannot offer much help..
And Collins said, at that time, he.
realized his atheism is dangerously thin..
He was quoting Chesterton, one of the great theologians.
in the 19th century..
He said, atheism is the most daring of all dogmas.
because it is the assertion of universal negative..
There's no God..
So he became a believer in two years,.
something written by C.S. Lewis and talking to other people..
And then finally came to faith in Jesus Christ..
Take notice of the acceptance speech..
The title is "In Praise of Harmony.".
Harmony between science-- he's a top-notch scientist--.
and faith, an evangelical Bible-believing question,.
these two can go together..
So these two, faith and reason, can go hand in hand..
And I hope this is your experience as well..
There's no need to just deny human reasoning when we come.
to faith in Jesus Christ..
So virgin birth is beyond our finite human reasoning..
It is beyond..
But that does not mean it's unreasonable..
That does not mean it's irrational..
When something goes beyond our finite reasoning,.
we do not need to conclude that this is irrational..
It can be-- probably the best term.
is trans-rational, beyond our rationality..
Our brain is finite..
Our mind is finite..
But God can go beyond all of this,.
surpassing our finite rationality..
In fact, when the angel announces the virgin birth.
to Mary, how did Mary respond?.
Luke 1, 29..
"Mary was greatly troubled by what the angel said to her.
and wondered what kind of greeting this might be..
But the angel said to her, 'Do not be afraid, Mary..

$^{361}$You have found favor with God..
You will be with child and give birth to a son..
And you are to call his name Jesus..
He will be great and will be called the Son of the Most High..
The Lord God will give him the throne of his father David..
And he will reign over the house of Jacob forever..
His kingdom will never end.'.
And Mary responded by saying, 'How will this be?.
How will this be since I am a virgin?'".
The angel continues..
"The Holy Spirit will come upon you, will have power on you..
So the Holy One to be born will be called the Son of God..
Even Elizabeth, your relative, is.
going to have a child, another miracle child,.
in her old age, which is John the Baptist..
She who was said to be barren is in her sixth month..
For nothing is impossible with God..
I am the Lord's servant, Mary answered..
Finally, she obeyed..
And she said, 'May it be to me as you have said.'".
Mary's faith, Mary's obedience, the first step.
to take in what the angel has said to her,.
what God is going to do through her life..
Now, here's the first takeaway, the doctrine of virgin birth..
First, Jesus came to this world in a humble manner..
That's why we come to God..
We come to Jesus..
We need to be humble as well..
Now, I don't have a picture here in Israel,.
in the town of Bethlehem..
Today, there is a church called Church of Holy Nativity..
Holy Nativity, meaning remembering the birth of Jesus..
You walk in that church, it used to be a really, really big door.
to welcome people..
But century after century, there's renovation,.
there's changes to the door or the structure of the church..
Now, it became a really, really small door..
And you have to duck to go in..
And the door is called the door of humility..
You come to see the birth of Jesus, the place he was birthed,.

$^{401}$according to traditions..
And your first action you need to do is to be humble..
And that's the first takeaway..
And God is infinite, and he's beyond our finite mind..
Yes, we cannot understand how it works,.
but God demands our humble trust..
So there's nothing irrational about Jesus' birth..
Yeah, it's trans-rational..
But I think the second takeaway is from another passage.
talking about the virgin birth of Jesus.
in the Gospel of Luke, chapter 2..
Now, you know what happened here is that Joseph and Mary,.
they're traveling from their workplace in Nazareth.
in the north and all the way to Bethlehem.
while Mary was still pregnant..
What happened was the Roman Empire.
imposed some sort of taxation census..
So everybody has to go to their hometown to register..
It's kind of an imperial policy..
So those who are being oppressed by Joseph and Mary,.
the normal, typical Judeans, they.
have to go back to the hometown..
This is how the story goes, chapter 2, verse 4..
So Joseph also went up from the town of Nazareth in Galilee.
to Judea, to Bethlehem, his hometown, town of David,.
because he belonged to the line or house of David..
He went there to register with Mary, who.
was pledged to be married to him and was expecting a child..
While they were there, the time came for the baby to be born..
And she gave birth to her firstborn, a son..
She wrapped him in cloths and placed him in the manger.
because there was no room for them in the inn..
Or I guess the better translation is the guesthouse..
No matter how you understand it, maybe there's.
a Motel 6 in Bethlehem, and it was full..
There was no vacancy for them..
Or they couldn't book the Airbnb at Bethlehem.
or whatever that is..
I think it can be translated properly as a guesthouse.
because in the typical Palestinian house, even today,.

$^{441}$sometimes-- so often, they have two floors..
And the upper floor, as you can see in this reconstruction,.
is the guest room, the bedrooms..
It's up there..
Down there, the lower level, is where the manger is,.
where the production happens..
They would cook, or they would have the animals in there,.
especially in wintertime..
They would keep them warm..
That's how a typical Palestinian house looks like..
And the guest room, the Gospel of Luke says, has no room..
Now, you can imagine in that sort of census.
across the Roman Empire, everybody.
is going to their hometown..
And it's so crowded..
There were other relatives, there.
were other friends who were staying at the so-called Airbnb.
on the second floor..
And Mary showed up with Joseph..
And he was about to give birth, right?.
So she went down to the first floor..
However you want to interpret it, in or guest room,.
it really doesn't matter because the message is the same..
He was born in the manger, down there, where the animals are..
Not in the palace..
That's another takeaway..
The Son of God came to our world, our broken world,.
in the manger..
We draw close to God every Sunday..
We better check our attitude when we come to worship..
Sometimes we come to worship as if we come to a concert..
Sometimes we come to worship as if we are watching a show..
There's always customer reveal at the end,.
because I'm the center..
I come to be entertained..
I come to be--.
somebody has to satisfy me, right?.
Well, we go to a restaurant after worship service,.
go to a restaurant if the service is bad, right?.
What do you do?.

$^{481}$You walk away, right?.
What do you do?.
You give them one star, right?.
But not so in the church..
Not so in worship..
Worship is all about exalting God..
And we come and we bow..
We need to learn the lesson of humility,.
because the Son of God came as a lowly, lowly man,.
born in the manger..
Next, we should consider what the Catholic tradition has.
to say about Mary..
That's another takeaway..
Now, the Apostles' Creek, as you can tell,.
is different from creeks or statements.
issued by the Catholic tradition..
In the Catholic tradition, in the general official teaching.
of the Catholic church, Mary is the co-redeemer..
Mary is the co-redeemer who brings the Messiah.
into the world, right?.
And Mary, accordingly, was sinless..
The perpetual virginity of Mary--.
you've probably heard of that--.
she remains pure in the rest of her life..
And some traditions in the Catholic teaching,.
Mary never died..
She fell asleep, and she was taken up to heaven..
Now, all of this are unfortunate,.
because as you can see, there's no biblical warrant.
for all this teaching or all this concept..
There's nothing intrinsically special about Mary..
Now, in Jesus' day, Mary was probably a high school girl--.
16, 17, or even 15, 14-year-old, right?.
This is a young girl, young Hebrew faithful girl..
The angel told her, Messiah will come through you..
And she responded with obedience, with faith..
So there's nothing intrinsically special about Mary..
It's not about Mary's status..
What is praiseworthy about Mary is not about herself..
What is praiseworthy is about her faith..

$^{521}$It's her obedience..
That is special..
Luke 1, 45-- Elizabeth was meeting with Mary..
Both of them are pregnant..
And Elizabeth, you remember, was pregnant with John the Baptist,.
the cousin of Jesus..
And then Mary was pregnant about to give birth to Jesus..
Both mothers met together..
And Elizabeth, the mother of John the Baptist,.
speaks about Mary in this manner--.
45-- "Blessed is she who believes.".
Blessed is she who believes..
It's her faith that matters..
It's her faith that is our example..
Now, I don't know your relationship with God recently..
Maybe God is calling you to do something,.
to witness to someone, to serve in a certain capacity..
No matter what God is calling you, maybe you're struggling..
Maybe you're waiting..
I'll tell you this..
God is advancing His plan through you..
God will..
But you're called to obey..
You're called to respond..
It's your faith that matters..
Sometimes we struggle with the calling of God..
Well, this is mundane, right?.
Mundane stuff I'm doing, maybe in my workplace,.
maybe in my family..
Just go through the routine day after day and day after day..
Nothing is really happening, right?.
But think about it..
She who believed, she said yes to God..
Yes, your will be done..
And God-- only then God can use us in our neighborhood,.
in our church, in our workplace..
We just say yes to God..
All the details, God will figure out..
But you need to trust in Him, right?.
Whether you're moving away, whether you're.

$^{561}$serving or stick to some sort of commitment, trust in God..
And God will work out the details..
All right, let's pray together..
Father God, we give you thanks for you.
are working in this world..
Sometimes we cannot see it..
And most of the time, you seem invisible..
But in scripture, story after story, page after page,.
you tell us you are faithful..
And you are actively doing something..
And you want to do it through us..
And we pray that no matter what we are struggling with,.
give us faith..
Give us the-- just to fully trust in you,.
that you will be working out your purposes in our lives..
You will be fulfilling your plan in our lives..
We pray that we will learn to be faithful..
We will learn to trust in you..
Thank you for being our faithful, faithful father..
Thank you for being our loving God..
In Christ, we pray..
Amen..
Well, God bless you..
And remember, happy--.
\newpage



\section{}
\label{sec:1_MORjvbY98}
\textbf{父親雅各的轉變 (創 32 章) - 陳崇基牧師}
\newline
\newline
連結: \href{https://youtube.com/watch?v=1-MORjvbY98}{\texttt{https://youtube.com/watch?v=1-MORjvbY98}} ~~~~ 語音日期: 2022-06-21
\newline
\newline
\hyperref[sec:Bxy4FSvUGwA]{\small{< < < PREV SERMON < < <}}
~
\hyperref[sec:index_chronic]{\small{[返順時目]}}
~
\hyperref[sec:index_scriptual]{\small{[返順卷目]}}
~
\hyperref[sec:5Ocf3zsSFnw]{\small{> > > NEXT SERMON > > >}}
\newline
\newline
$^{1}$祝父親節快樂.
今天是父親節.
今天我們說一個單元.
不說一些系列.
我也不繼續以前的系列.
說一個跟父親節有關的故事.
在我們社會的氣氛下.
父親節是沒有什麼人理會的.
相比於母親節是非常盡色的.
我預備的時候.
讓我的小朋友看到這張照片.
我找到六月有什麼特別的節日.
是國際的節日.
其中一個節日是國際的碎星日.
6月30日 你有沒有聽過.
沒有 沒聽過.
大部分人都沒有聽過.
沒有什麼人理會.
小朋友問我為什麼要搞這些.
就好像父親節一樣.
父親節好像碎星日一樣.
誰理會.
除非女兒帶你去公園.
或者你買一輛新車.
這樣就不同了.
但通常都沒有人理會.
特別是跟母親節相比.
但是今天怎樣也要父親節.
我們也應該說說父親節的事.
不如我就用一個聖經的故事.
跟大家說.
首先我想問問大家.
你覺得父親節給母親節訊息.
是什麼原因.
有什麼原因.
是不是因為男人不驅小節.
還是太多失敗的例子.
父親節失敗的例子.
還是男人比較短命.
大家覺得怎樣.

$^{41}$有誰覺得是第一個.
男人不驅小節.
所以我們不太在意父親節.
有沒有人跟我們慶祝.
有誰.
有 這些不驅小節.
有誰覺得太多父親節失敗的例子.
所以父親節沒有人理會.
有沒有.
沒有 不敢舉.
第三個是誰覺得男人比較短命.
所以我們不太在意.
父親節有沒有人慶祝.
有沒有.
OK.
這些可能都是一些事實.
父親節沒有人理會.
但是如果你回到聖經的話.
有很多父親的例子.
當然這個是聖經當時的文化.
以男為上的文化.
以父權為上的文化.
但是你看到聖經裡面.
從頭到尾都有很多父親的故事.
有一些聖經的故事.
或者聖經的事蹟.
是一些很有人情味的事蹟.
在創世記裡面.
其中一個最有人情味.
或者記載得最多的.
就是一個父親的故事.
雅各的故事.
雅各 你記得.
阿布拉罕 以撒雅各.
第三代的族長.
在聖經裡面有記載很多他的故事.
或者這麼說.
聖經裡面對他的刻畫.
是非常細緻的.
除了可能大衛 彼得 新約等等不同的人物.

$^{81}$雅各是特別多的.
當然因為雅各就是以色列國的先祖.
如果你說阿布拉罕是的話.
雅各更加是.
因為雅各所生出來的.
十二眾子.
然後成為以色列國.
以色列的民族.
他的名字由雅各變成以色列.
所以為什麼這麼重要.
如果你看看聖經裡面.
如何去刻畫雅各這個人.
你就大概get到.
從頭到尾都說.
雅各是一個怎樣的人.
而且上帝在他的生命裡面.
做了些什麼.
25章 尤節.
這裡剛剛開始出世的時候.
聖經裡面有一些典型的場面.
典型的刻畫.
是叫type scene.
說到一個人出世.
或者這麼說 形容一個人出世的時候.
就是形容這個人的character.
到底是怎樣的 他的本性是怎樣.
這裡說到雙胞胎.
他的哥哥 以素.
和雅各出世的時候.
隨後以素的弟弟也出世了.
因為雙胞胎.
但是他的手抓住了他哥哥.
的腳跟.
因此就為他起名叫做雅各.
你見到希伯來文.
經常都會玩一些字.
Jacob就是雅各.
腳跟就是阿Keth.
所以那個字有些諧音.
就說這個雅各.

$^{121}$是一個很喜歡搶東西的人.
或者一件事.
他自己控制了.
你記得雅各的故事.
我們曾經提過.
第一次.
他在他哥哥以素最軟弱的時候.
用一碗紅豆湯.
來搶到什麼.
丁權.
那個長子的名分.
就是有double portion的名分.
然後第二次.
就是他.
假扮以素.
然後去到.
加夫.
就拿到那個.
的blessing.
拿了所謂的patriarchal blessing.
就是在富家裡面.
他有的祝福.
他有的權柄.
兩次都是不擇手段.
這個其實.
可以這樣說是生意人.
為你仕途.
這個就是雅各.
但想不到因為他.
搶了以素的.
最後的blessing之後.
以素就起殺機.
要去殺一個雅各.
殺他的親兄弟.
於是他就要逃難.
雅各就逃亡.
離開了自己的本家.
就去投靠他的舅父.
叫做拉班.
去到另外一個地方.

$^{161}$去到.
另外一個異地.
在那裡他就.
見到一個女孩.
叫做拉傑.
在那裡就.
愛上了拉傑.
29章20節.
就因為愛上了.
拉傑的緣故.
他就幫拉班.
服侍了他,做了苦工7年.
如同幾日.
聖經這樣說,就是因為愛他那麼深.
誰知你記得那個故事.
洞房花燭夜之後.
就發生什麼事.
原來那個不是拉傑.
原來拉傑的姐姐.
就是尼雅.
我們很多時候.
看雅各的故事.
就當他是一個愛情故事.
說他深愛拉傑.
然後服侍拉班.
答應他服侍7年也好.
他都當幾日.
如果你不當一個愛情故事看的話.
其實是當一個報應來看.
或者可以這樣去讀.
以往他是不擇手段.
樣樣東西都好像.
自己掌控了.
掌控自己的名份又可以騙到.
祝福又可以騙到.
用自己不同的方法不同的手段.
可以拿到自己最想要的東西.
誰知去到拉班的時候.
強中自有強中手.
他被人找了笨.

$^{201}$結果就先取尼雅.
然後就被拉班找了他笨.
就要再服侍多他7年.
才可以得到拉傑.
他心裡面最愛的女人.
強中自有強中手.
以為自己人定勝天.
誰知他的舅父.
比他更厲害.
比他更高昌.
這個是上帝對他的管教.
上帝對他的訴讀.
然後我們在29至30章.
就見到眾子的誕生.
11個的兒子.
十二支派裡面11個.
就是在這兩章三章裡面交代.
11個兒子出世.
就是將來十二支派的先祖.
你看這些故事的時候.
我們很多時候會留意.
爭風吃醋不同的老婆.
不同的侍女.
誰去生誰.
但有一句說話很重要.
在這段的記載裡.
就是甚麼說話.
就是耶和華洗他的妻子生.
哪一個賜給他兒子.
或者哪一個賜給雅各.
能夠在異地裡面.
他的產業都能夠這麼豐厚.
是上帝.
沒有其他原因.
不是因為他厲害.
不是因為他自己手段高強.
而是上帝去以恩典去賜給他.
在這裡大概中.
他的故事中段的時候.
你就見到上帝想在他身上做些甚麼.

$^{241}$或者在以色列的國父身上.
想做些甚麼.
在這個父親身上.
就是告訴他.
你之所以能夠成家.
不是因為你手段高強.
或者你有強勢.
而是耶和華去賜福給你.
在當中做父親的.
你有沒有想過你第一次.
和你兒女的第一次.
有很多第一次.
第一次產房接生的時候.
抱起你的兒女.
第一次和他玩騙.
教他走路.
教他踩單車.
或者甚至我們比較年長的.
見到他成家立室.
有沒有想過這些第一次.
是一回甚麼事.
雅各這個父親的故事.
就是說不是己令.
是上帝的恩典.
完全是上帝的恩典.
你才能夠成家.
我才能夠成家.
有一個想破碎雅各.
破碎以色列祖父的一個經歷.
就是不是靠他己令.
能夠成家成國.
他謙卑下來.
不是靠你成功.
是因為上帝的恩典才能夠成功.
講起謙卑.
我們做父親的.
我們很多時候都會有些.
令我們自己謙卑的經歷.
我講一個.
大家認識我這麼久.

$^{281}$都知道我對自己要求很高.
這件事是好還是不好.
是好事是嗎.
不好事是甚麼.
Percy你覺得.
有難有不難.
那些跟我一起生活的.
就很辛苦.
我覺得真的.
可能其他家外的人.
他跟我接觸的時候.
覺得他都很容易相處.
但家裡就未必.
一定每次都是這樣.
我記得五六歲的時候.
我兒子學琴.
我記得有一次.
鋼琴老師來家裡上課之前.
他要練琴.
那時候五六歲.
像鬥陣一樣.
他就要練琴.
他練來練去都練不到.
我說怎麼會練不到.
我就坐他旁邊.
我說你彈五十次一定練到.
沒有可能練不到.
這個世界有這些東西.
你肯做就行.
我這樣說的時候.
現在我當然氣定神彈跟你說.
但那時候我是失控的.
我自己是很大聲罵.
罵到一個地步.
我自己都有一點害怕.
你有沒有試過你發脾氣的時候.
你自己都覺得很害怕.
你怕你自己.
那時候我走開.
走開繼續彈.

$^{321}$我也懂一點音樂.
我聽到他彈錯.
然後我就在走廊另外一邊.
我就大聲罵他.
我說為什麼那個Bb會彈錯.
你再彈一次.
結果那次之後.
我自己靜下來.
我發現我自己失控.
發現我自己脾氣太臭.
我就祈禱.
我就跟上帝說.
我都要謙卑下來.
那也是有一點困難的.
因為我要跟兒子道歉.
很正式跟他道歉.
我求上帝給我恩典可以做.
結果晚上我就跟他說.
我很對不起.
爸爸跟你發脾氣.
你可不可以原諒我.
我不知道你有沒有試過這樣做.
跟一個五六歲的小孩.
去跟他說.
解釋得清清楚楚.
我剛才做了什麼.
是得罪他.
是不應該這樣.
你可不可以原諒我.
去禱告去認罪.
要很大勇氣.
那次之後.
我就對他有一點轉變.
我就覺得.
我正在想的就是.
到底我是想他做我心裡面的事.
去達到我的要求.
還是應該是我引導他.
去達到上帝要在他心裡面做的事.
兩件事很不同.

$^{361}$我自己都有一個這樣的經歷.
道完業之後.
一個這樣的操練.
到底我作為一個父親.
我要摸到他一個地步.
是要合我心意.
還是我是一個渠道.
讓上帝去塑造他.
那次經歷之後沒多久.
我就跟鋼琴老師說.
他們都不學了.
鋼琴老師就很失望.
他說為什麼不學.
其實我跟鋼琴老師解釋.
我說我沒錢.
沒錢.
但其實我覺得.
他這樣學下去還是不行.
所以你看到他一直都沒有彈鍵盤.
可能是這個原因.
可能而已.
所以他打鼓.
就不是我教的.
因為我知道如果坐在他旁邊.
我教的話.
那又是死的.
所以我就彈開.
好像我女兒有句金句.
Chill out.
Take a chill pill.
冷靜點.
為父的經歷.
可以是一個謙卑.
我們自己.
讓我們知道上帝在做什麼.
不是靠我們自己.
那套故事都是一樣.
到第三十一章.
三十一章這裡發生什麼事.
他就成家了.

$^{401}$有很多的財產.
然後上帝要他回到迦南地.
他就火速的離開了.
因為他叔叔對他非常差.
三十一章二十二節.
到了第三天趕路的時候.
有人就告訴拉班.
他說亞國原來是偷走離開了.
拉班就火速的去追他.
追了七天.
然後去到基列山的時候.
就追上他了.
這一次是.
另一次亞國遇見撒基的時候.
二十四節.
夜間上帝就跟他的叔叔說.
報夢給他.
他說你要小心.
不要向亞國說好說歹.
原來又是上帝去保守.
上帝去帶領著他.
以致他又逃離了撒基.
然後下一次.
就有一個更大的威脅.
就是他回到迦南地的時候.
最大的威脅就是.
面對那個曾經想殺他的爾素.
爾素帶著幾百人.
就朝他的方向來的時候.
三十二章第二七節.
亞國就很懼怕.
而且受煩.
這個第一次直接告訴我們.
亞國當時心裡面在想什麼.
他就跟他的同伴的人.
同行的人.
他們的羊群,牛群等等.
他就說如果爾素來擊殺我們其中一隊的話.
剩下另一隊可以逃脫.
他就將他的財產分成兩隊.

$^{441}$以免爾素來到就一鋪過.
什麼都沒有了.
這是一個最大的危機.
一個最大的威脅的時候.
他就去將他的東西.
打理好.
用他自己的己力.
去分成兩隊.
然後他就發現.
還是搞不定.
祈禱.
在十字.
他向耶和華祈禱.
這是其中一段.
他說.
你向僕人所施的一切慈愛信神.
我一點也不配.
我先前只是用我一斤將過約旦河.
但現在我的財產已經有兩隊了.
求耶和華去脫離他哥哥的仇.
脫離仇家爾素.
因為很怕他來殺我.
將我所有的兒妻都會殺盡.
在這段時間.
聖經的刻畫是很細緻的.
幾乎是鏡頭慢下來.
是有一個慢動作.
去形容雅各發生了什麼事.
亞爾哲.
他夜更起來.
就帶著妻子.
兩個婢女.
十一個孩子.
就過了雅各渡口.
他帶著他們.
送他們過河.
他所有一切也都過河了.
只剩雅各一人.
只剩下他一個人.
在這裡你不能不想想.

$^{481}$到底他在做什麼.
這是不是一個為父的榜樣.
好像很不負責任.
過河就要去奸人地.
但他叫所有奸人過河.
然後發生什麼事.
他自己就留下了.
可能是.
可能是他恐懼太大.
太怕.
所以他就未過河轉.
只剩下雅各一人.
在這個時候.
就是聖經裡面最著名的一個故事.
就是雅各和天使摔跤的記載.
三十二章這麼長的故事.
我們不能夠逐節逐節說.
但這裡說到什麼呢.
他有一個人來跟他摔跤.
跟他打架.
直到黎明的時候.
那人就是天使這個使者.
其實就是神之記.
見到他不能夠勝過雅各.
就摸了他大腿窩一下.
雅各的大腿窩就在那個時候就流傷.
那人就說天快亮了.
讓我走.
不要再跟我打.
他說你不給我祝福的話.
我就不會讓你走.
這個也是雅各的本性.
是要搶住.
拿住.
去爭取住.
上帝給他的祝福.
那人就跟他說.
有一個名字的對話.
名字就是權能的一個對話.
有名字就是有權能.

$^{521}$你叫什麼名字.
你知道我的名字是什麼.
意思就是上帝將他的名字留下.
上帝將他的名字留下.
就是你不能夠濫用我的名字.
或者你可以這樣看.
你不可以濫用上帝的權能.
不是因為你自己多厲害.
不是因為你自己有權有勢.
你就可以濫用上帝的名字.
上帝的祝福.
上帝的權能.
於是他就將他改名.
變成以色列.
Israel就在這裡.
以色列的國父.
然後最後他受傷.
雅各就將那個地方的名字.
改成比魯伊林.
Peniel.
意思就是Face of God.
這段聖經很奇特.
有人說這段聖經是每一個男人.
應該要經歷的.
是屬靈的黑夜.
要在夜間裡面.
和上帝去糾纏.
到底上帝想做些什麼.
這裡是勝不過他.
這裡肯定不是說雅各大力一點.
可以勝過上帝.
不可能.
因為你看到他最後的時候.
摸他一下.
他就扭傷了.
這似乎告訴我們.
是雅各一貫的作風.
他一貫的本能.
一個很典型的爸爸.
是人定聖天的一種本能.

$^{561}$每樣東西都好像他自己可以控制.
他是一個強人.
雅各是一個強人.
Strong man.
但正是這個黑夜裡面.
上帝要做什麼.
就是要他倒下.
就是最關鍵的時候.
最大衛脅的時候.
他就是要他跨低入泡.
他有一點像今天.
一個強人.
一個像什麼都能搞定的爸爸.
突然間就進了ICU.
有一種這樣的經歷.
心靈黑夜的經歷.
就是上帝要在他身上去做一些事.
33章.
這裡就要見耳掃了.
帶著四百人一齊來到的時候.
雅各有什麼反應呢.
雅各就將他的孩子.
跟著莉雅和拉傑的兩個婢女.
將他們的兒女分散了.
然後就叫兩個婢女和他的孩子走在前頭.
莉雅和她的孩子就跟在後面.
最後是誰.
最後是拉傑和約瑟.
因為這兩個人.
拉傑你記得等了十四年.
十四年的苦工才可以換取拉傑回來.
他所生的唯一的兒子.
是約瑟.
他就將他放在最後.
他自己卻走在他們前面.
一連七次的伏在地.
才靠近 靠近他的哥哥.
他的排隊是別有用心.
最深愛的就放在最後.
你怎麼看他多麼軟弱都好.

$^{601}$這個時候似乎他有些改變.
可能跛了腳 有些改變.
他自己卻走在他們前頭.
再不在後面.
走在整個家的前面.
開頭到現在.
一個生意佬 是不是.
一個唯利是圖的人.
現在危機來了.
他就走在全家的前頭.
七次的低頭 為家人去解困.
再不是為了自己的命.
我想他可能.
這個細節裡面.
你見到有一個這樣的.
所謂的transformation.
雅各可以成家是因為他知道.
因為上帝的恩典.
我一點也不配得.
現在就走在眾人的面前.
他知道他自己要為父.
去保障家人.
這些很仔細的描述.
我不知道家人怎麼看他.
聖經也沒有說.
但可能大家可以想像.
把家父跛腳.
走在面前.
去迎接這個最大的危機.
不知道後面的家人怎麼看.
父親節的時候.
我想起以前中學.
有一篇文章.
中學的時候讀的.
是朱自清寫的背影.
不知道你有沒有印象.
朱自清說到他爸爸.
和他父親的關係的時候.
說到他爸爸送他離開南京.
去北京上大學.

$^{641}$他就見到.
和他爸爸一起在月台那裡.
他爸爸問他要不要吃柑吃橙.
於是他就爬過月台.
爬到軌道.
然後再爬上來.
為了他的兒子.
送他一跌的時候.
就給他幾個橙.
背影這篇文章.
就是寫到朱自清.
怎麼看到他爸爸的背影.
怎麼去餵著兒女.
做一些小的動作.
這些是一些很仔細的描述.
我想這裡也是一些很仔細的描述.
雅各的背影.
可能在他兒女的眼中.
我不知道他們看到了什麼.
但這裡的刻畫.
是說到一個老父.
父親去保護自己的兒女.
這些實在是天經地義的事.
但你看完整個的故事.
你就會發現不是雅各的保障最重要.
是耶和華的保障.
由頭到尾都是耶和華自己的恩典.
他能夠勝過每一個危機.
能夠成家.
能夠在一個最大的威脅裡面.
他能夠保得住.
是因為上帝的恩典.
整個故事不是雅各有多厲害.
整個故事是上帝是成家者.
上帝是你和我的成家者.
各位父親.
是上帝的恩典.
你可能真的很厲害.
但你不要忘記.
不是上帝保守著你.

$^{681}$我們早就玩完.
你能夠抱起你的小孩.
能夠帶著你的小孩長大.
上帝的恩典.
最終你要知道.
是上帝的主權.
上帝的恩惠.
不是你.
我家族認識一個朋友.
認識了很多年.
由他中學開始.
我們都見到他是一個音樂很厲害的人.
他一生的夢想就是搞音樂.
搞他喜歡的音樂.
年輕的時候已經很有恩慈.
他和他爸爸很大爭拗.
他爸爸是很典型的華人.
我們很實際.
他爸爸也是大力反對他搞音樂.
非要他做工程師.
非要他讀工程學院.
他就不肯.
結果兩父子的關係都是非常差.
雖然這個爸爸信了主.
但到今天還是不欣賞他.
我聽說這個人已經四十多歲.
成家了.
偏偏就是他爸爸到今天還是不欣賞.
好幾年前.
他拍了一套紀錄片.
一套短的紀錄片.
是講移民的.
因為他是ABC.
講一些華人的經歷.
我們夫婦就去看他上映.
完了之後有個茶會.
我們就和他聊天.
他就和我說.
他說他今天.
他爸爸一點都不欣賞他做的事.

$^{721}$你想一下.
這個人已經成家了.
並且這個是一個很好的弟兄.
很有呼召感.
很清楚上帝要在他身上做些什麼.
是 他可能不是飛黃騰達.
可能不是真的很多錢.
但我見到這個弟兄.
很清楚自己上帝要他做的事.
可以和上帝交仗.
頂天立地.
偏偏生命裡面就是他爸爸不肯放手.
不肯放手.
是一個很自然的表達.
作為一個華人的父親.
我們現在想保護我們的兒女.
天經地義.
當然.
給他最好的環境.
入最好的學校.
搬去最好的地方.
這些全部都是我們作為父親.
我們想去保護他.
想去培養他.
想使他在一個最好的環境裡長大.
但我們一定要問清楚自己.
最後最後不是要按我的主意.
是要按天父的主意.
天父在他身上要做些什麼.
我想這個問題才是最打緊的.
最打緊不是我們的理想.
或者我們的想法.
而是上帝對他的想法.
上帝對他的成長.
要想塑造他成為一個怎樣的人.
我們自己不過是一個工具.
在基督教的倫理裡面.
有一個字叫做vocation.
呼召.
vocale.

$^{761}$calling.
在拉丁文.
就是有vocation的一個字.
什麼叫vocation.
vocation.
你問一個人有沒有calling.
有沒有呼召的時候.
你怎樣測試一個人有沒有calling.
就是那個人說.
這個角色不是為了人工.
是為了自己的成就感.
是有一種超然性的transcendent的本質.
就是要對上帝負責.
超然性的.
去榮耀上帝.
最終的問題.
有沒有calling有沒有vocation.
就是他知道要向上帝負責.
要討厭上帝.
唯一要討厭上帝.
他就要盡忠職守.
這個就叫做vocation.
呼召.
我是Daddy.
我是做父親的.
是一個vocation.
是一個calling.
這個calling.
是有限期的.
只有兩個可能性.
第一個可能性.
那個限期就是我過期了.
我要離開.
完結這個vocation.
第二個限期可能是.
他要離開.
讓我先離開.
完結這個vocation.
長或短都好.
那個calling是出於上帝.

$^{801}$那個calling是神的旨意.
追打的.
我們看我們的耳裡.
追打的最重要.
是不是他能夠遵從主的旨意去走.
做一個頂天立地的人.
也不是千萬不要弄錯.
只是去養育一個兒女.
做一個nice person.
這麼簡單.
這個不是基督徒的calling.
基督徒的calling是我成為上帝的工具.
我去塑造我的兒女.
以至他能夠向上帝去教長.
新教,延教.
樣樣做齊.
為了什麼呢.
就是為了養育他.
能夠帶他去上帝的面前.
這個是我們的calling.
上帝在他心靈裡面.
能不能夠作主.
我能不能夠在我這個calling裡面.
你都可以盡忠職守.
不是為了benefit.
transcendent的.
相投意理的.
這個calling.
或者看回這個雅各的故事.
似乎都有一個這樣的方向.
因為說很多東西不是那些兒女有多懂事.
是上帝在雅各身上做了什麼.
在以色列的身上.
在雅各的身上做了什麼.
上一代有一個很出名的精神科醫生.
叫做Frank Pittman.
Frank Pittman說過一句話.
他說教育兒女.
並非是一個完美男人要做的東西.
教育兒女.

$^{841}$卻是要去承傳一個男人的功績.
教養兒女最後的成果.
不是子女.
而是我們作為父母.
上帝要在我們身上要做什麼.
求主你不問我們.
我們為我們所有的父親去祈禱.
為我們的家長去禱告.
求主幫助我們給我們一個sense of calling.
我們只不過在上帝的僕人.
在我們有限的支連.
有限的vocation.
做好我們本分.
帶這個生命去到上帝面前.
同時在這個工程裡面求主是改變的.
最主要不是我們的子女.
改變的是我.
我們得到祈禱.
主要我們多謝你.
因為你的愛,你的恩典從來沒有離開過.
我們每一個為父的.
甚至為母的.
我們都是破碎的人.
但我們見到上帝你的恩典.
你的信神.
一點我們都不配.
求主你在我們生命當中.
去作你自己的聖聖的事.
使我們更加.
能夠去討你的喜悅.
成為你自己的勇人.
在我們這個vocation裡面.
盡忠職守.
遵行你的旨意.
祈禱奉尊名給我們.
\newpage



\section{啟示錄 6:12-8:5}
\label{sec:5Ocf3zsSFnw}
\textbf{亂世啟示:第六至第七印 (啟示錄 6\_12-8\_5) - 陳崇基牧師}
\newline
\newline
連結: \href{https://youtube.com/watch?v=5Ocf3zsSFnw}{\texttt{https://youtube.com/watch?v=5Ocf3zsSFnw}} ~~~~ 語音日期: 2022-06-24
\newline
\newline
\hyperref[sec:1_MORjvbY98]{\small{< < < PREV SERMON < < <}}
~
\hyperref[sec:index_chronic]{\small{[返順時目]}}
~
\hyperref[sec:index_scriptual]{\small{[返順卷目]}}
~
\hyperref[sec:X2LCQ0cGbnk]{\small{> > > NEXT SERMON > > >}}
\newline
\newline
啟示錄 6:12-8:5
\newline
\begin{longtable}{cl}
\hline
\hline
章節 & 經文 (和合本修訂版)\\
\hline
6:12 & \begin{tabularx}{0.7\textwidth}{X} 羔羊揭開第六個印的時候，我看見地大震動，太陽變黑像粗麻布，整個月亮變紅像血， \end{tabularx} \\ \\ \relax
6:13 & \begin{tabularx}{0.7\textwidth}{X} 天上的星辰墜落在地上，如同無花果樹被大風搖動，落下未熟的果子一樣。 \end{tabularx} \\ \\ \relax
6:14 & \begin{tabularx}{0.7\textwidth}{X} 天就裂開，好像書卷被捲起來；山嶺海島都被移動離開原位。 \end{tabularx} \\ \\ \relax
6:15 & \begin{tabularx}{0.7\textwidth}{X} 地上的君王、臣宰、將軍、富戶、壯士，和一切為奴的、自主的，都藏在山洞和巖石穴裡， \end{tabularx} \\ \\ \relax
6:16 & \begin{tabularx}{0.7\textwidth}{X} 向山和巖石說：「倒在我們身上吧！把我們藏起來，躲避坐寶座者的臉面和羔羊的憤怒； \end{tabularx} \\ \\ \relax
6:17 & \begin{tabularx}{0.7\textwidth}{X} 因為他們遭憤怒的大日子到了，誰能站得住呢？」 \end{tabularx} \\ \\ \relax
7:1 & \begin{tabularx}{0.7\textwidth}{X} 此後，我看見四位天使站在地的四角，執掌地上四方的風，使風不吹在地上、海上和各種樹上。 \end{tabularx} \\ \\ \relax
7:2 & \begin{tabularx}{0.7\textwidth}{X} 我又看見另有一位天使從日出之地上來，拿著永生神的印。他向那得到權柄能傷害地和海的四位天使大聲喊著， \end{tabularx} \\ \\ \relax
7:3 & \begin{tabularx}{0.7\textwidth}{X} 說：「你們不可傷害地、海和樹林，等我們在我們神眾僕人的額上蓋了印。」 \end{tabularx} \\ \\ \relax
7:4 & \begin{tabularx}{0.7\textwidth}{X} 我聽見以色列人各支派中受印的數目有十四萬四千； \end{tabularx} \\ \\ \relax
7:5 & \begin{tabularx}{0.7\textwidth}{X} 猶大支派中受印的有一萬二千；呂便支派中有一萬二千；迦得支派中有一萬二千； \end{tabularx} \\ \\ \relax
7:6 & \begin{tabularx}{0.7\textwidth}{X} 亞設支派中有一萬二千；拿弗他利支派中有一萬二千；瑪拿西支派中有一萬二千； \end{tabularx} \\ \\ \relax
7:7 & \begin{tabularx}{0.7\textwidth}{X} 西緬支派中有一萬二千；利未支派中有一萬二千；以薩迦支派中有一萬二千； \end{tabularx} \\ \\ \relax
7:8 & \begin{tabularx}{0.7\textwidth}{X} 西布倫支派中有一萬二千；約瑟支派中有一萬二千；便雅憫支派中受印的有一萬二千。 \end{tabularx} \\ \\ \relax
7:9 & \begin{tabularx}{0.7\textwidth}{X} 此後，我觀看，看見有許多人，沒有人能計算，是從各邦國、各支派、各民族、各語言來的，站在寶座和羔羊面前，身穿白衣，手拿棕樹枝， \end{tabularx} \\ \\ \relax
7:10 & \begin{tabularx}{0.7\textwidth}{X} 大聲喊著說：「願救恩歸於坐在寶座上我們的神，也歸於羔羊！」 \end{tabularx} \\ \\ \relax
7:11 & \begin{tabularx}{0.7\textwidth}{X} 眾天使都站在寶座和眾長老，以及四個活物的周圍，俯伏在寶座前，敬拜神， \end{tabularx} \\ \\ \relax
7:12 & \begin{tabularx}{0.7\textwidth}{X} 說：「阿們！頌讚、榮耀、智慧、感謝、尊貴、權能、力量都歸於我們的神，直到永永遠遠。阿們！」 \end{tabularx} \\ \\ \relax
7:13 & \begin{tabularx}{0.7\textwidth}{X} 長老中有一位回應我說：「這些穿白衣的是誰？是從哪裡來的？」 \end{tabularx} \\ \\ \relax
7:14 & \begin{tabularx}{0.7\textwidth}{X} 我對他說：「我主啊，你是知道的。」他向我說：「這些人是從大患難中出來的，他們曾用羔羊的血把衣裳洗得潔白。 \end{tabularx} \\ \\ \relax
7:15 & \begin{tabularx}{0.7\textwidth}{X} 所以，他們在神寶座前，晝夜在他殿中事奉他；那坐在寶座上的要用帳幕覆庇他們。 \end{tabularx} \\ \\ \relax
7:16 & \begin{tabularx}{0.7\textwidth}{X} 他們不再飢，不再渴；太陽必不傷害他們，任何炎熱也不傷害他們， \end{tabularx} \\ \\ \relax
7:17 & \begin{tabularx}{0.7\textwidth}{X} 因為寶座中的羔羊必牧養他們，領他們到生命水的泉源；神必擦去他們一切的眼淚。」 \end{tabularx} \\ \\ \relax
8:1 & \begin{tabularx}{0.7\textwidth}{X} 羔羊揭開第七個印的時候，天上寂靜約有半小時。 \end{tabularx} \\ \\ \relax
8:2 & \begin{tabularx}{0.7\textwidth}{X} 我看見那站在神面前的七位天使，有七枝號賜給他們。 \end{tabularx} \\ \\ \relax
8:3 & \begin{tabularx}{0.7\textwidth}{X} 另有一位天使拿著金香爐來，站在祭壇旁邊；有許多香賜給他，要和眾聖徒的祈禱一同獻在寶座前的金壇上。 \end{tabularx} \\ \\ \relax
8:4 & \begin{tabularx}{0.7\textwidth}{X} 那香的煙和眾聖徒的祈禱從天使的手中一同升到神面前。 \end{tabularx} \\ \\ \relax
8:5 & \begin{tabularx}{0.7\textwidth}{X} 天使拿著香爐，盛滿了壇上的火，倒在地上；就有雷轟、響聲、閃電、地震。 \end{tabularx} \\ \\
[1ex]
\hline
\hline
\end{longtable}
$^{1}$我們很光榮見到大家在現場.
能夠跟大家一起崇拜.
亦都在網上.
在芝加哥以外的地區歡迎你.
我們繼續我們講道的系列.
就是講忠末論和啟示錄.
講忠末論的時候.
我們講到一些必要性的東西.
不是一些想聽的東西.
為什麼呢.
就像剛才麗倩所分享的一些事項.
這些都不是我們希望.
願意或者想聽的東西.
這些不是我們想聽的東西.
但是忠末的東西.
聖經裡面.
我不知道大家有沒有想過.
其實大部分聖經裡面的東西.
都是一些不是我們想聽的東西.
但是我們必須要聽的東西.
是不是.
大家知道認識我的話.
你知道我一向睡覺都有問題.
我不是睡著我是沒有問題的.
但是睡了之後我會醒的.
那個是問題.
我睡兩小時.
兩小時鬆一點我又會醒.
然後又睡不著.
所以我掙扎了這個問題都挺久.
結果上次我看醫生的時候.
我就跟醫生說.
我有些這樣的問題.
他看到我檢查過之後.
他覺得你都很健康.
OK的.
好像沒有什麼了.
我就說不是.
其實我有一個問題.
我在掙扎的.

$^{41}$就是我發現.
如果我再年輕一點的時候.
我就OK的.
因為年輕一點.
好像橙一樣.
可以拿橙汁.
還可以滴出來.
現在年紀開始大了.
就發現不行了.
我要好好保養我自己.
所以我就跟醫生說.
這樣搞法不行.
這樣我會很快死的.
我真的這樣跟他說.
因為你知道美國醫生.
一定要跟他說得大聲一點.
他才聽得到.
結果他就說.
這樣嗎.
不如這樣吧.
我試過有些病人開過一種藥給他.
是有效的.
不過療癒的.
他問我試不試.
我說試一下.
結果他就開給我.
我就試了幾天之後.
我就習慣了.
我覺得是OK的.
起碼我可以睡到.
起碼五個小時六個小時左右.
就是連續這樣睡.
所以我就跟他說.
跟他解釋了之後.
他就開了藥給我.
他跟我說的時候.
我就發現他跟我說一些.
是我必須要聽的東西.
不是一些我想聽的東西.
他也忠告我.

$^{81}$牧師你也開始要留意.
不然你玩不了多少年.
是一些必要性.
需要聽的東西.
中末論.
Ascotom.
其實就是告訴我們.
我們需要去看些什麼.
是我們必須要看的.
我記得90年代的時候.
那時候我在溫哥華.
有機會在教會事奉.
因為有一個會友.
他的爸爸是微信主.
叫我們傳道人去探望他.
我們會友當然也叫他肯去給你探望.
我們就說好啊.
我就去探望他.
我探望他90年代末的時候.
在溫哥華.
你知道那裡的香港人很有錢.
是吧.
起碼溫哥華有兩三間屋.
香港可能有幾棟樓.
我跟他談.
我們弟兄姊妹跟他談.
他真是.
他口氣多大.
很難怪的.
他有錢又有權.
又有勢.
有地位.
他每天都是出去喝茶.
他為什麼要聽你的話.
我又明白.
我們這樣的窮傳道.
爛名一條.
你說什麼.
他的態度不是很好.
跟他談這些.

$^{121}$他說你教徒.
什麼叫感恩.
什麼叫奉獻.
我自己白手興家.
我用不著靠你.
最後不是很好收場.
最後探訪的時間.
我們說不妨礙你了.
不好意思.
可能你說的東西不是你喜歡聽的.
希望將來有機會再探望你.
他在那種境況裡面.
不想聽一些他必須要聽的東西.
我們每個人都是脆弱的.
我們總會面對.
中末的問題.
中末論不是一些抽離的東西.
他告訴我們.
有一些我們必須要聽.
一定要聽.
一定要打大耳朵.
睜大眼睛去看.
想清楚.
到底這件事跟我們有什麼關係.
我們看過啟示錄.
一張又一張的.
告訴我們.
到底中末會是什麼事.
我們看七印.
七印這裡說到.
所謂中末的徵兆.
signs of times.
就是現在我們所活著的時間.
就是一個last chapter of history.
有戰禍.
有經濟的動盪.
有瘟疫.
有迫害.
對教會的迫害.
有天地的徵兆等等.

$^{161}$然後有一段間奏interlude.
十四萬四千人出現.
然後第七印就出現了.
這裡七印很特別.
就是只剩下半個小時.
在天上.
我們會繼續看下去.
但是你看到有一個pattern.
在任何聖經裡面.
特別是新約.
說到中末的時候.
都會有一個這樣的形式.
就是一些危機出現.
然後危機就會去到一個高潮.
那個高潮.
就會是一些.
可能天上的一些兆頭.
就會出現.
是一些我們人不能夠去控制的東西出現.
然後最後上帝的審判就會來到.
我們看過馬達方案第十四章.
Oliver Discourse.
耶穌說到.
他回來的日子會是怎樣.
好像是主後六十年代尾.
七十年的時候耶路撒冷被毀的時候.
那個情況.
一樣都是說這個東西.
民要攻打民.
國要攻打國.
好像第一義人那樣.
然後有其他的經濟的動盪.
有饑荒.
然後有人出賣信仰.
或者很多人在艱難的時候.
他就離開主等等.
迫害的時候.
然後也有天兆出現.
然後主說他就要隨時回來.
同一個Panel.

$^{201}$甚至整本啟示錄的結構.
我想告訴你.
是應該這樣看的.
七人其實說的話.
和七記號.
七碗.
七號.
說的話其實都是一樣.
是同一個主題.
就是上帝終末的時候.
要去審判活人死人.
這件事一定會發生.
那件事會是怎樣呢.
七印就在打開的時候.
就告訴我們.
世界的終末.
你的終末會是怎樣.
然後七號又再說同一個主題.
不過是再激進.
或者再緊張一點.
這樣告訴你.
然後七號七碗.
都是在說一個Cycle.
就好像你聽一首交響樂曲.
不斷的有一個Theme.
一邊Play一邊Play的時候.
就告訴你.
這件事是多麼的重要.
所以我們由第一印.
說到現在.
我們就去到第六印.
第六印說什麼呢.
六章第十二節.
他說高揚.
這個被殺復活的高揚耶穌基督.
揭開第六印的時候.
他就看見地大震動.
太陽就變黑.
然後月亮變紅像血.
天上的星辰墜落在地上.

$^{241}$好像無花果樹被大風吹搖一樣.
落在未熟的果子一樣.
天就裂開.
像書卷被捲起來.
山嶺海島都被移動.
離開本位.
地上的君王.
臣宰.
將軍.
富戶.
壯丁.
一切圍爐的.
自主的.
都藏在山洞和岩石的懸崖.
的月石月裡面.
向山和岩石說.
倒在我們身上吧.
把我們藏起來.
躲避.
助補助者的念念.
和高揚的憤怒.
因為他們遭憤怒的大日子到了.
誰能暫立得住呢.
在這裡說了一大堆的東西.
我們可以叫做天兆.
即是一些天上的天兆.
天上可以看到的日蝕月蝕等等.
對古代甚至對我們今天的人來說.
都是一些我們不能控制的東西.
特別對古代的人來說.
是一些特別在不利的時候發生.
不適合那個日子.
有月蝕.
或者他不是預計得到那個日子.
到了有月蝕的時候.
對他來說是一個空兆.
這裡就用一些這樣的語言.
或者這樣的意象.
來說上帝的審判.
他的公義的審判.

$^{281}$必定會到.
還有一些什麼的預兆.
我們再說多一次.
這些不是預先寫好的天月報告.
這些在第一世紀已經有發生.
第一世紀文獻有記載.
特別是耶路撒冷70年.
被滅之前已經出現.
並且每一代的信徒.
甚至我們今天這一代的信徒.
我們都會看到這樣的signs.
如果你看看舊約的話.
舊約都已經有這樣的一個pattern.
譬如在以賽亞書第十三章.
上下文說什麼.
說巴比倫帝國要傾覆.
以賽亞先知預言將來.
巴比倫打壓上帝子民的角度.
必定會受審判.
那個情況會是怎樣.
你看到是差不多類似的情況.
十三章第十節.
天上的星獸都不發光.
太陽一升起就變黑.
月亮就不放光.
日蝕月蝕.
然後必定要有.
我必因邪惡.
即是巴比倫帝國的邪惡.
來審判這個世界.
因罪孽來懲罰這些凶惡人.
止息狂傲人的狂妄.
禁止殘暴者的傲慢.
巴比倫為列國的榮耀.
迦勒底人所誇耀的華美.
必像上帝所傾覆過的.
索多瑪奧摩賴.
什麼時候發生的.
巴比倫帝國.
覆亡的時候.

$^{321}$一夜之間.
主前大概五百多年的時候.
歷史有講.
我們也知道.
聖經也是這樣預言的.
所以這裡告訴我們.
這種模式.
當你看到這些徵兆.
你不需要再想.
到底有多近.
主的日子就一定會來.
並且是隨時來的.
我們回到第六人.
這裡說什麼.
山洞和岩石縣.
這些也是很本土的描述.
當人們有橫的時候.
當時巴勒斯坦地的人.
一定會走去哪裡.
就是走去山的洞穴裡.
在那裡躲起來.
然後說到什麼人.
七個階層的人.
七個字眼.
地上的君王.
這個字很明顯就是代表.
羅馬帝國裡面的蓋薩皇帝.
和他的貴族.
臣宰.
在羅馬帝國底下.
有些所謂臣宰.
那些叫Klein Kings.
大希律就是其中一個.
是代他管.
好像特首那樣.
管一個民族.
或者管一個地方.
然後在他下面.
就有將軍.
軍長.

$^{361}$有庫戶.
有壯丁.
特別是每一個羅馬城裡面.
有錢的公民.
我們會叫他.
新約的時候.
會叫他壯丁.
一切為奴的.
自主的.
上帝第一次來.
是以耶穌基督.
十字架的方式來.
是以恩典來.
但如果講到上帝再來的時候.
耶穌基督再來.
就會是公義的審判.
就是全天下的人都逃不了.
全天下的人都不能夠避免.
有錢的.
沒錢的.
在上掌權的.
還有甚至是做奴隸的.
沒有一個人能夠逃避.
你和我都不能夠逃避.
講到終末的時候.
就講到甚至一些勇者.
壯丁.
軍長.
有權有勢的.
都逃避不了.
我講回剛才那個九年代的時候.
在加拿大溫哥華.
探望那個朋友.
記得我第一次探望他.
結束之後我說什麼.
我希望再有時間探望你.
結果真的有時間探望.
但我再探望他的時候是ICU.
我去到ICU的時候.
潘若良人.

$^{401}$當然我都是蒙恩德救的人.
我不能夠說死了你.
當然不是.
我都是蒙恩德救.
我是算什麼.
我都是爛命一條.
但我很強調的一件事就是.
都是潘若良人.
當時風光的時候.
他可以不聽一些他必須要聽的東西.
到他ICU的時候.
我們一班大英哲會去探望他.
我們照樣講同一樣東西.
我們很脆弱.
我們需要上帝救恩.
當我們面對終末的時候.
你可以暫時不想.
但你不能逃避.
你儘管遲延你要考慮的東西.
但你不能夠逃避.
其實你想想.
福音起碼是一回事.
信耶穌.
其實就是我們放下身段.
我知道我自己不行.
我知道我是脆弱的.
他肯去聽.
他在病床上接受了.
這個上帝救恩.
上帝的恩典.
我們每一個都需要.
終末就是我們必須要知的東西.
必須要去面對的東西.
不是你想面對的東西.
你和我都不想.
但你一定要.
在第六印講到上帝的審判終末來到的時候.
一個印一個印打開.
就告訴我們我們需要悔改去歸向上帝.
第六印之後.

$^{441}$你記得這裡第六印問一個問題.
誰能夠暫立得住.
誰能夠逃避.
誰可以逃避.
基督再來公義審判的大子.
誰能夠暫立得穩.
第六印打開之後.
就是一個所謂的間奏.
我們看第七章第一節.
此後我見到四位的天使站在地的四角.
即是全地.
執掌地上四方的風.
使風不吹在地上.
風這裡啟示錄.
或者天啟文學應該是有傷害的意思.
叫他不要吹在地上海上和各種樹上.
所以沒有風了.
我又看見另一個天使從日出之地東邊出來.
他拿著永生上帝的人.
他把那得到權柄能傷害地和海的四位天使.
大聲喊著說.
你們不可以傷害地海和樹林.
在我們上帝中僕人的額上蓋了印.
我聽見以色列各支派中受印的數目有十四萬四千.
這裡很明顯是在說一個間奏.
即是中國還沒到之前.
就是天使出來.
在上帝的僕人的額上蓋印.
給他們一個mark.
有多少人呢.
有十四萬四千人.
這裡說到mark of god.
即是上帝給他的僕人的一個印記.
我們說啟示錄.
我知道有很多人亂解.
這裡順便提一提.
有些人說那個受的印記.
十三章十六節.
他就是受的印記.
受即是羅馬帝國.

$^{481}$代表撒旦勢力的一個權勢.
有使眾人無論大小的自主的為奴的.
在他右手或在額上打了一個印記.
這樣除了那有印記的.
有受的數目的.
即是羅羅羅.
就不能夠做買賣了.
順帶提一提.
這個受的印記不是我聽過有人說是打疫苗.
因為打在手上.
所以.
I don't know.
這樣也有的.
但你見到是斷章取義.
因為印記打在哪裡.
手 或 額.
有沒有看漏.
額.
有沒有人打疫苗打在頭上.
I don't get it.
Anyways.
很明顯.
第六印之後的Interlude.
說14萬4千人是受印.
這個和受的印記是相反的.
即是14萬4千.
說到的就是上帝的印記印在神的子民身上.
說起14萬4千.
這個英文很出名.
因為耶和華見證人說14萬4千是實數.
所以真正得救的人只是14萬4千.
不會是14萬4千零一.
一定是14萬4千.
如果你這樣解釋天啟文學一定是撞板.
因為不是說實數.
Please.
不要當我們現代人的文學去看.
OK.
你看這些天啟文學.
不要當一個會計師去看.

$^{521}$因為不是實數.
不是說平衡.
是說meaning.
symbolic meaning.
14萬4千.
很肯定不是實數是代表什麼呢.
為什麼不是實數.
首先第一.
到第一世紀的時候.
你們知道.
到第一世紀不再有十二支牌.
因為北國的十支牌一早已經滅亡了.
The Lost Ten Tribes.
已經再找不到了.
雖然今天有很多story說怎樣找回那個.
No way.
The Lost Ten Tribes.
剩下的是猶大國.
或者他包著的西元.
剩下的是Judah.
一直被擄回歸的人當中.
剩下的只有這麼多.
這些所謂漁民.
Remnant.
剩下來的神的子民.
十個支牌已經告吹了.
一早在亞述帝國的時候.
只剩下南國猶大支牌.
但是.
在舊約的先知書裡面.
雖然十個支牌不變是Lost Ten Tribes.
但經常都說到將來有一個應許.
就是上帝一定會復興十二支牌.
十二支牌從某一個角度來說.
就會被復興起來.
什麼時候呢.
就是舊恩來到的時候.
利塞亞來到的時候.
所以記得耶穌來到的時候.
有多少門徒.

$^{561}$十二個是嗎.
是不是我們小組教會那種方程式.
一早有十二個就是最perfect的number.
因為一個組長take care十二個人是ok的.
十三個就不行了.
是不是這樣計算.
不是的.
為什麼一定要十二.
因為在講以賽亞書.
或者其他先知書裡面所講的利塞亞的盼望.
就是他來到的時候.
舊恩來到的時候.
就會復興十二支牌.
上帝講過的東西一定會算數.
耶穌來到有十二個門徒.
初期教會.
少林寺第一章你記得.
猶大出賣主.
自己失去了他十二個應該有的位置.
第一章發生什麼事.
就一定要算回第十二個.
馬提亞.
馬提亞在少林寺轉位.
在新約的歷史裡面.
我叫他做蛋散.
又沒有講馬提亞有什麼故事.
有什麼威水史.
沒有的.
為什麼一定要找他出來.
就是因為十二.
就因為number是十二.
特別在路加.
到史唐傳.
就是在講.
上帝的應許一定會實現.
利塞亞來到的時候.
他一定會復興真正上帝的子民.
就是有十二.
湊夠十二.
因為上帝的計劃.

$^{601}$要湊夠十二才可以再開始.
所以史唐傳第一章.
就要交代一些我們覺得.
跟我們不相大雅的東西.
但是在舊恩的歷史.
上帝的應許裡面.
就是非常之重要.
所以一定要湊回十二.
十四萬四千.
就是十二乘十二乘一千.
啟示錄裡面.
二十四位像.
十二個舊約.
十二個新約.
代表整個上帝的子民.
是蒙受上帝的恩典.
才能夠拿到這個恩典的人.
才能夠在上帝公義的審判.
來到的時候.
能夠站立.
這個恩又不是一個物質的恩.
是一個靈性的孝.
是一個聖靈的蓋恩.
你得到舊恩的時候.
你的蓋恩.
問題就是這樣.
你是不是真的那一位.
在十四萬四千裡面的其中一位.
就是上帝臨到的時候.
你已經是預備好的.
上帝臨到的時候就是第七印.
第八章第一節.
高揚見開第七印的時候.
天上寂靜約半小時.
安見那站在上帝面前的七位天使.
有七支的號角賜給他.
另一個天使就拿著金香爐.
這個金香爐代表什麼呢.
就是眾聖徒祈禱.
獻在寶座面前的金壇上.

$^{641}$拿香的煙.
和眾聖徒的祈禱.
從天使的手中一同升到上帝的面前.
天使拿著香爐.
然後炭上的火.
倒在地上.
天上的景況就是這樣.
然後將這些香炭放在地上.
然後就有不同的天兆和閃電等等.
就是上帝臨到的時候.
你記得西大山也是這樣的情況.
你記得我們講第五印的時候.
就說曇下有殉道者的靈魂.
他們求上帝來申冤.
這個禱告.
到第七印我們發現.
就出現了.
就是上帝必定會按照殉道者的祈求.
他會為他們申冤.
上帝的公義一定會臨到.
主必快來.
如果我們看舊約歷史.
就是教導我們這件事.
你今天未必可以參透.
縱帝國或者地上的權勢.
改朝換代的事情.
但是在舊約裡面.
已經給我們一個歷史.
讓我們能夠看到這個圖案.
邪惡的國度.
邪惡的勢力.
代表撒旦的勢力.
亞述國曾經張牙舞爪.
最後都是覆亡.
按照先知所說.
巴比倫剛才提了.
其他的國度都是一樣.
現在在第一世紀.
第一代的讀者.
看見啟示錄的時候.

$^{681}$他們當時是受到馬帝國的迫害.
受到他們的打壓.
這裡就告訴他們.
上帝是聽到他們的禱告.
上帝是聽到我們今天.
去呼喚他行公義.
去出手這個禱告.
不一定是按照我們想要的時間表進行.
但是上帝說.
必定按照他的議程來進行.
上帝的公然的審判.
必定會來到.
天上寂靜半小時.
第七天.
什麼意思.
我想兩約中間有一個很好的解釋.
就是天上有停頓.
寂靜一段時間.
兩約中間有一些猶太的文獻.
說到上帝工作之前.
是有沉默的.
等於軍隊出戰之前.
都會有一段沉默的時候.
寂靜的時候.
一段樂曲突然間停頓.
其實期待的是什麼.
就是下一個音符出現的時候.
那個音符在末世論中末論.
就是一個高潮.
就是末日來到之前.
天地會有寂靜.
因為上帝將會工作.
如果這樣看的話.
十四萬四千就不是一個間奏.
因為在第七天之前.
其實是一個呼召.
你是不是十四萬四千裡面的其中一位.
那你說.
那就是信耶穌就行了.
就是等上天堂.

$^{721}$似乎又不是.
如果你看十四萬四千的記載.
那裡說到猶太之派.
什麼之派.
有多少個國家加起來就是十四萬四千.
那裡說到長老中有一位回應我.
他說你見不見到這十四萬四千是穿著白衣的.
他們是從哪裡來的.
約翰就說.
我主日你是知道的.
然後這位長老就回答他.
這些人是從大患難中出來.
他們曾用羔羊的血.
把衣裳洗得潔白.
講到舊恩.
恩憐.
但如果留意.
舊約兩約中間.
怎樣去形容洗衣裳這件事.
洗衣服這件事.
是什麼時候洗衣服的.
是上帝的軍隊打完仗.
回來迎宗的時候.
他們就需要自結.
要把衣裳洗好.
這裡形容的是上帝的軍隊.
得勝再回來的時候.
這一班人就是十四萬四千.
那些忠心為主打仗.
不會因為世界威權去威逼利誘而放棄的.
至死忠心的那班人.
為主作戰的那班人.
那班人就是上帝的子民.
十四萬四千.
是因為上帝的道.
他們所做的見證.
他們回來上帝的迎宗.
他們就等候什麼.
天上的正正半小時.
第七天.

$^{761}$上帝的公然就臨到.
第七天打開.
就是七號.
第七號打開就是七碗.
第七碗打開.
最後就是最後的陣戰.
然後新天新地才會到.
十四萬四千.
就是預備好.
忠於主命的.
肯說和撒旦去交戰.
藉著什麼方法交戰呢.
藉著上帝的道.
藉著我們的見證.
去和撒旦交戰.
我們看最後一段聖經.
就是在十二章.
這裡雖然不是說十四萬四千.
但都是異谷同工.
說到教會和撒旦的爭戰.
上帝的子民和撒旦的爭戰.
不過聽見在天上有大聲音.
我臣的救恩.
他的國度.
所立的基督的權柄.
現在來到了.
因為那個在我們上帝面前.
晝夜空攻我們的弟兄的.
這位撒旦.
已經被率下.
弟兄勝過那條龍.
撒旦.
此恩膏揚的權.
是因自己所見證的道.
雖然至於死.
他們也不惜自己的性命.
所以諸天和住在其中的.
我們都快樂的去讚美.
地和海又和了.
因為上帝的公義已經臨近.

$^{801}$並且魔鬼知道自己的日子.
事後不多.
就氣憤憤的去到地上打退.
這裡說到.
這些弟兄姊妹.
是因著高揚的血.
得到見證.
並且因著自己所見證的道.
不惜一切的去跟隨著.
那些就是願意與撒旦爭戰.
十四萬四千.
神的子民.
我知道我們說到殉道的時候.
可能我們警方.
我們北美.
或者不同的地方.
比較自由的社會.
不在一些創啟地區的信徒教會.
我們可能覺得有些利地.
什麼殉道.
我等一下就出去吃飯.
殉什麼道.
下個星期我又來教會.
可能是我們身處的環境.
我們不知道.
但你知道第一世紀是這樣.
他們是躲在聚會裡.
唱歌也不能大聲一點.
對他們來說當然是一種安慰.
但我相信這裡不是只說利地的事.
什麼叫殉道.
什麼叫十四萬四千.
就是你忠於去跟著時刻的準備.
在我們這個自由社會裡.
豈不是我們更加需要付出更多.
我們說中國大陸的教會被迫害很慘.
他們很忠心.
Yes 他們是.
我們在自由社會.
我們豈不是要付出更多.

$^{841}$我女兒最近預備考試.
因為她已經大學第四年.
她快要考一個.
我們伊利諾州的教師的牌照的試.
原來聽她說.
全美國這麼多州裡面的小學教師.
考一些執照的試.
原來我們伊利諾州是最難的.
居然有這樣的天方夜譚.
教小學.
她給我看一些題目.
都很難 我也回答不到.
因為她是什麼都會.
什麼都要懂.
我懂的那一塊就是幫到她的.
古代歷史你給我吧.
我告訴你什麼事.
真的很難.
但我從來沒有見過她在四年讀書.
是這麼不鬆懈.
因為她就是要預備大考.
那種心情.
是吧.
時刻要去準備.
我們作為信徒.
是不是又鬆懈了.
你想想你.
你是不是放鬆了.
是不是很開心.
啟示錄就是告訴我們.
面對終末的時候請你準備.
不能鬆懈.
因為魔鬼的時間知道自己少.
如何在我們生活裡面.
付出更多.
比創啟地區的弟兄姊妹付出更多.
來見精神.
最近這兩個月.
我處理一個夫婦關係的問題.
花了我很多時間.

$^{881}$兩夫婦不明白為什麼.
不知道是不是疫情搞得大家很煩躁.
幾乎有家暴.
見了他們好幾次.
一開始我都沒有信心.
我唯獨是祈禱.
見了一次之後.
兩次.
終於大家開始有一點共識.
原來大家的溝通非常有問題.
兩夫婦在家裡.
一個講中文.
另一個講法文.
兩者完全溝通不到.
很多事情有很多誤會.
很多事情就吵起來.
花了很多時間.
最後講的都是一些很淺白的東西.
有多難.
不用讀音樂.
就是在講基督教中間.
我們夫婦越近上帝.
放下我們身段.
以主為我們的主.
不是再爭誰贏誰輸.
其實很簡單.
你也會說.
只是我比較專業的意思是.
他們在吵架的時候.
我又不會和他們在一起.
我只會調停他們.
但都是一些很淺白的道理.
以主為我們的主.
這就是聖經中的核心.
告訴我們一些中末的事情.
請你想清楚.
如果主是我們的主.
我們想清楚我們的中末的時候.
現在我們有的問題.
就可以放下.

$^{921}$有很多爭拗我們就可以放下.
因為我們最終要誰贏.
基督要贏.
不是我要贏.
也不是他要贏.
是我們兩個都放下.
然後我們說基督你要贏.
請你來到我們中間.
來到我們作主.
生命其實不是這樣.
我們想起中末.
我們作為門徒.
不再鬆懈.
以主為我們的中心.
最後我想邀請你.
做一個這樣的祈禱.
無論你面對著什麼處境.
什麼將來.
有什麼決定你要做.
有什麼問題你正在面對.
你跟上帝說.
主我要你的心意在我生命裡面.
象現出來.
給我一個僕人的取態.
去聽你的話語.
去遵主你的命.
我要主在我身上能夠得勝.
你立刻這樣祈禱.
這些祈禱很靈驗.
你平時那些清單的祈禱.
通常都不靈驗.
因為那些是為我自己得勝而禱告.
為主得勝而禱告.
你試試這樣禱告.
看你生命會怎樣.
讓羔羊在你生命裡面得勝.
我們同心得悠祈禱.
主愛多謝你.
因為你對我們這群罪人的恩典.
你對我們這群這麼破碎的人.

$^{961}$你願意羔羊的血來洗淨我們.
羔羊的十家去求恩.
在我們生命裡面主愛求你.
讓主的道能夠在我們生命當中得勝.
主愛求你.
讓我們能夠成為你的僕人.
我們這麼祈禱 放主美金.
\newpage



\section{創世記}
\label{sec:X2LCQ0cGbnk}
\textbf{創世記第6章 (6.6) - 背景釋經 [方舟,再創造 1] (陳崇基牧師)}
\newline
\newline
連結: \href{https://youtube.com/watch?v=X2LCQ0cGbnk}{\texttt{https://youtube.com/watch?v=X2LCQ0cGbnk}} ~~~~ 語音日期: 2022-06-24
\newline
\newline
\hyperref[sec:5Ocf3zsSFnw]{\small{< < < PREV SERMON < < <}}
~
\hyperref[sec:index_chronic]{\small{[返順時目]}}
~
\hyperref[sec:index_scriptual]{\small{[返順卷目]}}
~
\hyperref[sec:zuvpXyLO2aE]{\small{> > > NEXT SERMON > > >}}
\newline
\newline
$^{1}$我們上幾次有講過聖經考古的背景.
或者一些比對研究的歷史.
現在我們可以回到聖經.
就是第六章羅亞方舟或羅亞洪水的故事.
這段聖經是很大的.
六章至到第九章才完結.
我們經常都會說.
創世記1至11章就是一個可以在內容上分出來的單元.
我們叫它做史前史或者起源史.
英文叫做Proto History或者Primeval History.
等等不同的字眼都是告訴我們這些是很遠古.
和起源有關的記載.
這些經文所講的東西都是和世界的開始或世界本身的秩序有關.
我們也說過不能夠當它是一個現代歷史的寫作.
也不能夠說它是一個科學的報告.
我們要去到明白這些經文.
1至11章特別.
這些所謂的Proto History 起源史.
是在近東的背景寫出來的.
我們已經說了很多次.
由第一章開始我們去解讀創世記六日七日創造的時候.
都已經不特別強調這件事.
不是一個現代的科學報告.
不會解答我們所有現代的問題.
記得嗎?.
是在一個近東的背景底下寫的.
而他們所關心的東西往往不是我們現代人所關心的.
所以我們要留意這個角度.
所以為什麼我們好像分差了.
去到講一些比對研究等等.
這些近東背景的資料.
為什麼不可以只看聖經呢?.
因為聖經不是直接寫給我們.
而是為著我們每一代的信徒.
包括我們這一代在生的信徒.
為我們而寫的.
兩個的分野很重要.
因為聖經不是直接寫給我們.
語言是希伯來文.
背景概念都不是我們現代人的背景和概念.

$^{41}$宇宙的觀念.
就是宇宙的構造等等.
都不是今天我們科學人士了解的.
而是近東的人所了解的.
所以Bible is written to them.
not to us.
而是written for us.
for us的意思是.
在近東的背景底下.
上帝啟示他的心意和教導.
是藉著近東的文化背景去彰顯出來的.
所以對第一代的讀者來說.
他們有那種普遍的概念.
或者那種文化的背景.
宇宙觀的背景等等.
他們就明白聖經的意思了.
他們是初代的讀者.
而上帝藉著他們有的.
普遍的一些觀念.
他們的語言.
他們的手法.
來啟示出他們自己的真理的時候.
那個真理我們要先探索了.
然後那個真理才會是.
超越所有世代.
超越時空超越地理的.
正如我們看第一章也是一樣.
看第一章我們不要當它是一個科學報告.
因為這樣就很倫潤.
引發很多不需要有的問題.
因為他主要想說的是who.
不是how.
不是這個宇宙怎樣變成.
這些是我們現代人經常會想的問題.
經常會問的問題.
他是說who.
就是說這個宇宙.
到底誰創造.
誰值得敬拜.
所以是一種神學的宣告.

$^{81}$這個誰是創造主.
誰值得敬拜.
who這個就是超越所有的世代.
所以我們今天看創世記第一章的時候.
我們也了解.
原來創世記第一章就是介紹創造主.
不是說創造界如何形成.
有些什麼規律.
或者有些什麼物理的原則.
能夠去造成它.
所以我們要了解這方面.
我們看每一章聖經其實都是一樣.
今天我們繼續看第六章.
第六章我們說到洪水的故事.
之前交代過了.
神的意志半變.
無論你是天使的解說.
或者是地上的暴君的解說也好.
總之那個的原因.
有洪水來毀滅當時已知的世界.
當時的No World.
已經宣佈了原因.
然後就去到講解.
或者啟示出一個故事.
這個羅亞方舟.
羅亞洪水的故事.
我們先看經文.
六章十四節.
耶和華神對羅亞說.
你要為自己用戈飛木造一艘方舟.
並在方舟裡面建房間.
內外要圖上匿稱.
方舟的造法是這樣的.
長三十爪.
闊三十爪.
高三十爪.
方舟上要造天窗.
再上是上一爪.
方舟的門是開在旁邊.
分上中下三層.

$^{121}$有什麼特別.
特別的地方是.
如果我們用一個現代人的看法.
我們很多時候就是去關注一些問題.
水力問題.
一些物流的問題.
有沒有可能全世界的動物都走進去.
這個方舟夠不夠大等等.
我們就引申了很多物流的問題.
到底安排會怎樣.
管理會怎樣.
有沒有可能等等的事情.
其實不需要.
因為古代的人.
近東的人寫洪水的故事.
聖經不是第一個.
也不是唯一一個.
也有很多近東的故事.
都是講到洪水的故事.
我們之前已經交代過.
其中一個最出名的就是.
Gilgamesh Epic Tablet 11.
第十一塊的泥板.
這個基加美斯.
所以我們可以用古代的方法去看.
因為他所記載的東西.
是有他的意思.
或者一些他想表達.
想強調的東西.
那些往往不是我們有興趣的.
但是我們需要給予經文去控制.
給予聖經他要講的東西.
去控制我們應該看什麼.
和應該問什麼問題.
不應該問什麼問題.
好了,這裡講到方舟.
方舟這個字.
在希伯來文是Tiva.
Tiva是什麼意思呢.
什麼叫Tiva呢.

$^{161}$其實這個字是一個埃及字.
埃及的借用字詞.
我不知道有沒有提過.
在摩西五經裡面.
有很多是埃及的借用詞.
特別在五經.
後期的希伯來文的經文.
就越來越少這些所謂的Egyptian loan words.
借用的字眼.
Tiva就是其中一個埃及的文字.
埃及文字.
是借用字詞.
什麼叫借用字詞呢.
比如說我是廣東人.
很清楚以前在香港長大的時候.
有士多這個字.
store.
其實士多是什麼呢.
就是雜貨店.
其實是借用字詞.
士多的中文不會有士多這個字.
只不過是英文的store.
convenience store.
音譯過來.
所以士多這個字就是所謂借用的字眼.
就是外語借用的字眼.
在五經裡面特別多這些埃及近東的借用字眼.
正如我所說後期的比較少.
如果你有機會看到出埃及記的時候.
就有越來越多這些借用的字眼.
這個就反映聖經寫作的背後歷史性.
他是在那段時間去寫的.
所以特別你去比對其他文獻的時候.
你就發現有些希伯來文原來本身是埃及文.
只不過是音譯過來.
正如我說的士多.
這個teva這個字在埃及的一些文獻出現.
通常是說一個盒子或者棺材.
或者一些四四方方的東西.
就借用這個字.

$^{201}$在舊約裡面出現了28次.
有兩次就出現出埃及的二章三或者五節.
那裡說到teva是做什麼的.
就是摩西在尼羅河.
法老要去殺害一些男嬰.
然後他的媽媽就將這個嬰孩放在一個teva.
一個籃子裡面套上拿清.
然後放他到河裡.
然後有很多上帝背後幕後做很多事情.
然後留下了這個所謂拯救的英雄.
這個國家的英雄就是摩西.
所以就用了teva這個字.
所以這裡說到洪水的故事的時候.
有心就用這個字.
26次在創世記的洪水的故事裡面去使用.
有兩次總共在舊約裡面28次.
創世記的故事裡面有26次之多.
很明顯就是那個作者.
他就是想提起這個羅亞方舟的故事.
是一個拯救.
一個救贖.
另外除了看teva這個字之外.
我們就看看到底怎樣去做.
怎樣做.
他說用哥菲木去做.
gover.
你看看不同的翻譯就知道.
我們不是很清楚什麼叫做哥菲.
中文譯做哥菲木.
什麼叫哥菲木我們根本就不知道.
只不過是音譯叫做gover.
我們就譯做哥菲木.
原因是什麼呢.
因為這個gover這個字.
在舊約裡面只出現了一次.
就是這裡.
我們看看英文怎樣翻譯.
NIV NLT.
New International Version.
New Living Translation.

$^{241}$就譯做cypress wood.
柏木.
這個肯定不是柏木.
因為柏木有其他字.
可以叫做柏木.
cypress wood是怎樣譯回來的呢.
就是譯者.
英文聖經的譯者.
他是按照上下文推理出來.
他覺得是一個船.
應該用一些硬的木.
所以就取了cypress wood.
其實NIV的note已經說了.
meaning uncertain.
就是說我們不知道那個意義.
最近代的學者甚至說是草.
哥菲其實不是哥菲木.
是哥菲的植物或者是草.
你再看這個新納斯本.
他用了一個希臘文.
兩個希臘文.
就去說哥菲木.
就是square timber.
正方的木.
這個又幫不了我們很多.
因為什麼叫正方的木呢.
正方的木可能又是按照上下文.
好像cypress wood那樣去推理.
覺得是用正方的木材.
去做這個方舟.
所以這些都幫不了我們.
也告訴我們.
似乎經文著眼的不是說一艘船.
因為在舊約的希伯來文.
是有船這個字.
or ani.
有船隻這個字.
如果他真的想描述一艘船的話.
他大可以用那個字.
ani或者aniuk.

$^{281}$就不需要用teva.
所以我們首先看一看這段經文的時候.
暫停下來的時候.
我們就想清楚.
我們平時構想的方舟.
就是一艘船.
甚至美國有一些很保守的信徒.
做了一個公園.
叫做方舟公園.
說是聖經的一比一.
去做出來的.
我們就加上很多的想像.
手船前面應該是尖尖的.
因為我們想的船就是尖尖的.
如果聖經作者真的想說一艘船.
他就可以不用teva.
或者可以用teva再加上ani或者aniuk.
這些字.
再加上這些木.
又不知道是甚麼木或甚麼植物.
所以他的著眼的地方.
不是我們今天所留意的.
我們可以試一下用一些比對的材料.
讓我們知道到底聖經的作者想做些甚麼.
或者他可能在回應當時的其他.
所謂一些洪水的故事的記載.
他去到做爭辯去到反駁等等.
可能是在做這件事.
正如我之前說過911這件事.
911發生之後.
有很多的陰謀論出現等等.
那些陰謀論很多時候都是錯的.
危言聳聽的.
然後最後可能有些作者.
有些歷史研究政治的學者.
就會寫一本書.
去到說911到底本身發生了甚麼事.
去到以正視聽.
有些人亂說話.
他就說出來.

$^{321}$所以方舟的故事.
你可以這樣角度去看.
正如我所說不是寫給我們.
直接寫給當時近東的人.
當時街上或者社會裡面.
可能流傳不同的洪水故事.
包括Gilgamesh Epic.
羅亞方舟這個故事.
創世記這個作者.
他想寫一些東西來回應.
有甚麼近東的材料可以給我們.
或者可能去做一些對比.
到底這個方舟是甚麼一回事.
我們看看近東的一些其他洪水故事.
說到也有方舟.
我們做方舟的記載.
在Gilgamesh的史詩.
記載方舟的尺碼.
神靈和洪水的英雄.
叫做Yunapesh.
跟他說要做一隻方舟.
這隻方舟是120爪乘120爪乘120爪.
或者120爪乘120爪正方形.
有七層.
等於像這個構圖一樣.
雖然沒有七層.
為甚麼要這樣形容呢.
這是近東背景如何幫助我們.
這七層的東西.
一層一層或者一個Cube.
120乘120乘120.
這裡很明顯.
在近東的宗教神話裡.
有這個叫做Sigure.
廟塔這回事.
之前我們也有提過.
在創世記第十一章一至九節.
提到巴別塔的時候.
巴別塔不是後代歐洲人所想像的故事.
其實不是歐洲的建築.

$^{361}$是巴比倫的建築.
是米索不達米亞的建築.
這些廟塔很多時候.
其實是一個樓梯.
很多時候形容它是塔頂通天.
最上面就是一個聖所.
一個Temple.
方便天上的神靈下來.
然後廟塔旁邊就會有一些廟宇.
人們就會在這裡供奉.
方便神靈下來接受供奉.
所以整個方舟.
在Gilgamesh Epic裡面說.
這個方舟就是要這樣的尺碼.
方舟就像一個廟塔.
這裡我有一個地圖.
就是在這本書很舊的.
1980年的時候寫的.
1990年的時候出版的地圖集.
Cultural Atlas of Mesopotamia and Ancient Near East.
Michael Routh出版了這個地圖.
我覺得這本地圖是很好的.
雖然已經絕版了.
但這個地圖有很多圖表.
幫我們了解近東的背景.
有一個構圖說到找到的廟塔.
在近東的兩河流域.
找到不同類型的廟塔.
也有不同的形狀.
所謂的Ziggurat.
不是金字塔.
金字塔是埃及的.
埃及的金字塔是墳墓.
即是皇陵.
所以裡面是安裝了法老.
有不同的通道和空間.
這些廟塔是實心的.
製作了一個框之後.
他就把泥和雜物放進去.
然後把它裝好.

$^{401}$一直這樣建上去.
所以裡面是沒有東西的.
正如我所說.
最重要的是上面.
所以這些廟塔就像樓梯一樣.
是方便把光放下來.
來接受人的敬拜.
這裡的地圖說找到一些廟塔.
有些廟塔是三層.
有些明顯是樓梯.
有綠色,黃色,紅色和藍色.
代表不同時代建造的廟塔.
例如這個綠色一層層的.
就是Wurth III.
吾爾第三王朝.
就是BC主前2100至2000年.
這個大概就是阿布拉罕出吾爾的時候.
吾爾就在下面.
然後有Urdur.
Urdur就是埃利都的古城.
有Uruk.
Uruk就是Gilgamesh的皇帝.
大概在三千多BC的時候.
他做皇帝的一個歷史人物.
在Uruk這個古城.
也有這種類型的廟塔.
綠色的在這裡.
然後去到Old Babylon的時候.
舊巴比倫帝國的時候.
大概2000至1700年主前的時候.
有黃色的.
然後再繼續下去.
有伊蘭和Cassi.
這個大概是中亞述或中巴比倫的時間.
這個大概1400至2000年.
這段時間大概是出埃及的時候.
然後就是Late Assyrian.
就是後期的亞述帝國的時期.
也找到藍色的廟塔.
有好幾十個這樣的廟塔.

$^{441}$我們現在已經知道.
並且文獻去到表述.
描述這些廟塔是塔頂通天.
是天和地之間的匯點等等.
這些我們去到十一章的時候.
說到巴比塔就會多一點.
在這裡為甚麼提出這件事呢.
就是這個洪水的故事.
在Gilgamesh Epic裡面所描述的.
所謂的方舟.
就是一個距離近東的人.
所謂廟的建築.
用的尺碼或形狀都是這樣.
所以你會看到.
有不同近東的洪水的故事.
都是記載有方舟的尺碼.
而且那個尺碼.
本身是和米索不達米亞的文化有關.
你問為甚麼是120乘120乘120.
因為米索不達米亞.
那個完美的數目是60.
我們今天一分鐘有60秒.
等於60進制的東西.
是米索不達米亞人來發明的.
是他們用計算的方法.
計算時間等等.
60是他們想出來的.
所以以前的人都很聰明.
所以在米索不達米亞的文獻裡.
有60進制的數目.
所以你會看到.
這裡表達的廟塔.
方舟就像廟塔一樣.
方舟是一種完美的比例.
你再比較一下.
有一位最近的學者.
把幾個洪水的故事.
裡面講到方舟的平面面積.
去做一個比較.
第一個就是阿達哈斯史詩.

$^{481}$我們曾經提過.
裡面講到神創造人類等等.
有很多東西是和創世記.
一字首一章是可以對比的.
裡面也提過一個洪水的故事.
這個洪水的故事裡面講到的方舟.
就是平面的面積.
是14400的平方爪.
形狀是圓形的.
然後是古格密斯時代.
剛才我們提過.
120乘120平方.
所以那個平方爪都是一模一樣的.
一樣的大.
不過形狀就不同了.
或者是一個廟塔形狀.
然後我們看創世記.
裡面所提過的羅馬的方舟.
居然平面的面積.
是差不多的和另外那兩個.
這個不奇怪了.
因為我剛才講過.
記載這個洪水的災難的故事.
聖經不是唯一一個.
近東的文獻也有提過.
所以你看到有些東西是可以比對的.
或者很相似.
平面面積一樣.
只不過是長方形.
因為方舟的面積是300乘50爪.
我們來看看挪亞方舟.
長闊高是300乘50爪乘30爪.
如果在近東的背景裡.
你就會明白為什麼是300.
為什麼是50.
為什麼是30.
這個數字.
我們之前講的家譜.
或者壽數等等.
都應該不陌生了.

$^{521}$因為米索不達米亞是60進制.
所以300就是60乘5.
50就是300的長度除以6.
所以有6.
或者60.
高是30爪.
這個很明顯就是60爪除以2.
所以也是運用一些數字.
去表達理想或者ideal.
等於Gigamish Epic裡面的方舟是120乘120乘120.
這個可以讓我們想一想.
到底挪亞方舟想記載的是什麼.
我特別有一個構圖.
不用我們平常會看到的構圖.
因為平常看到的構圖.
是我們現代人扼了很多想像.
覺得是一艘船.
正如我說的Téva.
他不是真的想說那艘船.
有船這個詞為什麼不用.
到底他想說什麼.
想比對什麼.
很明顯他是想回應一些當時的.
洪水的故事.
這裡值得我們想一想.
如果你喜歡我們的影片.
就記得要Like, Comment和Share.
拜拜!.
\newpage



\section{}
\label{sec:zuvpXyLO2aE}
\textbf{比對研究 4: 必要, 發展, 方法   -  陳崇基牧師}
\newline
\newline
連結: \href{https://youtube.com/watch?v=zuvpXyLO2aE}{\texttt{https://youtube.com/watch?v=zuvpXyLO2aE}} ~~~~ 語音日期: 2022-06-24
\newline
\newline
\hyperref[sec:X2LCQ0cGbnk]{\small{< < < PREV SERMON < < <}}
~
\hyperref[sec:index_chronic]{\small{[返順時目]}}
~
\hyperref[sec:index_scriptual]{\small{[返順卷目]}}
~
\hyperref[sec:HciKujZi6ec]{\small{> > > NEXT SERMON > > >}}
\newline
\newline
$^{1}$上幾次我們談到比對研究.
是一些工具讓我們宏觀一些.
去看看背後的聖經背景掙扎的問題.
我們在創世記由第一章到現在.
不時會談一些更多的背景.
因為我們開始明白聖經不是天跌下來的.
有它的背景和處境.
這不會令到聖經的神聖.
聖經華語的權威性受到文化約束.
但是我們知道背景的時候.
上帝的華語往往是藉著一些本地的.
一些獨特的處境.
一些社會的狀況.
文化的狀況或者思潮去啟示出來.
所以希伯來聖經也是這樣的.
舊約聖經包括新約聖經.
在不同的時代處理不同的問題.
掙扎不同的問題.
聖經很多時候和這些思潮對話甚至批判.
所以我們從第一章到現在第六章.
就是洪水的故事.
我們先不進入很多經文裡面很細緻的東西.
我們先交代一些背後的東西.
我們上一次提到.
說在十九世紀開始.
當世界開始發現古近東的背景.
聖經的世界的時候.
人們就開始有不同的問題產生出來.
其中一個問題就是為何聖經會好像.
一些近東的故事.
例如洪水故事.
為什麼我們要看比對研究的方法呢?.
因為洪水故事不能夠避免.
要和其他近東的洪水技術做一些比對.
例如很有名的Gilgamesh Epic.
第十一塊泥板.
翻譯出來之後轟動全世界.
到今天還在研究.
到底希伯來聖經創世記.
那個洪水的故事.

$^{41}$和這些近東的故事.
其中一個就是Gilgamesh Epic.
其他還有的.
為什麼會這麼相似呢?.
你看到有17點.
敘事的脈絡其實都是非常相似的.
所以我們需要掙扎這些問題.
上次我們說到在1902年的時候.
就有一個分水嶺的講學.
這個講學就是Frederick Deluge.
他所說的.
這個叫做Bible of Babel.
巴別語聖經.
已經百多年了.
就是因為這個講學.
德國的學者Deluge.
他主張某一些的學說.
某一些的假設.
引起了很大的爭論.
因為當時的德國.
他自己的民族.
正在面對很多要成為強國.
或者民族主義上將等等.
他就有這些影響.
結果我們上次說.
其實其他學者也是這麼說的.
現在已經100年之後.
再回顧就更加明白.
當時為什麼他會說一些這樣的東西.
認為德國民族是很好.
然後以色列的民族.
或者聖經.
舊的聖經.
希伯來聖經.
猶太教是一些後期的劣貨.
一些比較低劣的產品.
近東找到這些泥板的時候.
發現原來更早期的.
這些是聖經的根源.
聖經是抄襲他們的.

$^{81}$或者是模仿他們的.
所以就引致一種這樣的錯誤.
或者一種不應該有的看法.
我們說過那種看法就是.
叫做Parallel Mania.
就是到處去找一些相似的地方.
找到的時候就會認為是巴比倫.
高超過希伯來聖經.
希伯來聖經是後期的.
所以就是比較差的東西.
也引致一些所謂的認信派.
Confessional.
我們今天叫他做保守派.
信聖經或者尊重聖經.
就非常之卻步.
因為很怕看到這些東西.
會不小心將聖經和其他東西相比.
時至今日都是這樣.
特別在西方教會也是這樣.
華人教會更加.
我們一看到這些.
聖經背景.
近東的神話的時候.
我們非常之抗拒.
因為我們覺得勝經是神聖.
是超越性的.
正如我所說.
我們了解那個方法的時候.
了解那個文化的Dynamic.
文化的動態的時候.
我們就不需要這樣看.
上次我們停在這裡.
其實在Delish.
1902年.
他講了這個講學之後.
有大概70-80年裡面.
學界都不斷有更新.
在方法論上是有成長.
有成熟的.
不斷有更多近東考古的發言.

$^{121}$當然.
這些神話的例版.
或者這些神話的文獻.
都是不斷慢慢翻譯出來.
其中一個很好的學者.
就叫William Howell.
William Howell在70年代的時候.
80年代的時候.
他就開始提出.
一種進路叫做背景的進路.
Contextual Approach.
或者有時候是一種對立的進路.
就是Contrastive Approach.
他就寫了幾本.
或者集合了幾本書.
學術的研究.
有三冊.
都是很舊的書.
但是都是很好的書.
因為他在講一種方法.
就是我們去看這些.
譬如巴比倫和以色列.
希伯來聖經之間的關係的時候.
我們不需要一定整天去想.
就是誰借用誰.
一定是後期的前者.
所以後期就是比較低劣一點.
前者就是先祖.
所以他們就比較厲害.
不應該這樣看.
因為文化是不會這樣運作的.
所以要將兩者都放在.
同一個近東的背景下.
就是大家有一個所謂Shared Environment.
一個分享的共通的環境.
一個文化的環境.
甚至是講到一些文學著作.
思潮.
想法的方法等等.
都是大家同屬於同一個背景下.

$^{161}$所以在這個方法裡面.
所謂Contextual Approach.
就是Comparative Study.
現在叫做.
就一定要看相似的地方.
相似的地方是我們預期的.
但是一定要留意.
有分別.
希伯來聖經所記載的可能是同一件事.
近東的文獻神話也同記載一件事.
但是大家的神學.
大家的解釋.
傳釋是很不同的.
所以一定要在一個宏大的文化背景下.
文學等等之下去看聖經.
我們就更加明白文化的動態.
希伯來聖經也是在文化的處境裡面.
去到反應.
去到作對話.
甚至作批判.
有爭辯等等.
所以有相似的地方.
這個是我們絕對會認同的.
有概念上一樣.
譬如我們曾經說過創世記第一章.
所謂Cosmology.
人怎樣看宇宙.
和我們今天21世紀的人怎樣看宇宙.
是完全兩回事.
以前是平面.
上面有個蓋.
蓋上面是水.
下面也是水.
世界本身是一個氣泡.
分開了水.
然後我們就能夠有這個世界.
這些聖經的記載.
創世故事的記載.
和近東的創世故事的記載.
都是有很相似的概念.

$^{201}$但聖經不是.
經文的主要目的.
不是去教我們近東的Cosmology.
不是教我們近東背景下.
近東人怎樣想事.
然後我們就照樣轉換今天.
我們就說這個是合乎聖經.
不是的.
因為那個聖經是用.
神啟示那個經文的時候.
是用近東背景.
人的文化的宇宙觀的概念.
大家當時明白的一種意識.
來啟示出上帝自己永恆的真理.
所以我們曾經第一次的時候都提過.
創世記的時候.
就是我們要記清楚.
想清楚.
聖經不是直接寫給我們的.
是為我們而寫的.
所以我們要明白那個背景.
共享文化背景的時候.
我們就知道.
他們這些文獻.
兩者是有相似的地方.
但是差異我們同樣需要去強調.
這個就是William Howell.
一直去到強調的事.
就是差異.
有什麼差異呢.
希伯來的文化.
就是一個.
獨神論.
然後呢.
一個神人之間的關係是完全不同.
整個世界的秩序.
有倫理上的秩序.
有order, moral order等等.
所以希伯來的聖經.
就是去到否定反對.

$^{241}$或者去辯別.
去爭辯這些思潮.
就告訴希伯來的讀者.
告訴聖經的讀者.
真正的神是怎麼樣.
神在世界當中在做什麼.
所以這個才是所謂.
聖經是為我們而寫.
所以我們今天能夠去應用出來.
所以文化是怎麼work的呢.
我們需要去了解.
坊間,今天的坊間仍然有很多人會這樣說.
聖經那個羅亞洪水的故事.
一定是抄襲人的.
後期的嘛.
真正那個故事本身.
是巴比倫的版本.
所以是照版主觀的抄襲.
所以沒有什麼洞見的.
這個看法雖然是很流行.
仍然是非常流行.
但是有一個很大的漏洞.
那個漏洞就是我們不明白.
近東的文化是怎麼運作.
或者文化一般性是怎麼運作.
如果你喜歡我們的video.
就記得要like, comment和share啦.
\newpage



\section{}
\label{sec:HciKujZi6ec}
\textbf{比對研究 5: 必要, 發展, 方法   -  陳崇基牧師}
\newline
\newline
連結: \href{https://youtube.com/watch?v=HciKujZi6ec}{\texttt{https://youtube.com/watch?v=HciKujZi6ec}} ~~~~ 語音日期: 2022-06-24
\newline
\newline
\hyperref[sec:zuvpXyLO2aE]{\small{< < < PREV SERMON < < <}}
~
\hyperref[sec:index_chronic]{\small{[返順時目]}}
~
\hyperref[sec:index_scriptual]{\small{[返順卷目]}}
~
\hyperref[sec:v_Mfq956F90]{\small{> > > NEXT SERMON > > >}}
\newline
\newline
$^{1}$第五部分 文化古文作需要了解.
不只是古代文化 現今文化.
如果我們想比對研究的時候.
我們知道文化的處境.
我們可能會明白多一點聖經和當時的文化的互動.
我們說過以色列的文獻 希伯來舊的聖經.
有它的處境 有它的文化 有它的文學背景等等.
這些都是和其他以色列以外的埃及 巴比倫 亞述所有的背景共通.
所以我們找到一些大家寫的東西 討論的東西 方法 手法.
是相似的 是絕對是我們預期當中要有的.
應該要有的 正如我所說 聖經不是從天而降的.
有它自己本身寫作的環境.
所以它有一些概念是大家差不多.
其實這個是我們預期會發生的.
不需要詫異的 但我們這個comparative study.
這個approach 比對研究 需要強調不單止是相似.
就是要有differences.
因為以色列的文獻和其他所謂外邦的文獻.
他們的討論是有很大的分別.
本質上有很大的分別 神論 獨一的神論.
或者人觀 世界觀等等 都是有一些對話的.
甚至倫理的教導 神與人之間的關係等等.
這些多數會反對或者提出另一個的alternative.
一個counter的option 或者一個polymics 一個爭辯.
所以在William Harrold和其他教學者開始去認識這一方面的時候.
開始在比對研究的方法上不斷向前.
去refine的時候就明白了.
我們不如舉一些例子 怎樣去明白一下.
一個文化的動態會是怎樣的呢.
不如我們舉一個現代一點的例子.
我在美國的 英文有時候會講這句話.
Eat, drink and be merry, for tomorrow we may die.
或者頭一句 Eat, drink and be merry.
不要理那麼多了 今天我們先行動 今天我們先開心.
因為我們反正都會死的.
英文人講這句話的時候.
很多人都會知道這句話 會懂得講.
但是他很少會想到底是引述哪一個人的.
那個典故是怎樣 我們有一些成語的典故是怎樣.
但是沒有人會理會典故是怎樣的.

$^{41}$也沒有人會在講這句話的時候刻意或者無意去引用典故.
這句話我們很多時候說.
本身出自一個希臘的哲學家叫做Apicurus.
Apicurus是第一個去講這句話的.
Eat, drink and be merry, for tomorrow we may die.
他是主前三百二百年的人.
他是一個希臘的哲學家.
所以他是第一個這樣講.
但是沒有人在引用他.
不單止是這樣 你再想多一步.
就算是Apicurus是第一個在文獻裡面.
在他的寫作裡面有提過這句話.
他都不是第一個去present(展示)這個idea(想法).
或者不是第一個去講人生的觀念的取態.
這個已經是約定俗成也好 或者已經是繪製人口也好.
或者已經在文化當中 大家共享的環境裡面.
已經是假設了 預設了.
只不過是有人講了出來而已.
formally(正式)講出來而已.
其實一向人都已經有這樣的看法的了.
所以今天聽到有人講這句話的時候.
如果你明白文化的dynamics(動態).
你就不會說他在抄襲誰 或者在引述誰.
他只不過是跟他共有的文化.
一些的意識去到一些的resonance(共鳴).
一些的迴響 有些的共鳴.
所以你這樣知道.
不是說聖經裡面有些東西很像近東的文獻.
所以聖經是抄襲他們早一點的.
不是這樣work(運作)的.
In fact(事實上)你看這個Epicurus(伯多利亞).
如果你記得在史蘭傳第17章的話.
保羅去到雅典的時候 他就跟哲學家去辯論.
其中一班就是Epicurean(伯多利亞).
這班人本身已經是很流行的了.
這個哲學的派系和Stoicism(斯多噶主義).
都是很流行的了 去到第一世紀的時候.
你記得三百多年前.
就是這個Epicurus(伯多利亞).
他開始講述 講學等等.

$^{81}$直到有些弟子 他們都有些這樣的哲學的主義.
保羅就跟他們去爭辯.
就說這個人講什麼 是不是外邦神明來的.
因為他在講耶穌和復活的福音.
所以不能夠這樣講 就是說總之講到任何東西.
Epicurus(伯多利亞)那句話就是在引用.
因為已經是含括在文化裡面.
大家都在裡面做一些回應.
甚至可能有人會說 耶穌在登山寶訓那裡說.
不如為明天憂慮 因為明天有明天的憂慮.
一天難處理 當就夠了.
有人說 這個就是抄Epicurus(伯多利亞).
希臘的哲學 是不是.
不是的 因為這些是很流行的人生的哲學.
人生的取態 是不是.
所以耶穌去到講這個說話的時候.
他只是在回應這個已經在文化裡面.
共有了很久的一種的意識 是不是.
所以他就會教門徒 我們不是說.
Don't worry, nothing to worry about.
我們就行啦 我們信神就行啦.
他不是在講這件事.
他在挑戰這件事 同時講出.
一個真正的門徒 要擔心的是什麼.
我們真正憂慮的是什麼呢.
憂慮的不是這些東西.
那些外邦人 你記得嗎.
外邦人所憂慮的東西.
耶穌說什麼 要憂慮的是.
求上帝的國和義.
我們每天擔心 要憂慮 要去到work on(工作)的.
就是上帝的國和上帝的義.
所以是一種文化的互動 是不是.
不是引述誰 抄襲誰 是不是.
這樣想太過幼稚.
再舉多一個例子 今朝有奏今朝在.
大家很明白這句話.
明日壽來 明日壽 是不是.
沒有人說這句話.
是說引自我找到的典故.

$^{121}$就是唐代的詩人羅忍.
沒有人會說.
我們現在是在引述他.
或者我們說一些類似的東西的時候.
我們就是在引述這個人.
這個文化不會這樣work(運作)的.
因為文化本身已經有很多這樣的意識在.
只不過有個唐代詩人.
他作了一首詩 說了出來.
再舉多一個例子.
就是911.
911是一件很大的事 是不是.
震驚世界 震驚美國的事.
如果你看看911之後寫的書.
或者拍的電影.
或者做一些藝術作品.
都會去詮釋911的事件.
這個大的災難.
到底發生了什麼事.
原因在哪裡.
aftermath(後面的數據)是怎樣.
中間的情節是怎樣.
如果你看看這些書.
或者紀錄片.
或者電影.
大致上它的脈絡.
大家都掌握得很好.
因為一些很基本的脈絡.
就是雙子大廈等等.
United 93.
但如果你看不同的詮釋.
不同的記載的時候.
他自然的手法就不同.
有些人做了很大的人像出來紀念.
有些國家或者美國的主義.
他就去表達出來.
譬如這個History Channel.
Dates that shape America.
他就有不同的角度去切入.
這個Commission Report.

$^{161}$一個美國政府官方的報告.
或者.
這本書.
Debunking Myths.
一些神話.
一些Conspiracy Theories.
有很多人引起一些陰謀主義的理論.
為什麼911會發生.
裡面有些什麼事.
他就去Debunk.
去De-Mythologize.
沒有這些事的.
有些人是這樣寫.
回應這樣的詮釋.
你這樣看的時候.
對一個災難.
不同的詮釋不同的回應.
或者有時候De-Myth.
有些陰謀論.
或者錯誤的看法.
他就改正.
不同的手法去回應.
那件事.
你這樣看就像.
羅亞洪水的故事.
和其他近東的.
洪水故事.
都是做一些類似的事.
不是誰抄誰.
又不一定.
最早期.
9月12號那年.
寫的書就一定是最準確.
有很多時候.
回想的時候.
後面寫的東西.
才是更加準確更加有權威性.
因為調查清楚.
反而早期寫的東西.
可能就是口齒薄皮.

$^{201}$或者以片開傳.
後期的可能更加.
是Reliable的報告.
所以不能夠說.
因為聖經是希伯來文.
所以給那些阿卡德文的文獻.
就一定是遲了.
所以他一定是Latecomer.
好像Deluge.
1902年Bible of Babel的講座.
這種看法.
完全是.
不明白文化.
是怎樣的.
一個共享文化的環境.
到底那個.
互動的關係.
是怎樣.
所以回到這裡.
比對的研究.
其實最主要不是去跟蹤.
誰先.
誰後.
先的超越一點.
後的低一點.
這個方法完全是錯的.
反而我們留意的是整體的環境.
有什麼相似.
和有什麼不同.
所以到70年代之後.
這個方法開始掌握了.
希望可以.
對我們聖經研究.
我們的焦點.
有幫助.
所以你看到Google Maps Epic.
Tablet 11.
跟羅亞洪水相比的時候.
脈絡當然是.
很相似.

$^{241}$正如你看911不同的書.
講論這件事.
脈絡都是一樣.
想起洪水就一定有方舟.
所以有什麼事發生.
然後水乾就會有什麼事發生.
所以在情節上.
有相似的地方.
絕對是我們會期望的.
反而可能.
聖經我們需要留意的.
就是它的差別.
它的獨特的地方.
它的不同的地方.
到底在哪裡.
我們可以用這個.
途徑去了解一下.
背後.
大家有一個所謂的911.
一個歷史的.
洪水的災難.
然後有不同的人.
有不同的版本.
巴比倫可能有Old Babylonian的版本.
有Standard.
之後的.
New Babylonian的版本.
亦都可以有亞述自己的版本.
以色列.
有希伯來.
聖經的版本.
黑人亦都有它的版本.
這些就是不同的文士.
不同的社會.
他們去到回應.
大體上是同一件事.
這兩本書.
我列出來的.
其實主要是在講.
當時近東的社會.

$^{281}$文書的文化.
Scribal culture 是怎樣的.
我們不明白的時候.
我們很多時候將現代人的Scribal culture.
我們的writing culture.
我們的書的文化.
又講版權.
不准抄.
抄了就是壞人.
抄了就是後者.
就是二次創作.
就是低等一點的.
不原創的.
近東的Scribal culture不是這樣的.
這個.
Vander Toon.
近東學者特別做了一個研究.
一個很基本的研究.
就是用一些.
近東的Scribal.
他們留下的泥板.
他們怎樣書寫 怎樣抄東西.
怎樣去流傳一些東西.
他們沒有講什麼侵犯版權.
沒有說.
會不會是.
去到照抄是不對的.
每個Scribal.
都是這樣.
流傳一些所謂的.
sacred的tradition.
這些可歌可泣的.
作品.
所以他們就在泥板裡面.
一直抄.
你這樣看的時候 有這樣的背景.
你就知道 希伯來聖經是怎樣的.
我們經常說希伯來聖經.
都是好像一個archaeological tale.
一層一層的.

$^{321}$有它的歷史.
之後的.
所謂editor.
編輯 可能會加一些東西.
update一些東西.
他們在Scribal culture裡面.
才能了解他們在做什麼.
不能用現代人的眼光來看.
另外一本比較新的.
就是John Wotten和Brandt Stendhal.
這個Laws of Scripture.
Laws of Scripture.
講到literary culture.
就是古代人.
這種書寫的文化.
所以講到這裡.
我們可以再問多一點問題.
就是.
特別是講到洪水.
不單止洪水.
有很多其他的事跡.
譬如亞述王 為什麼亞述王死.
獵王劫上下 都有不同的傳息.
對近東的文獻也有不同的傳息.
所以講到洪水的時候.
有什麼的.
dynamics需要留意.
就是希伯來聖經和其他的.
近東的文獻.
都是分享著一個傳統的背景.
有不同的傳息.
有不同的技術.
有不同的切入點.
最著名的.
最多找到copy的.
就是這個Gilgamesh Epic的Akkadian version.
應該是出自巴比倫這個地方.
因為在近東很多其他的地方.
都找到Gilgamesh Epic的.
例版或是例版的碎片.

$^{361}$巴比倫就在.
米索不達米亞的南邊.
這裡就是烏魯克城.
其實這個Gilgamesh Epic就是在講烏魯克王.
他的冒險的故事.
他的豐功偉績等等.
裡面有很多很精彩的細節.
其實他就是在找著.
長生不老 不想死的Obey.
所以就攀山涉水.
有很多的經歷.
很精彩的.
巴比倫也找到亞述的一些例版.
甚至到客人的地方.
喀塗沙 京都.
都有客人去抄寫.
或者去講述烏魯克王Gilgamesh Epic的故事.
甚至在密吉都.
以色列的米吉都.
也找到Gilgamesh Epic的碎片.
所以看到有不同的迦南地的文士.
也知道或是抄寫這些故事.
最著名當然就是阿卡利安.
最多最完整的版本就是他.
但這裡很有趣.
我就引述Jeffrey Teague.
裡面講到一個比較技術性的.
講到這些例版.
或者文學作品之間的關係.
講到這個literary borrowing.
我們經常用borrow這個字.
一用borrow的時候.
我們就用現代人的版權問題.
我抄你 你抄我.
後期抄早期的等等.
他講到這個Gilgamesh Epic的不同的例版.
他就提到這個黑人的version.
亞歇提的version.
他發現有趣的地方.
黑人的version就是縮減了或是省略了.

$^{401}$巴比倫版本的一些重要情節.
省略了什麼呢.
就是講到烏魯克城的偉大.
城牆等等.
黑人的版本就將它省略了.
增長了什麼呢.
增長了國王去征服黎巴嫩森林.
Lebanon Forest的旅程.
Gilgamesh Epic裡面.
其中一個很重要的情節.
就是烏魯克王去打敗守護Lebanon Forest的守護神.
叫做胡娃娃.
很精彩的將她打敗了.
殺了她.
這個version很奇怪.
講到Lebanon的征討旅程.
那個冒險的故事.
非常詳盡.
甚至加了很多不同的細節.
問題就是這樣問.
T記就問.
到底是黑人的文字.
抄襲了早期的巴比倫的阿卡德.
Arcadian的version.
還是黑人的文字有另一個我們還不知道的版本.
那個版本可能早過我們現在知道的Arcadian的version.
希望你弄清楚.
不是一定是黑人抄襲巴比倫的版本.
所以黑人是後期.
這個黑人也不後期.
這個黑人的版本是Lebron的時候.
1500-1200BC的時候.
還是黑人有另一個版本我們現在找不到的.
而巴比倫和黑人都是抄襲那個版本.
他們故事的根源是從那裡來的.
只不過我們還沒有找到.
我們還不知道有.
這裡很有趣.
為什麼他提這個黑人的版本呢.
就是有趣的地方.

$^{441}$就是他刪減了巴比倫烏魯克的故事的細節.
不說那個城多偉大.
而詳述的就是那個黎巴嫩森林之旅.
如果你看地圖你就知道.
因為很可能在黑人.
黑屠殺在這裡.
黑人的版本的Gilgamesh Epic.
他們對有什麼有興趣呢.
巴比倫在烏魯克的城有多偉大.
他們不在乎.
他們在乎的是近距離他們的黎巴嫩森林的旅程.
才是最吸引他們的讀者.
所以他們留下的就是一個比較詳盡的.
Lebron Conquest的故事.
Lebron就在這裡.
就是拿近東世界裡面最好的木材.
香柏木就在這裡.
你記得所羅門王建造耶路撒冷聖殿的時候.
你回去看那個細節.
就是問誰拿木材.
就是問推羅西頓.
泰西西頓的Hiram.
希蘭王.
就是拿什麼呢.
巴比倫的香柏木.
然後將它運到約柏港.
然後找牛車也好.
找人也好.
拉上去耶路撒冷建第一個聖殿.
所以巴比倫.
一個黎巴嫩的森林.
對於客人來說就大感興趣.
因為他們最近.
和他們最有興趣知道的一個故事.
所以無論如何.
你講文學背後的關係的時候.
就不一定是先有巴比倫.
然後有他們的版本1234.
然後客人抄襲他們.
其實可能大家都是源自一個更早的版本.

$^{481}$這個原則.
你就明白了.
在一個這樣的大文學文化裡.
大家在assume911裡面.
不同的著作.
不同的傳識等等.
你就需要明白.
是不是就好像Delos.
1902年八十多年前.
那種幼稚的說法.
巴比倫先.
所以就正一點.
然後希伯來聖經就後一點.
所以是差一點.
劣一點.
在這裡你又需要留意.
聖經的傳統.
強調什麼呢.
你留意聖經.
由第一章到第十一章.
到阿巴拉罕出現的時候.
都是強調.
阿巴拉罕是米索不達米亞的人.
阿巴拉罕是什麼人.
阿巴拉罕是南米索不達米亞的吾爾人.
即是他在那裡出來的.
整本聖經經常都有留意.
甚至五經以外都有這樣說.
約書亞記.
第十四章.
這裡最後結束了.
所謂事件之約.
那個講話.
約書亞講什麼呢.
對中百姓講.
耶和華神是如此說.
古時你們的列祖就是阿巴拉罕拿學的父.
他拉.
他住在哪裡.
大河那邊.

$^{521}$across the river.
那個river是什麼river.
就是Euphrates.
兩河流域.
米索不達米亞.
一條最大的河.
上面就是Tigris.
底格拉斯河.
然後這條就是所謂大河.
就是Euphrates.
幼法拉底河.
他是哪邊的人.
即是他們來到事件.
入主迦南.
然後說.
記得你們的先祖在哪裡.
大河那邊.
我們回去查一下原文.
約書亞這句說話.
其實開始的時候不是講古時什麼的.
第一句說話是.
住在哪邊.
住在大河那邊.
所以強調列祖的原居地.
他是什麼人.
米索不達米亞.
當然是不是說.
阿巴拉罕出吾爾的時候.
他帶著一些禮板.
或者等等不同的記載.
這個不能夠證明.
就算是口傳文化.
All culture.
他都會知道有什麼故事等等.
洪水故事.
我們講到創世記1至11章.
經常都提到米索不達米亞.
為什麼.
因為1至11章是非常濃厚.
有米索不達米亞的意識和思潮在裡面.

$^{561}$那種思想.
洪水也是.
所以這個比對研究.
去到逼我們去想.
到底共享的文化是怎樣.
所以我們不需要驚訝.
為什麼這些洪水故事和近東的故事這麼相似.
根本應該是相似的.
正如你拿起兩三本不同911的講論書.
只不過詮釋是很不同.
所以這個.
下一次我們繼續講.
這個洪水故事的詮釋上.
到底有什麼distinctive.
有什麼獨特的地方.
如果你喜歡我們的YouTube頻道.
就記得Like, Comment和Subscribe.
\newpage



\section{4:25-5:2}
\label{sec:v_Mfq956F90}
\textbf{Relationship in the Church (Ephesians 4\_25-5\_2) - Steve Tan}
\newline
\newline
連結: \href{https://youtube.com/watch?v=v-Mfq956F90}{\texttt{https://youtube.com/watch?v=v-Mfq956F90}} ~~~~ 語音日期: 2022-06-27
\newline
\newline
\hyperref[sec:HciKujZi6ec]{\small{< < < PREV SERMON < < <}}
~
\hyperref[sec:index_chronic]{\small{[返順時目]}}
~
\hyperref[sec:index_scriptual]{\small{[返順卷目]}}
~
\hyperref[sec:w05Jl15bCEk]{\small{> > > NEXT SERMON > > >}}
\newline
\newline
$^{1}$Hey brothers and sisters, good to see you guys again..
Especially the college adults,.
it's good to see you guys back..
Oh yeah, I have been cutting my hair for the last two years.
and I kinda messed it up, so I gotta start over again..
Yeah, so I did save a lot, but I appreciate it..
Let's bow our heads, let's get ready for God's message..
Father, we thank you.
that you are the everlasting God,.
that you have so much.
that you probably need to take care of, Lord..
But yet, you would spend Sunday,.
a day that you command us,.
where we need to pay attention to you, Father..
Lord, it's for our own good and it's for your glory, Lord..
And we pray that everyone here,.
may your spirit be with us,.
may you open our hearts, Lord, to receive your message..
In Jesus' name we pray, amen..
Well, a couple of time ago,.
I probably told you guys that I had a family feud.
between my uncles, you know, who, over some inheritance,.
it's not like 100 million or anything,.
it's like, it's like child change, you know..
Because of my grandmother died, you know..
It's feel like these relationships are broken right now,.
because of greed and pride..
They're still not on talking terms as of now,.
as I know..
But I just want you guys to think about the church..
How are relationships broken or strained in the church?.
Think about that..
What caused these broken relationships?.
What caused these strains?.
Well, I can tell you, I mean, based on my experience,.
probably a lot of you guys have been through it too..
A lot of things could be like gossiping, you know,.
being arrogant, being prideful of yourself,.
or even it could get to a point like being unable.
to share power between the leaders and the congregants,.

$^{41}$right, how things should be done..
These could cause a lot of disunity..
But the reality is that, as we are in the fallen world,.
I mean, the way how we do things,.
I mean, a lot of times it leads to our,.
I mean, prideful arrogance, which leads to sin, you know,.
which cause friction among the body of Christ..
That God, I mean, that God,.
he did everything he could to establish this church,.
this body..
Therefore, do we just accept the way it is?.
Well, the world is broken, right, it's fallen..
The way how we act, is this the way we should accept.
this reality?.
Absolutely not..
Because God expects more of us than this..
So today, I want to.
talk about the relationship in the church..
How do we conduct ourselves?.
In the past semesters, I talked about how did Jesus Christ.
establish this church, right?.
How, at the end, we have to be united.
and be mature inside the church,.
but how do we relate to each other?.
Today, I want to talk about that..
(papers rustling).
This is the thought, yeah, how we conduct ourselves..
Okay, the first point I want to talk about is like,.
do not give the devil a foothold..
Do not give him a stronghold in your life..
So, starting on Ephesians chapter four, verse 25 to 28,.
it says, "Therefore, each of you must put off falsehood.
"and speak truthfully to your neighbors..
"For we are a member of one body..
"In your anger, do not sin..
"Do not let the sun go down while you are still angry..
"And do not give the devil a foothold..
"Anyone who has been stealing must steal no longer,.
"but must work together, do something useful.
"with their own hand, that they may have something.

$^{81}$"to share with those in need.".
Now, right here, Paul starts with,.
"Therefore," that means there's a relationship.
from a previous passage to this one right here..
Yes, the previous passage, which is chapter four,.
verse 17 to 24, talks about the general principle.
of the difference between the old man,.
which is your old identity, your pagan identity,.
comparing it to the new identity, which is in Christ..
So, it's a very general comparison, telling us.
we must put on our new man, our new person,.
get rid of the old person..
Right here, Paul is giving us some specific guidelines.
of how we should conduct ourselves among brothers.
and sisters in the church..
So, let's take a look at this..
Here, Paul gives a list of actually,.
a list of vice and virtues, of what we shouldn't do.
and what we should do..
So, virtues are things that kinda like unite us,.
and in contrast, vice are things that disrupt.
the unity of the community..
So, the first thing that you see in verse 25.
is basically, it says, you gotta put off falsehood..
The first vice is lying..
This is the mark of the old person,.
which is corrupted by sin..
But, Paul says, this is not enough..
We can't just stop lying..
You have to replace it with a virtue,.
and that virtue right here is that we must speak truthfully..
But, why?.
I mean, the motivation is that because we,.
as the second part, it says, we are all members of one body..
We are part of the body of Christ..
We have to put on this new person here..
So, besides that, speaking the truth to each other.
is also a direct quote that Paul uses from Zacharias..
It's a old, Zacharia, the book of Zacharia.
is a Old Testament book..

$^{121}$It talks about what God's promised about the new Jerusalem..
So, these are, this is what happened..
During that time, they were allowed to come back.
during the Persian Empire..
They have been exiled for 70 years..
They came back, and God's telling them,.
you guys will inherit Jerusalem again, right?.
Since Paul is using this quote here,.
it seems like he's talking to us..
We're the recipient of that promise..
And, these are the promises that the prophets.
and apostles are making known to us right now..
I mean, us in terms of what I'm talking in their context.
to the first generation Christians.
that no one knows until they reveal it to them right now..
So, it is significant for us..
It's just that we are the recipients of this..
We have to behave accordingly to who we are..
Therefore, speaking the truth is a sign of this new person.
because God is truthful..
As we move on to verse 26,.
the second thing is sinning in your anger..
That's the thing Paul wants us not to do..
We should stop sinning when we're angry..
Paul is also quoting, actually, from Psalms 4..
You know, he's quoting about Psalms David wrote,.
claiming that he is being slandered.
because of crimes that he did not commit..
He is sure of that..
Therefore, what does he do?.
He cries out to the Lord,.
and his anger turned to joy and peace..
Now, this peace and joy,.
despite the situation you're in,.
is a mark of a new person..
Now, as we continue to take a look,.
it seems like anger, seems like it's permitted..
So, is anger in general okay, we're asking?.
Well, I would tell you being angry for the wrong reason,.
of course, is obviously not right..

$^{161}$However, can we be angry for a righteous cause?.
Like, if you see oppression,.
or you see injustice, of course we can..
There's such thing as a righteous anger, right?.
God has righteous angers..
But I could tell you, rarely our angers are righteous..
Even if it is, do not let it boil..
So, because as the second part says, right,.
do not give the devil a foothold, right?.
Do not let the sun go down while you are still angry..
So, you gotta get rid of it quickly..
So, because our natural,.
we still have our old person within us, you know?.
It's suddenly kinda like, let us feel like,.
oh, whatever we always experience, our anger,.
must be righteous, but other people's anger.
is always like a tantra, bad temper, right?.
So, they must be wrong, I must be right..
But I'm telling you guys, brothers and sisters,.
better assume the worst in us..
When you have anger, I mean, better assume.
it might mostly not be righteous,.
because it could be damaging to relationship..
Let me give you an example..
For me and my son, right, I told you guys,.
I mean, my son is here, but you know..
During this time, I feel like he's in a terrible three,.
right, whatever I tell him to do, he does not do..
If I tell him not to do it, he will do it..
Even when I feel like I have to talk to him,.
I yell at him, whatever, I feel like I'm right..
But in the heated moment, I could tell you,.
my discipline will not be the best,.
because I'm very emotional at the time..
Like, it's quite personal, right?.
Let's talk a little bit about more.
that we could relate to the body of Christ..
Like, I remember, not remember,.
it's still happening right now, right?.
The last couple years in Hong Kong,.

$^{201}$there's some social unrest, right?.
I'm not talking about commenting.
on anything political right here,.
but more of like, during that unrest,.
during these couple years,.
I mainly look at how we respond as Christians, right?.
Especially, they kinda, what I saw,.
is that there became two factions, right?.
One has a different, has one viewpoint,.
the other has a different viewpoint..
The thing is that, what I saw is that,.
despite that emotion, that anger,.
at the end, they were at each other's throat,.
and they did not talk to each other, right?.
I even see that from some of my own brothers and sisters,.
and man, that was tough to see..
What I'm trying to tell you is that,.
even with the best intention,.
our emotion will cloud our judgment..
So, as I said, if you continue to look at verse 27,.
it seems to suggest that anger is not caused by the devil,.
but Satan will use it against you.
to disrupt the body of Christ..
But let's come back to this issue in a moment..
Let's go to verse 28..
The third vice is stealing,.
and the virtue is not just to not steal anymore,.
that's not enough, Paul is saying..
We need to work in order to have something.
to share with others in our community,.
that's what he's saying..
Instead of having something that you have negative,.
we just don't turn it to zero,.
we don't just neutralize it, as a Christian..
A lot of things, we have to have a plus,.
from negative, not just a zero, but a plus..
And right here, as you see the word work,.
in this context, it means to labor extremely hard.
to the point of fatigue..
So we have to do our best with our abilities,.

$^{241}$with our mind, to work, so why?.
It says so we can have something to present to others,.
especially the ones in need..
Why?.
Because we have to make contribution to the body of Christ,.
for the sisters..
Now, this letter didn't address any specific problems.
towards the Ephesian church, mainly,.
but it's also a circulating letter..
I mean, other churches will read this..
These vice might be like a general practice.
that Christians, during that time,.
after they have been converted to the new person,.
they might continue to be engaging.
in some of these activities..
I mean, stealing, along with other things,.
are probably still acceptable during that time..
Why?.
They don't have social programs.
during that time in the Roman world..
I mean, you have to make a living somehow, right?.
But the thing is that, what Paul is saying,.
but we have to get rid of those things..
If you are gonna be a new person,.
you have to get rid of these things..
These are not acceptable..
Because these things are causing harm.
to brothers and sisters, right here..
Now, let's get back to verse 27..
This kind of summarizes the points right here..
Giving a foothold, right?.
I could tell you, right,.
Satan do not want anyone to be Christians.
or to live in a Christian manner..
It means when we give him a foothold,.
that means we do not reflect the new person..
We do not reflect Christ-like behavior..
Our lies, our angers, our stealing, you know..
This is what the devil needs..
He need you to do these things.

$^{281}$in order for him to infiltrate the church..
I mean, one of the reasons why relationship is so difficult,.
because Satan attacks them..
He wants to undermine what God has created..
So, what in the world is really a foothold?.
Take a look..
Imagining yourself being chased by a bad person,.
and you get to a point, you get to your house,.
and you close the door and shut 'em off..
But no, this person is able to stick one foot in the door..
And think about it, when his foot is here,.
you can't close that door completely..
So, basically, this is what the devil is doing every time..
When you can't close the door,.
he's able to come in, give you a whisper..
Hey, do this, this is good, stealing is good..
You're just doing a little stealing, it's okay..
It's nothing, you're not gonna hurt anyone..
He's continuing to influence your life right here.
with that foothold..
So, in conclusion, brothers and sisters,.
we need to get rid of the devil's foothold.
at the church door..
Therefore, we gotta produce these virtuous behavior.
that will serve others in the community..
This is a practical step to continue.
to maintain this unity in Christ..
The second point I wanna talk about.
is don't grieve the Holy Spirit,.
which consists of verse 29 to 32..
Do not let any unwholesome talk come out of your mouth,.
but only what is helpful for building others up.
according to their needs..
That it may benefit those who listen,.
and do not grieve the Holy Spirit of God,.
with whom you were seated for the day of redemption..
Get rid of all bitterness, rage, and anger,.
brawling and slander, along with every form of malice..
Be kind and compassionate to one another,.
forgiving each other, just as in Christ God forgave you..

$^{321}$Now in the second point,.
after talking about these lies, these angers, these stuff,.
Paul seems to returning to the first topic of speech..
The first topic of speech is speaking the truth in love..
Here, Paul's telling us not to engage.
in unwholesome or harmful speech..
But the thing is, our culture makes that a norm..
It kinda infiltrates into our church..
Think about it, here's some examples..
I always see, especially with ones with a consumer mindset..
Wow, our church, look at that nursery, man, it's so small..
You know, I wish there was a slide in there..
It would make the nursery way better..
You know those Sunday school teachers?.
It's not passionate at all..
If they were more passionate, we have better curriculums,.
I think my kid would join, so I can't come to church..
Oh, my church doesn't have a coffee machine, you know?.
I can't get free coffee, man, that sucks..
Or come on, we need some sofa here, you know?.
I mean, it's always consuming, like what we don't have,.
what we don't have, this or that, you know?.
This is that consumer mindset..
Another thing is that, I would say, I don't know,.
I hear this a lot from the Chinese churches..
If I have a problem with someone,.
I have a problem with person A right here,.
I will not talk to them,.
I cannot confront them directly..
It's always like, hey, do you know what person A did to me?.
Hey, do you know what person A did to me?.
But it's never confronted, face to face, you know?.
You never resolve the problem..
And one more thing I wanna say is that.
we always like to criticize our leaders and pastors..
Only if the pastor could do this or that, right?.
Only if the pastor could preach better..
Dude, I'm telling you guys,.
the pastor is just a human being like us, right?.
He has his talent, he has his shortfall also..

$^{361}$I can't expect a Pastor Mark to do everything perfectly.
and everything good..
That's why you guys are here..
The body of Christ gave you guys give, right?.
We all have some special gifts,.
so you can't expect the leaders and the pastor.
to do everything and you're just sitting here.
complaining this and that..
Now, instead of engaging in these harmful speech,.
engage in helpful speech, right?.
Speech that is beneficial..
We need to build up our community according to our needs..
Well, I would say this is what we should do..
We gotta think about how we speak to each other,.
especially how we cannot harm other with these speeches..
Watch out for our tones, for our demeanor,.
so as we talk to each other, be wise and sensitive..
Learn when and where is the right place.
to say what you want to say..
How about encouraging each other?.
I feel like we as Asians, we like to criticize,.
but how about saying good job?.
I'm glad you did a good job, right?.
Nice try..
We are all ever asked,.
hey, how are you guys doing these days?.
How are you guys doing emotionally, right?.
I mean, these are things we have to do more,.
something more positive to lift up our brothers and sisters.
with our speech..
I mean, in this case, if we look at this context,.
verse 29 flows well with verse 28,.
because both of these verse.
is to seek the benefit of others..
Verse 28 is to seek the benefit of others with our actions,.
and verse 29 is to seek the benefit of others.
with our speech..
These are really essential things.
for the good and the growth of the church..
As we are obedient to scriptures,.

$^{401}$the concern of the growth and the good for the church.
should be natural to every one of us.
who profess faith in Jesus Christ..
So think about it..
Next time you see someone getting baptized,.
are you happy for that person?.
You're happy because that person.
pledged obedience to Jesus Christ,.
and Satan's kingdom is crumbling..
Are we happy when someone else is shining up here.
besides us?.
Are we happy when we see someone.
is maturing in their Christian life?.
We should be happy..
These are natural for us..
It should be natural for us..
Let's move on to verse 30..
Besides trying to maintain this unity,.
verse 30 also gives us another motivation.
to achieve that,.
which is to not grieve the Holy Spirit..
We should not make the Holy Spirit sad or feel sorrow..
Now, the Holy Spirit of God right here.
is a link to Isaiah, the book of Isaiah,.
in Old Testament, chapter 63, verse 10..
I mean, right here, the Israelites are constantly sitting..
Isaiah is a prophet, he's trying to tell them,.
you better come back to God,.
if not, bad things are gonna happen..
So right here in this chapter,.
he's reminding the Israelites,.
how God saved you guys from the Egyptians.
and brought you into the wilderness,.
yet you guys are unfaithful..
And God became their enemy and told them,.
you guys, this first generation of Israelites,.
will not enter the promised land..
Hey, this same idea is applied to this verse.
of God's new people, the body of Christ..
Hey, we have been sealed with the Holy Spirit..

$^{441}$It's dwelling within us..
We have been blessed with every blessing from Christ,.
and we are on to full redemption.
when all things are under the authority of Christ,.
when God bring that to one day..
No, this taste, the taste of this glorious day,.
theoretically, should be seen in the church..
It should be..
But Satan is doing everything to undermine it..
But guess what, brothers and sisters?.
It's useless..
Yes, we seem to be stumbling back and forth, right?.
We probably seem to be losing a battle there and here, right?.
But we didn't lose the war..
Because the Spirit is greater than the devil..
For the things that we still behave like the old person,.
the devil's foothold will still be on us..
That grieves the Holy Spirit..
So we need to behave like the new person..
Build up the church, and when we build up the church,.
it simultaneously tears down the devil's stronghold..
Let's move on and continue to verse 31 and 32..
The next two verses is a summary of the way.
we should live by not grieving the Holy Spirit..
If you take a look at 31, it gives a broad list of vice..
It's talking about bitterness, rage, and anger..
These are very internal and personal..
And from there, it goes to brawling and slander,.
to verbal and outward..
And this results to every form of malice, of wickedness..
But in turn, Paul is saying right here in verse 32,.
we should be kind, compassionate to one another,.
forgiving one another..
And Paul is saying, the motivation is found, right,.
at the end of verse 32, just as in Christ God forgave you..
So basically, this is the higher order of command,.
of encouragement, you know..
God demonstrate forgiveness,.
so you must also demonstrate forgiveness..
Paul is referring back to his theology.

$^{481}$of what God has done for us in the first three chapters.
of Ephesians..
What Christ has done for us,.
what God is currently doing in our lives,.
then we should be motivated to live the right way..
So, if we're not motivated, this is how it would look like,.
I would say, if we grieve the Spirit..
You're in a relationship, a lot of us are couples,.
if you're not in, probably you will have an idea..
Relationships are complicated..
As I know a person who lives in the state,.
recently, the wife.
found out that he was committing pornography.
since whenever, you know..
And the wife found out,.
he told the congregation and the leaders,.
and this person had to step down..
I'm so glad the wife did that,.
because it made this person, this friend,.
came to his sense..
He still feel like, it did not hurt anyone,.
but myself, it's okay..
But it hurt his wife,.
and other people in the congregation.
once they found out,.
because he deceived others in the church..
So, the wife is obviously mad..
I would say the main emotion is sad..
But the thing is that,.
the wife is still there,.
but the wife says,.
"I don't trust whatever you're saying right now.".
You know?.
So, at nine o'clock,.
the computer is gonna be on a password..
And they have to go to counseling, right?.
The wife is really mad..
If I could tell you one good thing.
that comes out of this, brothers and sisters,.
if the wife or the husband is still grieving,.

$^{521}$that means they're still in it, you know?.
They're still in this fight with you..
If the Holy Spirit is still grieving,.
they're still in this fight..
If one day your wife or your husband.
or your girlfriend or a boyfriend said,.
or they don't have any emotion,.
"Oh, I don't care..
"I don't care about that person anymore.".
Then you worry..
Then I'm gonna tell you,.
your relationship is over..
Do you wanna get to that point?.
I don't think the Spirit will get to that point.
in terms of our theology that we know..
But grieving God's Holy Spirit is just like that, right?.
Let's not get to that point, right?.
And the last part I wanna talk to you about,.
the last point is point three, is imitate God..
This is actually the conclusion.
Paul uses to end this passage..
Now, after he's giving us.
Roddy's specific points of we should not steal,.
lie, or anger, he sums up his main point right here..
He says, "Follow God's example..
"Therefore, as dearly loved children,.
"and walk in the way of love,.
"just as Christ loved us and gave himself up for us.
"as a fragrant offering and sacrifice for God.".
As we identify as God's children,.
we are to walk in love, Paul is saying..
And where's that motivation coming from?.
In verse two, it says,.
"Just as Christ loved us and gave himself up for us.
"as a fragrant offering and sacrifice to God.".
We are to do the same, we are to have the same attitude..
We are to imitate God..
We are to walk by love by displaying sacrificial love.
to others, to our brothers and sisters..
But you're saying, you might be thinking,.

$^{561}$it sounds impossible, right?.
But I want you to remember two things,.
my brothers and sisters..
Number one, the Spirit is within us..
He will complete the sanctification work within us..
Therefore, we could speak truthfully to each other..
Therefore, we can quickly subdue our anger..
Therefore, we can supply others.
with our work to supply their needs..
I mean, these are a couple of important virtues.
that Paul is talking about to build up the church..
But the general principle Paul is saying is that.
the Spirit is greater than the devil..
We can overcome any other sins.
that prevent us from maintaining this unity in Christ..
The second thing I want you guys to understand is that.
this is the expectation of a normal church..
This is what a church should look like, Paul is saying..
And God commands it..
If God commands it, that means it's doable..
And God will give you the resource to complete the job..
Now think about a guy tearing down buildings..
In 2004, during Hurricane Katrina,.
I had a chance to go down to New Orleans..
My main job of ministry is to tear down houses.
because it was flooded and they need to tear down the walls..
All I do is need a sledgehammer..
It's not like rocket science..
I never majored in any sort of engineering..
I know how to break stuff..
It's so easy to break stuff, right?.
We're so natural..
Breaking things is so easy..
I've been breaking stuff..
I have a good timeline..
I can break things and not get in trouble..
All right, good..
But think about building something..
Building something, you really need to go to school for..
You really need to know what you're doing..

$^{601}$Your measurement has to be correct..
Even if you layer these bricks,.
these bricks have to be very careful in building..
Exactly what it is..
Satan, stronghold in our lives,.
we could destroy things so easily..
But maintaining this unity in the body of Christ is hard..
It's really hard, brothers and sisters..
But I want you to do one thing today..
I want you to think about behavior that you could improve.
and do better to maintain the unity within the church..
One thing, I can't tell you to do a million things..
I mean, we have a lot of work, but work on one thing..
What is one thing that you could do better.
to contribute to the church to maintain this unity?.
Hey, since we know we are God's children.
and we know what he did for us,.
we therefore live like a princess.
or a prince of the king of king..
Live like your truthful identity, brothers and sisters..
We have to choose to imitate God.
by walking in love towards each other,.
towards our fellow brothers and sisters in Christ.
with our speech and actions..
So we could maintain this unity.
and at the same time be in a manifestation.
of what the glory of God would look like.
when the church operates normally..
Let's pray..
Father, we thank you..
We thank you for the job of unity.
is already done through you..
Thank you for giving us some specific instruction.
of how we can continue to maintain this unity, Lord..
Father, as you gave us all the tools already,.
the spirit is within us..
Help us continue to tear down.
those strongholds from the devil..
Those footholds are very damaging, Lord..
I pray nothing more, Lord,.

$^{641}$that we do not demonize our brothers and sisters.
for any other reasons, Lord,.
because our main identity is in Christ,.
not because I am a Chicagoan.
or I am from Hong Kong, I'm Chinese, or whatever..
That's secondary..
Our citizen in heaven is primary, Lord,.
and help us hold dear to that,.
knowing that we might have differences,.
but help us love each other..
Help us have that sacrificial love, Lord,.
toward each other even though we might not have,.
we have different points of view, Lord..
Help us walk like Jesus, Lord..
Help us be like Jesus..
We thank you for today and the message..
In Jesus' name we pray, amen..
\newpage



\section{3:16-19}
\label{sec:w05Jl15bCEk}
\textbf{Why are we here? (John 3\_16-19) - Rev. Sam Johnson}
\newline
\newline
連結: \href{https://youtube.com/watch?v=w05Jl15bCEk}{\texttt{https://youtube.com/watch?v=w05Jl15bCEk}} ~~~~ 語音日期: 2022-06-28
\newline
\newline
\hyperref[sec:v_Mfq956F90]{\small{< < < PREV SERMON < < <}}
~
\hyperref[sec:index_chronic]{\small{[返順時目]}}
~
\hyperref[sec:index_scriptual]{\small{[返順卷目]}}
~
\hyperref[sec:GttbrJ8u_y4]{\small{> > > NEXT SERMON > > >}}
\newline
\newline
$^{1}$Thank you, I'm excited to be here today..
Always love being here at UV and enjoyed the worship..
Thank you to the singers and musicians who led us, that was really excellent..
And this morning we're going to be talking about why we're here in general, why we're.
here..
Not just the big existential question, we're not really going to address that..
That's a great question and just as a spoiler, the Bible has some really good answers to.
that question as to why we're here on earth, why we exist, what our purpose is..
We're going to talk more specifically about why we're here today on a Sunday, why we're.
in church, why are we here in church, why do we gather together, why do we do this and.
what some of the reasons are for that..
And we're going to look in the third chapter of John..
As we, before we read this, I'm going to remind us what's going on in the third chapter of.
John, in the third chapter of John, Nicodemus, who's a Pharisee, he's a part of the nobility.
and also the religious leadership in Jerusalem..
He comes to Jesus and he comes at night because he doesn't want anyone to know that he's talking.
to Jesus because Jesus is not in favor with the religious leaders, with the Pharisees..
And he comes to Jesus to ask him these spiritual questions..
So Jesus begins to explain to him about salvation, about how it's the new birth, it's like being.
born again spiritually..
Jesus talks about the Holy Spirit, he talks about the kingdom of God, he explains some.
really deep stuff to Nicodemus..
And in the middle of this passage, there is John 3, 16, right, the verse that we all know.
so well, where Jesus is talking some about the gospel and God's love for us..
And so that's the, we're going to look at John 3, 16 to 19 today, these verses that.
are right in the middle of this passage, that's the context of what Jesus is talking about..
But we know this as kind of a summation of the gospel and that's exactly what it is..
But today we're not just going to look at that piece, we're going to look at what it.
means for us in the church, John 3, 16 to 19..
So I'm going to pray real quickly and then we'll read these verses together..
God, I just pray that you would be present with us this morning, that you would, through.
the power of the Holy Spirit, really illuminate the scriptures to our hearts, that you would.
speak to us and that we would hear what you have to say, that people would not hear my.
words or my ideas, but that we would hear from you this morning and that your word would.
shape us and change us like it's supposed to and that we would go out from here or from.
wherever we are watching this and be different, be changed and be more like who you want us.
to be, God, in Jesus' name, amen..
All right, so let's read these verses..
John 3, 16 to 19, it says, "For God so loved the world that he gave his only son, that.
whoever believes in him should not perish but have eternal life..

$^{41}$For God did not send his son into the world to condemn the world, but in order that the.
world might be saved through him..
Whoever believes in him is not condemned, but whoever does not believe is condemned.
already because he has not believed in the name of the only son of God..
And this is the judgment, the light has come into the world and people love the darkness.
rather than the light because their works were evil.".
Now, this is an important question for us to consider, why are we here?.
Why do we come to church?.
Why is it a part of our lives?.
And if you are not familiar with church or you're not familiar with faith, just initially.
it seems like a strange proposition, right?.
That there would be a day of the week where we would get up and go to a lot of effort,.
you know, we take a shower, get dressed, do all those things and get up early, get in.
our car and go to a place and gather with other people..
And honestly, sometimes even the people we gather with, to outsiders it doesn't make.
that much sense, right?.
Because we, a lot of times we're very different..
A lot of times our church gatherings we gather with people who come from different backgrounds.
and different experiences, even different cultures..
And so, why do we do that?.
Why is it important?.
On the way here this morning, I was a few minutes late getting here because I drove.
past Sauk Stadium, it's right down the street..
And there was a game today, I'd forgotten they had a home game today, that's why I was.
a few minutes late, because I got caught in some of that traffic..
But that's something that we all understand, like why do you go to a baseball game?.
Well, because it's fun and you know, they have overpriced food that's really good and.
so that's why you go, you go to enjoy it and have a good time..
But why are we here?.
Why do we do this?.
Or why would you even, if you're part of the virtual audience, why would you say, "Hey,.
I'm gonna make an appointment in my week, set aside some time to make sure that I tune.
in and I listen to some brothers and sisters in Christ worshipping God and I worship with.
them and I listen to the word of God being talked about.".
Why do that?.
It doesn't make a lot of sense on the face of it..
And one of the things that happens sometimes in church is we just kind of take that question.
for granted..
We don't really think about it or think about the answers to it, we just take it for granted.
and we just show up without ever asking it or we just engage in church..

$^{81}$If you're watching the live stream, maybe just engage without ever asking it..
Until the day that we find out we don't really have a good answer and then we just stop showing.
up..
There's no reason to..
And we know that that actually happens..
It happens to people in the church..
The data is really clear on that..
In the decade prior to the COVID-19 pandemic from 2009 to 2019, Pew Research reported that.
church attendance dropped in North America, in the US, by 7\% during that decade..
And the Barna Research Group, which is another group that does research amongst church goers,.
they reported that during the pandemic, 32\%, a third of people who identified as practicing.
Christians, so these were people who self-identified as serious about faith and about church, related.
that they had stopped attending church either in person or virtually, stopped totally..
They weren't going in person, but they also weren't connected virtually..
They were just not engaged with church..
And these statistics are even more stark for young people..
In 2019, the Barna Group published research that said that among 18 to 29-year-olds who.
grew up in church, 64\% were withdrawing from church involvement as an adult after being.
active as a child or a teen..
So clearly, this is an issue, at least in the US and North America, that people at some.
point decide there's not really a reason for me to be there..
I'm not sure why I'm there..
And so they stop showing up..
They stop being engaged..
And that's why it's really important for us to consider this question thoughtfully and.
biblically, and that's what we're going to try to do this morning..
And I want to start by being honest and saying that there are a lot of problems in church,.
and we don't do well to try to ignore those or to try to cover them over..
In our modern church, especially here in North America, there's a lot of problems..
There's a lot of mess..
I mean, there's a lot of good reasons to want to disengage and to say, "I'm not sure that.
I want to be there because of some of the things that go on.".
At times, instead of reflecting the character of Jesus to the world, the church kind of.
stinks, kind of smells bad of selfishness and hypocrisy or dysfunction and conflict.
or legalism or materialism and chasing after money, even abuse..
Those are things that happen in the church..
They shouldn't, and they're dishonoring to God and they're heartbreaking, but they do..
And so we understand why there might be moments for people in church and for us to say, "I.
kind of want to disengage because of what I'm seeing.".
We can all find reasons to withdraw..

$^{121}$So what reasons are good enough to keep us here?.
Are there reasons that are good enough to keep us here?.
As we think about this initially, you know, a lot of stuff might come to mind, and some.
of those reasons might be partial answers to the question..
Some of them are a part of the true answer to why we're here, but some of them are just.
incorrect and wrong..
There's a lot of reasons that we come that may not be healthy, and some of them are tied.
to wrong assumptions about what church is..
And some of those things that we might initially think of, you know, there's obviously sociological.
reasons that we can attend church..
Experts even who are not believers agree that it's good for our mental health and our sense.
of well-being and purpose in life to have kind of the purpose that faith gives us and.
to find community in church..
So there's sociological reasons..
There's emotional reasons..
Again, according to the Barna Group, most people that regularly attend church report.
that they experience enjoyment at church and cite this as a main reason for attending..
Church can make us feel good and help us to get through the week, and we can enjoy the.
worship or we can be entertained by things that go on at church..
Or maybe we just really like to see our friends..
We get excited about seeing friends and acquaintances that we have at church..
Well, those things are wonderful..
Those are great things, but that's not the primary reason we're here..
That's not the heart of why we come..
Then there are some of us who engage with church, whether it's in person or virtually,.
we do it just because it's a habit, right?.
We've always done it..
It's a part of our week..
It's a part of our routine..
We just do it out of habit..
At times, some of us probably have participated in church out of a sense of obligation..
I grew up in church, so that's been true for me at different times, especially as a young.
person..
I did it because I thought I was supposed to, and I thought it was expected of me, and.
perhaps that's one of the reasons that we participate..
We feel like our parents or our family members or other friends or people around us or even.
the leaders of the church are really expecting us to be a part, and so we do it out of obligation..
And there can be cultural reasons, too..
Sometimes we feel like that it's just what we do, what people of our cultural class,.
social class, our economic class or our group, that's what we do..

$^{161}$Even some people are engaged in church because of their political affiliation, not because.
they know Jesus, but just because they think it's what you do when you're a certain type.
of person in society..
And then there's spiritual reasons, too..
We can participate in church because we're trying to justify ourselves in some way, because.
we're trying to show that we're a good person, and maybe we're trying to convince ourselves.
of that..
Maybe we're trying to, we know some stuff's not right in our life, and so we're trying.
to be like, "Well, I want to be good, and so I guess that means I go to church.".
Or maybe we're trying to convince someone else..
Maybe we're trying to show other people that we're a good person..
We're like, "Yeah, I'm someone who goes to church..
I mean, I'm a solid person..
I'm a churchgoer.".
Or maybe we're even trying to convince God..
Maybe we're trying to earn something with God by showing up at church, and by saying,.
"Yeah, God, I'm a good person..
I've been going to church, so shouldn't you answer my prayers, or shouldn't you care about.
me more, because I've been trying to do this.".
The reality is none of these reasons are really adequate to explain why we're here, and why.
we should be here, or to keep us here in the future..
Some of them are true and helpful, like the fact that we can find enjoyment, and encouragement,.
and support, and community..
That's true, and that's a part of the true reason we're here, and that's a beautiful.
thing..
But some of them are wrong..
If we think that going to church justifies us before God or anyone else, that's wrong..
That's just not true..
And even the best of these answers are really just byproducts of the true answer..
The enjoyment, the blessing, the community we find, that's really a byproduct of the.
true answer..
It's not the answer in full..
So to really find the ultimate answers, we need to go back to God's Word and look at.
this..
That's what we're going to do together this morning..
And here in verse 16, we see the first part of this answer for us..
We're here because God is there..
God exists..
That phrase is something that the late theologian Francis Schaeffer used to say, I think he.
actually taught a book, "The God Who Is There," and he spoke about the God of the Bible as.

$^{201}$the God who is there, the God who exists..
And this is really big and foundational..
This passage in verse 16 begins with the words, "For God, for God so loved the world, for.
God.".
And the existence of God is assumed as an established fact here..
And this truth is an essential part of why the church exists and why we're gathered here.
today because God exists, God is there..
If we're unsure about this, then it's certainly easy to be confused about the need to be here,.
to show up, to be connected to a church or the need to gather together if we're not sure.
if God exists..
The whole Bible begins with a statement about God..
It says, "In the beginning, God created the heavens and the earth.".
And from this verse and a lot of other verses throughout scripture, it's clear that God.
is the central character of the narrative of scripture..
He's also the creator, he's the first cause of the universe, and ultimately that tells.
us that this is God's universe and this is his earth and this is his story..
He created us in love and he has lovingly allowed us to take part in his story, but.
we're taking part in God's narrative, not the other way around..
God's not a little piece that gets added on to our narrative..
That's usually how we view it in life, that we're, it's about our narrative and God, we.
ask God to come take part in what we're doing..
But it's actually the opposite, that's true..
This universe is God's narrative and out of God's love, he allows us to take part in his.
story..
But his existence is the bedrock of reality in this universe and it defines everything.
else..
It defines our existence..
We exist because God exists..
So I want us to think about this for a minute here, spend a little extra time talking about.
this because God's existence is really foundational to a lot of the other reasons that we're here,.
that we gather together..
And it's also something that's really attacked in our modern culture and a lot of us have.
probably experienced moments of doubt or skepticism or wondering or questions about the existence.
of God or his nature and that's okay..
Honestly, that's a healthy thing for us to do..
Acknowledging those doubts and those questions allows us to wrestle with them and to take.
them to God, to examine God's word and to actually find answers..
So that's really healthy for us..
If we never acknowledge those doubts and those questions, then the questions never get answered.
and eventually we would just walk away assuming that there are no answers to those questions.

$^{241}$and that's not true..
There are..
God has answers for us and so it's healthy for us to wrestle with those..
So if you're someone today who says, "I'm kind of wrestling with that a little bit,".
with is God really there?.
Does he exist?.
Is he real?.
What is his nature like?.
I want us to just talk briefly about this, okay?.
And I want to point out that if we're struggling with the existence of God and we're saying.
maybe God isn't there, then there are actually three really big questions that we have to.
answer if we're trying to escape the truth of God's existence..
And the first one is if God isn't there, if there is no God, then why are we here?.
How do we explain the presence of really anything else in the universe?.
There's really not a naturalistic evidence-based explanation for the beginning of the universe.
because even though we can say, "Well, we can look at evidence from the very early history.
of the universe," we have no evidence of what came before..
It's a physical impossibility for us to know what came before because there was nothing.
before the universe existed..
And so even all of the best explanations of the best scientific minds in our world when.
it comes to how the universe came into existence go into metaphysical territory..
People talk about maybe there was oscillating universes or some type of multiverses or all.
these things like that, but those are metaphysical explanations..
There's no evidence, no scientific evidence we can have that any of that's true..
So even the best scientists and the smartest people in the world can't say what began the.
universe, what started the whole thing..
The Bible tells us that God did..
And there are other good reasons to believe that what the Bible says is true, but the.
Bible tells us that God did..
So that's something that points to the existence of God..
It's a proof of God's existence, the fact that the universe is here itself, the fact.
that we exist and everything else that we see is here..
If God did not exist, where did those things come from?.
The other thing that we have to think about is if God doesn't exist, where do we get our.
innate sense of morality of right and wrong?.
And why is it that us as humans, even those of us who claim no religious affiliation or.
belief at all, are still really convinced about certain things that they believe to.
be right or wrong, certain things that they believe to be just and unjust..
All of us as humans have an innate sense of morality..
It's shaped by our experiences and it's different for different people, but everyone has that.

$^{281}$sense..
And again, where does that come from?.
We can't explain that without the existence of God..
We can't just say, "Well, it's there because it's something that humans developed over.
time to help their species survive.".
Well, that's really not a good explanation because a lot of what we believe about justice.
and injustice and a lot of our innate sense of morality transcends our desire, our instinct.
for survival..
Throughout human history, people have been willing to give up their life for something.
that they believed was right..
And so it's not just something that we have developed in order to survive as a species,.
it's deeper than that..
It's more important to us than that..
And it's there in all of us..
We all believe some things are right and wrong, even if we don't believe in God or any form.
of faith..
And the third thing that we have to think about that shows us that God exists is just.
all of the order and complexity and intelligence that we see in our universe..
Our universe is full of just amazingly complex systems, from the macro on the galactic structure.
level all the way down to the micro, our cells and our DNA..
And we've seen over the last few years how chaotic things can become, even because of.
microorganisms, right?.
There's a lot more complexity and a lot more intricate, amazing intelligence involved and.
information involved, even on a cellular level, than any of us can even kind of wrap our minds.
around..
Where does all that come from?.
Again, this is something that points us to the fact that God exists..
Now we might say, "Okay, well, those things do point us to the fact that God exists, but.
how do we know that the God of the Bible is the God that we're talking about?".
And one of the things I would point us to there is that the Bible actually is provable..
The Bible talks about things that we can investigate..
The Bible has historical and archaeological things that we can look at and go back and.
say, "Okay, this person was talked about..
Was this a real person?.
This place, was this a real place?".
And we see that as we do that, the Bible is corroborated by those sources of evidence,.
by history, by archaeology..
The Bible has prophecies about people that were, where things were predicted in the future,.
and then those people actually lived and did those things..
Those things came to pass..

$^{321}$The Bible's internally consistent..
It agrees with itself..
Obviously there are lots of things that are difficult to understand in the Bible that.
we have to study out, but the Bible has an internally consistent voice and witness to.
us..
The text is consistent..
Over time, it has stayed consistent..
Honestly, the existence of the church throughout church history is one of the proofs that the.
Bible is true..
The fact that the church has continued to persevere and thrive in the space of so much.
opposition throughout history..
You may say, "Okay, those things are helpful, but I'm still not, I don't know, you know,.
I kind of want like a hundred percent proof that God is there.".
If that's what you're seeking, then I would say, "I don't think you're going to find that.".
And here's why..
Because the God who reveals himself in the Bible is a God of love, not a God of force.
and coercion..
God invites us into a love relationship with him..
He doesn't say, "Hey, I'm going to give you so many, so much evidence, and I'm going to.
beat you over the head with it until the only thing you can do is choose to be in a relationship.
with me.".
God doesn't do that..
God doesn't force anyone to believe in him..
But he says, "Hey, I'm going to give you enough to quiet your fears, and then I'm going to.
invite you to trust me and to love me.".
That's the God who is there..
And ultimately, that's why we're here today..
We're not here in church just because it's a fairy tale, but we keep showing up because.
we haven't figured it out..
We're not here in church just because it's good for our mental health to be around each.
other and to believe in something..
That's great that it is..
It's a byproduct of the truth..
But we're here because this is real, because God is real..
He really exists..
He's really there..
And that's the foundation of why we're here, because God is there..
Next thing we see in this passage is that God loves us..
We're here because God loves us..
As we read further in verse 16, Jesus, and remember, this is Jesus talking to Nicodemus..

$^{361}$So when we're reading verses 16 to 19, Jesus himself is speaking these words..
Jesus tells us that not only does God exist, but also that God is a loving God who loves.
us specifically..
And this is such wonderful news..
Just knowing that there's an all-powerful creator God, even if we believe that, that's.
not so much of a comfort, because what if God is angry or mean or vindictive?.
But the Bible tells us that not only is God there, but he also loves us..
And he loves specifically..
And what I mean by that is that God doesn't just love the earth, love us, love everybody..
He does..
He loves globally..
But he also loves specifically..
He loves you..
He loves me..
And this really has enormous implications for our worth and our value as humans, our.
lives..
We are created in God's image, but we're objects of his love..
The fact that we're valued by the most powerful being in the universe means we have worth.
and value that no one else or no thing else can take away..
Our lives matter and God loves us..
And that's true for each one of us, because God loves you specifically, not just globally..
His love is directed towards you..
And as Pastor Mark was mentioning earlier, when we go through those challenges of life,.
those uncertainties, those job challenges or trouble in school or things that we can't.
figure out or issues in a relationship with someone else and we're hurt or we're confused.
or we don't know what to do, we need to know that God loves us and loves us specifically.
and cares about our lives..
And we see that that is why we are here today in church..
We're not just here out of habit or routine..
We're not just here because we don't know what else to do on Sunday morning..
We're here because God loves us..
He loves you..
He loves me..
We're objects of His love..
And that's why we show up..
That's why we're here..
Next we see in this verse that God has revealed Himself to us..
God has spoken..
We see this in verse 19..
It tells us that the light has come into the world..

$^{401}$The statement makes it clear that God has shined on us..
He's revealed Himself to us..
God has revealed truths about Himself, about the nature of reality, about the significance.
of what it means to be a human, and about how He wants us to live in relationship with.
Him and with other people, with the people around us..
God has spoken and revealed that stuff to us..
He's revealed these things in a long love letter that He wrote to us through human authors.
that we call the Bible or the Word of God, and we call it that because it is spoken by.
God to us..
And He has told us who He is..
He's told us who we are..
He's told us what is good and what is true in life, what is false and evil..
He's described to us His good and beautiful design for our world and for our lives in.
it, and He's invited us to live in right relationship with Him and with the people around us..
He's described what that relationship looks like..
He's been honest with us..
He's assured us that if we follow His good and beautiful path for our lives, we'll experience.
deep fulfillment and health in our relationship to Him and others, even though we'll still.
live in a broken and fallen and at times painful world..
But He's also told us that if we reject His overtures of relationship, we'll experience.
separation, and that this separation will lead to unhealthy, self-destructive behavior.
that's damaging to us and the people around us..
God's been honest..
He's told us the truth..
He's revealed that His Word is vibrant and alive with contemporary truth..
It means something for us today..
It's without error, and that we need to interact with it..
We need to read it..
We need to study it..
We need to seek to understand what it means, even across the gaps of language and culture.
and history that exists between us in this current moment and the time that it was written..
That's why it requires listening and study..
That's why He's instructed us to talk to each other about it and to listen to it, proclaim.
and explain through preaching..
And then we're supposed to be changed and shaped and formed by it..
We're supposed to align our lives to God's truth as it's revealed to us..
We're supposed to let God's truth and God's Word change us, and that's why we're here,.
because God has revealed Himself..
He has spoken..
He has something to say to us, and it's significant, and we're supposed to be formed by it..

$^{441}$The next thing we see in these verses is that we're here because there are heavy realities.
in the universe..
This passage reminds us of that..
Verse 16 tells us there's a very real danger that we could be in a state of separation.
from God and perish, as the word that's used in that verse, for eternity..
The relationship between humans and God was broken by our rejection of Him, and because.
of that, both spiritual and physical death are realities..
That was the thing that God told the first humans in the garden was, "Hey, if you break.
our relationship, there will be death..
Death will become a reality.".
And since then, it has been a reality, both physically, that none of us live forever,.
but also spiritually, that we can be separated from God..
That's a heavy reality..
Verses 18 and 19 in this passage make it clear that good and evil are real..
Darkness and light exist in our world, and God's one day going to judge between good.
and evil in our world, in our lives, and in our own actions..
Now sometimes that's comforting when we think about God judging evil in someone else's life,.
but that's heavy when I think about it in my own life, right?.
It's a heavy reality..
It's a sobering truth..
That means the choices that we make in life do matter, and our decisions, they either.
move us towards God or away from God, towards what He's spoken and revealed that is good.
and right, or away from those things, towards life or away from life, towards deeper relationship.
with Him or towards separation..
That's what God has revealed to us, and that's heavy..
And that's one of the reasons we're here in church, not just for cultural reasons, certainly.
not for political reasons..
We're here because this stuff is real and it matters..
These realities are heavy, and we should care about them, and we do, and that's why we take.
part in church, and we gather together, and we listen to God's Word, because the choices.
that we make have consequences, because there is real good and evil in our world..
Life's not just some type of game that's a collection of inconsequential choices..
These things really matter..
Next thing we see in this verse is that we're here because of Jesus..
God sent a rescue in verses 16 and 17 to tell us that God sent His Son into the world in.
order that the world might be saved through Him, that whosoever believes in Him should.
not perish but have eternal life..
Jesus came to rescue us, to save us from sin and death, to forgive our past sins, to reconcile.
us and bring us back into right relationship with God..
And verse 16 tells us that God provided that rescue because of the greatness of His love.

$^{481}$for us..
God loved us so much, and that's why He sent the rescue..
And this is really, really good news, and that's part of the reason, that's the biggest.
part of the reason that we are here today, because Jesus, God sent Him to rescue us,.
to save us..
Even though we were in a state of continuing to reject God, we were saying, "God, I'm not.
interested, I want to continue to go my own way, and I want to continue to live in selfish.
and self-destructive ways," but God said, "No, I love you anyway, and I'm going to reach.
into your world and rescue you.".
It's an act of God's grace, and that word means unmerited favor, it means we didn't.
deserve it, but God sent the rescue anyway..
We don't perform work to earn it, to earn the rescue that God sent, we simply accept.
it by faith and receive it, receive God's grace..
And that tells us something really big about why we're here in church..
It's not because we're trying to earn favor with God, because we can't, and it's not because.
we're good people..
Sometimes people inside and outside church think that people come to church because they're.
good, and that's not true..
People even say like, "Oh, he's a good guy, he goes to church," or "She's a good girl,.
she's there every week," right?.
No, we don't come here because we're good, we come here because we're not, we come here.
because we're imperfect and we know that, and we need Jesus and we need rescue, we need.
saving, and that's why we're here..
Because we need to be changed, we need help, and because of God's grace and mercy through.
Christ that we are saved and rescued and redeemed and forgiven and healed and reconciled, and.
that's why we're here..
We're here because we're made new in and through Jesus..
If we accept God's offer of rescue, that gift, we're restored to right relationship with.
God, and then God begins to work in our heart, and we begin to change..
Verse 16 tells us that we will have life, and this reminds us that Jesus said that he.
came so that we could have life more abundant and fulfilling and full..
God's word tells us in other parts of the scripture that in Christ we are made into.
a new creation, and that's a huge part of why we're here, because we want to be made.
new..
Apart from Christ, before relationship with God, most of us were probably selfish and.
had a lot of things that were unhealthy in our lives..
I know that's true for me..
Now that I know Jesus, I'm still selfish and have a lot of unhealthy things in my life,.
but God is at work slowly stripping some of those away and helping me to be more like.
Christ, helping me to love others more, helping me to love my wife more and love my children.

$^{521}$more, helping me to love my community and my neighbors more, helping me to care for.
others in a deeper way, helping me to learn to pursue after and connect with God in a.
different way..
We call it in church, we call that process sanctification..
It's just a church word that means that God is changing you and making you more like Jesus,.
and that's why we're here, because we want to be made new..
We want to hear God's word proclaimed to our hearts..
We want to change from our selfishness and all the mess that we used to be into a reflection.
of Jesus..
That's why we're here, because we're made new in Christ..
And we're here because we're responding to God..
In verse 18, we see that God is calling us to respond to his love and his rescue with.
belief or faith..
Whoever believes in him is not condemned..
God's calling us to respond with belief..
And the Bible is clear that God wants us to respond to his love..
We respond to his love as we return it by obeying the truth that he's revealed to us..
We return his love and respond to his love by making God the object of our praise and.
worshiping him, something that we did this morning, something that we do in church..
We respond to his love in community..
God's made it clear that we need the support and encouragement and the help of one another.
and a group to get through life, and certainly to get through life as a follower of Jesus.
is not easy..
There's a lot of challenges..
There's a lot of setbacks..
There's a lot of things that hurt us..
We need each other..
And so we respond to God's love by saying, "Okay, I'm going to love the brothers and.
sisters that God has put around me.".
We support and care and pray for one another, even as we struggle with the painful realities.
of living in the world that we live in, where things are broken and where there's conflict.
and there's hatred and there's all these problems..
We respond to God and his love as we serve side by side and we go out and serve in our.
families but also in our communities, as we reach out to our neighbors, as we practice.
hospitality, as we do the things that the people of God have been known for throughout.
the history of the church, caring for the poor and needy and caring for those who are.
suffering and those who are vulnerable and those who are weak..
Those are all things that we do as a response to God's love..
And we do those together in a group with other Jesus followers..
That's why we're here, because we're responding..

$^{561}$We're not doing this to show off or to get anything..
We're doing it because we're responding to what God is already doing in our life..
And so we're loving each other, we're caring for each other, we're really listening to.
each other, we're supporting, we're praying, we're going out together to serve and to care.
and to try to impact the world around us..
We're telling the story of Jesus and the good news of how he can change lives..
Those are things that we're doing together because we're responding to God..
We're worshiping him, we're praising him together because we're responding to him..
So those are the reasons, the real reasons that we're here..
Like I said, it's not because it's a fairy tale, it's not because we don't have anything.
else to do or we're just used to it..
These things are real and God's clear about these things in his word..
And as we close, I just want to remind us that the history of the church is full of.
people who knew why they were here, who knew why they were a part of the church, why they.
were participating..
They were convinced of the realities of these things that we've discussed this morning..
They were certain of God's existence and of his revelation of truth in his word, the fact.
that God had spoken and revealed himself..
They were totally convinced of the gospel, the good news of God's love and grace and.
forgiveness through Jesus..
And because of this, they were consistently and faithfully engaged in church..
They were gathering together, preaching and discussing God's word together, fellowshipping.
with one another, worshipping God side by side, remembering Christ together, serving.
their neighbors together, sharing the good news with those outside the fellowship and.
caring for those in need..
And through the history of the church, people have withstood unspeakable persecution and.
life-altering challenges..
From the first century believers who were, because of their faith, were thrown before.
lions in the Coliseum, to believers in the last century who were imprisoned and tortured.
in oppressive regimes that didn't want people to believe in God, to believers that in this.
very moment are imprisoned or suffering consequences or even their lives are endangered because.
of their faith..
And throughout that history, there's this testimony of people all along the way, loving.
one another, caring for one another, loving their neighbors, serving the poor, taking.
care of the fatherless and doing all of these things because they knew why they were here..
They were certain of the truth of God's word..
And we may not share those experiences and that's okay if you're sitting there going,.
"Wow, I haven't gone through that stuff.".
That's okay..
I'm grateful for that too..

$^{601}$But we should share those convictions because we have the same reasons for being here that.
every other believer has throughout the history of the church..
It's not just a cultural thing, it's not a political thing, it's not just because we're.
trying to be good people or we're trying to show off, it's not because we're trying to.
earn something with God or it's just our Sunday morning habit..
It's not just because we enjoy it or it feels good, even though those are great blessings,.
it's not just because we want to see our friends, even though that's fantastic..
But it's because this is real, God is real, his love for us is real, his mercy and salvation.
and forgiveness, those things are all real..
It's because through Jesus we can really and actually be changed..
We don't have to be the people we used to be..
We can be something different..
We can love and obey God the way that we should and we can love and serve others the way that.
we should..
And that's why we're here..
And if that's why we're here, then that will keep us here through all the changes and transitions.
in life as we grow older and we go through different transitions in school and we're.
in college and then we're out of college and we're working and we have families and we.
have children and we are caring for relatives and life is crazy and we're moving and all.
that stuff will continue to be here because we know why we're here..
Because God is there, he really is at work in our lives, he really does love us and he.
really can change us and he has called us to be together..
Let's pray..
God thank you for these truths that we see in your word today..
God I pray that you would help each one of us to establish these things in our heart.
so that we know why we're here..
Help us not to be here for unhealthy reasons that won't really sustain us..
Help us to be convinced of the reality of your truth, of your existence, of your love,.
of your redemption, of the fact that you can change us and that you want us to come together.
to worship and to fellowship and to support and care and serve others together God..
Help us to be convinced of those things..
And give us perseverance in all the challenges of life to continue to follow after you, to.
continue to show up and to continue to really love one another side by side and to serve.
our communities and our neighborhoods and our city with the love of Christ..
God we thank you that you are here with us this morning and we thank you for our brothers.
and sisters in Christ who are here both in this room and in the virtual broadcast..
God we thank you for those who are here..
We pray that you would keep us here and continue to add to our number God..
In Jesus name, amen..
Thank you very much Stan for your good word..

$^{641}$God bless you..
\newpage



\section{約翰福音}
\label{sec:GttbrJ8u_y4}
\textbf{酒喇, 係咁先喇 (約翰福音 2: 1-11) - 陳德晶傳道}
\newline
\newline
連結: \href{https://youtube.com/watch?v=GttbrJ8u_y4}{\texttt{https://youtube.com/watch?v=GttbrJ8u\_y4}} ~~~~ 語音日期: 2022-06-28
\newline
\newline
\hyperref[sec:w05Jl15bCEk]{\small{< < < PREV SERMON < < <}}
~
\hyperref[sec:index_chronic]{\small{[返順時目]}}
~
\hyperref[sec:index_scriptual]{\small{[返順卷目]}}
~
\hyperref[sec:ufhLD6E1n5U]{\small{> > > NEXT SERMON > > >}}
\newline
\newline
$^{1}$大家好 大家午安.
不知道大家在公寓之後有什麼消遣.
我沒得出街 看多了電視劇.
發現這個世界各地很多人一放工的時候.
都很喜歡去喝兩杯.
疫情之下這個情況更加嚴重.
我留意到網上顯示.
英國,美國,台灣,香港都有關注.
疫情之下喝多了的情況.
其實說真的 這幾年.
朋友和我自己都是這樣 喝多了.
甚至去學一下品酒.
到底是因為大家壓門撩賴 沒得旅行.
還是借酒消愁 逃避現實 遺老送別.
我希望大家心中有數.
既然大家有一份興趣.
不如和大家說說酒吧.
要在這裡先提一提.
聖經裡面提醒我們不要醉酒.
不過是容許我們喝酒的.
很多處的經文都有提及.
其中有一段很有趣的.
就是《約翰福音》二章一至十一節.
「迦南婚宴」.
有趣之處是在這個婚宴裡面.
是「酒吧!先這樣吧!」.
事情發生在耶穌見完拿達葉之後的三天.
婚宴在耶穌的成長地拿撒勒.
北邊的十四里的一個地方.
叫做迦南去舉行.
耶穌和祂的母親和門徒都獲邀出席.
我們其實不太清楚.
為什麼這個主人家和耶穌有什麼關係.
或者會不會因為拿達葉是當地人.
所以請了那些門徒呢?.
這些我們都不知道.
不過我只知道一件事.
「沒有酒了!先這樣吧!」.
在現代的婚宴裡面.
沒有酒這個情況是不可能發生的.

$^{41}$更加是當時更加不可能的.
因為婚宴是擺放了七天.
必須要預備足夠的酒去宴客.
沒有酒.
不單止令親友掃興.
更加令主人家非常非常沒有勉氣.
會將一個本來是祝福的場面.
變成一個怨聲四起.
每個人都會跟你說.
「先走!先這樣吧!下次不知道還會不會來」.
不過希望不用再有婚宴.
可能瑪利亞在當中都不能忍心.
看到主人家沒有面子.
所以就跟耶穌說.
他們傾談了一點.
就吩咐傭人將水倒入石缸裡面.
水就變成酒.
這個神蹟的過程很簡單.
但你要知道約翰福音並不多這個神蹟.
約翰精挑細選了七個.
而水變酒卻是第一個.
你見到很有心思,很有義氣.
而11節指出的目的,意義.
是要顯出耶穌的榮耀.
而另一方面.
你見到約翰在結構裡面.
都有一個精心的佈局.
原來這裡是一個扇形結構.
顯現出酒裡面藏著不同的層次和韻味.
不如讓我們慢慢探索一下.
第一個層次我進入到最深入的層次.
第六節.
聖經說.
「照猶太人潔淨的規矩.
有六口的石缸擺在那裡.
每口可以乘兩至三桶水」.
這個扇形結構的中心.
通常很多人都會覺得是一個重點.
而這裡.
你見到很細緻的描寫.

$^{81}$那個缸.
但為何那麼重要?.
重要在哪裡?.
如果作為一個品酒的人.
懂一點釀酒.
我會說.
其實這裡是在交代釀酒的材料和過程.
不少人喜歡在網上學習浸煤酒.
你都會經歷過一件事.
看就容易.
做就很難.
我學品酒.
都試過自己去釀酒.
真的比想像中更難.
釀酒.
最基本.
你需要一些含糖分的原材料.
經過壓榨.
或者一些蒸煮.
然後加入一些酵母.
釀去發酵.
而發酵需要一個時間.
不是一個很短的時間.
需要一些長的時間.
更加需要很適合的環境去儲存.
慢慢去醞釀出來.
才可以釀出一瓶酒.
但耶穌在這裡.
你見到只有水.
即使缸裡面沒有天然酵母.
但他一點都沒有含糖分的原材料.
所以這次釀酒就好像.
創世時在水面之上.
無中生有的創造一樣.
再看.
如果從釀酒的角度.
適當的溫度,濕度,沒有菌的環境.
是一個不能或缺的地方.
即使猶太人認為石頭是潔淨的.
所以用石頭來造的缸.

$^{121}$是可以承載潔淨禮儀的水.
但大家都知道.
心中有數.
這六個水缸簡直是空虛雲動.
淤棉黑暗.
但耶穌在當中是怎樣?.
掌握著這個空間,這個環境.
還有.
發酵需要時間.
而且是長的時間.
但耶穌在這裡是怎樣?.
是傾刻之間完成了的.
根本沒有可能的.
但在耶穌的手中.
就是這樣發生的.
無中生有.
突破空間,時間.
你能否開始感受到.
這瓶酒有多好,有多榮耀?.
穩藏著的是耶穌那份神聖.
祂就是無中生有的全能,創造主.
我們再感受多一層的層次.
二章四至五節.
我們會發現第四節.
耶穌回應母親的說話都頗刻薄.
但大家要知道.
在原文裡面.
婦人其實是對女性的一份尊重.
一份尊稱.
就好像我們現在.
Dear Lady.
是一份尊稱.
耶穌在十字架之上.
都是這樣稱呼的.
所以我們知道.
不單是尊稱.
更加是用了神子的身份.
去說出來.
「我與你有什麼相干呢?」.
好像想扔下母親.

$^{161}$但這裡不是割蓆.
在原文直譯.
比較好的說法是.
「這事跟我和你有什麼關係呢?」.
耶穌在這裡.
不是對事對人.
只是在提醒母親.
焦點不在酒裡面.
而是在時間.
或更貼近原文.
在時辰.
在時辰尾道.
《張瀚呼音》裡面.
時辰尾道出現了九次.
指向十字架之上.
耶穌得榮耀的時間.
所以耶穌在這裡指出.
祂自己有權柄.
去請祂的母親.
要馴服祂的權柄.
跟祂的方法,時間.
第五節.
祂的母親不單不覺得冒犯.
更加好醒.
她聽得明白了.
所以她就馴服耶穌的權柄.
並且叫她的用人.
跟著耶穌的方法走.
第七至八節.
耶穌讓走師的方法技巧.
是多麼厲害.
完全意想不到.
很有榮耀在當中.
正如創世時一樣.
只是出口.
說有就有.
滿有權柄在當中.
你能否品到.
祂走的底蘊.
正正就是耶穌的神聖.

$^{201}$滿有權柄.
祂是滿有權柄的造物主.
我們再品多一層.
愛一層.
二章一至三節.
這裡說的.
在交代讓走的目的和心意.
我們知道.
耶穌不是基於母親的請求.
但祂為甚麼要讓呢?.
是一朝二節.
這裡重複了「筵席」這個字.
眾所皆知.
親臨的筵席.
送走給人.
就是想祝賀.
耶穌讓走.
送給人.
當中極具心意.
帶著滿有祝福.
事實上.
在舊約要說走.
正正是代表.
神親自臨到賜福人.
在《若爾書》三章十八節.
都會這樣說.
「到那日.
大山要流出甜酒.
小山要流出奶子.
猶大的溪河都有水流」.
是一份祝福的圖畫.
而你再看看那些酒缸.
在第六節.
你會看到有兩至三桶水.
在《聖經》裡面指出.
其實是二十至三十加侖的水.
如果一瓶紅酒有750ml.
我沒有計錯的話.
應該有606至909支酒.
多厲害.

$^{241}$你見到多厲害嗎?.
耶穌這份人情.
簡直是極大的豐富.
這樣去送給大家.
你感受到嗎?.
當你看到這麼豐富的時候.
你感受到嗎?.
你有沒有一種置身於花園.
在伊甸園創世時的感覺?.
正在經歷創造主不斷的賜福萬物.
更豐富的賜福比一.
看著形象.
新養眾多.
比你自養大地.
管理萬物.
賜予我們更多的食物.
你感受到嗎?.
在這杯酒裡.
那份味道.
正正充滿著耶穌.
賜福.
豐富賜福的創造主人.
耶穌所養的酒樓.
那份層次豐富.
韻味十足.
令人聯想創世的時候.
並且耶穌就是創造主.
無中生有.
沒有權柄.
為了要豐足的賜福給我們.
我們在美酒當前.
我們應該有什麼回應呢?.
我猜大家第一個反應.
就是很讚嘆.
但是你看看第十一節.
約翰指出.
門徒即時的反應.
就是相信.
相信耶穌就是創造主.
而相信往往都會和馴服放在一起.

$^{281}$應該這樣說.
相信很多時候.
在心裡面.
或者只是嘴上說.
我相信你.
但是如果馴服.
就是一個很誠實.
去反映真偽的行動.
而你看到七,八節.
這班傭人就完全去馴服.
他的行動很誠實地告訴你.
去馴服他.
我猜這件事情.
你想一想清楚都覺得奇怪.
不是很合常理.
耶穌這樣的做法.
這個方法是很不合理的.
如果我作比較.
我會質疑.
或者那些傭人.
也會在當中質疑.
不過在這裡.
約翰告訴我們.
無論怎樣.
他中間有沒有和耶穌討論.
最後他們選擇了馴服.
當他馴服的時候.
你就會看到.
耶穌就是滿有權柄的創造主.
而這個神蹟就是這樣發生的.
相信馴服之外.
相信有另一個同義詞.
就是「接受」.
你記得的.
《約翰福音》一章十二節說.
凡接待或新譯本的譯法就是「接受」.
就是信他明的人.
9-11節.
這裡的本言直.
虛顯出來.

$^{321}$就是甚麼都不知道.
所以說那些傭人為甚麼這麼慈善.
為甚麼新郎會這樣.
而這個新郎更特別.
現值的主角應該是新郎.
但這裡卻是無言以對.
好像傻傻的.
其實他們也在表達一件事.
其實那些酒是百百得來的.
他們沒有預計過.
沒有預算過.
是出乎意料的一份禮物.
所以我估計不用再多說下去.
你也知道的.
那個主人家一定會立即去接受.
讓所有人去享受這份豐富的祠福.
美酒就擺在面前.
我估計誰人見到這樣的情況.
這麼虛幻.
起碼六百多瓶的酒.
擺在你面前.
而且是好的酒.
你一定會相信的.
你一定會馴服的.
你一定會在當中去享受的.
不過這幾年.
在很多人的面前.
是甚麼都沒有的.
是黑暗.
是空虛混沌.
淤棉,黑暗.
你叫我怎樣相信?.
又為甚麼要用這樣的方法?.
你的時辰未到.
那何時到?.
你叫我怎樣唇氣?.
怎樣馴服?.
更加整體整個環境.
你只會見到有奏坐在當中.
我不知道哪裡有祝福.

$^{361}$你叫我怎樣去享受?.
其實面對大家的質疑.
我很坦白地說.
我不懂得.
我頂多只可以.
好像門徒一樣.
在旁邊慢慢觀察,脈觀.
正所謂作品,回說話.
我會脈觀一些作品.
看看迦南分院的畫.
也會聽聽廣東歌.
最近有一位廣東歌.
福音詩歌的重要人物離世.
他是翁惟雲.
他有一首歌.
令我悼念他的時候.
拿出來聽.
他的作品是這樣說的.
他說:在這裡天天享繁榮.
時民從前誰會預測.
靜悄悄三兩破雨雪.
勤勞做工不計雨或晴.
經過夏秋冬風季眾無情.
轉眼叢林猶言猶新.
在這裡感嘆天錯外.
原來根已種植土裡.
以前的香港.
起初的創造.
還是迦南分院裡的石缸.
其實我們都看到一樣東西.
都是一樣的.
什麼都沒有的.
在那個時間.
是不知道發生了什麼事.
有什麼變化在當中.
但是轉眼間.
就變成風俗的似乎.
聽姐妹.
當你這樣看的時候.
我自己看下去.

$^{401}$我真的心裡覺得.
我仍然要相信.
我相信這個創造主.
以前能夠做得到.
今天或將來都能夠做得到.
他那份祝福.
沒有離開過我們.
不過這首歌.
是九七之前的作品.
應該是八十年代的作品.
我想很難令我們有共鳴.
不如我們說一下酒吧.
如果一個超凡的釀酒師.
你都知道什麼叫超凡.
就是釀酒的方法.
或者那些技巧.
時間的掌握.
空間的掌握.
往往都令你覺得很荒謬.
是你猜不到的.
我比較熟悉的是清酒.
大家都會聽過二角三.
或者稱為23\%的精米步合.
這些酒.
這樣的標榜一定是好東西.
因為在表達.
其實在清酒裡面.
米用白米來釀製.
做原材料.
白米表面有蛋白質和雜質.
需要用一些打磨.
或者稱為精米.
精米越少留下.
代表酒越來越香.
有點像威士忌.
放得久就越好.
不過我試試.
你看看這瓶酒.
這瓶酒的標榜.
就是十角.

$^{441}$100\%精米步合.
如果我擺在你面前.
你會不會喝呢?.
但很好奇.
我看著這個介紹.
100\%.
我都很好奇.
到底行不行.
但我以前喝過這個酒造的酒.
我知道這個釀酒師很厲害.
很有熱誠.
他的方法.
他的技巧.
我永遠都覺得很讚嘆.
為什麼可以這樣?.
我見過的.
他用同一種米.
同一種釀製的方法.
但他只是把溫度改變.
他可以做出最dry和最甜的酒.
很離譜.
當我還在半信半疑.
我就嘗試去重複去釀了一瓶這樣的酒.
我就買了回來品嘗.
結果是怎樣?.
我猜你都猜到.
我好像進入了一個稻米田.
旁邊有些花香.
慢慢飄過來.
但最重的是一份清香的米味.
不知為何.
在我整個喉嚨.
那種餘韻.
不用吃飯.
很棒.
好像浸在了米水當中.
很棒.
我猜耶穌留下這麼多酒.
其實都可能想教我們一件事.
真正.

$^{481}$我們真的什麼都沒有.
或者我們真正覺得時間未到.
不斷去罵不同的人.
但是一杯酒.
讓我們深入了解一下.
很具體地讓我們知道.
原來去到這裡.
走入受賞.
受更受.
當然我們不要喝那麼多酒.
但若然你真的借酒消愁.
很愁的時候.
就一杯喜,兩杯慈.
但你不妨嘗試慢慢去嗅.
嗅一下.
然後嗅.
在當中有多少層次.
多少韻味.
同時你又要想想.
耶穌這個釀酒師.
他那份創造的奇妙.
他的技巧.
多出乎你的想像.
雖然乾.
但是又很順滑入口.
雖然苦.
但是又微微出現甜味.
有不同的香氣.
百花盛放在你整個人的身上.
那份餘韻十足.
令你感覺有一份幸福的感受.
見大家都好像差不多了.
我又望一下手錶.
是喝一杯的時候了.
先走了.
先這樣吧.
下次再見啦!.
謝謝大家收看,再見!.
\newpage



\section{希伯來書}
\label{sec:ufhLD6E1n5U}
\textbf{帶著信靠的順服 (希伯來書11:8-10) - 黃嘉漢牧師}
\newline
\newline
連結: \href{https://youtube.com/watch?v=ufhLD6E1n5U}{\texttt{https://youtube.com/watch?v=ufhLD6E1n5U}} ~~~~ 語音日期: 2022-07-06
\newline
\newline
\hyperref[sec:GttbrJ8u_y4]{\small{< < < PREV SERMON < < <}}
~
\hyperref[sec:index_chronic]{\small{[返順時目]}}
~
\hyperref[sec:index_scriptual]{\small{[返順卷目]}}
~
\hyperref[sec:ndrseOs8_m4]{\small{> > > NEXT SERMON > > >}}
\newline
\newline
$^{1}$各位成會的弟兄姊妹好.
我們今天可以透過互聯網.
和大家一起去聆聽上帝的話語.
學習上帝的話語.
記得在我初信主的時候.
有一個字常常在我一些基督徒的導師.
或者一些牧長身上會出現的.
那是什麼字呢?就是信服了.
對於一個初信主的少年人來說.
信服非常簡單.
就好像聽爸爸說.
聽媽媽說.
爸爸說怎樣就怎樣.
媽媽說怎樣就怎樣.
當然我是一個少年人.
當時的我小時候很簡單.
打下球 讀下書 和朋友一起.
是這麼簡單.
沒有什麼其他特別的事情.
我講的是三十多年前的我.
今天我已經是一個中年人了.
我的身份變了.
我面對的經歷的事情都不同了.
作為一個牧師.
大教會.
在人生的路上都經歷了三十多年.
作為一個基督徒.
面對著很多的事情.
有家庭 有工作.
有很多面對這個世界上各種的挑戰.
在這個世界上各種黑暗.
所有事情都在我面前.
我已經不是一個少年人一樣.
這麼單純 這麼簡單.
因為我們周遭的世界很複雜.
我所顧慮的事情都會多了很多.
當我有著這樣的身份.
再進入去想回馴服這個字的時候.
突然之間有些事情出現了.
有兩個問題在我腦海裡出現了.

$^{41}$第一個問題就是.
當然我都知道馴服神糾纏.
聽爸爸的話 聽媽媽的話.
聽他的話就是這樣.
但其實什麼叫真真正正馴服神呢.
在我們生活範疇裡.
怎樣才叫做馴服神呢.
我開始去問.
是不是馴服神就是等於.
背幾個金蓋就等於馴服神呢.
馴服神是不是等於回教會就等於馴服神呢.
究竟什麼叫做馴服神呢.
還有一樣東西就是.
不知道是我熟悉的傳統.
還是在我的經歷當中.
在我初順主義的時候.
覺得馴服這個字很沉重的.
好像就是無可奈何被迫要跟隨神.
就好像爸爸叫你做一樣東西.
你沒得不做.
你反抗後果會很嚴重.
好像馴服神是一種無可奈何的體會.
但三十幾年想住這麼久.
是不是真的這樣呢.
是不是呢.
我們一起看看聖經好不好.
如果可以的話.
我們一起翻開今天的聖經.
就是希伯龍書的第十一章.
第八至十字.
這段經文其實是寫給一班.
有猶太背景的基督徒.
讓他們在一個羅馬的.
大約在一個羅馬的位置的.
教會信徒當中.
給他們一些鼓勵和探望.
為什麼需要鼓勵和探望呢.
因為在他講的時候已經感受到.
外來都有很多的壓力.
外來就是羅馬的所謂的.

$^{81}$面信主義的幫人.
有些外邦人給他.
有些外邦政權給他一些壓力.
就算是猶太人.
他原本的弟兄.
他這個猶太教.
相信猶太教的人.
都因為他們的一班人信耶穌.
就把他們推起了.
很多的壓力.
他們很想放棄這個教信的.
在這樣的情況之下.
希伯龍書的作者.
有人說不如可以叫他.
希伯龍書的.
好像一個牧者那樣.
寫一封信.
傳給當時.
位處於羅馬一帶的教會.
當時沒有MP4.
沒有現在那樣的.
那麼先進的科技.
所以就有一篇講章.
他用每一個小朋友看.
大約十多二十人左右.
就讀出來.
就成為鼓勵.
就是今天的希伯龍書.
而在今天我們所讀的希伯龍書.
特別是希伯龍書的作者.
是透過第十一章.
引用了很多舊約斯汀的人物.
去鼓勵當時的讀者.
能夠對神有信心.
活出跟從神的生命.
然後我們透過今天的經文.
幫助我們去回想兩個問題.
第一個問題.
究竟什麼是推動信徒.
信服神的動力.

$^{121}$第二個問題.
我們要怎樣去信服神.
使這個信服.
不是成為一個強迫.
又或者是一個空洞的屬靈口號呢.
當我們看完經文的時候.
我們發現希伯龍書的作者.
是引用了亞伯拉罕.
作為一個信靠信服上帝的例子.
在經文中說.
亞伯拉罕因著信.
蒙召的時候就遵命出去.
往將來要得遺孽的地方去.
出去的時候還不至往那裡去.
經文所說的蒙召.
是說亞伯拉罕蒙上帝的呼召.
上帝要指示他一家人去到新的地方.
在《創世記》第十二章.
第一至第三節說得很清楚.
是這麼說的.
上帝叫亞伯拉罕.
當時還沒改名.
亞伯拉罕就說.
你要離開本地本族父家.
往我所指示你的地區.
我必叫你成為大國.
我必賜福給你.
叫你的名為大.
你也要叫別人得福.
為你祝福的.
我必賜福於他.
那救助你的.
我必救助他.
地上的萬族都要因你得福.
亞伯拉罕.
因著上帝的呼召.
他的一家就成為了一列.
萬族因他們從神而來的福.
而得到福氣的一群.
祝福他人的人.

$^{161}$他們擁有了一個新的身份.
和新的使命.
亞伯拉罕聽到神的呼召的時候.
他就用信心而靠神.
信服神.
就遵命地踏上這條信心的旅程.
就算一家都不知道自己將要去到哪裡的時候.
他們都願意跟隨神引領而往前走.
其實對於我們來說.
這已經是一個很難的事實.
明白的事實.
如果我今天我要去航空公司.
買一張機票.
要去一個地方.
櫃口的人一定會很自然地問我.
先生請問你打算去哪裡.
如果我告訴你.
我不知道.
你隨便跟我買吧.
那個人跟我買機票.
他都覺得很頭痛.
那你真的去哪裡.
又或者當我們去計劃.
我們自己的將來的時候.
譬如我想搬屋.
我都會認定.
我要住在哪一個區份.
認定了.
然後怎樣.
去找屋.
找了屋之後.
我就在這裡買屋或者租屋住.
怎會有一個人這樣.
走在這裡的屋地.
然後又去紮營.
這裡就是我家.
然後又走到別處.
又去紮營.
這裡就是我家.
對於我們來說.

$^{201}$這是匪夷所思的一種想法.
但是亞伯拉罕.
因著聽到上帝的聲音.
他信服上帝的聲音.
信服上帝的引領.
跟著他去做一些.
你和我有些想不通的事情.
從希伯來書的作者的角度來說.
他認為亞伯拉罕和他的家人.
可以去到地獄的那一天的時候.
真是神引導他們的.
而且在引導他們的時候.
他們還有一件事.
就是他們不投身當地的文化.
成為當地的人.
他們沒有接受到迦南的宗教.
他們沒有取迦南的女子為妻.
他們卻是要成為一個寄居者.
就好像在異地居住象棚一樣.
我稱他們為異鄉客.
即是說他們來到這個城市當中.
擁有著一個特殊的身份.
這個特殊的身份就是.
祝福他人的身份.
他們祝福他人的身份.
因著這個身份的時候.
來到一個新的地方.
他沒有放棄上帝給他的身份.
投入成為一個異族當中.
他要成為一個異族當中的異鄉客.
我們這樣說.
亞伯拉罕和他的子孫.
都沒有容許自己安頓在那些.
安慰世俗的城市.
繁華的城市.
大家還記得.
亞伯拉罕當初的侄子羅德跟著他.
最後羅德和亞伯拉罕分家去了哪裡.
索多瑪 蒙哥拉.
對於我們熟悉經歷的人來說.

$^{241}$知道這兩個是一個罪惡之城.
但是這樣.
它也是一個很繁華的城市.
你可以說.
亞伯拉罕沒有去到.
這個世俗的城市.
他卻是選擇等候.
那個上帝永恆的城.
臨在來到.
他等候的那個有根基城.
是神所經營的.
神所建造的.
你可以說.
他們雖然會在一個充滿著.
外邦神的迦南文化當中.
但是他們.
卻繼承著獨一真神呼召他們.
去領受從神而來的福氣.
又將這個福氣帶給他人的時候.
擁有出一個新的身份.
就是一個祝福他人的異鄉客的身份的時候.
所以他們願意離開一切世俗的大城.
卻一心等候神所設計.
所建造的城.
看到這裡我們有沒有發現到呢.
希伯來書的作者是透過.
引用亞伯拉罕的生平.
來嘗試提醒多其時的信徒.
他們要以亞伯拉罕作為榜樣.
亞伯拉罕是因著信教神的恩召.
人神的呼召.
而信服地活出一種在世而不屬世的生活.
好使自己經歷從神而來的福.
又將這份福氣祝福他人.
其實希伯來書的作者.
他這個提醒我們在其他新約的經文當中都找到.
這個就是什麼教導呢.
就是一個在世而不屬世的教導.
是在這裡的.
如果可以這樣說.

$^{281}$信徒是因著聽到神對他們的呼召.
呼召有幾個層面.
第一層的呼召就是神呼召我們得到救恩.
第二層呼召.
神呼召我們成為聖別成為神的子民.
熟齡的生命.
第三層的呼召就是我們一般所說的神.
在一些處境當中呼召我們成為某種的人.
為了成就某種的事.
第三層呼召可以包括.
我們可以是職場的身份.
或者是教育的身份.
又或者是可以在家庭當中.
作為父母作為兒女的身份.
這些身份是神給我們的.
每個人都有些不同的.
不會每個人都做牧師傳道.
也不會每個人都做基督徒會計師.
也不是每個人都會結婚生子.
但是因著有著上帝給我們的.
召命和身份的時候.
我們就按著這個召命.
去活出上帝的生命.
我這樣說的時候.
就好像不太有意義.
所以當上頭聽到神的召命的時候.
神召我們成為祂子民的時候.
當祂召我們的時候.
不單是神召我們出於祂的愛.
我們本為罪人.
但是神的愛的緣故.
呼召我們領受這個救恩.
成為祂的兒女.
蒙祂所愛.
也同時領受了祂要透過我們.
使萬人經過我們.
能夠體會到上帝的福氣.
那個使命.
使人能夠體會到神的愛.
正因為這樣.

$^{321}$我們願意信服神的召命.
這個召命因為出自於.
是神親自叫我們.
是一個愛的呼喚.
不單是叫我們接受這個救恩.
同時也透過我們的見證.
使世上的人.
能夠因著我們.
可以得到福音的好處.
所以我們生活的方向.
因著這樣的時候.
就不再安頓自己.
在世俗的思想上.
安頓的意思是什麼呢.
安心立命.
就是說我不是叫大家.
不要打工.
不要結婚.
不要帶孩子.
不要讀書.
我不是這個意思.
而是說.
如果我剛才所說的各種事情.
賺錢也好.
又賺健康也好.
這些都不是錯.
但是如果我們的心念.
只是為了追求這件事.
就以為得到的時候.
我的心安頓的時候.
這樣就錯了.
因為那些事物是短暫.
就好像索多瑪摩拉一樣.
可以沒有了.
但是上帝要我們.
不要將自己安頓在這些事情上.
也不要在世俗上的邪惡當中.
受到煎熬.
在這個世界上有邪惡的價值觀.
教義諫論就好像在彼得慈書面.

$^{361}$二章十一到十二節所說.
親愛的弟兄呀.
你們是黑女.
是忌拒敵.
不是以黑客眼.
我勸你們要禁戒肉體的私慾.
這私慾是與靈魂爭戰的.
你們在外邦人中.
應該品行端正.
叫那些誹謗你們是作惡的.
因看見你們的好行為.
便在監察的日子歸榮耀吉神.
當基督徒願意放下世俗的價值觀.
和不安頓自己.
在世俗當中所給予的權力.
金錢名譽和短暫的健康上.
來求生命的滿足的時候.
當我們願意不去跟這些事情的時候.
這些都是為什麼會發生.
因為我們願意聽到.
聽到耶穌基督對我們的呼召.
來跟從我們跟從耶穌.
等候上帝將來所經營所建造的城.
即是我們在世當中.
我們在異鄉下.
我們等待一個城.
地上我們有生活應然.
會過不同的生活.
但我們等待一個城.
是什麼城呢.
就是在世俗的21章1至4節.
因為我看希伯來書.
我看了很久.
我想了很久.
究竟一個城是指什麼.
當我再想想的時候.
可能在差不多時期.
或是在後一點時期.
有一個城.
真是在新的聖經裡說得很清楚.

$^{401}$就是在世俗剛才所說的.
21章1至4節.
在新天新地當中的新耶路撒冷.
經文說.
我又看見一個新天新地.
因為先前的天地已經過了.
海也不再有了.
我又看見新城.
聖城.
新耶路撒冷.
由神那裡從天而降.
預備好了就如新褲裝飾.
整齊等候丈夫.
我聽見有大聲音從寶座出來.
說:看那神的障目在人間.
他要與人同住.
他們要做他的子民.
神要親自與他們同在.
做他們的神.
神要擦去他們一切的眼淚.
不再有死亡.
也不再有悲傷.
悲哀.
哭耗.
疼痛.
因為以前的事都過去了.
這個聖經挺好的.
就是說基督徒在這個在世當中.
我們一直都等待什麼.
我們以為財產.
資源.
學業.
健康能夠帶給我們生命的豐盛.
到頭來我們都逃避不到悲哀.
哭耗和疼痛.
但是怎樣才得到.
從神而來這份福氣呢.
就是當我們自己聽到上帝的呼召.
領受了神而來這份恩典.
成為了兒女.

$^{441}$在我們的身份當中去見證上帝.
也將上帝給我們的福.
和人分享賜予的時候.
我們彼此都會經歷到那份.
充滿著喜樂.
沒有哀哭.
沒有疼痛的經歷.
當然.
我們不是每天都經歷到.
但是常過.
什麼時候會完完全全地經歷到呢.
就是新城耶路撒冷尾的時候.
我們每一個都能夠.
因著自作孽.
救恩而來.
能夠領受到.
所以其實我們在世的人.
我們在追求一個真正的家鄉.
是天上的家鄉.
就是我們能夠完全經歷到的福.
也是我們盼望別人.
都可以和我們一起經歷到.
這樣的一種福氣.
如果這樣說的時候.
原來我們.
在這一刻.
我們要去信服神.
不只是聽父親.
好像聽父母說這麼簡單.
我們信服神.
就是我們自己的心靈願意信服神.
也都願意將信服神.
將這個福氣帶給他人.
和信服神帶著盼望.
等候著這個新城來到.
在這段日子當中我們有所盼望.
所以對我們來說.
如果我們這樣相信的話.
而相信這是基於神的愛.
而不是一種勉強的神的愛.

$^{481}$呼召我們才可以去經歷這份福氣.
也都去信服在這個引導之下.
去將福氣帶給他人的時候.
其實這個影響了我們幾個同事的生活.
因為有三個應用.
是值得你和我今天去注意的.
首先.
如果我們知道.
我們在等候著一個.
神所建造的城.
新城耶路撒冷的時候.
我們就不會將現時眼中.
所看到的一切看得太重要.
當然我們會留心.
我們日常的生活.
所以我和大家說.
你仍然要用心讀書.
工作.
照顧自己的身體.
照顧自己的家庭.
做一個好的公民.
這些全部都不是錯的.
我們的工作.
家庭財富健康都是有需要的.
這些需要是真實的.
但是我們不會將自己.
安頓在這些事情上.
我們不會用這些事情來安身立命.
就是說這些事情.
我們知道有時會沒有了.
有時會過了.
但是我們知道.
有的感恩.
沒有的我們知道.
我們不是等候這些東西.
我們一生的目標與追求是什麼.
是追求上帝的福氣.
輪到在我們身上.
然後將這個福氣帶給他人.
正因為這樣追求的時候.

$^{521}$我們就不會將眼前的事情.
得得失失.
看得太重.
今天很多人.
看著眼前的得失.
心裡面很痛苦.
基督徒一樣.
但會想說.
我們在尋求的是誰.
我們在尋求的是上帝自己.
所以容許上帝提醒我們.
我們今天得與失.
重點是在於上帝.
而施下他的恩典福氣給我們.
而叫我們將這個福氣帶給他人.
同時我們是因著聽到神的呼召.
而領受到他的救恩訊息.
信靠神.
而且也因著他的呼召的緣故.
我們能夠活出經歷神福音的大能.
從而去信靠這個位愛我們的神.
因著這份信靠.
我們就有能力推動我們去信.
是愛推動我們去信服.
這是第二件我自己要留意的事情.
信服神.
不再是一種強迫我們自己做的決定.
而是因為神的恩典的呼召.
所產生出來的動力.
推動我們.
所以我們在祈禱當中的時候.
我們要學習尋求神的引導.
使我們能夠聽到神.
透過他的話語提醒和聲明的引導下.
去信服神.
我們要學習信服神的真言.
信服不再是一個強迫.
信服是上帝的愛.
所以我們要學習聆聽.
學習在我們這麼繁忙的生命當中.

$^{561}$上師靜下來.
在禱告裡面.
在神話語當中.
讓神的話語提醒我們.
叫我們知道我們今天所面對的事情.
有沒有上帝的指引.
其實話語師好像很抽象.
其實一點都不抽象.
當你將你心裡面的想法.
將你心裡面的願望.
或者將你心裡面的一些難處.
放在祈禱的時候.
我們千萬不要像水蛇出現一樣.
說一說然後啞門就收聲.
就好像遞一張清單給上帝.
拜託你幫我搞定它.
應該反過來.
當我們將心裡面的受苦.
或者心裡面的請求.
和神說完之後.
我們要說一句.
主呀 僕人在此.
我們敬聽.
我們在此僕人敬聽.
學習聆聽神.
學習從神的話語裡面.
讀經當中看回神的心意.
聖經裡面有很多神的心意寫在裡面.
神的願意透過他話語裡面提醒.
有些提醒是直接到不得了.
我們源源易得.
其實神的心意很清晰.
有時當我們去尋求的時候.
一起去祈禱的時候.
神會馬上將他的心意顯現出來.
亦都有弟子們互相鼓勵人士.
我的願意呢.
各位.
要活著一種.
忘福的生命.

$^{601}$我們需要打靠.
需要去聽神的愛.
和愛當中的挑戰.
當我們用信心信靠.
信徒神的時候.
我們不單止要放願意.
放低在世上的價值觀.
不再以價值觀成為.
我們生命的目標.
而且呢.
我們要看清楚.
那些都是短暫的.
先要知道.
這是短暫.
而且正因為知道短暫的時候.
我們更願意追求一個.
脫離罪惡的生活.
使我們的生命成為.
見證神的生命.
使別人能夠看到我們.
就見到上帝的福氣.
各位.
我們是帶著盼望等候.
神的神.
新城耶路撒冷的降臨.
我們帶著盼望.
為什麼這麼重要呢.
因為這個世界太黑暗.
因為很多的戰爭.
使人失去盼望.
人變得得過且過.
人也變得自私.
因為既然世界這麼邪惡.
我為什麼要做好呢.
但是如果告訴你.
你今天所活著的.
不是為了今時.
而是為了將來.
那個永恆的生命.
那個新城.

$^{641}$而那個新城已經說了.
是有盼望的時候.
那我們願不願意.
為著上帝開好事情的路呢.
有人說得很好.
黑暗最黑的那一刻.
就等於泥濘將要降臨.
地圖也是一樣.
在一個黑暗的邪惡面當.
我們要信心信靠.
用信服態度去見證神.
我們是神福音的見證人.
不單是用口來傳播音.
也透過生命.
去將神的福氣.
帶給他人.
好的行為.
好的見證.
是一個使人.
很具體看到福音的見證.
教會.
口要傳福音.
也要手去做福音的工作.
使人.
能夠得到.
看到福音的好處.
各位.
我們是否要因著信靠神的恩兆.
而信服地去活出.
在世而不在世的.
在世而不屬世的生活呢?.
我們要使自己和別人.
能夠受到.
神賜給我們的福氣.
請問.
我們你和我.
又會不會願意.
花時間學習聽神的聲音.
帶著信靠的心.
活出信服的生命.

$^{681}$使自己與別人.
都蒙受神賜給你的福呢?.
所以.
各位聽證會的朋友.
今天的訊息.
給了一個.
很簡單.
但有很大挑戰性的.
一個提醒.
我們願不願意.
因著信靠.
而去信服神.
因著信靠.
就去活出一種.
信服神的生命.
就是一種在世而不屬世.
離棄邪惡.
將上帝的福氣.
他的福音.
帶到你周遭身邊的人.
的一群.
信主.
以神賜福給大家.
我背台詞了.
\newpage



\section{}
\label{sec:ndrseOs8_m4}
\textbf{Apostles' Creed \#7 (Holy Spirit) - Rev. Mark Chan}
\newline
\newline
連結: \href{https://youtube.com/watch?v=ndrseOs8-m4}{\texttt{https://youtube.com/watch?v=ndrseOs8-m4}} ~~~~ 語音日期: 2022-07-19
\newline
\newline
\hyperref[sec:ufhLD6E1n5U]{\small{< < < PREV SERMON < < <}}
~
\hyperref[sec:index_chronic]{\small{[返順時目]}}
~
\hyperref[sec:index_scriptual]{\small{[返順卷目]}}
~
\hyperref[sec:8rUfl6H7c2c]{\small{> > > NEXT SERMON > > >}}
\newline
\newline
$^{1}$acceptable in your sight in Christ and pray. Well if you have your Bible or Bible app.
please turn to the first very first book of the Bible the book of Genesis we're.
going to continue our Apostles Creek series and we come to the third section.
of the Apostles Creek I believe in the Holy Spirit. Now the first first.
occurrence or first appearance of this spirit can be found in the first book of.
the Bible the book of Genesis. Are you ready? Genesis 1. In the beginning not in.
the PowerPoint but in your Bible. In the beginning God created the heavens and.
the earth now the earth was formless and empty darkness was over the surface of.
the deep and the Spirit of God was hovering over the waters. The Spirit of.
God is hovering over the water. That's the first encounter of the Holy Spirit. I.
think correctly translated to should be the Spirit of God instead of the wind of.
God. Hebrew sounds pretty interesting because the pronunciation.
gives you the meaning of the word sometimes. So in Hebrew spirit you have.
to repeat after me okay even you have your mask on even better if you have your mask on..
Spirit in Hebrew is "Ruach". Yeah exactly yeah exactly that is the Spirit that's the.
wind of God the life force of God. He was the first created agent of the entire.
universe. Last Tuesday there was a big news around the world. Have you ever.
heard of it? The James Webb Space Telescope sent back the first set of.
images from the deep space. You know how the story goes right in I think the end.
of 2021 NASA launched the James Webb Telescope into deep deep space and then.
travel for you know some time and then last Tuesday beaming back the first set.
of image and what we see in the images are lights probably the very first light.
since Big Bang. Those lights the James Webb Telescope captures probably.
travel already traveled 13 point some billion years to reach the telescope the.
camera so it was very very amazing it's just mind-blowing to think about how big.
the universe is. Hundred years ago we thought there was only one galaxy in the.
cosmos now we know there were so many so many so many others galaxies unnamed.
galaxies haven't seen before right this is how big and how vast the universe is.
and behind it is the Creator and behind it is the Spirit of God so we come to.
the third part the Holy Spirit so we have two two questions to tackle it.
says I believe in the Holy Spirit first question is who is the Holy Spirit.
second what does the Holy Spirit do in your life if you profess to be a.
question and you believe in God if you believe in Jesus the Holy Spirit is in.
you and it should be working in your life in your heart so first question is.
who is the Holy Spirit the Holy Spirit is the third person of the Trinity the.
third person of the Trinity now please understand the Bible is actually a.
library there were 66 books in this library the whole Bible is a huge.
library there are many different genres different types of literatures different.
stories and poetries and poem and genealogies and records of history and.
all that stuff all right so the Bible is actually a library from the first book.

$^{41}$Genesis to the last book that we have the book of Revelation in the New.
Testament all together there were so many different voices and as God.
unfolds his revelation to mankind to all of us it unfolds it progressively so in.
the Old Testament we don't know much about the Holy Spirit because it's not.
unfolded yet until we come to the New Testament and we see more clearly who.
Holy Spirit is all about and who he is now the clearest statement about the.
identity of Holy Spirit is found in John 14 as we read this passage I want you to.
answer this question how many persons how many persons are there in the.
passage right John 14 all right now now about the context okay what is what's.
going on here Jesus is speaking to his disciples it was the last Passover meal.
last dinner together the Jewish Passover before he goes to the cross and die on.
the cross for the sin of humankind and to give us the salvation right so before.
before the Good Friday he was speaking to the disciples around the table around.
the Passover table and he was about to leave and he was about to be crucified.
for our sins and will be arise again from the tomb and will be ascended to.
the Father reigning as the King of Kings so before all of this events Christ.
Jesus was speaking to his disciples teaching them what it means to follow.
him this is the last evening Jesus spent with disciples around table of Passover.
so preparing the disciples to continue his work once he left once he ascended.
to the Father and reign forever in heaven the work of mission has to be.
continued with his disciples his disciples need to pick up the baton and.
go on with what is left to be done about the mission of Jesus so 1415 here he.
says Jesus speaking to the disciples speaking to the church speaking to us if.
you love me you will keep my commandments and I will ask the Father.
and he will give you another helper to be with you forever even the spirit of.
truth whom the world cannot receive because it neither sees him nor knows.
him you know him for he dwells with you and will be in you now how many persons.
are there how many persons are there if you love me right speaking to the.
disciples me being Jesus himself right and then the Father and he says I will.
send another helper and other assistant so to speak and other helper to be with.
you at least three persons are being spoken off here well right minus the.
disciples so the three personas three persons here now first of all I want to.
clarify love in the Bible is usually not a feeling love in the Bible usually is.
not a sentiment where our cultural script is teaching us love is all about.
you know how you feel about yourself how you feel about your lover or somebody.
else right it's always the emotional kind of you know feeling that is.
defining love right but love is not generally not a sentiment in the Bible.
and what is it the love in the biblical culture is the term love is often used.
between a great king and a smaller king when they make covenant together when.
they make treaty to treaties together or diplomatic relationship they enter into.

$^{81}$that relationship that agreement together and they sign some contracts.
right there were some terms in the covenant there's some terms in the.
international treaty and then you get a copy I get a copy you put your copy in.
the temple I put my copy in my temple and then the terms of the covenant.
often uses the term love that you have to love me and I have to love you now.
think about in political terms right if I'm the United States enter into some.
sort of international diplomatic relationship with another country right.
they're using the terms love right a king is loving another king and then a.
king is loving back the other king you think love in this sort of relationship.
is about romantic feelings no right they never really fall in love with each.
other right but love means in this kind of context loyalty political alliance.
political allegiance it's kind of like when you when you sing national anthems.
right that's sort of you know sentiment it's a political allegiance so Jesus.
said if you love me doesn't mean oh I feel Jesus you know so warm and so.
comforting well of course that's one of the consequence right but mainly it's.
all about political alliance you are my ally and we're gonna work together you.
are my political partner and you got to pledge your allegiance to me and I'm.
faithful to you so if you love me you will keep my commandments and I will ask.
the Father he will give you an other helper to be with you forever now you.
notice there were at least three persons here and other helper speaking of the.
Spirit the Holy Spirit to be with you and the adjective describing the Spirit.
is Spirit of truth the biblical truth the gospel truth he will tells you he.
will tell you the truth of the gospel he will speak to you the truth of God.
whom the world cannot receive and he will remain with you and in you so here.
is one of the clearest statements about who the Holy Spirit is it's the third.
person of the Trinity with the Father and the Son and you all of a sudden you.
remember the song how great is our God Spirit comes second right right how does.
it go father spirit son right well you say that's not right right it should be.
father son spirit right shouldn't it what do you think did Chris Tomlin.
make a mistake no I think I think it's poetry right so.
spirit can come second in that song but in the Bible the unfolding revelation as.
the grand story unfolds we have the Father and the Son Jesus and then the.
Holy Spirit who came another clear statement is Matthew 28 Matthew 28 the.
last word Jesus spoke to his disciples before he goes up into heaven in.
Ascension 28 18 Jesus came and said to the disciples all authority in heaven.
and on earth have been given to me go therefore and make disciples of all.
nations this is what the church should be doing all of us should be working.
together to bring people into a living relationship with Jesus who makes.
disciples of all nations baptizing them in the name of the Father the Son and.
the Holy Spirit in teaching them to observe all that I have commanded you.

$^{121}$and behold I will be with you always to the end of the ages you see here name of.
the Father the Son and the Holy Spirit name is a singular or plural in Greek it.
should have been the names of the Father the Son and the Holy Spirit here you.
know there's no textual variant or anything like that all the manuscripts.
are saying in the name singular of the Father Son and the Holy Spirit right so.
it's only one name but there are three persons being named the Father the Son.
and the Holy Spirit in the old days when you know people trying to try to.
understand what the Bible says they use special terms so one of the old terms.
going back to several hundred years ago they use the word God had you know God.
had God had doesn't mean God has a head God has the old theological term back in.
the days of John Wycliffe you heard about Wycliffe Bible Translation Society.
is still doing some translation work around the world so John Wycliffe was a.
guy in the 14th century several hundred years ago he came up with the term God.
had God had actually means theologically is the essence of God God himself the.
self of God God had and he was trying to use this term to try to capture what is.
going on here there is no grammatical errors in Greek however the statement.
there in the name singular of Father and Son the Holy Spirit is the theological.
statement speaking to us with just one God but there are three persons in that.
God had right not going back to Chris Tomlin song there's word God had right.
God have three in one do you realize that there was pretty strange to our.
contemporary language right we don't use God had anymore right we don't pray to.
the God had and and all the other stuff but you know you need to understand.
there's a theological background to it God had three in one so the name singular.
of the Father Son and Spirit so right here the concise statement of what the.
doctrine of Trinity is is this so you repeat after me the Father's God the Son.
is God the Spirit is God but the Father is not the Son the Son is not the Spirit.
and the Spirit is not the Father there's not three gods but one God but one God.
in three persons you got it all right let's do it together the Father's God.
the Son is God the Spirit is God and the Father is not the Son the Son is not the.
Spirit the Spirit is not the Father how many gods are there there's only one God.
however one God in three persons right one personal God and three different.
persons it's pretty hard to for us to understand because we're talking about.
the God had the self of God right God self is way beyond our rationality.
right it's not irrational it's trans-rational right it's beyond our.
intellect I think Augustine puts puts it pretty well he used a human analogy of.
love to explain the mystery of Trinity he's a thing about love right there are.
always at least three aspects of love the person who loves the person who is.
loved and love itself all right pretty interesting but no one is greater than.
another or the sum of the other two the lover the one being loved and the love.
itself right God is like a community right so he created us that's why we.

$^{161}$have community now who is the Holy Spirit is the third person of the.
Trinity now it comes to us what does he do what does it do in us now I think so.
often in our contemporary debate or a struggle with understanding who who the.
what the Spirit does sometimes we so bogged down into some sort of tunnel.
vision that the Spirit all he does or explicitly mainly what he does is to.
give us the gift some spiritual supernatural gift of speaking in tongues.
right have you ever heard of you know people who speak in tongues speaking in tongues.
being a sign who we do have okay all right that's cool someone pointing at.
you speaking in tongues meaning you know there was really such such such thing.
someone who is gifted to speak some of the languages maybe known language or.
unknown language back in the days when I was in seminary I had three roommates.
Korean roommates we are all in the same class same school I remember one of my.
friends he speaks tongue speaks in tongue he speaks in tongue every single.
morning when he does the devotion the quiet time so you know I know a little.
bit what Korean sounds like all right so definitely he was not speaking in.
Korean he was speaking in some non Korean language right not English right.
not Chinese not some other languages so so one day I'm not on the door and ask.
him what are you doing he said oh I'm having a communion with God right.
speaking in tongue speaking in some other mysterious language right some.
spiritual language now that could be one sign of the work of the Holy Spirit but.
I think we are so truncated sometimes or we do all the works of the Holy Spirit.
into one thing right because it's so explicit so we center on it oh you have.
the Holy Spirit you have to speak in tongue I don't think that's that is.
really biblical teaching the work of the Holy Spirit is so embracing there were.
so many things we can experience the Holy Spirit in different ways and let's.
go back to the word about Holy Spirit when Jesus was speaking to his disciples.
he says I will send or the Father will give you and other helper.
procreators procreators it's being translated in many different ways helper.
counselor advisor strengthener encourager advocate or any other many.
different translation but what is clear is that the Spirit is he not it the.
Spirit is he as a person he speaks he relates and he must be obeyed he must be.
loved and we lay the back along with the Father and the Son now so often we think.
of the Word of God or the work of the Spirit we often think of it in some.
truncated way we always think that okay the Word of God is simply just you know.
sprinkling you know rose petals along the way and you know tells you every.
single morning how great you are how wonderful you are and fearfully made you.
are and all that stuff and to comfort you and to be with you you know it's.
just repeating the same message message every every every single day now.
nothing wrong with that nothing wrong with that but you know if you think.
about it that's if that is what God does right if you reduce the work of God the.

$^{201}$work of Holy Spirit to just this tiny bit of praising you or uplifting you or.
building up yourself right you get numb after one or two days right have you.
ever relate to someone or your lover or your boyfriend or a girlfriend or your.
spouse every single morning we just praise you we just honor you and and.
tell you how good you are and how wonderful you are you know you get bored.
after a day or two right well give it a month or two you definitely get bored I.
think I think the Bible present to us the Holy Spirit God himself to be on the.
side of who you could become more than the side of who you are right now you.
want a God that is on the side of who you could become more so than who you.
are right now if you reduce what God does in your life simply frees you into.
what you are right now it will stunt your growth as a man or woman God.
doesn't want to do that the Spirit doesn't do that so we come down to what.
the Spirit really what what he really does in your life is to transform you to.
change you gradually sometimes radically to become who you could become not who.
you are right now so I want you to listen to the Spirit the Spirit is.
talking to us every single day we make in us every single day reminding us the.
biblical truth of the gospel reminding us the the word of the scriptures so.
that we can be growing into what we could become Ephesians 5 talking about.
this the work of the Spirit now listen carefully 515 look carefully how you.
walk not as unwise but as wise making the best use of time because the days.
are evil therefore do not be foolish but understand what the will of the Lord is.
do not get drunk with wine and beef but be filled with the Spirit talking to one.
another addressing to one another in Psalms in hymns in spiritual songs.
singing making melodies to the Lord with your heart giving thanks always and.
everything to God the Father in the name of our Lord Jesus Christ what does the.
Spirit do it's an analogy between wine and the work of the Holy Spirit now we.
all get drunk at times right I mean we if you never maybe you will okay you.
know I'll get getting drunk is being under the heavy influence of alcohol of.
the wine you're not in control anymore or at least you are not primarily in.
control the alcohol is taking control over you that is getting drunk right how.
do you tell someone's getting drunk you know what is the biggest influence at.
that point for that person right is it wine or himself or herself the analogy.
is the work of the Holy Spirit is giving yielding to the control and influence of.
the Holy Spirit more and more in our lives when the Spirit reminds you.
something you yield to it you obey it that is getting drunk right getting.
influenced by the Spirit and that should be our Christian spiritual journey yes.
we've we walk and we stumble and we fall at times but you know in the long in a.
big scheme of things we should be giving more influence to the Spirit instead of.
our sinful self right have the Spirit has the Spirit talked to you recently I.
think the Spirit is very active reminding us right how to live a life.

$^{241}$that is pleasing to God how to live out the Christian character right more and.
more influence of the Holy Spirit on our lives and we walk walking with the.
Spirit right walking according to the will of the Spirit let me give you a.
note another analogy I was sitting here singing songs and and all that and I.
remember the first time I taught someone how to drive a stick shift how to drive.
a stick shift back in the more than 10 years ago I was in Israel and we got a.
rental car and most of the rental car back then was manual transmission so I.
was the only one who could drive it right so I was with my team members so.
we checked into a Airbnb hotel motel or whatever so we stayed there and then in.
the evening after sunset one of the one of the girls in my team everyone is is.
it everyone is a driver everyone has driver license everyone knows how to.
drive but only the automatic transmission right that sucks right so.
so the girl girl asked me hey you know what I want to learn how to drive a.
stick shift I said yeah sure because it's a rental car right you break it you.
know you can return it right so so I go out with her and some of the members and.
we just drive around in the parking lot first gear second gear first gear second.
gear and let's start the engine and start again start engine start again no.
over the entire process you know almost three hours I was circling in the.
parking lot so I'm gonna in parking lot and then uh you know the key for her to.
learn manual transmission what's the key it's to listen to me right just do.
exactly what I say right well there are times you know her no automatic.
transmission mind kicks in right she would drive the stick shift like an.
automatic transmission right I said no that's not right no you have to correct.
yourself right you can't think in a different way so in three hours I.
successfully I taught her how to drive a manual transmission - how wonderful I am.
right three hours right and then we came back from from Israel and then she.
actually brought a stick shift vehicle and she enjoys it right our relationship.
with the Holy Spirit it's pretty much the same listen to the Spirit teaching.
us correcting us guiding us changing us then our mindset will change our.
behavior will change our hands and our feet and our eyes coordination will.
change because the Spirit's voice is to condition us so that we can follow the.
goodwill of the Father so let's do that I have an ending prayer right here we're.
gonna end with this prayer together all right I want you to recite this prayer.
with me all right in your heart or verbally and let's do that then we'll.
close in prayer dear God fill us with all joy and peace as we believe in you.
Holy Spirit empower our lives to produce your fruit of love joy peace patience.
kindness goodness faithfulness gentleness and self-control let's all.
pray together Father we give you thanks because the Spirit dwells in us because.
of the Spirit speaking and teaching and guiding us 24/7 Lord our prayer is that.
we'll be sensitive to your voice to listen and you to your voice thank you.

$^{281}$for remaining in our lives and transforming ourself and we pray that.
gradually in our Christian journey we'll experience you more and more may the.
Holy Spirit continue to be the growing influence in our lives teaching and.
guiding our ways..
\newpage



\section{詩篇}
\label{sec:8rUfl6H7c2c}
\textbf{創造之歌 (詩篇19) - 陳崇基牧師}
\newline
\newline
連結: \href{https://youtube.com/watch?v=8rUfl6H7c2c}{\texttt{https://youtube.com/watch?v=8rUfl6H7c2c}} ~~~~ 語音日期: 2022-07-19
\newline
\newline
\hyperref[sec:ndrseOs8_m4]{\small{< < < PREV SERMON < < <}}
~
\hyperref[sec:index_chronic]{\small{[返順時目]}}
~
\hyperref[sec:index_scriptual]{\small{[返順卷目]}}
~
\hyperref[sec:KZ1AIqUYQds]{\small{> > > NEXT SERMON > > >}}
\newline
\newline
$^{1}$我們回到正經,詩篇第19篇,我們繼續思考詩篇的說話,.
我們這個系列會說幾篇詩篇,上一次說到137篇,.
我知道網絡會斷了,希望有時間可以錄,.
137篇斷了一半,後面那半沒有了,不過未有時間錄,.
不過這次我們看詩篇第19篇,我們先來讀,.
「諸天術說,神的榮耀,窮莊全陽他的作為,天天發出言語,夜夜傳出知識,沒有話語,沒有言詞,人也聽不到他們的聲音,他們的聲音傳遍全地,他們的言語傳遍地極,神在他們中間為太陽安設了帳幕.」.
這句是詩篇第19篇,頭一半,這裡每時每刻告訴我們,原來整個宇宙在說話,在對我們發出言語,.
但這些言語很奇怪,他說無言也無語,但繼續說話,很明顯不是我們人說的人文語,.
但他很清楚發出一些訊息,就是說到上帝的榮耀,上帝的權能在哪裡看到?.
在這個世界宇宙裡面看到,我們每踏足這個世界,每一步,都是屬於上帝的領域,.
上帝自己的物業,想清楚,我們每一個平方吋,都是屬於上帝的,.
在我家附近,最近有一間房子出賣,他賣出來放在市場,我和我兒子八卦,我們去看看,.
不是因為我們有錢買,我們沒有錢買,但我們看到他弄了這麼久,翻新過,我們進去看看,.
進去的時候我們發現沒有人,下午的時候,開放式屋的時候,走進去,.
但那種感覺很新鮮,因為是一間空的屋,一間完全裝修好,裝修得很漂亮的一間小屋,.
但進去的時候,我和我兒子說快點,因為趕時間要走,我們進去,上到樓上,再下來,看到有幾個房間等等,.
有個感覺就是那個地方不是屬於我的,是屬於別人的,和我自己進屋是不同的,和我自己住的地方是不同的,.
因為ownership不是我的,其實我們每一步的時候,都應該有一個這樣的感覺,.
感覺就是這個是上帝的property,這個宇宙是上帝所創造,是屬於他的,沒有話語,沒有言辭,.
人也聽不到他人的聲音,但是每一天去到述說,時時刻刻的去傳揚,告訴我們原來這個宇宙是在說,.
上帝的存在,上帝的glory,上帝的榮耀,等於你去到看一幅很漂亮的畫,很漂亮的藝術品,.
藝術品是在跟你說背後那一位藝術家,他的偉大,他的智慧,詩篇十九篇,.
他提醒我們,雖然耳聽不見,但是卻是在跟我們說上帝的榮耀,上個星期二大家知道有一單大新聞,.
跟太空有關的,大家有沒有留意,大家沒有留意這些,大家都是想著吃的,有什麼好吃,有什麼好喝,.
有什麼新的餐館開等等,上個星期二有一件很大很大的新聞,所謂震驚全球的一單新聞,.
就是2021年12月的時候,美國太空總署發射了一個望遠鏡,上到太空,那個望遠鏡叫做James Webb Space Telescope,.
他就彈了這個東西上去,然後一直向深的太空進發,然後上個星期二,第一次傳送幾張相片回來,.
大概幾張,五六七張,回來地球,然後太空總署NASA就收到,然後就把它公佈出來,.
這個是有史以來第一次看到那個深空裡面,以前沒有看過的,一百年前我們所知道的宇宙,.
只有一個銀河是Galaxy,現在當然我們越來越知道多,但是這個James Webb Telescope射上太空的時候,.
他傳送第一個人類從來以前沒有見過的相片,就是左邊那張相片,.
就是告訴我們這些光線其實已經走了13個billion years,我不知道多少個零,你自己畫一下,.
有多少個零,相片裡面看到的光幾乎是宇宙開始的時候所謂大爆炸,開始的時候看到的第一道光,.
你見到右下角的圖表,就告訴你以往有不同的望遠鏡是射過上去太空的,.
這次是所謂看到Big Bang之後會出現的光,航行了13個billion years,.
現在終於被我們肉眼能夠看到,在相片裡面能夠拍攝出來,.
見到宇宙我們越來越去探索的時候,就是越來越大,越來越無限,.
讓我們看到宇宙是多麼的浩瀚,CP19篇雖然是寫在古代的時候,.
但當時詩人已經看到這個宇宙,是述說著一樣東西,是上帝的榮耀,上帝自己的偉大,.
19章第四節,第四節到第六節,那裡說什麼呢?.

$^{41}$他說太陽,神把他們中間為太陽安設了帳幕,太陽如同新郎出洞房,.
猶像勇士歡歡喜喜的跑路,他從天的這邊出來,繞到天的那邊,.
沒有什麼可以隱藏,得不到他的混亂,用了太陽,用了幾節的聖經,.
去到說太陽的可見程度,或者可以感受的程度,每個人都會感受到太陽的猛烈,太陽的光,.
當然這個和當地巴勒斯坦地很有關係,特別是夏天,去到以色列的時候,.
7月8月,是沒有雨水的,整個天都是藍色的,和我們芝加哥有很多不同,.
我們幾乎有時候一年裡面,有幾個月是不見天日的,除了我們不出街之外,.
太陽很少見到的,因為多雲,或者有雨,好像今天這樣,我們就見不到太陽,.
但是在以色列,你逃避不了太陽,因為太陽出來的時候,知道它日落,.
我們都會感受到太陽的猛烈,用太陽來述說,說上帝自己的存在,.
上帝的權能,是每個人都可以見到的,為什麼不用其他真琴異獸呢?.
一些很奇特的生物,海洋裡面有很多,我最近看了一個YouTube,說八爪魚,有多厲害,.
裡面有很多很多的Nuance,很通明的一個活物,為什麼不用這些真琴異獸來告訴我們上帝創造的奇妙,.
和上帝創造主的奇妙,他的智慧,原因在哪裡呢?.
就是因為太陽人都見到,特別在以色列地,特別在巴士那裡,.
我們先來搞一些翻譯,第五節,太陽如同新郎出洞房,這個我覺得中文翻譯很奇怪,出洞房,.
可能有些現代的猶太人的傳統可以幫到我們,因為洞房這個翻譯了,或者英文ESV等等,.
翻譯成Pavilion,或者Tent,其實是什麼意思呢?.
現代的一些猶太人,他們會有一種叫做Hooper的東西,其實所謂洞房的字背後,希伯來文就是Hooper,那個Tent,.
這個猶太人現在已經有很多年的傳統,就是他們行婚禮的時候,.
新人和那個拉比,或者主婚人,就會在那個Hooper底下,在那個帳幕的底下進行婚禮,.
禮成之後,那個新郎就會帶著新娘離開那個Hooper,離開這個所謂的帳幕,.
就代表他們一個新的家庭出現,一個新的創造,所以如果十九篇這樣看的話,.
就是太陽出來的時候,就好像新郎離開這個Hooper一樣,是一個New Creation,一個新的Creation,.
就是說每一天都是上帝自己新的創造力,如果你這樣看,詩篇十九篇,.
就是在說一些不是我們現代人很多時候想的,一種對宇宙的觀念,我們經常覺得宇宙是一個機器,.
一直按照他的萬有引力,或者其他的principle,他會繼續去run,好像一個機械,一個engine,.
但是詩篇十九篇就是形容這個宇宙,就是有上帝去掌管和帶領著,有上帝去involve的,.
他去介入的,一個這樣的cosmos,一個這樣的宇宙,用太陽去導出上帝自己的榮耀,.
我相信也有另外一個意思,為什麼用到婚禮Hooper,因為在古代的婚禮裡面,永遠都是公開的,.
古時的婚禮永遠都是公開,不需要有我們今天現在搞婚禮的時候,我們有guest list,.
有些人請,有些人不請,不請就怕不請他,得罪大家,請了他又怕他不給禮金等等,.
我們不會這樣,古時的時候婚禮就不會這樣,古時的婚禮整條村都知道的,.
誰結婚,誰有組織新的家庭,整條村就being invited,就會被歡迎來參與婚禮,.
所以換句話說,因為說到太陽,每個人都知道的,就好像婚禮每個人都知道一樣,.
就是你沒理由會miss那個婚禮,你想像一下在那些小的村落裡面,農村,.
怎麼可能有喜事的時候,你是不知道的,沒可能,因為你一定會去那個banquet,.
一定會被歡迎,你可以參加,這裡就告訴我們上帝所述說的榮耀,是不能夠不知的,.
沒理由你聽不到,沒理由你留意不到,就是這個意思,詩篇十九篇,.
我記得好幾年前去網上聽一個英國的兩個學者去辯論,辯論上帝的存在,宗教等等信仰,.

$^{81}$其中一個學者就是一個英國的大拉比,叫做Jonathan Sachs,幾年前他剛剛離開世界,.
另一個跟他debate的,就是Jonathan Sachs,就是英國的大拉比,代表猶太教,.
另一個就是英國的一個生物學家,或者進化論的學家,叫做Richard Dawkins,.
他就認為是無神的,一個有神論,一個沒有神論,兩個去到debate的時候,.
我覺得很有趣的一個地方,當然是很熱烈的debate,但我其中一個對話我覺得很有趣,.
那個對話就是Rabbi Sachs跟Richard講,他說,religion is music, but you are tone deaf,.
Dawkins就回答他,他說,yes, it's true, I am tone deaf, but there is no music,.
這個對話值得我們去想一想,到底有沒有music,還是我們是tone deaf,.
詩篇十九篇,就說有一個music不斷的奏響,there is no way you miss, no way你會聽不見,是無法推諉的,.
只要你豎起耳朵聽一下,你就知道上帝的榮耀,這個在聖經的概念裡面,.
就叫做普世的啟示,universal revelation,即是每個人都可以知道,在你心裡面知道,.
上帝的存在,在見到這個宇宙的浩瀚的時候,你會感受到上帝的存在,.
universal或者general revelation,普遍的啟示,詩篇十九篇,頭一半就是講這件事,.
然後就講到另外一種的啟示,我們或者可以叫它做特別的啟示,特殊的啟示,special revelation,.
不是每個人都可以有的,只有屬於相信他的人才能夠得到這些啟示,.
上帝的子民才能夠得到上帝給他們這些訊息,.
看看你能不能找到,三對並聯的詩句,第七節開始,.
耶和華的律法傳遞,使人素成,耶和華的法道確定,使圓滿人有智慧,.
有一對的詞,Torah,律法,法道,testimony,當然有很多不同的字,都是指到舊約聖經,希伯來聖經,.
但他們可能是同意,但都有他自己特別的強調的地方,.
耶和華的分辭,是三個的名詞,正直,分辭是正直的,使人心快活,.
耶和華的命令,第四個,清潔,使眼目明亮,耶和華的典章真實,傳言公義,.
驚威耶和華是純潔的,傳道榮元,找不找到三對在哪裡,律法和法道,.
分辭與命令,典章和所謂驚威耶和華,Fear of the law,律法是甚麼意思,.
有時我們中文翻譯不是很好,我們想起法律,開車,就捉你,.
Torah是甚麼意思,Torah其實就是guidance的意思,即是指示我們的意思,.
是一個像指南針那樣,使我們人心能夠蘇醒,能夠回轉,.
法道就是一種警惕的意味,即是前面有危險,就使人愚昧,.
不知道危險的,不知道試探的人,不知道會跌倒的人,他就得到智慧,.
這個就叫做testimony,是法道,分辭就代表一種從上而來的command,.
一種authority,命令亦都是同一樣,但很清楚地告訴你,不會出爾反爾,.
上帝是對我們生命裡面有著這種吩咐,這種命令,應該很清楚,.
即是在近東裡面所有人都是有神論,很多時那些神明是不會講清楚他們自己想怎樣,.
因為他要發脾氣,就發脾氣,好像我們人一樣,但是你打開上帝的律法,.
耶和華的律法的時候,他就樣樣東西都跟我們講清楚,大的原則會講清楚,.
大的誡命會很清楚地點醒我們,耶和華的典章講到他很真實,很堅定的,不會改變的,.
耶和華的話語和典章,代表上帝真真實實將他的心意告訴我們,聽的人就能夠得到好處,.
最近我教我的兒子駕駛棍波車,因為我覺得這些,我知道現在通常都買EV,.
下一代都是買EV,電子車,應該都不用學了,但我就覺得我們的家族,.
有些Lost Art是不可以不懂的,我就教他駕駛,我每晚都出去教他駕駛,.

$^{121}$你知油費多貴,現在好像跌了一點點,好多了,如果你懂得駕駛棍波車,你會知道,.
你會浪費很多油,因為他不停地丟油,你知不知道怎樣可以成功地教到一個人駕駛,.
當然要有一個好教師,但我想那個Key都是,那個學車的人,他剛剛聽完,.
有很多東西我會清楚講出,但去到路面的時候,他自己就要去揣摩,.
我會給他大的方向,哪裡前面轉右,哪裡轉左,所以他學得都幾好,.
他學得都OK,不要覺得好像很危險,我太太就會覺得很危險,她從來都不敢教,.
我就去教她,但我教了她一兩次,她開始明白到,到底一波起步,二波起步,三波是會死貨的,.
但她肯聽,所以慢慢她就覺得原來不是那麼難,這個Key不是聖靈在我們生命裡面做的工作,.
當我們信主的時候,我們得到這個特殊的啟示,一種特別的啟示,上帝會用他的話語來提醒我們,.
當然不是我們平日生活裡面任何很細微的決定,他都會指明我們不需要,.
但有很多大的原則就已經足夠給我們,就好像我教車,.
我會告訴一些大的原則,有危險的地方在哪裡,我會給他有法道,有命令,讓他知道怎樣去開好車,.
事實上我們基督徒的生活都是一樣,我不知道你最近你自己的生活是怎樣,.
基督徒的生活,可能你都是懵懵懵,或者有些事情忘記了做,是時候去考牌,.
是時候去想清楚上帝對我們的吩咐是怎樣,這裡當然不是單單說你聽我說就可以,.
他告訴我價值在哪裡,What is the value,value很多時候是給我們的,.
是對我們有好處的,如同金紙的寶貴,大量的精金更加寶貴,比蜂蜜或者比蜜棗更加甘甜,.
基督徒的生活,我們跟神的關係應當是這樣,我們會漸漸覺得是一些很寶貴的東西,.
經歷上帝給我們的指示,對我們的指導,慢慢我們會覺得是甘甜的,.
開頭可能有一波二波死火,很笨拙,但慢慢他就會說原來都不是那麼差,其實都很enjoy,.
我兒子前晚跟我說一樣東西,他說I enjoy it,不愧是enjoy it,因為車如果不死的話,可能都是你繼承了,.
舊車一輛,如同金紙那麼寶貴,大家在美國的移民,我們知道三藩市又叫舊金山,.
特別是上一代的移民更加清楚,為什麼叫舊金山,因為十九世紀中的時候,有很多人尋金,.
1849年的時候正式掀起尋金業,當時我以前在北加州住,有很多地方都是紀念這些事跡,.
哪裡曾經尋金,哪裡曾經有人去掘金等等,你記得以前去掘金的時候,那些相片或紀錄,.
採金,掘金,不是在山上真的有大塊金,好像金磚那樣拿走,那些是拍電影才有,.
是要用沙拿盤子,慢慢篩,然後才有一兩滴出來,然後才集合到一些所謂金紙,.
很難的,要很多耐性,我覺得這個比喻非常好,上帝的話語,我們要得著寶貴,得著公榮,.
同樣我們需要做什麼,很多時候要花耐性,很多時候都需要給功夫,有血有汗,你才能夠拿到,.
不是說30秒就下載下來,不是說15分鐘,甚至3分鐘,看完一個影片,我們所有東西就掌握到,.
不,你需要去揣摩,你需要去學習,需要有功夫,需要有技術,在我們經常快走的文化相反,.
但你所得到的東西就是屬於你的,你在上帝的話語裡面所揣摩的所有得著,就是給你生命裡面的好處,.
得著的啟示,是一些很寶貴的東西,請你不要忘記,不要因為疫情去打斷我們很多時候教會的生活,.
什麼時候案例就重新上,然後我們有很多擔心,我們有很多東西要擔心,.
生命裡面是需要去尋急上帝的道,上帝的律法,你要我們生命去到經歷上帝的話,就需要下功夫,要給耐性,.
《詩篇》十九篇,由第一段講到創造界在說話,普世的啟示,然後第二段講到創造主的說話,特殊的啟示,.
他叫律法,十一至十四節,就講到被造者在說話,就是私人在說話,這個就是我們的回應,是你的回應,.
十一節至十四節,並且你的僕人也藉著這些特殊的啟示,剛剛講完,上帝的律法分辭,.
謹守這些就得到大的賞賜,誰能知道自己的錯誤,求你赦免我隱而未見的過失,.
求你攔阻你僕人,不犯任意妄為的罪,不許他們克制我,我也可以完全不犯大過,.

$^{161}$願我口中的言語,心裡的意念,在你面前蒙月臨..
這個很特別,因為如果你習慣看詩篇的話,通常詩篇那些請求的詩歌,都是請求開始,.
通常所有的詩篇都會有這樣的格式,詩篇十九篇就是相反,先有讚美,先有敬拜,.
然後到第三段的時候,就由我們自己去思想,自己去悔改,.
詩人見到上帝的榮耀,在宇宙裡面彰顯出來,見到上帝的話語,他的訓詞,在聖經裡面彰顯出來,.
現在他就去悔轉,這個很特別,除了格式之外,倒過來,.
就是見到上帝的話語奇妙,宇宙浩瀚的時候,他不是發出讚美,他是發出請求,求上帝去赦免他,.
其實這個很有道理,當我們見到上帝的偉大的時候,我們其中一個反應,應當就是去反省一下,.
我們有多麼渺小,我們有多麼多的錯實,我們有多麼的broken,有不足,誰能知道自己的錯實,誰會知道,.
這些是一些rhetorical question,一些修辭上的問句,意思就是說答案其實大家都知道,問來的時候其實就是要提醒我們,.
誰會知道自己的過失,沒有人知道自己的過失,你和我都不能去找清楚我們到底有什麼的不是,.
我們不知道的,無論我們做了多久的基督徒也好,做了多久的人也好,我們總有一些是我們不知道的,原來那些不知道的就是我們的不是,.
謙卑在上帝面前,我們去承認原來我們有的過失,不單止我們明做出來,明講出來,而是隱而未見,看不見,.
並且我們的罪成,我們人的本成,他去到反省度的時候,他說他會犯任意妄為的罪,.
剛才我們守詩歌都是一個很好的提醒,我們來到上帝面前,聽見他的律法,來到敬拜的時候,我們有沒有想清楚我們是不是謙卑,我們自己下來,.
這裡多少次講到他形容自己是僕人,兩次,你的僕人,你的僕人,肯聽悔憤,肯聽受教,見到上帝的榮耀,聽見上帝的律法對我們說話,.
我們有的回應就是去受教,原來我口中言語心裡面的意念在你面前蒙越南,acceptable,這個經常在五經裡面出現,.
就是那些古意識的人,他們拿著牛,羊,去到上帝面前都去獻祭,無論你獻什麼祭都好,贖罪祭,感恩祭,凡祭等等,.
acceptable,蒙越南,獻上的時候,上帝是歡迎的,你都去接受,或者我們今天來到上帝面前敬拜,或者每一個禮拜來到上帝面前敬拜,.
我們有沒有這樣去自問,你有沒有這樣自問,那個時候我們來到教會,我們總會帶著顧客心態,今天音樂好不好,歌剛剛好聽,.
這當然是很基本的,偏偏我們最基本的東西就未能掌握,今天那個program行不行,是不是給五顆星,你今天來到敬拜,.
如果你真的想去到遇見你的主,你的上帝的話,不是就這篇詩篇最後那幾節說,你的心態,你的態度,才要調好,.
如果你不調整好的話,我告訴你,你來多少次教會,都不會對你有大的賞賜,你想你心靈裡面得到易主的話,.
上帝我們求你,你都去謙卑我們,讓我們知道我們是你的僕人,主你對我們說話,並且你幫助我們,.
去誓聽你自己的言語,請在我們心裡面有很多我們不能夠看見的東西,主我們有個過失,我們求你赦免,.
求你都去引導我們,行一條易路,讓我們作為你喜悅的,我們的敬拜,我們所奉獻的,夢你的月亮,祈禱奉傳明禁..
\newpage



\section{}
\label{sec:KZ1AIqUYQds}
\textbf{比對研究 6: 基迦美斯史詩 \#1  -  陳崇基牧師}
\newline
\newline
連結: \href{https://youtube.com/watch?v=KZ1AIqUYQds}{\texttt{https://youtube.com/watch?v=KZ1AIqUYQds}} ~~~~ 語音日期: 2022-07-19
\newline
\newline
\hyperref[sec:8rUfl6H7c2c]{\small{< < < PREV SERMON < < <}}
~
\hyperref[sec:index_chronic]{\small{[返順時目]}}
~
\hyperref[sec:index_scriptual]{\small{[返順卷目]}}
~
\hyperref[sec:ExG_dJbvHck]{\small{> > > NEXT SERMON > > >}}
\newline
\newline
$^{1}$所以我們不需要怕這些事情,正如我們之前說過.
有相似的地方不等於抄襲.
不等於哪個B一定抄A,或者遲來的一定是錯的.
或者遲來的一定是劣貨.
早來的就一定是好一點,不一定的.
正如我所說,最好就想一個現代的例子.
譬如我之前曾經提過911.
911剛剛發生之後,有很多人寫書.
特別是一些陰謀論.
有很多這樣的報導,有很多這樣的故事出來.
令到大家對這件事很混亂.
過幾年之後才有人寫一本真正的歷史書.
去到議政試聽.
所以那本書才是權威.
而不是剛剛出現的頭一兩個月的亂說的理論.
所以我們比對研究就是在做這件事.
將聖經放回在它的背景下.
我們不妨把這一點再拿出來.
就是一些經外文獻,一些歷史的背景.
文化地理或者考古的背景.
可以幫助我們明白聖經背後的世界.
以致我們看聖經的時候不會那麼利地.
因為聖經也是在回應當時的一些故事和說法.
所謂的議政試聽.
來講出神所啟示的版本.
那個解釋才是最有權威的.
特別是講到洪水的故事.
當然不只是單單是洪水故事.
聶王記上下也有很多其他的故事.
譬如亞述帝國來入侵等等.
巴比倫帝國也來入侵.
先知也會講出一些議政試聽的說法.
神的啟示,講出真正神在做什麼.
所以我們看看背後的世界是怎樣.
這次我們特別去看一個史詩.
叫做《基加尾獸》.
這個表曾經出現過.
之前打錯了這個字.
現在是11的,應該是11.
這塊泥板.

$^{41}$叫做《Mission Epic》第11塊.
有很多相似的地方.
相似的地方,無論你信聖經也好.
不信聖經也好.
你都不可以想成什麼呢?.
就是抄襲.
因為抄襲是我們現代人版權的概念.
古代的文士沒有這樣的概念.
完全不會問這個問題.
因為他們寫作的文化是不同我們的.
所以我們不要將我們現代.
甚至西方的文化扼到古舊文獻上.
包括聖經,舊的聖經也是一樣.
如果這樣做的話,我們就非常幼稚.
或者我們根本不明白古代世界的文書文化是怎樣.
我們以為我們那個才對.
我們的看法才對.
所以比對研究其實是好的.
幫助我們知道敏銳於其他古代的文化.
曾經第一次兩次的時候.
我用了一個比喻.
我們看聖經的時候就像去旅遊一樣.
去一個從來沒有去過的地方.
從來沒有接觸過的文化.
你去那裡看看風土人情.
要一個旅遊有意義的話.
就一定要放下我們自己的一些成見.
放下我們一些現代人有的社會習俗,語言等等.
我們去到另一個地方的時候.
我們就要學習去聽.
到底他們在做什麼.
然後我們才明白他們發生了什麼事.
如果不是的話,我們就以為我們自己很厲害.
我們自己是superior.
這些其實是文化帝國主義.
所以比對研究或者甚至讀經這個動作.
本身就是放下我們自己的議程.
來進入聖經的世界.
在聖經的世界裡面.
我們聽清楚神對當時初代的讀者講什麼.

$^{81}$怎樣以正視聽也好.
那些神學的訊息才能夠應用在我們身上.
我們大茶經,主日學.
我們一向都是在做這件事.
不過我們更加刻意去做.
所以我們看到Gilgamesh Epic的時候.
我們不妨去看看.
到底這個Gilgamesh Epic是講些什麼.
我就將Gilgamesh Epic很粗略地翻譯了.
就是這個十一塊泥板.
整個故事翻譯了.
為什麼是十一塊泥板呢?.
因為是由第一塊泥板開始Gilgamesh Epic的史詩.
就是講到烏魯克王.
就是南米索不達米亞.
大概三千四千BC的一個皇帝.
一部所謂小說.
一部史詩.
就是讚揚他多少豐功偉績.
做了多少冒險.
探險的一些故事.
很刺激的一些故事.
其中有一些故事好像上次提過的.
就是他去黎巴嫩的森林採木頭.
去打贏守住黎巴嫩森林的神靈胡阿華.
他和他的摯友在那裡.
有一場很精彩的爭戰.
然後打贏了他.
然後回到這裡.
這些古代的史詩篇常都會講這些事情.
因為在歌頌一些英雄的時候.
除了講他的戰績.
他有多大力.
有多神奇之外.
或者甚至講到他.
Gilgamesh epic裡面講到.
Gilgamesh這個王的時候.
三分二是人.
三分一是神.
三分一是人性.

$^{121}$講到他有一些神聖之外.
也都會講很多這些所謂的旅行故事.
就是他去了哪裡遠行.
然後回來.
因為古代的人很少會遠行.
是那些偉大的人才會去遠行.
去了哪裡有旅行.
有一些探險.
有一些冒險.
然後回來這裡.
所以這個Gilgamesh epic.
就是由第一塊泥板.
一直講Gilgamesh.
他的一些經歷等等.
然後在不同學者有不同看法.
在歷史裡面某一個階段.
這個Gilgamesh epic的故事.
就繼續擴大.
然後有些文字加了第十一塊泥板.
不同學者有不同的意見.
認為這塊泥板.
這個洪水的故事.
第十一塊泥板.
是什麼時候加進去的呢.
有些人說是Old Babylonian的時候.
大概1800-2000BC左右.
有些人說是New Babylonian的時候.
新巴比倫帝國的時候.
就是舊約聖經裡面的那個.
利物浦格尼撒王那段時間.
主持六世紀.
無論如何.
在這個故事的發展裡面.
這個史詩繼續多了內容.
然後是一種家傳戶曉的史詩.
等於我們今天問一些坊間上的人.
不同的文化不同的社會.
他都有一些家傳戶曉的小說故事.
一些民間的傳說等等.
歷史的故事等等.

$^{161}$都會有的.
所以在以前近通的社會裡面.
有這樣的動態.
文化的動態.
其實是很自然的.
我們在比對研究裡面.
希望給大家.
起碼有一個品嘗.
起碼知道古代的世界.
不是很多時候我們現代人這樣想的.
是很悶的.
我們覺得古代的世界.
是什麼都沒有的.
沒有我們這麼精彩的.
其實他們古代的世界都很精彩的.
在他們的語言上面.
有不同的語言多元語言的語境.
有不同的人可能會不同的語言.
阿卡利安 埃及文等等.
有不同的文化.
也有不同的動態.
大家會互相的去影響.
或者互相的去對話等等.
所以是一個很熱鬧的世界.
聖經的世界其實就是一個很熱鬧的世界.
不是我們現代人一定這樣想的.
我們才是精彩的.
他們很沉悶 沒什麼可以做的.
他們有很多的文學著作.
還有很多的學術的思潮.
世界的觀念等等.
有很多的衝擊在那裡.
聖經就是在這種環境裡面去寫成的.
它的context是這樣的.
所以我們看基亞美修史詩的時候.
特別看到這個洪水的故事的時候.
這只是眾多洪水的故事.
其中的一個版本.
我們知道聖經 創世記就是其中一個版本.
所以我們看看.

$^{201}$我將它粗略的中文翻譯出來.
希望這次我們一口氣.
一兩個的課堂裡面.
我們可以看完這個Tablet 11.
Gilgamesh Epic.
我們看的時候 我想問一個問題.
到底裡面的神學是怎樣的.
裡面怎樣去描述神.
當然 巴比倫或者近東的世界.
就是多神.
你就問一下 他們都有他們的神論.
不只是基督教 聖經才會有它的theology.
到底這些巴比倫的宗教有什麼神論呢.
他們民間信仰裡面所信的那些所謂的神靈.
到底他們的本性是怎樣的呢.
所以我們希望花一段時間去看.
這個是Tablet 11.
本身第十一法例版.
都已經是一個可以自立的文學單元.
這個是一個故事.
講到一個洪水生還的英雄.
叫做Unapishtim.
我把它譯成Udapishtim.
名字很長.
意思就是說.
He who finds life.
就是那個找到永生者.
為什麼會在Gilgamesh Epic裡面提到Unapishtim.
有這樣的一個名字呢.
因為整個Gilgamesh Epic.
由第一到第十二法例版.
都是在講烏魯克王.
去找到長生不老之藥.
我們以前也有懷疑文化.
遠古的時候也有這樣的傳記式的故事.
人要面對死.
這是遠古以來到現在.
都是在面對的存在的問題.
我們都要面對死亡.
所以Gilgamesh Epic為什麼這麼出名呢.

$^{241}$就是一個很精彩的故事.
在講Gilgamesh.
去很多地方征討.
去很多地方探險.
就是要找到不死.
想不死.
怎樣才可以不死.
怎樣才可以繼續活下去.
因為面對死亡的時候很辛苦很慘.
因為死了一切都化為烏有.
所以這個洪水故事為什麼插入Gilgamesh Epic.
因為這個烏納匹斯.
他就是生還了這個洪水.
然後得到永生.
很奇怪.
我們看看這個故事.
這裡話說之前.
十一塊泥板.
之前有十塊泥板.
這塊是第十一.
這裡是第十一.
之前Gilgamesh已經走到地的極處.
他都還沒有找到不死之虞.
仍然要面對死亡.
死亡非常真實.
他就去到一個地方.
他就見到一個人.
叫做烏納匹斯.
他就見到這個人是不死的.
和眾神並列.
為什麼他可以這樣做.
於是就在第十一塊泥板這個故事裡面.
Gilgamesh就問這個洪水生還的英雄.
叫做烏納匹斯.
問他到底有什麼秘訣.
他居然可以得到永生.
然後他就講了一大輪.
講到他為什麼可以得到永生.
這裡我們首先開始讀.
雖然比較長一點.

$^{281}$但是我們分開兩三次這樣看.
基加美修就對這個烏納匹斯說.
就是對這個遠居於世界極處的烏納匹斯.
他說我見到你烏納匹斯.
你的四肢沒有什麼奇特.
和我一樣.
你一點都不是奇特的.
你和我一樣.
你和我一樣善戰.
Gilgamesh在頭一兩塊泥板.
就講到Gilgamesh是多麼善戰.
他的力量多大等等.
他見到這個烏納匹斯.
和他差不多都是一樣.
但是你看到他們兩個人有什麼分別.
烏納匹斯就是和眾神並列.
但是Gilgamesh就不是.
他在找這個不死之羊.
他就說.
當然這裡有很多細節.
那些泥板斷了.
學者會爭論到底那些字眼是什麼.
有時候我們看一兩行一兩句.
不是那麼清楚的時候.
不要緊的.
最重要是看到一個大體的故事.
然後看這個故事的時候.
看看它的神學是什麼.
還有他們是怎麼解釋世界的事.
特別是為什麼會有洪水發生.
這個世界的命令到底是怎麼樣的.
所以他就和烏納匹斯說.
請告訴我.
當你求永生的時候.
你為什麼會呢.
你尋找生命的時候.
為什麼你會加入了這個.
Ranks of Gods.
眾神的行列.
就是你現在有如神一樣.

$^{321}$為什麼呢.
烏納匹斯就回答他.
對基加美修說.
基加美斯就說.
讓我告訴你一個秘密.
就是眾神的秘密.
讓我告訴你.
他就說到一個地方叫做Sarupa.
Sarupa是一座你都知道的城.
在右法拉底河的河邊.
那個城是一個亂古的城.
眾神也是住在這個.
Sarupa這個地方.
就說回往事.
歷史曾經這些大神.
就決意來降下大洪水.
Sarupa在哪裡.
說到烏納匹斯團隊本身作王的地方.
就是這個地方.
Sarupa.
現在也是一個考古的徒幼.
曾經有少少的法律工作.
烏納匹斯就在下面.
所以也是南米索不達米亞的城邦.
說到這個烏納匹斯團隊.
就是一個王在Sarupa作王.
他就說眾神曾經決定.
降下大洪水.
為什麼呢.
說到眾神.
眾神曾經向他們的父神.
叫做Father God.
就是阿魯.
和他們的謀士的戰神Eli.
你知道多神.
不同神就管不同的東西.
等於我們民間信仰華人的.
土地就管土地.
修成的就管修成等等.
月亮的就管月亮等等.

$^{361}$他們也是這樣的.
另外一些的神.
管運河的.
Enlugi.
其中一個的首領.
叫做Ia.
Ia也和這些眾神一起起誓.
他就將.
向著一間茅屋.
去複述他們的計劃.
他們這些大神的計劃.
Ia就是.
其中一個所謂的厭子.
就是他透露的秘密.
眾神一起誓.
我們一定要消滅人類.
為什麼呢?.
通常在米索不達的神話裡面.
眾神想要毀滅人類.
原因就是.
人太多.
生養太多.
然後他們太吵.
令眾神不能睡覺.
我曾經說過.
都是同一個術語.
表達的方法.
眾神不能睡覺.
他們就很生氣.
發脾氣想殺人.
Ia就很有趣.
他就去複述眾神的計劃.
他怎麼複述呢?.
他不是直接跟那個人說.
因為直接跟那個人說.
就是犯了他的起誓的約.
出爾反爾.
但他就不是直接說.
他對一間茅屋說.
茅屋茅屋.

$^{401}$長長的.
叫他留心聽.
聽什麼呢?.
他就說.
他就對著一面牆說.
大聲到一個地步.
屋裡面的人都可以聽到.
又不是真的所謂的爆竄.
洩露秘密.
但又給了拜他那個.
拜伊雅的那個.
他說書魯柏的人.
Man of Sleeper.
這個烏巴圖圖的兒子.
這個作王的.
Lunar Pitch.
他說你要拆掉這個.
這個房子.
就做一艘船.
所謂做一艘方舟.
就棄你的財物然後尋生命.
做方舟救生命.
將所有的活物的種子.
就帶上去船裡面.
然後你要將要做的船.
它的尺碼.
它的尺寸要量度.
闊度長度要相等.
就將它蓋上.
蓋一個蓋在上面.
就好像隱藏的深淵.
然後呢.
Lunar Pitch Team就聽到.
這個.
洩露的秘密.
他說我完全明白.
我就對我所拜的那個主神.
伊雅說.
他說你的命令.
我主.

$^{441}$我剛說過的.
我會忠心的來執行.
但是我要.
怎樣回答城裡面的.
群民和長老呢.
因為如果那些人.
看到我拆了房子又在做船的話.
那些人好奇地問.
為什麼你會做這些事.
伊雅就獻計.
她的神靈就跟她說.
你就說我是她的僕人.
我是伊雅的僕人.
年輕人.
你要這樣對他們說.
就說什麼呢.
他說這個恩尼不喜歡我.
另外一個大神恩尼.
大家起誓的時候.
其中一個叫恩尼.
恩尼是大神裡面最大的.
他最有能力的.
就說這個恩尼不喜歡我.
所以我不能住在你的城裡面.
可能我就不再.
搭在城裡面的乾地.
這裡屬於.
恩尼爾的.
Sarupa的地方等等.
我就要去到.
水的陰.
深淵的地方.
就跟我的主.
伊雅同在.
因為我不會.
似乎這裡說的是她不會拜.
這個Elite.
她只是拜這個E.
Ea.
這個我的伊雅.

$^{481}$就會大大的祝福你.
這個Elite.
就會大大的祝福你.
有雀鳥.
有漁獲,有很多收獲等等.
早上有餅蛇.
晚上有糕.
然後有大地的曙光.
然後就開始造船.
地上的眾民就圍繞著.
這個洪水的英雄.
有些木匠就拿著斧頭.
然後有些斬茅草的就拿著刀.
有些工人等等.
將房屋拆掉.
因為那時候到現在都是.
南緯素不大文化.
那些屋很多時候是茅屋.
在河邊長的很高的草.
他們斬了下來.
然後就將它織成.
不同的籃或者.
甚至房屋的牆都是這樣造出來的.
他們就將房屋拆掉.
然後就造成繩子.
然後就拿一些磁鐵來.
那些防水的磁鐵.
黑黑的那些.
就塗上去等等.
窮人就帶不同的物資來.
第五天.
我就為方舟做了框架.
On the fifth day I laid the ship's framework.
一整英畝就是方舟的面積.
120爪就是它的長.
它的高度.
120爪就是它的周邊.
你記得我上一次說過.
米索不達米亞的一些.
方舟就是洪水的船.

$^{521}$的尺碼是一個正方體.
120乘120再乘120.
的尺碼是一個正方體.
我就畫了它的輪廓出來.
線條也畫出來.
就將它分成六層.
七層.
分成九層.
就像一個四方格.
我們那些Tic Tac Toe.
就是那些.
過三關的.
我們那些.
九宮格.
每一層就像九宮格.
就將它分開了.
每個房間裡面有房間.
將水塞在它的旁邊.
水不會進水.
塞住它.
然後有一些欄杆等等.
這些都是細節.
然後就拿了很多的.
匿青.
3600兩乘3.
的匿青.
又是60的單位數.
倒在爐裡面.
匿青要燒起來.
然後才可以用.
柏油也拿了這麼多.
然後油我就帶.
供應給那些.
搬籃子的人.
等等.
然後就拿了很多的.
匿青.
油就有多少.
兩次的3600兩油.
那個船夫.

$^{561}$就儲存起來了.
所以有些是.
整方舟的.
妙術.
然後就有一些物資.
就想帶上去.
船裡面.
然後就為了工人.
有很多人在那裡一起.
做船.
就載了很多.
工牛.
然後每天都要載羊.
又為了他們.
預備一些啤酒.
大麥製造的酒.
葡萄酒.
甜蜜棗的酒.
等等.
就供應給那些.
工人去喝.
他們喝很多.
他們在元旦的時候就舉行.
然後就.
打開油分配給他們.
到第七天.
總之很多天之後.
就做好這艘船.
是非常困難的.
我相信他指的是.
工程.
他們就.
帶著船的.
頭尾的甲板.
等等.
然後就上船了.
有三分之二的船的新高.
然後就把貨物載到船上.
有多少銀子呢?我都載到船上.
有多少金子呢?我也都載到船上.

$^{601}$有多少活物呢?我也都.
帶到船上.
我把我的家人和我的親友全部都帶到船上.
草原上的動物.
野獸.
工匠.
那些Civilization.
我們講創世記第四章的時候也有講過.
這些是那些.
米索不達米亞.
人類文明的.
最重要的支柱.
鐵匠,工匠,石匠等等.
Craftsman.
這些都是留給.
人類文明的.
然後就搬進去.
這艘船上.
如果你喜歡我們的影片.
就記得要Like,Comment和Share啦.
\newpage



\section{}
\label{sec:ExG_dJbvHck}
\textbf{比對研究 7: 基迦美斯史詩 \#2  -  陳崇基牧師}
\newline
\newline
連結: \href{https://youtube.com/watch?v=ExG_dJbvHck}{\texttt{https://youtube.com/watch?v=ExG\_dJbvHck}} ~~~~ 語音日期: 2022-07-19
\newline
\newline
\hyperref[sec:KZ1AIqUYQds]{\small{< < < PREV SERMON < < <}}
~
\hyperref[sec:index_chronic]{\small{[返順時目]}}
~
\hyperref[sec:index_scriptual]{\small{[返順卷目]}}
~
\hyperref[sec:bjjiDZD_qQI]{\small{> > > NEXT SERMON > > >}}
\newline
\newline
$^{1}$我們繼續第十一發言例版的故事.
這個故事也很長.
之前我們說要整方舟 他們已經整好了.
這裡說他上船 關上門.
船的船夫 田馮的工 有個船夫.
船夫是Pusu Amuri.
有個船夫上了去駕馭那艘水上的舟.
他放棄了他的宮殿和財產.
黎明的光出現了.
他看到一朵烏雲從天邊升起.
雲中有一個阿達神.
他是雷神 風雨神.
在米索不達米亞的雷神阿達.
ADAD.
他就是風雨神.
在迦南地 你記得嗎.
誰是風雨神.
就是巴黎Bell.
所以不同地方有不同的風雨神.
但他做的事都是一樣.
因為他管風雨的.
就是治風雨的.
阿達就在打雷.
然後有些小的神就在他前面前行.
前行的時候好像在打仗一樣.
雨就來了.
然後就將其他的神.
就將這些堤壩撕毀.
米索不達米亞是一個平原來的.
很多時候那些農田.
就會有一些水道是人工造的.
要擋住水 或者將水引去不同的地方.
所以這些堤壩就決裂了.
然後就將這些堤壩全部爛掉.
開始有洪水來.
然後有另外一些的眾神.
這些叫做阿魯納.
Anuna Guts.
他們就帶著火炬來.
這些都是一些比較詩意的描述.

$^{41}$其實有雷電來的.
他們就帶著火炬來.
這些閃光就點燃了大地.
阿達的大能就全遍了朱田.
本身是光明的地方.
他都變成黑暗了.
很好笑的.
一邊打雷.
有光 有雷電.
但是全部都變成黑暗了.
全部都變成混沌了.
他就把大地擊打.
好像把瓦窩打裂一樣.
一日之久就有暴雨 暴風.
這些飛快的吹來.
洪水也湧來.
好像戰事一樣.
去到掠過人民.
因為下大雨 傾盆大雨的緣故.
人們都見不到身邊的人.
彼此不能看見.
你見到他的描述.
描述到那些水 雨是多麼大.
這個時候就把鏡頭從地上轉到天庭上.
眾神就因為洪水而懼怕.
他們就往後退縮.
直奔到天庭最高處.
就是阿努的地方.
阿努就是天上的神.
眾神就被描述成狗仔一樣.
很害怕.
蹲在外面.
這個伊什塔.
伊什塔就是其中一個生育的女神.
她就像分娩的婦人一樣在尖叫.
很痛苦.
聲音甜美的另一個神就是貝萊特.
她就大聲的來到新人.
本身是很美好的.
很動聽的聲音.

$^{81}$但是因為暴風雨來到.
大洪水來到.
連眾神都買她怕.
他們都很害怕.
都管理不了.
都搞不定她.
哪一天會否化成無誘呢.
化成烏有呢.
你可以想像眾神在那裡說這樣說那樣.
在那裡吵架也好.
在那裡投訴也好.
在那裡申訴也好.
在那裡申淫也好.
在那裡說什麼呢.
之前你記得剛開始的時候.
他們的氣勢是降洪水.
要去消滅人類.
所以你可以想像.
每個人都在眾神的天上縮上去.
很害怕的時候就在那裡叫.
為什麼那天我在眾神的會中說出惡言呢.
說肯去同意去到有洪水.
現在不單止是人死.
現在我們都害怕.
我們眾神都害怕.
我怎麼會在眾神的會中出惡言呢.
又要致命地攻擊我們的人民等等.
所以大家就很害怕.
見到下面發生什麼事呢.
就像魚一樣.
鏟鱗遍滿大海.
即是那些人都死了.
遍滿了水面.
阿魯那眾神就與他同哭.
即是那些神害怕到一個地步.
他們都做不到任何事情.
然後就一直哭泣.
因為他們都搞不定洪水.
六日七夜.
然後風吹著.

$^{121}$又有風了.
然後洪水暴風就把大地移為平地.
到第七日.
洪水暴風和洪水就停了它的毀壞.
他們發起的就好像一場生死的搏鬥.
大海現在平靜了.
風雨就停止了.
洪水也停止了.
講回船上的角度.
剛才我們講到洪水來到.
人很多被毀滅.
然後天上的眾神就縮了上去.
因為他們搞不定洪水.
他們被威脅.
然後他們很懼怕.
現在我們講的是船上.
就是那個Unapitched.
就看到直靜一片.
全人類變成黏土.
因為每個人都像黑一樣死了.
地景平坦.
然後我就打開艙門.
日光照在我的面.
然後我就縮著.
因為我也很害怕.
坐在哭泣.
然後淚流滿面.
我就看到船去到哪裡.
去到當時的世界.
當時的世界是平的.
他們的宇宙觀是平地.
去到世界平地的邊緣.
你記得我講過巴比倫的平面地圖.
去到邊緣的有什麼.
就是所謂的宇宙海.
Cosmic Ocean.
已經是地的極處.
就接壤著大海.
The Sea, The Ocean.
然後有幾長.

$^{161}$12倍的60米等等.
不同的單位.
都是60進制的單位.
然後就看到一座山升起.
船就停在Mount Nimish上面.
然後就陸陸的停在山上.
或者山區裡面.
船不准動.
Mount Nimish就不是創世記講的Mount Elravel.
或者Mountains of Elravel.
山區.
阿拉納山區.
這個創世記講的.
但是Gilgamesh Epic第十一塊泥板.
是講Mount Nimish.
尼穆斯山.
大概在哪裡呢.
都是在阿拉納山區附近.
阿拉納山是主後十一二世紀才出現的一個傳說.
說羅亞方舟停在那裡.
這個傳說.
是在土耳其東邊.
這個就是創世記講的Mount Elravel.
但是Mount Nimish就在下面.
大概在這個地方.
在亞述的北邊附近的地方.
怎樣都好.
這些地理的標籤.
都是讓我們看到.
在米索不達米亞的地極的地方.
當時他們想像的地圖.
這裡是平地.
一個圓碟.
外面的地方就是宇宙山.
Cosmic Ocean.
頂著什麼.
頂著上面一個蓋.
那個叫做Rakia.
在希伯來文.
就是中文翻譯做窮倉.

$^{201}$窮倉就是頂著上面的大水.
頂著上面的大水.
讓大水不會下來.
所以他們講的這些地標.
都是當時他們的地極之處.
船就停在那個地方.
不准動.
第一天 第二天.
尼穆斯山上面就停了這部船.
不准它浮動.
第三 第四天都是這樣.
第五 第六天也是這樣.
到了第七天.
我就放甲子出去.
甲子就來往反感飛翔.
因為它找不到乾地.
就回到船上.
然後我就放燕子出去.
燕子又來往反飛翔.
然後找不到乾地.
它就回來.
第三隻放出去的就是烏鴉.
烏鴉飛出去.
看到水退了.
烏鴉就吃了一些東西.
盤旋.
甚至有糞便.
就是說有乾地.
它就不再回來.
然後呢.
就是說有乾地.
洪水退了.
這個所謂的船長.
Ona Pitchfield.
她就把船上所有的東西.
帶出來 釋放出來.
放了出來到地的四角.
在山上就見到一座壇.
她就放置了七個有七個祭壇的器皿.
在祭壇上.

$^{241}$然後上面就堆滿了一些爐位.
一些香.
她就燒著它.
然後眾神靈聞到香味.
就好像蒼蠅一樣.
就湧了來祭物面前.
或者祭壇面前.
為什麼呢.
你想一想到底神學是什麼.
然後讓眾神到近前來.
但是有一個神不要讓他來.
那個就是Enlil.
為什麼呢.
因為一開始的時候說什麼.
說要降洪水的時候.
眾神裡面最大的那個.
有這個主意的就是Enlil.
其他眾神就說.
這個Enlil不可以來吃這些祭物.
為什麼呢.
因為他蠻不講理.
就降下洪水.
也打算去滅掉所有的人.
突然間這個時候.
Enlil就出現的時候.
他就見到有架船在那裡.
他就很生氣.
對眾神滿懷憤怒.
為什麼呢.
因為Enlil的主意.
就是要滅掉所有人.
為什麼會有人生還呢.
Enlil就問.
不應該有任何人是囂傳的.
然後另外一個神明就叫做Ninuta.
他就準備開口.
就對戰士Enlil說.
知道為什麼嗎.
除了Ia之外.
誰會計劃這些事.

$^{281}$誰做得出這些事.
因為Ia就像全知的神一樣.
很有智慧的神.
所以有幾個神就出現.
Ia就說為什麼Hunapishteam.
或者這個船會令到有人生還呢.
原來入數的就是Ia.
Ia就準備開口.
就對這個戰士Enlil說什麼.
就罵他.
去斥責他.
戰士你是眾神的聖者.
你為什麼會蠻不講理.
你會導致洪水氾濫.
然後懲罰眾人.
罪人的罪.
懲罰有過犯的人.
應該寬容.
免得有人被斷絕.
忍耐.
不要讓他們全部被滅.
為什麼要降洪水.
Instead of your bringing on the flood.
你應該做什麼.
不如叫獅子來咬死人類.
那就不會那麼嚴重.
不會那麼嚴重.
可以叫浪來.
可以讓饑荒發生.
土地被摧殘.
可以降溫.
叫人類死一點.
為什麼要降洪水.
因為洪水是大家最怕的東西.
你記得在近東的宇宙觀裡面.
未創造的時候.
大水是最大的威脅.
任何災難都不能比.
因為水再淹沒這個世界.
就什麼都沒有.

$^{321}$你可以有獅子.
又可以有浪等等.
但是洪水就很大件事.
將整個創造淹沒.
眾神都會害怕.
眾神都會怕大水.
這個Ian就繼續說.
揭開眾大神的秘密不是我.
我讓人見到夷夢.
這裡就說回亞特哈西斯.
亞特哈西斯就好像.
阿特哈西斯的epic裡面.
他也是洪水的生還者.
和這個Una Pachim都一樣.
所以不同的故事有不同的名字.
亞特哈西斯聽見眾神的秘密.
就是Ian透露給他們的秘密.
現在他生還了.
就應該為他們訂一個計劃.
然後發生什麼事呢.
因尼爾就走上船.
他就拖著Una Pachim的手.
將他帶上來.
也將他的妻子抱過來.
妻子就跪在他面前.
這個Elie就知錯了.
摸一摸他的頭和尾巴.
然後就站在他們夫婦中間.
就祝福他們.
祝福他們什麼呢.
這裡就最後了.
接著那些我就沒有了.
最後Elie就說什麼呢.
他說以往到現在.
Una Pachim只不過是一個人.
凡人一個.
現在Una Pachim和他的妻子.
將會變成好像我們眾神一樣.
Una Pachim就可以住在哪裡.
住在Far off的地方.

$^{361}$At the source of the rivers.
就是去到這個越之極遠處.
河流的水源頭的地方.
就是眾神所住的地方.
這樣他們就把我帶到.
越之遠的極處.
然後就住在河流的源頭.
我們暫時把這個故事講到這裡.
之後發生什麼事呢.
之後就是Una Pachim.
講一些他解釋.
和Gilgamesh Epic再一個小結.
然後那個故事就說.
Gilgamesh就得不到這個永生.
最後還是要死.
他是魔頭.
我們看到這裡.
大家覺得這裡的神學是什麼呢.
你看這些東西.
好像看一些華人民間的信仰.
裡面有的故事.
即觀音的故事.
或者土地神的故事.
神靈的故事等等.
看這些近東的神話.
有一點像是看那些東西.
為什麼呢.
因為是多神論.
多神論.
我們可以再下一次.
我們才去總結一下.
有系統地總結一下.
但這裡起碼可以留意到什麼.
忠神是怎樣的.
忠神大家的意思是不一樣的.
你罵我我罵你.
這個世界誰是主宰.
沒有主宰的.
整群神都是自己吵架.
而且是猜不到他們會做這些事.

$^{401}$因為沒有解釋原因.
沒有道德的原因來解釋.
為什麼會有洪水降下來.
另外神學.
有什麼神學呢.
就是米索不達曼的神學是怎樣的.
他們想像這些眾神的時候.
這些眾神遇見洪水是怎樣的.
遇見洪水.
哪一張.
哪一張說他們很害怕.
有洪水來的時候.
他們好像狗一樣.
是畏縮的.
他們都怕怕.
因為他們都搞不定洪水.
因為我們講過創世記的第一章.
創造的背景的時候.
就是講創造的故事.
很多時候近東人用的概念.
就是誰可以將大水分開.
使寒地露出.
這個所謂創造.
能夠做到這個的神靈.
那個才是值得敬拜的.
所以為什麼創世記第一章.
又是講到大水.
又是講到真正的創造主.
唯一的創造主.
就是他可以去克服大水.
因為大水都要聽他的話.
事就這樣成了.
但這些眾神.
在米索不達米亞文化裡面.
就不是這樣的.
他們都搞不定.
他們都不知道發生什麼事.
好像那些婦人.
他們這樣形容.
去到新淫等等.

$^{441}$然後就覺得.
哇 這件事我們都搞不定.
我們都會死的等等.
所以你用這樣的背景.
看看聖經裡面有沒有這樣的描述.
跟這些造成一些爭辯.
告訴他們真正的神是怎樣.
真正的神遇見洪水的時候.
應該是怎樣.
所以這裡是神學.
另外一點 神學.
就是講到洪水完結的時候.
有一個憲制.
Gilgamesh Epic Tablet 11.
就說這個Unapeached Team.
就是這個所謂的船長.
就是洪水的生還者.
泊了一艘船.
然後他就把船裡面的東西拿出來.
釋放他們.
動物等等.
然後他就見一座壇.
放置一些器皿.
然後就燒香.
燒香給眾神.
有祭物給他們.
眾神聞到香味的時候.
就像蒼蠅一樣.
就湧去這些祭物那裡.
湧去這些sacrifice那裡.
你可以想像像蒼蠅一樣.
就滾到祭肉那裡.
其實是一些很現實的描述.
因為近東的廟宇.
到了憲制的時候.
很多時候都會發生這樣的情況.
像蒼蠅一樣黏在上面.
是一些很骯髒的描述.
為什麼會形容他們是蒼蠅呢.
在這裡我們可以做一些對比的文學.

$^{481}$Comparative Literature.
就是說其他的洪水的故事.
我們希望將來可以看到.
譬如亞達哈斯史詩.
形容洪水的情況.
差不多是這樣的.
洪水來到.
眾神知道要死了.
他們很害怕.
然後洪水過了.
洪水的所謂英雄亞達哈斯.
就會捉壇然後去憲制.
憲制的時候.
眾神就像蒼蠅一樣.
滾到祭物那裡.
亞達哈斯史詩裡面.
或者其他的洪水故事.
就會形容他們很餓.
為什麼會餓呢.
原因就是這些眾神.
好像人一樣.
在米索不提米亞的文化裡面.
他們都是有需要的.
和人一樣.
你記得我們曾經提過.
亞達哈斯史詩的時候.
為什麼眾神要做人類.
因為那些小的眾神.
低等的眾神.
做事很累.
因為他們要供應給大眾神.
要耕田 要有食物.
給大眾神.
讓他們享受等等.
他們就投訴.
然後大眾神就說.
不如我們做人類出來.
然後他們就做了七對人類.
做了一群人出來.
叫他們服婦.

$^{521}$那些人類就按時按後.
初一十五就去燒香.
和去獻祭肉.
給上面的神靈.
讓這些神靈能夠吃飽.
這裡洪水發生什麼事.
人沒有了.
所以寺廟也沒有了.
人們就去不了寺廟去獻祭.
所以久了的時候.
這些眾神吃少一餐.
再吃多一餐.
就開始肚餓.
他們肚餓的時候就很辛苦.
就開始明白.
為什麼要降洪水.
我們現在後悔了.
因為降了洪水.
沒有人燒香給我們.
沒有人煮飯給我們.
我們就很肚餓.
幸好這個伊勒救了Unapeach.
於是在洪水之後.
他就再去獻上這些sacrifice.
這些祭肉.
然後這些神就好像肚餓的蒼蠅.
一路湧過去.
你看到這裡有個神殼.
就是人和神靈.
是一個互相互補的關係.
就是神又需要你.
你又需要那個神靈.
所以你就定時定候去供奉祂.
祂定時定候去祝福你.
英文可以叫做symbiosis.
就是一種相棲的關係.
神在米索不達米亞的眾神觀裡面.
多神觀裡面.
不是自有永有的.
祂是有需要的.

$^{561}$好像我們人一樣.
所以人就一定需要供奉祂.
所以我們看了幾樣.
第一 眾神有沒有同一個意志.
沒有的 大家吵架.
第二 洪水來的時候.
這些神的反應是怎樣.
大水來的時候是最大的威脅.
他們都會害怕.
第三 就是神和人之間的關係.
眾神和人的關係.
是一種互相互補的關係.
你就去想一下.
在這樣的背景下.
你就可以看到.
不單止是羅亞的洪水的故事.
就導出了另外一種.
自然不同的神學的一種theology.
也導出了另外一種.
完全自然不同的人與神之間的關係的圖畫.
看看創世記.
第六章 這裡做了方舟.
看哪個好呢.
第八章第一節.
水勢洶湧在地上有一百五十天.
然後整個故事的轉折.
八章第一節.
神紀念羅亞和方舟裡面一切的走獸大地.
神使風吹過.
整個故事有沒有說神是怎樣害怕.
沒有的 每樣東西都在祂的控制下.
這篇是第29篇.
大水氾濫的時候.
要說災雀為王.
世界發生什麼事都好.
神都不會怕.
因為神掌管所有的東西.
連大水都在祂的管理下.
大水都是要聽祂的話的.
受祂控制.

$^{601}$所以這裡給我們看到.
創世記想寫一個.
以正視聽的洪水的故事.
來改正當時的人怎樣看曾經發生過的一件事.
歷史裡面有的洪水.
藉著這個事件.
來導出一種完全不同的神學.
再想一想.
除了一神和眾神的分別.
還有神在洪水裡面的反應.
眾神就害怕 肚子餓.
夜禍就災雀為王.
第三就是神與人之間的關係.
這裡說一點點.
可能下一次再說.
我們在六章開始的時候.
花了不少時間去處理神的兒子的問題.
處理什麼問題.
為什麼要有洪水.
就是為什麼要有洪水.
在這些米索不達米亞的古代故事裡面.
沒有什麼解釋洪水為什麼會來.
剛才我開頭的時候已經說了.
大概很多時候原因都是因為眾神分不清.
這些神就像人類一樣.
有需要 太吵的時候睡不著.
人太多了 他們太吵睡不著.
他們就安息不了.
於是他們就要滅感人類.
以致他們有得意睡眠.
為什麼會有洪水.
就是因為他們睡不著.
你猜都猜不到為什麼會有這樣的原因.
這是什麼原因.
但是在創世記就很清楚說明這個原因.
是一個道德的原因.
因為人的罪惡 反藍.
然後他們敗壞了全地.
地上的行為都敗壞了.
不單止敗壞他們的生命.

$^{641}$敗壞了所有的東西.
地上充滿了什麼 暴力.
第十一節 地上充滿暴力.
這裡有幾個字 最後那句.
地上充滿暴力只有三個字.
在希伯來聖經裡面.
三個字 第一個字就是充滿了.
phai ma le 即是it is filled.
然後phai ma le ha ares ha mas.
ha ares 即是地 全地.
ha mas 就是violence 暴力 粗暴.
或者殺害無辜生命.
其中一個意義在這裡.
ha mas 這個字.
如果你可能聽過現代的新聞.
說到以巴衝突的時候.
在加沙的地帶.
就有一個政黨叫做ha mas.
ha mas 即是那些恐怖份子.
用了一個這樣的名字.
我不知道他刻意這樣用還是怎樣.
但是一個很適合的名字.
因為他們做的全部都是暴行.
去殘害平民百姓等等.
ha mas 就是這個字.
在六章第十一節 創世記就是這樣形容.
他發現了什麼.
發現只有ha mas.
feel with ha mas.
有很多的暴行出現.
所以這些人都不需要上帝去毀壞他們.
他們自己已經毀壞了.
血肉之軀在地上的行為都已經被毀了.
其實他們按了那個按鈕.
是一種self-destructive button.
self-destruction button.
自己毀滅自己.
所以上帝說 不行 我公義的審判一定要來到.
所以六章這裡我花了很多時間討論.
為什麼呢.

$^{681}$因為聖經不單止花了一章時間去討論.
這個故事和其他洪水的故事去解釋洪水的原因.
都是特別鮮明的.
解釋的時候就是說.
這些人是行惡事.
所作所想的盡都是惡事.
地上最大惡極.
耶和華見到的時候.
他就很遺憾.
他不是突然間發怒.
他很遺憾因為大地被這些罪和暴行破壞了.
所以他說要按那個reset button.
上帝說要按那個reset button.
等於我們的電腦按那個reset button重新來過.
所以聖經有很清楚的一種道德的原因.
去交代為什麼會有洪水發生.
又是一種議程時停的味道.
因為很多其他的洪水故事.
都講這些原因.
眾神做的事情.
完全是不make sense.
所以羅亞洪水故事就是寫出來.
上帝啟示去教導人.
為什麼什麼是道德.
什麼叫做行義.
以至不會遭到上帝公義的審判.
所以是一個公義的故事.
講到上帝的公義.
講到上帝怎樣去維持著整個世界的秩序.
這個世界的秩序是一種道德性的秩序.
所以呈現了一種很不同的神學.
亦是一種很不同的神與人之間的關係.
亦是一種很不同的世界觀.
我們下次繼續有系統一點去處理.
這些種洪水的故事.
和創世的故事有什麼鮮明分別的地方.
這個比對研究.
因為我們看的時候.
這些故事去到浪費去比較的時候.
我們不可以只說有什麼相似的地方.

$^{721}$因為相似的地方很多時候都是很表面.
你看到這些故事的脈絡.
其實你閉上眼都猜到.
一定會是這樣發生的.
一定會有這樣的情節.
相似沒有問題.
最重要的是我們看到什麼.
就是那個theological differences.
即是聖經原來是和當時世界.
很熱烈的對話.
這些洪水的故事不是單獨的天跌下來的.
是有它的context.
去到以正視聽.
去到教導人真正的神學.
教導人真正的觀念.
好.
如果你喜歡我們的YouTube頻道.
就記得Like, Comment和Subscribe.
謝謝大家收看 再見.
\newpage



\section{詩篇}
\label{sec:bjjiDZD_qQI}
\textbf{詩前想後 - 詩篇一 (研經) - 陳崇基牧師}
\newline
\newline
連結: \href{https://youtube.com/watch?v=bjjiDZD_qQI}{\texttt{https://youtube.com/watch?v=bjjiDZD\_qQI}} ~~~~ 語音日期: 2022-07-19
\newline
\newline
\hyperref[sec:ExG_dJbvHck]{\small{< < < PREV SERMON < < <}}
~
\hyperref[sec:index_chronic]{\small{[返順時目]}}
~
\hyperref[sec:index_scriptual]{\small{[返順卷目]}}
~
\hyperref[sec:B4uBevTGH_s]{\small{> > > NEXT SERMON > > >}}
\newline
\newline
$^{1}$我們這次就講詩篇第一篇.
Again 我們在講的是.
這個希伯來文可能呢.
挖深一點.
我們可能對詩篇或者經文的意義.
會更加好的掌握.
就不會用我們那種的意境去看.
所以我們在這裡希望能夠做到這方面.
我相信這些詩篇是歷久常新的.
什麼時候看呢 感受也不同.
當我們受苦的時候看也不同.
當我們孤單的時候看也不同.
但是無論如何 在詩篇的解釋裡面.
我們試試多做一點功夫 看一些背景.
然後做一些圓文.
我們對詩篇的認識會更加不同.
這次我們看第一篇.
Again 我們有三行.
有Hebrew.
然後有Hebrew的讀音.
然後有英文和中文的翻譯.
Ashrei haish asher lo halak.
ba asad risayim.
這裡是Ashrei haish.
其實就是Blessings.
就是祝福的意思.
滿有祝福的意思 即是眾數.
通常英文的翻譯都是.
Blessed is he.
Blessed is the man.
或者是the person.
我們中文的翻譯就是.
有福的人.
這個也是登山寶訓的背景.
Ashrei這個字.
然後翻譯成希臘文.
新約的時候就說.
貧窮的人有福了.
清心的人有福了等等.
Ashrei haish.

$^{41}$這個字有一個中末性的意思.
中末性的意思就是.
Askeological.
即是有將來性.
因為在希伯來聖經裡面.
Ashrei這個字.
通常不是說一些.
現在已經很亨通.
現在已經很順利的狀況.
而是最終會得到賞賜.
最終將來一定會.
得到上帝的祝福.
所以是有一個中末性.
有將來性在裡面.
所以無論是新約的.
說到拜福.
或者這裡說到Ashrei.
有福的人有福.
都是有一個中末性.
即是有將來性.
換句話說.
現在那個人.
那個詩人未必一定是一個最好的狀況.
往往不是的.
往往都是一個很困難的時間.
但是因為上帝給他們盼望.
他們可以有所謂的Ashrei.
有一個祝福.
可能你這樣去意譯.
就是最終得賞賜的人.
就是什麼人呢.
Ashrei Ha'ish.
Ashelo Halak.
不幸在.
什麼呀.
這是第一動詞.
Halak 走路.
或者生活的意思.
Ba'a Sa'ad 就是那個council.
就是那些人.

$^{81}$什麼人.
Ressaib 就是惡人.
我們譯作wicked.
有時英文中文也叫惡人.
在聖經裡面特別的箴言.
或者一些教導性的詩篇.
經常都會說到Ressaib.
和他的相反就是Satikim.
就是義人.
惡人和義人.
惡人就是那些不會看上帝的秩序.
對世界的秩序.
對道德的秩序.
那個的旨意.
這些人是不顧上帝的.
給人間的那種order.
那些moral order.
義人就不是了.
義人就是會去到.
體察到上帝在世界裡面的order.
就是勤力的.
在這個箴言裡面.
勤力的人.
親手作供的人.
他一定會得到上次.
還掛得掛還抖得抖的意思.
就是有一種moral order.
wisdom的order.
義人就是不會看這件事的.
他們說有福的人是什麼呢.
No halak, but asad.
Ressaib.
不會拒絕走在這些計謀的路上.
然後第二句說話.
就把他調過來.
你看到有一個副詞.
拒絕的分詞.
No halak.
然後asad.
in the bard的意思就是.

$^{121}$in the.
asad resaib.
然後就倒過來.
不是走動詞.
動詞放在最後.
然後先放那句說話.
或者叫prepositional phase.
就是前置詞的分句.
Uvr derik hatayim lo amad.
就是in the way.
in the road的意思.
derik這個詞經常出現.
就是走路的意思.
就是那條道路.
或者是說到我們生活的模式.
lifestyle.
就是生活的習慣等等.
derik.
hatayim就是由hata這個字出來.
我們也叫sinner.
或者是罪人.
就是不行公義的.
去到行出罪惡的.
lo amad.
他就不站在那裡.
所以我們經常看中文聖經也好.
英文聖經也好.
我們都留意到三個動詞.
第一個動詞就是走.
然後就是站.
站在不站在.
罪人的道路.
Uvr mashaf lesim lo yasaf.
這裡又是另外一個prepositional phase.
就是前置詞的分句.
Uvr masaf就是in the seat.
就是坐的位置.
或者住的地方.
今天在以色列.
有一個小鄰村.

$^{161}$叫做Mosaf.
就是用這個字.
Mosaf就是settlement的意思.
就是那個社群.
大家一起耕種.
在村落耕種.
他們叫做Mosaf.
在以色列裡面有很多小的Mosaf.
Mosaf什麼的.
就是代表大家走在一起.
有一個社群.
所以就特別指到關係.
lesim marcus.
用口去欺壓人.
或者說一些狂妄的話.
等等.
lo yasaf就不住在那裡.
在那些Mosaf.
lo yasaf.
又是同一個字.
這個yasaf加一個m.
在前面.
就變成一個名字.
就是一個settlement的意思.
lo yasaf.
他就不會留在那裡.
不會settle在那裡.
不會定居在那裡.
或者不會和這些人打交道.
所以有三個動詞.
halak amat和yasaf.
三個的lo.
lo lo lo.
然後這裡指一些.
不跟隨上帝的話的人.
這個是詩篇第一節.
就像打開門一樣.
整篇的詩集.
就是邀請我們去敬拜神.
打開詩篇.

$^{201}$詩本的時候.
詩集就是一些詩歌.
來敬拜上帝.
如果要敬拜上帝的話.
整篇詩集的編輯.
最後的編輯就告訴你.
先打開第一篇第一節.
看了第一篇.
你就知道到底什麼人.
可以來到上帝的面前.
就像詩篇第24篇.
誰能登上王的山.
誰能到上帝面前敬拜他.
到聖殿去敬拜他.
就是這個人.
就是ashrey haish.
ashelo halak等等.
所以第一篇就像打開門.
告訴我們什麼人.
才能來到神的面前.
第一篇一向在解經歷史.
或者在教會歷史.
都是很重要的.
因為第一篇和第二篇.
其實在某些抄本的傳統裡面.
其實是放在一起的.
第一篇和第二篇一起.
所以第一篇就像一個大門口.
邀請一個人進來上帝的聖殿.
之前他需要想一想.
到底誰能夠去敬拜他.
這個抄本是一個中世紀的拉丁抄本.
很出名的.
13,14世紀.
主後13,14世紀.
你看到在抄本的詩篇的時候.
第一篇是用紅色墨水去抄的.
這個墨水是用紅色墨水去抄的.
這個抄本的開頭.
你有沒有看到拉丁文的幼福字.

$^{241}$一個大的B字.
當時這些字母會有一些裝飾.
特別的裝飾給他.
第一個就是B.
然後有一個大衛在彈琴.
然後有EATUS等等的字眼.
Bietus 幼福的人.
然後他就繼續抄下去.
用意想說出的就是.
詩篇第一篇是一個很重要的一篇詩篇.
其他的都不會像這樣抄.
第一篇他就故意這樣抄.
證明他掌握到第一篇想做些什麼.
第一篇是整個詩集的介紹.
等於我們看書的時候開始有序言.
序言很多時候我們不看的.
序言就是告訴你作者想說些什麼.
作者最主要的目的是什麼.
所以就好像preface introduction.
詩篇第一篇我們都很熟悉.
其實是屬於某一種的詩篇.
這裡我有一個簡圖.
就是將150篇的詩篇.
就是他的number 他的數字列出來.
然後按照不同的顏色.
或者不同的數目.
我們就將他group在一起.
將他組別在一起.
你看到不同的顏色就代表.
不同的genre 不同的題材.
詩篇有不同的題材.
藍色的就是哀歌.
這些25篇 第9篇 第3篇等等.
很多都是開頭有的.
這些都是叫做lament.
lament就是一個格式.
總是這樣開始的.
我們就知道他是一個哀歌.
然後就有贊美詩.
贊美詩也有他自己的典型格式.

$^{281}$有所謂的信心詩.
Trust, Song of Trust.
這些就是91篇 63篇等等.
君王詩篇就是89篇 第2篇.
第一篇 第二篇 你記得.
登基的詩篇.
就是慶祝大位的君王.
或者皇朝的君王.
他能夠登基作王 代表上作王.
這種royal songs.
然後在這邊有敬拜.
敬拜和贊美差不多.
不過有時候有多了一些元素.
可能會解釋清楚所謂的呼召.
敬拜的呼召.
call to worship.
那個原因是什麼.
你要敬拜神 因為什麼.
因為神的恩典很大.
神的命是大的等等.
有個恩惠在那裡.
但是有相似的地方.
然後感恩.
感恩就是私人感受過 經歷過上帝.
然後回頭去多謝上帝.
然後中間這裡就是所謂的Wisdom Songs.
就是這些智慧的詩篇.
119篇就是其中一個.
有時候叫做Torah Songs.
就是托拉詩篇.
專門講述上帝的道.
上帝的律法的奇妙.
19篇我們要記得.
第一篇就是這裡.
算是一種叫做Wisdom Songs.
Wisdom Songs有什麼特色呢.
彭真講到Wisdom 智慧的時候.
永遠都是two ways.
就是他在你面前放了.
只有兩個選項.

$^{321}$一是生命的道.
一是死亡的道.
一是通往上帝的道.
一是遠離上帝的道.
Only two ways.
就是逼你去選擇 逼你去選擇.
在你面前只有一個分岔路.
左和右.
到底你要走哪一邊呢.
這些Wisdom Songs 智慧詩篇.
很多時候就是想做這件事.
去教導我們 去勸服我們.
要去走上帝的道.
詩篇第一篇就有這樣的特色.
只是two ways.
一是上帝的道.
一是惡人的道.
就沒有選擇了.
沒有其他中間線去選擇.
所以詩篇第一篇可以算是一個Wisdom Songs.
By the way.
詩篇詩集150篇.
有五本書.
有五卷的.
1至41 42至89.
就完了卷三.
90開始就卷四.
到最後等等.
每逢卷一和卷二.
每卷中間的銜接地方.
很多時候就有Wisdom Songs.
如果你回去查一下.
即管回去查一下.
一是就是Wisdom Songs.
一是就是君王的詩篇.
所以就好像讓我們看到最後的編者.
將那些讚美詩編在一起.
譬如有些人將小敏的詩歌.
集在一起.
他就別有用心地將不同類型的詩歌放在一起.

$^{361}$可能很多時候是要一起看的.
看完一篇不單止要看那一篇.
要看前後.
有時候是有互相呼應的.
一張第二節.
(粵語).
就是合在一起.
可以有一種Contrast.
一種對立.
對立是什麼呢.
上面說什麼呢.
不從我人的計謀.
不暫罪人的道路.
不坐涉萬人的座位.
但是他做什麼呢.
除了不不不之外.
不做那些事.
不參與.
不去跟他們打交道.
(粵語).
這個就好像剛才寫了一個Bate這個字.
(粵語).
現在就不是In the way.
就是In the Torah.
(粵語).
他的Passion.
他的Desire.
他最喜愛的東西.
令他最快樂最滿足的東西.
是Beth Torah Adonai.
就是耶和華的律法.
還有什麼呢.
在他的律法裡面.
(粵語).
加了兩個字都是同一樣.
加了什麼呢.
後面一個叫做.
Third Person Masculine Singular.
的Suffix.
字尾 His.

$^{401}$加了His 加了一棟.
(粵語).
就是Meditate.
思想.
很多時候Meditate的意思.
不是我們沉默地在想的意思.
在當時.
(粵語).
其實就好像今天那些.
如果你有去以色列.
去西場.
西場是一個露天的會堂.
所謂的哭場.
現在不叫哭場了.
叫西場 Western Wall.
那些猶太人穿黑白衣的.
正統的猶太人.
他們的頭髮會動來動去.
其實傳說就是說.
以往他們一提到.
(粵語).
因為他們不會讀(粵語).
因為上帝的名字是不能讀出來的.
所以我們就用主的名字(粵語).
每逢他們讀到神的名字的時候.
他們就會鞠躬.
但他們每逢去讀經 去祈禱的時候 都會發聲.
自己任人做 自語這樣 默默自語這樣.
所以meditate就是這個意思 一直說 一直說 不斷的去meditate.
meditate的習慣是怎樣.
Your mom dates by day or daily for Lila day and night by day and by night.
我們也叫做晝夜思想.
不是每分鐘都在做這件事 而是習慣地.
打不掉的習慣 不斷地做.
Your mom for Lila.
這是一篇第二節.
然後繼續說下去 說到有福的人到底是怎樣.
He will be for Hayah.
A chaperone attend battle.
He will be for Hayah.

$^{441}$Hayah就是to be的意思.
He was or he is 差不多是這樣的意思.
他是一個怎樣的呢.
Key就是一個英文是light的意思.
彷彿 偶像.
Key是什麼呢.
Is就是樹.
後面有一個形容詞 形容這棵樹.
Is是什麼呢.
Chateau.
我們說栽在溪水旁.
就好像樹栽在溪水旁.
其實是transplant的意思.
帶有這樣的味道.
Transplant是什麼意思.
人工地把一棵樹 大家種過樹 種過植物都知道.
把它移植到別處.
所以是移種 移栽的意思.
移栽由一個地方 把它抽起.
然後栽在另一個地方.
這個叫做transplant.
這個是移植的意思.
那棵樹好像移植了去哪裡呢.
Alpagama形.
溪水旁所謂的.
Alpagama.
這個字很特別.
Alpagama 等一下再說.
為什麼不用那哈呢.
因為希伯來文裡面.
是有自然的河道 自然的河水.
這個字的 那哈.
為什麼用Alpagama呢.
因為這裡說到水道.
Canal 人工的Canal.
Canals 就是耕種梟灌用的水道.
移植到水道旁邊.
然後阿Share去形容它.
那是怎樣呢.
移植到這些水道那裡.

$^{481}$Puyo y temperto.
它的果子.
Puyo y temperto.
創作的第一張石頭出現了.
它的種子 它的果子.
It will give 或者它會帶出來.
它會結出來.
果子它會結出來.
它的果子也會結出來.
在什麼呢.
Betot in his time 或者in its time.
這個很有意思.
Betot 你看到後面這個字尾.
就是his.
It 就是time的意思.
in his time.
他也是his time.
his 是誰呢.
都是一個私人.
有福的人的時間.
還是另一個時間呢.
耶和華的時間.
是不是.
還是外面的季節的時間.
希伯來文還沒有正式說.
就是要你去想一下.
要你去思想一下.
如果這個人.
我們是被移植到這個水道的.
我們會結果子.
但是什麼時間我們才會結果子.
in his time.
whose time.
是誰的時間.
是不是.
可以是神的時間.
不等於現在馬上會有.
然後繼續說.
就說葉子形容這棵樹.
Fa-alehu.

$^{521}$就是it's leaf.
就是這棵樹的葉子.
是不是.
Fa-alehu.
Lo-yikbo.
就不會枯乾的.
Lo-yikbo.
說到它是常綠的一棵樹.
Fa-ko-ashe-yaseh.
Yas-li-ha.
所有他作的.
Asah 這個字.
也是創世的.
第一張我們經常碰到.
之前碰到.
就是創造.
就是做事的意思.
就是doing things.
會prosper.
會興旺.
會順利.
有時候會覺得這個意思是這樣.
Yas-li-ha.
順利.
這個第一章第三節.
所以就說到.
有福的人.
好像樹一樣.
栽在溪水旁.
On ta-dik-sih-ho-git-go-tze.
然後凡所做的事.
都會興旺.
當然也都是在說時間的問題.
現在興旺.
還是什麼時候興旺.
在誰的時間.
In his time.
這裡說到溪水旁的樹.
就好像我們看CB23篇一樣.
會不會我們現代人.

$^{561}$會想起這種.
北美某些很興旺的公園.
國家公園的衣景.
很多水.
很多樹.
這些就是evergreen.
就是上春的樹.
是不是就是這個意思呢.
Again 我們看回本身的背景.
應該不是這樣.
因為巴士坦地沒有這樣的景色.
也沒有這樣的大河.
或者可能是這樣.
這裡就是隱基底.
隱基底是在曠野的.
隱基底.
也有山在這裡.
這裡是中央山地.
沿海平原一直上來.
大概二千多呎的中央山地.
然後這裡一直下去.
就是一個曠野的地方.
連今天的衛星圖片都可以看到.
沒有綠色的地方.
隱基底在哪裡呢.
隱基底在這裡.
就在死海的左邊.
西邊的海岸.
如果你zoom in 你就知道隱基底是什麼地方.
就是在這個山上.
有一個水泉.
剛剛上面有一個水泉.
今天如果你去到這裡泊車.
可以走進去.
是一個自然交流徑.
可以hiking 進去.
裡面有一個瀑布.
一個很小的瀑布.
這是隱基底.
隱基底就是大衛逃避蘇羅追殺的時候.

$^{601}$逃避蘇羅追殺.
在地獄山走難離開的時候.
或者在這個山地離開的時候.
下來下面隱基底就躲起來了.
然後有這樣的故事.
講到蘇羅等等.
就是在曠野當中.
全部草原組畢生的.
沒有水的地方.
剛剛有一個洗的清泉.
他就躲在這裡.
你看到地勢.
這裡有條類似河.
很幼的河流下來.
隱基底就剛剛有水落.
然後樹就在水的旁邊.
剛剛這樣長出來.
不是常春的.
周圍沒有樹沒有草.
它的根要長得很深.
然後才能夠支取這些水.
然後才能夠在懸崖邊.
在這些所謂的瓦地.
阿拉伯文叫瓦地.
就是有時候會乾的河床.
才能夠根直下去.
可以在那裡生存.
有幾乎是囂傳的意味.
如果不是靠著它的根.
再下去的話.
生下去的話.
它就一定會死.
詩篇第一篇可能有這樣的意境.
或者這裡.
瓦地窟.
這裡有一個很舊的.
很遠古的修道院.
叫聖George Monastery.
今天仍然在運作.
瓦地窟在哪裡呢.

$^{641}$瓦地窟在北面一點.
死海一上去之後.
沒多久就是瓦地窟.
什麼叫瓦地窟呢.
我剛才說了.
瓦地窟就是那些會乾水的地方.
偶爾會有春雨.
然後才會長出一些.
你見到有一點綠色.
這個就是我們說過詩篇23篇的死淫幽谷.
這些山崖的地方.
耶路撒冷在哪裡呢.
耶路撒冷在上面.
耶路撒冷夏利利哥.
就會經過瓦地窟這個地方.
這個地方就是有這樣的景色.
你見到剛剛下面有些水.
偶爾會乾.
雨季的時候會有些水多了.
在分水嶺那裡流下來.
樹栽在溪水旁.
就是這個意思.
這些就是詩篇第一篇的樹.
它的根要生下去.
然後才會有水.
另外一個景色.
詩篇第一篇.
這個聖George Monastery.
這個修道院.
然後你見到有水嗎.
緊緊有水.
一直流下去.
然後旁邊有些樹在那裡.
這個詩篇第一篇.
不是那些常春的很大的河.
不是壯闊的河.
而是緊緊足夠.
然後你的生命力可以在那裡生存.
為什麼呢.
因為In the Torah.

$^{681}$上帝的說話.
詩篇第一篇.
溪水旁.
我們譯做溪水.
自然就被人想起是自然的支流.
河水自然的支流.
溪水.
但是我剛剛說過了.
Pargee這個字.
Palak.
那個字.
有自然的河水.
自然的支流.
這些字是希伯來文的.
但是沒有用到.
這個字如果你trace.
有很多希伯來文字.
和Arcadian的字是相似的.
甚至是借過來的.
所謂的Lone Word.
廣東話有時候會說.
拗叫.
為什麼大家有拗叫.
Argue.
其實就是Argue.
Argue就是英文.
因為有不同東西的文化影響.
我們就用了這個字來拗叫.
Lone Word就是這個意思.
借用的詞.
其實是一些外語.
詩篇第一篇.
溪水的字.
Pargee Mayim.
來自Arcadian的Paku.
Paku是什麼呢.
如果你查一下.
Chicago Assyrian Dictionary.
我們曾經講過CAD.
CAD的第P的Volume.

$^{721}$因為是Pagu.
P開始.
你查一下62至64版.
就講到什麼呢.
就講到這個字Pagu是什麼意思.
就是Canal.
是人工的運河.
Irrigation Ditch.
就是一些灌溉的水溝.
這裡有個地圖.
是米索不達米亞.
一個地方Nepal出來的.
是一個泥板的地圖.
有個建築師.
或者一個管轄的市長.
就畫了一個地圖.
在Nepal這個地方.
你見到右上角.
有一條Pagu.
就是這樣下來.
然後繞個彎.
這個是主要的Pagu.
然後有其他.
河道.
是人工的.
全部挖出來.
這些Ditch.
挖出來.
做什麼呢.
這些是農田的地方.
就是要努力.
有血有汗地挖.
在平原挖這些Pagu出來.
Pagu就是這個意思.
Canal.
然後這裡有不同的.
視型文字.
中間就是最大的忽然.
建築師或者畫圖的人.
將它畫成一個村落的中心.

$^{761}$然後有Pagu經過.
圓形的地方其實就是村落.
如果翻譯出來.
就是那個村落.
那個村 什麼名字等等.
然後就講到他們之間.
有什麼關係.
這些Pagu.
很多時候是運河.
他們做商業的貿易.
以物換物等等.
就是藉著這些Pagu.
然後這些Pagu也是去灌溉周邊的田.
所以說的是Pagu Mayim.
你記不記得.
Pagu Mayim.
Ki'i Sato 二席.
Alpagay Mayim.
Palak.
就是Canal 運河.
是人造的.
人造出來不是自然的.
應該是有這樣的意思.
如果你這樣看的時候.
《詩篇》第一篇的意思是很不同的.
你要被二席.
說真的我們每個信主.
都是一個被二席的過程.
由我們本身的地方.
紮根的地方.
二席到哪裡.
上帝的道.
In the law, in the Torah.
然後我們之所以能夠得到資源.
得到灌溉.
人工是要下功夫的.
那個看法很不同.
如果跟一條自然的河一樣.
剛剛我們那個地方.
另外一個景色就是這樣.

$^{801}$你看到有樹嗎.
有隻小驢在那裡.
這個就叫做Pagu.
Pagu Mayim.
這些就是Pagu Mayim.
很明顯就是人工造的.
挖挖挖出來.
然後再在溪水旁邊.
有個人.
他特意這樣做.
我們當他是神.
特意這樣做.
然後將那個人的生命再再溪水旁邊.
然後挖了這些溝渡.
讓他有灌溉.
他紮根在這個礦業的地方.
這個詩篇第一篇.
在大英博物館有一個圖像.
我記得有一次去的時候.
在倫敦.
我就看到這個景象.
有一個亞述王的皇宮.
裡面的壁畫.
你看到亞述王在上面.
舉起手.
敬拜的姿勢.
然後這裡有一個大的道.
上面有個香壇.
有國的香壇.
然後有一個大的道.
他上到山上.
然後這裡有一個有柱的建築物.
是皇宮或者是一個聖廟等等.
然後有樹在那裡.
你見不見到柏固.
就是那些運河水道.
有的.
上色了 本身有色的.
藍色的地方就是所謂的柏基瑪燕.
這些叫做溪水旁.

$^{841}$運河的水旁.
就栽了這些樹.
然後這個王就在那裡敬拜.
在皇宮裡面去敬拜.
司彬第一篇似乎是在說這件事.
通常有花園.
如果一個人就夜思想.
又說要去做律法.
生活多艱難都好.
他住在哪裡.
就在皇宮裡面.
他就好像一棵樹在柏固.
這些水道的旁邊.
然後按他的時候結果子.
按他的時候有順利.
在那裡敬拜.
司彬第一篇就是在說這件事.
你看回司彬裡面的世界.
其實都是一樣.
凡是有樹的地方.
就在哪裡呢.
司彬五十二篇第八節.
至於我的詩人說.
我就像神殿中.
Bavate Elohim.
Biote Elohim.
就是Temple或者是Home of God.
Bavate Elohim.
我就像神殿裡面.
Biote Elohim.
Bate Elohim.
茂盛的橄欖樹一樣.
因為我倚靠上帝的慈愛.
直到永永遠遠.
凡是有樹的地方.
都是在神殿裡面.
你記得所羅門的聖殿.
都有很多樹的裝飾.
九十二篇第十二節.
義人要興旺如種樹.

$^{881}$生長如禮巴聯的香柏樹.
他們在什麼地方.
或者應該是移植於.
移植於耶和華的殿中.
又是這個字.
當然是在說種樹.
Shetolim.
就是Being Transplanted.
Shetolim.
Bavate Adonai.
Bavate Elohim.
Bavate Adonai.
Shetolim.
你記得司彬第一篇說什麼.
Shetol.
Alpagai Maim.
又是同一個字.
在哪裡呢.
被移植到耶和華的殿中.
又在那裡發旺.
在上帝的聖所的園中.
凡是有水有樹的地方.
都是在上帝的皇宮裡面.
在上帝的家裡面.
在上帝的聖殿裡面.
有一條河.
46篇.
就是眾多之流.
洗神的城就是耶路撒冷石安.
充滿快樂者城就是至高者住的聖所.
有河.
然後以西結最後的40-48章.
說到新的耶路撒冷的衣帳.
說到河的兩岸.
有樹在生長.
然後有葉子不枯幹.
果子不斷在聖所裡面流出來.
所以有河有水的地方.
都是說到聖殿所在的地方.
所以這裡說的是什麼呢.

$^{921}$詩人說.
他晝夜思想耶和華的律法.
晝夜思想.
他的喜愛就是耶和華的律法.
他就好像什麼呢.
好像Keyesh.
Shetol.
Alpagai Maim.
就是水道旁邊移植了的樹.
如果他活在耶和華的律法裡晝夜思想的時候.
其實他就是在神的同在裡.
還有沒有其他地方會是這樣呢.
不會的.
全宇宙裡只有一個地方.
是有水.
樹.
上帝的聖殿.
上帝的聖所.
原來在想上帝的律法.
晝夜思想.
想想怎麼去形容.
想想怎麼去解釋.
活在上帝的律法.
這是我們習慣的時候.
原來詩人邀請我們.
就是活在上帝的同在裡.
生活可能很艱難.
不過.
這個時候就可以享受到.
在上帝的同在裡.
並且這些人最終會得到福氣.
Ashray.
繼續看下去.
第四節.
Locan.
又是一個對立的.
Locan Hashayim.
就是惡人.
Wicked的人.
Locan Hashayim.

$^{961}$Kiim.
又是一個對立的.
但是他們會怎樣呢.
Kamuz.
Ashay.
Tetefnu.
Ruach.
Kiim就是一個反轉.
然後Kamuz.
你記得我們之前看.
Iyan是怎樣的.
Kaish.
這裡就不是了.
Kamuz.
音也差不多.
是諧音.
Kamuz.
他們就不是Ish.
他們就是Muz.
因為這些Wicked的人.
Rashayim.
這些不看上帝在眼中的人.
Kamuz.
他們像Kamuz.
不是Kaish.
Kamuz.
就是匡批的意思.
匡批就會被風吹散.
Ruach這個字又出來了.
就是不應該翻譯成Spirit.
不是姓寧.
不是零的意思.
是風的意思.
這是另外一個意境.
匡批被風吹散.
這些在農耕的社會裡面.
經常會出現.
在最高的地方.
那些播羊,播菊的地方.
羊菊的地方.

$^{1001}$通常是比較大風比較高的地方.
這些農夫去到收割的時候.
就把這些穀物,大小物.
把它們養起來.
然後把它們壓碎.
再把它們養起來.
風吹過的時候.
這些東西會被吹散.
就是那些Muz.
那些匡批.
然後第五節.
Au ken lo yakumu, visaim ba mishpat.
這樣.
Ken就是Yes的意思.
So或者In this way.
他們是什麼.
lo出現了.
但不是二人的lo.
二人的lo就是不幸在惡人的道路等等.
他們是站立不住.
Yakumu.
他們不能夠站立.
就是惡人不能夠站立.
就是在審判的時候.
Mishpat.
上帝最後審判的時候.
所以整首詩篇.
由第一個字.
Ezra.
就是一個Ecological的Direction.
一個綜合性的Direction.
就是仰望將來.
這裡是仰望將來.
惡人.
很猖狂很哀悼.
但是Mishpat的時候.
他們就是lo yakumu.
這就是Visaim.
就是到審判來到的時候.
他們就站立不住.

$^{1041}$然後.
Vakhatayim.
罪人.
Bakadat Satyakim.
Bakadat就是罪集的意思.
就是Congregation.
我們教會說的罪集.
我們的會眾.
Assembly的意思.
Satyakim.
我們說過這些罪人.
還有Visaim.
Bakadat.
相反的就是Satyakim.
就是Satik這個字.
詩篇23篇你記得.
說過什麼嗎.
牧者引我走.
為自己的名引我走.
二路.
二路的意義字.
Satik.
Satikim.
就是這些義人.
義人就是看上帝的眼中.
去愛上帝的律法.
去奉上帝去敬畏上帝.
這些人在他的會中.
因為到審判來的時候.
上帝罪集他的義人.
罪人就不能夠站在義人的會中.
最後.
Ki yodaya adonai derik satikim.
fa derik visaim tove.
因為什麼呢.
最後第六節就解釋了.
最後一句.
Yodaya adonai.
adonai這個字又出現了.
就是耶和華知道.

$^{1081}$Yadah這個字出來.
Yadah就是認識的意思.
或者你可以譯作.
當然他有不同的彈性.
要我們思想才知道.
或者應該是上帝去看顧.
看著.
看著什麼.
derik satikim.
義人的生活.
義人的每一步.
義人的derik.
derik那些路.
但是呢.
另外一些derik.
就是derik visaim.
就是那些惡人的道路.
惡人的生命.
惡人的將來.
他們引去的方向.
tove.
perished.
avat.
就是會滅亡.
我們這樣看的時候.
有一點背景.
多一點希伯來文.
等下再聊多一點文法等等.
可能給你新的角度.
或者擴大一點.
你再看.
詩篇第一篇.
脈絡都很清楚.
abba.
有福的人.
ashrey haish.
比較reshaim.
reshaim就是那些惡人.
有福的人.
ashrey.

$^{1121}$low low low.
但是kiim.
betorat adonai.
hephso.
他所喜愛的.
就是耶和華的律法.
然後b.
再比喻.
有福的人.
他就像什麼.
移植到水道的旁邊.
按他的時候結果.
他是誰.
主的時間.
耶和華的時間.
像一棵樹.
kiim.
對比.
另外一個比喻.
就是惡人.
reshaim.
就是被風吹散的康培.
惡人不像者.
又是low這個字.
第五節.
他們卻是kiim.
上面有kiim.
下面又有kiim.
說到reshaim的時候.
就是kiim.
好像什麼.
康培.
khamus.
對比一下.
kiish.
和khamus.
字母一樣.
是k.
最後那個.
是sadiq.

$^{1161}$中間是差iron和mem.
故意有個諧音.
kiish和khamus.
他們像康培.
最後.
二人和惡人的比較.
惡人的道路就會滅亡.
tove.
有解經者就說.
這個詩篇.
第一個字.
就是alif的字.
但最後這個字.
就是taf的字.
因為alif是第一個字母.
taf是第22個.
最後的希伯來文字母.
整首詩篇就是給我們挑戰一個邀請.
你要經拜耶和華.
你是否做一個有福的人.
是否活在上帝的聖殿裡面.
和祂同在.
現在可能是一個比較艱難的時候.
不過在耶和華的道裡面.
耶和華的律法.
在耶和華的律法裡面.
就好像一棵樹.
在上帝的聖所裡面.
紮根.
一直到那個地方.
用功夫的.
用很多時間的.
將來一定會有上次.
有福的人.
將來一定會按他的時候去結果.
如果你喜歡我們的影片.
就記得要like, comment和share啦.
\newpage



\section{18:16-33}
\label{sec:B4uBevTGH_s}
\textbf{我們與義的距離 (創18\_16-33) - 朱莉傳道}
\newline
\newline
連結: \href{https://youtube.com/watch?v=B4uBevTGH_s}{\texttt{https://youtube.com/watch?v=B4uBevTGH\_s}} ~~~~ 語音日期: 2022-07-22
\newline
\newline
\hyperref[sec:bjjiDZD_qQI]{\small{< < < PREV SERMON < < <}}
~
\hyperref[sec:index_chronic]{\small{[返順時目]}}
~
\hyperref[sec:index_scriptual]{\small{[返順卷目]}}
~
\hyperref[sec:QzV41v7JIro]{\small{> > > NEXT SERMON > > >}}
\newline
\newline
$^{1}$頂支配平安.
今天想和大家一起分享的題目.
叫《我們與惡的距離》.
這個名字其實是2019年所播出的一套台灣電視劇的名字.
當時播出引起了很大很大的討論.
因為這套劇有別於過往.
我們所熟悉的台灣那些情情踏踏的電視劇.
這是一套社會寫實劇.
是討論一個無差別殺人事件之後.
整個社會的實況.
什麼是無差別殺人事件呢?.
就是沒有一些特定的殺人動機.
或者殺人對象.
而殺人之後亦都.
那個兇手是沒有任何益處的.
而很多時候最後我們發現那個兇手是有精神病的.
在台灣這幾年其實發生了好幾宗.
是引起全世界關注的事件的.
例如在捷運裡面的殺人事件等等.
其實我想身處在美國的大家.
對這些事情都有很多的感觸.
而這套劇特別的地方是什麼呢?.
是它不只是從受害者或者受害者家屬的角度來出發.
它甚至是很特別地選取了從加害者.
就是兇手.
加害者家屬.
甚至為兇手辯護的律師的角度.
亦都會討論到在社會裡面有同樣.
精神病的病患者.
和他們家屬的角度.
去討論這個這麼複雜的社會議題.
所以這套劇的名字叫《我們與惡的距離》.
的確.
我們很容易會定性.
兇手就是惡.
我們都不否認這件事.
但在同時的裡面.
其實當媒體.
例如大力去起底.
那個兇手的各樣的家庭背景.

$^{41}$他的學歷.
他所經歷的事情等等.
或者是坐在家裡電視機面前.
或者在電腦面前.
在社交平台很容易留言說.
這個兇手該死等等的.
社會普羅大眾.
其實這套劇想探討的.
是我們這一班人.
好像與我們無關的這一班人.
到底與惡有沒有距離.
當然今天.
我不是想和大家來個劇評.
有機會的話我們可以一齊討論一下.
這些劇.
今天仍然是想和大家分享聖經.
我們一齊來看《創世記》18章16到33節.
《創世記》18章16節.
三人就從那裡起行.
三人是在說前面的上文.
是在討論一些訪客.
來探望亞伯拉罕.
以致亞伯拉罕和他們有下變的事情發生.
三人就從那裡起行.
亞伯拉罕也與他們同行.
要送他們一程.
而話說.
我所要做的事.
豈可瞞著亞伯拉罕呢.
亞伯拉罕必要成為強大的國.
地上的萬國都必因他得福.
我眷顧他.
為要叫他吩咐他的眾子和他的眷屬.
遵守我的道.
秉公行義.
使我所應許亞伯拉罕的話都成就了.
而話說.
所多瑪和阿摩拉的罪惡甚重.
聲聞於我.
我現在要下去察看他們所行的.

$^{81}$果然進帳.
那達到我耳中的聲音一樣嗎.
若是不然.
我也必知道.
22節.
二人轉身離開那裡.
向所多瑪去.
但亞伯拉罕仍舊站在耶和華面前.
亞伯拉罕近前來說.
無論善惡.
你都要剿滅嗎.
假若那城裡有五十個二人.
你還剿滅那地方嗎.
不為城裡這五十個二人.
饒恕其中的人嗎.
將二人與惡人同殺.
將二人與惡人一樣看待.
這斷不是你所行的.
審判全地的主豈不行公義嗎.
耶和華說.
我又在所多瑪城裡.
見有五十個二人.
我就為他們的緣故.
饒恕那地方的眾人.
27節.
亞伯拉罕說.
我雖然是灰塵.
還敢對主說話.
假若這五十個二人短了五個.
你就因為短了五個.
毀滅全城嗎.
他說.
我在那裡若見有四十五個.
也不毀滅那城.
亞伯拉罕又對他說.
假若在那裡見有四十個.
怎麼樣呢.
他說.
為這四十個的緣故.
我也不做這事.

$^{121}$亞伯拉罕說.
求主不要動怒容我說.
假若在那裡見有三十個.
怎麼樣呢.
他說.
我在那裡若見有三十個.
我也不做這事.
31節.
亞伯拉罕說.
我還敢對主說話.
假若在那裡見有二十個.
怎麼樣呢.
他說.
為這二十個的緣故.
我也不毀滅那城.
亞伯拉罕說.
求主不要動怒.
我再說這一次.
假若在那裡見有十個呢.
他說.
為這十個的緣故.
我也不毀滅那城.
又說.
與亞伯拉罕說完話就走了.
亞伯拉罕又回到自己的地方去了.
經文明顯是在說.
關於亞伯拉罕為所多瑪城代求的描述.
但如果不是這樣好聽.
其實就是跟神討價還價的一段記錄.
就好像一個遊客.
去香港女人街買東西的時候.
在那裡講價.
而這種講價式的代求.
其實是很匪夷所思的.
另一方面更加奇怪的就是.
為何亞伯拉罕要為所多瑪求情呢.
所以今天我嘗試從兩個角度.
一起去討論這個代討.
第一方面 內容.
另一方面 原因.

$^{161}$我們先說這個禱告的內容.
我們留意到.
原來經文有六次.
亞伯拉罕開價來向神求.
我們先不要討論數字.
即是50,45,40,30.
但你留意到是有一個模式的.
整段經文的記述.
是非常像跟隨一個方程式一樣.
我們先留意一下模式.
首先.
第一件事我們留意到的是甚麼呢.
就是神沒有還價.
每一次亞伯拉罕開價.
很有趣地要求就接受.
當然如果我們真的去遊客區買東西.
我們真的這樣去討價還價.
我們當然希望小販接受我們的條件.
但無論如何神竟然是.
很順攤可以這樣說.
就接受了亞伯拉罕開出的條件.
另外內容亦是非常公式化的.
假若有多少個義人就怎樣怎樣.
很明顯.
義.
Righteousness.
是這一段的關鍵詞.
我們要再仔細地看的話.
你會發現.
一開始的時候.
亞伯拉罕是問了神兩個問題.
正正亦是整段記錄最核心的問題.
為甚麼或禱告的內容.
是圍繞著甚麼呢.
就是二十三節的.
無論善惡你都要剿滅.
和二十五節的.
亞伯拉罕問上帝.
審判全地的主.
豈不行公義.

$^{201}$所以無論是禱告的內容.
是圍繞著「義」這個字.
這個Righteous.
或者Righteousness這個字.
而亞伯拉罕來到同神.
在討價還價的過程裡面.
其實他都是圍繞著兩個問題.
就是到底神的審判.
是不是公義的.
亦都是.
你是一個Righteous Judge.
當你這樣評判的時候.
其實你是不是Righteous.
當然.
我們覺得很匪夷所思.
為甚麼亞伯拉罕可以質疑上帝呢.
但其實他沒有質疑上帝.
因為他馬上回答了.
這是一個Rhetorical Question.
這只不過是一個反問句.
他有答案的.
他說.
「將義人和惡人同殺.
將義人與惡人一樣看待.
這斷不是理所行的」.
所以儘管亞伯拉罕提出這個疑問.
但其實他馬上回答了.
他否定了神不會是這樣.
我們退回前一點.
我們想想亞伯拉罕是否否定.
審判的必要性.
是否說神去審判.
就是違反公義.
就是沒有大愛.
是否這樣呢.
又不是.
因為如果我們看前一點.
在十八章二十節到二十一節.
其實當時經文是這樣說的.
「又話說所多瑪和·俄摩拉的罪惡甚重.

$^{241}$聲聞於我」.
或者用另一些翻譯.
「呼冤的喊聲」.
或者「控告他們的聲音」.
即是「聲聞於我」是甚麼.
就是一些控告所多瑪和·俄摩拉的聲音.
已經達到神的面前.
十八章二十一節.
「我現在要下去.
察看他們所行的果然盡像也.
達到我耳中的聲音」.
一樣吧.
所以神是一個回應來的.
神的回應是甚麼.
回應有些人.
來控告所多瑪和·俄摩拉.
當然是甚麼人來控告所多瑪和·俄摩拉呢.
就是一些喊冤的聲音.
是一些被所多瑪和·俄摩拉的罪惡.
所影響的那些人的聲音.
控告所多瑪和·俄摩拉的聲音.
達到神的面前.
於是神要去察看.
到底是否真實.
其實這件事不陌生的.
我們想起甚麼.
《創世記》第四章提到.
阿伯被戈隱所害的時候.
他的血亦有聲音.
是從地裡向神哀告.
所以公義的主.
是必須要為受害的.
受冤的人伸冤.
這樣才是符合神的公義.
所以阿伯不是質疑神應否審判.
阿伯質疑甚麼呢.
質疑如果二人和五惡人.
一同在所多瑪被滅亡的話.
神會如何無差別地對待二人和惡人.
這件事就違反神的公義.

$^{281}$所以神的公義的核心是甚麼呢.
是否神會無差別地對待二人和惡人.
以至最後令到神的公義受虧損.
這是阿伯提出的疑問.
正如我剛才所說.
他是否定的.
我們說說數字.
為何是50,45,40,30等等.
為何一開始是說50.
其實解經家有很多猜測.
其中一個可能的解釋是甚麼呢.
因為在舊約的《先知書》有提到.
在《阿摩斯書》有提到.
一個小的城市可以有100人.
即100人的地方可以稱為一個小的城市.
50人就是一個小的城市的一半.
所以大概是說50是一個社群.
一個社群的一半.
就是50.
最後為何阿伯可以停在10呢.
而不是停在1呢.
你想想其實更加戲劇性.
如果阿伯提出.
如果有一個二人.
神是否不會毀滅這個城市呢.
神說好.
一個二人.
我都毀這個一個二人的遠古.
不去毀滅這個城市.
這樣更加戲劇性.
更加凸顯到神的慈愛.
但不是.
阿伯是停在10而不是1.
為甚麼呢.
因為10仍然是一個群體.
1就不是.
所以50到10的範圍.
其實是說甚麼呢.
就是到底一個群體.
一個意義的數字.

$^{321}$它能否發揮影響力.
反過來.
用另一個說話來說.
到底在一個城市裡.
二人的影響力大.
還是惡人的影響力大.
好像一個天秤.
如果一個城市只有100人.
52人就是50惡人.
但它一直遞減.
甚麼意思呢.
惡人的數字是越來越多.
所以天秤是越來越側向惡人那邊.
而二人的數字又越來越少.
本來50.
變成二人是45.
更加少一點就是40.
30這樣遞增的話.
二人是越來越佔少數.
那麼天秤就越來越不平衡.
反過來說.
在整個城市都惡貫滿盈的時候.
到底佔少數的二人.
這個正義的大多數.
能否起關鍵的作用.
能否逆轉城市的命運.
這個數字.
是要去解釋這件事.
到底一個正義的大多數.
或者有多少正義的人.
可以在一個惡貫滿盈的城市裡面.
起到一些改變和逆轉的作用.
常識告訴我們是不行的.
正如我們廣東人都有句諺語.
一粒老鼠屎.
掉進一棵竹裡面.
可以壞掉整棵竹.
或者如果你試過穿白色衣服.
然後吃薯條.
茄汁不小心弄到白色衣服的話.

$^{361}$你就知道很麻煩.
要漂白一些骯髒的東西.
永遠是難過一些小小骯髒的東西.
去弄壞所有好的東西.
所以.
到底一個少數的二人群體.
一個正義的大多數.
它的存在是否有作用的呢?.
它的存在有多少的力量.
可以平衡和改變結果呢?.
如果這段經文.
在這段經文讓我們看到什麼呢?.
這段經文.
原來神是說肯定的.
如果因為五十個.
神肯定可以不毀滅這個城市.
一直一直去到十個.
神都沒有還價.
神都肯定亞伯拉罕.
和亞伯拉罕說:Yes! OK!.
代表什麼呢?.
神肯定原來二人的群體儘管微小.
其實是有拯救的力量.
是可以勝過周遭更加龐大的邪惡力量的.
是可以改變一個沉淪的地方的.
所以.
如果我們真的明白這件事的話.
對於我們今天身處的社會.
當我們發覺邪惡的力量.
彷彿已經完全壓倒一些良善的力量的時候.
我們很容易會絕望,會灰心.
我們覺得沒有事情可以做.
但是亞伯拉罕和上帝的對話.
確定了一個事實.
十個.
只要十個二人.
都可以改變.
一個 Righteous Majority.
不是一個.
是一個 Community.

$^{401}$儘管他是一個 Majority.
但他是一個 Community.
他可以改變那個地方的命運.
好了.
我們討論完內容.
為什麼亞伯拉罕要為所多瑪代求呢?.
我們很容易.
第一個 Instinct.
我們的想法是什麼呢?.
哦.
因為他為了他的侄子羅德的緣故.
如果大家有.
我們一會兒都會看.
如果大家有印象.
羅德他的侄子.
亞伯拉罕的侄子.
就是在所多瑪的.
是不是因為親情的緣故.
顧念他的侄子.
所以他為所多瑪代求呢?.
但又不合理.
因為如果亞伯拉罕是為了羅德的緣故.
不如直接開口.
索性跟神說.
神啊 你為羅德的緣故.
來拯救這個地方.
他不是.
而且不是說.
只要有一個二人都可以.
所以.
亞伯拉罕明顯不是因為羅德的緣故.
來為所多瑪代求.
我們說回一點背景.
為什麼羅德.
會在所多瑪.
我們回到《創世記》第十三章.
在《創世記》第十三章提過.
當羅德和亞伯拉罕為了分地.
有一些討論.
最後羅德是怎樣?.

$^{441}$羅德選擇了所多瑪.
因為所多瑪是非常美麗.
十三章第十節.
「羅德舉目看見約大河的全平原.
直到所以.
都是滋潤的.
那地在耶和華未滅.
所多瑪 阿摩拉爾先.
如同耶和華的原子.
也象埃及地」.
所以.
羅德是因為什麼?.
是因為貪戀.
或者因為喜歡.
所多瑪的美地.
而選擇遷移到那裡.
但其實十三章.
作者已經留下了一些提示給我們知道.
那裡已經提到.
「那地在耶和華未滅.
所多瑪 阿摩拉爾先」.
甚至在十三章十三節提到.
「所多瑪人在耶和華面前.
罪大惡極」.
而所多瑪更加是什麼呢?.
所多瑪其實是在迦南境內的.
而大家都應該有印象.
亞伯拉罕應許的地方是什麼?.
就正正是迦南地.
神應許亞伯拉罕所得之地.
就是迦南地.
換句說話來說.
所多瑪其實是亞伯拉罕的敵人.
亞伯拉罕要得到迦南地.
必須要攻陷所多瑪.
才可以得應許地.
如果是這樣的話.
明明所多瑪是罪有應得的.
而且所多瑪是亞伯拉罕的敵人.
為什麼亞伯拉罕要為所多瑪代求呢?.

$^{481}$香港人最有名的是什麼呢?.
就是喜歡吃花生.
既然是這樣的話.
亞伯拉罕最有理由是什麼呢?.
就是在旁邊吃花生,手指子.
活該吧.
就應該有這樣的結局.
為什麼亞伯拉罕要為所多瑪代求呢?.
我們回到十八章前面.
我們看十七到二十節.
「而我所要做的事豈可瞞著亞伯拉罕呢?」.
這個更有趣.
原來亞伯拉罕知道.
神要滅所多瑪,俄姆拉這件事.
是神自己說的.
是神主動告訴亞伯拉罕的.
為什麼要告訴亞伯拉罕這件事呢?.
十八章十八節.
「亞伯拉罕必要成為強大的國.
地上的萬國都必因他得福.
我眷顧他.
為要叫他吩咐他的眾子和他的眷屬.
遵守我的道.
秉公行義.
使我所應許亞伯拉罕的話都成就了」.
又說「所多瑪,俄姆拉的罪惡甚重」.
又再重複說「他們罪有應得」.
所以原來神預先告訴亞伯拉罕.
神自己的計劃.
為什麼呢?.
因為十八章十八節.
「亞伯拉罕必要成為強大的國.
地上的萬國都必因他得福」.
原來這件事令我們想起什麼呢?.
想起亞伯拉罕的呼召.
在《創世記》十二章一到三節.
詳細的我不跟大家看了.
大家可能都有印象.
當時亞伯拉罕從吾爾出發.
被神呼召.

$^{521}$那個呼召的內容是什麼呢?.
就是叫萬國得福.
所以這裡十八章又再重複說這件事.
地上的萬國要因他得福.
反過來說.
亞伯拉罕是神賜福萬國的途徑或者管道.
怎樣叫萬國得福呢?.
十九節.
「我眷顧他」.
或者其實應該說.
「我揀選他」.
即是神揀選亞伯拉罕.
「為要叫他.
吩咐他的眾子和他的眷屬.
遵守我的道.
秉公行義.
使我所應許亞伯拉罕的話都成就了」.
怎樣叫神所應許亞伯拉罕的話都成就了呢?.
就是當亞伯拉罕吩咐他的眾子和他的眷屬.
遵守神的道.
秉公行義.
又是二字.
換句話來說.
神的計劃是什麼呢?.
亞伯拉罕和他的群體.
這裡是吩咐他的眾子和他的眷屬.
不只是亞伯拉罕一個人.
甚至不是他直系的兒孫.
還有他的眷屬.
整個亞伯拉罕群體.
要怎樣呢?.
要秉公行義.
當然要遵守神的道.
神的道其實就是秉公行義.
他們要成為一個義的群體.
要成為一個關鍵.
但是少數的義的群體.
在萬國的面前展現神救恩的計劃.
所以為什麼亞伯拉罕要為索多瑪代求呢?.
因為這就是亞伯拉罕的使命.

$^{561}$其實亞伯拉罕正在回應神給他的使命.
為什麼神要告訴亞伯拉罕.
他要毀滅索多瑪阿摩拉呢?.
因為神當亞伯拉罕是他救贖計劃裡面.
一個很重要的夥伴.
所以神告訴亞伯拉罕.
而亞伯拉罕亦在回應神當初給他的使命.
希望叫萬國得福.
怎樣叫萬國得福呢?.
就是亞伯拉罕和他所屬的群體.
在一個沉淪的世界裡.
在一個罪惡滿盈.
甚至是大多數的萬國的面前.
他們要成為一個正義的大多數.
儘管是異的少數.
但他們卻是關鍵的少數.
這是舊約創世記給我們看到的.
到了新約 原來是一致的.
我們今天被稱為領受了大使命的基督徒.
我們要做什麼呢?.
我們要使萬民.
馬太福音28章告訴我們的大使命.
要使萬民作主的門徒.
所以我們也在領受同樣的命令和呼召.
在萬民面前.
而萬民絕對是大多數.
但我們這群屬神的人.
就要成為一個正義的大多數.
去完成神的任務.
儘管這個任務是不可能的.
因為我們是少數.
但當神以我們為祂的救贖計劃的夥伴.
神用我們.
我們是可以改變世界的.
這就是一個異的群體.
儘管是少數.
但它可以發揮很大的力量.
去逆轉他們身處的地方的命運.
我再說.
神是確定的.

$^{601}$神沒有還價.
每一次阿爸,阿嬸開價.
神是認同的.
說OK OK OK.
問題是什麼是異呢?.
遵守神的道.
秉公行義.
什麼是秉公行義呢?.
當然可以有很多的答案.
我們可以去短算,長算.
在社會裡面有伏羨.
甚至我們信靠神.
我們奉獻等等.
我們可以有很多很多的答案.
但是《創世記》的十八十九章.
我們剛剛讀完這段經文的上文和下理.
其實是有一些很驚訝.
但很具體關於異的答案.
到底什麼是異呢?.
因為時間關係我們看不完.
我們只看十九章.
我們看《創世記》十九章一到第九節.
當阿爸和耶和華的協議.
在十八章完結.
鏡頭一轉.
所多瑪發生什麼事?.
十九章第一節.
「兩個天使晚上到要所多瑪.
羅德正坐在所多瑪城門口.
看見他們就起來迎接.
面伏於地下拜.
說:我主啊.
請你們到僕人家裡洗洗腳.
住一夜.
清早起來再就.
他們說:不.
我們要在街上過夜.
羅德徹徹地請他們.
他們者才進去.
到他屋裡羅德為他們預備筵席.

$^{641}$烤模烤餅.
他們就吃了.
他們說沒有躺下.
所多瑪城裡各處的人.
連老大少都來圍住那房子.
呼叫羅德說:.
今天晚上到你這裡來的人在哪裡呢?.
把他們帶出來.
任我們所為.
羅德出來.
把門關上.
到眾人那裡.
說:眾弟兄.
請你們不要做這惡事.
我有兩個女兒.
還是處女.
容我領他們出來.
任憑你們的心願而行.
只是這兩個人.
既然到我舍下.
不要向他們做什麼.
眾人說:退去吧.
又說這個人來寄居.
還想要作官嗎?.
現在我們要害你.
比害他們更甚.
這段經文.
可能大家都有聽過.
覺得.
羅德為什麼要這樣做人父親.
為了接待客女.
為什麼可以推自己的女兒出去.
或者我們都有聽過.
這個「任他們所為」.
其實或者「任我們所為」.
其實是在說一些同性戀行為等等.
我們可能在這段經文裡.
都會有很多討論.
甚至我相信Mark Sir.
他這麼創世記的專家.

$^{681}$可能都會有一些的資料給大家.
但我今天想.
用一個角度去討論這段經文是什麼呢?.
是很明顯有個對比的.
就是.
這裡對比是什麼?.
是羅德怎樣接待客女.
和所多瑪的眾人.
怎樣接待或者不接待.
甚至加害客女.
是很明顯有個對比的.
我們看到羅德怎樣.
Go all the way.
甚至負上代價.
推自己的女兒.
我們雖然今天不處理這個問題.
為什麼要推自己的女兒出來.
但至少我們可以肯定.
他為了要保護客女.
甚至負極重的代價.
他都願意.
反過來.
反過來所多瑪的眾人是什麼呢?.
他們不單不接待客女.
他們加害客女.
這裡提了兩次.
第七節.
羅德說:眾弟兄,請你們不要做這惡事.
即是這件是惡事.
不接待客女.
甚至任他們所為是一件惡事.
然後第九節.
所多瑪的眾人自己確認.
眾人說:退去吧.
有說者個人來寄居.
還想要作官.
現在我們要害這個害字.
其實亦是惡的字.
你要做這個惡事.
所以我們看到這裡是一個接待的技術.

$^{721}$講羅德怎樣接待客人.
這段技術.
甚至是負上極沉重的代價.
很熱情地招呼這兩個客人.
有趣的地方是什麼呢?.
原來十八章的頭.
亦有另一段技術.
因為時間關係.
今天沒有辦法跟大家詳細地講.
但如果大家有興趣.
可以回去看那段經文.
那段經文是講什麼呢?.
是阿巴拉罕怎樣接待客女.
所以換言之.
是有兩段接待客女的經文.
包圍著.
剛才我們讀的.
《阿巴拉罕為所多瑪代求》的經文.
而這兩段《阿巴拉罕接待客女》.
和十九章的這段《羅德接待客女》.
是非常之近.
非常之類似的.
果然是兩叔侄.
一方面他們都非常積極.
你留意得到.
他們很詳細地描述.
他們怎樣去接待客女.
怎樣大排筵席.
怎樣跪下幫他們洗腳等等.
要求,懇求他們.
去到他們的家.
然後另一方面他們是怎樣呢?.
他們其實是意料之外的.
他們不是知道.
收到電話,收到電郵.
他們有客人要來.
他們是意料之外的.
但是當客女意料之外地來.
他們都擺上最好的款待.
那個客女.

$^{761}$而羅德更加明顯的這段經文是什麼呢?.
是為了接待.
甚至將自己置於危險的境地.
其實他不只是推自己的女兒出去.
其實你留意得到.
在第六節.
他將自己都放在危險的境地.
他將門在後面關上.
然後自己對著那班來勢洶洶的所多瑪眾人.
當然有人可能會問.
是否因為他們知道.
天使是上帝.
所以當然是殷勤呢?.
但是起碼羅德這個技術.
他們似乎是不知道的.
為什麼呢?因為如果羅德知道的話.
根本就不會關上門.
在他後面.
不會在後面.
自己來到一呼當關.
萬呼莫敵地去撞上.
甚至叫天使你來.
但是似乎是不知道的.
所以.
有趣的地方是什麼呢?.
到底什麼是秉公行義呢?.
在討論意義的這段經文的前後.
都是講接待.
都是講怎樣去盡力.
甚至乎代價地接待客人.
而這就是義.
成為所多瑪眾人的相反.
因為所多瑪眾人不單不接待客人.
甚至加害客人.
這成為他們定罪的原因.
義.
等於接待客女.
Hospitality.
這不單是舊約.
不單是十八十九章講的.

$^{801}$如果我們跳到新約.
新約有新約的揣釋.
但是非常類近的.
我們試試多讀兩段經文.
希望會更加清晰.
第一段.
希伯來書十三章.
希伯來書十三章二到三節這樣講.
「不可忘記用愛心接待客女.
因為曾有接待客女的.
不知不覺就接待了天使」.
正正是講這段.
講《挫勢記》十八十九章的這段.
「你們要紀念被捆綁的人.
好像與他們同受捆綁.
又要紀念遭苦害的人.
想到自己也在獄身之內」.
客女是和什麼人相反呢?.
客女也是被捆綁的人.
客女也都同義詞.
或者同類.
是被捆綁的人.
和遭苦害的人.
這些弱勢的群體.
作為神的群體.
作為屬神的人.
是要接待,紀念,紀念.
另一段經文.
《馬太福音》二十五章.
三十一到四十.
比較長.
但是絕對值得我們慢慢讀.
《馬太福音》二十五章三十一節這樣講.
「當人子在他榮耀裡.
同著眾天使降臨的時候.
要坐在他榮耀的寶座上.
萬民都要聚集在他面前.
他要把他們分別出來.
好像牧羊的分別綿羊山羊一般.
把綿羊的安置在右邊.

$^{841}$山羊在左邊.
於是黃向那右邊的雪.
你們這無我父賜福的.
可來承受那創世以來.
為你們所預備的禍.
因為我餓了.
你們給我吃.
渴了 你們給我喝.
我作客裡.
你們留我住.
我赤身露體.
你們給我穿.
我病了.
你們看顧我.
我作監裡.
你們來看我」.
三十七節.
「二人就回答說.
主啊 我什麼時候見你餓了.
給你吃 渴了 給你喝.
什麼時候見你作客裡留你住.
或是赤身露體給你穿.
又是什麼時候見你病了.
或是在監裡來看你呢.
黃耀回答說.
我實在告訴你們.
這些事你們既做在我.
這弟兄中一個最小的身上.
就是做在我身上.
又是接待客裡.
耶穌稱這些為二人.
當祂審判的時候.
當祂坐在祂榮耀的寶座.
祂作為王的時候.
祂這樣說.
作在弟兄中一個最小的身上.
就是作在王身上.
很可惜.
《創世記》十九章.
後面的故事告訴我們什麼呢?.

$^{881}$所多瑪連十個二人都沒有.
十個二人都沒有.
羅德也沒有發揮.
他作為二人的影響力.
他只是僅僅跟他太太.
跟兩個女兒得救.
最後天火焚城.
當然神沒有忘記二人.
神在漠人的會中.
將二人救出來.
因為神始終是一個紀念二人的神.
在今時今日.
我們的亞伯拉罕在哪裡?.
我們的羅德又在哪裡?.
到底我們這群屬神的群體.
是不是一個異己群體?.
有沒有成為關鍵.
發揮力量.
大少數的異己群體呢?.
有一句名言這樣說.
The only thing necessary.
for the triumph of evil.
is for good men to do nothing.
邪惡德性的唯一條件.
是當善良人,當好人.
袖手旁觀.
所以今天.
我們不是問我們與惡的距離有多遠.
或者我們要反過來問自己.
我們離異的距離有多遠.
無論我們身處哪裡.
是香港.
是美國.
是世界其他地方.
加拿大,澳洲.
我們這群屬神的人.
有沒有完成我們從神領袖的使命.
在一個萬國.
但是卻是罪惡滔天.
罪惡勢力強大.

$^{921}$而且是majority的一個這樣的世界.
我們可以發揮一個righteous.
但majority.
卻是有影響力的一個小數群體.
去改變,逆轉我們身處群體的命運.
我們一起祈禱.
在我祈求你加增二人的力量.
在一個不容易的世代.
讓我們學習.
怎樣發揮你讓我們領受.
從你而所領受的異議力量.
關顧我們身邊的人.
特別是那些被綑綁,遭苦害的關鍵.
卻是受虐的群體.
以致我們能夠見到.
你的救恩計劃成就在這個地上.
我們很需要你的憐憫.
與我們同在.
祝福我們.
奉耶穌基督的聖名.
\newpage



\section{}
\label{sec:QzV41v7JIro}
\textbf{Apostles' Creed \#8 (Church and Eternal Life) - Rev. Mark Chan}
\newline
\newline
連結: \href{https://youtube.com/watch?v=QzV41v7JIro}{\texttt{https://youtube.com/watch?v=QzV41v7JIro}} ~~~~ 語音日期: 2022-08-02
\newline
\newline
\hyperref[sec:B4uBevTGH_s]{\small{< < < PREV SERMON < < <}}
~
\hyperref[sec:index_chronic]{\small{[返順時目]}}
~
\hyperref[sec:index_scriptual]{\small{[返順卷目]}}
~
\hyperref[sec:BAdEYU5b_F0]{\small{> > > NEXT SERMON > > >}}
\newline
\newline
$^{1}$Alright, today we're going to work through the Apostle's Creed and we come to the third.
section of the Creed..
As I said, this is a road map of what Christianity is all about..
If you've ever been wondering, okay, what Christians should believe and what the Bible.
says about the Christian faith, this is a good road map..
Twelve articles, twelve statements, English and Chinese I'm giving to you..
And we have talked about I believe in God, the God who exists..
We don't believe in the existence of God..
That's not what the church is all about..
That's not what the Bible is all about..
It's about the God who exists and the God who relates to us, God who loves us, God who.
cares for us..
Alright, so we believe in God, that's the first article..
And then the second section is about Jesus Christ..
There's a whole second section about what He did in history, what He did in our lives,.
right, to save us from sin..
And He is now reigning, seated at the right hand of the Father..
And then the third section, the yellow section, we have been working through since last Sunday..
And we believe in the Holy Spirit, the triune God being with us..
And then today we're going to work through the other two articles, the Holy Universal.
Church and then I believe in the Communion of Saints..
Very interesting statement, Holy Universal Church..
Now if you know a little bit about Latin, behind the word universal, it was Catholic..
So there are different translations of the Apostles' Creed that says, "I believe in the.
Holy Catholic Church.".
Now Catholic does not mean Roman Catholicism, does not mean the Roman Catholic Church..
Catholic is an old word..
It means universal..
It means worldwide..
It means there is something about the Church, God's community, that extends beyond time.
and generation and geography..
And we are all bound together as one family, one Church, capital C, Church..
Holy Catholic Church, Universal Church, right?.
I wonder if you have a thought about there is a special relationship between Christians,.
between those who believe in God, those who profess the name of Jesus..
There is something special connection about, you know, believers in Christ..
Two years ago when the pandemic began, we, you know, full scale, we go online..
And I find it very interesting and largely inspiring to be able to connect with a lot.
of different people around the world through the online streaming and online Zoom meeting.
and all that stuff..

$^{41}$As you know, our YouTube channel was very small..
It was under 1,000 subscribers..
And because of the pandemic, now we have, what's the newest number?.
13K?.
Something like that..
Yeah..
So, we have a lot of connections with other people, you know, from Asia, from West Coast.
and East Coast, even from UK..
I just got to know someone from UK a couple of weeks ago, and she was attending my class.
online..
And then we were connecting and I understand more about her career, what she has been doing.
and she understands more about our church..
In any case, I think there's a special connection that we have as brothers and sisters..
The most common titles, most common labels for believers in the Bible from Old Testament.
to the New Testament is brothers and sisters, family, right?.
And that's what this article is all about..
I believe in the Holy Universal Church..
And I think one of the passages that talks about the unity of believers that extends.
beyond time and geography is Psalm 133..
Psalm 133..
Let's read verse 1..
Psalm 133..
A song of essence of David..
How good and pleasant it is when brothers, of course, including sisters, living together.
in unity, dwelling together in unity..
There is something good about believers getting ready and gathering together, like what we.
do today, singing together, worshipping God together, praying together, caring for one.
another, like what we do in the afternoon..
We have the family gathering, the general assembly to take care of some of the businesses.
and some of the things that are happening in the church..
And we got a photo booth back there..
Thanks to the team who put this up and we can take some pictures for memories..
You know, it's all about the family and brothers and sisters living together, dwelling together..
That is something good in the sight of God..
When God sees God's people getting ready, getting together to worship, to have fellowship.
together, there is something good about just the gathering..
Well, according to the words of Moses in Deuteronomy, he says there are three times in a year that.
the people of God, especially their representatives, the males, back in the ancient society, it.
was the male dominant or patriarchal society..
Of course, that includes the whole family, right?.

$^{81}$But the head of the family is the patriarch, is the male..
So three times in a year, all your males should appear before the Lord your God in the place.
where He chooses..
There were three times, the Feast of Unleavened Bread, which is in our calendar, April or.
May, the Passover, the Good Friday and Easter season..
That's the first time..
Second time, 49 days, seven weeks later, the Feast of the Weeks, the Pentecost..
About May or June, they need to go to Jerusalem again..
And then the third time is the Feast of Sukkot, the Tabernacle, the Booths..
It was about the fourth time, September or October..
They shall appear before the Lord to worship, to take the pilgrimage from where they live.
to the place God has chosen, which is Jerusalem..
And they shall not appear before the Lord empty-handed because they are coming to worship..
Three times in a year..
Now you can imagine in the old days when you had no iPhones, no GPS, no cars, and nothing.
like that..
If you live in Galilee, which is in the north, you see the little freshwater lake up in the.
north where the beginning point of the arrow is..
If you live there, you have to travel three times a year to Jerusalem..
That is about 90 miles of distance..
In the old days, the average walking speed of a human, guess what it is?.
How many miles per hour can you walk?.
Average?.
Three, right?.
Yeah..
Well, give or take two or three, right?.
Two or three miles per hour..
And you have to walk 90 miles..
How many days would that take?.
That is about, if I'm right, about, I'm math, five to six days..
Eight hours a day you just keep walking, right?.
Keep walking and keep walking until you reach your destination, Jerusalem..
It takes about six or even seven days..
And think about how much preparation one has to make to worship God to come to Jerusalem,.
right?.
It takes intentional sacrifice of time and energy and all this, of tons of preparation.
to be in the presence of God, right?.
I believe this passage talks about the people of God need to appear before the Lord in Jerusalem..
It's a reminder for us..
You know, today, as we come before a corporate worship, together, we need to prepare ourselves.

$^{121}$as well, right?.
Even in our modern days, when we prepare for a trip, right?.
I'm leaving for Tennessee for a family trip tomorrow..
Guess how much preparation there is, right?.
Guess how much preparation there is..
The end of August, next month, I'm leading about 10 people to visit the land of Israel.
to study together..
And guess how much intentional preparation we have to make, right?.
A couple of days ago, I received two SIM cards from a company where I need to rent two local.
phones in order to have better communication in Israel..
And we make months of preparation..
Think about what Deuteronomy says..
Think about what Psalm says..
We have to prepare ourselves to worship..
Now, notice the title of the Psalm, 133..
The Psalms of Ascent..
There are 15 Psalms that has the same title, Psalm of Ascent..
Psalm of Ascent means rising up, going up, right?.
Going up..
It comes from the word even Israeli Jews today would use, aliyah..
When they say aliyah, they talk about an immigration process..
Most Jews or Israelis, they are living outside of the land of Israel..
They live in diaspora in Europe or in America..
If they decide to immigrate to the land of Israel, the modern nation of Israel, the whole.
process of immigration is called aliyah..
Now, what does it mean?.
Ascent means going up..
Aliyah means going up, right?.
And why?.
If you check the geography of the land of Israel, Jerusalem is in the mountain area..
It is more than 2,000 feet above sea level..
So wherever you come in, north, south, east, or west, you always go up..
You always go aliyah..
Even today when they immigrate to the nation of Israel, they are going through aliyah..
Have you done your aliyah yet?.
Something like that..
This is a beautiful picture of what worship is all about..
Going up, looking up..
You come here, you don't just look at each other horizontally..
It's a vertical relationship..

$^{161}$There are many different images of Christian journey or Christian life you find in the.
Old Testament or especially in the New Testament..
You can find the image of disciple..
The word disciple means student, right?.
It's a lifelong learning process..
You can find the image of temple..
The Christians are together like stones..
We are living stones..
We are building a temple where God will dwell in the midst of us..
There are different kinds of images, but I think one of the most neglected images that.
a Christian life is all about is the word pilgrimage..
Pilgrimage..
It's a journey..
It's aliyah..
It's going up to somewhere..
It is going to a certain end point or destination..
That may be new to you, but I think Psalm 133 is a song of a sense, one of the fifteen..
It's reminding you and me that worshipping God, relating to God is a journey..
It's an aliyah..
It's a pilgrimage..
You are going somewhere and the pilgrimage, the journey is a sacred journey..
Verse 2, it says, "It is like when people coming together to Jerusalem, worshipping.
together, brothers and sisters.".
Verse 2, it says, "It is like good oil, precious oil poured on the head and running down the.
beer of Aaron down upon the opening or the collar of his robes.".
Yikes..
I don't think it's a beautiful picture..
So you can imagine oil coming down here and then there's a long beer..
I think it's pretty ugly for us as modern people..
But in the ancient time, when they think about Aaron being anointed with so much oil and.
good oil, good precious olive oils, pure oil, they are thinking of the anointing of the.
high priest..
The high priest is the representative of all of Israel, all the people of God..
And the high priest will go into the holy of holies of the temple of Jerusalem once.
a year, only one time a year..
And he is allowed to go into the holy of holies to worship God..
And he, Aaron, the high priest, is actually the representation of every single Israelite..
Every single brother and sister is all upon the high priest..
So anointing him, meaning this is the most sacred climax of the worship life of the whole.
nation, of the whole people of God..

$^{201}$When the high priest being anointed, meaning that we can come before the presence of God..
We can come before God himself and we can worship him..
When Aaron enters the holy of holies, the whole people of God enters with him, go into.
the temple..
And please notice that..
It's always corporate..
Worship is always collective..
In the Bible, from the first book, Genesis, all the way to the last book of the New Testament,.
Revelation, when it talks about worship, it's not just about singing..
It's not just about gathering..
It's a collective experience..
It's a corporate thing to do..
It's never your own, your individualistic exercise..
Now of course you can worship God with your guitar, you can sing to God in your bedroom,.
you can do it by yourself..
That is of course authentic worship, right?.
You can do it yourself..
But when the Bible talks about worship, it's often referred to as something corporate,.
something we can see each other, to come together..
And that is what true worship is all about..
Now I think one of the cultural scripts of our days, cultural scripts, what I mean, cultural.
scripts meaning our culture has different grand narratives, grand ideologies, grand.
ideas..
The ideas are so much affecting how we think and how we act..
But I think one of the cultural scripts today, which is also a cultural right, is that you.
are on your own..
And you can define yourself and you are not related to anyone..
You decide for yourself, your future, your career, your identity, or what not..
You can do it yourself..
It's all DIY, right?.
It's one of our cultural scripts in America..
Individualism..
Who is going to define you?.
Who is going to define your identity?.
You, the person in the mirror..
That is one of the cultural rights..
You singular can change the world..
A couple of months ago I attended a graduation ceremony..
It was not too long ago..
I pay attention to the speech being said there, although nobody paid attention to that..

$^{241}$But I pay attention to what the speaker is talking about..
In the commencement speech, playing the whole entire cultural script, "Graduates, you are.
important..
You can change the world..
Graduates, you are good..
You can change the world all by yourself..
Individualistic.".
It's one of our cultural scripts..
Some years ago, actually some decades ago, Mother Teresa was being interviewed by a British.
journalist called Malcolm McElrich..
Now in the entire interview, it was done in 1968..
And because of that interview, Mother Teresa became internationally known by many other.
people..
It was only because of that interview..
Everybody after that knows Mother Teresa..
And in their conversation, in the interview, it's very interesting..
There was one point..
Mother Teresa said to Malcolm, "Malcolm claimed to be a Christian, but he was never willing.
to join the church..
But he says he's a Christian.".
So Mother Teresa said to Malcolm, she said, "Malcolm, you're a good man..
Why don't you go all the way and become a church member?".
And then Malcolm replied, he said, "Well, to answer you, Mother, I guess God sees that.
I am a good man..
And God needs some good men outside of the church as well as inside of the church.".
And to that, Mother Teresa replied with three words, three English words..
And guess what they are?.
Mother Teresa said, "No, he doesn't.".
And to that, Malcolm could not reply anymore..
And in his autobiography, Malcolm writes, "I could not answer that argument, so I became.
a church member.".
No, you are not the most important person in the world..
We are not..
God is..
God is most important..
And we are His servants..
And God doesn't call us to be successful..
God calls us to be faithful..
Faithful with our time and our gift, our talents, our treasure..
That's all God asks us to do..

$^{281}$We are not the most important people in the world..
But when God sees His people, He sees a family, a church, a family..
Verse 3, "When God's people together, worshipping Him, it is as if the dew of Hermon, which.
is the highest mountain in Israel, the dew, the moisture, falling on Mount Zion, which.
is Jerusalem..
For there the Lord bestows His blessing, even life forevermore.".
Dew of Hermon is the cold moisture, the highest mountain, 9,000 feet above sea level..
Snow-capped mountain..
Moisture coming down to bless the land so that the land will have life, especially in.
summer months, about April to September..
There's no rain in the land of Israel..
It's so dry and so hot..
There's no rain during the summer months..
And in order to survive those months, you need the dew of Hermon..
The morning dew, the mist, the hazy thin moisture that we see early morning, every early morning.
before sunrise..
It is life being sustained by the moisture..
So you see the metaphor here..
When the people of God gather together to worship Him, collectively, cooperatively,.
it is like life we can experience..
Again you believe in Jesus Christ, you become a Christian, right?.
And you have eternal life, right?.
So often we mistakenly thought that the eternal life can be experienced individualistically..
No, the Bible says no..
The Apostles' Creed says no..
We can only experience eternal life, the fullness, when the people of God gather together, sisters.
and brothers, getting ready to worship God..
And that is the beautiful picture about corporate worship..
I think the pandemic has done quite a bit damage to Christian life, to church life,.
to our own life..
Because it forces us to isolate ourselves, right?.
To go through, I mean to abide with the culture script of individualism..
That we are okay as long as we are on our own..
No, the Bible says no..
We are being created as someone connected with the bigger family, a bigger church..
I understand in this generation, many organized religious institutions fell, leaders fell..
I have talked about it in the previous sermon..
Please don't put your faith in any human beings, myself included..
We are all broken..
We need to be redeemed, every one of us..

$^{321}$I like what Martin Luther said..
He said, "Every Christian is 100\% sane, being redeemed by God, forgiven by God..
But he is also 100\% sinner.".
We are all broken and we need the redemption of Jesus Christ offered by the cross..
We need the gospel, all of us..
We need the gospel..
We need the new life, the eternal life of God..
So don't put your faith in human beings..
And don't put your faith, by all means, in any organization, any religious institution..
They are broken as well..
They are broken as broken as we are..
Put your faith in God..
Put your faith in Jesus Christ..
So back to the beautiful picture of corporate worship..
I think the Psalms invites us to re-imagine, to re-envision what worship is all about..
What life in God is all about..
When God's people getting ready together, right, to worship God together..
Ask yourself next time you go to worship..
Ask yourself next week when you come in, alright..
How do you get ready?.
Your body, your mind, and your spirit, right..
Even I as a pastor, sometimes I go to another church to worship because of traveling..
I go to another church to attend the worship service..
I always ask myself before I step into the worship service, before the service starts..
I always pray a simple prayer..
I pray to God, "Prepare my heart..
Prepare my mind..
Prepare my body so that I can meet God in His presence..
I can worship with other believers.".
When I read the Bible together, when I think about faith, think about this, you know, God.
thing, right..
Over and over again I want God to shape my mind, right, and shape my heart..
You know, this psalm is telling us we are coming together as a community..
We are never isolated individuals..
Come together as one people, living together, worshipping God together..
And of course now our high priest is our Lord Jesus Christ..
He is the perfect high priest and He offered Himself as the perfect sacrifice..
To the holy of holies, in Jesus, in Christ, we can come before the very presence of God..
And we can come before God and when we worship Him together, God sees something good in us..
There is something pleasant in His sight..

$^{361}$Let's all pray together..
Father, we give You thanks because You are our good God..
And thank You for Your grace that formed this community, formed Your church, formed Your.
community that is to worship..
And we pray that we never lose sight of the fact that we are worshipping You as one people..
We are worshipping You in Christ and with other sisters and brothers..
And thank You for the opportunity, for the blessing that we can experience when we worship.
You together..
And pray that You will continue to help us and remind us to prepare our hearts and our.
mind and our body so that we can worship You and the worship act itself is pleasing and.
is good in Your sight..
And thank You that we can be Your children, Your sons and daughters..
In Christ's name we pray..
Amen..
[BLANK AUDIO].
\newpage



\section{}
\label{sec:BAdEYU5b_F0}
\textbf{群體之美 (詩 133) - 陳崇基牧師}
\newline
\newline
連結: \href{https://youtube.com/watch?v=BAdEYU5b_F0}{\texttt{https://youtube.com/watch?v=BAdEYU5b\_F0}} ~~~~ 語音日期: 2022-08-02
\newline
\newline
\hyperref[sec:QzV41v7JIro]{\small{< < < PREV SERMON < < <}}
~
\hyperref[sec:index_chronic]{\small{[返順時目]}}
~
\hyperref[sec:index_scriptual]{\small{[返順卷目]}}
~
\hyperref[sec:pMRJMgqJITI]{\small{> > > NEXT SERMON > > >}}
\newline
\newline
$^{1}$今天我們看另外一篇詩篇.
我們一直在看每一篇詩篇.
我們一直在挑選不同的詩歌來看.
詩篇有133篇.
聖經這樣說.
大衛的上行之詩漢.
弟兄和睦同住.
是何等的善 何等的美.
留意有時候詩歌是有標題的.
有title 有superscript.
有些title是特別有意思的.
特別是在這裡說到上行之詩.
Song of Ascent.
是上行 上山的意思.
在耶路撒冷是二千多呎的海拔之上.
在山區.
所以無論什麼方向去耶路撒冷.
都是要爬山 都要上去.
所以這些詩歌都有它自己特定的背景.
背後有一個字.
Alaa 或者 Aliyah.
今天現代的以色列人 猶太人.
如果他們散居在世界各地不同的地方.
他們想移民回到現代的以色列國.
Modern Israel 的話.
他們整個移民的process.
就叫做Aliyah.
Aliyah是什麼意思.
就是上 高up的意思.
所以這個很有意思.
去到耶路撒冷的時候永遠都是高up.
大家熟悉聖經的話就知道.
有《上行之詩》這個title的詩歌.
有十五首.
120至134.
我們正在看的是尾二的那首.
133首詩篇.
這裡告訴我們.
這些以色列人.
可能是他們歷世歷代.

$^{41}$他們會上耶路撒冷朝聖.
上耶路撒冷去敬拜的時候.
他們會累積了一些這樣的詩歌或者歌曲.
在他們上路的時候.
上耶路撒冷去敬拜的時候.
他們可能就會唱出.
因為在摩西律法裡面.
《申命記》十六章十六節.
那裡說什麼.
他說你們所有的男丁.
男丁就是代表整個家庭.
或者有時候帶著家人.
要在隨孝節.
就是Passover.
大概是現在的四五月的時候.
七七節.
就是Pentecost.
或者Shavuot.
就是他們修葛節.
就是七七四九一之後.
Pentecost第五十日.
七七節那裡.
他們也要上去耶路撒冷.
然後在Sukkot.
駐彭節的時候.
大概現在九月十月的時候.
一年的三次.
要上到耶路撒冷所揀選的那個地方.
有耶路撒冷名字居住的那個地方.
來去朝見他.
也不可空手來去敬拜上帝.
今天我們的科技很發達.
我們上行.
或者去一個地方旅遊.
都要做很多準備.
雖然我們已經有很多科技.
昨天我剛剛收到.
就是我們八月底的時候.
我們有些弟兄姊妹會.
我們一起會去以色列兩個星期.

$^{81}$我昨天才收到那張SIM卡.
因為要做一張SIM卡.
要去那裡有一個本地的電話.
是方便一點的.
都等了兩個星期我才收到.
有很多東西要準備.
你想想你去旅遊的時候.
有多少東西要收拾.
你想像一下以前古以色列的時代.
舊約的時代.
他們需要預備的東西是多少.
需要準備的東西是多少.
如果你想像一下你在加利利居住的話.
在北面.
在加利利湖附近.
你要上行去耶路撒冷的話.
你要走差不多五至六天的路程.
中間要停不同的村莊等等.
要借宿等等.
要做很多很多不同的準備.
這十五首上行之詩.
可能就是他們上行去耶路撒冷的時候.
盛傳下來的一些詩歌.
好像我們有很多詩歌.
我們歷史歷代都會唱或者會記得.
會使用的.
上耶路撒冷去敬拜.
聚集.
他說何等的善 何等的美.
我相信我們疫情到所謂結尾了.
還是反反覆覆.
大家還是有不同的憂慮.
我想這種聚會.
已經打斷了.
對不對.
或者我們養成了一些不同的習慣.
聚會是不是一種失落了的藝術.
或者操練.
Lost art.
一種Lost discipline.

$^{121}$做什麼你都要聚會.
為什麼原因你都要做.
詩篇133篇或者上行之詩.
就是給我們另外一個角度來看.
作為一個基督徒.
作為一個敬拜上帝的子民.
到底最重要的是什麼.
上行之詩.
Aliyah.
Go up.
去敬拜上帝.
去崇拜真神.
聖經裡面有很多不同的圖畫.
來形容一個Christian life.
一個journey.
可以是一個Disciple.
Disciple意思就是說.
門徒是門生學生的意思.
有基督的新娘,新婦.
代表預備好出嫁等待.
聖殿.
我們基督徒是活石.
能夠組成一起.
建立一個宮殿.
讓上帝能夠在我們中間同在.
等等這些不同的圖畫.
第一個我們很少留意的就是.
Pilgrimage.
朝聖.
上路.
這個是比較不是靜態的.
動態的一種畫面.
來形容我們作為一個基督徒.
是應該是一個怎樣的群體.
你看看詩篇133篇.
那裡說兄弟姐妹同住.
一起來去守節.
去敬拜上帝.
去耶路撒冷去敬拜上帝.
在上帝所揀選的地方,時間.

$^{161}$那裡一同聚集.
是河南的善和美.
你留意那裡不是說兄弟姐妹一起同住.
是因為血統一樣.
是因為種族一樣.
是因為我們語言一樣.
或者食物一樣等等.
這些很多文化的標籤.
那裡說到什麼呢.
那裡說到一些神的子民能夠一起聚集.
上帝捍衛美好的事.
很多都不是.
原因.
為什麼捍衛善,捍衛美.
都不是一些猶太人.
他們最著緊或者最引以為傲的身份的標籤.
為什麼捍衛善,捍衛美呢.
Good and Pleasant.
為什麼上帝會這樣看.
是因為他們敬拜的緣故.
是因為他們聚集一起去敬拜.
一起去敬拜同一位的上主.
這個就是捍衛美,捍衛善.
第一字.
同住的同字.
是一下.
Together.
一起走在一起.
這個字是說到一起工作,一起打仗.
一起同一個目的.
一起傳福音等等.
一起事奉等等.
這個捍衛善,捍衛美.
如果你記得有一樣東西叫做寫開古卷.
寫開古卷就是一幫人.
他們離開了耶路撒冷.
他們不滿耶路撒冷.
然後他們就搬到一個地方叫做昆瓦.
在曠野裡面.
他們就一起聚集.

$^{201}$組成了另外一個群體.
在那裡敬拜神.
在那裡抄寫聖經等等.
他們就叫他們自己什麼呢.
Yahad.
Yahad Community.
就是一個合一的一個Community.
一個立約的群體聚首聚集的時候.
就是一個這樣的社團.
捍衛美,捍衛善.
Good.
Talk.
你記得勝經裡面第一次提到.
Good.
是什麼時候.
創造的時候.
上帝創造萬有天地萬物的時候.
他捍衛好.
是就者人成了.
God saw it was good.
居然我們原來聚集在一起的時候.
兄弟姐妹聚集在一起.
同心敬拜的時候.
上帝就好像看著他原本的創造美好一樣.
就是看到我們一起.
走在一起的時候.
就是捍衛美好.
在他眼中.
是Good.
創世的第二章.
你記得有個不好的.
是什麼.
人獨居不好.
人孤獨是不好的.
因為人被創造就是有個Community.
上帝的救恩來到的時候.
就是要創立一個新的Community.
一個孤獨的人.
或者孤立的信徒.
是不好的.

$^{241}$Not good.
所以上帝要有教會.
才能夠救.
用一個很簡單的比喻.
以前經常用的.
就是木柴.
木柴要堆在一起才能夠燒得紅.
火才能夠燒得黃.
把木柴放在旁邊的時候.
可以燒就燒.
沒多久它就會熄滅.
比喻我們每個人都是廢柴.
但是我們走在一起的時候.
就燒得黃.
燒得光亮.
教會 弟兄姊妹.
同住一起去敬拜.
捍衛美 捍衛事.
一起去參與一個教會.
加入一個教會.
去commit to一個教會.
一個local的faith community.
一個信仰的團體.
一個教會的家庭.
是好事.
所以為什麼我們今天有家事大會.
不是因為我們沒事做.
所以加一些這樣的環節.
而是因為這是我們的家.
這個不是一個酒店.
而是屬於我們的一個地方.
我想在現代的社會裡面.
現代的文化裡面.
其中一個方話.
one of the lies.
一個方話.
是我們有時候會buy in的.
一個方話是什麼呢.
就是個人主義.
individualism.

$^{281}$就是我們自己單獨的就ok了.
單獨我們就可以改變世界.
你是多麼重要.
你是一個individual.
你是多麼的尊貴.
你是一個單一的單元.
所以你自己一個已經很重要了.
你自己一個已經夠了.
這個是one of the lies.
如果你去一些畢業禮.
我早前去了一個畢業禮.
裡面的畢業的致詞.
結語的致詞.
都是說一樣東西.
you can change the world.
you are important.
但是如果你細心想一想.
這種文化劇本叫做cultural script.
其實是一種方話來的.
誰說你是很重要的.
上帝才是最重要的.
我們只是他僕人.
我們單一孤獨是不行的.
因為我們是屬於一個群體.
屬於一個group.
去commit to a group.
以前在幾十年前.
大家可能聽過一個名字叫做Mother Teresa.
不難修女有聽過.
無論是年輕的也好.
年長的也好.
都有聽過Mother Teresa.
Mother Teresa在1997年的時候.
她離世了.
她在1968年的時候.
有一個訪問.
英國有一個很重要的傳媒工作者.
那個人叫做Malcolm Rookeridge.
Malcolm Rookeridge是一個英國.
曾經在BBC做過很久的一個journalist.

$^{321}$他在1968年的時候.
就和Mother Teresa做了一個個人的訪問.
那個訪問.
是因為那個訪問.
才使得Mother Teresa揚名國際.
之前是直直無名.
沒有人知道誰是得難修女.
在Calcutta做了多少事情.
1968年他去訪問她的時候.
才使得Mother Teresa成為一個國際的名人.
很多人就明白.
她在Calcutta做過些什麼.
有多重要.
在訪問裡面有一個對話很有趣.
得難修女就和Malcolm說.
Malcolm, you are a good man.
Why don't you join the church?.
You have a question.
Why don't you join the church?.
因為Malcolm本身自己都是professor的question.
但他就不加入教會.
Mother Teresa這樣問他的時候.
Malcolm就回答他.
他說.
我想在上帝的眼中.
我是一個好人.
在教會裡面.
上帝需要一些好人.
但在教會以外.
outside the church.
上帝都需要一些好人.
所以他就說不加入教會.
得難修女用了三個字來回答他.
你猜猜是什麼字.
三個字 英文字.
No, he doesn't.
就是說上帝不需要你在外面做一個好人.
Malcolm聽了得難修女的回答之後.
他就無言以對.
在他自己的自傳裡面這樣說.

$^{361}$他都回答不了Mother Teresa.
於是他就加入了教會.
You are not important.
You need to join a group.
join a church.
因為只有在一個群體裡面.
才見到上帝的心意更加明顯.
神的計劃裡面.
其中一個最重要的一環.
就是教會.
舊約就是以色列的民族.
上帝的子民.
舊約都算是教會.
新約就更加了.
你看看新約裡面有多少的書信.
是和教會生活有關.
和教會的一個群體有關.
但是所有這些書信都是寫給眾數群體的.
不是寫給一個人.
大部分都是寫給整個的family.
整個的spiritual family.
一個屬靈的群體.
一個的教會.
《詩篇》第183篇第二節.
說到群體的時候.
他說一起敬拜的時候.
就好像那貴重的油.
飄在亞倫的頭上.
流到胡蘇.
又流到他的衣襟.
他說像亞倫被高納作大祭司的情景.
你知道每一個群體都有代表性的人物.
每一個群體.
每一個時代.
都有他最代表性的人物.
近日大家知道.
我們有很多都是香港背景的人.
香港在MIRROR的演唱會發生不幸的事故.
大家都很緊張.
為什麼.

$^{401}$因為在這個時代.
特別是香港後創商的一個時代.
就有一個代表性的人.
或者代表性的一班人.
他們就代表了他的心聲也好.
代表香港的一些情懷也好.
總是有一個代表性的人物.
所以這些代表性的人物.
或者代表性的事件.
有一些不幸的事發生.
一些大的事故.
就成為整個集體的一個經歷.
一個群體的一個經歷.
以前的我們比較年長一點的.
你記得了.
Beyond的黃家駒.
他開演唱會的時候.
頭兩天經常都.
最近這個MIRROR.
開演唱會頭兩天經常有這樣的小意外.
我就經常想起什麼呢.
就想起黃家駒.
Beyond黃家駒.
或者想起柯基哥哥.
或者那些梅艷芳等等.
不同的代表性的人物.
那些盛極一時的時候.
有一種集體的代表.
對於這些人來說.
夠了一件事.
他們回望最代表性的一個人物.
或者一件事情.
就是亞倫.
大祭司被高瀾.
特別第一次被高瀾的時候.
你記得摩西帶領以色列人出了埃及之後.
設立了回望.
所以King Crown那些東西.
那些回望的工具.
器皿.

$^{441}$全部東西都高瀾了.
然後潔淨了.
用水潔淨.
用油潔淨等等.
到最後最高潮的時候.
就是什麼.
你還沒看到第八第九章.
就是要高瀾大祭司.
因為那些回望都是工具.
工具來的.
是媒介來的.
誰才能夠代表我們.
能夠去到上帝的面前.
能夠去獻祭.
就是大祭司.
大祭司一年只有一次.
進到他的致政所.
來代表整個的以色列人民.
所以他們回想的時候.
就發現這個就是最高潮的時候.
就是我們一起來敬拜的時候.
就有如這個最高潮.
高瀾這個High Priest.
最神聖最secret的時候.
原來一起敬拜.
一起去所謂同住.
一起去耶路撒冷.
來去敬拜上帝.
就是一件這樣的事.
你還沒看到第九章.
這裡最高潮的時候發生什麼事.
阿倫最亂 遵了護目.
摩西和阿倫出來.
然後為百姓祝福.
為什麼說這裡是高潮.
在這裡才看到.
耶和華的榮光.
向全體的百姓去顯現.
有火降臨.
然後燒毀了.

$^{481}$焚燒了這些祭物.
全體百姓一起見到的時候.
就面服於地來歡呼喝采.
這是最高潮最climax的時候.
神的榮光降臨.
是怎樣降臨.
《詩篇》133篇.
弟兄姊妹合一.
一起去想去敬拜的時候.
就見到上帝的榮光彰顯出來.
第三節又好比.
黑門的金爐降災石安山.
在那裡有耶和華所命定的福.
就是永遠的生命.
黑門山是以色列最高最高的山峰.
九千多呎.
有另一個別名叫做白頭山.
我們老了的.
有些更加年長的人.
就叫白頭了.
為什麼呢.
為什麼叫白頭山.
因為是長年積雪的.
在這個地方長年積雪.
就是很多時候.
南風吹下來的時候.
就可以滋潤整個的以色列地.
特別每年夏季的時候.
四月到九月初的時候.
最炎熱的夏日.
最乾旱的夏日的時候.
那些農作物.
因為在四月到九月的時候.
是不會有雨的.
既然沒有雨.
那些農作物就很容易死.
很容易會枯乾.
會渴死.
除非有朝早.
天矇江.

$^{521}$你有沒有試過晨運天矇江的時候.
出過去後一.
一陣的薄霧.
那些就叫做甘露.
那些甘露是從哪裡來的.
最重要的地方就是在黑門山吹下來.
滋潤石安野老山.
聖文集體敬拜.
原來就是經歷到.
在夏天乾旱的時候.
生命的降臨.
永生.
這個也是我們的問題.
我們說信耶穌.
我們就有永生.
這個就像個人主義.
我信了.
我就有永生.
OK了.
詩篇133.
整個聖文在一起的時候.
才經歷到有永生.
耶和華所命名的福就是永遠的生命.
才能夠經歷到上帝所給的永生.
一起敬拜的時候.
才可以見到.
原來這個是神聖的時候.
這個是享受永生的時候.
你今天來到敬拜.
我想我們也要問一問我們的問題.
到底我們的期望是什麼.
Right?.
這個期望就很直接的去影響.
到底我們拿些什麼走.
很多時候都是聽一下演講的講員講得好不好.
講得好就給他五顆星.
講得不好就不要了.
這個都是期望的一部分.
但真正詩篇133講的期望是什麼.
就是你要知道這個時候.

$^{561}$敬拜的時候.
是一個神聖的時刻.
這個時候.
是一個永生的時刻.
你只能夠經歷永生.
Fullness.
滿足地能夠經歷到什麼叫做Eternal Life.
Eternity.
能夠一起去見到大家.
去到聚集的時候.
詩篇133就是.
這看為好.
和Sacred 是Holy.
我最近在英文堂講的一個系列是《思路信經》.
我自己以前也沒有留意過.
《思路信經》差不多有2000年的一個傳統.
教會裡面所應信的這些條目.
大概有12條.
我最近才留意到原來有兩條.
足足有兩條.
12條很短的條目.
就是我信上帝創造天地萬物主.
我信基督等等.
我信聖靈.
原來兩條的條目.
都是和Community有關.
我信聖而公之教會.
聖徒相同.
兩個的Article.
聖而公的教會.
英文有時候翻譯做Catholic Church.
I believe in the Holy Catholic Church.
Catholic Church不是說到羅馬天主教.
因為Catholic這個字是大公的意思.
就是說是超越時空.
超越世代.
超越地理的.
我信聖而公之教會.
就是說當我們聚集的時候.
我們相信上帝是我們的宗教.

$^{601}$是一個Sacred的Gathering.
是一個Eternal.
Eternity.
Taste of Eternity的Gathering.
我信聖而公之教會.
我信聖徒相同.
Communion of Saints.
被耶穌基督所拯救的人.
我們就能夠成為.
有一種特別的Relationship.
特別的關係.
詩篇133篇.
弟兄姊妹.
一齊來敬拜.
是一個何等善 何等美的事.
在上帝眼中.
無論你怎麼看.
在上帝眼中是一件Good thing.
Pleasant thing.
我們現在有的大祭司.
當然不是阿倫.
因為他所預表的就是耶穌基督.
獻上完美的祭.
我們的主耶穌基督.
已經是最完美的大祭司.
帶領我們能夠來到上帝的面前.
信了耶穌的我們成為上帝的子民.
在基督裡面就是一個真的耶路撒冷.
在基督裡面就能夠與上帝遇見.
但這個絕對不是一個個人.
孤立的人.
能夠享受的.
只有教會能夠聚集在一起的時候.
才能夠見到上帝的榮光.
在我們中間.
求主去調教一下我們的預期.
我們的Expectations.
如果你來到教會去敬拜的時候.
你不單只需要身心靈準備.
當然.

$^{641}$但我想問你.
問清楚到底你自己的問題是什麼.
到底你想得到什麼.
我每次來教會崇拜的時候.
無論是這個教會.
或者我去第二個教會.
我自己都調教一下我的Expectations.
我自己的預期.
我的預期.
是我希望能夠在這裡事奉上帝.
我希望能夠在這裡經歷到永生.
詩篇133.
作為一個鼓勵.
我們同心得道祈禱.
主我們多謝你.
因為你的恩典.
你將我們構成一個屬靈的家庭.
主我們成會教會.
就是你自己榮光能夠降臨的地方.
當我們來到聚集.
在基督裡面來到聚集.
我們就享受到.
上帝你自己所命定的福.
永遠的生命.
主我們感恩讚你.
因為這個敬拜.
我們每一個主日的敬拜.
我們實在可以這樣經歷.
求你自己聖靈來提醒.
幫助我們.
見到弟兄姊妹的時候.
我們一齊同心來到敬拜的時候.
我們要見到的是你的同眾.
要享受到的就是你自己的同眾.
主我們感恩.
我們待會下午都會有一個家事大會.
都是詩篇133篇裡面所講的.
弟兄姊妹能夠同住.
一起關心教會.
為教會禱告.

$^{681}$主我們感恩.
我們能夠有這樣的機會.
求你都祝福每一個本地教會的家庭.
讓弟兄姊妹同心來到聚集的時候.
你自己的榮光能夠彰顯出來.
\newpage



\section{約翰福音 4:43-54}
\label{sec:pMRJMgqJITI}
\textbf{信你就奇 (約翰福音 4\_43-54) - 陳德晶傳道}
\newline
\newline
連結: \href{https://youtube.com/watch?v=pMRJMgqJITI}{\texttt{https://youtube.com/watch?v=pMRJMgqJITI}} ~~~~ 語音日期: 2022-08-10
\newline
\newline
\hyperref[sec:BAdEYU5b_F0]{\small{< < < PREV SERMON < < <}}
~
\hyperref[sec:index_chronic]{\small{[返順時目]}}
~
\hyperref[sec:index_scriptual]{\small{[返順卷目]}}
~
\hyperref[sec:KfYBQZvmZ4Y]{\small{> > > NEXT SERMON > > >}}
\newline
\newline
約翰福音 4:43-54
\newline
\begin{longtable}{cl}
\hline
\hline
章節 & 經文 (和合本修訂版)\\
\hline
4:43 & \begin{tabularx}{0.7\textwidth}{X} 過了那兩天，耶穌離開那地方，往加利利去。 \end{tabularx} \\ \\ \relax
4:44 & \begin{tabularx}{0.7\textwidth}{X} 因為耶穌自己作過見證說：「先知在自己的家鄉是沒有人尊敬的。」 \end{tabularx} \\ \\ \relax
4:45 & \begin{tabularx}{0.7\textwidth}{X} 到了加利利，加利利人都歡迎他，因為他們也上耶路撒冷去過節，曾經看過他在節期間所做的一切事。 \end{tabularx} \\ \\ \relax
4:46 & \begin{tabularx}{0.7\textwidth}{X} 耶穌又到了加利利的迦拿，就是他從前變水為酒的地方。有一個大臣，他的兒子在迦百農病了。 \end{tabularx} \\ \\ \relax
4:47 & \begin{tabularx}{0.7\textwidth}{X} 他聽見耶穌從猶太到了加利利，就來見他，求他下去醫治他的兒子，因為他兒子快要死了。 \end{tabularx} \\ \\ \relax
4:48 & \begin{tabularx}{0.7\textwidth}{X} 耶穌對他說：「若不看見神蹟奇事，你們總是不信。」 \end{tabularx} \\ \\ \relax
4:49 & \begin{tabularx}{0.7\textwidth}{X} 那大臣對他說：「先生，求你趁著我的孩子還沒有死就下去吧。」 \end{tabularx} \\ \\ \relax
4:50 & \begin{tabularx}{0.7\textwidth}{X} 耶穌對他說：「回去吧，你的兒子會活！」那人信耶穌所說的話，就回去了。 \end{tabularx} \\ \\ \relax
4:51 & \begin{tabularx}{0.7\textwidth}{X} 正下去的時候，他的僕人迎面而來，說他的兒子活了。 \end{tabularx} \\ \\ \relax
4:52 & \begin{tabularx}{0.7\textwidth}{X} 他就問甚麼時候見好的。他們對他說：「昨天下午一點鐘熱就退了。」 \end{tabularx} \\ \\ \relax
4:53 & \begin{tabularx}{0.7\textwidth}{X} 他就知道這正是耶穌對他說「你的兒子會活」的時候；他自己和全家就都信了。 \end{tabularx} \\ \\ \relax
4:54 & \begin{tabularx}{0.7\textwidth}{X} 這是耶穌從猶太回到加利利後所行的第二個神蹟。 \end{tabularx} \\ \\
[1ex]
\hline
\hline
\end{longtable}
$^{1}$大宇姐妹朋友,開始的時候想考考大家.
有一條線是近年香港人很希望自己身上發生擁有的.
你知不知道是什麼呢?.
可能你會想,沒有的呢,就即時想到.
有?有點怪,什麼來的呢?.
或許美國的大宇姐妹不太熟悉香港的特色.
所以你想不到,其實這件事很簡單.
就是快插棒上面那一條線.
為什麼呢?因為這條線是一個Sign.
是顯明你沒有中招.
你是否很想擁有呢?.
或者你想想香港人多想擁有,是不是?.
我對這個Sign就有很深刻的體會.
因為我中過,我身邊的人都中過.
那時候我們中招的時候,天天都很渴望.
天天都在測試,都很渴望一件事.
有一條線,要出現.
我最慘的是,在第七天的時候.
仍然是沒有的.
那一刻,你知不知道.
我就在心裡想,我是不是可以放監的.
我什麼時候才可以放監的.
很沮喪,很想哭出來.
我想其實今時今日.
有另一個Sign.
我們在這個荒謬,混亂,失控.
令人很辛苦,無望的世界裡面.
這個Sign,我們同樣很渴望它出現.
不過可惜等來等去都見不到.
很沮喪,信祂才奇怪.
有一個這樣出現.
這個Sign是什麼呢?.
就是神蹟.
因為神蹟在原文的意思就是.
Sign,一個色諾在當中.
它的目的就是要顯明.
耶穌的身份就是神的兒子.
是神自己.
在《約翰福音》裡面出現過17次.
今天很想和大家去看《約翰福音》4章43至54節.

$^{41}$這個段落裡面.
這個Sign出現過兩次.
甚至約翰在最後的時候.
54節,他特意要補充.
是第二次的Sign.
這個神蹟出現.
所以你會看到,其實在這一段裡面.
神蹟這個Sign是一個重點.
但到底神蹟怎樣發生呢?.
我們一起看看.
如果你看下去,你就會發現.
其實發生得很奇怪的.
很怪的.
耶穌極少為自己作見證.
但一做的時候,就很「咩」.
「咩」到不得了.
44節,他就說.
「先知在本地是沒有人尊敬的」.
尊敬或者叫尊貴.
是一個心態,一個態度.
不少的學者都指出.
他在批評當時將信心.
單單建基於一些神蹟上的人.
對耶穌來說.
簡單來說.
他們對耶穌是信祂才會信.
不可能會發生在那裡.
但最怪的是什麼呢?.
就是他竟然正正就是要去.
這一群人的地方.
去本地加利利的地方.
從哪裡去呢?.
這個地方你很熟悉的.
撒瑪利亞.
你有印象.
《約翰福音》四章.
撒瑪利亞的出現.
就是耶穌和撒瑪利亞婦人.
見面有對談的時候.
之後撒瑪利亞婦人.

$^{81}$就向當地的人介紹耶穌.
和祂做見證.
當地人相當歡迎耶穌.
這個情況像什麼呢?.
為什麼我說怪呢?.
你想想.
你現在在A公司做得好好的.
每個人都很欣賞你做的一切.
而且也有很多機會給你去發揮.
但是B公司呢.
那些人是看不起你的.
信你才行.
就是會很想你快點死掉.
不斷想方法去針對你.
好了 我問你.
你會不會現在辭職.
由A公司走去B公司.
當然不會.
是吧.
這麼笨的嗎?.
但是你看耶穌.
祂只是留了兩天.
在一個祂受歡迎的地方.
就走去一個不受歡迎的加利利.
你說是不是怪呢?.
但是為什麼要這樣呢?.
我猜祂一定有些事情想做.
或者祂想去證明一些事情.
45節 祂果然是這樣.
耶穌去到加利利.
證明了情況變了.
當地人的心態.
從不去敬畏祂.
尊敬祂.
變成了接待祂.
新日本的譯法就是「歡迎」.
由一個不尊敬去歡迎祂.
因為他們曾經上耶路撒冷.
去過各節.
見過祂所行的一切事.

$^{121}$當中你見到.
這節經文說.
重點是什麼?.
「體見」.
在《楊漢福音》.
「體見」這個字出現了28次.
比其他的福音書更多.
是一個重點.
原來我們會見到.
「體見」很重要.
因為「體見」神跡.
就會顯在我們的眼前.
我們繼續看下去.
46節.
耶穌又來到一個地方.
你很熟悉的.
加利利的加勒.
加勒這個地方.
令你怎樣也會想起一些事.
但是約翰爾這裡更加厲害.
他明知怕你忘記.
他補充了一句說話.
他說從前.
變水圍走的一個地方.
提醒你的.
特意好刻意告訴你.
耶穌在當中行神跡.
並且好像意味著.
耶穌應該會再有神跡在這裡發生.
所以鏡頭一轉.
就有一個大神出現.
大神原文的意思是.
一個皇室或者屬於皇帝的.
Royal的人.
可能是分封的.
王希律安提柏的神父在當中.
無論如何.
Royal這個字.
你一聽到都覺得很尊貴.
但是尊貴得來又怎樣?.

$^{161}$都不能夠逃避.
一些荒謬,混亂,不如意的事情.
發生在自己的心上.
這次就來到他的遺址.
他的遺址.
去重病.
睡危在當中.
當面對親人睡危.
當然很迫切.
很熱誠.
你見到相當熱誠的.
他就由加勒林.
走了十七公里.
走到去找耶穌.
如果你看回52節.
更加厲害的.
其實那個路程真的相當厲害.
是走了一天一天.
是吧?.
一個很長的旅程.
特意要長途跋涉去找耶穌.
但是尊貴一件事是什麼呢?.
不是去.
這樣走過去.
而是你面對一個誠意拳拳.
走了一天路的人來到.
也都見到他的親人就快死了.
給了你 你會怎樣做?.
給了我們怎樣都會同情他.
更加耶穌在當中.
理應很欣賞他 很關心他.
但你看看耶穌又說什麼?.
又是一個很敏感的對話.
「若不看見神跡奇事.
你們總是不信」.
用不著這麼敏感吧.
雖然這裡的敏感.
說明是一個眾數.
理無應該也是針對當時的本地人.
不敬不信的那一群人.

$^{201}$但在當中也不用這麼敏感吧?.
我就去想一下.
我的為人其實都是一些.
很有衝勁 但昂直的人.
所以很多朋友都會擔心我不知道.
往往會串通我.
很敏感地對我說話.
其實想做什麼?.
提醒我.
我這樣想的時候.
我們再看看大神的言行.
聖經說.
「去求他下去醫治他的兒子」.
他的心態只是一樣.
按照自己的方法.
「耶穌 你跟我下去」.
是自己作主.
但你不是忘記了.
神跡是誰做主導的?.
神跡是要顯明.
耶穌是神的兒子.
是由神作主的.
一切都是跟隨耶穌的意思.
不是跟隨人的意思.
所以耶穌站在那裡.
似乎也在提醒或挑戰大神的信心.
信心在《仰古音四章》出現過七次.
單單這個段落也有三次.
這是很明顯.
在這段經文裡.
很想我們知道的重點.
信心其實是要說對象.
即是你信任哪一個.
而不是說你信會發生什麼事.
結果會是怎樣.
或是會否跟隨你的安排和預期.
所以如果這樣說.
如果你想將神跡顯明在你眼前.
其實你很需要一件事.
就是你要信心的對象對得準.

$^{241}$這樣神跡就會顯在大家的眼前.
到了49節.
大神的回應是.
「先生求你.
趁著我的兒子還沒有死的時候就下去」.
先生是相當尊重的.
但是這裡又是再重覆一次.
「要下去」.
似乎他很心急 很迫切.
聽不明白耶穌提他的說話.
怪不得 又是耶穌古怪地.
又是來一次很刻薄地命令他.
「回去吧 你的兒子活了」.
耶穌雖然刻薄 很難聽.
但很有用.
大神終於在當中聽得明白.
他沒有再死纏爛打.
他終於明白一件事.
是神說了算的.
是去命令人的.
是去發施號令的.
是去說怎樣就怎樣的.
但是耶穌這一句說話.
聖經學者告訴我們.
除了是一種命令之外.
其實是有一種應許.
還有帶著一些憐憫的說話.
所以在當中大神知道.
並且放下自己.
沒有再按自己的意思.
去決定相信耶穌是可以的.
他就馴服他.
所以在五十節他說.
「那人信耶穌所說的話 就回去了」.
跟隨耶穌的話 回去了.
所以在當中他放下了自己的心急.
願意去相信 馴服.
神蹟就顯在他的面前.
故事來到這裡還有下文.
補充了一段.

$^{281}$一個僕人去迎接一個大神.
他們之間的對話.
你看看 五十三節在當中.
他說「去便知道.
這正是耶穌對他說.
你的兒子沒了的時候.
他自己和他全家就都信了.
哇 大團圓結局.
在當中的時候.
你就會想起香港的電視就會說.
我們一起歡喜 一起去BBQ.
但是你看到這裡很有趣.
耶穌的反應.
或者應該說這個大神的反應.
不是去BBQ 不是去開筵席.
卻是再說多一次.
去告訴所有人.
耶穌在當中行了神蹟.
而且他是甘心命底地信去就竊.
意思不是疑惑.
而是信耶穌你就是有神蹟發生.
不過你可能會問.
信你才是.
我以為今天聽到會有些意見.
每次聽到這個教會這個網台.
總是有些意見在當中.
我都以為今天可以聽到一些意見.
讓我可以知道.
我在一個這麼荒謬.
這麼混亂的世界裡面.
這麼辛苦的時候.
可以等到一個神蹟.
可以知道可以怎樣去做.
除了不用向現實低頭跪下之外.
還有什麼出路.
你可能都會這樣想.
但是怎樣呢.
我講完那些話.
都不知道媽媽是女人嗎.
信耶穌得永生.

$^{321}$信耶穌有奇蹟.
浪費我的時間嗎.
其實我很明白大家在想什麼.
因為當我認識了經之後.
我都在問自己.
我是不是在浪費時間.
到底有什麼是他.
我用不著去轉變.
但是我作為一個傳道人.
不是叫人信做什麼.
我經常想我們最後的信仰.
不是唱出來的.
而是要去堅持.
我更深信一件事.
就是聖賢是有他的能力.
文本是會說話的.
我們有疑惑.
不是主耶穌的說話有問題.
而是我未明白他的說話.
所以我反覆去想一下.
耶穌幾句的說話.
他為什麼要這麼刻薄.
想說什麼.
我發覺都真的不是簡單的.
他似乎是一步又一步再一步.
去幫這個大神.
越來越有信心.
讓我再去highlight出來.
幫大家在荒謬不如意的世界裡面.
我們可以越來越有信心.
第一步.
耶穌在這裡.
令人由不敬到歡迎他.
不論是加里尼還是大神.
起初的時候.
你不尊敬他,不相信他.
但是之後他能夠歡迎他.
主動去找他.
為什麼?.
你看到是吧?.

$^{361}$你記得的.
他肯去看見.
我在牧養的過程裡.
經常都看到一些很crazy.
很混亂,很荒謬.
不如意的事情發生.
我這個星期都處理了不少.
其中一個要處理一個家庭的需要.
爸爸媽媽的兒子和我和貞德哥醫生一起談.
談了一個小時.
竟然是不歡而散.
家人很不好意思.
覺得要搞我.
而且大家肚子餓了.
就抓了他去吃飯.
在飯菜之間都很特別.
媽媽就開始說.
這幾年發生了很多事.
我整個家庭都很混亂.
我自己也是.
很辛苦.
神派了你.
輔導員,醫生.
幫助我們一家.
如果不是這樣.
我相信我們一家人.
是一個又一個跳下去的.
在最混亂的時間.
她就說.
她說在最混亂的時間.
反而是令她和她的關係越來越好.
她的丈夫在旁邊.
我看著他.
就有點心動.
我記得見過一次.
很好笑.
竟然太太稱讚她的丈夫.
她的丈夫說.
我結婚十幾年.
我第一次聽到你稱讚我.

$^{401}$我聽完.
我都有點無本度.
我發覺很有趣.
原來在艱難的日子.
當我們不只是看著艱難.
我們去看見一些恩典.
或者神做的事.
是可以不同的.
神的態度不同.
我一路回家的過程.
雖然都很累.
或者覺得還未解決的事情.
不知道可以怎樣.
都會有點擔心.
但是我記得.
神怎樣一步一步地去幫助這個家庭.
我就知道神還在.
祂的神跡是沒有停過.
只是我還未看到.
今時今日.
世界發生了很多事.
我們打開電視.
無論你看ABC,NBC,CBS,CNN還是Fox.
還是Levinson.
很多資訊在當中.
但是你看到什麼?.
你看到的是荒謬,混亂,失控,失望.
但是你在當中.
又看不看到神在做事?.
你看不看到當中有神跡?.
很多人遇到神跡.
都不覺得是神跡.
根本就不會尊重我們.
但是當人好像吉利利和這個大神一樣.
肯去看神跡.
就會發生在我們每一個人的眼前.
好了.
看見是第一步.
第二步就是我們要相信.
這兩步經常都連著一起發生.

$^{441}$大神有他自己的要求.
有他自己的想法.
認為這樣才可以有神跡.
但是耶穌很主張的說話.
就是在挑戰他.
神跡是要神作主.
人要做的只是一件事.
相信.
相信他就可以了.
我再一次補充.
因為我們通常都不記得.
信心不是說大小,多少.
而是說對象.
你信任誰?.
簡單來說.
你有沒有對準神去信他?.
這幾年我在香港經常都要做一件事.
就是吃一個farewell飯.
變成一個日常的生活.
每一次farewell的時候.
我們都會很自然地想起當年.
每個人都會說當年這件事情.
如果可以這樣發生就好了.
這件事情.
我想其實應該要這樣發生.
這樣就行了.
之後說完一大堆.
我們就越來越失望.
因為沒有一件事.
跟我們這樣想,這樣的安排.
就能夠發生.
我想起很多年輕人.
跟我會說一句話.
他說沒有希望就沒有失望.
但我很想提醒很多年輕人.
如果你沒有望.
你會不知道做什麼.
好像盲人一樣.
你望不到東西就不敢做.
不懂做任何事.

$^{481}$所以始終要做的事情.
但是望什麼?.
我們記得一首歌.
「當轉眼仰望耶穌」.
你要望準,望實耶穌.
你就能夠相信.
當你相信.
神蹟就會在你的眼前發生.
到最後.
我們能夠看見,能夠相信.
下一步就是我們要願意去信服.
這一步也不是容易的事.
當我們很心急,很迫切.
最有安全感的事情.
就是我們要捉緊自己的信念和想法.
這樣就沒事了.
我很想別人跟我的做法.
但是看著別人.
不是跟你的做法.
萬條司你.
愛你不理.
那一刻會怎樣?.
「頂」我就火都來了.
你看著.
好像耶穌一樣.
你就真的大件事了.
那個大神.
其實他可以安慰那個大神.
「沒事的 我會救你的兒子」.
其實你說這個就可以了.
但是偏偏耶穌.
真的不是你想的東西.
一問一答.
還要命令你.
為了什麼?.
就是要這個大神.
願意放下他的心急.
讓他可以信服.
相信和信服.
Trust and obey.

$^{521}$永遠都是一對.
所有歌曲都可以唱.
我想起一件事.
我想起我是一個沒有車牌的人.
但是認路挺好的人.
在香港是死不了的 沒有車牌.
我認路是怎樣的.
我都覺得我自己在腦海當中有個GPS.
我是會走的.
很有方向感的.
我有一次去趕時間.
當然是跳上的士.
去一個遠的地方.
當我看著司機去繞路的時候.
我明明知道他在繞路.
我跟他一樣的.
我火都要來.
然後把身體靠向前.
就想出聲罵他之際.
我突然間看到一件事.
我看到那個Google Map.
Google Map上面有一條紅線.
我看一看地圖.
我就知道.
那條紅線.
就是原本我想走的那條路.
紅線代表什麼呢?.
是實情.
我那一刻.
哭了一口.
很安心.
順服那個司機繼續載我.
當我們相信.
他選的那條路.
是有他的心意.
他有他獨特的眼光.
或者那條路在那一刻.
是應該最行得通的一條路.
如果我們知道他這些的判斷.
我們可以怎麼做?.

$^{561}$就是我們願意去相信.
相信神會繼續在這裡帶領我們走.
我願意大家藉著今天的道路.
我們可以繼續去學習.
看得見神的神蹟.
能夠真心去相信.
並且順服.
那一刻我們會和這個大神一樣.
可以經歷耶穌這一份神蹟.
充实充足的一天就完成了!.
\newpage



\section{}
\label{sec:KfYBQZvmZ4Y}
\textbf{Apostles' Creed \#9 (Life Eternal) - Rev. Mark Chan}
\newline
\newline
連結: \href{https://youtube.com/watch?v=KfYBQZvmZ4Y}{\texttt{https://youtube.com/watch?v=KfYBQZvmZ4Y}} ~~~~ 語音日期: 2022-08-16
\newline
\newline
\hyperref[sec:pMRJMgqJITI]{\small{< < < PREV SERMON < < <}}
~
\hyperref[sec:index_chronic]{\small{[返順時目]}}
~
\hyperref[sec:index_scriptual]{\small{[返順卷目]}}
~
\hyperref[sec:AlQ26GIFgJs]{\small{> > > NEXT SERMON > > >}}
\newline
\newline
$^{1}$to finish the long series on Apostles Creed..
All right, Apostles Creed..
This is a long series..
I think this is part number nine or 10..
We began with the series and saying that.
the Creed of the Apostles are like.
roadmap of Christian faith..
I remember when I was still in college.
and there were a bunch of guys and girls,.
brother and sisters, we grew up together.
in a local church and we were still.
in the college fellowship in my home church..
Had quite a close friends and we grew up together..
And back then we were all solid Christians,.
standing together, serving together,.
at least judging from what we said.
and how we act back then in those years..
But as years go by, some have decided to quit..
Some have decided to drop out of Christian faith altogether..
But as I continue to discover,.
some of these friends quit.
because they reject the Christian faith..
And as I continue to discover,.
they have some sort of truncated version.
or truncated understanding of what.
Christian faith is all about..
For instance, I don't know what got into their thinking..
One friend believes that Jesus is just a man,.
just a typical human being,.
and teaching mankind how to be good for some reason..
But I find it interesting..
And then the sort of Jesus that he believes in.
is simply a moral teacher..
So it is understandable why there's no point.
to continue to follow Jesus, right?.
Because if Jesus is just a moral teacher,.
and what the heck does he do?.
What the heck does Christianity do to you, right?.
There are many moral teachers in the world or in history,.
right, and Jesus is one of them..

$^{41}$And again, as I continue to discover,.
this is sort of a truncated, reduced understanding.
of what Christian faith is all about..
Another friend of mine drops out of the Christian walk.
because life gets tough..
Don't get me wrong, life does get tough..
Different circumstances or difficulties.
happen along the way, and then sometimes we lost focus..
And I think this friend of mine definitely presupposes.
that following Jesus means nothing.
but everything will go well, right?.
Now, if that version of Christianity shattered.
by different difficult circumstances in our lives,.
then definitely you will drop out, right?.
Now, again, I wanna focus on saying.
that this sort of understanding is problematic.
because following Jesus is being thought of.
as the skies are always blue and joy without sorrow.
and all that..
And I tell you this, these kind of thinking.
are nothing but heresy, nothing but heresy, right?.
We are being blown by different cultural winds of our day..
And the culture script has always been saying to us,.
hey, your life will go well because God is good..
If your life doesn't go well,.
then there's some problem with God, right?.
There's some problem with your faith or anything like that..
So you will automatically drop out, right?.
In our days and age, we love to pick and choose..
We do, we love to pick and choose, right?.
I do that too, right?.
We go on Amazon, AZ, whatever that is, right?.
Recently, I was buying some item off Amazon.
and it took me quite a few minutes to pick and choose..
I look at the customer view.
and I want to find the best deal I can get, right?.
It's for my benefit, right?.
That sort of mentality, if we bring to faith,.
it is problematic..
And we think that sometimes we are tempted to think.

$^{81}$that the Bible, the Christian faith is like Whole Foods,.
right?.
Ever been to a Whole Foods?.
You just go in and, you know, to a section..
You can pick the salad you want..
You pick up some items or drop some items, right?.
Not interested, you just decide not to buy it..
And I tell you, this is the biggest lie.
Satan is telling us..
In the Bible, particularly in the Psalms,.
in the Old Testament, you see the poet, the believers,.
they struggle with their life circumstances..
They struggle with suffering..
They suffer from pains and sickness,.
but they never drop out..
The only thing we know they are not dropping out.
because they, it's the fact that they are praying..
The fact that they are struggling with God..
They struggle mightily with God..
When life is good, they praise God..
When life is good, they give credit to God..
But when life goes through a valley,.
they lament, they struggle, they complain.
until their faith is refined.
and their relationship with God is so polished..
Have you ever seen a couple who are married for many years?.
Have you ever seen one of those couples, married couples?.
Married for so many years, okay?.
If you know them well, you know that their relationship.
is no more the kind of, you know,.
rosy, romantic stuff, right?.
They are, they passed that stage..
I mean, they still have romance..
They still have romance, but it's a different kind, right?.
The beginning phrase of the relationship.
of a married couple, probably just puppy love.
or something like that, right?.
It's very warm..
I'm not going to talk down on those experiences..
Those are great experiences..

$^{121}$However, when the couple are committed to walk together,.
you find that their lives and their relationship.
is mature, right?.
And there you go, in years of experiencing.
the relationship together, walking together,.
you find husband and wife,.
they do have, you know, they do have mistakes..
They sin against one another..
They oppose one another..
They fought, right?.
They argue..
However, you see also the maturity of love,.
which means forgiveness, right?.
They learn to overlook each other's shortcomings..
And the relationship is refined and continue to grow..
And that sort of relationship is a great picture.
of our relationship with God, right?.
Our spiritual journey with God is exactly like that..
We grow in the midst of difficulties..
We grow in good times and bad times..
And we grow also in our ability to discern,.
to evaluate things on earth,.
biblically and spiritually..
And we grow in our ability to take on things that matter..
As my relationship with God continues to grow,.
I know that there is something that really matter,.
eternal value, something of eternity,.
something that does not change and will continue to grow..
And that is the most important thing.
that matters in my life..
You know, I grew up, as I said,.
I grew up in my good years, you know, in college..
I lost some friends, you know,.
but what the hell, I don't care, right?.
I know, I get to know other new friends..
I get to know other people, right?.
It's like, it's like traveling on a train, right?.
There are some people who enter into the car.
and then some, you know, along the way, they exit the car..
A few of them stay together until the destination..

$^{161}$You know, people come and people go,.
but what is most important in your relationship with God.
is that you need to discern and you need to stick to.
what is most important..
What really matters in your life, right?.
Romans 8..
It gives us the biggest possible roadmap.
of our Christian journey..
This is the biggest, like 50,000 feet above bird eye view.
of what Christian journey is all about,.
what God is doing in your life..
Romans 8, 29..
For God knew his people in advance..
He chose them to become like his son,.
so that his son, Jesus Christ,.
would be the firstborn among many brothers and sisters..
And having chosen them, he called them to come to him..
And having called them to come to him,.
he gave them right standing with himself..
And having given them right standing with himself,.
God also gave them his glory..
Now, this is a bird eye view.
of what your relationship with God should look like..
All right?.
Everything aside, all right?.
Oh, whether I have a good time or bad time,.
whether I have good buddies to walk with me,.
it's not even in this picture..
God knew his people in advance,.
chose them to become his son..
Now, we can call this sanctification..
You know, when you started to follow Jesus,.
in definitely in different circumstances,.
some people come to Jesus in their difficult times..
Some of them, you know, mountaintop experience,.
they come to know Christ himself..
The process is still the same..
The process of sanctification,.
sanctifying your life, right?.
God chose you..

$^{201}$To be like his son..
It's a long, lifelong process.
that you will continue to become more and more like who?.
Well, not your old self..
Have you ever checked the mirror recently?.
Look at yourself..
Have you changed?.
You know, a couple of years ago, many years ago,.
definitely you change, right?.
You know, I got more, you know, gray hair.
and, you know, wrinkles in here and there, you know?.
No matter what I do to myself..
But, you know, I continue to grow older..
But in the process that God has launched in your life,.
when you decide to profess that Jesus is my Lord,.
Jesus is the Christ, Jesus is my savior,.
and he died for my sin,.
and he rose from the dead for my hope,.
I kick in, God kick in..
The lifelong process of sanctification.
is a process of sanctifying you,.
making you like who?.
Like his son, right?.
Hopefully at the end of the process,.
you will become more like his one and only son,.
Jesus Christ..
So that his son will be the firstborn,.
the one who get the glory among the sisters and brothers..
And God chose them and he called them to come to him..
Having called them,.
he also gave them right standing with himself..
Now that is what we call justification..
You are justified..
You became like Jesus Christ..
God sees you just like he sees his one and only son,.
because you believe in the salvation, right?.
You believe in Jesus Christ..
And that this process is not a long-term process..
It's a punctual, one-time act..
God acts in your life,.

$^{241}$and you became his child..
You became his son or daughter..
That is justification..
That you stand before God and declare sinless..
You declare totally forgiven..
And that is eternal life..
That is a one-time punctual event..
When you are in Christ,.
God sees you as he sees his perfect son..
Now, can you imagine the timeline, okay?.
Timeline of the Christian faith..
All the way to the left, all the way to the left..
All right, your left..
It's justification, all right?.
You became a Christian..
You became God's daughter or God's son..
That is one-time justification..
You stand justified,.
all because Jesus Christ has done for you.
on the cross and in the empty tomb..
And then the timeline goes,.
the lifelong process of sanctification,.
that you will become more and more like his own son..
Just like this..
You start with justification, boom..
And then there was up and down, up and down, you know,.
but in general, you're going up..
In general, you're going up to heaven or the new world..
That is the lifelong process that we are in right now,.
process of justification..
And then at the end of it is glorification..
Notice what he said, Paul said at the last phrase,.
"Having given them right standing justification,.
"God gave them his glory.".
One day when we see Jesus face to face, all right,.
either Jesus come back in his second coming.
or you go to Jesus, all right,.
because of whatever, you know, you die, okay?.
You die..
And then your spirit will go, go to the presence of God..

$^{281}$And one day, you know, the last two chapters of the Bible,.
the new world, the new heaven and new earth,.
when heaven and earth joined together.
in the state of glorification, the perfect state,.
there's no pain, there's no sickness,.
there's no homework, all right?.
There's no frustration anymore..
Then you kick in the new world,.
the new heaven and new earth..
Now, there are different,.
different funny understandings of what heaven is like,.
you know, what new heaven and new earth is like..
Some people would think,.
"Oh, probably we just sing and make music 24/7," right?.
That's kind of boring, right?.
Especially if you are not musically talented, right?.
Unless, I guess, you know,.
in the last process of glorification,.
you resurrected and you got the glorified body.
just like Jesus had on the third day,.
then you probably become the best bass player.
or whatever, right?.
But still, I mean, making music 24/7 is kind of boring, right?.
But you know, you need to understand,.
you need to understand new heaven and new earth.
is Genesis 1 and 2,.
we fulfill, you know?.
Fulfill right before your eyes..
Remember what Genesis 1 and 2 is all about?.
When God created humanity,.
His image bearers, all of us,.
God said to them,.
"Go and take control and rule the world for me..
"You represent my image, my glorious image on earth,.
"and you're gonna bring my order to the rest of the cosmos.".
That is Genesis 1 and 2, right?.
Of course, it didn't go well, right?.
Genesis 3, man proclaim himself to be God, right?.
And kicks in the entire grand story of the Bible..
But God never forsaken His world,.

$^{321}$His creation, our creation, you know, our world..
He set himself a program to redeem the whole world,.
to bring humanity back to Him.
through His son, Jesus Christ..
And now we are part of it, right?.
So the new heaven, new earth, what will it look like?.
It's not just playing music, saying,.
"Hallelujah 24/7.".
"No, no, I don't wanna go there.".
It's gonna be a glorified state.
that will bring God's holy, sacred, beautiful order.
to the rest of the cosmos, all right?.
I don't know if there's a space exploration.
in heaven and new earth, but whatever..
There's no death, there's no pain, there's no sickness..
All you do in this life, think about it..
You know, you all go to work, right?.
Think of moments in your work.
that you find so much fulfillment..
You find so much fulfillment..
What is that?.
Ask yourself, what is that?.
That after, you know, your labor, your sweats,.
you work on something and you finish something..
You finish a product, you finish maybe homework, okay?.
Maybe your assignment or something..
And then you get a sense of satisfaction or fulfillment..
That is who you are, right?.
You're so gifted in this area, no matter what that is..
That sense of fulfillment multiplies it infinitely..
That is what new heaven and new earth will look like..
That you will be exploring and extending.
the order of God to other areas..
Take control of other areas..
You know, I always remind some of our accountants.
or finance people in my church..
I always ask them, you know, have you ever realized,.
you know, God's original creation.
is the creation that has potential to multiply..
Just as when you read through it, you know,.

$^{361}$seed within seeds, God creates trees and the plants.
and the animals and they can multiply..
Things are meant to be multiplied by human hands..
You know, I think last hundred years,.
the biggest discovery for ancient history.
is that when you pick some, when you take some wild grain,.
take some wild grain, and then you work it.
and you cultivate it and you take the seed out of it,.
human hands, plant it, and then it grows..
The biggest discovery is that through a human process,.
wild grain becomes better, right?.
You know, the rice, the white rice that we eat, okay?.
I don't eat that much now, but why is it so soft?.
Why is it so rich?.
Because it's process and process of human effort.
being added into the creation, and the creation multiplies..
The value becomes higher..
Cultivated grain is more expensive than the wild grain, right?.
Because there is human process involved in it..
And that is the process of multiplication..
That is what God originally meant us to do..
The original creation, Genesis 1 and 2,.
and the last two chapters, the new heaven and new earth,.
it's not a static concept..
It's the holy, sacred, dynamic..
We, as we deem people in Christ,.
we will participate in the new heaven and new earth.
to explore and to invent, you know,.
to get to know this amazing creation.
and to get to know God himself, and God is infinite, right?.
So I hope I will continue to teach in heaven and new earth..
And you will continue to do Bible study.
in heaven and new earth..
And whatever you do, and God made you to who you are, right?.
Whether you do a labor work or those,.
some guys doing PhD researching or something, right?.
Making molecules, those, I don't understand those stuff..
But we will continue to engage in the creative process,.
you know, bringing God's order to the rest of the cosmos..
That is the bird eye view of what it is,.

$^{401}$what it is all about to follow Jesus..
Justification and lifelong process of sanctification,.
and then finally, the glorification, all right?.
So don't get, don't retire yet..
All right, new heaven, new earth is not for retirees..
We continue to work without frustration, right?.
And again, the apostles,.
we end with Romans 8,.
is telling us that once we kick in.
to that relationship with Jesus, it's a dance with God..
It's a dance with God..
We dance with the music of the Holy Spirit..
And sometimes the music sounds bright and happy..
Sometimes the melody sound sad, lamentation..
That is the long process of ups and downs.
in our sanctification, but it's an eternal dance..
It's an eternal dancing with God..
A long-term process, a long-term journey.
of discovering God and being transformed.
into His Son's image..
To say to a 15-year-old that you are okay.
with the way you are is the utmost disrespect to humanity..
Let me say that again..
To say to a teenager or younger or older, or even you,.
you are okay, perfectly okay with the way you are..
It's disrespect of God's creation..
You are not okay with the way you are right now..
Agree?.
You cannot avoid the need and the joy of transformation..
God is transforming you and shaping you..
That is the mark of maturity..
The mark of maturity is to say to yourself in the mirror,.
I'm not okay with the way I am today.
because the Holy Spirit is shaping me constantly..
And I need to let God's grace to work in my life..
Let the Word of God continue to shape me.
and challenge me to shape me,.
to transform me into the image of His Son..
And the Apostles' Creed is challenging us to aim for that,.
to aim high in the kingdom of God..

$^{441}$You are not okay with the way you are today, no way..
You will be transformed and let God transform you..
I believe in the Holy Spirit..
That's the first section of the Creed..
I believe in the Holy Spirit..
Have you realized this is the new creation?.
When the Spirit comes upon you,.
when the Holy Spirit is making you God's children,.
you are a new creation..
The Holy Universal Church, the Holy Catholic Church,.
the Chameleon of Saints speak to us.
that once the Holy Spirit is taking you.
into the kingdom of God,.
the church becomes your new family..
And that new family never change..
The earthly family will change..
The configuration of it, you know,.
it will not last forever, right?.
Family relationship will not last forever..
You understand?.
Someday, my son is not my son, right?.
And I'm not his father, someday, okay?.
This family, earthly relationship will not last forever..
Right?.
To make it last forever as if it really lasts forever.
is actually idolatry..
To make it idol in our lives..
It will change if we transform so much.
that it will be glorified, right?.
And the fulfillment that we have in the family relationship.
will exceed how we feel..
Will exceed the sense of fulfillment we have right now..
In the new day, in the new heaven and new earth,.
it will be so glorified..
So the church, the Universal Church,.
the capital C Church is the new family..
And then the Chameleon of Saints,.
I think we talk about it, is the new relationship..
When we are in Christ, okay?.
When we are in Christ,.

$^{481}$whether you're a Christian struggling with your faith.
so much or some other, you know,.
having a mountaintop experience,.
praise God every single day..
If we are in Christ, we have this new relationship,.
the Chameleon of Saints, like togetherness, right?.
Somehow being united together..
As we worship today, I don't know if you have noticed.
one of the songs, One Voice, right?.
Now, I didn't plan that,.
and the worship leader picked that song..
The song talks about some sort of unity..
That kind of unity goes beyond our ethnic unity,.
our political identity..
It goes beyond that..
In Christ, we find that new communion, new unity..
The forgiveness of sin, the new relationship with God..
Chameleon of Saints is the horizontal relationship,.
human being to human being, question to question,.
new chameleon of saints..
And the forgiveness of sins is vertical..
So we look upon God and say, "We are justified,.
"not because of our own effort, our own work,.
"but what Christ has done in your life.
"that made you someone who is new before God,.
"as if Jesus Christ, the one and only perfect Son.".
And then finally, I believe in the resurrection of the body,.
the new existence..
God kicks in the new existence in our lives.
when we profess the name of Jesus..
You know, you realize that?.
Yeah, I know my body continues to decay,.
your body too, right?.
We get old and we get sick, we get COVID..
We are weakened, right?.
You know, it used to be I could do 70 pushups,.
just like that..
Or you could do that..
I can't do it anymore..
Right, now I even have struggle with 20, right?.

$^{521}$We decay..
However, we have the hope of getting the new body,.
the new existence, just like the third day of Jesus Christ..
He walked out of the empty tomb..
He conquered death for us..
He conquered sin and darkness and Satan for us..
No more forces of darkness,.
no more forces of deterioration can go upon you..
And this is the new existence that we have..
And we have that hope in future..
Whether we see Jesus face to face.
or in the new heaven and new earth,.
we'll get a new body..
A lot of people think that, okay, once we die,.
our spirit will go like floating around, you know,.
just being creepy and all that..
No, that's not the biblical picture..
The Bible tells us that we are going somewhere,.
eschatology, going some end,.
and the end is the glorified body..
All God's people will share in Jesus Christ..
And lastly, the Creed says,.
I believe life everlasting..
That is the new world..
Eternal life is not something.
that we experience after we die..
Eternal life is the experience that we have once.
you come to know Christ..
When you come to know Jesus, that is eternal life..
And God is with you..
And you live in the presence of God..
That's why the Bible continues again and again.
telling you, you need to be obedient..
You need to listen to the word of God..
You need to walk with the spirit.
so that you can live in the presence of God, right?.
And that is conditional..
I know God's love is unconditional in some sense,.
but it's also conditional..
If you want to experience the power and the joy,.

$^{561}$the peace, the hope that we have in Jesus Christ,.
well, there is condition..
We need to be committed to worship Him..
John 10, 10..
Let me close with this verse..
Jesus is talking to the Jewish people in the New Testament,.
to the people of God..
The thief comes only to steal and kill and to destroy..
I, Jesus said, came that you may have life.
and have it abundantly, abundantly, you know?.
Like overflowing in your life, right?.
All the good things of the kingdom,.
the joy, the peace, the patience, the hope,.
and all those good stuff will be so overflowing.
in the cup of your life..
You cannot bear to contain them..
It has to overflow out of your life.
and you need to share with other people, right?.
You need to, you know, have other people experience it..
It's shareable, right?.
And it has to be shared..
There is life abundant..
In Old Testament term, this is New Testament,.
Gospel of John, in Old Testament term, it's shalom..
Shalom..
Even the Israeli Jews today will meet each other,.
see each other, and they will say the first word,.
not hi, not hello, but shalom..
Shalom..
What is shalom?.
Shalom is not merely the absence of war..
It's not just the absence of frustration,.
absence of death, all the bad things..
It's not just the absence of the negative things..
Shalom is positively means all the great things.
that we have in the kingdom,.
all the great attributes of God that we can share,.
we can experience, that is shalom, right?.
Quite meaningful, you know, to greet each other,.
to say shalom..

$^{601}$I wish you well..
Your body, your spirit, your mind, your emotion,.
everything in your life will be so blessed.
that you can have shalom, right?.
That is a foretaste, a foretaste of the new world,.
the new heaven and new earth..
Well, I pray that God will bless you in such a way.
that you will continue to experience the foretaste.
of the banquet, foretaste of the great things to come,.
the greatest thing to come..
Nothing the world can offer you that,.
only in Christ we have this hope..
Let us all pray together..
Father, we give thanks because you are our great God..
And thank you for the hope that we have in the kingdom,.
the eternal hope that we have,.
that that hope challenges and conquers death,.
and conquers the dark power of Satan and sin..
And I pray that this hope will be,.
will sink in into the very core of our being,.
that we'll continue to look forward.
to this new heaven and new earth, this life eternal..
And thank you that we are justified in Christ..
And thank you for your spirit being sent to be with us,.
to dwell in us, to sanctify us,.
to help us to walk and follow you..
And I pray that one day we look forward.
to the day of glorification,.
that we can experience kingdom in the fullness..
And right now, no matter where we are,.
no matter what circumstances we are facing,.
we get the foretaste of the glorious heaven..
That someday we'll see you face to face,.
and your people will continue to worship you, to serve you..
And thank you for the church that we have,.
the community of faith that we share,.
the communion of saints..
And I pray that this sense of unity.
will continue to grow your church, to grow your family,.
that your church will be your people,.

$^{641}$empowered to glorify your name..
In Christ's name, we pray, amen..
Well, lastly, two minutes, I want to..
\newpage



\section{}
\label{sec:AlQ26GIFgJs}
\textbf{Psalms for Every Season (Psalm 63) - Rev. Mark Chan}
\newline
\newline
連結: \href{https://youtube.com/watch?v=AlQ26GIFgJs}{\texttt{https://youtube.com/watch?v=AlQ26GIFgJs}} ~~~~ 語音日期: 2022-08-23
\newline
\newline
\hyperref[sec:KfYBQZvmZ4Y]{\small{< < < PREV SERMON < < <}}
~
\hyperref[sec:index_chronic]{\small{[返順時目]}}
~
\hyperref[sec:index_scriptual]{\small{[返順卷目]}}
~
\hyperref[sec:DvMyIFbhzjs]{\small{> > > NEXT SERMON > > >}}
\newline
\newline
$^{1}$- Some weeks ago, oh, I forgot to turn it on..
I'm sorry, yeah..
You know, this movie is about emotions inside out,.
and quite a few weeks ago, I remember Speaker Brian Ame.
was with us, and he talks about the emotional will, right?.
Remember that?.
Yeah, how we should articulate our emotion correctly,.
and how we process our emotions, okay?.
So here's my question to you this morning..
Is where are you gonna put yourself in the past week.
on the emotional will?.
Where are you gonna put yourself?.
Okay, so pick a color, pick a color, or pick one emotion..
Well, you have some stronger emotion,.
so-called positive emotions, happy,.
you know, feeling strong and pump up, right?.
And then on the other side is feeling sad and depressed,.
bored, exhausted, lonely..
And then down there, there's fear, the emotion of fear..
And then on the upper right is feeling mad, you know,.
feeling mad about something, okay,.
or feeling mad about someone, all right?.
So take a look at the emotional will..
Where are you gonna put yourself?.
Where are you gonna put yourself in the past week?.
All right?.
So how many of you put yourself in the fear section?.
Fear section..
Wow, you guys are so bold, okay?.
No one's in fear, all right?.
How many of you, just take a sample,.
how many of you put yourself in the mad emotion?.
The yellow..
Paul, you were mad, right?.
Yeah, mad about someone, probably..
(laughs).
Or mad about yourself, right?.
How many of you feel the sad part?.
Depressed a little bit?.
Okay, all right..

$^{41}$How about on the left side, the strong side,.
the happy side?.
How many of you fall into that quarter?.
No?.
Okay, those of you who didn't vote,.
those of you who didn't vote,.
you're probably being emotional less,.
which is not good, right?.
Which is not good..
All right, so I'm sure the past week.
probably very hectic for you..
Whatever happened in your life journey,.
we need to understand,.
as we look at one of the Psalms today,.
we have been looking at the book of Psalms,.
using different Psalms, different songs.
in the book of Psalms to study..
You need to understand emotion is a window..
It's a window to your soul..
It's a window to something inside of you..
Now, this is very interesting.
because you look at the book of Psalms,.
there were 150 Psalms recorded,.
being collected together..
And those Psalms are actually prayers of people to God,.
words of humans to God..
And at the same time, we collected those Psalms,.
and those Psalms are collectively God's word to us, right?.
So 100\%, they are prayers to God, right?.
From a human poet, from a human player,.
playing a song, playing a prayer to God..
At the same time, it is God speaking to us..
If you run out of materials for your quiet time.
or devotions, you always can go back.
to the book of Psalms, right?.
You never get tired of thinking about,.
pondering on the words of the Psalms..
So, you look at the 150 Psalms,.
you find all this kind of colors,.
and all this kind of emotions..

$^{81}$And that's why I want to ask you how you feel today..
How are you feeling in the past week?.
Let's process your emotions, right?.
Life is hectic, right?.
You have to take care of many, many different things,.
and you have to tackle different challenges and issues..
But it's time for you to come before the presence of God.
to process your emotions..
How do you feel today, right?.
How do you feel today?.
And let us go to Psalm 63..
Now, as we read through some of the verses of this Psalm,.
I want you to understand and hear.
what David is trying to say, his prayers,.
his words to God, and how did he feel..
A Psalm of David..
And he was in the wilderness of Judah..
"O God, you are my God..
"Earnestly I seek you..
"My soul thirsts for you..
"My body longs for you..
"In the dry and withered land, where there is no water,.
"I have seen you in the sanctuary,.
"and beheld your power and your glory..
"Because your love is better than life,.
"my lips will glorify you.".
The first few verses that we have just read.
are the prayers of David..
And he made this prayer, he made this Psalm,.
he composed this Psalm, where?.
In the wilderness, right?.
In the wilderness of Judah..
Now, if you go back to the background of the Psalm,.
you will understand wilderness is actually part.
of the landscape of Israel..
It's part of the landscape, a default geographical setting.
of any Israelite living in the land of Canaan..
They know where the wilderness is..
You climb the mountain and you descend the other side,.
you will see wilderness..

$^{121}$The land of Israel is very interesting..
It's very engaging..
Next week, I'm leading about 10 people.
going to the land of Israel..
It's gonna be interesting..
I want to tell you something, it's very interesting..
When we drive from west to east,.
we drive west to east,.
you see a total different two worlds..
All right, let me show you how..
The land of Israel, going from left to right,.
from west to east, is actually a mountain.
climbing all the way up to almost about 2,000 feet.
above sea level, 2,500 feet..
That's where Jerusalem is..
Can you see that?.
Can you make it out?.
And then after Jerusalem, you continue to drive,.
and then you will continue to descend.
in terms of elevation..
You continue to descend and you can hear it,.
you can sense it because your ears pop, right?.
Because you are descending so much.
to the lowest point on earth,.
1,400 feet below sea level..
It's very interesting to drive from left to right.
and from west to east,.
because when you pass Jerusalem,.
there's a tunnel in Mount of Olives..
You go through that tunnel,.
before you went through the tunnel,.
everything looks kind of green..
After you went through the tunnel,.
you went to the right side,.
the right side of the blue line,.
where the water shadows..
After you pass that, everything looks brown and yellow,.
because there's no grass growing.
on the other side of the,.
the right side of the watershed..

$^{161}$Now, what is important is the watershed, the blue line,.
because the rain cloud is moving from left to right,.
from west to east..
And when the rain clouds move through the land,.
especially in the rainy season, the winter season,.
about October to March,.
it will rain on the left side of the watershed..
And then after the rain clouds go through,.
went through the watershed line,.
there's no more rain..
So after that watershed,.
there's no more rain on the right side of the watershed..
So let me show you two pictures..
You'll see the difference..
The left picture is the left side of the watershed..
You see kind of green, right?.
Most of the time, there are greens..
And then after you move through the,.
past the line of watershed, past Jerusalem,.
you will see everything is dry, right?.
Everything is so dry..
It's like desert..
That is wilderness of Judah..
So for Israelite, for the people of God,.
in the Old Testament, they know what wilderness is..
They climb the mountain and they will see wilderness..
They will go into the wilderness..
And David, as the title of the psalm says,.
he wrote the psalm in the wilderness of Judah..
Now, because wilderness is kind of default setting.
for David and for the people of God,.
wilderness becomes a metaphor, a metaphor of life..
It represents testing and tough times.
you have to go through in your life..
Challenges and trials and temptations..
And from the landscape, you see the soul scape..
What is going on in your life, right?.
Perhaps today, you are going through some wilderness..
Perhaps recently, life has been tough for you..
Life has been challenging to you..

$^{201}$And because of the background of the psalm,.
we understand this wilderness..
It's unavoidable..
There's no way you can avoid it, right?.
You lost a job, that's a time of wilderness..
You lost someone you love, that is time of wilderness..
You lost some of the good things in life,.
that is a time of testing, time of wilderness..
Now, the most important question is,.
if you are going through wilderness, what do you do?.
How do you maintain and sustain yourself.
through wilderness, through the dry.
and weary season of your life?.
And here, the psalmist David said,.
"Oh God, you are my God..
"Earnestly I seek you, my soul thirsts for you,.
"my body longs for you..
"In a dry and weary land of wilderness, there is no water.".
What is he doing?.
You know, something is not right in his life..
And in fact, the context tells us,.
if you continue to read Psalm 63, he was on the run..
David was a fugitive..
You know some stories of David.
in the book of First and Second Samuel..
He got a lot of enemies trying to kill him,.
pursue him, and take his kingship and take his power..
So he was on the run..
And where are you gonna run when there is trouble?.
Where are you gonna run when you are pursued by enemies?.
You run into wilderness..
It is a land of fugitives,.
land of escape,.
land of going through some tough time..
Wilderness becomes a metaphor for all of us..
And he was on the run, and he was praying the psalm.
in the wilderness where there's no water..
It was time of testing for him..
And of course, if you're in the wilderness,.
the first thing you want is water..

$^{241}$You need water..
When you are thirsty, you need a glass of water..
I go to Israel quite often,.
and the most important thing you need to bring.
when we are on a day trip,.
we leave our hotel or Airbnb very early in the morning,.
and then there's a day trip for us,.
and we won't come back until probably evening,.
after sundown..
The most important thing in our backpack.
is the bottle of water..
You probably have to drink at least one to two liter.
of water every single day..
Dry and wiggly land..
The first thing you want, your body thirst for, is water..
But David is doing something very interesting here..
He transformed his physical experience.
into a spiritual experience..
He take that as a spiritual transformation moment..
It's for him to yearn for God..
Because if you know,.
if you have gone through life for some period of time,.
whether you're going through some really tough time,.
or your lack of money, you worry,.
or you deal with a lot of negative emotions,.
you need to understand that.
in terms of the physical body,.
the water can do something to your body,.
can crunch your thirst..
But there's emptiness in your life..
Water cannot crunch..
There's only one thing that can crunch the thirst,.
crunch your yearning..
And Psalm 63 says, "It is God.".
God himself..
Only God can fulfill and make you satisfied.
from the inside out, from your heart..
So seek God and thirst for God..
That is what David does in the wilderness..
Notice what he says, "You are my God.".

$^{281}$My God..
Now usually the English pronoun,.
we don't think it is important..
But in this Psalm, verse one,.
"God, you are my God.".
God of mine..
That personal pronoun is important..
Because for David, God is his God..
His personal God..
He has a vibrant, personal, living relationship.
with God himself..
So God for David is not someone else's God..
It's not some other pastor's God..
It's not some other family member's God..
It's not some other friend's God..
Why so often we associate with God.
and our relationship with God.
is by other associations, right?.
I have a good friend who is very good Christian,.
who is very solid Christian,.
who is very, very solid believer..
And because of him, I kind of close to God, right?.
But it's got your God..
It's got your personal, relatable being..
Self-sufficient in the wilderness..
David said, "You are my God.".
God and me is a personal relationship..
I'm not sure where you are in your relationship with God..
Can you really call God your God?.
Can you really call Jesus your savior,.
your Lord and your leader?.
Yours, not someone else's..
Personal relationship in the wilderness..
He remembers his relationship with God..
And that's the thing, that helps him..
I want to challenge you..
You probably heard of Jesus..
You probably accept the gospel,.
you know, here and there for some time..
But can you really pray 63 verse one?.

$^{321}$"Oh God, you are my God..
"And Jesus, you are mine.".
Right?.
And secondly, what does he do?.
He recalls his past experience with God..
He continues his prayer..
He says, "I have seen you in the sanctuary,.
"in the temple, in the temple of Jerusalem.".
David used to be a worshiper, regularly, weekly,.
on Shabbat, on Sabbath, Sunday..
He goes to the temple, he goes to the court,.
worshiping God with the community, right?.
He remembered that in the wilderness..
I really appreciate the comment made previously.
about the need for community..
Need for a friend, right?.
You know, in my circle, recently there was a big news.
about some people being lured.
into a human trafficking trap..
You know, in my circle, there's some people from Hong Kong..
They are being tempted by some human trafficker.
to go to Cambodia or Thailand.
to make some big money, right?.
And they're being tempted.
because there's so lack of money, right?.
They wanna make some quick money, quick income,.
and that's why they got themselves into the trap.
and they followed the head of the human trafficking scheme.
and went to Cambodia,.
and they never able to return, they disappeared..
It's a horrible news..
There's no bottom line in morality..
But what they did was that they were in a time of weakness,.
time of wilderness, so to speak,.
and they are being tempted to get some quick money.
but got themselves into huge trouble..
What David did in the wilderness.
is that he recall something.
that he had experienced with God..
Memory defines us..

$^{361}$Think of your memory, your spiritual memory..
What do you remember?.
Reading the Bible, God speaking to you,.
the Holy Spirit touching your life..
You know, here and there,.
there must be some sort of spiritual encounter with God..
Some people would encounter God with a different ways..
Some people would encounter God.
in a totally different means,.
and some people would cry for three days..
Some people would just feel it, and life goes on..
But those are real experience..
Those are real spiritual encounter with God,.
and that goes into our memory,.
and this is what exactly what David is doing..
He recall in his memory, in the midst of wilderness,.
in the midst of tough time..
Do you do that?.
Do you do that when life gets difficult?.
When life gets tough?.
You recall memories with God, right?.
I wonder if you have ever lived with someone.
who has suffered Alzheimer's disease..
Have you ever?.
For some years, my mother-in-law.
was suffering Alzheimer's disease..
Now, Alzheimer's disease is painful to see..
It's painful to experience with her,.
because at least physically,.
I mean, her organs or body functions.
are continually deteriorating..
But I think the most painful thing.
to see her going through suffering Alzheimer's disease.
is the loss of memory..
She lost much of her memory,.
and she lost herself..
She doesn't know who she is..
Not to mention where she is, right?.
I remember that one time, she always go out for a walk,.
and there was one afternoon, I came back home,.

$^{401}$and we came back home, and we don't find her..
So we're kind of very desperate..
And then I drove my car around the neighborhood,.
every single block, searching for her..
And all of a sudden, on Western Avenue, on Western,.
I found a cop, and there she is..
The cop is talking to her, but she doesn't know English..
She doesn't know where she is..
And then, thank God, we found her..
And then I talked to the cop, and he says,.
"My mother-in-law, okay, I want to pick her up.".
So he was kind enough, and no question asked,.
and we bring her home..
It's painful to see that she just don't know where she is..
She don't know who she is..
And after she came home, a couple hours later,.
we asked her, "What happened?.
"What happened, Mom?".
She cannot recall..
And I told her, "We were searching for you..
"You got lost in the neighborhood,.
"and some of the neighbors encountered you.
"and called the police..
"That's why you came home.".
And she cannot remember, recall any of this experience..
I think God is merciful to those people,.
because they just don't know who they are..
And they even cannot sustain their belief..
If you ask her questions about faith,.
she cannot respond positively..
Memory loss..
Now, what I'm trying to say is,.
memory is a powerful thing..
Memory defines us..
When you meet someone who's new to you,.
meaning a new friend, what do you ask?.
How do you know someone?.
Hearing the story, right?.
If you want to get to know someone.
or introduce yourself to someone else, what do you do?.

$^{441}$You talk about yourself..
You talk about your story, your history,.
your memory, where it came from,.
and what did I do, and what happened to me.
in the past few years, or whatever..
It's your memory that connects with other people..
And memory is a powerful thing..
And in the wilderness, Psalm 63,.
David is recalling his past experience with God..
I have seen you with my own eyes..
The sanctuary, the temple, the glorious temple,.
the sacrifice, the songs that I hear,.
the choirs that are singing the praise and worship..
Other people of God in the community..
That's why his spiritual life is vibrant..
Now, how do you sustain yourself?.
How do you even grow in the wilderness?.
The answer is to remember..
To memorize what God has done in your lives..
So often we think when life is not going well,.
only if God can turn around my bad situation,.
then everything will go well, right?.
This is our naive prayer..
Now, of course we want to get healing, right,.
when we're sick..
Of course when we lack of money,.
we want God to shower us some cash, right?.
Turn around our circumstances, turn bad into good,.
and then we'll praise God, we'll thank God, right?.
No, of course, we'll pray that way, right?.
We'll pray in that orientation..
But don't you know, at least from the Old Testament,.
or from our real life situation,.
when people find themselves, everything is going well..
It's going well..
They forget about God..
For the Old Testament, story after story,.
the people of God, when God bless them,.
they forget about God, right?.
They forget about Him, and they follow the idols..

$^{481}$They follow money..
They worship gold and silver..
This is the irony in our relationship with God..
However, David is not praying.
that God will take him out of wilderness..
In the wilderness, he recall his past experience..
He recall his experience of worshiping God in the temple,.
worshiping God with other people..
And guys, this is how you keep yourself,.
sustain yourself in the presence of God,.
in the love of God..
It does not merely depend on your circumstances..
It does not depend merely on your external condition..
It is right here, the soul scape..
It's not the landscape that defines you..
It's the soul scape..
It's your relationship with God Himself..
I remember I had a really good friend,.
and his life really inspired me..
What it means to be a Christian..
You know, he's a very humble man, very quiet man,.
but his faith is rock solid..
I remember in his last days,.
I went to ICU to visit him..
He couldn't talk..
He could only say a few words..
But his faith is rock solid..
Now that inspires me because,.
not because he was so blessed, right?.
Not because everything is going well with him..
On the other hand, he's not going well, right?.
He was suffering from cancer..
But in his whole life,.
he becomes inspiration for many of us,.
because even in times of wilderness,.
he lives in the presence of God..
He and God have a relationship..
Have a close, personal relationship..
And that inspired me because of his faith..
He reminds me there is some other reality in the world..

$^{521}$It's not just something we can hear.
or we can touch right now..
It's not just here and now..
There is another world..
There is another hope that transcends everything..
That inspires me..
The goodness of God..
Now we will continue to study the psalm,.
but I want to leave you this challenge today..
Is God really your God?.
Or you just follow your friends?.
Yeah, we have a good time, right?.
But God is just someone else's, not yours..
You are my God..
Earnestly I seek..
And if you find yourself in times of wilderness,.
I pray that your life will continue.
and completely transform..
It's not just a quick fix..
All right, everything goes well..
And then I forget about God..
And life goes on..
God is taking the landscape into your soul scape..
He's working something transformative in your spirit.
so that you will draw closer and closer to God..
Let's all pray together..
No matter where you are.
and what is going through in your life..
I pray that you take a moment before God.
to reorient yourself..
God, thank you..
You are the God who listens to our prayers..
Thank you for all the emotions that we encounter..
All the emotional highs and lows.
that you allow to have happen in our lives..
For our experience to transform us,.
to bring us closer to you..
I pray especially for those sisters and brothers.
who are going through some tough time of wilderness.
where there's no water and they're exhausted..

$^{561}$I pray for your strength..
I pray that your presence will be real to them..
You will walk through the wilderness with them..
You will be with them..
You will let them sense your presence.
as something so real that we can take it.
and call you our God, our personal God..
And thank you for being our great and faithful Lord..
I pray that every step of the way in this coming week,.
you will grant strength and grace to your people..
That your church will be built up.
and your people will be closer to you..
In Christ's name, we pray, amen..
In Christ's name we pray..
Amen..
\newpage



\section{}
\label{sec:DvMyIFbhzjs}
\textbf{曠野的記憶 (詩 63) - 陳崇基牧師}
\newline
\newline
連結: \href{https://youtube.com/watch?v=DvMyIFbhzjs}{\texttt{https://youtube.com/watch?v=DvMyIFbhzjs}} ~~~~ 語音日期: 2022-08-23
\newline
\newline
\hyperref[sec:AlQ26GIFgJs]{\small{< < < PREV SERMON < < <}}
~
\hyperref[sec:index_chronic]{\small{[返順時目]}}
~
\hyperref[sec:index_scriptual]{\small{[返順卷目]}}
~
\hyperref[sec:X4wQCaypggM]{\small{> > > NEXT SERMON > > >}}
\newline
\newline
$^{1}$我每幾個星期都會看不同的詩篇.
讓詩篇來對我們說話.
我記得好幾年前有一齣電影叫做《Inside Out》.
大家有沒有看過?.
有可能是講情緒的.
即是喜怒哀樂.
不同的角色代表不同的情緒的起點.
我亦都記得不久之前有一位英文的講員.
跟我們講這個叫做Emotional Wheel.
即是情緒的圓圈.
情緒的車輪.
我就找不到一個好的中文翻譯.
找不到一個好的中文翻譯.
我只找到英文.
用不同顏色代表不同的情緒.
我不知道你過去一個星期會放自己在哪裡.
甚麼顏色?.
給大家一點時間看看.
你是哪一種顏色?.
有紅色,有藍色,有紫色,有黃色.
我不知道你是哪一種顏色.
或者在圓圈裡面你會選擇哪一個?.
選定好嗎?.
OK?選定.
那你問一下為甚麼?.
為甚麼你會感覺到自己做了甚麼?.
或者感覺到自己做了甚麼?.
為甚麼你現在的心理狀況,情緒狀況是這樣?.
是否因為小朋友激氣?.
或者公司,學校有很多濕濕的東西要處理.
所以令你很疲憊等等.
這個觀景幫助我們.
去看看我們現在的心境是怎樣.
往往很多時候上帝是藉著我們的情緒.
去跟我們說話.
我們的情緒其實就是一個鏡子.
是我們心靈的鏡子.
Mirror of our soul.
或者是一個Window of our soul.
就是一個窗戶.

$^{41}$讓我們能夠看到我們裡面是甚麼.
上帝很多時候是藉著這個窗戶.
去跟我們說話.
如果你回到詩篇的話.
150首歌是嗎?.
你幾乎找到所有人類.
所有你能夠經歷到的.
或者感受到的一些情緒.
喜怒哀樂全部都在這些詩篇裡面出來.
很奇怪.
如果你這樣想的時候.
詩篇其實就是人對神說的話.
是祈禱來的.
同時也是上帝對我們的啟示.
也是上帝對我們說的話.
如果你翻開詩篇.
你會發現這些詩人.
即是這些信徒.
他們去祈禱的時候.
不是經常天色上藍才禱告.
很多時候是在掙扎.
在投訴.
在軟弱裡面.
在很多申訴裡面.
他甚至會害怕.
甚至會很憤怒.
這些情緒都會在詩篇裡面找到.
我不知道你今天是甚麼情緒.
希望上帝藉著這個視窗.
去跟你整理一下.
到底上帝和你的關係.
我們先禱告.
上主我們多謝你.
因為你是那一位對我們說話的神.
不單止你用你的話語對我們說話.
也藉著我們的心情.
藉著我們每一天有的經歷.
你對我們啟示.
你去到和我們同行.
打開你的話語.

$^{81}$讓你自己的說話.
再一次牧養我們.
帶領我們.
祈禱奉主明教.
我們看詩篇63篇.
大衛在猶大曠野的時候作了者詩.
神啊你是我的神.
切切尋求你.
在乾旱疲乏無水之地.
我的心渴想你.
我的肉身切勿離.
我在聖所曾如此瞻仰你.
為要見你的魂力和你的榮耀.
因你的慈愛比生命更好.
我的嘴唇要頌讚你.
有些詩篇是有標題的.
有個title.
title是甚麼.
大衛在猶大曠野作了者詩.
在曠野的時候.
寫出了這首詩.
大家可能知道.
我們經常說.
曠野其實是巴勒斯坦地.
以色列地.
是一個devil.
是一個本身有的physical setting.
所以曠野總是在舊約和新約聖經裡.
是不能夠逃脫的一部分.
有舊約聖經裡的祈禱.
有信徒有不同的經歷.
有不同的事蹟.
往往都是有曠野在那裡.
猶大曠野到底是甚麼意境.
你試想像一下.
想起聽到猶大曠野的時候.
你會想到甚麼的picture.
在巴勒斯坦地.
是山多平地少.
你由西至東這樣去的話.

$^{121}$由地中海.
一直上山.
上到耶路撒冷的地方.
就是Watershed的地方.
那個分水嶺的地方.
然後到右邊一直降下去.
甚至去到sea level之下.
即是地球最低最低的地方.
就是死海的地方.
過了死海的平原的曠野.
然後再上來.
為甚麼會有曠野呢.
因為那些雨雲是從左至右這樣移動.
當地中海吹雨雲過來的時候.
特別是雨季.
10月底至3月底左右的時間.
Winter的時候.
就是Rainy Season的時候.
那些雨雲就會由左至右這樣移動.
當雨雲由左至右移動的時候.
沿海平原.
在分水嶺左邊的.
以西的.
那些就會有雨水降下去.
過了分水嶺之後.
去到曠野.
就不會再有雨水.
所以為甚麼巴勒斯坦地.
在耶路撒冷.
你就總會很近曠野.
因為走過去就是曠野.
過到這邊.
就是可以說是比較綠一點的.
給大家看兩張照片.
在分水嶺的左邊.
就是經常看到比較綠一點的地方.
在右邊.
你幾乎是找不到有任何的綠草的.
所以.
Wilderness.

$^{161}$永遠都是.
在舊約聖經裡面.
一個Default.
曠野是避不開的.
你想避也避不了.
就因為這個緣故.
聖經經常用曠野.
去做一個人生的寫照.
在生命裡面所遭遇的事.
往往可以用曠野去比喻.
問題就是這樣.
在曠野裡面怎樣去Survive.
怎樣去生存.
或者甚至在曠野裡面怎樣去存在.
《詩篇63篇》.
它就是在曠野裡面寫的一首詩篇.
「神呀你是我的神」.
「我切切地尋求你」.
「在乾旱疲乏無水之地」.
有大曠野.
「我渴想你」.
「我的心渴想神」.
曠野裡面當然你最渴求的.
就是水.
因為那裡是乾旱疲乏無水的地方.
但大衛做了一樣東西很有趣.
他就將他肉身所渴無味.
水.
他就將它轉化過來.
成為一個Spiritual的Yearning.
一個屬靈上的渴求.
因為我們知道肉身上.
我們Thirsty.
我們口渴.
喝水可以解決.
但在心裡面.
在生命裡面.
你的饑渴.
你的Thirst.
是不可以只喝一杯水就解決的.

$^{201}$必須要上帝給到他滿足.
留意他在這裡祈禱.
他說「你是我的神」.
我的神.
通常這些pronoun我們都不理會.
因為叫代名詞.
我的神.
為何這裡這麼重要.
因為大衛和上帝之間的關係.
是一個你與我的關係.
一個Personal的關係.
不是別人的關係.
你留意這裡「你是我的神」.
不是說你是我爸爸的神.
或者我其他弟兄姊妹.
我認識的朋友的神.
是我自己的神.
不是別人的神.
我不知道你認識了上帝多久.
你與上帝的關係.
能不能夠說出「你是我的神」.
還是你只是跟著風.
別人做那樣東西.
你就做那樣東西.
別人這樣行徑.
你就這樣行徑.
真正的信仰就是「我的上帝」.
這是很大的分別.
特別是在曠野的時候.
逃亡的時候.
在詩篇下面.
他正在逃亡.
因為有人追殺他.
但在曠野裡他沒有忘記.
就是上帝和他之間的關係.
你想想曠野.
就是一些磨練的日子.
Trial and Temptations.
你想想在舊約裡面.
新約裡面.

$^{241}$有些什麼事跡.
是與曠野有關.
Wildness.
以色列人出了埃及.
在曠野四十年.
上帝要去磨練他們.
耶穌傳道之前.
他要去四十天的曠野試煉.
曠野就是一個試煉的時間.
我不知道你今天是不是在一個曠野的時間.
大衛說「你是我的神」.
他還有做了一件事.
使他在曠野可以生存.
可以survive and thrive.
可以存在.
「我在聖所曾如此瞻仰你」.
「為要見你的能力和你的榮耀」.
「因你的慈愛比生命更加好」.
「我的嘴唇要去頌讚你」.
「還活著的時候我要舉起我的手來敬拜你」.
在曠野是帝國.
他做了什麼.
記憶.
他記起以前的事.
我曾經做過一些事.
在聖殿裡面.
在耶路撒冷他曾經敬拜過上帝.
他回憶的是什麼.
他回憶的是他曾經有的spiritual experience.
在曠野補水的地方.
他回憶的就是以往他曾經經歷過的experience和上帝.
他們的夢.
怎樣在曠野裡面能夠生存能夠存在.
我們經常都以為.
總之上帝帶我離開曠野.
轉變我現在的情況.
使他更加好.
或者將他扭轉過來.
化腐朽為神奇.
我們就覺得這樣我們就讚美了.

$^{281}$當然這個是應該的.
我們都會這樣祈禱.
你病的時候你希望不病.
健康.
你失業的時候希望有職業.
你讀書很困難的時候.
你當然希望每樣東西都可以順利.
我們覺得如果轉化過來的時候.
我們就可以生存.
就可以去絕狀.
但這裡相反.
不知道你有沒有留意.
他說什麼.
他在曠野裡面.
他記憶回憶以往的事情.
我不知道你有沒有留意過.
可能有些人.
有了上帝的祝福的時候反而就是亡因.
對不對.
事實上舊約裡面以色列人都是這樣.
當他們住了進去迦林地的時候.
安居了.
博業了.
上帝在哪裡.
他們就去拜賈神.
他們去拜偶像.
所以不是說那個情況轉了就搞定.
往往情況改善了.
我們就去拜偶像.
換句話說.
上帝對我們來說.
那時候就是工具.
用完我們就即棄.
但我這裡很謹記.
上帝是他的神.
在曠野裡面.
他不只是求離開曠野.
他見到的是什麼.
上帝是他的上帝.
是他的神.

$^{321}$他學想他.
並且想起以往跟隨上帝.
他們的夢.
能夠維持住.
就在他曠野裡面都可以存在.
告訴我們記憶是多麼重要.
你想想記憶是多麼重要.
你想清楚.
記憶就是我們的定義.
就是我們的身份.
你要去認識一個人.
或者介紹你自己.
去認識一個新的朋友.
你說什麼.
就是你的記憶.
你告訴他你的故事.
我在哪裡來.
我曾經經過什麼.
My story.
你如何認識一個人.
你聽他的故事.
你聽他有什麼經歷.
有什麼經歷.
你就可以認識他.
Relationship.
就是從記憶來的.
我們的身份就是從記憶來的.
我不知道你認識一些人.
或者有沒有跟一些.
Alzheimer 的人相處.
我們以前叫他老人痴呆.
現在我們叫他另外一樣.
Alzheimer.
有沒有試過跟這些人相處.
很多年前我的外母.
就是有Alzheimer.
你知道痛苦的地方在哪裡.
痛苦的地方不只是肉身.
肉身當然很痛苦.
因為他不斷地衰敗.

$^{361}$不斷地扭壞.
身體有些器官不適.
但最掙扎的.
其實就是沒有記憶.
沒有記憶.
不記得所有事情.
或者應該記得的不記得.
不應該提的他就提.
整個人亂了.
你說他是誰.
可能自己也搞不清楚.
最後那幾年.
他曾經在芝加哥.
住過一段時間.
出了去散步.
又不懂得回來.
以前認識的人.
他不認得.
不過感恩他最後的時候.
他也記得我.
化虎令狗.
你想想記憶的Lost.
就是Lost yourself.
就是沒有了你自己.
迷失了你自己.
到最後的時候.
可能連信心的能力也退化.
信仰的能力也退化.
我相信上帝特別眷顧這種軟弱的.
但是重點在哪裡.
我想告訴你.
記憶是很強大的東西.
我們的Spiritual life.
我們的崇年生命.
是靠什麼來建立.
不是經歷多些高山的情況.
或有什麼復興等等.
是靠Look into the past.
上帝怎樣去恩代過.
記起上帝的信心.

$^{401}$過去現在將來.
這些Memory都是累積回來的.
所以為什麼在經過圖的生活裡面.
經過圖的Journey裡面.
經常會去到Accumulate.
去到累積一些東西.
去Deposit一些東西.
每一個敬拜的經歷.
都是一個儲蓄的經歷.
讓我們將來可以Recall.
好像學一個新的語言一樣.
學一個新的技巧一樣.
是Little and often.
你慢慢去到累積的時候.
你的生命才可以去到存在.
謹記上帝在你身上做過的事情.
謹記在你身上上帝有的信心.
第三章第七節.
因為你曾經幫助過我.
我並且要在你翅膀的任何地方.
記起上帝曾經怎樣保護過他.
帶領過他.
這個使他能夠在曠野無水之地.
能夠繼續走下去.
曠野你走不下去.
我們總會有曠野.
你今天可能是青草園.
明天可能就會有曠野.
翅膀的陰霞是什麼意思.
翅膀的陰霞.
在舊約裡面.
這個說話不時會出現.
出埃及的.
拯救意識的人出埃及之後.
他說你看不看到.
我在埃及人當中所行的歧視.
你們既然看見的話.
你就知道我好像鷹一樣.
背你們在翅膀上.
帶你們歸給我.

$^{441}$Save you to be with me.
Return to me.
翅膀的陰霞.
其實是一個神學術語.
告訴你.
上帝曾經恩代過你.
拯救過你.
就好像在翅膀的陰霞裡保護你.
摸了四十年曠野之後.
帶著意識的人.
臨到應許之地的邊緣的時候.
快要過河.
去到耶利哥.
快要入主加林的時候.
他再回憶上帝在他們身上做過的事.
說了什麼.
耶和華在曠野之地看顧你們.
看顧他.
就是意識.
鷹這樣攪動巢穴.
使初鷹在上面飛翔.
展開雙翅.
接住初鷹.
背在兩翼之上.
耶和華也是照樣獨自去引導.
並無外邦的別臣與他同在.
耶和華使他的指紋.
請辭在地的高處.
四十年曠野完了.
一想起.
回憶.
記起的是什麼.
不是曠野有多辛苦.
不是曠野有多難.
你想想你的人生.
過了一段時間.
或者可能過了很久.
如果你回想起的事情.
就是有多苦的時候.
你想想那個生命.

$^{481}$有多麼乾.
因為你想想你人生裡所想起的.
只不過是什麼.
你所記載的.
或是主旋律的.
你所想起的.
只不過是你的失敗.
考試失敗.
戀情失敗.
喪偶.
失去了我們最愛的人.
失去了我們最希望有的情況.
甚至可能婚姻破裂.
如果你想起的.
只不過是這些事情的時候.
你想想生命是多麼乾.
如何能夠存在.
摩西和以色列人再做一些回憶的時候.
他們就想起什麼.
翅膀的任務.
就是上帝如何帶領他們經過.
你想想反過來看.
如果看回我們一輩子的時候.
我們想起的.
主旋律不是那些失敗的.
失去的.
而是想起上帝如何拯救.
上帝和我們同在.
你想想你的命.
你的生命.
生命的最後回憶.
是上帝曾經在曠野裡帶領回憶.
大衛在詩篇63篇.
同樣一樣.
我要在你翅膀任下歡呼.
上帝的拯救.
值得我們的敬拜.
值得我們的歡呼.
並且給予我們生命有力.
我希望你的記憶是一樣.

$^{521}$我希望你的主旋律是一樣.
你有的記憶.
最主要的東西.
所裝載的東西.
不只是我們生命裡有幾幅.
曾經跌過多少次.
失敗過多少次.
上帝在當中.
他們有的恩典.
這個才是最重要的.
翅膀還可以提到什麼.
你記得初夏那段日子.
他們弄藥櫃的時候.
沒有提到翅膀.
就是弄上帝的藥櫃.
就是上帝的寶座.
裡面有兩塊的法板.
寶座的上面.
是有二基羅柏.
兩隻天使.
兩隻天使是有翅膀.
去遮蓋著藥櫃.
翅膀就想起上帝的致勝所.
Holy of Holies.
在上帝的同在裡.
我要在你翅膀底下.
去歡呼.
縱使在曠野裡.
猶如進入致勝所.
進入上帝自己的同在裡.
見到上帝的榮耀.
大衛所記載的.
在曠野裡.
為什麼可以形容.
所謂什麼.
請持在高處.
可以請持在高處.
就是他來到去敬拜.
想起上帝的恩典.
上帝的榮耀.

$^{561}$記得上帝的同在是沒有間斷的.
詩篇六十三章成為你和我的禱告.
我不管你剛才選了什麼顏色.
總會高呼一些.
可能覺得不是太好的情況.
讓你的記憶去支持你.
上帝的同在.
沒有間斷.
我們唱最後一首的詩歌.
我們敬拜上帝.
這首詩歌就像六十三篇一樣.
告訴我們.
如何可以在曠野裡.
你仍然可以請持在高處.
讓你的主旋律記起上帝和你的同.
無論你今天什麼情況都沒有.
什麼心情都有.
這個視窗就是上帝對你說話.
我們唱這首詩歌.
《天父》 詞:陳家麗 曲:陳家麗.
天父 誰聽我禱告.
我軟弱 似我力量.
當我困倦時 請聽我傾訴.
每一天都我自虧.
因你是我的避難所.
你是我力量 是我的拯救.
給我盼望 住一點豐富.
靠著你 我一生變得善良.
你是我力量 是我的盾牌.
你是我盼望 是我的惜望.
我在你陰霾 每天高唱.
你是我永不改變.
神啊 請你堅固我.
我在你陰霾 每天高唱.
你是我永不改變.
神啊 請你堅固我.
我和岑德都期盼.
希望這段時間能給你一段時間.
去到神的面前.
告訴祂你的心情.

$^{601}$讓上帝藉著這個時槍對你下水.
岑德,謝謝你.
讓我們能夠去忘記你的榮耀.
在基督的福音裡.
就看到你自己的愛字.
將你的榮耀彰顯.
將你的愛彰顯在十字架上.
十字架上的救恩.
在自己墳墓裡的見證.
給我們新的修行.
給我們一個超越任何境況的盼望.
神啊,求你祝福自己的恩典.
林道慧不會入圍.
讓你的同在繼續背我們向前.
在你的翅膀底下.
在你自己榮耀的底下.
讓我們知道.
你自己的同在沒有間斷.
為了每一位弟兄姊妹.
無論生命裡,學業裡,職業裡.
任何的境況.
他們遇見不同的難處和抗疫.
主要是你成為他們的力量.
背負他們的經過.
讓他們知道你是那位完全信實的神.
祈禱奉傳明教.
\newpage



\section{詩篇}
\label{sec:X4wQCaypggM}
\textbf{聖經地理 EP2 - 地理與釋經 - 詩篇125-1 (陳崇基牧師)}
\newline
\newline
連結: \href{https://youtube.com/watch?v=X4wQCaypggM}{\texttt{https://youtube.com/watch?v=X4wQCaypggM}} ~~~~ 語音日期: 2022-08-23
\newline
\newline
\hyperref[sec:DvMyIFbhzjs]{\small{< < < PREV SERMON < < <}}
~
\hyperref[sec:index_chronic]{\small{[返順時目]}}
~
\hyperref[sec:index_scriptual]{\small{[返順卷目]}}
~
\hyperref[sec:b63slc4ANbo]{\small{> > > NEXT SERMON > > >}}
\newline
\newline
$^{1}$我們先講講聖經地理和《釋經》有什麼關係.
我不想停留在一個概念性的層面.
雖然概念也很重要.
但是我們先給大家一個樣板.
到底為什麼要認識聖經地理呢.
我們先看一段經文.
就是詩篇第125篇.
上行之師,倚靠耶和華的人,就好像石安山.
安穩坐鎮,永不動搖.
縱山怎樣圍繞耶和華也照樣圍繞他的百姓從今時至永遠.
惡人的象,必不在二人的土地停留,免得二人伸手作惡.
耶和華求你善待行善的心裡正直的人.
至於那偏行彎曲道路的人.
耶和華必將他們和作惡的人一同驅逐出去.
願平安歸於以色列.
這篇詩篇可以做一個很好的導演.
就是為什麼要知道聖經地理.
對解釋這篇詩篇有什麼好處呢.
首先第一樣.
大家知道120篇至到134篇15首的詩歌.
叫做上行之師.
Shia lam ma lu.
什麼意思呢.
這個字其實就是Alaa,就是向上的意思.
高up的意思.
上行之師就是向上走的意思.
甚至有些中文翻譯叫做登階之師.
就是走樓梯上樓梯的詩歌.
大家可能聽過一些說法.
耶路撒冷南面入聖殿的時候.
有15級樓梯上去.
或者拉比文學說到不是聖殿山以外的15級樓梯.
而是以色列院裡面.
內院裡面進入聖殿的時候有15級樓梯.
這15首的上行之師就是利未人的師班.
每一級走上去的時候唱一首.
第一級唱120篇.
上行之師第一篇.
然後唱到第15篇134篇.
但是這個有地理的因素在裡面.

$^{41}$上行之師為什麼是上行呢.
很明顯就是去耶路撒冷敬拜神.
耶路撒冷就是石安.
所以你聽見舊約裡面說.
耶路撒冷其實就是石安.
其實就是石安山.
兩個的名字其實都是指同一個地點.
石安山有754米高.
即是海拔以上.
或者2400多尺海拔以上.
我們可以看一下Google map.
看看有什麼幫助.
Google map.
我們下去耶路撒冷的時候.
我們去到今天的所謂Old City這個地方.
Old City這個地方.
白色這個地方.
看到嗎.
我們將它3D這樣看的時候.
今天你去到耶路撒冷.
所見到的那個城牆.
就是一個16世紀的時候.
才興建的.
即是幾百年前才興建.
在奧特曼時期.
即是土耳其奧特曼帝國.
Ottoman Empire.
那時候才興建的耶路撒冷.
這個城牆.
所以這個城牆不是舊約聖經.
或者不是新約耶穌時代的城牆.
怎樣也好.
這個地方其實就是.
所謂舊約新約所指的耶路撒冷.
當然不同時代.
耶路撒冷的面積有小也有大.
按照一個城市的發展.
人口多了.
人史又少了.
多了又少了.

$^{81}$所以在不同的時間裡面.
耶路撒冷的城的邊界.
會大或者會小.
但是耶路撒冷在哪裡呢?.
在海拔的二千多呎以上.
所以如果你要去耶路撒冷的話.
你從哪個方向去.
上行之師.
那些朝聖者.
在耶路撒冷以外.
或者甚至加尼尼.
上面加尼尼海這裡.
他們去耶路撒冷的時候.
你從甚麼方向.
大致上都是上山.
所以為甚麼上行之師.
就是上去耶路撒冷朝聖的師過.
甚麼叫上去.
因為地理上確實是上去.
你記得.
如果我們轉過來這邊.
你見到有耶利哥.
耶利哥這裡.
上面就是耶路撒冷.
有東山有西山等等.
我們等一下再說.
這裡就是海拔以下的地方.
你見到旁邊就是死海.
約旦河的水流到死海.
就沒有地方可逃.
一定在這裡蒸發.
所以死海的水就越來越鹹.
因為水蒸發了.
就留下鹽在下面.
所以人們去死海玩的時候.
就是去那裡浮.
你不懂游泳都可以.
你下去就會浮.
所以死海的地方.
就是在海拔以下好幾百米的水平.

$^{121}$所以是全地球最深的地方.
算是地面最深的地方.
就是在死海這個地方.
死海旁邊就是耶利哥.
所以你見到.
Google map都可以看到.
耶利哥就是在這個平原上面.
死海的平原那裡.
然後你上耶路撒冷.
就是一直去到上山.
那個好撒瑪利亞人的比喻.
大家都記得.
耶穌說什麼.
有一個人從耶路撒冷下耶利哥.
不是上 反過來是下.
因為耶路撒冷是比耶利哥更加高的.
所以由什麼方向.
上去石安山或者去耶路撒冷.
都是上山的.
所以就有所謂上行之師.
今天以色列的猶太人.
如果有些散居在不同地方的猶太人.
他們說要移民回去以色列.
今時今日的現代的以色列國.
這個移民的過程.
都是叫上去.
Aliya.
Aliya其實就是阿拉這個字變出來.
所以他們說整個的移民的過程.
就叫做Aliya 上去的意思.
就是上去耶路撒冷.
所以這些上行之師.
絕對當然就是摩西五經裡面所說.
一年有三次不同的節期.
猶太人特別是南丁.
就要上到去耶路撒冷.
這些上行之師大概就是.
他們上路的時候.
去到耶路撒冷的時候.
那個過程裡面會唱的一些詩歌.

$^{161}$如果你看一下耶路撒冷的地理.
你就知道.
這個石安山.
耶路撒冷.
或者最初的時候叫大衛城.
因為大衛攻取了耶布斯人的石安.
Stonehall Zion.
他就將它命名為City of David.
大衛之城.
大衛城.
所以指的都是同一個地方.
大衛城大概就是2300呎海拔以上.
你看到四周的東南西北.
不同的山都是比大衛城高還是低.
大衛城有2300呎海拔以上.
北面的聖殿山就是2400多呎.
所以比它高.
右邊東面的橄欖山.
就是2500呎高.
南面的有一個聯合國的辦公室.
今天還在.
南面的有一些比較高的山.
或者有一個地方叫Hespalmanay.
很多時候人們會去那裡拍照.
比較高一點的地方.
拍回去耶路撒冷.
有南面拍照去耶路撒冷的一個地方.
一個公園.
都是2536呎.
所以比大衛城石安比較高.
接著你看到西面.
左邊就是西山.
今天的舊城.
耶路撒冷舊城Old City.
有一個地方叫Mount Zion.
以前的名字亂了.
所以本身的地理叫Western Hill.
就是西山.
你看到有多高.
2541呎.

$^{201}$所以四面的山都是比耶路撒冷高.
如果你看耶路撒冷的地勢.
這裡是Google Earth.
你都可以看到.
這個綠色的地方.
就是大衛城.
原本的石安就是在這裡.
然後剛才白色的那個紙作.
就是奧特曼時代的耶路撒冷.
四四方方的.
一個公里.
One kilometer.
乘以一個公里.
很小的地方.
這個就是所謂的耶路撒冷.
三邊是有山谷的.
黃色這條線就是恩聯谷.
紅色這條線就是喀倫谷.
如果你站在這個地方.
站在綠色的上面.
就是所謂的大衛城的地方.
大衛的皇宮大概會在這裡.
因為聖殿在上面.
黏著聖殿的南面的.
就是大衛的皇宮.
所羅門的皇宮等等的地方.
所以如果你站在這裡的話.
你周圍前後左右.
東南西北這樣看一看的時候.
你就會發現上面聖殿山.
就是今天這個金頂清真寺的地方.
主要是在第七世紀六百多年的時候才建的.
金頂清真寺.
比這個地方是比較高的.
北面就是聖殿山.
有多高呢?.
2446呎.
海拔以上.
然後旁邊這邊.
你會看到白白的地方.

$^{241}$這些山峰就是監覽山.
Mount Olives.
如果你轉過來的話.
Mount Olives上面.
很多時候我們去耶路撒冷的時候.
我們拿到耶路撒冷最好的拍照的地方.
就是這個地方.
就是Mount Olives.
Mount Olives比這裡的金頂清真寺更高.
海拔以上2500呎.
甚至更高.
所以東面比石安更高.
北面的金頂清真寺.
也是比大衛城更高.
西面就是這裡.
這裡.
轉過來.
在這個恩怨之谷上面.
恩怨之谷上面.
你會看到有這個所謂的西山.
今天叫做Mount Zion.
石安山的地方.
比大衛城石安更高.
南面這裡.
剛才我說到公園.
聯合國的地方.
他的辦公室.
也是比石安原本的大衛城.
耶路撒冷更高.
所以你看到四面環山.
都是包圍著耶路撒冷.
所以你看到東南西北.
這樣包圍著的時候.
你就明白了.
這個詩篇到底想說什麼.
\newpage



\section{詩篇}
\label{sec:b63slc4ANbo}
\textbf{聖經地理 EP4 - 地理與釋經 - 詩篇125-3 (陳崇基牧師)}
\newline
\newline
連結: \href{https://youtube.com/watch?v=b63slc4ANbo}{\texttt{https://youtube.com/watch?v=b63slc4ANbo}} ~~~~ 語音日期: 2022-08-23
\newline
\newline
\hyperref[sec:X4wQCaypggM]{\small{< < < PREV SERMON < < <}}
~
\hyperref[sec:index_chronic]{\small{[返順時目]}}
~
\hyperref[sec:index_scriptual]{\small{[返順卷目]}}
~
\hyperref[sec:USIUOn0X4jM]{\small{> > > NEXT SERMON > > >}}
\newline
\newline
$^{1}$我們繼續講詩篇125篇.
我們回到剛才所講的那一節.
就是寫125篇的人一定是熟悉耶路撒冷的地理.
他才能夠寫出來.
這個就是耶路撒冷.
回到這裡就是看到石安大衛城的地方.
中山怎樣環繞耶路撒冷.
東南西北去到抱擁著石安.
就引伸出耶路撒冷的保護.
依靠耶路撒冷才能像山崗一樣不會動搖.
然後他繼續講.
惡人的權柄惡人的國度.
惡人欺壓義人的國度.
必定不會留在義人的土地.
以免義人會跟他一起作惡.
所以這裡的詩人就說一定要趕絕這些惡勢力.
這些反對上帝的國度 勢力 君主等等的權勢.
然後耶和華說 耶和華求你善待行善心裡正直的人.
因為這些人才能夠永存.
因為有耶和華去環繞著他們.
以至那偏行彎曲道路的人.
耶和華必將他們和作惡的人一同驅逐出去.
在耶路撒冷這個地方.
在石安這個地方 敬拜神的地方.
只有義人才能夠進來.
詩篇由第一篇開始到現在都是這樣說.
要敬拜耶和華 要親近他.
這些人一定要全心全意去依靠耶和華.
這些人才能得到平安.
得到上帝所賜的福.
就是要依靠耶和華的那一群人.
我們繼續看下去.
我們發現這篇詩篇說到平安.
在上行之詩篇常提到平安.
平安從哪裡來.
平安就是在耶和華所住的地方.
就是耶路撒冷 大衛城石安山.
去敬拜他聖殿的地方.
他就會有平安的源頭.
真正的平安是在耶路撒冷那裡出來.

$^{41}$耶路撒冷是怎樣的.
這篇詩篇有提過.
當然有很多其他詩篇特別上行之詩都有提過.
耶路撒冷的原文是Yerushalayim.
很奇怪 希伯來文的運作.
字尾im其實是雙數字尾.
或者是眾數字尾.
譬如水.
希伯來文的水是Mayim.
又是Yim.
是結尾的.
又是雙數字尾.
當然有時候這些文法純粹是比較生硬的.
沒有什麼意義.
但在耶路撒冷似乎有這樣的意味.
因為說到耶路撒冷.
Yerushalayim這個字.
這個城的名字.
為什麼是雙數.
或者眾數.
因為耶路撒冷其實是在兩個山上的城市.
一個叫東山.
一個叫西山.
一個叫Eastern Hill.
一個叫Western Hill.
為什麼呢.
因為你看回這個.
耶路撒冷的地理.
你就會看到.
白色這個是什麼 現在的Old City.
綠色這個地方.
已經有一條線在這裡伸出來.
就是叫做Eastern Hill.
而這一邊上面.
我們剛才在視頻裡說.
360度那裡.
就是這個Western Hill.
比這個大衛城更高.
這個叫做西山.
耶路撒冷.

$^{81}$無論哪個時代.
它的大規模.
有大有小.
但都是坐立在這兩個山上.
第一個就是Eastern Hill.
東山.
東山最高的地方.
天然最高的床岩.
山崗最高的地方.
就是這個金頂清真寺.
這個就是東山最高的峰頂.
這個金頂清真寺.
這個金頂.
就是在這個峰頂建的.
所以第一聖殿.
或者第二聖殿的致勝所.
應該是在這個金頂下面.
也都按照以前中東.
一直以來對聖地的概念.
是經常都不改的.
一層又接一層.
之後建下去的廟或者會堂.
或者其他的殿.
甚至之後的穆斯林所建的清真寺.
都是建在神聖的地點上.
那些叫做Continuity of Sacred Space.
即是那個地方歷史以來.
已經是一個神聖的地方.
在歷史的後期.
不斷都會在那個地方.
是不會改的.
有一個保守性.
所以我們在地理上.
我們知道有很多地方.
以前舊約裡面提過.
新約裡面也提過.
現代的找出來的.
都是承接著這個歷史.
所以我們那些聖經的地圖.
不是靠猜測的.

$^{121}$是背後有一個科學.
有一個研究.
然後我們知道大部分舊約.
新約裡面的地點到底在哪裡.
這個是後話.
但是耶路撒冷.
為什麼要用Selium.
或者是雙數美.
或者是眾數美.
因為在兩個山上面.
Double Hill.
或者City of Two Hills.
Eastern Hill和Western Hill.
除了這件事.
我們也可以看到.
在一些聖經的地圖.
你也可以看到.
很多時候.
耶路撒冷的平面圖.
是這樣畫的.
這個就是耶穌時期的耶路撒冷.
為什麼會這樣.
大衛城長長的一條.
然後上面就是聖殿山.
原因就是.
好像左邊這個地形圖.
你就知道.
下面Eastern Hill.
你看到其他的線.
是比較密的.
比較傾斜.
所以山頂就是長長的山脊.
Eastern Hill最高的地方.
就是這裡.
就是Mount Moriah這個地方.
就是摩利亞山.
就是獻以殺的地方.
這個名字.
摩利亞.
這裡也是第一聖殿.

$^{161}$第二聖殿.
至聖所聖殿的地方.
所以在Eastern Hill上面.
是比較長的.
所以我們有時候看聖經的地圖.
雖然是平面的.
沒有地勢的圖案.
讓我們知道哪裡高哪裡低.
但你就知道為什麼大衛城會這樣畫出來.
然後到這個上城的地方.
就是這個Western Hill的地方.
今天叫做Mount Zion.
石安山的地方.
其實石安山原本就是在這裡.
在東山的上面.
所以你可以.
從這個地形.
就可以了解.
你本聖經裡面.
有的聖經地圖.
原來就是.
為什麼會這樣畫出來.
因為耶路撒冷就是.
City of Two Hills.
City on Two Hills.
或者.
耶路撒冷的另一個意思.
就是講到它字的名字.
剛才我們講到地形.
Eusulam.
City on Two Hills.
它的名字.
大家很清楚了.
就是City of Peace.
City of Shalom.
創世記十四章十八節.
這個亞伯蘭.
打贏了縱王.
救了羅德.
記得嗎.

$^{201}$回來的時候他就見到墨基西德.
當時還是迦南營的地方.
這個墨基西德叫做什麼.
Malik Shalem.
就是King of Peace.
撒冷王的意思.
撒冷就是Shalem.
Shalem.
王Malik.
就是撒冷王.
所以這個地方叫做撒冷.
Shalem就是平安的意思.
為什麼會是耶路撒冷呢.
可能呢.
就是來自.
一個阿卡德文.
叫做Ru.
你懂一點點希伯來文的組合.
有時候加一個字首.
加一個Y或是一個YOR在前面.
就是叫做Yerushalem.
Yeru.
Yerushalem.
Ru是什麼意思.
在阿卡德文裡面.
就是城市的意思.
所以Jerusalem.
就是Shalom.
City的意思.
平安的城.
所以你看到在詩篇125篇.
他就說到.
去敬拜耶和華.
去到石安山的時候.
才會有平安.
會有Shalom.
Shalom說過很多次了.
在希伯來的概念裡面.
Shalom.
平安.

$^{241}$不是無災無劫無戰爭那麼簡單.
而是整個生命的慷態.
就是耶穌說.
我來時叫人得生命.
我來時得得更豐盛.
那種就是Shalom.
不單止是無災無劫.
並且有著上帝的祝福.
上帝的同在.
無論什麼情況都好.
有上帝的同在.
這就是真正的Shalom.
所以一篇上行之詩.
說到上去耶和華石安山去敬拜的時候.
你會發現原來有很多背景.
特別地理.
這篇詩篇的信息就會豐富很多.
\newpage



\section{}
\label{sec:USIUOn0X4jM}
\textbf{聖經地理 EP5 - 地理與釋經 - 中間之地-1 (陳崇基牧師)}
\newline
\newline
連結: \href{https://youtube.com/watch?v=USIUOn0X4jM}{\texttt{https://youtube.com/watch?v=USIUOn0X4jM}} ~~~~ 語音日期: 2022-08-25
\newline
\newline
\hyperref[sec:b63slc4ANbo]{\small{< < < PREV SERMON < < <}}
~
\hyperref[sec:index_chronic]{\small{[返順時目]}}
~
\hyperref[sec:index_scriptual]{\small{[返順卷目]}}
~
\hyperref[sec:dVqF8V7wdDo]{\small{> > > NEXT SERMON > > >}}
\newline
\newline
$^{1}$我們又再來另一個部份,就是講到地理和《釋經》有什麼關係.
有很多《聖經》的經文,如果你知道地理的時候.
那個意義,神學的信息就會更加清楚和更加表明出來.
在這裡我給一些經文給大家去嘗試去掌握一下.
並且有些地理上的動態可以跟大家解釋.
然後看回這些經文的時候,意義就會非常不同.
以色列,巴斯坦地是古代世界的中心.
Center of the world.
《聖經》有些經文是講這件事.
地理當然也是講這件事.
所以有些概念我們可以嘗試去掌握.
以色列有很多地名.
不同的人在歷史裡面都會叫以色列這塊地方有不同的名字.
有時候叫它以色列,Israel.
或者叫它巴斯坦.
巴斯坦主要是在第二世紀的時候.
羅馬帝國為這塊地起名.
用非利士這個字叫這塊地的名字.
叫它巴斯坦.
Palestine.
之前就是非利士.
所以這個名字就流傳到現在.
叫做巴斯坦.
Israel, Palestine.
有些人會叫它聖地.
聖地旅行團.
聖地學習團.
去聖地旅遊.
叫做Holy Land.
這個有宗教性的意思.
之前那兩個就是比較政治性的名稱.
另一個名稱就是聖地.
Holy Land.
為什麼呢.
因為在耶路撒冷.
是世界上三大獨神宗教信仰的基地.
三個最大的獨神宗教.
Monotheism.
Monotheistic Faith.
都是以耶路撒冷為它的發源地.

$^{41}$或者基地.
哪三個呢.
基督教當然.
舊約新約.
整個的新舊約傳說.
基督教Christianity.
當然就是以耶路撒冷為它的基地.
猶太教.
只是認定舊約.
所謂希伯來聖經.
是他們的經典.
是他們的正典.
對猶太教來說.
這個耶路撒冷.
就是他們的Holy Land.
第三個當然是穆斯林.
或者伊斯蘭.
剛剛我們看到.
之前看過一些清真寺.
這個是主後之後.
伊斯蘭的入侵之後.
Islamic Conquest之後.
大概主後六百多年.
七世紀的時候.
侵佔了這個地方.
然後有很長的伊斯蘭傳統.
所以在這個地方.
就會建立了金頂清真寺.
或者Al-Aqsa Mosque.
其實意思就是說.
是最遠的清真寺.
因為他們的發源地本身在沙利阿拉伯.
在Mecca這個地方.
在Medina這個地方.
在現在的沙利阿拉伯.
在南邊.
但是在一些穆罕默德的傳說裡面.
可蘭經的傳說裡面.
就是穆罕默德在夜間.
有一個異象.

$^{81}$就是來到最遠的清真寺.
拿了經.
拿了他們的經典.
所以就在這裡叫Al-Aqsa Mosque.
就是The Farthest Mosque的意思.
就是最遠的清真寺.
所以在這裡.
就是穆斯林.
伊斯蘭認為是一個聖地.
所以有很多人就叫聖地.
聖地你都要問清楚.
對誰來說是一個聖地呢.
哪個宗教.
第四個.
就是Levant.
意大利文來的.
就是以前那些歐洲人.
來探索聖地的時候.
探索巴勒斯坦地的時候.
他們起了這個名字.
叫做Levant.
Levant是什麼意思呢.
就是Rising的意思.
這個名字這個標籤.
就有西方的味道.
因為什麼叫Rising.
什麼叫日出之地.
就是從西方.
西方就是.
歐洲.
歐洲的大陸.
去看這個.
意大利在這裡.
然後他們看過來.
這個巴勒斯坦地.
他們就說.
這個是日出之地.
東方之地.
這裡就很有歐洲的.
政治的標籤.

$^{121}$但是更加好的.
一個地理的說法.
其實就是.
叫做中間之地.
Land between.
這個是最忠誠的.
純粹是講地理.
就是那個地的.
動態.
大家之間的關係.
為什麼叫做中間之地.
或者叫做地橋.
因為.
很明顯.
這個就是當時.
古代世界的中間的地方.
小亞細亞.
歐洲.
就是在他的左上角.
亞洲.
西亞.
米索不達米亞.
就是在他東邊的地方.
在他左下角.
就是非洲.
這裡看到長長的綠色的就是尼羅河.
然後下來下面的平地.
就是分散了很多支流的尼羅河.
三角洲.
所以這個地方.
就是這三大洲的.
古代世界.
三大洲的中間的地方.
中心的地方.
是一個地橋.
因為你要在三大洲.
大家彼此通往的時候.
就要經這個地方.
這個就是叫做.
Geography.

$^{161}$就是地理.
Geographical dynamic.
就是地理的動態.
是怎樣.
這個圖就更加清楚.
給你看到.
迦南地.
舊約裡面叫做迦南地.
他就是三大洲.
中間的地方.
這個綠色的地方.
在.
1920年代30年代的時候.
有一個學者.
在芝加哥大學.
他研究中東.
近東文化歷史.
他就提出.
就說.
這個綠色的地方.
就是剛剛夠雨水.
能夠耕種.
能夠安居的地方.
剛剛好是夠水的.
因為有河道.
因為有雨水.
他就說.
這個綠色的地方叫做Fertile Crescent.
新月沃土.
所以由米索不達米亞右邊.
米索不達米亞.
其實就是.
兩河流域的意思.
這個是希臘文.
就是說兩條河下面.
就是底格拉斯河.
然後左邊.
最盡的這邊就是.
埃及.
上面就是今天的土耳其.

$^{201}$以前就叫做亞納托里亞.
或者叫做小亞,細亞.
中間就是迦南地.
在這個地圖.
這一處.
你能夠觀察到的就是.
在整個的.
這個.
古代的世界裡面.
古代人類最早文明.
開始的地方.
在米索不達米亞.
在整塊新月沃土.
的整個世界裡面.
最好的地方.
都是在兩個的盡頭.
第一個就是.
米索不達米亞.
因為長久都有清水.
有河水能夠耕種.
在左邊更加.
就是埃及.
有利羅河.
三角洲.
你今天.
去利羅河看.
在衛星圖片.
你都能看到.
是幾個肥肉的地方.
全部都是綠色.
你見到利羅河的支流.
旁邊是什麼.
其實全部都是農田.
全部都是農田.
平平的農田.
所以是一個.
很值得耕種.
很適合安居的地方.
米索不達米亞.
吾爾哈蘭等等.

$^{241}$巴比倫都是一樣.
有大的河流.
米索不達米亞.
和埃及.
都是有的.
唯獨是沒有的.
就是迦南地.
我們繼續講下去.
地理的情況.
就是沒有大的河流.
最大的那條是什麼.
約旦河是怎樣來的.
就是由旦這個地方.
上面就是黑門山.
黑門山就是.
長久有積雪的.
因為很高.
所以在黑門山.
這個地方.
Mount Hermon.
水就會.
溶了流下來.
然後流經旦.
這裡有水泉.
然後流下去.
加里里海,加里里湖.
清水的.
但希伯來文是沒有分別的.
海和湖是一樣的.
所以是加里里海,加里里湖.
加里里湖南面再出來.
就是約旦河.
一條很幼的線.
一直流經峽谷.
一直流下去.
就是死海.
所以在.
迦南地.
最大的河就是約旦河.
所以在迦南地生存的話.

$^{281}$沒有大河給你.
你只能靠.
雨季的時候.
冬季的時候有雨水.
吹過來的時候.
就灑雨水在地上.
以致你能夠在這裡耕種.
所以這個地方.
就是一個.
世界的中心.
所以這個地方.
就是一個世界的中心.
迦南地.
是世界的中心.
有什麼.
你明白了地理的.
開放的時候.
你就可以明白聖經裡面所說的.
到底神做了什麼.
我經常和初信的人.
他們不是很明晰的背景的人.
就會和他們解釋.
為什麼會在耶路撒冷開始.
耶穌會在耶路撒冷成就救人.
原因就是.
世界的中心就是.
這個地方.
聯絡所有.
大的.
綠洲.
都是在.
三個大的大陸.
非洲,歐洲,亞洲.
都是藉著.
迦南地來聯絡.
這是世界的中心.
世界的中心.
\newpage



\section{}
\label{sec:dVqF8V7wdDo}
\textbf{聖經地理 EP6 - 地理與釋經 - 中間之地-2 (陳崇基牧師)}
\newline
\newline
連結: \href{https://youtube.com/watch?v=dVqF8V7wdDo}{\texttt{https://youtube.com/watch?v=dVqF8V7wdDo}} ~~~~ 語音日期: 2022-08-25
\newline
\newline
\hyperref[sec:USIUOn0X4jM]{\small{< < < PREV SERMON < < <}}
~
\hyperref[sec:index_chronic]{\small{[返順時目]}}
~
\hyperref[sec:index_scriptual]{\small{[返順卷目]}}
~
\hyperref[sec:PjaM3O4lbt8]{\small{> > > NEXT SERMON > > >}}
\newline
\newline
$^{1}$有了這些地理知識就可以去解經.
解什麼經呢?.
譬如這一段.
《申命記》第四章第六節.
講什麼?.
他說摩西和那些以色列人.
他們快要進入迦南地.
在哪裡講的?.
《申命記》是在耶利哥對面的摩亞平原.
你都可以跟他們講話.
摩亞平原在哪裡?.
就在耶利哥附近.
尼泊爾山.
摩西死的地方.
在哪裡?.
就在這裡.
我剛剛點了出來.
尼泊爾山.
這條是約旦河.
摩西最後的講話.
在《申命記》中揭露出來.
就是跟這些以色列人說.
他們已經是第二代了.
因為他們已經離開埃及.
去了西乃的曠野.
埃及出來了.
然後他們經過西乃半島去到西乃山.
然後在曠野四十年.
第一代已經過去了.
他們繞過伊東地.
摩西的半島.
摩亞地.
然後已經來到尼泊爾山的山腳.
這裡就是摩西跟他們第二次.
跟他們再申報.
重申上帝的律法.
《申命記》就是這個意思.
第二條法律.
整本《申命記》都是摩西的講章.
跟這些以色列人說什麼呢?.

$^{41}$他們準備進入迦南地的時候.
左邊.
準備進入主迦南.
他們在這個.
死海上的.
摩亞的平原.
跟他們說什麼呢?.
他們說你們要謹守遵行上帝的約.
這就是你們在萬民面前的智慧和聰明.
為什麼特別提到萬民眼前?.
當然他們見證就是在眾民之間.
所有外邦的民族之前.
但如果你有了地理的動態.
你就知道.
在迦南地講萬民面前.
就真的是地理的現實.
因為萬國都會經過那個地方.
三個大陸的人.
去第二個地方都會經過迦南地.
在萬民之前.
你們就是成為智慧和聰明.
那些萬民聽見這一切的律例.
就會說這個大國真的有智慧有聰明.
哪有一個大國有神明.
與他們相近的東西.
我們的神.
在我們求告他的時候.
與我們相似的東西.
我們就說我們是神.
我們是神.
在我們求告他的時候.
與我們相近.
哪一個大國有公義的律法典章.
就像我們今天在你們面前.
所頒布的這一切的律例.
沒有一個神會這樣做的.
跟你講清講楚.
他的心意.
只有一個會說是這樣的.
如果你跟手進行的話.

$^{81}$你就成為萬國的見證.
因為你就是在世界的地橋上面.
中間之地.
去見證著神的律法是多麼奇妙.
萬國眼前.
所以你這樣明白.
這個就真的在講.
宣教的見證的地理的現實.
你站在中間之地.
就在萬民的面前.
然後以西結書第五章.
這裡就先知去罵以色列人.
這個耶路撒冷的居民.
說他們背道.
背約背道.
就是耶路撒冷.
我曾將耶路撒冷安置在列國之中.
列邦都在他的四周.
這個又是一個地理現實.
地理的現實.
他們確實就在列邦的四周.
四周就是列邦.
埃及在下面.
然後米索不達米亞.
亞述.
黑人.
今天的土耳其是黑人帝國.
巴比倫等等.
這些大國都是環繞著他們四周.
他們就在世界的中央.
再看下去.
以西結書第38章第12節.
這裡講到以色列的受敵.
就是哥格.
與馬國等等.
他們不詳細解釋.
以色列民住在地的高處的百姓.
他們就是住在地的高處的百姓.
因為中文不能翻譯出來.
我就不明白.

$^{121}$為什麼中文會翻譯成地的高處.
因為Tabu Ha'ales.
什麼叫Tabu Ha'ales.
Tabu其實就是土耳其的意思.
人的中心就是土耳其的地方.
地的中心.
應該是翻譯成住在地的中心中間.
所以新譯本就翻譯出來了.
住在世界中心的人民.
以色列民住在耶路撒冷.
住在迦南地.
就是世界的中心.
Naval or Center of the World.
是不是.
英文翻譯出來了.
和學本就翻譯成高處.
不是很對.
有了地理知識就明白了.
在做什麼.
以色列民就是萬國的見證.
萬國的光.
在十六世紀.
一五幾年的時候.
這幅畫很有名.
你去以色列某些博物館.
都會解釋耶路撒冷在世界的位置.
都會用這幅畫來解釋.
這幅畫就像一朵花.
然後有三塊花瓣.
然後就解釋這三塊花瓣是什麼.
周圍就是海.
有船等等.
這裡就是Europe.
紅色這塊就是Europe.
綠色這塊就是亞洲.
Asia.
然後下面這塊黃的啡的就是非洲.
中間如果你看得清楚的話.
就是Jerusalem.
寫著耶路撒冷.

$^{161}$耶路撒冷整個畫出來.
這裡就告訴你.
在中世紀.
或者有史以來.
以色列.
耶路撒冷.
就是世界的中心.
救恩就是從這裡出來.
直到地極.
耶路撒冷直到地極.
亞洲從亞洲開始直到地極.
美洲在左下角.
又說古代的世界.
人類開始的地方.
文明開始的地方.
就是在這個地方.
然後這個地方的中間.
地橋.
也是耶路撒冷的交通地.
創世記十二章.
我們看看.
這個很熟悉的經文.
對亞伯蘭人說.
你要離開本地本族.
即是米索不達米亞.
米索不達米亞.
他離開了不爾.
不爾就在這裡.
南米索不達米亞.
當時文明最富庶的地方.
有河水.
有耕種.
有很多平地.
有很大的麥田.
很大的米田.
這些就是古史會升的地方.
一定會升的地方.
他就說要離開本地本族.
本家.
他就先去了哈蘭.

$^{201}$然後就去了加拿大.
往我指示你的地方去.
是不是.
他知道要去加拿大.
為什麼要去加拿大.
他還說在這裡.
使你成為大國.
必賜福給你 使你名為大.
你要使別人得福.
為你祝福的我必賜福給你.
我必去咒助你.
這是納粹的語言.
說我一定要保護你.
地上的萬族都要因你得福.
這句說話就有意思.
你看完地圖就知道.
是不是.
因為地上的萬族就是要經過這個地方.
對不對.
迦南地.
再看下去.
史唐傳.
在五秦節的時候.
五秦節的時候.
就說什麼呢.
就是聖靈降臨.
然後在耶路撒冷居住的.
這些歸信耶路撒冷.
歸信猶太教的人.
他們上了聖殿敬拜.
他們做什麼.
他們就聽見.
聖靈的降臨.
然後使徒就用他們的說話.
來跟他們傳道.
他們說什麼.
這些猶太人.
他們就說.
帕提亞人.
馬代人 以蘭人.

$^{241}$和住在米索不達米亞的人.
這些人是什麼人.
米索不達米亞的人.
帕提亞.
就是在這個地方.
米索不達米亞上面.
這個地方.
馬代人也在這裡.
以蘭人就在米索不達米亞.
波斯灣的右邊.
這個地方.
還有住在米索不達米亞的人.
他們就回到.
耶路撒冷朝聖的時候.
聖靈降臨.
他們就知道.
東方這邊的人代表著.
他們就來到耶路撒冷.
聽見聖靈降臨.
然後就回到猶太.
猶太就是耶路撒冷這個地方.
還有什麼人.
猶太人.
然後加帕多加.
本都.
亞細亞.
佛雷加 龐菲利亞.
這些人在哪裡來.
在地圖上.
今天土耳其的地方.
龐菲利亞 本都.
亞細亞.
亞洲大陸.
他們又來到耶路撒冷.
去守節 然後聖靈降臨.
然後他繼續說.
埃及的人.
並靠近古利內.
利比亞一帶地方的人.
咦.

$^{281}$他說的是什麼 埃及.
利比亞就在這裡.
古利內.
你記得古利內人西門被迫去背耶穌的十字架.
他在哪裡來.
在利比亞的古利內.
就是埃及的.
左邊 西邊.
然後他繼續說什麼.
僑居的羅馬人.
他就說到.
這對人.
這裡 意大利.
羅馬.
來朝聖的人.
說了十五種人.
這裡繼續說下去.
就是克里特島.
克里特島.
還有阿拉伯人 南邊.
沙地阿拉伯人.
這些人做什麼.
他們都聽見.
用我們的.
香談.
講論上帝的大座.
你看到那個圖畫是怎樣.
是整個地球的人.
當時的.
世界的人.
都聚集在耶路撒冷.
這些就是那些.
猶太教去到朝聖的人.
代表著世界的萬國.
十五種人.
他們來到耶路撒冷.
這裡就是福音.
發源的地方.
由耶路撒冷到猶太.
然後到撒瑪利亞.

$^{321}$直到地極.
今天就去到世界各地不同的地方.
美國 中國等等不同的地方.
福音繼續擴散出去.
就是從這個地方開始.
這個不單止是一個使命.
一個宣教的現實.
也是一個地理的現實.
由這裡開始.
所以我們看到.
以色列是中間之地.
其實就是一個世界的.
匯點.
等於大城市.
不同的公路.
轉折的地方.
不同的公路.
交叉的地方.
匯點的地方.
交換的地方.
匯聚的地方.
這是交通的中心.
無論是行軍.
或者是商貿.
不同的人來人往的地方.
就是迦南地.
就是以色列這個地方.
世界的中心.
所以有一個聖經地理學者說得很對.
他說在.
這個以色列的地方.
最好做些什麼呢.
就不是建地上的國度.
因為你看看以色列.
由頭到尾.
古代至現代的歷史.
很少有一個時窗.
好像有一個小小的時間.
是有獨立的國度.
你記得舊約裡面.

$^{361}$曾經有沒有.
有 為什麼.
因為是所羅門.
但是不夠八年.
他就敗家.
國度就沒有了.
所以只是一段很短的時間.
能夠曾經獨立過.
地上的國度.
然後到耶穌之前.
瑪加比的時代.
大概167BC.
到60多BC.
一百年左右.
瑪加比的國度.
又是大概只有一百年的歷史.
是他們獨立.
自己有自己的國度.
到今時今日.
現代的以色列國.
1948年復國.
到現在大概70多年.
在這麼悠長的歷史裡面.
很少機會.
這個地方有自己獨立的國度.
因為通常都是.
埃及的現身.
因為埃及霸這個地方.
希望能夠霸佔.
其他更好的地方.
所以是一個踏腳石之地.
在米索不達米亞的政權.
亞述.
巴比倫.
他們想拿這個.
中間之地.
為什麼 因為他們想要落埃及.
要攻佔埃及.
最好的地方都是這兩個.
甚至上面.

$^{401}$土耳其.
和印度.
和埃及的帝國.
彼此相爭在哪裡相爭.
就是爭這一塊地.
兵家之爭的地方.
就是這個地方.
因為這是世界流通商貿最發達的地方.
人來人往最好做什麼呢.
聖經地理學家說.
最好建造的不是地上的國度.
所以神都不是要建一個地上的國度.
是由這個地上的國度.
建一個無界限的天國.
所以最好在這裡建造的是什麼呢.
他就說.
是港台.
是一個傳播信息的地方.
傳福音的地方.
這個就是中間之地.
有的地理.
有它的.
strategic的地方.
有它的策略性的地方.
你這樣看.
其實屬命上都是一樣.
上帝在這個地方成就舊人.
然後差派他的門徒.
去到地極.
有了地理.
就幫到我們了解很多.
這些聖經的章節.
\newpage



\section{11:19-26}
\label{sec:PjaM3O4lbt8}
\textbf{The beauty of ordinary faithfulness (Acts 11\_19-26) - Rev. Sam Johnson}
\newline
\newline
連結: \href{https://youtube.com/watch?v=PjaM3O4lbt8}{\texttt{https://youtube.com/watch?v=PjaM3O4lbt8}} ~~~~ 語音日期: 2022-08-30
\newline
\newline
\hyperref[sec:dVqF8V7wdDo]{\small{< < < PREV SERMON < < <}}
~
\hyperref[sec:index_chronic]{\small{[返順時目]}}
~
\hyperref[sec:index_scriptual]{\small{[返順卷目]}}
~
\hyperref[sec:f_Inn31K7Tc]{\small{> > > NEXT SERMON > > >}}
\newline
\newline
$^{1}$Well, thank you, James, and thank you for leading us in worship..
James and Annie really enjoyed that..
And I'm excited to be here, always excited to be at Urban Voice.
and really enjoy being here with you all..
Glad to be back again this morning..
We're going to be looking in Acts chapter 11..
So if you want to turn there.
in your Bible or in your Bible app, Acts chapter 11,.
we're going to be looking at verses 19 to 26 and looking at a special group.
of believers in this passage and in our world..
The modern world that we live in, a lot of times what we see.
is that our world really values and celebrates and craves.
things that make a lot of noise..
We we get excited about recognition and praise and fame and celebrity..
Big, big accomplishments..
That's what's important to us..
And our culture really values that we value chatter and things that go viral..
Those are the things that we like..
And a lot of times we value that even over content..
And this kind of milieu, it's less important.
what is being done or said and more important how many people hear it or see it..
We kind of end up prizing noise over quality a lot of times..
And sometimes even the most foolish or ridiculous moments.
can be praised and seem important or special.
just if enough people see them or notice them if they get enough views..
And it's not just noise..
Our world also assigns value based upon quantity..
We look at things that are small, someone who's doing something small.
with small numbers, and we think that that is not very important..
We might think it's inferior or even useless at some times..
We look at something big that's being done with big numbers,.
and we think that that's a big deal..
And so a lot of these cultural voices are constantly telling us.
to only do big things, to go big or go home..
We're we're told that we should only ever dream big dreams.
and we should only give our time to something if it's going to be noticed..
And sometimes we feel like if we're not getting noticed,.
if people are not seeing what we're doing,.
if we're not being recognized, if we feel unseen,.
then we feel worthless..

$^{41}$We feel like we're failing. We feel like we are losing..
We feel like if we're not accomplishing something big.
and earth-shaking and revolutionary, then we're getting nowhere..
And if no one knows our name or no one sees us, then we have no worth..
And culture often tells us these things explicitly and implicitly..
And these things give us the impression that.
living an ordinary life is at best, you know, kind of.
failing to be what you should be or could be,.
and at worst, maybe just foolish and not even a good idea..
But the problem here is that we lose sight of the value.
of the ordinary, of ordinary life and of ordinary obedience to Jesus..
A lot of times we've forgotten how much we need lives of everyday usefulness,.
and we're ignorant of how profoundly these lives really weave together.
the fabric of our society and our world and support everything that our world contains..
And sadly, a lot of times this is just as true in the church.
as it is in culture outside of the church..
In the church, a lot of times we've lost our vision.
of the beauty of ordinary faith and the beauty of ordinary life,.
the beauty of ordinary obedience to Jesus..
And this scripture is important because it helps to refocus our perspective..
The reality is that the Bible consistently throughout the scripture,.
it consistently affirms the value of an ordinary life, of obedience to God.
and how essential that is in the kingdom of God..
And it elevates this to a place of prominence and says that it's important.
and it's valued by God..
It's honoring to God himself..
Some of the biggest things that have happened in the kingdom of God.
have always taken place through the lives of unseen, anonymous,.
ordinary followers of Jesus..
And that's what this passage points us to..
So let's read these verses together..
We're going to read Acts 11, verses 19 to 26..
Now, those who were scattered because of the persecution that arose over Stephen.
traveled as far as Phoenicia and Cyprus and Antioch,.
speaking the word to no one except Jews..
But there were some of them, men of Cyprus and Cyrene,.
who on coming to Antioch spoke to the Hellenists also,.
preaching the Lord Jesus and the hand of the Lord was with them..
And a great number who believe turned to the Lord..
The report of this came to the ears of the church in Jerusalem,.

$^{81}$and they sent Barnabas to Antioch..
When he came and saw the grace of God, he was glad and he exhorted them all.
to remain faithful to the Lord with steadfast purpose,.
for he was a good man, full of the Holy Spirit and of faith..
And a great many people were added to the Lord..
So Barnabas went to Tarshish to look for Saul..
And when he had found him, he brought him to Antioch for a whole year..
They met with the church and taught a great many people..
And in Antioch, the disciples were first called Christians..
Now, this passage is about some very ordinary believers..
In fact, they're so ordinary that we are told basically nothing about them..
There's no names. There's no genealogies..
There's no educational background, no social information about these believers,.
no positions of rank within their culture or their society..
They're they're simply referred to as some of them.
with that them being followers of Jesus, saying some of these followers of Jesus..
Yet their nameless testimony speaks so powerfully of God's work.
in their lives and through their lives, God's work in their community..
God does something huge and powerful in Antioch through this group of nameless.
believers. And their testimony speaks to us powerfully.
about the beauty of lives lived in ordinary, faithful obedience to Jesus..
So as we consider this passage, first, I want us to think about three points.
of what ordinary faithfulness is, talk about what it is..
And then we're going to think about four points of what it does in our life.
when we live it out, try to answer those two questions..
What is ordinary faithfulness to Jesus and what does it do in our lives?.
So first, let's think about what it is..
What is ordinary faithfulness?.
The first thing I want us to see is that ordinary faithfulness.
is quiet devotion to Jesus..
It's it's devotion to Jesus and our relationship with him..
Now, what is devotion to Jesus? What does that mean?.
What does that look like?.
Well, devotion is the manifestation in our actions of the love.
and focus of our hearts..
For example, I love my wife.
who is here today with me,.
and I seek to manifest that love by acts of devotion..
So I should devote my time and my attention to her..
I should serve her. I should listen to her..

$^{121}$I should sacrifice for her. I should prioritize her..
I should prioritize our relationship by making time for her..
I should commit to value what she values..
The things that are important to her should become important to me..
I should love those that she loves and care for what she cares for..
I should even build my schedule and structure.
the rhythms of my life around her and our relationship..
If I'm doing that, then it'll be clear to her that I'm devoted to her..
It'll be clear to other people that I'm devoted to her..
So that's just kind of a simple example.
of what devotion looks like that we all should be able to grab onto and understand..
The same is true of our relationship to Jesus..
Devotion to Jesus looks like giving him the best of our time and attention..
It looks like spending frequent, consistent, undistracted time.
connecting with Jesus through God's word in prayer and in worship..
It looks like prioritizing our relationship with Jesus over other things in our lives,.
listening to his voice, serving him with our whole heart,.
being willing to sacrifice and put away things in our life.
that are going to pull us away from him, pull us into sin.
and seeking to commit ourselves to value what he values..
The things that are important to Jesus should be important to us..
We should love the things that lift him up in our lives and draw us closer to him..
And we should stay away from stuff that takes us away from Christ..
And we should build our schedules and the rhythms of our life.
around connection to God, making sure that we have time to connect with God..
That's devotion to Jesus..
And the reason I call this ordinary faithfulness quiet.
devotion is not because it's silent or noiseless..
There are times when our devotion may be full of loud worship,.
but I call it quiet because it should be done.
not for recognition, not to be seen by others..
It's not done to be noticed..
It's done out of love for Jesus..
It's not performative faithfulness to make myself feel better.
or to gain acknowledgement from other human beings or even to earn favor from God..
It is simply a faithful obedience.
and devotion that is motivated by love for Jesus..
That's pretty easy to understand, but not always easy to do..
It requires God's grace at work in my heart, requires the activity,.
the Holy Spirit in my life, and it requires commitment day in and day out..

$^{161}$It requires ordinary faithfulness..
The next thing that ordinary faithfulness is.
is quiet obedience to God's commands..
Ordinary faithfulness is obeying the things that God has told us to do.
and seeking to live out the character that God has instructed us to live out..
I'll give us another simple example..
Think about the fruits of the spirit in Galatians five..
You know, the list here, love, joy, peace, patience, kindness, goodness,.
faithfulness, generosity, self-control..
Think about that list now and ask yourself, do we live out those character traits?.
Are those the types of things that characterize our relationships at home.
with our family or at work with our coworkers, our boss, in our community.
with with the people that live around us or in our church?.
Are those things that characterize our relationships?.
That's something that God has commanded us to do, to live out those values..
Speaking of church, are we obeying God's commands to be part of church?.
Are we engaged with the community of the people of God?.
Are we involved in a way that is honest and transparent in that community?.
Are we opening ourselves up to relationship with other believers?.
Are we loving and depending on that community?.
These are things that God has instructed us to do..
Are we loving God with all of our heart?.
Are we putting nothing else before him?.
Are we loving those around us as much as we love ourselves?.
Obviously, there are a lot more to God's commands than that..
But but that's a great place to start with just some of the simple things.
that we know God has told us to do..
That's what it means to obey God's commands..
Are we?.
To love God with all of our hearts, to love our spouses, our children,.
our parents, our brothers and sisters, our neighbors, our family in the church,.
to give, to serve, to uplift, to care for, to be peacemakers, to love truth,.
to stay away from sin..
These are things that God has commanded us..
And scripture is full of instruction about the kind of lives.
that God wants us to live..
And we should be seeking to live that way, to obey God, God's commands..
And again, I call this obedience quiet, not because it's silence at silent..
At times it may be bold, but it's quiet because it's obedience..
That's motivated by love for Jesus..

$^{201}$It's empowered by God's work in our lives..
It's not about us or recognition..
It's not something that we serve to be seen..
It's not for an audience other than Jesus..
And the honest truth is it's not even possible to live out that life of obedience.
to God's commands if God's not already at work in your heart,.
if you don't already have relationship with him..
Because once again, it's easy to understand but hard to do..
It requires God's grace to love others..
It requires God's work to help us forgive..
It requires the activity of the Holy Spirit..
It requires commitment day in and day out and faithfulness..
The other thing I want us to see about what ordinary faithfulness is,.
ordinary faithfulness is the proclamation of God's message..
The things we've talked about so far about ordinary faithfulness are really.
attitudes and actions..
But it's more than attitudes and actions, there's also a verbal component that we.
have a story to tell..
And yes, our actions speak and they're essential to that story that our actions.
honor God, but we also have words to tell..
The scripture is clear that we should pair our actions with words that describe.
and explain and communicate the good news of the gospel, the good news of hope,.
forgiveness, grace, and salvation through Christ..
And we proclaim God's message as we share the grand redemptive narrative of the.
Bible, God's loving creation of the universe, the relationship between God and.
humans that's broken because of human sin, then God's loving rescue and.
restoration that God sends through his only son, Jesus, the free gift of.
salvation that comes by God's grace and through faith in Jesus, God's ongoing work.
of loving restoration and recreation that he's doing in our world, in the universe,.
through the coming of his kingdom..
That's the story that we have to tell..
We can tell that grand, the big picture story, the grand narrative of redemption..
But we also get to proclaim this message as we share the redemptive narrative of.
God's work in our own lives, God's loving creation of me and the broken relationship.
between God and me because of my sin, God's rescue of me through Jesus,.
the ongoing work of restoration and recreation that God is doing in my life.
and in my heart and changing me..
We all have a personal version of this narrative to tell..
And as we tell these redemptive stories, whether it's the grand narrative or our.
personal narrative, even if we're just telling a little bit, even if we're just.

$^{241}$able to share a small piece of that story of what God's done in our life or a small.
piece of God's love for someone else, and even if it's just a snippet at a time,.
as we tell those to the people around us, we obey God by proclaiming the message.
he's given us..
And that's what we see in this passage..
These believers, these ordinary faithful believers, they were just talking to people.
and telling the story of God's love and redemption..
So those are some of the things that ordinary faithfulness is in our lives..
And I want us to think now about some of the things that ordinary faithfulness does..
At first, we see that ordinary faithfulness points to the beauty of the gospel..
Now, the followers of Jesus that are referred to in this passage,.
as we talked about earlier, they're ordinary, they're normal people,.
they're not wealthy or upper-class people..
And we can assume this because verse 19 tells us they came to Antioch because they.
were fleeing the persecution in Jerusalem..
And that's the persecution that started when Stephen was martyred and the people.
stoned him and Saul, who would later become Paul, was watching over their coats..
That's story in Acts 8.3..
So the people that are fleeing this persecution, if you think about any time.
of social upheaval or crisis, the refugees are usually the ones from the lowest rungs.
of the economic ladder, right?.
People with more wealth and status, they have more ways to kind of stay safe and.
to insulate themselves..
But the people who don't are usually the poor people, the people who don't have.
as much..
And so they're the ones who are forced to flee..
So that's probably who these people in Antioch were..
They were not the aristocrats from the church at Jerusalem that had to flee.
to Antioch..
They were just normal families..
And we know they were families because in Acts 8.3, it talks about the fact that Saul,.
when the persecution started, he was going house to house and dragging out men.
and women..
So we know these were men, women, and families, just normal people who loved.
Jesus, who were living their lives..
And because of this persecution, now they have to flee to another city,.
to Antioch..
That's the picture we have here..
They're not wealthy or powerful..
They're not influential..

$^{281}$They don't have status..
And yet, God uses them in an amazingly powerful way..
This passage tells us that God uses them to bring a great many people into the.
family of God..
It tells us that the hand of the Lord was with them..
God uses them to do something amazing and powerful..
So what's so amazing about these people?.
Were they amazing communicators?.
Were they really compelling in the way they shared the gospel?.
What was it about these people?.
Well, we know ultimately the answer to that is that God was at work,.
the Holy Spirit was at work in their lives..
This was the work of God..
But I think we can also see in this passage that a part of what was amazing.
about these people was they were living lives of faithful and beautiful and,.
yes, ordinary devotion to Jesus that aligned perfectly with the words that they.
were proclaiming..
They weren't known for being important or great or viral..
They weren't known for being influencers..
They were just known for living these God-honoring lives of faithfulness and.
beauty and obedience to God..
And we see this testimony reinforced throughout the history of the church..
For every name that you might think of that's great in the history of the church,.
for every Martin Luther or Billy Graham or whatever name you might think of,.
there are hundreds of thousands of ordinary faithful believers who lived their.
lives, loved their families, loved their spouses, loved their children,.
served their neighbors..
They helped the sick and helped the poor, and they told people the story of Jesus..
And many of them even suffered in ways that we can't even really wrap our minds.
around, but we don't know their names..
But their lives of ordinary faithfulness glorified God and brought people to God.
and brought people into the kingdom..
That's the history of the church..
That's the history of the church right now..
The modern moment of the church, all of the growth and all the exciting things.
that God is doing are in places where you don't know the names of those people..
There's not pastors in villages in China or in Africa or in the Middle East that.
most of us would know their names..
They're not famous..
And yet, they're living lives of faithful obedience to God, and God is using them to.

$^{321}$do amazing and powerful things..
The goodness of the character of God that we read about throughout scripture was.
reflected in the lives of these followers of Jesus..
Things like reconciliation, forgiveness, peace, the fruits of the spirit that we.
talked about a moment ago, the shalom and wholeness of living out the values of the.
kingdom of God, the manifestation of true agape love, something that our world is so.
hungry for, the wonder of being full of grace but also being full of truth and.
speaking the truth, obeying God's commands, living a life of righteousness.
before him, all these things are what define these ordinary believers..
Their lives pointed to the beauty of the gospel..
That's what ordinary faithfulness does, and that's just as true for us..
This can still happen when we live a life of ordinary faithfulness..
Our ordinary, everyday, faithful obedience to God points to the beauty of the gospel..
It strips away all the noise of ourselves and our ambition, our pride,.
our branding, our mess, and it brings to the forefront the simple,.
beautiful things that make a life of following Jesus an incredible experience.
of both growth through difficult trials and also joyful rest in his strength and.
who he is..
It points to the wonder of the good news of God's grace, the good news of the.
gospel, the good news of the love of God..
Ordinary faithfulness points to the beauty of the gospel..
Ordinary faithfulness also puts Jesus at the center..
This is one of the coolest things about this passage to me..
In verse 26, we see that the disciples were first called Christians in this.
Antioch church, and this statement is a little bit astonishing in itself..
Why were they not called that before?.
But it tells us the defining characteristic of this movement taking place in Antioch.
was Jesus Christ..
Isn't that cool?.
Jesus was the thing that outsiders noticed and remembered and gave credit to..
That's because this work was a movement of God through ordinary believers..
So there was no celebrity preacher they could name this movement after..
There was no wealthy donor that was driving things that they could name this.
movement after..
There was no local landmark that they wanted to name it after..
There wasn't an affiliation to describe this group, no river or park or anything.
else..
This wasn't called the Olive Tree Mission or the Broad Street Gathering..
The only description that outsiders could think of to describe what was going on was.
Jesus..

$^{361}$I think that's so cool and so powerful..
And I think it's a challenge to us..
What if our local churches left that kind of impression?.
Jesus was at the center of this church and he was at the center of this movement..
And when people on the outside looked at it, Jesus was the only thing that they saw.
and God received the glory..
And that's why they began to call these people Christians because they were named.
them after Jesus..
This is what ordinary faithfulness does..
It puts Jesus at the center..
It lifts him up..
And here's a reality that's hard..
It's hard for me..
If we are living a life of ordinary faithfulness, a lot of times we're going to feel small..
We're going to feel unseen..
But we have to learn to embrace this reality rather than fight against it..
Because when we are made small, then God is made big in people's eyes..
When we are obscured, then God is made obvious..
And that's actually what we want..
What a beautiful way to glorify God and put Jesus at the center by being faithful in those.
small, ordinary, obscure ways..
And that's what it did in this passage..
Next thing I want us to see is that ordinary faithfulness brings us purpose and peace and.
freedom..
It brings something beautiful to us..
Learning to embrace ordinary faithfulness gives us a healthy sense of purpose in our.
daily lives..
You know, a lot of us, or at least me, maybe I'm just talking to myself here, but a lot.
of us awake each day both overwhelmed and underwhelmed with purpose..
We're overwhelmed because we want to accomplish so much..
We want to do such great things..
We're pressured and crushed by the weight of all the stuff we have to do and how long.
our to-do list is..
We know we want to make great strides and we want to have a lasting impact and we want.
to be impossibly productive and get so much stuff done..
And, you know, we want to make astounding profits in whatever business endeavors we.
have..
We want to give back to the community..
And of course, we want to leave a legacy for others..
There's so much that we want to do that it's just overwhelming..

$^{401}$But at the same time, the relentless demands of our calendars also leave us underwhelmed.
because even though we have these lofty ambitions, our days are full of mundane tasks..
And we have to do all these mundane things in our families, at work, around the house,.
in our communities..
So we're underwhelmed because we feel like we spend so much time doing things that don't.
matter or maybe they're even useless..
And we wonder if we're wasting our time or sometimes we wonder if we're wasting our lives..
And living a life of ordinary faithfulness speaks to both of those emotional quagmires.
and gives us truth..
First of all, living a life of ordinary faithfulness to God relieves the pressure of being overwhelmed.
by questions of purpose..
We don't have to live our day with the weight of all that we must accomplish..
We don't feel like we need to change the world and leave our mark in just 24 hours..
We can live the day with a simple understanding of what is significant and important..
Be devoted to Jesus today..
Obey God's commands today..
Live out his truth..
Share his message..
It's that simple..
We can know that our purpose is actually modest and clear for today and for this week and.
for this year..
And further, we can learn to rest in the truth that with the power of God at work in our.
life that purpose is actually attainable..
By God's grace, we can actually be faithful in the small things that he's given us to.
do..
And this day has purpose and it has value..
And I don't have to feel rushed to accomplish everything at once because there is time for.
what's truly important..
There is time for ordinary faithfulness to Jesus..
And second, living a life of ordinary faithfulness imbues even the smallest aspects of my life.
with deep purpose..
I'm reminded here of this brother Lawrence who was this medieval Christian monk and he.
wrote this classic called Practicing the Presence of God..
And in it, he says that when he was doing his work in the kitchen, he had to work for.
decades in the kitchen peeling potatoes and doing these mundane things..
It says that he having accustomed himself to do everything there for the love of God.
and with prayer upon all occasions for his grace to do his work well, he found everything.
easy during the 15 years that he had been employed there..
So this man spent decades of his life peeling potatoes and doing menial kitchen tasks..
But he found a way to see this work as honoring and pleasing to God and he found a way to.

$^{441}$connect with God while he did it..
So he can peel potatoes with purpose and significance, then surely we can be faithful at our workplaces.
or in our homes and see it as honoring to God..
We don't have to be underwhelmed by mundane tasks..
We can find purpose in glorifying God even as we vacuum the house or finish up those.
reports for our boss or run errands or help a neighbor install a new thermostat..
When we're preparing a meal or helping a classmate or cleaning a messy floor or even changing.
diapers, we can do those things as acts of devotion to Jesus because we love him and.
want to love others who he loves..
We can do these things as acts of obedience because God commanded us to serve others and.
give to others and care for others..
We can find purpose in even the smallest tasks in our lives..
And embracing ordinary faithfulness can also give us peace..
You know, a lot of us grew up on Disney movies and so we heard a lot of things like if you.
can dream it, you can do it or there's no limit to what you can accomplish if you put.
your mind to it and work hard..
Nothing wrong with those things, but sometimes those things, even though they were meant.
to boost our confidence and inspire us, they built into an expectation that we started.
to feel like we were actually supposed to change the world and change everything..
Maybe we feel like we're supposed to start something really big, some big company that.
makes a lot of money or that changes everybody's lives or we feel like we should, you know,.
save the world, clean all the carbon out of the atmosphere and save all the animals everywhere..
We feel like we've got to do so much..
And even in the church, we feel this way..
We feel like we need to embrace all the causes and support all the needs and right all the.
wrongs and sometimes we're frustrated because we haven't made more progress..
Why haven't we fed all the hunger yet or whatever it is that we want to do?.
Those are good things and it's good to have big dreams and big ideas, but sometimes that.
becomes a pressure that begins to crush us..
And if we're honest, that unspoken expectation can become very heavy, so heavy that we can't.
live up to it..
And when we figure out we can't live up to it, a lot of times we're discouraged..
It can damage our faith and cause us to despair..
It can even cause us to give up on trying to make a difference or even give up on our.
faith and our trust in God..
The reality is we were never supposed to change the world..
Only God can do that..
Only God can do the big things..
We can't save anyone..
We can't even change our own hearts..

$^{481}$God is God..
Jesus is the Savior of the world and we are not..
And there's actually freedom in knowing that..
Freedom as we get out from under the crushing burden of pressure and expectation that we're.
supposed to fix things, freedom as we understand that we're not responsible to do something.
that we never could have done in the first place..
I can't do it..
There's freedom in saying those words and in understanding that, that God created us.
with limitations and that's okay..
But simply acknowledging our inadequacy isn't enough because it might lead us to despair.
if we don't also at the same time see God's power..
It's not just enough to see our weakness, we need to see God's strength..
And when we do that, we find peace..
Peace as we see that God can change us, he can change our world and he is at work to.
do both..
Peace as we realize that by God's grace we can be used by him to take part in his global.
restoration project that is his kingdom and is accomplished through his church, through.
us..
Yes, only God can do big things, but he can use his people doing small things to do a.
big thing..
He can use us..
Our actions matter and there's so much peace and freedom in knowing that..
And then we see that ordinary faithfulness is fruitful beyond our imagination..
What fruit comes out of this passage?.
It's more fruitful and healthy and impactful than anyone could have imagined..
First the church at Antioch becomes the first multi-ethnic church..
Before this the church was an assembly of Jewish believers..
In Acts 13.1 we read the people who are part of the leadership of this church and it's.
a list of people from North Africa, from all over the Mediterranean, very different diverse.
backgrounds and diverse people groups..
This church becomes the first multi-ethnic church..
It's a major turning point in history, in Christian history to go from a Jewish background.
movement to a multi-ethnic background movement..
God does something amazing through this church..
But also this church becomes the main missionary sending church of the New Testament..
If you go on through the book of Acts you see in Acts 13.2 and 3 the church at Antioch.
commissions Paul and Barnabas, prays over them and sends them out on the first missionary.
and church planning journey..
In Acts 14.26 to 28 Paul and Barnabas come back to the Antioch church to give a report.
of what God did on this journey..

$^{521}$In Acts 15.40 Paul and Silas are sent out again and commended by the Antioch church..
In Acts 18.22 Paul once again returns from his church planning and strengthening journey.
and comes back to the church at Antioch..
And these passages show that to Paul and to the Antioch believers themselves this church.
shared in the fruit of the missionary journeys..
God used this church to be a part of planning churches in Philippi, Corinth, Thessalonica,.
Ephesus, Galatia..
All these things we read about came out of the lives of these ordinary nameless unknown.
believers..
They did something small in obedience to God and God used it to do something really, really.
big..
So how do we respond to this passage?.
I think this passage is both a challenge and an encouragement to us..
For some of us it's an encourage, it's a challenge maybe to reorient some of our priorities.
towards lives of devotion to Christ and obedience to God's commands and sharing his message.
and doing those simple and ordinary things that maybe we've been too busy to do or we've.
been too distracted or we've been focused on other things..
Maybe we need to reorient towards those things..
But for all of us I hope this passage is an encouragement, it's an encouragement as we.
deal with the day-to-day tasks of our lives and we find ourselves confronted with the.
reality that our reach is shorter than our grasp, we're not able to accomplish all that.
we want to and we feel small and insignificant, we feel unseen in those moments, we're unsure.
if we're having lasting impact, sometimes we're afraid that we're failing or we're doing.
something futile..
This passage assures us that if we're living a life of obedience to God we are not unseen..
Our days are not futile..
By God's grace we can accomplish more than we know and our contributions even though.
they may seem small, they're used by God to do something big in our world..
God is powerful and great and good and he is at work and so if we'll live this life.
of ordinary faithfulness to him he'll use it in a powerful way..
I want to close just by encouraging us to remember there's nothing wrong with dreaming.
big dreams or wanting to do something big for God or do something big in your life,.
that's fine and God might have that for you but here's the reality, living for Jesus is.
certainly a lot more than ordinary faithfulness but it is never less..
No matter how much God has for you in your life it will never be more than that moment.
of devotion to him..
Then it will never be more than that obedience to his commands and that sharing his message.
with other..
It may turn into something really big, a lot more than ordinary faithfulness but it'll.
never be less than that..

$^{561}$Let's pray..
God thank you for the truth of your word, thank you for this passage and these believers.
and what a testimony they are to us..
We don't know any of their names or their histories or their stories but you use them.
in such a powerful way..
They were ordinary faithful people that just honored you with their lives and God I pray.
that that would be true of us, I pray it would be true of me..
That I would be devoted to you, that I would obey your commands, that I would share your.
message with others and that you would get all the glory for everything that you do in.
my heart and life..
In Jesus name..
\newpage



\section{}
\label{sec:f_Inn31K7Tc}
\textbf{上行之詩 (詩 122) - 陳崇基牧師}
\newline
\newline
連結: \href{https://youtube.com/watch?v=f-Inn31K7Tc}{\texttt{https://youtube.com/watch?v=f-Inn31K7Tc}} ~~~~ 語音日期: 2022-08-30
\newline
\newline
\hyperref[sec:PjaM3O4lbt8]{\small{< < < PREV SERMON < < <}}
~
\hyperref[sec:index_chronic]{\small{[返順時目]}}
~
\hyperref[sec:index_scriptual]{\small{[返順卷目]}}
~
\hyperref[sec:FkGSJO7nqgY]{\small{> > > NEXT SERMON > > >}}
\newline
\newline
$^{1}$真的非常好.
OK.
今天我想總結詩篇的系列.
最好我們看看上行之詩.
Song of Essence.
所以我們此前看過不同的詩篇.
看過不同的上行之詩.
我們想這次總結一下.
上行之詩是什麼意思呢?.
15首詩歌從詩篇120篇到134.
15首.
可能是世世代代舊約神的子民.
他們去到worship去到敬拜.
因為他們的律法裡面說.
每一年要上去耶路撒冷.
無論你住哪一個地方.
都要三次.
第一次是Passover的時候.
我們的Good Friday Easter的時候.
第二次就是七七四日.
九日之後.
就是五十日.
五秦節.
就是他們的修國節.
然後就到秋季.
就是到駐彭節的時候.
大概是我們西曆九月尾.
十月的時候.
你看到這三個節日裡面.
其實包括了整年.
他們的農耕的.
他們的所謂的Fiscal Year.
三次的.
不單止自己要在家裡.
要去紀念去敬拜.
而且三次也要上到耶路撒冷.
去敬拜上帝.
所以因為這樣的緣故.
可能就是一直累積一個傳統.
當他們上行.

$^{41}$去到耶路撒冷的時候.
他們就需要唱一些詩歌.
做一些祈禱.
預備他們的心.
所以由120篇至到134篇.
可能就是15首的詩歌.
他們上行的時候.
他們會唱.
或者你也可能是.
他們去到耶路撒冷.
就有拉比文學告訴我們.
上去聖殿的時候.
進了耶路撒冷的City Limit.
然後去到Temple.
進去Temple的時候.
就有15級樓梯.
1960年代.
1970年代的時候.
那些以色列的人.
就去到法國那個地方.
在聖殿山.
就是今天的清真寺下面.
那一面.
在法國的時候.
就發現有15級的樓梯.
每一級的樓梯都是雙倍的.
兩級兩級.
但是有15級.
所以今天的猶人.
很多時候就去到那個地方.
在聖殿山的南面.
就好像當時耶穌的時候.
猶太人第一世紀的時候.
他們上到去耶路撒冷朝聖的時候.
真的進到聖殿的門口.
所謂進到Court的門口.
就在那15級裡面.
可能就是那些利未人的師班等等.
他們就在那裡唱這15首詩歌.
每一級就唱一首.

$^{81}$就是一個禮儀.
所以可能在一個建築上面.
和一個敬拜的禮儀上.
一種合併.
所以在上行之詩.
似乎就是特意編排.
以致讀上行之詩.
用上行之詩.
由120篇到祈禱.
到唱這些詩歌的時候.
就仿似你心裡面那種旅程.
一種spiritual.
一個in your spirit的一種journey.
就是繼續去耶路撒冷.
去到神的面前去敬拜.
所以看幾節聖經.
我們由120篇開始.
你就會感覺到120篇開始.
那個詩人在哪裡呢.
第五節這裡說.
禍災.
祈禱那種哀歌.
他說.
慘了這節.
我寄居在米切.
住在基達帳棚之中.
我與那恨惡和平.
不喜歡沙窿的人.
許久同住.
我願和平.
當我發言的時候.
他們卻要爭戰.
你看到米切和基達.
無論在哪裡都好.
米切大概是現在黑海南邊.
就是現在的土耳其北邊的地方.
基達大概是今天的敘利亞.
沙地阿拉伯北邊的沙漠.
這些地方.
無論什麼地方都好.

$^{121}$你很肯定知道.
那個地方遠離耶路撒冷.
它在基達的帳棚.
在米切這個地方.
是遠離耶路撒冷的.
他就覺得那裡很慘.
因為那裡的人不會說上帝.
不會看上帝長眼.
他們不要和平.
要爭戰.
這個詩人.
一個以色列人.
或者之後的猶太人.
總之是神的子民.
他有些原因.
他就是在exile(逃離)裡面.
在流放的裡面.
不是在自己的家鄉.
他真正家鄉所屬的.
就是耶路撒冷.
就是堅實的面.
但他有些原因.
他就是在米切.
在這些所謂流放的.
這種情況.
流散的情況.
如果你這樣看的時候.
這些詩篇.
其實跟我們有關.
因為我們今天.
無論在裡面.
或者在外面.
我們都總覺得有一種流散的感覺.
流亡的感覺.
我記得剛開始來美國的時候.
那時候已經是90年代.
頭幾年.
我們在教會團體.
大家聊天的時候.
有些來得久了的.

$^{161}$來美國的弟兄.
剛剛就說了一些國家大事.
他們說的時候.
就是My country.
我的國家.
那時候剛來到美國的時候.
你不是很感覺到.
為什麼你會說是我的國家.
因為你總覺得.
你在異鄉.
你來到這裡.
就覺得.
見到不同膚色的人.
可能你會說他是外國人.
但其實你是在外國.
你才是外國人.
但你明明心理的張力.
你在異鄉.
流落的時候.
你總是不覺得那是你的家鄉.
那不是我的家鄉.
無論你在裡面的景況.
那種心境也好.
Soulscape.
或者外面的位置.
那個Landscape.
都是有一種的Diaspora.
Exile.
就是流亡的感覺.
流亡就是說還沒回到家.
如果你來了美國很久.
你就會說My country.
你已經落地生根了.
但如果你還沒來很久的話.
你還是會覺得.
不上不下.
到底你屬於哪裡呢.
有一種流散的感覺.
上行之詩.
其實就是寫給流散的人.

$^{201}$流散的人.
他們希望能夠回到家.
去見上帝的面.
夜路上行.
所以第一首120首.
就是告訴我們.
這個詩人在異鄉的時候.
他有什麼感覺.
平安.
而平安是從哪裡來的呢.
繼續看下去你就會知道.
121.
第二首的上行之詩.
標題看清楚.
上行之詩.
聖經真的有這個標題.
第二首.
我要向山舉目.
我的幫助從何而來.
我的幫助從造天地的耶和華而來.
你會感覺到.
詩人開始走近.
家裡.
開始所謂回家.
走近夜路上行.
我要向山舉目.
我的幫助從何而來.
他望見群山.
山脈的時候.
他就很嚮往要回家.
為什麼你說快要回家.
因為夜路上行是在海拔2000多呎.
的群山裡面.
所以你去夜路上行.
一定會見到山脈.
中山.
山脈.
他見到中山.
他就向那裡舉目.
因為他知道.

$^{241}$耶和華在那裡.
耶和華是創造天地.
並且在那裡能夠賜力量給他.
所以他很希望.
唱到121篇的時候.
他就離開他自己流散的異鄉.
慢慢慢慢他向夜路上行邁進.
他就見到群山的時候.
他就很嚮往能夠回到家.
122號.
22號 第三首.
第三首.
又是上行之詩.
大衛上行之詩.
是不是大衛寫的.
總之就是Song of Ascent 第三首.
我喜樂因為人.
對我說.
我們到耶和華的殿去.
耶路撒冷.
我們的腳.
站在你的門內.
耶路撒冷被建造.
與和連絡.
整齊的一座城.
單看這幾節.
你就知道他已經到達了.
我很開心.
因為人對我們說.
Let's go to the temple.
一起去耶和華的殿.
快要到的時候.
Let's go there.
耶路撒冷我們的腳.
你看到他之前望群山.
但現在他望見的是自己的腳趾.
自己的腳.
站在城門之內.
通過不同的城門就進入了耶路撒冷.
第三節.

$^{281}$耶路撒冷被建造.
與和連絡整齊的一座城.
意思就是說.
連絡整齊.
我想很客氣.
這樣譯法.
其實是很擠迫的意思.
因為以前舊的城市.
是很擠迫的 街道很窄.
水洩不通.
很擠迫.
大家會撞到肩膀.
走進去的時候.
是很擠迫的.
很緊密的.
耶路撒冷就是一座這樣的城市.
一座這樣的市鎮.
進去的時候.
他就終於來到.
所以斯人.
由外地.
異地.
他就到了目的地.
站在耶路撒冷的城門之內.
然後我們看1,2,3.
到1,3,4最後.
就是跟著那十二首詩歌.
就是在目的地裡面.
做的事情.
在聖殿裡面去敬拜上帝.
去到石安.
Zion.
就是耶路撒冷.
他真正得到平安.
上帝給他的力量.
最後第134篇.
最後一首第15首.
第15首.
上行之詩.
也是title.

$^{321}$來.
稱頌耶和華.
夜間樹立在.
耶和華殿中.
耶和華的僕人.
那些的師班啊 祭司等等.
當向聖所.
舉手.
稱頌耶和華.
願做天地的耶和華從石安.
賜福給你.
有什麼feel呢.
完成了行程.
去了耶路撒冷.
然後到最後第15首.
第134首.
就是那些祭司跟他說.
願做天地的耶和華.
從石安耶路撒冷賜福給你.
他要離開了.
要上機 要走.
完成了整個旅程.
15首詩歌.
你就發現是一個所謂pilgrimage.
一個所謂朝聖.
一個屬靈的旅程.
一個實體的旅程.
一個完完整整的旅程.
上行之詩.
告訴神的子民.
無論身居哪個地方.
其實我們都是on the journey.
我們的life應該是on the journey.
上行意思就是不會settle down.
去旅行.
你有沒有去過旅行.
有啊 一定有.
covid之後你很想去.
你會找一個目的地.
可能中間有些終站.

$^{361}$但那個終站不會是你安居的地方.
如果你安居.
I think you defeat the purpose.
有沒有人去旅行.
你目的地在哪 我不知道.
我去到哪 那個就是目的地.
那就很廢了.
你總要有個destination.
你不可以用一個return to determine.
你做的事.
你最後當然有個final destination.
但你要去哪一處.
那些終站你是不會settle down.
今天我講這個主題是挺好的.
因為你看到我一身的打扮.
就知道我已經準備好了.
我平時都不會那麼casual.
當然平時都很casual.
我都不會打領帶等等.
但今天更加casual.
因為完了之後我們就上機.
你知道收拾東西.
旅行收拾東西.
我相信其他隊員都.
很多東西收拾.
我昨晚都收拾到.
差不多一半.
我昨晚都收拾到.
一點多兩點.
原來那麼多東西收拾.
想著想著.
我又那麼差 收拾那麼久.
又怕忘記帶這個.
又怕忘記帶那個.
今早起床我再check check.
其實最主要都是兩樣東西.
就是什麼.
passport.
還有呢.
credit card.

$^{401}$最重要有身份.
最重要有政治身份.
最重要還有錢.
其實就是這兩樣東西.
所以收拾的時候.
我覺得有些東西都不用帶.
不用帶那麼多.
你上行的時候.
你就是希望.
能夠輕便.
因為你永遠不會小孩.
我不知道有沒有想過這樣的picture.
其實就是一個基督徒的人生.
in fact其實我們本身的人生.
本身人生就應該是這樣.
基督徒更甚.
你想想.
在舊約新約.
剛才說食物.
整個食物鏈都說出來了.
從頭到尾.
舊約新約.
你不要說食物.
說旅程.
故事經文有多少.
幾乎處處都是.
第一本書.
創世記.
亞伯拉罕.
離開吳爾.
離開他的家鄉.
因為上帝呼召他你要走.
你住得很舒服不過你要走.
因為有個mission給你.
他就去了加拿大.
是一個上行.
一個journey.
不是去settle down.
不是去到落根.
永遠沒有落根.

$^{441}$你看那些族長的故事.
亞伯拉罕爾薩雅.
but not at home.
回家了.
但好像還是沒有回家.
你有沒有這種感覺.
回到自己的地方.
家裡.
城市.
但這個城市總是.
不會令你滿足.
總是有一樣東西.
是差一點.
the destination 仍然是.
yet to come.
目的地仍然未到.
亞伯拉罕爾薩雅.
族長的故事.
出埃及機.
更加是.
以色列人住了多少年都好.
四百年圍爐之地.
在埃及.
上帝都要叫他們走.
他們都很不捨得.
在那裡很舒服.
樣樣都習慣了.
上帝說你要走.
摩西就帶了以色列人離開埃及.
然後去曠野.
曠野四十年就是.
其實一直上行.
雖然他們比較陰公.
不聽話所以四十年在那裡繞一繞.
到了進入了迦南地.
好像是settle down.
開始建立國度.
索羅大衛所羅門等等.
但那個國度不會長久.
分裂成為南北國.

$^{481}$然後北國亡.
然後南國亡.
他們又是on the journey.
他們流放到Exile to Babylon.
去了一個不是屬於.
他們應該留的地方.
在那裡大約七十年.
然後上帝就說是時候了.
你們可以回去.
他們又去回歸.
你們這樣看的時候.
整個舊約的聖經.
整個舊約的進程.
都是一個journey.
因為神的指紋.
做人.
本身就是一個journey.
本身我們就是要.
往一個.
目的地去行.
如果你忘記了目的地.
你就很多時候會搞錯.
其他東西.
想清楚你的終結.
你的final destination.
上行之師.
就是說.
上帝的心.
上帝的生命.
就是我們每一個要走向.
的旅程.
生命就是一個journey.
無論你今天年輕.
或者年長.
都是we are on the go.
on a journey.
搞清楚你們要去哪.
最近我們夫婦.
幫一個年輕人.
十幾歲.

$^{521}$初中生.
十幾歲.
土生土長.
想幫他.
最近.
情緒很困擾.
覺得自己有depression.
覺得.
他突然意識到.
近這一年.
突然一個十幾歲的年輕人.
意識到.
生命很短暫.
最終我們會離開.
we will pass away.
於是他想起.
他想起這件事的時候.
他就覺得很困擾他.
想起人生的短暫.
生命要面對死亡.
走不掉.
這個問題.
我不知道是不是疫情.
COVID.
令到人開始想這件事.
或者令到人.
不能避免.
以往可能生活好好的.
世界好像.
很有秩序.
可能不需要掛問這些.
但是可能因為這個原因.
他就問多一些.
十幾歲.
十幾歲通常是想什麼.
十幾歲通常是開開心心.
玩.
花父母的錢.
買你喜歡買的東西.
去你想去的地方.

$^{561}$十幾歲通常都是冒險者.
那些會take risk.
所以為什麼.
十幾二十歲.
開車的時候.
因為撞車的是他們.
不代表年老的不開快車.
年少的開快車的機會.
一定是大過.
年長的.
還有他們魯莽的程度會多一些.
不代表年長不魯莽.
我們都可能會魯莽.
但是年輕的都會開快車.
會瘋了.
三個字就是不怕死.
他們怎會想死.
十幾歲的年輕人.
十幾歲的年輕人很有趣.
很有趣.
我們text幫他.
開解他.
我相信.
這個十幾歲的年輕人.
他心智是.
很成熟.
很mature.
十幾歲.
他就會想.
Where is the final destination.
如果final destination.
即是死.
Oh my god.
那我怎麼辦.
你不想這件事的時候.
你怎會evaluate.
去評估一下.
你現在走的方向.
你是為了什麼而去生存.
為什麼我說他成熟.

$^{601}$因為成熟.
可能比所有年長的人成熟.
因為我們最終想到的.
通常是什麼.
退休.
我們未求的就是想退休.
不用工作.
然後就不斷工作.
誰知道到我們退休的年齡.
老了身體又不好.
我們想退休去哪裡玩.
哪裡吃.
誰知道又走不動.
有錢又吃不到那些東西.
吃了那些東西.
吃了之後就要去ICU.
你又不夠膽吃.
我們想的終點.
我們想的終點就是什麼.
retirement.
我們說retirement真的好.
你都沒有想過retirement.
it's a risk.
你沒聽過很多retired的人.
沒工作.
they get depressed.
因為價值都沒有了.
我還生存在這裡幹什麼.
這個年輕人有什麼maturity.
就是他way beyond.
你retirement.
他不會說我們六十幾歲的時候.
六十幾歲我們要退休.
SSN又不夠錢.
等等.
他不是想那些東西.
他去想reality.
final destination.
最後的目的地.
你想清楚的時候.

$^{641}$你才可以評估到底我現在是為了什麼.
for what.
I'm on the journey.
回到剛才的比喻.
你去旅行的時候.
你不會沒有final destination.
到哪裡去就是目的地.
這個就很廢.
人生其實都是.
想清楚你的final destination.
我去開打他的時候.
他不回教會.
我們就認識了一間教會.
我們就.
在第二個state.
我們就介紹了他.
去那間教會.
那裡有個youth group.
有個青年的團體.
他很乖.
開始有一點.
我們為他祈禱.
希望他能夠.
be enlightened.
想清楚.
你試想想基督教信仰是很relevant的東西.
因為死你死定了.
如果那個死定了.
是你的final destination.
這樣就很可憐.
這樣就很廢.
上行之師.
on the journey.
去找什麼.
誰是我們的final destination.
聖經.
告訴我們人生最終的目的地.
是上帝的心.
是.
in the presence of God.

$^{681}$在上帝的同在裡.
舊約我們剛才說了.
新約也有.
這樣的概念.
新約我們選了幾段經文.
彼得前書第二章.
十個字.
他說親愛的弟兄姊妹.
你們是.
什麼.
不是根深蒂固.
落地生根的人.
你們是客類.
是寄居.
you are alien.
不是外星人.
大家那個是綠卡.
叫你alien.
不屬於這個地方.
不過被你拘留.
是要寄居的.
彼得說什麼.
我勸你們要禁戒肉體的情慾.
這情慾是和你.
神的心靈.
神進入了你的心靈.
你的靈魂不斷的增淺.
你們活在外邦人中間.
要品行斷情.
要讓那些人.
雖然毀謗你們是作惡的.
會因看見你們的好行為.
而在監察的日子.
就是上帝來審判的日子.
歸榮耀彼實.
看不看到你的身份.
你的身份.
不是美國的公民.
不是中國的公民.
不是任何一個地上國度的公民.

$^{721}$這不是你的身份.
你最重要的身份.
是外星人在地球上.
你去某個地方.
你是去某個地方的.
客類寄居的.
輕便的.
外星人在地球上.
是去歸的.
輕便的在旅行的.
去最後的目的地.
那就是最終的那一天.
你要去見上帝的那一天.
彼得前書很有趣.
第二章在這裡.
到第五章最後那幾節.
他說什麼.
彼得寫話.
他說我簡單寫了這封信.
寫了五章.
勸你們.
又證明這恩是真實的恩典.
我們務要在這恩典上.
站立得住.
不要跌低站立得住.
在巴比倫與你同盟.
揀選的教會.
也都向你們問安.
願平安.
歸於你們所有在基督裡面的人.
在巴比倫.
為什麼是巴比倫.
彼得肯定不是巴比倫.
因為大概當時的時候.
主後六十幾年的時候.
就是彼得殉道之前.
在羅姆.
怎麼會寫了巴比倫.
這個其實是一個很長久的.
猶太人的傳統.

$^{761}$因為在舊約裡面最出名巴比倫的就是.
將猶太人.
去到流放.
exile他們.
去迫害他們.
所以對彼得來說.
對當時第一世紀初代的基督徒來說.
巴比倫就是一個的code name.
就是一種政權在迫害他們.
那個其實就是.
不是真的巴比倫.
是羅馬.
羅馬帝國.
在迫害教會.
或者用不同的方法去引誘教會.
或者去到打散教會.
巴比倫就是一個代號.
是猶太人會這樣說的.
每逢是一些敵對教會.
敵對上帝國度的.
那些就是巴比倫.
所以你看到作者彼得.
是在流亡裡面.
他還沒有回家.
哪裡有他家.
當然不是巴比倫.
耶路撒冷.
真正的耶路撒冷就是基督徒.
你想想新約.
除了這裡還有很多.
其實你翻開新約第一本書.
你就知道.
Christian.
信主的人.
信耶穌的人.
其實就是on the journey.
你記得耶穌對門徒說的第一句話嗎.
跟隨我.
follow me.
當他在教會裡召門徒的時候.

$^{801}$他說跟隨我.
follow me的時候.
他literally.
實體他真的說了.
you follow me.
behind my back.
你跟隨我.
always on the journey.
最後一段聖經.
希伯來書第十二章.
這裡又是說journey.
因為之前十一章.
所謂英雄榜.
一直都是說舊約裡面的.
所謂信心的英雄.
神的子民.
亞伯拉罕.
以撒等等.
就是說他們在找一個eternal city.
那個eternal city不是在地上.
而是在天上將來要來.
他們個個都沒有家.
真正的家在天上.
在上帝裡面.
你查一查第十一章.
偏偏沒有說的是什麼.
就是聖殿.
沒有說聖殿.
因為聖殿對猶太人來說.
其實就是.
很多時候令他們停下來.
他們以為有個building就好了.
有個這樣的物業就好了.
是多麼輝煌的東西.
我們就停下來.
我們就沒有上路.
十一章偏偏沒有說建聖殿.
猶太人最夢寐以求的就是聖殿.
就是他們能夠坐立.
能夠settle down.

$^{841}$偏偏第十一章.
希伯來書就沒有說.
就是說個個都在跑步.
你有沒有跑馬拉松.
在跑步.
你在去一個destination.
在衝線去到一個final outcome.
然後第十二章就說.
所以我們有這樣多的見證人.
好像在運動場上看著我們.
跑的時候如同雲彩圍繞著我們.
我們就應該做什麼.
要輕便一點.
卸下我們什麼.
一些纏累我們的罪.
忍耐的心.
去奔跑在我們前面的道路.
然後他就說耶穌基督.
耶穌基督也是這樣.
我們的主.
他也是仰望前面的終點.
為我們信心創始成中的耶穌.
擺在前面.
ahead of him.
他的起龍.
衝線那種的起龍.
輕看了現在有的羞辱.
甚至是忍受十字架的痛苦.
如今他去哪了.
他真的坐在上帝的同在裡.
坐在上帝的右邊.
這就是一個progress.
一個journey.
我不是很確定.
你信了主這麼久.
有沒有想過.
到底你是在去哪裡.
Where are you going.
你的destination.
你的終點.

$^{881}$你的目的地在哪裡.
你有沒有想過.
如果只是死硬兩個字的話.
我就覺得很廢了.
但如果way beyond.
是去到上帝的裡面.
是一個永恆的裡面.
你現在走的道路.
或者你做的選擇.
就應該是不同的.
你的家庭的選擇.
婚姻的選擇.
你對你兒女的選擇.
不只是她去哪間學校.
不只是我住在哪個地方比較好.
那些都重要.
但不是最重要的.
不只是我的身體要保健等等.
這都是中站.
中間的站.
不是最終的目的地.
你有沒有想過.
你的目的地在哪裡.
上行之師就告訴我.
We are on a pilgrim.
我們中間的一些年長的基督徒.
可能就聽過.
以前在17世紀的時候.
有一個人叫John Bunyan.
他寫了一本小說.
叫做Pilgrim's Progress.
Pilgrim's Progress是一本經典.
17世紀已經面世.
然後有不同的翻譯.
中文也有 英文也有.
叫做《天路歷程》.
整本小說講什麼.
就是講一個基督徒.
信主之後.
他往中站.

$^{921}$最終的目的地.
天國的旅程.
中間他預見的是什麼.
在中間的路途.
他在小說裡 故事裡.
就預見有很多物質的引誘.
可能有些同伴.
會引誘他停下在某個城.
就不要繼續上路.
或者會有些同伴.
可能跟他們開始的時候.
是同路一起走的.
誰知道最後他就掉下了.
然後就離開了.
到底你是一個主角.
那個基督徒.
你會怎麼做呢.
他就講他的struggle.
很多時候可能有懷疑.
或者前面有一些威嚇.
或者一些恐懼等等.
他有沒有停下來.
這個故事.
就是翻山越嶺等等.
到最後他就要過死亡關.
死亡關.
Value of death.
他講了什麼呢.
你看到Value of death.
你就是不想過.
所以有些人就不想.
停下來.
不去就行了.
或者走分岔路去別的地方.
不要想最終的終站.
或者終站beyond的final destination.
是beyond death的final destination.
面對Value of death的時候.
他講什麼呢.
Going back is nothing but death.

$^{961}$To go forward is fear of death.
But life everlasting beyond.
I will yet go forward.
OK 你回頭.
It's nothing but death.
You still die.
You still end.
But to go forward is fear of death.
But life everlasting.
回頭是死路一條.
但前面死亡的威嚇.
和跨越死亡的永生.
那個就是destination.
我仍要往前走.
但我們每一個都是.
想清想楚.
你人生.
特別你基督徒的生命.
Where.
Ultimately where are you going.
Right.
不要停下來.
繼續高呼.
可以同心低落.
主我們多謝你.
因為你是那位在我們旅程上.
給我們恩典力量的神.
還有我們每一個屬於基督的人.
我們都是on the journey.
主我求你remind我們.
提醒我們不要停下來.
也不要將一些次要的東西.
變成我們的終點.
在你的裡面.
在上帝的生命.
真正的你的同在.
是我們的終點.
主因為這個終點.
因為這個終末.
教我們如何評估我們現在有的選擇.

$^{1001}$我們做的決定.
主讓我們敬拜你.
去事奉你.
在我們這個journey裡面.
我們繼續跟隨你.
讓耶穌成為我們的主.
帶領著我們走.
讓我們最終來到你面前的時候.
我們知道我們就屬於你.
求你給我們恩典.
給我們力量.
給我們每一位.
叫我們上上落落.
高山低谷的時候.
有你跟我們同在.
帶領我們去最終.
你的榮耀裡面.
我們這樣祈禱.
奉主庇救.
\newpage



\section{3:16}
\label{sec:FkGSJO7nqgY}
\textbf{Hang Tight (John 3\_16) - Joyce Tam}
\newline
\newline
連結: \href{https://youtube.com/watch?v=FkGSJO7nqgY}{\texttt{https://youtube.com/watch?v=FkGSJO7nqgY}} ~~~~ 語音日期: 2022-09-05
\newline
\newline
\hyperref[sec:f_Inn31K7Tc]{\small{< < < PREV SERMON < < <}}
~
\hyperref[sec:index_chronic]{\small{[返順時目]}}
~
\hyperref[sec:index_scriptual]{\small{[返順卷目]}}
~
\hyperref[sec:ahiCSxfoOC8]{\small{> > > NEXT SERMON > > >}}
\newline
\newline
$^{1}$Hello, my name is Joyce. I am one of the ministers at UV. I live in Arizona and I mostly serve.
the church online. And I fly back to Chicago once a quarter to serve in person. In 2008,.
the U.S. economy started to decline. That summer, our family went to China to do a short-term.
mission with our home church. We returned to Chicago in August, and next month, September,.
my husband was laid off. Since then, he was out of work for the next 11 months. The year.
after, 2009, we went to China for another short-term mission trip. While we visit families.
and friends in Hong Kong, my husband got a job offer. So after he brought us back to.
Chicago, he flew back to Hong Kong for the new job. He only had 100 Hong Kong dollars.
in his pocket for that month until he got his first paycheck. My daughter and I stayed.
behind in Chicago for another three months because of her school. We have about a month.
of mortgage and living expenses in the bank. We live in Hong Kong for two years, and we've.
moved to Arizona in 2012. In September 2014, we got a job offer from a bank in Hong Kong.
and we came back to Chicago for a friend's wedding. A day before the wedding, our whole.
family went for a physical checkup because we were in the middle of changing insurance..
A day after the wedding, my husband was admitted to the hospital. Because the blood test showed.
his blood count was dangerously low, his hemoglobin level is about one-third of a normal person..
They did some tests on him, and the tests showed that he had five ulcers in his stomach..
He was bleeding nonstop. We stayed in the hospital for a blood transfusion, and after.
he was discharged, we went back to Arizona, and he needed another blood transfusion. Unfortunately,.
there was a month between the two insurances that we didn't have insurance, and it was.
that month. So the medical bills pretty much emptied our bank account. And with all the.
chaos, I barely survived my second semester of seminary study. I can go on and on with.
similar stories like this, but I think you get the idea. There were times in my life.
that was very, very difficult. But standing in front of you today, when I look back, I.
have nothing but thankfulness. I thank you, God, for His grace. I thank God for the many.
trials in my life. I am blessed because I'm a Christian. How so? How could that be? I.
thank you, God, for the many trials in my life. How so? How could that be? Today's message,.
hang tight. Hang tight means to continue doing something or staying somewhere, although it.
is difficult or unpleasant to do so. Let's pray. Lord God Almighty, blessed be the name.
of the Lord. We come before you and ask for a teachable heart. May the Holy Spirit speak.
to each one of us as you pleased. Your name be glorified. In Jesus' name we pray. Amen..
Today's passage, John 3, 16. Shall we read it together? "For God so loved the world that.
He gave His one and only Son, that whoever believes in Him shall not perish but have.
eternal life." Many of you, or probably all of you, are very familiar with this verse..
But maybe it is not a bad idea to take another look, or you just have no choice since I'm.
the speaker. I'm the one who chose the verse. As you see, this is not a long verse. However,.
it is so rich in meaning and significance. So let us unpack this verse together. In this.
verse, there are seven important truths, truths that you and I need. The first truth is that.
we need, number one, God, the ultimate authority. Two, love the world, the greatest action..
Three, that He gave His one and only, one and only Son, the best gift. Four, whoever.

$^{41}$believes in Him, the widest coverage. Five, believes in Him, the only way out. Six, shall.
not perish, the holiest salvation. Seven, eternal life, the priceless gift. God, the.
ultimate authority. From the very beginning, chapter 1 in Genesis, there are 11 God says..
The authority of God is clearly reflected in the command spoken from His mouth. It is.
also recorded in the Old Testament that God called Abraham with authority. In Exodus,.
God also played an authoritative role in the exile and return of the Israelites. When it.
comes to the New Testament, we also see that God had given authority to Jesus. "All authority.
in heaven and earth has been given to me." And later, authority was passed on to the.
disciples. More importantly, Paul mentioned in his letter that the authority from the.
Holy Spirit is fulfilled the mission in the church of God. God is the ultimate authority..
He has the authority over the earth, over the authority of your life, over the challenges.
in your life. Two, love the world, the greatest action. God loves the world. The word love.
consists of thoughts, wills, and emotions. But we must understand God's love for the.
world is not limited to thoughts, wills, and emotions. God's love is also an action. Christ.
died for us. In this, the love of God was made manifest among us, that God sent His.
only Son to this world so that we might live through Him. For reason why loving the world.
is the greatest act? Because this is a sacrificial love. On the cross, the Father makes the sinless.
Son the object of the Father's wrath for our sake that in Jesus we might become the righteousness.
of God. Someone took the action to pay for your wrong. Someone took action to bear the.
pain. Someone loves you more than He loves Himself. Three, one and only Son. The Father.
one and only Son, the best gift. Eternal life belongs to God. The Son of God also has.
eternal life. Eternal life is supernatural. So unless God gives it to man, there's no.
way for man to have it by himself. John 10, 28, "I give them eternal life and they will.
never perish and no one will snatch them out of my hand." The only begotten Son is the.
best, ultimate, and most priceless gift. Life, you will never receive another gift in your.
life that can compare to it. There were times in my life that were so difficult that I can.
barely breathe. There seems to be no way out and I don't even have money to pay the mortgage.
next month. I used to joke with my brother's sisters in the church like, "If I don't have.
money to buy grocery, I'll come to your home for dinner." But the gift of, the thought.
of, the knowing of Jesus carried me through. You may be a professional. You may feel like.
you have control over your life. You might feel confident or you may have tons of money,.
but in the end, you cannot save yourself. Four, whoever the widest coverage. You must.
know that the Jews did not think that God loved the world because they believed that.
God only loved Israel. But it was Jesus' word that showed God's love is so long, wide, and.
deep that it covers everyone in the world. All people in the world can be saved and all.
people can have the life in the Lord. Jesus' teaching challenged the Israelites to humble.
their hearts and their status as God's chosen people and deprive them of the privilege of.
thinking they had God's blessings. Jesus transformed people's understanding of God's love and.
salvation. We are covered. You and I are within the coverage..
Five, believe in Him, the only way out. Believe in Him, describe people who are willing to.

$^{81}$accept God's love. God loves the world, but if people don't accept the Lord Jesus and.
the gift from God, they can't have eternal life at all. Believing in Jesus is not just.
about cognition in the head, but believing in Jesus is about believing in the head and.
relying on the heart and following in the action..
John 5, 24, "Truly, truly, I say to you, whoever hears my word and believes Him who sent me.
has eternal life. He does not come into judgment, but has passed from death to life.".
To believe is to connect, to hang tight our lives in Jesus, to transfer our sovereignty,.
to change our values, and to allow God to renew our hearts and guide our lives from.
now on. You must believe in Him because He is the only way out..
There is a verse that I love since I was in high school. Proverbs 3, 5-6, "Trust in the.
Lord with all your heart and do not lean on your own understanding. In all your ways acknowledge.
Him and He will make straight your path.".
Trust in the Lord with all your heart and do not lean onto your own understanding. In.
all your ways acknowledge Him and He will make straight your path..
Sixth, "Shall not perish the holy salvation." The word perish, destroyed in the Bible means.
to be destroyed, to be regarded as useless, to be killed, to go to hell, to be lost, to.
be declared, sorry, to be put to death..
Charles Spurgeon, one of the most influential English preachers in the 19th century. He.
once said, "Morality may keep you out of jail, but it takes the blood of Jesus Christ to.
keep you out of hell.".
Elaine Pagels, she is a professor of religion at Princeton University, the world's leading.
biblical literature research expert. In 1987, her son Mark died of terminal illness at the.
age of six, followed 15 months later. Her husband died in a climbing accident..
In 2018, she released a book called "Why Religion?" This book was written after the death of her.
son and her husband. Elaine tried to explore the teachings in the Bible from the religious.
studies she was familiar with, which could help her face the death of her son and her.
husband and overcome the struggles of extreme pain..
Although Elaine studied the Bible all her life, she was not a Christian, and even denied.
Christianity and religion all her life. Ironically, at the end she had no choice but to admit.
that the grace from God brought her out of the shadows. Although Elaine Pagels experienced.
grace from God and the death of her son and her husband, in the end she still chose to.
define God as the object of her study, not her savior. Her knowledge of the Bible is.
so extremely rich, and she has received many, many awards and respect in the academic world,.
but none of these can save her..
How about you? Do you truly believe He is the God of your life, or is it just a brainful.
of Bible knowledge? "Shall not perish," the holy salvation, is the ultimate goal of God's.
love for the world. Anyone can know God, but unless they confess their sins and believe.
in their hearts, they cannot escape eternal destruction..
Eternal Life is the Priceless Gift.
John 11, 25-26 Jesus said to her, "I am the resurrection and the life. Whoever believes.

$^{121}$in me, though he die, yet shall he live. And everyone who lives and believes in me shall.
never die." Do you believe this?.
Eternal Life, the Priceless Gift. Have you received it? You claim it. Open it. The understanding.
of Eternal Life helps us live our lives beyond the here and now. It empowers us to live life.
with a wider kingdom's perspective. We can claim victory over comes the world..
How can we still sleep like a baby when you lose your job, when there's no money in your.
bank account, or the number in your bank account continues to drop? How can we finish the race.
when everything seems to go against you? And how can we testify his name in our lives when.
there is suffering, when there is continuous challenge, struggle?.
God, the ultimate authority, loves the world, the greatest action, one and only Son, the.
best gift, whoever the widest coverage, believes in him the only way out, shall not perish.
the holiest salvation. Eternal Life, the Priceless Gift..
Seven Truths. Seven truths of the gospel we Christians need to hang tight because eternal.
hope is a big deal..
In Pastor Mark's video, he mentioned Matthew chapter 6, 31 to 33. What is kingdom living.
like? What should disciples, followers of Jesus think of every single day?.
Seek ye first the kingdom of God. Understanding who we are is the foundation of this worthy.
walk. Let's not live like we don't know God. Act or live like a Gentile..
Anyone with kingdom perspective can conquer the storm in life. Trials shall pass. Christ's.
followers can rest peacefully in the storm because he saves, because God is with us..
Hang tight. Trust in the Lord. Pray without ceasing. Don't let go. Let's pray..
Redeeming God, thank you for salvation, full and free. You are rich in mercy and grace..
You have promise and eternity spent with you by faith in Jesus Christ. Nothing can separate.
us from you. Please grant us kingdom perspective when trials come. We don't worry about tomorrow,.
what we eat, what we wear, because there is something more important. May your name be.
glorified in our lives. In Jesus' name we pray. Amen..
[BLANK AUDIO].
\newpage



\section{何西阿書 5:15-6:6}
\label{sec:ahiCSxfoOC8}
\textbf{主愛又如何? (何西阿書 5\_15-6\_6) - 吳君才牧師}
\newline
\newline
連結: \href{https://youtube.com/watch?v=ahiCSxfoOC8}{\texttt{https://youtube.com/watch?v=ahiCSxfoOC8}} ~~~~ 語音日期: 2022-09-05
\newline
\newline
\hyperref[sec:FkGSJO7nqgY]{\small{< < < PREV SERMON < < <}}
~
\hyperref[sec:index_chronic]{\small{[返順時目]}}
~
\hyperref[sec:index_scriptual]{\small{[返順卷目]}}
~
\hyperref[sec:T9LnkydmX7E]{\small{> > > NEXT SERMON > > >}}
\newline
\newline
何西阿書 5:15-6:6
\newline
\begin{longtable}{cl}
\hline
\hline
章節 & 經文 (和合本修訂版)\\
\hline
5:15 & \begin{tabularx}{0.7\textwidth}{X} 我要去，我要回到原處，等他們自覺有罪，尋求我的面；急難時他們必切切尋求我。 \end{tabularx} \\ \\ \relax
6:1 & \begin{tabularx}{0.7\textwidth}{X} 來，我們歸向耶和華吧！他撕裂我們，也必醫治；打傷我們，也必包紮。 \end{tabularx} \\ \\ \relax
6:2 & \begin{tabularx}{0.7\textwidth}{X} 過兩天他必使我們甦醒，第三天他必使我們興起，我們就在他面前得以存活。 \end{tabularx} \\ \\ \relax
6:3 & \begin{tabularx}{0.7\textwidth}{X} 我們要認識，要追求認識耶和華。他如黎明必然出現，他必臨到我們像甘霖，像滋潤土地的春雨。 \end{tabularx} \\ \\ \relax
6:4 & \begin{tabularx}{0.7\textwidth}{X} 以法蓮哪，我可以向你怎樣行呢？猶大啊，我可以向你怎樣做呢？因為你們的慈愛如同早晨的雲霧，又如速散的露水。 \end{tabularx} \\ \\ \relax
6:5 & \begin{tabularx}{0.7\textwidth}{X} 因此，我藉先知砍伐他們，以我口中的話殺戮他們；對你的審判如光發出。 \end{tabularx} \\ \\ \relax
6:6 & \begin{tabularx}{0.7\textwidth}{X} 我喜愛慈愛，不喜愛祭物；喜愛人認識神，勝於燔祭。 \end{tabularx} \\ \\
[1ex]
\hline
\hline
\end{longtable}
$^{1}$是,再說一次.
平安不是必然的.
大約在三個月之前.
我年輕的侄子突然不舒服.
氣喘,胸悶.
結果去了急症室.
之後被確診患有.
罕有的癌症.
而且發現是末期的.
他的病情急起急轉直下.
七天之後.
不治離世.
最後要白頭人送黑頭人.
有些親友不期然會懷疑.
為何這些不幸的苦難.
會淋到信主的人家庭當中呢?.
主的愛是如何理解.
苦難與主的愛又有何關連呢?.
今早我願意重溫主的愛.
另類的表達方式.
對我們有何提醒.
我以經文.
是《河西亞書》第五章.
十五節到六章六節.
請打出經文.
我稍後會.
分段時會多說一些.
大家看一看.
用一分鐘左右.
我先簡介一下《河西亞書》.
《河西亞》先知是公元前.
七百二十一年.
撒瑪利亞陷落之前.
那段紛亂時期的先知.
他是繼先知阿摩斯之後.
在北角以色列傳上帝的信息.
他特別關心百姓拜偶像.
和不忠於上帝的行為.
他大膽地以自己.
跟一個不忠貞的婦人結婚.

$^{41}$給以色列人那種背信.
正如他妻子戈滅對他的不忠貞.
上帝的子民同樣不忠.
背棄了上帝.
因此上帝的懲罰臨到以色列.
不過到最後.
上帝仍然對他的子民.
顯出無比信實的愛.
他要贏回以色列.
跟自己恢復正常的關係.
正如上帝的心意.
上帝的愛.
在《河西亞書》第十一章第八節.
那幾句話.
可以反映出來.
我用現代中文譯本譯出來的意思就是.
以色列人.
我怎能拋棄你們呢.
我怎能離開你們.
我的心不忍這樣做.
因為我多麼愛你們.
上帝的這種愛.
真是那種信實不變的愛.
神的愛是愛人到底的愛.
剛才你所看到的經文.
我們今早會盡心思想一下.
一些事情.
就是.
神的愛是怎樣.
人對主的愛又應該怎樣回應呢.
用三個方面的問題.
大家思想.
我們先有個簡單的祈禱.
天父有時我們不明白你的愛.
但是你的愛是信實.
是慈愛不變.
求你藉著今早.
在《河西亞書》.
我們看到你的心意.
也更願意愛你愛人.

$^{81}$祈禱奉主耶穌基督的名.
阿們.
第一個我們要思想的.
就是苦難.
痛苦.
可以彰顯主的愛嗎?.
你看到經文.
那裡說.
我要回到原處.
讓他們自覺有罪.
尋求我眠.
他們在急懶的時候.
必切切尋求我.
原處是指天上.
整句的意思是表示.
父神離開他的百姓而去.
其實是屬於暫時的離開.
目的是要讓他們.
能夠認罪悔改.
然後恢復尋求他.
其實人往往是自我中心.
自以為是.
事實上憑自己的能力.
用自己的思想辦法去行事為人.
很多時候我們會偏行己路.
這個其實是人性的真實寫照.
但是當急難苦楚不幸的事.
相繼臨到的時候.
人的驕傲,自意,自信.
才迫著要放下.
或者慢慢消失.
在痛苦,問題,難處當中.
原來人才會軟化下來.
再想想.
思想一下信仰.
再歸於神那裡.
事實上父神往往是藉著.
在你人生當中的逆境,不幸,折磨.
或者難處來管教.
造就他的子民.

$^{121}$所以管教的目的不是破壞.
而是更加促使人要親近神更多.
更加謙卑下來.
回轉歸向父神.
因為你肯認罪.
因為你肯謙卑在神面前.
其實是表明了你是否定自己過往.
始憑自己的能力.
而進而認識到不能夠再專注自己了.
不可以再靠己意了.
這些是不對的.
然後歸到神那裡.
這種整個生活態度和方向的扭轉.
正是悔改.
真實的悔改是會帶來切切尋求主的心.
是以父神作為你的人生主宰.
並且以主的旨意作為生活的指引.
你的人生觀,價值觀.
是扭轉定位在神那裡.
記得我在紐約的牧會的時候.
有人介紹一個住在附近的.
一個末期癌症的病人來.
差不多有兩三年的時間.
我也接載和陪同這個人.
去接受電療,化療.
這位男士因為這樣的緣故.
藉著這段經歷去信了主.
他後來去分享見證.
自己說.
父神原來是藉著這個癌症.
這個不治之症的苦難,不幸.
使他思想人生.
使他思想到信耶穌的問題.
父神是透過這個癌症的煎熬.
使他折服,降服,歸順了主.
並且切切地尋求神.
他信主之後.
變得有盼望.
雖然他知道自己死亡的日子是接近.
但是他心中那種牢苦掛慮.

$^{161}$都懂得交託給神.
他不再像以前那樣消極頹廢.
他說這是他未信主之前所沒有的.
在我的牧養當中.
曾經看到不少的家庭.
或者夫妻的關係出現了問題.
但是當他們肯讓神介入的時候.
神幫助的時候.
很多人都有感而發地說.
他說要改變這個這麼困難的局面.
人真的不能夠做什麼.
唯有好好地依靠神.
真的.
弟兄姊妹.
我們可以看到.
今天很多的環境.
很多的人士.
似乎壓迫著我們.
尤其是在近三年這個疫情的日子當中.
有些人真的忍不住.
但是請你留意.
主的愛.
原來是往往藉著苦難痛苦而彰顯的.
在你的難處逆境當中.
苦痛當中.
原來神的大能的手.
可以化身成為奇妙的祝福.
今天我們要確實知道.
神是明白你和我的限度.
祂永遠不會容許任何困難或者苦難.
是超越過我們的力量或者才剛所能承擔的.
我們要謹記.
神是信實的.
必不叫你們試探過於所能受的.
因此弟兄姊妹.
各位朋友.
如果你今天仍然身處在苦難逆境.
飽受試煉的考驗和折磨.
我請你仍然好好依靠主.
抬頭仰望神.

$^{201}$連續地繼續鼓起勇氣.
因為我們的父神知道.
你可以承受多少壓力.
你的極限是多少.
我們改變用另一個角度來看.
我們思想到主的愛.
原來是透過今時今日你所遭遇的一切.
這些苦難,這些問題,這些不幸.
磨練你,造就你.
甚至管教你.
為了使你更親近主.
歸向神.
並且切切尋求你.
當你肯謙卑回轉.
向父神支取力量的時候.
你就會更加體會到.
神的愛,神的恩典.
原來是何等的奇妙.
祂從來沒有將試探或者考驗.
比神的力量,比主的愛更大.
有一篇短文.
叫做《主耶穌的保證》.
是英文來的.
譯了中文之後也很感人,很激勵人心.
不知道你們能否看到.
聽我細細讀出.
「你永不會遇到恐懼是主不能克服的」.
「你永不會面對敵人是主不能打敗的」.
「你永不會進入戰役是主不能勝利的」.
「你永不會有需要是主不能應付的」.
「你永不會遇見試探是主不能勝過的」.
「你永不會背負重擔是主不能承托的」.
「你永不會面對難題是主不能解決的」.
「你永不會有觸磅是主不能拆解的」.
「你永不會有時段是主不會照顧的」.
「你永不會有時刻是主不再拿捏的」.
不一樣,我們看到神無比的奇妙作為.
第二個問題我們要想.
主的愛有甚麼應許呢?.
在第六章,《學西夏書》第六章一到三節.

$^{241}$是這麼說的.
「奈巴,我們歸向耶和華」.
「祂撕裂我們也必醫治」.
「祂打傷我們也必前果」.
第二節,過兩天.
「祂必使我們蘇醒」.
第三天.
「祂必使我們興起」.
「我們就在祂面前得以存活」.
第三節.
「我們務要認識耶和華」.
「竭力追求認識祂」.
「祂出現,學如神光」.
「祂必臨到我們像金魚」.
「像滋潤田地的春雨」.
一到三節是以法蓮和猶大的「華」.
後幾節,第四到六節是耶和華的答覆.
主的愛原來是透過四項動作.
下一章是說的.
是撕裂,醫治,打傷,前果.
這四種方法來成就.
主的審判好像猛獸一樣損傷人.
但是神的公義當中仍然有慈愛的成分.
祂必使人恢復健康.
撕裂是野獸對掠物的一個行動.
用爪,用牙去損傷.
通常野獸先弄死小動物.
然後再撕開來吞吃.
但是我們耶和華的刑罰.
是不會將人置於死地的.
因為父神的目的不是毀壞,而是醫治.
父神打傷人,是不會使人受傷致命的.
而是在傷處前果.
使人得到完全的痊癒.
請注意,全能全知的神.
祂有奇妙有慈愛的方法.
因為若果不撕裂,不破壞.
就不可能有改造和創新.
當你看第二節的時候.
你看到內容.

$^{281}$這個破壞的工作相當之嚴重.
幾乎是瀕臨死亡,置於死地.
好像要完全除滅一樣.
根本沒有恢復的可能.
但是父神仍然有慈愛奇妙的工作.
因為祂使人「蘇醒」.
「蘇醒」是救活的意思.
甚至第三天使人「興起」.
得以存活,即是死不去.
古時的異教迷信觀念.
認為人死了之後.
第三天的靈魂才會脫離身體.
所以在三天之內.
如果不復活就會永遠死亡.
這種說法當然不是這本經文的內容.
但是大家可以理解.
三天,兩天的時間.
始終都是一個很短暫的時間.
換句話說,沒錯.
父神是藉著不幸困境難處去撕裂,搞破壞.
去重重的,又重重的打擊,打傷.
但是主的愛卻是奇妙的去醫治,去懲割.
使人得以重建,造就,康復,恢復過來.
仍然得以存活.
所以父神的撕裂,打傷.
始終都是兩,三天.
是屬於短暫.
最後人在父神面前仍然得以存活.
在神面前仍然蒙大福.
得主的眷由.
所以這裡說:我們務要認識耶和華.
竭力追求認識他.
甚麼叫認識主呢?.
認識主就是要敬愛他,敬畏他.
多年前在香港.
有一位弟兄從事工程.
生活富裕不錯.
他有回教會,間中有事奉.
不過信仰對他來說只是生活的一部分.
直到有一次他發生嚴重的車禍.

$^{321}$去了急症室之後轉送到ICU.
醫生跟家人說.
他有嚴重的受創傷.
他生存的機會是三成左右.
於是教會發動弟兄姊妹.
為這位弟兄祈禱.
祈禱得很狠.
祈禱的力量真的大.
「褔腎」念文奇蹟真的出現了.
他始終渡過危險時期.
而且慢慢地好轉.
三,兩個星期之後.
可以轉到普通的病況.
自此之後這位弟兄的生命有極大的改變.
他自己的人生方向也有扭轉過來.
因為他認為是「褔腎」的念文.
自己才可以撿回一條命.
所以他定義今後要為主而活.
生命為主所用.
出院之後經過調理,康復.
他申請進入神學院.
接受裝備.
幾年之後他畢業.
然後他成為傳道人.
專注地為主作供.
其實「褔腎」今天也會在你和我身上.
往往透過一些苦難.
透過一些人生的衝擊.
要甦醒你.
要救活你.
更加要使用你.
所以我們今天若經歷了很多苦難.
其實是幫助我們更認識神.
更敬愛主.
更能夠好好地事奉主.
主的愛臨到真的好像曙光.
好像春雨.
曙光,神光是何等的耀眼.
人人都可以看到.
是很真實可靠.

$^{361}$春雨滋潤大地.
更使人得著豐富.
以及昌盛繁榮.
不單是屬於物質方面.
亦是心靈的加力.
使人能夠被主的愛和恩典.
得著潤澤和滋養.
影子們,各位朋友.
在苦難當中.
在困境,失意,絕望當中.
在你認為我偷偷碰著黑的時段.
請你用屬靈的眼光去看.
父神原來是一而再.
透過這些來撕裂,破壞,打傷,打擊你的情況.
我們可以看到一樣東西.
就是神要你更盡心認識祂.
所以祂藉著這些東西.
是有奇妙的安排.
根本一切的苦難,一切的煩惱.
只不過是兩三天.
是短暫的,是自輕的.
相對於將來無比的榮耀.
是不是算不得什麼呢.
一位姐妹因為一次意外.
搞到之後要坐輪椅.
有一次她淚流滿面向神禱告.
主啊,如果我是一直健康.
不用坐輪椅.
我就可以替你做很多的工作.
父神出人意外的平安和想法.
神的意念,不同人的意念.
比有限的人超越更多.
祂向這位姐妹表達清晰地說.
有很多人肯為我忙碌工作和生活.
但是很少的人願意肯花時間來親近我.
做我親密的朋友.
聽姐妹,你對主在你人生片段當中的撕裂.
打傷,打擊.
有沒有體會到原來有父神慈愛的依治和果前呢.
為要使你和我更認識神.

$^{401}$為要使你和我成為祂親密的朋友.
建立一個親切的關係.
有一首古舊的詩歌.
很直接地帶出.
或者你也聽過.
歌詞是這樣說.
前要的是福祉.
今要的是主.
前要的是感覺.
今要主懷與.
前要的是恩賜.
今要賜恩主.
前我尋求依治.
今獨要主自己.
我們經驗到主的愛背後.
那份依治,包摘,前果.
加力量.
以及更認識主的應許.
第三個問題.
我們要思想.
人對主的愛.
會不會忠誠堅定呢.
就是《河西亞書》第六章四到六節.
我讀出來.
主說:伊法蓮啊.
我可向你怎樣行呢?.
猶大啊.
我可向你怎樣作呢?.
因你們的亮善.
如同早晨的雲霧.
猶如速散的甘露.
因此我直先知.
懇佛他們.
以我口中的話.
殺戮他們.
我施行的審判.
如光發出.
我喜愛亮善.
不喜愛濟滓.
喜愛認識神.

$^{441}$勝於凡濟.
在巴勒斯坦這些地區.
是很炎熱,乾燥.
雨水,露水很快就會被蒸發.
《河西亞》才知.
就用這個現象來比喻以色列的人.
他們的義和對神的愛是不能夠持久.
如早晨的雲霧.
和速散的甘露.
因為他們的悔改.
甚至對主的愛.
都是很膚淺和短暫.
難怪父神要差遣忠信的先知.
在歷世歷代進行懇佛和殺戮.
懇佛原文的意思.
有刁鑽,破碎,開鑿這樣的意思.
經過這些過程就會建造得更好.
殺戮是表示滅絕,除盡,淨盡,清除這樣的意思.
也是父神藉著先知.
和藉著人生苦難的遭遇來管教.
聖化屬於他的每一個子民.
另一個說法有人說.
懇佛殺戮其實是指神譴責的意思.
往往父神是透過先知.
用他的話語來表達他的心意.
神的話語真的很有功效.
沒錯,神口中的話是給人生命.
但是當你不信服.
你背叛,遺命的時候.
卻使人與神的關係疏離隔絕.
這正是神的施行審判如光如閃電所發出的果效.
關於歸根究底.
神的懇佛,殺戮甚至譴責.
並不是要除滅.
而是要施恩.
然後要祝福.
所以正如第六章第一節所說.
我們對神的愛的回應.
不要那麼膚淺,飄忽或短暫.
第六節提及到神喜愛良善.

$^{481}$良善的細節或作連述.
原意是包括人對神的愛.
有行為表現出來.
以及對神的忠心.
第六節所提及的.
正如先知撒母耳所說的.
聽命,誠於獻祭,順從,誠於供養的自由.
有類似的涵義.
向神獻祭,服侍神,事奉主.
如果你沒有內心那種誠實,良善,忠誠.
神決定不會接納.
真正認識神的人.
乃是全心全意歸向神.
敬愛神,敬畏神.
不是只有外表認識的事奉.
否則你的事奉,你的獻祭.
只是有口無心.
並沒有持久的信靠,信服.
這樣神會唾棄.
靜姐妹.
我們真的要深切掌握這個真理.
盡心,盡意,盡性,盡力地回應主的愛.
今天我們在這個疫情當中.
無論在線上線下.
我們為主,我們敬拜神,我們愛主.
是否有種盡心,盡意,盡性,盡力的心態.
還是好像很多人說.
我很疼愛神的,我有保留的.
甚至我們很分心的.
我們只是做自己喜歡做的事.
愛自己做的事.
我們是有選擇性的擺上.
未算是盡呢.
真實的盡心,盡性,盡意,盡力.
不是只有口或只有知.
因為你和我都知道.
神是看人內心.
神知道我們裡面是怎樣.
但願我們不是表面,膚淺或短暫.
而是經得起考驗.

$^{521}$有聽命,順從的意願.
有行動印認出來.
最近我看一本孫教士的傳記.
拍出來的那本書.
就是《山雨,富能人身傳》.
這本書看的時候.
我深深感受到什麼叫盡心,盡意,盡力,盡性的愛神,愛人的表現.
這個富能人是一個孫教士.
他在1906年的時候.
他在英國倫敦皇家學院畢業.
準備做工程師.
但他又有天賦的音樂才華.
他彈了一首傑出的鋼琴.
他20歲居然可以在倫敦舉行第一次鋼琴演奏會.
但神在他身上有特別的計劃和呼召.
兩年後.
他夢召要去到一個人不想去的地方.
就是中國.
他放棄了人們以為很有前途的工程師的工作.
很高薪厚祿的未來.
他去到一個人看來很赤柏的地方.
因為他打算前往中國西南.
不毛之地的雲南.
而是專門向一些少數民族.
你見到這班是綠蜀族的族人.
向他們傳呼音.
他很艱苦的做開荒佈道的工作.
很辛苦.
過著極其簡樸和辛酸的生活.
他的足跡遍及雲南西北各地.
差不多三十年的時間.
在中國邊陲之地和土人做朋友.
他又學他們的綠蜀語.
少數民族語言.
他又用他們的語音.
編撰一些綠蜀語的新約聖經.
其實他經歷不少的患難險阻.
多次曾經被人殺.
瀕臨死亡的邊緣.
但仍然堅持的事奉.

$^{561}$直到後期.
他因為腦弱疾.
在1938年死於雲南寶山.
享年52歲.
其實他很艱苦經營和傳道.
足足12年才有第一間綠蜀人的教會成立.
之後神奇妙的.
不少綠蜀人和其他當地的土人.
分別信了主而成立了自養的教會.
在中國西南的地方.
很多間聖靈工作.
直到1949年中共的政權之後.
教會的情況才有改變.
當我閱讀這本山語.
關於116年前的宣教士的傳記.
這個庫林人的生命和傳道表現.
令我自己有不少反省和學習.
原來一個真正愛神愛人的人.
是怎樣毀身.
怎樣捨己.
怎樣亡我.
怎樣具體地在日常生活的細節上.
表現出盡心盡意盡性盡力愛主愛人.
庫林人在他生命最後的日子.
他用了很多的時間來禁食祈禱.
他說當他重讀約翰·維斯里的書的時候.
那時候正是他離世之前一個月.
他的日記是這樣寫.
他說我常思想極少數人.
能夠像約翰·維斯里一樣.
藉著嚴格的自律.
自律是時下已不流行的.
是準備妥當.
終生為神所用.
其實不單止約翰·維斯里的生命.
甚至庫林人的一生.
確確實實是終生為神所用.
他一生是為主消盡法綱.
弟兄姊妹.
如果作為基督徒.

$^{601}$你和我都要想一想.
我們的自律的情況是怎樣.
我愛主愛人愛到什麼地步.
我願不願意終生為主所用.
求主幫助我們透過何西亞書.
在這個疫情底下.
在這個這麼不穩定的時勢裡面.
我們仍然能夠明白.
和自己調整生命有三方面.
第一.
體會到主的愛是藉著苦難痛苦而彰顯的.
所以我們能夠接受.
和從苦難痛苦當中成長.
以及經歷主的愛更多.
第二.
要深刻經歷主的愛.
原來是有醫治.
包摘.
前果.
以及加力量的應許.
第三.
要對主的愛.
更忠誠堅定.
能夠盡心盡義盡力盡聖的愛神愛人.
好好地在生命和生活上見證主榮耀神.
讓我們誠心低頭.
結束祈禱.
多謝主藉著先知何西亞的提醒.
讓我們體會到神下原來是往往藉著苦難痛苦.
來彰顯你的愛.
幫助我們透過苦難痛苦.
有成長經歷你的愛更多.
更幫助我們深刻經歷神你的愛.
原來有醫治包摘.
果前.
以及加力量給我們的應許.
所以我們憑你的恩典.
繼續活得精彩.
主下我們更需要回應你的愛.
就是能夠盡心盡義盡聖盡力的愛你愛人.

$^{641}$好好地在生命和生活上榮耀見證你.
祈禱奉主耶穌基督的名.
阿們.
\newpage



\section{20:20}
\label{sec:T9LnkydmX7E}
\textbf{What Does God Really Care? (Reflection on Hezekiah's Achievements - 2 Kings 20\_20)}
\newline
\newline
連結: \href{https://youtube.com/watch?v=T9LnkydmX7E}{\texttt{https://youtube.com/watch?v=T9LnkydmX7E}} ~~~~ 語音日期: 2022-09-13
\newline
\newline
\hyperref[sec:ahiCSxfoOC8]{\small{< < < PREV SERMON < < <}}
~
\hyperref[sec:index_chronic]{\small{[返順時目]}}
~
\hyperref[sec:index_scriptual]{\small{[返順卷目]}}
~
\hyperref[sec:aMelPey4yl8]{\small{> > > NEXT SERMON > > >}}
\newline
\newline
$^{1}$Wave hands to your friends, yeah..
(laughs).
Alright..
Well, one last announcement is the registration, I think..
Oh, okay, alright, we have our continuous daily devotion..
Make sure you subscribe to it,.
join our social media group,.
and read a small portion of the devotional material.
to help you walk in the way of God, okay?.
So the theme is very interesting..
It give us today our daily bread, okay?.
Eating is so important for many of us, right?.
Especially Asian and Chinese, right?.
We greet each other, "Have you eaten yet?".
(laughs).
"Have you eaten your lunch or your dinner?".
We're so concerned about our own feast, right?.
Our own eating, our own stomach..
But the Bible has a lot to say about eating, right?.
The Bible has a lot to say about our food..
And in fact, the food, the meals that we have,.
are often tests and temptations to all of us..
So I hope this daily devotional.
will help you and challenge you..
Today is the last day for registration.
to be the voting member..
I have explained to all of you,.
the system at UV is kind of a system of democracy..
Alright, so all the voting members registered to be active,.
and you can nominate each other,.
and then elect each other to be the new,.
to become the core team of the new term..
Our Angus congregation, although it's smaller,.
we are not small, right?.
We are not small, as I can see..
There are, I think five, right?.
Five nominees coming out of you guys, alright?.
Maybe 10 out of the 10 o'clock service,.
but five out of you guys..
This is so important because our pastoral staff.

$^{41}$are working alongside with our elected.
core team members to lead the church..
So for the next three years,.
we're gonna generate a new leadership team,.
and we're gonna meet on a monthly basis.
for one hour or so to discuss church business.
and church direction and a lot of things..
We come to the church meetings,.
whether it's worship service or fellowship,.
or later you go to have your feast tonight at Mears..
You know, you need to understand,.
sometimes we show up to eat, and then we'll be gone, right?.
But there are tons of things behind the scene..
People are working so hard to make sure.
the schedule is right, and make sure all the tools,.
all the elements are there in order for us to enjoy..
So you need to understand, church is not just.
what you can see, it's not just one hour.
or two hours a week..
There are 24/7 behind the scene,.
people are doing things to prepare stuff for us..
So the core team leadership is actually in charge.
of all, many of the ministry.
and all the works behind the scene..
So this is so important..
So register if you have not..
If you have registered, praise God, okay?.
If you have registered, after today,.
we will enter into the second phase of our election,.
is the phase of nomination..
So there will be three Sundays,.
October 9, 16, and 23..
We're going to generate a list of nomination..
So once you register, you will be able to nominate people.
from the list of qualified nominees, or potential nominees..
You're gonna nominate among each one of us..
Whoever got the highest number of nominations,.
will go into the third phase, which is election..
All right, election is on the Sundays.
of November 6, 13, and 20..

$^{81}$So for three Sundays, we're gonna cast our vote, all right,.
to the list of nominees..
And then after the three Sundays of election,.
November 6, 13, and 20, we're gonna generate.
a new group of qualified members,.
and the results will be out on 20th of November..
All right, got it?.
I like this system..
I like this system, because we always have room.
for new members, for new leaders, for new ideas..
We need all of you to serve together..
In fact, the Bible tells us we need to do.
our best to serve God, to use our time,.
to use our talent, and to use our treasure.
to serve God's kingdom, all right?.
So three Ts, your time, your talent, and your treasure..
In my years in the past few churches, including this one,.
I'm so amazed by how God has raised up so many people,.
brothers and sisters, to do the little work behind the scene.
and to make the gospel possible, make the church possible..
Each one of us is gifted in many different ways, all right?.
I want to remind you, the church is composed.
of different members..
All of us are different, totally different, right?.
You don't want a church that, you don't want a church,.
the church body that has only one guy,.
like everyone looks like Mark Chan..
That would be the worst church in the world..
Everyone is unique, everyone is different, right?.
So use our talents and discover our talents..
We started our worship service 1130 until now..
There were many people working behind the scenes..
Some were doing the PowerPoint,.
some were doing the OBS,.
some were doing preparation and the worship team,.
they practice beforehand and come and set up.
and lead us into the presence of God..
So there were a lot of things..
We can use our gifts and use our talents, all right?.
So these three phrases, please remember today, September 11,.

$^{121}$you're going to register,.
this is the last Sunday to register, all right?.
All right?.
All right, so next one..
Oh, courting position, yes, there you go..
All right, so there will be a fellowship dinner today, right?.
Praise God, praise God, praise God for good food, okay?.
But I hope we have good time of fellowship as well, all right?.
It's not just about eating, it's about being together, right?.
And then what you do as giving, remember my three Ts?.
The third T is treasure..
You are not just giving for something unknown,.
you're giving to the ministry..
We are trying to build up our ministry,.
that our ministry is not just here and in the neighborhood,.
it's not just that..
Of course, this is our first priority..
As a community church,.
we wanna serve the people in our neighborhood,.
but our ministry goes to, I don't know, globally, right?.
Some of our YouTube videos, some of our teaching.
will go to somewhere else in the world.
and bless other people..
So as you give, remember this,.
you are participating in the global work of Jesus Christ..
All right?.
There was something I remember, I wanna remind you..
I forgot, yeah..
Shouldn't be important, but anyways, yeah..
So continue to give and give faithfully, all right?.
So today is my substitute preaching.
and let me give you a five minute recap.
of what I did in the past two weeks..
And then I give you the highlight.
or my highlight of the whole trip..
We went to different sites and different places in Israel..
If you have seen some of the videos,.
you should, first of all, you should remember,.
we often have some sort of preconception.
of what the land of the Bible, Israel, is, right?.

$^{161}$Now, if you ask your neighbors,.
if you ask your neighbors,.
what do you think of the land of Israel, right?.
What kind of image, what kind of impression.
come into your mind when people mention Israel, Palestine?.
Okay, talk to your neighbor, go to your neighbor..
What are some of the first impressions that you had.
with the land of Israel?.
What are some of those?.
Yeah, what were some of those that came to your mind?.
What was that?.
Here, here, this corner..
What you guys talk about?.
(laughs).
Dusty, dusty and sandy, right?.
Yeah, right, right, right..
What else?.
How about that?.
Huh?.
War, yes, war, right?.
Yeah, it's always a missile and bombs and all that, right?.
That's all you got, right?.
That's all you got..
What else?.
I'm sorry?.
Promised, promised land..
Promised land, did I get it right?.
Okay, it's good, good..
That is correct..
(laughs).
Well, is it sandy, is it dusty?.
Some areas are dusty, right?.
Just like Chicago..
Some areas are really dusty..
It was very interesting, I meet some friends.
or people, strangers on the street,.
and sometimes we exchange conversation, right?.
And then they ask, "Where you from?".
And I ask, "Where you from?".
"Oh, you from Sweden or from Germany.".

$^{201}$And then I say, "I'm from Chicago.".
And their impression of Chicago,.
it's very dangerous place..
It's probably the most dangerous place in the world..
(laughs).
That's true, right?.
In fact, yesterday, if you heard about the neighborhood news,.
I got, I think it happened yesterday or Friday..
Yeah, it happened Friday..
There was shooting in the Kenwood High School neighborhood.
in Hyde Park..
You heard about it?.
Because you are numb..
Because you are numb,.
because there's so many different shootings, right?.
You just don't care..
I don't care..
It's not special anymore..
And in fact, Friday, there was a shooting.
near Hyde Park, Kenwood Academy in Hyde Park..
And I talked to some of the parents,.
who's children are students of Kenwood..
And they told me there was a soft lockdown in Kenwood..
Kind of scary, right?.
It's the beginning weeks of the school year..
And already it's happening, right?.
But you talk to some of the foreigners,.
some of the people who've never been to Chicago,.
and then you try to correct their preconception, right?.
You try to correct them..
And I don't come from,.
there were not bullets flying 24/7 here, right?.
You try to tell them, you know,.
I'm from Chicago, I'm normal..
My street, my neighborhood is actually pretty safe, right?.
I live in McKinney Park,.
which is right next to Bridgeport, where we're at now..
I try to correct them, you know, it's not what you think..
Now, why did they think the way they do?.
Because of the media, right?.

$^{241}$Because of the media..
Now, you turn on your TV news,.
you turn on any news channel,.
you always get the sensational stories, right?.
Because what is newsworthy.
is always the dangerous stuff, right?.
The exciting stuff, the sensational stuff..
You don't get any news from me, right?.
Am I newsworthy?.
No, I'm not, right?.
I'm just a normal guy, ordinary, just like you, right?.
I don't read your story on, you know, CNN.
and all the major news channels..
All we got from the news channel is always a,.
well, one of the perspectives, right?.
One of the blind spots..
What is newsworthy is often sensational..
That's why when they hear about Chicago,.
they always think of the gangsters,.
the shootings, right, the killing,.
and all that stuff..
So I try to correct that..
And I go to the land of Israel for the past two weeks,.
and we had that kind of reorientation too..
Now, of course, I've been there for many, many times,.
and I don't need that..
But once we arrive, once we arrive,.
all of the Team Israel people,.
they tell me, "Wow, this is Israel?.
"Is that right?".
I said, "Yeah, you are in Israel..
"Check your GPS..
"You are in Israel.".
(audience laughing).
"You are not in Sahara Desert..
"You are not in some, you know, no man's land..
"It's actually pretty urbanized, right?".
I keep telling them, actually,.
Israel is a mini United States, right?.
So we land Israel on Monday, two Mondays ago,.

$^{281}$after 13 hours of flight..
So we land in the Ben Gurion Airport in Tel Aviv,.
or near Tel Aviv, and then we check out our rental cars..
So I drive one car and the other guy, Ronald,.
drive another car..
So we head out to the north side of Israel, to Galilee,.
right?.
It's a beautiful place..
It's the freshwater lake,.
the Sea of Galilee, Lake of Galilee,.
is a freshwater lake at the lowest point on planet Earth..
It's a beautiful place, right?.
Besides the heat wave, right?.
It's a beautiful place..
And they took a lot of pictures..
They have fun..
And then we check into our Airbnb..
We stayed there for less than a week, four to five days..
And then we drove all the way down to the middle.
of the country, which is Jerusalem..
Now, once we arrived in Jerusalem, it was a Thursday..
It was a Thursday..
Thursday night, we arrived and checked into our second Airbnb.
and then we settled in..
And then after an hour, I told the team,.
"Hey, let's meet on ground floor..
Because we have two different Airbnbs.".
I said, "After an hour or so, you settle down.
and we're gonna meet on the ground level.
and then we're gonna walk the street, okay?".
'Cause we checked in the Airbnb,.
and which is located right at the heart of the city,.
which is Ben Yehuda Street..
Ben Yehuda Street..
If you Google that, you'll find out Ben Yehuda Street.
is one of the busiest streets in the city of Jerusalem,.
in the new part of the city of Jerusalem, not the old city..
Anyways, it's very close to the old city..
So we went to the ground level and I said,.
"Hey, it's about time to go for dinner.".

$^{321}$Right?.
So they walked the street.
and there's light rail and all that..
And then shops and restaurants are all open..
And Thursday, I mentioned Thursday, we arrived.
because Thursday, Thursday evening,.
is the busiest day of the week..
Busiest day of the week..
Why?.
Because Friday is the Shabbat..
So it's the day before Shabbat, right?.
It's like our Friday night..
So their Thursday night is a very busy night..
And everyone is, you know,.
all the restaurants have tables and chairs out there.
on the sidewalk..
Good food, you know, good drinks,.
and people are chatting and having a good time..
And one of the team member told me,.
"This feels like Europe.".
"This feels like...".
Is that right?.
Europe is like that, right?.
You just chill out in the evening, nice and cool..
I said, "Yeah, exactly.".
Right?.
And then we went to the Western Wall, the Wailing Wall,.
and to check out, you know, all the activities there..
And it was crazy, so busy..
Because some of the Jews,.
especially the religious observant Jews,.
they would go to the Western Wall almost every day.
to pray and to read scripture..
So we had a good time..
We went to different site,.
and to so-called retrace the steps of Jesus,.
the steps of the Old Testament stories.
or New Testament stories..
Every place that we go is a storied place,.
storied place..

$^{361}$It means place full of stories,.
place full of traditions, histories, and legends..
Almost every site..
You know, I used to study in Israel for quite a bit,.
and I stay in the Old City..
So I go out to the Old City almost every day,.
just to walk and to get grocery or stuff like that..
Every stone has a story..
Every corner of the street has a story..
Has a story that goes back to thousands of years ago,.
go back to the Old Testament time..
And you can relate to what has happened in the past.
to what is relevant to all of us today..
And today, I would like to share one of the stories.
or one of the sites that I went to,.
which is one of the highlights, all right?.
Now let me read the scripture first..
2 Kings 20, 20..
20, 20, 2 Kings..
Now, this is the conclusion of the reign of King Hezekiah,.
one of the Old Testament kings,.
one of the kings reigning in the southern country of Judah,.
right around Jerusalem..
So this is the conclusion, the concluding note..
So after chapters of information,.
history about the reign of Hezekiah,.
the writer of the Bible says,.
as for other events of Hezekiah's reign,.
all his achievements and how he made the pool,.
the water, the water pool,.
and the tunnel, the water tunnel,.
by which he brought water into the city,.
are they not written in the book of Kings of Judah?.
Now, the last question that he asked.
was a rhetorical question..
But what I mean by rhetorical question.
is that the answer is expected, of course, right?.
There were other events,.
many other events and achievements of King Hezekiah..
Are they not written in the book of Kings of Judah?.

$^{401}$The archive, the government files,.
you know, the, you know, how many gigabytes,.
you know, in today's term, are all stored.
in the government archive of the day..
And that archive is called the archive of Kings of Judah..
Right?.
So there are many other things.
that the biblical author can talk about,.
about King of, about King Hezekiah..
And he said, I can talk so much,.
but they are all written in the government archive..
And I didn't talk about it, right?.
So in the previous chapters,.
he talked about what King Hezekiah has done in his reign..
Some of the highlights in the history.
of the reign of Hezekiah..
Now, Hezekiah is a good king..
He's a good king..
Now, why do I have to mention the good king?.
Because in the book of Kings, as you can remember,.
there were many, many kings that we quoted..
Northern country, Israel, the southern country, Judah..
The book of Kings is all about the history.
of the royal people, the kings and the queens.
and the elites of the country..
Most of the kings are bad kings..
If you talk about the northern king,.
the kings of Israel, the northern country,.
none of them is good..
None of them is good..
Everyone has gone astray from the faith.
and let the whole nation to worship the false god, Baal,.
worship the false gods..
None of them is good in the northern country..
That's why they went into exile..
And then he continued to talk about the southern kings,.
the kings of Judah..
There are many of them because the history is,.
of course, is longer than the northern country..
The southern country, Judah, also has many kings,.

$^{441}$but many of them are bad kings, all right?.
And here and there, there are some good kings..
Hezekiah is one of them..
Now, how do you define a good king, right?.
Now, I told you the bad kings,.
all they did was to worship the false god.
and lead the people to worship the false god,.
to leave the faith, the true faith of Yahweh..
The good kings will turn around,.
turn around the whole nation.
and have the whole nation worship Yahweh, the true God,.
the one and only God, the creator God, the covenant God..
So Hezekiah is one of them..
And in fact, in the early reign of King Hezekiah,.
he exercised and implement.
some of the all nationwide religious reform..
So he broke down all the false god shrines.
and false gods temples,.
some of the orders, rightful orders,.
he broke them all down..
So he got all, you know,.
all encompassing policy implement in his reign..
In the early time of his reign..
So that was King Hezekiah..
So there are many things the biblical author.
talks about this good king..
But one of the biggest crisis.
that Hezekiah faced when he became king.
was the incoming invasion of Assyrian empire..
Now, back then, one of the superpowers in the whole earth.
is the Assyrian empire..
Now, if you learn a little bit about Assyrian empire,.
you have never heard of, you know,.
some of the cruelest or the most militarized nation.
on earth, Assyrian empire..
They go to different places to conquer different people.
and to do some of the most evil things,.
you know, to oppress the conquered people.
and to destroy the cities and all the towns..
And the Assyrian empire was growing.

$^{481}$in the time of Hezekiah..
So one of the major crisis Hezekiah had.
was the growing Assyrian empire.
that would threaten the security.
of his small nation Judah..
So he did many things, Hezekiah,.
in order to defend his country, all right?.
Now, facing, you know, a Goliath, right,.
Assyrian empire coming at your footstep,.
you know, you gotta do a lot.
because you're a small nation, right?.
You got to spend a lot of time and a lot of effort.
in your military defense, right?.
Jerusalem is a very small place..
Jerusalem is also not a land of,.
that has a lot of resources..
So he did a lot of things, Hezekiah,.
in order to defend his country..
And one of the things that we can retrace his steps.
is the Hezekiah tunnels..
It is a almost 2000 feet long tunnel.
that brings water from the water source,.
the only water source in the city of Jerusalem,.
Gihon Spring, all the way,.
dug this tunnel 2000 feet underground.
into the mountain,.
all the way, bringing water from the Gihon Spring.
to the heart of the city, inside the city..
You know, in the old days,.
usually the water source or the spring of the city.
is located outside or at the foothill of the town, right?.
The town, when we build a town in the ancient time,.
we want to build a town on top of the mountain.
for defense, right?.
But you know, there's always pros and cons..
You build high, you don't get water, right?.
Where's the water?.
The water's coming from the underground water table, right?.
Have you ever dug well?.
You dug deep, right?.

$^{521}$You dug deep into the ground..
Why?.
Because the water table is underground, right?.
So you have to dug like 50 feet down, 20 feet down,.
in order to reach the water table in the limestone..
Same case in Jerusalem..
The water table is down there..
So imagine, you know, you got a hill, right, or a ridge..
City of Jerusalem, the city of David,.
is on top of the hill,.
and the water source is the Gihon Spring..
It's right at the foothill..
Now, how do you solve the problem of military defense, right?.
When the Assyrian Empire,.
when the Assyrian soldiers come to attack the city,.
they will surround the city, right?.
And they will tap into your water resource.
and won't let you use the water, right?.
So in the normal day, when there is no military attack,.
you can go out of the city wall, out of the city gate,.
and you go all the way down to the foothill,.
and then you have a bucket,.
and then you fetch water from the well, or from the spring,.
and then you go back uphill, right?.
Now imagine, day by day, you have to do that, right?.
In the morning, because you gotta make coffee,.
so you gotta have some water, right?.
So early in the morning, usually the women do the job,.
all right, the men stay back on the lazy boy,.
and wait for the women to bring water to home..
So in the morning, and probably in the evening,.
you know, the women are very busy going out of the city gate,.
out, you know, from the mountain top,.
go to the foothill to fetch water back home..
In the normal day, you can do that, that's fine..
But in order to prepare for the ongoing,.
or the upcoming Syrian attack,.
Hezekiah dug this tunnel, 2,000 feet long, is that long?.
Pretty long, right?.
All by human hands, can you imagine?.

$^{561}$All by human hands..
And today, we can go to the Hezekiah tunnel,.
and I tell you, this is the highlight of the whole trip..
Everyone thinks so, everyone thinks so..
Unless you are claustrophobic, kind of,.
you know, you can't go in..
We have a couple of members who cannot go in,.
so they decide not to go, which is good, which is good..
If you get panicked in the tunnel,.
there's no way I can rescue you, right?.
There's no phone signal, there's no Wi-Fi in the tunnel,.
so you have to walk through the tunnel one way,.
you know, from the Gihon Spring..
You walk in with your flashlight, right?.
You gotta bring your flashlight,.
and then the size of the tunnel,.
remember, I'm five nine, I have to duck most of the time..
I have to duck..
If you're really short, that's fine..
I have to duck most of the time..
And you have to walk through the tunnel.
from the starting point to the end point, 2,000 feet..
How long does it take?.
In the water, you wait in the water..
There's fresh water coming out of the Gihon Spring..
It still works, the spring is still working today..
You know, it's just amazing..
For 3,000 years, it's fresh water.
coming out of the Gihon Spring..
So we walk in the dark with a flashlight.
for almost 45 minutes, sometimes an hour,.
and sometimes longer because I stop,.
and I explain to the team right inside the tunnel..
And think about it..
You know, this Heisekai Tunnel was the greatest.
water engineering project in his day, in the entire world..
You don't find anything like this.
in other parts of the country,.
or in other parts of the planet Earth..
It's the greatest, smart, and genius idea..

$^{601}$They have two teams dugging together,.
you know, trying to join each other inside the mountain..
2,700 years ago..
By what?.
By sound..
When they dug in the limestone,.
that you can hear the other team, what they're doing,.
and are we coming closer?.
That's why the tunnel is not straight..
It's going like S-shaped, you know,.
zigzag into the mountain..
It's one of the greatest water projects.
in the world at Heisekai's time..
So Heisekai did that in order to defend the city..
Now back to the scripture..
Did the scripture talk about it?.
Yes, it does..
But right at the end, it's only one phrase..
All the other events of Heisekai's reign,.
all his achievements, and he mentioned this..
He made a water pool, and a water tunnel,.
and he brought water into the city from the Gihon Spring..
Outside of the city, into the city,.
in order to prepare for the Assyrian siege..
Now, for such a great water engineering project,.
it's very interesting, the Bible only has one statement..
All the stories of Heisekai, recorded in the Bible,.
are not about his human effort..
Because if you read the story,.
you understand eventually the Assyrians came.
and besieged the city..
And almost Jerusalem fell because of the Assyrians..
But you read the record of King Heisekai,.
at the day of distress, he went to the temple..
He went to the temple of Jerusalem to pray..
To pray for deliverance..
To pray that God would help him and to deliver him..
Now what a contrast..
Because tons of money and tons of effort.
go into ducking the water tunnel..

$^{641}$Right, the greatest water project ever in his day..
But even that system, even that water project.
did not save him, did not save his nation..
And Heisekai's story is all about the salvation of God..
Only God can save..
Humans cannot save..
Have you ever thought about it?.
Like the days of Heisekai, he spent tons of effort.
in order to defend himself, to prepare for coming crisis..
We do that too, right?.
How much time have we invest, think about it..
How much time have we invest into our 401k,.
into our pension, in thinking about our future?.
Thinking about how we prepare for, I don't know,.
retirement or days of distress, right?.
How to prepare for any crisis that we can think of..
The medical insurance in the United States,.
it's such a mess, right?.
But we have to spend so much energy, so much worries.
in order to prepare for that, right?.
Human efforts..
The story of Heisekai, the biblical record,.
give us another perspective..
Yes, you should prepare..
Yes, you should do all that you can.
in order to defend yourself..
But ultimately, you should not forget,.
it is Yahweh who saves, it is Jesus who saves..
I'll tell you the story and then I'm gonna end..
Years ago, I was in Vancouver..
I was a young pastor..
And you know, sometimes we do a pastoral visit.
to different families and different people..
I remember it was in the late '90s..
And I was in Vancouver, and there were many rich people.
in Vancouver..
Oh, they're so rich..
You know, in fact, I was a young pastor.
and I was taking care of a youth group..
And one of my students, he lives alone in a big mansion..

$^{681}$It's bought by parents from Hong Kong,.
but they don't stay there..
Only the kids stay there, right?.
17-year-old, 18-year-old, crazy..
I went to his house..
I really like to go to his house.
because you walk into his mansion,.
the huge living space is occupied by a snooker table..
You know the billiard?.
The snooker table?.
You know the pool table?.
Pool that you have a stick and then you poke the ball..
You know that table?.
But the huge size ones, right?.
The huge size..
So we go to his house and we play snooker, right?.
We play billiard..
It's just that rich, right?.
I remember one time I was arranged to go.
to one of the other businessmen..
Went to visit him..
And we went to his house and talked to him about gospel.
to share with him what Jesus has done for us..
You know how he look at me?.
I still remember how he respond to me.
because he's richer than me..
He's probably, I guess he's more good looking than I am..
He's successful..
He got the family..
He got all, everything intact, right?.
I mean, he said, "I have nothing to worry about..
"Don't talk to me about death and stuff.".
Right?.
What Jesus has done for, to him, Jesus is a poor man..
Right?.
To him, the gospel is, what, it sucks, right?.
It sucks..
Because compared to him, he has the best of the world..
Right?.
So I talked to him and kind of being looked down.

$^{721}$and kind of being, you know, he didn't yell at me.
but I can tell he's pretty, the despising tone and look..
And I left, right?.
I left..
And that was one of the worst moment in my ministry..
You know, I have doubt, what am I doing here?.
(laughing).
Some years later, some years later,.
you know, one of his family members called us and said,.
"Hey, can you visit my dad?".
Which is that guy, the businessman..
About three or four years later, I visited him again..
He had cancer, last stage..
We remembered the contrast..
We remember the contrast, his tone and his look..
Before, he was successful and despising..
Two or three years later, because he had been diagnosed.
with last stage of cancer, liver cancer,.
his tone totally changed..
His perspective totally changed..
At that point, we explained to him the gospel..
We explained to him what Jesus has done for us..
We explained to him the hope that we have,.
beyond death, right?.
As you can tell, when crisis hit,.
you cannot be proud anymore, right?.
You cannot be prideful of your achievement,.
of your human effort, because someday we're gonna fade away..
All we can do, the best we can do right now,.
is to delay our response..
Don't think about it, right?.
Well, we don't have to think about it..
Have you ever thought about tomorrow.
you're gonna be diagnosed with cancer?.
You won't, right?.
All you care about is, okay, today, tonight,.
we're gonna have a big feast, right?.
Yeah, we're gonna eat, and then we're gonna sleep well,.
and tomorrow we go to work, right?.
All we need to think about is this, right?.

$^{761}$To delay our response to the most important things.
in our lives..
Do we have hope beyond death?.
Jesus saves, and only Jesus saves..
That is what the story of Hezekiah is all about..
You can spend all your budget in your military defense,.
in your health insurance, in your achievement..
You can plan your education, you can plan everything, right?.
But only Yahweh saves..
Don't forget about that..
The salvation of Jesus Christ is so great of a salvation.
that we can have life beyond death,.
and this is eternal life..
Right?.
Yes, we need to prepare, we need to go to work,.
but don't forget, we need to worship..
We need to worship God and follow Him..
We need to give thanks, right?.
We need to do the will of God, because we belong to Him,.
and we know only Jesus saves..
All right, we're gonna end here,.
and we're gonna pray together..
May this be a reminder to all of us.
how great a salvation that we have in Jesus..
Let us all pray..
Father God, we give you thanks for saving us,.
for redeeming us..
Yes, all the human achievement that we can do,.
yes, we need to be responsible with our lives..
We need to plan, we need to arrange,.
we need to implement our responsibilities,.
but ultimately, let us not forget that only you save..
And thank you for your great love..
Thank you for your grace..
And be reminded that we are your people..
Be reminded, we need to praise you..
We need to give thanks to you, to follow you..
Thank you for being our leader..
Thank you for being our Lord..
And lead our life in the coming week..

$^{801}$And let us not be forgetful that we can stand in your grace.
to relate to you in our family, in our campuses,.
in our work, we're gonna glorify you..
In Christ's name, we pray..
Well, thank you very much, and I wanna see you next week,.
and let's worship together..
Let's worship God and serve him together..
God bless you..
\newpage



\section{}
\label{sec:aMelPey4yl8}
\textbf{佳節倍思恩(一): 住棚節(上) - 陳崇基牧師}
\newline
\newline
連結: \href{https://youtube.com/watch?v=aMelPey4yl8}{\texttt{https://youtube.com/watch?v=aMelPey4yl8}} ~~~~ 語音日期: 2022-09-24
\newline
\newline
\hyperref[sec:T9LnkydmX7E]{\small{< < < PREV SERMON < < <}}
~
\hyperref[sec:index_chronic]{\small{[返順時目]}}
~
\hyperref[sec:index_scriptual]{\small{[返順卷目]}}
~
\hyperref[sec:CndBkmupUxY]{\small{> > > NEXT SERMON > > >}}
\newline
\newline
$^{1}$我發現我預備的講章沒有什麼吸引力.
因為我沒有講吃東西.
我們去玩很多時候都不理會文化背景.
社會裡面的東西.
總之有東西吃就行了.
我們就是這樣的人.
你知道在以色列很有名的.
就是那些從中國大陸來的旅客.
那些導遊最大的抱怨是什麼.
就是那些人去到不知道為什麼中國大陸.
有時候可能都不是教會的人.
是非教會的人.
他就供了去以色列.
然後就在那裡玩兩個禮拜.
當地的導遊對中國大陸來的人.
最大的抱怨是什麼.
就是上車就睡覺.
下車就拍照.
拍完照就上車又睡覺.
然後導遊說什麼他就完全不理會.
然後回去拍了很多照片.
一百多張二百多張幾百張.
然後他就說拍了什麼.
他所關心的就不知道是什麼.
然後就在哪裡有好東西吃.
他才覺得是心滿意足.
我覺得這個文化很奇怪.
我相信聖地不可以這樣玩.
你可能去夏威夷可以這樣玩.
夏威夷也有很多歷史.
你應該去看一下.
你去其他地方可以這樣.
但是在那裡是不可以的.
因為我相信那地方有很多豐富的資訊.
豐富的資料.
我相信你去兩個禮拜.
等於你讀幾堂神學的課程.
我意思是真正的神學院的課程.
當然你一定要跟一個好導遊.
這個世界有很多好導遊.

$^{41}$我只是其中一個.
說起這件事.
我都有點後悔.
有一天我們去了納沙拉.
納沙拉有很多山路.
路很窄.
我們兩輛車.
我一輛然後Ronald一輛.
我走錯路了.
我走進去一個很窄的巷.
剛剛一輛車是下不了車的.
是這麼窄的.
我以為那條巷是可以出去的.
誰知道是出不了去的.
在一個45度的斜坡.
我們要掉頭.
那次真的很驚險.
車輪都出煙了等等.
真的很厲害.
我非常欣賞Ronald的車.
是可以OK.
我們成功脫險.
我最大的後悔不是走錯路.
我最大的後悔就是我們卡住的時候.
我們響了輪胎的時候.
我沒有拍攝那些東西.
我應該下車裝好鏡頭.
然後錄下最realistic的情況.
所以這個世界有很多好導遊.
有些人是不會走錯路的.
我們這次開始去拿一些題材.
在聖經裡面講到聖經的節日.
你可能留意過.
讀舊約新約的時候.
經常都會提到一些festival.
一些holiday.
一些節日.
聖經裡面其實有很多.
我們不講聖地了.
我們講聖經裡面的時間.

$^{81}$一年裡面有很多不同的節期.
春夏秋冬都會有.
其實對我們有什麼意思呢.
除了在歷史裡面有發生過.
記載了不同的事蹟.
其實對我們有很多提醒.
你想起節期的時候.
佳節.
佳節 裴司欣.
你想起佳節.
一些holiday.
一些重要的日子.
你會想起什麼.
想起什麼日子.
最近有勞工日.
勞工的假期.
勞工節.
911.
特別是我們美國.
如果是紐約的朋友.
一定是最大的印象.
最深的印象.
911的時候可能有不同的紀念項目.
Memorial Day.
接下來的有什麼節.
感恩節.
感恩節 沒錯.
感恩節.
又是吃的時候.
年底就是聖誕節.
這些節日很多時候就是一些marker.
一些標籤.
去提醒我們某些東西.
我們不要水過鴨背.
然後就過了去.
有一些reminder.
都是一些重要的節日.
譬如有什麼reminder.
對你來說可能你不會miss的.
生日.

$^{121}$你有沒有想過生日蛋糕是從哪裡來的.
有沒有想過.
我們就知道總之是買了生日蛋糕.
大概中世紀的時候.
在德國有班人.
開始為小朋友生日的時候.
做蛋糕.
然後越做越甜.
然後來到美國就更加甜.
所以birthday cake是這樣的來由.
去到慶祝某一些東西.
你想起節日的時候.
festival.
holidays.
或者一些important dates.
我們就一定想起某一些reminder.
多想一想聖經.
聖經裡面有一些節日.
一些festival.
都是有這樣的目的.
在《五經》裡面申命記第十六章.
它說你們所有的男丁.
當然當時的社會是男人.
擔起頭家.
in fact 今天應該也是.
男丁他們就是代表整個家庭.
他要在三個節日裡面.
一年三次.
去到耶和華所揀選的地方.
去見神.
去敬拜祂.
這是特別的敬拜.
除了每個禮拜有一個shabbat.
shabbat去到安息日去會堂敬拜.
或者去聖殿敬拜等等.
一年三次.
特別的節日.
一些holidays.
一些festival.
不可以空手去見耶和華.

$^{161}$各人要按照自己手中的能力.
召耶和華.
你上帝所賜你的福.
去奉獻禮物.
三個什麼節日呢.
第一個.
除孝節.
Feast of Unleavened Bread.
其實是passover.
因為passover是一天.
完結之後.
就是一個禮拜七天的.
就是叫做除孝節.
他們吃的東西就是不會有yeast.
不會有孝.
就是紀念什麼.
他們出埃及的時候.
以色列人出埃及.
是趕著出去的.
趕著出去的時候.
你做那些所謂麵包的時候.
做你的食物的時候.
就不可以等它發酵.
你做麵包你知道的.
要等一等.
然後等它發酵.
出埃及的時候.
是夜間出埃及的.
上帝拯救他們.
以恩典拯救他們.
然後擊殺法路.
擊殺埃及.
是沒有時間的.
是趕著走的.
所以passover除了吃完羊羔.
七天他們就去紀念.
unleaven bread.
就是當年的事蹟.
當年上帝的拯救.
他的salvation.

$^{201}$然後第二個日子就是七七節.
七七節就是passover之後.
七七四十九日.
七個禮拜之後.
七七節.
或者有不同的名字.
shut for all.
或者收割節等等.
新約裡面就叫做pentecost.
pentecost就是fifty的意思.
七七四十九日.
第五十日.
那個又是另一個收割節.
就是五旬節.
這個又是要去紀念的.
紀念上帝給我們的土產.
給我們的供應.
然後第三個節日就是駐棚節.
feast of tabernacle.
或者feast of booth.
或者希伯來文叫做sukkot.
那又是七天的.
又是今時今日都是一樣.
是一些很大的節日.
很隆重的慶祝.
三次的提醒他們要做些什麼.
要去到上到耶和華所揀選的地方.
以前就是耶路撒冷.
今天就在基督裡面.
上到耶路撒冷.
他們就去守節.
並且帶出奉獻.
去到按照自己的能力.
然後去到供應上帝的國度.
他的事工等等.
總之就是一種敬拜的reminder.
一年裡面有三次.
除孝節pass over大概是什麼時候.
我們基督教年曆的good friday.
和easter的時候.

$^{241}$那個就是pass over 除孝節.
七七節五重節.
再數七個星期.
就是pentecost.
使徒恆傳裡面提過.
然後駐棚節就是現在大概九月尾.
十月初左右.
如果你check一下今天的calendar的話.
大概就是十月頭.
即是快到了.
他們就在Sukkot.
在以色列這個地方.
我們今天看第一個節日.
我們嘗試看這個節日的其中一段.
一部分.
到底這個節日是用來做什麼的呢.
剛才我解釋了pass over.
解釋了七七節.
pentecost.
然後feast of tabernacle.
feast of booth.
Sukkot.
到底是用來做什麼的呢.
在五經裡面利未記23章.
就這樣吩咐以色列人.
神的子民要做什麼呢.
你要吩咐以色列人說.
七月的十五日.
就是feast of tabernacle 駐棚節.
要向耶和華守這節七天.
你們要做什麼呢.
specifically 具體的跟他說.
你們要住在棚裡面七天.
凡以色列家裡面出生的.
都要住在棚裡面.
好叫你們的後代.
你的小孩.
你的下一代.
知道我領以色列人出埃及的時候.
曾使他們住在棚裡面.

$^{281}$我是耶和華.
你們的神.
駐棚節是什麼時候呢.
十月初.
或者九月尾.
就叫他們做什麼呢.
他們就一定要住在棚裡面.
整家就住在棚裡面.
換句話說.
他們進到家林地的時候.
他們自己的房屋是石砌出來的.
但是到這個駐棚節的時候.
九月尾.
十月初.
他們農曆的七月十五日.
他們就要七天.
在棚裡面居住.
紀念什麼呢.
remind他們什麼呢.
出了埃及之後.
四十年他們走礦業.
走礦業住什麼地方.
沒有石屋住.
在周圍見到水的地方.
他們就會去.
在礦業的地方.
他們就住棚.
有如我們今天camping.
有誰喜歡去camping的呢.
有.
住camping很過癮吧.
那些沒有舉手的請不要舉手.
camping.
就是一些很temporary.
臨時.
紮的東西.
沒有人去camping.
就想著.
我在這裡留三年半載.
很少會這樣.

$^{321}$都是暫時暫居的地方.
很輕便的.
他們以前走礦業的時候.
就正是這樣.
今年你check check日期.
大概是十月九至十六日.
就是今年的sukur.
在以色列這個地方.
就會有些人去.
在outside.
去做紮營.
做一些營出來.
做一些所謂棚.
他們叫做sukur.
我記得在.
今年我們去聖地.
是covid以來.
疫情以來第一次.
我有一點心大心細.
是不是這個時候最好.
不過我聽見朋友去了.
我就發現是ok.
因為everything is open.
所以我們就今年.
疫情以來第一次.
還有很多次的.
我希望.
但疫情之前.
我們2017年的時候.
我記得那時候我和我老婆.
我就和call team說.
不如讓我們放一放假.
我們想走一走.
休息一下.
他們很照顧我們.
儘管走開.
我和我太太就去了以色列.
我們兩個人去.
2017年.
那時候我都沒有想過什麼timing.

$^{361}$什麼時候去最好.
我call team說ok.
我們就選一個時間.
照顧小朋友.
然後我們就說.
ok我們去兩個星期.
然後我們就買機票.
訂airbnb等等.
我們去的時候.
是九月底.
那兩個星期.
我以為是最好的時間.
因為我去到當然要四處走走.
我經常和我老婆說.
這次你真的執到寶.
有一個隨身最好的導遊.
他跟著你.
誰知道他覺得不是這樣.
他覺得我們在趕他走.
無論如何我們去兩個星期.
我們投得去的時候.
我以為是一個最好的時間去.
誰知道去到第一個星期.
原來就是擁潔坡.
就是贖罪日.
是他們以色列國裡面.
最大的一個high holiday.
贖罪日就是他們刻苦幾生.
然後全世界都躲在家裡.
樣樣都關了.
所以在擁潔坡.
你知道星期五新丹之後.
他們又沒有車.
街上很少人等等.
如果是擁潔坡.
一年的一天.
九月底十月初.
贖罪日.
在《五經》裡面有提過.
他們要刻苦幾生.

$^{401}$然後到悔改.
去到想清楚下一年的方向.
又是一個reminder.
這個我們今天不說了.
但無論如何我們去了擁潔坡.
誰知道擁潔坡真的什麼都關了.
這張照片.
大概是早上.
到中午的時候拍的.
交通燈就閃了.
我站在路中心拍這張照片.
完全不用看車.
因為沒有車.
除了小朋友可能踩著單車在玩.
擁潔坡全部寂靜.
我說真的慘.
所有東西都關了.
吃屎都關了.
不過有的.
去基督徒區就有.
因為基督徒不會有猶太節日.
我們去到擁潔坡.
慘了這次.
幾天之後就是Sukkot.
交了幾天之後就是Sukkot.
Sukkot又是很多東西要預備.
雖然是高度假期.
但都有很多interruption.
有些地方有些公園沒有開.
因為Sukkot的緣故.
我就覺得原來我都沒有留意.
所有東西都關了.
但可能對我來說都是一個good break.
我們很安靜地過.
但我們經歷了Sukkot.
這個高度假期.
豬棚節.
他們在豬棚節的時候.
他們會在外面搭棚.
你見到這些臨時搭出來的棚.

$^{441}$如果你是handy一點的.
你就可以的.
如果你本身不是很handy.
即是搭釘搭東西的時候.
不適合你的話.
那就很麻煩.
是一個很重要的日子.
所以他們都非常重視.
他們會在Panio.
然後他們就見這些棚.
很有趣.
What does it mean?.
這就是我們想.
這次第一次去到講一講.
這些所謂Sukkot.
Sukka他們叫.
到底有什麼意思呢?.
這些都是你可以在Sukkot的時候.
即是豬棚節的時候.
在以色列都能夠見到.
這裡是西牆的廣場.
他們就拿那些棕樹.
即是剛才惠仙說的棗樹.
那些葉子.
那些枝那些框.
他們就去弄一些帳棚出來.
然後那星期.
很多時候他們就在帳棚底下.
去吃東西.
三餐.
然後有時候有些人甚至睡覺.
在這些Tent裡面去睡覺.
是什麼叫Reminder?.
什麼叫Reminder?.
他們這些猶太人的做法.
有時候都會很多時候討論.
這些Sukka.
棚子應該怎樣做.
有一個故事是這樣說的.
以前曾經有兩個很出名的拉比.

$^{481}$律法教師.
他們有一次因為Sukkot的時候.
他們就在坐船.
在旅行.
因為Sukkot來了.
所以我們要建棚.
船上他們就去建棚.
其中一個拉比就在船頭.
在頂上.
就做一個棚出來.
我不知道他找什麼做.
因為做了之後.
晚上因為有風來的時候.
那個棚就塌了.
另一個拉比就跟他說.
為什麼你做一個棚這麼容易塌.
那個拉比就回答.
他說這個Exactly.
就是我們做棚的時候的要求.
做了之後是會塌的.
是很容易塌的.
因為那件事是Temporary.
等於我們去露營.
我們露營的螢幕.
都是很Temporary.
In fact 他們的律法裡面.
就是說我們做棚的時候.
就是不可以太實淨.
因為鑄棚節要紀念什麼.
就是他們一代40年.
在曠野.
上帝帶領他們走過.
鑄棚.
那個棚.
Sukkot.
其實代表我們人生.
就是一些很Temporary.
很暫時.
很暫時就會經過.
很快我們就會經過我們人生.

$^{521}$鑄棚節每一次來到的時候.
他們住在棚裡面.
去吃喝 去唱歌.
去到有團年等等.
吃飯.
都是remind他們.
生命裡面我們就是.
永遠都不會真的落地生根.
Life is transitional.
永遠都是過渡.
永遠都是過客.
無論你今天住的地方.
你選擇的世界都好.
都是很容易就會倒.
Life is fragile.
很容易爛.
很容易就過去.
最近我去了一些婚禮.
以前的學生.
現在長大了.
結婚.
見到很多很久沒見的人.
我們一整晚.
見到朋友say hi的時候.
經常都有支歌唱.
有句對白.
就是Time flies.
時間真的過得很快.
大家說起.
特別是他見到小朋友.
以前是鬥釘那麼小.
或者上次見他.
在IG或者Facebook見到他的照片.
鬥釘那麼小 手抱.
今天見到他.
他已經可以到處走 去搞事.
問我.
我的兩個小朋友.
原來他那麼大.
我以為他還在高中.

$^{561}$他說不是.
明年我女兒大學畢業.
他要找工作.
你想想.
我昨晚想起.
我的雪櫃.
有他們小朋友的照片.
我說那不是妹釘嗎.
今天不可以叫妹釘.
今天已經少女了.
快要叫她女人了.
Life is transitional.
朱鵬子就是remind我們這件事.
他說他很叫你世世代代.
紀念你出埃及的時候.
上帝帶領你.
住過棚.
不是permanent的東西.
但是上帝帶領著.
如果是這個緣故.
他remind我們人生是多麼快會過去.
你就應該想一想.
到底你應該怎樣計劃人生.
在上帝裡面.
到底你應該做些什麼.
既然生命會過得這麼快.
時間會飛逝.
到底我們應該在神的面前.
怎樣聽他旨意.
我希望提醒你.
Life is not permanent.
是transitional.
在上帝面前我們應該怎樣立志.
我們同心協力去祈禱.
上主我們感恩.
因為你是那位.
帶領我們走過曠野.
帶領我們走過高高低低的上帝.
神我們去紀念你.
你去remind我們.

$^{601}$提醒我們.
人生就是一些很快的過渡.
暫時的.
並且非常脆弱.
朱鵬的意思.
讓我們謹記.
既然生命是這麼短暫.
時間會很快地飛逝.
我們怎樣把握這個時間.
我們的生命來事奉你.
為了一些有永恆價值的事.
我們去憂去忙.
求你幫助我們.
幫助教會.
能夠在你的reminder裡面.
更盡心盡力.
愛你跟隨你.
\newpage



\section{}
\label{sec:CndBkmupUxY}
\textbf{The role of Nations \ and  Neighbors in a growing church (Luke 10) - Keith Draper}
\newline
\newline
連結: \href{https://youtube.com/watch?v=CndBkmupUxY}{\texttt{https://youtube.com/watch?v=CndBkmupUxY}} ~~~~ 語音日期: 2022-09-30
\newline
\newline
\hyperref[sec:aMelPey4yl8]{\small{< < < PREV SERMON < < <}}
~
\hyperref[sec:index_chronic]{\small{[返順時目]}}
~
\hyperref[sec:index_scriptual]{\small{[返順卷目]}}
~
\hyperref[sec:IDFYoAmcEPE]{\small{> > > NEXT SERMON > > >}}
\newline
\newline
$^{1}$All right, I want to say thank you to Kelly for doing that really off the cuff, being.
able to speak directly from the scriptures there..
It's a long passage, it includes a number of different things, and you might think,.
I hope he's not going to talk about all of that as we're here today, and I'm not..
I thought maybe as I come, often when I am here, some of the illustrations I use will.
be of my family, and most of you maybe have never seen my family because I'm traveling.
to different churches, but I have ten children, three of them are now married, and I have.
two grandchildren, so this is my wife and my family..
So when you hear my illustrations about big family and craziness of that, now you'll get.
a sense of that, and that's out on our front porch here in the West Suburb..
So anyway, that was free, that wasn't really part of the sermon, just a little bit of back.
story on me so you'll know me a little bit better..
But I really want to come to you today and bring you greetings as my brothers and sisters.
in Christ who are here at Urban Voice..
I'm so glad to rejoin you today on this journey you're on, and to really encourage you from.
the Bible as fellow followers of Jesus Christ in Chicagoland..
I confess, though, that I am far less aware of you and what you're up to these days, as.
I have devoted much more of my time lately to the many new people that are being displaced.
here from places like Afghanistan, Ukraine, and now Venezuela and Colombia..
Though from being with you when you started as a church, I know that you're approaching.
an age when inertia begins to set in, an inertia that causes churches to focus more inward.
than outward, to be more concerned about maintaining than growing the church..
As a church grows older, it takes more and more people in the congregation to see one.
new follower of Jesus Christ..
And just like you and me, if your church is not growing, you're declining..
So this really begs two very crucial questions that I'd like to ask you today..
What does growth look like for your church?.
And number two, how can I keep my church growing?.
To begin to answer these very important questions, from almost 40 years of missionary experience.
that I've had here in Chicagoland, this is what I would say to you..
Hear me in this..
There are a lot of good things that you can do as a church, but you must keep your focus.
on the greats..
You must keep your focus on the greats..
Let me repeat that to you..
There are a lot of good things that you can do as a church, but you must keep your focus.
on the greats..
So what are the greats?.
Well, I'm referring to the great commandment and the great commission..
So let me just quote them for you so that both of them are in your mind here..

$^{41}$The great commandment is what we just heard from Luke chapter 10, and it's to love the.
Lord your God with all of your heart, soul, mind, and strength..
And the second is likened to it, to love your neighbor as yourself..
The great commission is at the end of Jesus' time on earth as he is ascended to heaven,.
and these are the last words that he has to say, "Go and make disciples of all nations,.
baptizing them in the name of the Father, the Son, and the Holy Spirit..
And lo, I'm with you always, even to the end.".
So I want you to notice with these two greats, the great commandment and the great commission,.
that there are two key words in them that begin with the letter N, the letter N, like.
M-N-O-P-N..
And the Ns that are here, the two words that begin with N, are numbers and noses, which.
sometimes in church growth we kind of get into that..
How many people did we have show up for service today?.
What does the number look like?.
But really what you see in these two greats is two other Ns, and that is the N of neighbor.
and the N of nations..
So I want to start by talking a little bit about growing the church through neighbors.
and nations..
Let me start with nations..
Focusing on nations seems like a pretty straightforward approach for a church, and maybe even a roadmap.
to try to track progress as a church..
It's almost like a continental approach of reaching others with each continent, like.
starting in Asia, America, Europe, and going to all of the continents..
Well, maybe not Antarctica..
I don't think we have really much mission work that needs to be done in Antarctica..
It does mean going after people in each country, though, to know Jesus when we take a nations.
approach..
But here's what's interesting..
Nations rise and fall..
Nations come and go..
Sadly, we see that even with Ukraine recently as a country's trying to take it over and.
basically absorb it..
So how do we fulfill our mission if we're looking at nations when that list of nations.
is changing constantly?.
The United States didn't even exist when the first church was born..
My point is this, how do we measure progress when our lists are ever-changing?.
To bring greater clarity to this, I believe it helps us to understand the meaning of nations.
in the Scriptures..
The phrase for nations is actually the phrase "ponta ta ethne," and it's a Greek word,.
which, or Greek phrase, which means "people groups.".

$^{81}$So if you take the Great Commission, it's to go and make disciples of all people groups,.
all ponta ta ethne, baptizing them in the name of the Father, the Son, and the Holy.
Spirit, teaching them to observe all that I've commanded you..
In Acts 2, we actually see one example of this as the church began..
There's a list of people who heard the gospel at Pentecost, and it included this list that.
is given, Parthians, Medes, Elamites, Mesopotamians, Judeans, Cappodiceans, people of Pontus, people.
from provinces across Asia, Phrygians, Pamphylians, Egyptians, Libyans around Cyrene, visitors.
from Rome who were both Jews and converts to Judaism, Cretans, and Arabs..
Now that's probably the list I should have given Kelly to read, but then she would not.
have liked me anymore today after reading that list..
So why such a comprehensive and particular list of the nations is given here?.
Here's the early church, the Spirit has come, and this group of people respond..
It's just a group of people..
But in the Scripture here, it's a particular list of particular people from particular.
places..
And that's because the church was birthed not just for the Jews in Jerusalem..
As the Spirit fell on the church, people groups speaking different languages heard the gospel.
in their own language..
They repented, and it says that there were over 3,000 people who were added to the church.
that day, at least 3,000 people..
And then it says they all went back to their homelands to tell others what they had heard,.
and a multi-ethnic church spreads to other ethnicities and other people groups as the.
gospel goes to the ends of the earth..
You know, yet people groups might seem to you, too, like a moving target..
Just as we talked about nations coming and going, people groups might seem like a moving.
target to you..
But I want to assure you from Matthew's gospel that it is a finite work, meaning it has a.
start and a stop to it..
In Matthew chapter 24, verse 14, it says this, "And this gospel of the kingdom will.
be proclaimed throughout the whole world as a testimony to all people, or all Pantetah.
ethnei, and then the end will come.".
So let me say this, one day our job will be done, and Christ will return to bring us to.
heaven to worship Him from all that have been reached by the church through the years..
And it says that they will be from all peoples, tribes, and tongues..
And this is given to us by a vision from John in the book of Revelation chapter 7, verse.
9..
I know I'm going to a lot of places in Scripture, so my apologies here, but then if you look.
in Acts 1735, you see that Paul tells the people of Athens how important proximity and.
purpose is to our mission..
He says, "From one man God created all nations, or Pantetah ethnei, of men and women, and.

$^{121}$determined the times and the places where they live.".
Why did He do that?.
It says, "So that they might know Him, though He's not very far from any one of us.".
So therefore, one key really to a growing church is keeping the nations in front of.
us, especially those nations who are in our own backyard..
To help the nations know Him, and especially nations that are near us in global cities.
like Jerusalem then and Chicago now, a church must have a clear vision that keeps nations.
as a priority in our work..
Think about it..
We even have states like Texas that are busing nations here regularly to us, over 500 over.
the past three weeks..
Could I ask you, is there a particular people group that God may be touching your heart.
to share and show the love of Christ to them?.
Could it be that together with brothers and sisters in Christ here that you do that?.
Nations will never know unless you go..
Go and tell them about Jesus, and some of them are your people in your own backyard,.
even your own neighbors..
So then I'd like us to go on to think a little bit about neighbors now..
I mentioned nations, I want to talk just a bit about neighbors..
The great commandment speaks to us of really having a sold-out kind of love, to love the.
Lord God with all of your heart, your soul, your mind, and your strength, and love your.
neighbors as yourself..
Well, I think we have a pretty good idea of how to love ourselves, but we don't always.
know who our neighbor is or how we're supposed to love them..
Jesus was asked this specific question, and his answer was quite telling in what we've.
seen here in Luke chapter 10 from our Bible reading today..
It's a story not of greats, but is a surprisingly interesting story of a good person, a good.
Samaritan who's normally despised by others..
Jesus tells the story to inform us who our neighbor is and how we might behave toward.
them, so pay attention here..
The broader context, which is why I had Kelly read the whole chapter, is the context of.
deployment..
Jesus is deploying his disciples at the beginning of chapter 10..
He's sending out a second wave of 72 followers in pairs..
Now that must be a growing church when you can send out 72 people and still be growing,.
right?.
But his purpose is so that they would go out to tell still others, so that they would believe,.
and he warns people who don't believe..
Then something really amazing happens here..
As they go, it says that even the demons obeyed them, and Jesus says that he saw Satan fall.

$^{161}$like lightning at the power of Jesus in them as they went..
You get that?.
As they were going, as they were out doing the work, God's power was with them, and even.
Satan fell at the act of God's people on the move, people going..
So Jesus prays for them, and he's rejoicing in their childlike faith and the obedience.
that they have to change their neighbors and even to defeat Satan in those actions..
But it's not happening by people who are wise and clever by the world's standards..
Jesus gives this parable then of the Good Samaritan as an example of what we're supposed.
to do..
Finally, the chapter ends with this story that seems unrelated about these sisters,.
Mary and Martha..
We have four characters in the story, especially the Good Samaritan story..
One of them is his own, a fellow Jew who's left for dead on the side of the road, and.
then there's this preoccupied pastor and deacon, that's my modern day descriptions of that,.
who are preoccupied..
They've got way too many things to do to step into this situation..
And then there is someone who is outside of their own beliefs, not one of their own, who's.
different in belief, and he's also different in culture from them, and he's even despised.
by them, but he proves helpful to this person..
And only one of these four characters is described as good, the good neighbor..
The despised Samaritan is the only one who is described as a good neighbor in this story..
And what does this good neighbor, this despised Samaritan do?.
Well, he really, I think, lays out four things that you see happening here if you just look.
at that passage..
It says that he felt compassion for the person that was hurting..
He saw a person on the side of the road nearly dead, and he felt compassion for him..
It says then he stepped into the world of that person to help them..
He met their urgent needs for help and for healing, and then he made plans for their.
longer term needs..
Sounds kind of familiar, doesn't it?.
Think about it..
What did Jesus do for us?.
He felt compassion for us hurting..
He stepped in our world to help..
He met urgent needs that we had, but he also made plans to address our longer term needs..
Jesus describes himself as a good neighbor, and he gives us the example of a despised.
Samaritan also being a good neighbor..
Now why is he a good neighbor?.
In the story it says because he showed mercy..
He crossed culture and language and religion to even an enemy by stepping in to help instead.

$^{201}$of stepping away..
So maybe to be a good neighbor we just need to do the same thing, right?.
Not so easy, is it?.
It wasn't easy for the lawyer either who heard this..
He asked the question, he said, not because he wanted an answer, but he wanted to justify.
himself..
He already had given the correct answer..
He had quoted the great commandment to love the Lord God with all your heart, soul, and.
mind, and love your neighbor as yourself, and he said that he had done this all of his.
life..
I've always loved the Lord God Almighty, and I've always loved my neighbor..
So Jesus obviously knew that he needed an illustration to make clear and more measurable.
where he was missing the mark..
So Jesus told the parable of the good Samaritan to explain how to be a good neighbor..
Then Luke follows with this seemingly unrelated story of sisters, Mary and Martha, as we close.
this chapter on deployment, on Jesus' disciples being deployed..
It says that Jesus is present with the sisters, and he's showing ways that they responded.
to his presence with them..
One is Martha, who welcomed Jesus, but then is more anxious about all the things that.
need to be done, all the things that need to be maintained..
The other is Mary, and Mary's concerned about sitting at the feet of Jesus and learning.
from him..
Jesus says Mary discovered the one thing worth being concerned about most, being with Jesus.
and learning from him..
So now how does all this tie together?.
Well, what does all of this mean for you and I today?.
Well you see, loving neighbors and nations starts by remembering our first love, Jesus,.
and then listening to what Jesus tells us to do, even helping people that are different.
from ourselves living in our own neighborhoods..
Yes, we cannot lose sight of our first love, and it's because of our growing love for Jesus.
that we see the nations and the neighbors around us in a more compassionate way, worthy.
of his love too..
We love because he first loved us, even when we were his enemies..
So what is God saying to you today?.
Maybe your love has grown cold for Jesus..
Why don't you confess that to him and ask him to rekindle your love for him?.
He can do that..
And if you love Jesus, then you should love your neighbor more, even those who are not.
like you..
I believe God brought your church here together and put you where you are to grow here and.

$^{241}$not decline, to thrive and not just survive..
Where do you see neglect of your neighbors and your nations around you?.
Where you see neglect of neighbors or where you see nations that are not being reached,.
maybe there's something you should do about it..
Personally, I want to ask you, what is your next step?.
Let's pray..
Father, thank you for the moments to be here..
What is a very long passage but has some incredible meaning for all of us as followers of you.
that have been deployed beyond what happened in that town of Jerusalem to all of the ends.
of the earth..
Thank you, God, that the church then saw their neighbors, those who spoke different languages.
and had different backgrounds, and pursued them..
And today we are beneficiaries of what the church did then to reach those beyond themselves,.
their neighbors and the nations that were in their midst..
Show us how to do that today, God..
Give us a vision for what you would have us do that isn't just about how many people.
sit here in a seat, but about the neighbors that you put around us and the nations that.
need to know you..
And we pray these things in Jesus' name..
\newpage



\section{}
\label{sec:IDFYoAmcEPE}
\textbf{Faith and time: Feast of Tabernacle 1 (Rev. Mark Chan)}
\newline
\newline
連結: \href{https://youtube.com/watch?v=IDFYoAmcEPE}{\texttt{https://youtube.com/watch?v=IDFYoAmcEPE}} ~~~~ 語音日期: 2022-10-08
\newline
\newline
\hyperref[sec:CndBkmupUxY]{\small{< < < PREV SERMON < < <}}
~
\hyperref[sec:index_chronic]{\small{[返順時目]}}
~
\hyperref[sec:index_scriptual]{\small{[返順卷目]}}
~
\hyperref[sec:eW87A4hZtZA]{\small{> > > NEXT SERMON > > >}}
\newline
\newline
$^{1}$Today, we're going to begin our series on faith and time..
Think about it, your Christian faith, your walk in Christ,.
has so much to do with how you use your time,.
how you mark your calendar, right?.
How you, you know, check the clock..
Faith and time, they go together..
Sometimes, we think that Christian faith is all about,.
you know, our verbal confession, what we say..
Do you believe in Christ?.
Yes, I do..
Are you a Christian?.
Yes, I am..
Verbal confirmation..
Sometimes, we stay on this level..
Hopefully, this will match what our heart says, right?.
Otherwise, it means nothing to all of us.
and means nothing to you..
But, you know, think deeply about faith..
Ask yourself whether your faith is biblical or not..
Because, you know, you can invent.
your own version of Christianity, right?.
There's so many examples in the Bible,.
or, you know, outside of the Bible, in church history..
People would create their own faith,.
create their own version of Christian faith..
And they end up, you know,.
deviating from the teaching of Christ..
That's why in Christianity, in biblical faith,.
it's the faith of book..
You know, have you ever wondered?.
There's so many religions out there..
They don't care much about, there's one book.
defining what they do and what they think.
and what they believe..
So many religions in the world,.
so many worldviews in the world,.
even a religious worldview,.
they don't really have some sort of standard, right?.
It's pretty much anything goes..
You like to define your faith,.

$^{41}$you like to set a standard for your faith,.
you know, it's all up to you..
It's like going to Whole Foods, right?.
You go to Whole Foods and you pick and choose, right?.
Whatever you want, or whatever you feel like..
Now, this kind of Whole Foods, you know, faith,.
has nothing to do with biblical faith..
It has nothing to do with the Bible..
So think deeply about it..
Even you profess to be a Christian,.
even you were baptized, right?.
Couple weeks ago, we just returned from Israel.
and some of my team members,.
they were baptized in the River Jordan..
You talk about holy water, it's right there..
You know, everyone wants to, you know,.
get baptized there, right?.
There's something special about it..
Even you do that, it really doesn't mean,.
or not necessarily mean that you follow Christ..
Think deeper in your faith, you know, think deeper..
Because it's dangerous to come to the Bible,.
come to God, as if we go to Whole Foods,.
to pick and choose, you know, according to our own choices..
That's why we need to think deeply.
about what faith really means..
Faith and time, I want you to think about in this series..
Faith and your time..
What does it mean to follow Christ.
when it comes to how you use your time.
and how you see time?.
So this series is on faith and time..
And we'd like to pick some of the passages,.
some of the content from the Bible.
that talks about biblical faith and time and festivals..
There are certain markers in the Bible,.
especially in the Old Testament and New Testament as well..
During the week, during the month, during the whole year,.
there are special markers, there are special dates.
during the year, during the month, during the week,.

$^{81}$that marks the grace of God,.
that tells the people of God, "You have to worship me..
"You have to remember something.
"about your Christian journey..
"You have to think about and to be reminded.
"your experiences of God.".
So think about dates, some of the dates.
that appear in the Bible that has to do.
with walking in the faith of God..
And think about time, we are beings in time..
We are all controlled by time..
We are all managing our time..
Think of birthday, all right?.
How many of you, whose birthday.
are in the month of September?.
Very important month, right?.
Very important month..
(laughs).
'Cause my birthday is in September, right?.
How many of you birthdays in, well, what's the date now?.
October, yes, October..
October, all right, all right, all right..
We got some hands here..
So birthday, you know, it's a marker.
in your journey of life, right?.
You mark something important..
You have done one year, right?.
You have gone through one year..
And what you look for the next year, right?.
So in birthday, sometimes we have birthday cake..
Do you know the origin, origin of birthday cake?.
Where did it come from?.
We all have birthday cake on our birthday, right?.
Where does it come from?.
Definitely not in times of Jesus..
Definitely not in the Bible..
Well, if you chase the history of birthday cake,.
it actually comes from the Middle Ages,.
when some of the people in Germany,.
they like to make cakes for their kids.

$^{121}$when they celebrate birthday..
And then, you know, slowly and gradually,.
that tradition builds up, right?.
In the very beginning of the history of birthday cake,.
birthday cake was not very sweet..
But as it imported to America, you know,.
we add more sugar, right?.
So that's why we have birthday cakes,.
very sweet birthday cake today..
Think about other times in the year..
Today is October..
November..
Thanksgiving, right?.
What do we do in Thanksgiving?.
Hot pot, right?.
Hot pot. (laughs).
Hot pot, right?.
Not turkey..
What is Thanksgiving all about?.
It's kind of like a national cultural story, right?.
How the pilgrims came in their first harvest,.
and then they celebrate the very first time.
God has provided for them, you know..
That's a kind of a cultural memory,.
or cultural identity, right?.
All Americans, all Americans,.
Asian Americans, all kinds of Americans included,.
we all share that kind of identity, that kind of story..
Christmas, right?.
December 25th, 24th,.
sign night..
It's markers of memory,.
to remember something that God has done, right?.
What else?.
Important days during the year..
Easter, Good Friday, right?.
You gotta remember that, right?.
Because those are the dates.
that reminds you the cross of Jesus,.
what he has done..

$^{161}$So think about time and your faith..
You don't wanna forget, right?.
For those who are married,.
you better remember the day of your anniversary, right?.
Sometimes I don't..
(laughs).
Forgot, yeah, I forgot, I confess..
Sometimes I even forget my own birthday..
So, I'm not very date-driven, you know?.
Sometimes I'm too busy..
But, you know, it's all because of that..
We're going through our daily life in busyness,.
you know, being distracted,.
and sometimes life becomes, you know, chaotic..
That's why there are fixed dates in the calendar.
to remind you in your faith journey..
So, the Bible has a lot of different dates.
that reminds the people of God.
and what following Christ or what following God.
is all about..
Deuteronomy 16, 16..
It says three times a year, three times..
What are the three times?.
Three times a year, all of you, represented by men,.
must appear before the Lord, your God,.
at the place he would choose..
You know, eventually it's Jerusalem, the temple..
At the Feast of Unleavened Bread..
First time..
The Feast of Weeks, second time..
And then the Feast of Tabernacles..
No one should appear before the Lord empty-handed..
You gotta bring something..
Your offering, your sacrifices, your attitude,.
your whole being, each of you must bring a gift.
in proportion to the way your God has blessed you..
So there are poor people, there are rich people,.
they come three times a year to the temple to worship..
Right?.
Now, of course, every single week, you know, on Shabbat,.

$^{201}$they will worship..
That is one of the most important marker of their calendar,.
a weekly reminder, the Shabbat..
And then three times a year,.
there are three great festivals,.
three great dates that they must remember..
It's kinda sound like our Good Friday,.
our Easter, and our Christmas..
Three times a year, they must come before the Lord.
to praise him, to worship him..
The first one is Feast of Unleavened Bread..
That equals to Passover..
Passover, one day, and then seven days afterward,.
they celebrate for seven days and eating unleavened bread..
Passover..
We'll talk about that more later..
The Feast of Weeks, that is seven weeks after Passover..
The 50th day after Passover is another day of harvest..
They celebrate the harvest of barley and wheat,.
and wheat and all that..
In New Testament term, it's Pentecost..
Pentecost means 50, right?.
So it's the 50th day after Passover..
That's the second time they must come.
before the Lord in the temple..
The third time is Feast of Tabernacles..
Tabernacles falls on the fall, the early fall..
If you check the morning calendar, if you Google,.
Feast of Tabernacle, 2022, or Sukkot,.
S-U-K-K-O-T, Sukkot, 2022..
It falls on, what's today?.
October--.
- [Audience Member] Second..
- Second, right?.
Second, right?.
Sorry, yeah, second..
The next week, actually, is Sukkot in Israel..
Sukkot..
It's the fall festival..
So three times a year, they must bring something,.

$^{241}$some gift to this important worship opportunity..
What do they do?.
Let me talk about the Sukkot starting today..
Let's talk about the third festival,.
because it falls right on next week, right?.
Early October..
Leviticus 23, 33..
The Lord said to Moses,.
"Say to the people of God, Israelites,.
"on the 15th day of the seventh month.
"of the Lord's Feast of Tabernacles begins,.
"it lasts for seven days..
"Live in the Sukkot, in our booths, Tabernacles,.
"for seven days..
"All native-born Israelites are to live in the booths.
"so your children and your children's children,.
"your descendants, will know that I had.
"the Israelites live in booths.
"when I brought them out of Egypt,.
"through the wilderness, when they live in the desert,.
"when they live in the wilderness..
"I brought them out of Egypt..
"I am Yahweh, your God..
"I am the Lord.".
Now, as I said, it falls on October 9th to 16th this year..
And I told you last time, my wife and I,.
some years ago, we went to Israel.
for our two weeks vacation,.
and didn't know what festival fall in that two weeks..
So the first week, it was the Yom Kippur,.
the Day of Atonement, when they confess their sins,.
so they shut down all the whole country, right?.
I told you that..
And then the following week,.
it's a totally different mood, you know, on the streets,.
because it was the week of tabernacles, sukkot..
It was a high holiday in Israel..
It actually is probably one of the most important festivals.
that they celebrate..
And so we did experience one of their sukkots that year..

$^{281}$They call it sukkot..
In Hebrew, that's the word for tabernacle, the tent..
So during sukkot, we walk around the city,.
and I find it really, really interesting..
Just like a non-Chinese walks into Chinatown, right?.
Walks into the community of most of the Chinese people..
They probably would be amazed.
by how we celebrate some of the dates, right?.
How we go about our daily lives, right?.
I used to have a Caucasian professor..
He told me the first time he went to Chinatown.
when he was a kid..
So he went to Chinatown, and he looked at one of the stores..
And the storefront, there's a water tank,.
and there's some fish in there..
So as a kid, he find it really amazing.
that the fish, live fish, is on display on the street, right?.
He walk into the store, and he find it so sad,.
because people are killing the fish..
So he was amazed..
As a white person, he had never experienced.
something like that..
So kind of a same feeling that I had in sukkot..
Finding the Jews, they were building different kinds.
of tents on the streets and in the balcony..
So this is some of the pictures I took during sukkot..
They built this kind of tent they call sukkah..
So how do you build a sukkah?.
You need four by eight, some of the four by eight,.
and then you nail them together..
And then on top of it, you put some branches.
or you put some plants to cover it..
It's not very permanent, right?.
It's not very strong..
Actually, if the weather goes bad, it may fall..
So on the street, they build all kinds of sukkah..
On the balcony, they build some sukkah..
So this is what they do, because that's what the Torah says..
You have to celebrate in the sukkah,.
in the tent, for seven days..

$^{321}$So what do they do?.
They build some tents outside..
It has to be outside, not inside your house..
So they will dine under the tent..
They will talk under the tent..
They will drink under the tent..
Now, it's very interesting, right?.
So we went to Western Wall..
That is the day before Sukkot..
So I saw some workers..
They cut some of the branches from the palm tree,.
and then they were building sukkah..
Very interesting..
Building sukkah on Western Wall,.
you know, on the Western Wall Plaza..
The Western Wall is actually an open-air synagogue,.
so people are meeting there..
So they're building all these kind of things,.
and it looks like a hut,.
and it's very flimsy..
But during seven days, they will eat and drink.
and even sleep in the sukkah..
How many of you like camping trip?.
You know, go out to camping, maybe in Michigan.
or in some other national park..
How many of you enjoy camping?.
How many enjoy camping?.
You do?.
Okay, then we'll have one hand only..
Okay, okay, another..
Okay, all right..
So you don't like outdoor activities?.
Why not?.
Mosquitoes, right?.
Yeah, and then bugs in, you know, bathroom, right?.
Yeah..
All right, so you guys are too spoiled..
Urbanites..
Too spoiled urbanites..
I used to like camping, too..

$^{361}$I used to like camping..
But as I get older and older, I find myself,.
I can't do this anymore..
Yeah..
But you know, when you go camping, you know,.
it's kind of like living in a sukkah, right?.
You sleep in there, you eat in there,.
you play games inside the tent..
It's kind of fun, right?.
When you were young, right?.
(laughs).
This is what they do for seven days..
And if you're not handy, this can be a challenge to you..
How many of you are really handy here?.
You know, nails, woodworking, and all that..
How many of you are really handy?.
Or me, okay..
(audience laughs).
Yeah..
But when you think about sukkah,.
it really doesn't require you to be handy..
Because in the Jewish literature,.
when they talk about building a sukkah,.
they actually dictate that you build a flimsy sukkah,.
temporary sukkah..
It's okay if it falls..
Because sukkah does not mean to stand forever..
They are to be taken down..
They are temporary dwelling place..
Why?.
In fact, there's a story in one of the Jewish literatures.
that talks about two Jewish rabbis..
They were debating how strong.
or how permanent a sukkah should be..
And the story goes, two rabbis were traveling during sukkah..
So they were not home..
They were on a ship..
They were on a, you know,.
on the Mediterranean Sea..
During sukkah..

$^{401}$And they find that, okay, it's time in the year.
to celebrate sukkah..
So they need to build something, a sukkah, on the ship..
One of the rabbis,.
he chose to build his sukkah on the deck of the ship..
And then after he built it,.
during that night, wind came and blew down the sukkah..
And then the two rabbis were debating,.
"See, how could you build such a flimsy sukkah?".
Right?.
The wind came and then it fell..
And the rabbi who built the sukkah on the deck,.
he said, "I built the sukkah the way it's supposed to be..
"That's what our law says..
"It has to be fragile..
"It has to be flimsy.".
The Jewish literature, it's very funny..
You know, talk about sukkah so much,.
you have to build three walls at least..
And then you have to build your roof with branches only..
You know, all sorts of details,.
or building code for the sukkah..
Now, why did they do that?.
Because sukkah, the tent, the tabernacle,.
is a symbol of life..
Life is transitional..
Life is like a shadow..
Here today, but will be gone tomorrow..
That's what sukkah is all about..
Because when the Israelites, when they came out of Egypt,.
they spent 40 years in the wilderness..
There were no permanent houses in the wilderness..
There's no stone building in the wilderness..
All they live in was temporary tents, right?.
And life is always transitional, right?.
Don't you think so?.
Sukkot is a reminder every year.
for the people of God to remind them.
how short their life is, how transitional our life is..
Don't you think so?.

$^{441}$Life is short..
Time flies..
I went to a wedding not too long ago..
I met some of the students I used to be counselor with..
You know, now they grew up..
They got married, right?.
And I met some of the friends,.
you know, some of the students that I used to be with..
Now they grew up, right?.
They're working, or they have their own career..
Now, I met different people,.
"Oh, long time no see," right?.
"Wow, long time no see, long time no see.".
You know, that's one of the tunes.
that I play the whole night, right?.
"Oh, long time no see, long time no see.".
It becomes automatic, right?.
Yeah, whoever you see, "Oh, long time no see," right?.
Because I haven't seen them for a long time,.
especially in the past three years because of the pandemic..
I haven't seen those friends for a long, long time, right?.
So, and the other tune that we play, or I play,.
most of the time during the night, during the banquet,.
was that, "Wow, your kid is so big now..
"Wow, you," talking to the kids,.
"Wow, you're so tall..
"When I saw you last time, you were like, you know,.
"a tiny, tiny being, right?.
"Now you're so big," right?.
Then they, and then we update each other,.
and then I told them, you know,.
my son and my daughters are in college,.
and they say, "Wow, dang, I cannot believe that.
"because, you know, in my mind,.
"your kids are just kids," right?.
I said, "They're not kids, they're adults now..
"They drive, they drink, they're adults.".
Life is fleeting, life is so short..
That is one of the things we hear over and over again..
Time flies..

$^{481}$Sukkot, it's a reminder.
that your life is not permanent..
Don't make it permanent..
There's no need to make it permanent..
But, you know, it's very interesting..
Sometimes we spend most of our effort.
to make life permanent..
Try to make it permanent..
Try to make it look like it lasts forever..
Sukkot tells you, "Don't pretend..
"You can't do that..
"Life is not permanent, it's transitional.".
And how much we invest in our lives,.
how much effort, how much thinking.
have we invest into just planning our life,.
building our life, it seems like our Sukkot.
would last forever..
But no, it won't..
Life is always transitional..
You know, one of the scriptures.
the Jewish people would read during Sukkot.
in the synagogues is the book of Ecclesiastes..
Ecclesiastes..
If you haven't read this book,.
the book of Ecclesiastes in the Old Testament,.
you should read it..
It talks about life under the sun..
It's meaningless..
Because of death..
And one day, we're gonna die, right?.
And the whole book of Ecclesiastes is to remind us.
where should we find meaning in our lives..
Let me close with this last scripture..
The Lord said to Moses, "On the 15th day of seven months,.
"the Feast of Sukkot lasts for seven days..
"You live in Sukkah for seven days, all of you,.
"so that your children and your children's children.
"will know that you used to live in Sukkah,.
"in wilderness for 40 years..
"I rescue you..

$^{521}$"I provide for you..
"I guide your way..
"Your short life is under the hand of God.".
Well, a reminder for us, life is transitional..
And how much effort should we invest.
into something really eternal, right?.
Something really eternal..
C.S. Lewis says something like,.
do not, something like do not,.
do not stress out for things that are not eternal..
Do not be so stressed out about things that don't last..
Change our attitude..
You know, adjust our orientation..
What is life all about?.
Especially we have the life of Christ..
How should we invest our time in this short, short life?.
Invest into something that will last..
Now, what will last in the Bible?.
Let me tell you this, okay?.
What will last in the Bible?.
I can think of three things.
that will last forever in the Bible..
First, God himself..
God himself is everlasting..
All right?.
Second, is the soul, human soul..
Human life will last forever, eternal life, right?.
Three, the gospel..
Last forever..
Worth our effort, worth our investment..
Let's all pray together..
Father, we give you thanks for the reminder.
that we are in Christ,.
that we have eternal life and eternal hope..
And thank you for your grace..
Thank you for the gift of eternal life..
Remind us once and again, again and again,.
that this life is short..
This life will not last..
Yes, we are all like Israelites living in tents.

$^{561}$in the wilderness..
When wind came, everything would be gone..
Yes, life is transitional,.
but help us to anchor our life in eternity..
Help us to fix our eyes.
in something that will last,.
to preach your gospel,.
to build up your kingdom, to serve you..
Thank you for being our eternal God..
In Christ's name, we pray..
Well, thank you very much and we'll see you next week.
and we'll continue to talk about faith and time..
God bless..
\newpage



\section{5:1-5}
\label{sec:eW87A4hZtZA}
\textbf{算得上的人生 (創5\_1-5, 21-32) - 朱莉傳道}
\newline
\newline
連結: \href{https://youtube.com/watch?v=eW87A4hZtZA}{\texttt{https://youtube.com/watch?v=eW87A4hZtZA}} ~~~~ 語音日期: 2022-10-10
\newline
\newline
\hyperref[sec:IDFYoAmcEPE]{\small{< < < PREV SERMON < < <}}
~
\hyperref[sec:index_chronic]{\small{[返順時目]}}
~
\hyperref[sec:index_scriptual]{\small{[返順卷目]}}
~
\hyperref[sec:W7L4nTwMQOU]{\small{> > > NEXT SERMON > > >}}
\newline
\newline
$^{1}$等子妹平安.
相信大家都會認同.
在這兩三個星期最舉世矚目的新聞.
必然是英女皇以利沙白二世的離世.
無論在英國.
甚至在北美或香港.
都有很多人.
從不同的途徑.
來向她表達哀思.
我特別留意人們是怎樣懷念她.
並不是來數算她所擁有的財產.
是世界排名第幾.
或者她所擁有的金銀珠寶的富存.
是多麼的足夠等等.
是有這樣的報導.
但並不是主流.
或者強調在她的任內下.
英國怎樣國力鼎盛.
大國崛起等等.
而失實.
甚至可能是相反.
但很多人.
是去表達甚麼呢?.
是表達她本身是一個怎樣的人.
去強調她的待人接物.
甚至是講到她和寵物的關係.
最重要,最多人提起的.
是她怎樣看待她的身份.
怎樣鞠躬盡瘁.
去到最後.
今日想和大家分享的一段經文.
在《創世記》第五章.
《創世記》第五章.
因為比較長.
所以我們不會看完.
我們選取一些經文來看.
在《創世記》第五章.
我們先看一到五節.
之後二十一到三十二節.
《創世記》第五章一到第五節.

$^{41}$《亞當的後代》記載下面.
當神做人的日子.
是照著自己的樣式做的.
並且做男做女.
在他們被做的日子.
神似乎給他們稱他們為人.
亞當活到一百三十歲.
生了一個兒子.
形象,樣式和自己相似.
就給他起名叫「失德」.
亞當生失德之後.
又再世八百年.
並且生兒養女.
亞當共活了九百三十歲就死了.
第二節.
以諾活到六十五歲.
生了瑪吐撒拉.
以諾生瑪吐撒拉之後.
與神同行三百年.
並且生兒養女.
以諾共活了三百六十五歲.
以諾與神同行.
神將他取去.
他就不再世了.
瑪吐撒拉活到一百八十七歲.
生了拉密.
瑪吐撒拉生拉密之後.
又活了七百八十二年.
並且生兒養女.
以諾共活到七百七十七歲.
就死了.
二十八節.
拉密共活到一百八十二歲.
生了一個兒子.
給他起名叫「羅亞」.
說:這個兒子必為我們的操作.
和手中的奴婦安慰我們.
這操作奴婦是因為要話咒坐地.
拉密生羅亞之後.
又活了五百九十五年.

$^{81}$並且生兒養女.
拉密共活了七百七十七歲.
就死了.
「閃含亞佛」.
相信大家都會覺得.
這是一段很特別的經文.
為什麼呢?.
因為是家譜.
我不知道大家讀經的時候.
遇到家譜的時候會有什麼反應.
通常呢.
要不就完全跳過它.
就算不跳過.
很乖地看完它.
還是快快地看完就算了.
為什麼呢?.
因為我們覺得這些人和我們沒有關係.
讀名字也需要一番氣力的這班人.
和我們現在今時今日所過的生活又有什麼關係呢?.
所以我們今天上師.
從《創世記》第五章的家譜.
一齊去看人生.
要讀明家譜.
其實有兩個需要注意的地方.
第一是家譜的格式.
第二是上下文.
我們先從家譜的格式.
透過《創世記》第五章.
我們希望可以掌握一些技巧.
聖經裡面的家譜.
是絕對不少的.
我們以後讀到這些家譜的時候.
原來其實我們有些東西可以留意的.
有些東西我們可以嘗試去理解.
到底經文的意義在哪裡.
而不是快快讀過就算了.
所以第一我們要留意的是家譜的格式.
如果看這段經文你會留意.
家譜的格式正如我們普通日常生活裡.
或者現實生活裡的家譜一樣.

$^{121}$是非常regular的.
是很follow一個formula.
很formulative的記載.
在《創世記》第五章的formula是甚麼呢?.
某某人活到多少歲?.
生了A.
某某人生了A之後又活了多少年?.
並且生兒養女.
某某人共活了多少年?.
就死了.
要留意得到這個pattern.
差不多每一個entry都是重覆這個pattern.
我們去想想.
用這個pattern去表達甚麼呢?.
人生最重要.
透過這個家譜讓我們看到的竟然是甚麼呢?.
是傳宗接代.
他沒有其他highlight.
他人生的highlight是甚麼呢?.
就是生了兒子.
用兒子的出生成為那個人的defining moment.
很有趣.
傳宗接代.
原來猶太人都很像中國人.
很著重傳宗接代.
我們一會兒去解釋一下.
為何這件事這麼重要.
第二件事要留意的.
就是他無論壽命有多長.
我們剛才一直讀下去.
你都留意得到.
多少百,多少百歲等等.
無論壽命有多長.
每一個entry,每一個人.
他最後都有三個字.
就是「就死了」.
無論他有多長的壽命.
當然這是有特定的歷史因素.
所以這班人是這麼長壽命.
但他們無論多長壽命.

$^{161}$都難逃一個事實.
是甚麼事實呢?.
死亡的事實.
死亡從何而來呢?.
我們在看《創世記》第五章.
其實不是很久之前.
就在《創世記》第三章.
或者第二章開始提到.
當時神吩咐人不要吃禁果.
結果吃禁果帶來的後果是甚麼呢?.
在《創世記》第二章.
神親自這樣說.
說會必定死.
但當然第三章我們看到的是甚麼呢?.
早晚吃了禁果.
但好像沒事發生似的.
但在第五章.
就將死亡的事實記錄下來.
死亡的結果.
沒有避開.
結果一樣臨到.
人始終要為犯罪付上代價.
所以家譜沒有將這件事.
記錄下來.
甚至告訴我們.
其實在死亡面前人人平等.
不在於他活了多長.
他們每一個人的人生總有長短.
雖然他的短比我們的長很多.
但最後還是要面對死亡.
所以重點是甚麼呢?.
就是如何活.
如何活一個與別不同.
又精彩的人生呢?.
我們就要看在家譜固定的格式.
很普通的裡面.
有一些不普通的地方.
有些人與別不同.
有些人是破格的.
在這樣的家譜裡面.

$^{201}$特別是《創世記》第五章.
我們看到有些人.
就是在這個格式裡面.
跳了出來.
我們看兩個人.
第一是爾諾.
21至24節.
第一個 爾諾.
21至24節.
爾諾活到65歲.
生了瑪圖撒拉.
爾諾生了瑪圖撒拉之後.
與神同行300年.
並且生兒養女.
爾諾共活了365歲.
與神同行.
神將他取去.
他就不再逝.
所以在這個很普通的中間.
他有一些很普通的地方.
到底他有甚麼普通的地方呢?.
第一 兩次提到他與神同行.
即是與神有一個很親密.
很相交的關係.
如果你從希伯來文的文法裡看到.
你會留意到他那種與神同行.
不是一時興起.
「我今天是主日」.
「我上完教會心裡很火熱」.
「我今天要與神同行」.
並不是這樣的.
他那種表達.
爾諾是一種很慣常.
很恆常.
一個習慣地去發生的一件事.
他與神來對同行.
但你必須要留意.
他與神同行不代表.
他是在深山裡居住.
他甚麼事都不做.

$^{241}$只是做屬靈的事.
是否這樣呢?.
有首詩歌.
不知道大家有沒有印象.
我們年輕時經常唱的青年聖詩裡.
有一首歌叫「與神同行」.
「與神同行我就遠離世界」.
是否這樣呢?.
當然那首歌有它的意思.
我們就不去討論了.
但他與神同行的同時.
要留意的是甚麼呢?.
他保持著那些正常的事.
他一樣生兒養女.
他仍然是生了他的兒子馬托薩拉.
成為他的定點時刻.
所以與神同行並不代表離地.
不代表只是做屬靈的事.
他仍然很踏實地過他的生活.
這是第一.
第二,他最後的三個字.
不是「就死了」.
他最後的三個字是甚麼呢?.
「他就不再逝了」.
或者是「他就不再了」.
他沒有跟隨那個格式裡.
那個結局裡的特定三個字.
我們不知道為甚麼.
神最後將他接去.
或者從其他經文我們看到.
的確神將他接去.
因為他討神的歡喜等等.
我們不確定為甚麼會這樣.
但無論如何.
我們看到的是.
以諾過了一個很不一樣.
很破格的人生.
這是第一個.
破格的人.
第二,一個破格的人生是誰呢?.

$^{281}$是羅亞28至32節.
我們看28至32節.
「拉麥活到一百八十二歲.
生了一個兒子.
給他起名叫羅亞.
說這個兒子必為我們的操作.
和手中的奴婦安慰我們.
這操作奴婦是因為瑤華咒助地.
拉麥生羅亞之後又活了五百九十五年.
並且生兒養女.
拉麥共活了777歲就死了.
羅亞五百歲生了閃含亞忽」.
又是了.
他是有很不同的地方.
但我們要看.
他不是沒有了那些普通的元素.
他全部都在這裡.
即是甚麼時候活到多少歲.
生了誰.
而且還要加生兒養女.
然後活到多少歲等等.
但是.
羅亞有特別的地方是甚麼呢?.
第一,他的命名很特別.
在家譜裡面.
沒有其他人像他一樣解釋他的名字為甚麼.
羅亞的意思是甚麼呢?.
正如他爸爸在29節這樣說.
「這兒子必為我們的操作和手中的奴婦安慰我們」.
即是帶來安慰.
羅亞的意思是帶來安慰的意思.
但是從哪裡帶來安慰呢?.
為甚麼需要安慰呢?.
29節下半節這樣說.
「這操作奴婦是因為瑤華咒助地」.
原來是從咒助中帶來安慰.
又將我們帶回伊甸園.
人類的犯罪現場那裡.
無論是死亡.
或者這些操作的奴婦.

$^{321}$其實都是人類犯罪的惡果.
但是要留意.
這裡提到的是神咒助地.
在《創世記》前面第三章亦是這樣說.
神不是咒助人.
是咒助地.
以致要人汗流滿面.
才能糊口.
即是工作的奴婦.
所以加普將這兩樣東西.
在我們面前.
一個是死亡的威脅.
一個是工作的奴婦.
但原來有些人是可以跳脫出來的.
我們看到以洛其中一個.
羅亞原來亦都可以帶來安慰.
從工作的奴婦中帶來安慰.
其實我不需要再講下去.
因為我相信我們都認同這個事實.
無論《創世記》第五章.
或者我們今時今日.
我們面對我們的人生.
一樣其實我們都面對著這些威脅.
死亡的威脅.
以致例如疾病.
為甚麼這個疫症.
我們要做一個大事.
為甚麼我們要這樣去面對.
因為會死.
因為死亡率的確是有存在.
另一方面.
我們都面對非常日常.
就是我們面對工作的奴婦.
其中一個我經常很喜歡引用的例子.
因為我覺得這個是非常諷刺的.
就是在很多全球的排名.
或者調查裡面.
香港往往是名列前茅的.
但香港到底有甚麼這麼厲害.
是拿世界第一的呢.

$^{361}$第一是工時最長.
全球裡面香港人的工時是最長的.
平均每個星期花50個小時.
同時香港亦是房價最貴的地方.
即是你工作這麼多.
然後你用你大半生所勞苦的結果.
就是為了一層樓.
生活費亦都是很貴.
但最諷刺的是甚麼呢.
原來香港人是全球最長命的.
其中一個群體.
所以非常諷刺.
你這麼長命原來就是為了工作很長.
以至去供你那層樓.
難獨難怪.
香港人亦是最不快樂的其中一個群體.
但聖經裡面可以告訴我們.
原來可以不一樣.
羅亞有甚麼不同的地方.
除了它的名字是不同之外.
它最後到32節.
「羅亞五百歲生吊閃含阿忽」.
這是第五章的最後一節.
是否代表羅亞不用死呢.
是否就是「The End of Story」呢.
但如果大家熟悉聖經的話就知道.
當然不是「The End of Story」.
羅亞是否沒有「就死了」這三個字呢.
原來羅亞的故事從來不是停在第五章.
第五章只是提到他三個兒子就停在那裡.
而且羅亞也很特別.
他提了三個兒子的名字.
為甚麼這個重要性有機會再講.
但去到第六章提到甚麼呢.
第六章第九到第十節.
《創世記》六章九到十節這樣說.
羅亞的後代記在下面.
羅亞是個二人.
在當時的世代是個完全人.
羅亞與神同行.

$^{401}$羅亞生了三個兒子.
就是「閃含阿忽」.
你又看到羅亞的後代記在下面.
但又是更加深奧.
不單提到他的兒子.
還提到羅亞與神同行.
他是一個二人.
在當時的世代是一個完全人.
有更多的細節提到他的人生.
那他是不是不用死呢?.
不是的.
第六章提到他與神同行.
後面繼續下去.
要去到第幾章才記載他故事的結束呢?.
就是第九章.
《創世記》第九章二十八,二十九節.
洪水而後.
羅亞又活了三百五十年.
羅亞共活了九百五十歲.
就死了.
有否留意.
又回到家譜的格式.
但中間發生什麼事呢?.
第五章的結束只提到羅亞生了三個兒子.
第六,七,八到第九章.
才提到羅亞故事的結束.
換句話來說.
羅亞的進入比其他人到場更多.
他在家譜裡.
他不只是兩句或三句的說話.
他有幾章的聖經去記載他的人生是怎樣過.
中間最重要的是什麼呢?.
第六,七,八,九記載的是什麼呢?.
洪水.
第九章二十八節.
洪水而後.
原來洪水的記載其實是夾在家譜裡.
而他在洪水裡.
羅亞最重要的貢獻是什麼呢?.
他按照神的吩咐做方舟.

$^{441}$是一個很不容易的世代.
當時那個世代是滿了強暴的世代.
但他是一個義人.
他是一個完全人.
他與神同行.
在日常的生活裡.
他形形異異.
但是他承擔了神給他的一個特別的任務.
是什麼呢?.
就是做方舟.
面對人的譏笑,怒罵.
他過了一個破格的人生.
承擔了神給他特別的使命.
所以原來可以這樣.
在一個不容易.
對我們來說也是一個提醒.
到底在日常生活裡.
在我們工作的路虎裡.
我們怎樣可以看見神的同在.
怎樣去承擔神給我們的使命.
以致我們都可以過一個破格的人生呢?.
這個可能我們有機會回到家裡.
聽完這篇道.
慢慢去思考的一個方向.
有兩個閱讀家譜需要注意的地方.
第一是格式.
我們看到格式是什麼.
特別有些人是破格的.
第二我們可以留意從家譜裡面是什麼.
上下文.
我們回到第五章的一開頭.
我們留意一下那個家譜的引言是怎樣的.
《創世記》第五章一到三節.
亞當的後代記載下面.
然後有個括弧.
「當神做人的日子是照著自己的樣式做的.
並且做男做女.
在他們被做的日子.
神賜福給他們.
稱他們為人」三括弧.

$^{481}$是三節.
亞當活到一百三十歲.
留意那個括弧其實是後加的.
是和合本加的.
你留意英文是沒有這樣.
可能是翻譯的時候覺得.
這個好像是一個加插的內容.
所以就將括弧這樣加了進去.
所以如果你將括弧.
拿走不讀括弧.
第一節亞當的後代記載下面.
第三節亞當活到一百三十歲.
其實是很順的.
但是就是這個括弧.
就是這個括弧才是家譜的重點.
非常非常重要.
為什麼呢?.
因為亞當的後代記載下面.
我們就會想.
亞當的家譜的源頭是誰呢?.
既然說亞當的後代.
那麼亞當後代的源頭.
那個最老祖宗是誰呢?.
當然是亞當.
但是聖經就是這麼奇妙.
它將那個源頭很清楚地告訴我們.
我們再讀一次括弧它怎麼說.
括弧裡面的內容是.
「當神造人的日子.
是照著自己的樣式造的.
並且造男造女.
在他們被造的日子.
神賜福給他們.
稱他們為人」.
原來真正亞當家譜的源頭.
不是亞當自己.
是上帝.
因為是上帝所造的.
是祂設計的.
所以不斷強調是神造人的.

$^{521}$而且如果我們回到《創世記》那裡說.
祂是按照祂自己的形象,樣式去造人的.
我們回到《創世記》.
第一章.
我們去看看.
既然是神造我們的.
而且是祂精心設計去造我們的.
祂不單有版權.
祂有話事權.
我們要怎麼過我們的人生.
其實要問我們的設計者.
正如最近iPhone又出了.
出新的iPhone了.
到底iPhone應該怎麼用呢?.
其實設計者或工程師是最清楚的.
我不會用iPhone去煮飯的.
不會用iPhone去按釘的.
因為它不是設計成這樣.
它的創造,它所做的.
是有特定的功能.
當你去發揮它的功能.
按照它的設計去用的時候.
就能夠將它的功能.
它的效率去發揮到最極致.
當神造人的時候是怎樣呢?.
《創世記》第一章26至28節.
「神說我們要照著我們的形象.
按照我們的樣式造人」.
這些字句很熟悉.
因為在第五章有出現的.
「使他們管理海裡的魚.
空中的鳥.
地上的生畜和泉地.
並地上所爬的一切昆蟲.
神就照著自己的形象造人.
乃是照著他的形象造男造女.
神就賜福給他們」.
這些用字在第五章有重複的.
「又對他們說要生養眾多.
遍滿地面治理這地.

$^{561}$也要管理海裡的魚.
空中的鳥.
及地上各樣行動的活物」.
所以清楚地告訴我們.
在《創世記》第一章.
神創造我們的時候.
原意就是要賜福.
怎樣能夠獲得神所本意.
要賜福的福氣途徑呢?.
神造人的時候.
祂的本意是什麼呢?.
原來祂在這裡有使命給我們.
祂不是只想玩玩.
造幾個人出來玩玩.
不是這樣的.
《創世記》第一章提到.
祂有使命給我們.
使命是什麼呢?.
「生養眾多遍滿地面治理這地」.
當然,祂有當時的環境因素.
所以要求早日生養眾多.
到現在我們要怎樣理解.
這是另一個話題.
但神給亞當夏娃的使命是什麼呢?.
是非常清晰的.
是什麼呢?.
「生育」.
而第五章.
我一開始提到.
第五章這個家譜的重點提什麼呢?.
「生育」.
原來第五章家譜不斷重複提到的是什麼呢?.
不單是生一個兒子.
還是生兒養女去強調這些東西.
為什麼要強調這些東西呢?.
你可能會說.
「不是的,家譜當然強調生育」.
「難道說其他東西嗎?」.
對不起.
家譜真的可以說其他東西.

$^{601}$所以我們對比另一個家譜.
剛好亦在第五章前面第四章.
有另一個家譜.
當然不是剛好.
是作者刻意安排.
你將兩個家譜放在一起.
你就會明白到底《創世記》第五章.
記這些東西的重點在哪裡.
我們看《創世記》第四章.
17至22節.
另一個家譜.
《創世記》第四章.
17至22節的另一個家譜.
17節.
「該人與妻子同房.
看妻子懷孕生了兒郎」.
有另一個家譜.
但這是「該人」的家譜.
「該人建造了一座城.
就按著他兒子的名.
將那城叫做爾洛」.
當然這個「爾洛」.
與我們第五章的「爾洛」不同.
「爾洛」生「爾拿」.
「爾拿」生「米胡雅利」.
「米胡雅利」生「馬圖薩利」.
「馬圖薩利」生「拉麥」.
「拉麥」娶了兩個妻子.
一個名叫「阿大」.
一個名叫「洗拉」.
「阿大」生「阿伯」.
「阿伯」就是「著張鵬」.
「牧養生畜之人的祖師」.
「阿伯」的兄弟名叫「尤伯」.
他是「一切彈琴吹簫之人的祖師」.
「洗拉」又生了「土伯」「戈隱」.
他是打造各樣銅鐵利器的.
我們回顧一下.
「戈隱」當然是亞當的長子.
但在第四章的前面提到.

$^{641}$他因為妒忌殺了他的弟弟「阿伯」.
所以「戈隱」就成為了一個分離的入口.
他的後代.
他的線條就記載在第四章.
然後亞當的另一個兒子「瑟特」.
他的後裔就記載在第五章.
是分開的.
但同樣兩個都是家譜.
同樣在四,五章相比.
即是競爭對比.
是可以對比他們的.
你會發現他們同樣都是生育.
都有提到子女等等.
但「戈隱」的家譜是怎樣的呢?.
是有各種成就的.
建一座城.
有甚麼彈琴吹簫.
做銅鐵利器等等的佐詩.
如果你將它放到現代.
你會留意到甚麼?.
你會希望在第四章的家譜.
我都希望我爸爸建一座城.
然後城的名字是我的名字.
我都希望我的子女.
彈琴吹簫等等.
銅鐵利器等等.
即是律師,會計師,醫師等等的專業人士.
我們都會希望是甚麼?.
其實我們希望會在第四章的家譜.
而不是第五章.
好像甚麼都沒有.
就是只講生育.
的確「戈隱」的家譜充滿了各樣的成就.
但其實我還未讀完.
如果你繼續讀下去.
《創世記》第四章23至24節怎樣說?.
「拉麥對他兩個妻子說」.
我再說「拉麥」是第四章的「拉麥」.
不是第五章的「拉麥」.
「拉麥對他兩個妻子說:.

$^{681}$亞大洗拉聽我的聲音.
拉麥的妻子誓聽我的話語.
壯年人傷我我把他殺了.
少年人損我我把他害了.
若殺戈隱遭報七倍.
殺拉麥必遭報七十七倍.
戈隱殺了他的兄弟.
但這裡說「拉麥」他的兒孫.
是比戈隱更加殘暴.
是,第四章的家譜充滿了各樣的成就.
但同時也充滿了很多的強暴.
那個罪惡是加倍.
這樣在家譜裡蔓延.
最重要的是什麼呢?.
留意一下.
原來第四章的家譜.
是沒有記下他們的年歲.
第四章的家譜是沒有記載這些人的年歲.
而聖經也沒有再提這些人.
聖經有再提戈隱.
但每一次再提都是負面的.
第四章這群人.
是算不上,計不上的人生.
是不蒙神紀念的人生.
我的意思不是說.
亞當的後代沒有什麼成就.
所以只記載他們的年歲.
彷彿有些事情要記下.
不是.
為什麼亞當的後代.
瑟特的後代第五章要記載這件事呢?.
其實是記載亞當的這條線.
這條瑟特的線是如何完成神的使命.
一代又一代的人去完成神當初給他的先祖亞當的使命.
生養眾多,遍門地面.
所以第五章不是說亞當的後代,瑟特的後代沒有什麼成就.
只是將他們人生最重要的那件事.
如何去滿足神給他們的使命這一點.
去記下來.
而事實上第五章他們每一個人的年歲.

$^{721}$都記得清清楚楚.
他們是被計算的.
他們是算得上的人生.
是夢神紀念的人生.
英女皇以利沙白二世.
中年96歲.
其實是今年年初才慶祝她登基70週年.
她也破了英國本身的紀錄.
是在位最長的君主.
甚至世界的近代史裡.
她也是在位最長的君主.
一個擁有這麼高財富,權並地位的人.
在她今年慶祝登位70週年.
致國民的一封公開信上怎麼說呢?.
她這樣說:.
「As we mark this anniversary」.
她說她今年去紀念這個週年.
「It gives me pressure to renew the pledge I gave in 1947」.
她說這是一個好的機會.
重新一次重新她的誓言.
在1947年當時她登基的那個誓言.
什麼誓言呢?.
「That my life will always be devoted to your service」.
我的人生是會貢獻來服侍她的國民.
最後那封信的署名.
寫著「Your servant, Elizabeth」.
寫著「你的僕人」.
或者「你的婢女,以利沙白」.
當然你可以說這是門面的東西.
這是官方的公關.
但某個程度上.
這也提醒我們.
我們未必有她這麼長命.
我們也應該不會像她有這樣的地位.
今受全世界上至國王.
元首下至平民的尊敬.
其實我們也一樣可以成為至高君王的僕人.
《賽世紀》第五章.
是她們以生育.
以完成神的使命.

$^{761}$成為她們人生的方向和重點.
她們不僅被神紀念連綏.
她們的名字也再一次出現在聖經裡面.
在《路加福音》第三章.
大家自己回去看.
《路加福音》第三章.
耶穌基督的家譜.
《賽世紀》第五章的每一個名字.
每一個名字都再次出現在.
《路加福音》第三章.
耶穌基督的家譜.
到底我們是哪個僕人呢?.
我們有沒有時時刻刻重新宣誓.
去奉獻我們的生命.
給哪個的服務呢?.
我們是否渴望有一天.
聽到我們的主人跟我們說.
有良善,有忠心的僕人.
可以來享受跟僕人一起坐寂的快樂呢?.
我們一起祈禱.
撐下去.
求你提醒我們.
雖然我們是面對很多的不確定性.
亦都是一個很不容易的世代.
但主幫助我們.
讓我們學校怎樣去做一個與神同行的人.
並不是每個人都要奉獻一生去服侍.
做傳道人,做使徒等等.
但在營營役役的人生裡面.
在勞苦的工作.
為了戶口而勞苦的人生裡面.
主要是怎樣與你同行.
以致能夠過得有意義.
甚至看見你給我們的特別任務.
去承擔你給我們的使命.
在這個人生裡面.
來服侍你,服侍人.
主幫助我們,提醒我們.
讓我們的人生在你裡面.
是被算得上.

$^{801}$是成就你福音的使命.
以致我們有一天見到你的時候.
能夠被你稱為良善有中心的僕人.
求你聽我們在你面前的禱告.
奉耶穌基督的聖名.
\newpage



\section{馬太福音 4:1-11}
\label{sec:W7L4nTwMQOU}
\textbf{曠野的吶喊 (馬太福音4\_1-11) - 張麗仙傳道}
\newline
\newline
連結: \href{https://youtube.com/watch?v=W7L4nTwMQOU}{\texttt{https://youtube.com/watch?v=W7L4nTwMQOU}} ~~~~ 語音日期: 2022-10-12
\newline
\newline
\hyperref[sec:eW87A4hZtZA]{\small{< < < PREV SERMON < < <}}
~
\hyperref[sec:index_chronic]{\small{[返順時目]}}
~
\hyperref[sec:index_scriptual]{\small{[返順卷目]}}
~
\hyperref[sec:yX0mME__Hqg]{\small{> > > NEXT SERMON > > >}}
\newline
\newline
馬太福音 4:1-11
\newline
\begin{longtable}{cl}
\hline
\hline
章節 & 經文 (和合本修訂版)\\
\hline
4:1 & \begin{tabularx}{0.7\textwidth}{X} 當時，耶穌被聖靈引到曠野，受魔鬼的試探。 \end{tabularx} \\ \\ \relax
4:2 & \begin{tabularx}{0.7\textwidth}{X} 他禁食四十晝夜，後來就餓了。 \end{tabularx} \\ \\ \relax
4:3 & \begin{tabularx}{0.7\textwidth}{X} 那試探者進前來對他說：「你若是神的兒子，叫這些石頭變成食物吧。」 \end{tabularx} \\ \\ \relax
4:4 & \begin{tabularx}{0.7\textwidth}{X} 耶穌卻回答說：「經上記著：『人活著，不是單靠食物，乃是靠神口裡所出的一切話。』」 \end{tabularx} \\ \\ \relax
4:5 & \begin{tabularx}{0.7\textwidth}{X} 魔鬼就帶他進了聖城，叫他站在聖殿頂上， \end{tabularx} \\ \\ \relax
4:6 & \begin{tabularx}{0.7\textwidth}{X} 對他說：「你若是神的兒子，就跳下去！因為經上記著：『主要為你命令他的使者，用手托住你，免得你的腳碰在石頭上。』」 \end{tabularx} \\ \\ \relax
4:7 & \begin{tabularx}{0.7\textwidth}{X} 耶穌對他說：「經上又記著：『不可試探主—你的神。』」 \end{tabularx} \\ \\ \relax
4:8 & \begin{tabularx}{0.7\textwidth}{X} 魔鬼又帶他上了一座很高的山，將世上的萬國和萬國的榮華都指給他看， \end{tabularx} \\ \\ \relax
4:9 & \begin{tabularx}{0.7\textwidth}{X} 對他說：「你若俯伏拜我，我就把這一切賜給你。」 \end{tabularx} \\ \\ \relax
4:10 & \begin{tabularx}{0.7\textwidth}{X} 耶穌說：「撒但，退去！因為經上記著：『要拜主—你的神，惟獨事奉他。』」 \end{tabularx} \\ \\ \relax
4:11 & \begin{tabularx}{0.7\textwidth}{X} 於是，魔鬼離開了耶穌，立刻有天使來伺候他。 \end{tabularx} \\ \\
[1ex]
\hline
\hline
\end{longtable}
$^{1}$每個星期早上回來.
我都會有一個很遺憾的事.
就是我不能夠跟每一位弟兄姊妹.
個別打個招呼.
因為時間匆匆.
同時網上也是.
沒辦法可以有這樣的時間.
跟大家打個招呼個別.
但是在這裡.
所以就趁這個機會.
跟每一位弟兄姊妹.
無論是現場的還是網上的.
願你平安.
我不知道過去幾個星期.
你聽到的時候.
有沒有一個這樣的感覺.
就是看著台上的Mark Sir.
他好像變黑了.
又變尖了.
可以說嗎.
看他站出來.
老實說.
其實我磅一磅自己.
都變輕了三磅.
磅可以騙人的.
Mark Sir沒有說.
磅可以騙人的.
但是褲子是騙不了人的.
褲子鬆了.
我身邊有幾位弟兄姊妹.
都告訴我.
他又變輕了三,四磅.
這些全部都是去聖地的團友.
大家不要誤會.
我們好像沒得吃.
有圖為證.
我們有時吃得也不錯.
但是說真的.
去聖地真的有點捱.
體力上的捱.

$^{41}$但是又捱得挺開心的.
所以輕幾磅.
其實我們當拿獎而已.
當初想到去聖地這個旅程的時候.
的確不單是我身邊很多團友都是.
最擔心的就是.
我捱不捱得到呢.
我們捱不捱得到呢.
特別想起一個地方.
我自己就最怕.
就是曠野.
因為那裡是很曝曬.
好像很無水.
還有很乾旱的地方.
究竟能不能走得到呢.
不過可以告訴各位.
在這一次去完聖地.
我一個最深的體會.
就是曠野.
我記得車一開到.
猶大曠野的時候.
我才知道原來曠野是這麼宏偉的.
你有沒有想過.
曠野很宏偉的.
很壯觀的.
給大家看看幾幅.
這個是昆蘭山洞.
曠野是一個這樣的地方.
不過這樣.
坐在冷氣裡面.
看出去是挺好看的.
好像我們現在這樣看.
是挺好看的.
挺宏偉的.
但是如果大熱天時走在上面.
是很要命的.
是很熱的.
是挺大的挑戰.
但是正如我剛剛說.
這一次去聖地.

$^{81}$讓我對曠野.
都有一個很新的改觀.
我回來.
如果你問我.
這一次聖地.
你最深印象在哪裡呢.
我可以告訴你.
就是曠野.
很多時候.
回來我還做了很多長夢.
都是曠野.
發了兩個星期.
為什麼這麼深感受呢.
不單是因為宏偉壯觀.
更加是我想起聖經.
想起我們人生.
我們今天的人生.
我突然間看到.
神原來藉著曠野.
對我們給了很多重要的人生啟示.
所以忍不住.
我就定了今天的講題是.
曠野的吶喊.
曠野大聲跟我們說話.
就是一個寸草不生.
寂靜無聲的曠野.
他真的向我們大聲說了.
最低限度.
我看到兩個人生的功課.
第一.
我們原來是一無所有.
好像曠野一樣.
我們是靠神的恩典而活的.
你有沒有留意到.
在舊約和新約聖經.
其實曠野是扮演很重要的角色.
經常都出現曠野.
你想想也是.
曠野是神祝福和磨練他子民的.
一個很重要的地方.

$^{121}$我們也很熟悉.
當日神拯救他的子民.
逃離埃及奴隸的生活.
前往應許之地.
中間走的就是曠野.
曠野什麼都沒有.
在這裡你找什麼吃呢.
曠野沒有光.
沒有民居.
你在哪裡有光照著你呢.
特別是在晚上的時候.
照著你可以行動.
但是就是神的恩典托著他們.
神賜嗎哪.
賜鵪鶉給他的子民做食物.
神賜雲柱,火柱.
好像大光燈一樣帶著他們前行.
要他們明白第一件事就是.
你們知道嗎.
其實你們本來就像曠野一樣.
一無所有.
是靠我的恩典.
你才可以存活.
這個風貨其實不是容易學的.
說就容易.
所以我們也看到.
聖經記載.
以色列民.
他們很久也學不到.
他們時時一不順利.
一有很難捱的日子.
他們就盲怨.
他們盲怨說.
我情願回到埃及做奴隸.
也比去曠野好.
他們又會願意去做偶像來敬拜.
代替了神.
代替了耶和華.
沒有水喝的時候.
他們要一再埋怨.

$^{161}$埋怨摩西,埋怨神.
他們越不相信神.
神就要他們越走長一點的曠野.
結果他們走了四十年.
真是不可思議.
而且不信的一代.
神不讓他們進入應許之地.
這是他們的功課.
但是大兄姐妹.
我們其實也要學這個功課.
這個功課是不容易的.
我們本來也像曠野一樣.
我們赤裸裸來到這個世界.
我們沒有什麼.
我們沒有.
我們一切都是神供應的.
不過我知道.
其實這樣說也很難接近.
因為一向都這樣想.
什麼我一無所有.
有沒有搞錯.
我有頭腦.
我有才能.
我一生用我雙手努力.
而賺取到我今天的成就.
沒錯.
也是.
我不可以否定這些成就.
不可以否定我們的能力.
但是不要忘記.
誰讓你頭腦.
誰給你才能.
甚至誰在你最難熬的日子裡撐過你.
我們信你的日子.
我們很難明白這個道理.
我們是靠神的恩典而活.
這是道理.
但是當一個大能打過來.
一個大病.
一個感情的衝擊.

$^{201}$家庭的破裂.
事業失敗.
政治動盪.
甚至是一個疫情.
我們大家都深深體會到.
原來我們是這麼無助的.
原來我們一下子所擁有.
我們以為我們擁有的東西.
一下子就被拆除.
被奪走.
我們面前的就是抗爭.
回顧我的人生.
我想很多事情.
我走過這些日子.
突然間發覺自己.
是一無所有的.
我知道在我們當中.
弟兄姊妹.
你們不少都走過這些日子.
但是也在這些日子.
我們學會謙卑下來.
發現一個道理.
就是我們是一無所有.
是神的恩典托住我們.
甚至那個人幫我.
那個人幫我.
也是神派來的天使.
來幫我渡過困境.
有一個聖經學者說得很有意思.
說得很妙.
他說.
In wilderness you have nothing.
but you lack nothing.
在曠野你一無所有.
但是你又一無欠缺.
什麼意思呢.
其實這個聖經學者說的時候.
他是assume了中間有一句.
就是因為你所有的.
是神給你的.

$^{241}$就正如《申命記》所說的.
摩西對神的子民說.
這四十年.
是耶和華你的神與你同在.
故此你一無所有.
In wilderness you have nothing.
but in God you lack nothing.
這個就是我今天想跟大家分享的.
曠野的第一個吶喊.
就是我本來一無所有.
是靠神的恩典而活的.
但是一個這麼loud and clear的道理.
為什麼我們好像有.
很多時候都聽不到這樣呢.
為什麼很多人都不肯接受呢.
是因為我們人類的受敵魔鬼.
就是從伊甸園開始.
已經是想我們不要相信這個道理.
他扭曲事實.
告訴我們.
其實我們什麼都有的.
我們不用靠神的.
這個是他想我們相信的.
但是美妙的就是.
神的兒子耶穌.
他親自來到.
他這個世界.
他甚至走去曠野那裡受試探.
他身上魔鬼的謊言試探.
而他的成功.
正是我今天想跟大家分享的.
第二個曠野的吶喊.
第一個是我本來一無所有.
是靠神的恩典活著.
第二個曠野的吶喊是什麼呢.
是耶穌受試探告訴我們的.
以下我想跟大家很快閃.
看《馬太福音》第四章一到十一節.
大家都很熟悉的一段聖經.
是說到耶穌在曠野.

$^{281}$午食夜禁食四十天.
受魔鬼的試探.
我們希望從這個快閃那裡.
看看這個一二三的試探.
究竟耶穌交給了一個什麼武器.
來勝過魔鬼呢.
是你和我也很值得學習的.
我們看第一個試探.
是將石頭變成食物.
我們都很熟悉.
當時耶穌被聖靈引到曠野.
受魔鬼的試探.
他禁食四十晝夜.
餓死了.
試探人的魔鬼進來對他說.
如果你是神的兒子的話.
你可以吩咐這些石頭變成食物.
就是這麼簡單.
可以馬上變.
耶穌你這麼厲害.
耶穌回答說.
經上記著說.
人活著不是單靠食物.
乃是靠神考慮所出一切的話.
在《申命記》第六章那裡.
請大家留意這裡的字眼.
耶穌在曠野受試探.
不是意外來的.
看到嗎.
是被聖靈引到.
是神自己精心設計的一個重要行程來的.
記住耶穌當時是100\%人性.
當然他有100\%神性.
但在受試探的時候.
他有100\%人性.
他四十天不吃東西.
我們四天不吃東西行不行.
我就不行了.
四十天就夜不吃東西.
餓死了.

$^{321}$魔鬼就在這個時候闖入.
你將石頭變成食物吧.
說起來很吸引的.
馬上就有東西吃了.
但是有什麼問題呢.
耶穌將石頭變成食物.
有什麼問題呢.
耶穌不是做過很多神蹟.
做食物的神蹟.
餵飽四千人,五千人.
問題就是.
當時將石頭變成食物的這個行動.
是代表了耶穌聽了魔鬼的意思.
不聽父上帝的意思.
我們看上文.
耶穌受曠野試探之前.
是在約旦河受浸.
當時神在天上說.
「這是我的愛子,我所喜悅的」.
「你要聽他」.
神向這個世界,向這個靈界說.
「這是我的愛子」.
「今天我就將我的愛子帶到曠野」.
「聖靈帶他來曠野受魔鬼的試探」.
「是我的計劃來的」.
但是如果你一下子就引誘了耶穌.
將石頭變成食物.
就是耶穌沒有堅持到這個使命.
承受這個試探.
祂失敗了.
如果耶穌將石頭變成食物.
是失敗了.
祂是向魔鬼讓步.
但是耶穌沒有這樣做.
我們都知道耶穌沒有這樣做.
祂清脆利落.
用生命的術語說.
人沒有著不是單靠食物.
是靠神口裡面一切的話.
是靠我天父的旨意.

$^{361}$不是靠這些眼前的滿足.
星期二晚上十點鐘.
我在我家預備講章的時候.
我突然間聽到.
「砰砰砰砰」.
十聲.
大約十聲.
我以為有裝修工在裝修.
想想不要這麼晚裝修.
有沒有搞錯.
我馬上走到客廳看.
有很多弟兄姊妹來過我家.
都知道我樓下有一個十字路口.
我看到十字路口有一輛車.
停在那裡.
沒有動.
當時十點鐘因為已經很晚.
街上也沒有什麼人沒有什麼車.
那輛車在那裡橫.
不是橫.
總之在十字路的中間停了.
沒有動.
那一會有另一輛車在他旁邊.
我想那輛車是問他.
然後打電話.
因為很快就有消防車救火車.
又有八輛閃著藍燈的警車來到樓下.
我才知道.
看到的是一個目擊.
應該是一個目標的擊殺.
有人尋仇.
我看新聞沒有看得這麼清楚.
沒有講清楚.
但似乎就是這樣.
海外網上的朋友.
你不要怕.
我們不是想有這些事.
而是你聽到這裡不用為我們擔心.
我們真的見怪不怪.
只是求神保守我們.

$^{401}$我們也盡量小心.
這件事我為什麼要跟大家講呢.
跟魔鬼說.
你將石頭變成食物.
有什麼關係呢.
樓下發生的不是賣雞蛋賣魚蛋.
是子彈.
但是同時.
撒旦魔鬼可以將任何東西都造成我們的試探.
在這一刻.
他可以使我很害怕.
樓下的恐怖.
現在之一崗的恐怖治安.
大家都感受到.
都會使我們試探.
懷疑神.
你究竟有沒有聽我們禱告.
你究竟有沒有能力保護我們的城市.
你究竟愛不愛我們.
我們也可以這樣想.
但是耶穌的回應.
使我們明白.
人活著不單只是靠當前的食物.
或者是平安.
而是更加靠神更高的旨意.
靠祂的話.
因此無論警方怎樣也好.
我們不要失去這個眼光.
我們要學耶穌.
怎樣也要相信神.
尊崇神.
不會做一些違背祂旨意的話.
的事.
第二個試探.
魔鬼帶耶穌進聖城.
在聖殿的山上.
對祂說.
如果你是神的兒子.
你就跳下去吧.
因為經上記住.

$^{441}$主要為你吩咐.
祂的使者用手托著你.
免得你的腳碰在石頭上.
耶穌對祂說.
經上又說.
不要試探主你的上帝.
請大家留意這裡的字眼.
魔鬼帶耶穌去聖城.
殿頂.
這些都是代表了一些.
他們當時覺得是很重要的.
權柄之手.
你跳下去吧.
父上帝必然會叫天使托著你.
是什麼意思呢.
就是說.
你可以挑戰一下上帝的界線.
耶穌.
可以挑戰一下祂.
祂未必不行.
大家有聽到伊甸園的聲音嗎.
那個果子真的神叫你不要吃嗎.
未必.
這個是同出一轍的.
請注意.
這個試探裡面是擊中了.
耶穌同時祂人性的弱點.
權柄往往令我們飄飄然.
耶穌當時的心裡面.
可能也會閃過.
對呀.
父上帝會向我讓多少步呢.
祂會不會讓我這樣做呢.
這樣跳下去.
祂會不會托著我呢.
弟兄姊妹我們要提防.
這些試探同樣會臨到我們.
我們在金錢上.
時間.
性.

$^{481}$權.
自由.
我們在這一切上面.
很多時候都明知道.
神是有祂的界線的.
我們都想去測試一下那些界線.
但是這些界線.
往往帶來幻想.
接著就帶來什麼事.
一團糟.
令我們不能反過來.
使我們人生失控.
帶來很深沉的內疚.
痛苦.
無助.
是因為我們測試了神的界線.
耶穌回應一針見血.
祂不接受這個試探.
祂拆穿了魔鬼的方形.
祂說不要試探主.
你的上帝同樣是生命計的.
無論吸引多大也好.
耶穌看到真正權柄的只有父上帝.
我要堅持尊崇祂.
我要堅持相信祂.
無論眼前的吸引是多大.
請大家也記住.
耶穌這個試探目的不是想指控我們.
我們今天看的時候.
祂不是想用這個試探來指控我們.
而是想勉勵我們.
祂要勉勵我們就是.
祂讓我們知道.
100\%的人性的祂.
是完全明白我們的軟弱.
我們想跌下去.
我們想踩下去.
祂是完全明白這種軟弱.
但是祂會告訴我們.
我們在哪裡跌倒.

$^{521}$祂就在哪裡扶持我們.
祂的十字架已經告訴我們.
祂來到就是這個目的.
我們在哪裡跌倒.
只要我們願意承認.
祂就在哪裡扶持我們.
祂也提醒我們.
我們的仇敵魔鬼.
是時時不手軟地試探我們.
想我們越過神的界線.
想我們試探神.
使我們失敗.
但是我們要和耶穌一樣記住.
我們要相信和尊崇主上帝.
我們就必定能夠勝過魔鬼.
第三個試探.
魔鬼帶耶穌上一座山.
最高的山.
把世上的萬國和萬國的榮華指給祂.
對祂說.
如果你俯伏拜我.
我就將這一切都給你.
耶穌說.
撒旦,我看出你什麼意思.
你快點走.
因為經常說.
當拜主你的上帝.
但要事奉祂.
請大家也留意下字眼.
最高的山.
top of the world.
萬國和萬國的榮華.
全都給了你.
條件就是你要拜我.
這個條件是很厲害的.
魔鬼的試探明顯是升級了.
他不用客氣說.
你不用靠神的.
靠我吧.
我坐了神的位置.

$^{561}$你靠我吧.
榮華本身不是罪.
榮華就是富貴,有錢.
這些本身不是罪.
但當榮華的條件.
是要我們敬拜事奉祂.
取代我們的神的時候.
這個是很大的罪.
就是我們時時說的.
我們是敬拜偶像.
這一次我們去聖地牛.
我們踏足過很多大希律王.
他的榮華.
他的宏偉建築.
是宏偉豪華的宮殿行宮.
遍佈在耶路撒冷城.
曠野瑪薩大.
凱撒利亞港.
等等等等.
很多他的宏偉建築.
是很豪華的.
但我們看到.
這些宮殿,這些行宮.
其實就是一個神主牌.
巨型神主牌.
告訴別人.
我追求這些.
我已經背棄了神.
我不再以耶和華是我的主.
希律.
大希律就是這樣.
不過弟兄姊妹.
今天榮華絕對不是.
只是皇帝專有追求.
其實你和我.
都有很多機會去追求.
很容易.
我們都好像大希律那樣.
建了一個神主牌.
或者我們騙自己看不到.

$^{601}$但你身邊的人看到.
看到你拜誰.
看到誰在你心目中是最重要的.
魔鬼看到.
神更加看到.
你對祂的忠誠.
耶穌對於這個試探.
一心見血地說.
當拜你的主.
單要事奉祂.
耶穌看得出.
正如祂曾經說過.
你的財寶在哪裡.
你的心就在哪裡.
你以什麼為主.
你的心就在那裡.
你日思想得是那些.
你已經敬拜那一樣東西作你的神.
耶穌說.
單要敬拜主你的上帝.
唯獨事奉祂.
這個是祂對撒旦.
這麼厲害的試探的回應.
同樣是.
我要堅定相信.
和尊崇我的神.
你不要跟我說這些.
說完了.
三個試探.
耶穌在曠野.
禁食四十晝夜.
受的三個試探.
我剛剛說.
我們今天再快閃看的時候.
我們看看耶穌在這裡.
勝過試探的武器是什麼.
我相信大家.
心水清的都看得到.
第一個試探.
將石頭變成食物.

$^{641}$耶穌說.
你人活著不是單靠食物.
是靠副上帝.
耶穌的意思.
是堅持他對副上帝的相信和尊崇.
第一.
第二.
當魔鬼要試探耶穌.
在墊頂跳下去.
顯示自己是神的兒子的時候.
耶穌堅持不要試探上帝.
堅持不要試探他的上帝.
是要說.
我要堅持.
我對上帝的相信和尊崇.
第三個試探.
魔鬼要祂俯伏敬拜祂.
將萬國的榮華給祂.
耶穌堅持說.
拜我的主上帝.
單單事奉祂.
是表示祂堅持對副上帝的相信和尊崇.
所以我們都看到了.
耶穌面對三個試探.
祂成功的武器就是.
祂無論如何都堅持相信和尊崇副上帝.
耶穌這個武器正是我今天想和大家分享.
曠野對我們第二個的吶喊.
人生充滿試探.
但是我們仍然相信和尊崇神.
而且是堅定的.
祂藉著穿草畢生.
寂靜無聲的曠野.
向我發出這兩個吶喊.
我今天早上在這裡就想和你分享.
勉勵你.
無論我們的日子是風盛滿載.
或者是乾涸無水.
都要記住.
我本來一無所有.

$^{681}$我們是靠神的恩典而活.
面對人生的試探.
我無論如何都堅持相信和尊崇我們的神.
我們祈禱.
神啊.
我們在耶穌這個三個試探的面前.
我們真的無聲可說.
因為裡面也說到.
其實你是多麼地認識我們每一個人.
你多麼地知道我們的軟弱.
所以你親自來到曠野受這個試探.
你勝過撒旦的權勢.
你更加是因此鼓舞我們在軟弱無助當中.
我怎樣能夠有個眼光看到.
你是那一位掌管一切.
你才是那一位負上帝.
教我們不要再做旁觀者.
其實我們就是主角.
曠野上面我們是主角.
怎樣去學習在這個曠野路裡面.
更加認識你.
還有你的大愛.
周正宇.
詩歌說.
因為你的愛.
你為我們所受的變相.
使我們得意志.
心觸我們的心.
我們今天也會來到你面前領受祝餐.
再一次藉著曠野這些提醒.
我們感激你對我們的愛.
我們真的願意再一次在你面前說.
我們忠誠於你.
一生不變.
我們無論日子如何.
堅持相信你是我們.
是我們的神.
我們堅持相信和尊崇你.
多謝聽我們的禱告.
奉救主耶穌的名.

\newpage



\section{15:1-19}
\label{sec:yX0mME__Hqg}
\textbf{Lightening Our Yokes (Acts 15\_1-19) - Brian Ng}
\newline
\newline
連結: \href{https://youtube.com/watch?v=yX0mME--Hqg}{\texttt{https://youtube.com/watch?v=yX0mME--Hqg}} ~~~~ 語音日期: 2022-10-19
\newline
\newline
\hyperref[sec:W7L4nTwMQOU]{\small{< < < PREV SERMON < < <}}
~
\hyperref[sec:index_chronic]{\small{[返順時目]}}
~
\hyperref[sec:index_scriptual]{\small{[返順卷目]}}
~
\hyperref[sec:R_BYGFCmbDs]{\small{> > > NEXT SERMON > > >}}
\newline
\newline
$^{1}$- Thank you for having me back here..
Thank you for having me back here..
The last time I was here was back in May,.
so it's been a little while,.
but the fact that I'm back here means that.
it was okay last time, so thank you for having me..
The last time I was here,.
I spoke on the topic of an emotive God,.
and it was about how God encourages us.
to express the full range of emotions.
that he also expresses, that he wants us to also express..
And this time around, Pastor Mark was nice enough to,.
again, extend freedom for me to choose what to preach on..
And usually I go through a process of prayer and meditation.
to figure out what it is God might be pressing.
on my heart to share,.
and I happened to be going through the book of Acts.
for my own devotion..
And at the time when I received the invitation,.
I was going through the passages that actually.
were going through today,.
and I just felt really strongly that this was the message.
that the Lord wanted me to share today..
So,.
if any of you have visited Thailand,.
you've probably done the very tourist thing.
of riding one of those elephants..
I don't know, anybody ever gone to Thailand here?.
Yeah?.
Okay..
I did many years ago,.
and when we see these animals in pictures and videos,.
most of our lives, we don't really get the true sense.
of how just amazingly powerful and large they are..
These elephants can lift logs with their trunks,.
they can crush cars with their feet,.
and yet somehow they submit.
to these very fragile human masters..
I always wondered how such big animals can be so compliant,.
and I eventually learned years later.

$^{41}$how they actually train them..
The way that they train them is that when they're very young,.
they chain them or they bind them,.
so when they even try to break free, it hurts them..
They're in a lot of pain.
if they try to escape their constraints..
And so, over the years, as they get older and older,.
these human masters start to take away these devices,.
but the elephants are so trained.
to think that they're trapped.
that they continue to think that they're bound,.
and so they operate in submission as a result..
The elephant imagery resembles.
how sometimes we can be bound spiritually.
by our early and external influences..
Whether it's from our culture, society,.
or our religious context,.
we get stuck and trapped in these artificial limitations.
that we put on ourselves.
that might not even be biblically prescribed..
Today's message is actually about that..
It's about reflecting on what maybe we have in our lives.
that are unnecessary spiritual restrictions.
that we place on ourselves..
And today's, oh, sorry, I forgot..
Yeah, so this was the Emotive God,.
and I was thankful to see.
that Pastor Mark actually reviewed it..
I somehow was on UB's website in August,.
and I was like, oh, he found a better graphic than me..
(audience laughs).
And here's one of the elephants, right,.
that you would see in maybe Thailand or Burma..
And today's passage is from Acts 15, one to 21..
I will read it, but it's also on the screen..
"Certain people came down from Judea to Antioch.
"and were teaching the believers..
"Unless you are circumcised,.
"according to the custom taught by Moses,.
"you cannot be saved..

$^{81}$"This brought Paul and Barnabas into sharp dispute.
"and debate with them..
"So Paul and Barnabas were appointed,.
"along with some other believers,.
"to go up to Jerusalem to see the apostles and elders.
"about this question..
"The church sent them on their way,.
"and as they traveled through Phoenicia and Samaria,.
"they told how the Gentiles had been converted..
"This news made all the believers very glad..
"When they came to Jerusalem,.
"they were welcomed by the church.
"and the apostles and elders,.
"to whom they reported everything God had done through them..
"Then some of the believers who belonged.
"to the party of the Pharisees stood up and said,.
"The Gentiles must be circumcised.
"and required to keep the law of Moses..
"The apostles and elders met to consider this question..
"After much discussion, Peter got up and addressed them,.
"Brothers, you know that some time ago,.
"God made a choice among you.
"that the Gentiles might hear from my lips.
"the message of the gospel and believe..
"God, who knows the heart,.
"showed that he accepted them.
"by giving the Holy Spirit to them,.
"just as he did to us..
"He did not discriminate between us and them,.
"for he purified their hearts by faith..
"Now then, why do you try to test God.
"by putting on the necks of Gentiles.
"a yoke that neither we nor our ancestors.
"have been able to bear?.
"No, we believe it is through the grace.
"of our Lord Jesus Christ.
"that we are saved, just as they are..
"The whole assembly became silent.
"as they listened to Barnabas and Paul.
"telling about the signs and wonders.

$^{121}$"God had done among the Gentiles through them..
"When they finished, James spoke up..
"Brothers, he said, listen to me..
"Simon has described to us how God first intervened.
"to choose a people for his name from the Gentiles..
"The words of the prophets are in agreement.
"with this, as it is written..
"After this, I will return and rebuild.
"David's fallen tent..
"Its ruins I will rebuild, and I will restore it,.
"that the rest of mankind may seek the Lord,.
"even all the Gentiles who bear my name,.
"says the Lord who does these things,.
"things known from long ago..
"It is my judgment, therefore,.
"that we should not make it difficult.
"for the Gentiles who are turning to God..
"Instead, we should write to them,.
"telling them to abstain from food polluted by idols,.
"from sexual immorality,.
"from the meat of strangled animals, and from blood..
"For the law of Moses has been preached in every city.
"from the earliest times,.
"and is read in the synagogues on every Sabbath.".
The passage begins on the backdrop.
of Paul and Barnabas' missionary journey in Acts 13 and 14..
They travel from Antioch to Cyprus,.
making their way from the eastern coast of the island.
to the western coast,.
where they then travel to the mainland of modern-day Turkey..
They mimic Jesus' strategy.
by first targeting the Jewish communities.
in these areas where they already have.
the religious background of the Old Testament..
They actually mimic Jesus' strategy.
by first targeting the Jewish communities.
that already have the religious background.
of the Old Testament..
They appeal to the Jews and the God-fearing Gentiles,.
and then connect Jesus as a promised Savior,.

$^{161}$referred to in the scriptures that they were all expecting..
Paul and Barnabas head through a number of cities.
before backtracking, so you can see that they go inland,.
they go all the way to all these cities,.
and then they kind of do a U-turn back,.
and they backtrack to Antioch..
And then Paul and Barnabas report.
all that God had done through them,.
and how he had opened the door of faith to the Gentiles.
when they returned..
In verses 1 to 4, after returning to Antioch,.
some people came from Judea, where Jerusalem is,.
the religious heart of Judaism,.
and contended that in order for new Gentile believers.
to be fully saved and accepted into the community,.
they needed to be circumcised..
Circumcision is a hallmark of Jewish believers..
In Genesis 17, Abraham is called by God.
to undergo circumcision as a sign of the covenant..
Paul and Barnabas argued against this,.
and then made their way to Jerusalem to clear things up..
And then in verses 5 to 11,.
what follows is another discussion,.
but this time against the believing Pharisees..
This is different from just some people,.
because here are the respected elders.
of the believing community,.
and it required a more delicate handling of the situation..
And then in 12, verses 12 to 18,.
Barnabas and Paul countered the idea..
They argued against the Pharisees.
with the explanation of the sign of acceptance.
of the believers to the provision.
of the Holy Spirit to Gentiles..
They detailed the signs and wonders.
God had done among them..
And the conversion story capturing Cornelius,.
the centurion and others among him,.
are one of the stories depicted earlier in Acts..
In verse 10, in verse 10 of the passage,.

$^{201}$there's a mention of putting on a yoke.
that neither the Jewish disciples.
nor the ancestors have been able to bear..
It refers to circumcision and the extensive set.
of Mosaic laws that the traditional Jewish community.
believe had to be followed to be part of God's community..
Though we do not have these rules in modern Christianity,.
I'd say that we have our own versions.
of heavy yokes that we place on ourselves..
This part of the passage I'm gonna cover.
later in the message,.
along with this..
So this, this is a yoke..
A yoke is a tool placed on oxen.
to pull a cart, a plow, or a wagon..
And it's made of solid wood,.
and back in the old days, it was really, really heavy..
Actually, if you look at my shirt here,.
this is actually called the yoke as well..
It's actually what holds and hangs the shirt on our bodies..
Think about what yokes you might have.
in your personal situations..
This yoke can come internally from ourselves.
and/or externally from those influences around us..
They come from areas like our families,.
our culture, and even our church traditions..
Recognizing some of these self-imposed yokes.
can help lighten what actually might be barriers.
for us to enjoy more of God's grace..
It's actually funny that Miroslav's sharing.
the passage on the yoke..
I didn't pay her for that..
It just happened to work out..
And when we accept,.
recognizing some of these self-imposed yokes.
can help us lighten what might be barriers.
and enjoy more of God's grace..
And when we accept Jesus into our lives,.
we go through a process of justification.
as we enter the family of God..

$^{241}$But just because we enter the family.
doesn't mean that we feel like part of the family..
There's a lifelong process of sanctification..
It's like marriage..
When you marry, you create a new family,.
and then you join your spousal's family..
And there's a period of acclimation.
to adjust to the new dynamics and rituals that you have..
Going through this kind of acclimation.
is like sanctification..
And the Biblically guided process of sanctification.
is becoming more like Jesus..
Though the Bible should be our guide,.
our societal, cultural, and familial influences.
inevitably affect how we think.
the spiritual growth process should develop..
This is true even for those who grew up in the church..
There are traditions that become embedded.
in what we believe we need to do.
that really are not Biblically prescribed..
They are more our preferences,.
and they become heavy yokes that we end up carrying..
These yokes can be individual,.
and they can be organizational..
I was born into the Chinese church..
My parents were young and were active.
in the Cantonese side of the church they were in..
And I was part of that in my earlier years..
Sunday services were very traditional..
All the songs came from the hymn books..
Everybody came in suits and dresses..
It was very, very boring for me..
Even though now that I'm older,.
admittedly, I find comfort in some of that..
In any case, as I started entering middle school,.
I started to attend the English side..
It was more of a Western disposition..
It was more casual, it was more young..
And attendees were mostly American-born.
or mostly raised in the United States..

$^{281}$Now, my normal attire every day,.
going to school through the week,.
were sweatpants and flip-flops..
And they were super comfortable,.
and every day, even into college,.
this was my daily uniform..
But my mom absolutely hates me wearing flip-flops..
And I see this as a generational phenomenon.
because I remember when I was dating Rebecca, my wife,.
and I would go and visit her grandmother,.
she would always warn me,.
"Do not show up in your flip-flops.".
Or else she would just get an earful afterwards..
And even when I visited her parents,.
she'd be like, "Make sure you don't come.
"in your flip-flops.".
And sometimes, I would listen..
But I think it's actually a generational thing..
One thing my mom disliked me more than wearing flip-flops.
was wearing flip-flops to church..
And I think there's this image for those.
that grew up in Hong Kong in the earlier years,.
and only young thugs ran around wearing flip-flops,.
(speaks in foreign language).
And these were the so-called (speaks in foreign language).
I would consistently argue with my mom on this topic.
whenever we would prepare to go to church on Sundays..
And some of you know my mother,.
so you're probably laughing, imagining her..
We'd go back and forth, and at the end of it,.
she would usually conclude that Jesus would disapprove..
Like that would be her way of trying to get me.
to stop wearing my flip-flops..
And being the young teen I was,.
I would then reply that Jesus can't disapprove.
because he only wore sandals, right?.
(audience laughs).
In fact, I visited my mom's church a couple of weeks ago,.
and I wore flip-flops to church,.
and I was attending the Chinese service with her..

$^{321}$And at the end of it, she's like,.
(speaks in foreign language).
And then, of course, I was like,.
I'm actually being more like Jesus.
because I'm wearing flip-flops..
So she didn't take that very well,.
but this is a funny example of how we sometimes.
might add extra add-ons to what is required.
to be Christ-like..
And this example might be funny and not harmful,.
but substitute the flip-flops for something else..
Jesus would not approve of blank..
Or you need to do more of blank to be more like Jesus..
A more serious subject could actually be.
where you should and should not go as a believer..
You shouldn't go to the bars anymore..
You can't stay out late on Saturdays anymore..
You should only go to church on Fridays..
That's the only place you can go to..
Or maybe it's the people that you associate with..
You should only hang out with Christians..
Or you should disconnect yourself from your friends.
who are not believers, even though you might be.
the only connection to Jesus that they have in their lives..
Or it's related to your spiritual practice..
Maybe it's you have to pray a certain way.
for a minimum amount of time,.
or you have to pray in tongues.
for your prayers to be effective..
You have to read the Bible through so many times in a year,.
or you have to go to church so many times each week,.
or you have to serve in so many different roles..
These can all be our modern yokes of circumcision..
They can be extra practices that we add that accumulate,.
and they become more and more burdensome,.
and we can suffocate under the weight of them..
We become very performance-oriented.
versus being very grace-driven..
And we should ask ourselves,.
are there certain practices that we have.

$^{361}$that might be individual yokes in our own lives?.
Praying is good, church is good, serving is good,.
but are we doing them from a sincere posture of faith,.
or are we doing them just to say that we did them?.
So these yokes can also be organizational..
We assimilate a lot of our own cultural traditions.
in how we do church..
They might seem natural and preferential.
because they're grounded.
in our very familiar cultural practices,.
but they should not be classed as biblically required..
These demands can actually overemphasize.
what is actually critical for a church,.
and can create organizational drag.
that prevents effective ministry, leadership, and fellowship..
In the Chinese church, these can be even more pronounced.
when we operate within the greater American environment,.
because the difference is so stark..
And I think one strong Asian influence.
that really pervades our perspective of how we do church.
is our very Asian emphasis on education..
It's a very notable virtue,.
but it can also dangerously border on idolatrous behavior..
Having lived in Hong Kong,.
I saw firsthand the intense pressure there is.
with education, and it starts at the earliest levels,.
preschool or maybe even when you're born..
It got to very absurd levels..
The competition is so fierce,.
and people have to apply their kids.
to try to get into school,.
to even apply for school as soon as they're born..
You can't apply a child for school unless they exist,.
and they do so right after they're born,.
because the lines are so long and it's so competitive..
They start doing interviews with three-year-olds.
to see if they're gifted,.
and I don't know what they can test.
with giftedness at that age,.
other than can they sit down and listen.

$^{401}$for hours to a teacher teach..
I'm sure many of you have your own stories,.
and you've seen this or you've experienced it..
A more extreme example that I have.
is a number of years ago,.
my friend was teaching in one of these.
international private schools in Hong Kong,.
and they do these tours..
Parents can come and they can check out the school,.
and one of these parents,.
or who they thought was a parent,.
was this lady coming on her own,.
and she was part of the tour,.
and it was found out later on.
that not only did she not have a child yet,.
she wasn't even married..
She was preparing that far in advance.
because it was so stressful to think about this topic,.
and they had to plan ahead because it's so competitive..
And educational craze and demand in Hong Kong.
has fueled an extremely lucrative industry.
in tutoring centers in Hong Kong..
We actually have some of these in Bridgeport..
Some of you might have your kids there..
I'm not gonna blame you..
But I don't know any other place other than Hong Kong.
where we have these tutor gods, these bozafsan..
They are individuals who have supposedly mastered.
the art of passing high school and college entrance exams,.
and they make millions of dollars..
They're celebrities..
Their faces are all over billboards and on buses,.
and they're on TV, and it's crazy..
But I believe this actually is grounded.
in our Confucian culture,.
the Confucian-based imperial civil exams..
There's this long tradition from the 600s to the 1900s.
where anybody in society, if they could pass this test,.
they can jump through the ladder in society.
and do very, very well..

$^{441}$And I've seen this historical influence.
in a lot of these Asian churches,.
in Singaporean churches, in Korean churches,.
in Chinese churches..
And it bleeds into what we require educationally.
for candidates that can serve in our churches..
I've seen some job postings for youth pastors.
where you shouldn't apply unless you have,.
unless you're pursuing a master's in divinity, right?.
Or you can only apply to be a pastor or senior pastor.
if you have your PhD or if you're getting your DEM in,.
right?.
Knowledge is a priority..
But being empathetic, pastoral in nature,.
a good leader, or even caring about people,.
that might not even be listed, right?.
Like that, there's a disconnect there..
These restrictive educational requirements that we have.
can be our organizational yokes..
And I wonder if Jesus applied to some of these jobs today,.
if he would be considered as the carpenter's son, right,.
as his main vocation..
We're also a culture of very industrious people, right,.
as immigrants or children of immigrants..
And this infiltrates how we do church sometimes..
Like there's a busyness that gets translated.
into how we do church..
And we have program after program,.
study after study, event after event..
It's like we're working our way towards salvation..
And Sunday Sabbath can actually feel not restful..
And that just isn't right, nor is it biblical..
I value education..
I value strong work ethic..
And these are attributes that I want my own children.
to learn..
Wisdom and hard work are praised virtues in the Bible..
Blessed is the one whose delight is in the law of the Lord.
and meditates on his law day and night..
Good planning and hard work lead to prosperity..

$^{481}$But hasty shortcuts lead to poverty..
We need to ensure that these virtues don't become.
organizational limiters, organizational yokes for us..
That we don't see them as a means and the measurement.
of what a greater church is built on..
Putting on these yokes causes us to lose focus.
and lose dependency on God's grace and mercy..
Substituting it with a sense of control.
that we can have ourselves..
Our motivation needs to be grounded in the community.
of faith and love that is authentic and intentional,.
the same way that Jesus is..
Are we caring for the brothers and sisters around us?.
Are we caring for the community,.
for the people in our community?.
Or are we creating a safe bubble for us to just be in?.
I know I'm ruffling some feathers here..
I'm causing discomfort for some folks,.
maybe most folks here..
And I would probably too if I was in the audience.
and I ended here..
'Cause it sounds like, hey, it sounds like Brian.
is saying we should have no yokes, no rules..
I'm not saying that at all..
I'm not advocating for liberation from any constraints..
Far from it, we need yokes..
Without rules and structure, there would be chaos..
And we as Asians, of course, like rules..
It's why I think distribution-wise,.
we tend to have more people in the areas of engineering.
or the sciences or accounting,.
because they're grounded in very clear rules..
But we need to remember that our strengths.
often have a mirror side to them..
They're usually our weakness as well,.
and we shouldn't be blindsided by them..
Sometimes adhering to the rules can be to a fault,.
where we stop questioning why they exist..
What I want to emphasize is that we need to know.
what the true rules are,.

$^{521}$instead of adding in our own salt and vinegar,.
adding our gacho, right, into what we expect.
church needs to be like and how we operate within them..
(silence).
We should see ourselves with a lightened yoke..
So, not with no rules, but with the lightened yoke..
The end of our passage ends with James, Jesus' brother,.
the leader of the Jerusalem church,.
that says that things shouldn't be made difficult.
for the Gentiles coming to Christ,.
but they should still be instructed to abstain from food.
polluted by idols, from sexual immorality,.
and from the meat of strangled animals,.
from blood, and from blood..
After everything I said before,.
some of you reading this will say,.
"Hey, this sure sounds like.
"we should still be doing a lot of stuff.".
And in the legal field, there is what's considered.
the letter of the law and the spirit of the law..
The letter of the law is kind of our direct,.
very superficial, quick reading of the law..
It's very literal..
But the spirit of the law is actually what the heart.
of the law is meant to promote..
Using maybe the speed limit, for example,.
we can see that when you learn to drive,.
you're taught that you should not exceed the speed limit..
But as you actually start to drive yourself,.
there's this idea of what's called the flow of traffic..
You follow the flow of traffic,.
and that flow of traffic might be faster.
than the speed limit, or it could be significantly slower.
than the speed limit, but you follow the flow of traffic.
because the heart of the law is actually to maintain safety.
for yourself and for everyone around you..
The biblical law is meant to serve as a guide.
for us to understand what it takes to be holy,.
and to understand we are unable to fulfill it ourselves..
Jesus said that he did not come to abolish the law.

$^{561}$or the prophets, he came to fulfill them..
God sent Jesus because we can't do it ourselves..
He came as a perfect being to follow the law,.
perfectly, and to highlight that it is by God's grace.
through Jesus that we can become saved.
and strive to be holy..
Following him is the way to go..
The end of the passage here in Acts is really prescribing us.
to operate with a mind to be holy,.
holiness that is spiritual, holiness that is sexual,.
holiness that is nutritional or bodily..
The instructions are still to maintain.
basic standards of holiness, but not to the degree.
where it becomes a checklist to save ourselves..
And you know what?.
Like, these are practices that non-believers can identify.
as distinguishing marks of the believing community,.
and why James still prescribes these rules.
for those to follow, because he mentions that.
it's been told for a long time.
that this is what the Jewish community does..
The Christian faith is an invitation from Jesus.
to become part of God's family..
It's also a lifelong journey to become more like him,.
to be set apart, to be holy..
The process is not meant to cause us.
to feel like we're drowning in rules,.
where in our individual and church lives,.
the rules have become our masters,.
and like the trained elephants,.
we're afraid to step out of bounds..
The rules themselves are guardrails for us,.
so that we don't recklessly charge ahead.
to our own harm or to others' harm..
Jesus provides us a lightened yoke, an easier path..
It is a path ahead that is full of forgiveness and grace,.
for he has bore the burden of the law..
We have to just open our eyes.
and see what the new reality actually is..
Again, I said, no, I didn't pay Mirs.

$^{601}$to say the verse that she did,.
but it's actually in my slide as well..
And Jesus says, "Come to me,.
"all you who are weary and burdened,.
"and I will give you rest..
"Take my yoke upon you and learn from me,.
"for I am gentle and humble in heart,.
"and you will find rest for your souls..
"For my yoke is easy, and my burden is light.".
Jesus does the heavy lifting so we don't have to..
Make sure to put on the lighter yoke,.
and maybe feel free to wear some flip-flops at church as well..
(audience laughing).
Thank you..
Let me just end in prayer really quickly..
Heavenly Father, we just thank you.
for Jesus, we thank you that he came to this world.
so that we can have a lighter yoke,.
that we will not feel crushed by the rules.
that we feel we need to follow to become saved,.
and that you extend that freely through your grace,.
through Christ's sacrifice..
Bless us as we leave here,.
and help us to reflect on what might be.
our individual yokes,.
and what might be our organizational yokes..
And let us feel freedom,.
freedom in your grace and your mercy..
In your son's name we pray, amen..
\newpage



\section{}
\label{sec:R_BYGFCmbDs}
\textbf{佳節倍思恩(三): 安息日  - 陳崇基牧師}
\newline
\newline
連結: \href{https://youtube.com/watch?v=R_BYGFCmbDs}{\texttt{https://youtube.com/watch?v=R\_BYGFCmbDs}} ~~~~ 語音日期: 2022-10-19
\newline
\newline
\hyperref[sec:yX0mME__Hqg]{\small{< < < PREV SERMON < < <}}
~
\hyperref[sec:index_chronic]{\small{[返順時目]}}
~
\hyperref[sec:index_scriptual]{\small{[返順卷目]}}
~
\hyperref[sec:DueqCDH4G_Q]{\small{> > > NEXT SERMON > > >}}
\newline
\newline
$^{1}$(請看片段).
我們繼續看這個系列.
我們看看在我們生命當中有什麼節期.
有什麼時間提醒我們.
這些屬靈的生命.
所以我們這次去看另一個節期.
上一次我們提過.
就是豬膀節.
Feast of Tabernacles.
剛剛這個星期以色列完結.
這個豬膀節.
今天我們要看的是賈伯.
賈伯這個應該很清楚.
因為是七天的其中一天.
這一段時間.
如果你在以色列的街上走.
我們就會覺得很神奇.
因為星期五的日落.
到星期六的日落.
就是他們的安息日.
就是一個星期的第七天.
所以他們會在街上沒有車.
或者沒有人來走等等.
因為在賈伯那裡的時候.
一切都是靜態的.
每個人都回家.
然後不做工.
也不打火.
所以不按一些按鈕.
或者電話會關掉等等.
或者走路不可以多於半英里.
如果那些是非常保守的人.
就會半英里之內.
只是半英里之內可以走路.
這些都會對你很神奇.
其實不只是在那裡.
美國也有不同的角落是這樣的.
如果你有去過紐約的布魯克林.
一些比較多猶太人住的地方.
或者我們芝加哥北面.

$^{41}$斯科奇.
也有很多不同的猶太人在那裡聚居.
他們就會實行這個安息日.
安息日.
顧名思義就是安息不做事.
所以有很多我們覺得.
現代人覺得比較神奇的東西就會出現.
譬如這裡有幾張相片給你看.
有一些大樓.
有一些高樓大廈.
就會有一些安息日的電梯.
Schochberg elevator.
這裡看到Schochberg elevator是什麼.
就會是計劃了.
以至星期五的新單開始.
電梯就會每層停一次.
為了什麼.
為了你不會打火.
因為你按鍵就會打火.
按鍵就會亮燈.
所以這些比較神奇的東西.
這個挺有趣的.
這個是酒店的Schochberg elevator.
它有兩部elevator.
第一部在右邊的elevator A.
你見不見到它想說什麼.
由下面basement二樓.
去到ground floor零.
然後它會逐層停一次.
然後它就會下來.
然後再倒過來.
如果你要四樓去一樓的話.
你就要繞過大轉.
另外這個elevator D.
就會跳上去五樓開始.
然後五樓到九樓就逐層停.
然後再倒過來.
這個挺有趣的.
In case of fire.
Do not use the elevator.

$^{81}$Of course.
火警的時候當然不要坐電梯.
但如果你真的坐錯電梯的話.
那又怎樣.
你坐了Schochberg elevator.
可能生命有危險.
因為它可能會把你載上去.
這些都是很有趣的東西.
就是講到Schochberg停止.
聖經裡面用了很多動詞和名詞來形容.
什麼叫做安息.
我們看其中幾段經文.
《創世記》第二章.
這裡第一次記載所謂和安息日有關.
《創世記》第二章第一次.
天地的萬物已經做齊了.
到第七日.
上帝做物的工已經完畢.
就是在第七日.
揭了他一切的工.
安息了.
上帝似乎給第七日定為聖日.
因為在這一日.
他揭了他一切創造的工.
就安息了.
但這段聖經.
就是在六日.
蒼世之後記載的.
你記得有六日.
有晚上有早晨.
因為他們計算一日的開始.
那時候是晚上開始計算的.
由晚上安息到第二天的早晨.
就有六日.
上帝就創造了萬日.
一切都是非常美好的.
一切都是非常好的.
然後就記載第七日.
二章第一節.
他說一切都做成了.

$^{121}$然後呢.
上帝就停止了他的工作.
是不是他累了.
是不是他太過辛苦了.
所以就像我們這樣.
上了一星期的班.
終於可以休息一下.
終於可以停下來.
是不是這樣.
似乎又不是.
他說停了他之前一切的工作.
兩次提到.
停了他一切的工作.
就是之前六日的工作.
就是因為這樣的緣故.
停止了.
然後有第七日.
於是猶太人的文獻.
或者舊約裡面的經文.
就會用這個作為一個基礎.
來宣佈有安息日的誡命.
安息日誡命其中一個最熟悉的經文.
就是在《出亞給記》二十章第八節.
這裡說了十誡.
第四誡.
就是說到安息日的誡命.
他說當紀念安息日首為成日.
六日你要勞碌你一切的工.
但第七日是向耶和華你上帝當首安息日.
這一日.
你和你的兒女瀑被生畜.
和城裡面寄居的客女.
無論何工都不可作.
因為六日之內.
耶和華做天地海和其中的萬物.
第七日就停了.
安息了.
所以耶和華賜福與安息日.
定為成日.
第五.

$^{161}$星期五的晚上.
開始了這個安息日.
一切都停工了.
一切都是暫止著.
好像在一篇樂章.
你聽一首歌.
有丘子芙.
停下來.
安息日會不會就是提醒我們.
在生命裡面需要有這樣的律動.
一個這樣的韻律.
六日.
你不是偷懶.
你要作工.
但到第七日.
你就要停下來.
這個誡命寫的時候就是基於《創世記》第二章所說的.
上帝就是這樣做事.
所以有這個誡命.
人就是怎樣.
好像上帝一樣.
就要去停止.
停工.
你再看回剛才那段經文.
第二章.
上帝停止了他一切的工作.
第七日安息了.
上帝就賜福給第七日.
定為成日.
祂祝福了那天.
然後.
Sanctify.
去將祂定為聖日.
就是將祂分別出來的意思.
很奇怪.
上帝開始做事.
創立萬有的時候.
祂沒有Sanctify任何東西.
沒有一個地方定為聖地.
沒有一個地點定為一個神聖的地點.

$^{201}$也沒有一個空間定為所謂Holy的.
第一樣東西 祂定為聖的是什麼.
Sanctify.
是時間.
不是Space.
而是Time.
Sacred Time.
什麼意思.
定為聖.
就是分別為聖.
Sanctify.
你知道這些字眼.
通常都是祭司的文件裡面出現的.
或者聖殿裡面的文件.
說到聖殿的器皿.
什麼金碗等等.
金燈台等等.
油啊.
全部都是物件.
就是將這些物件.
將這些東西.
所謂定為聖.
Sanctify.
將它Set apart.
就是說這些東西特別不同.
這些東西特別跟其他東西不同.
有聖與俗之分.
這個很奇怪.
上帝第一樣東西定為聖的.
是時間.
在生命裡面一段時間.
不是一些物品.
現代人有一個弊病.
就是我們沒有.
時間上沒有分別.
我們一個星期七天.
每天24小時.
我們都捍衛一樣的.
Sacred.
每樣東西都一模一樣.

$^{241}$每時每分都是一樣.
時間只不過是我們的工具.
我們的工具.
我們想達至某一個目的.
我們希望有生產.
我們就用時間.
來達到我們的目的.
我們希望能夠達到某一些的目標.
我們就用時間.
來成為一個途徑.
It's a means to some other end.
這個安息人.
就將它反過來.
原來時間不只是工具.
不是我們用的東西這麼簡單.
上帝定為聖.
在第七天捍衛神聖.
Sacred time.
在你生命裡面.
有沒有時間是Sacred的.
神聖的時間.
在這段時間.
跟其他時間是完全不同的.
這個時間.
這個period.
是完全跟其他東西不同的.
這個很有趣.
有一個很有名的猶太人的拉比叫做Abraham Herschel.
他寫了一本書.
幾乎是經典.
在猶太人的文獻.
甚至基督教的文獻.
是講Sabbath.
這個拉比是20世紀的時候.
很重要的拉比.
他寫Sabbath的時候.
是寫什麼呢.
他說.
There is a realm of time.
where the goal is not to have but to be.

$^{281}$not to own but to give.
not to control but to share.
not to subdue but to be in accord.
在生命裡面有很多不同的領域.
他說時間.
講到Sabbath的時候.
上帝定位性的東西.
那段時間.
那段sacred time.
不是要我們去佔領的.
也不是我們要去控制.
而是我們要去share.
在這段時間.
這段sacred time.
is to be.
在那裡存有.
不是佔有.
什麼意思.
妹妹在聖經裡面.
提到安息日的時候.
Sabbath.
當然是停工.
應該是跟我們之前那六天.
做的事情有所不同.
有什麼不同.
妹妹提到Sabbath的時候.
就是跟什麼有關.
譬如這段聖經.
利未基.
23章第三頁.
他說六天你要作供.
第七天就是.
Sacred.
Sabbath.
聖安息日.
當有聖會.
Holy Assembly.
你們什麼供都不可作.
這是在你們一切的住處.
香燕華守安息日.

$^{321}$耶和華的節期.
就是你到了日期.
要宣告為聖會.
前後就說到一些.
Festival的時間.
一年裡面有不同的時間.
甚至一個星期.
有不同的時間.
就是定為聖的.
做什麼人.
聖會.
就是說他們聚集.
會去敬拜.
回到創世記第二章.
到底說到Sabbath的時候.
說到安息日的時候.
是什麼意思.
很多時候我們.
現代人看的時候.
我們就覺得.
可能他做完事.
累了.
所以他們休息一下.
對不對.
好像我們六天做完事很辛苦.
我們就停下來.
絕對當然不是說上帝累了.
上帝不像我們.
你看回那個.
形容創造世界的時候.
那些字眼.
有什麼特別.
如果你換了你的角色.
你是當時的以色列人.
你是一個祭司.
你會聽到這些字眼的時候.
想到什麼.
上帝說什麼.
他一切都做起了.
如果你是一個祭司.

$^{361}$你就會想起.
做會望.
會望就是上帝住在我們中間.
上帝的同在.
是一個標記.
昭和禁忌39章這裡.
章望就是會望.
一切的工就這樣完了.
凡耶路亞話所吩咐摩西的.
以色列人就照樣做.
一切的工作都完了.
這句話一次又一次.
在建造會望的時候出現.
所以如果你用這樣的字眼.
去看回創世的第一章.
上帝創立這個世界的時候.
其實他就創立這個世界.
好像創立一個會望.
去到為自己建一個聖所.
一個聖殿.
你這樣看創世的第一章.
他就不是純粹說一些.
星球,九大行星.
全部放在一起.
不是一種純粹物理上.
物質上的擺設.
而是一個聖所.
整個宇宙.
整個的Cosmos.
Universe.
原來就是上帝的聖所.
宇宙就好像上帝的聖殿.
萬有上帝所安排的東西.
無論是天上的太陽,月亮.
地上的走獸等等.
都好像聖所裡面.
會望裡面的家具.
一個會望有的.
就是有很多不同的器具.
不只是一個空的場所.

$^{401}$裡面有什麼.
金燈台,陳設餅.
不同的擺設.
全部東西就像摩西他們.
做會望的時候.
做Temple的時候.
這個東西放在哪裡.
那個東西放在哪裡.
然後這個東西放在外面.
有些東西放在裡面.
就是這樣安排.
這樣Arrange了.
完了之後就怎麼樣.
上帝的榮耀就住在他們中間.
這個就是會望做的過程.
如果這樣看整個世界.
整個宇宙創世的第一章.
就是整個宇宙.
就是上帝擺設了所有的東西.
然後他就住在中間.
所以聖經裡面提到安息這件事.
不是因為累了就休息.
而是上帝擺設好所有東西的時候.
就要運作.
整個世界開始運作的時候.
他就住在裡面.
所以第七天.
其實就是登基的時候.
就是上帝真的作王,作主.
他要去運作.
Operate這個世界.
舉個例子.
疫情開始的時候.
大概三年前.
因為家裡要買一部新電腦.
因為要做那些Online的東西.
就買了一部Desktop.
然後有很多線.
有很多鍵盤.
或者有很多零件.

$^{441}$Hard Drive等等.
就要全部裝起來.
有鏡頭.
又有不同的Screen等等.
搞了我很久.
幾乎搞了幾天.
做完之後怎樣.
那些線全部駁好.
我會不會說.
我們可以打工了.
停工.
駁好所有東西的時候.
你就做什麼.
你就touch部機.
你就開始做什麼.
就是工作.
開始Operate.
或者好像搬屋.
搬到一個新的空屋.
你進去可能是廚房.
把碗碟放在裡面.
爐頭放在裡面.
所有器具全部都搞好了.
你就做什麼.
你就不會罷工.
反而是相反的.
你去Operate.
是時候要開工了.
是時候要做事了.
就是這個意思.
創世記第一章.
說到上帝安息的時候.
其實上帝擺好了整個宇宙的東西.
Rest.
安息.
不是停止不做事.
相反是反而去到工作.
上帝作王.
去到控制這個世界.
去到管理這個世界.

$^{481}$這個和我們現代人想安息這個字是不同的.
我們想起安息就是累了不做事.
但是在聖經裡面說到安息的時候.
就不是這回事.
譬如約書亞記第21章44字.
這裡你猜猜哪個字是和安息有關.
耶和華指著他向列祖所起的誓.
所應許的一切說話.
他們洗耳識的人.
四境平安.
他們一切仇敵裡面沒有一個人在他面前.
可以站立得住.
哪個字是和安息有關.
四境平安.
四境都安息.
這裡說的是什麼呢.
約書亞記入主迦南的時候.
以色列人入主迦南.
然後他們有很多仇敵.
有很多混亂.
然後他們擺平這些混亂的威脅的時候.
發生什麼事.
就來到這裡.
就是四境平安.
上帝使他們四境平安.
能夠有安息.
仇敵沒有一個人在他面前能夠站立得住.
能夠有安息是什麼意思.
是不是說我們坐在那個懶人.
坐在那個沙發.
然後我們就休息.
不是.
四境平安的意思就是說角度已經建立了.
現在是時候去運作.
是時候去運作.
用這樣的概念.
我們可能明白創世記第一章想說什麼.
萬物被安排好.
上帝安息了.
上帝作王.

$^{521}$他開始運作這個世界的時候.
這個世界就進入一個新的order.
一個新的秩序.
一個新的rhythm.
所以安息日和敬拜就有關係.
叫以色列人守安息日.
並且守聖會.
其實意思就是說.
我們認定.
在這個sacred time的時候去認定.
真正operate這個世界.
真正控制我們這個世界.
是上帝.
That's why we have holy assembly.
That's why人會停下六天的工作.
放下我們所控制的東西.
今天在這個時間.
我們就知道原來上帝才是眾生.
上帝才是我們整個宇宙的眾生.
不是我.
好幾年前有一個拉比.
猶太的拉比.
他是一些解答問題的主持人.
在節目裡有些人會寫信給他.
然後他就會用信來回應.
告訴他.
你應該怎樣去招架你現在生活裡的困難.
有一次他收到一封信.
那封信大致上寫了這件事.
那封信就說.
敬愛的拉比.
敬愛的老師.
我很不開心.
我老公不聽我的話.
我子女也不聽我的話.
我的工作也沒有進展.
他說了一大堆.
我時時都不如意.
我的心情也很抑鬱.
你可不可以幫幫我.

$^{561}$如果你是那拉比你怎樣回應他.
猶太人的拉比就很妙.
好像耶穌和法利賽人對話.
答了又不答.
經常有這樣的事.
那這個拉比又怎樣回應.
他收到一封信.
他就將每一句的第一個字圈起來.
然後再摺起那封信寄給我.
那是什麼字呢.
就是我字.
I feel depressed.
I don't feel okay.
I am not happy.
I lack this or that.
安息人.
守聖會.
有個sacred time.
意思就是說.
世界的中心宇宙中心不是我.
是上帝.
上帝管理我的生命.
上帝管理我的世界.
每逢這個時間.
所謂的sacred time.
敬拜的時間.
其實就是一個reset button.
好像你用別部電腦用得久.
你就要關了它.
shut down.
然後再reboot.
是一個reset button.
我們來到上帝的面前.
我們知道祂才是世界的中心.
祂才是宇宙中心.
不是我.
安息人怎樣實踐.
在新約的時候.
這個安息人慢慢慢慢.
由第七天就轉為第一天.

$^{601}$為什麼呢.
就因為七天的第一天發生什麼事.
基督從死裡面復活.
七天的第一天.
清晨的時候.
基督從墳墓復活.
七天的第一天.
祂成全了舊人.
祂做了祂一切要做的事.
祂復活的時候是七天的第一天.
就是新的秩序開始.
在基督裡面.
孟乖.
在新約教會到現在.
就認定七天的第一天.
就是那個reset button的時候.
就不再是七天的第七天.
是七天的第一天.
上帝藉著祂的兒子.
從死裡面復活.
一個新的creation.
一個新的開始.
這個或者應該提醒我們.
今天所謂的守聖會.
我們去敬拜的時候.
我們worship的時候.
是應該帶著什麼心態.
有兩樣實踐的事.
我們現在已經試過了.
我提出兩樣.
第一樣.
六天之後有第一天.
六天之後有第二天.
六天之後有第三天.
這個would.
其實我們都需要放在時間表裡.
當然我現在跟你說.
其實都是跟我說的.
因為大家知道我那時候不休息.
但其實我也休息的.

$^{641}$我休息你不知道為什麼.
我來到敬拜.
很多時候.
人會覺得你只是來工作而已.
我可以有這樣的心態.
但我盡量都不會有這樣的心態.
我今天來敬拜.
是跟大家一起事奉.
守聖會.
Holy assembly.
雖然我有些事情要做.
但我是一個敬拜.
一個事奉這樣的心態.
我不是當一個task.
而是當一個.
This is a sacred time.
是一個守為聖日.
一個令人不同的一段時間.
很多時候敬拜的時候.
我都聽到一些道不是很中聽.
包括自己那些.
但我不覺得有關係.
因為那個point.
sacred assembly.
是跟上帝遇上.
怎樣跟上帝遇上.
我校正了我的心態.
我就能夠跟上帝遇上.
所以今天.
是應該來到敬拜去事奉.
不是純粹坐下來.
我們什麼都不做.
喝杯茶就算了.
六天.
然後又一天.
六天之後又一天.
一個小的習慣.
去到影響我們跟上帝的關係.
第二.
除了我說.

$^{681}$這個神聖的時間.
應該是重新在我們心中reset.
應該是我們做好準備.
用心做好準備.
第二.
就是每一天.
我們都應該有一個小的安息.
在我今天的beat.
我寫了一個人.
叫做Brother Lawrence.
勞倫斯帝兄.
幾百年前.
他有一個小的修士.
他就在法國巴黎.
一個修道院打工.
去洗碗.
做廚房的事.
他就很struggle.
每一天.
我做的事.
好像沒有什麼意思.
然後他就領會了一個秘訣.
怎樣活在上帝的同在裡面.
怎樣調整自己.
能夠知道上帝是跟他同在.
就是他洗碗的時候.
很忙的時候.
他就會take一分鐘去禱告.
早上是這樣.
中午也是這樣.
下午他也是這樣.
做一個小的祈禱.
他就說.
上帝我是為了你.
因為愛你而做這件事.
洗碗的時候做.
切東西的時候他做.
忘戀的時候他這樣做.
他說這個就是他的秘訣.
怎樣在小事上.

$^{721}$能夠知道上帝是跟我同在.
上帝跟我是present.
求主幫助我們.
reset一下我們.
自己忙亂的生活.
我們請敬拜的女兒上來.
我們就去唱一首詩歌.
讓這首詩歌成為你和我的祈禱.
這個禱告是.
跟這位創造天地萬物.
去控制一切的上帝.
我不知道你心裡面有什麼煩惱.
剛才敬拜開始的時候.
也有問這個問題.
可能你連問這個問題.
停下來想一想.
也沒有這樣的時間.
這是一個sacred time.
就是我們可以去reset我們自己.
唱一首詩歌.
《我愛你》 詞:陳家麗 曲:陳志遠.
藏我在你的思旁.
藏在眼裡你用無聲說我說話.
白日太陽必不傷我.
夜裡月亮也不害我.
此刻是有耶穌在我的補償.
再一次呈現我身於你面前.
唯獨你永遠化此生是我歡愉.
在世間再沒有別人.
可以奪去我的思想.
我的一切全屬於你.
我連一切也都歸你.
耶穌 耶穌 藏我面前.
命吩咐你的使者保護我.
用你的雙手托著我.
親眼看我每個記載 每個痛覺.
耶穌 我只要你耶穌.
好 深呼吸.
上主我們在我們宇宙的中心.
你是各位創造萬有的神.

$^{761}$創造我們都在我們生命當中.
作安排 做管理的神.
主我們求你.
求你在這個時候.
讓我們能夠活生生.
在生活的網路.
在世間上有很多令我們分曉的事.
上帝我們求你.
讓我們安靜 讓我們平安.
讓我們在你自己的主權底下.
能夠有安息.
主你做好.
讓我們在你面前.
聽你的說話 去親近你.
主為我們中間每一位.
甚至網上的人.
離開芝加哥不同的地方.
每一位的朋友.
主我們求你.
順天的平安與我們共聚.
讓我們認定你是那位管理萬有的神.
主你的自然.
能在我們生命當中.
願我們在你面前.
多謝各位 上帝祝福各位.
\newpage



\section{}
\label{sec:DueqCDH4G_Q}
\textbf{聖經地理 EP11 - 地理與釋經 - 中央山地 - 1 (陳崇基牧師)}
\newline
\newline
連結: \href{https://youtube.com/watch?v=DueqCDH4G-Q}{\texttt{https://youtube.com/watch?v=DueqCDH4G-Q}} ~~~~ 語音日期: 2022-10-21
\newline
\newline
\hyperref[sec:R_BYGFCmbDs]{\small{< < < PREV SERMON < < <}}
~
\hyperref[sec:index_chronic]{\small{[返順時目]}}
~
\hyperref[sec:index_scriptual]{\small{[返順卷目]}}
~
\hyperref[sec:Pgm9dzwLc3c]{\small{> > > NEXT SERMON > > >}}
\newline
\newline
$^{1}$我們現在不如講一些.
同一個家園地有不同的地帶.
我們逐個地帶來講一講.
在這些地區裡面.
我們能夠了解到到底有什麼動態.
地理上有什麼動態.
以及幫助我們怎樣去看聖經.
首先來這個地圖.
我們先放大.
這裡先關掉爐子.
我們講過這個家園地是中間之地.
聖經裡面也有這樣的字眼來形容.
這個以色列巴勒斯坦地.
聖地這個地方.
由北面這個黑門山.
直到南面的死海.
就是聖經裡面所說的應許之地.
家園地這個地方.
我們說這個家園地長長窄窄的.
由海邊到約旦河和死海的地方.
其實開車一個多小時.
所以是一個很窄的走廊.
左邊就是海洋.
再右邊就是沙漠.
沙地阿拉伯的沙漠.
中間這個家園地有幾個地區可以分到.
第一個地區就是所謂沿海平原.
加密山以南.
一直到西大半島的北面.
都是沿海平原.
沿海平原再上一點.
就是一些低山地.
英文叫rolling hills.
大概200米高.
比較小的山.
還有一些山谷比較寬.
然後上到這裡.
我們看到有耶路撒冷.
巴黎行.
以至下面有希伯倫.

$^{41}$這裡叫做中央山地.
耶路撒冷會在中央山地的分水嶺.
過了耶路撒冷就是這裡.
就是猶大曠野.
或者變雅敏曠野.
下來這裡就是死海.
所以如果你這樣算.
北與南的話就有不同的地區.
事實上聖經裡面也有這樣的分區.
聖經不同的字眼不同的經文.
都有說到這個地方有不同的字眼去形容它.
第一個經文想起的就是約書亞記第十章.
這裡說到約書亞帶領以色列人.
去到入主迦南.
他們決定性地擊敗了迦南裡面的人.
他講到什麼呢?.
全地的人.
然後有山區.
尼納夫.
低地.
和山坡的人.
這裡有四個.
首先有四個地理的名稱.
第一個就是山區.
haha.
就是the hill的意思.
或者the hill country.
就是說山的區域.
然後就是Hanageth.
現在比較新的譯法就是用那個音.
Hanageth.
就是the negeth.
negeth.
其實就是south的意思.
以往的和合本的譯法就是南地.
因為確實真的是南地.
所以首先這兩個地方.
就指出有個山區.
還有山地.
還有negeth.

$^{81}$然後第三就是這個Shephelah.
這個地方.
Hashephelah.
這個地方.
Shephelah.
或者Shephal.
其實是降低的意思.
就是低地的意思.
所以其實一個正確的譯法.
應該是現在這個和合本的修訂版.
所譯的低地.
這個就是Shephelah.
以往舊的譯法.
就是譯得不是很準確.
就是地理上的字眼.
又譯了做高原.
因為你聽到高原的時候.
你不會想起一些小的山.
一些很高的地方.
平原的地方.
海拔是比較高一點的.
你會想到高原.
不會想到低地.
所以這兩個譯法.
你一定要二取其一.
地理上正確來說.
就一定是低地或者低山地.
英文就是low hill.
或者是西山地.
foothill這樣的意思.
然後就用一個比較籠統的字.
來講山坡的人.
山坡.
Hashadot.
其實按照不同的經文.
其他的經文.
這個Hashadot.
Hashadot.
其實就是Mibbar.
就是曠野的意思.

$^{121}$Mibbar就是曠野.
所以很多時候.
山坡其實是講到.
中間山地的分水嶺的右邊.
東邊的山區的地方.
就是一直下去死海的地方.
剛才那四個地理的名稱.
記不記得.
首先有什麼.
山區.
就是我所講的.
耶路撒冷.
百里行到下面的希伯倫.
上面一直去到事件的地方.
這裡長長窄窄的一個山.
就叫做山區.
然後第二個.
就是Nagev.
Nagev就是下面這裡.
如果找Bezhevel出來.
就是別士巴.
見到嗎.
別士巴.
聖經裡面經常有句話叫彈到別士巴.
彈在哪裡.
Tal'din.
彈就在這裡.
就是黑門山下面.
加利海上面.
彈這個地方.
黑門山山腳.
彈到別士巴就是由北至南.
這個其實一般來說就是迦南地的.
南北的邊界.
所以Nagev就是這個範圍.
圍著這個別士巴這個範圍.
越南的水越少.
所以這裡比較乾.
雖然久不久都會有雨水可以種蜜田.
但是大部分時間雨水不是很充足.

$^{161}$也很難預計到.
所以這個地方大部分時間都是木洋的地方.
連今天都是一樣.
Bezhevel附近比較乾.
Nagev就是這個意思.
南地在南邊.
第三個形詞就是低山地.
希望以後我們知道的時候.
這個翻譯應該是做低地.
這個低地的意思.
在地理上的位置.
就是在這個沿海平原.
和這個山地的中間.
Bezhevel這個地方.
你看到這裡的山都是比較低一點.
這些山比較低一點.
一直向右.
之前都有提過.
就會上到耶路撒冷,百里行等地方.
中央山地.
就會去到海拔比較高的地方.
一直下來的時候那些山和谷就會比較大一點.
直到下到沿海平原.
所以聖經裡面都有這樣的了解.
聖經的經文都有這樣的了解.
不是亂寫的.
然後就是山坡.
為什麼叫山坡呢.
因為比較籠統的形容就是這個地方.
因為由這個分水嶺.
一直下來到死海.
就是一個斜的山坡.
這個地方也是Mibbar.
就是這個抗疫的地方.
抗疫我們之前提過.
這個山地之後.
雨水很多時候不會到的地方.
所以去年的旅遊歷史歷代.
都是一個木洋的地方.
乾旱的地方.

$^{201}$和下面的Nagev都是一樣.
這裡就是其中一段經文.
講到這些地理.
將這些地理分了.
將整個聖經的地圖.
分開了不同的區域.
有這個地理的分析.
分開了不同的區域.
有中間的山區.
然後南面的Nagev.
然後去到比較西一點.
去到所謂的西山腳.
Shafila這個地方.
然後再過到這邊.
講到這個山坡.
這個就是約斯亞.
進去的時候.
經文這樣記載.
我們試一下.
只看這個Hahar.
Hill Country這個地方.
給一些相片給大家看.
這裡就是Hill Country.
這裡是北面.
這個以法連山地和馬來西亞山地.
比較北一點的地方.
這個地方.
就是以法連山地和馬來西亞山地.
這個地方.
首先拿一個Shakhtyam出來.
試一下.
這個.
史劍.
史劍就是他們.
進入了迦林的時候.
一個很重要的地方.
然後還有什麼地方.
Shiloh.
就是停泊會幕的地方.
你看到這裡幾點.

$^{241}$由這個.
耶穌的平原.
一直向南.
這個史劍.
士羅.
然後中間有伯特利亞.
等等.
然後就到耶路撒冷.
這個.
伯特利亞.
還有這個.
希伯倫.
你看到幾乎是一條線.
幾乎是一條線.
這個就叫做Hahar山地.
Hahar山地.
中央山地.
你看到這些山谷都是比較窄一點.
下去看看.
在耶路撒冷附近看一看就知道.
是吧? 還有橄欖山等等.
都是比較窄的.
這些的.
山谷都是比較窄一點.
是吧?.
然後.
你離開了山區之後.
山谷都會比較.
闊一點.
這裡一直向西的時候就很闊.
去到伊什圖.
或者伯氏墨等等.
這個就是山區.
我們為什麼要認識呢?.
因為山區.
是一個嘉林地的中心地帶.
以色列人進入了嘉林的時候.
他們的.
中心地帶.
即是他們起初.

$^{281}$進入嘉林的時候.
他們所聚居的地方.
都是在這個山地.
因為這樣的緣故.
所以有很多時候他們去形容.
一些事物.
或者他們表達神學的時候.
都是和山地有關.
即是.
會說上帝是.
盤石.
什麼是盤石?.
不是一塊塊的盤石我們可以扔的.
拿得起扔得到.
那些不是盤石.
盤石是一些大的山岡.
那些的.
石灰岩就是一層一層的.
很遠古的時候.
已經形成了.
所以一層一層的緣故.
那些水就會.
涉入岩石裡面.
因為岩石本身也是.
會滲水.
所以很多時候山腳就會.
出一些水泉.
那些水泉都是一些自然的水.
所以.
在聖經裡面.
盤石的時候.
是一種.
賜予生命的.
一些地理的名詞.
一些地理的元素.
在這些山區裡面.
有這樣的山與谷.
很多時候.
不單止盤石.
是給人可以去到.

$^{321}$有保護.
甚至.
是給人有生命.
所以這些山的地方.
很多時候找到水泉等等.
這個就是我剛才所說的.
以法連山地.
有史以來.
都有.
累積了很多人工.
所造成的一些改變.
這個山本身.
就不是這樣一級一級的.
只不過.
以往的古人.
聖經以往的時代.
他們在這些山區裡面聚居的時候.
就把這些山.
做成一些梯田.
等於我們中國南方.
其實都有很多這樣的地方.
因為山多平地少的地方.
特別.
就需要在這些山那裡.
挖一些梯田出來.
這個就是北面的山區.
我們看看南面的山區.
南面的山區.
我們可以看這個.
這個就是.
梳列谷的山區.
你見到這個梳列谷.
本身是一些很窄的山谷.
有高山.
有低谷.
但是在這些山上面.
也很清楚.
你見到有些梯田在那裡.
這裡是一個很漂亮的.
郊野公園.

$^{361}$遊人可以在這裡.
去遠足.
走走這些山徑等等.
上面你見到有不同的梯田.
就是上面.
梯上面的表面.
就是一些窄的田園.
很多時候最適合.
栽種的就是一些果樹.
你見到有些不同的果樹.
橄欖樹.
無花果樹等等.
或者是葡萄園.
葡萄的.
原子.
就在這些梯田上面.
去建造.
這個就是一個很漂亮的.
在這些梯田上面去建造.
這裡甚至見到.
有一間屋.
有一間古舊的城樓.
有些人可能住在這裡.
周邊就是他們的.
田園.
梳列谷.
就是其中一個.
山區.
中央山區.
去西面的其中一個山谷.
叫做梳列谷.
\newpage



\section{}
\label{sec:Pgm9dzwLc3c}
\textbf{聖經地理 EP12 - 地理與釋經 - 中央山地 - 2 (陳崇基牧師)}
\newline
\newline
連結: \href{https://youtube.com/watch?v=Pgm9dzwLc3c}{\texttt{https://youtube.com/watch?v=Pgm9dzwLc3c}} ~~~~ 語音日期: 2022-10-21
\newline
\newline
\hyperref[sec:DueqCDH4G_Q]{\small{< < < PREV SERMON < < <}}
~
\hyperref[sec:index_chronic]{\small{[返順時目]}}
~
\hyperref[sec:index_scriptual]{\small{[返順卷目]}}
~
\hyperref[sec:26WK5R9cMYM]{\small{> > > NEXT SERMON > > >}}
\newline
\newline
$^{1}$Sawick Gorge就是其中一個中央山區.
下到西邊的其中一個山谷.
叫做Sawick Valley.
在哪裡呢?.
你見到我點了出來.
但我只能點一點.
一個谷當然不是一點.
Sawick Valley.
你見到這裡有一個山谷.
然後這個山谷一直上去.
然後上去就去到哪裡呢?.
就是去到耶路撒冷.
所以這個就是Sawick Gorge.
在整個山區有不同的山谷下來.
有好幾個.
其中一個最重要的就是Sawick Gorge.
如果你甚至看看今天的路線的編排.
很多時候都是跟著這樣的谷去跟.
這個就是Sawick Gorge.
38上來然後上到1號公路.
1號公路就是一個很重要的幹線.
很重要的公路.
就是由這個沿海平原上去耶路撒冷.
跟著1號公路一直上去.
就去到這個山谷的.
所謂山地的中間最重要的地方.
就是整個的分水嶺.
然後就去到耶路撒冷這個地方.
所以跟著1號公路.
你一路開車的時候就會看到.
由平原一直開車上去.
平原一直跟著公路.
你就上到山區.
其中一個的谷就是剛才所說的.
這個Sawick Gorge.
整個的自然的山谷.
這個地方在哪裡呢.
拍照這個地方.
這個地方就是近耶路撒冷.
就是跟著耶路撒冷.

$^{41}$耶路撒冷出一點點.
找一下耶路撒冷.
耶路撒冷出一點點.
就有一個郊野公園.
所以這裡是山谷以上的地方.
你會看到有這些梯田.
梳理谷在聖經裡面提過.
史詩紀第十六章.
說這是以後參孫的故事.
有座梳理谷愛上的一個女子叫做大莉拉.
其實她就是菲律賓人.
梳理谷.
梳理是什麼意思呢.
Sawick.
其實就是葡萄藤的意思.
就是種這些葡萄藤.
所以這個名字很適合這個地方.
改名就是因為這個地方.
很多時候種這些葡萄園.
等下我們會看多一點葡萄園的相片.
所以在梳理谷.
你會看到有梯田.
有果樹.
你會看到這些梯田是怎樣造成的.
需要花很多人工.
將本來是一個自然的山谷.
然後就挖了一級一級.
然後剩下的石頭.
就不扔掉那些石頭.
那些石頭就放在每一個梯田的邊邊.
然後放在那裡.
就能夠將這些泥土保留.
因為由高處沖下來的泥土.
這些很重要的肥沃的泥土.
就會承載在這些梯田上面.
所以你看到一級兩級.
一級兩級.
三級四級五級.
一直上去.
所以花很多人工才能夠造到這些梯田.

$^{81}$要好幾年的時間.
不久之前在那撒勒這個地方.
那撒勒.
那撒勒當然不是中央山區.
但也是一個山地的地方.
在加利利的山地.
那撒勒.
在這裡.
就是耶穌長大的地方.
你放大的時候.
這個就是一個山區的地方.
剛剛在山脊的後面.
就是一個小鎮耶路撒.
那麼多年都沒有變過.
都是一個在一些大的公路.
外面的地方.
一個支流的地方.
所以是一個小的小市鎮.
我記得曾經說過.
這個Villamares.
就是這條黃色線.
就是由大馬士革.
一直下來到加利利海的北邊.
然後經過哪個地方.
在海邊.
耶穌傳到的地方.
就是Capernaum.
這個加帛隆.
然後加帛隆繞道過來.
繞過來.
因為這裡有些山谷比較進修一點.
比較崎嶇一點.
所以那條國際的商路.
就會繞過這個山崗.
然後走過來這邊.
一直向沿海加密山這個地方進發.
然後穿過加密山.
就到了沿海平原.
一直下去下面就是埃及.
這個那撒勒.

$^{121}$就是在這個大的公路的旁邊.
所以是一個小鎮.
在那撒勒好幾年前.
去到做法國的時候.
有一個地方現在叫做那撒勒村.
Nazareth Village.
是一個基督教的事工的地方.
他們買了這個地方.
然後做了一些法國的工作.
找到一些舊約時期有的梯田.
這裡就是他們那撒勒村莊.
還沒成立的時候.
開始想去那撒勒村莊.
是做什麼的呢?.
是一個很好的事工.
是一班事工.
他們將那撒勒山邊的地葉.
重建成第一世紀的村莊.
將它重建出來.
遊人可以去那撒勒村莊.
可以參觀一下.
然後可以知道.
到底第一世紀的田園是怎樣運作.
所以是一個很有意思的工作.
今天還在那裡.
所以他們剛開始去開墾這個土地的時候.
開發的時候.
他們就找到一些舊約時期的梯田.
看到嗎?.
這裡有一支梁度的尺.
考古學家經常都會動在那裡.
一格就等於多少半米.
有時候是一米.
就知道這裡有多高.
他們發掘的時候.
就看到有些石頭放在一起.
這些很明顯就是人工留下來的.
上面也是.
所以他們就在這裡發掘.
了解那時候的梯田是怎樣做的.

$^{161}$將這些實物找回來的時候.
重構的時候就更加清楚.
到底是怎樣去做梯田.
你看到這個山坡.
由上到下.
有幾級的梯田.
如果將它重構的時候.
就會是這樣.
剛才梳理局看的梯田都是一樣.
本身就是一個斜坡.
然後農夫就在斜坡那裡.
挖出了.
作出這些硬的石頭.
然後將這些石頭保留.
所以第一步是怎樣做.
首先在床岩.
Back Rock那裡.
斜坡那裡就挖起這些石頭.
用鐵鎚慢慢挖起.
然後就形成台階.
一個樓梯.
然後挖出來的石頭.
正如我所說的.
是不會扔掉的.
大的石頭就特別放在樓梯邊那裡.
以致泥土.
侵蝕了的泥土.
風化了沖下來的泥土.
這些都是肥肉的泥土.
它就會盛載了.
一代又一代.
一年又一年.
盛載了這些很好的土壤.
所以第一就是先挖了這些樓梯.
第二就是將這些石頭推起來.
做一個護土牆出來.
一個擋土牆出來.
然後保存了這些泥土.
第三就是在梯田上面可以種植.
種一些果樹或者葡萄園等等.

$^{201}$所以格式或者形狀就是這樣.
我們看回這個拿撒克村莊.
在這些梯田附近.
就找到這些走柵欄.
什麼叫做走柵欄呢.
上面你看到的是一些葡萄園.
葡萄園就是葡萄的枝.
需要一些棚來依靠.
它們就長在這些棚裡.
每逢年底的時候.
他們要去收割這些葡萄.
將這些葡萄一抽一抽放在走柵欄的上面.
你看到這裡挖出了一個正方形的走柵欄.
然後他們放了這些葡萄在這裡的時候.
旁邊有個洞.
有時候會搭一些棚.
讓人可以拉著.
然後他們就會用腳.
是赤腳的.
來踩這些葡萄.
以致這些葡萄就能夠變成葡萄汁.
為什麼要赤腳呢.
因為他們不會穿鞋子.
草鞋也不會穿.
因為葡萄的核是苦的.
所以踩的時候.
將葡萄的果實變成果汁的時候.
就不能踩爛那些種子.
所以每個人都要赤腳.
所以幾個人就赤腳在這裡踩.
踩了之後就變成果汁.
在這裡你看到有一個小的坑.
這些葡萄被踩過之後.
這些果汁就會從坑裡滴到這個地方.
這裡就是裝起這些葡萄汁.
他們就會裝在啞瓶裡面.
或者裝在啞罐裡面.
然後裝起這些葡萄汁.
就會變成將它再發酵.
很多時候就會變成葡萄酒.

$^{241}$在這裡更近的地方.
你將葡萄放在這裡.
踩了之後.
\newpage



\section{}
\label{sec:26WK5R9cMYM}
\textbf{宣教路 - Steve Tan}
\newline
\newline
連結: \href{https://youtube.com/watch?v=26WK5R9cMYM}{\texttt{https://youtube.com/watch?v=26WK5R9cMYM}} ~~~~ 語音日期: 2022-10-25
\newline
\newline
\hyperref[sec:Pgm9dzwLc3c]{\small{< < < PREV SERMON < < <}}
~
\hyperref[sec:index_chronic]{\small{[返順時目]}}
~
\hyperref[sec:index_scriptual]{\small{[返順卷目]}}
~
\hyperref[sec:mXwPLGGYFwA]{\small{> > > NEXT SERMON > > >}}
\newline
\newline
$^{1}$影子們今天真的很開心有機會在這裡.
和你們分享我們的歷程.
我們一家下個月15號可以啟程去台灣.
我們很開心.
好像Mark Sir這麼說.
但同一時間很傷心.
因為我自己兩歲來到芝加哥.
在這裡住了36年.
所以很多親人.
很多朋友都已經.
離開了.
所以要從頭開始.
多謝今天教務同工給我們機會在這裡分享.
但是先給個警告.
因為廣東話不是我用得最好的話.
在台上說話.
所以有時有些說話或者粗魯點.
你見諒.
我們開始.
我們宣教歷程.
我們想和你說.
我用了至少10年時間去預備.
這裡說你要注守.
你要知道我是神.
這10年.
這個經文常常在我腦袋.
為什麼是這個經文呢.
我同一時間用這個經文.
讓你看到我們.
如何經過這10年去預備.
特別是在疫情的時間.
這個經文是如何幫助我們.
但是我們看看.
在詩篇46篇.
為什麼上帝會和以色列人說這句話呢.
是不是一個很安靜.
已經平定的時間去說呢.
你要安靜.
通常我們見到.
我們說靈修你要安靜.

$^{41}$去同神有一個好的時間.
有機會有時也會用這一節.
但其實我們看一看.
因為時間的緣故.
我很簡短地去解釋.
你見到46篇.
一至三節.
其實第三節那裡.
其實你見到最主要是有天災的.
但是第一節就說.
你不需要怕.
因為耶和華是我們的避難所.
是我們的力量.
是我們隨時的幫助.
所以有天災.
你是不需要怕的.
和以色列人說.
4至7基本上.
你在第六節見到.
它的總結就是.
有戰爭的.
打到來.
但是耶和華都說.
你不需要怕.
因為上帝是住在這個城市.
上帝的城是什麼.
耶路撒冷.
他有什麼時候住在那裡.
你都不需要怕.
這個城是一定不會倒的.
所以開頭第一部分是有天災.
第二有戰爭.
最後他在那裡再描述一件事.
第八和第九節就是.
他有能力可以停止戰爭.
如何停止.
他可以毀滅這些弓箭.
可以銷毀這些盾牌.
最後的高峰就在第十節.
你要駐守知道我是神.

$^{81}$為什麼在一個戰爭的背景.
他會說你要駐守.
其實他和兩個方面的人說.
第一和敵軍說你要駐守.
為什麼要和敵軍說你要駐守.
他要讓敵軍知道.
你怎麼打都沒用.
其實我是一定贏的.
因為我是上帝.
你不要浪費心機.
第一和敵軍這麼說.
第二和儀式里的人說.
你不需要怕.
你是一定贏的.
因為我是你的上帝.
我是創天造你的上帝.
你不需要怕.
你看看你認清楚.
我幫你做這件事.
你認清楚誰是上帝.
所以這個駐守就是.
你們醒醒吧.
你們認清楚.
你要知道.
這個知道不只是我們讀神學.
在聖經里看到上帝是怎樣.
不是.
這個知道通常是經歷到上帝的.
經過這十年.
我經常用這個經文.
詩篇寫六篇.
經歷到上帝.
他是真正的上帝.
我想趁這個機會和你們分享.
我們這十年的歷程.
第一.
我第一次去台灣是2009年.
因為我有個弟兄在香港的婚禮.
令我有機會去台灣.
但是我去台灣.

$^{121}$因為我想知道上帝在亞洲做什麼.
唯一一個地方可以有這個聯繫.
就是在台灣.
我不認識這個牧師.
這個叫林維光牧師.
在台灣已經做了20年.
他是宣教士.
我去到他接納我.
他最後那一個星期.
我跟著他.
他去哪裡.
他做什麼我都跟著他去做.
我看到宣教士真的很辛苦.
由早上八九點.
我跟著他到晚上兩點鐘.
他在家裡繼續什麼都做.
他跟太太說.
Steve 你不如過來和我一起去事奉.
我心想.
我剛剛畢業.
我有小雞蛋.
我哪裡有資格和你去事奉.
那時候我單身.
我說不行.
回來了.
但是我太太那時候.
我們還沒結婚.
我還沒去約會.
我2008年.
我太太已經奉獻了去全世界事奉.
關什麼事.
但是關事因為.
我2011年和她結婚.
因為我沒想過和她結婚.
因為我有一次說過.
我忘記後面說過.
誰和你結婚了就死定了.
因為那時候我們性格不是很好.
我們對立.
就是我.

$^{161}$結婚一年後.
我們去2012年.
我去宣教的很大型的開會.
我心想.
不可能.
我不會去宣教.
這件事不適合我.
但不是因為講完說得多好.
因為去到一個部分.
所有不同宣教的機構都在那裡.
他們在那裡招攬人.
我就問.
你們做宣教士有什麼條件.
沒有.
其實很多時候你想用你自己厲害的部分.
你可以用那部分去幫助宣教.
我說不是要做牧師還是怎麼樣.
他說不是.
最好有神學的背景.
去讀書.
我叫他這樣.
因為以前我每次宣教士回來教會.
都是牧師.
他說牧師不太適合我.
之後我說原來是這樣的.
我同一時間.
我說上帝.
我開口.
但你開路給我.
你給我一個感動.
同一時間在那裡.
我也認識了另外一個宣教士.
在那個會.
也是在高雄來的.
他也說.
我之後去探望他.
他也說.
Steven出來了.
其實我們台灣高雄很需要人.
你可不可以過來一起和Steven.

$^{201}$唯一一個宣教士.
叫我們去的也是在台灣.
之後我們繼續籌備.
馬上我老婆辭工.
我很害怕.
因為我們沒了一半的工資.
我們的收入.
之後14年.
我老婆也說.
輪到你去讀書了.
有機會開始慢慢去修煉.
我心想.
我很害怕.
為什麼.
如果無敵神學院知道我是一個什麼人.
他一定不會收我.
因為我這個人.
是一個很大的罪人.
這樣.
最後我老婆也幫我填表.
我抵上去.
原來她收我.
我去神學院讀了.
在神學院讀了半讀.
同一時間做全職.
很辛苦.
我輪到我老婆.
她差不多畢業了.
17年.
她就說.
阿明.
阿思.
輪到你去辭工了.
我回家工作.
我說不是吧.
辭工.
因為這件事對我來說是很艱難的.
如果你是男人.
你會知道我明白什麼意思.
我的工作是我的身份太大的一部分.

$^{241}$我辭工那一天之後.
我說我所有的名譽.
我的聲望好像都沒有了.
我還要照顧兩歲大的女兒.
還要讀書.
真的很難過這一關.
有點不開心.
同一時間我去到神學院.
我以為她會上帝的知識很多.
但不是的.
她上帝的知識太大.
我都不知道.
她學到聖經的東西太多了.
同一時間.
我和我女兒一起過這段時間.
看著她幫她讀書.
和她一起過這段時間.
兩歲.
其實是我人生中最開心的時間.
因為我不能拿回這段時間和她陪她.
所以你現在問我多一次.
我不會再.
讓我選事業和照顧女兒.
我一定會照顧女兒.
這不用說.
同一時間.
我去了神學院.
其中一件事不是因為我有什麼碩士學位讀到.
最主要.
2017年.
我是一個神學生.
我是一個小學生.
團體導師.
但我常犯罪.
最慘的是犯完罪你才去團體做導師.
你試過這樣嗎.
有一天去到一個地步.
我真的受不了.
我說上帝.
我真的不想再這樣.

$^{281}$我是不是真正的基督徒.
我有點懷疑.
直到有一天.
我正在讀詩篇.
《沙爾篇》.
戴衛描述他自己犯罪的過程.
原來有人明白我犯罪的感受.
所以.
同一時間.
我不知道為什麼去到《約伯記》最後那部分.
約伯會審問神.
上帝為什麼我會受苦.
上帝就是上帝.
他的答復就是.
他不告訴你.
但他會問約伯.
如果我告訴你你知道嗎.
你明白嗎.
我創天造地的時候你去哪裡.
約伯不敢說話.
我就是同一時間.
同一個人的經歷.
我說上帝.
我不需要知道任何答案.
我要知道.
我要知道你與我同在.
你說你真的愛我的.
你救贖我.
你是愛我的.
我是你的兒子.
是夠了.
求你給力量我去相信這件事.
我做完那一個祈禱.
我真的看到.
有一把黑雲黑煙在我身上.
走出來了.
我從此之後.
撒旦沒有再來找我這五年.
我每天會做些什麼呢.
我擦牙之前.

$^{321}$我在那塊鏡子上畫.
我是上帝的愛子.
他很愛我們.
我的身份就是這個.
這件事每天都幫我去.
預備我自己的心.
在世界的屬靈的伸展.
這樣.
2019年我們就畢業了.
然後我們去了世界的屬靈的機構.
去參加這個機構.
同一時間參加這個機構的時候.
我們就.
2019年的6月.
之後開始籌款.
我去到2020年3月開始疫情.
大概籌了三成.
之後就疫情.
怎麼辦呢.
沒有人知道怎麼辦.
這個真的是很大的考驗.
將來是怎麼辦.
沒有人知道 看不到.
但是我可以做些什麼.
糟了 台灣又要排隊了.
但是台灣有不同的.
全世界都在經歷疫情.
除了台灣.
但是我繼續在那裡.
在芝加哥籌款.
同一時間學國語.
去到2021年大概5月的時候.
我就有機會去.
2021年去台灣.
突然台灣開始有疫情.
我說有沒有搞錯.
但是一件好事就是.
我去到一個地步.
我籌了八成.
我找不到人.

$^{361}$因為我不認識人.
我從小在一個教會長大.
我認識的人就是這麼多.
我還可以找到誰.
好就是.
無端端有一個老牧師.
在左邊白色衣服的老牧師.
他90歲了.
他打電話給我.
其實他經過機構找到我.
他說有沒有宣教士去台灣.
我們是唯一一個宣教士.
在中美和東邊.
我們說你有沒有興趣過來紐約.
去探訪我們教會.
我說怎樣都好.
現在都不知道.
可以去哪裡籌款.
過去了.
在牧師那裡住.
還有去認識他.
他原來1960年.
在台灣那裡做宣教士.
我真的很佩服他.
現在90歲了.
依然在那裡做牧師.
今年我再去探望他一次.
剛好8月尾.
同一時間因為他要下來.
因為他找到新牧師在這個教會.
很小的教會.
其實三四十人.
八人教會.
在農村的地方.
我去到我們和他吃飯.
或者怎樣和他聊天.
認識一下大家.
他們很多時候都說.
多謝你們.
你們改變我生命.

$^{401}$我和我老婆說.
我做什麼.
我同一時間.
都有這種想法.
見到不同的人.
上帝擺在一起.
真的有不同的後果.
現在他們是我們其中一個.
支持我們的教會.
去到7月之前.
今年1月.
我去到那裡.
我都有一點擔心.
但上帝都和我說.
如果你記得《馬太福音》六章.
最後那裡.
你不需要憂慮.
那些雀鳥.
那些花.
是誰照顧的.
是上帝.
你不需要擔心.
我說上帝.
很難的.
我現在真的在這個情況.
我真的很難不擔心.
但他和我說.
你相信我吧.
我相信了他.
之後繼續去.
後來暫時有份工作.
在培德補習.
去幫街坊申請福利.
有機會用我的國語.
這樣.
直到7月台灣開了.
台灣終於開了.
等了我兩年多.
台灣開了.
我們8月就開始申請了.

$^{441}$8月1日.
第一個星期.
申請了.
怎樣呢.
他說要三個月.
很久的.
你知道很多人在排隊.
我說不要緊.
我和辦事處說.
台灣辦事處說.
朱子嘉哥.
那好吧.
你知道.
我.
接著.
他打電話.
他寫了個電郵給我.
2號.
他說你們的簽證可以了.
下個星期4天來拿.
8號.
這麼快.
一個月就搞定了.
我和所有的頂線說.
我們拿到簽證了.
下個星期可以拿了.
很高興.
真的很高興.
但是.
下個星期2.
6號打電話給我.
不可以.
你們.
我的上司要見一下你們.
因為你們在大陸出生.
他直接和我說.
我說.
我明白.
這個時勢.
我明白.

$^{481}$沒辦法.
這件事我改變不了.
我大陸出生.
是大陸出生.
雖然我來了36年.
我入籍都差不多30年了.
我說好吧.
算了.
去見一下他.
我星期五去見他.
星期五他在那裡.
和我面試一下.
他說你知道.
你現在這個時勢.
你也是在中國出生.
我知道.
這麼久關我什麼事.
但是同一時間他又問我.
你們去台灣做什麼.
你們有沒有經驗.
那些東西周圍問.
我就照樣回答.
我當然沒有發脾氣.
同一時間.
我就說.
來的週末.
10號 11號.
我就說.
我真的很累.
我等了兩年多.
現在我真的很擔心.
很擔心.
因為我覺得這個人.
是令到我.
可以一無所有.
我這麼多年的部署.
他可以說不批評我.
我說我怎麼和弟兄姊妹.
鄉親父老叔伯兄弟交代.
我真的很害怕.

$^{521}$我真的害怕到.
我說同一時間.
我放下我的子女.
我經常在29街Huston.
這個公園.
因為我經常和上帝對話.
我說上帝.
祂令我記不起.
《馬太福音》第六章.
你不需要怕.
我搞定了.
既然是這樣.
上帝.
我拜託你這個.
15號星期四.
奉耶穌的名.
拜託你給我簽證.
我不再玩了.
我真的這麼和他說.
我真的做了祈禱.
然後他立意打電話給我.
你教會是回哪裡的.
你給我電話.
他要讓我知道.
他打電話給教會.
證明我是不是真的.
還有神學院.
大禍.
又來了.
13號.
我把資料都給他了.
13號.
14號.
他無端打電話給我.
你知道15號.
我祈求這個土.
他說我再上去見你們.
這樣.
我上去了.
這次他說不好意思.

$^{561}$你知道我們.
處長叫我們去.
查得緊實一點.
因為中國那些.
明白的.
然後他給了簽證給我.
然後他跟我說.
其實處長和我們.
有幾個都是基督徒.
回到我們的鄉下.
幫我們傳這個福音.
我們早說.
大佬.
真的.
所以你要明白.
在這個環境.
我有什麼可以做.
我做不到任何事.
知道嗎.
所以我明白.
你要住守.
你要知道我是上帝.
我很明白.
我那一刻.
拿了簽證.
我所有的壓力都沒有了.
我的血壓高都沒有了.
上帝想我說什麼呢.
上帝和我說.
我要住守.
我要安靜.
你要明白我是真的誰.
我才是上帝.
這十年他要將我所有.
我覺得重要的身份.
破碎了.
我的身份最重要.
是我再記錄.
要將我所有的知識.
我覺得是行得通的東西.

$^{601}$毀滅了.
他要讓我知道.
他才是上帝.
他會幫我搞定.
他真的讓我看到.
所以這十年.
他真的預備了我們.
讓我們知道.
誰是上帝.
才帶我們出去工廠.
最後我想和你說這一句結尾.
他說.
你想找到安全.
在上帝那裡.
你要先不要在世界那裡找到安全.
沒有東西在世界會給你安全.
上帝才會給你的.
我做基督徒.
上帝都說.
你要獻上所有.
我其實覺得.
如果我不是獻上所有.
我其實不是真正做基督徒.
那我不如不要做了.
我不想做一個半假半恐怖的基督徒.
所以上帝說.
你要出去做宣教.
我馴服.
我想做真正的基督徒.
最後就是我們在高雄.
台灣做什麼.
最主要是做.
一個出世的嬰兒.
要餵奶.
什麼都不懂.
我們想去服侍.
特別是初生的基督徒.
我們想陪著他們成長.
幫他們長大成人.
做一個真正成熟的基督徒.

$^{641}$多謝各位弟兄姊妹.
和教牧童工.
給我們這個時間分享.
在這個時候.
不如我們為Steve和Shirley.
和他們整家人禱告.
Shirley我們請她上來.
我們都為宣教士.
當他們快要離開.
要去台灣的時候.
我們要禱告.
我們同心單祿祈禱.
上主我們多謝你.
因為你是那位.
願意去猜探的上帝.
你在我們的生命當中.
你不斷地作事.
不斷地感動和帶領.
在聽見Steve和Shirley的經歷.
他們在你的裡面.
在聖靈裡面.
他們有的感動.
一直沒有磨滅.
而是你不斷地.
每一步來帶領.
為他們開門.
為他們關門.
主我們獻上感恩.
你在他們的生命當中.
在他們心頭裡面.
所擺的志向.
給他們的感動.
一直沒有磨滅.
直到今天.
主要他們要去到.
預備去到台灣.
去宣教工廠的時候.
神我們求禰.
在每一方面來顯出.
禰自己是那位.

$^{681}$不需要令他們擔心.
禰是那位願意去供應.
願意去為他們預備的神.
主我們求禰上帝.
每一步來帶領他們.
特別是他們帶著兩個小的.
一家四口.
他們要去到台灣.
一個新的工廠.
新的環境.
主禰的適應.
我們求禰親自與他們同在.
無論在語言上.
文化上.
社會上.
環境裡面.
禰每一方面來加力給他們.
求禰讓他們在工廠上.
他們更多的學習.
經歷禰自己.
是那位加力給他們.
成為他們的避難所.
去拆解他們的上帝.
主我們知道宣教士.
去到工廠的時候.
會有不同的困難.
也都需要吃苦.
主我們求禰去建立他們的心智.
叫他們能夠去吃苦頭.
叫他們遇見困難的時候.
不會退縮.
反而來到去親近禰.
在禰身上來到去支取力量.
主我們求禰都去保守.
讓他們去到工廠的時候.
和其他的同工一起來到去.
彼此的配搭.
彼此的合作的時候.
上帝禰自己親自與他們同在.
讓聖靈藉著他們.

$^{721}$帶領人來認識禰.
藉著他們將福音.
在台灣這個地方.
有很多人沒聽過耶穌.
不認識耶穌的.
他們能夠去認識基督.
主我們求禰讓教會.
我們本地的教會.
不同各地的支持他們的教會.
我們在禱告上面.
來到去繼續的去支持住他們.
用我們的民安.
用我們的禱告.
用我們的支持.
來到去使他們的工作.
能夠得到免禮.
我們使他們在宣教的一程裡面.
能夠更加體會到上帝禰的大能.
上帝禰是值得我們敬拜.
值得我們盡心盡力去服侍的.
主也幫助我們教會.
能夠成為一個猜探.
宣教的教會.
幫助我們每一個.
能夠成為一個福音的管道.
將你自己的祝福.
你自己的恩典.
能夠向其他人展示出來.
我們這樣祈禱.
奉主耶穌基督的名字.
再次多謝Steven和Shirley.
\newpage



\section{}
\label{sec:mXwPLGGYFwA}
\textbf{聖經地理 EP7 - 地理與釋經 - 交通-1 (陳崇基牧師)}
\newline
\newline
連結: \href{https://youtube.com/watch?v=mXwPLGGYFwA}{\texttt{https://youtube.com/watch?v=mXwPLGGYFwA}} ~~~~ 語音日期: 2022-11-01
\newline
\newline
\hyperref[sec:26WK5R9cMYM]{\small{< < < PREV SERMON < < <}}
~
\hyperref[sec:index_chronic]{\small{[返順時目]}}
~
\hyperref[sec:index_scriptual]{\small{[返順卷目]}}
~
\hyperref[sec:PWN_xBip_GQ]{\small{> > > NEXT SERMON > > >}}
\newline
\newline
$^{1}$我們講到迦南地的地理和聖經地的地理.
我們在不同方面不同角度去看.
所以這次繼續講聖經地理.
這個地理告訴我們這個地圖是以色列.
以色列或者巴勒斯坦地.
或者我們叫它中間之地.
在海洋地中海和沙漠中間.
約旦河東這個地方.
過一點就是今天的沙地阿拉伯北邊.
就是沙地阿拉伯的沙漠.
是一個像隧道一樣的走廊.
一條橋.
最現在的以色列國.
最寬就是115公里.
72英里.
所以由左至右海邊開車去到死海.
由這邊開車去到這邊.
都是一個多小時.
公路來說是一個多小時.
所以很窄.
一條地橋.
由最上面的北黑門山.
到下面今天這個地方紅海另外一個灣.
這裡是420公里.
263英里.
所以你看到是一個長長的.
左邊就是地中海.
右邊就是東邊的沙漠.
叫做Eastern Desert.
所以在海洋和沙漠中間.
這裡有的氣候都是屬於地中海氣候.
主要來說地中海氣候.
就是夏天.
大概4月至9月.
夏天是乾和熱的.
10月尾至明年3月4月.
就是雨季.
比較冷的.
所以有濕有冷.
大概有兩個季節.

$^{41}$春夏秋冬.
感受上只有兩個季節.
一個是夏天是乾的.
一個是冬天是冷的.
是濕的是雨季.
所以這個就是迦南地.
這裡有另外一個地圖.
由打直地分到迦南地.
由左至右.
有不同的區域.
或者自然地理的部分.
自然地理的區域.
剛剛說了地中海.
就是Rainy Sea.
上岸之後.
這裡有綠色的一段.
就是平原.
Coastal Plain.
就是沿海平原.
上面的一隻角.
就是今天的以色列第三大城市.
叫做海法市.
上面就是加密山的地方.
所以這個沿海平原綠色的地方.
由北至南.
就是一個平原的地帶.
所以比較廣闊一點.
行軍行商路線會比較方便一點.
然後我們慢慢向東邊.
就是向這邊.
由左至右.
由西邊至東邊.
過了這個沿海平原.
這個地方很重要.
在聖經裡面經常提及.
叫做低山地.
低山地意思就是說.
它不是平原.
那麼廣闊.
但也不像山地那麼多山和谷.

$^{81}$那些山比較矮小一點.
還有山谷會比較寬闊一點.
這個地方叫做低山地.
低山地就是夾在沿海平原和山地中間.
山地就是由上面.
這個是加密山的地方.
加密山上面有一片綠色的地方.
就是以色列的平原.
或者以色列谷.
這裡一直下來叫做hill country.
或者central hill country.
或者中央山地.
這個地方最高大概有海拔二千多呎.
耶路撒冷就在這裡附近.
然後過了這個山脈.
中央山脈.
然後下去就會是一個曠野的地方.
因為這裡就是分水嶺.
在中央山脈.
然後過了中央山脈的分水嶺之後.
就不會有雨水.
所以天然來說.
就一定會有曠野.
這裡就是.
《史篇63篇》所說的.
乾旱疲乏無水之地.
猶大的曠野就在這裡.
然後水平線一直下降.
就去到死海的平原.
或者叫做阿拉巴列谷.
這個列谷是很多年前.
地質學的時間.
有百萬年前.
兩塊大的板塊撞在一起的時候.
大家有些張力.
然後就撞起來.
就造成了這個列谷.
這個就是所謂的約旦河谷.
由上面這個黑門山.
溶雪下來.

$^{121}$是水源的地方.
然後經過這些河流.
下到加里尼湖.
加里尼海.
然後加里尼海的南面.
就會出來.
就是經過這個河.
叫做約旦河.
就下到下面的死海.
死海很明顯的原因.
就是因為進去之後.
就死到一條了.
因為水下去之後.
就不會再分到其他地方.
只會蒸發.
然後過了這個列谷.
或者死海的列谷.
約旦河列谷.
上來就好像對面的山地那麼高.
有二千多呎.
這個我們叫做Transjordanian Plateau.
或者是約旦河東的高原.
所以兩邊都差不多高.
中間就有一個死海的列谷.
約旦河谷分開.
這個主要是由左至右的區域.
這些是天然造成的地理地帶.
有不同的山.
有不同的谷.
這麼多年都是聖經地理的桌面遊戲.
好像大富翁.
大家玩過大富翁四遍.
最開頭那邊一定是最便宜的.
最後那邊是最貴的.
所以就有一個桌面.
定了裡面所發生的歷史.
裡面有的神所做的事蹟等等.
上下.
北邊和南邊就有一些.
小一點的區域.

$^{161}$北邊就是加利利的區域.
有下加利利.
比較低一點.
還有上加利利的地方.
高山多一點的地方.
中間這裡就是加利利湖.
也是全地球最低的淡水湖.
全地球海拔最低的淡水湖.
再去右邊就是約旦河高原的伸展上去.
北邊.
我們今天叫高山高原.
以前聖經就叫做巴山地.
巴山牛等等.
詩篇有時候會提過.
舊約裡面有提過.
南邊下來.
雨水就越少.
所以一直下來就是去到.
這個地方.
這個就是主要的一些.
就是地理上的一些分區.
有了這個桌面遊戲.
我們就知道.
怎樣去處理.
裡面的一些經文.
或者裡面的動態.
我們先再講一下交通的問題.
我們上一次講過.
迦南地以色列巴山地.
是整個近東的世界.
古代的世界的中心.
那個地橋.
Land between.
歐亞非的中間之地.
這個地方既然是中間之地.
所以有很多的路會經過.
很多的路經過.
有藍色這裡一條路.
上面北邊.
我沒有翻譯到.

$^{201}$但是從這個.
很遠的米索不達.
再來到這裡.
就是大馬士革.
或者大馬式.
然後接著藍色這條線.
就是跟著一些自然.
通常這些路線都是跟著一些.
自然的山嶺.
山脊這樣走.
所以從南邊下來經過這個下所.
再下來.
加里尼湖的.
左上角.
岸邊.
加里尼湖的西南的.
湖岸.
這裡就有一些.
重要的市鎮.
其中一個.
就是加巴隆.
就是耶穌.
傳道的開始.
基地的地方.
在加里尼湖的左上角.
西南角的岸邊.
是一個碼頭.
一個漁民的碼頭.
然後繼續下來.
經過這個耶斯列平原.
就去到這個米吉多.
就是米吉多.
大家可能聽過.
在啟示錄裡面.
亞米吉多頓.
有一個傳奇.
就是從那裡出來的.
因為那裡經常打仗.
因為這是一個很重要的關口.
米吉多.

$^{241}$過了這個加密山.
就是Aphrodite.
中文叫阿弗.
在經文裡面.
也有不同的經文出現過.
然後下來.
下面有一個小的海港.
叫做約帕港.
在舊約裡面.
和新約.
比特.
看到有一塊布.
掉下來.
有一些不潔淨的食物.
神的使者叫他吃.
在約帕港這裡發生.
然後下去下面.
有加薩.
再下去就是經過.
這個沙漠的地方.
然後就去到.
利羅河的沙漠洲.
就是埃及.
有一個很重要的.
一個研究.
是講到這個迦南地的.
交通.
我在這裡就將.
那個研究的.
結論在這裡和你分享.
首先我們知道.
這些路線.
等等.
講到古代的路線.
古代路線.
我們想一下當時的人.
他們的步行速度是怎樣.
就不像我們今天.
有很大的公路等等.
一天一個人步行.

$^{281}$大概如果你真的.
這樣八小時十小時這樣步行.
大概你可以cover到.
二十英里.
當然要視乎地形.
地勢.
如果太崎嶇就會慢一點.
如果是沿海平原.
當然是一條直路這樣去.
就會走快很多.
在波斯時期我們不同的文獻.
我們會見到一個普通的旅客.
平均他每天.
步行大概十八英里.
在不同的文獻.
特別是希臘的歷史學家.
講到波斯帝國的一些細節的時候.
從一個京都.
叫波斯波里斯.
穿越整個波斯帝國.
去到西面的外禽海.
大概需要.
只是走路.
需要三個月的時間.
大概九十天的時間.
我所指的地方在哪裡呢.
在波斯波里斯.
即是今天的伊朗的地方.
大概在波斯灣.
這裡上面.
就是波斯波里斯.
經過米索不達米亞.
再去到今天的土耳其.
最西面的地方.
就是外禽海和希臘半島的世界.
由剛才波斯波里斯.
即是波斯帝國的中心的地方.
去到西面.
今天的土耳其.
西面這個地方.

$^{321}$大概需要三個月的時間.
就這樣步行.
九十天.
但是騎馬當然是很快.
在一些波斯帝國的文獻.
或者有關他們一些交通的信件等等.
我們知道九天或十天.
不同的馬去接力的時候.
好像我們華人.
中國以前一些古代的方法.
有驛站.
不同的驛站去接力.
然後送錶.
送信件.
就會很快去到另外一個地方.
大概九至十天.
走路就三個月.
然後第一世紀.
有一個歷史學家是猶太人.
叫做約瑟夫.
他說由加里.
南部去到耶路撒冷.
看到在哪裡嗎.
加里南部.
去到下面耶路撒冷.
還需要多少.
需要三天.
騎馬或騎小驢.
可能要兩天.
所以想起約瑟瑪利亞.
聖誕的故事.
他們在那撒勒.
去到哪裡.
巴拉里海.
那撒勒本身.
約瑟在那裡打工.
應該在那裡打工.
我們知道那些背景.
所以他們趕去那裡.
因為天下要報名上賊.

$^{361}$羅馬帝國的時候.
他們經過山地.
或者經過約旦河谷.
然後上去耶路撒冷附近.
再下去巴拉里海.
巴拉里海這個地方.
所以你看到.
很艱辛.
大概要三天的時間.
他們應該沒有小驢坐.
因為比較窮.
一天可以騎在馬或驢背上.
大概25英里.
如果沒有什麼困難.
駱駝的話.
每天可以走25英里.
並且可以承載很多東西.
所以人類明白.
如何去信仰駱駝的時候.
這個時候大概是.
主前一千年的時候.
開始已經可以很準確的.
或者很掌握到.
如何去信仰駱駝.
開始可以用駱駝了.
減少旅客要經過的.
旅行的時間.
這個研究.
就是一位學者.
在90年代初的時候.
叫做David Dorsey.
他做的.
專門研究古代以色列.
迦林地這個地方的路線.
有什麼公路.
或者支流路線.
大小路線等等.
他就.
開著電單車.
周遊列國.

$^{401}$在以色列這個地方.
80年代90年代的時候.
就周山跑.
然後就做了他的博士論文.
就是講.
The Highways of Ancient Israel.
他就用這些地勢地形.
去重構.
以往在鐵器時代或者舊約時代.
一些主要的幹線和支線.
他就重構了245條.
鐵器時代的主要支流路線.
最多交通的地方.
就自然.
剛才我們所講的.
就是沿海平原的地方.
沿海的路線.
沿海的平原.
他總共找到62條.
是中向的.
由北至南.
南至北.
包括國際路線.
和其他分支的道路.
但是道路網絡.
是伸延到整個迦林地.
不同的地方.
有9條路線.
就是從沙龍.
這個地方.
通往到北面的.
Becca Valley.
這個地方在哪裡呢?.
Sharon.
沙崙玫瑰.
沙崙玫瑰就是.
沿海平原.
就快到米吉多.
就快到.
加密山這個地方.

$^{441}$這個北面的沿海平原.
叫沙崙平原.
再上到黎巴嫩附近的.
Becca Valley.
他找到9條.
不同的路線上去.
左邊中間右邊.
等等.
他也找到33條國際路線.
International Route.
這些就是.
連接沿海的路線.
和約旦山地的北部.
所以在這些地方.
他也找到大小的路線.
你看這些數字.
就知道有很多.
這麼窄的地方.
居然有這麼多的道路.
34條中向的路線.
是穿過哪裡呢?.
猶太地.
或者薩瑪利亞山地.
薩瑪利亞山地就是中央山地的北面.
猶太山地就是.
中央山地的.
山脈的南面.
這個山地的地區.
5條南北的路線.
是在巴西部和低山地.
低山地就是這個地方.
然後14條本地橫路的路線.
是穿過加利利的.
14條穿過加利利.
這麼小的地方.
29條在薩瑪利亞附近.
然後59條是橫向的.
在猶太山地.
在這裡我問了其中一段說什麼呢?.
他說如此廣泛的國際跨區的路線.

$^{481}$不單見證了迦林地作為古代世界主要的迴路中心.
也見證了在舊約的時候.
在鐵器時代.
聖地的通行頻率.
太多人經過.
In and out.
又進又出.
並且告訴我們.
這個地方地橋就是當時世界的中心.
所以說到迦林地.
我們想起地理的動態.
就是communication.
是通話.
交通.
交通來往.
東南西北等等.
等於你在一些大的城市.
在廣州.
在上海.
在香港.
在北美的芝加哥.
這個地方.
等等.
這些大城市通常都是一些匯點.
東南西北周圍有不同的交通會經過.
所以這個就是迦林地的本質.
好,有望無.
\newpage



\section{}
\label{sec:PWN_xBip_GQ}
\textbf{聖經地理 EP8 - 地理與釋經 - 交通-2 (陳崇基牧師)}
\newline
\newline
連結: \href{https://youtube.com/watch?v=PWN-xBip-GQ}{\texttt{https://youtube.com/watch?v=PWN-xBip-GQ}} ~~~~ 語音日期: 2022-11-08
\newline
\newline
\hyperref[sec:mXwPLGGYFwA]{\small{< < < PREV SERMON < < <}}
~
\hyperref[sec:index_chronic]{\small{[返順時目]}}
~
\hyperref[sec:index_scriptual]{\small{[返順卷目]}}
~
\hyperref[sec:wlVkY6Ioxpo]{\small{> > > NEXT SERMON > > >}}
\newline
\newline
$^{1}$我們繼續看中間之地的交通,世界書樓的地方,會點的地方.
我們可以處理一下經文.
以賽亞書第十九章,這裡選了一些跟路線和道路有關的經文.
十九章,以賽亞書,講到將來耶和華救恩的日子會是怎樣呢?.
他說,到那一天,必有一條從埃及通往亞述的大道.
亞述要進埃及,埃及的人要進入亞述.
埃及人要與亞述人一同敬拜耶和華.
所以講到埃及和亞述.
到那一天,他就用了道路來作為一個比喻.
亞述就是這個地方,北面的米索不達,和埃及.
必定有一條大道通往大家.
就是說,兩邊的帝國不再打架,不再爭,其實就是通往大家.
講到將來的終末的時候,會發生的普世宣教現象.
所謂宣教的意象.
將來這些外邦會有一條大道通過去,大家會認識真神,耶和華.
以賽亞書第九章,那裡講到什麼呢?.
那受過痛苦的必不再見憂暗.
從前上帝使西布倫地和拿忽他尼地被藐視.
滅後將使玄海的路若但河外,外邦人的加里尼得著榮耀.
在黑暗行走的百姓看見大光.
住在死暗之地的人有光照耀他們.
這個就是以賽亞書.
看到終末來到的時候,利賽亞的日子,神的救恩臨到的時候.
會發生什麼事呢?.
上帝會使西布倫地和拿忽他尼地以前被藐視.
滅後在玄海的路.
在拉丁文來說就是via maris.
玄海的路,這是一條國際的路線.
剛才介紹了,米索不達米亞通到大馬士革.
然後到加里尼,經過加帕隆,再到米吉多.
然後玄海路線一直到埃及.
這就是聖經所說的玄海的路.
是一條實際的路,不是想像中的路.
而是在迦林地最快的路,最多支流的路.
剛才David Dawsey所說的研究.
最多的中路,中向的,南北的,都是在哪裡呢?.
Coastal Plain,就是在玄海平原.
他說將來在這條路上會看到大光.
約旦河外,這個地方.
約旦河之外,這個地方.

$^{41}$外邦人的加里尼,是中間之地的中間之地.
整塊的迦林地就是中間之地.
但在中間之地裡面也有很多十字路口,大大小小的.
其中一個就是加里尼這個地方.
雖然這裡山峰比較險峻.
但這條藍色的國際路線.
就會經過加里湖,左上角.
加勒比隆這個地方,然後繼續走國際路線.
再說說西布倫地和拿忽他利地在哪裡呢?.
我左邊這裡有一個十二支派的分地圖錄.
當然我們不計以東,不計摩亞,不計阿們.
我只是將它放在後面,這些打斜的字,不是十二支派.
但其他十二支派,我們就將它放在上面.
約書亞記裡面摩西五經,約書亞記和聶王紀.
等等,他們說到他們住在哪裡,支派分了甚麼地.
將它放在地圖上,這樣就很清楚了.
西布倫地和拿忽他利地,就是在上加里尼這個地方.
為甚麼這麼重要?.
因為西布倫地和拿忽他利地在舊約歷史裡面.
曾經發生一些事,對以色列人的故事有很大的影響.
發生甚麼事呢?大概在主前七百多年的時候.
主前七百多年,那時候已經分了南北國.
你記得所羅門聯合王國之後,上面就是耶羅波安.
下面就是羅波安南國,Jeroboam,耶羅波安在上面.
北國十個支派,下面就是猶大支派為主.
分了南北國,那時候大概主前八百九百年左右.
所羅門之後,分了南北國,上下南北國.
但到了公元前大概八世紀末的時候.
即是公元前七百多年的時候.
亞述帝國在那時候是他們最興旺的時候.
霸權最大的時候.
亞述在公元前七百多年的時候就開始進佔迦南地.
他們進佔的地方就在哪裡?就在北面.
因為亞述帝國行軍的話,他就要在北面下來迦南地.
他最終的目的當然就是要去埃及,要去佔領埃及.
但一定要佔領埃及的話,就必須要拿兵家必爭之地.
就是迦南地中間之地.
所以他下來的時候,首當其衝的第一班人.
第一個地方受到亞述國的侵佔的就是拿伐塔尼和希伯來.
所以所謂第一班地老的人.

$^{81}$在以色列的歷史裡面被外族攻陷了.
然後將一些人殺害,將一些人擄掠,然後帶走等等.
第一次就是在這個地方發生.
主前大概七百多年的時候.
所以以賽亞書就在講這個背景.
因為以賽亞講預言的時候,以賽亞先知到事奉的時候.
都是這段時間,大概在公元前七百多年的時候.
他就知道這個背景,因為以賽亞書的背景就是亞述帝國.
如果你翻一下以賽亞書裡面的一些細節.
有很多時候都是講到亞述帝國的軍事入侵.
那些人就很害怕,才叫他們不要拜人偶像,要依靠耶和華.
耶和華就會將這些受敵擊退,救恩就在耶和華那裡.
所以在西伯倫地和拿伐塔尼地是第一班被藐視的以色列人.
很慘,被擄,被流放,被放逐.
但是將來神要眷顧以色列人的時候.
黑暗中的百姓就會見到大光.
死刃裡面的人就會見到有光照耀在他們身上.
西伯倫和拿伐塔尼地.
所以你記得這段經文就是在馬太福音第四章裡面出現.
就是耶穌來到,耶穌聽見施洗約翰下了監.
就退到加利利去,他就由大地上回到加利利.
就離開了拿撒勒,拿撒勒就在這個地方.
然後在西伯倫和拿伐塔尼的邊界上.
他搬了地方,搬了去哪裡?.
加利利湖附近.
這時應驗才知以賽亞所講的,就是剛才那段經文.
在在死刃的會有光照耀他們.
從那時開始,耶穌就傳起道來了.
就是在加伯隆這個地方.
天國近了,你們應當悔改.
真是天國臨到,直至神的兒子彌賽亞來到.
所以回應了這個應許,這個承諾.
西伯倫地,拿伐塔尼地,在加伯隆的地方.
就見到有大光照耀在黑暗裡面被流放的子民.
神的恩典眷顧,會臨到.
這裡就是加伯隆的上空.
你見到就在加利利湖的旁邊.
我們由北面望下去.
這個加利利湖,這裡有碼頭,有漁船可以停泊等等.
裡面有些考古的發掘.

$^{121}$當時第一世紀,一些村莊和房屋等等.
是這麼大,一個很小的市鎮.
但是這個地方是在沿海的路線,國際的路線之上.
由上面這裡,高瀾高地,一直下來.
剩下這個關口,這個坳口,是最方便的.
去到一個國際的行商,行軍的路線.
就是經過加伯隆,然後下來莫大拉.
然後經過當時的加利利省的首都,叫做提比利亞.
Tiberias.
一直向南,來到沿海路線,Vienna Marist.
經過米吉多,加密山,然後到下面的沿海平原.
所以你見到,這些全部都有地理的實際動態.
這個是加伯隆,Capernaum.
帶大家去看一看.
這裡門口寫著Capernaum,Town of Jesus.
耶穌的市鎮,也有寫著.
這個地方現在是屬於天主教的物業.
因為他們很多年前就買了這個地方,然後開始發掘的工作.
所以這裡有個修道院,天主教的修道院.
你進去入場,就可以看到發掘了的東西.
有一個主侯四世紀的幽魂教會,猶太人的會堂.
主侯四世紀,即是公元後第一世紀.
這個就是幾百年後的遺跡.
我曾經說過,就是這個聖地的繁殖.
在聖地這個地方,近東這個地方.
有些地方是做了聖堂,是一些神聖的地點.
那個地點很少會變的.
因為近東的文化是很傳統的文化.
那個地方曾經是一個殿堂.
之後世世代代上去,雖然被拆毀,他們再建一個新的.
其他人去到那個地方,都會在那個地方,在那個點裡面去建上去.
所以有一個Continuity of Sacred Space.
不會亂放的,他們會知道那個傳統,然後承繼那個傳統.
特別是那個Sacred Space,Sacred Place的傳統.
所以你會看到這個白色的石灰岩所做的猶太人的會堂.
很大,有三道門,然後進去有很多的柱子.
有一個大的讀經室,聚會的地方等等.
雖然這個是主後四世紀,主後四百多年才建上去.
但是在這個四世紀的猶太人的會堂.
是建在一個一世紀的地基上.

$^{161}$所以你看看這裡,白色這個就是四世紀.
但是他建在下面一個黑色的,是當地的本地的石,玄武岩的石.
這些叫Berserk Stone,建在這個整齊的地基上.
以前那個可能是黑色的,小一點的.
但是那個地基是保留了.
所以之後的人都會在那個地方再建一個猶太人的會堂.
所以這個就是加勃隆的地方.
所以我們看到這個地方其實就是一世紀的地方.
另外一個在加勃隆找到的東西,我沒有放上去.
我們看回這裡就知道了.
剛才我們看這個主後四世紀的加勃隆的會堂,猶太人的會堂.
對面那幾間屋就是一個像是一個鑊形,我們叫UFO.
的一個天主教教堂.
為什麼會像是一個鑊形升起了呢?.
因為這個鑊形下面就找到一些早期基督教的建築,一些遺跡.
這個遺跡似乎就建在彼得的家附近,或者甚至是彼得的家那裡.
彼得之家.
福音書不是跟我們說過耶穌在加勃隆的時候,他回過家.
馬可福音有提過,他回家在哪裡?.
在這個考古的法國裡面,看到這個地方是特別.
又是sacred space的continuity.
以前可能是一個很簡陋的茅屋.
但是慢慢變成一個早期的基督教,大占庭時代的教堂.
所以你這樣估計的話,彼得之家可能就在這裡.
真的是彼得的家,就是在這裡.
我們主要今天想說的是什麼?.
就是沿海的路線,經過這個地方.
所以如果你有這樣的地理背景的時候,你就會發現這些不是亂來的.
在那佛塔里和西布倫地會看到大廣.
耶穌就在這裡傳道,就是在Capernaum這個漁港.
首先去到趙門徒,然後就傳道,就講天國的信息.
天國近了.
現在在衛星圖片都看到一個白色的猶太人會堂的遺跡.
下面就是一個第一世紀的地基.
下面是一個八角形的天主教教堂.
下面就是彼得之家.
最主要就是想觀察到這個地方是一個交通的要道.
所以耶穌選擇Capernaum是有他的意思.
不是亂來的,為什麼不選擇別的地方呢?.
他偏偏選擇這裡.

$^{201}$不單止當時歷史,這個地方是一個重要的行商路線.
在地理的處境裡面是一個行商的路線.
謝謝大家收看 再見.
\newpage



\section{}
\label{sec:wlVkY6Ioxpo}
\textbf{佳節倍思恩(四): 逾越節  - 陳崇基牧師}
\newline
\newline
連結: \href{https://youtube.com/watch?v=wlVkY6Ioxpo}{\texttt{https://youtube.com/watch?v=wlVkY6Ioxpo}} ~~~~ 語音日期: 2022-11-10
\newline
\newline
\hyperref[sec:PWN_xBip_GQ]{\small{< < < PREV SERMON < < <}}
~
\hyperref[sec:index_chronic]{\small{[返順時目]}}
~
\hyperref[sec:index_scriptual]{\small{[返順卷目]}}
~
\hyperref[sec:PVQ3X0QwXNk]{\small{> > > NEXT SERMON > > >}}
\newline
\newline
$^{1}$讚美你 祈禱奉尊明教.
好 今天我們繼續看這個系列.
就是看聖經的節期.
每逢講到節期的時候.
很多時候我們都會想起.
跟一些故事.
或者歷史的事蹟.
或者食物有關.
譬如我給你猜猜.
或者你想一想.
到新年的時候我們華人吃什麼.
年糕 沒錯 年糕.
或者.
蠔豉發菜那些東西.
我小時候最怕吃那些.
吃一天OK.
第二天又吃 第三天又吃.
吃到初七 那我就覺得很辛苦.
譬如還有什麼節日.
跟食物有關.
冬至.
湯圓.
還有呢.
糯米 還有呢 吃火鍋.
吃火鍋.
然後還有什麼節日.
感恩節呢.
是吃火鍋.
端午節吃什麼.
就不是吃火鍋.
我們只會吃火鍋.
講到聖經的節日.
其實都有很多跟食物有關.
譬如我們這次講.
預節.
預節跟什麼有關.
食物.
羊我們想到.
我們第一樣東西就會想到是羊.
事實上聖經裡面.

$^{41}$第一次守預節的時候.
就是跟羊有關.
在出埃及記你記得.
當耶和華要降十災.
最後那一災.
就是.
擊殺埃及長子的一災.
是上帝的審判.
去到這個.
所謂的強國.
然後以色列人住在埃及裡面.
他們是需要被拯救出來.
於是上帝就跟他們.
設立這個預節.
十二章第三節.
那裡講到.
預節的安排.
你們要去做什麼.
然後.
吃完這一餐飯你就要走.
十二章第三節.
他說你要告訴.
以色列傳媒眾.
本月初十.
各人要按照家庭.
娶羔羊.
一家娶一隻.
然後他就講到.
羔羊要被烤.
然後要在家裡.
吃.
然後就要將宰殺了的羔羊.
羔羊的血.
就要塗在哪裡.
門尾上面.
以致上帝的審判.
臨到國度的時候.
上帝就會.
越過.
門尾上面有血的.

$^{81}$家庭.
是一個拯救的故事.
他說你要這樣去吃羔羊.
然後.
你們要腰間束帶.
腳要穿著鞋.
這個很有趣.
因為猶太人.
起碼以色列人.
進到屋裡的時候.
他們就會脫下鞋在門口.
然後.
要上路的時候.
他們就會.
戴皮帶.
他們出門的時候.
很多時候都會戴皮帶然後穿鞋.
準備走的.
吃羊羔.
手裡面還要拿著杖.
特別是講到那些.
Walking stick.
你要行路的時候拿著那一張.
快快地吃.
然後還交代了很多細節.
我們沒有時間講.
吃無烤餅.
為什麼是無烤.
因為如果你做麵包的話.
你知道你等麵粉發酵.
要等很久.
要等好幾個小時.
然後你才能夠.
烘到麵包.
因為趕著走就不可以了.
所以他們在這個節日.
就要去吃無烤餅.
紀念上帝拯救他們.
離開埃及.
拯救他們.

$^{121}$以致他們不能夠不會被審判.
這是耶和華的預言.
Passover.
越過.
越過他們的門.
或者去到.
保護他們.
每年西曆的時候.
三四月.
就是猶太人的.
Passover的時候.
正是春分.
之後的.
第一個新月.
所以.
也是一個.
季節與季節之間過渡的.
一段時間.
當時特別講到.
如果你說耶穌的時候.
當時有很多散居.
流散在世界各地的.
猶太人.
他們就會回到耶路撒冷.
然後就會守節.
當時的人會很擠迫.
很裂漏.
這是一個最大最大的節日.
有點像我們華人的.
新年.
華人的一些重要的慶典.
或者好像.
美國這裡.
Thanksgiving的時候.
最忙的地方是什麼地方.
機場餐館.
華人餐館.
很多人走來走去.
猶太人就是這樣.
猶太人就會回去.

$^{161}$以前沒有一架飛機.
他們就會趕路.
去到耶路撒冷.
所以我們想起.
預節的時候就想起.
高陽.
剛才歌也有唱.
上帝帶領我們.
預月經過.
類似這樣的意思.
但你要知道在兩月中間的時候.
就是舊約結束了.
到新約還沒開始的兩三百年.
當時的猶太人.
他們想起預節的時候.
他們就不只是想起高陽.
他們想起整個package.
整個上帝的拯救.
會有些什麼呢.
你想一下有些什麼.
除了羊之外.
他們想起一連串的東西.
有摩西出現.
因為上帝呼召摩西.
去帶領他們.
好像一個牧羊人帶領他們離開埃及.
然後當然就是我們有所謂.
sacrificial lamb.
預節的高陽.
預節的高陽的血.
無烤餅.
還有想起什麼.
就是水的神蹟.
過紅海.
過完紅海.
整個package.
就看完一個曠野.
這麼多年上帝帶領他們.
是走過的.
在曠野又給他們嗎哪.

$^{201}$讓他們每一天能夠有東西吃.
什麼叫嗎哪.
manner.
大概可能聖經的描述.
是一些圓圓的像餅的東西.
其實我們都不太知道.
因為那個故事告訴我們.
嗎哪的名字是怎樣來的.
就是上帝降了一些.
好像麵包的東西給他們.
然後他們就不知道是什麼.
他們就說.
Mazet.
希伯文就是.
什麼東西.
什麼來的.
所以就叫做嗎哪.
所以其實我們不知道是什麼東西.
但總之他們就吃了這麼多年.
上帝就沒有間斷過的.
來供應他們.
我們說兩樣中間.
耶穌之前來到的時候.
他們就想起一連串的.
這些所謂預節會發生的東西.
或者pass over的時候.
他們會想起的東西.
一路由過紅海到離開埃及.
到這個曠野的時候.
他們就會想起.
有很多的故事.
甚至會加上一些故事.
加上什麼故事呢.
就來到耶穌來之前.
那些猶太人就會想像.
因為聖經確實有說.
他們過完紅海.
過完四十年的曠野.
進到迦南的時候.
嗎哪就不降下來.

$^{241}$他們入主迦南之後.
去到他們的家.
真是去到目的地.
就不再有那些什麼東西降下來.
然後他們就吃迦南地的土產.
所以他們就會想像.
天上不再降嗎哪.
為什麼呢.
因為原來摩西有條鑰匙.
當以色列人入主迦南之後.
他就將天上的嗎哪倉庫關了.
就不再降嗎哪.
所以他們就繼續想.
如果上帝真正的救恩.
最終極來到的時候.
應該會是怎樣呢.
就是摩西會天上的倉庫.
就拿著鑰匙再開.
然後就會降下來嗎哪.
然後我們就可以吃那些嗎哪.
就是天上降下來的糧食.
以致我們能夠有終極的救恩.
天上降下來的天糧.
Spiritual food.
如果是這樣想的時候.
我們就明白.
聖經裡面提到御節的時候.
那時候跟耶穌有關.
我們看一段經文.
我們只是選好幾節來看.
因為比較長.
約翰福音第六章.
約翰福音第六章第四節.
那時猶太人的御節近了.
可能還沒到或是已經過了.
總之是附近.
猶太人的過節.
耶穌舉目看見一大群人來.
就對他們說.
我們去哪裡買餅給他們吃.

$^{281}$當時他們是加利利.
很多山頭 曠野的地方.
耶穌就說.
那麼多人跟著我們.
那怎麼辦.
你記得那個故事.
就是五餅而魚的故事.
每每我們看到五餅而魚的時候.
我們就會想起.
上帝的供應.
或者上帝給人的恩典是足夠的.
因為五餅而魚就餵飽了五千人等等.
但你留意到.
特別是約翰福音.
約翰福音才會這樣記載.
因為五餅而魚這個故事.
四本福音書都有提.
但只是約翰福音.
是加了這句第四節.
他就在留意.
這個五餅而魚的故事.
神跡的故事.
是何時發生.
過世的時候.
附近的時候.
才會發生.
約翰福音很多時候.
講到上帝藉著耶穌所做的事.
那些的做事.
那些的偉大行為.
上帝要做的工作.
其實是一個的證明.
一個的路牌.
換句話說.
譬如你去餐館.
你就打開餐館的菜單.
你看到有很多相片.
他就告訴你.
這款菜是什麼名字和什麼樣子.
那個就是證明.

$^{321}$證明是指出現實.
你不會拿著菜單來吃.
你會告訴侍應.
我要這個.
到他拿來的時候.
就是那一碟菜.
貨不對辦.
你說和那個證明是不同的.
菜單就是那個證明.
在約翰福音裡面.
耶穌做的事.
一舉一動.
都是一個的證明.
五餅而魚就是一個的證明.
point to something beyond itself.
不是那個實體.
而是指向實體的東西.
五餅而魚這個神蹟.
約翰就說.
想清楚.
不只是供應.
每天吃東西這麼簡單.
耶穌在做的一件事.
是最終極想做什麼.
約翰福音第六章當然很長.
我這裡純粹給大家一個outline.
大概五段.
首先就是五餅而魚.
第二.
發生什麼事.
是故意鋪排的.
不一定是時間上這樣發生.
但記載耶穌做的事的時候.
第二個故事他們說到的時候.
就是耶穌在加利利海平靜風浪.
都不拉監.
吃完東西然後又說平靜風浪.
你記得我剛才說什麼.
猶太人想起御節的時候.
近了第六章第四節.

$^{361}$他就想起整個package.
想起食物.
天上有糧食給我們.
然後想起過紅海.
上帝把水分開.
或者上帝把水control.
把那些混沌.
邪惡的東西來control.
然後他就提到耶穌平靜風浪.
所以耶穌平靜風浪又不是單單.
我們生命裡面有風浪.
上帝就平靜了.
當然是這個意思.
但也是一個sign.
告訴我們耶穌在做什麼.
這裡想起整個御節出埃及的故事.
然後下面三段都是說話.
如果你回去看六章約翰福音.
很長.
看完會拗頭.
總之耶穌和猶太人辯論.
大家有很多對話.
說什麼呢.
就是說食物.
因為人們看到耶穌五餅如魚.
這個人經常請吃飯.
當然經常去找他.
結果他們真的去找他.
坐船去找他.
找到他了.
他就跟他們說.
你可不可以給我們糧食.
耶穌就跟他們有很多的debate.
有很多的對話.
然後有些人就發怨言.
有些人聽不明白就跑了.
他說這個人說話太爛了.
我們接受不了.
然後有很多本身跟耶穌的人.
就馬上走了.

$^{401}$因為他都不再請我們吃飯了.
就不要跟他做朋友了.
就是這個意思.
然後耶穌最後的時候就round up.
說一件事.
就跟12門徒說.
你們走不走.
門徒就說.
你有生命的道.
我們還要跟誰.
你有永生.
我們還要跟誰.
類似一個這樣的對話.
我們看幾節.
讀了這17節.
這就是大家在debate的時候.
就是那些猶太人.
他們試了一次五餅如魚.
哇 都不用做事.
然後這個人請我們吃飯 多好.
先找他.
找到他了.
然後就跟他說.
喂 你可不可以給我們這些東西吃.
你可不可以再降這些糧.
你記得整個預節的package.
就是有摩西開上面的bank.
然後降嗎哪下來.
這個會不會就是那個呢.
讓我們不用每天做事.
然後耶穌就跟他說.
他說不要為那會壞的食物而操勞.
但in fact.
在出阿俱記裡面.
說到那些嗎哪掉下來的時候.
如果那些人.
將它儲起的話.
先拿多一點.
明天我怕不夠.
然後我就雪起它.

$^{441}$然後第二天我再加熱它.
差不多這個意思.
其實是會壞的.
他就說不要為那會壞的食物而去操勞.
是你可以賺錢.
Sorry.
但不要為那東西而操勞.
為什麼而操勞呢.
存著永生的食物.
這食物就是.
耶穌說我人子賜給你們的.
你要找那些食物.
是什麼食物.
然後那些人聽不明白.
34節.
他們就對他說.
主要是請上上每日.
打謝糧給我們.
每日你都給我好嗎.
我每日都想你給我飯吃.
記得在出阿俱記的時候.
他們走過鄺氣40年.
真的每日給他們.
就算到Sabbath來之前.
都給兩份他們.
等到Sabbath的時候就不用煮飯.
40年上帝就是這樣去供應他們.
但這裡.
耶穌就有少少.
用他們所熟悉的故事.
來講一些再深一點的東西.
他說你吃完每一日的東西.
你都要再吃.
40年都是這樣吃.
事實上那些以色列人.
吃了40年.
都是死在曠野.
怎樣都是死硬.
為什麼要毀這件事.
這麼粗魯.

$^{481}$耶穌講一些更加重要的食物.
6章第35節.
他說我就算生命都要.
來到我這裡的.
終不饑餓.
信我的永不乾旱.
我給食物和水給你們.
這些是永遠的食物.
這個就很出名.
在約翰福音裡面有七句.
七次.
耶穌講什麼.
我是.
I am.
他說我是門.
我是羊的門.
等等好幾次.
這是其中一次.
I am the bread of life.
是永遠的生活.
他說我就是生命的糧.
來到我這裡的話.
你就不會饑餓.
信我的話就不會乾旱.
這個當然一定是比喻.
因為我們會吃東西.
吃食物.
我們不會信食物.
在你面前的飯都是吃的.
你不會說我信你阿凡.
我很行的.
我信了你就可以.
這個肯定就是在講比喻.
我們不會信麵包.
我們會吃麵包.
耶穌就是這個意思.
我就是生命的糧.
信我的就不會饑餓.
下一節.
又繼續debate的時候.

$^{521}$就講什麼.
他說我就是從天上降下來.
生命的糧.
上面的倉庫打開.
降下來有馬欄.
人若吃這糧.
就永遠活著.
我為世人的生命.
所賜下的糧.
就是我的肉.
這個更難.
什麼.
他說他降馬欄.
降馬欄.
其實他就說.
我就是馬欄.
他說不只是給你一碟飯.
他說我就是飯.
我就是飯.
你信我的話.
你吃我的話.
你就能夠永遠活著.
在以前農業的社會.
跟我們有一點不像.
農業的社會做什麼的.
農業的社會很清楚.
他們吃任何東西.
吃的東西都一定要犧牲了.
你才能吃到.
舉個例子.
我們今天拿個漢堡包去吃.
有什麼死了你才能吃到漢堡包.
以致你能夠不餓.
能夠活著.
牛.
除非是那些.
那些不知道是什麼.
可能麥當勞就是這樣.
蒼蠅都不吃.
牛.

$^{561}$菜.
那塊菜.
一個植物死了.
然後放在你面前.
夾著漢堡包.
你才可以吃.
麥.
大麥小麥.
做出來的麵包.
是由什麼來的.
是由田產出來的.
是那些麥子死了.
然後才煮到在你面前.
然後你吃了.
你就可以活著.
以前農業的社會.
很清楚這件事.
我們說吃羊羔.
吃任何的東西.
幾乎都是有些東西死了.
有些東西犧牲.
以致我們能夠得到生命.
剛才我們說的.
全部是還的節期.
全部都是有東西死的.
新年不要說這些.
比你更差.
但每樣東西都死了.
放得上去的.
都是已經死了.
以致你能夠有生命.
耶穌就是在說這件事.
耶穌就說.
我的死.
使得你有生命.
每一天.
你能夠有糧食.
都是因為有生命犧牲.
以致你能夠有生命.
如果說到永生的話.

$^{601}$就要先是我為你犧牲.
以致你不會受審判.
以致你能夠得救.
我們等下會搜聖餐.
聖餐.
就是從這一類的經文.
引申出來.
就是耶穌的死.
每時每刻我們拿聖餐的時候.
想起祂的犧牲.
餅和葡萄汁.
大麥小麥的死.
葡萄樹的死.
能夠祝福我們有生命.
這才是我們要操勞的.
最重要我們要知道.
我們是領受上帝的救人.
開始的時候.
我們提醒大家.
在感恩節裡面.
為什麼去操勞.
為什麼去說.
Friend God.
最重要的.
是天上國下來.
生命的.
耶穌基督自己.
我們請勁八隊上來.
唱這首歌的時候.
我們在最後.
第一節的聖經.
是跟預知有關.
跟耶穌的聖餐有關.
他說既然你們是無孝的民.
在說教會.
在說御節.
無孝的民.
舊孝隨政.
好使你們成為新團.
我們御節的羔羊.

$^{641}$基督就是被殺.
犧牲的.
我們來守這個節.
其實他在說聖餐.
不只是猶太人的節.
守這個節的時候.
我們就是.
作無孝的病.
孝代表罪.
代表一些不好的東西.
除去的時候.
領受真正的生命的良.
以致我們能夠永遠活著.
我們同心大道.
上主我們多謝你.
因為你是那一位.
帶領我們愉悅的羔羊.
你就是那個.
過世的羔羊.
天上降下的力量.
叫我們能夠得到永生.
上主我們感恩.
去讚給你.
因為你為我們預備的救恩.
你的愛子就是.
御節的羔羊.
在十家上.
為我們流血寫生.
救我們離開罪惡.
離開審判.
離開死亡和黑暗.
主帶領我們這一刻.
去領受這個.
代表你同在的.
這個餅和這個盤.
提醒我們.
你救恩的貴重.
主當我們領受的時候.
我們知道你就是天上降下來的.
生命的力量.

$^{681}$叫我們永遠不記憾.
主求你的救恩.
能夠在我們流血的.
這個餅和這個盤.
能夠在我們內心裡面.
在我們生命最深的地方.
去滿足我們.
祈禱奉主的名.
我們先唱這首歌.
《在困苦艱難》 歌手:陳昇.
在困苦艱難.
仍牽迫肚肚要安靜.
知道你是神無罪羔羊.
曾被失怨犧牲.
為救贖流血捨身.
令人越死亡.
越掛慌亂.
神恩守不可一錯.
成功是實在和.
潔淨一生無錯.
自己都不怕痛苦.
願救恩歸於.
坐在寶座上的耶穌.
願救恩歸於.
坐在寶座上的耶穌.
願救恩歸於.
坐在寶座上的耶穌.
願救恩歸於.
坐在寶座上的耶穌.
令我越死亡.
越掛慌亂.
神恩守不可一錯.
成功是實在和.
潔淨一生無錯.
自己都不怕痛苦.
我們有三張聖誕套餐.
如果方便的話.
你都準備好的話.
拿一條套餐.
然後回到我們的座位.

$^{721}$我們一起應酬.
如果你行動不太方便的話.
可以叫身邊的弟兄姊妹幫你拿.
所有我們信耶穌的人.
心靈的靈.
一個代表.
一個象徵.
讓我們一起應酬主的病.
也要和領主的病.
我們同心得道.
主愛多謝你給我們生命.
永遠的生命.
永遠天上生命的靈.
感恩你是那位.
在我們生命裡面.
帶領我們如意的神.
主愛為了每一位弟兄姊妹.
可能我們要經歷過的禍困.
要面對很多的危疾.
甚至要面對瘟疫.
主愛你知道你的愛.
你自己的拯救.
超越這一切的.
主愛求你讓我們信靠你.
叫我們能夠去領受.
你天上的律.
領受我們耶穌基督.
我們的主.
你知道你的愛.
永遠和我們同在.
我們無所懼怕.
主叫我們養手去養我們.
叫我們能夠在每一天.
靠著你的恩典.
去明白你是那位同行的神.
我們祈禱防住病毒.
願上帝祝福各位.
希望下星期能和大家一起.
事奉敬拜.
\newpage



\section{}
\label{sec:PVQ3X0QwXNk}
\textbf{Faith and time: Passover - Rev. Mark Chan}
\newline
\newline
連結: \href{https://youtube.com/watch?v=PVQ3X0QwXNk}{\texttt{https://youtube.com/watch?v=PVQ3X0QwXNk}} ~~~~ 語音日期: 2022-11-13
\newline
\newline
\hyperref[sec:wlVkY6Ioxpo]{\small{< < < PREV SERMON < < <}}
~
\hyperref[sec:index_chronic]{\small{[返順時目]}}
~
\hyperref[sec:index_scriptual]{\small{[返順卷目]}}
~
\hyperref[sec:E_n46N8Pd_8]{\small{> > > NEXT SERMON > > >}}
\newline
\newline
$^{1}$Thank you very much for coming,.
and we're gonna worship God together,.
and we're gonna turn to our series,.
which is the festivals,.
the important days in the Old Testament..
So we talk about the Sukkot, Tabernacle..
We talk about Shabbat..
Then today, we talk about another important festival.
that talks about in the Bible..
So let's see..
Let me start with these questions, okay?.
Now, we often know there's a relation or association.
with some festival or holidays with some food, right?.
When you think of some holidays or some festivals,.
you always think of some sort of food, right?.
Food items related to the festival..
Thanksgiving, what do you think of?.
- Turkey..
- Yeah, turkey, right?.
- Hot pot..
- Hot pot, yes, hot pot, yeah..
So when Thanksgiving comes,.
and you're gonna prepare some food.
that is related to the festival,.
Christmas, what do you eat?.
- Hot pot..
- Hot pot, yes, hot pot, yeah..
How about the New Year's?.
- Hot pot..
(congregation laughing).
- How about the Chinese New Year?.
- Dumplings..
- Dumplings, yes, dumplings..
Okay, finally, something other than hot pot..
How about Valentine's Day?.
I asked my wife, what do we eat on Valentine's Day?.
- Steak..
- Steak, oh, okay..
I thought you were gonna say hot pot..
(laughing).

$^{41}$All right, so the food items related to the festival.
or the holidays of the calendar,.
and we think about what the Bible says about Passover,.
what do you think of?.
Food item..
What food items are related to Passover?.
Yeah, unleavened bread, yeah..
And lamb, yes, the lamb, right?.
In fact, one of those big items.
that they will prepare for Passover,.
at least for the Old Testament time,.
the ancient Israelites,.
they will do prepare lamb, Passover lamb, right?.
So one family will take one lamb,.
and they will consume it,.
and they will offer the rest of it to God..
And this is right at the first Passover,.
right at the eve when they leave Egypt,.
when God is going to judge the nation,.
judge the nation of Egypt,.
to destroy Pharaoh and his empire..
So judgment of God will come..
And then right before it comes,.
the 10th plague, plague number 10,.
which is judging the firstborn sons.
of Egypt, right before the divine judgment.
comes to the land of Egypt,.
God said, "You're gonna prepare a meal,.
"because you're gonna go, you're gonna go in a hurry..
"You're gonna go in a hurry in the time of judgment,.
"and I'm gonna save you out of death and of judgment.".
So it says in Exodus chapter 12,.
"And now the whole community of Israel,.
"when they were still living in the land of bondage,.
"in the land of slavery in Egypt,.
"on the 10th day of this month,.
"each family, you would choose a lamb,.
"young goat, for a sacrifice,.
"one animal for each household.".
And this is the instruction..

$^{81}$Be fully prepared, fully dressed,.
wear your sandals, carry your walking stick in your hand..
Now, that's interesting,.
because the ancient Israelites,.
they go into someone's home,.
they will take off their shoes, just like us, right?.
Most of the Asians, we come to someone's house,.
or go back to your home,.
you take off your sandals, you take off your shoes..
Here, when they eat inside the home,.
they need to put on their shoes, right?.
They need to get prepared,.
and dress themselves, fully dressed,.
so that they are ready to go..
They are eating in a hurry, with urgency,.
because that night is gonna be the Passover of the Lord..
Now, what does it mean?.
They would kill the lamb,.
and then they would take some blood,.
and smear on the doorpost, or the lintel of the door..
And then, when the angel of judgment.
passes through the entire land of Egypt,.
it will see the doorpost,.
it will smear with blood, and it will pass over,.
it will skip that house..
And that is the story of salvation..
In times of judgment, the angel of the Lord,.
the judgment of God, come to the land,.
and they will be safe, the Israelites will be safe..
So it talks about blood, talks about judgment,.
talks about protection, talks about redemption..
God's gonna save his people.
out of the land of darkness, out of the land of slavery..
So they get prepared..
So this is the first Passover,.
and then year after year,.
they're gonna celebrate this festival,.
which is one of the most important festivals.
of all other holidays,.
because it marks the most important acts of God,.

$^{121}$that they are being saved from slavery, from bondage,.
and they are being redeemed..
So year after year, Jews,.
wherever they live, in diaspora,.
in different parts of the world,.
they will go back to Jerusalem,.
go back to the temple, especially in Jesus' time..
It was one of the highest holiday of the year,.
and right about the time of spring,.
March and April,.
and they will come to Jerusalem for a week to celebrate..
And in that week, the population of Jerusalem.
will skyrocket two times or three times of their population..
And they will come around the city,.
and then they will go into the city,.
and they will arrange a big meal,.
which is the Passover meal, sacrificial lamb..
They're gonna remember what God had done.
in their history, in their life..
Now, we talk about the Passover,.
thinking of food item being the sacrificial lamb..
But the Jews, the Israelites,.
when they think about the Passover,.
they think about what God had done in their lives..
It's not just the lamb..
They think about the whole package,.
the whole package of God's redemption..
So let's do some review..
When you think of the book of Exodus,.
when the Israelites came out of Egypt,.
what are the events follow after Passover?.
What are the events?.
Think about it..
Huh?.
- Crossing the Red Sea..
- Yeah, crossing the Red Sea..
Crossing the Red Sea, and then after what?.
What happened?.
Crossing the Red Sea, where did they go?.
Did they go to Chicago?.

$^{161}$- The wilderness..
- The wilderness, yes..
We have a Buddhist student here..
We got all the answers..
So after crossing the Red Sea, God split the Red Sea..
There's a water miracle,.
and then they go into the wilderness for 40 years, right?.
And then they receive the Ten Commandments.
and all that stuff..
And then in the wilderness,.
who is going to provide for them?.
What do they eat?.
There's nothing in the wilderness..
What do they eat?.
- Manna..
- They eat, yeah, they eat manna, right?.
Manna..
So when they think about the Passover,.
they think about the story of Exodus,.
they don't just think about just one lamb and that's it..
They think about a whole series of events,.
the entire process that God will bring about.
the salvation to his people from the start to the end..
So from the start, Moses is gonna come, right?.
And lead the nation as a shepherd,.
and lead the entire people out of Egypt..
Sacrificial lamb, we talk about it..
And then for one week,.
they will celebrate with unleavened bread, right?.
Because you're in a hurry, right?.
You don't have time to wait for the yeast.
to take the effect, right?.
So they're gonna eat unleavened bread..
And the,.
you know, if you have baked bread before, right?.
You know how many hours or how much time you have to wait.
for the yeast to take effect, right?.
The dough, you know?.
It's very time consuming..
So when they are leaving Egypt in a hurry,.

$^{201}$God is saying to them,.
"You're gonna eat unleavened bread..
"Don't wait, you gotta go..
"You gotta go by my command.".
And then they go to the shore of Red Sea,.
and then God split the water, conquered the water,.
and they go through it on dry ground,.
going to the wilderness,.
and the people grumble,.
complain about,.
"Hey, is God gonna provide for us what we need?".
Right?.
And then through Moses,.
God sent food from heaven, manna, right?.
What's manna?.
What is it?.
Now, according to the description in the Bible,.
in the book of Exodus,.
it's something like a round stuff, you know?.
It's like a biscuit or something, right?.
Well, the point is we don't know..
We just don't know what it is, right?.
In fact, the biblical story talk about the manna..
The Israelite did not know what it is..
So when the first time, the very first time,.
very first day, God sent the manna to the field,.
they go collect the manna,.
what did they say?.
They say, "Ma set, ma set.".
Ma set is in Hebrew..
It means what it is..
Ma is what..
Set is this..
And what the heck is this?.
So that's why you got the name manna, right?.
Ma set, manna, manna, all right?.
So the most important point is we just don't know what it is..
But when they grumble, when they complain about it,.
God gives them food, rain from heaven,.
raining down from heaven, manna from heaven..

$^{241}$And then as the story goes from the Old Testament,.
you know, go back all the way to Exodus.
and all the way down to the kingdoms.
and all the way down to the time of Jesus..
And then the Jews, the ancient Israelites,.
will think about and remember Passover.
and adding a lot of stories to it..
And one of the interesting stories they add.
around the time of Jesus is that they,.
because according to the story of the Old Testament,.
when they went through the wilderness for 40 years,.
God gave them manna day after day, right?.
And then right after they crossed the River Jordan,.
they entered into the promised land..
In the book of Joshua, it says, there's no more manna..
It stopped raining from heaven, no more manna,.
because they're gonna eat the produce.
from the land of Canaan..
So some of the Jews around the time of Jesus,.
they would say, you know, because the manna stopped,.
because there was a bank, there was a treasury in heaven,.
raining down manna to all of them..
And when they entered the land of Canaan,.
the bank is closed, the treasury, the storage is closed..
And who closed it?.
Who has the power to close it?.
Moses..
Moses has the key..
So he opened the treasury for 40 years,.
raining down manna to the people of Israel in the wilderness.
for 40 years..
After they entered Canaan, Moses closed it..
So they eat the produce of the land..
So when the Jews thinking about the story,.
they have the legend, they have a legendary telling.
that someday in the future, when God's salvation.
is gonna be fulfilled ultimately,.
then Moses will come again and open that bank..
It's like an ATM, open it again,.
and manna will come from heaven again..

$^{281}$Now you think about the whole story,.
you can tie different things together,.
and you can think of some of the description.
that the gospel talks about Jesus..
Talk about Passover, talk about what Jesus.
will ultimately do to fulfill the salvation..
Now let's take a look..
The gospel of John, John chapter six..
Now remember what I said about what we thought about.
the entire Passover, the entire Exodus story,.
the entire package, the manna and all that stuff..
Read the story, you will read it in new light..
The gospel of John chapter six..
It was nearly the time for the Jewish Passover celebration..
And Jesus saw a big crowd of people.
coming to look for him..
And then he said to Philip, one of his 12 disciples,.
"Where can we buy food for them?.
"How do we feed these people?".
There's so many of them..
So you know how the story goes,.
it's the five loaves and two fish..
Five loaves and two fish..
It's one of the most well-known Jesus miracles of all..
Because all four gospel, Matthew, Mark, Luke and John,.
talk about this story, five loaves and two fish..
We think about this little story,.
five loaves and two fish, it's a very nice story,.
it's a very nice children's story, kid's story..
I'm sure they all heard of it in our Sunday school..
And we think about the provision of God,.
God provides and Jesus provides..
But here, there's something very interesting..
Sixth four, John at the statement,.
Matthew, Mark and Luke did not record the story.
the same way..
John, at one special statement,.
it was the time of Passover..
It was the time of Passover that the five loaves.
and two fish story happened..

$^{321}$Now in the gospel of John, all the great deeds,.
all the great acts of Jesus are signs..
John called them signs..
Now what is a sign?.
What is a sign?.
When you're driving on highway, for example,.
you drive on the highway and on the right side,.
you see a sign, Chicago, 21 miles..
The sign is not Chicago..
So you don't hit the sign..
The sign is pointing you to the reality in the future..
The sign is pointing you to something else.
because signs is always pointing beyond itself..
So all the great acts of Jesus are actually signs.
and telling you there was some deeper meaning to it..
There was some reality the sign is pointing to..
You go to a restaurant, you pick up a menu,.
you see picture of fish..
You don't eat the menu..
You see the menu is the sign..
The menu is the sign for the reality..
You eat the fish that is cooked by the restaurant..
The reality comes to your table and then you partake it..
You take it, you take it in..
Menu is set aside..
Now this is what John is doing..
John is recording this story in a special way.
to tell you what it means..
So if you look at the entire chapter six of the gospel of.
John, it's a very long chapter..
We're not going to study it verse by verse..
But I'll give you the outline here..
This is John six..
There are a few things that you need to understand..
There were a couple of stories..
And then there were a bunch of things that recorded.
about the dialogues or the debate between the Jews.
and Jesus and then finally his disciples..
Now what kind of stories are being recorded here?.
Of course the five loaves and two fish, right?.

$^{361}$Jesus feeds 5,000 people..
Nice story..
But following this story,.
John recorded that Jesus calms the sea of Galilee..
There was a storm story..
Now why put it in such a way, you know, back,.
front and back, you know, together?.
Remember Exodus..
Remember the Passover, the whole package..
What kind of miracles are being recorded?.
The Red Sea miracle..
God controlled the sea..
God conquered the sea so that his people are safe..
What a miracle..
And then feeding the 5,000,.
feeding the people in wilderness..
The Exodus story is all about, you know,.
raining down manna..
God is going to provide for them in the wilderness..
Even they complain about it, even they grumble about it..
For 40 years, God is going to give them breath from heaven..
So John's putting different things together..
It's like a TV show producer, you know,.
orchestrate everything, different stories,.
pack it together so that you know who Jesus really is.
and what Jesus has done for you..
The food miracle..
And then the water miracle..
So Jesus is doing something that the Jews,.
when they heard about this story,.
they understand what it means..
There are different things here, feeding 5,000,.
the water miracle, people grumbling..
Now, you can see part three, four, and five.
are all about talking..
There's so much talking..
You know, first time I confess, I read John 6..
I just don't know what the heck are you talking about?.
They ask Jesus some questions,.
they debate about the meaning,.

$^{401}$and then Jesus taught them different things..
You know, but the entire John 6.
is all about in the wilderness.
when people are grumbling, complaining..
In fact, some of the Jews heard about what Jesus said,.
they left..
Originally, they were disciples..
They follow Jesus, but when they heard about.
some of the teachings of Jesus in chapter six of John,.
they didn't understand and they were confused..
That's why at the end of the story, they left..
They complained and then they left..
Then toward the end, Jesus round up the whole thing.
and asked the 12 disciples, "Are you gonna leave too?".
Remember one of the disciples, Simon Peter?.
He said, "You have the eternal word of God..
"Who shall we follow?.
"Who else shall we follow?.
"We're gonna follow you.".
Now, let's take a look at some of the verses.
and see what the story means about Jesus..
627..
Now, when they're talking about the food,.
remember the first story, Jesus fed 5,000..
Some of the people are following,.
big crowd are following him..
Try to use your imagination..
If you're one of those who experienced.
the feeding of 5,000,.
what would be your immediate response?.
If you go to a restaurant with your friend.
and your friend pay for your dinner,.
he said, "It's gonna be free..
"I'm gonna pay for you.".
What would be your immediate response?.
Wow, okay, wow..
The next day, what would be your immediate response?.
You're gonna call your friend,.
"Hey, let's go out for dinner again.".
Yeah, right?.

$^{441}$'Cause every meal, you're gonna pay for it, right?.
So this is how the people react..
So they were looking for Jesus.
after the 5,000 feeding miracle..
They were looking for Jesus,.
just disappear and on the other side of the sea..
So they ride on the boat and then they come to Jesus.
and say, "Hey, can you give me more?.
"Can you give me more?.
"'Cause you just gave me one day of the daily bread..
"Okay, that's good, all right?.
"How about tomorrow's dinner?.
"Or the dinner's after tomorrow, right?".
So Jesus was telling them,.
"Do not be so concerned.
"about perishable things like food..
"Spend your energy seeking the eternal life.
"the Son of Man can give you.".
Perishable food..
Don't be so concerned about things that perish..
Now think about the Exodus story..
When man came, the first week,.
some of the ancient Israelites, as the story goes,.
they collect more than they need..
And they put it back into their tent in their home..
And then the next day, it mowed..
It mowed..
They didn't trust God..
They didn't trust God enough for daily bread..
So they collect maybe two days, three days, right?.
Just like going out with your generous friend,.
you're gonna order more..
And then as we talk, you're gonna put it in your to-go box..
But it mowed and it rot..
Perishable things..
Jesus said, "Don't be so concerned.".
Yes, you need to plan, you need to be responsible,.
you need to work..
However, don't be so concerned about perishable things..
Aren't we all do that?.

$^{481}$Aren't we all like that?.
We're too concerned about things that will perish, right?.
You're gonna store up your money, store up your food,.
because we worry, right?.
Many of us, Chinese heritage,.
food is so important to us, right?.
Food is too important for us..
Jesus says something is more important..
Seek the life, seek the food.
that the Son of Man can give to us..
34, the crowd said to Jesus,.
"Sir, give us that bread.".
How about every day, all right?.
Pay for our dinner every day, all right?.
So we don't have to work for it..
Thinking about Moses, opening the treasury..
Hey, can you be that Moses?.
Can you just be that Moses,.
and every day we're gonna have bread on the field,.
we don't have to work for it..
And Jesus said, "Don't be so concerned about it,.
"because I'm gonna give you food that will last,.
"that will not perish.".
35, Jesus replied, "I am the bread of life..
"Whoever comes to me will never be hungry again..
"Whoever believes me will never be thirsty.".
Jesus is saying, "I'm not just giving you food.".
He said, "I am the food.".
That's kind of astonishing, right?.
When you have dinner with your friend,.
your generous friend,.
your generous friend, all of a sudden, if he say to you,.
"I am not the dish in front of you, I am the dish.".
You say, "This guy's nuts.".
Give you a plate of fish, I give you fish,.
however, I am the fish, man..
What are you gonna say?.
You say, "Jesus, what are you talking about?".
In fact, that's what the people react..
They say, "That's crazy.".

$^{521}$Now, of course, the language is metaphor, right?.
When Jesus said, "I am the bread of life,.
"you believe in me, you will be safe,.
"you have eternal life.".
It's metaphor, because we eat bread and we drink water..
We don't believe in bread, we don't believe in water..
It's metaphorical..
The bread is Jesus himself, right?.
As the story goes, 51, he said, "I am the bread of life.".
Jesus said, "I am the bread of life.
"who came down from heaven, from the treasury of heaven..
"I'm not just the one, I'm not like Moses.
"with the key open the depository in heaven,.
"treasure in heaven to rain down food for you..
"I am the food, I'm the living bread,.
"came down from heaven, anyone who eats this bread.
"will live forever.".
Now, what is this bread?.
My flesh..
We cannot take a bite of your arm or something..
It's a metaphor..
Now, people in the ancient days,.
especially in the agricultural world,.
where you have to farm and you have to harvest,.
you get your food..
People understand almost everything they eat.
is something that has died, right?.
You think about it, even today,.
when you go to a fast food restaurant,.
you get a hamburger,.
you eat something that has died, right?.
What has died in order to serve you the hamburger?.
Wheat and barley made into flour,.
and flour baked into bread, right?.
The lettuce, tomato, pickles, the plants,.
somehow they die being harvested..
Serve within the hamburger, the patty..
The animal die so that you can protect the flesh,.
and you get the hamburger, right?.
Unless you're eating a hamburger totally engineered,.

$^{561}$you know, genetically, you know, like a McDonald, right?.
Even the flies gonna eat it..
Everything that we eat entails some death..
Something has died in order to serve us food.
to keep us going, and people understood it..
If we continue to live, other living things need to die..
If those things don't die, what will happen?.
You die..
You gotta have something else die for you.
so that you can live, right?.
Jesus is using this metaphor..
He said, "I am the Passover lamb..
"I am the manna..
"I am the water..
"I die so that you can live.".
Our daily bread needs, you know, daily things to die..
Eternal bread give us eternal life..
It's not just daily life, eternal life..
Jesus is the Passover..
He died so that we can have life,.
and that's why almost all the songs that we sing in worship.
have to do with the death of Jesus,.
His sacrifice, His sacrificial act, the Passover..
He sacrificed Himself so that we can live..
He died on the cross, and He shed His blood.
so that we can live, we can have eternal life..
Well, we're going to observe the communion today..
I want you to understand all the things.
that we talk about about the Passover.
are signs pointing to the reality..
The Son of God died on the cross because He loves you..
He loves me so that we can be redeemed..
We can live forever..
As we take the cup and the bread, the communion,.
there are also signs pointing to the reality..
If we partake it, if we take the bread,.
that means we're receiving Jesus, we believe in Jesus..
We're going to have eternal life..
That's going to change the whole world for you..
Whatever we work for, wherever we are,.

$^{601}$we're going to change because our identity change..
We are the sons and daughters of God..
We belong to the kingdom..
We belong to that big story,.
you know, from creation to redemption.
to all the way to new heaven and new earth..
We are just a little story..
Be part of that big story..
Partake Jesus, and we're going to enjoy.
the eternal relationship with our Father..
Let us all pray together and prepare ourselves.
as we take the communion..
All right..
Father, we thank you for the death of your Son,.
for His sacrifice, for what He has done on the cross..
Yes, because of His sacrifice,.
because of the Passover lamb,.
that we have the eternal life..
And let us hear your word again..
I am the bread of life..
We will not go hungry eternally because we partake in you,.
because we believe in you..
And thank you for your sacrifice..
Thank you for your grace that we can become your children..
As we take the bread and take the cup,.
remind us our relationship with you..
Remind us your so great of a salvation,.
the eternal life that we have..
And help us orient ourselves toward you.
so that we are not too concerned about our daily bread,.
but we are more concerned about your kingdom,.
our relationship with you..
In Christ's name, pray..
All right, if you're ready, you can go back, I guess, right?.
\newpage



\section{}
\label{sec:E_n46N8Pd_8}
\textbf{The Greatest Commandment (Matthew 22: 34-40) - Brock Lockenour}
\newline
\newline
連結: \href{https://youtube.com/watch?v=E_n46N8Pd_8}{\texttt{https://youtube.com/watch?v=E\_n46N8Pd\_8}} ~~~~ 語音日期: 2022-11-13
\newline
\newline
\hyperref[sec:PVQ3X0QwXNk]{\small{< < < PREV SERMON < < <}}
~
\hyperref[sec:index_chronic]{\small{[返順時目]}}
~
\hyperref[sec:index_scriptual]{\small{[返順卷目]}}
~
\hyperref[sec:042xzVMibbM]{\small{> > > NEXT SERMON > > >}}
\newline
\newline
$^{1}$Thank you..
Okay..
So I have a friend named Grant,.
and I'm guessing most of us have a friend like this.
in your friend group..
Grant is that one person in your friend group.
who is heads and shoulders, smarter and more successful.
than everybody else in your group..
If you know who I'm talking about for yourself,.
think about them..
Grant is a genius..
Grant is very influential in his career..
And this guy,.
we always joke that if you wanna feel bad about yourself,.
you should Google him.
because the number of podcasts he's on.
and the number of interviews he's been given.
and the number of lists in magazines.
that he's made it onto is a little sobering.
when you're like, "Oh, I'm about his age, that's great.".
And the reason Grant is successful and well-known.
is because Grant is an entrepreneur..
So he opens businesses, he runs small businesses,.
and he's been doing this since he was 16,.
since he was maybe 14..
I am not kidding..
When he was in his teenage years,.
this guy opened a thrift store in Hyde Park.
that is still running to this day..
He's also written books..
He's made it into magazines..
And right now he runs a company.
that has more than 50 employees..
It's a little insane..
Now, one thing Grant always talks about.
when you talk to him about his businesses.
is he'll talk a lot about projects.
that they will dump or not do.
because of the fact.
that it doesn't fit their mission statement..

$^{41}$And this is something he says he's been doing.
ever since he started opening that thrift store.
when he was in his teenage years..
So Grant would tell you that when you open a business,.
your mission statement is the very foundation.
of what you do..
And you don't just have a mission statement.
because it's fun to put there..
You don't have a mission statement.
just because you need something to pitch to investors..
You have a mission statement.
because when an opportunity pops up,.
you can compare it to what you first set out to do.
and decide whether or not it's something.
you're going to do based on that..
So if you're opening a burger restaurant.
and your mission statement.
is to provide the best dining experience.
with the most unique burgers in Chicago,.
you're probably not going to jump.
on the chicken sandwich trend when it comes around..
You're probably not going to have a Taco Tuesday special.
because it doesn't match your mission statement..
You're supposed to stay on point..
Everything you do is born out of that mission statement..
It's a foundational piece..
It's like the cement upon which you build a tower..
And he tells me that every single one of his businesses,.
past and present, has always and will always.
have a very strict and defined mission statement.
because otherwise you're going to end up wasting your time..
So I want us to think today.
about what is your mission statement?.
Now, I know people are not like businesses..
We don't sit down and decide that we're going to.
build a person and then decide what they're about.
and what their core values are..
Most of us probably don't have an index card.
with our mission statement written down on it..
I'm guessing most of us don't have an elevator pitch.

$^{81}$for who we are as a person..
That we can just recite at any moment, right?.
But the thing is that each and every single one of us.
has a mission statement,.
whether or not we understand that we do..
Even if we don't think we do,.
there is still something that is that foundation.
that drives every action you do,.
every day that you spend your time,.
it drives the decisions you make..
So I want you to pause for a minute.
and think about what that would be for you..
What is the thing that drives you?.
What gets you out of bed every single morning?.
If I were to distill every action and motivation.
and every task you do and the way you spend your time.
and the people you hang out with,.
if I were to distill that.
into some kind of a mission statement or a purpose,.
what would that be for you?.
Think about that for a minute..
Maybe your mission statement is,.
I want to provide for my family, right?.
Maybe it's, I need to make enough money.
so that I can buy a house.
and have enough money for groceries.
and send my kids to college..
Maybe that's your mission statement..
Maybe you're in college right now.
and your mission statement is to either survive.
or get out with good grades.
and be in a position to get a good paying job..
Maybe you're in high school.
and you want to please your parents with good grades..
Whatever it is,.
each of us has a mission statement right now..
And if you can't define it,.
that's probably a good indicator.
that you've just kind of been going with the flow.
and maybe the thing you're building your life on.

$^{121}$is being developed without you knowing it..
That's not a good policy..
Today, we're actually going to discuss a time in the Bible.
when Jesus was asked.
about what the greatest moral commandment is..
And Jesus gives an answer.
that I think fits into this mission statement box.
really perfectly..
And his answer tells a lot about why he's there..
And it also shows us a lot about.
why we make the decisions we make..
And whether you're in this room as a believer,.
whether you're in this room as somebody.
who's grown up in the church.
and you're not sure if you're a believer,.
maybe you're here or you're watching online.
and you're not a Christian, but you're intrigued..
I want you to be thinking.
about what that mission statement is for you..
And as we're talking about what Jesus endorses,.
think about how that compares and contrasts..
We'll do that as we go along..
Let's pray that God would bless our time today..
God, we come to you today.
knowing that there is a lot.
of really difficult things in the world..
We know there's a lot wrong with existence right now,.
from natural disasters to oppression to injustice..
God, we know and are fully aware of the fact.
that we don't live in a world that is easy.
or oftentimes fair..
But God, it is all the more because of that.
that we thank you for your goodness..
And we thank you that you love us,.
even though we don't deserve it,.
even though it's not something that we earned,.
it's not something that we should have..
And honestly, the fact that a God like you cares about us,.
even in the slightest bit, should astound us.
and should gobsmack us into silence for minutes at a time..

$^{161}$We love you, God..
We ask that as we explore this topic,.
you would be moving in our hearts,.
whether we already believe in you.
and call you our God or not,.
would you be working in us.
so that we might experience your love.
and so that we might be moved in a certain way.
so that we know you more and are able to have lives.
that please you and look like you even more..
We pray this in your name, God..
Amen..
Now, the passage we're gonna be in today is Matthew 22..
So you can open up your Bibles.
or turn on your Bibles if that's your speed.
and turn to that passage..
I'll be in the, oh gosh,.
I can't remember if I'm in the ESV or the NIV today..
You choose..
But before we read the passage,.
I wanna give us a little bit of background.
into where we're dropping in..
So back in Jesus' day, right,.
the temple was the epicenter of religious life..
And since everybody in that day and age.
kinda had to be religious,.
the temple was kind of the center.
of everybody's life in general..
So I want you to think about a sprawling religious building.
and I want you to think about people.
who are there all the time, okay?.
You would pass by this place on your way into the market..
When you're going about your dealings and visiting friends,.
you would pass by this place..
You would see people going in to make sacrifices,.
going in to talk to the Pharisees..
You would be here often, right?.
And so at that day and age,.
there were people whose entire job it was.
to teach the law and the scriptures and the prophets.

$^{201}$at this building, at this temple..
And these Pharisees or these lawyers or Sadducees,.
whatever you wanna call them,.
they would be out in the temple courts preaching,.
but most of the time,.
they're doing it because they want a following..
They want popularity..
They want political power..
At that day and age,.
political power and religious power.
were basically the same thing..
So these guys are out there.
and they're preaching and telling people.
how they should live their life,.
what the scriptures mean,.
and they're usually doing it.
because they want to have a following..
They want people to listen to them..
They want people to know their name and respect them..
They want to improve their rank, right?.
So these guys are essentially people.
who just want the clout..
If they were alive today,.
they would probably be Twitch streaming..
That's a bit of a stretch, but trust me..
They do this kind of stuff..
They're in it for their name..
They're in it for the political power..
And another way that they decide.
and they realize they can get that power.
is by debating people..
So let's say you're Bob the Pharisee.
and you out there.
and you're preaching to a bunch of people.
and people are following you..
And then Dave over here.
sees that you have such a following..
Well, if Dave can call you out on something,.
if Dave can debate you.
and make you look like an idiot,.

$^{241}$Dave is going to get more popular than Bob..
Dave is going to get those fans.
because he's shown everybody.
that Bob doesn't know what he's talking about..
So this kind of thing was something you would hear.
and see all the time..
You'd have crowds of people coming to these people.
and listening to them..
Now, Jesus, at this point in his ministry,.
Jesus does have clout..
Jesus is very popular..
People like this guy..
People are murmuring about maybe he's the Messiah..
Maybe this guy is the one.
who is fabled to come and save our people..
They're listening to him..
And you know what that does?.
That paints a target on his back..
That makes people want to seek him out and debate him,.
make him look dumb..
If they can just trap him in some kind of a mistake.
or make him say something that isn't actually biblical,.
or maybe it's unpopular culturally,.
they can take that popularity and power for themselves.
and invalidate him..
So this is the situation.
and the sort of arena that Jesus enters.
when he's in the temple courts and he's teaching..
So we're gonna start in verse 34.
and we're gonna go to 40..
So hearing that Jesus had silenced the Sadducees,.
that's basically a different political party.
or different religious party,.
the Pharisees got together.
and one of them, an expert in the law,.
the smartest among them,.
decided that he was going to test Jesus..
And he tested him with this question..
(clears throat).
Teacher, what is the greatest commandment of the law?.

$^{281}$Now, maybe this guy is thinking.
that he can get Jesus to say something wrong..
Maybe he's thinking he can get Jesus.
to say something unpopular..
Whatever it is, this guy does not have positive intent..
This guy ain't being sincere..
He's trying to trip Jesus up..
Here's how Jesus answers..
And most of us probably know this passage pretty well..
Jesus replied, "Love the Lord your God.
"with all your heart, with all your soul,.
"and with all your mind.".
This is the first and the greatest commandment..
And the second is like it..
"Love your neighbor as yourself.".
All the law and the prophets.
hang on these two commandments..
Jesus is actually quoting Deuteronomy here..
So, excuse me, he is pulling his answer.
from scriptures that they all would have known..
So he's saying, "Love the Lord your God.
"with all your heart, soul, and mind,.
"and then love your neighbor as yourself.".
That's the greatest commandment..
So Brock, you're talking about mission statement..
Why do you turn to this passage?.
I understand that it's not as if Jesus says,.
"Hey guys, this is my mission statement.".
This is not what's written on Jesus's.
about page on his website, okay?.
I recognize that he doesn't actually come out.
and say in his many words,.
"This is my mission statement," or "This is my purpose.".
But the reason I like to look at it through this lens.
is because it does seem to be like something.
that is very important and something Jesus sees.
as really meaningful..
What do I mean by that?.
This story, most of you probably know Jesus answering,.
"Love the Lord your God with all your heart,.

$^{321}$"soul, mind, and strength," right?.
And then love your neighbor as yourself..
It's because this story seems to be.
in three of the gospels, right?.
It seems to be in three different accounts.
of Jesus's ministry..
We have it in Matthew, right?.
We have it in Mark and we have it in Luke..
We're looking at Matthew today,.
but it's also in these other accounts of Jesus..
So when you see things that are in multiple gospel accounts,.
you should probably take notice..
So here are two options when we see multiple accounts.
of this story, right?.
We say either all three of these writers.
found this as important.
and wanted to put it into their book..
Or there's a second option.
that maybe this happened multiple times..
Maybe there were three instances.
where Jesus gave this answer..
So let's look at each of the actual scriptures..
I won't go into detail..
You can look at them yourself in your own time..
There's a Wikipedia page that lists all of them in one..
You don't have to go to several different webpages..
Matthew, the one we just looked at,.
has a lawyer testing Jesus..
It explicitly tells us that the Pharisees.
are conniving together and they're plotting something..
In Mark, it's a little different..
In Mark, there is a sincere scribe..
It says that there's a scribe,.
so somebody who writes down.
and his whole job is to copy the scriptures..
It says that in Mark, there's this scribe guy.
who's kind of sitting on the sidelines.
and he's like, "Huh, this Jesus guy,.
he's answering really well..
He's giving good answers.".

$^{361}$So then he asks Jesus,.
"What's the greatest commandment in the law?".
Not to test him,.
but because that's what he's working on 24/7..
He's like, "What is the most important part.
of this whole thing that I'm transcribing all day?".
And in that account, Jesus says to the scribe,.
"Hey, you're close to the kingdom of God, man..
Keep on, you're doing great.".
And then in Luke, it's similar to in Matthew.
where people are trying to test him,.
but Jesus doesn't actually answer himself..
He says, "Well, what do you say it is?".
He turns a question into a question..
He says, "You answer this for me.".
So what is it?.
You could believe that all three of these authors.
remembered this instance differently,.
that it was one event.
and three different writers remembered it differently..
That's a possibility..
But what I believe, what I think seems to be the case.
since there are so many different,.
interesting, unique details,.
is that I think Jesus used this answer often..
It seems that Jesus had this answer prepared.
for when people asked him,.
"What is the greatest moral commandment?.
What is the most important of these things?".
You kind of get the sense that Jesus had this answer ready..
So it seems important to him..
And I think it fits the idea of a mission statement.
really, really well..
And we're gonna have to go look through.
what he's actually talking about here now..
First of all, Jesus says,.
"You are to love God with all of your heart.".
That's your emotion..
That's the things that you feel..
You are to love God with all of your soul,.

$^{401}$all of who you are..
Think about the way that you define yourself..
I am Brock, I'm a graphic designer..
That's my identity, right?.
Jesus is saying your soul, all that you are as a person,.
all of your identity and all of who you are,.
God is worthy of that..
All of your mind, your intelligence..
How many of you have ever found a new hobby.
and you've gotten so into it.
that you learn the minutia and the details.
and you're on the subreddit looking at it every day.
and you are learning so much about this hobby or this thing..
And then eventually you just kind of lose interest.
because you've beaten it, you've conquered it,.
you've figured it out, right?.
Jesus is saying that God is worthy.
of all of your intelligence..
There is no point where you can exhaust that.
in your exploration of him..
Your strength, this isn't in the Matthew version.
we just read, this is in some of those other versions,.
so we're gonna include it as well..
This is your passion..
This is the oomph with which you live your life..
When you're motivated to do something.
and it's outside of the ordinary,.
it's not what you just do with your everyday life..
When you put forth effort towards something,.
that's your strength..
Jesus is saying that your heart, your soul,.
your mind, your strength, all of it.
is worth being thrown at God..
God is worthy of all of it..
Now I want you to think about your mission statement.
from before..
Even if you don't have it completely filled out,.
maybe you're just like, ah, I kind of seem security driven..
Maybe I seem a little bit driven by my own fame.
or popularity, or maybe I seem money driven.

$^{441}$or career or status driven, whatever it is..
Ask yourself, is the mission statement.
that you currently live your life by.
worthy of all of that?.
Most of us recognize when a person.
puts all of themselves towards a pursuit..
Sometimes we applaud it, right?.
You get basketball stars whose only goal is to win..
And when we have them on our fantasy draft, we like that..
But in real life, we don't really believe.
that that's the healthiest thing, right?.
We look at that and go, that's a little whack..
To put all of yourself into one thing,.
whether it is success, whether it's status,.
is gonna leave you in a bad spot..
Jesus is saying and affirming that all of you.
is worth throwing towards God..
That if you're looking for something.
that is meaningful enough for everything in your life,.
it's God..
I want you to think about if your current mission statement.
lives up to that..
So love God with everything..
Here's the second part..
Jesus says, love your neighbor as yourself..
Now this is interesting, right?.
He's asked what is the greatest commandment?.
And he gives two..
He doesn't leave this out..
That's weird, right?.
He could have just left it at love God with everything.
and that could have been it, but he doesn't..
He goes into more detail..
He includes and love your neighbor as yourself..
Now, this is the part where if you're not a believer,.
if you don't actually consider God as your God,.
you don't look at Jesus and say, yeah, that's my Lord..
You probably would sign onto this..
You would probably say, hey, I can get behind that..
Right?.

$^{481}$I might not like a lot of the rules.
or a lot of the things that tells me not to do in the Bible,.
but this part, this seems good..
This seems fine, right?.
I can sign off on this..
I can say that this is a good thing,.
even if I don't vibe with everything else in the Bible..
Here's the reason Jesus says.
that they are intertwined though, right?.
Jesus says that these two commandments are very intertwined.
because you are not loving God.
if you are not loving your neighbor..
I'll say that again..
You are not loving God if you are not loving your neighbor..
And conversely, you cannot love your neighbor.
if you are not loving God..
You cannot love your neighbor if you are not loving God..
I'll say those two one more time..
You are not loving God if you're not loving your neighbor..
End of the story..
You cannot love your neighbor if you aren't loving God..
These two things, these two concepts are so intertwined.
that you can't have one thing without the other..
I can't say to God, you're really great..
I like that you died for me..
I'm gonna make you my God,.
but I really, really hate other people..
I'm just not gonna deal with them..
You can't do that..
And Jesus is also putting forth the idea here.
that you can't actually love your neighbor as yourself.
if you are not first loving him..
This is an interesting idea, right?.
For those of you who are in Summit, you might remember this..
I do, as Pastor Mark said, I do youth ministry here.
and we work with students in the church.
and students in the neighborhood..
And this past semester, we actually went over this topic..
We went over this verse.
and we were specifically looking at the idea.

$^{521}$of loving other people as yourself, right?.
Now, where we're at in Summit,.
where we're at with the kinds of discussions.
we have in our youth group,.
it's really interesting because rather than preach,.
we just kind of discuss it and we talk about it..
And we're at a spot where a lot of people.
are willing to talk about the things.
that maybe don't make a ton of sense to them,.
or they're willing to say, this doesn't make sense to me..
That seems like craziness..
Let's talk about why that is..
And when we were talking about this idea,.
we had a young woman in our discussion who is whip smart.
and she's awesome, but she immediately locked on.
to some idea about this concept..
She said, okay, loving your neighbor as yourself..
And as we're talking about what that actually means,.
as we're talking about what that actually looks like,.
she goes, that is a terrible life policy..
She identifies, she identified that this was not going.
to put you ahead, right?.
Loving your neighbor as yourself,.
giving them your spot in line,.
preferring the things they want over what you want,.
giving your time and money to other people.
instead of using it for yourself..
Like to the point where you would actually prefer.
other people's good over your good all the time..
That does not make sense..
Can you even build a life on that?.
Can you actually get anywhere with that?.
Aren't you always going to be left in need or wanting.
if you actually do that?.
And here's the thing, she's right..
She's really onto something there..
There is no reason to love other people.
more than you love yourself.
if you do not have God in the equation, right?.
If you are a just random person here by accident,.

$^{561}$making it alone and making it in your life alone.
and you're watching out for number one,.
there really isn't any reason.
to value other people so highly..
There's definitely a reason to care for other people..
There's a reason to be kind for them..
There's a reason to love them, sure,.
but to love them more than you love yourself,.
love them as yourself, what?.
That makes no sense..
This commandment only makes sense.
if there is first a God who has loved you so much.
that it would be crazy not to love your neighbor.
as yourself..
If the God of the Bible, right?.
If the God that is Jesus actually loves you,.
if he actually died for you,.
if he actually cares for you,.
if he's actually watching out for you so much.
that there will never be anything you need.
that he will not provide,.
whether that's physical, mental,.
whether that's financial,.
whether that's emotional or relational,.
if there really is a God out there.
who cares that deeply about you,.
it would be crazy not to love your neighbor as yourself..
If you've been given everything,.
why would you not want to share that kind of everything.
with other people?.
These commandments are intertwined..
For one final time,.
I want you to think about the mission statement.
that you thought about at the beginning..
I want you to think about what your life is defined by..
An easy way to do this is to think about your time..
How much time do you spend doing this thing.
versus that thing?.
Or how much of your thought life.
do you devote to something?.

$^{601}$Is your thought life always focused on the next promotion?.
Are the things coursing through your mind.
how you're going to get to the next step financially,.
that next stage?.
Are you consumed with another person.
and their thoughts of you?.
Is that where you get yourself worth?.
What are the things that drive you and motivate you?.
What is your mission statement?.
The reason I wanted to talk about this today.
is because what Jesus is putting forth.
is the only mission statement,.
the only foundation that is going to be satisfying..
If you live your life Saturday to Saturday,.
exhausted throughout the week,.
just looking for the next moment.
that you're gonna get for yourself..
If you live your life looking for purpose or meaning.
and feeling like there's just none of it..
If you live your life feeling lonely.
and exhausted of being lonely,.
what Jesus is saying here and putting forth.
is something that is going to matter to you..
It's something that's gonna fill you up..
And think about what your life would look like then..
You're not only fulfilled,.
you're also fulfilling others..
You're showing them,.
there's a God who loves you..
And I love you because he first loved me..
I get to show you that I can value you.
more than I value myself,.
because God valued me more than he valued himself..
Everybody's life has some sort of foundation,.
some sort of mission statement..
Your actions don't happen randomly..
There is a value behind what you do..
And what Jesus is saying is that.
if you want to be fulfilled,.
if you want to be happy,.

$^{641}$if you want to know that other people are also loved,.
if you want to be able to love other people.
more than yourself,.
you have to know that first you are to love God.
with everything..
There is not a lot of things in this world.
that are actually worthwhile..
The older you get, the more you realize that..
So much of what you do and care about now is dust later,.
or doesn't matter years down the road..
But what Jesus is saying is that this does..
And if you can't say with honesty.
that God has changed your heart to the point where.
you are able to love other people more than yourself,.
he's saying that you can have that,.
that all you have to do is love him first..
Because the truth is he's loved you before that..
Let's pray..
God, we thank you that these are not just words.
you're telling us,.
that it's not just some commandment you're giving to us.
and saying, good luck with it..
God, we thank you that this is something that Jesus did,.
that this is something you did first..
That Jesus himself loved God with all of his heart,.
soul, mind, and strength, all of his emotion,.
all of his intelligence, all of who he was,.
all of his passion..
And that because of that,.
he was also able to love his neighbor.
more than he loved himself..
That when he's in the garden of Gethsemane,.
on his face, sobbing and sweating blood,.
because he's terrified of the cross,.
that he prays to you and says, your will be done,.
and comes to the conclusion that he still loves his neighbor.
more than he loves himself..
Thank you that rather than running,.
he turned and he walked towards the cross of his own will,.
because he loves us..

$^{681}$God, would you give us a heart like that?.
Would you help us to live for something.
that actually matters instead of the rat race.
that this world is?.
Instead of all the things we could pour our heart and mind.
and soul and strength into that is going to go away soon..
God, give us the wisdom to see that this.
is how we should be spending our life.
and that this is something that is going to fulfill us..
And give us the conviction to follow through on that..
Not just to believe it's the case and say, that's nice,.
and I'll do that on Sundays,.
but to actually wake up every morning,.
so thankful and grateful that you love us..
\newpage



\section{5:1-9}
\label{sec:042xzVMibbM}
\textbf{一個人在 (約 5\_1-9) - 陳德晶傳道}
\newline
\newline
連結: \href{https://youtube.com/watch?v=042xzVMibbM}{\texttt{https://youtube.com/watch?v=042xzVMibbM}} ~~~~ 語音日期: 2022-11-14
\newline
\newline
\hyperref[sec:E_n46N8Pd_8]{\small{< < < PREV SERMON < < <}}
~
\hyperref[sec:index_chronic]{\small{[返順時目]}}
~
\hyperref[sec:index_scriptual]{\small{[返順卷目]}}
~
\hyperref[sec:i2GG1qfPho8]{\small{> > > NEXT SERMON > > >}}
\newline
\newline
$^{1}$大家祝福平安.
相信每個人都知道.
物想退收是一個很寶貴的屬靈傳統.
但是我們往往很少會走去物想退收.
反而近年我們很多人都會說Me Time.
每個人都會走去一些Me Time的時間.
因為可能可以令我們離開一些煩擾.
一個人原來都可以盡興.
沉默看星 聽到月兒呼應.
繼而平靜 到放一人之境.
嘩!這種感覺多棒.
大家也可以更加想像.
如果那一刻耶穌一起在.
陪你去聽風 看月光.
甚至和你有些對話.
那一刻多棒.
其實在《約翰福音》五章裡面.
都有一個一人之境.
有一個一個人在的境況.
我們一起去再看看.
《約翰福音》五章一朝九節.
這個人的處境又是怎樣呢?.
經文一開始不是提一個人.
反而是很詳細地說這個人所在的地方.
經文第一節是這樣說的.
「這是以後東正猶太人的一個節日.
耶穌就上耶路撒冷去.
在耶路撒冷靠著陽門有一個寺子.
在白蘭話叫做「北士大」.
旁邊有五個男子.
這裡躺著黑眼的.
腐褪的血氣.
富幹的許多病人」.
「北士大」這個詞.
其實是在說這個地方.
在耶路撒冷聖殿山東北面的一個小谷.
附近是聖殿的陽門.
大家聽到這個名字都猜得到.
是被那些憲制的羊牛出入的後門.
但其實這個名字本身有一個意思.

$^{41}$很有意思的.
是一個叫「傾流之處」.
說到傾流之處.
你就會猜到這個詞.
就跟它的中文譯名不同.
「北士大」完全不同的.
是一個相反的.
它說的這個谷.
其實形勢是一個天然的儲水池.
很大的.
不小的.
其實在當時.
在耶路撒冷裡面一個很主要的儲水池.
考古發現這個池.
好像一個日字的形狀.
一個很大的長方形.
中間有一個狼子分割了.
變成兩個小字.
所以就有五個狼子在當中.
約翰提到一個讀者.
要留意一件事.
要用希伯來文去理解「北士大」.
原來希伯來文的意思是「憐憫自惡」.
所以你明白.
在第三節指出.
在這裡躺臥了一些黑眼的.
腐腿的.
血氣枯乾的.
有很多病人.
在這裡做什麼?.
等.
等天父傾流出來的那份憐憫.
這裡很詳細講到這個人所在之處之後.
鏡頭就慢慢放到這個人的身上.
在第五節.
他沒有講到這個人有什麼病.
按照他的回應.
你能夠知道.
他起碼看得到水洞.
又能夠走得遠一點.

$^{81}$而且能夠生存這麼多年.
所以你應該猜到.
他不是黑眼.
又不是腐腿.
或者血氣枯乾.
甚至可以講得清楚一點.
他應該不是絕症.
或者是重病.
事實上經文很特別.
他沒有關心過他有什麼病.
他只關心他在哪裡.
他在哪裡躺臥.
在哪裡住的意思.
或者他關心的.
是他在那裡的時間.
已經38年.
到底為什麼會講時間和地點.
和這個人有什麼關係.
有什麼特別.
我們很想知道.
但未答之先.
這一刻耶穌就出現.
很特別.
病人無數.
重病的也不少.
但是耶穌偏偏只聚焦.
專注在這個人的身上.
莫非這個人好像其他神職的主角一樣.
在大叫.
耶穌救我 耶穌救我.
我猜你也看到.
不是這樣的.
情況完全不是這樣.
這一次是耶穌很主動提出要幫助他.
他很簡單做了三個動作.
針對同樣是他的時,地,人.
他先做什麼呢?.
他看見他躺臥的地方.
之後他知道他病了很長的時間.
於是他就問這個人.

$^{121}$你要得全愈嗎?.
這個問題其實是明知故問的問題.
但耶穌為什麼要問這個問題.
要對這個人問這個問題.
不如我們聽聽這個人如何回答.
第七節說.
先生 水凍的時候會有人把我放在池子裡.
我靜去的時候就有別人讓我先下去.
經文除了很客氣地稱呼耶穌之外.
這個人再一次提到時間.
而且有兩個提到時間.
他說水凍的時間和我靜去的時間.
應該在說自己一直錯失了一些時機.
帶來一些無奈,無望.
因為相傳在這個靈魂之屋裡面.
會有天使在按時攪動池水.
當水凍之後.
只有一個人能夠早點進去.
無論是什麼病都能夠得到全愈.
他錯失了這個時間.
所以很無奈,無望.
他也提到第二點.
就是地點.
他說在這個池子裡.
這個池子可能代表他容身的地方.
也可能代表他依賴的方法.
一直是靠這個池子.
他才能夠得救.
或者甚至可能.
同樣代表他容身的地方.
和依賴的方法.
除此之外 他也提到人.
但他指出沒有人能夠幫助他.
反映著他沒有人關心.
極度無助.
提到人另一個很特別的地方.
你剛才聽我讀的時候也聽到了.
在短短一節的經文裡面.
他重複了三次說「我」.
不斷聚焦在自己的世界.

$^{161}$自我的中心裡面.
在當中拼命一切.
但是也沒有用.
我們可以想像一下.
去那個時空.
一個人在等.
一個人的希望.
開始時滿懷幻想.
正所謂人生沒有一個的夢想.
和鹹魚沒有什麼分別.
但是經歷無數次的失敗.
一等就三十八年.
等到自己真的變成一條鹹魚.
沒有無奈,沒有望,沒有助.
只有這個錢給他容身.
是唯一可以依靠的方法.
不過在心裡想.
況且做到一個已經習慣了做乞丐.
懶做工.
你說他正在逃避現實的世界也好.
或者很害怕面對一個新的生活也好.
他也不會再有力.
去面對其他的生活模式.
當你聽到這樣的形容.
我不知道你有沒有.
腦海當中有沒有想到一些人.
或者一個人.
其實也很普遍.
美國有一個心理學家叫做Martin Solomon.
他在1967年做了一個很經典的實驗.
他分別有兩批狗.
放在不同兩個籠子裡面.
當中籠子裡面有一個按鈕.
這個按鈕有什麼用呢?.
就是會有警報.
警報聲之後.
就會有些電擊給這些狗.
而第一批狗當他按到這個按鈕的時候.
就能夠停止這個電擊.
但是第二批狗.

$^{201}$即使有這個按鈕.
但是不能夠阻止這個電擊發生.
這個實驗做了很多次.
發現有一樣東西很有趣的.
重複了這麼多次之後.
第二批狗.
即使開了那個閘.
那個籠的閘.
但是第二批狗.
它不再想任何的逃走方法.
它只會在那個籠子裡面坐以待斃.
甚至它其實有機會走.
它也不走.
不知道為什麼會這樣.
會在當中很痛苦地去滲入.
Martin Sullivan.
人類也有同樣的觀察.
他說如果有一個人在一件事.
或者在當事件裡面.
去重複,持續去受到一些打擊.
痛苦的話.
是會很無奈,無望,無助,無力.
即使可以改變.
有改變的機會.
它也不會掙扎,選擇.
不會有任何想去改變這件事.
只會在這裡等.
在這裡等死.
所以他給了這個現場一個名字.
叫做Learn Heaviness.
學習得來的絕望.
我在《木戀》裡面.
其實也見過不少的人.
有些人因為失業,失戀,長期病患.
情緒病,隱蔽青年.
坐牢,吸毒,賭徒,走鬼.
或者因為社會的狀況的轉變.
都陷入了這個學習得來的絕望當中.
有什麼可以讓我們在絕望裡面有出路呢?.
我猜還是要回到聖經.

$^{241}$我們看看耶穌怎樣幫助這個人在絕望.
第八節耶穌就跟他說.
「起來!拿你的玉子走吧!」.
「起來!」.
你聽到耶穌在用一個命令的語氣.
三個動作都是一個命令的語氣.
顯明出祂是有神的兒子這個權柄.
「起來!」.
不僅是提醒我們祂是神的兒子.
更加好像提醒我們.
「不!現在就是那個時機!」.
「起來!」.
你可以有起身的希望,選擇.
耶穌似乎在幫助這個人.
針對這個人.
去對時間的迷信.
祂說「你要拿起你的玉子」.
其實耶穌在這裡.
大可以只會說「拿起玉子來」.
但是你明顯看到.
祂刻意加了「你」這個字.
似乎是在回應那個人.
不斷提醒「我」.
在幫助這個人.
讓他感受到一件事.
去耶穌這個人.
是在乎「去」.
祂自己有價值.
祂會否得到憐憫.
不需要再一次只困在自己的世界.
你可以「滴起心肝地」.
帶來改變.
耶穌再一次針對那個地方.
祂說「起來!」.
是要離開這個磁鐵.
不用再貼著那張床.
就是一直去那個容身之處.
那個依賴.
能夠幫助他得救的方法.
耶穌在這裡讓他知道.

$^{281}$其實方法不只是一種.
應該這麼說.
方法就在耶穌裡面.
其實馬丁·塞勒曼.
也有同樣的觀察.
說絕望的人.
通常都差不多有相同的思維.
在時間方面.
往往會想到一些很短暫的事情.
會變成一個很永恆化的.
他會常常想.
我永遠都不會得到出路.
我永遠都會是這樣的狀況.
面對地.
他的處境同樣是這樣的.
他會將很局部的處境.
變成一個很全面的處境.
他會不斷去看.
即使我有方法.
但也沒有什麼用.
是幫不了我的.
或者他說.
既然沒什麼用.
那我為什麼要去治療呢?.
在這裡等死就算了.
或者我死了也好過了.
往往會局限了自己得救的方法.
只能看到眼前那個方法.
或者自己的處境.
也在另一方面.
說到人.
這是最重要的.
當絕望的人.
往往他會對人的信念.
很有自己的看法.
往往會很自責.
自我中心.
他認為沒有人會再關心我了.
所有問題都是我自己扭來的.
不關其他人的事.

$^{321}$但是我現在很大問題.
甚至這樣說.
覺得自己毫無價值.
我死了也好過了.
簡單而言.
你會看到.
通常絕望的人.
都會在時,地,日裡面.
不斷轉後角線.
他指出.
Solomon 指出一樣東西.
只要我們在形字上扭轉這些思維.
其實是可以得救的.
我們不需要再在糾纏.
在時,地,日的方面.
我們是可以得救的.
但你再看聖經.
很特別.
耶穌很厲害.
祂比Solomon早一千多年之前.
就已經理解到.
並且幫助這個在絕望裡面的人.
讓他有不同的信念.
我剛才也有提到.
我陪伴了不少陷在絕望當中的人.
沒養那麼多年.
就常常都見到這些人.
其實陪他們一起面對絕望.
我也沒有什麼方法.
甚至這樣說.
越陪就越絕望.
當我真是一個人.
在默想這段經文.
我投入自己.
是像一個人一樣的時候.
其實我心裡面的那份絕望.
突然間都覺得很被激勵.
很感動.
我真的哭了出來.
我看到的是.

$^{361}$普天下 病人無數.
但主在際邊 陪我寂寞.
那份耶穌的主動.
很感動我.
我以為我沒得可以幫.
或者我都幫不到人.
很無奈.
但原來耶穌很主動地走過來.
很想幫我.
很想幫更多絕望的人.
舉個很簡單的說話.
其實我底下不太明白.
我聽不明白他說什麼.
為什麼這麼簡單.
但不知道為什麼有一份動力.
很想試一下起來.
試一下再來.
試一下用多一點力.
去幫一些絕望的人.
但很奇怪.
你再看看那些猶太人.
他們只是掛著一些規矩.
但是我知道.
讀了這麼多年神學.
很知道這個安息日.
那個意義.
就是在乎完美.
在神創造天地的時候.
很想完美.
他看到自己完美就可以休息.
他也很想那些人完美.
而耶穌在這裡很在乎我.
很在乎這個人.
他要這個人.
要我.
都能夠得著完美.
他很在乎我健不健康.
不單只是心態.
是我的心靈健不健康.
但是當我還不知道可以怎麼做的時候.

$^{401}$他再走過來跟我說.
你已經得了痊癒.
不用再犯罪.
恐怕你做完更加厲害.
這個說話好像帶來一份關心.
其實我也不知道自己在做什麼.
我還有什麼可以幫助其他人.
或者幫助自己.
但是好像這個說話再一次提醒我.
不是的.
不要再問自己行不行.
自己的世界裡面是怎樣.
不要在這裡自憐自卑.
可以的.
走出來就可以了.
我相信這一刻.
你未必會有我的感受.
我其實真的很相信.
Me Time是可以幫助我們每一個人.
所以我也很想這一刻.
我特意說一段.
是很想這一刻.
試一下你有Me Time的時間.
你試一下可以有一人之境.
你試一下安靜.
默想一下這段經歷.
你試一下去想像.
自己身處在那個現場.
耶穌很主動跟你說.
跟你聊天.
你有什麼感受.
你有什麼回應.
我們安靜一下.
我相信這一刻.
你會有我的感受.
我其實真的很相信.
Me Time是可以幫助我們每一個人.
所以我也很想這一刻.
你試一下有什麼感受.
你試一下可以有一人之境.

$^{441}$你試一下去想像.
自己身處在那個現場.
讓我們一起有個祈禱.
你是很樂意.
傾流你的那份恩典給我們.
即使我們在絕望當中.
你很樂意主動幫助我們.
我再請你.
求主憐憫.
施忍.
\newpage



\section{}
\label{sec:i2GG1qfPho8}
\textbf{聖經地理 EP9 - 地理與釋經 - 地景 (陳崇基牧師)}
\newline
\newline
連結: \href{https://youtube.com/watch?v=i2GG1qfPho8}{\texttt{https://youtube.com/watch?v=i2GG1qfPho8}} ~~~~ 語音日期: 2022-11-15
\newline
\newline
\hyperref[sec:042xzVMibbM]{\small{< < < PREV SERMON < < <}}
~
\hyperref[sec:index_chronic]{\small{[返順時目]}}
~
\hyperref[sec:index_scriptual]{\small{[返順卷目]}}
~
\hyperref[sec:JhyniyhyQUs]{\small{> > > NEXT SERMON > > >}}
\newline
\newline
$^{1}$迦南地的地景是怎樣呢?.
它的地勢等等.
我們希望用一個短的時間來整理一下.
左邊是一個做得很好的.
雖然是英文的地圖.
由南至北有一條紅線.
由西至東有一條淺藍色的線.
為甚麼畫兩條線上去呢?.
由南至北紅色這條線.
做一個橫切面的畫.
將它斬開.
看看它的高高低低.
就是上面這個圖.
由北面開始.
Upper Gallery 上加利利的地方.
高山比較多.
然後就下來.
我剛才提過下加利利.
Lower Gallery.
就是那撒勒的地方.
那撒勒在這裡.
它的山峰高度大概是.
不夠二千呎高.
這裡就是Lower Gallery.
大概在那撒勒附近.
這個地方.
塔伯山等等.
塔伯山不是尼泊爾山.
這個就是Lower Gallery.
在那撒勒附近.
有一個平原.
一個大闊的山谷.
叫做耶斯列谷.
耶斯列谷.
在這裡.
Jesville Valley.
這是一個通往西東.
一個很重要.
很方便的路線.
這裡就是Jesville Valley.

$^{41}$就是耶斯列平原.
然後一直向南.
就開始往中央山脈.
就是這地方.
今天有一個很大的.
Palestinian的城市.
在West Bank.
Palestinian Territory.
巴士丹地區.
這裡就開始.
上到差不多一千呎以上.
然後一直向南的時候.
就去到.
以巴魯山.
和基里森山.
Mount Gerizim.
你記得這個地方.
就是事件的地方.
這裡就是約書亞入主了迦南之後.
他們最後就去到事件的地方.
在兩個山頭上.
來到立約和耶和華.
有祝福.
有奏坐的說話.
在約書亞記第24章.
這個地方.
如果你今天看過地景.
你都看到是一個十字路口.
Shakim在哪裡呢?.
我們先找Shakim.
就是事件.
事件就在這裡.
今天就是Nablus.
另外一個大的巴士丹的城市.
你看到.
你這樣看的時候.
Shakim就好像一個T字路線的中間.
這個是一個十字路口.
又是一個中間之地.
然後又是一個T字路線.

$^{81}$又是一個中間之地.
的中間之地.
最多交通經過.
南北東西.
前面我們看到有地中海.
地中海.
然後過了沿海平原.
再進來的時候就是薩馬利亞山地.
然後我們就看到有.
以巴盧山.
和.
這個基利森山.
就是結在.
這個事件的.
一北一南.
Galaxy.
Mount Galaxy.
基利森山.
看到嗎?.
我們zoom out.
一點的話你會發現.
由北至南一直下來.
就是去到這個事件.
然後北面就是.
以巴盧山.
有一個找到一個鐵器時代.
舊約時代的.
約書亞時期的大概1200BC.
一個祭壇的地方.
這個相信就是.
約書亞記.
裡面說到他們.
在以巴盧山建祭壇.
的位置.
然後這裡就是基利森山.
在這裡他們立約.
基利森山今天.
在山頭上還有.
那些叫做薩馬利亞人.
有2000年的歷史.

$^{121}$還有一群薩馬利亞人.
自己一個村莊.
他們自己內族通婚.
現在只有800-900人左右.
他們就在這裡.
世世代代.
在這個山頭上.
建立他們的地方.
有他們自己的宗教傳統.
薩馬利亞人就.
只是接受.
摩西五經.
還是接受他們那個版本的摩西五經.
不接受整個的.
希伯來聖經.
我們叫的九六聖經.
所以這裡又是一個.
十字路口.
來到基利森山和以巴盧山.
一直繼續下去.
你會看到耶路撒冷.
伯利恆.
這個所謂的Hill Country.
北面的Hill Country.
就是馬來西亞和以法尼亞的山地.
下面.
南面就是.
伯利恆附近.
耶路撒冷附近.
這些的.
很大的山地.
一直下去你會看到整個的山脈.
去到下面就是希伯倫.
希伯倫.
族長時代.
伯利恆以撒亞國很多時候就是希伯倫.
這個地方做的基地.
事實上伯利恆買第一塊.
的墳地就是在希伯倫.
是不是.

$^{161}$然後過了這個中央山脈.
再下去的時候.
再向南就下去.
一千多呎的.
海拔以上.
這裡就叫做Negev Basin.
就是這個南地的盤地.
這裡最重要的地方.
就是Besheva.
就是.
這個.
Bessabala.
所以這裡就是比較乾旱一點.
就叫做Negev.
然後再下去就是Negev Highland.
再下去下面.
就是比較高山一點的.
南地.
如果你看聖經的話.
你會想起這種地理.
譬如.
在這個.
創世記第十二章.
第十二章.
你記得嗎.
就是亞伯蘭回應上帝的呼召.
在萬國得福.
去到中間之地.
他就照著耶和華的吩咐.
羅德也帶著他.
然後出了哈蘭的時候.
他就帶著妻兒.
他們的產物等等.
往迦南地去.
他們就來了迦南地.
經過那地.
到了第一個地方是什麼.
就是事件的地方.
那時迦南人住在那地.
然後上帝就跟他們顯言.

$^{201}$耶和華跟他們顯言.
他就捉了一座壇在事件.
然後再從那裡.
又遷到去伯特利東邊的山.
支達象棚.
西邊的伯特利.
東邊的艾.
就是在伯特利和艾中間.
那裡又為耶和華捉了一座壇.
求告耶和華的名.
宣講耶和華的名.
後來就去哪了.
南地.
你記得嗎.
進了迦南地之後在事件捉了一座壇.
在伯特利捉了一座壇.
伯特利和艾中間捉了一座壇.
然後就去到南地.
如果你看地圖.
就是由北至南.
事件.
伯特利.
如果找回伯特利在哪.
北部.
就是在這裡.
伯特利和艾中間.
所以他跟著中央的山脈.
一直向南走.
下去.
Negev.
南地.
別士巴這個地方.
由北至南一直下來.
他跟著這個分水嶺.
一直到半遊牧民族.
經過不同的地方.
所以你見到他.
就是這樣去走.
去到Negev.
由南至北.

$^{241}$見到上帝賜給他將來的.
這塊地.
這個紅線.
然後我們就看.
藍線.
藍線就是由西至東.
西至東.
你見到.
這裡就讓你看到.
由地中海開始.
然後就是阿什特.
這個地方.
然後再進入一點.
就是到低山地.
就是Ekron.
以勒倫.
其中一個非利士人的.
大的城市.
Ekron.
再上到去就是.
這個地方.
很多時候.
三武爾的上.
四章五章.
很多時候在這個地方發生.
然後經過.
然後再上去.
這個中央山地.
你見到分水嶺就在這裡.
就在耶路撒冷附近.
去到耶路撒冷.
然後二千多呎的海拔以上.
再下去的時候就去到死海.
就一直滑下去.
然後這個地方就寸草不生.
因為雨雲.
由西面吹到東面的時候.
雨水都是下在.
分水嶺.
以西的地方.

$^{281}$左邊的地方.
以東的地方.
就不會有雨水.
就變成有大礦業.
一直下去到死海的.
峽谷.
然後再上去約旦河東的.
高原.
你見到Mount Nebo.
尼泊山的地方.
如果你看看這個立體的話.
大概你都.
get到是這個意思.
由地中海.
一直過來這裡.
你見到農田都是比較大塊.
因為是平原.
這裡簡直是機場.
今天在特拉維夫的機場.
然後你經過另一個大城就是基適.
再過來這邊.
剛才我們提到白事墨.
這個就是低山地的地方.
低山地.
就是中央山地.
和沿海平原.
中間的.
gateway.
關口.
門口就在這裡.
所以你記得.
亞西加.
大衛和哥利亞打仗.
大家去到.
單對單搏鬥的地方.
綠色這個地方就是伊拉谷.
伊拉谷.
你見到我們現在.
往東看的時候.
就是這個.

$^{321}$猶大的山地.
中央山地.
你再看多一點.
當然你在伊拉谷.
看不到穿不過這個山.
這個山頭後面.
就是什麼地方.
死海.
猶大礦業.
回來這裡就是中央山地.
然後再下來.
這個就是赤斐拉.
這個.
低山地的地方.
然後再下來就是平原的地方.
為什麼會在.
亞西加附近的伊拉谷.
大衛和哥利亞.
猶大人和以色列人.
和非利士人會有衝突呢.
因為這個就是低山地.
就是那個gateway.
上面就是以色列人.
入主迦南聚居的地方.
下面平原包括.
加特人.
哥利亞.
哥利亞就在.
加特這個地方出來.
這個平原.
就是非利士人的地方.
上面就是以色列人和猶大人的地方.
所以他們經常有衝突的地方.
就是在低山地.
這個就是一個緩衝區.
那些約櫃的故事.
都是在這些地方發生.
因為兩者兩個的.
勢力就在中間.
所以很多時候會有衝突.

$^{361}$所以你這樣看的時候.
你會發現原來這個地景.
可以幫助我們看到聖經裡面.
發生了什麼事.
非利士人在沿海平原的地方.
最好的地方.
恆商發達的地方.
都在這裡.
以色列人就在山.
他們要發展.
要去到擴充的時候.
就一定要下來下面.
然後非利士人就和他們有很多這樣的張力.
這個就是.
讓我們一個地景.
可以明白到聖經裡面的一些經文.
\newpage



\section{}
\label{sec:JhyniyhyQUs}
\textbf{Cultivating Gratitude (Psalm 100) - Rev. Sam Johnson}
\newline
\newline
連結: \href{https://youtube.com/watch?v=JhyniyhyQUs}{\texttt{https://youtube.com/watch?v=JhyniyhyQUs}} ~~~~ 語音日期: 2022-11-16
\newline
\newline
\hyperref[sec:i2GG1qfPho8]{\small{< < < PREV SERMON < < <}}
~
\hyperref[sec:index_chronic]{\small{[返順時目]}}
~
\hyperref[sec:index_scriptual]{\small{[返順卷目]}}
~
\hyperref[sec:xHuMbJI2NdE]{\small{> > > NEXT SERMON > > >}}
\newline
\newline
$^{1}$Thank you, James, and thank you, Jason, as well, for leading us in worship..
I'm excited to be here and grateful to be here with you all, and always really enjoy.
being at UV, so glad that I get to be with you today..
And as Jason was reminding us, we are headed into the holidays, and the first kind of American.
holiday we have this time of year is Thanksgiving, which is a great holiday..
My daughter and I, as we were driving here, we were talking about the things we enjoy.
about Thanksgiving..
And there's lots of obvious things, like the food, and spending time with family, those.
things are great..
But it really is a good time for us to think about being grateful, and practicing giving.
thanks and practicing gratitude..
So that's what we're going to talk about today..
I thought since we're headed towards this time anyway, we should talk a little bit about.
what it means to be grateful and to cultivate gratitude in our hearts..
And gratitude is really an important thing for us..
It's a God-honoring value that Scripture commands in a lot of different places..
The Scripture tells us to practice gratitude..
First Thessalonians 5.18 says it real plainly, real clearly..
It says, "Give thanks in all circumstances, for this is the will of God in Christ Jesus.
for you.".
Can't get much more clear than that..
Give thanks for everything in all circumstances..
But even if that wasn't enough, just the command of Scripture, we know that practicing gratitude.
has a lot of obvious benefits in our lives..
There's a lot of scientific studies have been done to show the benefits of gratitude..
It improves our physical health, our psychological health, it strengthens our relationships,.
it helps us to be more resilient as people, it gives us a better view of ourselves, and.
it even helps us to sleep better..
So there's all these reasons that we should be learning to be more grateful people..
And yet, it's not necessarily something that is easy for us..
A lot of us struggle to practice gratitude..
I know I do..
It's not something that we just naturally do..
And I think the reason for that goes back to who we are as human beings..
The fact that we're fallen, we're broken, we're flawed, and we live in this broken world..
And so there are things in our world that disappoint us..
And we disappoint ourselves..
And things in life don't go how we want them to go..
And because we're broken and flawed human beings, we're inherently selfish..
We kind of know how we want life to go, and we want it to be focused on us and good for.

$^{41}$us..
And when that doesn't happen, we are frustrated, or we're disappointed, or we're hurt, or we're.
confused, and it's hard for us to be grateful..
That's one of the reasons that gratitude just doesn't naturally come out of us most of the.
time..
And some people may, their personality, their gifting, maybe they're just more grateful.
than others..
So if you're one of those people, that's beautiful, and that's fantastic..
But for others, you may say, "Well, it's not easy for me, and I don't find myself being.
grateful very much..
I don't even find myself saying thank you or things like that.".
I think for all of us, we can learn to be more grateful..
Even if you're already a grateful person, there are more ways you could practice gratitude..
But if you're someone who struggles with that, then certainly you could grow in that..
And the point is that because it's not natural to us, we have to cultivate it..
That's why we're talking about it today and saying cultivating gratitude, because it's.
not something that's just going to naturally spring up in our life..
We have to cultivate it, just like you cultivate plants..
You have to plant the seeds, you have to water and care for the plant and protect it from.
things that would damage it so that it will grow and bring fruit..
And that's how gratitude is..
We have to plant the seeds, we have to frequently water the plant and allow it to grow in our.
life and bring fruit of thanksgiving and gratefulness in our life..
So we're going to look today at Psalm 100..
It gives us a really beautiful template for thinking about gratitude..
If you'd like to look at Psalm 100 in your Bible or your phone app, we're going to read.
through the entire Psalm, but don't be scared, it's only five verses..
Some Psalms, if we were going to read through the whole thing, I would just read the Psalm.
and then we'd go home..
But this one's short, so we can take it..
This is Psalm 100, it says, "Make a joyful noise to the Lord, all the earth..
Serve the Lord with gladness..
Come into his presence with singing, 'Know that the Lord, he is God..
It is he who made us and we are his..
We are his people and the sheep of his pasture..
Enter his gates with thanksgiving and his courts with praise..
Give thanks to him, bless his name..
For the Lord is good, his steadfast love endures forever and his faithfulness to all generations.".
Now, hopefully, just reading that Psalm encourages us a little bit and fills us with a little.
bit of gratitude..

$^{81}$But this Psalm has some really practical clues in it that help us think about how to cultivate.
gratitude in our own lives, how to grow gratitude in our lives..
The first thing I want us to see is that we cultivate gratitude as we praise God for the.
wonders of his creation..
We cultivate gratitude in our lives as we praise God for the wonders of his creation..
You can see this in verse one..
It talks about how the whole earth is making a joyful noise to the Lord..
God's power and creativity, his wisdom, his care, his beauty, those things are displayed.
in nature and in the things that he has made..
Verse three reminds us that he's the one who made us..
He's the one who made our world..
So that reminds us that we should experience some gratitude as we look at this amazing.
world that God has made and as we look at our own lives and the amazing things that.
God has done in our lives..
Verse five reminds us of God's goodness displayed in his love for his people and throughout.
the generations of human history..
And those are important things to think about..
If we're having trouble being grateful, we say, "I really struggle with gratitude because.
I feel disappointed or I feel hurt or I feel confused about things in the world.".
A great place to start is looking at the wonders of the world that God has made..
Psalm 19 says that the heavens declare the glory of God and the sky above proclaims his.
handiwork day to day, pours out speech, and night to night reveals knowledge..
There's no speech nor are there words whose voice is not heard..
It's the first three verses of Psalm 19..
And that passage reminds us that no matter where you are on the earth, no matter what.
your language or cultural background is, you can look at God's creation and experience.
the wonder of it and you can also understand that God is big and powerful and good and.
full of beauty and full of wisdom just by looking at God's creation..
And that's a healthy thing for us to do..
Honestly, I believe that if we can't acknowledge the goodness of God in creation, then we're.
really practicing ingratitude..
We're really being ungrateful..
If you think about someone giving a special and unique and valuable gift to someone else,.
if someone receives that gift and feels no gratefulness for it whatsoever, then they're.
practicing ingratitude, right?.
And God in the world that he has made has given us this amazing gift of beauty and of.
wonder and we need to acknowledge it..
We need to acknowledge him..
Romans 1 talks about this when it talks about how God has made plain himself in nature..
Verse 21 in Romans chapter 1 says this, "For although they knew God, they did not honor.

$^{121}$him as God or give thanks to him.".
It's talking about people who looked at what God had made, they could see the reality of.
God through what he had made, but they chose not to acknowledge him or give thanks..
And that's ungratefulness..
And that's obviously not where we want to be..
That's not what we want to live out and practice..
We have to guard our hearts against this..
And one of the ways that we guard our hearts against that is just by making sure that we're.
aware of how amazing God's world is and how beautiful it is and all the things that God.
has made..
So there's different ways we can do this, obviously, but I kind of put it into categories.
as I was thinking about it..
These are different ways that I see some of the wonder of God's creation..
First of all, we see the aesthetic beauty of God's creation, right?.
And most of us are probably very familiar with that..
That's when we look at something and we say, "That is beautiful.".
Maybe it's a sunset..
Maybe it is where you go to the lakefront and we're looking out at the water..
Maybe you get out of the city and spend some time in the mountains or you're in the forest.
and you just say, "Wow, this is beautiful.".
And when you experience that, the aesthetic beauty of what God has made, you should feel.
a sense of gratitude and say, "I'm grateful to God..
Look at what he has done..
This is just amazing and beautiful and totally gratuitous.".
That's one of the things I love about nature..
God didn't have to make it that beautiful..
God didn't have to make it that....
I mean, it just does things to our heart and soul when we connect with it, when we see.
it, when we experience it..
God didn't have to do that..
He didn't have to use as many colors as he did..
And yet he did that to display himself and to display his creativity and his wisdom and.
his power and his love for us..
And so our heart should see that and respond and say, "God, I praise you..
This is amazing..
What a beautiful sight that I'm seeing or what a beautiful day it is.".
And Jason's right..
We get a lot of winter in Chicago..
But even that has its own beauty, right?.
Today is a nice day..

$^{161}$It's cool and crisp out there, but it's still beautiful..
The snowy days are beautiful in their own way..
Then it gets kind of dirty after that..
But we can enjoy God's beauty through nature..
Even just looking out our window into the yard, we can enjoy that..
The other way that I really try to reflect on God's goodness in creation is kind of maybe.
more from a scientific perspective, the scientific beauty that we see in creation..
What I mean by that is if you start to think about the intricacy and the complexity of.
the world that God has made, if you learn a lot about the biology of our own bodies,.
you start to say, "Wow, this is really an amazingly intricate and complex thing that.
God has done in us.".
Or if you learn a lot about the universe and the cosmology of the universe kind of on a.
macro scale, it's amazing..
It's mind-blowing..
It's more than we can wrap our heads around, the power of God that's displayed in that..
And those things, again, should cause us to say, "Wow, I'm grateful to God..
I praise God because you didn't have to do this..
You didn't have to show off like this..
You didn't have to make it this great, and yet you did..
And it's worthy of our praise..
Our gratitude belongs to the Lord.".
Then the other category that's really big for me is we can really be in awe of the human.
beauty of creation..
And I think that's significant..
And I don't just mean human physical beauty, although God created that too..
But I mean all the things that it means to be a human being that God made about us that.
we get to experience..
The fact that we get to have friendships, we get to enjoy laughing together over something.
that's funny, we get to taste food together and enjoy that experience, we get to listen.
to music, and we get to have love and relationships where we can't really, as human beings, we.
can't really describe, even from kind of a scientific perspective, we can't really describe.
why love is such a big deal to us, why we feel the way we do about a family member or.
a close friend..
Again, God didn't have to put that much beauty into us..
He didn't have to give us those things, but he did..
He gave us these friendships that we can't even really fully express and relationships.
with people that love us and that we love in ways that we were not even able to describe..
God made something really beautiful when he made humans, and the Bible's clear about that..
When God makes everything in the book of Genesis, humans, us, that was the masterpiece..
It was the last thing and the best thing that God made..

$^{201}$And we should be grateful for that..
As we look around and we see people and we see the way that they love one another, we.
see the beautiful things that happen through relationships in our lives, we see the things.
that other people create and that other people do..
We can be grateful to God for that wonder of his creation, the human element..
It's something that we can be in awe of and we should be..
The reality is, as we look at the wonder of what God has created, we start to sense our.
own smallness and it fills us with gratitude..
Hopefully we've all experienced that on some level..
If we haven't, then we're not looking, we're spending too much time looking at the sidewalk.
or something like that..
But all of us should be looking at the world God has made and saying, "Wow, I'm grateful.
to see this and it gives me a sense of how small I am and my place in God's world and.
I'm grateful that I get to be here.".
That doesn't mean that there's not challenging things there are, and that doesn't mean there's.
not things that we see that are hurtful..
There certainly are, but there are also things that we can see, even in difficult times in.
life that cause us to be grateful, that cause us to be grateful that we get to be in God's.
world..
I remember one time when I was in college, I went to college in Hawaii and I was snorkeling.
on the big island of Hawaii, this place called Kealakakua Bay..
I was pretty far out, it was a beautiful reef, probably about 20 feet deep, the water..
And then the drop off was over there where it just dropped off, it got down to maybe.
30, 40 and then just dropped off into the dark blue water..
If you've ever snorkeled in a place like that above the reef, the drop off is just really.
dark and you just can't see what's down there..
So I was having a great time looking at all the fish and the reef and it was colorful.
and amazing..
And I looked over to my left, I was kind of going this way, the beach was over here, drop.
off's over here..
I look over, I see something moving over by the drop off, it's kind of silver..
So I stop and observe for a minute and it's this big, it's coming up towards the reef.
so I can see it better and better as it's coming toward me..
It's this big tiger shark, probably 18 to 20 feet long, huge shark..
And the weirdest thing was what I remember feeling so distinctly in that moment..
I had my fins on and I felt like, you know, when you have fins on you feel like you can.
move pretty good in the water..
And I was just watching this thing cruising up and it didn't even look like it was moving.
its tail hardly and it was still moving like five times faster than me..
And it wasn't even trying, it looked like it was drifting and it was moving so fast..

$^{241}$And I just had this overwhelming feeling of being so vulnerable and out of my element.
because I felt how slow and incapable I was of going anywhere and I could see this giant.
shark like, it was easy for him to cruise through that water..
And obviously I went into the beach and got out of the water..
But it was one of those moments where I really felt my own smallness..
And when we look at God's creation enough, we'll start to feel that and say, "Wow, I.
can feel my own smallness and it makes me so grateful that I get to be a part of this.
world that God has made and I get to experience it and I'm grateful for what he's done.".
Next thing we see here in Psalm 100 is the gratitude that we can have for God and the.
beauty of redemption..
We cultivate gratitude as we praise God for the beauty of his redemption..
In verses three, four, and five of this Psalm, we're reminded that God is good and loving..
It tells us in verse three that he's the one who made us, we belong to him, we're his,.
we're his people, we're the sheep of his pasture..
In verse five tells us that the Lord is good, his steadfast love endures forever, his faithfulness.
to all generations..
God is full of care for his people..
He's full of love and grace and mercy..
God has done something in redeeming us that we, is beyond our capacity to describe and.
honestly beyond our capacity to really comprehend or even fully understand what we have received.
from God..
But if we will begin to think about it and consider it and reflect on it, our hearts.
will begin to be full of gratefulness..
Think about some of the things that it means for you, you personally, to be redeemed by.
Christ..
God has forgiven your sin..
He's covered your mistakes, those things in your past that you wish you could erase..
God has covered those things..
You may still have to deal with some of the consequences of them, but you have no guilt.
or shame from them..
God has erased them..
Ephesians one, verses seven and eight says this, "In him we have redemption through his.
blood, the forgiveness of our trespasses according to the riches of his grace, which he lavished.
upon us in all wisdom and insight.".
That's your story if you have a relationship with Jesus..
All those things are forgiven through his grace and he has lavished his love on you..
God has healed our hurts and scars..
We know that because we live in this broken world, all of us are harmed by different things..
We have traumas..
We have all of these scars in our life..

$^{281}$Through Christ, God can and does heal those things..
God is at work helping us to be made whole, to restore us, to reconcile us to other people,.
to bring us back to health in our relationships..
That's a beautiful thing that we can be grateful for..
God's rescued us from sin and death..
We don't have to be afraid of death..
We don't have to be afraid of sin..
Jesus defeated those things..
God's promised us a beautiful future..
First Peter one says this, "Blessed be the God and Father of our Lord Jesus Christ according.
to his great mercy..
He has caused us to be born again to a living hope through the resurrection of Jesus Christ.
from the dead to an inheritance that is imperishable, undefiled, and unfading, kept in heaven for.
you.".
That's your story through redemption..
You have a living hope that cannot be touched by anything in this world..
It will not be destroyed or diminished in any way..
It's kept in heaven for you..
That's your story because of redemption..
God's given us a new start in life..
That's one of my favorite parts of the gospel that we can be made into a new creation in.
Christ..
We know that 2 Corinthians 5:17 talks about that, "Therefore if anyone is in Christ, he's.
a new creation..
The old has passed away, behold the new has come.".
That's something to be grateful for..
That's something to praise God for, something to give thanks for..
We can be something new..
God gives us peace and joy and hope..
Romans 15, 13, I love this verse..
It says, "May the God of hope fill you with all joy and peace in believing so that by.
the power of the Holy Spirit you may abound in hope.".
Those are some amazing promises, amazing promises..
Type of stuff that we could jump and shout and yell on the street because of how powerful.
and beautiful that is..
The type of stuff that we should certainly be grateful for..
They come to us because of redemption, because of what Jesus has done in our lives..
That's something that should, as we look at it, as we think about it, as we reflect on.
it, it should fill our hearts with gratitude..
It's a way we can cultivate gratitude in our lives is looking at the redemption that God.

$^{321}$has given us..
We need to focus on the beauty of what He has done for us..
Sometimes when we think about scripture, we think about redemption in this grand narrative.
and God saving all people..
That is true..
That's the beautiful grand narrative of redemption..
Redemption also has a personal narrative that God saved you..
God's at work in your life and He's healing you from your hurt and from your past and.
He's making you a new creation and He's making you more like Jesus..
So when you start to think about this and be grateful for it, don't just be grateful.
that God is saving people around the world..
Yeah, be grateful for that..
But bring it home to your own heart and say, "I'm grateful that God saved me..
I'm grateful that I'm not who I was..
I'm grateful that even though I'm not yet who I want to be, that God is moving me towards.
that and He's doing something beautiful in my life.".
And you can be personally grateful for redemption..
When I think of this, one of the illustrations I think of is like a really skillfully crafted.
work of art, if you think about something like that..
Especially like if you think of a paint on canvas type of work of art..
If you go down to the Art Institute, they have this giant, some of the Renaissance art.
are like huge, like 20 feet wide and 10 feet tall, huge, huge..
You think about how long it would take to make a work of art like that, all the layers.
of paint that have to go on there, the type of vision that the artist has to have when.
they start to even know what they want to get onto the canvas, but then also the dedication.
they have to get it on there day after day for years they have to work on something like.
that..
And you also have to know that even the most precise artist is going to make some mistakes.
and they're going to get some lines that are a little bit off..
And so then they've got to go back and touch that up a little bit with another layer of.
paint..
It's an amazing accomplishment to bring something like that to completion, a work of art like.
that..
And yet that's what God is doing in our lives..
God has this vision for your life that's not just a one week vision, it's, there's a Christian.
song from when I was in high school, I've always remembered this lyric, it says, "I.
have a plan for you, it will only take a lifetime to see through.".
That's God's vision, it's for your whole life..
God sees the whole thing, he knows what he's doing in your life..
He's putting on layer after layer of paint to craft you into who he wants you to be..

$^{361}$And the whole time there's mistakes being made, right?.
Not by God, but by me..
God doesn't make mistakes, but I do..
I get outside the lines all the time..
And God in his providence and his sovereignty and his grace and mercy and love, he finds.
ways to craft even my mistakes into something beautiful and something that's gonna be for.
my good and for his glory..
But here's the crazy part, God's not just doing that for me, he's doing it for everyone.
in the world who has a relationship with Jesus..
He's doing it for all these people at the same time..
What an amazing thing it is..
And when you stand in front of one of those giant works of art, for me, I'm always kind.
of in awe and I think, "Wow, whoever made this, like I said, I give glory to them for.
what they have done.".
And so when we look at God's work in redemption, we should stand back and say, "Wow, this is.
amazing..
I can't even believe God is able to redeem all the mess in my heart, much less everybody.
else's heart, and then bring us together in all these different ways and then fix all.
of our mistakes even when we keep making them.".
And it brings glory to God because of what he's doing in redemption..
So that's a great place to cultivate gratitude..
Say, "I'm gonna grow gratitude in my life by thinking about what God's done for me,.
what he's doing in our world.".
And then the last big way I want us to think about cultivating gratitude is as we praise.
God for the majesty of who he is..
You know, the first two, praising God and cultivating gratitude as we look at the wonders.
of his creation, as we look at the beauty of redemption, those things kind of have to.
do with God's work in our world..
But here's the crazy thing..
Scripture is really clear about this..
God is worthy of praise even if he never did any of that, just because of who he is..
And we can cultivate gratitude in that way as well, just because of who God is..
Verse 3 here of Psalm 100 says, "Know that the Lord, he is God.".
Obviously there's other statements in that Psalm, but that statement is enough by itself..
Know that the Lord, he is God..
Even if there's nothing else, God is worthy of our praise because of who he is..
Our gratitude belongs to him..
Jehovah, the word in verse 3 there, the Hebrew word for God, means the self-existent one.
or the eternal one or the existing one..
That's the picture that we have of God in Scripture..

$^{401}$Completely apart from the wonder of creation and the beauty of redemption, God is still.
worthy of our praise..
And so when we contemplate who he is and reflect on who he is, our hearts should be full of.
gratitude..
Think about some of the attributes of God that Scripture revealed..
God's power, which is unlimited..
God's goodness..
It always blows my mind when I think about how powerful God is and yet he is also good..
If God was powerful and not good, he would be completely terrifying..
God is not just powerful, he's also good and he loves us..
His wisdom and knowledge, Romans 11.33 says, "Oh, the depth of the riches and wisdom and.
knowledge of God, how unsearchable are his judgments and how inscrutable his ways.".
God's wisdom and knowledge is completely beyond our comprehension..
God's holiness, his love for us, Psalms 100 talks about it, verse 5, "The Lord is good,.
his steadfast love endures forever.".
God's love for us, his grace, his mercy, God's justice, God's sovereignty, God's immutability..
That just means the fact that God doesn't change..
He's so right that he never needs to change because he's not wrong and he's not going.
to be wrong..
He never changes because he doesn't need to change..
The eternal nature of God, he's always existed, his omnipresence, he's in every place, his.
transcendence, his infinitude, which is crazy because when you think about the attributes.
of God and all the things that make up who God is, he's just, he's loving, he's all these.
things, you have to also remember that not only is he those things but he is those things.
in infinity and he is the very definition of those things..
We only know love because God is love..
We only know justice because God is justice..
We only know mercy and grace because God is mercy and grace..
God's perfection, not only is he those things in infinitude but he is also those things.
in perfection..
He is perfectly loving..
He is perfectly just..
He is perfectly holy..
He is perfectly good..
When we think about that stuff, we should say, "Wow, I want to praise God for who he.
is..
I want to be grateful for who he is.".
And God is all these things completely by himself..
He self-exists and is completely self-sufficient..
Back to this three, verses 13 and 14, this is when Moses and God are having this conversation..

$^{441}$God's come to Moses, said, "Hey, I want you to go to the people and go to Egypt and lead.
the people out of Egypt and I'm going to give you what to say.".
Moses is not wanting to do it so he's talking back..
"Then Moses said to God, 'If I come to the people of Israel and say to them, 'The God.
of your fathers has sent me to you,' and they asked me, 'What is his name?.
What shall I say to them?'".
Verse 14, "God said to Moses, 'I am who I am.'.
And he said, 'Say this to the people of Israel, 'I am has sent you.'".
Now, it's not quite as powerful in English as it is in Hebrew, but verse 14, God is really.
letting Moses know in no uncertain terms..
He's saying, "You want to know what my name is?.
You want to know who I am?.
You want to know who to announce?.
Announce the eternal, existent one, the one who has always existed and always will exist.
and who needs no one else and is totally self-sufficient and is without any degree of diminishment.
or need.".
That's what God's saying to Moses in that verse..
"I am the existent one..
I'm the one who has always been, always will be..
I need nothing..
Everything else came from me.".
That's the God that we serve..
That's the God of the Bible..
And he's worthy of our praise and our thanksgiving and gratitude just because that's who he is..
If you've ever been in a situation where you kind of experienced something that was really.
transcendent, gratitude was probably a part of that situation..
I know in times of my life when I've experienced something transcendent, gratitude just comes.
from nowhere..
One time when I was a youth pastor many years ago, I was in a service at a camp..
The presence of the Lord was really present..
We got done..
We did worship at the end..
It was time to go..
I told everybody to be dismissed, to go back to the cabins because we were supposed to.
go to bed..
And no one would leave the room because of God's presence..
And so everybody stayed..
So we started doing worship again and we just kept the service going and ended up being.
there much, much longer..
And God was just working in powerful ways..

$^{481}$Kids got saved that night, came to know Jesus..
It was really a transcendent thing..
It wasn't anything that I was doing or anybody else was doing..
God was at work..
And I was full of an overwhelming sense of gratitude just to be there, to be present..
It's not a lot of transcendent experiences in my life..
But another one that I can think of that's really powerful to me, every time one of our.
children has been born, I'm there with my wife and I see one of our children born, been.
there for all five of them..
And I'm reflecting on just the promise and the potential of a human life and what it.
means that this new life is taking their first breath in the world..
And I get to see it, get to see them take their first breath and then announce themselves.
with screaming and crying..
And it's always like just a beautiful moment..
And there's just emotion that I don't have to manufacture..
It just is overwhelming..
But there's also overwhelming gratitude..
Because I'm so grateful for that moment, for that life, for what God is doing..
And when we spend some time in God's presence and experience Him and start to see how transcendent.
He is, that's what happens..
We don't have to try to make ourselves feel grateful..
It just comes out..
We're overwhelmed with gratitude for who God is..
So I want to encourage us to seek to cultivate gratitude in our lives, to take our eyes as.
much as we can off of ourselves and put them on God and what He is doing..
And there's a lot more ways that you can practice it, but I would encourage you, if you don't.
know where to start, start with these three things that we talked about today..
Cultivate gratitude by praising God for the wonder of His creation..
Cultivate gratitude by thanking God for the beauty of His redemption in your life..
And cultivate gratitude just because God is..
And He is so amazing and so transcendent..
I'm going to close just by reading this psalm one more time and letting its words fill our.
hearts..
"Make a joyful noise to the Lord all the earth..
Serve the Lord with gladness..
Come into His presence with singing..
Know that the Lord, He is God..
It is He who made us and we are His..
We are His people and the sheep of His pasture..
Enter His gates with thanksgiving and His courts with praise..

$^{521}$Give thanks to Him..
Bless His name for the Lord is good..
His steadfast love endures forever and His faithfulness to all generations.".
Let's pray..
God, we just want to thank You for who You are..
We can't even describe it..
We can't even express who You are, God..
We are so small in Your presence and yet, as the word says, by Your mercy we're not.
consumed..
You still love us..
And You choose to come and be involved in our lives..
And God, we thank You for redemption..
That even though we were broken, fallen sinners living in a broken, fallen world, You wrapped.
Yourself in our flesh and You loved us enough to come into our world and to begin to heal.
us one heart at a time and to reconcile us with You and to bring us into right relationship.
with our Creator..
God, we thank You for that..
We thank You for this beautiful world that You have given us and allow us to live in..
And God, I pray for myself and for all of us, especially over the next few weeks as.
we have this time of thanksgiving, that You would help us to cultivate gratitude to begin.
to see how good You are in our lives..
And even in the most difficult moments, You are still at work..
\newpage



\section{}
\label{sec:xHuMbJI2NdE}
\textbf{Faith and Time: Shavuot - Rev. Mark Chan}
\newline
\newline
連結: \href{https://youtube.com/watch?v=xHuMbJI2NdE}{\texttt{https://youtube.com/watch?v=xHuMbJI2NdE}} ~~~~ 語音日期: 2022-11-28
\newline
\newline
\hyperref[sec:JhyniyhyQUs]{\small{< < < PREV SERMON < < <}}
~
\hyperref[sec:index_chronic]{\small{[返順時目]}}
~
\hyperref[sec:index_scriptual]{\small{[返順卷目]}}
~
\hyperref[sec:LC2O0QPKB1U]{\small{> > > NEXT SERMON > > >}}
\newline
\newline
$^{1}$And here is it, let's welcome Pastor Mark for today's sermon..
Hey, good morning guys..
I had a really good Thanksgiving week and we ate a lot in Toronto..
We spent five, six days, right?.
From Monday to Saturday..
Right, how many days?.
And we ate a lot of good food and we spent a great time there..
Nine of us went there..
So we did a lot of shopping..
I started to think in the third day or fourth day,.
I started to think my main job is actually just driving and paying,.
using my credit card to pay all the good stuff, right?.
So, you know, it just got so good, right?.
So Thanksgiving is just over and it's a holiday.
or it's a time of the year to remind us God has been good to all of us,.
you and me, and we need to give thanks, right?.
There are some calendar markers along the year to remind us.
what God is doing in your life.
and when to respond positively to God's good work..
And today we continue this series to talk about the festivals,.
the special days in the Bible..
And we talk about Passover, we talk about Sukkot,.
now we talk about Pentecost, what we call Shofoot..
Shofoot simply means weeks in plural, so many weeks..
So in the Old Testament calendar, it's seven weeks,.
seven Sundays after the Passover..
So it's the 50th day is the day of Shofoot, the day of Pentecost..
That's why we call it Pentecost in Greek, 50, right?.
It's a day that is special in the sense that there is nothing.
related to God's miracles in the history of His people in Israel..
We talk about Sukkot, right?.
Tabernacles, it has to do with God's miracle works.
in the wilderness, right?.
God is guiding His people out of Egypt and through the great.
and fearful wilderness for 40 years and giving them manna, right?.
Giving them food and water, miracle works, you know, for 40 years..
That's why they celebrate Sukkot and they make temporary shelters.
and they live in it for a week to remember what God has done.
in their history, in their lives..
And then the second example I give you is that most of the.

$^{41}$biblical calendar, most of the biblical festivals have to do.
with God's miraculous work in the lives of the people..
And Sukkot is one of them, Tabernacle is one of them,.
and Purim is another one..
Purim, if you remember, is a holiday that is set aside.
by the end of the book of Esther, right?.
When the Jews were separated or scattered across the Persian Empire,.
they were persecuted and God, you know, saved them and rescued them, right?.
And then the third one is mentioned in the Gospel of John,.
chapter 10, is called Hanukkah, right?.
Hanukkah is around the time of Christmas..
Sometimes we call it Festival of Lights..
It has to do with the Jewish people in the year of 160-something BC.
when they fought against the Greek Empire who persecuted them,.
oppressed them, and they won..
They won the battle and they gained their political independence..
That's why they cleansed the temple and they rededicated the temple..
They call it Hanukkah, right?.
The fourth festival we talk about is the Passover..
Passover, it has to do with Exodus, right?.
The 10 plagues, God defeated the oppressor Pharaoh.
and saved his people out of Egypt..
Now, all these examples have to do with God's great works..
You know, he has done some signs and wonders to save his people..
But Pentecost almost has nothing to do with miraculous work of God..
Now, let's take a look at this verse from Leviticus 23..
It says, "From the day after Sabbath,".
now, the previous context was talking about Passover, Passover..
"When Passover is done, from the day after the Shabbat,.
count off seven full weeks and then keep counting until the day after Sabbath,.
which is 50 days later..
Then present your grain offering to the Lord,.
whatever your harvest is, your first fruit, you present to God in the temple.".
So, it has no miracles..
It has to do with farming, harvest, you know, agricultural festival..
50 days after the Passover, the Pentecost,.
they bring in the wheat harvest and being reminded God is the provider..
God is faithful..
Now, as I said, there's nothing to do with the saving works of God..
No exodus, no wilderness, no miracles in the Bible, but they celebrate it..

$^{81}$Now, between the Old Testament and the New Testament,.
there are about a few hundred years in between..
We call it intertestamental period..
After the Old Testament, before Christ came in the New Testament,.
there were about a few hundred years..
The Jewish people in the land of Palestine,.
they would merge some of the dates of festival together..
So, in between the Old Testament and New Testament,.
the Jewish people would take the day of Pentecost to Sheffield.
as remembrance of God's giving the law to His people..
Now, it's about the same time, right?.
Because if you remember the people of Israel,.
they came out of Egypt at the twilight of Passover,.
and then exactly two months later, they arrived at Mount Sinai,.
and God gave them the law, right?.
Part of the law is Ten Commandments..
So, it's two months after Exodus, they came to Mount Sinai,.
they received the law, right? They became a covenant people..
They became people of the law..
And then the Sheffield O is 50 days after Passover,.
so they kind of merged that, right? It's around the same time..
So, they would use this time to remember God giving them the law,.
the revelation..
That's why in Pentecost, in between the Old Testament and New Testament,.
they sometimes call it Pentecost, but also called the Festival of Revelation..
God is giving them the law, His regulations, His will to His people..
So, it makes sense around this time they merged these things together..
But there's one note we need to make..
What is the law? What is the Torah?.
Now, this becomes problematic for us,.
because in modern translations, usually we translate the word Torah,.
which is Hebrew, into law or regulations in English, right?.
Or, you know, in some other modern translations..
You get the sense that law has nothing to do with relationship, right?.
You see a stop sign, you ran a stop sign, right?.
And then you got pulled over..
There's nothing relational about it, right?.
The police officer will give you a ticket, right?.
Or it hurts your wallet, right?.
You got fined, but there's no relational dynamics involved, right?.

$^{121}$You just violate some regulations..
You see the speed limit, 55 miles per hour..
You know, you went over it, and too bad if you got pulled over..
You know, it hurts financially, but there's nothing relational about it..
But in the Bible, Torah is not impersonal..
It is relational..
It has everything to do with your relationship with God, right?.
A few weeks ago, I was in Kansas City, and I was taking one of the seminars..
So we were talking about how we should translate the word Torah,.
because definitely it does not mean the modern sense of law..
How do you translate Torah?.
I think one of the examples given by the instructors, the professors,.
is really appropriate..
He says, you know, Torah in the Bible, law and regulations in the Bible,.
is like how, you know, if you're a parent, you tell a child,.
"Don't touch the stove," all right?.
"Don't touch the stove.".
You know, one-year-old, two-year-old, "Don't mess with the stove," right?.
Now, this is like the law, right?.
For a two-year-old, that makes sense, right?.
You don't touch it..
But when the child grows up, if it's simply,.
"Remember, just stay away from the stove," for the rest of his life,.
I think he obviously does not understand what it means..
When the parents tell the child, "Do not touch the stove,".
what does it really mean?.
It means use the stove properly, right?.
So when the child goes up, when the boy goes up,.
he knows how to cook, right?.
He will approach the stove and cook a good meal, right?.
That's what it means..
That's the essence of don't touch the stove, right?.
So as you grow up, you understand, you know, so-called the Torah..
Actually, you have to interpret it and apply it with wisdom..
As you grow up, you need to understand this is the essence of the law..
Now, think about Ten Commandments..
Ten Commandments..
I'll give you examples..
"Thou shalt not murder," right?.
Well, what does it mean?.

$^{161}$Does it mean, "Okay, I just don't kill anyone," right?.
"I don't murder anyone, so I'm fulfilling the law," right?.
No. No..
That's not what it means.
because it means to nurture a life-giving society,.
a life-giving culture..
That's why in other regulations,.
when they unpack what it means "Thou shalt not murder,".
they will say you have to protect, you know, the innocents..
You have to protect other people who are weaker than you.
because God loves every life, every human being..
"Thou shalt not commit adultery.".
What does it really mean?.
You cannot imagine a husband going home every night and say,.
"You know, I did not commit any extramarital affair..
I love you," right?.
No. That's not what it means..
It means culturing or nurturing a good relationship,.
a good family, a good marital relationship, right?.
That is what the commandment means..
It's relational..
"Thou shalt not steal.".
That does not mean, okay, I don't take anyone else's goods.
or I don't take their stuff..
No, it means to promote generosity..
What I have is not just fulfilling my needs and my wants..
I will give. I will give back to God.
because God is giving you abundantly..
That's why you make your offering..
That's why you want to use your resources to bless other people..
That goes beyond just calculating, okay,.
"Oh, it has an offering, you know, 10\% and all that stuff.".
You know, that simply means that we do not understand the Torah..
Just like the example, "Don't touch the stove.".
So if you think about the Torah,.
you think about the law of God given at the time of Pentecost..
You understand there is something we need to, you know,.
interpret and we need to understand, right?.
It's relational and there is some essence you need to understand..
That's why today in Jewish synagogues,.

$^{201}$during the festival of weeks, the Pentecost,.
they will read two passages from the Bible..
Now, one of them, of course, is what I have already explained,.
is 10 Commandments..
They will read this portion of Scripture.
to remember God has revealed Himself..
God has revealed His character by His law..
That's why as the people of God, we need to follow.
and to understand and apply the law of God in their lives.
to live out the life of God..
That's why they recite the 10 Commandments in the synagogues..
And another passage they will recite during this festival,.
festival of Pentecost or Shavuot,.
they will recite the Book of Ruth..
The Book of Ruth..
And why?.
If you remember the story of the Book of Ruth,.
it was the time of judges.
when the Israelite society was in chaos..
And Ruth was a Moabite..
And she left her own country.
and confessed faith in the true God Yahweh..
And she lived in the land of Israel..
And she was poor, right, with her mother-in-law..
They went through a lot of difficulties.
trying to settle down in a foreign land,.
in the land of Palestine..
Now, what does it have to do with Pentecost or Shavuot.
or, you know, the time of harvest?.
Well, it happened that in Chapter 2,.
Ruth the girl met her future husband, Boaz,.
in the time of Shavuot, in the time of harvest..
Chapter 2, "One day Ruth the Moabite said to Naomi,".
her mother-in-law, "Let me go out to the field,.
the harvest field, right, that people are harvesting, right?".
This is the time of Pentecost. This is the time of Shavuot..
"Let me pick up the stocks of grain.
left behind to anyone, anyone who is poor,.
and is kind enough to let me do it.".
So she went. And Naomi said, "Yeah, go ahead.".

$^{241}$And as it happened, by chance, by God's providence,.
she found herself working in the field that belonged to Boaz,.
the relatives of her father-in-law..
So it just happened she ended up in the field of Boaz.
in the time of harvest, in the time of Pentecost..
So the book, Ruth, the story of Ruth,.
is actually a harvest story, kind of in late spring,.
you know, in our calendar, May or April,.
for a foreigner to be able to settle down.
because of Boaz's generosity..
He is good to her and to her family..
He's generous enough to let her pick up the crops.
so she can live, she can flourish..
That is why the Jewish synagogues will recite the book Ruth.
in times of the festival of weeks,.
festival of Pentecost,.
because one of the most important themes in this season,.
in the festival, is how to practice generosity..
Remember, that's about the time.
the people of Israel had their harvest..
They got their paychecks..
It is right at this time they remember a story like this,.
to be generous because God is generous,.
to be good to others,.
to be kind to others..
Now I told you I spent a week in Toronto, right?.
Now people ask me, "What did you do in Toronto?".
It's the same thing, right?.
Typical stuff, right?.
We just eat, eat, eat, and eat, right?.
And then shop, and shop, and shop, right?.
I spent some time with some of the friends..
Some of them I've known for years, for years..
I met this--I had dinner with one of my friends from Hong Kong..
I haven't seen him for almost 30 years, almost 30 years..
So he just moved with his family from Hong Kong,.
moved to Toronto..
He got a job there..
And you probably know some of the politics of today, right?.
Canada has this lifeboat programs,.

$^{281}$you know, for people from Hong Kong..
They kind of like a refugee status.
because of the political change or political hardship in Hong Kong..
So this friend of mine is part of that wave of immigrants,.
wave of people that left Hong Kong..
They went to Toronto a couple months ago.
and tried to settle down, you know..
Now, so often, we visitors, we are travelers,.
you know, especially we are holding U.S. dollars..
We go in Toronto, "Wow, everything is so cheap..
Everything is so cheap.".
Every meal is like, "Dude, I can eat 10 meals here..
I'm using U.S. dollars because U.S. dollars is so strong, right?".
And we go buy stuff, right?.
We go to the shopping mall or the outlet..
"Wow, everything is so cheap..
Oh, man, this is such a great trip," right?.
But you talk to the locals, you hear their stories..
It's the other side of the coin..
They had difficulties..
With all the inflation after COVID,.
with all the financial hardship they are going through, you know,.
I think as we spend good time in Toronto,.
spend good time as a traveler, as a tourist,.
I think we need to open our ears at least to know some stories, right?.
To listen to their stories..
Can you imagine someone who has been his whole life in Hong Kong.
and he and his family has to be re-rooted in a foreign land, right?.
It's a tough life, you know?.
You use U.S. dollars in Toronto, it's a good deal, right?.
But for those who earn money, try to settle down there,.
it's a totally different story, right?.
It's kind of like a story of Ruth, right?.
You know, I was talking to a couple who just spent a couple weeks in Hong Kong.
after the service, the previous service..
They were talking to me..
They spent two weeks in Hong Kong..
They said the whole society has changed..
It's not the one they used to know..
They told me they stayed in an area near a university in Hong Kong.

$^{321}$and they go to the subway every day, right?.
And the police officers were picking on all the college students,.
all the young people, for no reason..
They said everything has changed because of politics, right?.
Think about people like this..
They have to leave..
It's not because it's a better life out there..
They got to go, right, for the kids or for the family or whatever reason..
They just can't stay there..
I was talking to someone in the morning..
He was updating me about what's going on in China..
You cannot imagine living in a place like this, you know,.
where COVID is almost over, but the government is not willing to step down..
It's not willing to change the policy, right?.
I just Googled, you know, BBC News, China COVID,.
protesters, you know, from Beijing to Chengdu, Chengdu,.
protesters openly urged Xi Jinping to resign over the COVID policy..
All over the place, you know, people are protesting..
You cannot imagine living in a place like this, and you have to go..
You go to somewhere else, right?.
So I met some friends in Toronto..
They have the same stories, right?.
Yeah, we have a happy time, but, you know, they have a tough life, you know..
That's why in the Torah, when it mentions Shavuot,.
the time of harvest, the time when you get your paychecks, right,.
it continues to talk about how we should live with all we have..
So we read part of the passage from chapter 23,.
the book of Leviticus..
Let's continue..
Verse 22..
Now it's talking about Shavuot, right, Pentecost..
You have to bring your grain offering to the Lord..
And then verse 22, it says, "When you harvest your crops of your land,.
do not harvest the grain along the edges of your fields,.
and do not pick up what the harvesters drop..
Live it for the poor and the foreigners, the aliens living among you..
Be kind to them. I am the Lord, Yahweh, your God.".
Now, listen to this law, this Torah..
It is essentially relational, right?.
It flows out of your relationship with God.

$^{361}$because God is giving you good things..
Now you can harvest, you can get your paychecks, right?.
And then from your blustered life,.
you flow out the kindness of God to other people who are needy around you..
Who is the first reader of this Torah, of this regulation?.
It's not the poor..
It's not the needy..
It's written for those who have a lot, right?.
Who has land, who has harvest, who has workers working in the field, right?.
23, 22, when you harvest your land, right?.
So you own the land, right? You have a farm..
You have harvest, you got your money, you got your paycheck, right?.
Don't just think about, "Okay, I want to buy this," or "I want to buy that.".
Just fulfill my needs and fulfill my wants, right?.
After Thanksgiving, after Black Friday, and then we think about,.
"Okay, Cyber Monday, what should I buy? What do I want?".
This regulation is telling you, as the people of God,.
we need to think about those around us who are in need..
Leave them..
The edges of your field, you know, when you cannot harvest everything,.
just leave it there..
Share with some other people, right?.
Leave it there..
When you drop something, leave it there.
so the people following you, behind you,.
just like Ruth, the girl, the moorbyte,.
leave it for the poor and the foreigners living among you..
Be kind to them..
Always think of when you harvest, when you get your paycheck,.
how can I be generous to those in need?.
Everything has to do with relationship..
Torah, the law, is imitating God..
God is generous..
Look at the universe around us..
God is generous..
He made a huge universe for us to explore, for us to enjoy..
This is the life of God..
That's why the life of God's people.
needs to be generous and gracious.
to imitate the character of God..

$^{401}$The last verse I want to share with you is Acts 2.1..
Now when we think about Pentecost,.
when we think about the festival weeks, the harvest,.
we always think of Pentecost,.
what happened was the Holy Spirit.
descending into the lives of people,.
and then there were signs and wonders, right?.
People speaking in tongues and all this..
Miracles happening in the early church..
Now listen to what it says, 2.1..
2.1..
When the day of Pentecost arrived,.
first phrase, first sentence..
A better translation is actually.
when the day of Pentecost was being fulfilled..
It's being fulfilled..
What is being fulfilled?.
It's the spirit, the essence of the festival,.
the harvest, and being generous..
Now we often, as I said,.
we often pay attention to the signs and wonders, right?.
But if you trace the story of Pentecost.
up to chapter 4, 4.34,.
it describes the life of the early church,.
the early Christians, the early Jewish Christians..
In the day of Pentecost, when it is fulfilled,.
when the Holy Spirit came upon His people,.
and they shared their things with one another, right?.
Those who were well-off,.
they shared their own properties..
They shared their own goods.
with those who are in need around them..
4.34, there were no needy people among them..
That is the miracle..
That is the true sign and wonder..
It's not speaking in tongue..
It's not just enjoying,.
okay, we have a high worship time, high life, right?.
It's the fact that they are being generous to one another, right?.
If you check the phrase, "no needy people among them,".

$^{441}$this phrase is in the book of Torah,.
the book of Deuteronomy..
When the people of God is really applying,.
you know, the Torah of God,.
being generous, living the life of God together,.
there's no needy people among them..
That's what the early Christian church experienced..
No needy people among them..
So I invite you to go back to the Bible.
as we check the--as we study the festival of Shavuot, Pentecost..
Think about how God has blessed you so much, abundantly,.
and how you can bless other people,.
and you can be generous to other people, right?.
Look around the world, you know, pray for those who are in need..
Pray for those who are under oppression..
You know, I hope you understand I'm not talking about early politics..
I'm talking about the kingdom of God,.
full of justice, compassion, and righteousness..
That's what God wills, right?.
That people will respect one another, to love one another,.
under the Torah of God, to participate in the life of God..
I'm talking about that..
I'm talking about heavenly politics,.
justice, and compassion, right?.
Think about, you know, past week,.
we spent a great time, Thanksgiving,.
we ate a lot, we probably threw out a lot of food..
Just learned plenty..
How we can bless other people and exercise generosity..
Let's operate together..
Father, we give you thanks because you have blessed us abundantly..
Thank you so much for your blessing..
The life that we have, the stuff that we have, is more than enough..
Thank you for being our good God and faithful God..
And most of all, thank you for giving your son, Jesus Christ,.
to give us eternal life, give us forgiveness of sin..
And boy, we pray for not only being thankful,.
pray that we can be kind and generous and compassionate toward other people..
Those who are heartbroken, those who live in places of oppression,.
those who are struggling with settling down in a foreign land,.

$^{481}$the needy people around us..
Open our eyes, God, that we don't just see and so focus on our needs and our wants..
Because you have indeed blessed us so much..
Open our eyes so that we can see the needs around us..
No matter if it's churches, it's ministries, it's other people that we know,.
help us to be generous..
Help us to give..
Help us to bless other people..
And we pray that we will be a people living under your Torah.
to know your will and apply your will in our lives..
And we pray that you will send us off this week.
so that we will go to our corners, go to our own places,.
our situation in our lives, that we can be a blessing to other people..
And what a privilege we can participate in the life of God..
What a privilege we can share your grace, share your blessing to other people..
May our lives be an instrument of your grace..
In Christ's name we pray. Amen..
Well, God bless you, and we'll see you next week,.
and continue to worship and learn about God..
All right, God bless..
\newpage



\section{}
\label{sec:LC2O0QPKB1U}
\textbf{佳節倍思恩(五): 五旬節  - 陳崇基牧師}
\newline
\newline
連結: \href{https://youtube.com/watch?v=LC2O0QPKB1U}{\texttt{https://youtube.com/watch?v=LC2O0QPKB1U}} ~~~~ 語音日期: 2022-11-28
\newline
\newline
\hyperref[sec:xHuMbJI2NdE]{\small{< < < PREV SERMON < < <}}
~
\hyperref[sec:index_chronic]{\small{[返順時目]}}
~
\hyperref[sec:index_scriptual]{\small{[返順卷目]}}
~
\hyperref[sec:SrXfS_omV5U]{\small{> > > NEXT SERMON > > >}}
\newline
\newline
$^{1}$現在將時間交給他.
大家好.
昨晚我剛從唐人街回來.
大概十點,十一點.
去了幾天.
很好玩.
好玩的每個人都這樣問我.
其實我主要目的是開車.
和給錢.
所以每次去商場.
我都坐在那裡.
我預備說一張.
然後時間夠了.
他們叫我要你的卡.
我就乖乖地去給錢.
很好玩吧.
哈哈.
不知道大家有什麼慶祝.
大家應該吃得很飽.
還有大家應該和家人或朋友聚首.
這是很難得的時間.
Thanksgiving.
是一些節日.
停一停低.
讓我們去明白.
想一想上帝在我們身上做過的恩典.
我們講開這個系列.
聖經裡面都有一些.
一年三百六十五日.
有一些節日.
一些標籤.
在時間上的標籤.
讓上帝的子民想起上帝的恩典.
這次我們去講這個節日.
是三大節日的其中一個.
之前我們講過Tabernacle.
Sukkot.
也講過什麼.
也講過Passover.
這次我們講的是中間的收割節.

$^{41}$或者五旬節.
這個沒有什麼人講的.
因為好像和一些大上帝的作為.
沒有什麼關係.
但是這個節日.
我在舊約裡面提到的時候.
就和上帝的律法有關係.
當然我們提到五旬節的時候.
一想就想起.
就是這個銷行傳第二章.
聖靈光臨五旬節的時候.
其實五旬節就是收割節.
一個農耕的節日.
我們叫它做Shavuot.
或者Feast.
Festival of the Weeks.
Shavuot其實就是星期眾數的意思.
Shavuot.
七個眾數的意思.
所以通常就是在Passover之後.
再數七個星期.
七個主日.
七七四十九日之後.
第五十日.
就是這個所謂Shavuot的時候.
就是這個所謂收割節的時候.
適逢剛剛小麥在巴勒斯坦地收割的這段時間.
這個Shavuot和其他的Festival有些不同.
你可以想一下.
其實我們每一年有不同的節日.
包括 Thanksgiving 或者接下來的 Christmas.
這些節日都有它特別的地方.
這個收割節和其他的 Passover.
Tabernacle 這些節日是有些不同的.
你想想有什麼不同.
我舉個例子.
有一些舊約聖經裡面提過的節日.
都是和上帝在歷史裡面.
曾經做過的事情.
一些神蹟.

$^{81}$一些奇事是有關的.
譬如駐彭捷.
就是去紀念神帶領以色列人經過四十年曠野.
他們住在彭裡面.
然後一路一路帶領他們進入迦南.
這個就是 Shavuot 或者 Tabernacle.
另外一個節日比較少提到的.
就是普爾節.
Purim.
Purim大概是我們西歷的一二月的時間.
Purim其實是以色列的 Halloween.
他們就扮鬼扮馬等等.
Purim 在哪裡提過.
就是 Esther.
這個以斯帖記.
提過上帝怎樣拯救自己的子民.
以致他們不會被殲滅.
這個就是 Purim.
所以現在的以色列人.
猶太人就會去會堂那裡.
大概一二月的時候.
Purim 的時候.
他們就會特意去會堂那裡.
然後他們就會讀.
The Book of Esther.
以斯帖記.
每逢讀到他們的仇敵哈曼這個字的時候.
他們一邊讀一邊讀.
由第一章第一節一直讀下去的時候.
讀到哈曼的時候.
他們就會哼.
因為是他們的仇敵.
所以他們就有這樣的節日.
然後他們就扮有不同的服飾等等.
這個就和上帝在他們的子民身上做過的事情有關.
另外就是修典節.
修典節在聖經裡面沒有提過的.
Hanukkah.
Hanukkah 大概是十二月的時候.
西曆十二月的時候.

$^{121}$大概是聖誕節的時候.
如果你去一些猶太人的社區.
你就會看到他們會.
所謂光明節或者修典節.
也是和上帝在上帝的子民身上做過的歷史的工作有關.
當時大概是耶穌基督來之前大概200年左右.
160多年BC的時候.
他們去到反擊猶太人.
他們反擊欺壓他們的希臘帝國的軍隊.
然後他們成功地爭取到獨立.
一個本身是underdog.
他們是被人欺負的.
underdog然後反擊他們贏了.
就有一個很大的慶典.
所以就有一個所謂修典節.
拿回他們的聖典.
有一點像日本打贏德國.
怎麼可以打贏.
最後入一球.
今天大家很高興.
世界盃.
underdog贏了的時候.
就有這個.
Hanukkah.
然後我們提過了.
就是講到Exodus出埃及.
但是Sherford Ho和上帝在歷史裡面的作為.
是沒有什麼關係的.
我們看一節聖經.
利美的23章第15節.
這裡說你們從安息日.
之前他就講了passover.
安息日的次日.
次日就數幾天.
七個的安息日.
第七個安息日的次日.
就共計50天.
所以為什麼叫Pentecost.
希臘文就叫50.
你們要將新的掃祭.

$^{161}$當時剛剛收割小麥的時候.
就獻給耶和華.
在你的住處要拿你的獻祭.
然後是初熟之物.
就獻給耶和華.
first food.
因為第一次的收成.
整年都沒有第一次收成.
你看到這裡沒有什麼miracles.
沒有什麼神蹟奇事.
只不過每年去到出糧的時候.
就是收割的時候.
有harvest.
他們就將這些掃祭.
或者grain offering.
將這些麥子來獻給耶和華.
是時候他們再去奉獻.
但是到了兩日中間的時候.
就是救藥完了.
新約還沒到的時候.
中間這段時間.
發生了很多事.
其中就是修奠節.
他們去反擊.
得到政治的獨立等等.
在這段時間.
猶太人就將這個收割節.
就跟上帝的律法.
上帝賜下的律法.
就拉上關係.
其實差不多.
因為他們離開埃及.
在修割的時候離開埃及.
整整兩個月之後.
他們就到了西乃山.
兩個月.
大概八個星期左右.
整整兩個月.
他們就去到西乃山.
然後上帝就頒佈他的律法.

$^{201}$給上帝的子民.
這個修割節大概是.
修割之後五十天.
所以大概這段時間.
他們就將兩件事merge.
然後就會認為.
這個時間他們紀念到什麼.
不是單單收割.
就是上帝賜律法給他的子民.
來引導他的子民.
來怎樣去生活.
所以今天也是.
在那些猶太人的會堂.
他們就會紀念到上帝.
將Torah.
來給上帝的子民.
一件很重要的事.
但沒有什麼.
崩山碎石的事情發生.
只是上帝將他的law.
他的律法.
Torah.
來給上帝的子民.
為什麼這麼重要.
我們需要講一講.
什麼叫Torah.
什麼叫law.
我們現代人的翻譯就很笨.
我們經常講law的時候.
就是講什麼.
是一些跟關係沒有什麼拉攏的東西.
譬如你開車.
你沒有停stop.
警察抄你牌.
你不會說得罪誰.
你只會得罪你的錢包.
或者你開車.
speed limit寫著65mph.
但你不會說有什麼關係的意思在那裡.
只不過是一個law.

$^{241}$一個regulation.
一個規則.
你不要過那個speed limit.
或者你過了speed limit就犯規.
是一些impersonal的東西.
就是一些跟關係沒有什麼拉攏的東西.
這個就是我們現代人的印象.
什麼叫律法.
我們中文譯Torah就譯做律法.
其實有很淪盡的.
有些不熟悉聖經的人會覺得.
法律.
其實不是真的法律.
Torah是什麼呢.
當這個修國節的時候.
我們紀念到上帝又賜下Torah.
賜下他的律例.
就不是我們現代人這樣想的.
純粹是一個rule.
規則.
應該是跟relationship有關的.
我記得兩個星期前我去上課的時候.
我們剛剛就去到.
有一個討論的主題.
我們就說到底Torah應該是怎麼翻譯.
怎麼去了解什麼叫做law of Yahweh.
什麼叫做上帝所賜的律法.
我覺得有一個提議挺好的.
就是他用一個比喻來說.
什麼叫做Torah.
你問一下你自己什麼叫做Torah.
上帝賜他的律法賜他的華語給我們.
是什麼意思.
其中一個比喻挺好的.
其中一個同學說.
如果你告訴一個小孩.
我們中間有很多很小的小孩.
Don't touch the stove.
一歲兩歲的時候.
Don't touch the stove.

$^{281}$或者你家裡的貓狗都可以叫他們.
Don't touch the stove.
不要搞亂我的東西.
什麼意思.
對於一兩歲的小孩來說.
Don't touch the stove.
但當小孩越來越大的時候.
如果他到十歲.
十五歲.
甚至到二十歲的時候.
他都會說這個就是Torah.
就是Don't touch the stove.
如果他還是這樣想的時候.
他就不是很明白.
那個說話Don't touch the stove.
那個所謂的律例的意思.
其實Don't touch the stove是什麼意思.
Use the stove properly.
Right.
Use the stove properly.
到他慢慢長大的時候.
他就明白原來這句說話.
開頭父母跟我說的.
不要碰那個爐頭.
其實意思是什麼.
是很好用那個爐頭.
到他十五十七歲的時候.
二十歲開始懂得做東西.
開始懂得去用的時候.
他就應該去到Use it properly.
煮餐好的.
這個就是Don't touch the stove的意思.
你明白我的意思嗎.
背後的精神.
慢慢長大的時候.
他就應該要get到.
如果他get不到的時候.
他就有些問題.
Don't touch the stove.
其實你看看其他的律法也好.

$^{321}$都是一樣的.
是要你去領略背後的意思.
那個精神.
你才知道到底律法是一回什麼事.
當五旬節來到的時候.
猶太人的傳統.
就會去紀念上帝頒佈律法.
其中一個最重要的律法是什麼.
十誡.
Ten commandments.
In fact, Ten commandments就是第一套.
上帝在西太山給他們的律法.
你舉例.
其中一些的律法.
不可謀殺或者不可殺人.
你不是說.
總之我不殺人就ok.
其實不是.
你看看其他的律法.
去到典章.
去unpack不可殺人.
Do not murder, thou shalt not murder.
的意思的時候.
其實不是這樣.
其實是保護所有生命.
保護所有的life.
因為上帝的子民應該是.
應該是推廣生命.
因為生命是神聖.
Do not commit adultery.
不要犯奸淫.
總之我不做那些事.
那我就叫做守了.
就叫做按照律法去生活.
不是.
你想想一個老公回到家.
每晚回到家.
我沒有犯奸淫.
所以我很愛你.
不會這樣算的.

$^{361}$應該是什麼.
還有他正面的精神.
就是怎樣去正面的去推廣.
去nurture.
去培養那個關係.
家庭的關係.
婚姻的關係等等.
Thou shalt not steal.
不要偷東西.
你不是說.
總之我不偷東西就ok.
我不拿別人的東西.
律法裡面.
更加重要的是什麼.
就是去慷慨的.
分享你的東西.
這個叫做不可偷渡.
不可偷竊.
不是說我不拿別人的東西就算了.
背後的精神就是.
我有的東西.
我應該不要在那裡算死.
然後跟別人分享.
其實最主要的精神是generosity.
你想想這些律法.
寶貴的地方就在這裡.
不要像我們小時候聽到.
Don't touch stove.
然後我們一輩子都不去touch stove.
其實精神在哪裡.
就是要我們去領略那個意識.
所以為什麼猶太人的傳統.
就很著重這個shuffle of the soul.
就是上帝給我們律法.
就是給我們神的子民.
按照上帝的美意.
去生活出來.
上帝要我們去保護生命.
推廣生命.
上帝要我們去推廣關係.

$^{401}$好好愛身邊的人.
上帝要我們be generous.
我們有的東西可以跟別人分享.
所以為什麼去到shuffle of the soul的時候.
大概是我們西曆四月五月的時候.
就是Pentecost的時候.
猶太人就會去到會堂.
他們特別會讀兩段聖經.
第一段聖經是什麼呢.
就是剛才我們在這裡看了.
就是十誡.
他們再想清楚到底十誡是什麼意思.
所以為什麼有時候他們會說.
shuffle of the soul其實就是.
不只是修國節.
就是啟示的節日.
上帝啟示他自己的美意.
啟示他自己的性情.
去告訴我們.
一個律法背後就是上帝自己的.
美善的本質.
上帝就是這麼慷慨.
你看上帝創造的宇宙.
是多麼無憾多麼大.
給你去享受.
上帝可以很吝嗇.
用到就夠了.
但他給我們的風景.
給我們這個宇宙去欣賞.
上帝就是一個這麼慷慨的上帝.
所以上帝說我是慷慨的.
你們就應該要去慷慨.
如果你這樣看律法的時候.
就是一些relational的東西.
是關係的東西.
你和上帝之間有的關係.
就應該在這個關係裡面流露出來.
那些律法其實是教導我們.
去引導我們.
怎樣去活出上帝自己的生命.

$^{441}$第一段十誡.
就是猶太人在會堂裡面.
五分之五的時候他們會讀.
第二段是什麼呢.
第二段就是路得記.
跟路得記有什麼關係.
你記得路得記的故事嗎.
路得記的故事.
是一個摩訶的女子.
一個外邦的女子.
她去投向以色列人.
投向耶和華.
然後她就搬了去迦南地.
然後裡面發生了很多不同的細節.
就跟修國節有關.
為什麼呢.
因為路德第二章.
就在修國節發生.
因為摩訶的女子路德.
對那沃米說.
她的家婆.
跟她說.
等我出去.
到拾取墨水.
什麼時候拾取墨水.
就是修國節的時候.
工人去收割莊稼的時候.
他們作為一個窮人.
路德和劉沃米.
就出去田間拾取墨水.
她就說.
看看我可以跟到誰.
誰可以蒙恩.
劉沃米就跟她說.
女兒你去吧.
這個年輕女孩就出去了.
就去到田間.
在修國的人的身後.
來拾取墨水.
因為她是窮人.

$^{481}$所以她拿不到正式的莊稼.
但剩下的東西她就可以撿起來.
撿起零碎.
以至回到家裡就有東西吃.
恰巧就遇見誰.
波瓦斯.
年輕女孩就遇見.
她本族的一個族人.
波瓦斯.
就去到恰巧地.
就去到她的田裡.
大家知道那個細節.
波瓦斯就很恩愛她.
然後就娶了她.
以至她的家族能夠留名.
發生的時間就是收割的時候.
拾取墨水的時候.
但最重要的.
不單單是一個拾取莊稼.
或者收割莊稼.
收割節的一個故事.
奴隸記是什麼故事呢.
其中一個最重要的主題就是.
慷慨.
Generosity.
就是波瓦斯對一個愛邦女子.
對一些窮人.
他就實行慷慨這個道理.
In fact 律法書裡面.
就是告訴那些有剩餘的人.
有收入的人.
有多的人.
他們需要做些什麼.
有多的不是說.
啊 有多的.
那我明天這個Cyber Monday.
我又多買些東西給自己.
我們就是這樣的.
我已經買了一個星期的東西了.
過去一星期.

$^{521}$我都快瘋了.
我一直預備這個講章的時候.
我就一直告訴自己.
不需要憂慮.
明天就Cyber Monday.
我有多錢.
不如我再買多些東西給自己.
我們滿足的就是我們想要的東西.
不是我們需要的.
我們想要的.
奴隸記.
剛剛在這個Cyber All發生的時候.
就是要慷慨.
去幫助身邊有需要的人.
剛才我說了多不多.
過去一個星期.
每個人都這樣說.
我們拿著美金去花.
我們覺得好像發達一樣.
樣樣東西好像不用錢一樣.
特別我的小朋友他們覺得樣樣東西都很便宜.
總之我就不教他們怎樣換算.
因為美金強.
上到去吃東西.
哇 這麼便宜的東西.
不吃不行.
樣樣東西都要吃.
這樣又要吃 那樣又要吃.
吃到你吐都要吃.
但我有另一個領會.
我有另一個領會就是.
因為我探訪了一些朋友.
有些朋友在香港去到多倫多.
在那裡定居.
和我們以前在香港來的那班.
來到美國也好 來到北美也好.
那些移民的心態是不同的.
我開車去探望我一個老朋友.
他們是那些類似舊生艇.
加拿大有些舊生艇的program.

$^{561}$來到去叫做舊仔.
叫做舊仔的香港人.
我們這些來北美這麼久.
已經落地生根的.
有時候我們有一點離地的.
我們不太清楚他們的生活狀況.
我們假設他們好像我們以前那樣.
其實你想想我們以前也是捱生捱死.
你來到一個異地.
又要被你自己本族的人欺負.
然後你做生做死.
我還記得我來到美國的時候.
我只是洗碗碗的.
洗到手都難.
凍水洗.
手機也沒有.
你就是一天十幾個小時在這裡洗碗.
你才可以有錢讀書.
其實很苦.
他們這班來到多倫多的.
他們很多故事都是這樣的.
我們就開心 我們就在這裡吃東西.
有些人要做幾份工作.
為了能夠拿到居留.
為了能夠有一個好一點的生活.
聽一下他們的故事.
其實有苦就是自己知道的.
特別現在他們有很多是live bowling.
你看看世界裡面有多少地方.
生活的人是多麼苦.
我們就拿著美金走去吃一樣的.
我們覺得很開心.
我不是叫你不吃 吃.
儘管吃.
但你聽一下他們的故事.
聽一下他們的故事.
中國現在搞成這樣.
剛剛早上我們大紀元節目聊開告訴我.
各大城市因為動態清零.
不知道叫什麼.

$^{601}$其實是鎮壓.
其實是控制人.
你想一下多麼苦.
多麼苦.
我們去旅遊.
我們要吃好的東西.
有時候是不是太離地.
聽一下他們的故事.
我在多倫多聽了好幾個故事.
都是我熟悉的人.
你就知道他們的生活是多麼艱難.
所以為什麼在聖經裡面.
看到路得記 修國節的時候.
都有提過一些跟窮人有關的東西.
有需要的人 有關的教導.
剛才我們看了利未記.
你再看利未記22章.
講到修國節的時候.
你們在自己的地修國莊稼的時候.
不可割盡田的角落.
也不可拾取莊稼所吊下來的.
要把他們留給窮人和寄居的.
因為我時也會說你的 上帝.
這句話是寫給誰的.
這段律法.
這個律例是寫給什麼人的.
是寫給有田的人.
不是寫給窮人的.
你不要以為是寫給有需要的人.
是寫給自己的田地.
自己有收拾莊稼.
有工人去幫你的那幫人.
他就說你們收取莊稼的時候.
不要收到田角也拿了.
你賺幾個開不要賺到盡.
田角的東西應該留給窮人.
寄居的人.
在你中間異地異鄉的人.
慷慨的跟他們分享.
這就是重點.

$^{641}$重點不是說我不拿別人的東西.
而是你剩下的東西.
應該跟別人分享.
你有剩餘的東西.
不要每次都想.
我又想要那樣.
我想想身邊的人有需要.
當然你可以做什麼.
你吃完東西可以做什麼.
你可不可以幫到那些人.
最終來說可能真的不行.
但我覺得在心態上我們可以聽聽他們的故事.
聽聽他們有什麼辛酸.
去關心他們.
你可以的話可以跟他們分享.
這就是Torah的精神.
我們可以去模仿一下上帝的生命.
上帝是一個這樣的上帝.
慷慨的恩待我們.
感恩節想起.
就是上帝給我們的福祿.
有好的東西.
每樣東西都供應給我們.
可能有很多我們還是欠缺的.
但我告訴你.
十該其九可能都是想要.
I want.
不是need.
不是我們需要.
上帝是一個這樣恩待我們的神.
他現在的律法.
就是告訴我們我們可以行出他的美意.
身邊的人.
或者一些需要的事工.
我們慷慨的去賜予.
這才是上帝的律法的精神.
最後一段經文.
使徒函傳第二章.
就是聖靈降臨的時候.
五旬節的那一日到了.

$^{681}$其實應該是更好的翻譯.
五旬節的日子.
實現了.
fulfill.
我們想起五旬節的時候.
就想起聖靈降臨.
然後他們就講方言.
然後有很多人信主等等.
這些神蹟奇事.
推下去使徒函傳.
講到五旬節之後發生的事.
第四章講到初期教會.
他們所謂凡物共融.
大家慷慨的去分享.
去祝祭身邊的人.
講什麼?.
四章三十四節.
他們中間會有一個缺陷.
查了整本律法書.
舊約.
這句話在申命記出現過.
中間沒有缺陷沒有窮的人.
No one is needy among them.
在初期教會.
修國節,五旬節來到的時候.
聖靈的工作到了.
他們就明白.
這是上帝賜予律法的時候.
我們要去實行慷慨之道.
使我們中間沒有缺陷.
這才是奇事.
這才是神蹟.
神蹟奇事不只是過紅海.
不是過藥彈河.
不是一些崩山碎石的事.
而是上帝賜予律法給他的子民.
聖靈在我們中間來行事的時候.
我們去行出上帝的美意.
有上帝的生命.
這才是奇事.

$^{721}$這才是最驚人的事.
多大的完美我們能夠.
仿效上帝的生命.
我們可以唱一首回應詩歌嗎.
剛才第二首那首.
供應者.
麻煩你上來.
唱這首詩歌的時候.
讓我們做一個好的回應.
這樣就過了.
需要的.
想要的.
都有了.
你想一想.
上帝的聖靈.
來到我們生命當中.
去行出.
我們能夠有這樣的特權.
可以行出上帝的生命.
數算一下上帝的恩典.
不只是我們.
好好地過生活就罷.
行出上帝的自己的生命.
成為上帝的子民.
想一想我們自己感恩.
自己落袋的東西.
是好.
但也想一想我們怎樣能夠.
去奉獻給人.
怎樣能夠去實行.
慷慨之道.
實行上帝的話語.
上帝的律法.
這才是神跡奇事.
聖靈的工作.
我們唱這首歌.
《天空降下雨水》.
作詞:陳家麗 作曲:陳家麗.
天空降下雨水.
我在半空.

$^{761}$小花汲取寂寞.
綻放色彩.
當失約與困苦.
你仍許沒變改.
令我深深感到痛.
曾帶我心灌自開.
盡量你才回答.
此刻我.
更加索性一次.
令我心喝彩.
祝福我在這生.
信心不會變愛.
懷著謙卑心願讚頌你愛.
神使我這生當然真.
作風滿滿傾城.
更擁抱我渡過寒夜.
神使我這生光年正.
任光照我方圓.
許多就做條著你可改成.
神願替萬物歌唱光年.
我們用心低頭祈禱.
因為你是那位公認滿滿的祝福我們的神.
主叫我們恩著你給我們的恩典.
恩著你給我們的福氣.
我們不單止自己去到感恩讚美你.
也讓我們生命能夠行出上帝自己的律法.
讓你自己的道在我們生命當中能夠實現出來.
使我們教會.
使我們自己的生命能夠成為其他有需要的人的祝福.
教導我們慷慨之道.
教導我們分享奉獻.
教導我們不單止去到.
只懂得去計算我們自己有多少.
想要多少.
叫我們學習我們自己的生命.
叫我們的生命能夠去祝福更多的人.
藉著我們使更多人得到你自己的恩典.
祈禱奉著名感恩.
好,以上就是我的祝福,希望下個星期能夠跟大家一起來到市中心.
\newpage



\section{}
\label{sec:SrXfS_omV5U}
\textbf{佳節倍思恩(四): 逾越節  - 陳崇基牧師}
\newline
\newline
連結: \href{https://youtube.com/watch?v=SrXfS_omV5U}{\texttt{https://youtube.com/watch?v=SrXfS\_omV5U}} ~~~~ 語音日期: 2022-12-06
\newline
\newline
\hyperref[sec:LC2O0QPKB1U]{\small{< < < PREV SERMON < < <}}
~
\hyperref[sec:index_chronic]{\small{[返順時目]}}
~
\hyperref[sec:index_scriptual]{\small{[返順卷目]}}
~
\hyperref[sec:kSJ3Fo7YGUs]{\small{> > > NEXT SERMON > > >}}
\newline
\newline
$^{1}$好 今日我地繼續睇呢個series.
就係睇下聖經ge 節期.
每逢講到節期ge 時候.
好多時我地都會諗起同一D 故事.
或者歷史ge 事蹟.
或者食物有關.
譬如我比你估估.
或者你諗一諗.
到新年ge 時候我地華人食咩.
年糕係咪 冇錯 年糕.
或者.
蠔豉發菜gwo D 咁ge ye .
我細個ge 時候最怕食gwo D .
食一日ok.
跟住第二日又食 第三日又食.
食到初七 咁我就覺得好辛苦.
譬如仲有咩節日呢.
同呢D 食物有關.
冬至.
湯圓係咪.
喂 湯圓.
仲有呢 糯米.
仲有呢 打邊爐係咪.
打邊爐.
仲有咩節呢.
感恩節呢.
係打邊爐係咪.
端午節食咩.
就唔係打邊爐.
我地淨係識打邊爐.
咁講到聖經ge 節期.
其實都有好多同食物有關.
譬如今次我地講.
預節.
預節同乜ye 有關.
食物.
羊我地諗到係咪.
我地第一樣ye 就會諗到係羊.
事實上聖經裡面.
第一次守預節ge 時候.

$^{41}$就係同羊有關.
係出埃及記你記得.
當耶和華要降十災.
最後gwo 災.
就係擊殺埃及長子ge 一災.
係上帝ge 審判.
去到呢一個ge .
即係所謂強國當時.
然之後呢.
以色列人住係埃及裡面.
佢地係需要被拯救出黎ge .
於是乎上帝就同佢地.
設立呢一個ge 預節.
十二章第三節.
gwo 度講到.
預節ge 安排.
即係你地要去做D 乜ye .
然之後呢.
食完呢一餐飯你就要走.
十二章第三節.
佢話你要告訴以色列.
傳會眾.
本月初十.
各人要按照家庭.
娶羔羊.
一家娶一隻.
然之後佢就講到呢個羔羊.
要去被烤.
然之後呢要係屋企裡面食.
然之後呢就要將.
宰殺jor ge 羔羊.
羔羊ge 血.
就要塗係邊度.
門尾上面.
以至上帝ge 審判.
臨到呢個國度ge 時候.
上帝就會越過.
有門尾上面.
有血ge .
即係家庭.

$^{81}$係一個拯救ge 故事.
佢話你要咁樣去吃.
呢個羔羊然之後呢.
要腰間束帶.
腳.
要穿著鞋呢個好有趣.
因為猶太人.
起碼以色列人.
入到屋ge 時候佢就會脫下鞋.
係門口.
然之後.
要上路ge 時候.
佢地就會.
戴皮帶係咪.
出門口ge 時候.
好多時都會戴皮帶然之後穿鞋係咪.
準備去走ge .
食呢隻羊羔.
手裡面仲要拿住杖.
特別係講到gwo D .
walking stick係咪.
即係你要行路ge 時候lor 住gwo 張.
快快地.
然後仲交代jor .
好多ge 細節我地冇時間講.
食無烤餅.
點解係無烤.
如果整個麵包ge 話你知道.
你等gwo D 麵粉發酵.
要等好耐ge 係咪.
要等好幾個鐘頭.
然之後你先至能夠.
烘到gwo D 麵包.
因為趕住走就唔可以.
所以呢佢地係呢一個.
ge 節日就要.
去食無烤餅.
就紀念返上帝拯救佢地.
離開埃及.
拯救佢地以至佢地.

$^{121}$唔能夠唔會呢係被審判.
這是耶和華的預言.
Passover.
越過.
越過佢地ge 門.
或者去到.
保護佢地.
每年西曆ge 時候.
三四月呢就係呢個.
猶太人ge .
Passoverge 時候就淨係呢個.
所謂即係春分.
ge 時候之後ge .
第一個ge 新月.
所以亦都係.
一個季節.
與季節之間過渡ge .
一段時間.
當時特別講到如果你講.
耶穌ge 時候.
當時有好多散居.
流散係世界各地ge 猶太人.
佢地就會返返去.
耶路撒冷.
然後就會守節.
當時ge 人會好擠迫好烈鬧.
你可以想像呢個係一個最大.
最大ge 節日.
好似我地華人ge 新年.
華人ge .
一D 重要ge 慶典.
或者好似.
美國呢度Thanksgiving.
ge 時候最忙ge 地方.
係咩地方.
機場餐館.
華人餐館.
好多人走來走去.
猶太人就係咁樣.
猶太人就會去返.

$^{161}$以前冇得搭飛機.
佢地就會趕路.
去到耶路撒冷.
所以我地諗起.
預節ge 時候就諗起高揚.
頭先歌都有唱.
上帝帶領我地.
預約經過.
Passover.
類似咁樣ge 意思.
但係你要知道係兩約中間ge 時候.
即係舊約完jor .
到新約未開始ge 兩三百年.
當時ge 猶太人佢地諗起.
預節ge 時候佢地就唔係.
諗起高揚.
佢地諗起成個package.
成舊ye 成個上帝ge 拯救.
會係有D 乜ye 呢.
你諗下有D 乜ye 除jor 隻羊之外.
佢諗起一連串ge ye .
有摩西出現.
因為上帝呼召摩西.
去帶領佢地好似一個牧羊人.
帶領佢地離開埃及.
然之後當然就係我地有所謂.
sacrificial lamb 即係呢個.
預節ge 高揚.
預節ge 高揚ge 血.
無烤餅.
仲有諗起乜ye .
就係水ge 神蹟.
過紅海.
過完紅海.
成個package.
就睇埋一個.
曠野.
咁多年上帝帶領佢地.
係行過ge .
係曠野又比佢地嗎哪.

$^{201}$等佢地每一日.
能夠有ye 食.
咩叫嗎哪.
manner.
大概可能聖經ge description.
gwo D ge 描述.
係一D 圓圓地ge 好似餅咁ge ye .
其實我地都唔係好知.
因為呢個故事.
話比我地聽.
嗎哪個名係點黎.
上帝降jor 一D 咁樣ge .
好似麵包咁樣ge ye 比佢地.
然之後佢地就唔知係D 乜ye 黎.
佢地就話Mazet Mazet Mazet.
Mazet 希伯來文.
What is this.
乜料.
乜ye 黎架咁樣.
所以就黎jor 呢個名叫做嗎哪.
所以其實我地唔知乜料黎.
但總之佢地食jor 咁多年.
上帝就.
冇間斷過ge .
黎到去供應佢地.
我講兩日中間.
耶穌之前黎到ge 時候.
佢地就諗起呢一連串ge .
所謂預月節會發生ge ye .
或者pass overge 時候.
佢地會諗起ge ye .
一路由過紅海到去.
離開埃及.
去到呢個曠野ge 時候.
佢地就會諗起.
有好多ge 故事.
甚至佢會加上一D 故事D .
加上咩故事呢.
就黎到耶穌黎之前.
gwo D ge 猶太人就會想像.

$^{241}$因為聖經確實有講.
佢地過完呢個.
紅海過完四十年ge 曠野.
入到迦南ge 時候.
就有講嗎哪就唔降落黎.
佢入主迦南jor 之後.
去到佢地ge 屋企.
真係去到目的地.
就唔再有乜料落黎.
然後佢地就食迦南地ge 土產.
所以佢地就會想像.
天上面唔再降嗎哪.
點解呢.
因為原來摩西有條鑰匙.
當以色列人入主迦南之後.
佢就將天上面.
嗎哪ge 倉庫就關jor .
就唔再降嗎哪.
所以佢地就繼續諗.
如果上帝真正ge 救恩.
最終極黎到ge 時候.
應該會係點呢.
摩西會天上面gwo 個倉庫.
就lor 條鑰匙再開返.
然後就會降落黎嗎哪.
然後我地就可以食gwo D 嗎哪.
就係天上降落黎ge 糧食.
以至我地能夠有終極ge 救恩.
天上降落黎ge 天糧.
Spiritual food.
但如果咁樣諗ge 時候.
我地就明白.
咦 聖經裡面提到御節ge 時候.
gwo 時同耶穌有關.
我地睇一段ge 經文.
我地只係揀好幾節黎睇.
因為比較長.
約翰福第六章.
約翰福第六章第四節.
那時猶太人ge 御節近了.

$^{281}$可能未到或者已經過jor .
總之係附近.
猶太人ge Passover.
耶穌舉目看見一大群人黎.
就對肥利講.
喂 我地去邊度買餅比佢地食.
當時佢地加利利.
好多山頭 曠野ge 地方.
耶穌就話 喂 咁多人跟住我地.
咁點呀 你記得個故事.
就係五餅而魚ge 故事.
每每我地睇到五餅而魚ge 時候.
我地就會諗起.
上帝ge 供應.
去到俾人ge 恩典係足夠.
因為五餅而魚就餵飽jor 五千人.
但係你留意到.
呢度特別係約翰福音.
約翰福音先至會咁樣記載.
因為五餅而魚呢個故事.
四本福音書都有提.
但淨係約翰福音係加jor 呢句第四節.
佢就係留意.
呢個五餅而魚ge 故事.
神蹟ge 故事.
係幾時發生.
Passover 附近ge 時候.
你都發生.
約翰福音好多時.
講到上帝藉住耶穌所做ge ye .
gwo D ge Acts.
gwo D ge Great Deeds.
上帝要做ge 工作.
其實係一個Sign.
一個路牌.
即係換句話講.
等陣你去餐館.
你就打開餐館Menu.
你見到有好多相.
佢就話比你聽呢隻ge 菜係咩名.

$^{321}$同埋係咩樣.
gwo 個就係Sign.
Sign係Something point to the reality.
你唔會lor 住個Menu黎食.
你話比個攝影聽我要呢一個.
到佢lor 黎ge 時候就係gwo 碟sung .
貨不對版.
你話同個Sign係唔同.
個Menu就係個Sign.
係約翰福音裡面.
耶穌做ge ye .
一舉一動都係一個Sign.
五餅二魚就係一個Sign.
Point to something beyond itself.
唔係個實體.
而係指向實體ge 東西.
五餅二魚呢個神蹟.
約翰就話諗清楚.
唔係淨係講供應.
每日食ye 咁簡單.
耶穌做緊一樣ye 係最終極.
想做D 乜ye .
約翰福音第六章當然好長.
我呢度純粹比一個Outline比大家.
大概五段.
首先就係五餅二魚.
第二 發生咩事.
係特登咁鋪排.
唔一定係時間上係咁樣發生.
但係記載返耶穌做ge ye ge 時候.
第二個故事.
佢地講到ge 時候.
就係耶穌係加里海平靜風浪.
都唔拉監.
食完ye 跟住又講個平靜風浪.
你記得我頭先講咩.
猶太人諗起御節ge 時候.
御節ge 時候近了嘛.
第六章第四節 約翰福音.
佢就諗起咩 成個Package.

$^{361}$諗起食物.
天上有糧食比我地.
然之後諗起過紅海.
上帝將水分開.
或者上帝將水Control.
將gwo D 混沌邪惡ge ye .
黎到去Control住.
跟住佢就提到耶穌平靜風浪.
所以耶穌平靜風浪又唔係淨係單單.
我地生命裡面有風浪.
上帝就平靜啦.
當然係咁ge 意思.
但亦都係一個Sign.
去話俾我地聽耶穌做緊D 咩.
你呢度諗起就成個御節出埃及ge 故事.
然之後下面三段都係說話.
如果你返去睇下六章約翰福音.
好長呀 你越睇越拗頭.
唔知佢講乜.
總之耶穌就同那些猶太人去辯論.
大家有好多對話.
講D 乜ye 呢.
就係講食物.
因為D 人見到耶穌 五餅二魚.
呢個人成日請食飯.
梗係成日搵佢.
結果佢地真係去搵佢.
由烈國營之後搭船走去搵佢.
搵到佢啦.
佢就同佢地講.
你可唔可以比D 糧食比我地.
耶穌就同佢有好多ge Debate.
有好多對話.
然之後有D 人就發怨言.
聽唔明就走鬼jor .
佢就話呢個人講ye 太爛.
我地接受唔到.
然之後有好多本身跟耶穌ge 人.
就la 一聲走jor .
因為佢都唔再請我地食飯.

$^{401}$就唔好同佢做朋友.
就係咁ge 意思.
然之後耶穌最尾ge 時候.
就Round up講一樣ye .
就同12門徒講.
佢話你地走唔走.
門徒就話.
你有生命ge 道.
我地仲要跟邊個.
你有永生.
我地仲要跟邊個.
類似一個咁樣ge 對話.
我地睇幾節.
讀jor 呢17節.
呢度就係大家有Debatege 時候.
即係gwo D ge 猶太人.
佢地試jor 一次五餅如臘.
嘩都唔使做ye .
然之後呢個人請我地食飯幾好.
搵佢先.
搵到佢啦然之後就同佢講.
喂你可唔可以比D 咁ge ye 比我地食呀.
你可唔可以再降呢D 咁ge 糧呀.
你記得成個預節ge Package.
就係有摩西開上面gwo 個Bank.
然之後降嗎哪落黎呀嘛.
呢個會唔會就係gwo 個呢.
等我地唔使每日做ye 呀嘛.
然之後耶穌就同佢講.
佢話唔好為那會壞ge 食物而操勞.
In fact.
係出《阿俠記》裡面.
講到gwo D 嗎哪跌落黎ge 時候.
如果gwo D 人呢.
將佢儲起ge 話.
拿多少少先.
聽日我驚唔夠呀.
跟住我就雪起佢.
然之後第二日我叮翻熱佢啦.
咁上下咁ge 意思.

$^{441}$會壞ge 其實.
佢就講緊咩呢.
不要為那會壞ge 食物而去操勞.
係你可以搵食.
Sorry.
但唔好為gwo 樣ye 黎到操勞.
為D 乜ye 而操勞呢.
傳住永生ge 食物.
呢D 食物就係.
耶穌話我人子賜比你地.
你要搵gwo D 食物.
係乜ye 食物.
跟住D 人聽唔明34節.
他們就對他說.
主呀請常常每日把這糧給我地啦.
每日你都比我好嗎.
我每日都想.
你比飯我食.
唔記得係出埃及記ge 時候.
佢地行過曠野40年.
真係每日比佢地.
就算到Shabbat黎之前.
都比兩份佢地架.
等佢Shabbatge 時候就唔使煮飯.
40年上帝就係咁樣去供應佢地.
但係呢度呢.
耶穌就有少少.
用佢地所熟ge 故事.
黎到講一D 再深D ge ye .
佢話你食完每一日ge ye .
你地都要再食.
40年都係咁樣食.
事實上gwo D 以色列人食jor 40年.
都係死係曠野.
點都係死硬呀.
點解要為呢樣ye 咁粗魯.
佢都講一D 更加重要ge 食物.
六章第三五節.
佢話我就使生命都要.
黎到我呢度ge .

$^{481}$終不饑餓.
信我的永不乾旱.
我比食物同埋水比你地.
呢一D 係永遠ge 食物.
呢個就係好出名係約翰福音入面.
七句七次耶穌講咩.
我是.
I am.
佢話我係門.
我係羊ge 門.
等等好幾次.
呢就係其中一次.
I am the bread of life.
Eternal life.
佢話我使生命ge 糧.
黎到我呢度ge 話.
你就唔會饑渴.
信我的話就唔會乾渴.
呢D 就係比喻.
因為我地會食D ye ge .
食食物.
我地唔會信食物.
係你面前D 飯都係食ge .
你唔會話我信你阿凡.
我好掂架我信jor 你就得.
佢肯定就係講緊比喻.
我地唔會信麵包.
我地會食麵包.
耶穌就係咁ge 意思.
我就是生命.
信我的就唔會饑餓.
五眼一節.
又繼續debatege 時候就講咩.
佢話我就是從天上降下來生命的糧.
上面個倉庫打開.
降落黎有馬欄.
人若吃這糧.
就永遠活著.
我為世人ge 生命.
所賜下ge 糧.

$^{521}$就是我的肉.
呢個重眼.
咩呀.
講馬欄.
其實佢就話.
我就是馬欄.
佢唔係只俾碟飯黎.
佢話我就是飯.
我就是飯.
你信我的話.
你吃我的話.
你就能夠.
永遠活著.
係以前農業ge 社會.
跟我地有少少唔似.
農業ge 社會做咩.
農業ge 社會好清楚.
佢食任何ye .
食ge ye 都一定要犧牲jor .
你先食到.
舉個例.
我地今日lor 個漢堡包食.
有D 咩死jor 你先能夠食到漢堡包.
以致你唔餓能夠活著.
牛.
除非你gwo D .
牛.
除非你gwo D .
那D 唔知咩黎.
麥當勞gwo D 就係咁.
烏鷹都唔食.
牛.
菜.
gwo 塊菜.
植物死jor 放你面前夾住漢堡包.
你先可以食.
大麥小麥整出黎ge 麵包.
係由咩黎.
係由田產出黎.
係D 麥子死jor .

$^{561}$先至煮到係你面前.
然後你食jor .
你就能夠可以活著.
以前農業ge 社會.
好清楚呢樣ye .
我地話食羊羔.
食任何ye 幾乎.
都係有D ye 死jor .
有D ye 犧牲.
以致我地能夠得到生命.
頭先我地講gwo D 全部還ge 節期.
全部都有ye 死.
新年唔好講呢D ye 啦.
大家猜過你咩.
但樣ye 都死啦.
放得上去ge 都已經死jor .
以致你能夠有生命.
耶穌就係講緊一樣ye .
耶穌就係話.
我ge 死.
使到你有生命.
每一日.
你能夠有糧食.
都係因為有生命犧牲.
以致你能夠有生命.
如果講永生ge 話.
就先係我.
為你犧牲.
以致你唔會受審判.
以致你能夠得救.
我地等間會搜聖餐.
聖餐.
就係從呢一類ge 經文.
gwo 度引申出黎.
就係耶穌ge 死.
每時每刻我地lor 聖餐ge 時候.
諗起佢ge 犧牲.
餅.
同葡萄汁.
大麥小麥ge 死.

$^{601}$葡萄樹ge 死.
能夠祝福我地有生命.
呢個先至係我地要操勞ge .
最重要我地要知道.
我地係領受上帝ge 救人.
開始ge 時候我地提醒大家.
係個感恩節裡面.
為D 乜ye 去操勞.
為D 乜ye 去到講.
Fan Guard.
最重要ge .
係天上降下來.
生命救流.
耶穌基督自己.
我地請我地敬拜隊上黎.
我地唱呢首歌ge 時候.
我地係最後.
我地第一節ge 聖經.
係同呢個預知有關.
同耶穌ge 聖餐有關.
佢話既然你地係無孝ge 面.
佢講緊教會.
講緊御節.
無孝ge 面.
舊孝隨正.
好使到你地成為新團.
我地ge 御節ge 羔羊.
基督就係被殺.
犧牲ge .
我地到守呢個節.
其實佢講緊聖餐.
唔係淨係猶太人ge 節.
守呢個節ge 時候我地就係.
作無孝ge 餅.
孝代表罪.
代表一D 唔好ge ye .
隨去ge 時候.
領受呢一個ge .
真正心靈ge 流量.
以致我地能夠.

$^{641}$永遠活著.
我地同心大道基督.
上主我地多謝你.
因為你係gwo 一位.
帶領我地.
愉悅ge 羔羊.
你就係gwo 個.
天上降ge 力量.
叫我地能夠得到.
永生.
上主我地感恩去讚比你.
因為你為我地預備ge 救恩.
你就係你ge 愛子就係.
御節ge 羔羊.
係十家上面.
為我地流血寫生.
救我地離開.
罪惡.
離開審判.
離開死亡同黑暗.
主帶領我地呢一刻.
去到領受呢一個.
代表你同在ge .
呢個餅同呢個盤.
提醒我地.
你救恩ge 貴重.
主當我地領受ge 時候.
我知道你就係天上降下來ge .
生命不可.
叫我地永遠不記憾.
主求你ge 救恩.
能夠係我地內心裡.
能夠係我地ge 心裡.
係我地生命最深ge 地方.
去到滿足我地.
祈禱扶住命.
我地先唱呢首歌.
在困苦艱難.
仍謙卑土高.
要安靜.

$^{681}$知道你是神.
無罪羔羊.
曾被失怨犧牲.
為救贖.
流血捨身.
為救贖.
流血捨身.
凝望月死亡.
月掛方岸.
神恩神不可欣賞.
成功是實在和.
跌剩一生無盡.
知己都不懼痛難過.
願救恩歸於.
坐在寶座上高昂.
耶穌.
願救恩歸於.
坐在寶座上高昂.
耶穌.
願救恩歸於.
坐在寶座上高昂.
耶穌.
願救恩歸於.
坐在寶座上高昂.
耶穌.
凝望月死亡.
月掛方岸.
神恩神不可欣賞.
成功是實在和.
跌剩一生無盡.
知己都不懼痛難過.
願救恩歸於.
坐在寶座上高昂.
願救恩歸於.
坐在寶座上高昂.
耶穌.
願救恩歸於.
坐在寶座上高昂.
願救恩歸於.
坐在寶座上高昂.

$^{721}$願救恩歸於.
坐在寶座上高昂.
願救恩歸於.
坐在寶座上高昂.
願救恩歸於.
坐在寶座上高昂.
願救恩歸於.
坐在寶座上高昂.
願救恩歸於.
坐在寶座上高昂.
願救恩歸於.
坐在寶座上高昂.
願救恩歸於.
坐在寶座上高昂.
願救恩歸於.
坐在寶座上高昂.
願救恩歸於.
坐在寶座上高昂.
願救恩歸於.
坐在寶座上高昂.
願救恩歸於.
坐在寶座上高昂.
願上帝祝福各位.
希望下星期能和大家一起事奉敬拜.
\newpage



\section{}
\label{sec:kSJ3Fo7YGUs}
\textbf{Reason For The Season - Rev. Mark Chan}
\newline
\newline
連結: \href{https://youtube.com/watch?v=kSJ3Fo7YGUs}{\texttt{https://youtube.com/watch?v=kSJ3Fo7YGUs}} ~~~~ 語音日期: 2022-12-20
\newline
\newline
\hyperref[sec:SrXfS_omV5U]{\small{< < < PREV SERMON < < <}}
~
\hyperref[sec:index_chronic]{\small{[返順時目]}}
~
\hyperref[sec:index_scriptual]{\small{[返順卷目]}}
~
\hyperref[sec:du2XrTgHsNA]{\small{> > > NEXT SERMON > > >}}
\newline
\newline
$^{1}$And let's welcome Pastor Mark for the sermon..
- Hey, good to see all of you, at least most of us, right?.
I know it's a very exciting game, the World Cup..
It comes down to penalty kicks,.
and I don't know, which team are you rooting for?.
Which team are you rooting for?.
Here..
It's over?.
It's already over..
So Argentina won, right?.
All right, all right, let's celebrate, let's celebrate..
Okay, anyways..
Well, good to have the update so we can pay attention,.
and we can concentrate of why we're here,.
because we're here to worship, right?.
So I would like to share something.
about the Christmas story recorded in the Bible,.
in the New Testament..
When we come to Christmas,.
when we come to this Advent season,.
there are many questions that we may ask about Christmas..
For example, especially Christmas,.
you know, it has inherited a lot of different traditions,.
in our Western world, at least, right?.
So one tradition is Santa Claus, right?.
Santa Claus..
Do you think Santa Claus has anything to do.
with the Bible Christmas story?.
What do you think?.
- No..
- Yeah..
(congregation laughing).
You know where Santa Claus come from?.
- Coca-Cola..
- Coca-Cola?.
(laughs).
You know..
- St. Nicholas?.
- St. Nicholas, yes, yes, you're right, James..
Yeah, St. Nicholas..

$^{41}$Is a bishop in today's Turkey..
He was a Christian leader and known for being kind.
and helpful to those who are needy..
Back in the days, AD 200, 300, something like that,.
or 200 to 300 years after Christ..
So it's an early church leaders..
One of them is called St. Nicholas..
Because of his kind work,.
you know, the tradition has developed.
that someone who is giving gifts to those who are needy.
in the history, in the church history..
So that's why we have Santa Claus, right?.
So although Santa Claus today has very little to do.
with the Christmas story,.
have very little to do with the Bible,.
but you need to understand the root of the tradition,.
the root of the legend comes from a bishop.
called St. Nicholas..
Secondly, why December 25th?.
Why?.
Have you ever asked this question?.
Now, according to the Bible,.
actually Jesus was born in the summer days..
You know?.
Actually in the summer days..
You know why?.
Because Luke 2 tells us the shepherds.
are grazing the sheep in the field of Bethlehem..
Now, according to the ancient way of working.
between farmers and shepherds,.
after the harvest of the barley and wheat,.
they would invite the shepherd to come with their sheep,.
and they will eat the stubble of the leftover.
from the harvest..
So they would let them into the field..
So that is right in the summer days, right?.
Right in the summer days..
'Cause harvest happens in April and May,.
and then the summer season kicks in..
The shepherd will come to the farming fields,.

$^{81}$to the land..
They will graze in the land of the farmers..
So there was a contract between the shepherd.
and the farmers, and they would let them in..
So with that background,.
we kind of know Jesus was actually born in the summer days..
All right?.
I'm wrong, okay?.
But why December 25th?.
Now, one theory, one theory,.
is that during the early church,.
when Christianity was being recognized,.
they found that on December 25th, or around this time,.
in the Roman empire, the Roman people.
would celebrate the sun god..
It's a sun god festival..
They would give gifts and lights,.
and something like that, like candles,.
to celebrate the arrival of sun god..
But when Christianity is being recognized properly.
in the Roman empire,.
the Christians, the Roman Christians,.
they took over the festival,.
and replaced it with the celebration of Christmas,.
the coming of true God, son of God..
So that's how it has developed,.
December 25th, instead of worshiping the false god,.
the sun god, instead of being controlled by him.
to worship him, the Christians were saying,.
let's celebrate the true God, son of God, Jesus Christ,.
who come into our world, right, to give us hope..
Not the sun god, not the false god,.
but the true God, son of God..
So with that background, you understand.
the first Christmas was actually counter-cultural,.
because the entire cultural world of the Roman empire.
was celebrating following the false god,.
you know, the sun, the moon, and the stars..
But the Christians come along and saying,.
hey, let's not do that, and replace it.

$^{121}$with the true meaning, with the true God, son of God..
So that's why December 25th, right?.
And another theory is how they figure out.
is 25th of December was the birth of Jesus Christ,.
that they count nine months from Passover..
Because according to one tradition,.
Mary was conceived by the Holy Spirit on Passover,.
but there's no other better day for her to be conceived,.
I guess, so they thought of the conception happened.
in Passover, and then count nine months,.
because in the ancient days, they believe after conception,.
there'll be nine months in the womb,.
and then giving birth to a child..
So they count nine months from Passover.
to end up on December 25th, right?.
But there are different Western and Eastern church,.
they have different traditions.
of using different calendars..
So sometimes they calculate the Passover.
would be on a different day..
So some of the Eastern Christians, Eastern Orthodox,.
they will not celebrate Christmas on December 25th..
They instead, they celebrate in January..
So if you go to a Eastern Orthodox church,.
you worship with the Orthodox communities,.
they would not celebrate on the 25th of December,.
they would do so on January 7th or 6th, something like that..
So there are a lot of questions.
we can ask about Christmas..
However, there are one question.
there is one question we have to ask.
is what does the Bible say about the Christmas story, right?.
So today I would like to turn to the Gospel of Luke,.
Luke chapter one, to take a few verses.
out of the first two chapters of the Gospel of Luke.
and see what the Bible writers want us to know.
about the Christmas story..
Of course, there are tons of things that we can ask.
about the birth of Jesus Christ,.
the scenario, the circumstances,.

$^{161}$what kind of diaper did he use,.
there are all sorts of different things we can ask,.
but what the Bible want us to know.
and don't want us to miss,.
we have to go back to scripture, right?.
So there are two Gospels.
that talk about the Christmas story, okay?.
Which one?.
The Gospel of Luke is the first one, it's one of them,.
one of the two..
The second one is Gospel of Matthew, right?.
Only the Gospel of Matthew and the Gospel of Luke.
talk about the Christmas story..
So Luke chapter one, here we go..
When Herod was king of Judea,.
there was a Jewish priest called Zechariah..
He was a member of the priestly order of Abijah.
and his wife, Elizabeth..
She was also from the priestly line of Aaron..
Now, Zechariah and Elizabeth were righteous in God's eyes,.
careful to obey all of the Lord's commandments.
and regulations..
They, however, had no children.
because Elizabeth was unable to conceive.
and they were both very old..
Now, notice how the story goes..
This is the stage setting first..
Tell us what happened in the first Christmas, right?.
So chapter one talks about the birth of John the Baptist,.
the cousin of Jesus, the pioneer of Messiah, right?.
The foreigner of Messiah..
So they talk about a couple, Zechariah and Elizabeth..
Now, how does the Bible describe them?.
When you read the Bible,.
sometimes you can pick up some details.
and ask the question,.
what kind of social statements are these?.
Social statement meaning statements of descriptions.
that talks about their social status,.
their position in the society,.

$^{201}$their ranking in the society..
What kind of people do they belong to?.
Zechariah and Elizabeth, old couple, no children, right?.
Now, in those days, in the biblical world,.
when you had no children,.
you cannot conceive, you don't have a child,.
especially you don't have a son..
You are considered being under curse..
You're being under curse..
If you're able to give birth, you have God's blessing..
But in those days,.
in the ancient society of honor and shame,.
not being able to conceive is considered as a curse,.
is considered as a shame..
It's a social shame..
So you don't have the social currency of boasting,.
hey, I have son, right?.
I have a daughter..
I have my own family..
I have a next generation..
In those days, you are considered a curse.
if you cannot have a son or daughter..
And also, they are hopeless..
Elizabeth and Zechariah as a couple,.
very old, hopeless..
So you have no hope of bearing another son.
because the day is coming, you're counting down,.
you're very, very old..
So when Luke talks about the first Christmas,.
he introduces us to the first main characters,.
Elizabeth and Zechariah, under curse, hopeless, very old..
Now, let's fast forward..
Chapter one, the rest of the chapter one.
talking about the birth of John the Baptist..
And then chapter two, we'll talk about the birth of Jesus,.
the real story of Christmas..
Chapter two, verse one..
At that time, the Roman emperor Augustus.
decreed that a census should be taken throughout the empire..
This was the first census taken.

$^{241}$when Quirinius was the governor of Syria.
and all the people returned to their own ancestral towns.
to register for this census..
So that's why we have the story, Mary and Joseph,.
they probably got a donkey..
I guess they didn't..
They didn't have enough money to rent a car or something..
So they travel all the way from the north, Nazareth,.
where they are, and then all the way down to the south,.
to the central hill country, south of Jerusalem,.
to Bethlehem, go back to their hometown..
And then the story goes, as you can imagine,.
they went back to their hometown.
in order to register for this Roman census..
They couldn't find an Airbnb,.
so they end up in a barn or something..
They end up in a fairly humble place..
That's how the story goes..
Now, notice when I read Luke 1, 5,.
the introduction of the Christmas story, John the Baptist,.
and then chapter two right here,.
talking about the birth of Jesus, the first Christmas,.
how many characters are being introduced?.
Or maybe I should ask, how many powerful characters.
are being introduced in those days?.
Chapter one, five, Herod the Great, right?.
Herod the Great..
He was a mighty great king..
He ruled the land of Palestine and Israel in those days..
He has so much power..
Actually, he had the absolute power in the region,.
in the eastern frontier of the empire..
Here, chapter two, and other two powerful characters.
are being introduced..
Augustus..
Now, in those days, the superpower of the entire world.
was the Roman emperor,.
and Augustus was the first emperor of Roman history..
And then the second character is Quirinius..
Quirinius is the prophet of Syria,.

$^{281}$which include Israel and Palestine,.
the governor of Syria, and Quirinius is the commander..
He's the governor of the entire province..
Now, as you can tell,.
the three powerful characters in those days,.
absolute power, super powerful figures,.
are being introduced in the first Christmas story,.
according to the Gospel of Luke..
Herod the Great, Quirinius, and then Augustus..
Now, think about it..
Now, if you walk into a bookstore, okay,.
I know bookstore is fading away, okay?.
Let's imagine you go to a library or go to a bookstore..
You go to the biography section,.
biography section..
What kind of biographies you will find?.
Or what kind of biographies will be the best sellers?.
Has to be a famous figure, right?.
Not too long ago, I was in the library..
I bumped into a book written by Hillary Clinton..
Biography of hers..
She's famous, right?.
Or she's powerful..
So these kind of biographies will be the best sellers..
In Jesus' day, it was the same way..
The powerful characters, the powerful people.
who have influence in their own world.
or in the lives of many other people,.
they have their own biographies..
Sometimes they have inscriptions..
Sometimes they have some public billboard.
to talk about their good deeds, to talk about their birth..
For example, Augustus, he was born on a certain day,.
and that certain day becomes a special holiday.
for the entire empire, for all the people..
And when Augustus was born,.
when people talk about his biographical information,.
Augustus was being called the son of gods,.
the savior of the world..
He bring peace, the Roman peace, the Pax Romana,.

$^{321}$because of the Romans, because of the Roman military might..
You will have peace, Pax Romana, Roman peace..
So when Augustus comes, when he was born,.
he brings the gospel..
You know, the gospel was a political term in those days..
The gospel meaning the news, the good news.
of the king being born, the kingdom being ushered in,.
being launched because of this figure..
So because of Augustus, you have joy and peace.
and good tidings and gospels..
And he is the son of gods..
Now, these kind of powerful figures.
will have their own famous biographies,.
the best sellers of those days..
People would read that..
People would listen to those stories and thinking,.
geez, they really control our lives..
They really are very powerful..
But according to the gospel of Luke,.
when he writes about the story of Christmas,.
these guys, they are not occupying.
the foreground of a stage..
They are not in the front of a stage..
They are in the background of a stage..
These guys who can do a lot of things.
to control the lives of the common people,.
they are in the background..
And who is in the foreground?.
The minor characters..
Zechariah, Elizabeth, John the Baptist,.
Joseph and Mary, and Jesus..
They should belong to the background..
But in the gospel of Luke,.
they are being brought forth to front stage..
Now let's take a look at the first Christmas story,.
chapter two, verse four..
And because Joseph was a descendant of King David,.
he had to go to his hometown Bethlehem in Judea,.
David's ancient home..
He traveled there from the village of Nazareth.

$^{361}$in Galilee in the north..
He took with him Mary to whom he was engaged,.
who was now expecting a child..
And while Joseph and Mary were there,.
the time came for her baby to be born..
She gave birth to her firstborn son, Jesus,.
wrapped him, and then put him in the manger..
Because there was no lodging available for them..
Now this is how the story goes, right?.
As I said, they came to their hometown Bethlehem..
They had no place to stay..
Now in this short passage,.
many times it talks about census, the Roman census..
Now it happened that we know.
the Roman emperor would call for census.
for the entire emperor once every 14 years..
So about 10 to 14 years, they would call for another census..
And the census will dictate the people.
to go back to their hometown..
Maybe they're working somewhere else,.
some other town, some other cities..
But according to the decree of the Roman emperor,.
they have to return to their hometown to register,.
to pay taxes..
Kind of like, you know, once a year we have to pay tax, right?.
In March, the tax day, the IRS..
Here, Luke is talking about the Roman IRS.
calling you to go back to your hometown to pay for taxes..
But a little bit different is that when they decree you.
or they dictate you to pay the Roman taxes,.
it is heavy tax, very heavy tax..
So whenever the Roman emperor called for a worldwide census,.
there was always unrest in the society..
People would fight back and people would revolt,.
something like that..
As you can imagine what kind of time they live in..
Joseph and Mary, they rush from Nazareth.
back to their hometown, Bethlehem..
And you see a lot of social unrest,.
just like what we see today..

$^{401}$It's very unsafe to travel..
They end up in Bethlehem..
And verse five says, "Now expecting a child," right?.
And you can imagine how difficult it is to travel.
with a pregnant woman, right?.
And the last verse here, it says,.
there was no lodging available for them..
Now, according to our modern imagination, as I told you,.
they went to Bethlehem and they couldn't find an Airbnb.
or a Motel 6, and then they end up in a barn or something..
But if you understand the background of the story,.
most of the Palestinian houses were designed.
as a two-flat house..
There were two stories..
The top upper level is the guest room.
or the living room or the sleeping quarter,.
or the bedrooms, it's upstairs..
And then downstairs is the manger.
or the kitchen or the area of productivity..
They cook, they sew, they work on their farm stuff.
and crafts in the lower level,.
where you have sheep and donkeys and goats and all that..
Now, if you can imagine, there were two stories of the house..
The upper level is actually the guest room..
Actually, that's what the gospel is telling us..
The word being used is a special word.
for guest room, for the upper room..
Remember Jesus had the last supper?.
Where did he have it with the disciples?.
The upper room..
The gospel used the same word..
It's the upper room that they had the last supper..
They had the Passover meal, the last Passover meal..
And the lower level is where the animals are..
Now, this kind of design is very practical..
You know why?.
Because when in the winter days or at night,.
who's keeping you warm?.
Now, in Chicago, it's getting really, really cold,.
getting really, really chilly here..

$^{441}$We have heater, right?.
We have water heater, we have the pipe heater.
to keep us warm..
But the Palestinian houses are designed in such a way.
as to keep warm of the people who stay.
or who sleep in the second floor..
The first floor, who's keeping you warm?.
The donkey, the animals, right?.
Anyone who have pets here?.
Dogs and cats?.
Do you sleep with them?.
You do, right?.
Why?.
They want to, right?.
Why, why do they want to sleep with you?.
Because you keep them warm, right?.
And that they keep you warm, all right?.
So that's the same way..
Lower level, where the donkeys are,.
where the horses or goats or sheep,.
they will keep the upper level warm, right?.
Hot air rises..
So that's how the house was designed in those days..
So if with that background,.
we can actually translate the last verse here,.
because there was no room in the guest room.
available for them..
Now, as you can imagine, the Roman IRS comes in,.
everyone has to go back to their hometown, right?.
And when you go back to your hometown,.
where are you gonna stay?.
Usually you stay with your relatives, right?.
Your forefathers or your family members, right?.
Or your distant family members..
You knock on their houses and say,.
"Hey, can I stay here?".
Right?.
Something like that..
Next week, I'm traveling to California.
after the Christmas service here..

$^{481}$So four of us will travel back to Northern California,.
where most of my in-laws are there..
You think we would have to rent an Airbnb?.
No, no, we knock on their houses..
And we stay with them, right?.
That's what happened with Mary and Joseph,.
because Joseph's hometown is in Bethlehem..
So they went back to their hometown..
So many people are being squeezed into the guest room..
And Mary and Joseph just happened to be at the doorstep..
They couldn't find any place on the upper level..
And by all means, Mary was pregnant, right?.
It's pretty tough for a pregnant woman.
to climb to the second level anyway, right?.
So they probably stay in the first floor, right?.
So there's nothing like, "Oh, they're so poor,.
"and we pity them," right?.
It's actually a story of inviting.
and welcoming the Holy Family..
This family, whoever that is, in Bethlehem,.
are actually inviting Joseph and Mary..
I can imagine they would go to the first floor.
and try to put things together, put things in order,.
so that Joseph and Mary can stay in a warm place, right?.
To wait for doing their things, their taxes in Bethlehem..
With that background,.
Gospel of Luke chapter two is actually a story.
of inviting Jesus into their home, into their house..
Now, we started with chapter one..
Now we're in chapter two..
And as you know, there are different parts.
of the Christmas story..
The shepherd welcomed Jesus..
They received the gospel..
Zechariah, Elizabeth, they saw God is speaking..
They saw God is working..
The history of redemption continues..
God show up again, speaks to them again..
Chapter two, this family in Bethlehem.
welcomed Joseph and Mary,.

$^{521}$and ultimately, the birth of Jesus..
It's a welcoming story..
When Jesus comes,.
do you have room for him?.
Are you inviting him to come into your lives?.
Now, as I told you, chapter one and chapter two,.
all this superpower, all the important figures.
are being pushed back into the background..
And what comes forth to the center of the stage,.
getting the spotlight, are the common people,.
the lower level, those who are ordinary,.
those who are being seen as unimportant people,.
or even shamed, being cursed..
Unimportant people are taking the center stage.
in the Christmas story..
Now, every year comes to the month of December,.
to the end of the year..
If you notice, Time magazine.
always elect someone who is important, right?.
Someone who is important..
To be the icon of the year,.
to be the person of the year, right?.
The icon of the year, this year,.
who is it, do you know?.
Michelle Yao..
Michelle Yao, right?.
You know Michelle Yao?.
Michelle Yao..
Y-E-O-H..
Yo, right, yo, or Yao, right?.
She has been an action figure,.
Asian actress, for so many years, right?.
How is she right now?.
60s, 60s, yeah..
She's being elected icon of the year..
And the person of the year in 2022, who is that?.
Zelensky, Zelensky..
Ukrainian prime minister, right?.
Because of him, he keep up the spirit of the Ukrainians,.
you know, to fight back of the Russian invasion..

$^{561}$Now, this kind of people are being regarded important..
Now, who among of us would be qualified.
to be called Time's magazine person of the year, right here?.
Who is it?.
Who would have the qualification to be like that, right?.
Nobody, probably, right?.
Nobody..
None of us..
So, according to the eyes of the world,.
we are not important..
And sometimes we feel that way..
Subjectively, personally, we think that people don't care..
We don't think people would care about our future..
What are we doing?.
But in the eyes of God, according to the gospel of Luke,.
you are important..
The ordinary people have the spotlight..
The ordinary people, the powerless people,.
as they receive Jesus,.
they would be regarded as valuable in God's sight..
It's valuable in God's sight..
Some years ago, I read something about Mother Teresa..
You know Mother Teresa, right?.
And she was famous, right?.
You know, she was not very famous in the beginning..
For many years, nobody knows about her existence..
It's only because in 1970s, I believe it's 1970s,.
when a BBC reporter went to Calcutta to interview her,.
to ask what Mother Teresa is doing.
on the streets of Calcutta..
What she was doing is to wash the bodies.
of those who are dying on the streets,.
bring them home so that they can die in dignity..
That is Mother Teresa..
Day after day, she was doing something like that..
So nobody knows, nobody pay attention..
Only one day, a BBC reporter was curious.
about what she's doing, and air an interview..
So the whole world know who Mother Teresa is..
But for many years, she was silent..

$^{601}$She was quiet, doing her work..
Mother Teresa once said, and I want to quote,.
"Not all of us can do great things,.
"but we can all do small things with great love.".
(audience member sighs).
Not all of us can do great things,.
but we can do small things with great love..
You see, that's the Christmas story..
Not all of us can do great things, right?.
Sometimes we just thought that our works are insignificant..
Nobody pay attention..
But we can all do small and little things with great love..
Couple weeks ago, I came into the church.
during the weekdays..
So as usual, I check out some of the mails that come in..
So I open some of the envelopes..
Some of them are junk mails, right?.
But I receive a card..
And then I notice, where is it from?.
It's from our state..
It's not from Illinois, it's from other state.
of the country..
So I was curious, who would write a card.
to us, right?.
So I open the envelope..
There was a card in there..
There was a small offering being made to our church..
And I read the content of the card..
It's addressed to the Yuvi family, to our church..
It's written from a very old couple, very old parents..
And they wrote to us and thanked our church.
for reaching out to one of his children..
Because they have been praying for the children.
for a long time, adult children..
Praying for them for a long, long time, for many years..
And because of our ministry, we received this family.
into Christ..
And I read that card, I read the message..
I finally realized, this is what our church.
should be doing, right?.

$^{641}$This is what we are here..
That's why we are here..
That's the purpose of our existence..
We're here, it's not because to please ourselves,.
but as instrument of the grace of God..
We can receive someone who has been lost for a long time.
into, back to the family of God..
And in fact, this is the spirit of Christmas..
Now in this Christmas, as we end,.
I want you to ask yourself, what can you do?.
What small things can you do with great love?.
Perhaps visiting someone, or attending to someone.
who is lonely..
Now ask about someone who has been forgotten..
Who's gonna do that?.
You are going to do it..
Small and little things with great love, right?.
But the ultimate question is,.
do we have room for that great love?.
Jesus Christ himself came into your world..
Open your home, open your lives,.
open your heart to receive it..
Because without Jesus in our lives,.
we would not be able to do small things with great love.
in the kingdom of God..
So think about some of the people you can reach out to.
in this season..
Think about it, all right?.
Pray for them and do something..
Doesn't have to be great things, small things,.
but with great love..
That's the closing prayer..
Father, we give you thanks because you see us as important..
Father God, we give you thanks.
because we are valuable in your sight..
And I pray that in this season,.
as we celebrate the birth of Jesus, our Lord, our Savior,.
we would remember that in your sight,.
in your vision, that we can be your instruments..
Help us to do small things, to bless other people,.

$^{681}$to be your instrument, to bring love and joy and peace,.
the gospel, the true gospel to our world..
Maybe someone that we know, maybe our friends,.
maybe writing a card, you know,.
send some words of encouragement to say hello,.
to connect, to reconnect..
I ask that you will use us as your instrument.
to bring forth your peace and your joy in this Christmas..
In Christ's name we pray, amen..
Well, God bless you and Argentina one,.
and we can have lunch, okay?.
I wanna see you next Sunday, 25th..
We're gonna have a banquet here,.
luncheon here with a special Christmas dinner..
\newpage



\section{}
\label{sec:du2XrTgHsNA}
\textbf{聖誕的原因 - 陳崇基牧師}
\newline
\newline
連結: \href{https://youtube.com/watch?v=du2XrTgHsNA}{\texttt{https://youtube.com/watch?v=du2XrTgHsNA}} ~~~~ 語音日期: 2022-12-20
\newline
\newline
\hyperref[sec:kSJ3Fo7YGUs]{\small{< < < PREV SERMON < < <}}
~
\hyperref[sec:index_chronic]{\small{[返順時目]}}
~
\hyperref[sec:index_scriptual]{\small{[返順卷目]}}
~
\hyperref[sec:LoAxQ1cr1f0]{\small{> > > NEXT SERMON > > >}}
\newline
\newline
$^{1}$今日ge 中文講完係Mark Sir 而家將時間交比佢.
大家好 今日好高興見到大家.
我知道今日係打世界盃.
而家打下下半場.
後面工作人員話比我聽.
2比0 阿根廷係贏緊.
ok 大家好安心.
可以聽到.
都唔使睇.
點都好.
我地繼續去迎接張琳琦.
張琳琦的Advent Season.
聖誕節就黎到.
我地思想一下聖誕節.
我地可以有好多問題問.
譬如聖誕節 D 舊ye .
即係有D 唔同ge symbol.
唔同ge 記號.
唔同ge 蠟燭等等.
係邊度黎ge 呢.
譬如有D 問題係聖誕老人.
呢舊ye 係邊度黎.
唔知.
聖誕老人係主後大概.
二百幾年 大概三百年.
有個人叫聖尼克勒斯.
呢個主教.
係一個基督徒ge 主教.
聽講佢就係.
係傳記裡面講.
就係一個好樂意去.
行善ge 一個人.
去對其他有需要ge 人.
好盡心去幫助.
正因為咁ge 緣故.
所謂聖誕老人就係從.
聖尼克勒斯慢慢慢慢.
黎到去.
演變出黎.
個燈有少少壞.

$^{41}$等間下午我地就會.
掉jor 佢個燈.
另外一個問題.
我去問 除jor 聖誕老人之外.
就係呢個聖誕節.
日子.
係點樣黎ge 呢.
咁.
十月廿五號 聖誕節.
Christmas day.
都幾distracting.
不過唔緊要.
聖誕嘛.
有少少斬燈係okge .
Thank you.
Anyways 聖誕日子.
係邊度黎架呢 有一個ge 說法.
就係話係.
早期ge 教會.
gwo 陣時.
以前gwo 個羅馬帝國.
羅馬帝國.
就有一個慶祝太陽神.
ge 一個節日.
而gwo 個太陽神ge 節日呢 就係arm arm .
係呢一個ge 時間.
咁所以到呢個ge 教會.
開始係歐洲.
或者西方ge 世界.
去到受到人.
重視ge 時候 佢地就lor jor 呢個.
節日呢 黎到去取締jor .
將耶穌ge 降生.
取締jor 呢一個.
ge 太陽神ge 節日.
咁當時佢地就會送禮.
咁然之後呢 教會ge 傳統.
將佢取締jor 之後 又會將.
一D 民間裡面.
有ge .

$^{81}$呢D ge .
做法 佢地會lor jor 佢出黎.
所以呢 如果係咁樣睇法ge 話.
聖誕節呢 係一個.
counter cultural.
celebration 黎ge .
即係話呢 世界點樣去.
慶祝 慶祝ge 就係太陽神.
假ge 神 但係呢.
基督教ge 教會.
就lor jor 呢個ge 節日呢.
黎到去.
慶祝真正ge .
救主 耶穌ge 理能.
呢個係另一個說法.
另外一個ge 說法.
點解會有聖誕節呢個節日.
就係.
初期教父.
或者初期gwo D ge church fathers.
佢地去計呢 耶穌.
出世gwo 日.
係幾時呢 佢當時就諗呢.
一定係passoverge 時候.
係童貞女瑪利亞.
感人.
所以一數就係.
九個月 佢當時係計九個月.
唔係十月懷胎 九個月懷胎.
所以十二月廿五號.
就係耶穌正式出世ge 時候.
咁所以有D 著作.
就講到討論到Christmas day.
係幾時 佢由passover.
一路計就計九個月.
因為咁ge 緣故.
所以就有十月廿五號.
但係當然.
另外一個問題就係 到底個passover係幾時.
咁所以有唔同ge 說法.

$^{121}$就講話passover前前後後.
爭少少.
咁所以今日就唔係淨係十月廿五號.
係某D 教會黎到慶祝Christmas.
譬如東方ge 教會.
Orthodox Church 佢地就會係一月ge 時候.
先會慶祝.
呢個Christmas day.
因為大家計個passovergwo 日開始.
就係有D 唔同.
但我地可以用呢D 咁ge 問題.
黎到去討論一下聖誕節.
但係一樣ye 一定要做ge .
就係睇下聖誕.
聖經係點講.
你諗起聖誕節 諗起乜ye .
Christmas day.
好似一D 童話故事咁.
即係有D 好可愛ge ye .
或者有D .
好安詳ge .
衣景咁.
聖經裡面係點講.
你越睇ge 時候.
你就越發現.
唔係一個咁簡單ge 故事.
我地睇幾節經文.
路加福音第一章.
第二章就係講到聖誕.
你記得係聖經裡面.
只係得兩本書.
兩本福音書.
係講到Christmas.
第一本.
係邊本呀.
馬太福音.
第二本就係路加福音.
Gospel of Luke.
第一章第五節.
當希律作猶太王ge 時候.

$^{161}$亞比亞班裡面.
有一個祭司.
叫做撒迦利亞.
佢ge 妻子.
係亞倫ge 後代.
就叫做Elizabeth.
伊利沙伯.
佢地兩個人在神ge 面前.
都是二人.
進行主一切ge .
誡命同埋條例.
會有可指責的.
指示佢地會有孩子.
因為Elizabeth係畢生奴.
兩個人呢亦都係年紀老邁.
呢度第一次.
黎到去講聖誕節ge 故事.
就介紹gwo 個ge 場景.
呢度講到大希律.
Hell the Great.
佢作王ge 時候.
就發生jor 一件事.
就係有兩個夫婦.
係年紀老邁ge .
係沒有兒子ge .
呢個係第一個ge 介紹.
第二個介紹呢就係第二章第一節.
黎到講開呢個ge .
聖誕節附近發生ge .
一D 事蹟.
第一節.
當那斯人子.
凱撒奧古斯督.
Augustus.
有意旨下來.
叫天下人民都報名上策.
呢個時候呢.
亦都係居里紐.
作Syria.
敘利亞巡撫ge 時候.

$^{201}$係頭一次去到行.
報名上策ge 事.
眾人就各歸各城.
報名上策.
我睇jor 兩節經文.
兩節經文就提到.
當時一D 歷史ge 人物.
歷史人物唔係聖經提過.
當然 聖經以外.
好多歷史著作.
或者傳記都會提到.
有邊幾個.
有大希律第一章第五節.
有奧古斯督.
有居里紐.
呢三個.
都係當時ge 歷史人物.
唔單止係咁.
亦都係當時有權有勢ge 人.
好大powerge 人.
甚至係有absolute power.
大希律.
當然就係管猶大地.
即係巴勒斯坦.
以色列呢個地方.
即係當時.
Roman Empirege .
東面ge 防線.
另外就係.
奧古斯督.
就係整個地中海世界統一天下ge .
emperor.
第三個就係居里紐.
居里紐就係.
管 頭先提到.
以色列巴勒斯坦地.
一個大D ge 省份叫做敘利亞.
Syria.
呢一個ge 官就係居里紐.
三個.

$^{241}$都係歷史人物.
三個都係有權有勢.
可以呼風喚雨ge 人.
當以前.
去到提到.
某D 人出世ge 時候.
特別係.
一D 歷史著作.
特別係名人.
先會有D 所謂birth story.
可以講.
即係佢出世ge 時候幾巴閉.
出世ge 時候有人提過.
出世ge 時候發生D 乜ye 事.
呢D 所謂birth story.
都係比一D hero.
或者比當時係最重要.
最重要ge .
譬如奧古斯督.
當時有absolute powerge ye .
係管理.
整個ge 地中海ge 世界.
係佢gwo D ge .
出生ge .
記載.
gwo D ge birth story.
講到佢邊年邊月出世.
講到佢出世ge 時候.
佢就帶黎福音文.
Gospel.
講到佢出世ge 時候就帶黎.
有大ge 喜樂.
joy.
good things.
佢出世ge 時候.
亦都係.
眾民 整個ge 天下.
裡面ge 人民.
ge 救主.
saviour.

$^{281}$亦都係當時帝國ge 主宰.
甚至有一D ge .
文獻話佢就係son of gods.
即係眾神ge .
兒子 天子.
奧古斯督.
亦都係gwo 一位帶黎.
新ge peace.
佢地叫做Roman peace.
Pax Romanum.
帶黎一個新ge kingdom.
一個新ge 國度.
所以奧古斯督ge 生日係每年都會慶祝.
等於一個今日.
名人.
我地每年都會去到.
記得 某年某日.
就係佢ge birthday.
就係佢ge 生日.
係一D 世界舞台上面ge 主角.
兩節經文.
我地睇jor 呢D .
重要ge 人物.
但好奇怪 路加介紹ge 時候.
呢D 人呢.
就有咩地位呢.
佢地就唔係舞台ge 中間.
佢變成jor 舞台ge 背景.
唔係foreground.
係background.
當時人人聽見呢個名.
都會諗起佢ge 事蹟.
諗起佢ge 故事.
佢地係世有權ge 男人.
有absolute powerge 人.
每每睇到佢地ge 時候.
佢地世界上面 梗係舞台ge 中間.
foreground.
路加福音.
去到介紹.

$^{321}$耶穌出世ge 事蹟.
講到施洗約翰.
作為佢ge 前鋒.
ge pioneer 講到耶穌ge 出世.
都係.
將到佢地調返轉.
路加福音第一章.
頭先我講過.
大希律作為背景.
有咩人.
有伊利沙伯.
伊利沙伯.
係.
兩夫婦同埋撒迦利亞.
兩夫婦都係.
畢生慾同埋係年紀老邁.
如果你將聖經.
你去查ge 時候.
你當呢D 咁ge statement係social statement.
即係話比你聽.
一D 描述.
係講到佢地社會ge 地位.
你就知道.
畢生慾同埋年紀老邁.
並非係一D 註腳.
其實就係講清楚.
佢地係社會.
係世界係一種點樣ge 人.
小人物.
呢個of course.
但同時.
年紀老邁.
甚至係畢生慾.
近東ge 社會.
古代ge 社會.
係受咒助.
cursed.
甚至係今日我地可以叫佢低端人口.
gwo D 人.
無權無能ge .

$^{361}$但係佢地係.
上帝面前.
作為二人.
即係上帝面前.
佢地.
係去依靠佢.
去信靠佢.
反而唔係一D backgroundge 人.
係路家ge 筆下就變成.
係主角.
神居然.
係向佢地.
顯現去啟示.
你睇到呢個故事.
繼續落去我地唔睇好多細節.
就係薩加利亞本身.
係一個祭祀.
輪到佢輪班.
佢能夠入到去聖所裡面.
去公職ge 時候.
無論係點.
登台.
或者打理聖所裡面ge 事務.
輪到佢班節ge 時候.
佢就入jor 去.
聖所.
然後佢裡面就上帝向佢顯現.
就係天使向佢顯現.
同佢話比佢聽.
而家有說話.
黎到佢臨到你.
有咩咁特別.
你記得我上個禮拜講過.
係舊約完結之後.
到新約未開始.
有幾多年時間.
大概有三百幾四百年時間.
三百幾四百年時間.
gwo 段時間叫乜ye .
我地叫silent period.

$^{401}$因為舊約最尾個先知係.
Malachi.
最尾gwo 個.
Malachi.
中文係咩.
係馬拉記.
馬拉記.
完jor 之後到新約未開始ge 時候.
有大概三百多年ge 時間.
係冇有上帝ge 啟示.
甚至你睇下當時gwo 幾百年.
猶太人ge 著作.
佢地都認定上帝好似冇講ye .
上帝冇錯之前有promisegwo D ye .
但係呢幾百年呢好似靜jor 落去.
沉默ge 時期.
所以到路加福音第一章.
天使突然間顯現ge 時候.
上帝講說話.
佢就講話.
上帝講說話.
然後佢出到黎ge 時候.
其他ge 人民.
見到佢入jor 好耐.
返出黎ge 時候佢就啞jor .
講唔到ye .
換句話講上帝對佢有啟示.
對佢再有說話.
呢件係好大件事.
所以路加福音第一章.
一開始就話上帝有啟示圖.
上帝ge 計劃繼續.
開展.
繼續去advance繼續向前.
雖然有沉默gwo 時候.
但上帝而家又開始講ye .
所以你睇第一章ge 時候.
你有咩背景你就發現.
原來好大好大ge 事.
冇講ye 好耐突然間.

$^{441}$而家又再講ye .
並且係同救主降生.
係有關.
路加福音第一章.
到第二章.
我地睇呢個故事ge 時候.
佢就講咩呢.
講jor 施西若汗ge 出生.
第一章.
第二章就交代耶穌ge 出生.
Christmasge story.
第二章第四節.
約瑟.
就從加里.
ge 那沙勒城.
上猶太去.
到了大衛的城.
叫做伯利恆.
因他本是大衛一族.
一家的人.
要和他所聘之妻瑪利亞.
一同報名上閘.
那時瑪利亞ge 新人.
已經從遼.
瑪利亞ge 產期已經到.
生jor 頭胎ge 兒子.
用布包起黎.
放在馬槽裡面.
因為客店裡面冇地方.
前前後後第二章.
第一節到而家.
有好幾次講到報名上閘.
四次.
講到報名上閘.
咩回事.
當時羅馬帝國.
要求那些人.
下聖旨.
然後回到他們原居地.
回到他們的家鄉.

$^{481}$去報名上閘.
和我們今日不同.
IRS就要報稅.
三四月的時候.
他們都是一樣.
十四年左右.
就來一次要去報名上閘.
所以我們發現呢個故事.
佢地呢D ge .
小民.
約瑟瑪利亞.
本身係係拿撒勒.
然後佢地因為.
接納jor 聖旨.
聖旨落jor 之後.
佢地就要返返去自己屋企.
換句話講約瑟係咩人.
佢係伯利恆ge 人.
佢鄉下係伯利恆.
所以佢就返返去報稅.
但係羅馬帝國.
一講到報名上閘.
報稅ge 時候.
就好大件事.
報稅ge 時候就係加稅.
好似我地最近芝加哥.
我地成日話地稅加.
地稅加大唔大件事.
好大件事.
所以個個都要去投票.
想去搏返.
低少少地稅.
當時你諗下.
今日ge 地稅再加多十倍.
就係當時羅馬帝國.
要下聖旨.
每人要報名上閘.
ge policy.
好大件事.
所以呢每逢.

$^{521}$十年半載.
佢地要報名上閘ge 時候.
好多地方就會有暴動.
因為一種強制ge 加稅.
其實跟之前所謂.
即係一D 好強硬ge .
政策.
同之前.
中國.
動態清零.
差唔多都係夾硬黎.
佢要你點.
你就要點.
不相伯仲.
每逢報名上閘ge 時候.
就係一D 好政治上.
社會上好不安ge 一段時間.
各地都會有暴動.
係個小故事裡面.
瑪利亞同埋約瑟.
就趕路從納沙.
係北邊.
落到去野口山附近.
ge 百里行.
趕路去到gwo 度ge 時候.
佢地就返返去佢地祖家.
到jor 百里行.
佢話黑店沒有地方.
我地諗起.
就係airbnb.
去到gwo 度我地.
冇酒店.
但係其實要check下.
真正ge 原文.
路加.
唔係講緊酒店.
路加如果要寫酒店.
佢有另外一字可以用.
你記得有個比喻叫做好撒瑪利亞.
有個人經過.

$^{561}$gwo 個地方.
見到一個人比人打到半死.
然後佢就去照顧佢.
帶佢去邊.
黑店.
真係有路上面.
一D 所謂holiday inn或者motel 6.
咁ge ye .
佢就帶jor gwo 個比人.
打到半死ge 人.
然後就比錢你都去照顧佢.
gwo 個叫做黑店.
呢度用ge 字其實唔係.
黑店ge 意思.
其實應該係客房ge 意思.
客房.
即係話約瑟瑪利亞返返.
去自己祖家ge 時候.
當然就係搵咩.
搵返自己親戚.
下個禮拜我會.
飛加州返返去.
所謂我ge home town.
咁我地使唔使住酒店.
我地唔使住酒店.
點解.
因為我in law係gwo 度.
我應該成閘gwo 度.
換人住一晚都得.
所以我就係咁ge 緣故.
我就返返去home town.
我就唔需要租地方.
約瑟瑪利亞都係一樣.
返返去自己home town.
可能係返返去自己.
一D 親戚ge 祖屋等等.
係gwo 處就去.
投靠佢地住住先.
因為報名上冊就要排長龍.
或者.

$^{601}$有好多ge 手續要做等等.
當時ge 房屋.
好多時兩層.
兩層ge 平房.
等於好似以前香港ge 丁屋.
上面就係客房.
下面就係馬棚ge 地方.
或者廚房.
紡織ge 地方.
生產ge 地方.
如果你咁樣睇ge 時候.
客房沒有地方.
ge guest room not available.
換句話講.
有新人擁埋一齊.
佢地可能冇地方比佢住.
於是乎有新人ge 瑪利亞.
就落jor 下面.
馬棚ge 地方.
係咪好陰公.
有D 陰公.
但都唔差ge .
因為gwo D 兩房ge 設計.
下面.
有D 牛羊等等.
其實就令到gwo 個ge 屋企.
會亂D ge .
以前就冇暖氣.
所以gwo D 馬牛羊.
佢就放jor 下面.
然後暖氣就會上升.
所以下面都係ok.
去到馬棚呢個地方.
佢地就係gwo 度.
生產.
可能我諗咁想像.
大老婆都爬唔到上去.
係樓下.
佢地可能執個地方.
然之後就比佢地係gwo 度.

$^{641}$生仔.
因為咁睇ge 時候.
呢個聖誕故事.
並非好似.
唔一定係好似.
我諗馬尼亞同埋.
約瑟.
冒住風雨冇人理佢地.
搵唔到Motel 6.
然之後就去jor .
一個馬棚.
然之後生jor 耶穌出黎.
未必係咁.
可能呢個故事其實係一個.
Welcomege 故事.
即係佢地去到伯尼行呢個地方.
客房上面冇地方.
然之後主人家.
或者佢ge 親戚.
執jor 個地方.
然之後比佢地係gwo 處.
能夠安睡.
產其到日.
耶穌就出世.
如果係咁睇.
路加方第一二章.
就係講緊一個Welcomege 故事.
有D 乜ye 人.
能夠Welcomejor 耶穌.
進入佢心靈.
有D 咩人可以接受到上帝ge .
福音 接受到上帝ge .
話語.
唔係gwo D 大名鼎鼎上面ge .
而係gwo D .
真正ge 小角色.
撒加尼亞.
伊尼撒伯.
約瑟 瑪利亞.
當時社會裡面.

$^{681}$最低端ge .
最無能ge .
佢地卻係.
見到耶穌.
見到上帝ge .
工作.
每年到年尾ge 時候.
時代雜誌.
Time Magazine.
都會選Person of the Year.
或者Icon of the Year.
你知唔知.
今年佢地選jor Icon of the Year係邊個.
楊紫瓊.
楊紫瓊 Michelle Yeoh.
你知唔知.
都ok.
因為我地亞洲人.
選jor 佢為我地爭光都好.
即係楊紫瓊都幾歲.
都好老.
60歲.
佢都打jor 江山好耐.
係電影界.
係一D 動作片.
等等.
得到呢個獎.
Icon of the Year.
Person of the Year.
即係今年.
所謂風雲人物.
烏克蘭總統.
點解.
因為俄羅斯入侵烏克蘭ge 時候.
佢帶領佢ge 人民去反抗.
Keep住gwo 個Spirit.
Ukrainiange Spirit.
呢D 都係一D 好重要ge .
Icon of the Year.
Person of the Year.

$^{721}$我地中間邊個可以選到.
Icon of the Year.
有冇.
如果我地係任何一個人.
係Person of the Year 或者 Icon of the Year.
就巴閉囉係咪.
可惜唔係呀.
冇乜人理我地.
可能你ge 你ge 你最近ge 感受都係咁.
冇人理你.
社會裡面可能唔係咁多人.
重視我地.
In fact我地就係一種咁樣ge 人.
但係路加.
話接受福音ge .
上帝ge 恩典.
臨到ge 呢D 生命.
就係呢D 咁ge 小國蟲.
約瑟瑪利亞.
撒迦利亞.
伊利沙伯.
一D 冇人理ge 上帝.
就臨到佢地.
等於我地唔使去崇拜gwo D .
有權有勢ge .
相反聖誕故事.
就話比我地聽縱使我地係小角色.
上帝.
仍然係我地.
身上可以共同.
你係上帝ge 眼中.
係一個好重要ge 人.
係社會裡面.
可能係低端.
係社會經歷上面.
可能你覺得冇人理你.
但係你見到呢D .
小角色第一二章出現ge 時候.
都係一D .
默默去事奉.

$^{761}$默默去忠心.
去謹守.
冇乜地位.
冇乜影響力.
上帝ge 恩典.
居然係臨到佢地.
如果你拍電影.
你就唔會搵呢D 咁ge 人拍.
你搵咩人拍.
戴希倫.
奧古斯都.
去做佢地ge 故事.
呢D 叫做Epic Story.
我地冇Epic Story.
但聖誕故事.
老家就話.
佢係呢D 咁ge minorge character.
可以見到.
耶穌神ge 恩典.
拯救.
能夠黎到.
寂寂無名ge 係呢D 小角色.
能夠見到上帝ge 共同.
我地ge 救主.
亦都係寂寂無名ge .
佢可以見到呢個世界.
我記得好幾年前.
我聽.
德蘭修女.
你知到邊個係德蘭修女.
德蘭修女都好出名.
其實本身佢一D 名都冇.
只不過英國BBC.
有個記者.
走jor 去嘉義國塔Calcutta.
黎到去訪問佢.
訪問一播出去之後.
成個英國西方世界就知道.
原來有個咁ge 人.
就係服侍gwo D 垂死ge 人.

$^{801}$係街垂死ge 人.
佢就執佢地返去.
洗乾淨佢個身體.
以至佢能夠有尊嚴去世.
德蘭修女Margaret Theresa就做呢樣ye .
Margaret Theresa講jor 一句說話.
今日印象仍然好深.
佢講乜ye 呢.
佢話Not all of us can do great things.
But we can do small things.
With great love.
Not all of us can do great things.
But we all can do small things.
With great love.
佢就係咁直直無名.
係嘉義國塔.
做修女去服侍.
係佢身上就見到上帝ge 工作.
上帝ge 恩典.
唔係我地個個都未必可以做到偉大ge 事.
佢就係咁直直無名.
係嘉義國塔.
做修女去服侍.
係佢身上就見到上帝ge 工作.
上帝ge 恩典.
唔係我個個都未必可以做到偉大ge 事.
但我可以做小事.
With great love.
係聖誕節.
平平無奇ge 人.
平凡ge 人.
先係主角.
佢係社會裡面.
上頭ge .
有權有勢ge 人.
佢地反而係背景.
Do small things with great love.
係聖誕節裡面諗一諗.
有D 咩可以默默ge .
去繼續耕耘去做.

$^{841}$花少少時間.
可能為一D 有需要ge 人祈禱.
可能寫張問候ge 卡.
去慰問一個有需要ge 人.
一個可能孤獨ge 人.
你去幫助佢.
幾個禮拜前.
我入黎教會.
weekday 時候入黎教會.
有時我地收到信.
個個巡禮就.
D 信就放jor 係出面.
咁我就開gwo 幾封信.
有封卡黎ge .
外州寄黎.
外州會有人寄黎ge .
唔係芝加哥寄黎ge .
芝加哥唔使寄黎啦.
外州寄黎.
gwo 張卡.
我打開用英文寫.
佢寫比我地教會.
好多謝我地教會.
因為教會最近幾年.
我地去接納jor 佢ge 兒子.
一對老夫婦.
年紀老邁ge 夫婦寫信黎比我地.
好多謝你.
因為接納jor 我個仔.
我為jor 我個仔祈禱好多年.
希望佢能夠返返去上帝gwo 度.
一個腐鹿撞棍咁.
佢就黎jor 呢度教會.
成為我一份子.
睇完之後我就發現.
呢個就淨係教會應該要做ge .
唔係大新聞係咪.
但係佢係呢個人ge 生命裡面.
去做jor 一D 上帝ge 工作.
呢個就係聖誕ge 故事.

$^{881}$聖誕ge 意義.
諗一諗.
你身邊有咩ge 人.
你係需要去到關心佢.
可能一個小ge 動作.
Do small things with great love.
但係最大ge 問題就係.
到底我地有冇邀請耶穌.
進入我地ge 生命.
我地可唔可以見到耶穌.
可唔可以讓耶穌進入我地ge 生命.
呢個聖誕節我希望你唔係淨係.
過一個好忙ge 節期.
可能又要去探親.
或者去做好多唔同ge ye .
要買禮物 要包禮物等等.
要佈置等等.
Welcome 耶穌.
進入你ge 生命.
近期.
你知聖誕節好多聖誕歌.
其中一首.
我特別被提醒.
一首.
細路仔ge 歌.
聖誕節有好多細路仔唱ge 歌.
其中一首.
我最近記起.
叫做張心給主.
中文翻譯叫張心給主.
英文係.
What can I give him.
我可以比耶穌D 乜ye .
我一個咁平凡ge 人.
可以為佢貢獻D 乜ye .
奉獻D 乜ye .
張心給耶穌.
我等下會播呢首歌.
我希望呢首歌成為你ge 提醒.
成為你ge 禱告.

$^{921}$記住乜ye .
Not all of us can do great things.
But we can do.
We can all do small things.
With great love.
但係要讓耶穌.
進入我地ge 生命先.
要讓我地ge 生命裡面有空間.
俾耶穌作我地ge 主.
呢首歌.
\newpage



\section{}
\label{sec:LoAxQ1cr1f0}
\textbf{A Child is Born - Rev. Mark Chan  有一嬰孩 - 陳崇基牧師}
\newline
\newline
連結: \href{https://youtube.com/watch?v=LoAxQ1cr1f0}{\texttt{https://youtube.com/watch?v=LoAxQ1cr1f0}} ~~~~ 語音日期: 2022-12-26
\newline
\newline
\hyperref[sec:du2XrTgHsNA]{\small{< < < PREV SERMON < < <}}
~
\hyperref[sec:index_chronic]{\small{[返順時目]}}
~
\hyperref[sec:index_scriptual]{\small{[返順卷目]}}
~
\hyperref[sec:lBh3rgPDiTo]{\small{> > > NEXT SERMON > > >}}
\newline
\newline
$^{1}$(speaking in foreign language).
In Christ's name we pray..
- Well, I wanna say Merry Christmas to all of you..
And I think like many of you,.
Christmas can be a crazy time..
It can be a very busy time..
And this Christmas season,.
I'm thinking of fathers, mothers,.
hardworking people in the restaurants.
who are running, operating in our society,.
making sure our loved ones are being taken care of..
And today I know many of you who are part of our church,.
you did your best to make this church a better place..
Especially today, those who are in the backstage,.
those who prepare the table,.
those who prepare the good food that we're gonna enjoy,.
and the people and the personnel on stage,.
the musicians and all the helpers..
I know this is a crazy, busy time,.
but it's also a time for us to slow down,.
to think about what this season is all about,.
what Christmas is all about..
I also think of many of you,.
or even many of us who are joining us online.
and cannot be on site..
I'm thinking of you because I know you're going through.
some crisis in your life,.
some troubles, some difficulties..
And you are one of those taking care of the sick,.
the dying in your family..
And know that God is with you..
This is a challenging time..
Maybe you are dealing with broken relationships..
Maybe you're dealing with a lot of uncertainty..
And many of us, this holiday season is not easy..
It's a tough season..
But this is the time we should turn to God..
We should turn our eyes to heaven.
because Christmas is a special time.
when we celebrate the birth of Jesus Christ,.

$^{41}$someone who is special..
In Jesus Christ, we have everlasting joy,.
not happiness..
We have everlasting peace, not temporary peace..
In Jesus Christ, we have hope, hope that is everlasting..
I know you probably listen to a lot of songs.
about Christmas on radio or on Spotify..
If you pay attention to some of these Christmas songs,.
Christmas hymns, some of them are hundreds of years old,.
very, very old..
And you listen to the lyrics..
All of these Christmas songs.
speak about the most basic human desire,.
the desire deep down in our heart..
We yearn for peace in our time of chaos..
We yearn for joy in our time of sadness..
We yearn for someone to be with us, and we feel lonely..
We yearn for hope when we see there's no end in the tunnel..
It is no wonder that in this Christmas season,.
the rate of suicide and depression increase..
It's no wonder..
This is a very depressing time for many of us..
But yes, even though we are in a difficult time,.
even though we see a lot of oppressions, injustice,.
darkness in our world,.
we are reminded on this Christmas Sunday.
that a newborn King was given to all of us..
(speaking in foreign language).
So in this Christmas,.
we wanna turn our attention to this child..
The kid, the infant was born for all of us..
(speaking in foreign language).
We're gonna ask the children's choir to come up here.
and they will help us to celebrate Christmas.
by singing two very old Christmas hymns..
(speaking in foreign language).
Well, let's give them a big hand again..
(audience applauding).
The first Christmas was recorded in the Bible,.
in the book of Matthew, Matthew chapter one..

$^{81}$This is how Jesus the Messiah was born..
His mother Mary was engaged to be married to Joseph..
But before the marriage took place,.
while she was still a virgin,.
she became pregnant through the power of the Holy Spirit..
And Joseph, to whom she was engaged, was a righteous man.
and did not want to disgrace her publicly..
So he decided to break the engagement quietly..
As Joseph considered this,.
an angel of the Lord appeared to him in a dream..
"Joseph, son of David," the angel said,.
"do not be afraid to take Mary as your wife.
because the child within her.
was conceived by the Holy Spirit..
And she will have a son, and you are to name him Jesus,.
for he will save his people from their sins.".
(speaking in foreign language).
Now, as you read this portion of the story of Christmas,.
this is a profound mystery,.
perhaps the most profound mystery of all..
God himself incarnate came to our world,.
fully divine, but fully human..
There's something beyond our rationality can explain..
This is the most profound truth in the Bible..
And when Joseph received the news,.
God is saying to Joseph,.
"My plan is different from your plan..
My will is higher than your will.".
Yes, for Joseph, if you can imagine,.
he was trying to build his own simple family, right?.
Have a good life, have a small family in Israel.
to farm or to be a carpenter,.
just want to live a normal life..
But God said to him, and to Mary, to this holy family,.
he said, "I want to send the Savior through your family.".
Now, this is also profound,.
because at times you see in this story,.
God will act in a way that is different from your plan..
You may have your own plan,.
and you may have thought your plan.

$^{121}$is the most perfect plan of all,.
but God shows up..
Christmas is telling us that God is close to us.
and God can do his work..
Something that is higher than your will,.
something that is different than your plan..
But again and again, God's way is different..
Though it's different, it is the best..
So Joseph obeyed God and took Mary,.
who was conceived by the power of God,.
and son was born..
And because of their obedience,.
God's saving plan for all of us sinners is fulfilled..
(speaking in foreign language).
the Lord. (speaking in foreign language).
(speaking in foreign language).
Now if you read this scripture again in 121,.
in the book of Matthew, it says,.
she will have a son and you are to name him Jesus..
Right?.
Name him Jesus..
His name will be Jesus..
(speaking in foreign language).
Now what does it mean?.
What does the name Jesus mean?.
Now in Hebrew, Jesus will sound like this..
Yeah, sure..
In Greek, it's Yesu..
In Mandarin, it's Yesu, right?.
In Cantonese, it's Yesu, right?.
(speaking in foreign language).
(audience laughing).
The name Jesus, what does it mean?.
It means saving, to rescue, to redeem..
That's why we call him savior..
Because the name Jesus, Yesu,.
it means to save us from something..
And the story says, Jesus is gonna save us from sin,.
from death, from the power of darkness..
Now as I said before, I know many of us.

$^{161}$are facing different challenges.
and seeing a lot of different forms.
of injustice and oppression..
In the time of Jesus, many people thought that way too,.
because they were under the shadow of the Roman Empire..
And the Roman Empire was an oppressive regime..
And today, many of us,.
certainly we see some situations similar.
to what they had before..
We see a situation that is oppressive.
in many other countries, or even in our society..
Because we see all the brokenness in our world..
We see sin and darkness..
Some of the people think if we have the right policy,.
if we have the right precedent,.
if we have the right political power,.
all of these issues will be resolved..
But the Bible cuts through all of these answers.
and telling us the root problem is sinfulness..
Is that we are sinful, all of us, living in darkness..
And we need a savior..
And only Jesus himself can save us.
from that kind of darkness..
Because until the power of sin is broken,.
all changes are superficial..
Only when God came to our world in our heart.
and start to transform us,.
then our relationships can be reconciled..
Our world can be changed..
And that's why we celebrate Christmas..
And perhaps the story of first Christmas.
remains the most profound mystery..
Because on the first Christmas,.
this baby, the newborn king, was born,.
fully human and fully divine..
And he is the bridge between us and God himself..
Christmas shows us God is near..
Now the year of 2022 is coming to a close..
Perhaps this year is tough and challenging for you..
It's very critical for you..

$^{201}$And my prayer is that we will all turn to God,.
knowing that only in Jesus, our newborn king,.
we can have true peace, true joy, and true hope..
So as we continue to worship, we will light the candle..
And as we light the candle,.
as we pass the light from one to the next neighbor,.
and then on and on, as we light up the room,.
I want you to say a prayer..
Say a prayer, first of all, for yourself..
Maybe you are facing some difficulties..
Maybe you are facing some issues.
that is beyond your ability to solve..
Say a prayer for yourself..
Let God be your God..
Then also say a second prayer.
for someone you know who is needy in this season..
Maybe a family member, maybe a friend or relative,.
maybe a neighbor..
Say a prayer for them..
May God be with them..
(speaking in foreign language).
So as we continue to worship, we will light the candle.
and we'll pass the light from one to the other..
I know you have some candles on your table.
and you can get that prepared..
Make sure you get the tray ready..
This is our tradition at UV in our church..
On Christmas Sunday, we will light the candle.
and pass the light..
And we will light up the room together..
And this symbolizes the newborn king who has come,.
who has come to save us..
He dispels all kinds of darkness.
because light always conquers darkness..
(speaking in foreign language).
♪ O holy night ♪.
♪ The stars are brightly shining ♪.
♪ It is the night of our dear savior's birth ♪.
♪ Long lay the world in sin and error pining ♪.
♪ Till he appeared and the soul felt its worth ♪.

$^{241}$♪ A thrill of hope ♪.
♪ The weary world rejoices ♪.
♪ For yonder breaks a new and glorious morn ♪.
♪ Fall on your knees ♪.
♪ Oh hear the angel voices ♪.
♪ Amen ♪.
♪ Divine ♪.
♪ O night when Christ was born ♪.
♪ O night divine ♪.
♪ O night, O night divine ♪.
♪ O night divine ♪.
♪ Truly he talks to one another ♪.
♪ His law is love and his gospel is peace ♪.
♪ For the slave is our God and in his name ♪.
♪ All our precious shall sing ♪.
♪ Sweet hymns of joy ♪.
♪ In grateful chorus raise we ♪.
♪ Let all we believe us praise his holy name ♪.
♪ Christ is the Lord ♪.
♪ Our praise he sings forever ♪.
♪ His power and glory evermore proclaim ♪.
♪ Glory evermore proclaim ♪.
♪ Evermore proclaim ♪.
(gentle music).
♪ Praise his holy name ♪.
♪ Christ is the Lord ♪.
♪ Our praise his name forever ♪.
♪ His power and glory evermore proclaim ♪.
♪ Glory evermore proclaim ♪.
- 1 John chapter four, for God is love..
God show us how much he love us.
by sending his one and only son into the world.
so that we might have eternal life through him..
This is real love, not that we love God,.
but that he loved us and sent his son as a sacrifice.
to take away our sins..
(speaking in foreign language).
This is love..
(speaking in foreign language).
There's no wonder why Christmas is a time.

$^{281}$to give gifts to each other,.
to exchange gifts with our friends..
I know we will have a game of white oven later,.
but let's hold onto this.
because I believe it will be worthwhile.
for us to ask ourselves the question..
Now, if you have the opportunity to ask God.
for a gift in this season, what would that be?.
If you pray to God, you have the opportunity.
to grant you one wish, one prayer, what would that be?.
What would that be considered a gift for you in this season?.
Perhaps you are thinking of a lot of different things,.
perhaps a better career, a better future,.
a better relationship, a better house, a better car,.
or whatever that is better than what you have..
And I believe that God can certainly grant.
these good things to all of us..
But here the Bible says,.
"God has already granted us the greatest gift.".
And this greatest gift is what we are celebrating today..
It's the one and only Son, Jesus Christ,.
His salvation, His one and only Son.
who will grant us eternal life, life indeed..
But this scripture tells us why does God do this?.
Why did God give us His one and only Son?.
Because of His love..
The greatest gift is given to us because of His love,.
because it is who God is and giving us the best,.
the best gift of all..
(speaking in foreign language).
\newpage



\section{}
\label{sec:lBh3rgPDiTo}
\textbf{6/2/19 主日講道\_耶穌天國的小故事-4\_寶貝與珍珠(陳崇基牧師)}
\newline
\newline
連結: \href{https://youtube.com/watch?v=lBh3rgPDiTo}{\texttt{https://youtube.com/watch?v=lBh3rgPDiTo}} ~~~~ 語音日期: 2023-01-02
\newline
\newline
\hyperref[sec:LoAxQ1cr1f0]{\small{< < < PREV SERMON < < <}}
~
\hyperref[sec:index_chronic]{\small{[返順時目]}}
~
\hyperref[sec:index_scriptual]{\small{[返順卷目]}}
~
\hyperref[sec:j1mkZMKAJog]{\small{> > > NEXT SERMON > > >}}
\newline
\newline
$^{1}$我們有聖經的話,我們回到聖經,好嗎?.
《馬太福音》第十三章.
我們今天看四節經文.
十三章四十四節.
《馬太福音》第十三章.
四十四節.
他說:「天國好像寶貝藏在地裡,人遇見了就把它藏起來,.
歡歡喜喜地變賣一切所有的,賣這塊田地..
天國又好像買賣人尋找好珠子,遇見一顆重價的珠子,.
就去變賣他一切所有的,買了這顆珠子.」.
我們繼續看比喻,耶穌基督說的比喻.
你要知道有很多的比喻,是嗎?.
有部分的比喻,耶穌特別說完之後,他會解釋的.
比如說殺種的比喻.
他就說每一樣的泥土代表什麼人,是嗎?.
或者有一個比喻叫做麥子,拜子的比喻.
就是麥田裡面有一些好像麥子的,好像weed的.
那些拜子出來的時候,很像樣的.
那個比喻他也有作解釋,也有一些比喻是他故意不給解釋的.
換句話說,那種比喻是有一點啟開性的,就是open end.
就是叫世世代代的人來聽這個比喻的時候.
特別是我們今天聽這個比喻的時候.
信息沒有變大,但是形容就有些不同.
所以今天說這個比喻,可能信息很清楚.
但是你拿走的東西,可能提醒你的東西是有些不同.
有人說信仰宗教是要安慰人的心靈的.
受傷者的心靈.
但是同樣也是折磨那些自以為安穩的.
這句話有沒有道理?.
耶穌說的比喻很多時候就是這樣.
不是一定你聽完比喻之後,你覺得很安心就走.
可能你自覺很安心,很安穩,很行的時候.
他就要去折磨你,或者不一定是折磨.
就是去挑釁你.
所以你看聖經裡面的一些比喻的時候.
你要留意這一點.
特別是你平時看聖經.
到底是為了純粹得到更加安心.
還是聖經在挑戰我們.
這裡很簡單,兩個比喻.

$^{41}$說到寶貝的比喻.
這些藏寶的故事,經常每個人都會.
歷世歷代的新聞都會說.
我講個例子給你聽.
幾年前在英國有三個學生.
他們去救世軍那裡買了舊梳化,二十多塊.
搬回家.
誰知道他們擺放的時候,清潔的時候.
就拿了很多信封出來.
在梳化底那裡.
一邊開一邊開的時候,原來是現金.
結果就有四萬塊現金在梳化裡面.
二十塊就換了四萬塊.
划算嗎?.
很划算吧,下次你去買梳化.
你記住看清楚.
特別是買舊梳化的時候.
原來最後拿到的信封.
就有那個物主的名字.
原來是一個九十歲的婆婆.
她坐了那麼多年,就放了很多現金在裡面.
這些寶物,藏寶物的故事.
應該是歷世歷代的人都很熟悉.
特別我們華人,最有錢的.
或者最有價值的東西在哪裡?.
在保險箱,藏在下面.
小心一點.
所謂保險箱,保險箱放一放,買一買.
不知道是什麼.
你知道最重要的東西當然在那裡.
這個比喻有一點這樣的情況發生.
相信也是當時一些很常見的現象.
有一個農夫在耕田或者在田地裡面工作的時候.
他找到一些寶物.
古今中外都是這樣的.
特別是動盪的地方.
很多時候有戰爭,有不安的地方.
人們就把那些東西放在自己的田裡.
掘個洞放在那裡,然後賣掉.
希望動盪過了之後,他就拿出來.

$^{81}$如果他不幸死了,又沒有交低的話.
那些東西就不見了.
所以後人可能會找出來.
這裡就是一個這樣的故事.
大概這個人按照當時的經濟狀況來說.
社會的狀況來說.
這個人應該是一個租田地的農夫.
就是一個窮的農夫.
自己買不到地方,那怎麼辦?.
就租了,租了塊田去耕種.
然後收成的出產有部分歸自己.
大部分歸業主.
就是這樣的情況.
買不到就租,你租的地方還是買的地方住?.
你還沒準備買的,就租吧.
租一個好的地方.
記住,租的地方要想清楚.
可能有些寶藏在裡面.
要花錢找找.
大概是一個窮的農夫.
有一天耕田的時候.
產一插進去,原來是硬的東西.
原來是有寶藏的,就把它收起來.
發現原來那塊地是有寶藏的.
就是說業主不知道.
業主知道怎麼會讓你找到.
業主買那塊地的時候他也不知道.
到租田的農夫.
突然間偶爾發現了的時候.
他就邊賣他所有的東西.
賣光他自己擁有的東西.
或者借錢.
他就為了要得到這塊田.
就是一個這樣的比喻.
下面那個比喻也很簡單.
就是講到不是一個窮鬼.
上面就講一個窮鬼.
下面就講一個商人.
他本身是周遊列國的.
去找一些價值連城的東西.

$^{121}$找珍珠.
珍珠你有沒有?.
說不定你帶了幾顆.
你知道珍珠是怎麼形成的嗎?.
就是有些蜆走進去.
然後就滲出一些東西.
然後就包了很多年.
形成了一顆珍珠.
現在有的珍珠很多時候是人造的.
人養的蜆做出來的珍珠.
以前一定是wild的.
一定是很難很難才找到.
所以是價值連城的東西.
有如今天講鑽石那樣.
周遊列國去找一顆很名貴的珍珠.
終於被他找到了.
然後他的回應就是變賣他其他的東西.
為了要得到價值連城的珍珠.
兩個比喻都是在講value.
就是人生裡面的價值.
你想想你人生裡面最有價值的是什麼?.
人生世界最有價值的是什麼?.
有些人說自由.
有些人說健康.
有些人說是關係.
可能都是.
這些都是人生裡面我們很寶貴的價值.
六四週年來到了嗎?.
六四三十週年.
大家記得六四三十週年.
今天六月幾號?.
二號是,兩天後.
六四三十週年將近.
有人說自由真的無價寶.
我看了一個六四死難者母親的紀錄片.
發生這樣的慘劇.
這樣的屠殺.
拆散了家庭.
然後留下了母親.
你想想母親多傷心.

$^{161}$社會的國家如果沒有自由.
沒有人權,沒有基本的人權的話.
她會很可憐.
所以這些都是人生裡面值得去想的.
值得去追求的價值.
你想想價值是什麼?.
人生裡面有什麼是值得你追求的?.
最近香港有修訂逃犯條例.
到處都有聯署.
在香港有很多不同的地方.
不同的群體.
當然比較少教會出聲.
但一些自己發起的群體.
自己發起的無論是平民百姓.
或者一些紀錄.
聯署聲明去反對這個法例.
一個惡法.
言論自由,新聞自由.
這些是不是我們人生裡面要追求的核心價值?.
Value.
你說到基督徒.
我們是擁抱天國的一班人.
這些價值是不是應該是我們的價值?.
還是我們只是純粹.
我們有上帝的救恩,我們等上天堂.
我們是很安穩的一班人.
到底什麼才是值得我們追求的價值?.
如果這些都是天國的價值.
換句話說,也是天國裡面有的永恆的價值.
兩個比喻就是想去挑戰我們.
有沒有想清楚.
我們人生追求的方向.
是一些真正值得我們追求的價值?.
兩個比喻裡面.
也有說到那個找到寶藏的人.
他有什麼反應?.
那個窮鬼農夫.
是歡歡喜喜的.
將他所有的東西變賣.
然後買那塊地回來.

$^{201}$可能一個更好的反應.
應該不是joy,不是歡歡喜喜.
上文下面應該是驚喜.
Surprise.
因為他沒有想過會有這樣的寶藏在地上.
但下面那個商人.
他是走遍全世界去找名貴的.
所謂的名珠.
名珠的一粒.
他就沒有說驚喜.
因為有什麼可以驚喜呢?.
你明知道有一顆火的名珠存在於世界裡面.
你千山萬水去找它.
到你真的找到的時候.
你就說我終於得到了.
就沒有驚喜了.
所以沒有說他心情上的反應.
因為說到擁有天國的人.
有兩種.
有一些是偶爾遇上天國的.
偶爾在人生裡面.
福音給他帶來驚喜.
沒有想過有人探望我.
跟我說福音我會相信.
沒有想過就是那天抽起一根筋.
我就可能去了教會.
我就可能聽到福音.
聽到主耶穌基督的救恩.
救恩就臨到我了.
從人的角度看好像很幸運.
好像很幸運.
我最近這兩個星期看了一個中東的新聞.
說伊朗.
當然伊朗最近都很緊張.
他說伊朗本身是一個回教徒的國家.
一個穆斯林的國家.
每個人都是穆斯林.
一輩子都是做回教徒.
那些人覺得.
開始有些人.

$^{241}$因為在一些動盪裡面.
他認為這個伊斯蘭教.
可能真的不是真的.
可能不是真的.
但是居然有一班人.
他去訪問他們的時候.
他們就是偶爾地遇上基督教的信仰.
你想想在伊朗.
在伊朗一個這樣的地方.
一個福音硬土的地方.
沒有人會去建立教會的地方.
就是一個女孩子.
她在心靈裡面覺得欠缺一些東西.
而伊斯蘭教是不能給她的.
她就突然間遇上一本新約聖經.
New Testament.
她就打開這本新約聖經.
就看到耶穌的故事.
她就說她Fall in love with Jesus.
看到耶穌的故事.
看到耶穌的如何和別人相處.
如何和別人回應.
說什麼道理的時候.
她一讀之下就放不下手.
就說這個耶穌真是一個世界的救主.
她就遇上了福音.
這是歡歡喜喜地.
Surprise.
也有另一種人擁抱天國的時候.
不是Surprise.
而是她的人生裡面.
她是想四海尋覓.
到底人生真正有沒有永恆的價值.
終於被她找到了.
她就立志跟隨耶穌到底.
我認識一個朋友.
有一次跟他聊過如何信主.
他就告訴我.
她是一個女孩.
一個姐妹.

$^{281}$她年輕的時候.
每天晚上都會去派對.
特別是星期六晚上就去派對.
大家都覺得心裡有些事情不舒服.
到底是為了什麼呢?.
沒有什麼滿足.
突然有一個朋友就邀請她去教會.
她就去教會check out.
她覺得去主教崇拜.
星期天早上去教會.
很悶.
牧師說話很悶.
但她又繼續去.
星期六晚上又要派對.
又要喝酒.
兩三點才回家.
第二天她又急起來.
又去教會.
又覺得很悶.
但她又繼續去.
去了一年多的時間.
一直沒有間斷的.
去一年多的時間.
然後突然間終於她明白.
這就是她一直在尋找.
在派對中找不到的東西.
就是救主耶穌基督.
千山萬水去尋找真理.
然後她就跟隨耶穌.
順了主.
長長一年的時間.
很難過.
因為你要去派對.
然後又要去教會.
去教會又很累.
牧師說話很悶.
然後又要睡覺.
完了之後又不知道得到什麼.
但她就每個星期這樣去.
結果她的生命就改變了.

$^{321}$兩種反應.
兩種擁抱天國.
進入天國的人.
一種是驚喜.
一種是四處尋找.
終於找到了.
比喻道出了兩種人.
也說到兩種人找到天國的時候.
他有什麼回應.
有什麼行動.
這個比喻只有兩句說話.
第一句,天國好長.
第二個比喻也是.
天國好長.
第二句就是回應.
說到那個反應.
不論你怎樣遇上天國.
無論你怎樣遇到耶穌.
你是驚喜還是.
終於找到那種.
反應都是一樣.
變賣所有的.
犧牲一切的.
因為他找到了.
他就想盡辦法去擁有.
對身邊的人來說.
你想想那個窮鬼農夫.
對身邊的人來說.
因為他找到寶藏.
他不會到處告訴別人.
別人問你幹嘛.
你借這麼多洞.
為了買那塊爛地.
你買都不買那裡.
你買downtown.
你買那些會升值的地方.
但他有一個秘密.
就是不能告訴別人.
這個會升值.
對其他人來說.

$^{361}$是一個很大的風險.
對其他人來說.
是不太合理的.
不說農夫.
說商人.
中文翻譯日本.
賣賣的人.
買賣的人.
Emporium.
Emporium你記得這個字.
如果你懂一點英文的話.
以前那些大百貨公司叫做Emporium.
就是說那些.
host seller.
就是批發商.
批發商有什麼特色呢.
就是賣不同貨品.
因為你不想將所有蛋都放在同一個basket.
這個是.
不是很聰明的投資.
沒理由為了一樣東西而去盡.
你當然將蛋放在不同的basket.
有一樣東西倒塌也好.
你的portfolio也保得住.
這個就是Emporium.
買賣的人.
merchant.
商人.
就好像你沒見過costco只是買一樣東西.
很多百樣東西.
但是他找到這一顆珍珠的時候.
這個商人.
連身份也變了.
就是不做批發商.
他只是為了這一顆珍珠.
對其他人來說.
是不太合理的.
因為這個孤注一擲.
將所有雞蛋放在一個籃子裡.
那你覺得.

$^{401}$這個商人.
他有什麼特色呢.
他有什麼膽的論.
農夫的比喻.
商人的比喻.
到底想說什麼.
肯定不是他想說.
告訴我們.
我們可以做一個不誠實的人.
我們應該路不拾遺.
找到寶藏.
然後收起來.
好像很貪心.
不是這樣的意思.
你知道比喻不是每樣東西都重要.
不是每個細節都重要.
也不是每個細節都想.
導出那個信息.
因為肯定耶穌不是說.
我們應該像這些人.
在你家裡.
找到寶藏.
不告訴你.
趕你走 買了你的地方.
然後我發達 你沒法發達.
他不是叫我們貪心.
不是叫我們不誠實去做人.
你不可以強迫每個細節都有.
那個訊息在那裡.
那個訊息是什麼.
得到寶藏之後.
到底你的行動是什麼.
我說過古今中外.
都有尋寶的故事.
古代也有.
聖經的時代.
耶穌的時代 也有很多尋寶的故事.
但是你看回.
這麼多尋寶的故事.
有一件事很重要.

$^{441}$就是逢講到.
有個人尋到寶.
都很強調.
他尋到寶之後.
他的生活.
是多麼好.
買車 買房子 奢華生活.
有賺到大錢.
那當然是強調.
他之後的生活有多奢華.
每個尋寶的故事都是這樣.
Treasure hunting stories.
今天的也是.
現代有人中了Lottery.
新聞報導.
新聞報導什麼.
就是他會買什麼.
怎樣買一間漂亮的房子.
怎樣用那些錢.
偏偏耶穌就講.
尋寶的故事.
一句都沒有講.
一句都沒有提到.
他之後有什麼奢華的生活.
著眼的地方是什麼.
是那個人找到寶藏的時候.
他.
作犧牲.
不惜代價.
去拿到它.
變賣一切的.
為了得著.
一個寶藏.
變賣一切.
這句說話.
經常有.
馬太福音19章.
你記得那個故事.
上下文講什麼.
有一個很年輕的.

$^{481}$很有錢的.
我們叫他二世祖.
他跟耶穌講.
老師.
我可以做什麼可以得到榮心.
是不是.
對答了之後.
耶穌就跟他說.
你還差一件事.
如果你想做完全人的話.
做一個不是完美的人.
意思是成熟的人.
做一個成熟的.
懂得選擇的人.
可以去變賣你所有的.
分給窮人.
就必得.
財寶在天上.
你還要你都去跟隨我.
變賣你所有的.
這句說話是代表什麼.
不是叫我們做窮鬼.
不是這件事.
而是要我們重新編排我們生活裡面的次序.
生活裡面的價值.
是有優與先的.
也有次要的.
要我們重新再reorder.
我們的priority.
變賣一切的.
感覺就是清盤.
意思是total reordering of your priorities.
你將你生命裡面的priorities.
全部再reorder.
因為你找到天國.
因為你得到了天國.
而這個應該是你不惜一切.
你要努力.
要去得到守住的東西.
Nothing stand in the way.

$^{521}$沒有一件事應該去阻止你.
或者減慢你.
或者令你分心.
變賣一切.
對這個財主來說.
這個少年的二世祖來說.
錢他的財產.
就是他的神.
這個神就阻住他.
你都去找耶穌.
耶穌說你要變賣你一切.
分給他.
財寶在天上.
你要跟隨我.
reorder你自己的priority.
因為你找到的東西.
是有永恆價值的.
我們剛才報告的時候.
不是說有護教學.
問道點答.
在護教學的文獻裡面.
基督教的護教學.
有二千年的時間.
有一個說法.
經常都出現.
就是有一個護教學的.
不可以叫它argument.
有一個護教學的.
比喻.
可以這麼說.
經常都會出現.
那個叫做Pascal's Wager.
這個叫做.
法式賭注.
Pascal 大家知道.
是一個很出名的十七世紀的.
哲學家.
你學電腦就知道.
以前有一種language叫Pascal language.
就是以他命名的.

$^{561}$Pascal 賭注是什麼意思呢.
Pascal 寫了很多東西.
是和基督教有關的.
他講了一句說話.
或者寫了一段說話.
他說.
假設在生命裡面.
只有兩個選擇.
第一個選擇就是神不存在.
就是沒有永恆.
就是神存在是永恆.
另一個選擇就是神不存在.
就是沒有永恆.
只有今世.
第一個選擇就是神存在.
沒有永恆沒有天國.
那就做一個wager.
就是說.
所謂下一個賭注.
你信上帝的話.
你信神的話.
你就有永恆的祝福.
是不是.
有永恆的東西.
但如果真的沒有神的話.
ok.
Pascal 的意思就是說.
損失其實都很少.
但是如果有神的話.
你所得到的.
就是永遠的東西.
不能夠估計的價值.
是priceless的.
所以他的意思就是說.
如果一個有理性的人.
就一定應該去選擇.
相信上帝.
當然這個不是一些.
airtight的.
護教學的論證.

$^{601}$我絕對不贊成我們這樣去護教.
因為你這樣講好像.
我信仰就是一個賭注.
買大小.
買魚蝦蟹.
信仰不是這樣當然.
但你要知道那個context.
為什麼他講這件事.
Pascal 去address的人.
我想我理解的.
他的背景就是.
他想address那些.
nominal questions.
就是說從小到大都在教會長大.
覺得天國這些沒什麼.
他想address的那班人.
就是他想要的.
他想address的那班人.
就是這種人.
這個不是一個邏輯的論證.
你知道不是風物的.
很神秘的一個邏輯論證.
但他想做什麼呢.
他想做一點遊說.
是一些emotional persuasion.
就是一些.
心理上的遊說.
你想想.
每逢我們傳福音的時候.
我們practically.
實際上.
其實是邀請對方.
邀請你暫時.
放棄你對人生的看法.
你暫時.
站在我這一邊.
預設一下.
基督教聖經裡面的信仰的世界觀.
到底是不是make sense.
你考慮一下.

$^{641}$聖經裡面的信仰大體上是不是.
合情合理.
放棄你自己那種本身有的觀念.
你站在我這一方.
傳福音談道.
很多時候都是做這件事.
就是你希望persuade到對方.
你明白.
人生.
應該好像聖經裡面的信仰.
告訴我們.
這樣去檢視人生.
人生應該是這樣的.
其實耶穌說很多東西都是這樣的.
只有兩個選擇.
你不要想有其他選擇.
沒得想的.
因為耶穌說什麼.
祂說你只有兩個選項.
就是永死.
是一種游說的方法.
我想Pascal的焦點是一樣.
是說大家對談的時候.
那種dynamics.
不是用理性去克服對方.
很多時候不是用理性.
講贏一個人信主.
而是那個人願意.
OK我放棄我對我自己本有的人生的觀念.
我們來看看聖經.
講人生.
聖經講世界.
是有什麼不同.
是不是合情合理.
如果合情合理的.
我們應該相信.
耶穌經常說只有兩道門.
一道是窄門.
一道是寬闊的門.
窄門小門.

$^{681}$不是每個人都喜歡進去.
但那道門是通往永生的.
那闊門是通往永死.
兩個選擇.
binary.
就是逼我們.
游說我們.
要選擇一個對的選擇.
要信服神.
還是信服人.
信服上帝.
真正在天的掌權者.
還是信服地上的掌權者.
那個政權.
你選擇.
誰合理.
你將你的房屋.
將你的生命建在盤石上.
還是建在沙土上.
都是一個binary的選擇.
不需要太多東西去混亂你.
就是希望你選擇.
一個永恆的選擇.
其實這個比喻也是一樣.
告訴我們有一樣東西.
是有永恆價值的.
無論你剛剛遇見.
剛剛遇見.
幸運.
還是你千山萬水找到.
今天終於找到了.
終於聽到原來耶穌基督對我的愛.
對我的死.
在十字架上的救恩.
去赦免我的罪.
給我有永生.
為我復活.
給我有永生的盼望.
咦 我找到了.
這個就是你應該要做的選擇.

$^{721}$binary option.
這個就是Pascal Wager的意思.
問題就是你想想.
你自己.
要聽清楚聖經裡面所說的東西.
耶穌基督說的那兩個比喻.
最有價值 最有永恆.
意義的東西.
而你選擇的時候.
是不是一個最合理的選項.
如果你不選擇.
你的損失將會非常之大.
Alright.
耶穌的persuasion其實就是.
做這件事.
我不知道馬太.
馬太福音 馬太社.
馬太是其中一個.
耶穌的門徒.
我不知道他是不是因為做稅利.
所以他每件事都很公公正正的.
每件事都很平衡的.
你知道馬太福音.
是五個部份組成的.
五個部份.
都是說到耶穌的教導.
第一個部份說到耶穌的倫理.
天國的倫理.
登山寶訓等等.
八福.
然後他就說到天國的宣講.
就是他差派門徒出去.
去教他們.
傳福音.
第三個部份就是說到天國的比喻.
第四個部份說到教會.
天國的群體.
第五個部份說到天國的將來.
將來會有審判 你準備好了嗎.
我們剛才看的比喻.

$^{761}$就是在第十三章.
就是剛剛在中間.
整本書的中間.
在這裡他就將七個的比喻.
混在一起.
你不要理會.
有些人在估計.
這個比喻本身在哪個故事裡面出現.
只有耶穌才會說出這樣的比喻.
總之.
馬太就將這七個比喻混在一起.
在這裡說.
比喻再加比喻.
再加比喻 問我們七次.
說七次的比喻給你聽.
就是要你去檢視一下.
到底你在天國裡面是甚麼位置.
你想清楚.
首先他說的殺種的比喻.
其實就是說到我們.
到底我們是不是一個好的土壤.
麥田拜子的比喻.
還是一個好的種子.
是一個好的麥子.
還是一個壞的麥子.
壞的幣子.
第三個比喻.
就是說到天國裡面的能力.
第四個也是.
這兩個我們都說過.
皆菜種的比喻 很小顆.
變成很大.
麵梟的比喻.
很少 藏起來了.
但全團發酵.
就是隱藏的天秘的能力.
當人們散佈在社會.
在世界裡面不同的地方.
他們講出公義 為真理發聲.
去表達上帝.

$^{801}$給他們生命的時候.
他們那種能力在哪裡.
你是不是這樣.
發現寶藏和尋見珍珠.
這個比喻.
就是說到.
如果你知道這是永恆的價值的時候.
你有沒有這樣的慧眼.
你會決心去追究.
變賣你一切.
你所有的優先.
去追究天國裡面的價值.
最後拉網上.
分魚 有些潔淨 有些不潔淨.
就是警惕我們自己.
到底我們是哪一種魚.
但你要留意.
在這十個比喻.
完結之後.
馬太又很工整 每次都加一個story.
而story是.
連接之前所講的東西.
天國的傳道.
然後教導.
然後神蹟.
就有天國的宣講.
然後耶穌和.
施洗約翰.
有個debate.
他是否fulfill所有舊約裡面.
先知的應許.
然後天國就到了 就講七個比喻給你聽.
下面的故事講什麼.
就是講到耶穌啟示.
告訴你 祂真正的身份是誰.
祂就是永生神的意志.
你不要.
這麼輕率去看耶穌.
耶穌是永生神的意志.
所以第十六章下面的故事.

$^{841}$講完所有的比喻.
其中一個故事你記得了 在蓋撒尼亞肥亞比.
他就帶著門徒.
他就問門徒.
到底我是誰 你覺得我是誰.
你認定我是誰嗎.
他說 你是基督.
你是永生神的意志.
在你那裡有永恆的救恩.
這個救恩是無價的.
你所成就的天國是無價的.
將來的天堂是無價的.
而你今天天國的生命.
進入我的生命裡面.
開始和我同行的時候.
這個都是無價的.
是一個value的statement.
如果你認定耶穌是基督.
是理賽華.
是上帝的救世主.
去救我們離開罪和死亡.
這就是一個榮耀.
無價永遠的救恩.
沒有一樣東西可以比較.
你的價值觀.
我的價值觀.
就是應該要這樣去重整.
將一些真正有永恆價值的東西.
放在我們的第一.
或者可以這樣問.
如果你.
你認定耶穌是基督.
你覺得他是基督.
如果你家裡大火.
你只有三分鐘的時間.
你進入火場裡面.
拿回你的東西.
你那三分鐘會拿回什麼.
你想想你會拿回什麼.
錢吧.

$^{881}$可能你床底下的現金.
你需要拿出來.
到底想救些什麼.
你心靈裡面的價值是些什麼.
耶穌都問這個問題.
阿根也是在這個故事裡面問的.
講完所有的比喻.
天外的比喻.
他問什麼.
他說有人跟隨我就要寫幾杯在他們的十字架裡跟隨我.
因為.
凡要救自己生命的就喪掉生命.
凡要喪掉生命的就必得著生命.
說什麼呢.
為什麼說的這麼玄呢.
他說凡要救自己地上的生命.
就會喪掉永恆的生命.
凡要.
凡要喪掉地上的生命.
先把它放下.
以上都是為主.
就必得著永遠的生命.
兩種不同的意識.
再講多一次.
凡要救你自己的生命.
你的一口氣.
你就會喪失永恆的生命.
凡要喪掉生命.
你的一口氣.
不要只是這一生的東西.
你就必定會得著永恆的生命.
人若賺得全世界.
配上自己的生命.
你的靈魂.
有什麼益處.
到底你最後賺些什麼.
人說拿什麼來換取生命.
真正永恆的生命怎樣去換取.
人只在他父的榮耀裡.
跟眾天使降臨的時候.

$^{921}$天國將到的時候.
那時候他會按照.
人的行為來作審判.
亦都是在講這件事.
有永恆的生命.
有永恆的生命的時候.
到底你找些什麼來賺到它.
可能我們這個時候.
就需要做盤點.
在生命裡可能要做些些審查.
計一計數.
計一計自己的帳戶.
耶穌就是帳戶.
如果你賺到全世界.
你失掉永生.
到底你最後賺什麼.
如果你還是左望右望.
猶豫不決的話.
你失去的.
可能你計都計不到.
信耶穌.
就是去顛覆我們的價值.
我們真正的價值.
是在天上.
問題就是這樣.
我們的回應是不是變在一起.
去重新命令我們所有的優先.
為了得到天國.
我們同心協力.
主我們求你來到.
做一次挑戰.
去問我們.
到底我們生命裡面最有價值的是什麼.
主天國既然進入我們人生.
既然你在我們人生當中.
出現.
我們做一次挑戰.
到底我們的生活.
是不是重新命令.
我們有的優先.

$^{961}$是為了你的永生的國度.
永遠的國度.
永恆的價值而去努力.
主我們求你再一次.
改變我們的價值觀.
幫助我們去認識我們的救主.
亦認識活著.
是為了你的.
主求你繼續去幫助.
生命是去向人們表達天國裡面的價值.
你要我們赦免的人.
我們願意去赦免.
你要我們去恩待身邊的人.
就算是受敵.
我們都可以去恩待.
因為這就是天國生命的價值.
你叫我們為真理.
為公義去發聲.
主我們願意去站出來.
因為這是天國裡面的原則.
神你要我們的生活有聖潔.
好像理性潔一樣.
主我們求你去reorder我們的priority.
由我們的生活帶回我們的家庭.
我們的工作.
是要有你天國裡面的生命.
主求你幫助.
我們用祈禱服全民革命.
主我們求你去reorder我們的生活.
是要有你天國裡面的生命.
主求你幫助.
我們用祈禱服全民革命.
\newpage



\section{撒母耳記下 24:1-25}
\label{sec:j1mkZMKAJog}
\textbf{蓋棺定論(撒母耳記下24\_1-25) - 朱莉傳道}
\newline
\newline
連結: \href{https://youtube.com/watch?v=j1mkZMKAJog}{\texttt{https://youtube.com/watch?v=j1mkZMKAJog}} ~~~~ 語音日期: 2023-01-02
\newline
\newline
\hyperref[sec:lBh3rgPDiTo]{\small{< < < PREV SERMON < < <}}
~
\hyperref[sec:index_chronic]{\small{[返順時目]}}
~
\hyperref[sec:index_scriptual]{\small{[返順卷目]}}
~
\hyperref[sec:ZVIv_jPkXmQ]{\small{> > > NEXT SERMON > > >}}
\newline
\newline
撒母耳記下 24:1-25
\newline
\begin{longtable}{cl}
\hline
\hline
章節 & 經文 (和合本修訂版)\\
\hline
24:1 & \begin{tabularx}{0.7\textwidth}{X} 耶和華的怒氣又向以色列發作，激起大衛來對付他們，說：「去，數點以色列人和猶大人。」 \end{tabularx} \\ \\ \relax
24:2 & \begin{tabularx}{0.7\textwidth}{X} 大衛對跟隨他的約押元帥說：「你來回走遍以色列眾支派，從但直到別是巴，數點百姓，我好知道百姓的數目。」 \end{tabularx} \\ \\ \relax
24:3 & \begin{tabularx}{0.7\textwidth}{X} 約押對王說：「願耶和華－你的神使百姓的數目增加百倍，使我主我王親眼得見。我主我王何必要做這事呢？」 \end{tabularx} \\ \\ \relax
24:4 & \begin{tabularx}{0.7\textwidth}{X} 但王堅持他對約押和眾軍官的命令。約押和眾軍官就從王面前出去，數點以色列的百姓。 \end{tabularx} \\ \\ \relax
24:5 & \begin{tabularx}{0.7\textwidth}{X} 他們過約旦河，在迦得谷中、城的右邊亞羅珥安營，與雅謝相對。 \end{tabularx} \\ \\ \relax
24:6 & \begin{tabularx}{0.7\textwidth}{X} 他們來到基列，到了他停‧合示地，又來到但‧雅安，繞到西頓。 \end{tabularx} \\ \\ \relax
24:7 & \begin{tabularx}{0.7\textwidth}{X} 他們來到推羅的堡壘，以及希未人和迦南人的各城，又出來，到猶大尼革夫的別是巴。 \end{tabularx} \\ \\ \relax
24:8 & \begin{tabularx}{0.7\textwidth}{X} 他們來回走遍全地，過了九個月又二十天，就回到耶路撒冷。 \end{tabularx} \\ \\ \relax
24:9 & \begin{tabularx}{0.7\textwidth}{X} 約押向王報告百姓的總數：以色列拿刀的勇士有八十萬；猶大有五十萬人。 \end{tabularx} \\ \\ \relax
24:10 & \begin{tabularx}{0.7\textwidth}{X} 大衛數點百姓以後，心中自責。大衛向耶和華說：「我做這事大大有罪了。耶和華啊，現在求你除掉僕人的罪孽，因我所做的非常愚昧。」 \end{tabularx} \\ \\ \relax
24:11 & \begin{tabularx}{0.7\textwidth}{X} 大衛早晨起來，耶和華的話臨到迦得先知，就是大衛的先見，說： \end{tabularx} \\ \\ \relax
24:12 & \begin{tabularx}{0.7\textwidth}{X} 「你去告訴大衛：『耶和華如此說：我向你提出三樣，隨你選擇一樣，我好降給你。』」 \end{tabularx} \\ \\ \relax
24:13 & \begin{tabularx}{0.7\textwidth}{X} 於是迦得來到大衛那裡告訴他，問他：「你要國中有七年的饑荒呢？或是你在敵人面前逃跑，被追趕三個月呢？或是在你國中有三日的瘟疫呢？現在你要考慮思量，我怎樣去回覆那差我來的。」 \end{tabularx} \\ \\ \relax
24:14 & \begin{tabularx}{0.7\textwidth}{X} 大衛對迦得說：「我很為難。我們寧願落在耶和華的手裡，因為他有豐盛的憐憫；我不願落在人的手裡。」 \end{tabularx} \\ \\ \relax
24:15 & \begin{tabularx}{0.7\textwidth}{X} 於是，耶和華降瘟疫給以色列。自早晨到所定的時候，從但直到別是巴，百姓中死了七萬人。 \end{tabularx} \\ \\ \relax
24:16 & \begin{tabularx}{0.7\textwidth}{X} 天使向耶路撒冷伸手要毀滅這城的時候，耶和華改變心意，不降那災難，就對那在百姓中施行毀滅的天使說：「夠了！住手吧！」耶和華的使者正在耶布斯人亞勞拿的禾場那裡。 \end{tabularx} \\ \\ \relax
24:17 & \begin{tabularx}{0.7\textwidth}{X} 大衛看見那在百姓中施行毀滅的天使，就向耶和華說：「看哪，我犯了罪，行了惡，但這群羊做了甚麼呢？願你的手攻擊我和我的父家。」 \end{tabularx} \\ \\ \relax
24:18 & \begin{tabularx}{0.7\textwidth}{X} 當日，迦得來到大衛那裡，對他說：「你上去，在耶布斯人亞勞拿的禾場上為耶和華立一座壇。」 \end{tabularx} \\ \\ \relax
24:19 & \begin{tabularx}{0.7\textwidth}{X} 大衛就照著迦得的話，照著耶和華所吩咐的上去了。 \end{tabularx} \\ \\ \relax
24:20 & \begin{tabularx}{0.7\textwidth}{X} 亞勞拿觀看，看見王和臣僕向他走過來。亞勞拿就出去，臉伏於地，向王下拜。 \end{tabularx} \\ \\ \relax
24:21 & \begin{tabularx}{0.7\textwidth}{X} 亞勞拿說：「我主我王為何來到僕人這裡呢？」大衛說：「我要買你這禾場，為耶和華築一座壇，使瘟疫在百姓中停止。」 \end{tabularx} \\ \\ \relax
24:22 & \begin{tabularx}{0.7\textwidth}{X} 亞勞拿對大衛說：「我主我王，你眼中看為好，就拿去獻祭。看，這裡有牛可以作燔祭，有打糧的器具和套牛的軛可以當作柴。 \end{tabularx} \\ \\ \relax
24:23 & \begin{tabularx}{0.7\textwidth}{X} 王啊，這一切，亞勞拿都獻給王。」亞勞拿又對王說：「願耶和華－你的神悅納你。」 \end{tabularx} \\ \\ \relax
24:24 & \begin{tabularx}{0.7\textwidth}{X} 王對亞勞拿說：「不，我一定要按價錢向你買；我不能用白白得來的東西作燔祭獻給耶和華－我的神。」大衛就用五十舍客勒銀子買了那禾場與牛。 \end{tabularx} \\ \\ \relax
24:25 & \begin{tabularx}{0.7\textwidth}{X} 大衛在那裡為耶和華築了一座壇，獻燔祭和平安祭。耶和華垂聽了為那地的祈求，瘟疫就在以色列中停止了。 \end{tabularx} \\ \\
[1ex]
\hline
\hline
\end{longtable}
$^{1}$大家好, 姐妹平安.
今天想和大家分享的題目是「蓋棺定論」.
可能有些年長的人會覺得.
大家是女士, 為什麼要用這些題目呢?.
但我覺得往往在一個喪禮.
是一個最好的陪伴.
因為在一個人的終點.
回望一回.
的確是一個很好的機會思考人生.
正如我們現在站在年尾或年初的時候.
當我們回望過去的一年.
我們有沒有一些事情想改變呢?.
或者我們有沒有一些教訓想總結.
再展望新的一年呢?.
今天想和大家看的題目是.
《撒慕爾記》下二十四章.
我們會花些時間看完整章聖經.
我們先來看一到十七節.
我們跳來看.
雖然比較長.
但我們也用些時間來讀這段經文.
《撒慕爾記》下二十四章第一節.
耶和華又向以色列人發怒.
就激動大衛.
是他吩咐人去數點以色列人和猶大人.
大衛就吩咐跟隨他的元帥約翰說.
你去走遍以色列眾支派.
從彈藉到別士巴.
數點百姓.
我會知道他們的數目.
約翰對王說.
無論百姓多少.
願耶和華你的神再加增百倍.
使我主我王親眼得見.
我主我王何必喜悅行者事呢?.
但王的命令勝過約翰和眾軍長.
約翰和眾軍長就從王面前出去.
數點以色列的百姓.
我們跳到第十節.
大衛數點百姓之後.

$^{41}$就心中自責.
禱告耶和華說.
我行者事大有罪了.
耶和華求你除掉僕人的罪孽.
因我所行的甚是原末.
大衛早晨起來.
耶和華都說臨到先知迦德.
就是大衛的先見說.
你去告訴大衛說.
耶和華如此說.
我有三樣罪.
隨你選擇一樣.
我好降與你.
於是迦德來見大衛對他說.
你願意國中有七年的饑荒嗎?.
是在你敵人面前逃跑.
被追趕三個月呢?.
是在你國中有三日的瘟疫呢?.
現在你要揣摩思想.
我好回覆.
那差我來的.
大衛對迦德說.
不心為難.
我願落在耶和華的手裡.
因為他有豐盛的憐憫.
我不願落在人的手裡.
15節.
於是耶和華降瘟疫與以色列人.
自早晨到鎖定的時候.
從彈藉到別時巴.
民間死了七萬人.
天使向耶路撒冷伸手要滅城的時候.
耶和華後悔.
就不講這災了.
他吩咐滅盟的天使說.
夠了 絕手吧.
那時耶和華的使者在耶布斯人阿魯拿的和場那裡.
大衛看見滅盟的天使就禱告耶和華說.
我犯了罪 行了惡.
但這群羊做了什麼呢?.

$^{81}$願你的手攻擊我和我的父家.
事實上.
《撒母耳記下》二十四章.
都有一個總結的意味.
一方面可以說是大衛生平的總結.
雖然大衛的死亡是記載在《列王記》上的.
但在《列王記》上明顯大衛只是一個過渡的人物.
主角已經換成他的兒子索納門.
所以大衛.
在舊恩歷史裡有這麼重要地位的人物.
他的蓋棺定論竟然就是這張聖經.
可以說是用一個污點.
一個數點文數野神發怒的事跡來作結.
所以非常奇怪.
當然我們可以說.
聖經沒有引惡揚善.
很真實地記載.
這是必然的.
因為我們看到聖經很多人物都是這樣.
兩邊平行地記載.
但除了《撒母耳記》下24章.
除了是大衛人生的總結之外.
亦是《撒母耳記》上和下這卷書.
如果把兩卷書放在一起.
一樣地看的話.
這卷書的總結都是在《撒母耳記》下24章.
而《撒母耳記》我們知道.
最重要的主題往往是問什麼呢?.
就是誰是以色列最理想的王.
Who is the ideal king for 以色列.
所以整個《撒母耳記》上.
既然不是撒母耳又不是蘇羅.
找到大衛了.
他應該是那個ideal king.
最合神心意的王.
但是傳書的最後那一節.
即是最後這一段24章.
好像推翻了前面的結論.
所以是非常奇怪的.
如果我們看前面一點.

$^{121}$其實《撒母耳記》下23章.
已經是不同的風景.
《撒母耳記》下23章第一節這樣說.
以下是大衛沒了的話.
耶西的兒子大衛得居高位.
是雅各臣所高的.
作以色列的美歌者.
然後一連串勇士的名單.
多威風 多蒙神祝福.
那個更似結尾.
為什麼要有24章呢?.
我們華人的父母很喜歡說一句話.
贏在起跑線.
但是如果輸在終點.
終結往往更重要.
所以我們去環繞這個問題.
來看24章.
怎樣總結了大衛.
甚至是搜尋「理想王」的過程.
《撒母耳記》整卷書的內容.
的訊息在這裡怎樣總結.
而大衛又如何.
為什麼如何.
也是How 一個Why.
仍然被認證為神最理想的王呢?.
首先我們必須承認.
這段經文是很難解釋的.
因為有很多的難題.
最明顯的兩個難解經文.
第一節.
「耶和華又向以色列人發怒.
就激動大衛.
使他吩咐人去掃點以色列人或猶太人」.
所以很奇怪.
為什麼始作俑者到底是神還是大衛呢?.
是神來激動大衛去掃點.
為什麼又去怪責大衛呢?.
這是其中一個很大的難題.
另一個難題.
我們可能也會覺得奇怪.

$^{161}$就是當大衛意識到自己犯罪的時候.
先見迦德奉神的命令去見大衛.
然後叫他選擇刑罰.
好像點心紙那樣.
好像你去喝茶點心.
你自己點吧.
很奇怪.
聖經從來沒有記錄到神給人選擇刑罰.
所以這段聖經的確有很多很多我們不明白的地方.
可能是什麼呢?.
可能我們放錯重點.
問錯問題.
很多的事情聖經沒有解釋.
但相反.
有一些地方聖經很詳細地記載.
甚至是我們想像不到.
不合乎比例地記載一些事情.
那些往往是我們才需要注意的地方.
但我們卻很多時候會忽略了.
例如這段聖經.
《撒慕爾記》下24章.
我們剛才讀到17節.
但18至25節才是危機如何解決.
其實解決的部分在這段.
應該是故事的高潮.
但這段聖經我們讀下去.
反而不太注重這段經文.
而且這段經文有些很特別的地方.
我們先讀18至25節.
18節.
當日迦德來見大衛對他說.
「你上去在耶布斯人阿勞那的和牆上為耀華築一座壇」.
大衛就照著迦德奉耀華名所說的話上去了.
就迎接出去.
面服於地向王下拜.
說:我主我王為何來到僕人這裡呢?.
大衛說:我要買你這和牆.
為耀華築一座壇.
使民間的瘟疫止住.
22節.

$^{201}$阿勞那對大衛說:我主我王你喜悅用甚麼就拿去獻祭.
看那這裡有牛可以作繁祭.
有打樑的器具和套牛的鴨可以當柴燒.
王阿這一切我阿勞那都奉給你.
有對王說:願耀華你的臣閱立你.
24節.
王對阿勞那說:不然我必要按著價值向你買.
我不肯用白德之物作繁祭獻給耀華我的臣.
大衛說:就用50赤克勒銀子買了那和牆與牛.
大衛在那裡說:築了一座壇獻繁祭和平安祭.
如此耀華垂聽國民所求的瘟疫在以色人中止住了.
所以瘟疫的危機是如何解決呢?.
當然是大衛聽神的話照著神的話去築壇獻祭.
但如果築壇獻祭是一個如此重要的解決方案.
為什麼不是詳細記載大衛如何築壇獻祭呢?.
相反只有25節很短的幾個字去形容這件事就算了.
相反這段是不成比例.
超乎我們想像的詳述大衛如何得到那塊地.
很奇怪.
另外很奇怪的地方就是角色調換.
有沒有留意地主亞勞那不斷獻上.
而大衛王座不斷推辭.
如果我們更加明白背景.
耶穌亞勞那其實是什麼呢?.
他可能是耶路撒冷的原住民.
當大衛攻佔了耶路撒冷之後.
他可能留下了或是俯首稱臣.
所以一個以色列王對著一個歸順了的手下敗將.
為什麼如此客氣呢?.
其中一句我們可以留意的這段話.
關鍵詞是哪個呢?.
是22節.
「亞勞那對大衛說:我主,我王,你喜悅用什麼.
就拿去獻祭」.
關鍵詞是「拿去」.
「Take」.
T-A-K-E.
「Take」.
亞勞那說:你隨便拿.
你喜悅什麼就隨便拿.

$^{241}$英文是說:你用眼睛掃過.
你看到一些Pleasing in your eyes的東西.
你都可以隨便拿.
為什麼他這麼信他呢?.
當然了.
因為他對面站著的那個.
或是正在和他對話的那個是誰呢?.
大衛是王.
他是王.
王就有權去管轄,去壓制.
在《撒母耳基相》第八章.
當時以色列人哀求先知撒母耳為他們立王.
那裡有一段經文.
提到撒母耳如何警告以色列人.
他們立王的要求.
其實是會最後自食其果的.
撒母耳如何警告以色列人呢?.
《撒母耳基相》第八章十一節.
「管轄你們的王必這樣行.
他必派你們的兒子為他趕車」.
其實這是「Take」.
「派」的意思是他拿你的兒子.
「跟馬奔走在車前.
又派他們作千夫長,五十夫長.
為他耕種田地,收割莊稼.
打造軍器和車上的器械」.
第十三節「必娶」.
「娶」又是「Take」.
「你們的女兒為他製造香膏,做飯,烤餅」.
「也必娶」又是「Take」.
「你們最好的田地,葡萄園,橄欖園.
次給他的臣僕.
你們的糧食和葡萄園所出的.
他必娶十分之一給他的太監和臣僕」.
「又必娶」又是「Take」.
「你們的僕人婢女建葬的少年王們.
你們的奴工他差役.
你們的羊群他必娶十分之一.
你們也必作他的僕人.
那時你們必因所損的王哀求」.

$^{281}$「耶和華」.
「耶和華卻不應運你們」.
原來一早在撒慕爾的上第八章.
撒慕爾已經警告.
你們要立王這件事.
只不過是製造一個王來管轄.
去「Take」你們所有的東西.
而的確,這就是人性的本性.
當我們有權力去做的時候.
我們就會為自己拿取.
同樣的道理.
其實大衛犯奸淫和謀殺的事件.
在撒慕爾的下十一章.
我們不會詳細地說.
但是當時發生什麼事.
如果有機會你回去看看.
就是他是神所選召的王.
但他也怎樣?.
他也用他的權力去「Take」.
「Take」那個必孕的妻子.
甚至間接地可以這樣說.
他「Take」奪走了烏尼亞的性命.
所以我們回到撒慕爾記下二十四章.
當數點人數的時候.
其實發生什麼事呢?.
我們再讀二十四章的第一到第九節.
「耶和華有向以色列人發怒.
就激動大衛.
使他吩咐人去數點以色列人和猶太人」.
第二節.
「大衛就吩咐跟隨他的元帥約翰說.
留意這和合本是譯作「大衛就吩咐」.
但如果你看英文譯本.
是「王就吩咐」.
「The King 就吩咐」.
事實上應該是譯作「王」.
所以第二節.
「王就吩咐跟隨他的元帥約翰說.
你去走遍以色列眾支派.
從彈藉到別氏巴數點百姓.

$^{321}$我好知道他們的數目.
約翰對王說.
無論百姓多少.
願耶和華你的神再加增百倍.
使我主我王親眼得見.
我主我王何必喜悅行者事呢?.
但王的命令勝過約翰和眾軍長」.
約翰和眾軍長就從王面前出去.
數點以色列的百姓.
然後我們跳到第九節.
「約翰將百姓的總數就告於王.
以色列拿到的勇士有八十萬.
猶大有五十萬」.
留意那個稱呼.
第一到第九節.
除了第一節有說到.
這個王是誰就是大衛.
他被神激動去做這件事.
但由第二節一直到第九節.
不斷稱呼的是什麼?.
就是這個王.
當然大衛是王.
這是事實.
但原來數點人數的問題出於哪裡呢?.
就是當大衛覺得自己是最王的時候.
我最有權力去拿取.
去剔.
我的權力是最高峰.
甚至是第四節.
王的命令勝過.
其實約翰有勸他.
他說「王的命令勝過」.
或者有些背景資料.
我們看到原來這個階段.
很可能是大衛王最權力最盛.
國力最盛.
版圖最大.
厲害了.
大衛的國.
但他是距離合臣心意.

$^{361}$可能是最遠的時候.
轉機在哪裡呢?.
他覺得自己是王.
很厲害.
可以做很多決定.
但錯了.
第十節.
大衛數點百姓之後.
就心中自責.
禱告耶和華說.
我恨這是大有罪了.
耶和華求你除掉僕人的罪孽.
因我所行的心思愚昧.
他稱自己愚昧.
原來撒母耳基相.
都同樣用同一個字「愚昧」.
來形容數來的失敗.
他承認.
他好像前任一樣的失敗.
第十一節.
大衛清早起來.
耶和華的話臨到先至加德.
就是大衛的先見.
這個名稱又再變回大衛.
當他再一次意識到.
他站在神的面前.
他不是那個王.
當他肯承認錯誤.
當他肯面對改變的時候.
結果我們去到後面.
剛才讀那個交易.
大衛和阿勞那的對話.
關於那塊田要怎樣交易的時候.
是反過來的角色.
是阿勞那一直說.
你是王.
我要奉給你.
我要獻給你.
但大衛就.
不行.

$^{401}$不可以.
我不可以這樣.
所以反而是他堅持.
不拿.
不取走阿勞那的地和牛.
廣東人說.
這麼大的甲拿.
已經送到上門.
但他寧願付代價.
也不是.
不想行使本應該有的王的權力.
因為他知道.
合神身的王.
不是拿取的.
我再看回中間.
講到點心紙的刑罰.
神給大衛選.
有三種刑罰.
大衛怎樣選.
24章14節.
大衛對迦德說.
我甚為難.
我願落在耶和華的手裡.
因為他有豐盛的憐憫.
我不願落在人的手裡.
這裡怎樣?.
他不是「滴」.
他是「課」.
他寧願落在.
他寧願將自己交在神的手裡.
當一權有權力的時候.
除了「滴」之外.
你以為你有很多選項.
但大衛這裡是否拒絕去做選擇.
為甚麼?.
他將權交.
其實不是拒絕做選擇.
是他將可以做選擇的權力.
怎樣?.
他交給上帝.

$^{441}$而且他知道上帝有豐盛憐憫的上帝.
他寧願選擇去相信神豐盛的憐憫.
當然在後面.
我們今天不會有詳細時間說.
上帝的確是為人後悔的上帝.
不是我們人那些說.
我們做錯了甚麼.
我們後悔的.
是怎樣?.
神是會為人改變他的心意.
從刑罰到私恩.
我們快有聖誕節了.
其實聖誕的意義也是這樣.
我記得我看過一個片段.
這些片段大家應該經常會看到.
就是玩團體遊戲.
玩信任遊戲.
Trust game.
一個人站高.
然後要跌下去一個位置.
然後他的同伴要去接住他.
這是一個信任的遊戲.
但有些搞笑的片段是甚麼?.
他不是向後跌.
他的同伴在這裡.
他應該向後跌下去.
讓他的隊友去接住.
不是,他跌出了方向.
他向前跌.
當然是搞笑.
但到底我們將自己交給甚麼人?.
我們將自己交給.
或者這樣說.
我們很習慣用「Tick」.
來證明自己的價值.
我們將自己「Tick」了很多東西.
然後證明了自己的價值.
我們有很多頭銜.
我們有很多資產.
我們好像有在世上很多的成就.

$^{481}$很多值得去記錄的事情.
同時我們也將自己交給甚麼?.
交給股市,樓市.
交給我們所擁有的資產.
我們的銀行,存設部.
交給我們的家人.
然後我們也像大衛一樣.
數我們的銜頭,我們的學位.
我們銀行存款有多少個零.
甚至我們的身體檢查的結果.
其實大衛提醒我們.
最安全的是將自己交在神豐盛的憐憫裡面.
即使神看起來不合理.
就好像上帝不合理的時候.
大衛仍然選擇去相信神仍然是有豐盛憐憫的神.
我們對這段經文仍然有很多不明白的地方.
我們不明白為何神會容許這些事情發生.
為何又激動又刑罰.
為何是人民來面對這些刑罰.
我們希望看到的是甚麼呢?.
看到大衛對神仍然有豐盛憐憫的相信.
而為何大衛在最後三母耳記下二十四章.
仍然承認被認證為神最理想的王.
是The Ideal King 仍然是David.
因為我們看到他是一個肯承認唯有上帝是王的人.
他發現原來他只有王的位置.
不是要剔他是要Fall.
他Fall不單是落在上帝的慈愛裡.
也落在上帝的權力下.
他承認唯有上帝是王.
在這裡我們想起另一個人物.
三母耳記的另一個人物.
在整本書裡出現的第一個主角.
我們最後二十四章大衛是主角.
但在三母耳記上第一章或第二章.
其實有另一個主角.
那個是誰呢?.
那個是Hannah.
如果我們將大衛和Hannah作比較.
你覺得沒有什麼可以比較.

$^{521}$Hannah的確是三母耳記上第一章的人物.
而且老實說她是典型的無線電視劇內容.
首先她不辱.
甚至三母耳記上第一章第五到第六節.
形容為「耶和華不使Hannah生辱」.
然後她的丈夫爾利加拿的另一個妻子.
又惹她發怒.
以致她不吃飯.
心裡受苦哭泣.
然後到侍來的時候祈禱.
甚至以利以為她喝醉酒.
Hannah又如何在神的面前傾心吐怒.
後來又去吃飯.
面上再不戴手用.
結果最後Hannah懷孕.
因為她曾經許願.
要將苦苦得來的兒子歸給神.
就是三母耳.
後來她就將他奉獻給上帝.
Hannah的故事是典型的電視劇內容.
Hannah和大衛似乎沒有什麼可以比較.
好像完全沒有什麼相關.
一個沒有任何功績的婦人.
一個以色列最偉大的王.
但其實她是一頭一尾.
在三母耳上和三母耳下的一頭一尾.
都告訴我們.
其實上帝心中最理想的王.
最合神心意的人生應該是怎樣.
最合神心意的王是怎樣呢.
就是一個承認唯有神是王的人.
就是一個唯有承認耶和華是王的人.
Hannah得到三母耳.
但她結果拿去奉獻.
在三母耳下我們看到.
大衛寧願自己家庭受苦.
也不想再看到人民受苦.
其實這種態度是基於一種意識.
就是唯有上帝是王.
儘管問題來源可能是上帝.

$^{561}$因為Hannah不辱.
聖經形容是神使她不辱.
而在三母耳下二十四章又形容.
大衛是被神激動.
我們不明白.
但的確神是有主權.
似乎祂可以使人不辱.
又可以使人生辱.
祂可以降災.
亦可以撕裂民.
我們必須承認.
上帝是有絕對的王權.
儘管我們不明白.
最合神心意的秘訣在哪裡呢.
三母耳記上第二章Hannah的禱告.
三母耳記上第二章第一節.
第一節:「哈拿圖告說:.
我的心因耶和華快樂,我的國因耶和華高舉,.
我的口向受敵張開,.
我的因耶和華的救恩歡恩,.
只有耶和華為聖,.
除他以外沒有可比的,.
也沒有盤石像我們的神.」.
神是不可比的.
第五節:「數來飽足的反作,.
用人求食,.
飢餓的再不飢餓,.
不生育的生了七個兒子,.
多有兒女的反倒衰微.」.
第六節:「要說使人死也使人活,.
使人下陰間也使人往上升,.
他使人貧窮也使人富足,.
使人卑微也使人高貴,.
他從灰塵裡抬舉貧寒人,.
從糞堆中提拔貧乏人,.
使他們與王子同坐,.
得著榮耀的座位.」.
這就是一個這樣的神.
他是隨己而行的.
他有絕對的主權.

$^{601}$他更加是無可比擬的.
他叫人死又叫人生.
他可以完全這樣去顛覆現實.
問題是我們是否肯承認他的主權.
是否肯將人生交在他的主權下.
儘管有時我們不明白.
而且相信他是有豐盛憐憫的神.
撒母耳記下24章第25節.
最後一節.
「大衛在那裡為耶和華築了一座壇,.
獻繁祭和平安祭,.
如此耶和華誰顧國民所求的,.
溫逆在以色列中就止住了.」.
最後的決定是什麼呢?.
其實真正的王從來都是耶和華.
是神聽了為那個地方的祈求.
是神洗溫逆在以色列中止住.
他仍然是有絕對王權.
天地萬物的唯一主宰.
他是上帝.
他是神.
撒母耳記提醒我們.
我們是否肯承認這件事.
是否肯將自己交給他.
或者這件事.
亦是我們在一年的結束.
展望新的一年.
都是一個提醒.
到底我們人生應該怎樣過呢?.
怎樣去過一個促成心意的人生呢?.
求主來幫助我們.
我們一起祈禱.
神下求你幫助我們.
憐憫我們.
在人生裡面有很多高低起跌.
就算我們回望在這一年裡.
可能我們都有很多不同的經歷.
當然有很開心的時候.
可能亦都有失望.
有受挫折的時候.

$^{641}$但無論如何.
主要幫助我們.
讓我們的眼光仍然對準你.
仍然去學習怎樣謙卑地在你面前.
與你同行.
在一個不確定的年代.
一個充滿挑戰危機的時代.
我們緊緊地抓緊你.
讓我們去承認.
我們有很多事不明白.
但我們仍然信服.
在你的主權底下.
因為我們真的知道.
唯有你是王.
主要求你祝福我們.
憐憫我們.
我們這樣禱告.
奉耶穌基督成名.
阿們.
\newpage



\section{13:3-5}
\label{sec:ZVIv_jPkXmQ}
\textbf{Stewarding the Gift of Power (John 13\_3-5, 12-16) - Tommy Lee}
\newline
\newline
連結: \href{https://youtube.com/watch?v=ZVIv_jPkXmQ}{\texttt{https://youtube.com/watch?v=ZVIv\_jPkXmQ}} ~~~~ 語音日期: 2023-01-10
\newline
\newline
\hyperref[sec:j1mkZMKAJog]{\small{< < < PREV SERMON < < <}}
~
\hyperref[sec:index_chronic]{\small{[返順時目]}}
~
\hyperref[sec:index_scriptual]{\small{[返順卷目]}}
~
\hyperref[sec:DsrtebXqSNI]{\small{> > > NEXT SERMON > > >}}
\newline
\newline
$^{1}$All right..
Good morning, happy new year..
Southwest Airlines, if you go to Dallas, Texas, there is a Herb Kelleher Way..
Not sure if you ever heard the name of Herb Kelleher..
He was the CEO and chairman and president of Southwest a number of years ago, a couple.
decades ago, retired..
When you think of Southwest Airlines, you think of one of the premier efficient customer.
service on time..
It messed up pretty badly, as Mark is a testament of..
And over 25,000 flights around the country were all messed up..
And the reason being why?.
Southwest got too confident..
They knew they were one of the best, best price, best efficiency, all of that stuff..
They started letting things go lax..
They started not to invest in different technology..
They started taking chances, thinking, you know what, we'll just keep riding..
We'll make sure we continue to make the most profit..
All the stuff that Herb Kelleher wanted to do, efficiency, customer service, on time,.
everything, all started being lost in the whole process because they all became about.
making money for the bottom line..
And as a result, they didn't have to upgrade any of the equipment..
And as a result, they paid for it as what we saw this past in the last couple of weeks..
And that's a topic that we're going to talk about today..
It goes to this whole theme of power, a topic that we never really talked about before..
How do we steward the power that God gives us?.
All of us have power..
All of us have influence..
All of us have gifts..
All of us have skills..
Financially, what we're able to do, our work, our status, our position in life, we all have.
power..
How do you then steward that power well?.
One culture and in all throughout history, we've always seen the use of power and the.
misuse of power..
We've seen this picture..
I've shown you this picture before..
And for those who are new, a lot of times if you go over to Russia as well, too, President.
Putin, when he meets dignitaries and people come in, he wants to show you his power..
And he wants you to come and meet him in this room..
That is a demonstration of power because he wants to show you that he's in charge..

$^{41}$Another thing, we've seen the misuse of power throughout history..
We see Adolf Hitler, who used his power to continue to ravage around all of the world.
and go against people groups and abuse and enslave them..
We've also seen this..
We celebrate the two-year anniversary of the January 6th insurrection at the White House..
We see signs like Jesus saves Christians who are saying that their pastors or churches.
and all that stuff who are going into the Capitol building and saying that the election.
has been rigged and who's sitting there going after the vice president, different state.
officials in order to really be able to allow a person who many believe is the king of that.
power to remain in power..
We look at culture as well, too..
Look at this..
When we think of the future of AI or artificial intelligence, pretty soon many of you guys.
who are in the tech field, many of your kids who will be in the tech field will have massive.
amounts of power..
We see cars that are now driven without hands free..
Everywhere you go from McDonald's in different places, kiosks, you go to the United Center,.
they don't have cashiers or airports, they no longer have cashiers..
Everywhere you go in, you sit there and you show them your phone and whatever you take.
now is charged on your phone..
And as more and more data begins to really be built, as more of these computer systems,.
artificial systems begin to understand your ways and gather more data, power begins to.
be built..
And that power through artificial intelligence can be used for good, but it also can be used.
for evil..
Social media, as we begin to see, can also be used for good and also used for evil as.
well, too..
We're seeing things like Twitter and Facebook, which are beginning to aggregate your behavior..
They're feeding information that caters to what you believe in..
As Christians who are working in technology, whether it's AI or in some of these technology.
companies, how do we think what is best for the world?.
We need Christ followers in these areas to steward that power well, to really be able.
to say what decisions do we have to make in our jobs, in these industries, to continue.
to create a world that cares and balances, but not be used and abused by different people..
Power is found in artificial intelligence..
It's found in social media..
Power can be used for evil, but power also can be used to inspire and to help..
We saw this picture in 9/11..
One of the most powerful images that you ever see, and I'm not even saying whether you're.
Republican or Democrat, and for those who are watching as well, too, this is a picture.

$^{81}$of President George Bush..
When 9/11 hit, President Bush went over to New York City and grabbed a microphone, and.
he reminded the people that he hears them..
He gave hope in the world that needed hope..
He's retired now..
I'm not sure if you know what President Bush is doing right now..
He does this for a living..
He paints pictures of immigrants, and he tells their story of hope..
If you go over to Bush Library, I have been there over in Dallas as well, too..
He paints images of each person, their face, and he tells the story of how they have overcome.
trials, and he gives us a picture of hope within that..
And one final thing, Chick-fil-A, the story of a company who uses their power for good..
That when you go through the drive-thru or go into the restaurant at Chick-fil-A, what.
is it that you hear from every single employee?.
"How may I serve you today?".
You see the story of Dan Cathy, who really believed as a Christ follower that he wanted,.
as he says, that we are here to serve..
We keep the needs of operators, their team members, and customers at the heart of our.
work, doing what is best for the business and best for them..
And as a result of that, in every single situation, we serve the customers, we serve our clients,.
we serve our staff employees..
How can we serve you?.
And they continue to build a brand, and every single time you go to Chick-fil-A, your orders.
are usually correct, you are treated well, and the customer service is phenomenal because.
it stands for the glory of who he works for..
As Christ followers, our topic today, I'm going to go back to this, is this..
That Jesus gives us a new framework of power..
That we are called not to give up our power, but we steward the power we are given for.
the gospel..
That's the big idea I want to get across today..
When I think of power, when I want you to think of power, the influence that you have,.
the skills that you have, the money that you have, Jesus gives us a new framework of power.
as we're going to look at John 13 today..
That as Christ followers, we are not called to give up our power..
That is never the case..
You don't give up your power..
But you steward your God-given power, your God-given skills, your God-given talents,.
your God-given resources in the form of money and status, and you steward that power we.
are given for the glory of the cross and for the gospel..
If you have your Bible, let's take a quick look over in John chapter 13..

$^{121}$We begin in verse 3..
John 13, verse 3..
Or actually, you know what, let's start in verse 2..
Going to supper when the devil had already put it into the heart of Judas Iscariot, Simon's.
son, to betray him..
Jesus, knowing that the father had given all things into his hands and that he had come.
from God and was going back to God, rose from supper..
He laid aside his outer garments and taken a towel, tied it around his lilies..
Verse number one, point number one..
Jesus knew what power he had and where his power came from..
Jesus, knowing that the father had given, I'm reading from the ESV, had given all things.
into his hands and that he had come from God and was going back to God, rose from supper..
Jesus knew what power he had..
Jesus knew who had given him the power..
As Christ followers, we know that Jesus gives us power..
We know that God gives us power..
We know the skills that you have, the talents that you have..
We know the position, the job that you have, how smart you are, the resources, financial,.
how rich and wealthy you are financially, is all given to us by God..
We know where that power comes from..
We know what power we have..
Now at the same time, what are you doing with it?.
We're reminded in John chapter one, automatically, if you look at John chapter one, story of.
John chapter one begins this..
In the beginning was the word and the word was with God and the word was God..
He was in the beginning with God..
All things were made through him and without him was not made any, was not anything made.
that was made..
In him was life and the life was a light of men..
The light shines in the darkness and darkness has not overcome it..
Way back in the beginning, John one, we were reminded and John gives us the entire backdrop.
of this is who Jesus is..
This is the power that Jesus had..
He was at the beginning..
John chapter two, we begin to see the story of what God has..
John chapter three, who does he encounter?.
He encounters Nicodemus..
Nicodemus comes up to him and says, "Rabbi, who are you?.
Are you the teacher that we were expecting?".
And that many of you guys are all familiar with John chapter three, the most famous of.

$^{161}$the verse..
What does he say to him?.
For God so loved the world that he gave his one and only son..
That John chapter one tells us and gives us an option of this is who Jesus is..
At the beginning of life, when all of this was not even created, there stood Jesus, the.
living word..
Jesus is the rabbi in John chapter three, in which Nicodemus comes to him..
Jesus, as Jesus tells Nicodemus, "I am the savior for God so loved the world that he.
sent me his one and only son, that whoever believes in me shall not perish, but have.
eternal life.".
And he says it to who?.
Pharisee, teacher of the law, a Jewish man..
A man who believes that his power is only focused on the Jews..
And Jesus tells him, "You're wrong, because when God sent me and my power, it was for.
the world.".
And story after story after story, if you look at John chapter four, chapter five, chapter.
six, chapter seven, you see the story of Jesus and the Samaritan woman..
Jesus in chapter five, Jesus and the Gentile official..
John chapter six, Jesus and the invalid..
John chapter seven, Jesus and the feeding of the 5,000..
Story after story after story of what it means to love the world..
What it means to use your power for good..
Jesus knew his power..
Jesus knew where his power came from..
Now the big question, point number two, is that during supper, when the devil had already.
put into it, Jesus, knowing that the father had given all things into his hands and that.
he had come from God and was going back to God, rose from supper..
He laid aside his outer garments and taking a towel, tied it around his waist..
And he poured water into a basin and began to wash his disciples' feet..
Point number two, once we know our power, once we know where our power comes from, point.
number two, you make a choice of what you're going to do with that power..
Jesus made a choice of what he's going to do with that power because we are called as.
Christ followers not to give up our power, but we steward the power we are given for.
the gospel..
And if we know what our power is and we know who gives us the power, what choice are we.
going to make to do with that power?.
He gets up from the meal..
What choice are you going to make for that?.
During in the last couple of weeks, I've gotten a chance over the Christmas break to sit down.
and I've gotten a chance to talk with two individuals..

$^{201}$Let me backtrack in that one time..
All of us as individuals, as Christ followers, make a choice each and every single day of.
what we're going to do with that power..
You're asked to steward that power..
You could take that power and continue to build up your own life..
That's easy..
You could take that power, enjoy a comfortable life..
That's easy too..
But how do you begin to use that power and make a difference in your life?.
One such example of that person that sticks out to me is right here in the Bridgeport.
area..
If you ever go down, there's a dental office right near Healy Elementary School..
For those who are watching, Bridgeport is this neighborhood that we're in..
There's a guy by the name of Dr. Sam Park..
Sam Park is a Korean dentist who is overseeing and is the lead dentist for Bridgeport Children's.
Dental..
Sam Park and I have known each other for a long, long time..
Sam Park is one of the most generous individuals as a Christ follower that I knew..
When I worked in the church, he was the elder of a church up north when I first met him..
If there is one person who teaches me so much about generosity, it's Mr. Sam Park, Dr. Sam.
Park..
Every single day, Sam walks in..
And every single day, he wants to create a dental office that really ministers and serves.
the community well..
Every single employee is tasked with treating every single one of your employees that as.
a Christian, that he wants you to experience the love of Christ through all of his employees.
as soon as you walk in..
Dr. Sam Park sits down with every single one of his employees..
And he sits there and he wants to talk with them about what their goals are, what they.
want to earn in life, and how he could serve them well..
Dr. Sam Park, when I first started Resource Global, we had no money..
Sam Park was one of the most generous individuals that I knew..
He knew that our organization had no money..
This was about 12 years ago..
And when he knew it was payroll time, he would always give me a call and says, "Tommy, how.
are we doing on payroll?".
And I would sit there and say, "Sam, things are a little short this week," or, "Hey, we're.
working through some donations and when you're first starting about 12 years ago, you're.
kind of always short.".
And right before payroll, about a couple days right before payroll came, I would open up.

$^{241}$my check in the mailbox and there was a check from Dr. Sam Park, every single payroll..
He showed me what it meant to really be able to be generous in everything that we do..
During Christmas break, though, I had two conversations with different people..
And both of them were very bad conversations because all of us as Christians are stewarding.
our power well or good, but also some of us choose this in a way..
And so one of the conversations I had was a phone call from one of my employees who.
told me, "Tommy, I just wanted to let you know that we finalized our divorce with my.
husband..
My husband has recently, he had cheated on her for about a couple times in the last few.
years..
And as a result now, he wanted to divorce her, but he wanted the house..
And negotiations were not going very, very well..
And it got to a point where he felt like as the male person in a Christian fellowship.
or in the marriage, the woman had to defer to him..
And he ended up getting the house..
And as a result, they had to tell their kids after Christmas, who had no idea that their.
parents were getting a divorce, that the wife, as well as the kids had to move and the husband.
was getting a divorce..
Second conversation I had was with another person who says, "Tommy, I would like to have.
a breakfast.".
And she got together and we sat down and I said, "How are things going?".
And she says, "I just want to let you know I separated from my husband..
My husband also is one of these individuals who really believes that in a marriage, the.
men are in charge and women submit to the authority of men..
And as a result, Tommy, I'm the primary breadwinner," this lady tells me..
And she is..
She makes millions every single year through her business..
But she says, "I have to come before my husband..
I pay for everything in our house, everything from electricity, mortgage, everything I pay.
for..
But he won't let me spend any money until I get permission from her because he's the.
man of the house..
When we raise our kids, when we disagree on what the kids need to do, I have to submit.
to him..
Our kids always have to always submit to him in anything he wants, even if they disagree.
because he's the man of the house.".
My daughter tried to commit suicide three times, Tommy, because she couldn't stand my.
husband anymore..
He was so domineering..
He was always telling her she wasn't good enough..

$^{281}$And so one day she wrote a note to him about why she wanted to take her life..
And she gave it to him and says, "This is how you made me feel.".
And he says, "I am the man of the house..
Biblically I am in charge.".
We have a choice of what we do with our power..
We have a choice of whether we use and steward our power for good or for bad..
God has given each one of us, young or old, certain skills, certain resources, all of.
that..
We sit here with certain power..
What is the choice that you're going to make to do with power?.
Jesus gives us this new framework of power..
Jesus who knew what power he had, he knew where his power came from, makes a choice.
to stand up, take off his robe, and begin to wash his disciples' feet..
Because you know why?.
He was not concerned about status stemming from his power..
He was not concerned about status..
And the disciples had to learn something..
In Luke 22, right?.
As they were walking over to dinner, you know what they were arguing about?.
Who is going to be the greatest in the kingdom and who is going to sit on the left hand and.
right hand of Jesus?.
That's the discussion that they were doing..
That's what they were arguing about..
They were concerned about their status..
And as Christians, one of the things that Jesus teaches us in this whole idea of the.
upper room, the last supper, is don't be concerned about your status..
He takes off his robe and he begins to take a basin of water and washes his disciples'.
feet..
The Roman author writes this and he says this, "Even today, many cultures preserve the.
sense of shame that accompanied the feet in Jesus' day..
The lowest part of the body, the one most covered in the grime of the world, far more.
grimy in Jesus' day than our modern sidewalks, and the most likely to smell a sweat and toil..
To touch another person's feet was an act of profound subjugation fit only for a slave..
The feet were the antipodes of power, kept away from the table, toward which all guests.
lean, accorded the least attention, and was the least dignified..
Within a normal feast, some household slave would have performed the foot washing..
Should the ritual of foot washing somehow be neglected by the servants attending the.
feasts, it would have been the job of the least powerful guest present, making a quick.
and entirely calculation status, to quickly take up the culture and perform the duty.".
Amongst the people in this room, who was going to be the least?.

$^{321}$Was it going to be Matthew, the tax collector?.
No, not Matthew, the tax collector..
Was it going to be the two, the fishermen that were sitting there, Peter, James, or.
John?.
It certainly was not going to be Judas Iscariot..
He had already sold out Jesus and he was hoping to grab all the money for himself..
That wasn't going to be him..
Who was the least amongst all these twelve?.
It was going to end up being Jesus..
Jesus makes the choice to take off his robe and start washing the feet of the disciples..
He takes off his outer clothing, he wraps a towel around his waist, and after that he.
poured water into basin and began to wash his disciples' feet, drying them with the.
towel that was wrapped around him..
And verse 12, "When he had finished washing their feet, he put on his clothes and returned.
to his place and he says to them, 'Do you understand what I have done for you?'.
He said to them, 'You call me teacher and Lord, and rightly so, for this is what I am..
I am your teacher..
I am your Lord..
But now that I, your Lord and teacher, have washed your feet, you shall also wash one.
another's feet..
I have set you an example that you should do as I have done for you..
Very truly I tell you, no servant is greater than his master, nor is a messenger greater.
than the one who sent him.".
We know our power..
We know where our power comes from..
We make a choice of what we do with that power..
We don't get caught up in the status of that power..
And in the process, we change the paradigm of how we perceive power and how people perceive.
power..
Jesus gives us a new framework of power..
And that framework of power is not to give up our power..
A lot of times people interpret that when he took off his robe that he was giving up.
power..
We do not give up power, but we steward the power that we are given for the gospel..
And that particular incident was so powerful, we still talk about it to this day..
We're still preaching to this day..
Universities, MBA programs, professors are still talking about what it means to be a.
humble servant, that one of the best ways of a CEO is to be the servant leaders..
Whether it's Christians or non-Christians, many of them will tell you, you have to be.
a servant leader..

$^{361}$Because this framework became so successful, it changed the trajectory of how we look at.
leadership and how we steward our power well..
Jesus gives us a framework of power..
We're not called to give up our power, but steward the power you and I are given for.
the gospel..
Let me give you some final thoughts, and I want to show you one video, okay?.
Number one, application time..
Many of us as young people, I say this to young people, technology has completely created.
the ability for all of us to be an influencer..
Some of you guys who are here, or those who are watching, technology has elevated all.
of our platform where any individual can now jump on YouTube or Facebook or Instagram and.
whatever it is, and create a video..
And you have the ability to be an influencer..
Many of you guys who are sitting here right now are talented and bright enough to be the.
CEO and executive of your company..
You will have more riches than I will ever imagine..
You will have more power than you will ever know what to do with it..
My biggest encouragement for you, make a choice of what you're going to do with that power..
For instance, a lot of times when money, I heard this from one donor who taught me well,.
and this donor says to me, Tommy, as you get older, and before you even become rich, figure.
out what you're going to do with your money..
A lot of times people always say, well, you know what, I'm going to keep saving my money,.
I'm going to make sure I have kids, I'm going to put away from savings, all of that stuff,.
and as I get wealthier, I'll start giving away and all that other stuff..
And he says, no, no, no, Tommy, before you even become rich, before you even get that.
great job, figure out and make a decision of what you're going to do with your resources..
Why?.
Because when you become wealthy, and when you get that power, you're going to be too.
tempted to serve yourself..
Make a decision in your head right now, what you're going to do with that money, what you're.
going to do with that power, what you're going to do with your life, what you're going to.
do with your skills..
Make a decision and commit to it..
Point number one..
Point number two, continue to learn to submit..
As you continue to grow in your power and your influence and what you're able to do,.
learn to continue to submit to the people around you, and also learn what it means to.
continue to grow in humility..
I say this one thing because I have this one girl in Jakarta, bright girl, very, very bright.
girl, very talented..

$^{401}$She's about 24 years old..
She even made it to Miss Indonesia..
She's pretty talented..
And in the midst of that conversation, you know what she says to me?.
She says, "Tommy, I hate when people say no to me..
I disregard anyone who says no to me.".
And we've had to have a conversation that sometimes as you grow in your power, you have.
to learn to submit and surround yourself with people who say no to you..
That sometimes in the midst of our power, you got to make a decision, but you got to.
continue to maintain that humility and continue to really be able to learn to submit to the.
people around you..
Number three, application..
Learning from Jesus doesn't guarantee you will make the right decisions..
It's the decisions that you have to make each day..
Why is that?.
That who was present at this last supper?.
You had Judas who literally watched what Jesus did, watched what it meant to be a servant,.
and what does he do later on that day?.
He sells out Jesus for pennies on the dollar..
He betrays Jesus..
Who's also sitting in that room?.
Peter..
And what does Peter do?.
He betrays Jesus..
Learning from Jesus does not guarantee you're going to lead the correct life..
It's a decision that you have to make each and every single day because we are written,.
we are living by the gospel..
The gospel guides us in all that we do, and there's a choice that we make each and every.
single day..
And the one final illustration, when you serve as a servant, don't ever expect anyone to.
return that favor back..
Do not have any expectations of what people can do for you..
Too often, a lot of times, when we do something for someone, what do we do?.
We expect it to be done in return..
That when you and I use our power, our God-given power, our God-given talents, and our God-giving.
resources, don't expect anything in return..
Because when we start expecting things to return, you're going to be sorely disappointed..
Serve out of the gratefulness of your heart..
Serve out of the glory of what the grace of God has done for you..
Serve out of the joyfulness of your God..

$^{441}$We understand the power that we have given..
We understand who has given us the power..
We make a choice each and every single day to steward that power..
We don't care about status..
And in the process, we change the paradigm of how people perceive power..
Because Jesus gives us a new framework of power..
We use that power and we ignore that power, but we steward the power that we are given.
for the gospel..
Let me show you one person who did it really, really well this past week..
For those who do not watch football, and for those who are watching American football,.
there was a situation that happened on Monday night football, where one of the players from.
the Buffalo Bills tackled one of the opposing players on the Cincinnati Bengals..
The guy stood up, the guy who tackled, stood up for a minute or so and then collapsed and.
went into cardiac arrest..
Game was canceled and we didn't know if the player was going to make it or whatever it.
is..
And all throughout the entire football community came and offered up prayers..
But out of all that prayer, one person, a Christ follower, who had a platform, set an.
example and showed us what it meant to do use his power well..
Phoebe, let's show that video..
>> Football gave me everything..
And I think even through the midst of absolute tragedy last night, I think you saw some of.
the beauty of football as well that it's brought us all here together..
This is a little bit different..
I've heard it all day, like thoughts and prayers..
You just heard Sheriff and Jonathan Allen say all we can do is pray for them..
I've heard the Buffalo Bills organization say we believe in prayer..
Maybe this is not the right thing to do, but I want to pray for them..
>> I'm going to pray for you..
>> What decisions are you going to make to impact your community, to impact this church, to impact your city?.
It's a decision that you're going to have to make each and every single day..
It's a decision that sometimes in the midst of understanding power you're going to make that doesn't benefit you..
It's not going to reward you in any way..
And that's what we're called to do..
We're called to be stewards of that power..
We're called to be sacrificed because we have the greatest example of the person who sacrificed his very life to die for our sins..
It's a choice that he had to make..
It's a choice that all of us as Christ followers who are saved by grace each and every single day, we make that choice to continue to be an influence in all that we do..
To live by integrity..
To embrace that power..

$^{481}$To understand that power..
To not care about status..
To love every single person and to serve the world as we know it..
Let's pray..
Father, thank you for these wonderful people in this wonderful church..
Thank you for bringing home Mark safely..
Thank you for those who are here for Christmas vacation..
We thank you for people like even Dan Orlovsky who uses his platform as a Christ follower to continue to show us and demonstrate what it is..
But dear Jesus, you are the model of what that power and sacrifice is..
You showed us not only what it meant to wash the feet of the disciples..
You gave us a new paradigm of power..
But you took it to a new extreme and died on that cross so that we are here..
And we still talk about that powerful servant attitude..
You give us inspiration and you give us hope..
Let us make those decisions every single day to honor you..
In Christ we pray..
Amen..
Thank you, Father..
Thank you so much, Tommy..
God bless you..
And I want to ask you for prayers as we meet today, this afternoon, with the old and new Quartermembers..
We're going to try the course for the next year or so and I pray that you will pray for our leaders as well..
Okay?.
So God bless you..
\newpage



\section{雅各書 4:13-17}
\label{sec:DsrtebXqSNI}
\textbf{計劃明天(雅各書4\_13-17)- 陳崇基牧師}
\newline
\newline
連結: \href{https://youtube.com/watch?v=DsrtebXqSNI}{\texttt{https://youtube.com/watch?v=DsrtebXqSNI}} ~~~~ 語音日期: 2023-01-17
\newline
\newline
\hyperref[sec:ZVIv_jPkXmQ]{\small{< < < PREV SERMON < < <}}
~
\hyperref[sec:index_chronic]{\small{[返順時目]}}
~
\hyperref[sec:index_scriptual]{\small{[返順卷目]}}
~
\hyperref[sec:daNqcFTl4KQ]{\small{> > > NEXT SERMON > > >}}
\newline
\newline
雅各書 4:13-17
\newline
\begin{longtable}{cl}
\hline
\hline
章節 & 經文 (和合本修訂版)\\
\hline
4:13 & \begin{tabularx}{0.7\textwidth}{X} 注意！有人說：「今天或明天我們要往某城去，在那裡住一年，做買賣賺錢。」 \end{tabularx} \\ \\ \relax
4:14 & \begin{tabularx}{0.7\textwidth}{X} 其實明天如何，你們還不知道。你們的生命是甚麼呢？你們原來是一片雲霧，出現片刻就不見了。 \end{tabularx} \\ \\ \relax
4:15 & \begin{tabularx}{0.7\textwidth}{X} 你們倒應當說：「主若願意，我們就能活著，也可以做這事或那事。」 \end{tabularx} \\ \\ \relax
4:16 & \begin{tabularx}{0.7\textwidth}{X} 現今你們竟然狂傲自誇；凡這樣的自誇都是邪惡的。 \end{tabularx} \\ \\ \relax
4:17 & \begin{tabularx}{0.7\textwidth}{X} 所以，人若知道該行善而不去行，這就是他的罪了。 \end{tabularx} \\ \\
[1ex]
\hline
\hline
\end{longtable}
$^{1}$歡迎大家,以後時間交給Mark Sir.
(字幕提供).
大家好,剛剛我做了兩個星期假,回來嘗試去休息一下,因為放假之後又要再放假.
大家都有詢問過,大家很關心到底那兩個星期假發生了什麼事.
也有很多弟兄姊妹很關心和禱告,在這裡特別多謝各位.
今天我想嘗試講一下我的經歷,然後看看有什麼可以學習到,我真的學習了很多東西.
大家知道我很少放長假,這次放長假是一個很大的經歷.
話說第一部分我先講我的經歷,然後再講一下我自己學到什麼.
希望在這個經歷裡面大家都能夠學習.
10月25日我們有聖誕節崇拜,我們完成崇拜的時候已經很快去收拾東西趕著走.
因為那天大概3點多我就要去機場,因為聖誕節崇拜,又要拿禮物,又要玩White Elephant.
要拿完禮物才走,我回去的時候已經差不多2點多,我是完全沒有收拾東西.
我一家大小有4個人,3個都已經收拾好東西了,我隨便找一些衣服,然後就丟到行李箱裡.
然後拿完證件,拿完信用卡,最重要是拿信用卡.
然後我就馬上叫了Ivan,就載了我們去機場.
原本我們想著可以私私然去,誰知道我們的親友就一味去Tag我.
大佬,你知不知道Southwest發生什麼事?.
因為我們搭Southwest的航空公司有很大的Meltdown.
我一邊檢查那個App的時候,就發現那個飛機都OK的.
他還延遲了半個小時,延遲了半個小時我就說好啊,因為我趕嘛.
我打算慢慢來,誰知道他的親友就一直在提醒我,你肯定你的飛機沒事的嗎?.
我一邊檢查的時候,他就一直延遲了15分鐘,再延遲15分鐘.
我說不行,然後我就上去電腦再檢查.
原來那個電腦上檢查的時間和我手機上檢查的時間是兩回事,差一點點.
我說不行,我就跟家人說,我想我們一定要早點去.
然後我們就去那個機場問清楚他,到底那部飛機是怎樣.
因為我原本是從芝加哥飛到菲尼克,然後菲尼克那邊不夠一個小時就要轉機.
就要搭Sacramento,因為第一部飛機的第一站就一直延遲.
所以我說如果趕不到第二部飛機,那我就很糟糕.
所以我們就叫了Ivan載我們去Midway的機場.
我們一下車的時候,我們去到那個機場已經大排長龍.
那裡起碼有幾百人,我們就馬上,總之有龍我就排隊.
排了的時候,我就發現原來所有的人都是嘗試去重新訂機票.
嘗試去再重新訂機票,因為每個人都是大混亂.
我們查了資料的時候,原來很多的飛機都是延遲的.
我們就排隊,誰知道那條隊很慢.
我們三點半鐘,十二月二十五號就去到機場.
就等到去到Counter那裡,Counter那裡只有幾個人在工作.
我們就一直在等,等到大概五,六個小時之後.
我們終於快要去到Counter,快要去到.

$^{41}$你知道我剛才說只有三,四個人在工作,那些Southwest都不夠人工作.
結果快要去到Counter的時候,已經看到Counter.
那個經理人就出來,他說不好意思,我們電腦崩壞了.
大家忍耐一下,我說你已經延遲了很久.
在那段時間我已經明白,我根本趕不上第二部Fenix飛到Segmento那部機.
因為那部就算怎麼延遲都好,都不會等到我,那時候已經八點,九點.
我已經在機場站了五,六個小時.
我就做一下Research,然後我就有一個想法.
如果不飛就沒有機會飛了,因為如果他這樣大混亂,Southwest meltdown.
如果回家的話,其實就不用去了.
我就想了想,然後我就做一下Research.
我兒女就跟我老婆說,我說我有一個瘋狂的想法.
在那個時候,很少一家人都這麼留心聽你說話.
那個是父親的時刻,我有一個瘋狂的想法.
我就聽,我的瘋狂的想法就是,我們先飛到Fenix.
然後在那裡我們才決定怎麼去Segmento.
我們飛到Fenix,因為我們有一個同工在Fenix,Joyce在Fenix.
我發訊息給她,可不可以在那裡寄居一晚.
我就跟兒女說,既然你都沒有去Las Vegas,不如我們去Fenix.
然後我們再租車,就開車去Las Vegas.
在那裡酒店留一晚,反正你都沒有去過拉斯維加斯.
去一次吧,因為我們兩夫妻都去過.
我說不如這樣,我們去租車,然後開五個小時車.
由Fenix開到拉斯維加斯,在那裡住一晚,住酒店住一晚.
然後我們再開九個小時車,再上去Segmento.
他們說,來吧,因為他們不知道後果是怎樣.
只不過我心裡盤算,我覺得這個計劃是可行的.
結果我們去到那個Countdown的時候,我們就跟那個小姐說.
不要理,我知道我們第二次租車就取消了,不過讓我去Fenix.
她也有想過,小姐說,你去到Fenix怎麼辦.
我說不要緊,我們會處理好,我們有個好朋友.
你讓我們飛去Fenix,結果她就把我們的票再重訂.
然後我跟她說,請確保我們五件行李全部標籤在Fenix機場.
我看著她標籤,問了她好幾次,她說就是去Fenix.
我說OK,我們就飛去Fenix,然後我馬上訂了租車.
但當我們去到機場的時候,你知道進到去西西那邊的機場.
已經大混亂,每個人都在等他們的飛機.
有些人就轉了幾次閘口,最後就取消了.
最後就回家,真是麻煩,我們進到閘口就繼續等.
我們上到機的時候,應該是大概12點鐘.

$^{81}$所以我們3點鐘去到機場的時候,去到閘口.
等待上機的時候,已經差不多12點鐘.
真是慘,因為我們沒有東西吃.
幸好我靈機一觸,還沒出門之前,我就叫他們帶好食物.
因為將會是非常走藍色的,去搭飛機.
真是走藍色,去到閘口的時候.
因為每個閘口都在等收回食物,有很多都是失敗.
結果那個人叫我們可以入閘,那些人鼓掌.
我沒試過在美國機場有人鼓掌.
我們就排隊上機,上到機的時候已經是12點鐘.
已經是午夜了,在那裡等了9個小時,終於上到飛機.
那個機師還說,這次為什麼可以飛呢?.
因為剛剛有一個飛機乘客願意去飛.
他願意去飛,他特意去飛.
所以剛剛湊夠人,你的飛機才可以起飛.
所以我們的方案就是這樣.
我的瘋狂的想法就是,先飛去菲尼斯.
然後那裡才決定怎樣搞.
你覺得這是不是很好的方案?.
其實不是最好的方案,因為是唯一一個方案.
唯一一條路我只能夠選擇.
所以我們就上飛機,飛到菲尼斯大概3個小時.
我們到了菲尼斯的時候,已經是當地的大概凌晨3點鐘.
所以我已經糾纏了12個小時左右.
到了機場,我覺得安心一點,就去拿行李.
等那些行李出來,哇,那麼少行李.
等了很久,都沒有我們的行李.
因為是全程的,因為是全程的.
還有其他乘客都是沒有行李.
所以我不知道是很生氣還是很好笑.
好笑的就是有人陪.
在菲尼斯的閘口,領行李的閘口.
那裡放了幾乎一千個的行李.
有些是比我們早一點到的.
我一邊看一邊看,我心想這次死定了.
真的沒有,我們又跑去排隊.
跑去排隊就是跟行李箱說,行李不見了.
結果他發現原來行李還在地下.
我只有一個標籤,只有五個標籤,五件行李.
拿著一張紙,他說你回去等一下,我們看看可不可以寄來.

$^{121}$原來那五個行李沒有跟我們去菲尼斯.
結果我們什麼都沒有,只有手提行李.
但幸好我寧記一足.
我們臨出門口之前,我說麻煩你們把底衣褲多撿一套.
OK,幸好我們就是多撿一套.
所以我們就在那裡可以暫時過一晚夜.
我們去到貨物車,然後去到租車.
已經三點半了,四點鐘AM.
整個人都殘了.
去到租車的停車場,為什麼沒有什麼車.
又等,原因就是其實每個人都在租車.
等到那個人終於來到我們面前.
我們是enterprise的,我們沒有理會車的尺寸.
給你一個free upgrade.
他就說了一堆我們有的車.
我一聽到camera就說給我camera.
我就上了camera.
終於由一個我不能控制的環境.
上到車我就可以自己控制.
我們就馬上飛車去到Joyce那裡.
幸好我們同工在那裡.
如果不是我真的想像不到有其他的乘客.
如果他們真的是出外去旅遊.
沒有地方住的時候.
你要住酒店,衣服又洗不到.
當然我都沒有衣服洗.
是怎樣搞的.
所以在Fenix那裡留了一天.
其實都是一個祝福.
我們去探望Joyce的家.
探望了Abby等等.
是一個驚喜.
我們喜歡.
睡醒了.
我們就去不同的地方買東西.
買什麼?.
買衣服.
買了很多衣服.
買到我今天回來發現我的抽屜不夠位置.
我就想要買一個新的抽屜.

$^{161}$我們就在Fenix留了一晚.
我們就開車去Las Vegas.
星期一晚上我們終於去到Las Vegas.
開了五個多小時車.
終於去到Las Vegas的酒店.
去到酒店.
一進去就排長龍.
又排長龍.
為什麼呢?.
因為前前後後聽見他們說話.
或者跟他們聊天的時候.
原來大家都是同命相憐的.
有很多人被迫留在Las Vegas.
等了一小時.
我才去到櫃面那裡.
拿到酒店的鑰匙.
是大酒店.
我們去到那裡終於可以睡.
我們就在Las Vegas.
都是一個驚喜.
正如我開頭說的那個瘋狂的想法.
他們可以去到拉斯維加斯.
我們就在那裡走走拍拍照.
去了Las Vegas.
在Las Vegas我沒有賭錢.
原因是什麼呢?.
因為整個旅程都是一個賭注.
所以是一個Gamble.
我們在Las Vegas.
結果我們就.
星期二下午.
我們就決定離開.
我們開了十個小時車.
由Las Vegas經過Bakersfield.
然後北上在I-5上到Sacramento.
開個人車很辛苦.
為什麼呢?.
因為遇到在加州西岸有大風雨.
所以整個車都是下雨開車的.
並且有很多地方是山路.

$^{201}$我已經老了.
有一點腦眼分化.
晚上看東西是看不清楚的.
我記得有一段路是山路.
我是看不到那些線.
因為雨太大.
因為我租的是Camry.
是新車.
很好開.
我開啟它的Lane Detection.
但Lane Detection是不可行的.
因為Lane Detection都看不到那些Lane.
很辛苦開了十個小時車.
結果我們終於去到Sacramento.
都是一個祝福.
安全抵達.
我們見到很多親友.
有些親友因為疫情的緣故.
我們幾乎六年沒有見面.
六年七年.
或者有些教會的弟兄姊妹.
去到北加州Bay Area等等.
都是一個祝福.
雖然很辛苦.
我講過很少放長假.
弟兄姊妹經常說.
不如你放長假吧.
你這麼辛苦.
我終於知道為什麼我不放長假.
現在放長假就變成這樣.
下次放長假真的要想清楚.
但很感恩.
見到很多很久沒見的朋友.
弟兄姊妹.
聽到一個這樣的經歷.
一個這樣的故事的時候.
或者自己再回想.
再去評估每一個步驟的時候.
我覺得有很多恩典.
因為和每一個弟兄姊妹見面.

$^{241}$去到reconnect或者new connection.
都是一個恩惠.
可能網上的一些可能是加州的朋友.
我們有機會在短短兩個星期見到面.
我覺得都是一個恩典.
有些是第一次見面.
有些是沒有見面很久.
我覺得是一個gracious moment.
很感恩.
但整個旅程裡面.
我不時想起.
這一段聖經.
就是雅各書第四章.
這裡和修本和合本的修訂版.
這樣說.
第二.
有人說今天或明天.
我們要往某城去.
在那裡住一年.
做買賣賺錢.
其實明天如何.
你們還不知道.
你們的生命是什麼呢.
你們原來是一片雲霧.
出現片刻就不見了.
雅各書.
大家可能很清楚.
是一本很practical的書.
只有五章.
雅各作為一個牧者.
在耶路撒冷的教會.
他寫給那些可能被迫害.
被教會迫迫.
然後他們散居在不同的地方.
小組的教會.
弟兄姊妹.
他寫了一封信.
或者集合了一些的港獨.
然後寄去給散居在不同的地方.
神的子民.

$^{281}$在第四章這個地方.
他就去針對某些人.
針對什麼人.
生意的人.
商業的人.
想做生意.
然後可以買賣.
可以賺錢.
這些在雅各書的世界裡.
這些很多時候都是一些.
因為有錢.
因為有資源.
所以是比較高傲的一些.
生命的取態.
很多時候沒有放上帝在眼中.
雅各在雅各書裡面.
寫了好幾句的說話.
或者好幾段的教導.
去對抗這些人.
這段經文說什麼.
我們要去某個地方.
我們要留在某個地方.
將來式.
我們要做生意.
我們要賺錢.
你有沒有留意他.
We will and we will.
雅各就是針對這些人.
因為見到自己有財有勢的時候.
他就做很多盤算.
做很多的打算.
就說We will do this or that.
將來.
好像是他能夠控制到.
雅各就是去教訓這些人.
告訴他們.
當你有資財的時候.
到底上帝是不是在你心靈裡面的圖畫.
當你去盤算.
下一步我要做什麼.

$^{321}$今天 明天我要做什麼的時候.
有沒有將上帝放在眼裡.
他說什麼呢.
你們完事一片雲霧.
出現一片黑就不見了.
給他一個當頭的棒.
你可以計劃.
不過你有沒有想過.
你自己生命裡面的本質.
其實你只是miss.
morning fog.
一片的雲霧.
最好你就是順服.
最好你就是謙謙卑卑.
告訴上帝.
其實你不是真正in charge.
你只不過是管理.
你是steward.
是管理.
不是去到擁有明天.
我們在拉斯維加斯的時候.
星光大道.
你見到有很多那些名人.
Walk of Fame.
我們見到Alphys Presley.
最近知道Alphys Presley的女兒.
唯一一個女兒.
突然去世.
很年輕.
生命就像雲霧.
我們有機會去到三藩市.
你知道三藩市去金門橋.
有沒有去過.
去金門橋不是經常拍到照片.
因為那裡很煩.
經常出霧.
霧一去就可以拍.
霧就是這樣.
你捉不到.
你是不可以擁有.

$^{361}$你怎可以擁有雲霧.
它出現一段時間.
然後很快就離開.
我們沒有絕對的control.
對我們的life對我們的生命.
是沒有絕對的control.
為什麼我經常想起聖經.
因為我真的沒有一件事可以控制.
我真的沒有一件事可以控制.
我控制到的是什麼.
就是租一輛車.
可以開車位.
但是如果天氣不好.
我都開不了車位.
有幾個moment.
我幾乎要停車.
因為我不能夠再去.
雨太大 太黑.
你不能控制.
生命裡面想清楚.
亞國其實就是想將現實.
現實去告訴我們.
有些事我們真的不能夠去控制.
芝加哥的大影子會知道.
我們的天氣.
是很反常的.
反常到什麼地步呢.
我們.
這個的.
picture.
沒有了.
剛才我有一個forecast的picture.
那個picture說什麼呢.
就是今天是50度.
明天是多少度.
3度.
今天是陽光普照.
芝加哥明天就會下雪.
就是天氣我們都不能夠去控制.
最近有些電影節目跟我們說.

$^{401}$連蛋的價錢.
我們都控制不到.
對不對.
最近我看到一個網上.
講的一幅相.
有個男孩拿著.
幾打蛋.
就因為情人節.
情人節送蛋是最好的.
經濟我們不能夠去控制.
生命裡面有很多的.
元素我們不能夠去控制.
當然你說.
亞國是不是.
什麼都不要計劃.
他不是說不計劃.
是需要計劃.
只不過我們的取態要想清楚.
你作為一個管家.
我作為一個管家.
我們需要想清楚我們的計劃.
這個是我們的責任.
但是我們不能夠有一個.
錯誤的取態.
他在這段聖經說什麼.
他說你們倒應當說.
主若願意.
我們就能活著.
也可以做這事和那事.
然今你們既然狂傲自誇.
凡這樣的自誇.
都是邪惡的.
所以人若知道.
該行善而不去行.
這就是他的罪了.
亞國書.
其實就是新約裡面.
有點像舊約裡面的.
智慧文學.
真言傳道書等等.

$^{441}$裡面有很多的教導.
其實就是給我們一些生命裡面的.
wisdom.
生命裡面的智慧.
其中一個最重要的.
亞國書.
智慧文學要說的.
就是生命裡面應該有的取態.
敬畏上帝.
謙卑在上帝面前.
明主權.
是上帝的旨意.
高過你的旨意.
是上帝的主權.
高過你的權能.
上星期一晚.
我們英文部開會.
我們同工說開一些趣事.
有什麼趣事呢.
就是我們每個人都是開車的.
我們曾經都教過人開車.
教人開車.
最害怕的是什麼.
就是那個人不聽你的話.
還有我們不是那些教車師傅的車.
我們不能自己踩.
我們就說開.
有些朋友我們去教開車的時候.
我們覺得很害怕.
我不知道你有沒有試過.
坐在乘客的位置.
那個司機開車的時候.
令到你很害怕.
這要看.
你開車的時候.
這要看.
如果你坐在我旁邊.
我開車的時候你害不害怕.
可能都害怕.
不過是另外一種害怕.

$^{481}$開得快的那種.
那種害怕.
但如果你想想.
是上帝開車.
是上帝開這部車.
你可能試過坐一部車.
那個司機是很順得過的.
是可以在那裡睡覺.
是可以安然.
又平安.
這就起不來我們心裡應該有的取態.
上帝是.
在司機的位置.
我們只不過是坐在車上.
你跟司機說.
Take me according to your will.
按照你的心意.
帶領我.
去到我的目的地.
因為我信得過你.
上帝的主權.
人的謙卑.
結束的時候.
我想用施普珍一句說話.
他說到上帝的主權.
上帝的sovereignty是甚麼.
施普珍.
這位已經百多年前的.
一位廣告家.
他說甚麼呢.
他說當你經歷試煉的時候.
上帝的主權.
就是你安然.
拯守的地方.
The sovereignty of God.
is the pillow.
upon which you can.
lay your head.
上帝的.
心意 上帝自己的.

$^{521}$微而 祂的主權.
帶領我們.
我們儘管謙卑在祂面前.
去信服.
讓上帝帶領我們心靈.
我不知道今天你自己經歷過的是甚麼.
可能未必好像.
我兩個星期長假.
所經歷的那麼多喜福.
亦都可能在你心靈裡面.
經歷很多不同的試煉.
很多trial and temptations.
記住上帝的主權.
去到帶領我們.
我們請離見八對再上來.
我們唱一首回應詩歌.
這首詩歌.
帶領我們來到主人面前.
讓我們知道上帝.
是我們心靈的主人.
讓上帝帶領著我們.
我們不會奏詐.
但是上帝要給我們.
信服的心.
我們的取態要一致.
要聽清楚.
上帝每一步的帶領.
以致我們的心靈.
走在上帝所開的路程上.
《上帝的心》 詞:陳家麗 曲:陳志遠.
上帝的心 詞:陳家麗 曲:陳志遠.
上帝的心 詞:陳志遠 曲:陳志遠.
《上帝的心》 詞:陳志遠 曲:陳志遠.
一步又一步 這是恩怨之路.
你愛你說 將我緊緊抓住.
一步又一步 這是盼望之路.
你愛你說 牽引我追著人生路.
你是我的主 引我走這一路.
高山或低谷 都是你在保護.
萬人中唯獨 你愛我認識我.

$^{561}$永遠不變的事 這一切都是祝福.
一步又一步 這是恩怨之路.
你愛你說 將我緊緊抓住.
一步又一步 這是盼望之路.
你愛你說 牽引我追著人生路.
一步又一步 這是恩怨之路.
你愛你說 將我緊緊抓住.
一步又一步 這是盼望之路.
你愛你說 牽引我追著人生路.
我們來到神面前 待會兒會祈禱.
給大家一點機會 將你的前路放在主的身上.
也求主給我們一個信服去帶領我們上去.
想向你道謝 因為你是我們人生路上的主.
縱使我們走的人生路上 可能有些迂迴曲折 高山低谷.
但我們深信你是那位 滿有美意 並且滿有權能 滿有主權的上帝.
生命中我們可能遇見不同的問題 未必能夠有答案.
但我們知道每一個彎位 有你自己引導.
每一個低谷 有你自己陪伴.
主要是為了每一位弟兄姊妹 我們在心裡可能我們籌算著 計劃著.
或者很關注某一些決定 或者將來我們有一些未知之數.
我們未知道下一個彎位會是怎樣的.
我們深信你自己那一位 是滿有主權來帶領我們的上帝.
並且你不會令我們走差路.
神啊我們求你 聖靈給我們一個更新的上帝.
去信服你自己帶領我們.
謙卑我們 讓我們聽清楚 聖靈對我們說的話.
引領我們走每一個的一段路.
每一個的彎路 每一個的低谷.
主知道你是那位願意 跟我們共行的聖人.
我們感激到 奉主的名.
願上帝平安 治理大事.
\newpage



\section{}
\label{sec:daNqcFTl4KQ}
\textbf{Becoming A Healthier Church 1 (Ephesians) - Rev. Mark Chan}
\newline
\newline
連結: \href{https://youtube.com/watch?v=daNqcFTl4KQ}{\texttt{https://youtube.com/watch?v=daNqcFTl4KQ}} ~~~~ 語音日期: 2023-01-24
\newline
\newline
\hyperref[sec:DsrtebXqSNI]{\small{< < < PREV SERMON < < <}}
~
\hyperref[sec:index_chronic]{\small{[返順時目]}}
~
\hyperref[sec:index_scriptual]{\small{[返順卷目]}}
~
\hyperref[sec:hnLQJ3HSDsE]{\small{> > > NEXT SERMON > > >}}
\newline
\newline
$^{1}$Let's do it together..
All right, let's welcome Pastor Mark..
- Hey, happy new year..
Happy new year to all of you..
And such a good day to be here..
I know it's snowing outside, right?.
It's still snowing outside?.
'Cause I've been here since 8 p.m..
So I don't know what's going on outside, you know..
Huh?.
- 8 p.m..
- 8 a.m., sorry, 8 a.m., yeah..
It's been a long day for me,.
and it's been a long week for me..
So good to see all of you..
And why don't we shake hands with our neighbors.
and say happy new year,.
or words of blessing to each other..
Find five neighbors, five neighbors around you..
Yeah..
All right, all right, yeah..
Well, there's some new friends in the midst of us,.
and we wanna welcome you to be in our church..
Let's round of applause for all the new friends.
who are visiting us the very first time, all right?.
(audience applauding).
So it's still new year..
You know, the atmosphere, the color, or whatever.
reminds us what it means to have new year,.
and what it means to have change..
What do we do as new year kicks in,.
and a new period of time, a new season in our lives?.
I think this is a great reminder for us to just rethink,.
you know, where are we going?.
So we may think of new year's resolutions, right?.
New year goals, all right?.
So let me do a survey here..
How many of you have your new year resolution,.
and one of your resolution is to lose weight?.
How many of you?.

$^{41}$Yeah, that reminds us, you know,.
this reminds us we have to do it every single year, right?.
Although we already fail, we already fail,.
but that's okay, you know?.
We continue to grow, right?.
Not in weight, but in our spirit, all right?.
You know, some of the new year resolutions,.
some of the new goals that we're gonna set for ourselves,.
or some of the new mindset we want to have..
So this new year will be reminding us, you know,.
doing something that will require us to change, right?.
So what it is for you, think about it..
And today, I would like to start a new series.
on the sermon, which is about church,.
about healthy church, healthier church..
How do we become a healthier church?.
Think of yourself, right?.
You have been a Christian for a long time,.
or maybe a short time, or perhaps you stumble and fall,.
stumble and fall, and along the way, you know,.
you still find your identity in Christ, right?.
And how to be a healthy Christian,.
what does it mean, healthy Christian, right?.
You may have your own, you know, imagination,.
what it means to be a healthy Christian,.
what it means to be a healthier church,.
but you know, we want to, you know,.
go into the word of God, and see what God has to say.
about being a healthier church, right?.
What does it mean to be a healthier Christian?.
Now, as we talk about, you know,.
new year's goals and resolutions,.
you know, we always think of, you know,.
there's some maybe poor habits we need to do away..
Some new habits we need to redevelop, right, to pick up..
Something like that, spiritual disciplines, right?.
You know, things like we have to change our lifestyle.
in order to be healthier,.
not just physically, but spiritually as well..
So think about, as an individual question,.

$^{81}$how we should change in such a way to be healthier,.
and then as a group, as a corporate, as a body,.
as a community, how do we grow together.
to be a healthier church?.
So that is the new topic we need to talk about..
Now, why do I want to talk about this?.
Because, you know, as you think about church.
as a community, as a family, we talk about relationship..
Unless our relationship is healthy,.
our church will not be healthy..
Unless your relationship with other people,.
Christians, non-Christians, brothers and sisters in Christ,.
you know, your siblings in your family,.
your parents or your children,.
unless your relationship is healthy,.
there's no way you can walk.
in a spiritual pathway of God, right?.
So we have to dig into how do we grow our relationship.
in a healthier way in order to be a healthier Christian.
as an individual and as a group, as a church..
In fact, if you read through the Bible,.
there are many places in the Bible.
that describes church as a family..
You know, in the first century, 2,000 years ago,.
in the time of Jesus, in the times of the first Christians,.
the early Christians, all churches are house churches..
So they gather together, not in a separate building,.
but let's say we are the sample here in the first century..
We're gonna gather together for worship,.
for fellowship on Sunday, on the Lord's Day,.
and we're gonna ask around different people,.
you have a house big enough for us to squeeze in, right?.
To sit together, to sing together, to pray together,.
and to study the word of God together is house church, right?.
It's pretty much the situation in many parts of China..
Right?.
Having churches not in a separate building, but in houses..
So when the apostle talks about church,.
he actually described the church as family,.
literally as a family,.

$^{121}$because the church will gather in your living room.
as a family, right?.
And all the members belong to the same family,.
the big family of God, right?.
Where you have, of course, your brother and sisters,.
someone who is directly related to you,.
but many of the other people.
who will squeeze into your living room in the house church.
are related in Christ, right?.
Sister and brothers in Christ..
So in the days of Paul, all churches are house churches,.
and the family is the basic unit,.
not only just the society,.
but for church, capital C, church..
That's why in this verse, Ephesians chapter two,.
it says you used to be Gentiles,.
not being part of the kingdom of God,.
not being related to the people of God,.
but now in Christ, you are no longer strangers,.
you're no longer foreigners, you're citizens..
You are the people, along with all the holy people of God,.
and you belong to the family of God..
You are members of God's family..
That's why we need to talk about family relationship, right?.
Because I really believe that.
if our relationships are strong,.
if we are closer, if we grow together,.
then our church will grow..
Now, if we don't grow in our relationship,.
if we don't grow mature in our relationship,.
our church will be stifled..
So we need to talk about relationship,.
just as Paul, writing Ephesians, talking about relationship..
Now, 2022, last year,.
there was a new survey done in United Kingdom, in England..
So the survey asked a couple of questions..
One of the question is,.
to ask different people on the street,.
and ask them, how many of you feel that the pandemic.
has improved your family relationship?.

$^{161}$How many of you feel that the pandemic,.
which lasts for how many years?.
Seems like forever..
Three, four years, right?.
Three years, right?.
How many of you feel that COVID,.
the pandemic, has improved your relationship,.
particularly your marriage or your family relationship,.
parents and children and all that?.
How many of you feel that COVID has,.
three years, have improved your relationship?.
And the result is, take a guess..
How many of those people, what's the percentage.
that would say COVID will improve our relationship?.
Exactly, one third only, one third only..
Is that exciting?.
Not very exciting, right?.
Pretty sad..
The second question is,.
they survey some of the largest family law firms in UK.
about the increase of divorce.
in the past three years of pandemic..
And these law firms submit their results,.
and there was increase in terms of cases in divorce..
Take a guess, percentage..
How many percentage increase because of COVID for divorce?.
50, any?.
I tell you, it's 95\%..
Is that exciting?.
Is that exciting for lawyers?.
Not for the families, right?.
Not for the marriages..
It is indeed very, very sad,.
because in reality, the pandemic.
has a lot of negative impact on our relationship..
Now, think about it..
And we go to the Word of God and say,.
what does God want, right?.
What does He require of us?.
How can we be renewed in our relationship, right?.

$^{201}$Because we are family of Christ, right?.
So I've chosen a passage today from the Bible,.
a very familiar passage..
And I wanna think about this passage.
and what it means to be a family in Christ..
What does it mean when Christ comes into our relationships,.
all of our relationships?.
Now, my relationship with my spouse, my children, if any,.
my brother and sisters, my in-laws,.
or my, any of our relationship in church.
or outside of the church, right?.
So let's think about it..
But as I read through this passage,.
I want you to ask yourself,.
as if you are reading this passage the very first time,.
as you listen to the words of this passage.
the very first time,.
how do you feel about this passage?.
How do you feel about it?.
All right, are you ready?.
You ready?.
I'm gonna read through Ephesians 5,.
verses 21 and following..
And further, submit to one another.
our reverence for Christ..
For wives, this means submit to your husbands.
as to the Lord..
For husband is the head of his wife,.
as Christ is the head of the church..
He is the savior of his body to church..
As the church submits to Christ,.
so you wives should submit to your husbands in everything..
25, for husbands, this means love your wives.
just as Christ loved the church..
He gave up his life for her to make her holy and clean,.
washed by the cleansing of God's word..
He did this to present her to himself.
as a glorious church without spot or wrinkle.
or any kind of blemish..
Instead, she will be holy and without fault..

$^{241}$28, in the same way,.
husband ought to love their wives.
as they love their own bodies..
For a man who loves his wife.
actually shows love for himself..
No one hates his own body, but feeds and cares for it.
just as Christ cares for the church..
And we are members of his body..
31, as the scripture says,.
a man lives his father and mother is joined to his wife.
and the two are united into one..
This is a great mystery,.
but it is an illustration.
of the way Christ and the church are one..
So again, I say, each man must love his wife.
as he loves himself..
The wife must respect her husband..
All right, we're done..
How do you feel about it?.
How do you feel about this passage,.
you know, reading it or listening to the word.
as if this is the first time?.
How do you feel about it?.
I tell you, as modern readers,.
there is something really hard to swallow, right?.
Some even find this passage oppressive,.
you know, because it seems like the passage.
is encouraging some sort of blind obedience, right?.
Especially for, on part of the wives, right?.
If you're a wife, if you're a female,.
you probably will read this passage as a modern reader,.
you will find trouble understanding what it means, right?.
Because we value 21st century, we modern people,.
we value equality, right?.
Husband and wives, you know, they are equal in status..
So you may feel that it is so old and traditional, right?.
Wives, you know, we believe are not subjects,.
but partners in marriage..
Amen?.
Amen, right?.

$^{281}$And we respect each other, there's no blind obedience..
And even the servants,.
the following context will tell us,.
even the servants need to be free, right?.
They need to be free men..
So how can this passage be related to us?.
That's the big question..
What is the passage really talking about?.
Now, I think you heard me before..
As we read the Bible, we need to understand.
the Bible is not written to us,.
but the Bible is written for us..
Got it?.
The Bible is not directly written to us, right?.
There's no way by saying this would devalue,.
you know, the sacredness of scripture..
Because as God revealed himself through scripture,.
he has to use some of the ideologies,.
some of the wordings that were familiar.
to the first audience, right?.
In order to communicate to them, right?.
Using the language or the cultural values.
that they are familiar with,.
and to communicate some eternal truth to them, right?.
So it's not directly written to us,.
it's directly written to them, using their language,.
using their categories to communicate.
some of the eternal truth that would be applicable.
for all cultures and all times and space, right?.
So it's not written directly to us, it's written directly.
to them..
So we have to understand the scripture.
or read the scripture as if we were.
first century Christians..
If we live in the Roman Empire, in the time of Jesus,.
in the time of Paul, then we probably will be better,.
on a better place to understand the real meaning.
of the scripture..
Now, as you can tell, you know, every society has their own,.
you know, understanding of the responsibility.

$^{321}$and roles of, you know, different members, right?.
Male are expected to do something, right?.
Or behave in such a way..
And female are expected to behave in some other way, right?.
Every society, you know, ancient or modern,.
there are always some sort of, you know, social categories.
to, well, you can say box in,.
but, you know, in some way it's to define who we are, right?.
Based on gender or based on, you know, different roles.
in the society, in the relationship..
Now, it is the same way in the first century,.
in biblical time, in the Roman Empire..
There was some sort of, you know,.
social understanding of how wife should behave.
and how husband should behave,.
how the master of the household should behave, right?.
It's just part of the social fabric of the time..
Now, reading this passage in the first century context,.
we can call it household codes,.
household regulations, household principles, right?.
I know you probably heard of a name called Aristotle,.
right, Aristotle..
Ring any bell?.
Aristotle is a great philosopher, a great teacher..
He wrote a lot of things about social relationship..
And in one of his writings called "Politics,".
Aristotle talked about household management..
And he said about regarding household management,.
there were three parts..
There were three different dynamics.
for every head of the household he has to manage..
The first one is the rule of a master over slaves,.
over the servants..
The second aspect is the rule of a father over the children..
And then the third one is a husband ruling over the wives..
Now, three categories or three aspects.
of a man's relationship in the entire household..
Now, this is called the household code, right?.
So that is the popular understanding on the street..
So if you go back to Rome in Jesus' day,.

$^{361}$in the first century,.
and you ask the people on the street, right,.
male or female, you ask them, "Hey, what's the role?".
"Oh, you are male.".
"What is your role and responsibility?.
"What is your duty in the household?".
And he would probably be able to cite you.
something like that..
Oh, there are three ways I need to relate.
to my household and manage them..
One is my servant, and one is my children,.
one is my wife, right?.
So they would be able to tell you that the popular idea.
that a male will rule over his household as a sovereign,.
almost like absolute authority..
And the male is ruling the subordinates,.
and the subordinates are the servants, right?.
The wife and the children..
So there are three parts, right?.
So the consensus is that the man is justified.
to rule over his household because wives, servants,.
and children are by nature inferiors, right?.
They are subordinates, they are to submit..
So the purpose of the household codes like this.
in Paul's day, in the first century,.
is to maintain a good order,.
a good hierarchical order in the Roman Empire..
The Roman society is peaceful, is secure,.
is all because all the household heads,.
the males, the sovereign in the household.
are managing the family properly..
Now with that understanding,.
you go to the book of Ephesians,.
you find something really interesting.
'cause Paul is almost citing or using some of these ideas,.
but he used it and twisted it in such a way.
and turned it upside down..
Now if you're familiar with the book of Ephesians,.
as Paul talks about the family church,.
the family relationship, he talks about three parts..

$^{401}$First, what we have just read,.
husband and wife relationship..
And then the second part, if you go on with the passage,.
he talks about a father and children..
And then he talks about what?.
Master and slaves, right?.
Master and servants..
But Paul is not just copycatting the ideas of his day..
He's using the idea of his day,.
the social understanding of his day,.
and put a new meaning in Christ..
Twist it and almost turn it upside down..
Let me tell you the context, okay?.
We just read 522 to 33..
If you keep on reading,.
he talks about other relationship, right?.
Now let me try to demonstrate..
If we were in the first century,.
hearing the words of Paul in the book of Ephesians,.
as if the first time, right?.
You're all familiar with the idea of a household.
from Aristotle, right?.
People on the street, they always assume that kind of idea..
The male head, all the other people in the household.
are subordinates, right?.
And then Paul comes in, and he's using a similar idea,.
and he talked to different parties in the household..
Now let me demonstrate how probably a male would feel,.
what he would feel like when he heard this word.
first time, first century..
522, Paul writes, "Wives, pay attention, listen up..
"You need to submit, right?.
"You need to do this and that, right, to your husbands.".
And then if you're a male in the congregation,.
hearing the words the very first time,.
what would be your response?.
I would imagine if I'm a male,.
and I hear this 522 to 24..
Wives, you need to submit to your husband..
I'll go, my response is, "Yeah, of course.".

$^{441}$Of course the wives have to submit to me, right?.
Because I'm the head of the household..
And the 525, Paul comes in and say,.
"But, but hey, husbands, the males, listen up.".
And there's a long verses, you know, 25 to 33,.
actually so much longer than the passage.
we've done for the wives..
Listen up, men..
Lord Jesus comes in..
He is the sovereign, not you, right?.
And this is what you do..
Remember what 25 to 23 says?.
Jesus died for the church,.
sacrificed himself for the church..
This is how you love your wife, right?.
And as a male reader, I would go, "What?.
"Why does it concern me?".
Because that's not what the slogan says, right?.
I mean, all the commands are given to the subordinates,.
right, to submit..
Not for me..
However, Paul turns it upside down and say,.
"Hey, hey, hey, hey, hey, men, slow down.".
Because there's a new sovereign coming into your household..
And you are now become a steward..
You're no longer the absolute authority..
You're accountable to Lord Jesus Christ..
And look at how he managed the household of Christ,.
the household of church..
He died and sacrificed, right?.
He loved selflessly..
And then six, one to three, children, right?.
Hey, kids, listen up..
You know, this is what you do..
You have to respect your parents and all that stuff, right?.
As a male reader, I would say, "Yeah, of course, right?.
"My kid has to submit to me.".
And then six, four, Paul comes in and say,.
"Surprise, you know, fathers, listen up..
"There's something for you to do.

$^{481}$"because you are now in the Lord..
"You are not the sovereign..
"Jesus is the sovereign.".
Same way for servants..
Servant, listen up, right?.
All the employees in your household,.
yeah, of course, they need to work hard and faithfully..
However, six, nine, Paul comes in and say,.
"Hey, masters, listen up as well..
"You men, you better behave.
"because Lord Jesus Christ is above you, right?".
There was a, can you see that?.
There was a, there was a,.
a fraudy reconstruction or reconfiguration.
when Christ Jesus comes into your family,.
into your church..
This is what we call subversive..
Paul is using similar familiar ideas,.
but he turns it upside down.
all because Christ Jesus comes in, right?.
Now, it's reorganization, right?.
Look, before I entered ministry,.
I used to work in a chemistry lab..
I was a sort of a biochemist..
I was working in a environmental lab for some time..
And I remember I was working for this lab.
and I was in the office.
and there were other co-workers beside me, right?.
And I still remember there's one co-worker.
I really dislike, all right?.
I really dislike, you know why?.
You know why?.
You know me, right?.
I work hard..
I'm a very good worker, all right?.
Right?.
And that co-worker that I dislike,.
he comes in late and he exit early..
He hardly do anything and he's lazy boy..
I was like, how is that possible.

$^{521}$that I'm a new worker.
and he has been working there for years?.
And then after some weeks,.
I learned about the relationship in the department.
and I realized he is actually directly related.
to the department head..
He's some sort of relatives to the department head..
I said, bingo, bingo, that's why..
That's why he hardly working, right?.
He's hardly working and he's getting paid, right?.
It's such an unfair situation..
Have you ever encountered a situation like that?.
Right?.
That's so unfair because they are blood with relatives,.
right?.
So of course he's protecting his cousin, right?.
And he's getting paid by just showing up for a little bit..
So I feel really oppressed, to say the least..
And then I remember one day,.
the CEO of the company,.
the head of heads, right?.
He fired the department head.
and he hired a new department head to replace him..
It's very interesting after that,.
the co-worker that I dislike,.
he start working hard..
It's very interesting, his life almost turned around..
It's because the department head now is new,.
is under new configuration..
There's a new head now, right?.
So he has to change his behavior,.
has to change how he works..
Now this is an analogy of Ephesians 5 and 6..
The household is working well according to Roman,.
Greco-Roman values..
However, when Christ Jesus comes in,.
everything will be reorganized..
Now he's the head..
Everyone is kept accountable under Jesus..
The values will be different, right?.

$^{561}$The order of the relationship has changed..
Just like this..
The left is the old order..
The male is up there above, right?.
Unchallenged or unchallengeable..
And then wife and children and servants,.
they are all subordinates..
And then there's a reorganization happening in Christ..
If you believe in Jesus,.
every single of your religious should change.
in a very fundamental way..
And that's the right side, the new order..
Christ is above all of them..
Male, female, children, servants,.
all are accountable to one sovereign, one Lord,.
one Lordship of Jesus Christ..
Someone is more powerful than anyone in the household..
That is the fundamental change.
that should happen in our lives..
If you believe in Jesus,.
when Jesus comes into your home,.
coming to our home, coming to our family,.
every relationship will be reordered in such a way.
that we, every one of us,.
will be accountable to the one true Lord, Jesus Christ..
So how does it work?.
How does it work?.
I remember, I don't know if you listen to radio..
Radio is an old thing, right?.
You radio program, do you listen to it regularly?.
Don't, or no, right?.
Sometimes I listen to radio programs..
I remember not too long ago,.
I was working in the church,.
it was very, very late, very late, almost 11 p.m..
And then I was driving home..
Takes about five to 10 minutes..
And then the radio was on, right?.
I didn't intentionally turn to one of the programs,.
but it was on..

$^{601}$And 11 p.m., almost midnight,.
what kind of program would be air in a radio station?.
You know, those songs, "Dedication Time," right?.
The caller will call in to the radio station.
and call the DJ and say something really touching, right?.
Expressing love and all that stuff, right?.
And then I didn't really pay attention.
until I was a couple blocks away from the church.
on the way home..
And I listened up and there was a middle-aged man,.
50-plus year old..
He call in, he call in to the radio station..
And he said to the DJ and saying,.
"Hey, you know, I have been married.
"with this woman of my love for 25 years,.
"but I have never,.
"I've never taken the opportunity to apologize,.
"to say a formal apology..
"So I want to dedicate this song to her.
"because I love her..
"For 25 years, I have been not really properly.
"related to her, and especially in our communication..
"And I want to dedicate this song to her.
"and I want to say sorry..
"And I want to change.".
You know, in those moments,.
you know, I almost crashed, you know,.
almost crashed because I was paying attention..
I was like, wow, you know, that takes a lot of courage..
Right?.
You know, as a male, I find myself,.
that's pretty hard to do, you know?.
Would you do that?.
You as male, as a man,.
you probably would not have the guts to do it, right?.
But I was listening to him and I listened to the song..
I forgot what song it is, but you know, anyways..
I said to myself, that is how Christ should change.
our family, our marriage, our relationships, right?.
Admitting our wrong, right?.

$^{641}$There's no sovereign, I cannot do something like apology.
because I'm ultimately accountable to Christ, right?.
That is the new order, right?.
That's how it works in our family, in our marriage..
That's how it works with our relationship.
with other brothers and sisters..
The old order would say, you know, I insist on my way..
The new order, when Christ comes in,.
there's forgiveness, there's understanding,.
there's proactive communication, right?.
This is how Christ should change our,.
your marriage, family, and relationship..
Well, let's work on that..
Let's really work on that..
It's a new year, new phase of our lives, work on that..
What does it mean for you when you say Jesus is Lord?.
How does it affect fundamentally and practically.
your relationship with other people, right?.
Let's all pray together..
Father, thank you for loving us..
That Jesus, your son and only son,.
would come and sacrifice himself on the cross.
to die for us, to save us and redeem us..
Thank you for your lordship..
It's not a lordship of power,.
but a lordship of love and grace..
And I pray that as we launch into a new year,.
your lordship will mean to us,.
fundamentally and practically,.
will renew and change our relationship..
We welcome you, welcome you to be our Lord,.
to be our savior..
Now pray that you will renew our relationships,.
whether it's a relationship with our spouses,.
with our brother and sisters, our in-laws,.
our family members, even our spiritual family members..
Now pray that your love will fill us.
and change us in such a way that we'll honor you..
And thank you for being our loving Lord..
And thank you for being,.

$^{681}$for the blessing to be your church..
And I pray that Jesus, you will be our Lord..
You will lead us and change us.
in such a way to honor your name..
In Christ's name, we'll pray..
Well, God bless you and happy new year..
And there is a bag of blessing.
we're going to distribute to all of you.
who haven't gotten one..
Make sure you get one, all right?.
I don't know what's in there,.
but it should be good things..
All right, and then this afternoon,.
we're going to have a fellowship meal together.
at Gene Fung's house church, all right?.
So you're all welcome to come..
And I don't know what kind of good food they prepare.
or we prepare, or we bring our own dish, right?.
I'm sorry, I'm sorry..
So it's a potluck, okay?.
It's a potluck..
So you bring something good.
and I'm sure it will be a great time.
of fellowship together, all right?.
I'll see you later..
\newpage



\section{以弗所書}
\label{sec:hnLQJ3HSDsE}
\textbf{同建健康的教會 1 (以弗所書 5) - 陳崇基牧師}
\newline
\newline
連結: \href{https://youtube.com/watch?v=hnLQJ3HSDsE}{\texttt{https://youtube.com/watch?v=hnLQJ3HSDsE}} ~~~~ 語音日期: 2023-01-26
\newline
\newline
\hyperref[sec:daNqcFTl4KQ]{\small{< < < PREV SERMON < < <}}
~
\hyperref[sec:index_chronic]{\small{[返順時目]}}
~
\hyperref[sec:index_scriptual]{\small{[返順卷目]}}
~
\hyperref[sec:5LAc82ueC_Y]{\small{> > > NEXT SERMON > > >}}
\newline
\newline
$^{1}$各位好 先跟大家拜個年.
我也有兩個小孩的 所以先謝謝大家.
新的一年.
不知道大家還在想什麼.
可能還有新年的氣氛.
新一年裡面有什麼計劃等等.
所以我也想在這裡有一個新的系列跟大家分享.
就是講到教會.
講到教會就好像講一個大家庭一樣.
如果新年的時候想一下有什麼四字詞語.
或者是揮春.
跟一家人有關的 有沒有.
我想了想也想不到.
有什麼揮春.
或者是四字詞語是跟一家人有關的.
合家平安 OK.
我想不到 中文開始差.
還有呢.
家肥屋潤.
嘩 這些真的是.
大家吃了很多東西.
一團和氣.
我也想到一團和氣.
真是聰明啊 你們.
所以我們想起這些話語的時候.
就想起教會其實應該是這樣的.
因為聖經裡面怎麼形容教會呢.
聖經裡面形容教會是一個家.
譬如已忽所書第二章第十九節.
他說什麼呢.
他說我們信了主的人.
在基督裡面的人.
就不算是那些所謂的外人.
不算是客類.
因為我們已經有一個所謂.
天國公民的一個身份.
屬於某一個國度.
屬於某一個國家.
這個國家當然就是上頭而來.
天上的國家 神的國.

$^{41}$就是跟聖徒去同國的.
是神家裡面的人.
Family of God.
God's Family.
我們屬於神家裡面的人.
所有本地的教會.
甚至整個普世的教會.
就是一個神家裡面的人.
Family of God.
所以在新約的時代.
保羅 彼得那段時間.
他們所有的教會.
當然跟我們認識有些不同.
本質上一樣.
認識有什麼不同呢.
就是他們是一些House Church.
就是某一個弟兄姊妹.
他們的家比較大一點.
譬如我們這裡.
圍內我們說.
誰家裡夠大.
我們就去那裡聚會.
主日他們就去那裡聚會.
這些叫做House Church.
所以保羅在這裡這樣說.
他說我們作為教會.
在基督裡面的人.
我們是一個Family.
其實是很有意思的.
因為這個Setting已經是一個Family.
可能在Living Room.
在大廳.
我們一起坐在一起.
就算是一個Family 一個家庭.
所以我們在這裡可以去想一想.
新的一年.
想一想一個比較健康的教會.
應該是怎樣的呢.
有沒有想過.
新的一年應該有一些.

$^{81}$新的目標是跟健康有關的.
大家是不是.
有誰是想減肥的.
不用舉手.
有 有人舉手.
很勇敢.
後面很多.
工作人員那裡.
當然坐在那裡.
要減肥.
有誰想過的.
有的吧.
每年都想過我們.
每年都想.
證明我們下一年又再想過.
去到今年已經初一.
我知道大家吃了很多東西.
應該都是失敗的.
不過不要緊.
我們繼續不要氣餒.
健康的人.
或者健康的生命.
總是有一些好的習慣要養成.
和一些壞習慣要除去.
教會其實都是一樣.
健康的教會.
有一些壞習慣.
我們需要去除去.
有些新的習慣.
好的習慣.
我們需要再pick up.
懶惰的.
我們疏懶了的東西.
我們要再pick up.
一個健康的基督徒.
也是一樣.
有一些懶惰的習慣.
忘記了.
忘記了的東西.
我們就要再pick up.

$^{121}$有一些好的習慣.
我們要再養成.
也是一樣.
我們作為一個family.
作為神的家.
我們都需要去想一想.
所以為什麼.
在我們今天開始.
我所說的這個series.
是和大家去想一想.
怎樣去應用.
我們作為一個family.
作為一個教會.
我們應該要做些什麼.
疫情總算過去了.
對不對.
總算過去了.
對於一些人來說.
疫情仍然持續.
不出街.
不敢去見人等等.
有時候未必是因為疫情.
可能是你需要的疫情.
你想有這個疫情.
那你可以避免很多其他事情.
真的.
疫情已經散去了.
那要想一想.
這個疫情對我們有什麼影響.
幾年的時間.
全球的疫情.
在去年2022年.
在英國做了一個研究.
說什麼呢.
就是說這個疫情.
問到底有多少個family.
有多少個人.
覺得這個疫情.
會令到他們家庭的關係.
人際關係.

$^{161}$有進步的.
有改善的.
因為疫情而有改善的.
你知道調查結果.
有多少人認為疫情.
是幫到大家的關係.
你猜猜有多少.
百分之多少.
10\%也好.
三分一.
換句話說.
三分一這麼多.
不是啊.
我不知道是不是泰國的錯誤.
總之他就說有三分之一.
只會有三分之一的人.
說疫情是幫到他們的關係.
那三分之二是什麼.
很可能是令到關係更加差.
和張力更加大.
換句話說.
這裡這麼多人.
我們將它分三楷.
大家做那個調查.
做完調查之後.
我們問的問題.
同一個問題就是.
疫情有沒有令到你的關係.
或者家庭的關係.
婚姻的關係.
是會有改善的.
三楷在這裡.
怎麼回答.
只有一楷說有改善.
這裡.
然後我們這兩楷就挖著它.
我們有沒有教錯.
是不是啊.
你還可以這麼幸福.
No way.

$^{201}$所以這是一個很sad的news.
想想你自己是不是這樣.
第二個問題.
這個UK的study.
就是說在UK.
在英國.
最大的那些家庭法律的公司.
他們的report.
就是他們的報告.
去年.
就是說離婚率.
是增加了多少.
猜猜增加了多少.
百分之幾多.
百分之九十五.
哇.
很開心.
他笑我開心.
離婚率增加了百分之九十五.
打官司都打到手軟.
換句話說.
我們用這裡做一個sample.
算了.
不要說這些.
新年快樂.
怎樣也好.
這個疫情對我們來說.
是很大的impact.
很大的影響.
無論你的工作也好.
你的個人的生命也好.
甚至家會也好.
所以這個時候.
我們想去想一想.
回到聖經裡面.
怎樣去看.
關係.
怎樣去看教會.
我就選了這段聖經來開始.
我還沒讀這段聖經的時候.

$^{241}$我想問你一個問題.
就是當我讀這段聖經.
從頭讀到尾的時候.
你有什麼感覺.
你的feeling是什麼.
你聽這段聖經.
可能是一段很熟悉的聖經.
又或者.
如果你當這段聖經.
是你第一次聽的.
第一次讀的.
你有什麼感覺呢.
我們開始吧.
已忽所書第五章.
第二十一節.
要傳敬畏基督的心.
彼此馴服.
作妻子的.
你們要馴服自己的丈夫.
如同馴服主.
因為丈夫是妻子的頭.
如同基督是教會的頭.
他又是這身體的救主.
教會怎樣馴服基督.
妻子也要怎樣馴服丈夫.
二十五節.
作丈夫的.
你們要愛自己的妻子.
如同基督愛教會.
為教會捨己.
以水藉著渡把教會洗淨.
使他成為聖傑.
好獻給自己.
作榮耀的教會.
毫無玷污咒文等類的缺陷.
以示聖傑沒有瑕疵的.
二十八節.
丈夫也要照樣愛妻子.
如同愛自己的身體.
愛妻子就是愛自己了.

$^{281}$從來沒有人肯護自己的身體.
總是保養愛惜.
正如基督待教會一樣.
因我們是他身體的肢體.
三十一節.
為這個緣故.
人要離開父母與妻子聯合.
依人成為一體.
這是極大的奧秘.
而我是指基督和教會說的.
然而你們每個人都要愛妻子.
如同愛自己一樣.
妻子也要敬重丈夫.
讀到這裡.
有什麼感覺.
我們做個調查.
有誰覺得感覺好.
感覺好.
男人.
我看不到女人感覺好.
家裡感覺好.
That's good.
很明白聖經的道理.
有誰不是感覺好.
不一定感覺不好.
可能不是感覺好.
讀完聖經.
有人嗎.
OK.
可能有.
但不敢舉手.
很可能.
你覺得不太合聽.
因為我們現代人.
我們說什麼.
平等.
Equality.
丈夫與妻子應該是平等.
應該是平起平坐.
大家應該是平等的地位和待遇.

$^{321}$這段聖經為什麼不合聽.
因為一味叫那個女人去馴服.
所以有些人都會這樣.
強行來解釋.
總之丈夫做什麼.
什麼都好.
妻子一定要馴服.
因為要馴服.
所以每件事都要聽話.
丈夫叫你洗廁所.
你就要洗廁所.
你不要去洗爐頭.
一定要叫你走左.
你不要走右.
我們現代人會覺得.
甚至有些人會覺得oppressive.
好像欺負女人.
不知道你有沒有這個感覺.
可能會有這個感覺.
這個是understandable.
好像很古板的一些鼓勵人.
盲目的去信從.
對不對.
我不是講了你心聲嗎.
講了你心聲.
看聖經應該怎麼看.
我經常都這樣說.
聖經不是直接寫給我們.
這段聖經不是直接寫給我們現代人.
他一定是針對當時的處境.
去講這番說話.
我不是說將聖經框在以前的社會裡面.
而是他藉著當時社會所了解的語言也好.
或者那些比喻也好.
那些說法也好.
用那種思想形態.
來溝通某一樣東西.
那樣東西.
上帝要啟示要溝通的.
就是藉著當時的處境來溝通出來.

$^{361}$然後我們捉到那個精義.
那個精義就是橫跨所有時空.
我應該這樣看聖經.
因為聖經不是直接寫給我們.
而是為我們而寫.
所以我們看完聖經也是一樣.
雖然這段聖經好像很熟.
為什麼好像很熟呢.
我做過很多婚禮.
我覺得那些新人都沒有想清楚.
他們覺得Sound good OK.
那OK的.
有些聖經我們婚禮的時候用.
當然要Sound good.
難道你婚禮的時候講那些舊約的擊殺經文嗎.
你當然找一段比較合聽一點.
適合那個處境.
所以很多時候婚禮都會用這段聖經.
但我想你今天看.
就好像第一世紀的婚禮.
讀這段聖經.
就不是二十一世紀.
如果第一世紀的耳朵.
去聽這段聖經的時候.
會有什麼的分別.
和你捉到什麼的.
那個是正義.
每一個社會都會對某一些家庭的關係.
社會裡面人際的關係.
有一些的明文也好.
非明文也好.
的規例.
有些可能寫出來.
有些可能大家明白的.
譬如例子.
給個例子.
我們華人的文化都是這樣的.
特別對男人.
或者對女人.
諸多要求.

$^{401}$華人的傳統文化.
女人多些要求還是對男人多些要求.
女人.
女人要什麼.
三從四德 沒錯.
男人就三妻四妾.
所以是兩回事.
華人文化都是這樣.
有一些家規.
或者有一些的code.
我們叫它做social code.
去處理和去形容.
或者甚至去.
限制了.
某一些的關係的動態.
回到第一世紀.
羅馬帝國的社會.
有什麼文明.
或者非文明的法規.
來處理一些家庭的關係.
那些我們可以統稱.
叫做household code.
就是家庭的法規.
國有國法.
家有家規.
家庭的法規.
我們叫它家庭的法規.
家庭的法規.
國有國法 家有家規.
他們的家規.
是怎樣呢.
所以我在這裡讀一段文字.
不是聖經.
是當時.
應該是耳熟能詳.
在羅馬逛街的時候.
在大城市逛街的時候.
你會聽到這首歌就會這樣唱.
或者有一種這樣的.
社會的價值.

$^{441}$或者所謂的household code.
我不知道你有沒有聽過一個人.
叫做亞里士多德.
不是竇麗士多德.
新年頭你們上了路.
亞里士多德.
亞里斯多德.
當時是一個哲學家.
有很多的teaching.
有很多的教導.
都是深入民心的.
亞里士多德說了一句這樣的說話.
聽清楚.
他說什麼呢.
他說說到household management的時候.
怎樣管治一個家庭.
這樣就分三部分.
他說第一部分是什麼呢.
就是主人怎樣對僕人的家規.
第二部分就是父親怎樣去管治子女的家規.
第三部分就是丈夫怎樣對管理.
管治妻子的家規.
第三部分就是丈夫怎樣對管理.
管治妻子的家規.
第三部分就是丈夫怎樣對管理.
管治妻子的家規.
第三部分就是丈夫怎樣對管理.
管治妻子的家規.
第三部分就是丈夫怎樣對管理.
管治妻子的家規.
三部分.
所以說到.
一連串就說到.
男人.
男人是一個家庭之主.
男人是家庭之主.
所以他就要管理什麼呢.
管理三件事.
當時羅馬的家庭.
特別大的家庭.

$^{481}$就會有三個方面.
關係的方面.
第一個就是主人和僕人.
第二個就是主人和父親.
和子女.
第三部分就是主人和妻子.
換句話說.
當時的所謂哲學.
當時的價值.
這些家規.
就是那個男人.
就是整個家庭裡面的.
sovereign.
主人.
最高的主人.
他有最高的.
authority.
最高的權威.
所以每個人都要聽他的話.
不單止這樣.
亞里斯多德寫這本書.
叫做Politics.
政治.
他講的政治不是我們今天很限制那種.
我們投票那些.
政治是講到.
整個社會裡面的.
order.
整個的規律.
整個羅馬帝國的規律.
就好像一個大的家庭.
哪一個是最頂的.
就是羅馬帝王.
然後羅馬帝王就處理.
他下面不同的關係.
他就是有特權.
有唯一的.
大權.
如果將這個廣義的.
帝國一個.

$^{521}$所謂帝國一個家庭.
縮小了變成一個家族.
一個家庭一個household的時候.
都是一樣.
你整個的household最高的那個是誰.
男士.
父親.
丈夫就是最大的.
甚至當時的.
羅馬的.
社會就說.
男人要好好處理自己的家庭.
處理這三個東西.
以至整個的.
國度整個的帝國.
才會有.
所謂國安.
才會有security.
那些叫做pax romana.
就是羅馬盛世.
羅馬盛世怎樣.
可以維持到.
就是所有這些人好好處理.
自己的household.
所以那個男人.
就是有這樣的特權.
去命令.
和去管治在他下面的人.
妻子 兒女.
僕人.
那麼.
家庭好的時候.
大家就好了.
大家就國安了.
舒舒服服.
可以生活下去.
保羅.
在二弗所說.
講到這些所謂household的relationship的時候.
現在聽清楚二弗所說.

$^{561}$就是有三個東西.
對嗎.
第五章第六章.
你回去看.
講這三個東西.
丈夫怎樣管理妻子.
父親怎樣管理兒女.
然後家裡面的主人.
怎樣管理他的僕人.
其實就是這三個東西.
但是保羅.
用一些當事人.
一聽是很熟的東西.
但是他就.
轉他一點.
等於一些人.
將一首.
當時很普遍.
一聽就明白的.
一聽就懂的一些歌曲.
就將它做一個remix.
然後聽起來.
好像很熟.
但是又很陌生.
保羅寫一封信.
給教會的時候.
你就發現原來有一些東西.
他做了.
是當時的人會覺得很突然.
和很詫異.
這裡給一個outline給你.
第五章第六章.
講到這些所謂household management的時候.
講這三個東西的時候.
保羅做了什麼.
我試一下去demonstrate.
當時的人.
我們回到第一世紀.
第一世紀婚禮的時候.
讀這段聖經會和我們今天.

$^{601}$再聽一段聖經有什麼不同.
第一世紀的人.
在他們耳中.
聽見這些教導的時候.
他們會有什麼反應.
首先五章二十二節.
他們說你們作妻子的.
wives.
listen.
當時大部分第一手的讀者.
去聽這段聖經的.
應該是男人.
男人.
聽到牧師在讀.
作妻子的.
你們凡事要順服等等.
男人會怎樣.
of course.
pritchard.
當然叫老婆聽話.
但到五章二十五節.
講了幾節.
講到妻子要做什麼.
of course 大家知道.
二十五節他說.
你們作丈夫的.
husbands.
listen up.
聽清楚我想講什麼.
你記得他講什麼.
二十五至三十三.
很長.
聽起來.
為什麼燒了我的答案.
就是要燒你的答案.
二十五至三十三節.
講什麼.
耶穌作為一個榜樣.
他怎樣去管理教會.
是為教會.

$^{641}$捨身.
愛教會愛到一個地步.
愛到底.
這個就是真正的主.
丈夫你怎樣去做.
如果你想想.
第一世紀的時候.
坐在婚禮上.
如果我是男人.
我真的不舒服.
平時我說了算.
無端端去到教會的時候.
那個牧師保羅.
講這樣的東西.
燒我的答案.
好了這個算了.
作丈夫的又要我一樣.
又要我那樣.
耶穌是主人.
然後六章一至三節.
好像有小孩在.
聽完.
兒女要做什麼.
要順從父母.
如果你是男人.
你怎樣聽這段世紀.
當然.
但是六章四節.
就講.
父親.
listen up.
要在主裡面去教養.
你的兒女.
不要激怒他們.
不要令他們煩惱.
要去到.
愛他們等等.
這裡又燒了他的答案.
六章五至八節.
作僕人的.

$^{681}$打工的.
你要聽聽話話.
要好好地做事.
忠心地做事.
男人聽起來的時候.
of course.
六章九節又燒他們的答案.
作主人的.
為什麼.
因為耶穌才是主.
你看到保羅.
用一些很熟悉的.
格式.
或者很熟悉的一些說法.
來講出在基督裡面.
如果耶穌來到家裡.
來到教會的時候.
他作主的時候.
會有很多不同.
right.
理所當然地.
當時的人聽起來.
講到household.
怎樣去處理關係的時候.
當然會講什麼.
作妻子的要做什麼.
作兒女的要做什麼.
怎知道他們在中間就隔了三格.
就講到.
當時的household裡面.
所謂絕對的主宰.
你要做什麼.
因為耶穌已經到了.
所有的.
關係都會再重組.
reorganize.
我記得以前.
在.
我還沒有傳道之前.
就在一間公司.

$^{721}$打工的.
一間實驗室.
在lab裡面工作.
我記得.
我很不喜歡一個同事.
坐在我旁邊.
為什麼呢.
因為你知道我這個人.
我這個人是膽小不偷吃.
我很勤力的.
那.
我也明白美國工作.
不可以太快.
有時候太快你做了別人的那份.
就不太好.
但我已經慢了下來.
我都很忠心地工作.
我就很不喜歡我旁邊的同事.
那個同事就是那些.
又遲來又早走.
然後.
天天就在那裡.
不工作.
我一開始進去的時候.
奇怪這個人為什麼可以做這麼多年.
然後我發現原來.
他的部門的頭銜.
是他的親戚.
原來是那些三姑六婆.
大家在那裡.
好來好去.
他就不用工作.
有一天.
上頭老闆.
CEO.
就辭退了那個部門的主管.
就辭退了.
第二天就換了一個新的來.
那天我回到公司的時候.
我就發現那個星期.

$^{761}$那個人我很討厭那個同事.
他突然間.
整個人變了.
你不相信有生命改變嗎?真的有.
整個人改變了.
為什麼?.
連坐姿都不同了.
平時很懶的.
這個叫做重組.
因為第二個主來了.
你就不要.
這麼架勢.
耶穌基督來到家庭.
來到教會.
其實都是一樣.
他將那個結構.
轉了一面.
因為轉了另外一個管理.
每一個人.
在整個的集會裡面.
包括那個主宰.
男的.
丈夫,父親,主人.
他的關係.
都完全不同.
因為所有的人都會被問責.
你不可以再懶.
所以你就問.
到底耶穌來到的時候.
對一個家庭.
就算所謂教會.
當時的教會就是family church.
對他們的關係應該有多大的改變.
新的主管到了.
就沒有人.
可以蹋蹋跳跳.
並且這個主宰.
是一位愛我們到底的主宰.
在當時的家規.
家庭的法規裡面.

$^{801}$對主宰是沒有什麼要求的.
對他的要求就是.
他要命令下面的人.
他老婆,他的子女,他的工人.
耶穌一到了.
整件事就會不同.
重組過的架構.
所有的人都在主的主權底下.
就有分別.
以往就是左邊那個.
男人都大了.
現在我們reorganize.
就變成怎樣.
所有人都在基督之下.
每個人都是keep accountable(保持管理).
這個就是一個新的家庭.
就是一個以主為至上.
基督為至上的教會.
這個就是一個健康的家庭.
這個就是一個健康.
或者合神心意的一個婚姻.
我不知道你作為一個人.
你的婚姻.
你的家庭.
會不會有這樣想.
需要時刻重組.
到底我在基督面前.
我是不是一個好的丈夫.
不是我太太面前.
不是太太面前那麼簡單.
我作為一個妻子.
是不是在基督面前.
是合他心意.
我做兒女的.
我是不是合上帝的心意.
等等我們需要再想一想.
怎樣在這個所謂.
department再reorganize的底下.
重組底下.
我應該怎樣去生活.

$^{841}$因為討好的最終是誰.
基督.
問責的是誰.
基督 耶穌基督.
我不知道你有沒有習慣去聽收音機.
有沒有.
我都很少的.
很多時候開車就會有收音機.
有時候收音機已經開了.
我記得有一晚.
我在教會做一件好事.
12點.
我就開車回去.
開車回去十多分鐘.
那個收音機就開了.
我很少會選擇什麼台.
剛剛那個台.
晚上有什麼節目.
清談那些.
點歌那些節目.
那些很肉麻的.
我都不是很懷疑.
聽著聽著的時候.
為什麼懷疑呢.
打電話上去.
他說.
這個男人五十多歲.
他打上電台跟那個DJ說.
我想點一首歌.
給我太太.
我們結婚25年.
我一直沒有跟她說過道歉.
但我今天想藉著這首歌.
跟她說道歉.
嘩 我開車的時候.
差點撞車.
很肉麻.
但很感動.
很感動.
是對的.

$^{881}$原來這麼久都沒有一句.
合心意的話說過.
他承認他自己錯.
一個男人可以五十多歲.
打個電話上去DJ那裡.
點一首歌.
然後向太太公開道歉.
我都未必做得到.
重組婚姻的關係.
我經常這樣說.
在婚姻裡面.
丈夫和妻子.
是沒有贏或輸的結果.
只會有贏或輸.
或者雙敗.
不會有贏輸的.
婚姻裡面.
不是我贏你輸.
或者你贏我輸.
從來不會.
要麼兩個一起贏.
要麼兩個一起輸.
這就是重組.
到底我在基督.
在耶穌裡面.
我是否可以忠心講道.
老規矩.
某君可能說了算.
可能以前是男人.
現在可能是男或女.
看情況.
我們正在去的方向.
就是上帝要重組在基督裡面.
重組我們所有的關係.
無論你是男還是女.
我們都是為了耶穌.
在基督裡面.
我們做一個中心的管家.
無論你是做丈夫也好.
父親也好.

$^{921}$媽媽也好.
妻子也好.
兒女也好.
整個關係都會改變.
以前的老規矩可能是.
我不想去溝通.
新規矩就是基督給我們的榜樣.
去溝通 學習去溝通.
學習去proactive去溝通.
這才是一個合成心意的家庭.
以至教會的家庭.
都會合成心意.
以前可能真的不想聊天.
我不知道怎樣去表達.
現在我就要學了.
因為這是新的order.
怎樣去proactive.
主動的去了解.
然後大家在愛裡面能夠相連.
重組.
原因就是.
我知道.
我是問責於.
這位愛我的上帝.
最近你知道嗎.
我女兒回家了.
跟我們住.
起碼會有一個學期.
因為她是教小學的.
所以她快要畢業了.
所以她就要回來.
做一個實習.
所以回到city.
回到Chicago.
她每天早上.
就吃兩隻蛋.
做早餐.
我也是吃兩隻蛋.
她吃兩隻蛋.
我吃兩隻蛋 不是成績吃蛋那些.

$^{961}$我們煎兩隻蛋.
她吃兩隻 我吃兩隻.
最近.
我聽到.
你知道蛋很貴嗎.
我聽到.
我老婆跟我女兒說.
她說女兒.
她很喜歡這些protein.
她每天吃兩隻蛋.
她說有沒有想過.
現在的蛋這麼貴.
你吃一隻好了.
她說這些.
是金句.
她不是跟我說的.
我偷聽聽到.
其實我覺得可能是跟我說的.
這叫做主動.
這叫做active listening.
我就.
對呀 這些蛋真的很貴.
她說完這句之後.
我就沒有早上吃蛋.
我吃奧妙.
我的心想 其實都ok的.
我在雪櫃看到蛋.
我都不再發育.
吃那麼多蛋 這麼貴.
真的很慘.
蛋這麼貴.
有時候吃蛋會便宜一點.
不是等到雞年才便宜一點.
吐龍舍馬.
羊後雞 七年了.
我的觀點是.
是active listening.
是需要大家去理解.
需要主動去理解對方的需要.
這個家庭 這是婚姻.

$^{1001}$才會合乘心意.
我們去做.
重組我們的家庭.
婚姻的關係.
以致我們的教會.
都會成為一個健康的教會.
我相信這句說話.
當我們教會 我們家庭.
無論是young family也好 其他family也好.
如果我們家庭是健康的話.
教會才會健康.
對嗎.
Let's work on that.
求主幫助我們.
我們唱最後一首的回應詩歌.
上帝一給我們智慧.
夫婦之間的關係.
我們和兒女之間的關係.
在我們主耶穌的主權底下.
我們再重組.
讓我們知道.
這不是上頭壓下來的命令.
因為這位主宰.
是愛我們的上帝.
我們唱這首詩歌結束.
然後我們才祈禱.
我們祈禱.
祈禱.
然後我們才祈禱.
《我找到你》 詞:陳零九 曲:陳零九.
從你的居所 我找到抱初之窩.
在你的國度裡 那句迷路與失去.
是你保護我 在你恩典中寄居.
是你幫助我 讓我可以跨過重陽深海.
從你的聲音 我找到世間歡欣.
在你的愛裡 永沒疑惑與唏噓.
是你安慰我 驟變的不必去追.
是你堅定我 讓我可以不怕洪濤掩蓋.
你是我的門石 是我千古的寶藏.
最堅固無比山寨 是我力量.

$^{1041}$我試過在這世逃生 贈我光 賜我熱戀.
唯獨你是信神到晚曬.
你是我的盤石 是我千古的寶藏.
每天長長釋放著 每刻似我的手.
在我心靈內 贈我光 賜我熱戀.
唯獨你是我神 我上帝.
我們同心得禮.
主我們再一次歡迎你來到我們中間.
在我們教會 在我們家庭 在我們婚姻.
在我們個人的生命裡面 你再一次作主.
你是我們唯獨的主宰 獨一的神.
主讓我們尊你偉大.
上帝在我們來到你面前 再重組這個關係的時候.
讓我們站在你面前 我們謙卑的.
做你的僕人 你所託付我們的.
每一個角色 無論是我們自己個人的生命.
在教會裡面的弟兄姊妹 團契裡面.
跟其他弟兄姊妹的相處.
還是自己的家庭 還有他們的婚姻.
跟兒女的關係 甚至可能有幾代的關係.
主我們都在你面前 知道你是我們唯一的主宰.
主讓我們聽見你的聲音 觸摸到你給我們的智慧.
讓我們每一方面 我們的管理是忠心的.
恩賢的來接受 接受你的管教 接受你的懷念.
主讓你自己的道 去提醒我們 去更新我們.
叫我們的生命 是一個健康的生命.
一個熟悉健康的生命.
叫我們的生命 是有著主你給我們的大愛.
你給我們的恩念.
我們做一個忠心的管教 討你的喜悅.
主保守我們眾多的關係裡 讓我們能夠見到上帝你自己的恩念.
見到上帝你的憐憫.
主一次把我們的心放在這裡.
我們求你充滿我們的心靈 成為我們真正的主宰.
我們願以基督的福祉 祈禱給我們.
\newpage



\section{彌迦書 4:1-5}
\label{sec:5LAc82ueC_Y}
\textbf{假使世界原來不像你預期(彌迦書4\_1-5)- 陳智祥傳道}
\newline
\newline
連結: \href{https://youtube.com/watch?v=5LAc82ueC-Y}{\texttt{https://youtube.com/watch?v=5LAc82ueC-Y}} ~~~~ 語音日期: 2023-01-30
\newline
\newline
\hyperref[sec:hnLQJ3HSDsE]{\small{< < < PREV SERMON < < <}}
~
\hyperref[sec:index_chronic]{\small{[返順時目]}}
~
\hyperref[sec:index_scriptual]{\small{[返順卷目]}}
~
\hyperref[sec:NqYX_zjUCDg]{\small{> > > NEXT SERMON > > >}}
\newline
\newline
彌迦書 4:1-5
\newline
\begin{longtable}{cl}
\hline
\hline
章節 & 經文 (和合本修訂版)\\
\hline
4:1 & \begin{tabularx}{0.7\textwidth}{X} 末後的日子，耶和華殿的山必堅立，超乎諸山，高舉過於萬嶺；萬民都要流歸這山。 \end{tabularx} \\ \\ \relax
4:2 & \begin{tabularx}{0.7\textwidth}{X} 必有許多民族前往，說：「來吧，我們登耶和華的山，到雅各神的殿。他必將他的道指教我們，我們也要行他的路。」因為教誨必出於錫安，耶和華的言語必出於耶路撒冷。 \end{tabularx} \\ \\ \relax
4:3 & \begin{tabularx}{0.7\textwidth}{X} 他必在許多民族中施行審判，為遠方強盛的國斷定是非。他們要將刀打成犁頭，把槍打成鐮刀。這國不舉刀攻擊那國，他們也不再學習戰事。 \end{tabularx} \\ \\ \relax
4:4 & \begin{tabularx}{0.7\textwidth}{X} 人人都要坐在自己的葡萄樹和無花果樹下，無人使他們驚嚇；這是萬軍之耶和華親口說的。 \end{tabularx} \\ \\ \relax
4:5 & \begin{tabularx}{0.7\textwidth}{X} 萬民都奉自己神明的名行事，我們卻要奉耶和華－我們神的名而行，直到永永遠遠。 \end{tabularx} \\ \\
[1ex]
\hline
\hline
\end{longtable}
$^{1}$請交給Mark Sir.
麻煩Carmen.
今天是初幾呀?.
初八.
初八,OK,OK.
還可以派年,是吧.
跟你身邊的人說新年快樂.
趁著還可以派年的時候.
那些未派利是的快點派.
OK.
未收利是的快點收.
很高興今天在中間能夠見到大家.
因為今天都是冰天雪地,是吧.
今天地面很滑.
還有開車不容易.
特別是那些很遠來的朋友.
再一次歡迎大家.
今天我們都是一位很遠來的傳道朋友.
跟我們做客席講完.
就是Joseph.
他多遠呢?.
就是在D'Ifill.
你在D'Ifill住嗎?.
他在Trinity附近.
現在Joseph就快完結他在三藩神童院的M.Diff的課程.
所以我們趁著有機會可以特別叫他來.
特別是新年的時候我們叫他來.
冰天雪地的時候.
他的講題我們不多講了.
《假使世界不像你預期》.
歌名來的.
經文就是尼加書第四章.
有聖經的話就給我打開.
尼加書第四章一至到第五節.
我們PowerPoint不會打出來的.
你拿回你的手機拿回你的聖經.
尼加書是舊約的一本先知書.
四章一至到第五節.
末後的日子.
耶和華殿的山必建立.

$^{41}$昭扶珠山.
高舉過於萬靈.
萬民都要留歸這山.
必有許多國的民前往說.
內伯我們登耶和華的山.
賓雅國神的殿.
主必將他的道教訓我們.
我們也要行他的道.
因為悔憤必出於石安.
耶和華的言語必出於耶路撒冷.
他必在多國的民中施行審判.
為遠方強盛的國斷定是非.
他們要將刀打成泥頭.
把槍打成鐮刀.
這國不舉刀攻擊那國.
他們也不再學習戰事.
人人都要坐在自己的葡萄樹下.
和無花果樹下無人驚嚇.
這是萬君子耶和華口所說的.
萬民各奉己神的名而行.
我們卻永永遠遠奉耶和華.
我們神的名而行.
我們將時間交給祖師爺.
我們歡迎各位媽媽.
很感恩今天能夠和大家.
在星期日早上在一起敬拜.
就像Mark Sir說的.
我是在D'Irviel開車過來的.
今天也有些.
幸好今天好像如期一樣.
不塞車.
就很平安地來到.
我們今天說的經文是.
尼加書的四章一到五節.
講題是.
假使世界原來不像你預期.
話說我聖誕節之後.
就和朋友去了洛杉磯旅行.
我們通常幻想的加州.
就是陽光,海灘,溫暖.

$^{81}$誰知我去到那兩個星期之後.
就發覺一下機就說.
天這麼陰暗.
完全沒有太陽.
去到整個星期.
風大到吹到頭快要掉了.
去海灘都去不到.
什麼都做不到.
其中同行我們有一個朋友.
他們就說.
為什麼加州會是這樣.
這個都不是我想像中的加州.
好像芝加哥這個時間還好一點.
他在旅行中一直說這件事的時候.
其實都影響到我的心情.
當加州原來不是像我們預期.
就會令到東西不太好吃.
東西不太好玩.
又想起塞車.
塞到我快要刺線.
我就說整個旅程都很低落.
最後令我覺得.
這個真的不是我想住的地方.
我想想如果一次這樣的旅行.
都可以這樣影響到我們心情的時候.
何況是我們每一天身處的環境.
假使世界原來不像我們預期的時候.
我們會怎樣選擇呢.
我們會選擇逃避.
說我不再住在這個地方.
還是我們繼續以一個很負面的心態來迎接呢.
在尼加書四章一到五節當中說的是.
當世界不像上帝預期的時候.
上帝怎樣來回應.
這段經文告訴我們.
上帝沒有離開我們.
上帝並沒有放棄我們.
當世界不如上帝預期的時候.
上帝要做的卻是要來恢復祂的子民.
恢復我們對祂的敬拜.

$^{121}$我們首先來看看尼加書是一本什麼書.
尼加書未必是大家最熟悉的舊約的一本書.
尼加書主要來說的是上帝的審判訊息.
整本書可以分成簡單來說.
三個大的部分.
第一章 第一到第二章.
第三到第五章.
和第六到第七章.
這三個部分都是以一個審判.
再有一個盼望的訊息.
是一個巡迴.
來宣告上帝的心意.
而我們今天看的第四章.
就是第二個部分.
我們今天看到的是一個盼望的訊息.
盼望之前是有審判的.
審判的內容是什麼呢.
或者說上帝究竟是要說.
以色列人做錯了什麼事呢.
我們可以繼續看聖經.
可以看一看尼加書的第三章.
第三章的第一節是這樣說的.
「我說:雅各的首領,以色列家的官長,你們要聽.
你們不當知道公平嗎?」.
即是說他們那裡沒有公平.
去到第九節我們跳一跳.
上帝繼續審判.
「雅各家的首領,以色列家的官長.
你們要聽我的說話.
你們厭惡公平.
在一切事上屈往正直.
以人血建立錫安.
以罪孽建造耶路撒冷.
首領為賄賂行審判.
祭司為顧家司悔憤.
先知為銀錢行占卜.
他們卻依賴耶和華說.
耶和華不是在我們中間嗎?.
災禍必不臨到我們」.
先知尼加在這一段經文.

$^{161}$主要針對的是以色列人在宗教上的敗壞.
而來說的.
他們這一班人.
這一班首領.
這一班祭司.
這一班先知.
杯稀耶和華.
敬拜偶像.
撇棄上帝的公義.
他們這一班宗教的領袖.
在做的是收人錢財.
幫人擋災的工作.
如果放回現在這個處境.
我們可以想一想.
如果我們走進去一間教會.
牧師就跟你說.
我可以幫你祈禱.
不過有一個先決的條件.
他就放了一個奉獻箱在這裡.
你知道牧師祈禱比較靈.
所以他就這樣提供給你.
他就說.
如果你給一百元.
你奉獻一百元.
我就幫你祈禱.
祈禱你身體健康.
一千元就幫你祈禱.
升職加薪.
一萬元就幫你祈禱.
六合彩.
你就有一個很好的回報率.
現在還是新年.
新春大減價.
買十次送兩次.
你就每個月都有一次牧師幫你祈禱.
這些教會變成什麼教會.
這些變成不是認識上的教會.
這些變成買賣的教會.
就好像當日的以色列一樣.
所以上帝說聖殿耶路撒冷.

$^{201}$由敬拜上帝的地方.
變成上帝厭惡的地方.
是尼加書第3章第12節.
上帝甚至說.
所以因你們的緣故.
錫安必被耕種成一塊田.
耶路撒冷必變為亂堆.
這殿的山必長成叢林的高處.
尼加預言.
耶路撒冷會變成一個亂堆.
將會被毀滅.
因為耶路撒冷不能再完成它的使命.
耶路撒冷代表的不再是上帝的榮耀.
而是代表虛妄的地方.
當上帝的殿不再成為敬拜上帝的地方的時候.
上帝說要將它毀滅.
對於當時的以色列人來說.
耶路撒冷殿被毀滅.
其實是一個很實在存在的威脅.
對於他們來說.
這不是一個假設或是很久之後才會發生的事.
當時以色列人面對的是亞述帝國的一些逼迫.
在公元前主前722年的時候.
亞述人已經攻佔了以色列的北角.
佔領了他們的首都撒瑪利亞.
他們下一個要打的就是耶路撒冷.
所以耶路撒冷他們是充滿一個威脅的.
我們好像看到一個充滿了紛爭.
好像毫無希望的情況的時候.
尼加卻宣告說.
上帝將要重建.
將要恢復它榮耀的這個信息.
上帝告訴我們.
當耶路撒冷不是好像他預期的時候一樣.
上帝說他不會袖手旁觀.
在第四章的第一節.
尼加說到末後的日子.
即是將來的時間.
耶和華殿的山必被建立.
超乎珠山.

$^{241}$高舉過於萬靈.
我們看到整段經文的語氣.
一個情況變得很不同.
第三章說的是一個好像末日般的信息.
要毀滅 完全沒有了.
也是第四章.
上帝說他要重新建立屬於他的結局.
重新建立屬於他的世界.
要按照他的心意重新建立他的殿.
耶和華殿的山必建立超乎珠山.
高舉過於萬靈.
在這裡說的話.
耶和華的殿.
即是耶路撒冷的那座山.
是要高過世界上一切的山.
超過一切的高度.
如果想像一下.
我們現在一抬頭一看.
就看到的不是只有高樓大廈.
而是一見就看到耶路撒冷那座山.
上帝的聖殿的時候.
那份榮耀 那份威榮.
他在說上帝的殿是要成為世界的中心.
勝過一切的勢力或者神明.
但上帝說他要建立他的殿.
不單要彰顯的是他的榮耀.
更加重要的是要重建那一班以色列人的心.
第四章是第一節的下半部分到第二節這樣說.
他說萬民都要留歸這座山.
必有許多國的民前往.
留歸.
我們有時會想像就像一條河流這樣湧下來.
但現在這裡卻是一些很反自物理的現象.
不是湧下來而是要湧上去.
所有的人 萬民.
即是所有的人就好像一條河流這樣湧上耶路撒冷山.
好像幾波巨浪一般要來覆蓋上帝的殿.
這個情景經常讓我想起.
那時候在香港 去到11月尾就會看到美國Black Friday的新聞.
令我印象很深刻的就是.

$^{281}$每逢Black Friday 前夕人們就會等在商店門口.
到一開門的時候 人們就會湧上去 衝上去.
他們就會去買他們想要的東西.
我們想像一下 如果人們是用這樣的心情.
去到上帝的殿的時候會是怎樣呢?.
尼加遜說的訊息就是說.
這班人等的不是要去上帝的殿當中買祝福.
也不是去上帝的殿當中買紀念品.
他們要去做的 要去尋求上帝的話語.
第四章第二節.
這一班盲民所有的人說.
來拜我們登耶和華的山 奔雅各上帝的殿.
因為當我們去到上帝面前的時候.
上帝會將祂的道教訓我們.
以至於我們可以走祂的路 可以跟隨祂.
因為悔憤必出於錫安 耶和華的言語必出於耶路撒冷.
唯有我們來到上帝面前的時候.
我們才可以認識祂的話語.
這個景象就好像昔日以色列人去朝聖那樣.
他們好像唱著上行之詩.
他們帶著一個很激動的心情來敬拜上帝一樣.
在詩篇第122篇.
這樣來形容他們的心情.
他們說.
人對我說 我們往耶和華的殿去 我就歡喜.
上帝要重見的是甚麼人.
上帝要重見的是一班真誠地來學悟敬拜.
學悟上帝話語的人.
這裡是上帝在應許我們.
他要來施行他的拯救.
我們很少將拯救和敬拜放在一起.
但其實對上帝正確的敬拜其實是一個救贖.
這裡的救贖是說.
上帝怎樣重新更新這一班人.
以至於這一班以色列人可以和上帝有一個正確的關係.
對於當時的以色列人來說.
耶路撒冷是一個充滿賄賂,虛假.
他們沒有辦法正式正確地認識上帝的地方.
以色列人是一班敬拜偶像的人.
但上帝在這裡說.

$^{321}$他要扭轉這個情況.
他要重見耶路撒冷成為敬拜的中心.
他要向人教導他的心意.
以至於人都可以遵行他的旨意.
上帝做的是要扭轉這一班敬拜偶像的人.
甚至領導萬國萬民.
都自願地來到上帝的面前.
去學悟上帝的話語.
這就是第一點.
上帝所要施行的救贖.
就是要讓我們可以真誠地去敬拜他.
和他有和好的關係.
在這段經文第三到第四節當中.
我們看到上帝要重新建立耶路撒冷城.
重新建立的這一班以色列人.
他甚至要在這一班人當中實現他的公義和和平.
在第三節這樣說.
他必在多國的民眾施行審判.
為遠方強盛的國斷定是非.
上帝要在人的面前施行審判.
列國都要來到上帝面前.
要上帝為他們斷定是非.
就是要上帝作為審判的人.
來審判什麼是對什麼是錯.
成為他們的仲裁.
他們國與國之間不需要再因為爭奪什麼資源.
或者什麼紛爭而來發動戰爭.
其實如果我們想一下這個世界.
很多時候很多紛爭的源頭.
有時候不單止是為了確保自己的利益.
很多時候有些紛爭是為了證明我們是對的.
就好像俄羅斯和烏克蘭的戰爭.
站在俄羅斯那一邊.
他們就覺得烏克蘭加入了NATO.
對於我們的領土作為一個威脅.
但是對於烏克蘭來說.
這是一個自保的行為.
當然他們當中有利益的紛爭.
但他們更加想證明他們的觀念是對的.
或者我們想起很多時候有些真橙的案件.

$^{361}$我們經常都喜歡看這些唐山風暴的事件.
這些唐山風暴的事件很多時候.
其實不是因為他分多一點.
他和我分得一樣多.
我們而爭取.
很多時候是因為爭取一口氣.
明明我才是照顧他的人.
為什麼我的弟弟分得和我一樣多.
為什麼我不分多一點.
為什麼他不分少一點.
我們經常都用自己的情來量度.
究竟什麼是對 什麼是錯.
但上帝在這裡說.
祂要恢復祂的主權.
要讓祂的公義.
可以成為一個絕對的標準.
祂要恢復祂創世時候的秩序.
因為上帝作王.
國與國之間不用再爭奪什麼.
所有的人 所有國家.
都要俯伏在上帝的主權面前.
他要按照他的公義的準則來管理.
這個就是上帝要恢復他的公義的方法.
而當上帝恢復了他的公義之後.
我們見到一個很好的和平的實驗.
在第三節的下半節.
他這樣說.
他們要將刀打成泥頭.
將槍打成鐮刀.
因為萬國都承認了上帝公義的準則.
以至於他們不需要再持有他們的武力.
不需要再持有他們軍事的力量.
上帝要帶來的是一個普世的和平.
因此他們可以將他們打仗用的刀.
和打仗用的槍.
變成一些耕種.
農耕可以用的泥頭.
和收割的時候可以用的鐮刀.
因為上帝彰顯了他的公義.
以至於這國不再舉刀攻擊那國.

$^{401}$他們也不再學習戰事.
因為上帝主權的實現.
上帝公義的彰顯.
他們不需要以武力再威脅別人.
保護自己.
上帝透過他自己的能力.
制止了紛爭.
終止了一切的戰爭.
帶來一個和平.
我們看到上帝在這裡恢復的是.
他的救贖.
他的公義.
和他的和平.
這三點是舊約學者Peri Yoder.
他的書當中指出.
是Shalom的三個部分.
我們應該大部分人都很.
對於希伯來文的字.
是很熟悉的Shalom.
我們很多時候將它翻譯成.
平安或者和平.
Peace to you.
但其實Peri Yoder在這裡指出.
他就說Shalom其實不單止是包括和平.
平安的意思.
更加是一個大的包圍.
可以說是一個完全完整.
滿足的生命.
是一個Fullness of life.
當世界不如上帝預期的時候.
上帝並沒有放棄這個世界.
而上帝更加要做的是.
要將他的Shalom.
他這份完整的生命賜給我們.
在第四節當中我們看到.
尼加在描繪這個平安.
這個Shalom很重要的一個部分.
第四節說到.
人人都要坐在自己的葡萄樹下.
和無花果樹下.

$^{441}$無人驚嚇.
這是萬君之言和親口說的.
上帝在這裡應許的是.
每一個人都可以享受他們土地的出產.
享受這份富足和上帝的供應.
過著沒有驚慌的生活.
但這個時候我們會問.
既然上帝已經應許了平安.
但為什麼這裡會強調無人驚嚇呢?.
什麼是他們成為驚嚇.
或者他們fear的源頭呢?.
在這個時候我想請大家一同再一次翻開聖經.
我們再看看聶王記下第十八章.
在聶王記下第十八章.
這裡說到亞述攻打耶路撒冷時.
有什麼回應呢?.
我們可以看看聶王記下第十八章第29到第31節.
第29節是亞述使者和以色列王說的話.
亞述王如此說:你們不要被希西加欺騙了.
因為希西加王不可以救你們脫離我的手.
也不要聽希西加王使你們倚靠耶和華.
耶和華必要拯救我們.
這城必不交在亞述王的手中.
你們不要聽希西加的話.
因為亞述王這樣說:你們要與我和好.
出來投降我.
你們投降了.
你們不再順依耶和華的時候.
你們各人就可以吃自己葡萄樹和無花果樹的果子.
喝自己這裡的水.
亞述王這裡做的是將這份驚嚇給了他們.
他說如果你們願意來投降的話.
你們就可以過著一些富足的生命.
你們就可以吃你們自己的無花果樹和葡萄樹的果子.
亞述王很喜歡用驚嚇作為戰爭的手段.
他們的政策就是一路打一路打.
打到人都害怕.
以至於他們不可以再去反抗.
這就是亞述帝國眼中的和平和平安.
其實如果我們回想.

$^{481}$現在很多國家都同樣用這些手段.
北韓用勞改營來嚇他們的國民.
伊朗就用他們的武力來鎮壓當時的示威.
這個世界運作的法則就是.
誰惡誰正確.
但上帝卻說他要實現的和平不是這樣的.
上帝所要彰顯的主權.
所要彰顯的能力是要除去一切的驚慌.
是神學家Cornelius Planting.
在一本書說.
Not the way it's supposed to be.
這本書說的是罪對於我們這個世界的影響.
他在這本書指出.
聖經當中的平安Shalom.
其實是指一個完全的狀態.
當人所有的需要都被滿足.
當神和所有受造之物.
在一個充滿公義滿足喜樂的關係當中.
對於他來說.
Shalom 平安是甚麼.
平安就是The way things ought to be.
就是上帝創造時候所定下的心意.
而不是因著罪而被扭曲的一個世界.
上帝說他要帶來一個真正的和平.
真正的平安.
他不是以驚惶或者武力作為基礎的平安.
而是要將人從驚慌當中拯救出來.
將恐懼當中拯救出來.
上帝要做的是要我們從罪當中得到拯救.
如果我們繼續來到看尼加書第五章的時候.
我們會見到這個應許的延續.
上帝應許說有一位王要由伯利恆而出.
在尼加書第五章第四節.
他說這位王要倚靠耶和華的大能.
並耶和華他上帝之名的威嚴.
牧養他的羊群.
他們要安然居住.
第五節說到.
這位王要作我們的平安.
這是尼塞雅的寓言.

$^{521}$這個尼塞雅的寓言告訴我們.
耶穌基督的降世將要為我們的生命帶來真正的平安.
我們很多時候對於平安和平的一個想法.
就是我們沒有戰爭,沒有威脅,沒有恐懼.
這就是平安和平.
我們過得很安逸.
我們有我們需要的東西.
我代代平安就是和平.
我們的和平可能是我的日子如何.
我的平安又如何.
但很多時候這份平安是否真的呢?.
回想在幾個月前.
當上海在封城的時候.
我外婆住在家裡.
住得很安逸.
她覺得很平安.
因為她不用離開家裡.
依然有足夠的食物.
有足夠的空間來走動.
但同一時間.
其實有很多人.
他們生活有很多負擔.
大家都知道封城的時候.
食物賣得很貴.
他們甚至不能外出.
其實都不知道在吃什麼.
他們只可以捱肚餓.
因為他們選擇了食物.
就可能沒錢交租.
對於我外婆來說.
那段時間她有她生活所需的一切.
食物或者住屋.
對她來說.
其實那段時間她過得很平安.
但對於某些人來說.
那段時間卻是一個災難.
這個災難是什麼?.
這個災難是.
這不是它應該是這樣的.
上帝所設計的世界.

$^{561}$上帝所預期的世界.
並不是這樣的.
但上帝卻要透過耶穌基督.
成就祂真正的平安.
耶穌基督帶來的平安.
是拯救.
是讓我們可以跟上帝一個復和的關係.
以至於我們可以成為祂的子民.
可以以一個正確的敬拜來朝見祂.
耶穌基督所成就帶來的平安是公義.
這份公義不是一個假象.
不是好像平靜.
不是好像什麼事情都發生這樣.
或者都沒有發生這樣.
不是為了滿足某些人的需要.
忽略了某些人的一份公義.
祂給我們一個公義的盼望是說.
當祂再來的時候.
祂會以祂的公義審判人一切的罪惡.
耶穌基督帶來的平安是和平.
是和平.
祂要將我們一切的驚惶都消除.
但祂帶來的和平不是透過妥協.
不是透過.
而是透過跟那惡者的爭鬥.
重新建立屬於上帝的國度.
恢復創世時的那份秩序和美好.
這就是上帝將要成就的一份平安.
當世界不如上帝所預期的時候.
上帝沒有停止你的工作.
當世界不如我們預期的時候.
我們又會怎樣做呢?.
在這段經文的第五最後一節當中.
我們看到是等候上帝應去實現的時候.
有兩種不同的人.
第一種人是萬民卻奉己神的名而行.
屬於世界的人沒有歸向上帝.
而是割除己義.
沒有遵行上帝的意思.
但這段經文給我們提醒.

$^{601}$有第二種人我們可以成為.
我們卻永遠奉耶和華.
我們上帝的名而行.
我想對於當時的以色列人來說.
他們是有一個選擇的.
他們可以選擇信亞述王所說的話.
還是選擇上帝的話語.
其實我想選擇相信亞述王是很簡單的.
只需要拋棄那位好像我見不到的上帝.
我不知道他是哪裡的上帝.
我只要信我眼睛所見到的東西就行了.
相信上帝對於他們來說是一個更加難的選擇.
因為他們需要有一個勇氣.
來拒絕亞述王的邀請.
我們今天同樣有這個選擇.
我們可以選擇我們眼睛所見到的現實.
還是我們可以選擇相信上帝的應許.
今天的現實是我們見到這個世界有很多挑戰.
有很多不公義的事情發生.
這個世界有很多方面不是我們預期的一樣.
我們見到大大小小的戰爭.
烏克蘭和俄羅斯在打仗.
南蘇丹有種族衝突.
緬甸,敘利亞,耶門,衣索比亞有內戰.
或者說近一點.
是美國2023年開頭.
1月到現在二十多天的時間.
已經有40宗的mass shootings.
有40人,有80人.
接近80人死亡,180人受傷.
面對這些不公不義.
面對這些罪惡的時候.
可能我們都會問.
上帝你的公義,上帝你的平安在哪裡呢?.
我們雖然眼見這些現實.
但現實同樣都是.
上帝的平安已經透過耶穌基督的降世實現了.
祂已經將這份平安賜給我們.
這個現實是無論我們今天見不見到.
上帝依然是我們的世界當中掌權.

$^{641}$在過去這幾年的時間我們面對病毒.
或者今天我們面對戰亂.
面對我們現在身處美國社會當中的矛盾,分裂.
或者是香港,中國一些政治環境的壓迫.
或者是人與人之間的一些衝突的時候.
我們的心會不會被這些恐懼所佔據了.
以至於我們見到的只有不平安呢?.
當我們見到世界並不如我們所預期,所期望的時候.
我們同樣有沒有這個勇氣.
來願意選擇相信上帝.
依然來見證上帝賜給我們的平安是真的呢?.
或者我們有沒有這份勇氣.
來宣告上帝在這個混亂當中祂依然作王.
耶穌基督的平安,耶穌基督的拯救.
要將這個世界恢復成為「The Things Ought To Be」.
就是上帝創造時所定下的心意.
假使世界原來不像你預期.
我們面對這個世界很多的不友善.
很多人拒絕上帝.
但是上帝依然提醒我們要按照祂的路而行.
要活在祂的心意當中.
在這些的境況當中.
上帝鼓勵我們更加需要有勇氣.
來永遠的事奉耶和華.
我們神的名義.
而不是被這些驚惶所佔據.
而是活在祂的平安當中.
我們一同有一個禱告.
天父上帝我們知道.
真正的平安來自於你.
唯有你的拯救,唯有你的公義.
唯有你彰顯的能力.
才可以帶來真正的平安.
就讓我們的眼目能夠專注於.
基督你所承受的平安當中.
在這個時代當中仍然能夠堅守對你的信靠.
專心等候你應許的實現.
求主你來繼續兼顧我們的信心.
在這個等候的時間仍然不忘記.
神你今天依然作王.

$^{681}$你依然掌管一切.
求主你繼續隨便我們禱告.
禱告大福.
靠著耶穌基督得勝的名給.
阿們.
多謝Joseph給我們的提醒.
剛才我們說到拜完.
代代平安.
希望我們不只是口頭禪.
我們知道平安是從耶穌基督而來.
剛才講完他提醒我們.
真正的平安在神.
真正的神.
多謝.
願上帝給大家平安.
下個星期能夠和大家一起來.
事奉一切敬拜.
多謝.
\newpage



\section{}
\label{sec:NqYX_zjUCDg}
\textbf{Building A Healthier Church 3: United for Mission - Rev. Mark Chan}
\newline
\newline
連結: \href{https://youtube.com/watch?v=NqYX-zjUCDg}{\texttt{https://youtube.com/watch?v=NqYX-zjUCDg}} ~~~~ 語音日期: 2023-02-06
\newline
\newline
\hyperref[sec:5LAc82ueC_Y]{\small{< < < PREV SERMON < < <}}
~
\hyperref[sec:index_chronic]{\small{[返順時目]}}
~
\hyperref[sec:index_scriptual]{\small{[返順卷目]}}
~
\hyperref[sec:Cv_wNlk__oQ]{\small{> > > NEXT SERMON > > >}}
\newline
\newline
$^{1}$Okay, Mark..
- Hey, good morning, everyone..
Why don't we say hello and shake hands to one another,.
call each other brother and sisters,.
remember what I said last week..
So find some people and shake their hands.
and welcome them to our worship community..
- Hello, Mark..
(audience chattering).
- All right, good to see you guys, good to see you..
Good to see you..
All right, so it's good to be here.
and today at the end of service,.
we have Communion Sunday..
It's the first Sunday of the month.
and we're looking forward to this time.
to take the communion elements,.
to remember what Christ has done for us..
And again, Pastor Mark and the announcer was called,.
Kelly, did she introduce herself?.
- I don't know..
- Okay, oh, you were there..
Okay, Kelly, thank you very much..
And there's also one more announcement..
There's a sister show show next week, right?.
So it's after our Sunday worship service,.
you guys are gonna eat together, hot pot, right?.
Hot pot together and make candles together..
Yeah, candles and then hot pot and then Super Bowl.
and then show show, whatever, okay?.
So it's gonna be a long day next week..
Anyways, but today let's turn our attention.
to the word of God..
We are on a series called Becoming a Healthier Church..
Healthier Church, I think that's what God desire..
God wants His church and local churches to be healthy, right?.
He wants you to be healthy too, right?.
So it's kind of related topic..
You think about how to be healthier..
Maybe it's your New Year's resolution..

$^{41}$Nobody cares about it now, right?.
It's already February..
So maybe it's one of your goals to be healthier..
What kind of things that you would do or change.
in order to be healthier?.
Now, again, a related topic is the church..
As a church, as a faith community, as a family of God,.
how can we be more healthy?.
And we spent two weeks to take a look.
at some of the passages from the letter of Ephesians.
in the New Testament..
Now we are turning our attention to another book,.
another letter written by Paul as well..
It's the letter of Philippians..
So if you have your Bible, you can turn to Philippians 1,.
how to be a healthier church..
Chapter one, above all, you must live as citizens.
of heaven, conducting yourself in a manner worthy.
of the good news, the gospel of Christ..
Then whether I come and see you again,.
or only hear about you, I will know that you are standing.
together with one spirit and one purpose,.
fighting together, contending together for the faith,.
which is the good news, the gospel..
So it's a short verse from chapter one.
of the letter of Philippians, written to the Philippians,.
written to the Christians in Philippi..
So today, I just learned some of the new here.
who are students, especially students.
at Moody Bible Institute..
So I prepared some pop quiz..
Okay, I prepared some pop quiz..
All right, so well, I think you will do well.
because you're Moody students..
Anyways, let's go over these questions together..
All right, it's a multiple choice questions..
As we begin to study the book of Philippians,.
I would like to begin with some of the questions.
to warm up your mind, all right?.
So are you ready?.

$^{81}$- Yes..
- All right, question one, all right?.
Question one, all right..
You ready?.
It's all about the epistle of the Philippines, all right?.
The church in Philippi was planned by, who planned it?.
All right, A, Apostle Peter, B, Apostle Paul..
A or B?.
Who said A?.
Who said B?.
Why you hesitate?.
(laughing).
Why you hesitate?.
All right, so it is..
It is planned by Apostle Paul, right?.
Why you hesitate?.
It's such a simple question, all right?.
So more difficult question to come, all right?.
All right, now this is a tough one, okay?.
This is a tough one, all right?.
On which of his missionary journey.
did Paul plan the church in Philippi?.
All right, A, first, B, second, C, third journey..
All right, think about it, okay?.
Work with your neighbors, work with your neighbors, okay?.
Ask your neighbors, do you know the answer or not, okay?.
On which of his missionary journey.
did Paul plan the church in Philippi?.
The first one, second, or third, okay?.
How many of you say A, the first one?.
Okay, I see a hand in the room, okay..
By the way, thank you very much..
They work so hard to make our worship service possible..
Sometimes there are technical difficulties, all right?.
You have any complaint, we invite you to join the team.
to help us out. (laughing).
Anyways, option one is A, is first missionary journey..
How many of you say B, second?.
All right, some said B..
What, you following others, or?.

$^{121}$Okay, are you sure James is right?.
All right, how many of you say C, the third journey?.
All right, so let's see the answer..
It is B, it is the second journey, okay?.
Second missionary journey, all right?.
All right, not bad, guys, not bad..
All right, those of you who are hesitating.
to raise your hand to answer the question,.
you have to answer this question, okay?.
Question three, does the church in Philippi.
support Paul's mission?.
Yes or no, okay?.
Yes or no, how many say yes?.
Okay, how many say no?.
That's why he wrote the letter, right?.
No, yeah, he did..
All right, all right..
So Paul did, Paul did receive offering.
and support from the church of Philippi, all right?.
So they're a pretty good church, very generous, okay?.
So last question, okay?.
Uh, there you go, okay..
This is a tough one, okay?.
Which of the following people were the leaders.
of the church in Philippi?.
Which of the following people were the leaders.
of the church in Philippi?.
A is, um,.
Euodia and Sinti, okay?.
Euodia and Sinti..
A, B is Alexander and Philip..
C is Timothy and Luke..
All right, how many of you say A?.
All right, all right, I see some hands, okay..
How many of you say B?.
B, names very familiar, Alexander and Philip..
By the way, they're father and son, yeah..
No, C, Timothy and Luke?.
Timothy and Luke, okay, good..
Brave hand, brave people..

$^{161}$The answer is,.
is A, all right?.
If you read chapter four, you find these two names show up..
They are actually having some problem in the church..
In any case, I want to warm you up with these questions.
to think about what the letter has to say about the church,.
about us as people in Christ..
So the church of Philippi was planted by the apostle Paul.
in his second missionary journey..
He went the first one, he went to what we call.
Asia Minor or modern day Turkey in that area.
from Antioch..
So he preached the gospel to different towns and cities.
and find some people who embrace the gospel.
and planted all these house churches in different towns..
And then they return to their home base.
in Jerusalem and Antioch..
And then as they pray about it,.
they felt that the Spirit is moving them.
to go on for an other journey..
So they started the second journey..
And the second missionary journey,.
they pass through what we call Turkey,.
you see the blue line,.
they were in Jerusalem and then they were back to Antioch..
And then from Antioch, they move west.
through Asia Minor..
And there are different things happen.
and they wanted to go to that town,.
but the Holy Spirit is not leading them to..
And then they moved on to another town.
and then move around..
As they move around, they move westward.
to a place or a port or harbor called Chorus..
So when they were there, they spend the night.
and God show Paul a vision..
And in that vision, a man, a Macedonian man,.
Macedonian man by race, Macedonia..
And in the dream, in the vision,.
this Macedonian man call on Paul and his team.

$^{201}$to come over, to cross the sea,.
the Aegean Sea to the sort of the peninsula of Greece..
So to call Paul to cross the sea,.
to come over and help him, all right?.
Now that's interesting,.
because if you know something about history.
or the context of the Bible,.
you know, the Macedonian, all right?.
Remember the name, Alexander the Great?.
Alexander the Great, ring any bell?.
Right?.
Fairly familiar name, a great military leader.
living in 300 BC or so..
And he was a Macedonian..
And Alexander the Great, he crossed the sea eastward.
and conquered what do they call Turkey.
and the entire Asia, the entire continent.
and brings with him the Hellenistic culture,.
the Greek culture..
So it's the Macedonian helping the rest of the world, right?.
Conquered the entire piece of land, entire continent..
And Alexander the Great conquered the world.
in order to help other people in the east..
Here, the vision is the other way around..
Paul saw this vision, the Macedonian man.
calling him to help him, bring the gospel over to help him..
So it's such a cultural reverse, right?.
So Paul follow God's guidance.
and he did cross the sea.
and eventually he came to a town called Philippi, all right?.
Actually in the book of Acts, recorded this incident..
So we, here we have Luke, the writer of the book of Acts.
and the apostle Paul and all the other co-workers,.
they bought a ship from Chorus.
and they sailed straight to another island.
and then the next day they landed.
to the other side of the land in Macedonia..
And from there, they traveled to different towns.
in which the big city called Philippi, all right?.
So they came to Philippi, which is the district.

$^{241}$or the province of Macedonia..
It's a colony of Rome, okay?.
It's a colonial city..
So we stayed there for several days..
Now, as you can tell, this is a major city in Macedonia..
It's a major Roman colonial place.
and it's an important city to the Roman empire..
Now, what is a colonial city?.
What is a colony?.
What does it do?.
All right, what does it do?.
What does a colonial city do?.
What is the purpose?.
What is the function?.
For the empire, for the imperial power,.
they established different places, foreign places..
They conquered the city and named that city.
and changed that city into a colonial port.
or a colonial town..
And a colonial place has a really, really important.
political social function..
For the capital of the imperial power,.
they established colonies around the world.
or outside of the border..
It's to promote their own power, their own culture.
and their own values..
So you live in Philippi as a Roman colonial city..
Your job or your colonial dream.
is to become a Roman citizen..
And by becoming a Roman citizen,.
your function in the empire, in a colonial city,.
is to promote the power and the value and the culture.
of the imperial power..
That's what colonialism is, right?.
We are talking about our world today..
It's a post-colonial world, right?.
Because it used to be the Brits,.
they conquered all the places around the globe..
And what they wanted to do is to promote their own culture.
and project their power in other places..

$^{281}$So you can imagine if you're a local person in Philippi,.
your lifetime dream is to become a Roman citizen.
and to promote the interests, the power,.
the propaganda of the imperial power, which is Rome..
And then if you understand this background,.
you should understand some of the special words.
that Paul used in the letter of Philippians..
So this is a special word alert..
I'll give you some context first..
Some of the pictures are taken.
belonging to the first century of Philippi, all right?.
Upper left is Roman forum,.
and lower right is Roman theater,.
which seats about 50,000 people,.
part of the imperial power projecting their own culture..
And lower left is Roman toilets, all right?.
So the Roman lifestyle,.
they bring with their own imperial lifestyle.
to different places in different city, all right?.
So if you never try one of those Roman toilets,.
you should, yeah, because this is what we do, all right?.
This is what we do today, thanks to the Romans, right?.
So these are the buildings and architecture.
of imperial power..
Whenever they go, they will establish their own symbols.
and their own lifestyles into the city,.
especially the city in the Roman East, all right?.
So they impose foreign power.
onto different towns that they conquer..
Now, here's a special word alert that I'll talk about..
We read this verses, 127..
Paul is telling the church in Philippi,.
"You must live as citizens of heaven,.
"conducting yourself in a manner worthy of the gospel.".
What's so special about this word.
is that behind this word is the Greek term..
Can you pronounce that word?.
Yeah, politiomai, right?.
Politiomai, right?.
Something like that..

$^{321}$Now, you don't have to be,.
you don't have to be fluent in English.
in order to understand what this word is all about..
What we translate in Chinese or other language.
or in English, live as citizen,.
conducting that verb, behind that verb, is politiomai..
And politiomai, you can tell it's related.
to some of the words that we use today..
Politics, right?.
Citizenship..
So Paul is using a Roman term, Roman colonial term,.
to tell us, to tell the church in Philippi.
what it means to be a Christian,.
what it means to be a follower of Christ..
Paul used this colonial word not for Roman citizenship,.
but for citizenship in heaven..
Heaven is citizenship..
He kind of used that common word.
that people were familiar with..
And turn it around, subverting it or remixing it,.
and put it in the context of the gospel..
When you came to Christ,.
your ultimate identity is not the earthly citizenship..
I know it could be a very relevant topic to all of us,.
right?.
Many of us, we probably,.
some of our life goals.
is to obtain the United States citizenship, right?.
United States citizenship..
Or any citizenship of any strong country, right?.
Big power, superpower..
We want to have that political right..
We want to have that political identity..
Paul is reversing it..
Paul is saying what you are seeking.
and what you are pursuing should be.
the citizenship in heaven..
So we are Christ's colonial city..
We are the heavenly colonial city..
Christ plant us in different towns,.

$^{361}$in different communities,.
in order to promote the gospel,.
promote the value,.
the interest of the one and only eternal kingdom of God..
It's not the Roman way that we're concerned about..
It is the Christ way that we should be concerned about..
Paul is using this term.
and reminding the people in Philippi,.
all of us in Christ,.
you are being called to the best citizenship of the world..
Better than the Romans,.
better than the United States,.
better than anything..
It is the kingdom of God..
That is our identity..
And that's what church is all about..
You must be political mind..
You must be the citizens of heaven..
Remind yourself what are you after in your life?.
What is your greatest interest?.
It is to promote the kingdom of Christ,.
promote the value of the kingdom..
Again and again, Paul's used this term.
in chapter three as well..
Same word, but you are,.
you're not Roman citizens..
You are citizens, political mind, of heaven..
Where the Lord Jesus lives..
And we are waiting for him to return as our savior..
And he will take our weak mortal bodies.
and change them into the glorious body like his own..
Using the same power with which.
he will bring everything under his control..
He's talking about the coming king, Jesus Christ..
The coming king of our kingdom..
And he is coming and we are waiting for his coming,.
the glorious coming of Christ..
And that's why we are going forward to heaven.
and we are the citizens of heaven..
To promote the interest and the power.

$^{401}$in the kingdom of Jesus Christ..
So if you ask the question,.
what does it mean to be a healthy church?.
What does it mean to be a healthy Christian?.
What does it mean?.
You think in the colonial way, the heavenly colonial way..
A healthy Christian, a healthy church,.
is most concerned about promoting the gospel..
And that's why Paul wrote the letter of Philippians..
Let me give you some suggestions as we close..
Some of the marks or some of the indicators.
of a healthy church..
All right, it's the outreach to other people..
Church is the only organization on earth.
that is concerned most about the non-members..
You see what I'm saying?.
Church is the only body on earth,.
or heavenly body, heavenly sent colonial body,.
that is most concerned about the non-members..
How many of you are Costco members?.
Yeah, are you proud of it, being a Costco member?.
Yeah, right?.
(congregation laughing).
You're proud of it, right?.
Yeah, because you need that membership card.
to go in to that big box, right?.
Without that card, you cannot go in..
And with that card, you can bring a lot of people, right?.
You can bring the whole church to Costco.
and we can all buy some cheap eggs..
You know, all the organizations on earth.
are concerned about their own members, right?.
But church is the other way around..
It's the other way around..
How do we become a part of the heavenly kingdom?.
It's that we should pay attention.
to those who do not belong to us..
Those people who have not heard of Christ,.
those people who need the gospel,.
who need to hear the good news..

$^{441}$Our ultimate concern and interest should be those people..
That's why church community, churches,.
should always be open doors, right?.
I know there were a lot of concerns these days.
about church safety, shooting, and all of that..
But let us be reminded that the church is on earth.
being planted from the heavenly kingdom above.
is to promote the gospel, the good news..
Some of the actions item I would like to suggest to you,.
and let's apply what we have learned today in your lives..
Number one, pray, pray in your own prayers..
Pray that you will be able to reach out to the non-members,.
people around you..
I know it's all good..
We have Christian brothers and sisters,.
we can hang out together, right?.
Last week, I spent hours in Moody Bible Institute.
to have fellowship with some of the brothers and sisters.
in the midst of us..
And it was such a great time we had together, right?.
We can pray together, we all talk in the same language,.
right, we all talk, we're in the same wavelength, right?.
And we enjoy it..
There are times like this, but our ultimate concern.
is to be speaking and talking and relating and outreaching.
to those who are different from us..
We are followers of Christ, but we want to reach out.
to those who have not heard of the gospel, right?.
Yesterday I gave rides to some of the students in the city..
There were some students that I had never met before,.
but I gave them a ride, it was a pretty long ride,.
two hours ride..
I drove them to some place for their exam,.
and then drove them back..
I really enjoy it..
To talk to them, to tell them what I am, right?.
So in our conversation, I met them and I introduced myself,.
and then they didn't know I'm a pastor..
So after some hours of being together in the car,.
and then we went to lunch, and then I finally said,.

$^{481}$"I'm a pastor.".
And then that girl was like, "Wow.".
Her eyes was brightened..
I don't know if it's for good or for bad,.
but, "Oh, you're a pastor.".
And then I got a nickname..
I got a nickname, she calls me Pastor Daddy..
I guess I'm so old now, right?.
The students are calling me Daddy..
I don't know if it's a good thing or not,.
but, you know, so I was able to just represent Christ.
as I relate to them, right?.
As I provide some of the services or help them out,.
this is what we should be concerning about..
So number one action item is that for you to pray.
that God will use you to outreach other people..
Second, second action item..
So funny, right?.
I didn't tell Rachel, but, you know, God bless you..
You got some sisters..
And number two, action item number two,.
is to invite others to join the fellowship..
You know, we got social activities.
and whatever we want to outreach to other people..
It's a great platform to invite other people to join us..
Number three, when other new people are coming.
to our fellowship or our dinners or our social activities,.
you know, we need to open our hearts..
We need to open our lives to invite them,.
to make them feel welcome, right?.
That's what the family of God should be..
That's what the church should be like..
So number one, pray that you will outreach to others..
Number two, invite others to join our community..
Number three, welcome them, right?.
Let's all pray together as we partake.
the communion, bread and wine together..
Let's all pray..
Dear Father, we give you thanks for calling us in Christ,.
that we have the privilege to hear the gospel.

$^{521}$and to believe in your name..
And we want to thank you for your love..
We want to thank you for your grace..
Because of the blood of Jesus Christ,.
we can become heavenly citizens..
And God, we pray that you will help us.
to live out our heavenly citizenship..
Let our concern, let our interest,.
is to promote your gospel..
And help us to reach out to others.
and use us to invite others to join your heavenly community..
And I pray that your love will fill us up..
That we'll be your channel, we'll be your instruments..
As we take the bread and the cup,.
remind us once again, once again,.
that you are our Lord, Savior..
You are our loving Lord and God..
Thank you for redeeming our lives,.
forgiving our sins, so that we.
are called brothers and sisters in Christ..
Thank you for your love in Christ's name..
The world is not a place for you to live..
\newpage



\section{約翰福音 2:1-11}
\label{sec:Cv_wNlk__oQ}
\textbf{你可以活得更豐富 (約翰福音2\_1-11)- 程子信牧師}
\newline
\newline
連結: \href{https://youtube.com/watch?v=Cv_wNlk-_oQ}{\texttt{https://youtube.com/watch?v=Cv\_wNlk-\_oQ}} ~~~~ 語音日期: 2023-02-07
\newline
\newline
\hyperref[sec:NqYX_zjUCDg]{\small{< < < PREV SERMON < < <}}
~
\hyperref[sec:index_chronic]{\small{[返順時目]}}
~
\hyperref[sec:index_scriptual]{\small{[返順卷目]}}
~
\hyperref[sec:qMFmuMv64OQ]{\small{> > > NEXT SERMON > > >}}
\newline
\newline
約翰福音 2:1-11
\newline
\begin{longtable}{cl}
\hline
\hline
章節 & 經文 (和合本修訂版)\\
\hline
2:1 & \begin{tabularx}{0.7\textwidth}{X} 第三日，在加利利的迦拿有一個婚宴，耶穌的母親在那裡。 \end{tabularx} \\ \\ \relax
2:2 & \begin{tabularx}{0.7\textwidth}{X} 耶穌和他的門徒也被請去赴宴。 \end{tabularx} \\ \\ \relax
2:3 & \begin{tabularx}{0.7\textwidth}{X} 酒用完了，耶穌的母親對他說：「他們沒有酒了。」 \end{tabularx} \\ \\ \relax
2:4 & \begin{tabularx}{0.7\textwidth}{X} 耶穌說：「母親，我與你何干呢？我的時候還沒有到。」 \end{tabularx} \\ \\ \relax
2:5 & \begin{tabularx}{0.7\textwidth}{X} 他母親對用人說：「他告訴你們甚麼，你們就做吧。」 \end{tabularx} \\ \\ \relax
2:6 & \begin{tabularx}{0.7\textwidth}{X} 照猶太人潔淨禮的規矩，有六口石缸擺在那裡，每口可以盛兩三桶水。 \end{tabularx} \\ \\ \relax
2:7 & \begin{tabularx}{0.7\textwidth}{X} 耶穌對用人說：「把缸倒滿水。」他們就倒滿了，直到缸口。 \end{tabularx} \\ \\ \relax
2:8 & \begin{tabularx}{0.7\textwidth}{X} 耶穌又說：「現在舀出來，送給宴會總管。」他們就送了去。 \end{tabularx} \\ \\ \relax
2:9 & \begin{tabularx}{0.7\textwidth}{X} 宴會總管嘗了那水變的酒，並不知道是哪裡來的，只有舀水的用人知道。於是宴會總管叫新郎來， \end{tabularx} \\ \\ \relax
2:10 & \begin{tabularx}{0.7\textwidth}{X} 對他說：「人家都是先擺上好酒，等客人喝夠了才擺上次的，你倒把好酒留到現在！」 \end{tabularx} \\ \\ \relax
2:11 & \begin{tabularx}{0.7\textwidth}{X} 這是耶穌所行的第一個神蹟，是在加利利的迦拿行的，顯出了他的榮耀來，他的門徒就信他了。 \end{tabularx} \\ \\
[1ex]
\hline
\hline
\end{longtable}
$^{1}$各位誠惠的弟兄姊妹和朋友,早安.
也很開心可以今天和大家分享神自己的話語.
多謝Mark Sir的邀請.
我是宣道會石安堂的Pastor Jackson.
今天和大家分享一個訊息.
活得更豐富.
新的一年,希望大家在主的裡面.
可以有更豐富的恩典.
不知道當說到豐富的時候.
你會聯想到些什麼.
我們中國人很喜歡吃東西.
我自己也很喜歡吃東西.
所以一說到豐富的時候.
我就想起一餐豐富的飯.
可能對有些人來說是一餐豐富的自助餐.
我記得我曾經和太太一起去澳門.
探訪在澳門住的舅父.
他帶我們去吃自助餐.
你知道澳門有很多酒店.
有很多自助餐.
看到的時候,很多東西吃.
我們不是經常吃,自助餐也很少吃.
所以看到的時候,很豐富.
舅父說這些很小事,很小事.
他說有另外一間我想帶你去.
但是因為你只有一天的時間.
他說免得浪費時間在吃自助餐上.
Mark Sir十一月放假也來了多倫多.
我們一起吃了一餐豐富的午餐.
加拿大有一間銀行.
那間銀行的標語真的很貼切.
它的標語就是.
You are richer than you think.
你比你想像中更富有.
挺好的,你比你想像中更富有.
人人都喜歡豐富.
人人的定義也很不同.
例如剛才說到在北美.
那種豐富,很多時候都是說選擇上.
我們怎樣可以選擇一個更好的.

$^{41}$就叫豐富.
但是對比起第三世界國家.
他們的豐富可能只是有一餐飯吃.
所以豐富的定義,人人都不同.
但是既然今天我們來到相聚一起的時候.
一起敬拜的時候.
我們就問聖經的豐富.
有什麼意思在裡面呢?.
在新的一年,我想藉著耶穌以水變酒.
一個這麼熟悉的神蹟.
和大家看看什麼是豐富.
第三天,在加里尼的迦那有一個婚宴.
耶穌的母親在那裡.
耶穌和他的門徒也被請去赴宴.
在這個故事裡面.
我們看到從頭到尾都有一個婚宴.
根據當時猶太人的傳統.
或者一些舊約聖經的記載.
我們都看到猶太人的婚宴.
不是只慶祝一天.
而是一個星期,大飲大食.
其實也不是很難明白.
雖然我們今天不會說有七天的慶祝.
但其實我們很多人都分開幾天.
在不同的地方慶祝.
好像我認識的一些弟兄姊妹.
他們在這裡先擺酒.
然後他們是在中國來的.
然後去中國又慶祝幾天.
然後後來去到香港.
又要和一些在香港的親戚朋友.
又要吃一餐.
所以你說猶太人七天的慶祝.
是不是很過分呢?.
也不是,我們都已經去了三四天.
不過是在不同的地方,不同的時間.
約翰福音裡面.
迦勒的婚禮是一個很特別的記載.
因為一方面.
其他的福音書.

$^{81}$馬太福音,馬可福音,盧迦福音.
我們找不到.
很特別,只有約翰福音才有.
而另外這個神蹟又是很特別的.
就是令到我們很多人都不明白.
耶穌的第一個神蹟.
約翰說這是耶穌第一個神蹟.
耶穌第一個神蹟.
不是醫病講鬼.
不是叫死人復活.
可能你和我來到去.
是耶穌的話.
我們第一個神蹟當然是想.
君天動地.
想多些人認識你.
但是這個神蹟又不是這樣.
而且這個神蹟.
再想真一點.
其實都很奇怪.
只不過是可以繼續開派對.
沒有酒了.
耶穌變一些酒出來.
大家嘰呀哈哈,很開心.
這樣又可以繼續.
所以究竟我們應該怎樣.
來理解這個故事呢?.
約翰福音的作者約翰.
其實是一位滿有豐富的情感的文學家.
為什麼我這樣說呢?.
因為他寫的故事.
有很多的比喻.
有很多的隱喻,對比.
等等這些手法.
我舉幾個例子.
第一張,約翰福音大家都很熟悉.
我們記得就是一個光明和黑暗的對比.
你記不記得?.
光照在黑暗裡.
黑暗卻不接受光.
一個光明和黑暗的主題.

$^{121}$其實這個主題在第一章之後.
稍後那幾章都有出現.
好像尼哥底姆的故事.
大家記不記得尼哥底姆是何時找耶穌的?.
是晚上.
尼哥底姆是晚上找耶穌.
而最後他是不認識耶穌的.
聖經沒有給他那個結論.
他不認識耶穌.
但到井旁.
第四章井旁的撒瑪利亞婦人.
大家記不記得撒瑪利亞婦人是何時遇見耶穌的?.
是正午,太陽最猛烈的時候.
最光的時候遇見耶穌.
所以我們很快就看得到.
她有光就認識神.
她在晚上找耶穌.
她是不認識耶穌.
然後第九章這個光明和黑暗的主題都在當中.
大家記不記得第九章天山的蠻人.
他在整個故事裡面.
當然他是唯一一個黑眼的.
但故事很諷刺的地方.
他是黑眼的.
但他認識著.
相反那些有眼睛的人.
那班法利賽人.
他們是黑眼的.
他們看不到.
不認識著.
所以我們見到光明和黑暗在楊方的裡面.
這裡那裡用不同的故事去表達出來.
另外《約翰福音》的第六章.
耶穌做完五餅二魚的神蹟.
然後跟門徒說.
我的肉是可以吃的.
我的血是可以喝的.
嚇死人.
即是嚇到人們.
所以人們就走了.

$^{161}$我們讀《約翰福音》的故事的時候.
我們不可以只是看表面.
看表面的時候我們就走寶了.
我們應該來看一下它背後的意思.
所以在聖經裡面.
當講到有婚宴的時候.
其實婚宴,筵席.
就是代表了神的國的降臨.
有婚宴,有筵席.
就是因為神的國已經來到了.
這個就是聖經的筵席的意思.
而那班信徒,神的子民.
因為神的國的降臨.
他們就要起來.
就要慶祝.
就要吃飯,喝酒,慶祝.
好像《賽亞書》25章第6節,第9節.
萬君的耶和華必為萬民擺設筵席.
有肥金與美酒.
到那日人必說.
看吶!這是我們的神.
我們向來等候他.
他必拯救我們.
這是耶和華.
我們向來等候他.
我們必因他的救恩.
歡喜快樂.
神來到了.
神的國到了.
所以我們要來擺設筵席.
然後在另外一些聖經裡面.
我們也見到有這個例子.
大家很熟悉的.
耶穌另外一個比喻.
《馬太福音》22章.
天國告庇一個王.
為他的兒子擺設娶親的筵席.
然後提示祿.
我們要歡喜快樂.
將榮耀歸給他.

$^{201}$因為羔羊的婚期到了.
他的新娘也預備好了.
天使對我說.
你要寫下來.
凡被請赴羔羊婚宴的人.
有福了.
所以作者約翰.
用了婚宴來做七個神蹟的開始.
不是隨便.
隨便.
想到甚麼就寫下去.
不是.
是刻意的.
是有目的的.
就是想告訴我們.
九曰裡面.
一直在應許神的國.
來了.
真正的新郎.
到了.
一直想的時候.
我都在想.
有點像甚麼呢?.
我想大家都看過奧運會.
每一場的比賽.
都是很精彩的.
不知道你們喜歡看甚麼比賽.
有人喜歡看游泳.
有人喜歡看跑步.
我也很喜歡看跑步.
我弟弟喜歡看健身.
我不知道大家喜歡看甚麼.
但無論甚麼比賽.
哪一場比賽.
總是有一樣東西.
在奧運會裡面.
是最吸引人的.
最多人看的.
大家都是拭目以待的.
那是甚麼呢?.

$^{241}$那就是開幕禮.
開幕禮將所有人的注意.
每個人都期待著.
每個人都想看.
當然很多開幕禮都很精彩.
但我不知道大家還記不記得.
當年這個開幕禮.
真是厲害.
2008年8月8日.
北京舉辦的奧運會.
一流在鳥巢舉行.
很久以前去過北京旅行.
我們都要特地來參觀鳥巢.
那天的奧運會.
中國的民防四寶功夫.
真是很漂亮,很精彩.
當時其他國家的總理.
總統,傳媒等等.
每個人都讚那場奧運會的開幕禮.
只是開幕而已.
但很重要,很精彩.
所以在《楊康》裡面.
七個神蹟裡面的第一個.
用婚宴來表達.
就是一個開幕禮.
告訴我們,告訴神的子民.
我們要歡喜快樂.
因為神的國來了.
神的時代,新的時代開始了.
彌賽亞已經來臨了.
經文接著一路來講.
這個神蹟.
在五節的第八節.
「他母親對傭人說.
他告訴你們什麼.
你們就做吧.
之後他人潔淨禮的規矩.
有六口石缸擺在那裡.
每口可以乘兩三桶水.
耶穌對傭人說.

$^{281}$把缸倒滿水.
他們就倒滿了.
直至缸口.
耶穌又說.
願再逃出來.
送給燕迴總管.
他們就送了去」.
當他們來到發現沒有酒之後.
故事裡面就講到有六口的石缸.
我找到一些最接近當時石缸的相片.
相信都是當時有可能用到的石缸.
我知道Mark Sir五月帶隊去以色列.
看看Mark Sir有沒有機會.
給大家看到這些當時的石缸.
當時猶太人的家裡.
都是有石缸去做潔淨禮儀.
主要是潔淨一些食具和洗手.
今時今日你去猶太人的餐館.
當然你不會見到石缸.
但你會見到一個洗手盤.
就是專門用來給猶太人做潔淨禮.
經文繼續仔細講那個神蹟.
不是只是說耶穌祈禱.
然後他們就有很多酒.
不是這樣的.
經文很仔細.
耶穌吩咐一班婚宴的工人.
將那六口的石缸裝滿了水.
照猶太人潔淨禮的規矩.
有六口石缸擺在那裡.
每口可以盛兩三桶水.
裝滿了之後.
然後又說倒.
裝到很滿.
甚至去到缸口.
然後經文繼續講.
當工人出去給賓客的時候.
拿些水,酒給賓客的時候.
燕梅的總管就跟新郎說.
人家都是先擺上好酒.

$^{321}$但你竟然留到現在.
竟然現在才拿些好酒出來.
我們今天講這個豐富.
聖經裡面的豐富.
我說如果作者約翰要表達神的豐富.
他會怎樣來表達呢?.
約翰用了婚宴來形容神的豐富.
就是神的國來到了.
所以神的子民要開心,要慶祝.
神來到我們的生命裡面.
帶給我們真正生命的豐富.
然後約翰藉著神蹟.
來講到神自己的豐富是很多的.
是滿舍的.
是比以前好.
神蹟裡面講到那六個石缸.
可以裝很多水.
水是滿的.
甚至是滿舍的.
去到缸口.
就是其實表達豐富.
負責燕梅的人說.
你把好酒留到現在.
新的酒取代了舊有的酒.
而新的酒是更加好的.
負責燕梅的人說.
你把好酒留到現在.
所以約翰講的.
這些就是神的豐富.
一路看的時候.
你會發現豐富不是只是多.
Quantity.
我們會說很多很多酒.
就是豐富了.
飲不完的酒就是豐富了.
可能對我們來說.
我想起以利亞仙姿.
大家記不記得.
去幫一個窮寡婦.
有一個桶油.

$^{361}$那桶油很少.
但可以不停用.
不停用不停都有的.
但這個用完.
用乾淨可能就沒有了.
所以約翰的故事不是這樣.
豐富不是只是多.
在約翰的眼中.
豐富是根深.
是全新.
是改變.
Quality.
不是Quantity.
是Quality.
因為大家留意.
耶穌不是以不好的酒.
甚至過期的酒.
變成好的酒.
我再說一次.
耶穌不是以不好的酒.
壞的酒.
變成好的酒.
耶穌是以水變酒.
是根深了.
是改變.
是全新的.
這個就是神給我們的豐富.
新的一年.
我們有很多很多的希望和願望.
有人希望減肥.
有人希望儲錢.
有人希望生孩子.
不知道你的希望.
是不是時時圍繞著Quality.
多一點.
我想要多一點多一點.
多一點錢.
多一點樓.
多一點衣服.
多一點新衣服.

$^{401}$多一點娛樂.
但作為基督徒.
我們不應該只是要多一點多一點.
我們要新.
我們要全新.
我們要根深.
這個新是指到我們生命內的根深.
而這個新比起世上追求數量上的多.
是更加好的.
因為這個全新.
這個根深.
是一個生命的素質.
是沒有其他人可以奪去的.
是聖靈的果子.
彼得候書第一章說.
有了信心又要加上得恆.
有了得恆又要加上知識.
有了知識又要加上什麼什麼等等.
不知道大家有沒有機會看過一本書.
這本書非常好.
其實是一本很舊的書.
這個新的版本.
書名叫做跟隨他的腳踏中行.
英文很出名.
In His Steps.
很精彩的一本小說.
如果你沒有看過.
希望你可以去看一下今年.
其實故事也很熟悉.
故事裡面說到.
應該說故事裡面說了一句說話.
大家是很熟悉的.
故事裡面說到牧師邀請一班人.
因為發生了一些事.
後來牧師就問教會.
有沒有人肯參與一個運動.
牧師就邀請大家.
你們有多少人肯陪我在這一年裡面.
認真問一條問題.
就是如果是耶穌.

$^{441}$祂會怎樣做.
你每做一件事都問這條問題.
如果是耶穌.
祂會怎樣做.
What would Jesus do.
這句說話就是從這本書.
In His Steps裡面出來的.
於是就一個人舉手.
兩個人舉手.
一個人舉手.
三個人舉手.
就有一班人就參加了這個運動.
當他們來到去參加了這個運動.
做每一件事都問耶穌會怎樣做的時候.
到最後他們發現.
他們自己改變了.
他們的生活豐富了.
最後每一個參加者.
完了之後.
他們發現他們的生命改變.
他們的家庭改變.
他們的教會改變.
他們的社區改變.
所以新的一年.
讓我們跟自己說.
我不要多.
我要新.
我要生命的更新.
我要生命的質素.
我要更加像主耶穌.
當你開始這樣祈禱的時候.
當你開始這樣立志的時候.
你發覺你生命.
自己開始給神改變.
你身邊的人.
也慢慢地來到改變.
很多時候我們經常都在祈禱.
我們希望坐在旁邊的那個改變.
坐在旁邊的太太.
丈夫.

$^{481}$兒女.
來到改變.
但我們經常都說.
神好像不太聽我們祈禱.
祈禱很久都見不到什麼.
你試一下.
你祈禱.
你說神求你改變.
當你改變的時候.
你身邊的人.
就開始慢慢地來到改變.
我記得有一次.
我踩單車回家.
當時我在左邊.
那個車在右邊.
剛剛在紅燈前面.
是紅燈.
所以右邊的車停在那裡.
我單車就慢慢地踩.
紅燈準備停下來.
說時遲那時快.
當時我和車已經停在那裡.
我就慢慢地開.
因為夏天很熱.
我就開了窗.
當我的車慢慢到達的時候.
那個人就突然對著我.
「哈遲」.
對著我「哈遲」.
我就驚了.
我就走前一點.
我還好像看到有口水.
濕漉漉的.
當時我心裡很不安.
很不平安.
因為COVID.
去年暑假的時間.
在那一刻.
我心裡有一節經文.
在我心裡.

$^{521}$那時候在讀彼得前書.
彼得前書裡面就說到.
有困難的時候.
都要選擇善.
善事的善.
於是我就等了幾秒鐘.
很自然.
我就不知道為什麼.
大家記不記得.
很自然地.
我就回頭看著那個人.
我就和他說.
Bless you.
留意我的態度.
Bless you.
不是這樣.
Bless you.
祝福你.
那個人都呆了一夜.
Thank you.
你知道嗎.
我說完這句話.
Bless you.
祝福你之後.
我內心裡竟然突然之間.
沒有那種搞錯.
或者我對COVID那種.
不知道有沒有中招.
那種恐懼的情緒.
突然之間沒有了.
這些混亂的思緒.
沒有了.
隨之而來.
給我有一個醒覺.
我發現原來每一件事.
我們都有一個選擇.
我們每一件事.
我們都有一個選擇的能力.
我們可以選擇以惡報惡.
也可以選擇以善報惡.

$^{561}$而基督徒.
你和我是有這樣的能力.
在我們裡面的.
所以當我們希望今天.
不是只是多.
而是生命的更新.
你發覺當你自己更新.
當你自己改變.
當神的靈在你裡面的時候.
你改變.
你身邊的世界.
身邊的人.
也都因為你的改變.
而他們改變.
這個就是這麼奇妙.
神蹟的豐富.
就是耶穌來到我們生命裡面.
所以豐富.
豐富不是更加的多.
是更新.
最後豐富.
是和別人.
和別人分享我們的生命.
經文裡面.
約翰給我們看到的.
是耶穌.
耶穌他自己的豐富.
在《二水變酒》這個故事裡面.
我們見到約翰一開始.
我一開始的時候就說.
我們不可以看表面.
要讀背後的意思.
所以經文裡面.
我們見到這些字眼.
就是第三日.
我的時候還沒有到.
顯出了他的榮耀來.
第三日我的時候還沒有到.
顯出了他的榮耀來.
這些其實都是說到.

$^{601}$主耶穌之後.
掌襲自家.
奉獻自己的生命.
而這個奉獻的生命.
就是一個豐富的生命.
在《約翰福音》十二章.
都有說到這件事.
耶穌回答他們說.
人子得榮耀的時候到了.
我實在告訴你們.
一粒物似不落在地裡.
死了仍舊是一粒.
若是死了.
就結出許多似粒來.
若有人服侍我.
就當跟從我.
我在那裡.
服侍我的人也要在那裡.
若有人服侍我.
我忽必尊重他.
所以從一開始.
耶穌豐富的生命.
耶穌的生命是豐富的.
耶穌豐富的生命.
就是告訴我們.
祂來了.
不是要人服侍祂.
當時的豐富.
就是有很多的僕人.
就是豐富.
因為多僕人,多奴僕.
就是代表了你是富有的.
你是高上階層的.
但是耶穌說我來了.
我不是要人服侍我.
而是我去服侍人.
我們見到耶穌.
是跟人分享自己的生命.
用神的愛去服侍.
而耶穌也是叫門徒.

$^{641}$來這樣做.
在我們教會裡.
有一對宣教士.
他們本來就在這個地方.
來去事奉.
他們在這個地方事奉.
事奉了大概五,六年的時間.
自己出錢,出力.
阿魯巴這個地方.
義務的性質.
作為丈夫.
在那裡很努力地事奉.
我舉幾個例子.
有多努力.
他平時在多倫多.
是不開車的.
很少開車.
但是他去到阿魯巴.
他們有一輛小貨車.
他要開車去到不同的超級市場.
接載那些小朋友.
因為他們的父母是超級市場的擁有者.
是老闆.
小朋友放學後.
沒有事情做.
他就開著這輛車載他們去教會.
去上學,上課.
另外.
這位弟兄平時不是很會說話.
他又要學習怎樣去講道.
當六,七年之後.
宣道會正式宣佈.
宣道會認識猜派出去的宣教士的時候.
大概在宣佈之前的幾個月.
這位弟兄很不幸.
發現了一個急性的癌細胞.
很可惜在幾個月之後.
被天父接了去.
在人看來真的很可惜.
一個這麼年輕的弟兄.

$^{681}$當我自己回頭看的時候.
我在他在醫院的時候.
我都有探望他.
在臨離開這個世界之前.
我們一起來聊天.
在他出事之前.
去到阿魯巴這六,七年的時間.
五,六年的時間.
我和他來說.
其實你這五,六年的時間.
你提早退休.
辭職不做.
然後去到阿魯巴這個地方.
其實你可以在這五,六年裡面.
繼續工作.
然後才去阿魯巴.
但是猜也猜不到.
這五,六年是他最後的人生.
但是他就將這五,六年.
最好的時間.
給了天父.
奉獻給天父.
奉獻給阿魯巴.
阿魯巴這幫弟兄姊妹.
當他離開了這個世界之後.
太太將一些骨灰.
放在教會那裡.
在阿魯巴教會的十字架下.
作為紀念.
阿魯巴的弟兄姊妹.
很感動.
所擺上的心機.
所擺上的時間.
所擺上的愛心.
這位太太.
今天也繼續在這個地方.
來事奉.
所以我們看到.
豐富的生命.
不是要人服侍自己.

$^{721}$豐富的生命.
好像耶穌一樣.
是服侍別人.
去服侍別人.
其實我再說一個真理.
奇妙的真理.
讓大家知道這個服侍.
當你去服侍人的時候.
就得到這個豐富的生命.
當你服侍人的時候.
就得到這個滿足.
我舉多一個例子做結束.
可能我們有些人.
曾經也進過醫院做手術.
做完手術之後.
我們也會怎樣呢?.
也會痛.
這裡痛,那裡痛,腳痛.
走不動,不想走.
但是醫生護士會怎樣跟你說呢?.
不行,不可以坐在這裡.
你坐在這裡就繼續痛.
坐在這裡沒有幫助.
所以你要繼續走.
你要走多一點.
你動多一點.
這樣就會好起來.
一走幾次.
身體就會慢慢好起來.
其實我們的服侍.
也是同一個道理.
你要人服侍你.
就好像你坐在這裡.
什麼都不動.
你不會生命滿足的.
你只會空虛.
你只會越來越埋怨人生.
你會跟人比價.
看一樣東西.
你看都看不順眼.

$^{761}$這個不喜歡,那個不喜歡.
為什麼這樣,那樣.
但相反當你開始去服侍.
當你參與的時候.
當你付出.
當你流汗的時候.
當你奉獻給神的時候.
當你工作的時候.
當你處處為人著想的時候.
當你打算不是拿人給的時候.
你會發現你的生命是很滿足.
你心裡面是很平安穩妥.
你知道為什麼嗎?.
因為這個就是神給人豐富的秘訣.
這個是一個屬靈健康.
唯一一個的途徑.
沒有其他方法.
我也想有其他方法.
不過沒有.
神給人豐富的秘訣.
唯一一個屬靈健康的途徑.
就是去付出.
所以新的一年我在這裡祝大家.
有一個欣賞佳音的生命.
在主裡面有一個更新的生命質素.
就是將你的生命奉獻出去.
付出去.
這樣你就會更加的健康.
更加的開心.
更加的平安.
因為神都是這樣將他的兒子給了我們.
耶穌也是這樣將他自己給了我們.
願主祝福大家.
我和大家一起去祈禱好嗎?.
親愛的天父我們多謝你.
多謝為了成會的弟兄姊妹院上的感恩.
在疫情期間.
我也有來看到他們的崇拜.
看到教睦同工他們在網上所擺上的.
真是一個很大的祝福.

$^{801}$不單止對成會一個小小的教會.
一個小小的地方.
藉著這個小小的地方.
祝福世界各地的人.
所以求主你幫助在教會裡面的弟兄姊妹.
讓我們都能夠好像成會那個理念一樣.
使用我們.
以致我們能夠將這個生命和其他人分享.
而當我們肯將自己和別人分享的時候.
我們發現我們裡面就有平安.
我們裡面就有豐富.
我們裡面就有滿足.
願你致福給我們.
加我們的力量.
去成為你自己在地上的手和腳.
可以代理在這個世界上.
好像主耶穌說.
我們能夠作光,作業.
聽我們禱告奉主耶穌基督.
得勝名字祈求.
阿們.
安穩.
\newpage



\section{17:1-7}
\label{sec:qMFmuMv64OQ}
\textbf{Before Whom I Stand (1 King 17\_1-7) - Tommy Lee}
\newline
\newline
連結: \href{https://youtube.com/watch?v=qMFmuMv64OQ}{\texttt{https://youtube.com/watch?v=qMFmuMv64OQ}} ~~~~ 語音日期: 2023-02-13
\newline
\newline
\hyperref[sec:Cv_wNlk__oQ]{\small{< < < PREV SERMON < < <}}
~
\hyperref[sec:index_chronic]{\small{[返順時目]}}
~
\hyperref[sec:index_scriptual]{\small{[返順卷目]}}
~
\hyperref[sec:8B7OdFQev9A]{\small{> > > NEXT SERMON > > >}}
\newline
\newline
$^{1}$guest speaker Tommy. Thank you very much. Very good job, very good job. Thank you.
very much. Hey let me ask you a quick question for those especially those.
actually doesn't apply to those who are online but more so in person. Who watches.
football here in this room? Oh am I am I on?.
Is it on?.
Good? All right perfect. Who watches football here in this room? All right not.
too much. All right but if you do come to the Super Bowl party all right I want to.
at least take one minute to give you some basic information so you can come.
and know what's going on. Here are the two teams that are playing today. The.
Philadelphia Eagles and the Kansas City Chiefs. The two quarterbacks Jaylen Hurts.
on the Eagles very strong Christian believers. Patrick Mahomes Kansas City.
Chiefs very strong believers as well too. All right your three guys that you guys.
need to pay attention to. There is Chris Jones their defense he's the guy on my.
left. Patrick Mahomes on my right. The middle guy with the C is Travis Kelsey.
who is the tight end. He's your best receiver on there. The other guys don't.
really matter so just pay attention to those guys. The third Jaylen Hurts is.
right in the middle. Then you have AJ Brown and then from that point up Miles.
Sanders to my right. The coach on that team is Andy Reid who is currently the.
coach of the Kansas City Chiefs but he was the former coach of the Philadelphia.
Eagles so that's why they're calling it that Andy Reid Bowl. All right so at.
least some basic information that when you're sitting there and you're talking.
with someone you're not asking so what quarter are we on and how many baskets.
have they made. All right so and then also you you know this I mean a lot of.
times we're Super Bowl everyone pays attention to the commercials and all of.
that stuff so enjoy it's the last game of the season but it's a good time of.
fellowship as well too. Well this morning we're you know you guys are stuck with.
me it looks like you guys are stuck with me once a month here and so you get to.
hear my thoughts every single month and every single year I'm not sure if you.
get a chance to do this every year I like to pick a character or I like to.
pick a book in the Bible and for that entire year that's who I study. Now at.
the same time I totally understand a lot of times there's these Bible one-year.
plans and all that stuff and I appreciate that I like going deep into.
something and so whether it's me a lot of times at one year I studied the life.
of David then I spent some time studying Solomon I actually spent about a good.
nine months studying John chapter 1 all the way to John chapter 7 to really.
fully understand it. You saw that with Joseph all of that I want to do that.
with Elijah and if you ever get a chance just spend dive deep into and say Lord.
what is it that you're teaching me through this character or through this.
book and you're reading the same verses over and over and over again. Unlike Mark.
who has to prepare a sermon once a week I get a chance I have the luxury.

$^{41}$of preparing it once a month and so for the whole month that's what I study so.
for the last month I've just been studying 1st Kings 17 1 through 7 and.
I'm always asking Lord what is it that you're asking me what is it that you're.
saying to me I just keep rereading and rereading and rereading because there.
are days you sit there and you look at the same passage for the 50th time and.
it's something you realize wait a minute I never saw that before it makes.
scripture come alive what as you begin to understand it so hey today as we dive.
into 1st Kings chapter 17 and look at the lives of Elijah I do want to add in.
one bit background information that I think sometimes it's very good to.
understand is I think a lot of times in order to fully appreciate the life of.
Elijah you've got to understand the historical background of what was going.
on prior before Elijah came there was kind of a stability probably not perfect.
but for over a hundred years Israel was essentially led by three kings you had.
Saul who was your first king and he was probably stable he was the strongest of.
everyone he was the most athletic everyone on paper he fit what a king.
should be he was that warrior that the people chose but the problem with Saul.
is he had character flaws he had character issues that pretty soon.
started affecting his leadership to the point where he was disqualified from.
leadership and God had to find a new leader and we know the story he finds it.
in this young shepherd boy in David when Samuel the Prophet goes to Jesse he's.
in there bring all your sons I want anoint the next king and Jesse forgets.
David because he's out in the field tending his father's sheep story goes he.
defeats Goliath and later on he is chased after Saul who tries to murder.
him all of that and he ultimately becomes king he doesn't become just any.
king he becomes the greatest King Israel besides Jesus later on becomes the.
greatest king that Israel ever sees he builds the first temple everything but.
he makes one fatal flaw what happens with him in Bathsheba and eventually.
with Uriah the Hittite completely destroys his kingdom his leadership and.
he now even though he regrets it there is ultimatums there is all of the.
repercussions to the mistakes that he makes and pretty soon it goes to his.
other son Solomon well Solomon unlike David and also Saul was the smartest of.
all three of them he merely asked God for wisdom and God gave him wisdom and.
he was the wisest of all the Kings but lust and women got in the way and pretty.
soon he was deterred because his focus once on other things besides pleasing.
the Lord and here we are at the end of Solomon's rule a civil war breaks out.
the stability that they have faced even though it was run by weird leaders here.
and there a civil war breaks out and there the nation is divided into a.
northern kingdom which is referred to as Israel where the story of Elijah takes.
place the northern kingdom and a southern kingdom referred to as Judah in.
the northern kingdom they had 19 Kings during that span and all of them were.

$^{81}$wicked all of them were wicked according to scripture it's sort of like what.
what happened during Joshua during the days of judges when Joshua passed away.
these judges took over and a lot of times one of the phrases that you see in.
judges over and over again is that they did evil in the eyes of the Lord same.
thing is happening here in the southern kingdom they had 17 Kings now out of the.
17 eight of them followed the Lord but nine of them were wicked they were.
taken over later on by the Babylonians and it becomes the setting for Nehemiah.
and Ezra because the walls of Jerusalem are destroyed and you remember Nehemiah.
he is the cupbearer for the king he has a burden to rebuild those walls and now.
we dive into first Kings chapter 17 verses 1 to 7 I don't need those slides.
anymore if you guys mind taking that off as well - first Kings chapter 17 verses.
1 to 7 now Elijah the Tishbit of Tishbi in Gilead said to Ahab as the Lord the.
God of Israel lives before whom I stand there shall be neither dew nor rain.
these years except by my word let's break it down in terms of this passage.
it's first a little bit now Elijah a lot of times now for those who I got a.
chance to meet with a bunch of different Moody students and sometimes you want to.
look at the original languages I still remember taking my two semesters at.
Greek and you're always looking at the original Old Testament language that was.
written in Greek I don't always like going there but there is one point I do.
want to make about that Elijah in the original Hebrew language because the Old.
Testament was written in Greek in Hebrew Elijah is translated as literally my God.
is Jehovah you have the L you have the job with the I in the middle literally.
my God is Jehovah that is what Elijah's is translate name is translated to.
Elijah the Tishbit of Tishbi in Gilead now where in the world is that if you.
ask the folks who wrote the original passage they don't know if you ask the.
people now who are scholars where this was they don't know they know it's.
somewhere in Gilead but do they know where it is not really it's probably.
some rural area small little town that no has been to my brother actually have.
a twin brother he recently moved from Chicago and he lives in Payson Arizona.
has any of you guys ever been to Payson Arizona it is two and a half hours away.
from the nearest airport it is an hour away from the nearest Costco the.
internet over there is 3g I went there to visit I'm sure what in the world did.
I walk myself into it feels like you've gone 10 15 years back I was trying to.
get an uber and they said there's no over I said how do people get around they.
said well there's this phone call of this guy that you call and he could come.
and pick you up I said really and so here Elijah comes from this small little.
town that no one the Tishbit of Tishbi in Gilead and he says to Ahab is the.
Lord the God of Israel lives there shall be neither dew nor rain these years.
except by my word so somehow this guy my God is Jehovah translated that's what.
his name means who comes from this rural area it's a meeting with the king how he.

$^{121}$gets the meeting with the king I have no idea somehow he gets a meeting with the.
king so he stands there and he just says in one word because we didn't know.
anything about Elijah he just makes his first appearance in chapter 17 and he.
says to him Ahab there'll be neither dew nor rain these years except by my word.
okay he I get the feeling he's kind of rough around the edges he kind of says.
whatever is in his mind you recognize some of those guys right I mean we see.
that even in scriptures we see that with the prophets Jeremiah was kind of rough.
around the edges Paul definitely was rough around the edges Paul said.
whatever was I I think Paul lacked a lot of that social grace a lot of that so.
the social interaction John the Baptist was definitely like that as well too and.
here Elijah kind of just says whatever is on his mind and he directly says to.
him you're not gonna get dew nor rain these years except by my word now why.
was this so important why was holding off dew and rain so important to the.
story to understand that you don't need to turn to that let me turn it read to.
you first Kings chapter 16 verse 9 first Kings 16 in the 38th year of as Asa king.
of Judah southern kingdom now Ahab the son of Omri began to reign over Israel.
northern kingdom where the story takes place northern Sunder Kingdom and Ahab.
the son of Omri reigned over Israel in Samaria for 22 years so in the northern.
kingdom Samaria is the capital now if you think through John chapter 4 Jesus.
and the Samaritan woman Jews don't go through Samaria you're beginning to.
understand why capital city is in Samaria and Ahab the son of Omri did.
evil in the sight of the Lord more than all who were before him so what you're.
saying there and remember we talked in background there was 19 kings in the.
northern kingdom all of them were bad there isn't a good one all of them were.
bad you're telling me Ahab was the worst of all of them he was the most evil of.
all of them that's it and as if it has been a light thing for him to walk in.
the sins of Jeroboam the son of Nebat he took for his wife Jezebel the daughter.
of Ethel king of the Sidonians wait Jezebel why is that important if you go.
and we do not have time to go through each of the kings in first Kings it.
always mentions the name of the king but never the wife never the Queen for the.
first time we hear Jezebel we see Jezebel the wife and the Queen is.
mentioned later on as we look throughout the stories Elijah in the next few months.
that we talk about this Jezebel plays a huge role because she's even more evil.
than Ahab I actually think she's the strings behind why Ahab's the person.
that he is Jezebel plays a key role here in this and so the author wants to tell.
us hey look Ahab was bad there is Jezebel lurking and she is pure evil and.
that's the story that we have and from that point on they went and served Baal and.
worshipped him he erected an altar for Baal in the house of Baal which he built.
in Samaria the capital all right why is that important what is Baal then? Baal is.
a god of lightning and storm and responsible for the fertility of the.

$^{161}$land they built Baal because they trusted and were gonna place their trust.
in Baal that he was gonna make their land fertile he was gonna bring rain he.
was gonna bring everything and they were all gonna benefit that and here Elijah.
comes and says really you're gonna put a god that you're gonna worship in the.
nation of Israel so that they could bring you fertility of land interesting.
I'm gonna bring my God and up you one and I am gonna withhold rain and dew.
until I say so let's see what your God does with that I think as Christ.
followers we tend to do the same thing don't we it's not like we have statues.
it's not like we built altars like the Asherah pole all that stuff I think.
sometimes we do have these little gods that we have in our life things that we.
place our trust in that we think will bring our success and achieve our goals.
maybe that's in the form of relationships if I could get that one.
relationship that guy or that girl or if I could get that one job or their career.
path that job that opens the door and that would achieve all my goals all.
right if I can have this much money and we establish these little gods we're.
still following Christ we still believe in Christ but God Christ and God becomes.
a little God amongst all these other gods that we have and we pretty soon.
start putting our trust in multiple things now here's the one thing it.
doesn't happen overnight it's not like one day we wake up one day and say all.
right Lord I followed after you but you're not it and so I'm gonna follow.
everyone in people and then as a result here we go it doesn't happen like that.
it happens in little decisions that we make in the adventures of our life let.
me give you an example I had a pastor friend who I loved dearly I really.
really respected he was so influential influential on my life growing up as a.
pastor he soon started getting frustrated at church politics it started.
happening he says I don't like this church politics I'm sick and tired of it.
then from that point I started becoming money problems we don't have enough.
money as a pastor and with four or five kids growing if I had enough money life.
wouldn't be so hard eventually he left the pastorate and church and then want.
to become a pastor he got a secular job in the corporate world he got more money.
and he started enjoying that money and he started realizing wow this money.
helps open doors for many different things and you begin to see pastorate.
didn't have money got tired of church politics from that point on now he has.
money and he's beginning to put his trust in that money now that money opens.
doors for him to do traveling to go to nicer restaurants all of those things I.
like this pursuit this is pretty good it opens doors on Sundays it's okay if I.
miss miss church service it's okay if I don't get fellowship and you start.
distancing away politics he soon gets frustrated with a lot of the right wing.
and all the other stuff so he gets frustrated and pretty soon he stops.
believing the Bible the Bible is no longer the central authority maybe the.

$^{201}$Bible was wrong maybe people aren't imperfect when he wrote the Bible and.
now he's completely gone I sat there with him not too long ago he says how do.
you still have hope because I lost hope a while ago it doesn't happen overnight.
we walk in there when we come to accept Christ with our focus we're gonna follow.
Christ we're gonna do all our but alongside of it different aspects in our.
life start getting in our way it starts affecting our life and so every single.
day right we have to and here's a good illustration we're so used to in this.
technology in Chrome and Google world and everything like that to enter our.
password login password and suddenly a notice pops up says would you like to.
remember this password yes because I don't want to sit there manual entering.
my password every day but yet at the same time we log in and everything's.
there we get our Google Chrome we get our Google Drive we get Dropbox.
everything like that but every single day we have to log in again every day we.
have to manually log in and choose to follow Christ and commit to follow.
Christ each and every single day because when we do not we start establishing.
these little gods and it begins to move us away from who our primary focus is.
life is not sitting there what these days I just realized chat GPT is not.
chat GPT where I sit there and say hey what's the sermon for first Kings 17 in.
two minutes it spits out my sermon it's not me sitting there and logging in my.
password and for the rest of my life it logs me in no every day it logs me off.
and every day I got restart the process I got redo the process I got recommit.
every single day now look at here's the one thing I forgot or actually not.
forgot but here's the part that I glance over look at verse 1 right here's what he.
says you're gonna have neither do nor rain these years and he says as the Lord.
the God of Israel lives underline this portion this next phrase I'm gonna read.
underline it highlight it because it becomes the theme of this passage and.
all of Elijah's life before whom I stand that's his thing that's his motivation.
it is the Lord before whom I stand and that is why I'm confronting you a how.
that is why I'm doing what I'm doing before whom I stand motivation informs.
action the actions of a Christian is good you could go to church and do all.
these things you go small group all of these things great.
what's your motivation Elijah's motivation is before whom I stand your.
motivation informs your action your commitment before the Lord if you.
believe truly that you have the same vision that Elijah has that I stand for.
here because I follow the Lord before whom I stand it changes how you approach.
work it changes how you approach your life as a spouse how you treat your wife.
how you treat your husband it affects you and how you treat your family if.
facts you in terms of how you treat people you don't like your boss your.
co-workers you don't like you excel at doing work every day you excel and not.
being lazy at work you excel at school you don't cheat at school because all of.

$^{241}$those things because my motivation is before whom I stand.
it's that your motivation is that my motivation or do we just go through the.
routines of action every single day Elijah my God is Jehovah comes out of.
nowhere from tish be the way knows where tish be is a tells a ham this is what's.
gonna happen why because my motivation is from the person before whom I stand.
keep going verse 2 and the word of the Lord came to him depart from me and turn.
eastward and hide yourself by the brook a chariot which is east of the Jordan.
you shall drink from the brook and I have commanded the Ravens to feed you.
there so he went and did according to the Lord the word of the Lord he went.
and lived by the brook chariot that is east of the Jordan and the Ravens.
brought him bread and meat in the morning and bread and meat in the.
evening and he drank from the brook and after a while the brook dried up because.
there was no rain in the land a couple things I want to say about this number.
one I was very intrigued by the Raven and I spent a lot of time trying to.
understand so how did the Raven bring him food did the Raven just sit there.
and just go on his shoulder and bring him food and all that so I had to.
understand that a little bit so pretty much I think one of the reasons why God.
was telling him about the Raven you know how Ravens bring food they sit there and.
they get wrapped the food and they fly to mountain cliffs and they put it in.
different holes and so I'm assuming that as different scholars are writing that.
Elijah just had to spot where the Ravens were coming every day what holes they.
were gonna put in and there his food was versus a Raven lands on his arm and.
brings the food and eats it and so there a lot of times that's how God is and.
also another thing that's important God opens and closes doors as he sees fit.
God opens and closes doors as he sees fit why in the world did he leave Elijah.
away was it because he wanted him to rest maybe did he want Elijah be.
protected because Ahab was gonna go after him maybe did he want Elijah to.
spend some time in solitude as he can do the work on his character his inner life.
maybe maybe have been combination of all three but God has a plan and God opens.
and chooses and closes door as he sees fit but yet at the same time God is in.
the midst of always working in your life God is always in the midst at pulling.
you at certain points in your life taking you away from things to continue.
to allow you to be in solitude to learn the inner workings the inner character.
the inner struggles your mental health your emotional health of your life he.
did that with Moses he did that with Joseph over in jail he did that with.
David when he was a little shepherd boy tending sheep fighting bears and lions.
he did that with Jesus for 40 days 40 nights he did that when Mary and Martha.
Martha had to learn that because Martha was so busy preparing for the arrival of.
Jesus that she forgot what was important God begins to work in the quietness the.
solitude of his life and begins to transform you in the person that he.

$^{281}$wants you to be take a quick look at the end of 17 the last verse we started the.
story by introducing you to Elijah the Tishbit of Tishbi from Gilead at the end.
of 17 the woman said to Elijah now I know that you are a man of God and that.
the word of the Lord in your mouth is truth he is no longer Elijah the Tishbit.
of Tishbi in Gilead but he is now a man of the Lord and all his word is truth.
from the Lord God begins to work in the transformation in all of our lives to.
shape us into the person he wants us to be but he leads us on that journey and.
too often is done in what he calls us to do but it is also done in solitude as he.
begins to work in all of our hearts to help you to become that person a couple.
applications as we wrap up number one God uses people we least expect to do.
his work God uses people we least expect to do his work we don't get a chance we.
don't get any understanding of who Elijah is he appears out of the blue and.
we immediately know that his name means my God is Jehovah we also know that he.
comes from a city or a town or rural area that we're not even familiar where.
he lives and he comes out of nowhere and God uses people we least expect to do.
his work and that means a lot of times never count yourself out never disregard.
yourself of how God can use you you may not be smart enough that's okay you may.
not be talented enough that's okay you may be shy you may be a little awkward.
socially awkward at times that's okay you may have baggage in your life all of.
those things that's okay never count yourself out never count yourself out.
because of your age definitely never count yourself out because of your sex.
because men and women could do wonderful things never count yourself out because.
of your skin color never count yourself out because God uses and worse in the.
people we least expect it to number two God's methods are unique and different.
God has us on each a unique journey all of our journeys are different our life.
sometimes is a culmination of bad or good that God weaves in and now.
continue to mold you and shape you into the person that he wants you to be our.
interactions our relationship with our family our siblings our friends our.
experience within church all of that because of God's methods are unique and.
different he may call you here in Chinatown next week he may call you to.
Hong Kong he may call you in the full-time ministry pull you out full-time.
ministry all of those things his methods are unique and different sometimes we.
will never understand it sometimes it and most times it will not go according.
to plan because but at the same time number three God's direction include.
God's provision God's direction includes God's provision because as much as he.
calls you he will provide now not in the way you want him to provide probably not.
in a very smooth way not in an easy way but God's direction includes God's.
provision just like the Raven comes just like the Brooke comes God provides.
probably not in the way we expect it to but he opens and closes doors as he sees.
fit and we have to be careful those doors be watching for those doors.

$^{321}$praying for those doors praying for open and praying for closed doors that he.
continues to make known those doors far for guide us in everything that we do.
and we look upon our entire journey we think that we're studying for one career.
well that door suddenly closed it moves to another career I have a friend whose.
house I was just at the other day I sat down with her and her husband after 25.
years in one company she was just informed this past week she's being laid.
off with no severance after 25 years it's the only thing she's done for 25.
years since graduating from college and she's wondering Lord what are you asking.
me to do because I am NOT trained in anything else and she's asking for the.
Lord to make it clear to her what she wants her to do you may be asking that.
same thing for those who are in college right now you're a senior in college and.
all your life you followed school a direct path and now you're about to.
graduate you don't know what that path is some of you are thinking about global.
missions or returning back home to your country or wherever you are or switching.
careers or even getting married you don't know you pray for open doors you.
pray for closed doors God uses people we least expect to do his work his methods.
are unique and different his direction includes God's provisions God's.
direction includes his provision for our motivation is before whom we stand that.
you and I are here because of one thing the Lord before whom we stand that is.
our motivation our motivation informs our actions it's why we do what we do.
it's why we treat people the way we treat people because of the motivation.
that we have for the Lord and you will see as we journey through Elijah will he.
be perfect no he will not be perfect but he will get to that point that he paves.
the way for John the Baptist who paves the way for Jesus Christ who paves the.
way for the gospel that you and I believe in and on that day later on when.
Jesus stands on the transfer on the mountain for his transfiguration there.
stands Moses and there stands Elijah Elijah the tishbit of tishbi from.
Gilead is standing right there God opens and closes doors as he sees fit and uses.
us in ways we can never ever expect let's pray father thank you for the.
wonderful story of Elijah a guy who seemingly came out of nowhere who was.
kind of rough around the edges and kind of says whatever is on his mind but.
somehow you used him and you will continue to use them on this journey.
because he really believed that he stood before you the Lord before whom I stand.
I hope that is a reminder for all of us that we have that same model in our life.
we have that same motivation in our life that motivation informs us in any.
everything that we do thank you for the life of Elisha thank you for the words.
of scripture for guiding us and every single day Lord challenge us to log off.
and log on and begin another day of following after you Christ would pray.
mark thank you.
\newpage



\section{箴言 3:5-12}
\label{sec:8B7OdFQev9A}
\textbf{我兒... (箴言3\_5-12)- 譚文倩莊傳道}
\newline
\newline
連結: \href{https://youtube.com/watch?v=8B7OdFQev9A}{\texttt{https://youtube.com/watch?v=8B7OdFQev9A}} ~~~~ 語音日期: 2023-02-13
\newline
\newline
\hyperref[sec:qMFmuMv64OQ]{\small{< < < PREV SERMON < < <}}
~
\hyperref[sec:index_chronic]{\small{[返順時目]}}
~
\hyperref[sec:index_scriptual]{\small{[返順卷目]}}
~
\hyperref[sec:vHt7tPIvb_w]{\small{> > > NEXT SERMON > > >}}
\newline
\newline
箴言 3:5-12
\newline
\begin{longtable}{cl}
\hline
\hline
章節 & 經文 (和合本修訂版)\\
\hline
3:5 & \begin{tabularx}{0.7\textwidth}{X} 你要專心仰賴耶和華，不可倚靠自己的聰明， \end{tabularx} \\ \\ \relax
3:6 & \begin{tabularx}{0.7\textwidth}{X} 在你一切所行的路上都要認定他，他必使你的道路平直。 \end{tabularx} \\ \\ \relax
3:7 & \begin{tabularx}{0.7\textwidth}{X} 不要自以為有智慧；要敬畏耶和華，遠離惡事。 \end{tabularx} \\ \\ \relax
3:8 & \begin{tabularx}{0.7\textwidth}{X} 這便醫治你的肉體，滋潤你的百骨。 \end{tabularx} \\ \\ \relax
3:9 & \begin{tabularx}{0.7\textwidth}{X} 你要以財物和一切初熟的土產尊崇耶和華， \end{tabularx} \\ \\ \relax
3:10 & \begin{tabularx}{0.7\textwidth}{X} 這樣，你的倉庫必充滿有餘，你的酒池有新酒盈溢。 \end{tabularx} \\ \\ \relax
3:11 & \begin{tabularx}{0.7\textwidth}{X} 我兒啊，不可輕看耶和華的管教，也不可厭煩他的責備， \end{tabularx} \\ \\ \relax
3:12 & \begin{tabularx}{0.7\textwidth}{X} 因為耶和華所愛的，他必責備，正如父親責備所喜愛的兒子。 \end{tabularx} \\ \\
[1ex]
\hline
\hline
\end{longtable}
$^{1}$今天中文堂的講演是Joyce傳道.
將時間交給她.
今天有機會在大家當中靜讀.
今天的主題叫《我而》.
經文在《箴言三章》五至十二節.
我先跟大家讀一次.
你要專心仰賴耶和華.
不可以靠自己的聰明.
在你一切所行的事上.
都要認定它.
它必指引你的路.
不要自以為有智慧.
要敬畏耶和華.
遠離惡事.
這邊依自你的肚臍.
滋潤你的白骨.
你要以財物和一切初熟的土產.
尊榮耶和華.
這樣你的倉房必充滿有餘.
你的酒窖有新酒盈日.
我而你不可輕看耶和華的管教.
也不可厭煩他的責備.
因為耶和華所愛的.
他必責備.
正如父親責備所喜愛的兒子.
如果.
你相信主已經一段時間.
我相信這段經文.
大家都很熟悉.
特別這個《箴言三章》五至十二節.
你要專心仰賴耶和華.
不可以靠自己的聰明.
在你一切所行的事上.
都要認定祂.
祂必指引你的路.
表面看來.
這段經文一點也不難明白.
似乎是說人要專心仰賴神.
不要自己冷醒靠自己的聰明.
認定神是我們生命的主.

$^{41}$如果人這樣做的時候.
神就會給我們一條出路.
一條方法或最好的結果.
作為一個基督徒.
特別我生了女兒之後.
這句經文就是我的座右銘.
我經常祈禱問神.
給我智慧如何養育我的女兒.
依照聖經的教導.
以致我不要得罪神.
特別是養小朋友的時候.
小時候有一樣東西很煩擾我.
就是每一次上學或外出.
他就會在抽屜裡找衣服.
我女兒找衣服的方法是這樣的.
找到一團亂.
如果你認識我.
我是一個很整齊的人.
這件事令我很苦惱.
有一次洗完衣服.
一起在房間坐在地毯上.
再摺衣服的時候.
我又想起這段經文.
神祈禱.
主啊 你一定要給我一些方法.
如何教導我的女兒.
要選衣服的時候.
我自己想.
可以這樣揭開.
不用這樣抄.
一直祈禱的時候.
期間我有一個意念.
意念就是.
當我們摺好衣服的時候.
我就一手把我女兒摺的衣服推倒.
我女兒當然嚇了一跳.
媽媽 你在做什麼.
我就告訴她.
每一次這樣摺衣服的時候.
把衣服弄亂.

$^{81}$她現在的感受就是媽媽的感受.
我女兒點頭.
告訴我她明白.
當時我女兒六歲.
我當時當然在想.
主啊 你真好.
你立刻給我一個方法.
可以教導我的女兒.
這句經文實在非常好用.
但問題是.
特別是我進神學院讀書的時候.
發覺很多時候.
我們都會把經文應用.
隨著我們自己的感覺.
或者表面的讀法來解經.
來應用.
我不能否定.
經文這樣看的時候.
對我們有幫助.
只不過問題是.
其實神是否藉著那段經文.
真的這樣跟我們說話.
正確明白神話語的方法.
看經文的時候我們要看文本.
我們要看看.
神在當時.
跟當時的人說什麼.
為什麼要跟他們說這些呢.
當時的全境.
神學的反思.
跟耶穌 跟風音有什麼關係呢.
跟我們或今時今日有什麼關係呢.
神藉著聖經讓人去認識祂.
彰顯祂自己.
人需要讀經文.
在祂的名義上.
按照聖經的原意去閱讀聖經.
如果我們不謹慎地花時間去閱讀的時候.
明白神到底藉著聖經.
要跟我們說什麼.

$^{121}$就會好像我當天這樣.
隨著我自己的心意去應用.
將神的話語基本上是降級的.
當作只是一個很小很容易的座右銘.
為什麼說這些呢.
因為我在想今天的經文.
我自己一邊看的時候.
我自己也很熟悉這段經文.
我覺得有時候因為我們太熟悉那段經文的時候.
就很容易掉以輕心.
所以我很想特別在講到之前提醒大家.
我們先做個低頭祈禱.
(低頭祈禱).
首先呢.
我很想和大家總結一下.
《箴言三章》是講什麼的.
《箴言三章》的第一節.
十一節和二十一節.
都用「我而」來開始.
三個「我而」將三章分為三個段落.
第一個「我而」是一至十章.
回應《箴言二章》所提到品格轉化之後的後果.
如果你有留意細心閱讀.
其實這段經文的一,三,五,七,九單數經文.
其實列出一些對神的孩子的忠告.
第一節.
「我而」不要忘記我的法則.
你心要謹守我的誡命.
第三節.
不可使慈愛誠實離開你.
要是在你頸項上.
刻在你心把上.
第五節.
你要專心養賴耶和華.
不可依靠自己的聰明.
第七節.
不要自以為有智慧.
要敬畏耶和華.
遠離惡事.
第九節.

$^{161}$你要以財物和一切初熟的土產.
尊榮耶和華.
兩個不要.
三個要.
兩個不可.
兩個你要.
一共有九個忠告.
雙數的經文.
二,四,六,八,十.
是說The Great I Am.
對孩子帶著鼓勵的條件.
因為祂必將長久的日子.
生命的年數與平安加給你.
第四節.
這樣你必在神和世人眼前蒙恩寵有聰明.
第六節.
在你一切所行的事上都要認定祂.
祂必指引你的路.
第八節.
祂必先醫治你的肚臍滋潤你的白骨.
第十節.
這樣你的艙房必充滿有餘.
你的酒窖有新酒盈日.
必生命長久.
必蒙恩.
必指引你的路.
必充滿有餘.
四個承諾.
九個應許.
第二個我兒在十一至二十節.
一開始就來一個很簡單的勸告.
提醒神的孩子.
管教當中的價值.
也再一次鞏固之前一至十節所提到的忠告.
第三個我兒.
二十一至三十五節.
被智慧塑造的品格.
生命有力不至跌倒.
不用怕.
晚上睡覺睡得好.

$^{201}$二十五至三十一節.
提到神的孩子不需要怕.
不要對人吝嗇.
若果你有能力的時候.
要盡能力去幫人.
不要推遲.
不要陷害人.
不要誣告人.
最後的三節.
就講到總結正直人和惡人最終的結果.
很快的.
剛才我跟大家總結了第三章.
經文所提醒.
我們再回來三章.
五至十二節.
再一次跟大家讀五至六節.
《真言三章》五至六節.
你要專心仰賴耶和華.
不可以靠自己的聰明.
在你一切所行的事上.
都要認定祂.
祂必指引你的路.
這兩節經文當中.
是講述三件事我們需要做的.
我都提及了.
我們要專心仰賴.
不可以靠自己的聰明.
要認定神.
很多時候我們做基督徒.
都會講很多基督徒的語言.
但其實.
你有沒有去想一下.
去研究一下.
究竟什麼叫專心仰賴神呢?.
什麼叫不可以靠自己的聰明?.
似乎應該是很容易的.
因為不靠自己聰明就是靠神.
但什麼叫靠神呢?.
認定神.
什麼叫認定神?.

$^{241}$一大堆基督徒時常掛在後面的詞彙.
其實是在說什麼呢?.
其實這兩節經文.
所講的這三件事.
要做的事.
答案是怎樣做呢?.
其實是在下面那幾節經文.
就是三章七至十二節.
不要自以為有智慧.
要敬畏耶和華.
遠離惡事.
這便醫治你的肚臍.
滋潤你的白骨.
你要以財物.
和一切初熟的土產.
尊榮耶和華.
這樣你的倉房必充滿有餘.
你的酒窖有新酒盈日.
我二.
你不可輕看耶和華的管教.
也不可厭煩他的責備.
因為耶和華所愛的他必責備.
正如父親責備所愛的兒子.
遠離惡事.
第二.
奉獻財物和一切初熟的土產.
三.
不輕看耶和華的管教.
不厭煩他的責備.
遠離惡事.
惡事在希伯來原文.
Raha是說有邪愛破壞或者傷害的意思.
箴言十四章十六節.
智慧人懼怕就遠離惡事.
願望人卻旁傲至此.
表面看來遠離惡事.
就是在說一些不敬畏神的人.
不應該做出一些傷害別人的事.
不僅是在行為上.
也在言語上甚至思想上.

$^{281}$創世記二十五章二十七至三十四節.
雅各趁著耳塑又累又肚餓.
就用這個熊頭湯換了他的名分和祝福.
這是一個很典型很容易明白的例子.
我想和你分享另一個見證.
我在亞利桑那認識一個太太.
她在十多年前被她丈夫拋棄.
他們分開之後她要匿躲.
當時她回教會很小.
她就祈禱求神.
讓我找到一間屋.
不僅是讓我自己住.
也可以開放給其他弟兄姊妹來聚會.
沒多久她真的找到一間合適的屋.
唯一的問題是那間屋的價錢超過了她的預算.
但她繼續祈禱反覆思想.
那間屋的位置和尺寸真的很適合.
她越看越喜歡.
雖然很掙扎.
但最後她也決定去前進.
去銀行封鎖那天.
她說去到門口.
突然間她的腦海出現了聲音.
問她是否真的願意為了這個夢想屋.
而放下她的一個熱誠.
這位太太是一個很有愛心.
很有善良招待恩賜的姐妹.
除了她想幫教會.
能夠買一間屋可以服侍弟兄姊妹.
她的另一個熱誠就是.
她很願意幫助一些有需要的人.
如果她用盡了她大部分的錢.
去買了這間屋.
其實她很難繼續去做招待或幫助其他人.
她站在門口反覆思量.
想來想去都覺得很不平安.
最後決定不買.
但決定不買之後.
她的心也很不安樂.
為何不安樂呢.

$^{321}$因為她想起如果不買這間屋.
賣屋的人又買不到她想買的屋.
第二她覺得在整個過程去找屋的時候.
代理人做了很多事.
她覺得這樣做對代理人很虧欠.
她就有個想法.
她就跟賣屋的人說.
不如這樣.
我把我本來買屋的錢借給你.
買你想買的屋.
你買到你的屋之後.
你才給我我的錢.
她也跟代理人說.
她說我真的不好意思.
你浪費了這麼多時間.
你有什麼損失.
用了什麼錢.
你讓我知道.
我願意還給你.
她跟我說.
Joyce.
其實我知道我是不需要這樣做.
也沒有什麼要求我這樣做.
她說但是.
我覺得這是對的事情.
是一件對的事情.
長話短說.
經紀也很快收到她的佣金.
事情有一個很圓滿的結果.
買不成屋.
沒有地方住.
這個太太就決定去參加短孫.
她參加了幾個短孫.
而在一個其中短孫的聚會裡.
她就認識了一個很好的弟兄.
今時今日這個弟兄就是她的丈夫.
他們結果在亞利桑那有一間.
我想比她之前想的夢幻屋更大.
還有一個客房.
他們家長期開放.

$^{361}$讓弟兄姊妹在那裡聚會.
而有很多人.
甚至不是基督徒.
都經常得到這兩夫婦的幫助.
根據經文原文的解釋.
遠離惡事.
其實不單止是說.
不做一些不好的事.
傷害人的事.
誣告人的事.
更加是說一個敬畏神.
尋求神心意.
聆聽神帶領.
信服神的人.
好像這位太太.
當她決定聆聽聖靈的聲音.
去作出決定.
甚至過程當中.
她自己是有吃虧的.
她甚至不知道.
之後沒有屋住的情況.
但她仍然選擇信服.
因為作為一個基督徒.
她知道她這樣做是對的.
第二.
奉獻財物和一切初熟的禱產.
很多年前.
我們在教會.
還跟弟兄姊妹.
有一個很深奧.
真的要把心心都投入.
分享的討論.
說的是奉獻.
其實應該是報稅前.
還是報稅後.
不知道你們有沒有曾經有過這樣的掙扎.
我們當時.
這個討論是很認真的.
現在聽起來好像很少.
但當時我們很認真地說了很久.

$^{401}$沒有答案.
但都挺有趣的.
回想起來.
到現在仍然歷歷在目.
當我去到神學院的時候.
我們也在課堂中討論.
大部分的神學生.
牧師,傳道人都沒有錢.
生活也很勉強.
我們都已經把我們的心門.
放上去奉獻給神.
去事奉了.
其實可不可以不受奉獻呢.
可不可以當我們的服事就是奉獻呢.
聖經裡面說到寡婦的兩個小錢.
不知道你們怎麼看.
耶穌對著銀庫坐著.
看眾人怎樣投錢入庫.
有好色財主往裡投了若干的錢.
有一個窮寡婦來.
往裡投了兩個小錢.
就是一個大錢.
耶穌叫門徒來說.
我實在告訴你們.
這窮寡婦投入倉庫裡的.
比眾人所投的更多.
因為他們都是自己有餘.
拿出來投在裡頭.
但這寡婦是自己不足.
把得一切養生的都投上.
神起了窮寡婦的兩個小錢.
因為她是把她用來生活的錢.
資源.
保障她將來的錢放上去.
相對其他人只是把自己剩下來的.
投入到那個箱子裡.
我們的教授當日.
鼓勵我們一班神學生.
要好好地操練金錢奉獻的功課.
操練將我們建立在金錢上的安全感.

$^{441}$人為自己保障將來的籌算.
換來降服.
認定神是我們的供應主.
若果做不到十一奉獻的時候.
或者你覺得真的很難領的時候.
他鼓勵我們嘗試從1\%.
2\%.
3\%開始.
不要放棄.
不要就這樣就算了.
求神幫助我們不要小看這個操練.
《箴言三章》五六節.
你要專心仰賴耶和華.
不可以靠自己的聰明.
在你一切所行的事上.
都要認定他.
他必指引你的路.
在奉獻金錢的敬拜動作中.
我祈求神幫助我們每一個人.
在奉獻的功課中盡忠.
在奉獻的功課中經歷神.
將貴重的,最好的擺上.
因為你願意神在你奉獻的動作中得著榮耀.
3.不輕看耶和華的管教.
不厭煩神的責備.
三章一節.
我而不可忘記我的法則.
你心要謹守我的誡命.
三章十一節.
我而你不可輕看耶和華的管教.
也不可厭煩他的責備.
三章二十一節.
我而要謹守真智慧和謀略.
不可使他離開你的眼.
三章十一節.
我而要謹守真智慧和謀略.
不可使他離開你的眼目.
不要忘記.
要謹守.
不可輕看.

$^{481}$不可厭煩.
要謹守.
不可使他離開.
你對上一次被神管教.
或者責備是何時.
或者想起管教的時候.
你第一樣想起的.
好像我這樣做媽媽.
就想起我的女兒.
或者你是做爺爺奶奶.
公公婆婆.
怎樣管教你的孫子.
甚至可能你有隻小動物.
上年年頭發生的事.
結果我和我的女兒鬧了一場大交.
其實從小到大.
我們很少有爭吵.
所謂的吵架.
其實主要是媽媽罵女兒.
這件事發生之後.
我一直很不開心.
我都沒有想過會不開心這麼久.
我不開心了很久.
因為我覺得我女兒的反應.
傷害了我.
再者我在她的行為上.
我也不覺得有什麼改變.
我自己心裡在想.
女兒明知她傷害我.
也知道我很不開心.
她也沒有改變.
即是她不太介意傷害她媽媽.
這件事過了之後很多個月.
其實我都一直很不開心.
大半年之後.
我女兒就找實習生.
有一次我開車回她的公寓.
她就跟我說.
媽媽我將會見工作.
她有很多問題要我回答.

$^{521}$我已經預備好了.
怎樣怎樣.
她問我有沒有興趣.
想看一看我的預備.
當然求之不得.
我說好呀好呀快點給我看.
我女兒就發了文件給我看.
我打開文件一邊看的時候.
我發覺我女兒長大了很多.
很成熟.
我自己覺得很安慰.
看到文件最後有一條問題.
就問她.
你怎樣處理衝突.
請舉例.
想不到我女兒.
就將那次的事寫下來.
她說她不太明白.
我當日為何發這麼大脾氣.
她說起初其實我嘗試跟我媽媽解釋.
但我發覺她都聽不進去.
我決定就閉嘴.
反過來聽聽我媽媽罵我甚麼.
當她媽媽繼續不停的哭泣.
她說她突然間明白了一件事.
她說她發覺其實她媽媽.
她媽媽很愛惜她.
她媽媽因為很擔心她.
很怕她有事.
所以反應這麼大.
怕她一個女孩晚上外出危險.
我女兒在這件事後做了一些事.
我自己沒有留意.
因為那次之後.
她每次外出都會跟朋友拍照.
發給我和她爸爸.
我當時還未看這份文件時.
我經常覺得我女兒有點奇怪.
為何經常發這張照片給我.
為何這麼有趣.

$^{561}$原來她是有原因的.
她很想我不擔心她.
不知道她跟誰人外出.
《歌羅西書》三章二十一節.
不要惹兒女的氣.
恐怕他們失了志氣.
回想起來.
我自己很放縱自己的血氣.
我的脾氣去管教和責備.
結果神在這件事上.
去管教和責備我.
藉著我女兒去找工作.
讓我看到自己做媽媽的自以為是.
以為我最適合.
以至我對著她發脾氣罵她是可以的.
讓我自己的情緒和脾氣.
主導了我的言行.
《心言三章》十二節.
因為耶和華所愛的,他必責備.
正如父親責備所喜愛的兒子.
正如我剛才所說.
我很希望能夠依照聖經的教導.
神的引領去做一個媽媽.
我看完我女兒的文件後.
明天我打電話給我女兒說對不起.
我女兒說:It's OK, Mommy.
我知道你很疼愛我.
《心言一七》.
驚威耶和華是智慧的開端.
如每人藐視智慧和憤悔.
我們真的不要做一個愚昧的人.
就好像一個不認識神的人.
藐視了耶和華的智慧和憤悔.
其實不論你是什麼身份.
你是什麼年紀.
只要你是神的兒女.
神一定必定管教責備.
因為這個是父親的愛證.
是父親的愛證.
是證明父親的愛.

$^{601}$是父親愛兒女的一個印證.
你要專心仰賴耶和華.
不要靠自己的聰明.
在一切所行的事上.
你都要認定神.
以至神可以指引你的路.
是真言智慧的提醒.
我兒.
你要遠離惡事.
要敬畏耶和華.
尋求神的心意.
你要聆聽神跟你說的話.
要信服.
我們要將貴重的最好擺上.
原因是我們要將榮譽再一次歸還給神.
願以神在當中得著榮耀.
最後不傾向耶和華的管教.
不要厭煩神的責備.
我們一起低頭祈禱.
親愛的主我們的神.
你的說話帶著改變生命的能力被釋放出去.
願人在有生的年日裡.
我們真的專心仰賴.
不要靠我們自己的聰明.
在一切事上我們都要認定你.
不傾向聖靈的提醒.
我們願意信服.
讓我們的生命.
你在當中得著榮耀.
祈禱奉主耶穌.
基督成名祈求.
阿們.
\newpage



\section{}
\label{sec:vHt7tPIvb_w}
\textbf{Building A Healthier Church 4: Spiritual Gifts - Rev. Mark Chan}
\newline
\newline
連結: \href{https://youtube.com/watch?v=vHt7tPIvb-w}{\texttt{https://youtube.com/watch?v=vHt7tPIvb-w}} ~~~~ 語音日期: 2023-02-20
\newline
\newline
\hyperref[sec:8B7OdFQev9A]{\small{< < < PREV SERMON < < <}}
~
\hyperref[sec:index_chronic]{\small{[返順時目]}}
~
\hyperref[sec:index_scriptual]{\small{[返順卷目]}}
~
\hyperref[sec:u_InczBESQY]{\small{> > > NEXT SERMON > > >}}
\newline
\newline
$^{1}$All right, and yeah, this is the address..
So let's welcome Pastor Mark..
- Good morning, everyone..
Good to see you all..
And today we're gonna continue our series on healthy church..
What does the Bible say about church being healthy?.
All right, so one big question in my mind.
is that maybe you should ask yourself,.
is that the longer you have been to church,.
the longer you have been the question,.
does it mean that you know what the church is all about?.
Yes or no?.
Maybe..
Maybe not, right?.
The longer you have been a Christians for years,.
sometimes we have our own concept or our own preferences.
that creep into our understanding.
of what the church is all about..
So this whole series,.
I hope that we have the desire.
to return to the Word of God,.
return to the Bible and see what God has to say.
about being a church, becoming a church..
It's not about your preference..
It's not about my preference,.
although that may be important,.
but that is not the most important thing..
The most important thing is for us.
to go back to the Word of God..
That is our final authority..
That is our highest authority,.
what the church is all about..
And as I said before, in this series,.
church is one of the most vital expressions.
of the kingdom of God..
If God comes to reign, reign in us, reign in our lives..
One of the most vital expression is the local church..
So you may have heard my story,.
how I was called to ministry, to church ministry..
Back then, I was in California years ago,.

$^{41}$and I was in my local church..
And my local church was without a pastor.
and without leadership..
And it was in the lack of the leadership,.
the lack of teaching, that the church,.
you know, suffered from time to time..
So that's how God moved me to this path of ministry..
Maybe I was praying that if God is calling me.
to the ministry, then I'm willing to answer the call,.
because I wanna be part of the movement.
to build up and strengthen the local churches..
And I will think that if the local church is strong,.
then God's name will be glorified..
On the other hand, if the local church is weak.
and sick and not very healthy,.
then God's name is gonna take a dent, right?.
That's why we need to go back to the word of God.
and see what the word of God,.
what the Bible has to say about being a church,.
especially the New Testament..
As you know, there are many, many letters.
in the New Testament that talks about how to be a church.
and what the church should look like, all right?.
So let our understanding of the church continue.
to be aligned with the word of God, the final authority..
All right, so as we do this series,.
sometimes I give you a pop quiz, all right?.
Since we have some students.
from the mighty Moody Bible Institute..
So yeah, so I'm sure you don't mind, right?.
So I prepared some pop quiz questions.
and hopefully these questions will get us thinking.
about what the church is all about in the Bible, all right?.
So first question, are you ready?.
What are your neighbors?.
All right, there are some trick questions here, all right?.
So let's warm up, okay?.
Are you ready?.
First question, what are your neighbors?.
All right, how many letters did Paul write.

$^{81}$in the New Testament?.
Four, 13, or 27?.
All right, how many, how many?.
Think about it, talk to your neighbors..
Talk to your neighbors..
Hey, do you know?.
How many letters did Paul write in the New Testament?.
All right, how many?.
Okay, how many of you say four?.
How many of you say 13?.
All right, some hands here, okay..
Looks like they're really good students..
27, how, 27?.
All right, it's not 27..
(laughing).
Well, you know why, right?.
By the process of elimination, right?.
How many books in the New Testament?.
27, there's no way Paul wrote all of them, right?.
Come on, yeah, give it to some other authors, right?.
Yeah, Paul wrote 13 letters, right?.
Okay, you can find it from the table of content.
of your Bible, all right?.
So it's actually 13, all right?.
Second question, which of the following letters.
was not written by Paul, okay?.
Galatians, Ephesians, 1 and 2 Thessalonians,.
first letter of Paul, and Titus..
Which one is not written by Paul?.
Okay, how many of you say Galatians?.
- D. - D, right?.
Wait, but it's the first letter of Paul..
You know, there's a first letter of John, you know that?.
It's written by John..
And the second letter of John is written by?.
John, right, and the third epistle of John.
is written by?.
John, right?.
So what, no, no, so first letter of Paul.
is not in the Bible, okay?.

$^{121}$It's not in the Bible, it's not in anywhere else..
So, all right, you got it?.
All right, these are easy questions, okay?.
Number three, all right?.
It's getting harder, getting harder..
Ready?.
All right, the verse, "Love is patient, love is kind,.
"love is not jealous, not boastful, not proud,".
can be found in which letter of Paul?.
How many of you say Galatians?.
It was a good saying, you know?.
It should be in the Bible, right?.
It should be in the Bible somewhere, right?.
How many of you say Galatians?.
How many of you say 2 Corinthians?.
1 Corinthians?.
All right, all right, it's 1 Corinthians, all right?.
1 Corinthians..
This is an easy question, now it's a hard one, okay?.
Now it's a hard one..
All right, ready?.
It's 1 Corinthians, the first letter written by Paul.
to the church in Corinth..
It's 1 Corinthians, the first letter written by Paul.
to the church in Corinth..
Okay, talk to your neighbors, talk to your neighbors..
What do you think?.
Yes or no?.
How many of you say yes?.
First, it's 1 Corinthians, come on, guys..
How many of you say yes?.
All right, some people are brave,.
and some of you say no..
Then what did he write?.
Why is it called first?.
It's actually no..
1 Corinthians is not the first letter.
he has ever written to the church in Corinth, no..
Why, how do we know?.
From the letter of 1 Corinthians 5, 9,.

$^{161}$Paul said, when he was writing the letter to Corinth,.
after four chapters, 5, 9, he said,.
"When I wrote to you before,".
before, so there's, I guess, Corinthians zero,.
right, beta, maybe, beta version,.
before what we call 1 Corinthians..
So there are different letters Paul has written.
to churches in his time,.
and teaching them what the church should be like, right,.
correcting them, teaching them, and all that stuff..
And what we have collect in our New Testament canon,.
the 27 books that we have in the New Testament,.
are just 27 of those..
And you know, talking about Pauline letters,.
only 13 of those..
So there are, there are maybe other different letters.
that Paul has sent to the church.
to teach them about doctrine,.
teach them about God, and teach them about church living..
So this help us to think about.
what is the purpose of church letters?.
What is the purpose of the letters,.
after letters we found in the New Testament, 27 books,.
right, many of them are called epistles, letters..
What is the purpose, and what are the purposes?.
It is actually for teaching and correcting,.
and addressing some of the issues, right?.
It's like your health report, you know?.
Anybody done a medical checkup lately?.
Huh?.
You're too young, right?.
Yeah. (audience laughing).
I got my report, and the doctor warns me,.
you know, there are some food items.
you should not put in your mouth..
Sugar, right?.
And you know, something you should not, you know,.
because, you know, my doctor said,.
you know, you, you, you, you, you,.
you know, you're not young anymore..

$^{201}$You better watch out, right?.
So it's like a, it's like a medical report, right?.
So the church letters can be seen as a report.
of the church health,.
and Paul, one of the letters' author,.
is trying to correct some of the problems,.
and you know, let's dive into the letter.
of 1 Corinthians, and you can see.
that Paul is jumping from topic to another topic.
to address some of the questions..
Now, the church letters can be called.
the occasional letters..
Occasional letters meaning, being,.
that the letter is being written.
to respond to some of the issues or questions,.
especially in 1 Corinthians..
As you can see, 7.1, regarding the questions.
you guys asked me in your letter..
Hey, so there was an other letter sent by the Corinth,.
by the Corinthian Christians to Paul in the first place,.
and Paul is responding to some of their questions..
So 7.1, 8.1, regarding the question.
about food offered to idols..
You know, in their background,.
there's religious background, pagan temple,.
and all that stuff, and you go buy meat..
The meat was already offered to the idols,.
so some of the people would not allow them.
to eat those meats because of conscience, right?.
They just don't feel comfortable, right?.
Kind of like some of us, we do have that background, right?.
I remember not too long ago,.
I received a call from a church member,.
and the church member asked me,.
because the Chinese Asian, pagan worship,.
and they have the chicken already offered to the idols,.
and then of course, after the idol is done with eating,.
and it's your turn to eat, right?.
It's your turn to eat..
And the church member asked me.

$^{241}$whether I should eat that or not..
What do you think?.
Should we eat it?.
Of course..
If it's cooked well, if it's made well,.
of course I would take it, right?.
But because there's some conscience issues, right?.
Not feeling comfortable about partaking.
into that sort of ritual, and then so she has some concern..
And these are some of the questions.
that the Corinthian Christians were asking Paul..
11.2, Paul said, "You are following the teachings.
"I pass on to you, let me expand.".
So 11, chapter 11 is all about public worship,.
and then 11.17, yeah, 11.17..
In the following instructions, Paul said,.
"I cannot praise you guys, for it sounds as if more harm.
"are being done than good in your public worship.
"when you meet together.".
So again, he's addressing some of the chaos.
that the Corinthian church was having..
Again, it's a corrective measure..
Paul is trying to correct them..
12.1, now dear brothers and sisters,.
regarding the question, your question,.
about special abilities, or what we call spiritual gifts,.
that the Spirit gives us..
So let me expand, let me teach you what it means..
So for a few chapters, 12, 13, and 14,.
Paul is trying to address this issue,.
this question of special grace gifts, or spiritual gifts..
And then 15.1, he talks about what he has preached to them,.
i.e., the doctrine, the belief of resurrection..
So chapter 15 and on, Paul is talking about the doctrine.
of Christ, you know, raising from the dead, resurrection..
So again, I want you to think with me.
that we need to return to the Word of God..
This is, at the end of the day, this is what matters..
It's not your preference or opinions.
about what church should look like,.

$^{281}$or whether I like this church,.
or like some aspect or dislike some aspect of the church..
These are not important at the end of the day..
The most important thing is the Word of God..
So let's return to the Word of God..
So today, as we go on to the series,.
let's take a look at the 12th chapter of 1 Corinthians..
So he's talking about spiritual gifts or special abilities.
the Spirit has given the church,.
spiritual gifts or special ability of the Spirit..
"I don't want you guys to misunderstand this,".
Paul said, verse two..
"You know that when you were still pagans,.
when you were worshipping the idols,.
no, the non-God, they don't exist,.
you were led astray and strapped along.
in worshipping speechless idols..
Those are nothing..
But you were worshipping the pagan idols,.
the Greco-Roman gods and deities..
Those are speechless..
They can't do anything, those idols..
But now you're different," verse three..
"So I want you to know that no one speaking.
by the Spirit of God will say, 'Jesus grants a curse.'.
And no one with the Spirit can say,.
no one can say, 'Jesus is Lord,'.
except by the Holy Spirit.".
Now, what Paul is saying is that if you're in Christ,.
if you believe the Lord Jesus Christ,.
and you're genuine, born-again Christian,.
you have accepted Jesus Christ as your Lord and Savior,.
your life has changed because the Spirit dwells within you..
Now you have the life of the Spirit,.
and the life of the Spirit will enable you.
to say, "Jesus is Lord.".
And if you will have the Spirit,.
you would not say, "Jesus grants a curse.".
Now, that's my translation from the Greek,.
because there's no verb in those two words in Greek..

$^{321}$No one speaking by the Spirit.
will say, "Jesus grants a curse.".
What does it mean?.
Now, remember, the Bible is not written to us, right?.
It's written for us..
It's written to the first century Christians..
Now, I explain a little bit the Corinthian Christians.
when before they came to Christ,.
they were in all sorts of pagan worship..
That's why Paul said, "When you were still pagans,.
you worshiped idols.".
Now, in Corinth or in any other big cities of the day,.
people worship all kinds of false religions..
And all the pagan religions have one common ground,.
one common element among all of the Greco-Roman religions.
is that people worship the idols.
because they want the idol to do something for them..
People use religions to achieve their own purposes..
So let's say you have a co-worker in your workplace,.
all right, you don't quite like the guy, all right?.
You know, in those days, in Paul's days,.
they go to the pagan temple, all right?.
You know, you talk to the pagan deity, the false deity,.
and say, "I don't like John, you know, I don't like John..
Please grant a curse, okay?.
Please grant something, do something bad, okay?.
Maybe have him fired, right?.
So I have a better day, you know, to go through.".
So this is what you do..
You go to the pagan temples.
because you use the name of the deity.
to achieve your own end, right?.
In some of the temples,.
actually we found curse tablets, you know, documents,.
you know, left behind by pagan worshippers..
And some would say, "So-and-so God, grant a curse.".
Grant a curse to whom?.
Someone you don't like, right?.
Maybe you don't like mother-in-law.
and you ask for a curse, right?.

$^{361}$It's really bad..
Maybe you don't like somebody, you know,.
or you just don't like it..
You know, you see what was going on.
in the popular religion of the day?.
Is that you use religion for your own end,.
for your own purposes..
And as you can imagine,.
some of those people from the pagan worship background,.
they come to the church and now they accept Christ, right?.
Now they probably would carry the same mentality.
and they want Jesus to do the same thing for them, right?.
Don't we all do that sometimes?.
And we need to repent..
Because we come to Jesus, we come to the true God,.
and we want this true God to be my servant,.
to do something for me, right?.
If he doesn't do something for me,.
then he must be false, right?.
This is the wrong mentality that we bring.
to the biblical faith..
And Paul is saying, "No, no, no, no, no, no, no..
"No, if you really have the spirit," Paul said,.
"you won't say, 'Jesus, grant the curse.'".
That means you don't use the name of Jesus.
for your own purposes..
You are not the Lord..
Jesus is the Lord..
You don't use Jesus for your own benefits..
Whether you want to get some more money,.
a bigger house, a higher status in your society,.
all that stuff..
Paul is correcting some of the wrong-headed mentality..
He said, "If you have the spirit, you say,.
"'Jesus is the Lord,' that means you surrender yourself..
"You surrender your will and your desire.
"to the true Lord, Jesus Christ himself,.
"and you are a servant.".
That's a totally different ballgame in the Bible,.
in biblical faith..

$^{401}$You come to worship God, to call him Lord,.
and you acknowledge his lordship..
And that's why we worship, right?.
That's why we worship..
That's why we worship every day, especially on Sunday..
We set aside this time because we want to acknowledge.
that God himself, Jesus, is the Lord, and I'm not..
He is in charge, and I'm coming before him..
I surrender myself to him..
And that's how Paul started with this whole chapter,.
talking about spiritual gifts..
He said, "There are different kinds of spiritual gifts,.
"but the same spirit is the source of them..
"There are different kinds of service,.
"but we serve the same Lord.".
God works in different ways, and it is the same God.
who does the work in all of us in the church..
A spiritual gift is given to each and every one of us.
so that we can help each other,.
or the common translation is for the common good..
Verse 11, "It is one and only spirit..
"He distributes all these gifts,.
"these special abilities to different members of the church..
"He alone decides which gifts each person should have.".
That is spiritual gifts..
And then, in the latter context,.
he talks about what kind of gifts..
So Paul kind of gives us a list..
First, the apostles, the gift of apostles,.
meaning planting churches or launching churches..
In our UV network, we belong to the Southern Baptist.
church planting movement, so we launched 2016..
And we got a lot of church planters in our network..
They are in the city, they are in different areas,.
they are starting different churches,.
and these people are doing what the apostles did..
You know, Apostle Paul, going from town to town.
and city to city to start churches,.
to begin churches, to launch church plants..
First gift is the apostle, and second gift is prophets,.

$^{441}$prophetic words..
Now again, prophetic words is not weather forecast..
It's not our modern meaning..
Prophetic, being prophetic, being a prophet,.
in Paul's days, in biblical terms,.
is speaking on behalf of the Word of God..
Based on the final revelation,.
the final authority of the Word of God,.
he proclaims the Word of God..
That's what the Old Testament or even New Testament.
in Paul's time still have prophets..
What they're doing is proclaiming.
the authoritative Word of God to the people, teaching them..
And the third gift is teachers..
Teacher, correcting and teaching.
the right doctrines of the church..
And then he lists out different kinds of gifts,.
but the whole point is, the whole point is,.
it's not an exhaustive, comprehensive list.
of spiritual gifts..
You can find different lists of spiritual gifts.
from 1 Corinthians, no, 1 Peter,.
and other church letters..
But the whole point is not giving you a comprehensive list..
The whole point is tell you what spiritual gifts are for..
The most important word in this passage.
is the word soul..
In order that..
The ultimate purpose..
What is it for?.
In order so that we can help each other or common good,.
to increase the common good of the church community..
That's what spiritual gifts are for..
And what the problems with the Corinthian Christians,.
if you read the whole letter,.
is that they regard spiritual gifts.
as something to build up their own.
high status in the community..
So the people will speak the tongue,.
and they think because we have this kind of explicit.

$^{481}$expression of the Holy Spirit's work in our lives,.
and we can speak the language of angel and all that stuff..
And that's why we are of higher status in the community..
Paul flipped it around and say,.
this is actually the last one..
You check the list of spiritual gifts,.
speaking in tongue is the last one..
The first one, or the first few spiritual gifts.
that Paul wanted to elevate,.
were actually has to do with communicating the word of God.
for the benefit of the whole church..
So Paul is trying to correct this problem and say,.
look guys, what is spiritual gift?.
And he used this special word, charismatic,.
to describe spiritual gifts..
What is it?.
Maybe the most simplest definition.
is freely given gifts, given without merit..
It's not because of who you are..
It's not because you're smarter..
That's why you get this gift and that gift..
It's not because you're better..
It's not based on your merit.
that the Spirit gives you certain kind of gifts..
It's freely given..
It's like salvation..
It's like our redemption in Christ..
It's freely given to all of us without merit..
I remember some years ago, I was reading this book.
and the author has a way to categorize,.
you know, grouping different spiritual gifts together..
And he mentioned four big categories..
The first one, he said there were different abilities,.
different spiritual gifts in the church.
that would celebrate the name of God..
Well, one of the most clear examples.
is those who are musically gifted, right?.
Right?.
Any of you have learned a music instrument before.
and you gave up, right?.

$^{521}$(audience laughing).
You do, right?.
Yeah, many of you guys are Asians, right?.
Must have learned some music talents, right?.
And then you finally give up..
You know, nothing wrong with you..
You're just not gifted..
(audience laughing).
Right?.
Maybe you're not disciplined enough..
That's why you gave up, right?.
Some of you pick up the instrument,.
they really learned well, you know?.
They didn't spend a lot of time, right?.
But they are naturally or spiritually.
given with that gift, right?.
That's one of the clearest examples of gifts.
that can celebrate the name of God..
Worship is one of them..
And the second category is communicate,.
communicating the word of God..
Some people are more gifted in teaching, right?.
Some of them are really gifted in, you know,.
doing study and clearly communicating, you know,.
the truth of God so that the church.
can get benefits out of it, right?.
Yeah, I know we have the open AI these days..
And, you know, one of my friends actually, you know,.
plop into the AI and ask the AI to,.
hey, can you write a devotional on this chapter.
or that chapter in the Bible?.
And the AI actually can spill out pretty good stuff..
You know, as you know, I'm doing studies, right?.
Old Testament studies in seminary, part-time..
And some of you guys encourage me, you know,.
because you heard, you know, I have to write a paper, right?.
So I actually try..
I actually try to, you know, plop into the AI and say,.
hey, I want this topic..
You know, can you write me something?.

$^{561}$And then I try it a couple of weeks ago..
Actually, the AI gave me some correct information.
with some wrong information..
(laughs).
With some wrong information.
about the Old Testament background..
So I correct the AI..
I said, this is wrong..
That's not from there..
You know, it's from here, right?.
So, you know, we live in the new generation.
and it's kind of interesting, you know,.
how we can use, you know, this tools, you know,.
to minister..
Although there's so much danger in it,.
but, you know, it's one of the tools, right?.
It's like computer, right?.
This morning I was talking to a brother.
who is in charge of all this IT tools of our church..
You know, you walk into church,.
there's so many things have to do with IT, right?.
You know, this thing, this thing,.
and, you know, computer and all that stuff..
Well, someone has to be gifted in also the IT stuff.
in order to make it work and make the church operate..
And some of the brothers and sisters.
are gifted in such a way, you know,.
that I'm not gifted in..
And they use it to build up the church..
The third category is caring..
Some people are really good, really good..
In hospitality..
You know, not, of course we need to love one another,.
right, no matter who, right?.
Sister or brother in Christ, we have to love one another..
But some people are really gifted in caring..
You know, go visit the sick and, you know, cooking..
Just prepare good food, right?.
You know, I'm not one of those..
You know, I cook..

$^{601}$Yeah, but the stuff that I cook, it sucks..
(audience laughing).
'Cause I have the instrumental mentality to food, you know?.
I said food, I just devour it,.
and then I'll go on to do something else, right?.
So I really don't know how to enjoy it, right?.
My son does..
Yeah, my son, Nathan, yeah..
He is a very, very good cook..
Sometimes I have to beg him..
You know, can you make a little bit more?.
Yeah, so I can have a portion, you know?.
Some people are really good at caring, you know?.
Just making a loving environment,.
creating, you know, relationships,.
so that the church can be built up..
Collaborating, leadership..
Did I mention that?.
Yeah, Paul mentioned that, the gift of leadership, right?.
People can put things together.
so that we can, you know, operate as a team, right?.
It's not like a car..
You know, you have different parts, right?.
The engine, the water pump, and the motor oil,.
the brakes, and all that stuff..
Everything has to work together.
in order for the car to bring you from A to B, right?.
You walk into the church,.
maybe when you walk into church for the first time,.
what was your impression?.
You walk into church, our church, UV..
- Cars, benches. - Yeah, okay, okay..
The cars, okay..
So you walk in, there's six or seven, you know,.
toy cars out there..
None of them is mine..
It belongs to the owner of the building..
This is toy car parking lot..
So you walk through the cars, right?.
Some people thought, you know, we are car dealers..

$^{641}$(audience laughing).
That'd be nice..
But, you know, you walk through, you know, all this stuff,.
and you come to our unit, right?.
What was your impression?.
It's a very nice place, right?.
The cafe is a very nice place..
I want to tell you the story.
because, you know, it's not me who designed it..
You know, if I design it,.
it's gonna be just folding tables, folding chairs, right?.
'Cause I'm a practical guy, right?.
When you find, you want coffee, here you go, you know?.
(audience laughing).
I've, you know, it's not in my radar, you know,.
how it looks and the design..
Because when we launched our church, Urban Voice, 2016,.
one of my core team members is a very gifted artist..
So she told me, I still remember, she told me,.
"You step away, pastor.".
(audience laughing).
You know, this sort of stuff, you know,.
it's you're not gifted, I'm gifted..
So let me design it, right?.
So she actually designed a lot of different things,.
you know, all the looks and feel, it's all from her brain..
It's all from her, right?.
It's not me..
Now, this is the conviction.
that the Bible wants you to carry..
It's that every single one of us.
needs to discover your own spiritual gifts..
You need to find out where you're gifted at..
And you need to discover it, you need to develop it..
You need to grow, right?.
I remember some time ago, I was reorganizing.
my study room, my home office..
And I found some of the old sermons that I had, you know,.
years ago, years ago..
I look at the script and I said, "Dang, it sucks..

$^{681}$"It's really bad sermon.".
You know, what was I doing, right?.
Because I was just, you know, first year in ministry.
and I have so much to learn, right?.
Well, even today, I look at my sermons,.
sometimes I told myself, "That's not gonna be good.".
Right?.
And why did I preach that, you know?.
But amazingly, you know, God used crappy sermons.
to accomplish his goal, right?.
It's not an excuse, but we need to study, right?.
We need to grow..
I need to grow and you need to grow as well..
Discover your gifts and use it, develop it, right?.
12.14, Paul continues to say,.
"Yes, the body has many different parts,.
"not just one part..
"If the foot says, 'I am not part of the body.
"'because I'm not a hand,'.
"that does not make it any less a part of the body..
"If the ear says, 'I am not part of the body.
"'because I'm not an eye,'.
"would that make it any less a part of the body?.
"No, every part is important.".
18, "But our bodies have many parts.
"and God, according to his sovereignty,.
"has put each and every part just where he wants it.".
Who's the subject here?.
God..
He has the sovereign will to do whatever he desires.
to give different members with different things, right?.
It's like a football team, right?.
Those wide receiver go to the end zone.
and got the touchdown, right?.
How did he do that?.
Of course, he cut the ball, right?.
But the ball has to pass to him, right?.
And whose credit should that be?.
You gotta have the quarterback good enough.
to pass the ball to him, right?.

$^{721}$And he's good enough to catch it, right?.
To score it..
It's like a team of people working together..
The Spirit gives us different things,.
abilities to build up the church, right?.
I'm going back to very beginning..
I was called to ministry and I answered the call..
Of course, I stumbled and fall..
I learned a lot throughout these years..
But there has to be a community working together, right?.
Now, Paul is using the analogy of human body.
to teach what church should look like, right?.
And you understand the culture of Paul's day..
Actually, the body image was often used.
by politicians in Paul's day..
There was a story in Paul's time..
The politicians were trying to appeal.
to workers who were on strike..
So there were a lot of different workers and slaves.
that were on strike because of uneven pay.
or unequal pay, treatment, and all that stuff..
So they were on strike..
So the Roman politicians comes in and say,.
"Hey, think about it..
"We are like a body, right?.
"And we, the leaders, the governing class,.
"the upper class, we are the stomach.".
And the speech actually goes like this..
"We are the stomach, we are the governing body,.
"and you workers are providing food for the stomach.".
Right?.
Using the body language to tell everyone,.
"Okay, don't go on strike..
"Go back to work..
"Because if you go back to work,.
"then we have something to eat.".
Right?.
Now, as you can tell, the analogy is working.
for those who are oppressing the people..
Right?.

$^{761}$The ruling class..
Paul is using the same image subversively..
He turn it around, turn the table..
In fact, the whole analogy,.
if you read the 12th chapter of 1 Corinthians,.
stomach is not there..
You know why?.
Because of the background of Paul..
He's not like a politician..
He's an apostle..
He wants the church to grow and to be strengthened..
And that's why Paul is saying, "No stomach.".
There's no ruling class..
Every single one plays its part..
And even the strong should value the weak, the despised..
Read the chapter on your own time..
That's how the church should be put together..
And the more I think about spiritual gifts,.
you know, sometimes we have the, you know,.
not quite right mentality..
We thought he has that ability..
I don't have that ability..
He got that tool..
I don't have that tool..
You know, it's not like that..
The more I think about it,.
spiritual gifts are more like opportunities..
Opportunities that we should seize.
and use for the glory of God..
Spiritual gifts are more like opportunity.
God has given to us..
So we talk about what are spiritual gifts.
and what are those spiritual gifts for?.
And what are they like?.
Opportunities..
God given opportunities for you and me.
to build up the church..
Now I told you last week,.
I went to my seminary for a whole week seminar.
and I present paper and some of the students present paper..

$^{801}$You know, the school that I'm going to is Midwestern..
I really like the slogan..
Can you see that?.
For the church..
Whatever we're doing,.
whatever lofty idea we're doing, you know,.
scholars, there's 30 different stuff in the Bible..
The ultimate purpose is for the church..
What a privilege we have the opportunities to do that..
Yesterday, there was a funeral,.
a memorial service for a beloved brother..
He's not coming to UV, but he goes to CCC South.
and known him for years, Dr. Dominic Ho..
And in the entire service, you know,.
different people have different reflections on Dr. Ho..
And I was really amazed by, you know, this man's humility..
You know, he's a wonderful doctor.
and is accomplished in his field,.
but it's very humble..
You know, I've known him for years and he's generous,.
he's humble, always there to help other people.
using his opportunities in his career, in the church,.
to honor the name of God..
You know, what an example..
God is giving us here and there different opportunities.
to glorify his name..
The beginning of a service,.
Miriam mentioned the revival at Ashbury University..
What started as a chapel service scheduled on February 8th,.
not too long ago, 10 days ago,.
and the service prolonged.
and different people, you know,.
they gave different expressions, right?.
Radical humility, confession,.
consecration, surrounding themselves to God..
You know, there are different people from, you know,.
distant community and travel all the way to Ashbury.
in order to experience the revival, right?.
Now, who made it?.
Who made it happen?.

$^{841}$It's not a human created thing..
It's a God sent..
Nobody planned for it, right?.
If you study the history of revival in the church history,.
nobody actually manipulate and plan for revival..
It's the spirit who sends it..
The spirit is at work in different brother and sisters.
so that the glory of God can be magnified..
And the most important thing.
is that the people are returning to God and His word..
That's the most important thing..
Yeah, the emotional, you know, experience,.
you know, those are valuable,.
but the most important thing is that we will return.
to God Himself and say, "Jesus is Lord..
"Only Jesus is Lord.".
And I'll do away all the pagan practices..
I'll do away all this non-biblical lifestyles,.
mentality, attitudes, and come to God Himself.
to experience this is the true revival..
The true revival is the people will finally.
and ultimately return to God, right?.
And let us do that..
Surrender our will to God Himself and say, "Jesus is Lord.".
Let's all pray together..
Father, we praise you..
We worship you because Jesus is our Lord,.
and He is the only true Lord..
And thank you for the opportunity we can serve you..
And thank you for the different abilities,.
different talents that we have.
so that we can use them to strengthen your community,.
your church family..
We ask that once again, your spirit will come upon us.
to open our lives and our hearts.
so that we'll be surrounding our will to you.
and be used by you, ready..
We're ready to be your servants..
We're ready to be touched and moved by your spirit,.
you know, to love and to serve and to worship you..

$^{881}$And I pray that your spirit will be at work in our midst..
Will be at work in our lives..
Help us to repent and return to you,.
to live under your word, to live out the gospel,.
live out the life of the spirit..
And I pray that you will help us to learn,.
to grow as a church community,.
that your body, the body of Christ,.
will continue to be strengthened.
and to give you glory..
In Christ, we pray..
Amen..
But again, your homework this week is to discover and to develop..
\newpage



\section{}
\label{sec:u_InczBESQY}
\textbf{同建健康教會 4:恩賜 - 陳崇基牧師}
\newline
\newline
連結: \href{https://youtube.com/watch?v=u_InczBESQY}{\texttt{https://youtube.com/watch?v=u\_InczBESQY}} ~~~~ 語音日期: 2023-02-20
\newline
\newline
\hyperref[sec:vHt7tPIvb_w]{\small{< < < PREV SERMON < < <}}
~
\hyperref[sec:index_chronic]{\small{[返順時目]}}
~
\hyperref[sec:index_scriptual]{\small{[返順卷目]}}
~
\hyperref[sec:zKpR0b97ynU]{\small{> > > NEXT SERMON > > >}}
\newline
\newline
$^{1}$中文堂同埋英文堂ge 講道呢就係Mark Sir.
加番時間比Mark Sir.
大家好 早晨.
好咁我地繼續係呢一個ge series呢叫做.
建立一個健康D ge 教會.
我地呢繼續去思想呢個問題.
有兩三個禮拜冇同大家係港台上面同你大家分享.
咁我arm arm 過去一個禮拜就去jor 神何縣.
交完paper咁我就禮拜五晚就返黎.
咁所以呢都有一D 好忙ge 時間咁.
跟住呢晏晝呢今日晏晝呢我就要去探我仔.
我仔冇見佢好耐呀唔知佢肥jor 定瘦jor .
定係冇咩點變呢咁所以呢今日晏晝都要zar 車落去.
Well anyways.
我地睇呢個ge 健康教會ge 題目.
今次係呢個講堂裡面係第二次講.
咁我地繼續睇呢個題目ge 時候呢我睇一段聖經.
係響哥林多前書第十二章.
如果你有聖經ge 話可以先打開去呢個地方.
點解我地睇呢個題目呢因為好多時呢有一個問題呢.
就係我地需要問自己.
你返jor 教會咁耐你覺得唔係返教會耐ge 話.
就知道教會應該係點架.
返教會越耐就越清楚教會應該係點.
唔一定呀係咪.
返教會幾耐上jor 幾耐.
你唔一定好清楚教會應該係點.
可能返得耐ge 更加唔清楚D .
因為我地自己對教會ge 概念.
我自己一D ge 價值觀可能都會影響我地.
點樣睇一個教會.
所以點樣都好我地都一定要返返去上帝ge 話語.
係上帝ge 話語聖經點樣去教導教會.
話俾教會聽教會應該點樣去生活.
教會點樣係一個健康ge 一個家庭.
係聖經講ge 唔係你同我去諗ge .
所以個權威始終都係響上帝ge 話語上.
唔係響我地.
唔係我地自己諗法亦都唔係我地自己ge 傳統.
更加唔係我地自己ge 意見.

$^{41}$最重要ge 就返返去聖經裡面.
到底聖經點樣去到講教會.
所以諗一諗你返到教會好耐又好.
或者唔係好多arm arm 開始返教會ge .
到底一個教會應該係點ge 呢.
最主要唔係你點諗.
係上帝點諗.
上帝ge 話語點樣去教導我地.
好咁我地開始ge 時候呢.
不如我地試下去到睇聖經裡面一D 經文之前.
咁不如我地做少少QA好唔好.
我地有個pop quiz比你.
有個測驗比大家.
準備好未.
咁我就準備jor D 講保羅書信ge pop quiz.
同你隔離gwo 個你partner問一問佢.
到底答案係乜ye 呢.
點解呢因為有好多新約ge 書信.
新約ge 一D 書卷都係講到教會應該係點樣.
好多都係出自呢個保羅ge 書信.
OK準備好未.
同你隔離gwo 個partner.
到底有幾多本書係保羅寫.
教會ge 書信有幾多卷.
OK你同你隔離gwo 個傾傾好嗎.
前後左右都得.
幾多卷.
知唔知幾多卷.
OK有邊個話係4卷.
4卷OK.
有邊個話係13卷.
13OK27呢.
你做考試ge 時候有一個技巧係做咩.
就係illumination位.
點解唔會係27.
新約已經27卷.
點可能佢27卷都佢寫佢玩曬咩.
咁所以只能夠係4或者13.
4定13.
係13架OK.

$^{81}$羅馬書 哥林多前後書.
加拉提書.
爾弗所書 肥納比書.
歌羅西書.
接受現在前後書.
提摩太前後書.
提多書同埋肥利門書等等.
加加埋埋呢13卷.
OK得唔得.
OK呀 知道答案.
有人錄低個答案.
要知架13卷.
傳統黎講我地話13卷.
呢一條.
邊一卷書以下邊一卷書唔係保羅寫架呢.
A係加亞太書.
B係爾弗所書.
C係貼哂我地而家前後書.
D係保羅一書.
E係提多書.
邊個答案 邊個唔係保羅寫架.
係咪保羅一書.
但係有約翰一書.
約翰一書係約翰寫架.
OK所以係保羅一書架.
OK呢個好清楚呀.
大家知道o下 冇保羅一書呢舊ye 架.
OK.
但係保羅寫gwo D 書信呢.
係冇歸納或者冇收係新約書信架.
知唔知呀.
我地唔係話佢13卷ge 寫jor .
但係其實佢寫多過13卷ge .
OK第三.
愛是恆久忍耐.
又要因此愛事不疾到不自誇.
不張狂 係從哪一卷書出黎架.
解太書 哥林多後書 哥林多前書.
或者爾弗所書 邊卷.
係呀哥林多後書.

$^{121}$前書 冇錯係前書.
我特登咁講 等你醒一醒先.
OK.
係啦 冇錯係哥林多前書.
OK仲有冇.
仲有.
係呢卷.
哥林多前書係咪保羅寫比哥林多教會ge 第一封書信呢.
是定否呢.
奇怪咁樣問呢D 咁ye .
哥林多前書冇理由前前書.
OK.
是定否.
哥林多前書係咪保羅寫比哥林多教會ge 第一封信.
有邊個話係架.
有邊個話唔係架.
有D 驚.
咩料呀.
唔係咩.
答案係唔係架.
點解呢 因為哥林多前書gwo 個話比你聽呢.
五章第九節佢寫一寫下ge 時候.
佢就話.
保羅講咩 佢話我之前寫信過比你地架.
I wrote a letter to you.
即係話有個哥林多前前書.
係咪.
真ge .
好過癮.
去jor 邊度gwo 本書.
唔見jor .
因為做最重要ge ye 已經收納.
係呢個新約ge 正典裡面.
所以寫多jor ge .
可能都唔係話最重要ge ye .
So that's OK.
哥林多前書五章九節講咩呢.
話我寫過信比你地ge .
之前 而家我先再寫比你.
所以我地而家有ge .

$^{161}$所謂first Corinthians.
哥林多前書.
起碼佢係第二封信.
係寫比哥林多教會.
點解需要諗呢樣ye 呢.
睇新約ge 書卷ge 時候.
所謂新約ge 書卷.
就係一D .
所謂occasional letter.
即係話一D .
去回應一D 情況.
回應一D ge 問題.
而去寫ge 書信.
好似你去寫一個電郵咁.
你係複一個電郵.
係咪 有人寫一個電郵比你先.
然後你去複一個電郵.
你複ge 當中.
你就去回答一D ye 或者回應.
某一個ge 問題.
回應某一個ge 處境.
所謂一個occasional letter.
等於你.
睇一個ge .
回覆ge email.
你冇本身個emailge .
但係你係gwo 個回覆ge email裡面.
你去嘗試去領略.
到底佢答緊D 咩問題.
回應緊咩處境.
我記得好幾日前.
我zar 車.
我老婆坐我隔離.
因為佢去一個.
所謂after school programgwo 度做.
有人寫敢唔敢有人打電話比佢.
佢就打電話比佢.
佢就回答佢.
咁就.
我聽完一大輪.

$^{201}$因為我只係聽到佢gwo 個對話裡面.
佢gwo 句.
佢收線之後我就同佢講.
我就終歸答jor .
你頭先gwo 個問題係咩問題.
我話第一呢個家長一定係好麻煩.
第二.
你gwo 個對話裡面你gwo 一部分ge 對話.
我只係聽到你一邊.
我話呢個家長應該係有D .
咁樣咁樣ge 問題.
然之後佢就氣定神壇.
咁樣同我講.
佢話其實唔係ge .
其實係咁樣咁樣咁樣.
咁佢同我分析jor 之後.
咁我起碼估中jor 一樣ye .
就係gwo 個家長好麻煩.
所以我地睇書信.
其實都係要咁樣.
所以特別睇哥林多前書.
或者保羅gwo D 書信全部都係.
try to address some of the problems.
教會裡面到底有咩問題.
佢點樣去糾正gwo D 問題.
點樣去處理同埋教導gwo D 問題.
我地比少少ge sample比大家.
你睇哥林多前書.
我知好多時我地好難one sitting睇晒.
成卷哥林多前書.
或者一卷長D ge 書信.
但係你睇下裡面.
字裡行間佢講緊咩ye 呢.
哥林多前書特別多.
譬如七章第一節.
佢話regarding the questions.
you asked.
即係話你曾經send個email比我.
而家我就答返你ge 問題.
寫jor 六章.

$^{241}$佢講jor 一大輪ye .
知道佢地ge 問題等等.
哥林多教會發生咩事.
而佢第七章佢就話.
你信上所提的事.
八章第一節.
講咩呀.
就另外一D 問題.
可能佢地信裡面.
曾經問過佢ge 問題.
祭過偶像的食物.
你地曾經係你emailgwo 度問過.
而家我回覆你.
呢封書信話比你聽點樣處理呢D 咁ge 事.
十一章第二節.
佢又跳另一個主題.
就係講到.
我要讚下你地.
我教了你咁耐.
你地凡事都記得我同埋堅守.
我曾經傳過比你地ge .
教導.
十一章第十七節.
又轉笠.
講咩呀.
我而家吩咐你地ge 話.
唔係要讚你 而係鬧你.
去責備你 因為你地敬拜ge 罪過.
係有損而無益.
佢又處理一大輪.
佢地公眾敬拜ge 問題.
十二章第一節.
弟兄姊妹 關於屬靈ge 恩賜.
我不願你們不明白.
到底乜ye 叫做spiritual gifts.
乜ye 叫做屬靈恩賜.
十二章開始.
佢就去處理呢個問題.
寫jor 好幾章.
十五章第一節.

$^{281}$弟兄們 你地要認清.
我先前傳給你們ge 福音.
唔好搞錯.
跟住佢就講復活.
呢個復活節成日會講 哥林多前書第十五章.
所以我地睇到呢D letter.
呢D 書信ge 時候.
本來其實係去.
回應緊一D ge 問題.
一D 可能.
教會裡面錯亂jor ge ye .
again我就係返返去開頭.
問大家gwo 個問題.
你唔係話你返教會耐jor .
你就知道教會係一回咩事.
唔一定 哥林多教會就係咁.
一個已經成立jor ge 教會.
但係裡面有好多ge .
混亂ge 地方或者搞唔清楚ge 地方.
所以哥林多前書或者後書.
都係寫比佢地.
黎到去處理某D ge 問題.
所以點樣都好.
我地都要返返去上帝ge 說話.
一個健康ge 教會.
一定要返返去上帝ge 說話.
去問清楚到底上帝對教會ge 要求.
係D 咩.
十二章第一節.
今日我地ge 主題ge 經文.
就係關於淑寧ge 恩賜.
各位大兄姐妹你地知道.
我想話比你聽乜ye 叫做淑寧ge 恩賜.
乜ye 叫做spiritual gift.
我不願你地不明白.
你地知道.
你地曾經做愛邦人ge 時候.
隨時被引誘.
受了迷惑去.
拜不能出聲ge 偶像.

$^{321}$所以.
我要你地知道.
被上帝ge 靈感動ge .
冇人能說耶穌.
發出奏詞.
若不是被聖靈感動的.
也沒有人能說.
耶穌.
就是主.
呢度劈頭一句.
佢就話要處理spiritual giftge 問題.
即係講唔到教會唔係好清楚乜ye 叫spiritual gift.
如果你睇字裡行間.
或者上文下理.
你知道佢地ge 教會.
佢地覺得某一D ge .
淑寧ge 表達.
某一D 淑寧ge 表現.
先至算係.
更加好.
更加淑寧ge .
一種叫做gift.
恩賜一種能力.
如果你睇上下文.
你會知道佢地覺得講方言.
就係最.
高超或者最.
好ge 現象.
如果一個人有聖靈ge 話.
佢一定要.
唔一定ge .
講方言gwo D 就最好.
所以佢就處理緊spiritual giftge 問題.
花jor 幾張聖經.
話比佢聽到底乜ye 叫做.
spiritual gift.
處理一下翻譯ge 問題.
中英文翻譯都有少少論盡.
呢度講乜ye 呢.
佢話你曾經.

$^{361}$未信主ge 時候.
你地就係拜gwo D .
不能出聲ge 偶像gwo D ge idols.
呢D 當然係當時ge .
背景甚至我地今日ge 背景.
都會係.
有好多人佢地係從一D 可能異教.
ge 風俗.
或者一D ge .
即係民間ge .
信仰gwo 種ge 背景.
而走出黎ge 而去信主ge .
咁所以佢地gwo 時係拜一D ge 偶像.
等等.
而保羅就同佢講呢D ge 偶像.
係唔能夠出聲.
係.
當時一D ge .
廟.
希臘神ge 廟或者.
羅馬神ge 廟.
佢當時就會搵到一D 呢佢地.
即係D 人去廟gwo 度拜.
gwo D ge 偶像點解呢.
因為好多時佢地會.
為jor 某一D ge 目的.
咁就會搵到一D ge 信.
或者一D ge .
碉板.
咁佢地就會叫某某ge 神明.
黎到去發出一個咒詛.
即係話.
譬如你公司有個.
人唔係你.
唔係幾順眼ge .
咁你就初一十五.
你就走去gwo 個廟gwo 度.
你就叫某某個神明.
不如發出一個咒詛.
走去燉佢冬菇又好.

$^{401}$乜ye 都好.
咁你知咩.
咁樣ge 宗教環境裡面.
就好多時佢地會覺得.
神明gwo 度做咩.
神明係聽你使冇錯.
神明其實就係.
就係你ge 工人黎ge .
發出一個咒詛等佢燉jor 冬菇.
咁我就唔使咁辛苦返工.
等等.
呢個就係gwo 種ge 背景.
咁所以如果你咁樣翻譯ge 話呢.
呢度我翻譯.
就係耶穌發出咒詛.
所以佢地咁樣ge 背景.
入jor 黎教會.
佢入jor 即係信jor 主.
信jor 耶穌之後.
佢地好多時就帶住咁ge 心態.
一種咁ge mentality.
就覺得耶穌應該幫一幫我.
所以我祈禱ge 時候就希望做咩.
就jor 咒gwo 人好.
其實就係當耶穌係乜ye .
係你ge 工人.
等於以前你ge 背景一樣.
你當邊個神明就係你ge 工人.
然後你就去做D ye 去整下佢.
你依家信jor 耶穌.
你就當耶穌又係咁樣.
所以保羅就糾正第一樣ye .
如果你真係有.
聖靈ge 話.
如果你真係信主ge .
就唔會講一種咁ge ye .
或者帶一種咁ge 心態.
Jesus grants a curse.
on somebody else.
喂.

$^{441}$耶穌可唔可以.
服事一下我.
等我達到我ge 目的.
No 保羅話唔係.
唔應該有咁ge 心態.
如果你真係信主ge 話.
你係講乜ye .
耶穌是主.
Jesus is lord.
只有耶穌先係主.
我係佢邊個.
我係佢ge 僕人.
我係比佢使用ge .
呢種ge 心態.
先至係正確ge .
You are a servant.
Jesus is lord.
搞清楚呢度.
搞清楚呢三節.
跟住佢就講屬靈ge 恩賜.
佢話恩賜有好多種.
確是同一位.
聖靈所賜ge .
事奉亦都有唔同ge 種類.
確是事奉同一位ge 主.
工作亦都有唔同ge 種類.
確是同一位ge 上帝.
在萬人當中運行萬事.
聖靈去到彰顯.
係國人ge 身上.
係教會ge 國人身上.
是要使人得益處.
這一切都係由唯一.
同一位聖靈所運行ge .
隨著自己的旨意.
留意上下文.
自己的旨意唔係我旨意.
係聖靈旨意.
去分級俾國人.
呢度就講到屬靈ge 恩賜.

$^{481}$原來係教會裡面.
聖靈就隨著.
佢自己ge 主權.
上帝照住佢自己ge 主權.
俾每一個人有唔同ge .
所謂屬靈ge 恩賜.
Spiritual gifts.
或者Grace gifts.
有唔同ge 翻譯.
或者唔同ge 能力.
甚至有D 現代D ge 英文ge 翻譯.
就會叫做Abilities.
有唔同ge 能力.
就去到構成一個教會.
而構成ge 方法.
係藉住聖靈自己ge 主權.
然後下文我唔睇咁多.
就係佢講到有唔同種類ge 恩賜.
佢就做一個列表.
列表裡面講D 乜ye 呢.
譬如有D 係做使徒.
建立教會ge .
有D 人真係去plant churches.
特別係我地UVge 網絡裡面.
有好多ge church planters.
佢地會去建立教會.
有好多ge church planters.
佢周圍去唔同ge 地方.
佢係去建立教會.
呢個第一個列表裡面有ge 恩賜.
第二個恩賜就係先知.
先知就唔係未book.
先知就唔係話聽日落唔落書.
聽日有幾多度等等.
聖經裡面講先知.
係權威ge 宣講上帝ge 華語.
係Spokeman.
Spokeperson of God.
即係上帝ge 代言人.
按照上帝ge .

$^{521}$按照上帝ge 啟示.
黎到去宗教宣講上帝ge 道.
呢個叫做先知.
第三就係教師.
教導人黎到.
上帝ge 真理.
然後有行神蹟.
醫治.
幫助人ge 都有.
有D 人特別有.
gwo 種能力去到幫助其他人.
去serve其他人.
然後佢地都講過.
有治理.
有疫力.
有D 人係教會裡面.
係有一種的才幹.
Spiritual gifts.
可以去帶領教會ge .
等等.
最後佢先講謊言.
當然保羅係特登ge .
因為教會D 人好嚮往.
佢地去講謊言.
保羅就話.
調返轉真係可以造就到教會ge .
好多時就係gwo D .
做教導呀領導呀.
去建立基督身體ge 恩賜.
當然呢D 咁ge .
屬靈恩賜ge .
清單冇一個係.
comprehensivege .
係新約ge 書信裡面唔同ge 地方.
都提過呢D 咁ge 所謂.
List of spiritual gifts.
List of talents.
List of special abilities.
冇一個list係comprehensive.
佢只係比.

$^{561}$一D ge 例子比大家.
最重要ge 就係話比你聽乜ye .
你睇完經文最重要gwo D 字係咩.
最重要gwo D 字.
上下文.
唔係列出到底我係有D 咩ge .
Abilities或者有D 咩Talents.
或者有咩Spiritual gifts.
最重要gwo 個字係邊度.
英文就係兩個字.
冇就係so.
so that.
in order to.
中文係要使人.
要使其他ge .
要使教會ge .
有D 乜ye .
益處.
common good.
舊D ge 英文翻譯就會益處.
so that.
能夠實現到.
common good.
整個ge 身體.
整個ge 社群能夠有好處.
呢D 又係希臘.
羅馬ge .
一D ge 字眼黎ge .
common good.
好多時佢地講到市政府.
或者某一個ge .
一個社群.
就需要有個common good.
保羅就借用呢D 咁ge 字眼.
黎到講.
我地作為一個教會.
就係為jor D 乜ye .
其他人.
要使其他人能夠得到益處.
所以佢反翻轉過黎.

$^{601}$因此係做咩ye .
因此唔係為自己得益.
亦都唔係為自己ge 地位.
亦都唔係純粹為自己ge 表達.
或者自己ge 經歷.
咁簡單.
最重要ge 係咩.
個方向就係in order that.
so that.
大家係能夠分享到呢個ge common good.
所以乜ye 叫spiritual gift.
俾個定義.
charismatic.
freely giving gift.
without merit.
唔係按照你ge 意識.
唔係按照你ge 好處.
唔係按照你ge .
比人地好D .
所以你得到.
呢D ge 素靈恩賜.
freely giving.
係上帝ge 聖靈.
親自按照.
佢自己ge 意思.
黎到去分比唔同ge 人.
freely giving.
邊個去決定.
神自己決定.
佢自己決定.
所以呢個冇咩好比較.
點解佢有gwo D 我冇呢.
我可唔可以有佢gwo setge gift.
普羅話唔需要.
佢特登用呢個字.
黎到去形容.
原來我地教會.
大家有唔同ge spiritual gift.
曾經我睇過一本書.
佢咁樣去分類.

$^{641}$即係當然我地有好多清單.
我地可以去處理係咪.
有D 咩素靈ge 清單等等.
有D 咩素靈ge 恩賜.
佢就咁樣分四個category.
第一個category就係celebrate.
有D 素靈恩賜係去到.
敬拜同埋去.
高舉上帝.
好簡單je 音樂.
就係其中一樣.
有D 人係有咁ge 恩賜.
有D 人你練極都冇咁ge 恩賜.
你有冇玩過音樂.
你細個有冇學過音樂.
實有你亞洲人.
實有你父母一定逼過你.
出過木頭笛都好.
但總係冇人咁叻ge .
唔知點解.
arm 呀.
唔使妒忌.
有D 人係好D .
你冇ok.
咁可能有其他恩賜.
有D 人係celebrate上帝ge 名字.
有一類就係communicate.
去傳遞.
上帝ge 真理.
教師呀.
所謂頭先講到先知呀.
等等.
呢D 都係communicate.
去到傳遞上帝ge 訊息.
第三樣啦.
就係caring.
去到維繫一個ge 社群.
gwo D 叫做幫助人ge 恩賜.
即係保羅其中一個ge example.
serving.

$^{681}$hospitality.
去到款待人.
有D 人真係特別有咁ge 恩賜.
你知唔知我都識煮ye ge .
但總係煮得唔夠人地好.
arm 呀.
即係我仔而家都.
睇先.
佢二十歲啦.
我仔二十歲.
佢煮ye 又好叻.
唔知點解佢咁有心機煮ye .
我老爸都唔係咁.
即係我老爸都唔係咁.
我老爸都唔係咁.
我食ye 只不過係一個工具性ge 食ye .
我食完跟住我就做ye .
我從來唔會enjoy.
煮ye 同埋食ye gwo 個過程.
in fact佢成日話我.
點解你食ye 咁快.
我食ye 快點解因為我要做ye .
I got something to do.
佢煮ye 好叻.
當然佢而家係大學.
我少D 見佢.
佢係屋企住ge 時候.
佢好多時朝早起身.
佢就自己煮一D ye .
咁我做老爸.
佢就唔煮我gwo 份.
咁arm 撞見佢.
不如你個part大少少.
煮埋我gwo 份.
有D 人真係煮ye 係特別叻.
有D 人呢.
就係collaborate.
或者你叫佢做領袖.
Leader.
即係點樣維繫大家整架車ge 零件.

$^{721}$點樣能夠run一齊去合作.
一齊去cooperate.
一齊去make things work together.
呢D 都係一D collaboratege .
一D ge 恩賜係唔同ge 時代上.
就係需要呢D 咁樣ge 人.
譬如我頭先同一位幫我地ITge 威大兄去傾偈.
我先發現到.
今日ge 教會要去operate.
要去運作ge 時候.
係好需要ITge tools.
ITge 電腦呀呢D 等等ge ye .
全部都係IT部門ge ye 黎ge .
呢D 都係collaborate同埋leadingge 呢一方面.
我地每一個都有某一D ge 恩賜.
有D 人可能多D 有D 人少D .
但係按照聖靈ge 主權.
黎到去分配ge .
做乜ye 呢.
for the common good.
所以每一個弟兄姊妹.
係教會係必需要.
去發現.
discover.
發現自己ge 恩賜.
我自己ge 角色.
並且要去develop.
去發展自己ge 恩賜.
每一個弟兄姊妹都係咁.
你做基督徒幾耐.
做基督徒好耐.
我一向有ge 動機.
去教會ge 動機.
除jor 敬拜等等.
我就係要係本地ge 教會裡面.
去搵到我自己ge 角色.
乜ye 角色都好.
如果唔係去教會.
我覺得冇乜人.
聖經確實係咁樣教導我地.

$^{761}$去到一個教會.
本地local churchge 教會.
我地就去到使gwo 個教會.
能夠係gwo 個自縫.
去使用我自己ge 恩賜.
arm 呀.
唔代表你.
有某一方面ge 恩賜.
你就一定.
我就ok啦.
我就唔需要成長 係需要ge .
我地都需要去到.
訓練自己去學習去develop.
係咪.
今時今日我地.
特別係我地.
呢個時代 譬如港獨.
港獨.
有D 人話.
有D 人特別有港獨恩賜.
你其實諗清楚.
可能開始ge 時候 港獨ge 時候.
佢都唔係咁.
D D 碰碰去到摸熟.
對唔對.
我記得有一次.
我屋企執ye .
搵返以前.
我一D ge 講章.
好多年前ge 講章.
我睇返ge 時候.
gwo 時用手寫.
我都唔知點寫D 咁ye .
sucks.
但你唔好.
唔好覺得過.
好奇怪 上帝又可以用.
一D gwo 時gwo 邊.
難度.
我唔記得講完之後有D 人覺得有用.

$^{801}$Well.
You know.
當然唔係要鼓勵大家.
講難度.
但我地都要啄磨.
成長需要去到學習.
今日仲有AI.
可以幫你講到.
大家有冇試過.
我試過.
gwo D 人.
有D 大研究師同我講.
你而家要讀書要寫paper.
搵D AI幫你寫.
我真係試過.
我就同佢講.
Write a paper on this topic.
咁佢出D ye .
錯ge .
我話比個AI聽.
呢個唔係gwo 度黎.
我估你個AI會.
學習一下.
到底呢個.
我地個個都係學習緊.
甚至D 所謂AI工具.
我地都需要去.
學習點樣去用.
我有個牧師朋友話.
佢試寫一個.
Devotional.
即係一個靈修ge 一個短文.
咁就叫gwo 個.
AI幫佢寫.
寫jor 出黎之後佢就放jor .
係佢ge social media.
我睇之後.
又ok喎.
真係好恐怖.
唔知點講.

$^{841}$但係我地都要識得去.
點樣使用.
但係.
當然你都唔會.
prefer AI同你講到.
你知唔知今日唔係AI同你講到.
係真人.
但係我個point就係.
我地都需要去學習.
我地每一個.
無論有邊方面ge 屬靈ge 恩賜.
我地都需要去學習.
點樣去build up the body of Christ.
for the common good.
呢個就係.
哥林多教會ge 問題.
佢認為呢D 屬靈ge 恩賜.
係比我自己抬舉我自己ge 地位.
我某方面ge 恩賜.
講方言也好.
我地就有個.
higher status.
保羅係將D ye 調返轉講.
十四節.
佢話yes.
身體.
原不只是一個ge 肢體.
而是有許多肢體.
假如腳話.
我唔係手.
所以唔屬於身體ge .
他不能因此就不屬於身體.
然後繼續講落去.
但現在.
隨住自己ge 意思.
將肢體一一.
安置在身體之上.
假如.
全都係肢體ge .
身體係邊處呢.

$^{881}$等等等.
就用身體呢一個ge image.
去到去.
表達清楚一個教會.
係點樣ge .
點解用身體呢個比喻.
因為當時.
gwo D ge 政客.
就鍾意用身體呢個比喻.
去到去.
講比市民聽叫D 市民.
係為著大家ge 大我.
為著大家ge 福利.
黎到去生活.
有一次.
一個咁ge 故事.
當時有一個城市.
D 工人罷工.
gwo D 羅馬ge 官員.
就出去為jor 平息.
罷工gwo 個strike.
咁佢就用jor 一個身體ge 比喻.
佢話我地就好似一個身體咁.
我地呢D 所謂ruling class.
我地呢D 咁ge .
管理ge 官員.
我地都係希望你.
唔好罷工繼續做ye .
因為我地就係gwo 個位.
你地就係gwo D 手腳.
你地手腳要繼續做ye .
以至我地ge 位能夠郁到.
你聽到呢D 比喻ge 時候.
你地知道佢講緊乜ye .
唔好罷工點解.
因為你服侍我.
保羅.
講身體比喻.
就剛剛相反.
每一個ge 部門.

$^{921}$都係重要.
唔係gwo D 領袖就係個胃.
Stomach.
in fact佢冇提到心脈.
反而係一D .
所謂缺欠ge .
所謂軟弱ge .
gwo D 我地.
更加需要去抬舉佢地.
去欣賞佢地.
呢個先至係基督教會.
存在ge 目的.
你今日ge 功課.
就係去搵你ge .
屬靈恩賜.
同埋去發展同埋去使用.
所以我地講jor 幾樣.
乜ye 係屬靈恩賜.
屬靈恩賜係為D 乜ye 存在.
for the common good.
為jor 教會.
為jor 其他人ge 著想.
恩賜似乜ye .
如果你咁諗ge 時候.
恩賜其實似機會.
opportunities.
唔係咁多係講緊.
到底我個能力.
qualification.
我ge 才幹.
係講緊opportunities.
講緊時機.
上帝比我地.
生命裡面唔同ge 季節.
就有唔同ge 時機.
gwo 個時機就要我地去把握.
gwo 個所謂spiritual opportunity.
spiritual opportunity.
spiritual gift 黎到去事奉.
最後兩個例子.

$^{961}$呢個我個神學院.
唔係D 咩名校.
但我幾鍾意佢一個ge .
即係口號.
for the church.
為jor 教會.
而存在.
我地去讀呢D 咁ge ye .
或者去study某一D ye .
課堂呀等等.
最終目的是做咩.
為jor 教會.
而存在.
呢個就係我地有ge 時機.
我地ge opportunity.
arm arm 上完堂gwo 個禮拜.
隔離一個後生仔.
即係成個class裡面.
我係最老gwo 個.
Gabriel.
即係佢廿零歲.
咁大家傾開計認識下ge 時候.
我發現.
咩叫時機.
佢同我講.
佢而家arm arm 開始又讀返.
我話點解停jor 呀你.
佢話因為呢.
年幾前呢.
佢有一個車禍.
一個好嚴重ge 車禍.
就令到佢有brain damage.
好大ge brain damage.
咁佢係ICU好長ge 時間.
然之後先至開始recover返.
我話你而家係咪fully recover.
佢話我仲未fully recover.
所以佢睇電腦.
佢係咁樣睇ge .
好辛苦呀.

$^{1001}$仲要睇gwo D 希伯來文.
gwo D 一點.
chicken scratch.
我話點睇呀.
我先realize.
我地大家坐係度.
能夠去study.
能夠去所謂發展我地ge 恩賜.
係時機.
opportunity.
咁我地就.
opportunity.
唔係個個都有咁ge opportunity.
我地就需要把握呢個時機.
黎到去試防治.
昨日有個安息禮拜.
何醫生.
安息禮拜.
係成南堂好多年.
一個好忠心ge 一個大兄.
我好感動.
佢個故事就係.
一個好humblege 人.
佢從來唔張揚ge .
唔張揚佢有幾叻ge .
我一聽佢ge 心平行坐.
原來佢咁犀利.
不過我地唔知je .
原來佢係gwo 行裡面.
佢係咁厲害ge .
一個醫生.
咁受人地去尊敬ge .
一個醫生.
但一個咁humblege 人.
就係用佢自己ge .
opportunity 黎到去.
試防治.
去謹守上帝俾我地ge 崗位.
讓我地做好我地ge 功課.
搵下你自己ge .

$^{1041}$屬靈ge 恩賜.
或者搵下你屬靈ge .
opportunity.
你點樣可以去.
貢獻去建立.
基督教.
一個健康ge 教會就係我地每一個.
照住上帝ge 主權.
俾我地ge 恩賜我地去使用.
可以同心得道.
主要多謝你俾我地.
呢個教會.
係教會當中我地知道.
聖靈係呢度ge 工作.
俾到我地每一個.
姐妹每一個弟兄.
我地有唔同ge .
才幹或者更加清楚ge .
就係有唔同ge 機會.
主要俾我地用我地ge 恩賜.
去建立你自己ge 教會.
for the common good.
讓你自己係我地中間.
你ge 名能夠被.
高舉.
去使用我地去發展我地.
使我地能夠去建立你自己ge 身體.
你自己名得榮耀.
祈禱奉主名救.
\newpage



\section{}
\label{sec:zKpR0b97ynU}
\textbf{比對研究 8: 巴力宗教 1 - 陳崇基牧師}
\newline
\newline
連結: \href{https://youtube.com/watch?v=zKpR0b97ynU}{\texttt{https://youtube.com/watch?v=zKpR0b97ynU}} ~~~~ 語音日期: 2023-02-24
\newline
\newline
\hyperref[sec:u_InczBESQY]{\small{< < < PREV SERMON < < <}}
~
\hyperref[sec:index_chronic]{\small{[返順時目]}}
~
\hyperref[sec:index_scriptual]{\small{[返順卷目]}}
~
\hyperref[sec:9YLRCR_mn34]{\small{> > > NEXT SERMON > > >}}
\newline
\newline
$^{1}$我們看巴黎宗教的背景.
我們以前看過比較研究.
曾經看過一些聖經背景的東西.
正如我所說.
看到聖經的背景.
其實可以幫助我們如何領略聖經的信息.
聖經對當時的人有什麼信息.
而我們領略了之後.
對我們今時今日有什麼當代的應用.
就會更加容易運用得到.
所以我們看看一個很出名的假神.
叫做巴黎.
在巴黎的背景下.
到底聖經如何去對抗.
或者那些所屬耶和華的人.
敬拜耶和華的人.
真心事奉耶和華的人.
他們如何去對抗這些巴黎的假神的影響.
這個情況.
其實就好像我們今天要面對不同的.
所謂今天的巴黎.
當然今天已經沒有人再說所謂拜巴黎.
但它所代表的就是一些.
嘗試去毀滅耶和華的僕人.
嘗試去毀滅耶和華的信仰.
毀滅聖經的信仰.
這些文化的意識等等.
或者這些勢力.
如何去對抗.
我們知道聖經是有它的背景.
很多時候我們就將聖經直接.
想應用在現代的處境裡面.
但我們說過.
如果聖經是這樣的話.
我們就將聖經抽離了.
就沒有了聖經本身的世界.
聖經第一代的讀者.
古代的讀者.
如何去理解那些經文.
我們首先需要去明白了.

$^{41}$這個絕對不是像這裡那樣.
希望這個圖畫可以幫助到大家.
不是說將聖經框在以往古代的世界.
而不能夠應用在我們現代的世界.
而是我們認識了以往近東的世界的時候.
我們知道聖經如何跟它對話和如何對抗.
以致我們能夠影響到那個意義.
而那個意義就能夠應用在我們今時今日的世界.
所以很簡單想介紹一下.
巴力是一回什麼事.
我們經常看到先知書或者以利亞的故事.
說到如何去跟巴力的先知去鬥法.
跟巴力去抗衡等等.
到底是怎麼來的呢?.
就跟大家去一個地方叫做烏加利.
烏加利就是今天Rashomwa這個地方.
就是今天在土耳其和敘利亞邊境的地方.
這個地圖就讓大家看到.
以色列就在南邊.
耶路撒冷這個地方.
而北邊一直藍色的海旁.
就是當時最大的海港.
在地中海東的世界.
Mediterranean East.
這個世界裡面的霸主.
貿易最大的海港的霸主.
就是這些海港.
譬如有推羅,泰爾,西頓,比布斯,比布拉斯.
上面烏加利這個地方.
烏加利這個地方在1920年的時候.
就找到很多跟巴力宗教有關的背景.
烏加利就是一個海港.
在烏加利的上面大概二三十公里左右.
就是這個地方叫做Mount Saffong.
Mount Saffong就是藍色的海岸.
叫做菲利基.
海岸一直伸延下去.
由北邊Mount Saffong一直伸延下去加密山這個地方.
這裡整片都是拜巴力的大本營.
所以你們記得以利亞的故事.

$^{81}$就是阿哈王跟誰結親.
就是跟西頓王的女兒夜西別結親.
所以巴力才知道和以利亞去到抗衡鬥法的地方.
也是在加密山這個地方.
所以這個地方在1920年的時候.
不單找到土優也找到很多文獻.
在這裡發掘之後.
就是古烏加利城.
有的城牆和裡面不同的建築物.
比如有什麼建築物呢.
這裡有一個很大的皇宮的地方.
就是西世代烏加利王.
在這裡作王.
是他們的宮殿.
而且也找到很多archives.
就是文件和例版.
他們用一些攝影文字寫的例版.
有一些他們的信件等等.
也有他們守護的城門.
有住的地方.
右上角這裡就是有Bell Temple.
這個所謂的巴力的大廟.
裡面也有一些祭司的房屋等等.
也有他們的大祭司的文獻寫下來.
當然多神論的世界裡面不止一個廟.
和耶路撒冷不同.
也有大滾神明的廟等等.
找到這些東西.
最重要的就是找到一些烏加利的例版.
正如我所說.
在攝影文字刻上去的例版.
燒乾了之後.
就會有一些文字在上面.
這個烏加利的例版.
我們為什麼需要知道呢.
因為你知道在舊約聖經裡面沒有提到有烏加利這個地方.
但是很多時候提到巴力.
巴力才知道.
巴力宗教.
巴力神明.

$^{121}$巴力和亞瑟拉等等.
這些多神論的背景.
在烏加利所找到的文獻.
就幫助我們去理解到底巴力是什麼一回事.
為什麼對當時的以色列人來說是一個這麼大的試探.
正如我所說.
這個平衡可以看到我們今天現時的世界.
到底有什麼是屬於所謂的巴力.
對我們在文化裡面.
思想上面是經常在操控我們.
經常在影響我們.
這些就是今天的巴力.
過去的巴力是什麼呢.
這些烏加利的例版.
有幾千塊左右.
發掘出來之後.
翻譯了之後.
就發現有三類這些文件.
第一類的文件就是宗教的文獻.
宗教文獻就是說.
講述他們的神話.
講述他們宗教的信仰.
他們所信的神明.
一些故事等等.
譬如有些什麼呢.
我這裡列出的就是一個叫做Bell Cycle.
就是巴力的循環故事.
Bell Cycle.
另外一個就是一些皇的史詩.
就是Quarter 或者 Arkhart.
這些史詩都是講到可能跟聖經裡面一些的寫作.
是有相似的地方.
譬如探討苦難等等.
當然這些宗教的文獻.
就是三類裡面是最多的.
當然在這個宗教的文獻裡面.
最出名.
最多人關注的就是第一個了.
就是Bell Cycle.
巴力的循環故事.

$^{161}$等等我們會看看.
或者之後我們會看多一點.
到底巴力在神話裡面怎麼形容他.
以至為什麼.
住在迦林地裡面的以色列人.
這麼容易受巴力宗教的影響.
這個第一類是Religious Text.
第二類就是Apostolary.
就是書信.
烏加列這個城市或者城邦.
是大概當時主前1500至1200年左右.
到1200年左右的時候.
就是被毀滅的.
所以有幾百年的時間.
有很多的皇在歷史歷代在管治這個海港.
這個大的城邦.
他們也寫了很多的信件.
譬如跟埃及王有一些互通的信件.
或者官員,人民寫了一些信函等等.
這些給我們看到他們當時的生活.
這是第二方面.
第三方面就是Administrative Text.
就是一些行政的文件.
這些行政文件有時候有一些明治的清單.
或者是合約等等.
所以找到這些很豐富的內容.
我們嘗試用這些內容.
給我們一些亮光去理解聖經的背景.
甚至幫助我們理解一些聖經的經文.
巴力的循環就是最著名的.
有時候我們叫這個故事.
在六塊的泥板上講述了巴力的故事.
有時候叫做巴力的循環故事.
Bale Cycle.
有時候叫做巴力神話.
Bale Myth.
或者有時候叫做Bale Anarch Myth.
Anarch就是巴力的妹妹.
或者又是他的情人.
他們叫做兩夫妻.

$^{201}$他們有一些神話的故事跟他們有關等等.
在這六塊泥板上.
Tablet 1-6.
記載了一個相當完整的故事.
講述巴力如何在眾神裡面崛起.
然後他最後贏得了王權.
成為至高者.
成為真正的God of Gods.
就是眾神裡面的神.
所以我們需要理解一下.
當時的人如何理解巴力.
以至我們可以理解到聖經裡面提到巴力的時候.
耶和華跟巴力的鬥爭上.
到底為什麼會出現這樣的故事.
或者這樣的記載.
你看看地圖.
你就知道烏加利亞就在最北邊.
以色列就在南邊.
耶路撒冷這個地方.
當時分了南北角.
阿哈的時候.
北角就是以色列.
下面就是猶大國.
但是整塊土地叫做大迦南地.
Greater Canaan這個地方.
之前我們提過這個地方是一個中間之地.
所以在文化上一定會有互通.
或者互相交流.
互相有互動.
所以看烏加利亞講巴力的故事.
這些泥板一定可以幫助我們去理解.
以色列或者猶大國.
以色列在迦南地這個地方.
他們所面對的一些文化的衝擊.
巴力給他們的衝擊.
一定是有用的.
並且我們知道烏加利文.
剛才我們所講的泥板等等.
烏加利文是屬於西面閃族的語言.
West Semitic.

$^{241}$這一堆.
敘利亞和巴勒斯坦地.
很多不同的語言.
包括希伯來文.
都是屬於West Semitic.
就是西閃族語.
西閃族語是什麼意思呢.
大概最簡單來說.
西閃族語就是.
很多時候他們的字根.
都是用三個字母來組成.
所以如果讀過一點點希伯來文就知道了.
大部分希伯來文的生字.
動詞.
都是用三個字母.
三個consonant來組成的.
這個就是西閃族語的特色.
並且我們知道在西閃族語裡面.
烏加利文的文字.
和希伯來文是非常相似的.
有很多的字眼.
甚至是一模一樣的.
只不過寫法是不同的.
他們就是用楔形文字.
然後後期的希伯來文.
就是用一些square script.
就是正方的字母.
所以這個對比.
是一定有好處的.
一定可以幫助我們理解.
聖經怎麼去對抗當時的文化.
\newpage



\section{}
\label{sec:9YLRCR_mn34}
\textbf{Art, Beauty, and God: Gen 1-2 - Rebecca Wu}
\newline
\newline
連結: \href{https://youtube.com/watch?v=9YLRCR-mn34}{\texttt{https://youtube.com/watch?v=9YLRCR-mn34}} ~~~~ 語音日期: 2023-02-28
\newline
\newline
\hyperref[sec:zKpR0b97ynU]{\small{< < < PREV SERMON < < <}}
~
\hyperref[sec:index_chronic]{\small{[返順時目]}}
~
\hyperref[sec:index_scriptual]{\small{[返順卷目]}}
~
\hyperref[sec:5SRk2JnEBBc]{\small{> > > NEXT SERMON > > >}}
\newline
\newline
$^{1}$Hi, everybody. Good to see you. Good morning. It's nice to be back at UV. Many of you may.
be here or maybe, maybe, maybe not here. I don't know. No, maybe not here. But I used.
to do these videos for UV during the height of the pandemic. And it's just like a one.
woman variety show where it's me and the camera. I'm just like, Hi, everybody. I'm Mr. Rebecca.
de la Homa. And I'll say it all in Cantonese. And I'll throw in some English in there. But.
today I'm going to be talking to mostly adults. So we're not going to play any games. But.
we're going to have some few fun interactive activities. And I probably will be less hyper,.
but just as interesting. All right. So this morning, we'll be exploring a topic that's.
really near, dear to my heart. As you can see on the slide, it's about art, beauty,.
and God. And I want to start by asking everybody a question. Gosh, it's so techie here. Last.
time I was here, I wasn't so techie. So what value do you place in art and beauty? So in.
other words, how important is it to you, like personally? And how important is it to society?.
It's a rhetorical question. So you can think about it. So I'll share with you what I often.
hear, especially from the Chinese culture that I grew up in. So art is way impractical,.
like not very useful. Ancillary. What? The painting is \$3,000. Why so expensive? No way,.
I can do that, close my eyes. Paint, paint, paint, all done. So, and beauty, of course,.
is nice, but who has the time for all that? We don't have time to make everything pretty,.
everything to look good. But if it's for skin care, we can do it. So the funny thing is.
that despite all these negative comments that I've heard or even have absorbed in our Chinese.
culture that I grew up in, and in culture in general, and society in general, we spent.
a lot of money and time and resources in artists, like BTS, Taylor Swift. I mean, they are artists..
Practicing their art form, performing on stage for tickets, sold out shows, multiple shows..
We spend money on other people's art, other artists that are celebrities. So the word.
celebrity means we celebrate them, right? But we don't really actually celebrate art,.
we just celebrate them. And especially here in the States, we have game shows and reality.
shows, especially for things like this, like America's Got Talent. But have you noticed.
when anybody goes on stage on America's Got Talent, it's like the same kind of talent?.
So we have someone's gonna sing, definitely gonna sing, maybe with a guitar, maybe with.
a piano, maybe both. Someone's gonna dance, someone's gonna do magic and wow the crowds..
So, what I'm seeing is they're doing these things, but you don't see someone stand up.
on stage and perform their talent in doing a math problem. How quickly can I solve this.
puzzle? Can I do a Rubik's Cube in less than a minute? Or how good am I at organizing spreadsheets?.
Those are not talents we celebrate, those are not talents we wanna compete for, but.
actually we do that, a lot of us, every day. That's what we practice, that's part of our.
work. So, we as the audience love it when other people display art and beauty, right?.
We love it. So, if we're gonna spend time to watch these shows and watch the next YouTube.
video and the next YouTube video for the next two hours of our lives and not sleep, or buy.
tickets to watch these singers, it would seem that we highly value the arts. We highly value.
art and beauty. But I don't know if we actually encourage that in ourselves or for others..
I wonder why. And for me, this topic is really near and dear because I am an artist. I make.

$^{41}$my own artwork, I lead art making sessions and workshops, I work with a community, I.
exhibit their art, and I'm an author. These are all artistic, creative things that I do..
And it's a huge part of who I am. But I am no superstar celebrity. So, does it still.
matter that I do this stuff? Is it worth it? So, I've struggled a lot to be able to put.
what you see on the slide here in terms of identifying myself as an artist. It's way.
easier to say I'm a teacher. Way easier when I was a business consultant to say, "Oh, I'm.
in consulting. All the aunties love it. Make so much money. Yeah, very good." In college,.
before college, I wanted to go to art school. No way is that happening. But I can do everything.
art on the side if I wanted, as long as it stayed a hobby. So, even after being a consultant,.
being an educator, et cetera, what do I come back to? I still come back to my passion..
And I can't deny that I feel my best when I make art and when I teach art. But even.
as I put these pictures up on this slide and talk about being an artist, it's still hard.
to reconcile these messages that I've heard either from growing up or messages from society.
and balance that with what I feel is very important and valuable. And kind of combating.
with what I am hearing in society saying, "It's really not that useful. It's shallow..
It's superficial. It's ancillary." So, kind of selfishly, not too selfishly, but also.
like altruistically, I want to explore what the Bible says about art and beauty, about.
this topic. So, essentially, I want to know what God thinks about this in hopes that I'll.
find affirmation in what I'm doing. Because I'll get affirmation about being an accountant,.
no problem. I'll get affirmation about being a doctor, a lawyer, no problem. All the affirmation.
in the world. And if I don't get affirmation from God, maybe he doesn't care about art..
Maybe he doesn't care about that stuff that I care about. Then maybe as I explore, I'll.
find another direction. So, today, I want to explore with you what I've discovered so.
far. So, we're going to look at art and beauty as a theme. I know, I said I'm a teacher,.
right? So, I got to teach you these things in this kind of way so that all learners are.
included, all diverse learners in all these ways are included. So, you'll see graphs,.
you'll see pictures, I'll be talking. And that's what I do. So, we're going to go through.
the Bible. And when I look at themes, it's a theme because we see it across the Bible..
And the Bible is one narrative, not little stories picked out and put in a collection.
like fables. It has a beginning and it has an end. So, we'll look at creation. We'll.
look at the flood. Exodus, when God's people were taken out of Egypt as slaves. We'll look.
at the Psalms. And we'll look at parts of the New Testament. So, we're going to start.
from the beginning, the very beginning. So, in creation, the first thing we read is what?.
In the beginning, God created the heavens and the earth. Now, the earth was formless.
and empty. Darkness was over the surface of the deep. And the Spirit of God was hovering.
over the waters. And God said, "Let there be light." And there was light. God saw that.
the light was good. And he separated the light from the darkness. God called the light day.
and the darkness night. And there was evening and there was morning, the first day. So,.
what I notice immediately here is that God is the original artist. He created something.
from nothing. So, if we look at art, you have a blank canvas and you put something on it..

$^{81}$You look at a sculpture, there's stone, there's clay. It was a lump and the artist makes something.
out of it. Taylor Swift went from no song to a song. So, something was formless and.
was given form. Something was empty and it was being filled just by his word. There's.
a vision and there's an order and there's a plan for his creation. And he filled it.
layer by layer. In Genesis later, it says, "And God said, 'Let there be a vault between.
the waters to separate water from water.' So, God made that vault and separate the water.
under the vault from the water above it. And it was so. God called the vault sky. And there.
was evening and there was morning, the second day. And God said, 'Let the water under the.
sky be gathered into one place and that the dry ground appear.' And it was so. God called.
the dry ground land and the gathered waters he called seas. And God saw that it was good.".
So when we paint, whether it's on a canvas or a piece of wood, whatever it is, when we.
paint, we start out with something empty. And we have to paint from the background to.
the foreground, layer by layer. And you have to let it dry. If not, all the layers start.
combining and mixing. It gets all mixed up. And you can't tell each layer apart. So God.
first paints this background of the heavens and the earth. Then he adds the water, the.
sky, the land, plants, vegetation, seed bearing fruits, the flora. Then he hung luminaries,.
lights, the sun, the moon, the stars. More layers. The sky needs birds. All the fall..
The sea needs creatures. And the land, the livestock, everything creeping and crawling,.
all the wild animals. Everything is there. And this is the beautiful painting God created.
with every detail from every vein of a leaf, every seed of a fruit, every crater in the.
moon, every feather in a bird, every fin on a fish, every spot, pattern on any creature,.
any animal. And I think we can all say, if we look at God's creation in nature, wow..
What a design and what artistry. There's form, there's function, and there's finesse. So.
even in today's world, scientists, designers, and artists try to mimic what God created.
just by his word. Airplanes mimic birds. They call it biomimicry. So we have creations mimicking.
each other. Camouflage, where's that from? Things that scuba divers wear, where's that.
from? And we try to mimic it or even capture it with photographs, but we really can't capture.
all its majesty. So up until this point in creation, the earth is being filled and filled.
and filled as the backdrop for the last creation, humans. It's the only creation that God created.
in his own image. The scriptures say, "Then God said, 'Let us make mankind in our image.
and our likeness, so that they may rule over the fish in the sea and the birds in the sky,.
and the livestock and all the wild animals and all the creatures that move along the.
ground.'" So God created man in his own image. In the image of God, he created them. Male.
and female, he created them. And the Lord formed a man from the dust of the ground and.
breathed into his nostrils the breath of life. And the man became a living being. I think.
we've established pretty much that God is an artist, is a designer, is a creator here..
And every artist makes their mark. There's some sort of mark when you create. When you.
make the art, you mark it. This is Picasso's signature. We all recognize it as his mark..
Not just his style, but his mark. Fancy Le Bouton shoes, underneath the red, ooh, a mark..
This is a fake pair. Mid-century modern furniture, my favorite. What's the mark? The angled legs,.

$^{121}$the streamlined lines. That's a mark. God left a mark on all of his creation. And the.
biggest mark he made was on man. So only man was sculpted and molded from the dust of the.
ground. And the only creation given the breath of life. His own breath of life. So as humans,.
we're created in God's image. In his image. So I have kids. They're created in the image.
of my husband. They both are boys, and they both look like him and his brother. And people.
love to tell me that. Wow, he looks just like your husband. You are so correct. But I know.
they are mine. Because I birthed them out. And in certain angles and in certain lighting,.
with a certain smile, they sort of resemble me like 2\%. So when they're created in our.
image, is it just what they look like? And of course we love that. We'll bring them..
Oh, look, this is my eyes. No. Oh, look, it's my nose. Also not my nose. So what I tell.
people is, yes, created in our image looks just like my husband and his brother. But.
all the insides are me. All that in there, me. So if we're created in God's image, we.
must share characteristics of our Father. And from what we just read in Genesis, in.
the beginning of the Bible, in the very beginning, God is a creator. God is an artist. So we.
must possess creativity and artistry. But really, what's the point of having that? We're.
going to go say something and make some flowers appear? I don't have that power. So I don't.
have exactly what God has, but an image of it. But what's the goal of it all? Why? He.
recreated everything, didn't he? What's the point? So as I keep exploring, we've laid.
out the foundation in creation, right? So maybe as we look at this really pretty picture.
from up in the mountains looking, this creation is for enjoyment. It's so we can delight in.
God's creation, so God can delight in his own creation. That must be it. But what happened.
after creation? Adam and Eve were sent out of the garden. And then after generations,.
God's people were being corrupt and wicked and evil. God thought they were so evil, he.
regretted he made them. I never think that about my boys. Usually. He wanted to wipe.
them out with the flood. I just want to send them to the corner. But there's this one person.
he wanted to save. And that's Noah. Only Noah was righteous in the eyes of God. So he asked.
Noah to do something. Hey, Noah, I'm gonna wipe out the whole world with lots of water.
like a flood, but I'm gonna save you and your family. And this is what you're gonna have.
to do. Build a big boat. Like a really huge boat. So in scripture, he tells him exactly.
like what to do. Make it out of cypress wood, put rooms in it, coat it with pitch inside.
and out. Who's ever built a boat before? I've been in a boat. I've never had to make the.
boat that keeps me alive. Oh, and then it has to be 300 cubits long, 50 cubits wide,.
30 cubits high. And make sure you have a roof on it. And below the roof, make sure there's.
an opening also one cubit high all around. Put a door on the side too. Make sure there's.
a lower, middle and upper deck. Oh, and after you're done with this big boat, put two of.
all living creatures in it, 'cause we're gonna save the moms and dads of all the animals..
Make sure you have every kind of food too, 'cause that's all you're gonna have to eat,.
you and your family. Whoa, dude. It's like getting a verbal blueprint of how to make.
this boat, but having no blueprint. So this picture of like, oh, there's A and there's.
A prime, there's B and B prime. And say, oh, don't forget C prime, prime, prime. And then.

$^{161}$this, I wouldn't know what to do. And I'm already good at math. I won't know what to.
do. No power tools. The scriptures say Noah did everything just as God commanded him..
Okay. So I don't know about you, but I wouldn't be able to know what to do about that. Can't.
even ask my spouse for help or my sons or anything, just me. But God knew, right? When.
he found Noah and asked him to build a boat, he knew he has craftsmanship. He knew he'd.
understand. So how does God use his talents? 'Cause this verbal blueprint is what designers.
do, is what architects do, is what contractors do, what builders do. Ultimately, Noah uses.
craftsmanship to obey the Lord out of obedience. So maybe that's like an objective of all this.
art and beauty stuff. To obey, to follow. Now, further on after the flood happened,.
let's look a little further for more examples. God's people were once enslaved by the Pharaoh..
And in Exodus, there is a different kind of ark that needs to be built. As God's people.
were saved by Moses, parting the Red Sea, going into this land, they were away from.
their country, they were away from a place. And God asked Moses, "Build a tabernacle.".
Tabernacle is a place where they can worship, a tent. It's very orderly. So if you look.
at the layout of a tabernacle, it's very symmetrical. There's lots of order. Who gets to be where.
and why. But inside this tent is an ark. Sort of like this. This isn't like the actual one,.
which is a model of it. But it's made of metal. It looks pretty intricate to me. You have.
to like know a lot of different things. And God said to Moses, "Oh, okay. I want you to.
make a tabernacle, a sanctuary, so I can dwell in it amongst my people. And make all the.
furnishings inside it, exactly like the pattern I'm about to give you." So pattern, design,.
blueprint. Patterns are what a lot of tailors and seamstresses use. You have to make it.
out of this wood, this high. Inside and out has to be pure gold. Put four rings for it..
Fasten it to four feet. Two rings on one side, two rings on the other. Oh, but don't forget.
the poles. Overlay it with gold. Insert the poles into the rings. Make sure you make an.
atonement cover of pure gold. Make sure it's this long and this wide. Wow. Man. Because.
if I were Moses, I'll be like, "What? How long is a cubit?" 17.5 inches long. How do.
you even measure a cubit? How do I make sure it's definitely that long, that wide, and.
that high? But don't worry. God has an answer. So in Exodus, we also read later, like a few.
chapters later, after all this direction, "Oh, I've already chosen people to take care.
of this. Don't you worry." God chose specifically, it says, "See, I have chosen Bezalel, son.
of Uri, the son of Hur, and the tribe of Judah. I filled him with the spirit of God, with.
wisdom, with understanding, with knowledge of all kinds of skills, to make artistic designs.
for work in gold, silver, and bronze, to cut and set stones, work in wood, and to engage.
in all kinds of crafts." He's like the Renaissance man. Like he can do everything. This is a.
true skilled artisan. Oh, and don't worry. I also appointed Aholiab, son of, I don't.
know how to pronounce this, Ahisamach, I don't know. But he's from the tribe of Dan, to help.
him. So Mr. Renaissance man isn't enough to make this sanctuary for the Lord. We need.
more skilled artisans. And don't worry, I've given ability to all skilled workers to make.
everything that I've commanded you. Wow. When's the last time you've heard of maybe a church.
or a place, an organization say, "Top priority, go get the artists." Number one, make sure.

$^{201}$we hire the artists and make sure we train them. Make sure they're skilled. Make sure.
they have all the support. I don't think that happens. Usually we go, "Make sure we hire.
the accountant. Make sure we have the CEO. Make sure we have the janitor." But that's.
not what's happening here. God is handpicking artisans out of everybody for very important.
work to build the tabernacle where he's going to dwell. The highest honor. But oftentimes.
these artisans, at least nowadays, are the outliers. When I say the word artisan, what.
do you think of? Bread. Art is no bread. Very good. I can't bake it. But it takes a long.
time to do it. I think of the really pretty baskets with the leather handles. Like, "Ooh,.
really pretty, really nice." I don't know how to weave it either. I don't even know.
how to sew my own clothes. I don't know how to do those things. I went to a wood shop.
at the park district. It took me three months to build a toy shelf for my child. Three months..
And then in two weeks, they were sick of it. I should have gone to IKEA. But I can't even.
build the IKEA shelf either. So God's saying, "I'm just going to pick these skilled artisans.
because they're my people of choice." And he treats them with priority and honor. Because.
this level of beauty, this level of art, this level of skilled craftsmanship is a level.
of excellence God deserves, wants, and requires. So art and beauty in this case isn't just.
for worship because the tabernacle is a place for worship. It's a form of worship. Because.
we want things to be excellent, high quality, and beautiful for God. And he even says that.
these people are blessed with this talent. Not everybody got it. And they've worked hard.
to develop this skill. Pure talent doesn't mean you can create. You need to practice.
and work hard. So these people are supported by everyone else in the community. These people.
were not picked. But I'm pretty sure someone's feeding these artists. So in this congregation,.
artisans were far from being ancillary. They were essential. The essential workers, basically..
Essential for worship. And I know nowadays in modern day church, we might not see worship.
in this way. Like I did not help Pastor Mark build this stage. I did not. I made the PowerPoint,.
but that's what I did. I didn't make the beautiful cityscape back there that Pastor Mark loves.
to build and create. When we think of the term worship, we usually think of music. So.
the songs that we sang, the music that we play on instruments, the DJ, which is pretty.
much all of Psalms in the middle of your Bible. That's what I was taught. I didn't grow up.
in the church, but when I went to my first Bible study, I couldn't figure out the order.
of all the books. And every time the teacher said, "Psalms, yes! Just open to the middle.
of the Bible, and I'll find it." And I always wondered why does it have to be in the middle..
But it kind of makes sense. Because, well, these are songs, right? And in the Bible,.
we see when music is used, when instruments are used, it's for times of jubilee, celebration,.
the trumpets of Jericho falling down the walls. All you got to do, make some noise with your.
trumpets. Wow. But also in Psalms, in some of them, if you see, like in the beginning.
of Psalm 5, it tells you what instrument to use. Because Psalms are for music directors..
It says, "Oh, use flutes with these. It'll be best." Wow. But what I want to focus on.
in Psalms is, yes, the melody instruments are nice, but really on the words, the lyrics..
And I don't think we hear the lyrics or know the lyrics or absorb the lyrics as well, unless.

$^{241}$we're pretty familiar with the tune already. Like the songs we sang today, this morning,.
I knew one of them and I had to learn the other two. And I love it when I can sing along..
That brings me a lot of joy. But when I already know the tune, when I already know the music,.
I can't focus on the words. Because not every song is about joy and victory and happiness.
and fa-la-la, tra-la-la. A lot of songs, even the Psalms, the lyrics are about sadness,.
about fear, about anger. But everything is honest, authentic from the heart. So when.
words are written, either in prose or in poetry, accompanied by music or not, I think the objective.
of it is to share, to express. And worship music and worship songs are like a staple.
in modern day church service. So it's always in the beginning, it's always at the end when.
you respond to the teaching, the message. At least every Sunday service I've ever gone.
to. But why? Because the words serve to communicate and connect with what you just learned and.
what you just heard from the scripture. And in planning, you call it a response song after.
the sermon. Because it helps us respond to the Word of God. Not just praise Him, but.
respond like a conversation between us and God. So recall a very moving sermon I listened.
to one time, had to do with like our wickedness. It wasn't by Pastor Mark, but he also gives.
good sermons. But this one particularly talked about our wickedness, our wicked hearts, and.
all of us needed God's mercy. It was like cutting deep. And the response song after.
the sermon was the tried and true Amazing Grace. And even non-regular church goers might.
know the tune. Amazing grace, how sweet the sound. And then we got to these words, because.
I already knew, I knew how to sing it. I knew the tune. Got to these words, we got to the.
words, I was lost, but now I'm found, was blind, now I see, and then tears are flowing..
I couldn't even control myself. And I felt it. And as much as I know it was the Holy.
Spirit resonating in me, it also reminded me of back in the day, because I'm kind of.
old, when you have to call the radio to pick a song. Radios have DJs, so it was not on.
demand. You had to wait for your song to be played on air, or else you can't hear it,.
unless you happen to have the CD or the cassette tape. And it never fails when you ask the.
DJ for a song, probably for your crush, for your boyfriend, whatever, or maybe someone.
had called in about a heartbreak or something, the DJ picks a song in response to that situation..
And every time I listen to the radio, when that happens, it never fails to articulate.
those feelings. And when you hear your feelings being reflected in music, in words, in a song,.
it resonates. It comforts. It knows. At least to me, I know that I'm not alone. And that.
resonance is very powerful. All right. So, so far in the Bible, we talked about art being.
used to delight, to enjoy, to obey, to worship, communicate, connect. But art's art, right?.
Art could be kind of like anything. So what makes art beautiful? Because beauty seems.
sort of subjective, doesn't it? Like this watercolor looks beautiful to me. Is it the.
most beautiful piece of art I've ever seen? Is it the most beautiful piece of art you've.
ever seen? Would you not call it beautiful, but maybe just pretty? That's nice. That's.
interesting. Oh, I'm sure they really tried. Looks like a blob. What is it trying to say?.
I mean, the colors are pretty. Oh, it's very vibrant. How do we exactly define beauty?.
What's beautiful? I mean, in Exodus, when God was telling Moses what to build for the.

$^{281}$tabernacle, oh, there was purple, there was blue, there was colors, there were linens,.
there were tassels, there was gold. Wow. All the finest materials, down to the color, down.
to the fabric, it was defined. That kind of definition was beauty. And maybe sometimes.
we have certain rules that we follow. Maybe not too as specific as the parameters for.
the ark and the tabernacle, but we follow some sort of rules. So certain styles are.
trendy while others are classic. One color is in, like now, like mauve is really in,.
like beige is really in, grays, muted tones. But there was a period of time when neon was.
very in, like in the '80s. But maybe that's out now, or maybe it's coming back. So we.
judge beauty often within certain parameters. And in mathematics, there's something called.
the golden ratio. And mathematicians have discovered it in architecture, artwork, flowers,.
and nature. And that has to do with the proportion of the length and the width of something..
And usually that proportion ends up being about 1.6. Because these things are nice to.
look at when the length and the width and the proportions are about 1.6. So we have.
a photograph a little bit in the foreground and a lot in the background. So essentially.
we really like looking at rectangles. And we think they're really pretty to the eye..
So for a person's face, that ratio means the length is about 1.5 times the width. So more.
like a rectangle. I did not measure my face today. But some of us have longer faces, some.
of us have rounder faces, some of us have heart-shaped faces, some of us have broader.
faces. Not beautiful? More beautiful? Some people have said, "Well, you have a very symmetrical.
face if you put a line down the middle of your nose. If one side reflects perfectly.
on the other side, like Angelina Jolie. Beautiful. But if not, like one eye is bigger than the.
other, which I have. Not so pretty. So then in my pictures I have to angle. I want to.
show you, look at these photos of all these people. Different standards of beauty across.
cultures, norms, countries, generations. Who's most beautiful? Sometimes, in some places,.
lighter skin is beautiful. Sometimes darker skin, straighter hair, curly hair, coily hair..
What if we go to different cultures? What if we don't show our skin? What if we paint.
our skin? We'll put markings on our skin, piercings on our skin. Is that still beautiful?.
I mean, the list goes on and on. Big eyes are more beautiful than small eyes. Snow white.
porcelain skin is really hit in Asia. So we have whitening creams, whitening masks, whitening.
this, whitening that. But you take that same Asian person and put them in a different context,.
tanner skin might be more in. Big eyes, small eyes, double eyelids, single eyelids, thick.
lips, thin lips, thick eyebrows, thin eyebrows. The list goes on and on and on and changes.
perpetually, time periods. At one point, fair and plump was really beautiful. And now it's.
skinny and dark. Sometimes it's voluptuous, like Beyonce. Sometimes it's not. And that's.
just for females. We haven't even talked about men yet. With this perpetually changing trend.
and beauty standards, what standards really stay true and can stand the test of time?.
So with all those different pictures, I think if I asked you who's the most beautiful, you.
might each pick a different face. Because it's based on preference. And beauty is in.
the eye of the beholder, isn't it? So if you were the beholder, you would pick what you.
prefer. But the ultimate beholder is God, is Jesus. And Jesus sees the beauty within..

$^{321}$In the New Testament, in the Gospels, we see Emmanuel, which means God with us. Jesus,.
God with us. Dwelling with us. So he basically kind of reset what's worthy, who's worthy,.
and where is worthy. He didn't go to the fancy places to see the fancy people. He didn't.
go to Michelin star restaurants. He went to where it was ugly and rejected. He met with,.
hung out with, cared for the dirtiest people. The prostitute who was criticized and judged.
for her profession, Jesus looked at her heart of faith. The tax collector seen as greedy,.
a liar, a cheat. Jesus looked at his heart for repentance and reverence for him. Even.
worship requirements kind of changed and got reset in the New Testament. In the Old Testament,.
we needed an ark to be like this. The tabernacle to be set up this way. We needed this, this.
high, this long, this wide, this color, that color. But those specific standards aren't.
so strict anymore because of Jesus. When Jesus died on the cross, there was like this earthquake..
In the Bible, it talks about a curtain tearing. The veil was torn from top to bottom, not.
bottom to top. And that veil was what separated people from God. Only the most holy and the.
most clean. And if you were to go in there, you'd die because you're not clean enough..
Jesus changed all that. His death changed all that. We don't need these specific strict.
guidelines of what to craft, what to make, even though it's beautiful. But the focus.
is on the beauty of salvation and redemption instead. So things on the inside. So we have.
tons of examples, I think, throughout the Gospels on how inner beauty was emphasized..
But actually, we see that God was after our hearts from the beginning, after all. So that.
theme is actually carried, not just in the New Testament where it seems like things have.
changed, there's a turning point, but it runs in both Testaments. So who knows what this.
is about? Yes, my son is raising his hand. Yes, David and Goliath. Thank you for the.
applause. Little David beats big giant Goliath. We know the story. My boys love this story..
And I remind them that there's more to this story because the person of David is a key.
character in the whole Bible. David was a little guy who beat the big giant with the.
tiny stone. Classic hero. But not without help. I always ask my boys, how did David.
do that? Was he the strongest? Yep, there he is shaking his head. Was he the tallest?.
Nope. How did David do it when no one else dared? Yes, Ethan. Yes, because God was on.
his side. Their answers are so pure, aren't they? And confident because God was with him..
Because God. Yes. And also because David. David believed and had faith. Because all.
the other people, yeah, God was still there. But they didn't believe. They didn't have.
a heart of faith. Okay, now this is the victorious picture here. Just rewind a little bit, just.
a teeny bit. Before this happened, David was actually anointed king by Samuel. Because.
Saul was not going to be king and Samuel had to go to David's house. David's dad was Jesse..
And David had tons of brothers, older brothers. Stronger, taller, maybe more handsome. So.
Samuel had to bypass all those brothers because he thought, oh, that one is probably king..
That one's called tall and strong. He's pretty handsome too. But God's like, nope, not him..
Keep going. Next brother? Nope, not him either. Man, none of them. Do you have any other sons?.
Yeah, we got like that. David, he's like out there somewhere taking care of the sheep..
Go bring him in. For real? David? All right. That's him. Wasn't even what Samuel thought.

$^{361}$either, let alone Jesse. Because David didn't fit the bill of what we have an image of a.
king. But God chose David for his heart. And look what he did in the battle with Goliath..
And even when David became king, he committed a great sin. He took someone else's wife..
I tell my boys, that's a big no-no. We do not take someone else's wife or someone else's.
husband. And then King David made plans for her husband to die in battle. That's like.
big no-no number two. And even amidst that grave sin, horrible sin, God looked at David's.
heart. So in God's eyes, the heart is huge, which puts a different perspective on beauty..
Because it's in the eye of the beholder, right? In David's case, I might look at him and go,.
"Oh my goodness, I can't believe he did that. I can't believe he did that with Bathsheba..
I cannot believe it." And God still sees his heart. Yeah. But isn't it true that God's.
the ultimate beholder? He's really the only one. So I don't know about you, but I can.
be kind of shallow when it comes to looking at couples and celebrity couples. I like to.
see pretty people with other pretty people. And when that doesn't happen, in my eyes of.
course, I tend to criticize. Like, "Oh, she chose him." Or, "Oh, he chose her." I know,.
I'm shallow. I admit it. Yep. Like Hugh Jackman. Okay. So as superficial as it may sound, I'm.
speaking the truth of what I think and what I believe, right? Just based on outward appearances..
But I fail to see in these celebrity couples is their inner beauty, which I would never.
be able to see, because all of my impression is based on their public image. So they may.
look really good, but you know what? Maybe they have the golden ratio face. I don't know..
But once you hear rumors of scandal, that pretty face is no longer pretty. Their beauty.
immediately plummets. But on the contrary, when you hear of a celebrity's loyalty, adopting.
dogs, rescuing babies, whoa, their beauty points skyrocket. They might not even have.
like a prettiest face, and it don't matter. All of a sudden, they're beautiful. Why does.
that happen, though? Do you think that inner beauty is visible on the outside? Does it.
change the person's face? Are things like kindness, humility, attractive? So attractive.
that it's more attractive and alluring than looks? So me and my husband have been married.
for a long time. And as I talk about being shallow about celebrity couples, I think about.
God looking at the heart, and then I have to ask myself, well, do I? Sometimes, sometimes.
not. In premarital counseling, the pastor asked us to read Peter, specifically the part.
that talks about wives, for me. First Peter, chapter 3, verses 3 to 5. "Your beauty should.
not come from outward adornment, such as elaborate hairstyles and the wearing of gold jewelry.
or fine clothes. Rather, it should be that of your inner self, the unfading beauty of.
a gentle and quiet spirit, which is of great worth in God's sight. For this is the way.
the holy women of the past, who put their hope in God, used to adorn themselves." I.
thought it was really funny and ironic, because all I could think about was my dress, my hair,.
my makeup, my bouquet, and all my accessories. But I get what the pastor was trying to say..
And, all joking aside, even though my nails were really important for that day, I understood.
what we're reading is important for all the days after the wedding day. And as time passes,.
even though I think my husband is super cute, it was not the only thing that attracted me.
to him. And as time passes, we're almost at 15 years, and I've noticed since I was a kid,.

$^{401}$we've aged. And our looks have changed. Even though Asian don't raisin and beige don't.
age, I look younger than my age, but I have aged. And my looks have definitely changed.
from when we were dating, and especially from our wedding day, when I had lots of people.
helping me look a certain way. So as that change has occurred, it's really those character.
traits that we have to be attracted to. Although sometimes those same characteristics can really.
drive us crazy, it's really what continually is refined in our marriage and really carries.
us through. And I have to agree that that beauty from within is indeed precious and.
unfading. So we've kind of hopped through the whole Bible a bit to explore the theme.
of art and beauty in the scriptures, and the theme seems to be everywhere, shown in different.
ways. And all of this exploring has really shown me how much of a master artist God is..
So in the art world, when you talk about master artists, people say Rembrandt, Monet, Dali,.
Picasso. And I don't know if you've done a lot of art, maybe you have a practice of.
making art, or you know artists, but for me and some artists, mostly many of the artists.
that I know, we can spend days, weeks, or months on a piece, whether it's a painting,.
a picture, editing a photograph, designing something, a poster, whatever it is. When.
we step back and look at it, we appreciate what we made. It's not bad, it's pretty good..
Keep fixing it. And when it's finally done, all done, maybe I'll put it in a frame, maybe.
I'll put it up somewhere, I celebrate the beauty of the artwork. And I call it good,.
like very good. I go rest, refresh, and realign. The next piece of art. To go back to creation,.
God saw everything he had made. And indeed, it was very good. So the evening and the morning.
were the sixth day. Thus the heavens and the earth, and all the host of them were finished..
And on the seventh day, God ended his work which he had done. And he rested on the seventh.
day from all his work which he had done. And God blessed the seventh day and sanctified.
it. Because in it he rested from all his work which God had created and made. So unless.
we're commissioned to do a specific piece, like a portrait or draw my old house or something.
like that, artists don't really need their art pieces to be beautiful to you or anyone.
else. Because it's meaningful and beautiful to themselves. That's why they do it. So when.
I go to the Art Institute and go, I gotta go see the Monet exhibit, so beautiful. Gotta.
go to Paris, go to the Louvre. When my husband and I went for our 10th anniversary to Paris,.
I've never been, went to the Louvre. And there was a huge line for Mona Lisa. But the painting.
is actually quite small. Like in the textbooks, like it feels like it could be really big.
because you only see it on the page. You get there, it's like, I'm like 12th row from like,.
where is it? It's like, just like a regular sized portrait. And everyone's staring at.
it, talking about it. I'm not really sure if the artist cares what we think about Mona.
Lisa. But we really care about talking about it. So if I were to put like examples from.
artists like Picasso, Michelangelo, Da Vinci, Van Gogh, Monet, Cicel, Andy Warhol, abstract.
painters, installations that have one piece of stone on a pedestal and you walk away,.
new installations. I think we would all agree that it's all art. And objectively, all impeccable..
But regardless, I don't think we would agree on which one's the best, or the most beautiful..
And the artist frankly doesn't care. The maker of each piece of artwork didn't create it.

$^{441}$for you or for me, but for himself. So the maker already knows the purpose and the value.
of every piece that anyone else outside would never be able to clearly see. The maker is.
the beholder of the beauty and the creative process it took to make that holds a lot of.
feelings and stories of inspiration behind every piece. Therein lies the beauty, the.
inner beauty that only the artist knows very well. God is our maker, our creator, our artist,.
we are his creation. We're his art piece. In the Bible there's this analogy of the potter.
and the clay and Jeremiah. Who's ever done pottery class here? Like two. Who's ever seen.
someone do pottery before? Oh, like four, five. Who's ever watched a YouTube video of.
pot? Oh, okay, maybe 10. So it's messy. It's messy. Here's what happened in Jeremiah. The.
word came to Jeremiah from the Lord saying, "Hey, arise and go down to the potter's house.
and there I'll cause you to hear my words." So then I went to the potter's house and there.
he was making something at a wheel. The vessel that he made of clay was marred in the hand.
of the potter. So he made it again into another vessel and it seems good to the potter to.
make. So when you start making pottery, you throw clay, it kind of looks like this. You.
have a big clump, you want to make it wet, you put it on the wheel and it spins. Oh,.
you better start holding that clay. You have to put pressure on it. Not enough pressure,.
everywhere. Too much pressure, clump. Your hands shape it with pressure, with angling.
and sensitivity. Why was he told to go see the potter and the clay? Then the word came.
to him after he saw this. So if you can imagine, this is what Jeremiah is seeing. "Oh, house.
of Israel, can I not do with you as this potter?" says the Lord. "Look at the clay is in the.
potter's hand, so are you in my hand, oh house of Israel." So very experienced potters can.
do this. Make cups, make vases, many shapes, many heights, no problem. But should it happen,.
like a scripture is saying, it's marred, you messed up, darn it. You can still make corrections,.
you can remold the clay, a different shape, a different vessel. Until the potter chooses.
to glaze and fire, complete creative control over what the final outcome would be. So to.
me, that means we, being the clay, are continually being molded by the potter, being God. So.
we have to trust his hands to make us into vessels he's pleased with. Because the clay.
can't really talk back. Can't be like, why'd you make me like this? I was really intrigued.
by this pottery metaphor and this analogy, so I found in Romans 9, Paul reminds the people,.
puts them in their place, like you know your clay, right? And he says, "Does not the potter.
have the right to make out of the same lump of clay some pottery for special purposes.
and some for common use?" So I have no idea how God is going to mold me or mold you, and.
mold us, and I'm not sure what type of vessel you would be or I would be. I definitely don't.
know fully how to define common or special use, especially in this context. But what.
I do know is that God is an artist who makes beautiful things. So from his hands, some.
of us were given the talents of artistry, and we need to use them in obedience, to worship,.
praise, and connect. Maybe you don't consider yourself an artist, a designer, a maker, a.
painter, a sculptor, a musician, a singer, an actor, a poet, whatever you deem an artist.
to be. But do know that you were created as an art piece, and you were created by an artist,.
by design, which means you were created to be creative. And we were all created to admire.

$^{481}$and nurture beauty, both inside and out. So I encourage you to pause more often, step.
back and appreciate not just God's artistry, but also your own. Notice what is around you..
Take time to value art, value beauty, and value yourself. Because we often spend a lot.
of time nurturing quantitative, measurable things. Like I'm on a diet right now, because.
Thanksgiving, Christmas, Chinese New Year really hit me bad. But how do I quantify my.
progress? Inches, pounds, I wish it weren't pounds, it would be just like milligrams,.
that would be way more success and achievement there. But we spend a lot of time doing those.
things. How many hours did I sleep? How many hours did I work out? How many hours did I.
work? Blah, blah, blah. Quantitative, measurable. But since art and beauty are so highly valued.
by God, how about we try to nurture that part of us too? So to be corny, WWJD. Who knows.
what that means? No one has ever seen that on a bracelet before. So when I was growing.
up in church, which I didn't grow up in church, people were wearing these bracelets that say.
WWJD. What would Jesus do? T-shirts, hats, bracelets, everywhere. I was like, what the.
heck is that? What would Jesus do? I don't know, die? I don't know. Why are you wearing.
this? But really it's to remind us in every case, what would Jesus do? And in our theme.
and topic today, I believe Jesus will look for the beauty in the ashes, beauty in the.
everyday things, and the beauty in everybody, the beauty in people. Because after all, that's.
what artists do. So I started off this morning wondering what value God placed on art and.
beauty in hopes to affirm myself as an artist. And I think I found what I was looking for,.
and I hope you found something too. Thank you..
Thank you very much. Thank you very much for insightful sermon. And I want to pray together..
Can we? Yeah. Let's all pray. Father, we give you thanks because you are the master artist..
Thank you for creating this universe in such beauty. Thank you for the sunrise. Thank you.
for all the things that we can see and touch. And thank you so much for making us in your.
image. And thank you for your reminder that we can be creative and we are made to be so..
And I pray that we will worship you in this way to appreciate what you have created. And.
most of all, appreciate what you do and who you are. Thank you for being our beautiful.
God and our mighty God in Christ. Thank you so much. Thank you so much for the insightful..
We should invite her to come back, right? Yeah. And such a great perspective. And by.
the way, I didn't make that cityscape. Irene did. Yeah. So praise God for that. And Willie.
did. We really appreciate the presence of artists. And as I told you a couple of weeks.
ago, people walk into this space, come to our cafe, they appreciate the beauty of it..
And it's not my job. I'm no artist. And because of some of the artists in our team, in our.
church who have created this beautiful space, right? Let's do so. Let's put to work and.
be more creative and appreciate God's creation. All right? So today we have the fellowship.
in the afternoon, right? And we're going to start. So you're invited to join us. All right..
So I want to see you next Sunday and God bless you and you're dismissed..
[BLANK AUDIO].
\newpage



\section{}
\label{sec:5SRk2JnEBBc}
\textbf{同建健康教會 5: 愛與恩賜 - 陳崇基牧師}
\newline
\newline
連結: \href{https://youtube.com/watch?v=5SRk2JnEBBc}{\texttt{https://youtube.com/watch?v=5SRk2JnEBBc}} ~~~~ 語音日期: 2023-02-28
\newline
\newline
\hyperref[sec:9YLRCR_mn34]{\small{< < < PREV SERMON < < <}}
~
\hyperref[sec:index_chronic]{\small{[返順時目]}}
~
\hyperref[sec:index_scriptual]{\small{[返順卷目]}}
~
\hyperref[sec:CFdO0Ddf8oE]{\small{> > > NEXT SERMON > > >}}
\newline
\newline
$^{1}$時間交比Mark Sir.
唔該曬啦Cindy.
好我地再睇呢一個ge topic.
我地一路都係講緊.
健康ge 教會.
第一個問題想問大家.
一個健康ge 教會當然就係.
一班健康ge 基督徒黎到組成.
一個健康ge 基督徒.
應該係點架.
好似medical report咁樣.
最近有冇睇醫生.
有呀.
我就去jor check up.
醫生就同我講你小心D .
有D ye 唔等你入口.
你就唔好放入口.
OK.
佢就話你都唔曬啦.
應該留意下你自己.
跟住你又訓得唔好.
加加埋埋好多ye 你就玩完.
佢就比個report比我.
OK.
有D ye 係需要承認.
我自己都有D ge 生活上面.
一D ge .
或者dietary上面.
一D ge adjustment.
食ye 少jor 呀等等.
食ye 少jor 唔係因為窮.
食ye 少jor 因為有D ye 係真係唔落得口ge .
咁就算啦.
咁我就唔食咁多gwo D 我鍾意食ge ye .
一個medical report.
聖經比我地.
教會基督徒.
都好多時有一個medical report.
去check check.
到底你係咪健康先.

$^{41}$check check到底你係咪healthy.
我記得係呢個series裡面問過呢個問題.
係咪信得耶穌來ge .
就知道.
好明白.
一個健康ge 教會.
一個健康ge .
受靈心靈.
就應該係點ge .
係咪信得耐就知ge .
唔係呀嘛.
可能係初信ge .
你更加清楚.
just shame on us.
有時係信得耐ge .
反而唔係好知道.
問一問.
健康ge 基督徒係點.
一個健康ge 教會應該係點ge .
ok 健康ge 教會應該係點呀.
舉個例.
一家好ge 餐館應該係點呀.
係咪個menu乜ye 都有呀.
係咪有成千樣menu.
你入去揀ge 時候.
唔知點揀ge .
有冇去過D 餐館呀.
百幾個menu item.
你知點揀.
揀jor 成15分鐘都未揀到.
又唔得定架喎.
唔得定gwo 個ge 餐館.
樣樣ye 都有齊.
包羅萬有ge .
有一間好ge 餐館.
另外一個例子.
你話你家.
你屋企.
係咪樣ye 都有呀你屋企.
但我問你呀.

$^{81}$如果你最近你屋企想買D 咩.
想到想買D 咩.
爭D 咩呀你屋企.
爭屋.
爭D 咩電器呀.
有D 需要用ge 用具呀.
有冇呀.
成日都有架係咪.
買幾多買唔完.
係咪你家屋企.
乜ye 都有.
就係一個健康ge 快樂ge 家.
英文好有意思呀.
英文係.
我記得gwo 陣時我地.
婚前輔導ge 時候我同我太太.
gwo 個counselor同我地講.
You are not building a house.
You are building a home.
House同Home有咩分別.
Home係講relationship.
House係講咩.
硬件.
有D 咩爐頭呀.
個廚房幾靚呀雲石做呀.
地板幾靚呀.
木ge 地板.
冷氣幾好呀.
但係裡面.
係點ge .
就睇你個home係點.
睇你個relationship係點.
其實教會都係一樣.
教會都係一樣.
呢個topic值得我地去諗下.
乜ye 係一間.
合上帝心意ge 教會.
again.
唔係將我地ge preference.
放係頭先.

$^{121}$可能人地問到你.
一個好ge 教會係點樣ge .
你話樣樣ye 都有囉.
有兒童士工.
有培基士工有小學士工.
有中學士工有大學士工.
青少年士工等等所有ye .
全部都有晒.
係咪一定係好ge 教會.
again we are not building a house.
we are building a home.
right.
again唔係我地ge preference.
係最重要.
係上帝ge 話語.
呢個先係最高ge 權威.
所以點解我地一定要返返去.
上帝ge 說話.
上一次我地睇jor .
係哥林多前書第十二章.
gwo 度講到話每一個.
係我地在座ge 只要我地.
係基督裡面ge 如果你係未信ge .
你係仲係摸索緊ge .
除外.
如果你已經信jor 主ge 話.
無論你信jor 幾耐.
聖靈呢按照佢ge 主權.
係比jor 你.
恩賜.
比你恩賜.
ge 原因同埋個目的.
就係為jor 造就教會.
係咪.
哥林多前書第十二章.
本來就係處理緊.
spiritual giftsge 問題.
每一個人都有spiritual gifts.
咁你地點樣去到使到呢D .
spiritual gifts去造就教會.

$^{161}$當然哥林多教會可能.
係新約教會裡面.
新約裡面眾多ge 教會裡面.
最豐盛ge 一間教會.
資源.
資材最豐盛ge 應該.
因為係哥林多前書第一章.
開始ge 時候.
本來開頭ge 時候通常呢佢就會.
話為教會係咩祈禱.
佢都已經講jor 句咁ge .
說話佢話因為.
你地呢.
在神裡面在主裡面.
凡事富足就.
具有各種口才.
各種知識.
ge .
好唔好呢間教會.
最sharpge 人都係gwo 度.
社會裡面.
教會裡面所有.
有口才ge 人有知識.
ge 人.
都係.
教會裡面.
但係呢個教會.
又係最混亂.
因為裡面.
有好多ge 人佢地追求呢D ge .
才幹同埋.
恩慈.
就覺得自己係十分熟練.
覺得我比其他人.
洞見更加好.
我經歷呢就比較深.
其中一個ge 恩慈.
大家知道就係講方言ge 恩慈.
我地唔係gwo 種ge .
傳統我地唔係gwo 種.

$^{201}$即係第一節會講方言gwo 種ge .
傳統我地可能未必會直接.
去即係處理到.
佢裡面即係.
處理講方言ge 問題.
但保羅提醒jor 乜ye 呢.
最主要個中心.
每一間教會都可以聽ge .
或者每一間教會都可以應用ge .
個原則就係乜.
佢具有超凡ge 能力.
OK 比人多.
知才比人多.
resources.
你地要知道保羅話係出自上帝.
係聖靈.
所賜比你ge 按照佢ge 主權.
按照佢ge 自由.
你都可以賜比你.
以致你能夠.
使用係教會裡面.
因此目的.
就係為jor the common good.
保羅就係十二章gwo 度講.
用身體.
黎做一個比方.
十二章最後.
我地繼續睇落去.
二十七節.
希盟士基督ge 身體係各自.
都士之體.
神係教會裡面所設立ge .
第一係使徒.
第二係prophet.
先至.
第三係教師.
其次ge 係行義能ge .
再次ge 係.
醫病ge 恩賜.
幫助人ge .

$^{241}$治理人ge .
最後係說方言ge .
佢特登將.
佢地最渴求ge gwo 樣ye .
最能夠表現到gwo 樣ye .
就係講方言.
佢就放係最尾ge .
唔係話gwo 樣ye 冇價值.
而係保羅話咩.
恩賜ge 目的.
係要去造就其他人.
去使其他人.
有common good.
能夠去分享得到.
如果係冇ge 話呢.
gwo 個恩賜可能就要諗一諗.
點樣去發揮出黎.
保羅就係話你有ge 恩賜.
就應當要去發展.
去運用ge .
最重要.
就係上帝要求你.
係要去使用.
你ge 恩賜.
去祝福其他人.
十二章.
最尾ge 時候.
睇落去就到第十三章.
十二十三章.
十二章第三節.
佢講完哂D 恩賜ge 時候.
佢話你地要追求更大ge 恩賜.
佢話哥林多教會.
你有好多恩賜.
好呀.
你去追求ge 恩賜好.
不過你爭jor 一樣ye .
就要追求.
最大ge 恩賜.
The best gift.

$^{281}$最大ge .
屬靈ge 恩賜係乜ye 呢.
你地要去追求.
就係我而家將呢個.
最妙ge 道去指示你.
十三章就係我地最熟悉ge .
一段經文應該.
佢話我若能夠說.
以前和合本就係.
萬人ge 方言係咪.
和修版.
和合本ge 修訂版.
就係我若能說人間ge .
方言甚至.
天使ge 語言.
有冇聽過天使語言.
應該有.
有冇.
通常我都聽jor 返唔到黎.
有D 人都係咁.
有D 咁ge 經歷.
near death experience.
聽到天使好美妙ge 聲音.
如果.
我能夠講到.
唱到呢D 咁ge 說話.
呢D 咁ge 壞話.
卻沒有愛.
我就成為明ge 羅.
響ge 拔一半.
如果我能夠有先知講道ge 能力.
可以明白.
所有ge 傲備.
各樣ge 知識有齊備ge 信心.
使我能夠.
而生添呢個信心.
實在係.
卻沒有愛算不得什麼.
如果將所有ge 財產.
黎到去救濟窮人ge .

$^{321}$又犧牲自己ge 身體.
讓人地去誇讚ge .
呢個翻譯比較好D .
犧牲自己.
使別人能夠誇獎我.
卻沒有愛.
與我無益.
卻沒有愛三次出現.
講jor 幾個ge .
example.
應該呢D 道理都好清楚.
唔需要點解釋.
一個教會.
有幾多ye 都好.
或者你要求幾多ye 都好.
卻沒有愛.
算不得什麼.
is nothing.
again係上帝ge 眼中.
唔係我地眼中.
我地可能覺得樣樣ye .
包羅萬有真係好.
樣樣ye 都.
tap tap掂.
樣樣ye 都有齊哂.
太好了.
但係聖經講咩.
卻沒有愛.
算不得什麼.
上帝往往.
睇ge ye 就同我地係唔同.
教會.
係過去幾年我相信我地都.
都好勞累.
我地作為同工.
包括好多.
好忠心去.
事奉ge 弟兄姊妹.
我地都會gui ge .
我都會gui .

$^{361}$因為捱jor 幾年.
我好想係度欣賞一D .
係我地台前幕後.
有好多ge 同工.
唔能夠一日個個都講.
但我真係好appreciate.
因為佢地一路都冇quit過.
係我地可以.
安坐係沙發.
或者訓係床上面.
我地去睇下youtube.
或者甚至唔睇.
因為佢地.
佢地一路冇quit過.
重啟之後.
我地都係.
不眠不休.
我地都係搏命咁做.
我地仲有好多弟兄姊妹.
係做外展添.
每次.
去到做一D 外展工作.
派下食物.
去接觸下一D 街坊.
一D 朋友.
真係好唔話得.
重啟之後.
你一冇人黎.
佢都照開.
arm 呀.
你可以選擇今日.
I don't feel like going.
或者我細路都唔理.
佢係咩度.
佢一樣係度.
我地gwo 幾個ge 義工.
我真係好appreciate佢地.
冇休息ge .
喂.
琴晚佢地仲係度避課.

$^{401}$講講我老婆.
我真係好唔話得.
我覺得佢.
好忠心咁事奉.
落晒力ge .
冇休息過ge .
你話成日叫我休息.
不如你叫佢休息.
係我.
咁多年同佢一齊事奉.
去教會.
我發現一個道理.
當我太太.
高興ge 時候.
教會應該係健康.
當佢唔高興ge 時候.
我就唔健康.
我去教會.
可能都係.
麻麻地.
我可以理解.
因為我地黎到教會.
我地好多時都有.
顧客ge 前設.
因為我地.
黎去接受服務.
我地黎到呢度.
我地去接受服務.
有邊個去餐館.
食ye 個侍應係唔理你.
你嬲唔嬲.
嬲係咪.
你點樣去教訓佢.
唔比貼士.
冇錯.
平時都唔比.
今次更加唔比.
平時有比.
算你啦.
而家生意好難做.

$^{441}$所以我而家.
去比貼士我比多jor .
因為我真係比多.
佢做得好唔好都好.
但係教會就唔同.
喂.
如果挑剔得多.
意見多多ge 話.
手指尾又唔願意去郁ge 話.
我又呼籲你.
唔好做一個明ge 羅.
響ge 笨.
因為呢個對你冇益處.
因為睇到嗎.
係上帝要求你.
呢個係你係上帝面前有ge 呼召.
追求更大ge 恩賜.
就係愛之道.
The way of love.
愛ge 道理.
所以係度呼籲你唔好做一個明ge 羅響.
不如.
你教會ge 時候.
你掛住口唇邊另一句說話.
How can I help you.
What can I do.
我自然識一個好好ge 長老.
佢真係每次黎教會.
佢都係咁講.
What can I do for you.
How can I help.
How can I pray for you.
我有咩為你祈禱等等.
起碼你掛住口唇邊一句咁ge 說話.
你會提醒你自己.
呢個係教會.
上帝ge 教.
我係其中一個上帝教ge 家人.
What can I do.
Instead.

$^{481}$我黎到係消費ge .
當然你係arm 黎教會.
或者你摸索緊唔知信仰係乜ye .
你想摸索探索下.
除外.
但我想話.
如果你係係基督裡面ge .
你係信jor 主ge .
你自然你ge 心態.
你ge 取態應該係完全不一樣.
Right?.
人見到你ge 取態都應該係不一樣.
人地會Wonder.
原來係咁樣ge .
唔似去餐館.
唔似去消費.
唔係Religious Consumerism.
Right?.
特別係我地中間信得耐ge .
有邊個信jor 十年有多.
係我見到ge .
真係唔少.
十年有多.
十年做到個牧師出黎.
係讀神學讀三四年.
信主五六年.
你就可以.
訓練到你自己.
成為教會領袖.
你教會.
追求.
更大ge 讀.
更大ge 恩賜.
愛之讀.
你教會.
How can I help you?.
How can I use my gifts?.
呢個就係Home.
就唔係House.
呢度我認識一位弟兄.

$^{521}$唔講名.
佢工作好忙.
成日要加班.
佢話佢出唔到力.
但幾次同佢傾偈.
無論係屋企同佢傾偈.
或者約佢黎傾偈.
佢都係咁問ge .
佢有咩可以奉獻.
有咩可以.
我慷慨ge 去幫到你.
好為呢D ge 生命去感恩.
白玫瑰gwo 次.
邊個講到話.
十一奉獻歲前歲後.
邊個講.
係.
十一奉獻歲前.
定係歲後.
佢話Interesting question.
你唔覺得呢個question其實好silly.
點解會咁問.
你訓練你細路仔.
我講過呢個比方.
訓練你細路仔.
佢開始識行ge 時候.
Don't touch the hot stove.
個熱ge 爐頭你唔好掂.
到大過ge 時候.
如果佢八歲十歲ge 時候.
佢都係唔掂個爐頭.
十七歲二十歲佢都係唔掂個爐頭.
三十歲都唔掂個爐頭.
你同佢講點解你唔掂個爐頭.
你由細到大教我Don't touch the stove.
你唔覺得好silly咩.
因為佢唔明白.
Don't touch the stove係咩意思.
Use the stove properly.
到你成熟ge 時候.

$^{561}$就自然唔需要問呢D 咁硬ge 規矩.
你需要有智慧.
因為你成熟jor .
歲前歲後.
點解會問呢D 問題.
I think it's a silly question.
I think it's a stupid question.
Okok.
因為你要個number.
咁咪比十分之一.
十分之一.
你鍾意有個number.
十分之一.
歲前歲後都好.
邊個少ge .
十分之一.
得唔得.
我咁樣答你.
咁你又唔需要掙扎咁多呢D 問題.
奉獻.
但係如果你話.
係咪真係有所謂奉獻.
如果你要歸納舊約裡面.
gwo D 奉獻ge .
所謂條例.
奉獻的律法.
其實係30\%.
有冇諗過.
在第一世紀耶穌ge 時代.
佢話猶太人係比30\%.
因為仲有.
有first harvest.
跟住有每年.
三次去朝聖ge .
等等加埋仲有比窮人.
gwo D 咁樣等等.
30幾\%.
美國歲高唔高.
高呀.
你試下去加拿大.

$^{601}$加拿大可能仲高.
Well, I think it's a silly question.
因為最主要ge 係乜ye .
係應該係.
你同上帝ge 關係.
gwo 種ge 愛.
黎到去奉獻出黎.
其實最主要唔係十一奉獻.
最主要係慷慨.
Generosity.
Generous.
係咩意思.
十一奉獻.
十一奉獻jor .
跟住gwo 九分之十.
十分之九.
90\%.
就隨得我用.
樣樣都係我自己.
腦一分一毫都係上帝所賜.
個point就係.
我地要.
Be generous.
That's the whole point.
唔係再問gwo D 細路仔要問ge 問題.
要長大.
我頭先問你.
信jor 住幾耐.
十年gwo D .
冇理由要再問呢個問題.
It's a silly and stupid question.
因為上帝講乜ye .
你係教會建立上帝ge 家.
就係一個愛之道.
愛之道.
The way of love.
十三章第四節就係七節.
愛是恆久孕來.
又有恩慈.
愛是不疾度不自誇不張狂.

$^{641}$呢度唔係講緊婚姻.
呢段聖經所謂偉大的愛篇.
呢段聖經如果你未曾.
睇過聖經ge 話.
你要知道呢段聖經.
係歷史裡面最偉大的.
愛篇.
讀到邊度去先.
不作害羞的事.
不求自己的益處.
不輕易發怒.
不計算人的惡.
不喜歡不義.
只喜歡真理.
凡事包容凡事相信.
凡事盼望凡事忍耐.
唔係講緊婚姻.
講緊教會生活.
講緊弟兄姊妹.
大家之間應該怎樣去相處.
呢個係最大的恩慈.
更妙的度.
就係愛之度.
因為有幾多恩慈有幾多資材.
有幾多programs有幾多活動.
如果沒有呢樣ge 話.
算不得乜ye .
So what.
你睇你間屋有幾靚.
值幾多錢.
好巴閉係咪.
你發現關係做得唔好.
搞得唔好.
係唔和諧ge .
家就不成家.
Who cares.
你住ge 三千呎.
九千呎.
八萬房洋樓.
I really don't care.

$^{681}$Carege 你夫婦.
你家庭.
你之間ge 關係.
如果關係好ge 話.
住九百呎有乜所謂.
Right.
Seriously.
食乜ye 有乜所謂.
好話你日日食.
蟹食龍蝦.
日日食海鮮.
So what.
係咪膽固醇咁高D .
要有愛.
你先至能夠成家.
上帝ge 家都係一樣.
愛是恆久忍耐又有恩慈.
Explanatory.
By the way 呢度雖然唔係講婚姻.
但係樣樣ye 都apply to 婚姻.
Of course.
恆久忍耐.
gwo 度冇一樣ye 係同.
Romance有關.
全部都係motion.
唔係emotion.
Motions of love.
恆久忍耐.
有恩慈唔嫉妒.
唔自掛唔張角.
唔作害羞.
唔求自己ge 抑阻.
先求對方ge 抑阻.
凡事係忍耐.
All motions.
你留意度形容愛.
有D 乜ye 係.
抗衡世界ge 文化.
世界好多時.
講緊ge 愛.

$^{721}$同呢D 往往係相反.
Quick temper.
求自己ge 抑阻.
輕易ge 法律.
計算人ge 惡.
琴晚我地家人.
去睇一套ge 電影.
因為我個女Rachel.
就好想.
佢成日都重睇一D 電影.
即係有D 電影.
睇jor 好多次佢都要再睇.
佢琴晚就話.
不如我地一齊睇呢套.
Cars我睇.
我話唔係呀.
睇Cars 1.
你知唔知幾時拍緊呢套ye .
2006年.
嘩.
十幾年前.
但係Remasterjor 4Kge .
睇Cars.
我一家人就睇Cars.
個故事你記唔記得係講咩.
故事就係講.
一個年輕ge 快車.
Lightning McQueen.
紅色gwo 隻.
佢就係逢跑必贏ge .
比賽一定係贏.
佢就好想贏一個大獎.
一味就想世界冠軍.
但係不幸.
佢就被迫就留落jor .
美國gwo D 小鎮.
可能開Highway經過.
有D 小鎮幾百人.
好細好細ge 小鎮.
佢就好想好想走.

$^{761}$gwo 個小鎮阻佢住.
阻佢實現自己夢想.
咁所以佢就留落jor .
係gwo 度但係gwo 個小鎮.
又好富人情味.
gwo 度大家.
雖然人唔多.
但係好富人情味.
點知係gwo 度.
就留落jor 係gwo 度ge 時候.
佢就遇上jor 一個係退休ge .
神車手.
退休ge 神車手.
gwo 個神車手係贏過好多個獎杯.
就係佢想贏gwo D 獎杯.
贏過好幾屆.
裡面其中一個對白我覺得好有意思.
就係呢個神車手.
就去同佢講.
When was the last time.
you care about something.
except yourself.
你諗下係上一次.
你關懷ge .
關心ge ye .
係同你自己無關.
係幾時.
愛之道就係咁.
不求自己的益處.
唔單止.
係婚姻係咁.
教會生活都係一樣.
係關係裡面.
弟兄姊妹之間ge 關係.
一個.
走埋一齊.
有ge 團契.
所謂聚會.
裡面有D 乜ye .
Sexually係好次要ge ye .

$^{801}$最重要係gwo 班人.
你走埋一齊ge 時候.
唔係求自己益處.
係愛對方.
係我地.
以後ge 世界.
我地更加需要去重整呢樣ye .
有一樣ye 令我覺得好詫異.
幾年ge 疫情.
可以完全侵蝕.
一個基督徒的生命.
疫情之前好好地.
疫情一黎到ge 時候.
佢成個信仰生命就打垮jor .
That really amaze me.
點可以係咁.
你同上帝ge 關係.
係一個永恆ge 關係.
因為一個疫情.
乜ye 都可以忘記.
I just don't get it.
或者一個疫情.
令到我地.
唔識控制.
我地ge 情緒.
我地需要學習.
學習乜ye .
本來就係話愛之道.
佢講jor 乜ye  愛係必要.
有恩切冇用.
有活動冇用.
你沒有愛算不得什麼.
愛ge 定義.
愛ge 長久.
愛是永不止息.
先知ge 講道.
都歸於無有.
講方言ge .
終必停止 知識也要歸於無有.
我地唔係gwo 種ge 傳統.

$^{841}$唔係靈恩ge 傳統.
However.
你有幾多ge 活動.
幾多ge 人 幾多ge services.
終歸於無有.
我地而家所知ge .
只有限.
先知講道ge .
有ge 恩賜ge .
都會係有限.
到那完全ge .
有限ge 就會消逝.
我地而家做.
小孩子.
心思像小孩子.
意念像小孩子.
長大成人.
用大子ge ye 去雕起.
而家我地有ge ye .
就係對住鏡子.
模糊不清.
到那時 我地就要面對面建築.
諗清楚.
目的.
最終係要做D 咩.
完全ge 來到位.
建築面對面ge 時候.
到底咩ye 係常存.
信 忘 如 愛.
第三樣裡面.
最大ge 係愛.
你諗下你花幾多心思.
幾多精力.
去做D 乜ye .
有冇花更多精力.
你去行.
愛ge 道.
呢個先係最大ge 恩賜.
結束ge 時候我想講一個例子.
就係最近係基督教ge 世界裡面.

$^{881}$有一個好多ge 消息.
呢間ge .
細細ge 大學.
叫做Asbury University.
係Kentucky.
就係2月8號ge 時候.
一個平日ge 早會.
一個chapel.
學生一齊敬拜 然後聽道等等.
點知完jor 之後.
gwo D 學生唔想完.
佢地就自發ge .
繼續聚會.
繼續去敬拜.
禱告 繼續去學習.
甚至去奉獻 去事奉.
等等.
一路一路咁樣去聚會.
24小時輪流ge .
唔停止ge .
已經三個禮拜.
最後校方話.
我地係時候停.
校長講ge ye 好有意思.
無論你點樣去.
詮釋呢件事都好.
係咪一個復興 所謂revival.
係咪聖靈所比ge 一個.
屬靈ge 奉獻都好 點樣詮釋都好.
校長講咩ye .
佢話people are hungry.
佢地想去到.
渴求上帝.
想去上帝ge 同在.
gwo 度面 去歸向上帝.
去敬拜上帝.
係咁樣ge 心態 咁樣ge .
情況 咁樣ge 現象.
你點樣詮釋佢都好 點樣解釋都好.
校長講咩呢.

$^{921}$佢見證到乜ye 呢.
見證到聖經裡面所形容ge .
聖靈ge 果子.
聖靈ge 果子九樣.
第一樣係乜.
人愛.
人愛喜樂和平忍耐.
恩慈 良善信神.
溫柔節制.
佢所見證到ge 就係.
呢D 大學生ge 生命.
能夠展現到呢一樣ye .
I think he is right.
唔係一咩emotional high等等.
我地好多時.
黎教會 最近聽到一個弟兄同我講.
嘩如果我地有個full band就好啦.
好激昂 好激苦.
好偉大ge 咁樣敬拜.
情緒好澎湃ge .
咁樣好唔好呀.
敬拜ge 時候 音樂ge 時候.
音樂敬拜 好唔好呀.
好 但係如果係顧客心態ge 我就唔做.
What is that for.
幾動人都好.
算不得什麼.
就係呢D 人ge 生命.
有冇去歸向上帝.
呢個先係最重要.
佢冇歸向上帝 而係你自己.
個人ge 生命裡面有重整.
係你婚姻ge 關係.
係你家庭ge 關係.
能夠有重整.
係你教會 工作.
甚至到整個社會.
能夠有重整 呢個先係最重要.
Right.
有幾多呢D 咁ge .

$^{961}$放幾多炮仗都好.
放幾多煙花都好.
算不得什麼.
最重要係你ge 生命.
有冇被改變過.
求主幫助我地.
The way of love.
花時間.
花精力去做呢樣ye .
Right.
邊個去負責.
你個人ge 生命.
係能夠去實踐愛.
你.
唔係教會.
邊個去負責.
你細路仔.
ge 個人生命.
你.
唔係教會.
係你先行.
教會去輔助.
邊個去負責 你有冇一個美好ge .
見證ge 生命.
你.
責無旁貸.
如果上帝ge 愛.
已經進入你生命.
Let it be.
讓上帝ge 生命去充滿你.
去改變去更新你.
以致你能夠成為一個.
成熟ge .
基督徒 健康ge 基督徒.
以致教會能夠健康.
可以同心大道.
祈禱.
主愛多謝你.
因為.
你係愛我地ge 神.

$^{1001}$縱使我地係破碎.
我地係軟弱ge .
但係因為你對我地ge .
大愛.
係基督ge 十教裡面去彰顯出黎.
主愛你ge 生命.
願意比你地度去.
傾住.
讓你ge 愛充滿我地.
以致我地ge 生命能夠去彰顯到.
基督自己ge 生命.
以致我地ge 教會能夠成為一個.
合理心意ge 教會.
求主你憐憫 由我地教會能夠去.
討你ge 喜悅 去榮耀神靈.
祈禱 奉主名.
好 原上祝福各位.
下個禮拜能夠同大家一齊去事奉.
謝謝收看.
by bwd6.
\newpage



\section{}
\label{sec:CFdO0Ddf8oE}
\textbf{Becoming A Healthier Church 5: Way of Love - Rev. Mark Chan}
\newline
\newline
連結: \href{https://youtube.com/watch?v=CFdO0Ddf8oE}{\texttt{https://youtube.com/watch?v=CFdO0Ddf8oE}} ~~~~ 語音日期: 2023-03-07
\newline
\newline
\hyperref[sec:5SRk2JnEBBc]{\small{< < < PREV SERMON < < <}}
~
\hyperref[sec:index_chronic]{\small{[返順時目]}}
~
\hyperref[sec:index_scriptual]{\small{[返順卷目]}}
~
\hyperref[sec:z2UTnfWH0W8]{\small{> > > NEXT SERMON > > >}}
\newline
\newline
$^{1}$Thank you everyone..
- All right, good to see you all..
And why don't we say hi to each other?.
Okay, just a little, yeah..
There we go, all right..
Good to be together as a church family.
as we continue to reflect on the topic of being a church,.
being a healthier church,.
being a healthy church that will bring glory to God..
Again, the adjective is important..
It's a healthy church..
It's not a good church..
Now, I don't like that adjective, good..
Is this church good church?.
You know, it's really subjective..
But you know, at least another adjective, healthy church,.
will bring us a little bit closer to what the Bible says..
Of course, the best adjective, I think,.
is biblical church, right?.
But biblical church, you know,.
it depends on how much do we know about the Bible,.
what the Bible says about church..
So I hope this series will continue to, you know,.
help you ponder on the question.
of what the church is all about, okay?.
Now, an analogy, an analogy..
If you go to a restaurant, you check the menu, right?.
And you're gonna order something from the menu..
Do you think a restaurant is a good restaurant.
because the menu has tons of items?.
No?.
Have you ever been to a restaurant like that?.
- [Audience Member] Oh, yeah..
(audience laughing).
- There you go, yeah..
You know, sometimes we go to a restaurant like that,.
and we probably have to spend 20 minutes or so.
in order to pick what you like, right?.
Too many menu items..
Sometimes it's good, sometimes, you know,.

$^{41}$well, what I'm saying is,.
what that restaurant is a good restaurant..
(audience laughing).
All right, because we're going online, right?.
Yeah, I don't take any commission,.
but that restaurant is a good restaurant..
But, you know, all I'm saying is,.
it's not necessarily so, right?.
You go to a restaurant,.
probably the restaurant has a few menu items,.
and they do really good..
They do really good quality food, right?.
And that is also a great restaurant, right?.
Some weeks ago, I was in California, right?.
And people always tell me, if you go to California,.
you gotta go to In-N-Out, right?.
You have to go there, right?.
So it's kind of like a pilgrimage..
You gotta go there and order or eat there,.
you know, order a hamburger, right?.
Yeah, those who are from California,.
you know, In-N-Out is a big thing, right?.
Anyways, I don't think In-N-Out is as good as Culver's here..
You think so?.
No?.
How many of you think Culver's is better than In-N-Out?.
- Yeah..
- Culver's, yeah, right, yeah..
Well, anyways, now my example is,.
you go to In-N-Out, you know, how many items do they have?.
Only a few, right?.
Only a few..
But they do it really, really well, right?.
It's a great restaurant..
Same for church..
No, a healthy church does not mean.
that you have all kinds of programs.
that you can think of, right?.
Oh, you can walk in the church.
and you can think of, you know,.

$^{81}$babysitting, preschool, Sunday school,.
kids' Sunday school, youth,.
college, young adult, couples,.
and senior adults,.
or all kinds of life stages are being covered..
Does it mean, does it mean if a church.
has all the menu items,.
does it mean the church is a healthy church?.
Not necessarily so, right?.
It would be nice to have those things..
But I want you to think about.
what a church really means, right?.
So there's nothing better to go to the Word of God..
As I said, as I emphasized before,.
the authority of God's Word..
You are Christian, how you live,.
and what you think about God,.
how you think about God,.
is not by your own preferences..
How you think about ministry,.
how you think about serving God,.
using your gifts in different areas of the church,.
is not according to your preferences..
You may have preferences,.
you may have your own opinions..
However, the most important thing is.
you go back and return to the Word of God..
That is the most important thing,.
because the Word of God is the ultimate standard.
and authority..
Now, we have been studying the book of 1 Corinthians..
If you have your Bible, turn to it..
We were at chapter 12..
And as you know, Apostle Paul, as a pastor,.
he wrote this letter to the church in Corinth.
and tried to address some of the problems they have..
Now, you think about it,.
among all the New Testament churches,.
you can think of church in Philippi,.
church in Galatia,.

$^{121}$church in Ephesus,.
anywhere else..
The church in Corinth probably is the Macca church.
of the first century..
Macca church doesn't mean a lot of people..
Probably it was a really, really small church, right?.
That's the same case for other churches.
in the New Testament era..
However, they are the best in terms of manual items,.
in terms of people being gifted and talented in the church..
In fact, the Apostle Paul said that is the case..
Right off the bat in chapter one, verse five,.
he said, "Through him, God has enriched you guys.
"in every single way you can imagine,.
"your eloquent words and all of your knowledge.".
So you can pick one guy or some person.
from the church of Corinth,.
you can tell, wow, he's so gifted, right?.
This person or that person is so gifted.
or probably in a very high level.
in terms of the social status..
So Paul actually praises the church in Corinth.
and saying, "You guys are so gifted and so talented.
"and so equipped with all kinds of words and knowledge.
"and talents and gifts.".
That's why in chapter 12,.
he come to address the issues arising from spiritual gifts..
He has to teach them and remind them.
what spiritual gift is all about..
As you know, the Corinthian church,.
they are gifted in some way.
that other churches probably don't have..
The speaking of tongue, right?.
They can speak the language of the angels,.
the heavenly language..
Nobody could understand unless there was someone else.
who is gifted how to interpret that kind of language.
into human language, right?.
So whether you believe it or not,.
that is being addressed in chapter 12..

$^{161}$So Paul is trying to tell them,.
"You guys are so spiritually gifted,.
"but you need to understand the ultimate goal.
"of spiritual gifts or how you being talented.
"or your resources, your spiritual gifts that you have..
"What is the ultimate purpose?".
The ultimate purpose, remember, is common good..
It is for common good..
You are gifted in a certain way,.
you are gifted in some other way..
You need to think of the image of body..
We are just parts of the whole body,.
making up the whole community..
So in chapter 12, toward the end,.
the apostle Paul said,.
"All of you together are the body of Christ,.
"and each of you is a part of it.".
Here are some of the parts that God has appointed,.
he has ordained for the church..
First, first set of gifts is apostles..
Second is our prophets..
Third are teachers..
And then those who do miracles..
Would be nice to have a church.
where some people can do miracles, right?.
Yeah, right?.
If something doesn't work, you know,.
(fingers snapping).
right, it works..
That's nice, right?.
Well, the church in Corinth,.
they do have some people like that..
Those who have the gift of healing,.
that would be nice, right?.
I'm not feeling good, you know, these days..
What if we have a spiritual doctor, right?.
(fingers snapping).
I'm okay..
It would be nice, right?.
Those who can help others,.

$^{201}$and those who have the gift of leadership,.
those who speak in unknown languages,.
the speaking of tongue..
Well, pay attention to how Paul.
organized or prioritize the spiritual gifts..
He put apostle, prophets, and teachers.
the very beginning of the list..
And then he pushes down the so-called spectacular gifts..
Those gifts that can show off, right?.
Those gifts that can express, you know,.
how powerful the presence of God is, right?.
Like speaking of tongue..
But he put that in the last, right?.
In the last set of the gifts..
So there's a priority, why?.
Again, it's the ultimate question is,.
what is the gift for?.
What are spiritual gifts for?.
What are your resources for?.
It's for the common good..
And he put apostles, those who establish churches,.
church planters these days,.
prophets who proclaim and speak the word of God.
in authority, teachers, teaching the church,.
teaching other people what to believe and how to behave..
These three gifts are put in the priority list..
And because they are used to build up other people,.
not for themselves..
Now, there were many, many different gifts.
and God is using those gifts.
so that we can build up the church..
Now, in my prayer, I mentioned.
there was a core team meeting today, right?.
We have a monthly core team meeting.
on the first Sunday afternoon..
So all the core team members are deacons..
We don't call them deacon,.
but we call them core leaders, right?.
Core team members..
We're gonna meet together for a couple of hours.

$^{241}$for guiding the church and leading the ministries..
If you sit down with the core team members,.
I thank God that we are being gifted in so many,.
why a variety of gifts, right?.
How many of you know that I'm not very good at finance?.
I'm really not that good, yeah, right?.
I'm really not that good, you know?.
I don't track the budget..
I try to, right?.
Do you guys know how to use Excel?.
Excel, yeah?.
So only a few of you?.
Oh man..
I use Excel, you know?.
I know how to use it, you know?.
The basic functions, you know,.
how to calculate the sum or summary and all the formulas..
I know how to do it..
But you know, I'm just not good at numbers..
I'm just not good at finance, right?.
And my wife is so much better, you know,.
handling finance..
So when we got married,.
she has always been the one who is,.
she's the watchdog, you know?.
She's telling me, "Hey, you overspend, you know?.
"You need to stop, okay?.
"Yeah, whatever you wanna buy,.
"but you know, you have to wait," right?.
She is the watchdog..
She key in almost every single transaction.
in a Excel worksheet, right?.
Some people are really gifted in that way, you know?.
They may not be the ones who are going upstage.
and doing something in public,.
but they are good..
They are really gifted..
So in our quality members,.
there are a couple of people.
who are really gifted in that way..

$^{281}$So they use the gifts for common good.
so that we can build a church together, right?.
You cannot imagine if our church,.
the finance and budget are managed by me..
That would be a disaster, right?.
It would be a disaster, right?.
And the accountants, those who are trained.
in handling numbers and finance,.
there we go, right?.
And some people are really good at IT stuff, right?.
The computer and all this stuff..
We rely on them in order to have a smooth surface, right?.
Smooth operation, right?.
Some people are really gifted in youth ministry,.
talking to young people, right?.
I'm not one of those, right?.
I cannot teach young people, right?.
Some people are really, really gifted in that way,.
and God is calling them to do those things..
Some people obviously are gifted in music, right?.
They can pick up an instrument.
and lead us in praises and in inspiration..
So God uses different gifts together.
in order to build up our church..
So my homework for you is that,.
remember I told you something about spiritual gifts, right?.
In the first round of our series..
Is to discover and to develop, right?.
Your spiritual gifts, right?.
So my homework for you is,.
after the service, you go to at least three people.
and tell them what I think you're gifted at, okay?.
What I think you're gifted at..
So not for yourself, you go to the other person.
and say, "Hmm, I think you're gifted at this or that.".
Right?.
All right, no pressure, no pressure..
It's just you think, right?.
Yeah..
The person may not think that way, but you know..

$^{321}$Well, it's one of the ways.
that we can encourage one another to discover..
Always tell people, you know,.
spiritual gifts are like pepper in the teeth..
Pepper in the teeth..
You know, have you had some like black pepper,.
you know, seasoned food,.
and you had it and it looks really gross, right?.
The pepper is between your teeth..
Nobody can see it..
You need somebody else to tell you,.
"Hey, you got some pepper in your teeth," right?.
Spiritual gifts work like this..
Nobody knows that you are gifted in this or that,.
unless other people tell you, right?.
That's what I'm saying..
Pepper in the teeth..
Somebody has to tell you..
Somebody has to confirm it..
The community has to confirm it..
Before I went to seminary,.
I felt called to ministry, right?.
I felt called, at least, you know,.
I feel that I am gifted teaching, right, in teaching..
So that is the direction I was going forward..
And I did talk to my elders of the church, my home church,.
and they have to come together and say,.
"Hey, Mark, you are, I think you're gifted..
"Yeah, but you still need training..
"Yeah, you still suck..
"But you still need training..
"You still need to be equipped, right?.
"But I think you, maybe you got the potential, right?.
"The trajectory, I can see that..
"So you have to work hard, too, right?.
"So, you know, in school, in seminary, I work hard..
"I really work hard..
"It doesn't mean I'm gifted..
"Then everything is smooth, right?.
"I don't need to do homework or, you know,.

$^{361}$"and all that stuff..
"I still need to work hard..
"I still need to develop it.".
That is the nature of spiritual gifts, right?.
So all the Moody students here,.
let me remind you, you know, study hard..
Okay, all right..
So what is spiritual gifts for?.
And he continues, Paul continues to talk about it..
He said, "It's not enough that you are gifted..
"It's not enough that you think you are talented..
"The most important thing is you should earnestly.
"desire the most helpful gifts..
"The most helpful gifts are those gifts,.
"those talents that you can use for the common good.
"to build up the body of the church.".
So Paul is saying, "Now, let me tell you.
"a way of life that is best of all.".
And that way of life is the way of love..
Take a look at this passage..
He said three hypothetical situations..
If I could speak all the languages of the earth,.
even the language of angels, but did not love others,.
I would be a noisy gong, noisy cymbal..
Second situation, if I had the gift of prophecy.
and I understand all the secret,.
I possess all the spiritual knowledge, mystical knowledge..
If I have such faith that I could move mountains,.
but did not love others, I would be nothing..
Not even something, nothing without love..
If I gave everything I have to the poor.
and even sacrificed my body, I could boast about it,.
but if I did not love others, I would have gained nothing..
What Paul is saying is without love,.
no matter how gifted you are, it means nothing..
Three times he said, "If I did not love,.
if I did not use my gifts and use my time.
and talents, my treasure for the common good,.
for the eternal cause, I would be nothing.".
Now, at this point, I wanna,.

$^{401}$I just wanna appreciate some of the people.
who have been with us in the last few years,.
especially the last few years..
As you know, because of the pandemic,.
the last three, four years, we're really, really tough..
We're really, really tough..
You know, I still remember in year 2020, in March,.
we came back here because of the CDC guideline..
Everything has to be locked down..
Everything has to be shut down,.
and we have to switch to online completely..
I still remember one weekday, weekday evening,.
we came with some brothers and sisters,.
and we tried to grab something,.
and then we hauled it back to my home..
And my basement became the workstation, right?.
The studio..
And we did the online service from there..
It's been so tough..
And a number of months ago, we opened, right?.
Now, what I wanna appreciate is all those sisters.
and brothers, the workers,.
who have been working tirelessly during this period..
Why I appreciate them most is that it was tough..
When you're doing something.
and you don't see the results immediately, right?.
When you're serving and you don't see people returning,.
it has been really, really tough..
That's why I wanna appreciate each one of them..
I can't name all the names,.
but I really appreciate those who are trying.
to get things going, right, with our limited resources,.
to get the worship going, to get the singing going,.
teaching online, preparing something backstage,.
figure out the IT problems, troubleshoot all the issues..
Right?.
Just cannot imagine..
This is crazy..
So I would say they are tired..
They are exhausted..

$^{441}$If you talk about burnout, you know,.
I think they are the first one to get burned out..
It's just not easy, right?.
Not easy..
So I really appreciate all of them..
They did not quit..
And they are with us until this day..
Really appreciate it..
Now, when we open, we are still working tirelessly, right?.
Just give you an office example, it's my wife, right?.
You know, she's a tough nut, yeah..
She's really, really tough..
You know, of course, when we switched to online,.
she got, she had no kids ministry, right?.
Everything is done by video and all that stuff,.
and she didn't do it..
But when we opened some months ago,.
she's the one and with some other sisters and brothers.
trying to, you know, we build and we open the kids ministry..
What I really appreciate people like this.
is that no matter what, they are there..
No kid, they still prepare, right?.
In fact, last night, we came home,.
and she was preparing for Sunday school..
Of course, with her team, when I said team,.
there's only one or two people..
But they are tirelessly working together.
in order to build a ministry..
So, you know, we, sometimes we throw up the church..
It's very natural for us to assume the role of consumers..
We come here, we enjoy it..
Do you enjoy today?.
Today's church?.
Eh, not really, right?.
Yeah..
Do you enjoy the program here?.
You know, that kind of mentality is customer mentality..
I mean, I can't blame you..
I mean, it's very natural, right?.
You show up, yeah, there's something here, right?.

$^{481}$You don't show up, eh, they're okay, right?.
I wanna say with apostle Paul, we need to do it with love..
We are being called together to constitute the church,.
to form the church body, the spiritual family of God..
Everyone has his part..
Everyone has her part to play, to give,.
to use our talents to build up the church..
That is the way of love..
And that's why 13.4 and following.
is the most famous love passage of all times..
Take a look..
Love is patient and kind..
Love is not jealous, not boastful or proud or rude..
It does not demand its own way..
It is not irritable..
It gives no record of being wrong..
It does not rejoice in injustice,.
but rejoices whenever the truth wins out..
Love never gives up, never loses faith..
It's always hopeful, endures through every circumstance..
That is love..
Remember, this passage is often used in Christian wedding..
This passage actually has nothing to do with marriage..
Remember that?.
The context I told you from chapter 12, now to chapter 13,.
Paul is talking about church life, the way of love..
Love is patient, love is kind, right?.
You come to the church, we live together.
as a church body, as a church family..
The best way to do it is not administration..
The best way to do it is not leadership..
The best way to do it is not we have spectacular programs..
The best way to come together as a church.
that is pleasing to God is the way of love..
That means when you show up, when you come here,.
you think of others..
That is what love is..
It's not self-seeking..
It's not about me, it's about you..
Remember I said, you know, this is,.

$^{521}$at least in the context, it's not talking about marriage..
But every single principle here.
actually is applicable to marriage, right?.
In all relationship, in all social gatherings,.
in all situations, love is the best way..
Now a couple days ago, or actually a couple weeks ago,.
I think last week, a week ago, yeah,.
my daughter suggested that we as a family,.
we're gonna watch a movie, all right?.
So what kind of movie that she likes to watch?.
Um, it's this one..
Cars..
No, as you know, have you ever watched this one?.
Cars 1?.
You did?.
Nobody?.
Oh, come on, this is probably one of the best movies, okay?.
So you should get together and watch this movie..
And it talks about Lightning McQueen, right?.
The race car..
And he always want to win the race, right?.
To win the cup, right?.
The what?.
The piston cup, right?.
The piston cup, whatever..
He just wanna be number one..
He just wanna beat other people,.
and then he's gonna win, right?.
And unfortunately, he end up in a very small town.
called Radiator Spring, right?.
Radiator Spring on Route 66..
And he got stuck there, but he quickly find out.
this town of not too many people,.
they are very simple people..
You know, they love the town..
They wanna build up the town..
They feel sad if the town is being abandoned,.
being in ruin..
So they love each other and care for one another..
It's a small family..

$^{561}$So they went to, so the story,.
a lot of the stories happened in that little town,.
Radiator Springs..
And then, eventually, Lightning McQueen,.
this young race car, met someone.
who has been retired,.
used to be very successful race car..
Dog, right?.
Dog..
And there's one line in this movie,.
I think it's very, very meaningful..
Because this old guy, you know, retired,.
and he look at the piston cups like nothing, right?.
It's just empty cup, right?.
In fact, he won a few of them, right?.
So for him, it really means nothing to him..
And this old guy, retired,.
taught Lightning McQueen with this line..
He said, "When was the last time.
"you care about something except yourself?".
When was the last time?.
You care about something except yourself..
That's a great line..
When was the last time you care about something.
except yourself?.
Like how you enjoy ministry?.
Like how you enjoy church?.
Like how you enjoy the family?.
Like how you get something out of it, right?.
When was the last time you turn the table around and say,.
you know, that is not important.
because that is not the way of love..
Love is not self-seeking..
I'm not number one..
Who cares?.
Whether I get something out of it, right?.
Whether I enjoy it..
That is not the point..
The point is whether I'm practicing the way of love..
Paul continues..

$^{601}$He said, "Prophecy, speaking in unknown languages,.
"speak special knowledge, will eventually become useless.".
It will be useless..
Eventually it will be useless..
In all the programs, fantastic stuff,.
you can make it happen in the church..
Eventually it will become nothing.
because they will not last..
They will not last forever..
They are not eternal..
However, the kids ministry, how exciting it is,.
how entertaining it is, how engaging it is,.
Paul is saying it will become nothing.
compared to what the gifts, the spectacular gifts.
the Corinthian Christians have,.
speaking in unknown languages..
Wow, how great that is, right?.
You know, when I was in seminary,.
it was in Southern California,.
I was roommate with three other Korean guys..
One of those guys, Korean guys,.
he said he's gifted in speaking in tongue..
Have you ever heard speaking in tongue?.
Every morning, he was doing his devotional, right?.
And he prays..
And nobody understands what he's saying..
But he is so close to God..
That's what, at least that's what he said, right?.
Now, I for sure, I for sure, he was not speaking Korean,.
even though I know nothing about Korean..
But those are not human language, right?.
But he was using that language to get close to God, right?.
He was gifted in such a way..
But Paul is saying, even so,.
it will become useless..
But love will last forever..
Love will last forever..
Even the gift of prophecy,.
only reveal part of the reality, right?.
No matter how great the teaching is,.

$^{641}$only, we only see partially the whole picture..
He said, when the time comes,.
that means when Jesus comes back,.
time of perfection comes,.
these partial things will become useless..
Useless..
Now, we need to mature..
When I get mature, and Paul is saying,.
"When I was a child, I spoke like a child..
"When I grew up, I put away childish things..
"And now we see things imperfectly,.
"like puzzling reflections in the mirror.".
But then, when Jesus comes, when the kingdom comes,.
we will see everything with perfect clarity..
Now what we see is incomplete,.
but we will know everything completely when the time comes..
Therefore, only something will last forever..
And Paul said, three things will last forever..
Surprisingly, nothing about speaking in tongues,.
nothing about your spiritual gifts,.
nothing about your talents..
Only three things will last forever..
Faith, hope, and love..
And the greatest of these is love..
So my challenge to you today is that you come to church,.
or even you're not going to church,.
you're not in the weekdays, right?.
You're not seeing other sisters and brothers..
How do you practice the way of love?.
Well, at least I would suggest to pray for the church..
When you come together, when you care for one another,.
love is not self-seeking..
It's to seek the benefit of the church, the common good..
Seek the benefit of other people..
So let's all pray and prepare our hearts and our mind.
and pray to God..
I spend some moment to pray that, you know,.
come to God and say, "Dear God,.
"let it not be about myself,.
"but about the kingdom, the church..

$^{681}$"Dear God, can you use me in such a way.
"to build up the body of Christ?".
Let's all pray together..
(pastor and congregation murmuring).
Father, we come before you.
and we realize that we are your vessels..
We are simply your instruments,.
that you use your sovereignty.
to grant us different gifts and talents.
and treasure and resources..
And I pray that we will be a church.
that will practice the way of love..
To care for one another, to put others before ourselves,.
to serve other people for the benefit of your church.
so that your name will be made famous..
That people will come and observe us.
as a church family, as a spiritual family..
They will be surprised by the way of love..
They will see us as a community.
that is different from the culture..
Yes, the world is teaching us and molding us.
in such a way that we always come as a customer.
to seek for our own enjoyment,.
to seek for our own benefit..
But we come before you at the foot of the cross..
We know that the best way, the eternal way.
that will last forever is the way of love..
Help us to care for one another..
Help us to pray for one another..
To build up your church so that your family will be strong..
Strong in love, close in love..
And I pray that you will continue to grow our church.
and grow in such a way that we will continue.
to be mature in the way of love..
In Christ name we pray, amen..
Well, thank you for coming..
And I pray that you will have a good time today.
and this week and God bless you in your work.
and your career and your family..
And make sure you pray for one another,.

$^{721}$especially those who are out of town this week or next week.
and we'll continue to serve one another in our church..
All right, God bless you..
\newpage



\section{使徒行傳 26:15-20}
\label{sec:z2UTnfWH0W8}
\textbf{非常任務 (使徒行傳26\_15-20) - 黃茂林傳道}
\newline
\newline
連結: \href{https://youtube.com/watch?v=z2UTnfWH0W8}{\texttt{https://youtube.com/watch?v=z2UTnfWH0W8}} ~~~~ 語音日期: 2023-03-07
\newline
\newline
\hyperref[sec:CFdO0Ddf8oE]{\small{< < < PREV SERMON < < <}}
~
\hyperref[sec:index_chronic]{\small{[返順時目]}}
~
\hyperref[sec:index_scriptual]{\small{[返順卷目]}}
~
\hyperref[sec:bEexjAddGwM]{\small{> > > NEXT SERMON > > >}}
\newline
\newline
使徒行傳 26:15-20
\newline
\begin{longtable}{cl}
\hline
\hline
章節 & 經文 (和合本修訂版)\\
\hline
26:15 & \begin{tabularx}{0.7\textwidth}{X} 我說：『主啊，你是誰？』主說：『我就是你所迫害的耶穌。 \end{tabularx} \\ \\ \relax
26:16 & \begin{tabularx}{0.7\textwidth}{X} 起來，站著，我向你顯現的目的是要派你作僕役，為你所看見我的事，和我將要指示你的事作見證人。 \end{tabularx} \\ \\ \relax
26:17 & \begin{tabularx}{0.7\textwidth}{X} 我也要救你脫離百姓和外邦人的手。我差你到他們那裡去， \end{tabularx} \\ \\ \relax
26:18 & \begin{tabularx}{0.7\textwidth}{X} 要開他們的眼睛，使他們從黑暗中轉向光明，從撒但權下歸向神；使他們因信我而得蒙赦罪，和一切成聖的人同得基業。』」 \end{tabularx} \\ \\ \relax
26:19 & \begin{tabularx}{0.7\textwidth}{X} 「因此，亞基帕王啊！我沒有違背那從天上來的異象； \end{tabularx} \\ \\ \relax
26:20 & \begin{tabularx}{0.7\textwidth}{X} 我先在大馬士革，後在耶路撒冷和猶太全地，以及外邦，勸勉他們應當悔改歸向神，行事與悔改的心相稱。 \end{tabularx} \\ \\
[1ex]
\hline
\hline
\end{longtable}
$^{1}$講到之後就有聖餐的時間.
以下時間交給Mark Sir去讀經.
大家好 節目早上好.
讓我們打開聖經.
Pastor Samuel和太太都在.
Nancy 歡迎你.
今天他們講到的經文就是史郎傳第26章.
15至20節.
如果你有聖經的話不妨翻開.
PowerPoint上面也有.
我先讀上帝的話語.
我說 主啊你是誰?主說:.
我就是你所逼迫的耶穌.
你起來站著.
我特意向你顯現.
要派你作執事,作見證.
將你所看見的事.
和我要指示你的事證明出來.
我也要救你脫離百姓和外邦人的手.
我差你到他們那裡去.
要叫他們的眼睛得開.
從黑暗中歸向光明.
從撒旦泉下歸向神.
有因信我.
得蒙謝罪.
和一切誠信的人同得基業.
阿基帕王啊.
我故此沒有違背那從天上來的意象.
先在大馬士革.
後在耶路撒冷和猶太傳遞.
以及外邦.
勸勉他們應當悔改歸向神.
行事與悔改的心相承.
我們請王傳德來歡迎他.
有太太好一點.
有些人獨居不好.
有些人一個出一趟都不好.
城會裡面來分享主人的訊息.
為這個聚會禱告了很久了.
上一次來的時候.

$^{41}$也是疫情還沒開始的時間.
轉眼就三年.
大家又成熟了.
大家有些面孔我都不認得了.
但是我很感覺到大家那種朝氣.
是不是.
年輕.
心境都年輕.
我都很年輕的其實.
我很想和大家分享今天的訊息.
就是非常任務.
什麼是非常任務呢.
不知道怎麼說.
我想起我從前有一個非常的任務.
我從前是做核數的.
進去核數師樓.
幾個月中間不是太忙.
有一天就收到一份特別的任務.
就有個上司召我進房.
交給我一個很厚的信封.
用一些很硬的雞皮紙.
寫了那個名字.
誰收.
用那些蠟印.
用那些wax的燒.
然後簽名交給我.
就說很認真地跟我說.
你記得要將這個信封.
交給這個名字上面那個人.
你不要想著回答人家幫你.
你又不要想著順道去做這件事.
你專門要做這樣的工作.
他簽回一個回調作實.
核對那個名字.
就是核對身份證.
然後你才回來.
拿了這個信封之後.
我就不可以休息.
上下網玩.
我馬上就要拿著這個東西.

$^{81}$去一個地方.
看一下地址.
ICAC.
就完了這個工作之後.
我就要第一時間.
call回公司.
告訴我的部門主管.
交條量.
然後親手交回調.
在上司的手中.
我這個case就搞定.
這個是我的非常任務.
我連想問一下.
究竟有什麼特別.
要我這個專業人士.
去做專業的送派服務.
後來我才知道原來那份東西.
就是關乎當時一個公職的高官.
以權謀私貪污的財政證據.
交給他那個人.
我也不知道是誰.
接一個非常任務.
就有非常的要求.
我很記得一件事就是.
由頭到尾我非常緊張.
我很怕搞錯了.
譬如忘記了對名字.
忘記了簽名.
或者那個人收不到.
保羅他都接受了一個非常任務.
他的非常任務.
有不同的機會都和人家複述.
剛才讀的聖經.
《士唐傳》的26章.
他有機會在一個見證的裡面.
和人家去講.
我覺得保羅這個是非常聰明.
又非常有效的.
人家都想你講了.
你還不講.

$^{121}$有時候我們傳呼音的機會.
什麼時候來呢.
就是人家都問你了.
你還不講.
你有沒有試過.
有時候要製造機會.
這些是要學的.
人家會送上門過來問你的.
保羅都遇到這個機會.
向這個分封的王阿基伯和他太太.
伯尼基士一齊來分享.
因為他藉著自己羅馬公民的身份.
在中國的訴訟裡.
他行使他的權利.
上訴該殺.
當時是可以的.
還要給他女婿.
有人保護他.
在他的分數裡面.
剛剛包括剛才Mark Sir所讀的聖經.
所以透過保羅他的見證.
我們又看到他怎樣看自己.
那個非常的任務.
首先我們看到.
他認定自己的任務是怎樣來的呢.
是一個更高的權.
更高的掌管者.
我們叫他上頭.
這個更高的上頭交給他的.
我的工作就是當時的部門負責.
那個高級經理.
僅次於合夥人的partner.
交給我.
我當然很害怕.
很嚴肅看待這件事.
所以同樣保羅也都很嚴肅看待.
主耶穌跟他說交代他的事.
這個經文裡面提到的.
主耶穌是他自己這樣說的.
主親口說.

$^{161}$我就是你所逼迫的耶穌.
從前保羅有眼不識泰山.
自以為是拼命來逼迫.
主耶穌那些的跟隨者.
那他一旦認清了身份之後.
他就反過來拼命為主做見證.
不怕別人取笑.
不怕別人威脅他的生命.
絕對是認真不苟且的.
我親眼都碰過很多信耶穌的朋友.
從前他是反對神的.
用基督徒來取笑的.
或者是做很多事情來傷害基督徒的.
但當他一信了耶穌之後.
他是改變過來的.
一定改變的.
他裡面就很認真去看待.
他信耶穌的這件事情.
我不知道我們在座之中.
有多認真看待你信仰的事情.
可能有時候.
有些人生活無聊.
沒什麼好做.
禮拜天那些人都還沒醒.
不如回來教會吧.
等那些人醒了才出街玩.
有些人這樣想的.
但也有些人好好地準備自己去敬拜神.
保羅很認真地看到.
他難得這個機會.
因為這個交給他的任務.
就是一個愛他的主.
很有權柄的.
沒有人不會不聽他的話的.
權柄這件事大家感受到嗎?.
做事的人很覺得.
有很多青年人還沒做事的時候.
不覺得什麼叫權柄.
出來做事之後就知道什麼叫權柄.
我可以捨死你的.

$^{201}$我可以捨得你被人調職.
我捨得你永遠升不到職.
甚至被人解僱.
你怕不怕?.
我出來做第一份工作.
我的管理員就嚇我.
你知不知道兩年之內就死了.
因為兩年一個最後評估.
能不能做到會計師.
就捨死了你.
你就慘了.
我記得有一個小女孩.
以前考上中學就要派位.
她家人有些問題.
明明公佈了結果.
就要在幾天之內去登記.
她家人就不知道有什麼事.
旅行之類.
錯過了登記日子.
回來就去校務處.
校務處就說.
沒有了.
你這個中學位很矜貴的.
沒有了.
多次跟進不得要領.
跟著那個小女孩和家人走著走著.
經過公園就瀑了在那裡哭.
哭得很厲害.
有位阿姨經過.
看到小女孩哭得這麼慘.
一邊安慰這個小女孩.
一邊又安慰她的家人.
不要擔心.
跟著就拿出電話.
你寫下這個電話.
下個星期一.
打這個電話.
叫辦事處.
看能不能幫你.
你記得打.

$^{241}$家人就.
我想小女孩有沒有哭也不要緊.
家人就覺得.
不如試試打電話.
到星期一.
真的打那個電話.
打電話就說說她的苦處.
聽到那個人都很慘.
一個小時內.
原來那間學校就打回來了.
他說妹妹.
王太.
你女兒的學額有了.
你再回來拿回來.
補辦手續.
後來才知道.
原來那天剛經過那個.
那個區.
那個是一個區監督.
剛剛有事經過那裡.
看到這個事情.
之後就背著卵告訴他們.
那些辦事的人聽.
你們要幫這個小女孩.
他們會不會不聽呢.
他們不會不聽的.
因為這個是什麼人告訴他們的.
是那個監督告訴他們的.
主宰他們的工作.
親愛的兄弟姐妹.
第一樣.
當我們說一個非常任務.
這個非常任務.
之所以非常是因為.
交這個工作給我們那個人.
是一個非常的身份.
第二.
主耶穌交代的是一個什麼任務呢.
我可以歸類三樣東西.
三樣東西就是這個任務的內容.

$^{281}$第一樣任務的內容就是.
你要做神的服事者.
剛才讀經文用的字眼.
叫執事.
大家不要誤會了這個執事.
是教會的那個執事.
城會好像用別的名字也不一定.
不知道叫什麼成員.
那個任務是其中一樣.
要他做一個執事.
其實只是一個服事者.
一個服務員.
還查那些聖經的字典.
強調這件事不是高級的.
是低級的.
我真的覺得.
主啊 你這麼辛苦救我回來.
你給我這麼不高級的工作.
弟兄姊妹.
我們去完成主給我們的任務.
其中一樣東西.
我們要習慣去看自己去服事神.
不要看自己太高.
可能這個不適合很多人聽也不明白.
但這個的確是主耶穌和保羅所說的.
但你想想.
保羅是什麼人士.
Walkman Horse.
他是很年輕的.
但是聖經說他.
在公會裡面一個成員.
還可以出名定案的意思是什麼.
就是簽名.
我以前的工作.
做財務主管那些.
就是我的簽名很重要.
公司千里追蹤我.
就是要我.
很重要你快點回來簽.
我記得有一年.

$^{321}$我就要教兒童聖經班.
公司就因為我緊要.
就是姑且這麼說.
有什麼事我要打電話給你.
那個時候是說30年前.
我竟然不聽.
那就慘了.
那一次後來被公司問責.
為什麼不聽電話.
很重要.
有個很大件事要交賬.
要付款.
要出貨等等.
很害怕.
我再說.
我們看自己是什麼身份.
我們只不過是一個僕人.
服侍的.
是主教給我們.
這個這麼重要的任務.
第一樣.
我們的心態看自己.
不要看自己太過高.
保羅是一個這麼重要的人.
當他信了耶穌之後.
他能夠放下自己.
公會成員的身份.
一個明日之星的身份.
你幻想一下.
有個參議員.
或者有一個眾議員.
30歲左右已經成名.
然後成為明日之星.
一下子他不做那件事.
他要跟隨耶穌基督.
他做耶穌基督的僕人.
第二樣主耶穌要告訴你.
你要做一個見證人.
什麼見證人.
做過見證人沒有.

$^{361}$想不想做.
做過沒有.
見證人很重要的.
還有保羅是做.
耶穌基督復活的見證人.
這個不是每個人都有資格做.
我沒有資格做.
我想我們在座之中.
應該大多數人都沒有資格做.
因為在第一世紀都很難找到這些人.
第一世紀他們都要猛金去掃票.
這些人符合這個資格.
才可以被推選為候補仕徒.
所以聖經裡面只有兩個做候補仕徒.
第一個落選.
選到的就是馬提亞.
所以他們很知道.
這個復活的見證人很重要.
也都很難得.
證人對一件案件.
關係重大.
你知道嗎.
在多倫多發生一個交通意外.
那時候還沒有興行車紀錄儀.
結果應該沒有負法律責任.
但是你知道交通意外的特點.
就是很多人說謊.
一剎那覺得真真假假.
假假真真沒有人知道.
又沒有紀錄儀.
那怎麼辦.
結果對於本來不應該負責的司機很不利.
好了 沒辦法.
突然有一天警察局打給司機.
他說有好消息.
竟然有個人主動聯絡警方.
說他可以做證人.
真是天大其事.
竟然有個人那麼主動.
然後他證明了之後.

$^{401}$警方就說決定不會起訴你.
你沒有這個法律的責任.
我就是那個受害者.
這個真的很切身.
我很記得 證人啊 證人啊.
神派你做證人真是太好了.
神為我預備了一個樂於助人的證人.
弟兄姊妹.
主耶穌基督復活了是事實.
但是能夠做復活的見證人.
可能只是那一班真的認識主耶穌.
真的知道發生什麼事.
如果他們是那個見證人.
又沒有義氣不去那怎麼辦呢.
就沒有見證人了.
所以弟兄姊妹.
我們如果是真心信了主耶穌.
跟了祂 知道主耶穌復活是一件什麼事.
真理 福音是一件什麼事.
主耶穌對你多好的事.
你又親身經歷的話.
如果我們那麼沒有義氣.
我們不肯去宣告.
我們就是連做證人都不肯做.
沒有討人厭的不做證人.
屈枉正直.
就是這樣的人搞成.
第三樣主耶穌跟保羅說.
你要做什麼人呢.
信息的傳遞人.
就是你看到什麼.
主耶穌說什麼.
那些道理你就說出去.
難不難呢.
哪有難呢.
養過撩哥沒有.
你沒有養過撩哥.
值得養一隻.
不是太貴的話.
我養過的人不用錢的.

$^{441}$他自己飛來的.
很聰明的.
我很記得.
我那時候還上學.
我教他叫我阿林.
我中文名字有個林字.
那他早上就很聰明的.
阿林 阿林.
什麼都行.
我又教他早晨.
阿林 早晨.
很聰明的.
我覺得撩哥比小朋友還乖.
你有沒有碰過很多小朋友.
其實不叫人.
不要說早晨.
叫你的名字都不叫.
題外話.
我做父母的時候.
我現在都是父母.
不過我的子女現在是成人.
我大女兒結婚了前年.
我很留意一樣.
她小時候.
要叫人.
叫人的名字叫早晨.
這個是福音秘訣.
你不知道的.
原來介紹自己的名字.
說到人的名字.
是令人深刻印象的秘訣之一.
你看主耶穌.
我們聖經裡面的神.
他怎麼介紹自己.
在舊約跟摩西說.
你要介紹我耶和華.
你不認識其他神的事也好.
你要認識神的名字.
名字就代表很重要的東西.
回到信息的傳遞人.

$^{481}$有些事你知道了.
你要將那個內容告訴別人.
這個是非常之重要.
你說那隻鳥哥.
當然了.
幾句東西當然學到了.
其實我也花了很多心機.
教他.
講完講完反反覆覆.
兄弟姐妹.
我們去查經.
我們又來聽道.
或者我們自己又會看看聖經.
又會讀書.
又聽視頻.
學了這麼多東西.
做什麼的.
就是被人覺得你好厲害.
感覺良好.
我相信不是.
是主耶穌.
是神的心意.
被我們裝備了上來.
就將一些重要的信息.
跟別人交流.
分享.
可不可以.
這個信息的傳遞人.
按著主耶穌和保羅說.
就是他這樣說.
你就要這樣去傳.
那就是說.
你要原汁原味.
將聖經的真理.
跟別人分享.
所以傳福音.
既容易.
又不容易.
為什麼呢.
那個信息沒有改變.

$^{521}$信息基本的框架是一樣.
框架就是.
這是一個因點.
是因為人有罪.
不能夠救自己.
而神是公義聖潔.
會追討罪.
神因著我們.
面對不到這個困境.
神就安排.
這位的救主耶穌基督.
去甘心樂意為我們犧牲.
凡是肯樂意.
願意接受.
主耶穌已經完成的救恩的.
憑著信心的.
就得到拯救.
是不是很簡單.
我幾十秒已經講完了.
但是那個內容再演繹.
進深過程.
聖經怎樣去講解.
我們繼續去學的時候.
聖經就要我們忠實地.
將這個內容講出去.
忠實地去介紹.
昨晚我正在吃飯.
突然有個電話打來.
我就扭轉了.
不認識的.
我就不聽了.
接著就把手機刪除了.
以前我都不會.
昨晚不知道做什麼事.
手癢.
接著他就留了個錄音.
他是一個問福音問題的朋友.
我真的想打自己兩巴掌.
看看怎麼辦.
我有兩個電話.

$^{561}$我兩個電話.
講了一個小時之後.
飯還沒吃的就算了.
接著還可以.
(音量問題).
每個20分鐘左右.
你聽兩個.
大家有興趣都可以聽的.
在芝加哥的電信中心.
那個知多點網上教室的頻道.
我想講的就是.
那個訊息沒有改變.
但是我們需要中心去傳揚.
不可以說一點不說一點.
說一點不說一點的福音.
就有什麼危險呢.
最近我教的跟進栽培進深課程的班裡面.
有一個同學就說.
老師有很多人在教會是基督徒.
但不像基督徒.
那個學生很激動.
我偶爾都要輔導一下他.
自從我上課為止.
但是他的心是好的.
我就告訴他.
其中有一課我說什麼呢.
我說有很多人是要澄清.
他的福音的觀念.
有很多人上教會.
為很多其他不等洗的原因.
當然神可能藉著這個.
他原來是不等洗的原因.
帶他認識神.
但是他不應該因為這個.
不等洗的原因.
藉到永恆.
都是為這個原因上教會的.
是不是.
我認識一些人很稀奇古怪的原因上教會.
我的兒子很頑皮.

$^{601}$我希望教會教好他.
又可以.
教會又可以.
在神的大能裡面.
他真的服.
有些人找不到丈夫太太去教會.
這麼不正確的目標來教會.
神都糾正他.
神都幫助他.
最後他就知道自己的原因上教會.
跟隨耶穌基督.
但是他不應該永遠都停留在那裡.
有不少人在福音的混淆.
不明白的裡面.
需要我們愛心忍耐去關懷他.
你一腳.
一針拆穿他是沒有用的.
等神用方法拆穿他.
有一天他讀聖經.
他聽到.
神自己光照他.
讓他看見.
原來這個福音.
那個福分寶貝是什麼.
很多人都不知道.
包括昨晚那位朋友.
他是因為怎樣回教會追尋信仰.
你猜.
他說黃傳都.
我睡不到覺.
他的原因就是睡不到覺.
他覺得這個信仰幫到他.
很多人有很多原因.
主耶穌強調.
他給我們那個福音.
真理的信息.
我們忠實去傳遞.
當然十八節那裡.
就不要漏了一樣東西.
主耶穌說.

$^{641}$是我要猜你到他們那裡去.
所以在主耶穌給保羅的.
最開始的吩咐就是.
有三種人你要做.
服事神的人.
見證人.
和忠實要傳遞這個信息的人.
而且你不可以只是遠距離的參與.
你不可以遠距離看一下這件事就算了.
你要自己親身去.
in person的.
投入的.
親自的去.
所以弟兄姊妹.
我們當中信主的.
鼓勵你.
做很多事都好.
就是要投入才行.
不要經常做一個外圍的人.
有些事做外圍的人.
做觀眾不過癮的.
不好玩的.
你一定要投入裡面才開心.
服事主更加.
投入在裡面那種興奮.
無可比擬.
相不相信.
你投入過你就知道了.
什麼叫第一個人信耶穌.
我聽過一個七十多歲.
做了這麼多年pharmacist.
退休了.
他分享他第一個人信耶穌.
多興奮.
我到現在都好記得.
跟著他就.
好想找機會和別人分享.
這個是他的心境.
有沒有什麼果效呢.
有沒有果效.

$^{681}$主耶穌和保羅說.
你其中的任務就是要使那些人眼睛得開.
好多人的屬靈眼睛是看不到.
屬靈的真相.
重要的東西.
你要負責去傳揚這個訊息.
令到那個人的眼睛得開.
屬靈的眼睛打開又怎樣.
就會知道.
信耶穌是寶貴.
未看到之前呢.
就好多時候我們會好灰心.
覺得他沒有反應的.
他不理睬我的.
我還好記得我為我爸爸的祈禱.
去信耶穌.
二十八年.
我爸爸信耶穌的時間.
是我做了傳道都一段的時間.
我曾經因為爸爸未曾信耶穌.
我都在掙扎.
會不會不是太好見證啊主.
不如等我爸爸信了.
我答應你我不逃跑.
我到時就奉獻好不好.
神有祂的時間.
神就真是好信實.
帶領我爸爸信主.
那些人眼睛得開好重要.
第二樣呢.
他們替暗投明.
由黑暗裡面歸向這位真神.
由撒旦的權下歸向這位愛他的主.
這個不是寓意故事.
今天有時好多故事靈異化了.
被我們誤以為.
屬靈的爭戰好像是一些.
好外圍的事.
不是真的那麼HIT到.
主耶穌其實在他的生平裡面.

$^{721}$都承認一樣東西.
真實性.
他是有一定的權勢的.
只不過他是被摔倒在地上.
他不可以對我們屬神的人.
令我們驚怕.
害到我們.
斷沒有甚麼能害我們的.
在路加夫的十章.
十九節那裡記載.
斷沒有甚麼能害我們.
嘩! 好厲害.
還說我們有權並踐踏蛇和蠍子.
撒旦的影響是很真實的.
如果你有傳福音.
你就會經歷得到.
你知不知道在現在還有些情況.
就是有些人想去除偶像是除不到的.
打擊偶像都不難的.
好多年前.
我就去巴西那裡做短宣.
在巴西做短宣的時候.
就留意到.
收到一個消息.
在巴西的聖保羅姆基院.
有一個叫聖約瑟士.
在聖約瑟士那裡.
有一個姐妹的媽媽信得住.
她就是.
她身邊有些基督徒幫她除偶像.
打擊不難的偶像.
打擊不難的那些東西.
他們很害怕.
那個姐妹的媽媽更害怕.
就是死了.
自己就決志信主除偶像.
誰知除不了偶像.
她還拖的.
是不是.
她這麼有法力保護自己.

$^{761}$她為什麼不會再進攻.
她真的很害怕.
我們短宣隊剛去到.
她就說請我們有幾個進去.
其實我們都不是什麼大隻佬.
你指望我五隻五.
這樣打擊那個偶像都很難.
我們靠著的是主耶穌.
我們進去就切切禱告.
第一件事就搞清楚.
她媽媽信得夠清楚.
你是不是真心真意.
有沒有藏起一些重要的東西.
秘密符咒 秘密偶像等等.
全部都供出來.
放在那裡我們祈禱.
還記得.
那時那班人裡面.
我最年輕.
你來.
給個錘你 戴上手套.
然後為你祈禱.
他們站在後面.
我就說好啊.
我祈禱.
我記得那個感覺是.
那個偶像一錘就粉碎.
我不用補飛.
那你說.
我練了什麼武功.
我自己都奇怪.
我自己最後的結論就是.
很明顯是神要粉碎她.
要彰顯她的大能.
很讚美神 榮耀神.
真是主要彰顯她的能力.
所以她的媽媽很開心.
知道主耶穌才是真神.
所以保羅.
他要傳揚一個訊息.

$^{801}$他要告訴亞基百王.
我不是普通的人.
屬靈或者宗教狂熱.
現在被人糟蹋.
我是奉耶穌的名去傳揚這個真道.
面對這些的事.
但是我們的主是真的.
是使人在黑暗裡面歸向光明的.
Amen.
不單止這樣.
好了 主的能力很大.
但是有一樣更加重要的是什麼呢.
就是祂的赦罪的能力.
是一種改變的大能.
這是保羅去見證到的一件事.
傳道者最大的福氣就是.
將這個生命的見證.
給過這個領受福音的人.
這個領受福音的人一生會改變.
就是他從前在罪裡面.
核制裡面.
有些人被罪的引好控制.
有些人被罪的習慣控制.
和家人的關係破損 破壞.
有人可以十幾年.
沒有和哥哥說話.
他信了之後才開始和哥哥說第一句話.
差點都忘記了上一次和哥哥說什麼.
是主耶穌基督改變人的生命.
使那些人罪得到赦免.
得到真正的平安.
很有趣 保羅說到這裡.
突然提到說.
這些人有一個特點.
是會與其他信的人同得基業.
我覺得要給我們留意.
我們看信仰是很個人的.
但有沒有想過.
信仰的美麗不只是個人的發揮.
信仰的美麗就是當我們在一起的時候.

$^{841}$你幻想一下.
如果我現在分享.
只是對著我太太.
或者對著Phoebe聊天.
是好 但只是很有限的層次.
但如果我們有一班屬神的人.
一起去聚集 一起去見證.
一起去唱詩 讚美神.
一起去享受那樣東西 好東西.
這樣才好.
所以你有好東西.
你有好東西吃.
你不要只在家裡吃.
你不要只收起來.
最好是怎樣.
一起享受 一起分享.
屬靈的福氣更加.
這樣分享的時候.
越分享越享用的時候.
會越多起來.
所以五餅二魚其中一個美麗是什麼呢.
就是在它傳下傳下傳下傳下就多.
這樣才美麗.
不是說突然間滿了.
整個郊外都很多東西吃.
你們來拿吧.
不夠它拿著那個餅.
給不給 給 只不給 給 只不給.
不是 給 給.
在這個過程裡面.
就將這個福氣傳開.
所以保羅去看到主耶穌的心意.
原來神的心意計劃.
是要他做這樣的服事.
最後保羅是怎樣回應的呢.
我只說一點.
主耶穌跟他說.
叫他這樣做.
二十節就說.
他先在大馬色.

$^{881}$後在耶路撒冷和猶太傳地.
以及外邦.
同樣使徒行傳一章八節.
不是這樣說的嗎.
從耶路撒冷.
猶太傳地和撒瑪利亞直到地極.
為什麼好像不合的.
起點不合的.
保羅不是應該飛一個旅程.
回到耶路撒冷開始嗎.
究竟應該怎樣理解那段經文呢.
那段經文強調一件事.
或者保羅領會了一件事.
你的起始點.
就是你現在這個起始點.
所以你不用飛去耶路撒冷.
你不用等Mark Sir去旅行的時候.
才開始在耶路撒冷開球.
做福音工作.
你現在在哪裡.
在芝加哥.
你的起始點就不是大馬色.
你不是耶路撒冷.
你是在芝加哥.
你和我的起始點就是在芝加哥.
或者有些人在網上.
我不知道你在哪裡聽.
如果你現在的起始點.
你在紐約.
你在香港.
你應該在你現在的起始點.
香港就去傳福音.
就去見證耶穌基督才對.
這個是主給我們的優先.
原來文字不是說做完了.
耶路撒冷才去猶太傳遞.
那裡用的原來文字是一個連接.
楷的字.
希臘文.
是同步同步同步的意思.

$^{921}$一起去做.
所以主耶穌基督.
很清楚將這個訊息.
這個非常的任務告訴他.
這個任務是什麼.
我們是一個報人.
我們是一個見證人.
我們要將真理認真去傳開.
這個任務交給我們.
我們的反應是立即去做.
我們在這個起始點去做.
你願不願意這樣做.
是真還是不是.
還是聽著主耶穌的說話.
最後.
保羅在大馬式呼召回應神.
我們在這裡去呼召回應神.
我給大家一些挑戰.
UV有什麼外展的工作.
你知不知道.
你會不會禱告去參加.
第二.
我剛剛才知道你們有一些弟兄.
有一個西班牙教會.
他們派食物去接觸一些華人.
你有沒有想過.
一有機會.
如果這些機會.
是可以有任何人去關顧的.
你會不會準備.
第三個挑戰.
每年你知不知道.
芝加哥短宣中心.
就是我服事的機構.
九月有一個騙傳運動.
很多時候有一天.
都會在你教會那裡做訓練.
但其他日子都會有機會傳福音的.
你有沒有想過.
祈禱一下.

$^{961}$會不會你可以參加一下.
按著你的能力.
可以的.
我做一個僕人.
我做一個見證人.
我做這個福音的訊息傳遞者.
願主使用你.
我們一起祈禱.
主耶穌基督懇求你的恩典.
激勵感動我們.
使我們再不是一個觀望.
坐在看台.
看別人來事奉.
看別人來為福音努力.
很欣賞拍掌的人.
主一樣.
我們都是在裡面敬祝.
為到神的榮耀.
神的賞賜.
神的喜歡.
來服事.
我們將我們眾人交在你的手.
主啊.
你得掌我們的心.
你將你的愛.
那個使命.
這個非常的使命.
交給我們.
願你恩代.
最聽我們的禱告.
奉靠主耶穌基督的聖名.
阿們.
\newpage



\section{}
\label{sec:bEexjAddGwM}
\textbf{比對研究 9: 巴力宗教 2 - 陳崇基牧師}
\newline
\newline
連結: \href{https://youtube.com/watch?v=bEexjAddGwM}{\texttt{https://youtube.com/watch?v=bEexjAddGwM}} ~~~~ 語音日期: 2023-03-11
\newline
\newline
\hyperref[sec:z2UTnfWH0W8]{\small{< < < PREV SERMON < < <}}
~
\hyperref[sec:index_chronic]{\small{[返順時目]}}
~
\hyperref[sec:index_scriptual]{\small{[返順卷目]}}
~
\hyperref[sec:2q17d1G4vqo]{\small{> > > NEXT SERMON > > >}}
\newline
\newline
$^{1}$我們來看看巴力的循環.
Bell Cycle.
正如我所說,六塊的禮板裡面.
很清楚地講述巴力的故事.
到底巴力神明是一回什麼事呢.
第一和第二塊禮板.
主要講什麼呢.
一個至高者叫做El.
El的翻譯就是神.
希伯來文也常常叫神是El.
El是一個名號,總之是形容神.
在Bell Cycle裡面,巴力的神話裡面.
首先開始有一個至高神El.
因為一些原因沒有解釋.
他偏愛一個海神,叫做Yam.
Yam.
聽到的時候就知道希伯來文是Yam.
Yam.
El很喜歡Yam.
但就允許巴力這個比較年輕的神明.
去挑戰Yam這個海神.
結果故事發展下去的時候.
巴力成功挑戰了Yam.
然後他就贏了眾神和眾人的王權.
贏了之後.
禮板的記載就是巴力要求建立自己的宮殿.
在哪裡建立呢.
當然是在聖山上面.
就是這個薩峰山.
巴力就在這個薩峰山上面.
去建立他的Temple Palace.
他的寺廟和他的宮殿.
所以這個Mount Saffong.
在聖經裡面出現過好幾次.
Saffong的意思其實就是北面的意思.
支北的意思.
Saffong就是北面的山.
所以音譯就是薩峰.
但其實是北面的山.
所以你剛剛看到地圖就知道.

$^{41}$就是在北面.
就是最北面的地方.
就是大迦林最北面的地方.
就是Mount Saffong.
就是眾神聚集的地方.
巴力在眾神中間.
就是稱神的地方.
Mount Saffong.
然後.
過了兩塊泥板.
到底三四塊泥板呢.
就是描述巴力如何鞏固自己的皇權.
以及如何完成.
他的寺廟的建築.
裡面就提到一些細節.
說到他如何在這個Mount Saffong上面.
建立了一個很偉大的宮殿.
然後他就接受敬拜.
到最後兩塊泥板.
兩塊泥板五和六.
就是故事的結語.
在這裡就講述巴力.
為自己的皇權.
去到辯護.
作守衛.
守護自己的皇權.
特別就是有一個叫做Maud.
這個Maud.
的冥神.
Maud其實就是Guard of Underworld.
就是在死人的地方.
或者下面陰間的地方.
那個死神.
翻譯的譯音.
就是Maud.
我們可以譯作Maud.
Maud.
這個名字也是.
Maud.
就是希伯來文的死亡.

$^{81}$所以.
烏加列文.
也是用Maud這個字.
來形容死亡.
死亡的神.
就是Maud.
他就挑戰巴力.
就殺死了.
巴力就死了.
死了之後.
阿勒特.
就是他的情人.
他的妹妹或者情人.
就為他復仇.
為他復仇之後.
他就殺敗了Maud.
Maud死了之後.
就令巴力又再重生.
又再復生.
所以這個就叫做.
巴力循環.
所以你想起.
聖經裡面有沒有一些地方.
跟這些是相似的呢.
我們說到.
聖經跟這些背景相似.
不是說聖經抄襲他們.
不是說聖經去到.
去到.
copycat.
去到.
重複他們.
要說的話.
而是某些東西.
在文化裡面已經根深蒂固.
一些的思想.
一些的概念.
譬如巴力神.
他的死.
然後他再復生.

$^{121}$然後下年又死又再復生.
就好像巴勒斯坦地的氣候一樣.
所以有些東西.
已經根深蒂固了.
所以不是說誰抄誰.
而是去到.
大家.
巴力這個神話.
聖經也是對當時的背景.
對當時的文化.
作出回應.
所以如果我們看了巴力循環.
巴力神話.
一些的內容.
我們很快地總結來說.
可能你也會想起一些聖經.
是神.
和巴力神明.
去爭鬥的.
一些的經文.
剛才也提過一點.
這些就是我們簡稱.
可以叫做podemic.
去爭辯.
這裡很簡單地列出某一些經文.
譬如詩篇29篇.
詩篇68篇.
還有遼王記.
17至19章.
這裡就是以利亞的故事.
甚至我們繼續延續下去.
以利沙的故事.
都是和巴力來抗爭.
所以我們認識的背景.
一定是對我們有好處的.
為什麼這個故事.
這些經文.
會這樣描述.
有沒有它自己的意思在裡面.
而那個意思.

$^{161}$一定要在背景裡面才能夠掌握.
如果我們.
第一步不去掌握背景的話.
我們就很容易將.
我們現在有的觀念.
有的字眼或意識.
我們就會先讀進去.
所以我們先.
理解了第一步.
他們怎樣去到.
初代的讀者.
初代的作者.
在巴力的神明的影響下.
他們怎樣去.
理解.
和怎樣去到.
為耶和華作爭辯.
以致我們今天能夠知道.
意義我們今天怎樣.
和不同的.
反對信仰.
或者反對.
聖經信仰的.
思潮可以作一些辯識.
和作一些爭辯.
巴力神是怎樣的呢.
這一塊.
大的石碑就是在烏加利.
這個地方挖出來的.
這裡就.
畫出巴力.
是一個怎樣的神.
巴力你看到.
這裡就是一個人像.
他最重要的地方.
當然是他的右手.
右手拿著一個棒.
一支棍.
或者一個.
武器.

$^{201}$就在他的頭上.
右手拿著.
他就很準備地.
隨時可以攻擊.
他的受敵.
然後他的左手就拿著什麼呢.
左手就拿著.
一個.
雷電.
或者有些解釋.
雷電.
和植物.
當然雷電和植物很有關.
因為和雨水有關.
所以拿著一個雷電.
電棒.
他的這塊.
這支電棒.
插在地上.
代表他的特性.
他戴著的.
就是一個牛角的.
一個.
頭盔.
然後.
他身上所穿的.
就是一個半.
到膝頭.
這麼長的.
一個.
裙.
就是一種戰士的.
裝飾.
這些的衣飾.
或者這些的.
代表.
就是告訴人聽.
巴力是一個年輕的.
很沒有力量的.
一個神明.

$^{241}$你看到他的腳.
不是直立的.
是向前走的.
換句話說.
他是很沒有力量的.
他可以向著他的受敵.
去作攻擊.
而他腳下面踩著的是什麼呢.
如果你隱約看到.
就是一些波浪.
所以他就代表著.
是剋聖水.
水的混沌.
大水等等.
就代表巴力.
其實是可以去到.
掌握著風雨.
掌握著雷電.
掌握著雨雲.
然後將雨刺在地上.
以致地上.
能夠豐衣足食.
他也有權能.
得勝了大水.
所以因為巴力的緣故.
這個宇宙.
就會有.
四季的循環.
就會有命令.
不會有混亂.
混沌已經被他戰勝了.
所以巴力.
就可以.
保證.
可以證實.
可以堅立.
這個世界有命令.
所以.
拜巴力的人.
其實是做什麼呢?.

$^{281}$我們就和當時的社會.
很有關係.
因為當時是農業的社會.
我們看看一些.
巴力神話裡面的.
片段.
這些片段.
看完之後.
我們就會想想.
到底聖經裡面有沒有類似的記載.
而那個記載的手法.
不單止是像.
也有很多分別.
本質上有不同的.
我們看看.
一些.
在巴力神話.
或者巴力的一些文獻.
裡面怎麼形容.
巴力.
其中一個.
巴力的.
名號.
當然以前的神明.
有很多的名號.
耶和華也是一樣.
巴力也是一樣.
巴力有很多不同的名號.
其中一個最多時會提到的名號.
就是叫做.
駕運者.
Cloud Rider.
或者.
他是駕駛著天雲的.
那一位.
這裡我就截錄了出來.
就是Bell Cycle的其中一段.
這段的上下文講什麼呢?.
就是講到巴力.
戰勝了.

$^{321}$殺敗了這個Yam.
這個海神.
混沌的神.
God of Chaos.
那怎麼形容呢?.
就是巴力.
拖走.
這個Yam.
將他肢解.
打完一輪之後.
巴力贏了.
毀滅了這個武士.
何川.
Yam.
Yam有時候叫做River.
然後就有一些.
大的聲音去讚美.
這個巴力.
或者去高舉巴力.
就說驅散.
驅散.
駕運者.
Rickhaw.
或者Rickhaw.
這個就是烏加利文.
英文翻譯成.
讀音的意思.
好,這個就叫做.
Cloud Rider.
駕運者.
另外一段就是Epic of Arkhan.
這個史詩.
這裡講到什麼呢?.
上下文就講到.
巴力死了.
死了之後就.
發生什麼事呢?.
其他的眾神.
因為巴力死了的緣故.
所以就沒有雨水降下來.

$^{361}$他們就在旱災的時候.
他們就去哀歎.
為巴力作哀歌.
他說.
七年巴力不在.
Seven years bail is absent.
八年.
那駕運者.
換句話說,沒有了另外一個意思.
意思就是說他也不在.
那駕運者不在.
是同一個字.
Cloud Rider.
他說沒有金路下來.
No dew.
也沒有下雨.
沒有深淵的漩渦.
沒有人搞動深淵.
也沒有歡迎巴力的聲音.
就是雷電的聲音.
沒有雨水的聲音.
等等.
這裡給我們看到.
另外一個章節.
就是講到Cloud Rider.
這個就是巴力.
第三段.
在另外一個神話裡面.
叫做Bear Fathers Above.
就是巴力生了一隻牛.
這個神話裡面.
有一段是講什麼呢?.
他說好消息.
我們要歡迎這個好消息.
歡迎他.
就是巴力.
巴力歡迎這個好消息.
巴力是大菌的後裔.
大菌的後裔.
因為巴力生了一頭牛.

$^{401}$一頭野牛.
這頭野牛是誰生的?.
是誰給的?.
是駕運者.
Rider of Clouds.
Bale the Almighty rejoices.
大能者巴力就歡喜了.
為什麼會說.
這麼奇怪.
為什麼會說他Fathers Above.
就是生了一隻牛.
這個就是.
當時農耕的背景.
就一定知道.
因為牛就代表大力.
生產力.
農作物的豐裕.
就是有頭牛.
所以說巴力.
就代表他生了一隻牛.
就是他能夠控制.
這個生產的生命.
所以就說他生了一隻牛.
但最重要的是.
他怎麼叫巴力.
Rider of Clouds.
所以他是駕運者.
他是去操縱天運.
那些黑運.
當你看到黑運來的時候.
有雨落.
是誰在巴力的宗教.
或者巴力的思想在迦南的文化裡面.
就是因為巴力.
去控制著這些東西.
他可以有天.
他可以帶天運.
所以因為這個緣故.
他也是那個.
在地上有雨水的神明.

$^{441}$在這裡另外一節錄一段.
就是Epic of Cota.
又是另外一個烏加利的史詩.
裡面提到.
又是汗災.
去到想找雨水.
他就說要擦汗地面.
就是想有巴力的雨水.
田地想有.
至高者的雨水.
這裡另外一個巴力的名號.
就是至高者.
他說.
在大地裡面多好.
就是巴力的雨水.
為了田地的緣故.
有多美好.
就是至高者的雨水.
這些就像寫詩.
你看到像詩篇一樣.
都是這樣.
一句一句.
然後一句對一句.
然後A句又對B句.
巴力的雨至高者的雨.
我們再看.
最後一段.
就是Bell Cycle.
就說到巴力.
得勝之後.
他就在聖山上.
去建立他的.
就是.
這個宮殿.
就說到.
巴力建立廟宇的時候.
他在廟宇裡面.
就想打開一個洞.
在宮殿裡面.
要開一扇窗戶.

$^{481}$巴力就在.
雲中打開裂口.
巴力就發出.
聖潔的聲音.
巴力就述說.
他嘴唇的事.
他是一個說話的神.
說什麼呢.
就是有聖潔的聲音.
以致大地會震動.
有地震出現.
所以呢.
他就是那位.
教運者.
將黑雲與雲帶來.
然後他來到的時候.
有很大的聲音.
然後這些聲音要使大地可以震動.
這個就是.
那個神話的意思.
因為他在雲層那裡.
有一個缺口.
他可以打雷下來.
就說到這些開口.
這些窗戶等等.
天雲裡面有些洞口.
巴力可以打雷.
降雷下去.
降火下來.
然後地帶震動聽見他的聲音.
有雨水出來.
\newpage



\section{}
\label{sec:2q17d1G4vqo}
\textbf{Becoming A Healthier Church 6: 1 Cor 14 - Sound Teaching  - Rev. Mark Chan}
\newline
\newline
連結: \href{https://youtube.com/watch?v=2q17d1G4vqo}{\texttt{https://youtube.com/watch?v=2q17d1G4vqo}} ~~~~ 語音日期: 2023-03-14
\newline
\newline
\hyperref[sec:bEexjAddGwM]{\small{< < < PREV SERMON < < <}}
~
\hyperref[sec:index_chronic]{\small{[返順時目]}}
~
\hyperref[sec:index_scriptual]{\small{[返順卷目]}}
~
\hyperref[sec:lL2aGhSKVWQ]{\small{> > > NEXT SERMON > > >}}
\newline
\newline
$^{1}$- So we have our pastor, Pastor Mark,.
who's gonna give our sermon today..
- Thanks, Annie, thank you..
Well, we continue to ponder on this important question,.
how to build, how to grow a healthy church..
Again, we were at studying the book of 1 Corinthians..
We were at chapter 12,.
and last week we studied the love poem of chapter 13..
So today, let us turn to 1 Corinthians 14, all right?.
So these three chapters are the Trinity chapters..
Three chapters together, 12, 13, and 14,.
addressing the use of spiritual gifts..
And actually, there was a bigger picture in mind..
Paul is talking about church as a human body,.
as an analogy, that we are all different parts.
of this whole body, building up the body of Christ..
And then 13, as you remember, the famous love poem,.
love is patient, love is kind,.
to help us to get a right orientation.
of how to use our talents and our gifts.
to build up the church..
And now we are at chapter 14..
And let me begin by reading the passage..
And there are two points in this chapter, chapter 14..
First is the importance of sound biblical teaching.
for growing a healthy church..
That's the first point..
The second point is the value,.
the value of sound biblical teaching..
As I have reminded you, what a church should look like.
and what a church should be,.
is not our preferences or our opinions..
It is not our suggestions that matter..
The end of the day is the word of God,.
the final authority of God's revelation, the Bible..
We come back to the Bible, we turn to the scripture.
and see what God has to say about our church,.
about church in general..
So 14, verse one,.
Paul is dealing with the use of spiritual gifts,.

$^{41}$especially the gift of speaking in tongues,.
speaking in unknown language..
14, one, pay attention..
Paul says, "Let love be your highest goal,.
"but you should also desire the special abilities.
"or spiritual gifts the Spirit,.
"the Holy Spirit gives to all of us,.
"especially the ability to prophesy..
"For if you have the ability to speak in tongues,.
"you will be talking only to God,.
"since people won't be able to understand you..
"You will be speaking by the power of the Spirit,.
"but it will all be mysterious..
"But one who prophesies,.
"who has the spiritual gifts of prophesying,.
"strengthens others, encourages others,.
"and comforts others.".
So Paul is making a contrast.
between the gift of speaking,.
well, how does the NLT translate it?.
Speaking in tongues, right?.
Or speaking in unknown language..
He's contrasting one set of spiritual gifts,.
speaking in tongues or in unknown language..
And then the other one is prophesying,.
speaking prophecy, you know, prophesying..
So from this passage,.
we understand not all spiritual gifts are equal..
Yes, the Holy Spirit will grant us different abilities.
or talents in the church,.
different members or different brother and sister,.
they have this set of gifts,.
or he has this set of abilities or talents,.
but not all spiritual gifts are of same importance..
Why?.
Now, some of the spiritual gifts are obviously stronger,.
according to Paul..
Why does Paul say that?.
Because the Corinthian Christians,.
they were desiring, they want,.

$^{81}$or they're yearning for some sort of spectacular gifts..
You know, speaking in tongues,.
or speaking in unknown language, right,.
mysterious language, seems to them,.
the Corinthian Christians,.
would be something they wanna strive for,.
they wanna pray for, and they want that kind of gifts..
And that sort of gifts probably look more spiritual to them,.
or look more supernatural to them,.
more spectacular to them..
But, you know, as Paul has emphasized.
throughout these chapters,.
the Word of God is the standard..
When we return to the scripture.
and see what the Bible has to say about spiritual gifts,.
even though we are not part of that charismatic tradition..
So what is Paul saying?.
Paul is saying, he has nothing against speaking in tongue..
He has nothing against that sort of spiritual gifts..
Some people will speak in tongue.
because the sovereignty of the Holy Spirit,.
God has granted them this sort of abilities..
But Paul seems to be doing something here..
He deprioritize, you know,.
try to put some of the gifts,.
especially speaking in tongue, to a lower priority..
He elevate some of the gifts.
the Christians in Corinth were not seeking,.
the gift of prophecy..
Now, what does it mean?.
Okay, let's continue..
Verse five, he says,.
"I wish you could all speak in tongues,.
"but even more, I wish you could all prophesy,.
"for prophecy is greater than speaking in tongues..
"Unless someone is here in the public service,.
"when someone is speaking in tongues,.
"someone interprets what you're saying,.
"so that the whole church will be strengthened.".
Now, as you can tell,.

$^{121}$he is definitely not against speaking in tongue in general..
But he said there was a greater gift,.
greater spiritual talent that has to be used.
when the church gathers, and that is the gift of prophecy..
Now, what is prophecy?.
If you remember, I told you before,.
prophecy is not weather forecast, right?.
It was snowing outside, right?.
Right, it was snowing outside..
Prophecy, in general, in the Bible,.
has nothing to do with weather forecast,.
whether it was snow at 3 p.m. today, 3.30 p.m. today,.
whether the forecast is correct or not..
It has nothing to do with it..
Prophecy, in the Bible, is not so much prediction..
It's about preaching..
It's about speaking the word of God.
according to the final standard of Scripture..
Authoritative preaching, you can translate that..
Prophecy is stronger according to Paul,.
because it trains your mind..
You know, when the questions in Corinth,.
when they gather, you can imagine,.
like, say, we are the church in Corinth..
Most of the people, they were looking forward.
to speak in tongue..
Somehow, this Holy Spirit touches them.
and give them the spiritual ability.
to speak in mysterious language.
so that they can draw closer to God..
But Paul is saying, by exercising those gifts,.
speaking in tongue, you are drawing closer to God.
only by yourself..
And that's why that sort of gift, speaking in tongue,.
is lower in priority..
What is higher in priority is that when the church gathers,.
preaching, biblical teaching, has the higher priority..
When you come to church,.
when you come to a worship gathering,.
Paul says you need to seek for the gift.

$^{161}$of biblical teaching,.
the gift of sound doctrine,.
the gift of training our mind..
Because some of the people, when they are preaching,.
when they're explaining the word of God,.
they're expounding on the word of God,.
everyone else would understand, right?.
See, for example here, I'm using a language.
that everybody can understand, right?.
Otherwise, your mind will not be trained,.
and you will not be able to understand it..
I am speaking a common language everyone can understand..
If you cannot understand, there is a radio to translate..
That's why we take this so seriously..
We prepare the script, we prepare the sermon,.
and we prepare some teaching,.
drawn from the Holy Scripture..
Because when the church gather together,.
Paul is saying we need to know.
and understand the will of God..
It trains your mind, doesn't it?.
Now, a question for you..
Do you think your mind, your thinking,.
matters to God?.
You think having an intellectual mind,.
a sound, rational mind, right thinking,.
about the teaching of God, about the world,.
about the Bible, about our lives,.
you think those things matter to God?.
They matter to God, right?.
They matter to God..
The ability to discern, the ability to think critically,.
the ability to examine things rationally,.
these things matter to God..
They matter to God..
So much that Paul is saying there's importance,.
utmost importance..
When the church gathers together,.
we need to understand the teaching of the Bible..
That's why he elevate one set of gifts over the other..

$^{201}$Speaking of tongues, that's nice, Paul says..
But this is for your own private fellowship with God..
Because nobody understand what you're saying..
When the church gathers, everyone has to understand.
what we are talking about, what we are thinking about..
So the ability to discern and ability to use our mind.
correctly, to think about things critically,.
these things matter to God..
These are not the monopoly of some of the scholars,.
some of the intellectualists among us, right?.
There are some PhD students here..
What they're talking about, we don't understand, right?.
(audience laughing).
What they're doing, no clue, right?.
But let me say this, right thinking,.
right thinking, in order to think critically.
and to discern things, truth and error,.
these are important things to God..
These are very important..
So not only biblical teaching,.
some biblical teaching is important..
Secondly, my point is, it is of value, all right?.
Everyone should watch this documentary on Netflix..
It's a horrible, horrible thing to watch..
I haven't watched it yet, all right?.
It's called "In the Name of God, Holy Betrayal.".
Now it's a very, it's a big hit these days.
because it was released March 3rd, only a few days ago..
It was released on Netflix..
It's a documentary series exposing the biggest cult.
in Korea, the Big Four..
Now, as I said, I don't have time to watch it yet,.
but I will watch it..
Have you ever watched it?.
Anyone?.
No?.
You're scared to watch it.
or you just don't have time to watch it?.
You should watch it, you should watch it..
Because this documentary series has many episodes..

$^{241}$It exposed the cult leaders in Korea..
It's very interesting, it's very intriguing.
that all these cult leaders,.
they were able to attract thousands.
and thousands of followers..
And according to this documentary series,.
this cult leaders in Korea, or even all over the world,.
because they plant churches, different groups,.
all over the world,.
they attract thousands and thousands of college students..
Can you believe that?.
I mean, listen, let me tell you one of the quote.
from this series..
One of the cult leaders used to say to his followers,.
and I quote, "Do you know who I am?.
"I bet you don't..
"Well, I am the Messiah..
"You don't need to search for God..
"If you can't see God, look at me.".
These are some of the things that the cult leaders.
used to say to his followers..
Thousands and thousands, tens of thousands of followers..
"Look at me, I am God.".
It's very interesting to me because.
I just cannot understand why this sort of teaching.
would attract such a following,.
even among college students..
Now, college students are supposed to be not so dumb,.
'cause they got into colleges..
They must have some intellectual ability to discern.
in order to get themselves into campuses..
It's just very intriguing to me..
Many, many of the followers of this cult.
in Korea and around the world,.
they would follow wrong teachings,.
obviously false teachings..
And I encourage you to watch this documentary.
to think with me, to ponder with me..
What is going on?.
Remember what Jesus said..

$^{281}$He said, "Love the Lord with all your heart.
"and all your mind.".
You got to use your mind when you follow Christ..
God expects us to use our mind..
And our mind with the right thinking,.
to think critically, to discern,.
is pleasing to Him..
It's pleasing to God..
To discern and to understand what is going on,.
what is true and what is not..
Christians, with all your mind,.
should be able to grow in this sense,.
to think logically, to think rationally,.
should make decisions wisely and critically..
This is all expected in our Christian discipleship..
Because our faith is not blind faith..
Our faith is reasonable faith..
It can be justified..
It's the right thinking that matters to God..
Now, as I read some of the basic information.
before I watched the documentary series,.
I quote some of the obviously wrong false teachings.
from all these cult leaders..
As I said, I find it very, very troubling..
Students, you know, college students.
would fall into those traps, right?.
Because if you think about it,.
if you have the mind to really study.
and to discern, to think critically,.
you quickly will come to some of the scriptures.
obviously stated in the Bible.
that you can quickly tell.
these leaders are not teaching sound doctrine..
For example, Matthew 24,.
when Jesus was talking about the last days,.
he told his disciples, "In the last days,.
"don't let anyone mislead you,.
"for many will come in my name,.
"claiming I am the Messiah,.
"and they will deceive many.".

$^{321}$See, if you use your mind and use your brain.
to just flip some of the pages in the Bible,.
you quickly can tell who is telling the truth.
and who is not..
How could this educate people.
by this kind of lies, right?.
I just cannot understand..
I mean, do you have an answer?.
I don't have an answer yet..
Now, for many people who observe.
this terrible things and horrible things happening.
in people who use the name of God to deceive people,.
some of the people who are observers,.
kind of outsiders,.
they probably will conclude that,.
okay, all right, they will just abandon faith altogether..
They will just abandon the Bible altogether.
or religion altogether in general..
Because, see, some of the cult leaders.
using the name of God, the name of religion,.
to deceive many people and do horrible things, right?.
They will conclude that this is no use..
I will just abandon God..
I will just abandon the Bible altogether..
But again, if you think about it,.
such conclusion is actually irrational, right?.
Because somebody misused the name of God.
doesn't mean God is useless, right?.
If you think clearly, I think the best way.
in observation of all these terrible things happening,.
the best way is not abandonment, but discernment..
It's not abandonment, it is discernment..
God is reminding you and challenging you to discern,.
to grow in your ability to think critically..
Now, returning to chapter 14 of 1 Corinthians,.
you notice there is a key term being mentioned,.
being repeated over and over and over again.
in the entire chapter of 14..
Let me give you an example, 14.12..
Paul is speaking about the gift of preaching,.

$^{361}$sound doctrine, and the gift of speaking in tongue, right?.
He elevated one over the other and put speaking in tongue.
to the least of the list of gifts..
The gift of preaching is much better..
Verse 12, he said, "The same is true for you guys..
"Since you are so eager to have the special abilities,.
"the spirit gifts, seek those that will strengthen.
"the whole church.".
He's talking about some of the talents,.
some of the gifts that has to be elevated.
because those gifts will strengthen the whole church..
The word strengthen or the verb strengthen.
repeats again and again, over and over again.
in the entire chapter, because Paul understands.
when we get together, we need to strengthen.
and build up the church..
Now, strengthening the word, many times in this chapter,.
literally, it can mean building up,.
building things together..
It's used all over the Bible, especially in New Testament,.
about the temple, the temple in Jerusalem..
Brick and mortar, people put things together..
Building up, strengthening..
Talking about palace, the royal palace,.
palaces of the king, houses that people live in,.
physical buildings..
He's using this word to describe when we get together,.
we need to use our presence and our talents.
and our resources to build up the church,.
to build like the temple, put bricks together..
When we come together to have fellowship,.
when we come together to worship together,.
to serve together, we are putting bricks on one another.
and build up the temple so that God can be.
in presence with us..
That is our goal..
And that's why there's value in understanding.
the better gifts..
The better gifts are those that will build up.
and strengthen the church..

$^{401}$You probably know what contractors do, contractors..
Have you ever done any renovation in your house?.
Repair jobs, right?.
Those are really, really tough jobs, right?.
I, in the past few years, because I got some major.
renovation in my house, I have to hire somebody.
to do the hot job, right?.
So I really respect those guys..
You know, they got their hands dirty.
and they got the blueprint and they build things for me..
Now, as you understand, as contractors,.
the building process is not an end in itself,.
but they always have the end goal in mind..
What contractors do is to build things.
so that when the final project is done,.
the people who live in that space,.
who use that space, will be benefited, right?.
They got the blueprint, they got the process,.
they got all the tools, and they know what the end goal is..
The end goal is to build a beautiful building.
and finish the whole project so that whoever use it,.
the users, the people who live in it, they will enjoy it..
That is the whole point of building the church..
So Paul is encouraging us to use your gifts.
to build the church..
Now, to summarize what the gifts are,.
I have a quick way to summarize it..
It's three Ts..
Your time, and your treasure, and your talents..
Remember these three words, all right?.
The first one is time..
What does it mean?.
When we gather, it takes time, right?.
There's no gathering of Christians,.
there's no worship gathering that would not take time..
You set aside a time, you allot this time slot,.
I'm gonna worship God with the people of God, right?.
I'm gonna show up and be present, bodily present, right?.
Time, your treasure, right?.
Giving generously to the ministry of God, right?.

$^{441}$To train up the people.
so that his people will follow his way..
Your talents, right?.
Now, I don't have to repeat all those stuff, right?.
Because we already talk about,.
we are gifted in different ways, right?.
Some of those are musically gifted,.
some of those who are gifted in some other ways..
We are all using the gifts, right?.
Time, treasure, and talents to build up the church, right?.
Urban voice, our church..
Let's build it so that God will get the glory..
God will get the glory..
I have some confidential pictures I wanna show you..
As I prepared a sermon, I found some of the old pictures..
As you know, Urban Voice UV,.
we launched our church as a church plan.
with just maybe eight or nine families together..
It was back in 2016..
So we were officially launched on the Easter Sunday, 2016..
So as I prepared a sermon,.
I found some of the old pictures, all right?.
So the old pictures were all taken.
when we first move in to this place..
When we first move in here, it had nothing..
It's just an empty box..
Who built it?.
People who have talents,.
people who are committed to build the church together..
And I wanna show you this picture, it's about six or seven..
And I want you to, well, who were with us.
in the very first Sunday, 2016 Easter Sunday?.
Who were here?.
Yeah, there are only a few, right?.
So praise God, these people,.
they were not here in the very first day,.
but they gradually, and God sent them here,.
and what a blessing to have all of you..
But some of these pictures, as I show you,.
you have been coming to this church.

$^{481}$to attend the worship service,.
to participate in our ministries..
Thank God for that..
You should know our environment a little bit,.
right, our facility..
So as I show you these old pictures,.
you tell me, where was this picture taken?.
Where, okay, exactly where this picture was taken?.
You ready?.
Where was this picture taken?.
Where?.
Where was the camera?.
Right, it was around here, right?.
So we took a picture all the way there..
Now it's changed a lot because we at this room,.
you know, this studio room, before COVID,.
before the pandemic, the room was not there,.
part of the building process, right?.
Next one..
Where was this taken?.
That's right, right here, right here..
Taken this direction, right?.
It's amazing to see how the builders, you know,.
build the church together, or build all this stuff, right?.
This thing is hollow..
There were holes in it, right?.
They built a frame, you know,.
I was there almost every single day,.
and I took some pictures, you know,.
just to document what God has been doing in our midst..
So this is the stage..
How about this one?.
Where was this taken?.
Yeah, it's the cafe area, right?.
Back then we had only this crappy, you know, cabinet, right?.
And this worker was doing the drywall,.
and before we put the wood panels on..
God built this church..
How about this one?.
Where was this taken?.

$^{521}$Any guess?.
It's outside, it's actually the same angle,.
the cafe before the wall, before the drywall was put up..
Right, so this wall right there..
It's outside, the picture was taken outside, right?.
It's amazing, it's amazing..
Every day I come here, you know,.
there's a process, more has been done, right?.
How about this one?.
Any guess?.
Where was this taken?.
Where was this picture taken?.
To what direction?.
Ah, that's tough, right?.
That's actually in the cafe area..
If you stand outside the nursery door,.
you take a picture this way,.
not all the way to the end of the cafe,.
there's a dining table there, right?.
There's a bookcase there, right?.
Remember?.
Do you remember?.
Okay, all you pay attention to is the food, right?.
The food and the drinks..
Seriously, this is taken from the end to the other end.
in the cafe area..
How about this one?.
This is tough..
It took me a long time to figure out where was this taken..
I didn't even remember..
It was actually taken from where James is..
Yeah, it's right there, it's this corner..
Can you tell?.
It's this corner, that's the back door..
I want to show you this picture because,.
just like what Paul says,.
the church needs to be built up, you know, brick by brick,.
dry wall by dry wall..
But Paul is not simply talking about hardware..
He's talking about the people, the group, the church,.

$^{561}$the community getting together..
All these pictures, I add this verse in there, 1412..
In the same way, it is true for you guys..
Since you are so eager to have your special abilities,.
the spirit gifts, seek those that will strengthen,.
will build up the whole church..
A stronger church is the foundation,.
and the foundation is our mind..
Our mind is oriented to the Word of God,.
a better understanding of faith..
Disciples should be thinkers..
You should be able to discern, to think,.
to debate, to question..
A good worship service is not one.
that would make you tear up..
All right?.
Although I was worshiping and James was leading,.
it was such a great time we can praise our God, right?.
Offer our worship to Him..
But a good service, a good worship, a good church.
is not one that makes you tear up,.
but one that challenges you to think deeply,.
to think more critically, to think better..
Think in terms of God's Word..
Think along the line.
with the biblical teaching of the Bible..
Let's all pray together,.
and let's continue to live in the truth of God,.
and the truth is we live in His love..
Let's all pray..
Father, we give you thanks because you are the one.
who built this church..
And thank you for the fact that the Holy Spirit.
is the builder..
The Holy Spirit gives us in many different ways..
So as we come before you, as we conclude our service,.
we ask that you would bring us into your presence again,.
and know that we are being saved by your grace,.
by your love, and your love would transform us.
to change our perspective, change our orientation,.

$^{601}$so that our lives, our gifts, our time, our treasure,.
and our talents will bring honor to you,.
to build up the church, to build up this community,.
so that your name will be lifted high..
Thank you, and we praise you..
Christ, we pray..
(upbeat music).
\newpage



\section{12:1-11}
\label{sec:lL2aGhSKVWQ}
\textbf{Final Days of Jesus: Bethany (John 12\_1-11) - Rev. Mark Chan}
\newline
\newline
連結: \href{https://youtube.com/watch?v=lL2aGhSKVWQ}{\texttt{https://youtube.com/watch?v=lL2aGhSKVWQ}} ~~~~ 語音日期: 2023-03-20
\newline
\newline
\hyperref[sec:2q17d1G4vqo]{\small{< < < PREV SERMON < < <}}
~
\hyperref[sec:index_chronic]{\small{[返順時目]}}
~
\hyperref[sec:index_scriptual]{\small{[返順卷目]}}
~
\hyperref[sec:zbYaP2QF_GY]{\small{> > > NEXT SERMON > > >}}
\newline
\newline
$^{1}$Pastor Mark for our seminar today..
- Welcome, morning..
And we have some special guests here..
Let me welcome them again..
They were in the 10 o'clock Kennedy service..
They are joining us in this worship service as well..
And those are some of the special students.
from Moody Bible Institute..
Let's welcome them, all right?.
Let's welcome them..
(audience applauding).
They're here for class assignments..
They will ask some questions,.
so don't give them a hard time..
Yeah, okay..
Let's be friendly to them..
Okay, they are Moody Bible Institute..
Okay, they're from Moody, all right?.
Name we can trust..
They're from Moody, all right?.
So let's see..
We start a new sermon series on Good Friday and Easter..
I think it's appropriate because as you can hear.
from the announcement, the Good Friday is approaching..
And I think it's appropriate and it is the best time.
to reflect on the last week of Jesus,.
the last days of Jesus..
So we turn to some of the episodes in the Gospels,.
Matthew, Mark, Luke, and John,.
and think about and reflect on what Jesus has done.
and what he's doing to our lives..
So as we begin the service, as we begin the sermon,.
I would like to do some quiz, all right?.
So again, some pop quiz, all right?.
Since we have some Moody students here,.
all the church people, you are not going to Moody..
This is the time to show them..
We know the Bible without going to Moody, okay?.
All right?.
(congregation laughing).

$^{41}$This is the time, okay?.
This is the time we're gonna show them..
We don't have to go to, well,.
we don't have to go to Moody.
and we know well the Bible, all right?.
So, okay, all right..
So let's do it, okay?.
Question one, which Gospel recorded the birth of Jesus.
like Christmas, which Gospels?.
A, Matthew and Mark, B, Luke and John,.
C, Matthew and Luke, D, the Gospel of Mark..
Which one?.
Well, talk to your neighbors..
C, who said that?.
Moody students?.
I want, yeah, yeah, I want a non-Moody student, okay?.
To answer this question..
Just show them..
You know, we know the answer..
Yeah, Matthew and Luke..
We know the Bible, all right?.
So it's getting tougher, okay?.
Quiz two, all right?.
Question two, which Gospel recorded childhood.
and years before Jesus' pulpit ministry?.
Matthew, Mark, Luke and John..
Which one?.
Huh?.
Can I have the,.
can I have the non-Moody student answer this question?.
Come on..
Come on, guys..
Well, no, no, no, no, no..
He went to Moody..
James is not qualified to answer the question..
Okay, here we go..
C, it's Luke..
Yes, yes, give her a hand, give her a hand..
(audience cheering).
All right, all right..

$^{81}$She answered all the questions..
No way, no way, all right?.
So, question three, which Gospel recorded crucifixion,.
death, burial and resurrection of Jesus?.
Which Gospel?.
Matthew, Mark, Luke and John or all four Gospels?.
What's your answer?.
Huh?.
Are you sure?.
That's your final answer?.
All right..
Yes, it's all four Gospels, okay?.
All four Gospels talk about the death.
and resurrection of Jesus Christ, okay?.
So pretty good, guys..
Pretty good, all right?.
Pretty good, but you can still go to Moody..
You get better..
Now, think about it..
There are some important days.
in the calendar of Christianity..
The Christian calendar will mark some dates.
as more important than others..
Either in practice or in theory,.
Christmas is one of them..
So that's why most of the churches.
will celebrate Christmas in one way or the other..
Good Friday, of course,.
Easter and some church traditions.
will celebrate, will emphasize Pentecost as well..
So these important Christian calendar events.
remind us what Jesus has done, right?.
The life of Jesus, the birth and his life.
and his death and resurrection..
And think about what modern people.
would like to hear about the story of Jesus..
How many of you have watched "Chosen"?.
"Chosen.".
Yeah, pretty good, right?.
Pretty good interpretation or rendering.

$^{121}$of Jesus' life..
And think about some other Jesus movies..
We modern people, modern audience,.
are most interested in some of the stuff.
that is not recorded in the Bible..
Like when Jesus went to kindergarten,.
did he study or he did miracle, right?.
And got all A's..
How did he relate to like Joseph and Mary,.
you know, his early parents..
Things like this really interests us..
In fact, in the last two verses of the Gospel of John.
says this, listen,.
this disciple, which is the author of the Gospel of John,.
is the one who testifies to these events of Jesus,.
has recorded them here in the Gospel..
And we know his account of these things.
is accurate, reliable..
Jesus also did many other things, many other things..
He said, "If they were all written down,.
"I suppose the whole world, the whole cosmos,.
"the whole universe, the vast universe,.
"could not contain the books written.
"about all the things about Jesus.".
He said that..
John is one of the disciples..
He probably heard of what Jesus did when he was young, right?.
What Jesus did before his public ministry, right?.
All the funny things that he did.
or he experienced and all that stuff..
You know, rabbi and students,.
they live together 24/7..
So there has to be a lot of different stories, you know,.
a lot of different things that can be said of Jesus..
But, you know, it's very interesting..
When we go to the Bible, the content,.
majority of them has to do with Good Friday,.
the death of Jesus and his resurrection..
Nothing else..
No more speculations..

$^{161}$Only focus or emphasis is right on his death,.
his dying and resurrection..
Now that sets apart Jesus from all other religious leaders,.
all other leaders in general..
You know, all the leaders, human leaders,.
they came to live a life so that we can see as an example..
We can follow their example..
They're good examples..
Think of the heroes of the country, of the nation,.
the founding fathers, right?.
The cultural heroes, all these great stories about,.
you know, human, you know, inspiring figures..
They're all, you know, living the life.
so that we can follow as example..
Only Jesus came to die so we can live..
That sets apart Jesus from all the other people..
That sets apart biblical faith, Christianity,.
from all other faiths..
And let's take a look at today's passage..
It's taken from the Gospel of John,.
and the story has to do with preparation.
for nothing but Jesus' death..
Again, the focus is on the death of Jesus..
Chapter 12..
Six days before the Passover celebration began,.
Passover again is around the time Jesus was crucified..
So it's one week before the Good Friday..
Jesus arrived in Bethlehem..
That is the home of Nazareth,.
the man he had raised from the dead..
Remember John, in the previous chapter, chapter 11,.
talk about the last miracle of the Gospel of John.
is raising Nazareth from the dead..
Verse two, a dinner was prepared in Jesus' honor,.
and Martha was there serving,.
and Nazareth was among those who ate with him..
He was just raised from the dead,.
and Nazareth was one of them joining the dinner..
Again, Martha, Mary, and Nazareth, they are siblings..
All right, so they are very close to Jesus, this family..

$^{201}$A big dinner is there,.
and then Mary took a 12-ounce jar of expensive perfume.
made from essence of nard..
Where's the nard from?.
It's from India, okay?.
It's not curry..
It's from India..
All the way from India and transport all the way.
through commerce in the Middle East..
It's a very rare product,.
very expensive product..
And Mary annoyed Jesus' feet with it,.
wiping his feet with her hair,.
and the house was filled with the fragrance..
"But Judah is a currier, the disciple,.
"who would soon betray," Jesus said..
"That perfume was worth a year's wages..
"It should have been sold,.
"and the money given to the poor.".
Not that Judah cared for the poor..
He was a thief,.
and since he was in charge of the disciple's money,.
he often stole some for himself..
Jesus replied, "Leave Mary alone..
"She did this in preparation of my death, my burial,".
and following resurrection..
"For my burial, you will always have the poor among you,.
"but you will not always have me.".
When all the people heard Jesus' arrival,.
they flocked to see him and also to see Nazareth,.
the man Jesus had raised from the dead..
Then the leading priest of Jerusalem.
decided to kill Nazareth as well,.
because many people had deserted,.
back then Judaism, deserted the synagogues,.
and they believed in Jesus..
So there's a short story, 11 verses..
It talks about the beginning of the final week of Jesus,.
and you can imagine, you know,.
Jesus was preaching the kingdom of God.

$^{241}$in the north of Israel, in Galilee,.
and slowly, step by step, he's moving away from Galilee,.
going down to Jerusalem in the south,.
and day after day, Jesus is talking to his disciples..
What is his purpose of approaching Jerusalem?.
And now, Jesus finally, almost there,.
almost there, approaching Jerusalem,.
and he was in Bethany..
And Bethany is a small village on the crest of Mount Olives,.
which is about maybe two miles away from Jerusalem..
So it may take about, you know, 15 to 30 minutes.
to walk from Bethany to Jerusalem..
And you can imagine the scene..
Many people from the villages,.
they jam into the little house of Mary, Martha, and Nazareth,.
and they wanna see Jesus..
Of course, they wanna see Nazareth as well,.
because they heard news of he's being raised from the dead..
And by the way, you know,.
and here is a big dinner was prepared,.
and right in the middle of the dinner,.
right in the middle of this huge banquet,.
Mary, she takes the most expensive perfume,.
the most valuable thing in her house, probably,.
and she breaks it..
And the perfume, the fragrance,.
permeates the entire house..
You know, you can imagine, this is a beautiful picture..
It's a beautiful scene..
And in this short story,.
there are at least two main characters,.
Judas and Mary..
Well, let's talk about Judas first..
Verse four, what did he say?.
In this beautiful scene, he said,.
"That perfume worth a year's wages,.
"and it should be so,.
"and the money we gain can be offered to the poor.".
And there's a footnote there..
He didn't really care about the poor..

$^{281}$He was a thief, and he was in charge of the church money..
He was in charge of the offering..
You know, when Jesus go from village town into city,.
preaching the gospel, preaching the kingdom of God,.
many people supported Jesus, right, men and women,.
and they offered the resources, and sometimes food,.
sometimes money, to support the ministry.
of the disciples and Jesus..
And Judas was, I think, the elected court team.
to take care of the finance, right?.
Maybe he was gifted, maybe he was an accountant..
I don't know, yeah, maybe different reasons,.
he was in charge of the money of the disciples..
And you can imagine, you know, with all the people there,.
this is something worth recording..
The gospel of John recorded this conversation, this dialogue,.
and Judas said, you know, these things..
Now, imagine with me a spectrum, a spectrum..
All the way to our left is someone or some people.
who oppose or reject the gospel..
And then slowly, you move to the right..
There may be doubters, right?.
They question about the faith, and probably sincerely,.
they wanna find out what Jesus is all about..
And then there are seekers, more serious about learning.
about the faith, and then finally,.
they became so-called believer..
Now, this is a simplistic spectrum, right?.
I'm not saying there's a difference between believers.
and disciples or servant, right?.
They are all the same..
In this scheme, however, you can imagine someone's moving.
from seeking the faith and then believing,.
and then being his disciples,.
and then being his followers, even serving, right?.
Now, in this scheme, where do you put Judas?.
Where do you put him?.
Where do you put him, the betrayer?.
Somewhere from the right all the way to the left, right?.
As time go by, he moves from the right.

$^{321}$all the way to the left..
And very interesting, right?.
Very interesting..
How does one, how does a person move.
from the right to the left, right?.
How does someone who has been a disciple for three years,.
he was personally called by Jesus himself, right?.
Judas was called to be one of the 12,.
personally called by Jesus when he began his ministry..
And Judas, among the 12,.
he actually experienced the power of the kingdom, right?.
The relationship with Jesus firsthand..
And he was initially trusted by the group.
to take care of the money..
And then eventually, Judas moved slowly to become a doctor..
And then finally, someone who rejects Jesus and betray him..
What is going on?.
And I think the root of it may not be money..
I think the root of it is self-interest, self-interest..
Now, Judas may be called a tip of the iceberg.
in the Jewish population of the day..
You know, the context of Jesus' story.
is Roman military occupation..
The Jewish people were suffering.
and they were being oppressed by the Roman empire..
So they were looking for some sort of Messiah..
Now, everybody was looking for a Messiah,.
you know, in those days and in today, right?.
We are all looking for some sort of salvation..
We are all looking for some solution.
to our supposed human problem..
In those days, the Jews were looking for a political solution.
to kick out the Romans,.
people that they don't like, oppressing them..
So this Messiah should come and to help us,.
you know, start a revolution or something.
and to do some change in our political situation.
so that we can have freedom, right?.
You know, most of the Jews would think.
that the immediate problem is the Roman empire..

$^{361}$This is our problem and we want political change..
Now, don't get me wrong, you know,.
there's so many social change,.
so many different political changes are good, all right?.
They bring about common good, you know, common grace..
Don't get me wrong..
I'm not saying, actually, you know, in recent history even,.
there were many social changes.
were initiated by Christian leaders, right?.
Like slavery..
Now, some of the social changes are good,.
but they are not the ultimate good..
They're not the ultimate good..
And Jesus come along and say,.
"Yeah, you can change from one political power to another,.
"but you are simply changing the external, right?.
"The outside, the external configuration, right?.
"Because you topple one regime,.
"you replace it with another, right?.
"You take care of one system.
"and then you replace it with another system,.
"and then it will go corrupt.".
And we are in the same situation, right?.
And Jesus come along saying,.
"Yes, I'll do some changes..
"I'll do some revolution,.
"but the revolution I'm bringing about.
"is not the external..
"It's from the inside out..
"The grace of God, the love of God, the kingdom of God.".
And this is what I'm coming to do..
I am coming to atone your sin..
This is the very root of a human problem,.
the heart of the problem, the heart of the issue..
It's our heart, right?.
It's the human heart..
We are corrupt and we are sinful,.
and we need God to redeem us..
And that's why the cross is necessary..
You know, I've searched in the church for some years..

$^{401}$You probably experienced church for quite some time..
And there are many different people that I encounter..
They come to the church, they come to faith.
with different reasons..
But my general conclusion is that most of the people,.
they don't ask enough questions..
They don't struggle with the cross enough..
They look at the church or look at Jesus from the outside..
Oh, he's such a good moral teacher.
and such a good example we can follow..
And I come to church, most of the people,.
with some other reasons than coming to faith..
There are some benefits coming to church, right?.
Don't you think?.
Maybe I can get connected with someone.
and then I'll get a new job, right?.
I come to church with some benefit too, right?.
Recently I hurt my knee and there are two PTs here,.
physical therapists, and then I get the service for free..
(audience laughing).
Just a little bit, right?.
Not to take advantage..
I mean, there has to be some benefit.
when we come to church or come to faith..
However, Jesus is challenging us,.
or the gospel is challenging us,.
to look at the real emphasis,.
which is the necessity of the cross..
Why is the cross necessary?.
Now we come to the second character, which is Mary..
Now Mary is also a disciple..
Jesus did something really extraordinary.
when he was preaching the kingdom of God..
As a rabbi, he actually had female disciples, right?.
All the women following him and supporting his ministry,.
and Mary was one of them..
So Mary, like Judas, she had witnessed the power of Jesus..
Remember before chapter 12, it was chapter 11..
That was the greatest miracle in the gospel of John..
Jesus raised Nazareth from the dead,.

$^{441}$brought his brother back to life..
So Mary has experienced God firsthand as well,.
experienced Jesus firsthand..
So she knew Jesus had come into the world for some purpose..
And that's why we have the story..
She broke the jar of perfume to wipe those on Jesus' feet..
To do what?.
Because she realizes Jesus comes to die, to die for us..
That's the only way to bring us back to God..
Jesus dies to save us..
Look at verse three..
It's a short verse,.
but it's like the camera zooming in.
to what Mary is doing to Jesus..
She did not hold back anything..
She gave the best gift, a costly gift, a perfume..
In those days, it probably could be sold for 300 denarii..
It's about 300 silver coins..
No, it's 300 silver coins..
It's equivalent to almost a year's salary,.
your whole year's salary..
What is your whole year's salary?.
What is yours?.
Well, tax day is coming..
So you should know, right?.
How much do you earn a year?.
A lot..
Well, maybe for some people a little poor..
A lot..
So she broke it and offered it to Jesus..
Now, in those days, when people had dinner,.
they did not sit up straight, right?.
They reclined on couches..
That's why their feet are sticking out, right?.
So you can imagine this crowded house,.
a cozy dining table,.
and then Mary came and she annoyed Jesus,.
his feet with the perfume..
And look at how the camera zoom in..
He not only annoyed it with the perfume,.

$^{481}$she was wiping his feet with her hair..
In those days, a lot different from today,.
women don't unbind their hair in public..
It is a disgrace..
It is humility, humiliation..
So you see what Mary is doing is.
she is taking the lowest, lowest, lowest possible place,.
humble herself to love Jesus..
You know, consider the entire story of the Gospel of John..
Jesus begins ministry in chapter one, all the way..
And there were seven miracles..
There were different teachings of Jesus..
And it goes on for some years, for almost three years..
And Jesus moved from Galilee in the North.
to Jerusalem in the South..
And he told his disciple many times.
that he came to die,.
to be betrayed,.
to be condemned,.
to be crucified..
And he said also on the third day, he will rise again..
Now, why do we need the cross?.
Why do we need the cross?.
The disciple didn't understand, right?.
They debate about it..
In fact, before Jesus went to the cross,.
they were fighting among themselves, right?.
Who's gonna be the highest?.
Who's gonna be the king when the kingdom comes, right?.
They thought Jesus was the political Messiah.
to solve their own perceived problem,.
the political problem, not the heart problem..
So the disciple didn't understand.
why the cross was necessary..
So here it is..
We need to think and struggle more about the cross..
About the cross..
Why is the cross necessary?.
For modern years, we as modern readers,.
sometimes we think of the cross like blood,.

$^{521}$very painful..
It's not according to our taste, right?.
And we often think that if God is all loving,.
why can't he just forgive, right?.
Why can't he just forgive and forget, right?.
And then we move on, right?.
We often think that way..
Just like your computer crashed, right?.
We format your drive, right?.
You buy a new computer and you start all over again..
Why can't, what's up with God, right?.
Why can't he just forgive and forget, right?.
And we often think in this terms..
Now I want you to struggle with me..
Why the Bible says the cross is necessary..
Jesus has to die, has to come to the cross..
Imagine if someone really, really wrong you,.
if someone really, really wrong you,.
the perpetrator, the oppressor,.
the one who sinned against you,.
can he just come to you and say, "I'm sorry,.
can we start all over again?".
Would you take that?.
I mean, someone really wrong you, okay?.
Someone really wrong you..
It's not trivial stuff, right?.
For trivial stuff, we can just forget and go on, right?.
If someone really wrong you,.
you know, took away your, maybe your name, your future,.
your happiness or your reputation,.
something really wrong has done to you, you know?.
There has to be a cost..
There has to be a debt, right?.
Either you pay for the cost,.
or the perpetrator pays for the cost,.
or you both share the cost, right?.
When my kids were younger,.
when they were, you know, they sit in the car.
and, you know, we driving,.
and I remember one day, they don't remember anymore,.

$^{561}$but they asked, you know, "What if we crash the car,.
right, damage the car,.
or somebody else damage our car, what will happen?".
I said, you know, "The insurance will pay for it.".
And then, back then, they were so much younger,.
they didn't understand, they were so happy..
"Oh, because somebody is paying for it,.
so that's cool, that's cool," right?.
What you don't understand is.
if there was some damage done to you,.
there has to be a debt that has to be paid for..
Someone has to pay for it..
Right?.
And again, if you think about most of the wrongs.
done to people, cannot be calculated in terms of dollars..
No matter how much compensation is being fought for,.
it's just not eloquent, right?.
Someone may have hurt you,.
or robbed you of something that is important..
You can't put a price tag on those things, right?.
Now, in this situation, imagine with me,.
you know, there are two options..
The first option is retaliate, right?.
You do the same thing to the person.
who did the same thing to you, right?.
But is that a good option?.
Maybe not, right?.
Because it's a vicious cycle, right?.
Evil, and then evil, and then evil..
Hurt people, hurt people, right?.
You heard about that?.
Hurt people, always hurt people..
That's the first option..
The second option, what may be called the redemptive option..
The redemptive option is going after justice and forgiveness..
Justice and forgiveness..
And this option is redemptive, is sacrificial,.
because the price has to be paid, right?.
Think about in very simple human term, right?.
Someone has done something wrong to you,.

$^{601}$someone has to pay..
Even you say, you know, I'm not gonna retaliate,.
I'm gonna forgive, you know, right away..
I'm gonna forgive, but you know, think about it,.
dynamic is that you take the pain,.
someone take the cost, someone pay for the cost..
That's why people said forgiveness is suffering..
It is suffering, because we suffer and take the cost..
You know, all of those redemptive elements.
in this kind of dynamics is sacrificial..
And think further,.
all redemptive actions are sacrificial and life-changing..
We parents, we know that, right?.
I don't know how many times I sacrificed.
in order to give life to somebody else, right?.
I reorder my life priority.
in order to give life to the other person..
If I don't do that, their life will be in ruin, right?.
Parents do that, lovers do that, friends do that..
Now imagine with me, you know,.
those things times a factor of infinity..
That God himself, who is totally holy and pure,.
and we violate evil, we did evil..
We sin..
This is what amazing about the cross,.
that God himself chooses to love us..
And that's why Jesus had to sacrifice himself.
to pay for the debt of sin..
The debt of sin incurred by us.
sinning against the holy, eternal God..
It's an eternal relationship being ruptured..
Jesus has to come and pay for the debt on the cross..
And when you truly believe that,.
when you really take this love into your heart,.
you should change from the inside out..
You should gradually change your perspective on things,.
your relationship with God, your serving,.
all that you do and all that you make.
in terms of decisions..
You will orient yourself toward the cross,.

$^{641}$you're toward the God who loves you..
Mary does this..
She realized the cross, the sacrificial death,.
the substitutionary atonement,.
someone paid for your debt..
That's the only way we can come back to God..
And this is her last chance to honor this death,.
to honor what Jesus has done for her..
Anoint his body, his feet with perfume..
Respond with love and sacrifice in a generous way..
When you respond to God, the love of God,.
with the same love,.
your life will be changed inside out..
We're gonna sing our last song, our response song..
I'll invite the praise team to come back up here..
I want you to, as we sing the song again,.
imagine what is going on in Bethany,.
what Mary did..
And think with me, and struggle with me,.
why the cross is necessary..
Why is it necessary for you?.
And you respond appropriately..
To the cross..
(gentle music).
♪ I don't know how ♪.
♪ To say exactly how I feel ♪.
♪ And I can't begin ♪.
♪ To tell you what your love has meant ♪.
♪ I'm lost for words ♪.
♪ Is there a way ♪.
♪ To show the passion in my heart ♪.
♪ Can I express ♪.
♪ How truly great I think you are ♪.
♪ My dearest friend ♪.
♪ Lord, this is my desire ♪.
♪ To pour my love on you ♪.
♪ Like oil upon your feet ♪.
♪ Like wine for you to drink ♪.
♪ Like water from my heart ♪.
♪ I pour my love on you ♪.

$^{681}$♪ If praise is like perfume ♪.
♪ Then I'll lavish mine on you ♪.
♪ 'Til every drop is gone ♪.
♪ I'll pour my love on you ♪.
Would you please rise?.
(gentle music).
♪ I don't know how ♪.
♪ To say exactly how I feel ♪.
♪ And I can't begin ♪.
♪ To tell you what your love has meant ♪.
♪ I'm lost for words ♪.
♪ Is there a way ♪.
♪ To show the passion in my heart ♪.
♪ Can I express ♪.
♪ How truly great I think you are ♪.
♪ My dearest friend ♪.
Let's sing this..
♪ Lord, this is my desire ♪.
♪ To pour my love on you ♪.
♪ Like oil upon your feet ♪.
♪ Like wine for you to drink ♪.
♪ Like water from my heart ♪.
♪ I pour my love on you ♪.
♪ If praise is like perfume ♪.
♪ Then I'll lavish mine on you ♪.
♪ 'Til every drop is gone ♪.
♪ I'll pour my love ♪.
♪ Like oil upon your feet ♪.
♪ Like wine for you to drink ♪.
♪ Like water from my heart ♪.
♪ I pour my love on you ♪.
♪ If praise is like perfume ♪.
♪ Then I'll lavish mine on you ♪.
♪ 'Til every drop is gone ♪.
♪ I'll pour my love ♪.
♪ On you ♪.
Let's go to prayer..
Father, thank you for your love, your great love..
Beyond all the words that we can explain and describe,.
it's incomprehensible..

$^{721}$Your grace abounds in our lives..
So we come before you, we want to respond..
Respond with love..
Responds with our open heart..
So let your grace and your love, the love of the cross,.
to fill us up to change and transform our lives..
Change our orientation, change our perspectives,.
change our direction..
How we make decisions, how we live our lives,.
is to respond to you as a sacrifice..
Thank you for the fact that your love never change..
And I pray that your spirit will move us in such a way.
that we'll love you back..
Thank you..
Thank you for being our loving God..
In Christ, we pray..
Will you dismiss, and I want to see you next week,.
and let's worship God together..
[MUSIC PLAYING].
(upbeat music).
\newpage



\section{約翰福音 12:1-11}
\label{sec:zbYaP2QF_GY}
\textbf{耶穌最後的日子:伯大尼  (約翰福音12\_1-11) - 陳崇基牧師}
\newline
\newline
連結: \href{https://youtube.com/watch?v=zbYaP2QF-GY}{\texttt{https://youtube.com/watch?v=zbYaP2QF-GY}} ~~~~ 語音日期: 2023-03-20
\newline
\newline
\hyperref[sec:lL2aGhSKVWQ]{\small{< < < PREV SERMON < < <}}
~
\hyperref[sec:index_chronic]{\small{[返順時目]}}
~
\hyperref[sec:index_scriptual]{\small{[返順卷目]}}
~
\hyperref[sec:5_8ZAI7O6Mc]{\small{> > > NEXT SERMON > > >}}
\newline
\newline
約翰福音 12:1-11
\newline
\begin{longtable}{cl}
\hline
\hline
章節 & 經文 (和合本修訂版)\\
\hline
12:1 & \begin{tabularx}{0.7\textwidth}{X} 逾越節前六天，耶穌來到伯大尼，就是他使拉撒路從死人中復活的地方。 \end{tabularx} \\ \\ \relax
12:2 & \begin{tabularx}{0.7\textwidth}{X} 有人在那裡為耶穌預備宴席；馬大伺候，拉撒路也在同耶穌坐席的人中間。 \end{tabularx} \\ \\ \relax
12:3 & \begin{tabularx}{0.7\textwidth}{X} 馬利亞拿著一斤極貴的純哪噠香膏，抹耶穌的腳，又用自己頭髮去擦，屋裡充滿了膏的香氣。 \end{tabularx} \\ \\ \relax
12:4 & \begin{tabularx}{0.7\textwidth}{X} 有一個門徒，就是那將要出賣耶穌的加略人猶大，說： \end{tabularx} \\ \\ \relax
12:5 & \begin{tabularx}{0.7\textwidth}{X} 「為甚麼不把這香膏賣三百個銀幣去賙濟窮人呢？」 \end{tabularx} \\ \\ \relax
12:6 & \begin{tabularx}{0.7\textwidth}{X} 他說這話，並不是關心窮人，而是因為他是個賊，又管錢囊，常偷取錢囊中所存的。 \end{tabularx} \\ \\ \relax
12:7 & \begin{tabularx}{0.7\textwidth}{X} 耶穌說：「由她吧！她這香膏本是為我的安葬之日留著的。 \end{tabularx} \\ \\ \relax
12:8 & \begin{tabularx}{0.7\textwidth}{X} 因為常有窮人和你們在一起，但是你們不常有我。」 \end{tabularx} \\ \\ \relax
12:9 & \begin{tabularx}{0.7\textwidth}{X} 有一大群猶太人知道耶穌在那裡，就來了，不但是為耶穌的緣故，也是要看耶穌使他從死人中復活的拉撒路。 \end{tabularx} \\ \\ \relax
12:10 & \begin{tabularx}{0.7\textwidth}{X} 於是眾祭司長商議連拉撒路也要殺了， \end{tabularx} \\ \\ \relax
12:11 & \begin{tabularx}{0.7\textwidth}{X} 因為有許多猶太人為了拉撒路的緣故，開始背離他們，信了耶穌。 \end{tabularx} \\ \\
[1ex]
\hline
\hline
\end{longtable}
$^{1}$各位好,歡迎各位,我們不如和大家旁邊的和左右的和他們打個招呼,歡迎..
又到這個時候,我們快要到壽苦節和復活節的時候,所以我們都應該思想一下耶穌基督的壽苦,他的十字架和他的救恩和三天之後的復活..
這個很重要,我們需要大家思考救恩的必要,為什麼需要十字架,和十字架對我們有什麼意義..
我們有特別的嘉賓,有六位學生,我們歡迎他們.他們來做課程,所以他們會問你一些問題,不要給他們熱時間..
正如我所說,我們快要到復活節和壽苦節的時候,壽苦節或十難節,那個星期五,我們當然有音樂崇拜,剛剛報告也有提到..
就是一個音樂的崇拜,我們思想到上帝給我們的救恩,十字架的救恩..
壽苦節很多時候叫做Good Friday,然後到復活節的時候,我們Easter就會有唐慶的時候..
所以我覺得我們可以在接下來這幾個星期,去思想一下,特別是我的講道裡面,就是講到耶穌最後的一個星期,或者最後的那幾天..
聖經裡面有很多不同的記載,如果你從來沒有看過聖經也好,你去新約的頭四本書,《馬太》,《馬可》,《路加》和《約翰福音》四本的福音書,你都會找到耶穌的生平..
所以這些是很容易找到的一些內容,亦都應該是大英姐妹,如果你是回教會一段時間,你看過聖經一段時間,應該你對這些內容都很熟悉..
所以今天既然是我們講這個題目的時候,不如我們給一個測驗給大家..
第一條問題,哪一本福音書是記載耶穌的出生的,即是所謂的聖誕節..
答案是什麼?應該是《馬太》和《路加》吧,對不對?很難嗎?.
應該是《馬太》和《路加》才有記載耶穌出生的事,《馬可》是沒有的,《馬可》一開播就講耶穌公開的傳道,已經成人了,耶穌三十多歲開始他的傳道最後那幾年,《馬可》就是講這個..
《約翰福音》孟車篇你都可以講,太初有道,但那個是講永恆的道,然後道成肉身來到我們中間,所以《約翰福音》就沒有怎麼講聖誕節的時候..
所以《馬太》和《路加》第一和第三本福音書才有提到耶穌的出生..
好,第二條,哪本福音書記載了耶穌公開傳道之前的事?《馬太》,《馬可》,《路加》還是《約翰》呢?.
可以跟你旁邊那個人談談嗎?.
誰?《馬太》,《馬太》,《路加》,沒錯,是《路加》,12歲耶穌去到聖殿的時候,不是玩失蹤的嗎?.
第二章,只是《路加》才有提過,其他都不感興趣,我們很感興趣的東西,福音書不是很感興趣..
第三條,哪本福音書記載了耶穌,中文翻錯了,是耶穌的受死,被釘,死亡,埋葬和復活的?.
哪本書記載了?誰說是E的?你們好害怕,是E的,不用怕,四本福音書都有記載耶穌的死和復活的,死和復活四本福音書皆有記載..
我們在基督教的所謂年曆裡面,有幾個重要的日子,聖誕節,壽難節,復活節,五旬節,這些都是所謂的紀念日,我們會記載的一些耶穌的生平和事蹟..
有很多其他的東西,相信耶穌是做了,但是新約的聖經是沒有記載,譬如耶穌小時候在幼稚園,他怎樣跟其他同學相處?.
耶穌有沒有溫書?還是他考試的時候可以用神蹟,然後拿到A呢?這些都是我們很喜歡揣測的東西.事實上,在聖經以外的一些記載,或者有些人寫的一些小說,或者你今時今日看一些有關耶穌的電影等等,這些東西都是會感興趣的..
他長大的時候到底會是怎樣的呢?跟瑪利亞,約瑟父母相處,跟兄弟姐妹相處等等等等,這些現代人我們都很感興趣的..
但是聖經卻是焦點不是放在那裡,事實上在約翰福音最後那兩節,最後那兩節,21章最後那兩節,講些什麼呢?.
其實耶穌發生了很多事,做了很多事,他說這門徒就是作者,約翰福音的作者,他說為這些事作見證,就是之前21章裡面記載了的東西,記載這些事我們知道他的見證是真的..
耶穌所行的事還有很多,其實還有很多關乎耶穌的事的,是不是?一天是不是吃三餐,他幾多磅,腰圍三圍是什麼等等,這些我們非常感興趣的東西,他帥不帥等等,這些都是我們非常感興趣的..
他說有很多事,還有很多,若是一一都寫出來,全世界都容不下所寫的書,所以有很多事,沒錯那些門徒可能聽過或者接觸過,可能知道的,或者耶穌曾經有講過的..
這些東西如果要記載的話,全世界都會放不進去,但是四本福音書,或者甚至整個新約聖經所記載的事,都是跟什麼有關?都是跟Good Friday和Easter有關的..
譬如你看約翰福音,約翰福音整本都是記載他最後那幾年,公開傳道,到他在耶路撒冷受害,被釘十字架,死而復活,二十一章聖經裡面,花了這麼多章節,全部都是記載最後最後那幾年..
馬可福音更加是,其他福音書都是一樣,所以我們很感興趣的東西,聖經卻不感興趣,聖經卻告訴我們最後那一個禮拜,最後那幾年,才是你要去放焦點的地方..
這個就是耶穌和其他的領袖,或者包括宗教領袖,不同的地方,所有的領袖都是來到這個世界上,他是活一個生命,以致你有一個榜樣,你去學校的..
你想起的英雄,想起的重要的人物,都是來這裡,他在這個世間上,能夠生活,以致他有一個好的生活,有一個好的榜樣,以致你能夠去學校他..
只有耶穌一個,來到這個世界,是死的,他的目的就是來這裡死,不是單單來這裡活,給我們一個榜樣而已..
當然這個是真的,他的生命是一個榜樣,但是他來的目的,不是單單來這裡活,他的目的,是死在十字架上,這個就是耶穌和其他的宗教信仰,千萬不要說,所有宗教信仰都是一樣..
不是的,只有耶穌基督,只有他一位,來到的目的是去死的,並且他所作的所有東西,說的所有東西,都是預備那一天,就在十字架上死的人..
今天我們來看一段經文,就是在耶穌最後的一個星期的時候發生的,約翰福第十二章,十二章第一節,我們先讀經文..
預節前六天,預節就是所謂Good Friday的時候,預節前六天,耶穌來到一個村莊叫伯大利,就是他使拉撒路從死人中復活的地方,有人在那裡為耶穌預備筵席..

$^{41}$馬大侍候,拉撒路也在和耶穌坐直的人中間,瑪利亞就拿著一根極貴的純拿達香膏,抹耶穌的腳,又用自己的頭髮來擦,屋裡面就充滿著香膏的香氣..
有一個門徒,就是那將要出賣耶穌的教育人尤大說,為什麼不把這香膏賣三百個銀錢來周祭窮人呢?.
他說,並不是關心窮人,而是因為他是個賊,又管錢囊,常偷取錢囊中所存的..
耶穌說,尤他吧,他者向告本是為我安葬之日,我死的那一天,留著的.因為常有窮人與你們一起,但是你們不常有我..
有一群猶太人,知道耶穌在哪裡就來了,不但是因為耶穌的緣故,也是為了看耶穌,使祂從死人中復活的拉撒路..
於是眾世之長和民間的領袖就相議,連拉撒路也要殺了..
因為有很多猶太人,為了拉撒路的緣故,就背棄了當時的猶太教,信了耶穌..
這段很短的故事,應該很多人都很熟悉的,故事發生在整個耶穌的生平,終於去到一個所謂的高峰,或者最後的一個禮拜,花最多的筆墨來形容的一個禮拜,.
The Final Week,就是耶穌終於快要來到耶路撒冷,耶路撒冷就是祂死的地方,.
祂去到一個村莊叫做伯大利,Bethany,Bethany和耶路撒冷,其實只不過是大概四五個字的路程,大概兩英里,很近很重要的地方..
很多很多的人因為耶穌的緣故,或者甚至因為耶穌做過的一些神蹟,例如之前那一張手,就是拉撒路從死裡復活的神蹟,也是在伯大利那裡發生的..
所以有很多人就,圍觀也好,或者感到很讚歎也好,他們就去到伯大利的地方,剛剛因為有一個中東的文化,就是有什麼客人來到,或者尊貴的貴賓到了的時候,他們就會大排筵席來歡迎他..
所以有很多的村民就走在一起,去到一間房子裡面,可能是一個很小很小的房子,只有幾百呎的房子裡面,他們就擠湧在裡面,然後就在那裡有一個很大的晚餐..
所以有一個筵席就為耶穌來去舉行,在吃東西的時候,中間的時候,瑪利亞就拿了他一瓶,我想大概12盎司左右,這麼小的一些香水的瓶子,那些叫做牙達香膏,是一些膏油來的..
他就打爛了,然後就將它抹在耶穌的腳上,那你就可以想像那些香氣會充滿整個房間,有誰喜歡吃榴槤的?.
那些香氣你拿一塊出來就夠了,喜歡吃的那些人就會擠湧在一起,那些不喜歡吃的就馬上驅走他們..
就是一個這樣的情況,那種香氣就會瀰漫著整個房間,其實就是在一個當時的社會,一個情境,是一個很美的情境來的,有香膏,有香水,來去抹在耶穌的腳上..
這個故事裡面,起碼有兩個重要的人物,第一個就是猶大,第二個就是瑪利亞,我們先看猶大,一個這樣的美景出現的時候,猶大說什麼呢?.
他說為什麼不將這些香膏賣掉,然後我們有三百兩銀錢,三百公斤錢,來去周濟窮人呢?.
這裡很合加勒人猶大的品性,在不同的科學技術裡面,加勒人猶大都是有這樣的恩賜,可以這樣說,素口很好..
他說這個香膏一看就知道值三百兩銀錢,換算以前,當代,三百兩銀錢大概有多少錢呢?.
就是一個普通的工人,他每年的薪金,每天大概有一兩銀錢,以致他能夠有戶口,有東西吃..
三百兩銀錢大概是一年的工資,所以他一看見的時候,他就發現,這個香膏或者這瓶香水,其實是值一年的年資,一年的年資多不多?.
很多,叫你奉獻一年的年資,可以嗎?不要吧,上帝不會這樣的,我們對上帝很有信心,我不奉獻,其他人一定有奉獻的..
我們一定會這樣說,一年的年資你想想,你算一年你的年資是多少?我們快要tax day了,大家應該在算錢吧,一年的年資..
他就把它打爛,然後就塗在耶穌的腳上,所以加勒人猶大一看見的時候,他就覺得這個非常之不妥,並且到預節或者good friday的時候,.
猶大人有一個習慣,就是他們會有一些額外的奉獻,會奉獻到耶路撒冷的聖殿,是周祭共人用的,.
每逢到大時大節,就是所謂的passover預節的時候,就是耶穌死的那一天,十字架,受苦節的那一天,good friday的那一天..
所以猶大人就已經心裡面有數了,為什麼不把它賣掉呢?他對馬尼亞的批評,賣了它,正好這個時候我們有慈惠奉獻,我們不如去到周祭共人吧..
當然約翰作者已經交代了,其實他心裡面在想什麼,by the way,在說猶大的時候,在新約裡面說出賣耶穌的加勒人猶大,其實他不只是取笑這麼簡單,.
其實是一種很意外很突然的描述,意思就是說,居然這個門徒是最後出賣耶穌的,我們知道最後的故事是怎樣的,猶大衛了幾十塊的銀錢,.
他就出賣了耶穌,把耶穌交給了當時的公安或者祭司長等等,這個就是猶大,我想想,猶大,如果我們想一個人和耶穌的關係,.
你想想一條線,最左邊的就是一個反對信仰的人,然後你慢慢慢慢去到中間的時候,可能他有些質問或者提問,懷疑等等,直到他變成信徒,.
再繼續去和信仰打關係,和上帝打關係,和耶穌繼續去認識多了的時候,這個人就不單止是一個所謂信徒,變成一個門徒,甚至可能一個僕人,一個領袖,.
在一個這樣的光譜底下,由你的左邊,最左邊,是一個反對信仰,反對主的,到你的最右邊,就是一個最親近主的,去服侍主的,在中間你覺得猶大應該放在哪裡?.
猶大應該放在哪裡?猶大是耶穌親自去呼召的門徒,十二門徒之一,猶大也是每一天要跟著拉比去到出行,去到不同的村莊傳道,.
甚至他是被信任去管理錢的,等於我們在我們的執事會,我們有一些比較因此這方面強一些的,他們就是管錢,.
他們負責教會的奉獻,到底為什麼會有一個耶穌親自呼召,並且三年跟隨耶穌的人,他慢慢會移動去到變成一個反對甚至出賣耶穌的人,到底中間發生什麼事呢?.
有什麼人或者有什麼事發生,會使一個人能夠從一個被呼召作門徒的,慢慢他對他自己可能對主有懷疑,並且到最後他是怎樣?.
他是反對的,不知道發生什麼事,會不會是因為錢的緣故,可能又不是,因為賣耶穌的時候都沒有講價,三十多萬錢,他就答應了當時的當局,祭司長就賣了耶穌,.

$^{81}$會不會是一些自我的目標,或者自我追求的一些動機,以致使他的信仰,他的生命變質,或者可能會,因為當時的猶太人,大部分的猶太人的群眾,.
他們想找一個彌賽亞,想找一個救主,都是一個軍事政治的一個救主,慢慢你跟隨耶穌的時候,就發現原來不對勁,他不是想做我心裡面做的事,.
我們每個人去信主的時候,我們都可能懷著一些不是太過純正的動機,這也沒問題,開頭的時候,有人來教會可能為了認識女孩子,有些認識不到就走,.
有人來教會可能為了有公,有網絡,有些社交的好處,等等不同的利益或者不同的益處來教會,.
開頭的時候是沒問題,但慢慢你久了的時候,你就需要去面對一個問題,就是到底這個主是怎樣的,這位耶穌,真正的這位耶穌應該是怎樣的,.
這個就是很多時候我覺得很多人對十誡,對耶穌的十字架,對耶穌的救恩,是我們想得不通的,想得不夠的,掙扎得不夠,.
對這個十誡的道理,我們不是很放心思,或者我們對這個十誡的道理,一個這麼難明的道理,我們掙扎得不夠,猶太會不會是這樣呢?.
可能好像當時的猶太人一樣,既然這個都不是我心目中想的救主,其實三十兩銀錢已經給多了,將它賣了,.
耶穌就說,耶穌一向的傳道,就是說我來到,要給的改變,不只是單單政治上的改變,.
猶太人就很想有政治上的改變,很簡單,當時的情況,就是他們在羅馬帝國底下,他們被軍事管理,他們是很不順,.
所以他們希望利塞亞尼做什麼呢?就是帶領一個軍事的革命,然後有一個社會或者政治上的改動,一個改革,.
這個就是猶太人一向想做的事情,事實上歷史裡面真的是這樣,第一世紀主要是七十年,主要是一百三十二年,兩次大的革命都是為了做這件事,他們就是想推翻羅馬帝國,.
你聽聽耶穌說的話,聽聽耶穌所做的事情,你就會了解他不是想做這件事,或者應該這樣說,不是說這些社會的改革一定是錯的,請你不要誤會我的意思,.
有很多社會的改革是好的,並且在歷史當中,就算是近代的歷史,有很多社會的改革都是基督教的領袖站出來的,譬如說廢除黑奴,.
等等這些所謂社會的改變都是有它的好處的,但是最多的好處,最極的好處只不過是短暫的,社會的改革,你只不過是將另外一個的權兵放上去,.
然後權兵也會腐敗,然後到下一次又會將另外一個權兵放上去,這個權兵也會腐敗,如果看通看透的話,耶穌就說我要做的事情是一個更不同,.
或者更深的一種改革,耶穌說我來為你贖罪,我去犧牲,是解決人類最基本最心臟的一個問題,就是我們的罪,我們與上帝的決裂,我們與上帝的分裂,.
如果這個不解決的話,一切都會解決不到,就是在你生命當中,inside out,由裡面出來的,你被上帝的愛,被上帝的救恩,上帝的道去拯救,去更新之後,你的生命開始改變,.
這個一個又一個的改革才是最基本的改變,耶穌說我們做的改革就是這種改革,由大可能他只不過是看到表面,正如我所說,我們對十字架的救恩,很多時候我們掙扎,真的不夠,我們真的沒有想清楚,.
我們想的東西,或者我們思考十字架的時候,我們只不過是很皮毛,我們只不過是很皮毛的時候,在想,他真的很偉大,耶穌很偉大,他受苦,他受痛楚,以致我們有個榜樣,這個不是錯,.
是太過表皮,裡面到底是什麼意思,由大錯失了這個機會,或者掙扎不夠,瑪利亞呢,瑪利亞就不同了,瑪利亞也是耶穌的門徒之一,雖然他不是並列於十二門徒,.
但是他是門徒裡面其中一個,因為耶穌做的一些事情,當時拉比通常都不會做的,就是有女門徒跟著,有些女門徒會跟著耶穌去服侍他,然後去支持他的事工,或者等等,奉上他的金錢等等不同的資源,就是想去支持拉比所做的事情,他所傳的,瑪利亞是其中一個,.
事實上我剛才提到,上一章,瑪利亞是看著他弟弟死而復活的,拉瑟魯,他也知道,似乎他也感受到,耶穌真正的目的是什麼,十二章第三節,.
這裡很仔細的說,瑪利亞他的行動,將最好的獻上,極貴的純拿達香膏,拿達在哪裡來的,純拿達香膏是當時的印度來的,那些植物,香草,是從印度運到中東的,.
所以是很貴重,很遠而來的香草,去到提煉出來的香水,香油,這裡當然沒有解釋,為什麼瑪利亞會有這麼貴的東西,應該是他們很窮,.
可能是他的家傳之寶,當時很多女士,女孩子,都會有他家庭的所謂一代傳一代的那些很貴重的東西,.
我記得我們結婚的時候,我老婆結婚的時候,出嫁那天,她媽媽就拿了一些從來沒有見過的東西來,她媽媽都很節儉的,.
從來都不穿金戴銀的,無端端那天這麼多的,我計算一下,我都覺得OK,.
溶了它,不是,要留給女兒,可能是一些家傳之寶,一些嫁妝,所以瑪利亞就拿了他最貴重最貴重的東西,一年的資金,.
他就拿出來,就是他整個的401k,整個的retirement的account,他就拿出來然後打爛它,在眾目睽睽之下,他就抹耶穌的腳,.
當時吃東西不是我們今天這樣坐起來,坐立,吃東西的,當時吃東西是餓著吃的,睡在那裡,很享受的,.
所以他的腳就凸出來,瑪利亞就在中間,他就打爛了他的香水瓶,然後抹在耶穌的腳上,並且用他的頭髮來刷,.
這個就是和我們今時今日的文化有些不同,以前的女人是不會披頭散髮的,一定會包著或窄起,.
披頭散髮其實就是特別在公眾上面,披頭散髮其實就是謙卑他自己,他拿了一個最低等的地位,.
來將這些香水抹在耶穌的腳上,如果你看整個的約翰福音的次序,.
你會知道耶穌開始傳道,然後他有幾年的時間,他就在加利利不同的村莊傳道和醫治等等,.
他有好幾次很明確的就是和門徒說,人子來要做什麼,不是每一個都很榜樣,人子來就是為你罪人而死,.
說過好幾次在加利利,他說我將會去耶路撒冷,做什麼呢,不是去做王,而是去受死,被出賣,埋葬,然後再復活,.
瑪利亞應該會明白就是這個東西,就是十字架的必要,為什麼一定需要十字架,.
我們現代人想起十字架的時候,我們只會想起什麼呢,我們想起血腥的東西,一些很痛苦的東西,一些刑罰,.

$^{121}$這些不適合我們的口味,如果用一個純粹的taste,純粹一個我們現代人的喜不喜歡,.
來衡量耶穌在十字架上所做的事情,I think we miss the point,很多時候我們搞不清楚十字架的道理是什麼,.
為什麼要人子,神的兒子親自在十字架上死,來去代罪,犧牲,然後才可以有這個救恩,.
因為我們現代人只會很表面的去想,如果上帝是一個慈愛的上帝,為什麼他不就這樣,寬恕就行,.
當作是忽略了他,罷袖,可不可以這樣,為什麼要搞一大輪,我想給大家一個小小的例子,.
讓大家去明白一下,想一想到底十字架的救恩對你是什麼意思,你想像一下,.
如果有一個人做了一些事情是得罪你的,損毀了你的,或者搶奪了你人生裡面一些很重要的東西,.
你想一下,應該會怎樣做,做一些很不對的事情在你身上,他可不可以就這樣說對不起,可不可以算數,.
可不可以,是不可以的,我不是說一些皮毛的東西,而是他真的去得罪你,無論是損毀你的名聲也好,.
或者去強暴你也好,或者使你失去了一些自由也好,可不可以就這樣就算數,當然不可以,.
如果不行的話,我們應該大概有兩個選項,第一個選項就是報復,使他得到他加在我身上的傷害,.
但如果第一個選項就是報復的話,這樣又不行,因為受害的人通常會加害他人,.
你只會繼續有一個邪惡的循環,惡再以惡去報復,以惡報惡,這個世界會更加醜陋,更加邪惡,.
你想一下第二個選項是什麼,第二個選項可能可以叫做一個舊熟悉的選項,.
舊熟悉的選項就是公義與寬恕,公義與寬恕是什麼意思,.
罪一定是有公價,罪一定是有後果,有consequence,或者有個cost,有個price,.
所有加害或者得罪的事情,最終一定會有規矩,有虧損,舉個例子,我小時候,.
他開始坐車的時候,他明白什麼叫汽車保險,他就問,如果你爸爸撞車,車壞了,.
怎樣修理,我說,insource會賠的,如果人家錯撞壞你的車,假設我是存保,.
那些虧損誰負責,我說,都是insource公司會給的,他們小時候就很開心,.
原來這麼好,原來有人賠,怎樣做都有人賠,他不明白,其實是有人給錢,.
每一個錯誤,每一個錯事或者罪,都一定有代價,有虧損,.
你想一下虧損裡面,一定是有一個人去承擔,一定有他者去承擔,永遠都是,.
以至公義能夠被完成,以至寬恕可以來到,一定有人pay for這件事,.
你想一下所有罪認成的後果,規絕,一定是有他者,有一者,來擺上犧牲,.
就算你說寬恕一個人也好,我不會報復,我就這樣吃了他,我忍辱負重,.
你就是那個他者,你就是去承擔了那個規矩,承擔了那個所謂罪的公價,.
上帝,你再想一下我們和上帝的關係,上帝是一個永恆無罪的一位上帝,.
聖潔全然無罪,當我們有罪的時候,我們和他的關係虧缺的時候,哪一個來對付,.
基督教的信仰就是告訴你,原來這位上帝願意犧牲來給這個公價,.
這個是人間能夠明白的道理,也是超越人間的道理,人間明白就是我剛才給你的例子,.
大家都明白了,虧損一定要被擺上,一定要pay for,然後才可以有所謂的寬恕,.
這個超越人間的道理,超越的意思就是說上帝他沒有做錯任何事情,也沒有任何責任要去負上那個規矩,.
但是他的兒子願意來在十字架上犧牲,去挽回我們和上帝永恆的關係,.
上帝親自犧牲,瑪利亞應該了解這個所謂代替的愛,代替的贖罪,.
當你了解這個愛的時候,當你了解這個救恩的時候,這個救恩才能夠在生命裡面產生應該有的效果,.
所有的愛,生命改變的愛,都是一種這樣的愛,是一種犧牲的愛,.
父母我們會為兒女犧牲,我們會為兒女擺上,這只不過是人間的層面,上帝擺上他自己的生命,.
去挽回我們,這個才是真正的愛,這個才是真正能夠改變你生命的愛,瑪利亞多貴也好,.
他說這個是值得的,值得去得到上帝的救恩,值得去看見十字架對他生命的後果,.

$^{161}$接受這個十字架,這個生命才能夠改變,我們結束祈禱的時候,.
讓我們給這個十字架的救恩再一次去提醒,求主讓這個十字架的愛,在你生命裡面能夠產生它的功效,.
主我們多謝你,因為你是那一位願意來犧牲,去拯救我們的神,.
十字架的救恩,我們人能夠明白,也是超越我們人世的道理,.
主我們今天來到你面前,我們要去認識這個十字架,認識你在我們身上所做的,.
你擔當我們的苦難,擔當我們罪的規矩,.
以致有公義,有赦免,我們能夠成為你自己的子民,主我們感恩,讚美你,.
主讓我們的生命向你敞開,去接受這個道理,接受你十字架的救恩,.
讓你自己的愛能夠去轉化,去改變我們,我們作為你的兒女,我們生命能夠被更新,.
在我們生命當中,我們所有的動機,所有的決定,能夠去回應你對我們的大愛,.
我們這樣祈禱,奉主命令,.
我願上帝祝福各位,希望下星期大家能夠一齊思想上帝自己的愛..
謝謝收看~.
\newpage



\section{}
\label{sec:5_8ZAI7O6Mc}
\textbf{比對研究 10: 巴力宗教 3 - 陳崇基牧師}
\newline
\newline
連結: \href{https://youtube.com/watch?v=5-8ZAI7O6Mc}{\texttt{https://youtube.com/watch?v=5-8ZAI7O6Mc}} ~~~~ 語音日期: 2023-03-25
\newline
\newline
\hyperref[sec:zbYaP2QF_GY]{\small{< < < PREV SERMON < < <}}
~
\hyperref[sec:index_chronic]{\small{[返順時目]}}
~
\hyperref[sec:index_scriptual]{\small{[返順卷目]}}
~
\hyperref[sec:WPZJt_sBFNI]{\small{> > > NEXT SERMON > > >}}
\newline
\newline
$^{1}$所以剛才略提過巴力神話為什麼會叫做巴力循環.
BioCycle.
為什麼是一個循環.
因為在這個故事裡面.
巴力永遠都不是最終可以得勝的.
因為他會被殺敗.
但是他的情人又為他復仇.
然後他又再復生.
所以在這個巴力神話裡面.
其實就可以看到巴勒斯坦地的農耕和氣候的循環.
每年到10月至4月的時候.
西曆10月到4月的時候.
就是雨季.
這個時候就有如巴力復生.
每逢雨季過了.
大概5月左右.
到8月底9月的時候.
就是巴勒斯坦地.
所謂巴力死亡的時候.
因為這段時間就不會有雨落.
盼望是什麼時候.
盼望就是希望雨季真的會動.
也就是說在巴力的宗教裡面.
在巴力的信仰裡面.
就是希望巴力再次復生.
因為這段時間最乾旱的時候.
5月至9月.
就是巴力死亡.
然後那個死神.
去到操控一切.
沒有雨雲沒有雨水的時候.
希望能夠有雨水回來.
如果你這樣看背景的時候.
之後我們可以再看多一點.
不過今天就看到這裡.
你可以看一下聖經是怎樣去回應這些文化.
在巴力死亡的背景.
《斯坦第68篇》第4節.
那句話是要向神去唱詩.
去歌頌他力明.

$^{41}$向那駕車經過曠野者去高唱.
他力明師爺會說.
你們要在他面前歡樂.
駕車經過曠野.
這句話很特別.
但其實可以翻譯成.
是教訓者.
或者應該翻譯成教訓者.
因為有其他抄本的字母是不同的.
有哪個字母是不同的呢?.
你看到這裡.
駕車經過曠野的這句話.
就是Law of Cat.
Baa-ra-foo.
Law of Cat就是駕駛.
駕駛者.
那個字母是什麼呢?.
Law就不用理會.
Law就是前置詞.
To the one, to.
unto的意思.
然後Rokef.
就是你拿走的響音.
就只剩下三個字母.
就是RK和B.
Baa-ra-foo.
這個就是曠野的意思.
In the 即是Baa.
Our vote 就是曠野.
就是駕駛經過曠野的那一位.
但其實有些抄本.
或者應該在烏加利的背景下.
其實你看到這裡.
如果翻譯成教訓者.
背後的原文.
其實差不多一樣.
除了尾三這個字母.
本身就是一個B.
但有些抄本.
就是寫成P.

$^{81}$所以其實讀音差不多.
Law of Cat.
Baa-ra-foo.
或者Pot.
所以幾乎一樣.
B就變成P.
其實嘉臨的語文.
希伯來文是其中一樣.
烏加利文是其中一樣.
還有很多西閃族的語文.
很多時候B和P是沒有分的.
所以很多時候搞亂了.
有時候本身是P.
有時候就說了B.
本身說了F.
有時候就說了V.
等於我們今天.
譬如英文聖經的名字叫Peter.
比特.
中文就叫比特.
但如果其他的名字是用B字開始的話.
我們都是用比的.
B和P我們是沒有分的.
如果你這樣看的時候.
就明白了.
為什麼會是教訓者.
是make sense的呢.
因為有烏加利的背景.
剛才我們看了一連串的.
Bell Cycle.
或者Bell的Epic.
這些史詩.
裡面所形容巴力是怎樣形容的.
就是這一句說話.
這個巴力的名號.
RKB.
一撇.
其實就是一個iron.
對希伯來文就是一個iron.
一撇.

$^{121}$然後是RPT.
Rage.
Pay.
Tough.
因為這個緣故.
所以我們make sense一點.
去將它譯作.
向上帝唱詩.
歌頌他的名.
向那家混者.
高歌.
他的名字.
不是巴力.
是耶和華.
你們要在他面前.
來歡樂.
所以這些背景.
可以暫時幫助我們.
以後我們會再看多一點.
這些烏加利的.
史詩等等.
其實已經翻譯了.
其中一個的edition.
其中一個的編輯.
當然翻譯了做英文.
剛才我們所翻譯的巴力神話的中文.
就是一些很.
粗略的翻譯.
但是.
這個其中一本.
叫做Rogaritic Narrative Poetry.
將剛才的巴力神話.
或者那些史詩.
將它翻譯出來.
最重要的文獻.
我們知道原來看完這些東西.
我們發現.
詩篇68篇.
想講什麼呢.
他就用當時的人.

$^{161}$聽起來.
教運者當然一定是巴力.
整個社會.
整個文化.
走在街上.
每個人都會說.
巴力神明給我們風雲風雨.
詩篇68篇.
或者其他的詩篇.
告訴會眾.
告訴讀者.
不是巴力.
真正的教運者是誰.
是耶和華自己.
這些就叫做爭辯.
或者.
好像我們今天所講的.
apologetics.
什麼叫護教學.
大家就知道護教學.
apologetics.
護教學.
其實就在.
比特前書第三章.
裡面出來.
這裡說什麼呢.
他說只要心裡面奉主基督為聖.
尊他為主.
有人問你心中盼望的.
理由.
你要常作準備答覆.
這裡就是.
apologer.
這個字.
就好像今天英文.
apology的意思.
說對不起.
道歉的意思.
我講過apology.
但其實apology本身不是這個意思.

$^{201}$真正apology的意思.
或者希臘文的背景.
其實就是.
在一個法院上面.
一個人被.
誣告的時候.
他要給一個apology.
一個答覆.
這個答覆就是一個.
defense.
去保護自己.
或者去辯護自己.
這裡說什麼呢.
他說當時的教會.
他們正在面對不同的誣告.
有人去.
誣告他們.
他們信仰的東西.
等等.
或者誣告他們.
彼得就說什麼呢.
要常作準備.
準備去答覆.
給一個apology.
不過我們的心態是什麼呢.
要以溫柔敬畏的心來回答.
要存無愧的良心.
在你們的.
在你們被誣告的賀詩上.
在那凌辱你們.
在基督裡面有很不幸的人.
會自覺羞愧.
他就說出.
真正的apology是什麼.
是我們生命的apology.
所以這個.
舊約聖經裡面.
所說的巴力和耶和華.
之間的爭辯.
其實是一種apology.

$^{241}$一種apologetics.
和新約的大同小異都是一樣.
我們需要作答辯.
這種叫做.
一種護教學.
護教學就是提出一種.
回答.
回答很多時候是要去否定.
這個世界所說的.
而去.
講清楚真正的信仰.
聖經的信仰.
是說什麼.
回答方面也有另外一方面.
回答apology.
也有一種.
所謂.
去揭示.
世界的.
謊話.
文化裡面.
給我們的謊話.
世界意識形態給我們的謊話.
expose那些lies.
cultural lies.
然後去到說出.
真理的美善.
提出信仰的美善.
這個就是我們每一個.
所以我們說每一個基督徒.
其實都是一個apologist.
一個護教學者.
當然講得清楚不是和人爭辯.
而是.
溫柔敬畏心回答.
並且有一個好的生命.
使對方能夠被感染.
甚至那些冒過我們的人.
能夠自覺收攜.
這裡就是.

$^{281}$一種.
我們可以應用到的東西.
所以以後我們可以多看一些巴黎.
宗教的東西.
然後我們再看看.
我計劃會繼續看.
這篇其他地方和.
以利亞的故事.
以致我們更加清楚聖經在做什麼.
我們應用到.
聖經給我們的.
信息.
\newpage



\section{13:1-15}
\label{sec:WPZJt_sBFNI}
\textbf{Final Days of Jesus: Example of The Cross (John 13\_1-15) - Rev. Mark Chan}
\newline
\newline
連結: \href{https://youtube.com/watch?v=WPZJt_sBFNI}{\texttt{https://youtube.com/watch?v=WPZJt\_sBFNI}} ~~~~ 語音日期: 2023-03-27
\newline
\newline
\hyperref[sec:5_8ZAI7O6Mc]{\small{< < < PREV SERMON < < <}}
~
\hyperref[sec:index_chronic]{\small{[返順時目]}}
~
\hyperref[sec:index_scriptual]{\small{[返順卷目]}}
~
\hyperref[sec:9Bv9h74L7iM]{\small{> > > NEXT SERMON > > >}}
\newline
\newline
$^{1}$you up to the service. Thank you guys. Oh, Mark..
I want to say hi to one another. Give a handshake or a high five, you know, to your.
neighbors. What a blessing to be here and to worship together..
That's it. That's it. OK..
As you can tell, I'm suffering from a very mild form of coming cold, so I usually.
don't wear a mask, but, you know, I wear a mask..
So people are pretty surprised this morning and asked me why I'm wearing a mask..
You know, as you can tell, you know, I don't want to spread germs..
So but it's very mild form of coming cold..
I'm coming to the end of it..
So so praise God for for the opportunity to to worship together because, you know,.
showing up here, you know, not being sick, not being sick enough to be to be absent..
It's a blessing, right?.
It's a blessing to sit with, you know, brother and sister in Christ to worship.
together and to think about the word, especially in this coming weeks as the.
Holy Week approaches us..
The Good Friday, the Palm Sunday, next Sunday and then Good Friday and Easter.
Sunday is the best time to reflect and to think about what the love of Christ.
really means for us..
And there's nothing that is better is to go into the word of God, to come to the.
scripture. And last week, as you remember, we were talking about the last days of.
Jesus, the last week of Jesus..
And we were at chapter 12 of the Gospel of John..
And we brought out a couple of points when we studied the story in chapter 12 of.
Gospel of John, the Mary of Bethany..
He annoyed the feet of Jesus with the most expensive perfume and she broke the jar.
and just started to annoy Jesus..
And the purpose of it is to prepare for the death of Jesus Christ..
So, in the perspective of Mary, she knew the cross was necessary..
So, there was the first point..
The first point is a big point..
The necessity of the cross..
She saw that Jesus had to die, had to die to pay for the sin, for sinners to.
accomplish the work of salvation on the cross..
So, the cross was necessary..
Jesus has to die..
That's why she prepared Jesus' body ahead of time..
And then, in the same story, we see another character, Judas, who betrayed.
Jesus, and we see the foolishness of the cross..
And certainly, for Judas, he could not see any wisdom from the cross..

$^{41}$And he was following a leader, the Messiah, for his political need to provide.
some sort of practical solution to his issue, to the issue of his day..
It's the political problem..
The Roman Empire was the political problem for the Jewish people..
So, they were looking for a political Messiah..
Of course, for Judas, this Messiah, Jesus, does not fit the profile that.
they were looking for..
So, there he goes..
He only saw the foolishness of the cross..
So, now we are on chapter 13..
Chapter 13..
Let me read the scripture..
It's a familiar story..
I'm sure you heard about it..
So, let's read John chapter 13..
"Before the Passover celebration, Jesus knew that his hour had come to leave.
this world and return to his Father..
He had loved his disciples during his ministry on earth, and now Jesus loved.
them to the very end..
It was time for supper, and the devil had already prompted Judas, son of Simon.
Zechariah, to betray Jesus..
Jesus knew that the Father had given him authority over everything and that he had.
come from God and would return to God..
So, he got up from the table, took off his robe, wrapped a towel around his waist..
He poured water into a basin, and he began to wash the feet of the disciples,.
drying them with the towel he had around him..
When Jesus came to Simon Peter, Peter said to him, 'Lord, are you going to wash.
my feet?'.
Jesus answered, 'You don't understand now what I'm doing, but someday you will.'.
'No,' Peter protested, 'you will never, ever wash my feet.'.
Jesus replied, 'Unless I wash you, you will not belong to me.'.
Simon Peter exclaimed, 'Then wash my hands and head as well, not just my feet.'.
Jesus replied, 'A person who has bathed all over does not need to wash except for.
the feet to be entirely clean..
And you disciples are clean, but not all of you.'.
For Jesus knew who would betray him..
This is what he meant when he said, 'Not all of you are clean.'.
After washing the feet, he put on his robe again and sat down and asked,.
'Do you guys understand what I'm doing?.
You call me teacher and Lord, and you are right because that is what I am..

$^{81}$And since I, your Lord and teacher, have washed your feet,.
you ought to wash each other's feet..
I've given you an example to follow..
Do as I have done to you.'.
This is the word of God..
Pardon my extra bass voice..
I know it's sexy, but....
Yeah, I struggle whether I should preach today, but, you know, I got a little bit better..
So thank God for answering my prayer..
Otherwise, I will call on some of you to substitute the preaching today,.
especially you, Moody students..
I'll call on you Saturday night and say, 'I cannot preach, okay?.
You should preach.'.
All right..
So God bless you today..
All right..
Now, back to the story, okay?.
The story..
Jesus washing the feet of the disciples..
Very familiar, right?.
There are two major points here I want to point out..
First one is that when we talk about foot washing,.
so often we can think of is this iconic painting, all right?.
What is the name of this painting?.
Last breakfast, lunch, or supper?.
Supper, right?.
Supper..
Da Vinci's..
So it was back in the 1600, I think..
Yeah, he drew this iconic picture..
And as you can see in this picture, 15th century painting,.
one long dining table, and Jesus is at the center..
Can you see Jesus?.
Yeah..
He's right at the center, and six disciples line up on his left,.
and then six others on his right..
So this painting reflects the dinner setting at the time of the painting,.
15th century, 16th century, the time of Renaissance..
European setting, European dining table, right?.
Something we are more familiar with, the dining setting..

$^{121}$But as you know, in Jesus' day, the setting of supper is quite different..
And you remember, as we read the story,.
this is probably the last evening of Jesus' spend with his disciples..
And as we know from the storyline,.
John 13 is the last evening, last dinner with his disciples..
And it was probably on Thursday evening, Thursday evening..
And the next day is so-called Passover and the Shabbat..
And Jesus and his disciples gathered together,.
probably in a pretty big room..
And the Bible says it's an upper room..
It's an upper room..
It's on the second floor..
So it's a big dining area..
And they set up for their last supper, Jesus' last supper..
And it was probably a formal Passover meal of Jewish people..
At this time, it's the Passover season,.
and people were getting ready in Jerusalem to celebrate the Passover meal..
The setting of Jesus' day is different..
What we call, when they sit together to have dinner,.
it's not a long dining table..
It's actually something called triclinium..
As you can tell from the word, triclinium is to recline..
And there are three tables or three couches reclined, forming a U-shape..
Now, back in the days of Jesus, under Roman occupation,.
many of the Jewish people actually welcomed some of the Roman lifestyle..
And this is one of them..
This is one of them, triclinium..
So one day, one meal a year, the most important meal, the Passover meal,.
they would dine in the Roman style..
And that is the Roman style they would imitate..
Triclinium is to have all the guests, all the people dining together..
They will recline together..
And then the food would be in the center..
You would recline on your left side, lying on your left side,.
and your right hand would be free to reach out to the food in the center..
All right?.
So it's a U-shape, and then the center is for the waiter or waitress or the busboy..
Right?.
They will serve the food, and then you ask for refill..
The servant will do it for you..

$^{161}$And your legs, as I explained before, your legs will be sticking out..
Right?.
So your legs, your feet, away from the food..
And you stick your feet out there as you enjoy the feast..
And then the servant in the midst of them will do the foot washing..
Will do the foot washing..
Now, something more about the triclinium..
Something more about the Roman way of dining..
Is that the seating arrangement, where you sit,.
where you sit around the triclinium, there's different levels of honor..
This is pretty much the same as a traditional Chinese or Asian culture..
I remember years ago, I visited southern China,.
and then we were meeting with some church leaders and city officials..
So I remember one evening, we went out to a restaurant,.
to a private room, a dining room..
And then we went to the dining hall, and there was a big round table, right?.
The Chinese table, usually round, right?.
So we went into the room, and back then I was a young guy..
Right?.
So I still remember I walked into that room with a bunch of elders.
and people who are much older than me, right?.
And do you know what you should do and how you should behave?.
You do not sit..
You are told where to sit, right?.
So the city officials, the big guys, right?.
They are probably the oldest person in the room, the elder..
He will sit first, and he will assign a seat, right?.
So, you know, beside him, right or left,.
you know, probably other governors or city officials or government workers..
And then come alongside, and then the last one being probably the youngest one..
So it is a culture of honor and shame..
So even where you sit, where you sit, it is a matter of honor or shame..
This is the same way in the Roman triclinium arrangement..
Every seat is prearranged from the elder to the least,.
from the greatest to the least, from the eldest to the youngest..
It's a culture of honor and shame..
So if we map out the triclinium, the host, okay, the host,.
probably a little bit small, the font, but that's okay..
The host will be on the left, sitting second to the left..
And the guest of honor on his left..

$^{201}$And the trusted friend, another guest of honor,.
usually the youngest, on the right side of the host..
And then the seating will continue in a clockwise direction,.
often from the older and then all the way to the youngest..
And then on the lower right, on our lower right,.
would be the seat of the least honor..
And usually the youngest person or the servant.
would be sitting in the lower right..
Now as you can tell, triclinium in the Roman culture.
is a symbol of political power..
It's all about domination..
It's all about who's going to sit at the highest honor,.
who's going to take the least honor..
The ruling class, the leaders, would exercise their right.
and their power over other people..
So you come into the dining room of Roman triclinium,.
you will feel that this is not just a meal..
You know, it's politics..
It's all about power..
It's all about dignity and about honor..
Now if you map this out in such a way.
that Jesus is the host,.
he would be sitting on the second to the left, right?.
He's the host, Jesus there..
And then probably on Jesus' left is Judah..
And then on his right is the youngest of the 12,.
Apostle John..
And then on the other side of the dining hall.
is Peter himself, Simon Peter..
Now why did I say that?.
Because there are different details.
in the Last Supper stories from all four Gospels..
There are different details that tell us.
actually this is how they arrange the seating positions.
in the Last Supper..
Let me give you an example..
Now remember this map, right?.
This is where you sit..
This is where Jesus is..
If you're Simon Peter, where are you?.

$^{241}$You're on the list, right?.
By the door..
So, you know, he is in the position of servanthood..
John, probably the youngest among the 12,.
and we know John died the last among the 12, right?.
He was also the author of the last book of New Testament,.
the Book of Revelation..
So naturally we would think that John is the youngest guy..
Judah is right next to Jesus and sharing the same bowl,.
dipping food into the bowl and enjoy the food together..
So Jesus and Judah's probably are very close to each other..
So if you remember this map,.
you remember this seating arrangement,.
let me turn to one of the verses in chapter 13,.
talking about the Last Supper..
You would immediately understand,.
all details here would immediately make sense..
The story goes, when Jesus says, "Someone will betray me.".
He said, "Among you 12, someone would betray me.".
Verse 22, the disciples look at each other,.
wondering whom he could mean..
The disciple Jesus loved,.
probably John himself, right?.
He's the beloved disciple..
He wrote the Gospel of John..
He's the youngest guy..
Now, how is he being described?.
The disciple Jesus loved John,.
was sitting next to Jesus at the table..
Simon Peter motioned to him, motioned to John and asked,.
who is Jesus talking about?.
But why didn't Peter go ask Jesus himself?.
Well, the only reason is because in the triclinium,.
he's the farthest guy away from Jesus..
He could not ask Jesus personally, individually..
So he has to motion to John..
John is right next to, actually in front of,.
on the right side, reclining in front of Jesus..
So he motioned to him and asked,.
who is Jesus talking about?.

$^{281}$So that the disciple leaned over to Jesus,.
remember they're reclining,.
and asked, "Lord Jesus, who is it?.
"Who is it going to betray you?".
Jesus respond, "It is the one to whom I give the bread.
"I dip in the bowl.".
And when Jesus had dip it, he gave it to Judas,.
not Peter, right?.
If he had to give it to Peter, he has to throw it over,.
right, give it next to Judas..
So immediately you understand,.
the gospel description reflecting the local Roman,.
or Roman Jewish, you know, high class triclinium dining,.
you know, fashion..
They're sitting around the U-shape, reclining,.
feet sticking out, right?.
And then the background give us, you know,.
some more, you know, some more ways to understand.
some of the biblical descriptions..
And I think this background brings about.
some of the surprises in the story..
Some of the big surprises..
Perhaps the story is too familiar to us,.
and that sense of surprise, that sense of, you know,.
ha-ha, you know, it's gone..
But let's revisit..
You know, in those days, everybody wear sandals..
Everybody wear sandals..
Anyone wear sandals today?.
Open toes..
You wear sandals and you go to the dining hall,.
from your home to the dining hall,.
your feet bound to pick up some dirt, right?.
Stinky feet, dirty feet..
And foot washing in those days.
is a necessary part of dinner arrangement..
You gotta have someone do the foot washing for you, right?.
Usually people of lowest ranking class,.
servants, slaves, women, or widows,.
they are being tasked to do such things..

$^{321}$They are the waitress, they are the busboys,.
or they are the servers..
So they are the foot washers as well..
Now, then you can understand why.
in some of the gospel stories,.
talking about in the Last Supper,.
the disciples were arguing against one another,.
who is gonna be the greatest,.
and who is gonna be the least, right?.
Remember some of those stories?.
You know, why did they argue about it?.
Well, there's a practical need to argue about it..
You know what I'm saying?.
Someone has to do foot washing, right?.
Now, Jesus is the host, He arranged the entire dinner..
What He missed doing, He didn't do one thing right..
Jesus didn't do one thing right..
He did not arrange a servant among them..
He did not hire a servant, right?.
He did get the Airbnb and get the dining hall ready,.
but nobody's gonna serve, right?.
So all the 12 disciples show up in the dining hall,.
what food is being served,.
but there's no one among us.
to do the most necessary part, right?.
Because our sandals, our feet,.
when we took off our sandals,.
our feet are full of dirt, right?.
And part of the manner of dining in Trichilium.
is to have your feet washed, right?.
Think about our culture, the dining custom..
When we sit down in, say, a Chinese restaurant, right?.
What's the first thing that is necessary.
when you sit down?.
Tea..
Somebody has to bring tea or ice water..
Somebody has to do that, right?.
You expect that..
It's kind of our cultural expectation or cultural values..
In Jesus' day, around the Trichilium,.

$^{361}$people would expect the first thing..
We come in, a servant will wash our feet, right?.
That's why they argue against one another..
Who's gonna do it now, right?.
Oh, Jesus didn't, forgot probably,.
forgot to book a maid servant, right?.
And nobody would be so disrespectful.
and say, "Teacher, you missed something.".
Nobody would do that..
That's why among the 12, they would have to argue..
They have to figure out,.
"Hey, when we got to have our feet washed,.
"who's gonna be the one, right?".
And then as they come to the table, surprise, surprise, right?.
Jesus put all the 12 in a clockwise position..
On this right is youngest disciple, John..
On his left is Judas..
And then clockwise, and then to the last one,.
is Simon Peter..
Now that's a surprise,.
because Peter has been the head among the 12, right?.
He is the best Sunday school student, right?.
He always answer the questions often wrong,.
but he's the most active student among the 12..
And he's actually the one of the closest to Jesus..
And surprise, surprise to Simon Peter..
He's sitting right next to the door, the least position..
The surprise is, Jesus is doing the foot washing..
Jesus got up and he did it..
And look at what Jesus did..
Jesus knew that the Father had given him.
authority over everything,.
that he had come from God and would return to God soon..
So Jesus got up from the table, from the triclinium..
He took off his robe,.
wrapped a towel around his waist,.
and poured water into a basin..
There, then he began to wash the feet of the disciples,.
and then dry them one by one with towel around him..
Now, as you can tell, I highlight all the verbs here..

$^{401}$If you read the Gospel of John from chapter one.
all the way to this chapter, chapter 13,.
you understand the narrator or the author,.
it's like a cameraman..
It's like a cameraman..
At this point, it was pretty fast.
from chapter one to chapter 12..
Boom, boom, boom, miracles, miracles, miracles, miracles..
And talking, talking, talking..
All the controversials about Jesus,.
all the teachings about Jesus..
One chapter after the other, pretty fast, pretty fast..
And then come to chapter 13, it comes to a slow-mo..
You know, just zoom in to the actions, right?.
You know, all the verbs, you know,.
just two verses, four and five..
He rose, he lay aside, he took up a towel,.
he wrapped around his waist, poured water out,.
and then washed their feet, and wiped dry their feet..
It is not domination..
It is not Roman politics..
It is servanthood..
The greatest becomes the least..
It's all about serving..
Now, do not miss, please do not miss.
the significance of this story..
It's not just a nice, warm story about servanthood,.
about Jesus washing the feet of a disciple..
It's not just an act of servanthood..
It foreshadows something else..
It foreshadows something that's gonna happen soon,.
which is the crucifixion of Jesus..
He took off his majesty..
He stood up, he took it off,.
and he took on the form of servant,.
and then he died on the cross..
He poured out his blood this time, pretty soon..
On Calvary, not water, but like water,.
he poured out his blood to be the least..
And remember in one of the seven words,.

$^{441}$one of seven last words of Jesus on the cross,.
one of them in the Gospel of John.
is that Jesus said, "I am thirsty.".
The living water become thirst..
So foot washing is not just about serving..
It points us forward, or what we call foreshadows,.
the cross..
Jesus washing the feet of disciples is to prepare them.
what is he doing ultimately on the cross..
That's the first point..
The second point is that foot washing points.
toward the way of disciples,.
the way of how disciples should live..
Now why do I say that?.
I mean, I think the point is obvious.
because in washing the feet of disciples,.
Jesus demonstrates so clearly the true meaning of service,.
or better, the better way to put it,.
is the true meaning of using his power..
We all have power..
We all have our rights, our privileges..
You know, especially in American culture,.
we emphasize power..
We emphasize our rights..
Nothing wrong with that per se,.
but as we follow Christ,.
Christ gives us another way of using our power,.
using our rights..
There are a few points we should point out.
in the last supper,.
especially in the triclinium arrangement..
First, is that Jesus puts Judas in the seat of honor.
on his left, on his left..
He took a piece of bread,.
dip it in the bowl,.
and give it to Judas..
Judas is not sitting just on his left.
in the seat of honor,.
but for the very first time,.
or very last time,.

$^{481}$Jesus gives him one less reason to betray him..
One less reason..
Loving the enemy..
Demonstrate clearly..
Second point..
I said John should be the youngest among the 12, right?.
Jesus put him first on his right..
The youngest, the least, is put on the first..
Third, Peter, far right, lower right,.
has been head of the 12,.
but now Jesus put him last..
The entire Jesus triclinium arrangement.
is some sort of anti-Roman political story..
Jesus said, "This is how you work in the Roman world,.
"but when you come around my table,.
"when you be my disciples,.
"we work in quite different ways.".
It's an upside-down world, right?.
Jesus turned the whole world upside down..
The head would be the last, right?.
The first would be the least..
The oldest would be the servant..
You know, what does it talk about, discipleship?.
What does it talk about our church, right?.
We have a short series on being a healthy church, right?.
I think the most healthy church,.
most healthy church culture.
would be a culture of Jesus' triclinium,.
the Last Supper..
When we come together,.
all powers are being reconfigured..
Everything has changed..
The configuration of structure,.
social, political power, spiritual power, or whatever..
When you come to the table of Jesus, everything has changed.
because He came to serve, not to be served..
There's a book, a very good book,.
written by Scott McKnight..
He used to be a New Testament teacher.
at Northern Seminary, not too far away from us..

$^{521}$He said, the title of the book is "A Church Called Tov.".
Okay, what is it?.
Tov is actually a Hebrew word, means good, goodness..
All right, good, tov..
So the book talks about a healthy church culture,.
a tov culture, a goodness culture..
How do we form a goodness culture,.
especially in our generation?.
When people look at church.
or established religion,.
organized religion,.
they often look at church in a negative way..
There's just so many horrible stories.
done in the name of God, in the name of the church..
It's just sobering to hear all the stories..
And Scott McKnight wrote this book,.
and he lays out four ways of power..
He said the first way is power over..
You have domination over other people..
You're in charge, and other people listen to you..
The second way is power to..
You have the power to influence other people..
And the third way is power with,.
to share the spotlight,.
to share power and your platform with other people..
Let other people speak, let other people lead..
Empower another person, right?.
Power with..
But the ultimate form of, but but call, exercise of power,.
Jesus' way, is power for..
That's what Jesus did..
He demonstrated to his disciples,.
to serve, but not to be served..
The greatest shall be the last..
Greatest shall be the least..
Now, some of the homework that I wanna give to you..
What does it mean to apply,.
to use our power and use it for, you know, power for?.
Couple of points..
First, forgiveness..

$^{561}$Have you ever thought about forgiveness as an act of power?.
You decide to forgive someone who has wronged you..
It is an act of power..
And biblical forgiveness is an exercise of power for..
You decide not to retaliate..
You decide to forgive, just like Jesus did..
Forgiveness..
Second, giving..
Giving is also exercise of power for, right?.
Instead of taking..
Taking, accumulating, is power over..
But in the biblical lifestyle, in discipleship,.
we ought to exercise and continue to exercise power for..
To give is an act of power for other people,.
for the church, for other ministries..
Third, serving..
Serving definitely, ministry definitely, is powerful..
Let me ask you a question, a very quick question..
Have you ever done anything to serve other people lately.
without telling anyone?.
Without telling anyone, even your wife or husband,.
your best friend, nobody knows..
You give, you do something for other people, you know?.
(audience laughing).
Powerful..
That is the way of discipleship..
That's how we should act..
That is the whole point of Jesus' foot washing.
for the disciples..
Let me close with one very short quote, okay?.
Very short quote..
Back in the 1700s, 1700s, 200 some years ago,.
there was a German missionary..
His name is Sensendorf..
It's not very well known..
Nobody would have interest to know about his life.
or what he did..
But there's one saying that he said.
that has been echoed.
or has touched so many lives..

$^{601}$To remind people what it means to follow Jesus Christ,.
what it means to be his disciples,.
what is the purpose of our life.
as we commit our life to Jesus Christ himself,.
to be Christ's follower..
What does it mean?.
What is the ultimate purpose?.
And I wanna end today's sermon with this quote..
It's a very short quote..
The quote said, "Preach the gospel, die, be forgotten.".
Be forgotten..
That is the ultimate goal of Christian discipleship..
Now, we don't have any problem.
with preaching the gospel usually, right?.
Probably we are afraid of doing so,.
but that's not a big issue, preaching the gospel..
And we usually have no problem dying, right?.
Someday you die, someday I die, right?.
You have to die, you know, someday..
Usually you have no problem with the second part..
Usually we have problem with the last part, the third one..
Be forgotten..
You wanna do something.
and you wanna prolong so-called your name, right?.
Your reputation or your power..
But this story in John 13.
tells the ultimate purpose for discipleship..
Follow Christ..
It's not about us..
It's not about us being served.
or us being, you know, feeling good..
It's about us doing something to respond.
to the grace of God with our life.
and ultimately be forgotten..
And that's okay..
Let's work on it..
Give our life back to God,.
God who save us, who redeem us..
And let's work for His glory,.
for God's name, not my name, not your name,.

$^{641}$not even our church's name, not our ministry's name,.
but for the name of God..
Let God's name be glorified..
Let Him get glory and we be forgotten..
Let's pray..
Father God, we give you thanks.
because you have demonstrated your love to us..
You are the one who is greatest.
and you willingly, sacrificially lay down your majesty.
to come into our midst,.
took on the form of humans,.
took on the form of servant, slave..
You pour out your life to save us, to redeem us,.
to bring us back to our Father..
And we want to give you thanks..
I pray that this is our prayer,.
to do the work, to preach your gospel..
Then we die and it's okay to be forgotten.
because ultimately your name is what matters.
in Jesus' name..
[MUSIC PLAYING].
\newpage



\section{以弗所書 4:11-16}
\label{sec:9Bv9h74L7iM}
\textbf{我健康,成長嗎?(以弗所書4\_11-16) - 吳君才牧師}
\newline
\newline
連結: \href{https://youtube.com/watch?v=9Bv9h74L7iM}{\texttt{https://youtube.com/watch?v=9Bv9h74L7iM}} ~~~~ 語音日期: 2023-03-27
\newline
\newline
\hyperref[sec:WPZJt_sBFNI]{\small{< < < PREV SERMON < < <}}
~
\hyperref[sec:index_chronic]{\small{[返順時目]}}
~
\hyperref[sec:index_scriptual]{\small{[返順卷目]}}
~
\hyperref[sec:k6VIyucQ7dI]{\small{> > > NEXT SERMON > > >}}
\newline
\newline
以弗所書 4:11-16
\newline
\begin{longtable}{cl}
\hline
\hline
章節 & 經文 (和合本修訂版)\\
\hline
4:11 & \begin{tabularx}{0.7\textwidth}{X} 他所賜的有使徒，有先知，有傳福音的，有牧者和教師， \end{tabularx} \\ \\ \relax
4:12 & \begin{tabularx}{0.7\textwidth}{X} 為要裝備聖徒，做事奉的工作，建立基督的身體， \end{tabularx} \\ \\ \relax
4:13 & \begin{tabularx}{0.7\textwidth}{X} 直等到我們眾人在信仰上同歸於一，認識神的兒子，得以長大成人，達到基督完全長成的身量。 \end{tabularx} \\ \\ \relax
4:14 & \begin{tabularx}{0.7\textwidth}{X} 這樣，我們不再作小孩子，中了人的詭計和欺騙的法術，被一切邪說之風搖動，飄來飄去。 \end{tabularx} \\ \\ \relax
4:15 & \begin{tabularx}{0.7\textwidth}{X} 我們反而要用愛心說誠實話，各方面向著基督長進，連於元首基督， \end{tabularx} \\ \\ \relax
4:16 & \begin{tabularx}{0.7\textwidth}{X} 靠著他全身都連接得緊湊，百節各按各職，照著各體的功用彼此相助，使身體漸漸增長，在愛中建立自己。 \end{tabularx} \\ \\
[1ex]
\hline
\hline
\end{longtable}
$^{1}$頂止無懼平安.
平安.
有一個人.
好幾年前.
信了耶穌.
但他撕不掉打麻將.
(笑聲).
經常攻打四方城.
(笑聲).
很少很少.
讀經祈禱.
亦間中回教會.
在九個月之前.
他的兒子.
三十幾歲.
突然呼吸覺得有問題.
胸又痛.
頭又痛.
去了急症室.
擾攘了好幾天.
做檢查.
最後發現他患上了.
罕有的癌症.
而且是末期.
在一個星期之內.
他就匆匆離開了世界.
這個突然的打擊.
喪子之痛.
令到這個人.
情緒很低落.
抑鬱.
難過.
傷痛不已.
好幾個月.
教母關懷和開解他.
慢慢才走出.
這個傷痛.
和抑鬱的陰霾.
之後.
他有少少的.

$^{41}$扭轉.
他要調整自己的人生.
未來方向是怎樣.
要更多依靠神.
親近主.
要更加積極回教會.
他願意做一個健康成長的基督徒.
亭姐妹.
做個健康成長的信徒.
仍然是你和我今時今日.
要呼應的.
這個亦是神的期望.
我們經常要檢視一下.
究竟我健康成長嗎.
我們有個簡單的祈禱.
天父監察我們.
聖靈提醒我們.
不住要健康成長.
求你帶領我們以下.
聽到講到的時間.
祈禱奉主耶穌基督的名.
阿們.
今日我願意.
我願意.
用.
以弗所書四章.
一到十六節.
作為彼此勉勵的訊息.
經文已經打出來了.
你看一看.
時間關係.
因為講的時候也會分別講.
跳得這麼快.
猜到這裡按不到那裡.
知道老人家手腳沒那麼靈敏.
一到十六節.
十一節到十六節.
提醒我們.
要自己.
追求靈命健康成長.

$^{81}$我今天會用.
兩大方面和大家思想.
第一要健康地善用恩賜.
第二就要.
有成長的表現.
要健康地善用恩賜.
用第十一節那裡來看.
首先我們要再想一想.
自己的恩賜和慮別.
在哪裡.
這裡說他所賜的有使徒.
有先知有傳福音的.
有牧師和教師.
有五樣東西.
你也聽過.
使徒原文是鳳差衍的意思.
使徒不只是一種恩賜.
更加是一種積分.
羅馬書十一章十三節這麼說.
除了十二使徒.
由主耶穌親自設立.
之外.
在馬可福音三章十三到十六節有記載.
還有聖靈所設立的使徒.
譬如保羅.
巴拿巴等.
使徒行傳十三章一到四節.
八節十四章三到六節.
按聖經所說.
五重節之後使徒的主要.
有三個任務.
第一是見證基督的復活.
使徒行傳一章二十節.
第二是建立教會.
以弗所書二章二十節.
第三是確定真理的標準.
加拉太書一章九節.
提摩太后書一章十四節.
彼得后書三章十五到十六節.
使徒也見過主耶穌的復活.

$^{121}$臨前十五章六到七節.
並且有權並行神蹟奇事.
哥林多后書十二章十二節.
事實上.
一些開方直堂.
創立教會的人.
大多數有使徒的恩賜.
他們有.
倚仗.
有膽識.
有使命感.
你認為在我們中間.
我們成會教會的陳牧師.
是不是這類呢.
第二先知.
在這裡所說的先知.
既然是基督賜給新約教會的.
五種恩賜之一.
當然就不是指舊約的先知.
你知道舊約的先知.
最明顯的一種工作.
就是照神的旨意.
講預言.
或者預言警告.
和勸導當時的百姓.
而新約的先知.
主要的工作.
按臨前十四章一節所說.
是作先知的講道.
作先知的講道.
主要的.
按臨前十四章三到四節所說.
是要造就.
安慰 勸勉人.
換言之.
是一種造就教會.
教導人 建立人.
的恩賜.
臨前十四章十二.
二十九 三十三節所說.

$^{161}$弟兄姊妹.
不妨你又檢視一下.
自己 或者我們當中.
是不是.
有些人.
很懂得關心人.
懂得安慰人.
也有教導造就的恩賜.
當你發覺.
或者你認識的人.
他們.
很有愛心.
用行動.
或者說話來幫助人.
使得人得無見能.
你就要想想.
這個也可能是恩賜.
第三.
傳福音的.
就是特別有傳福音的負擔.
和恩賜的人.
你見到.
這類人很有.
心去傳福音.
去分享福音.
通常這類人.
有某種吸引力.
而且很容易.
甚至很喜歡去接觸人.
他們很熱情.
聖經裡面.
使徒.
譬如彼得.
保羅大多數有.
傳福音恩賜.
但是有一個人.
聖經說稱為.
傳福音的肥利.
他不是使徒.
他只是七個執事裡面一個.

$^{201}$《素能傳》.
21章8節這樣說.
一個.
執事.
一個資深的信徒.
一個堂委.
或者call team.
核心領袖.
甚至團契職員.
原來都可以成為.
福音的得力助手.
你有沒有善用.
這個恩賜呢.
第四.
恩賜是牧師.
原文是牧羊者.
主色書說牧師.
這個字不單止有餵養羊群的意思.
也有看顧照顧羊群.
給牠們關懷鼓勵安慰.
牧師是神所賜給教的.
一種恩賜和職範.
是要負上很大.
很重要的.
屬靈責任.
所以弟兄姊妹.
你要為牧師.
好好的祈禱和用行動.
來支持牧師.
完成牧師的任務.
第五.
是教師.
是教導人明白聖經有解經恩賜.
能夠造就信徒.
保羅巴拿巴在未受聖靈.
參見之前和安提拉教會.
裡面的西面路球.
馬廉.
同被稱為.
先知和教師.

$^{241}$今天在教會裡面.
一些喜歡訓練信徒.
教導人的.
一些事奉的工作.
應該都是有教師的恩賜的.
一個能教導人.
明白神的話語.
堅固信徒.
以及激發他們去愛主的.
這些都是教師的恩賜.
我相信.
在我們中間在座的.
或者在網上的.
弟兄姊妹.
你一定擁有.
你體會到.
神賜給這種恩賜積分的原意嗎?.
其實.
在聖經裡面.
在哥林多前書十二章八到十節.
更加提及聖靈賜給.
有一些人.
有智慧的言語.
知識的言語.
信心 醫病的恩賜.
行義能 辨別諸靈.
說方言.
翻方言的恩賜.
而羅馬書十二章六到八節.
講到.
有作則事的.
治理的恩賜.
勸化施捨.
憐憫人的恩賜.
弟兄姊妹.
這些的恩賜.
你有沒有想過.
原來都是可以.
發揮出來.
能夠健康地.

$^{281}$去教會的呢?.
我們看第二點.
健康地發揮恩賜.
的目的在第十二節.
為要成全聖徒.
各盡其職.
建立基督的身體.
那個.
的成全.
有個裝備的意思.
裝備得很好.
為神所用.
今天.
我們會發現到.
很多人有恩賜.
但卻糟蹋了埋沒了.
這個令神很傷心的.
是有一些.
不健康的阻滯.
或者我稱之為障礙.
我用信徒.
甲乙丙丁來講一講.
他們可以代表四類的人.
第一類.
甲.
是一個姊妹.
以前是信民間信仰.
或者佛教.
信了耶穌.
說我信.
但是.
你卻發覺原來.
他沒有將.
耶穌.
擺在自己人生的主權裡.
他接受耶穌只是救主.
不當耶穌是自己人生的主宰.
所以他仍然用以前的生活方式.
處事為人的心態.
實際上.

$^{321}$他未完全將自己的人生主權.
交給神管理.
他未被主得著.
或者反過來講.
他自己仍然未能.
完全得著.
因為.
他未曾真正看到.
基督是我人生至寶的.
他的心眼只是其他事物.
屬世的那些.
他根本.
完全未肯接受.
聖經的教導.
或者讓神完全掌管.
所以講恩賜這件事.
他是完全不理的.
信徒月.
他是一個.
專業人士.
他很聰明.
可以說.
好像小世界一樣.
也好像以前大老倌.
梁醒波所講的.
操世界要聰明.
就當玩把戲.
他遊戲人間.
樣樣.
頭頭是道.
不過他也是一個很好的家長.
愛子女 愛家庭.
他一邊說.
愛主 其實他一邊.
都是繼續.
貪戀世界的事物和活動.
他很在意.
一些東西.
其實外面的.
名利 虛榮.

$^{361}$世界很多的玩樂.
牽引他很大.
所以有時.
去了那些都未必.
著重教會那些事物.
或者恩賜那些事物.
其實你和我都知道.
人若愛世界.
愛父的心.
就不在他裡頭.
是不是.
雅各更講得.
嚴重.
如果你和世俗為友.
就是.
與神為敵.
這麼嚴重.
他都不理這麼多.
所以有些進途.
也是真的 這個世界.
有時做事.
我真的經歷很多人.
人在江湖 身不由己.
你不是這麼容易做.
除非你裡面有種.
能量或者智慧.
所以.
很多時候我們見.
這類人有很多版本出現.
我們有時稱.
這類人叫一腳踏兩船式的信仰.
骨子裡對聖經真理.
或者對恩賜.
是否未必有興趣.
在心靈的深處.
其實我們又想回.
我們何嘗不是這樣.
我們有時不肯.
將自己的前途.
或者自己一切.

$^{401}$好好奉獻給神.
剛才也有宣傳.
宣佈一下你的奉獻.
甘心樂意.
很多人都未必這麼甘心.
你的時間才剛.
是否甘心.
會捨得好為主呢.
你會善用你這一點嗎.
信徒.
丙.
丙是一個長者.
你聽到有些長者.
他信耶穌.
我信.
不過有時信完之後.
我老了我要事奉.
你不要找我.
你找年輕一點的.
我讀書少.
你不要搞我.
很多人.
類似都是這樣.
不少的信徒.
其實是貪圖安逸.
不肯付出.
怕麻煩 怕困身.
真正的原因.
原來是怕為主吃苦.
怕吃虧.
怕辛苦.
更怕付代價.
特別現在後現代.
這個世代.
很多人要自我感覺良好.
我要對感覺.
我要享樂.
要開心.
搞到我麻煩.
不開心的那些不要搞我.

$^{441}$恩賜事奉的那些.
不要搞我.
信徒丁.
就是另一個人.
年輕一點的.
他其實對自己.
沒有什麼信心.
也對神沒有什麼信心.
他有時信不過自己.
可以承擔任何職務.
不可以的.
但另一樣深層的深處.
他不信任人.
那個牧師這麼懶.
都不聯絡我.
這個傳道人又怎樣.
用一種陰謀或懷疑.
來看一些教會的事.
看的未必是完全正面.
對一些人.
我不理那麼多.
你要我事奉嗎.
什麼恩賜嗎.
理你都傻.
我都不想盡心的事奉.
這些是一些你生命上.
卡住你成長.
或者善用恩賜的障礙.
ABCD.
你見到沒有沒有.
這四個沒有.
其實是你要檢視一下.
究竟你不能夠善用你的恩賜.
是不是因為更深層的原因.
一個老信徒在祈禱會裡面.
突然說出一番說話.
哎呀 神啊.
我要用十言玄的樂器來讚美你.
參加祈禱會的人都奇怪.
明知他根本不會彈任何樂器.

$^{481}$究竟是什麼意思呢.
這個老人家.
老信徒就說.
我們敬拜讚美神的時候.
是用我們雙手去服侍神.
用我們雙腳.
走在神的旨意裡面.
用我們的口.
為神的恩典作見證.
用我們的心來愛神.
這樣.
兩眼 兩耳 兩手.
兩腳 口和心.
加起來就是十樣東西.
這就是十言樂器.
弟兄姊妹.
我們的五官.
我們的手腳.
心口.
各有功能.
為神的緣故.
你願不願意很健康地.
運用發揮出來呢.
我們看另一樣.
第二方面.
我們要有成長的表現.
成長的呼聲.
在第十三節.
直等到我們中人在真道上同歸於一.
認識神的兒子.
得以長大成人.
滿有基督長成的新娘.
在另一個譯本.
是說我們將上帝的兒子.
有一致的信仰和認識.
我們將長大成人.
達到基督成熟的程度.
史圖保羅.
著意地提醒我們.
每一個基督徒.

$^{521}$都要長進 要長大成人.
要達到基督所盼望的水平.
那種程度.
不要再小孩子.
不要那麼靈敏 幼稚.
第十三節.
顯示一個特徵.
如何檢視自己.
有沒有長大成人 成熟呢.
今天有時我們都很遺憾.
或者難過.
有些人信了主十年.
八載 二十年.
三十年.
其實靈敏都是很幼稚.
有些人請他祈禱.
牧師不要叫我.
我雖然信了主二十多年.
我都不懂祈禱.
我聽到 嚇 張開嘴巴.
當然公眾祈禱是一個挑戰.
但私下祈禱.
生活又如何呢.
有沒有追求長進的心呢.
這裡提出一樣東西.
就是你有一致的信仰.
和認識和其他.
的長進信徒.
一致.
這個水平.
今天我們求主幫助我們.
究竟我們是怎樣一樣的人.
我們又看看.
長進成長.
有些特徵.
這裡.
十四到十六節有說.
第一樣.
健康成長的特徵.
就是有定力.

$^{561}$不再作小孩子.
中了人的詭計和欺騙的法術.
被一切異教之風搖動.
飄來飄去.
隨從各樣的異端.
今天.
很多.
的人.
面對不同的衝擊.
在第一世紀的時候.
保羅的時代出現了若干異端.
諾斯底的主義迷惑人.
今時今日.
亦都有些異端.
邪教.
為何人間.
最近.
Netflix紀錄片裡面.
以神之名.
信仰的背叛.
你有沒有看過.
他又說.
南韓有幾個自稱為.
救主先知的教派領袖.
犯罪的行徑.
其中有.
攝理教的部分.
引起特別香港人.
關注更多.
因為三個星期前.
藝人方力申.
他的女朋友.
Mabel 葉萱.
就是以受害人的身份.
控訴遭受.
教主鄭明識.
性侵的經歷.
Mabel 在香港.
遇到.
攝理教.

$^{601}$然後信了教.
異端害人.
異端對正統和基督教信仰構成.
混亂和威脅.
其實在美國 在香港.
仍然都有這些異端很活躍.
你也聽過.
比如夜華的見證人.
守望台 沒數聖徒教會.
即是摩門教.
全能神教會.
統一教這些.
向一些人下手.
兩個多星期前.
有個老姐妹打電話給我.
牧師 那些甚麼.
夜華見證人.
是甚麼.
即是守望台.
我的護理員.
他們的福音.
好熱心.
一三告假.
走去傳福音.
那些異端比我們正統的基督徒還積極.
異端.
搶人.
基督徒 正統的人.
就放人 甚麼都不理.
今天 無可否認.
作基督徒真的不容易.
在信仰的路上.
我們外面有撒旦.
有世界罪惡的誘惑和攻擊.
內部有自己的肉體.
和私慾的困擾.
和身體健康的問題衝擊.
很難令你有定力.
有時我們會受到.
父神的管教和試煉.

$^{641}$所以我們很容易.
沮喪.
容易缺乏喜樂.
我們甚至軟弱到一個地步.
失去對神的恩典.
救勇的信心.
失去了定力.
失去了盼望和勇氣.
曾經有些人對信仰.
以為.
可能我當走的路.
就快要跑盡.
想放棄自己的信仰.
放任自己.
我不知道你今天.
你的屬靈情況落在甚麼光景.
或者是否在不冷不熱.
的信心狀態.
但是希伯來書.
二章一節.
特別提醒我們.
要有定力.
所以我當閱法正從.
所聽見的道理.
恐怕我們隨流失去.
一旦我們不正從.
去聽和行主的道.
你和我活在這個.
隨流失去的世界裡.
就已沒有定力.
唯今之計.
你真的要抓緊真道.
正從閱法.
聽到的.
道理.
這樣才不會被世俗.
或者屬肉體的潮流所捲起.
昔日以色列人.
因為不恆心守神的約.
希伯來書八章九節說.

$^{681}$所以.
神也不理他們.
所以如果我們沒有定力.
不恆心守神的約.
不正從去行.
和聽道.
難怪你的祈禱.
好像覺得神不聽.
為何神好像不理我一樣.
撲問弟兄姊妹.
不要以為.
你所聽到的.
真道.
都和自己沒有益.
或者沒有影響沒有關連.
不是.
每一次聽道.
包括今天你聽道.
你要學習用你的信心去調和.
所聽見的道.
雅各書四章一到三節提醒.
我們這樣做.
雅各書四章十七節更加說.
如果你知道.
要行善.
卻不去行.
這就是他的罪了.
這麼嚴重.
求主幫助我們.
能夠依靠神更多.
抓住神.
回到他的愛和恩典裡面.
好好的有定力.
所以我們要謹慎.
上上留意.
保守自己和神有緊密的關係.
學習.
和其他肢體.
趁著今天.
要天天彼此相勸.

$^{721}$即是互相守望.
互相提醒.
互相獻告.
這樣才有信心的定力.
時間很快.
第二樣要有愛心.
唯有用愛心說誠實話.
凡事長進.
念於元首基督.
這裡提醒我們.
要有愛心.
一個健康長進的基督徒.
一間健康興旺的教會.
每一個成員.
必須要有愛心.
按這裡來說.
其中一樣.
就是能夠以愛心說誠實話.
不只是說.
阿兒淺美奉承的說話.
也不是風花雪月的表面話.
更加不是說.
是非八卦閒言閒語的閒話.
是用愛心.
正直勇氣說.
真心的話.
誠實的話.
誠經真理的話.
當然了.
誠實的話.
是說真相.
要過濾一下.
所說的是不是真的有愛心.
是不是有建設性.
可以造就到人呢.
幫到人呢.
如果不是.
禁戒不說.
這是杜絕教會裡的.
閒話,傳說.

$^{761}$和觀弄是非.
這個德性的方法.
教會裡.
有時很多傳來傳去.
我聽到一些東西.
我都會去背著.
哦,那些都不是這樣的事.
怎麼會這樣傳呢.
健康合神心意的表現不是.
有位姐妹.
在病危當中.
遺漏就死了.
丈夫把握.
這個仍然清醒的時刻.
跟她說.
愛心的誠實話.
去求自己的妻子.
為著她自己過往的表現.
要求她的妻子.
描述她.
這個姐妹在虛弱的病中.
仍然憤憤然.
堅持不肯.
為著多年的積怨.
原諒她的丈夫.
夫妻關係.
家家有本難聞.
很難一時三刻說得清楚.
但她真的很生氣.
當時牧者也在場.
幫看.
哎呀,姐妹.
你都快見主了.
你不要輸人.
你自己也很辛苦.
但這個姐妹病得來.
仍然很有氣質.
No,我不肯.
牧者等了一會.
聖靈給她有智慧.

$^{801}$姐妹.
你把你的右手.
按在你的胸裡面.
左胸.
然後拉著這個姐妹的左手.
拉著手背.
然後.
連連親吻.
她的手背.
一邊禱告.
一邊.
求神.
融化這個姐妹的心.
讓她堅韌的心得到改變.
親吻了兩三下.
奇妙的事就發生了.
她整個人.
變得軟化下來.
眼睛慢慢滴下淚水.
然後慢慢跟老公說.
老公.
我原諒你了.
我都愛你了.
老公被她一說就很激動.
抱著她.
兩個人相擁哭泣.
這個是.
釋放的淚.
是喜樂的淚.
更加是搖樹的淚.
沒多久之後.
這個姐妹就安然去世.
這個是真人真事.
記錄了愛的威力.
搖樹的威力.
真人的.
愛心並且說誠實話.
意義及搖樹實在是基督徒.
重要的功課和有力的見證.
不過很多人.

$^{841}$以前是.
佳偶.
為何今時今日.
變成怨偶呢.
住在同一屋簷下.
卻不能好好交流.
包容.
就是因為沒有愛.
就是因為沒有搖樹.
就是不肯溝通.
就是不肯用愛心說誠實話.
和真心話.
求主幫助我們.
一間.
有健康的教會.
一間.
一個健康的基督徒.
是用愛心說實話的人.
時間所限快快跳第三秒.
是.
是合適地聯絡.
彼此相助.
建立自己.
這裡那節經文.
意思就是.
要我們.
互相聯絡.
當肢體發揮功用的時候.
整個身體.
就會愛中漸漸長大.
時間所限.
我先問你三個問題.
你自己想一想.
第一,你肯不肯和人聯絡.
你肯不肯和弟兄姊妹.
真心地交談.
還是.
封閉自己.
不肯開放.
第二個問題.

$^{881}$在你心底的深處.
你肯或不肯.
各按各職.
彼此相助.
肯發揮自己.
的恩賜才華.
第三,這麼多年來.
你自己想一想.
有沒有建立到自己.
原因在哪裡.
John McArthur.
曾經寫了一本書.
Mask of a Healthy Church.
即是健康教會的標誌.
提出健康教會.
有十二樣的標誌.
其中有好幾樣涉及.
合適地聯絡有關.
譬如,快快說.
你是否一個活躍的.
教會會員.
你肯不肯.
彼此.
去關心.
你肯不肯熱愛教會.
這個家.
你自己肯不肯改變一些意願.
你肯不肯有犧牲的心態.
提出這幾樣東西.
其實都很重要.
如果你肯.
你就可以和肢體有更好的聯絡.
投入衛生.
做到發揮功能.
健康的成長.
今天仍然有些人覺得.
我做好自己就算了.
不用理會其他人.
我自己可以不用.
和其他肢體聯絡.

$^{921}$我也沒有興趣.
理會別人那麼多.
自己那麼多麻煩的事.
這類人.
很不幸.
卻在教會裡.
成為信仰的.
獨行俠.
老死與人不相往來.
但這不是.
神的心意.
求主幫助我們能夠.
健康善用恩慈.
發揮恩慈.
在長進的事上.
有美好的定力和見證.
更加有愛心的流露.
合適的和弟兄姊妹.
好好的回答.
我們誠心低頭祈禱.
神啊.
你希望我們能夠健康.
好好的善用恩慈.
你更加希望.
我們好好的彼此.
與人聯絡.
榮耀你的聖名.
去除我們的難處.
\newpage



\section{}
\label{sec:k6VIyucQ7dI}
\textbf{UV 匯客室- 朱莉傳道}
\newline
\newline
連結: \href{https://youtube.com/watch?v=k6VIyucQ7dI}{\texttt{https://youtube.com/watch?v=k6VIyucQ7dI}} ~~~~ 語音日期: 2023-03-30
\newline
\newline
\hyperref[sec:9Bv9h74L7iM]{\small{< < < PREV SERMON < < <}}
~
\hyperref[sec:index_chronic]{\small{[返順時目]}}
~
\hyperref[sec:index_scriptual]{\small{[返順卷目]}}
~
\hyperref[sec:hADCTqthLUc]{\small{> > > NEXT SERMON > > >}}
\newline
\newline
$^{1}$Hello 各位弟兄姊妹.
我是Joyce 傳道.
今天很高興可以在香港.
去面試朱莉傳道.
不如朱莉傳道你先介紹一下自己.
Hello 各位UV的弟兄姊妹.
我相信應該不會太陌生.
因為之前在網上.
我們有神交過.
Joyce是我認識了很久的朋友.
在芝加哥的時候.
今天在香港見面.
我們剛才也吃了粥.
油炸粿等等.
哈哈哈.
是啊.
很棒.
朱莉傳道之前在芝加哥的三一神學院讀書.
但是後來她就回來了香港.
現在.
不如我直接問你.
你在哪一個教會做傳道.
還有你回來了多久呢?.
我回來了應該也有七年了.
我在香港佐敦的一間教會.
叫做基督中心堂.
在那裡服事.
另外兼職教神學.
也會在其他教會服事.
教書或者是講道.
嗯.
你回來了香港.
嗯.
為什麼會回來了香港呢?.
一方面那邊的讀書告一段落.
然後也覺得離開家裡很久.
因為我之前的學位.
也是在美國讀的.
所以就離開家裡.
而我家一直都在香港.

$^{41}$所以就離開家裡很久了.
也覺得是時候要回來負一些責任了.
所以也覺得很希望人生的下半場.
至少在現在.
會在香港一個熟悉的環境.
來服事.
所以就回來了香港.
至於為什麼我會蒙召呢?.
其實要講的話就要講得很早期.
我信了主沒多久.
神就已經召喚我了.
當時我在下令會裡面.
因為唱一首詩歌.
叫做《主令我何往必去》.
是一首很青年性歌的時代.
暴露年齡系列.
其中一句就說.
主令我何往必去.
這一句我就很感動的.
我覺得這也是我的心願.
但是中間就兜兜轉轉了很多年.
因為當時我只是中學生.
後來就讀大學.
讀完大學又做了一段時間的事情.
但是神都沒有將這個呼召減退.
我特別感動的是在華人的群體.
在中國大陸的.
當時.
所以那一次一蒙召.
唱那首詩歌的時候.
我就傻傻地在詩歌本的下面.
自己寫了五個字.
寫著作中國傳道.
我還保留著那本書.
當然還認得自己的字.
所以應該也不能離開了.
所以就算我後來大學在美國讀的.
但是我一直的心願都是在.
中國或者香港這些地區服侍.
所以讀完大學做了一段時間的事情.

$^{81}$當時在Silicon Valley.
就回到香港讀神學.
誰知道在香港讀完第一個神學學位之後.
神又再給我這個負擔做神學教育.
所以我又再去了芝加哥.
再繼續進修.
最後就回到香港.
也是一個很特別的時機.
回到香港.
所以回看之後.
很多事情的確不是我計劃的.
但是就是這樣一步一步神帶領我.
嗯.
OK.
我想問一下你.
你的服侍生涯裡面.
有些什麼是令你.
就是.
bring you the biggest joy 的呢?.
Biggest joy 是嗎?.
當然最開心是與人同行.
我覺得是一個privilege.
就是你可以在某一個階段.
在某一個人有需要的時候.
你可以跟他一起去經歷神.
一起陪伴他經過.
我覺得這個是事奉裡面最大的喜樂.
就是與人同行這個privilege.
嗯.
可以參與在不同的人的生命裡面.
是.
OK.
如果在.
說起侍工的話.
Biggest joy.
但是另一方面.
最難的是什麼.
就是最大的challenge.
你面對.
是什麼呢?.

$^{121}$當然是一個biggest joy的flip side.
就是.
當與人發生很多衝突的時候.
你能夠與人同行是很大的喜樂.
但是當你與人有衝突.
當你發覺對某些人失望.
或者是受到傷害的時候.
就會發覺也是最大的challenge.
就是最大的hurdle.
就是在事奉裡面.
所以.
歡喜也是人.
受到傷害.
讓你失望也是人.
但是傳達人的工作就是對人的工作.
所以是必須要克服的.
嗯.
我記得我曾經在一個moment.
我也有懷疑過.
我也有想過.
到底其實做傳道是不是一個keep the distance.
就是與人同行.
因為你是一個privilege.
去參與其他人的生命.
嗯.
但是同時.
你也會因為做一個dumpster.
就是很多人經歷的喜怒哀樂.
特別是一些哀或埋怨.
一些憤怒.
一些不開心的時候.
他也會向你講.
通常是比較負面的.
是的.
所以你就會承受了很多這些.
就是人的幽暗,黑暗的一面.
到底我應該怎樣呢?.
我應該怎樣回應呢?.
是keep the distance.
去pray safer.

$^{161}$去keep me safer.
還是應該去更多的信任等等呢?.
後來我想了這個問題一段時間.
後來我覺得.
作為傳道人.
authenticity.
就是being authentic.
是一個很重要的element.
所以我就想.
還是要有authentic love.
就是要對人有基本的誠信.
對人有信任.
儘管可能受到傷害.
儘管可能會覺得.
甚至會被人欺騙.
但是我也覺得.
這樣的付出是值得的.
做在神的身上.
一來祂會保護我.
二來我相信是有獎勵的.
就是ultimately.
嗯 OK.
你回到香港已經一段時間了.
我想大家都知道.
特別這幾年香港有很多轉變.
你覺得在香港的服侍的模樣.
近幾年特別有哪三樣.
你覺得有多大的改變.
影響教會事宮呢?.
首先我又不是很覺得.
我都是近這大半年的看法.
我又不覺得這些情況.
是很unique to Hong Kong.
我又不覺得香港有什麼很特別.
現在傳頓是很難.
現在在香港服侍.
因為我之前也說過一篇.
也給我自己提醒.
其實歷世歷代的信徒.
都面對很多的challenge.

$^{201}$在他們不同的年代.
在不同的地方.
其實他們都面對很similar的事情.
所以我又不覺得.
因為現在可能政治環境.
或者社會氣氛各樣東西變化.
例如你說其中一樣東西.
是很divisive.
就是那些很分裂.
因為政見.
甚至因為疫苗的看法.
大家就變成撐得很狠.
沒有common ground等等的狀態.
但其實我回頭看.
美國也是這樣.
不只是香港是這樣.
其實美國也同樣面對很similar的situation.
所以我又不覺得特別unique to Hong Kong.
當然香港教會也要面對這些問題.
需要思考前路等等.
但是一方面不覺得很unique to Hong Kong.
另一方面我也覺得是的.
作為在這個時候在香港的墓會.
是需要去認清一個事實.
就是不同了.
我想就算不是墓會.
是在香港生活.
要認清一個事實.
不要整天都覺得以前不是這樣的.
以前我可以這樣的.
你就會不斷回頭看.
就會很痛苦.
我就會覺得.
Let's face reality.
Let's face forward.
無論怎樣.
無論環境怎樣變化.
其實整個工作沒有停止.
我們教會發覺.
仍然有新朋友.

$^{241}$有新的需要.
不同的人在不同的階段.
都一樣有那些的需要.
我仍然要attend to 他們的需要.
其實有些東西是沒有變的.
但是在一個新的環境裡面.
怎樣去侮辱他們.
怎樣去餵養他們.
這個是我需要去思考.
需要去figure out 的東西.
但是你說在大環境.
是不是真的有這麼大的變化呢.
其實我覺得面對的現實就是不同了.
可以有多少空間就做多少的事情.
甚至可不可以創造一些空間.
用創意一些的方法.
來做一些事情.
所以.
我也覺得.
因為剛剛.
可以在更個人的層面去分享.
剛剛和Joyce去看一齣電影.
我們兩個.
我發覺其實.
你說宗教是受到影響.
受到一些壓制.
這個是事實.
但是不止宗教的.
其實在例如文藝創作.
這些的範疇.
其實他們也面對一個收窄了的空間.
嗯.
也會想到到底會不會觸碰到紅線.
但是其實在現在香港的文藝創作.
或者戲劇拍戲.
音樂創作那方面.
你就看到.
這些香港本土的人.
就嘗試發揮創意.
是.

$^{281}$在有限的空間裡面.
仍然可以說一些信息.
那些信息仍然可以.
是啊.
仍然可以說一個香港的故事.
是.
仍然可以觸碰到香港本地人.
現在面對一些問題.
他們的心靈.
就在那些位置做.
嗯.
就是嘗試去發揮創意.
仍然在某些空間去做一些.
嘗試去做一些事情.
所以就是.
有變.
但是也有不變.
是.
就是時候去發揮創意.
是.
去更加細心去敏感.
到底現在人的需要是什麼.
然後嘗試去捉摸他們.
嗯.
我突然間想到一個很實際的問題.
譬如在新冠肺炎期間.
很多東西會是在網上.
是.
就是現在慢慢再開放.
其實教會會不會面對一個問題.
就是很多人說.
我寧願繼續在網上.
不回來聚會.
教會怎麼去繼續那個模樣呢?.
當然了.
我剛剛也教完一科神學.
我發覺我有一半的學生都是prefer.
長期是zoom的.
也不開鏡頭的.
所以剛剛交完功課.

$^{321}$其實我也沒有見過那個學生.
我只知道他的名字.
我一方面要尊重他們的preference.
另一方面在教會的模樣層面.
我就發覺.
其實他們有時候不是真的那麼死板.
我就不想出去了.
只不過可能在過去這幾年裡.
習慣了這個模式.
也覺得維持得到一個safe的.
好像一個洞.
讓他可以在那個很舒服地去.
比較comfortable對他來說.
對那個人來說.
可能他比較安心一點.
比較舒服一點.
也未必是因為疫情的原因.
可能他本身已經不是很舒服去跟別人交往.
現在他有這個選擇就這樣.
所以我一方面尊重.
但另一方面我發覺.
只要你reach out to them.
例如有時候我們碰到一些會友.
他回來教會可能有些事情.
他回來拿東西.
你只要跟他多聊幾句.
不如我約你吃飯.
找一次崇拜完之後約你吃飯.
其實他們接觸的.
都很歡迎的.
如果他本身跟你有.
有關係的話.
是的.
或者他願意讓你接觸的話.
其實可能你只要多走一步.
他可能會比較容易回到社群.
而我也覺得這個是有需要的.
人始終是需要有交流.
需要有眼神接觸.
需要有個人觸感等等.

$^{361}$可能就是教會的教育同工.
比之前要更加多主動去reach out to.
特別可能沒有回來聚會的弟兄姊妹.
是的.
當然我們也發覺年齡層是很有趣的.
我們有些接觸一些長者.
他們反而是很願意要出來.
很快當疫情比較減緩.
他們可能是第一批回來教會聚會的那一批.
所以有時候不是關於健康疫情的原因.
可能是他們的情緒.
他們的性格等等.
都會有影響的.
嗯.
OK.
是.
OK.
你有一個問題想問一下你.
作為一個傳道人.
你怎樣平衡自己的需要.
嗯.
你家裡的需要.
和時工.
知道時工很多時候都很忙.
你怎樣平衡你的生活呢?.
嗯.
這個我會很坦白的.
我不會回答模特兒的問題.
因為我的確覺得.
可能在這個時代.
對傳人的要求.
再不是那些24小時on call.
當然我也試過在深夜收到電話.
需要照顧一些事情.
需要幫一些人.
那些當然是緊急事件等等.
但是我也覺得需要有一個平衡的.
所以例如.
身處在香港就有這個優先權.
我們很容易去到其他國家旅行.

$^{401}$所以也會安排一些時間和家人去旅行.
嗯.
也會有些娛樂.
去看電影.
好像剛才看電影.
是啊.
約朋友吃飯.
也會覺得.
會畫一個邊界.
是需要off.
需要在這時候.
我要離開教會.
我要回家報到.
要讓媽媽見一下.
等等.
我也會畫這些邊界.
我也覺得是很重要的.
因為你的健康的關係.
不單只是在教會裡.
只是事奉.
是和家人.
和朋友等等.
也會讓你可以走得更遠.
所以在這方面.
我也覺得我有一個很好的平衡.
好.
最後想問一下.
你是我們的弟兄姊妹.
可以怎樣為你禱告呢?.
嗯.
繼續為在香港的教會.
我是有一個小小的想法的.
但是這個的確是.
可能是一個夢想.
不知道有沒有機會實現.
我想教希伯來文.
因為我在神學院裡.
也有機會教希伯來文.
但我想試一下教平信徒.
希伯來文.

$^{441}$特別在香港教會.
我們也在預計.
未來到底有多少空間.
去做多少事情.
會不會可以將神學教育.
普及到一個地步.
是平信徒也有某一個程度的裝備.
以致可以預備他們.
去面對未來的挑戰.
我在這方面有興趣.
我就想會不會有.
可以幫助弟兄姊妹去掌握原文.
原文的教育不只是神學院.
不只是傳道人.
而是平信徒.
或者是某些聖經的工具.
讓他們自己去做一些更加深入的研經.
以致他們以後可能也可以自己.
甚至自己餵養到身邊的人等等.
所以我也是有這個.
當然課程的課程規則.
就要重新設計等等.
所以也有想過這方面的事情.
我本來沒有想問.
但你這樣說我想問一下你.
為什麼這個原文和你剛才說的.
幫助平信徒.
怎樣面對有什麼關係.
面對將來或者難處.
有什麼關係呢?.
當然有人這樣說.
我們基督徒就是.
People of the book.
或者People by the book.
就是聖經.
所以我們信仰很重要的部分.
或者基礎就是聖經.
所以到底我們有多掌握聖經的真理.
我們有多容易或者多深入.
去明白聖經對我們去面對.

$^{481}$也不只是香港.
在全世界面對不同的日常挑戰.
或者是時代的挑戰等等.
都是很重要的.
所以我就想會不會在這方面.
讓他們有裝備.
有刀叉去嘗試去解構.
去更加明白聖經.
更加掌握得好聖經.
就會對他們站得穩一點.
分辨能力也強了.
我覺得現在這個世代.
實在太多不同的資訊.
到底怎樣分辨.
這個也是我這幾年的學習.
就是看到很多事情發生.
或者每一天都有很多新聞.
到底應該怎樣去看這些新聞.
我的立場我的看法應該怎樣.
就嘗試用聖經的訓練去看去分辨.
嘗試用一個律師的角度去想.
所以你如果不掌握好聖經.
就很容易會被這些東西擺動.
所以這是其中一件事.
我覺得在這個時代是很需要的.
我們繼續在討稿裡紀念你.
也紀念你在香港的復仕.
多謝 多謝大家.
多謝Julie 遲些再見.
再見.
\newpage



\section{17:8-16}
\label{sec:hADCTqthLUc}
\textbf{Spent And Empty (1 King 17\_8-16) - Tommy Lee}
\newline
\newline
連結: \href{https://youtube.com/watch?v=hADCTqthLUc}{\texttt{https://youtube.com/watch?v=hADCTqthLUc}} ~~~~ 語音日期: 2023-04-03
\newline
\newline
\hyperref[sec:k6VIyucQ7dI]{\small{< < < PREV SERMON < < <}}
~
\hyperref[sec:index_chronic]{\small{[返順時目]}}
~
\hyperref[sec:index_scriptual]{\small{[返順卷目]}}
~
\hyperref[sec:QvSDeFIblPQ]{\small{> > > NEXT SERMON > > >}}
\newline
\newline
$^{1}$back. All right, now let's welcome Tommy Lee..
Good morning everyone. All right, Kelly emailed me this past week and says, "Hey Tommy, we.
have a special announcement we're going to add in addition to the staff, but secretly.
I think she says you have 25 minutes." I think in her own way she says, "Make it short, don't.
bore everyone, and you have 20 something minutes because we have some translation." So we're.
going to dive in, go in and look at the passage and see what it is and get out of there so.
we can have our special announcement as well too. But I got the message from Kelly, don't.
bore everyone. Thank you. Jesus walks in and he enters Jerusalem and everyone's cheering.
him, but boy do they turn on him pretty quick. Boy do they turn on him pretty quick. They.
had an expectation of who Jesus was, what he was able to do, and the problem is he didn't.
meet their expectations. And as a result, the people ended up doing what they did for.
him and he ended up crucifying. And in the long run, it became something that opened.
the doors for what we are experiencing today. We look at the story of Elijah today. Elijah.
is the one who John the Baptist was the new Elijah who continued to preach Jesus into.
our lives, but Elijah's life also helps us to understand what that grace means in.
our life. A couple months ago I was here and we looked at 1 Kings chapter 17 verses 1 through.
7. And during that period of time, if you don't, so I'm just going to quickly go.
through it. We looked at the life of Elijah and in the original language, the name Elijah.
simply means "My God is Jehovah." That's what Elijah means, "My God is Jehovah.".
He came from, he's the Tishbit of Tishbi in Gilead. That's what 1 Kings tells him..
He comes from this obscure little city in Tishbi and scholars didn't really know.
where it was, so it tells us that it's somewhere in Gilead. It's this rural city. And so.
Elijah of Tishbi in Gilead comes from obscurity. And this guy, this rural guy who comes from.
obscurity ends up talking with the king, King Ahab of the Northern Kingdom. If you were.
here I explained to you during that period of time, there's a Northern Kingdom and.
there's a Southern Kingdom. The story of Elijah takes place in the Northern Kingdom..
And he says to this king, King Ahab, there's 19 kings all throughout the Northern Kingdom.
in its history before it's taken over. Ahab was the worst of all the kings. Jezebel later.
on, his wife, the queen, is worse than Ahab. So this couple is the worst that any of the.
regime in the Northern Kingdom ever saw. And Elijah says to King Ahab, "Neither dew nor.
rain except by my word." Why? Because Ahab was worshiping the god Baal, the god of rain,.
of storm, of weather, fertility, all of that stuff. Because King Ahab believed that if.
he put all his trust on this king or this god Baal, he was going to produce fertility..
Fertility produces food, produces money, produces flour, all of that stuff, and it was going.
to help him accomplish his goals. And he turned his eyes on the Lord. And so Elijah comes.
and says, "Neither dew nor rain except by my word." And he tells us, and that becomes.
the theme of what we talked about, "Before whom I stand." That was his motivation. That.
was the theme, "Before whom I stand." Later on, God ushers him at the end of the verse,.
that last part of chapter 17 in verse 7, he ushers him away to Cherith. And there he doesn't.
stay, but he rather lives in Cherith for a period of time. God takes him away from seclusion..

$^{41}$There God protects him. There God continues to work in him. And during that period of.
time, he's fed by ravens, and he is fed by the brook in Cherith until it dries, which.
leads us to today's passage in verses 8 through 16. "Then the word of the Lord came.
to him, 'Arise, leave Cherith, and go to Zarephath, which belongs to Sidon, and dwell.
there. Behold, I have commanded a widow there to feed you.' So he arose and went to Zarephath..
And when he came to the gate of the city, behold, a widow was there gathering sticks..
And he called to her and says, 'Bring me a little water in a vessel that I may drink.'.
And as she was going to bring, he called to her and said, 'Bring me a morsel of bread.
in your hand.' And she says, 'As the Lord your God lives, I have nothing baked, only.
a handful of flour in a jar and a little oil in a jug. And now I am gathering a couple.
of sticks that I may go in and prepare for myself and my son that we may eat it and die.'.
And Elijah said to her, 'Do not fear. Go and do as you have said. But first make me.
a little cake of it and bring it to me. And afterward make something special for yourself.
and your son.' For thus says the Lord, the God of Israel, 'The jar of flour shall not.
be spent, and the jug of oil shall not be empty until the day that the Lord sends rain.
upon the earth.' And she went and did as Elijah said. And she and her husband ate for.
many days. The jar of flour was not spent, neither did the jug of oil become empty according.
to the word of the Lord that he spoke by Elijah." If you're following after today, we are.
reading from the ESV version, 1 Kings chapter 17, verses 8 through 16. If I literally had.
to title this sermon, I would literally rename it 3 by 2. And it literally means this. There.
are three things that Elijah needs to know. Three things Elijah needs to know. Two things.
that he has to overcome. Three things that Elijah needs to know. Two things that he has.
to overcome. Let's talk about those three things. Number one, Elijah needs to know that.
the Lord sees him and knows where he's at. Number two, the Lord knows where he's going..
Number three, the Lord knows what he needs. Three things. Number one, the Lord knows and.
sees him and knows where he's at. Number two, the Lord knows where he's going. Number three,.
the Lord knows what he needs. Verse 8, "Then the word of the Lord came to him." The Lord.
sees you. The Lord sees me. The Lord sees you. There is something so powerful about.
knowing that we are seen by God. There is something powerful knowing that the great.
omniscient God who controls all of the universe sees you, has a plan for you, and has a unique.
plan in your life. It gives us dignity. It gives us worth. It makes us feel special..
We've all learned that ever since we were kids. We've all wanted to seek approval, and.
we've all wanted to be seen by our parents. I have a little girl, her name is Kara. Kara.
is four years old. And in this day, right now, she goes to Puy Taks school, and there's.
an Easter program. And Kara is in the school assembly, and from that point on, she sits.
there, and then later on when they call her, she gets up to the stage. And immediately,.
you guys probably know if you guys have kids what that means. She's looking around. And.
you know what she's looking around. She's trying to figure out where her parents are..
Why? Because she wants to know that we're watching her. She wants to know that we're.
looking and seeing at her. After the program is done, she runs, and Kara's really quiet.

$^{81}$for the most part, pretty shy. But for this time around, she runs and she looks for us,.
and she runs to us. And immediately, her first words were, "Did you see me? Did you see me?.
Did you see me?" I said to her, "Of course I saw you." She looks at me and says, "No,.
no, did mommy see me?" Because she doesn't really care what I say. She always cares about.
what her mom says to her. "Did my mom see her?" We want to be seen. We want to feel.
dignity. We want to feel worth. Story in Acts chapter 3, Peter and John. Rarely do you ever.
get Peter and John together. It's a great team up that they have there. Maybe this is.
the second story in Acts chapter 3. A man laying from birth was being carried, whom.
they lay daily at the gate of the temple. That is called the beautiful gate to ask alms.
of those entering the temple. Seeing Peter and John about to go into the temple, this.
man laying from birth asked to receive alms. In verse 4, here's a powerful verse, chapter.
3, verse 4. Peter directed his gaze at him, as did John. It says, "Look at us." He directs.
his gaze at him. You guys, all of us know what that means. When we drive off and we.
enter off and exit off Highway 55 into Chinatown, when we go and exit off Cicero to Midway Airport,.
when we get off a highway, what do we always see? We see someone begging for money. We.
see people sitting there banging their buckets asking for money. And immediately, you know.
what all of us do? We sit there and we're stopped at a red light and we pretend like.
we're looking at our phones or we're looking straight ahead because we know what that means..
If we make eye contact, they're asking us for money. And as Asians, we don't want to.
say no, so we don't make eye contact. But here, Peter and John, there's a man begging.
for money. Peter and John sits there and directs his gaze at him and he looks at them straight.
in the eye and says, "Look at us." He gives a man who is lame, who is begging for money,.
dignity. He lets that man know, "I see you." The Lord sees us. The Lord knows where we're.
at. Verse 9, "Arise, go to Zarephath, which belongs to Sidon, and dwell there." Second.
point is that Lord knows where we need to go. He has a plan. He opens and closes doors.
as he sees fit. Joshua chapter 3 is one of the most amazing directions. When Joshua takes.
over the Israelites in chapter 3, they are ready. They've lived in the wilderness for.
40 years. And on one day, God says to Joshua, "Tell the people we're going to the promised.
land. Pack up." And after 40 years on that day, they pack up and they're ready to go..
Can you imagine what they're feeling on that day knowing that after 40 years, they're packing.
up? They don't know what they're going to deal with. They don't know what dangers are.
in store. They look at the Jordan River. What in the world do we do there? But one of the.
most powerful verses in John chapter 3 is the Lord God says to Joshua, "Tell the people.
to follow, to follow the ark of the covenant." Why? Because you've never been there before..
You've never been there before. Because the Lord knows where you need to go. All of us.
who are sitting here, I think a lot of us are older, there's not too many college students.
here. All of us have struggled with our career. You may have graduated and you're trying to.
figure out what's next in your life. You may be in a job right now and you just kind of.
just don't like your job. You may be watching right now and says, "Man, I'm trying to figure.
out what's next." But there's no clarity. And you're asking, "Lord, what is it that.

$^{121}$you want me to do?" And here's an illustration, right? I had a friend, Ian. I went to Washington.
University in St. Louis. Ian was a pre-med major. And Ian, all throughout freshman year.
and all the way to senior year, he was studying to become a doctor. That's what his parents.
wanted. And that senior year, he was ready to take MCATs and something just didn't sit.
well with him. He realized, "I don't want to be a doctor. My parents want me to be a.
doctor. I don't want to be a doctor." You know what he wanted to do? He wanted to do.
film. He really loved doing film. And so a lot of times he started doing film on his.
hobby and everything like that, entered the film festival at Washington University in.
St. Louis, won first place in the film festival. And he made the decision, because he felt.
like the Lord was tugging on his heart, to leave medicine, disappoint and frustrate his.
parents, and go into film school at Florida State. He graduates from Florida State with.
film. And from that point, he takes all these little odd-end jobs. Doesn't make any money..
And so he's trying to figure out, "Lord, you brought me out of here to figure and do.
film and now I can't find a job. Is this really what you have in store for me?" But he kept.
pressing on and pressing on. And pretty soon, he got hired to do this startup program that.
just got released, America's Next Top Model. And so he started filming that. And later.
on, he got another job offer in Los Angeles. They had two shows that they wanted him to.
do. They were just starting up shows. It was CSI or Criminal Minds. And he'd go pick.
your pick. And that grew into what he's doing today. And he kept telling me, "Tommy,.
I didn't know what to do. I just kept praying. Help me to spot what you are asking me to.
do, Lord. Help me to continue to see what you are showing me, not my will, but what.
yours." The Lord opens and closes doors as he sees fit. My prayer and your prayer in.
everything that we do is always saying, "Lord, may we spot those open doors. May we spot.
those closed doors. May we be aware of it according to your own time. But may we continue.
to take myself, my goals, what I want out of it and continue to seek you. And may you.
make that aware of it every single point in time." The Lord sees you. The Lord knows.
where you are. The Lord knows where you need to go. But thirdly, the Lord knows what you.
need. "Arise, go to Zarephath, and behold, I have commanded a widow there to free you.".
He knows what you need to do. Now, it may not be what you want. Really? A widow? Seriously?.
Cherith, you're going to go to Cherith and you're going to go to Zarephath. You know.
how that's a hundred miles. You want me to walk a hundred miles out in the open? I hid.
in Cherith this entire time. I was hiding in Cherith away from Ahab because you remember.
this, I threatened the king. He's not sitting around waiting for me. He's looking for me..
So now you want me to walk a hundred miles to Zarephath? Yes. And there, who's going.
to take care of me? A widow. Really? A widow. The Lord knows what you need. It may not be.
what you want. It may not be what you think you need. But he knows what you need. The.
Lord knows where you are. The Lord sees you, number one. The Lord knows where you're going..
It may not come the way you want it to go. It may take some time to figure out what those.
open and closed doors are. And he knows what you need. But at the same time, you're going.
to have to overcome two things. You're going to have to overcome first impressions and.

$^{161}$expectations. And secondly, you're going to have to overcome what's physically impossible..
Let me talk about that. So he arose and went to Zarephath. And when he came to the gate.
of the city, behold, a widow was there gathering sticks. And he called to her and says, "Bring.
me a little water in a vessel that I may drink." And as she was going to bring it, he called.
to her and says, "Bring me a morsel of bread in your hand." And she says, "As the Lord.
your God lives." Not even as the Lord my God. As the Lord your God. Really? She's not even.
of the faith? "As the Lord your God, I have nothing big." Really? "Only a handful of flour.
in a jar and a little oil in a jug. And now I am gathering a couple of sticks that I may.
go in and prepare it for myself and my son that we may eat it and die." Really, Lord?.
You're going to ask me to walk a hundred miles? A widow's going to take care of me? And that.
widow is not part of the faith? That widow has only a handful of flour in a jar, little.
oil in a jug, and she is going to sit there, make the remaining of the food that I'm supposed.
to be eating with her servants and her son so she could go in and die. Yes! Yes! That's.
what I want you to do. So along the path of following God, you're going to have to overcome.
expectations and first impressions. You're going to have to trust that whatever you see.
before you is going to be more than you ever imagined. And sometimes when we have first.
impressions, we struggle with it. All of us love food. Hey, look, how many of our friends.
sit there on Instagram and have food collections where you take care of food? And a lot of.
times when we travel to all these different restaurants, right, we sit there and we think.
if the menu is kind of expensive, it's very nice, the food must be very good. And if we.
think it's a dumpy restaurant, the food must not be very good. But as many of you guys.
know, especially in Asia, we have all these hawker stalls where they have some of the.
most amazing food in these hawker stalls. For some of you who live—I'm 46 years old..
I think Mark is a little bit older than I am, but I grew up here in Chinatown. And back.
in the days in Chinatown, there was this—Mark, I'm not sure if you remember, there was this.
old rundown restaurant called Golden Gate Restaurant. Golden Gate Restaurant was this.
hole in the wall where these old Chinese people, including my dad, would go every day. And.
for a dollar, you bought a cup of coffee and some toast, and you bought dim sum for like.
a dollar fifty. I still remember going with my dad all the time on Saturday, Sunday. He.
would sit there with his friends, there's cigarettes smoking everywhere, and my brother.
and sister and I would sit there in our own little table and eat dim sum for a dollar.
something. It was the oldest, smelliest, dirtiest place you ever went to, right under Poy Tak.
Center, it's no longer there, but it had the best dim sum you ever had. But strictly.
by going in, you would never know. Because sometimes first impressions, sister says,.
"It doesn't allow God to work. Don't judge a book by its cover." It's a whole idea..
And when we go into situations, we have to sit there and overcome expectations, overcome.
first impressions to see how God works. And then you're going to have to overcome physical.
impossibilities. We see it all throughout scriptures, don't we? Feeding up to 5,000,.
where Jesus takes some bread and some fish and was able to feed the 5,000. We see it.
in the walls of Jericho, where the Israelites go and the walls of Jericho come crumbling.

$^{201}$down. We see it with the story of David and Goliath, where tiny David, where all of the.
Israelite army are running away. But young David is able to defeat Goliath. You're going.
to have to overcome physical impossibilities to continue to follow after him. But yet,.
at the same time, he says, verse 13, "And Elijah says to her, 'Do not fear.'" And here's.
the key phrase as we begin to wrap up. "Go and do." Do not fear. Go and do. Why was Elijah.
able to do that? Because there's two things I think Elijah was able to figure out. Elijah,.
the Tishbit from Tishbi in Gilead, whose name literally means "My God is Jehovah," knew.
God's word. He knew God's word. He knew God's promises. But he also remembered what God.
did for him in the past. He rescued him out of Ahab and into Cherith. He watches the Lord.
provided for him through ravens and through a book and now sitting there in this situation,.
he knew that God was going to watch him. Probably not in the way he expected to do, so he's.
able to go and do. You and I, in the midst of following after God, in the midst of seeking.
God's will in our life, we're going to have to overcome physical impossibilities. We're.
going to have to overcome first impressions because God sees us. Now, along the journey.
of life, he's going to call us to do something, which sometimes we may be scared to do or.
we're going to face situations where we don't know what to do. And he's going to call us.
to trust in him because he will provide us according to his need. He calls us to go and.
do. But we're able to trust him if and when we continue to know God's word and we continue.
to remember what God did in the past and cling upon those memories. I leave this with one.
final story. I did youth ministry for about 18 years. I had this one student who grew.
up, I knew him all the way from junior high to high school to college, and when he graduated.
from the University of Michigan, he got married at a young age, 24 years old. He and his wife.
moved back to Seattle, and after being married at age 26, 27 I think, one day she had a headache..
He was at work, he was an optometrist at Costco, and she had a headache. And the headache got.
so bad she went to the hospital by herself and she collapsed. And she had went into a.
coma and she never woke up. She had a brain aneurysm. And he never got the chance to say.
goodbye. He rushed over to the hospital and by then she was pretty much gone. And I asked.
him, I said, "What kept you going during that time? You never had a chance to say goodbye,.
you never even, you walked out of there that day, you never saw her again." She was also.
an only child, so her parents had to deal with it. And he says, "I had to choose to.
trust God's word in my life, and I trusted and remembered the memories of how he worked.
in our life, in our marriage, even though it was a short time. I had to trust that..
I had to choose to trust that. God knows where you are. God knows where you're going. God.
knows what you need. But you're going to have to overcome first impressions. You're going.
to have to overcome expectations. And you're going to have to overcome physical impossibilities..
The God that we love and choose, whose people in Israel turn on him a couple days later,.
who cries and yells to God in the Garden of Gethsemane, "Take this cup away from me,".
which God chooses not to do, who dies on the cross, continues to be our inspiration of.
why we do what we do. And that's why as we enter Good Friday, as we enter Easter, we.
do have that confidence to overcome. We do have that confidence to overcome physical.

$^{241}$impossibility, expectations, because Jesus showed us on the cross of what that meant..
He showed us what it meant to deal with the pain and the punishment so that you and I.
can be sitting here this morning. Let's pray. Father, thank you for the story of Elijah..
The story of Elijah teaches us three wonderful things and three things that we need to overcome..
Thank you for the stories that you give us through Scripture that guides us on where.
we must go, what we must do. Help us as we continue to go forth in our life to seek you.
every step of the way. Christ, we pray..
[BLANK AUDIO].
\newpage



\section{申命記 8:1-20}
\label{sec:QvSDeFIblPQ}
\textbf{當紀念主 (申命記 8\_1-20) - 陳智祥傳道}
\newline
\newline
連結: \href{https://youtube.com/watch?v=QvSDeFIblPQ}{\texttt{https://youtube.com/watch?v=QvSDeFIblPQ}} ~~~~ 語音日期: 2023-04-03
\newline
\newline
\hyperref[sec:hADCTqthLUc]{\small{< < < PREV SERMON < < <}}
~
\hyperref[sec:index_chronic]{\small{[返順時目]}}
~
\hyperref[sec:index_scriptual]{\small{[返順卷目]}}
~
\hyperref[sec:aVFFH_WArz0]{\small{> > > NEXT SERMON > > >}}
\newline
\newline
申命記 8:1-20
\newline
\begin{longtable}{cl}
\hline
\hline
章節 & 經文 (和合本修訂版)\\
\hline
8:1 & \begin{tabularx}{0.7\textwidth}{X} 「我今日所吩咐你的一切誡命，你們要謹守遵行，好使你們存活，人數增多，可以進去得耶和華向你們列祖起誓應許的那地。 \end{tabularx} \\ \\ \relax
8:2 & \begin{tabularx}{0.7\textwidth}{X} 你要記得，這四十年耶和華---你的神在曠野一路引導你，是要磨煉你，考驗你，為要知道你的心如何，是否願意遵守他的誡命。 \end{tabularx} \\ \\ \relax
8:3 & \begin{tabularx}{0.7\textwidth}{X} 他磨煉你，任你飢餓，將你和你列祖所不認識的嗎哪賜給你吃，使你知道，人活著，不是單靠食物，乃是靠耶和華口裡所出的一切話。 \end{tabularx} \\ \\ \relax
8:4 & \begin{tabularx}{0.7\textwidth}{X} 這四十年，你身上的衣服沒有穿破，你的腳也沒有腫。 \end{tabularx} \\ \\ \relax
8:5 & \begin{tabularx}{0.7\textwidth}{X} 你心裡要知道，耶和華---你的神管教你，像人管教兒女一樣。 \end{tabularx} \\ \\ \relax
8:6 & \begin{tabularx}{0.7\textwidth}{X} 你要謹守耶和華---你神的誡命，遵行他的道，敬畏他。 \end{tabularx} \\ \\ \relax
8:7 & \begin{tabularx}{0.7\textwidth}{X} 「耶和華---你的神必領你進入美地，那地有河流，有泉源和深淵的水從谷中和山上流出。 \end{tabularx} \\ \\ \relax
8:8 & \begin{tabularx}{0.7\textwidth}{X} 那地有小麥、大麥、葡萄樹、無花果樹、石榴樹，那地也有橄欖油和蜂蜜。 \end{tabularx} \\ \\ \relax
8:9 & \begin{tabularx}{0.7\textwidth}{X} 那地沒有缺乏，你在那裡有食物吃，一無所缺；那地的石頭是鐵，山中可以挖銅。 \end{tabularx} \\ \\ \relax
8:10 & \begin{tabularx}{0.7\textwidth}{X} 你吃得飽足，要稱頌耶和華---你的神，因為他將那美地賜給你。」 \end{tabularx} \\ \\ \relax
8:11 & \begin{tabularx}{0.7\textwidth}{X} 「你要謹慎，免得忘記耶和華---你的神，不守他的誡命、典章、律例，就是我今日吩咐你的。 \end{tabularx} \\ \\ \relax
8:12 & \begin{tabularx}{0.7\textwidth}{X} 免得你吃得飽足，建造上好的房屋，住在其中， \end{tabularx} \\ \\ \relax
8:13 & \begin{tabularx}{0.7\textwidth}{X} 你的牛羊增多，你的金銀增多，你擁有的一切全都增多， \end{tabularx} \\ \\ \relax
8:14 & \begin{tabularx}{0.7\textwidth}{X} 於是你的心高傲，忘記耶和華---你的神。他曾將你從埃及地為奴之家領出來， \end{tabularx} \\ \\ \relax
8:15 & \begin{tabularx}{0.7\textwidth}{X} 曾引領你經過那大而可怕的曠野，有火蛇、蠍子、乾旱無水之地。他也曾為你使水從堅硬的磐石中流出來， \end{tabularx} \\ \\ \relax
8:16 & \begin{tabularx}{0.7\textwidth}{X} 又在曠野將你列祖所不認識的嗎哪賜給你吃，為要磨煉你，考驗你，終久使你享福。 \end{tabularx} \\ \\ \relax
8:17 & \begin{tabularx}{0.7\textwidth}{X} 你心裡說：『這財富是我的力量、我手的能力得來的。』 \end{tabularx} \\ \\ \relax
8:18 & \begin{tabularx}{0.7\textwidth}{X} 你要記得耶和華---你的神，因為得財富的能力是他給你的，為要堅守他向你列祖起誓所立的約，像今日一樣。 \end{tabularx} \\ \\ \relax
8:19 & \begin{tabularx}{0.7\textwidth}{X} 你若忘記耶和華---你的神，隨從別神，事奉它們，敬拜它們，我今日警告你們，你們必定滅亡。 \end{tabularx} \\ \\ \relax
8:20 & \begin{tabularx}{0.7\textwidth}{X} 耶和華在你們面前怎樣使列國滅亡，你們也必照樣滅亡，因為你們不聽從耶和華---你們神的話。」 \end{tabularx} \\ \\
[1ex]
\hline
\hline
\end{longtable}
$^{1}$(麥克風轉動).
大家早安.
是,很高興今天有機會回到UV.
跟大家一起來聚會.
今天我們看的經文是.
看《申命記》第八章1到20節.
我會一邊說一邊讀這段經文.
今天的講題是要紀念.
當紀念你的主.
說到紀念.
我經常會想紀念的反而詞是什麼呢?.
你記得的東西.
反面來說就是你忘記了一些東西.
這是我們生命中經常發生的事.
做學生,我們知道.
上學背生字或數學公式.
我們都背得很辛苦.
但又經常都忘記了.
作為家長,可能要記小朋友今天參加什麼活動.
他什麼時候要考試,什麼時候要上學.
都很難記.
或者上班或事奉的時候.
我們經常都會問.
今天我要開什麼會呢?.
有什麼我今天要在限期之內做完的事呢?.
我們用盡全力來記住我們要記的東西.
但我們經歷的就是一次又一次的.
會忘記,會遺忘了.
而在《申命記》第八章這段經文.
相當同樣提醒我們說.
我們要去紀念它.
我們要記得它的作為.
我們要記得要遵守它的律法.
上帝知道我們的記性有限.
所以它提醒我們.
它們知道我們經常都要一個提醒.
就是提醒我們說我們是屬於它的子民.
我們需要進行它的旨意.
在這段經文當中我們分為兩個部分.
我們首先來看第一部分.

$^{41}$是《申命記》第八章的一到十節.
如果有聖經的話我們可以一起來翻開來看.
《申命記》第八章第一節.
我今日所吩咐的一切誡命.
你們都要謹守遵行.
好叫你們存活,人數增多.
且進去得耶和華向你們列祖.
起誓應許的那地.
你也要紀念耶和華你的神是曠野.
引導你的這四十年.
是要撫練你,示現你.
要知道你心內如何.
肯守他的誡命不幸.
他撫練你,任你饑餓.
將你和你的列祖所不認識的嗎哪賜給你吃.
使你知道人活著不是單靠食物.
乃是靠耶和華口中所出的一切話.
這四十年你的衣服沒有穿破.
你的腳也沒有腫.
你當心裡思想耶和華你神管教你.
好像人管教兒子一樣.
所以你要謹守耶和華你神的誡命.
遵行祂的道,敬畏祂.
因為耶和華你神領你入美地.
那裡有地有河,有泉有源.
從山谷中流出水來.
那地有小麥,大麥,葡萄樹,無花果樹,石榴樹,橄欖樹和麥.
你在那地不缺食物,一無所缺.
那地的石頭是鐵,山內可以滑動.
你吃得飽足,就要稱頌耶和華你的神.
因他將那美地賜給你了.
在這段經文第八章的第一節.
摩西向以色列人發出一個很清楚的信息.
我今日所吩咐一切的誡命.
你們都要謹守遵行.
以至於你們可以存活.
以至於你們可以進入耶和華向你們列祖.
應許的那個地方.
以至於你們可以在應許之地來享受上帝的祝福.
上帝的誡命很清楚.

$^{81}$他應當要以馴服回應上帝.
但是在第二至第五節當中.
我們也看到很有趣的一件事.
耶和華說你們要記得我.
但他說不是要靠你們自己的能力來記得我.
而是你們要靠著他自己的作為來謹記上帝.
在第二至第五節當中.
上帝幫以色列人反思的是他們四十年來.
在曠野當中的一個經歷.
如果我們想曠野的話.
其實曠野的經歷對以色列人來說是一個屬靈.
或者是他們出埃及旅程的一個低潮.
他們進入曠野是因為上一代以色列人的不信.
以至於他們錯失了可以進入應許之地.
進入迦南地的這個機會.
他們因為看到迦南人的強大.
他們看到那地的居民強大很高.
他們就覺得很害怕.
而忘記了那位帶領他們出埃及.
帶領他們過紅海.
帶領他們戰勝埃及一起軍兵的上帝.
他們因為眼見的一個危險而忘記了上帝.
所以在曠野當中.
這一批以色列人等了四十年.
這一批聽眾.
《申命記第八章》的這一批聽眾.
其實是一批他們可能有記憶以來.
都是活在曠野當中的一批以色列人.
這一批以色列人和上一代人不同.
他們沒有經歷過上帝帶領他們過紅海.
沒有經歷過他們在埃及遺奴之地的經歷.
也沒有親眼見證過上帝和他們立約的那一個時刻.
但是這一批以色列人.
就是有機會見證上帝應許實現的一個地方.
在第二節到第五節當中.
我們見到上帝說你要紀念耶和華.
你的神是曠野引導你的這四十年.
上帝是再一次和以色列人重申.
你們就算在曠野當中受苦.
你們以為上帝拋棄你們.

$^{121}$但是上帝其實沒有.
上帝這一批以色列人.
他們當然很想就這樣可以走進應許之地.
就這樣可以拿到上帝的一個祝福.
但是上帝說他要按著他自己的心意和計劃來引導他們.
以色列人在曠野當中的經歷.
其實不是為了上帝的好處.
也不是為了單單來懲罰以色列人他們的不順服.
上帝說他要把以色列人放在曠野當中.
是要讓以色列人自己心裡可以清楚地知道.
心裡清楚地明白他們的心在哪裡.
他們究竟願不願意來遵守上帝的誡命.
在曠野當中,當他們沒有食物的時候.
上帝給了他們馬納.
這個是你和你列祖都不認識的事情.
這個是對這一代以色列人特殊的啟示.
當他們沒有水的時候,上帝供給他們.
即使當他們認為受苦的時候.
上帝都在保護他們.
他不讓他們的衣服磨損,也不讓他們的腳變腫.
上帝就像一個父親一樣.
來關心以色列人,管教他們,教導他們.
教他們知道生命當中最重要的不是食物.
也不是物質,也不是他們眼見的一些巨人或者威脅.
上帝要教導他們的,就是在第四節當中說.
第三節當中說,他要知道人活著不是單靠食物.
是靠耶和華口中一切的說話.
如果以色列人他們要求的只是一些物質上的需求.
神說我可以很輕易地來滿足你們.
你要食物我可以給到你.
你要水我可以給到你.
你要衣服我都可以給到你.
但是我唯一給不到你們的,就是你們的一個馴服.
這個是你們自己需要來學習的.
神要告訴他們,這一位為他們供給一切的上帝.
這個沒有離開過他們的上帝.
才是他們生命當中最重要的東西.
而在第六節到第十節當中.
上帝不單止來回憶過去為以色列人的作為.
他更加來向以色列人保證.

$^{161}$他昔日對他們的應許終會實現.
在第六節到第十節我們可以見到.
可以見到,穆斯在說.
他們將要進入的地是多美好的.
那裡有河有泉有源.
在山谷當中流出水.
有小麥有大麥有葡萄樹無花果樹石榴樹橄欖樹和麥.
那個地不缺食物,一無所缺.
我們見到他們等候這個應許實現的時候.
有兩個很重要的基礎,就是第六節當中.
他們可以進入應許之地.
是要靠著遵行耶和華的道.
要敬畏耶和華.
對於以色列人來說.
上帝這個應許是對他們至關重要的.
在《創世記》十二章當中.
我們已經見到上帝向以色列人的祖先亞伯拉罕來應許.
說你們以後的子民將要來進入這個應許之地.
當穆斯帶領以色列人離開埃及的時候.
他們同樣帶著的是對上帝這一份的應許.
雖然他們上一代人不信服.
令到這個承諾,令到這個應許被推辭.
但是經歷曠野漂流的四十年之後.
上帝跟他們說,神的應許依然存在.
這一塊地依然是你們的.
這一塊迦南地,這一塊有地有河有水有大麥有小麥有葡萄樹無花果樹.
不缺食物的這一塊地,仍然是你們的.
只要你們願意跟隨,只要你們願意遵守上帝的誡命.
而不是好像你們的列祖一樣.
事實上這一代的以色列人已經見證了上帝的應許是怎樣來實現的.
在《民數記》二十一章當中.
我們見到上帝已經帶領以色列人戰勝迦南人的兩個王.
以色列的兩個支派,約旦河的西邊已經來到去駐紮.
已經領受了他們應許的土地.
其實整個以色列的群體,這一代的以色列人.
其實他們見證了上帝的應許是怎樣來實現的.
他們見證了上帝的大能.
怎樣來幫助他們戰勝一些他們過往認為他們不可以戰勝的人.
當波西來宣講這篇信息的時候.
他們處於約旦河東的摩亞平原.

$^{201}$摩亞平原他們離開迦南地,只差了一條河的地方.
我們如果這樣想的話,我們可以這樣有一個畫面.
就是當波西來宣講這篇信息的時候.
其實以色列人他們是看到迦南地的.
他們是看到他們將來要住的地方是怎樣的.
他們在曠野當中等了40年,就是等那個地方.
他們列祖,由亞伯拉罕開始等了400多年.
所以要進去的就是那個地方.
上帝的應許離開他們很近.
他們知道他們快要可以進入到去.
他們知道他們將會見證上帝應許實現的人.
因此摩西很著重地提醒他們.
當你們進入這一塊地的時候.
你們要記得耶和華.
在第十節摩西強調.
你們在這塊土地吃得飽足.
你們吃飽滿足了你一切需要的時候.
你們就要去稱讚耶和華,你的上帝.
因為祂將那塊地賜了給你們.
耶和華上帝才是你們安居樂業.
得享繁榮的原因.
你們要敬拜祂,你們要馴服祂,你們要去紀念祂.
我們看到第一到第十節其實是一個很美好的應許.
上帝在回憶祂和以色列人之間的關係.
祂怎樣在文約當中保護以色列人,供給以色列人.
祂將來也會怎樣用信實來回應他們.
但是在第十一到第二十節的時候.
我們看到摩西的語氣變了.
祂說的不再是一個應許,而是一個危機.
祂在警戒以色列人不要忘記了你的上帝.
我們一起來看第十一到第二十節.
《申命記律》八章十一節.
你要謹慎免得忘記耶和華你的神.
不受祂的誡命,典章,律例.
就是我今天所吩咐你的.
恐怕你吃得飽足,建造美好的房屋居住.
你的牛羊加多,你的金銀增添.
並你所有的存都增加.
你就心高氣傲,忘記耶和華你的神.
就是將你從埃及地遺奴之間領出來的.

$^{241}$引你經過那大而可怕的曠野.
那裡有火蛇,蠍子,乾旱,無水之地.
祂曾為你洗水,從堅硬的盆石當中流出來.
又在曠野將你列祖所不認識的嗎哪賜給你.
是要撫練你,試煉你,叫你終究享福.
恐怕你心裡說,這貨才是我的力量,我的能力得來的.
你要紀念耶和華你的神.
因為得貨才的力量是祂賜給你的.
為要堅定祂向你列祖所立的約,像今天一樣.
你若忘記耶和華你的神,隨從別神.
事奉敬拜,你們必定滅亡.
這是我今日警戒你們的.
耶和華在你們面前怎樣使列國的民滅亡.
你們也照樣滅亡.
因為你們不停從耶和華你們神的話.
在這段經文當中我們見到摩西警告以色列人.
你們不要忘記了上帝.
你們不要不守祂的律例,典章,誡明.
你們不要忘記我今日跟你說過的一切說話.
在第12到第13節當中我們見到摩西來解釋.
以色列人為何會來忘記上帝.
他們忘記上帝不是因為他們再一次見到一些危險.
或者一些威脅他們沒辦法來戰勝.
他們忘記的原因正正是因為他們領受了上帝的祝福.
他們見到他們的牛羊家多.
他們的金銀真添.
他們全部所擁有的東西都變得越來越多的時候.
他們就會心高氣傲忘記耶和華你的上帝.
當以色列人他們生活安逸.
不需要再為他們生活所需的東西來受煩.
定居在迦南地上帝賜給他們這塊美地的時候.
以色列人就會自高自大.
他們忘記上帝.
他們陷入的是一個陷阱.
他們陷阱是因為這些繁榮,這些豐盛.
都是他們努力所得來的結果.
以至於他們忘記了他們可以有這些祝福.
其實全部都是因為上帝的恩典.
或者簡單來說摩西在警告以色列人.
當上帝的應許得到實現的時候.

$^{281}$就是這班以色列人最有可能忘記上帝的時刻.
當他們被生活所擁有的豐盛來滿足的時候.
他們的眼目就會轉來上帝.
他們專注在這些安逸當中.
我們見到以色列人的歷史其實充滿了這些遺忘.
這些陷阱摩西給他們的境界是真實存在的.
在《約書亞記》當中我們見到.
以色列人渡過約旦河的時候.
上帝就要叫他們在急隔的時候建立一個石堆.
就代表上帝好像再一次帶領他們出紅海一樣.
很奇蹟般地帶領他們過了約旦河.
以至於他們可以領受上帝的祝福.
這些石頭作為以色列人的記號.
讓他們可以記得上帝的能力,上帝的應許.
怎樣來跟他們同在的.
但即使有急隔這些石頭的記號.
我們見到當這一批跟約書亞同在的長老,領袖離開的時候.
以色列人就已經很快地忘記上帝.
我們不用看很遠.
只是看《約書亞記》之後的時詩記.
我們已經見到以色列人忘記了上帝.
他們陷入了一個危險.
就是當他們領受了這個土地的時候.
他們就轉來上帝.
他們只有在陷入危難的時候才會呼求上帝.
而當他們的上帝令他們的生活再一次重回正軌的時候.
他們又再一次地忘記,再一次地轉來.
這件事令我記得.
是我中學時經常發生的一個經歷.
當我生病的時候.
我就會向上帝有一個這樣的禱告.
我就說「上帝,你醫治我吧」.
「我知道你是這個創造天地的真神來的」.
「我保證如果你治好我,我得了醫治的時候」.
「我就會記得向你禱告的」.
這些禱告其實都很百事百靈的.
通常過了兩三天之後就康復了.
但是我康復之後.
我就再也沒有祈禱了.
這件事都很百事百靈的.

$^{321}$上帝在哪裡呢?.
在醫治好我之後,我已經忘記了祂.
其實我們很多時候都有這樣的心態.
當我們遇到危險,需要幫助的時候.
我們就很想尋求上帝.
因為只有上帝是會解決我們生命當中問題的一個主.
但是對我們來說,上帝是不是只是這樣呢?.
是一個解決師,是一個幫我們解決問題的人,神.
這就是摩西警告我們的東西.
摩西來提醒我們說.
我們要紀念上帝,要記得遵守祂的誡命.
不要享受完上帝的祝福.
經歷完上帝的應許大能之後就忘記了祂.
為什麼記得對我們來說是這麼重要?.
其實記憶不是我們個人層面上的一個認知.
不是說我們記得就好像我們知道.
記憶更加是一個群體性的經歷.
我們看到摩西其實是整段經文當中.
其實跟以色列人說.
你們要記得的不是你們個人的一些和上帝的經歷.
不是你們個人當中我們怎樣見證上帝的作為,這麼簡單.
他要叫以色列人記得的是,要記得你們的歷史.
要記得你們作為一個群體.
你們怎樣一起來經歷,經驗上帝大能的時刻.
例如他們是曠野受到保護.
或者他們見到出埃及的大能.
這些是一些以色列人集體性的回憶.
這是他們共同所擁有的一些東西.
這些傳承,這些結構來告訴他們.
以色列人究竟是怎樣的.
以色列人是因為和上帝之間有這個盟約.
所以以至於他們這個群體.
可以經歷到上帝的保護,上帝的幫助.
上帝的恩典,上帝的應許,怎樣來實現.
社會學家就將這些現象,將這些共同的回憶.
給了一個很美麗,很fancy的詞語.
他就叫做社會的記憶.
這個是說一個群體.
他們是因為有一些共同的回憶.
共同的經歷,共同的經驗.

$^{361}$而來告訴他們這個群體究竟是怎樣.
而來定義這個群體有甚麼獨有的價值.
有甚麼獨特的身份.
以色列人他們經歷的是上帝的保護.
而這個共同的回憶對他們來說.
就是他們經歷,他們作為上帝子民.
他們與上帝在盟約之間的關係.
他們應該怎樣來回應.
他們作為上帝子民.
這個就是他們獨特的身份.
在第19到第20字當中.
摩西最後來警告他們說.
如果你們忘記了上帝.
你就必定會滅亡.
你如果忘記了上帝.
隨從必實,事奉敬拜.
你就會滅亡.
在第20字當中他們說.
如果他們忘記了上帝.
他們不再知道上帝是誰.
同時他們拋棄的是.
其實是自己的身份.
因為耶和華說.
我怎樣在你們面前使列國的民滅亡.
你們也照樣滅亡.
以色列人忘記上帝.
就是說其實他們不再成為上帝的子民.
他們就要好像列邦一樣.
來敬拜他們的偶像.
他們忘記了自己的身份.
忘記了自己的使命.
忘記了他們是上帝聖潔的國民.
他們有這個責任給外邦人.
來見到屬於上帝的人是如何與眾不同.
這就是為何忘記對於以色列人來說.
是這麼危險的一件事.
這段經文為我們指出.
忘記或遺忘的一個危機.
但近年科學家就指出.
其實忘記或遺忘.

$^{401}$其實對於我們是有好處的.
在時代週刊的一篇文章當中.
就說遺忘其實是我們大腦的一個應對機制.
甚至有一本書就說.
忘記是有好處的.
忘記或遺忘有些東西.
只不過是我們大腦嘗試為我們清理.
哪些資訊是不重要.
或者哪些資訊是重要的.
以至於我們可以刪除一些不重要的資訊.
以至於我們可以學一些新的東西.
例如我們不會記得.
我們昨天在Costco收銀.
或者收銀的人穿的衣服是甚麼顏色的.
這些資訊對我們的生命其實是不重要的.
如果我們只記得這些這麼無聊的東西.
我們哪有時間去背誦經文.
但這可能正正是我們忘記上帝的原因.
以色列人忘記上帝.
離棄上帝.
轉向敬拜和事奉其他偶像.
是因為他們覺得.
外邦的偶像或者和其他國家結盟.
對於他們的生存.
對於他們的生活更加息息相關.
他們覺得這些國家更加強大.
更加可以保護他們.
他們忘記了上帝.
是因為他們覺得上帝對他們的生命不重要.
上帝在他們的生命當中沒有任何的關係.
我們今天是否同樣地忘記上帝呢.
我們看到使徒保羅在他傳福音建立教會的時候.
即使面對被人擊打,被人羞辱,被人威脅.
在各樣的情況下他仍然忠於上帝.
在《肥納比書》第三章裡面提到.
保羅重新回憶.
上帝如何將他由一個迫不及待的教會的法利賽人.
變成一個忠誠的僕人.
他記得的是上帝如何令他的生命得到更新.
他說「我將萬事都當作有損.

$^{441}$因為我認識我主耶穌基督為自保」.
他記得他自己一個新的身份.
他是作為一個天上的公民.
他的眼光是要專注於上帝.
是要視基督為他生活當中最終極的喜樂和榮耀.
他的生命是要被基督的復活的大能改變.
以至於保羅可以記得.
是因為基督的復活.
他的生命才得到轉化.
他才有力量去為上帝繼續服侍他.
保羅的經歷提醒我們.
其實紀念上帝不是說我們只是背經文或背上帝的名字那麼簡單.
紀念上帝不是一個頭腦的認知.
紀念上帝更加重要的是我們的行為.
紀念上帝紀念他的作為.
是告訴我們我們是誰.
我們今天有責任來做些什麼.
正如以色列人要記得和他們立約的上帝.
要帶領他們出埃及的上帝.
我們今天記得的是和我們立約的上帝.
我們的目光要專注的是耶穌基督.
他在十誡上的犧牲.
他成就的救恩.
作為基督徒的我們.
我們這個群體其實最重要的一個特徵.
或者我們共同所有的一個經歷回憶.
就是耶穌基督為我們的犧牲.
我們是一個因著基督的死和復活而被更生的群體.
正如我們剛才報告當中說到.
其實今天是教會傳統年曆的一個中樹主日.
或者中枝主日.
其實就是御府期的最後一個主日.
是為了慶祝耶穌基督.
他就是改為上帝應許的彌賽亞拯救者.
來為我們犧牲.
耶穌基督是站在路上.
他是以一個謙卑的身份來走進耶路撒冷.
最後一步一步走到十誡當中.
作為基督徒我們很知道.
或者我們很記得耶穌基督的犧牲.

$^{481}$或者我們甚至背得出.
就是說耶穌基督是十誡傾出的寶血來潔淨我們的罪.
但是今天我們的眼目.
我們的行為.
我們的心思意念.
有沒有被基督的犧牲所改變呢?.
我們有沒有去到基督所想的.
基督所行的.
基督所愛的呢?.
今天耶穌基督的復活.
有沒有真正改變我們生命當中的一切.
有沒有令到我們可以更加願意來服侍他.
我們一齊來禱告.
在這個時刻.
讓我們可以再一次來到.
讓我們的眼目專注於基督你所承受的救恩當中.
讓我們在這段時間.
再一次來紀念你的作為.
紀念你的救恩.
求神你繼續來到去跟隨我們.
幫助我們.
讓我們的心思意念都被你所改變.
行基督所行的.
愛基督所愛的.
我們感恩感恩地禱告.
阿風靠主耶穌基督得勝的名求.
多謝Joseph提醒我們.
\newpage



\section{}
\label{sec:aVFFH_WArz0}
\textbf{比對研究 11: 巴力宗教 4 - 陳崇基牧師}
\newline
\newline
連結: \href{https://youtube.com/watch?v=aVFFH_WArz0}{\texttt{https://youtube.com/watch?v=aVFFH\_WArz0}} ~~~~ 語音日期: 2023-04-08
\newline
\newline
\hyperref[sec:QvSDeFIblPQ]{\small{< < < PREV SERMON < < <}}
~
\hyperref[sec:index_chronic]{\small{[返順時目]}}
~
\hyperref[sec:index_scriptual]{\small{[返順卷目]}}
~
\hyperref[sec:_Z2VGiey2J8]{\small{> > > NEXT SERMON > > >}}
\newline
\newline
$^{1}$我們一路在看比對研究.
比對研究其實是將正經和其他的文獻.
做一個 side by side 的比對.
看看中間有什麼相似和分別的地方.
我們之前一直在說巴力的宗教.
巴力的背景.
我們曾經提過巴力的大本營.
或者最多我們今天能夠知道巴力宗教.
或者巴力的神話.
巴力神明的地方.
他所出土的文獻就在哪裡呢?.
就是烏加利的大海港.
就是塞浦路斯的東邊.
這個地方.
眾神所匯集的山就是 Mount Saffong.
就是薩峰山.
這個就是大迦林地的最北邊.
也是在這些大的海港裡面.
一個比較北一點的地方.
在烏加利找到一些文獻.
讓我們知道巴力的神話是怎麼樣的.
我們做比對研究的時候.
有一件事情我們可以做的就是看聖經.
同時.
看看聖經背後.
他所預設的這些意識形態.
這些思想.
很多時候就是這些經外的文獻能夠呈現出來.
有了這些經外文獻來看聖經的時候.
或者 read the bible in light of those documents.
我們就可以更加清楚聖經裡面的一些字眼.
還有一些故事.
一些主題.
為什麼會這樣記載這樣寫的呢?.
我們不妨去看看.
跟巴力宗教很有關的一些經文.
這個首當其衝一定是以利亞的故事.
就是聶王記第十七章第一節開始.
大家很清楚這些故事了.
我們從小到大在教會聽到的故事.

$^{41}$都是講到以利亞是多麼英雄.
是不是?.
我們不妨用一個比對研究.
用背景的資料.
來看看有沒有一些亮光讓我們看聖經.
裡面的一些字眼.
會不會跟巴力的文獻有一個呼應.
呼應的目的.
就是去做一些所謂的辯護.
或者去抗衡.
去講清楚.
巴力不是神.
耶和華才是神.
第十七章第一節.
聶王記上.
他說「住在激烈的提斯比人以利亞對阿哈說.
我指著所事奉永生的耶和華.
以色列的神起誓.
這幾年我若不禱告.
或者我若不說話.
必不降露水也不下雨」.
這裡沒有什麼特別.
不是以利亞的禱告.
或者他才說的話.
就不會有雨水.
不會下雨.
但是為什麼會提這兩樣東西.
露水和下雨.
我們畫了比對研究.
我就選了一段巴力的.
烏加利的文獻.
提到巴力.
這個就叫做Epik Akkad.
這裡就提到一個烏加利的英雄.
英雄的王.
是一個公義的王.
我們不妨看看.
這一段文獻.
到底想講些什麼.
我們讀一讀這一段Epik Akkad.

$^{81}$下面粗略的中文翻譯.
就說「但奈爾·拉比烏人.
在浩熱的當中.
當時有旱災.
他在形容旱災.
在浩熱的當中.
你去召喚雲彩.
將它召喚過來.
扣它過來.
讓雲在當空的夏天下雨.
然後甘露把露水灑在葡萄之上.
七年巴力不在.
旱災造成饑荒.
七年巴力不在.
八年嘉運者不在.
C句A句B句互相呼應.
所以B句很多時候就短一點.
巴力不在就是嘉運者不在.
七年八年他不在.
還有一個很長的時間.
巴力不在.
嘉運者不在.
之前我們提過嘉運者.
這個綽號.
Wakaf Arapu.
這個已經在詩篇或其他地方有提過.
聖經是爭拗.
真正的嘉運者.
真正掌管雨雲的.
在耶和華不是巴力.
七年八年巴力不在.
嘉運者不在.
沒有甘露沒有傾盆大悅.
沒有深淵的漩渦.
沒有巴力的聲音.
Welcome voice.
巴力的聲音.
沒有了.
這裡有兩個字我們首先要留意.
讓雲在夏天下雨.

$^{121}$在烏加列文裡面.
Ugaritic.
就是Mata.
甘露.
在烏加列文裡面.
就是Tau.
他說.
沒有雨也沒有甘露.
兩個字.
並列在.
同一句的說話裡面.
沒有雨也沒有甘露.
這裡很清楚.
巴力就是掌管雨水和甘露.
是不是.
所以這裡很清楚.
嘉運者掌管雨雲.
也掌管Mata and Tau.
就是下雨.
雨水和甘露.
順帶一提這個Danel.
就是Danel.
是什麼人呢.
其實聖經裡面有提過.
以西結書十四章十四節.
其中一節經文.
整章裡面也有提過.
這個所謂Danel.
中文的翻譯或是英文的翻譯.
都是譯作Danel.
我們一想起Danel.
我們就想起那個被擄.
流放到巴比倫.
那個書裡面的主角.
Danel.
但這裡的原文.
馬蘇拉抄本.
或是某一些的抄本.
其實串法就是.
Daniel.

$^{161}$這個串法.
有別於.
平時的.
但爾里.
但爾里書裡面的但爾里.
那個串法.
你會看到有分別.
Daniel.
中間有個y.
在下面.
然後有個sere 有個長的e音.
Daniel.
如果以西結.
真的想說但爾里的話.
他可能就不是.
為什麼會串這個.
所以大部分的.
聖經學者.
今時今日都會認為.
這個但爾里其實就不是.
但爾里書.
那個但爾里.
在烏格麗文獻裡面.
告訴我們那個公義的王.
那個Danel.
有名於.
他去.
照顧孤兒寡婦.
行公義的一個王.
這個叫做但內爾.
如果你這樣看.
提起羅亞.
但內爾.
和約伯這三個人.
其實就吻合十四章裡面.
所說的.
這三個人是什麼呢.
他們只能因自己的義.
救自己的命.
為什麼要提這三個人.

$^{201}$因為這三個人有一個共通點.
三個都不是以色列人.
羅亞約伯.
和但內爾都不是以色列人.
下面有一個.
學術的文章.
你可以參考一下.
怎樣也好.
這裡說到但內爾.
所以我們.
回到剛才的.
烏格麗文獻.
我們說.
最主要我們想留意的.
就是雨水.
和金露.
是同一句說話.
如果翻開.
到列王記上十七章第一節.
介紹.
以利亞的時候.
以利亞和阿哈說什麼呢.
這幾年我約伯討個.
不靠我口中的說話.
必不降露水.
也不降雨.
希伯來文是什麼呢.
必不降露水.
或者金露.
就是Tau.
Tagalamut.
Tau.
T和L.
音標就是A.
A音Tau.
就是露水英文翻譯做Dew.
也不下雨.
就是Mata.
Mata.
就是M T R.

$^{241}$音標就是用了A音.
這個A音.
英文是Rain.
這個絕對不是.
一些意外的事.
因為.
以利亞是來抗衡.
巴力的.
告訴以色列的民眾.
告訴人們.
真正駕運者.
掌管雨雲的是耶和華不是巴力.
同時他們第一句說話.
跟阿克說的是什麼.
露水和雨水.
Tau和Mata.
回到剛才那個.
烏加利的文獻.
那個史詩.
Mata和Tau.
是同一句說話出來.
對於一個.
拜巴力的人或者有巴力背景的人.
他就一定聽到.
什麼呢.
就像雨一樣簡單.
聽到是一個抗衡巴力.
神話.
抗衡巴力的神學.
一個宣告.
在希伯來的.
文字裡面.
有很多不同的字.
去形容雨水.
或者水.
在這裡我再將.
加林地的龍耕的循環再列出來.
由1月至12月.
你會看到.
每一年的循環.

$^{281}$由9月開始.
9月尾開始.
駐彭徹的時候.
就是第一個雨季出來.
雨季來到了.
下了一陣雨.
然後農夫就可以耕地.
將泥土翻鬆.
然後就可以播種.
殺種的比喻就是這裡.
然後再有冬雨.
這些種子就可以發芽.
然後如果你下了8吋的雨.
就應該有大麥.
收成.
雨節到了.
然後下了12吋的雨水.
小麥就成功可以收成.
這個雨季過了之後.
夏天就會到了.
夏天就沒有雨水.
天色一片青藍.
是嗎?.
那靠什麼來維持住.
那些果樹啊.
農作物的.
能夠結果呢?.
到年尾的時候.
再雨季來的時候.
能夠收割的時候.
葡萄啊 無花果啊 石榴啊.
橄欖啊 蜜棗啊等等.
靠什麼.
雨季結束之後.
最後可以收成呢?.
靠什麼可以熬過.
這個熱的夏天呢?.
就是靠.
淘(淘).
露水.

$^{321}$所以在這裡說淘和麥他.
淘的意思是.
5月至9月的時候.
是需要的甘露.
甘露是什麼呢?.
就是清晨的時候.
空氣裡面的露水.
它凝固.
然後變成液體.
然後落在青草上.
我們晨早出去的時候.
天朦剛的時候.
就經常在一片草地裡面.
看到這些露水.
這就是淘.
這就是非常重要的.
如果沒有了它們.
就不會有好的收成.
甚至是寒季.
再一年發生的時候.
連這些淘和甘露.
都會沒有了.
所以在五經裡面說到沒有甘露的時候.
其實就是一種.
聖若裡面的.
咒詛.
有甘露就是若裡面的.
祝福.
麥他就自然是在說雨季.
冬天.
到春天.
所下的雨.
提到這兩個字.
當然就是反映出.
迦林地本身的.
農耕的循環.
有麥他.
也有淘.
但最重要在比對研究裡面.
可以讓我們看到的.

$^{361}$其實就是在說抗衡巴力的神學.
真正控制麥他.
和淘.
整年的水分.
整年的雨水.
包括了甘露.
在誰的手中呢?.
就是在耶和華永生之主的手中.
若果不是靠著.
先知的說話.
先知的禱告.
這些水分.
是不會下降的.
以致.
就會有寒季出現.
但是人又要去.
歸向耶和華.
然後去認定.
耶和華是那位教養者.
去控制.
控制著.
我們生命的.
露水和雨水.
那位的神.
去敬拜他 去事奉他.
這才是真正的信仰.
這就是.
十七章第一節.
比對研究可以給我們一些亮光.
以後我們繼續看下去.
\newpage



\section{}
\label{sec:_Z2VGiey2J8}
\textbf{Resurrection Gospel: John 20 - Rev. Mark Chan}
\newline
\newline
連結: \href{https://youtube.com/watch?v=_Z2VGiey2J8}{\texttt{https://youtube.com/watch?v=\_Z2VGiey2J8}} ~~~~ 語音日期: 2023-04-18
\newline
\newline
\hyperref[sec:aVFFH_WArz0]{\small{< < < PREV SERMON < < <}}
~
\hyperref[sec:index_chronic]{\small{[返順時目]}}
~
\hyperref[sec:index_scriptual]{\small{[返順卷目]}}
~
\hyperref[sec:PhsM03F2OIM]{\small{> > > NEXT SERMON > > >}}
\newline
\newline
$^{1}$All right, I'll hand it over to Mark..
Well, thank you, everyone, for coming..
And as Kelly told us last Sunday,.
it was our Easter Sunday and also our seventh anniversary. Right..
So we can say everybody, everybody,.
happy birthday to each other because we are launching.
the seventh year as a church, as a church plan, as a church family..
So it's a happy day. And so last week we celebrated Easter Sunday..
And I want to tell you.
sometimes or most of the times.
we do treat Easter as an afterthought,.
a kind of a footnote,.
because we often feel like the center of Christianity,.
the center of being a Christian is the cross..
Whatever Jesus did for us on the cross, he died for us sacrificially..
He died for us redemptively and to rescue all of us as sinners..
So sometimes we focus on the cross and there's nothing wrong with that..
There's everything right about it, because biblical faith does center.
on the on the cross, on the accomplished work of Jesus Christ..
And that's why we celebrate and worship together on Good Friday..
Right. Were you here?.
And that's why we also celebrate Easter..
However, sometimes I get the feeling that, you know, after Easter,.
after we talk about resurrection and then we get it done.
and then we just move on. Right..
And then we focus on, OK, the love of God, the love of Jesus on the cross..
Though there's nothing wrong with that..
However, if you flip the pages of the Bible,.
New Testament, particularly all four Gospels,.
all the New Testament books, especially the letters to the churches,.
especially Paul's letters to the churches..
They often remind the early believers,.
the early church to center their beliefs, not just on the cross,.
not just on the love of God, but on the resurrection of Jesus Christ,.
the son of God..
I don't know if you pay attention to that..
If you check all of the early church documents..
Including Apostles Creed,.
we recite that last Sunday on Easter Sunday.
as we worship together..

$^{41}$The Apostles Creed has 12 articles..
And whoever draft that those 12 articles.
and summarizing the entire biblical faith in 12 statements,.
they did not forget to add, I believe in what?.
The resurrection..
The resurrection of Jesus Christ..
I believe in that..
And that's why we have hope..
And that's why our faith is founded on solid rock..
So that's why I want to start this mini series on living the resurrection.
to remember what resurrection means for us..
And don't just think that this is a footnote of Christianity,.
just a footnote of the biblical faith..
Because again, as I said, all four Gospels and all the epistles,.
all the letters to the churches, especially Paul's letters,.
never stop at the cross..
They move on..
Yes, on the cross, God redeems us..
Jesus died for us, sacrificially for mankind, for sinners..
But all four Gospels, all Paul's letters.
go to a great length, emphasizing.
the fact of empty tomb, the fact that Jesus' tomb is empty.
and the resurrection of Jesus Christ..
The early church did..
And so we should. Right..
So there are two points I want to share with you..
And those are very, very important points for this.
first part of the sermon..
But before we do that, let's read the scripture..
And I picked this passage, this story..
It's a long, kind of long 18 verses.
record of the first Easter Sunday..
I want you to go through this 18 verses, this story,.
this historical account with a mind that,.
you know, you all go students, right?.
You're good students..
As good students, we pay attention to details. Right..
So we're going to read through this 18 verses..
And I want to I want you to pay attention to some of the details.
and ask the question why the details are there..

$^{81}$All right..
Gospel of John, chapter 20, verse one..
Early on Sunday morning, while it was still dark,.
Mary Madeline came to the tomb and found that the stone.
had been rolled away from the entrance..
She ran and found Simon, Peter and the other disciple,.
the one whom Jesus loved should be John..
As you said, they have taken the Lord's body of the tomb.
and we don't know where they have put him..
Peter and the other disciples start off for the tomb..
They were both running and the other disciple outran Peter..
Probably John is the little John and he's younger..
So he's he he's outrunning Peter and he reached a tomb first..
So he stood and look in to the tomb and saw the linen.
wrappings lying there..
But he did not go in..
Then Simon, Peter finally came, finally arrived..
Trying to catch his breath and went inside the tomb..
He also noticed the linen wrappings lying there.
while the cloth that had covered Jesus head and was.
it was fold up and lying apart from the other wrappings..
Then the disciple who had reached the tomb first.
also went to the tomb and he saw and believed..
For until then, they still had not understood the scriptures.
that said Jesus must rise from the dead..
Then they went home..
Mary, however, was standing outside the tomb crying.
and as she wept, she stood and looked inside the tomb..
She saw two white robe angels, one sitting at the head,.
the other at the foot of at the foot of the place.
where the body of Jesus had been lying..
Dear woman, why are you crying?.
The angels asked her, because they have taken away my Lord,.
she replied, and I don't know where they have put him..
She turned to leave and saw someone standing there..
It was Jesus, but she did not recognize him..
Dear woman, why are you crying?.
Jesus asked her, who are you looking for?.
She thought he was the gardener..
So she said, if you have taken him away,.

$^{121}$tell me where you have put him..
I will go and get him..
Mary, Jesus said, she turned to him and cried out,.
"Rabboni," which in Hebrew means rabbi, teacher..
"Don't cling to me," Jesus said,.
"for I have not yet ascended to the Father,.
but go find my brothers and tell them,.
I'm ascending to my Father and your Father,.
to my God and your God.".
Mary Magdalene found the disciples and told them,.
"I have seen the Lord.".
Then she gave them his message..
Eighteen verses..
And there are two points here..
A very important point..
The first one is the historical reliability,.
the historical authenticity of the first Easter Sunday account.
that we found in the Bible..
The historical authenticity..
When you read the record, when you read the story,.
there is a, you find marks and marks of authenticity.
that is a true report..
There's a true news report..
And the second point is the new creation has begun.
on Easter Sunday..
Now, first point, historical authenticity..
Now, like good students, right?.
You're good students and there are many good,.
moody students here, all right?.
I'm going to give you a quiz..
(laughs).
Just kidding..
(laughs).
As you read the story, pay attention to some of the details..
Pay attention to some of the details..
What are some of the details that would bear the marks.
of historical authenticity?.
Now, what do I mean by that?.
What do I mean by that?.
If you know there's an event that happened.

$^{161}$and you turn on the TV or you go to YouTube,.
you read some news report..
You probably read the version from CNN, NBC, CBS,.
or any news channel..
And you find different versions of the same event, right?.
Yes, the details are different..
Some of the minor details are different..
However, these reports have the marks of authenticity..
The marks, what marks them authentic is not.
because all these accounts are exactly the same..
On the other hand, it is because the exact details.
are all different, they bear the marks of authenticity..
The Gospel of Matthew, the first gospel.
of the four gospels, said who first witnessed.
the empty tomb was Mary Magdalene and the other Mary..
The gospel mark, Mary Magdalene, and there's another Mary..
By the way, there were a lot of people called Mary.
back in the first century..
So you go on the street and you call Mary,.
there are probably 10 or 20 people responding to you..
Hey, yeah, right, Mary's here..
And gospel mark, Mary Magdalene and Mary the mother of Jesus.
and other women were there..
Gospel of Luke, the women, and also Mary Magdalene.
and other named women were there..
Gospel of John, we just read, only Mary Magdalene was there.
because the camera is on Mary Magdalene, on her only..
Now, all four accounts are quite different,.
but because they are different,.
that's why they can be historical reliable, right?.
Authentic..
Now, pay also attention to the two pieces of linen cloth..
There's one wrapping the head of Jesus' body,.
the other piece of cloth wrapping his body..
So there are two pieces..
In fact, archeologists tells us.
what we found in the first century tomb,.
all the first century tombs around Jerusalem,.
when they wrapped the body, prepared the body to be buried,.
they have two pieces of cloth,.

$^{201}$one for head and one for the body..
So as you can tell, the account was very, very authentic,.
reflecting the custom of the day..
One more detail..
Both pieces of linen cloth were neatly folded up..
You don't see any signs of disorder,.
or the tomb being robbed, or the tomb being sabotaged..
Everything is in order..
You look into, they look into the empty tomb,.
everything is in order..
How many of you make your bed after you wake up?.
How many of you?.
How many?.
One..
Wow..
We don't, two, okay..
What minorities?.
You don't?.
I don't make my bed..
So in some sense, if someone walked into your room right now,.
they probably would think this is in chaos, right?.
Or probably some thieves just went there,.
just robbed your bedroom, right?.
But on the other hand, the tomb of Jesus was in order..
Everything is in order..
Now, aside all these details,.
I'm sure you pick up the other detail..
The other detail is the appearance of women..
The women are the first witnesses.
of the resurrection of Jesus..
And what is authentic about it.
is that if you understand the first century context,.
the background of first century,.
women were often looked down by men..
No surprise, right?.
It was a patriarchal culture..
You know, all through the years,.
I have collected some of the Jewish writings.
that reflect the men's attitude toward women..
And I want to share with you..

$^{241}$There were a few of them, all right?.
So I want you to read this, not in the Bible,.
not in the Bible, but it's outside of the Bible,.
outside scripture..
And this statement or this sentence, this writings,.
are reflecting the patriarchal attitude,.
perspective on women..
You ready?.
First one..
This is what the rabbi said, okay?.
This is what the rabbi said..
"So long as a man talks too much with a woman,.
he will bring to himself three consequences..
These three bad consequences is first one,.
he, the man, will bring trouble on himself.
if he talks too much with a woman..
Second is, waste his time being spent on studying the Torah.".
Right?.
You want to do Bible study than talking to a woman..
Number three is that you will end up in hell, right?.
Gehenna, right?.
Era of Gehenna..
Now what kind of attitude is that?.
The man are looking down the woman, right?.
Second, again, this is what the rabbi said,.
or the Jewish rabbi of Jesus' day..
He said, "Because women are gluttonous,.
they like to eat, they like to eat a lot,.
for she is suspect of uncovering the instant part.
of her female friends to find out what she is cooking.".
Now that kind of attitude is reflected.
by men who wrote this thing,.
and how they see women, right?.
Because women, they like to cook,.
and they like to eat,.
and they have no self-control..
They don't even lack of self-control with what they cook,.
but they are curious about what other people,.
their neighbor, are cooking..
Number three, Josephus, a first century Jewish historian..

$^{281}$He said, "Let not a single witness be credited,.
but three or two at the least,.
and those such whose testimony is confirmed.
by the good lives..
But let not the testimony of women be admitted to the court.
on account of the levity and boldness of their sex.".
There you go, you can see how men see women..
By the way, that's not my perspective, all right?.
That is the perspective of the first century people,.
especially Jewish society..
So again, if you look at the women.
of the story of the first Easter Sunday,.
you would not expect women to be there..
If someone would fabricate a story.
and make you believe that Jesus did rise from the dead,.
he would not write such a account, right?.
He would say the first one who witnessed the empty tomb.
must have been one of the disciples, one of the males,.
because their testimony has to be.
more credible than the women..
Now, reverse that question..
You're gonna ask, why did the gospel writers,.
all four of them, Matthew, Mark, Luke, and John,.
they recorded the way they did,.
that the women were the first witnesses of the empty tomb?.
Well, very likely, it did happen that way..
So they cannot lie..
They just have to report the truth, right?.
Now, I'm sure if you read the empty tomb story,.
the first Easter Sunday story,.
you probably get the feeling.
that the women were going in and out,.
nobody's tracking them, they're not afraid of anything, right?.
And the male disciples, where were they?.
Where were they when Jesus died?.
They were hiding, they were hiding..
Now, what's going on here?.
Nobody care about the women going in and out of the tomb.
and checking on the tomb?.
I remember some years ago,.

$^{321}$I heard a testimony from a Christian missionary.
who went to Middle East for some mission work..
And he was telling us that it was time of war..
So both sides are fighting,.
with the guns and tanks and all that..
And the defensive side, all the men were hiding.
because they were scared..
But the missionary was telling us that, you know,.
from eight to five in the morning, in the afternoon,.
the women are going on the street and going to supermarket.
and get their vegetables, get their food,.
and then go back home and then they cook..
They bring the kids outside and then they come back in, right?.
So this is kind of the context of the gospel account..
The men were hiding because the government,.
the Roman Jerusalem government,.
were looking for the Jesus followers.
and persecute after them, right?.
They want to pursue them, they want to arrest them.
because they're starting a movement.
that will rock the boat of the society..
They would disturb the peace of the society..
So they were looking after and pursuing the male disciples..
While the women, nobody cared, right?.
Nobody cared because women,.
in those days, they have the anti-woman attitude, right?.
So you look at the empty tomb account,.
the first Easter Sunday account,.
this sort of cultural anti-woman attitude is not there.
because the first witnesses of the empty tomb.
were the women..
In fact, if you track the story of Luke,.
after the women discover that Jesus rose from the dead,.
they went back to the male disciples..
And the male disciples, how did they respond?.
The story sounds like nonsense to the men..
So they did not believe it..
Don't think for a minute, don't even think about,.
okay, because they're ancient people, 2,000 years ago,.
they didn't have iPad, they didn't have iPhone,.

$^{361}$they didn't know anything better..
We know science, we know how the world works, right?.
And don't think that because they're ancient people,.
so they must have believed the resurrection..
No, it says Luke 24, 11..
The story sounds like nonsense to the men,.
to the male disciples..
So you get the feel that none of the disciples is ready.
for the resurrection of Jesus..
Now, to tell you the truth, in Jesus' day,.
most of the Jews believe in bodily resurrection..
They believe the physical resurrection of the human body..
Because there's a hope in Jewish theology in those days,.
that in the eschaton, in the final day,.
when everything, God will set everything right.
in the final day, the eschaton,.
the last day of human history,.
God will come and his glory will descend,.
and everything will be set right..
There will be justice, there will be peace.
and righteousness and love, and all that stuff..
And when everything will be consummated,.
the new creation will finally come..
And they all believe that..
Most of the Jews, they will believe.
and look forward to that day,.
and all the dead will rise from the dead..
So they have some sort of idea of resurrection..
But what caught them surprise.
on the first Easter Sunday morning.
was that Jesus rose from the dead first,.
before everyone else..
Jesus is the first man who conquered the tomb,.
and he was resurrected from the tomb,.
and his tomb is empty..
And what caught them surprise was this..
Hey, we were looking for some day.
that God will consummate every single thing.
in our universe, the new heaven and new earth will descend,.
and glory will fill the entire world, and entire cosmos..

$^{401}$But hey, now Jesus rose from the dead,.
and that is surprising, because we didn't expect that..
I remember when my kids were younger,.
they go to elementary school,.
and there were some days they set aside,.
it's called the pajamas day..
You know pajamas day?.
They wear pajamas, and then they go to school,.
and then the teachers will also wear pajamas.
and go to school, right?.
And they kind of have fun, right?.
You know I think about it,.
the first Sunday was kind of like that..
It's like a pajamas day..
Okay, so imagine with me, imagine with me,.
in the early morning, six o'clock, six a.m.,.
the sun rises, breaking the dawn,.
six o'clock, and where are you?.
You're still in bed, right?.
You're still in bed, and you're still in your pajamas..
And you hear knocking on the door..
Six o'clock, you're still in bed with your pajamas..
Six o'clock, somebody knocking on your door loudly..
And from the outside, you hear a voice,.
hey, this is Joe, Joe Biden..
I'm at your door..
I have coffee and I have donut ready for you..
Please let me come in, you open the door..
And how will you respond?.
How will you respond?.
You probably will brush your teeth and wash your faces,.
and you're still in your pajamas,.
and you open the door, right?.
And you hear, hey Joe, right, can you help me with taxes?.
You know, I need some help, right?.
You invite him to come in,.
he is the biggest man in this country..
And coming, knocking on your door..
Now this is kind of like the first Easter Sunday morning..
The disciples were not quite ready..

$^{441}$They were still in their pajamas..
They just don't know what to do yet..
The first Easter Sunday was kind of like that..
Jesus rose from the dead unexpectedly and surprisingly..
And that's why in the New Testament, in the early church,.
they were trying to sort out.
what is going on in the empty tomb..
What is going on in the bodily resurrection of Jesus?.
When Jesus met them, talked to them about the kingdom of God,.
they were trying to figure out what is going on.
because yes, we're looking forward to some day,.
the final day, everybody will rise from the dead..
But however, Jesus is the first guy who rise from the dead..
And that's why the New Testament uses some terms.
to try to understand what is going on.
with resurrection of Jesus Christ..
They describe Jesus Christ,.
one of the example is first fruit of the harvest..
Now if you're farmers, you know how it works..
The first produce harvest from your field.
will guarantee the rest of the field.
will have harvest for you..
And that's why the first fruit.
will guarantee the rest of the fruit..
First fruit is the resurrection of Jesus Christ..
You see him, he rose from the dead, his tomb is empty,.
and he has defeated death, decay, and destruction..
He has defeated evil and will guarantee the finishing off.
of everything God is going to do..
If you read the Bible carefully,.
actually you don't have to read much..
If you read the first two chapters of the Bible,.
Genesis one and two,.
and you read the last two chapters of the Bible,.
which is Revelation chapter 21 and 22, right?.
First two chapter and last two chapters,.
you figure out the entire Bible is a grand story..
It's a grand narrative that God created the world all good..
That is the original creation..
And sin came into the world and evil destroyed the world.

$^{481}$and twisted everything else, right?.
So partially we can see the original creation.
is still good today..
However, there were a lot of things,.
you observe the universe are being twisted..
You see pollution, you see racism,.
you see all kinds of different issues with our world..
It's good and bad, it's a mixed bag..
But God is going to bring everything to the finish line,.
which is the new creation..
You read the last two chapters of the Bible,.
Revelation chapter 21 and 22..
You see that it's not just us escaping from the earth,.
escaping from this physical world universe..
It's that the, it's on the other hand,.
it's the heaven descending to the earth..
It's like our Lord's prayer..
You know how do we say that?.
Our Father in heaven, hallowed be thy name,.
thy kingdom come, the future kingdom will come..
Your will will be established on earth as it is in heaven..
This is kind of like a vision, a future vision.
of what God will be doing..
And what God will be doing is that he will bring heaven.
and earth together and the new earth.
and new heaven will arrive..
And now with the resurrection of Jesus Christ,.
all of a sudden the future kingdom launch into our world..
The new chapter has begun..
The new creation has become..
And we'll see this new creation being fulfilled finally.
in the final day..
However, in the Jesus resurrection,.
we see that this chapter has already opened..
This chapter has already launched..
So my second point is that not only did Easter Sunday.
account in the gospels were authentic historically,.
the new creation in the first Easter morning.
was that the new creation has come..
It has arrived..

$^{521}$What we looking forward to the new kingdom,.
the new reality has already broken into our lives,.
broken into our world..
Easter Sunday is the beginning of that new creation,.
heaven and earth..
But those of us, I know not many of us here who are older,.
we should remember there's an old hymn called In the Garden..
Now I don't know if that rings any bell,.
but you check In the Garden hymn, you check the lyrics..
The hymn says, you know, me and Jesus walking in the garden..
Everything is so nice..
You know, the birds are singing and the nature is so nice..
The weather is so nice, unlike Chicago..
And everything is so warm and it's like a spring..
And I'm walking with Jesus in the garden..
It's a song, it's a hymn that communicates,.
that reflects the communion we have,.
the fellowship that we have with God, right?.
But if you think about it further,.
you know, the Bible or what Jesus has done for all of us,.
is not just going back to the good old days..
It's not going back to the first chapter of the Bible,.
going back to Eden..
It's not back to paradise..
It's a linear development that we look forward.
someday salvation will come to all of us..
And that salvation is not just our soul or our spirit.
being released or escape from this physical dying body..
What God is doing, the creation and redemption project,.
is to rescue everything in our cosmos..
Not only our soul will be escaping.
from this physical bodies, this dying body,.
and then we're just kind of flowing in the heaven.
and sitting on the cloud and that's it..
That's only the first part..
The second part, the final part, is the physical redemption..
That is the final destination..
That is the final finishing off.
of God's entire project of redemption..
Redeeming all of us, even our physicality,.

$^{561}$our physical body..
Parents with young kids probably have this sort of experience.
and I remember when I was younger and our kids were younger,.
sometimes we plan for a long road trip..
So we drive from Chicago to,.
let's say for example, we go to Disney World..
Disney World in Florida..
And how many hours of driving that would take?.
Many, many hours..
So on the long road trip, we'll just get everything packed.
and then put in a minivan and then we'll just drive off..
And then there's always many stops in between..
So as today, you go on the highway.
and then sometimes you need a stop in between,.
bathroom break or lunch break,.
and then we go to McDonald's..
Our kids will go into McDonald's and as you can tell,.
McDonald's is a very nice place,.
and they go in and they get the hamburgers,.
they get fries and happy meals.
and almost all the McDonald's have playland.
attached to the restaurant..
And then the kids will, after they eat,.
and then they'll go in and play and take off their shoes.
and just play in the playland..
And there are times that daddy, mommy are saying,.
"We gotta go because we are not there yet.".
And they were crying and then yelling at us.
and almost we have to abuse them.
to grab them out of the playland, right?.
Because we are not there yet, right?.
They are so happy because of the playland,.
they thought this is the Disney World,.
this is the whole story..
No way, I said, "Son and daughter, we gotta go..
"Come on, man, we gotta go..
"Wake up, we have a final destination to hit.".
Now sometimes I feel that Christian faith.
and Christian life is like that.
because we stop at the playland,.

$^{601}$just the grass, just God loving us..
Yeah, good, that's very nice..
However, these two stories telling us.
the new creation yet to come and has already here, right?.
Someday it will bring to the finish end..
So don't think for a moment that God merely redeemed us.
so that our souls will go to heaven.
and we're sitting on a cloud..
No, this is not the salvation that we have in the Bible..
Otherwise, there's no need for empty tomb..
There's no need for resurrection of Jesus Christ.
in his empty tomb..
The best is yet to come and it has already begun..
It is so revolutionary to the first century Christians.
that they even changed the day of their worship..
For Jewish Christians, they used to worship.
on the seventh day, the Shabbat of the week,.
the seventh day, the last day of the week..
However, when Christians came to faith,.
they understood the reality of resurrection..
They changed the worship day to first day of the week,.
Sunday of the week, that they worship Jesus Christ.
because he has begun the new creation..
This is why when the Bible talks about resurrection,.
it's not an afterthought..
It's the core of our belief..
It's the foundation that we have hope..
Yes, we live in a sad world..
This is a broken world in many different ways..
You look around our cosmos, our world,.
you see, hey, everything is not working right, right?.
There are issues, there are problems,.
there are twists and distortion that we see..
However, it's the linear hope that someday God will come.
and Jesus will bring about the new heaven and new earth..
And in the empty tomb, now we see that..
The new creation, the new chapter has begun..
When you accept Jesus Christ, this new reality has begun..
You have a new life and this new life will consummate,.
will finish off someday in the future..

$^{641}$And who guarantees it?.
It's the empty tomb of our Jesus Christ..
I wanna close with the last verse of this Easter story,.
John chapter 20, 18..
It says, "Mary Magdalene found the male disciples.
"and told them, 'I have seen the Lord..
"'I have seen the Lord..
"'The new life has begun in all of us in Christ.'".
You are being sent, just like Jesus sent Mary,.
the first witness, to go out to her world.
and to tell of the resurrection of Jesus Christ,.
the new life that we have in Christ..
How do we do that?.
You know, in our daily life, you have your own role..
Maybe you're a housewife,.
maybe you're a worker in your workplace..
Maybe you're someone who thought.
that you're not significant enough..
Well, we're all are very small figures in our own world,.
but we make the resurrected Jesus Christ.
visible by our lives..
Talk in a new way, in a biblical way..
Relay in a biblical way..
Relay to other people..
Care for other people in a new way.
because the new life has already begun..
Shall we do that?.
Shall we do that this week?.
Bring that message of resurrection,.
that hope, that new life, into your own life..
Relay to your spouse, relay to your kids,.
your neighbors, in a totally different, brand new way.
because that is what Jesus' resurrection is all about..
Yes, the best is yet to come,.
but it has already begun..
And let it begin with our lives..
Let's all pray together..
Father, we praise you because of the resurrection.
that we have in Jesus Christ, our Lord..
As we reflect on the teaching of resurrection,.

$^{681}$the accounts of resurrection,.
praise you, thank you..
Thank you for the arrival of the kingdom,.
the new creation has begun..
And we want to offer our praises and our worship to you..
And may our lives be visible to other people,.
the reality of the empty tomb..
In Christ's name we pray..
[BLANK AUDIO].
\newpage



\section{約翰福音}
\label{sec:PhsM03F2OIM}
\textbf{復活的福音:約翰福音 20 章 - 陳崇基牧師}
\newline
\newline
連結: \href{https://youtube.com/watch?v=PhsM03F2OIM}{\texttt{https://youtube.com/watch?v=PhsM03F2OIM}} ~~~~ 語音日期: 2023-04-18
\newline
\newline
\hyperref[sec:_Z2VGiey2J8]{\small{< < < PREV SERMON < < <}}
~
\hyperref[sec:index_chronic]{\small{[返順時目]}}
~
\hyperref[sec:index_scriptual]{\small{[返順卷目]}}
~
\hyperref[sec:SK0VEvCIu6I]{\small{> > > NEXT SERMON > > >}}
\newline
\newline
$^{1}$咁今日ge 講到呢中英中英文堂呢都係Mark Sirge .
咁而家將時間交比佢.
早晨各位.
大家好.
我要同大家講Happy Birthday.
因為呢我地成會七年.
所以我地呢今個主日呢上個禮拜我地arm arm Celebratejor .
我地就開始七歲.
你唔介意ge 話你同隔離gwo 個講Happy Birthday.
有冇邊個今日生日架.
冇啊OK.
咁即係我地大家教會一齊生日啦係咪.
All right very good very good.
我地上一次呢係復活ge 主日.
咁所以我地有堂慶啦同埋我地去Celebrate我地ge Easter.
咁但係好多時呢我地好多時呢.
就會留低jor 停jor 係呢一個地步.
跟住我地覺得上個禮拜係Easter Sunday.
然之後呢我地就繼續Move on 我地就唔理啦.
我已經過jor 復活節啦我地可能你地有細路ge .
我見到你好多相片走去lor 雞蛋走去做Egg Hunt等等.
Well 我地好多時呢就停留jor Easter Sunday.
Celebratejor 然之後呢我地Put it behind.
然之後我地就唔記得到底發生咩事.
好多時呢我地諗我地基督教信仰ge 時候.
諗起聖經ge 信仰ge 時候我地都有咁ge 感覺ge .
我地好多時就係會諗十字架.
The Cross.
咁呢個冇錯ge .
上帝ge 愛係十字架上面表明出黎係咪.
耶穌幾多係十字架上面為我地犧牲呢個冇錯ge .
不過好多時呢我地淨係諗起十字架.
就係耶穌為我地犧牲就停jor 係gwo 度.
如果睇下四福音書.
瑪太馬可路加同埋約翰.
再睇埋新約ge 書卷ge 話.
你會知道基督教信仰唔係淨係基於上帝ge 愛ge .
往往初期教會只能夠可以傳揚主ge 福音.
能夠周圍去建立教會並且歷世歷代係不敗ge .
唔係淨係因為上帝ge 愛.

$^{41}$而係因為上帝ge 兒子復活.
復活ge 主日唔係單單一.
52個禮拜裡面一個ge 主日.
唔係單單一個ge 節期.
我地守完就完啦.
而係一個ge 基礎.
正因為耶穌係死裡面復活.
所以點解初期教會能夠有咁ge 勇氣.
能夠有咁ge 能力.
可以去周圍.
穿州過省.
去到去傳揚主ge 福音建立上帝ge 教會.
所以呢係今個主日開始.
我想同大家諗一諗.
到底復活節係一回乜ye 事.
你話唔係喎 復活節過jor  仲講復活節做咩.
就係因為復活節完jor .
我地先應該講復活節.
我地先諗起.
以新約裡面講到我地基督ge 信仰.
gwo 個ge foundation 基礎.
就係耶穌基督ge 復活.
所以今日好簡單.
我想有兩點係同大家分享.
但係未分享之前.
我地先讀聖經裡面記載.
第一個復活ge 主日.
到底發生咩事.
係咪好有feel.
過jor 復活節我地仲係度睇復活節ge 經文.
你睇下新約裡面好多時就係唔放手.
唔係耶穌係十字架上死就完.
係有主從死裡面復活.
我地睇睇呢個故事.
約翰福音第20章第一節.
七日ge 第一日清晨.
天還黑ge 時候.
沒大拉的瑪利亞來到墳墓.
看見石頭已從墳墓挪開了.
就跑去見西門彼得和耶穌所愛的那個門徒.

$^{81}$呢個應該係約翰.
對佢地說.
有人從墳墓裡面把柱移走了.
我地不知道佢地把他放在那裡.
彼得和那門徒就出去往墳墓去.
兩個人同跑.
那門徒比彼得跑得更快.
呢個可能約翰比彼得跑得更快.
可能佢後生D .
應該佢應該後生D .
後生仔應該跑快D .
後生仔就後面.
到了墳墓低頭往裡面看.
看見西瑪布還放在那裡.
只是沒有進去.
西門彼得個老D ge .
隨後ge 終於喘晒氣咁到.
進了墳墓看見西瑪布放在那裡.
又看見耶穌ge 果頭根.
沒有跟西瑪布放在一起.
係另一處捲起ge 摺起jor ge .
然後先到墳墓ge 那門徒約翰.
亦都進去.
他看見就信了.
他們還不明白聖經所說.
耶穌必須從死人中復活的意思.
於是兩個門徒回自己的住處去了.
瑪利亞卻站在墳墓外面哭.
他哭ge 時候低頭往墳墓裡面看.
看見兩個天使著著白衣.
在安放耶穌身體的地方坐著.
一個在頭一個在腳.
天使對他說.
夫人你為什麼哭.
他對他們說.
因為有人把我主移走了.
我不知道他們把他放在那裡.
說了者說.
他轉身來看見耶穌站在那裡.
卻不知道他是耶穌.

$^{121}$耶穌問他.
夫人你為什麼哭你找誰.
瑪利亞以為他是看完子的.
就對他說.
先生若是你把他移了去.
請告訴我你把他放在那裡.
我去把他移回來.
耶穌對他說.
瑪利亞.
瑪利亞轉身來用希伯來話對他說.
拉玻璃.
拉玻璃就是老師拉筆的意思.
耶穌對他說.
不要拉住我.
因為我還沒有升上去見我的父.
請你到我的弟兄那裡去告訴他們.
我要升上去見我的父.
也是你們的父.
見我的上帝也是你們的上帝.
穆達拉的瑪利亞就向門徒報信.
我已經看見了主.
他又把主對他說的話告訴門徒.
這個就是第一個主日.
清晨發生的事.
今天我想和你分享兩點.
第一點.
如果你去到作為一個好學生.
我們都是好學生.
我們看一段文字的時候.
我們學一下怎樣去看它.
看這些復活的記載.
復活的主日的故事.
裡面的技術裡面的細節等等.
我們發現這些復活的主日的記載.
都是有歷史的寫實性.
就好像你看一個新聞報道.
上兩個星期.
我們一個多星期前.
這附近有大火.
在阿脫那裡有大火.

$^{161}$你看看不同的新聞報道.
那些細節都不一樣.
但是所報道的東西都是有共通點.
你就發現有件事發生了.
你看復活的記載其實都是一樣.
有歷史的寫實性.
這是第一點.
第二點就是復活的清晨.
復活的主日是一個新的創造的開始.
第一點.
就是那個所謂的Historical Authenticity.
你看那個復活的記載的時候.
你會找到一些證據.
或者找到一些記號.
你會發現那些不同的人的記述裡面.
那些細節是很不同的.
馬太,馬可,路加和約翰.
四幅音書.
記載同一件事.
都是有些不同.
如果你看看馬太的福音.
去記載復活節的清晨.
第七天的頭一日.
清晨的時候發生了什麼事.
有誰在呢?.
就是有馬尼亞.
和另外一個馬尼亞.
當時很多時候女孩子都叫馬尼亞.
你生女兒就叫馬尼亞.
你在街上叫馬尼亞.
有十幾個人認你.
這是一個很常見的名字.
但是馬太福音就記載了.
有莫大拉的馬尼亞和另外一個馬尼亞.
然後馬可福音就記載了.
有莫大拉的馬尼亞.
和一個雅國的母親.
叫什麼名字?.
叫馬尼亞.
然後有另外一個女人.

$^{201}$然後去到勞加的時候.
又是記載有莫大拉的誰?.
馬尼亞.
沒錯.
還有其他的女人在當中.
約翰福音剛剛看的那個故事.
就只有莫大拉的誰?.
馬尼亞.
沒錯.
又是馬尼亞.
但只是看到馬尼亞.
這些不同的四個的記述.
四份的所謂報章.
去記載這件事情的時候.
細節都是不盡相同.
但有一些共通點.
有什麼共通點呢?.
他們記載了一個空墳墓.
並且那個空墳墓.
不是有盜賊來搶東西的.
你有沒有留意剛剛看那段經文的時候.
有一個細節.
是作者特意去留神.
要我們知道的.
就是他說有兩條的細麻布.
一條是裹頭巾.
一條是包屍身的布.
有兩塊的.
我們找到第一世紀的墳墓.
是這樣撞當時的屍身.
是有兩塊布.
一塊布是蒙頭的.
一塊布是蒙屍身的.
他怎樣記載呢?.
他說裹頭巾和細麻布.
就是放在各一處.
然後怎樣記載呢?.
是摺起了的.
很整齊的摺起了.
你今天早上起床的時候有沒有摺床?.

$^{241}$沒有吧?.
如果有人現在進入你的睡房的時候.
可能他會覺得有賊仔進去.
因為太亂了.
有誰摺床的?.
Very good.
我不會摺床的.
但我老婆會摺床.
你看到這個空墳墓的記載.
他特別講到什麼呢?.
就是兩塊的毛巾.
兩塊的細麻布.
很整齊的放在那裡.
那個墳墓是沒有被盜獵過的.
全部東西都是in order.
很整齊的放在那裡.
這是第一.
第二.
四個福音書記載第一個復活的主人.
都有記載哪一個是第一個的見證人.
什麼名字呢?.
瑪麗亞.
是女人.
這個很奇怪.
如果你想去fabricate一個故事.
在第一世紀的背景裡面.
如果你想去作一些東西.
以致有公信力的話.
你就大概不會這樣去記載第一個的復活清晨.
因為第一班見證人就是女人.
不是男人.
以前第一世紀猶太人的社會.
上到法庭的.
或者做到法庭上見證人的.
全部都是男人.
女人是不可以上法庭的.
因為女人說的話是不可信的.
你今天是不是這樣想?.
所以你最近有沒有收到那些Jury duty的信?.
以前第一世紀女人就不會收到那些信.

$^{281}$因為全部都是照男人去做法庭的見證人.
或者在法庭裡面工作.
女人說的話大概是不可信的.
因為你知道當時是男尊女卑的社會.
Patriarchal society.
所以男人是尊貴一點的.
我自己過去這麼多年.
我搜集過一些當時的一些人寫有關女人的東西.
我覺得這幾句是很表達到當時.
一般來說是怎樣去看待女人.
或者女人所說的話.
我和你分享一下好不好?.
拉比說什麼呢?.
智者曾經說.
男人和女人談話太多的話.
就會有三個後果.
第一個後果就是他會自己帶來禍患.
第二就是浪費時間.
忽視了言讀律法的時間.
第三就是最終會接受地獄.
對不對?.
如果和女人說話太多的話.
不如多點茶經吧.
OK 第二句.
很有回響.
他說女人是貪吃的.
因為女人是涉嫌揭開別人的論社的instant pot.
想知道別人煮什麼.
OK.
這些都是猶太拉比的說法.
猶太拉比就說這些.
其實他們都很官人於微.
他們知道女人的特性是怎樣.
但他在說什麼呢?.
其實就是女人的天性.
你弄清楚.
不是我這樣看的.
我很尊重女士.
我不是這樣看.
是以前的猶太拉比.

$^{321}$所以做Civus.
第一世紀的歷史學家.
猶太人.
他說見證人的旨在不可受理.
但要三個或最少兩個.
其實提摩太前遊書都有說.
要兩個或三個.
而這些見證人.
是要有好的生命和見證.
但不可承認女人的正宮.
因為女人是率性和魯莽.
誰贊成?.
OK.
所以你可以看到.
如果第一個的副節.
當他們如實去報道的時候.
他們就沒得選擇.
他們一定要告訴你.
第一班的見證人.
就是女人.
所以你就可以看到.
這些所謂的技術.
我們做一個好學生.
那些技術是那種.
我們叫Authentic.
很實在的.
真是這樣報道.
其他的男人都去哪了?.
全部都躲起來.
In fact.
路加福音第十四章.
都有這樣的味道.
有些所謂看不起女人的態度.
就是那些女人去見到空墳墓的時候.
然後她們回來.
和那些藍門徒.
全部都是躲起來的那些藍門徒.
來解釋的時候.
他們說什麼呢?.
這些使徒就以為.

$^{361}$這些女人所說的話.
就是無稽之談.
就不相信.
Nonsense.
或者以前和學本也做.
是胡言.
胡說八道.
所以你看到第一個主人.
就發生這樣的事.
很奇怪.
女人就出去.
她們想打理.
耶穌安葬的事.
因為之前的安息日太急了.
搞不定.
因為安息日不可以做事.
所以七天的第一個早上.
早上她們就去.
帶香膏等等.
去預備那個師生.
然後每樣東西都預備好.
搞定了.
第一班去到的就是女人.
你會發現.
四方的記載.
復活節的清晨.
女人都是自出自入的.
男人就因為怕迫害.
然後她們全部都躲起來.
這個其實很符合.
中東的文化.
我記得有一次聽一個.
去中東宣教的.
傳教士.
他就回憶中東.
某些地方.
去打仗的時候.
戰火連天的時候.
他說當然.
他就是躲在家裡.

$^{401}$但他家裡望出去的時候.
他發現全部那些.
所謂守衛的男人全部都躲起來.
但是每一天.
從早到晚那些女人.
都是照樣出街買菜.
然後回家煮飯.
為什麼呢?.
因為那些男人不看女人在眼內.
對不對?.
她們可以.
手無殺雞之力.
當然可能你們.
比很多男人更愛宰雞.
但你看到男人.
怎樣看女人在這樣的環境裡面.
所以你看到.
這些節節記載的時候.
你就會發現.
好像他們.
找到這個空墳墓的時候.
不單止是如實的記載.
也會想起.
到底為什麼會這樣發生.
好像大家不是很預備好.
當時的猶太人.
他們是相信身體復活.
大部分的猶太人.
除了殺到國人之外.
法利錯人等等.
他們相信有身體復活.
他們相信.
但身體的復活.
是何時會發生呢?.
在當時猶太人的神學裡面.
就是將來的.
那個所謂末世.
末日的時候.
那個所謂Eschaton.
最終最後一天的時候.

$^{441}$才會發生.
所以他們就沒有想過.
突然間中間會有第一個人.
耶穌基督.
救世主 Messiah.
會先復活.
然後他的復活會帶來其他人的復活.
所以大家不是很.
明白.
不是很準備好.
我相信你看.
復節的記載都是一樣.
這個也是一個如實的報道.
不是說先前.
到第七天的早上.
我們記得.
星期五他剛死.
但之前拉比說過.
第三天復活我們在等.
我們在等待交定鬧鐘.
我們知道是復活了.
其實大家都好像.
不是很聚財.
因為.
怎會有這樣的發生.
是很驚訝的.
我記得小朋友小的時候.
他們有時候上學.
有些Pajamas Day.
小朋友有嗎?.
穿睡衣上學.
我第一次接觸的時候.
奇奇怪怪.
為何會穿睡衣上學.
原來整間學校.
由大到小都是穿睡衣上學.
他們就很開心.
因為他們不需要.
怎樣.
裝漂亮.

$^{481}$他們就拿著毛公仔.
然後就上學.
有些這樣的事.
你想想復活的主日.
清晨的時候.
其實有少少像Pajamas Day.
你想想.
早上六點鐘.
太陽剛剛升起的時候.
剛日出的時候.
你在哪?.
在床上嗎?.
還在睡覺?.
突然間.
你聽到有人敲門.
突然間有人敲門.
很大聲的在門口跟你說.
This is Joe.
Open the door.
Joe Biden.
美國的總統.
突然間六點鐘.
都未夠.
他就來探訪你.
I got coffee and donut.
Open the door.
然後你會怎樣做?.
你會有什麼反應?.
你會叫他過主?.
你立即叫他擦牙.
洗臉.
擦眼屎.
穿著睡衣.
然後就想開門.
就是一種這樣的情況.
你不是很準備好.
突然間這樣發生.
耶穌基督的復活.
那天的清晨.
就是這樣的情況.

$^{521}$不預期的.
但卻這樣發生.
跟我們神學的預期.
都不同.
因為我們說最後一天才會發生.
為什麼突然間這個時候發生.
這個就正是新約.
開始要去到.
處理的一個問題.
就是耶穌第一個.
從身體復活的.
以至將來.
所有人身體復活的一個證據.
一個保證.
新約裡面有不同的字眼.
來形容這件事.
叫做first fruit.
就是說一個.
初熟之果.
耕種的時候.
第一批的收成.
或者出來的土產.
就是保證你將來整塊田.
就會有土產.
耶穌基督復活.
為什麼像初熟之果.
就是第一批的成果.
是保證和保證.
將來的成果.
這個救世主.
大家記錯了.
打敗了死亡.
打敗了罪惡和邪惡.
打敗了撒旦.
他是保證將來.
一定會得勝的.
原來上帝.
要做的事.
不只是單單十字架上的兒子.
為我們死為我們復活.

$^{561}$為我們犧牲這麼簡單.
原來上帝一直以來想做的事.
就是想新天新地.
能夠駕臨到我們中間.
就像《使徒文》裡面的祈禱一樣.
我們在天上的父.
願人都尊你名為聖.
願你的國降臨.
願你子兒行在地上.
如同行在天上.
如果你是熟悉聖經的話.
其實也不需要太熟悉.
你只是看四章的聖經就夠了.
《創世記》一二章.
上帝原本的創造.
聖經最後那兩章.
啟示錄二十一,二十二章.
上帝最後的創造.
新的創造.
新的新天新地.
是何時會到.
猶太人就覺得將來一定會到.
但怎知耶穌突然復活的時候.
從死裡復活.
就保證將來.
新天新地就會駕臨.
並且現在已經到了.
所以如果你這樣看的時候.
確實第二個point就是.
不單止那個技術.
是有歷史的時間.
技術是有歷史的寫實.
復活的主人.
也是一個新的一章開始.
新的創造開始.
那個新的創造將來.
會完完全全來到.
但現在已經展開了.
在哪裡展開? 在耶穌基督的復活裡面.
展開了.

$^{601}$所以整班新約的教會.
所有的信徒.
他們就快點去處理一下.
猶太人耶穌基督.
要在我們身上 要在這個世界裡面做些甚麼.
原來整個的.
聖經裡面告訴我們.
是一個創造.
和一個新的創造.
一個大的故事.
上帝曾經創造過.
那個原本美好的創造.
今日仍在 雖然是有罪惡.
有邪惡 去到天污了.
或者拆毀了 有很多東西扭曲了.
但仍然是美好的.
而且上帝不斷來到更新.
將來的新創造.
新天新地一定會來到.
你想想 最後最後那兩章.
那裡說到.
天堂會怎樣來.
天堂是從天降下來.
聖聖耶路撒冷.
Holy Jerusalem.
The New Jerusalem.
是會從天 Heaven 下來到地上.
天和地就合一.
上帝原本創造.
就有天和地.
一直去到最後.
ascension的時候 末日來到的時候.
最後一天來到的時候.
上帝所有的工作.
去救贖和創造的工作.
就會完結.
我們何時知道.
這件事是否會有保證 就是耶穌基督.
第一個從死裡復活.
來到去保證.

$^{641}$所以整個聖經是一個直線.
由開始原創.
到上帝新的創造.
就在耶穌基督.
的墳墓裡面.
去誕生出來.
由祂出墳墓那一天.
就展開了.
一個新的創造.
這個就是我們有的盼望.
為甚麼.
我們今天見到.
原本上帝創造的東西.
是扭曲的 但我們仍然可以讚美.
仍然可以敬拜 仍然可以跟隨耶穌基督.
仍然有力可以向前行.
原因就是知道.
大結局.
一定會臨到.
而那個大結局.
初常的大結局.
能夠讓我們見到.
get a glimpse.
就在耶穌基督的身體復活.
那個空墳墓裡面出來.
在我們中間有些比較老的信徒.
我們有時會想起.
一些舊的詩歌.
舊的聖詩.
我不知道你記不記得.
有一首詩歌叫做在花園裡.
in the garden.
那首歌講到.
人和神的契合.
是一首很好的詩歌.
那首歌如果你看看.
歌詞就講到.
他和耶穌在花園裡面.
能夠同行一起行.
幾乎像拍拖一樣.

$^{681}$在花園裡面.
去到同行 然後有鳥叫.
然後很美景.
跟上帝一個.
fellowship 一個關係.
一個很親近的關係.
我那時聽那首歌的時候.
我都發現.
如果只有那首歌.
是有點不妥的.
因為上帝的救贖.
不是要我們回到.
以前的伊甸園.
上帝的救贖.
是向前的.
是要我們去到將來.
新天與新地.
請你不要搞錯.
基督教信仰.
不是只叫我們靈魂得救.
然後這個身體.
一個血肉的身體.
是一個敗壞的身體.
好像一件衣服一樣.
我們就脫下.
然後我們上到天堂.
那裡就是我們最終極的天降.
基督教信仰.
不是停在那裡.
基督教信仰.
是講身體的救贖.
即是說這個物質的世界.
我們物質的身體.
是會被救贖.
這些在我們series之後.
我們仍然會講.
但我想講清楚.
就是剛才我講的那個大的故事.
就是天與地最後臨到的時候.
上帝的榮耀住在中間的時候.

$^{721}$那個是yet to come.
還未臨格的結局.
那個真真正正上帝整個的project.
在世界裡面要行的事.
就是要我們的身體.
都會救贖.
我們身體沒錯.
我們會敗壞.
我們每一日隨著年月歲月.
是不留人.
最近你知道我媽來了.
我媽剛來到.
她一落機第一句跟我說什麼.
她說Mark Chan你為什麼不染頭髮.
我一看鏡子.
我說不是啦.
不是真的這樣.
可能燈光的問題.
但我們是會扭壞.
我們會扭壞.
但上帝在神的兒子耶穌基督的復活裡面.
就告訴我們原來我這個會扭壞的身體.
不是脫了它丟了它.
而是上帝會更新.
將來一定會有身體的復活.
其實這個史道信經裡面.
史道信經是一個很有智慧的總結.
總結整個基督教信仰所有基礎的東西.
其中一條就是我信身體復活.
你有沒有留意.
我信身體復活.
不是一種幻象.
而是第一個復活的精神告訴我們.
原來真正目的地在將來.
我們的救贖在將來.
所以縱使今天這個身體會扭壞.
不怕.
因為將來這個取死的身體一定會在榮耀裡面復活.
小朋友小時候.
大家在我中間有些家長.

$^{761}$可能都有這樣的經歷.
就是你可能計劃一個很長很長的路程.
大家一家人上一輛車.
然後我們去一個很遠的地方.
比較遠的地方我們去玩.
打個比喻去迪士尼.
誰去迪士尼.
去科羅達.
科羅達那個比較大.
去科羅達要開多久車.
很多個小時車.
你去一個很長很長的路程.
然後你開車的時候.
因為小朋友小.
你就一定要讓他們有時候下車停一停.
他們要去廁所也好.
要去吃東西也好.
大人也要.
可能大人的需要更加大.
我們就停下來.
我們就停下來一個地方.
最方便停下來的是什麼.
麥當勞.
我們就在麥當勞買東西吃.
麥當勞很好.
麥當勞旁邊有廁所.
有一個所謂的playland.
那個遊樂場.
我們就在這裡買一點東西吃.
然後他們就去.
就去那些playland那裡玩.
他們玩到樂而忘返.
樂而忘返.
這裡真的很好玩.
叫他走的時候他不肯走.
然後你抓他走.
幾乎要abuse他.
你要扯他走然後他就哭了.
他不明白這是什麼.
他說這個是playland.

$^{801}$你還沒到迪士尼world.
你最終的目的地還沒到.
你就當這個就是你的目的地.
怎麼搞的.
你就想那些酒店關了那麼多次.
幹什麼.
你旁邊的playland已經玩夠了.
這個就是我們很多時候.
了解信仰的問題.
我們就好像小朋友一樣.
我們覺得這個是playland.
上帝在十字架上.
為我們成就的救恩使我們罪得赦免.
以致我們罪得赦免.
ok了我們靈魂得救就搞定了.
這個就是我們的playland.
但真正目的地.
是還沒到.
真正的escathon來到的時候.
我們更加會surprise.
我們更加會.
見到更大的榮耀.
因為我們不是純粹.
被上帝救贖.
而是我們靈魂上天堂.
然後我們飄來飄去就算了.
這個是希臘哲學.
這個敗壞的世界.
我們離開我們.
in fact大部分.
初期教會的異端.
都是在說這個.
這個物質的世界是邪惡的.
扭曲的.
我們這個身體是邪惡的扭曲的.
沒有用的.
甚至今天所有的意識形態.
都很多時候是初期教會的翻炒.
是不是.
這個身體是取死的身體.

$^{841}$有什麼意思.
我們怎樣去define它.
叫它gender也好.
什麼的定義也好.
其實不重要的.
因為這個身體沒有用.
是不是.
No 上帝說.
他原本創造的東西.
他會再recreate.
在耶穌基督的復活裡面.
那個新的.
新天新地.
is yet to come.
and it will come.
第一個復活的主人.
是一件這麼大的事.
所以為什麼新約的教會.
見到一件這麼大的事的時候.
連他敬拜的那一天都轉了.
本身是猶太人基督徒的傳統.
就是18那一天.
第七天.
七天裡面的第七天.
但現在轉了一個禮拜的頭一天.
七天的第一天.
他們來到敬拜.
因為他們知道在基督的復活裡面.
他們知道將來會有這樣的反應.
所以為什麼.
在聖經裡面.
從頭到尾特別是新約.
提到我們基督徒有盼望.
基督徒我們有動力的時候.
原因是什麼.
不是純粹十字架感動我們這麼簡單.
是復活的大能臨到.
在新天新地.
那個大結局已經揭開了.
來到我們中間.

$^{881}$所以為什麼我們會有盼望.
每一次提到復活的時候.
就告訴我們.
我們仍然有工作要做.
並且我們參與上帝的工作.
那個工作一定會有大團圓結局.
最後結束的時候.
20章第18節.
這個故事最後說什麼.
耶穌說.
瑪利亞你去.
你去做什麼.
瑪利亞不是說.
耶穌不是說瑪利亞.
你叫那些藍的門徒來.
說清楚發生什麼事.
他說瑪利亞你儘管去.
因為是一個委任.
一個差派.
差派他去哪裡.
去報信 報什麼信.
就是上帝這個創造的計劃.
一個復活的計劃.
是一定會完結的.
在哪裡展開了.
在耶穌基督的復活展開.
沒大拉的瑪利亞就向門徒報信.
我已經看見了主.
他又把主對他講的話.
去告訴他們.
願意讓我的生命成為一個復活的生命.
我們就去.
回去你的小角色.
回去我的小角色.
做什麼.
我已經看見了主.
怎樣能夠讓人看到主是復活.
人看到我的生命.
起碼看到多少.
上帝的兒子身體復活.

$^{921}$這個真相.
我們看看在哪裡.
主多謝你讓我們盼望.
多謝讓我們在這個人世之間.
我們雖然有軟弱.
有破爛.
這個世界被扭曲.
但是我們卻看見.
你自己救贖的工作.
是剛剛展開.
在我們的主耶穌基督.
從墳墓裡面出來的那一刻.
展開這個新的創造.
主我們仰望你.
因為你將來會來到的新天新地.
已經展開.
主願意這個新的生命.
這個新天新地.
這個真實.
在我們生命裡面展開出來.
主教導我們今天.
你拆散我們.
去回我們的工作.
去回我們的角色.
大的角色也好.
小角色也好.
讓我們能夠讓人看到復活的主.
復活的大能.
賜給我們勇氣.
賜給我們盼望.
祈禱逢轉命球.
\newpage



\section{}
\label{sec:SK0VEvCIu6I}
\textbf{比對研究 12: 巴力宗教 5 - 陳崇基牧師}
\newline
\newline
連結: \href{https://youtube.com/watch?v=SK0VEvCIu6I}{\texttt{https://youtube.com/watch?v=SK0VEvCIu6I}} ~~~~ 語音日期: 2023-04-22
\newline
\newline
\hyperref[sec:PhsM03F2OIM]{\small{< < < PREV SERMON < < <}}
~
\hyperref[sec:index_chronic]{\small{[返順時目]}}
~
\hyperref[sec:index_scriptual]{\small{[返順卷目]}}
~
\hyperref[sec:9WOLJvd45fU]{\small{> > > NEXT SERMON > > >}}
\newline
\newline
$^{1}$我們現在看以利亞背景.
上一次我們提到以利亞背景.
我們需要了解當時最大的假神.
最大的宗教,文化.
對當時的政局的影響有多大.
最大的宗教,影響力.
就是巴力這個假神.
在巴力的宗教背景下.
我們了解到為何聖經會這樣記載.
以及以利亞會遇見不同的情況.
以利亞有很多不同的故事.
包括之後的以利沙的故事.
為何會有這樣的神蹟故事.
有這樣的主題在裡面.
我們需要先了解背景.
我們提過巴力的神話.
巴力循環.
我們提過巴力當時最大的假神.
其中一個的綽號是駕駛者.
就是Rider of the Clouds.
他掌控著帶來雨水的雨雲.
所以叫駕駛者.
他也操控當時的雨水.
所以當時人們一定要去拜巴力.
因為拜巴力的時候就得到巴力的權能.
能夠有巴力的祝福.
能夠巴力降雨給他們.
供應他們的需要等等.
同時也控制了金路.
金路在巴勒斯坦的氣候是很重要的.
最難熬的幾個月.
寒季夏天的時候.
就是需要有金路.
Morning Dew.
以致農作物能夠得到收割.
能夠成熟能夠收割.
他也操控著生與死的權能.
巴力的故事.
起碼巴力的故事的循環.
每年雨季的時候.

$^{41}$就等於他復生10月至4月.
每年到寒季的時候.
沒有雨水的時候.
他就死亡.
但是他死亡只是幾個月.
之後他再打敗死神.
然後他再復生.
10月來到的時候.
雨季又會到.
就代表巴力又會再復活.
所以這個也是我們需要了解的.
我們提過這幾點.
巴力最重要的權能的彰顯.
就是有雨水有金路.
掌控生死等等.
你這樣想的時候.
在這個背景下.
你就想一下.
以利亞的故事.
是不是故意說一些.
或者觸碰一些這樣的主題.
是不是刻意將這些故事.
經歷來鋪排.
以利亞所經歷的.
他作為一個神人.
作為先知.
他有的那些遇見的情況.
那些神跡的故事.
是不是跟這些事.
都有關係呢.
所以為什麼我們需要了解.
巴力的背景.
我們上次17章開始.
其實17章.
開始的時候.
只不過是以利亞突然間出現.
之前到底發生什麼事呢.
16章 列王記嘗.
16章 這裡就交代.
阿哈王登基的時候.

$^{81}$做以色列王.
說些什麼呢.
跟17章很有關係.
17章其實就是一種回應.
是什麼回應呢 我們先看看.
他說.
猶大王 即是南面的王.
阿薩第38年.
暗利的兒子阿哈登基.
作以色列王.
阿哈在薩瑪利亞.
作以色列王有22年.
在薩瑪利亞.
就是北國的首都.
在哪裡呢.
就在這裡.
我們記得.
看過一些地理.
就告訴我們.
北國在中央山地的上面.
北面.
南國猶大就在耶路撒冷.
作為京都的.
猶大的山地.
所以在北國.
的歷史裡面.
由耶路撒冷開始.
到阿哈.
搬過幾次首都.
最後第三次.
搬到薩瑪利亞.
在薩瑪利亞.
他就作王了.
然後就交代.
他的政績.
這些都是列王記的.
格式.
這些所謂的.
Royal Record.
就是一些.

$^{121}$國王.
或者.
皇室的.
紀錄.
暗利的兒子阿哈.
行爺華.
比以前的所有的王更加嚴重.
他就犯了.
利伯的兒子.
耶羅伯安.
就是第一個北國的王.
所犯的罪.
患當作是小事.
就娶了西頓王.
海巴力的女兒.
夜死別.
為妻去事奉巴力.
敬拜他.
不單止這樣.
就建巴力廟.
在廟裡面.
為巴力去祝壇.
阿哈也做.
亞瑟拉.
你記得我們在背景上講過.
在巴力神話裡面.
他是母神.
就是母神.
他所做的就是.
耶和華以色列的上帝法老.
比他以前所有的以色列王更甚.
阿哈的.
兒子伯特利人.
希伊娜.
從修.
耶利哥.
納根基的時候就喪了長子亞比蘭.
安門的時候就喪了.
幼子西閣.
正如耶和華直至亂的兒子.

$^{161}$約書亞所說的話.
在這裡就記錄了.
用一個很典型的格式.
每一個王.
記錄了他什麼時候作王.
然後就作一個評估.
就是說這個王的.
他的政績.
他的政策等等.
很多時候作者就會說.
這個王.
犯了什麼錯.
所以為什麼到最後.
那個國家就覆亡.
所以這個就是.
歷史書.
或者列王記有的.
特別的一個format.
或者按照猶太人的.
聖經.
猶太人的Torah.
他就沒有歷史書.
這一部分.
所以我們今天叫歷史書.
他們叫做先知書.
列王記上下.
就是他們所謂的先知書.
就是前先知書.
我們今天叫的先知書.
以賽亞 耶利米 俄西亞.
等等.
那些叫做後先知書.
為什麼叫做先知書.
Prophetic Books 明明是講歷史.
因為這些歷史的記載.
不是純粹記載一些事實.
而是去評估.
和作一些先知的宣告.
因為先知就是.
站在神.

$^{201}$與神的子民中間.
他很多時候就是做那個.
prosecutor.
就是那個檢控官.
就是神和以色列人中間.
是有約的.
每逢這些以色列人.
走在正路上.
才會被.
宣告.
但更多時候.
神的子民走來.
猜答錯的時候.
先知就好像那個檢控官.
拿著上帝的約書.
去到神的子民面前.
來指責他們.
希望他們能夠悔改.
來迴轉.
所以我們每個人都需要先知.
在我們心靈裡面需要先知.
國家也需要先知.
社會裡面需要有先知的聲音.
等等.
所以這個就是一個很典型的.
其實是一個先知的.
一種評估.
在約的.
關係下.
作出一些評價.
阿哈王比其他王更甚.
因為這個緣故.
所以在.
申命記經常說到.
約書的時候.
上帝與以色列人立約.
與神的子民立約.
神的子民.
走得正站得正的時候.
就有祝福.

$^{241}$但如果他們走猜答錯.
特別是敬拜別神的時候.
就會有咒詛.
咒詛在申命記最後幾章.
講到咒詛的時候.
其中一種咒詛.
就是沒有雨水.
就是上帝不會再按照他的恩典.
按照他的信實.
降雨在地上.
因為申命記開頭也說到.
這個迦南地.
其實是一個.
靠著天上來的雨水.
來維生.
上帝說.
如果你是忠於我所約的話.
我就會降雨.
如果你不忠於我的約.
我就會將雨水收起.
因為上帝才是操控雨水的那個.
這個情況.
在阿哈裡面.
就更甚了.
因為他都事奉巴力.
又將夜洗別取過來.
夜洗別從哪裡來的呢.
西頓王的女兒.
她是西頓的公主.
西頓就在這裡.
之前我們講過.
如何知道巴力的宗教.
是講什麼的.
背景是從哪裡來的.
就是在這些海灣.
一連串的有心的海灣.
我們叫做菲尼基拉.
那裡有很多很出名的海港.
其中一個海港.
我們知道有巴力的神話出土.

$^{281}$就是烏加列.
這個地方.
烏加列.
就是在菲尼基地最北面的一個大海港.
下面南面.
就是我們經常在舊約聖經裡面會看到的.
兩個海港.
就是西頓和推羅.
在以利亞的故事裡面.
也有提到.
他在阿哈.
在這裡.
以色列的國度.
和京度.
然後他就和.
西頓的公主.
去成親.
在舊約近東的背景裡面.
成親.
或者.
娶過.
其他的公主過來.
有些是我們華人.
以前的西施.
等等的故事.
娶了過來的原因.
主要是外交.
和經濟.
所以.
如果你看地圖.
北角.
以色列在這裡.
他們就需要.
如果要將.
更多的收入.
更多的錢.
能夠去分享的話.
他們就需要.
北角就需要上面的海灣的錢.
所以和他們結親.

$^{321}$是很自然的事.
非常自然.
你要發展.
你要擴張你的以色列.
事實上阿哈就是.
列王裡面.
擴張的.
國度最多.
最大的一個國王.
他和他父王.
暗利.
特別是暗利也是很出名的.
所以在京外文獻.
都有提過暗利和阿哈.
特別是亞述的文獻.
為什麼他們這麼富強.
因為他們就要靠.
這些海港.
在以色列這邊沒有什麼大的海港.
沒有天然深的海港.
上面深的海港就是在菲利基.
所以他們就是.
菲利基的人或者在海港的人.
就是去操控著.
操控著.
東地中海的.
所有的貿易.
既然沒有海灣.
自然就需要去結親.
就要和西頓結親.
當然不同的時代.
有不同的政策.
在這個時代的政策似乎就是.
你和他結親的時候就要將他的宗教.
帶進來.
所以他就.
很出力的將巴力的宗教.
西頓啊.
推羅啊.
菲利基的這些假神.

$^{361}$帶進來以色列國.
\newpage



\section{6:33-34}
\label{sec:9WOLJvd45fU}
\textbf{First Things First: Matthew 6\_33-34 - John Yi}
\newline
\newline
連結: \href{https://youtube.com/watch?v=9WOLJvd45fU}{\texttt{https://youtube.com/watch?v=9WOLJvd45fU}} ~~~~ 語音日期: 2023-04-24
\newline
\newline
\hyperref[sec:SK0VEvCIu6I]{\small{< < < PREV SERMON < < <}}
~
\hyperref[sec:index_chronic]{\small{[返順時目]}}
~
\hyperref[sec:index_scriptual]{\small{[返順卷目]}}
~
\hyperref[sec:CsnEAjBlXuk]{\small{> > > NEXT SERMON > > >}}
\newline
\newline
$^{1}$(audience applauding).
- Thank you so much for having me back..
It's good to see everybody..
And actually the group is larger.
than the last time I was here..
I didn't look back behind me..
We know when I was sitting in the front row.
to see what was happening behind me..
But one thing I noticed that's kind of interesting.
about churches is that the most desirable seat in a church.
is usually in the back row..
I don't know why that is..
And the reason I think I was thinking about that is,.
you know, this morning when I saw Rachel,.
she was wearing this Cornelius Street shirt.
and I don't know anything about Cornelius Street..
She shared the reference with me about Taylor Swift..
And the sermon's not about her, so don't worry..
But I thought about, you know, when I was younger,.
I used to go to a lot of rock concerts too..
And I know it's still the case.
that the most desirable seat, the most expensive seat,.
the best seat in the house is in the front row, right?.
You wanna be as close to the action as possible..
You don't wanna have to look at.
what's happening through binoculars..
I did that once at a theater, at a live theater,.
and I was like, this is the worst experience ever..
You know, I'll never do this again..
If I can't get a good seat,.
I'm not gonna go to the theater..
So I would encourage you,.
if you're used to being in the last row,.
check out the action in the front row, okay?.
At least try it, okay?.
So let's go to God's word first..
And the passage we're looking at.
is a very familiar passage..
And it's in Matthew 6, chapter 33 to 34..
And this is what it says..

$^{41}$But seek first his kingdom and his righteousness.
and all these things will be given to you as well..
Therefore, do not worry about tomorrow,.
for tomorrow will worry about itself..
Each day has enough trouble of its own..
So I wanna thank you for coming together as the UV family..
You know, there's a lot of things.
you could probably be doing this afternoon,.
but you are here..
And I believe that you are at the right place.
at the right time, and that God has brought you here.
because not only is the Lord here.
to receive worship from you,.
but he has something he wants to give you..
He has something he wants to give you..
And so I hope we'll approach our hour or so together.
with this attitude of seeking Lord,.
what is it you have for me today?.
So let's pray..
So we do ask you, Lord, what is it you have for us today?.
We have tried to bring you our worship,.
bring you our very best, even in our tithes and offerings,.
in the way that we serve one another.
and serve our community, we do it for you..
We are bringing our best to you..
But Lord, thank you that you are a God.
who is open-handed, that you're not stingy,.
but you are generous..
In fact, we say you show us amazing grace.
because you give us better than we deserve..
And so we ask you, what do you have for us today, Lord?.
Not only are we your servants, but we are your children.
before a wonderful heavenly father..
And we are ready to receive from you..
Open our ears, our hearts are open too..
And we pray that you will give to us.
so that we can bring you even more glory..
We pray in Jesus' name, amen..
So you know, friends, it's no secret,.
we've been living in a age, a time, a season of life.

$^{81}$that is so filled with problems, so many stresses,.
so many troubles, worries, so many things.
that can cause fear..
And you know, I'm not gonna tire you with the list.
of all the things that we're sharing together..
And as much as the season has all these problems.
and troubles, all these things that might cause us.
to worry and be afraid, that doesn't even count..
Like, what are you experiencing?.
What are your unique sets of problems?.
Problems, needs, desires?.
What's stressing you out?.
What's causing you to lose sleep?.
What's causing you to wring your hands.
and wonder what's next?.
And I wanna tell you, Jesus has an answer for you..
It's in verse 34..
And it says, Jesus said,.
"Therefore do not worry about tomorrow,.
"for tomorrow will worry about itself..
"Each day has trouble of its own.".
Doesn't that bless you?.
That's enough, right?.
We can just pack it up and go home now..
'Cause Jesus said, "Don't worry.".
You know, by itself, this verse actually sounds pretty cold,.
brutally cold..
Like, Jesus is, it sounds like he's saying, like,.
"Come on, just get over it..
"Don't worry.".
But see, it's not a verse that just stands by itself.
or a sentence that stands by itself..
It is in the context of something bigger..
And like any good Bible student,.
we should be concerned about context, right?.
We already know how dangerous it is.
to hear a piece of news out of context.
and think, "Oh, I understand what's going on.".
But we may have no idea what's actually going on,.
what the person actually meant when they said it, right?.

$^{121}$But in context, we get the picture.
of what is the person saying?.
So we don't wanna just go with a soundbite from Jesus..
We wanna hear what is his message for us today?.
So we gotta look at the verse before..
And it says, "But seek first the kingdom of God,.
"or seek first his kingdom and his righteousness.
"and all these things will be given to you as well.".
So let's kind of unpack this going backwards.
as we've already done, starting with verse 34,.
and then go to verse 33,.
which is really gonna be our key text today..
With all that we're experiencing,.
verse 33, Jesus tries to center us..
He tries to prioritize us..
And even though he starts by saying,.
"Seek first his kingdom and his righteousness,.
"I know everybody's attention went to,.
"and all these things will be added to you," right?.
So let's get that out of the way.
so that you're not distracted by the whole time.
wondering what are all these things.
that he's gonna give me, right?.
So let's go through it..
And again, in context,.
we are in the midst of Jesus' Sermon on the Mount..
It began in chapter five..
And in the beginning of chapter five,.
it starts with Jesus promising.
that God our Father has a response..
He's not sitting idly and quietly.
and hiding in the shadows,.
but he has a response for those who are poor,.
those who mourn, those who are gentle,.
who seek righteousness or justice,.
who need mercy, who desire to be pure in heart,.
who desire peace, who are persecuted..
God has a response..
God has something for us..
We find that in the Beatitudes, as we call them..

$^{161}$And then there's this big section.
that covers from chapter five through chapter six,.
and it has to do with our spirituality,.
our relationship with God..
And so he covers the law,.
spiritual disciplines like piety,.
generosity, fasting, and prayer..
He talks about how should our attitude be about wealth..
So spiritual disciplines..
And then finally, he covers basic necessities..
Basic necessities..
Concerns about food, water, clothing,.
survival, a longer life..
And are any of these of concern for you?.
Are any of these things that keep you up at night?.
You know, most of us in America,.
you know, we don't miss any meals..
We're not, unless, I guess, unless we're on a diet.
and we choose to miss a meal..
Nobody is really going thirsty..
You know, water is so readily available, fresh, clean water..
And then nobody's walking around naked.
unless they decide to do that, right?.
It's not for lack of clothing..
And yet there are people all over the world.
who are struggling with these concerns every day..
And not to mention all the other ones,.
starting from the beginning of chapter five..
And I just did a quick Google search.
and I just asked the question,.
is food insecurity better or worse?.
Is it getting better or worse today?.
And I was surprised with all the attention to ending hunger,.
that according to Google,.
food insecurity is getting worse in the world..
And then I thought, you know, I'll ask one more question..
Is abject poverty, that's extreme poverty,.
is it getting better or worse in the world?.
And I was surprised again..
It's getting worse..

$^{201}$Approximately 10\% of the world is hungry..
About 10\% of the world is facing abject poverty..
Wow, this is a word from Jesus..
And I think this speaks not only to us.
and all the needs and worries we have, but to the world..
So let's be ready to receive this..
So Jesus says in chapter six, verse 25 to 32,.
he says, don't worry about your life..
You know, I should ask, how many people remember the last,.
who was here the last time I came to visit?.
Any hands?.
Only Pastor Mark, everybody else?.
Okay, just a few people..
So I wanna ask our worship leader,.
do you remember what I preached on last time?.
Okay, I'm just kidding..
I had to look it up.
'cause I couldn't remember what I preached on..
I just, I look it up because I wanna make sure.
I don't preach the same message again.
and just because that sometimes happens, right?.
But I preached on a text that was earlier,.
right before this text..
And I didn't intentionally choose this text today,.
thinking, oh, you know,.
because I preached on the section before,.
it'd be natural to preach on..
It just, this is something that the Lord.
has just laid on my heart..
And, but the last time I preached,.
I talked about how our heavenly Father.
knows our needs before we even ask..
That was basically the gist of the message..
And that Jesus is saying,.
don't worry about your life and what you need..
Don't be like the Gentiles who are fixated.
and consumed with the pursuit of these things..
All these things that cause this universal worry.
and stress on our lives..
Jesus was speaking to a largely Jewish audience..

$^{241}$So it would have been,.
he's saying, don't be like the Gentiles,.
like the rest of the nations or the pagans,.
the idol worshipers..
Don't be like them who worship dead gods..
Today, if he was saying it,.
it might sound more like,.
you are worshipers of a living God..
So behave like people who have a living God,.
who actually hears what you need,.
but also already knows what you need..
And because he loves you and he is gracious,.
he's already at work..
He's already on the job,.
taking care of what we need..
God's got this..
Know with certainty,.
God loves you and he's taking care of you..
So don't be distracted by these things..
Don't be distracted by concerns for things.
that our loving and capable father.
can already take care of..
Instead, seek first the kingdom of God..
Seek first his kingdom and his righteousness.
and all these other things will be taken care of,.
will be provided for you..
So what does it mean then to seek something first?.
I mean, this might seem really like kind of pedantic,.
you're like, why say it?.
But I think it needs to be said.
'cause it's in the text, right?.
What does it mean to seek something first?.
It has to do with, it speaks with,.
what is your greatest passion and desire,.
your most important, highest priority pursuits?.
What are you intentionally chasing after.
with the most determination and fixation?.
Think about that..
What are you chasing after?.
With everything you've got,.

$^{281}$in your heart you've said to yourself,.
I must do this, I must have this,.
this is worthy of whatever it takes..
What is that thing?.
And so many of us, especially we got a young audience here,.
we're thinking about our education and our careers,.
finding the love of our lives,.
building a family, paying the bills..
Maybe it goes a little farther.
and you're consumed with, I must be wealthy one day..
I must have this kind of house or this kind of watch..
I must have these things..
Is that what consumes your whatever it takes?.
I must have this, I must do this..
Jesus says, seek first his kingdom and his righteousness.
so that everything else will be provided for you..
And so what is the kingdom?.
The Bible and the New Testament talks about the kingdom.
in so many ways, Jesus talks about it often in parables..
And he says, the kingdom is like yeast..
It permeates and it penetrates..
And it's not something you can simply point at.
and say like, there it is..
Or look, here it is..
Because the kingdom of God is also in our midst..
It's like the air that we breathe..
It's present even when we don't recognize it,.
it's present even when we refuse to see it or acknowledge it.
even when we insist on denying it,.
the kingdom of God is there..
We behold the miracles of God and the kingdom is there..
The kingdom of God grows and bears fruit.
and the kingdom bears the fruit with or without us..
With or without us..
The kingdom might seem small to some.
and so hard to notice,.
but the kingdom of God is also so big,.
so much bigger than anything we could imagine..
You know that we're told the kingdom of God.
is evidenced by the side effects..

$^{321}$When we see certain characteristics,.
we know that the kingdom of God is there,.
just like when we see the wind blow,.
we know the wind is there..
Or when we see the leaves move, we know the wind is there..
Or when we see certain signs in the weather,.
we know that seasons are changing..
And then we're also told we have to receive it like a child.
and unless someone is born again,.
they cannot see the kingdom of God..
Jesus, John the Baptist and the apostles.
spent a lot of time talking about the kingdom of God..
So what is the kingdom of God?.
I like easy..
So I'm gonna boil it down,.
maybe as much as that you can possibly boil something down..
Somebody once said, this isn't mine,.
but when I heard it, I said, I'm gonna use that.
because that makes sense to me.
'cause I'm not smart, that's easy..
You know what it is?.
They said, whenever you see God get his way,.
you're seeing the kingdom of God..
Did you get that?.
Whenever God gets his way,.
it's the kingdom of God you're seeing..
So when you make a decision, I'm not gonna do it my way,.
but I'm gonna do it God's way,.
that's the kingdom of God at work..
When you see somebody choose God and not something else,.
we're seeing the kingdom of God at play..
When the king gets his way, we see his kingdom..
There are literally thousands, no, millions.
of books, articles, sermons, lessons.
that have been taught on this and written on this,.
but it's when God gets his way..
Are you seeking first that God would get his way.
in your life?.
'Cause that's where it starts..
God, I am surrendering right now..

$^{361}$I'm not gonna hang on to anything,.
I'm letting it all go and saying, have your way in me..
That's hard..
It's hard to let it all go and say,.
I'm gonna seek your kingdom first..
And then, are you praying that people around you,.
your friends, your family, your boss, your neighbors,.
are you praying that the people around you,.
your church, would have God's way,.
that they would choose God?.
I mean, isn't that what we're praying.
when we're praying for our brother.
who doesn't know Jesus yet?.
God, would you save him?.
And what is that?.
Have your way, God, with my brother..
And then how about the rest of the world?.
When we see wars and rumors of war.
and we're praying for what's happening here or there,.
when we're hearing about people who are hungry.
or naked or who are exposed and vulnerable.
and we're praying for all the things.
that are going on in the world,.
is that not God?.
Would you please have your way?.
Do what only you can do..
That's the kingdom of God..
So seek first the kingdom of God and his righteousness..
And for most of chapter five, Jesus talks about the law..
The law, the rules which make up.
the measuring stick of righteousness..
And the rules which govern people's lives.
and their relationship and their standing with God..
So you know the Jews, they lived and breathed the law..
They lived and breathed the law..
They were fixated on it, they were just engrossed in it..
They were obsessed with the law of God..
613 of them in the Old Testament..
Doing everything possible to not violate any laws,.
to obey and observe the laws..

$^{401}$And the law was also a heavy burden.
because it affected everything in their lives..
It had something to say about every aspect of their lives..
And so I don't know about you, when you think of the law,.
even the law as you understand it in the Bible,.
do you feel like, man, you know,.
I'm pretty good at this obedience thing..
Tap yourself in the back..
Good job, John..
Or maybe you're the kind of beats yourself up every day..
And you're counting all the ways that you've failed,.
all the ways you've disappointed God,.
you've broken the law..
Jesus has a bomb to drop on us..
I wanna say it's like a cluster bomb..
'Cause there's at least three things,.
three explosions that are gonna go off in our heads,.
in our hearts..
And the first one is this..
That when it concerns the law, Jesus says,.
"Do not presume," and he says this back.
in chapter five, verse 17,.
"Do not presume I came to abolish the law and the prophets.".
So it's a heavy burden..
Jesus came to save us..
Take off our heaviest burden, Lord..
Take the law, would you just eradicate the law?.
If we don't have the law,.
then man, I wouldn't feel guilty anymore..
I wouldn't have to constantly be measuring myself..
But guess what?.
Jesus says, "I didn't come to get rid of the law..
"I came, didn't come to abolish the law.".
Then he gives us bomb number two..
And he says, "I came to fulfill the law..
"There is no one else who has made such an audacious claim..
"I am here to completely fulfill this law..
"Obey it perfectly, observe it perfectly.".
So don't even try to mess with the law..
This law, which is the air and the atmosphere.

$^{441}$of the Jewish experience, the religious leaders,.
the teachers of the law,.
they led and they maintained the standard..
They knew God was holy..
And they also knew, they did know.
that the law was virtually impossible to keep..
So you could say they would even mess with God's law.
by creating more rules and laws around the law..
So kind of figuring, you know,.
if mom says, "Don't touch this little screen here,".
then what do you do?.
You put a fence around it..
So there's no way you can reach that..
There's no way you could violate that rule..
So they'd create these fences, these hedges around the law..
The religious leaders interpreted the law for everybody..
The teachers drew the boundaries around it..
The judges determined how to interpret it..
The priests set the bar, the standard of obedience.
and what righteousness looks like..
And even though we tend to remember them.
for their hypocrisy, I submit to you.
that better than anybody else,.
they tried to obey the law..
They were the best law keepers..
That is until Jesus came..
As far as human effort goes,.
nobody was more successful when it came to the law..
But here is the biggest bomb of all..
Because we do look up to certain people, right?.
Like, do you look up to Pastor Mark?.
And, you know, when it comes to,.
maybe you're like, you know,.
I don't wanna be wealthy like Pastor Mark,.
but I do wanna be spiritual and godly like Pastor Mark..
Maybe that's, you know, how you look at it..
I wish I dressed better like Pastor Mark does..
See, we hold people up as like,.
gosh, that's something to aim for..
I wanna be like that..

$^{481}$And the Jewish people basically held up the priests,.
the teachers of the law, the Pharisees and the Sadducees.
as these are the people we're gonna emulate..
And what does Jesus say?.
In verse 20 of chapter five,.
"For I say to you, unless your righteousness surpasses.
"that of the scribes and Pharisees,.
"the experts in the law,.
"you will not enter the kingdom of heaven.".
Did you hear that?.
Unless you are more righteous than Mark Chan,.
you will not enter the kingdom of God..
Unless you are more righteous than Billy Graham,.
you will not enter the kingdom of God..
You can put anybody's name on that list,.
any great Pharisee, any great teacher of the law,.
and if your righteousness doesn't surpass that, you've failed..
Not that if you will just meet it,.
but even if you can't exceed it,.
if you can't surpass them..
So, you know, one of the most righteous Jews.
in the New Testament is the Pharisee, Saul..
Later, after an encounter with Jesus,.
his name changed to Paul,.
and again, one of the most righteous men on the planet.
says this about his former self..
And it's in Philippians chapter four..
I don't know if you wanna turn there,.
but I'll read it to you,.
and you know, maybe that'll make it easier to even listen..
He says in verse four,.
"If someone else thinks they have reasons.
"to put confidence in the flesh,.
"that's human ability for righteousness,.
"I have more.".
I have more confidence..
"I'm circumcised on the eighth day of the people of Israel,.
"of the tribe of Benjamin, a Hebrew of Hebrews,.
"and so far,".
do you notice everything on that list.

$^{521}$is nothing that he had any control over?.
Did you notice that?.
Circumcised on the eighth day..
So, I mean, like no baby says,.
"Dad, we got an appointment, eight days,.
"time for circumcision.".
And he had no choice..
There was nobody who said,.
"You know what?.
"Let's make sure that Saul is born to Israel.".
He didn't control that either..
That he's born in the tribe of Benjamin..
He didn't choose that..
Hebrew of Hebrews, he didn't choose that..
But then in regards to the law, he says,.
"I'm a Pharisee.".
As for zeal, that means passion..
I persecuted the church..
As for righteousness, based on the law, I am faultless..
That's what Saul, Paul says about his former self..
And yet, what does he conclude about this old self.
that looked so enviable to everybody else around him?.
Like, "I wanna be Saul.".
He says, "But whatever were gains to me,.
"I now consider them loss for the sake of Christ..
"What is more, I consider everything a loss.
"because of the surpassing worth.
"of knowing Christ, Jesus my Lord,.
"for whose sake I have lost all things..
"I consider them garbage,.
"that I may gain Christ and be found in him,.
"not having a righteousness of my own.
"that comes from the law,.
"but that which is through faith in Jesus Christ.".
Not the law, but a righteousness that comes.
through faith in Jesus Christ..
The righteousness that comes from Jesus Christ,.
the righteousness that comes from God.
on the basis of faith..
I want to know Christ, yes,.

$^{561}$to know the power of his resurrection.
and participation in his sufferings,.
becoming like him in his death,.
and so somehow, attaining to the resurrection from the dead..
Paul says, "I'm not gonna pursue the righteousness.
"of even the holiest of men.".
I'm not gonna accept that as the standard..
Why?.
Because that is not God's standard..
I will pursue God's standard of righteousness..
And the standard of God's righteousness is Jesus himself..
He says, "I'm not gonna put my faith and trust in garbage,.
"even though it's shiny and pretty..
"I'm gonna choose Jesus.".
Seek first his kingdom and his righteousness..
Some of us might be thinking,.
man, you know, poke your elbow at the person next to you.
and say, "I'm better than you..
"I'm more righteous than you.".
That's not good enough..
You're selling yourself short..
The righteousness you've set for yourself,.
whatever standard that is, if it's not Christ, it's garbage..
Simple as that, it's garbage..
So get ready for a little heresy, okay?.
So if you've got stones, please try not to throw them..
And let me just say, I don't, this is not heresy,.
but you might think it is until you let me finish, okay?.
So please let me finish..
So you know, if you've been going to church.
any number of times, if this is the second time.
you've attended Urban Voice Church,.
then you will have already heard.
what is the only way to eternal life?.
What is the only way for salvation?.
What is the only way to God?.
Anyone?.
(audience laughing).
I hope you're just being shy, but it's faith in Jesus Christ.
that's as simple as it gets, right?.

$^{601}$It's trusting Jesus, believing Him..
John 3:16, "For God so loved the world.
"that He gave His only Son,.
"that whosoever shall believe in Him.
"will have eternal life.".
Well, they will not perish, but have eternal life, right?.
So I'm gonna submit to you that there's one way.
one other way, one other way..
I mean, I affirm Jesus is the way to the Heavenly Father,.
but I'm gonna say there's one other way..
And guess who tells us what the other way is?.
Jesus..
So Matthew 19, verse 16 through 22,.
you might wanna turn there just to say like,.
man, I'm gonna prove John wrong today..
There is no other way but Jesus..
And then someone, you might know this story, it says,.
"And someone came to Him and said,.
"Teacher, what good thing shall I do.
"that I may obtain eternal life?".
I.e., or in other words, salvation..
That I may have heaven..
And He said to him, that's Jesus said to him,.
"Why are you asking me about what is good?.
"There's only one who is good,.
"but if you wanna enter life, keep the commandments.".
That would have been the perfect time for Jesus.
to put another exclamation point.
on salvation by faith alone through Christ alone, right?.
He could have said, "Trust me, believe me,.
"put your faith in me.".
But He doesn't say that..
He says, "Obey all the commandments.".
Now, what a normal person would have said,.
well, then, I've already violated,.
even before this conversation, I've already violated.
all the commandments..
And what we should say is, okay, is there another way?.
That's what should have happened, but in this story,.
the man who's asking the question,.

$^{641}$when told just obey all the commandments, what does he say?.
Okay, I did that..
I've obeyed them all since I was a kid..
And you can imagine anybody who's standing around listening,.
they're like, "Oh my gosh, this guy is so unaware of himself.
"has no idea..
"He is living in a fantasy world, he's delusional.".
But it says, "Jesus loved him, though.".
And then he said, "You got one more thing.
"that you're lacking..
"So sell everything you have and give it to the poor.
"and then follow me.".
And there was the answer right there, right?.
Basically, it should have begged the question,.
Jesus, what other way is there?.
Because this is an impossible way,.
the law is an impossible way..
And then Jesus would have just said, "Then follow me.".
And that's what he says to us..
There is another way, but it leads to death.
because the law won't save you there,.
will only accuse you, will only confirm.
you're a sinner, that you deserve the judgment.
and wrath of God, but instead choose me..
Hang on my coattails and faith.
and let's walk into glory together..
So Paul, again, arguably the most righteous man.
walking on the face of the earth, chooses Jesus..
Arguably the most successful law keeper,.
perhaps in the history of the world,.
says, "I will not choose that garbage..
"I will not choose what looks like my success story..
"I'm gonna choose Jesus.".
And brothers and sisters,.
Jesus says, "Seek first the kingdom,.
"his kingdom and his righteousness..
"Choose him, choose Jesus,.
"God's standard of righteousness.".
And then we can say, "Therefore," in verse 34,.
"all these things that consume our thoughts,.

$^{681}$"that distract us from pursuing.
"and keeping the main thing, the main thing,".
Jesus says, "Don't worry..
"Don't worry about tomorrow,.
"for tomorrow will worry about itself..
"Each day has enough trouble of its own.".
You don't need to worry about those things.
because you're seeking first his kingdom.
and his righteousness and all these other things.
that try to distract us, that try to consume our thoughts..
"I've got it, I'll take care of it..
"Don't worry.".
So friends, let's put our confidence.
in Christ and his kingdom..
Let's seize the righteousness that Jesus is,.
the righteousness that God offers to us as a gift..
Let's receive him in faith..
And if you have chosen any other way.
to gain the approval of God, to win your ticket to heaven,.
then let's throw that in the garbage can where it belongs..
Paul did..
And let's choose Jesus..
Let's bow our heads..
I'm gonna give you a moment to respond to the Lord..
I'll start with that first question..
What is it that we pursue with all we have?.
What is our highest priority desire?.
Let's just confess it to the Lord..
(soft piano music).
And put that, and maybe there's more than one,.
maybe there's dozens of things..
And put that in one hand, and then in the other hand,.
the kingdom of God and his righteousness..
(soft piano music).
I'm just asking you, in your own heart,.
just weigh those two things out..
Jesus said, "What does it profit a man.
"to gain the whole world and yet lose his soul?".
(soft piano music).
And as the Lord leads you,.

$^{721}$would you empty one of those hands?.
Lay it at his feet..
And tell him what you're gonna do..
(soft piano music).
And you can pray alongside me,.
or you can just keep praying in your own words..
Father, we acknowledge right now that you are a holy God.
who has an impossible standard of what is holy.
and what is righteous..
And first, we just ask you to forgive us.
as we confess that we have missed the mark..
We have not lived up to your standard..
And even those days that we may have fooled ourselves.
into thinking that we are holy enough,.
good enough, righteous enough, Lord, it was a lie..
Forgive us for trying to earn our way,.
earn our way to your approval..
(soft piano music).
Forgive us that we have sought our kingdom and not yours..
We have chosen my way, not the Father's way..
And forgive us that we have not trusted you,.
we have not believed you,.
that we have not embraced your son, Jesus,.
that we have not clung to him desperately..
But today we put all of our desires,.
all of our pursuits, our worries, our fears, our concerns,.
and we drop them at your feet,.
the feet of Jesus who says,.
"Come, all who are weary and heavy laden,.
"and I will give you rest.".
The one who invites us to take the heavy burden.
off our shoulders and that you promise to trade with us.
a light yoke, a light burden, an easy burden..
Lord, we receive that..
Thank you for taking the load off of our shoulders..
And now as we live our lives,.
seeking that you get your way,.
and trusting in your son, Jesus, the righteousness of God,.
(gentle music).
Lord, let us experience the joy of following you,.

$^{761}$the joy of seeing you provide everything that we need,.
coming to you with open hands,.
holding nothing back from you,.
as you pour out on us the blessings of heaven..
We trust your promises..
We trust your word..
We trust your son..
Have your way with us..
God, would you help us?.
We already know we're not able..
We cannot obey your law..
No matter how hard we try, we keep failing..
Lord, would you help us?.
And would you put our sights on Jesus.
who has already fulfilled your law perfectly?.
And we're not gonna try to win our own salvation.
'cause Jesus has already done it..
Thank you, Lord Jesus, that you went to the cross.
perfect, spotless, sinless,.
and you became sin for us and crucified our sin,.
paid the debt that we owed for our own sin..
We love you and we are so thankful to you..
Lord, as we go from this place,.
I pray you put a bounce in everybody's step,.
knowing that there's a lighter load on their shoulders,.
that they are not walking hopelessly,.
but now we are walking with so much hope,.
so much promise,.
Lord, we are walking in the goodness and grace of God..
Thank you so much, we pray in Jesus' name, amen..
and I hope..
\newpage



\section{哥林多前書}
\label{sec:CsnEAjBlXuk}
\textbf{復活的福音:哥林多前書 15 章 - 陳崇基牧師}
\newline
\newline
連結: \href{https://youtube.com/watch?v=CsnEAjBlXuk}{\texttt{https://youtube.com/watch?v=CsnEAjBlXuk}} ~~~~ 語音日期: 2023-04-24
\newline
\newline
\hyperref[sec:9WOLJvd45fU]{\small{< < < PREV SERMON < < <}}
~
\hyperref[sec:index_chronic]{\small{[返順時目]}}
~
\hyperref[sec:index_scriptual]{\small{[返順卷目]}}
~
\hyperref[sec:wxXYrlF2bG8]{\small{> > > NEXT SERMON > > >}}
\newline
\newline
$^{1}$交比佢.
AZ等節目好.
我地繼續一個topic.
係講到復活.
復活節過jor .
不過仍然要講復活.
因為復活節或者復活ge 訊息.
復活呢件事.
身體復活.
基督ge 復活.
係我地信仰ge 一個核心.
唔係一個footnote.
好多時我地當jor 佢係一個footnote.
咁我地有好多ye 我地就搞唔清楚.
所以點解我地係呢幾個禮拜.
我地會去到討論下呢一個ge 課題.
諗一諗.
上帝係基督裡面ge 死而復活.
對於我地生命有D 乜ye ge 鼓勵.
對於我地生命裡面有D 乜ye ge 形容.
呢D 唔係純粹將來會發生ge ye .
因為將來個目的地.
我地去到新天新地.
我地有一個復活榮耀ge 身體.
好似基督自己有榮耀復活身體一樣.
係應該定決定我地今時今日.
點樣睇我地ge 經歷.
我地呢個人生.
走ge 方向一定要睇最尾gwo 個目的地.
搞清楚目的地.
我地而家就知道我地唔會咁容易走失.
所以呢度有三個point.
我呢就想同你黎到去思想一下.
第一個point就係復活ge 真實性.
第二個point就係復活ge 身體.
我地思想一下.
一個復活ge 身體.
係一回咩事.
點樣係一個復活ge 身體.
第三方面.

$^{41}$復活帶比我地有D 咩ge 得勝.
所以第一個point係復活ge 真實性.
呢度好明顯講到復活ge 時候.
其中一段最長ge 經文就係哥林多前書第十五章.
你知道哥林多前書由頭到尾.
有十六章聖經.
裡面處理好多教會ge 問題.
或者教會對保羅ge 提問.
所以保羅就寫信作為一個牧者.
黎到話比佢聽唔同ge 事情應該點樣處理.
譬如婚姻點樣處理.
單身ge 點樣處理.
性倫理點樣處理.
佢地中間有分黨或者有不和.
甚至有D 上到去法庭gwo 度互相去告狀等等.
呢D 全部都係哥林多前書.
前面ge 章節一路黎到去討論.
所以保羅係所謂哥林多前書.
其實係一個Q and A黎.
到最尾ge 時候佢似乎佢佈局.
就係提到.
或者討論到一個最重要ge ye .
最尾時候通常都係講D 最重要ge ye .
十五十六章.
最尾兩章.
十五章佢就整章長ge 篇幅黎到去討論呢個問題.
呢度其中一節十二節.
好清楚話比我地聽.
保羅可能聽到.
或者有D 人向佢詢問.
就係係哥林多教會裡面有一D 信徒.
佢地係唔信基督復活ge .
沒有死人復活的事.
你知道哥林多處於一個咩地方.
係希臘半島.
所以佢地好受當時世界裡面.
文化裡面ge 希臘思想影響.
其中一個ge 希臘思想就係柏拉圖主義.
柏拉圖係早前五六八年.
一個好出色ge 哲學家.

$^{81}$大家聽過啦.
蘇格拉底,柏拉圖同埋亞歷斯多德.
呢D 咁ge 希臘哲學有一個共通點.
特別柏拉圖.
柏拉圖就係話.
呢個世上我地捉到摸到ge ye .
包括我地物質上ge 肉身呢D ye .
都係唔會永恆ge .
呢D ye 甚至可以係邪惡ge .
係唔好ge ye 黎ge .
Metal.
Body.
唔係好ge ye .
Someday我地就會係一個監獄裡面.
被釋放出黎.
因為希臘人ge 思想.
就係靈魂.
係永遠ge .
靈魂係永遠會落去ge .
肉身呢.
就係唔好ge ye .
咁所以呢係.
福音去到哥林多ge 時候.
就對佢地有一D ge 衝擊.
呢個柏拉圖主義就係保羅去到想.
去改正或者去到.
挑戰ge 一個ge 思想.
呢個柏拉圖主義.
我想講呢D ye 唔係話D 好心ge ye 係咪.
以前你讀大學ge 時候讀Philosophy 101.
都講到呢D ye 係咪.
而家就係Philosophy 101.
呢柏拉圖主義之後.
演變成乜ye 呢.
就係洛斯底主義.
你話咩黎架.
以前我大學冇讀過.
確實你真係冇讀過.
呢個其實係初期教會.
面對一個最大ge .

$^{121}$第一個最大ge 異端.
就係洛斯底主義.
其實係點樣黎ge 呢.
就係從柏拉圖主義呢一個ge .
主流gwo 度分支出黎ge .
洛斯底主義就係話.
我地信jor 耶穌ge 時候.
我地都唔需要理呢D 咁ge 所謂bodyge ye .
呢D matterge ye .
因為我地最終.
真係得救ge 係點樣.
我地脫離呢一個.
取死ge 身體.
掉走佢.
係咪.
信仰呢就係脫苦海.
係咪.
脫離呢一個ge .
物質ge 環境.
呢個邪惡ge 世界等等.
然之後我地去到somewhere.
呢個就叫做.
救恩.
保羅就去抗衡呢樣ye .
如果睇下新約ge 書信.
有好多都係受呢D 咁ge .
思想ge 洪潮黎影響.
其實甚至今日都係一樣.
係咪.
今日係唔同ge .
一D 所謂cultural ideas.
係美國也好.
係唔同ge 地方也好.
係我地ge 文化裡面就話比我聽.
body.
有乜所謂.
我地咪亂黎.
係咪.
隨便去定義.
或者.

$^{161}$我地唔覺得呢樣ye .
係.
係俾我地identityge .
唔係我地身份ge 一部分ge .
我地可以隨意黎到去將佢定義.
呢D 其實都係.
其實好巧好得意.
所以我就.
同D 人講.
特別D 神學生.
點解要讀教會歷史.
雖然教會歷史好悶.
係咪.
有好多ye 其實係翻炒.
係當時翻炒.
你諗o下2000年.
都冇點變過.
係咪.
律師主義影響當時ge 教會.
亦都影響我地今日.
好多時唔小心ge 時候.
我地就會諗.
信耶穌係一回咩事.
lor 個護照上天堂咋嘛.
係咪.
我地係靈魂咁樣.
我地走jor .
跟住我地就坐係舊魂上面.
然後係度默想.
呢個就係叫救恩.
呢個唔係基督教ge 信仰.
所以點解保羅係十五章gwo 度.
就好著名ge .
好詳細ge .
黎到去交代.
甚至去同佢地作一D ge 爭辯.
話基督ge 復活.
係實在ge .
十五章第一節.
佢話大宣節目.

$^{201}$我希望你記住所傳比你地ge 福音.
所領袖ge 要堅持ge 福音.
你地如果持守呢個福音ge 話.
必因福音而最終能夠得救.
佢講緊ge 得救唔係就係靈魂出竅.
然後去到天堂就算數.
佢講緊我地ge 身.
與我地ge 心.
都會得救.
心靈與身體.
都會得救.
呢個就係基督死而復活ge 福音.
佢話我所領袖ge .
並傳比你地ge .
乜ye .
最重要ge 訊息.
保羅你諗下可以講好多教義.
可以講好多ge 道理.
可以教你點做好人.
點樣做善行等等.
但係佢話最重要.
the most importantge 訊息係乜ye .
基督照聖經記載.
為我地死了.
然後被埋葬了.
又照聖經ge 記載.
第三天復活了.
又向彼得顯現.
向十二使徒顯現.
後來一次向五百位ge 大兄姐妹顯現.
其中許多依舊仍然健在.
有些已經安息了.
隨後佢向亞國顯現.
向眾使徒顯現.
到最後.
就向我.
我就是一個未夠產期出世的.
即係一個早產ge 一個英雄.
點解呢.
因為其他跟過耶穌幾年的門徒.

$^{241}$幾年先能夠成為使徒.
三年.
而家我地十月懷胎.
其實三年懷胎.
三年先搞到一個門徒出來.
佢地就做使徒.
保羅沒有跟過耶穌.
所以佢就說.
我見到復活ge 基督.
我見證到福音.
但我又不像佢地跟過.
所以我就是未夠產期.
我就出世.
不過上帝給我恩典.
我能夠成為使徒.
其實呢.
我就不配.
因為我曾經迫害過教會.
保羅講了一大堆講什麼.
其實講出一個當時的信經.
當時初期教會.
有一些信經.
他們在崇拜裡面用.
或者他們在聚會裡面會用.
佢就將這個信經.
再一次講出來.
這個就是最重要的信息.
所以你想想上一個星期.
我不是講過話.
第一個見證耶穌復活的是什麼人.
莫大拉的什麼.
瑪麗亞.
還有誰呢.
還有另一個女人.
叫什麼名.
瑪麗亞.
還有另一個女人叫什麼名.
瑪麗亞.
是的.
沒錯.

$^{281}$為什麼在這裡不講.
因為他在這裡講的是一個公信力的信條.
當時的文化來講.
就是主要是男人.
做一個法庭上的見證人.
是不是.
所以這裡講的和四方講的.
就有一點點分別.
但是其實都是一樣.
就是復活的真實.
到他交代了這個之後.
他就去到第二個的要點.
你可能會問.
喂.
家裡是復活的.
是不是.
基督的復活.
然後代表我們身體會有復活.
那麼到底我們身體復活會怎麼樣.
其實這個也是哥林多教會人所問的問題.
保羅已經預先講了.
他說或許有人會問.
死人怎樣復活呢.
死了.
復活之後有什麼身體呢.
怎樣的身體.
想一想.
到底復活的身體是怎麼樣的.
喂.
聖經裡面.
基督復活之後.
有些少記載.
講到基督復活了榮耀的身體是怎麼樣的.
這是一個很有趣的問題.
是不是.
我相信不單止當時的人會問.
今時今日我們都可能會問.
這麼神奇.
這裡是大家想像力.
是不是.

$^{321}$我們有時候講到所謂的終末論的時候.
我們就會想.
到底復活的身體會是怎麼樣的.
我問你.
你想像一下.
到底將來復活的身體.
如果你是基督裡面.
復活的身體.
不是靈魂出去了.
然後就永遠永遠在天堂.
在上面浮浮浮.
新天新地來到的時候.
基督再來的時候.
我們就有一個榮耀復活的身體.
一個榮耀的身體.
那個身體應該是怎麼樣的呢.
你覺得會怎麼樣.
你覺得會有咒文嗎.
你今天有沒有咒文.
你想不想復活身體有咒文.
不想.
OK.
Alright.
我以前去以色列的時候.
經常去一些店鋪.
人們經常會介紹一種商品.
就是 Ahava.
就是以色列裡面最出名的護膚用品.
護膚用品.
有沒有用過.
聽說挺好用的.
OK.
我當然沒有用.
所以人們經常問我.
會不會有些人用了會好一點.
不知道.
你試一下.
可能日本那些好一點.
不過聽說那些死海產品都是很好的.
你想象一下.

$^{361}$那些護膚品.
其實就是在做一件什麼.
去減慢我們的 decay.
Right.
咒文肯定會有的.
咒文永遠都會繼續少還是繼續多.
有沒有人會繼續少的.
千年樹人.
一定會繼續多的.
殘缺的地方一定會繼續多的.
衰壞的東西一定會繼續多的.
其實那些護膚品就是將過程減慢而已.
Right.
是不是.
無論如何都會多的.
Right.
希望將它減慢.
復活身體沒有咒文是不是.
OK.
你說譬如.
我們到了這個年紀.
我們開始沒有頭髮.
你那個復活的身體會不會有頭髮.
會不會復活的時候突然間頭髮都彈出來.
那你以後就要帶一把梳子了.
煮菜的時候.
想象一下到底什麼是一個復活身體.
整段的聖經這裡.
三十五節開始.
保羅就是去到交代復活.
不是復甦.
不是 resuscitation.
大家記得有一樣東西叫做DNR.
是嗎.
Do not resuscitate.
Right.
就是到你很不濟的時候.
Do not resuscitate.
不要再救我了.
讓我舒舒服服.

$^{401}$Right.
不要再電我了.
因為充電了又電我.
充電了又電我.
不要吊著我的命.
Do not resuscitate.
保羅就是說.
It's not resuscitation.
Resurrection.
不是 resuscitation.
他用了這個比喻.
當時人明白.
現在也有人明白.
他說無知的人.
你們想一下.
你復活的身體會是怎樣.
你們所種下的種子必須要死.
然後再會生出來.
你們種下的.
並不是那將來長成的形體.
僅僅是種子.
種下去的僅僅是麥子或其他的種子.
有沒有種過東西.
你種了麥子下去.
或者種了種子下去的時候.
譬如你種橙樹.
你將橙裡面的核.
是嗎.
種子.
你放進泥土裡面.
就等於死了.
然後生出來的是什麼.
就不是橙的種子.
形體一定是不同的.
他就用一個這麼.
平日人可以去了解.
可以去觀察到的.
再去說到復活.
橙的種子種出來的.
就是橙的樹.

$^{441}$是兩種形體完全不同的.
我們的身體也是一樣.
有如種子被埋葬.
死去.
然後我們再復活過來的時候.
就是完全.
a different kind.
完全不同的形體.
不同的形體.
好像你心想的那樣.
你種下去的是很多皺紋的.
上來的時候.
沒有皺紋.
你也想.
不知道是不是.
聖經裡面沒有怎麼說過.
只是說什麼.
It's totally different kind.
It's gonna be different.
It's gonna be totally different.
四十二節.
他說死人復活也是同樣的道理.
種下去的是會扭壞.
我們這個會腐朽的身體.
復活的.
就不會再扭壞.
種下去的是羞恥的.
復活的時候就是榮耀.
種下去的是軟弱的.
復活的時候就是剛強.
種下去是血肉之軀.
復活的時候就是屬靈之軀.
你在這裡搞清楚他說什麼.
Different kind.
由dishonor.
變成有honor.
Glory.
軟弱的會變成剛強的.
然後種下去的是血肉之軀.
Our natural body.

$^{481}$Matter.
抓到摸到的.
然後復活的時候就是一個spiritual的body.
要講清楚.
什麼叫做屬靈的軀體.
屬靈的軀體不是沒有形質的軀體.
聖經裡面說.
Spiritual.
屬靈的.
其實是聖靈所驅動.
的東西.
那個形容詞spiritual.
不是說沒有任何物質.
而是driven by the holy spirit.
怎麼說也是一個totally different的body.
這裡夠少少材料.
讓我們想一想.
因為基督的復活.
帶給我們整個人的得救.
整個人身心與靈.
我們的靈魂.
非物質的.
以及我們的身體.
是屬物質的.
都會傳言得救.
變成榮耀的身體.
夠材料讓我們想一想.
到底將來.
身體復活.
是一回什麼事.
預備的時候我看了兩個的見證.
第一個見證就是Johnny Ellison Tartar.
這位姐妹大家可能都很熟悉.
因為她已經很多年全福音.
Johnny Tartar.
她年輕的時候.
還是青少年的時候.
她有一次跳水的意外.
跳水跳下去的時候.
她有意外.

$^{521}$然後她下身頸而行.
就是癱瘓.
在過去50年.
都是坐輪椅.
特別在過去這麼多年裡面.
她寫了很多書.
以及去很多地方.
來鼓勵那些殘障的人士.
有一次有人問她.
她說你覺得將來復活的時候.
你的身體會是怎樣.
她說將來復活的時候.
有榮耀的身體.
我拿到這個榮耀身體的時候.
我想我應該會有一對.
是永恆而至的.
復活了的.
我能夠立刻在輪椅起來.
然後立刻.
我就會跪在基督面前來敬拜.
這個可能都是我們的想法.
她年輕的時候.
或者她今生今世.
她身體.
她的甚至心靈.
有的傷痕.
復活的時候.
能夠完全的去拆去.
完全被殲滅.
另一個見證我聽到的.
也都是一個殘障人士.
Chatelle Kunick.
她在1980年代.
有一個女士.
一個姐妹.
也是在Johnny Tartar的.
時光裡面有份的.
一個女孩.
出生的時候就是腦殘.
她從來沒有走過路.

$^{561}$同一個差不多的遭遇.
有人問同一個問題.
她說Chatelle你.
復活的時候.
你覺得.
一個榮耀的身體會是怎樣.
她的回答就是.
她說不一定有對身的腳.
而是復活的時候.
她一個穿梭機的輪椅.
連線坐在輪椅.
但是.
是可以穿梭太空的.
不受限制.
兩個不同的對復活的景象.
一個可以說是傳統的.
傳統我們就是想.
現在這個取捨的身體.
到復活的時候是完全不同.
是沒有傷痕帶進去的.
到我們去新天新地的時候.
沒有任何創傷.
今生我們有的創傷.
無論在肉體上.
或者可能有很多時候.
心靈上是更重的創傷.
這些是不會帶去新天新地.
因為是榮耀復活的身體.
但是我們有一個.
榮耀復活的身體.
但是另一個版本.
另一個的理解.
就是說.
是的.
我們可能都會帶著一些創傷.
進去新天新地.
但是進去的時候.
到我們到了新天新地.
我們的身體得到復活.
榮耀的身軀的時候.

$^{601}$可能都會有一些傷痕帶著.
但是意義會是完全不同了.
這些都是我近年開始去.
新的一些神學的問題.
我去想.
到底聖經神學.
怎樣去處理這些問題.
這個是近十年二十年左右.
才有的題目.
就是殘障研究.
Disability studies.
就是對於一些殘障的人來說.
到底福音對他們來說.
有些什麼意思.
因為柏拉圖主義.
就是很有意思.
不要這個殘障的身體.
但是真正的福音告訴我們.
是有一個榮耀的身體會復活.
種下去的.
和生長出來的.
是完全不同的.
這個才是全人的得救.
本來就是去抗爭這件事.
兩個不同的復活的景象.
哪一個是呢.
還是兩個都是.
兩個都有可能.
為什麼我會留意這件事.
因為我在PhD的一個課程裡面.
我其中一個同學是坐輪椅的.
有一隻狗帶他來.
很不簡單.
看東西看得不清楚.
然後.
經常要坐輪椅等等.
他也在讀新學.
很佩服他.
到底上帝的大能應該是怎樣.
你聽到兩個的version.

$^{641}$都是有上帝的大能.
一個.
就是取消我們所有現在的創傷.
然後到新的身體.
來到的時候.
完全沒有傷害.
另一個上帝的大能.
就不是一種好像法術那樣.
而是我們經歷過一些跌破.
我們留下過一些傷痕.
仍然會帶進去的.
不會好像粉筆字那樣寫.
寫了之後抹掉.
上帝的大能是可以超過.
甚至是帶進去的一些傷痕.
復活的大能.
我們應該怎樣去思考這件事.
在以色列有一個博物館.
一個猶太人的博物館.
他們專門在這裡搜集.
納粹德國的時候.
有很多猶太人.
其實是會拉小提琴,中提琴等等.
因為可以隨身攜帶.
在納粹的時候.
你知道有很多猶太人.
幾百萬的猶太人.
被拉進去集中營.
然後在集中營裡面處死.
留下來的就是一些很殘破的小提琴.
這個猶太人的博物館.
專門收買這些剩下來的.
爛的小提琴.
有些是被蟲咬住了.
有些是脫頭脫骨.
破舊了的小提琴.
我們一般人來說.
我們就會覺得.
這些小提琴已經破舊了.
應該扔掉它.

$^{681}$因為都再沒有價了.
沒有價錢標記.
不如買一個新的.
我們是這樣想的.
特別是壞人.
扔掉它.
買一個新的.
但是這個博物館專門做這件事.
他們的信念是什麼.
這些破舊的小提琴.
雖然沒有價錢.
但是有價值.
沒有價錢標記.
但是有價值.
有價值.
所以他們找了很多工匠.
來修理這些小提琴.
然後讓它能夠出聲.
然後用這些舊了.
能夠被所謂復活的小提琴.
去彈奏一些音樂.
錄製成為CD.
去籌錢等等.
他們叫做violins of hope.
可能這個圖畫.
是更加讓我們明白.
到底什麼叫做死裡面復活.
基督復活.
身體復活是一回事.
不一定好像魔法一樣.
所有我們現在有的傷痕.
全部一筆勾銷.
這樣才是上帝的大能.
上帝的大能可能就是.
不是一種魔法.
而是在傷痕依舊當中.
變成祝福.
這個就是救贖的大能.
最近有一個牧者.
在網上寫了一個帖子.

$^{721}$他講到是鄭秀文.
鄭秀文拿獎大家知道.
她經歷過抑鬱.
丈夫出軌.
傷痕累累.
但很奇怪.
觀眾對她的喜愛.
是有增無減.
因為牧者這樣去解釋.
他說原因是什麼.
原因是因為她.
沒有遮掩過她的傷痕.
是一個有血有肉的藝人.
幾年前你記得她有一首歌.
叫做《我們都是這樣長大的》.
就引起很多人共鳴.
不是一些情情踏踏的歌.
歌詞裡面講什麼.
就是我們人世之間.
我們會有恐懼.
有傷痕.
有心碎.
有唏噓.
成年有成年的唏噓.
每年我們這樣適應.
這樣老去的時候.
但我們卻在這個過程裡面.
我們能夠被歷練.
可能這幅圖畫.
才是真正復活救贖的意思.
傷痕可能依舊.
但是.
有上帝救贖的大能.
我想我們應該有時間.
可以聽一聽這首歌.
OK.
可以聽一聽這首歌.
但是導出一件事.
就是.
我們人生裡面.

$^{761}$在世裡面.
我們都會帶著傷痕.
會帶著一些痛苦的經歷.
鄭秀文真的很不簡單.
很不簡單.
我曾經跟你說過.
我有時候自己祈禱的時候.
我只會為兩個基督徒的藝人.
藝人禱告.
第一個就是.
鄭秀文.
第二個是關心妍.
原因不是因為其他不重要.
而是我只認識這兩個人.
但無論如何.
她說到.
生命裡面有痛傷.
有創傷.
可能神的大能.
真的告訴我們.
真正復活的救贖.
不是要去遮掩.
也不是像粉筆紙那樣.
抹掉就算了.
而是可以帶到榮耀的一個地步.
可以見證上帝在我們身上有的醫治.
而醫治了之後.
不是說我們以前有皺紋.
有傷口.
現在我們復活了.
我們就什麼事都沒有.
我有時候都很驚訝.
復活的基督.
First Easter.
他復活的身體.
為什麼不會好像.
好像嬰兒一樣.
皮膚那麼滑.
有人放第20章.
當他向門徒顯現的時候.

$^{801}$向多瑪顯現的時候.
他說把你的指頭伸在這裡.
看一看我的手.
伸過來.
看看我的肋旁.
你看看有沒有傷口.
不要疑惑.
總要信.
最後一點.
就是基督復活的特性.
在我們人生的塵世之間.
我們會拿到不同傷痕.
我們會有不同的創傷.
這些都是咒詛.
罪和死亡.
殺蛋可以造成的.
甚至是剛才我們所說的殘障.
一些磨不滅的殘障.
我覺得復活的特性.
就不一定是一筆勾銷.
可能會是.
但可能不單止是一筆勾銷.
可能是一種轉折.
可能是一種轉化.
如果你這樣想的時候.
可能我們人生裡面.
很多時候我們盡力.
我們就是將這些傷口埋葬.
不想去面對這些傷痕.
因為覺得這些傷痕是沒有用的.
取死的身體.
但基督復活就告訴我們.
居然未必一定是完全消失.
而是一種轉折.
可能這種的力量更加大.
這種真正才是復活的彈力.
最後這段聖經說什麼.
他說告訴你一個秘聞.
就是我們都要改變.
什麼時候改變.

$^{841}$就是主在來的時候.
號角聲響起的時候.
一個扭壞的身體就會改變.
必扭壞的變成不扭壞.
死的會變成不死的.
這一切發生就是應驗聖經所說的.
死亡會被吞滅.
他說死亡你得勝的權勢在哪裡.
死亡你的毒勾在哪裡.
在復活的基督裡面得勝.
你看清楚.
他不是說死亡的毒勾是什麼.
他問這個問題不是.
你的毒勾在哪裡.
死亡與死亡.
是你的毒勾在哪裡.
你的毒勾在哪裡.
不能夠再威脅我們.
這個將會復活的身體.
有血有肉的.
我們的人生.
可能會因為自己或者他人.
加給我們的傷痕.
甚至可能是殘障.
但到基督復活的時候.
相反的這些傷痕這些創傷.
可以在新天新地裡面.
能夠被肯定下來.
不是否定.
基督的復活不是否定我們今時今生.
我們遇見的不同跌碰.
或者留下過的傷口.
不是否定而反而是肯定.
十八章最後一節.
所以親愛的弟兄姊妹.
我告訴你們你們要堅定不移.
因為知道在主裡面.
我們所作的一切.
是不會白費.
失敗有傷痕.

$^{881}$你曾經被加害過的.
也可能會有傷痕.
但在主的愛裡面.
在基督的復活裡面.
就可以有被轉化.
讓我們的生命.
我們可以讓別人看多一點.
這種復活轉化的生命.
我們同心低頭祈禱.
上主我們多謝你.
因為你自己復活救贖的大能.
讓我們能夠見到我們今生.
一切所面對.
所經歷的.
主在我們中間有很多.
我們經歷的一些創傷.
曾經有過一些傷口.
主原來你自己復活的大能.
你自己的救贖.
不只是單單要我們一筆勾銷.
奮不自磨過就算.
原來在你自己死而復活的能力當中.
能夠將這一切轉化.
上帝我們禱告.
來到你面前.
讓你看到我們傷痕累累的生命.
我們會說.
是.
我們曾經受過傷.
在罪,在邪惡.
在咒詛裡面.
我們都有跌破.
但你卻沒有否定.
你自己的救贖是去肯定.
在基督的十九號.
在你自己空的墳墓.
榮耀復活的救贖.
主我們感恩我們能夠有份.
求你讓我們這個.
已經在生命裡面展開了復活大能的生命.

$^{921}$在基督裡面有的生命.
能夠展現出對你的讚歎.
對你的讚美.
我們感激到很大.
願上帝祝福各位.
希望下星期能夠大家一起來敬拜.
\newpage



\section{9:1-9}
\label{sec:wxXYrlF2bG8}
\textbf{The Gospel of Resurrection: Act 9\_1-9 - Rev. Mark Chan}
\newline
\newline
連結: \href{https://youtube.com/watch?v=wxXYrlF2bG8}{\texttt{https://youtube.com/watch?v=wxXYrlF2bG8}} ~~~~ 語音日期: 2023-05-02
\newline
\newline
\hyperref[sec:CsnEAjBlXuk]{\small{< < < PREV SERMON < < <}}
~
\hyperref[sec:index_chronic]{\small{[返順時目]}}
~
\hyperref[sec:index_scriptual]{\small{[返順卷目]}}
~
\hyperref[sec:upXqQwuuS0M]{\small{> > > NEXT SERMON > > >}}
\newline
\newline
$^{1}$And right now we'll have past Mike for a message..
- Hey guys, if you want to find one word.
that would sum up the Bible, Old Testament and New Testament,.
one word, okay, what would that be?.
Sum up the entire Bible, right?.
How many of you say God?.
That's kind of lame, right?.
But what I'm going to suggest to you.
in the ultimate sense is resurrection..
It's this word, resurrection..
It's the victory of death over sin.
and over death over Satan in the empty tomb of Jesus Christ..
And that's why we are trying to think about,.
trying to apply what this teaching is all about,.
the resurrection of Jesus Christ..
I know it's not Easter, it's Easter behind our back.
and we celebrate Easter Sunday some time ago,.
but I want to argue and I want to suggest to you.
that the doctrine of resurrection,.
the teaching of resurrection, the reality of it.
is the core, the foundation of our faith..
From Old Testament to New Testament,.
this is the word, resurrection, the victory of Jesus Christ..
You know, I am preparing for our Israel trip.
at the end of May..
We have 26 people joining us..
10 of them are from this church.
and 16 of them are from all over the place,.
from Hong Kong, from West Coast or East Coast, anyways..
So 26 of us..
So I'm preparing some materials.
and I just stumble upon some of the books I've read before.
and I didn't pay attention to it.
when I talk about resurrection..
And these are the two books,.
written, both of them, written by Jewish scholars..
And they are Jewish, they are not Christians..
So there's no intention for them.
to promote the Christian faith,.
to promote the fact that Jesus' tomb is empty..

$^{41}$The first one on top is a top-notch Jewish scholar,.
a historian, he passed away not too long ago,.
Gisela Vermers..
Gisela Vermers has been a historian for a long, long time.
and well-respected in the Jewish world,.
in fact, in the world..
And he specialized in early century.
of the land of Israel..
And he said there's no way that the disciples.
would make up a story of empty tomb..
There's no way psychologically, historically,.
there's no way that they would expect a resurrection,.
expect an empty tomb..
So there was no purpose for them to fake one..
And down below is Shimon Gibson..
He's also a historian, not a Christian,.
and that's why I quote them, right?.
That's why I quote them..
These are not Christians..
These people do not believe Jesus is the Son of God..
But however, in their investigation,.
in their historical studies,.
they find that the tomb is indeed empty..
And how do you explain that?.
There's no way you can explain that by science.
or by history unless you adopt the biblical explanation..
The biblical reason for the empty tomb is,.
unless we adopt, look at the quote below,.
we adopt a theological one, which is the biblical reason,.
the resurrection of the Son of God..
And Shimon Gibson is very interesting..
He end this book by saying,.
I'll leave it to you guys, all right?.
This goes beyond our human proofs..
But do you want to accept what the Bible says?.
That is faith, right?.
So it's not blind faith..
It is based on historical reality..
And that's why in the Bible, it tells us.
resurrection is the core and foundation.

$^{81}$of everything that we believe, the entire Bible..
Paul said, if Christ has not been risen,.
all our preaching is useless and your faith is useless..
I like the SV translation, in vain, you know?.
Everything, every single thing that we do,.
that we talk about, that we believe in, is in vain..
It's useless, it's stupid..
It's just nonsense if Christ has not been risen..
But he has been raised..
He was resurrected from the dead..
And give us new chapter, give us new hope..
And I have been trying to focus on the fact.
that God is trying to bring the new heaven.
and new earth to us..
And he began the new creation project.
in the resurrection of Jesus Christ..
Not someday in the future,.
but already it has broken into our lives.
if we accept Christ..
That's why the Bible keeps on talking about,.
in Christ, the old is gone, the new has come, right?.
Speaking of new chapter, speaking of new beginning,.
you know, I hope you feel that way recently,.
'cause spring has finally come to Chicago, right?.
The long-awaited spring has already come..
And for my personal life, it has also come,.
a new chapter, a new, sort of a new beginning,.
and it feels good about it..
You know, I'll tell you what,.
in the past term, a few months,.
has been really, really tough..
You know why?.
Because Rachel has come back to Chicago.
for student teaching,.
and she has to go to Healy School to do the teaching..
And then she has to show up 7.30 in the morning,.
on time, right?.
And, you know, being a good father,.
I give her a ride every single morning,.
every single, well, one or two,.

$^{121}$I give them by mom, right?.
Okay, all right, good for her..
You know, if you know me enough,.
I'm not an early riser..
I've never been..
You know, the most efficient hours of my day,.
I don't know about you..
How many of you are early riser?.
Kelly, you are?.
What time do you wake up, usually?.
- [Kelly] I'm a fatherly..
- Okay, okay, adult, right?.
So, she is early riser..
Make you one..
And, you know, I've never been..
The most efficient hours for my day.
is 9 p.m. to 2 a.m..
Anyone amen to that?.
There you go, right?.
That's my most efficient hour..
I can get so much done at night, right?.
Night hour..
So, when she came back, when Rachel came back,.
I had to drive her every single morning..
It's killing me..
I'm sure it's killing her as well..
But finally, last week, last Friday,.
it was the last day of her work..
So, praise God..
Hallelujah, hallelujah..
So, I feel kind of good..
I don't have to, you know, think about this anymore,.
or at least for the time being..
I don't have to do it anymore..
So, you know, I'm so glad to drive her to Starbucks.
to get some, you know, bonus, right?.
Some sandwich and some coffee in the last week,.
and that's good..
So, that's a new beginning..
A new attitude, right?.

$^{161}$New attitude of morning, to at night, to my life, right?.
That is what the resurrection should bring us..
A new beginning..
The new creation has begun.
in the resurrection of Jesus Christ..
That's why the Bible talks so much about resurrection..
Many, many times, all four Gospels,.
all the New Testament letters,.
at least they talk about it explicitly,.
or they stress on it as the foundation of our faith..
Why do we follow Jesus?.
Because we are not following a dead man..
We are following a living God, a living person,.
Jesus Christ himself..
And today, we would like to take a look.
at one of the stories of the encounter of the risen Christ,.
and Paul's, the apostle Paul's own encounter.
of the risen Christ..
Now, according to the timeline of the Bible,.
Paul seems to be, the apostle Paul,.
seems to be the last person.
who has witnessed the risen Christ..
And we found this record in the chapter nine.
of the book of Acts..
We're gonna read through some of the verses here,.
the story, the historical record,.
and then there are a couple of questions.
I want to quiz on you..
All right, nine one..
Meanwhile, Saul was uttering threats..
With every breath,.
and was eager to kill the Lord's followers..
So Saul, Saul is Paul, Paul is Saul,.
two different languages..
He went to the high priest,.
and he requested letters addressed to the synagogues.
in Damascus, asking for the cooperation.
in the arrest of any followers of the way,.
the way of Jesus, the way of the Lord..
He found there..

$^{201}$He wanted to bring them, both men and women,.
back to Jerusalem in chains..
As he was approaching Damascus on this mission,.
a light from heaven suddenly shone down around him..
He fell on the ground and heard a voice saying to him,.
Saul, Saul, why are you persecuting me?.
Who are you, Lord, Saul asked..
The voice replied, I'm Jesus, the one you are persecuting..
Now get up and go into the city..
You will be told what you must do..
The man with Saul stood speechless,.
for they heard the sound of someone's voice,.
but saw no one..
Saul picked himself up off the ground,.
but when he opened his eyes, he was blind..
So his companions led him by the hand to Damascus, the city..
He remained there blind for three days.
and did not eat, did not drink..
Now this is a story, and it's such a dramatic story,.
and there are two points here..
And first of all, two questions first..
What was Paul planning to do in Damascus?.
To write a letter or to persecute Christians.
or to meet the high priest?.
What is it?.
To be, right?.
To persecute Christians..
All right, that's easy, right?.
The harder one here..
Where did the raising Christ appear to Paul?.
On the way to Jerusalem, to Samaria, or to Damascus?.
Which one?.
See, right?.
Damascus, right?.
Damascus..
Not Jerusalem, away from Jerusalem,.
away from the hill country of Samaria,.
coming on his way, approaching Damascus..
The two points here..
First, this story, we see the raising Christ.

$^{241}$in an unworthy place..
Unworthy place, unexpected..
Second, Christ, who is resurrected,.
for an unworthy person..
Two points..
Now I think we need to understand this passage..
We first need to understand the background of Paul..
Who is this guy?.
Who is this Paul or Saul, right?.
Now I think one of the ways we can understand him.
is that he is kind of like the super orthodox Jew,.
in his days, right?.
Super orthodox, super kosher, right?.
And following the law straightly,.
following the traditions of Judaism straightly..
And today we can kind of picture him.
as some of the ultra orthodox people, right?.
These are called the ultra orthodox Jews, right?.
Where can you find them?.
Some in Skokie, right?.
Some in Brooklyn..
And of course, many of them in Jerusalem, right?.
And they are the super conservative communities,.
especially they are following straightly the Shabbat,.
the Torah, going to synagogues on time,.
wear black and white all the time..
Isn't that nice, black and white all the time?.
You don't have to worry about your outfit.
every single morning..
You just open your closet, everything is the same, you know?.
'Cause they put the energy in.
following the way of the Torah, right?.
And Paul is 100\% or even 200\% of orthodox Jew..
Now, if you go to Jerusalem today,.
say, I usually drive my own car, rental car..
If I drive at the wrong time,.
in the wrong neighborhood, in one of those neighborhoods,.
on the wrong time, like Shabbat,.
I can get stoned easily.
'cause you are violating the most sacred law, right?.

$^{281}$So, you know, we got to know our way, right?.
We can't just drive into the religious orthodox community.
without knowing the timing, right?.
So you can get stoned, you know?.
You got to be very, very careful..
These are the people like Paul..
And Paul in his day, you need to understand,.
he was one of the super Pharisees, right?.
Super Pharisees..
And the Pharisees trying to do their best.
to follow the letters of the Torah..
Now, sometimes we got it wrong..
Sometimes we thought, okay, these people,.
they are doing some work righteousness..
They're trying to follow the Torah and fulfill the Torah.
so that they can get so-called saved in Christian terms,.
so that they can be considered the chosen people..
No, I think we got it wrong..
It's the other way around..
In Paul's day, Judaism means.
some of the people, the Jews, the orthodox Jews,.
they consider themselves as chosen people..
And because of this identity,.
that's why they show you who they are..
That's why they observe the law..
That's why they do the law,.
not because they want to be counted as the people of God..
Oh, there are one or two laws I violate,.
so I'm not the people of God..
I have to fulfill those..
No, it's the other way around..
They are chosen people..
That's their identity..
So they don't want to compromise their identity..
That's why they want to do the law..
They set up different boundaries around them,.
time boundaries, geographical boundaries,.
to not to compromise their identity..
Who do you eat with?.
That's one of the boundary,.

$^{321}$one of the cultural religious boundary..
How to worship, what time to go to worship..
In Paul's days, Pharisees like Paul,.
believe that if everyone observe the law completely,.
then the kingdom of God will arrive..
That's why they want to follow the law..
That's why they're zealous and enthusiastic.
about doing the law..
Because this is not an individual thing to do..
Everyone, every Israelites, Jew,.
if we all follow the law,.
then the kingdom of God will be ushered in..
There's only one way..
We follow the law together..
So not only I follow the law,.
I want to make you to follow the law..
And you are zealous about following the law,.
me zealous about following the law,.
all of us, the entire community, the entire people,.
then the kingdom of God will come..
That is our hope..
That is the means..
Now this explain, this kind of background will explain.
why Saul or Paul would persecute the Christians..
The early Christians in the first century, who were they?.
Ethnically speaking, who were they?.
Jew..
They were Jewish people..
They were Jewish people 100\%, ethnically speaking, right?.
And they recognized Jesus as the Messiah..
Now to Paul as the Pharisee, this is troubling..
That's not good, right?.
For one thing, you believe in Jesus the Messiah.
being crucified on the cross..
That is a shameful thing, right?.
That's a stumbling block..
And secondly, there's no way that he rose from the dead..
There's no way that this man being crucified on the cross.
and his tomb was empty..
There's no way you can get around that..

$^{361}$That's why Paul was trying to stop this Christian church,.
you know, Jewish Christians,.
or what we call today Messianic Jews..
Those people who are Jewish.
and recognize Jesus as the Messiah..
And Paul was trying to do his best as a Pharisee.
to stop this so-called virus from spreading out, you know?.
Remember the outbreak of COVID, right?.
Everyone go to the computer, check what?.
Check the world map, right?.
And see the red dots, you know, spreading around,.
you know, different countries and different regions..
And, "Oh, it's coming to my place..
"It's coming to our region.".
And we are scared, right?.
You know, most of us are Asians..
We are scared..
Well, Asians are easily being scared, right?.
We are free of death..
That's why we pay attention to those dots, right?.
And the spread of the virus..
And Paul in his day as a Pharisee,.
as a super culture rabbi, he was so concerned about this..
That's why the story of chapter nine happened..
What is Paul trying to do?.
You look at this map..
Where is Paul going?.
He begins from Jerusalem, he goes north..
And then most likely he avoid,.
not the most straightforward path to go north.
because who are those people in the north?.
Samaria, right?.
The Samaritans, right?.
You know, Samaritans and Jews,.
you know, a little bit background,.
you know, they are like,.
they don't like each other, right?.
They don't eat with each other..
They're actually hostile to one another..
So he probably scroll around.

$^{401}$the northern Samaria kill country,.
and then go up north to that, you know,.
little lake there, the Sea of Galilee,.
and then go northeast and climb the plateau.
of what we call Golan Heights today,.
and go into a Gentile, totally Gentile territories..
And the foothill of Mount Hermon, can you make it out?.
Mount Hermon, he is on his way to Damascus..
It's, that is a total of 150 miles.
and six days of walking nonstop..
Now, if I were Paul, I have super culture,.
you know, Orthodox Jew, super culture Rabbi,.
I would be so, I'd probably not be walking..
I'll be running, you know, I'll be running..
I'll be running for a few, you know, marathon races..
Let's try to get to Damascus..
But your question is, why Damascus?.
Why Damascus?.
Why not other cities?.
Why not Chicago, right?.
Why Damascus?.
Now, remember, when the early church begins to be persecuted.
by the authorities, Roman authorities.
or Jerusalem authorities,.
the Christians were scattered everywhere.
'cause they are trying to, you know,.
run away from troubles, right?.
And bringing with them the gospel, right?.
That's what the book of Acts is all about..
When persecution came to Jerusalem,.
to the Jewish Christians, they scattered off.
to different towns and villages or cities..
And Damascus is at the crossroad of civilizations..
It's the international center of the day..
If you go to Damascus, then you can go right,.
you go to today's Iraq, Mesopotamia,.
a lot of Jews over there since exile of Babylon..
And then if you turn left from Damascus,.
you go to what they call Turkey, Asia Minor,.
and then Europe, right?.

$^{441}$And a lot of Jews in Diaspora,.
they're scattered into different synagogues..
And you know, they said Damascus is like New York..
It's today's New York..
It's a crossroad of everything..
It's an international crossroad..
That's why the apostle Paul,.
well, back then he was not the apostle,.
the super-cultural rabbi Paul..
He went to Damascus, that's why he went there..
Because he wanted to do,.
probably he wanted to engage with some preemptive strike.
against the early Christianity, right?.
That is how you stop Christianity from spreading, right?.
That's how you stop the so-called virus from spreading..
You go to the frontier and try to persecute them.
and bring them back to Jerusalem..
That is the mindset of Paul..
Now, you need to also understand.
the Jewish super-cultural rabbi's mindset..
Paul, what was Paul thinking?.
And in some way, in the first century,.
Judaism, what Paul believed in,.
he was advocating some sort of, you know,.
Second Temple Judaism or First Century Judaism..
And there are three T's in Judaism, in Paul's day..
The first T is time, second is Torah, third is temple..
Time, Torah, and temple..
Now, what do I mean by that?.
If you follow the law, you need to observe the time..
Time of Shabbat, especially, right?.
It's time to worship,.
and you have to follow the law, straightly..
Torah, of course..
And then, most importantly for this story, is temple..
The most sacred place on the planet Earth,.
the entire universe, is one place,.
and you can find it on the map..
You can point it to a GPS coordinates..
That is where the temple is..

$^{481}$That's where the holy of holy is..
Even today, Orthodox Jews,.
they pray toward some direction, right?.
Not to Chicago..
They pray facing Jerusalem, right?.
They pray and praise,.
they live their life toward that direction..
You go to Western Wall,.
why do the Orthodox Jews would go to Western Wall.
more than any place else?.
You know, Western Wall is simply an open air synagogue..
It's like open air church..
So people just go there to read scripture,.
to read Torah and pray..
That's why they go to the wall,.
because the wall is closest that they can get today,.
to the holy place, the most holy place,.
you know, where the golden dome of the rock has mark..
That's why they wanna go to Western Wall, right?.
For Paul, the most sacred place on Earth,.
the sacred spot is Jerusalem..
Now with that in mind, you can understand,.
Paul is not simply a Roman citizen.
traveling from one region to another, right?.
It's not like us traveling from Chicago to San Francisco,.
or Chicago to Indianapolis..
It's not from one region to another region..
Nothing is neutral for them..
Jerusalem is the most sacred place..
And the further away from Jerusalem,.
the secular you get..
So there is a sacred and profane difference..
For Paul traveling from Jerusalem to Damascus,.
he's actually traveling from a sacred place,.
the most sacred place, to the territory of the enemies..
Where is Damascus?.
Damascus is in what we call Golan Heights or beyond, Syria..
In the Old Testament, it's Aram..
If you know a little bit about the Old Testament history,.
you know Arameans, people from that region,.

$^{521}$and the Israelites in Israel or in Judah,.
they were always enemies..
So for Paul, for this super-kosher rabbi,.
he was actually going from a friendly place.
to the enemy's territory..
And that's why it is totally unexpected.
for Paul to meet and encounter the living God..
On his way to Damascus, there was no way..
And that's why for three days he did not eat.
or he did not drink..
It's just earth-shaking for him..
It shake up his whole world..
And he was dumbfounded and silent for three days..
He's trying to figure out what the heck is going on..
I did not expect God to show up on the way to Damascus.
in the territory of the enemies..
I expect God to show up where?.
In the temple, right?.
In Jerusalem, in the most sacred place on earth, right?.
And the resurrected Christ is telling us.
that he is no longer bound by territories..
He is no longer bound by geographical boundaries..
Everywhere he is, because he has ascended.
to the right hand of the Father..
He is the resurrected Christ..
He is the living God..
He is not bound by anything else on the Google map..
You can meet God here,.
as well as you can meet God in Jerusalem..
I remember some people who joined my Israel trip..
Sometimes we go to the Western Wall, right?.
Everyone wants to go to the Western Wall, right?.
At least you need to take a picture there..
And they kind of find themselves in a situation..
Oh, this is the temple..
And they say, "Hey, Pastor Mark, please,.
"this is the temple..
"This is the most holy place.".
And I said, "No, no, this is not the holy place..
"The holy place is in the person of Christ, amen?.

$^{561}$"You no longer worship Christ on the GPS coordinate..
"You worship Christ in the living person..
"The dead, once dead, and now risen, Christ himself.".
And that's why in this story,.
the risen Christ is showing up.
in an unworthy, unexpected place..
And second point is that the risen Christ.
show himself, demonstrate himself to an unworthy person..
Now, if you ask your friends,.
you know, if you have some neighbors,.
you may be able to, well, most likely,.
you can kind of put them on a kind of spectrum..
How close are they to God?.
How close are they to the gospel?.
Some of your friends or some of your people.
that you know or neighbors,.
you may, at least in human speaking,.
they're so far away from the gospel, right?.
Maybe they're even hostile to the gospel..
And some of them may be more interested in faith.
and asking you questions.
and you answer some of the questions..
You know, you could put them on a spectrum..
You know, humanly speaking..
And then you think about it..
Paul is actually the least expected person.
to come to Christ, right?.
He's again, super culture rabbi, right?.
Super culture, orthodox Jew..
How would you expect him to come to Christ?.
And this story, quite unexpectedly and surprisingly,.
telling us the persecutor of the church coming to faith.
because the power of the risen Christ..
And there is a second point..
God's saving grace in an unworthy person..
Essentially, you think about it,.
that applies to all of us, right?.
We're not more worthy than other people to come to Christ..
It's all about saving grace..
If you track the story of the book of Acts,.

$^{601}$you can start with chapter seven,.
a little bit before the story that we read..
Paul's conversion, the way to Damascus, chapter nine..
Chapter seven already talks about God showing himself,.
his glory to many different people..
And that is the long sermon of Stephen, the first martyr..
And in that long sermon in chapter seven,.
he's actually doing something like a summary.
of the entire Old Testament..
So we start with Abraham,.
that the glory of God showed to Abraham.
in a foreign place, in Mesopotamia..
Where was Abraham from?.
Where was he from?.
Country of birth..
Mesopotamia, Iraq..
In some sense, Abraham is Iraqi..
He was not straightly an Israelite..
He's a foreigner, but God showed himself to him..
And then Stephen surveyed the Old Testament story.
and saying that in wilderness,.
God showed himself to Moses and to the Israelites..
And he talks about different episodes.
of the Old Testament story,.
and God will show them the glory.
and be present with his people..
And interestingly, what is not mentioned.
in Stephen's sermon was Jerusalem..
He never talked about Jerusalem, the first temple,.
the glory of God filled the temple.
when it was devastated in the time of Solomon..
He never mentioned that..
At the end of his sermon, his punchline here,.
the most high God does not live in temple.
made by human hands..
He was probably, he was speaking this sermon in Jerusalem..
He was probably pointing to the temple..
And that's why the Jews wanted to kill him, right?.
You speak bad words about our temple..
Are you kidding me?.

$^{641}$Right?.
And then they stone him..
And then chapter eight, Philip sent to Samaria..
Samaria, the enemy's territory..
Samaritans, unworthy of Christ..
And they came to Christ..
And then the end of chapter eight,.
he was being sent to a eunuch.
and somebody with defect in his body..
And he's unworthy to go into the temple,.
going to the presence of God to see God..
But God show himself to this unworthy man..
And then chapter nine is Paul's conversion..
Chapter 10, Peter is being sent to Caesarea Militima..
Caesarea was the Las Vegas of the first century..
And he met a guy called Cornelius, a Roman,.
someone who is enemy to the Jews..
And he came to faith..
He was unworthy..
And everyone is unworthy..
But the living Christ show himself to them..
And story after story proves one point,.
that God's glory, God's grace is extended to all,.
is not bound by human boundaries..
And God's resurrection power can be displayed.
in every single one of us, even you and me..
And that's why Paul and his chapter,.
chapter 15 of first Corinthians,.
he says, "I am not the worthy one to receive this grace,.
"special favor on me..
"I work harder than anyone because I feel unworthy,.
"because I sense, I feel this grace is so real to me.".
The living Christ showing himself to me..
And made me a new man, a new beginning..
In Christ, we do have new beginning..
Paul is the proof..
And you are the proof..
We are the proof..
In Christ, we should be new,.
in some foundational way..

$^{681}$I'll close with one story..
I know a sister, a sister in Christ in Israel,.
some years ago..
And she has been coming to church for worship.
very faithfully..
And she is married to a Jewish man..
And as you know, most of the Jews,.
they don't accept Jesus as Messiah..
This is, you kidding me?.
This is a joke, right?.
This is something block..
But eventually, you know,.
I spent some time in Israel, in Jerusalem,.
and I often go to that church every single week..
And then eventually, her husband comes to the church..
And according to her husband,.
why is he attracted to the church?.
He's a Jewish man..
And he said, "Very simple..
"Because I see the differences.
"the gospel makes in my wife's life.".
And he said, you know, every single week.
after she goes to church,.
the husband feels like she's different..
She's like some sort of new person..
Joyful, more positive, you know,.
dealing with life issues, you know..
A foundationally different way..
And, you know, cooking is better, at least for the day..
So he says the difference, you know..
I see, you know, that's the whole point of the story..
Sometimes we come to church, we feel like,.
oh, this is such a burden, right?.
A checkbox, okay, I'm done, okay?.
Okay, I go on my way Monday to Saturday, right?.
Okay, you know, forget about it..
This is what the difference should make in our lives..
A new beginning, right?.
A new pursuit..
Back in our old self, back in our old days,.

$^{721}$what do you pursue about?.
What do you think about most of the time?.
You know, lust, material things?.
What do you concern about?.
What do you worry about?.
You know, the living Christ who died and rose again.
should give us a new start.
so that we can work toward the new world,.
which is the new heaven and new earth..
We pray that we can do that and make that decision today..
Think about what resurrection will make.
a foundational difference in your life,.
your attitude, your direction, right?.
The object of your pursuit..
Everything will be ruled by Christ.
and we are working toward that goal, right?.
Let's all pray and close..
And we're gonna sing our last song,.
"You Are My No No,".
and I ask that you praise Him to come up here..
Let's all pray..
Father, as we close,.
we ask that your Spirit will remind us.
the old things that are still, you know,.
putting a hose on our lives,.
making us to look back all the time..
And I pray, I pray that you will make us new..
In Christ, who has died and rose again,.
He is the one who give us hope..
He is the one who give us a new chapter.
and help us to be in Christ..
The old has gone and the new has come..
And praise you, Father,.
because we believe and we follow a living Christ..
Show yourself, show your glory.
in this place, this unworthy place..
Show your glory and your grace to this unworthy person..
And thank you for your grace..
In Christ, we pray..
[BLANK AUDIO].

\newpage



\section{使徒行傳 9:1-9}
\label{sec:upXqQwuuS0M}
\textbf{復活的福音:使徒行傳 9\_1-9 - 陳崇基牧師}
\newline
\newline
連結: \href{https://youtube.com/watch?v=upXqQwuuS0M}{\texttt{https://youtube.com/watch?v=upXqQwuuS0M}} ~~~~ 語音日期: 2023-05-02
\newline
\newline
\hyperref[sec:wxXYrlF2bG8]{\small{< < < PREV SERMON < < <}}
~
\hyperref[sec:index_chronic]{\small{[返順時目]}}
~
\hyperref[sec:index_scriptual]{\small{[返順卷目]}}
~
\hyperref[sec:ixAcYGJ3Y7Y]{\small{> > > NEXT SERMON > > >}}
\newline
\newline
使徒行傳 9:1-9
\newline
\begin{longtable}{cl}
\hline
\hline
章節 & 經文 (和合本修訂版)\\
\hline
9:1 & \begin{tabularx}{0.7\textwidth}{X} 掃羅不斷用威嚇兇悍的口氣向主的門徒說話。他去見大祭司， \end{tabularx} \\ \\ \relax
9:2 & \begin{tabularx}{0.7\textwidth}{X} 要求發信給大馬士革的各會堂，若是找著信奉這道的人，無論男女，都准他捆綁帶到耶路撒冷。 \end{tabularx} \\ \\ \relax
9:3 & \begin{tabularx}{0.7\textwidth}{X} 掃羅在途中，將到大馬士革的時候，忽然有一道光從天上下來，四面照射著他， \end{tabularx} \\ \\ \relax
9:4 & \begin{tabularx}{0.7\textwidth}{X} 他就仆倒在地，聽見有聲音對他說：「掃羅！掃羅！你為甚麼迫害我？」 \end{tabularx} \\ \\ \relax
9:5 & \begin{tabularx}{0.7\textwidth}{X} 他說：「主啊！你是誰？」主說：「我就是你所迫害的耶穌。 \end{tabularx} \\ \\ \relax
9:6 & \begin{tabularx}{0.7\textwidth}{X} 起來！進城去，你應該做的事，必有人告訴你。」 \end{tabularx} \\ \\ \relax
9:7 & \begin{tabularx}{0.7\textwidth}{X} 同行的人站在那裡，說不出話來，因為他們聽見聲音，卻看不見人。 \end{tabularx} \\ \\ \relax
9:8 & \begin{tabularx}{0.7\textwidth}{X} 掃羅從地上起來，睜開眼睛，竟不能看見甚麼。有人拉他的手，領他進了大馬士革。 \end{tabularx} \\ \\ \relax
9:9 & \begin{tabularx}{0.7\textwidth}{X} 他三天甚麼都看不見，也不吃也不喝。 \end{tabularx} \\ \\
[1ex]
\hline
\hline
\end{longtable}
$^{1}$多謝Carmen提醒我們教會有什麼重要的事情.
特別今天會有大會.
我們會繼續去看復活的信息.
講到復活節.
講到resurrection.
講到聖經裡面怎樣提到復活這件事情.
是我們基督教信仰的核心.
沒有了這件事 有神的愛也沒有用.
沒有了這件事 有上帝的恩典也沒有用.
因為這個復活節.
或者復活的事實.
是我們生命裡面一個很重要的基石.
沒有了這件事就什麼都沒有.
聖經新約裡面也是這樣講.
是事實上這樣強調.
說上帝的復活.
在耶穌基督的復活.
是給我們一個新的開始.
因為上帝自己將來會帶來的新天新地.
新的耶路撒冷.
是從耶穌基督復活的那一刻.
來加入的.
加入了你和我的生命.
所以我們在基督裡面也是有一個新的開始.
講到新的開始.
我們這裡芝加哥的天氣又開始回暖了.
我們在芝加哥真是很有福的.
四季都看到.
很多時候一個星期有四季.
我們真是很好.
講到春天的時候就是一個新的開始.
我們也知道在我們中間的課堂.
或者我們英文課.
也有些朋友可能他們的家人.
或者他們自己是畢業或者升學等等.
都是一些新的開始.
你知不知道我的女兒畢業了.
她5月14日畢業.
就是媽媽那天畢業.
所以那天我就不能夠和大家一起.

$^{41}$因為不知道為什麼大學不喜歡在母親節那天.
畢業禮.
真是麻煩.
所以我們一家就要去到Champaign那裡.
看她的畢業禮.
都是一個新的開始.
特別是一個新的開始不是畢業禮這麼簡單.
因為她過去一個學期.
她都已經回到芝加哥.
她因為要做教師.
所以她就在一間小學那裡.
去做Student Teaching.
就是做實習.
因為她一向都不是早起的人.
整個學期.
過去整個學期.
她是晨運這麼早起床.
六點多就要起床.
有誰六點多起床?.
哇 這麼多.
哇 這麼厲害.
我們約晨運.
有誰不是六點多起床?.
是很難的.
OK Good Amen.
我也不是.
我們兩個都不是早起的人.
因為你熟悉我的人.
我是夜鬼.
我最有效率的時間.
是晚上九點至凌晨兩點.
因為我要載她.
你知道我很喜歡女生.
我載她上班.
真是要命.
所以我六點多就要起床.
說真的就算她開車.
你也不放心.
或者她開車.
她也趕不上.

$^{81}$她在Healy小學上班.
拍邊.
除非有電單車給她.
因為這個緣故.
所以她頭一天的學生教育.
我就要載她.
我覺得很辛苦.
很辛苦.
她也很辛苦.
我更辛苦.
OK 我更辛苦.
我載完她之後.
簡直殘了我整個人.
所以很辛苦.
但她在過去的學期.
她也堅持.
也挺厲害.
所以有時候她要多花點錢.
再去買一杯Starbucks.
OK Fine.
It's OK.
剛好這個星期她結束了.
沒有一天遲到.
很厲害.
其實有幾天差點遲到.
因為我開車.
所以沒有遲到.
怎樣也好.
很辛苦.
但你看到她結束了這個學期.
終於完結了.
踏入一個新的章節.
她很開心.
完結了一個所謂的使命.
或者一章.
是一個很開心.
踏入新的一章的時候.
就不再回頭看.
當然.
將來去找一份工作.

$^{121}$一樣是這樣.
一樣這樣做.
但怎樣也好.
由舊去新.
是一種新的感覺.
一種新的經歷.
其實耶穌基督的復活.
就是想帶給我們這件事.
往往我們很多時候忘記了.
沒有被提醒.
我們做著做著基督徒的時候.
我們就忘記了這件事.
你知道耶穌基督的復活.
是一件很重要很重要.
很大件的事.
在聖經整個的救恩歷史裡面.
是一個中心.
十字架之後.
一定要有空墳墓.
空墳墓就是.
將整個的歷史.
分為兩半.
一個是舊的.
一個是新的.
有基督復活的緣故.
所以整個的世界就不同了.
所有屬基督的人.
他的方向也都不同了.
我正在預備去以色列.
因為五月尾會帶頸尖之妹去以色列.
我有一天就去到.
留意到.
之前沒有留意過這本看過的書.
下面那本我提過了.
講到復活的.
兩個的作者都是猶太人.
兩個都不是信耶穌.
所以為什麼要引用他們.
那些人就覺得.
賣花當然是賺花香.

$^{161}$聖經這樣講嘛.
當然聖經作的也可以.
我們有很多很瘋狂的想法.
我們覺得.
福音書裡面或者聖經裡面.
新約裡面講到復活的時候.
一定是講到賣花賺花香.
我就特別要引用這些作者.
或者學術的研究.
兩個都是猶太人.
兩個都是不信的.
第一個就是基辛格.
基辛格是一個很重要的.
已經過去世了.
很重要的一個猶太的識經者.
和一個猶太的歷史學家.
他說什麼呢.
他說沒有可能門徒是偷屍的.
偷屍不是偷秘密.
偷屍體.
沒有人偷屍.
沒有可能.
因為從他的歷史的角度來看.
從心理的角度來看.
那些門徒根本沒有可能會想出.
或者虛構一件這樣的事.
因為沒有任何動機.
因為沒有任何期望.
耶穌會首先復活.
但是誰知道.
新約的歷史的記載.
就是說他們就發現.
這個墳墓是空的.
耶穌是第一個從死裡復活.
並且帶給我們將來.
一切在基督裡面的人.
有身體復活的盼望.
下面Simon Gibson不說了.
這裡你可以自己參考一下.
他就說用歷史的角度.

$^{201}$科學的角度根本無法解釋.
總之那個墳墓就是空的.
就是空的.
除非你接受聖經的解釋.
就是上帝使神的兒子從死裡復活.
基督復活很重要.
他說不多前書第十五章.
上次提到最後結論的時候.
保羅說.
若基督會有復活.
我們傳的便是謊言.
你們信的也是謊言.
useless.
我很喜歡ESV的翻譯.
他說in vain.
所有東西都是徒勞的.
你今天為了基督所走的路.
信耶穌的信仰.
一切都是徒勞的.
如果基督沒有復活.
基督既然有復活.
所以一切都是有意義的.
不會是空跑的.
謊言不會是徒勞的.
是有內容.
並且是有盼望的.
聖經裡面強調耶穌基督復活的事實.
和我們生命當中有的意義.
一次又一次在新約聖經裡面.
數不盡的這麼多次提到的.
新約的書信.
福音書裡面的記載.
所以今天我們看一段的故事.
是講到基督復活的顯現.
史羅行傳第九章.
看完的時候我想問你兩個問題.
我們留心看.
第九章第一節.
蘇羅仍然向主的門徒.
口吐威嚇兇殺的話.

$^{241}$去見大祭司.
求文書給大馬識.
或者大馬士革的國會堂.
若是找到信奉者道.
即是那些基督徒.
無論男女都准他捆綁.
帶到耶路撒冷.
蘇羅就上路.
將到大馬識.
忽然從天上發光.
四面照著他.
他就突然在地.
聽見有聲音對他說.
蘇羅蘇羅你為什麼逼迫我.
他說主啊你是誰.
主說我就是你所逼迫的耶穌.
起來進城去.
你所作的必有人告訴你.
同行的人站在那裡.
說不出話.
聽見聲音卻看不見人.
蘇羅從地上起來.
睜開眼睛.
竟不能看見什麼.
有人拉著他的手.
領他進大馬識.
三日不能看見.
也不吃也不喝.
這裡說到大馬識的顯然.
復活的主和保羅.
來到顯然.
第一個問題.
準備好了嗎.
第一個問題.
保羅去大馬識其實計劃了什麼.
他是寫信去大馬識.
還是迫害基督徒.
還是見大祭司.
他在耶路撒冷出發.
去大馬識.

$^{281}$還沒到大馬識的時候.
他想做什麼.
去大馬識抓那些基督徒.
以至他們能夠被帶回耶路撒冷.
然後他們受迫害.
嚴刑敲打等等.
所以是迫害.
第二個問題.
復活的基督徒在哪裡向保羅顯然.
是往耶路撒冷的路上.
往撒瑪利亞的路上.
還是往大馬識的路上.
是什麼.
是施而為.
大馬識的路上.
去大馬識的路上.
所以這段聖經說到他上路去大馬識的時候.
他就見到復活的主.
然後這個本身迫害教會.
迫害基督徒的使徒保羅.
以前叫蘇羅.
蘇羅或者保羅.
不同兩個語言.
然後他信了主.
然後他成為了上帝自己的使徒.
這裡有兩點.
第一點.
就是復活的主.
是一個根本不會預期不配的地方出現.
第二點.
就是復活的主向一個不配的生命顯然.
保羅在大馬識的路上.
見到復活的主.
榮耀的主.
向他顯然.
這是一個歷史的事實.
也沒有什麼人會懷疑.
並且在自古到今.
也有很多人記起這件事.
有很多畫家.

$^{321}$一些很有名的作品.
都是以這個主題來畫出來的.
你可以看到左邊這幅.
1600年的時候已經有的油畫.
裡面提到.
或者用了一個題材.
就是他上大馬識的時候.
突然有榮光向他顯然.
復活的主向他顯然.
然後他就從馬上跌下來.
然後那種混亂.
這件事到底有什麼我們需要留意.
第一我想需要留意的是.
到底保羅是一個怎樣的人.
他本身是一個怎樣的人.
他的家底或者他的背景是怎樣的.
保羅在不同的地方提到.
說他是一個法理錯人.
是嚴守上帝的律法的.
他是一個所謂超正統的猶太人.
我覺得今天如果你要畫圖.
保羅是一個怎樣的人的時候.
你就要想這些人.
這些叫做超正統的猶太人.
在布魯克林紐約或者斯科奇.
更多當然在耶路撒冷.
這些超正統猶太人.
日日夜夜他們所穿的衣服就是黑與白.
每天在家裡他們就是去祈禱.
然後去讀經.
在以色列的政府或者以色列的環境裡面.
他們是有政府養的.
所以每天就是去到嚴守他們的律法.
好不好這樣做人.
OK喔.
不用做.
不用做啊要做啊.
你去禮拜的.
每天都要去的.
不用做.

$^{361}$好像你上班一樣.
他們就去西場或者去不同的會堂.
他們就去讀經去祈禱.
飲食方面當然他們也是超正統.
所以會按照摩西律法的吩咐.
等等等等.
這些就是所謂當時的保羅.
如果我們明白.
法利賽人到底會發生什麼事的時候.
可能會多一點點明白保羅的行動.
或者訴羅的行動.
我們千萬不要搞錯.
有時候我們誤會.
就是他們守律法.
是因為他們覺得守盡律法.
他就可以得救.
用基督徒的術語來說.
我們做不夠.
我們要做夠他.
守夠安息人.
守夠食物的條例.
Culture等等.
我們才可以叫做神的子民.
才能夠做到神的子民.
做不到的時候我們就很缺陷.
我們就不夠了.
我們就做不到上帝的子民.
No No No.
不是啊.
是相反的.
相反就是因為他們.
認為他們是神的選民.
所以他們去嚴守律法.
你明白我的意思嗎.
不是因為他們做不夠.
不是一種Work Rights.
大部分的猶太人.
特別我講第一世紀的時候.
訴羅保羅等等.
他們不是為了守律法而去得救.

$^{401}$而是相反因為他們已經是立約的子民.
起碼他們覺得他們是.
因為他們血統上是.
我們就是那種族.
或者我們就是那班的選民.
所以我們的身份.
就應該履行這些律法的生活和責任.
是反過來的.
以前的法理賽人.
當時的猶太教裡面.
其中一個的學派.
其中一個的支派 中派.
法理賽人.
他們守律法是因為他們是神的子民的身份.
並且他們相信.
或者那些真的超正統的法理賽人.
他們相信如果全民皆守律法的時候.
所有人 我們每一個都是守律法.
所有的猶太人.
我們同種族的人.
我們的骨肉之親.
全部都是去守律法的話.
神的國才會來到.
所以他們說叫做大發熱心是什麼意思.
不單止他們要守律法.
也是大家一起守律法.
見到其他的自己人.
不去守律法的時候就很大件事了.
為什麼.
因為你阻礙地球轉.
上帝國要臨到.
為什麼你不守安息.
為什麼你不做神的子民應該做的事.
這個叫做大發熱心去守律法.
就是因為他們盼望著上帝的國臨到.
上帝的國臨到的時候.
全國的世界就會改變了.
新天新地就會來到.
然後身體也會有復活.
所以你明白 為什麼保羅這麼緊張.

$^{441}$這麼大發熱心.
不是因為他覺得他是上帝的子民.
他們那班人是上帝的子民.
我們就應當如此行.
並且我們要大發熱心.
所有其他我們同樣種族的自己人.
我們都要去守律法.
有這樣的背景 你想想.
突然間有基督教會出現.
初期的教會是什麼人多.
猶太人 或者我們今天叫做Messianic Jews.
那些信耶穌是彌賽亞是救主的.
那班猶太人叫做Messianic Jews.
突然間有教會出現.
自己人 他們相信耶穌.
相信耶穌是救主彌賽亞來到.
在十字架上受苦.
這個已經是承諾了.
因為我們的救主沒理由這樣.
所以對猶太人來說是一個絆腳石.
更加就是他們信了耶穌之後.
他們繼續去擴張.
他們傳一個死而復活的耶穌.
對於當時的超正統猶太人.
或者保羅自己來說.
這個更加不能夠去接納.
所以你這樣的背景就明白.
為什麼這個訴羅不是因為撩交打.
他不是因為真的想去殺人.
而是他為了上帝的國大發熱心.
為了他自己本身所信的東西.
縱使他現在知道是錯誤.
他去大發熱心.
所以他盡他的能力.
就是什麼.
要去停止教會的擴張.
這個是史萊姆傳第九章.
如果你知道他做了什麼的時候.
你知道之前發生什麼事.
第七章就是史提反.

$^{481}$第一個猶太的基督徒.
是被殉道的.
人們用石頭打他.
原來保羅又在中間.
所以他是其中一個大佬.
他去允許這件事發生.
但是八章交代了其他東西.
第九章他就出發.
你看著地圖.
他從耶路撒冷向北邊走.
當然他不會直接向北邊走.
因為北邊是什麼人.
北邊是撒瑪利亞人.
懂一點背景的話.
總之撒瑪利亞人和猶太人是不合的.
不會一起吃飯.
也不會見到大家.
也不會去大家的地方.
所以作為一個超正統猶太人.
他就繞過北邊的撒瑪利亞山地.
然後上到上面.
北邊的加利利海附近的加利利.
然後再攀上去.
我們今天叫做高蘭高地的地方.
高蘭高地攀山越嶺.
沿過黑門山.
然後一直向東北邊.
去大馬士革.
就是今天敘利亞一個很大的城市.
全程大概150英里.
一六天時間.
我相信他應該很趕路.
很趕路.
但我有一點點知道保羅的身世.
或者保羅自己是法治創人的背景的時候.
你就知道他一心一意想做什麼.
去大馬士革.
奇怪為什麼去大馬士革.
為什麼不在耶路撒冷.
耶路撒冷也有初期教會的信徒.

$^{521}$因為你知道在史坦傳的交代裡面.
那個技術裡面.
就是初期教會.
史提反殉道之後.
就四散.
因為大受迫害.
所以那些基督徒就到處跑.
跑去避難.
或者跑去其他地方.
他就帶著福音去其他地方.
你知道大馬士革是什麼地方.
大馬士革是當時的紐約.
當時的紐約.
是世界的通道.
當然加拿大已經是世界通道.
但去到大馬士革是一個大的城市.
一個在中東東南西北的通路.
如果你到了大馬士革.
你就可以去到今天的伊拉克.
當時其中一個文化的中心.
就是米索不達米亞.
如果轉左的話.
你就去到今天的土耳其.
就是小亞細亞.
或者去到歐洲.
那是一個很重要的十字路口.
在地理上來說.
可以這麼說.
大馬士革就是地極的開始.
記得史坦傳第一章.
耶穌吩咐他們門徒.
從耶路撒冷去哪裡.
猶太傳遞.
撒瑪利亞直到地極.
所以他們四散的時候.
這些基督徒相信有很多去了大馬士革.
去一個很重要的城市.
然後他們帶著道就去其他地方.
所以如果你想一下.
你是一個法理財人.

$^{561}$一個大發熱心的人.
你盼望上帝的國會臨到.
復活的世界會臨到.
你第一步要做什麼.
你就不是留在耶路撒冷.
你就要趕緊的.
好像要去.
剛剛COVID的時候.
人們不是很留意COVID的擴張散佈嗎.
經常看著世界地圖.
哪裡擴張到哪裡.
快要來到美國了.
為什麼不封關.
我們那時候很害怕.
這個就是保羅當時在想的.
那些基督徒到處擴散.
好像病毒一樣.
不行.
一定要去哪裡.
一定要去他們最前線.
最front line的地方.
來將基督教阻截.
你不阻截的話.
就很大件事.
因為上帝的國是不會臨到.
我們知道.
為什麼這麼遠的水路.
要走到那裡.
在保羅的mindset.
或者在他的心理裡面.
在他當時的概念裡面.
你知道猶太教.
可以這樣總結.
當時的猶太教.
起碼法理上.
可以有三個T來總結.
第一就是時間.
第二就是律法.
第三就是聖殿.
什麼叫時間.

$^{601}$就是說他們守安息日.
是很重要的.
節日是很重要的.
這個不只是他文化身份的一部分.
也是他認為要嚴守律法的時候.
然後帶到上帝的國.
第二個不需要說了.
就是律法.
摩西律法.
我們今天的舊約聖經.
希伯來聖經.
或者那些口傳的律法.
今天到現在都有很多這樣的口傳律法.
在聖經裡面沒有說的.
比如潔正詞等等這些.
舊約聖經哪有說的.
但是從哪裡來的.
是拉比的口傳的傳統.
那些tradition那裡來的.
第三個更加重要的.
就是聖殿.
聖殿在哪裡.
耶路撒冷.
在他的對世界的理解裡面.
最聖潔最神聖的地方在哪裡.
就是你給他地圖.
他就可以GPS讀到給你聽.
這裡就是最聖潔.
哪裡最聖潔.
就算世界各地不同的猶太人都是這樣.
他向著耶路撒冷祈禱.
向著耶路撒冷祈禱.
我記得有一次坐飛機的時候.
有些團友問我.
正在去塔拉雲湖.
他說為什麼他們在飛機上走的時候.
不需要再sit belt.
他們就走去後面廁所那裡.
就是空位的地方.
他就向著某個方向去祈禱.

$^{641}$或者等著上機的時候.
他們又向著玻璃去祈禱.
其實他就向著耶路撒冷祈禱.
為什麼向著耶路撒冷祈禱.
而不向芝加哥.
那有什麼好呢.
就是pothole.
pothole多.
tax多.
corruption多.
這三個最出名的.
因為這裡不是聖城.
這裡不是神聖的地方.
耶路撒冷就是上帝所居住的地方.
temple.
是最聖潔的地方.
最近最近向著的地方.
這個就是最神聖.
神所在的地方.
神在那裡出現.
神在那裡顯現.
所以如果你這樣想的時候.
你就發現在保羅的世界裡面.
其實有幾個圈.
第一個圈最中心的圈.
就是聖殿.
知聖所.
然後外面那個圈.
就沒有那麼神聖.
猶太傳地.
再外面一點.
就更加是世俗.
撒瑪利亞.
不同種族的人.
跟我們不合的那群人.
再外面一點就是那些所謂的外邦人.
更加跟我們是拉不上關係的.
那些就是profane.
由sacred to profane.
就在一個這樣的mindset裡面.

$^{681}$你就可以理解到.
當保羅離開耶路撒冷.
他不是一個羅馬的公民.
羅馬公民就是出境入境.
他拿著護照去別的地方.
好像我們去別的地方.
別的國家.
保羅不是單單一個這樣的心態.
一個正統的猶太人.
他是離開最神聖的地方.
去到敵人的地方.
敵人的地方.
每個地方都有一個故事.
每個地方都有它的story.
特別是聖經的地方.
大馬式是什麼地方.
在舊約到新約.
和大馬式的人.
或者亞蘭人舊約.
一直都沒有好過.
經常都打仗.
所以保羅是正在進入一個敵人的地方.
他更加不會預期在那裡見到上帝.
你想想保羅如果去越來越遠的時候.
到時候祈禱他會向哪裡祈禱.
不是向前面.
是向後面.
因為耶路撒冷在他後面.
神的榮耀沒有理由在那裡出現.
但是就是復活的主基督.
在那裡向他顯現.
你仔細想一想在基督的復活之後.
其實有很多這樣的故事都和這些有關.
譬如路加福音第24章.
兩個門徒傷心欲絕.
他們覺得耶穌死了.
不會復活的.
怎麼辦.
他們就離開耶路撒冷去二馬五師.
二馬五師是耶路撒冷的西邊.

$^{721}$離開耶路撒冷.
誰知道走著走著的時候.
誰和他們走.
就是耶穌和他們走.
耶穌居然不是在耶路撒冷和他們出現.
一個榮耀的上帝不是在耶路撒冷出現.
而是在四周出現.
在更加不會預期祂出現的地方.
祂就出現了.
所以這是第一點.
第二點就是一個根本不會預期.
在他們生命裡面出現的一個人.
就是保羅自己.
剛才我說了一點保羅的背景.
大概我們都猜到了.
沒理由這個人會相信耶穌.
沒理由.
in fact.
當保羅真是信主之後.
第九章之後.
史瓦蘭傳.
他去到耶路撒冷的教會.
那些人信不信他.
比我都不信.
那裡的人都在顫抖.
因為這個應該是艷仔.
滲進來這裡想做些什麼.
結果因為有一個門徒叫阿拉利亞.
他親自去接他進來.
然後巴拿巴親自接他進來.
款待他進來.
他才能進到教會.
很自然.
你想想你身邊的人.
可能有很多都沒什麼機會信主.
你覺得對嗎.
你可能會有一群朋友.
你可能會這樣分級數.
可能對信仰有些開放.
有些興趣.

$^{761}$OK 他們可能會近一點.
我們很自然會這樣分.
有些完全是和你反對的.
對信仰是反對的.
叫他教會永遠都不會來.
那些根本沒可能.
這個故事發生.
偏偏相反.
在一個不配的生命裡面就出現.
扭轉了這個人的一生.
也扭轉了基督教歷史.
上帝救贖歷史的所有一切.
新的一章就開始.
在這個根本是不配的人身上.
復活的主就向他演.
所以為什麼保羅最後說他不配.
如果你看看小林傳的脈絡.
剛才我們交代了一點.
第七章說什麼.
第七章說史提反的殉道.
史提反殉道有什麼好說.
其實不是殉道那一刻有東西說.
而是講道那部分有東西說.
因為是一篇很長很長的道.
很長的道.
但是說什麼呢.
很長的道.
你不用回去聽.
我今天告訴你一分鐘.
史提反由頭到尾說.
其實就是一個救恩歷史的故事.
就是說上帝在救約整個歷史裡面.
在哪個地方出現過.
他首先在米索不達邁瓦講.
在吾人呼召阿巴罕的時候.
榮耀的上帝向他顯現.
然後在曠野那裡向摩西和馬爾穆文顯現.
然後一直講下去.
但是很奇怪.
史提反就是沒有說什麼.

$^{801}$神的榮耀在耶路撒冷出現.
這個他隻字不提.
在猶太人眼中.
這個是夢寐以求的good old days.
希望能夠回去第一聖殿被獻的時候.
獻殿的時候.
所有人獻殿的時候.
那種威榮.
上帝的榮光充滿整個的聖殿.
耶路撒冷的聖殿.
然後地都震了.
上帝的榮耀淋到我們中間.
這個是good old days.
但是史提反卻沒有說.
還加了一句什麼呢.
他說你知道嗎.
至高者其實是不用住在人手所造的殿.
其實他就指著那殿.
所以那些正統猶太人.
保羅包括.
是不是很討厭這個人.
講哪一套講到這樣呢.
不單止是一篇爛道.
當然是對我們挑釁的道.
於是他們就拿起石頭就扔死他.
但是你看到史提反所講的是什麼.
至高者.
神的榮耀.
不是限制於在耶路撒冷.
在你心裡面.
那個所謂sacred place.
四周圍都可以顯現.
因為基督已經復活.
第八章.
講到費力.
去到傳道.
去哪裡傳啊.
不是耶路撒冷.
叫他去撒瑪利亞.
仇家的地方去傳道.

$^{841}$然後有些人信了.
然後又叫他去哪裡.
去南邊.
然後那個是誰.
埃塞俄比亞的人.
那個是一個大權的太監.
這個故事很奇怪.
有什麼好講的呢.
是不是埃塞俄比亞特別一點.
最特別是什麼這個故事.
是太監.
太監.
太監就是身體有殘缺的.
他們是不可以進聖殿的.
也不可以近聖殿.
因為身體有殘缺.
但是卻信主受洗.
不用回耶路撒冷.
他繼續回去.
所以有些不配的人.
一直一直來信主.
到第十章.
就是彼得.
彼得就去到蓋薩尼亞.
一個羅馬人的海港.
就遇見哥尼留·庫內利安斯.
這個人.
這個是羅馬的公安.
羅馬公安.
和當時的猶太人也好.
基督徒也好.
都不是很合得來.
有什麼事.
我們不會派出所有人去喝咖啡.
彼得就被上帝差派去.
去到最大的派出所.
就去找哥尼留.
就是帶他去信主.
你看到一個一個不配的生命.
是見到上帝的榮耀.

$^{881}$其實想想我們自己也是一樣.
我們都是難民一條.
有誰想過我們配受上帝的恩典.
但是上帝他向我們顯現.
所以保羅最後說什麼.
他說我是不配.
是未到產期而生的人.
使徒當中我根本就不配做使徒.
曾經迫害過教會.
今天我為什麼會在這裡.
為什麼會格外勞苦.
因為上帝的恩典.
基督的復活就是開展了一個新天新地的工作.
新的創造就來到我們中間.
臨到我們是產生什麼的分別.
產生他追求上的分別.
這個保羅以前是嚴守律法.
並且是以為用他的方法能夠帶上帝的國來到.
基督的復活就扭轉一切.
然後基督就給他一個新的使命.
復活的基督在我們生命裡面.
有沒有開到新的一章.
我們敢不敢給復活的基督.
在我們生命裡面開一個新的一章.
因為我知道我們正在經歷很多不同的患難.
不同的問題.
但是基督的復活.
你想一下如果是基石的話.
就應該會起碼方向上已經扭轉了我們.
叫我們追求的是什麼.
以前可能追求的就是可能是物慾.
可能是一些世界上很掛慮我們的東西.
情慾的東西.
祈禱的都是那些.
希望上帝給我這樣給我那樣.
但是基督的復活應該給我們一個新的方向.
就是向著復活的主.
可以為他做什麼.
我記得以前在以色列.
我認識一個.

$^{921}$在耶路撒冷認識一個姐妹.
那姐妹信了主.
她的丈夫是猶太人.
你知道猶太人是不會進教會的.
對耶穌和基督教會都沒有什麼好感.
都是彈開的.
誰知道跟她結了婚.
Anyway.
她說很有趣.
她的丈夫為什麼可以帶她回教會.
Eventually帶她回教會.
原來.
她丈夫說.
就是這個太太.
她每個星期去完教會之後.
她很爽.
她的生命不同了.
她覺得教會真的可以做到一些事情.
在她老婆身上.
Something different.
Right?.
我不知道我們是不是這樣.
我們是不是去完教會之後.
我覺得.
終於有個差事辦完.
太好了.
我Go on with my own life.
真是麻煩.
又要去教會.
然後又要聽這個人講道.
我們去敬拜上帝.
是應該給我們一個.
生命上的一個轉變.
多多少少.
基本上.
有一個轉變.
他就是看到他太太.
回來的時候.
照顧孩子.
煮飯.

$^{961}$每樣東西都特別好一點.
沒有煮東西吧.
這個就是.
我們生命裡面.
應該有一種.
可以說服到人的.
印象.
說服到人.
英國有一個護教學家.
叫做Alister McGrath.
他是專門在科學護教.
他說什麼呢.
他說我們作為福音的代表.
作為全福音.
我們不是.
只是叫人們決志填那個科目.
That is not the point.
那個point是什麼.
他說我們作為基督徒.
作為福音的代表.
就是要給人們.
我們的目標.
是要給人一個.
對福音有說服力的生命.
convincing testimony.
人們看到我們的時候.
或多或少都會看到.
基督復活的真相.
是應該是這樣的.
讓我們的生命.
被基督的復活.
有多少繼續改變.
好嗎.
讓我們能夠convince其他人.
這位復活的主.
是一位活著的主.
我們同心抵禮基督.
好嗎.
上主我們多謝你給我們復活的盼望.
因為耶穌基督是一位活著的主.

$^{1001}$是因為生命當中活著.
主你求你幫助我們.
我們今天來到敬拜你的時候.
我們將我們一切的疑難.
一切的憂慮.
冒出一些問題.
很多時候我們未能夠解決.
但你卻告訴我們.
我們生命已經開展了新的一頁.
新的一章.
是因為基督的恩典.
復活的盼望.
主差引我們去到不同的地方.
讓我們去服侍你.
去代表你.
我們這樣祈禱奉祝明天.
剛才我們報告的時候.
說到有一個報道會.
我剛才沒有說到.
我最重要的目的是做什麼.
不是叫人去缺資.
有些人可能你想請他來報道會.
他就灑手露頭.
為什麼呢.
因為很多時候我們給人的印象就是.
你去簽名.
你舉起你的手.
我想告訴你.
你叫你的朋友來報道會的時候.
我是不會叫他做這件事的.
你說不是吧牧師.
你這麼不熟練.
我就是不會.
因為這不是重點.
重點是什麼.
是我們應該有一個.
有信心的印象給人.
所以你可以告訴你的親友.
不要緊.
你不會被迫去做某件事.

$^{1041}$享受一個音樂的晚會.
聽到主的福音.
我希望他心靈能夠得到改變.
希望我們祈禱裡面去紀念.
這個報道會.
上帝祝福各位.
\newpage



\section{}
\label{sec:ixAcYGJ3Y7Y}
\textbf{比對研究 13: 巴力宗教 6 - 陳崇基牧師}
\newline
\newline
連結: \href{https://youtube.com/watch?v=ixAcYGJ3Y7Y}{\texttt{https://youtube.com/watch?v=ixAcYGJ3Y7Y}} ~~~~ 語音日期: 2023-05-06
\newline
\newline
\hyperref[sec:upXqQwuuS0M]{\small{< < < PREV SERMON < < <}}
~
\hyperref[sec:index_chronic]{\small{[返順時目]}}
~
\hyperref[sec:index_scriptual]{\small{[返順卷目]}}
~
\hyperref[sec:hS4BoW8e_54]{\small{> > > NEXT SERMON > > >}}
\newline
\newline
$^{1}$於是我們就知道17章第一節發生了什麼事.
就是一個先知的回應.
上帝回應阿訶王做的事.
以利亞在基列住的提比斯人.
提比斯的字大概意思是寄居者.
不過我們音譯就叫做提比斯.
其實就是寄居者.
在基列的寄居者 以利亞.
他奉差派對阿訶王說.
就是上帝說我收起怒水 收起雨水.
這個我剛才說了 生命的背景.
就是說這些就是咒詛.
當我們不在上帝的約裡面.
不忠於祂的時候.
很多時候約裡面所講清楚的咒詛就會淋到.
你不再活在耶和華的同在下.
神的同在下 你失去了上帝的同在.
在以前迦南地的人來說.
最重要的一件事.
沒有上帝的同在的話.
就是沒有怒水也沒有雨.
所以有幾年的時間.
這樣的情況就自然會有旱災出現.
耶和華的話就淋到以利亞.
首先上帝跟先知說.
沒有水 旱災來到.
不過我會養活你.
他說你離開這個地方往東去.
當時他在伊斯利國.
跟阿訶講完話.
大概應該是在京都.
他就叫他去約旦河東的地方.
就是他回到他本來來的地方.
基列就在這裡.
所以他就離開了應許之地.
講清楚了.
約旦河其實就是迦南地東邊的邊境.
叫他去約旦河東的基立溪旁.
叫Reef.
不知道在哪裡.

$^{41}$不過總之就是約旦河東邊.
這個最為重要.
因為他就離開了迦南地.
離開了迦南地還有沒有上帝同在.
應該是沒有.
不過在先知上.
在以利亞的身上發生的就是.
上帝居然在他的境外.
就是在迦南地以外.
他仍然是控制著水和食物.
換句話說他是誰.
他是全地的主人.
這裡已經講清楚了.
他不是巴力.
巴力就有地頭.
譬如烏加拉.
在薩峰山.
Mount Saffong.
但是耶和華沒有的.
去到哪裡都好.
他都是可以供應.
還有是操控著水和食物.
所以有一個這麼可愛的故事.
就是赫那溪裡面的水.
已經吩咐了烏鴉在那裡供養你了.
上帝就用不同的方法來供養他.
於是以利亞就去了.
遵照上帝的話做.
就住在約旦河東的基納溪旁.
早上有餅和肉.
晚上也有餅和肉.
然後他也喝溪裡面的水.
直到喝乾為止.
這裡就講清楚了.
上帝控制了所有的東西.
這是第一個故事.
以利亞.
然後那裡的水乾了.
發生什麼事呢.
雨沒有再下在地上.

$^{81}$耶和華就臨到以利亞.
他說你起來.
這次去哪裡.
就去西頓的薩拉法去.
薩拉法.
西頓的地方.
住在那裡.
這裡奇怪了.
他本身是這裡的人.
奉差派做先知去對付.
以利亞就說.
上帝就跟他說.
你住在約旦河東.
以色列的認許之地之外.
但我在那裡供養你.
我很信實.
我仍然去控制所有的東西.
然後那個故事完了.
就叫他去哪裡.
就叫他去西頓的薩拉法.
西頓附近的薩拉法.
居然叫他去敵邦的故土裡面.
進入他的故土.
這裡肯定就是以色列的認許之地以外.
叫他去西頓的地方做什麼呢.
在那裡他就吩咐了一個寡婦來供養他.
他就起來去了那個地方.
他也在境外.
所謂迦林地的境外.
上帝很奇妙地又去供應他.
在城門口他就遇見了一個寡婦去收拾.
然後就跟他說.
你用器皿給我喝水.
他就拿了水給他.
又呼喚他.
你可不可以拿些餅給我吃.
現在是旱災.
所以那個寡婦就跟他說.
我指著永生的耳環.
你的神和氣勢.

$^{121}$我沒有餅.
潭內也只有一把麵.
瓶裡面只有一點油.
看我去撿兩斤柴.
回來我和我的兒子做餅.
我們吃了之後就等死了.
因為剩下很少很少東西.
已經山窮水盡的時候.
以利亞就跟他說.
不需要怕.
你照著我的話去做.
正如之前以利亞照著耶和華的話去做.
耶和華就恭應他.
現在他就吩咐寡婦.
用他的信心照著他做.
先給我一個小餅.
拿給我.
然後就為你和你的兒子做餅.
因為什麼.
神說什麼.
潭內的麵必不用盡.
瓶裡面的油必不短缺.
直到耶和華始於降在地上的日子.
就是說這段時間裡面.
一定會給你小麥.
大麥.
麵粉.
橄欖樹.
出橄欖.
收成了.
有油了.
有橄欖油用.
是不是?.
這些全部都是甘地的氣候.
或者農耕的 cycle(循環).
直到耶和華始於雨降在地上.
就是用神蹟去給他油和麵粉等等.
他就照著以利亞的話去做.
正如以利亞照著耶和華的話去做.
是不是?.

$^{161}$塔羅和以利亞和他家中的人就吃了許多日子.
潭內的麵果然沒有用盡.
瓶裡面的油也沒有短缺.
正如耶和華.
藉著以利亞所說的話.
這裡提了什麼呢?.
一段故事我們可以不理會背景.
照樣可以形容.
上帝很信實.
去到公認等等.
裡面提到雨水或者麵粉,油等等收成的東西.
我們大可以不理會.
也可以立刻形容.
總之就是上帝的信實.
但是有沒有背景可以幫助我們多一點理解經文.
一些更加重要的主題呢?.
巴力的背景.
巴力宗教.
跟這些有什麼關係呢?.
我們看一段巴力的文字.
這個史詩.
是講到一個烏加尼的王.
叫做King Curta.
所以這個史詩可以叫做King Curta的史詩.
這個史詩上下文大概是說什麼呢?.
雖然有一些子女行家.
因為禮版已經沒有了.
但是我們在上下文裡面推理.
我們大概知道.
King Curta.
就是國王.
他患了病.
患了病也有旱災和饑荒出現.
有饑荒有旱災出現.
在巴勒斯坦的稀有裡面是怎麼樣的呢?.
就是因為沒有水.
這段史詩講什麼呢?.
就是祈禱.
天上的大神.
其中一個.

$^{201}$很多神.
其中一個神叫做El.
或者其餘的另一個神.
祂說什麼呢?.
哪個神也好.
不重要.
說什麼才重要.
說什麼呢?.
從第二節開始.
祂說要察看大地和天空.
要察看地的遠極處.
有沒有墨紙呢?.
能不能找到墨紙呢?.
因為有旱災.
那麼祂說要仰望誰呢?.
看看地上.
希望有巴力的雨.
就是地上因為沒有水.
希望巴力會降雨.
因為沒有降雨的時候又有旱災.
所以祂想找什麼呢?.
巴力的雨.
地上也是一樣.
希望有至高者的雨.
巴力的雨.
祂就說.
對大地有多麼多好處.
因為有雨下就會有收成.
是不是?.
對於田地來說.
至高者的雨水是多麼好.
所以這裡提到雨水.
也提到墨紙.
還有什麼呢?.
祂說巴力的雨實在是好.
因為為了什麼是最好呢?.
就是耕地的墨紙.
就是泥溝裡面的墨紙.
田間裡面的小墨罐.
那些小墨水漲出來的墨水.

$^{241}$祂說祂們舉手.
農夫也舉手.
大地好像現在是呼喚著.
很渴望.
很仰望著天上的神能夠降雨.
農夫也是一樣.
向天上去仰望.
仰望誰呢?.
就是巴力.
巴力另外一個綽號就是大菌的僕人.
大菌你記得在聖經舊約.
譬如《士師記》.
參書人故事有出現過.
達加恩.
就是那個….
其實祂就是墨紙的神.
墨田的神.
祂說為什麼要向他們呼喚呢?.
為什麼要呼喚大菌的僕人巴力.
能夠給他們水.
能夠給他們水之後.
能夠有墨水出現呢?.
因為糧食.
從儲存的地方.
都已經吃光了.
油.
酒也都沒有了.
酒在皮袋裡面已經乾了.
皮袋都乾了.
當時釀酒就是放在皮袋裡面的.
桶裡面也沒有油.
全部都用盡了.
如果你這樣看的時候.
一個這樣的背景.
你就留意到.
我們剛剛說的兩個故事.
都是有這樣的主題的.
包括這個寡婦的故事.
正如我所說.
你將所有東西.

$^{281}$全部背景拿走.
都能夠應用的.
是不是?.
就是上帝信實.
是不是?.
但是在這裡我們就知道.
為什麼這個故事會這麼寫.
還有為什麼主要.
講出幾件事情.
就是麵粉.
是不是?.
還有.
瓶裡面的油.
是不是?.
油用不盡.
麵粉用不盡.
為什麼呢?.
因為耶和華.
不是巴力.
是耶和華.
去操控著.
掌管著.
這些的.
墨紙.
還有田間裡面的出產.
就不是巴力.
為什麼我們知道這樣呢?.
因為如果你這樣去對比之下.
你就知道.
巴力就是一個這樣的神.
就是管理雨水.
管理墨紙.
小麥罐.
是不是?.
糧食.
食物.
酒.
還有油.
是不是?.
跟當地的農耕.

$^{321}$很有關係.
以利亞的故事.
就想講什麼呢?.
起碼開頭這兩個故事.
就是告訴我們.
不是巴力去控制我們.
要的麵粉.
油.
墨紙.
也不是他去控制雨水.
是耶和華去控制.
我們回到剛才巴力的背景.
我們再想想.
有了這個背景的時候.
我們就了解了.
為什麼那些小小的神蹟的故事.
是這樣發生.
因為上帝藉著以利亞的經歷.
來講出.
祂是一個比巴力更大的神.
來證明巴力是假的.
祂才是真實.
我們看到兩個故事.
現在想下去.
看看有沒有其他以利亞的故事.
跟這些主題有關.
相信這些背景.
對比能夠幫助我們了解.
這個故事真正的意思.
\newpage



\section{17:13}
\label{sec:hS4BoW8e_54}
\textbf{Go and Do: 1 Kings 17\_13 - Tommy Lee}
\newline
\newline
連結: \href{https://youtube.com/watch?v=hS4BoW8e_54}{\texttt{https://youtube.com/watch?v=hS4BoW8e\_54}} ~~~~ 語音日期: 2023-05-09
\newline
\newline
\hyperref[sec:ixAcYGJ3Y7Y]{\small{< < < PREV SERMON < < <}}
~
\hyperref[sec:index_chronic]{\small{[返順時目]}}
~
\hyperref[sec:index_scriptual]{\small{[返順卷目]}}
~
\hyperref[sec:UpJWRYxxqyc]{\small{> > > NEXT SERMON > > >}}
\newline
\newline
$^{1}$Help me.
Sound good.
All right.
Thank you very much.
Mark thank you very much.
for inviting me back again.
You know.
It's been a couple of weeks.
since I was last here.
And I will probably say.
After I spoke to you guys.
I have been literally.
been on the road.
for the last couple of weeks.
I think I've been more.
in hotel rooms.
and traveling on planes.
than I actually have been at home.
And so it is good to be back here.
at least for the next week.
that I'm off to Seattle this week.
that I'm back in Texas.
and Los Angeles.
for the next couple of weeks.
before finally being home again in June.
We've been starting this series.
on the book of.
the first Kings.
where we looked at the character Elijah.
Elijah when we first began.
in first Kings 17 1-7.
We saw that Elijah literally means.
that our introduction of Elijah was this.
We know that his name literally means.
My God is Jehovah.
DL means My God.
Jah means Jehovah.
I is the connecting point.
My God is Jehovah.
We knew that Elijah is a Tishbite.

$^{41}$from Tishbe in Gilead.
What does that mean.
Merely that.
they have no idea where he lives.
It's somewhere in Gilead.
and he lives in obscurity.
And so this guy who comes from a rural town.
who's kind of rough on the edges.
as you will see.
whose name literally means.
My God is Jehovah.
comes over to King Ahab.
And King Ahab is the king of the northern kingdom.
If you guys remember the back story.
There is a northern kingdom.
during that period of time.
and a southern kingdom.
The southern kingdom is where.
Jerusalem, Nehemiah, Ezra.
All of those guys take place.
Elijah takes place in the northern kingdom.
And in the history of the northern kingdom.
there are 19 kings who rule.
And none of them are good.
All of them are bad.
King Ahab is the worst of all of the kings.
And then there's Queen Jezebel, his wife.
who a lot of times.
when you look at scripture.
Queens are rarely ever mentioned.
Here's a king and here's a queen.
But Jezebel is mentioned.
You know why.
Because Jezebel will turn out to be the worst.
of all of the kings.
including Ahab, including the 18 kings.
She's this bad.
Well, Elijah suddenly gets.
a meeting with the king.
How a guy from Tishbe.

$^{81}$gets a meeting with the king.
who comes out of obscurity.
I don't know.
but he gets a meeting with the king.
And he says that neither dune nor rain.
except by my word.
And he says this in verse 2.
As the Lord the God of Israel lives.
before whom I stand.
That becomes the theme of Elijah's life.
That becomes the theme of this particular passage.
That's Elijah's motivation.
before whom I stand.
That's why he does the things that he does.
He tells them that there will be.
neither dune nor rain.
And just like that.
right in the middle of the action.
God takes him away to Cherith.
And now he lives in Cherith.
He doesn't just stay there.
He lives there.
And he's fed by ravens.
And he's staying there.
until the brook suddenly dries up.
And suddenly in the last time that was here.
which will eventually lead us.
to what we're talking about today.
The word of the Lord came to him.
and says arise go to Zarephath.
So you've been to Cherith.
And now you're going to Zarephath.
which belongs to Sidon.
and dwell there.
Behold I've commanded.
a widow there to feed you.
All right so God you want me to.
When I go against Ahab.
You want me to go to Cherith.
I went to Cherith.

$^{121}$It's nice.
I know you're protecting me.
Ravens fed me.
Brook fed me.
So now I'm happy.
Now you want me to go to Zarephath.
Zarephath literally is a hundred miles.
You want me to walk a hundred miles.
Do you not realize.
that Ahab is coming after me.
Do you not realize my life is in danger.
Yes.
And by the way.
which I never told you this.
Because I found out.
Right before.
Right after I came and spoke.
Zarephath is the birth town.
In the hometown of Queen Jezebel.
So you want me to walk a hundred miles.
Be exposed to danger.
And you want me to go to Jezebel's hometown.
Yes.
And someone's going to be there.
to take care of you.
Great.
He walks a hundred miles.
He risks the danger.
And there he has seen a widow.
Who when he meets in verse 12 says.
As the Lord your God lives.
Here's what she says to him.
He meets a widow.
As the Lord your God lives.
She's not even of the faith.
She doesn't even belong in the faith.
Your God.
I have nothing baked.
I only have a small handful of flour.
And a little oil.

$^{161}$And a jug.
So I have a little baked.
I have a little flour.
I have a little oil.
Really.
And now I am gathering a couple of sticks.
That I may go and prepare it.
For myself and my son.
That we may eat it and die.
Lord.
That's what you're sending me to.
You made me walk a hundred miles.
To the birthplace of the Queen Jezebel.
Who is the worst of everyone.
To meet a widow.
Who has no flour.
Little oil.
And she thinks she's about to die.
Yes.
And as you remembered.
For those who are here.
For those who are not here.
Quickly we talked about three things.
Three things that you need to know.
And two things that you need to overcome.
Number one.
The Lord sees you.
The Lord sees Elijah.
The Lord knows where Elijah is.
He knows that he's been in Tishb.
He knows that he went to see Ahab.
He knows that he was a chariotee.
He knows that he's in Zarephath.
The Lord sees you.
Number two.
The Lord knows where you need to go.
You need to go to Zarephath.
You just need to go.
It reminds us of Joshua chapter three, right?.
It's after forty years in the wilderness.

$^{201}$To roam around the wilderness.
One day the Lord says to Joshua.
Tell everyone to pack up.
After forty years.
Pack up.
We're going to cross the Jordan River.
But follow me.
I love this verse.
Follow the ark of the covenant.
Because you have never been there before.
And you don't know where to go.
Follow the ark of the covenant.
Because you don't know where to go.
Because you have never, ever been there before.
Not only does the Lord see you.
The Lord knows where you need to go.
But the Lord knows what you need.
Now it may not be what we think we need.
It may not be our expectations of what we need.
But the Lord knows what we need.
He encounters a widow.
Doesn't have a lot but just enough.
He encounters a widow.
And there's two things.
That you need to overcome.
You're going to have to overcome.
First impressions or expectations.
That when we think that.
When we follow the Lord.
There's all these expectations.
Of what God should provide for us.
If we're faithful.
If we're faithful to you.
Here's what I expect out of you.
And we have these expectations.
And when these expectations.
Don't come through.
We get a little frustrated.
We get a little annoyed at God.
And we soon lose trust in God.

$^{241}$So you have to overcome.
These first impressions.
This widow and these expectations.
And trust that God.
Will open those doors as you see fit.
Mark mentioned that.
Some of you guys are graduating.
I was just talking with Rachel.
As she's graduating as well too.
All your life from first grade to second grade.
From third to fourth.
From seventh to eighth.
From eighth to freshman year.
Then from that point you go to college.
You know where you need to go.
Bam, bam, bam, bam, bam.
Upon the graduation from college.
It will be the first time.
College graduates look at their life and say.
There are five roads I could take.
There are multiple job offers.
There are master programs that I could take.
I don't know where to go.
And we think that if we take one passage.
If we think that's what God is calling us to do.
Everything will be okay.
And God will provide.
It may not happen.
So you have to overcome.
First impressions or expectations.
And then you're gonna have to overcome.
Physical impossibilities.
He, Elijah, had to overcome.
Physical impossibilities.
He had to overcome the very fact.
That this particular widow.
Didn't have very much.
But then he says to her.
Which I love and which a lot of times.
We actually didn't have enough time to spend on.

$^{281}$Because so today I want to spend time.
On this one verse.
Really unpack it.
And see what it means to me.
He says to her.
Go and do.
You're gonna have.
And here's a.
I'm losing it.
Oh Lord, here we go.
Here we go.
Elijah says to her.
Do not fear.
Go and do as you have said.
Even though you have seen.
All of these impossibilities.
Even though you have seen.
All these expectations.
That are out the window.
He says to the widow.
Go and do.
And that's a calling for us.
As Christ followers.
For us to go and do.
Three things that we're gonna talk about today.
What does it mean for us.
As we look at this particular phrase.
For us as Christ followers.
To go and do.
What does it mean for us not to go and do.
What does that mean for us not to go and do.
What do we do.
As Christ followers.
Three things we're gonna talk about.
What does it mean for us to go and do.
Number two, what does it mean not to go and do.
And what does it mean for us to actually do.
Those three things.
I told you the last three weeks.
I have literally been on the road.

$^{321}$Went to Minnesota for spring break.
With the family.
Went over to Romania.
To the city of Bucharest and Cluj.
Then from that point.
I went over to Prague.
Over the Czech Republic.
And I ended up over in Orlando.
With a bunch of our cohort.
From around the different worlds.
When I went to Romania.
I arrived in Bucharest.
The capital city of Romania.
On Orthodox Easter.
Where the city was completely empty.
It was completely deserted.
Why.
Because it is a huge national holiday.
Did you know in Romania.
90\% of all the people.
Who live in the country of Romania.
Will say they're Christians.
Out of that 90\%, 80\% of them.
Are Orthodox Christians.
Romania as a country.
Is going through a huge rebuild.
Like when I stepped foot in Romania.
I've never been to Romania before.
Some of you guys.
Who have probably studied.
Eastern European history.
That's the only time.
I've ever heard about Romania.
It was embedded in communist regime.
All the way to 1989.
It was ruled by this general secretary.
By the name of Nikolai Ceaușescu.
And he ruled it with an iron hand.
And after his overthrow of the regime in 1989.
Everything started rebuilding.

$^{361}$But it got to a point.
If you know anything about.
What was going on in Romania.
This guy Nikolai Ceaușescu.
Was so obsessed.
About creating his own kingdom.
He ignored the basic necessities.
Of the people there.
He ignored the health of the country.
He ignored the needs.
Until it got to be so much.
They soon overthrew him.
On one of those days.
We visited his former palace.
It is now the parliament.
For where the congress.
And the senate takes place.
When you go into this former palace.
It is humongous.
We walked around.
It took me 25 minutes.
Just to walk around the outer perimeter.
Because it's Orthodox Easter.
So there's no security guard.
I'm sitting there.
Why is there no security guard.
Around the parliament.
I had no idea.
And so you had to walk all the way.
And we kept getting lost.
We finally ended up in this place.
To take a tour of this place.
It is so big.
It is the second biggest.
Government building.
Next to the pentagon.
So of all the different.
Government buildings.
Around the world.
This palace was the second biggest.

$^{401}$They told us that it takes.
\$500,000 a month.
To power this thing up.
To pay for all the utilities.
And you know what was ironic.
It was empty.
75\% of the building.
Is completely empty.
It was deserted.
And they said unless.
Congress and Senate meets.
It's just empty all the time.
Do you know any groups.
Who would like to rent this room.
It just was crazy.
And here's the one thing.
That's happening in Romania.
We went to Bucharest.
And later on we went to Cluj.
More on Cluj in a little bit.
If there's 90\% of people in Romania.
Who says they're Christians.
Why do they have such an issue.
With corruption in Romania.
Because a lot of times.
Back in the days of the communist regime.
Many people had to trust God.
For basic necessities.
They had to trust that God would provide.
They couldn't trust that.
The government would provide.
Because the government.
Was taking all their money.
So a lot of times.
Many of the older baby boomers.
Lived in this regime.
Where their basic needs were not met.
When communists fell in 1989.
Democracy came.
You started bringing a whole new generation.

$^{441}$Who did not live in that regime.
So now you know what happened.
They started attending church.
But now because of those opportunities.
They didn't need to depend on God anymore.
They don't need church anymore.
I'll go on Sundays.
Because that's what is expected.
I've always lived.
I could control my own destiny.
And did they have mentors.
To mentor them as they were growing up.
No.
Because this is the first generation of Christians.
Who now are living in a democratic society.
That is different from their parents.
Their parents couldn't be the mentors.
Because they couldn't understand it.
So suddenly you have all these Christians.
In Romania trying to figure out.
What it means to integrate their faith.
We then went over to Cluj.
Cluj is the second biggest city in Romania.
Cluj was filled with entrepreneurs.
Was filled with university setting.
People with new ideas.
Wanting to change their city.
But at the same time.
Many of the older folks.
There's a huge divide between the church.
And the business leaders.
As they're trying to figure out.
So a lot of times you have these Christians.
Who are living in Romania.
Trying to figure out how.
And what it means to integrate their faith.
From Romania I went over to Prague.
And in Prague we gathered.
This is my second year there.
We gathered all the wealthiest.

$^{481}$European Christians from all of Europe.
Who gathered in one place.
To talk about collaboration.
Because the perception is.
The church is dying in Europe.
But that's not so.
There's so much rebuilding in Europe going on.
And all these wealthy people.
Marketplace leaders are all gathered together.
To come and say.
How do we pull all of our money together.
To engage and see and collaborate.
That's the key thing.
Collaborate in what was Europe.
So from that point on to Prague and to Europe.
And two more stories.
And it ties into what we're talking about.
We then get to Orlando.
And there we gathered.
All of our young marketplace leaders.
From Hong Kong, from Indonesia, from Malaysia.
All of that.
And we have a bunch of different speakers.
We have one speaker.
Who said that who is on his way.
To minister to business leaders.
And CEOs in Hong Kong.
He was going to meet with the CEO of.
The President of Pakistan.
Oh by the way.
He was the former president.
Of a major organization.
A Christian organization.
A 30 million dollar organization.
Whose leader and founder.
Fell from grace.
But fell from grace in a very powerful way.
Because of sexual sin.
And this individual.
Who was the president.

$^{521}$Of this major organization.
Shared about how the entire ministry.
Completely fell and disintegrated.
Because of the sin of one person.
We heard the pain.
As he learned how to rebuild his life.
We also heard the story of Larry Smith.
Larry Smith is the CEO of Tokyo Electron.
He shared his story.
That when his kids graduated from college.
He and his wife sold everything they had.
They sold their mansion.
And they moved to a 300 square foot home.
In a tiny town.
Full of drug addicts.
Full of homeless people.
All of that.
So he could minister to them.
And learn how to build bridges with them.
Why did I tell you all these different stories.
Because each one of those individuals.
Whether you're here or listening today.
Are doing what God has called them to do.
Has continued to be so molded.
And so focused in their calling.
They are walking each day to make a choice.
Because they realize they are stewards.
Of what God has given them.
One speaker said this.
I love what he said.
I wrote it down.
That when you and I are born.
You look like your parents.
But when you die.
You look like your decisions.
When you and I are born.
You look like your parents.
But when you die.
You look like your decisions.
What are the decisions.

$^{561}$That you and I need to make.
To continue to serve and honor the Lord.
In everything that we do.
What does it mean for us to go and do.
Acts 1A tells us.
But you will receive power.
When the Holy Spirit has come upon you.
And you will be my witnesses in Jerusalem.
And in all Judea and Samaria.
And to the ends of the earth.
For the longest time when I read this.
I always thought it was just the Great Commission.
We always knew it was the Great Commission.
This is where God wants you to be.
Where's your Samaria.
Where's your Judea.
Where's your all of these different things.
Where's the end of the earth.
I understood it as the Great Commission.
We all see it as the Great Commission.
One pastor one time says to me.
Tommy, have you ever thought.
That besides this being a Great Commission.
How would your perspective change.
If you look at Acts 1A.
Not just as a commission.
But as a declaration.
Because when you look upon it.
As a declaration also.
It changes your perspective.
It creates urgency in what you need to do.
What do I mean.
It says that when you will receive power.
You're definitely going to receive power.
When the Holy Spirit has come upon you.
You will be my witnesses.
It's not a choice.
You're going to be my witnesses in Jerusalem.
You're going to be my witnesses in Judea.
You're going to be my witnesses in Samaria.

$^{601}$Samaria Jews don't go to Samaria.
You will be my witnesses.
In places that you don't want to be.
And you will be my witnesses.
To the end of the earth.
When we see it as just a declaration.
There is no choice.
There is an urgency for what we must do.
It changes our perspective.
It allows us to look at our life.
And what we're doing right now.
And realize the call of the Great Commission.
This declaration that Jesus is coming.
Means that the Holy Spirit is in each one of us.
And we are called to make a difference in our work.
In our family, in our community, in our church.
Not only from the outside, but from the inside.
It's how we live.
It's how we treat people.
It's how we think, all of it.
Because we are witnesses.
Because the Holy Spirit is upon us.
That's what we're supposed to be doing.
It means that we go into places.
Encounter people that aren't Chinese.
It means that we go to people.
Encounter people that are different than us.
In their skin color, in their beliefs.
In their sexual identity.
We engage all those issues.
Because we're Christ followers.
That's what it means to go and do.
That's what it means to be witnesses.
In everything that we do.
To go to places.
That are different than us.
My former boss.
The president of a major seminary.
Once said to me, "Tai.
All we need to do.

$^{641}$Is to teach people to study the Bible.
That's all you need to do.
That's it".
And I once had to say to this guy, "I disagree with you".
And I'd love to hear your thoughts later on.
I disagree with you.
Helping people studying the Bible is very key.
But I said to him, "A lot of times.
You have to understand this.
When I see you, you are a white Caucasian man.
President of a seminary.
Where you earn mid-six-figure salaries.
When you first started in a church.
You came into a big church.
Also earned six-digit salaries.
When you left the seminary and went to a church.
Another church hired you.
And gave you a big six-digit salary.
You never had to worry.
About racial discrimination.
You never had to worry about money.
You never had to worry about.
Any of those different things.
But some of us do.
Some of us have to deal.
Some of the issues in culture.
We're going to have to deal.
With racial discrimination.
We're going to have to deal with money.
And so a lot of times.
And for all of us.
As Christ followers.
Dealing with studying the Bible is great.
It's very great.
But we as Christians.
Have to learn how to engage the culture.
We have to learn to address the culture.
We have to learn to be.
What it means to really be able to be.
Engaging with the Christ-filled mind.

$^{681}$That's what it means to take the gospel.
To all the ends of the earth.
What does it not mean.
Here's what it does not mean.
It does not mean that you live.
Like a nomad like John the Baptist.
And be in need and give up everything.
Sometimes I think we as Christians.
Say, "I'm not dealing with this.
"Because it means that I live this poor life.
And it's like, no, it doesn't.
For some people, they've chosen to do it.
But for us as Christ followers.
It does not mean that we give up everything.
Like John the Baptist.
But here's what it does mean.
It means that you do not get to.
Pick and choose when to follow Jesus.
It does not mean that you get to.
Pick and choose when to follow Jesus.
We pick and choose.
When we pick and choose.
When to follow Jesus.
It doesn't lead to transformation.
That means that every single day.
When you and I wake up.
Our mentality is.
We are a steward of grace.
We are a steward of the grace.
That God has given us.
We are a steward of God's will in our lives.
And how do we use those talents and strengths.
To continue to make a difference.
To live a life that glorifies Him.
In every single way.
Luke 14, 26 tells us.
If anyone comes to me.
And does not hate his own father and mother.
And wife and children.
And brothers and sisters.

$^{721}$Yes, and even his own life.
He cannot be my disciple.
And as an Asian person.
A lot of times there's a lot of honor and respect.
That we have toward our family.
And when we read this, we say.
Whoa, whoa, whoa, you want me to hate my mom and dad.
That's not very nice.
What it means is that your dedication.
Your focus, your frame of mind.
Your focus of mind.
Is so strong every single day.
That when it is compared to your relation.
With anything else.
You are hating your mother.
You are hating your father, your wife, your children.
Your own brothers.
That's the focus that we need to have.
We don't need to pick and choose.
We don't even get to pick and choose.
It's our focus every single day.
Why?.
Because the Holy Spirit is within us.
And when the Holy Spirit is within us.
We are called to go and do.
When the Holy Spirit is within us.
We go and do to go to Samaria.
We go and do over in Jerusalem.
We go and do over in all the ends here.
That's what we're called to go and do.
And that means every single day.
We are light on the city of a hill.
To make a difference in our own lives.
In the lives of other people.
But the problem is in Christianity.
We've created a Christianity.
And I'm not sure for those who are watching.
We've created a Christianity.
Where no one wants to be part of us.
We've created a Christianity.

$^{761}$Where people are sick and tired of us.
We've gotten too cute.
Where we sit there and look at pastors.
These days who love their platforms.
I was literally on Facebook the other day.
And I came across these ads on.
How do you help your pastors in church.
Snippet their sermons.
And create instant reels.
And create their platform.
We've gotten to a point where we're too cute.
That we love to be influencers.
And we want to create these platforms.
We've gotten too cute.
In the sense that a lot of times.
Here in the US.
We've started fighting over dumb things.
We want to know who's right.
On the political spectrum.
We want to know who we're supporting.
And within the church.
That has completely divided us.
We want to know that we're right.
And you better believe.
There's no longer dialogue.
And conversations and arguments.
All of these things.
Because everybody wants to be right.
And the rest of the world.
Has been watching us.
And they are looking at us and say.
If that's the way the church is.
If that is the people.
And the leaders that we encounter.
If that is the hypocrisy that we see.
If that is what we're dealing with.
With pastors who want to create their platform.
With people using each other.
With people misspending money.
With people fighting all the time.

$^{801}$Then why in the world.
Do I want to be part of that church.
Because as Christ followers.
We do not have a singularity of mind.
To know that every single day.
We're going to deal with different people.
And every single day.
Through the power of the Holy Spirit.
We are called to take his gospel.
To every part.
And to encounter people.
Who are different to us.
That has different opinions.
And how do you show God.
In a loving way.
That's what we do.
So what do we do.
Couple of things as we wrap up.
Couple of things.
How has God designed you.
How has he gifted you.
What are the skills that he has given you.
We're not all called to do the same thing.
We're all called to have the same focus.
But we're not called to have to do the same thing.
Peter was not sent to the nations.
Paul was sent to the nations.
But did Peter have to learn.
The importance of what it meant.
To reach out to Gentiles.
He did.
Joseph was called to be a city official.
To be the chief of staff.
To really be able to nurture.
During that period of time.
Nehemiah was called to rebuild.
The walls of Jerusalem.
At a later age.
Before that he was called to be.
A faithful cupbearer to the king.

$^{841}$Zacchaeus of all of those guys.
Was not like Matthew.
Called to leave his role as a tax collector.
And follow Jesus.
And be one of his disciples.
He was called to be a chief tax collector.
And stay at his job.
And be faithful at his job.
What do we do.
It comes down to how God has created you.
And designed you to be.
It comes down to how God has gifted you.
And the passions that he has given in your heart.
We're all called to take the gospel.
To the ends of the earth.
We are not called to pick and choose.
When we follow God.
But how that looks like.
Is going to be different.
For each person.
And it's going to involve your time.
It's going to involve your outlook.
It's going to involve your motivation.
And it's going to involve your money.
It's going to involve all of those things.
As we seek to follow God.
In all that we do.
That's really what we're called.
When Elijah says to the widow.
Go and do.
God says the same thing for you.
Go and do.
Because if our motivation.
Is the grace of the cross.
What does it mean for you to go and do.
What does it mean for you to go and do in your work.
What does it mean to go and do in your family.
What does it mean to go and do.
To engage with people who are different than you.
We're called to go and do.

$^{881}$We're called to engage with people who are different.
Who think differently.
Who act differently than we are.
That's what we're called to do.
How that looks.
Is going to be different for each person.
Father, thank you for the words of Elijah.
You call us to go and do.
Where do we go and do.
How do we go and do.
We know one thing, Lord.
We don't get to pick and choose when we go and do.
But may we be aware of open doors.
And open opportunities.
May we be aware of how this affects all of our life.
Both our inner character and our outer lives.
How it affects our money, our calling.
Our decisions that we make.
Our motivation, Lord.
Help us to understand what that go and do means.
Help us as we reflect upon that phrase.
Wrestle with that word, those words.
And help us to see you in the midst of it.
Christ, we pray, amen.
\newpage



\section{1:26-38}
\label{sec:UpJWRYxxqyc}
\textbf{Unseen Service (Luke 1\_26-38) - Sam Johnson}
\newline
\newline
連結: \href{https://youtube.com/watch?v=UpJWRYxxqyc}{\texttt{https://youtube.com/watch?v=UpJWRYxxqyc}} ~~~~ 語音日期: 2023-05-16
\newline
\newline
\hyperref[sec:hS4BoW8e_54]{\small{< < < PREV SERMON < < <}}
~
\hyperref[sec:index_chronic]{\small{[返順時目]}}
~
\hyperref[sec:index_scriptual]{\small{[返順卷目]}}
~
\hyperref[sec:_7dHuWNR_kw]{\small{> > > NEXT SERMON > > >}}
\newline
\newline
$^{1}$him. Thank you and I'm excited to be here today. I love being at UV. I get to see.
old friends and new friends every time I'm here and it's really good to be.
back. I haven't been here for a few months. It's been a while but I was here.
years ago for a long time for a period of years was here and always a part of.
the UV English service so it's really good to be back with you. This morning.
we're gonna look at the story of a mother. It is Mother's Day as we've.
already noticed and discussed this morning and that's a special day..
Hopefully you've already had a chance if your mother is living. Hopefully you've.
already had a chance to connect with her and offer some words of appreciation and.
gratefulness and encouragement to your mom. My wife is here with our kids..
We have five kids and our oldest is 11 so she's a very busy mom and my mom I.
come from a family of 10 kids so I had nine siblings. I've seen a lot of.
moms up close who had a lot to do and seen the service that they give and.
that's what we're gonna talk about today how that can be a model for us. I know.
not everyone has a mom who models God's love for them. Sometimes our.
relationship to our mom can be painful and if that's your experience that's.
fine that God can heal that and God can can bring us through that and even the.
things that we're gonna talk about today are things that apply really to all of.
us and so we're gonna be looking in the first chapter of Luke. We're looking at.
this very famous mom in Scripture, Mary the mother of Jesus and in this passage.
we see the beginning of the story of Jesus's physical manifestation and his.
birth here on earth. Now most of us associate this story with Christmas..
We're in Luke chapter 1 we're gonna be looking at verses 26 to 38 and a lot a.
lot of times we think about this story in terms of Christmas and it is that's.
okay it is a part of the Christmas story but this is also the beginning of not.
just Jesus's birth but the beginning of Mary and of what we see of her.
throughout the story of the Gospels and so today we're gonna think about her.
life a little bit and some of the things that we can learn from her life and the.
theme of what we're talking about that we see modeled in her is this this heart.
attitude this posture of unseen service and we see this in Mary's life and if.
you had a mom who loved you especially if you had a mom who loved Jesus and.
loved you I'm sure you saw this in your mom's life. Moms model this so well for.
us this this attitude this posture of unseen service. So let's think about Mary.
for a minute. Here she is a seemingly unremarkable young woman. She's a peasant.
woman. She doesn't have any wealth, any authority, any position, any status. She's.
from one of the most forgotten small villages in Israel, Nazareth, and she's.
going to become known for being the mother of Jesus but she's also going to.
be known as we're considering today for being unknown. There it's conspicuous her.
story is conspicuous because of how little we know about her even though she.
plays this really important role we really don't know a lot about her and.

$^{41}$the the biggest picture we see of her is at this point when when Jesus is born.
and then after that throughout the Gospels there's only really small pieces.
of her story that we see and that tells us that so much of her service to God.
was unseen so much of her service to God was invisible to us and honestly.
invisible even to the people in her life. We don't know how much time she spent.
caring for Jesus or all the things that she did as a mom to care for him when he.
was an infant or when he was a toddler or when he was a boy and eventually a.
young man. We don't know all of those things that she did. We know she did.
those things because that's a part of human life and she was a mom but we.
didn't see them they're not recorded for us. We don't know all the prayers that.
she prayed over her son. We don't know all the tears that she cried. Mary we.
know from the scripture it talks about her reflecting and considering things.
in her heart. She was deeply spiritual. She knew God. She was obedient to God..
She was a servant of God and she watched Jesus be persecuted and she watched all.
the things that Jesus went through in life. That must have been deeply.
emotionally wounding to her but we don't we don't see that. That's not a part of.
the story that we're given and that's one of the exemplary things about her.
that we can learn from is that she plays such an important part but she's so.
humble and she's content for her service to God to be unseen and that's why we're.
gonna look at her story today. Since it's Mother's Day she's a great model for us.
and like I said so many moms model that for us as well but I want to be clear at.
the beginning here we're not saying that Mary and her unseen service is a great.
model for moms although that is certainly true. We're saying she is a.
great model for all of us because being willing to serve and to serve God and to.
be obedient even in ways that no one notices or sees or thanks us for is not.
just something for Mary to do it's not just something for moms to do it's.
something for every believer every follower of Jesus. We should all have.
this attitude this this posture of heart that we see in her. So that's what we're.
gonna consider today. So let's read this passage together beginning in chapter.
26 this is Luke 1 26 to 38. In the sixth month the angel Gabriel was sent from.
God to a city of Galilee named Nazareth to a virgin betrothed to a man whose.
name was Joseph of the house of David and the Virgin's name was Mary and he.
came to her and said greetings Oh favored one the Lord is with you but she.
was greatly troubled at the saying and tried to discern what sort of greeting.
this might be and the angel said to her do not be afraid Mary for you have found.
favor with God and behold you will conceive in your womb and bear a son.
and you shall call his name Jesus he will be great and will be called the son.
of the Most High and the Lord God will give him the throne of his father David.
and he will reign over the house of Jacob forever and of his kingdom there.
will be no end and Mary said to the angel how will this be since I am a.

$^{81}$virgin the angel answered her the Holy Spirit will come upon you and the power.
the Most High will overshadow you therefore the child to be born will be.
called holy the Son of God and behold your relative Elizabeth in her old age.
has also conceived a son and this is the sixth month with her who was called.
barren for nothing will be impossible with God and Mary said behold I am the.
servant of the Lord let it be to me according to your word and the angel.
departed from her the first thing we see in this passage about Mary's heart.
attitude her posture is that she has this posture of submission and surrender.
and that's so important for us if we want to be willing to enter into unseen.
and unacknowledged service to God we have to have a hard attitude of.
submission and surrender it's right there in verse 38 Mary says behold I am.
the servant of the Lord let it be to me according to your word now Mary was a.
young woman again there's so much we don't know about her but we know she.
must have she was a human being just like us so she must have had a lot of.
plans and dreams and hopes for her future just like every young person every.
young person has things that they want to do and ideas in their head of what.
they want their future to be like and whether it's relationships or something.
they want to achieve or something that they want to be a part of everybody has.
those dreams so Mary would have had those as well just like any other young.
person and yet she hears this she hears God's truth communicated to her from.
Gabriel and what does she do she releases her will and she surrenders to.
God's Gabriel's announcement is what we would call a watershed moment or a.
turning point in her life she's never gonna be the same after this it marks a.
pivot a change in the trajectory of her life she may have been heading in one.
direction before and now forever after her life's gonna be totally different.
she's going in a different direction and she knows that she understands that and.
she says that's okay I'm gonna submit to what God is doing in this moment and I'm.
gonna say God let my life be what you want may it be to me according to your.
word she surrenders to what God is doing and so that's something that we need to.
live out in our lives as followers of Jesus we need to be willing to submit.
and to surrender to what God is doing to God's truth obviously that's a difficult.
thing submission is a scary word for us just hearing that word is scary because.
it means a loss of control it means a loss of what I want my priorities and.
objectives and so I want to remind us that as we apply that in our lives what.
Mary was doing what God calls us to do is to submit to God's Word and submit to.
God's truth we're not simply submitting for the sake of submission and that's.
one of the reasons that word is so scary for us in our culture a lot of times we.
say well I've seen people who have submitted to the wrong thing and then.
they've end up being harmed or abused by that posture of submission right that.
causes us to be very wary of the concept of submission and that's true and we are.

$^{121}$to be wise and discerning because we're not to submit to everything we're.
actually supposed to resist evil remember James 4 7 submit yourselves.
therefore to God and resist the devil and he will flee from you so there's a.
lot of evil and injustice in our world that we are not supposed to submit to.
we're actually supposed to resist but the thing that we are to submit to is to.
God's truth and to what God has to say about our lives and that's something.
that if we will submit to it it is good and beautiful and right for us to submit.
to God and if we think about that for a moment it actually makes sense right.
because if God is who he says he is if he knows everything and he knows.
everything about our life and our heart and our future and if he's also who he.
says he is that he is perfectly good and loving then it makes a lot of sense to.
submit to his plan to say God here's what I would like to do here's how I see.
this working out in my future but I'm gonna trust that you have a better plan.
than I do I'm gonna trust that you know more than me I'm gonna trust that your.
desire for my heart is better even than my own desire for my heart and I'm gonna.
submit to you it's actually a good and beautiful and right thing for us to.
submit to God because he is the definition of good and his will is the.
definition of right and his expressions are the definition of beauty and so if.
we will submit to what he's doing then we will experience those things in our.
life that goodness that righteousness that beauty but we need to be clear on.
that when we're looking for that posture of submission where are we submitting.
we're submitting to God to what he's saying Mary was right to submit to.
Gabriel because Gabriel was communicating God's truth to her and.
then what comes out of that what fruit is produced in our lives when we begin.
to submit to God the outcome is that God has been able to work his good and.
beautiful and right plan in our lives for his glory and for our good Mary's.
submission meant that she was able to play this part in God's grand redemptive.
narrative of every human that was going to be saved through Jesus now that's a.
beautiful and mind-boggling thing right we think about all of the people from.
every tongue tribe and nation throughout all the generations of the earth that.
have come to relationship with God through Jesus Christ and through his.
sacrifice on the cross and through the redemption that Jesus brought and Mary.
got to be a part of that story if she had not submitted and surrendered to God.
she would not have been a part of that story now obviously in that moment it.
may not have felt great to her to say I'm gonna give up all of my plans and my.
dreams and my desires my ambitions and and obey God and she probably had some.
some conflict within her some second thoughts but because she did that she.
got to be a part of something that was beautiful and amazing and so far beyond.
herself and that's true for us as well it's not easy to surrender to what God.
is doing in our lives sometimes it can be difficult to submit but as we do that.

$^{161}$we get to be a part of something that's good and beautiful and way bigger than.
us and one of the things that's clear in scripture is that God always works.
always for his glory and for our good anything that God allows in to our life.
is only there because it's either gonna glorify him and make his name great and.
cause people to see him for who he is or it's gonna be good for us maybe not good.
in the way that we imagine when we think about what's good for us we think.
comfort blessing success you know material things that's what we think.
right sometimes I think God says hey you know what's good for you at times is.
adversity and challenge because that brings growth or to go through something.
difficult because you're gonna learn a lesson you wouldn't learn any other way.
but we can trust that God's always at work for our good even if it doesn't.
feel great in that moment and that's the fruit that was produced in Mary's life.
that's the fruit that will be produced in our life if we have that posture of.
submission to God God's glory will be evident in our life God will be.
glorified and God's plan will be at work in our life and we'll get to see God's.
good plan for us and for those around us the other thing we see here is that this.
posture of unseen service is a posture of the background vocalists throughout.
the rest of the Gospels Mary kind of fades from our view that she's there in.
different places she reappears briefly when she and Joseph take Jesus to.
present him to the Lord at the temple that's in Luke chapter 2 when her and.
Joseph are looking for Jesus when they leave Jerusalem after the past Passover.
that's later in Luke chapter 2 she's there at the wedding in Cana that's in.
John chapter 2 she was certainly there when Jesus was rejected at Nazareth.
that's in multiple Gospels Matthew 13 is where it's at in Matthew she's visiting.
Jesus during his ministry that's recorded in Luke 8 and also in the other.
Gospels and John 19 tells us that she was at the cross when Jesus was.
crucified so she's there she obviously was a true believer in what Jesus was.
doing and she was following him she was a follower but we don't see much of her.
her story really fades she fades into the background just like a background.
vocalist if you ever see a singing group or maybe you see some a favorite musician.
or artist on TV that you really enjoy most of especially on TV but even at a.
concert if there's screens that are showing them most of the shots are of.
the artist right the main person that's mostly you're gonna see and occasionally.
you'll see the background vocalist standing off to the side and they're.
usually you know snapping and dancing sometimes doing their thing but you just.
see them for like a half second and then it goes back to the main person right.
and that's the role that Mary took she had this posture of fading into the.
background now from a human perspective that doesn't seem fair right because we.
say well you know the main singer may be a great singer but the background.
singers they can sing great too why is it fair that we're always looking at the.

$^{201}$main singer and from a human perspective that's that's actually true none of us.
as humans should should necessarily be thrust into the limelight or celebrity.
that's that's not something that that God ordained for us because we're to.
love and respect one another we're all made in God's image so we're not to.
Center ourselves but when we think about Jesus it's different right because Jesus.
wasn't just another man even though he was fully human he was also the Son of.
God and he was also here to do something that was so much bigger than anyone.
could comprehend Jesus actually reminded Mary of this in the wedding at Cana in.
John chapter 2 verse 4 when Mary comes to him and she's saying Jesus my friends.
are having a wedding they ran out of wine like what can you do and Jesus says.
hey it's not the right time yet he reminds her hey I there's something.
bigger going on here now Jesus demonstrates his love and care for her.
and that he goes ahead and does the miracle it turns the water to wine but.
he was reminding her in in verse 4 of chapter 2 there that he's saying there's.
a bigger picture here that you don't see and that was true of Jesus he really was.
supposed to be centered because he was the Son of God and he had come to redeem.
the souls of humanity and so everyone else was supposed to be in the.
background his purpose was so much bigger than Mary or anyone else he was.
fulfilling God's redemptive plan and he was doing something for the good of all.
of us and Mary recognized this and she was content to fade into the background.
we see it there in verse 35 in verse 34 Mary says how will this be in verse 35.
angels answers her and says the Holy Spirit will come upon you in the power.
of the Most High will overshadow you God's power will overshadow you you will.
be you will fade you will recede it's what that language is communicating and.
we need to apply that in our lives we need to have that posture as well.
remember Jesus taught us to pray in the Lord's Prayer in Matthew chapter 6 verse.
10 he taught us to pray your kingdom come your will be done on earth as it is.
in heaven God made your purposes come to the.
forefront and may mine recede and fade John the Baptist said one of the John.
the Baptist said a lot of really profound stuff he didn't say there's not.
a lot recorded of what he said but it's all super profound but in John 330 we.
have this he said he must increase but I must decrease what a beautiful.
statement of what it means to be a disciple follower of Jesus is that Jesus.
gets bigger in people's view and I get smaller and Mary modeled that for us she.
faded into the background what about the fruit that comes out of that what fruit.
is produced when we fade into the background when we fade into the.
background God's narrative is centered what God is doing becomes the main focus.
of the people around us but here's the cool thing is that we still get to see.
what's happening and we get to watch it up close I had a friend in another state.
that worked in a big music venue for a long time from like the 70s to I remember.

$^{241}$having a conversation with him in the early 2000s he was telling me about this.
and he was talking about all the bands that came through the venue over the.
years and big-name bands that you know the biggest bands in North America would.
come through and he was there for all the shows but he wasn't in the audience.
he was on the stage now he was like a tech so he wasn't performing he was.
making sure everything else worked but he got to see every show from right up.
close and he got to see a different view than most people saw because he was.
backstage instead of in the audience and that's one of the perks of being in the.
background you're not visible but you get to be a part of something from a.
perspective that no one else sees and if you're willing to be in the background.
of God's work in your life yeah no one's gonna be showering you with praise and.
and your service to God might be unseen and unrecognized but you will get to see.
some miracles up close you'll get to see some amazing things happen from just a.
few feet away I can't even imagine what that was like for Mary for the disciples.
for the other followers of Jesus to see the things that Jesus did right there up.
close firsthand it's why their lives were radically transformed right because.
for them it wasn't something abstract it wasn't a story they heard from someone.
else they saw people get healed they saw dead people raised to life with their.
own eyes they were right there for it and that's one of the beautiful things.
about being willing to fade into the background as you get to be a part of.
that and this is true as a parent it's true for moms it's true for dads if you.
try to make it about yourself and center yourself in your kids life and what's.
going on you're gonna create all types of problems but if you'll fade into the.
background and say I'm gonna love this young person I'm gonna help them to.
follow Jesus I'm going to help them to walk into what God has for them you'll.
get to see some really neat things happen you'll be a firsthand witness to.
those things it's true in disciple making if you're trying to help someone.
else follow Jesus and you make it all about you and you're trying to.
communicate all the things that you know and you are really trying to make that.
person follow you they're gonna be an unhealthy disciple but if you can fade.
into the background and help them learn to follow Jesus and you come alongside.
them and support them and encourage them and continually point them towards.
Christ you're gonna see a miracle of transformation in their life it's true.
when you think about praying for someone else interceding on someone else's.
behalf you're gonna pray stuff for someone if you love someone you say I'm.
gonna pray for this person you're gonna pray things they're never gonna know.
about and no one else is ever gonna know about no one's gonna know that you were.
praying for for them to come to know Jesus or that you were praying for them.
to experience blessing in their life or you were praying for them to be healed.
in some way they may never know no one else may ever know but when God answers.

$^{281}$that prayer you see the miracle and you get to see it up close because you were.
willing to be in the background and it's true in all forms of unseen service that.
honors God it's not easy to do something obscure that no one is noticing and to.
do it without being acknowledged and without being thanked but one of the.
beautiful rewards is that you'll get to see God do some powerful things and.
you'll see them up close and then we see here that Mary her unseen service this.
is a posture of trust and that's what we've got to learn if we're gonna be.
willing to serve God in a way that doesn't center ourselves it's not about.
us that's not focused on me then we've got to trust that God's Word God's truth.
God's way is best we see that here in her story verse 29 says but she was.
greatly troubled at the saying and tried to discern what sort of greediness.
might be verse 30 and the angel said to her do not be afraid Mary for you have.
found favor with God Gabriel comes and he begins to explain God's truth to her.
and her initial reaction is fear she's saying I don't get it this is terrifying.
I'm seeing an angel she's totally afraid but what happens by the time we come.
down to verse 37 and verse 38 Gabriel in verse 37 says for nothing will be.
impossible with God Gabriel saying hey I know you're scared he even says that in.
verse 30 do not be afraid for you found favor with God then he gives her the.
message and then at the end he says I know this is tough to believe but.
nothing's impossible with God God can do this and then verse 38 what does she say.
she says I'm the servant of the Lord let it be to me according to your word she's.
saying okay I'll I'll believe you I'll go with it I'll obey we go through that.
same process in our life too right when we initially encounter God's truth.
God's Word it challenges us sometimes it scares us right and we say God how could.
that happen how could that be possible just like Mary was saying this is there's.
no way and I could I could have a child I'm not even married I've never been.
with a man and there's no way that that child could be the Son of God in the.
side it's not even possible and God may bring you to a place in your life where.
he's he's asking you to do something and you said this is not possible you say.
God you don't understand it's not possible for me to talk to that person.
at work about Jesus like I don't know how to do that I'm not good at that.
stuff and that person doesn't even really like me so it's not even.
possible say God it's not possible for me to find time in my week to serve you.
or to connect with this other person I know you've put that on my heart but.
it's just not possible I don't have the time God it's not possible for me to use.
my gifts in church or somewhere else to honor you because you know I'm I'm.
really shy and I'm uncomfortable doing that God it's not possible for me to.
make a big move in my life even though I know this is something that you're.
putting in front of me it's scary we all come to to God's truth and are.
challenged and frightened by at different points in our life just like.

$^{321}$Mary was but what bridges the gap for us is we have to trust God that's what.
Mary did it's not that between verse 30 and verse 38 she understood the whole.
process because she did she didn't say okay now I get it it all makes sense.
yeah no she still doesn't understand how it's gonna work because Gabriel told her.
she was gonna have a child a super that was gonna be supernaturally placed in.
her body and so she's saying I still don't understand any of this so verse.
38 is an expression of trust not an expression of oh I get it now I can go.
for it I'm comfortable with it she's not saying that she's saying okay I trust.
you it's a posture of trust and that's true for us as well we need to trust.
what are we trusting in we're trusting in God we're trusting in his love for us.
and his goodness as I said earlier God's only gonna allow things in our.
lives that are for our ultimate good because they bring him glory and because.
they're at work to benefit us even if we don't see the benefit in in the moment.
of those circumstances now that's hard to believe sometimes because we go.
through stuff we say I don't see any way this could be for my good so we bridge.
that gap with trust say God doesn't feel good doesn't look good doesn't smell.
good right now but I'm gonna trust you that this is good that you're at work.
that you still love me that you're still good and that what you're doing in my.
life it's gonna bring you glory and it's gonna be best for me and when we are.
able to do that when we're able to trust in God's love we have peace that even.
our unseen service is valuable that's one of the tough things about unseen.
service it's one of the tough things about serving God and obeying in a way.
that no one sees because if you do something and everybody in the room says.
you did a great job and we're so grateful you did it you already feel.
great about that you you don't ever ask yourself well did I do a good job.
because you just heard affirmation from everybody right when you do something.
especially when you pour your heart into it you say I really want to honor God.
with this choice or with this ministry or I want to honor God by the way I'm.
caring for this person in my life and you pour your heart into it you do it.
with all that you have and you use all the creativity and the wisdom and the.
energy you have to do that and then you walk away and no one says anything no.
one says thanks no one says good job no one says that was a blessing to me then.
you do have those questions Wow was that just a waste of time why did I even do.
that why am I even here what what's the purpose of this but if we're trusting.
in God's love we have the peace that says yeah that that service was valuable.
even if it was unseen because God loves me he loves the people in my life he.
wouldn't call me to do this and pour my heart into it unless it mattered to him.
and if it matters to him that's all I need whether he uses in a big way or.
small way that's all I need and when we trust in God we see that God uses that.
unseen service in a way that's beyond what we can see or imagine Mary is a.

$^{361}$background figure in the story of Jesus but she's not forgotten she's not.
invisible here we are talking about her life 2,000 years later so her role was.
important she got to play that part in the grand narrative of redemption that.
we were talking about earlier I'm sure she had no idea in verse 38 when she.
said okay God let it be to me according to your will she had no idea what that.
meant she had no idea that Jesus was the Messiah she had no idea that God was.
going to use Jesus to bring people from every tribe and tongue and nation into.
relationship with him and to give them a pathway to heaven she had no idea that.
we would be sitting here today because of Jesus there's no way she could have.
imagined that but she obeyed she surrendered she trusted and that's true.
for us too especially in those areas of unseen service we have no idea how God.
wants to use that we have no idea the ripples that will go out from that or.
how God can use that to change someone else's life or to change someone else's.
life who is two or three people removed from the person in front of us that.
we're struggling to love and so when we have that posture of trust and we say.
God I'm trusting you that your work is going forward that you're loving me that.
you're even in a difficult situation that you're doing this for my good and.
for your glory God can do something huge out of that so much bigger than we'll.
ever know so that's my desire and prayer for us today now unseen services it's.
not easy it's not easy to be a mom wasn't easy for Mary to play the role.
that she played it's not easy for any of us to serve God and in an obscure.
anonymous way and honestly if you don't know Jesus it's impossible if you don't.
have relationship with God there's really no way you can truly love and.
serve someone else in the way that you should without getting some type of.
reward or getting something back so if you're sitting here today and you say.
this makes no sense to me it's impossible understand that you have to.
have relationship with God through Jesus Christ that's why Jesus came to the.
earth it's because we're all sinners were all broken and he through his life.
and then his death satisfied God's justice and God's holiness so that we.
could have relationship with God if you're gonna try to love someone else.
this way without relationship to God it's not gonna work but if you do have.
relationship to God then God wants you to love and serve others in this way.
it's not easy but we're all called to it it can be painful at times it can mean.
pouring ourselves out pouring out our energy our resources our gifts in.
anonymity and no one might see it no one notices all that we're giving but if.
it's done in obedience to God and that's the the motive of our heart if we do.
these things if we have a posture of submission and surrender if we have a.
posture of fading into the background if we have a posture of trust then God can.
use it in an amazing and powerful way and that's what I want to encourage us.
with today I want to salute and commend all of the moms who who model this for.

$^{401}$us and who do it so well day in and day out I want to commend this example that.
we have in Mary and I want to challenge myself and all of us to live into this.
to be obedient to God even when no one's watching to do it for him to surrender.
submit to his will to let ourselves not be centered or platform but to fade into.
the background and then to trust that he's going to do something with it let's.
pray God we thank you for this example we thank you for moms that give us this.
example God I pray that we would be grateful and gracious today as we.
connect with our moms and just encourage them and express gratitude for for the.
way that they loved us but God I pray that you would help us to grow in this.
every one of us help us to grow in serving you obedience to you surrender.
being willing to to fade and let you be the forefront of our lives God help us.
to grow and trusting you help us to be beautiful unseen reflections of who you.
are so that the world can see you and so the people can be drawn to that and be.
saved in Jesus name amen.
[MUSIC PLAYING].
\newpage



\section{申命記 4:9}
\label{sec:_7dHuWNR_kw}
\textbf{謹記這一口氣 (申命記 4\_9, 23) - 張麗仙傳道}
\newline
\newline
連結: \href{https://youtube.com/watch?v=_7dHuWNR-kw}{\texttt{https://youtube.com/watch?v=\_7dHuWNR-kw}} ~~~~ 語音日期: 2023-05-16
\newline
\newline
\hyperref[sec:UpJWRYxxqyc]{\small{< < < PREV SERMON < < <}}
~
\hyperref[sec:index_chronic]{\small{[返順時目]}}
~
\hyperref[sec:index_scriptual]{\small{[返順卷目]}}
~
\hyperref[sec:GOoaQbj9lhw]{\small{> > > NEXT SERMON > > >}}
\newline
\newline
申命記 4:9
\newline
\begin{longtable}{cl}
\hline
\hline
章節 & 經文 (和合本修訂版)\\
\hline
4:9 & \begin{tabularx}{0.7\textwidth}{X} 「但你要謹慎，殷勤保守你的心靈，免得忘記你親眼所看見的事，又免得在你一生的年日這些事離開你的心，總要把它們傳給你的子子孫孫。 \end{tabularx} \\ \\
[1ex]
\hline
\hline
\end{longtable}
$^{1}$然後中文課.
今天就是黎先傳道講到這裡.
現在時間交給黎先.
Hello 各位好.
喜歡那個生物認領柱嗎?.
想不想試一下?.
哈哈哈.
那個生物認領柱的盒子只有這麼大.
但如果你有些東西大一點的.
例如你的外套.
或者你的小朋友不見了.
你也可以找Amanda.
哈哈哈.
好各位母親節快樂.
今天我知道大家都很有慶祝感覺.
因為我知道不僅是母親節.
還有畢業禮.
好像很多.
是不是?.
好像有我們的牧師去了畢業禮.
所以其實今天是很多慶祝.
很開心的一個日子.
所以我今天預備講到的時候.
我預備今天的講到的時候.
我自己的心情也很輕鬆起來.
我想一下.
不如我也輕輕鬆鬆地跟大家講一整句話.
不過雖然是輕鬆.
但絕對不是不重要的.
而且我覺得今天要講的東西也是超級重要的.
先講一個故事.
我不知道你聽過沒有.
從前有一位學業和事業都有成的少女.
二十多三十歲.
她有四個男朋友.
哈哈哈.
人家是這麼挑的.
不是四個老公.
是四個男友.
第四個呢.

$^{41}$她是.
第四個她是最喜歡的.
OK.
聽一下為什麼.
她喜歡到經常都買.
為什麼?.
經常都買一些漂亮的衣服給她.
穿得她很型格.
很瘋狂.
很貴氣.
經常跟她去吃美食.
韓日,意大利,法國.
中國海底撈.
她都會混合的.
總之就是跟她一起.
每次都是跟她去A級的享受.
不是A4的享受.
是A級的.
OK.
第三個男朋友.
她也是非常喜歡的.
而且這個男朋友很有成就.
所以這個女朋友很喜歡帶他出去介紹一下人家.
她是什麼博士呢?.
介紹給人家認識.
別人稱讚她.
她的心很灰心.
但是她對這個男朋友的心裡面.
也有些顧慮.
就是她經常怕他丟下她.
OK.
第三個男朋友.
第二個男朋友.
她也是很喜歡的.
而且這個男朋友一向給她最大的安全感.
這個男朋友很親切.
很體貼.
還有很能忍受她.
所以有什麼事情不妥的時候.
她當然會先找他.

$^{81}$她當然會找他.
她當然可以幫助他.
這是第三個男朋友.
第四個.
倒過來.
這是第二個.
真的會聽我說話的.
OK.
這是第一個男朋友.
最後一個.
這個少女是最不上心的.
但是這個男朋友是超級忠心的.
這個女生從她年輕成長讀書.
出來工作.
一直都很支持她.
不過這個少女不知為什麼偏偏對她最冷淡.
最少時間跟她在一起.
除了自己很搞不定的時候.
都會找她.
說到這裡.
問問現場的姐妹.
有沒有一個像你的男朋友呢?.
你喜歡哪一個呢?.
但是這樣.
弟兄不用生氣.
因為這個故事絕對不只是給女朋友說她有幾個男朋友.
其實可以倒過來看.
因為有一個男朋友.
有一個男生.
他有四個女朋友.
這樣數下來都可以的.
OK.
不過你很想知道故事發展如何.
下集.
很不幸.
這個少女突然患了蟲病.
她知道自己日子無多了.
她心裡真的很無助.
她回望她一生.
她這四個男朋友.

$^{121}$她想到.
這幾個一生都跟我這麼親密的人.
我問問他.
可不可以陪我去到哪裡.
於是她問第四個男朋友.
他說我一輩子都這麼愛你.
我給你最好的你.
最後給你最好的你吃.
我真的看你最重要的了.
現在我走了.
你可不可以陪我.
我走了.
大家知道是什麼意思吧.
現在我走了.
可不可以陪我.
答案就是.
斬釘截鐵.
當然不行.
一句話都不說.
這個男朋友就走了.
你可以想到.
其實這個答案也不是很離譜.
我怎麼陪你走.
走去哪裡.
但是.
的確是令這個少女心如刀割.
因為這個男朋友她太重視了.
第三個男朋友.
她走去問他.
他說我一生都這麼愛你.
我很引以為榮.
我事事都很proud of你.
所以經常在別人面前稱讚你.
高舉你.
現在我走了.
你可不可以陪我.
誰知道他的回答更厲害.
比之前那個更厲害.
他說.
當然不可以陪你.

$^{161}$人生這麼短.
你走了.
我很快會找第二個.
是不是真的要命.
這個時候這個少女心真的沉下大海.
第二個男朋友.
跟她說.
我一生人當中.
你真的對我最好.
從我小學到大學.
我出來工作.
你真的對我無微不至.
我對你很感激.
但是現在我走了.
你可不可以陪我.
這個男朋友跟她說.
我真的很愛你.
我一生裡沒有什麼不給你.
我盡力幫你.
但是這次我做不到.
我可以陪你走到你的墓地.
但是我不可以再走下去.
說出來也是很傷心.
而這個少女的心真的很傷心.
她真的想不到.
一生她倚靠的第二個男朋友.
在這個時候也這樣跟她說.
好了.
她還沒想到最後那個男朋友.
就是第一個.
誰知道這個就走過來跟她說.
主動過來說.
她說.
我是會陪你走的.
無論你去到哪裡.
我都會陪你的.
這個少女看一看她.
嚇了一跳.
因為其實她很久沒有看他了.
你知不知道.

$^{201}$很久沒有看他.
看一看她.
瘦成這樣.
好像形容不良.
為什麼會這樣呢.
這個少女在這個時候.
看到她這麼瘦.
又看到她說這些話.
真的心裡很難過.
眼淚都流出來了.
她說是我錯.
其實你對我一生.
你一直都不離不棄的.
但是我呢.
就時時對你這麼冷淡.
整天不理你.
所以她自己也很傷心.
故事結束了.
究竟這裡說什麼呢.
如果是你的話.
你會怎樣呢.
正如我說的.
其實這個故事.
不單是關於姐妹.
其實也是關於弟兄.
關於我們每一個人.
這個故事.
其實是給我們很好的醒覺.
第四個男朋友是什麼呢.
很明顯是一個象徵.
是你的身體.
是你最愛的.
在你一生裡.
你真的給他最多的關注.
給他穿得漂亮.
又看得到.
給他吃最好的.
因為你也會選擇最好的來吃.
哪有人會選擇不好的來吃.
那麼.

$^{241}$你一生裡.
你真的對他是最好的.
但是當你離開的時候.
你的身體會馬上離開你.
是一個母親.
沒有第二句話說.
是嗎.
第三個男朋友是你的財產.
你所擁有的.
有你的擁有.
可以說是你的名利.
你的成就.
你的金錢等等.
是值得你驕傲的.
你也忍不住有時候.
跟別人說.
我有這些東西.
但是很可惡的就是.
當你離開的時候.
這些東西是會跟著別人去的.
不會再是你的.
第二個男朋友.
你猜是什麼呢.
是你的家庭.
是嗎.
是一生裡最關心你的.
是一生裡願意給你最多的.
你最有難處的時候.
他會幫你的.
但是這個第二個男朋友.
無論怎樣也好.
當你離開的時候.
他可以送你去墓地.
接著他不可以走下去.
沒有辦法可以走下去.
而這個第一個沒有體面.
這個這麼瘦的男朋友.
這個少女是最不注意他的.
最冷淡的男朋友是什麼呢.
是什麼呢.

$^{281}$就是你的靈魂.
如果你是主耶穌.
其實你也是在說那些東西.
是你的靈魂.
一生裡你和我.
都是花很多時間.
在前面那三個男朋友身上.
問心真是.
我們是花很多時間.
在前面那三個男朋友身上.
我們嚴重的忽略.
第一個男朋友.
就是我們的靈魂.
現實就是.
只有他.
只有這個第一個男朋友.
這個靈魂.
不單止陪你走到人生盡頭.
而且會繼續走下去.
無論你是走上去還是走下去.
他都和你在一起.
靈魂是和你在一起.
直到永遠.
我想問問.
其實以前聽過這個故事.
請你舉高你的手給我聽聽.
這樣也不錯.
第一次將這麼有趣的故事.
告訴大家.
這是一個虛構的故事.
但我覺得是一個.
很發人深省的故事.
今天只有一式尚存的.
你和我在這個世界.
這個故事對我們.
是一個當頭棒的提醒.
我們有沒有太輕視.
我們的靈呢.
這個真正永恆的我們呢.
說到這裡.

$^{321}$我馬上想起我小時候.
第一次接觸聖經.
聽到的第一句話.
我想給大家看看.
這句話我不知道你們聽過沒有.
耶和華上帝用地上的塵土做人.
將生命之氣吹進他們的鼻孔.
這人就成了有靈的活人.
在《創世記》第二章第七節.
是我小時候在基督教學校讀.
老師在聖經堂裡面說.
上帝創造天地的時候.
當他說到這句話的時候.
我覺得很震撼.
雖然當時我很小.
很不懂事.
但我覺得這句話也很震撼.
因為上帝做天,做地.
做很多動物,植物,海洋生物.
但沒有一件事.
他是會這樣做的.
是吹吹嘔吼氣.
使他們成為有靈的活物.
難怪我小時候.
總覺得我比我的狗和貓厲害.
雖然有些東西是比我厲害的.
牠跑得比我快.
跳得比我高.
這些東西我不及牠.
但我比牠聰明很多.
為什麼呢?.
有什麼不妥.
我會向上帝求.
牠會不會求呢?.
貓,狗.
我害怕的時候.
我會告訴上帝.
我害怕.
我看到大自然這麼雄壯.
這麼偉大的時候.

$^{361}$我會震驚,震撼.
我不知道貓,狗會不會.
我不知道牠吠的時候.
是否代表牠很震撼.
我不知道.
但我相信是不同的.
因為上帝吹了一口氣在我們身上.
我們是有靈的活人.
甚至我們做錯事的時候.
我們良心會不安.
我們有道德的.
無論你多麼疼愛你的狗.
我知道這裡有很多愛狗的人.
KK,旺旺,Pepper.
很多好狗.
但無論這隻狗多麼厲害.
牠沒有這一口靈氣.
但你的子女呢?.
有.
好像中國人常說.
除非你生的叉燒.
如果你生的子女是一個人.
牠就有靈氣.
好像你,我.
好像歷世歷代的人一樣.
因為上帝吹了一口靈氣在他們身上.
只可惜我們的人生裡.
很多時候我們都像這位少女一樣.
我們太過顧惜我們頭三百個男朋友.
身體,擁有,家庭.
是代表一些親密的關係.
而忽略了靈魂的地位.
牠的重要.
想想我們真的.
我們超級滿足我們的身體.
但很多時候我們超級輕視.
忽略我們的靈魂.
是本末倒置的.
在聖經裡.
神告訴我們這是本末倒置的.

$^{401}$但這也是撒旦的陰謀.
從創世記開始.
就是這樣的陰謀.
他要我們忽略神給我們一個獨特性.
這個靈.
因為這個靈使我們與上帝能夠溝通.
能夠連接.
使我們成為神的形象.
撒旦就是故意攻擊這個東西.
使我們一生都把這個靈放得很低.
結果就是我們一生的精力.
我們為肉體,為我們生活.
我們很拼命.
但我們沒有重視到我們靈魂的需要.
也因此而沒有孕育到我們與我們的主.
創造主的關係.
不過這個被我們忽視的靈.
它真的會時時不離不棄地跟我們說話.
它會跟我們呼喊.
我相信你跟我一樣都會體會到.
很多時候晚上,深夜的時候.
或者半夜起床上廁所.
回到床上還沒睡著的時候.
這些時候.
因為你已經沒有這麼忙了.
沒有為你的身體去做什麼.
去忙的時候.
你靜下來的時候.
你覺得你的靈會跟你說一些.
你白天沒有想到的東西.
你只不過說說你睡著了.
所以又是沒有重視這個靈.
但是.
很微妙的就是.
在我們人生低谷的時候.
什麼低谷呢?.
就是.
或者你有些事情很挫折.
很大大的失敗.
又或者你跟一些人的關係突然間.

$^{441}$撞擊得很厲害.
你很傷心.
無論那個是誰都好.
又或者你是.
你是.
發覺自己身體有些問題.
你.
很震驚.
不知道怎麼辦的時候.
在這些的時候.
你會聽到靈的聲音.
又再大聲的跟你說.
你的靈跟你說什麼呢?.
就是告訴你.
不是靠你自己.
你回到上帝那裡.
祂會給你盼望.
祂會給你希望.
祂會給你力量.
祂會給你智慧.
不是靠你自己.
在這些難關.
在這些難處裡面.
又或者你甚至掙扎在大病裡面.
這些難處裡面.
你反而會大徹大悟.
你反而會靈魂蘇醒.
在這個時候.
咒詛.
在你人生的咒詛.
說我會變成祝福.
因為你開始重視你的靈.
有人說得很好.
他說.
淚水洗過的眼睛.
比以前看得更清.
眼淚洗過的眼睛.
就是你在困難裡面.
比從前更看得清.
而我加一句就是.

$^{481}$痛苦.
使我們的靈魂蘇醒.
我很喜歡.
大家都很熟悉的一位作家.
C.S. Lewis 著名的作家.
神學家.
他在《Problem of Pain》.
《痛苦的奧秘》這本書裡面.
他說過一句很出名的話.
他說.
在我們的享樂裡面.
享樂.
享樂不是壞事.
不過他形容.
在我們的享樂裡面.
神會用很微聲對我們說話.
在我們的良心裡面.
神會用正常的聲量.
對我們說話.
告訴我們.
這件事不對.
那件事對.
這件事應該做.
那件事不應該做.
這是正常的聲量對我們說話.
而在我們的痛苦患難裡面.
神會用抗音器的聲音.
麥克風的聲音.
對我們說話.
這是什麼意思呢.
其實很多時候就是.
往往在困難.
痛苦患難裡面.
我們的靈魂甦醒.
我不想大家.
我不想我自己.
不想大家真的會走到這樣的時候.
才靈魂甦醒.
但是很諷刺的就是.
很多時候都是.

$^{521}$不過這個也是給我們一個很好的提醒.
我們不需要走到患難的日子.
我們才靈魂甦醒.
今天是母親節.
我忘記了.
真心的祝願.
每一位母親.
希望你今天能夠和你的家人.
有一個很快樂的慶祝.
很多教會有傳統.
就是在母親節.
是送禮物.
送花給母親的.
神會教會好像沒有這個傳統.
但是我告訴你.
神會教會有一個很重要的傳統.
就是教睦在講台上.
一定會今天給一些.
很真誠的神話的免禮給你.
讓你也知道.
作為一個父母.
你如何得到從神而來的智慧和亮光.
所以今天講的不只是母親.
也是給父親的.
因為到父親節的時候就沒得聽了.
就像媽媽經常說的.
父親節沒人記得了.
今天我給大家.
各位父母從神話語的免禮.
就是在《申命記》第四章那裡.
第四章我給大家看兩節.
這兩節的說話都是.
神藉著他的僕人摩西對當時的父母.
語重心長的說話.
這第九節就說.
請聽.
你只要謹慎.
父母,你只要謹慎.
因勤保守你的心靈.
免得忘記你親眼所看見的事.

$^{561}$又免得你一生.
這是離開你的心.
總要傳給你的子子孫孫.
接著23節.
你們要謹慎.
免得忘記耶和華.
你們神與你們所立的約.
為自己挑剔偶像.
就是耶和華.
你神所禁止你做的偶像.
這是23節的.
大家看這兩節聖經的時候.
你會發覺有些字是重重複複.
很鄭重的.
謹慎,因勤.
這是很鄭重的態度.
還有心靈.
其實跟我們剛才說的靈魂是同一個字.
還有就是免得忘記.
提了三次在這兩節聖經.
免得忘記.
接著.
不要忘記傳給你的子子孫孫.
不要忘記神所立的約.
不要為自己挑剔偶像.
我希望藉著這兩節聖經.
送給父母衷心的勉勵.
很重要的勉勵.
不過我想說一說.
申命記.
如果你剛回到教會.
你對申命記不知道是甚麼.
我也想說幾句.
聖經66卷書.
舊約39卷.
頭頭曲5卷就是.
創世記,出埃及記,利未記.
創世記,出埃及記,利未記,民數記,申命記第五卷書.
創世記是說上帝創造天地萬物.
但是人墮落了.

$^{601}$人在墮落的世界中.
揀選了一群人作為他的子民.
對亞伯拉罕揀選他的族為他的子民.
而接著的四卷書.
出埃及記,利未記,民數記,申命記.
這四卷書都是說這群子民.
他們很不幸地成為奴隸.
受了很多苦.
上帝將他們在奴隸的生活中救出來.
帶他們進入應許之地.
給他們祝福.
就是這樣.
而申命記是最後一卷.
差不多走完出埃及的路.
進入應許之地的時候.
神藉著摩西的僕人.
對這群子民說了一些很重要的話.
也是我今天想給大家的一些美麗.
熟悉聖經的人.
讀過申命記的人都會很明白.
這群選民.
以色列民.
他們第一代是很失敗的.
他們剛出埃及的那群人很失敗.
因為他們神那麼大能大力.
帶領他們出埃及.
但他們在途中有很多很多埋怨.
經常忘記神的恩典大能和拯救.
所以他們時時都埋怨.
因為沒有飲食.
不合胃口.
看不到領袖.
他們就很驚慌很埋怨.
馬上做偶像.
但是神很容忍他們.
一次又一次地給他們機會改過.
而申命記第四章.
在即將進入迦南應許之地的時候.
神命令摩西跟新一代說.
你們爸爸媽媽.

$^{641}$即是第一代.
是很失敗的.
他們時時都在埋怨懷疑我.
但是你們不要再這樣了.
你們要怎樣呢.
就是接著說到這幾句話.
你們要好像聽摩西說.
你們怎樣去對我.
就是你們要殷勤.
免得忘記.
我怎樣對你們.
今天要跟大家說.
摩西這一段的勉勵.
是因為每一位信主的人.
包括你跟我.
還有今天在場的父母.
網上的父母.
我們都是藉著耶穌基督.
成為了神的子民.
我們在我們的信仰旅程裡面.
其實我們都是.
好像以色列人一樣.
我們是做過奴隸.
我們都是在罪的綑綁奴役裡面.
但是神祂自己救我們.
脫離這些奴隸的生活.
脫離罪的綑綁.
進入應許之地.
進入祂祝福的境界.
我們成為子民.
神的子民.
其實是受祝福的子民.
但是我們很多時候.
都有很多失敗的地方.
我們人生遇到挑戰的時候.
我們都會好像以色列人一樣.
埋怨上帝.
我們都會去弄很多偶像出來.
所以今天摩西這兩句話.
絕對不是只是對於當日的子民說.

$^{681}$各位弟兄姊妹.
是對於我們今天每一個人說.
是對於今天每一個父母說的.
我和大家看這兩節聖經的時候.
我也想說出兩點.
第一點就是請你記得.
請你謹記.
將上帝在你身上所做的.
傳給你的子女.
這就是第四章第九節所說的.
進入應許之地的父母.
你們要謹慎保守你的心靈.
你的靈魂.
你的心靈.
這是你和主可以溝通的.
這一個方面.
你的靈魂.
你的心靈.
你要謹慎保守它.
免得忘記將你親眼所看的東西.
免得忘記將這些東西離開你的心.
總要傳給你的子子孫孫.
究竟免得離開.
當時是指甚麼呢.
我和大家看了第四章九節.
至四章二十三節中間一大段.
十節至二十二節那裡.
其實是在說.
上帝在昏黑,默暗,幽暗裡面.
從火焰那裡對以色列民說話.
對於這班不忠不信的人說話.
我現在賜十誡給你們.
和你們立約.
你們是我的子民.
你不要忘記這一點.
這就是你們不要忘記的這一點.
今天的我們.
今天的父母.
或者你沒有去過聖地.
或者你去過聖地也好.

$^{721}$你不一定去過河列山.
就是頒布十誡的西乃山.
但是你去過也好.
你未必見到這個.
聖經說到昏黑,默暗,幽暗的景象.
你未必見到.
不要說未必.
直接是未見到.
就是神在火焰當中對你說話.
但是其實你和我.
見過一些類似的東西.
就是當我們由罪的捆綁.
脫離罪的捆綁.
進入神應許之地的時候.
我們見過黑暗的場面.
也就像昏黑,默暗,幽暗中.
神對我們說話.
就是耶穌基督釘十字架.
在各個地方釘十字架和我們說話.
耶穌用寶血和我們立約.
當時神用十誡和我們立約.
但是耶穌基督用寶血和我們立約.
我們是屬於神的子民.
我們不可以忘記.
我們要切切記住.
不要因為生活中很多很多事情.
忙這個忙那個.
其實是忙一輩子都不會忙完的.
不要因為忙這些.
而忘記將這些東西傳給你的子子孫孫.
不要忘記這件事.
這是你重要的責任.
我們要切記.
使他們靈魂肅成.
第二.
要謹記你是立約的子民.
一生尊神為聖.
就是這節聖經.
你要謹慎免得忘記耶和華和你立約.
當你忘記的時候.

$^{761}$你就會雕刻這些偶像.
其實偶像是耶和華絕對禁止的.
當時在以色列民中間.
他們真的在那段聖經.
十二節十章十二節.
繼續這樣說.
當時的以色列民.
當摩西離開的時候.
他們真的做了很多很多偶像.
他們有男的神.
有女的神.
有野獸的神.
有飛鳥的神.
有星象的神.
他們拜很多神.
因為他們以前接觸過外邦人.
他們是拜這些神的.
所以他們一不信耶和華的時候.
一堆人有懷疑的時候.
他們就拜這些神.
我的問題是.
現在想問大家.
我們今天有沒有偶像?.
我們今天立約的子民.
我們有沒有好像當日的以色列民一樣.
犯了這個毛病.
我們建立很多偶像?.
老實說.
我們今天的日子.
是比以前的日子物質富裕一點.
大部分人都是.
但是物質富裕的日子.
也是我們最受試探的日子.
因為肉體的東西.
身體上的需要.
真的永無止境.
是忙到你不可以停下來.
但是我們往往就是在這些時候.
將這些成為了我們的偶像.
這些本身是好的東西.

$^{801}$但是我們會將它成為偶像.
所以我們將靈魂.
放到我們都不認得的地方.
各位父母.
我想.
你也去想想.
當你今天很努力地為你的子女安排好的教育.
好的節目.
為他們安排很多很好的東西的時候.
你有沒有用足夠的時間.
好好地教導你的子女去認識神呢?.
或者你有沒有好好地和他們分享.
你在神裡面的經歷呢?.
神怎麼拯救你.
怎麼祝福你呢?.
你有沒有教他們去感恩.
愛神愛人呢?.
各位父母.
每一個星期日.
你回來教會.
就像現在這樣敬拜.
很開心見到大家.
但是每一次你回來敬拜之前.
今天回教會敬拜.
是不是你的選項呢?.
今天只做了這個.
好,來吧.
誰知道平時有多幾樣.
就不來了.
就去了第二樣.
會不會是這樣呢?.
各位父母.
你有沒有嘗試過努力.
在家裡為子女做一些家庭的敬拜或者靈修呢?.
或者不可以每一天做.
有沒有試過一個星期做一次呢?.
有沒有試過在家裡培養一個敬拜神.
神在你生命裡是很重要的.
這樣的實習呢?.
各位父母.

$^{841}$願藉著以上摩西對以色列人說的話.
勉勵你.
你將子女.
你關心你的子女.
不單是身體方面.
更加是她的靈魂方面.
你小心.
如果你不做這些.
你今天提供給你的子女.
一切的好東西.
都會成為她生命的偶像.
因為越注重肉體的需要的時候.
就會越忽略神.
而這些東西.
往往使你的子女的成長.
他們會越來越變得迷失.
你買一個好的電話.
一個很好的ipad.
一個很好的電腦給她.
這些會成為她的偶像.
如果你不將她帶到神那裡.
相反.
如果你撫養她一生.
你是著重她的靈魂的需要.
你著重帶領她去尊崇敬拜神的時候.
我告訴你.
你一輩子都不用擔心她.
因為她在神的話語裡成長.
最後我想跟大家說說.
我自己的經歷.
今天是母親節.
我也想起我的媽媽和爸爸.
在物質上.
我們的年代絕對沒有現在那麼好.
大家都知道.
我們這些年代是很窮的.
每個人都捱苦的.
我們一家九個兄弟姐妹.
都在窮苦裡長大.
沒有什麼物質的豐富的供應.

$^{881}$但是我爸爸媽媽.
留給我最豐富最寶貴的遺產.
是什麼呢.
就是在他重病.
在醫院留院的日子裡.
他信了耶穌.
因此爸爸在最殘弱的日子.
他帶了我回教會.
爸爸去世了.
媽媽繼續帶我回教會.
帶我回去.
那個吹口靈氣進我身上.
使我成為有靈的活人.
那位上帝那裡.
祂使我靈魂甦醒.
我一生都可以體會到.
與主同行的那種福氣和力量.
我們一起祈禱吧.
神啊我們為了這一口氣.
讚美你.
你是這樣對我們.
我們是受造物.
我們算得了什麼.
但是你竟然愛我們.
你給我們靈氣.
使我們成為完全不同的受造物.
更加當我們在罪裡面墮落.
離棄你的時候.
你拯救我們.
神啊雖然我們今天.
在我們人生裡面.
我們都有很多的衝擊.
很多的難處.
甚至有很多弟兄姊妹在這裡.
他們是忍受著很多的患難.
但是當我們知道這一口氣的時候.
我們的生命有了新的結果.
就是神你顧念我們.
你重視我們.
而且你拯救我們.

$^{921}$有一天我們會回到你那裡.
而當我們回去之前.
求你教我們明智的.
聰明的好好地生活.
在世界上生活.
敬拜你.
重視我們的靈.
我們的靈魂蘇醒.
我們也使我們後一代靈魂蘇醒.
求聖靈幫助.
奉救主耶穌的聖名.
阿們.
願神的說話.
祝福大家.
\newpage



\section{}
\label{sec:GOoaQbj9lhw}
\textbf{比對研究 14: 巴力宗教 7 - 陳崇基牧師}
\newline
\newline
連結: \href{https://youtube.com/watch?v=GOoaQbj9lhw}{\texttt{https://youtube.com/watch?v=GOoaQbj9lhw}} ~~~~ 語音日期: 2023-05-20
\newline
\newline
\hyperref[sec:_7dHuWNR_kw]{\small{< < < PREV SERMON < < <}}
~
\hyperref[sec:index_chronic]{\small{[返順時目]}}
~
\hyperref[sec:index_scriptual]{\small{[返順卷目]}}
~
\hyperref[sec:i8guCKlYag0]{\small{> > > NEXT SERMON > > >}}
\newline
\newline
$^{1}$我們說比對研究可以幫助我們了解聖經的背景.
我們看了好幾集,或者說了好幾次.
都是和巴力的神話,或者巴力的文獻有關.
我們說在烏加利這個地方,在以色列北面這個地方.
找到一些泥板,這些泥板讓我們更加清楚知道.
拜巴力是一回什麼事呢?以及巴力的神話,巴力的地方.
有什麼影響,或者有什麼主題和巴力宗教有關.
所以我們就靠著這些所謂經外的記載.
這些文字,可以讓我們看得舊約裡面的一些聖經.
可能會更加清楚.
我們一直以來都是在看,用這些巴力的文字.
來看以利亞的故事.
以利亞的故事,大家很清楚有幾章,列王記上17至19章.
就是和巴力很有關.
所以我們就說,將這些文字做一個背景.
嘗試去理解一下,到底以利亞的故事是在說什麼.
我們說巴力的神話.
和以色列巴勒斯坦地的龍更氣候的循環很有關係.
記得我們說過巴力的循環,或者是BioCycle.
或者是巴力神話.
其實就是迦南地的龍更氣候的循環.
每次巴力再生的時候,或者是有活力的時候.
就是雨季的時候,它就會賜雨給地上.
當下季來到了,沒有了雨水,就代表巴力自然是死了.
就不會有雨水,也不會有土產,就會有死亡而來.
所以這個cycle是和巴勒斯坦地的龍更息息相關的.
開始的時候,我們說過巴力這個風雨的神.
就是加雲者.
他控制著雨雲的神明.
他的死對頭有兩個.
在巴力的神話裡面,第一個就是Yam.
就是海神,或者是混沌的神.
你記得烏加利的文獻裡面說到巴力的死對頭.
叫做Yam.
或者是Y-A-M-M.
Yam就是海神或者混沌的神.
在希伯來文裡面,講到地是空虛,混沌,雨面黑暗,水等等.
這些都很多時候用同一個字,Yam這個字.
只不過在創世記裡面,沒有了神話的味道.
只不過是上帝創造的,或者聽上帝控制.

$^{41}$被上帝所管治的,其中一些大地裡面的東西.
所以有些這樣的字眼,其實是有互通的.
所以我們曾經看過巴力的神話裡面.
巴力的其中一個頭號敵人就是Yam.
他將Yam或者海神混沌打敗了.
然後他就帶來order,帶來世間上.
或者農耕或者氣候裡面的生產的規律.
巴力就能夠帶來,有雨水的時候就有雨水.
沒有雨水的時候,只有一段時間.
然後巴力就會再去給雨水給地上.
以致人們可以耕種,然後有土產.
另一個頭號敵人就是Mod.
就是Mod.
Mod在烏加利文,或者甚至在希伯來文.
是死的意思,所以是死亡.
所以這個就是God of Mod,或者死亡的神.
我們中文譯音可以只把音譯出來,Mod.
其實就是死亡的意思,無論在烏加利文或者希伯來文.
裡面都是一樣.
好,我們在巴力神話裡面,我們看看巴力和Mod.
死亡之神Mod.
能夠去看一看,有些文字去表達.
這個死亡之神是怎樣的.
還有巴力的神Bale.
怎樣去跟他交戰.
我們好像之前提過的循環.
每逢到沒有雨水的時候.
就是這個旱季來到的時候.
巴力好像被打敗一樣.
然後他就死在,死了,然後就下到了Mod的陰間.
因為這個Mod是死亡之神,也是陰間之神.
直到這個Mod死亡之神被打敗.
然後巴力才重新復生過來.
所以我們看一看這個文字,就講到Mod是怎樣的.
我們就抽取幾段經,這些文字.
就去到講出到底這個Mod是怎樣的.
在一段文字裡面,這個死亡之神Mod.
就跟這個Anart,就是這個阿納.
阿納就是巴力的妹妹或者老婆.
總之就是一男一女.

$^{81}$就跟他說一件事,他說他是怎樣的呢.
Mod就說他的胃口經常都很餓.
想吞噬人.
要吞噬什麼呢?.
就是地上的眾生.
所以你看到這個死神是怎樣去形容他.
就是死亡就是跟他掛鉤.
就是他有死亡的能力.
他吞噬人的形容,其實只不過是說什麼.
只不過是說這個死亡一定會影響.
或者會去到吞噬每一個地上的人.
所以這個Mod就成為死亡之神.
成為巴力,雨水之神,生育之神,生產之神的對頭.
所以巴力就代表生命.
而死神就代表死亡.
除了這一段講到他的胃口.
要去吞噬人的生命之外.
也講到什麼呢.
這裡就很生動地形容這個Mod神,死亡之神是怎樣呢.
他說他一片嘴唇是向著地.
另一片嘴唇是向著天.
你看到他的嘴巴很大.
舌頭伸向天上的眾生.
其實去到巴力循環的這一段.
就講到巴力被死神吞噬了.
巴力終於死了.
沒有雨水.
所以這個死神Mod就得逞了.
巴力必定會進入他的肚腹.
被他吞噬了.
然後他就落到Mod的口中.
所以巴力就喪生,死了.
因為他是一個死亡之神.
所以近著他的時候.
近著這個Mod的時候.
就會很危險.
因為巴力循環的一些神話的內容.
就會說Mod這麼恐怖的時候.
他是代表死亡的權兵.
你近著他就會很危險.

$^{121}$其餘一段就是這一段.
就是KTV1-4.
就是說如果你太近他.
一定會有死亡.
你一定會死.
因為你太近他的話.
不要太近他.
以免他當你成為他口中的其中一隻羊羔.
你就好像小羊一樣.
在他的鱷的底下被壓碎.
他就好像很大隻的野獸一樣.
這個死亡之神就會將你吞噬.
所以在一些神話的記載裡面.
巴力神話裡面的記載.
就說到這個死亡之神是很厲害.
很有權兵.
連巴力都被他吞噬.
你記得我們曾經提過巴力循環的時候.
每逢漢季完的時候.
完的時候,巴力怎麼可以再復生呢?.
你記得我提過嗎?.
就是巴力死了.
然後他的伴侶阿納為他報仇.
然後就將死亡之神摩特殺死了.
於是就將巴力從陰間的死亡領域裡釋放出來.
以致雨水又再帶到大地.
\newpage



\section{}
\label{sec:i8guCKlYag0}
\textbf{Beholding Your Glory - Rev. Mark Chan}
\newline
\newline
連結: \href{https://youtube.com/watch?v=i8guCKlYag0}{\texttt{https://youtube.com/watch?v=i8guCKlYag0}} ~~~~ 語音日期: 2023-05-22
\newline
\newline
\hyperref[sec:GOoaQbj9lhw]{\small{< < < PREV SERMON < < <}}
~
\hyperref[sec:index_chronic]{\small{[返順時目]}}
~
\hyperref[sec:index_scriptual]{\small{[返順卷目]}}
~
\hyperref[sec:OQpeQNO5TLA]{\small{> > > NEXT SERMON > > >}}
\newline
\newline
$^{1}$And here's the announcement..
Let's welcome Pastor Mark for today's sermon..
Welcome everyone..
And again, I want to remind you on your way out,.
there will be a table, there will be some products, CDs and media products featuring sun music..
And please go pick it up. It's not free..
You have to pay for it..
And I think it's worth for you to support a ministry like this, right?.
To have strong and solid spiritual hymns and contemporary songs, worship songs, serving for the church..
And we want to praise God and thank God for using Brenda in the past many years and writing so many songs..
And Cantonese songs or English songs or Mandarin songs, right?.
Mandarin songs. Wasn't it a beautiful song?.
What was it?.
The Mandarin song I'm talking about. Yeah, the third one..
Okay, those of you who don't know Chinese, it's okay..
It's still beautiful, right?.
It's still beautiful..
And we want a ministry like this to grow and to work with the church and support the church.
so that churches around the world can sing praises to our Lord, right?.
Amen?.
So please do support the Sun Music ministry..
By the way, they have a YouTube channel and go Google it..
Sun Music, S-O-N, Son of God, S-O-N, God's Son, Jesus Christ, S-O-N Music, okay?.
So go and subscribe, all right?.
Go and subscribe, please. Go and subscribe, all right?.
And support her ministry and it's really a bless to have her in the past weekend..
And we got five people last night to come forward and made the decision to come to know Jesus..
Amen to that, right?.
Thank God for that..
And today, I know my time is very limited and I'll make it short..
And I want to start with this question..
It's a big question..
I wonder if you read the Bible, you try to understand the Bible,.
have you ever asked yourself a question, "What does God really want?".
What in the world does God desire the most?.
You know, look at Genesis 1 to the last chapter of the Bible, Revelation 22..
There was a theme underlying all of these books..
Now, there are 66 books in the Bible, 39 books in the Old Testament, and 27 books in the New Testament..
But among all of the books from Genesis to Revelation,.
there is an underlying central theme that runs through every single book, every single passage..

$^{41}$And what does God really want?.
And that theme that underlines and links every single thing together,.
one of those themes is that God wants to make His home among us..
He wants to live among us, among His people..
Whenever we worship, whenever the church gathers,.
we sing praises to the name of Jesus, we exalt the name of the Lord..
When we worship together, God lives among us..
And that's one of the major themes underlying the entire Old Testament and New Testament, entire Bible..
Do you realize that even though there are 66 separate books in the collection,.
what we call Bible, is actually a big story, a big narrative,.
big narrative to introduce to us who God really is..
I want you to pay attention in the next five to seven minutes as I cut this sermon into a very condensed form..
I want to give you a very, very quick summary..
What the Bible is all about, what God really wants, is to live in your life..
God wants to fill your life with His love, His joy, and His peace, and His kingdom..
That's what God really wants. That is what God really desires..
God is not some kind of transcendent concept out there somewhere and doesn't care about you..
God cares about you deeply. God wants to make His home in your heart..
And He wants to make His home in His people, in the churches, in all of God's people..
I want you to pay attention to three circles..
And these three circles is a very, very condensed summary of what the entire Bible is all about..
And you ever use a napkin, a piece of napkin?.
You can put the entire Bible on a piece of napkin..
This is a napkin theology, all right? Napkin theology..
You can just do it during lunch or during dinner..
On the dining table, take a piece of napkin, and you can summarize the entire, entire story of the Bible..
And this three circles is, you have to imagine, the first circle is on your upper left,.
is that God created this world good and orderly..
God created everything, the entire cosmos, the entire universe, out of love..
Do you realize that God is Trinity, right? Three persons in one God, right?.
That's kind of interesting. That kind of defiles our logic, right?.
There is no way for us to wrap around what is going on..
God is one, but He's also three. God is also three, but one, right?.
Isn't that interesting, right? But that is what God in His essence..
He is not a lone person. He is a community, one community, united together, and love one another..
The Father loves the Son, and the Son loves the Spirit, and the Spirit loves the Son and the Father..
And there's a mutual community of the God Himself..
God is three in one, Father, Son, and Spirit..
How many of you have children? How many have children? Yeah..
A few of you, yeah. You should understand, right?.

$^{81}$First, when we first gave birth to the first child, what did we do? What did we do?.
We prepared the crib, we prepared the diapers, we prepared the bottles, we prepared everything..
And we set aside a room for her, and we set aside so many things, and put the child in your first priority..
At least, you know, ten months or so before the birth of the child, right?.
And why do we do that?.
And you have to know, parents are doing this, not only they want to set up the physical space, the house..
God, parents, parents want to do a home..
You know, house and home are different, right?.
We have some brothers and sisters moving around in the summer, in the past year,.
moved from house one to another one, and they bought another house,.
or they rent another house and moved to another place..
Those of you who are students, you move out of your dormitory, and then move to another apartment,.
and then you move from here and there, right?.
Those are house stories, just physical space..
Now, this is what's interesting to us, because we opened the Bible, Genesis 1..
God created the heaven and the earth, the entire cosmos..
And so often, we treat the story, Genesis 1, as a house story..
We think that God is simply just creating all the physical structures, right?.
That's the planet Earth, and you've got different planets and different stars, and sun and moon, and all that stuff..
It's all about material that we think about..
But Genesis 1, the Bible tells us, is not a house story..
Genesis 1 is a home story..
God wants to create humankind so that He can enter relationship with those created in His image..
That is not a house story..
That's a relational story, right?.
Much like parents preparing the arrival of the child..
You know what the child needs?.
It's not the physical space. It's not just that..
It's not just buying stuff for him or her..
It's the relationship..
You want a child to be born in a home that is warm and welcoming..
That is, you find relationship in it, peace in it, love in it..
That's what Genesis is all about..
That's what God's creation is all about..
God created you and me because He wants to enter into a relationship with us..
So that we can worship Him, and we'll be His people..
And God's life will dwell among us..
It will dwell in your heart..
So God created everything good..
And as you can tell, if you trace the entire story of the Bible, things didn't go well..

$^{121}$Quickly, humankind rebelled against God..
They want to say to God, "You are no longer God. I am the God.".
"We are the gods. We are the little gods. We are the bosses.".
"And we take charge of our own life. You create it for us.".
"No, you don't be our God. We will be our own God.".
And that's why they were in exile..
And if you trace the entire story of the Bible,.
God is not giving up His desire to dwell among you..
To dwell in your life..
You trace the story of Genesis, Exodus, and all that..
All the Old Testament books..
God is trying to make His dwelling..
God is trying to make His home in His people..
So after Eden, you got Abraham..
And God called Abraham out to form a nation out of him..
So that through this nation, Israel, everyone in the world will be blessed..
In the sense that everyone, every nation, will be brought back to God..
To His fellowship so that God can live among them..
So that when Israel was saved by grace from Egypt,.
God was telling them, "I'm going to walk with you.".
"I'm going to be present with you.".
"That's why I have my tabernacle among you.".
The book of Exodus..
Chapter after chapter..
Talking about how to build a tabernacle..
A tabernacle is simply a portable tent..
A portable home..
Because they were in wilderness for 40 years..
And God said, "I'm going to be with you.".
"I'll be your God and you will be my people.".
And then finally they arrived at the promised land..
The book of Joshua..
They settled down in the promised land and God fought the battles for them..
And let them settle down in His good land..
And finally, ultimately, they arrived at Jerusalem..
And King David brought the Ark of the Covenant to the capital, Jerusalem..
And then his son, Solomon, built the first temple..
And you read the account of the building and dedication of the first temple..
You find that God's glory filled the earth..
He filled the entire temple because God finally made His home in His people..

$^{161}$It was so glorious and everybody saw His glory..
And of course things didn't go well..
Soon after, the kingdom split..
Right?.
And they abandoned God and they followed the false gods..
They worshipped idols..
And then finally they brought the end to their nation..
So Israel was gone..
And then Judah was gone..
And then the entire people of God was in exile..
However, God did not give up..
God dwelt among the exiled people in Babylon..
After 70 years of exile, they returned back to the land of promise..
They returned to Jerusalem..
And God continued to dwell with them in the second temple..
And after decades and centuries, God finally showed up..
It's Jesus..
Jesus Christ, the Son of God..
He came and dwelt in His people..
So that's the first circle..
That's the first circle..
And we sinned and we did rebellious things..
And we're going to set up ourselves as gods..
And that's why the world has so much suffering..
The world has so much brokenness..
I read an article not too long ago..
It says in the 21st century, mental health has become a crisis and is on the rise..
Now this is nothing new to you..
And the article tried to capture what is going on in our broken world because of sin..
The author said it's brokenness..
You observe around our world, you see brokenness everywhere..
Family is fragmented..
We live in an age, the author said, an age of alienation..
We are aliens to one another..
We are separated. We are broken off from one another..
And most of all, broken from God, our Creator..
And that's why Jesus came..
Jesus came and offered a way so that we can come back to the home of God,.
to the relationship of God..
And whosoever repent and whosoever believe in the name of Jesus,.

$^{201}$our life will be restored..
And that is what the gospel is all about..
And because of the gospel, we are recovered, we are reconciled with our Creator,.
the good home, the owner, the Lord of the home..
And that is what the gospel is..
Gospel of John 1.14.
"So the Word became human and made His home among us..
He was full of unfailing love and faithfulness..
And we have seen His glory, the glory of the Son of God,.
the glory of the Father's one and only Son.".
Pay attention to those words..
"The Word became human and made His home among us.".
By the way, the original text, the original language,.
"made His home among us" is actually tabernacle..
He made His tent, His home in the midst of His people..
That is what the gospel is all about..
If you yield to God, if you repent and return to God,.
and say yes to God, then God will dwell in your life..
And you will see His glory..
And the glory we see is on the cross..
He came in humility and He died in your place..
The righteous one became sin for all of us..
So that He could bear all the wounds and brokenness.
and all the sins on the cross..
And by the cross we are redeemed..
And that's what the gospel is all about..
God wants to dwell in your life, to write His story in your life..
Back in the 1960s, and I used this illustration last night..
Back in the 1960s, the Russians,.
they were successful in sending a man to outer space..
So they sent a man on a rocket, and a man went to outer space,.
and then returned..
So as the Russian returned, he talked to the news outlet..
In the news conference, he said, "Oh, we have been to outer space,.
and we didn't find God, or we didn't find heaven..
So God does not exist.".
That was his conclusion..
Or maybe God did exist, but He doesn't care about us anymore..
He's not there anymore..
So he made this news, this kind of conclusion in the news conference..

$^{241}$Now back then in the 1970s, '60s, in England,.
there was a literature professor, a Christian theologian..
His name is C.S. Lewis..
And C.S. Lewis has been a very good writer, a very creative writer..
And he responded to the Russian, and he said,.
"So this guy being sent to outer space, and then returned,.
he said this guy is looking for God in the outer space..
It's kind of like Hamlet looking for Shakespeare in his own world.".
You know Hamlet and Shakespeare?.
How many of you have heard of Shakespeare?.
You did, right? Yeah..
How many of you read some of the works of Shakespeare?.
Okay, wow, you guys are really educated..
(Laughter).
And one of the main characters in one of the plays.
written by Shakespeare is Hamlet..
It's a character full of life, right?.
A character inside the story of Hamlet,.
and he had to make a lot of decisions, life and death,.
be or not to be, that's the question, that kind of thing, right?.
He went through so much ordeal..
And Hamlet, think about it, in the story of Hamlet,.
in the world of Hamlet, he has no way to get to know Shakespeare..
Because if you read the entire story, the entire play of Hamlet,.
you won't find Shakespeare in Hamlet's story..
So C.S. Lewis was saying,.
"If the guy, the Russian, went to outer space, and then returned,.
he couldn't find God,.
it's kind of like Hamlet looking for Shakespeare in his own world.".
The only way Hamlet can get to know Shakespeare.
is for the author himself to write himself his story.
into the story of Hamlet..
That is the only way..
That is the only way for the character to get to know the author.
is for the author to take the initiative.
to write himself into his story..
You know, if you think about it, what C.S. Lewis has said,.
it's actually the gospel..
The gospel tells us God enters our world in Jesus.
and meets us where we are..

$^{281}$He came to live among us, and God literally writes his story,.
writes himself into your story..
Jesus is willing to humble himself to be known by you.
so that you can have a relationship with God..
And that's what God desired from the original plan..
God's desire has always been to live among us.
so that we can worship him..
John 1:14, "So the Word became flesh,.
became human just like us, and made his own among us.".
I tell you, the world needs to see the glory of God..
How do they see it?.
See through you that the life of Jesus can live inside of you..
And the church has a lot we need to do..
There are so many things we need to do together..
When we worship together, when we live in Christ together,.
we learn how to love one another,.
we learn how to fellowship with one another,.
to love each other as Christ loves us..
In this community, then God will be shown..
Last night we had a sister, Lily..
She was a devout Buddhist..
And it's just amazing to hear her story again and again,.
how she came to our church and how she came to know Christ.
because of the community that live out the gospel..
And let's do that..
When we worship, when we gather,.
when we form a community in Christ,.
Jesus lives among us and God made his home in the midst of us..
Let's all pray together..
And I ask the praise team to come here to sing the last song as a response..
Father, we give you thanks because of your desire to be with us..
And you humble yourself in your son, Jesus Christ,.
to write yourself in our story..
Thank you for being our living God..
May our community continue to display your glory,.
to tell the world that you are real,.
to tell the world that you are in our midst..
Thank you. In Jesus' name, amen..
[BLANK AUDIO].
\newpage



\section{}
\label{sec:OQpeQNO5TLA}
\textbf{看見...你的榮耀 - 陳崇基牧師}
\newline
\newline
連結: \href{https://youtube.com/watch?v=OQpeQNO5TLA}{\texttt{https://youtube.com/watch?v=OQpeQNO5TLA}} ~~~~ 語音日期: 2023-05-22
\newline
\newline
\hyperref[sec:i8guCKlYag0]{\small{< < < PREV SERMON < < <}}
~
\hyperref[sec:index_chronic]{\small{[返順時目]}}
~
\hyperref[sec:index_scriptual]{\small{[返順卷目]}}
~
\hyperref[sec:jG3lPKnsQeA]{\small{> > > NEXT SERMON > > >}}
\newline
\newline
$^{1}$各位早晨.
今日我地講到敬拜.
因為呢一個禮拜都係講敬拜.
音樂的敬拜.
敬拜音樂的報道會等等.
我唔知你呢你呢個教會有幾耐.
就算你係今日第一次踏入教會.
或者第一次所謂接觸聖經呀.
聽緊耶穌基督的信仰呀等等.
你可以問一個問題.
一個所謂大ge 問題.
就係到底上帝想要.
教會成立.
或者教會成立.
就係到底上帝想要乜ye .
上帝的desire.
上帝的盼望.
上帝的渴望.
上帝的渴望係乜ye .
佢需要ge 係乜ye .
或者佢最想.
發生ge ye .
係D 乜ye ge 事情.
係我地未開始敬拜ge 時候.
之前呢我地就同D .
弟兄姊妹傾偈.
就講開呢BB仔出世ge 事.
我地知道有D .
已經出世啦亦都有D 就黎出世啦.
幾好呀係咪.
多D 黎呀多D 黎嘛D 時候.
細路仔出世我就諗起.
gwo 時.
我第一個ge 細路出世ge 時候.
咁就有預產期啦當然.
咁就懷jor 佢.
十個月係咪.
十個月唔係我懷呀當然.
見住佢懷.
即係我老婆懷住佢十個月.

$^{41}$十個月裡面呢我地就去.
預備屋企.
預備屋企ge 房.
家屍.
預備買D 乜ye 買D 乜ye 啦係咪.
準備D 乜ye .
等等.
有好多ye 要預備.
係個house裡面.
但最重要.
預備ge 係乜ye 呢.
唔係個house.
係個home.
大家明白house同home有分別.
house只不過係D .
架槍.
你要煲D 乜ye .
你要整D 乜ye .
有個鑊鏟係度.
有D 咩工具.
或者需要個細ge 床比佢.
或者需要car seat比佢等等.
呢D 全部都係housege items.
唔係home.
你問我一個做父母ge .
我最希望.
我ge 仔女出世誕生之後.
我會做乜ye .
not a house.
a home.
一個家庭.
一個和諧ge 家庭.
一個愛ge 家庭.
如果你咁樣問ge 時候.
你會發現其實上帝都係想做呢樣ye .
我地好多時.
睇聖經我地真係水過鴨背.
我今日想同你.
睇晒成本聖經.
由頭到尾睇.

$^{81}$我地想搵一搵.
諗到底上帝想做D 乜ye .
創世記第一章.
天地萬物.
被創造.
上帝做天做地.
做海做世界裡面.
所有ge 萬物.
第一天有光有暗.
第二天有天空有海.
第三天有乾地.
第四天有太陽有月亮有星宿.
第五天有天上ge 飛鳥.
海裡面ge 魚.
第六天所謂最高潮.
就係創造地上ge 萬物.
同埋人類.
按照上帝ge 形象.
你都可以做男做女.
做出人類出黎.
點解要咁樣做.
因為萬物裡面.
萬物之靈係人類.
只有人類先能夠.
同上帝有一個獨一ge 關係.
係咪.
呢個關係.
其他動物係唔會有ge .
我知道我地中間有好多係.
狗奴或者貓奴.
我知道佢地係好可愛.
但佢ge 關係.
只能夠去到.
某一個地步.
而我地同上帝ge 關係.
佢係可以好親近好親密.
但點解要.
創造天地萬物.
最重要ge .
上帝唔係做一個house.

$^{121}$唔係做一個物質ge 世界.
然之後有晒所有ge 結構.
所有ge 家屍.
所有ge 裡面有ge .
嘉環.
然之後就掟D 人係度.
上帝唔係要做一個house.
上帝要做一個home.
我同你講過.
以前有一個好.
有智慧ge counselor.
同我講過.
佢話幾時係最好去到有細路.
幾時.
佢話如果.
兩個人相愛到一個地步.
佢地愛人咁滿舍.
好似兩個杯.
佢地能夠滿舍.
滿到滿日.
能夠倒出黎ge 時候.
呢個時候就最好有新生命.
誕生ge 時候.
因為gwo 個新ge 生命.
就可以同你分享到.
你所滿出黎ge gwo D 愛.
其實home.
上帝創造萬有.
都係一樣.
上帝係三位一體.
我地解你解佢都解唔清.
講理論又好講邏輯又好.
咩叫三位一體.
其實最簡單就係上帝.
就係愛.
有冇諗過神是愛.
我地成日講神是愛係咩意思.
神是愛就因為上帝.
本身ge 本體.
就係一個愛ge 本體.

$^{161}$愛ge 本體本身愛ge 時候.
就一定係多過一個ge .
所謂位格.
多過一者.
愛者.
維繫愛者同被愛者.
中間ge 愛.
一定係有呢幾樣ye .
或者可能三樣.
上帝本身ge 本體.
就係一個愛ge 本體.
當佢創造呢個世界ge 時候.
唔係純粹set up一個house.
係set up一個home.
係與人類.
同在.
上帝想做D 咩ye  上帝想住.
係人類ge 中間.
同我地建立關係.
係一路以黎有ge .
心意黎ge .
由第一章開始就係咁.
我頭先講到話.
第六日ge 時候.
所謂最高峰就係創造人類.
其實gwo 個都唔係最高峰.
最高峰係咩.
就係我地覺得好悶ge 第七日.
神.
天地萬物都做齊了.
佢就停止jor 佢ge 工作.
然之後佢就安息了.
安息了.
佢就係聖經ge 文化或者語言裡面.
其實係上帝.
坐係聖殿裡面ge 意思.
整個世界就好似.
一個大ge home.
係上帝.
所創造ge 神聖ge home.

$^{201}$一個ge 家庭.
佢最後第七日.
就係坐係.
佢ge home裡面.
佢去作王.
係人類中間.
佢去到同佢地有.
fellowship.
能夠有建立關係.
只有人類 只有你.
同上帝.
先可以有gwo 種ge 關係.
只有我地 先能夠同上帝.
有gwo 種ge 關係.
呢個創世記開頭ge 時候.
話比我地聽 到底上帝要D 乜.
上帝ge 心意就係要.
住係你ge 心靈裡面.
住係佢子民ge 心靈裡面.
呢個先就係最重要.
上帝想要發生ge ye .
所以創造一切ge 時候.
創造一個home.
到第二章.
我講成本聖經.
而家第二章.
創世記第二章 伊甸園.
伊甸園我地諗諗下.
都係好似我地北美ge .
房屋後面gwo D ge .
後yard 咁樣.
間唔中我地去剪下草 間唔中我地去種下ye .
No 我考你.
伊甸園向邊ge .
門口 伊甸園有門口.
你都知o下嘛 你後yard都有門口.
伊甸園ge 門口.
向邊.
向東.
你返去check下 向東.

$^{241}$創世記第三章.
有兩幾路柏守住門口.
係咪 伊甸園東邊.
點解向東.
所有ge 聖所.
係聖經世界裡面.
都係向日出之地.
向東 伊甸園裡面.
有D 咩 有金 有銀.
有寶石 有好多.
好美ge 植物.
有樹.
有水.
河水.
呢D 全部ge 意境.
都係聖所ge 意境.
係咩地方.
有金有銀 然之後又有天使.
然之後又有.
花草樹木 有河水.
有水源.
係咩地方.
你翻下聖經.
聖所 聖殿.
就會有呢D ye .
所以諗返起.
原來上帝創造一切ge 時候.
佢就係住係.
人類ge 中間.
人類就係去敬拜佢.
人類就係咁樣去事奉佢.
敬拜.
原來就係天職.
你唔好以為你入黎.
我地唱下幾首歌.
ok 等暈jor 個時間.
15分鐘.
17分鐘.
好似呢D 聚會要造勢咁.
我地warm up jor 個場先.

$^{281}$跟住最尾就點呀.
有個人出黎講道.
呢個先係戲肉.
No.
敬拜.
係我地ge 天職.
上帝要住係.
人類中間.
代表佢榮耀能夠充滿係人類中間.
佢榮耀能夠.
淋到佢ge 子民.
所以由頭到尾.
開始ge 時候上帝ge 爭議就係咁.
但係大家知道整個.
聖經ge 故事就唔係咁樣.
完美地發生ge 係咪.
因為人.
半變.
人半變.
人去得罪上帝.
人想離開上帝.
人想話我做上帝.
你唔好做.
結果呢 佢地就被流亡.
人類就被流亡.
離開jor 異地.
就缺少jor 分裂jor .
上帝ge 關係.
但上帝.
如果你係上帝 你會唔會咁做.
Game over.
然之後reset.
Wipe outjor 佢.
算啦.
如果係咁樣ge 話.
如果呢個上帝係咁樣ge 話.
我都唔想信呢個上帝.
聖經就話比我聽呢個上帝.
係有憐憫.
係有恩憐.

$^{321}$創造jor 人.
一個最perfectge home比佢.
佢半變.
離家出走.
上帝去邊.
上帝當然去等佢.
我唔止比喻你.
或者應該更加.
講得更加貼切.
如果你繼續去到追蹤.
上帝去尋找人.
上帝去搵返迷失jor ge 人.
唔願意住係佢ge home.
鍾意出去流亡ge .
鍾意出去訓街.
上帝就要去.
尋索呢D 人.
去搵返佢地.
然之後去到.
救贖返佢地.
挽回返佢地返去.
創世記一路開始第三第四章.
就係講緊上帝有D 乜ye 計劃.
去到成就佢.
一直以有ge 心意.
佢ge design就係要住係.
人ge 生命當中.
住係佢ge 子民.
如果你咁繼續落去你fast forwardge 話.
你就去到.
上帝呼召阿巴拉罕出黎.
上帝叫阿巴拉罕出黎做咩.
佢話你移民.
而要居去一個地方.
就係當時世界ge 中心.
我要使你成為大國.
以至你能夠.
將我ge 計劃.
帶到去萬國.
Right.

$^{361}$你就能夠成為不成.
你就成為祝福.
然之後萬國都能夠.
離返上帝ge 控.
所以你有阿巴拉罕.
有以撒有亞國然之後有以色列國出現.
點解會揀選以色列國.
因為以色列國.
唔係淨係操身敗格.
係操身.
for 格.
係for上帝自己本身有ge .
Home plan.
佢ge Home plan就係要想萬民.
能夠黎到佢面前.
同佢建立返呢個關係.
能夠同佢係埋一齊.
能夠去敬拜.
如果你trace呢個storyge 時候.
你好快就去到.
出埃及.
呢班神ge 子民.
或者阿巴拉罕ge 後裔.
佢地聯絡到去一個流亡ge 地步.
就係圍爐之地.
做奴隸.
係埃及.
上帝用恩典拯救佢出黎.
唔係因為佢地好.
係因為上帝.
應許過話我會救贖.
我ge 人民.
所以佢就係埃及.
水深火熱ge 當中去拯救佢地出黎.
恩典拯救佢地出黎.
然之後做D 乜ye .
我地諗諗下.
我地成日覺得出埃及就係.
Liberation.
我地就係自由.

$^{401}$我地就唔使人管.
唔係喎 出埃及ge 故事唔係咁.
離開jor .
克制佢地ge .
Oppress佢地ge .
法老王.
歸向一個新ge 王.
就係耶和華自己.
就係上帝自己.
出埃及機一個好大段.
好大段交代.
就係上帝創造會幕.
佢話俾摩西聽.
你做一個會幕.
做一個營.
其實會幕就好似.
鍾意呀.
有邊個唔鍾意.
唔鍾意ge 你就好難經歷到.
到底乜ye 教會幕.
我都唔鍾意.
會幕係乜ye .
係一個帳幕.
係一個帳幕上帝.
set up係佢ge 人民ge 正中間.
係以色列國.
周圍十二支派.
就係gwo 附近去扎營.
圍繞住中間gwo 個ge .
所謂Tabernacle.
gwo 個ge 會幕.
上帝要住係佢地裡面.
然之後帶住佢地.
不斷走唔同ge .
旅程.
係曠野裡面俾ye 佢地飲.
俾ye 佢地食.
會幕就係一個home story.
就係上帝話我要住係你地中間.
雖然你地鍾意流亡.

$^{441}$但唔緊要我同你同在.
你係我ge 子民.
我係你ge 神.
呢個就係Tabernacle.
所以佢地每逢set up一個地方ge 時候.
佢地就set up個Tabernacle先.
因為上帝要住係佢地中間.
然之後佢地要再起行ge 時候.
佢地就有其他人gwo D 聖職ge 人員.
佢地就將個會幕收起黎.
我地就去另一個地方.
好似hiking咁.
然之後去到第二個地方佢地又停留.
gwo 度有水源 佢地就停留係個樹.
一段時間.
然之後一步一步ge 上帝帶領住佢地.
我要帶你返去屋企.
gwo 個就係加林地.
加林地就係摩西五經完jor ge 時候.
然之後約書亞記.
開始.
就係上帝話我帶你入去應許之地.
Promised Land.
就係所謂我地講ge Holy Land 聖地.
係gwo 度入jor 去之後.
上帝話唔緊要我同你地去爭戰.
我為你打勝仗.
打贏所有ge .
壓迫你ge 仇敵.
然之後我就會set up.
我ge home again.
係你中間.
結果有段好大段ge 故事係舊約.
就係講到.
佢地點樣入去加林入主加林.
然之後有好多仇敵去到壓迫佢地.
佢地有好多軟弱.
有好多fall and stumble.
然之後佢地又跌低又遠離上帝.
但上帝冇離開過佢地.

$^{481}$到最後立jor 一個旺叫做大衛.
大衛就話係時候.
上帝使我地.
國境太平.
我地過以色列家.
終於可以落地生根.
呢個時候.
要做D 乜ye .
就要將個約櫃帶黎到.
京都耶路撒冷.
將個約櫃.
Art of Covenant.
個約櫃其實係乜ye 黎架.
係上帝gwo 個ge 座駕.
係上帝gwo 個.
上帝gwo 個ge 龍椅.
上帝gwo 個ge .
Air Force One.
將佢拍翻入黎.
黎到首都.
耶路撒冷.
將個約櫃帶到黎.
耶路撒冷ge 時候.
大家歡呼喝彩.
又係敬拜ge 故事.
因為上帝一路有ge 心意.
就係要同佢ge 子民同在.
然後佢ge 子民回應佢.
用感恩黎到回應佢.
獻上佢地ge 祭.
然後去敬拜上帝.
敬拜佢地真正ge 天父.
約櫃黎到.
然後下一任ge 王.
就係所羅門.
所羅門就建聖殿.
聖殿建好ge 時候.
你記唔記得發生咩事.
gwo 個係The Perfect Day.
Of the Old Testament.

$^{521}$成本舊聖經裡面.
最所謂高潮ge .
起碼人眼中最高潮ge .
就係上帝永要住係聖殿裡面.
我地一諗起聖殿.
我地就諗起好似好奇怪.
ge 一個building.
其實係一個homege building.
上帝要make himself a home.
gwo 個home就係一個temple.
就係一個聖殿.
所以上帝充滿呢個聖殿ge 時候.
gwo 時就濃煙密佈.
然後有好多榮耀.
連祭司都入唔到去.
呢個叫聖殿.
呢個叫上帝一路有ge 心意.
就係住係佢自己ge 子民裡面.
佢成為佢地ge 上帝.
佢地去到敬拜.
去事奉佢.
當然好景不常.
你記得gwo 個故事.
所羅門之後就係分裂jor 南北角.
然後佢地有大段好混亂ge 時間.
上帝就係南北角.
佢地就係南北角.
有段好混亂ge 時間.
上帝又派先知去同佢地講.
喂大佬上帝已經對我地好好.
上帝已經俾我地好多好多ge 恩典.
俾我個perfect home俾我地.
點解我地仲係係度嗌交.
點解我地仲係reject佢.
或者否定上帝.
要住係我地中間.
唔想上帝去到作我地王.
先知一個一個ge .
被派去.
就係去話比佢地聽.

$^{561}$你回歸 你要返返黎.
上帝同我地立ge 約.
同我地守ge 約.
咁我地先可以.
繼續可以係個home.
但結果佢地無.
北角被滅.
南角亦都被滅.
主前586年.
巴比倫軍隊黎到上帝ge .
homege 時候.
係耶路撒冷ge 聖年.
佢地去射箭.
掟石.
用火焚燒.
就被毀.
上帝gwo 個ge 架.
就玩完.
因為無我食.
我地而家去到邊度.
係聖經.
聶王記下最尾gwo 幾章.
好快.
聶王記下最尾幾章.
因為上帝玩完.
如果你係神.
你會點做.
game over.
唔好搞.
但上帝唔係咁.
上帝依然同.
gwo D 流亡jor ge .
神ge 子民.
同在.
佢個座駕無因為.
聖殿被毀.
而被毀滅.
佢座駕反而去jor 邊度.
去jor 被擄ge 子民當中.
以西結書第一章.

$^{601}$上帝ge 榮耀仍然ge .
係佢自己.
流亡ge 子民裡面.
佢地係黏埋一齊.
有70年左右ge 時間.
佢地係過住異鄉流亡ge 生活.
然之後上帝.
又帶返佢地返黎.
返返去耶路撒冷.
去到重建返gwo 個ge 上帝.
gwo 個ge 聖殿.
gwo 個ge home.
就冇以前咁好.
冇咁多榮耀.
但係都係一局ge home.
上帝要同佢自己.
子民一路黎到同座.
教導佢地.
派先知係佢中間.
引導佢地.
以至佢地真係做返.
神子民應該做ge ye .
就係敬拜上帝.
好奇怪.
gwo D 被擄回歸ge 子民.
開始聽話.
冇咁多偶像.
不過仲係好遠.
過jor 好幾個世紀.
好幾個世紀.
你成本舊約聖經睇完.
今日好快.
15分鐘你睇哂成本舊約聖經.
完啦.
但係你記得舊約聖經裡面.
成日有預言.
先知會講預言.
先知話上帝.
有ge promise 佢ge 承諾.
係一定會實現.

$^{641}$等.
係有根據.
因為上帝就係一個首約ge 上帝.
等jor 好幾百年.
結果呢.
真係有個先知到.
gwo 個就係施洗羊.
John the Baptist.
我地最近英文堂.
開始團契.
我地就查經.
我地查馬可夫.
馬可夫第一章你記得講咩.
一開波就有.
所謂施洗羊上台.
佢就話你呢你呢.
上帝ge 救主.
上帝差黎ge 僕人.
要駕臨.
我唔係gwo 個.
我做佢先鋒.
我係做個announcer.
我去報告下介紹下佢.
然之後佢就逃.
呢個就係耶穌基督.
佢話科罕第一章.
第十四節.
佢話道成了肉身.
住在我門中間.
匆匆足足ge .
有恩典有真理.
我地親眼見過.
佢ge glory.
佢ge 榮耀.
正是父獨生子耶穌基督.
ge glory.
親手捉過佢.
親手見過佢摸過佢.
我地見過佢ge glory.
所以點解我地敬拜佢.

$^{681}$因為佢住係我地中間.
by the way 佢住字.
就係Tabernacle呢個字.
就係會幕.
耶穌基督黎到我地中間.
道成了肉身.
好似我地一樣咁樣.
有血有肉ge 軀體.
一個ge 載體.
載住乜ye 呢.
唔係罪.
唔係苦難.
真理.
與恩典.
住係我地中間.
to make it a home.
係我地中間.
每當上帝ge 子民.
肯聽上帝說話.
肯放低自己.
自己唔做上帝.
上帝話.
上帝你想點.
Yes.
我願意.
你要我敬拜你.
願意.
你要我赦免.
我願意.
你要我接納.
Yes 我願意.
點解.
唔係因為你強權.
我一定要咁聽你話.
因為上帝愛我地.
愛到一個心存.
我地要去回應佢.
每一個信服ge 回應.
係一個Thank you.
ge 回應.

$^{721}$去彰顯你自己榮耀.
榮耀係充滿恩典.
充滿真理.
有冇諗過.
每一個人生係真係有十字架.
就係一個願意.
去寬容.
願意赦免.
願意為其他人行多一步.
ge 呢個生命.
一個愛ge 生命.
一個有家庭.
Homege Story.
當你咁做ge 時候.
你就會知道上帝住係你中間.
你諗過每個教會.
每個群體每個家庭.
只要我地知道有十字架.
係我地中間ge 時候.
就會充充足足有恩典.
同有真理.
耶穌基督就Show up.
但相反ge 時候.
就唔係耶穌基督Show up.
就俾魔鬼留地步.
點解.
因為由頭到尾.
創世ge 第一章.
上帝有ge 心意.
就係要一鬼.
永恆ge 愛.
住係你生命裡面.
使你有平安.
有和諧.
你同上帝講.
I do.
上帝先會住係你中間.
你就見到上帝ge Glory.
又話第一章.
到城你獄真.

$^{761}$住係我地中間.
充充足足有真理.
教會如果真係明白.
整個大ge 故事.
就係耶穌基督住係我地中間.
一個敬拜ge Story.
人黎到ge 時候.
就會見到耶穌基督係真真實實ge 主宰.
我地請敬拜隊載上黎.
帶領我地敬拜.
記住我地去到敬拜.
回應我地ge 位主ge 時候.
係一個Home Story.
上帝住係.
佢ge 子民敬拜ge 歌聲裡面.
上帝住係佢ge 子民聚集ge 地方.
一個充滿有愛.
有恩憐有真理.
就有上帝生命.
係我地中間.
求主與立我地敬拜.
(下集待續).
\newpage



\section{}
\label{sec:jG3lPKnsQeA}
\textbf{People Who Worship God (Psalms 95) - Rev. Mark Chan}
\newline
\newline
連結: \href{https://youtube.com/watch?v=jG3lPKnsQeA}{\texttt{https://youtube.com/watch?v=jG3lPKnsQeA}} ~~~~ 語音日期: 2023-05-29
\newline
\newline
\hyperref[sec:OQpeQNO5TLA]{\small{< < < PREV SERMON < < <}}
~
\hyperref[sec:index_chronic]{\small{[返順時目]}}
~
\hyperref[sec:index_scriptual]{\small{[返順卷目]}}
~
\hyperref[sec:wLYSrSGINOw]{\small{> > > NEXT SERMON > > >}}
\newline
\newline
$^{1}$from the word of God..
- Thank you, James, thank you..
Well, as you can see, I'm ready to fly,.
and the plane going to Israel for me.
is around five o'clock, I guess..
I haven't checked, but we'll be there..
And there was a connecting flight in Vienna,.
and we're gonna stay there, lay over for three, four hours,.
play some music in Vienna,.
and we're ready to go to Tel Aviv, all right?.
So yeah, thank you for praying for us,.
and I wanna say a note of thanks to the church.
and the leaders and the pastoral staff, especially,.
because all of the things that they are doing,.
trying to fill the gap while I'm gone,.
and it's very helpful,.
so that I can be sent to go to Israel.
to lead a group of people..
And this round, we have 27 people in total..
Only 10 people are from our church,.
and all the rest of the 17 are from Hong Kong, Canada,.
a lot of different places, yeah..
Actually, a few of them already arrived,.
so thank God, all right?.
Thank God they already arrived safely, all right?.
I just got a text not too long ago..
One of our members forgot to bring her credit card,.
and she had to go back home and grab the credit card.
before going to the airport..
Yeah, that is very important..
Trust me, going to Israel, you only need two things..
One is the valid passport, US passport,.
valid US passport or visa,.
and second thing is a credit card, all right?.
That's all you need, all right?.
You don't have to bring anything, just two things, right?.
Just a passport and a valid credit card,.
and then you'll be good, yeah..
And yeah, I'm looking forward to this trip,.
and by the way, I'm not a tour guide..

$^{41}$I don't charge, I don't get any tips, you know?.
I don't live off--.
- You should..
- I should, right?.
I should, yeah, I should, I know, but I don't take it..
I don't take it..
I do it for, you know, as a service,.
as a teaching service to other people..
So God has blessed us with some resources,.
and we can share with some other sisters and brothers.
in Christ, some other people who are interested..
In fact, this round, we have 27 people..
One person is a non-Christian,.
and he's just interested to go to Israel, right?.
Who's not interested to go to Israel, right?.
Yeah, it's exotic, exotic place, yeah, exotic place..
Yeah, it's gonna have a lot of fun..
I'm gonna buy you something, okay?.
I'm gonna buy you some ajava hand cream, okay?.
So you come back Tuesdays, I'll have a big bottle.
of ajava dead sea, dead sea extract hand cream, all right?.
How about that?.
Anyone?.
No, you don't want it?.
Do you want it or not?.
All right, all right, okay, I'll bring you something..
I'll bring you something, all right?.
I'll bring you something we can share, okay?.
Lotion is welcome, all right?.
Anyways, I may bring you a bottle of water,.
you know, holy water from Jordan River..
How about that?.
So yeah, right?.
And we can take that in the chameleon, right?.
Everyone can take a sip..
It's good for your health, yeah..
Anyways, we'll continue with our topic on worship..
We talked a little bit about worship last week.
when Brenda was here, Sun Music was here..
And I think it's appropriate to continue to think.

$^{81}$about what worship is all about..
And I think in our generation, worship is a lost art..
You know, if you think about it,.
when you come to church to worship,.
when you come to worship today,.
what did you do to prepare yourself?.
What did you do to prepare yourself to worship?.
All right, talk to your neighbor..
How did you prepare for today's worship?.
What did you do?.
What exactly did you do, you know, okay?.
Maybe a cup of coffee,.
maybe, did you brush your teeth?.
Yes, I hope so..
You change, and you walk into a place of worship,.
and you wanna offer praises to God..
And as I said, worship is a lost art today..
We take things for granted..
We often approach worship with a mindset.
that we don't need to prepare..
There's nothing for us to do.
to prepare ourselves to worship..
When we think of preparation,.
we think of someone else is doing the preparation, right?.
The worship band, the worship people,.
they gotta prepare, right?.
Did they practice?.
Well, you can tell, right?.
They did practice, right?.
Couple days ago, they practiced here,.
and that's why we have a worship set, right?.
And in the backstage, they got the lyrics on,.
they got the cameras on, they got everything,.
PowerPoint and stuff like that,.
and the temperature of the room is just.
to a comfort of my level..
There are many things we can think about,.
preparation of worship, but so often,.
we shift the focus to someone else.
doing the preparation for us..

$^{121}$And for us to come to worship,.
we often think that we don't need to do anything to worship..
You need to understand, some of the ancient architecture,.
when the ancient people, they built a temple,.
a place of worship, a sacred place, or ancient church,.
they always have stairs going up to a temple.
or a church building..
And so often, you realize the architecture.
is a little interesting, because the steps going up.
to the place of worship, the sacred place,.
is often uneven steps, some longer, some shorter,.
and then some longer, some shorter..
What are they trying to do?.
They're actually trying to design the building.
in such a way that when people, worshipers,.
approach the sacred place, they will approach.
with care and slowly..
And that's why the steps are uneven at times..
Those are called processional steps..
So you have to be careful..
Not every step is even..
And you have to take care when you go up to the temple.
or go up to a sacred place..
It's a matter of preparation..
I have been to Israel so many times, since 2008..
Sometimes I go there two times a year,.
sometimes once a year..
I was fascinated by another culture,.
not so Western culture,.
how they approach this sort of sacred place,.
worship place..
I remember one time, I was, a couple years ago,.
before the pandemic, I took a sabbatical.
a couple weeks off from the church,.
and I went with my wife to Israel,.
and we were on the Temple Mount..
You know the Temple Mount, the golden thing?.
The golden dome?.
Yeah, the Temple Mount..
So we went up to the Temple Mount,.

$^{161}$and then, of course, it's a beautiful place, right?.
And then we were trying to take pictures,.
and my wife was next to me,.
and I put my hand on her shoulder,.
and then trying to take a selfie, right?.
And then the security guard hit my hand..
And say, "You don't do this here.".
I said, "What?".
He said, "You don't display affection or intimacy.
"on the Temple Mount.".
I said, "What?".
You know, I've never been to Israel..
It was our first time going together as a couple, right?.
I've never been to a person I can put my hand around, right?.
So I just want to seize the moment, right?.
So I just take a selfie, and hit my hand really hard,.
and then say, "Give me your camera, or give me your phone..
"Delete that picture..
"Delete it right here, delete the picture.".
I said, "What?.
"Okay, don't shoot me, all right?.
(audience laughing).
"I'll delete that picture..
"Too bad, right?".
And there was another time,.
you know, I was doing some YouTube video and all that..
I'm doing my selfie, right?.
I have my selfie stick,.
and just outside of an ancient church building,.
I'm trying to record what this is,.
and give a short introduction,.
and then slowly I walk from outside.
into the ancient church building..
I walk in, and it was so quiet, because no one was there..
And I keep filming, you know, the ceiling,.
and all the side wall, and the church environment..
And then a monk came after me and hit my head..
(audience laughing).
Hit my head right here..
I said, "What is going on?".

$^{201}$He said, "You don't wear hat in the church, all right?.
"You take off your hat.".
I said, "Okay, man, what's the big deal, right?".
It's very interesting, right?.
You go to a Christian church in Israel,.
you take off your hat,.
but you go to a synagogue,.
or Jewish, like Western Wall, Jewish sacred place,.
you have to put on your hat, or a kippur, right?.
The kippur, the Jewish, right?.
The Jewish man..
You know, my point is,.
I'm just fascinated by, in other culture,.
how they prepare themselves to worship, right?.
Now, of course, this sort of things,.
we modern people, or Western people,.
we try to laugh at, right?.
"What's the big deal, right?.
"What's the big deal?.
"I can wear a tank top, right?.
"To the church, right?.
"Yeah, right?.
"It's democracy..
"It's freedom.".
No, you don't do that, right?.
So, as I prepare the people, 27 of us, to go to Israel,.
we had some Zoom online sessions.
to prepare ourselves to go..
And one of the most frequent asked questions is,.
you know, what to wear, right?.
What to wear?.
You know, like this, what I'm wearing,.
I can go to anywhere, sacred place, synagogue, churches,.
any place, I can go, right?.
But you can't wear tank top..
You can wear shorts, go into a church or a synagogue..
So, it's a matter of preparation, right?.
Sometimes we've lost that..
We lost that kind of, you know,.
the art of preparing ourselves to worship..

$^{241}$Now, if you think about preparing for worship,.
preparing ourselves to go into the presence of God,.
there are many scriptures you can think of..
Probably you can think of what Jesus said.
in Sermon on the Mount..
In the sermon, Jesus said,.
"Before you offer sacrifice in the temple,.
"you have to be in a reconciling with your neighbors.".
Right?.
You gotta do something to prepare yourselves..
And one of the things you need to do.
is to be at peace with your neighbors,.
as much as possible, right?.
As much as possible..
The best of your ability to keep peace with other people.
before you approach God to worship..
Now, one of the Psalms,.
there are 150 Psalms in the Book of Psalms..
One of the Psalms that taught us about worship is 95..
Okay?.
Psalm 95..
It's equally divided..
Very short, 11 verses, Psalm..
First half, one to seven..
First half right here..
It's about who to worship and why we worship..
And the second half is about the voice of God, right?.
Now, for worship song, it's very typical..
It's very typical to begin with some call to worship..
Call to worship..
Now, what call to worship means.
is that some sort of phrases or some things.
that the Psalmist singing.
or calling attention to the people of worship..
Call for attention..
It says what?.
Come, let us sing to the Lord..
Let us shout joyfully to the rock of our salvation..
Let us come to Him with thanksgiving..
Let us sing Psalms of praise to Him..

$^{281}$Can you trace how many,.
can you point out how many verbs there are?.
Come, right?.
Literally it means walk, right?.
Walk, physically you walk..
You move yourself to some place..
And then singing, shouting, right?.
Shouting in both Psalms is often victory, war cry..
Victory..
God is the victory..
God has won the war for us..
And it requires you to have some energy, right?.
Shout, right?.
It's an active verb..
It's the active things that we do..
And the call to worship so often.
will tell us why do we worship..
The reasons for worship..
Why do we come here?.
Have you ever asked yourself this question?.
Why do we come here, right?.
Why do we come?.
It's more than a routine..
It's more than a spiritual habit.
that all Christians should have..
It's so much more than a ritual that we do..
It says we worship God who is great..
And in fact, the greatest..
For the Lord is the great God,.
great King above all gods..
He is the greatest..
He holds in his hands the depths of the earth,.
the highest of the mountains,.
the mightiest mountains,.
and the sea belongs to him,.
for he created it..
His hands formed the dry land too..
The depths of the earth,.
the lowest part of the world,.
which is the Dead Sea, by the way,.

$^{321}$in the land of Israel..
And then the highest of the height, right?.
The highest mountain..
And then the sea..
So he's trying to capture the entire cosmos,.
which is under the creative power of God..
The sea, in an Eastern culture,.
Oriental culture,.
is the untamable forces..
When you think of sea, the ocean,.
it's always unconquerable..
It's untamable by man, right?.
A couple days ago,.
a deep sea research company.
produced a 3D scan of the Titanic, right?.
The unsinkable vessel..
What year did it sink into the ocean?.
Atlantic?.
1912..
That's like more than 100 years ago, right?.
So this event, this tragedy,.
has fascinated many people for so many years..
And finally, we got some 3D scans..
I think the company captured about 700 images.
from a deep sea scan..
And it took about 200 hours.
for computer to compile them.
and create it into a 3D scan..
That is the Titanic..
1912, right?.
It's haunting images.
at the depth of the ocean..
With today's modern technologies,.
we still feel that the sea, the ocean,.
is untamable, right?.
It's just haunting images, right?.
Yeah, can you,.
just can't imagine if you suffer.
from this sort of tragedy, right?.
Crazy..

$^{361}$And this is a haunting image..
So much more realistic than this,.
Jack and Rose, yeah..
My point is, the creation is untamable, right?.
For even for modern culture,.
or for the culture of the Bible,.
the sea, the oceans, the highest of the mountain,.
the depths of the earth,.
we cannot conquer it,.
but God holds it in His power, right?.
God holds it in His power..
All the forces that threaten our well-being.
is all under God's power..
I attended a few funerals recently..
One is younger than me..
The other one is older than me..
Even today, with our modern technology, right?.
Healing takes time,.
and sometimes it doesn't happen, right?.
No matter how much money you have,.
how much technology you have,.
no, we feel, we still feel,.
or we have to face the final enemy,.
the forces of death, right?.
But the Psalm says, "God holds everything in His power,.
"and He has power over you, power over me.".
The Psalm talks about the power of God, the Creator..
And to many people, and to maybe your theology,.
your understanding of God,.
is that, okay, this powerful man is like some old guy,.
some old white guy out there, right?.
It's distant from us, it's like an absentee father..
He created the world, boom, okay, it's right there,.
and that machine will continue to work.
according to His physical principles and all that..
And then God, where is God?.
You know, God has nothing to do with us..
You know, kind of deistic, you know, deism..
It's not atheism, but deism, right?.
Some sort of deistic concept of God, and He's not here..

$^{401}$But the Psalm says otherwise..
The Psalm says God has everything to do with us..
Listen to what the Psalm says, verse six and seven..
"Come, let us worship and bow down..
"Let us kneel before the Lord, our Creator, our Creator,.
"for He is our God.".
We have a relationship with Him..
And we are the people He watches over..
He takes care of us..
We are like flock under His care, right?.
Notice how the verbs of worship going from action verbs,.
shouting, making praises, and dancing and all that,.
and finally come to stillness, bowing down,.
kneeling before Him, right, on our knees..
Number seven is a familiar picture, right?.
God is His shepherd, and we are His flock, His sheep..
Now, again, most of us, modern readers of the Bible,.
we don't know how to appreciate this kind of imagery..
We always think of sheep and lambs, right?.
What do you think of?.
Food on the table, right?.
Cook and done well, right?.
Right, in a pot, and then we can consume it..
You know, in the Oriental or Palestinian culture,.
or nomadic culture, you rarely sacrifice a lamb, right?.
Just like you don't sacrifice a dog, right?.
You don't..
You pet it, you take care of it, right?.
They don't, they don't sacrifice a lamb..
Very rarely that happens..
Maybe once a year, in Passover, right, with deep meaning..
To the first audience, to the Oriental people,.
shepherd and sheep, it's a tight relationship, right?.
Shepherd and the sheep, because the sheep listens.
to the voice of the shepherd only..
No one other than the shepherd..
There was a news story some years ago..
The Israeli army took over or captured some Palestinian lands.
and then they raided and they got all the properties.
they seized them, and one of those is a huge flock.

$^{441}$of sheep, hundreds of sheep..
So the Israeli soldiers put all the sheep into the ranch..
And how many of them?.
200 or 300 sheep in one ranch..
And then a certain Palestinian shepherd came.
to the Israeli army and asked the soldier.
who's guarding the door of the ranch,.
and the shepherd said, "Can I have my sheep back?".
And the Israeli soldier, "How many sheep do you have?".
And the shepherd said, "I have 20, 25, 25 sheep.
"in the ranch.".
And this Israeli soldier, a young soldier,.
laughed at the Palestinian old guy, the shepherd,.
and said, "Well, you can try..
"There are 300 sheep here..
"Go ahead and find your own sheep..
"Everyone looks the same, right?".
And the shepherd walks to the door.
at the entrance of the ranch,.
and he started to raise his voice..
It's like, "Come, come, come.".
Something like, "Yalla, yalla," right?.
You know what yalla is?.
Yalla is hurry up, Arabic, yeah..
Yalla, yalla..
And then, three, four minutes later,.
25 sheep came to the door,.
and the shepherd said, "These are my sheep.".
It's very interesting..
All the rest of the sheep,.
they're just doing their dumb things..
They don't listen to the voice of this shepherd,.
but this 20-some sheep followed the voice of the shepherd..
It's a true story..
(stammering).
The sheep will listen to the voice of the shepherd..
Remember Jesus said that, right, John 10?.
"I am the good shepherd,.
"and the sheep listens to my voice.".
And that's why this image of shepherd and sheep.

$^{481}$is a tight relationship..
It's a relational metaphor,.
and that's why at the end of verse seven,.
it goes into another theme,.
not unrelated, it's actually very related,.
talking about the voice of God..
If only you would listen to the voice of God,.
His voice today..
The Lord says, "Do not harden your hearts.
"as Israel did at Meribah,.
"as they did at Masar in the wilderness..
"For there your ancestors, your fathers,.
"tested me and tried my patience,.
"even though they saw everything I did..
"For 40 years, for one entire generation,.
"I was angry with them,.
"and I said, 'They are a people whose hearts are hardened,.
"'turns away from me..
"'They refuse to do what I tell them to do.'.
"So in my anger, I told and oath,.
"'They will never enter my place of rest.'".
Now in the beginning, I said Psalm 95 had two parts, right?.
It's actually a pretty strange psalm, right?.
If you write a worship song today,.
you probably would end your song,.
like the chorus or the ending, on a positive note, right?.
You kind of wrap around, going back to, you know,.
what you call it, tonic, right?.
Yeah, you always end the song nicely..
Psalm 95 doesn't do that..
It opens with a joyful sound, making noises..
"Come on, guys, make some noises, shout for joy.
"for our maker is here, we're gonna worship him.".
But then the second half, the ending,.
ends with a serious warning, a serious lesson, right?.
We often don't think a song is a good vehicle of education..
We don't do that because a song.
usually is three, five minutes long, it's very short..
But it can be a powerful vehicle.
for the message of God..

$^{521}$When we think of the Old Testament stories,.
or the entire Old Testament scripture,.
39 books, if I asked you what stories.
would tell you the tension between God and Israel,.
what are some of the stories?.
The stories in the wilderness,.
when God saves his people by grace,.
out of the land of slavery, out of the land of Egypt,.
and he walked with them and led them.
through the wilderness for so many years,.
but after they came out of Egypt,.
they complained and they grumbled against God,.
and they distrust what God is doing..
And God is patient, God is always long-suffering..
Those are the stories you can think of..
And in fact, that's what Psalm 95 does,.
to bring us back to the history of God's patience..
Now, so often we read through the Bible,.
we always get the impression,.
"Okay, God is so mean, right?.
"How can you have so mean of God, right?.
"And he doesn't like the people,.
"he yelled at the people," and all that, right?.
But if we read scripture like that,.
I think we are doing some selective reading,.
selective understanding..
If you read the Bible, at least the Old Testament,.
or maybe you add the New Testament,.
the entire Bible, 66 books,.
you realize that the story is not about God's wrath,.
it's about God's grace, how God is so patient.
with those who are rebellious..
We rebel against him all the time,.
but God is always patient, right?.
And Moses is always the mediator between the two, right?.
And God is always faithful..
And that's what Psalm 95 wants to get at..
It is one of the best psalms on worship..
The ending of Psalm 95 does not cater to our preference..
We always think, "Okay, we come to the face of God,.

$^{561}$"we come to worship God in the presence of God.".
When we worship, yeah, it has to be,.
but you come out of it, you come out of worship.
feeling peaceful, right?.
Feeling a little bit comfortable, right?.
That is our misunderstanding..
Sometimes the best worship is that we come before God.
and we are warned..
We receive a serious warning to our lives..
But we stop listening, then we'll go astray, right?.
And we'll end up in curses, in death,.
and separation from God..
Psalm 95, second half, the last half,.
is telling you, you must direct your attention.
to the voice of God, because God is a speaking God..
God still speaks, speaks to your heart through the spirit,.
and speaks to you in warning and in teaching..
If we refuse to accept or listen to the voice of God,.
then we suffer the consequence..
But God is always the gracious one..
He wants us to listen, right?.
Sometimes the best worship doesn't cater to our preferences..
It caters to God's preferences..
That your attention is always on God and on God alone,.
nothing else..
God is God speaking to you, and you are ready to listen..
Now C.S. Lewis says something really interesting.
and fun to say about all things, about all things..
I really started to enjoy his writings..
He compares true worship with a pair of good shoes..
A pair of good shoes..
Let me give you that quote..
He said, "A good pair of shoes.
"are the ones you don't notice..
"The perfect church worship service.
"would be one we were almost unaware of..
"Our attention would have been on God alone.".
That makes sense..
You're wearing a good pair of shoes now..
Do you pay attention to them?.

$^{601}$No..
A good shoe will not get your attention, right?.
But it will carry you through, right, while you're walking..
But you don't pay attention to them.
because you're unaware of the existence of shoes..
A true worship is kind of like that..
Unaware of everything else, right?.
But we are fully aware of the presence of God..
So I invite you today to polish your craft of worship..
Psalm 95 is not written to clergy..
It's not written to the worship band..
It's not written to those who are in charge of the worship,.
the worship personnel, the religious workers.
who prepare the worship service..
It's written to all the people of God..
It's written to you..
How do you craft your art of worship.
so that your worship is pleasing to God?.
How do you prepare your heart, right?.
You know, the most terrible sins.
that the Israelites committed in the wilderness.
were not breaking any laws..
They were not doing horrible crimes..
Do you get that?.
They never do really horrible things..
They were not criminals..
The most horrible thing, the most terrible thing.
that the Israelites, the people of God did.
was being hard-hearted, not willing to listen..
They're not willing to open their hearts..
So how do you prepare your hearts?.
How do you prepare your art of worship?.
It's to pray..
Sometimes a simple prayer,.
to pray that the Spirit will speak to you, right?.
Even for myself, sometimes I attend worship service..
Someone else is speaking, right?.
I pray that God will speak to me, right?.
That I will have a softened heart to come before Him,.
to acknowledge, to be totally aware of His presence..

$^{641}$Let's all pray together..
And I ask the worship band to come back up here.
and sing our last song..
And let the song be our worshipful response to God..
And most importantly, as Psalm 95 says,.
come before God with listening ears..
Let our heart be open to Him..
Let's all pray..
Father, we give You thanks for the fact.
that You are worthy of our worship..
Yes, You are worthy, because You are our Creator,.
our Redeemer..
We praise You, we make song, we shout..
For You, and joyfully we do so,.
because You are worthy of our praises..
And yes, we ask that Your Spirit will open our hearts,.
humble us, so that we'll listen to You..
Thank You for being our God, for Your love, for Your grace..
In Christ, we pray..
(upbeat music).
\newpage



\section{詩篇}
\label{sec:wLYSrSGINOw}
\textbf{敬拜神的子民:詩篇95 - 陳崇基牧師}
\newline
\newline
連結: \href{https://youtube.com/watch?v=wLYSrSGINOw}{\texttt{https://youtube.com/watch?v=wLYSrSGINOw}} ~~~~ 語音日期: 2023-05-29
\newline
\newline
\hyperref[sec:jG3lPKnsQeA]{\small{< < < PREV SERMON < < <}}
~
\hyperref[sec:index_chronic]{\small{[返順時目]}}
~
\hyperref[sec:index_scriptual]{\small{[返順卷目]}}
~
\hyperref[sec:4x5CA4Q65_w]{\small{> > > NEXT SERMON > > >}}
\newline
\newline
$^{1}$我地交比佢.
大英姐妹好.
我係準備好啦.
arm arm 有個隊友就話.
趕唔切你呀.
因為唔記得帶credit card.
所以返jor 去lor credit card.
好緊要呀lor credit card.
最緊要有護照同埋credit card.
仲有credit card credit limit要幾萬蚊.
咁就ok啦.
一定會返到黎ge .
首先多謝大英姐妹.
同埋特別教育同工.
有個好ge 配搭.
以致我就能夠帶大英姐妹去Israel.
咁我地帶佢幾多個UV呀.
十個係咪呀.
十七個係世界各地唔同地方去.
咁已經有一D 團友已經到jor 啦.
安全到jor 所以佢地報jor 平安.
咁我想多謝教會同埋同工ge 配搭.
就係我地能夠去祝福到其他人.
大家知道我唔係做呢行架.
我唔收小費ge .
我唔收小費我去到目的.
希望大家充充份份有一個.
豐盛ge 一個旅程.
That's my purpose.
我ge job就係做呢樣ye .
咁我唔介意去流多少少汗ge .
咁但係我就唔係做呢行食飯.
我唔係靠呢行食飯ge .
咁但係可以有一D 資源同人地去分享.
呢個係一件樂事.
我地教會能夠祝福更多ge 人.
咁今日我地睇一個主題.
咁就係敬拜.
上個禮拜我地睇多jor 一D 敬拜ge 經文.
同敬拜有關ge 題目.

$^{41}$記得Brenda黎jor .
咁我地就開始黎到講一講敬拜.
乜ye 叫做worship.
你入到黎教會.
你呢個主日崇拜.
到底一回咩事.
我地值得去思考一下.
呢個係我地基督徒應有ge 習慣.
一個正常基督徒應有ge 習慣.
每個主日我地去崇拜.
咁到底崇拜係一回咩事呢.
詩篇第95篇.
我覺得呢篇詩篇係眾多ge 敬拜ge 詩歌.
你知詩篇裡面有150首ge 歌.
呢D 都係以色列人舊約ge 時候.
神ge 子民用黎去敬拜.
其中一首講到敬拜呢.
就係95篇.
我地先讀.
來呀我地要向耶和華歌唱.
向拯救我地ge 盆石歡呼.
用詩歌向佢歡呼.
因耶和華是偉大ge 上帝.
是超越萬神ge 大君王.
地的深處在佢手中.
山的高峰也屬佢.
海洋屬佢.
是佢造的.
罕地也是佢手造成的.
來呀我地要.
俯伏敬拜.
在做我地ge 耶和華面前.
跪拜.
因為佢係我地ge 上帝.
我地係佢草場ge 百姓.
係佢手中的羊.
惟願你地今天聽佢ge 話.
你地不可硬著心.
像在米利巴.
就在曠野馬殺的日子.

$^{81}$那時你地的祖宗.
視我探我.
四十年之久.
我厭煩那世世代說.
這時心裡模糊的百姓.
竟不知道我的道路.
所以我在路中起誓.
他們斷不可進入我的安城.
呢篇詩篇有D 古怪.
如果你揀一首歌.
去到敬拜去唱.
好似頭先幾首歌去敬拜.
去到唱出讚美詩詞.
通常都係end with.
結束的時候都係比較正面少少.
呢首詩歌就去到呢度.
就冇了.
你唔好返去搵12節13節.
就呢度完.
有D 古怪.
我地就睇下呢篇詩篇.
就俾我地聽敬拜係一回咩事.
今日黎到教會坐係度.
我地大家係度敬拜.
係到底一回咩事.
第一個問題想問大家ge 今日黎呢度ge 時候.
你做jor 咩準備功夫呀.
有咩準備功夫你做jor 呀.
食早餐係咪.
飲咖啡呀等等.
到底有咩準備功夫我地要做.
我覺得我地一個現代ge 世代.
教會 基督徒.
我覺得呢個敬拜呢回事.
好多時係一個lost art.
一個已經失去jor ge 藝術.
我地彷似呢黎到主日崇拜ge 時候.
我地好似係冇咩準備.
你都可以講有D ye 準備.
但通常我地講緊準備ge ye 都係.

$^{121}$人地做俾我地ge .
係咪.
有冇早餐呀 有冇咖啡呀.
邊個煲呀.
gwo D 早到gwo D 人煲係咪.
我黎到我expect有D 咩.
就有杯熱咖啡.
冇理由呢度凍冰冰ge 隔夜咖啡係咪.
下次你唔會黎啦係咪.
仲有D 咩要準備.
gwo D 敬拜隊呀.
音樂呀.
上黎做ye gwo D 人呀.
或者後面後台要做ye gwo D 人呀.
全部所有ge preparation都係佢地做ge .
powerpoint出得arm 唔arm 呀.
等等呀 有冇打錯呀.
有D 乜ye 我地需要準備.
還是我地黎到敬拜ge 時候.
我地只會覺得人地需要準備.
我地唔需要準備.
你知唔知古代gwo D 建築呀.
去一D 聖堂或者去一D 聖殿呀.
或者教堂呀.
gwo D ge 樓級呀.
係好多時一長一短一長一短.
我地叫佢做processional steps.
點解.
係咪特登整到你撲親咁呀.
係咪裝你呀.
唔係呀就係要人.
去聖殿或者去敬拜ge 時候.
佢要小心每一步.
係減慢佢速度.
係一個preparation黎架.
佢slow down.
你去到上帝面前ge 時候.
你去敬拜ge 時候.
唔係一D 等閒事.
因為you are approaching a sacred space.

$^{161}$arm 呀.
呢D 可能都係lost art.
我地去以色列之前就有D ye 準備.
裝束問題.
因為去某D 地方.
係唔能夠露咁多肉架.
各位女士.
現代教會ok.
我地都算係唔係ok.
當然唔好太少.
即係大家都好規矩.
即係費事有D 分心ge 情況出現.
咁所以係囉.
係咪.
姐妹講架.
唔係男人講.
我地就講裝束.
因為去到某D 地方你係唔能夠著短褲.
你唔能夠著背心.
男人都唔得架.
去到某D 地方.
你就要帶埋頂帽.
但係去某D 地方就要除頂帽.
好奇怪.
上幾個月去一個基督教教堂.
因為我要拍ye .
好辛苦我做ye .
仲要做埋義工.
自己拍埋ye 自己錄埋ye .
去到gwo 個教堂.
跟住gwo 個看更.
gwo 個修士.
一ye 打到我頭gwo 度.
除帽.
基督教教堂就要除帽.
但係去猶太人ge 會堂就要帶帽.
咁咩極煩呀.
arm arm 先去完猶太人ge 會堂.
咁就帶住頂帽.
或者gwo D keeper.

$^{201}$跟住去gwo D 基督教ge 古教堂.
跟住你又要除低頂帽.
係有D ye 要prepare.
有冇諗過聖經裡面都有講.
你敬拜ge 時候.
係需要有D preparation要做.
你記得耶穌.
登山補佛.
佢話你地去獻祭ge 時候.
當時仲有聖殿去獻祭.
獻祭之前你要做咩.
要同gwo 個同你有拗叫ge 論者.
盡可能你都去和好.
你去和好.
你唔可能你都去和好.
係咪.
大家和解jor 你先至好去獻祭.
即係搞掂jor .
你ge 關係先.
你先至可以approach上帝.
當然我地知道世界唔完美.
有D 人.
有D 情況有D 處境唔能夠.
但係個精神你知道嗎.
個精神gwo 個spirit.
就係我地需要有preparation.
relationalge preparation.
physicalge preparation.
等等.
盡可能ge 我地同人地和解.
盡可能起碼我釋懷jor .
我地到上帝面前.
呢D 都係我地需要有ge 準備.
咩叫做敬拜.
詩篇95篇.
分開兩半.
頭一半頭七節.
跟住gwo 半就係七節到.
七節下半節.
至到十一節.

$^{241}$係一D 好典型ge .
敬拜ge 詩歌.
點解話典型.
因為你睇詩歌ge 時候.
通常呢D 敬拜詩歌.
佢都會有一個所謂叫做宣召.
一個call to worship.
一個叫醒你.
一個提醒你.
係時候啦我地要去敬拜.
你睇下呢度第一節第二節第三節.
佢話.
來呀我地要向耶和華.
歌唱.
要向拯救我地ge 盆石.
歡呼.
來到他面前.
要以詩歌向他.
歡呼.
你諗唔諗到呢段頭一半.
有D 咩動詞.
唱歌歡呼.
讚頌.
singing shouting.
呢D 全部都係要energyge ye 黎.
換句話講.
黎到神面前唔係一個被動ge .
情況.
一個主動.
ge 參與係咪.
call to worship.
好多時都話比我地聽.
the reasons.
乜ye 原因要去敬拜.
點解要去敬拜.
搞清楚乜ye 叫做.
敬拜搞清楚乜ye 叫做.
原因點解要敬拜.
三四五節.
因為.

$^{281}$你見唔見到gwo 個字因為.
因為.
因為耶和華.
係最大ge 神.
耶和華係唯一ge 神.
地ge 深處.
係佢手中 山ge 高峰.
也屬佢 你見到寫得好靚.
地ge 深處.
山ge 最高峰.
由最低去到最高.
然之後.
海洋 旱地.
中間.
都係佢手所做.
整個世界 整個宇宙.
都係佢所做.
你睇頭一半.
你諗起咩經文.
創世記第一章.
創世記要講咩神學.
就係上帝創立.
所有ye .
你見到呢個世間 呢個宇宙.
仍然係有規律ge .
仍然係有規則ge .
有設計ge  邊個去創立.
邊個去design 係上帝自己.
特別講到海.
頭先講到地ge 深處.
山ge 高峰.
海.
海對當時.
古代ge 人黎講.
或者甚至今日.
都係一D 唔能夠戰勝ge ye .
唔能夠controlge ye .
untainable.
冇得去conquer佢ge .
因為海就係咁.

$^{321}$海洋 正如我地今日所講.
稀山磨稀水.
前兩日我睇一個新聞.
佢講到鐵達利號.
鐵達利號 幾時沉.
記得 鐵達利號.
1912年ge 時候.
鐵達利號.
我地今日現代ge 人.
諗返D 鐵達利號.
呢個所謂unsinkable boat.
unsinkable ship.
我地都好fascinated.
一艘好大好大ge .
郵輪.
係海洋gwo 處去沉jor .
然之後拍埋電影出黎.
等等.
有好多好浪漫ge 故事.
前兩日睇一個新聞.
講到有一個.
深海ge 研究.
ge 公司.
深海研究公司.
佢地就去到終於做jor 第一個.
鐵達利號沉船.
海底ge 掃描.
3D scan.
咁就係用jor .
70萬張ge 圖畫.
scan出黎ge 圖畫.
然之後再組合出黎.
用電腦用jor 200個小時.
先能夠砌返呢一個.
ge 鐵達利號ge .
即係殘骸.
你諗下幾咁.
難去到做ge 一件事.
1912年.
沉過呢個船.

$^{361}$成100年有多.
我地而家先至可以點呀.
用科技將佢.
唔係打撈出黎添.
將佢scanjor 佢.
大家知唔知呀.
你見到呢D 咁ge image.
係好huntingge image.
即係好.
深寒ge .
一D ge 圖畫.
同呢個好唔同呀.
Jack and Rose.
到今日我地都覺得海洋.
係唔能夠.
tameable.
欺山磨欺水.
但係呢D ye 係上帝.
ge control.
你諗下我地心靈裡面有D 乜ye .
係untameable.
最近我去jor 幾個喪禮.
有後生ge .
有年長ge .
你話咪有D untameable.
ge 威脅.
一D ge 勢力.
係影響住你同我.
但係上帝.
係最大ge 上帝.
佢有權能.
包括控制埋呢D ge .
勢力.
呢度就講到.
點解要去敬拜.
因為講到上帝ge 權能.
佢係創造主.
當然我地好多時.
諗到創造主好大力.
好多能力呀等等.

$^{401}$我地總會覺得呢個神.
佢係一個上帝.
係一D 同我地距離好遠ge ye .
可能好耐之前.
佢曾經創立jor 呢個世界.
跟住佢就係absentee father.
即係佢就唔再.
同我地一齊.
唔再係度同我地冇乜關係.
可能我地都會咁諗.
或者呢個culture要我地咁諗.
但係呢篇詩篇講咩.
佢話唔係喎.
你唔好忘記.
第六節講咩.
佢話來呀.
我地要俯伏敬拜.
做我地ge 夜幻面前.
跪拜.
因為又因為.
因為他是我地的上帝.
我地是他.
草場上面ge 百姓.
是他手中的羊.
你留意呢度.
佢越黎越.
第一節ge 時候佢咪話.
你呀我地要去敬拜.
要去歡呼要去到歌場.
然後慢慢慢慢好似呢個動詞.
越黎越靜落黎.
靜到乜ye 地步呢靜到跪拜.
跪係船.
You bow down.
靜止.
認識到上帝ge 偉大ge 時候.
我地就係佢面前去靜止.
第七節.
佢話神係我地ge 牧者.
我地.

$^{441}$係佢ge 羊.
我地係佢ge 子民.
我地又係現代人.
我地睇唔明呢D ye .
我地諗起羊ge 時候.
諗起支竹羊腩煲.
即係死jor ge ye .
我地好少見到生ge 羊.
生ge 狗呀貓.
我地成日見.
但生ge 羊我地好少見.
以色列zar 車經過.
羊呀羊呀好似全世未見過羊.
因為我地太少見.
我都諗起羊同埋牧者.
係乜ye .
係一種好親密ge 關係.
好似你養ge 狗同埋貓.
佢識認你.
佢就唔認其他人.
十幾二十年前.
有一單新聞.
係巴勒斯坦地發生.
就係以色列軍隊.
就走去一個巴勒斯坦ge 村落.
咁就去到.
即係攻佔jor 佢.
lor 晒裡面ge ye .
其中就lor jor 幾百隻羊.
咁就gwo D 軍隊就趕晒D 羊.
走jor 去一個羊欄.
係個ranch 裡面.
咁有一個巴勒斯坦ge 牧羊.
牧羊人.
佢返到黎唔見晒D 羊.
佢有二十幾隻羊.
唔見jor .
咁佢就走去gwo 個羊欄gwo 度同gwo 個.
以色列ge 士兵講.
士兵就守住gwo 個羊欄.

$^{481}$gwo 度有幾百隻羊.
gwo 個牧羊人就同佢講.
可唔可以lor 返我D 羊.
咁gwo 個年青ge .
以色列士兵就同呢個老人家講.
佢話.
你有本事你就lor 返你gwo 二十幾隻.
呢度成堆羊 隻隻都係羊.
咁gwo 個牧羊人就話.
你等我試下.
佢就企係羊欄ge 門口.
佢就係度講一D 好簡單ge 說話.
come come come 咁晒下.
然之後幾分鐘之後.
屬於佢gwo 二十幾隻羊.
就走到佢面前.
其他羊就唔理佢.
好奇怪.
耶穌講咩.
羊係聽牧羊人ge 聲音.
佢真係認到.
個主人牧羊人.
gwo 把聲.
我地是他.
草場上ge 羊.
係他手所做ge .
草場上面ge 百姓.
係佢手中ge 羊.
係要聽佢ge 聲音.
聽到上帝ge 聲音.
並且肯聽上帝ge 聲音.
所以點解呢.
第七節開始佢就講聲音.
邊個ge 聲音.
你睇下.
唯願我地今天.
聽他的話.
你唔不可硬著心.
像在米利巴.
就是在曠野馬殺的日子.

$^{521}$那時你們祖宗視我探我.
四十年之久.
我厭煩那世代.
這時心裡模糊的百姓.
竟不知道我的道路.
所以我在路中起誓.
他們斷不可進入.
我的安城.
我頭先開頭ge 時候講.
話呢篇詩篇係一篇好奇怪ge 詩篇.
係咪.
就係咁樣完.
就正係咁樣完.
開始就係歡呼.
我好似開過Party咁.
你閉嘴.
你冷靜下.
你聽.
聽.
係一個警告.
好嚴肅ge 教訓.
好多時我地睇詩歌.
我地都唔覺得詩歌係.
所謂的教育機構.
好多時我地覺得詩歌唔係一D .
教導我地ge 媒體.
因為詩歌有幾長.
一首詩歌大概.
都係三至四分鐘.
五分鐘.
一首比較長ge 詩歌.
我地好多時覺得詩歌咁短.
係冇咩教導ge 功能.
諗起詩歌.
我地諗起都係抒情ge ye .
我唔記得我年輕ge 時候.
都老.
有D 後生ge 都唔知道有D 咁ye .
以前年輕ge 時候.
我地zar 車係用cassette帶.

$^{561}$有用過.
我地年輕ge 時候.
我地鍾意gwo 首歌.
我地就dubjor 佢.
成個鐘都係唱gwo 首歌.
因為我地唔想翻帶.
翻帶好煩.
鍾意gwo 首歌.
點解呢?因為想聽返.
諗起詩歌.
當然唔會諗起D 好複雜ge .
lecture或者課程.
但詩歌係一個好重要ge .
教導工具.
所以點解有150首詩歌.
並且裡面有好多詩歌.
都係要去到.
做一D teaching.
呢一首就係.
講到話神ge 子民.
同上帝ge 關係.
就講到咩馬殺的日子.
米利巴.
點都好佢講咩ye .
講緊以色列人出jor 埃及之後.
上帝帶領住佢地.
藉住摩西帶領住佢地.
佢地就去.
厭煩上帝.
點解唔比水我.
點解唔比ye 食比我.
跟住入到去.
加cum 地.
去窺探gwo 個地ge 時候.
又返返黎又話唔可以入去.
所以上帝就係.
曠野ge 時候佢就話.
好 你地唔聽我講.
你唔想入我ge 安息.
佢就話好 你就係曠野.

$^{601}$40年.
第一代出埃及.
就係曠野.
就流離jor 40年.
點解會睇呢樣ye .
因為呢個係.
一個最典型ge .
係舊約整本聖經ge 故事裡面.
講到叛逆.
講到硬著心腸ge .
一個故事.
你諗起咩叫做硬著心腸.
猶太人.
猶太人係硬著心腸.
硬頸ge .
你就諗起呢D 咁ge 故事.
馬殺的日子 即係試探ge 日子.
等等.
所以呢首詩歌完結ge 時候.
佢話乜ye  上帝話你地.
唔能夠進入.
我應許之地.
因為你唔聽我.
說話.
你可能會.
睇完呢篇詩篇.
你會有個結論.
呢個上帝真係好鬼死惡.
惡死到不得了.
如果你睇完聖經.
或者睇完舊約.
你會覺得呢個上帝.
淨係無故去發脾氣.
無故去發嬲ge  去激打人ge .
咁可能我地睇錯.
因為你如果睇聖經.
你淨係睇到gwo D ye .
會唔會有D .
選擇性閱讀.
其實睇返整本ge 舊約.

$^{641}$整本新舊約都好.
你睇到一個好patientge 上帝.
一個好有忍耐.
好有寬容 其實係一個好有愛ge 上帝.
詩篇95篇都係一樣.
係 縱使我地.
黎到敬拜ge 時候.
我地需要一個warning.
一個嚴肅ge .
責備.
我覺得咁樣先係一個好ge 敬拜.
我見到上帝 遇見上帝.
然後比佢可能當頭棒喝也好.
或者好嚴肅ge D 提醒.
咁樣先至係一種真正同上帝ge 遇上.
真正同上帝遇上未必一定係好溫馨.
好舒服.
好多時係當頭棒喝.
E3篇詩篇就係咁樣完結.
我地ge Attention.
我地ge 焦點.
唔好放係任何ye 身上.
放係上帝自己身上.
一個英國ge 文學家.
我覺得佢講ye 好有趣.
有幾個著作.
佢講到敬拜.
一個真正ge 敬拜.
一個同上帝.
創造天地ge 主.
愛你ge 上帝係你ge 牧羊人ge 上帝.
去到相遇一個敬拜.
真正ge 敬拜.
True worship 係點樣.
True worship 好似一對好ge 鞋.
一對好ge 鞋.
咩意思.
佢話一對好ge 鞋.
係你唔會留意到.
你唔會留意你著緊一對好ge 鞋.

$^{681}$唔會係你ge 對鞋上面.
True worship 係咁.
好似一對好ge 鞋.
正如好似一個好ge 閱讀一樣.
你唔會留意眼鏡arm 唔arm .
有冇老花等等.
夠唔夠光.
你係進入閱讀裡面.
就好似你著住對鞋.
進入一個步行ge 功能.
佢係一個perfect church worship.
完美ge 敬拜.
係一個敬拜過jor 之後.
unaware.
唔會驚異到.
唔會留意到.
不知不覺地就到過.
因為.
你ge 焦點.
係係上帝身上.
你唔會留意到底首歌唱得好唔好.
你唔會留意gwo D 傢俬arm 唔arm .
燈光好唔好.
因為你全神貫注ge .
就係上帝.
好似一對好ge 鞋.
結束ge 時候.
比我地一D 提醒好嗎.
咩叫敬拜.
諗一諗點樣預備.
你ge 敬拜.
去同上帝.
面對面D 度相聯.
詩篇九十五篇.
寫比每一個.
人ge 子民.
唔係gwo D 美人呀.
或者gwo D 聖職ge 人員.
寫比每一個神ge 子民.
點樣讓我地ge 敬拜.

$^{721}$能夠更加不知不覺地.
係到過.
因為你ge 專注.
全部放係上帝.
自己身上.
預備下我地ge 心靈.
以色列人同上帝.
係一代裡面ge 關係.
四十年.
唔係話因為佢地大奸大惡.
你留意到佢地.
冇咩大奸大惡.
冇犯過咩crime.
但佢地.
最horriblege ye 係乜ye .
唔係做jor D .
咩刑事ge 事.
而係佢地心硬.
你知心硬就無得救.
唔肯去聽.
唔肯去到return.
最好ge 預備.
黎到敬拜 黎到神面前.
甚至你每一日.
你返開聖經ge 時候.
同上帝建立關係ge 時候.
最好預備就係.
你先黎到上帝面前 我謙卑度過.
你知我係會心硬.
人地提醒我.
專業人士提醒我 朋友去提醒我.
我都唔會聽.
gwo 個程度就好煩.
就係上帝話厭煩gwo 個世代.
就係佢偏唔聽.
靠聖靈你都可以.
軟化你ge 心.
今日你聽見上帝ge 聲音.
就不要硬著心.
聽清楚神對你講ge ye .

$^{761}$知道上帝ge 道理.
我同心低頭祈禱.
然後會回應思過.
上主我地求你今日對我地講說話.
求你今天讓我地.
係聖靈ge 裡面.
同你相遇ge 時候.
你俾我地一個渴慕你.
肯去聽你ge 心像.
俾我地呢個心能夠聽見你自己的信念.
俾我地一個ge 提醒.
讓我地今日唔係硬著心黎到你面前.
讓你係我心中你去工作.
帶領我地同你面對面你去相遇.
Yes 你就係我地ge 上帝.
唯一ge 上帝.
係我地能夠全神貫注.
在你自己身上.
係蒼天造地.
蒼天造地創造我地.
帶領我地ge 牧者.
我地向你歡呼.
向你敬拜.
祈禱奉全命感.
(音樂結束).
\newpage



\section{}
\label{sec:4x5CA4Q65_w}
\textbf{比對研究 15: 巴力宗教 8 - 陳崇基牧師}
\newline
\newline
連結: \href{https://youtube.com/watch?v=4x5CA4Q65_w}{\texttt{https://youtube.com/watch?v=4x5CA4Q65\_w}} ~~~~ 語音日期: 2023-06-03
\newline
\newline
\hyperref[sec:wLYSrSGINOw]{\small{< < < PREV SERMON < < <}}
~
\hyperref[sec:index_chronic]{\small{[返順時目]}}
~
\hyperref[sec:index_scriptual]{\small{[返順卷目]}}
~
\hyperref[sec:kj0zBwcuLqE]{\small{> > > NEXT SERMON > > >}}
\newline
\newline
$^{1}$除了死神的形容之外,還提到一個很重要的戰爭.
那個最重要的戰爭就是死亡之神被粉碎.
我們講兩段,其中一段巴力的文字叫做epic of Cotter.
上文下理講到國王有病,整個地上有瘟疫,有疾病,有死亡,有旱災出現等等.
這些有旱災,有瘟疫,有疾病,有死亡,這些在神話中或中東都會放在一起討論.
因為很簡單,有旱季,沒有雨水,沒有生產就自然會有飢餓.
會有瘟疫出現,會有病出現,最後會有死亡.
所以這些主題經常會放在一起討論.
在上文中,講到地上沒有食物,因為沒有雨水,所以會有飢荒出現.
因為沒有食物的緣故,在烏加利眾神會中,最大的神叫做EL.
他會造一個女神出來,叫做Shataka.
造了這個女神出來,讓這個女神把疾病和死亡拿走.
以至烏加利的王能夠健康一些,能夠有康泰.
以致王平安之後,又有雨水和食物出來.
這一段講什麼呢?.
Shataka要做什麼呢?.
他要將死亡粉碎,Shataka可以將死神打敗.
Shataka離開了神的重天庭,進入了王的房子.
進入了王的房子,然後進入了王宮,代表他粉碎了死亡,然後給了王健康.
當時的中東的王,就是神所設立的.
保守著國家平安,有雨水.
他可以代表眾民向巴力禱告,有雨水,有出產,以致大家有豐收等等.
所以王很重要,王健康也很重要.
所以EL就做了Shataka,來戰勝死亡.
之前我們提過巴力的循環,神話.
說在漢季完結的時候,巴力被摩特吞併.
然後他的老婆Anart出來復仇,將摩特殺死.
這一段最清楚,這裡提到Anart,這個女神.
抓住死神摩特,用劍劈開他,然後用西子去篩選他.
就像用菊來打散他.
所以死神雖然有這麼大權,但在循環中,他永遠不會有絕對的權力.
因為最終Anart將他打敗,然後Bale巴力在陰間再被復生.
我們說起Maud,有什麼名號,或者他出名是什麼呢?.
我們再看一段,然後才進入以利亞的故事.
我們在看的時候,以利亞的故事是這樣寫出來.
交代一些這樣的情節,為什麼呢?.
就是我們比對研究,我們拿經外的文字.
來做一些對比的時候,給一個背景的時候.
我們就更能夠清楚為什麼這些故事會這樣發生.
除了我們說如此真的發生過之外.

$^{41}$到底為什麼聖經作者會這樣堆砌以利亞的故事.
其實想做什麼呢?.
我們繼續看下去,這個Maud的名號是什麼呢?.
其中他有兩個相同的名號,就是死亡和解散.
Death and dissolution.
死亡和人死,就是說家庭有解散等等.
所以他的名號就是這樣.
我這裡就引用一本大典.
美國這本大典就是很多近東的學者.
寫有關近東的宗教,特別是一些神和魔鬼.
Deities and Demons.
這本書就成為了一本大典,一本百科全書或者字典.
這本是頗新的書,有一段文章.
就是講到Maud這個神.
除了講到他的名號是Death and dissolution之外.
就是說Maud是拿著喪親的權杖.
也是拿著寡婦的權杖.
這裡就很奇怪,為什麼會講這個呢?.
喪親就很明顯了.
因為他吞噬了巴黎,沒有雨水.
然後就沒有土產,然後有饑荒.
有饑荒出現就有病.
人就會軟弱,然後就會死亡.
所以Maud也有這樣的權柄.
就是他拿著親人的生命.
所以他有一個所謂的Bereavement.
就是喪親的時候哭泣的權柄.
還有他有什麼呢?就是寡婦的權杖.
為什麼會有寡婦呢?.
因為男人死了,女人就變成寡婦.
所以他們不單是喪親,他們也自己變成寡婦.
所以寡婦的出現,或者有寡婦的出現.
為什麼呢?在巴黎的神話裡面.
就是因為Maud是掌管著人的生命的.
吞噬人,被人死亡的.
所以就會有寡婦出現.
所以寡婦,我們一想到寡婦的時候.
不單是想到近代社會裡面的寡婦.
是社會裡面最低下階層的,或者無依無靠的人.
其實跟這些神話很有關係.

$^{81}$就是這個Maud就是造成寡婦的.
是不是?誰得罪他.
這個Maud就會將他變成寡婦.
所以他擁有著很大的生死之權.
這個就是Maud.
所以在巴黎神話裡面.
除了交代巴黎戰勝Yam之外.
即是雲頓大水之外.
也有很多章節特別來交代巴黎和Maud之間的爭戰.
因為巴黎要去到掌管雨水.
掌管著生命的時候.
他就一定要打敗Maud.
Maud未能夠打敗,生命始終是很脆弱的.
所以一年又一年又一年的循環裡面.
巴黎就不斷要去到對付死亡.
除了大水之外就是死亡.
我們剛才說一直看巴黎的文字的時候.
就想一想以利亞的故事.
以利亞的故事.
為什麼會交代這樣的主題.
我們從17章第一節開始.
以利亞的故事是先知被公認.
上帝說你去到一個溪旁.
然後烏鴉將肉帶給你.
然後你可以喝溪裡面的水.
縱使在旱災裡面.
沒有雨了.
耶和華說我會供養你.
所以耶和華才有給水和食物的權力.
他才是掌握生命的唯一的神.
不是巴黎.
第二個故事是說什麼.
第二個故事就是說.
以利亞去到西頓.
撒亞法這個地方.
就去找一個寡婦.
然後因為當時的旱災太嚴重.
剩下的油和麵粉是很少的.
以利亞的故事.
17章列王記上17章中間的故事.

$^{121}$就說到上帝去到公認一切的需要.
連寡婦和她的兒子都能夠養著他們.
先知也能夠養著他們.
只要憑信心相信.
這個真神耶和華是掌管一切的.
到第三個故事.
17節.
他說這事之後.
第三個故事.
那婦人就是那個寡婦.
就是那家裡面的女主人.
她的兒子就生病了.
病得很重.
並且沒有氣息.
死了.
婦人就對以利亞說.
那個神人我跟你有什麼關係呢.
你竟然到我這裡來.
使上帝來記起我的罪.
以致我的兒子死了.
這些都是很近東的概念.
因為有些咒詛臨到.
死亡臨到.
就代表我一定是和那個神明之間有些不好的關係.
得罪了他.
想起我的罪.
所以我的兒子現在要得到下場.
有咒詛來到.
以利亞就跟他說.
將你的兒子交給我.
以利亞就從婦人懷中.
接過孩子來.
抱他上到她所住的頂樓.
放在自己的床上.
以利亞就將這個兒子接過來.
一個死了的人.
他就叩高耶和說.
耶和說我的上帝.
我寄居在這個寡婦的家裡面.
你卻降災給她.

$^{161}$使得她死了.
於是三次伏在孩子的身上.
叩高耶和說.
耶和說我的神.
求你使這孩子的生命歸回給她.
耶和說聽見.
以利亞先知的呼求.
孩子的生命就歸回她了.
她就活著.
死然後重生.
死而復活.
以利亞就將這個孩子.
從樓上抱下來.
進了房間.
就給他的母親.
他說看你的兒子活了.
婦人對以利亞說.
現在我知道你是神.
耶和說.
你口中所說的話是真的.
為什麼會記載這個故事.
除了是自然地發生.
沒有幾年的下雨.
然後沒有出產.
就會有饑荒.
然後死了.
其實記載這個故事.
想做些什麼.
由第十章第一節開始.
以利亞就出來.
和以色列國.
特別是阿哈旺.
宣告一件事.
就是耶和說才是真神.
不是巴力.
阿哈是引進巴力.
將菲利基的宗教.
他的假神.
引進入薩馬利亞和以色列國.
以利亞在做什麼呢.

$^{201}$其實是在做一些所謂的護教學.
一些所謂的文化護教學.
Cultural Apologetics.
就是說這個文化告訴我們.
巴力就是掌管所有的.
雨水 食物.
還有什麼.
生命.
他打敗了摩特.
以至他能夠降雨.
但是這裡說的是什麼.
這個故事一個又一個的故事.
就是將一些巴力神話裡面的主題.
針對那些主題.
然後有些神蹟發生.
這些神蹟就是指向.
耶和說才是賜雨水.
賜生命過的神.
\newpage



\section{17:17-24}
\label{sec:kj0zBwcuLqE}
\textbf{Death in The Family: 1 King 17\_17-24 - Tommy Lee}
\newline
\newline
連結: \href{https://youtube.com/watch?v=kj0zBwcuLqE}{\texttt{https://youtube.com/watch?v=kj0zBwcuLqE}} ~~~~ 語音日期: 2023-06-05
\newline
\newline
\hyperref[sec:4x5CA4Q65_w]{\small{< < < PREV SERMON < < <}}
~
\hyperref[sec:index_chronic]{\small{[返順時目]}}
~
\hyperref[sec:index_scriptual]{\small{[返順卷目]}}
~
\hyperref[sec:7Sx7wB7NbaI]{\small{> > > NEXT SERMON > > >}}
\newline
\newline
$^{1}$All right..
Here we go..
Well, good morning, everyone..
For those, I recognize a few faces here this morning..
And for those who are here from college and all that, my name is Tommy Lee..
And just a little bit about me, I'm born and raised here in Chicago..
I'm 46 years old..
My parents immigrated here to Chicago in 1972..
So I lived in the Chinatown Bridgeport area all my life..
I still remember when Bridgeport was, I was one of only five Chinese students who had.
ever been to a college..
And I was a freshman in high school..
I still remember when she first started..
I went to Washington University in St. Louis..
I thought I was going to be a doctor..
Because a lot of times, when you have Asian parents,.
my parents gave me a choice..
Doctor, lawyer, engineer..
Pick one..
WashU gave me a full ride,.
so I guess I'm going to be a doctor..
After a first semester, I realized,.
yeah, I'm not going to be a doctor..
I was going to be a lawyer..
Many of you guys are familiar with Hebrews 11, the Hall of Faith..
It talks about all of the Old Testament heroes that the writers of Hebrews tells us belong.
into this Hall of Faith..
It gives us the names in order belonging to the Hall of Faith..
There's Abel..
There's Enoch..
Noah was in the Hall of Faith..
Abraham was in the Hall of Faith..
Isaac was in the Hall of Faith..
It talks about Jacob..
It went down to Joseph and Moses..
It even goes to the people who crossed the Red Sea..
Rahab..
And then the writer of Hebrews says, "For time would fail me to tell of Gideon, Barak,.
Samson.".
How in the world Samson made it into the Hall of Faith?.

$^{41}$I will never know..
Japheth, David, Samuel, and the prophets..
Missing from that list is Daniel..
Another omission for some reason, Daniel didn't make it to this list..
Solomon did not make it to this list..
The brightest of all of the kings of Israel, but yet at the same time, he had a lot of.
heart motives..
He could not overcome the stain of lust..
Solomon was omitted, and Elijah was not named in this list..
Why?.
Think through a little bit about that..
In verse 35 though, as he finishes this list, he does mention one thing..
He says, "Included in this list are the women who received back their dead by resurrection.".
Clearly referring to the story that we're going to look at today with the widow's son..
So now you're sitting there, you didn't tell me, you forgot to mention Elijah, but you.
pinpoint the action or the event in which he raises the widow's son from the dead, Elijah.
raises, and you add that story in there..
What were you trying to say to us?.
I have some theories, and we're going to go back to it at the end of the sermon once we.
look at this particular story..
And here's the big ideas we look at verses 17 to 24 today..
Our main idea, our big idea is this..
Miracles unfold in the hands of humble hearts..
God uses and performs miracles in the hands of people with humble hearts..
Miracles unfold in the hands of humble hearts..
Let's take a look at 1 Kings chapter 17..
And he writes this in verse 17, "After this, the son of man, the mistress of the house,.
became ill, and his illness was so severe that there was no breath left in him.".
Verse 17, it tells us, "After this, the son of the woman.".
After this, what happened in verses 1 through 17?.
For those who don't remember, let me just give you a quick recap beginning in verse.
17 to where we are right now..
Remember the first appearance of Elijah..
Who did we say Elijah was?.
His name really means, "My God is Jehovah.".
El meaning God, Jah meaning Jehovah, I connects the two, Elijah..
My God is Jehovah..
Where is Elijah from?.
He was the Tishbit of Tishbi in Gilead..
What does that mean?.

$^{81}$Well, Tishbi is somewhere in Gilead..
We know where Gilead is in the map at that point in time..
But no one knew where Tishbi was..
This Tishbit from Tishbi in Gilead, what does that mean?.
He came from obscurity..
No one came..
So this guy, this Tishbit of Tishbi, where no one ever knew where he comes from, gets.
and has a meeting with the king..
It happens over in the first three verses, King Ahab of the northern kingdom..
Post David, post Solomon, the whole thing went to disaster..
You know, soon separate into northern kingdom, where Elijah, the book of Obadiah, Queen Jezebel.
is at, and then the southern kingdom, where Israel, where Mark is right now, a lot of.
times Israel, Jerusalem, all those guys, that's where Nehemiah, Ezra, all of that takes place..
But we're focusing on King Ahab of the northern kingdom..
Well, this guy from Tishbi gets a meeting with King Ahab, and he says this to him, "Neither.
dew nor rain except by my word.".
Because you people are worshiping Baal, the god of fertility..
Let me then show you how powerful my God is..
There will be no more dew, there will be no more rain, nothing, except by my word..
And he tells him, here's my motivation, "It's before the God before whom I stand.".
And then later on, as we end up in verse 6, God at the climax of this event, he ushers.
and brings him to Cherith..
Cherith, solitude..
Cherith, he's by himself..
He's fed by ravens, all of that..
He's fed by the brook, water by the brook..
But the drought is three and a half years, and soon the brook dries up..
And God says to Elijah, "Go to Zarephath..
I now want you to go from Ahab to Cherith to Zarephath.".
Well Zarephath is a hundred miles away, so this guy has to walk a hundred miles..
King Ahab and soldiers are after him..
So then he not only goes, walks a hundred miles, he goes to Zarephath, which is the.
home of Queen Jezebel, who is married to King Ahab, who is the worst of all the worst..
And God then says to him, "I want you to go, and a widow will meet you.".
A widow?.
Yes, a widow..
Shows up to the widow's house..
The widow tells him, because she's gathering sticks, "I have no food..
I have a little flour..
I have a little oil..

$^{121}$I'm going to go bake some bread so that you, my son and I, will eat it and die.".
And Elijah's looking at him, he says, "Really?.
Of all the people in Zarephath, you made me walk a hundred miles to meet a widow who's.
going to tell me that she has no oil, no flour, and she's now going to make food to die?".
Yes, yes, that's what I'm going to tell you, because God knows what we need..
God knows what we need..
And the story goes that you're going to have to continue to do all these things and go.
and do..
And the second thing that we talked in sermon is not only does God know where you are, he.
knows where Elijah was at the beginning from chariot..
He wants him to go to Zarephath..
He then knows exactly where he needs to go..
You go to Zarephath, and he knows what you need..
I'm going to give you the widow..
It may not be what you want..
It may not be what you think you need, but it's what you need according to what I think.
that you need..
That's what God is calling him to do..
And then lastly, that I want you to go and do..
And that's what he says to the widow..
But now this thing comes, and the story goes in verse 17, after this, the son of the women,.
the mistress of the house, after all this thing happens, he becomes ill..
And his illness was so severe that there was no breath in him, and he dies..
And in verse 18, the widow says to Elijah, what have you against me, O man of God?.
You have come to bring my son to remembrance and to cause the death of my son..
That's what she believes..
A couple things I want to say about this..
Number one, the woman has a wrong perspective on sin..
Somehow or another, when this thing happens, she believes that her sin is being punished.
by God..
Now here's the one thing as Christians that we have to believe, right?.
We know that this woman has some sin..
We don't exactly know what this sin is..
We do know that the sin is eating her up..
There must be some guilt within that sin that keeps popping up, that is at the forefront.
of what she's thinking..
And eventually, for some reason, as we read, we infer that she thinks that eventually this.
sin will come in the open..
But when her son dies, she immediately thinks, God wants me to remember the sin and causes.
my son to die..

$^{161}$Couple of things I want to say about that..
Sin is present in the world, yes..
There are implications to your sin..
There are consequences to your sin..
There is also sin in the world..
That's why diseases, natural disasters, all of that, because of the presence of sin in.
our world, bad things happen..
So there are consequences and ramifications to our sin..
There is sin in the world because of the fall of Adam and Eve..
But here's something very important..
God will not kill other people or do harm to other people just to punish your sin..
God is not going to kill the widow's son just to punish the widow..
There are sin, there are ramifications to sin, there's sin around the world, but God.
is not a sadistic God..
That's something that the widow needs to understand..
She has a wrong perspective on sin..
Not only that, what happens?.
She's blaming Elijah because she's hurt..
Because sometimes we say and do things when we are hurt that we regret..
Sometimes we say things when we are hurt, when we are angry, that we can't take back..
As Asians, for most of us here as Asians, we don't know how to deal with anger..
One day, a lot of times, I may need you to help me understand..
As an American Asian born here in the U.S., sometimes I may not understand, but I find.
with my parents, they're very private about everything..
They just keep everything inside..
They never open up..
I've never understood why..
My mom always just says, "No, no, it's no problem..
You just got to put your head down.".
And they never share about their feelings..
In all my life, my parents have never shared what they were thinking about..
They have never been open and they never want to share with anyone..
Everything is a private matter..
So eventually one day when all of this pressure builds up and once in a while you just explode..
And here a lot of times we see the widow gets so frustrated, she says things to Elijah that.
she probably would forget..
She takes out her anger at Elijah..
Let me tell you a story..
Have you ever been to a bad funeral?.
Now funerals are kind of bad, but have you just ever been to a funeral that just went.

$^{201}$completely berserk?.
I went to one funeral one time..
It was a very close friend and she had about six or seven siblings..
And then the siblings did not get along..
The mother had passed away and now the father had passed away..
The older sister gets up and she's supposed to do the eulogy..
She gets up just like this and she starts doing the eulogy..
And so she starts talking about the dad..
Then suddenly she just went off course..
She started blaming all of the siblings about how no one took care of the dad, all the problems.
with each of the different siblings, and how she was the sister who took care of him when.
all the other sisters were not there..
Well, one of the brother-in-laws stood up in front of everyone and says to the sister-in-law,.
"I am so sorry for this..
She's not mentally right.".
So he gets up and he's trying to grab and pull her off stage..
And then she starts running around this funeral home that we're in trying to read her eulogy..
Next thing you know, the sisters stand up and they're all trying to corner her and she's.
running around this funeral home finishing up her eulogy so everyone knows how bad the.
sisters were..
And then suddenly they just start swearing at each other..
And we're all sitting there watching..
And you're sitting there, "This is terrible..
This is terrible..
Everyone's swearing at each other..
The casket goes out..
They're swearing at each other as the casket's going out.".
It was interesting..
I bet you as you look back, as you can imagine, the family was never the same..
As you can imagine, there was things said that you just regret..
There was things said that you wish at the key of the moment that was never said but.
can never be taken back..
I remember my mom..
My mom is, I told you a lot of times, she's not a very open person..
My mom could be a hard person and she keeps all these anger things built up..
And my mom a lot of times got into fights with my younger sister so much and they would.
say things that I'm sure they would regret..
It's been 20 years now..
I have not seen my sister..
She moved out, got into it with my mom, moved out, went, disappeared somewhere in New York.

$^{241}$and cut off ties with my mom..
She felt like as brothers, me and my twin brother will always side with my mom..
I'm sitting there, I said to her, "You know, she's like this..
You have to learn to forgive her.".
Because a lot of times she struggles with so much anger but my sister couldn't do it.
and she just couldn't stand it about my mom..
And she disappeared for 20 years..
My mom tried to call her..
My mom tried to apologize for once in my life, I've never seen that..
My mom tried to tell my sister that she took back everything she said..
My sister never returned her phone call..
It's been 20 years..
So sometimes in the course of life when we get angry, we will say things that we regret..
We will say things that we wish we could take back..
And after all this Elijah and all this stuff and she sees this widow, the son's widow,.
widow's son dies..
And the widow takes out her frustration..
You know what Elijah's response was in verse 19?.
And he said to her, "Give me your son.".
Elijah never told her that things were going to be okay..
Elijah never rebukes her..
Elijah doesn't argue with her..
Elijah doesn't tell her and reason with her, "Hey, I'm not at fault here, all right?".
And help her see what she was wrong in her thinking..
Elijah doesn't even speak scripture to her and tell her to trust God and place your faith.
in God and pray to God he'll save you..
He doesn't do all these things..
For some of us, a lot of times we've been in this situation, he doesn't even say, "Hey,.
how are you feeling?".
"Well, how do you think she's feeling?".
He also doesn't say, "I understand.".
I think all of us have had to deal with people who are hurt..
Death of a beloved one, loss of a job, circumstances that happen..
And sometimes when we're counseling different people, we want to solve their problem..
We want to tell them it's going to be okay..
That's not what they want to hear at that time..
We want to tell them you got to pray to God..
That's not what they need to hear at that time..
We want to tell them, "Hey, look, God is in control.".
Or we sit there and say, "How are you feeling?".

$^{281}$It'll get them even more mad..
One time I said to an individual who was diagnosed with cancer, I said, "How are you feeling?".
She says to me, "How do you think I'm feeling?".
Never said that ever again because I realized that's probably not the best thing to say..
Never even say, "I understand how you're feeling," because sometimes you don't understand..
You know what Elijah did?.
He sat there and listened..
He took it all in..
He doesn't sit there and give her all the solutions because you know and I know sometimes.
when you're really mad and you're really angry, you know what you did is wrong..
You know what you're supposed to do..
And you're not looking for people to tell you what to do..
You're looking for people to listen to you..
I had two close friends..
I had one individual who was going to retire early at 47 years old..
She was a paralegal and she had saved up her money at 47, 48..
Her goal was to retire early and she bought her house, everything..
On the week she was supposed to retire early, she went in for a checkup because her insurance.
was about to be over and the doctor told her she had stage four cancer..
And now for the last since leaving her workplace, she's been in chemo, radiation every single.
week..
And we talk at least once a week and half the time, 80\% of the time, it's her expressing.
anger and frustration over the situation and she's a Christ follower and she sits there.
and tells me how she's frustrated..
She doesn't know why this happened to her and what can you say but listen to her..
I had another person, my wife's aunt, who is not a believer..
Her husband went to bed that night..
I remember that because we saw him at a family party and it was just before Thanksgiving..
He said, "We'll see you at Thanksgiving.".
He went home that night and he had a stroke and he passed away..
And she calls me up because she knew that somehow or another I was kind of a pastor..
How do you deal with all this stuff?.
And for subsequent times, once a month I have checkup on her..
She's not a believer and she tells me all her pain and all I could do is listen to her..
Because sometimes in the midst of all this stuff, you know what they need?.
Is it for you to give them advice?.
Is it for you to tell them God is there with them?.
Listen to all this stuff but just be there to listen..
Elijah took the brunt of this widow's pain because this widow said some things that she.
probably will regret if she looked back..

$^{321}$And he just listened to her and he says one thing to her..
"Give me your son.".
And he took him from her arms and he carried him up into the upper chamber where he lodged.
and laid him on his own bed..
And here for the next couple of verses, I'm going to show you Elijah's prayer because.
I find this prayer is so powerful..
We'll glide over it but in the future I was thinking maybe we could just dive into it.
a little bit more..
And he cried to the Lord, "Oh Lord my God..
Have you brought calamity even upon the widow with whom I sojourned by killing her son?".
He brings the boy to his bed and he cries out to the Lord, "Oh Lord my God..
What have you done?".
Because he doesn't know..
He doesn't know why God killed her son..
He doesn't know the answers to all these things..
And when you have a deep enough relationship with the Lord your God, you're able to come.
to God with your pain, you're able to lay it before him and say, "Help me understand.".
Because I don't know..
Why have you done this?.
Why is this happening?.
I don't know..
Because sometimes understanding and following the Lord, we don't know why certain things.
happen..
But then he also intercedes for the boy..
He takes the widow's son and he intercedes..
He comes to the Lord and says, "I don't know.".
But then he struts himself upon the child three times and cries to the Lord, "Oh Lord.
my God..
Let the child's life come into him again.".
He intercedes on behalf of the widow's sons..
He wrestles with God..
He prays specifically for the boy's life to be saved..
But he's not afraid to pray for the impossible..
Remember when we talked about Zerah for some of those guys who were here and for one of.
the things..
One of the things that we talked about in that sermon is you have to overcome physical.
impossibilities..
You have to believe that God sometimes in his realm will work miracles in your life..
Elijah was not uncomfortable with praying for impossibilities..
He prayed for a miracle..

$^{361}$But he also realized God may not answer his prayer..
And it's okay..
You're seeking God..
You're wrestling with God..
You're coming before him praying for answers..
You're asking for miracles to happen..
You're asking him to do the impossible..
But you have to trust that he may or may not answer that prayer based upon that will..
And you have to be okay with that..
See the widow's son who expresses her anger..
We see Elijah listening to what he has to say..
He doesn't tell her what to do..
He brings her, lays him on the bed..
He's also wrestling the same thing because he doesn't know..
Comes before the Lord but he intercedes on behalf of the widow's son..
And the Lord listened to the voice Elijah..
And the life of the child came into him again..
And the boy revived..
And Elijah took the child and brought him down from the upper chamber into the house..
He delivered him to his mother..
And Elijah says see your son lives..
And look at verse 24 if you get a chance to underline verse 24..
And the woman said to Elijah, now I know you are a man of God and that the word of the.
Lord in your mouth is truth..
Here's what she didn't say..
Now I know Elijah you are a powerful man..
Now I know Elijah I can trust you..
Now I know Elijah man I need to worship you..
And that you have all this power that we can depend on..
That's not what she says..
She says now I know that you are a man of God..
And that the word of the Lord in your mouth is truth..
When people look at you, do they see something different in you?.
Do they see that it is God working in your life?.
They're not just seeing you but they see beyond and see God in the presence of your life..
That's the goal isn't it?.
That's the goal of why we do what we do isn't it?.
That when people when we do good and we allow and intercede for others, when we serve others,.
when we help others, all of these things, it's not because people say look at him..
Look at me..

$^{401}$Look at how great I am..
Look at how perfect..
No, no, no, no..
Our goal as stewards of the gospel is to say look at how great God is..
Look at how amazing God is..
And Elijah had this what we call a confident humility..
He was confident in what God could do..
He was confident in what God called him to do..
And he was humble enough to know that it was God working through him..
And everything that he did was to bring glory to God in all that we did..
If you guys are familiar with the name Tim Keller, Tim was a mentor of mine..
I love Tim Keller..
If you are not familiar with Tim Keller, I encourage you to Google and look him up or.
read one of his books..
Tim Keller was a man, if you ever met him, he was a pastor in New York, started becoming.
famous with his famous book Reasoning with God at age 53..
And as a result, we will look upon, he died two weeks ago..
When we remember Tim Keller, we will see him as such an influential leader who changed.
the church, the modern church as we know it today, that we will compare him to what Billy.
Graham did back in the 70s and 80s is what Tim Keller brought to us today..
But on the side note, as a result, I didn't know Tim very well..
I got a chance to oversee his Los Angeles conference, his big Chicago conference, and.
we interacted with each other very much so..
Tim was one of the brightest individuals you'll ever met..
Tim was a pastor, Tim was a scholar, Tim was a theologian, Tim was a professor, Tim had.
a chance, he read so much, he thought so much, he engaged so much, but the man was humble.
in a way that I've never seen a person be..
Tim got famous when he was a little bit older, age 53..
One of my friends asked Tim one time and says to Tim, "Tim, how does it feel being famous.
and wanting everyone to sign your books and all that?".
And Tim says to the woman, he says, "You should never ever feel comfortable that people worship.
you..
You should just never do.".
Tim was a person who dialogued in the day when, I was just talking with Rachel about.
this, in the day when people were just funny about politics and race and all the other.
stuff, Tim would just listen and dialogue all the time..
Especially in the last three years of his life, people would just criticize Tim because.
he says, "You're not playing on one side, you're not on the right, you're on the left,.
so which side are you?.
Pick a side.".

$^{441}$And Tim would never engage in those debates..
Princeton Seminary gave him an honorary doctorate, but Princeton Seminary had to rescind it because.
a lot of the students and faculty and some of their donors says, "Why are you giving.
this person who believes in marriage and who doesn't believe in LGBTQ, we're not giving.
him, we're not giving him this doctorate.".
So they had to rescind the doctorate from him..
Tim wrote a kind, kind note that says, "Thank you for even considering me and giving me.
this honor.".
When people criticized him on social media, you would never hear Tim get on social media.
and criticize him back..
Many times Tim would just call that person and says, "I want to understand your views,.
help me understand.".
One of the things that someone said to Tim, the most impressive thing about him was you.
never came away exclaiming, "Isn't Tim Keller amazing?".
But when you looked at Tim Keller, all he wanted you to know is, "Isn't Jesus amazing?.
Isn't Jesus wonderful?.
Isn't Jesus beautiful?.
Isn't Jesus captivating?".
That's our life as Christ followers, isn't it?.
We live like Elijah in which when miracles happen because of what God is able to do through.
us, that they're going to look at our life and say, "Wow, look at God..
Look at what God has done through him and shown him.".
Because that's what we're supposed to be called to do as Christ followers..
So let's go back to Hebrews 11..
Why wasn't he in the hall of faith?.
And why instead are we told in verse 35 that the dead, women of the dead, and it says this,.
"Women received back their dead by resurrection.".
I really believe that the author of Hebrews was trying to honor Elijah well by not talking.
about Elijah, but rather about what he did..
Let me explain this..
Elijah was all about others first..
Elijah was never about himself..
Elijah never did anything to give him the platform, but it was always about what he.
did..
My theory is that a lot of times the author of Hebrews honored Elijah well by not talking.
about Elijah, but rather about what he did..
Because Elijah was never about himself, but always about other people..
That is the conclusion of chapter 17..
Four weeks in chapter 17..
We learn that God could use me wherever I come from..

$^{481}$He used Elijah to Tishbit from Tishbit in Gilead, a place of obscurity, a rural person.
that no one ever knew..
God could use me wherever I come from..
I am here by God's timing..
God used Elijah and called him to Ahab, he called him to Cherith, he calls him to Zarephath,.
all of these things..
I am here by God's timing because God sees me wherever I am..
Elijah goes and goes wherever he is called to do..
He obeyed him to go to Cherith..
He obeyed him to go to Zarephath..
He seeks his will..
Before whom I stand is his motivation as we learn in chapter, in verse 2 of chapter 17..
He accepts what he is given..
He accepts the widow..
He expects, he accepts whatever happens with the widow who had no flower, who had no oil,.
who thought she was about to die..
He intercedes for others..
He wrestles with God in an authentic way and he gives God the glory but never takes on.
the glory..
Miracles unfold in the hands of humble hearts..
Elijah tells the widow go and do..
Go and do..
Go and do, go and do because miracles unfold in the hands of humble hearts..
Go and do..
Let's pray..
Father, thank you for the life of Elijah..
Thank you for all that you have taught us through this chapter 17 full of riches of.
the adventures with Ahab, with Cherith, with Zarephath, with the widow's son, and ultimately.
a lot of times you reminded us this morning it is not about us and it will never be about.
us, but how God you work through and in us in our role as stewards to give you the glory.
in all that we do..
May we learn from the journey of Elijah..
May we learn what it means to intercede for others, to die for ourselves so that we may.
serve others every single day, but to give all glory and credit to you..
As we go forth, may you guide us every step of the way..
May we go forth on this day in all that we do in our work, in our school, even for some.
of us who are on a summer vacation to honor you every step of the way..
In Christ we pray, amen..
Mark is back next week and have a good day..
Thank you..

$^{521}$[BLANK AUDIO].
\newpage



\section{8:21-22}
\label{sec:7Sx7wB7NbaI}
\textbf{Cost of Following Jesus (Matthew 8\_21-22) - Rev. Mark Chan}
\newline
\newline
連結: \href{https://youtube.com/watch?v=7Sx7wB7NbaI}{\texttt{https://youtube.com/watch?v=7Sx7wB7NbaI}} ~~~~ 語音日期: 2023-06-13
\newline
\newline
\hyperref[sec:kj0zBwcuLqE]{\small{< < < PREV SERMON < < <}}
~
\hyperref[sec:index_chronic]{\small{[返順時目]}}
~
\hyperref[sec:index_scriptual]{\small{[返順卷目]}}
~
\hyperref[sec:TVT11XSJIWY]{\small{> > > NEXT SERMON > > >}}
\newline
\newline
$^{1}$Well, so happy to see all of you, and as I said, I came back last night about 9 to 10.
p.m..
And safe and sound, good to see my family, good to see all of you..
And again, we went as a group, 27 people, 12 of them from this church and 15 of them.
from all over the place, Hong Kong, West Coast and East Coast..
And what a blessing to have all of us together for the past two weeks and to walk the land,.
actually..
This is one of my favorite pictures, and I have tons of pictures, and I want to just.
give me some time to reorganize them..
And starting next Sunday, I'll design a series and introduce to you what happened there..
This is one of the favorite pictures I have because we do have, this round, we do have.
a couple of semi-professional photographers, and they did take some good pictures..
And yeah, I don't know what I'm doing with my right hand, but in any case, this is one.
of the places that we went to in the past two weeks, and it was an additional place.
because we found some time..
So I talked to the driver, "Hey, can we just take some off-track places?".
And so we went..
So we were on that day, on this day, we were going from Jerusalem, which is high on the.
mountain, and all the way to the east, to the Judean wilderness, all the way to the.
Dead Sea and Jericho and all that..
Remember that, right?.
The road from Jericho down to, the road from Jerusalem down to Jericho..
And there is a road, right?.
And today is the Highway 1 in Israel..
So we were passing different places, and I talked to the tour guide, and I talked to.
the driver, "Hey, do we have time to go over there?".
And we made a left turn into this off-track places..
And the place was called Wadi Q, Q-E-L-T. Wadi is in Arabic..
It means some dry valley..
Usually it's just seasonal runoff will go after the, go into the valley, but usually.
it's dry..
And what you see on the left is actually the Wadi Q, and it's one of the valleys coming.
out of the mountain of Jerusalem all the way to the Dead Sea Valley..
So it's kind of exotic place..
Yeah, because we were in this big bus, right?.
27, almost 30 people..
And then we would drive, I wasn't driving..
I wanted to drive, but the driver was one, just amazing driver..
And he drove the bus on the cliff of Wadi Q. And we can see, you know, down the window,.
we can see the bottom of the valley..

$^{41}$And it's just amazing view there, right?.
On the left, lower left, you see a monastery..
The monastery is called St. George Monastery, and it's a monastic community that existed.
since AD 400..
So it's been around for 1,600 plus years..
Generation after generation, there are Christians, Eastern Orthodox Christians who gather there.
and to live a contemplative life and keeping the tradition of Christianity..
So it's just amazing..
And this picture did not show up, perhaps it did..
We were standing on, right on the top of the cliff and looking down to Wadi Q and St..
George Monastery, and on the other side of the valley, we saw a shepherd grazing the.
sheep on the cliff..
And there were probably like hundreds of sheep and running up and down the cliff and, you.
know, finding food, right?.
And this is actually Psalm 23..
The Lord is my shepherd, and I have lack of nothing..
And He leads me into green pastures and quiet waters..
And what do you think of when you read words like this in Psalm 23?.
Green pastures and quiet waters..
We always think of the Scottish Highland, right?.
Where everything is green, right?.
This is actually Psalm 23..
Green pastures..
Can you see anything green in there?.
Maybe a little bit, right?.
Maybe a little bit..
But that is the green pastures of Psalm 23..
Just enough for you, right?.
If you follow a wise shepherd, the wilderness is the place for grazing sheep..
You don't graze a sheep in your olive grove, right?.
In your wheat farm, right?.
Because those are for surplus, for humans, right?.
So animals, sheep and donkey and goats, they go out to the wilderness, which is just east.
of Jerusalem..
Often we read of land without water, land without water and wilderness..
This sort of picture is the wilderness..
And you think about Psalm 23..
Green pastures is the little green stuff underneath of the rocks..
The little green stuff, you know, between the cracks and sheep can go out and find them,.
right?.

$^{81}$And then they can get just enough, the essentials to be fed on..
That is the theology of Psalm 23..
You see quiet water?.
No, just seasonal runoff, right?.
As I said..
But if you follow a wise shepherd, you get just enough water for you to live..
Just enough for you to survive and even thrive in a land like this, right?.
It's beautiful..
So, you know, in order to understand some of the passages in the Bible, I think it's.
so essential for us to walk the land..
One of my professors who passed away some years ago, Anson Rainey, he's a professional.
biblical geographer..
He always, you know, when you take his classes, we always go out to the field, places like.
this..
We don't sit in the classroom..
We don't sit in an air conditioner, air conditioner room to listen to a lecture..
Most of the time we walk out there..
And sometimes we walk with him..
And Rainey would say, you know, this is the only institute, only classroom that we have.
a one-to-one map..
Indeed, it was one-to-one, right?.
Every step of the way is not the Bible map that we found in the Bible, but it's actually.
walking the land, right?.
So he always emphasized one thing is that we have to read the Bible with our feet on.
the ground..
Our feet on the ground, right?.
Not in the ivory tower, but with our feet on the ground, right?.
So yeah, it's beautiful..
I haven't been to Waterkill for years..
I've led tours for the past over 10 years..
And usually because we want to hit different sites, different spots, and we want to go.
out there and to see as many places as we can..
But it was such a unique opportunity..
We can just pass through the wilderness and spend extended time there..
So there I was actually reading Psalm 23..
It's not the Scottish Highland, but God is always faithful, even in wilderness, right?.
Amen to that, right?.
Standing on the ground as we read the Bible, right?.
So yeah, that's what, this is one of my passions..
And I want to thank the church for letting me and letting the pastor to go away for a.

$^{121}$little bit and to be a blessing to many other people..
And most of the people, we go out to different sites and their first response has always.
been, "Wow, ah, yeah, I got it now.".
Right?.
Yeah, you got it now, right?.
So walking with our feet on the ground as we read the Bible, right?.
So perhaps next year you can go with me, right?.
(speaking in foreign language).
Okay, we can go anytime..
I always tell people you can go, you know, next week..
As long as you got a passport and you got a visa card, that's all you need, right?.
But I did bring something back this time..
So as a tour guide, I got some free gifts..
So I got this mineral hand cream from a hover factory..
Yeah, so all the way from Dead Sea, right?.
So you can get a little bit after the service, just like a million, you know?.
(audience laughing).
Share, okay, share..
Don't drink the bottle, all right?.
Don't drink the bottle, all right..
So it's good to be here and I want to reflect on a short passage today..
And I want to tell you, we, in order to understand the message, the impact of the passage, we.
have to understand the context..
That's why we have this theme of walking the land..
You need to understand the biblical background of scripture in order to appreciate what Jesus.
has said..
And this passage is Matthew chapter eight..
Let's read that first..
When Jesus saw the crowd around him, he taught them, the disciples, and he instructed the.
disciples to cross to the other side of the lake..
And this lake is called Sea of Galilee, right?.
The Lake of Galilee, Sea of Galilee..
So they crossed to the other side, away from the crowd, and one of the teachers of religious.
law said to Jesus, "Teacher, Rabbi, I will follow you wherever you go.".
But Jesus replied, "Foxes have dens to live in and birds have nets, but the Son of Man.
has no place even to lay his head.".
Then other of his disciples said, "Lord, first let me return home and bury my father.".
But Jesus told him, "Follow me now..
Let the spiritually dead bury their own dead.".
Now first of all, this is a passage, one of the passages about discipleship, about the.

$^{161}$cause of following Jesus..
I must admit the very first time when I was much younger, when I first read this passage,.
I was kind of confused..
This is actually one of the tough sayings of Jesus, one of the most difficult to understand.
sayings in the Bible..
And what he said is that, especially the latter part of it, verses 21 to 22..
Other disciples said to Jesus, "Yeah, I want to follow you, but first let me return home.
and bury my father.".
But Jesus said, "Yeah, go ahead and bury your father and follow me.".
No, he didn't say that..
He said, "Follow me now..
Follow me now..
Let the dead bury their own dead.".
Now that's what I call really confusing, right?.
It's really tough..
What is going on?.
I just want to bury my father, right?.
So how long would that take?.
Maybe a day or two, right?.
You think about modern burial practices..
Can we just do that, right?.
And then I'll follow you?.
Because by all means, in Jewish culture, burying one's father is one of the most honorable.
thing for someone to do it, right?.
Especially for son, right?.
In that patriarchal culture, you have to honor your parents, right?.
By the way, that's one of the Ten Commandments, right?.
To respect and to honor your parents..
And what is going on here?.
And I think in order to appreciate the impact of what Jesus has said, you have to understand,.
again, it's the context, the background of the Bible..
The Bible is not something, you know, out of the blue, like a heavenly book..
It's dropped down from heaven to earth, and then we can read it..
And the message and the languages are magical..
No, the Bible has its own context, right?.
In order to understand what the Bible says, we have to understand the background..
And here we go..
We have to understand in Jesus' time, burial is a huge thing, just as today, right?.
No matter you like someone or you don't like someone, if you know him and he passed away,.
you probably have to change your clothes, all black, and you go to the funeral home.

$^{201}$just to, you know, pay the final respects, right?.
It's a huge thing, but in Jesus' time, it was a bigger, huge deal..
So in the Jewish culture, in Jesus' time, they would do this in their burial..
So from left to right, when someone died, what to do with the body?.
What to do with the dead ones?.
First, because in the ancient days, they didn't have refrigerator and ways or technology to.
preserve the body, and it gets stinky really quickly..
That's why someone died, and then they have to bury the body almost ASAP, almost immediately..
They have to put the body in the tomb, and this is the typical Jewish ritual of burying.
someone..
So first of all, someone died, all the way to the left..
You're going to go into the burial, almost within 24 hours..
And then you kick in a seven-day mourning period..
It's called Shavuot..
Shavuot means seven in Hebrew..
So seven days, they would mourn at home, almost privately..
And then, after the seven days of Shavuot, they would go into Shaloshim..
Shalosh is three..
Shaloshim is 30..
Thirty days of mourning, and then some of the members of the family and extended relatives,.
they will come to your home and visit you, bring you food or bring you gifts and all.
that stuff, and 30 days..
And then you wait for about a year..
So you put the body in the tomb, and then it will take time for the flesh to declay,.
right?.
To be all gone..
And then almost one year later, you will return to the same tomb and to gather the bones,.
right?.
If I die, you bury me in the tomb, right?.
And then one year later, you come back and you don't see me anymore..
You see my bones, right?.
And that's what they do..
It's the second burial..
Second burial..
And they put the bones into some stone boxes called asherurim..
And these boxes, you can still see all over Jerusalem..
And this is one of the pictures of the first century tomb, and they call Kokkim tomb..
Now this is what I said in the timeline..
Once the body is there, once someone dies, and you have to prepare the body, and you.
bring the body into the chamber of the tomb, okay?.

$^{241}$By the way, this tomb is right next to a coffee shop..
Yeah, it's very nice..
You can have a coffee and a croissant right in front of the tomb..
Interesting..
It's on Mount Scopus in Jerusalem..
Anyways, you see the entrance of the tomb, right?.
Back in the days of Jesus, 2,000 years ago, the tomb is dated to the time of Jesus..
It was covered..
It was a cave, actually..
So now it's open air..
Anyways, it was covered 2,000 years ago, and the ceiling of the cave did not collapse..
So you go into the entrance..
Can you see that?.
There are three side benches, right?.
And then you bring the dead body on the bench, and you put spices and whatever you need to.
prepare the body for..
And then you're going to prepare it well, and then you're going to stuff the body into,.
you know, see the small holes?.
Those are called coke..
K-O-K-H..
You stuff the body in there, and then one of the cokes, or one of the kookim, or kookim.
is plural, coke is singular..
So as you can see, there were a lot of kookim, right?.
Because most of the tomb is family tomb, right?.
So my father used it, I used it, my son used it, my grandson will use it..
So there were a lot of people being buried in the same place, in the same tomb, because.
they belong to the same family..
So that's why they were, it's like a refrigerator, all right?.
You put the body in the kookim, and then what do you do?.
What do you do?.
You go home for seven days of mourning, and then 30 days of mourning, and then?.
And then you wait, right?.
For one year..
And the body stay in the tomb, in the kook, or kookim, and then let the flesh decay, and.
when everything is gone, you go back one year later for the second burial, which is the.
ossuary, right?.
Or you see, one year later, you open the tomb again, go inside, and gather the bones into.
boxes like this..
Some of them are decorated with some geometric design..
Some of them are bigger, some of them are smaller..

$^{281}$And the length of the stone box is the length of your longest bone..
Where is your longest bone in your body?.
Right here, what do you call it?.
- Femur..
- Femur, right?.
Femur, yeah, this one..
So for me, I probably need a bigger box..
Anyways, it's just a matter of, it's a business of recycling, right?.
Because they are doing a lot of different stone work and stone building, especially.
Herb the Great, he built or expand the temple in Jerusalem so big, and that's why there.
were so many leftover stones, right?.
And they make vessels out of the stones, and one of the vessels is the ossuary..
So they will go back one year later and gather all the bones and put them all in the stone.
box..
And then you put the box by the entrance or by the, just on the side of the chamber..
So the dead upon dead will be gathered in the tomb..
And then the new dead will be buried in the same tomb, and then you go through the same.
cycle, three days and, what, seven days and 30 days, and one year, you go back for a second.
burial, gathering the bones and put in the stone boxes..
So with that background, you can actually understand what happened, first of all, on.
Good Friday..
All right, Good Friday..
Good Friday, when Jesus died, his body has to be buried, buried, right, immediately,.
right?.
And remember what time Jesus died on Friday?.
Before Shabbat, right?.
Before Shabbat..
Once Shabbat kicks in, you cannot do any work..
You have to wait for another 24 hours, right, to continue to work..
So that's what happened on Good Friday..
Jesus died, and quite unexpectedly, and the women disciples, or the disciples, would put.
his body into one of these tombs, and the Bible says it's a new tomb, a new tomb for.
the important people..
So they would put his body on the bench, and because Shabbat is coming, and they could.
not finish the work, right?.
They got to get out of the tomb, because it's Shabbat, right?.
You got to purify yourself in order to prepare yourself for Sabbath of 24 hours..
So they exit the tomb, and then when did they come back?.
After Shabbat, right?.
After 24 hours..

$^{321}$And then on the first day of the week, it's the first day of the work..
So they return, the women return to the tomb, and they try to, probably they would ask someone.
to open the tomb again, and come back, and just finish off whatever they did not finish..
It was kind of in a hurry, so they didn't finish all the ritual of mourning and burial..
So that's what they were planning to do..
Returning to the tomb for the first burial, put Jesus' body into what?.
One of the kukims, right?.
Stuff him..
Stuff him into one of the holes, and wait for another year, and come back for asurery..
That's what happened on Good Friday..
With that picture, or with this timeline, you would understand what, oh this is one.
of the diagrams, this is better, right?.
The entrance is on top, or it's middle of the picture..
You remove that big stone, and you get into the cave..
You have the center chamber, and then you have kukims..
In any case, with this timeline, you can understand what Jesus said..
He said what?.
First of all, and other disciples said to him, "Lord, first let me return home and bury.
my father.".
But Jesus replied, "Follow me now, let the dead bury their dead.".
That's what happened here..
Now I want to ask you a question..
With that background, I explain..
This and other disciple who came to Jesus and said, "Hey, first let me return home and.
bury my father.".
For this disciple, assuming his father just passed away, when did this disciple approach.
Jesus with that request?.
When?.
Is it before first burial, or before second burial?.
There's no way for this guy to have any room to come to Jesus and follow him and say, "Hey,.
first let me go back and bury my father.".
There's no way that he can find time before the first burial..
You understand what I'm saying?.
So it has to be this guy, when he approached Jesus, he already done his first burial..
And he's waiting for another year, right?.
Because the second burial has to be done, right?.
So when this and other disciple approach Jesus, it has to be between the first burial and.
second burial..
So what he meant when he said to Jesus, "First let me return home and bury my father.".
He was referring to the second burial, the ossuary..

$^{361}$So he is asking Jesus, "Hey, I want to follow you, but hey, let me have some time..
And let me get the ossuary done first, the second burial.".
So essentially, he was asking Jesus for a year time..
"I want to follow you, but next year, because I want to get something done, right?.
The second burial..
It's not done yet.".
Do you understand what I'm saying?.
There's no way his father just died and then he would come and follow Jesus, right?.
This is the most dishonorable thing to do, right?.
His father died, indeed, it's pretty sad..
But he has already done his first burial because first burial happened almost immediately,.
right?.
There was no time for this guy to approach and follow Jesus..
And this, another disciple actually is between the first and second burial and asking Jesus,.
"Hey, can I have some time before I follow you?.
Let me first go back home and bury the dead.".
Now in rabbinical tradition, for all the sons, especially the firstborn son, there are three.
rituals in their lives that are most important..
Think about your life, in your lifetime..
What are the most important things that you could have done?.
You could have done, right?.
For the Jewish tradition, the first most important thing, most important ritual for all Jewish.
men, the first one is circumcision..
It's done on the eighth day after the child is born, right?.
So it's done..
That's the most important ritual they have to go through..
The second most important moment in their life or ritual in their life is wedding..
Get married..
That is the biggest day..
And the third one, actually according to ancient literature, Jewish literature, is the second.
burial..
It's the ossuary..
Because you will finish the duty, the honorable duty of the firstborn son..
So it's not really the first burial that is important..
The most important thing is wait for another year until you come back to the same tomb.
and gather the bones of your father into ossuaries..
That is the most important ritual and most important social honors one has to gain..
Now with this background, you understand the impact of the passage..
You understand what Jesus is trying to say..
Because the entire Bible from Genesis to Revelation, no one can claim the authority higher than.

$^{401}$your parents..
Remember I said one of the Ten Commandments is honor your father and your mother, right?.
Your parents, in this kind of culture, is the most important people on earth..
You have to honor them..
You have to pay respect to them, right?.
From the first page of the Bible to the last page, nobody can claim the authority higher.
than your human parents, except God himself..
Only God can step in and say, "I have a higher authority..
You honor me above all others, even your parents, even your grandparents.".
So as you can see, what Jesus is claiming is actually the authority of divinity..
He is not a normal guy..
He has power and authority and honor much higher, infinitely higher than any human authority,.
including your parents..
Remember I said I quote this line to you before?.
Listen up..
This line says, "Either Christ is Lord of all, or he is not Lord at all.".
Right?.
A.B. Simpson, founder of Alliance Church..
Christ is all of all, otherwise he is not Lord at all..
So my point is, there are two basic points..
Let me go back to this slide..
This passage is the priority of following Jesus..
Following Jesus, honoring him, worshipping him and serving him is above all..
Above your own emotional stuff you have to deal with, above your parental authorities,.
above everything else..
And Jesus is the highest because he is God himself..
He is the Lord of everything..
He is over everything of your life..
It's the priority of following Jesus..
And secondly, it's the timing of following Jesus..
I gave you the chart..
I gave you the timeline of a typical Jewish burial in Jesus' day..
And I already explained how it is for this and other disciples to approach Jesus..
This has to be after the first burial and before the second burial, the austerity..
So this and other disciples is asking Jesus for a different timing of following him..
He said, "Yes, I want to follow you, but let me delay for about another year and then.
I will decide to follow you.".
And Jesus said what?.
"Let the dead bury the dead.".
What does it mean?.

$^{441}$With that tomb picture in your mind, you understand what Jesus said..
Because inside a family tomb there were so many dead people, right?.
In the past, your grandparents, the great-grandparents, the bones were there..
The dead were with the dead in the tomb..
And Jesus was saying, "Hey, you need to follow me..
You have to follow me now.".
And also, if you understand the story flow, the story plot of the Bible, the Gospel of.
Matthew specifically, we read chapter 8..
Chapter 8, Jesus already started his prophet ministry and he was on his way to Jerusalem.
and die on the cross, crucify, die, be buried, and rise again..
How much time does he have?.
How much time does Jesus have between the disciple approaching him to the time of Jerusalem,.
his death?.
Less than a year..
As you can tell, this another disciple did have the intention to follow Jesus, but he.
did not understand the time is running out..
If you don't follow Jesus now, you will never be able to follow him..
You see what I'm saying?.
Because he will be gone..
If you really need to follow this rabbi, if you really need to follow Christ, you've got.
to do it now because according to the story of the Gospel of Matthew, Jesus had less than.
a year on earth..
There's no time for you to wait for another second burial to be done, anything to be done,.
any stuff that you have to do..
Just a reminder for us, if we follow Jesus, we better do it now because he claims the.
highest priority..
And also, your time is running out..
You have to follow him..
Following Jesus is the biggest deal..
There are many matters in our lives are big deals, but there is no bigger deal than following.
Christ..
There's no one greater than Jesus himself in terms of authority..
So Jesus said, "Let the dead bury the dead.".
He's actually poking fun of the complicated burial ritual the Jewish had..
It's too complicated to follow..
It's actually challenging us to lay aside anything that is stopping us from following.
Jesus today, at this moment, not tomorrow, not next month, not next year, because who.
knows, right?.
In our past two weeks, I was with the 27 people on the trip and I think by the end of first.
week, some of the members are from Hong Kong, they received really shocking news because.

$^{481}$one of their relatives suddenly had a stroke and she's a young mother, 30-some years old,.
had a stroke on the street, didn't have any sign, and she had a stroke, a massive stroke,.
and she was immediately hospitalized in ICU, 30-some years old, and we were praying for.
her on our trip almost every single day, all of us, praying for her..
And I think she's still alive in ICU, but in a very critical condition..
What I want to tell you is that if you want to follow Jesus, this is the time..
Not wait until something has been fixed, right?.
Not wait, like, "Okay, let the church be better and then I'll follow Jesus.".
Or, "Let me take care of my relational issues and my work or my career, and even I won't.
follow Jesus until I retire.".
If that is so, I don't think this Jesus is what I want..
I don't think this is the Lord of all..
Jesus is not the God worthy of my following..
If I have to wait and He will become my footnote of my life..
No, Jesus is the content of your life, not the footnote..
We are the footnotes..
So I want to remind you, following Jesus is now, not tomorrow, not next year..
Let's all pray together..
Pray that we will return and repent from our delay, delaying our decision in ways that.
you can think of..
You know, sometimes we have reservation of following Jesus because this or that, no matter.
what is this or what is that, Jesus has the highest authority and following Him is now..
Father, we come before you, we realize that your Son, Jesus Christ, is the Lord of all..
And I pray that at this moment, all the brothers and sisters, Christ, all the friends, we do.
have some stuff we have to deal with..
And perhaps some of the issues in our lives, our career, our work, our study, almost anything.
we can think of can be the barrier, the barrier of following you right now..
And I pray that you will listen to us, listen to our prayers..
Help us see that you are the God worthy of our praise, worthy of our following..
And I pray that at this moment, your Spirit will cleanse us, lay aside all the issues,.
worries and reluctance that will come and follow you right now..
In Christ's name we pray..
Amen..
(dramatic music).
\newpage



\section{羅馬書 12:1-2}
\label{sec:TVT11XSJIWY}
\textbf{生命的祭 (羅馬書 12\_1-2) - 譚文倩莊傳道}
\newline
\newline
連結: \href{https://youtube.com/watch?v=TVT11XSJIWY}{\texttt{https://youtube.com/watch?v=TVT11XSJIWY}} ~~~~ 語音日期: 2023-06-13
\newline
\newline
\hyperref[sec:7Sx7wB7NbaI]{\small{< < < PREV SERMON < < <}}
~
\hyperref[sec:index_chronic]{\small{[返順時目]}}
~
\hyperref[sec:index_scriptual]{\small{[返順卷目]}}
~
\hyperref[sec:DoxRaGGvrWQ]{\small{> > > NEXT SERMON > > >}}
\newline
\newline
羅馬書 12:1-2
\newline
\begin{longtable}{cl}
\hline
\hline
章節 & 經文 (和合本修訂版)\\
\hline
12:1 & \begin{tabularx}{0.7\textwidth}{X} 所以，弟兄們，我以神的慈悲勸你們，將身體獻上當作活祭，是聖潔的，是神所喜悅的，你們如此事奉乃是理所當然的。 \end{tabularx} \\ \\ \relax
12:2 & \begin{tabularx}{0.7\textwidth}{X} 不要效法這個世界，只要心意更新而變化，叫你們察驗何為神的善良、純全、可喜悅的旨意。 \end{tabularx} \\ \\
[1ex]
\hline
\hline
\end{longtable}
$^{1}$開始之前 我想問問大家.
你們有沒有人每個星期為我們.
星期日服侍的弟兄姊妹祈禱.
Thank you.
特別是為了港元 有沒有?.
如果你有 我很想多謝你.
因為每次講道的時候 我都很戰戰兢兢.
莫說我預備了很多東西給你們.
我預備了幾個月的經文.
正是上個星期三一神學院校長.
包惠君牧師所講的經文.
求主憐憫.
我.
Seriously 其實很多謝包牧師.
幫我打了底.
希望今天我們再講羅馬書.
十二章一至二節的時候.
大家對當中的訊息能夠有更深入的反思.
重要的經文講兩次.
這個暑假我會分開三堂.
跟大家講羅馬書十二章當中的訊息.
第一堂就是今天.
就是羅馬書十二章一至二節.
生命的一章.
羅馬書十二章一至二節.
生命的祭.
到七月頭就會講生命的慈.
講神的恩賜三至八節.
第三堂就在七月底.
九至二十一節講到生命的智.
講智力去愛.
羅馬書十二章一至二節.
可能你們當中有些人跟我一起長大.
是我大學生團契的時候的一個團訓.
現在的團契已經不流行唱團歌背團訓.
很明顯這是我的自殺式.
順便揭露年齡分享.
想當年每個星期三的團契.
我們都會站起來去背團訓.
唱團歌.

$^{41}$我想起每一次背的時候.
我們大家都是很有使命感的.
我自己在想.
我每個星期背的時候.
重新去立志.
告訴神我願意將我的生命放上.
願意獻上我自己給神.
分享之前我們先做個祈禱.
出口主我們的神.
恭敬的將以下時間放上.
你恩高諧至的口.
讓你自己的話語.
藉著人的口帶著能力被釋放出去.
聖靈在人的心裡面工作.
用心去聆聽.
你要跟我們說的話.
這樣的祈禱.
奉耶穌基督的聖名祈求.
《羅馬書》在聖經裡面是一本很重要的書.
可能你聽我這樣說的時候.
可能在想.
誰不知道媽媽是女人嗎.
聖經裡面有哪本書是不重要的呢.
我想講重要的意思.
是《羅馬書》對一些影響世代的神學家.
是有些很震撼性的影響.
神學家 哲學家奧古斯丁.
他在386年的時候.
因為讀《羅馬書》13章.
13至14節而悔改信主.
後來他成為一個很出名的講道者.
我們讀神學院的人.
沒有一個學生是不會讀到他的著作.
他的著作是很舉足輕重的.
神學家 宗教改革家.
馬丁路德.
在1515年11月教授《羅馬書》的時候.
指出《羅馬書》的重點.
不是講一個公義的神要去審判不義的人.
而是神要因為人的信心去稱人為義.

$^{81}$這個教義也是去改革運動的一個動力.
在《羅馬書》12章1至2節之前.
我很想和大家講一講.
隨便講一講《羅馬書》1至11節.
大約是講了些什麼.
總括來說是在講基督教的基本教義.
也在講無論是猶太人.
在當時地位似乎高一點.
巴閉一點.
或者是外邦人地位低一點的人.
都是因為信基督而被稱義.
而不是因為遵守一些律法和教義.
來到《羅馬書》12章的結構上.
12章的開始是將1至11章所講的.
由一些教義轉向我們怎樣去應用.
怎樣去應用解釋福音的真義轉向怨念.
教導信徒在每日行事為人上.
我們怎樣實踐恩信的道理.
就像我剛才所講的.
今天的講題叫《生命的祭》.
我想先邀請大家和我一起讀這節經文.
如果你想當年和我團契的時候.
歡迎你背著.
1 2 3.
所以弟兄們.
我以神的慈悲勸你們.
將身體獻上.
當作活祭.
.
2008年開始.
我們一家隨從我們以前的教會.
一起進入中國大陸傳福音.
每個暑假.
差不多廁止之後.
我們每個暑假都會進入中國.
但自從2019年之後.
我們就不可以再進入.
但我心裡經常記掛著.
在中國大陸的弟兄姊妹.
當地的教會.

$^{121}$很想有機會再去探望他們.
很想問問他們你們最近怎樣.
還可不可以聚會.
有沒有堅守信仰.
同樣12章的一開始.
保羅在這裡.
.
出於神的慈悲.
去慰問弟兄姊妹.
出於他對弟兄姊妹的愛.
勸勉信徒.
要怎樣活出信徒的生命.
.
我為主作酬道的勸你們.
既然蒙召行事為人.
就要與你們所夢的呼召相請.
聖經裡面的教導.
從來都不是只是叫人閱讀.
或者讓人只知道一些教義.
教導的目的是要我們實踐.
看聖經是要我們更加認識神.
將身體獻上當作活祭.
上個星期巴木斯都有少許提過.
舊約的獻祭.
普遍是講被殺的生畜.
被揀選.
預備.
被殺.
奉獻.
一個取悅上帝的儀式.
而當中有很多獻祭的種類.
凡祭.
數祭.
平安祭.
贖罪祭.
贖遣祭.
有很多預備.
很多規矩.
很多步驟.
很多禮儀.

$^{161}$缺一不可.
保羅在十二章一至二節裡講到活祭.
與上面講的那些很不同.
但在講活祭之前.
我很想先跟大家講一講身體.
上個星期都有講到.
身體這兩個字不是講一個身體.
是眾數來的.
Multiple.
每個人的心思意念意向.
所以不只是講那個Physical Body.
不只是講身體這麼簡單.
肉體這麼簡單.
根據原文的解釋.
是很多個身體.
一個信仰的群體.
將身體獻上.
成為一個獻祭擺上.
我在預備的時候.
我在想什麼時候我們會這樣.
我在想信主的時候就是這樣.
正正我們是全人擺上的一個動作.
我們決志信耶穌.
承認我自己的過犯.
降伏在神的面前.
宣告我的生命不再屬於我自己.
我不再是我自己的主人.
決意認定耶穌作為我的個人救主.
從今以後活著的不再是我.
乃是基督在我裡面活著.
一個回應福音的動作.
羅馬書8章11節.
然而使耶穌從死人中復活的.
神的靈若住在你們裡面.
那使基督從死人中復活的.
也必藉著住在你們裡面的聖靈.
使你們必死的身體又活過來.
一個會死的身體.
因為耶穌而活過來.
將身體獻上當作活祭.

$^{201}$羅馬書12章.
這裡說的是甚麼.
人之所以可以成為活祭.
最重要的部份.
不是我們多麼愛神.
願意擺上成為活祭.
因為願意擺上的前設.
是耶穌自己首先擺上.
不是我們有多麼愛主.
多麼屬靈.
是一個多麼好的基督徒.
這不是重點.
重點是耶穌首先擺上.
成為最完美的獻祭.
成就救恩.
主耶穌的寶血清洗我們的罪.
拯救我們.
讓我們重新成為聖潔.
被稱為義.
是耶穌先成為祭身擺上.
人才可以成為活祭.
羅馬書10章10節.
我們憑著這旨意.
藉著耶穌基督僅此一次獻上他的身體.
就得以成聖.
想深一層.
想當年我自己在背誦經文的時候.
只是在想自己有多有使命感.
決意要做一個好的基督徒.
擺上自己.
那個想法是由我自己出發的.
自我感覺良好.
但其實想當年我根本不明白.
重點並不是在我身上.
重點是其實神有多麼愛我.
差派他的獨子耶穌擺上.
是奉乃是理所當然的.
正正是因為人明白耶穌的擺上.
跟隨主的人.
對福音最理所當然的回應.

$^{241}$好像上星期包牧師提到.
對神感恩.
最合乎情理的回應.
不要效法這個世界.
只要心意根深而變化.
這個世界當今的時代.
有很多自己的思想和生活模式.
是我們這群蒙恩得救的信徒.
必須要辨別清楚.
所以不要效法這個世界的意思是.
叫我們小心.
基督徒你要小心.
我們不要讓這個世代去塑造我們的生命.
耶穌負上一個很沈重的代價去拯救我們.
最近在一個飯局裡面.
有一個朋友提到.
因為他們兩個女兒的關係.
他經常會看迪士尼頻道.
近期他發現這些節目裡面.
越來越多不傳統的家庭出現.
傳統的意思是一個爸爸一個媽媽.
和兒女所組成的家庭.
他說這些節目裡面的家庭.
或者他們家庭裡面的互動.
和一般的家庭好像沒有什麼分別.
只不過不一定有一個爸爸一個媽媽.
他說他發現這些兒童節目.
好像在前沿物化地將不正常的關係正常化.
幾年前有一單新聞.
在加拿大有一個46歲的爸爸.
他有七個小朋友.
他說發現他其實是一個女孩子.
他開始的生活模式就是用一個女孩子的身份去生活.
他太太受不了.
結果他們分開了.
他就離開他太太.
離開他七個小朋友.
離開他家之後.
這個六尺四寸高的男人.
穿裙子.

$^{281}$戴頭箍.
然後在網上為自己賣廣告.
希望有家庭能夠收養他.
他自稱是一個八歲的女孩子.
你真的猜不到.
真的有人收養他.
收養了他之後.
他被一對老夫婦收養了.
自此之後他和這對老夫婦的孫子孫女.
玩遊顏色煮飯仔.
有一天其中一個孫女.
跟這個46歲自稱八歲的小女孩說.
她說我想做姐姐.
我不想做妹妹.
因為這個女孩子七歲.
她說你可不可以再小兩年.
不要做八歲.
做六歲吧.
這個男人說好啊.
沒問題.
他就做了一個六歲的小女孩.
今時今日他仍然和這對老夫婦的家人一起住.
最後一個例子.
我想起有一次我去接我女兒放學.
看到其他小朋友在等家長接.
我就嘗試要友善.
跟旁邊的女孩說.
你媽媽來接你嗎.
她看著我沒有說話.
後來我回到家.
我女兒才告訴我.
這個女孩沒有媽媽.
她有兩個父親.
我便想起從此之後我不可以以為每個小朋友的家裡都有一個父母.
因為那個組合可能是父父 母母 父母 母母 母母.
或者只有爸爸 只有媽媽.
或者沒有爸爸也沒有媽媽.
這是我們現今世代家庭關係的狀況.
將來還會出現多少個組合呢.
其實我覺得有時候也估計不了.

$^{321}$我們下一代的家庭觀念定位在哪裡.
我們現在活在一個什麼樣的世代.
作為基督徒的你 你會怎樣去回應.
作為基督徒父母的你.
當你的小朋友來到你面前去問你一些這樣的問題.
你又怎樣去回應.
若果今時今日你仍然只是專注在你的小朋友.
如何催谷他們去讀好書 進入一間好的學校.
甚至是一間名校 出來工作 賺多少錢.
有人跟我說過.
一年班最好讀二年班甚至是三年班的課程.
琴棋書畫.
不要吃壞東西 一定要吃有機食品.
我想在這裡跟你說.
想你再想想.
我不是在這裡嚇你 因為我自己也是一個媽媽.
有時候我的女兒來跟我聊天.
跟我說一些我聽都沒聽過的議題的時候.
我都不懂得怎樣去回應.
我心裡面在祈禱 主啊求你幫助我.
讓我有智慧 起碼不要亂說話.
大部分的時間我自己可以在這裡聆聽.
嘗試瞭解 知道他在外面在聽什麼 學什麼.
不知道你們最近有沒有上我們UV的YouTube頻道.
我在香港的時候特地和朱麗全導演做了一個訪問.
沒看的話很鼓勵你上網再去看.
不是我做訪問做得多好.
你看的時候就知道我做得多不好.
很明顯這個不是我的強項.
朱麗全導演在訪問里給大家很多提醒.
其中一段說到基督徒是People of the Book.
我們是屬於聖經的人.
他有一個心願 教導平信徒去讀希伯來文.
將神學普及化.
讓基督徒有某個程度的裝備.
應付未來的挑戰.
他盼望能夠幫助信徒.
自己可以深入研讀聖經.
甚至可以慰養自己身邊的人.
提摩太后書三章十六節.

$^{361}$聖經都是上帝所默示的於教訓.
督責使人歸正.
教導人學而都是有益的.
教訓是要帶領你走應該走的路.
督責是要告訴你走歪了.
使人歸正是要帶你走一條對的路.
教導人學而是要幫助你不要再走他一條錯的路.
基督徒的信仰根基是神的話語.
有父母 即不止一個父母.
經常會來問我.
作為一個基督徒的父母.
我們如何去幫助 去回應 去教導.
我們今時今日小朋友所面對的種種事情.
林林種種我們自己都不清楚的信息.
現今做父母到底我們應該如何教他們.
親愛的弟兄姊妹.
面對將來的不明朗.
時代的挑戰.
你要記住神的話語.
是我們最大的幫助.
你看聖經的時候.
可能你有時覺得我都不知道看些什麼.
或者有些悶.
或者好像看完我沒懂些什麼.
但我很想告訴你.
日子有功.
儲下儲下.
你天天看聖經.
那些就是你的彈藥.
那些就是你打仗的彈藥.
唯有親近神.
藉著聖靈.
幫助人心意根深而變化.
幫助我們去辨別這個世代.
靠著聖靈重新規劃這個世界墮落的思想.
在動蕩不安的世代裡面.
我們仍然可以站立得穩.
我想起很多年前看了一個劇集.
其中一幕我很深刻.
講到有個珠寶商.

$^{401}$他的朋友就問他.
其實你如何去學分別鑽石的真假.
他說他小時候.
他爺爺會給他鑽石玩.
過了很多年之後.
他爺爺就鑽了假的鑽石進去.
他說一看我就看到那些是假的.
弟兄姊妹.
你掌握真理嗎.
純粹屬於一年.
兩年.
三年.
五年.
甚至十年二十年.
我們信了很多年的基督徒.
你對真理有多掌握.
你在你的生活裡面.
你每天的生活裡面.
你花多少時間去看聖經.
去閱讀神的話語.
羅馬書十二章一至二節.
所以弟兄們.
我以神的慈悲勸你們.
將身體獻上.
當作活祭.
是聖潔的.
是神所喜悅的.
你們如此事奉.
乃是理所當然的.
不要效法這個世界.
只要心意根深而變化.
叫你們察驗何為神的善良.
順存可喜悅的旨意.
生命的祭.
耶穌已經首先擺上.
一次過成為最完美的祭.
你們怎麼回應.
我們教會的弟兄姊妹.
作為一個信仰的群體.
我們怎麼一起去回應.

$^{441}$我很想問候你們.
你們在做什麼.
你們有沒有信耶穌.
信到滑了牙.
上次看聖經是什麼時候.
神的話語在你們的生命里.
有沒有份.
生命的祭.
重生得救的生命.
願意活出一個.
與舊恩相請的生活.
全人擺上.
回應我們一點功勞都沒有.
而得到的舊恩.
定義如此事奉.
正正是因為.
這是對我們領受福音的.
每一個基督徒.
最理所當然.
感恩合乎情理的回應.
求主幫助.
面對將來的不明朗.
時代的挑戰.
我們每個人.
個人所面對的掙扎.
願意在神的面前.
承認我們無能為力.
我們需要神.
基督徒要抓緊神.
我們要努力讀神的話語.
因為唯有靠著神的話語.
幫助我們.
有辨別示我的能力.
靠著聖靈重新規劃.
這個世代的墮落思想.
在動蕩不安的世代.
我們每一個基督徒.
能夠站立得穩.
中國人有句話很好.
台上一分鐘.

$^{481}$台下十年功.
求主幫助我們.
不要做一個.
沒有餡的基督徒.
我們一起低頭祈禱.
主.
我們每一個人.
在禰面前的光景是怎樣.
禰是最清楚的.
主求禰繼續引領我們的生活.
禰幫助我們.
不要忘記禰的恩惠.
叫我們回應禰的愛.
全人擺上.
當作活祭.
求聖靈提醒.
幫助我們緊緊緊.
在禰的話語上.
讓我們在動蕩不安的世代.
裡面.
站立得穩.
在這裡我也很想特別.
為一些正在掙扎.
軟弱的弟兄姊妹禱告.
掙扎的找到出路.
軟弱的得著主的同行同在.
一些生活得不開心.
有很多事務纏繞.
主是禰陪伴.
禰給智慧.
傷痛憂傷的.
禰抱緊.
禰安慰.
主求加添力量.
生活越難過.
越混亂.
越壓迫.
追求從禰話語.
支取力量的心更強.
叫我們來到禰的台前.

$^{521}$可以成為一個土地喜悅的活祭.
你也是一個聽禱告的神.
心存感恩.
不配的禱告.
奉主耶穌基督的聖命祈求.
Amen.
多謝Joyce給我們提醒.
\newpage



\section{}
\label{sec:DoxRaGGvrWQ}
\textbf{比對研究 16: 巴力宗教 9 - 列王記上 18 - 陳崇基牧師}
\newline
\newline
連結: \href{https://youtube.com/watch?v=DoxRaGGvrWQ}{\texttt{https://youtube.com/watch?v=DoxRaGGvrWQ}} ~~~~ 語音日期: 2023-06-17
\newline
\newline
\hyperref[sec:TVT11XSJIWY]{\small{< < < PREV SERMON < < <}}
~
\hyperref[sec:index_chronic]{\small{[返順時目]}}
~
\hyperref[sec:index_scriptual]{\small{[返順卷目]}}
~
\hyperref[sec:2hbMIJp1DE8]{\small{> > > NEXT SERMON > > >}}
\newline
\newline
$^{1}$因為我們以往試過看17章.
我們來到18章 所謂的以利亞故事17至19章.
主要有三章聖經.
我們看完17章 現在看18章.
我們一邊看一邊想一想.
那些經文背後會不會有一些背景.
我們可以在巴力的文獻或資料.
幫我們看得更加深入一點.
17章的時候我們已經看了三個故事.
其實三個故事所講的主題都是巴力神話.
或者巴力的宗教所宣揚的東西.
死而復活 有油 土產 有水 有食物等等.
這些全部都是巴力所推廣的東西.
你拜巴力就會有這些東西.
寫的時候這個故事去和巴力做一些爭辯.
就是說真正的供應者 真正的教運者.
死而復活的神是耶和華自己.
我們來到列王記第18章.
18章 開頭有一段這樣的說話 這樣的記載.
我們一邊看的時候想一想.
我們曾經看過一些巴力的經文.
或者巴力的資料 可能跟這個有關.
他說過了許多日子 到了第三年.
耶和華臨到以利亞 他說你去.
讓阿哈看見你 我要降雨在地面了.
耶和華就去 要讓阿哈見到他.
那時撒瑪利亞的饑荒非常嚴重.
阿哈召了他的管家 俄巴底.
俄巴底是上敬畏耶和華的.
耶洗別殺 耶和華先知的時候.
俄巴底就將一百個先知藏著.
每五十個藏在洞裡面 拿餅和水供應他們.
阿哈就對他的管家俄巴底說.
我們要走遍這地道 一切的水泉旁或者溪旁.
或者可以找到青草 救活我們的馬和女子.
免得喪失一切的生畜.
於是兩人就分頭行事 巡查.
阿哈獨自走一路 俄巴底就獨自走另外一路.
這裡發生什麼事呢 幾年之後.
結果這個旱災 因為沒有雨水的緣故 失修.

$^{41}$就造成了饑荒.
所以在這個時候 饑荒非常嚴重.
嚴重到一個地步 連這個什麼 撒瑪利亞.
撒瑪利亞是什麼地方 就是北角的京城.
撒瑪利亞 加上山在上面.
所以來到北角的中心地帶.
應該是最富強的地方 最多土產.
最多米糧的地方 結果都令京城有饑荒.
非常嚴重 於是他們做了什麼呢.
交代了這個俄巴底是什麼人.
說要洗別殺這些耶和華先知的時候.
他就很忠心地去養活他們.
又是一個養活的故事.
有水有餅 好像第一個故事 十七章.
以利亞就是耶和華養活先知.
這裡就是耶和華籍住俄巴底去養活先知.
俄巴底就聽阿哈說 說什麼呢.
整串的說話就是我們要去找東西.
去找水泉 找溪蚌 找青草.
為了救我們的牲畜和馬等等.
所以因為旱災 因為失收 為了饑荒的緣故.
所以就需要分頭去查 找找到底有沒有水等等.
如果你記得我們曾經提過巴力的神話.
在巴力的循環裡面.
沒有水的時候 發生什麼事.
就是有人來到黃河 或者是獅人.
去做一些愛歌 希望能夠見到巴力的雨水.
這個到處去找巴力的水 找巴力降的雨.
在我們曾經看過的epic of crota.
那裡說什麼呢.
他說王有病 因為旱災的緣故等等.
令到整個國家沒有水.
因為這個緣故 故事在天庭裡.
最至高神L 或者其他的神.
發出一個命令 去找食物和找水.
所以你看到這裡說的主題都是一樣.
這裡說到旱要去巡查 測看大地和天空.
測看地的懸處 找找有沒有墨紙.
找找地上有沒有巴力降的雨水.
到田野裡 最高者的雨 巴力的雨.

$^{81}$一個的lament 一個的哀歌.
巴力的雨水對地有多好.
因為有雨水就有田地 有田地就有出產.
所以最高者的雨水是何等好呢.
實在是好 因為耕地的墨紙.
泥溝裡的墨紙 田間裡的小灌墨.
能夠發芽 能夠有墨紙 能夠有結果的結水.
這個是因為巴力給的水 巴力降雨.
他們就舉手 農夫也舉手 向天去呼求.
大菌的僕人 巴力其中一個的名號.
除了是教運者等等.
就是大菌的僕人 大菌是什麼呢.
達加在史詩紀念出現過.
大菌神 Samson 參孫的故事.
有一個神明是菲利斯人所拜的其中一個神明.
就是大菌 大菌就是菊類的神.
就是米的神 大小麥的神.
所以他們就向他呼求.
因為什麼呢 糧都吃光了 酒也用光了.
桶裡的油也全部用光了.
你看看這段話 再看剛才.
我們之前17章都有提過.
整個故事其實是圍繞著這樣的主題來寫的.
整個以利亞的故事.
就是表達清楚 哪一個才是有能力來降雨.
結果不是巴力 巴力是廢的.
耶和華才是供應者.
因為他所有的先知都被耶和華保守著.
藉著俄巴底 藉著烏鴉來養活他們.
所以誰才是養活者 降雨者.
給食物和水 就是耶和華.
第七節 兩個人巡查的時候.
就像巴力的神話一樣 現在沒有水 沒東西吃 沒雨.
所以就到處去找東西.
所以這裡有一個故事就是到處去找東西.
誰知道找到什麼呢 就看到以利亞.
看到真正可以神藉著他來降雨 給食物的.
那個就是真正的先知.
他說看見以利亞就遇見他.
俄巴底認出他來 就面服於地.

$^{121}$說我是我主 以利亞對他說 我是你去.
告訴你的主人 看啊 以利亞在這裡.
看這段經文的時候不妨留意一下有什麼keyword.
有些很重要的字常常出現.
俄巴底就說 僕人犯你什麼罪呢.
你竟把我交在阿哈的手裡 使他殺我.
我知永生的耶和華 你的上帝起誓.
哪一邦哪一國 我主都派人去找你.
若他們說不在這裡 他就叫哪邦哪國的人起誓說.
實在找不到你.
現在你說 你去告訴你的主人說.
看啊 以利亞在這裡.
恐怕一離開你 耶和華的靈就將你提去了.
到我不知道的地方去.
這樣我去告訴阿哈的時候.
阿哈就找不到你 然後他一定會殺我.
但是你要知道 僕人是自幼去敬畏耶和華的.
好 那我們繼續.
夜洗別殺耶和華先知的時候.
俄巴底繼續說.
我把耶和華一百個先知藏著.
每五十個藏在一個洞裡 拿餅拿水來供養他們.
這個也是其中一個主題.
難道沒有人把我做的事告訴我主嗎.
現在你說 你去告訴你的主人.
看啊 以利亞在這裡 祂一定會殺我.
以利亞說 我指著所事奉永生的萬軍之耶和華起誓.
我今天要讓阿哈見到我.
於是俄巴底就去應見阿哈 告訴他這事.
阿哈就去見以利亞.
有什麼常常出現的字呢.
當然主題一直扣著的就是水和食物的問題.
如果留意的話 聖經很長 其實也很沉悶.
一直有些對話.
在對話裡你講我 我講你等等.
有幾個字很明顯是常常出現的.
就是主字.
我主 你的主.
總共出現了六次.
在單單這個故事裡.

$^{161}$俄巴底這個故事.
並且在這裡這麼多張的聯王記.
17至19張.
沒有用過我主這個字.
只是在這個短的故事裡出現了六次.
這裡highlight了給你.
第七節.
你是我主阿多尼 以利亞媽.
然後以利亞就跟他說.
告訴你主人阿多尼卡 你的主人.
說什麼什麼.
第十節.
我主阿多尼都派人去找你了.
你去告訴你的主人阿多尼卡說什麼.
十三節.
難道沒有人把我做的事告訴我主阿多尼卡嗎.
十四節.
你去告訴你主人阿多尼卡說.
這個adorn這個字.
就是主人的意思.
主人的意思.
作者就將這個字.
adorn這個字.
後面加一個字尾.
希伯來文加字尾通常是加一個屬系的意思.
我的主.
adorn.
adorn卡.
卡.
你的.
單數你的.
然後adorn我的.
my lord.
adorn卡.
your lord.
adorn我主.
阿多尼卡.
你的主.
為什麼要這麼沉悶呢.
一直說來說去.

$^{201}$就是說這個.
想帶出什麼呢.
似乎這裡說出主權的問題.
想帶出一個主題.
就是之後的故事.
想交代的.
就是那個Lawsuit的問題.
誰是你的主.
誰是你的主人.
誰是你所真正敬拜的.
為什麼會這樣說呢.
因為.
去到下一個故事.
就是這個加密山的故事.
加密山的故事.
就是在說誰是真正的主人.
以色列人到底是不是還在騎牆.
左右擺定.
不知道拜耶和華好還是拜巴黎好.
一左一右.
不知道怎麼決定.
在這個俄巴底的故事裡面.
就是先來預備.
那個經文.
那個主題.
將它帶出來.
帶出什麼呢.
就是誰才是主.
所以似乎這裡就是.
以利亞和俄巴底的對話裡面.
其實不是多餘的.
其實就是要俄巴底.
去到走前多一步.
去認誰是他的主.
是不是.
都很奇怪.
因為一開始已經說得很清楚.
俄巴底是一個很好的人.
一個敬畏神的人.
還有保守著供養著.

$^{241}$真正的先知.
他所服侍的.
為他打工的就是阿克王.
是不是.
地上的主就是阿克王.
似乎以利亞就叫他.
去講清楚.
到底俄巴底的主人到底是誰.
真正的主人是誰.
為什麼這樣說.
因為這裡一開始就說.
我主我主這個字的時候.
帶出這個主題.
到最後的時候18章第11節.
他說現在你說.
你去告訴.
你主人說.
看啊 以利亞在這裡.
俄巴底這樣說.
最後這一句.
這句就很有趣.
就說以利亞在這裡.
就是跟阿克說.
跟我的主說.
你去告訴你的主人 以利亞跟他說.
俄巴底就要對他的主人阿克說.
以利亞在這裡.
看啊 以利亞在這裡.
只有兩個字 希伯來文就是.
Hiney Eliyahu.
Hiney就是surprise的意思.
或者behold.
要留意著.
所以只是一個座語詞.
以利亞在這裡.
其實就只有一個字.
或者.
Elijah is here.
只是Eliyahu.
所以其實.

$^{281}$只不過是在說什麼呢.
就是現在你說.
你去告訴你的主人阿克.
看啊.
Eliyahu.
其實是一個字.
Hiney Eliyahu.
Eliyahu是什麼意思呢.
其實就是以利亞這個名字.
以利亞這個名字.
所以很多時候.
英文也好中文翻譯也好.
就是翻譯成以利亞在這裡.
是不是.
Elijah is here.
因為希伯來文可以很短.
而且沒有一些.
很明確的動詞.
或者verb to be.
所以我們就加上去.
Hiney Eliyahu.
就是behold.
Elijah is here.
這個暗示了.
已經在這裡.
所以通常我們翻譯.
就是看啊 以利亞在這裡.
但是我們細想一下.
Elijah這個字.
其實除了以利亞這個名字.
以利亞是.
有什麼意思呢.
就是.
El.
El是什麼意思呢.
就是My God.
我的神.
Yahu就是Yahweh.
耶護華的意思.
所以可以讀成.

$^{321}$Hiney Eliyahu.
就是看啊.
Elijah.
就是說.
My God is Yahweh.
Elijah My God.
Is Yahweh.
所以似乎在挑戰.
這個俄巴底.
就是去他的主面前.
他的主人面前.
就是告訴他.
誰才是他的神.
所以.
如果他照直的說.
告訴阿哈.
看啊.
就是Hiney Eliyahu.
其實他就是在說.
看啊 我的神.
是耶護華.
或者說以利亞在這裡.
或者有兩種不同的意思.
在這裡.
這個就很吻合.
之前我們提過.
就是常常出現什麼字.
就是主的意思.
真正拜的主是誰.
而這裡很有趣.
就是用Elijah.
以利亞這個名字.
想讓俄巴底.
我的神.
是耶護華.
不是巴黎.
所以.
十八章開頭的時候.
一直在說什麼.
在說真正的主.

$^{361}$是誰.
那個主題.
就是一種.
耶護華.
有獨一的主權.
只是敬拜耶護華.
唯一一個.
不同巴黎.
巴黎的宗教.
或者巴黎我們看過的神話.
巴黎的宗教.
是多神的宗教.
就是Polytheism.
多神的宗教.
有什麼特別的地方.
就是很多個主可以放在一起.
沒什麼所謂.
沒有誰是唯獨一個.
最多一個大一點.
巴黎比阿塞拉大.
巴黎比達加大.
巴黎比L大.
或者我們敬拜的.
是拜巴黎多一點.
拜其他東西少一點.
但沒有一種獨一的.
排他性的唯一的.
這些主權.
出現.
你看到十八章第十八節.
其實都是在說這件事.
以利亞說什麼.
看到阿哈的時候.
他說以色列遭殃的不是我.
而是你和你的父家.
因為你們.
你起了耶和華的誡命.
去隨從巴黎.
這個中文.
翻譯不了.

$^{401}$其實是中數.
Abba Aleem.
The Bales.
有個leem.
im在後面.
就是一個陽性的.
中數.
Lim.
就是I am.
至尾.
其實中數就是巴黎.
中巴黎.
因為你離去了耶和華的誡命.
你去隨從中巴黎.
有很多不同的巴黎.
當然這個.
Abba Aleem.
在其他舊的經文.
譬如.
斯基亞等等.
都是說到很多的巴黎.
這個就是.
多神的.
宗教有的特點.
你可以有很多很多的主.
很多很多.
敬拜的對象.
效忠的對象.
現代都是一樣.
你可以有很多很多不同的巴黎.
不要緊的,大家是不會排斥大家.
但是去到.
這個以利亞的故事.
去到中間這一章.
很明顯想帶出.
耶和華是排他性的.
神明.
他不是普通的.
一個眾神的其中一個神.
你拜他的話.

$^{441}$就不可以拜其他東西.
你只能夠拜他.
其他東西就要被廢.
所以有一個exclusive.
主題.
出現.
所以我們處理了.
第十八章開頭的地方.
還沒有到加密山.
下一次我們就會處理加密山.
但這裡很清楚.
這個故事串連起來.
圍繞著.
多神論.
多神的宗教巴黎的主義.
去爭辯出.
一件事.
就是耶和華唯一一個可以敬拜的.
排他的.
exclusive的.
獨一的主宰.
如果你看17-19章.
其實.
列王記上17-19章.
是一個.
風景的單元.
因為他寫的時候.
很明顯是跟著一個.
脈絡.
而且每一個故事是跟著一個大的骨幹.
來進行.
這裡是怎樣進行呢.
如果你這樣列出來的話.
你就更加清楚.
不是一些單獨的故事.
隨便寫出來.
然後炒在一起.
而是很按部就班的來講出某一些東西.
而且是越來越顯出.
耶和華是一個獨一的主.

$^{481}$還在一個巴黎的環境.
裡面.
他是一個獨一的主.
這個脈絡.
就是宣告.
首先開始有一個宣告.
然後有一個旅程.
一個動向.
一個movement.
一個人去哪裡.
或者從哪裡去哪裡.
然後有兩次的相遇.
第一次的相遇.
是一些角色的相遇.
然後也有另外一個相遇.
又見到另外一個角色.
然後就有一個神蹟.
來作結束.
神蹟最後的時候.
就會在一個主角身上.
就會帶出一種逆轉.
一種改變.
或者是一種conversion.
所以你見到.
17章第一節開始.
是發生什麼事.
就是以利亞有一個宣告.
17章第一節.
不妨看看.
就是這裡.
以利亞宣告什麼.
對阿哈說.
就是說.
我所事奉的.
永生的耶和華.
以色列的上帝.
起誓這幾年不禱告.
必不降露水.
也不下雨.
第一個宣告.

$^{521}$宣告了.
然後他就離開了以色列.
去到基列.
去到基列的西部.
然後有第一個相遇.
就是烏鴉.
然後供養著以利亞.
故事沒有停止.
繼續下去.
講另外一個故事.
就是另外一個相遇.
就是叫他去到西頓.
見到一個寡婦.
然後呢.
就跟著有一個神蹟出現.
當然寡婦用不盡油.
也是一個神蹟.
但是一個很重要的.
一個神蹟.
就是這個episode.
這一集.
結束的時候.
就是寡婦的兒子死而復活.
是不是.
然後有一個改變.
就是寡婦.
就認清.
要說.
說話的能力.
先知所說的.
這個就結束了第十七章.
然後第十八章.
又來一個宣告.
就是剛才我們看了.
就是耶和華自己宣告.
宣告什麼呢?.
你去找阿哈和以利亞說.
於是他就來到以色列.
就是另一個movement.
本身他應該在寡婦的地方.

$^{561}$就是在北面的.
西頓這個地方.
你記得第一個movement是怎樣的.
他去了基列.
離開了以色列.
然後上去西頓.
找了一個寡婦.
現在他終於回到.
以色列.
所以另一個movement.
另一個動向就是.
來到以色列.
然後有第一個相遇.
就是和俄巴底.
大家聊天.
在對話裡面我們剛才處理了.
就是我主你主.
我主你主六次.
其實想帶出一個主題.
到底誰才是真正的主呢?.
故事沒有在這裡結束.
到下一個相遇.
就是以利亞.
面對著阿哈.
然後就去到.
加密山上面的所謂神跡.
就是上帝降雷.
顯明上帝能夠打雷.
能夠降雨.
就不是巴力.
這個我們下一次處理.
然後到最後有一個逆轉.
以色列人說耶和華是神.
耶和華是神.
然後阿哈也有多少的逆轉.
然後就結束第十八章.
然後就結束第十九章.
第十九章也是類似這樣的骨幹.
來進展的.
所以我們處理.

$^{601}$七至十九章.
這裡就給我們.
起碼一幅地圖.
知道這三章的經文.
這些故事原來是這樣串連起來的.
想說一件事.
就是去爭辯耶和華是.
獨一的主宰.
不是巴力.
也都不能夠把巴力混合在一起.
你要拜的話.
只能夠拜耶和華.
如果不是.
你就只是拜巴力.
巴力就加耶和華.
所以聖經很清楚的說.
沒有這些所謂的.
syncretism.
就是宗教混合的主義.
沒有我們好像.
華人民間信仰.
耶穌可以放在其他主那裡.
又可以拜.
或者耶和華.
耶穌和金錢又可以一起拜.
只能夠拜.
一個主.
一和真正的神.
謝謝大家收看 再見.
\newpage



\section{1:7}
\label{sec:2hbMIJp1DE8}
\textbf{Wisdom and Fear of God (Proverbs 1\_7) - Rev. Mark Chan}
\newline
\newline
連結: \href{https://youtube.com/watch?v=2hbMIJp1DE8}{\texttt{https://youtube.com/watch?v=2hbMIJp1DE8}} ~~~~ 語音日期: 2023-06-20
\newline
\newline
\hyperref[sec:DoxRaGGvrWQ]{\small{< < < PREV SERMON < < <}}
~
\hyperref[sec:index_chronic]{\small{[返順時目]}}
~
\hyperref[sec:index_scriptual]{\small{[返順卷目]}}
~
\hyperref[sec:F8HOtvMKDlc]{\small{> > > NEXT SERMON > > >}}
\newline
\newline
$^{1}$- Hey, welcome to all of you..
And some of our folks are actually out of town,.
so we have less people here, but that's okay..
We're gonna worship together,.
and I pray that you will open the Bible.
and turn to the book of Proverbs,.
and we're going to think about what it means to be a father..
Today's Father's Day, right?.
Yeah, happy Father's Day to all the fathers here, okay?.
That's me..
(audience laughing).
And then soon to be father, happy Father's Day..
I wanna start with a question..
What do you think of when you see this word, father?.
Father..
What kind of image do you have in mind?.
Do you think of father, fatherhood?.
Do you think of courage, loyalty, faithfulness?.
Responsibilities, loyalty, wisdom, perhaps?.
Do you think of that, those terms?.
You do?.
All right..
But most of the time, I think, sadly,.
we are in a crisis of fatherhood,.
and absentee father, probably you can think of that.
when you see father, or think about father,.
abusive father, detached father,.
or angry father..
I don't know what you think of,.
especially in the Asian context..
Sometimes we think of all sorts of negative images.
related to fatherhood..
According to some of the numbers,.
24 million children in United States,.
they grow up without a biological father..
So that's one in three, right?.
That's a pretty big number, right?.
One out of three children, you know,.
growing up without a biological dad in their home..
That is the crisis of fatherhood,.

$^{41}$and that leads to all sorts of different crises.
about your personal development,.
about our personal identity,.
about how we anchor who we are with our fathers..
So that's a huge crisis,.
and that's why we need to talk about this, okay?.
Today's Father's Day, and I would like to think.
a little bit more about what it means to be a father,.
and hopefully this will not be applicable.
only to fathers, but to all of us..
Think of the fathers in the Bible..
Any negative examples in the Bible?.
Abraham, Isaac, Jacob?.
They're not colorful or beautiful figures most of the time..
They have their own weaknesses and brokenness, right?.
Think of Jacob, right?.
He favored one of his sons, right?.
Joseph, and brings about some sort of family emergency.
or family crisis..
He almost breaks apart the entire family.
because of his favoritism on Joseph,.
destroying the family of God..
Without the grace of God at the end of the book of Genesis,.
the family cannot be safe, right?.
So negative image of father..
And then probably you can think of some positive image.
of a father from the Bible, right?.
Joseph, New Testament, Jesus' earthly father, right?.
Being described as righteous, right?.
Obedient to Torah..
Father Abraham, again, he believed in God.
and God credited him righteousness, right?.
He trusted in God..
So there's some mixed complexity.
with all the fathers in the Bible..
But I think this is good news..
The good news is that even with our own brokenness,.
I'm somebody's father and I struggle with my own witnesses..
I struggle with my, you know, there's all sorts of things..
Brokenness, in order to be a father,.

$^{81}$I try to be a father that is pleasing to God..
And I think this is good news when we see, you know,.
all sorts of different positive fathers.
or negative fathers in the Bible..
And the grace of God is always sufficient..
It's always abundant..
And only if we have the courage,.
only if we have the humility to come before God.
to be shaped by the grace of God..
So today, I hope the message will not be.
just applicable to the fathers, but to all of us..
Now, if you turn to your Bible,.
in the middle part of the Bible, there are five books..
Job, Psalm, Proverbs, Ecclesiastes, and Psalms..
These five books, what we call together,.
is the body of literature..
We call them wisdom literature..
Wisdom literature..
And each of these books has its own way.
to show the life of wisdom..
What does it mean to live wisely?.
And in the context of today,.
father is often teaching the son,.
like in the book of Proverbs, Ecclesiastes,.
Psalm of Psalms, teaching the son and the daughter.
how to behave, how to live a life of wisdom..
So there we go..
We have a huge body of literature.
right in the middle of the Bible,.
showing us what it means to live a life of wisdom..
Will you consider yourself wise?.
How many of you would consider yourself,.
I have some wisdom?.
You know, how many of you?.
All right, you guys are very humble..
You know, this is interesting.
about wisdom literature in the Bible..
You know, wisdom literature, if you read carefully,.
if you think about it carefully,.
it is different from all other different writings.

$^{121}$of the Bible, like the prophetical books..
Some of the writings are written in such a way.
that truth comes from above..
Thus says the Lord, right?.
The prophet said..
Thus says the Lord, this is what Yahweh said..
So you have to obey, right?.
Wisdom literature is a little bit different..
Actually, in terms of direction,.
wisdom literature from below..
It's all, it often gives us some sort of,.
you know, from observation to insight..
So the wise, the wise teacher in the book of Proverbs,.
for example, he moves from observation to truth..
He observed there's a divine order in the world..
Of course, the universe is complex and just,.
you know, finely tuned..
It's amazing..
But the wise teacher moves a little bit forward.
and says not only the universe is designed.
and created by God, and it's amazing by itself,.
he perceives there's an order in the universe..
And that order is not just a physical order,.
but a moral order, a spiritual order..
And that spiritual order, he perceives and he observed..
And then he got some insights and spiritual insights.
and lessons to how to live a godly life..
That is what wisdom is all about..
Wisdom in the biblical sense is practical skill,.
how to live a beautiful life,.
how to live a life, a godly life that is pleasing to God..
That is wisdom in the Bible..
So because you guys are all wise,.
I want to give you some questions, okay?.
So I'll give you a picture,.
and you tell me what kind of insight.
you can get out of it..
All right, you ready?.
All right, let me show you this..
So when you see something like this,.

$^{161}$okay, when you see something like this,.
what kind of insight can you get?.
Where you see this picture?.
Home Depot, right?.
A door hinge, right?.
A door hinge..
What kind of practical wisdom can you get out of it, right?.
It's kind of strange, right?.
Why you show me this picture on the worship?.
And, right?.
What do you see?.
The wise teacher will come and observe something,.
and how the door hinge works,.
and he will understand the order of life,.
you know, spiritual lesson, so-called..
Let me show you..
Proverbs 26, 14..
As the door swing back and forth on his hinges,.
so the lazy guy turns over in the bed, right?.
That's what he's getting at, right?.
From observation to truth, remember?.
From observation to insight..
From observation of how the world works..
And then he gets some spiritual lesson.
that is applicable to all of us..
You gotta work hard, otherwise you will be like a hinge,.
moving left and right, and you are not moving anywhere..
And your life will not be.
contributing to the community..
That's what the wise is getting at, right?.
So from observation to truth, okay?.
You get what it means, right?.
You get what I mean..
So I give you a second picture,.
and then you tell me what kind of insights.
or spiritual lessons you can get, okay?.
Now this second picture I took two weeks ago,.
and I was in Israel, and as you know,.
I was leading about 27 people for two weeks in Israel..
So we were in Bethlehem, and as you know,.

$^{201}$we tour around the country,.
and sometimes we need to do shopping, okay?.
We need to do shopping..
So we were in a shop,.
and it was owned by a Palestinian Christian family,.
and we spent a little bit of time in there,.
and they got some souvenirs, some olive wood,.
craft, all sorts of things..
And don't ask me, I didn't buy anything for you..
But they all bought something, right?.
So by the time I told the group,.
okay, we gotta go, right, we gotta go..
The time's getting late..
You know, we need to move on to the next stop, right?.
So everyone is exiting the shop.
and ready to get onto the bus, right?.
And then right at that moment,.
a street merchant ambush us.
and came with some sort of, you know,.
souvenir items and trying to sell those.
to some of the people who got extra money..
So that delayed us from getting onto the bus..
So right there, I was trying to tell them,.
you know, we gotta go, we gotta go, you know?.
You know, we're running out of time..
And then I took this picture..
(audience laughing).
So they were on the ground,.
they were on the ground with all this scarf.
and, you know, textile products,.
and then bookmarks, you know, fairly cheap..
And I think the merchant, you know,.
got some of our weaknesses..
You know, he's like, oh, I'm gonna sell you,.
like, you know, it was originally \$10 for five,.
and then now it's \$5 for five, okay?.
And then all this, you know, team members,.
as you know, they got delayed, you know?.
We actually, we are frozen in time..
And then they are trying to get all of these things..

$^{241}$It's actually very nice, it's very nice..
It's just not, this is not me..
So, this is the picture I took, okay?.
So, from observation to truth,.
what kind of insight did you get?.
Immediately, I thought of one verse in the Bible.
that talks about this picture, exactly, right?.
It's this one..
2014 Proverbs..
The buyer heckles over the price,.
saying it's worthless, it's worthless..
Then, brag about getting a bargain..
(audience laughs).
So, the wise teacher is making some observations.
about human nature..
He's trying to describe human nature, right?.
What is the human nature that he describes?.
He observes humanity..
We are all beings of contradiction, right?.
When they were looking for the different colors.
of textile products, this is not good,.
this is not good, this is not good, this is not good..
And, by the way, it's not just ladies, men too..
Yeah, this is not good, this is not good,.
this is not worthless..
They try to heckle the price, right?.
And after they got it, they got onto the bus,.
they say, "Oh, this is so good..
"This is so nice.".
Now, what the Book of Proverbs is trying to get at.
is that we tend to be prival..
We tend to be proud..
That's our natural tendency..
Our natural tendency has always not been humility..
You don't have to work on yourself,.
you don't have to work hard in order to be a prideful person..
You are a prideful person, right?.
And that's why our inner self is often contradicting.
ourselves..
So that is what the wise teacher is trying to do..

$^{281}$From observation on earth, and then he gets.
some spiritual lesson to understand who we are.
and understand the way of wisdom..
So there are two big points I wanna show you..
First is biblical wisdom is understanding.
the order of creation..
There is a certain order, spiritual order,.
moral order in this universe..
And the wise teacher, the wisdom literature,.
is trying to get at, is to observe carefully.
and then get some spiritual lessons out of it..
And the second point, which is important,.
is that wisdom begins with obeying the voice of God..
So let me show you the first point..
Understanding the creational order of God..
I'll show you this picture..
What do you see here?.
Stone fallen, right?.
Stone fallen wall..
It's very simple, right?.
And what do you get?.
Again, observation to insight, right?.
2430..
Here the teacher says, "I walked by the field.
"of a lazy person, the vineyard of one with no common sense..
"I saw that it was overgrown with all sorts of nasty weeds.
"and the walls were broken down..
"And then I look and then I reflect upon it.
"and I learn this spiritual lesson..
"A little extra sleep, a little more slumber,.
"a little more folding of hands to rest,.
"then poverty will come..
"It will pound on you like a bandit..
"It will come to you and attack you like an armed robber.".
So that's what he's trying to get at..
If you slumber a little bit more,.
if you're lazy a little bit more,.
you should know that life always works this way..
That is the principle of life..
He passes by a lazy person's vineyard..

$^{321}$You know how many years do you need.
to grow a mature vineyard?.
How many years of hard work?.
What, six years?.
(audience laughing).
20 years would be very mature, very mature, very valuable..
You know, you go to a wine store, you get some good wine..
Those are from a very, very old,.
a very, very mature-grown vineyard..
But what I'm telling you is at least six years,.
24/7, 365 days a year..
You have to work on the soil, you have to work on the wall,.
and you have to retain the soil, the good soil,.
and then you have to fertilize it, you have to water it,.
take care of it, prune it, and grow it..
For six years of sweat and labor..
And then you get a vineyard..
But if the vineyard is in the lazy person's hand,.
the vineyard will be ruined..
You'll find weeds..
You'll find your vineyard being destroyed..
Just naturally being destroyed..
Doesn't mean there's enemies trying to attack it,.
trying to break down the walls..
It's just naturally..
Stone will always fall down..
Stone from the stone wall will always fall down.
and the wall will be broken,.
and you lose all the soil and all your produce..
So the wise teacher is a spiritual teacher..
He's not a natural scientist..
He observes how the world works and he learn a lesson..
If you're a little bit lazy in your spiritual life,.
in your relationship with God,.
if you stop working, there's always forces of chaos.
that will destroy you..
That is the natural order of things..
And the order will naturally give way to chaos..
The wise teacher is saying, in your life,.
the world has always been this way, order to chaos..

$^{361}$Always order to chaos..
Unless you work hard and you work against.
the forces of darkness, forces of chaos..
You put the fallen stone back on the wall..
You input, you work and you labor..
You sweat, repair and maintain your life..
Otherwise you will not enjoy the fruit of your labor..
Spiritual life, or life in general, does not come easy..
That's what wisdom is trying to say to you..
There's always an insistence of emphasis of diligence..
You know, wisdom does not come easily..
It can only be earned..
It can only be earned by your hard work..
Spiritual maturity, character formation..
You want to be like Christ..
You want to be making advances in your spiritual journey..
It doesn't come naturally..
What comes naturally is ruin, chaos..
If you fold your hands, fold your arms,.
naturally you will be destroyed..
And the wise person is saying,.
you need to put sweats and tears into your spiritual life..
Think about it..
I've been a Christian for a long time..
Whenever I got lazy, I don't do anything about it..
I don't work against my own weaknesses, my own old self..
All sorts of sinful nature will come out..
You have to guard your heart..
In fact, the Old Testament and New Testament,.
the entire Bible often advises and command us.
to guard our heart..
And guarding is a military term..
It's a military term..
You stand on a tower to guard against the enemies..
They will come any time to destroy you.
and your life will be ruined..
And you have to guard your own life..
And how do you do that?.
That is wisdom calling..
Calling you to form your character,.

$^{401}$to mold yourself so that you are on the right track.
in your spiritual relationship with God..
So wisdom is understanding the creational order.
and then you learn some spiritual insight out of it..
For those of you who have the blessing.
to be born and to grow up in a Christian environment,.
I want to tell you, you cannot take it for granted.
because the forces of darkness will destroy you..
It will ruin you unless you work really, really hard..
That is Proverbs 24..
So second point is, wisdom is listening and obeying God..
Now don't get me wrong..
Wisdom literature is not about self-help, self-improvement..
It's about drawing us closer to God himself..
So Proverbs 19.27..
It says, "If you stop listening to discipline, my child,.
"you will stray from the words of wisdom,.
"words of knowledge.".
Words of discipline, think about it..
The Hebrew word means,.
actually it can mean chains, you know, boundaries..
Something you can think of,.
but it's not some comforting words..
Words of discipline..
It's gonna be tough to listen to,.
but you have to listen to it..
You have to be obedient..
You have to be open to discipline..
Strayed, it means unable to walk in a straight line..
Have you ever got drunk?.
Yeah, you did, right?.
You had before, right?.
You got drunk and you cannot walk in a straight line..
You just wander, wander, you waste your time and energy,.
and then you probably will fall, right?.
Fall and stumble, right?.
Strayed is unable to walk straight..
It's actually describing people who got drunk..
But what it says is,.
if you stop listening to words of wisdom,.

$^{441}$even tough words, discipline,.
then your life will be in ruin..
Now, we often misunderstand spiritual growth.
in some sort of rosy picture..
We always think, okay,.
how do you grow your relationship with God?.
Okay, how do you grow?.
Well, we're talking about going to a retreat, right?.
To St. Louis, going to a retreat,.
you know, going somewhere else..
And probably there will be times.
we will do some, you know, devotion, right?.
Quiet times..
When you think of quiet times, spiritual reflection,.
what kind of picture you have in mind?.
You probably have this kind of picture in mind, right?.
Ah, so nice, you know, go to a lakeside.
and you just sit there and open up your Bible.
and, you know, you turn to some of the passages.
that may be your favorite passages, right?.
You don't want to turn to any passages.
that will yell at you, right?.
That will discipline you..
You don't turn to that, right?.
Because, you know, we're in the retreat, right?.
We want to feel nice..
So, oftentimes we think of spiritual growth.
or spiritual relationship with God in this sort of picture..
I want to tell you today, this is a wrong picture..
You raise it..
It doesn't happen that way..
Spiritual growth, your spiritual life, your character,.
it doesn't grow this way..
This is a lie..
It doesn't work that way..
Now, think of spiritual growth not in this sort of picture..
Think of spiritual growth in this picture..
This is spiritual growth..
This is character formation..
This is spiritual formation..

$^{481}$I remember years ago, I was in a region college in Vancouver.
and I was taking MDiv, right?.
I was preparing myself for ministry..
And I have so many classmates, other seminarians,.
they were taking the program of a spiritual theology, yeah..
Spiritual theology or spiritual formation..
I never heard of any one of those.
would talk about spiritual growth this way..
I never..
To sweat?.
I mean, to labor?.
I mean, to work hard, right?.
Difficult?.
Being tough?.
I never heard of that..
You know, I came out and I really have a big question..
You know, how do you grow yourself?.
How do you grow in your character, in your spiritual life?.
They always think this way..
Oh my God..
Oh my God, I need to grow..
Okay, I need to go away and to reflect, right?.
On the calm, you know, weather..
It doesn't work that way, guys..
I'm not saying there's no room for your spiritual growth.
to have some quiet times and spiritual reflection..
We all do, right?.
Put out 15 minutes a day, 10 minutes a day..
It's necessary..
But most of the time, it's like this..
The Bible is very clear to all of us, you know?.
And the spiritual teacher is showing us.
that character growth, spiritual formation,.
requires hard work..
There's no shortcut..
If you forget about it, if you fold your arms,.
your life will be in ruin..
And your relationship with God will go down the pit..
Spiritual growth requires sweat and tears.
by a hardworking farmer..

$^{521}$Proverbs 1:7, it says, "Fear of the Lord.
"is the foundation of true wisdom,.
"but fools despise knowledge and discipline.".
My child, my child, listen..
When your father corrects you,.
what you learn from them, from wisdom,.
will crown you with grace, and you will have honor..
Now, how do you earn a life of wisdom and life of honor?.
You've got to listen..
And listening is hard work..
You've got to work against your own natural tendency, right?.
When someone says some words of discipline,.
correction to me, you know, I don't have to work on it..
I will naturally resist, right?.
I will naturally resist..
But you've got to work against that tendency,.
that sinful tendency, in order to get that honor,.
to grow your life and grow closer to God..
So to all the fathers today, or soon to be father,.
I want to challenge you and want to remind you.
to grow in the way of wisdom..
And that's going to be hard work..
There's so much work in there, right?.
And work hard to build a godly character..
Because as father, I think fatherhood.
is the embodiment of spiritual wisdom..
And there's no excuse..
There's no excuse for us to say,.
oh, no, I'm not a good Christian, that's okay..
You know, that sort of excuse will only lead to.
ruin and chaos and disorder..
Do your best and run in the way of wisdom..
Even though it will require you to work hard..
And to all of you, to each and every one of you,.
I want to remind and challenge you to grow spiritually..
And what I meant is not that rosy picture..
No, no, no, no, no, no, erase that..
This one, I want to remind you,.
we need to work hard to build godly character..
To build a community of God, to build a church..

$^{561}$This is the picture..
And that is the picture of honor..
And that is the future of honor..
So let's work on it and walk in the way of wisdom..
Let's pray and recommit ourselves.
before the God of wisdom..
It is his pleasing will, his pleasing desire.
for you and for me to grow in wisdom..
Let's all pray..
And I'll ask the team to come back up..
Father, we come before you and humble ourselves..
And may your spirit humble our hearts,.
knowing that we are sinful.
and we resist the word of wisdom..
And I pray and I beg for your grace..
I pray that you will change us and mold us.
so that we'll be teachable..
Our spirit will be obeying to the voice of you..
Speak the word of wisdom in our lives.
and grow us and mold us in the way of wisdom.
so that we can grow closer and closer to you..
Perhaps we need to get back to reading your word.
and just work hard on understanding your word..
Help us, Father..
Perhaps we need to get on our knees to pray,.
to pray that we will grow in holiness,.
grow in godliness, and to humble ourselves before you.
and acknowledge that you are our God..
(soft piano music).
So I pray that you will continue to help us to grow,.
even though it will take hard work..
We will not be lazy..
We will cling to your mercy.
and to grow spiritually and wisely.
so that your name will be honored in our lives..
In Christ, we pray..
\newpage



\section{箴言 1:7}
\label{sec:F8HOtvMKDlc}
\textbf{智慧與敬畏(箴言 1\_7) - 陳崇基牧師}
\newline
\newline
連結: \href{https://youtube.com/watch?v=F8HOtvMKDlc}{\texttt{https://youtube.com/watch?v=F8HOtvMKDlc}} ~~~~ 語音日期: 2023-06-20
\newline
\newline
\hyperref[sec:2hbMIJp1DE8]{\small{< < < PREV SERMON < < <}}
~
\hyperref[sec:index_chronic]{\small{[返順時目]}}
~
\hyperref[sec:index_scriptual]{\small{[返順卷目]}}
~
\hyperref[sec:K5EeMHs5F9w]{\small{> > > NEXT SERMON > > >}}
\newline
\newline
箴言 1:7
\newline
\begin{longtable}{cl}
\hline
\hline
章節 & 經文 (和合本修訂版)\\
\hline
1:7 & \begin{tabularx}{0.7\textwidth}{X} 敬畏耶和華是知識的開端；愚妄人藐視智慧和訓誨。 \end{tabularx} \\ \\
[1ex]
\hline
\hline
\end{longtable}
$^{1}$有時間交回給Mark Sir.
謝謝你揭曉.
今天是父親節.
不要緊,不用傷心.
父親節快樂.
我特意跟你說的.
特意跟你說.
今天我們想.
父親節的時候.
我們說一些應節的事情.
就是說父親.
我不知道你們想起父親.
這個形象的時候.
想起什麼形容詞.
大家想一想.
在說father的時候.
想起什麼.
可能是想起一些很有責任感的人.
或者是一些很忠心的.
委託於他.
他一定會執行的人.
一家之主.
家庭之主.
希望是這樣吧.
但會不會想起.
父親的時候.
會想起一些負面的形象.
absent father.
缺席的父親.
或者是一些detached的.
人在心不在.
一種這樣的形象.
或者可能想起一些.
可能會發脾氣的人.
其他沒有什麼情緒.
會不會你會這樣想呢.
很多時候想起父親的時候.
就會想起一些比較負面的事情.
我相信這一代.
我們真的要好好思考.

$^{41}$到底作為一個父親.
有什麼重要的角色.
我這樣說.
不只是我們是生物的father.
而是我們整個神的子民.
都需要去想一想.
到底什麼叫做parent.
什麼叫做father.
按照美國的一些調查.
現在你知道有.
在美國的小孩子.
有2400萬的小孩子.
在家裡沒有父親.
就是說三個裡面有一個.
家裡面沒有一個的父親.
這個是一個父親的危機.
父親危機大家可以想像.
可以引申出多少identity的問題.
多少身份的危機.
就是對一個人的成長.
對一個人的發育.
對一個人的發展.
會有多少的危機.
所以我們要想一想.
到底什麼叫做一個父親.
聖經可以給我們一些指引.
所以到父親節.
我想不只是吃喝玩樂這麼簡單.
我希望我們有一些挑戰.
聖經裡面有很多父親.
當然因為聖經是一個.
Patriarchal Society裡面.
所有的媒介.
所以聖經是一個以父權為主的.
一本的寫作.
裡面講到很多的父親.
也有很多失敗的父親.
亞伯拉罕有沒有.
亞伯拉罕 以撒雅各.
三個最重要的所謂開國之父.

$^{81}$都有他們自己失敗的地方.
譬如雅各.
你記得雅各.
創世記的那個雅各.
偏愛其中一個兒子.
然後導致家破人亡.
幾乎不是上帝在恩典.
那個家就散了.
所以有很多失敗父親的形象.
但同時也有一些正面的形象.
又是講亞伯拉罕 以撒雅各.
他們也有一些正面的東西.
或者上帝在恩典裡面.
來拯救了他們.
正因為這個緣故.
我覺得這個是Good News.
雖然我們每個作父親的人.
我們都有破碎的地方.
有很多脆弱的地方.
對不對.
有很多我們不事的地方.
這個是上帝給我們的恩典.
去提醒我們.
其實可以我們去倚靠這位上帝.
我們可以活出那種.
父親應有的勇氣.
和父親應有的謙卑.
父親應有的生命是怎樣呢.
所以今天我想和大家去想一想這方面.
在聖經裡面有五卷書.
剛好在中間.
白詩針全國.
是叫做智慧文學.
智慧文學可能我已經講過很多.
不過我想今天Center.
講一件事.
就是智慧文學和我們所謂作父親的.
會有什麼的點 重點.
我們今天可以留意.
白詩針全國.

$^{121}$很多時候這些智慧文學.
都是一個父親去教一個小孩.
對不對.
Father teaching the son.
或者一個Teacher teaching the student.
所以在一種這樣的背景下.
在一種這樣的脈絡下.
去寫出來的.
譬如.
箴言 傳道書.
甚至啞歌.
都是一個爸爸去教小孩.
啞歌可能更特別.
可能是教女兒.
這種智慧文學其實就是.
想去到講出.
人生命裡面.
他繼續發展 繼續成長的時候.
他需要的是什麼.
所以一個慈父就去教一個小孩.
或者一個教師去教一個學生.
當然在以前的社會裡面.
所有的教育都是在家庭裡面發生.
沒有說要上學的.
因為家就是學習的地方.
所以一個教師或者一個父親的角色.
來講出一些所謂的Wisdom.
一些智慧.
什麼叫做Wisdom.
很簡單.
其實就是一些應用性的.
一些技巧.
怎樣活到一個合神心意的生命.
那些practical skills.
到底是什麼.
所以智慧文學.
不同的部分.
會有不同的角度去看.
智慧文學有一個很有趣的現象.
和其他的書卷.

$^{161}$或者和其他的寫作方法.
是有一點不同的.
他不是直接去告訴你.
智慧文學往往就是.
想去到拋磚引玉.
去告訴你他observe了什麼.
然後在世情裡面.
他有一些觀察.
對人性的一些觀察.
然後在人性觀察裡面.
他會得到一些所謂的insight.
生命裡面.
生活裡面有的insight.
因為這樣的緣故.
我覺得有一個吸引的地方.
見到一些日常生活見到的東西.
然後那個智者.
或者那個父親.
就會想多一點東西.
他想多一點.
到底在上帝裡面.
有什麼值得我們學的東西.
About我們的life.
About我們的spiritual lesson.
我知道大家都很有智慧.
我們想考考大家.
好嗎.
譬如你看到這件事的時候.
你會想起什麼.
你們在哪裡見到的.
Home Depot見到的.
家裡也有.
門教.
門教你領略到什麼.
聲音.
智者都有看到門教.
以前的門教當然不是這樣.
但也有門.
智者在針圓裡面說什麼.
那個懶惰人在床上.

$^{201}$就好像門的門教一樣.
他在說什麼.
他說我們在床上不想起床的時候.
我們轉右轉左.
一直都是原地踏步.
好像一個教一樣.
所以你今早就好像一個教一樣.
轉來轉去.
關掉了alarm.
然後又再snooze過alarm.
這個就是智者想說的.
就是人生裡面應該要怎樣做.
就是需要勤奮去作功.
是吧.
不能走捷徑的.
有很多東西不能走shortcut的.
他就領略了這樣的道理.
你明白我的意思嗎.
由觀察observation.
就去到一個的insight.
所以每次你看到門教的時候.
就想起這句.
針圓十六章十四節.
想起我們不要這麼門教.
ok.
起身去作功.
這樣你的生命才能夠祝福到其他人.
是吧.
因為很多時候智者都在說.
不只是你自己這麼簡單.
在群體裡面.
你能夠祝福到其他人.
所以為什麼要勤奮作功.
不是為了自己的錢包多一點.
或者錢多一點.
而是一個義者.
一個智者.
在一個群體當中.
來生活的時候.
他應該要做什麼.

$^{241}$就是不要做一個門教.
懶惰的人就在床上.
好像軸心一樣.
他是從一個observation.
然後有一個insight.
第二張相片.
我想show給你看的就是.
你和我領略一下有什麼insight.
這張相片就是上兩個星期.
在以色列的時候.
我們去了一個地方.
在帛琉.
當然我們去遊學.
一定要買東西.
我們在那間店買完東西之後.
那裡有很多橄欖木.
那些產品等等.
食物等等.
買完之後我們就出來.
出來我們準備上車.
誰知道沿途就有一些兜售東西的人.
就截擊我們.
然後他就說有什麼便宜的東西等等.
就有一個這樣的情況出現.
他說很便宜.
那些絲巾.
然後有那些書籤.
有些有路等等.
很厲害.
我說走吧走吧.
我們要走了.
誰知道好像拜年一樣.
走都走不了.
因為那個人.
那個商人.
只是露面的.
除了我之外.
那個商人就抓到我們這個心理.
他就說便宜一點.
真的很便宜.

$^{281}$其實那些東西真的很便宜.
如果我的行李夠大.
我也想買兩件.
所以那些人不怪他們.
是應該的.
你觀察到一些東西.
一個觀察.
你領略到什麼.
領略到什麼.
趁便宜就買.
做一個好的消費者.
知者有句這樣的說話.
我看到這樣的情況的時候.
我就馬上想起.
這一句說話.
這一句說話很少我們會看.
但是在智慧的文學當中出現的.
知者說什麼呢.
他說買物的.
不好不好.
他這樣說.
及至買了之後.
他就自誇.
就是說買的時候.
不要不要這個太貴了.
這個不行.
不值這麼多.
但是買了之後.
他上車的時候就說.
哇多漂亮.
你看看.
這個就是我們.
有應該可以觀察到的東西.
大家可能也是這樣.
我覺得這個很有Answer.
雖然不是任何一個景點.
但是我就真的想起.
這句說話.
到底知者想做什麼.
知者是觀人於微.

$^{321}$所謂觀人於微.
他領略到人性是怎麼樣的.
Human nature.
大家如果有經驗.
有智慧.
人生裡面有多少智慧的時候.
有些人我們會形容他.
越人無數.
對不對.
有些人他看到一個人的時候.
很多時候都看對.
就類似這種情況.
這個知者就observe.
那個Human nature.
發現原來人類.
人是一個很多時候會矛盾的人.
我們心裡面有一種矛盾.
他就在說用買物這件事.
買東西購物這件事.
那種的Contradiction.
We are beings of contradiction.
我們很多時候人性就會這樣.
原因在哪裡.
他在描述和觀察人性.
原因在哪裡.
就是人會自掛.
我們不用教我們都會自掛.
對不對.
總是覺得我們自己對.
總是覺得我們自己是對的.
很難去學到謙卑.
對不對.
這個就是他describe humanity的時候.
去到observe人性的時候.
他領略到的所謂的insight.
Wisdom是什麼.
智慧.
就是我們work against.
我們自己那種自然的趨向.
這裡說自然的趨向就是我們會跨.

$^{361}$一定會的.
如果你不慎防你自己生命的方向的話.
你就會一定走向跨那裡.
不是在上帝面前謙卑.
這裡不是給什麼購物的judgment.
而是那個智者有所領悟.
這個就是我所說的智慧文學要做什麼.
就是要領略和觀察人性.
然後看出那個道理.
所以有兩點我想說的.
第一點就是什麼叫Wisdom.
Wisdom就是那個智者.
他會看到世情.
世間宇宙.
是一個order在那裡.
說的是creation.
說的是創造.
但不只是創造多偉大這麼簡單.
而是在創造界裡面.
是一個所謂道德的規範.
一個道德的規律.
懶的人就會肚餓.
勤奮的人他就會有貢獻.
這些都是一些所謂moral order.
moral的principle.
智者就是去洞悉這件事.
洞悉了之後.
他能夠有生命裡面的insight.
這個叫做Wisdom.
第二點我想說的.
什麼叫智慧.
智慧就是他明白.
最終是需要聆聽上帝的話.
這個才叫做真正的Wisdom.
給多一個圖片給大家.
這裡你看到什麼.
破爛了的牆.
石牆.
破爛了.
有這麼簡單嗎.

$^{401}$非常簡單.
但智者說了什麼.
第十四章.
他說我經過懶惰人的田地.
走過無知人的葡萄園.
他說看他長滿荊棘.
然後野草遍地.
石牆也都塌陷.
我就留心去觀察和思考.
我看見什麼.
I see.
And I receive insight.
災瑞遍事.
打頓遍事.
抱著傷臂躺臥遍事.
你的貧窮就如盜賊一樣.
貧窮就像拿盾牌的人.
好像士兵一樣來到.
這個應該大家很熟悉.
就是又是observation.
to insight.
觀察.
然後他有一些領悟.
就是他經過一個懶惰人的田地.
看見什麼.
就是石牆倒塌.
葡萄園.
你知道一個葡萄園要多少年才能夠成熟.
能夠出到那些葡萄.
才能去超級市場買到那些葡萄.
要多少年來培養它.
幾年至六年.
起碼幾年至六年.
要很悉心的培養它.
才能夠有一個葡萄園.
但落在一個懶惰人的手裡面.
去看管的時候.
你就見到生了野草.
然後石牆又會倒塌.
他就領略到什麼.

$^{441}$因為一個智者.
他不只是單單去觀察一些自然的規律.
然後放在實驗室去做實驗.
他主要想講什麼.
就是見到一個natural order.
天然的規律.
然後領略到一個所謂moral order.
一個生命裡面.
道義上的一些規律.
他經常講我們生命.
我們生命總會有一些勢力.
是想去破壞我們.
一定會.
你不需要去找那些東西.
它自己會來.
就好像石牆一樣.
石牆上的石頭.
一定是會向下跌的.
跌了之後會上回來.
一定會永遠都是我們的生命.
是由order變成chaos.
你想想是不是這樣.
我們當然不是單單只是你執不執房的問題.
我們的生活.
個人的生命.
你的家庭.
你的婚姻.
所有事情.
事業都是一樣.
永遠都是有一個這樣的order在.
如果你不做任何事的話.
它永遠都是會end up什麼呢.
是chaos.
除非你有那個領悟.
你會明白.
阿何一定要有input.
一定要去勤奮的去改善它.
以致那些的chaos不會再影響我的生命.
將那些塌了的石牆推上去.
將倒塌了的東西我們再推上去.

$^{481}$全部都是需要我們的labour.
全部都是需要你去肯做.
並且下手下腳去做.
spiritual growth都是這樣.
屬靈的生命都是一樣.
你就這樣把它放在那裡.
它是不會長的.
你說不長ok我沒事.
不是的 智者說不單止不長.
你還會長野草出來.
你的生命會倒塌的.
因為它們這些影響你的勢力.
破壞你的勢力.
不斷會來破壞你.
你一定要去落力的.
來給input下去.
生命的成熟.
我們的品格的塑造.
全部都不是自然而然的.
一定要用心去落下去.
有hard work.
要流汗甚至流淚.
這個就是智者所領悟不到的.
wisdom.
如果你心靈裡面越來越有wisdom的時候.
你應該留意到這件事.
並且你應用到在你心靈裡面.
以至你和上帝的關係.
你生活裡面每一方面.
你需要管理的東西.
包括你自己個人self leadership.
你自己的脾性.
你自己的所言所行.
你和人的關係.
你的配偶.
你的兒女.
你的家庭.
所有東西都是一樣.
我自己做父親的人.
我自己都很清楚.

$^{521}$當我不check住我自己生命的方向的時候.
很多時候.
就是一些破壞的勢力做出來的.
是我不需要給任何力量的.
它自然就會出來的.
給我們所有的父親.
或者所有的人.
一個很大的提醒.
第二方面.
他就說到最主要不是做一個好人這麼簡單.
原來就是我們和上帝之間.
有一個約的關係.
covenant的relationship.
就是我們需要在上帝面前.
來領略那個教訓.
以致我們能夠忠心來活.
十九章二十七節.
那個父親和兒子說什麼.
他說我兒.
如果你停止聽.
智慧的話.
你就必定會偏離.
如果你不肯去聽.
如果你硬著心的話.
你就一定會偏離.
again這個又是人性的order.
是我們需要work against的.
這裡說到你偏離智慧的言語教導.
其實那個字是discipline.
discipline有誰想聽.
沒有人喜歡聽.
你有沒有一些好的朋友.
一個好的智者在你心靈裡面.
很多時候你找他.
你不是很想找他.
甚至想避開他.
因為他說的話不是很合聽.
但是我們的人性在哪裡.
這裡說就是忠言逆耳.
明明是想你好.

$^{561}$我們偏是不聽.
這裡就叫做教導的意思.
discipline不是一些安慰的話.
等於你去看醫生.
為什麼這麼多人不喜歡看醫生.
因為醫生很多時候跟你講真話.
糟了你有cancer.
你喜歡聽那句話嗎.
不是很喜歡聽.
不是那些咖啡店.
聽聽音樂.
大家談的那些.
你去找他.
他照完你一輪.
然後他說一些.
總是不好的消息給你聽.
醫生牙醫我們最不喜歡.
但是如果你不聽的話.
你就一定會走歪.
這裡說什麼.
偏離道路.
原文有個意思就是說.
是不能夠行直路.
就好像一個醉了走的人.
他走不到直線.
這樣的意思.
如果你停止聽智慧的訓誨.
你就一定好像醉酒一樣.
你就走很多冤枉路.
然後你就走錯路.
跌下山都不知道.
很多時候我們想起.
我們跟上帝之間的關係.
所謂的spiritual growth.
就是你跟上帝之間有一個弱的關係.
我們很多時候都想起.
我們培養跟上帝的關係.
很多時候我們想起這種畫面.
對不對.
退休會.

$^{601}$我們坐在湖邊.
然後有一杯好咖啡.
然後我們有一本聖經在我們面前.
然後我們就默想一下.
很多時候我們想起這個圖畫.
我想告訴你這個圖畫是錯的.
這個是我們浪漫化了.
浪漫化了我們生命的成長.
生命的成熟.
真正的圖畫應該是好像這樣的.
這個才叫做spiritual growth.
是要流汗流血流淚的.
要下苦功的.
這些才叫做spiritual growth.
請你把它抹掉.
不要以為特別那些從小到大在教會長大的.
我們去什麼退休會.
我們很開心.
那些叫做我們去增加跟上帝的關係.
或者去增加我們的character.
我們的character formation.
我們的spiritual formation.
我們所謂的生命的成長.
大哥不是這個.
這樣是不會成長的.
It's not possible.
智者說什麼.
懶惰的人就會生命虧損.
勤奮的好像這樣.
你肯下功夫的.
肯去聽的.
肯去學的.
肯去謙卑下來的.
這個才會是一個祝福的道路.
請你不要看重.
多下一點力.
所以為什麼在開頭的浸賢的時候.
他說什麼呢.
他說敬畏耶和華是智慧的開端.
慾望人是藐視智慧和是忿悔.

$^{641}$聽清楚不是慾望人忘記了智慧和忿悔.
是藐視.
藐視.
Despise.
我現在要你聽父親的教誨.
因為如果你聽的話.
就會成為你頭上恩惠的華冠.
頸上的項鍊.
你就會有尊榮.
你的生命才能夠活出上帝給你的榮美.
在這裡對每一位的父親.
做一個挑戰.
也給我自己一個挑戰.
You need to work hard.
to build your character.
You need to work really really hard.
這樣你才會有一個鎮定的角色.
你是這樣.
我也是這樣.
是沒有藉口的.
沒有什麼藉口.
因為藉口我們就說.
我們都不是一個很好的基督徒.
It's ok.
I think you missed the point.
The point就是我們的生命一定要有那種formation.
不是在混亂裡面.
而是在order裡面.
以致我們能夠祝福其他人.
如果我們經常想著.
我也是這樣.
我也是很水皮的.
然後上帝也不care.
上帝care.
你去教養兒女.
你在家庭裡面有的責任.
有的使命.
上帝care.
很關心.
過程是交給我們.

$^{681}$你也去負責的.
結果不一定是我們能夠掌握得到.
我們不能夠控制結果.
兒女最後會走到哪裡.
這個不是我們的control.
但在我們的control.
在我們的management底下.
就是我們的人生.
上帝要我們是勤力的去做.
用功夫去做.
我相信我不是只在說父親.
是在說我們每一個.
你們做妻子也好.
做媽媽也好.
做姐妹也好.
做一個人就應該要在上帝面前.
去想清楚.
怎樣去build我們自己的character.
這些不是容易的東西.
有時候要流汗.
有時候要流淚.
有時候要聽中言.
雖然是易耳.
但我們需要聽受上帝的教導.
果然你要聽父親的教誨.
這樣你的生命才會有尊榮.
我們請我們敬八台上來.
我們同心得由祈禱吧.
上主我們多謝你.
因為你是那一位真正的智者.
主叫我們在智慧當中能夠去長大.
叫我們在生命裡面.
能夠更多屬靈的屬天的智慧.
主更加希望你能夠給我們.
一個馴服的耳朵.
讓我們能夠聽你自己的教導.
並且遵行去遵守.
讓我們敬畏你.
尊你為榮.
以致我們生命.

$^{721}$可以在這個智慧的開端上.
起到行第一步.
我們祈禱奉主名.
我們一起起來吧.
\newpage



\section{8:31-32}
\label{sec:K5EeMHs5F9w}
\textbf{Dwelling in The Truth (John 8\_31-32) - Sam Johnson}
\newline
\newline
連結: \href{https://youtube.com/watch?v=K5EeMHs5F9w}{\texttt{https://youtube.com/watch?v=K5EeMHs5F9w}} ~~~~ 語音日期: 2023-06-27
\newline
\newline
\hyperref[sec:F8HOtvMKDlc]{\small{< < < PREV SERMON < < <}}
~
\hyperref[sec:index_chronic]{\small{[返順時目]}}
~
\hyperref[sec:index_scriptual]{\small{[返順卷目]}}
~
\hyperref[sec:hjFzQBCexic]{\small{> > > NEXT SERMON > > >}}
\newline
\newline
$^{1}$Thank you, James..
And I'm excited to be here..
Glad to see everyone here..
I'm always honored to be able to be at Urban Voice and to be here again and to see old.
friends and new friends..
And excited to open the word with you all today..
So we're going to be looking in John chapter 8 today..
If you want to turn there in your Bible or on your Bible app or whatever you're using.
here, John chapter 8, we're going to be looking at verses 31 and 32, these kind of famous.
verses where Jesus is talking about truth..
And these are important verses for us..
Culturally, we live in a time where people talk a lot about truth and think a lot about.
truth, but it's hard for us to define it..
It's hard for us to figure out what we mean when we say the word truth..
And Jesus has some really profound things to say here that are really good for us to.
understand, especially in the church..
So we're going to look at these today..
John 8, 31 says this, "So Jesus said to the Jews who had believed him, 'If you abide in.
my word, you are truly my disciples, and you will know the truth, and the truth will set.
you free.'".
What's happening in this passage, Jesus, earlier in the passage in verse 12, he declares.
that he is the light of the world..
He reveals that to the people that are listening to him..
And immediately, the Pharisees who are listening accuse him of lying about himself in verse.
13..
So Jesus begins to explain to the Pharisees and everyone else who's listening that they're.
misunderstanding the kingdom of God..
And that happens in verse 14 on down to verse 29..
And then verse 30 tells us that as a result of what Jesus is saying, a lot of people who.
are listening begin to believe in him..
They begin to place their trust in him and have faith in him that he is who he says he.
is, that he is the Messiah, the light of the world as he proclaimed himself to be..
Then verse 31, he begins to speak to those people, the people who have begun to believe.
in him, the people who are following him..
And he says, "If you abide in my word, you are truly my disciples, and you will know.
the truth, and the truth will set you free.".
That means these words are spoken to people like us, people who are followers of Jesus,.
who have already made the decision to believe in him and follow him..
And Jesus is giving us some profound, significant truth about what it means to be a disciple..
Now obviously, in our culture, at times we have a strange relationship with what we call.

$^{41}$truth..
Most of us want to define truth for ourselves, and so we make a lot of statements like, "That's.
true for you, but not true for me," or we say, "Live your truth," things like that..
But even as we seek to make truth personal and self-defined, we keep running into these.
kind of unyielding boundaries of reality..
So then what we do sometimes is we relabel the boundaries in order to kind of fit our.
thinking..
We say, "Well, maybe if we change the names and move the goalposts, then it'll make sense.".
So now we have terms that culturally we use like misinformation and disinformation because.
there's so much falsehood..
There's so many things, so many labels that have been changed that people get confused.
about what is true, what's right and wrong..
Yet even in that, we can't decide really what disinformation is or who gets to determine.
what the standard of disinformation is, and we argue about that..
We're desperately trying to figure out what's true, but a lot of times we're also desperately.
trying to make truth in our own image..
And we want the truth to match our desires and our preferences and usually to excuse.
our shortcomings, to excuse our mess..
But the reality that the Bible puts in front of us is that we don't live in a world of.
our own creation..
We don't get to make truth in our image because we live in a world that's created by the almighty.
creator God, and he's the one who's defined reality and truth itself according to his.
own character..
And as much as we might try to escape that, we can't..
We can't escape that pull of his reality, the world that he has made..
It's like gravity, it keeps pulling us back down..
So what are we talking about when we say the word truth?.
How are we going to define that?.
How are we going to define that today?.
I think the simplest and most accurate way to think about the word itself is to say that.
truth is something that corresponds to reality..
So a simple illustration of this is if I were to say that I was wearing red leather pants.
today, we could safely say that's an untrue statement, right?.
Because I'm not wearing red leather pants, I'm wearing some type of cotton polyester.
blend that's a tan color, right?.
So that's kind of an obvious illustration..
If I say something that doesn't correspond to the reality, then we can say, well, that's.
a false statement, that's not true..
Well, if truth is something that corresponds to reality, where did reality come from?.
Where does ultimate reality come from?.

$^{81}$Well, the Bible tells us that ultimate reality comes from God..
That God's the one who created the universe, created everything that we see..
He is not only the creator of ultimate reality, but he's the only one who really understands.
it and comprehends it and encompasses it, right?.
Because God has all knowledge and all power in the universe..
So God's really the only one who can know what ultimate reality is..
He's the only one who can look at all of it and see all of it..
At very best, our versions of reality and our visions of reality are incomplete at very.
best..
Sometimes if we're honest, we can admit that sometimes they're just plain wrong..
Sometimes we get totally the wrong idea about what is real and important in life..
But at very best, what we see as real and important in life is gonna be incomplete because.
we cannot see all of reality..
We can't see all of truth..
We do not know all the things there are to know..
But God does..
He created reality..
He created the universe..
He created morality and what's right and wrong..
Those things are defined by him..
And so when he makes a statement and says, "Hey, this is what's right..
This is what's good for you..
This is the way to go.".
He actually knows what he's talking about because he possesses all the knowledge and.
all of the wisdom and he can see it from every angle and he can make a true statement..
And we can't do that because our knowledge is narrow..
Our wisdom is narrow and our vision is very closely defined..
We can only see things from our own perspective..
And that reminds us that when we think about truth, there really isn't any such thing as.
my truth and someone else's truth, right?.
There's really only God's truth..
God's the only one who can say what truth is because he's the only one who can see it.
all..
We can have our versions but ultimately, only God's truth is gonna be truth..
Only God can possess that..
And Jesus begins to speak about this in chapter 8 here..
Jesus starts to talk about how we, if we are gonna be a disciple, a follower of Jesus,.
then we need to understand his truth and how that takes place, what that looks like in.
our lives..
So the first thing I want us to see that Jesus is communicating here is that if we're gonna.

$^{121}$dwell in the truth, then we've really got to soak in God's word..
And in verse 31, there's a couple of key words here..
He says, "If you abide in my word, then you are truly disciples of mine.".
The word in the original language for abide means to remain or dwell or to continue to.
be present in a place or to sojourn there or another meaning which I really love is.
to endure..
So he's saying if you stay in my word, if you're always there, if you're dwelling there,.
if you're living there, if it's your home, then you're a true disciple..
And one of the cool things about this word, like I said, I like that meaning of endure..
It reminds us that there are times when it's not easy to remain..
It's not easy to stay there..
It's not easy to continue on..
And that's what Jesus is telling these followers who have just taken the step of faith..
They said, "Jesus, we believe..
What does it mean to follow you?".
He says, "You've got to remain..
You've got to stay..
You've got to continue to dwell there.".
The other key word in this verse is the word logos, the word for word..
When Jesus says, "If you abide in my word," that word in the Greek means like an idea,.
a concept, a teaching, a doctrine, a saying..
But in the book of John, particularly, it also refers to Jesus himself..
And I think that gives this verse a really beautiful double meaning..
Jesus is saying, "You've got to dwell and remain and abide in my teaching, but you've.
also got to dwell and remain and abide in my presence.".
There's two things that are being said in this verse..
Jesus is saying that you've got to understand the truth that I've communicated to you..
You've got to learn those things and you've got to abide in them, but you've also got.
to abide in relationship..
You've got to be in my presence with me..
And as we said earlier, when he says that, when we think about what truth is and we say,.
"Okay, what does it mean to remain in God's word?.
How are we going to know the truth about reality?".
The clearest implication for us, living 2,000 years after Jesus was on the earth, is that.
we need to spend time in God's word, the Bible..
Now for these people that Jesus was speaking to on this day, he said, "If you abide in.
my word," that could have meant following Jesus around and listening to what he was.
saying..
But we're not in that position, right?.
We're not in that position, right?.

$^{161}$Jesus is no longer physically present on the earth in a physical body..
So we can't say, "Okay, well, to dwell and remain in Jesus' teaching means I've got to.
go to where he is and listen to what he says," because his physical presence has left the.
earth..
So what it means for us is that we have to spend time in the Bible..
And we know that God has told us in the New Testament that the Bible, scripture, is his.
word spoken to us, just like if he were speaking in a conversation..
So the Bible is God's revelation to us..
It's God revealing himself, God saying, "Hey, this is who I am..
This is what's important about the universe..
This is what's important about how you live your lives..
This is how I want you to treat other people..
This is how I want you to care for others.".
All those things are spoken to us in scripture, and that's what it means for us to abide,.
to remain in Jesus' words, his teachings, his ideas, and also in his presence, is to.
be immersed and saturated and soaked in the scripture..
And I chose the word "soak" intentionally because it gives us the picture of something.
that surrounds you and you're just in it all the time..
And that's what our relationship to scripture should be like, especially in the world that.
we live in..
There are so many things that challenge our faith..
There are so many things that are difficult..
There are so many things that we face in any given week that can really discourage us or.
harm us or knock us down..
And it's not going to be enough to say, "Well, I read a Bible verse once," or, you know,.
"I remember something that I heard one time in church," or "One day a week, I'm going.
to try to connect with God.".
We need to be soaking in it..
We need to have it all the time..
It needs to be a life-sustaining thing that is constantly coming into our lives..
And so we need to be spending time in God's Word..
We need to be soaking there..
We need to saturate our lives, be reading it, be considering it, reflecting on it, meditating.
on it, saying, "What does this mean?.
What is God saying to me?".
We need to do this, where we get together with other believers and talk about it and.
say, "What does this verse mean, and how can we encourage each other to follow it, and.
how can we take care of one another and support one another?".
That's all a part of remaining in His words, in His teaching that He talks about in verse.
31..

$^{201}$We need to be alive in that relationship with Christ so that we're in His presence, and.
the Holy Spirit's working to illuminate God's Word to our hearts..
Now the obvious implication of soaking in God's Word is that to dwell in the truth,.
we have to know the truth..
And in verse 32, where it says, "And you shall know the truth," that word in the original.
language "to know" means that you perceive or understand or become acquainted with something,.
but it also communicates a knowledge that's grounded in personal experience..
So that's communicating to us that there's going to be a deep knowledge there that involves.
the experience of our lives, the things that we go through, not just kind of a head knowledge.
where we're like, "Yeah, I heard that," or "I read that.".
It's gonna really impact our day-to-day..
And first of all, the first step of walking in truth and knowing truth is to just become.
aware of it..
Like we just said, we've got to engage with it, we've got to read it, we make ourselves.
aware of it..
But it has to go beyond that..
If we are gonna dwell in the truth, we have to submit to the truth..
You shall know the truth, and the truth shall make you free..
That word "know" is not just saying you know it in terms of awareness..
A lot of times when we use the word "know," we mean awareness..
Like you might say, "Did you know that Chicago's the third largest city in the U.S.?".
I'm not saying, "Do you know that?".
Like, "Have you experienced it?".
I'm saying, "Are you aware that Chicago's the third largest city?".
It's two very different things, even though we use the same word..
And so when this says, "You shall know the truth," it doesn't mean you're aware that.
Jesus loves you, but it's saying, "Have you experienced the love and grace of Christ?".
Two very, very different things..
Knowledge that is more than awareness..
In Romans 1:19-21, Paul talks about people that know, it says, "For what can be known.
about God is plain to them, because God has shown it to them for His invisible attributes,.
namely His eternal power and divine nature, have been clearly perceived ever since the.
creation of the world in the things that have been made..
It's there without excuse, for although they knew God, they did not honor Him as God or.
give thanks to Him.".
Paul's talking about people who know that God exists, they know He's there, they're.
aware of God, but they're not submitting to Him..
They're saying, "Okay, I understand that God is there, but that's not going to change anything.
in my life.".
So if we're going to dwell on the truth, if it's going to be our living place, and as.

$^{241}$Jesus says in verse 31, if we're going to be a true disciple, then we've got to submit.
to it, not just know it..
We've got to say, "Okay, God, I hear what you're saying..
I may not really like it that much, but I'm going to submit, I'm going to agree with you..
I'm going to actually start to live that out and apply it in my life, even if it makes.
me a little bit uncomfortable, even if it's different than how I wish it would be.".
And as we submit to it, we've also got to respond..
That's a part of that knowing that is grounded in personal experience..
First John 3, 18 and 19 says this, "Little children, let us not love in word or talk,.
but in deed and in truth..
By this we shall know that we are of the truth and reassure our heart before Him.".
It's not enough just to talk about it..
It's not enough just to give mental assent to it..
We have to live it out..
We have to respond..
When God challenges our heart and says, "Hey, here's some of my truth, and this is what.
it would mean in your life," we have to say, "Okay, I understand that..
I know it..
I want to submit to it, and then I'm actually going to respond.".
So like tomorrow, I'm going to do something different than what I did yesterday..
That's the response..
There's the awareness of the truth..
There's the submission to the truth where we say, "Okay, I'm actually going to submit.
to it, apply it to my life.".
But we still haven't responded until we've done it..
And that's when the truth changes us..
That's when it starts to change who we are..
Another thing that's a huge part of the process of truth changing us and dwelling in the truth.
is that we have to reject lies..
And that's kind of an obvious statement..
It's like, "Well, obviously," right?.
But this is really, really tough for us, especially if you've been in church for very long..
You may know a lot about the Bible..
You may know a lot about God's truth..
You may even submit to it and be trying to respond to it..
But sometimes we're stuck in life because we're still holding on to lies..
God comes along and he says, "Hey, I want to put something great in your hand..
I want to use you for me, or I want to bless you in some way, or I want to do something.
in your life.".
And he can't put anything there because we're gripping tightly to something that's not true.

$^{281}$that we're still holding on to..
And we've got to reject those lies and let go of them so that we can move forward with.
God..
And understand this..
This is really important to understand..
The devil cannot stop anything that God is doing..
Sometimes we get confused and we start to think that God and the devil are like rivals..
They're not, okay?.
God is infinitely more powerful than every being in the universe, including the devil..
So the devil has no equality with God in any way..
The devil hates God and he wants to mess up anything God's trying to do..
But he can't stop God..
He doesn't have the power of God..
And so the way that he tries to stop God and mess with his plans is just by lying to us.
and convincing us not to cooperate with what God is doing..
The devil can't stop it..
He can't fight a head-to-head battle against God and win, and he knows that..
But he can lie to us and get us so distracted or so discouraged or so deceived that we.
head in a completely different direction and then it starts to thwart the plan of God in.
our life..
And so that's what he wants to do..
His main tactic in our lives is to lie to us..
And so it's really important for us to identify those lies and reject them..
To take them captive when they come into our mind, literally to say, "That's not true and.
I'm not going to believe that..
I'm not going to let that idea just roll around in my head.".
And then replace those lies with God's truth..
Say, "You know what?.
Not only am I not going to believe that, but I'm going to believe what is true.".
All right, here's a simple illustration for you to think about..
God is great..
God possesses all power in the universe, okay?.
That's a truth that we probably, most of us could easily agree with..
It's throughout scripture..
And we give mental assent to it..
We may even submit to it and say, "Yes, God is great..
That's why I'm a Christian..
That's why I'm here in church, because I believe that God is great.".
I'm trying to respond to that, right?.
But here's a lie that comes against that truth..

$^{321}$I have a problem in my life and I feel like the problem is going to destroy me or it's.
going to, you know, I'm super discouraged about it because I feel like there's no solution,.
right?.
When I'm thinking that way, I'm starting to believe that in some way God is not great.
enough for that problem..
I'm starting to believe that, yeah, God may be great, but like this is a really big deal..
God is great, but he's never met my boss..
So you know, we're starting to believe that way..
And so we have to reject that lie and say, "You know what?.
Yes, my boss is really challenging to deal with," or, "Yes, this situation is a mess,".
or "Yes, I've made some bad choices and created some chaos in my life.".
But the truth is that God is bigger than this..
And so I'm not going to listen to the lie that says, "Well, my life's going to be a.
disaster because of what happened or because of the mistakes I made or, you know, I'm going.
to fail at everything I try.".
I'm not going to listen to those lies because I know that the truth is that God has a power.
that's greater than anything else and he's at work in my life for his glory and for my.
good..
We have to reject those lies, say, "I'm not going to believe that," and replace them with.
God's truth..
Another simple one that's really huge is grace in our life, right?.
A lot of times, even though we give assent to the truth that, "Yes, Jesus loves me..
I'm saved by grace..
My sins are forgiven..
I don't, my mistakes are not held against me," we can believe that, but lies still come.
in and accuse us..
And in our mind, we feel so attacked and we feel like, "But I'm a miserable person and.
like I just fail so much and I'm not this enough or that enough or fill in the blank.
enough.".
And so we're discouraged because of that..
It's the same thing there..
We have to say, "Yes, I am certainly imperfect and I need God, but God's grace covers my.
sin..
I'm forgiven..
I'm not defined by those mistakes.".
As Pastor Mark read the passage earlier, "Nothing can separate me from the love of God.".
Those are truths and we need to bring those into our mind and say, "I'm not going to believe.
the lies that I'm just a failure and that I'm worthless and that everything's going.
to turn into disaster in my life..
Those things are not true..

$^{361}$I'm going to believe what God has said about me, that He loves me, that I'm forgiven, I'm.
accepted..
He has a plan for my life and yes, it's not complete yet..
I'm still a work in progress, but God is going to bring that to pass..
He is going to complete it.".
That's what it means to begin to reject lies..
If we're going to dwell in the truth, we've got to know the truth..
We've got to soak in it..
We've got to submit to it..
We've got to respond to it, but we also have to start saying no to these lies that get.
thrown at us by our own desires, by the culture around us, and by the enemy..
The devil wants to throw those lies at us..
They're fiery darts that he shoots at us to try to bring us down..
Then what I want us to see is what comes out of this..
If we begin to dwell in the truth in the way that Jesus is talking, what comes out of it.
in our lives?.
What's the fruit of it?.
Well, the first thing that we see Jesus point out here is that as we dwell in the truth,.
we live as true disciples..
We're really true followers of Jesus if we're remaining and living in His truth..
You know what that's saying is that the staying power of our faith in Jesus is not a product.
of the ease of our pathway as a disciple, but it's a product of our connection to Jesus.
and God's truth..
What I'm saying there is the reality and the truth of your followership of Jesus will not.
be a product of how hard or how easy your life is..
It will be a product of how close you are to Jesus..
If you're close to Him, you can go through a lot of really difficult things..
You can come out on the other side and be healthy and be joyful and be full of all that.
God has for you in your life..
Jesus here is not saying when He says, "If you abide in my word," think about that statement,.
"If you abide in my word.".
That implies that it's not easy to abide in His word..
It implies that there's going to be some challenges..
It implies that some people are not going to abide in His word and they're going to.
turn back..
It implies that some people aren't going to follow His teaching and some people aren't.
going to understand..
That statement in and of itself is saying, "Hey, it's not easy..
It's not always going to be fun, but if you will remain, if you will dwell in His teaching,.
in His word, in His truth, then you'll be a true disciple.".

$^{401}$I think all of us desire to be a true disciple of Christ..
That's why we're here..
Jesus is saying this is the way to do it, to remain, to dwell, to live in His truth,.
in His word..
And the other beautiful thing that we see is that as we dwell in the truth, we experience.
freedom..
Jesus makes that clear in verse 32, "You shall know the truth and the truth shall make you.
free.".
Think about that..
When we know the truth, when we know the truth about who we are, who God made us to be, the.
fact that we're forgiven, we're accepted, that we have a place in God's kingdom, God.
wants to use the gifts and the talents and the skills that He's given us to love and.
help others and to glorify Him..
When we know that truth, what a free place to be..
That can set us free from so much self-doubt and self-loathing and fear and worry..
When we know that God's power is greater than our sin, and we begin to experience the freedom.
of getting away from the things that hold us back, the hang-ups, the addictions, the.
struggles, and God gives us victory, and we're getting closer and closer to Jesus, and we're.
walking away from our old self and saying, "Yeah, that used to be who I was..
I used to struggle with that..
I used to get angry all the time..
I used to have these problems, but I don't because Jesus is changing me..
I'm being transformed from the inside out.".
That's a place of beautiful freedom..
A freedom from bondage, those things that have harmed us in the past..
And ultimately, it's a freedom to get closer to Jesus to know more of His truth..
It's like a spiral, except inverted in a healthy way, right?.
A lot of us sometimes spiral in unhealthy ways..
I do that plenty of times..
I start thinking about, "Oh, I messed this up," and then I realize, "Oh, I also messed.
that up.".
And then pretty soon, I'm like, "Oh, I messed everything up.".
We start spiraling down, right?.
But this is like a spiral going up in a healthy way, where as we know God's truth, and we.
get free from lies and unhealthy things, now we're able to believe more of the truth, and.
submit to more of the truth, and respond to more of the truth, and then we're able to.
learn something else that God wants to do in our lives, and submit to that, respond.
to that..
And pretty soon, we're getting closer and closer to Jesus..
We're getting healthier..

$^{441}$God's doing a lot of just beautiful and amazing work in our lives..
And that's what happens..
As we dwell in the truth, we experience freedom..
And all of us want that..
All of us need that..
And the other beautiful thing about this is that our world needs these type of people.
that we're describing, right?.
Think about people that you work with..
Think about people that you go to school with, or that you are around, or even people in.
your own family..
How much do some of those people need someone in their life who understands that they are.
loved by God, and who is able to love other people well, and who is not listening to all.
the lies of culture, and all the self-condemning things, but is able to see themselves in a.
healthy way, and respond to others in a healthy way?.
How much do the people around us need that?.
So as we dwell in the truth, God begins to use us..
So that's my prayer for myself..
It's my prayer for all of us..
As we follow Jesus, as we become true disciples, that we would increasingly grow in embracing.
God's truth..
We need to soak in God's word..
If you're not doing that, maybe that's a starting place for you to say, "Okay, I'm.
familiar with the truth a little bit..
I know some stuff, but I don't really spend that much time there.".
Maybe you need to start soaking..
Maybe you need to kind of get immersed in it, spend some more time there..
Maybe submission is the starting place for you..
Maybe you say, "Well, yeah, I've been in the church a long time..
I know a lot of stuff about scripture.".
But maybe there's just some areas where you know what you're supposed to do, but you're.
just not really doing it..
You haven't submitted to it..
You say, "Yeah, I know that's true, but I don't really want to mess with that part of.
my life.".
Maybe submission is the place for you..
Maybe it's responding..
Maybe you say, "I know a lot of truth, and I agree with it..
I've submitted to it, but I just, I try, and I don't really follow through..
It's hard for me.".
Maybe you need to respond..

$^{481}$Maybe you need to reject lies..
A lot of us get stuck there where we're trying to live out the truth, but we just keep holding.
on to some of those lies about who we are and about our own weakness and how we can't.
do anything..
So we just get stuck..
Maybe you need to start rejecting some of those lies..
But may all of us really grow in living as genuine disciples and being set free by the.
truth that God has for us..
Let's pray..
God, I thank you so much for the truth..
I thank you so much for who you are, that you, as the creator of the universe, you,.
the one who possesses all wisdom and knowledge, you care about us, and that you're willing.
to speak to us, and you're willing to tell us, "Hey, this is the way to go..
This is the way to be healthy..
This is the way to live your lives.".
Because without you, we wouldn't know that..
We'd just be wandering around in the dark trying to make ourselves happy and make ourselves.
feel good..
And most of us would be becoming increasingly miserable..
So God, I thank you that you have spoken to us..
You've given us truth in your word..
And God, I pray that you would help us, especially right here in this room, as your followers,.
as disciples, help us to dwell in it, to remain in it, to soak in it, to spend our time in.
it, help us to submit to it, to respond to it, to apply it to our lives, help us to reject.
the lies that come against the truth and to replace them with your truth in our mind over.
and over again, as many times as we need to, to begin to believe the things that you've.
said about us and about our world and about our lives, and help us to get free from things.
that hold us back, help us to walk in freedom so that we can go boldly and confidently into.
the world and love others and point people to the freedom that we've found in Jesus..
God, we just thank you and praise you for who you are and pray that these things would.
be more and more evident in our lives..
In Jesus' name, amen..
Amen..
[MUSIC PLAYING].
\newpage



\section{18:1-4}
\label{sec:hjFzQBCexic}
\textbf{Are You Limping Through Your Faith (1King 18\_1-4, 17-40) - Tommy Lee}
\newline
\newline
連結: \href{https://youtube.com/watch?v=hjFzQBCexic}{\texttt{https://youtube.com/watch?v=hjFzQBCexic}} ~~~~ 語音日期: 2023-07-08
\newline
\newline
\hyperref[sec:K5EeMHs5F9w]{\small{< < < PREV SERMON < < <}}
~
\hyperref[sec:index_chronic]{\small{[返順時目]}}
~
\hyperref[sec:index_scriptual]{\small{[返順卷目]}}
~
\hyperref[sec:NGEVacpuPjM]{\small{> > > NEXT SERMON > > >}}
\newline
\newline
$^{1}$Amen..
- Good morning, thank you very, very much..
All right..
Well, happy July 4th..
Mark, it's good to have you back from Israel..
Last time I was here, you were missing,.
and it is always good to be back here at UV..
We've been in this whole series.
about the life and times of Elijah,.
and one of the things I encourage you.
that I've really enjoyed doing is studying,.
taking a character, taking a book of the Bible,.
and just digging down into it..
We spent the last three, four times I've been here.
talking about 1 Kings chapter 17..
There is nothing like sitting there.
and just reading the same verses over and over again,.
and thinking about it throughout the week,.
and saying, "Lord, what is it.
"that you're trying to teach me?".
Today, we're gonna look at 1 Kings chapter 18,.
but before we do, let me just give a summary.
of where we've been for those who are new,.
for those who've been with us,.
just a refresher, and as we look at Elijah,.
we started off in 17.1, where we talked about,.
literally, Elijah means, his name means,.
my God is Jehovah..
Elijah is a Tishbite of Tishbe in Gilead..
What does that mean?.
He literally came from obscurity..
No one knew where Tishbe was..
They knew where Gilead was,.
but they didn't know where Tishbe was,.
a small, little city, and so he came from obscurity,.
and somehow this prophet who came from obscurity.
got to the capital city of the northern kingdom of Israel.
in Samaria, and he gets a meeting with King Ahab,.
and he tells him, "Hey, look,.
"there's gonna be a drought..

$^{41}$"There's gonna be a long, long drought,.
"and actually, that drought will last about three years.".
And he says, "Neither dew nor rain,.
"except by my word," and the motivation.
in the early parts of 17 is,.
"the God before whom I stand.".
That becomes a theme in his life,.
a theme in this chapter of Elijah's life..
It's the God before whom I stand,.
and just like that, as he confronts Ahab,.
he gets ushered away to Cherith,.
a place where he is in solitude..
He's fed by ravens, and he is also,.
drinks from the brook over there..
Well, in time, he doesn't just stay there for a short time..
He stays there for a long time,.
and now, later on, the brook dries up,.
and then now the second part of 17,.
God says to go pack up, leave Cherith, and go to Zarephath,.
and when he goes to Zarephath, he says,.
"Hey, look, I see you..
"I know where you're gonna go.".
It reminds us all the time that God sees us.
wherever we're at, and God knows where we need to go..
Sometimes we may not know where to go,.
but God knows where we need to go..
He also knows what you need, 'cause he says,.
"When you get to Zarephath,.
"you're gonna meet a widow there," and he does..
He meets a widow there, but the problem.
with this widow is she's collecting sticks,.
and she says, "I only have a little flour..
"I have a little oil, and I'm collecting sticks.
"so I can go in and make some bread,.
"and then from that point on,.
"my son and I are gonna eat it and die.".
Really, that's who you want me to sit there?.
Yes, because sometimes God knows what we need,.
and what we need is different than what we think we need.
or what we want, but God knows what we need..

$^{81}$Along the way, just like the story with the widow,.
we're gonna have to overcome first impressions.
or expectations..
We're gonna have to overcome physical impossibilities..
It's literally a widow who is poor,.
who's gonna feed him, who thinks she's gonna die.
because she only has a little flour and a little oil,.
but somehow God makes that oil and that flour.
just keep going and going and going,.
but yet at the same time, this is the last story.
that we looked at when I was last here..
Sometime later, as he stays with the widow,.
the widow's son dies, and the widow blames Elijah,.
"That you brought this upon my family.".
Elijah grabs the son, he brings it, and there we see,.
we looked at Elijah's prayer to the Lord,.
a genuine prayer where he wrestles with the Lord.
because he doesn't know what happened..
He sits there and he yearns and he prays.
in a genuine fashion, the fashion that God wants us.
to wrestle with him in our own prayer life..
He prays the impossible, but he also knew.
that God could say no..
That's what it means to walk with the Lord..
There is a deep relationship that you have.
with your prayer life..
It's not just sitting there in your prayer.
and just asking God to provide all the things.
that you want, but you're wrestling..
There's a natural, there is a genuine relationship.
with the Lord, and God chooses to heal the son..
And I love what verse 24 and 17 tells us..
Now I know that you are a man of God,.
and that the word of the Lord in your mouth is truth..
She never says to him, "Now I know that you are.
"a great man, now I know that you are a powerful man.".
That's not what she says to him..
She says, "Now I know that you are a man of God,.
"and that the word of the Lord in your mouth is truth.".
Because everything that Elijah does in 17,.

$^{121}$and we're gonna see it in 18, is as if he is a servant.
and a steward of what God has given him..
And that's what a servant does..
He gives light, he gives credit to his master..
I know that you are a man of God..
We talked about the theory behind Hebrews 11..
Some of you guys who may not be familiar with Hebrews 11.
is the Hall of Faith..
It talks about all these Old Testament characters,.
and you know who is not mentioned in the Hall of Faith?.
Elijah is never mentioned by name..
There is one verse that does describe Elijah..
It's this, in verse 35 of Hebrews 11,.
"Women receive back their dead by resurrection.".
One of the things that we talked about.
when we were last here is why in the world,.
in the Hall of Faith, Elijah plays such an important role.
that he was never listed in the Hall of Faith..
Now, I do wonder if the writer of Hebrews,.
in order to honor Elijah, never mentions who Elijah is,.
but rather what he does..
Because Elijah was never about himself,.
but always about being a steward,.
always being a spouse, servant,.
always giving credit to what God is doing..
Which leads us today to the confrontation in Mount Carmel.
in 1 Kings 18..
And there today, we're gonna look at five questions.
that emerges from these in this chapter,.
as we look at 1 Kings 18..
And let me read it to you..
It's gonna be a long one,.
so we're gonna just go on little different passages.
in different verses, and try to get through it.
as soon as we can in the next 20 minutes..
18, verse one..
After many days, the word of the Lord came to Elijah..
In the third year, saying, "Go, show yourself to Ahab..
"It's been three long years of this drought,.
"and I will send rain upon the earth.".

$^{161}$So Elijah went to show himself to Ahab..
Now, the famine was severe in Samaria,.
the capital city of the northern kingdom..
And Ahab called Obadiah, who was literally the chief of staff.
over Ahab's kingdom, and called Ahab,.
who was over the household..
Now, Ahab was feared the Lord greatly..
And when Jezebel, Queen Jezebel,.
cut off the prophets of the Lord,.
Obadiah took 100 prophets and hid them by fifties.
in a cave, and fed them with bread and water..
Let's go to verse 11..
Obadiah now says to Elijah,.
"And now you say," Obadiah says to Elijah,.
"Go tell your Lord..
"Behold, Elijah is here..
"And as soon as I have gone from you, Elijah,.
"the spirit of the Lord will cure you..
"I know not where.".
Because he literally went from Samaria.
to Cherith and Zarephath, and this entire time,.
they're all looking for him, but they never found him..
And so they're always wondering.
where this guy disappeared to..
And so when I come and tell Ahab,.
and he cannot find you again,.
you know what he's gonna do?.
He's gonna kill me..
"Although I, your servant,.
"have feared the Lord from my youth.".
And Elijah says, "As the Lord of hosts lives,.
"before whom I stand.".
Same theme as in the beginning of chapter seven,.
"Before whom I stand.".
That is a constant theme of Elijah's life..
"I will surely show myself to him today.".
So Obadiah went to meet Ahab and told him,.
and Ahab went to meet Elijah..
When Ahab saw Elijah, Ahab said to him,.
"Is it you, you troubler of Israel?".

$^{201}$Verse 18, pay attention to this..
"And he answered, 'I have not troubled Israel,.
"but you have, in your father's house,.
"because you have abandoned the commandments of the Lord.
"and followed the bells..
"Now therefore send and gather all Israel to me.
"at Mount Carmel and the 450 prophets of Baal.
"and the 400 prophets of Asherah,.
"total of 850 people who eat at Jezebel's table.".
So Ahab sent to all the people of Israel in verse 20.
and gathered the prophets together at Mount Carmel..
And Elijah came near to all the people and says,.
"How long will you go limping.
"between two different opinions?".
First question that we see.
comes in this very verse that I just read..
"How long will you be limping.
"between two different opinions?".
The first question that emerges, I ask you,.
"Are you limping through your faith?".
That's the first question that emerges..
Are you limping through your faith?.
Three characters and three perspectives arise..
There's a character Ahab, as we see in chapter 18..
Ahab, we don't even know where his faith journey is,.
how he got out to worshiping Baal, all of that stuff..
We have an idea it really is through the influence.
of Queen Jezebel..
Ahab, out of the 19 kings in the northern kingdom,.
they said Ahab was probably the worst..
Queen Jezebel actually ends up being the main villain,.
the main person who disrupts everything..
She was the worst of even all of the kings..
Now, at the same time, we do have some idea.
that at one point or another, Ahab was following the Lord..
Here's why..
When Ahab dies, his two sons take over..
His two sons, the first son is,.
I'm gonna butcher this, so I've been trying to practice this.
for the whole weekend, I still butcher it..

$^{241}$Ahaziah, all right?.
Ahaziah is one of the boys..
His name literally means "owned by Jehovah.".
Interesting, "owned by Jehovah".
is what his first name's son means..
The next one is Jerhoram,.
which means "Jehovah is exalted.".
So why in the world are your kids' name,.
it means "owned by Jehovah" and "Jehovah is exalted,".
if at one point or another, you weren't following the Lord?.
And somehow or another, Queen Jezebel, who you married,.
really wanted to make Baal the state religion,.
and as a result, suddenly, all the prophets were killed,.
and now there's 450 prophets of Baal,.
400 prophets for Asherah..
There's also Obadiah, the next character that we see..
Obadiah is the chief of staff, as we mentioned..
He's the one who oversees the entire household..
He's the one, when Jezebel was trying to kill.
the 100 prophets, he rescues 50 and puts 50 in each cave,.
total of 100..
Well, but this time, he is scared of Ahab..
When he confronts Elijah, he says,.
"Hey, look, if you send to me and Ahab can't find me,.
"he's gonna kill me.".
And then you have the third folks..
Those are the Israelites..
The Israelites were also scared of Ahab and Queen Jezebel..
But then also, it's been a long three years, hasn't it?.
The drought has taken for three years..
God has not responded in those droughts..
And probably the Israelites says,.
"Hey, look, if our God doesn't respond,.
"then let's give Baal a try..
"Let's give the God of fertility,.
"the God of harvest a try..
"What can we lose?.
"Our God's not paying attention to us as well, too.".
When we allow our fear, our pursuit of other things.
become our focus, we find ourselves limping.

$^{281}$through our faith..
When we allow our fear, when we allow our pursuit.
of other things become our focus,.
we find ourselves limping along.
and limping through our faith..
I'm not sure for some of you guys who are listening.
on screen or those who are sitting here,.
have you ever been canoeing before?.
Canoeing is fun and you're constantly going downstream..
It is hard to canoe upstream..
You're constantly going downstream..
There is a downward pull on all of our lives.
as Christ followers to pull us downstream..
There is a downward stream, there is a downward pull.
on all of our lives to pull us downstream,.
to affect our faith and to allow us.
to live in a state of limbo..
It's like an ocean current as well, too..
You ever sit in an ocean current?.
If you ever sat in an ocean current for 30 minutes,.
you just sit on the sand right near the edge.
where the ocean's coming..
If you weren't paying attention 30 minutes later,.
where would you have gone?.
A different place because the ocean current.
would have moved you from one place to another.
and you would have never known..
In our work culture today,.
here's a term that also describes where we are..
If you are working right now,.
Harvard Business Review calls this the quiet quit..
It's not the great resignation, it's the quiet quit..
I'm not sure if you ever heard that term as well, too..
Post-COVID, it's been very, very popular..
So literally, a lot of times,.
workers in corporations realize that I hate my job,.
I hate my boss, I don't like what I'm doing,.
but until I find a new job, I'm not gonna quit..
I'm just gonna sit here and collect a paycheck,.
do minimally what I have to do,.

$^{321}$and that is the quiet quit..
It has affected so many corporations at this time..
I'm not gonna be dumb enough to quit.
and not get the paycheck, but I will get the paycheck..
I will do as minimal as I can..
You tell me what to do, I will do it..
I am not gonna go above and beyond..
I am not gonna do anything more,.
and anything that you ask me that's not part of my job,.
I'm not going to do it..
My whole goal is to collect the paycheck..
I wonder how many of us approach.
the Christian life like this,.
that we approach it with this quiet quit mentality,.
where we sit there and we sit and we said,.
"Hey, look, we're gonna go through the routines..
"We're gonna come on Sunday, maybe once in a while,.
"if I have nothing to do,.
"I will attend a fellowship or a Bible study,.
"but I'm not gonna go above and beyond.".
I bet you too many of us will probably not sit there.
and say, "We're gonna give up Christianity..
"God is wrong, I don't believe it..
"Now, I don't think we're gonna go that far.".
But we go in a state of what we call lukewarmness,.
where we sit there, we go through the motions,.
we attend service, and then we're done,.
and we say we're still a Christian.
and we attend the church and we're religious..
I bet you that's where many of us are..
That is the quiet quit, that is the lukewarmness.
that Revelation warns us about..
Now, here's a quick question is,.
why do we end up in that place?.
What are the motives that causes us to go there?.
Let me actually go through some of those things..
Maybe sometimes what causes us to go through the motions.
is churches become a routine..
Maybe some of us have grown up in the church.
all our lives and all that other stuff,.

$^{361}$and ever since we're younger, a lot of times,.
we've gone through the same thing..
We've lost that personal attraction,.
we've lost that personal motivation,.
and it has become a routine..
Maybe for some of us,.
we've grown disenchanted by the church..
We've grown frustrated by the church..
Christianity is not seen in a very good light these days..
Because of politics, because of race,.
because of how pastors act or ministry leaders act,.
we're not always seen..
We're seen by the world as haters..
We don't love other people, we don't care about other people,.
and we all make it about building our own empires..
Christianity today is not seen in a positive light,.
and that's why sometimes you get frustrated.
and disenchanted by the church..
It's not routine, we're disenchanted..
Maybe a lot of times you're just distracted.
by what the world has to offer..
And you're sitting there, "Hey, look,.
"I'll do this church thing, I'll come in,.
"but at the same time, I gotta worry about money,.
"I gotta worry about career,.
"I gotta worry about all these things..
"But I'll come on Sundays, I'll give you that.".
And we live in a state of going through the routines..
We live in the state.
of what Revelation's called being lukewarm..
For us as Christians, let me give you.
a couple of signs of being lukewarm..
Let me give you a couple fast ones.
and see if any of these applies to you..
Signs of being lukewarm,.
you have an inconsistent and diminished prayer life..
You struggle with having regular communications..
You struggle with genuine engagement with the Lord..
If you do pray, it's, "Lord," you go through rhythms.
and you have a script or everything like that,.

$^{401}$but you've failed like what Elijah said,.
of wrestling with the Lord, crying out to the Lord,.
talking with the Lord..
An inconsistent and diminished prayer life..
Signs of being lukewarm also shows.
that you have a worldly mindset and values..
You sit there and say, "Lord, I totally understand.
"what you're saying, but Lord,.
"look at all these different things..
"Look at the need to have money and power and goals.
"and all of these different things.".
I literally was talking to a woman.
who lives in Australia, in Perth, Australia,.
and she works at Deloitte,.
and she's literally gonna be a partner over at Deloitte,.
but she works from 8 a.m. to 9 p.m..
every single day in Australia..
And she's a Chinese-Indonesian,.
and she says to me, "Tommy, I'm trying to figure out.
"that as I get promoted to be partner,.
"how do I still maintain a vibrancy.
"in my walk with the Lord versus getting caught up?.
"I'm already working from 8 a.m. to 9, 10 p.m..
"How do I continue to maintain a diligent,.
"a godly mindset, and never about going up.
"the corporate ladder?.
"What does it mean to have a worldly mindset and values?".
We find ourselves struggling with all of that..
Signs of being lukewarm is we neglect.
our spiritual disciplines, the importance of Bible reading,.
meditation, fasting, serving others..
Other signs, lack of convictions and moral compromise.
about our character, about doing what is right,.
about humility, because the things of God matters to us..
And when we struggle with lukewarmness,.
ah, yeah, sure, sure, fine..
But what it means to serve others,.
to die to ourselves every single day,.
because when you're a motivated Christian,.
you're constantly asking those questions.

$^{441}$each and every single day..
When we struggle with lukewarmness, we struggle..
We have a lack of Christian community,.
because we just don't see the importance of that anymore..
We see ourselves a lot of times just sitting there,.
and we don't need to share,.
we don't need to do anything accountable,.
we just do whatever we want..
We don't see the importance of fellowship,.
of people holding us accountable..
And for those who are sitting on the screens at home,.
and for those who are outside in the country,.
I totally understand, but for those who are in Chicago.
sitting there, you are the one.
that may be tempted to do this,.
'cause you're sitting at home and says,.
I will just watch, I'll attend virtual church,.
but I don't need any community,.
'cause I don't need anymore,.
I don't need anyone accountable..
When we live in a state of lukewarmness,.
when we're sitting there stuck,.
stuck in being lukewarm,.
and limping along, we don't see the need for community..
Are you stuck, are you limping along in your faith?.
And the second thing it tells us.
in terms of the verse is this,.
if the Lord is God, follow him,.
but if it is Baal, then follow him..
Make a choice, make a decision..
Will my gods accomplish what I want them to do?.
Because if it is not fully focused on our God.
and our Savior, we all have other gods..
Now I'm saying that you have statues.
and everything like that,.
but the gospel and the Bible tells us.
that we have gods or we have built altars.
in our hearts that distract us away.
from focusing on one true God..
For the Israelites at that time,.

$^{481}$Baal was the god of power and sexual potency..
Asherah was a god of fertility,.
and for them, the Israelites,.
and all the people in Israel,.
and King Ahab himself was so focused.
on sex, power, harvest, wealth,.
that became their god..
So if God is not your god,.
then what are the gods that you're chasing?.
If God is not your god,.
then is it wealth that you're chasing?.
Is it family that you're chasing?.
Is it relationships that you're chasing?.
Is it goals or dreams or comfort that you're chasing?.
And if it is, then you better take a hard look.
and ask yourself, if I chase after those things,.
will my god, will that god,.
money, wealth, family, relationship,.
will that god accomplish what I want them to do?.
And if the answer is yes in your mind,.
then there's no need for you to be here listening at all,.
because you've made up your mind,.
so no matter what I say, it won't matter to you..
But you better make sure that will accomplish.
what you need them to do..
Have you ever realized that some of the most unhappy people.
in life are wealthy people?.
And many wealthy people who spent their entire life.
just attaining and getting wealthy.
and finally they get to the end of it,.
and they have all the money that they want,.
they're lonely and they're realizing this is it?.
You have to ask yourself,.
will that god lead me to where I want to be?.
Now, if the answer is no,.
and you want to begin to explore the true god.
in the form of God and Jesus Christ,.
and see where that path leads you,.
then it goes to the third thing.
that we're gonna look at the question..

$^{521}$If the Lord is God, follow him..
But if thou, then follow him..
And the last part of 21,.
and the people did not answer him a word..
Third question..
If you're stuck and you're looking through your faith,.
and you realize that the god that I'm chasing.
is not gonna get me to where I want to be,.
and I want to begin to understand what the true God.
or the form of Jesus Christ and God helps me to do,.
are you going to do anything about it?.
What are you gonna do?.
What steps are you gonna take to do it?.
If you realize, hey, look, the god that I'm pursuing,.
the money and all that stuff, that's not the answer..
Lord, I want to pursue you..
What are you going to do about it?.
Because a lot of times, too often,.
we sit there, when we're sitting there and lukewarm,.
we know that we should be doing something about it..
We just don't have the motivation to push ourselves.
to do it, and that's where we continue to get stuck..
I love what my friend Dave Evans,.
he's a professor at Stanford, says..
And he says, "These days, meaning is the new money.".
We're not at a place where people are motivated.
by earning more money and all that other stuff,.
that people in this generation.
want meaning in their life..
And if meaning is not found in money or power.
or relationship, and it's found in the true God.
in what scripture calls us, as we pursue God,.
we will find meaning, what are you going to do about it?.
Many of you guys may not be familiar with basketball,.
but you will understand this story..
Many of you guys did see the last dance with Michael Jordan..
And when I say Michael Jordan and Chicago Bulls,.
all of you guys are familiar with it..
They had six championships..
They were one of the greatest team.

$^{561}$in NBA basketball history..
They had the first three championships,.
and then Michael Jordan retired,.
and then he came back for the last three championships..
There was an interview, the former coach,.
Phil Jackson, did on a recent podcast I was listening to..
I was very, it just caught my attention..
Phil Jackson says to the team, he says,.
"When Michael came in, we lost in the championship,.
"the Eastern Conference Championship..
"And then the following year, we all came back..
"And I brought the team together,.
"and at that point in time, we brought in this player,.
"little crazy, his name was Dennis Rodman..
"He was a little nuts, but we brought him to the team,.
"and a lot of people did not like him..
"So I brought the team together.
"right before the season began..
"And I said to the team,.
"how bad do you want to win this championship?.
"Because if you want to win this championship,.
"you're gonna have to do something about it..
"And if you want to win this championship,.
"you're gonna have to accept people.
"who are different than you to be part of this team.".
And Phil Jackson says, I pointed to Dennis Rodman,.
and says, "You see this guy?.
"Dennis over here is a little strange..
"Dennis over here will probably go to Las Vegas.
"sometime during the season..
"Dennis over here will probably not show up at practice,.
"and he'll just do his own thing..
"But Dennis over here promised that at game time,.
"he will be here at every single game.
"and play his tail off to help you win the game..
"Can you accept someone different than you.
"to be part of the team?.
"Because you have to do that.
"if you want to win the championship.".
So for you and I as well, too,.

$^{601}$is if the money, the power, the family,.
the relationships mean nothing,.
and you're sitting there and says,.
"That's not gonna give me fulfillment or meaning..
"And pursuing God is gonna give me the meaning.".
And you're willing to try that?.
What are you gonna do to pursue him?.
What changes are you gonna do?.
For the people, they cannot give him an answer..
What about you?.
Elijah said to the people,.
"I, even I, only am left to profit the Lord..
"But by all his profits are 450 men," in verse 22..
"Let two bulls be given to us.
"and let them choose one bull for themselves.
"and cut it in pieces and lay on the wood,.
"but put no fire on it..
"And I will prepare the bull and lay on the wood.
"and put no fire to it..
"And you call upon the name of your God,.
"and I will call upon the name of the Lord..
"And the God who answers by fire, he is God.".
And all the people says, "It is well spoken..
"We love that question, we love that.".
Let's go to verse 30..
"Then Elijah said to all the people, 'Come near to me.'".
Because now the bells are doing all this,.
nothing happens, right?.
They sit there and they call upon their gods.
and all that stuff, nothing happens..
They start cutting themselves,.
they start dancing, nothing happens..
And Elijah says, "It is now my turn.".
Verse 30, Elijah said to all the people,.
"Come near to me.".
And all the people came near to him,.
and he repaired the altar of the Lord.
that has been thrown down..
Elijah took 12 stones,.
according to the number of the tribes of the sons of Jacob,.

$^{641}$to whom the word of the Lord came saying,.
"Israel shall be your name.".
And with the stones, he built an altar.
in the name of the Lord..
Couple of things,.
if you feel like you're limping your faith, number two,.
if you feel like your God.
is not gonna accomplish what you need,.
and you're willing to make a decision to change it,.
you're gonna have to repair and rebuild the altar..
Repair and rebuild the altar..
Let me quickly go through what that means,.
that when you and I repair and rebuild the altar.
in our lives, what we're saying is.
we put the priority back on God..
God, you are, I am rebuilding the altar in my life.
so that it is dedicated to you..
What does that tangibly mean?.
That means you're gonna have to reflect upon.
what expected my relationship with God..
You're gonna have to look at some of the barriers.
that can and will affect my relationship with my God..
You're gonna have to look at some of the issues.
in your heart..
You're also gonna have to ask yourself,.
who has God made me to be?.
What gives me meaning in my life?.
And who do I need to help in my process?.
The altars in our life are the altars.
that God places in your heart..
We can always fake actions..
We can't fake what's in our hearts..
In order to rebuild the altar in our life,.
you will have to ask yourself,.
where is the state of my heart?.
What barriers are in my heart?.
What are the barriers that are affecting.
my relation with God?.
What are the barriers that are killing,.
or the issues that are affecting my heart?.

$^{681}$Who has God made me to be?.
What gives me meaning in my life?.
And you're gonna have to tangibly do all of that..
And he repairs the altar,.
and then he makes a trench above the altar.
as great as wood contain two seeds of seed..
And Elijah put the wood in order.
and cut the bowl in pieces and laid it on the wood,.
and he says, "Fill four jars with water.
"and pour it on a burnt offering.".
And on the wood..
And he says, "Do it a second time.".
And they did it a second time..
And he says, "Do it a third time.".
And they did it a third time..
And the water ran around the altar.
and filled the trench also with water..
Here's my last point,.
and we're gonna look at Elijah's prayer..
At the time of the offering of the oblation,.
Elijah the prophet came near and said,.
and here's Elijah's prayer,.
"O Lord God of Israel, Isaac and Israel,.
"let it be known this day that you are God in Israel.
"and that I am your servant.
"and that I have done all these things at your word..
"Answer me, O Lord, answer me,.
"that this people may know that you, O Lord, are God.
"and that you have turned their hearts back.".
Fifth question, we've gone through four questions..
Are you limping through your faith?.
And if you are, what are the certain gods.
in your life that you have?.
Do you know if those gods are gonna accomplish.
what you need, do what you want?.
And if it's not, what's your alternative?.
Is it found in the Lord your God, in Jesus Christ?.
Thirdly, if it is found in the Lord Christ,.
what are you gonna do about it?.
Then from that point of forth,.

$^{721}$how do you rebuild your altars?.
How do you repair your altars?.
And the last one, what does the prayer of Elijah teach me?.
Prayer of Elijah teaches me three things..
O Lord God of Abraham, Isaac and Israel,.
let it be known this day that you are God in Israel..
Prayer of Elijah teaches us that God is supreme..
God is the God of Abraham, God is the God of Isaac,.
God is the God of Israel..
Let it be known this day that you are God,.
you are the one true God..
Second thing that the prayer of Elijah takes place.
is that I am your servant..
It shows us our position in relation to God,.
that he is, we are his servant,.
and that I have done all these things at your word..
And the last one, answer me, O Lord,.
answer me that this people may know that you, O Lord,.
are the God that you have turned their hearts back..
Look at the key verse in this one.
as we wrap up this one too..
Look at the key thing..
It's not that, Lord, you're God and I am your servant,.
and force the people to turn back on him..
No, no, no, no..
Lord, work in it that you are the one.
that turns their hearts back to you..
It is not the people turning their hearts..
Lord, it is you who turns their heart back..
Then the fire of the Lord fell.
and consumed the burnt offering,.
and the wood and the stones and the dust,.
and licked up the water that was in the trench..
And when all the people saw it,.
they fell on their faces and said,.
the Lord, he is God, the Lord, he is God..
What does this story remind you of?.
Does it remind you of the gospel?.
Does it remind you of the sacrifice of Jesus?.
That God could have easily punished all humanity.

$^{761}$because of their sin, and the judgment of sin,.
he could have literally leveled the fire.
and brought his entire judgment on the people..
But he doesn't bring his judgment on the people..
He brings his judgment on the sacrifice..
And it is the sacrifice who takes the blunt of the flame..
For us, as we look at the cross of Jesus,.
it is not humanity and those who have sinned.
that is the recipient of God's anger and God's judgment,.
but rather he sends Jesus who becomes the sacrifice,.
and he sacrifices the sacrifice..
And through all that, through the sacrifice,.
it turns our hearts back to him..
That's the gospel message..
That's what the prayer of Elijah teaches us,.
a foreshadowing of Christ..
Some of us here today may be stuck.
in limping through our faith..
Some of us here may be wrestling through.
what are the true gods in our life..
We may be living this state of mediocrity,.
of lukewarmness..
We go through the flow, we go through the motions.
of coming to church, but there's things in our life.
that's tugging in our life..
You're gonna have to make a decision..
Who's gonna be the god of your life?.
And if Jesus is going to be the focal point,.
the god in your life, then you're gonna have.
to make a decision..
What are you gonna do about it?.
Are you gonna repair and rebuild the altar in your life?.
Are you gonna sit there and reflect upon.
what God's meaning in your life,.
what God's purpose in your life is?.
Are you gonna look at some of the barriers in your heart?.
And are you gonna understand and reflect upon.
the road that God tells us through Elijah,.
that God is supreme in our life,.
that we are a servant of Christ,.

$^{801}$and it is the gospel that draws us over to him..
That's what 1 Kings 18 teaches us..
That it is the gospel, as God draws us to him,.
we are his servant..
We are Elijah, continue to proclaim him as a servant..
Let me pray..
Father, thank you for your works..
Thank you for the, man, the life of Elijah.
as we study so much in chapter 17 to 18..
Thank you for how even through Elijah,.
you spoke and showed Jesus Christ through us,.
that you brought the sacrifice and your judgment.
was not placed on humanity, but placed on the sacrifice..
And it is because of that sacrifice, we are drawn to you..
May we continue to walk in your ways,.
especially those of us who are sitting right here,.
limping through our faith..
In Christ we pray, amen..
(gentle music).
\newpage



\section{12:1-2}
\label{sec:NGEVacpuPjM}
\textbf{The Perfect Sacrifice (Romans 12\_1-2) - Joyce Tam}
\newline
\newline
連結: \href{https://youtube.com/watch?v=NGEVacpuPjM}{\texttt{https://youtube.com/watch?v=NGEVacpuPjM}} ~~~~ 語音日期: 2023-07-11
\newline
\newline
\hyperref[sec:hjFzQBCexic]{\small{< < < PREV SERMON < < <}}
~
\hyperref[sec:index_chronic]{\small{[返順時目]}}
~
\hyperref[sec:index_scriptual]{\small{[返順卷目]}}
~
\hyperref[sec:bp9nD3hO7O0]{\small{> > > NEXT SERMON > > >}}
\newline
\newline
$^{1}$Okay. Okay..
Before we begin, I'd like to ask everyone a question..
Do any of you pray for the brothers and sisters.
who serve on Sundays, especially this weekend?.
If you do, I want to thank you..
Because every time I give a sermon,.
I am filled with fear and anxiety..
Let alone today I'm preaching in English..
Lord, have mercy on me..
Today I hope to discuss Romans chapter 12, 1 to 2 with you..
And I encourage everyone to reflect more deeply on the message within..
Let's pray..
Dear Lord, our God, we thank and praise you..
We dedicate this hour to you..
Anoint the lips of your child..
Let your words be spoken with power and authority..
May the Holy Spirit work in people's heart as they listen to you..
In Jesus' name we pray. Amen..
The book of Romans is a very important book in the Bible..
Well, you might say, like, every book in the Bible is important, right?.
By important, I mean that the book of Romans.
has had a profound impact on generations of theologians..
The theologian and philosopher Augustine in 386 AD.
repented after reading Romans 13, 13 to 14,.
and later became a renowned preacher..
His writings still carry great weight today..
The theologian and religious reformer Martin Luther in November 1515.
taught the main point of Romans is not about how a righteous God punishes the unrighteous,.
rather that God declares people righteous through their faith in him..
This doctrine was the driving force behind his reform movement..
Before we dive into Romans 12, 1 to 2,.
let me briefly summarize what was discussed in Romans chapter 1 to 11..
In summary, Paul explained the basic teachings of Christianity,.
emphasizing that both Jews, who seem to have a higher status and were proud,.
and Gentiles, who seem to have a lower social status,.
were both God's people justified through faith in Christ,.
not by observing the Old Testament law or rituals..
Now, as we come to Romans 12, verses 1 to 2,.
serve as a turning point from the previous chapters..
The apostle Paul transitions from explaining the true meaning of the gospel.

$^{41}$to teaching believers how to live out their justified lives through faith in their daily conduct..
Title for today, "The Perfect Sacrifice.".
Let's start by reciting Romans 12, 1 to 2..
"I appear to you, therefore, brothers, by the mercies of God,.
to present your bodies as a living sacrifice,.
holy and acceptable to God, which is your spiritual worship..
Do not be conformed to this world, but be transformed by the renewal of your mind,.
that by testing you may discern what is the will of God,.
what is good and acceptable and perfect.".
When I read this passage, I could easily relate it to how Paul felt at the time..
Since 2008, our family, along with many brother and sister,.
have gone on mission trips to China..
We almost go every summer..
But we haven't been able to go since 2019..
But my heart always yearns to see my brothers and sister in China..
I always think about, like, how they are doing..
I really hope to have the opportunity again to visit them,.
ask them about their well-being, their gatherings,.
and whether they are still holding on to their faith..
Similarly, in the beginning of chapter 12,.
the apostle Paul, out of godly compassion and love for the brothers and sisters,.
Paul encourages believers on how to practically live out life as a Christian..
Because the church was divided by their differences,.
they forgot they are following the same God..
In Ephesians 4.1, it says,.
"I therefore, a prisoner for the Lord,.
urge you to walk in the manner worthy of the calling to which you have been called.".
The teachings in the Bible are not just meant for reading or acquiring knowledge..
The ultimate purpose of these teachings is to know God and to put His words into practice,.
to present your bodies as a living sacrifice..
In the Old Testament, sacrifices were generally about selecting, preparing, offering,.
and literally sacrificing animals to please God in religious rituals..
There were various types of sacrifices, such as burnt offering,.
grain offering, peace offering, sin offerings, and trespass offerings,.
each with its own preparations, certain elements, rules, procedures,.
and they each had their specific purpose, all of which are necessary..
However, what Paul mentions here is different..
It is a living sacrifice..
Before we talk about living sacrifice,.
I want to first talk about the bodies mentioned in the Scripture..

$^{81}$The word bodies in the original text refers to the whole person..
This includes the body, the mind, spirit, and thoughts..
Furthermore, the word bodies is not one physical body, but rather many bodies,.
a community of faith, a community of faith offering themselves as one sacrifice..
When a person believes in Jesus, it is indeed an act of offering..
A person who accepts Jesus as their Lord Savior,.
acknowledge his or her sins, surrenders himself before God,.
declares that their life no longer belongs to themselves..
They are no longer their own masters,.
and they ultimately choose to recognize Jesus as his or her own personal Savior..
From then on, I no longer live, but Jesus lives in me..
This declaration is a response to the Gospel..
Romans 8:11, "If the spirit of him who raised Jesus from the dead dwells in you,.
he who raised Christ Jesus from the dead will also give life to your mortal bodies.
through his spirit who dwells in you.".
To offer our bodies as a living sacrifice,.
the sacrifice mentioned in Romans 12 is made possible.
not by whether a person is willing or not to offer themselves,.
but by the prerequisite that Jesus first offer himself as the perfect sacrifice,.
accomplishing salvation..
It is the precious blood of Jesus that cleanses us from all sins,.
enables people to be born again, become holy, and be called righteous..
It is because Jesus first offered himself as a sacrifice.
so that people can become a living sacrifice..
Hebrew 1010, "And by that will we have been sanctified.
through the offering of the body of Jesus Christ once for all.".
Reflecting back on the days of my college years,.
when I memorized Romans 12 1-2,.
I was only thinking about how mission-driven I felt.
and how determined I wanted to be a good Christian,.
offering myself as a sacrifice..
The motivation behind offering myself came from my own selfishness..
It makes me feel good to think this way..
However, I didn't truly understand the focus is not on me..
It is not about how much I love God.
and willing to offer myself as a living sacrifice,.
but on how much God loved me that he sent his one and only son, Jesus Christ,.
to be the sacrificial lamb so that I can live..
To offer our bodies as a living sacrifice is a true and proper worship.
because it is a proper response to Jesus' sacrificial offering..

$^{121}$For those who follow the law, it is the most natural response to the gospel,.
a response of gratitude to God,.
and a reasonable response of someone who understands the grace once received..
Verse 2..
"Do not be conformed to this world, but be transformed by the renewal of your mind.".
This present age has its own thoughts and lifestyle,.
which believers who have received God's grace must discern..
Therefore, the meaning of "do not be conformed to this world".
is that we must be cautious and not let these trends and values worldview shape our lives..
Recently, during a gathering, a friend mentioned that because of his two daughters,.
he often watches Disney Channel..
Lately, he noticed an increased portrayal of non-traditional family relationships in the program..
According to him, the portrayal of these families and their interaction in these programs.
are no different from traditional families,.
except for the fact that they don't necessarily have a father and a mother..
They seem to subtly normalize non-traditional relationships..
A few years ago, in Canada, a 46-year-old father with seven children.
discovered that he was actually a woman..
His wife could not tolerate his lifestyle, so he left his wife and his seven children..
After leaving his family, this 6'4" tall man with his hair tied up and wearing a custom-made dress.
transformed himself into an 8-year-old girl..
He advertised himself online and sought adoption..
Not having to wait long, he was adopted by an elderly couple..
He now lived with them and behaved like one of their grandchildren..
He and her grandchildren spent time coloring, playing with the toy kitchen, playing toys..
One day, one of the grandchildren of this elderly couple wanted to be an older sister..
So she asked this 46-year-old man to be two years younger..
So this 6'4" tall little girl changed his age again,.
played the role of a younger sibling, and settled at the age of six..
I remember a long time ago when I was picking up my daughter from school..
I saw a little girl waiting for her parents..
I casually asked her, "Your mom is coming to pick you up?".
She looked at me and did not answer..
Later, I found out that she didn't have a mom because she has two fathers..
When I returned home, I thought for a long time..
I realized that I should never assume that every child's family has a father and a mother.
because there are so many combinations nowadays,.
such as two fathers, two mothers, one father, multiple mothers,.
one father, one mother, or no parent at all..
This is the current situation of our world's family relationships..

$^{161}$I don't know how many different combinations will emerge in the future..
And I often wonder, where does the next generation's concept of family stand?.
What kind of world we are living in right now?.
As a Christian, how would you respond?.
As a Christian parent, how would you address the questions your children ask?.
If you are still solely focused on how to push your children to get into a good school,.
studying advanced courses at an early age,.
engaging in various activities like music, painting, chess, calligraphy,.
eating organic food, don't eat any fast food,.
I think you need to reconsider..
I'm not trying to scare you, for I am a mother myself,.
but I genuinely want to know,.
when was the last time you considered your child's character?.
The world urges us to put all this effort into making our child the smartest,.
the most successful, most high-achieving individual..
But sadly, we turn a blind eye to their heart..
As Christian parents, we need to pray..
We need to seek God for wisdom..
We need to listen. We need to reconsider..
Sometimes when my daughter talks to me,.
she brings up topics that I never thought about before..
And I don't always know the answer or know how to respond..
In those moments, I pray in my heart, asking the Lord for wisdom,.
hoping at least not to say anything wrong..
Most of the time, I'm speechless..
I'm just there to listen and try to understand more..
I don't know if you have watched the UV YouTube channel recently,.
for I have done an interview with Pastor Julie when I was in Hong Kong..
If you haven't watched it, I encourage you to watch it if you understand Cantonese..
In the interview, Pastor Julie provided a few reminders..
And one of them was that Christians are the people of the book..
Pastor Julie said she has a desire to teach lay people Hebrew.
and make theological education more accessible.
so that Christians are equipped to face future challenges..
She hoped to help believers deepen their study of the Bible.
and even be knowledgeable enough to feed those around them..
2 Timothy 3:16 says, "All scripture is breathed out by God.
and profitable for teaching, for reproof, for correction,.
and for training in righteousness.".
Teaching shows you the path to walk on..

$^{201}$Reproof shows you where you have gotten off the path..
Correction shows you how to get back on the path..
Training in righteousness shows you how to stay on the path..
The foundation of Christian faith is God's Word..
Christian parents have asked me how we can respond to current issues in various matters..
Unfortunately, we have to admit that we just don't have all the answers..
So how can we teach our children in this day and age?.
My brothers and sisters, when facing the uncertain future.
and the challenges of this world,.
remember that only God's Word is your greatest help..
When you read the Bible, you may not understand or feel its significance..
But day by day as you read, study the Bible,.
you are accumulating your knowledge..
You are allowing God to speak to you,.
allowing the Holy Spirit to transform your heart..
And sure, you will slowly store up ammunition for the battle..
Only by drawing near to God,.
allowing the Holy Spirit to renew and transform our hearts,.
give us the ability to discern good and evil..
Can we stand firm in this turbulent and unsettled world?.
I remember a scene from a drama series I watched when I was young..
A person asked his jewel of friends how to distinguish between real and fake diamonds..
The person explained that when he was young,.
his grandfather would let him play with real diamonds..
As he grew older, his grandfather added fake diamonds into the mix,.
and he could immediately tell the difference..
Do you have the ability to discern the truth, brothers and sisters?.
For those of you who have believed in the law for maybe two, four, or more than ten years,.
how well do you grasp the truth of the Bible?.
How much time have you set aside in your daily life in God's Word?.
Romans 12, 1-2..
I appeal to you, fearful brothers,.
by the mercies of God to present your bodies as a living sacrifice,.
holy and acceptable to God, which is your spiritual worship..
Do not be conformed to this world, but be transformed by the renewal of your mind,.
that by testing you may discern what is the will of God,.
what is good and acceptable and perfect..
The perfect sacrifice..
Jesus has already offered himself as the perfect sacrifice,.
accomplishing it once and for all..

$^{241}$How will you respond?.
How should we as brother and sister, a community of faith, respond together?.
A living sacrifice refers to the regenerate and save life.
that is willing to live a life worthy of salvation received,.
offering oneself fully in response to the grace of salvation that we do not deserve..
This is indeed the most reasonable and appropriate response to the gospel we have received..
As Christian, we need to hold on to God tightly and strive to study his Word diligently,.
for it is good through God's Word that our minds are renewed,.
giving us the ability to discern what is good and pleasing to him,.
that by testing you may discern what is the will of God,.
what is good and acceptable and perfect..
Let us bow our heads..
Lord, we thank you for loving us and for the sacrifice of Jesus Christ..
You know us better than anyone else..
Please continue to guide our lives and help us not to forget what you have done for us..
May we respond to your love by offering ourselves as living sacrifices..
May the Holy Spirit remind and help us to root ourselves in your Word,.
so that we may stand firm in this turbulent and restless age..
I also want to pray specifically for our brothers and sisters who are struggling and feeling weak,.
that those who are struggling may find a way out,.
and those who are weak may find your companionship and presence..
For those who are sad and burdened, may you grant them wisdom to face their challenges..
For those who are wounded and sorrowful, may you embrace them and provide comfort..
Lord, please strengthen us so that the more difficult, chaotic, and oppressive our lives become,.
the stronger our hearts will be in seeking strength from your Word..
As we, the church body, come before you, that we may be a living sacrifice that pleases you..
Thank you, Lord, for you will never abandon those who follow you..
And you are a God who listens to prayer..
We humbly come to you with gratitude..
In the name of the Lord Jesus Christ, we pray. Amen..
[BLANK AUDIO].
\newpage



\section{羅馬書 12:3-8}
\label{sec:bp9nD3hO7O0}
\textbf{生命的賜(羅馬書 12\_3-8)- 譚文倩莊傳道}
\newline
\newline
連結: \href{https://youtube.com/watch?v=bp9nD3hO7O0}{\texttt{https://youtube.com/watch?v=bp9nD3hO7O0}} ~~~~ 語音日期: 2023-07-11
\newline
\newline
\hyperref[sec:NGEVacpuPjM]{\small{< < < PREV SERMON < < <}}
~
\hyperref[sec:index_chronic]{\small{[返順時目]}}
~
\hyperref[sec:index_scriptual]{\small{[返順卷目]}}
~
\hyperref[sec:hrYJaJP5Fzk]{\small{> > > NEXT SERMON > > >}}
\newline
\newline
羅馬書 12:3-8
\newline
\begin{longtable}{cl}
\hline
\hline
章節 & 經文 (和合本修訂版)\\
\hline
12:3 & \begin{tabularx}{0.7\textwidth}{X} 我憑著所賜我的恩對你們每一位說：不要把自己看得太高，要照著神所分給各人的信心來衡量，看得合乎中道。 \end{tabularx} \\ \\ \relax
12:4 & \begin{tabularx}{0.7\textwidth}{X} 正如我們一個身子上有好些肢體，肢體也不都有一樣的用處。 \end{tabularx} \\ \\ \relax
12:5 & \begin{tabularx}{0.7\textwidth}{X} 這樣，我們許多人在基督裡是一個身體，互相聯絡作肢體。 \end{tabularx} \\ \\ \relax
12:6 & \begin{tabularx}{0.7\textwidth}{X} 按著所得的恩典，我們各有不同的恩賜：或說預言，要按著信心的程度說預言； \end{tabularx} \\ \\ \relax
12:7 & \begin{tabularx}{0.7\textwidth}{X} 或服事的，要專一服事；或教導的，要專一教導； \end{tabularx} \\ \\ \relax
12:8 & \begin{tabularx}{0.7\textwidth}{X} 或勸勉的，要專一勸勉；施捨的，要誠實；治理的，要殷勤；憐憫人的，要樂意。 \end{tabularx} \\ \\
[1ex]
\hline
\hline
\end{longtable}
$^{1}$上星期Mark Sir帶大家去耶路撒冷大衛城的希西加隧道.
今天我帶大家上水面.
重返地球.
去看羅馬書十二章三至八節.
是羅馬書系列第二講.
生命的刺.
講屬靈的恩賜.
先找書吧.
上次跟大家講羅馬書十二章一至二節.
生命的制.
保羅以神的慈悲勸勉弟兄姊妹.
將原本已經死的身體.
因為耶穌的犧牲.
活過來的生命立志獻上.
如此是奉乃是理所當然.
因為這是反映人真正明白福音.
合乎情理的回應.
基督徒要小心.
不要讓這個時代去塑造我們的生命.
求聖靈幫助我們心意根深而變化.
有淑齡的眼光去明白神的心意.
成就神的旨意.
今天我們繼續羅馬書十二章一至二節.
接下來的三至八節主題.
生命的刺恩賜.
我們不如先一起去讀這個經文一次.
羅馬書十二章三至八節.
我憑著我所賜我的恩對你們每一位說.
不要把自己看得太高.
要照著神所分給各人的信心來衡量.
看得合乎中度.
正如我們一個身子上有好些肢體.
肢體也不都有一樣的用處.
這樣我們許多人在基督裡是一個身體.
互相聯絡作肢體.
按著所得的恩典.
我們各有不同的恩賜.
或說預言.
要按著信心的程度說預言.
或服事的.

$^{41}$要專一服事.
或教導的.
要專一教導.
或勸勉的.
要專一勸勉.
施捨的.
要誠實.
自理的.
要殷勤.
憐憫人的.
要樂意.
我們先低頭祈禱.
親愛的主我們的神.
聖經都是你所默示與教訓.
督責.
使人歸正.
教導人學義.
求你的話語帶領人走要走的路.
知道我們走歪了.
走回一條對的路.
引導我們不會再走錯.
你似因讓跟隨你的人.
生命建基在你的話語上.
祈禱交托.
奉主耶穌基督的聖名祈求.
阿們.
羅馬書十二章三節.
那裡這樣說.
我憑著所賜我的恩.
對你們每一位說.
不要把自己看得過高.
要照著神所分給各人的信心來衡量.
看得合乎忠道.
前幾年我回香港.
在我住的地方.
我特意去了樓下的街市.
去買兩條藤條.
其實藤條對我來說都挺有感情的.
陪了我很多年.
我有規定的.

$^{81}$那條藤條要粗粗的.
如果你去過市場買藤條.
你知道有些是彎的.
我找的是要直的.
我盡量找一條粗粗的要直的.
最好藤條尾那裡.
有條膠索帶.
方便我把它掛起.
打者愛也.
想當年被人打的時候.
打我那位仁兄說.
我也是為你好.
因為愛錫你才打你.
我在寫講章的時候.
我想起藤條背後的愛.
曲調雖異 功妙雖同.
第三節這裡.
保羅的教導勸勉.
教會裡面屬靈恩賜.
配搭之前.
要先提到.
人要藉著自己所得著的恩典.
然後才講之後那幾節經文.
先打好底.
起樓打好樁.
然後才一層一層地砌上去.
記得《羅馬書》十二章一開始.
第一節提到保羅是以神的慈悲.
而這裡是講藉著神的恩典.
作出勸勉教導.
根據神學家John Piper的解釋.
恩典是在講神對罪人的良善.
而恩賜 憐憫.
是在講神對受苦的人.
受苦的人的良善.
憑著我所賜的恩.
保羅根據《使徒行傳九章》記載.
他是一個迫害基督徒的法利賽人.
在大馬色路上遇見耶穌.
得蒙恩典.

$^{121}$神既往不咎呼召他跟隨.
神既往不咎叫保羅這個罪人.
時時存著感恩的心.
轉化成一個活著不再是我.
乃是基督在我裡面活著的門徒.
在羅馬書12章3至8節的第三節一開節.
提到他所承受的恩典.
就是要提醒信徒要像他一樣.
先要記得你今時今日所有的事.
包括你的恩賜.
都是神所賜的.
這個認知驅使人.
將自己放在一個正確的位置.
最近我們教會有一個茶經.
婦女茶經叫做.
正正經經茶茶經.
打屁股我竟然說不出來.
我們每一次茶考的時候.
我們都會至少讀三個不同的版本.
今天我也很想和你們讀兩個不同的版本.
《新漢語聖經》羅馬書12章3節.
我藉著神賜我的恩典.
對你們各人說.
不要對自己評價過高.
超過所應有的評價.
都要按著神量給各人的信心.
忠肯地評價自己.
不如問問我們的同學.
這裡有哪一個生字是重複出現的.
評價.
不要對自己評價過高.
超過所應有的評價.
忠肯地評價自己.
再看多一個版本.
英文的.
《新美國標準》.
我藉著神賜的恩典.
對你們各人說.
不要對自己評價過高.
超過所應有的評價.

$^{161}$不如問問我們的同學.
在這篇文章中.
我們有幾個不同的版本.
我們會用三個不同的版本.
來評價自己.
我們會用三個不同的版本.
來評價自己.
我們會用三個不同的版本.
來評價自己.
我們會用三個不同的版本.
來評價自己.
我們會用三個不同的版本.
來評價自己.
十二章四至五節.
正如我們一個身子.
尚有好些肢體.
肢體也不都有一樣的用處.
且讓我們許多人.
在基督裡是一個身體.
互相聯絡.
作肢體.
我想這段經文.
大家都很熟悉.
甚至沒有這個概念.
在教會裡.
教會是一個整體.
一個身體.
我們每個人的不同部分.
是肢體.
要知道這段經文的背景.
當時的信徒有很大很大的分歧.
例如可不可以吃祭物.
守節旗.
而在這個群體裡.
亦有很多是外邦人.
是大多數.
我們可以利用當時在羅馬.
很流行的言詞.
去比喻教會肢體和身體.
身體是應該互相配搭.

$^{201}$不是互相攻擊.
我們再看另一段經文.
哥林多前書十二章十二節至二十七節.
就如身體是一個.
卻有許多肢體.
身體的肢體雖多.
仍是一個身體.
基督也是這樣.
無論是猶太人是希臘人.
是遺老的是自主的.
都從一位聖靈受洗成了一個身體.
並且共享這位聖靈.
身體原不只是一個肢體.
而是許多肢體.
例如腳說我不是手.
所以不屬於身體.
他不能因此就不屬於身體.
例如耳朵說我不是眼睛.
所以不屬於身體.
他也不能因此就不屬於身體.
假如全身是眼睛.
慶覺在哪裡呢.
假如全身是耳朵.
嗅覺在哪裡呢.
但現在神隨著自己的意思.
把肢體一一安置在身體上了.
假如全都是一個肢體.
身體在哪裡呢.
但現在肢體雖多.
身體還是一個.
眼睛不能對手說我用不著你.
頭也不能對腳說我用不著你.
不但如此.
身上的肢體人以為軟弱的.
更是不可缺少的.
身上的肢體.
我們認為不體面的.
越發給他加上體面.
我們不雅觀的.
越發裝飾得雅觀.

$^{241}$我們雅觀的肢體自然用不著裝飾.
但神配搭這身子.
把加倍的體面給那有缺陷的肢體.
免得身體不協調.
總要肢體彼此相顧.
假如一個肢體受苦.
所有的肢體就一同受苦.
假如一個肢體得榮光.
所有的肢體都一同快樂.
你們是基督的身體.
並且各自都是肢體.
都從一位聖靈受洗.
成了一個身體.
並且共享這位聖靈.
身體原不只是一個肢體.
而是許多肢體.
並且各自都是肢體.
保羅用這個比喻去強調.
一群信徒為了共同利益.
需要共同努力的必要性.
攀石專家Tommy Cogwell和Kevin Georgeson.
在2015年用了19天.
成功徒手攀爬全世界最困難的優山美地.
Yosemite National Park裡面.
有一個Captain.
裡面的Don War.
如果你們有機會去過加州.
去過這個地方.
我想你們知道我在說哪裡.
我看了他們的紀錄片.
其中一幕令我印象很深刻.
就是Kevin用了6天.
這個就是Kevin.
多次嘗試都過不了.
難度5.14D grade的Pitch 15.
還有32個Pitch.
我想我形容的Pitch好像是32個點.
如果你玩攀石.
或者你會知道我說甚麼.
相對他的拍檔Tommy.

$^{281}$不單止過了Pitch 15.
他還要16,17,18.
這樣一天一天向山頂進發.
多次又多次的失敗.
令Kevin很沮喪.
甚至懷疑他自己的能力.
他知道自己可能在拖累Tommy.
另一邊Tommy雖然一個Pitch一個Pitch.
可以繼續上去.
但當他知道原來Kevin還在15那裡.
不斷掙扎的時候.
他都很不好受.
在這6天裡面.
很多人都勸Tommy.
他說很明顯你的技術.
比起Kevin好很多.
他說不如你自己繼續上去.
去完成這個壯舉.
不要再等他了.
長話短說.
最後Tommy沒有因為他的技術比較強.
旁邊的人告訴他.
他比起Kevin更有能力去攻頂.
冒著他可能不能成為這個壯舉的結果.
回到Kevin身邊為他打氣.
結果前後6天.
你試想想,困在這塊石的這個位置6天.
才成功過了這個Pitch 15.
最後他們能夠一起完成.
徒手攀爬這個L-Captain Don Wall的壯舉.
一直看這個紀錄片的時候.
我一直在想.
若果教會裡面的弟兄姊妹.
都有Tommy和Kevin的團隊精神.
我們一起去為神的國去努力.
我們互相扶持前進.
當我們其中一個遇上低潮的時候.
強一些過一個願意放下自己的利益.
陪伴,鼓勵弱一些的那方.
不離不棄和他們一起同行.

$^{321}$教會又怎會不是一個很吸引的地方.
能夠成為一個彰顯神的地方.
基督女士一個身體.
互相聯絡,作肢體.
弟兄姊妹的相處.
特別很多時候我們一起事奉.
因為我們不同的性格.
處事的態度.
甚至是說話的態度.
競爭.
我們普遍都不願意承認自己的驕傲.
完全掩蓋了保羅在第三節所提到的.
我們每一個人都是一樣.
我們都是蒙受了神恩典的罪人.
事奉的時候.
其實重點並不是你喜不喜歡對方.
而是你看不看到神要藉著我們的不同.
因為他的關係.
我們有同一個目標.
所以我們可以一起同工.
彰顯救恩的大能.
成就神的旨意.
就像我剛才所說.
我們大家都很熟悉教會是一個整體.
一個身體的概念.
很清楚我們身體裡的不同部分.
肢體.
但很多時候聖經的道理.
我們真的只有一個肢.
我們都不願意對付自己.
因為實際上.
今天活著的仍然是我.
仍然是我在我裡面活著.
其實我們會不會和當日保羅所怨恨的那群信徒分別不大呢.
十二章六至八節.
按著所得的恩典.
我們各有不同的恩慈或說預言.
要按著信心的程度說預言.
或服事的要專一服事.
或教導的要專一教導.

$^{361}$或勸勉的要專一勸勉.
施捨的要誠實.
自理的要殷勤.
憐憫人的要樂意.
來到十二章六至八節.
保羅就更加提到七個恩慈如何運用.
或說預言是傳遞信息的恩慈.
能夠鼓勵,造就,安慰教會.
並且能夠對不認識神的人介紹認識神.
遵而敬拜神.
先知說預言這個恩慈.
很容易令人想起微卜先知.
但其實真正的意思是說一個傳遞信息的人.
一個messenger.
傳遞的信息裡面可能涉及將來的事情.
讓神的群體可以根據這些知識.
指導教會的使命.
或服事的.
服事基本上是指充當為人執行任務的中間人或傳遞人.
在四福音和使徒行傳當中.
事工服事經常是指為他人準備或提供膳食的人.
我想一說的時候.
大家馬上想起我們教會有一隊軍隊.
時時提供我們很多美食.
教導的在某種意義上.
其實所有信徒我們都應該互相教導.
但另一個角度來說.
並非所有人都是老師.
而保羅這裡更加是說一個專業的角色.
當保羅說上帝賦予他們傳道的能力的時候.
有教導恩慈的人.
就應該在傳講神的服事裡面.
使用這種能力.
賦予神給他們的恩慈.
說「勸勉的勸勉」是說.
擔任鼓勵,安慰,勸勉弟兄姊妹.
在某種情況下如何表現.
在生活中或甚至獨習.
施捨的.
施予,分享.

$^{401}$這個動詞不單是意味著要幫忙.
更加是分享訊息,智慧,贈送禮物.
或分享物資和金錢.
也可以是分享福音.
當中有慷慨,真誠的意思.
「自理的」.
「自理領導」是要謹慎地從事教導和祝福的活動.
要關心他人的利益.
最後一個「憐憫人的」回到前面的論點.
所有信徒都是上帝憐憫的對象.
這意味著我們要學習和敵人和好.
鑒於保祿在前面論證所提到憐憫的概念.
很難避免我們思想憐憫這個恩賜時.
是上帝給我們每一個信徒的禮物.
因此我們也要同樣去實踐憐憫別人的恩賜.
保祿在這六節經文中解釋了七個不同的恩賜.
當中「服侍,教導,勸勉」都在說要專一.
要專一服侍,專一教導,專一勸勉.
意思是在這些服侍中我們要專注.
實事求是,不要加鹽加醋.
「施捨」要誠實.
是在說贈送,分享的時候.
我們不僅要慷慨,也要真誠.
要動機單純.
我想起我生命中有些人.
每當他們幫別人或施與的時候.
他們真的不要求任何回報.
我很羨慕他們的心.
因為坦白說,我自己做不到.
我留意自己的時候.
我發覺我對別人好的時候.
其實我也希望別人稱讚我.
希望有回報.
但當我思想這個恩賜的時候.
我就在想,雖然我是這樣.
但不代表我要停下來.
我和我老公說笑.
至少我找一些我覺得他會回報.
或多謝我的人去做.
去實踐這個恩賜.

$^{441}$不要就這樣停下來.
「自理的要殷勤」.
是在說做領導的要謹慎.
要顧及其他人,要勤力.
憐憫別人的要樂意.
就是在說我們苦咒其他人的時候.
不僅要樂意.
是要帶著很歡喜快樂的心情.
甚至做完之後.
不要到處唱,或埋怨.
三節經文,七個恩賜.
保羅是否要在這裡.
要我們搞清楚我們有哪一個恩賜.
然後我們就各就位,對號入座.
這就叫做互相聯絡作肢體.
建立基督的身體.
和你們分享一個我自己的個人見證.
2008年,第一次去中山短孫.
我在福音隊裡面負責敬拜隊.
去之前的三年,我已經沒有大敬拜.
因為我媽媽去世之後.
我是處於一個情緒很低潮的地方.
我覺得我暫時要停下所有的事奉.
總之要休息.
當牧師找我去帶隊的時候.
其實我是很不願意的.
但又覺得沒得選擇.
那我繼續去.
我們的敬拜隊是由四個不同教會的弟兄姊妹組成.
第一次會面的時候.
牧師就介紹我認識和我一起大敬拜的一位姊妹.
她走進來的時候,黑衣黑褲,沒有笑容.
我當時在想,怎麼這麼不友善.
之後去中山之前,我們有多次的練習.
每次的練習都一點都不合適.
這件事從來沒有發生過.
我每次和別人練敬拜.
我覺得很容易,大家很有chemistry.
但不知為何總是和他不合適.
但我當時心想.

$^{481}$大家都去不同的教會.
將來都沒什麼機會見.
省下連朋友都不用做.
算了,去完短孫就算了.
幾個月之後,我們終於去了中山.
去到一個農村,第一場敬拜.
聚會三個小時前,已經開始有些村民來到.
我當時很驚訝.
你知道北美不遲到都很happy.
三個小時前坐在下面看聖經和看聖詩.
我當時已經覺得很感動.
時間到了,音樂一起進入敬拜.
無神神,我們整隊team的chemistry很棒.
真的很棒,很合適.
村民手舞足動.
音其實未必合適.
我們也有些外省人在當中.
我們當中有些是廣東話歌.
他們也不懂唱.
看到他們左右擺拍子也不對.
但這個敬拜卻是我多年來第一次被匯中敬拜的心.
反過來引領我大敬拜的一個敬拜.
當我很享受這個敬拜的時候.
同時我聽到神跟我說話.
祂跟我說.
不是你想不想大敬拜.
不是你喜不喜歡你隔壁的那位姊妹.
因為it was never about you.
我只需要你將屬於我的人帶到我的面前.
就在那一刻,我突然間很清楚知道自己的位置.
我們那次大敬拜,不單是坐滿了人.
整個教會是站滿了人.
唱到站在村口的人都聽得到.
聲音大得整個教會震動.
之後七場的敬拜我們都很合適.
那一年的報道會,我們做的報道工作有百多二百人信主.
之後有機會我再跟這個姊妹坐下來一起分享的時候.
很奇妙.
原來同一時間,神跟她說了同一番話.
原來我們各自都以為之後大家都不會做朋友.

$^{521}$想不到神在那個時候跟她說了同一番話.
再說一個笑話給你們聽.
一年之後我們再回到教會大敬拜.
教會有兩層.
上面沒有人站著,下面更爆.
原來那個教會是危樓.
所以我們唱歌的時候震動是真的.
是真的震動.
他們說過兩年我們應該拆了.
這個經歷從此扭轉了我大敬拜.
事奉的態度.
羅馬書十二章三節.
我藉著神賜我的恩典對你們各人說.
不要對自己評價過高.
超過所應有的評價.
都要按著神量給各人的信心.
忠肯地評價自己.
羅馬書十二章三至八節.
若果沒有了第三節.
就真的好像保羅只是叫人清楚知道自己的恩賜.
然後對號入座.
以至事奉能夠事半功倍.
建立基督的身體.
但是因為第三節.
重點不是只是鼓勵你找到自己的恩賜.
然後運用恩賜建立身體.
因為建立基督的身體.
運用恩賜的前設.
是說人要學孝主耶穌基督去謙卑自己.
放下自己.
信徒和基督復活的生命緊密相連.
而主基督耶穌服侍的榜樣.
其實你想一想.
從來都不是在說他喜不喜歡.
甚至他面對著被遺棄.
肉身受苦.
甚至不是在想千萬不要這麼辛苦去考量.
他只是專注他父神託付給他的使命.
去謙卑自己.
取了奴僕的形象去完成託付.

$^{561}$恩賜服侍也不是在說教會很需要幫忙.
請你回來看看你有什麼適合做的.
我感恩在2008年這個短宣的時候.
神要我明白我的恩賜.
我的服侍.
是要將我破壞了服侍的一些想法放下.
定睛在耶穌的榜樣裡謙卑下來.
和弟兄姊妹一起去同工.
因為當我這樣做的時候.
我就能夠看到神的大能.
我也能看到屬於神的人回到神的面前.
我們一起去蒙受神的恩典.
生命的詞.
求主更新我們的心.
我們對付自己的驕傲.
在祂的恩典裡.
我們時時都存著一個感恩的心.
察驗神的心意.
我們各按其職.
我們一起努力.
互補長短.
建立基督的身體.
被祝福見證神的大能.
我們一起低頭祈禱.
願意每一個跟隨你的人.
藉著運用他們的恩賜.
為你建立.
更認識你.
好像剛才的詩歌那樣說.
我們向著前方.
我們擁抱使命.
將福音廣傳出去.
祈禱奉主耶穌.
基督的聖名祈求.
阿們.
完主的平台.
\newpage



\section{}
\label{sec:hrYJaJP5Fzk}
\textbf{4/1/18 復活節崇拜話劇--轉機,轉機 Easter Sunday Skit}
\newline
\newline
連結: \href{https://youtube.com/watch?v=hrYJaJP5Fzk}{\texttt{https://youtube.com/watch?v=hrYJaJP5Fzk}} ~~~~ 語音日期: 2023-07-13
\newline
\newline
\hyperref[sec:bp9nD3hO7O0]{\small{< < < PREV SERMON < < <}}
~
\hyperref[sec:index_chronic]{\small{[返順時目]}}
~
\hyperref[sec:index_scriptual]{\small{[返順卷目]}}
~
\hyperref[sec:mbI7RGfwLmY]{\small{> > > NEXT SERMON > > >}}
\newline
\newline
$^{1}$[音樂].
是今天復活節所表達的那種訊息.
就是為了你在復活節當中的得勝.
得勝死亡的勝利.
得勝罪的權柄.
你去犧牲了自己.
我們為了十家獻上感恩.
我們為了第一個復活節的早上.
我們獻上感恩.
在這份墓裡復活的時候.
當你去表現出你得勝死亡的權能.
就是得勝死亡的權柄.
得勝罪的權柄.
去到得勝所有我們的破碎.
我們祈求你這個權能的彰顯.
去到改變我們.
去更新我們每一個.
我們祈求在基督裡面.
我們每一個新造的人.
當我們今天去遇見你自己復活的大能的時候.
我們在這個禮拜開始.
我們的生命都得到改變.
我們的心境會改變.
我們的行為會改變.
我們也會轉化.
我們也會去愛其他的人.
也會讓我們這樣愛他們的時候.
讓他們的生命得到更新.
我們在生命當中.
我們需要主理復活的能力.
藉著我們來接觸更多身邊的人.
我們來祈求你.
繼續在我們生命當中來作供.
改變我們以至我們能夠改變我們身邊的世界.
放耶穌基督的名頭.
這件是後悔.
我揹著後悔20年.
我真的去到哪裡我都揹著後悔.
這件是後悔.
我揹著後悔20年.

$^{41}$我真的去到哪裡我都揹著後悔.
這件是遺憾.
我揹著遺憾.
沒有10年.
都有8年.
不揹著它.
我周身不聚財.
這件呢.
這件是內疚.
究竟.
我何時開始.
帶著內疚呢.
我都不太記得了.
不過無所謂.
又不是重.
就帶著它.
到處去.
還有它們.
是痛恨.
和埋怨.
你知道的.
你當然知道.
我.
一向都是靠自己的.
反正.
我早就這樣了.
沒有人幫.
沒有人養.
沒有人關心.
沒有人陪.
和你說.
我想的.
我有的.
你聽才行.
神經病.
又是我自己和我自己說話.
我想怎樣.
我只想有多個人.
幫忙.
拿走這件行李.

$^{81}$我很累.
我真的很累.
這麼多年.
你都沒有照顧我.
有沒有.
我要求.
都不是我過份.
你讓我做.
你讓我做的事.
我都做了.
一步一步都可以.
你讓我相信你.
你說你會留心聆聽.
你說.
把肚子交託給你.
在這裡.
你看到嗎.
什麼.
你替我辛苦.
No.
你不要可憐我.
就算多重多累.
我一個人.
一個人都.
拿得起.
就算我一個人.
一個人湊大個人都可以.
有什麼難倒我.
還帶了一件行李.
休息一會.
用不著這麼多行李.
又一件.
感恩.
其實很多人.
都忘記感恩.
排排坐.
都不明白那些人.
帶這麼多行李出門幹什麼.
行李都掉了.
帶得多就掉得多.

$^{121}$你看.
每天都有十幾件.
我的行李在這裡.
失物認領處.
沒有東西.
等你主人回來收拾他們.
就有家了.
不過.
不是每一樣東西都失去了.
都找得到.
這兩件就快成為.
鎮店之寶了.
這件叫忠誠.
這件叫信任.
客人這麼多年.
都沒有回來找他們.
可能根本不知道.
掉了.
又或者明明知道.
根本都不在乎.
這兩件呢.
不在這裡.
這件和這件.
就是大老闆留下的.
就是親情和健康.
這樣都可以留下.
我有電郵給他.
他說健康先放在這裡.
遲下有機會.
才回來拿.
先講清楚.
沒有人擔保健康.
還在這裡.
這件親情怎樣.
他說不要老婆.
親情不要.
他要我們幫幫忙.
重新收拾.
大哥.
重新收拾不是我部門的.

$^{161}$關我甚麼事.
還嫌我不夠多工作.
這件呢.
厲害.
友情.
已經絕版.
就這樣不見了.
我不會太可惜.
拜託.
這麼貴重的東西.
小心一點.
這件呢.
哇.
充滿著歷史價值.
我記得.
我第一天來Lost and Found的時候.
就見到它.
主人不知道有沒有遺言.
掉了.
看一看.
這件呢.
扁嘴了.
看一看.
最初的內心.
名字很靚.
看看主人.
叫甚麼名字.
是我的名字.
我的東西.
我不見了這件東西嗎.
甚麼時候.
在哪裡不見的.
喂.
你出聲當幫忙.
我只是得到支持.
Hello.
算了.
你又怎會認我.
這麼多年來.
我又是自己跟自己說話.

$^{201}$這麼多年來.
我Lost了.
不過這些東西.
友情.
親情.
愛情.
夢想起散 希望 目標.
每樣東西我都有份.
我快連自己也Lost了.
我當然想找回這些東西.
哪裡不想.
信你.
我老嗎.
我很相信.
萬事都是由你掌管.
我也想交給你.
我問過你很多次.
Show me.
Show給我看看.
找你幫手.
我有叫過你幫手找的.
如果你想幫我的話.
你可以.
你可以贏.
我傷心又失望.
都不知道多少次.
我已經很累了.
我有時想.
你可能忙.
忙得連我都Lost.
但你都不知道.
我一直等.
一直等.
希望有一天可以找回這些東西.
但是等下等下.
又好像等到一天.
我終於找回這些東西的時候.
已經太遲了.
因為我才發現.
原來我曾經.

$^{241}$全都Lost了.
好了.
我畢竟是一個網絡紅人.
你們認識我嗎.
不認識我不要緊.
今天在Follow這個ID.
包裹你們走向.
潮流的尖端.
我呢.
今天已經在Facebook.
Line.
WhatsApp.
WeChat.
Snapchat Instagram.
Post了給大家看.
我今天會到Po .
讓我的Fans.
朋友和Followers.
都可以來一睹我的真面目.
不知道站在哪個位置.
最容易讓他們找到我呢.
喂.
不如你告訴我站在哪裡最好.
OK 不說算了.
自己和自己說話.
喂.
不要呀.
在做甚麼呀.
有沒有興趣一會去喝杯東西呀.
你沒空嗎.
你這麼忙嗎.
好吧.
遲些再約.
拜拜.
又是我呀.
沒有呀.
看看你一會有沒有空.
你移民了嗎.
現在三更半夜.
對不起.

$^{281}$早點睡吧.
拜拜.
現在.
注意了就去改電話.
我呢.
就不信我一千多個contact.
會找不到一個人.
來我的直播機.
驚喜呀.
是我呀.
在做甚麼呀.
你進了醫院.
做這件事.
陪你.
是呀.
我也很想過來陪你.
但我剛剛出了埠.
是呀.
真是不好意思.
是呀.
我回來馬上來陪你.
是早日康復.
差點中招.
算了 打給他吧.
喂.
你之前不是叫我教你化妝嗎.
擇日不如撞日.
就一會兒.
甚麼 警局.
真壞.
讓我改電話.
我呀.
你們呢.
之前不是經常叫我跟你們回教會.
去拜一下神嗎.
是呀.
今天好有mood.
好有興致.
去認識一下你的神.
不如呢.

$^{321}$你們過來載我吧.
崇拜.
快要開始了.
好了 看你們這麼有誠意.
給我地址.
我可以飛過來.
我們等一會 拿支筆.
好了 你說吧.
喂.
沒電.
Hello.
你現在看到了.
不是我不找你.
是你不給我機會.
我都很忙的.
你以為真的坐著.
等你來like我嗎.
反過來.
我就不信你會坐著.
等我來like你.
還有.
你那些followers.
跟我那些.
相差很遠.
經常說一套做一套.
甚麼.
是呀.
他們之前有叫我去那個.
甚麼.
報道會.
我那天真的沒空.
那天有個.
ig live show.
你想到的ig live show.
重要一點嗎.
甚麼.
是呀 他們之後.
還有叫我去.
pop up picnic.
那天九十幾度.

$^{361}$人都融了.
我去了.
中了暑 就幫不了.
然後.
他們就.
沒再找我了.
那你還想我怎樣.
我連電話都主動.
打了給他們.
是有誠意就過來.
載我回教會.
我現在要做的事都做了.
真的不虛把關.
不虛把就.
你以為有補嗎.
下次我連電話都不聽你們的.
慘.
快點.
差點就知道在哪.
我自己找.
我呢就不怕坐飛機.
不過最怕.
就是這件事.
不明白嗎.
金屬探測器.
弄到我的寶貝.
都打開了給人看.
那怎麼行.
甚麼來的.
沒有.
我沒有.
甚麼來的.
沒有.
我少少成績感.
這個.
全城祈禱用事.
冠軍.
這個.
四十日.
禱告達人.

$^{401}$這個.
很厲害.
看到嗎.
見歷史紀錄.
七七四十九.
金石紀錄保持者.
就是我.
OK.
少少成績感.
我很信祈禱.
祈禱的力量.
真的很大.
我很信祈禱.
祈禱最重要的是甚麼.
這個.
信心.
你一定要有很多信心.
另外還有甚麼.
就是這個.
恆心.
祈禱當然要有恆心.
教會有甚麼.
祈禱會.
你們一定會見到.
教會無論有.
大小的事務.
所有弟兄姊妹.
都希望我幫忙去祈禱.
因為他們說.
我祈禱的頭髮.
好像特別漂亮.
OK.
都幾好.
不過.
直至那一次.
師母病了.
我們大家都很誠心.
去祈禱.
祈求上帝去醫治她.
我.

$^{441}$我還要負責.
帶著一大群弟兄姊妹.
去負責祈禱.
但結果.
結果怎樣.
師母都沒有回來.
自從那次之後.
就再沒有人叫我.
去祈禱用事了.
再沒有人叫我.
去幫她上祈禱了.
我一個人說.
我的祈禱.
很漂亮.
你知道的.
是你給我力量的.
你看.
我有用的.
我有聽的.
你說的.
我們祈求.
我們敲門.
你就會應門.
但為甚麼.
為甚麼你要將你對我的承諾.
收回.
你叫我做祈禱用事.
我現在變不像了.
我像個後輩多些.
你以前不是這樣的.
你和我很有默契.
我們一起.
我們做了很多好成績.
為甚麼你現在會變成這樣.
為甚麼你會一時一樣.
你說吧.
你告訴我吧.
Hello.
我這樣說.
太過份了.

$^{481}$不過算了.
你不會理我.
我現在的祈禱都沒有把握.
我真的很累.
我很累.
你知不知道.
我很想你像以前一樣.
不如這樣吧.
你好像以前一樣.
Show給我看.
你這個無限的大能.
你和我說話.
就像無限的信心.
和恆心.
你看.
\newpage



\section{20:20}
\label{sec:mbI7RGfwLmY}
\textbf{Exploring the Holy Land: Hezekiah Tunnun (2 Kings 20\_20, 2 Chronicles 32\_30) -  Rev. Mark Chan}
\newline
\newline
連結: \href{https://youtube.com/watch?v=mbI7RGfwLmY}{\texttt{https://youtube.com/watch?v=mbI7RGfwLmY}} ~~~~ 語音日期: 2023-07-18
\newline
\newline
\hyperref[sec:hrYJaJP5Fzk]{\small{< < < PREV SERMON < < <}}
~
\hyperref[sec:index_chronic]{\small{[返順時目]}}
~
\hyperref[sec:index_scriptual]{\small{[返順卷目]}}
~
\hyperref[sec:LITCDKMXv4k]{\small{> > > NEXT SERMON > > >}}
\newline
\newline
$^{1}$email, so yeah..
So today we have Pastor Mark to give us a message today..
- Hey guys, we are starting a new series,.
and this series I want to connect the Bible and the land..
Now you have heard me saying the Bible.
is not written to us, right?.
It's written for us, alright?.
Now why do I say that?.
Why do I say that?.
Because indeed the Bible was an ancient document,.
and it's not directly written to us,.
but it's written for us..
That's why we need to understand the Bible.
in its own context..
So I want to connect the Bible and the land..
Sometimes we read the Bible, we read scripture,.
without the context, without the background..
And that's why we miss so much of the nuance.
of what the word of God will say..
So I would like to explore the whole land with you,.
at least virtually, virtually we can do that,.
and think about some of the Bible passages,.
some of the words of God..
And today we are going to a place called Hezekiah Tunnel..
Hezekiah Tunnel..
And one of those passages that mention this tunnel.
is found in the book of 2 Kings, chapter 20, verse 20..
Let me read that first..
The rest of the events, the account in Hezekiah's reign,.
including the extent of his power and how he built a pool.
and dug a tunnel to bring water into the city.
are recorded in the book of the history.
of the Kings of Judah..
Hezekiah died and his son Manasseh became the next king..
That is toward the end of the kingly account.
or royal record of King Hezekiah..
Now he's a good guy, and that's why we want to spend.
some time with what he did and what it means to us today..
There are three points I want to go through..
The first one is to understand who Hezekiah is.

$^{41}$and the record of his reign..
That's the first part..
The second part is the importance of faith.
from a biblical perspective..
And third point is the call for stewardship,.
biblical stewardship..
Now first one is the record of Hezekiah..
Now if you read 1 and 2 Kings, it's original one book,.
it records king after king what they did.
and sort of a political history, sort of a political book,.
talking about the kingdoms, right?.
Remember David, well Saul was the first king,.
and then David and then Solomon,.
and then the United Kingdom under these three kings.
crumble after Solomon..
And then it was split into North and South, right?.
The Northern Kingdom is called Israel..
The Southern Kingdom is called Judah..
So then they parallel and develop,.
and there are so many different kings.
in the Northern Kingdom and the Southern Kingdom..
And 1 and 2 Kings are tracing the political history.
of what happened to the people of God in those days..
So there were about 19 kings in the Northern nation.
of Israel, 19 kings, and all of them,.
all 19 of them are bad guys..
And then the Southern Kingdom had about 20 kings,.
only eight of them are good guys,.
and all the rest of them are bad guys..
Now this scripture, this passage that I just read.
is the conclusion of Hezekiah's reign..
Now the Bible in 1 and 2 Kings.
spend about three chapters, long chapters,.
to record what is important in the political history.
of King Hezekiah..
And this is toward the very, very end,.
the last couple of verses..
He said what?.
"The rest of the event in the reign of Hezekiah,.
"including the extent of his power,.

$^{81}$"and how he built a water pool.
"and dug a water tunnel to bring water into the city,.
"are recorded in the book of the history.
"of the kings of Judah..
"And then he died, and his son became the next king.".
Now what happened in the reign of Hezekiah.
was that he had a big problem,.
and the problem is called Assyrian problem..
Now if you know, imagine a big map,.
the Middle East, the Near East..
Right at your upper right-hand corner.
is the big guy, the Assyrian Empire,.
with mighty military power..
And they're expanding to the west..
Expanding to the west,.
and the ultimate goal is to conquer Egypt..
In order for Assyrian Empire to conquer Egypt,.
he has to pass through one land bridge,.
the only land bridge of the world..
It's the land of King..
That's where Hezekiah was, Jerusalem was there..
Israel, Palestine..
It's the land bridge for all three big continents, right?.
Europe, Asia Minor, and then Africa, and then Asia..
Now he had to pass through the land of Kingdom.
in order to go to his ultimate destination,.
which is Egypt, right?.
So he had to pass through that land,.
but if the local kings in the Middle East.
if the local kings in the land bridge.
obey him and follow the empire of Assyria,.
then you're okay, right?.
You just pay more taxes, right?.
If you rebel against the Assyrian Empire,.
the Assyrian Kingdom, the military might.
will come down upon you and will crush you..
Now in those days, the military strategy.
for conquering a city is to lay siege around the city..
You know what lay siege around the city is?.
They surround the whole city,.

$^{121}$and then they trap you inside the city wall, right?.
And then until you, staff you to death.
or thirst you to death, right?.
You can't come out, you can't go in, right?.
And that's what Assyrian military is trying to do..
To Hezekiah in Jerusalem..
So what Hezekiah did was he tried to block.
all this small water sources,.
all of this mountain spring, small springs,.
around the city of Jerusalem in the hill country..
And then he shut himself into his own city..
And Hezekiah's strategy is so that the city.
is prepared for the siege with enough food and enough water.
so that the Assyrian soldiers surrounding Jerusalem,.
hopefully they will get tired.
and then they will go back home, right?.
So as long as you make sure you are sustainable,.
the city is sustainable, then you're okay..
And that is the strategy of Hezekiah..
I think that's the only way, right?.
'Cause you cannot stand against a monster, right?.
You gotta prepare yourself..
So that's what Hezekiah did..
He built a new water tunnel to bring water into the city..
Now back in those days, usually when you need to get water,.
you had to go outside of the city wall to the valley..
'Cause usually the spring is in the valley, right?.
It's outside the city wall and then you come back, right?.
In time of peace, you're okay, right?.
You pass through city gate two times a day,.
go out and go in..
You get enough water, you come back..
But in times of war, you cannot do that..
That's what Hezekiah did..
He tried to block all this little water springs.
in Jerusalem area and then shut himself into his own capital.
and built a new plumbing to bring water into the city.
so that they can sustain the Assyrian siege..
Now I have three minute video..
I wanna show you what this tunnel is like..

$^{161}$Now as I said, the background of the Bible,.
the context of Bible is very concrete..
You can actually walk it..
You can actually see it and touch it..
So Hezekiah tunnel is one of them..
Let me try to show you the three minute video..
(crowd chattering).
- My right is Gihon Springs..
Gihon Springs right here..
(speaking in foreign language).
- Oh, so cold..
(crowd chattering).
(speaking in foreign language).
(crowd chattering).
(speaking in foreign language).
(crowd chattering).
(speaking in foreign language).
(crowd chattering).
(speaking in foreign language).
- This is the Hezekiah tunnel.
and we are walking through it.
and we are using a 360 camera.
so that you can see the environment..
And this is a man-made tunnel, 2,000 feet long.
and done by two different teams..
And one team starting from the spring,.
from the origin of the water..
And the other team is coming from the port of Siloam,.
south of the city of David..
And two teams listening to one another.
and they are built through the mountain..
And as you can see, the size of the tunnel is a body..
And they dig and dig and dig.
and they listen to each other and the noise of the digging..
And they try to meet together..
And at the end of the engineering project,.
two teams did meet together..
And it was a celebration..
And the water tunnel is through.
and the water from the Gihon Strait.

$^{201}$can go down to the south of the city of David.
so that people can get the water.
without going out of the city wall..
(speaking in foreign language).
(people chattering).
- Sorry about that..
This is my bad..
And now, the tunnel is too long..
(speaking in foreign language).
(people chattering).
(speaking in foreign language).
- Well, it is very interesting, right?.
You can actually walk the tunnel.
that is recorded in the Bible..
And especially you have a 360 camera..
And thanks to Hale for lending the camera..
I think in the YouTube world,.
we are the first one recording the tunnel.
using a 360 camera..
Yeah, praise God..
Anyways, makes you wanna go, right?.
Makes you wanna go..
It spent about 40 minutes to walk the entire tunnel..
Total pitch darkness..
And it was built in Hezekiah time, 700 BC..
That's when Hezekiah had a crisis..
If you see the diagram here,.
you see that the Gihon Spring,.
which is now flowing,.
even now, there are clean water coming out of the mountain.
from Gihon Spring..
Gihon Spring is on the right..
You see that blue line?.
That blue line, S-shaped, curved blue line.
is the line of Hezekiah Tunnel..
I was kind of walking the first third of the tunnel..
So you start from Gihon Spring on the right..
And then the entire tunnel is about 2,000 feet long..
And you walk the entire thing with a flashlight..
And you went through the tunnel to the end,.

$^{241}$and you come out to the open air again,.
which is on the left..
The left is Siloam Pool..
Now this is a mighty thing that Hezekiah did in his time..
And let me tell you this..
If you understand a little bit about the background,.
you know this tunnel, 2,000 feet long,.
a body width, right?.
If you're too big in size, you cannot fit in, right?.
And good luck to you, right?.
And then water flowing from the spring.
all the way to the Siloam Pool inside the city..
This water tunnel, this plumbing system.
was the greatest achievement in his time.
in the entire world..
You go back to 700 BC,.
the greatest plumbing system was done by Hezekiah Tunnel..
You can't find it in Rome..
You can't find it in Greece..
You can't find it in Iran or Iraq..
The greatest plumbing system of all time.
during Hezekiah's time was found there..
And we were walking through it, right?.
The greatest plumbing system..
Probably the plumbing system is better.
than those in your homes..
The best plumbing system..
The best achievement..
And because of that,.
now one of the questions is why the tunnel was not straight?.
Why wasn't it straight?.
From left to right,.
the shortest distance should be a straight line, right?.
Through the mountain, right?.
You just hill through the mountain, straight line..
Now what happened was that.
because during the Assyrian crisis,.
they had to build this plumbing system ASAP,.
as soon as possible..
So he divided into two teams of workers..

$^{281}$They're trying to dig from,.
one dig from the spring, downhill..
The other team, team B, is digging from the water pool.
at, on your lower left, all the way uphill..
And they're trying to hear the noise of digging..
Trying to discern where the other team is..
So they're trying to dig through the mountain..
That's why it's curving around inside the mountain..
Now this is a great achievement.
because, no, you didn't have any GPS, right?.
You didn't even have flashlight back then..
You only have oil lamp, right?.
You have oil lamp and candle,.
and then you have to go through the mountain,.
totally dark, and try to hear the noise.
so that you can meet together at the midpoint..
You see the midpoint?.
They're trying to dig in, you know,.
with some excitement, right?.
You want to get to work, right?.
Because the Assyrians are coming..
Crisis is coming to us..
So they're trying to meet together..
So the middle point is kind of like,.
there are some false starts,.
there are some false places.
they have to correct themselves..
So that's why the tunnel is not straight..
And then toward the end, there is an inscription..
About 150 years ago, there's a high school boy.
who skipped school that day, 150 years ago..
And then he was swimming in the pool..
And then he started to say to himself,.
"Okay, let me go inside the tunnel.".
So he swam, swimming, swimming,.
and swim a little bit into the tunnel,.
about six meters deep inside the tunnel..
And then he slipped..
He drowned a little bit, and then he got up..
And then he wiped his face,.

$^{321}$and he saw a stone inscription on the side of the tunnel..
And that inscription can be translated,.
and it is kind of, you know,.
two feet by one foot,.
a stone inscription..
And if you translate it,.
what it says is the story of digging through the tunnel..
Let me translate it..
The tunnel, this is the story of the tunnel..
And they are digging against each other.
while three cubits was left..
The voice of a man called to his teammate,.
his counterpart, and there was a sada,.
whatever that is, right?.
Probably a crack in the rock..
And on the right and on the left,.
on the day of the tunnel,.
the stone covers struck each other.
toward his counterpart, $\times$against X,.
and flow water from the source, the Gihon Spring,.
to the port for 1,200 cubits..
It's about 1,900 feet long,.
and six cubit high in terms of the height of the tunnel.
over the head of the stone covers..
Now, I don't know how to translate.
to express the excitement,.
because this is a historical record, right?.
So all you can do, all I can do,.
is add exclamation mark, right?.
You add a lot of exclamation mark..
This is the greatest achievement of all time..
It is the water tunnel is being done, right?.
So the high school boy who skipped school that day.
found this important inscription..
So I guess one of the lesson is it's okay to skip school..
You found something really, really important, right?.
Anyways, this is what the Bible recorded..
This is the water tunnel, the greatest achievement of all..
And second king did not fail to mention it..
And that's the first point..

$^{361}$But the second point I want to bring out.
is the importance of faith..
Let me try to understand the biblical perspective.
on the life of Hezekiah..
And how does the Bible see your life?.
How does Bible evaluate your life?.
How does the Bible evaluate the reign of Hezekiah?.
Let's take a look at this passage again..
2020, second Kings..
The rest of the events in Hezekiah's reign,.
including the extent of his power,.
how he built this and that water tunnel,.
recorded somewhere else in the book of history.
of the kings of Judah..
Now this face, the rest of his achievements,.
is often repeated throughout the book of first and second Kings.
for each king, his political history..
At the end of it, the book of Kings always concludes.
the rest of the events of this king or that king..
Now in terms of Hezekiah, the Bible sees,.
did record the great achievement, the water tunnel..
Again, it's the best plumbing system in his time.
in the world..
However, what is the key word here?.
The key word is the rest..
Now if you record, if you are the court historian,.
if you're going to write a book on Hezekiah,.
what would you spend most of your details.
or most of your content on?.
You're gonna talk about this tunnel..
You're gonna talk about the engineering project,.
the water engineering project, right?.
Now if you have read some of the political history.
or autobiographies of politicians,.
you will open the book and you find,.
okay, this president or that politician.
has done something great, right?.
That's what biography, royal biography does..
To praise a guy or a lady,.
how, what great achievement he has done or she has done..

$^{401}$All kinds of autobiographies,.
Obama, Hillary Clinton, or whatever big politicians,.
when they write their own biography,.
when they write their own story,.
they would lift up or exalt.
the great things they have done..
Here, second Kings has done something totally different..
The greatest project, yeah, I'll mention it,.
but it's the rest of the archive..
You can dig it up, right?.
I save it on another hard drive, right?.
But you can reshift the data if you want,.
if you're interested..
But I'm not interested in it, right?.
It's in the rest of his deeds,.
the rest of his achievement is recorded in the other book..
That's what I'm talking about..
I wanna talk about not his great achievement..
I wanna talk about something else..
Now, this is the biblical perspective on someone's life..
Think about your life..
If you're gonna write a book,.
a biography of yourself, your history,.
what would you put in it?.
You probably would think, naturally,.
we would think that we'll put our great achievement, right?.
Which school I went to, I went to Harvard,.
I went to the best school of the world,.
I got the best education, right?.
I graduated with, I don't know, straight A, 4.0, right?.
And I went to this company or that company,.
Fortune 500 company, I got so much achievement..
And I'm gonna put it in my book, my reign..
But that is not the perspective of the Bible..
The Bible pay attention to something else..
Yeah, you have your own achievement, that's fine..
I'm gonna talk about it in one verse..
The rest of the chapters from 2 Kings 18 to 20,.
what does the book of Kings talk about?.
And think about it..

$^{441}$If you understand the perspective,.
how would the Bible, how would God sees your life?.
How would God evaluates your life?.
The rest of his deeds, mentioned somewhere else..
But what I wanna focus on,.
what the Bible wants to emphasize is this..
At the beginning of Hezekiah's reign, his record,.
18 verse three,.
Hezekiah did what was pleasing in the Lord's sight..
Just as his ancestor, his father David, had done..
He removed the pagan shrines, smashed the sacred pillars,.
and cut down Asherah posts..
Those are the pagan idols across his country..
He broke them and he removed them..
Hezekiah trust in the Lord God of Israel..
There was no one like him among all the kings of Judah,.
either before or after his time..
He remained faithful to the Lord in everything..
And he carefully obey all the commands.
of the Lord had given Moses..
So the Lord was with him..
What the Bible sees is his faith,.
his spiritual life, his relationship with God..
When things are going right in his life, God walks with him..
When he faces crisis, a Syrian crisis,.
or whatever crisis you are facing,.
he depend on the Lord, he trust in the Lord..
That goes into the record of first and second kings..
That's what God wants to see..
That's what God want to focus on..
You see what I'm saying?.
Your construction, your achievements,.
maybe they are important..
You know, in Hezekiah's day,.
he had to devise some political strategy..
He has to achieve something in order to defend the city..
But at the end of the day, who save him?.
Yahweh save him, God save him..
When crisis hits you, you can do all the preparation..
This is part of our faithful stewardship, right?.

$^{481}$We have to be responsible..
We have to go to a doctor..
We have to find the best medical treatment..
We have to find the best way to manage our lives..
But the end of the day, do not forget.
as what happened to Hezekiah..
He did not save himself..
The water tunnel, as mighty as it was, did not save him..
At the end of the day, he went to the temple.
and prayed for salvation..
Prayed that God will help him..
And then God help him..
At the end of the day, the credit goes to the creator,.
his Lord and savior, no one else..
You see what I mean?.
You can manage your life..
You should..
You may have a little bit achievement here and there,.
here and there, and maybe those are important..
But when God sees your life, okay,.
I put one verse, summarize it..
All your achievement are becoming footnotes..
The main content of your life is your spiritual life..
It's your faith..
Whether you are depending on God in times of crisis,.
whether you do not forget God in times of blessing,.
it's your relationship with God..
That is all that matters at the end of the day..
That is God's perspective seeing you..
Not how the culture sees you, not how the world sees you,.
not how you see yourself..
It's how the Bible sees you..
That brings to my third point, it's the call for stewardship..
God has entrusted us with different things,.
resources, time, our money, all kinds of resources.
He has entrusted to us..
That's why we have this word, theological word, stewardship..
Comes from an old English word, steward..
A steward is a manager of the house..
To use a modern analogy, it's kind of like a bank manager..

$^{521}$When you go to the bank, you deposit your money.
into your bank account..
And the bank is entrusted with your money,.
and the bank, hopefully, the bank will faithfully.
and wisely manage your money..
In times of need, you go to the bank.
and you withdraw some money, right?.
It really does not belong to the bank manager..
The bank manager is managing someone else's resource,.
your resource..
This is the same way our relationship with God and you..
God entrusts something in you, in your life..
And you're kind of like a bank manager..
You're a steward..
And the only thing required for you is faithfulness..
Whether you're using all the resources as much as possible.
to serve God and to worship Him..
That is biblical stewardship..
And do not forget that..
And that's what Hezekiah's story is all about..
Why didn't the Bible, I actually want the Bible.
to talk more about the Hezekiah tunnel, right?.
I want a whole chapter talking about it..
However, the biblical author did not care..
And think about it..
In the great scheme of things,.
God does not really care about your earthly achievement..
What He cares is how you manage it to glorify Him..
How you build your life on the word of God,.
in obedience to Him..
That is what matters at the end of the day..
You are God's steward and trusted with all sorts of things..
How are you managing it?.
If God calls you to give an account today,.
how would you respond?.
Today..
Will God seize your life saying,.
"Hey, you are being faithful..
"You're a good and faithful servant..
"And you're coming to my kingdom.".

$^{561}$Or maybe God will see you, you are wicked..
You wasted all of my resources..
You are not a faithful manager..
That is biblical stewardship..
I officiate a lot of funerals..
I've served in the pastoral role.
for probably close to 30 years..
And I have done so many, so many funerals..
I remember one year, in my early year,.
I have to officiate 40 funerals in a year..
40 funerals..
Sometimes two funerals a week..
Especially in winter time in Chicago..
More people die in the depressing time..
If you go to a funeral, if you go to a funeral,.
well, I mean, it's not a nice thing to say,.
but you should go to funerals more..
At 10 funerals..
You know why?.
Because there's a part in the funeral service,.
there's a part called eulogy..
Eulogy means speak good, right?.
Speak good, it's a Greek word..
Speak good of a person, right?.
Eulogy is, you know, reciting the history.
or remembering the history of someone,.
the deceased family member..
All the eulogies I've heard.
have nothing to do with the deceased achievements..
Nobody will come up to the podium in the funeral service.
and say, "Oh, he got a lot of money, you know what?.
"He got a lot of investment..
"He's so great in terms of money..
"Look at how many figures he has in the bank account.".
Nobody will say anything like that..
Most of the eulogies will talk about.
the character of the person..
How he has been a blessing to other people..
The family, the community, the neighborhood, right?.
The church, his spiritual life..

$^{601}$Most of the things mentioned in the eulogies.
are those things, the spiritual things, so to speak..
Your relationship with God.
and your relationship with other people..
Whether you are walking in the way of God.
or you are loving other people just as God loves them.
and you try your best to do it,.
those are the things going to the account..
Those are the things going to God's record of your life..
Don't lose that focus..
Let's pray..
Father, we wanna build our lives on your grace.
and thank you for calling us..
Thank you for entrusting so many wonderful,.
great things in our lives..
Yes, we do have some achievements..
We exert great efforts in building our lives,.
but we pause right now.
and ask you to help us see our lives in your perspective,.
to evaluate all our doings, all our efforts.
in light of your eternal kingdom..
And I pray, I pray that in times of blessing,.
we do not forget you..
We thank you, we praise you..
And in times of crisis, we don't forget you,.
knowing that you are our rock and you are our salvation..
Thank you for your deliverance, thank you for your healing,.
thank you for your blessing that we enjoy..
So in times of crisis, help us to focus on you.
because in the end of the days,.
it's our faith in you that matters..
Pray that you will help us to focus on your kingdom..
\newpage



\section{1:1-18}
\label{sec:LITCDKMXv4k}
\textbf{Count It All Joy (James 1\_1-18) - Caleb Lee}
\newline
\newline
連結: \href{https://youtube.com/watch?v=LITCDKMXv4k}{\texttt{https://youtube.com/watch?v=LITCDKMXv4k}} ~~~~ 語音日期: 2023-07-25
\newline
\newline
\hyperref[sec:mbI7RGfwLmY]{\small{< < < PREV SERMON < < <}}
~
\hyperref[sec:index_chronic]{\small{[返順時目]}}
~
\hyperref[sec:index_scriptual]{\small{[返順卷目]}}
~
\hyperref[sec:2mwrzSe3J14]{\small{> > > NEXT SERMON > > >}}
\newline
\newline
$^{1}$Good morning..
Okay, the mic is working..
Today we will be.
in the book of James..
So if you would all open with me your Bibles to the book of James..
It is closer to the back following.
Hebrews. So I'll give you a moment..
If you do not have your Bibles with you or on you.
we have it on the screen here..
So we will be reading from its first chapter.
verses 1 to 18 and then we'll open.
in a word of prayer..
James chapter 1 verses 1 to 18..
James, a servant of God and of the.
Lord Jesus Christ. To the 12 tribes in the dispersion.
greetings..
Count it all joy my brothers when you meet trials of various kinds.
for you know that the testing of your faith produces steadfastness.
and let steadfastness have its full effect that you may be perfect and complete.
lacking in nothing..
If any of you lacks wisdom let him ask God who gives generously.
to all without reproach and it will be given him. But let him.
ask in faith with no doubting for the one who doubts is like a wave.
of the sea that is driven and tossed by the wind. For that person.
must not suppose that he will receive anything from the Lord. He is a double.
minded man unstable in all his ways..
Let the lowly brother boast in his exaltation.
and the rich in his humiliation because like a flower.
in the flower of the grass he will pass away. For the.
sun rises with scorching heat and withers the grass. Its flower.
falls and its beauty perishes. So also will the rich man fade.
away in the midst of his pursuits. Blessed.
is the man who remains steadfast under trial.
for when he has stood the test he will receive the crown of life.
which God has promised to those who love him. Let no one say when he is tempted.
I am being tempted by God for God cannot be tempted with evil.
and he himself tempts no one. But each person is tempted when he is.
lured and enticed by his own desire. Then desire when it has conceived.
gives birth to sin and sin when it is fully grown brings.
forth death. Do not.

$^{41}$be deceived my beloved brothers. Every good gift and every perfect.
gift is from above coming down from the father of lights with whom there is.
no variation or shadow due to change. Of his own will.
he brought us forth by the word of truth that we should be a kind of.
first fruits of his creatures..
Let's pray..
Father all praise all glory all.
honor be to you. For righteousness and justice.
are the foundation of your throne. The heavens.
declare your righteousness God and your people have seen.
your glory..
In the coming moments we ask that you would speak that your word would.
be illuminated for us. Guard the.
ears of your people oh God that their feet would not be stumbled..
Would you give them hearts that are renewed so that.
their hearts would yield to you and your spirit..
Call us to be still oh God so that we may again.
find our refuge in you. All this we pray in.
Jesus' name amen. Okay..
Let me go back to the beginning..
That's too far..
So I think you'll find.
that one of the most pressing questions that appear before.
us in our daily lives is the question what now?.
After all our lives are constantly changing.
because time is linear it moves in a linear fashion.
we are continually picking up new responsibilities..
A few of us here are in school..
Some of us most of us are in the workforce.
and there are a few of us here that are new.
parents. Those of us who are not going to be new parents we have.
friends who are being new parents. All new life stages.
all new circumstances everything is.
kind of inescapable. And with every change.
in our life we're faced with the question what.
now? What should my life look like? What is expected.
of me and what do I do now?.
And this question of what now becomes even more pressing when the.
changes in our lives are not just confusing or new.
but bad. Financial.

$^{81}$strain, relational strife,.
deep, deep loss..
All of it is a heavy, heavy weight on our shoulders..
And to further complicate things you'll find that there is.
more to this question of what now?.
It's a Sunday morning. You guys are here. You chose to be.
here. And so the question for you is.
as people of God how are you expected.
to live? If you look again with.
me to James chapter 1 verse 1 you'll find that the.
audiences of this letter are the 12 tribes of.
the dispersion. James a servant of God of the Lord Jesus.
Christ to the 12 tribes of the dispersion. These are the.
Jewish Christians living outside of Palestine. Those living in a foreign.
and famished land. They are displaced..
They're poverty stricken. They are starving from persecution.
and social exile..
In circumstances like that it's.
almost natural for hardship to bleed into character..
It's almost expected. Under this sort of.
pressure the cracks, our cracks begin to show..
They widen to reveal already present.
flaws in our character. And if we're being honest with ourselves.
we can all recognize some inadequacies.
in ourselves that are covered up and usually hidden behind our strengths..
Like we see in verse 14..
Let me go to verse 14 real quick..
Verse 14..
What threatens to drag us away can so often come from.
the desires within us that seek to produce sin and subsequently.
death. It says, "But each person is tempted when he is lured and enticed.
by his own desire. Then desire when it has conceived gives birth to sin.
and sin when it has fully grown gives birth to death.".
Life is hard. Life is hard.
and it makes us hard. And yet, and yet we find.
that James heaps instruction upon instruction.
upon these Jewish Christians..
He gives instructions that call for those of the.
dispersion to behave contrary to their circumstances..
To go against the natural flow of action..

$^{121}$Look at verse 2. What does it say?.
It says, "Count it all joy, my brothers, when you meet trials of various kinds..
Consider it joy when you face trials." James makes.
this point. Let this be an occasion for joy..
Joy..
This is a strange and unexpected response to suffering..
It almost feels inappropriate.
in some ways. But don't misunderstand..
Don't look at this and think to yourself,.
I gotta pretend that things are okay..
This is not a forced joy, but an eschatological.
joy. A joy that comes from looking beyond..
And seeing something joyous. It's a response of joy..
If you look at verse 2, James calls for his brothers and sisters.
to consider it joy when they face trials. And the reason.
the reason to consider such things joy, such hardships.
as joy is because the testing of their faith produces.
steadfastness. Note that this.
testing they speak of is not to determine.
whether or not there is faith in you. It's instead.
here to refine faith that is already present. And in this.
refining of faith, steadfastness is produced..
You'll find that only through this refining.
process would this facet of steadfastness.
this endurance and perseverance exist within.
believers. But.
with all this talk about steadfastness.
you'll find that this is not the goal. Not just for the.
sake of being a steadfast person. Steadfastness itself.
is not the goal. Instead, if you look with me.
to verse 4, you'll see that the goal.
and desired outcome of suffering trial is something else..
Right? And let steadfastness have its full.
effect that you may be perfect and complete, lacking.
in nothing. The goal is not a specific virtue..
It's not a particular fruit. The goal.
of steadfastness is your perfection..
Your completeness. Perfection.
is central to James' writing. It's this picture of wholeness.
a finished and flawless.

$^{161}$work. We see this idea of perfection.
in verse 5 through the contrasting of two men..
Let's take a look. The first person.
the first person spoken of by James embodies.
wholeness. They are lacking in nothing because they ask in.
faith and they receive wisdom from God. The first person is.
wise. But the second? The second person is.
like a wave of the sea, blown and tossed.
by the wind. They are double-minded.
and unstable, doubtful of the security of God's promises for them..
And they're torn apart with inner conflict. They wander.
on and off every road they can see. In everything.
they do, they are fractured and inconsistent. They're not.
whole and they are not wise..
In all of this is the implication that God desires.
for his followers to live true to their titles..
Followers. What do followers do? They follow. God desires for.
his followers to follow him. And more than that, it's implied.
it's implied that God desires for those who follow him.
to be refined, to be perfected, and to be completed..
To be wise and be affected by.
wisdom. Not in part, but in whole..
For the collective voice of scripture calls for you to not just be.
followers, but imitators of Christ..
That your whole being and every aspect.
of you would be refined and conformed into the likeness.
of the rock that is higher. And that from this.
completeness, your actions would be good and pleasing.
to God. See,.
this that we talk about, this being wise,.
being perfected, this result of suffering under trial..
It is a continuing of the work of God. The call to count.
all joy in poverty and humiliation is a command that is.
founded on the finished work of Jesus. You,.
having been saved, you are called away from being fractured.
and inconsistent. You are called toward being whole.
with every aspect of your being..
The reality we see in James is that.
faith has an effect on.
life. And that the culmination.

$^{201}$of Christian life is a whole.
and perfect conformity to the likeness of Christ..
Why?.
Verse 9, "Why should the lowly brother boast.
as lowly? Because his position.
in Christ is not lowly, but exalted. He is not poor,.
but rich. Why then should the rich.
boast in his humiliation? Because his richness.
comes from his inheritance from heaven, not.
the collection of dust on this earth..
Why should those undergoing trials and temptations.
consider it joy? Because in the.
trying of your faith, you and those around you become witnesses.
to the work of God that is being done in you..
You see with your eyes that faith is present in you. And now.
you're given this opportunity to exercise the faith.
that is in you by placing weight on it, by trusting.
it. Do you see?.
Do you see here? The command here.
to count it pure joy is not another burden.
for you to carry on top of the pains of life..
It's an instruction given to you so that you can live..
Blessed is the man who remains.
steadfast under trial for he has, when he has withstood the test,.
he will receive the crown of life which God has promised to those.
who love him..
And so,.
looking at all this,.
do you want to be wise?.
Do you want your life to look like?.
My brothers and sisters, in.
sickness, loneliness, bereavement, and disappointment,.
don't keep your eyes fixed on the circumstances of life..
Don't keep your eyes fixed on the desires.
of the flesh. Instead, look.
towards the object of your faith, the good and perfect gift.
of God to you, his son Jesus. And see.
that the promises of God are fulfilled to you in Jesus, in Christ..
Know that the work of Jesus will not be undone, and that.
the Spirit is in this very moment making you more and more.

$^{241}$like Christ, and consider it all.
joy. The task laid before.
you is this. Joyfully lean into.
the work of God that is being done in you,.
and continue in that way..
Become holy like Christ..
That was a little short, but let me.
close us in a word of prayer. Would you bow your heads with me?.
Father, we give you thanks.
because you have been with.
us every single second of the day. There is not a moment.
that we are away from you or apart from you. You have.
seen us and known us, and before we were, you were..
So, Lord, ask that you would continue to.
keep your hand on your people, on these brothers and sisters of mine,.
that you would reveal to them how deeply you love them,.
that they would start there, not looking at their own selves or the things.
around them, but God, let them look at you first and foremost..
Would you make them rely on you, depend.
on you for everything in their life, that they might be a reflection of the way.
that you provide for them? Remove from them.
all these obstacles that hinder them, Lord, and keep.
them firmly in your hand. Give them a.
stronger spirit, a motivation to grow,.
and help them be good representatives of who you are..
We pray this in Jesus' name. Amen..
Thank you very much..
\newpage



\section{羅馬書 12:9-21}
\label{sec:2mwrzSe3J14}
\textbf{生命的致:致力愛 (羅馬書 12\_9-21) - 譚文倩莊傳道}
\newline
\newline
連結: \href{https://youtube.com/watch?v=2mwrzSe3J14}{\texttt{https://youtube.com/watch?v=2mwrzSe3J14}} ~~~~ 語音日期: 2023-07-25
\newline
\newline
\hyperref[sec:LITCDKMXv4k]{\small{< < < PREV SERMON < < <}}
~
\hyperref[sec:index_chronic]{\small{[返順時目]}}
~
\hyperref[sec:index_scriptual]{\small{[返順卷目]}}
~
\hyperref[sec:qy6pJ6fRtpk]{\small{> > > NEXT SERMON > > >}}
\newline
\newline
羅馬書 12:9-21
\newline
\begin{longtable}{cl}
\hline
\hline
章節 & 經文 (和合本修訂版)\\
\hline
12:9 & \begin{tabularx}{0.7\textwidth}{X} 愛，不可虛假；惡，要厭惡；善，要親近。 \end{tabularx} \\ \\ \relax
12:10 & \begin{tabularx}{0.7\textwidth}{X} 愛弟兄，要相親相愛；恭敬人，要彼此推讓； \end{tabularx} \\ \\ \relax
12:11 & \begin{tabularx}{0.7\textwidth}{X} 殷勤，不可懶惰。要靈裡火熱；常常服侍主。 \end{tabularx} \\ \\ \relax
12:12 & \begin{tabularx}{0.7\textwidth}{X} 在盼望中要喜樂；在患難中要忍耐；禱告要恆切。 \end{tabularx} \\ \\ \relax
12:13 & \begin{tabularx}{0.7\textwidth}{X} 聖徒有缺乏，要供給；異鄉客，要殷勤款待。 \end{tabularx} \\ \\ \relax
12:14 & \begin{tabularx}{0.7\textwidth}{X} 要祝福迫害你們的，要祝福，不可詛咒。 \end{tabularx} \\ \\ \relax
12:15 & \begin{tabularx}{0.7\textwidth}{X} 要與喜樂的人同樂；要與哀哭的人同哭。 \end{tabularx} \\ \\ \relax
12:16 & \begin{tabularx}{0.7\textwidth}{X} 要彼此同心，不要心高氣傲，倒要俯就卑微的人。不要自以為聰明。 \end{tabularx} \\ \\ \relax
12:17 & \begin{tabularx}{0.7\textwidth}{X} 不要以惡報惡，眾人以為美的事要留心去做。 \end{tabularx} \\ \\ \relax
12:18 & \begin{tabularx}{0.7\textwidth}{X} 若是可行，總要盡力與眾人和睦。 \end{tabularx} \\ \\ \relax
12:19 & \begin{tabularx}{0.7\textwidth}{X} 各位親愛的，不要自己伸冤，寧可給主的憤怒留地步，因為經上記著：「主說：『伸冤在我，我必報應。』」 \end{tabularx} \\ \\ \relax
12:20 & \begin{tabularx}{0.7\textwidth}{X} 不但如此，「你的仇敵若餓了，就給他吃；若渴了，就給他喝。因為你這樣做，就是把炭火堆在他的頭上。」 \end{tabularx} \\ \\ \relax
12:21 & \begin{tabularx}{0.7\textwidth}{X} 不要被惡所勝，反要以善勝惡。 \end{tabularx} \\ \\
[1ex]
\hline
\hline
\end{longtable}
$^{1}$早晨 弟兄姊妹.
上次我跟大家分享了《羅馬書》第二講.
《生命的刺屬靈恩刺》.
有個朋友上網去看.
他說他一看就以為我講《生命的腸》.
(笑聲).
我都笑了很久.
他又笑了很久.
我亦都很想在這裡多謝一些網上的朋友.
很關心我經常做什麼奇奇怪怪的.
其實我的喉嚨是有點患窒.
但是醫生給我的藥有很多副作用.
所以我不可以吃.
所以暫時你們經常聽到我奇奇怪怪的時候.
就需要大家忍耐.
今天來到《羅馬書》十二章.
講到系列的最後一講《生命的智力愛》.
愛到盡了.
我們先低頭祈禱.
你的順服 聖靈.
事常為我們禱告.
主耶穌你的話語是我們腳前的燈.
路上的光.
求你謙卑人的心領受你的話語.
祈禱交托奉主耶穌基督的聖命祈求.
第三講《生命的智力愛》.
在《羅馬書》十二章九至二十一節.
不如我們先一起讀好嗎.
第九節.
愛不要虛假.
惡要厭惡.
善要持守.
愛弟兄要彼此相親相愛.
尊重人要彼此爭先.
殷勤不可懈妥.
要鄰里火熱.
常常服侍主.
在盼望中要喜樂.
在患難中要忍耐.
禱告要專心致志.

$^{41}$聖徒有需要要供給他們.
要潔誠接待旅客.
要祝福那些迫害你們的人.
要祝福不要咒詛.
要與喜樂的人一同喜樂.
與哀哭的人一同哀哭.
要彼此同心.
不要心高氣傲.
都要輔咒卑微的人.
不要自以為聰明.
對任何人都不要以惡報惡.
眾人看為美好的事.
要留心思想.
可能的話.
要盡你們所能與眾人和睦相處.
親愛的弟兄們.
不要自己報復.
寧可給神的憤怒留地步.
因為經上記著.
主說.
報復在我.
我必報應.
不但如此.
如果你的仇敵餓了.
就給他吃.
如果渴了就給他喝.
因為這樣做.
就是把炭火堆在他的頭上.
不要為惡所成.
反而要以善成惡.
看看羅馬書十二章九至十三節.
這短短五節經文.
當中有十三個提醒.
提醒我們要去愛.
但不單止去愛.
是不虛假的愛.
不要虛假的愛.
Be without hypocrisy.
惡要厭惡.
善要持守.

$^{81}$我們要彼此相親相愛.
但說的相親相愛.
不是單單說相愛那麼簡單.
是說兄弟如手足.
弟兄姊妹那種愛.
尊重而彼此爭先.
爭先的意思是.
我們好像差不多和對方比賽.
誰快一點去尊重對方.
殷勤不可懈妥.
要火熱.
是在靈里要火熱.
我們要服侍主.
是常常服侍主.
要有喜樂.
特別在盼望里.
在等候里.
我們要有盼望.
要忍耐.
特別在患難的時候.
遇到困難的時候.
我們要有忍耐.
禱告要專心.
聖徒有需要.
要供給他們.
最後那個說要揭城.
接待旅客.
意思是一些離開家裡的人.
我們要接待他們.
來到《羅馬書》十二章的段落.
第九節一開始.
它說到愛不要虛假.
《和合本修定版》說.
愛不要虛偽.
《NASB》說.
愛不必虛假.
ESV就很直接地說.
讓愛真實.
是真誠的.
什麼是虛偽.

$^{121}$什麼是虛假.
意思是說.
無論我裡面是怎麼樣.
最重要的是.
我外表的表現.
是怎麼樣給你看.
簡單一點來說.
就是表裡不一致.
基本上我們每個人都很渴望.
有人來贊賞.
欣賞和接納.
我們都很希望.
別人對我們有好的印象.
所以很多時候.
我們敬意也好.
不敬意也好.
都會在思考.
我走出來的時候.
我是要怎麼樣表現.
怎麼去展示我自己.
怎麼去照顧我自己.
另一方面.
我很小心.
不想讓你看到.
我一些醜陋的東西.
一些不好的東西.
因為我不想有負面的批評.
我不可以讓你給我少看一眼.
你每次見到我.
我走出來都是整齊整整.
好好看看.
甚至可能讓你有點羨慕.
哥林多前書13章3節.
我若將所有的財產.
救濟窮人.
又犧牲自己的身體.
讓人誇贊卻沒有愛.
仍然對我無益處.
阿婆過馬路.
你拖她過馬路.

$^{161}$你朋友經濟有困難.
你去幫他.
你被人欺負.
我兩腳插刀.
在所不辭.
但是最重要有人看到.
最重要有人知道.
表面我是為了你.
但其實是為了我自己.
動機一點都不純正.
保羅在這裡說.
若我們的動機不純正.
愛人不是出於真誠.
genuine love from the inside.
其實一點益處都沒有.
馬太福音15章8至9節.
這百姓用嘴唇尊敬我.
他們的心卻遠離我.
他們把人的規條.
當作教義教導人.
他們拜我也是枉言.
耶穌斥責發理菜人和民事.
因為他們是假冒為善的人.
表面做足了.
其實只有嘴巴.
我們繼續看下面的經文.
14至16節.
要祝福那些迫害你們的人.
要祝福不要咒詛.
要與喜樂的人一同喜樂.
與哀哭的人一同哀哭.
要彼此同心不要心高氣傲.
禱要苦咒卑微的人.
不要自以為聰明.
祝福你的敵人.
不要詛咒他們.
當身邊的人開心的時候.
你和他一起去開心.
當他傷心失意的時候.
你要陪伴他.

$^{201}$甚至和他一起哭.
我們要有團隊的精神.
要互相扶持.
接著17至19節.
對任何人都不要以惡報惡.
眾人看為美好的事.
要留心思想.
可能的話要盡你們所能.
與眾人和睦相處.
親愛的弟兄們不要自己報復.
寧可給神的憤怒留地步.
因為經上記著主說.
報復在我,我必報應.
被傷害被陷害.
不要想辦法去報復.
我們多留心欣賞好那方面.
如果在你範圍以內有能力的話.
盡量和人和睦相處.
不要想著別人對你不好.
經常要想著究竟我要做些什麼.
去補救他.
因為在神的內部.
這件事與你無關.
因為上帝說.
我會做出判斷.
我會處理.
又要講一下我們姐妹們.
正正經經查查經.
最近我們查經查到.
《馬可福音》15章.
講到耶穌被冰冰打他的頭.
吐口水.
玩他.
跪下拜他.
耶穌不僅沒有鼻子.
也沒有還手.
有位太太說.
我就不會了.
誰動手打我.
我一定還手.

$^{241}$今天這個太太也在我們當中.
我不是想拜她上台.
而是她說完之後.
我想這是我們很正常的反應.
我沒有做錯.
你乾嘛打我.
不關我的事.
你冤枉我.
你煽我.
我要為自己平反.
二十二十一節.
不但如此.
如果你的仇敵餓了.
就給他吃.
如果渴了就給他喝.
因為這樣做.
就是把炭火堆在他的頭上.
不要為惡所成.
反而以善勝惡.
聖經這裡告訴我們.
如果你的仇家.
你的敵人餓了.
你要請他吃東西.
他口渴.
你最好給他一杯水喝.
你這個慷慨的行為.
是會嚇到他.
因為你比他更厲害.
你殺他一個措手不及.
不是在他不行的時候.
你多踩一腳.
反而是在他不行的時候.
你幫他一把.
《針言》二十五章二十一至二十二節.
你的仇敵餓了就給他吃.
如果渴了就給他喝.
因為你這樣做.
就是把炭火堆在他的頭上.
耶和華必回報你.
這句經文也有這句.

$^{281}$就是把炭火堆在他的頭上.
什麼意思呢.
在我查考這段經文的時候.
我找到一些解釋.
我覺得很有趣.
他說把炭火堆在他的頭上.
意思是你用善良去殺死你的敵人.
英文是叫.
Kill enemy with kindness.
要通過將敵人變成我們的朋友.
來去消滅他.
這個解釋來自詩篇一百四十篇的十節.
他說願他們被雕在火中.
火炭落在他們身上.
意思是.
這個將會是最好的懲罰.
想想羅馬書十二章九至十一節.
這裡說的提醒.
勸告.
勸勉.
你會不會覺得.
其實有沒有可能做得到呢.
不知道你會不會覺得.
有時做基督徒也不是那麼容易.
好像很難.
我相信我們每一個人.
都是很想土神起樂.
我們也不是想效法這個世界.
但在真實社會里.
其實真的很難.
我們不得不承認.
很多時我們都做不到.
所以保羅在羅馬書十二章的開頭.
一至二節勸勉信徒.
要將身體獻上.
當作活祭.
是說你要將你每一日的生活.
睡覺,吃飯,上班,跟人相處.
好像一個祭物一樣.
帶到主的台前.

$^{321}$擁抱神.
讓神幫助你不要效法這個世界.
求神幫助你心意根深而變化.
以至我們的生命從內到外發生改變.
叫我們有熟悉的眼光.
能夠分辨到在神眼中捍衛善惡的事.
來過生活.
成就神的心意.
不要忘記我們每個人.
都是活在神的恩典當中.
我們這樣追求過生活.
是神喜悅的.
也是對我們所領受的這個救恩.
最理所當然,感恩,合乎情理的回應.
回到這一段經文的第一節.
其實後面的很多節.
是由這一節開始.
我之前有講過.
好像一個地基一樣.
愛不要虛假.
Love be without hypocrisy.
平時我和我老公.
我們很喜歡互相稱贊.
玩那種.
偶爾他在開車扔彎的時候.
我就會贊他.
哇,爹地,這個彎你扔得實在很漂亮.
或者當他投訴我做一些不好的事.
我就會說.
是的,你的確做得比較好.
你的方法實在非常好.
我真的要向你學習學習.
我老公的回復通常都是.
你好假.
You are so fake.
我不是一個很會說大話的人.
我不是說我不會說大話.
我是說我不是很會說大話.
莫說我自己對自己是有要求的.
對著別人我希望我是一個很真誠.

$^{361}$我是real的.
我是genuine的.
我不想是一個fake的人.
有一次我老公沒空.
所以由我拿車去換油.
去做all change.
我已經很久沒有做這個動作了.
因為通常是他做的.
當售貨員告訴我換油要37塊錢的時候.
我其實嚇一嚇.
我只認為是二十多塊.
哇,37塊這麼貴.
要完成一些文件工作.
要去一個等待室.
我是帶著這個心有不甘的心情.
走去這個等待室的.
我心裡不停在想.
什麼時候加價.
為什麼現在換油這麼貴.
在想的時候.
我看著前面有個snack table.
有些零食和飲品.
我決定要拿幾包.
和拿一瓶飲品去平復我的心情.
拿下一包.
我突然間心裡出現一個意念.
我要拿一塊7塊錢.
因為我覺得應該是二十多塊.
加上稅是30塊左右.
每包東西當作1塊錢.
我應該要拿7包才夠.
我就拿了一包一包.
拿了一塊7包.
坐下.
發揮我強勁執拾的本能.
我就好好地將這6-7包的snack.
放在我的腰帶最底最深處.
因為我不想被人看到.
再想想.
等下我付錢的時候.

$^{401}$打開袋子的銷售員可能會看到.
所以我再找些東西鋪在上面.
再把銀包放在頂上.
以免我付錢的時候.
就被銷售員看到.
我覺得很得意地坐著.
拿了一塊7塊錢.
一看上面就有個監控鏡頭.
有個監控鏡頭.
我告訴你.
我第一個念頭不是怕被人看到.
我第一個想到的是.
神在看著我.
問我你在做什麼.
然後我自己也在問自己.
你在做什麼.
那些snack我不吃的.
不是我喜歡吃的snack.
就是放在家裡.
別人來的時候給你吃的.
那為什麼還要拿呢.
不夠兩秒.
我擠起來.
去到那張桌子上.
放下所有東西.
再坐下.
我真的有想過.
如果有人走過來問我.
你在做什麼.
我決定要把這件事跟他說.
或者你會想.
幾包小吃而已.
用不著說得那麼嚴重.
《路加福音》十六章十節.
人在最小的事上忠心.
在大事上也忠心.
人在最小的事上不義.
在大事上也不義.
坦白說我相信基督徒.
如果在這些所謂的小事上.

$^{441}$小罪上.
覺得很小事沒什麼所謂.
慢慢慢慢.
我們會越來越對不同size的罪.
沒什麼敏感.
愛不要虛假.
Be without hypocrisy.
之前我說過.
我會要求自己做一個很真人.
我要對人很真誠.
I don't want to be a fake.
這件事發生的時候.
其實我在亞利桑那的一個教會.
帶著一個門徒訓練.
作為他們的導師.
我覺得我有必要將這件事情跟他們分享.
認罪.
悔改.
在承認我也有軟弱的時候.
另一方面我也不想留地步給撒旦.
我要將我的罪放在神的光照底下.
罪要見光死.
在神的面前.
其實我們也沒什麼好藏起來.
因為神徹徹底底.
內內外外認識我.
知道我的軟弱.
知道我的罪.
知道我的心.
想不到我將這件事.
跟我們門徒訓練的這群太太分享之後.
他們排著隊.
一個一個認罪悔改.
大家輪流分享.
我都貪心過的軟弱.
最近做了些什麼得罪神.
整組人一起認罪悔改.
約翰一書一章九節.
我們若認自己的罪.
神是信實的.

$^{481}$是公義的.
必要赦免我們的罪.
洗淨我們一切的不義.
雅各書五章十六節.
雅各更加強調互相認罪的重要性.
你們要互相認罪.
互相禱告.
使你們可以得意志.
愛不要虛假.
來到神的面前.
我們必須承認我們每一個人都是罪人.
我們都是在接受同一個恩典.
其實我跟你平起平坐.
而互相認罪悔改這個動作.
是罪人在神的恩典底下.
願意真真實實面對自己.
放下自己.
來到神面前.
去得著赦免.
被釋放.
得自由.
得意志.
弟兄姊妹能夠互相鼓勵.
我們選擇一起宣告我們這個信仰群體.
不要在這個世界的思維方式裡面去生活.
我們要為耶穌而活.
是的,你可能覺得做基督徒很難.
因為整本聖經.
有不單止在這裡所提到的十三個提醒.
甚至還要說到謙卑自己.
打不還手.
但其實我們可以回到原本.
在神面前全人擺上.
讓神幫助我們.
由心裡面去發出改變.
不是你可以做到多少.
或者你能否做到.
是神可以直接改變你的生命.
來去做.
讓人看到我們在信一個多麼強大的神.

$^{521}$因為他可以幫助我們.
不去隨著這個世界的標準來生活.
可以做一個真真實實的人.
表裡一致真誠的基督徒.
真誠地去愛.
其實你想想是很清新的.
也是這個世界所需要的.
如果你聽完我這個分享.
只是在想.
Joyce原來是這樣的.
我只可以說.
對不起,我令你失望.
我的確不是你想象中那麼好.
個人來說.
作為一個基督徒.
我希望我自己在人前人後是沒有分別的.
我不希望台上台下的我.
是不同的.
莫說我不想是一個表裡不一的傳道人.
因為如果我是這樣的時候.
代表其實我不愛你們.
智力愛.
我們所信的耶穌的愛.
是真的愛.
是一個愛到底的愛.
智力愛到底的愛.
思想的時候.
或者你回到家裡去想的時候.
想想你信主之前.
到你信主之後認識的這個耶穌.
在你的生命裡面.
作出了什麼改變.
或者當你這樣思想的時候.
你也可以為自己祈禱.
求神赦免我們的罪.
真是願意放下我們自己.
好好地跟隨祂.
被神去改變.
羅馬書十二章一至二節.
生命的祭.

$^{561}$將身體獻上當作活祭.
這個動作應該由一個領受了福音.
心意根深而變化的生命.
最理所當然感恩合理的回應.
十二章三至八節.
生命的祠.
屬靈恩賜.
保羅關心的是信徒.
在神的恩典裡面去評價自己.
借著運用恩賜建立教會.
建立基督的身體.
今天十二章九至二十一節.
生命的祭.
智力愛.
因為主耶穌基督真真實實去到盡的愛.
叫人擺上心意根深而變化的生命.
借著恩典我們能夠去反對這個世界的醜惡.
智力去真真實實地愛.
不做hypocrite.
借著耶穌的愛.
我們去在外面跟人相交.
羅馬書十二章一至二節.
所以弟兄們.
我以神的慈悲勸你們.
將身體獻上當作活祭.
是聖潔的.
是神所喜悅的.
你們這樣做是理所當然的.
不要效法這個世界.
要心意根深而變化.
叫你們測驗何為神的善良.
順存可喜悅的旨意.
這個世界總會逼使我們.
去隨著這個世界去隨波逐流.
基督徒唯有在神的裡面.
生命被根深作出變化.
我們才可以真真正正去明白神的心意.
真真正正能夠經歷神要給我們的祝福.
你信不信神可以這樣去祝福你.
你信不信神可以藉著你願意的擺上.

$^{601}$去改變你的生命.
以至面對這個世界不同的東西的時候.
你不只是覺得束手無策.
覺得做一個基督徒很慘.
有很多東西綁著.
你怎樣去應對.
你夠不夠膽試一下.
我們一起低頭祈禱.
求你幫助我們願意將生命獻上.
當作活祭.
不是我們有什麼東西可以擺在你的台前.
而是耶穌首先為我們死.
讓我們有生命.
我們是要回應你的愛.
求你轉化我們的生命.
我們贊美你的聖明.
借著恩典叫人的價值.
建基在你的救恩上面.
弟兄姊妹謙卑服事.
建立基督的身體.
幫助我們學校你去愛.
自力愛.
go against這個世界.
去愛你所愛的.
主.
願我們的生命能夠彰顯你自己的大能.
你當受敬拜.
奉主耶穌基督的聖明祈求.
阿們.
\newpage



\section{1:19-27}
\label{sec:qy6pJ6fRtpk}
\textbf{Being A Doer, Not A Hearer (James 1\_19-27) - George Crawley}
\newline
\newline
連結: \href{https://youtube.com/watch?v=qy6pJ6fRtpk}{\texttt{https://youtube.com/watch?v=qy6pJ6fRtpk}} ~~~~ 語音日期: 2023-08-01
\newline
\newline
\hyperref[sec:2mwrzSe3J14]{\small{< < < PREV SERMON < < <}}
~
\hyperref[sec:index_chronic]{\small{[返順時目]}}
~
\hyperref[sec:index_scriptual]{\small{[返順卷目]}}
~
\hyperref[sec:LDb_e4fz_Ck]{\small{> > > NEXT SERMON > > >}}
\newline
\newline
$^{1}$All right, good morning..
Hello everybody..
God is good..
And all the time..
Amen..
I didn't know if I was going to have to explain that..
Okay, awesome..
Sweet..
So, I want to thank God and I want to thank you guys for having me here, allowing me the.
opportunity to come before you and just share a message, bring you God's word..
Today, we'll be in James chapter 1, starting from verse 19 through 27..
I listened to the live stream last week and I was excited to see that Caleb did James.
1 because I love the book of James and I figured this is great..
I can continue on this and pick up on where he left off..
So, thank you Caleb for your word last week..
I was deeply encouraged by that..
Before we get into the word, I'll ask if you guys would just pray with me real quick..
Holy Father, our Lord, our God, thank you..
Thank you for your faithfulness..
Thank you that you are a good father..
Thank you that you loved us enough to send your son to die in our place, that we might.
have hope in this world, that we would have peace, that we would have comfort in the midst.
of a world that is full of uncomforting things..
We thank you for waking us all up this morning, acknowledge that not everyone did..
We thank you for the life that you give us..
We thank you that in you we live, breathe, and have our being..
So thank you God for this life and the eternal life and hope we have in your son, Jesus Christ..
I pray that you would open our hearts, that we would receive your word, that we would.
respond to your word in such a way that Jesus Christ would be glorified in and through our.
lives..
So be with us today, teach us today that it wouldn't just be my words, but that it would.
be your words simply spoken through me and to me..
And so all these things I pray in the precious name of Jesus Christ..
Amen..
Okay, so as I said, I'll be in James chapter one, starting from verse 19 through the end.
of the chapter..
It says this, "Know this, my beloved brothers, let every person be quick to hear, slow to.
speak, slow to anger..
For the anger of man does not produce the righteousness of God..
For put away all filthiness and rampant wickedness and receive with meekness the implanted word,.

$^{41}$which is able to save your souls..
But be doers of the word and not hearers only, deceiving yourselves..
For if anyone is a hearer of the word and not a doer, he is like a man who looks intently.
at his natural face in a mirror..
For he looks at himself and goes away and at once forgets what he was like..
But the one who looks into the perfect law, the law of liberty, and perseveres being no.
hearer who forgets, but a doer who acts, he will be blessed in his doing..
If anyone thinks he is religious and does not bridle his tongue, but deceives his heart,.
this person's religion is worthless..
Religion that is pure and undefiled before God the Father is this, to visit widows and.
orphans in their affliction and to keep oneself unstained from the world.".
So brothers and sisters, today I want to focus in on this passage and specifically the verses.
in the middle, 22 through 25..
And by God's grace, do the best I can to draw out what it means and what it would look like.
for us to become doers of the word of God..
So I mentioned I was excited about Caleb preaching James because the letter of James is one of.
my favorite books in the Bible because of the impact it's had on me, my faith journey,.
my walk with Christ, and because the book itself is just so full of practical wisdom.
for how to live out our Christian lives..
There's a lot of really good allusions to the Psalms, to the Proverbs, even to the teachings.
of Jesus that James uses in his writing..
And we'll focus on some of that today..
But from the writings of the four gospels, we know that James was a brother of this dude.
named Jesus who maybe you've heard of him..
It was kind of a big deal a few thousand years ago..
Of course you guys have heard of Jesus..
But the interesting thing about James is that even though he was Jesus' own brother, he.
didn't believe in him..
Initially he doubted the things that Jesus taught, the things that Jesus did..
And there's a passage in the Gospel of Mark that tells us Jesus' family actually thought.
and said that he was out of his mind for the things that he was doing, having all these.
people come to the house and hear him preach and teach..
So they thought he was out of his mind..
His brothers doubted him..
But eventually something happens to James following the crucifixion and the burial of.
Jesus that radically changes his outlook on Jesus..
He goes from a denier, a skeptic, into a believer, a denier to a follower, a detractor to a leader.
of the early church..
Whatever it was that happened, whatever it was that transformed James, it was so convincing.
that he was even willing to die for his newfound belief in his brother, the Lord Jesus..

$^{81}$And we know this because according to Jewish historian Josephus and a few other early church.
historians, James was martyred for his faith, for not denying faith in Christ..
When I think about this, you know, going from the brother of Jesus who doesn't believe in.
Jesus to a believer, a leader in the early church and even a martyr for the faith, I.
can't help but think that such a drastic transformation to me is one of the most compelling pieces.
of evidence for the resurrection of Jesus Christ..
You look at James' transformation, you look at Paul's transformation..
What could cause someone to change so drastically?.
And I've heard it put this way, so I want to ask you guys, how many of you have siblings.
in here, brothers or sisters?.
Okay, obviously I know you guys do..
So I'm going to point you out..
Rachel, think about this now..
What would it take for you to be convinced that Nathan is the son of God?.
No way, right?.
Now we're talking about Jesus here, like, you know, I'm sure James grew up with him.
as a brother and I'm sure as a brother he wondered, like, why doesn't Jesus ever get.
in trouble with, you know, mom and dad?.
He's such a goody two-shoes..
But still, he didn't believe he was the Messiah..
He didn't believe he was God..
But something so drastic happened, he began to believe and would even die for this belief.
that Jesus truly was God in the flesh..
And truly, what else besides seeing a man once dead, buried in the tomb for three days,.
alive in front of your very eyes, bearing the marks of crucifixion, what else could.
cause such a drastic transformation in a person?.
So you see, James' encounter with the living, resurrected, eternal word of God changed James..
He was a man who once only heard the teachings of Jesus and doubted them to a man who believed.
the words and acted upon them..
In the same way, James in this passage is encouraging his audience, his audience then.
and his audience today, encouraging us to not only hear the word of God, to also respond.
to it..
So that's what I will be encouraging for all of us today..
So I mentioned that James has a lot of allusion to other passages of Scripture, to Jesus'.
own teaching, and Jesus himself taught this point about being a doer and not a hearer..
At the end of the Sermon on the Mount, one of his most well-known, powerful sermons,.
which 2,000 years later, still to this day is shaping the ethics and morality of our.
world..
And at the end of this massive sermon, you know, Matthew chapter 5 through 7, he ends,.
he gives an audience this warning..

$^{121}$He says this, "Not everyone who says to me, 'Lord, Lord,' will enter the kingdom of heaven,.
but the one who does the will of my Father who is in heaven..
On that day, many will say to me, 'Lord, Lord, did we not prophesy in your name and cast.
out demons in your name and do many mighty works in your name?'.
And then will I declare to them, 'I never knew you..
Depart from me, you workers of lawlessness.'.
Everyone then who hears these words of mine and does them will be like a wise man who.
built his house on the rock..
And the rain fell and the floods came and the winds blew and beat on the house, but.
it did not fall because it had been founded on the rock..
And everyone who hears these words of mine and does not do them will be like a foolish.
man who built his house on the sand..
And the rain fell and the floods came and the winds blew and beat against the house.
and it fell..
And great was the fall of it.".
You can hear how James is echoing the words of Jesus, hearing and doing..
So the point that he's making is pretty clear..
Just because you hear the words of Jesus doesn't mean you are truly a disciple of Jesus..
Just because you attend church regularly or listen to sermon clips and podcasts doesn't.
really mean you truly know and are known by Jesus..
There's more to this faith of ours than hearing the word..
Let's look, for example, at the disciple who betrayed our Lord..
Judas heard just about as much of Jesus' teaching as anyone else..
He heard more of it than we ever have..
But how much good did that do Judas just to hear the teachings of Christ?.
It didn't really do him a lot of good because, as James puts it, he didn't receive the implanted.
word with meekness..
Judas was no doer of the word..
Although he heard it, he was not moved by it..
He experienced no repentance, no change..
We can see that he still had this love of money that was strong enough to cause him.
to betray his brother, to betray his rabbi, his Lord for 30 pieces of silver..
There was no obedience flowing from Judas' heart, no love of God that would cause him.
to follow God..
And James says that such people are only deceiving themselves..
And as easy as it is for us to look back and kind of single out Judas, you know, and pick.
on Judas, I think it's pretty easy for us to fall into the same deception..
How easy is it for us to go to church on Sunday, get a good word in, praise God, feel really.
good about ourselves, and then we go back on Monday to a week of inaction or indifference.
to the things that we were so encouraged about the day before..

$^{161}$How often do we quench the Holy Spirit and stop the word from truly penetrating our hearts.
in such a way that it then plays out in our lives?.
It's pretty easy to hear a sermon about controlling our tongue and then leave and find ourselves.
gossiping or slandering somebody or speaking harshly to somebody..
It's pretty easy to hear a sermon on, let's say, James chapter 1, be slow to anger, and.
on the way home, driving through traffic, someone cuts us off, and then we quickly find.
ourselves angry..
And I speak from experience..
My wife is always reminding me about how I drive..
So it's easy for us to find ourselves hearing these things, sermons and scriptures that.
address the hate in our heart, the pride in our heart, love of money, love of the things.
of the world, and by the end of it, feel very motivated, very encouraged..
But do we ever actually experience any change beyond that emotion, beyond that feeling and.
that agreeing?.
A lot of times we hear these sermons and we read the scriptures and we agree wholeheartedly,.
"Yes, that's true..
Yes, that's so good..
I should make that change in my life.".
And then two, three, four weeks later, do we actually make that change?.
And even as I'm speaking, I'm feeling convicted because the other day I found an old notebook.
of mine..
It has pages and pages full of sermon notes for the past three years, and I thought to.
myself, "When is the last time I even reviewed this?.
And when's the last time I actually wrote sermon notes on Sunday and then a week, two.
weeks later, went back and checked in to make sure I was still doing those things?".
So the question I'm asking is, do we simply hear or do we also do?.
And this topic, this scripture, another reason I love James so much is because I can attest.
to this kind of lifestyle very personally..
Just telling you a little bit about myself, for a majority, 20 plus years of my life,.
I was nothing more than a half-hearted, lukewarm hearer of the word..
You see, even though I grew up going to church, I would go to church on Sundays, youth group.
on Wednesdays, prayer night on Fridays, it didn't really make a difference to me..
You see, my dad is a pastor..
I was hearing the word constantly, if not daily, and yet it never impacted my heart..
I had a lot of knowledge of God in my head, probably not as much as I thought I did, but.
it didn't actually affect my heart..
It didn't really cause any change..
So by the time I went to college, I was living in complete rebellion against God..
I was a cheater, a liar, a thief, a drunkard..
I tell people this, if you go to the book of Galatians chapter five, there's this beautiful.

$^{201}$list of the fruit of the spirit, love, joy, peace, patience, goodness, kindness, faithfulness,.
gentleness, self-control..
That was not who I was..
If you go a few verses back, Paul gives you another list, and he calls it the deeds or.
the works of the flesh, envy, dissension, rivalry, drunkenness, immorality..
If you read through that list, just about every single one of those describes my life..
That's who I was..
That's how I was..
And the saddest part of that is that if you would have asked me at that time in my life.
when I lived like that, "Are you a Christian?".
I would have been able to look you dead in the face, completely serious, and tell you,.
"Of course I'm a Christian..
Of course I'm a Christian..
I grew up going to church..
My dad's a pastor..
I got baptized when I was eight..
John 3:16, "For God so loved the world.".
I know that one..
Because of those things, because I heard those things, I was so self-deceived, as James describes..
I thought I was saved..
And it wasn't until I was about 22 years old that God opened my eyes and showed me just.
how deceived, just how lost I had been all my life..
Because there's more to being a disciple, being a follower of Christ, than hearing His.
Word..
There's even more to it than simply agreeing with His Word..
There's another piece to it where you live out His Word, where you humbly submit to His.
Word and follow His Word..
There's another passage in James that's not part of this section, but I recall reading.
Scripture when I finally decided, "Okay, my life is like totally off here..
I need to get into God's Word and learn where I went so wrong.".
As I'm reading through James, he says something along the lines of, "You believe..
Good for you..
God is one..
Even the demons believe and they shudder.".
And that really hit me because there's a difference in hearing and believing the Word of God and.
doing and responding in humility to the Word of God..
And so my life is a great example of somebody who was self-deceived..
Never in my entire life did I have any desire to be anything more than a hearer of the Word..
But in light of all that Christ had done for me at the cross, how could I not respond to.
the gospel?.

$^{241}$How could I not be forever changed, transformed from the inside out?.
Hearing His Word week in and week out, failing to respond to it, neglecting to put it into.
action, it's as foolish, James says, as looking in a mirror intently, really staring, analyzing.
yourself in a mirror, and then walking away and forgetting what you look like..
And I'm glad that I have my wife, Kiamia, to help me with this kind of thing because.
I'm very forgetful and I have like little to zero fashion sense..
So if I look in a mirror to decide, "Does this outfit look good?.
Does this match with this?".
I look in the mirror..
I think it looks good, but maybe I go and I change and I put something else on, forgetting.
how the outfit matched..
If it wasn't for my wife, I would be walking around Moody's campus looking ridiculous because.
I just, I just, I can't do it..
So I thank God..
Thank you, Kiamia..
But that's what it's like..
It's foolish..
If you, let's say you went out to eat, you know, you went out to dinner with some very.
important people in your family, in your job, and you had the best salad of your life, and.
you went to the bathroom and you looked in the mirror and you realized, "How long have.
I had this piece of lettuce stuck between my teeth?.
Have they all seen me with this lettuce in my teeth?.
Why did no one tell me?".
And then you walk out of the bathroom, but you don't remove the lettuce..
You just forget it's there..
You look silly..
You look foolish..
In the same way, James is saying it's foolish for us to be little more than hearers of the.
Word..
So again, my encouragement is that we would not simply and solely be hearers of the Word.
of God..
Let us become doers..
Let us be moved like James to respond to the Word of God, to turn from whatever sin may.
be hindering our walk in relationship with Christ..
Let's apply the things that we're learning..
I know you guys here are learning some really good things..
I've listened to the sermons..
I've been here before..
I know Pastor Mark..
He's a man who loves the Lord..

$^{281}$He loves you guys..
He loves the Word of God..
So I know you guys are getting really good teaching..
Let's respond to it..
Let's act in response to the teaching, being not just hearers, but hearers and doers of.
the Word of God..
When we become doers of the Word, when we put the principles of God into practice, we.
show that our profession of faith is genuine, and we show ourselves to not be deceived..
So what would it look like?.
What does it look like to be a doer and not just a hearer of the Word?.
Thankfully James in his practicality throughout his letter, he gives us a few examples that.
we can go off of..
And so he says this, "If anyone thinks he is religious and does not bridle his tongue.
but deceives his heart, this person's religion is worthless..
Religion that is pure and undefiled before God the Father is this, to visit orphans and.
widows in their affliction, to keep oneself unstained from the world.".
So what James is saying is that it's pointless, it's worthless to think that you're religious,.
to think you're something, to think you're right with God and you're so spiritual and.
you're so mature, or to tell people that and only be a hearer of the Word, not backing.
it up with action..
It's one thing to talk the talk, it's another thing to walk the walk..
And if you look in the Gospels, this is what Jesus is constantly confronting the Pharisees.
on..
You know, they have so much knowledge of the Word of God, and he'll ask them, "Have you.
not read?.
Did you not know?".
He's being rhetorical because if there's anyone who had read the Word of God, who knew what.
the Word of God said, it was the scribes, it was the Pharisees, the teachers of the.
law, the experts..
But even though they knew it, they chose not to do it..
They chose to do their own thing, follow their own tradition..
And so Jesus is constantly critiquing them..
So one of the examples James gives is controlling our tongues, exercising wisdom in what we.
choose to say and what we choose not to say..
If we can't keep ourselves from slandering others, gossiping, crude, inappropriate jokes,.
you know, everywhere we go, we're cussing here and there..
One, it affects our witness, our testimony to the world that we're disciples of Christ..
And so there, you know, I can tell you how much of a Christian I am, and I want to tell.
you about Jesus and how much he's changed me, but at the same time, I'm going to tell.
you some inappropriate jokes or gossip about this other person that we know or, you know,.

$^{321}$use foul language..
It's pointless..
It's worthless, James says..
All true religion, all true faith in Christ is marked by a changed heart, a heart that.
responds to the word of God with love and with obedience, a heart that is moved to care.
about the things that God cares about..
And as you read through Scripture, I mean, from the beginning, Genesis to Revelation,.
you see clearly that one of the things that God cares about is helping the helpless, being.
a defender of the defenseless, a father to the fatherless..
And so showing compassion on those on the margins of society is one of the things that.
God cares about..
And as those who want to hear God's word and respond to God's word, James is saying that.
we also should care about those like orphans, like widows, who are oftentimes on the margins.
of society..
James also tells us that true religion is marked by a heart moved to repentance, a turning.
away from and separating oneself from sin, beginning to hate the things that God hates,.
controlling the tongue, helping the helpless, avoiding the wicked ways of the world..
Three really great examples James gives, but there's so much more to being a doer of the.
word than just those three examples..
Because in general, being a doer of the word is a heart posture of humility that says to.
God, "I hear your word and I agree, I believe your word, help me to respond, help me to.
apply this to my life, help me to respond to the conviction of the Holy Spirit by agreeing.
with you and doing the things that you want me to do, turning from the things that you.
want me to turn from.".
That is what it means to be a hearer of the word, is responding in such a way to the word.
of God, whether it's in your private time reading scripture, whether it's a sermon you.
heard, a podcast you heard, applying to the lessons and principles in such a way that.
Jesus Christ is glorified in the life you live..
That's what it means to be a hearer and a doer of the word of God..
So again, over and over and over again, I say this, I encourage you and I challenge.
you guys with two things here that we can evaluate for ourselves in being doers and.
hearers of the word..
The first thing is this, ask yourself and be honest with yourself, how have I done this.
past month, the past two weeks with responding to the word of God, whether it was a sermon.
you heard here, a podcast clip, your own personal devotion time, scripture reading, how have.
I done in responding to what God was teaching me in that word?.
How have I responded?.
Did I simply hear it?.
Did I simply read it?.
Did it sound good to me?.

$^{361}$Did I just acknowledge the truthfulness that was in the word?.
Did I just gain knowledge about Jesus and his will for my life or did I act upon that.
knowledge?.
Did I step in alignment with the will of God for my life?.
Was I a doer?.
Did I act upon this word?.
Did I put any points of application into practice in my life?.
I'm sure you guys know all about, you know, the end of a sermon, three points of practical.
application..
I told you I have a notebook full of them and I don't always go back and look at them.
and evaluate with myself or with others how I'm doing in applying those to my life and.
so that's my first encouragement is to just evaluate, you know, how have I been doing,.
putting God's word into practice in my life?.
Have I been humble in responding to the Holy Spirit's conviction?.
Did I follow the ways in which God prompted me to be more like his son and less like the.
world?.
And get honest, get brutally honest with yourself as you ask the question and regardless of.
how pretty the answer is or isn't, let's resolve and dedicate ourselves to be doers of the.
word and not just hearers..
Again, I know personally how easy it is to live a life only being a hearer and being.
content with that..
Let's challenge ourselves to be hearers and doers..
And the second thing I ask, the second encouragement and the challenge is this, simply because.
I don't know a lot of you guys very well, I do know some of you, but from personal experience.
I know that you can never really be too sure, I ask this question..
To everyone here, have you ever truly responded to the word of God?.
Have you ever truly been a doer of the word or are you merely a hearer?.
Church, please listen to me..
Being a hearer, only hearing the word of God, it's possible that you could fool other people.
and have them thinking, "Man, you know, look at George..
He sure knows his Bible and he really gets into the worship and he raises his hands and.
yeah, that's a great Christian there.".
But if I'm only a hearer, it doesn't matter what I know about God or what I can tell you.
about God if I don't act upon that..
And so the truth of the matter is that there will be many, many hearers of the word of.
God in hell..
There will be many church attenders in hell, many Bible memorizers, 10\% tithers, worship.
leaders, pastors, kids, pastors in hell because not everyone who calls out, "Lord, Lord,".
is truly doing the will of God..
And what does Scripture teach us the will of God is to believe in his son, to trust.

$^{401}$in Jesus, to repent of our sin and turn to him..
James puts it this way, "Therefore, put away all filthiness and rampant wickedness," that.
is repent, look at the sin in your heart, respond to the conviction that God is stirring.
up within you about the things in your life you know are wrong and turn from those sins..
And then James says to do this, "Receive with meekness," that's to receive with humility.
the implanted word which is able to save your soul..
And what is that?.
But the gospel of Jesus Christ, the gospel of grace that you're a sinner and there's.
nothing you can do to save yourself..
So Christ did it all for you..
He lived the perfect life..
He truly was the son of God despite what his brothers might have at one time thought..
And being the son of God, he went to the cross humbly, willingly, dying on the cross, paying.
the price of your sins, the price you couldn't pay so that by believing in him, trusting.
in him, turning from the world, turning from sin, turning to him, you would be saved, you.
would have a promise of eternal life..
This is the word of God..
This is the implanted word which when received with meekness and humility is able to save.
our souls..
So again, in light of this, the question that stands is, are you, am I, are we doers of.
the word of God or are we hearers only?.
The last thought I want to leave you guys with is again, looking back to the example.
of James and even the example of Paul, if you're familiar with his transformation, their.
lives were undeniably radically impacted by the grace of God once they believed in the.
Lord Jesus Christ..
And I'm not saying that all of us will have such a drastic conversion experience as Saul,.
as Paul..
I personally don't know anybody who walked around blind for three days with scales on.
their face, but spiritually speaking, we will have that transformation because it's the.
same Holy Spirit that changed James' mind, the same Holy Spirit that changed Paul's mind..
Scripture even says the same Holy Spirit which raised Christ from the dead lives inside of.
us and will cause transformation inside of us so that we can go from those who only hear.
the word and are unaffected at the core of our being to those who hear the word and humbly.
receive it and respond to it in our lives..
James and Paul, their lives bore fruit which testified to the fact that they were not merely.
hearers of God's word but doers who acted and responded upon it..
So the only answer, the only question that I can really leave you with is this, are you.
a doer of the word of God?.
Amen?.
Amen..

$^{441}$Thank you..
[BLANK AUDIO].
\newpage



\section{6:1-5}
\label{sec:LDb_e4fz_Ck}
\textbf{Exploring The Land: Little Town Jericho (Joshua 6\_1-5)- Rev. Mark Chan}
\newline
\newline
連結: \href{https://youtube.com/watch?v=LDb_e4fz-Ck}{\texttt{https://youtube.com/watch?v=LDb\_e4fz-Ck}} ~~~~ 語音日期: 2023-08-08
\newline
\newline
\hyperref[sec:qy6pJ6fRtpk]{\small{< < < PREV SERMON < < <}}
~
\hyperref[sec:index_chronic]{\small{[返順時目]}}
~
\hyperref[sec:index_scriptual]{\small{[返順卷目]}}
~
\hyperref[sec:NkubOgp9wRQ]{\small{> > > NEXT SERMON > > >}}
\newline
\newline
$^{1}$to join us and we're going to have some fun time together..
Exploring the land, we're going to, I guess,.
finish our series on this city called Jericho..
So today we're going to talk about Jericho.
and I have two main points..
One main point is I'm going to explain to you.
what Jericho is all about..
It's a little town, not Bethlehem,.
but little town of Jericho..
And then the importance of little things, all right?.
And I'll get to that..
So first of all about Jericho..
Jericho is a very well-known city,.
well-known story in the Bible,.
and who would you associate with when I talk about Jericho?.
Who is the guy who conquered Jericho?.
Who is the military general?.
Joshua, oh yeah, right, there you go, Joshua, right?.
Joshua led the Israelites to conquer the very first city.
as they entered the land of Canaan,.
and the land of the life of God, right,.
in the presence of God..
So the first city that they conquered or took over.
was Jericho, and it's a very famous city..
Not only is this a very well-known story in the Bible,.
the kids know it and the adults know it,.
but not only that, it is actually in human history,.
in the secular world history,.
Jericho is one of the most famous towns in the world.
because it is being known as the oldest city in the world..
Oldest city..
Now what I meant was humans, before they settle down.
and build their own city, for years, for centuries,.
there have been hunters and gatherers,.
and they go from oasis to oasis,.
from water source to another water source or food source,.
and they don't have their own home,.
until they somehow, they got together.
and organized themselves and built their own city..
And the first city that they built.

$^{41}$by ancient people was Jericho..
You're going back to 10,000 BC..
That's kind of crazy, right?.
Those who went to a few museums,.
you read some signs and you find some display.
that are so old, right?.
Very, very old..
So Jericho was the oldest city in the world..
And then you have to understand why they gathered together.
to organize the first urban society in Jericho,.
because Jericho had a spring..
And the spring, later tradition associated with Elijah..
It's called Elijah's spring,.
because there was a story in later scripture,.
talks about the prophet Elijah,.
he did a miracle to heal the water of the spring..
So that's why later traditions associate Elijah,.
the prophet, with the spring in Jericho..
In any case, that is the oldest city in the world,.
and also because there's Elijah's spring..
And the spring is still working, it's still operating today..
So this is Jericho, and think about Jericho,.
especially for those who grew up in the church..
You pick up the children Bible,.
and sometimes you find some pictures in the Bible,.
talks about Joshua and the Israelites conquering,.
taking over the city of Jericho..
You will find some pictures like this..
Big city wall fell down, right?.
As they shout out, cry out the victory,.
they are blowing the trumpet,.
and then the huge wall of Jericho,.
all of a sudden they fell,.
and then they took over the city..
It's a dramatic, wonderful story to remember, right?.
It tells you the power of faith..
If you really believe in God,.
if you really believe in the promises of God,.
and you act by faith, then the biggest city,.
the biggest wall, the tallest wall,.

$^{81}$will fall down in front of you..
Now, but as I emphasize one thing,.
we have to read the Bible story in its own context first..
All right, if you go back to the time of Joshua,.
it was about 1200, 1300 BC..
What kind of city was Jericho?.
It was actually a very, very small town..
This is the aerial picture of Jericho..
You see the curvy road on the lower left?.
Above that curvy, modern road,.
is Jericho..
Ancient town, the oldest town of Jericho..
And it's not that big..
It was, according to the most accurate estimation,.
it was just one acre..
How big is one acre?.
Any good student here?.
One acre..
How big is one acre?.
It's smaller than a city block in Chicago..
Smaller than one city block, right?.
It's a very small town..
It's about 200 feet by 200-some feet, right?.
It's actually smaller than Chinatown..
That is Jericho..
It's actually a small town, right?.
So by some, most of the most reliable calculation,.
in the time of Joshua, Jericho had about.
probably 200 to 300 people, and that's it..
It's a little town, but a strategic place..
You know why?.
Because if you check the Bible map,.
you understand it is guarding the eastern entrance.
of the land of Kingdom..
That is Jericho..
And Jericho is a strategic place,.
although it's a small, small town..
Now that may be surprising to all of us,.
because I grew up in the church,.
and I read Children's Bible,.

$^{121}$and Jericho was being imagined.
as some sort of very, very big city,.
and the wall is like the Great Wall of China..
And then the Israelites circle around the city,.
and then every single day,.
and then the last day they circle around it for seven times..
And then finally the wall just fell down, right?.
It's a very dramatic story,.
and we exaggerate the scene by later tradition,.
by later imagination..
You go back to the Bible,.
the Bible is quite careful in describing.
some of the cities of those days..
For example, after they took over Jericho,.
the next city to be taken over was Ai, Ai..
Ai, Ai..
And you see the description in 7.3..
When the Israelites returned,.
the spies returned and informed Joshua,.
they said, "There's no need for all of us.
"to go up and attack Ai..
"It won't take more than 2,000 to 3,000 men to attack Ai,.
"since there were so few of them in the city of Ai..
"Don't make all of our people struggle to go up there.".
All right, so take it easy, guys,.
because Ai is such a small town, right?.
It will only take 2,000 to 3,000 people..
And by the way, the word "thousand".
should be translated as "military unit.".
So the people are saying to Joshua,.
"Well, we checked out Ai already..
"If you want to take it over, you just need, what?.
"Two to three football teams, right?.
"And that's all you need..
"You will take over Ai.".
But of course, they fail at the end..
However, you see the description..
The Bible actually is very careful.
to describe some of the targets,.
Joshua and the Israelites,.

$^{161}$as they conquer the land of Canaan..
The second town is Hathor..
Hathor, 1110 in Joshua,.
it talks about Hathor killing the king and captured the city..
Hathor had one time been the capital.
and the head of all these kingdoms in the land of Canaan..
So Hathor is a big, big city..
That's what the Bible is telling you..
It is the head of all the kingdoms,.
head of all the cities,.
one of the most, the biggest city in the land of Canaan..
In fact, history tells us this is Hathor, the aerial picture..
You see the upper city.
with some buildings on it and trees on it,.
and then the upper right,.
the flat land with some farms,.
that is the lower city of Hathor..
Altogether, the size of Hathor is 200 acres..
How big was Jericho?.
How many acre?.
Jericho..
One, yes, one acre..
Hathor is 200 acres, 200 times of Jericho..
So chapter 11 talking about the taking over Hathor..
In fact, Hathor is the biggest site in Israel today..
It's a huge city in times of Joshua..
Now, as I said, the Bible is very careful.
to describe the cities..
And it's a small city, Ai, it tells you,.
there are only a few people there..
And Hathor, it tells you it is the head of the kingdom..
However, with that information in mind,.
you understand when the book of Joshua talks about.
the Israelites taking over different cities,.
it's very interesting..
For Hathor, the biggest city in time of Joshua,.
the book of Joshua only uses 50 acres..
15 verses to describe the battle..
They did take over Hathor, but the Bible only spent.
about 15 verses to talk about the battle of Hathor..

$^{201}$However, Jericho, the Bible used entire chapter,.
chapter six, to describe the first battle of Jericho..
If you talk about the previous context, chapter six, right?.
The previous context would be chapter one to five..
All talking about preparation of taking over Jericho..
So the Bible at least used six chapters,.
six chapter long content to talk about.
the first battle of Jericho..
The two spies went in there and they met someone.
called Rahab, the prostitute, right?.
And the prostitute Rahab changed side..
And she sided with, she abandoned the kingdom of God.
and went with Yahweh, the God of Israel..
Spent six long chapters to talk about.
one battle, the first battle..
We're gonna read some of the verses.
and let me pick chapter six, the first five verses..
Let's take a look..
Why does the Bible spend so much time.
and so much content on describing the battle of Jericho?.
Such a small town, one acre..
Now, the gates of Jericho were tightly shut.
because the people were afraid of the Israelites..
No one was allowed to go out or in..
But the Lord said to Joshua, "I have given you Jericho..
"It's king, all his strong warriors..
"You and your fighting men should march around town.
"once a day for six days..
"Seven priests will walk ahead of the Ark of the Covenant,.
"each carrying a ram's horn..
"On the seventh day, you are to march around the town.
"seven times with the priest blowing the shofar, the horns..
"When you hear the priest give one long blast.
"of the ram's horns, have all the people shout.
"as loud as they can..
"Then the walls of the town will collapse.
"and the people will charge straight into, up to the town.".
Now, this is the first description.
of the battle of Jericho..
If you check out the entire chapter six of Joshua,.

$^{241}$there are only three parts..
The first part is God's instruction.
to the Israelites, to Joshua..
That is what we have just read..
And the second part is the men of Israel.
executing the plan..
And then the third part is kind of a conclusion,.
talking about the victory and the salvation of Rahab,.
the king and the knight in the town..
Now, we're not going to talk about the details here..
You can go home and read about the battle of Jericho..
What I want to say is this is the first battle,.
first battle of a small town, right?.
And it's very interesting when you check out.
the book of Joshua, what we expect the book of Joshua.
to emphasize is strategy, military resources,.
how they prepare military strategy-wise,.
tactic, how to take over the first town, Jericho..
Last night I was with my family..
We went to this restaurant to have dinner together.
because my wife, it was her birthday yesterday..
Anyway, so we went out last night,.
and she's not here, that's okay..
So we were waiting for the table, right?.
And there were so many people in the restaurant..
So we went down a couple blocks.
and went into a used bookstore, used bookstore..
So we went into the bookstore, and there were so many books.
there, and I went past the section on military history..
There were so many different books.
talking about military history..
World War I, World War II, and then the Civil War,.
Vietnam War, and all this kind of famous battles in history..
So out of curiosity, I picked one of those,.
and then I just skimmed through it..
And then it talks about who are the heroes,.
and most of the heroes are the military generals..
And most of the military history genre.
would emphasize on military tactics, right?.
How they move around and maneuver all these resources.

$^{281}$and take over some of the targets..
These are what we expect from the military history genre..
You take a look at the Book of Joshua..
It really doesn't do that..
It's very strange..
It's a strange book, right?.
Some people would expect the Book of Joshua.
because it talks about battles and military stuff..
It has to be a man's book, right?.
The men like to do that, right?.
Men like military stuff, don't you?.
Right, all the men here?.
You love fighting..
You love wars and battles, right?.
The dramatic battles, right?.
The important battles..
Not so much for women, maybe it's just overgeneralization,.
but that's the way it is..
You take a look at the Book of Joshua..
It's very disappointing..
It doesn't really talk about military tactic.
and all this geography and all this military strategy.
and all this stuff..
They don't..
It's a very strange book..
Now let me tell you some of the data..
When the Book of Joshua talks about taking over Jericho,.
it uses 102 words only..
102 words..
For such an important city, the first city, 102 words..
However, in the previous chapter, chapter two,.
talking about Rahab, how she changed side.
from the Canaanite gods to Yahweh, the true God of Israel,.
chapter two uses 86 words to talk about a woman.
turning to the side of Israel,.
believing in God, right?.
And then make a big statement about the God of Israel.
being the true God..
86 words..
So it's about the same portion, same amount of words.

$^{321}$talking about Rahab's story, same amount of words.
talking about the battle of Jericho..
And in fact, in the portion that talks about Rahab,.
the prostitute, chapter two uses five verses,.
five long verses giving Rahab some of the most important.
lines in the Old Testament..
Some of the direct speech that Rahab said.
to the spies of Israel was actually statements of faith..
A theologically sound, solidly scriptural words.
that she talks about how the Lord is the true God,.
the Exodus, how the Lord has given the land to Israel.
and the Lord is the true God..
And all that stuff, five verses,.
is a public theological declaration.
spoken by Rahab, the prostitute, the Canaanite prostitute..
You skim through the book of Joshua,.
it's very disappointing from a military,.
typical military genre perspective..
The heroes, the heroes,.
the heroes are not what we expected,.
not Joshua, not this fighting man of Israel,.
but a woman, prostitute, Rahab..
That is what the Bible says..
And try to emphasize the importance of small characters.
and the importance of small deeds done by them..
They matter to God..
And in fact, they are more important.
than the military fighting and dramatic content..
The Bible put our focus on something.
that is more important in the sight of God..
The small characters, Rahab, all these unknown,.
unnamed people,.
they believe in God, they act in faith..
That matters to God..
So there's, this is the importance of small deeds,.
that's my second point..
In the book of Joshua, the small characters matter to God,.
matter to the biblical history,.
matter to what God is doing on earth..
You and me, we are the small characters..

$^{361}$We are not the famous people, right?.
We don't have a lot of followers on our IG, right?.
On your TikTok, right?.
Nobody take notes of you..
But you matter to God..
The small characters doing small deeds of faith..
Couple months ago, in June,.
we took a family trip to New York..
So we went to this military museum.
called Intrepid Museum in New York..
It's a fascinating museum..
So we went into a retired submarine,.
and then we went to a retired aircraft carrier,.
and it's just amazing, there's so much to see..
There's so much more to see than a film museum..
So I really enjoyed that day,.
and then I went into this gallery of aircrafts,.
and it has a one-to-one model of space shuttle..
It's fascinating, right?.
We don't use space shuttles anymore, right?.
It's a retired program..
So I was kind of taking pictures and reading.
all this information about the space shuttle program.
and all that, and then it just dawned on me.
there were two space shuttle disasters that happened.
in the history of the program,.
both related to space shuttle accidents..
Now, back in 1986, those who are older like me,.
would remember back in 1986,.
NASA space shuttle Challenger.
exploded seconds after launching,.
only 73 seconds after lifting off..
And the disasters claimed the lives.
of all seven people in the aircraft..
All people on board died, including a teacher,.
the first civilian in the space program..
And after the disaster, investigation revealed.
what was wrong..
It was actually a small, old rain seal.
that got brittle because of the low quality.

$^{401}$of the old rain seal..
It becomes brittle in cold temperatures..
And that small, old rain fail,.
and then the entire aircraft, flames broke out,.
and then the spacecraft just explode in the air,.
right in front of the eyes..
Exploded, and then disintegrated..
That was back in 1986..
The other big disaster was 2003..
Columbia, and other space shuttle..
Now, this time, the accident did not happen.
when it takes off..
It actually took off, and then went into space,.
and did its thing, and then on the way back.
to the atmosphere..
We entered the atmosphere, the space shuttle.
disintegrated into pieces in the sky of northern Texas..
Again, killing all seven crew members..
Investigation revealed what's wrong..
It's actually a small piece of foam insulation,.
fell off, and then it damaged the left wing,.
and then, as the spacecraft continued to go.
in supersonic speed, it disintegrated into pieces..
All people died..
Now, the long investigation of the entire.
space shuttle program tells us there was something wrong.
actually from early on..
It's that the crew members or the engineers.
trying to put together the aircraft,.
and program it, and all that,.
they actually made a lot of small mistakes,.
or small getaway..
Some of the lower quality stuff that they used,.
or they didn't pay attention too much..
And the engineers thought, "Oh, we could get away.".
And they did..
And time after time, they got away,.
but at the end, disasters happened..
They called the entire process of,.
"It's okay, we can get by.".

$^{441}$That sort of mentality, that sort of stuff that they did..
The entire process they called normalization of deviance..
Deviance is deviating from the standard practice, right?.
It's only a small deviation from the standard practice..
Not quite up to par, but a little bit less deviance,.
but they normalize it..
It's okay..
We make excuse..
All the small things,.
and the entire aircraft disintegrated..
I think this is a perfect example for all of us,.
because if you think about it, life,.
it consists of many mundane things,.
many routine things..
Sometimes we speak of spiritual discipline,.
spiritual habit, right?.
All of those things are really, really small..
All of these changes that God expects us to do.
are actually very, very minor, right?.
In the big scheme of things, right?.
In this generation, there's a phrase called.
to be voked, W-O-K-E, right?.
To be voked, what does it mean?.
To be voked means to be aware of.
important issues in the world..
Racial justice, social justice,.
those are the big things that we need to pay attention to..
To be voked, that is what it means, right?.
There's nothing wrong with that..
We should concern about justice issues..
We should, the big issues..
But you know, think about it..
Life often operates.
on small details..
The small and mundane things,.
they all matter to God..
You think about your character growth,.
your spiritual attitude toward other people,.
your attitude toward things entrusted to you..
Don't we all find small excuses?.

$^{481}$We all do, right?.
I do, right?.
And we always say to ourselves,.
what's the big deal, right?.
If I miss this, if I don't do that, right?.
If I'm less than faithful,.
I kind of deviated from the standard,.
that's all right, right?.
We're still alive, right?.
What's the big deal, God?.
What is the big deal?.
I can do it later..
I can procrastinate, right?.
Because there's no big deal..
The story of Jericho tells you,.
it's a big deal..
The small deeds are big deals..
Church life, the standard practice.
of a normal Christian life,.
how we relate to other people..
When we are on the wrong, we apologize..
We ask for forgiveness..
Those are really small things, right?.
It's not dramatic..
It's not something that will change the world,.
but they matter to God..
The story of Jericho tells you life,.
it's like a cathedral..
A cathedral is built by stone after stone.
and stone after stone..
All these pieces have to fit together.
in order to build a big building..
Your life is like that..
If you're missing one of those,.
someday it will collapse..
Someday it will crumble..
The growth of your spiritual life is the same way..
It's the small habit add on small habit.
add on small habit..
And the habit will shape you,.

$^{521}$will mold you into the image of God..
Without those habits, without those small deeds,.
your life cannot be built, right?.
Let's be reminded..
In the sight of God, small deeds matter..
The greatness or the maturity of your spiritual life.
or your character as a whole.
is measured by many, many small deeds..
And what God expects us to do.
is to be faithful with all the small deeds..
And it will add up..
It will pave you a path to be successful.
in the sense that to be spiritual,.
to be God-pleasing, to be Godly.
in your character growth..
Let us all remind ourselves and take a look at yourself..
What kind of small things am I giving excuses to lately?.
That I tell myself, what is the big deal?.
Chill out, God..
Can you chill out for a little bit?.
I just don't want to do that, right?.
Those are really, really important.
even though they are small..
Just like Jericho..
What is your Jericho?.
Let's all pray together..
Father, we come before you and.
we know that all the details of our lives.
are important to you..
And sometimes they are even more important..
They're not as big issues..
We just pray that your spirit will remind us all..
Help us to see that we need to be faithful.
not only with the big stuff,.
but the small details..
(piano music).
Remind us and help us to do some self-examination..
To look into our lives..
Are there any habits that we have forgotten?.
Are there any small deeds that we have.

$^{561}$given excuses to?.
Procrastinate, delay?.
Help us to be pleasing to you..
Thank you for being our God..
Help us to grow in such a way.
that will build up our character..
That our being will be shaped by your grace..
And help us do all the small deeds faithfully..
So that our life can be built up one stone at a time..
One layer at a time..
Help us to grow into the image of your son..
Remind us how we speak..
Every single word is important..
Remind us all of our meditation,.
our thought life,.
all the details are important..
How we meditate on your word,.
how we reflect on spiritual insights..
Help us build up our habits,.
our lives so that we can be pleasing to you..
In Christ we pray..
\newpage



\section{}
\label{sec:NkubOgp9wRQ}
\textbf{6/9/19 主日講道:神竟然比我和祂講數 (張麗仙傳道)}
\newline
\newline
連結: \href{https://youtube.com/watch?v=NkubOgp9wRQ}{\texttt{https://youtube.com/watch?v=NkubOgp9wRQ}} ~~~~ 語音日期: 2023-08-13
\newline
\newline
\hyperref[sec:LDb_e4fz_Ck]{\small{< < < PREV SERMON < < <}}
~
\hyperref[sec:index_chronic]{\small{[返順時目]}}
~
\hyperref[sec:index_scriptual]{\small{[返順卷目]}}
~
\hyperref[sec:T08JyDUvISQ]{\small{> > > NEXT SERMON > > >}}
\newline
\newline
$^{1}$.
弟兄姊妹好,今天開始我想跟大家做兩分鐘的練習.
就是你和你旁邊的弟兄姊妹說一說.
你知道聖經裡面究竟有多少卷書呢?.
新約幾卷,舊約幾卷,還有你最喜歡哪一卷呢?.
兩分鐘.
.
隔離前後都隨便找個人說說好嗎?.
(閒聊).
差不多了吧?.
(閒聊).
好了,我們回來了.
我不知道你得到什麼答案.
但是我知道有些人一定會說.
麗仙,拜託你問問題不要這麼低調.
但是每天都有人說.
我只會創世記和詩篇.
還有其他.
聖經好像有些阿太的書.
有白記和阿太福音.
所以無論你是對聖經很熟悉的.
或者不熟悉的.
我都想給你一點面子.
如果你很熟悉的話.
究竟聖經對你來說是頭腦知識呢?.
還是能夠改變你的生命呢?.
如果你對聖經是很不懂的.
我勉勵你,今年不如努力一點.
無論你是否熟悉,我都勉勵你們.
不如你們更加認識聖經.
我們教會有探索課程.
在很多團契都有查經的時間.
不如今年決定參加某一兩個.
我們一起在聖經的真理中成長.
老實說,一個人在家裡讀聖經.
很多時候都很難明白.
我們實在需要有人帶領.
需要大家彼此勉勵.
才會有更加進步.
而來了成會教會半年半的時間.

$^{41}$我一直都帶著一班.
很負責任的查經組長一起查經.
每一個月我們會預查一次.
預查了之後,他們就會回到團契那裡.
帶大家查經.
我們覺得這樣的查經是非常有益的.
不單止是更認識神.
我們更加認識自己.
而且我們發覺在人生裡.
這些話語對我們有很多提醒和勉勵.
早期我們查了幾本書.
《雅各書》,《爾弗所書》.
這兩本書都是新約的.
而最新鮮熱辣的.
我們剛剛查完的.
就是剛才Cindy也有提到的.
不知道她是否知道.
就是《馬拉基書》.
《馬拉基書》是舊約聖經.
如果你熟悉聖經.
馬拉基書是舊約聖經裡.
小先知書裡最後的一卷.
很多人都知道.
不過更加有很多人.
一提到馬拉基書就想起這件事.
每個人都是這樣的.
馬拉高書.
不要緊,你說馬拉高書.
什麼都好,我不介意.
但最重要的是.
我希望你真的明白.
究竟馬拉基書或馬拉高書.
是在跟你說什麼.
神藉著這本書跟你說什麼.
而且我告訴你.
你讀過馬拉基書都知道.
不是像馬拉高書那樣鬆化.
不是那麼軟熟.
因為馬拉基書裡的道理.
其實是很硬心硬肺的.

$^{81}$你讀下去就明白.
為什麼硬心硬肺.
但無論如何,貫穿這本書.
我看到一個令我驚嘆的精髓.
這是今天我想跟大家分享.
就是神竟然讓我們人跟祂講數.
馬拉基書就是一卷這樣的書.
神讓人跟祂講數.
什麼叫講數呢.
講數包括投訴.
大家爭辯,講條件,解釋.
包括這些東西.
目的是什麼呢.
是要取得公道.
就像在家裡.
夫妻之間,大家都會講數.
你還是這樣,我就怎樣怎樣.
如果父母和子女之間.
你還是玩遊戲機,打機.
不做好功課.
我以後怎樣怎樣.
我們在公司.
有時候都會跟上司講數.
我們在教會的弟兄姊妹.
在祭奉的時候都會跟上司講數.
有時候跟牧師講數.
講數是想取得公道.
神在馬拉基書上.
讓人跟祂講數.
我記得去年在港台上.
Mark Sir有一個詩篇系列.
他說原來我們在苦難裡.
我們可以向神質疑,埋怨.
我們可以跟祂喊叫.
聽完這個道之後.
有不少的姐妹跟我說.
原來可以這樣.
真的使我向神開了眼界.
原來可以這樣.
我可以告訴你.

$^{121}$馬拉基書也是讓我們向神開了另一個眼界.
就是神竟然這麼大胸襟.
讓我們跟祂講數.
整本書都是圍繞祂的六個講數.
一般理解我們都知道.
講數是人與人之間.
高高在上的神.
祂什麼都是對的.
為什麼祂會讓我們跟祂講數呢?.
為什麼祂會讓我們跟祂頂撞呢?.
但是聖經告訴我們.
我們的神的確是這樣的神.
在馬拉基書.
神讓人跟祂講數.
我特別看到兩方面.
我今天想跟大家分享.
第一.
神讓人跟祂講數之後.
祂是很耐心地去解釋.
跟你解釋.
我今天給你一個例子.
在馬拉基書裡面.
第一個講數就是.
神的子民以色列.
跟神說.
神你不愛我們.
可以這麼說.
馬拉基書第一章.
一開始神說.
我一直都愛你們.
不過如果你看中文聖經.
我要警告你一件事.
因為中文聖經是翻譯得非常不對的.
它寫我曾愛你們.
好像現在不再愛.
但是在希伯來文裡面.
是說我一直.
以前愛你們.
現在一直都是愛你們.
為什麼神的子民說.

$^{161}$神你不愛我呢?.
其實是有原因的.
因為當時神的子民.
他們是經過了亡國.
70年.
做奴隸.
已經是很慘的了.
大家可以想像.
之後.
雖然70年之後.
波斯王改變了奴隸的政策.
放回這些以色列人.
回到自己的國土.
雖然他們回到自己的國家.
可以重建聖殿.
重建圍牆.
但是回到國家100年之後.
覺得國家.
經常都是很垃圾.
不像以前那麼風光.
聖殿很垃圾.
農作物不好.
國家不富強.
大家的生活都非常艱難.
所以他們埋怨.
神,你都不愛我們.
如果愛我們.
我們怎麼會搞到這樣的田地.
這是馬拉的說話.
這是一個人對神的投訴.
但是在這裡.
神的回答是很奇怪的.
我讀給你們聽.
大家看一看螢幕.
(文字).
是不是.
看起來一點都不明白.
你記住問題就是說.
當時以色列人.
向神投訴說.

$^{201}$你不愛我們.
神說了這番話.
祂說.
你說我不愛你.
不如你看看.
你們隔離的那一族.
就是以素的那一族.
又稱為以東.
你看看他們怎麼樣.
我洗他們的山頂芳良.
我將他們的地業交給抗日的野狗豺狼.
他們想重建.
我都拆毀.
聽起來很殘忍.
為什麼神這樣回答他們.
說不愛他們呢.
但是我們不要忘記.
在公義聖潔的神的眼中.
這是犯罪的人.
應該有的懲罰.
神是不讓他們重建.
神是不再理會他們.
不再理會他們.
說起來.
其實世人都犯了罪.
無論是雅各,以素.
就是以色列或者以東一族.
因為你和我.
其實我們都是充滿罪孽.
我們其實都應該受這個懲罰.
但是這段話告訴我們.
神告訴向他投訴的以色列人.
我愛你們.
因為就算你們多壞.
就算你們不是比以東人好很多.
你們一樣這麼壞.
但是就算你們這麼壞.
我也願意給你們機會回轉.
其實是你們不配得.
完全是我的恩典.

$^{241}$以色列你埋怨我不愛你們.
但是你們有沒有想過.
你有沒有敬畏我.
我一直以恩慈地帶領你們.
但是你們有沒有敬畏我.
沒有,你們竟然帶入很多的偶像.
你們生活腐化.
所以你們得到亡國的教訓.
你們還不知道.
今時今日還向我埋怨.
不過你看不看到.
即使如此.
你亡國,我有給你重建的機會.
你們埋怨現在回到國家.
好像不是這麼輝煌風光.
但是你睜大眼睛看看.
旁邊的以東.
他們想重建都沒有機會.
以色列你不要搞錯.
你們不是比他們好.
事實上你和他們一樣是罪惡滿盈.
在我眼中全部都看到.
不過是我曾經揀選你們.
而且我和你的祖宗亞伯拉罕.
摩西亞國納納戈約.
所以我對你的愛.
你要知道是一個蒙約的愛.
是基於我的揀選.
不是你配得.
所以有聖經學者說得很好.
神對以色列的愛.
不是基於人感情的愛.
因為人的感情可以升可以跌.
神對以色列的愛是蒙約的愛.
有個學者說.
選擇與不選擇.
我對你的愛是基於我揀選了你.
不是因為你值得.
不是因為我對你的情緒上的回應.
神說我對你的愛是蒙約的愛.

$^{281}$所以我一直都愛你們.
縱使你們曾經離棄我.
縱使你們曾經幾富化.
縱使你們拜幾多偶像.
但是你離棄的日子.
你離棄我的日子.
我沒有離棄過你們.
我好像一個慈愛的爸爸.
由你們在這些患難中成長.
以致成熟.
究竟你們看不看到我對你的愛.
我想我一直說的時候.
做父母的特別明白.
有時你看到兒女都會跟你說.
爸爸媽媽你不愛我.
我想有這個你又不給我.
我想有那個你又不給我.
但是你知道那些不是對他好的.
你知道有時他們在所謂的受苦中.
反而使他們成長成熟.
打少一點機.
對他們來說是好事.
你心知肚明.
神對我們的愛就是這樣.
神說你不要懷疑我對你的愛.
這是神對當日向他投訴的以色列民.
說的一番話.
弟兄姊妹.
我們今天不是舊約的時候的人.
我們不是舊約的時候的以色列人.
但是我們都因為藉著耶穌基督的寶血.
神給我們立了一個新約.
我們都在這個盟約的愛裡面.
你和我都是被神揀選.
我們因著信耶穌成為了神的子民.
記住這個愛不是因為你和我特別好.
而是因為神揀選我們.
是因為祂的恩典.
那你有沒有懷疑過神對你的愛呢?.
前天我們才有姊妹旋.

$^{321}$每一個月我們一次聚會.
所以每一位姊妹都很珍惜回來.
只不過往往到那一天.
很多時候這個姊妹有事.
那個姊妹有事.
有時因為病.
有時因為家裡有突發事.
突然間就不能理會.
這些只是一個例子.
告訴我們人生真的很無可預測.
有很多的患難.
一個大浪這樣衝過來.
有時你覺得人生是很無情.
很無情.
你在那些境況裡面.
在那些困境裡面.
我相信你和我也有同樣的感受.
究竟神在哪裡呢?.
神你究竟愛不愛我呢?.
我們都會向神投訴.
但是神藉著《瑪麗基書》再一次提醒我們.
你要堅信我對你蒙約的愛.
我是不會變的.
在衝擊的日子.
你對我有信心.
你要相信我恩守愛護.
我會一步一步地帶領你.
你只要堅定相信.
你一定會體會到.
從我你而來.
是天上而來的祝福.
剛剛說了.
神被人投訴.
第一個例子就是《瑪麗基書》.
第一天.
他們說你不愛我.
還有的.
當以色列民族的同神說話.
神你不公義.
在《瑪麗基書》有很多篇幅.

$^{361}$神的子民投訴神不公義.
例如.
你看看那些聖職人員那些祭司.
他們獻祭的時候犯罪.
你都容忍的.
他們說神不公義.
但是在這個時候.
神很耐心地跟他們說.
你不要以為我不公義.
有一天我會用火和漂泊的繭.
來煉證這些聖職人員.
好讓他們悔改回轉.
他們的錯處.
在我眼中.
我全部都看到.
你不要以為我不理.
什麼叫火和漂泊的繭.
這是一個比喻.
是指神會用一些嚴厲的教訓.
有一天臨到這些聖職人員.
在中間.
使他們明白.
使他們悔改.
神的審判對他們.
其實是想他們悔改.
神的子民對他另一個控訴就是.
好人和壞人.
行善行惡都沒有分別.
我們時時都看到惡人得福.
這是另一個控訴.
但是神耐心地解釋說.
好人說的話.
我全部都聽著.
聽在我的心裡.
好人的善行.
我都記在我的紀念冊裡.
惡人是會受到審判.
審判那天.
渣都沒有剩下.
因為他們不敬畏我.

$^{401}$因為他們作惡地對撞我.
這是神在這裡對以色列的投訴.
神你不公義.
他的回應.
弟兄姊妹.
你和我今天在生活裡.
有沒有受到不公義的對待呢?.
我相信一定有.
有時我們會.
因為這些不公義的對待.
無論是在工作的地方.
在家裡.
在教會.
在什麼地方.
我們受到不公義的對待.
是很激氣的.
公義這件事其實是很重要的.
是令人氣憤難平.
有時你不平.
還會令你掉進憂鬱的深谷裡.
或者在不公義的日子裡.
我們都會埋怨神.
就好像以色列人一樣.
但是在《馬拉基書》裡.
神向我們說得很清楚.
祂的胸襟是大的.
你向祂投訴吧.
你投訴吧.
祂會接受的.
但是祂也盼望你.
聽到祂這些話.
公義在神的眼裡.
是極之重要.
公義是神的個性.
所以.
神一定會處理不公義的事.
你要相信祂.
你要保持一個心.
去敬畏,尊崇祂.
有一天.

$^{441}$你會見到神公義的臨到.
到目前為止.
我們看到《馬拉基書》很特別.
神給人投訴祂.
說你不愛我.
你不公義.
之後就耐心解釋.
但是我們不要太得意.
以為只有我們投訴祂.
因為《馬拉基書》裡.
也有神向祂的子民投訴.
說你們也有什麼不對.
我舉一個例子.
就是憲制的問題.
神說祂的子民對祂極不尊重.
在憲制的事情上.
祂的子民完全是亂來的.
不過我還沒說到.
神如何投訴他們的時候.
還沒說到.
他們究竟犯了什麼事的時候.
我讀兩段聖經給你們聽.
你們會看到.
他們如何回應神對他們的投訴.
他們的態度很差.
神說:我既為父親.
尊敬的我在哪裡?.
我既為主人.
敬畏我的在哪裡?.
你們卻說:我們在何事上.
藐視你的名?.
神說:你們將污穢的食物.
放在我的壇上.
且說:我們在何事上污穢你?.
另外一段.
神說:你們要轉向我.
我就轉向你們.
你們卻問:我們如何在事轉向?.
神說:人豈可奪取神的物?.
你竟奪取我的公物.

$^{481}$你們卻說:我們在何事上.
奪取你的物?你的公物?.
神說:你們在當立的十分之一.
和當現的公物上.
你們看這些紅字.
是以色列民對神指責他們的回應.
他們投訴神的時候.
投訴得很過癮.
但當神投訴他們的時候.
他們突然就裝作反叛青少年.
我哪有做過這些事?.
我何時有做過?.
其實這就是他們的態度.
當時的憲制問題是什麼呢?.
當時神的指紋.
他們對憲制非常藐視.
因為在《神比摩西律法》裡.
清楚地說明.
你們的憲制.
無論是牛羊還是飛鳥.
都是要無殘疾的.
因為我是聖潔的神.
你們不可以獻上有殘疾的.
但當時他們的憲制.
將盲眼的,跛腳的,有殘疾的,病的.
這些牛羊或者飛鳥獻上.
聖經裡說得很清楚.
而且他們是經常的.
不是偶然發生的.
神很氣憤.
他還說了一句話.
如果你送東西給省長.
你的省長.
你也不會送這些東西吧?.
為什麼你會送這些東西給我?.
這是神對他的回應.
我不知道他們當時.
因為生活情況不好.
所以他們發難.
送這些東西去憲制.

$^{521}$又或者有可能.
是因為他們曾經被擄過去巴比倫.
波斯這些國家裡.
他們見過很多其他國家的神明.
他們就覺得.
耶和華我們的神.
也是其中一個神.
所以他們對神的敬畏.
是降格的.
所以他們的憲制是很藐視.
這是當時他們的行為.
而且不單是平民這樣做.
祭司也這樣做.
大家都這麼壞.
所以神氣憤到一個極點.
他說了一句話.
你們再這樣.
我會在你們中間.
有一個人關上電門.
就是說禮堂的門.
關上,鎖上.
不讓人進來.
不用你們的憲制.
神說了這句話.
我根本不需要你們的憲制.
我不是等你們這些來吃飯.
我等這些來看你們對我的敬畏.
但你們竟然給我這些東西.
其實這裡對我們今天.
也有很大的提醒.
說真的.
母親節.
你會送爛的手袋給媽媽嗎?.
不會.
你送禮給上司.
會不會拿回家吃.
那盒朱古力.
吃了幾顆.
送給他?.
不會.

$^{561}$或者你想送他即食的.
但會不會送爛的給他?.
不會.
這些是很荒謬的.
這些是很不敬的.
但是.
神在這裡所說的.
以色列人的憲制.
就是類似這樣.
他們是送單眼的.
跛腳的.
病了的.
一些生畜.
飛鳥給神.
很明顯.
他們對神非常不尊重.
所以,弟兄姊妹.
這裡給我們也有很大的提醒.
我們今天.
來到神的殿裡面.
我們獻上的是什麼呢?.
我說的不只是金錢.
正如保羅說的.
你們要把身體獻上.
當作活祭.
是我們整個人.
究竟我們獻上什麼給神呢?.
金錢上.
我們的事奉上.
我們是不是獻上單眼的.
跛腳的.
殘疾的呢?.
這個值得我們反省.
而神對他們的投訴.
憲制上的投訴.
不只是在質量上.
或者獻上這些爛東西.
而且還是在質量上.
大家都經常聽到.
正如Cindy也提到.

$^{601}$九一的時候.
神特別的鄭重對以色列文化.
你們要獻上十分之一的奉獻.
十分之一的奉獻.
是將你的修成獻上十分之一.
這是一個責任.
主要是供奉聖殿和聖殿負責的人.
所以是利微之派.
所以獻上十分之一.
是你們對神的一個.
宗教上的一個責任.
但他們沒有.
還有,鄭華夏說.
神要求人奉獻十分之一.
但其實《舊約聖經》不只是十分之一.
十分之一只是硬性規定.
你要十分之一的奉獻.
來供奉聖殿的運作.
但其實在十分之一以外.
很多時候.
以色列文的感恩.
或者贖罪.
是另外的貢獻.
當他們感恩的時候.
或者當他們要求贖罪的時候.
他們是會再有奉獻的.
但當時以色列文是怎樣呢?.
他們的感恩祭.
他們的贖罪祭.
他們不獻上應該貢獻的.
這是神要指責他們.
簡單來說.
他們無論是十分之一.
還是當獻的公物.
他們都很不願意.
送一些隨便的.
最便宜,最不像樣的.
有時候還不送.
這是他們當時的問題.
神要指責他們這些問題.

$^{641}$神指責他們.
這些冷淡藐視的奉獻行為.
所以神說了這句話.
「你們奪取我的物」.
人豈可奪取神的物呢?.
你們竟然奪取我的公物.
因為你們不奉獻.
其實你們不奉獻.
就是你們是奪取了神有的東西.
人豈可奪取我的公物呢?.
你們通國的人都奪取我的公物.
所以我的咒詛.
我的懲罰要臨到你們身上.
當時的咒詛懲罰是什麼呢?.
神是洗到黃蟲.
割在他們的土產上.
葡萄樹早生產而失落.
這個時候.
以色列人他們深深體會到.
神的懲罰臨到.
因為當時的農作物生產極之惡劣.
但他們竟然不知道.
原來這是神給他們的懲罰.
要神對他們說明.
這是我給你們的懲罰.
所以你們要記住.
我的指紋.
不要只是我愛你們.
你們要愛我.
尊崇我.
是你們的神.
你們要作出應有的奉獻.
而且是從心裡面作出.
這是神對他們的講述.
講到這裡.
你會覺得.
又講自己的呃.
牧師傳道.
久不久在台上都會講錢.
因為久不久就提醒大家.

$^{681}$如果不是大家弟兄姊妹合作.
我們就不能發工資.
又講錢了.
我告訴你.
教會是要講錢的.
但我不是這樣看.
我不是說.
我們講錢是因為發不發工資的問題.
因為教會的事工.
有沒有錢運作的問題.
而當我們講錢的時候.
是再一次提醒我們.
哪一個是主.
我們有沒有敬畏我們的主.
就好像當日的以色列民.
他們沒有敬畏他們的主.
當神的懲罰臨到他們的時候.
他們才知道.
他們不是他們擁有的主.
那個才是主.
各位弟兄姊妹.
今天我們講錢.
我們在教會裡面講錢.
都是因為再一次提醒我們.
哪一個是我們的主.
我們有沒有敬畏我們的主.
不要以為這些東西是我們雙手賺回來的.
我們就是這些的主.
我們要明白.
神是一切的主.
不過.
在神講的懲罰裡面.
要臨到的時候.
懲罰之先往往都是神的恩典.
所以在《馬拉基書》裡面繼續這樣講.
神說我再給你們機會.
你們要將當那十分之一獻上.
使我家有糧.
以此試試我是否為你闖開天上的窗戶.
傾福於你們.

$^{721}$甚至無處可容.
萬國必稱為你們有福的.
因你們的地必升為喜樂之地.
如果你們真的這樣敬服我.
各位弟兄姊妹.
神的懲罰審判對以色列再一次告訴我們.
其實.
祂是想他們回轉.
而神的懲罰之先往往.
是一個恩典的提醒.
你們要回轉.
你知道我一直講的時候.
我希望什麼.
我好希望我們成會教會.
不用回轉.
我好希望我們今天就已經懂得.
怎樣尊崇我們的主.
在奉獻上.
在事奉上.
我們每一日.
每一時候.
我們都是獻上最好的.
我們知道祂是我們的主.
我們敬奉祂.
以至我們能夠一起這樣.
受到神從天上傾下來.
傾下來的福氣給我們.
以至無處可用.
《馬拉基書》裡面還有很多很好的教訓.
不過因為今天的時間關係.
我就講到這裡停.
但是在這裡我想和大家有一個回頭的總結.
特別是對你.
如果你是剛剛回到教會.
你是剛剛聽聖經的人.
我希望這個總結再一次的幫助你.
明白得到的就是.
《馬拉基書》再給我們看到.
我們的神是一個怎樣的神.
可以說我們對神開了一個眼界.

$^{761}$就是整本書都告訴我們.
神是有這樣的胸襟.
讓我們和祂講述.
高高在上的神.
竟然這麼不可思議.
讓我們和祂講述.
而講述的目的.
是讓我們明白祂對我們的愛.
是不會變的愛.
是蒙約的愛.
不是因為你不好.
祂就不愛你.
你不好,祂仍然愛你.
但是祂一直呼喚你.
你要回轉,你要回轉.
神的講述.
另一個目的.
不是要懲罰我們.
而是要我們明白我們自己的問題.
讓我們有機會去回轉.
以致我們能夠體會.
祂從天上傾倒給我們的福氣.
在我們的人生裡面.
在我們人生的各方面.
講完《馬拉基書》.
我想跟大家有個預告.
就是在六月開始.
城會的柴經班就開始一個新的書信.
就是一個和《馬拉基書》很不同的.
性質的書信.
也是一個很重要的書信.
因為這個書信可以說是.
基督教很多很多重要的神學.
都在這裡是一個基礎.
羅馬書.
在這裡我特別邀請各位弟兄姊妹.
一個月一次.
就是每逢第四個禮拜五.
七點半.
我們有這個柴經班.

$^{801}$歡迎你回來.
和我們一起在真理上面成長.
祈禱.
你的胸襟是這麼廣闊.
即使我們有時候.
很多時候都因為罪而頂撞你.
即使我們真的很多時候都不敬.
教我們知道.
教我們再一次的深信.
你是很愛我們.
教我們以愛去回敬你.
你不需要我們做什麼.
給什麼錢.
其實你需要的是.
我們這個真心對你的回應.
你是體重.
主求你幫助我們.
求你教我們以一個回敬你的心.
經歷你從天上給我們的福氣.
求你祝福我們每一個人.
求你祝福我們教會.
能夠榮耀你的名.
多謝你聽我們的禱告.
奉救主耶穌的名.
\newpage



\section{2:7}
\label{sec:T08JyDUvISQ}
\textbf{Learning to Love Dependence (Genesis 2\_7) - Sam Johnson}
\newline
\newline
連結: \href{https://youtube.com/watch?v=T08JyDUvISQ}{\texttt{https://youtube.com/watch?v=T08JyDUvISQ}} ~~~~ 語音日期: 2023-08-14
\newline
\newline
\hyperref[sec:NkubOgp9wRQ]{\small{< < < PREV SERMON < < <}}
~
\hyperref[sec:index_chronic]{\small{[返順時目]}}
~
\hyperref[sec:index_scriptual]{\small{[返順卷目]}}
~
\hyperref[sec:1rUpXzwI0NI]{\small{> > > NEXT SERMON > > >}}
\newline
\newline
$^{1}$her today, Sam..
(audience applauding).
- Good morning..
I'm honored to be here and excited to be here..
Privileged to be with you in worship today..
Thank you to the worship team for leading us..
I really enjoyed that time of worship in the Lord..
And this morning we're gonna talk about something.
that is really important for us,.
especially in the modern world that we live in,.
and that is learning to love dependence,.
learning to love dependence..
And that's a difficult thing for us..
Dependence is not a cultural value.
that we are used to embracing..
We live in a culture that really values,.
especially in the United States,.
really values individualism and independence,.
self-reliance, self-sufficiency..
We all feel like that's how we should live.
and that's how we wanna be..
We feel like we don't want to need anyone else..
We don't want to have to rely on anyone else..
We wanna be strong and get through things on our own..
But the reality is that's not a biblical value..
And in fact, what the Bible teaches us is the opposite..
The Bible teaches us that we are hopelessly dependent.
on God and that we're also,.
God made us to be interdependent on one another.
in the church..
And so our cultural instincts in the world around us.
kinda fight against that..
We say, no, I don't wanna be dependent on God..
I don't wanna be dependent on others..
I wanna do everything on my own..
But we have to learn to embrace the fact that we're human.
and that God made us with limitations.
and that we need him and that we need one another..
So that's what we're gonna speak about this morning..
We're gonna look at this verse in Genesis chapter two,.

$^{41}$verse seven, just one verse..
And if you're not familiar with what's going on.
in Genesis two, this is when God is creating..
And so in Genesis one, God's created the universe.
and everything in it..
He's created the earth, he's created the animals.
and the plants that are on the earth..
And in Genesis two, Genesis one is like an overview.
of creation, it gives all the different things God created.
and the days and when they were created..
And then Genesis two is like zooming in on day six,.
which is when humans were created.
and giving us more detail about when we were created..
And so in verse seven, it tells us, speaking about us,.
it says, "Then the Lord God formed the man of dust.
"from the ground and breathed into his nostrils.
"the breath of life, and the man became a living creature.".
That's our origin story right there..
The origin story of humanity,.
not as glamorous as we think, right?.
We didn't get bitten by a radioactive spider.
or come from an alien planet..
We got made out of dust and dirt..
That's our origin story, that God formed us from dust.
and he literally breathed life into us..
And not just life, but our soul.
and everything that it means to be human..
We know in, earlier in Genesis 1, 26 and 27,.
it says that God made us in his image..
And one of the things that we need to understand about God.
that helps us to understand ourselves is who God is..
And one of the most significant aspects.
of Christian theology is the truth.
that God is completely distinct.
from all of the created order..
God is in a class to himself..
Everything else that has been created,.
everything that we interact with,.
everything in our world was created or made..
The natural things were made by God..

$^{81}$Buildings and things like that were built by humans..
But everything we see and interact with.
was created at some point.
and it had a beginning point in time.
and it will have an end point in time, including us..
But that's not true of God..
God is completely different.
than everything in the created order..
We see in Genesis that God, in the beginning,.
God, he already existed before the beginning of the universe..
God was pre-existing, he was eternal..
He needs nothing from no one to exist..
That's who God is..
And that's important to help us understand who we are..
Because God made us in his image,.
but not the other way around..
We are like God in some ways..
God made us in his image..
So God gave us a soul, this spiritual part of us.
that can have relationship with him..
God gave us some of the things that we see in our own lives..
We see the appreciation for beauty and creativity.
and a desire for relationship..
Those are reflections of God..
God obviously loves beauty and he's very creative..
And he is a relational being..
Even in the beginning of Genesis there,.
we see God as father, son, and the Holy Spirit.
in the act of creation..
God's a relational being..
And we are too..
We're made in God's image..
We're like God in some ways..
But God is not like us in any way..
And so we have to understand,.
when we start to think about who God is,.
God is completely beyond us..
And we are finite creations of the infinite God..
God is unlimited in every way..
He has all knowledge, he has all power..

$^{121}$His attributes are unlimited..
He is not just loving, but he is perfect love..
He's the definition of what love is..
That's how we know love, because of God..
He is just..
That's how we know what justice is..
Because God is the definition of perfect justice..
He is holy..
That's how we know what's right and holy in the world,.
because God is holy..
He's the definition of those things..
God is infinite in every part of who he is..
But we're not..
He created us as finite creations..
We're a part of the created order..
Now we're a very special part of the created order..
God created us as the crowning achievement of his creation..
He gave us, again, his image.
and things that no other part of his creation has..
We have a soul, we have the ability.
to make moral choices for right or wrong..
We have the ability to have relationship with God.
and relationship with one another..
Those are really unique things..
But we are not God..
We are not infinite..
We are finite in our capacity..
And even in the attributes of God's character.
that God calls us to reflect,.
we reflect those things in a finite way..
God's love is perfect..
We love imperfectly..
And we should seek to love like God..
God calls us to that throughout scripture,.
but still, our love is imperfect, it's incomplete..
God is perfectly holy..
Our holiness is always flawed..
We are finite..
And a part of our finitude, the finite part of us,.
is that we have limitations..

$^{161}$In verse seven there that we just read,.
God formed humans from the dust of the ground..
How did he form us?.
He formed us in the way that we.
see one another in this room, right?.
God gave us spiritual bodies, physical bodies..
He gave us spiritual life as well,.
through the universe and do different things..
He gave us this physical container for us to live in..
And that physical container has some limitations, right?.
One of the biggest ones is we can only.
be in one place at one time..
As much as we would love to be in multiple places at once.
and feel like we could get so much more done,.
God gave us that limitation..
We can only be in one place at a time..
But there's a lot of other limitations too, right?.
We get tired..
All of us have to sleep..
We can work hard, we can hustle, we can do a lot,.
but then we have to rest..
That's just a reality of having a physical body..
And what I want us to see is that our limitations.
are actually a part of God's good design..
Sometimes we fight those limitations..
We say there's something wrong with me.
because I can't do more and I can't do.
all the things I wanna do and what's wrong with me?.
But the reality is our limitations.
are a part of God's good design..
God gave us those limitations on purpose.
because we need them..
And one of the biggest reasons that we need them.
is to make us depend on him..
We're limited creations that are dependent on God..
And one of the marks of God's good design in our lives.
is the limitations that he gave us.
that force us to rely on him..
I'm gonna let you in on a secret here..
Most of you probably already know this,.

$^{201}$but some of you may not..
But here's the secret..
It took me a long time to learn in my life..
But the secret is that you cannot have it all in life..
All right?.
And our world, most people in the world.
don't seem to know that or maybe they know.
but they just don't believe it..
And they think they can have it all..
And so at every stage of life,.
we are trying to live into that,.
that we can have it all..
So when we're young and we're a student,.
we're trying to get the best grades.
and also do all the extracurriculars.
and maybe some of the sports.
and also be the best rated gamer.
and be a great child to our parents..
We're trying to do it all.
because apparently we think we can.
and apparently we think we have to, right?.
Then we get a little bit older and we're working.
and we're trying to have it all in our career.
but also live this really cool life.
where our social media reels are full of great pictures.
and vacations and restaurants and beaches.
and we're trying to do it all..
And then if we have a family,.
we find ourselves trying to still have it all in career.
and everything else but also be a great parent.
and a great spouse and we're trying to do that..
And the reality is we can't..
God never intended it for us to have it all..
We are limited creations..
And if we're gonna live in a healthy way,.
then we have to embrace that and say,.
you know what, I'm dependent on God..
Think about what we see in this passage..
First, we're dependent on God for our physical life,.
for our very breath..

$^{241}$God breathe the breath of life into humanity..
And we could say, well, this is just talking.
about the first man, the first woman,.
it's not talking about us..
But the Psalms speak very clearly to the fact.
that God forms each one of us before we are born..
God, it says, God knits us together with his own fingers..
God is in there in the mess of your DNA.
and not just in the DNA..
The DNA is the physical part, right?.
God formed you how he wanted you to be,.
with your face and your height.
and your physical characteristics and all that..
But not just in that, God forms you in your emotions,.
in your personality, in your gifts, in your talents,.
in all those ways..
God forms you just like he formed the first humans..
So we're dependent on God for our physical life..
We're dependent on God for our emotional life..
We're dependent on God for our spiritual life..
All of those things, those dreams and desires.
and ambitions that we have,.
ultimately, we're dependent on God for those things..
James chapter four talks about this..
It says, come now you who say today or tomorrow.
we'll go into such and such a town.
and spend a year there and trade and make a profit..
Yet you do not know what tomorrow will bring..
What is your life?.
For you are mist that appears for a little time.
and then vanishes..
Instead, you ought to say if the Lord wills,.
we will live and do this or that..
As it is, you boast in your arrogance..
All such boasting is evil..
What is that passage saying?.
It's saying, look, you don't know it,.
but you're dependent on God,.
even for your future plans and the things.
that you think you're gonna do next week or next year..

$^{281}$If God doesn't sustain your life and your emotions.
and sustain you spiritually,.
you're not gonna do those things..
That's what that verse is reminding us of..
We, as human beings, we don't like dependency..
We try to convince ourselves.
that we're not actually dependent..
We love feeling like we're in control,.
even if it's just an illusion..
We're very good in the modern world.
at building a lot of structures around us.
that insulate us from difficult events.
and maintain the illusion.
that we're in control of our lives..
But we're really not..
We're dependent on God.
to get us through the day physically..
We're dependent on God emotionally..
When we go through times, as Pastor Mark was mentioning,.
where we are grieving or we're experiencing loss.
or we're hurting or something unexpected happens.
or we have a disappointment, we lose a job,.
we have a relationship where someone is unkind to us,.
we go through those things,.
or at least when I go through those things,.
I'll speak for myself..
Maybe it's not true for you, but it's true for me..
Sometimes I go through those things.
and emotionally I just get really wrecked.
and I realize how close I am all the time.
from just really struggling emotionally.
and getting stuck.
and just being in a really unhealthy place..
So if I think, oh, I'm doing great because I'm so strong,.
I'm lying to myself..
I'm dependent on God all the time for his strength.
and for his comfort and for his peace.
and to sustain me emotionally and certainly spiritually.
for the relationship that I have with God..
That's not gonna happen at all unless I'm depending on him..

$^{321}$If I'm relying on myself saying, yeah, I'm a great person.
and I know what I'm doing.
and I'm out here doing all these things because of me,.
I'm gonna be very spiritually unhealthy..
We have to get past our sinful desires.
and the lies of our culture that it's all about us.
and we need to rely on ourselves.
and we need to say, you know what?.
I am a limited creation and I need God in my life..
I need him to keep me going..
I pray this prayer pretty much every morning..
I prayed it this morning 'cause I was tired..
Just my wife and I, Jessamine is my wife's name..
We have five kids..
Our youngest is one.
and so he had trouble sleeping the last few nights..
So I've missed some sleep..
So I was tired when I woke up this morning.
and I prayed this prayer that I pray every morning.
before I even get out of bed..
I just say, God, I need you today..
I'm really dependent on you..
I need you physically to get me through the day..
Just give me the strength and the energy.
to get through this day..
I need you emotionally and spiritually.
to give me the patience that I need.
to deal with situations that are gonna arise today.
or people that are gonna give me difficult interactions.
or whatever it is..
I need you to get through this very day..
That's a healthy place for me to be..
It's a healthy way for me to start my day..
Then we see in Genesis as well.
that we are limited creations.
that are interdependent on one another..
Again, that's a part of God's good design..
God gave us limitations.
that not only make us dependent on him.
but they make us need other people..

$^{361}$And we're relational beings.
and God created us for community.
or another word that we could use is family.
with one another..
Genesis 2, 18, just a few verses after this in chapter two..
God looks at the man..
This is before he makes the woman..
God looks at the man and says this..
The Lord God said, it is not good.
that the man should be alone..
It's not good for us as humans.
to be isolated and on our own..
As much as we think it's good,.
we think oh, but then I would never be disappointed.
if I didn't have anybody else that I had to trust.
or no one would ever hurt me.
if I could just isolate myself.
from these other people who are being mean..
But God said, hey, it's not good for you to be alone..
You are relational..
You need other people..
And in Genesis 2, 24,.
God institutes this covenant of marriage..
He begins the first family between Adam and Eve..
But he doesn't just begin the first family..
He begins the first community..
Before there was Adam, there was an individual..
After that, there was community..
There was a group..
God made us to be in community, in family..
And we need that..
And we need to embrace that..
And one of the ways we embrace that is right here..
The church is a community that God gave to us.
for our own benefit..
And it doesn't matter if you're married or single.
or young or old or even what your background is..
God in the New Testament talks about that,.
how there's neither Greek or Jew,.
circumcised, uncircumcised,.

$^{401}$barbarian, Scythian, slave, free,.
but Christ is all and in all..
That God brings us together in the church..
It doesn't matter who we are or what our past was..
And he says, you're in a family now..
And so you can rely on each other and care for one another.
and depend on each other..
And again, the human part, our sinful nature says,.
I don't want that, I don't wanna need other people..
I don't wanna need people in my church..
Sometimes we don't even wanna need our own family members..
I don't wanna rely on my family..
I'm not even sure I like them..
But God says, you need those people..
I made you that way on purpose because you need help..
And if you didn't have the limitations,.
you wouldn't go to anyone else..
You'd just try to do it all on your own..
That limitation is a good thing for us that God gave us.
so that we would rely on others..
And I just wanna say real quickly here.
that those limitations are not necessarily sin..
Look at the order of Genesis, what's happening here..
Genesis one, the creation..
Genesis two, like I said, it zooms in.
and gives us this kind of zoomed in view of humans.
and the creation of humans..
Then Genesis three, what happens?.
Genesis three is where sin enters the world..
So think about that..
Genesis two, we're created as limited beings.
and there's no sin in the world yet..
It means our limitations are not sin..
And a lot of times we feel like,.
oh, well, I just can't do enough and be enough.
because I'm broken, there's something wrong with me.
and I guess I just need to try harder..
I need to figure out a way to be better..
And we treat our limitations like there's some type of sin.
or brokenness..

$^{441}$God never said that..
Now, our limitations can lead us to sin.
if we don't respond to them correctly..
If we hit the wall of our limitations and say,.
you know what, I need help.
and so I'm gonna turn to God for help,.
that's exactly what God wants us to do..
But if we hit the wall of our limitations and say,.
no, I'm gonna push through and I'm gonna find a way.
and I'm gonna do this on my own,.
then yeah, our limitations can lead us to sin..
But our limitations are not sin..
And so when we look at our life and say,.
I just can't do it and I can't say yes to everything.
and I can't keep up with it all.
and I can't take it all on, that's okay..
That's okay..
God didn't make you to do all those things..
God made you to be able to do a limited number of things.
each day in a limited amount of time.
before you have to sleep.
and you can only be in one place at a time,.
you can't meet every need,.
you can't be friends with everybody,.
you can't do everything that you want to do..
And that's tough for us in the modern world, right?.
In the past, maybe it was easier.
because you didn't see pictures all the time.
of what everybody else was doing,.
so you just lived in your community.
and you felt like, hey, I'm doing okay,.
I'm here in my community with my family,.
with my people, with my church,.
we're doing what God has for us..
It might have been a little bit easier.
to kind of keep your eyes.
on what God was doing in your life..
But in the modern world, with the internet and social media,.
we're constantly seeing what everybody else is doing.
and so we have this desire.

$^{481}$to always feel like we should be doing more,.
like we're not doing enough,.
or that we're not enough..
Like, you know, I'm not good enough.
because I see everybody else.
that seems to be better than me.
or more talented or their life is more exciting..
Those are not the things that God has told us..
God, it's clear here that God created you.
to be who he wanted you to be..
He intentionally gave you the talents that you have,.
but he intentionally gave you the limitations you have..
And he wants you to be dependent on him.
and dependent on others in a healthy way.
and to not try to do it all..
Our limitations are not sin..
And I want us to see here that dependence on God.
is cooperation, it's not resignation..
Sometimes we might be mistaken.
and say, okay, if I'm just gonna depend on God,.
that means I'm just gonna kinda resign.
and be like, I can't do it..
And so whatever happens in life is just gonna happen.
and I'm gonna give up because I'm overwhelmed, right?.
But that's not what the Bible teaches us..
It's not a posture of resignation where we just kinda quit..
The New Testament's clear that we are still.
supposed to be actively engaged in seeking to know God.
and seeking to draw near to him,.
seeking to surrender to his truth.
and seeking to obey the commands of scripture.
and glorify God in our lives..
So we're not just quitting,.
but we are cooperating with what he is doing in our life.
and allowing him to lead and depending on him.
and being honest with ourselves and saying,.
God, I can't make myself..
It's one of the interesting things about this passage,.
Genesis 2, seven, it gives us that origin story,.
that moment of formation,.

$^{521}$but the reality is God continues to form us.
throughout our lives, emotionally, spiritually..
And if we're a follower of Jesus,.
if we're someone who by grace and through faith.
has put our trust in Christ,.
then God is going to form us as a new creation.
into the image of Jesus..
He's gonna make us look more and more like Jesus.
and that's an ongoing process throughout our life..
And so we're supposed to be engaged in that..
We're supposed to be surrendered to that..
We're not supposed to just give up..
Dependence is a posture of cooperation.
where I admit my limitations and my need.
and I rely on God's work in my life.
and I cooperate with all of my mind and my soul.
and my strength with what God is doing..
It's not giving up..
Ephesians 2, 10 says,.
"For we are his workmanship created in Christ Jesus.
"for good works, which God prepared beforehand.
"that we should walk in them.".
It's a beautiful verse..
And it tells us that God has some good stuff for us to do..
But who is the focus in that entire verse?.
God..
We're his workmanship created in Christ Jesus,.
which God prepared beforehand..
It's all about what he's doing, not about what we're doing..
There are good works that we're gonna do,.
but it's through relying on him..
And that means we're actively involved, but we're dependent..
One of the beautiful fruits of dependence.
is that it can give us peace..
As I said earlier, having limitations is not sin..
And understanding that we can't do it all.
is actually a very peaceful place to be..
When we get to that point where we say, you know what?.
I'm not gonna be able to please everybody in my life..
I'm not gonna be able to commit to everything..

$^{561}$And instead, I'm gonna depend on God and say,.
God, what's really important for me here?.
What's important for me in this job?.
My boss keeps saying that I need to do this.
and meet these standards..
What do you have for me here?.
What am I supposed to be doing?.
And I'm gonna live into that..
That's a much more peaceful place.
than just trying to please all the people around you.
who are constantly saying, we wish you'd do this,.
we wish you'd do that, we want you to be more,.
we want you to give more..
It's a peaceful place when we can say, you know what?.
I was involved in this at work,.
I was involved in this other thing,.
someone asked me to volunteer over here,.
but I can't do that all..
I can't do all that stuff because I'm a limited human..
I'm gonna seek God about what God wants me to do,.
and then I'm just gonna pick one thing,.
or I'm gonna pick two things,.
or I'm gonna pick three things..
We have different capacities at different times in our life,.
but I'm gonna match what I'm doing to my capacity,.
and then I'm not gonna worry about the rest..
I'm gonna admit that I'm limited..
That's a place of peace..
And those things change in time of life..
I remember when Justin and I got married,.
before we had kids, I was serving as a youth pastor,.
and we did a lot of stuff in ministry,.
'cause we had more time than we had kids,.
and when our kids were little,.
they just got dragged along to everything,.
and did it all with us, and we still did quite a bit,.
but not as much as we had been doing,.
and as our kids have gotten older, that's changed..
It's not that we're doing less,.
we're actually doing more than we ever did,.

$^{601}$but it's just in different ways..
We invest more in our kids,.
and we still invest a lot in ministry,.
but not as much as we used to,.
'cause I can't have it all..
I can't be the dad that God wants me to be,.
and also be the number one preacher guy.
who's putting out stuff all the time..
I couldn't be that guy anyway..
I'm not talented enough to be that guy in the first place,.
but even if I was trying to be that guy,.
I wouldn't have enough time to be the dad.
that God wants me to be, and so I have to choose..
I have to say, God, what is important?.
What do you want me to do?.
You want me to be a husband and a dad first,.
and so I'm gonna put my time there,.
and I'm gonna give the time that I can give to ministry,.
but I'm okay with the fact that there's some things.
I'm gonna say no to, and there's some things.
that other people are gonna say,.
you know, if you would do this,.
it would really work out great for you,.
and I'm gonna have to say, that's probably true,.
but I can't, I don't have enough time,.
because God's given me other things to do,.
and that's a part of our life..
Our capacity changes..
I'm older now than I used to be..
I used to run a lot faster than I do, jump higher..
I played sports, all that stuff..
A few weeks ago, Zara, she's seven, she's sitting there..
We were at the park about two weeks ago.
playing a little family baseball game on the Diamond,.
just kind of hitting the softball and running the bases,.
and I was playing with him and running,.
and she goes, "Dad, I've never seen you run that fast.".
(laughing).
But it's true, I don't run nearly as fast as I used to,.
but I also don't run that much anymore,.

$^{641}$'cause it's a different time of life,.
and that's okay, that's a part of our limitations..
Aging is a part of our limitations..
And just like I said earlier,.
our limitations are not sin, but they can lead us to sin..
Some people in midlife, they hit that wall of aging,.
and they say, "No, I'm gonna fight against it,.
"and I'm gonna go out and blow all my family's money.
"on a new car, and I'm gonna leave my marriage.
"and do all this.".
It can lead us to sin, right?.
But if we're honest with our limitations, we say,.
"Hey, I'm not as fast as I used to be.".
Things are different now, that's okay,.
I'm gonna rely on God, and let him equip me.
to love the people in my life in different ways..
Maybe I can't do all the things I used to do,.
but God can make different ways for me to minister.
and love and care for people..
Then we can maintain health,.
and that's what's really important..
Embracing our dependence gives us a lot of peace,.
when we're able to say, "You know what, I can't,.
"and that's okay.".
That's where God wants me to be..
Embracing our dependence gives us freedom,.
when we're able to say,.
"You know, that gap between the desire to have it all.
"that we feel and that the world tells us,.
"and then what we can actually achieve,.
"that gap can give us a lot of stress,.
"and a lot of anxiety and a lot of worry.".
When we realize, "Here's what I wanna do,.
"and here's what I expected that I was gonna be able to do,.
"and here's what I thought I was gonna be able to do,.
"and here's the reality of what I can achieve.".
As a limited person, with a certain number of hours a day.
that can be in one place at a time,.
this is really as far as I can get..
And there's always a gap there..

$^{681}$We always want more than we can get to..
Every day is like that for me..
I'm a very optimistic person,.
so I start every day thinking I'm gonna get.
like 20 things done, and usually I get three to five..
And sometimes, if I'm not relying on the Lord,.
I'll get really stressed, 'cause I'm like,.
"I didn't get 20 things done, and now what am I gonna do,.
"and now I've got so much to do tomorrow,".
and I get full of stress and worry and anxiety..
Embracing our dependence gives us freedom.
when we say, "You know what?.
"I can do what God enables me to do..
"I'm gonna depend on Him to make sure.
"that I'm doing the right things,.
"and the rest of it, I'm gonna let go of.".
There's a lot of freedom in that..
Even our mission in the church,.
to make disciples of all nations,.
way bigger than anything we could ever achieve.
on our own, right?.
That is such a huge mission..
We need to embrace our dependence on God and say,.
"God, yes, we want this..
"We wanna serve you..
"We love you..
"We wanna be your people..
"We wanna make disciples..
"We wanna reach our community.".
But we can't..
These obstacles are way bigger than we are..
We need you to show up..
We need you to be at work..
And when we embrace our dependence on God,.
we're able to be free from the pressure.
of filling roles that were never ours to fill,.
trying to do things that we were never supposed to do.
in the first place..
We can be free from the stress and worry of anxiety,.
of carrying burdens that we were never supposed to carry..

$^{721}$Think about that in your life, in my life..
Again, I'll speak about myself..
A lot of the times when I get stressed.
and anxious and worried,.
it's because I'm trying to, in some way, know the future..
Was I ever supposed to know the future?.
No..
That's a human limitation..
We do not have the ability to know the future..
God does, but we don't..
So when I try to take that on,.
I'm definitely gonna get stressed out..
There's no way I can know the future..
Or I try to control other people..
Or I wanna have control over something like my health.
or someone else's health..
God never intended for me to have that ability..
And so if I'm taking that on,.
I'm just gonna stress myself out and frustrate myself..
There is freedom in just saying,.
God, this is a you thing..
This is not a me thing..
I can't..
You know, Matthew 11, 29 to 30 says this..
Take my yoke upon you and learn from me,.
for I am gentle and lowly in heart,.
and you will find rest for your souls,.
for my yoke is easy and my burden is light..
Jesus said that..
Now how could Jesus say that?.
How could he say his yoke is easy and his burden is light?.
'Cause this is the same Jesus that also said,.
if you wanna be my disciple,.
you have to take up a cross and follow me..
This is the same Jesus that also said,.
in this world, you will have troubles..
So when Jesus said his yoke was easy and his burden was light.
he wasn't saying there's not gonna be any problems,.
there's no adversity if you follow me,.
it's smooth sailing for the rest of your days..

$^{761}$That's not what he's saying..
What he's saying is if you follow me and you do it right,.
the yoke is easy and the burden is light.
because you're not carrying it, he's gonna carry it..
But for that to happen in our lives,.
we have to embrace our dependence on him..
And the last thing I wanna point us to.
is that embracing our witness is a dependence to others..
In a world that so values individualism.
and hustle and self-reliance,.
being someone who says, you know, I can't do it all,.
I need help of other people, I have to rely on God,.
that in and of itself is a witness,.
that in and of itself is counter-cultural,.
that in and of itself is something.
that people take notice of..
And they say, this person's different.
'cause everybody else is trying to do it all on their own.
and trying to project an image.
that they've got it all together..
And so if we are willing to say, you know what,.
I don't have it all together,.
I actually need things in my life,.
I'm a dependent person, I'm a limited person,.
that's a witness to what God is doing..
Now earlier I said that God continues to form us.
throughout our lives..
In Jeremiah chapter 18, in the first four verses,.
there's this beautiful illustration of the potter.
that's forming the clay, right?.
And it talks about how the clay.
has to surrender to that process,.
embrace being formed by the potter..
The clay doesn't get to say to the potter,.
I think this is how you should mold me..
The potter is the one who is in control of that process..
And that process of forming and shaping and removing,.
that's a painful process, right?.
Think about what happens..
The clay gets formed, it gets shaped,.

$^{801}$parts of it are chopped off, removed,.
and then what happens, if you know anything about pottery,.
what happens next is it goes in the fire..
And it gets heated really hot..
And a lot of times it goes in the fire multiple times..
But then what comes out on the other side.
is something that's beautiful,.
something that's valuable, something that's useful,.
and something that if it's really well made,.
it reflects the craftsmanship of the maker..
In our house we have some really nice, fine dinnerware.
that my wife's grandmother got us when we got married..
We never use it, but it's really pretty..
It sits in this hutch next to our dining room table..
And every time I look at it, I'm always amazed.
because it's so beautiful..
The craftsmanship in it is really intricate and amazing..
And I'm always kind of in awe of whoever made it.
because that's, I know I could not do that..
And that's what can happen in our lives..
If we're willing to surrender to being dependent on God.
and our need for him and allowing him to form us.
and shape us and mold us,.
and even through taking us through difficult things.
like going in the fire, maybe even more than once,.
when we come out on the other side,.
there's something in us, in our life,.
that's valuable and useful and beautiful,.
and it's glorifying to God..
And the people around us who observe it,.
they say, "Wow, that person is really different..
"There's something really beautiful in their life.
"and something that testifies to someone else.".
That's another beautiful thing about dependence.
is when other people see it, they know it's not us..
They don't look at us and say,.
"All this stuff going on in your life.
"is because you're so great and you're so smart.
"and you're so tough and you hustle so hard.".
They have to look at our life and say,.

$^{841}$"There's something else going on here.
"'cause it's not you," right?.
And that's what we want..
We want to be in that place where we are reflecting.
God's work in our life, where we're dependent on him,.
where we're experiencing the peace and freedom.
that comes from that,.
and that's a witness to the people around us..
And that only happens as we learn to embrace.
our need for God and the need for the people around us,.
that God created us with..
We are limited creations..
That's okay, God made us that way on purpose,.
and we need him, and he wants us to turn to him.
in our lives and to rely on him..
Let's pray..
God, thank you so much for your truth.
and the way it speaks to us..
There are times when it's tough to hear..
We don't always want to hear that we are dependent.
or that we have need..
We'd rather just hear that we're great.
and that we have it all together..
But the reality is we do need you..
You formed us from dust, from nothing,.
and you breathed life and vitality and a soul into us..
But we did not make ourselves..
You formed us, and you are still making us..
And so God, please help us to embrace our need for you,.
to be honest with ourselves about our limitations.
and the fact that we can't do it all,.
to embrace the interdependence that you gave us.
in the church, the people around us that we need,.
to rely on you and others,.
to help us live into who you want us to be.
and who you've called us to be..
God, thank you that you are loving us.
and you are at work in our lives to form us and shape us.
into something that's way more beautiful.
than we could imagine or we ever would have planned for..

$^{881}$And I pray that you would continue that.
even this week in our lives..
Give us comfort, give us peace, give us hope..
Help us to rely on you..
\newpage



\section{約翰福音 14:1-6}
\label{sec:1rUpXzwI0NI}
\textbf{新天新地:人生終末的再思 1 將來的父家 (約翰福音14\_1-6) - 陳崇基牧師}
\newline
\newline
連結: \href{https://youtube.com/watch?v=1rUpXzwI0NI}{\texttt{https://youtube.com/watch?v=1rUpXzwI0NI}} ~~~~ 語音日期: 2023-08-14
\newline
\newline
\hyperref[sec:T08JyDUvISQ]{\small{< < < PREV SERMON < < <}}
~
\hyperref[sec:index_chronic]{\small{[返順時目]}}
~
\hyperref[sec:index_scriptual]{\small{[返順卷目]}}
~
\hyperref[sec:p_FHxvDNGe8]{\small{> > > NEXT SERMON > > >}}
\newline
\newline
約翰福音 14:1-6
\newline
\begin{longtable}{cl}
\hline
\hline
章節 & 經文 (和合本修訂版)\\
\hline
14:1 & \begin{tabularx}{0.7\textwidth}{X} 「你們心裡不要憂愁；你們信神，也當信我。 \end{tabularx} \\ \\ \relax
14:2 & \begin{tabularx}{0.7\textwidth}{X} 在我父的家裡有許多住處；若是沒有，我就早已告訴你們了。我去原是為你們預備地方去。 \end{tabularx} \\ \\ \relax
14:3 & \begin{tabularx}{0.7\textwidth}{X} 我若去為你們預備了地方，就必再來接你們到我那裡去，我在哪裡，叫你們也在哪裡。 \end{tabularx} \\ \\ \relax
14:4 & \begin{tabularx}{0.7\textwidth}{X} 我往哪裡去，你們知道那條路。」 \end{tabularx} \\ \\ \relax
14:5 & \begin{tabularx}{0.7\textwidth}{X} 多馬對他說：「主啊，我們不知道你去哪裡，怎麼能知道那條路呢？」 \end{tabularx} \\ \\ \relax
14:6 & \begin{tabularx}{0.7\textwidth}{X} 耶穌對他說：「我就是道路、真理、生命；若不藉著我，沒有人能到父那裡去。 \end{tabularx} \\ \\
[1ex]
\hline
\hline
\end{longtable}
$^{1}$現在將時間交給他.
多謝Cindy.
今日我們嘗試去看一些新的題目.
有一個題目我們很多時候會想到的.
天生地.
一些綜活的事情.
人生裡面最末的時候 最後的時候.
或者世界最後的時候會發生什麼事.
這些問題都是很切身的問題.
特別是我們生命裡面遇見不同的難處的時候.
無論是生命裡面一些的失去.
或者特別是親友過去離開世界.
剛才我們祈禱的時候特別紀念到兩位.
在我們中間的大英姐妹家庭.
我們有時候需要去想一想這些題目.
雖然很多時候我們會覺得.
這些問題不是我們經常會朝思暮想的.
但是有一些題目是需要去想的.
是需要去思考的.
聖經裡面有很多地方.
講到這些所謂的終末的教導.
人生的終末 世界的終末.
很多時候不是一些很抽象的理論.
就是說聖經裡面的作者.
保羅等等.
他們講到一些人生的終末的時候.
不是純粹是茶餘飯後的時候.
想一想不如思想一下這些東西.
然後寫一些這樣的東西.
很多時候這些終末的教導.
都是出自一些實際的處境.
譬如教會裡面有一些的試煉或者困難來到.
或者有一些教會裡面的成員.
他們的家人去世了 過世了.
到底他們應該怎樣去思考這個問題.
很多時候聖經就是在這些處境裡面.
寫出一些所謂終末的教導.
講到天家 講到新天新地.
New Heaven and New Earth.
到底一回是甚麼事.

$^{41}$如果你記得有些地方講到主在來.
譬如《貼薩雷加前書》第四章.
講到主在來的境況是怎樣.
絕對不是一些抽空去思考這些問題寫一下.
原因是因為開始的時候.
他們中間有一些的家人.
一些的親友離開了世界.
本來就去動筆去寫一些這樣的東西.
告訴他們我們終末裡面有甚麼盼望.
或者前書第十五章.
講到死人復活的事 基督復活的事.
很多時候我們復活節的時候就會看這些經文.
原因不是因為想抽象地.
來去教導這些題目.
而是他們中間有一些這樣的問題.
而他們中間有些問題是因為.
他們中間確是有人已經離開了世界.
那麼他們會去哪裡.
我們又會去哪裡.
所以這些都是一些實際處境裡面出現的終末的教導.
再講譬如啟示錄.
啟示錄講到新天新地最後那幾章聖經.
為甚麼會講這些東西.
因為初期的教會.
他們面對著很多迫害.
甚至有人是殉道離開這個世界.
在一個帝國的迫害.
在一個帝國的陰影下.
他們怎樣可以得勝呢.
這些都是一些很實際的處境.
等於我們今天坐在這裡.
我們都面對著一些很實際的處境.
我們的親人離開.
或者我們的親人重病.
我們開始要被迫去想這些東西.
所以亦都是講到這些終末的事情的時候.
千萬不要覺得這些是不會發生在我們身上.
雖然我們不會經常想.
不過是必須要去思考的.
這次我們第一次去看約翰福音第十四章.

$^{81}$耶穌怎樣講呢.
他說你心裡面不予憂愁.
你們信上帝也當信我.
在我父的家裡面有許多住處.
若是沒有我就早已告訴你們了.
我去完是為你們預備地方.
我若是為你們預備地方.
就必在來接你們到我那裡去.
我在那裡叫你們也在那裡.
我往那裡去.
你們知道哪條路.
多瑪就說主啊我不知道你去那裡.
怎麼能知道哪條路呢.
耶穌對他說.
我就是道路真理生命.
若不借著我沒有人能到乎那裡去.
很熟悉的經文.
很多時候在安息禮拜會用的.
這段經文其實是在一段很大段的部分.
就是約翰福音十三至十七章.
那段大的部分是所謂叫做farewell discourse.
就是耶穌臨走之前.
他講了一段話.
或者一段很長的一篇講道.
在哪裡講呢.
最後聖餐的所謂上樓房.
upper room的地方.
去講這段話.
你記得第十三章約翰福音.
就是時候已經到了.
耶穌在晚餐的場合.
他就拿起盤子.
然後脫了衣服.
然後去為門徒洗腳.
在最後的晚餐裡面.
他就開始去跟門徒講解天國的事.
講解將來會發生甚麼事.
這個叫做道別的講話.
一大段的在聖餐的樓房裡面去講.
講完之後他們就完結聖餐.

$^{121}$他們就發生了甚麼事呢.
十八章第一節.
完了十三至十七章.
十八章第一節.
耶穌說完這些話.
講完這些discourse.
這些sermon.
就跟門徒離開那個地方.
過了吉倫谷.
Kidron Valley.
然後過到耶穌上東面的一個原址.
就是赫西瑪利園.
跟門徒一起在那裡.
然後他就被出賣.
他就被帶到耶路撒冷.
被審判之後.
就傷了十字架.
死了他.
然後三天之後就復活.
復活之後.
也是升天回到庫那裡去.
所以這一段的說話是非常重要的.
Farewell Discourse.
其實在聖經裡面.
所謂的farewell.
臨別抽波.
臨別講一句說話.
臨別講一些東西.
在聖經裡面有很多的.
我想問你.
有誰在舊約裡面你記得的.
是臨別抽波講很多東西的.
有誰.
最出名的那個一定是摩西.
是嗎.
你記得.
申命記是嗎.
申命記就是齊集了.
其實申命記是甚麼呢.
是摩西的精選.

$^{161}$摩西港島的精選.
全部收集在一本書裡面.
告訴私人我要離開了.
之後那一章怎樣算.
之後你們會怎樣算.
約書一下就接手.
所以在一種這樣的transition過渡的底下.
去講一些很重要的說話.
Farewell Discourse.
耶穌在這裡就是講這件事.
我們想起這些臨別會.
鎮宴.
臨別講的說話.
通常都是一些很重要.
或者我們可能聽完之後.
我們一輩子都會記得的東西.
你想起一些人離開之前.
他們講的說話.
或者做過的事.
不是一定是去世的時候.
可能是一些場合.
生命裡面一些的離開.
譬如有很多人.
香港有很多人移民.
經常都說要寬送.
或者有時沒有寬送.
那種場合所講的一些farewell的說話.
甚至你可以想起離世之前.
都是講一些很重要.
印象很深刻的一些說話.
是預備所謂身後事.
或者預備將來會發生甚麼事.
在聖經裡面這些farewell discourse.
就更加重要.
耶穌這幾章所講的東西.
都是一樣.
在這裡講些甚麼呢.
第一點可以想起的.
就是主講清楚.
他預備他的父家有地方.

$^{201}$在頭這四節裡面講到.
心裡不憂愁.
憂愁的原因不是因為他們很清楚.
耶穌之後那一天.
在耶路撒冷他會受害.
然後上十字架.
那種慘痛的經歷.
因為他當時的門徒都不是那麼清楚.
他們憂愁的原因就是.
夫子即是他們的老師要離開.
為甚麼要離開.
到底他會去哪裡.
他跟了他那麼久.
到底這個時刻為甚麼要講一些臨別的說話.
可能他們有些confuse.
有些不是很清楚.
有些迷糊.
所以他就講這段說話.
他說他會prepare.
rooms.
in the father's house.
在耶穌的時代.
很少有一些獨立的房屋.
獨立的房屋.
譬如我們芝加哥有很多獨立的房屋.
single houses.
一間一間分開的.
在以前的世代.
耶穌的時代是不會的.
大部分的房屋都是一件一件的黏在一起.
然後形成一個很大的家庭.
大家一個compound.
一個complex.
因為當時按照文化的慣例.
每逢兒子要結婚.
到了適婚年齡的時候.
談婚輪嫁.
大概當時.
今時今日很多時候.
我們叫他遲一點結婚.

$^{241}$以前十多歲就來了.
十多歲就是成人了.
今天就不是了.
今天可能四五十歲都還沒有成人.
十多歲就要成人了.
要擔委投家.
到了適婚年齡的時候.
兩方面的父母.
就會為他安排.
安排這對青年的男女.
就會形成一個家庭.
談婚輪嫁.
兩邊的家庭談好之後.
就會發生什麼事呢.
兒子.
通常在另一條村.
跟那個女孩談好了.
跟那邊的家庭談好了.
禮金等等全部都搞定了.
他就會回到他的父家.
來做房子.
就在他爸爸的家.
或者他長大的家旁邊.
他就會建一些房子的牆.
或者一些結構.
就會為他這個新的將來的家庭.
來建造一個新的樓房.
是黏著那個父的家的.
大概一年之後.
因為談好之後就訂婚了.
你記得約瑟瑪利亞就是這樣.
然後一年之後.
他就會再回去.
然後就會大肆慶祝.
然後接那個新娘.
就回到自己的家.
回到他的父家.
所以那個新人.
就會住在那間新建造的房子.
黏著那個父家.

$^{281}$然後所有的兒子都會這樣做.
一個又一個.
然後一代又一代.
然後就變成一個很大的家庭.
甚至可以變成一個小的鄉村.
中間很多時候有一個天井.
這些很熟悉吧.
以前華人的文化都是這樣的.
有一個天井.
大家都是一起住的.
這種就是猶太人的文化.
如果你有這樣的文化.
你就會發現耶穌說的話.
就是有些味道.
他快要離開了.
耶穌要死在十架上.
死後復活然後升天.
回到他的父家.
所謂天家.
他就說我現在要走.
但你不用怕.
因為我走.
不代表我拋棄你.
我走.
其實是回去父家.
去到為你預備地方的.
之後會回來接你.
對嗎.
之後一定會回來接你.
這段話.
不只是道別這麼簡單.
不只是farewell.
是一個promise.
是一個承諾.
一個將來必定會發生的承諾.
你想起farewell的時候.
很多時候不會這樣說的.
對嗎.
有沒有人說移民離開了.
離開香港.

$^{321}$然後跟他的朋友說.
我要走了.
不過我會回來接你.
怎麼會呢.
回來跟你吃飯.
又會.
但不會說我一定會回來接你.
這裡耶穌在說的就是.
在那個文化裡面.
用一個這樣的意境.
你要說到主離開.
不是拋棄.
不是abandonment.
而是promise.
promise他會再回來.
在主裡面.
就會接你回去天家.
這裡在說的是一個家庭.
我為你預備一個家庭.
這個家庭不只是空間這麼簡單.
是一個住得到的地方.
也是一個持續的關係.
是一個relationship.
是一個family.
對嗎.
我們想起家.
我們當然會想起地上的家.
對嗎.
門徒我相信你也會想起.
到底我們屬於哪裡.
哪裡可以棲身.
哪裡可以給我立足.
地上的家.
耶穌說.
不會永恆的.
in that sense是不會永恆的.
因為地上的家一定會有風吹雨打.
就會毀滅.
人事變遷也會散.
對嗎.

$^{361}$生老病死也可以受到威脅.
最近大家看到夏威夷的大火.
對嗎.
那些家園.
一把火就全部都燒了.
連火也不能控制.
對嗎.
我昨晚預備的時候.
我看了一個社論.
那個社論是CNN那裡發佈的.
那個社論的title說什麼呢.
他說.
Hawaii fire shows that Hawaii paradise was a dream.
大家去夏威夷玩過.
Paradise.
世外桃源的地方.
百萬富翁千萬富翁.
走去弄一間屋子.
好像大富翁.
有時候去渡假.
夏威夷給人的感覺.
都是一種世外桃源的地方.
我們介紹朋友的時候.
朋友說你夏威夷來.
真的好.
夏威夷的地方是世外桃源.
不代表那個人是世外桃源.
但是那個風景很美.
世外桃源是一個幻象.
一個幻形.
一場火就告訴你.
在地上有多好的家園.
一把火就可以消滅.
一有生老病死.
人事變遷.
就一定會煙消雲散.
哪裡真的有一個家.
那裡的家.
耶穌說我為你預備一個家.
這個家是不會被毀滅的.

$^{401}$這個家是不會過去的.
好像你.
有沒有試過你出去旅遊.
或者出去玩.
你去了幾個星期.
你很恨的是什麼.
是啊.
為什麼呢.
你去夏威夷玩.
為什麼想回家.
為什麼不住在夏威夷.
因為始終不是你的家.
農場不及狗竇.
你回家.
你那張爛鬼床.
跟你在酒店睡的那張.
差多遠.
沒有人做房間服務的.
但你總是覺得那裡舒服一點.
你睡的爛位置始終都是你的.
這就是home的意思.
不是一個house.
是一個home.
耶穌說我為你保證.
我會回來接你.
這個就是天教.
第二方面他不僅給我們天教.
還有什麼呢.
回家.
回家的道路.
進入四章五節.
多瑪做多肯文.
他說耶穌我們不知道你去哪裡.
我們怎麼知道走哪條路.
是不是.
你又走.
那我們怎麼去找你呢.
要去找你.
就需要知道走哪條路才行.
耶穌對他說.

$^{441}$我就是道路真理生命.
若不直著我.
不會有人能夠輔拿你.
因為提到好幾次.
是道路.
the way.
你去哪裡.
我們怎麼去找你.
從哪條路.
耶穌說我就是那條路.
你知道以色列這個地方.
是世界的中心.
非洲,歐洲,亞洲.
都是連結在以色列這麼小的地方.
長長窄窄的.
就因為在世界的中心.
所以有很多路都會經過這個地方.
其實你住在加拿大.
在以色列地.
其實就是住在路邊.
你家大概不是在路邊.
路邊是什麼意思.
譬如美國的那些高速公路.
有很多很重要的高速公路.
中間intersection的地方.
你會不會看到一間樓在那裡.
然後車輛到處走.
以色列加拿大就是一個這樣的地方.
世界溝通,交通的中心.
所以經常都會說路.
猶太人的神學.
都受這件事影響.
猶太人怎麼去想像他們的生活.
生活就是走.
walking.
所以猶太人有一個教導.
一個傳統.
一個典籍.
叫做halakha.
halakha就是halak這個字.

$^{481}$halakha就是walking.
going and coming.
halakha是說什麼.
就是他生活裡面應該有什麼規則等等.
是走路的意思.
做人就等於走路.
生活就等於走路.
有一點像我們中文.
行為.
行事為人.
都是說走的.
在以色列的人更加會意識到這件事.
所以這裡說道路.
起碼有兩詞.
在整本《約翰福音》裡面.
提到道路.
the way.
the road.
the route.
只有兩詞.
這裡.
和第一章.
第一章開播說什麼.
就是施捨約翰.
施捨約翰就是.
先鋒.
就是預備道路.
給主來的.
所以施捨約翰說什麼.
我就算再擴野呼喊的聲音.
修直主的道.
為什麼.
因為主要來這裡.
救世主.
那位彌賽亞.
上帝因為愛他.
差遣彌賽亞這個救主來.
就要為他預備一個道路.
等於替他鋪一個紅的地毯.
這個貴賓終於.

$^{521}$主國終於要來到了.
施捨約翰就是做這件事.
他說我預備主的道路.
修直主的道路.
以至主能夠.
能夠走到來.
來到我這裡.
正如以賽亞先知所說的.
如果回去查以賽亞書.
又經常講道路.
特別是以賽亞書後面那一段.
四十章至六十章.
經常都提到道路.
道路 道路.
我們只看一節.
什麼道路.
那裡講到那些以色列人.
他們是被擄.
沒有了他們的家園.
流亡於海外.
在巴比倫的時候.
上帝應許.
他說我一定會回來找你.
一定會施恩給你.
一定會再愛你.
上帝的工作會再來到.
上帝的工作來到的時候.
會做些什麼呢.
五十一章.
耶和華的蚌蜱.
它的power.
會再一次.
好像上古的世代一樣興起.
好像出埃及的時候.
那個世代一樣的興起.
會使得海和深淵的水乾固.
然後海裡面的深處會開路.
使救贖的民能夠走過.
耶和華救贖的民必定會.
做什麼.

$^{561}$從無家的地方.
流亡的地方.
他們會歸回.
來到耶和華的硬石岸.
他們會再回來.
會再回到家裡.
流亡回歸的路.
其實是講上帝的救恩.
其實意識到我們每個人都是在流亡的人.
縱使我們地上有所謂的家.
但是這個家是不會永久的.
哪個家才會永久呢.
上帝所預備的.
耶穌就說我為你預備家.
不單止是這樣.
為你預備回家的道路.
You can find your way home.
你留意耶穌在這裡說的.
我是道路真理生命.
他說我是道路.
我不是只指一條路給你.
我自己本身就是道路.
我想這個所謂的中末思考.
其實是幫助我們.
告訴我們我們有一個.
最終的一個家我們要去.
一個目的地.
我們人生道路上.
如果我們知道我們的目的地.
Destination.
Make a huge difference.
當然聖經裡面說到中末的時候.
我們將來會去哪裡.
回什麼家.
就不是叫我們純粹在這裡死等.
等我們將來去到就好了.
或者叫我們望天打卦.
你要明白將來的目的地.
是決定你今天的選擇.
如果你將來不知道目的地去哪裡.

$^{601}$你今天的選擇很多時候就會迷失.
或者浪費.
我記得好幾年前.
我有一個Mini Cooper.
我很喜歡開那個Mini Cooper.
因為我覺得很好.
你有沒有開過Mini Cooper.
很好玩的Mini Cooper.
沒有.
你自己試一下.
我記得有一次.
買了一個Mini Cooper沒多久.
我們的車要服務.
我就拿了去他的工廠.
因為服務兩天.
那個車行很好.
借了一個Low-NOW給我.
借車給我用.
借一個新車給我.
我已經開的算是新車.
然後我上了一輛再新一點的車.
我一觸碰到那輛車.
真的很好.
你有沒有想過那種人車合一的感覺.
兩座位 兩門.
那輛是Turbo的.
他借了一輛這樣的車給我.
我馬上就飛回家.
我跟我老婆說.
這輛車很棒.
出去遊車河.
我抓了他上去遊車河.
很開心開這輛車.
我上了車之後.
我插上了車.
我就說我們走吧.
他就問我.
我去哪裡.
我說沒有目的地.
我就繼續開車.

$^{641}$上了Highway.
我開的時候.
我就意識到.
十分鐘之後.
五至十分鐘之後.
我都不知道去哪裡.
車油又貴.
然後我又開回家.
這些叫做浪費油.
你都不知道自己在上哪裡.
你是不是浪費油.
生命裡面我們的道路.
如果我們知道我們的目的地.
目的地是決定我們現在轉左或轉右.
或者我們有什麼的抉擇.
終末的思考.
想到人生的終末.
其實幫助我們.
不是說我們死等將會去那裡.
不只是這麼簡單.
而是決定我們今天所有的選擇.
等於你去遠足.
你去爬山.
如果你想一下.
你在山頂先望一望整個的路程.
用一個高的.
上面的一個角度來看下面.
你知道是很迂迴曲折的.
你只是在下面去走.
你只是留意前面那幾步.
你不會享受整個的旅程.
因為有時候跟一些年長的弟兄姊妹去交談的時候.
其實是很有益的.
因為他們很多時候會告訴你.
為什麼當時不早點順序.
為什麼當時不花多點時間去事奉.
追趕的是其他的東西.
為什麼不花多點時間去認識主.
那些不只是後悔.
其實是一些智慧的領悟.

$^{681}$去領悟到我們人生裡面的道路.
應該是決定於我們將來要去的地方.
去想將來的地方既然是這樣的話.
我們今天應該要有一個怎樣的生活.
是智慧的領悟.
我記得我做過很多喪禮.
我記得有一次最大的影響就是.
我記得是二零零幾年的時候.
我忘記是哪一年.
就是有一個喪家不是來教會的.
他們需要找一個牧師幫他們主持.
我幫他們主持.
喪禮那個時候要問的就是心平行坐.
我說不如你家屬寫一個心平行坐.
因為你認識他.
他寫了之後才交給我.
我通常會查一查.
我看一看那個心平行坐.
我最大的影響就是.
一個我要跟他主持喪禮的人.
跟我出生年齡是一樣的.
那時候是我年輕的時候.
我三十多歲的時候.
很大的影響.
我送一個跟我同年的人離開.
不要說送一個比我更年輕.
好像昨天我們的喪禮一樣.
送一個這樣的人離開.
我還記得那天喪禮.
是沒有人來的.
只有一兩個人.
因為是一個新移民.
是台山人來的.
廣東人.
來到這裡.
你知道這些移民生命很苦.
一味在這裡工作.
一天做到走.
然後做到自己身體已經衰敗.
突然間有惡性癌症出現.

$^{721}$很快他就離開了.
連他的女兒都沒有來.
很傷心的一個喪禮.
我還記得那天去送殯要上山.
還是大風雪.
在芝加哥大風雪.
只有我和那個審判員.
我們就去看著他下葬.
下葬的時候還是很棘手.
因為大風雪.
是不是.
棺木又放得不正.
然後我們要幫忙用2x4去弄.
很棘手.
但為什麼對我最大影響.
其中一個最大影響.
我會想一下到底我圓場的時候.
我想怎樣.
如果你想想一下.
你自己圓場的時候應該是怎樣的.
你自己應該怎樣想.
就是聖經裡面所說的終末的教導.
沒有終末的目的地.
沒有終末的從上而來的角度.
你就有很多事情會搞錯.
昨天我們和Leah做過安息禮拜.
當然她三年多.
她和絕症要交戰.
是很苦的.
但聽回她的故事.
她居然可以做這樣的祈禱.
祈禱上帝可不可以用我的癌症.
都可以見證到神.
說真的.
我是一個健康的人.
我都未必可以做到這樣的禱告.
但你可以在她身上見到.
將來的終末的盼望.
是實在的.
是拿不走的.

$^{761}$天家的盼望拿不走.
在絕處裡面仍然可以逢生.
那種盼望可以從哪裡來.
就是因為看到The End.
然後決定到今天應該怎樣去做決定.
這個PowerPoint的背景.
是星期一晚上拍的.
其中一位弟兄.
我們在Leah的家.
他已經很平安地離開了世界.
我們在家探訪的時候就在等.
等那些處理遺體的人來.
有兩三個小時的時間.
其中一位弟兄是幫忙做這件事的.
他在外面等車的時候.
等人來的時候.
他就拍了這張照片.
他發了給我們那個群組.
我們幫他負責那個安西大拜的群組.
天家到底是什麼樣子.
可能比這個更美.
必定比這個更美.
終末的思考.
不是一些抽空的想像.
是嗎?.
無論你今天是信主還是不信主.
是必須要想的.
因為你的目的地.
決定你今天所有的抉擇.
所有的方向.
我們練了一首歌.
我們希望這首歌能夠幫助大家思考.
我們自己的頌目.
只是音樂.
我們將歌詞打出來.
希望你聽這首歌的時候.
你看見這些歌詞.
思想一下.
我們自己的終末是怎樣.
你希望它的延長是怎樣.

$^{801}$你希望人生的終末.
最後的時候.
是怎樣的呢?.
天家對你有什麼意思.
對你今天每一個抉擇.
應該有什麼影響.
你希望它能夠延長.
你希望它能夠延長你的生命.
在我們生命當中.
在我們生命的道路裡.
你都可以敬拜,讚美,事奉你.
祈禱奉主名.
願上帝賜給各位平安.
現在散會.
\newpage



\section{2:1-13}
\label{sec:p_FHxvDNGe8}
\textbf{Show No Favoritism (James 2\_1-13) - Caleb Lee}
\newline
\newline
連結: \href{https://youtube.com/watch?v=p-FHxvDNGe8}{\texttt{https://youtube.com/watch?v=p-FHxvDNGe8}} ~~~~ 語音日期: 2023-08-21
\newline
\newline
\hyperref[sec:1rUpXzwI0NI]{\small{< < < PREV SERMON < < <}}
~
\hyperref[sec:index_chronic]{\small{[返順時目]}}
~
\hyperref[sec:index_scriptual]{\small{[返順卷目]}}
~
\hyperref[sec:S12GkifPujs]{\small{> > > NEXT SERMON > > >}}
\newline
\newline
$^{1}$Good morning, everyone..
It's a frightening privilege to be here with you guys today..
Before I begin, let's open in a word of prayer..
Would you bow your heads with me as we pray?.
Heavenly Father, you have been so, so good to us..
And sometimes, Lord, these words come out of us a little contrite..
But God, you are the God, the creator of heaven and earth..
From your word, there was light..
And from your word, there was life in us..
And so we come before you by the blood of your son, Jesus..
Come humbly before your cross and ask that you would hear us, that you would know us.
and see us, that you would give us hearts that are softened to you..
Ask that you would use this moment as we proceed to open our eyes to your word, to help us.
hear what it is that you have to say to us..
For myself, Lord, I ask that you would guide my words, that you would help me speak good,.
sound doctrine, that you would protect the ears of my brothers and sisters from anything.
I say that is not of you..
So we thank you, Lord..
We praise you in Jesus' name..
Amen..
So for those of you who weren't here the last time I spoke, we are in the book of James,.
which is a letter from James to the Jewish Christians of the dispersion, which are the.
Jewish Christians scattered away from their homeland..
And what we find in James are instructions, instructions given to these Jewish Christians.
on how they should be living..
They're given so that James' brothers and sisters would live wise lives, lives that.
are reflective of the gospel that they have received..
And what we find in James is that wisdom is not just knowing a lot about stuff, a lot.
about stuff that's about life, but rather wisdom must also include practical application.
of what we know about our God into every area of our lives..
This is a theme that will come up again and again..
James writes with this idea in mind, that followers of Christ can't just be hearers..
No, not just hearers..
They must also be doers of the word..
So if you have your Bibles with you, I ask that you would open with me to the book of.
James..
The passage is also going to be on the screen..
I'll flip through the slides for you..
Today we will be reading from James chapter 2, verses 1 to 13..
Again, that is James chapter 2, verses 1 to 13..

$^{41}$Let's read..
My brothers and sisters, believers in our glorious Lord Jesus Christ must not show favoritism..
Suppose a man comes into your meeting wearing a gold ring and fine clothes and a poor man.
in filthy old clothes also comes in..
If you show special attention to the man wearing fine clothes and say, "Here's a good seat.
for you," but say to the poor man, "You stand there or sit on the floor by my feet," have.
you not discriminated amongst yourselves and become judges with evil thoughts?.
Listen, my dear brothers and sisters, has not God chosen those who are poor in the eyes.
of the world to be rich in faith and inherit the kingdom He promised to those who love.
Him?.
But you have dishonored the poor..
Is it not the rich who are exploiting you?.
Are they not the ones who are dragging you into court?.
Are they not the ones who are blaspheming the noble name of Him to whom you belong?.
If you really keep the royal law found in Scripture, love your neighbor as yourself,.
you are doing right..
But if you show favoritism, you sin and are convicted by the law as lawbreakers..
For whoever keeps the whole law and yet stumbles at just one point is guilty of breaking all.
of it..
For he who said, "You shall not commit adultery," also said, "You shall not murder.".
If you do not commit adultery but do commit murder, you have become a lawbreaker..
Speak and act as those who are going to be judged by the law that gives freedom because.
judgment without mercy will be shown to anyone who has not been merciful..
Mercy triumphs over judgment..
Now the amazing thing about James is that his instructions are very, very straightforward..
Now look, my brothers and sisters, believers in our glorious Lord Jesus Christ must not.
show favoritism..
To his brothers and sisters, James says, "Since you believe in Jesus, you must not show favoritism.".
It's right there..
If you are a believer in Jesus Christ, you have no business showing favoritism..
So what is favoritism?.
Does anyone know?.
Can anyone give me a definition?.
Okay..
Well, I'll give you one..
Favoritism is this practice of special treatment, okay?.
Special treatment of valuing things or people differently..
And what often comes along with this word is that favoritism is unfair..
I think most of us here have experienced favoritism or even shown favoritism at some point in.
their life..

$^{81}$You know, sometimes it's the small things..
Maybe it looks like going...it might look like certain lunch suggestions being more.
acceptable and then you never get to go where you want to go..
It might look like the favorite food of your brothers and sisters constantly being eaten.
night after night because your parents love them very, very much and so you're stuck eating.
that food..
Other times, it might look like conversations that are cut short or even avoided so that.
you can spend more time talking to your closest friends..
It might look like excess friendliness to the professor who grades your papers..
It might look like friendliness towards the person who writes your paychecks..
But let's let James show us what favoritism looks like..
I realize now that I did not keep going with the slides so you can read along with me..
But that's okay..
We're going to look at James chapter two..
We're going to look at what favoritism looks like..
We're given this scenario in verses two to four, right?.
Look with me there..
Two to four where we find two men joining believers in their place of meeting..
For if a man wearing a gold ring and fine clothing comes into your assembly and a poor.
man in shabby clothing also comes in, and if you pay attention to the one who wears.
the fine clothing and say, "You sit here in a good place," while you say to the poor man,.
"You stand over there or sit down at my feet," have you not then made distinctions among.
yourselves and become judges of evil thoughts?.
There's one man..
He's dressed in fine clothes..
His fine clothes are bright and shining, right?.
That's the description that's given about him..
He wears this gold ring that lets everybody know..
It announces how high his socioeconomic status is..
The other is described as this poor man, a man in filthy old clothes..
In contrast to the man in fine clothing, this man's unwashed, right?.
And his clothes are mismatched..
It is easy to tell what rung of society he is on..
Favoritism in this case looks like offering a good seat to the man wearing fine clothing.
while telling the poor man to stand there or sit where people are resting their feet..
It's this honoring of the well-off man for his appearance, for his wealth, and this dishonoring.
of the poor man based on his lack..
Doing this, James writes, is not only favoritism that believers must not show, it's also discrimination.
against yourself..
And it makes you into a judge with evil thoughts, into a judge that judges poorly..

$^{121}$Remember that the people James writes to are a scattered people..
These are people who live in a foreign land..
They are suffering under discrimination..
They are greatly in need..
And so the judgment that is being made towards the poor man and the treatment he's given.
in this what-if scenario is the very same treatment and judgment that they are receiving..
This is important to remember because we're so quick to think, of course, favoritism is.
wrong..
Anyone can tell you that..
You should not discriminate..
And you know what?.
It's not that difficult to avoid these things..
We think about that so easily and so quickly..
But think about these..
Think about how helpful it would be to a poor group of people if someone with wealth and.
resources were to join them..
Think about the extra burden that will be placed on this group of people just from having.
one more poor person join them, one more mouth to feed when you yourself have nothing and.
are struggling to feed yourself..
You see, favoritism is tempting..
Favoritism is very tempting..
And James writes this to his brothers and sisters because this is a real and present.
problem..
He talks about this because it's a real flaw in our character in Christians that keep us.
from being whole and complete in our entire being like Jesus..
So then, how does James tackle this problem?.
If you look with me to verses 5 to 7, I'll actually turn the slide this time..
It says, "Listen, my dear brothers and sisters..
Has not God chosen those who are poor in the eyes of the world to be rich in faith and.
to inherit the kingdom He promised to those who love Him?.
But you have dishonored the poor..
Is it not the rich who are exploiting you?.
Are they not the ones who are dragging you into court?.
Are they not the ones who are blaspheming the noble name of Him to whom you belong?".
He starts to listen, right?.
"Listen," James says, "your attention right here, your attention on this," which is a.
rhetorical question, rhetorical meaning that this is a statement of truth that is disguised.
as a question, all right?.
"Has not God chosen those who are poor in the eyes of the world to be rich in faith,.
to inherit the kingdom?".

$^{161}$And reading this, the response of James' brothers and sisters is, "Yes, of course, God has chosen.
those who are poor in the eyes of the world.".
They can't say otherwise because, you see, these are Jewish Christians..
The history of their people, of those who are tied to the line of Abraham is just filled.
with God's concern for the lowly and the needy..
The literature of Israel is filled with this theme of God's love for the poor, for the.
widowed..
God's blessings are even promised to those who care for the poor, who see the poor and.
reach out to them..
Furthermore, what's even more than that is that the brothers and sisters in James, the.
audience to whom he writes, they are living proof of God's doing..
God has chosen them, even though they were poor in the eyes of the world, to be given.
faith and to be heirs of God's kingdom..
Instead of treating them poorly based on external appearance, God's attitude towards the lowly.
and the needy is that of compassion, of mercy, and of love..
So yes, it's true..
God cares for the poor in the eyes of the world..
So then, to dishonor the poor is to go against God's attitude, to go against the will of.
God, to represent the opposite of God's will..
What we find here is that believers in Jesus must not show favoritism because God does.
not show favoritism..
Believers must not show favoritism because God does not show favoritism..
God does not make judgment based on outward appearance..
He looks at the heart..
And so, James reminds his brothers and sisters of God's attitude towards the lowly and the.
needy..
He reminds them of God's attitude towards them..
And by keeping this in mind and remembering God's compassion, believers can overcome this.
temptation to show favoritism..
But James does not just leave it there..
He doesn't just leave that be..
The issue with favoritism isn't just that it fights against God's compassion..
You see, favoritism is a sin..
This is the second point that we find in James, that believers in Jesus must not show favoritism.
because favoritism is sin..
Let's look at verses 8 to 11..
If you really keep the royal law found in Scripture, love your neighbor as yourself,.
you are doing right..
But if you show favoritism, you sin and are convicted by the law as lawbreakers..
For whoever keeps the whole law and yet stumbles at just one point is guilty of breaking all.

$^{201}$of it..
All right?.
It's pretty straightforward..
Believers that keep the royal law, right, they are doing right..
They're doing well..
But here's another question..
What is this law?.
What is this law that he speaks of?.
This is the sovereign and royal law of God..
It points to the law of Moses that is declared again by Jesus and fulfilled in Jesus..
It points towards the commandments found in Leviticus chapter 19 of how God's people are.
to be holy, how they are to be set apart for him, and it points towards what Jesus declares.
to be the second most important of the commandments, to love your neighbor as yourself..
In Leviticus 19, there is command after command given on how to treat those around you..
Anyone and everyone you come into contact with in God's eyes is your neighbor..
You see, favoritism is sin..
Favoritism is transgression because both the poor man and the rich man are your neighbor..
It's not enough to just keep the commandments of God in some ways or even most ways..
Under the law, it is not good enough to almost be good enough..
To stumble at just one point, at just one point, is to be guilty of breaking all of.
the law..
The law requires absolute perfection, and we're not perfect..
Nobody here can say that they are a perfect keeper of the law..
If you understand this, then you are on the same page as James..
Nobody can get to heaven through their perfect law keeping..
The statement of absolute requirement of the law is placed here to remind James' brothers.
and sisters of who they are because they're not people who have perfectly kept the law..
They are people who have had the law perfectly kept for them..
They're not people who are crushed under the judgment of the law that condemns..
No..
James' brothers and sisters are those that are under the law that liberates..
What does James want his brothers and sisters to do?.
Well, we see in verses 12 to 13, 12 to 13, what not showing favoritism looks like..
This is what not showing favoritism looks like..
It looks like speaking and acting, not as someone who is under the law that condemns,.
but as someone who is under the law that gives freedom, as someone who is under grace..
It's the same reason that was given before in verse 5, because Jesus Christ did not show.
favoritism..
He did not judge you based on your outward appearance..
Instead, he saw you and he saw me in our helpless state..

$^{241}$And in that moment where we expected judgment, where we deserved judgment, we were instead.
met with a deep and overwhelming compassion..
Brothers and sisters, favoritism is such a poison to God's people..
It damages relationships and it holds us back from spiritual growth without us even noticing..
Because of the sin of favoritism, too many of God's people have found their church to.
be one of the loneliest, most disconnecting places they have ever been..
To be surrounded by so many people while being unseen and unheard is just heartbreaking..
And this cannot be allowed..
We as believers of Jesus Christ cannot allow those who God has placed in our surroundings.
as our neighbors to remain separated and disconnected from the body of Christ..
Believers in Jesus Christ, you are under grace..
You were shown mercy..
Your sin was washed away once and for all..
And not only that, your identity is now more than a worldly measurement of ability or accomplishment..
You were chosen, you were elevated, and now exist as an heir to the kingdom of God..
And having been freed by this law that liberates, live as a child of God, as an heir to the.
kingdom of God..
Let your entire being be affected by the mercy of Christ and do not judge your neighbor by.
their outward appearance..
Let's reflect on James' instruction..
Do not show favoritism because Christ did not show favoritism..
Do not show favoritism because favoritism is sin..
Instead, show mercy and show compassion..
See your neighbor the way Jesus sees them, the way that Jesus saw you, and approach your.
neighbor just as Jesus approached you in his compassion..
Let me close us in a word of prayer..
Would you bow your heads with me?.
Heavenly Father, again we are here before you because of your mercy and your grace..
Because you reached out to us and you pulled us close to you..
You decided for whatever reason in your mind that you wanted us..
And so to an unlovable people, you loved us and you made us lovable..
And so I pray in this moment, Lord, for my brothers and sisters here, that you would.
lay your hand on them and that you would remove the calluses that they have formed..
You would remove the calluses from their hearts and that you would renew their hearts and.
their minds to give them new eyes and new ears, to renew them, Lord, to display before.
them the glory of your cross and to overwhelm them with your love so that they might in.
turn be representatives of who you are, that they might live as people who are under grace..
And so help them find peace in you, Lord, help them find rest in you..
And help them live every single part of their life affected by what you have done for them,.
Lord..

$^{281}$We pray this in Jesus' name, amen..
[BLANK AUDIO].
\newpage



\section{啟示錄 21:1-4}
\label{sec:S12GkifPujs}
\textbf{新天新地:人生終末的再思 2 將來新天地 (啟示錄21\_1-4) - 陳崇基牧師}
\newline
\newline
連結: \href{https://youtube.com/watch?v=S12GkifPujs}{\texttt{https://youtube.com/watch?v=S12GkifPujs}} ~~~~ 語音日期: 2023-08-21
\newline
\newline
\hyperref[sec:p_FHxvDNGe8]{\small{< < < PREV SERMON < < <}}
~
\hyperref[sec:index_chronic]{\small{[返順時目]}}
~
\hyperref[sec:index_scriptual]{\small{[返順卷目]}}
~
\hyperref[sec:mdfK7VGG0u4]{\small{> > > NEXT SERMON > > >}}
\newline
\newline
啟示錄 21:1-4
\newline
\begin{longtable}{cl}
\hline
\hline
章節 & 經文 (和合本修訂版)\\
\hline
21:1 & \begin{tabularx}{0.7\textwidth}{X} 我又看見一個新天新地，因為先前的天和先前的地已經過去了，海也不再有了。 \end{tabularx} \\ \\ \relax
21:2 & \begin{tabularx}{0.7\textwidth}{X} 我又看見聖城，新耶路撒冷由神那裡，從天而降，預備好了，就如新娘打扮整齊，等候丈夫。 \end{tabularx} \\ \\ \relax
21:3 & \begin{tabularx}{0.7\textwidth}{X} 我聽見有大聲音從寶座出來，說：「看哪，神的帳幕在人間！他要和他們同住，他們要作他的子民。神要親自與他們同在。 \end{tabularx} \\ \\ \relax
21:4 & \begin{tabularx}{0.7\textwidth}{X} 神要擦去他們一切的眼淚；不再有死亡，也不再有悲哀、哭號、痛苦，因為先前的事都過去了。」 \end{tabularx} \\ \\
[1ex]
\hline
\hline
\end{longtable}
$^{1}$早晨 各位姐妹.
今日繼續睇新天新地.
今日睇第二講.
我地講一D 人生中末裡面.
我地需要思考ge 問題.
我上一次講過呢話呢D ge 題目呢.
唔係一D 好抽象.
或者對我地呢完全冇關係ge .
非常切身ge 一D 問題.
特別係環ge 世界.
我諗環世界裡面呢有一樣ye .
其中一樣ge 教育呢係非常之缺乏.
咩教育呀.
生死教育.
生死教育有冇聽過.
去到思考生與死ge 問題.
特別係人生中末ge 事情.
其實呢個ge .
呢一代啦.
已經開始去意識到.
即係需要做一D 咁樣ge 教育.
生死教育.
咁點解我地咁疏忽呢樣ye 呢.
原因好簡單因為我地好唔想提.
因為我地ge 常規裡面.
成日去話比我地聽.
我地係唔需要面對死亡.
可能我地需要面對死亡.
但係人地需要面對死亡.
就唔係我要面對死亡.
生死教育就係想做呢樣ye .
聖經都有生死教育.
但係要知道點解我地唔想提.
其實唔係淨係華人ge 文化.
係唔鍾意提.
其實好多文化都唔想提.
美國有一個ge 醫生叫做Paul Blain.
如果你想知道痛苦.
Pain.
呢一類ge 事情.

$^{41}$同我地人生有咩關係.
信仰點樣去幫助我地.
你一定要睇Paul Blain.
呢一個醫生所寫.
點解呢因為Paul Blain.
呢個美國ge 醫生.
佢頭一半ge 職業.
行醫ge 生涯.
佢頭一半係響印度.
一D 好荒蕪ge 地方.
去到做醫生.
Paul Blain花jor 佢自己上半生.
就係響印度.
去到服侍一D 窮鄉僻壤ge 地方.
醫學醫療設備係比較簡陋ge 地方.
係呢D 人群當中.
佢就去做一個醫生.
佢話gwo 度ge 人.
因為佢接觸多ge 時候.
佢話佢地呢D 咁ge 人.
對於生死離別生老病死等等.
係非常之熟悉.
因為每一日呢佢地都要接觸.
但係Paul Blain下一半ge 生涯.
佢就返jor 黎美國.
返黎美國繼續行醫ge 時候.
佢就發現美國.
或者現代ge 世界.
或者一D urbanizege 地方.
就完全係一個世界.
Paul Blain其中一寫到.
話美國係一個點ge 地方.
佢話係美國我碰到一個社會.
呢個社會係不惜代價去避免痛苦.
病人ge 生活.
佢話係美國病人ge 生活.
佢話比任何世界一個地方.
病人ge 生活更加舒適.
但係呢度ge 病人.
似乎係應對痛苦ge 時候.

$^{81}$係不堪一擊.
所受ge 創傷係比其他人更加大.
好有意思呢句說話.
我地有最舒適ge 醫療.
最好ge 醫療.
全世界你搵醫療.
應該美國都數一數二.
但係我地卻係不堪一擊.
其他D 窮鄉僻壤ge 地方.
冇咩醫療設備.
佢地卻係對死亡.
對於傷痛.
對於創傷.
有另外一種ge 體會.
其實呢個唔係新ge 發現.
幾百年前有個神學家叫嘉爾文.
十六世紀ge 時候.
佢曾經講過一句咁ge 說話.
佢話我地每一日經營ge 一切.
就好似係地上面為自己創造永生一樣.
長生不老一樣.
如果我地見到一具ge 遺體.
屍體.
我地或許會對生命.
轉瞬之間一個ge moment.
我地會覺得.
我地覺得生命裡面.
我地做少少思考.
人生真係好短暫.
但只要轉身我地ge 視線一離開.
我地就有一種可以繼續活下去.
係唔需要面對死亡.
一種中斷ge 一種咁ge 思想.
固定jor 係我地腦海裡面.
當你唔再睇見一個dead bodyge 時候.
你好快你就返返去.
原來繼續去活下去.
係一個常規.
係永遠可以繼續活下去.
聖經點樣講生死.

$^{121}$生與死ge 教育.
呢個就係人生綜合裡面我地需要思考.
人生裡面.
係有一個目的地.
有一個destination.
我地對呢個destination越多了解ge 時候.
我地先至會明呢我地而家呢個journey.
現代人有一個世界觀.
雖然我講D 世界觀好似好抽象ge ye .
等陣間我會講少少哲學ge ye .
但係我知大家唔鍾意哲學.
That's ok 講少少je .
小小個.
咩叫世界觀.
世界觀就係一個world view.
即係我地點樣去睇世界.
睇事物睇人生最終極ge 問題.
呢D 都係世界觀.
世界觀就係其實一早ge 信仰.
一早以信心去到作出ge 宣告.
唔係一D 科學呀或者數學呀.
或者任何ge 理據去到得到ge 結論.
相反係一D ge 前設.
一D 我地ge 預設.
一D presupposition.
咩叫世界觀.
就係一早ge 一D 咁ge 觀念.
係我地先有預設ge .
以至我地用呢D 咁ge 預設ge 思維.
黎到去解釋世界其他ye .
ok呀.
現代人有個咩世界觀.
自然主義.
咩事呀咩自然主義.
naturalism.
意思就係話我地人生裡面所有ge ye .
就一定係要用呢.
現世物質自然ge 方法黎到去解釋ge .
其他的解釋呢我地係唔會接受.
即係有D 超自然ge 解釋.

$^{161}$有D 超自然ge 原因我地係唔會去考慮.
呢個就係我地ge 前設.
好多時現代人ge 文化就係咁樣.
但係聖經ge 文化.
聖經ge 世界觀卻係唔同ge .
聖經ge 世界觀係點樣ge 世界觀.
聖經話比我地聽.
呢一個ge 世界係上帝.
呢個自然以外.
超自然ge 一位.
黎到去創造呢個世界.
有始亦都有終.
所以我地人生亦都係有始亦都有終.
呢個ge 世界觀係點黎ge .
就從啟示黎ge .
Revelation.
聖經裡面講譬如啟示錄.
啟示係咩意思.
啟示呢個字其實就係.
將個蓋打開ge 意思.
Uncover.
即係本身你唔知道裡面係咩打開.
你煲湯煲緊咩呀.
你個蓋cum 住.
唔知道.
你將佢打開其實就係一個啟示.
上帝係聖經裡面啟示ge .
就係另一個ge 世界觀.
唔係一個純粹自然.
可以解釋一切ge 觀念.
而係有上帝以外.
呢個世界以外ge 解釋.
今日我地睇呢段經文.
《創世記》第21章.
我先讀.
我又看見新天新地.
因為先前ge 天和先前ge 地已經過去了.
海也不再有了.
我又看見聖神新耶路撒冷.
由上帝那裡從天而降.

$^{201}$預備好了就如新娘打扮整齊.
等候丈夫.
我聽見有大聲音從寶座出來說.
看呀!上帝的帳幕在人間.
他要與他們同在.
他們要作他的子民.
神要親自與他們同在.
上帝要拆去他們一切ge 眼淚.
不再有死亡.
也不再有悲哀 哭號 痛苦.
因為先前的事都過去了.
那在寶座上的.
佢話.
看呀!我要擺一切都更新了.
太后說.
你要寫下來.
因為這些話是可信靠.
係真實.
太后對我說.
成了我是阿拉伯.
我是俄梅雅.
我是起初.
我是終結.
我把生命的泉水.
白白的切割那口渴的人海.
得勝的.
要承受這些威脅.
我要作他們的上帝.
他要作我的兒子或者兒女.
這裡是聖經最後的兩章.
二十一章和二十二章.
最後的兩章說到將來的新天新地.
我們去的目的地.
人類歷史.
人類的世界.
將會去到的目的地.
就是將來會發生的新天新地.
怎樣描述呢?.
我們不可以說很多細節.
不過你大概可以get到.

$^{241}$有新天.
有新地.
新的天地.
有聖城.
耶路撒冷.
另外一個的城市.
有生命河.
有生命樹.
等等.
繼續看下去二十二章.
周座不再有.
神住在他的子民中間等等.
這些元素.
讓你想起聖經裡面其他什麼地方.
讓你想起其他什麼地方.
新天新地.
有天地.
對嗎?.
有生命河.
生命樹.
不再有周座.
上帝住在人中間.
就是創世記.
創世記頭三章.
就是說到上帝創造原本的世界.
原本的宇宙.
聖經好像一個bookend.
就是書架.
書架頂著兩邊那些.
那些bookend.
一個在開頭.
另一個在最後.
是首尾呼應的.
這個就是聖經所啟示的.
截然不同的世界觀.
就是怎樣看人類的歷史.
有始亦都是有終.
就是來幫助我們去詮釋.
去解釋我們生命裡面所面對的一切.
最近弟兄姊妹都很忙.

$^{281}$因為我們幫一些弟兄姊妹.
他們剛剛親人離開.
是很艱難的時候.
見到一些我們平時常規裡面見不到的東西.
是不容易的.
先不說有些家庭他們的至親離開.
這個非常艱難.
我們周邊的人去做關心.
做一些照顧的工作.
其實都是非常之勞累.
到底怎樣去解釋這些.
聖經的啟示.
就是說這個世界裡面.
到底有什麼問題.
而上帝要做什麼.
這一段聖經.
我想說兩點.
第一點就是他講到新天新地.
必定會來到.
新天新地.
必定會來到.
我們是向著這個目的地進發.
當然中間裡面可能有不同的事發生.
人死了之後.
在基督裡面.
我們的軀殼暫時放下.
我們叫塵歸塵 土歸土.
我們的軀殼會暫時放下.
但是人的靈魂就會去哪裡.
果然是你.
雖然沒有了軀殼.
是你來的.
靈魂.
就會回到天下.
我們所謂天堂.
Heaven.
但是不是永久.
不是永久.
我們的靈魂去到天下的時候.
絕對不是永遠就是這樣.

$^{321}$我們印象中或者想像中.
就飄來飄去.
在舊魂上面等等.
在天下不是故事的終結.
聖經講什麼.
聖經說還有一個目的地在去.
還有一個大結局在去.
那個結局就是新天新地的來到.
就是耶穌基督再來的時候.
將萬物來到去更新.
在基督裡面的信徒.
就會身體復活.
我們所謂的放下的軀殼.
會扭歪.
但是到耶穌基督來的時候.
新天新地到的時候.
我們就會得到一個復活榮耀的新軀.
這個新軀是怎樣的.
有沒有咒文.
有沒有一些我不喜歡的東西.
你看看耶穌基督的故事.
復活的故事.
你就略知一二.
耶穌基督的身體.
他從墳墓復活過來的時候.
他的軀體是能夠被認到的.
門徒可以認到耶穌.
雖然他很詫異.
可能不知道是不是他.
但最後他認到.
然後那個軀殼是一個物質的軀殼.
好像我們現在這個軀殼一樣.
是抓到和摸到的.
絕對不是飄飄飄不屬物質的.
但是他是一個榮耀的新軀.
一個榮耀的軀殼.
這個就是基督裡面給我們的認許.
所以要提到什麼呢.
他說有看見不是新天.
這麼簡單.

$^{361}$不是純粹新天.
而是新地球.
整個聖經的故事.
整個聖經的世界觀.
就說上帝首先創造天與地.
而且縱使人類的罪惡破壞世間.
或者上帝所創造的東西.
最終有一天這個天與地也可以走在一起.
上帝會更新所有的一切.
所以為什麼叫做新天和新地.
新不是舊這個或者現在這個.
是沒用的.
而是到新天新地來到的時候.
就是一個更新.
一個renew.
是什麼意思呢.
打個比方.
我的電腦就是用一個聖經軟件.
好幾年前.
我用的那個版本是第六版.
第六 number six.
因為需要做多點東西.
我就給點錢.
然後就upgrade它.
就變成第七版.
第七版和第六版有什麼分別.
一個舊一個新.
是嗎.
六的有很多bug.
有很多程式上的問題.
然後更新了第七版.
功能就多了.
或者它的運作就沒有毛病.
新天新地和我們現在這個所謂先前的天地.
就是一個這樣的關係.
當上帝來到的時候.
耶穌再來的時候.
就是一個新天新地.
因為先前的天和先前的地已經過去了.
先前的天地就是我們現在這個天地.

$^{401}$現在這個軀殼.
現在我們所見到的一切.
這些都是美善的.
因為上帝創造萬有的時候.
他看見是好的.
It is good.
And it is very good.
縱使有罪惡破壞這個世界.
有邪惡來扭曲這個世界.
但依然它的本相是美好的.
但當耶穌再來的時候.
現在這個所謂先前的天地.
就會過去了.
這個版本是夠了.
上帝會將它更新.
使它完美.
這個不是抽空的教導.
是和我們非常切身.
因為當我們明白目的地.
我們就明白今天.
我們應該有什麼方向.
第一點.
新天新地必定會到.
第二.
就是耶穌基督一定會得勝.
耶穌基督一定會得勝.
第一節最後一句.
海也沒有了.
這麼奇怪.
新天新地不會有海.
那就是不能住海景.
什麼意思.
現代人我們想起海.
就覺得很美.
很喜歡海景.
很喜歡看日落游船河.
但古代的人.
聖經時代的人.
海是什麼呢.
海是邪惡的根源.

$^{441}$因為對古代的人來說.
他們住在陸地.
海是很難參透的.
裡面有什麼呢.
是混亂.
是人類不能住的地方.
特別是啟示錄當中.
海不單止代表一些我們不能夠參透.
和混沌的東西.
在啟示錄裡面.
海是代表什麼呢.
我們選一節經文.
因為提到海有很多節經文.
十二章第十八節.
這裡有一個異象.
有一個vision出現了.
十二章最後那一節.
哪時龍站在海邊的沙灘上.
我又看見一隻獸從海裡面上來.
有十隻角七個頭.
十角上面帶著十個冠冕.
七個頭上面有蝕族的名號.
十二章最後就出現大龍.
十二章發生什麼事.
其實就是一個爭戰.
一個的battle.
在天上面有的battle.
就是正邪之間的一個battle.
然後這隻大龍就被擊敗.
然後它就跌落地上.
跌落地上它沒有放棄.
它繼續追殺上帝的子民.
追殺當時的教會.
去到迫害他們.
但上帝又保守住它自己的子民.
然後這隻龍就站在海邊.
它就叫它的助手出來.
有什麼助手.
你看看.
助手就是有一隻獸.

$^{481}$從海出來.
它有十個冠冕七個頭.
意思就是說它很有power.
我們不用打開報紙.
不用看報章.
我們去猜猜現在什麼政權.
代表誰.
不需要.
這個不是啟示書寫的目的.
也不是這種寫法的目的.
當時啟示錄發生什麼事.
寫給當時的讀者.
是因為當時第一世紀的教會.
第一世紀末的時候.
他們在經歷很艱難的時候.
水深火熱的時候.
是羅馬帝國在迫害他們.
打壓他們.
他們不知道在哪裡可以有家.
哪裡可以立足.
啟示錄就是寫給他們.
要去鼓勵他們.
要告訴他們天上到底真實發生什麼事情.
以至他們能夠做一個得勝的基督徒.
有艱難來到的時候.
他們說不需要怕.
因為上帝已經贏了.
今天這本啟示錄.
不是預先寫好的報章.
而是跟我們很切身.
當我們遇見很多不同的壓力.
生命裡面我們有很多的艱難.
或者被人追趕.
這些都是啟示錄裡面所代表的.
那個獸或者龍所代表的.
就是那些帝國的壓力.
當那些壓力追趕著你的時候.
你怎麼辦.
啟示錄就是說.
上帝已經得勝.

$^{521}$他會保守.
去帶領著他的子民.
今天這些壓力都不需要多說.
你個人生命裡面.
可能遇見很多不同的迫害.
令到你喘不過氣.
上個星期我看了一些新聞.
說到加拿大.
在Northwest Territory.
Yellowknife這個地方.
Yellowknife是.
你想看北極光就去Yellowknife.
去一些很北的地方.
去看北極光很美.
Wildfire.
多少萬人疏散.
2萬多人疏散.
這個世界發生了什麼事.
有這麼多邪惡的事.
聖經的世界觀就是說.
是這個世界被邪惡所破壞.
不過上帝最終會得勝.
很切身的.
你跟我今天可能都遇見很多掙扎.
或者甚至見到一些很不幸的事.
Cancer.
腫瘤.
有不同的壓迫.
不同的苦難.
或者有一些未知之數.
你搞不定的.
發生在我們自己身上.
或者我們所摯愛的人身上.
聖經的世界觀.
要想說什麼.
不是純粹抽空的哲學.
聖經的世界觀告訴我們.
有邪惡不過上帝.
最終會彰顯他公義.
海也沒有了.

$^{561}$什麼意思.
就是說新天新地當中.
不會再有這些邪惡的事發生.
因為上帝會住在中間.
上帝會滅絕所有的罪惡.
這個就是我們每一個人.
希望能夠去到的地方.
每一個人憧憬.
希望能夠住在的地方.
聖經就說.
我給你一個promise.
上帝給我們一個應許.
這個新天新地上帝.
會滅絕所有的邪惡.
歡迎你.
在基督裡面能夠來到.
剛才我們看這些經文.
我們看的時候覺得很奇怪.
什麼出現怪獸等等.
我不知道誰喜歡看戲就看.
喜歡看Sci-Fi.
Sci-Fi其實是一些很深的電影.
很深的道理.
其實啟示錄或者啟示錄這種寫法.
我們叫做天啟文學.
其實是一種好像Sci-Fi.
好像科幻這一類型的寫法.
原因是什麼.
因為這種寫法.
很多時候就是被一些.
被很多惡勢力去打壓的一班群體.
他們會想像這些惡勢力.
就是一些很凶悍的怪獸.
羅馬帝國外表可能是非常堂皇.
很華麗.
但內裡其實是什麼.
是暴力是暴權.
天啟文學或者啟示錄這一類的東西.
其實就是被我們盼望.
就是叫我們的生命.

$^{601}$你看到這些邪惡勢力的時候.
但是你不會被這些生命.
不會被這些勢力所騎劫.
我們不會恐慌.
因為上帝是話事.
剛才說到Sci-Fi.
有誰記得一些童話故事.
Fairy Tales.
記不記得.
這些是華人文化裡面很少的.
為什麼要提這些東西.
我們有很多小朋友.
家長會跟他們看一些Fairy Tales.
有沒有.
為什麼要跟他們說這些東西.
這麼恐怖.
但你知道Fairy Tales想做什麼.
這種寫法.
Fairy Tales裡面有怪獸出現.
有巫婆出現.
但最終會發生什麼事.
就是那些怪獸和巫婆等等的惡勢力.
就會被消滅.
Fairy Tales就是教小朋友.
告訴他們.
他的盼望他的信靠可以在哪裡.
你說說這些東西.
嚇怕小朋友.
有一個很出色的文學家.
他說童話故事是什麼.
童話故事不是告訴小朋友怪獸是存在.
童話故事不是告訴小朋友這個世界有怪獸.
童話故事不是想做這些東西.
因為小朋友本身已經知道怪獸是存在.
童話故事是告訴小朋友怪獸是可以被殺掉.
於是小朋友看完童話故事應該明白的是什麼.
就是他有勇氣可以面對這個世界.
啟示錄其實就是想做這件事.
二十一章第六節.
寶座上的說.

$^{641}$成了我是阿拉伯我是俄梅亞.
我是起初的我是結中的.
我要把生命的泉水白白賜給那口渴的人.
得勝的就會承受這些威脅.
我要作他們的上帝.
他要作我的兒女.
你知道英文字有多少個字母.
二十六個是不是.
希臘文呢.
二十四個.
第一個字母是Alpha.
Alpha.
最後一個是Omega.
俄梅亞.
即是說上帝是掌管整個的歷史.
掌管我們整個的生命.
由起初到終結.
由起初的天地到將來的天地.
整套劇都在他身上.
是他掌管是他導的.
是他導演的.
我老婆很喜歡看那些警匪片.
我經常都很Wonder.
為什麼她不喜歡看這些.
因為很多警匪片都是變態的.
即是有個瘋子.
抓了一些人也好虐待他們也好.
那些情節就很緊張.
有很多Suspense.
因為那個惡人是無法無天.
為什麼可以這麼無法無天.
那個時候因為她電視的那個地方.
就是我工作的那個地方.
我就經常聽到她有些反應.
她說這樣也可以.
但這些警匪片或者這些偵探片.
最終做什麼呢.
就是說公義會彰顯.
最後會有Justice.
你就大快人心.

$^{681}$有時候她還會重看.
為什麼會重看.
你都知道結局了.
劇透了 為什麼要再重看.
我們先想想.
拍得好的這些劇集.
或者拍得好的電影.
可以讓你重看.
縱使你知道結局.
就證明它真的可以.
你看的時候是津津有味.
縱使你知道結局是如何.
你仍然在劇情裡面是Enjoy的.
這個就是聖經.
終末的啟示要做的事.
我們知道結局.
同樣我們知道中間.
雖然生命裡面可以有很多曲折.
但是我們仍然有盼望.
得勝的.
就可以承受新天新地為賢.
在苦難的事件裡面.
我們知道上帝的工作.
繼續就進展.
上帝要帶領我們去的目的地.
仍然是一樣.
新天新地他的工作必定會成全.
這個就是終末的盼望.
讓這個盼望成為你的鼓勵.
如果你看見這個劇透.
最終會發生什麼事.
你就應該今天.
如果你生命裡面有一些傷痛.
我們會有安慰.
這個是聖經裡面所講.
屬天的安慰.
亦都讓我們生命裡面.
可以有使命有Purpose.
即是我們和上帝一起同工.
將天地帶來.

$^{721}$新天新地要來到.
今天不知道你心情是怎樣.
新天新地.
可以因為見到這個最終的結局.
你今天離開.
可以帶著信心.
可以帶著勇氣.
上帝是信得過.
可以同心低頭祈禱.
上主我們帶著我們的問題.
帶著我們的憂慮.
生命裡面可能有很多不同的掙扎.
不過我們知道.
你給我們新天新地的盼望.
縱使我們在生活裡面.
我們見到有死亡.
見到有疾病.
見到有很多.
不是很妥的地方.
有很多憂慮.
上帝讓我們知道.
你是那一位仍然帶領著歷史.
仍然帶領著我們生命的神.
主我們今天來到信靠你.
這個盼望.
讓我們知道這個盼望.
去信靠你的時候.
我們可以超越死亡.
超越一切壓迫我們的勢力.
因為你是得勝者.
主要叫我們都能夠成為你的得勝者.
讓我們知道耶穌基督在十字架.
空墳墓裡面所有的承傳的救恩.
成為我們今天生活裡面的鼓勵.
我們能夠有使命.
帶著使命和你同工.
當我們憂傷難過的時候.
我們能夠將教徒.
知道你是看管.
控制.

$^{761}$控制一切.
在我們生命裡面我們能夠去相信信靠你.
參見我們離開這個地方.
我們帶著上帝你新天新地的盼望.
走到人群裡面.
走到我們家庭裡面.
將新天新地的一切美善.
能夠帶到地上.
主繼續幫助我們.
讓我們能夠信心仰望.
祈禱奉著名.
願上帝祝福各位.
能夠有平安.
離開這裡的時候.
帶著上帝給你的能力.
跑去見證上帝.
\newpage



\section{14:1-7}
\label{sec:mdfK7VGG0u4}
\textbf{New Heaven and New Earth 1 - (John 14\_1-7) - Rev. Mark Chan}
\newline
\newline
連結: \href{https://youtube.com/watch?v=mdfK7VGG0u4}{\texttt{https://youtube.com/watch?v=mdfK7VGG0u4}} ~~~~ 語音日期: 2023-08-29
\newline
\newline
\hyperref[sec:S12GkifPujs]{\small{< < < PREV SERMON < < <}}
~
\hyperref[sec:index_chronic]{\small{[返順時目]}}
~
\hyperref[sec:index_scriptual]{\small{[返順卷目]}}
~
\hyperref[sec:OzeSIq8OBsY]{\small{> > > NEXT SERMON > > >}}
\newline
\newline
$^{1}$- Thank you everyone..
Mark..
- Good morning to all of you..
And we're gonna turn to the word of God today..
And we're gonna start a new series..
And this new series of messages.
is on new heaven and new earth..
New heaven, new earth..
If you know the Bible a little bit,.
you should know that the Bible has quite a bit.
to say about our future, the conclusion..
Just imagine there's a song or there's a great musical piece..
When you're playing along and listen along the song.
or the musical piece, it's going somewhere, right?.
It's going somewhere, it's going to some conclusion.
or some coda..
This is what the Bible says..
The Bible, the first two chapters is on the first creation..
And the last two chapters is on the final,.
new heaven and new earth, the new creation..
So the Bible, if you read from cover to cover,.
you should know there's a grand story..
It going from the first creation,.
which is our current life and creation,.
and it's going somewhere in the final destination..
It's like going on a journey..
You're taking different turns.
and different detours in between..
However, there was a final destination going toward..
That is what the biblical story and biblical worldview..
We need to understand this so that we know.
where our story fits into God's whole grand story..
And in the middle of it, first creation,.
and then the last creation, the new creation,.
in between is what Jesus Christ did for us..
Jesus came to human history and he died on the cross,.
sacrificed himself to redeem us.
so that we can be the people of God..
And then he rose from the dead,.
and then he ascended to heaven.

$^{41}$and seated at the right hand of the Father.
and reigning as our heavenly King..
In between is Christ..
And we are going toward the final conclusion.
of this huge, grand story..
So the Bible has a lot to say about new heaven.
and new earth, our future..
The Bible has many different passages.
that has valuable teaching, very good teaching,.
and to remind us where are we going..
I don't know if you've, sometimes we think about.
questions of ultimate purpose..
You know, where am I going?.
Where am I going?.
And where will I go when I pass away?.
Or some of the friends that you may know.
have already passed away and left us..
And where are they going?.
Does the Bible have anything to say about our future?.
Right?.
Couple, in the past three weeks,.
our staff and our church have been pretty busy,.
you know, trying to get some, prepare for funerals..
And it's not a very nice thing to do..
I mean, it's not pleasing to do,.
but, you know, we come alongside with the grieving families.
and try to help them out..
And we have couple of funerals.
and for people who are dearest to us..
And we were in the funerals and we kind of said goodbye.
or better, see you later to those we love,.
but they passed away before us..
So, you know, I try to reflect on, you know,.
what the Bible has to say..
And why does the Bible teach us, you know, those things?.
What happens when you die?.
And where is the whole creation?.
Where is the whole world heading to?.
So today we're gonna turn to the first passage,.
which is in the Gospel of John, the Gospel of John..

$^{81}$Okay, let me read that first..
This is what Jesus said to the disciples..
He said, "Do not let your hearts be troubled..
"Trust in God, trust also in me..
"There is more than enough room in my Father's home..
"If this were not so, would I have told you.
"that I am going to prepare a place for you?.
"When everything is ready, I'll come and get you.
"so that you will always be with me where I am..
"And you know the way to where I'm going.".
Now, Thomas said, one of the disciples,.
"No, we don't know, Lord..
"We have no idea where you are going..
"So how can we know the way?".
And Jesus replied, he told him,.
"I am the way, the truth, the life..
"And no one can come to the Father except through me.".
Now, this short passage is actually part.
of the long passage, or long sermon,.
message Jesus told his disciples..
And that long, long portion in the Gospel of John.
is from chapter 13 to 17..
Sometimes we call this long portion.
as a farewell discourse, farewell discourse..
It's Jesus' farewell discourse..
He's trying to say farewell to his disciples..
And you remember this was the last night.
that they had the last supper together..
And after that, in 18.1, so after chapter 13 to 17,.
many different things put together..
All the things Jesus has said to his disciples,.
his last night on earth, farewell discourse..
And then after all of this, 13 to 17, 18.1,.
it says, "After saying these things,.
"Jesus crossed the Kidron Valley with his disciples.
"and entered into a garden..
"And that garden, as we know, is Gethsemane..
"And he was arrested, betrayed,.
"and he was taken to Jerusalem, and he died on the cross.".
So the little passage that I have just read from John 14.

$^{121}$was part of this entire farewell discourse..
I think about it, when you have friends.
who are moving somewhere else,.
and you have a farewell dinner, right?.
What did they say, and what have they done,.
usually are very, very impressive to you,.
and impressive in your mind,.
and you probably would remember the moment.
or the environment that your friend.
has said farewell to you..
Or before someone passed away,.
usually the moment that he or she passed away..
There are times that the last moment of his or her life,.
he or she may have said something.
or may have done something.
that you will remember the rest of your life..
Now, this is kind of the farewell discourse.
we quoted in the Gospel of John..
And Jesus said many things to his disciple,.
and one of the things he said was to prepare a home.
for you and for me in heaven..
So there are a couple of points.
I wanna point out from this passage..
The first point is Jesus prepares our home in heaven..
Let's take a look at 14, one, and one, two, three..
Don't let your hearts be troubled..
Why?.
Because Jesus is about to leave them, right?.
But Jesus said, "Don't let your hearts be troubled..
"Trust in God, trust also in me..
"There's more than enough room in my Father's home..
"If there were not so, would I have told you.
"I'm going to prepare a place for you?.
"When everything is ready, I will come and get you.
"so that you will always be with me where I am.".
Now, in the time of Jesus, 2,000 years ago,.
in the first century, usually there were.
not many single houses..
You know, in Chicago, you drive around.
in different streets and roads,.

$^{161}$you see a lot of different single houses, right?.
How many of you live in a single house?.
Single house, right?.
How many of you live in a townhouse.
where you share a wall with other people.
or condo, right, or dormitory, right?.
Not many single houses in Jesus' day..
Most of the houses, most of the rooms.
are connected and attached to one another..
What is happening is that usually,.
according to the custom, is that when the son.
is ready to get married, usually around 18 or 19 year old.
or even earlier, younger, when the son is ready.
to get married and the parents will arrange the marriage,.
in those days there were no, you know,.
just dating and all that stuff,.
and the parents do all the things for the children,.
especially for their son..
So when the son is getting ready to get married,.
two families or two pairs of parents will come together.
and they talk about the terms or bride prices.
and all that stuff, and they sort out everything.
and they arrange the marriage for the young man.
and young woman..
And the son is getting ready and he will go with the bride.
and he will go with his father to the village of the bride..
You know, set out all the terms and sort out.
all the arrangements for marriage,.
and then the son will return back to his father's home.
and beside the father's house,.
he will build another small house..
And usually it will take about a year..
And then after that year, they engage for one year.
and then the son will go to the village.
and will receive the bride back home..
So usually the father's house,.
it has attached to the son's house,.
and then generation after generation,.
all these houses are connected together..
And then in the middle of it,.

$^{201}$probably is a courtyard or garden or something,.
and they share space in the middle..
Now this is quite similar to some Asian culture, right?.
Where all the houses are connected, right?.
You got the grandfather, you got the father,.
and then you got the son and the grandson.
and the great-grandson..
They are all living together.
and sharing one space in between, in the middle of it..
Now this is the same thing for the Palestinian culture..
When the son is getting married.
or getting ready to take the bride,.
he would take one year to prepare the house.
so that the bride will have a home..
Now with that kind of background,.
you understand what Jesus is trying to say..
He said, "I am leaving you now, temporary..
Just for the time being, I will not be with you..
However, I will come back and I will take you..
I will take you and I will bring you to my father's home,.
to my father's house, where I have already prepared.
a room for you.".
With that said, no, this kind of farewell discourse,.
this kind of farewell goodbye speech.
is not a matter of abandonment..
Jesus is not saying, "I'm abandoning you.".
Jesus is saying, "I am making you a promise,.
just like a groom promising a bride..
I will come and take you..
I will prepare a home, a heavenly home,.
where nothing can destroy, a perfect home for you.".
Last week, I was in Kansas City and I was in seminary.
and we have an intensive seminar..
So I was in Kansas City in the Airbnb for four or five days.
and it was pretty nice, pretty nice..
We booked an Airbnb and I shared the Airbnb.
with another pastor's friend..
And however, everything is clean,.
but I still don't like it..
I still don't like it..

$^{241}$A few days on, I kind of miss my home, right?.
I kind of miss my own bed, my own mattress..
Have you ever felt that way?.
You miss home, right?.
You're homesick, right?.
Yeah, of course you miss your family members,.
your sons and daughters and your wife and all that, right?.
This is the sense that within us,.
that we're yearning for home..
We get homesick..
Jesus is saying, "I am preparing a home for you.".
And that home is gonna be different..
It's in heaven and it's perfect..
And there's no pain and there's no tears..
There's no suffering in that home, right?.
We feel homesick from time to time..
And we know what that feeling is like..
A couple of weeks ago, as you know,.
Maui, right, Maui, the island in Hawaii,.
was destroyed by wildfire, big fire,.
displacing thousands of people..
And I was reading some news articles.
and some discussions online,.
and I ran into this opinion..
And the headline says,.
Maui fire shows that Hawaii paradise was a dream..
Have you ever been to Hawaii?.
Paradise, right?.
Paradise..
You know, a lot of billionaires,.
they just put a house there and, you know,.
build a house there, right?.
It's kind of like monopoly, right?.
Very nice place, I'm gonna build a house there, right?.
The fire consumes everything..
And the fire shows us this paradise on earth.
is simply a dream..
Where is our true home?.
Jesus said, "Your true home is not here..
"It's in heaven..

$^{281}$"It's the home I prepare for you.".
He said to his disciples, he said to us,.
"No matter how comfortable we are in our own homes,.
"they are not permanent..
"They're not eternal..
"Your loved one will pass away and will destroy the home..
"Your family may last for decades..
"However, at the end of the day,.
"the final enemy comes knock on your door..
"And his name is death..
"And your home will not last.".
The home in Hawaii won't last..
Our home won't last..
The only home that is prepared in heaven will last..
And that's where my yearning is..
That's where our desire should be..
So Jesus prepares a home for us..
And also, second point, is Jesus is the way home..
He not only prepared a home for us, but he is the way..
He's not just pointing us to the right way to go home..
He himself is the way home..
Verses five..
Thomas, one of the best Sunday school students, I guess,.
and he was curious..
After hearing Jesus' farewell and goodbye speech,.
Thomas said, "No, we don't know, Lord Jesus..
"We have no idea where you're going..
"And how can we know the way?.
"How can we go after you?.
"How can we follow you?".
Jesus told him, "I am the way, the truth, and the life..
"No one can come to the Father..
"No one can go home except through me.".
Here it says the way home, right?.
The journey or the route that will bring you home..
And you need to understand,.
it's a little bit geography here..
People living in Israel, especially in the ancient world,.
the land of Israel is the center of everything..
It connects Africa, Asia, and Europe..

$^{321}$Canaan, Israel, is the center,.
the heart of the entire world, the ancient world..
And therefore, there were many different roads.
and highways would pass through the land of Israel..
In fact, people living in Israel,.
when they imagine life, when they imagine daily living,.
they will talk in terms of walking..
How you live is how you walk,.
because every day you are walking..
Every day you are living right next to.
a huge highway intersections, right?.
Everyone has driven through downtown Chicago, right?.
There's so many different highways and exit, right?.
That is what land of Israel is like..
You live in Israel, you are like living right next to.
the intersection of different highways..
So when they talk about living, they talk about walking..
How you walk, walk in a way that is worthy of the gospel..
That means live in a way that will bring honor to Christ..
Roads and walking,.
many, many different passages talking about that..
Especially in the gospel of John,.
there were two places that talk about the way..
The first one, what the first one is here,.
Jesus said, "I am the way.".
And then chapter one, in the very beginning.
of this gospel of John, Jesus or his cousin,.
John the Baptist, the pioneer,.
also talk about the way, chapter one, 23..
John the Baptist, remember he is the pioneer.
of the Messiah, right?.
He's gonna introduce the Messiah to the world..
And John the Baptist start baptizing people by the river..
And John said, in the words of Prophet Isaiah,.
he said, "I am the voice shouting in the wilderness.".
To clear the way for the Lord's coming..
So John the Baptist is kind of like.
rolling out the red carpet.
for this honorable guest to come in..
And he is Jesus Christ, the Messiah..

$^{361}$John the Baptist is trying to clear the way.
in the wilderness so that the Lord Jesus can come..
He's the hero of the gospel, the kingdom..
Finally, the king is coming,.
he roll out the red carpet and this is him..
Now, he's quoting from the book of Isaiah..
And this has happened, the book of Isaiah.
talks so much about the way, the way for the Lord,.
the road for the Lord..
There are many different passages in Isaiah,.
especially the second half of Isaiah,.
chapter 40 and the following,.
talks so much about the way of the Lord,.
the way of the Messiah..
And one of the verses is 51, nine, Isaiah chapter 51,.
talk about wake up, wake up..
It's a prayer, a prayer that pray to the true God,.
the Lord, the God of Israel, clothe yourself with strength..
They're trying to wake up, wake God up..
God, wake up, this is the time for your salvation..
And rouse yourself in the days of old..
The old days was the first Exodus,.
when the people came out of Egypt,.
God defeated the oppressor and then saved his people..
And the Israelites became his people..
That's when, that is the days of old..
The one who dry up the sea, you see that?.
The first Exodus, they crossed the Red Sea on dry ground,.
making a way of escape, a way through the depths of the sea.
so that your people can cross over..
Again, this is a prayer,.
prayer that God will redeem his people,.
just like in the days of old..
In the days of old, the first Exodus,.
God save his people and bring them.
and prepare the way for them, they can cross the sea..
And now this is a time, God, you need to do it again..
And John the Baptist is quoting from the book of Isaiah.
and say, I'm gonna introduce you to the salvation.
that Jesus the Messiah gonna bring to us..

$^{401}$And it's gonna like the days of old..
And finally, people will return..
Those who have been ransomed by the Lord will return..
And they will enter the holy city of Jerusalem.
singing with everlasting joy..
There is a way home..
And Jesus said, amazingly, he said, I am the way..
I'm not pointing you the right way..
I myself is the way..
And Jesus said, he is the way, the only way we can go home..
The teaching in the Bible concerning the future,.
concerning the new heaven and new earth,.
is for us so that we know where we're going..
Imagine you're on the journey..
If you don't know your destination,.
if you don't know where you're going,.
you're wasting your time, right?.
The more you know about your destination,.
the more you understand where you're going,.
will determine how you make choices today..
That is the biblical teaching of the final things..
Trying to reveal to us what our final destination is.
so that in the detour and the turns in our lives,.
we know we won't get lost..
I remember some years ago, I was driving a Mini Cooper..
I really like driving Mini Cooper..
I sold it because of mechanical problems and all that stuff..
So I had a Mini Cooper for about seven years..
I really enjoyed it..
Have you driven a Mini Cooper before?.
No? You should try it..
It's like driving go-kart..
It's like driving go-kart all the time..
So I really enjoy it..
So I remember one day,.
well, right after we bought the Mini Cooper,.
and we have to bring it back to the dealer,.
to the shop for oil change or something like that..
And then it took about two days,.
the mechanic says,.

$^{441}$"It will take about two days for us to fix it.
"and get it ready for you.".
And then here's the loaner..
The loaner..
And I got a loaner from the car dealer..
It's another Mini Cooper..
No, it's brand new Mini Cooper with turbo..
So as I said, I really like driving Mini Cooper..
It's like you and the car, there's no distance..
I am the Mini Cooper, right?.
You know, if you haven't,.
if you've never driven a Mini Cooper, you should try it..
You know, it's really fun..
So I drove that car, I was really excited..
I got this loaner, I got this turbo, two-seaters,.
brand new Mini Cooper, and I drove home..
And then I talked to my wife and said,.
"Hey, I got a new Mini Cooper for two days..
"And let's go have some fun.".
So I drag her to the Mini Cooper and we drove off, right?.
I drove off, right?.
On the highway..
So it's kind of like, it's like driving a bullet..
So much fun..
And then after three minutes, she asked me,.
"Where are we going?.
"Where's the destination?".
I said, "I don't know.".
And then I thought to myself,.
"This is wasting gas.".
(audience laughing).
"Wasting money, actually.".
So then we drove back home..
Now, when you don't know where your destination is,.
it's no longer fun, right?.
Because your destination determines your decisions,.
determines the direction of your journey..
The Bible reveals to us our final home..
And we are all in Christ and heading that home..
No matter how comfortable we are here,.

$^{481}$this is not our final resting place..
Our final home is in heaven,.
in the new heaven and new earth..
With that said, it should apply to us.
how we're gonna make choices,.
how we're going to make different arrangement in our lives,.
how we prioritize all these different choices in our lives..
Going to a final home..
I remember one day, many years ago,.
when I was early 30s,.
that's some years ago, not too many..
(laughs).
Remember, well, I told you before,.
in my previous church where I was serving,.
we had a lot of funerals every year..
I remember one year I officiate 40 funerals, one year..
Almost like every single week, sometimes two times a week..
I was kind of tired.
because I have to do funerals all the time..
And not only people from my church,.
but there were many different people I needed to help out.
and to serve, they live in the neighborhood..
They live in the community.
and they want the pastor to help them out..
So there was one year I just officiated.
so many different funerals..
And you're just kind of numb at times..
But I remember one day, this family walked into my office.
and they said, "Pastor, can you help us.
"to prepare this funeral and officiate this funeral?".
I said, "Of course, of course, at your service.".
And then usually in the arrangement of funerals,.
I will ask the family to prepare eulogy..
Do you know what eulogy is?.
It's a life story, right?.
A life story of the deceased..
So I used to ask the family members,.
"Can you prepare the eulogy?".
And you write it out and then I'll double check it.
and then we're good to go..

$^{521}$So they did write down the eulogy for the deceased..
And they hand it to me the next day..
And I kind of check it, spell check and all that..
And this struck me that I know the life story.
of this person I'm going to officiate a funeral for..
He was born the same year as I was..
Same birth year..
I was in my early 30s..
And it really impacted me because it was the moment.
for me to ask the ultimate question..
What am I doing here, right?.
What am I doing?.
If I die today, where am I going?.
Am I sure where I'm going?.
Remember the funeral, the deceased was early 30s,.
same year, and read his life story.
and officiated the funeral..
In that funeral, it was so sad.
because there were only two family members show up..
Only a handful of people in the funeral..
No one mourned for him, no one really cared..
Even his daughter is not there..
So I still remember that funeral happened.
in the harsh winter of Chicago..
So we did the funeral and then we went out to the cemetery.
to perform the burial service..
And it was a big snowstorm..
And it was just crazy..
You know how crazy the harsh winter in Chicago, right?.
It was snowing so much and then the workers, of course,.
stick the ground six feet below.
and prepared the place for his final burial..
And then the workers had to lower the casket into the vault..
And it was so windy and so much snow..
And there was so much trouble trying to fit the casket.
into the vault, right?.
Even I and the funeral director have to help out..
So we grab some two by four and try to make sure.
the casket will fit so they can close the lid..
It was tough..

$^{561}$It was so, so sad..
And the grief of the family..
I think the most important thing for me.
is that what am I doing here, right?.
Am I using my resources faithfully?.
Am I being faithful to the call of God?.
What is my vocation?.
What is God calling me to do?.
And with my limited time on earth,.
what am I doing with the resources he has given me,.
the gifts that he has granted to me?.
What am I doing with all this stuff, right?.
And someday I will be home..
Someday I will see him again..
I will see Jesus, our Lord again..
And how can I answer?.
How can I reply to him?.
What have you done with your life?.
As I told you two weeks ago,.
we officiate a funeral for Leah, my dearest friend..
I've known Leah for years, since she was in college..
And it was just terrible suffering.
that she has gone through in the past three, four years..
She was discovered with, diagnosed with.
a fairly aggressive form of brain tumor..
And the last three and three plus years,.
she has been battling with cancer,.
you know, on and off, on and off..
And it's just, can't imagine what she has been going through..
But if you listen to her story,.
she actually prayed that, "God, can you use my suffering?.
"Can you use my suffering to tell your story,.
"to testify the goodness of God,.
"even in the midst of trouble?".
Now that kind of prayer, I tell you,.
honestly, I cannot make that prayer, right?.
From where she get the strength, right?.
From where she get the strength to pray like that, right?.
It has to be a hope beyond death..
Has to be the reality of the final home..

$^{601}$She knows where she's going..
And she has that hope anchor in her life..
And she could pray like that,.
even in the midst of most difficult trouble..
How can I use my time to witness to God?.
The future reality..
A couple of weeks ago, it was on Monday, yeah..
So I got a call that, you know,.
Leah in her final moment..
So I got a call and I rushed to her home..
And she passed away peacefully at home..
And there were different brothers and sisters were there..
We were with the parents,.
and actually the parents were here today, this morning,.
worshiping with us..
Very difficult time for them..
And that Monday morning when she passed away, we were there..
And her body was there..
And we wait for hours for the funeral workers.
to pick up the body..
So very, very difficult..
But we take comfort, the fact that she's now home..
In her final home..
And that's the home we are all heading to..
And during those few hours we were waiting.
for the workers to pick up the body,.
one of the brothers who has been in contact.
with the funeral home, he was waiting outside..
And he took this picture..
It's actually the background of the PowerPoint..
It's a beautiful picture..
I'm sure the heavenly home is so much more beautiful.
than this..
That's what Jesus is talking about..
I prepare a home for you..
And that's where your hope is..
What we live according to this hope..
We keep in mind we are heading there..
And we are on our way..
Every moment in this life counts..

$^{641}$We are kept accountable before God..
Use our time, use your resources, use your gifts and talents.
to serve God..
Because someday we're gonna give an account..
Someday we're gonna be home..
Let's all pray..
Let's all pray together..
Father, we give you thanks because you prepare a home for us..
That our Lord Jesus not only died for us on the cross,.
only rose from the dead..
He triumphed over death..
And he rose from the dead and give us this hope,.
this eternal hope..
And we hear his promises today..
That our Lord Jesus, you have prepared a home for us.
in heaven..
And we're all heading there..
And Father, may our desires and our direction.
are anchored toward that hope..
And I pray that we'll build our lives.
according to this final destination we are heading to..
And I pray that you will help us to live in wisdom,.
to use our resources faithfully,.
so that in our earthly period of time,.
we'll bring you honor..
Use our lives to serve you..
Good time or bad times, may you reign in us..
In Christ, we pray..
(gentle music).
(upbeat music).
\newpage



\section{帖撒羅尼迦前書 4:13-18}
\label{sec:OzeSIq8OBsY}
\textbf{新天新地:人生終末的再思 3 將來的君王 (帖撒羅尼迦前書 4\_13-18) - 陳崇基牧師}
\newline
\newline
連結: \href{https://youtube.com/watch?v=OzeSIq8OBsY}{\texttt{https://youtube.com/watch?v=OzeSIq8OBsY}} ~~~~ 語音日期: 2023-08-29
\newline
\newline
\hyperref[sec:mdfK7VGG0u4]{\small{< < < PREV SERMON < < <}}
~
\hyperref[sec:index_chronic]{\small{[返順時目]}}
~
\hyperref[sec:index_scriptual]{\small{[返順卷目]}}
~
\hyperref[sec:d5kiiPL31Z0]{\small{> > > NEXT SERMON > > >}}
\newline
\newline
帖撒羅尼迦前書 4:13-18
\newline
\begin{longtable}{cl}
\hline
\hline
章節 & 經文 (和合本修訂版)\\
\hline
4:13 & \begin{tabularx}{0.7\textwidth}{X} 弟兄們，至於已睡了的人，我們不願意你們不知道，恐怕你們憂傷，像那些沒有指望的人一樣。 \end{tabularx} \\ \\ \relax
4:14 & \begin{tabularx}{0.7\textwidth}{X} 既然我們信耶穌死了，復活了，那些已經在耶穌裡睡了的人，神也必將他們與耶穌一同帶來。 \end{tabularx} \\ \\ \relax
4:15 & \begin{tabularx}{0.7\textwidth}{X} 我們照主的話告訴你們一件事：我們這活著還存留到主來臨的人，絕不會在那已經睡了的人之先。 \end{tabularx} \\ \\ \relax
4:16 & \begin{tabularx}{0.7\textwidth}{X} 因為，召集令一發，天使長的呼聲一叫，神的號角一吹，主必親自從天降臨；那在基督裡死了的人必先復活， \end{tabularx} \\ \\ \relax
4:17 & \begin{tabularx}{0.7\textwidth}{X} 然後我們這些活著還存留的人必和他們一同被提到雲裡，在空中與主相會。這樣，我們就要和主永遠同在。 \end{tabularx} \\ \\ \relax
4:18 & \begin{tabularx}{0.7\textwidth}{X} 所以，你們當用這些話彼此勸勉。 \end{tabularx} \\ \\
[1ex]
\hline
\hline
\end{longtable}
$^{1}$好 各位電視節目的觀眾 我們繼續今天的主題.
過去的三講.
兩講.
這次是第三講.
我們講到新天新地.
聖經裡面啟示了一個世界的觀念.
或者怎樣去看我們人生的觀念.
不是一個沒完沒了不斷循環的觀念.
看人生或者看世界 看歷史.
都不是一個沒完沒了一直在循環.
好像一首樂章一樣.
總是有一個目的地要去的.
那個叫Coda或者conclusion.
一定會是一條直線.
雖然可能有時候會重複一些東西.
但到最後會有一個結論.
這個就是聖經告訴我們有的世界觀.
就是新天新地.
聖經由開始的第一頁.
到它最後的那兩章.
就是好像那本Book N那樣.
告訴我們整個的時間表.
由頭到尾上帝是帶領著整個歷史.
而在中間發生什麼事.
中間就是耶穌基督來到.
祂的誕生 祂傳道.
然後在耶路撒冷上面.
祂寫在十字架上面.
為我們罪人擔當我們的罪.
然後第三天.
埋葬之後 祂就復活了.
墳墓是空的.
然後呢 神的兒子也都升天.
現在坐在天上面來作夢.
將有一天 祂會再回來.
所以這個聖經告訴我們.
這個世界觀.
是和世界裡面所說的世界觀是不同的.
向著一個結論.
向著一個意義來走向.

$^{41}$而到新天新地來到的時候.
就是基督再來的時候.
人在基督裡面就會復活了.
我們就會有一個不壞的榮耀的身體.
然後新天新地就會降臨.
好像新的耶路撒冷聖城一樣.
是沒有罪 沒有邪惡的一個地方.
沒有痛苦 沒有疾病 沒有哭號的一個地方.
那個地方就是我們正在去的目的地.
所以聖經裡面說很多終末的事情.
對我們人生是很有關係.
因為我們不是在這裡沒完沒了的.
去做我們一生裡面經營的事.
而是朝著一個方向.
朝著一個新天新地來去.
這次是第三講.
我們會看一段經文.
就是貼在我現在的前書.
這段經文我們看的時候.
我只有兩個重點.
第一個重點就是.
說到這位的王者會再來.
就是耶穌基督我們的主.
這個死而復活.
並且在天上來作王的主耶穌基督.
祂會再來的.
並且聖經宣告是肯定的.
聖經的啟示.
上天給我們的啟示.
揭開這個奧秘.
就是這個第二次來臨.
這個將會再來的君王.
是肯定的.
第二點就是.
祂的再來是越來越接近.
越來越近.
我們先看經文.
貼在我現在前書第四章.
抵門至於已睡了的人.
我們不願意你們不知道.

$^{81}$恐怕你們憂傷.
像那些沒有子望的人一樣.
既然我們信耶穌死了復活了.
那些已經在耶穌裡面睡了的人.
上帝也必將他們與耶穌一同帶來.
我們照著主的話.
告訴你們一件事.
我們這活著還存留到主來臨的人.
絕不能在那已經睡了的人之先.
因為召集另一法.
天使長的呼聲一叫.
上帝的號角一吹.
主必親自從天降臨.
那在基督裡面死了的人.
必先復活.
然後我們這些活著還存留的人.
必和他們一同被提到雲裡.
在空中與主相會.
這樣我們就要和主永遠同在.
所以你們要用這些話.
這些話彼此勸勉.
這段話是保羅寫信給當時的教會.
在帖撒雷加的一個城鎮.
一個羅馬的大城寫給他們的.
這個初生的教會可能有很多的struggle.
特別他們可能中間有些親人去世.
或者有些人離開了.
他們就憂傷難過.
這是難免的也是自然的.
但是保羅在這封書信.
帖撒雷加前書.
之後還有一本後書.
前書他講什麼呢.
講到第四章的時候.
他就說提一提這個很重要的事.
那以色說了的人.
我們知道他將來會發生什麼事.
而且我們還存留的.
如果我們有福氣.
耶穌是在我們有生之年回來的話.

$^{121}$我們會發生什麼事.
這是很重要的一段話.
這些的終末的思考.
或者終末的教導.
不是一些抽象的東西.
是和我們今時今日面對的事有關.
生活裡面可能有憂傷.
會有難過.
甚至會有困境.
不過這些終末的教導.
是給我們生命裡面有盼望.
有力量.
當你看到目的地的時候.
你的人生的方向就不會這麼容易陣腳大亂.
當我們面對有不同的險阻.
不同的眾叛親離等等.
只要我們見到最終的終點.
我們今時今日就會有所不同.
第四章第十五節這裡說什麼.
他說我們照主的話告訴你一件事.
主的話是誰.
就是主耶穌基督說過的.
這裡就是主耶穌曾經說過的東西.
然後初期的信徒.
初期的使徒.
他們記起那一件事.
我不知道他真正在引用的是哪一道.
可能是鑒探山的港道等等.
你記得面對耶路撒冷的時候.
耶穌進入耶路撒冷之前.
他在鑒探山上說過一些有關主在來的事情.
特別是主後70年.
會有什麼動盪的事情會發生.
然後有戰亂.
有饑荒等等.
羅馬軍隊會打過來.
然後用這個情況.
他又說了主的在來.
這種預言.
聖經裡面有很多這樣的應許.

$^{161}$由舊約至新約.
神的承諾.
我不喜歡說預言這麼多.
因為現代人一聽到預言.
我們就覺得是預先寫好的報紙.
但聖經裡面所謂的預言.
其實是promises.
其實是應許.
或者是承諾.
舊約裡面有很多對將來的承諾.
對將來的應許.
是在歷史裡面發生的.
改朝換代.
或者地動山搖.
世事有變遷.
你去查一下舊約裡面的歷史和背景.
你就知道有很多重要的應許.
都已經實現了.
最重要的應許.
實現在誰身上呢?.
就是整個人類歷史中間這一段.
就是耶穌基督的來,死,復活和升天.
這些在舊約裡面有的promises.
都在耶穌基督的身上實現.
舊約,新約.
但在整本聖經裡面有一個預言.
或者一個應許.
是還沒有實現的.
並且一定會實現的.
那個是什麼應許?.
就是耶穌的再來.
耶穌的second coming.
將會實現.
並且必定會實現.
所以這裡說主的話.
去告訴你一件事.
這件事一定會發生.
將來一定會發生.
你怎麼知道?.
因為你看一下上帝有沒有承信.

$^{201}$你去查一查舊約至新約.
有沒有一些東西是上帝曾經說過.
所謂預言應許過的.
有發生過的.
所以上帝的character.
是可信的.
並且上帝的啟示告訴我們.
second coming.
耶穌的再來.
是必定會發生.
這個certain是肯定的.
這是第一點.
照著主的話告訴你一件事.
就是second coming的時候會發生什麼事.
我們看看他怎麼去描述.
這個所謂second coming是怎麼樣.
還有這個second coming.
是和我們之間有什麼關係.
這個描述主的再來的情況是怎麼樣.
他說末後的時刻.
就是說人類歷史的終結.
這個樂章的終結的時候.
到底會發生什麼事.
他說在基督裡面的人就會復活.
這裡有一個次序.
15節.
我們照著活著還存留到主來臨的人.
絕不會在那些睡了之人之先.
因為召集令一發.
天使掌的呼聲一響.
上帝的豪國一吹.
主必親自降臨.
那些在基督裡面死了的人必先復活.
然後我們這些活著還存留的人.
必和他們一同被提到魂裡面.
在空中與主相會.
這樣我們就會永遠與主同在.
這裡要搞清楚時間線.
誰會先復活.
睡了的人.

$^{241}$就是說那些已經離開了我們.
在教會離開了我們.
在基督裡面他們首先離開了.
這一班人就是所謂已睡了的人.
他們好像睡著了.
睡著了就會醒了.
他們好像暫時在睡覺.
所以首先要復活的就是那些已睡了的人.
就是已經安息主位的人.
然後第二班是什麼人.
就是我們還存留的人.
現在我們還活著的這班人.
還存留到主再來臨的時候.
但是誰先呢.
就是他們先.
然後我們是跟著的.
我們所謂第二組的人.
然後會有什麼事情發生呢.
就是豪國的聲音.
上帝皇者再來的德勝的聲音.
這些是打仗德勝的海歌.
就好像一隊軍隊打完仗之後.
海船回歸回到他們的京都的時候.
那種豪國聲和德勝的聲音.
主就會親自從天降臨.
發生什麼事呢.
當主來到的時候.
豪國聲德勝的聲音吹響的時候.
睡了的人起來.
然後我們活著的人.
也都會一同加入他們.
被提到雲裡面.
在空中與主相遇.
然後我們就與主永遠同在.
這兩句話我們需要解釋一下.
第一句就是被提到雲裡.
特別這裡有很多華人的信徒.
特別是老信徒.
我想我們以往年輕的時候.
我們都聽過所謂的災前被提論.

$^{281}$Pre-tribulational rapture.
被提的字就是Rapture.
在這段聖經裡出來的.
被提到雲裡.
大多數舊有的一套.
就是我們比較老一點的.
我們接受了這一套的觀念.
這一套觀念是什麼意思.
就是說基督再來的時候.
第一次他再來.
就是在說2000年以前.
人類歷史.
然後他再來第二次的時候.
首先有兩個階段.
有一個階段是沒有人知道的.
他來的時候.
被提的時候我們突然間就會消失.
所以這裡坐的全部都是信主的.
說一聲我們就不見了.
接著就會七年大災難降臨在世界上.
就是前所未有的大災難.
然後會非常慘痛.
到七年大災難完結之後.
主真正的再來就來了.
他就是榮耀的來到.
然後所謂收拾殘局.
然後帶來永恆的國度.
那些千禧年我們不鬥了.
我們不談了.
有一點複雜.
舊有的一套就是說這件事.
被提到雲裡面.
我們就是所謂脫苦海.
然後七年大災難才來到.
災前被提就是這件事.
In any case.
我想這段聖經是不是真的說這件事.
首先搞清楚.
這件事叫做在空中與主相會.
什麼意思被提.

$^{321}$然後在空中與主相會.
其實是一尊有名詞.
形容當時羅馬帝國.
那些軍隊將軍帝王帶著軍隊去打仗.
打完仗回來海船回歸的時候.
他進入京都之前.
京都的使鎮的長老.
他們就會組織一個歡迎的隊伍.
然後他們出城.
在路途中間與德勝的軍王迎接他.
所謂相會就是這個意思.
然後相會之後.
慶祝之後.
他就帶他回去京都.
所以那些城中的長老.
就會一出然後就回頭.
迎接一個軍王入城.
迎接一個德勝主.
一個王者入城.
是一種這樣的儀式.
一種這樣的習俗.
在羅馬帝國特別多.
剛才我說的貼薩萊加這個地方.
是一個大城.
是軍隊駐紮的地方.
所以他們應該很清楚.
保羅在這裡說的話.
與主相會.
一講到相會的時候.
就想起這種的習慣.
這樣的儀式.
如果你了解和想一想.
這些當時的儀式.
其實都很合理的.
有大官貴人.
或者有重要的人來到的時候.
你應該是走去迎接他的.
德勝我們教會請了一個特約的講員.
他飛到來我們芝加哥的時候.
我們會不會跟他說.

$^{361}$不如你自己坐車吧.
我們會在教會見你的.
你應該不會這樣.
禮貌上.
不想他來也好.
你都應該去迎接他.
上個星期我剛從加拿大回來.
上了一個星期課.
星期五我飛回來的時候.
你知道是誰來迎接我嗎.
是CTA來迎接我.
原因是因為我不重要.
電話都沒有通過.
總之你自己回來自己吃東西.
然後你繼續工作.
但如果你是大官貴人的話.
就不同了.
就是這種相會的儀式.
保羅在說.
Meeting together.
新約裡面另外兩處用了這個字.
整本新約翻譯了.
只有三次用了相會的字.
這是第一次.
第二次在哪裡呢.
就是史蘭傳第28章最後的時候.
那裡講到保羅.
他上京上告該殺的路途.
到最後.
conclusion的時候.
他快要去到羅馬.
那裡講什麼呢.
他說羅嗎哪裡的弟兄們一聽見我們快要來到的消息.
就到了哪裡呢.
阿比烏市和三館來到迎接我們.
愈發壯膽.
我們就進了京都羅馬城.
哪個字是和相會有關.
就是迎接.
迎接這個字就正是同一個字.

$^{401}$與主相會.
這幫人做什麼.
羅馬的教會弟兄.
他們就像接機一樣.
他們就離開羅馬的京城.
然後去到一個地方叫阿比烏市.
或者三館等等.
這個地方無論如何.
距離羅馬有多遠呢.
40個mile.
你今天開車.
高速公路都要45分鐘.
這幫羅馬的信徒聽見.
保羅和他的隊伍快要來到羅馬.
他們就特意有一個歡迎的.
接送的小組.
他們就出去.
羅馬城以外40多英里.
特意在那個地方.
去到相會.
去迎接他們.
然後呢.
迎接之後就去哪裡.
不是離開羅馬.
然後他們就帶著保羅回到羅馬.
這個是史林傳的結論.
其實是一個很妙的結論.
就像我講的一個惡章的一個coda.
最後最後的時候會發生什麼事.
這裡講到不是羅馬王進入京都.
而是保羅帶著福音.
帶著耶穌的見證.
帶著王者進入當時帝國最中心京都的地方.
這個就是史林傳最後的結論.
基督耶穌的福音.
雖然會在艱難的困所裡.
但是會植入羅馬的京城.
你要信福音.
是有這樣的power.
你信了福音的時候.

$^{441}$應該有這樣的改變.
不單止改變我們自己的心靈.
可以扭轉整個的帝國.
相會就是這樣的意思.
迎接保羅人之後回到京都.
就是海船回歸.
第三次用到相會這個詞.
就在馬太福音第25章.
那時候上文夏利講了什麼呢.
就是十個童女的比喻.
你記得嗎.
十個童女的比喻.
就是十個伴娘.
當時的伴娘.
他們就在自己的鄉鎮裡等新郎回來.
為什麼呢.
因為新郎通常是在另一個村子裡娶新娘.
他一早就去那裡.
然後接了新娘回來.
通常都要搞到很晚.
好像我們華人拜年一樣.
拜完年就走不了.
走了兩三個小時.
然後又坐下來打了兩圈麻將.
然後又吃了一些東西.
吃完東西又再吃.
然後搞幾個小時你才能走.
當時巴勒斯坦的婚禮就會是這樣的.
新郎走了去.
跟他的鄧穿石.
接新娘回來的時候.
通常都是晚一點才回來.
凌晨的時候才回來.
十個童女就睡了.
因為等了太久了.
這次要加班了.
結果發生什麼事.
第六節.
半夜有人就喊叫.
他說新郎來了.

$^{481}$你們出去迎接他.
伴娘你起來了.
起來.
因為新郎快來了.
你們去等他還是怎樣.
你們出去與他相會.
去迎接他.
是同一個字.
相會了之後.
迎接了之後去哪裡.
不是離開.
而是回家.
然後又再大事的慶祝.
直到清早的時候.
這個就是那個意境.
相會.
保羅在說什麼呢.
用一種這樣的迎接君王的儀式.
迎接一個海船回歸的將軍.
王者.
駕到的時候.
這個就是主在來的時候.
會發生什麼事.
已睡了的人.
暫時離開了我們的人.
好像睡覺一樣.
當基督再來的時候.
號角聲吹響.
他們首先會復活.
被提到雲上.
我們這班仍然存留的.
我們也會加入他們.
不是去脫離這個地方.
而是我們空中相會.
將基督帶回來.
然後下來這個地上.
新天新地.
就可以與基督一起來作夢.
這個就是聖經裡面說得很清楚.
終末.

$^{521}$The End.
終末給我們的應用.
為什麼寫這段話.
不是一些抽象的教導.
是因為在貼上教會.
他們正在掙扎這個問題.
有人離開了我們.
他們將來會怎樣.
我們呢.
我們將來會怎樣.
終末的教導.
就是告訴你.
你的終極目的地.
你越想清楚你的終極目的地.
你現在走的路.
就不會走冤枉路.
你走的每一步就會有意思.
你想像一下.
如果沒有基督再來.
如果歷史只不過是不斷循環.
不斷無意義地重複的話.
我們的人生會怎樣.
我們很多華人.
都可能有這樣的觀念.
那就認命.
生命就是這麼苦.
完了就沒有了.
到最後你很辜負的時候.
你就認命.
這個就是世界給我們的觀念.
但聖經的終末的教導.
就是揭開一個新的啟示.
不單止說這個生命裡面.
是有終結.
並且你和我.
是會去到某一個的目的地.
而那個目的地就是基督的再來.
基督得勝的凱歌會就位.
這個就是用聖經的角度.
來思考你人生的終極的時候.

$^{561}$你想一下有多大的分別.
想像一下那個意境.
就是基督已經得勝了.
他一步一步的已經逼近我們這個地方.
一直慢慢來的.
已在路上.
一個王者.
一個已經得勝的主在海船回歸.
正在海船回歸.
每過一日.
他就步近多一日.
每過一日就代表我們會快見主.
多一日.
會counting down.
如果你這樣看的時候.
生老病死.
可能有另外一個完全不同的角度.
主在來帶給我們的.
就是這個所謂的盼望.
他必定會到.
並且越來越近.
就教我們今天生活.
我們會有盼望.
會有活力.
今天有的時間.
不如就為主打拼.
因為他的再來.
是越來越接近.
肥立比書第一章.
用這章做結束.
這裡是肥立比書另外一篇的.
保羅寫的書信.
他說到他因為福音的緣故.
他被困所.
有些人說肥立比書其實是在羅馬.
他上京的時候.
然後他就被house arrest.
保羅說什麼.
他說因為我活著就是基督.
死了就有益處.

$^{601}$但是我現在在肉身裡面活著.
仍能有工作的成效.
我就不知道怎樣挑選.
因為正處於兩難之間.
我情願離開這個世界.
與基督永遠同在.
因為這個是好得無比.
但兩難之間是怎樣.
我現在仍然活著.
我仍然有工作的成效.
我仍然會為主打拼.
不知道要選哪樣.
兩樣都有意義.
但更加好的就是什麼.
與基督同在.
永遠好得無比.
所以死了就有益處.
這個生死觀.
或者生死的教育.
怎樣看我們老去.
怎樣看我們生老病死.
怎樣經歷.
是完全另一回事.
不是一種認命.
而是有盼望地我們步向將來的結論.
最近我們有些大兄姐妹.
因為我們搞一些安息禮拜等等.
當然是一些很困難的時間.
但也有一些我們休息的時候會談的東西.
有些人說將來.
因為我們每個人都要面對這些.
福地 墓地 土葬 火葬等等.
我感受到一件事.
我們基督去面對.
我們可以從從容的去面對.
沒錯我們會有憂傷會有難過.
這是必定的.
但怎樣可以從從容的面對.
因為我們知道destination.
知道終點.

$^{641}$而終點就決定我們今天的生活.
今天的取態.
生老病死是必然的.
好像保羅這裡所說.
可以從從容的面對.
是有困難但仍然可以有盼望的.
望向前.
搞安息禮拜的時候.
大家在說現在原來有很多很有創意的安葬方法.
我在研究的時候.
我發現原來土葬不只是火葬.
有些非傳統的.
有些可以葬在樹上.
將你的骨灰放在樹上.
當然越大棵樹越貴.
越小棵樹就越便宜.
有些.
剛才說這件事的時候.
有我想起好幾年前.
我有一個很好的朋友.
他特意來香港.
來美國.
送他其中一個家人離開這個世界.
他告訴我怎樣葬.
原來他走去那些太空發射火箭的場地.
去看安葬禮.
安葬禮不是向下是向上.
那些骨灰放在火箭上射上去.
原來有些太空葬禮.
Space Burial.
如果你想搜索的話.
Space Burial.
當然不是很急著說這些.
這個我不收佣金.
我自己想認識多一點的時候.
這個Space Memorial.
原來有不同的套餐.
如果你射上去然後掉下來融化的話.
最便宜了三千元.
但最貴的多少呢.

$^{681}$一萬多元.
一萬多元就射你上去深太空.
當然我不會說大家.
不吉利.
但你覺得我是不是應該用這樣的方法.
有誰是贊成的.
我還沒問過我家人.
有誰是贊成的.
你贊成嗎.
我知道你不想來探望我.
射我上太空.
我就想告訴你.
基督給我們的盼望.
是我們可以從容的去面對.
死亡心知.
死了就有逆轉.
你知道以前有一個聖經教授.
叫做Eugene Peterson.
他以前在加拿大溫哥華.
維珍神學院.
當時我讀神學學位的時候.
他就在那裡教.
他就將整本聖經翻譯成現代的英文.
就是一種paraphrase的翻譯.
意譯的譯本.
他怎樣譯這段聖經.
菲勒比斯第一章.
他說兩難之間.
保羅說我活著就是基督.
死了就有逆轉.
他意譯成怎樣.
他說life versus even more life.
他說原來我們人生不只是life versus death.
這是世界給我們的世界觀.
沒有意義的.
基督的復活.
基督的再來.
就告訴我兩個選擇.
就是life versus even more life.
當你離開這個世界.

$^{721}$與基督同在的時候.
even more life.
I can't lose.
因為我在基督裡面.
聖經告訴我們.
我們到主再來的時候.
就有一個這麼榮耀的盼望.
這個榮耀盼望不是我們死等.
等去那裡.
而是將這個盼望帶到我們中間.
保羅為什麼要講這些.
這種終末的教導.
就是因為我們會面對孤獨的時候.
會面對憂傷難過的時候.
但是保羅說什麼.
我不願意你們不知道.
免得你們傷心難過.
將會有死亡的人.
Yes我們當然會傷心難過.
是應該的.
但是不會沒有死亡.
現在我們可能有些被困的時候.
但不應該是到.
減弱我們對天國的使命感.
因為這個王者越來越近.
我們應該把握我們的時間為主打平.
縱使被囚的保羅他做什麼.
寫信.
寄信出去.
去勉勵教會.
去為他們每天晝夜的去禱告.
為主繼續打平.
或者將來我們面對一些不確定性.
你看看如果最終的終點是一個好得無比的人.
一個盼望.
好得無比的人.
我們今天.
豈不是生活應該更加有力.
我們唱最後一首詩歌.
讓這首詩歌的歌詞帶領我們.

$^{761}$去認識終末的盼望.
因為終末的盼望和我們息息相關.
我們請敬拜陀上來.
然後我們先結束祈禱.
《終末的盼望》.
作詞:陳小霞 作曲:陳小霞.
誰發現美麗宇宙同時光明共創.
誰歇斯西斯懷疑秘密的亂莫屬.
問最終的答案.
我真的不知道.
誰正預計著世事以後變幻結果.
神已畫過在往後每事兼局如風.
當記錄已照你話語.
從天邊降臨家著滿座.
神啊 到哪一天.
這世界最終必須結束.
你已愛在生生國土.
是你的國度.
神啊 世上需要你的我應該怎樣數.
我信你有天別再分.
是我的盼望.
到那天天地不再.
世間光輝不再複來.
我的心一再仰望你.
到世界的終結.
我的心一再仰望你.
到永遠.
主愛多謝你 因為你賜給我們一個復活終末的盼望.
一個榮耀的盼望 就是我們的主耶穌基督.
為我們的罪而死 死在十字架上.
第三天復活 並且現在在天上.
升天在神的右邊掌權的主宰.
祂一一的報敬.
我們這位的王者.
愛我們這位的君王.
一步一步的接近我們.
主讓這個盼望成為我們今天的鼓勵.
我們憂傷難過的時候.
讓我們知道我們不需要沒有指望.
反而我們更加篤信這個指望.

$^{801}$讓這個指望成為我們今天向前的方向.
我們能夠今天繼續去事奉你.
主叫我們在生命裡面有不同的難關.
有不同的困境.
我們知道這個生命.
我們的人生不是一個沒完沒了的無意義的循環.
一切都會向著這個榮耀的君王在內.
這個警方來推進.
主教導我們來相信.
教導我們這個永恆的盼望.
成為我們今天的鼓勵.
讓我們孤獨的時候.
我們知道有主與我們同在.
並且我們將來也會與主同在.
永永遠遠.
當我們有困難的時候.
我們知道你是我們的幫助.
叫我們今天帶著你給我們的使命.
我們事奉你.
將來有一天你再來的時候.
在空中與你相會.
見到你自己得勝.
聽見你的海歌.
在我們生命裡面能夠奏響.
主求你自己得勝.
你而到我們的人生.
教導我們仰望去信靠你.
祈禱奉主名.
願上帝祝福各位.
願上帝的盼望帶著你們一路離開.
下個星期希望能夠與你們一起事奉.
\newpage



\section{22:1-14}
\label{sec:d5kiiPL31Z0}
\textbf{New Heaven and New Earth: The Coming Banquet (Matthew 22\_1-14) - Rev. Mark Chan}
\newline
\newline
連結: \href{https://youtube.com/watch?v=d5kiiPL31Z0}{\texttt{https://youtube.com/watch?v=d5kiiPL31Z0}} ~~~~ 語音日期: 2023-09-05
\newline
\newline
\hyperref[sec:OzeSIq8OBsY]{\small{< < < PREV SERMON < < <}}
~
\hyperref[sec:index_chronic]{\small{[返順時目]}}
~
\hyperref[sec:index_scriptual]{\small{[返順卷目]}}
~
\hyperref[sec:ovpilkoc6io]{\small{> > > NEXT SERMON > > >}}
\newline
\newline
$^{1}$Pastor Mark for our spiritual today. Hey guys, good morning. Uh, one wish it has to one another.
some of the neighbors, uh, right next to you. Um, good to be, uh, together as the people of God to.
worship together. And, um, yeah, today, uh, we're going to, uh, talk about this topic. Uh, we're.
going to continue this series on new, new heaven and new earth. Try to picture what heaven looks.
like. Try to picture what heaven is like. Right. And, um, there were a lot of different images.
found in the Bible, especially in the new Testament that talks about what our future home.
looks like. I will like the, uh, the weekly devotional is tied to our preaching series..
Uh, it's on our future home. Uh, last time we were talking about John 14 and, uh, Jesus was saying,.
I'm going to, uh, prepare homes for you and a place for you in heaven, in my father's house..
And we're going to continue this topic because the Bible has so much to say about our future home..
In the old Testament, in the book of Isaiah, it talks about the future kingdom, new cosmos,.
new creation is going to be without pain and without tears, without crying, without all this.
suffering. And it will be a righteousness dwelling in the new heaven and new earth..
That's the book of Isaiah. And then in the book of Revelation, the last book of the Bible also.
talks about the new heaven and new earth. And we'll get to that in the future. But today we're.
going to understand one of the images of new heaven and new earth that Jesus gave. Uh, it is.
about wedding, a wedding banquet. All right. Um, I'm sure you have been to a wedding banquet, right?.
Haven't you? To a wedding banquet, to a wedding feast, uh, go to a ceremony and all that. And.
Jesus has something to say about it. So there are three points I want to cover today. It's the.
importance of the banquet in Jewish wedding, especially. And the second point I have is the.
invitation to the, to the wedding banquet. And then the third one, how do we prepare to be part of.
the wedding banquet in heaven? The first thing is the background. We need to understand how.
important this for the Jewish people in Jesus' days, how they see, uh, the Jewish, uh, wedding,.
you know, uh, the entire, uh, Jewish wedding and marriage custom can be understood in three stages..
The first stage is, uh, what we call the arrangement. Now in the, uh, in the old days,.
in, in the first century, uh, people would have their marriage arranged when they were about,.
for guys, uh, about 17, 18 years old. They're going to build a family together. And girls.
probably would get married about 14, 15 earliest. So it's pretty, pretty, uh, early compared to our.
modern, modern lifestyle. Uh, well, when you are 17, 18 year old, you know, you're building a family,.
you're an adult, right? Today is kind of different, right? You know, men who are probably 50 years.
old, they're still adolescent. Uh, they're still teenagers. Uh, today is different, but, uh, in.
the first century, uh, before you're 20 years old, you're, you're, you're going to get married and,.
uh, you're going to build a new family. And that is what the, uh, parents of both sides,.
they will come together and do some arrangement. And the arrangement is very, very serious because.
in the first century in the Jewish society, marriage is a transaction. It's an economic,.
financial transaction, especially for those who are, uh, from the rich families, uh, those who.
are from the priests or the sacred families. Uh, they will, uh, have a lot of terms to lay out,.
especially for both sides of the parents. They will come together and talk about money. Now, uh,.
today, how much do you spend average on a wedding? How much spend on a wedding?.
You remember? How much do you spend on your wedding? Today, probably, uh, I think the low,.

$^{41}$the low end is about maybe 50,000. No? 50,000. Yeah. Uh, something like that. Right? Well, the,.
the, the sky is a ceiling. Uh, it's a lot of money. So in the old days as well, when they talk.
about marriage, when they talk about wedding, they need to talk about finance and usually, uh,.
because it was in the patriarchal society, usually you think the man's side going to pay for it or.
the, uh, the, uh, the bride's side going to pay for it. Who's going to pay for it in Jesus day?.
Usually the man's side going to pay for all the major expenses. Right? Now, if, if, if I live in.
that kind of society, I better not have my son get married. Uh, my daughter is okay. Yeah. Um,.
so this is the arrangement side. And after they lay out all the terms and then, um, they were.
going to the second stage, be troval, which is kind of like today's engagement. But in those.
days it was a little bit more serious. You remember Joseph and Mary, Jesus, father, mother,.
Irving, father, mother. Yeah. They were in, in this stage, they were engaged and legally speaking,.
they were already husband and wife. So, uh, so it was kind of tricky when Mary report to Joseph,.
Hey, I got pregnant by the Holy spirit. Uh, it's going to big bomb for, for Joseph. Right? So he.
has to go to the court, uh, probably some local court to get some certification to a divorce..
Right. But he didn't, right. He didn't. Anyways. Uh, so this, this is the second stage and it's.
kind of like today's, uh, engagement and from, uh, our, our first part of this series, I talk,.
I talk about the farewell discourse of Jesus. Remember the book of John chapter 14 to 17,.
there was long passage talking about, uh, recording Jesus' speeches to the disciples. He said,.
I'm going to leave you now, but I'm not going to abandon you because I'm going back to the.
father's house to prepare a home for you. And this is what the, uh, the, the groom will do.
in this stage after they were engaged. After they enter into the second stage of betrothal,.
the groom, the man will, will promise the bride, I'm returning to my father's house and by my.
father's house, I will build another house. I will prepare a place for you. And, but I will come back.
usually one year later to take you. I will come back for you. This is the background why Jesus.
was talking, uh, in such terms to his disciples. I'm going to prepare a place in the father's.
house for you. There were many places for you, but I'm not going to abandon you. I promise you,.
I will come back and we'll take you. And the, and the groom also beside, you know, preparing a place.
in the father's house is to make some deposit, usually in some valuables, maybe a wedding ring,.
maybe, maybe there were some properties, uh, some commodities that would, uh, give the bride.
some sort of promise, some sort of security. Um, so, uh, not only building the father's house,.
but also, you know, deposit, depositing some gifts as a security deposit, so-called. Um,.
you remember the farewell discourse. Actually, that's what Jesus did. He said, I'm going away,.
return to my father's house to prepare a place for you. And also there was some other sermons,.
some other things that Jesus said. If you check one of the key words.
occurring in the farewell discourse in those chapters, one of the key words always pop up.
is giving. I'm going to give you this, and I'm going to give you that. For example, he said,.
I'm going to give you the Holy Spirit and the Holy Spirit, the counselor will be with you..
So you will not feel abandoned. God is with you. The Spirit of God will be with you..
In chapter 14, Jesus said, I leave you my peace, peace that I grant and I give to you..
In this world, you shall have tribulation, remember, but you will have peace..

$^{81}$And that is the background of Jesus' farewell discourse. That is the second stage..
And the third stage of the wedding custom is about one year later, the house is ready,.
and the groom will go to the bride's village and to take the bride. And of course, in this stage,.
wedding, wedding ceremony, and then of course, come to the conclusion in the banquet,.
the focus is on the couple, right? Focus is always on the couple. Actually today as well,.
you know, some weeks ago, I was officiating a wedding for a couple, a lovely couple. And.
we were about to start because everybody gets so, you know, so, so, so,.
they want to make sure their wedding, wedding service and wedding program goes well, right?.
So everyone is so tight. Everyone is talking about, okay, we have to do this,.
we don't get it wrong and make sure the order of service is correct, right? And I was,.
I was staying there and a couple of photographers, I'm sure they're not very experienced,.
new young people doing photography for the wedding. And they come to me and say,.
because I'm the officiant, right? And the girl said, "Hey Pastor, your shirt has some wrinkles..
That is not right. Yeah, can you cover it or something like that?" Right? And you know,.
as an experienced pastor, you know, I officiate a lot of weddings. I, you know, I gently remind her,.
I said, "Nobody would pay attention to me because everyone's attention is on the couple, right?".
So lo and behold, after the wedding, you know, I got all the photos, right? I got all the pictures,.
very beautiful pictures on the bride and the groom. None of the picture has me on it..
Nobody actually pay attention to the pastor, you know, but that's okay. I officiate it..
So everyone focus is on the couple. It was so in Jesus' day, but of course the,.
when the groom going to take the bride to his home, there was a big processional going on,.
right? All the people and all the bridesmaid and all the groomsmen, all the relatives and friends,.
they were in the processional. And when the bride and the groom about to enter, return to home,.
in their hometown, in their home village, they will have an other team of bridesmaid coming out.
to meet them, right? Remember the parable of 10 virgins? That's what happened. The 10 virgins,.
they prepare oil and they oil them and they were about to receive the new couple into their home.
and they will enter into the big banquet hall and there will be a big celebration..
And in this stage, the last stage, probably, the banquet stage and the host family will prepare.
all the food and all the wine and all of the items that is needed to give honor to the new couple..
That is the background of John 2. Remember the miracle of Jesus turning water into wine,.
right? That is not some magical thing that Jesus did. When they were running out of wine,.
in those days, as in today, in those days, it was a big social shame. And Jesus did that miracle,.
turning water into wine is to save them from social shame. He displayed His glory,.
at least in the sense that would give the host family the social honor that is needed..
But remember that miracle is only the first miracle recorded in the gospel of John,.
right? First of the seven. And remember what John says? Something along that line of,.
you know, this is the first miracle, but it is a pointer to something more glorious coming..
And the gospel of John is the second part. It's the most glorious of all miracles..
It's Jesus dying on the cross, feeling no shame, but glory, glory and honor from God,.
because He loves us so much. He was willing to die. And that is the most glorious moment.

$^{121}$in the gospel of John. But you remember the, if you keep that in mind, the cultural background,.
the context, you understand what is going on, right? It's not an isolated story..
For those who are invited to the banquet, they had already RSVP. They were already invited..
They got the invitation. And as I said, the host family will prepare everything that is ready.
for all the guests to come. And of course, if we enter into the banquet, if we go to the banquet,.
what do we need to do? What do you need to do when you go to a wedding banquet?.
You need to bring some money, right? You better bring some money, right? Some gift money, right?.
Because this is a social honor, right? You don't go there empty-handed, empty stomach,.
and then you devour the food and then you leave. No, that's not how you, that's not how it goes,.
right? You need to bring some gifts, some monetary gift so that you can honor the host family, the.
couple. Without doing that, it will be a dishonor. It will be a shame. And of course, on top of that,.
if you are invited and you RSVP and you miss the meeting, you did not show up, this is a shame to.
the host family. It's pretty much what we experience today, right? But with a, you know, a bigger.
magnitude in those days, in Jesus' day. So the wedding banquet is very, very important. That is.
the conclusion of the entire custom. You got to do a big banquet at the end. Because this is so.
important in Jesus' day in their culture. That's why Jesus picked on some of the images to describe.
the future kingdom, the future heavenly home. When we go there, when we're looking forward to go.
there, Jesus said, it's kind of like, very much like a heavenly banquet. I'm sure you've been to.
a wedding banquet. Multiply it infinitely. That is the glory of the future heavenly banquet..
And Jesus used this image or the New Testament Bible uses many of this wedding banquet images.
to describe our future home. And one of the passages today, we're going to read that quickly,.
is talking about the invitation to the wedding banquet. Matthew 22. Let me read that..
Jesus also told them other parables. He said, "The kingdom of heaven can be illustrated.
by a story. A story of a king who prepared a great wedding feast for his son..
When the banquet was ready, he sent out his servants to notify those who were invited,.
but they all refused to come. So he sent other servants to tell them, "The feast has been.
prepared. The bulls and the feline cattle has been killed. They are ready. Everything is ready. Come.
to the banquet." Now what is going on here? This story parable talking about the king preparing a.
banquet for his people. Those who are invited, they should have the royal honor to attend. And.
they have already RSVP, right? They reply to the email or click the link, you know, "I'm going to.
come." Right? "Come in." And all that. But it's very strange. These sort of people, they are not.
willing to come. They are not willing to attend. Now in those days, in the ancient days, there was.
no refrigerator. Once you cook something, you got to consume it. Right? You have to have it..
You got to finish it. Because you can't put it back into the fridge and then bring it out again.
next day or two and then put it in the microwave. No, you can't do that in the ancient Palestinian.
society. So you have to consume it. That's why the king prepare everything for the banquet and.
quickly send out his servants to the nearby neighborhood and say, "Hey, come on. Guys,.
this is ready. And you need to come right now." Usually it happen in the evening, at night..
Now pay attention to who are being invited here. Who have RSVP here?.
It says those who are already invited. To notify those who were invited, RSVP, right?.

$^{161}$It's very strange. They promised to come, but they never come. They're not willing to come..
What are the excuses do they have?.
The story continues. It just said, but the guest, he had invited, he had RSVP,.
ignored them and went their own way. One to his farm, another to his business. Other people,.
other guests, seized the servants, the messengers, and insulted them and killed them..
So the king sent out another team of servants and trying to urge them, "Come now. This is the.
moment. You accept the invitation. You need to show up." The excuses are two. One's going to the.
field, to farm. The other one's going to do his business or do some trading. And these two excuses.
are lame excuses. You know why? Because I told you some of the background. The banquet happens.
at what time of the day? In the evening. At night, right? You never go to, even today, you never go.
to a wedding banquet in the morning, right? Usually it's at night. At night, there's no opportunity.
that you have to go to your farm, right? There's no way you can do some trading, right? We work.
nine to five. It's after five, right? There's no way you can do all these things and excuse yourself.
from attending the royal banquet. The excuses reveals their idols, and the idols is in their.
heart. Essentially what they're saying is, "The banquet is not important. It's not worth it..
I'm going back to my work. I'm going back to my farm." You understand when we are confronted.
with the heavenly invitation, when the invitation is sent to you by God, usually our excuses.
reveal our priority. What's most important? What is worthy in your mind? It is always revealed when.
we try to respond to the invitation of the kingdom. When the gospel is presented to you,.
you can think of a lot of different ways, excuses, or reasons to reject or refuse it..
But you need to understand in the Bible, those excuses, those reasons, are usually indicators.
of what is truly your God. What is God to you? What is most important to you? Now I admit to.
all of you, myself included, we are all idols, manufacturers. Remember, I think 400 years ago,.
I think 400 years ago, 400 or 500 years ago, there was a reformed theologian called John Calvin..
John Calvin describes human nature, human beings, as idol manufacturers..
We manufacture our own idols. We create our own idols. And we put idols on our exotic place,.
and we worship it. And we spend all of our time consuming our energy. And that is your God. That.
is your idol. When I prepared this sermon, I was reminded of some of the books that I read from.
an atheist called John Paul Sartre. John Paul Sartre was a French existentialist..
Now what does it mean? Existentialism, simply put, means I want to do whatever I want..
My freedom is not, there's no boundary. There's no accountability in my freedom. That is the true.
freedom that I have. And John Paul Sartre was saying that. And many of his novels talk about.
the death of God. That God is dead. You know, we don't have to regard God as anything, anyone.
worthy in our life. You just erase it from your life. You know, reject him from your life. And.
John Paul Sartre said, "Belief in God makes my morality impossible." My morality of existentialism..
I want to be free. I don't want to be kept accountable. If God exists, He will keep me.
accountable. That's why if I believe in God, it will make my morality impossible, Sartre said..
To which a Dutch theologian replied, "Sartre's morality makes the non-existence of God a necessity.".
It is a necessity for an atheist, for an existentialist, that God does not exist..
It's not because it has been proven that God does not exist. That is not the conclusion,.

$^{201}$not the last part of the equation. That the non-existence of God, that God does not exist,.
God is absent, is my presupposition, not my conclusion. It is my worldview. It's my own.
belief that will etch out God. That is a necessity. I have to etch out God so that I can have my own.
life. That is what atheism in his version. Now think about it, this is not just philosophy..
That is essentially what Jesus said in the parable. Those people who have RSVP, they reject.
to attend the banquet. It's not because they cannot attend the banquet. Essentially they will not..
It's not because we cannot come to the kingdom. Standing in front of the foot of the cross,.
at the foot of the cross, it's not because we cannot accept the invitation..
The invitation of love and grace is not the ability. It is our will. Whether you will come.
or not. Now it's explained in the first part. In the ancient society, in Jewish society in Jesus'.
day, being invited and RSVP and you are willing, you are not willing to show up, that is an insult.
to the host family. Remember Jesus' parable? It's not just one guy, any dude trying to create or.
set up a banquet. It's not a certain dude on the street. It's the king. The king is preparing a.
banquet for his son. Because his son is getting married. You follow that line, you understand.
what the big deal about this banquet. It's the king's son's marriage, wedding banquet..
Now the son is getting married, that means the son will inherit the kingdom. He will have the.
kingship. The son is getting married, he will be the king of the kingdom. Not showing up.
does not mean social shame alone. It means treason. It means rejecting the kingship..
That is treason. 22.7. The parable goes on. The king was furious and he sent out his army.
to destroy the murderers and burn their town. And he said to his servants, the wedding feast is ready.
and the guests I have invited are not worthy of the honor. So go out to the street corners and.
invite everyone you see. So the servant did that and they find everyone they could find, good or.
bad alike and the banquet hall was filled with all the people. Now why Jesus told this parable at this.
point? If you understand context, as we read the Bible, we have to understand context, right? It's.
not to isolate the story. This is chapter 22. What happened in the previous chapter, chapter 21?.
Jesus riding on the donkey entered Jerusalem as a humble king, as the loving king. Matthew 21..
He is the king. He is riding the donkey entering the holy city. He is the true king. So he went.
into Jerusalem and he cleansed the temple and there were so many debates and controversies between.
Jesus and the authorities of Jerusalem, the Jewish people. This parable is ironic because this.
parable is trying to condense the entire Old Testament story into 10 verses. What happened.
in the entire Old Testament story, Jesus said, is that the king has already, has been inviting.
his people to come to his feast to enjoy the time together, to enjoy his food. But his people have.
been refusing to come and kill all the servants and those of the Old Testament prophets,.
from the first one to the last one, John the Baptist, they killed them all because they are.
not willing to honor the son's kingship. And finally they killed the son and crucified him.
in Jerusalem. Now this is the last week of Jesus' early earthly life. This parable is trying to.
tell us the entire Old Testament story and it ends really bad because in AD 70,.
the Romans did come and they came to destroy the temple, Jerusalem. And the temple crumbles.
and the people scattered. Jesus is trying to tell us the entire Old Testament story, God and his.
people. And because they refused, because they reject the gospel invitation, now the door is.

$^{241}$open to other people. Who are they? The Gentiles, right? The Gentiles would come..
All the nations will come because of Jesus, because of the grace..
Then they come to the banquet. That is the invitation. If you reject it,.
not because you can or cannot accept it, because you will not accept it,.
the consequences are there. AD 70. That's what happened to the Jewish people..
Now the parable could have ended at this point, verse 10. But Jesus continues. Now this is very.
interesting because you find similar parables from other books of gospels. But Jesus put on.
a twist here. He continued with verse 11 and he talks about how do we prepare to go into the.
banquet, to be part of this banquet. Verses 11 and following. But when the king came in to meet.
the guest, he noticed a man, there was a dude who was not wearing the proper clothing for the wedding..
He asked, "Friend, how is it that you are here without the wedding clothes?" But the man had no.
reply. Then the king said to his servants, "Bind his hands and feet and throw him into the outer.
darkness where there will be weeping and gnashing of teeth. For many are called, but few are chosen.".
Now I've talked about the wedding banquet. Usually it lasts for almost seven days. Can you believe.
that? If you have a wedding banquet in the Jewish society on Jesus' day, you have to host a banquet.
seven days. This is crazy. I attended one of my friend's wedding. The groom is from the Vietnam.
family background. They did host for 10 days. It's quite a bit of food. A lot of guests coming in.
and out, in and out. But when you attend a banquet, what do you need to do in terms of preparation?.
At least you need to dress up a little bit, right? You can't just go with a fit flop, right?.
That would be a dishonor to the guests. It's the same way in Jesus' day. They need to dress up..
Now there are different interpretations of why this dude without the proper clothing..
Perhaps one of the most convincing arguments is that the point is not how he ended up this.
way. The point is what his clothing, the proper wedding clothing symbolize. What does it mean?.
Now you flip the pages of the New Testament. You'll find proper clothing to enter the kingdom.
is the robe of righteousness, the white robe that is provided by Jesus. He died on the cross. He.
sacrificed himself. He redeemed you. And so to speak, you have a new cloth on you. God is not.
looking at you. If you are in Christ, if you believe the gospel, God is no longer looking at.
you, you know, in your filthy way. You're not a sinner anymore. God look at you just as he looks.
at his own son, his perfect holy son, Jesus Christ. That is what we call justification by faith, right?.
Justify fully righteous standing in front of God. That is probably what the wedding.
cloth is all about. In any case, this parable add a twist at the end that you have believed in Christ..
You supposed to be sitting in the banquet, but you are not with proper clothing..
That means your relationship with God is a little bit shaky..
Maybe you really don't know Christ. You have not really accepted him as your Lord and your Savior..
Now there are phrases like outer darkness, you know, gnashing and weeping..
Now, what does it mean? If you turn to other passages in the gospel of Matthew, you would.
not be surprised. Same phrases are being used as so often in the context of going into entering to.
future kingdom. And you find those places, you understand Jesus is saying,.
I don't really know you. Who are you? You go out to the outer darkness..
Gnashing of teeth, right? Weeping. You are not going into my banquet because you do not know me..

$^{281}$You have not really put in your faith in me. Now this parable ends with this twist..
It's a great reminder for you and me that we come to the church and we even think that.
we are sons and daughters of the King. Are we really? Right? Are we really?.
In the gospel of Matthew, it talks about believing in Christ is much more than lip service.
in the Word of God. To listen and to walk, right? Of course, God is not looking for a perfect man or.
woman. We can never be perfect in this life. In this side of heaven, we cannot be fully sanctified,.
100\% sanctified, right? We are clothed by Christ's righteousness. That's why we can enter.
into kingdom. But foundationally, in the bottom of your heart, you have this inclination,.
the willingness to say to God, yes, you are my Lord. I come before you constantly, regularly..
I want to turn my life to you. I want to obey you. Yeah, help me. Help me. That is your most.
fundamental change when you come to Christ. And that is the robe of righteousness. That is what.
this parable is all about. You see in the banquet, you better be making sure that you are part of the.
honored guest, properly clothed. The future kingdom is like a banquet. If we refuse the invitation,.
it's essential, it's essentially, it's like you can try it. Stand at the foot of the cross.
and you look at the cross. Jesus sacrificed himself and shed his blood for you..
That you are unworthy. You don't deserve it. But you stand in front of the foot of the cross..
You're saying to the Jesus who died for you on the cross, you said you're not worth it..
Come on. You know, I don't care. This is essentially what the refusal of the kingdom.
is like. How can we turn away such a great salvation, such a great, wonderful grace, right?.
The parable reminds us we better be sure right now, right? It has been said that.
there will be three wonders when you go to heaven. Three wonders when you step into heaven,.
you know, the new world, the new heaven and new earth. The first wonder is that.
you will see many in the kingdom whom you did not expect to see. You step into the kingdom, hey,.
why are they here? I didn't expect that. But they were here indeed because of the grace of God..
That's the first wonder. The second wonder is that you do not find many whom you expect to see..
I thought they should be here, but they're not here. Why? Because essentially they rejected the.
king. That's the second wonder. The third wonder is the greatest wonder of all. When I step into.
the kingdom, the greatest wonder should be I find myself there with my robe of righteousness,.
credited to me because of what Christ has done for me and knowing I don't deserve..
And that is the greatest wonder. And friends, I want to end this note just as Jesus ended..
The invitation is out there. You are our SVP. You are our SVP by your faith in Christ..
Make Jesus Christ your Lord today. Make him your leader, the leader of your life..
Bow before him and say, "My heart is right here. Please be my Lord. Guide me and help me listen to.
you. Help me obey you." We're going to end today with a chameleon. We'll do it once a month..
And we'll put it right at the end today. And there's a long tradition in the church that.
chameleon is a foretaste of the future heavenly banquet. And let us pray and prepare our hearts..
This is a foretaste. It's like an appetizer. The future kingdom, the banquet, the perfect,.
glorious banquet is ready for us. Take it. Let's all pray. Lord Jesus, we ask that you will make us.
one as your church until the final victory of the second coming of our Lord Jesus Christ..
We want to give you thanks because we are invited. We can be at your table. We can feast at your.

$^{321}$banquet table. All honor and glory to you. And once again, we pray that.
that Jesus Christ, our Lord, will be the Lord of all. He is the one.
who has the lordship over everything in our lives..
We're going to school, but our school life is under your lordship. Please help us to do so,.
to glorify you. Our family, our marriage, our church, our ministry, everything, our career..
You are the Lord. You are the Lord of all. And I pray that today we'll accept this invitation.
to taste, to foretaste of your glorious, heavenly banquet. Thank you for the invitation..
And we're going to accept it with our will. In faith, we put our trust in our dependence.
on our Lord Jesus Christ. All glory to you in Christ's name we pray. Amen..
We have a communion table back there. If you're ready to take this foretaste, the appetizer.
of the future banquet, you can go ahead and go to the table and get the packet..
And we're going to partake together..
So,.
so,.
but this is the foretaste of the kingdom, the kingdom banquet. Let's partake the bread together..
And the cup as well..
Let's pray a final time. Father, once again, we give you thanks because of the great love that.
you have for us. We are standing at the foot of the cross and we are amazed and wonder by your love..
Thank you for accepting us. Thank you for, for, for your righteousness, that we are fully justified.
because of what your son has done for us. And thank you for this hope of future banquet..
And I pray that our life will turn to you, will obey you, will listen to your word..
In Christ name we pray. Well, God bless you all and grant you peace in this Labor Day weekend..
Right? Yeah. Thank you for coming. And I hope to see you next week. And we're going to have some.
normal Bible study. We're going to have some fun as a fellowship. Okay. You're invited,.
right? But you have to pay for it, pay for your own. We're going to barbecue next..
\newpage



\section{馬太福音 22:1-14}
\label{sec:ovpilkoc6io}
\textbf{新天新地:人生終末的再思 4 將來的婚宴 (馬太福音 22\_1-14) - 陳崇基牧師}
\newline
\newline
連結: \href{https://youtube.com/watch?v=ovpilkoc6io}{\texttt{https://youtube.com/watch?v=ovpilkoc6io}} ~~~~ 語音日期: 2023-09-05
\newline
\newline
\hyperref[sec:d5kiiPL31Z0]{\small{< < < PREV SERMON < < <}}
~
\hyperref[sec:index_chronic]{\small{[返順時目]}}
~
\hyperref[sec:index_scriptual]{\small{[返順卷目]}}
~
\hyperref[sec:D2gsbWPv_jE]{\small{> > > NEXT SERMON > > >}}
\newline
\newline
馬太福音 22:1-14
\newline
\begin{longtable}{cl}
\hline
\hline
章節 & 經文 (和合本修訂版)\\
\hline
22:1 & \begin{tabularx}{0.7\textwidth}{X} 耶穌又用比喻對他們說： \end{tabularx} \\ \\ \relax
22:2 & \begin{tabularx}{0.7\textwidth}{X} 「天國好比一個王為他兒子擺設娶親的宴席。 \end{tabularx} \\ \\ \relax
22:3 & \begin{tabularx}{0.7\textwidth}{X} 他打發僕人去，請那些被邀的人來赴宴，他們卻不肯來。 \end{tabularx} \\ \\ \relax
22:4 & \begin{tabularx}{0.7\textwidth}{X} 王又打發別的僕人，說：『你們去告訴那被邀的人，我的宴席已經預備好了，牛和肥畜已經宰了，各樣都齊備，請你們來赴宴。』 \end{tabularx} \\ \\ \relax
22:5 & \begin{tabularx}{0.7\textwidth}{X} 那些人不理就走了，一個到自己田裡去，一個做買賣去。 \end{tabularx} \\ \\ \relax
22:6 & \begin{tabularx}{0.7\textwidth}{X} 其餘的抓住僕人，凌辱他們，把他們殺了。 \end{tabularx} \\ \\ \relax
22:7 & \begin{tabularx}{0.7\textwidth}{X} 王就大怒，發兵除滅那些兇手，燒燬他們的城。 \end{tabularx} \\ \\ \relax
22:8 & \begin{tabularx}{0.7\textwidth}{X} 於是王對僕人說：『喜宴已經齊備，只是所邀的人不配。 \end{tabularx} \\ \\ \relax
22:9 & \begin{tabularx}{0.7\textwidth}{X} 所以你們要往岔路口上去，凡遇見的，都邀來赴宴。』 \end{tabularx} \\ \\ \relax
22:10 & \begin{tabularx}{0.7\textwidth}{X} 那些僕人就出去，到大路上，凡遇見的，不論善惡都招聚了來，宴席上就坐滿了客人。 \end{tabularx} \\ \\ \relax
22:11 & \begin{tabularx}{0.7\textwidth}{X} 王進來見賓客，看到那裡有一個沒有穿禮服的， \end{tabularx} \\ \\ \relax
22:12 & \begin{tabularx}{0.7\textwidth}{X} 就對他說：『朋友，你到這裡來怎麼不穿禮服呢？』那人無言可答。 \end{tabularx} \\ \\ \relax
22:13 & \begin{tabularx}{0.7\textwidth}{X} 於是王對侍從說：『捆起他的手腳，把他扔在外邊的黑暗裡；在那裡他要哀哭切齒了。』 \end{tabularx} \\ \\ \relax
22:14 & \begin{tabularx}{0.7\textwidth}{X} 因為被召的人多，選上的人少。」 \end{tabularx} \\ \\
[1ex]
\hline
\hline
\end{longtable}
$^{1}$許多元主早安.
我們一起低頭祈禱.
讓我們今天來到.
敬拜去事奉你.
我們感恩因為你.
給我們恩典在基督裡面.
能夠成為你自己的子民.
你住在我們生命當中.
你在這個地上繼續建立天國.
一個超越死亡.
的一個盼望.
主讓我們繼續仰望你.
在我們生命當中.
在教會裡面有憂傷難過的弟兄姊妹.
主我們再一次去教導.
求你親自設平安和安慰給他們.
也有一些是有病痛的.
主我們求你親自去看顧和醫治.
在我們的生活當中.
我們也遇見很多不同的挑戰.
我們記起一些弟兄姊妹.
他們生命裡面可能有很多失去.
經歷一些難處.
甚至可能是失業.
失去一些至親.
主我們求你賜給他們力量.
讓我們知道在你的恩典裡面.
我們知道你給我們的力量是足夠的.
今天我們教會.
來到你面前我們求你.
在上帝的話語當中.
你自己對我們作提醒.
求聖靈來引領我們.
祈禱逢主名.
好那我們繼續去看這個主題.
這次我們第幾講.
第四好像是.
第四.
我們講到《馬太福音》第22章.
我們講到將來的天價.

$^{41}$將來的國度.
上帝的國度.
已經展開了在基督裡面.
在我們生命裡面展開了.
但仍然是.
未完全實現.
將來一定會實現.
這個就是耶穌基督給我們的盼望.
因為這個盼望的緣故.
所以我們今天能夠繼續向前.
講到我們人生的終末的時候.
有時候是我們的目的地.
來決定我們今天的路向.
(聲音不清楚).
(笑聲).
可能聽到.
AI聽到我說話.
anyway.
well anyways.
那個將來的天國的盼望.
去到決定我們今天.
行走的方向.
和我們今天所做的決定.
回到聖經.
舊約新約.
你都會看到.
這個講到天價.
講到天國.
在舊約裡面講到天國.
上帝會抹去人的眼淚.
會不會有痛苦.
他所答應.
給我們的國度一定會臨到.
在聖經最後.
也都講到將來的天國.
同樣用回這樣的字眼.
來形容.
在天國裡面.
有所有的萬國萬民.
能夠進入上帝自己的國度.

$^{81}$在榮耀的裡面.
去到和上帝一起同在.
有一個畫面.
在聖經裡面經常會用的.
去形容天國.
就是我們日常生活裡面.
我們經常會接觸的一些畫面.
就是婚禮.
婚姻.
將來的天國.
好像一個婚禮的筵席一樣.
羔羊的婚宴.
所以今天我們去看一個比喻.
我們去看一下.
聖經裡面怎樣講這個天國.
還有怎樣準備我們自己.
進入這個天國.
有三點.
第一樣我想解釋一下.
當時耶穌的時候.
那個婚姻.
那個婚禮.
結婚是一回什麼事.
為什麼這麼重要.
既然這麼重要.
所以為什麼主耶穌去到教導將來的天國的時候.
都是用一些這樣的畫面.
這樣的意象.
然後我們講到那個邀請.
邀請我們去喝酒.
去到婚禮.
然後.
我們現在可以做什麼的準備.
三點.
第一樣我們先解釋一下背景.
當時的猶太人.
耶穌時代的猶太人.
他們談婚論嫁大概可以分成三個階段.
第一個階段就是.
談婚.

$^{121}$去到做一些.
做一些安排.
之前講過.
很多時候都是兩個家庭.
一起去談.
一對男女.
到底是對不對.
是一些人生的大事.
大到一個地步.
有些拉比的律法.
說就算安息日的時候.
都可以談婚論嫁.
安息日的時候.
有很多事情不能做.
但是.
談婚論嫁仍然是可以的.
所以有什麼事情還沒講完.
安息日那天都可以.
去講一下嫁娶的事情.
當然他們當時最關心的是什麼.
就是門當戶對.
你知道什麼叫門當戶對嗎.
就是兩家人走在一起.
男女到底是不是合得來.
是不是.
不只是他們兩個合不合得來.
而是兩個家族到底是不是配.
特別是一些大戶人家.
或者是有高階級的人.
譬如說祭司等等.
利未人等等.
他們很多時候都要很關心.
就是我們這樣結婚.
到底合不合得來.
所以很多時候不只是男女的抉擇.
是兩個家庭.
甚至是家族.
來談的一件事.
當時是父權社會.
一定是重男輕女.

$^{161}$不代表聖經是重男輕女.
但當時的社會是重男輕女.
所以婚姻是這麼大件事的時候.
你猜婚事出的錢是誰給的.
男家還是女家.
是男家給的.
猶太人.
當時的猶太人.
男家給的.
所以大部分最大的開支.
錢都是新郎那邊.
是父家去給錢的.
所以講完這些東西.
才可以講下一步.
你知道我有兒子有女兒.
我覺得不錯.
如果我兒子結婚.
應該用美國那一套.
應該是女生給的.
如果我女兒結婚.
就應該用聖經那一套.
所以.
Anyway.
過了談婚之後.
就到訂婚.
什麼意思呢.
就好像今天的 Engagement 訂婚.
但當時就更加嚴肅.
當然.
因為當時訂婚了.
就一定是已經有法律的責任.
名份上就已經是夫婦.
你記得約瑟瑪利亞就是這樣.
訂婚之後.
兒子.
曾經講過.
會回到父家那裡預備地方.
但同樣他會做多一件事.
就是去做一些所謂 Deposit.
或者一些禮品.

$^{201}$你可以看.
禮品.
你可以叫他做訂婚的信物.
無論是戒指也好.
或者一些首飾也好.
總之是一些 Property.
一些.
有價值的東西.
他就會放在女家那裡.
然後就說一年之後就會回來接你.
並且我已經預備了地方.
我們講過耶穌說.
他離開.
去到父家那裡預備地方.
就是一個這樣的意境.
不單止是這樣.
你記得耶穌在臨別的時候.
臨別的講話.
Farewell Discourse.
好幾張13至17張約翰福音.
不單止是講到.
我預備地方.
也講到我給你很多東西.
如果你留意約翰福音第13至17張.
那裡整段的講話.
最後晚餐耶穌說什麼.
我去父家那裡.
為你預備地方.
一定會回來接你.
並且我會賜給你.
很多東西.
第一樣他說賜給的.
就是聖靈.
我會將聖靈賜給你們.
聖靈保衛師就與你同在.
一路帶領著你.
這就是一個 Deposit.
還有賜給你什麼呢.
賜給你平安.
世上裡面有苦難.

$^{241}$不過在我裡面有平安.
我留下平安給你們.
我賜平安給你們.
以致你知道在世上裡面.
你怎樣去生活.
這些都是所謂的 Deposit.
一年之後他就回來.
去接他們.
這就是一個訂婚的階段.
第三個階段就是婚禮.
加上婚宴.
婚禮和婚宴.
一年之後.
新郎回去接新娘.
當然新娘有很多準備.
很多打扮.
當時都是一樣很緊張.
很緊張.
特別是新人自己很緊張.
然後就會接新娘.
回到他們的村.
或者回到他們的副家.
這個都是一些.
非常熱鬧的情況.
和非常重視的情況.
前陣子我去證婚.
婚禮.
婚禮很緊張.
比熟人還要緊張.
很緊張.
大家很緊張.
不知道為什麼這麼緊張.
不知道是不是因為做久了.
所以不覺得緊張.
那對男女.
攝影師.
一男一女.
開始的時候就跟我說.
牧師你這樣不行的.
衣服皺皺的.

$^{281}$你應該整好一點.
我就跟他說.
我語重心長的跟他說.
其實沒有人會理我的.
他說你這樣拍照不行的.
我說你聽.
不會拍到我的.
我就跟那對攝影師這樣說.
因為攝影師很緊張.
結果就結束了婚禮.
我又收回照片.
真的被我批准.
我批准了什麼呢.
就是只有一張照片有我的.
並且我不是主角.
我只是在背景.
婚禮就是很大的一件事.
所以他們會.
大排筵席.
招呼.
十個同女的比喻.
將新娘帶回村子.
帶回她的父家.
然後就進入筵席.
筵席也是非常重要的一部分.
不是華人才.
重視吃東西的.
猶太人.
都是一樣的.
特別是當時社會的背景.
大排筵席.
其實很少機會.
有得大魚大肉的.
婚禮或者婚宴.
就是一個最重要的時間.
大家被邀請的人.
就能夠大魚大肉.
主人家一定會預備.
足夠的酒菜.
因為婚宴還沒擺之前.

$^{321}$他就會發出所有邀請的帖.
人家一定要.
預備好酒菜.
不夠就很醜.
那些就叫做.
social shame.
你記得.
加拿大婚宴就是這樣.
耶穌去加拿大婚宴.
記不記得.
瑪利亞也在那裡.
是很樂面的情況.
不夠酒.
酒是最重要的.
今時今日都是.
open bar.
任喝.
大家就很開心.
雖然你未必喝.
你不喝多少.
不過他說open bar就很好.
所以你看回那個故事.
耶穌是.
拯救那個家庭.
以致他們不會被落面.
是給他們honor.
你記得他將水變酒.
但是那個神跡.
告訴你.
神跡告訴你.
耶穌說沒有人知道.
不過有一些榮耀.
我再彰顯的時候.
那些才是真正的榮耀.
不是單單水變酒.
不是social honor那麼簡單.
真正的榮耀是何時來.
就是約翰福音最後.
他上十字架的時候.
那個才是真正的glory.

$^{361}$所以那個是婚宴.
宴席.
通常猶太人的宴席.
是不是很好玩.
今天也不會擺七天.
並且那七天.
大家也可以放假.
主人家一定會放假.
當然親朋戚友就不會.
在七天裡面.
也有一些人會出出入入.
可能有些朋友會去享受.
這些盛宴等等.
被邀請的.
是一個社會的榮譽.
今時今日也是.
人家擺酒請你.
是不是榮譽.
是吧.
你應該去的.
但沒有空手去.
如果不是你就很醜.
你應該要給禮金.
給多少都要.
要送上禮金.
被邀請的.
一定不可以缺席.
因為缺席.
就是不給面子主人家.
這些.
大家社會裡面的.
社交的東西.
大家應該知道.
就算人不到.
禮也要到.
這個就是當時.
第一世紀猶太人.
的.
叫做社會慣例.
社會慣例.

$^{401}$宴席是非常重要.
宴席.
可以這麼說.
是最重要最後最後的.
我記得那次.
去一個人的婚禮.
結束後.
他們去哪裡擺酒.
去Fogo D'Italia擺酒.
那次很開心.
我沒試過去一個婚禮.
吃那個飯堂.
吃到不想走.
無論如何.
是一個非常重要的.
一個禮儀.
正因為這麼重要.
難怪.
難怪耶穌說到.
我們將來最後的天家.
是用婚言.
來形容.
馬太方第22章.
說到.
婚言的.
invitation.
我們先讀.
耶穌用了一個比喻來跟他們說.
天國很像一個王為他兒子.
擺設.
娶親的宴席.
他打發.
僕人去.
請那些被邀請的人.
來赴宴.
他們卻不肯來.
王又打發別的僕人說.
你們去告訴那被邀的人.
我的宴席.
已經準備好了.

$^{441}$牛和肥畜已經宰了.
各樣都齊備了.
你們來赴宴吧.
這裡說的是什麼人擺酒.
一個王為他的兒子.
娶親來擺酒.
大家可以想像得到.
那不是一條村的問題.
這是整個國家.
整個國度.
一個最大的派對.
一個的宴席已經準備好了.
他就派人出去.
以前巴勒斯坦地.
第一世紀的時候沒有雪櫃.
你煮了東西一定要吃.
對嗎.
煮了東西你不能夠保存.
因為沒有那種保存的科技.
沒有那些工具.
所以立刻就要去食用了.
所以他就打發這些僕人出去.
叫他們是時候入席了.
留意這裡是哪裡.
誰人是被邀請去那個宴會.
那已經被邀的人.
就是說叫那些人來的時候.
那些人已經是RSVP了.
那些是VIP來的.
他們答應了會來.
但是那個宴席煮好了的時候.
那些人就趕緊出去.
然後就叫他們來.
但奇怪的是他們就不肯去.
剛才我解釋了那個背景.
你不肯去.
說了去不去.
這些就是不給面子.
今天.
我們都是這樣看婚禮的.

$^{481}$以前是更加重.
今天也是很重.
昨天我和太太聊天.
我說如果有人請你做.
主伴兩.
May of Honor.
你應該怎麼回應.
你是要去的.
除非是唐人新慶.
COVID.
不去就行.
May of Honor.
你想不想做.
人家把邀請給你的時候.
你就應該要去.
對不對.
如果有人請你做.
Best Man.
你是不是應該要去.
不去嗎.
兄弟.
無所謂.
總之有東西吃有酒喝就算了.
那種是Honor.
Social Honor.
社會學家叫做Social Currency.
你一定要去.
答應.
如果你不去的話.
就是來不及去.
如果去的話.
就是要下面.
這班人被邀請的.
他們就是不去.
什麼原因.
第六節.
那些人不理就走了.
一個.
到自己的田裡面.
一個說去做買賣.

$^{521}$其餘的.
抓住僕人凌辱他們.
然後將他們殺了.
再派僕人去.
催逼那些人.
都是缺席的.
到自己的田地.
然後去做買賣.
兩個你一聽到就是很廢的藉口.
為什麼.
因為婚宴是什麼時候發生的.
在一天裡面哪個時間.
通常是晚上.
你都沒有去過婚宴吃早餐.
是晚上.
晚上怎麼會去耕田.
晚上怎麼會去做買賣.
所以聽到這些根本是很廢的藉口.
今天.
等於今天.
天國的邀請.
給了你的時候.
你就應該去到回應.
但很多時候我們都是這樣的.
我們都會.
用一些很廢的藉口.
我們說耶穌很好.
大部分人聽見耶穌的時候.
都是覺得耶穌好.
很少說耶穌.
搞錯了.
這些大道.
教人做一些邪惡的事情.
沒有的.
通常一聽到耶穌的時候.
我們都會說是OK.
是好.
不過我就是不能夠接受.
耶穌在我的生命裡面.
去到說了算.

$^{561}$還要每一方面他要作主.
我就受不了.
所以說耶穌是好.
OK.
不過都不是很行.
因為我不想改變.
因為主來到我的生命裡面.
我就要慢慢.
一路一路.
我所有的每一方面都要被改變.
因為他做了我的主宰.
對嗎.
他上了我的車.
他就是開車的那個.
我不再是開車的那個.
由他去到作主宰.
所以很多時候.
我們的心理都是一樣.
但我想告訴你.
這裡有兩個藉口.
其實耶穌點出我們有時的藉口.
就是我們反映了.
我們心裡面其實有一些偶像.
那些偶像不一定是放在神主位上的那些.
是我們生命裡面一些的idol.
我不能夠花時間去敬拜上帝.
為什麼.
我會做一些更加忙的事情.
原因不是因為你時間少一點.
人家時間多一點.
原因就是我的優次的問題.
就是我心裡面有偶像.
我和你都是一樣.
我們都是一些偶像的製作人.
我們很喜歡製作一些偶像出來.
並且我們的興趣.
我們想做的事情.
很多時候都是一些可能.
表面上是非常光彩的事情.
沒有錯的.

$^{601}$在社會裡面是值得稱頌的.
但是我們就拿出來做偶像.
因為我們是偶像的創作者.
說到底骨子裡面.
其實就是我們不能夠接受.
一位主來到我們心裡面.
要再整頓我們的東西.
我們就作一些藉口.
就等於這裡被邀請的人.
一個說要去田裡面夜晚.
夜晚不行.
我要去買賣.
在我預備的時候.
記得以前看過一本書.
是一個.
這裡講一點哲學.
我知道大家不喜歡哲學.
但是看過一本書.
是一個法國的哲學家寫的.
那個人叫做沙特.
沙特的一生.
就是去反對上帝存在.
他是一個無神論者.
所謂存在主義者.
什麼叫存在主義者.
其實不是很深的東西.
存在主義者就是說.
總之我的自由就是沒有限制的.
我喜歡做什麼就做什麼.
我喜歡怎麼定義我自己.
就怎麼定義.
是沒有任何限制.
沒有任何束縛的.
如果上帝是存在的話.
就會束縛我的自由.
所以沙特寫很多小說.
還有寫很多的作品.
都是講這件事.
就是上帝死了.
其實不是真的上帝死了.

$^{641}$而是當他死了的意思.
因為如果上帝在我生命裡面的話.
我真的不行.
因為他又要拿著.
我又不想給他拿著.
所以沙特曾經講過.
他說.
Believe in God makes my morality impossible.
如果信仰上帝的話.
就會使我的道德生活不可能.
使我整個生命不可能.
Impossible.
所以為什麼上帝是不存在.
於是當時大概二十世紀中的時候.
有一個德國的神學家.
他就這樣回應.
他說沙特的道德.
是使上帝不存在成為必然.
Necessity.
正因為他這樣去定義自己.
所以一定上帝是不能夠存在.
不是上帝存在已經是一個理性的結論.
而是一個前設 一個開始.
上帝是不能夠存在.
以至我自己的生活能夠繼續.
問題不是.
這裡被邀請去赴宴的人.
不是他們能不能夠去.
而是願不願意去.
當然他們當然能夠去.
並且必須要去.
但是他們不願意.
不是Can or Cannot.
是Will or Will not.
我們面對上帝其實都是一樣.
不是You can serve God or not.
而是You will serve God or not.
在近東社會裡面.
被請去不來的.
本身已經是對主人家侮辱.

$^{681}$剛才說的這個比喻.
耶穌說的.
不單只是一個普通人.
是一個王為他的兒子來成親擺酒.
不赴宴不單只是不給面子的問題.
王的兒子擺酒.
結婚擺酒.
冠軍說什麼.
這個王子是一個天國的國度的繼承人.
他是一個新的主宰.
不赴宴不赴會的話.
就等於不承認他的王權.
Does it make sense?.
不是單單不給面子的問題.
是Treason.
是叛國罪.
所以這個王怎麼做呢.
第七節.
王就帶路了.
發兵來除滅那些兇手.
去銷毀他們的城.
於是王對僕人說.
起燕已經齊備了.
只是所邀的人不配.
所以你們要往大路口上.
凡看到的就邀請他們來赴宴.
那些僕人就出去到大路上.
凡遇見的不論善惡都招聚了來.
然後整個的Banquet Hall就滿了人.
這裡說什麼.
如果留意《馬太科學》的上下文.
這裡第22章.
21章發生什麼事.
耶穌騎驢進入耶路撒冷.
以一個王者身份.
一個以愛以謙卑來到耶路撒冷的王者.
他進入耶路撒冷.
然後發生什麼事.
有一連串的事情.
不可能不發生.

$^{721}$這是耶穌生平最後的一個星期.
發生的事.
他進入了耶路撒冷.
就潔淨聖殿.
然後你記得他講過很多其他東西.
還說不結果的無花果樹要被咒坐.
其實是指什麼.
指他自己的子民耶路撒冷的子民.
是不願意結果治.
他進入耶路撒冷的時候.
人有沒有歡迎他.
進去的時候.
隨從.
然後他旁邊來的那些會眾.
那些門徒等等.
他們歡呼喝采.
或者有些種樹枝等等.
進入城之後就不是另一件事.
進入城的時候.
裡面的權貴去反對他.
潔淨聖殿.
然後有很多的爭論.
爭論他有沒有王權.
所以每一個和薩都哥人.
法利賽人等等.
不同的故事都是在說這件事.
他們就是不要這個王.
來做他們的君主.
所以這個婚宴的比喻.
說得正確一點.
其實是一個風雨.
諷刺.
諷刺什麼呢.
整個的比喻的故事.
只有那幾個字.
大概七至八個字.
就是在說什麼呢.
在說耶路撒冷.
猶太人.
整個舊約的歷史.

$^{761}$說出來.
他們不斷去殺害這些先知.
因為王者不斷派僕人來.
那個僕人告訴他們.
你們已經知道那個宴會.
已經準備好了.
我準備了一個天國的宴會給你們.
那個宴會已經準備好了.
你們來.
你們被邀請.
儘管進來.
他們就不去.
不肯.
will not.
並且將這些僕人去殺害.
最後連他的兒子也殺了.
所以王就大怒.
第七節.
發兵來到去攻打這個城市.
這裡其實是在說.
主後第七十年的時候.
猶太人的聖殿.
就是被羅馬軍隊來去拆毀.
也在說.
這些本身已經被邀約的人.
他們本身就是天國的繼承人.
在這個筵席裡面是有份的.
但是他們不肯進去.
天國的門就向他們散.
然後天國福音的門.
就向其他願意進來的人.
去大開.
包括信主.
肯信主的猶太人.
小部分.
更多的就是外邦人.
這班人本身不是被邀約的.
不過最後是因為恩典的緣故.
他們就進來了.
這個比喻本身應該可以在這裡完結的.

$^{801}$總之就是外邦人.
他們恩典能夠進入這個宴會.
very nice.
但是耶穌沒有停.
第十一節.
就說第三個的大點.
就是如何get ready.
要prepare什麼.
第十一節.
耶穌就繼續說.
他說在這個宴會裡面發生什麼事.
王就進去見賓客.
看見那裡有一個沒有穿禮服的.
就對他說.
朋友你到這裡來.
怎麼不穿禮服呢.
那人無言可答.
王就對侍從說.
捆起他的手腳.
將他吊在外面的黑暗裡面.
在那裡哀哭之前.
因為被召的人多.
損上的人卻少.
但都很奇怪.
正如我所說.
那個比喻可以停在第十節.
但耶穌再加第十一節.
就是說那些所謂外邦人.
本身沒有被邀約.
但是他因為恩典的緣故被邀.
能夠進入這個婚宴的時候.
居然有什麼事發生呢.
王去見賓客.
就好像今天的新人.
到最後吃完飯.
他們去哪裡呢.
去敬酒.
到處去敬酒.
剛才我說的是不是吃七天的.
敬酒.

$^{841}$敬到你真的.
廁所沒時間去.
到處去見很多賓客.
因為星期一有些來.
星期三有些來.
星期四有些來.
然後走.
他們就經常去做這件事.
王就進去見賓客.
看見一個沒有穿禮服的.
這裡有很多不同的解釋.
有些人說為什麼會無端端有個人.
沒有穿禮服呢.
有不同的解釋.
第一個解釋就是.
可能主人家.
特別是以前.
一些的富戶.
他會為那些賓客進來的時候.
門口就給他們一些衣服.
一些禮服.
等他們進來的時候.
他們就可以再穿上去.
去那些宴會好不好.
就是說你穿短褲T恤就搞定了.
他去到門口就穿西裝.
給你了.
那你就可以坐在宴會裡面.
能夠去到享受.
這是其中一個解釋.
另一個解釋可能就是.
總之那個人沒有穿禮服.
你去宴會.
你都整齊一點.
對不對.
你沒有踢到拖鞋.
給面子也好.
怎樣解釋都好.
很可能.
就是說這個禮服代表什麼呢.

$^{881}$最重要不是中間發生什麼事.
為什麼他沒有穿禮服.
而是禮服代表什麼.
禮服應該代表在新約裡面.
信主的人.
是有衣袍.
白色的衣袍.
在馬太福音裡面.
經常都說到一些天國的宴會.
天國的附會.
十個童女等等.
就是說人進到這個禮堂的時候.
他需要去準備.
誰為他準備.
是基督的十字架的救恩為他準備.
你進到這個宴會.
天國的宴會.
能夠進去的.
就是因為上帝已經為你預備了衣袍.
所謂恩順清儀.
上帝看你.
不再看你本身.
你衣衫藍柳那個.
因為你是一個罪人.
但因為耶穌在十字架上為你犧牲.
你一個愛的救恩.
就為你穿上一條新袍.
這件新袍就可以帶你進去這個宴會.
馬太福音裡面前前後後經常提到一件事.
就是有些人進了天國.
以為自己進了天國.
但卻最後在外面哀哭截止.
這句話經常出現.
十個童女比喻.
五千個二千一千僕人的比喻.
或者登山寶訓.
也有說.
說什麼呢.
進到天國的時候.
有些人說主啊主啊.

$^{921}$我們為你做了很多事.
但主現在跟他說什麼.
我不認識你.
I don't know you.
終末的時候.
這裡就有一個驚喜.
被召的人多.
選上的人.
他就說少.
在這裡說到最後.
一個的婚言.
拒絕赴會的.
那些人當然活該.
本身他們可以承繼的.
他們卻不肯.
被恩典能夠邀請進入婚宴.
是一個很大的祝福.
如果我們都不肯的話.
就等於什麼呢.
等於我們.
你試一下.
站在十字架面前.
那個為你死的耶穌.
你不肯去接受.
其實就是看著這個十字架.
和這個愛你的.
犧牲為你去救你的救主.
你跟他說.
You are not worth it.
你沒有價值.
I don't care.
這個就是宴會裡應該有的準備.
如果你明白上帝的愛.
就不能夠不接受.
你這個恩典.
是一個美好的示威.
對嗎.
最後我想說的是.
有人說.
你進入天國的時候.

$^{961}$會有三個的wonder.
你有沒有想過.
你去天國.
進入天國的時候.
會有三個的wonder.
三個什麼wonder.
第一個wonder.
就是你以為.
有些人.
進不去天國.
但他沒有進去.
第二個wonder.
就是以為.
有些人進不去天國.
但他們在那裡.
第三個wonder.
就是我竟然在天國.
上帝給我的恩典.
基督給我的義抱.
我不知道哪個wonder會hit到你.
我希望是.
第三個.
居然我可以在那裡.
上帝給我的恩典.
宴會.
我能夠參與.
因為基督為我所做.
我們結束的時候.
我們用聖餐結束.
聖餐在一個.
基督教的歷史裡面.
教會的歷史裡面.
有一個很長遠的傳統.
聖餐是一個的頭台.
將來天國宴會的頭台.
那個appetizer.
餅和杯.
他告訴我們.
在十字架面前.
耶穌對我們的愛和恩典.

$^{1001}$是有多大.
不是一些等閒事.
我們同心協力祈禱.
然後我想請每一位.
已經接受了耶穌.
我們去領受這個聖餐.
上帝我多謝你.
多謝你對我們的恩典.
對我們的簡說.
我們能夠認識你.
為我們所作的.
你所犧牲.
你流出的寶血.
使我們能夠得到潔淨.
讓我們能夠被邀請.
進入這個天國豐盛的宴席.
主要是仰望將來會去到的宴席.
將來會去到的天國裡面的宴會.
今天主要是求你準備我們.
我們小信的今天來到你面前.
求你.
去到加添我們的信心.
可以有軟弱的主.
求你在你的恩典裡面.
帶領著我們.
但最骨子裡面要我們知道.
是因為這個恩典.
我們能夠進入天國的宴席.
讓我們拿起餅的杯.
去明白你的恩典是多麼的大.
讓你十字架的九恩.
去改變我們.
我們有幾張聖餐桌.
如果你方便的話.
很清楚你和耶穌的關係.
請你拿一個餅.
回去你的座位.
請你把聖餐的杯子放在你的座位.
我們求你,作為我們的主宰.
讓我們聽見你的話.

$^{1041}$也能夠進行你的動力.
在我們個人的生命,在我們的家庭.
在我們所接觸的一切上.
主要是願你自己,做我們的主人.
去帶領,去管理我們.
我們感激你,奉主命.
願上帝賜給大家平安.
下星期繼續跟大家一起去敬拜.
拜拜!.
\newpage



\section{}
\label{sec:D2gsbWPv_jE}
\textbf{比對研究 17: 巴力宗教 10 - 列王記上 18 - 陳崇基牧師}
\newline
\newline
連結: \href{https://youtube.com/watch?v=D2gsbWPv-jE}{\texttt{https://youtube.com/watch?v=D2gsbWPv-jE}} ~~~~ 語音日期: 2023-09-09
\newline
\newline
\hyperref[sec:ovpilkoc6io]{\small{< < < PREV SERMON < < <}}
~
\hyperref[sec:index_chronic]{\small{[返順時目]}}
~
\hyperref[sec:index_scriptual]{\small{[返順卷目]}}
~
\hyperref[sec:EH56cMFj_rc]{\small{> > > NEXT SERMON > > >}}
\newline
\newline
$^{1}$我們繼續處理這個《獵王機殤》.
《獵王機殤》第十八章.
就是這個很出名的加密山的故事.
上次我們說過17至19章的故事.
其實是一個.
你看完這麼多的故事.
其實它的編排是有它的脈絡在那裡的.
它的脈絡是說過.
就是由一個宣告開始.
然後有一個旅程.
一個movement.
然後有兩個encounter.
兩個相遇.
然後就有一個神蹟出現.
然後就有一些改變.
那些看到神蹟的.
或者接受神蹟的人.
就會有一個改變.
我們看17章的時候.
這樣去順序的時候.
你就發現.
其實18章的脈絡.
是非常相似的.
所以這些這樣的聖經的作者.
他們去到佈局的時候.
或者將這些故事.
放在一起的時候.
它的arrangement.
它的安排.
其實是有它的意思的.
所以我們來到18章的時候.
我們之前說過.
俄巴底.
然後.
以利亞和阿哈.
去到有一個confrontation.
就是有一個相遇.
然後就去到加密山上面.
之前也曾經提過.
加密山的地理.

$^{41}$我們就看看一些經文.
上下文.
遼宏紀.
其他的經文.
看看這個加密山.
為什麼會這麼重要.
18章20節.
就是阿哈派人.
到以色列的眾人那裡.
召集了先知.
上加密山.
Mount Karmel.
加密山我們之前曾經提過.
加密山就是以色列的海岸.
一直由南至北上去.
就有一個這樣的章在那裡.
加密山不是一個山.
加密山是一個山脈.
或者一個山區.
一個山地帶.
這個就是在以色列的北角.
就是薩瑪利亞的京城.
一直上去的時候.
一直向左上角.
伸展到海的.
這個地方就是今天以色列的.
其中一個大的城市.
叫做海法.
海法市.
海法.
這個加密山就好像插水一樣.
一直打斜插水.
到地中海.
在沿海的平原.
到達加密山的時候就停止了.
然後過了加密山之後.
就去到所謂的耶斯列平原.
拿撒拉就在上面.
這個整個平原就是東西的路徑.
東西的交匯點.

$^{81}$然後這個海岸線又繼續.
在阿阿科這個地方.
一直上去.
就是藍色的大地方.
我們叫做Felicia.
就是菲利基.
這裡有很多.
深的海港.
這裡就是已經去到今天的黎巴嫩的國境裡面.
但是這裡在舊約的世界裡面.
就是全部大的海港.
貿易的地方.
你就可以看到這個北角.
以薩瑪利亞為京都.
阿哈所在的地方.
跟這個菲利基.
推羅西頓.
是有邦交.
是很正常的.
如果我們知道一點點地理的話.
這裡不詳細說.
這些海港.
平地的地方.
是很少的.
這個是很正常的.
因為有深的海港的話就代表那裡比較進修一點.
所以他們平地是比較少的.
一些地理的常識.
就告訴我們.
這些海港.
雖然他們有貿易的好處.
但是他卻沒有農業的生產力.
是有限的.
因為他們的農地比較少一點.
所以這個北角.
跟他相.
大家做一個外交的時候.
其實是大家互相.
互保互利.
因為在北角.

$^{121}$這個地方就沒有一些大的海港.
這裡有一個叫.
Cesarea Martima.
這個在新約的時候才有的.
是大希律所建的一個人工的海港.
但是以前是沒有的.
以色列一直有的海域.
都是沒有深的海港.
但是他有的是比較多一點.
的農業的生產力.
比如大麥小麥等等.
果樹等等.
就是以色列的農產品.
所以跟他有建交.
是大家互相互利.
我就拿到.
你們的貿易的好處.
然後我又將我的田產.
可以供給你們.
這個就是聖地裡面.
一些地理的.
dynamic.
一些彼此之間.
互動的一些關係.
所以他就召集了去.
加密山.
加密山很明顯.
就是北面.
菲利基.
和南面以色列.
一個交界的地方.
特別是去到所羅門之後的時候.
就更加.
菲利基的地方.
和北國就更加近.
為什麼呢?.
因為回到所羅門的記載.
當他執政的時候.
他做了什麼呢?.
他就是跟.

$^{161}$北面的希蘭王.
有一個邦交.
我們看看這裡.
所羅門做聖殿和皇宮.
兩座殿宇的時候.
用了20年.
推羅王希蘭.
就在這裡.
以利亞的時候.
後一點的時候.
自然是西頓.
比較旺盛一點.
比較發展大一點.
但是不同的時期.
海港大家彼此有所謂的爭雄.
有時候是你大一點.
有時候我小一點.
有時候大家的商貿.
或者海港的生意會比較多一點.
在所羅門的時候.
就是推羅的時候.
他就跟希蘭.
有邦交.
他說曾經希蘭王.
推羅王.
就跟.
所羅門所有的資助的.
香柏木.
松木金子.
就將他運到來.
幫所羅門建這個聖殿.
所羅門交易上給他什麼呢?.
所羅門就將.
那把加里尼地的.
20座城就吸了.
吸了希蘭.
希蘭就從.
推羅出來.
然後察看所羅門吸他的城鎮.
看就看不上眼.

$^{201}$看不上眼.
他就說我兄啊.
你是跟我邦交的.
你給了什麼城鎮給我.
所以他就叫這些城鎮.
20座的城鎮.
加里20座城鎮.
就名叫加布納地.
直到今天.
重給.
所羅門120塔克蓮德的金子.
這個跟他邦交的.
就覺得很不順.
因為他給了20座加里尼地的城.
給他.
加里尼地在哪裡?.
加里尼地就在這裡.
他就給了20座.
可能隨便給了20座.
給推羅的希蘭王.
這個.
加布納在哪裡呢?.
卡布.
好幾年前在這個地方.
就有一個發掘.
就發現一些.
考古的跡象.
應該是.
不只是考古跡象.
留下的遺產.
而是一些.
推羅西頓的.
物質的文化.
就是說在以色列的境內.
居然.
有一些.
菲利基的.
物質的文化.
證明這個地方.
在所羅門的期間.

$^{241}$大概是主前.
十世紀.
主前九百多年的時候.
卡布這個地方.
就有.
土優.
和這些.
考古的跡象.
代表這個地方是一個城鎮.
當時是有人居住.
和有貿易的來往.
等等.
有些學者認為這個地方.
就是聖經裡面所說的.
希布納這個地方.
卡布這個地方.
卡布這個名字在哪裡呢?.
推羅王看不順眼的時候.
他就說這個地帶.
就叫做.
加布納.
什麼意思呢?.
希伯來文就是.
加布納地.
地帶.
地的意思.
地帶的意思.
卡布.
可能解經者的推斷.
就是由兩個字.
而去組成的.
這兩個字是什麼字呢?.
就是cur.
cur就是有餘的意思.
彷彿.
包.
包就是希伯來文裡面.
一個反義詞.
but的意思.
no的意思.

$^{281}$negative的意思.
詩篇裡面有不同的經文都提過.
用過這個字的包.
卡布納.
就是.
加布納的意思.
什麼意思呢?.
就是他看到這個所羅門給他的城的時候.
他下來視察的時候.
這個推羅希蘭王就說.
卡布.
Aris Kavu.
就是說這個地方.
等於廢墟.
或者廢物.
所以這個所羅門之後.
這個卡布納地.
就屬於.
推羅或者西頓.
就是這個菲利基的地方.
如果你這樣看的時候.
你就會發現.
這個菲利基的地方.
是有推羅,西頓,比布斯,比布洛斯.
這些地方.
這些海港.
但是到所羅門的時候.
佛利基.
這個菲利基的文化.
或者他的影響.
已經來到下面.
就是因為所羅門已經給了一些城.
給推羅王.
所以就屬於.
菲利基的地帶.
所以加密山就更加近.
如果你去過聖地的話.
去加密山.
其中一個山峰.
這裡有一個修道院.

$^{321}$就是這個卡米爾.
的修道院.
就是一個以利亞的.
教堂在這裡紀念這件事.
在上面一看的時候.
你會發現.
左邊或者.
西邊就是地中海.
然後在你後面.
就是.
以前的北國.
薩馬利亞的山地.
在北邊就是.
菲利基的海港.
就是說一看.
在北邊的時候.
看到拿撒勒的左邊.
這個海岸線的地帶.
就是卡夫爾這個地方.
所以加密山.
真的在.
菲利基和以色列.
中間的地方.
所以在.
加密山上面有一個.
這樣的爭鬥.
決定到底耶和華.
是神還是巴力.
為神.
在地理上.
是絕對合理的.
對不對?.
如果你看看這個十八章.
繼續記載的時候.
剛才我們看了二十節.
去到加密山這個中間的地方.
南北之爭.
上面是巴力.
下面就是耶和華的信仰.
是不是?.

$^{361}$以色列國.
那麼你會知道二十一節是想說什麼.
以利亞和百姓說什麼呢?.
他們說你心慈意義要到什麼時候.
如果耶和華是上帝.
就當順從耶和華.
如果是巴力就當順從巴力.
他們就一言不談.
你想想在山上可以講這個故事的時候.
你就覺得.
地理上和整個的環境.
其實都是在說.
心慈意義是什麼意思.
你看看南邊.
南邊就是耶和華的信仰.
北邊就是什麼?.
巴力.
心慈意義到什麼時候.
就是在上還是下.
前還是後.
左還是右.
在這個地方就要決定.
整個的以利亞的故事.
其實就是告訴我們.
聖經的信仰是一種獨特的信仰.
神就是唯一的神.
要拜上帝的話.
就只能夠拜上帝.
如果你拜巴力的話.
你可以加上耶和華也可以.
你拜巴力的話.
拜其他東西的話.
你也可以拜耶穌.
但是在聖經的信仰裡面.
就是一個排他的信仰.
你一定要不可以心慈意義.
一定要決定到底.
哪一個才是真誠.
如果耶和華是真誠的話.
你就只會順從耶和華.

\newpage



\section{2:14-26}
\label{sec:EH56cMFj_rc}
\textbf{A Living Faith (James 2\_14-26) - Caleb Lee}
\newline
\newline
連結: \href{https://youtube.com/watch?v=EH56cMFj-rc}{\texttt{https://youtube.com/watch?v=EH56cMFj-rc}} ~~~~ 語音日期: 2023-09-13
\newline
\newline
\hyperref[sec:D2gsbWPv_jE]{\small{< < < PREV SERMON < < <}}
~
\hyperref[sec:index_chronic]{\small{[返順時目]}}
~
\hyperref[sec:index_scriptual]{\small{[返順卷目]}}
~
\hyperref[sec:p__HCvfoaaM]{\small{> > > NEXT SERMON > > >}}
\newline
\newline
$^{1}$Good morning, everyone..
Okay, it works..
That's good..
Looks like people are doing well..
As always, it is a great privilege, an incredible privilege to be here to see you all this morning..
So let me first open us in a word of prayer..
If you would bow your heads with me, let's pray..
Heavenly Father, our God and our Father, what is man that we can stand before you in this.
moment that you would hear us?.
Only by your son, Lord, can we stand before you righteous and justified, seen not as enemies.
but as friends and as children, as heirs of your kingdom..
And so, God, I ask in this moment that you would be with us, that you would make your.
presence known to your people here, that you would magnify yourself..
Would you guard them against anything I say that might not be of you?.
Would you use what I say and help us see your word as you have spoken it to your people,.
Lord?.
This time is yours..
These people are yours and we pray all this in Jesus' name..
Amen..
Okay..
So, as many of you are aware, just looking at the slide, we have been in the book of.
James and we will be continuing our reading of the book of James, which is a letter from.
James to the Jewish Christians of dispersion..
And what we find in James is that the theme of James is followers of Jesus are to be affected.
by their faith, not just hearers of the word of God, but also doers of the word of God..
And this is James' call for his brothers and sisters, for them to live wise lives, lives.
that are reflective of the gospel that they have received..
Last time, we saw James address the way in which the Jewish Christians were to treat.
the poor and the needy that come into the midst of their gathering, how they are not.
to make distinctions among themselves based on outward appearance or seek to curry favor.
of those who are rich in the eyes of the world..
Instead, they are to show mercy and compassion, to love their neighbors just as God had loved.
them in their need..
Rather than seek favor from man, they are to look to God and God alone for their needs..
You see, the apostle James doesn't want God's people to live double lives..
He wants them to be whole, perfect, and complete, to be rid of the flaws in character that come.
from keeping parts of themselves separate from the effects of the gospel, and in doing.
so become more and more like Jesus..
Yet James' instructions, okay, his instructions are not just given because he wants you to.
be more Christ-like..

$^{41}$That is incredibly important, but not just that..
He gives them because it is a matter of life and death..
Having said that, let's look at today's passage..
If you have your Bibles with you, I invite you to turn with me to James chapter 2, verses.
14 to 26..
Again, that is the book of James, its second chapter from verse 14 to the end of verse.
26..
It's on the slide here, and I'll just read it for you guys, okay?.
Here we are..
James chapter 2, verses 14 to 26..
What good is it, my brothers and sisters, if someone claims to have faith but has no.
deeds?.
Can such a faith save them?.
Suppose a brother or a sister is without clothes and daily food..
If one of you says to them, "Go in peace, keep warm and well-fed," but does nothing.
about their physical needs, what good is it?.
In the same way, faith by itself, if it is not accompanied by action, is dead..
But someone will say, "You have faith, I have deeds..
Show me your faith without deeds, and I will show you my faith by my deeds..
You believe that there is one God, good..
Even the demons believe that, and shudder..
You foolish person, do you want evidence that faith without deeds is useless?.
Was not our father Abraham considered righteous for what he did when he offered his son Isaac.
on the altar?.
You see that his faith and his actions were working together, and his faith was made complete.
by what he did..
And a scripture was fulfilled that said, "Abraham believed God, and it was credited to him as.
righteousness, and he was called God's friend.".
You see that a person is considered righteous by what they do and not by faith alone..
In the same way, was not even Rahab the prostitute considered righteous for what she did when.
she gave lodging to the spies and sent them off in a different direction?.
As the body without the spirit is dead, so faith without deeds is dead..
We start here..
We start here in verse 14, with this question placed right before you..
What good is it?.
What good is it for someone to claim faith but exhibit no deeds?.
The scenario is given here, where there is someone who needs clothes and food..
James asks his audience, he asks them this question, "What good is it to pray and give.
blessing, right?.
Go in peace, keep warm and well fed without providing food or clothes..

$^{81}$What good is it?".
And so the implication here is that it is of no good..
It is not enough..
Before the blessing was pronounced, the person was in need of clothing and food, and after.
the blessing, that person is still in need of clothing and food..
The need is unfulfilled..
This is, to James, something that is significantly wrong..
Something is incredibly wrong with the faith of someone who only says, "Go in peace, keep.
warm and well fed," and does not provide clothes or food..
You can see for yourself just how wrong James considers it..
If you look at verse 14 with me, in the same way, faith by itself, if it is not accompanied.
by action, is dead..
James calls this faith dead, useless and dead..
No, it is of no good if someone claims to have faith but has no deeds..
Their faith is dead..
We see here in verses 14 to 17 a warning..
Do not confuse faith professed with faith practiced..
In other words, don't mistake just saying, "I have faith," with faith that results in.
action..
By itself, faith professed does not save..
For the person who says, "I have faith," will also be moved by their faith..
They will not just be hearers of the Word of God, they will be doers of the Word of.
God..
And this is something James understands..
Something that you and I need to understand..
It is a living faith that saves..
Okay, let me repeat that for you..
It is a living faith that saves..
Let's not miss this frightening statement from verse 14..
You see this one?.
The one pretending to be a question?.
What good is it, my brothers and sisters, if someone claims to have faith but has no.
deeds?.
Can such a faith save them?.
A faith that makes claims but is without deeds, can that faith save?.
Again, it's implied, right?.
It's implied..
No, that faith cannot save..
Faith without action is dead..
And a dead faith cannot and will not save you..

$^{121}$This I'll repeat for you as well..
Faith without action is dead..
And a dead faith cannot and will not save you..
17 to 19..
We see here that James quickly moves into an exchange between himself and a person who.
is disagreeing with him..
The disagreeing man tells him, "You have your deeds, I have my faith.".
And what is being said here is, "Will your deeds save you or will my faith save me?".
He's challenging him..
And this is a very real situation, right?.
Because you, you having come to faith, having received your salvation, you didn't receive.
it because you earned it, right?.
This is an absolutely true thing..
And James does not dispute this fact..
You having come to Christ did not receive your salvation because you earned it..
He does not dispute that salvation is in fact a supernatural work of God..
But James does dispute something..
Look at his response with me, okay?.
"Show me your faith without deeds and I will show you my faith by my deeds..
You believe that there is one God, good..
Even the demons believe that and shudder.".
Firstly, faith is not just knowledge of what is and is not true..
It goes further than that..
Just knowing something isn't enough to save you..
Believing that there is one God does not save you..
In fact, look, James makes it clear that demons know even more than you do that the God of.
the Bible is the God of heaven and earth and this knowledge does not save them..
In fact, it makes them shudder..
They tremble in fear because their knowledge does not change their behavior..
All right..
Secondly, looking at this, there's an error in this argument that's being made against.
James..
The error made against James is separating faith and works..
The error is thinking that saving faith can be divorced from merciful works..
Action and works aren't something additional to living faith..
Let me ask you a quick question, okay?.
You guys will ace this..
How many of you right now are breathing?.
Right?.
All of you..

$^{161}$Everyone here is breathing..
And so is your breathing an addition to your living?.
Is breathing something extra that you are doing?.
No, right?.
No..
Breathing is a function of living..
If you want to continue to live, you need to continue to breathe..
And so when we look at works the way that James looks at works, we see that works are.
as breath to those that follow Christ Jesus..
The action following faith proves that it is a living faith..
And this is not a new reality about faith..
This is the way that God has ordained living faith since the very beginning..
If you look with me to verses 20 to 24, you'll see that James begins to talk about someone.
named Abraham, the father Abraham, who is the patriarch of these Jewish Christians,.
all right, from the line of faith..
And we find that Abraham is considered righteous for what he did because of his interaction,.
because of the interaction between his faith and his actions..
Abraham's faith was alive..
It led him to obey God, to offer his son Isaac on the mountain where God's goodness and grace.
provided a substitute in Isaac's place..
And through this cooperation of his faith and his works, Abraham's faith was made complete..
It was able to fulfill its purpose..
Living faith is shown by its very nature..
Living faith is demonstrated through action..
When James writes in verse 24, look at verse 24 with me, he says like, he says something.
here, not by, oh, is 24 on the screen?.
No, it's not..
Let me, yes, okay, there we go..
He says something here..
You see that a person is considered righteous by what they do and not by faith alone..
This faith that he talks about, that he refers to, is a faith that is dead with no ability.
to save because it's a living faith that saves, a living faith that takes action..
Verse 25, just like Abraham, Rahab the prostitute was considered righteous because of her faith.
and her actions..
These figures that find themselves in the honored line of David stand before God through.
a living faith..
It wasn't a faith without action that saved them..
It was a faith completed by action..
And in the same way, faith without action cannot save you..
You can't just say, I believe in Jesus' death and resurrection and then carry on with.

$^{201}$your life as if nothing has changed..
Action apart from faith cannot save you..
You know this to be true..
You cannot earn your way to heaven..
Your actions that come out of a personal standard of righteousness will always fall short of.
God's perfection..
Only a faith accompanied by action, a living faith, can save..
Having said all that, like this brings us to an incredibly important question that you.
must answer..
Do you have a living faith, a faith that compels you to not just be hearers of the Word of.
God but also doers of the Word?.
Do you have this faith that can and will save you, a faith that lives?.
Consider this carefully because you are not meant to be dead..
My brothers and sisters in Christ, you are meant to be alive..
Your body is not without a spirit..
You are meant to be the dwelling place of the spirit of the living God..
While you and I were still enemies of Christ, He looked and He saw you..
He saw you in your disregard and opposition to Him and still He reached out to you..
While you and I were dead in our sin, the Word of God who was with God and was God,.
Jesus Christ, embraced death so that He could pull you out of death and into life..
On a wooden cross, Jesus placed Himself under the punishment you rightfully deserved and.
bought forgiveness for all your wrongdoing..
Every sin you have and will ever commit, in glory, Jesus rose from the grave to give you.
life in abundance..
Do you see?.
Do you see here?.
You aren't supposed to be like you were before..
You can't be like you were before..
Only the dead are comfortable in their coffins..
You are meant to be living children of God with a living faith..
And James tells you up front how this is supposed to look..
The book of James, right, as we've been going through it, it gives instruction after instruction.
on the way that Christians are supposed to live, the way that Christians should relate.
to those that are around them..
The verses right before today's passage spoke on how the Jewish Christians were to treat.
the poor and the needy that came into the midst of their gathering, how they were not.
to make distinctions based on outward appearance or to seek favor from those who are rich in.
the eyes of the world..
So we see in James that works of mercy are the natural and fullest expression of living.
faith..

$^{241}$Acts of love, of goodness, of charity, acts that provide..
This is James's understanding of living faith..
In this crowd, are you a husband, a wife, a father, a mother, a son, or a daughter?.
If that's you, then be merciful, compassionate to your mother, your father, your son, your.
daughter, your neighbors..
Be compassionate and merciful..
I understand that none of this is easy..
It is not easy to be merciful or compassionate because we are not an easily lovable people..
Because our unmet needs, they cause the cracks and the flaws in our character to widen, and.
then out of that comes sin..
So as we close up our time here, continue to fix your eyes on Jesus and remember what.
he has done for you, how he loved you to make you lovable, and how he redeemed you from.
the pit..
In that same way, look upon your neighbor, those that God has placed around you, and.
be both hearers and doers of the word of God..
Live as a living children of God with a living faith..
Be whole and complete..
Let me close us in a word of prayer..
\newpage



\section{尼希米記 1:1-11}
\label{sec:p__HCvfoaaM}
\textbf{靈命省察:四個重要的問題 (尼希米記 1\_1-11) - 張麗仙傳道}
\newline
\newline
連結: \href{https://youtube.com/watch?v=p_-HCvfoaaM}{\texttt{https://youtube.com/watch?v=p\_-HCvfoaaM}} ~~~~ 語音日期: 2023-09-13
\newline
\newline
\hyperref[sec:EH56cMFj_rc]{\small{< < < PREV SERMON < < <}}
~
\hyperref[sec:index_chronic]{\small{[返順時目]}}
~
\hyperref[sec:index_scriptual]{\small{[返順卷目]}}
~
\hyperref[sec:1Q8_M16V0ZI]{\small{> > > NEXT SERMON > > >}}
\newline
\newline
尼希米記 1:1-11
\newline
\begin{longtable}{cl}
\hline
\hline
章節 & 經文 (和合本修訂版)\\
\hline
1:1 & \begin{tabularx}{0.7\textwidth}{X} 哈迦利亞的兒子尼希米的言語如下：亞達薛西王二十年基斯流月，我在書珊城堡中。 \end{tabularx} \\ \\ \relax
1:2 & \begin{tabularx}{0.7\textwidth}{X} 那時，我有一個兄弟哈拿尼，同幾個人從猶大來。我問他們那些被擄歸回、剩下殘存的猶太人和耶路撒冷的情況。 \end{tabularx} \\ \\ \relax
1:3 & \begin{tabularx}{0.7\textwidth}{X} 他們對我說：「那些被擄歸回剩下的餘民在猶大省那裡遭大難，受凌辱；耶路撒冷的城牆被拆毀，城門被火焚燒。」 \end{tabularx} \\ \\ \relax
1:4 & \begin{tabularx}{0.7\textwidth}{X} 我聽見這話，就坐下哭泣，悲哀幾日，在天上的神面前禁食祈禱， \end{tabularx} \\ \\ \relax
1:5 & \begin{tabularx}{0.7\textwidth}{X} 說：「唉，耶和華－天上大而可畏的神，向愛你、守你誡命的人守約施慈愛的神啊， \end{tabularx} \\ \\ \relax
1:6 & \begin{tabularx}{0.7\textwidth}{X} 願你睜眼看，側耳聽你僕人今日晝夜在你面前，為你眾僕人以色列人的祈禱，承認我們以色列人向你所犯的罪；我與我父家都犯了罪。 \end{tabularx} \\ \\ \relax
1:7 & \begin{tabularx}{0.7\textwidth}{X} 我們向你所行的非常敗壞，沒有遵守你吩咐你僕人摩西的誡命、律例、典章。 \end{tabularx} \\ \\ \relax
1:8 & \begin{tabularx}{0.7\textwidth}{X} 求你記念所吩咐你僕人摩西的話，說：『你們若犯罪，我就把你們分散在萬民中； \end{tabularx} \\ \\ \relax
1:9 & \begin{tabularx}{0.7\textwidth}{X} 但你們若歸向我，謹守遵行我的誡命，你們被趕散的人雖在天涯，我也必從那裡將他們召集回來，帶到我所選擇立為我名居所的地方。』 \end{tabularx} \\ \\ \relax
1:10 & \begin{tabularx}{0.7\textwidth}{X} 他們是你的僕人和你的百姓，是你用大力和大能的手所救贖的。 \end{tabularx} \\ \\ \relax
1:11 & \begin{tabularx}{0.7\textwidth}{X} 唉，主啊，求你側耳聽你僕人的祈禱，聽喜愛敬畏你名眾僕人的祈禱，使你僕人今日亨通，在這人面前蒙恩。」我是王的酒政。 \end{tabularx} \\ \\
[1ex]
\hline
\hline
\end{longtable}
$^{1}$麻煩Carmen.
大家看我上來不覺得出奇.
不過出奇的是今天Marcel去了哪裡呢?.
其實他沒有躲起來.
大家看到節奏的話都知道.
他說他這星期去了Edmonton.
加拿大一間教會.
做考古的workshop.
他說「我只懂這些東西」.
所以今天我在這裡.
各位聽眾.
一年容易又中秋.
中秋快到了.
9月29日.
這個月底MetalSup就要來了.
但是我加一句.
一年容易又一歲.
同意嗎?.
真的很快又一歲了.
我想問你什麼時候生日呢?.
你不如說一句.
一年容易又一歲.
我想問你什麼時候生日呢?.
你不如告訴隔壁的人.
不管他是你老公老婆.
都告訴他.
提一提大家.
你幾月生日呢?.
或者你忘記了.
還有就是.
我想問問大家.
你記不記得你屬靈生日是什麼時候呢?.
有些人馬上想想.
你是受浸那天.
或者是你缺自信耶穌那天.
或者你是回歸神那一刻.
無論如何.
你的屬靈生日是什麼時候呢?.
這個懂不懂說呢?.
其實我最想問的就是.

$^{41}$大家.
你覺得你今天.
你的屬靈生命是怎麼樣的呢?.
我為什麼問這個問題呢?.
大家都知道.
我們有二四週五茶茶經.
我們在查尼希米記.
當我預備.
尼希米記.
第一堂.
我讀到的時候.
我心裡有一點激動.
為什麼呢?.
因為很有趣.
那一章下半段是一個尼希米的祈禱.
而這個祈禱.
我一直讀的時候.
一直發覺.
它一直在提醒我.
我的靈命是怎麼樣的.
這個短短的祈禱.
給我幾個很重要的.
靈命的醒察.
所以我今天都很想.
不單是那晚回來的弟兄姊妹.
有二十位.
因為我們時間趕了.
我沒有時間詳細地.
和你們分享這個.
我不想你們錯過了.
我也不想再再每一位錯過.
這麼好的祈禱.
包括網上的.
我也不想你們錯過這麼好的祈禱.
我們要說一下簡單的背景.
尼希米記究竟在哪裡呢?.
你一定要先知道.
舊約的聖經歷史.
創世記是一個頭.
尼希米記.

$^{81}$差不多是去到尾.
是舊約的歷史.
創世記神創造天地.
尼希米記是說神的子民.
是被老回歸.
是這一段日子.
在這一段日子裡面.
當時我們熟悉聖經的.
大家都知道.
當時神的子民以色列.
雖然是被巴比倫.
是老了一段很長的時間.
但是當波斯帝國.
滅了巴比倫的時候.
波斯王一上場.
就很奇怪.
採取懷柔政策.
大家知道什麼是懷柔政策.
就是對你好一點.
讓你可以對我有用一點.
懷柔政策就是.
他不再讓以色列人.
做他們國家的俘虜.
放他們回去.
他們猶大的國土.
耶路撒冷這些地方.
本來是很開心.
本來是很興奮的.
但是回國之後的情況很慘.
這個就引導.
我要和大家說.
今天這個祈禱.
尼希米的祈禱的一個背景.
我們先看看聖經.
請聽我很快地讀.
阿加利亞的兒子.
耶利米的言語如下.
阿達斯西王二十年.
基斯留約.
我在書山城堡中.

$^{121}$那時我有一個兄弟哈拿尼.
和幾個人從猶大來.
我問他們那些被老歸回.
剩下殘存的猶太人.
和耶路撒冷的情況.
你們回到耶路撒冷.
怎麼樣.
他們對我說.
那些被老歸回剩下的漁民.
在猶大省那裡遭大能.
受凌辱.
耶路撒冷的城牆被拆毀.
城門被火焚燒.
我聽見者說.
就坐下哭泣.
悲哀幾日在天上的上帝面前.
禁食祈禱.
這個就是.
祈禱的背景.
我們看看他怎麼祈禱.
說.
耶路說天上大而可畏的上帝.
向愛你.
遵守你誡命的人.
首約施持我的上帝.
願你睜眼看.
則耳聽你僕人今夜就夜.
在你面前.
為你眾僕人以色列人的祈禱.
承認我們以色列人.
向你所犯的罪.
我與我家.
我與我父家都犯了罪.
我們向你所行的.
非常敗壞.
唯有遵守你.
吩咐你僕人摩西的誡命.
而你卻不敢演章.
叩你紀念所吩咐你僕人摩西的話說.
你們若犯罪.

$^{161}$我就把你們分散在萬民中.
但你們若歸向我.
謹守遵行我的誡命.
你們被趕散的人.
雖在天涯.
我也必從那裡將他們召集出來.
帶到我所選擇.
立為我名居所的地方.
他們是你的僕人.
和你的百姓.
是你用大力和大能的手所救贖的.
唉.
主啊.
求你質以聽你僕人的祈禱.
聽喜愛敬畏你命中僕人的祈禱.
使你僕人今日亨通.
在這人面前蒙恩.
我是王的酒精.
尼希敏.
沒錯.
雖然回國是一件很興奮的事.
但是.
尼希敏聽到的.
就是一些這麼不好的場面.
你聽到他說.
在那裡.
被老龜狐盛夏的漁民.
在那裡怎樣.
遭大能受凌辱.
耶路撒冷的城牆被拆毀.
城門被火焚燒.
還有很多外族人的攻擊.
這裡沒有寫出來.
大家都應該明白這個心情.
因為上個禮拜.
香港黑雨.
水災得多厲害.
地鐵那裡都是水浸了.
好像浴缸一樣.
大家都覺得很擔心.

$^{201}$馬上問香港的家人.
因為我的家人在那裡住.
或者如果國內有什麼事.
大家都很擔心.
很擔心.
經常去慰問.
因為那是我們的國家.
我們的故土.
我們家人在那裡.
很憂心.
這個就是尼希米的心情.
尼希米當時的緊張.
是有原因的.
雖然.
尼希米在波斯那裡.
是做酒政.
是一個很高的官.
皇帝很信任的.
他其實是舒適豪華.
很尊貴的身份.
但是他始終是猶太人.
他心裡面是.
顧念猶太國土的事.
所以一聽到這個危機.
他就說.
我聽見這個危機.
我就.
我一聽見這個危機.
我就坐下哭泣.
為什麼不去那段.
我就坐下哭泣.
悲哀幾天.
在天上的上帝面前.
禁食祈禱.
這個是尼希米當時的心情.
為什麼.
我講這個祈禱.
和大家有什麼關係呢.
為什麼我今天.
要特別和大家講這個祈禱呢.

$^{241}$就是說我們可以在這個祈禱裡面.
看看我們今天的靈命.
究竟是怎麼樣.
有很多的提醒.
起碼有四個提醒.
第一個提醒就是.
究竟我當神是誰.
為什麼這樣說呢.
尼希米當神是誰.
一開始的話.
耶和華天上.
大而可畏的上帝.
向愛你守禮誡命的人.
守施慈愛的神.
我們知道.
是這樣嘆息.
他的心情很沉重.
但是在這個沉重的心情之下.
尼希米稱神為什麼呢.
天上.
大而可畏的神.
耶和華天上.
大而可畏的神.
他以這個為他的祈禱開始.
耶和華在.
以色列人的字典裡面已經是.
創天造地.
大而可畏的那一位.
耶和華在以色列人的歷史裡面.
更加是他呼召他們的先祖.
亞伯拉罕.
揀選了以色列做他的子民.
和他們立約的神.
而歷史一次又一次的證明.
耶和華永遠都守約的.
守施慈愛.
對他的子民.
不像當時其他的神明.
是很專制.
很自私的.

$^{281}$耶和華時時都守施慈愛.
對他的子民很寬容.
你看得到嗎.
你今天看得到嗎.
這個是尼希米心裡面的神.
所以他一禱告開始.
即使他心情多沉重也好.
他首先稱呼.
你是天上大而可畏的神.
你是守施慈愛.
你就是守施慈愛.
你是守施慈愛的神.
不是漂亮.
而是真誠的說話.
他真的這樣覺得.
而說出來的.
我們問問自己.
我們祈禱的時候.
有沒有一開始就對神有這種尊崇呢.
你有沒有呢.
很妙.
我不知道你有沒有聽過一些.
弟兄姊妹這樣說.
我聽過很多次.
總括來說.
我今天跟神說.
為什麼我祈禱了這麼久.
神也不出手.
如果是這樣.
你想我怎樣呢.
這些人祈禱其實也很好.
他真的當神是老友.
我們當神是老友是可以嗎.
可以的.
我們當神是老友是可以的.
甚至很多詩篇告訴我們.
你在質問神.
質疑神也可以.
你可以當神是你的爸爸.
神想我們這樣.

$^{321}$神的確想我們跟他有這麼親密的關係.
當他老友.
當他爸爸.
你可以向他投訴.
但是無論怎樣也好.
你是要有一個這樣的心懷在.
你是尊崇敬畏他.
你是以尊崇敬畏他的心.
你不是尊崇敬畏他.
你是尊崇敬畏他的心.
你是以尊崇敬畏他的心.
當他爸爸當他老友.
甚至質疑他.
不失去尊崇就好像尼熙米一樣.
他當時的心情很沉重.
他真的有一些質疑.
但是他仍然是說.
天上大而可畏的神.
守約斯前外的神.
我想起一句很有意思的說話.
是一個已故的牧師和作者.
很出色的牧師和作者.
A.W. Tozer.
他說過一句這樣的說話.
What comes into our mind.
when we think about God.
is the most important thing about us.
即是說.
當我們想起神的時候.
我們的腦海馬上出現什麼呢.
其實就反映我們心裡.
最重要的是什麼.
你代入自己.
當我一想起神的時候.
我就想起祂打救我.
即是反映.
當你是有很多的憂慮.
很多的困難.
如果你一想起神的時候.
你就想起祂審判你.

$^{361}$即是說.
你做了很多不好的事.
你有很多內疚在.
很怕神.
或者當想起神的時候.
你什麼都想不到.
即是說.
神在你心中根本沒有地位.
你都經常想不到祂.
沒有想過祂.
各位弟兄姊妹.
究竟你當神是誰呢.
往往在我們禱告.
第一句那裡就反映出來.
不是假的.
你.
究竟當神是誰呢.
我們今天的祈禱很多時候.
聖神是神.
天父 主.
不是說長就好.
但是這幾個.
這些明清的背後.
你有沒有真的當祂是天上.
大而可畏的神.
守施慈愛.
拯救你 憐憫你的神呢.
這個是最重要.
不是你的稱呼怎樣重要.
而是你的心裡面有沒有.
這種尊崇.
所以我覺得耶利米這個禱告.
一開始就給我一個靈命很重要的提醒.
就是.
我們有沒有.
我們的禱告 我們的靈性生活.
有沒有例行公事.
草草收場.
還是我們時時.
首先對這位神.

$^{401}$尊崇敬畏呢.
我覺得絕對是值得我們去思想.
經威耀霸.
是智慧的開端.
聖經一直強調.
這個也是生命的祝福.
願大家都去審查這個問題.
我當神是誰.
第二.
稱呼完了.
耶利米給我們第二個的審查是.
我是不是看罪很少呢.
今天黎師兄好像很通俗.
有時不通俗.
好像說不到心裡面的說話.
我看是不是當神很少呢.
我們來看看他怎樣說.
他說 神啊.
願你睜眼看.
則以聽你僕人.
今日晝夜在你面前.
頌僕人 以色列人的祈禱.
承認我們 以色列人.
向你所犯的罪.
我與我父家都犯了罪.
當耶利米稱呼.
元大而可畏.
首約施賜我的神.
他其實是應該有些很重要的東西.
祈求的 為耶路撒冷.
但他沒有祈求.
他馬上做了一件什麼事呢.
就是認罪.
很奇怪 我不知道你覺不覺得奇怪.
如果我們有什麼事.
我們當然要求神.
還首先認罪.
我想仍然的是這一句.
耶利米知道.
聽祂禱告的.

$^{441}$那一位是誰.
是大而可畏的神.
我首先要在你面前.
思想到我今天.
落到這樣的田地.
是我們犯了罪.
我們全國都犯了罪.
而且耶利米不敢稱高.
我自己不是為人犯罪.
是認罪.
如果我們做傳道人牧師.
不在為你們認罪的時候.
我們就真的太精高了.
耶利米在這裡說.
我自己都犯了罪.
他明白到在這個時候.
什麼都不重要.
不要先求.
首先要在神面前.
Be honest.
我要誠實認自己的罪.
這個是勇氣.
我有一張很漂亮的沙發.
上星期日我們完結了會議.
90多度.
Ivan坐在我那張沙發.
我們去一個弟兄姊妹的家.
他們搬回來.
Mark Sir,Kaleb和林廣.
在新水新旱搬回來.
那幾張椅子能坐多少人.
我數過應該舒舒服服.
能坐五個人.
坐到五個人.
令我想起一件事.
就是.
我看過一個廣告.
是這麼說的.
他說.
有一個人去買沙發.

$^{481}$進去.
那個售貨員告訴他.
這套沙發.
能坐五個人.
是沒有問題的.
那個人就說.
糟了 你真的考起我.
我怎麼找到五個沒有問題的人.
來坐呢.
不知道大家明不明白.
嚇死我了.
我忘記了這張沙發.
這個漫畫的意思就是說.
是什麼意思呢.
其實.
我們每一個人都是有罪.
不是我們旁邊的人有罪.
我們說認罪的時候.
我們是認自己的罪.
我記得很多時候.
跟弟兄姊妹在茶經班.
聊認罪的時候.
大家說說有什麼例子.
有什麼罪要認呢.
時時都聽到兩樣罪.
說謊.
貪心.
很多人都是說這樣.
我們的罪是不是這麼簡單.
他們說謊.
說謊貪心.
不是窮凶極惡.
大奸大惡那種.
是說謊貪心.
但是你有沒有想到.
我們有很多很多的罪.
自私自利.
妒忌的火.
惡毒的心.
無情冷漠.

$^{521}$有沒有這些罪.
是不是說謊和貪心.
你有沒有呢.
我就有.
我相信.
我們受這個世界的道德文化影響.
我們的生活模式.
也有很多的罪.
這個世界已經走得越來越偏離神的正道.
我們怎樣賺錢.
怎樣花錢.
怎樣用我們的時間.
我們的性關係.
我們還在乎神說什麼.
聖經教導什麼.
實在我們生活裡面.
已經積聚了很多的罪.
如果我們.
習慣了看罪.
是很小的事的時候.
就是我們遠離神.
越遠的時候.
越和神有隔的時候.
你能不能給我們.
很好的提醒.
危機當前.
我們求神保護.
首先的.
不要不想起.
你要在神面前認罪.
不過我在這裡想.
特別補充一句.
我的意思不是說.
你今天有危機 今天有大能.
就是你犯了罪.
我強調的就是.
我們需要時時在神面前.
審查我們的罪.
不過到現在.
或者你都還沒明白.

$^{561}$究竟你犯了什麼罪.
李希敏在這裡認什麼罪.
去到我們第三點.
就是給我們一些.
靈命的提醒.
我覺得這一點是非常重要的.
就是我們有沒有.
忠於與神所立的約.
首先我想問每一位基督徒.
在座每一位信了主的.
你知不知道.
或者你記不記得.
你和神是有約的呢.
其實.
你知不知道記不記得.
立約這件事.
在聖經裡面是非常重要.
說幾天都說不完.
但是我很想.
把握這個機會和大家重溫.
立約是一回什麼事.
其實約是一件.
非常美麗的故事.
在舊約聖經那裡.
說到.
約就是天上大而可畏的神.
祂給沒有機會的人.
一個重返.
回到祂的機會.
說到人類的始祖.
在罪裡面墮落的時候.
他們失敗了.
失敗了一直到如今.
當時他們被趕逐.
出伊甸園.
他們不能翻身.
他們只有在罪裡面.
在那些痛苦和破碎裡面.
折騰.
這個我們在舊約聖經那裡.

$^{601}$都看到.
但是就在這個時候.
慈愛的神給人翻身的機會.
這個機會.
就從立約開始.
最初我們知道.
神呼召一個在古近東的.
一個族長.
大家都認識的亞伯拉罕.
呼召他出來.
揀選他的族裔.
和他們立約.
後來稱他們為以色列.
在地上贖他的子民.
這是很重要的一件事.
因為.
神和這群子民的立約.
沒有機會.
以前是和我完全分割.
斷絕關係的了.
因為罪 但是現在我和你立約.
是建立一個特殊的關係.
那是什麼關係呢.
神用了當時.
近東的一個.
立約.
來表現這個關係.
就是宗主國和藩屬國.
宗主國簡單來說.
是強國.
可以吞併你的國家.
藩屬國就是弱國.
是被你吞併的國家.
當時有這樣的條約.
宗主國就是大亞哥.
保護小國.
神用了這個立約.
來表達他和人之間.
想建立的一個新的關係.
有同的地方.

$^{641}$有不同的地方.
就是神是宗主國.
我就是那個強國.
你們是弱國.
你們是無能 要人幫助.
要人保護的.
同的地方就是.
我會保護你.
但是你要聽我的話.
但是不同的地方就是.
世上的宗主國.
是以自己利益為主的.
就是蝦併吞.
蝦蝦霸霸吞併.
但是.
神和以色列所立的.
這個宗主.
和藩屬國的約.
是以外.
是以人的利益為出發點.
簡單來說 神這一方.
是祝福這個族裔.
後裔興旺.
好像天上的星.
地上的沙那麼多.
更加救世主會在.
你這個族裔那裡出來.
萬國.
都因這一位後裔.
而得福.
而人這一方 以色列做什麼呢.
不用做什麼.
他們只是受祝福.
他們不用付出什麼.
唯一的條件.
就是要忠誠.
忠誠於.
耶和華給你的.
典章 律例 誡命.
正如我們.

$^{681}$剛才所說的.
為什麼這麼重要呢.
因為神要他們遵守.
這些誡命的律例 典章.
要向世界聲明.
這班人是屬我的子民.
是聖潔的子民.
所以神很重視.
這件事.
而神也強調.
一件事.
我和你立約不單只是.
宗主 藩屬國.
更加是因此而有了.
一個前所未有的關係.
我要作你們的父.
你們要作我的子民.
是這樣的約.
神在當初已經.
和以色列納了這樣的約.
總之就是約.
是建基在.
耶和華的恩典上面.
因為人.
這一方是不配得.
也沒有付出什麼.
約也是建基在.
忠誠.
這兩個字非常重要.
Faithful 耶和華和子民.
都要忠誠.
耶和華對他忠誠 他要的子民.
都要忠誠 或者這樣說.
你就明白當時他們.
尼希米他們犯了什麼罪.
他是.
你是守施慈愛的神.
但是我們就沒有.
遵守你的.
律例 典章 誡命.

$^{721}$但是求你.
紀念.
你獨人摩西的說話.
我要跳過這裡.
因為時間比較緊張.
我想跟著和大家說.
這些立約和我們有什麼關係呢.
我們不是.
以色列人 更加不是.
二千多年前的以色列人.
如果我們這樣想的話.
我們就漏了一件很重要的東西.
就是神與人的立約.
不是停在以色列那裡.
而是你和我.
今天信了耶穌的人.
絕對是立約裡面.
重要的角色.
主耶穌基督雖然是生於以色列.
但是他用.
寶血來拯救的.
不只是以色列.
是拯救世界上每一個.
相信和跟隨他的人.
好像你和我一樣.
每個星期在這裡我們是守主餐.
記不記得Marshall經常讀的聖經.
就是我的血.
是為你們流出來的.
這個杯是我的血.
所納的新約.
是藥來的.
是為你們的罪而流出來的.
所以.
沒錯 因為耶穌基督.
為我們流出的血.
神給我們納了一個新的藥.
各位弟兄姊妹.
因著我們和父神.
和我們的神建立一個父與子的關係.

$^{761}$我們可以回到祂那裡.
有意思的就是.
今天我們.
如何守那些誡命.
還記不記得耶穌說.
誡命的總綱就是.
你要盡心盡性盡義盡力.
愛主你的神.
其次就是愛人如己.
愛你的靈性如同自己.
神就將很闊.
很艱深的誡命.
做了一個總綱.
今天我們有沒有活出.
愛神愛人的生命.
就是我們有沒有忠誠地.
去守這個藥.
所以我們在這裡.
找回之前.
黎英米禱告帶來的幾個靈靈聖策.
第一個提醒我們.
我們要謙卑地給予神.
應當有的尊崇和敬畏.
第二 我們要向神.
闖開我們的修脈.
承認自己的罪.
第三 我們要忠誠守神.
和我們納藥.
我們以為這段聖經說完了.
但是原來還有一件事.
還有一節聖經我們不可以漏.
還有一節是很重要的.
唉.
又哀了.
主啊 求你.
則以聽你獨人.
我黎希米的祈禱.
聽喜愛敬畏.
為名眾僕人的祈禱.
有沒有留意.

$^{801}$這裡 最後的這一句.
黎希米說的就是.
不單是他一個人祈禱.
原來還有很多人祈禱.
究竟有多少人.
我不知道.
因為聖經沒有說清楚.
但是.
重要的就是.
有很多人.
在這裡祈禱.
這群是忠誠死守.
至死.
都要為主忠心的人.
你有沒有很奇怪.
為什麼一波消亡.
一上場.
就可以讓他們回歸呢.
沒錯 是懷柔政策.
但是你有沒有想過.
是神聽了很多忠心的人的祈禱.
以致讓他們回歸呢.
藉著波消亡這個政策.
讓以色列回歸呢.
只可惜回歸之後.
他們當時仍然面對.
受靈辱 城牆被毀.
城門被燒 這些事還沒有搞定.
可能.
神仍然.
繼續要聽很多人.
堅守忠誠的祈禱.
今天你和我.
不是在亡國被擄的時候.
丁家.
但是今天你面對.
這個政治環境.
或者神的子民.
的群體 好像城會教會那樣.
或者你會看到很多問題.

$^{841}$你都會好像黎希米那樣.
嗨 圓 又嗨.
但是不要忘記.
在嗨裡面.
你要忠誠.
為主忠心.
洗手去祈禱.
用禱告來緊守到最後一刻.
你有沒有這樣做呢.
黎希米的.
禱告裡面.
最後告訴我們 有一班人.
你願意是其中一個呢.
說到這裡.
我很想和大家分享.
我年輕的時候.
忠誠洗手的.
立志和經驗.
雖然時間不夠.
我也很想說完 幾分鐘.
我記得我十幾歲的時候.
我信了主.
其實我那時候.
信主的時候.
我是一個迷洋.
我已經信了主.
離開了教會.
但是有一晚的報道會裡面.
神的話淋到我身上.
我決定回到我以前離開了很久的教會.
但是很有趣.
我回到那個教會的時候.
教會並不是我想像中那樣的.
為什麼呢.
施班子輝和大家練歌的時候.
我坐在那裡.
後面我聽到不斷有人.
在那裡說施班子輝的壞話.
不斷批評他.
說他悶 戴得不好.

$^{881}$穿什麼衣服 什麼都在吵.
我從來沒有聽過.
這麼惡意的批評.
不要怕 施班子輝真的有很多.
前美.
當時我真的聽到很多這樣的.
這樣也有的 教會.
還沒有說完.
主的崇拜開始的時候.
我練了詩就要獻詩.
十點鐘開始主的崇拜.
我坐在那裡 沒人來.
只有我一個人.
今天我獨奏嗎.
施班練歌的時候.
十個八個人.
主的崇拜一開始.
只有我一個人坐在那裡 我真的嚇得發抖.
誰知道十點多的時候.
有幾個又回來.
總之三五個人每一次都有獻詩.
就是這樣的教會.
就是這樣的地方.
還有我們執事會開會.
開會的時候.
坐在那裡的那班生意人.
不是生意人不好.
但是那班比較世俗.
通俗一點的生意人.
就是一邊開會一邊吸煙.
這樣.
教會裡面.
沒有什麼委身的屬靈領袖.
茶經 主學.
都是牧師一個人帶的.
而跟我當時.
決志信主 回教會的.
有另外一位同學 跟我差不多一代.
他也早就走了.
他已經走了成心.

$^{921}$去了一個比較好的大教會.
但是我當時.
決定.
我仍然的死守在這裡.
牧師雖然講到都很悶.
但是這位牧師很有愛心.
他默默耕耘地在教會.
很感動我.
我也見到有幾位.
默默事奉.
很忠心事奉的信徒.
包括那個被人批評的指揮.
我見到.
這樣的榜樣.
我更加看到聖經那裡.
講到耶穌基督為我們的罪.
忠心的死守.
他用禱告和行動.
在世界上.
為我受苦受死.
使我可以.
從罪的危機得到轉機.
是這幾樣東西.
令到我決心死守.
就在這段日子裡面.
我留下在這個教會.
轉眼後來成為.
這個教會其中一個年輕的領袖.
看起來.
神.
就是這樣用.
我們的忠心為他死守.
各位弟兄姊妹.
你有沒有想過.
你願意為神忠心死守呢.
即使你面對面前.
你覺得有很多問題.
你覺得有很多不滿.
但是你願不願意為神忠心死守呢.
這個答案是很重要.

$^{961}$是神很在乎.
不是神只是問.
黎希米和問我.
問大家每一位.
我們祈禱.
藉著黎希米的禱告.
帶給我們靈敏的醒察.
我們真的要敬畏.
你是大而可畏的神.
給你應有的尊重.
我們真的要重視.
我們在生命裡面犯罪.
得罪你.
我們要事事的.
真誠的去求你赦免.
even before.
我們去求任何的事情.
開始.
我們要忠誠的.
謹守和你納個約.
而且是至死忠心.
主我們願意這樣.
當我們做不到的時候.
求你讓事事的提醒我們.
教我們真的能夠活出.
是能夠.
令你得到榮耀.
的一群的子民.
奉救主耶穌的聖名.
我們禱告.
\newpage



\section{1:3-14}
\label{sec:1Q8_M16V0ZI}
\textbf{We Are Worshippers (Ephesians 1\_3-14) - Sam Johnson}
\newline
\newline
連結: \href{https://youtube.com/watch?v=1Q8_M16V0ZI}{\texttt{https://youtube.com/watch?v=1Q8\_M16V0ZI}} ~~~~ 語音日期: 2023-09-19
\newline
\newline
\hyperref[sec:p__HCvfoaaM]{\small{< < < PREV SERMON < < <}}
~
\hyperref[sec:index_chronic]{\small{[返順時目]}}
~
\hyperref[sec:index_scriptual]{\small{[返順卷目]}}
~
\hyperref[sec:Z86XRROS_74]{\small{> > > NEXT SERMON > > >}}
\newline
\newline
$^{1}$I'll give real briefly a little bit of background.
of this passage here..
This is the letter that Paul wrote to the Ephesian church.
and to the churches in the region of Ephesus..
And if we think about that church.
and what was going on in that region,.
he was writing to a group of believers.
that were in an urban context..
They were a group of believers that were in a pagan context..
Ephesus was, throughout the known world at that time,.
was known for the temples that were there..
And so it was a place where people were.
practicing a lot of unhealthy and sinful things.
that were not in line with God's word and God's truth..
This is the type of church that Paul is writing to here..
So obviously, he's writing to us..
Those are things that we can identify with,.
being in a context like that..
And he begins, very interestingly enough,.
in this letter, with worship..
So we're gonna look at this passage in verses three to 14..
Blessed be the God and Father of our Lord Jesus,.
who has blessed us in Christ with every spiritual blessing.
in the heavenly places, even as he chose us in him.
before the foundation of the world,.
that we should be holy and blameless before him..
In love, he predestined us for adoption to himself.
as sons through Jesus Christ,.
according to the purpose of his will,.
to the praise of his glorious grace,.
with which he has blessed us in the beloved..
In him, we have redemption through his blood,.
the forgiveness of our trespasses,.
according to the riches of his grace,.
which he lavished upon us in all wisdom and insight,.
making known to us the mystery of his will,.
according to his purpose, which he set forth in Christ..
As a plan for the fullness of time,.
to unite all things in him,.
things in heaven and things on earth..

$^{41}$In him, we have obtained an inheritance,.
having been predestined according to the purpose of him,.
who works all things according to the counsel of his will,.
so that we who were the first to hope in Christ.
might be to the praise of his glory..
In him, you also, when you heard the word of truth,.
the gospel of your salvation, and believed in him,.
were sealed with the promised Holy Spirit,.
who is the guarantee of our inheritance,.
until we acquire possession of it.
to the praise of his glory..
So this letter begins with this expression of praise..
It follows the form of a Hebrew barakah, or blessing,.
and it's different than the way that Paul.
normally begins his letters..
He normally begins with kind of a greeting,.
a thanksgiving, and a prayer..
And one of the things I like about this opening.
is it's kind of an outburst of praise..
It seems like Paul began to give thanksgiving.
and praise to God, and then he just couldn't stop,.
so he just kept going..
It's made of one really long sentence of 202 words..
It's just an outburst of praise..
And it begins with the word blessed,.
to denote the exaltation of God,.
and then proceeds to focus on all that God.
has done through Christ, and the blessings.
that we as believers have received in Christ,.
and God's plan for the future in Christ..
And some theologians believe that this may have functioned.
as an early Trinitarian hymn, that the early church.
would have used it in that way..
But this expression of praise also sets the tone.
for the entire rest of the letter..
And in this letter, Paul is speaking to the church.
about the church, about what it's like.
to have a healthy church, and all these things.
that are so important to us, and he sets the tone.
for the entire letter with this opening.

$^{81}$of verbal praise and worship..
And I think there's a lot for us to learn here..
There's kind of a template for us to think about worship,.
and about verbal praise and worship,.
and about how important it is,.
and how essential it is in our lives..
It's an essential part of our gathering.
when we get together like this,.
but it should also be a part of just the rhythm.
of our daily lives..
We should be praising and worshiping God verbally..
It's an essential part of our lives as believers..
And I'm using the term verbal praise and worship.
intentionally to kind of make a distinction there,.
because we know Romans 12 one talks about the fact.
that our whole life is supposed to be worshiped to God,.
and that is certainly true,.
and that's how we should be living..
We should be worshiping God in our work,.
in our relationships, in every part of our life..
But this morning we're specifically talking.
about verbal praise and worship,.
when we speak or sing praise to God..
And it's clear in this passage that that's supposed.
to be a part of who we are as believers..
It's supposed to define us in a lot of ways..
I'll also be using the term genuine worship.
just to remind us that our verbal praise and worship.
can be hollow and insincere..
Just because we're verbally praising God.
doesn't mean it's genuine,.
doesn't mean it's coming from our heart..
And we know that that's actually one of the things.
that Jesus often chastised the Pharisees about,.
because they were observing God's law.
and doing things externally that God had asked them to do,.
but it wasn't genuine, it wasn't coming from their heart..
And so in the church we desire to genuinely worship God.
from a place of spirit and truth.
that the scripture talks about..

$^{121}$And Paul gives us a really beautiful template here.
for how to think about what that looks like in our lives,.
what it looks like in our gatherings..
The first thing we see here is that we worship.
and praise God..
Now that seems like a very obvious point, right?.
But it's actually important for us to think about..
Who are we worshiping?.
And Paul makes it clear in this passage.
that we are worshiping the almighty creator God of the Bible..
He starts talking about that right away..
This expression of praise is Trinitarian..
So look at verses three, he talks about God.
and the Father of the Lord Jesus Christ..
In verse five, he's talking about.
what God has done through Jesus..
In verse 13, he's talking about the Holy Spirit.
that has sealed us..
So he's making it very clear who he's talking about.
and what God he's talking about..
The Trinitarian God of the Bible..
This worship is gospel-focused..
Verses five to 12 are all about what God has done.
through the gospel, through the good news of Jesus,.
that Jesus came and by his grace,.
he saved us, we've been redeemed,.
that he has given us forgiveness in verse seven..
And in verse eight, he has lavished on us.
that forgiveness in all wisdom and insight..
He's uniting all things in Christ..
It's really focused on the gospel..
That's a great template for us to think about.
in our worship..
And one of the things we see here in this expression.
of praise that Paul gives us is that genuine worship.
is not ambiguous about who God is..
And that's actually important in our culture..
We say, okay, we praise and worship God, that's obvious..
We don't even have to talk about that..
But in our culture, sometimes we do.

$^{161}$because people worship a lot of things.
that are not the God that Paul is talking about here..
And not just the obvious things like materialism.
or money or power or success..
Yes, people worship those things and orient their lives.
towards those things and follow after those things..
But even at times in the church,.
we can begin to make God in our own image,.
to try to come to God in worship on our terms..
And we can't do that..
God is a God of grace and mercy and love..
But he is God..
So when we approach him in worship,.
we approach him on his terms..
And we worship God for who he is,.
not for who we want him to be..
And we worship God for who he has revealed himself to be..
And we find that in his word..
So if we are worshiping a God or a belief in our mind.
that is not God as he has revealed himself to be,.
that is just God as we want him to be.
or God as we feel like is comfortable.
or palatable for us culturally,.
or we say, well, this makes sense to the people around me,.
so this is how I think God is,.
then we're not really worshiping God.
in the way that he calls us to worship him..
No God's made in our own image..
That's one of the 10 commandments, right?.
And that's one of the first cornerstones of genuine worship..
If we're gonna worship and praise God,.
we need to make sure that we are worshiping him.
for who he is as he's revealed himself to be..
The next thing we see here is that we worship and praise God.
because he is worthy..
And there's two aspects of God's worthiness.
that are really clear in this passage..
The first is God's majestic greatness..
Paul points us towards that..
He talks about God's eternal nature in verse four,.

$^{201}$how God chose us before the foundation of the world..
He talks about God's love for us in verse four,.
that what he did was in love..
He talks about God's omnipotent will in verse five,.
that God is all-powerful,.
that his purposes in the world cannot be thwarted.
and will not be thwarted..
He talks about God's wisdom in verse eight..
He talks about God's redemptive plan.
to heal all brokenness in verse 10..
That's a beautiful expression of God's greatness..
He talks about God's goodness in verse three,.
in verse six, in verse eight, in verse 11,.
verse 13, verse 14, throughout the passages,.
talking about God's goodness..
And when we genuinely encounter God's greatness,.
then worship becomes spontaneous..
We don't have to convince ourselves to worship God..
When we have an encounter with God's power,.
it's kind of like if you think about seeing something.
that's just kind of really awe-inspiring or amazing.
if you stand next to the Grand Canyon.
or you go to Niagara Falls or something,.
and it's not hard to just feel excited about it..
It's not hard to turn to the person next to you.
and say, "Wow, I can't believe this.".
You know, it's not hard to have something to say about it..
It's spontaneous..
And that's true when we are beholding God.
and really encountering God's greatness..
It's gonna come out..
It's spontaneous..
We're gonna begin to worship..
These attributes that Paul describes in this passage.
of who God is, those things never change,.
regardless of our circumstances,.
regardless of what's going on in our world..
Those things never change..
And that means that God is always worthy of our worship,.
every day of our life, no matter what's happening..

$^{241}$Even when we're going through really difficult times,.
God is worthy of our worship..
Even when the world is falling apart and things are crazy.
and we're frightened, God is worthy of our worship..
Even when we feel pain, God is worthy of our worship..
And we know that we have that testimony in scripture.
of the saints that worship God through lament.
as well as through joyful praise..
So verbal praise and worship doesn't just have to be.
a joyful expression..
It can also be an expression of lament..
But we can always worship God,.
regardless of what we are facing in life..
And we're called to do that..
And if we're gonna do that,.
if we're gonna encounter God's greatness,.
even when we're going through difficult things,.
we have to take our eyes off of our circumstances.
and begin to look at God,.
begin to orient ourselves towards God..
And one of the things that's true in my life,.
and I think it's true for a lot of us as believers,.
is that if we're struggling to worship,.
if we find ourselves in a place where we're saying,.
I just don't feel it, I don't feel in church,.
I don't feel it anywhere else, it's just not there..
Yeah, I can say the words, but it's not meaningful to me..
If we're struggling to worship,.
it's often because we're not seeing God's greatness..
We don't see how great God is..
We don't see how powerful he is..
We're not looking at that vista that's in front of us..
Because like I said, if we're seeing that,.
then it's easy for us to begin to worship..
The other aspect of God's worthiness that we see here.
is God's amazing grace..
Throughout this passage, it's talking about God's love,.
God's lavish forgiveness, his grace, his mercy,.
his wise and holistic and good plan for our good,.
for our physical, our spiritual, emotional good,.

$^{281}$and not just for us personally,.
but for the redemption of all things..
How God is gonna unite all things in Christ..
God's plan to glorify himself through his people,.
through us, God's gonna use us as a part of his plan,.
God's destiny for our future..
And if we're in Christ, if we have a relationship with him,.
that means that we are experiencing God's grace..
We can't say, well, I know God's got a grace,.
but there's no grace in my life..
If you are in Christ, if you have a relationship with him,.
you already have experienced his grace.
and you are experienced it moment by moment, day by day..
And so that's something that God is worthy..
We can praise him because of that..
When we encounter the depth of God's grace,.
we start to see how much we need God..
And we see how, we become aware.
of how much he's done in our lives..
And that's a part of seeing God and his worthiness.
and who he is, is saying, okay, this is the reality of me,.
this is the reality of who I am,.
both my sin, all my mess, all my junk,.
but also just the ways I fall short,.
the ways I'm not who I wanna be,.
the ways that I am not gonna get where I want to go..
I need God so much..
And when we see that, we begin to say,.
wow, God has given me so much grace..
God has forgiven me..
God has changed me..
God is at work in my life..
God is filling in all these gaps..
We can look around the world.
and see others who are struggling,.
and instead of having a heart that says,.
well, I'm just better than that person,.
so I can feel good about myself,.
I can feel proud about that,.
we can look at someone else and say,.

$^{321}$you know what, I would probably be in the exact same place.
if it wasn't for God's grace and God's work in my life,.
and that can cause us to worship..
Honestly, we don't know what we don't know.
about our destiny apart from God's grace..
It's true of all of us..
We really don't know where we would be without God's grace.
if God hadn't saved us, forgiven us,.
brought us into his family,.
began to transform us and change our desires.
to honor Christ..
If God wasn't doing that,.
we really have no idea where we'd be..
And so again, I'll say that if we're struggling to worship,.
a lot of times it's because we're not seeing.
the depth of our need for God.
and the depth of his grace for us..
If we really are seeing how much God has forgiven us.
and how much we need him.
and how much he is at work in our life day by day.
and all of the wonderful things he is doing for us,.
then we're gonna begin to worship..
So what we see in both of those areas,.
looking at God's greatness.
and looking at God's grace in our lives,.
is we see that if we will look at God and who he is.
and behold him, then we'll begin to genuinely worship..
Next thing we see in this passage.
is that we worship and praise because it is good for us..
It's so healthy and good for us.
to take our focus off of our circumstances.
and our mess and our struggles and our challenges.
and turn towards God and look at who he is.
and see his greatness and see his grace..
It's so healthy for us..
Psalm 147, the first three verses say this..
Praise the Lord for it is good to sing praises to our God.
for it is pleasant and a song of praise is fitting..
The Lord builds up Jerusalem..
He gathers the outcast of Israel..

$^{361}$He heals the brokenhearted and binds up their wounds..
What do we see in that passage?.
We see people who are brokenhearted,.
who have wounds, who are outcasts,.
people who, things are difficult in their lives,.
and it says it is good for them to sing praises to our God..
It's pleasant, a song of praise is fitting..
Ephesians chapter five, later on in this letter.
that Paul has written, he tells the believers.
to address one another in psalms and hymns.
and spiritual songs, singing and making melody.
to the Lord with your heart, giving thanks always.
and for everything to God the Father.
in the name of our Lord Jesus Christ..
It's so healthy and good for us to praise God.
and to do that in song like we just did,.
but also just to do it in speaking with one another.
and encouraging one another, talking about God's goodness.
in our lives, talking about how we see God at work.
in someone else's life, how we're grateful.
for what God is doing in their life,.
how we're grateful for the way God is using them.
or how we are seeing God taking them through something.
or supporting them in some way..
Those are all ways that we can express worship to God.
and it encourages one another and it encourages us..
When we're in some of our most difficult moments.
and life is terrible and we feel like everything.
is going wrong and we've lost our job.
or we have experienced a relationship.
where someone has said terrible things about us.
or maybe we've experienced a loss,.
someone that we love has passed away,.
whatever we're going through, even in those moments,.
if we're in Christ, then we are loved,.
we're redeemed, we're forgiven, we're accepted.
and we have a future hope, we have a blessed inheritance.
like this passage talks about and we have.
an eternal destiny of victory and of fellowship.
with God in heaven..

$^{401}$That's the truth about everyone in this room..
If you're in Christ, that's the truth about you..
And so even if you're going through something.
really difficult, I'm not talking about saying,.
well, I can just ignore my problems and wish them away..
No, life is challenging, it's hard and that's real..
And there are times when we feel that brokenness.
so heavily and we're discouraged..
I'm not saying that we don't..
As believers, we feel that..
That's why, like I said earlier,.
there are all the scriptures of lament.
because life comes in on us just like anybody else..
But even in those moments, we know that we are loved.
and forgiven, we're redeemed, we're accepted,.
that God is at work in our life for his glory,.
for our good, that we have a future hope,.
that we have this destiny of victory.
and fellowship with God in heaven..
And we can worship and it's good for us to think about that..
In the darkest moments of our life,.
when we're really struggling, it's good for us to say,.
yes, this is hard and it hurts and I'm being crushed by it,.
but God loves me and I'm forgiven and he's redeemed me.
and even all my mess, he's at work.
and he's making something good out of it..
He's transforming me and I'm not gonna be.
in this place forever..
God's gonna change me into something new.
and that's gonna happen in this life,.
but ultimately, God's gonna take me on to heaven.
in the future..
And so even the pain that I'm feeling right now.
is a temporary pain..
It's good for us to know that..
It's good for us to say that..
It's good for us to say that to one another..
It's good for us to say that to ourselves..
And we do that when we worship..
The songs we were singing this morning,.

$^{441}$singing about how we have a savior who's mighty to save..
That's what we're doing..
We're saying, hey, life is tough,.
but Jesus is bigger and he's powerful and he's mighty.
and he can save us out of all the challenges of our life..
It's good for us to say that to each other..
It's healthy..
It encourages us..
It keeps us going..
Worship and praise is so good for us emotionally..
It's so good for us physically..
It's so good for us spiritually..
And then the last thing we see here.
is that we worship and praise.
because it's a witness to others..
Some of the earliest second century texts.
talking about the early church.
agree that the worship gatherings of the early church.
that happened on Sunday were first and foremost.
what they called Eucharistia or Thanksgiving..
That was what defined those gatherings.
was that they were about thanksgiving to God..
Worship and praise to God and talking about who God is.
and how good he is and exalting him.
and talking about his work in our lives.
has always been a defining characteristic.
of the church of Christians from the very, very beginning..
In the second century,.
there was this pagan writer named Celsus.
who didn't like Christians..
He was a critic of the church..
He was writing to someone else.
and he commented in this letter.
that he didn't like Christian worship.
because the chants and the songs that they sang.
were unfamiliar to his ears,.
but they were also so beautiful.
that he felt that they dulled his reasoning faculties.
and made him more open to the Christian arguments..
From the very beginning,.

$^{481}$we as a people, followers of Jesus,.
have been known as people who worship.
and sing and sing beautifully..
That doesn't mean that we all necessarily.
know how to play instruments or even sing on key,.
but it's beautiful because it comes from our hearts.
and it's genuine and it's worship to God..
That's been a part of who we are from the very beginning.
and it's a part of who we are now..
In Acts 2, verses 46 and 47, it says this..
Day by day, attending the temple together.
and breaking bread in their homes,.
they received their food with glad and generous heart,.
praising God and having favor with the people.
and the Lord added to their number day by day.
those who were being saved..
That was the very first church..
Every day, they were getting together,.
doing different things and one of the things.
they were doing was praising God..
And one of the powerful things about worship.
is that genuine worship and praise to God,.
it can't be duplicated or counterfeited in the world..
We see that in American culture..
We have this beautiful tradition of gospel music.
in the US that comes from the African American church..
Well, there's a difference between someone.
singing that song in church.
and someone just singing it on the stage.
at an award show who doesn't know Jesus..
It's two totally different things..
You can't counterfeit genuine worship.
and the world around us notices that too..
It's a witness..
When we are genuinely worshiping God,.
they say that's something beautiful.
and that's an expression that you don't find anywhere else..
It can't be duplicated..
And genuine worship and praise to God.
is always focused on God, not on us.

$^{521}$and that's another reason it's a witness..
Genuine worship and praise to God.
is some form of saying, listen, this is who I was.
and then this is who God made me or is making me to be..
This is my present situation.
but this is what God has for me..
This is what my future used to be.
but now because of Jesus, this is what is in my future..
Those types of statements define our worship.
where we're saying, look, it's about God..
God did this..
It's not about me..
The focus is on him and how good he is.
and how great he is and all that he's done in my life..
That's what we see in these verses three to 14.
of Ephesians one..
And that type of worship is always a witness to other people..
The people around us, they may think it's strange,.
they may think it's unusual but they'll notice it.
and people will say, that's different..
What are you talking about?.
Why are you talking about it?.
Why are you joyful?.
Why are you talking about your God.
even when life is terrible and I can see that it's terrible.
and you're going through a time of lament.
but you're still worshiping in that?.
Those things are all a witness to the people around us..
So I wanna close out with application for us..
Like I said, this passage gives us.
this beautiful expression of worship..
It's a great template for us to think about.
and it makes it clear that this should be a part.
of our lifestyle and clearly in the church,.
just like we worship this morning,.
but it should be a part of our lifestyle every day..
We should verbally be praising God,.
talking to ourselves, nothing wrong with talking to yourself.
if you're gonna preach the gospel.
or talk about who God is to yourself..

$^{561}$That's a healthy, good thing..
We should be doing that,.
especially when we're going through difficult times..
We should be telling ourselves, you know what?.
It's hard but God's at work in my life..
God loves me, I'm forgiven, I'm redeemed..
That should be a part of our life rhythm..
But I know it's not always easy..
So how do we do that?.
How do we make that happen?.
I want us to think about something for just a minute..
To illustrate this, we can take a mental excursion real fast.
so let's pretend, you know, if you go down Halsted here.
to 31st and then you just go east on 31st,.
as far east as you can go, you get to the lake.
and there's the beach park there, right?.
Everybody been to 31st Street Beach?.
And it's a really pretty beach park..
There's that point out there where the marina is..
You can walk out there and they have kind of.
these big boulders for breakwater..
You can look back towards the city..
If you haven't done it, you should do it sometime..
It's a really, really cool spot..
It's a great spot to take pictures of the skyline..
But if you go out there and if I asked each one of you.
to describe what you would see when you're there,.
your answers would vary depending on how you were oriented..
If you're looking straight east from out there,.
all you're gonna see is water..
You're just gonna see the lake, which is beautiful..
That's all you're gonna see..
If you're looking back to the west,.
then you're gonna see the city skyline..
You're gonna see Lakeshore Drive..
You're gonna see something totally, totally different.
depending on how you're oriented..
And I think that's really instructive for us in worship,.
especially when we're struggling to worship..
We're saying, I know I need more worship in my life..

$^{601}$I wanna express praise to God,.
but it's not coming naturally right now..
I don't feel it..
We need to think about what are we oriented towards..
And sometimes in life, we need to turn the eyes of our soul.
away from our problems, away from our success,.
away from our failures, away from our relationships,.
away from our entertainment,.
away from all our social media feeds,.
and say, I'm gonna orient myself towards God..
I'm gonna look at God..
I'm gonna behold God..
And we do that as we look in His word..
We do that as we look into our hearts.
and recount God's grace in our life..
And we say, you know what?.
This is what God has done for me,.
and this is what God has done in the past..
Or we look at the lives of other believers around us.
and say, I see God doing this in someone else's life,.
and it's a beautiful thing..
I see something encouraging happening.
with a brother or sister..
When we do those things, we begin to orient ourselves.
towards looking at God..
And when we look at God, then we can begin to worship..
The scripture talks about this..
Psalm 75, one says, we give thanks to you, O God..
We give thanks for your name is near..
We recount your wondrous deeds..
Psalm 124, one says, if it had not been the Lord.
who was on our side, let Israel now say..
And so I would challenge and encourage all of us.
to do that and to seek ways to do that in our lives..
Begin to orient our soul towards looking at God..
There's a lot of stuff for us to look at in the world..
Way too much..
We can't take it all in..
And sometimes we need to turn away from those things.
and orient ourselves towards God..

$^{641}$And if we'll do that, it'll be so good for us..
It'll be life for us..
It'll be life-giving..
It'll be a witness to those around us..
And ultimately, it's an identity that God has given you.
as one of his children, is to worship him as a worshiper..
That's who we are..
We are worshipers..
Let's pray..
God, we thank you for your greatness..
We thank you for your grace..
We thank you for who you are,.
all that you have done in redemption..
We thank you for all that you have done in our lives,.
and we are undeserving..
We need you desperately, God..
And God, you love us and have redeemed us and forgiven us..
I pray that you would help all of us.
to orient our souls to look at that,.
to look at you, to look at what you've done for us,.
to look at your grace, and may worship come out of us..
In Jesus' name, amen..
(gentle piano music).
(piano music).
(gentle music).
\newpage



\section{創世記}
\label{sec:Z86XRROS_74}
\textbf{上主的護理 (1) 細想約瑟的故事 1 (創世記 37) - 陳崇基牧師}
\newline
\newline
連結: \href{https://youtube.com/watch?v=Z86XRROS-74}{\texttt{https://youtube.com/watch?v=Z86XRROS-74}} ~~~~ 語音日期: 2023-09-19
\newline
\newline
\hyperref[sec:1Q8_M16V0ZI]{\small{< < < PREV SERMON < < <}}
~
\hyperref[sec:index_chronic]{\small{[返順時目]}}
~
\hyperref[sec:index_scriptual]{\small{[返順卷目]}}
~
\hyperref[sec:_iX85FoHSZc]{\small{> > > NEXT SERMON > > >}}
\newline
\newline
$^{1}$為我們講道 現在將時間交給他.
早晨 各位姐妹 今天下雨.
希望大家來到沒有問題.
我們開始一個新的系列.
叫做Providence 上主的護理.
God's Providence.
我們會看很多的聖經.
特別是看創世記約瑟的故事.
Joseph的Story.
希望這個系列能夠幫助大家.
思想一下我們人生每一個階段.
我們經過的每一步.
很重要 聖經裡面有一個很寶貴的教導.
我們將它囊括起來.
就叫做Providence.
Providence是什麼意思呢.
Pro-Videre.
Pro 即是說.
Ahead的意思.
預先 之前 事先.
Videre.
To see的意思.
拉丁文來的.
Seeing Vision.
這個教義 或者聖經裡面的教導.
告訴我們上主 上帝.
是去到看著 預先.
看到我們走的每一步.
這個教義是一個很寶貴的教導.
告訴我們上帝不單是創造萬有.
他也是護理著 去到take care著.
所有的一切.
天下的萬無 whatsoever.
所有都在上帝的視覺裡面.
或者在上帝自己的管理裡面.
這個就叫做Providence.
所以我們去思想一下.
這個topic和我們有什麼關係.
神的護理 上帝的護理.
很多時候舊有的翻譯 中文的翻譯.

$^{41}$Providence翻譯成天明.
有沒有聽過 天明.
因為有一點像我們華人.
有一個很重的意識.
就是天明 宿命.
就是每一件事都set死了.
等於我們廣東話會說什麼.
整定.
什麼叫整定.
吃多少 穿多少 是怎樣.
整定的.
吃多少 穿多少 整定的.
就是每一件事都安排了.
命運裡面安排了等等.
比如說緣份的事不可以強求.
為什麼.
整定了.
你急都沒用.
整定的就有些事.
有一種這樣的概念.
或者我們廣東話說.
arm arm .
arm arm 這個詞語很好用.
arm arm .
就是這樣arm arm .
arm arm 可以說.
少少 數量.
arm arm 到位.
或者天下萬無.
你所遇見的事.
arm arm 善 arm arm .
不過我想說說.
聖經裡面所說的.
跟華人的傳統的思想.
有些不同.
因為我們說整定 命運.
命生成是這樣.
等等.
跟聖經裡面所說的不同.
因為命運不會跟你說話.

$^{81}$是不是.
整定的事.
或者所謂的命運.
是不會跟你有任何互動的.
但是聖經裡面說.
上主的護理就不同了.
上主是跟這個世界會有互動的.
會大家.
有彼此的互動.
上帝是會參與在世界裡面.
並且參與在你生命裡面.
對不對.
有時候直接.
有時候是間接的.
有時候是明明的.
有時候呢.
或者更多時候是暗地.
無論是明的 或者暗的.
聖經裡面都去到.
很直接的告訴我們.
我們的生命是在上帝自己的看顧裡面.
在上帝自己的.
providence裡面.
聖經裡面當然有些直接的.
教導 直接的經文.
說到這件事.
詩篇139篇 大家很熟了.
我認識 我坐下起來.
我休息的時候.
我做工的時候 起來.
你都曉得.
從遠處你見到.
知道我的意念.
我走路忙的時候.
我得閒的時候.
休息的時候 你都細察.
你也心知.
我一切.
whatsoever 我所行的.
這些是很直接的經文.

$^{121}$我們不會翻開這些經文.
為什麼呢.
因為聖經文化 希伯來文化.
甚至到今天都是.
很少會像我們.
比較西方.
現代人這樣想的.
我們就想point for.
1234 bullet point.
什麼叫上主護理.
1234這樣列出來.
將它抽空 或者將它抽象地.
有些概念.
去表達出來.
希伯來文化 聖經文化不是這樣的.
到今時今日都是一樣.
拉比 你問他一些東西.
怎樣解釋這個概念.
他會怎樣回答你.
Let me tell you a story.
聖經裡面都有.
這些東西.
先知拿丹去指責.
大衛 君王犯罪的時候.
他怎樣做的.
他不是一棍子.
這樣下去.
他講一個story給他聽.
講一個比喻給他聽.
耶穌的教導當然也是這樣.
所以是一種story telling的culture.
講一個故事給你聽.
讓一個故事.
讓你能夠咀嚼一下.
讓你明白.
裡面所講的一些很重要的.
很寶貴的教導.
很重要的一些概念.
所以為什麼我們做父母.
特別小孩子小的時候.

$^{161}$要跟他讀那些.
兒童故事 讀那些好的故事.
因為我們每個人都是tell a story.
hear a story.
learn from a story.
所以在這個系列裡面.
我們要多看一點聖經 看一點經文.
這些經文都跟故事有關.
所以今天我也跟你講一個故事.
就是約瑟的故事.
37章我們開始講.
就是約瑟的故事開始的時候.
這裡講什麼呢.
他說這個是雅各的事蹟.
約瑟十七歲和他哥哥們一同牧羊.
他是個少年.
父親的切丕拉和色柏的兒子們常在一起.
約瑟把他們的惡行報給父親.
以色列愛約瑟過於其他兒子.
因為約瑟是他年老時候生的.
他就給約瑟造了一件長袍.
這個是新的翻譯.
因為我們知道七十四本書說是彩衣.
希伯文本身就是彩袍的意思.
長的袍 總之就是漂亮的袍.
總之他有名牌的東西穿.
他哥哥就沒有.
全部名牌的東西都買在他身上.
第四節.
哥哥們見父親愛約瑟過於他們就恨約瑟.
不與他說友善的話.
然後經文下一節就講到兩個夢.
記得嗎.
有兩個夢.
約瑟做了兩個夢.
第一個夢是有關收割的時候的和菌.
把那些東西摘起.
然後他說全部哥哥的和菌就向他下拜.
第二個夢就講到太陽,月亮和十一顆星星.
就向他下拜.

$^{201}$對嗎 兩次.
很多時候兩次都是重要的東西.
所以約瑟整個故事裡面.
很多時候夢境都出兩次.
從不同的人都好.
總之是兩個夢.
兩個dreams.
這些dreams很多時候在以前的社會.
或者以前的文化.
以前的語言.
是帶著一種預言性在那裡.
就是預計將來會發生什麼事.
當然我們知道.
我們已經看完整個約瑟的故事.
或者聽過.
37至50章.
最後約瑟結束的時候.
他雖然經過很多苦.
但最後他成為一個很重要的人物.
救了他的家庭.
這個就是開頭.
如何去介紹約瑟這個故事.
14節.
以色列就是雅各.
對他說.
對約瑟說.
你去看看你哥哥是否平安.
羊群是否平安.
栽糊內告訴我.
於是他派約瑟出希伯來谷.
在南邊.
約瑟就往士劍去了.
在北邊.
由南走到北.
有人遇見他.
看他在田野走迷了路.
那人問他說.
你找什麼.
他說我找我哥哥.
請告訴我他們在那裡放羊.

$^{241}$那人說.
他們已經離開這裡走了.
我聽到他們這樣說.
我們往多丹去.
多丹在士劍再北一點.
約瑟就追趕他哥哥.
在多丹找到他們.
接著發生什麼事.
18節.
哥哥遠遠看見他的時候.
趁他還沒有就緊的時候.
就圖謀地去謀殺他.
他們彼此說.
看那做夢的人來了.
The dreamer.
現在來把我們把他殺了.
丟在一個坑裡面.
就說有猛獸把他吃掉.
我們且看他的夢將來會不會實現.
裡面聽見.
裡面是大兒子.
12個兒子.
裡面聽見就要救約瑟脫離他們的手.
說你們不可以害他的性命.
裡面又對他們說.
不可流他的血.
可以把他丟在這礦的坑裡面.
不可以下手害他.
裡面這樣說是什麼意思.
是想救約瑟脫離他其他哥哥的手.
將他交還給他父親.
約瑟終於到了.
23節.
到了他哥哥們的那裡.
他們就削去他的外衣.
就是他身上那件長袍.
他們抓住他.
把他丟在坑裡.
那坑是空的.
裡面會有誰.

$^{281}$你覺得那些兄弟想怎樣.
他們想將他謀殺.
不是純粹將他教訓一下.
因為憎恨到他入骨.
來到的時候就想將他殺死.
然後將他的屍體丟在坑裡面.
坑裡面沒有誰.
或者就算裡面大兒子Rubin.
想去救駕.
他都可能會死.
為什麼.
明明講清楚那坑沒有水.
在那裡.
如果將他丟下去.
誰會幫他.
叫天不應 叫地不聞.
然後發生什麼事.
丟了約瑟下坑之後.
那些兄弟居然有胃口吃飯.
他們坐下吃飯.
有一群以實瑪利人從基列來.
基列在東邊.
用駱駝帶著香料與香末藥.
他們要去哪裡.
去埃及.
猶大突然間出了一個主意.
跟其他哥哥說.
我們殺了弟弟.
遮掩他的血有什麼好處.
我們將他賣給以實瑪利人.
不要下手害他.
因為我們的弟弟是我們的骨肉.
其他兄弟就聽他說.
那時有些米電的商人從那裡經過.
他們就將約瑟從坑裡面拉上來.
他們就用20塊的銀紙.
20塊.
就將約瑟賣給以實瑪利人.
他們就把約瑟帶到埃及去了.
你以為完了嗎.

$^{321}$還未完.
29節.
回到坑旁.
這個不知道去哪了.
可能要照顧羊.
羊可能跑了.
可能要把羊抓回來.
剛剛好他不在.
回到坑裡面.
看見約瑟不在坑裡.
就撕裂自己的衣服.
回到他的兄弟那裡說.
孩子不在了.
我在那裡才好呢.
於是他們就宰了一隻公山羊.
拿了約瑟的那件外衣染了一些血.
派人把長袍先送回他爸爸那裡.
說我們發現這件長袍有血的.
你來認一認.
是不是你兒子的外衣呢.
然後我們skip下去第36節.
鏡頭轉過去就是米電人.
就將約瑟賣到埃及.
給了一個法老的官員.
一個護衛長叫波提法.
約瑟的故事當然很長.
37章至50章.
有14章的聖經是創世記最後的.
是很迂迴曲折的一個故事.
很多時候我們看這個故事的時候.
當然我們如果真的會看故事的話.
你想看故事的話.
其實這個故事是很好看的.
但是怎樣去解釋一些東西呢.
這個就是巧妙的地方.
希伯來人的文化是怎樣呢.
Let me tell you a story.
有一些很重要和你有關的.
屬靈的教導.
怎樣告訴你呢.

$^{361}$不是list個form.
point form給你.
而是去到給一個迂迴曲折的故事給你.
當然我們先劇透.
約瑟被賣到埃及發生什麼事.
經過很多苦難很多波折等等.
最後成為了埃及的一個大官.
然後當地有大饑荒.
然後雅各一家又被逼到埃及去吃糧.
然後故事最後有一個所謂好的收場.
輾轉之下兄弟和約瑟就復和.
整個家庭就被救回來.
到最後的時候.
50章第20節.
約瑟就說什麼.
跟他的哥哥說.
從前你們的意思是要害我.
但上帝的意思原是好的.
我們剛才看第37章.
有很多很多的細節.
這裡有兩個point值得留意.
第一.
在上帝的providence底下.
上帝的護理底下.
第一每個人都有自由作主.
37章我們剛才看了很多不同的經文.
家父做過什麼.
派約瑟去視察一下.
他們就平安.
他絕對有自由去做一個這樣的選擇.
派約瑟去探望一下.
保佑平安.
第二.
他的兄弟.
由希伯文谷上到北面事件.
他們做什麼呢.
放羊.
可能那裡找不到水.
不夠水不夠草.
他們再往北一點.

$^{401}$往北一點就會有多一點雨水.
他們一直往北走的時候.
他們是不是在做他們自己的決定.
非常自由.
free choice.
也是按照他們當時的環境去做一些決定.
你記得那個坑有沒有水.
坑有沒有水.
沒有水.
坑沒有水是什麼意思.
一年裡面什麼時間.
夏季最後的時候.
因為他們最喜歡掟坑.
雨水掉下去他們就儲起雨水.
然後夏季的時候沒有雨水.
四五個月.
他們就舀水來用.
無論是放羊需要喝的水.
或者人需要喝的水.
他們就會舀水.
到了夏季最後的時候.
就舀到底.
沒有水.
夏天最末的時候.
他們就上去越來越北面放羊.
這個是牧羊人應該要做的事情.
絕對有自由.
弱勢做了什麼.
他遵命在希伯倫谷往北上去.
先去哪裡.
事件.
然後再往多丹去追他的哥哥.
他有沒有自由.
他不是機械人.
他一直跑上去就是想找他的哥哥.
就是為了做什麼.
為他家父報一個平安.
然後兄弟遠遠見到.
弱勢來的時候.
他們做了什麼自由的選擇.

$^{441}$他說這個Dream是你的.
不如我們將他殺了.
兄弟是要去加害於他們弟弟弱勢.
他們要不要負責.
要負全責.
故事的道德責任.
永遠都是放在兄長上.
不是放在誰身上.
不是放在家父身上.
不是因為家父叫了弱勢上去.
所以你錯了.
也不是放在弱勢的身上.
可能弱勢很討人厭也好.
很抵得人愛也好.
那不是重點.
將弱勢苦害的道德上的責任.
是放在他的兄長身上.
不是雅各不是弱勢.
也不是其他故事裡面的人物.
所以每個人都有自由自主的決定和選擇.
但第二方面.
在際遇上也有很多意料不及的東西.
有什麼意料不及.
弱勢去了士劍找他的哥哥.
你記得走迷了路.
蕩失了.
浪費了拖延了某些時間.
然後發生什麼事.
剛剛有個老人鴿.
剛剛才有個老人鴿又怎樣.
老人鴿又八卦聽到那些哥哥.
離開士劍去哪裡.
去多丹.
剛好遇到老人鴿.
然後可以指示他離開士劍去多丹.
然後他追上他的哥哥.
兄弟遠遠看到弱勢正在來的時候.
本想做什麼.
就是將他殺死.
不知為何抽起一根筋.

$^{481}$裡面的大哥.
不知為何很好人.
想救他.
結果弱勢被吊在一個乾坑裡面.
誰知道在他們吃飯的時候.
遠遠看到有一班商人經過.
剛好那些商人去哪裡.
就去埃及.
不是去他們家.
去他們家是大件事.
就是要去埃及去別的地方.
剛好這班遠遠看到的商人.
他們就是賣貨.
拿著他們的財貨.
去埃及做買賣.
也不是說要賣奴隸.
剛好遇見他們的時候.
他們就說不如將弱勢賣給這些商人.
也剛好是猶大出的主意.
我相信不會是他們開幾個小時的執事會.
然後做一個最好的決定.
沒有計劃.
預料不到的情況下.
他們見到這些商人猶大就出了主意.
不如將他賣給埃及.
預料不到的當中.
弱勢就剛好暴死在坑裡面.
在白羅娛的聖經裡面.
很多時候會用一個字.
給讀者有些警覺.
讀故事的時候.
就算讀children's story也是一樣.
有些signal給你.
讀者告訴你.
有些事情你要留意.
那個字叫Hinay.
中文翻譯就是Horner.
有時候Horner的意思就是始料不及的意思.
預料不到.
surprise.

$^{521}$起碼在人的角度是surprise.
在這個故事裡面出現了好幾次.
甚至有個老人家遇見弱勢.
Hinay.
他就在田野裡迷了路.
那個人就問他你在找什麼.
然後他就指示他.
我剛剛聽到八卦.
19節.
兄長彼此說.
Horner他們也預料不到.
弱勢會來通風報信去問過平安.
是雅各加夫幾天前.
打發他去找他的哥哥.
Hinay.
居然這個造夢的人會來.
25節.
他們坐下吃飯的時候.
弱勢在哪裡?.
在坑裡面.
在那裡呼叫.
在那裡很慘.
Hinay.
Horner.
有一群商人從基列尼.
拿著他們的貨財往埃及去.
29節.
裡面回到坑裡面.
Hinay.
Horner.
弱勢不在坑裡面.
他就撕裂自己的衣服.
我剛才說聖經很多時候是說一個故事.
來交代或者教導一些很重要的原則.
很重要的功課.
就是用一個這樣的故事.
如果你細讀的話.
你會發現到底他想說什麼.
就是我們的主題.
Providence.

$^{561}$上主在幕後做一些護理.
Caring.
你可以將這個故事再列出來.
你就會發現.
所有東西都是我們所謂的剛剛好.
Right?.
整定.
其實就是Providence.
Right?.
對嗎?.
每個人都有他的自由去做任何的選擇.
但是在他做選擇的裡面.
也有一些始料不及的事.
好像背後有一隻手.
是看不見的手.
無形的手.
Invisible hand.
是帶領著整個的所謂棋局.
Right?.
大家有沒有下棋?.
有沒有試過和電腦下棋?.
和電腦下棋.
我下棋很差的.
所以我最討厭下棋.
以前和我爸爸下棋我最討厭.
因為你每走一步.
你覺得你那一步走得很精準.
但其實你就落了他的陷阱裡面.
Right?.
因為我看到的只不過是開頭那幾步.
但他看到的是什麼?.
他看到整局.
比我看得更加遠.
所以你和電腦下棋的時候.
很難贏的.
但是你有沒有自由?.
你絕對有自由.
和上帝相比的話.
我們是一些很差的棋手.
我們看到的非常有限.

$^{601}$上帝看到的是全盤.
全盤的棋局.
我們說人生如棋.
說得沒錯.
人生真的好像下棋一樣.
上帝的護理.
這個訊息就是告訴我們.
OK 人生如棋.
有很多事情我們做.
憑我們自己有限的看法.
我們做很多決定.
但整個棋局在上帝的眼中.
Providence.
Seeing ahead of time.
He saw it ahead of time.
關於上主護理的教導.
可能我們可以用一個Point form.
我們都是喜歡這樣.
不妨用一個Point form.
試試列出來.
我們繼續看約瑟的故事的時候.
可能我會發現更多.
起碼在37章開頭這章說什麼呢?.
上帝是用插不透的智慧指揮著.
和安排著萬事.
大小的事情.
時間.
際遇.
你想想連約瑟迷了途.
浪費了不知多少時間.
正是因為他浪費了時間.
他才死不去.
如果他早點去的話.
商人經過的時候.
可能他都沒有力氣去叫.
就是剛剛好.
他吊在坑裡面.
兄長吃飯的時候.
就遠遠見到.
不時間性的掌握.

$^{641}$誰去掌握.
不是故事裡面的角色.
而是在背後有人指揮著.
和有人安排著.
第二 神沒有認同過他兄長的罪行.
兄長是要為自己的惡行負上全盤責任.
神的護理.
上帝的Providence.
沒有抹殺過人的自由.
和人的道德責任.
包括我們每一個的行動.
我們每一個的禱告.
我們每一個的祈願.
沒有抹殺.
反而得到確立.
這個神不神奇.
這個就是神秘的地方.
神也會向所有人去問責.
因為你是自由的人.
你是會做抉擇的人.
上帝的護理.
上帝在背後的安排萬事.
沒有抹殺過你要負上的責任.
神是會用人的美善來達到他的計劃.
人做得好當然會達到他的計劃.
神會用他的僕人去做他的事.
但同樣 神也用人的惡行.
可以凌駕他們.
他們照做惡行.
但最後也會達成他的最終計劃.
第四.
第五.
最終目的.
上帝的護理是做什麼的?.
是使愛神的人.
我們能夠得到安慰和平安.
說得對.
如果我們認識到.
這個教導.
Providence.

$^{681}$我們不要說它是天命.
因為聖經從來沒有說過我們是靠運氣.
命運安排.
讓事情整定.
機緣.
這麼簡單.
因為命運是麻木的.
我們不是被一些麻木的命運去擺佈.
我們不是被一些不知道發生什麼事的.
啞的機緣去支配.
上帝是會說話的.
上帝是會行事的.
人類歷史裡面他會參與.
他會推進自己所做的事.
在教會歷史裡面有很多前人.
他們曾經去思考這個問題的時候.
或者去掙扎這個教導的時候.
他們很多時候還會產生一些信條.
我們叫信條.
Statement of Faith.
或者一些條文.
其中一個很著名的條文叫.
Westminster Confession of Faith.
他說到Providence of God的時候.
他說什麼呢?.
為什麼叫.
不要緊,這個是整定的.
說到這個Westminster Confession of Faith.
他說什麼呢?.
為什麼叫Westminster?.
因為在倫敦有一個地方叫Westminster Abbey.
一個大的教堂.
在大約17世紀的時候.
當時的教會就說.
不行,我們要將我們的信仰.
我們的神學整頓一下.
整理一下.
怎樣可以有系統的去整頓出來.
所以就弄了這個Westminster Confession of Faith.
一個很長的,差不多一本書的.

$^{721}$一個神學的總結.
整本聖經怎樣去討論.
教導某一些的主題.
說到Providence的時候.
他說什麼呢?.
他說偉大的上帝創造萬有.
藉著祂至高的智慧.
至聖的護理.
根據祂無誤的預知.
以及上帝的自由.
不變的旨意.
他確實的去維護著.
指導著.
安排著.
並且統管一切受造之物.
包括他一切從最大到最小的行動和事物.
為要使自己的榮耀,智慧,能力,公義,良善.
因為得到稱讚.
上帝創造萬有.
當然他是凌駕於我們受造物.
凌駕於這個創造.
但不等於上帝不在.
上帝不是抽離的.
上帝仍然是主動,直接,間接,明的,暗的.
去指揮著萬有.
這個就叫做Providence.
我預備的時候.
記起以前小時候.
我兒女幾歲的時候.
大家做父母應該很有同感.
就算不做父母的都應該明白.
小時候做爸爸的人.
經常跟他玩一個遊戲叫做Hide and Seek.
如果你有經驗.
做過爸爸跟小朋友玩Hide and Seek的時候.
你就知道Hide and Seek其實是很考技巧的.
為什麼呢?.
因為我絕對有這樣的能力.
將我自己完全隱藏.
他永遠都找不到我.

$^{761}$我記得有一次我這樣跟他玩.
我躲在那裡很久.
他找不到我,他就哭了.
然後我才了解原來Hide and Seek不是這樣玩的.
跟他玩的話.
是需要將自己隱藏.
但最後.
我可以被他找到.
可能有點像上帝跟我們之間的關係.
有時候我們經歷不同的際遇.
或者不同的高山低谷.
我們真的覺得有時候上帝在哪裡.
上帝在搞什麼.
但想想我們以前的故事.
我們以前的經歷.
上帝是不是完全隱藏.
會不會找到他.
創世記念37章.
上帝雖然有時候隱藏.
但不代表他消失.
上帝的護理.
給我們安慰,給我們平安.
我盼望今天你離開的時候.
你思考一下你過去.
上帝似乎隱藏.
但你卻看到時間,際遇.
有時候剛剛好.
上帝在幕後指揮著.
Yes.
我就有鼓勵.
我有信心.
我們每個人都經歷過.
前面是怎樣呢.
我不知道.
我是一個很差的騎手.
我可以知道多少.
我可以知道多少步,看到多少步.
不是很多步.
不過最重要的是什麼.
整盤棋局在上帝的手裡.

$^{801}$他得到多少.
因為你是偉大的主.
你有至高的智慧.
而且你對我們生命當中.
世界裡面一切的護理.
是至聖,至高的.
主要你也知道你自己的預知.
你的指導,你的安排.
你統管管理一切的.
是沒有錯誤.
包括一些最小的動作.
看似是一些很等閒的事.
或者一些很大的決定.
主要是讓我們的心在裡面有安穩.
因為我們知道我們可以去尋求你.
你也知道我們的生命不是被一些無情的.
啞的一些命運去擺佈.
縱使我們未必每樣都能夠全盤解釋.
但我們知道你自己至愛的護理.
你的引導是帶領著我們.
終瞻榮耀是歸給你的.
帶領著我們的生命讓我們繼續信靠你.
認識你,祈禱服從你.
願上帝和大家同在.
賜給大家平安.
\newpage



\section{}
\label{sec:_iX85FoHSZc}
\textbf{比對研究 18: 巴力宗教 11 - 列王記上 18 - 陳崇基牧師}
\newline
\newline
連結: \href{https://youtube.com/watch?v=_iX85FoHSZc}{\texttt{https://youtube.com/watch?v=\_iX85FoHSZc}} ~~~~ 語音日期: 2023-09-23
\newline
\newline
\hyperref[sec:Z86XRROS_74]{\small{< < < PREV SERMON < < <}}
~
\hyperref[sec:index_chronic]{\small{[返順時目]}}
~
\hyperref[sec:index_scriptual]{\small{[返順卷目]}}
~
\hyperref[sec:mJHLbAHSJBg]{\small{> > > NEXT SERMON > > >}}
\newline
\newline
$^{1}$如果你看十八章繼續記載的時候.
這裡說什麼呢?.
二十一節.
他說心慈意悅到什麼時候呢?.
他們就一言不達.
以利亞就跟百姓說.
他說作耶和華先知的子子只剩下我一個.
巴力就有四百五十個先知.
現在以利亞就說.
這個故事大家很熟悉了.
大家拿一隻牛肉出來.
然後切成小塊.
不要點火.
然後放在柴上.
也不要點火.
就求告你們神明的名字.
我也求告耶和華的名字.
應允禱告降火的就是上帝.
百姓就回答.
這個是很主要的.
以利亞就跟巴力的先知說.
因為你們人多.
你們先挑選一隻牛肉.
而預備好了之後就求告你們神明.
但是卻不願點火.
OK.
所以在鬥法裡面.
在加密山上.
大家拿出一個 sacrifice(犧牲).
一個祭出來.
哪一個神去降火.
哪一個神去燒盡祭物.
就代表祂是真神.
是不是?.
這個應該是不陌生的了.
因為在舊約近東的世界裡面.
很多時候那個獻祭是用來做什麼的呢?.
那個獻祭是用來討好神明的.
譬如說神明發怒.
你就獻祭.

$^{41}$如果神明肯收你的祭物的話.
就代表祂跳起信了.
祂就不再發怒了.
所以在近東的異教文化裡面.
巴力如果是生氣的話.
為巴力獻祭.
希望巴力收這個祭.
那麼祂就不會再發怒了.
就會祝福我.
所以神明和人之間是有一個 symbiosis(相伴).
就是大家互保互利的一種關係.
你知道的.
上下文的故事已經告訴我們了.
到底現在哪一個神明發怒.
巴力的宗教.
巴力的先知.
會想什麼呢?.
就是巴力.
不是很生氣的.
為什麼呢?.
因為有幾年沒有下雨.
巴力到底去了哪裡呢?.
是不是?.
所以在巴力先知的思維裡面.
你就可以知道.
以利亞所給的主意.
其實是很 make sense(合理)的.
當然這個時候我們需要去.
所謂 appease(討好)那個 angry god(憤怒的神明).
就是去討好他.
所以你看到他們很順耳.
他們聽話.
然後就說這是一個好的主意.
他們就將他們的牛毒預備好了.
從早上到中午就求告巴力的名下.
巴力就說求你應允我們.
卻沒有聲音.
也沒有回應.
他們在所捉的壇上四處崩跳.
到了午靜的時候.

$^{81}$太陽最猛烈的時候.
以利亞就嘲笑他們了.
他們大聲一點.
因為如果他們是神明的話.
或者他們在默想.
或者他們正在忙.
或者他們在路上.
或者他們在睡覺.
他們應該現在要醒了.
是不是?.
大家再大聲去求告的時候.
按著他們的儀式.
用刀刺割自己.
直到渾身流血.
中午過了之後.
他們狂呼亂叫.
直到獻晚祭的時候.
都沒有聲音.
也沒有回應.
也沒有理睬.
是不是?.
在這裡可以留意什麼呢?.
如果留意原文的話.
你記得我們剛剛看過二十一節的時候.
以利亞在加密山上跟他們說什麼.
他說你心慈意義到什麼時候.
如果南邊是耶和華信仰.
是不是?.
是上帝的話.
如果他是上帝就當順從耶和華.
是不是?.
如果北邊的話.
就應該當順從巴力.
他們就一言不達.
是不是?.
心慈意義是什麼意思?.
就是從這個Pasach這個字來的.
Pasach這個原文就是.
Poschim.
一個participle.

$^{121}$就是說你們要心慈意義.
就是擺動到什麼時候.
然後到了二十六節.
剛剛看完.
就是巴力的先知.
他們從早上到下午.
求過巴力的名的時候.
他們在做什麼呢?.
沒有聲音沒有回應的時候.
他們就在所觸的祭壇四處做什麼?.
又是Pasach這個字.
是不是?.
Paschal這個字.
就是他們崩跳.
或者如果你的翻譯是做得 consistent(一致)一點.
一致一點.
整張聖經裡面.
出現Pasach這個字.
只有這兩字而已.
第一個就是形容那些百姓.
心慈意義.
第二個就是形容那些巴力先知.
四處跳.
所以這個.
小說聖經作者有些嘲諷性.
是不是?.
那些百姓四處崩跳.
要到什麼時候呢?.
跳上跳下.
是不是?.
巴力.
巴力的耶和華.
事實上他們就是這樣跳來跳去的.
有時候跳到北面.
有時候跳到南面.
有時候跳到這裡跳到那裡.
他們就好像那些假先知那樣.
四處在Pasach.
是不是?.
四處在Pasach.

$^{161}$是.
這個你看完文的時候.
你就會發現.
如果翻譯得比較 consistent(一致)的時候.
其實你可以去做一個 note(筆記).
就是有一個 footnote(筆記).
有一個筆記在這裡.
就是你們心慈意義.
其實你們崩跳到什麼時候.
彈出彈入要到什麼時候.
這裡講到巴力不在.
沒有聲音.
沒有回應.
在巴力的神話裡面.
我們之前看過很多.
巴力的神話裡面確實有講到巴力不在.
譬如有一個巴力的詩句.
這麼說.
就是說 Anart.
就是他的所謂情人.
Anart.
就想去找巴力.
找巴力做什麼呢?.
他希望找巴力去做一些事情.
特別是降雨.
但是去到巴力的皇宮.
或者巴力的廟的時候.
是怎麼樣呢?.
當這個 Anart 去到這個廟的時候.
就被他的使者.
巴力的使者告訴他什麼事情.
你就會看到這裡一至五行.
就說什麼呢?.
這個 Anart 去到的時候.
他就說巴力在他的家嗎?.
或者這個 Hade.
Hade Hada.
神明在他的家嗎?.
他在家嗎?.
是不是?.

$^{201}$然後巴力的使者.
Led.
他們的僕人.
就這樣回答 Anart.
說巴力不在家.
Hada 神明不在家.
所以在這些神話裡面.
都有這麼說.
巴力是 Absentee God.
就是有時候他不知道跑到哪裡去了.
所以如果你看看.
聖經裡面記載.
以利亞所說的話.
沒有聲音沒有回應.
其實是熟悉巴力宗教的人.
其實就是在用他們的術語.
或者他們那種巴力的神學.
來去嘲諷巴力.
確實有時候真的不在.
不知道跑到哪裡去了.
是不是?.
如果只是看這裡.
你就會看到.
他可以是 Absentee God.
但是再回到.
迦南地的農耕的循環的時候.
你就會發現.
每逢到夏季的時候.
巴力確實是不在家的.
跑到哪裡去了?.
就跑到了 Underworld 裡面.
是不是?.
跑到了 Underworld 裡面.
他就不見了.
為什麼呢?.
因為他被 Yharnam.
就是 Chaos.
或者被死神 Mordek.
所吞併或者暫時打敗了.
如果不是他復活.

$^{241}$回頭的話.
就不會有雨水了.
所以一年又復一年.
巴力就不見了神.
不是不見人.
不見了神.
所以就用巴力的神話.
或者巴力的神學.
來做一些護教.
Apologetics.
或者一些 Sceptile.
就是一些諷刺的嘲諷.
以利亞就在做這件事情.
你看看如果你靠巴力的話.
就會是這樣了.
有時候他真的不在.
是不是?.
每逢到夏季的時候更加不會在.
你看看這個故事繼續下去的時候.
沒有聲音沒有回應的時候.
他們在四處彈來彈去.
四處崩跳.
到了正午的時候他就嘲笑他們.
以利亞說什麼呢?.
他說大聲呼喊.
因為他們神明在默想.
或者不見了人等等.
他們大聲求告.
按著他們的儀式就用槍.
用刀用槍刺傷自己.
直到渾身流血.
很多時候我們會想起這裡的時候.
就會覺得很血腥.
但是在近東的宗教裡面.
或者近東的一些宗教的儀式.
當你很趕急的時候.
很危急的時候.
或者很傷痛的時候.
很多人都會做這件事情.
為什麼呢?.

$^{281}$就是為了使神明去回應你.
在舊約的不同的經文.
都有提到這種的 practice(練習).
這些這麼奇怪的 practice(練習).
耶利米斯四十一章第五節.
那裡是講到主前586年的時候.
尼泊爾加尼沙巴比倫王.
將第一個聖殿.
耶路撒冷的聖殿毀滅了.
然後大家很傷痛.
在北國的地方.
有80人.
從士堅,士羅和撒瑪利亞.
他們來到耶路撒冷.
鬍鬚剃去.
衣服撕裂.
身體劃破.
手裡面拿著掃祭和乳香.
來到耶路撒冷的殿.
那個殿已經拆毀了.
在殿的位置上.
已經是灰燼了.
就想在那裡去祭祀.
四十七章第五節.
這裡就講到.
亞什基倫或者加薩.
就是在佛利時亞.
非利士地的地方.
非利士地在哪裡呢?.
就是在這裡.
就是迦林地沿海平原的南面.
加薩就在下面這裡.
亞什基倫.
這些都是剛剛的經文提到的.
他說.
加薩成為江突.
亞什基倫歸於無尤.
平原上所剩下的.
是怎麼算呢?.
就是死剩的人.

$^{321}$這裡又發生什麼事情呢?.
巴比倫王.
尼布格尼薩.
巴比倫君.
來到巴基斯坦地的時候.
他就剷除了很多這些城鎮.
包括了非利士人的地方.
剩下來的人.
他們就是受到糟蹋.
以及被毀壞.
平原所剩下的人.
就是死剩的那一群人.
南面平原的.
非利士地的.
你們割畫自己要到什麼時候呢?.
他們很傷痛.
想神明來回應他們.
但是這些神明當然都是假的神明.
所以有一些這樣的 practice(練習).
是一些不會 surprise(驚訝)我們的.
不是驚訝的事情.
但是更加讓我們看到的就是.
在巴力的文獻裡面.
提到一種這樣的哀悼.
就是當巴力不在的時候.
沒有雨水的時候.
旱災的時候發生什麼事情呢?.
我們看看巴力循環 bell cycle(鐘循環)裡面.
就是這一段了.
我們很粗略地中文翻譯出來.
全能的巴力死了.
死去的王子大地的王.
他離開了.
然後那些神明和人.
為了巴力不在.
他死去.
下到了 Underworld 陰間.
被 Mort.
死神所吞併之後.
那些神明為了這位大能者巴力死去.

$^{361}$然後傷痛.
這裡就講到.
他.
就是最高的神.
叫做 L.
用石塊來刮傷.
刮損皮膚.
用刀片的傷人.
他刮臉頰和下巴.
刮他的手臂.
刮他的胸膛.
如同刮田地一樣.
刮他的背脊.
如同刮山谷一樣.
你可以想像.
就是一個農業的文化.
沒有雨水的時候.
農夫照樣去刮田地.
但其實在做什麼呢?.
這裡的神話就說.
他們很哀悼巴力不在.
希望能夠刮傷自己.
使巴力能夠回來.
希望巴力.
因為敬拜者.
這樣的緣故.
用刀刺傷自己.
直到渾身流血.
然後希望在大聲疾呼的意識裡面.
巴力終於希望能夠回應.
當然他們所預期的回應是什麼回應呢?.
就是巴力的聲音.
巴力的聲音是什麼意思呢?.
就是雷電的聲音.
打雷.
下雨.
地震的聲音.
當大風暴來到的時候.
從地中海.
西邊.

$^{401}$一路刮風過來的時候.
你就可以看到.
打雷打雨.
然後黑雲滿佈的時候.
就將雨水落在地上.
有雷電交加的時候.
地也震動了.
這種就是他們所預期.
希望能夠巴力這樣去回應.
當然我們知道.
故事是怎麼發展的.
我們下一次再繼續看下去.
\newpage



\section{}
\label{sec:mJHLbAHSJBg}
\textbf{God's Providence 1 (Genesis 37) - Rev. Mark Chan}
\newline
\newline
連結: \href{https://youtube.com/watch?v=mJHLbAHSJBg}{\texttt{https://youtube.com/watch?v=mJHLbAHSJBg}} ~~~~ 語音日期: 2023-09-26
\newline
\newline
\hyperref[sec:_iX85FoHSZc]{\small{< < < PREV SERMON < < <}}
~
\hyperref[sec:index_chronic]{\small{[返順時目]}}
~
\hyperref[sec:index_scriptual]{\small{[返順卷目]}}
~
\hyperref[sec:Fs4lhTNde4E]{\small{> > > NEXT SERMON > > >}}
\newline
\newline
$^{1}$- Hey guys, we will turn to the book of Genesis chapter 37..
If you have your Bible,.
turn with me to the book of Genesis chapter 37..
And today we are going to begin a new topic..
This topic is called the providence of God..
So we're gonna read the story of Joseph..
Now, anyone has any idea what the word providence means?.
What does it mean?.
It's actually two words together, right?.
Combined compound word, right?.
Meaning pro, ahead or beforehand,.
and then seeing, where we get vision, right?.
That is the providence of God..
And this is a precious teaching found in the Bible..
At times we might think or mistaken providence as fate.
or some sort of blind destiny..
Everything is set, past, present and future,.
and everything is predetermined in your life, in our life..
Everything is pre-decided mechanically..
So we're just going through our journey of life.
and every day is already preset..
Sometimes we have this wrong idea,.
but in the Bible, providence refers not.
the blind mechanism of the universe or in your life..
Providence means God is caring and seeing ahead of time.
in whatsoever comes to pass, whatsoever..
The people that you meet, the coffee that you had today,.
driving here, right?.
Some more people got in your way, right?.
You probably yell at some people..
You come here and join the worship..
Everything whatsoever comes to pass..
That is the biblical definition of providence..
It's not a mechanical, you know,.
everything decided and predetermined kind of fate.
that we're going through our lives,.
but God is seeing ahead of time and caring for you..
So that's what it means, providence..
And there are a lot of different Bible verses.
that has to do with providence..

$^{41}$For example, Psalm 139, one of the favorite Psalms, right?.
"Oh Lord, you have examined my heart.
"and know everything about me.".
Everything, see the key word here?.
Everything..
You know, when I sit down or stand up,.
when you sit down, take a break,.
when you're standing up, doing something, working,.
Lord, you know my thoughts, even I'm far, far away, right?.
When you're traveling, we are not home..
Lord, you see me when I travel, when I rest at home..
That's NLT, by the way..
And you know everything I do..
That is providence, isn't it?.
Right?.
Many scriptures teach the providence of God.
or the doctrine of divine providence,.
that your life is under the beforehand seeing of God,.
beforehand caring of God..
Whether you believe it or not, right?.
Now there are also different Bible verses.
we're not going to dig into them,.
but this is something unique about the biblical culture.
or the Hebrew culture..
Even today, when you ask a rabbi a question,.
usually you don't get a direct answer..
The rabbi probably will respond, "Let me tell you a story.".
That's interesting, right?.
It's kind of different from our Western.
or urban kind of linear, you know,.
we all want to outline things.
and bullet points and all that, right?.
We're more abstract in our thinking,.
but in the Hebrew culture or Bible culture,.
they will tell you a story, right?.
Remember the prophet Nathan,.
when he was about to confront David.
for the great sin that he committed,.
what did Nathan do?.
(audience murmurs).

$^{81}$He told him a story, told him a parable, right?.
And Jesus did a lot of parables, illustrations, right?.
Daily life illustrations.
in order to flesh out some of the truth, right?.
So today we're going to read story, all right?.
Big chunk of it, right?.
Biblical story of Joseph..
You ready?.
Let me tell you a story, all right?.
Chapter 37, verse two..
This is the account of Jacob and his family..
When Joseph was 17 year old,.
he often tended his father's flock..
He worked for his half brothers,.
the sons of the father's wives, Belhar and Silpah..
But Joseph reported to his father.
some of the bad things his brothers were doing..
Jacob loved Joseph more than any of his other children.
because Joseph had been born to him in his old age..
So one day, Jacob had a special gift made for Joseph,.
a beautiful many color robe..
But his brothers hated Joseph.
because their father loved him more than the rest of them..
They couldn't say a kind word to him..
This is interesting..
The Hebrew here is they cannot say shalom to him..
So they can't say good morning or hi to him..
Can't say a kind word to him, can't say shalom to him..
And then the following context reports.
about the two dreams that Joseph had..
And we're not going to go into it,.
but the main idea is that someday,.
the dreams are foretelling the future, right?.
Ahead of time, that Joseph will be greater.
than all his family members, his father, his mother,.
and all the brothers..
And that's why the brothers hated him, right?.
Imagine your little brother telling you,.
"Someday you're gonna worship me.".
No way, right?.

$^{121}$Your little brother or your younger sibling.
saying to you, "Someday I'm going to have.
"all your investment, all your bank accounts.".
No way, right?.
That's why the brothers hated him..
And then verse 14..
Now what happened here is Jacob said to Joseph,.
"Go and see how your brothers.
"and the flocks are getting along..
"Then come back and bring me a report.".
So Jacob sent him on his way,.
and Joseph traveled to Shechem.
from their home in the Valley of Hebron..
When he arrived there, a man from the area noticed him.
wandering around in the field or countryside..
"What are you looking for?" he asked..
"I'm looking for my brothers," Joseph replied..
"Do you know where they are pasturing their sheep?".
"Yes," the man told him..
"They have moved on from here,.
"but I heard them saying, 'Let's go to Dothan.'".
So Joseph followed his brothers to Dothan.
and found them there..
Verse 18..
When Joseph's brothers saw him coming,.
they recognized him in the distance..
As he approached, they made plans to kill him..
"Here comes the dreamer," they said..
"Come on, let's kill him.
"and throw him into one of the cisterns..
"We can tell our father a wild animal has eaten him..
"Then we will see what becomes of his dreams.".
Right?.
21..
But when Reuben pulled up the scheme,.
he came to Joseph's rescue..
"Let's not kill him," he said..
"Why should we shed any blood?.
"Let's just throw him into an empty cistern.
"here in the wilderness..

$^{161}$"Then he'll die without our laying a hand on him.".
Reuben was secretly planning to rescue Joseph,.
and return him to his father..
So when Joseph arrived,.
his brothers ripped off the beautiful robe he was wearing..
Then they grabbed him and threw him into the cistern..
And the pit was empty, and there was no water in it..
Then, just as they were sitting down to eat,.
how can you eat with your little brother in the cistern?.
Crying out..
They're having fast food here..
They look up and saw a caravan of camels in the distance.
coming toward them..
It's a group of Ishmaelite traders.
taking a load of gum, balm, aromatic resin.
from Gilead down to Egypt..
Judah said to his brothers,.
"What will we gain by killing our brother?.
"We have to cover up the crime..
"Instead of hurting him,.
"let's sell him to those Ishmaelite traders..
"After all, he is our brother,.
"our own flesh and our own blood.".
His brothers agreed..
So when the Ishmaelites, who were Midianite traders,.
came by, Joseph's brothers pulled him out of the pit.
and sold him to them for 20 pieces of silver,.
and the traders took him down to Egypt..
Now, sometime later,.
Reuben returned to get Joseph out of the pit..
When he discovered that Joseph was missing,.
he torn his clothes in grief,.
and he went back to his brother and lamented,.
"The boy is gone..
"What will I do now?".
Then the brothers killed a young goat.
and dipped Joseph's robe in his blood..
They sent the beautiful robe to their father.
with this message, "Look at what we found..
"Doesn't this robe belong to your son?".

$^{201}$And then we skip to 36..
Meanwhile, the Midianite traders arrived down in Egypt,.
where they sold Joseph to Polypher,.
an officer of Pharaoh, the king of Egypt..
Polypher was captain of the palace guard..
Now, this is a long story, right?.
Whole chapter 37, pretty much we read every single verse..
It clearly tells us the older brothers of Joseph.
were ill-intent..
They wanted to kill Joseph..
Remember the pit, the cistern?.
What happened to it?.
Was there any water?.
No, there was no water, right?.
According to the Palestinian climate,.
this should be the end of summer..
No water at all..
They wanted to kill him,.
but they changed their mind and agreed with Reuben..
And Reuben, all of a sudden, he intervened and said,.
"Hey, let's put this boy into the pit," right?.
Where there's no water..
And he actually intended to come back and to rescue Joseph..
However, as the story goes, some traders,.
some international traders, Ishmaelites or Midianites,.
they came by and they decided,.
according to Judah's suggestion,.
sell Joseph to the trader and bring him down to Egypt..
As you know, the story of Joseph runs very, very long,.
right, 37, chapter 37 to all the way to the end,.
chapter 50, end of the book of Genesis..
And we all know the end, right?.
Well, let me spoil it, yeah..
Chapter 50, verse 20, finally,.
even the brothers and Jacob,.
the whole family came down from Canaan to Egypt.
because of famine, right?.
And just eventually, Joseph rose from the pit.
to being the second in command in the nation of Egypt..
And he was kind of a savior..

$^{241}$He saved his whole family..
And eventually, all the brothers reconciled..
And Joseph said to his brother, "That is the end,.
"almost the end of the story," chapter 50, 20..
"You guys intended to harm me,.
"but God intended it all for good..
"He brought me to this position.
"so that I could save the lives of many people.".
That's just a long story..
We will go through it, you know, chapter by chapter..
But this is how the Bible works..
Tell your story, a very complicated, very detailed story.
so you can chill upon it..
You can return the story again and again.
and think about the topic here, the theme here,.
the providence of God in your life..
Providence of God in your life..
Somehow God is seeing or caring for you way ahead of time.
before you can see how your life unfolds..
God sees ahead of time..
So there are actually two points I want to bring out..
First, if you read the story, just read chapter 37,.
what I want to point out is everyone in the story.
is free to make choices, right?.
Jacob, the father, he probably worries about, you know,.
the flock and the brothers,.
and send his beloved son, Joseph, to check on them, right?.
Nobody forced Jacob to make any choice..
He was a free agent..
He chose to send Joseph, right?.
Completely free..
So he made that choice..
And then the brothers were tending the flocks in Shechem..
Shechem, by the way, is in the north..
So he started, Joseph started in the south, Hebron,.
all the way up, probably walk three, two days,.
three or four days, to reach up north in Shechem..
And as I said, it was probably late summer.
because the pit was dry..
And then the brothers made their own choice as free agents..

$^{281}$They moved from Shechem to another place called Dothan,.
which is further up north..
Probably they were trying to find water for the sheep..
But in any how, they were free to do so, right?.
Nobody forced them..
They made their own choices..
And then Joseph, the beloved son,.
he obeys his father and goes to Shechem..
And they're kind of wandering, he got lost..
And then finally he reached the brothers at Dothan, right?.
Joseph made free choices as well..
And then the brothers were held fully accountable.
for victimizing Joseph..
The story lays blame on brothers,.
not Jacob, the father, not Joseph, the beloved son, right?.
No matter how Joseph is lovely or not,.
he's getting along with the brothers or not,.
or he's just a very wicked person or not, right?.
Good or bad..
The story lays blame only on the elder brothers, right?.
So everyone in the story is free to make choices..
However, some things happen in the story.
out of expectation..
Nobody would have thought how the story unfolds..
For example, Joseph goes to Shechem,.
try to find his brothers, right?.
He's lost..
Probably he has wasted some time, or a day or so,.
wandering in the wilderness, in the field..
And then surprisingly, a dude passes by..
And this dude, no name, right?.
Unnamed character in the story..
Just happened he heard what the brother said..
Remember what the brother says?.
Let's go out of Shechem and go to Doton, right?.
And then just, he heard about it..
He heard about it, and then surprisingly,.
he found Joseph getting lost, right?.
This young boy getting lost in the field..
And then surprisingly, he goes to Joseph and ask,.

$^{321}$"Hey, what are you doing?.
"What are you looking for?".
Right?.
And then Joseph reply, "I'm looking for his brothers.".
"Oh, I heard about your brothers.".
You know, they just went from Shechem here.
to someplace else, to Doton..
So some things happen in the story.
out of expectation for anyone, anyone in the story..
Secondly, the brothers..
Now I talk about Joseph, right?.
Wandering, wasting some time..
And then the brothers, the brothers meant to kill Joseph..
No question about it, right?.
But surprisingly, Reuben, the firstborn, intervenes..
He intervenes, "Don't kill him, don't commit murder.".
And Joseph end up in the pit, right?.
Now, why did Reuben think so?.
I don't know, I don't know..
Probably just, you know, spirit of the moment, right?.
Let's not kill him, let's not kill our brother..
Throw him in the pit, right?.
So the brothers, to them, something surprising happened..
Three, the cistern, the pit was dry, the end of summer..
If Joseph was left in the pit for too long,.
he would have died, right?.
Even he asked for help, nobody would have heard him.
because he would be too weak, right?.
The cistern was dry..
Nobody planned that, right?.
But something out of expectation happened..
And then finally, some of the merchants came by, right?.
Who were those guys, right?.
Medianites or Ishmaelites or whatever, right?.
Some international traders came by and said,.
didn't really talk to the brothers,.
but the brothers were having the, you know, fast food lunch..
And they saw these guys going down to Egypt..
And out of expectation, Judah did not plan this..
And he had an idea, let's sell our brothers, Joseph,.

$^{361}$to the traders, and the traders brought him down to Egypt..
So unexpectedly, Joseph did not die in the pit..
Eventually sold to Egypt..
Now, the Hebrew Bible has a special word in the story.
to catch our attention..
And usually this word is translated,.
behold or look, right?.
Some of the English translations, to my surprise,.
misses that..
And it's a word called hinei..
Hinei in the Hebrew word means.
trying to catch your attention, right?.
And they look, behold..
And in this story, in chapter 37,.
hinei happens again and again..
Let me give you some example, 15, verse 15..
A stranger, a guy, found Joseph and hinei..
He was wandering, he got lost in the field..
And the stranger asked him,.
and there was some, you know,.
correspondences between them..
Verse 19, the brothers said to one another, hinei..
Actually, the brothers did not expect.
Joseph would show up, right?.
'Cause Joseph is the beloved son..
He didn't have to work, right?.
He's always on the lazy boy playing games, right?.
He's beloved by the old father, right?.
Didn't expect Joseph would come and see how they're doing..
Hinei, he comes the dreamer..
And then they try to scheme to kill him..
25, the brothers sat down to eat their meal..
As they raised their eyes, hinei..
Raised their eyes and looked..
Hinei, look, behold, a caravan of Ishmaelites,.
the traders, coming from Gilead..
Gilead is up north, and then they're going down to Egypt..
Now, these guys, these traders,.
were not slave traders, right?.
With all of the commodities they are trying.

$^{401}$to bring to Egypt, they have no interest.
in slavery..
Judah got the idea and sold Joseph to those traders..
And then Reuben came back in 29..
Now Reuben returned to the cistern,.
and hinei, look, surprisingly,.
Joseph was not in it, so he torn his clothes..
Surprise, surprise, surprise, surprise..
At least there are four times here,.
hinei, the worst chosen..
So in the entire chapter, 37,.
especially the latter half,.
talking about the selling of Joseph.
or whatever he encounters,.
you get the idea that everything somehow is orchestrated..
Everything is working together.
down to the finest detail..
Joseph got lost in the field..
He wasted a day or two,.
and that's why he did not die in the pit..
And he was sold to Egypt for a greater purpose..
Providence, by means of a beautiful story,.
it tells you that providence includes getting lost..
Have you ever got lost?.
Providence includes getting lost and wasting time,.
wasting time..
Have you ever, you know, stand in line in DMV?.
That's like the most frustrating hell on earth, right?.
Providence includes wasting time..
Now, how many of you are good chess players?.
How many of you are good chess players?.
I'm very bad..
I just don't know how it works, right?.
When I was younger, I hate playing chess with my dad.
because he's much better..
He's a good chess player..
And I hate playing with him because I,.
well, you know, playing with my dad,.
a better chess player, and I'm so bad,.
I am absolutely free to make any move..

$^{441}$But my opponent often puts me into some sort of situations,.
some sort of positions that I have to respond.
with another move, right?.
Some sort of moves..
But I find that quickly, my moves lead to further.
his strategy to have him win..
And he, of course, he wins at the end..
Or have you ever played chess with computer?.
AI, can you beat it?.
Probably not, right?.
Compared to God, we are all like bad chess players..
We're really, really bad..
Because we can only see a step or two ahead..
But God see the whole thing..
God see the whole game..
But there's a difference..
No, me and God, you and God, in his providence,.
there was a difference..
Because God's moves in your life.
are not mechanically programmed..
He's relating to you, interacting with you all the time..
And another difference is God is not in your life.
in order to defeat you..
That is not the intent of our Karen father..
Our Karen father is driven by his compassion for us..
His loving plan for you and for the world, right?.
Don't you find this doctrine, providence,.
comforting, encouraging, especially in times of difficulty?.
That you are not under the control of any blind faith..
You're not under control of any blind destiny.
that is predetermined..
Yes, there are things out of your control..
There are things that you pray for,.
you hope for, didn't happen..
But your life, the Bible says,.
is not under some sort of blind faith..
The school that you're going to,.
or the school that you apply and rejected you, right?.
The girl you try to approach, rejected you..
I had that a lot when I was younger..

$^{481}$The people that you try to relate to,.
give you a cold shoulder, or whatever you encounter..
Don't you find it comforting.
if you know the biblical teaching of providence?.
Right?.
Now in the history of the church in the last 2000 years,.
many people, theologians,.
have tried to understand this mysterious,.
mysterious but comforting doctrine..
And they try to take through different places of the Bible.
and try to summarize it, right?.
And one of the great summaries of providence.
is found in something called.
Westminster Confession of Faith..
Now, what is Confession of Faith?.
It's actually a kind of a systematic summary.
of biblical teaching..
And this document is a very comprehensive document.
about all things related to the teaching of the Bible..
And it was drafted in the 17th century.
in the beautiful Westminster Abbey in London..
In chapter five, article one,.
here is what it says about providence..
Now I want you to pay attention to some of the words, okay?.
Let me read that..
Five one, Westminster Confession of Faith..
God, who created everything, also upholds everything..
He directs, regulates, governs every creature,.
action, and thing..
From the greatest to the least..
By his completely wise and holy providence,.
he does so in accordance with his infallible foreknowledge.
and the voluntary, unchangeable purpose of his own will..
All to the praise of the glory of his wisdom,.
power, justice, goodness, and mercy..
God orchestrates from the greatest.
to the least..
Everything is under his sight..
Nothing is out of his vision, his providence..
Here, this summary is only one paragraph.

$^{521}$out of the many paragraphs that talks about God.
and his providence in this document..
Here it says, God, the creator,.
transcends the creation..
He's beyond us..
But notice this, transcendence means his works.
are mysterious at times, right?.
Because he's beyond us..
He's outside time and space..
This is our God..
This is the God of the Bible..
He created us all..
He orchestrates everything in our journey..
Because he's beyond us, sometimes his works are mysterious..
But not irrational, because transcendence.
does not mean his absence, right?.
His beyond us does not mean God is absent..
He's here..
He's here in your life..
And secondly, from this paragraph,.
we notice, or chapter 37 of Book of Genesis,.
God never condones evil and sin..
We are all accountable, just as the older brothers of Joseph,.
accountable to God..
But God can use the good deeds of the people.
to accomplish his will..
He, God, God can use the evil actions of the people.
to accomplish his holy plan..
God can use both..
Use our good works and the evil plans and actions of others..
Number three, God orchestrates, sometimes directly.
and often indirectly, all the things.
from the greatest to the least in our life..
But he's not inactive..
He is not silent..
He is always there..
And this paragraph says there's a purpose to it..
The purpose is for the praise and worship of God..
If we really understand and really sink in.
this teaching of providence, we should find comfort.

$^{561}$and peace, even in times of difficulties and uncertainty..
Yes, we all have our moments that we doubt.
what is God doing and where am I heading, right?.
Especially when we think of our future, right?.
But the teaching of providence encourages us,.
even invites us to join the plan of God,.
to participate in the advancement of God's plan,.
preaching the gospel, building up the church,.
serving in obedience, working with God,.
cooperating with God..
He's orchestrating everything in my life,.
and that's why I can surrender my will to him..
I can pray to God, God use me, right?.
Guide me, help me follow you, help me obey you, right?.
All the things, what does it say?.
From the greatest to the least,.
by his completely wise and holy providence..
We're gonna end today with our response song,.
and I wanna invite the worship team to come up here..
Great is our God, and because he's so great,.
that's why we bow down and fall down before him to worship..
That's why we can lay our lives before him,.
because he is greatest, and he is holy and good..
We can trust in him..
No matter life throw at you, you can still trust..
Let's worship God, and then we'll close in prayer..
♪ Come let us praise him ♪.
♪ Let us kneel at the throne of our God ♪.
♪ Through his son our salvation was bought ♪.
♪ With mercy and grace ♪.
♪ Come let us bow down ♪.
♪ In his hands are the depths of the earth ♪.
♪ With one voice we proclaim his great worth ♪.
♪ Lord our God ♪.
♪ And we will seek him ♪.
♪ Our rock of salvation ♪.
♪ Morning by morning ♪.
♪ With thanksgiving we come ♪.
♪ And we will bow down ♪.
♪ With creation we cry out ♪.

$^{601}$♪ In daylight and darkness ♪.
♪ We sing to the Lord ♪.
♪ Great is our God ♪.
Let us pray..
Father, once again, we give you thanks.
because everything in our life,.
from the greatest to the least,.
is under your careful watch..
Thank you for being our Lord and our Savior..
Thank you for this comforting truth.
that we can lay our lives before you..
God use us, send us..
Send us to where you want us to be,.
to be your instrument..
And in times of difficulties and uncertainties,.
we pray that we will trust in you..
Help us to know that you are our greatest God.
and in your providence,.
your will will be done in Christ's name..
God bless you..
- Peace Church, and I hope that you can join.
our Bible study on the Gospel of Mark later..
And let's get together and study the word of God.
and get to know God better, all right?.
God bless you..
(gentle music).
[MUSIC PLAYING].
\newpage



\section{創世記}
\label{sec:Fs4lhTNde4E}
\textbf{上主的護理 (2) 細想約瑟的故事 2 (創世記 39) - 陳崇基牧師}
\newline
\newline
連結: \href{https://youtube.com/watch?v=Fs4lhTNde4E}{\texttt{https://youtube.com/watch?v=Fs4lhTNde4E}} ~~~~ 語音日期: 2023-09-26
\newline
\newline
\hyperref[sec:mJHLbAHSJBg]{\small{< < < PREV SERMON < < <}}
~
\hyperref[sec:index_chronic]{\small{[返順時目]}}
~
\hyperref[sec:index_scriptual]{\small{[返順卷目]}}
~
\hyperref[sec:Jzl6PyqrviE]{\small{> > > NEXT SERMON > > >}}
\newline
\newline
$^{1}$.
Anyways 我們繼續看這個主題叫做Providence.
上帝的護理或上帝的庇護.
什麼意思?Providence.
我們曾經解釋過這個字.
即是事先看見的意思.
即是上帝在你心靈當中大大小小的事情.
祂已經一早看見.
並且祂是去照顧著的.
這個是聖經裡面一個很寶貴的教導.
當我們心裡面有驚慌.
或者我們不知道下一步要怎樣走的時候.
我們會有安穩.
因為我們知道上帝有Providence.
在過去我們心靈裡面經歷過的事.
我們怎樣迂迴曲折都好.
我們都知道原來是有上帝自己的帶領著.
我們不是被盲目的一些命運去操控著.
或者擺佈著我們.
Providence 聖經裡面非常重要的一個教導.
所以我們繼續看這個約瑟的故事.
今次我們第二次看到約瑟的故事創世記念.
39章.
上一次我們看37章.
今次我們看39章.
39章開始.
一個很熟悉的故事.
我們不如整段來讀一讀.
39章第一節.
約瑟被帶下埃及去.
有一個埃及人波提法.
是法老的官員是護衛長.
他從那些帶約瑟下來的二十萬人手中把約瑟買了.
約瑟在埃及主人的家中.
耶和華與他同在.
他是一個通達的人.
他主人見耶和華與他同在.
就見耶和華使他手裡所辦的事都盡都順利.
第四節.
約瑟就在主人眼前蒙恩.

$^{41}$事後他主人主人派他管理家務.
把一切所有的都交在他手裡.
自從主人派約瑟管理家務和他一切所有的.
耶和華就因約瑟的緣故賜福那埃及人的家.
凡家裡面田裡面一切所有的都蒙耶和華賜福.
波提法把他一切所有的都交在約瑟的手中.
除了自己所吃的食物.
其他一事都一概不知.
然後第六節.
下半節.
這裡是我的翻譯.
因為原文是有個break在這裡.
其實原文作者就說要留意.
現在有一個新的段落.
現在你要留意.
約瑟是英俊健美.
這裡很有趣.
其實意思就是說.
首先他說他身材很靚.
就是他去健身.
也都說他樣貌很靚.
這個中文翻譯就是英俊健美.
留意了.
他是英俊健美的約瑟.
這時以後.
約瑟主人的妻子以沐宋情給約瑟.
說你與我同寢吧.
約瑟拒絕與他主人的妻子說.
看下一切家務.
我主人一概不知.
他把所有的都交在我手中.
在他家中沒有一個比我更大.
除你以外.
他也沒有留下一樣不教給我.
因為你是他的妻子.
我怎能行這大的惡.
得罪上帝呢.
他天天這樣對約瑟說.
約瑟他不聽從他.
不與他同寢.

$^{81}$也不和他在一起.
第十一節.
有一天約瑟進屋裡去辦事.
家裡沒有一個人在那屋子裡.
婦人就拉著他的衣服說.
你與我同寢吧.
約瑟把衣服留在他手裡.
逃出外面去了.
婦人看見約瑟把衣服留在他手裡.
逃到外面.
就叫了家裡的人來對他們說.
他帶了一個希伯來人到我們這裡.
戲弄我們.
他到我這裡與我同寢.
我就大聲喊叫.
他聽見我放聲大喊.
就把他的衣服留在我這裡.
逃出外面去了.
婦人把約瑟的衣服放在身邊.
直等到他主人回家.
就用這樣的話對他說.
你帶來這裡的希伯來僕人.
進來要調戲我.
我放聲大喊.
他就把衣服留在我身邊.
逃到外面去了.
主人聽見他妻子說.
對他說者說.
你的僕人就是這樣對我.
就非常生氣.
約瑟的主人就把他抓起來.
關在監獄裡面.
就是黃的囚犯被關的地方.
於是約瑟在那裡坐牢.
但耶和華與約瑟同在.
向他示恩.
使他在監獄長的眼前蒙恩.
監獄長就把監獄裡面所有的囚犯.
都交在約瑟手下.
在那裡一切事都由他處理.

$^{121}$任何交在約瑟手中的事.
監獄長一概不察.
因為耶和華與約瑟同在.
耶和華使他所做的都順利.
一個很熟悉的故事.
就是約瑟結果被埋到埃及.
然後他本身是一個亞國家裡面.
最被愛寵的一個兒子.
變了一個奴隸.
然後好好的.
結果就是由低做起.
但因為這個惡毒的女人的緣故.
就被引誘.
然後就弄到他一身癌.
沒有犯大錯也好.
他仍然守住他的誠信.
我們很多時候看這個故事.
我們都會當為一個道德的故事.
應該是有這樣的意思的.
一個年輕人.
去到一個異地.
一個異鄉.
他可以抵受著這些人遊.
這裡經文有沒有說過.
波提法的妻子漂亮嗎.
說不定她不漂亮.
怎樣都好.
有沒有說她年輕還是年長.
說不定她年長.
誰知道.
怎樣都好.
這裡講到他受的試探.
性的試探.
做一個很有誠信的人.
不願意犯大罪.
去得罪上帝.
但這個故事我們直接看.
應該是一個好故事.
我相信大部份的兒童聖經.
都會包含這個故事在裡面.

$^{161}$會包了這個故事在裡面.
但如果你看清楚的時候.
好像又不是很合理.
如果純粹是一個道德的故事.
因為你應該不會這樣教兒子.
你不會跟兒子女兒說.
告訴你故事.
臨睡前.
主角是一個很好的男孩.
品學兼優很乖的.
還很帥.
最後被一個女人害.
然後就弄得一身癌.
完結了.
你不會這樣說一個小朋友的故事.
如果這樣寫一個小朋友的故事.
或者一個道德的故事.
看完之後就睡覺.
你睡不著.
怎麼睡得著.
因為他二人受害.
所以你這樣看的時候.
這個故事不只是這麼簡單.
不是這麼簡單.
約瑟是一個好兒子.
好兒子會有好報.
這裡不是這樣說.
好兒子好像不是一定有好報.
當然最後他有好報.
但為什麼有一個這樣的故事.
我們說過希伯來人的文化.
不是我們西方那種.
或者現代那種.
我們列出一點一點來解釋.
說到上帝的護理.
Providence的時候.
很多時候是用一個Story.
一個故事.
來向你敘述.
所以上一次我也提了.

$^{201}$為什麼兒童的故事的書.
如果買一本好的話.
買一批好的話.
是很好的.
對小朋友來說.
是非常幫助他們.
起碼人格上的成長.
好的故事書.
以前我太太在.
我們讀神學的時候.
我們已經結婚了.
我讀她也讀.
她讀了一堂很有趣的.
她自己也覺得很有趣.
每次完結的時候.
晚上她做功課的時候.
我都八卦去看她做什麼功課.
那個是什麼課程呢.
就是講.
Children's books.
就是兒童故事的神學.
那個很有趣.
因為我成長的時候.
沒有人跟我講這些故事.
看來看去都是TVB爛鬼東西.
怎麼會學這些有道德精神的故事呢.
所以她讀了一堂的時候.
我覺得她很有興趣.
她就說在過去的幾十年也好.
幾百年也好.
有很多最好的兒童的故事書.
是什麼樣的故事書.
還有怎樣衡量這些故事書.
所以說回故事書.
這裡好像不是這麼簡單.
不是純粹去說.
好男孩就會有好報.
應該怎麼看.
這裡有兩點.
第一點.

$^{241}$講到Provenance的時候.
講到上帝的護理.
是超越我們人生的苦難.
為什麼會這樣說.
如果看整段的經文.
大概可以分成三段.
剛才在第六節我將它分開.
事實上可以這樣分開三段.
第一段.
就是講到約瑟在波提法的家裡面.
他有亨通.
凡他所做的事主人都交給他.
然後有一句話經常出現.
就是耶和華與他同在.
耶和華使他順利.
進度順利.
這個其實就是詩篇第一篇裡面所講.
一個倚靠神的人.
就是進度順利.
他遠離惡事.
然後倚靠上帝.
聽耶和華的律法.
進度順利.
然後第六節下半節到十九節.
就是交代他在波提法的妻子的手下.
發生了什麼事.
你留意這裡經常提到.
手和衣服.
那件衣服本身是誰的.
是約瑟的.
當然他已經沒有那件彩衣.
但他有衣服.
他就拉他的衣服.
想去引誘他.
然後就搶了他的衣服.
就經常說著留在他手裡.
那個淫婦的手裡.
他衣服留在淫婦的身邊.
經常提到這件事.
每逢聖經裡面提到.

$^{281}$在手裡的時候.
你可以說是.
將它翻譯成為power.
under her power.
在他主人的妻子的權力下.
他是不能走的.
是無助的.
已經是一個奴僕了.
已經是幾無助了.
誰知道在波提法的妻子的手下.
在他的權力下.
受他控制.
這個情況不可以比這個更壞.
你想想他本身是一個愛子.
幾乎天天都在家裡不用工作.
然後就被人糟蹋.
丟到坑裡.
還是不夠差.
然後還被人賣到埃及.
你會看到這個故事一直下降.
現在還落在一個女人的手下.
被他控制.
被他誣告.
終於誣告成功.
20-23節第三段.
他們被抓了.
然後坐牢.
但又有那句話.
而我為他同在.
而我為他所做的盡都順利.
我們說到providence的時候.
我們經常要處理一個問題.
處理一個問題就是上帝的主權.
因為說到上帝的護理.
先見.
seeing ahead of time.
就一定是說到上帝的主權.
他的sovereignty.
還有人間裡面.
發生一些不如意的事.

$^{321}$苦難問題等等.
聖經如何化解這件事.
神的主權和人的經歷的際遇.
聖經當然不是說.
我們是一個機械式的.
沒有自由沒有選擇的一種生活.
我們每個人都有我們的自由.
上一次說過.
我們每個人都有我們的責任.
但是每逢有一些不如意的事.
苦難的事.
臨到的時候.
到底聖經如何化解這個問題.
當然聖經沒有一個完整的.
或者一個完全解釋所有問題的答案.
我們可以說.
我們不能夠接受的就是.
也不能夠接受的.
就是一個串線.
串林的上帝.
是引致或者驅使.
cause.
導致這些苦難.
臨到我們的.
聖經是不會贊成的.
或者我們已經去到語言的限度.
人的語言不能夠去描述.
上帝的主權.
還有人間發生的事.
有我們的自由有我們的責任.
然後有苦難出現.
去到幾乎我們已經去到語言的盡頭.
我們不能夠說.
上帝是導致或者驅使這件事.
我們只能夠說.
上帝容許.
或者容忍.
tolerate.
或者allow.
這些不好的事.

$^{361}$邪惡的事.
發生在我們的身上.
約瑟.
在這個淫婦的手下.
我們不能夠說.
上帝促使這件事發生.
He did not cause it.
極其量我們只能夠說.
上帝是allow.
去到容許這些事.
發生在他身上.
聖經我們說過了.
包括聖經本身的文化.
從來都沒有一個很緊密的.
神密的一種理論.
來扔給我們.
將上帝的主權和人所經歷的苦難.
兩者之間的張力完全化解.
是沒有的.
事實上過去這麼多年.
哲學的思想都沒有一個人能夠化解.
是不是.
說真的.
化解了又如何呢.
解釋給你聽.
為什麼有些陰陋的事發生在你身上.
或者濕滯的事發生在你身上.
So what.
是不是.
為什麼你會生病.
好啊幫你解釋了.
那又如何.
你仍然會有病.
是不是.
聖經沒有這麼神密的哲學理論.
或者甚至外面有任何一個哲學理論.
來解釋上帝的主權和人間苦在的問題.
相反聖經做了什麼呢.
聖經就是用故事.
來跟你說這件事.

$^{401}$因為故事你可以再聽再看.
你再嘗試去領略裡面的一些細節.
所以它就將弱勢的故事帶出來.
我們大家都知道story.
是有它那個解說的能力.
是多過理論的.
是不是.
Tell me a story.
Let me tell you a story.
故事裡面.
再幫助你去明白.
去解釋到底發生什麼事.
所以看這三段的時候.
其實可以這樣加個標題.
什麼標題.
一至六節上半節.
就是上帝同在.
然後中間那段.
就是所謂月缺的時候.
沒有說上帝同在.
你有沒有留意.
當他們遇見這個苦難.
遇見這個隱憂的時候.
沒有說耶和華與他同在.
第三段.
你回去皇的監獄裡面亨通.
又說回神耶和華與他同在.
似乎告訴我們有一樣東西.
即使有難.
上帝好像隱藏了.
神不在.
或者應該說得清楚一點.
神好像不在.
好像不在.
是不是.
上帝仍然在庇護.
故事就告訴我們一個活生生的證明.
I tell you a story.
約瑟的故事.
就是告訴你上帝.

$^{441}$即使在他有難的時候.
彷似是隱藏.
找不到上帝在哪裡.
但上帝仍然在.
你有沒有留意.
整個故事的脈絡.
約瑟的地點是有改變.
我們經常留意.
總之有陰陋的事情發生在他身上.
但當故事由開頭到最後.
39章開始到39章尾的時候.
那個位置是鮮明有分別的.
他moving on.
因為39章開始的時候.
他在波提法埃及人的家裡面.
到最後的時候.
約瑟在哪裡.
在王的家裡面.
王第一次出現.
The King.
當你繼續講下去的時候.
你就發現原來上帝的計劃.
是繼續推進人.
這個只不過是一個插曲.
或者只不過是其中一集.
將來會發生什麼.
這39章已經給了一些clue.
一些顯示給我們.
約瑟是moving on.
你要繼續明白約瑟的故事是怎樣.
你就繼續看下去.
40章41 42.
但很明顯他在向前推進.
去到王的監獄裡面.
大家知道我們已經劇透了.
最後約瑟是做了什麼.
是王的大臣.
怎樣可以做到王的大臣.
39章就給你一個movement.
上帝的護理是繼續向前走.

$^{481}$即使好像他是隱藏.
但上帝的計劃是繼續advance.
第二點.
上帝的護理的時候.
亦都講到約瑟人生裡面的浮沉.
地位高與低.
為什麼這樣講呢.
我們試一下.
不要只看39章.
我們試一下抽出來.
看一看附近到底發生什麼事.
由37章開始約瑟的故事.
到現在我們看39章.
37 38 39.
到底有什麼事發生.
37章你看到.
有他的長袍或是彩衣.
本身是亞國給他最愛的兒子.
有件彩衣長袍.
然後他就被奪去他的長袍.
然後亞國最後被騙.
有些兄弟騙他爸爸.
就把那件彩衣給他.
有血沾上的.
你就看到是一種honor.
本身是有honor.
有件長袍有件彩衣.
然後就變成了屈辱.
羞辱.
shame.
然後我們知道到41章的時候.
他又有一個榮耀.
42章就講到他從監獄裡面出來.
然後法老給他一件衣服.
他就做了大臣.
由榮譽變成羞辱.
又由羞辱變回榮譽.
38章了,這裡就很奇怪了.
突然間插一個猶大的故事.
猶大和他媽的故事.

$^{521}$這個很奇怪,我們經常都struggle.
為什麼37章結束了,為什麼不直接跳去39.
為什麼要插一個38呢.
原因就是整個創世記特別是最後幾章.
就要交代到底長子的名分是給誰.
如果你看之前那幾章的話.
你就知道亞國有多少個兒子.
12個位.
頭3個已經是DQ了.
裡面因為跟他父親的妾去到同寢.
想拿家產,想secure那個家產.
裡面已經是DQ了.
然後西勉和利美.
你記得有一個故事,是很殘忍的故事.
就是他們去到事件的時候.
將那個城洗劫.
是兇殘的兩個兄弟.
所以第二第三個兒子.
老二老三都已經再DQ了.
因為他們是不應該得到繼承者的權利的.
disqualify.
第四個是誰.
就是猶大.
所以你記得賣藥室那一幕,37章.
之前我們上個星期講過.
是誰出主意的.
是猶大.
為什麼猶大出主意.
因為四公子知道應該到我了.
應該是我分相分的時候.
而這個穿著彩衣的那個.
就是在那裡也問也沒有.
因為家父是有資格去到隨便挑選.
或者按照他心意去挑選.
誰去做繼承人.
第四個兒子應該就是穩坐的那一位.
所以他就說不如將他賣了.
因為十二個兒子裡面最憎藥室的.
應該就是第四個兒子.
所以就加插了猶大這個故事.

$^{561}$有些細節我們不講.
猶大就去了另外一個地方.
然後有一個媳婦叫做他媽.
這裡有一些語句或者一些出現的物件.
又是在之前那一張和之後那一張出現.
他媽因為不同的緣故.
她就做了寡婦.
然後因為利未人那種制度.
你要為你亡夫來留名.
所以猶大應該給他們其他的兒子.
來去許配給他.
然後跟他們結婚.
然後為他的名.
亡夫的名去留下.
但是猶大自私不肯.
大概是這樣.
但最重要我想你留意的就是.
他媽就脫去她的.
母法之際.
她不知道怎樣去為她亡夫來留名.
她就脫去.
又是衣裳.
就是她寡婦的衣裳.
然後她就扮成一個妓女.
最後她就生了兒子.
因為猶大不肯給她.
你見到他媽.
又是由榮譽變成羞辱.
又是由寡婦扮成一個妓女.
然後最後她就得子.
地位又由低變回高.
最後他媽生了誰.
是法勒斯.
法勒斯繼續生下去的時候.
第十代就是大衛君王.
所以由榮譽變成羞辱.
又由羞辱變回榮譽.
三十九章我們今天看這章.
再講約瑟.
又是四公子和約瑟之間.

$^{601}$到底誰是有繼承權呢.
再講約瑟.
你見到約瑟整個故事裡面.
又是經常出現衣服.
鏡頭經常放在衣服裡面.
原因是三十七,三十八,三十九都是講衣服.
以前的社會文化.
衣服是甚麼.
衣服是地位的象徵.
特別是穿一件漂亮的衣服.
以前窮的古代的社會.
不是每個人都有穿Money.
不是每個人都有名牌的衣服.
名牌衣服只不過是百分之一的人有穿.
穿著那件衣服就代表他的地位和權力.
尊榮是最高的.
衣服就是有一個這樣的.
所謂社會裡面的功能.
或者社會的標記.
八九章的故事.
如果你這樣看的時候.
不是只是講好男孩就有好報.
而是講我們人生.
講約瑟的人生.
是有浮沉.
由他最高的尊貴有件衣服.
到他沒有了那件衣服.
然後上帝的計劃繼續推進的時候.
他又有那件衣服在穿.
這個尊貴告訴我們.
告訴你.
生命裡面一定是有浮沉.
不過你不是被一些盲和啞的命運去戲弄.
約瑟不是.
你也不是.
你不是被一些盲目的化學作用.
或者一些物理的基建.
在這個世界裡面.
整定你一些東西.
然後玩死你.

$^{641}$不縱使你未必可以解釋所有的東西.
但是這個故事想告訴你.
在上帝的呼應下.
你可以放心.
上帝的計劃裡面.
你未必看到全套.
不過有他開車的話.
你就不用怕.
你有沒有教過人開車.
我教過兒子開車.
開棍波車.
我跟他說你怎樣都要學棍波車.
棍波車好啊.
沒人搶.
我記得一開始教他開棍波車.
我心都寒.
你教人開車.
雖然你沒有腳掣.
你都踩下去.
就因為你知道那個新手.
他不夠經驗.
慢慢他夠經驗的時候.
你就很安樂.
現在他開始有經驗.
坐他車的時候我就很安心.
通常.
這個就是Providence.
如果你說你的生命裡面.
你的體理人生.
只不過是一些很隨機的東西.
沒有意義的東西.
完全撞在一起.
你的工作遇見不同濕濟的東西.
你家庭你自己身體.
你自己的人生的旅程.
遇見很多不同的障礙.
或者浮與沉.
越圓越缺也好.
Providence就告訴你.
你不是一個命運所玩死你的.

$^{681}$是一個人生.
Yes我們未必看得完.
不過.
Providence告訴我們.
你坐那個車.
是最好的司機.
會帶你去到他要你去的地方.
目的地.
我不知道你有沒有在生命裡面.
有一些禱告的交戰.
在禱告的交戰.
我相信我們每個人都有.
過去大概一年左右.
我自己心裡面都有很多交戰.
當然我經常都有.
不是一些很特別的東西.
但過去一年我經常在交戰.
苦討幾個月的東西.
我經常在臨婚的時候.
我都會跟上帝說.
上帝你在搞甚麼.
我只是做一個很簡單的禱告.
而那個禱告不是給我的.
我沒有任何得益.
我是為教會禱告或其他人禱告.
我是沒有得益的.
我實在很多個月.
好像跟空氣說話.
有些東西是不在我的控制之中.
正因為我知道是不在我的控制之中.
所以我才努力去禱告.
為甚麼要跪在上帝面前祈禱.
因為我知道我不是在控制之中.
你是開車的那個.
為甚麼不帶我去一個我想去的地方.
我相信你一定有試過這樣的交戰.
上帝你在搞甚麼.
我試過有一段時間.
我簡直想放棄.
因為如果你不交戰.

$^{721}$那你不用有張力.
不過我要講回.
每個星期三晚我們有祈禱會.
其實對我來說是一個很好的提醒.
每個星期有一個小提示.
一個溫馨的提示.
因為看到弟兄姊妹.
我們網上祈禱會.
交心.
大家去禱告 大家紀念.
很難得.
每次去完我都被提醒.
禱告要行前.
我未必看到很多部.
但我知道開車的那個是誰.
有一個基督教的作家叫做Philip Yancey.
Philip Yancey寫了很多文章很多書去處理.
苦難和上帝禱告不認魂和上帝的主權.
那種交戰的一些作品.
其中一段他這樣說.
上下文他講到耶穌在赫西瑪利園的祈禱.
耶穌在赫西瑪利園的祈禱大家知道.
是第一個天父吳英允的祈禱.
吳英允他的心意的禱告.
因為耶穌想求什麼.
求天父救他脫離死亡.
但天父他做什麼.
沒有照他的意思.
照著天父自己的意思.
這段他說什麼.
他說有多少次我祈求一件事.
他得另一件事.
大家應該有同感.
你求A為什麼會出了B,C,D.
他說只有上帝才有資格回應禱告.
只有上帝才有qualification去應運禱告.
有時候我們必須要放手.
並且相信上帝.
他能夠將我們以己為中心的祈禱.
轉化成想像不到的更大的計劃.

$^{761}$當耶穌能夠救他脫離死亡的天父祈禱的時候.
他沒有得著哪個拯救.
相反他得著的是一個更大的拯救.
世人能夠得救.
他沒有在死裡得救.
耶穌卻在死裡得到世人的得救.
約瑟的故事其實都是一樣.
發展到這裡.
同一道理.
上帝沒有立刻救約瑟離開埃及.
但上帝最終藉著他救了以色列人.
甚至當時饑荒的人.
是上帝有一個更大的計劃在他身上.
神的providence就是教我們.
Yes你可以安心.
因為上帝有一個bigger plan.
也教我們如果今天月圓的話.
去教你.
不單止你可以放心.
你更加可以奉獻自己.
去靠上帝.
更加值得你去愛上帝.
做他的僕人.
我們結束祈禱之前.
我們唱我們的回應詩歌.
我希望回應詩歌.
能夠再一次帶我們去到上帝面前.
因為唱的時候.
這些禱告.
我們有的跟上帝說的話.
讓你知道上帝在他的providence裡面.
在他的照管下.
護理下.
他有一個更加美的計劃.
一個更加大的計劃.
《上帝的禱告》.
他靠你翅膀任行.
在你光中不必見怕.
圍獨敬畏你的方向得以靠.
別錯自以來轉向你.

$^{801}$雄心即將發達至善.
是我轉做著為王.
願天歡喜.
願你想快樂.
願你每次上台.
但求我主的必不受愧.
是你甘願與我自然一切.
我們跟心底留結.
主我們多謝你.
因為你是那位照管一切的神.
天下萬物.
所有我們所經歷的際遇.
都在你自己的愛護.
你自己的恩典當中.
上帝再一次提示我們.
提醒我們.
我們每一個禱告.
每一個的際遇.
都有你自己更美的計劃在當中.
主我們短時未必能夠見到.
在你面前艱難的時候.
有什麼出口.
但我們相信神你依然同在.
神你依然與我們同行.
主叫我們來依靠你.
叫我們來相信你.
教導我們如何去服侍.
去敬拜你.
聽我們有聲無聲的禱告.
因為你是我們最美善.
全能的神.
我們人生在你的手中.
沒有走錯路.
主一靠你繼續帶領我們.
祈禱奉傳給大家.
願上帝賜給各位平安.
\newpage



\section{3:1-12}
\label{sec:Jzl6PyqrviE}
\textbf{Words Of Life (James 3\_1-12) - Caleb Lee}
\newline
\newline
連結: \href{https://youtube.com/watch?v=Jzl6PyqrviE}{\texttt{https://youtube.com/watch?v=Jzl6PyqrviE}} ~~~~ 語音日期: 2023-10-03
\newline
\newline
\hyperref[sec:Fs4lhTNde4E]{\small{< < < PREV SERMON < < <}}
~
\hyperref[sec:index_chronic]{\small{[返順時目]}}
~
\hyperref[sec:index_scriptual]{\small{[返順卷目]}}
~
\hyperref[sec:VkqqimXzBU4]{\small{> > > NEXT SERMON > > >}}
\newline
\newline
$^{1}$Good morning, everyone..
It is great to be here on another Sunday morning with all of you to worship with you, to take.
communion with you, and to just enjoy the presence of this family that we've been adopted.
into..
If I could, let me open us in a word of prayer..
Would you bow your heads with me as we pray?.
Our Heavenly Father, our Lord, and our God, we come before you today not by what we have.
done, not by our own merit, but because you saw us first and you loved us first, Lord,.
and you have brought us into your presence..
And so, God, this Sunday morning we give to you, I ask that you would just calm our hearts,.
Lord, and give us minds that are focused..
Would you clear away the clouds of sleep, of the troubles of this week?.
Would you focus our attention upon you?.
And so for that, Lord, I ask that you would be with me, that you would speak through me,.
and would you protect the ears of my brothers and sisters as they listen to what I say about.
your word?.
Would you help them remember your word and your word alone?.
All these things we pray in Jesus' name, amen..
Okay, so we'll be continuing our reading of James today, this letter from the Apostle.
James to the Jewish Christians of the dispersion, where he calls for his brothers and sisters.
to live lives according to the faith that they profess..
To be more than just believers in name, but also believers in their living..
In James' words, Christians cannot just be hearers of the word of God, you and I must.
also be doers..
Last week, we saw in chapter 2, verses 14 to 26, that it is a living faith that saves..
That a faith which is without works, without any impact on a person's life, is like trying.
to live without breathing..
It's a life that's incomplete, missing an essential part, an essential aspect, and it.
is one that cannot sustain you..
It's wrong to live in this way, and it's impossible to live in this way..
The only outcome from a faith without works is death..
And so, James calls for his brothers and sisters to have a complete faith and be complete people..
He wants his fellow believers to be perfect, affected by their faith and transformed so.
that their entire being both agrees and works in cohesive wisdom and righteousness..
Toward that end, we are brought to the next part of his letter to his displaced brothers.
and sisters..
If you have your Bibles with you, I invite you to turn with me to the book of James,.
it's third chapter, verses 1 to 12..
Again, that is James chapter 3, verses 1 to 12..
If you don't have a Bible with you this morning, they'll be right here..

$^{41}$So let's read, all right?.
James chapter 3, verses 1 to 12..
Not many of you should become teachers, my fellow believers, because you know that we.
who teach will be judged more strictly..
We all stumble in many ways..
Anyone who is never at fault in what they say is perfect, able to keep their whole body.
in check..
When we put bits in the mouths of horses and make them obey us, we can turn the whole animal..
Or take ships as an example..
Although they are so large and are driven by strong winds, they are steered by a very.
small rudder wherever the pilot wants to go..
Likewise, the tongue is a small part of the body, but it makes great boasts..
Consider what a great forest is set on fire by a small spark..
The tongue also is a fire, a world of evil among the parts of the body..
It corrupts the whole body, sets the whole course of one's life on fire, and is itself.
set on fire by hell..
All kinds of animals, birds, reptiles, and sea creatures are being tamed and have been.
tamed by mankind, but no human being can tame the tongue..
It is a restless evil, full of deadly poison..
With the tongue, we praise our Lord and Father, and with it, we curse human beings who have.
been made in God's likeness..
Out of the same mouth come praise and cursing..
My brothers and sisters, this should not be..
Can both fresh water and fresh water be the same?.
My brothers and sisters, can a fig tree bear olives or a grapevine bear figs?.
Neither can a salt spring produce fresh water..
Do you want to be what James calls his brothers and sisters to be?.
A person who is wise and complete and whole?.
A person perfected by their faith?.
Here in James chapter 3, verses 1 to 12, we see the centrality of the tongue in being.
perfect..
He starts off this chapter by advising caution when considering a position of teaching among.
believers, right?.
Verses 1 to 2, "Not many of you should become teachers, my fellow believers, because you.
know that we who teach will be judged more strictly..
We all stumble in many ways..
Someone who is never at fault in what they say is perfect..
It would keep their whole body in check.".
In speaking about strict judgment, we begin to think about the standard of God..
It's the standard of perfection..

$^{81}$And what James states is that among the many ways in which we fail to keep this standard,.
we most commonly and often fail through what we say..
We are prone to mistakes, and one great maker of mistakes is our tongue, to the point where.
James says that perfection in speech translates to perfection in character..
There's a correlation being drawn between a self-controlled tongue and a self-controlled.
body..
So you see in verses 1 to 2 that the tongue is important, that the tongue holds a key.
place in the direction of your living..
In order to help us understand this, James gives us two illustrations in verses 3 and.
4, a horse and a boat..
There we go..
If we put bits into the mouths of horses so that they obey us, we guide their whole bodies.
as well..
Look at the ships also..
Though they are so large and are driven by strong winds, they are guided by a very small.
rudder wherever the will of a pilot directs..
That is a different version, but here we have two examples..
One is a living creature, the other is a man-made vehicle..
These are large things..
Horses and boats move with immense power..
Have you ever seen a horse?.
Anyone?.
They're big..
Okay, I don't know how else to describe them, but they are much bigger than I pictured them.
to be..
When we quantify the strength of modern-day machinery, what do we use?.
What sort of word do we use?.
We use the word "horsepower" because horses are this measurement of strength..
We tell people, "This machine has this many horses' worth of power.".
By using a bit, which is a metal rod that fits across the inside of a horse's mouth,.
we're able to steer the horse..
We're able to cue and direct the horse in the direction that we want..
A small tug lets the horse know which direction it should turn..
The entire body of the horse is guided by this bit, and it's the same with a boat..
We're given this imagery by James of a large boat driven by strong winds, although they.
are so large and are driven by strong winds..
It's a large boat..
This boat is able to withstand the force of nature..
It takes the power of these air currents to be pushed across the waters..
And yet, using a rudder, a pilot is able to influence and change the direction of that.

$^{121}$force..
The boat is steered by the will of the pilot..
And so we're given these examples to help us understand the effect of words..
Because the tongue holds a key place in the direction of their living..
It is the means by which even a body of great size is directed and controlled..
It is a small member of the body, but it makes great claims, which turns the body in unholy.
and ungodly directions..
If you would look with me at James chapter 3, verse 5, it's on the screen right here..
Likewise, the tongue is a small part of the body, but it makes great boasts..
Consider what a great forest is set on fire by a small spark..
The tongue boasts..
It makes claims, it oversteps, and it speaks in ill ways..
We see in verse 5 through the idea of boasting that this small part of the body is not leading.
its bearer in the good directions..
There is not a proper helmsman at the wheel of this ship..
Instead of works of mercy and compassion, there is only death and destruction..
What a great forest is set on fire by a small spark..
You know, since earlier this year, I think it was around March and April, there have.
been a lot of wildfires consuming Canada..
And I think Chicago got a brief taste of this smoke and ash..
It drifted in over from the north..
And now what I read, this was in August, around August, I read was that the wildfires had.
burned up 37.8 million acres of land..
That is a lot of land..
And something interesting to go with that is that a good portion of these fires were.
started by dry thunderstorms..
It's this lightning produced by the dry environment that sparks these small fires..
Fires which spread across dried out vegetation in the blink of an eye, blazing into these.
massive uncontrollable conflagrations..
They're uncontrollable..
They grow faster than they can be put out and they destroy and displace with no regard.
for anyone and anything..
James describes the tongue as a fire, not just an object that steers, that changes your.
direction, but a fire that wreaks havoc on everything around it..
The tongue is also a fire, a world of evil among the parts of the body..
It corrupts the whole body, sets the whole course of one's life on fire..
Look at that statement in the first sentence of verse six..
The tongue is a world of evil..
The present reality of wrongdoing in this world and all the unrighteousness of that.
world occurs and it finds its representation in the tongue..

$^{161}$The tongue is the part of our bodies which is used by God's enemy to do harm..
James points at it as a manifestation of evil that corrupts the whole body, the entire body,.
the part which directs the rest of our body towards a life of evil and destruction..
It is set on fire by hell..
These are harsh words, scary words..
James looks at the tongue as a power for evil..
And so our tongues, holding a key position in the direction of our living, more often.
than not directs us towards evil..
Have you ever been in a situation where you found yourself processing hurt or anger?.
We're human beings..
Of course we have..
Maybe in those situations we find friends or family to vent our thoughts to..
Maybe we even sit and mutter to ourselves going over and over again and again how the.
offense happened..
But in that process we find that we work ourselves up into a rage..
We get angrier and angrier or more and more upset as we hear our own words, not just our.
spoken words, but also our thoughts as we revisit them over and over again..
And then sometimes we even begin to take it out on the people around us and transgress.
against our brothers and sisters..
The tongue, our words, influence us..
They stain and they spread..
The tongue corrupts the whole body, causing our hands and feet to stumble..
And there is terrible news beyond this..
If you would look with me to verses seven and eight..
All kinds of animals, birds, and reptiles, and sea creatures are being tamed and have.
been tamed by mankind..
But no human being can tame the tongue..
It is a restless evil full of deadly poison..
Think about how mankind has transformed all the things around us to suit our purposes..
The multitude of animals that have been domesticated for our benefit..
Consider the blessing and command of God to Adam and Eve to subdue the earth and to have.
dominion over the fish of the sea and the birds in the sky and over every living creature.
that moves on the ground..
We have dog owners in our crowd and cat owners..
These animals are now used to living inside our homes..
Our food and our clothing come from cattle, from sheep, from domesticated pigs..
Going even further than animals, the very land on which you now, I guess you are not.
standing, you are sitting, it's been shaped by mankind..
You live in Chicago, we're here in Chicago, and this place actually used to be a swamp..
And yet, even with all this accomplishment and ability, no human being can tame the tongue..

$^{201}$It says so right there..
No human being can tame the tongue..
While the things that inhabit this earth can be tamed, the tongue cannot be tamed by man..
This is terrible and sobering news..
Instead of us controlling our tongues, our tongues control us..
Our tongues don't just direct us towards evil..
They are themselves a restless force of evil that continuously fights to corrupt every.
part of us..
They're like a poison that infects everything we do and everyone we interact with..
And this, this is wrong..
James makes it clear to you and me that this is not right at all..
The entire letter of James reminds us that those who are believers in Jesus Christ are.
not people of death, but people of life, living children of God with a living faith who are.
called to be both hearers and doers of the word of God..
In verse 9, in verse 9, we are presented with this reality of what comes out of the mouth.
of men..
With the tongue, we praise our Lord and Father..
With the tongue, we curse human beings who have been made in God's likeness..
These are contrary actions..
Praises and curses strive towards separate purposes..
And not only this, the curses are directed at those made in God's likeness..
We are people who are in the image of God, a faint reflection of the one with infinite.
worth who is worthy of all praise..
So curses against these human beings made in God's likeness slander the very God who.
was just being praised..
And these come out of the same mouth, the mouth of those supposed living children of.
God with living faith..
This should not be..
James, in his usual fashion, asks rhetorical questions again..
If you would look with me to verses 11 and 12..
Can both fresh water and salt water flow from the same spring?.
My brothers and sisters, can a fig tree bear olives or a grapevine bear figs?.
Can both fresh water and salt water flow from the same spring?.
The implied and obvious answer is what?.
No, obviously not..
One when drank will satisfy thirst, the other will deepen thirst..
These are different waters..
They are not the same..
Can a fig tree bear olives or a grapevine bear figs?.
Again, the answer is no, right?.

$^{241}$These are not things that can happen..
To borrow a line from one of my classmates, we are proved through what we produce..
Fig trees being fig trees bear figs and olive trees bear olives..
They are not the same plant, the same with a grapevine..
Grapevines produce grapes because they are grapevines..
They bear their corresponding fruit..
So why are praise and cursing both flowing from the same source?.
How can both death and life be produced by the same tongue?.
What does that prove about the one who is speaking?.
James makes this clear to you and to me, that this is wrong..
It should not be and it must not be any longer..
You and I, being people who follow the Lord Jesus Christ, who profess faith in Jesus,.
must become living children of God with living faith..
Words of evil and death have no place coming from such people..
So then, knowing that the tongue holds a key place in the direction of our living, that.
our tongues direct us towards evil and that children of God should not be guided by or.
speak evil, what should you do?.
Is the solution to work hard to control our tongues?.
Is it to, from this point on, having heard what you heard, put all your effort into gaining.
control over our words?.
Obviously not..
If you look back at verse 8, no human power can tame the tongue..
It doesn't matter how much effort you put in..
You cannot force fresh water to flow out of a saltwater spring..
So what's the solution then?.
Do you want to be a person who is wise, complete, and whole?.
A person that is directed by righteousness rather than an evil tongue?.
If that is the case, then position yourself..
Position yourself to be shaped and led by the words of God rather than the words of.
men..
Because it is not a word of life that comes out of you that transforms you..
It is words of life that are given to you that transform you..
And God has given you his word to transform you..
James himself pulls from wisdom literatures such as Psalms and Proverbs..
And here is an example..
Never mind..
Where is it?.
One second..
Okay, guess it's not there..
If you have your Bibles with me, you can turn to Psalms..

$^{281}$The book of Psalms, its first chapter..
You'll be looking at verses one to three real briefly..
It is not on the screen, and that's okay..
I'll read it out for you, all right?.
Psalm chapter one, verses one to three..
Blessed is the one who does not walk in step with the wicked, or stand in the way that.
sinners take, or sit in the company of mockers, but whose delight is in the law of the Lord.
and who meditates on his law day and night..
That person is like a tree planted by streams of water, which yields its fruit in season,.
and whose leaf does not wither..
Whatever they do, prospers..
That's the example I give to you..
If you look at this verse, you'll find that there's this person here that's mentioned..
A person who resists evil and flourishes like a tree planted by streams of water..
And they do so by meditating on the word of God day and night..
Do you know what it means to meditate?.
It's not this concept of clearing your mind that we so often think about..
You're not emptying..
You are filling your mind..
In scripture, to meditate is to mutter, to murmur to yourself, and to speak the word.
of God as you read it over and over..
It's not just a mental effort either..
It's this engaging of your entire person, an exercise of your lips, your tongue, and.
your teeth..
To meditate is to engage fully and to truly take in the word of God..
This person here who is described, the one who is blessed, who does not walk in the ways.
of sin, is not directed by evil, and instead meditates on God's law day and night..
This is Jesus..
Meditate on God's word just as Jesus does..
And this is not new..
All right?.
This is not new at all..
This is the same truth over and over in the letter of James..
Our actions and behaviors are not the reason for our faith..
They are the result of our faith..
In this chapter that we have just looked at, chapter three, James calls for believers to.
correctly display the state of their hearts through what they say..
And so, having heard all of this, having seen the importance of the tongue in shaping the.
way in which you live, while also knowing that this is a daunting fight, one that you.
cannot win on your own, allow the word of God to direct the way in which you speak..

$^{321}$Be diligent..
Be a diligent hearer of the word of God so that you become a diligent doer of the word.
of God..
And by doing so, do not prove a bitter and critical heart through bitter and critical.
words..
Instead, prove to be a living child of God with a living faith through living words..
Let me close us in a brief word of prayer..
Would you bow your heads with me as we pray?.
Father, you are the one who gives life..
You are the breath in our lungs..
You are the one who spoke first and created all of creation..
And so, Lord, I ask in this moment that you would speak to your people, that you would.
speak to my brothers and sisters, that your word would be the first thing upon their heart,.
the thing that shapes them, because it is your righteousness that we cling to, your.
work that we cling to..
And so I ask that you would continue to place your hand upon them, that you would remove.
the calluses from their heart, that you would give their ears the ability to hear your word,.
their eyes the ability to see your word..
Lord, would you place your spirit in them anew that they might come to know you in deeper.
and deeper ways..
We pray all this in Jesus' name, amen..
\newpage



\section{}
\label{sec:VkqqimXzBU4}
\textbf{弟兄中的一個最小 - 譚文倩莊傳道}
\newline
\newline
連結: \href{https://youtube.com/watch?v=VkqqimXzBU4}{\texttt{https://youtube.com/watch?v=VkqqimXzBU4}} ~~~~ 語音日期: 2023-10-03
\newline
\newline
\hyperref[sec:Jzl6PyqrviE]{\small{< < < PREV SERMON < < <}}
~
\hyperref[sec:index_chronic]{\small{[返順時目]}}
~
\hyperref[sec:index_scriptual]{\small{[返順卷目]}}
~
\hyperref[sec:lg0Kj2q_1qU]{\small{> > > NEXT SERMON > > >}}
\newline
\newline
$^{1}$我可以感到很高興..
上週,有數十萬名耶穌基督徒逃亡到地牢中..
耶穌基督徒,你們在哪裡?.
他們在那裡..
他們的家庭人有很少的食物和水,.
我們不知道有多少人死亡..
因為他們的家庭人很少,.
他們的家庭人很少,.
他們的目標是教會當地難民自己建立超級Adobe,.
說起居民和病人,.
我亦想起一個沒有笑容的小朋友,.
他因為做手術,不是第一次來,是第二次,.
他出了病沒有笑容,.
長期大叫,大叫,.
大叫叫喊的那種叫喊,.
但這次我們有一個Child Life Specialist,CK,.
她穿著紅衫的姊妹,.
她用盡各種方法跟他們玩遊戲,.
跟小朋友哄他們開心,.
當這個女孩子第一次笑出來的時候,.
我們整個護士房的人爭先擁後要進房去看她笑,.
我又想起一個雙腳炎熟畸形,.
腐爛的大嬌,她的腳是腐爛到一個地步,.
我覺得如果不看人,你會以為是一個死人的腳,.
手術一次又一次她都未能完全康復,.
很多痛苦,她又是沒有笑容,長期很失落,.
治療師本來只顧小朋友,.
但她就開始撩大嬌一起玩Domino,.
結果大嬌開始不單笑,.
她還要玩到拍檯拍燈,.
指手畫腳對著我們的治療師說Kurdish,.
如同她明白他說什麼,.
我們就好像在說笑,.
大嬌好像不停在爆粗口,.
臨走之前那幾天,.
我和其中一個翻譯出街,.
他突然轉過來問我和Ivy,.
我很想問你們,.
為什麼你們願意出錢出力來到我們這個地方,.
和我們這些素未謀面,.

$^{41}$不同種族,不同言語,.
甚至不同宗教的人當中去服侍他們,.
我當時第一個反應,.
機會來了,.
我立即和我的同工說福音,.
我們說我們是基督徒,.
我們是基督徒,.
我們的門徒奉父子聖靈的名,.
給他們施洗,.
凡我所吩咐你們的,.
都教導他們遵守,.
看那我天天與你們同在,.
直到世代的終結..
不是開玩笑,.
這三個星期真的很辛苦,.
每天差不多直落十四至十六個小時,.
大約五個小時,.
但團隊的熱誠去做這件事,.
是很令人振奮的,.
我們的團隊是一起捱,一起笑,.
我在思想時,.
我很感恩我是一個基督徒,.
最後很想大家在禱告裡,.
紀念這個Yashidi的群體,.
因為離開之前,.
我們聽到難民營將會被肅逝,.
這些Yashidi會被散落在社區裡,.
其實社區的人是不想的,.
因為他們都會普遍覺得Yashidi來到他們的群體當中,.
就會分薄了他們的資源,.
也紀念基督教的很多,.
其實我們去到都接觸了一些基督教的團體,.
當中做的事工,.
不單止是醫治疾病,.
其實他們也有很多不同的技能去教他們,.
也有療癒,.
因為大概九年十年之前他們被屠殺,.
很多人失去家人,.
或者現在本來很親密的家庭,.
散落在不同的地方,.

$^{81}$有很多心靈的創傷,.
當然最重要的是我們希望將福音帶到他們當中,.
馬太福音25章40節,.
我實在告訴你們,.
這些事你們做在我弟兄中,.
一個最少的,.
就是做在我身上了,.
我們一起低頭祈禱,.
主耶穌你是我們的主,.
是我們的王,.
你施行醫治,你施行拯救,.
求你幫助我們對難民的處境,.
有更多的同情心,.
你軟化我們的心,.
並且幫助我們跟隨你的帶領,.
為他們尋求正義和憐憫,.
我們祈禱戰爭,貧窮和侵犯人權的行為能夠結束,.
主耶穌我們祈求為難民提供生活所需,.
工作,食物,乾淨的水和安全的住宿,.
我們請求你祝福他們所居住的家園,.
讓他們能夠找到平安和歡樂,.
並且建立新記憶,.
主耶穌我們祈求保護那些特別脆弱的人,.
保護他們免受剝削和虐待,.
殘障的人士,孕婦,媽媽,.
行動不便的人,老年人,小女孩,.
和家人失散或成為孤兒的兒童,.
主阿彌撒尼,願人都清理明日性,.
願意流離失所的耶穌能夠得著福音,.
享受真正在神家裡的平安,.
祈禱奉主耶穌基督的聖名祈求,.
散會之後在外面會有天倫基金的一本書,.
我會放在這裡,大家可以多點認識這個時宮,.
願神祝福大家..
多謝Joyce給我們的分享和挑戰.
我們下次見 拜拜.
\newpage



\section{}
\label{sec:lg0Kj2q_1qU}
\textbf{方舟考古問題逐一破障 (2023-10-2 鄺牧道頻,城匯網台聯播)}
\newline
\newline
連結: \href{https://youtube.com/watch?v=lg0Kj2q_1qU}{\texttt{https://youtube.com/watch?v=lg0Kj2q\_1qU}} ~~~~ 語音日期: 2023-10-04
\newline
\newline
\hyperref[sec:VkqqimXzBU4]{\small{< < < PREV SERMON < < <}}
~
\hyperref[sec:index_chronic]{\small{[返順時目]}}
~
\hyperref[sec:index_scriptual]{\small{[返順卷目]}}
~
\hyperref[sec:55qaAhCWIcM]{\small{> > > NEXT SERMON > > >}}
\newline
\newline
$^{1}$Hello Mark Sir 你好啊.
很開心見到你啊 我們的認識.
就是因為這個方舟 差不多12年後了.
那時候YouTube,Social Media都是文字會多一點.
但是現在就可以直接溝通了.
都很開心 Mark Sir.
都真是我想訪問你的.
Mark Sir你在考古上面都有很好的經驗.
其實關於方舟我想你都很熟悉的了.
直接講吧 因為這次.
就因為應金使團再去推動籌款去發現方舟.
它那條數加起來200萬港幣.
如果這件事值得做的 那就應該做.
應該花錢 但是我們12年前.
一幫教務都聯署打了很大的問號.
他們做的無論研究的手法,宣傳的手法等等.
Mark Sir那時候在這本《神師評辯》.
或者是整個你在那個Blog的時候.
都有很多發表 關於考古方面的關注.
所以找Mark Sir來很開心.
就是很希望你澄清一些事情.
我和Mark Sir合作了 我問Mark Sir問題.
Mark Sir可以回應 大家很清楚.
究竟那件事是發生什麼事.
第一我想問 其實去找方舟的遺骸.
或者是那些古蹟 能不能夠找到呢?.
可不可以回答一下.
這個很好的問題 有很多教會的弟兄姊妹.
可能想到這個方舟,羅亞洪水 都會第一個想到的.
就是找到一艘木船的殘骸.
當然有很多坐在教會裡面的弟兄姊妹.
很多時候我們接觸這個問題的時候.
譬如考古 找一些東西去證明.
特別是找方舟去證明.
我們都會比較純真一點.
我們會比較直接一點.
這個其實反映 我覺得這麼多年來反映.
華人教會的那個屬靈的觀念.
我們不是很講理性上的思考.
我們通常都是心領神會就算了.

$^{41}$所以有時候我去到教會跟別人分享的時候.
人們都會說 找到這些東西就好了.
因為會證明聖經是真確的.
這個是一個很簡單的想法.
但是如果你講到一些歷史裡面的問題.
和考古的證據.
那是不是就是這麼簡單的呢?.
我們首先要除去那種反智的.
我覺得屬靈的東西是不用想的.
屬靈的東西是不需要去嚴肅的思考的.
所以你剛才講的那本書.
我們那個名字叫做《神思明辨》.
是真的需要神思明辨的.
聖經裡面有很多的地方都告訴我們.
我們需要去用我們的智慧.
需要我們的知識來做分辨.
去做一些的分析.
這個和聖靈的工作是有關的.
所以我覺得屬靈的成長.
屬靈的生命是需要有這個理性的.
所以講到這方面的時候.
我就想奉勸大家.
是絕對不是這麼簡單的.
首先我覺得第一樣就是.
我經常都跟別人說.
甚至我那個網誌裡面的那個博客.
和那本書裡面的那些文章.
其他的那些作者都是這樣說的.
首先我們需要明白那個方法.
我經常都是這樣跟別人說的.
找方舟不是在山上找的.
是在山下找的.
為什麼這樣說呢?.
因為如果你說有一個這麼大的洪水的話.
這個洪水造成的破壞會是非常大的.
所以正因為這個緣故.
如果真的有一個這樣的洪水.
總共有一年的時間.
在聖經裡面所計算的時間.
羅亞的洪水總共有一年的之長.

$^{81}$所以地下上面的東西.
全部的城市一些的土優都會被破壞.
會被毀滅的.
是同一時間同一年.
所有這些全球的土優.
都會被毀滅的.
這個就是我的意思.
如果你要找羅亞洪水的證據.
考古的證據.
最實在的證據就是在山下找的.
在地上.
在不同的城市不同的古城.
有沒有這樣的洪水層.
洪水破壞的這些這樣的.
我們叫破壞層.
就是destruction layer.
是沒有的.
是沒有一個全球.
同一時間.
所有的土優.
都呈現一個同樣的有淤泥.
有破壞.
有被淹沒的這種跡象.
幾年前我的家裡的土庫.
水浸.
我都有痕跡.
我區裡面的房子.
全部地下都有這樣的痕跡.
那些叫洪水層的破壞.
考古很多時候就是需要做這件事.
因為我們知道每一個土優.
都是一層一層的文明.
一層一層的文化去疊上去的.
如果突然間有一個洪水.
去到消滅了毀滅了它的話.
它一定會留下很明顯的一些痕跡.
事實上在美國一些很機要.
很保守的.
也都想去找方舟去證明洪水.
洪水的一些機構.

$^{121}$他們都到處去找的.
他們甚至在大峽谷那裡.
他們甚至會爭論說大峽谷.
就是一個洪水所造成的.
這個叫flood geology.
這個不被學術界所認同的.
因為這些就是我們要去.
扭曲那些資料.
扭曲這些geological data.
以致fit到我們本身已經預設了的故事.
所以連一些這麼機要的美國組織.
弄出一個方舟.
然後想去找這些證據.
他們都去大峽谷那裡找.
所以你說如果他們說洪水真的這樣發生的話.
正如經文字面的解釋.
所有的土優都會呈現破壞的跡象.
這個在考古歷史裡面是完全沒有的.
另外一個難度就是.
羅亞到底是什麼時候的人.
我們首先要知道什麼時候洪水發生.
然後我們才可以找到證據.
你說洪水是什麼時候發生.
你才要計算時間.
怎樣在經文裡面拿到這些資料.
我們能夠推算羅亞的時代是什麼時候.
這個是不可能的.
因為1-11章的創世記.
他寫的目的 家譜等等.
他的目的的功能不是讓我們現代人去推算一個很準確的歷史.
所以你就呈現一個很大的問號.
到底找到什麼時間的木頭才可以.
那個洪水是什麼時候發生.
你才要決定.
譬如這個等等不同的機構.
開始發現的時候是11-12年前.
他們說大概是3500BC左右.
3000多BC.
我一聽到這個消息的時候我就說這個不行.
因為這個太早了.

$^{161}$太近我們了.
因為3000多BC是什麼.
Early Bronze Age的時候.
只是看中東.
就是有一個很普遍的城鎮化的urbanization出現.
我們可以在土修找回3000多BC至4000BC.
到底發生了什麼事.
是一些城市繼續建造.
然後有很多人聚居.
然後發展了很多不同的古代城市.
所以你不可以去算羅亞是什麼時候.
你就更加找不到任何實質的證據.
你不可以在山上找到一個木頭就可以.
我們在聖經裡面的年表.
我們最能夠計算到的最遠.
最多都是阿巴拉罕.
阿巴拉罕大概2000BC.
之後羅亞等等.
我們是計算不到的.
因為聖經是沒有這樣的功能給我們做.
所以首先第一件事.
找不到一個普遍的洪水層.
也計算不到羅亞是哪一個時間.
你不可以在山上找到.
並且你再去想一件事.
是有關歷史的.
現在我們有的土耳其阿拉伯山的傳統.
是主後九世紀十世紀才呈現的.
所以距離耶穌的時代已經有1000年時間.
越後期的傳統就越不可靠.
所以傳統又去不回羅亞的時期.
我們又不知道羅亞是屬於哪個時期.
你怎麼可以在山上找到古墓呢.
我們先用古墓的方法去計算.
去鑒證年份.
這已經是一個很大的問題.
其實這也是很常見的事.
我在家裡拆你的房子.
你的房子是用古墓來做的.
那是不是你的房子是非常遠古的.

$^{201}$千多年的房子呢.
不是的.
因為你是用古墓來做的.
所以這件事我們不去問的時候.
我們很多時候只是一個很簡單很幼稚的觀念.
我們不知道自己在做什麼.
大家討論一下Mark Sir有什麼想問的問題.
或者為了觀眾可以多一點去了解.
其實找不到一個普遍在全球的土優裡面有的洪水痕跡.
其實這個背景就是有一個叫做全球化的洪水.
有些人的看法就叫做局部的洪水.
我想會不會現在主流的考古界.
或者是色經界都是同意局部的洪水比較合理一點呢.
因為局部的洪水是一個很大的地區的佛.
整個米索不達米亞或者南米索不達米亞.
比如吳爾古城等等.
是找到很多洪水破壞的淤泥的層面.
就是有一個很厚的淤泥是洪水所造的.
這個是常識的,因為在兩河流域.
Tigris River和Euphrates.
經常都會有洪水泛濫,春季的時候.
所以這些洪水的故事為什麼會集中在這個地方出來呢.
比如古近東的泥板,聖經的其中一個記載.
為什麼會在那裡出來呢?.
因為跟他們的環境有關.
因為有些石化大山,然後有一個洪水層.
然後洪水層之上又繼續有連續的不同的文明的建築在上面.
所以你就可以挖下去一個土丘的時候.
你就可以看到整個的歷史.
整幅的歷史.
Regional flood是絕對有的.
這個有一個很好的問題.
就是剛才你問的.
為什麼在考古界裡面.
認真做考古的人.
從來都不會去嘗試去證明.
諾亞洪水或者全球洪水這一回事.
因為兄姐們需要去用理性去想清楚.
考古不是純粹去找一塊東西.
然後聲稱那塊東西就是.

$^{241}$我也可以隨便去找一隻杯子說是耶穌用過的.
吃最後晚餐.
第一個問題就是問你怎麼去證明.
你用什麼標準去證明.
我可以隨便去廚房拿一隻杯子說是耶穌用過的.
我不能夠隨便去找一個地方說.
這裡的聖經裡面發生過的事.
我另一個牽扯的問題就是.
我們怎麼看聖經.
我們看聖經的時候是不是看得小心.
在四章四字中說方舟停在阿拉納山上.
這是重文翻譯.
但原文是.
.
所以方舟不一定是停在山頂之上.
是山區的山腳.
山區的其中一個谷.
可以是半山也可以是山頂.
什麼都可以.
它是說一個群山.
一個山區.
所以不能夠說因為停在阿拉納山上.
而且很後期的時候不可靠的傳統說這是阿拉納山.
所以上面就有方舟.
不可以這樣.
因為看聖經已經看漏了.
我想問一下.
如果總結他們找方舟.
找木結構.
去證明聖經是真的.
他想做這件事.
但他做這個考古的時候.
完全背離了考古的原則.
地點又不清楚.
時間究竟他找什麼年代的又不清楚.
所以其實是很難成立.
我想問Mark Sir會不會回應一下.
現在他根本12年前也好現在也好.
都說要弄一隊考古隊上去.
其實當然剛才說的如果大家聽得清楚.

$^{281}$根本就是不可能去做得到.
又不知時又不知地.
但他們這樣去做.
他們建基什麼.
他們的過程大概是怎樣.
以致我們是不贊成這件事.
當然我不能夠去判斷到人的動機.
但我會覺得做到出來的是什麼.
做出來會fit到預設的解釋.
預設了對聖經的閱讀.
對聖經的interpretation.
很多時候我們想這樣做的時候.
我們不是在捍衛聖經.
我們其實是在捍衛我們對聖經的解釋.
大家看節目要分清楚.
我對聖經的解釋.
你對聖經的解釋.
不等於聖經的權威.
對不對.
我十幾年前對聖經某個經文的解釋.
和我今天對聖經的經文解釋都會有出入.
但我對聖經的interpretation.
不等於聖經本身的權威.
所以他們說去證明聖經.
證明聖經什麼.
證明我對聖經某一個解釋.
聖經可以有其他解釋.
大家都可以是一個合情合理.
甚至更加完備的解釋.
去找一個方舟.
只不過是去證明某一種對聖經的解讀.
不等於是證明聖經.
所以我們都想澄清.
不是因為我們不相信有方舟這件事.
我們是相信聖經的記載是有方舟這件事.
當然Mark Sir其實鼓勵大家聽一下.
城會網台的講座.
都講了很多 我都有聽.
比如說創世記方舟背後.
那個作者寫有神學的意義在後面.

$^{321}$但是有神學意義就不代表那件事是虛構的.
我們相信 當然我們相信.
起碼我和Mark Sir.
或者很多我們對這件事有存疑的人.
反而會相信有一個局部的洪水.
發生過的事也不是沒有洪水這件事.
但是問題是我們在講.
他的方法怎麼去做考古.
可不可以Mark Sir補充一下.
要怎樣才可以在過程中證明一個考古.
發掘到的東西 當他找到一塊木頭.
人家都會說這件事.
或者找到很多塊木頭.
過程是要怎樣才可以有一個認受性.
以致覺得這件事是真的.
因為考古的過程都是這樣的.
時間會有多長呢.
沒錯 有時候要證明一個考古的發現.
有時候需要多年的.
但是有一個方法或者過程是必須要經過的.
就是叫做peer review.
就是同行的平衡.
同學的學術平衡.
我曾經和一些牧者解釋過這件事.
那些牧者居然誤解了我的意思.
所以我和你談過就叫做peer review.
不是這樣的.
peer review是要寫一些學術的日記.
報告你的發現.
報告你的核心資料.
然後給其他這方面的專家.
因為不是每個人都可以在同一個圖標上發掘.
但最重要是圖標發掘了之後.
你需要報告出來.
你不能夠說我只是在新聞會發佈就算了.
我覺得在華人的信徒或者保守的信徒.
去到遇見一些好像很驚人的發現的時候.
我們經常都是譁眾取寵.
可能動機是很好.
我們想捍衛聖經.

$^{361}$但是我們不能夠這樣非理性地.
我們就說這個就是那個發現.
或者出埃及那條路就是那裡了.
或者為了人造衛星拍到.
有很多譁眾取寵sensational的說法.
所以一個步驟一定要做的.
就是要peer review.
就是他一定要寫一個objective的report.
去告訴學術界.
然後由學術界其他人來作criticism.
然後去給一些evaluation和review.
這個一定是需要做的.
這個也是理性的一個思考的過程.
以致我們不需要有盲點.
另外一樣東西就是.
找到很多木頭.
carbon date(加工)它.
屬於古代的木頭.
正如我所說.
你不能夠說只有山上有塊木頭就行.
你下面的洪水層你要match到.
所有的data一定要match的.
你不能夠說上面.
總之3000年BC的時候就有一艘船在上面.
然後下面沒有任何洪水的遺跡.
這個就很奇怪了.
沒可能的.
正如我剛才所說.
開始的時候.
找洪水的證據.
考古證據.
不是在山上找的.
是在山下找的.
山下找不到的話.
在山上找到什麼都是沒用的.
剛才你剛才說的.
你簡單的幫我summarize.
就是何時何地.
何時你都還沒決定.
你怎樣能夠供出證據.

$^{401}$就變成龍門就沒有擺了.
我剛開始找到的就是3500BC.
然後就有另外一個人上去看過.
原來是12000BC.
就算是12000BC.
對一些不認識的人來說.
一定很遠古.
一定好像很可靠.
其實我們人類歷史裡面.
考古的歷史.
一萬幾年二萬年BC.
有什麼不奇怪的.
也有很多的土幽.
有些新石器時代的土幽.
都是沒有洪水層的.
所以這樣就能證明.
我們這樣想太過武斷.
Mark Sir說得很好.
其實你找到一件東西.
你是要配合當地的考古場地的土幽.
去配合.
看看那件東西在哪個位置.
所以我覺得其實這件事.
很多鄧炳姐妹對考古認識很弱.
我不是認識很多.
但我都會看到.
很大問題就是.
他將那件東西拿了出來.
就去發佈.
其實你根本嚴格來說.
你是不可以拿走在考古場地.
你拿走在考古場地.
你做了什麼.
有很多污染.
有很多已經改變了的情況.
你不可以這樣做考古.
但是我們看到.
他做考古很有趣.
隨便拿那些東西出來.
第一,移動考古場地.

$^{441}$第二,他說Claim那個考古場地.
找到了.
但是又沒有人可以上去.
又要透過那些富爾德族人去帶.
又很神秘.
其實是沒有其他人可以找得到.
加上那時候他出第一套.
羅亞方舟經濟啟示.
第二套.
他還在說兩個場地.
完全不同.
我們很客觀的.
這些是一些事實.
不是我們三哥八祖說他是這樣的.
他做這件事.
使得我們.
為什麼會這樣.
還有一件事我想說.
剛才Mark Sir說.
你找考古的研究.
除了說PIA的考試.
還要是一段時間確認了.
但是其實永遠在這些學術裡面.
你是有一個機會.
可能有些東西還沒找到.
以致你決定的時候.
我們說80\%,90\%.
其實很少一點.
是真的這件事.
但是永遠有機會.
之後可能找到一些新的東西.
但不是這樣.
去反證他.
你是有一個這樣的情況.
這些很常見.
這個就是同行的審核是重要.
因為你不能夠說你覺得是就是.
因為這個不是我們平時靈修的結果.
平時靈修你覺得沒所謂.
沒什麼大問題.

$^{481}$但是你講一些客觀的科學的找尋.
你就需要經過一些正確的客觀的.
科學的同行審核的過程.
Mark Sir其實在整件事.
還有沒有補充.
我想你補充一下.
考古做到什麼呢.
你經常都說的.
很簡短的給大英姐妹知道.
其實考古可以做到什麼.
以致我們是不是可以用在全呼音那裡.
這樣去用呢.
我想第一件事.
考古可以幫我們看聖經再認真一點.
正如我們看洪水的故事.
首先第一件事.
我們不可以說那裡講到全地.
所有東西.
All the earth.
我們就說一定是我們今天想的全球.
因為你看看創世記.
創世記在藥室的故事.
講到饑荒遍滿全地.
同一句話.
All the face of the earth.
All the earth, all the earth.
41章5657字.
講到饑荒.
沒有人去證明.
饑荒是有澳洲的人上藥室拿東西吃的.
或者北極有人.
海斯基沒有人去拿東西吃.
不過只不過是regional.
一個大地域的影響.
因為是一個修辭.
沒錯.
一個rhetoric.
聖經裡面經常用一些修辭.
來形容一些.
用一些普世的修辭來形容一些regional的.

$^{521}$一些大的災難.
所以就逼到我們去想.
既然考古沒有這樣的證據的話.
我們是不是看聖經的時候.
看得太快了.
看得不夠仔細.
譬如有人說.
人子的日子就好像羅亞的日子.
這個耶穌說的.
他們說一定是全球的.
因為人子來到的時候就是全球的事.
但如果看清楚馬太福音24章那裡說.
人子的日子好像羅亞的日子.
到底他在比較什麼.
他在比較你看看上文.
就是你不知道哪個時候.
他就會到.
好像下洪水不知道的時候.
人家都不會預期洪水的時候.
忽然就到了.
人子回來也是一樣.
所以他比較的不是範圍.
而是那個突然.
那個驚訝.
他所比較的就是羅亞的日子就是這麼驚訝.
赤外駕取的時候洪水就到.
毀滅就到.
然後人子來.
同樣也是這個日子.
是突然間到的.
上文說的不知道哪個時候就會到.
所以比較的不是範圍.
而是突然的洪水.
這個就是我們讀經識經的時候.
我們很多時候很大膽.
我們不要想清楚.
你再想想考古可以做什麼.
如果沒有錯的話.
有些傳呼音的時候.
已經發掘了的東西.

$^{561}$是可以證明聖經的歷史性的.
是真的.
譬如有個Tel Dan inscription.
達成的石碑.
有Bed David.
House of David.
這個字眼在那裡.
很明顯大衛和他的皇朝就有歷史性.
這些可以絕對在傳呼音裡面告訴你.
然後有個以色列石碑.
一個埃及王打迦南的時候.
他留下的戰地.
他就說到以色列被摧殘了.
絕子絕孫.
這個是什麼意思呢?.
這個吻合時思記的時候.
那班以色列人上到山區.
只不過是一個人民沒有王.
在埃及的石碑.
或者埃及的中央電視台.
都是這樣報導的.
有班以色列人在山區裡.
他們就是一班民眾.
證明了聖經的歷史性.
或者彼拉多的石碑.
克薩尼亞的石碑.
證明了這些人的歷史.
是歷史的人物或者歷史的事件.
這些其實已經確立了.
經過學者評論.
這些是絕對可以用的.
但是說到證明聖經的時候.
證明了聖經的歷史性.
不等於證明了聖經的神學的真理.
證明了大衛是一個歷史人物.
你證明不了大衛的品格.
和神的關係.
他寫的詩篇.
如何套用他和上帝的關係.
這些是無法證明的.

$^{601}$就算你找到一塊石頭.
在曠野找到摩西道指一遊.
證明了摩西是史上最古老的.
你能否證明摩西是一個很偉大的民族英雄.
是無法證明的.
所以我們要對考古這件事.
不要一種迷思.
主要的目的可以做到.
沒錯有證明.
但做到的只不過是有限.
這也是很好的.
因為對不認識聖經的人.
有很多人認為聖經是失憂記.
傳播是虛構的.
這些證據其實可以拿出來.
可以談到的時候.
和別人分享對話的時候.
告訴他們聖經裡面有很多歷史性的文件.
但他們有自己的限度.
主要考古是做什麼呢?.
其實是給我們一個背景去讀經.
給我們一個context.
因為我們知道聖經不是直接寫給我們.
是為我們而寫的.
我們就知道初代的讀者.
古以色人也好.
或者初期教會的基督徒.
第一世紀的人.
聖經直接寫給他們的時候.
他們是如何去領受那個意義.
以至聖經為我們而寫的那部分.
在他們的背景下.
所表述的那些屬靈的真理.
是如何應用在我們今天21世紀.
所以我們讀書的時候.
一定需要明白背景.
問背景的問題.
什麼時候呢?.
聖經描述的時景是什麼時候.
那時候有什麼政治.

$^{641}$或者背景的資料.
我們能夠知道的.
文化的背景等等.
我們就能夠明白當時人怎麼看.
然後我們今天教導,講道,查經.
我們自己看聖經.
我們才會明白今天對我們有什麼意義.
這才是考古最大的貢獻.
很少人真的.
絕對沒有一些嚴謹的考古的隊伍.
是為了去證明一件事.
而去挖掘一些東西.
因為每一個考古的工作.
都會有一個計劃書.
計劃書的議程.
是需要一種學術,理性的進路.
通常一個考古正確的計劃.
就是他想明白該段歷史.
裡面一些的資料.
所以為什麼我們會去剔那個勢.
就不是說我們為了找哥利亞的骨頭出來.
然後我們到處亂找.
或者我們想找最後晚餐.
耶穌用的那一隻杯.
這不是一個正當有的.
合情合理的一種考古的方法.
是.
多謝Mark Sir.
其實都很精簡.
他講了考古可以做到什麼.
希望大家讀者,觀眾都聽得出.
考古能夠做到什麼.
以及不做到什麼.
還有哪些考古是不OK的.
尤其是.
我想有個小總結.
講到關於全副音.
我自己都是很重視全副音.
全副音你說用考古來證明一些東西.
OK.

$^{681}$就算你是找到.
我們都聽得出.
其實是不是完全可以證明呢.
尤其是現在在講.
影音視團說找方舟.
你在未證明到.
我們退了一萬步.
當你要證明他.
但你未證明到已經發佈了.
已經說找到了.
這個就是一個很大問題.
是不是.
不要說要很嚴謹去看他那件事.
行不行.
其實看回無論神仙名變.
或者我們回憶那時候.
其實很多人都勸過他們.
尤其是當中有一篇.
Dr.Snelling.
Andrew Snelling.
他們影音視團去發佈的時候.
在之前找過這位.
想他背書幫忙支持.
但他已經勸他.
你的證據還沒有.
你先不要發佈.
但他們幾天後就發佈了.
變成了Dr.要去寫篇文.
要澄清我們都覺得很有問題.
這個也是那時候.
我們一班傳道人.
除了說注重學術.
但譬如我自己注重傳呼音.
你好像搞這些東西.
有很多人聽到.
好像很驚訝.
找到方舟了.
叫做信了耶穌.
但他的根基真的很淺.
還有對一些所謂.

$^{721}$對這些有認識的.
或者比較重視信仰的理性部分的朋友.
他聽起來更加不相信.
你這班人.
你的根據是這樣的.
那你叫人怎樣信呢.
如果說傳呼音都不合格.
只會使那件事更加差.
所以這是我們很大的關注.
我知道現在他們再要發動這件事.
這條片我們相信都是希望.
橫教會的弟兄姊妹.
要小心行事.
正確認識應該認識的東西.
差不多了.
如果大家有其他好的回應.
想多說一點.
我想Mark Sir都很願意多說一點.
我們打算不要太長.
所以我們時間到了.
下次我們再談.
多謝大家.
下次再談.
OK 拜拜.
\newpage



\section{}
\label{sec:55qaAhCWIcM}
\textbf{God's Providence 2 (Genesis 39) - Rev. Mark Chan}
\newline
\newline
連結: \href{https://youtube.com/watch?v=55qaAhCWIcM}{\texttt{https://youtube.com/watch?v=55qaAhCWIcM}} ~~~~ 語音日期: 2023-10-11
\newline
\newline
\hyperref[sec:lg0Kj2q_1qU]{\small{< < < PREV SERMON < < <}}
~
\hyperref[sec:index_chronic]{\small{[返順時目]}}
~
\hyperref[sec:index_scriptual]{\small{[返順卷目]}}
~
\hyperref[sec:ZcJI3G9ScsI]{\small{> > > NEXT SERMON > > >}}
\newline
\newline
$^{1}$- It's so good to see you all on Marathon Sunday..
Yeah, and you got here okay?.
Traffic was okay, right?.
Yeah, all right..
You know why people wants to run marathon in Chicago?.
Know why?.
Why Chicago is one of the favorite cities.
in United States or even the world for marathon runners..
You know why?.
Well, I was talking to someone who is running today..
He should be still running now.
and as part of our UV family..
So he got the lottery to be able to join.
the Chicago Marathon and this is very interesting, right?.
It's not because you wanna run.
and then you can join the Chicago Marathon..
You gotta be in the lottery..
When you are picked, then you can run, right?.
You know about that?.
You don't know, you don't run, right?.
Yeah, Chicago is actually one of the favorite cities.
for marathon runners and I was talking to this brother.
who is still running today at this moment.
and he told me last year he was running in Atlanta..
Yeah, and he didn't have a good experience.
because in the 26 miles of running,.
there were so many sections of the journey.
was there was nothing there..
He was like, he told me he was like running in empty space..
Right, and no one is cheering him on, right?.
In Chicago is different..
Whether the people on the street are faking it or not,.
right, they're cheering you on, right?.
And there's so many things to see, right?.
You see the high rise and all that..
So Chicago is a very, very interesting place.
and especially for marathon runner..
And it got me thinking, actually this is our life journey.
in light of God's providence.
because God and his people are cheering us on..

$^{41}$Once we know what providence of God really means,.
we are being cheer on and we are encouraged.
every step of the way because we are not running.
in the empty space, right?.
We're not running in the wilderness all the time.
because God is always there on the sign line.
and trying to cheer us on, right?.
This is what providence of God really means..
So this is the second talk of God's providence.
reading Joseph's story..
A pop quiz for you..
What does the word providence mean?.
What does it mean?.
Can you talk to your neighbors?.
What does it mean basically?.
Providence, this word, providence..
What does it mean?.
Talk to your neighbor..
Basically, what does it mean, providence?.
- Provision..
(audience laughing).
- Okay, so here you go..
Providence means pro, vision, right?.
The word, that's how we got the English word vision,.
meaning that God sees everything in your life,.
whatever comes to pass in your life.
from the greatest to the least.
is under the love and care of our God..
So the Bible is very clear about the providence of God.
in your life and in my life..
In God's sight, there was no accident..
In God's sight, there was only incidents..
Got it?.
All these incidents that happened to you.
and you experienced, the people that you have met, right?.
The encounters that you experienced,.
it's all in God's sight, in the presence of God,.
because he's timeless, right?.
God is outside time, it's beyond time and space..
And for him, God looking at you,.

$^{81}$is that there was no accident, but incident,.
an incident under his love and care..
And that's why this doctrine is so important..
And the way the Bible presents and teaches this doctrine.
is through storytelling, right?.
The story in the Bible, how God act in history,.
how God act in the life of his people,.
is the same that God is doing today..
God is also active in your life,.
but teaching you through the means of storytelling..
So the story of Joseph is one of those.
complicated story to teach us this precious truth..
So today is number two,.
talking about the divine providence..
So we're coming to chapter 39..
Okay, let me read this scripture..
This is again, a story..
Let's see what we get from this,.
as we explore together what divine providence.
really means in your life, in my life..
39 verse one, "When Joseph was taken to Egypt.
"by the Ishmaelite traders,.
"he was purchased by Potiphar, an Egyptian officer..
"Potiphar was captain of the God for Pharaoh,.
"the king of Egypt..
"The Lord was with Joseph,.
"so he succeeded in everything he did.
"as he served in the home of his Egyptian master..
"Potiphar noticed this and realized.
"that the Lord was with Joseph,.
"giving him success in everything he did.".
Verse four, "This pleased Potiphar,.
"so he soon made Joseph his personal attendant,.
"his assistant..
"He put him in charge of his entire household.
"and everything he owned..
"From the day Joseph was put in charge.
"of his master's household and property,.
"the Lord began to bless Potiphar's household.
"for Joseph's sake..

$^{121}$"All his household affairs ran smoothly,.
"and his crops and livestock flourished..
"So Potiphar gave Joseph complete administrative.
"responsibility over everything he owned..
"With Joseph there, he didn't worry about a thing.
"except what kind of food to eat.".
That's interesting..
Continue with 6b,.
"So Joseph was a very handsome and well-built young man,.
"and Potiphar's wife soon began to look at him lustfully..
"Come and sleep with me, she demanded..
"But Joseph refused..
"'Look,' he told her, 'my master trusts me.
"'with everything in his entire household..
"'No one here has more authority than I do..
"'He has held back nothing from me except you,.
"'because you are his wife..
"'How could I do such a wicked thing?.
"'It would be a great sin against God.'".
Verse 10..
Let's see..
"She kept putting pressure on Joseph day after day,.
"but he refused to sleep with her..
"He kept out of her way as much as possible.".
Verses 11 and following..
"One day, however, no one else was around.
"when he went in to do his work..
"She came and grabbed him by the cloak,.
"demanding, 'Come on, sleep with me.'.
"Joseph torn himself away,.
"and he left his cloak in her hand.
"as he went from the house..
"When she saw that she was holding his cloak.
"and he had fled, she called out to her servants..
"'Since all the men came running, look,' she said,.
"'my husband has brought this Hebrew slave here.
"'to make fools of all of us..
"'He came into my room to rape me, but I screamed..
"'Then he heard me scream..
"'He went outside and got away,.

$^{161}$"'but he left his cloak behind with me.'".
16..
"She kept the cloak with her until her husband came home..
"Then she told him her story,.
"'That Hebrew slave you have brought into our house,.
"'try to come in and fool around with me,' she said..
"But when I screamed, he went outside.
"and left his cloak with me..
"Paul was furious when he heard his wife's story.
"about how Joseph had treated her.".
Verse 20..
"So he took Joseph and threw him into the presence.
"where the king's prisoners were held,.
"and there he remained..
"But the Lord was with Joseph in the prison.
"and showed him faithful love..
"And the Lord made Joseph a favorite.
"with the present warden..
"Before long, the warden put Joseph in charge.
"of all of the other prisoners.
"and over everything that happened in the prison..
"The warden had no more worries.
"because Joseph took care of everything..
"The Lord was with him and caused everything Joseph did.
"to succeed.".
That's an interesting story, right?.
And a familiar story indeed..
And we know that Joseph started out in chapter 37.
as a favorite son of Jacob of Israel..
And he was unfortunately sold to Egypt..
But down in Egypt, Joseph was doing well.
because God was with him..
And he was in the house of Paul the first,.
one of the captain or one of the officials.
of the house of Pharaoh..
And then here's a story goes,.
and the wife of Paul the first tries to seduce Joseph,.
because again, I think who can resist Joseph, right?.
He was a well-built young man, right?.
He goes to the gym and is very young.

$^{201}$and good looking young guy..
And this wife of Paul the first trying to,.
you know, get him to go to bed with her..
And, but Joseph maintained his integrity.
and he was committed not to sin against God..
And he fled from temptation..
Now this is a self-contained story..
It means when you, you know, you cut the story out,.
you know, chapter 39, pretty much is self-contained, right?.
You can put the story in a separate book..
A children's Bible, right?.
Or a children's story book, right?.
You just tell the story about this young man,.
lonely, away from his parents, his family,.
in a foreign place and fighting against sexual temptations..
And notice the details there..
You find that the wife of Paul the first.
tried to seduce him day after day after day, right?.
Who can resist that?.
But Joseph did..
He was worth a noble character and a man of virtue,.
a man of moral example..
And that's why this story is featured.
in almost every children's Bibles, right?.
Every single children's Bible has this story.
and as a self-contained story..
But you think about it,.
the story is, it doesn't seem to be.
just a simple, you know, moral story..
There were a lot of stories in the Bible..
If we take it artificially,.
or just, you know, read the text,.
it seems like it's a moral story, right?.
Teaching a good example, right?.
Joseph is one of them..
But if you dig deeper,.
it doesn't seem like it's just a story of good man,.
good young man resisting temptation,.
although that is true..
But, you know, remember how the story turns out..

$^{241}$The result of Joseph's integrity was really bad luck, right?.
He was being imprisoned and accused, right?.
Kind of a fake me too moment, right, for Joseph..
And he end up in prison, right?.
And that is not a good turnout, right?.
But some of you are parents with young children..
I cannot imagine the bedtime story would go like this, right?.
You tell a story to your kid, you know,.
your son and your daughter when they were younger..
See, once upon a time, there was a boy,.
very nice looking, he was a good boy,.
and maintain his integrity, and he end up in prison..
Good night, and good night, kids..
You wouldn't do that, right?.
This is how the story functions, right?.
But we want to cut that out,.
and we put it in the self-contained book,.
and tell the story..
So we have to dig deeper, and what does it mean, right?.
How the Bible puts it together.
so that we get what it really means..
Now, a lot of stories in the Bible.
are not just stories of good examples,.
so that we can copy and follow..
A lot of the stories in the Bible are actually God's story..
What is God doing in this man's life?.
And that story tells of his character, God's own character,.
who God really is..
So here, I have two main points..
One, as we dig deeper into the story,.
we find that the divine providence triumphs over even evil..
And second point is providence moves in silence..
First point, perhaps you need to understand this story.
is to understand the overall flow of the chapter..
Now, this chapter has three paragraphs..
Outline the passage, you find there were three paragraphs,.
1 to 6a, 6b to 19, and 20 to 23..
It's funny, interesting, in the first part of the story,.
1 to 6a, Joseph was in poverty's house,.
but there's a phrase that is repeated..

$^{281}$The Lord was with Joseph, Yahweh was with him,.
and giving him success..
And then the second part of the story, 6b to 19,.
you find that this phrase is absent..
There's no obvious presence of God,.
even when Joseph was in trouble, right?.
In fact, the cloak, the clothe that Joseph was wearing.
that Joseph is wearing, remember the story of Joseph.
has to do with what he wears all the time, right?.
Chapter 37, he got the many colored cloth.
from his loving father, right?.
And then chapter 39 right here,.
his clothes is in the hand of Paul's wife..
Now the word hand can mean power and authority..
So in the middle part of the story,.
Joseph, unfortunately, is under P.W.'s hands,.
Paul's wife, okay?.
Paul's wife, P.W..
Seizes the clothe and kept it under her guard.
to accuse him, right?.
She screamed and he went outside.
and leaving the clothe with P.W..
And in this story, Joseph is under the hand.
and the power of P.W..
Where's God?.
God is not there..
At least it's not explicitly mentioned..
And then you go to the third part, 20 to 23..
Joseph is in the prison, the royal prison,.
but then the phrase shows up again..
The Lord was with Joseph, giving him success..
Now, when we talk about divine providence,.
we have to deal with the inevitable tension.
between God's sovereignty, God's control in your life,.
and also human suffering..
Now, a lot of people do not believe in the Bible.
or resist the idea of the existence of God.
because they have seen so much suffering in the world..
There are so many troubles in the world, right?.
You read the recent news,.

$^{321}$what happened in Israel, Gaza..
Hundreds of people died, innocent people..
And to a lot of people, this is enough.
for them to say, "God, if you exist,.
"why all the troubles, right?.
"Why all the suffering?".
Now, this story tells us, or trying to.
sort out the tension between God's sovereignty.
and human suffering..
Now, the Bible never said,.
the Bible never says that God caused PW to tempt Joseph..
The Bible is clear on this..
Although the Bible, there's no anti-philosophical answer.
to the problem of evil..
In fact, there's none, right?.
If you read all the works of all the great philosophers.
who are trying to figure out or sort out.
the sovereignty of God and human suffering,.
there's no airtight philosophical rational answer.
to explain the two..
Never..
Even, you get that, even you have a rational understanding.
or comprehensive answer to why suffering happened to me,.
why suffering happened to other people,.
even you get the answer, the Bible says, "So what?".
Because you are still experiencing suffering, right?.
You still see suffering..
In fact, we don't need one..
The Bible is clear..
The Bible uses stories to try to explain that..
And as we all know, human language has its limit..
And when you hit that limit, storytelling comes..
Because story has the power to explain,.
more so than statements that we make..
So the language is important..
The Bible is clear that God does not cause suffering.
to happen..
Perhaps there's no better term..
We should say God allows or God tolerates evil.
or suffering to come to pass..

$^{361}$At best, we must say God permits or allows.
some bad things happen to us..
Joseph begins with a position of honor,.
being loved by his loving father..
And now he come to the bottom of the valley..
That's a story trying to tell us..
But he is here in the position of shame at the end of chapter.
or in the middle part of the chapter,.
is to prepare him for the next thing to happen..
God's providence is still there when he faces difficulties,.
temptations, and challenges..
Indeed, the power of God triumphs over all of this evil.
things in his life..
So that's how the story goes, showing you who God is..
He can triumph over all the troubles..
And that is true for Joseph..
That is true for you and true for me..
For a little while, Joseph suffered..
He did, he suffered..
But at the end of the episode, God raised him up again..
That is what God is doing in the life of Joseph..
And that should be also the case for all of you,.
for all of us..
So God triumphs over evil and suffering..
But you take deeper into the story,.
the providence of God actually moves in silence..
You have to pay attention, look closely to this chapter..
And if you divide this chapter as I did.
into three different paragraphs,.
you can also subtitle each part as so..
The first part, Joseph in providence's house,.
God is there, clearly, explicitly stated..
6B to 19, in the hand of P.W.,.
God seems to be silent..
God seems to be absent..
20 to 23, again, God is there..
You see, in this chapter of Joseph's life,.
God is there in the beginning, in Egypt..
And sandwiched between it is God seems to be silent.
in times of difficulties..

$^{401}$Is God really absent?.
That's the question the chapter's trying to pose to us.
as readers..
Is Joseph merely under the mercy of P.W.?.
Is he really?.
Or is not?.
As I try to understand this providence,.
this teaching of providence,.
trying to understand what God is doing.
in this episode of Joseph's life,.
I must admit at first, I don't quite understand..
But then I noticed some of the details.
that I didn't see before..
Okay, this is what I noticed..
I noticed a subtle changes,.
subtle progression of where Joseph is..
He started off in poverty's house in Egypt..
Then bad things happened to him..
And then the location changes to the prison..
That is so bad, right?.
However, the word king shows up the first time.
in Joseph's story..
He end up in the prison of the royal house, yes, indeed..
However, it is the prison of the kings, the pharaoh..
Now, why is this location change significant?.
As we all know, the whole story of Joseph,.
we know Joseph is moving closer and closer and closer,.
one step at a time to the royal court of the king of Egypt..
And God is preparing him,.
even in the midst of misfortunes happen along the way..
God allows or tolerates something bad happened to Joseph,.
but he is slowly, gradually moving Joseph closer.
and closer to the king of Egypt, to pharaoh..
At the end, as we know,.
Joseph was being placed there in the courthouse of pharaoh.
in order to save the whole family of Israel.
and indeed plus many other people who suffer from famine..
So God has a bigger plan, right?.
God can use evil and indeed triumphs over evil.
to accomplish his will..

$^{441}$God seems to be absent,.
but that does not mean God is not working behind the scene..
Right?.
As you observe your life,.
as you think about the journey of your life,.
all the twists and turns,.
if you think about it in light of God's presence,.
what a huge difference the people that you have met,.
the accidents that happen in our lives,.
but in God's sight, there's no accident..
It's just incident..
What is God doing in your life?.
Why is God silent or seemingly silent, seemingly absent,.
especially in times of difficulty, times of uncertainty?.
In the past few months, I've been struggling with.
just so much struggle in the presence of God..
I prayed for something and it didn't happen..
And there was a spiritual battle inside of me..
God, what are you doing?.
What the heck are you doing, God?.
Right?.
And there was so many sleepless night.
and I struggled with my prayers, my wish, my desire,.
not happening according to my will..
And then it turned out, praise God, right?.
It did happen..
And of course, according to God's good will..
Philip Yancy, ring any bell?.
Christian thinker, Christian writer..
He has written many, many books about sufferings,.
many, many books about divine providence..
And in one of his stories,.
Yancy talks about his struggle with divine silence,.
that God doesn't seem to be working..
God doesn't seem to be moving..
When he pray so hard, diligently every single day,.
he feels like God is ignoring his prayers,.
denying his request..
And one of those moments, he wanted to give up, right?.
And then he thought of something in the Bible..

$^{481}$And that something in the Bible, he thought of,.
was Jesus in Gethsemane..
Remember the last night of the life of Jesus,.
he was in the garden, he was praying so hard..
He prayed that the Father would take away.
the cup of suffering, death on the cross,.
soon to happen from him..
He was praying so hard in the garden in Gethsemane,.
that prayed that God would remove that cup..
And God denied his prayer..
It is so-called the first prayer in Jesus' life,.
or the only prayer in Jesus' life,.
that God did not answer the way his son prayed..
His son prayed, finally,.
"Not according to my will, but your will.".
So Yancy wrote this paragraph..
He said, "How many times have I prayed for one thing,.
"only to receive another?.
"God alone is qualified to answer prayer..
"And there's some time we have to let go and trust God..
"Who can transmute our self-focused request.
"into an unimaginably larger plan?.
"When Jesus prayed to the one who could save him from death,.
"Jesus did not get that salvation..
"Indeed, Jesus got the salvation of the whole world.".
That is amazing..
God did not answer his son's prayer the way he wanted..
But God has a bigger, bigger, larger plan..
Yet Jesus was not saved from death..
But in Jesus' death, he saved many others in the world..
That is the divine providence of God..
Look at our lives, right?.
There are many requests, there are many prayers.
maybe delayed, or even God said explicitly, "No.".
But we have to remember, who got the bigger plan in mind?.
Not us, right?.
We can only see a few steps ahead of us,.
but God sees everything..
Let me end with the Confession of Faith of Westminster..
Again, this document, a theological document,.

$^{521}$written in 1646, 17th century..
But it has been one of the best treatment.
of all the doctrines in summarizing.
all the important teachings of the Bible..
In chapter three, article three, five, article three,.
it says, "God uses ordinary means.
"to work out his providence day by day..
"But as God pleases, God may work without,.
"beyond, or contrary to those human means.".
God can do that..
God can work beyond those means that you can think of..
God can even work contrary to all those means,.
to all those instruments that you can think of..
In fact, if you think about it, that is the gospel..
That is the good news we found in the Bible..
The providence of God transcends human suffering,.
the suffering of Jesus on the cross..
Human evil, God can triumph over it..
When Jesus died on the cross, it was the moment,.
the cosmic moment of injustice..
But God transcends the means,.
perhaps the most devilish means, the crucifixion..
His son died unjustly on the cross..
But God triumphs over all and accomplishes his will,.
his holy purpose for saving us sinners from sin, from death..
There are moments in our lives, God may be silent..
He may be quiet, but God is not absent..
There are moments in our lives,.
we may experience God as if he's a passive God,.
detached from us..
But the story tells us who God really is..
He's still working, still working behind the scene..
Look at the cross, look at the crucifixion.
of his own one and only son..
He's still there..
He's still there in the life of Joseph..
He's still there in your life..
Let's all pray together..
(gentle piano music).
Father, as we try to sort out the mystery.

$^{561}$of your sovereignty,.
the mystery of all the happenings in our life,.
I pray that if we're going through a season of success,.
we don't forget, we don't forget your hand,.
your invisible hand is still working in our lives..
So we come before you to give thanks..
Thank you, because it's not by our own ordinary means.
that we can achieve all of this..
The fact that we enjoy peace,.
a good season in our lives,.
it's because of your grace..
But we are going through some valleys,.
some dark moments in our lives..
Help us to remember that you are still working.
behind the scene..
You're moving slowly,.
mysteriously, but faithfully in our life..
Help us to trust in you..
Help us to continue to put our faith in you,.
knowing that you are the God who sees all things..
We still wanna praise you,.
we still wanna give thanks to you,.
knowing that we can trust in you..
In Christ, I pray..
[BLANK AUDIO].
\newpage



\section{創世記}
\label{sec:ZcJI3G9ScsI}
\textbf{上主的護理 (3)細想約瑟的故事 3 (創世記 40-41) - 陳崇基牧師}
\newline
\newline
連結: \href{https://youtube.com/watch?v=ZcJI3G9ScsI}{\texttt{https://youtube.com/watch?v=ZcJI3G9ScsI}} ~~~~ 語音日期: 2023-10-11
\newline
\newline
\hyperref[sec:55qaAhCWIcM]{\small{< < < PREV SERMON < < <}}
~
\hyperref[sec:index_chronic]{\small{[返順時目]}}
~
\hyperref[sec:index_scriptual]{\small{[返順卷目]}}
~
\hyperref[sec:s6E5zSGVp5c]{\small{> > > NEXT SERMON > > >}}
\newline
\newline
$^{1}$弟兄姊妹早安.
今天是馬拉松.
昨天我們一班弟兄姊妹有幾個.
包括我老婆和我女兒.
我們去了年輕家庭的退休會.
很熱鬧的.
所以我們昨晚大概12點多才回來.
我們去了一晚.
今天我們又去.
去探探他們.
然後今晚又再回來.
雖然我不是跑馬拉松.
但我都走很多.
不少mile away.
說到馬拉松.
你知不知道為什麼芝加哥馬拉松這麼出名.
每個人都想來這裡跑.
你知道嗎.
我知道你不跑.
你跑不動.
但跑動的很喜歡來芝加哥跑.
為什麼.
我剛剛和一個參加馬拉松的弟兄聊天.
剛好是星期五.
他說為什麼這麼喜歡來這裡跑.
因為他跑過其他城市.
譬如Alanta,Seattle.
他說跑得很悶.
他跑很多段路.
沒有人去打氣.
然後旁邊一路跑的時候.
全部都是野草.
非常荒涼.
芝加哥不是.
他說很多人打氣.
又給他水喝.
有招吃.
你吃招吧.
繼續跑.
然後很多人去拉隊.

$^{41}$去鼓勵他.
他說特別開心.
在這裡跑.
為什麼芝加哥跑馬拉松.
要抽籤這麼厲害.
你知道嗎.
你以為你想跑就跑嗎.
要抽籤的.
其他城市不用.
其實你想想.
上帝在我們心靈裡面的帶領都是一樣.
無論你跑過任何荒野.
每一個路徑.
都是有上帝去打氣.
這就是Providence.
就是上主對我們的庇護.
你想想你走過不同人生的旅程.
不同的章節的時候.
如果有一個人在你附近.
有群人在你附近.
去cheer you up.
去幫你打氣.
那個感覺是多麼不同.
所以聖經裡面說到這個Providence.
是一個很重要.
很寶貴和很切身的教導.
我們說聖經裡面.
是用故事來介紹上帝.
你知道上帝會是怎樣.
你知道上帝在你心靈裡面在做什麼.
你不一定全盤有答案.
不過起碼你要知道這個上帝.
在你身邊一路帶領你.
一路去和你打氣的上帝.
到底是一個怎樣的上帝.
上帝就說你打開我所說的話.
我曾經做過的事.
我當日是這樣做.
這就是當日的我.
今日的我都是一樣.

$^{81}$這個就是用一個story telling的方法.
來告訴我們.
Joseph Story.
上帝是一個怎樣的上帝.
而不單單在約瑟的生命裡面是這樣的.
在你生命裡面也是一樣.
今天我們看到約瑟的故事.
第40章和第41章.
我們看了兩講.
這次是第3講.
第40章我們開始.
這時候埃及王的司酒長.
司善長.
得罪了他們的主埃及王.
法老對司酒長和司善長.
兩個官員法老.
把他們關在護衛長府內的監獄裡面.
就是約瑟被囚的地方.
護衛長把他們交給約瑟.
約瑟就侍候他們.
他們就被關了一段日子.
關在監獄裡面的這兩個人.
就是司酒長和司善長.
在同一個晚上各自做了一個夢.
每個夢都有自己的解釋.
到你早晨.
約瑟來到他們那裡看他們.
他們就甚憂愁.
他就問在監獄裡面所關閉的兩位官員.
他們今天為什麼面帶愁容呢.
他們對他說.
他們做了一個夢.
每個人為我們解夢.
約瑟就對他們說.
解夢不是出於上帝嗎.
請你把夢告訴我.
這段聖經很長.
然後他們就各自對約瑟說.
我做了一個什麼夢呢.
首先是酒長Cupbearer.

$^{121}$他說我夢見什麼.
三支葡萄枝.
然後他用手zar 爛葡萄.
把這些葡萄放在杯子裡面.
把杯子遞給法老.
然後約瑟就跟他解夢.
他說三天之後.
你就會復職.
然後這個善長.
就是煮飯的那個The Butler.
他就覺得這個解得很好.
因為你以前試過拜假神偶像的時候.
去廟裡求什麼上上下下簽的.
有沒有試過.
人家投了上上下下簽.
你為什麼要求上上下下.
這個The Butler就是這樣.
旁邊有一支上上下下簽.
他說我也說個夢給你聽.
他說我頭頂上有三個籃子.
三個籃子裡面都有法老吃的餅.
然後有些鳥來剁食.
約瑟就將他解罵.
他說三天之內你就人頭落地.
原來是這樣.
這一段故事可以暫時看到這裡.
第四十章.
就是講到兩個官員的故事.
然後兩個不同的解釋.
講什麼呢.
這裡告訴我們上帝的護理是怎麼樣.
上帝的護理包括當時的世事.
這兩個官員.
他們的命運.
他們所發生的事.
甚至是法老的埃及王的朝廷.
上帝的護理都在那裡.
世事裡面發生什麼.
雖然不是直接和約瑟有關.
不一定直接和神的子民有關.

$^{161}$不過上帝的護理仍然在那裡.
留意一下約瑟的故事.
或者整本的創世記.
有些事都是打孖.
約瑟的故事.
約瑟發了多少場夢.
兩個夢.
開頭三十出獎.
兩個夢.
然後有什麼打孖.
兩個兄弟打孖.
約瑟和猶大.
整本創世記其實都是這樣的.
很多事都是打孖.
創造的故事也是打孖.
然後亞伯拉罕兩次失誤跌倒.
又是打孖兩次.
所以很多事是兩次兩次.
代表是一個很確定的事.
發生了一次.
然後再發生另外一次.
在這個故事裡面.
其實都是一樣.
你將它列表的話.
酒掌的夢.
三支葡萄枝.
然後有法老的酒.
酒掌的手.
然後他將杯奉上給法老.
右邊就是善掌的夢.
三個攬在頭上.
法老的餅.
善掌的頭.
然後有飛鳥將這些餅吃掉.
在埃及或者古近東的文化裡面.
夢是很重要的.
你不要拿這段聖經說.
今天我發了什麼.
昨晚我發了夢.
你不要亂來.

$^{201}$因為當時的文化.
當時的宗教.
他們的思維.
就是用夢境來掌握世事的.
那些神明很多時候都說自己是報夢的.
所以你不要亂想.
你昨晚發了什麼夢.
你記不記得也好.
不要緊的.
當作發一場夢.
近東的人就不是這樣想的.
所以他們很多時候.
有一些所謂夢境的解說.
甚至有些書.
是幫宮廷的術士去做解夢的.
有一點點像民間信仰.
就是迷信的那些.
看相那些.
其實是什麼呢.
是數據來的.
就是某年某月某日.
或者哪個月的初一十五發生了什麼事.
他們就記錄下來.
然後越來越多這樣的個案.
他們就用這些個案來做一個基礎.
去如何去解釋.
發那些夢.
有一點點像今天的維基百科.
他們就想去維基百科查一查.
如何去為這麼黃的人去解夢.
這兩個官員.
酒長和膳長.
當然被關在監獄裡面.
他們沒有圖書館.
又沒有wikipedia.
那怎麼辦.
所以他們就非常之群怒.
神就偏偏安排了.
就在這個時候.
約瑟就在他們中間.

$^{241}$約瑟說了什麼呢.
四十章第八節.
他就跟他們說.
他說解夢的.
當然啟夢的一定是上帝.
但是解夢的也是上帝.
在這裡整個故事.
約瑟的故事.
三十七章到現在.
約瑟第二次口中說上帝.
第一次就是波提法的妻子去引誘他.
我豈能作這事去得罪我上帝.
第二次.
說出上帝.
在他口中.
說到上帝是控制一切.
啟夢和解夢.
結果剛才那個夢境已經解釋了.
一個就是三天之後復職.
然後三天之後.
另一個就是人頭落地.
結果是怎麼樣呢.
聖經就這樣交代.
到了第三天的時候.
就是法老生日派對的時候.
他擺設筵席.
然後所有的人都來到他面前.
然後正如約瑟所說的.
就是這樣發生.
但是司酒長就忘記了約瑟.
這些叫做飲水不思源.
曾經約瑟去幫過他.
但是他忘記了提貴約瑟.
但是我們看到整個故事.
雖然約瑟的故事從頭到尾.
都好像撲朔迷離和迂迴曲折.
但是你知道聖經裡面的擺放和設計.
就比你的故事裡面.
雖然好像有很多的混亂.
但是很按照步驟進行.

$^{281}$是有結構的.
是有程序的.
不是亂發生的.
雖然故事裡面有很多巧合.
也有很多起起跌跌.
解夢或者他遇見一些難處被引誘.
然後被誣告.
但是好像每件事都是整整有條的.
這個就叫做Providence.
你有沒有寫日記的.
有吧.
有沒有寫日記的.
沒有吧.
你們不寫日記的.
你查一下你的日記.
可能你回頭看的時候.
可能那天發生了這件事.
然後那天又發生了這件事.
這個season又發生了那件事.
好像沒有什麼聯絡.
尤其是約瑟的故事.
你翻查的時候.
好像有很多波折.
但是裡面好像很有條理.
我想問你.
你有沒有試過在街上看到一個車.
然後它越過很多障礙.
左穿右插等等.
然後去到目的地.
你會有什麼結論.
你會結論說.
這個車的性能不錯.
小心開車.
但是它沒有插.
這麼厲害.
可以這樣越過障礙.
然後沒事.
到了目的地.
你的結論是什麼.
你的結論會說.

$^{321}$這些工程師真厲害.
那個計算做得這麼好.
當然你其中一個可以有結論.
但是如果只有這個結論的話.
我會做一個片面的解說.
因為你一定會結論什麼.
有人在後面.
有一個mind在後面.
有一個好的司機.
即使你說那些自駕車.
都是有一個mind在背後.
有AI.
AI也是有人寫的.
都要寫到他可以學習.
所以你不能夠說.
有一個機器人很厲害.
然後背後那個誰.
誰在後面.
那個不能夠打掉的.
甚至那個應該更加重要.
你明白我的意思嗎.
你不可以只是說.
只有一個機制.
是很厲害.
然後你忘記了.
Agency背後那個.
誰在做所有的.
那個更加重要.
約瑟的故事.
Joseph Story.
就是說到這件事.
善長走長的生命.
甚至埃及朝廷裡面發生的事.
上帝控制著.
掌管著所有的東西.
將萬事都在他自己的看管底下.
上帝對萬事的護理.
對世上的國度.
甚至對世上的氣候.
我們繼續看下去的時候就知道.

$^{361}$舊年度的走長.
得到一個上上千不所謂.
有一個好的解釋.
但他卻忘記了約瑟這件事.
都是在上帝自己的保佑底下.
四十章.
就說到上帝對萬事的護理.
四十一章就說到上帝對他子民的護理.
上帝對他自己的People.
就是約瑟所代表的.
他怎樣護理他呢.
你記得.
那個走長忘恩.
喂.
那是怎樣呢.
四十一章就去交代這件事.
為什麼上帝沒有即時叫走長記起約瑟對他的恩惠.
原來故事不是在那裡畫一個句號.
只不過是一個逗號.
到四十一章你就發現了.
四十一章第一節.
過了兩年.
兩年.
法老做夢了.
又做夢了.
法老就做了兩個夢.
就做了一個什麼夢呢.
大概意思就是.
他站在尼羅河邊.
見到有七頭的母牛.
在尼羅河那裡突然上來.
這兩隻母牛很漂亮的肥牛.
然後在蘆葦草那裡吃草.
這個景況挺好的.
對當時的人.
甚至對今天的人都是.
有肥牛.
然後隨後就有另外七隻牛又在尼羅河那裡上來.
但就不是了.
那些是瘦的牛.

$^{401}$不可以吃火鍋.
七隻母牛就站在河邊.
掌上是醜陋的.
肌肉是瘦肝的.
不行.
然後突然法老就醒了.
為什麼會發生這樣的事.
然後那七隻醜陋的.
吃了他的肥牛.
哇 嚴重了.
第五節.
他又睡了.
又做了另一個夢.
又打麻將了.
第二個夢.
所以是很確實會發生的.
夢見什麼呢.
就是有麥穗.
七個麥穗.
又肥又佳美.
這個用來做麵粉最好.
是吧.
漂亮的麵粉.
然後又有七條的麥穗.
是細弱的.
被東風吹乾了.
居然那些細弱的.
又把那些肥的.
飽滿的穗子.
就把它吃了.
法老突然就醒了.
發了一個不太好的夢.
到了清晨的時候.
發生什麼事.
到了清晨.
到了早上.
第八節.
法老心裡就不安了.
但是又沒有人.
可以在那個computer解夢.

$^{441}$怎麼搞呢.
在這個時候.
就是這個酒掌.
就記起了.
兩年之後.
他就稟告法老說.
兩年前發生了什麼事呢.
在監獄裡面.
有一個無私自通的.
藥室.
他居然在上帝那裡有啟示.
然後事情卻是這樣發生.
十四節.
於是法老就派人召藥室來.
連忙的就把他到監獄那裡提出來了.
甚至裝得很漂亮.
就告訴法老.
這個人就在你面前了.
法老就對藥室說.
我做了一個夢.
其實是兩個夢.
沒有人能夠解說.
你聽了夢的.
就能解說的.
十六節.
藥室就回答法老.
這不在乎我.
又說上帝了.
必應允法老平安.
三次了.
說到上帝同在.
上帝實在.
解夢.
兩個夢都是有完成的.
七隻肥牛代表什麼呢.
就是七個豐收的年章回到.
七隻瘦的牛.
就是代表.
七條瘦的墨水.
或者七隻瘦的牛.

$^{481}$就是代表豐收之後.
七年豐收之後.
就會有七年的方年.
不僅如此.
藥室還提出.
我們應該做一些財政預算.
因為七年的豐收來了.
那怎麼預備再來七年的方年呢.
但要儲一儲錢.
做財政大臣.
未來十四年.
怎麼熬得過呢.
因為要證明上帝.
是控制著當時的氣候.
人間裡面.
世界裡面的氣候.
甚至埃及帝國的命運.
都是在上帝的手中.
這個時候就是最重要的時候.
因為正是上帝.
眷顧和扭轉藥室的生命.
也是扭轉神的子民.
在低谷裡面的生命.
我們很快地看.
42節.
於是法老就脫下戒指.
戴在藥室的手上.
然後穿什麼呢.
衣服.
衣服是代表謙卑和尊榮.
現在終於有細麻布.
穿在藥室的身上.
然後給他一個輔助駕車.
然後就將他升職.
他就成為埃及裡面.
一個很重要的人.
大家記得.
藥室開頭的時候.
就有兩個夢.
那時候是什麼時候發生.

$^{521}$藥室十七歲的時候.
當時他被寵愛.
然後他就下了低谷.
這個時候.
就是藥室三十歲的時候.
等了多久.
十三年.
差不多兩個十四年.
所以有人說.
藥室的人生.
在聖經裡面所交代的.
分開三段.
第一段是被寵愛的.
出生至十七歲.
有坐衣著.
然後他就被賣了.
就是他低谷的時候.
十七至二十九歲.
第三段.
就是reversal.
逆轉.
上帝將他逆轉.
就是三十歲開始.
很有趣.
我覺得勝經裡面.
那些刻畫的角色.
人生很多時候都是分三段.
摩西也是分三段.
或者可能在生命裡面.
都是一樣.
曾經被寵愛.
有專為.
但可能所謂的混濟.
有些左濟.
有些低谷.
然後上帝又將你逆轉.
我結束的時候.
我想說這個.
回到這個.
Wesminster Confession of Faith.

$^{561}$我們上一次看過第一節.
第二節.
這次我們看第三節.
第三節說什麼呢.
他說上帝雖然在一般的護理之中.
會使用方法.
means.
ordinary means.
上帝去護理和庇護世事.
人間.
你生命.
是會用不同的means.
但如果上帝看為好的.
上帝也可以不受.
這些方法.
約束行事.
上帝可以不用他們.
上帝可以超越.
這些人間的方法.
也可以抵觸.
這些人間的方法.
如果你反過來想想.
上帝在你生命裡面的庇護.
其實都是一樣.
他不受你所想的方法去行事.
他可以用不同的所謂human means.
但也可以transcend human means.
甚至可以contradict human means.
生命裡面我們遇見很多不同的東西.
但上帝可以用不同的方法.
也可以不用這些方法.
這就是他絕對的主權.
當然他不能夠說.
那我下一天.
明天.
下個季度會是怎樣.
我們不能寫包單.
但可以寫包單的是什麼.
是我們在上帝自己的庇護當中.
你想想其實這就是耶穌基督福音的奧妙.

$^{601}$福音有什麼奧妙.
上帝可以用人的方法.
但卻不受約於人的方法.
他可以超越.
也可以抵觸.
十字架的恩典其實都是一樣.
上帝可以用不同的方法.
讓他的兒子為我們犧牲.
但他也可以超越.
甚至抵觸這些方法.
今天我們在基督裡面可以相信.
這不就是上帝自己的福音.
十字架的救恩.
人以為可以剷除它.
甚至撒旦以為可以剷除.
誰知道上帝就是用這個方法.
超越.
甚至抵觸.
以致我們能夠得到永生.
我們能夠聚得社民.
在說到福音的時候.
不只是單單說到.
上帝好像有一樣東西在自己的計劃裡.
就是這麼簡單.
原來這就是我們能夠相信耶穌.
福音的奧秘.
福音本身就是上帝的庇護.
你能夠信主.
你能夠成為神的兒女.
都是上帝自己的庇護.
結束的時候我想問你.
最近你自己的季度是怎樣.
你的季節是怎樣.
或許你最近的季節.
是很順利的.
一切安好.
那時候我們祈禱會裡面都會這樣說.
有什麼為你祈禱.
我們通常都掛在口邊.
都沒有什麼特別.

$^{641}$什麼意思.
就是沒有破沒有爛.
No news is good news.
如果是這樣的話.
我們值得感恩.
有什麼值得為你祈禱.
感恩.
我能夠享受上帝給我的平安.
第二.
如果你最近的季度.
有很多阻滯.
甚至在低谷.
這個故事就告訴你這個上帝是怎樣.
昔日他是這樣.
今天在你心裡.
都是一樣.
我唱最後一首詩歌.
這首詩歌成為我的禱告.
然後我們才在基督裡面再祈禱.
我和心底留祈禱.
主愛多謝你.
我們生命在你的任下.
生命就是一個馬拉松.
你自己的庇護.
在我們身邊.
每一個路程.
為我們打氣.
我們感恩.
因為你是那位看顧.
護佑我們的神.
因為耶穌基督的福音.
顯明你對萬事的看顧.
對我們生命的看顧.
主給我們這個信念.
給我們一個信心.
我們繼續走我們前面每一步.
無論我們順利的時候.
我們就去感恩讚美你.
我們走在低谷裡面.
在我們憂困的當中.

$^{681}$我們依靠你.
因為我們知道你在我們身邊和我們同行.
感恩讚美你.
祈禱付出給他.
願上帝祝福各位.
下星期見.
\newpage



\section{}
\label{sec:s6E5zSGVp5c}
\textbf{比對研究 19 巴力宗教 12 - 列王記上 19 - 陳崇基牧師}
\newline
\newline
連結: \href{https://youtube.com/watch?v=s6E5zSGVp5c}{\texttt{https://youtube.com/watch?v=s6E5zSGVp5c}} ~~~~ 語音日期: 2023-10-13
\newline
\newline
\hyperref[sec:ZcJI3G9ScsI]{\small{< < < PREV SERMON < < <}}
~
\hyperref[sec:index_chronic]{\small{[返順時目]}}
~
\hyperref[sec:index_scriptual]{\small{[返順卷目]}}
~
\hyperref[sec:iNlk2JOLf1Q]{\small{> > > NEXT SERMON > > >}}
\newline
\newline
$^{1}$再處理第十九章.
由加密山到賀列山.
即是西大山.
十九章第一至第二節.
再看.
就是阿哈將以利亞一切所做的.
以及他用刀殺眾先知一切事.
告訴皇后耶洗別.
耶洗別由派使者到以利亞那處.
明天若在這時候.
我若不使你的性命像.
用那些人的性命.
總之那些被殺的巴力先知一樣.
願神命重重的來懲罰我.
然後這一段我們大家很熟悉的.
因為我們看這段聖經的時候.
我們有很多個人的感想在裡面.
個人的感動在裡面.
就是他害怕的時候.
這位神的僕人就起來來逃命.
就直走到猶大的別士巴.
當時已經南北各分了.
他們當時就在加密山.
鬥完法之後.
你記得十八章最後的時候.
阿哈就回到耶斯列的皇宮.
Jezreel這個地方.
在一個平原上.
這裡也有些考古的法國證明.
當時阿哈是一個居住的地方.
有些遺跡.
以利亞就跟著他.
在這個地方就是十九章開始的地方.
耶斯列的地方.
所以他生命受到威脅的時候.
他就一直向南走.
當時已經分了南北各.
北部的京都就在撒瑪利亞.
另一個皇宮就在耶斯列.
到了南國的時候.

$^{41}$首都就是耶路撒冷.
然後他就繼續向南走.
就到了南國最南的邊界的城市.
就是別士巴.
走哪條路我們不是很清楚.
應該就是照著聖地的地理.
就是中央的山區.
一直跟著山脊下來.
由耶路撒冷到希伯倫.
再到下面的南地.
就是非常乾燥的地方.
沒有水,像沙漠一樣的.
礦業的地方,來到別士巴.
來到別士巴他怎麼辦呢?.
原來有一個僕人是跟著他的.
他就把這個僕人留在別士巴這個地方.
他自己在礦業那邊走了一天的路程.
躺在一個駱騰樹下.
然後就在那裡求死.
他就說:耶護我救了.
求你取我的性命,因為我不比我的先祖更加好.
躺在駱騰樹下之下.
他就睡著了,然後有一位天使就拍拍他.
然後叫他起床.
頭旁就有一些烤了的餅和水.
他吃了,喝了之後又再躺下.
然後另一個使者又回來了.
第二次又拍他.
就說:起來,你要吃東西,要喝水.
於是他又再一次吃了東西,喝了東西.
然後像著飲食的力量就走了四十周夜.
到了上帝的山,就是荷里山Mount Horror.
或者另一個名字叫西大山.
當時以往的地點很多時候都多過一個名字.
西大山荷里山就是其中一個地點.
這個地方就是當時摩西到羅律法的地方.
也是在那裡看到上帝顯現的地方.
不知道為什麼他要回到西大半島的西大山.
這是傳統的西大山,雀波木沙.
很多時候人們到這個地方,就會由埃及出發.

$^{81}$然後坐巴士到這個地方.
在山腳那裡,日出之前就會上山.
上到這個山上就看日出.
很漂亮的,很多遊人都會做這件事.
上去大概要三四個小時的時間.
一路上山,遠足走上去.
大概日出五點左右,有人就會坐下來.
然後看著東邊,看著太陽升起.
是一個很漂亮的地方.
西大山,正如我所說.
這個地方就是上帝和以色列民第一次立約的地方.
他們出了埃及之後,經過了曠野,中間有水等等.
然後就去到西大山.
當然西大山在哪裡,有很多不同的理論.
但最實際的理論都是在這個地方.
因為按照聖經的地理資料和傳統等等.
走三天的路就可以去到西大山.
當然他們走得比較長.
但無論如何,帶著整個的以色列魔神來到西大山.
Bumhussa.
在這裡有法板,拿到一個回望的圖案.
上帝也在眾民面前,在山上加運了以利亞.
然後他出現了.
所以以利亞逃奔到這個地方.
按照一些考古的背景,我們知道.
當時的以色列人或猶大人.
他們都會有一些所謂的朝聖跡象.
朝聖是什麼意思呢?.
就是有些地方曾經是神聖的地方.
例如Bumhussa,Mount Sinai.
所以中間有一些驛站會停下來.
也有一些以色列人留下來的遺跡等等.
所以我們知道這裡大概有一條路.
以色列人或猶大人可能會久不久.
或某一些小撮的猶大人或以色列人會去到朝聖.
所以以利亞走這條路回來這裡.
絕對不是很罕見或很奇怪的事.
無論如何他都會回來這個地方.
為什麼他不去耶路撒冷呢?.
因為在北國首都就是撒馬利亞.

$^{121}$剛才我們所估計的,他應該是經過耶路撒冷.
南國和猶大國就不是把巴黎引入的國家.
起碼南國仍然是真正的信仰.
敬拜耶和華.
聖殿,唯一一個聖殿.
唯一一個正統的聖殿.
就是在耶路撒冷.
為什麼以利亞會不去耶路撒冷呢?.
然後他走到洛尼亞,別氏巴,是不是他太害怕呢?.
還是他去到西乃山有另外一些目的呢?.
還是他回到他的先祖那裡.
回到最終第一次立約的地方.
在喀麗山西乃山.
他去尋根.
當以色列國的國民.
拼棄了耶和華.
不要耶和華的時候,歸向了巴黎.
是一個很大的信仰危機.
上帝不做任何事.
他自己又覺得自己被追殺等等.
是不是因為這個緣故呢?.
所以有很多的psychological reading(心理學讀取).
到底為什麼他會下來這個地方.
或者他想回到上帝和以色列民所立約最根本的地方.
回到西乃山這個地方.
來問個究竟.
所以在中間的路程裡面,四十天.
有兩個使者和他相遇.
這個表我們回看的時候.
我們發現這些故事都有一個結構.
就是耶洗別首先有一個宣告.
然後離開了以色列.
有一個movement(行動).
之前兩個故事的pattern(模式)都是這樣的.
所以他就離開了以色列.
然後有一個天使來拍拍他.
給他水喝,然後給他吃東西.
然後有另外一個天使又拍拍他.
所以有兩個encounter(碰撞).
好像以往的pattern(模式)一樣.

$^{161}$這些分佈一樣.
故事的進展差不多都是這樣.
然後去到十九章第九至十八節.
就是十九章整章的高峰的時間.
然後才會有一些改變.
好,來到這裡,就是九節.
(念念不忘).
這裡很奇怪.
去到西乃山,賀列山的時候.
他就站在山上.
然後有很多神蹟.
裂峰,地震,火.
然後又說不在那裡.
火之後就有微小的聲音.
怎樣看這段記載都好.
作者似乎都是去到描述.
以利亞好像另一個魔世界一樣.
你記得在西乃山第一次.
出埃及的十九章二十章等等的經文.
那裡有一些上帝顯現的現象出現.
就是有黑雲,有風,有地震,有火等等.
然後耶和華就說話.
耶和華就跟神的子民說.
我是耶和華,你們的神.
曾經從埃及地領你們出來.
然後頒布十誡給他們.
所以有這樣的主題又出現了.
上帝顯現的場景.
有火,有風,有地震等等.
這些天象會出現.
然後又有曰的出現.
因為以利亞說什麼呢?.
以色列人背棄你的曰.
說回曰,西乃山第一次立曰.
就是摩西出埃及來立曰.
然後大發熱心.
你記得整個西乃山的故事.
第一次摩西在西乃山的故事.
也是說大發熱心.
就是說耶和華是一個忌邪的神.

$^{201}$他是不信任的,會起疾到的.
因為他是唯一的上帝.
他的子民跟他立了約之後.
救了他出來之後.
如果他不敬拜他.
他是不信任的.
大發熱心.
這裡就說到以利亞大發熱心.
所以這些字眼都被我們想起的時候.
想起第一次.
第一次西乃山發生什麼事.
第一次摩西大連以色列人民去到西乃山.
跟耶和華立約.
然後耶和華在山上顯現的時候發生什麼事.
所以又是回到山上.
在耶和華的面前.
但奇怪的地方就是.
有風,但耶和華不在風中.
風下有地震,但不在地震之中.
地震之後有火,即是雷電.
雷電,但耶和華也不在其中.
火之後就有微小的聲音.
我們先處理微小的聲音這個字.
這個很特別.
微小的聲音,就是三個稀白內字.
Kol Dmamad Dka.
Kol就是聲音的意思.
Voice的意思.
Kol.
Voice.
Dmamad就是寂靜的意思.
即是安靜的意思.
Dka就是比較薄弱的意思.
所以有些英文的翻譯.
我覺得也不錯.
Voice of sheer silence.
或者Voice of whisper.
即是好像突然有大風過了之後.
你聽到的聲音.
風的尾巴很靜.

$^{241}$風停了,那種聲音.
很有趣.
有聲音出現.
這個應該是我們預期的.
正如我剛才所說.
第一次在西大山.
最後最重要出現的是什麼?.
是Voice.
上帝說話.
然後頒布十誡等等.
但我們正如預期.
過了這些崩山碎石的場景.
就應該有Kol出現.
但這裡形容Kol是Dmamad Dka.
Kol Dmamad Dka.
即是Voice of sheer silence.
過了這些安靜的東西.
然後就有這樣的聲音.
怎樣去處理呢?.
我們說一些巴力的神話.
我們可不可以去做一些比對呢?.
巴力是一個風雨神.
所以在之前曾經提過.
巴力的神話.
說到巴力是一個拿著火.
雷電插著地.
然後他右手拿著雷電他的武器.
然後劈到地上.
因為他是風雨神.
所以在巴力的神話中.
亦都這樣形容巴力.
巴力出現的時候會發生什麼事呢?.
當然有火有地震有大風.
最重要的是有什麼呢?.
就是有巴力的聖音出現.
這裡我揀選了一段出來.
就是巴力循環巴力神話中的其中一段.
這裡說到什麼呢?.
就是巴力出現的時候的情況.
他說巴力裂開天雲.

$^{281}$即是在雲層裡面開一個洞.
然後他發出神聖的聲音.
然後巴力的嘴唇就說話.
他神聖的聲音就使大地震動.
這個場景有點像荷里山上的情況.
天雲開然後有聲音出現.
那些是什麼聲音呢?.
巴力神話中說的聲音.
很明顯就是在說雷電.
雷電交加大風的時候的聲音.
另外這裡我亦都留意.
我所點出的.
就是聲音這個字.
就是call.
在烏格麗文烏格利文裡面.
是call.
他神聖的聲音是call.
voice.
他的聲音.
他講說話的聲音.
即是他雷電的聲音.
打雷降雨的聲音.
我們曾經提過.
希伯來文和烏格麗文是非常相似的.
因為大家都是西閃族的語文.
只不過字體寫出來是不同的.
但有很多字其實是同組的字.
這個call字就自然是.
我特別點出這個字.
因為這個正是剛才我們所提到的.
那個是哪個字呢?.
就是聲音這個字.
call這個字.
所以作者似乎很清楚.
巴力出現的時候是什麼情況.
在一個迦南地文化裡面.
很清楚巴力宗教.
他們就知道巴力出現的時候.
就是什麼.
有一個call聲音出現.

$^{321}$打雷.
閃電.
大雨的聲音.
等等.
這個就是巴力出現的情況.
回到以利亞這個故事.
似乎想做一件事.
不是我們所預期的.
做些什麼呢?.
你見到有烈風大作.
但耶和華說他不在風中.
風之後有地震.
耶和華說不在其中.
地震之後有火.
耶和華說不在其中.
這是經文裡面一次一次強調的.
就是不在.
但我們所預期的其實是什麼.
有烈風大作崩山碎石之後.
耶和華說在風中.
為什麼呢?.
因為巴力應該是這樣的.
證明耶和華是強過巴力的話.
應該是說.
所刮的風會更大.
所給的地震應該是更大.
比巴力更大.
地震之後有火.
應該是那種火那種雷電.
是比巴力所劈下來的雷電應該更大.
但偏偏這裡的經文怎樣形容呢?.
不在其中.
火之後有微小的聲音.
這裡似乎作者.
在一個巴力的環境下.
他就去到強調一件事.
就是耶和華出現的時候.
不像巴力出現.
為什麼呢?.
因為巴力就會在風,火,地震的時候出現.

$^{361}$這段經文.
這段文獻.
說到他有神聖的聲音出現.
那種聲音使地大震動.
天雲裂開.
有火降臨.
這就是巴力.
但在以利亞的故事裡.
似乎告訴我們.
耶和華不像巴力.
甚至應該這樣說.
耶和華不是巴力.
不像人去到預期.
更大風,更多火,更大的地震.
全部耶和華都不在其中.
這裡有一點點顛覆了.
我們對一個風雨神.
降雨去到颳風的真正的神.
有一種不同的預期.
不在其中.
就代表以利亞以為.
我好像巴力一樣.
會更加大的情況出現.
更加大風,更加大火.
更加多這樣的神蹟.
誰知不是.
他說他不在其中.
最後才有微小的聲音.
這裡似乎告訴我們.
不單止耶和華不像巴力.
耶和華本質都不是巴力.
他的本性和巴力本質上完全不同.
真正出現的時候是何時出現.
似乎是微小的聲音出現.
如果我們看看這段記載的結構.
19章9至18節.
場景是在山頂山洞.
然後說的華臨到以利亞.
問他一件事.
以利亞你在做什麼.

$^{401}$回答.
我為耶和華萬軍之神大發熱心.
因為以色列背棄了你.
現在只剩下我一個.
然後耶和華說.
然後有四樣東西.
風吹過不在風中.
地震經過不在地震之中.
火不在火中.
最後有微小的聲音.
如果這樣列出的時候.
似乎結構上和故事的下半部一樣.
下半部是說什麼.
又是回到場景.
又是山洞.
又說有林多他.
我們覺得很沉悶.
為什麼會這樣.
因為經文的編排是這樣.
問他問題.
又問你在這裡做什麼.
同一句說話.
回答以利亞又是同一段說話.
他背棄你的約.
我為他大發熱心.
為上帝大發熱心.
只剩下我一個.
然後又是耶和華說.
然後又有四樣東西.
他說你回去.
你去告哈舍.
做阿萬王.
然後又告耶和.
做以色列王.
然後再告以利沙.
去接續你作先前.
這三樣東西.
一個是外邦的王.
哈舍.
另一個就是阿哈的其中一個將軍.

$^{441}$Jahu.
然後到以利亞.
其實這三者.
都是在列王記下面的內容.
去交代他們被上帝用.
去殲滅.
這個拜巴力的阿哈的家族.
哈舍.
就是亞蘭王.
所以之後我們看20章開始.
有亞蘭王和以色列王.
亞蘭王和阿哈.
去征戰.
然後擊敗他們等等.
然後到耶和.
耶和將耶洗別.
阿哈抄家.
然後就是以利沙.
以利沙去對抗巴力的先知等等.
然後第四樣.
上帝就說有七千人未向巴力屈膝.
如果回到剛才那裡.
其實這四樣東西.
就是一個平衡.
耶和說有四樣東西.
但最後有他的聲音出現.
他在聲音裡面.
然後再說.
耶和說命令他.
一二三.
我的計劃會繼續.
我的工作會繼續.
然後我也在以色列保守住七千人.
未向巴力屈膝.
就代表上帝的說話.
一直以來他的聲音.
call.
一直沒有失敗.
上帝的說話.
上帝的言語.

$^{481}$先知所說的.
上帝所認他的.
沒有離開過.
沒有失敗過.
是會繼續.
然後下面就看到.
七千人未向巴力屈膝.
所以一個故事.
完了之後.
以利亞就.
他就從原路回去.
所以一個重新.
再去到.
重新宣告.
上帝的話語.
是沒有被廢棄.
因為上帝會保守住.
保守自己的說話 保守自己的子民.
\newpage



\section{15:1-7}
\label{sec:iNlk2JOLf1Q}
\textbf{The Power of Storytelling (Luke 15\_1-7, Matthew 18\_1-6,10-13) - Rebecca Wu}
\newline
\newline
連結: \href{https://youtube.com/watch?v=iNlk2JOLf1Q}{\texttt{https://youtube.com/watch?v=iNlk2JOLf1Q}} ~~~~ 語音日期: 2023-10-17
\newline
\newline
\hyperref[sec:s6E5zSGVp5c]{\small{< < < PREV SERMON < < <}}
~
\hyperref[sec:index_chronic]{\small{[返順時目]}}
~
\hyperref[sec:index_scriptual]{\small{[返順卷目]}}
~
\hyperref[sec:RamdTJjsZh8]{\small{> > > NEXT SERMON > > >}}
\newline
\newline
$^{1}$and I'm so glad to be invited back again..
Last time I talked about art and beauty.
and took far too long..
And this time, I know I'm supposed to stay.
within 25 minutes..
So I'm gonna do that..
That's why I have notes and now I have a clock.
and I have all these things to help me stay on time..
All right..
So it's always a pleasure to be here and a delight..
Today, I'm gonna be talking about the power of storytelling..
But before we begin, I wanna make sure.
that we talk just a minute.
about why stories are so compelling..
So what do we get out of stories?.
What's the purpose?.
What purpose do they serve?.
So, one is that it's for enjoyment..
So it entertains us..
It lets us escape from our day-to-day stresses..
Like this book, "Harry Potter,".
which maybe some of you have read or read the sequels,.
watched all the films, went to the theme park,.
bought an owl and are like crazy fans..
But pure fun and entertainment..
You didn't buy an owl?.
Oh yeah, me too, me too..
And then like a little wand and everything.
and a lightning tattoo..
(audience member laughing).
So other times, stories are there to educate us.
so we can learn something new that can improve our lives..
One book that's been on the New York Times bestseller list.
is this one by Oprah Winfrey..
So anything by Oprah Winfrey.
usually gets on a bestseller list..
This book is about trauma and how it lives in our body.
and how we can heal from it..
Another reason we love stories.
is that stories connect us with other people..

$^{41}$So it helps us engage,.
has like an emotional attachment to it..
Perhaps in reading someone else's story,.
we can see a mirror to ourselves.
and know that we're not alone..
This book is called "Crying in H Mart,".
and it's about a daughter.
who takes care of her really sick mom.
and is really overwhelmed..
And one day while grocery shopping,.
can't hold it in anymore..
So for those of us who have been there,.
we know that when we're so stressed,.
even the most mundane task of grocery shopping.
can feel very daunting or triggering..
So here we have all these compelling reasons.
to love stories, to read stories, to listen to stories..
But why, Rebecca, are you here talking about stories?.
So you've been asked to come,.
and as a guest speaker,.
I feel like it's really important.
that I talk about things that I know..
And Pastor Mark talks about things that he knows..
And I know a lot about stories.
because I'm an author and I self-published my book..
And through that process,.
I learned a lot about editing,.
what it means to tell a story well..
And I wear a lot of different hats professionally..
So as an educator, as a workshop facilitator,.
if I think about it, story is really the common denominator..
And everything we do,.
when there's a story attached to it,.
we learn so much better,.
and the lesson is far more meaningful..
And I've witnessed when people,.
listeners are glued to a storyteller.
because they're just so good..
And that's just like human-made stories, right?.
Like stories that we write as humans..

$^{81}$So imagine how much better it could be.
when it's a God-made story..
And that's what we're gonna talk about today..
So today, the goal is really to see.
the power of storytelling through Jesus, who's the goat..
So if you don't know what goat means,.
it means greatest of all time..
And if you don't think Jesus is the goat storyteller,.
we're gonna go through it,.
and hopefully by the end,.
you're gonna see why he is the greatest of all time..
So today's story is the parable of the lost sheep..
Maybe it's familiar to you, maybe it's not..
Regardless, I hope that you'll see.
that through this simple story,.
that Jesus is really compelling all of us.
to listen well and to learn..
So a parable in its definition is a short story.
that teaches us a spiritual truth..
So parable rhymes with fable..
It's similar to a fable..
Maybe you've read it or heard the story.
of the tortoise and the hare or the lion and the mouse..
These teach a moral..
In a parable, it teaches through a metaphor, usually,.
and it's to apply to our spiritual lives..
So if you can turn to your Bibles if you have them,.
or on your phone, or read along with me,.
we're in Luke chapter 15, verses one through seven..
This is the NIV translation..
Now the tax collectors and sinners.
were all gathering around to hear Jesus,.
but the Pharisees and the teachers of the law muttered,.
which means that they're grumbling,.
this man welcomes sinners and eats with them..
Then Jesus told this parable,.
suppose one of you had 100 sheep and loses one of them..
Doesn't he leave the 99 in the open country.
and go after the lost sheep until he finds it?.
And when he finds it, he joyfully puts it on his shoulders.

$^{121}$and goes home..
Then he calls his friends and neighbors together.
and says, rejoice with me, I found my lost sheep..
I tell you, in the same way,.
there will be more rejoicing in heaven over one sinner.
who repents over 99 righteous persons.
who don't need to repent..
So in some translations,.
the word parable means picture story,.
sort of like a comic book or a graphic novel in modern time..
So I'm gonna lead you through a picture story.
of this very parable..
So here we have, right, here's a pasture, grassy field,.
here's the setting, and then the owner of the sheep goes,.
hey, go eat some grass..
Sheep go graze..
They go graze, they walk around, they roam around,.
they come upon maybe some trees..
The owner doesn't count, sort of like on a field trip..
Oh no, one is missing..
Where's that one?.
Maybe that one is scared, but I also have to leave the 99.
in order to find that one or I can't..
(laughs).
Maybe the sheep is calling for the owner..
The owner finally finds the sheep,.
puts it on his shoulders, is so happy,.
and says, hey everyone, I found my lost sheep..
Everyone, let's go party..
So that's my retelling of the story in pictures..
Now if we think about stories,.
I want us to look at this quote..
It says, the best stories are the stories.
that any and everyone can understand and relate to..
So it's a quote from a pretty smart person,.
and that person was my editor for my book that I wrote..
It's called Moving in Forever,.
and it's a picture book about hospice care..
So it's a heavy topic, and this is what she would tell me.
all the time when we were working on my book,.

$^{161}$because it's a book that's very personal to me.
about my college roommate who died..
So I wanted to write every single detail down.
from my perspective, 'cause I felt that was necessary..
And of course, my perspective.
and all those details are really valid,.
but if I'm to publish this book to make it meaningful.
for other people, other families, other children,.
I have to edit it..
I have to make it so that any and everyone.
can understand it and relate to it,.
'cause otherwise, the impact will not be there..
So this is what we're gonna do today..
We're gonna do a Bible study on this parable,.
but through the lens and eyes of an editor/publisher/writer..
So we're gonna decide as editors,.
if this story is up to snuff, is it gonna be a hit?.
Is it gonna be a New York Times bestseller or not?.
And how you do this is you have to look at it.
with some level of scrutiny..
So when my editor looked at my manuscript,.
she made sure that I was okay.
with whatever she's gonna say to me..
We have to look at it with some level of inspection,.
just like your writing teachers at school.
might be looking at your essays.
and your papers with inspection..
What do they look for?.
What do editors and publishers look for?.
First, it's like the structure..
Does it have good bones to a well-written story,.
beginning, middle, and end, and everything in between?.
And does it have the components of a well-written story?.
And it's those two parts that give you.
a story that's good enough, a good story..
But we're not here in the Bible to look at good stories..
We're here to look at great stories..
We're here to be compelled and convicted..
So these stories need powerful storytelling,.
good and strong storytelling,.

$^{201}$so that it's powerful..
Otherwise, it's just good enough..
And Jesus was never just good enough..
So let's look..
Let's see if this parable of the lost sheep.
survives our inspection..
So does it have a structure?.
Does it have the bones?.
Okay, yes, as a beginning,.
we are oriented to the setting where we are..
The sheep are grazing..
There is an owner or a shepherd..
And in every great story, there's a problem..
The problem is, oh no, one goes missing..
We don't leave the problem hanging..
We have to do something about it..
There has to be an action..
What are the characters gonna do?.
In this case, they're gonna leave the 99.
and go find that one..
You keep finding, you keep finding..
And as listeners and readers, we want a resolution..
Is there one?.
Yes, finally found the lost sheep..
Yay!.
And what is the ending about?.
It's what do the characters do.
after the problem's been fixed?.
So now everyone celebrates with the party..
Okay, the bones are there, check..
Are the components there in the parable of the lost sheep?.
Is it a good story?.
Well, we have the characters, got that..
And then we have the setting..
The main plot really is the lost is found..
Classic, straightforward, and the message is.
the lost matters even if it's only one..
And the voice that Jesus uses.
within this parable of the sheep is pretty firm..
And diction is just a fancy word to say.

$^{241}$the choice of words that the writer.
or the author chooses to use, or the storyteller..
So in this case, Jesus uses suppose one of you..
Well, doesn't he?.
And joyfully finds a lost sheep..
So he chooses these specific words.
to tell this specific story to this specific audience..
So we got the structure, we got the components,.
but do we have good storytelling?.
Or else it's just maybe B plus, right?.
And our Jesus is not a B plus student..
He's like A plus, plus extra credit,.
plus like principles, that's like.
everything valedictorian, right?.
So we have to have, and those of us who are Asian here.
and felt that academic pressure, you would understand..
Good storytelling, as my editor has told me,.
and what she has learned from the entire industry.
and profession of copywriters, are these aspects,.
these areas..
So you have to have a trusted teller..
So whether it's the writer, the author,.
or the person speaking, do people care.
about what this person's saying?.
Is this person worthwhile, right?.
If Oprah Winfrey were to tell you a story,.
I'm sure it'd be packed..
But I don't know, maybe not..
If it were someone else, I don't know..
Then is there enough drama?.
'Cause if there's not enough drama,.
we're not gonna be captivated, we won't be invested..
We actually like the up and down of all the emotions..
Is it relatable?.
Can I identify with it?.
Does it matter to me personally?.
And does it matter to anybody else?.
Is it simple?.
So the most powerful stories.
are the simplest stories, actually..

$^{281}$And every word and every sentence progresses the narrative..
So we don't waste time, we don't waste space..
Agency..
Agency is, does the reader or the listener.
have the power and freedom to make meaning on their own?.
So when I was trained to be a teacher.
for elementary education, in class, in grad school,.
I was told, "He who does the work does the learning.".
So I could teach children all they need to know..
I can tell them exactly what they need to know..
And they only need to memorize it and regurgitate it..
But that type of meaning, that type of learning.
is not meaningful to them..
It's much better when they connect the dots,.
when they can make meaning out of it,.
when they can have a personal connection.
to whatever we're learning..
Okay, so now we're gonna see,.
does the parable fit good storytelling aspects.
as we see it in the copywriting world?.
All right, so, trusted teller..
Well, yeah, it's Jesus..
People are lining up to listen to this man, right?.
God in human form..
Yes, we trust this teller..
We care what he says..
Check..
Drama, lost and found, lots of emotion..
Fear, joy, worry, celebration..
That's a lot..
Is it relatable?.
Now back then, herding sheep is pretty common knowledge,.
but just in case it's not, to these Pharisees,.
actually in the rest of the chapter in Luke,.
Jesus tells other lost and found stories,.
like the parable of the lost coin..
Have you ever lost money?.
The parable of the prodigal son..
Have you ever lost a child?.
So there is a lot of relatability here already..

$^{321}$And is it simple?.
It's, I don't, if you count the lines,.
it's maybe four to five lines of the story,.
you can't get simpler than that..
And every single piece moves the narrative..
And is there agency?.
So here is kind of tricky, because in the Bible,.
we don't know if, what these Pharisees.
or these religious leaders said.
after listening to the story..
The Bible doesn't say, necessarily..
They didn't say if they stood up and clapped,.
or they stood up and left..
It doesn't say..
But that's the beauty of agency, isn't it, for us,.
because we have the ability to make meaning.
out of it for ourselves..
So, looks like, check, check, check, we have it all..
Now I'm gonna share with you another quote.
from a very wise person..
"Sometimes we need to zoom out a bit.
"to see the bigger picture..
"Perhaps this is a story within a larger story.".
You might be thinking, wow, what a wise thing to say..
And you're wondering, who said it?.
Well, I said it..
I say it to my kids all the time,.
'cause as you know, I'm a teacher, I'm an artist,.
I'm an author, so I'm pretty much like a teacher nerd mom.
to my children, and I wanted to show you.
how we at home read the Bible and read stories.
and look at it with inspection..
So on the left here, you'll see a picture of our bookshelf..
And it has just a collection of Bible stories..
And these are like one book of one story,.
all designed for kids..
And on the right, you'll see a stack.
of what I call the Bible Bibles..
So it's not that the other ones aren't real,.
but these are the Bible Bibles we use.

$^{361}$for reference and quiet time..
And the very top one is mine, and I tell the kids.
that that's what adults read..
There's no pictures in it, there's lots of footnotes.
and lots of words and really tiny text..
And then I say that, but then actually I read.
all of the other ones way more frequently.
than I read my own..
And what we do, let's say for example,.
on the left side, I pull out a Noah's Ark story,.
and I have a few copies of them from different authors,.
different publishers, to see how the story changes.
based on who's telling it..
The crux of the story is the same,.
but maybe one is only pictures..
Another one, maybe they take out a little bit,.
but we don't know that until we go to the right.
and pull out a Bible Bible and look at.
what is the bigger story, what is the bigger picture?.
'Cause these are not standalone stories,.
is what I wanted to make sure my kids understood..
Noah's Ark is not a standalone story..
Easter is not a standalone story..
They all fit into a larger context,.
a larger story inside the Bible itself..
And it's the same for this parable of the lost sheep..
So it's part of a larger story,.
but it really depends on how far you zoom out..
So we have the parable of the lost sheep..
If we zoom out maybe 25\%,.
there's actually another story that's being told.
of Jesus teaching..
It's Jesus's teaching story..
And if we zoom out even more,.
all of Jesus's teaching story are in the New Testament..
We zoom out even more, well, it's really.
the huge story of the Bible from creation.
all the way to redemption and us reuniting with God..
But how far do we zoom out?.
For today, we're just gonna zoom out like 25\%..

$^{401}$We're just gonna zoom out to Jesus teaching.
just a little bit, right?.
So if we zoom out, we'll see.
that the lost sheep parable here.
doesn't essentially change, but how he tells it changes..
So these are the components of the parable.
of the lost sheep, right?.
We already talked about this..
And if you look outside,.
like before he starts telling the parable in Luke,.
there's actually more characters to this story..
There are religious leaders, there are sinners,.
there's lots of people here..
And as the setting and other components add on,.
you'll see this metaphor built up..
It's not just about lost and found..
It's about addressing the leader's biases..
It's about saying not just one sheep matters..
It's saying, hey, we don't need.
to have a person look a certain way, be a certain way.
to be able to enter the kingdom of heaven..
All we need is belief and repentance..
Now, even though within the parable of the lost sheep,.
Jesus' voice is firm, I wanna say that outside of that,.
that storyline above it,.
he's actually pretty clever and witty..
So the Pharisees and the religious leaders.
are basically grumbling..
They're like, oh my goodness, look at that Jesus..
He's like talking to sinners and like eating with them..
Basically, ew, gross, right?.
Like how horrible, dirty..
Jesus doesn't really like respond to that directly,.
but you can see by the words he chooses..
He just tells them a story..
He's like, well, there's just gonna be more rejoicing.
for that one sinner who repents.
than all 99 righteous people who don't need to repent..
Link, link, like who are these 99 righteous people?.
And I found that really genius.

$^{441}$because he didn't really have to say anything..
And it's like, mic drop..
And it reminds me of my grandma.
'cause she used to do that to me..
So she used to tell me these stories.
in order to get me to do or not do something..
A typical story is, it's all in Chinese,.
but in English it goes,.
once there was a child who swallowed a watermelon.
and swallowed all the watermelon seeds..
And two days later, a watermelon grew out of his head,.
period..
And this is what she would tell me often.
whenever I'm eating watermelon or anything that has seeds..
But she would change this story ever so slightly..
So she'll change it like this..
There once was a little child.
who swallowed so many watermelon seeds.
that a watermelon grew inside of her..
So there's no more room in her stomach for other food..
Oh, and by the way, I made noodles..
This is a response, an indirect response to my question of,.
can I have more watermelon?.
So I just find Jesus here, not exactly like my grandma,.
but just so clever in how he answers.
where I'm sitting there as a child thinking,.
oh no, I should not have the watermelon anymore.
'cause I want the noodles..
So I wanna say that Jesus takes these stories.
and kinda changes them,.
but maybe not exactly like my grandmother.
who's just getting me to not eat more watermelon..
And these stories are true.
'cause I've swallowed some seeds.
and when the watermelon doesn't grow,.
I realize the story's fake..
But I want you to see how Jesus tells the same parable,.
but in a different way,.
and if you can spot the difference..
So we're gonna go to Matthew chapter 18,.

$^{481}$verses one through six in IV..
I am going to read them and you can read along with me..
So at the time, the disciples came to Jesus and asked,.
who is the greatest in the kingdom of heaven?.
He called a little child to him.
and placed the child among them..
And he said, truly I tell you,.
unless you change and become like little children,.
you will never enter the kingdom of heaven..
Therefore, whoever takes a lowly position of this child.
is the greatest in the kingdom of heaven..
And whoever welcomes one such child in my name welcomes me..
If anyone causes one of these little ones,.
those who believe in me to stumble,.
it will be better for them to have a large millstone.
hung around their neck.
and be drowned in the depths of the sea..
I don't know if Jesus actually said it like that,.
but that's how I hear it..
Skipped a few verses, Matthew 18, 10 to 13..
See that you do not despise one of these little ones,.
for I tell you, their angels in heaven.
always see the face of my Father in heaven..
What do you think?.
If a man owns a hundred sheep and one of them wanders away,.
will he not leave the 99 on the hills.
and go to look for that one that wandered off?.
And if he finds it, truly I tell you,.
he's happier about the one sheep.
than about the 99 that did not wander off..
In the same way, your Father in heaven.
is not willing that any of these little ones should perish..
So I'm getting the message here..
And I've never even tended sheep..
I don't even enjoy touching them..
But I'll tell you as a parent,.
losing a child is one of the worst possible nightmares..
And when my family were in seminary,.
we lived in graduate housing,.
I thought I lost my eldest son, Tyson..

$^{521}$He was only not yet even three years old.
and we're out in the back, up on a hill playing..
I turned around, picked up my younger son, Ethan,.
who was barely one at the time..
And in a split second, Tyson was gone..
I could not see him anywhere..
And this is in Deerfield..
It's not like there's like crowds and like buildings..
This is like Deerfield, suburban Chicago..
But he's nowhere to be seen.
and it felt like time stood still forever..
And I could hear my heart beating so fast.
and it was like coming out of my body..
So I had to decide, do I leave Ethan to go find Tyson?.
Which is hard to do 'cause he's barely one..
So I left him with other children.
and told the children to watch him like a hawk.
while I look for Tyson..
And I'm screaming and screaming..
And I'm so loud that all the other seminary students.
drop their books, come out of their dormitories,.
into the field going, what's going on?.
I lost Tyson..
And then everyone's looking for Tyson..
We're at the soccer field, the football field,.
the playground, inside the slides, in the tunnels,.
in the forest, everywhere, in the crosswalk,.
everywhere, no sign of him..
I am starting to panic..
The worst like comes to mind..
Oh my goodness, he was abducted..
I don't know by what, a coyote?.
I'm not sure..
But in my head, it was like worst possible thing..
I go around the apartment building and I see him there..
He's just sitting on the step waiting for me..
I run to him and I grab him.
and I'm rejoicing inside to tears..
I pick him up and I just run around everywhere..
I found him, I found him, I found him..

$^{561}$And everyone is celebrating with me..
So if you ask me, was I glad and happy.
that Ethan didn't wander off?.
Of course..
But was I happier that I found Tyson?.
Most definitely, right?.
So even though I've never heard a sheep,.
a tended sheep or even one touch them,.
this story I get, I can relate to it..
I identify..
Because you see, the parable of the lost sheep,.
as told in Luke, did not do that to me..
I didn't think of that..
Because there was not that whole section.
about Jesus talking about how important children are..
Right?.
It didn't hit me in the same way..
So if we look at this,.
the components within the components of Mark,.
you'll see that he actually uses a different voice.
because his disciples are basically bickering..
Who's better?.
And I experience that all the time with my children..
Who's better?.
I'm better, I'm better..
No, I'm better..
No, you're not..
Yes, I am..
The point is, there's a bigger plot..
He's settling the rivalry between them with this parable..
And he's saying, all you need to do is be like children..
But I wanna also highlight,.
he's pretty casual when he tells the parable,.
but he's pretty disciplinary.
when he tells them about children..
And here's like the tone,.
unless you change, you're not gonna get this..
And it's about my father and your father..
In Luke, Jesus didn't do that..
Right?.

$^{601}$He didn't say your father, my father..
He used heavens..
There was no personal relationship there..
And also in Luke, he says,.
suppose one of you owned a hundred sheep..
In this one, he said, suppose a man..
Well, why?.
Why are there these differences?.
Because it's a different audience..
The disciples are not wealthy people..
They would never own a hundred sheep, right?.
But why does that matter?.
It matters a lot because within that context,.
it tells a storyteller how to modify the story.
to make it impactful for whoever's listening..
Right?.
And then later he was also saying,.
you just have to be, you just have to be like children..
So in that context,.
I don't think that Jesus is particularly condemning,.
but really teaching..
So if we take a look at those two zoomed out storylines,.
side by side, let's compare them..
And I wanna focus a lot on the audience..
So we've talked about a bunch of these already before,.
but I wanna focus on the audience..
In Luke, there's these Pharisees.
and these religious leaders, right?.
And as much as he's being clever and witty,.
kind of tongue in cheek,.
he's basically trying to show them.
that the gates of heaven are welcome.
to people they don't expect..
And the leaders wanna gatekeep it,.
but Jesus is saying, no, as long as there's repentance,.
there's acceptance..
And for some of us here,.
maybe we're grumbling about someone or something..
I know I sure am a lot of times.
when I think those people can't have access.

$^{641}$to the kingdom of heaven..
Those people, the conniving people,.
the self-righteous people, the rude people,.
the exclusive people..
And then I'm thinking, oh, wait a minute..
Am I being exclusive here?.
So there's a compelling moment.
and you're convicted, right?.
Now in Mark, even though Jesus is directly talking.
to the disciples and teaching them, disciplining them,.
he told them to be like children,.
even though they were acting like kids, ironically..
But I don't think he was condemning his disciples either,.
but really helping them to solve their problem.
and their debate..
And it's subtle, but actually not so subtle..
If they're just like children,.
their father's not gonna let any of them perish..
They're just chill out..
And for some of us,.
maybe we think someone else is not deserving.
or we're competing with those around us..
Because after all, that's what society tells us to do..
But I wanna mark that actually,.
that in Mark, Jesus is also talking to children.
'cause he invites a child in their midst..
So the children might not be the active, direct listeners,.
but they're listening..
As you know, if you ever have kids,.
you can say one thing and they hear it.
and they'll repeat it..
But you'll say 10,000 other things.
that you want them to say and they won't..
Unless that's just my kids..
And I wonder when these children are listening.
to what Jesus is saying to the disciples,.
if they feel comforted..
Because as children, they have no power..
They have no choice..
But what Jesus is saying is,.

$^{681}$they're so important that all of you need to be like them..
And maybe some of us,.
we don't relate to judging Pharisees..
We don't relate to maybe disciples who are quarreling..
Do we relate to children?.
Do we feel helpless?.
Do we feel confused?.
Do we feel lost?.
Do we feel inexperienced?.
But yet, our God tells us.
that children are the most beautiful thing..
And what we see here is that bottom row, agency..
We get to have that agency to look at these stories.
in a zoomed out way..
So if we're to be like Jesus here on earth,.
healing the sick, finding the lost,.
how does this story look to us today?.
So who is the audience today reading this story.
and the story outside the story?.
The audience, meaning you and I here right now..
So ask yourselves, what context are you currently in?.
Are you in a place of a lot of turmoil, decision, chaos?.
Are you in a state of jealousy or unforgiveness?.
How does that affect how you hear the story?.
We're not just one of these sets of audiences forever..
We can be all of them at different times in our lives..
So maybe there is a lost sheep in your life..
Maybe there's a need to reconcile that..
Or maybe you think you deserve more than that person..
Or could that lost sheep be actually you?.
Are you feeling lonely and scared or helpless?.
But remind yourself of what the parable says,.
regardless of audience and context,.
that each one matters..
Every single one, you, me, every single one of us matter..
So in summary, the goal was to help you know.
the power of storytelling, which I hope I made that point,.
and that you realize that Jesus is the goat..
If not anything else, Jesus is the goat..
But I also want you to know that this is one story..

$^{721}$This is the lost sheep, one story,.
that is inside like a zoomed out story, right?.
That is inside an even bigger story..
So I am really excited and looking forward.
to continuing diving into these stories with you.
as we move along..
Next week, I'll be here..
We'll talk about another story..
And so far, that's the end for me today..
But if you have any stories you wanna share with me,.
you wanna hear more of my stories,.
feel free to find me and contact me..
But that's all I have for you today..
[BLANK AUDIO].
\newpage



\section{}
\label{sec:RamdTJjsZh8}
\textbf{以巴衝突的反思 (1) - 陳崇基牧師}
\newline
\newline
連結: \href{https://youtube.com/watch?v=RamdTJjsZh8}{\texttt{https://youtube.com/watch?v=RamdTJjsZh8}} ~~~~ 語音日期: 2023-10-17
\newline
\newline
\hyperref[sec:iNlk2JOLf1Q]{\small{< < < PREV SERMON < < <}}
~
\hyperref[sec:index_chronic]{\small{[返順時目]}}
~
\hyperref[sec:index_scriptual]{\small{[返順卷目]}}
~
\hyperref[sec:Xbkzaf9FOTY]{\small{> > > NEXT SERMON > > >}}
\newline
\newline
$^{1}$好 大家早安.
我們暫時今天去做一點點的detour.
我們不說providence.
我們說說以巴的衝突.
大家在新聞上和傳媒上應該接觸了很多這樣的資訊.
我相信我們作為基督徒.
我們都需要去了解一下.
到底這個地方是發生了什麼事.
在我自己這麼多年對以巴的接觸.
不單止有時候會帶團.
甚至裡面去做一些工作等等.
我想在這裡和大家分享一下.
因為我認為有很多時候.
每逢以巴有衝突的時候.
在傳媒或者在基督教的圈子裡面.
就會有一些不同的講論.
我認為有兩套的講論.
第一套的講論是一些比較helpful.
比較有幫助的進路.
另一套就是很多時候.
特別是傳統教會的圈子裡面.
我認為是很harmful.
沒有什麼建設性.
甚至是有傷害的一種進路.
來討論以巴的問題.
為什麼需要講一件事情.
除了在傳媒裡面.
我們已經不能夠避免地.
我們要接觸很多這樣的新聞.
這些資訊.
甚至在教會的圈子裡面.
我們接觸到這些不同的講論.
我相信我們作為基督徒.
在世界上我們都需要去了解一下.
用一下我們的理性來思想一下.
到底這件事是一回什麼事.
或者有什麼叫做我所說的helpful的comments.
或者一些harmful的comments.
每逢講到以巴的時候.
一定會觸碰很多不同的意見.

$^{41}$我相信我今天講這篇.
在網上也有不同的迴響.
這些是預料到的.
以巴裡面有很多不同的爭執.
不同的甚至矛盾的說法.
但是都需要我們去做一些反思.
正如我所說.
有一些的講論是幫助比較大的.
有些的講論我認為幫助是比較少一點的.
甚至是有害的.
每逢以巴爆發的時候.
很多人就會很快用聖經.
或者用一些末世的語言.
來解釋當前的一些局勢.
有很多時候這些坊間裡面所流行的說法.
我們需要問一個問題.
到底是不是嚴謹的.
是不是有一些地方不是很仔細去想清楚.
我覺得作為基督徒.
我們手拿著聖經.
我們是不是認識聖經.
是不是真的明白聖經.
然後是不是將我們一些一知半解的看法.
然後套上去一個這麼複雜的情勢.
所以我們會處理一些經文.
然後在這些經文裡面.
我們去想一想.
到底會不會和現在地上發生的事情有關.
還是沒有關.
這個就是以色列.
加沙的地圖.
在左下角就是加沙strip.
或者是加沙走廊.
加沙corridor.
左下角長長的.
在海邊的那個地方.
就是今天我們在新聞裡面看到.
以色列和巴勒斯坦地.
或者以色列和加沙之間的一些衝突.
這個地方加沙strip.

$^{81}$是大概360多個平方公里.
有多大.
其實是香港的三分之一.
大家有些是香港來的.
香港的三分之一.
裡面有230萬巴勒斯坦人住在那裡.
另外右上角那個地方.
好像蝴蝶的一半.
那個地方叫西岸.
叫west bank.
為什麼叫西岸.
因為是約旦河的西岸.
約旦河Jordan river的west bank.
那裡是中央山地.
中間你看到是耶魯沙坎這個地方.
在左下角加沙strip這個地方.
過去十幾年是受一個政府去管治的.
這個政府的政權叫做哈瑪斯.
哈瑪斯.
他是一個武裝的革命.
武力革命的一個政權.
一個恐怖的組織.
他誓要去反對以色列的軍事佔領.
所以很多時候在一些報章或者一些新聞裡面.
我們經常會看到.
為什麼他們會射一些火箭過去以色列的南邊.
譬如Bezhara.
Bezhara這個地方.
我記得2014年的時候.
我在以色列.
我剛剛就去Bezhara那裡做了一些工作.
正好這個時間就是2014年.
哈瑪斯又是射火箭過去以色列的南部.
所以我們整個星期就在Bezhara.
我們天天聽到這些火箭.
他們detect到這些火箭的警報.
然後我們就躲在防空洞.
所有的活動都是暫時去.
各自暫時取消的.
就是一個這樣的情勢.

$^{121}$現在都是一樣.
在這個地方主要傳媒在說的.
你看CNN CBS等等.
在說這個左下角的一個地方.
我們先不說太多所謂政治的東西.
我們先處理一下一些坊間裡面比較流行的.
一些比較popular的theories.
我們看看聖經裡面到底是怎麼說的.
我大概總結了.
有五個坊間裡面我們經常接觸到.
說到以巴的言論或者一些說法.
第一個就是猶太人就是以實的後裔.
阿拉伯人就是以實瑪利的後裔.
特別我想斟酌的是後面那一句.
第一句後面那一句.
就是阿拉伯人Arabs.
就是Ishmael的後裔.
這是第一個的命題.
經常在坊間裡面我們會接觸到.
第二個就是阿拉伯人就是伊斯蘭教徒.
因為伊斯蘭教徒的緣故.
所以他們有比較多的印象.
他們是一些激進的分子或者恐怖分子.
這是第二個命題.
第三個命題就是我們坊間裡面經常接觸到的.
就是迦林地就是上帝無條件地賜給以色列人或者猶太人.
第四個就是神的計劃當中.
有兩個的peoples兩個的子民.
一個是他所揀選的猶太人.
另外一個就是外邦人包括我們.
這個我們經常坊間裡面聽到.
第五個就是以巴的衝突就是末世的徵兆.
這個更加多聽見了.
所以我們今天未必有時間去處理五個.
希望能夠處理到頭三個.
第一個命題我們需要去想一想.
就是阿拉伯人是以實瑪利的後裔.
這個這樣的印象.
因為在坊間裡面我們講到以巴的報道的時候.
我們經常聽見就是千年的恩怨.

$^{161}$這是幾千年解決不了的恩怨.
你知道以實瑪利是誰.
阿拉伯人有起碼兩個兒子.
第一個就是以撒是撒拉所生的.
那個是繼承產業繼承聖約的.
那個就是Isaac.
第二個兒子就是夏甲所生的.
就是叫做Ishmael.
以實瑪利.
以實瑪利不是被揀選去承受聖約或者產業的.
所以他最後被拆走離開了迦南地.
我們在不同坊間裡面或者後期的傳統.
我們就會說這些阿拉伯人.
就是Ishmael的Descendants.
是他們的以實瑪利的後裔.
正因為這樣的緣故.
所以我們就說現在的以巴衝突.
以色列人就是猶太人.
然後巴勒斯坦人就是阿拉伯人.
所以猶太人和阿拉伯人是世仇.
因為阿拉伯人是甚麼人.
是甚麼時代的人.
是主前二千年的人.
你說是不是世世代代這麼厲害這麼多的恩怨.
距離現在多久.
有四千多年的時間.
我們老公和老婆吵架.
吵了五分鐘之後你都不記得吵甚麼.
何況吵了四千年.
所以我們有一個坊間的一種印象.
是幾千年的恩怨.
所以是難分難解的.
到底是不是這樣.
以實瑪利.
以實瑪利的一些經文.
原來上帝仍然是祝福以實瑪利.
在創世記21章說.
上帝聽見孩子的聲音.
這個就是以實瑪利.
當他小的時候.

$^{201}$下格帶著他走路.
然後上帝的使者就從天上呼叫下格說.
下格你在這裡不要怕.
上帝已經聽見孩子在那裡的聲音.
他們當時在曠野.
起來把你孩子扶起來.
用你的手握著他.
因為我必使他成為大國.
所以上帝仍然是有祝福去祝福.
以實瑪利的後裔.
在創世記後面的一些章段.
就去解釋.
有上帝祝福的應許的實言.
創世記25章.
然後他的後代就住在沙特阿拉伯半島的北邊.
Northern Arabian Desert.
然後就是一些幼木民族.
你就會想.
一定是所有阿拉伯人.
就是從以實瑪利所生出來的.
但是你要明白.
創世記裡面有記載.
有其他的民族.
是住在Northern Arabian這個地方.
譬如哪一個呢.
阿寶艾罕有另外一個設施.
叫做基土拉.
在不同的家譜裡面.
記載基土拉的後代.
也是住在Northern Arabian.
另外呢.
以素.
雅各有兩個兒子.
媽仔你記得嗎.
以素和雅各.
以殺有兩個的媽仔.
以素和雅各.
以素也是住在以東.
甚至是和Northern Arabian這個地區.
是同一個地帶.

$^{241}$裡面也有提過不同的人.
其中一個人種.
叫做阿瑪利人.
你說說來說去很悶.
很亂.
總之意思就是.
在所謂Northern Arabian.
或者所謂阿拉伯人的先祖.
那個地方.
其實已經有不同的人種了.
不同的人種.
不只是只有以實瑪利的後裔.
所以你這樣看的時候.
阿拉伯人是不是就是以實瑪利的後裔呢.
No.
就算你說.
有部分阿拉伯人.
是以實瑪利的後裔.
但你都不能夠以偏概全地.
你就說所有的阿拉伯人.
就是以實瑪利的後裔.
所以你不要用一種這樣的坊間.
或者很popular的一種看法.
來看現在以巴的衝突.
為什麼不可以這樣看呢.
如果你這樣看的時候.
我覺得這種角度.
不單止對現實是沒有什麼幫助.
甚至是有害的.
因為你這樣看的時候.
你就當作一種民族.
就是看它為負面.
不可以說抹黑.
但起碼有一種妖魔化.
他們是阿拉伯人.
他們是以實瑪利的後裔.
所以他們是沒有光環.
而另外一邊就是有光環.
因為他們是被揀選的民族.
你看看創世記.

$^{281}$創世記神學講得很清楚.
上帝選擇阿巴拉罕和他的家族.
是做什麼的.
是要使萬國能夠得福.
所以以實瑪利也是受上帝的祝福.
其他民族也受上帝的祝福.
整本的創世記由頭看到尾.
所有在以色列附近的萬國萬族.
包括埃及.
都是受到上帝自己的不成.
你這樣看的時候.
你就不會說將某一種民族.
我們將他categorize.
然後覺得他是差一點的.
另外一種可能就是有光環.
你說不是.
伊斯蘭的傳統.
說他們自己是會教徒.
說自己是以實瑪利的後代.
可蘭經或者穆罕默德的其他一些遺訓.
他們這樣claim.
他們這樣聲稱.
不代表是有負歷史的.
可蘭經說.
阿巴拉罕獻以實的時候.
其實不是獻以實.
其實是獻以實瑪利.
那你說你信誰.
可蘭經說我就不是那麼拜.
他們想將他們的identity.
想將他們的名份.
和阿巴拉罕其中一個很出名的兒子.
以實瑪利來認同.
是有他們當時的原因.
後期的原因.
但是歷史裡面是不是真的這樣.
所以你搞清楚.
不要將某一種民族.
我們按照坊間的一些看法.
我們就覺得他們是差一點.

$^{321}$另一種民族就是好一點.
另外一個要處理有關以實瑪利的經文.
就是十六章十二節.
就是講到以實瑪利之後會變成怎樣.
他必將野奴.
他的手要攻打人.
人的手也要攻打他.
他必尚與他眾弟兄作對.
你這樣看.
你就發現他好像經常與人打架.
很野蠻.
這個翻譯是一個比較負面的翻譯.
原文可不可以只是這樣翻譯.
我又想將他.
其實你可以這樣重新翻譯.
這個重新翻譯.
是正面的翻譯.
或者中性的翻譯.
這是一個猶太學者的翻譯.
怎樣的翻譯呢.
他說應該是.
他必將野奴.
即是說他是一個自由的遊牧的民族.
他的手要與人.
with everybody.
with him.
with his kinsmen.
意思是甚麼呢.
即是說他是一個自由的民族.
和他周邊的民族進行貿易.
進行合作.
然後他比其他遊牧民族更加大的更有能力.
如果這樣轉過來翻譯.
你就不會說那些Ishmael.
就是一些很野蠻的人.
周圍與人作對.
只不過是說他與其他人合作.
然後他的權勢可能比較更加大.
這個也是聖經裡面我們處理得比較小心.
正如我所說的創世記的神學.

$^{361}$亞伯拉罕的家族.
上帝所呼召和揀選出來的.
使萬族能夠得福.
以色列人能夠成為祝福.
以色列人周邊的人包括.
以實瑪利人.
都能夠得福.
所以要搞清楚.
不要用一個坊間的千年恩怨.
以巴的衝突是甚麼時候的恩怨.
如果要數的話是近代的恩怨.
大概一百多年左右的恩怨.
所以你寧願用一個比較有幫助的方法.
去想一想到底近代發生了甚麼事.
特別是1947年聯合國.
他們通過了.
說要在巴勒斯坦地位猶太人.
去建立一個猶太國家.
就是今天的以色列.
開始有很多不同的戰爭和爭論.
如果你回到聖經的時候.
如果亂用聖經.
我就認為這個不是很負責任.
和不是很有幫助.
第二個問題我們可以去處理的.
阿拉伯人就是伊斯蘭教徒.
我們以印象中.
多謝我們西方的傳媒.
印象中阿拉伯人是恐怖份子.
有沒有這樣的印象.
如果你沒有的話.
你就很清心了.
但我們很多時候會有這種錯誤的印象.
我們就覺得阿拉伯人.
一定是很恐怖的人.
我曾經聽過一個人說.
他在學阿拉伯語.
或者在學希伯語.
開始去讀那些希伯語.
或者那些阿拉伯語的時候.

$^{401}$他就說笑.
他聽出來的聲音.
好像那些恐怖份子說話.
那你廣東話都很沒禮貌.
廣東話真的出名沒禮貌.
很刺激.
所以用廣東話吵架是最好的.
因為殺傷力最大.
你聽到台灣人說國語.
很溫柔.
是吧.
起碼我的印象是這樣.
你不能夠這樣說.
這樣說就是告訴我們.
我們被西方的傳媒.
或者偏袒一方面的傳媒.
我們很多時候會有一個.
可能錯誤或不是很對的印象.
雖然有很多阿拉伯人.
是穆斯林.
但你知道在阿拉伯的人口裡面.
也有很多數以萬計甚至百萬計的基督徒.
是Arab Christians.
是阿拉伯族的基督徒.
所以Arab這個字.
阿拉伯這個字.
其實是一種文化.
意思就是說指到所有的人.
他們的第一語言.
是Arabic.
是阿拉伯語.
阿拉伯人之間.
他們有一種共通的文化.
共通的一般的一種聯絡.
所以他們大家都是成為Arabs.
不代表你可以混為一談.
Arabs有一點點像Asian.
亞裔.
亞洲人.
你想想亞洲人是怎樣的.

$^{441}$Asian.
我們不是一個單一的種族.
是有很多不同的群體不同的種族.
混為一談.
每個人都當是同一的.
早前你記得楊紫瓊嗎.
楊紫瓊.
楊紫瓊不是我像嗎.
中國有些官方傳媒居然說.
他們為中國人爭光.
關什麼事.
人家馬來西亞人來的.
Asian.
但你不能夠混為一談.
對嗎.
Arabs其實都是一樣.
是一個單一的種族.
有些可能是藍眼睛.
有些皮膚可以黑一點.
有些皮膚淺一點.
有些是介乎兩者之間.
大部分的阿拉伯人沒錯是Muslim.
但有成千上萬的阿拉伯人.
是基督徒甚至是福音派的基督徒.
所以我們不要想得太過簡單.
或者只是用一些我們平時流行的.
街邊我們聽到的一些訊息.
我們來看以巴之間的衝突.
很多時候我去以色列的時候.
我都帶弟兄姊妹去西岸.
因為我們加薩路不可以去.
加薩路是封鎖的.
但我們進去West Bank.
我記得有幾年.
我和弟兄姊妹安排了.
他們住一些巴斯坦人的家.
在那裡借宿.
這些真的是Airbnb.
因為你和他一起住.
你和他一起住的時候.

$^{481}$你才發現.
當我們和其他人有種誤解的時候.
為什麼呢?.
因為我們只有一些很抽象的印象.
我們沒有Faces.
我們沒有Names.
我們不知道那個是誰.
我們很多時候就是用一種.
坊間或者傳媒給我們一種印象.
來了解對方.
但當你住在他家的時候.
你發現原來有巴斯坦的基督徒.
他們有家.
有兒有女.
有孫.
幾代同堂.
你住在他們裡面.
住在他們家裡.
他們熱情地招待你.
主人家寧願自己少吃一點.
他們都願意去服侍你.
你想想你和他們聊天.
他們其實都是和我們一樣.
是普通人.
都是想落地生根.
想平平安安.
你都想平安.
你都想建立你自己的家庭.
能夠落地生根.
只不過在以巴的雙方.
雙方都有一些激進的分子.
或者激進的派系.
政黨等等.
兩方面都有.
腐敗.
有自私.
有貪婪的政權.
你和這些人聊天.
無論是猶太人或者是巴斯坦人.
其實生活都是很無奈的.

$^{521}$夾心人.
夾在中間.
You cut in the crossfire.
當兩邊有槍戰的時候.
你剛剛就在中間.
這個情況就是一個這麼苦的情況.
一個政權.
哪一邊的政權都好.
都不會代表到.
所有政權底下的人.
哪一個是我們的政權.
哪一個是我們的政黨.
政黨做的事.
不一定代表你做的事.
對嗎.
所以不能夠去到混為一談.
甚至我們應該去審視一些西方的傳媒的時候.
我們都要用一個批判一點的眼光.
批判一點的眼光.
第三方面.
就是迦南地.
就是迦南地是上帝無條件的賜給猶太人以色列人.
因為時間關係我們只是暫時處理.
下一次我們再繼續.
聖經是從哪裡來的.
在創世記十二章.
耶和華呼召亞伯蘭的時候.
要離開本地本族父家.
去到我所指示你的地方.
我必使你成為大國.
我必賜福給你.
使你的名為大.
你要使別人得福.
我祝福的我要祝福他.
就在你的我必就在他.
地上的萬族都要因你得福.
我們經常聽見這段經文的時候.
我們就說.
上帝一定是寫了包單.
一個無條件的地契.

$^{561}$給了亞伯蘭和他的後裔.
是嗎.
大家可能有聽過這樣的故事.
很多時候我去到以色列巴士坦地的時候.
都會有些人和我這樣說的.
特別是一些亞歷克森.
他會和我說什麼.
他說有些很好心的人.
在美國去.
基督徒.
他們就去到那個地方.
他就和一些巴士坦的基督徒說.
他說你們走吧.
搬去別的地方.
這個地方是屬於猶太人.
如果有人這樣和你說.
你有沒有家.
去敲門.
他說麻煩你搬.
因為有個民族比你早到.
這個地方是屬於他們.
你搬不搬.
你怎麼搬.
你說搬就搬.
是嗎.
用這樣的說法.
I am sure it is well intended.
I am sure.
I am sure it is well intended.
但是.
是不是有幫助呢.
前設就是.
舊約的聖經.
是以色列猶太人的地契.
I don't think so.
結束的時候讓我一分鐘.
來解釋一下.
然後我們下一次再繼續.
先講一講聖經的文法.
講一講聖經的文法.

$^{601}$中文譯做我必.
上帝說必定會或將會.
文法是imperfect.
或者future tense.
好像寫了包單.
但是你裡面.
你發現有幾個動詞.
不是future tense.
是一種文法裡面.
sorry 要講這麼悶的話.
文法裡面叫coordative.
什麼叫coordative.
是你可以譯做.
是讓我.
給我或者等我的意思.
我舉個例.
我和我老婆說.
我必定會.
幫你掃地的今晚.
那我想寫包單.
因為將來式.
但如果我用coordative的話.
勸告式的話.
我不知道勸告式是不是最好翻譯.
因為可能性.
他表達的可能性.
潛在性不一定會發生.
我就會和我老婆說.
什麼coordative.
等我幫你掃地今晚.
那就不是必定掃.
因為可能性.
等你.
讓我.
意思就是說.
我要和你一起partner.
一起做.
如果你這樣看的時候.
這個就更加make sense.
在阿巴拉罕整個故事.

$^{641}$上帝說讓我使你成為大國.
等我能夠成為.
等我使你的命為大.
等我使你成為祝福.
使萬國能夠得到祝福.
要等阿巴拉罕和上帝.
一齊來到.
阿巴拉罕要信任上帝.
信靠上帝.
不用自己的方法.
去達到這個可能性.
potentiality.
所以不是一種無條件.
unconditional.
神有一方.
人也有一方.
整個阿巴拉罕故事其實就是在說這個道理.
阿巴拉罕很多時候.
他會用自己的方法.
自己的human means.
但是.
很多時候用了自己的方法.
就會成為一些障礙.
如果他相信上帝的話.
那件事.
就會好辦很多.
coherent.
是不是一個永遠的契約.
永遠的地契.
下回分解.
在聖經裡面有很多不同的.
應許.
最重要的.
我們要看到的.
在基督裡面得到實現.
土地的應許.
國度的應許.
祭司聖殿的應許.
全部都在基督裡面.
是應驗的.

$^{681}$我們就是那一班.
在基督裡面.
有這個應許的人.
我們就是truly.
是被祝福的那一班.
上帝祝福我們在基督裡面.
給我們有的應許.
是多的時.
你的生命.
有多少上帝的blessing.
土地就能夠成為其他人的祝福.
最後的時候.
用一段很簡單的聖經.
他說那麼多的後書一章二十節.
這裡說什麼呢?我讀出來.
他說神的應許.
不論有多少.
在基督裡面都是是的.
yes.
所以我們藉著他說.
阿們amen.
使上帝能夠得到榮耀.
我們暫時在這裡結束.
謝謝.
上帝我們在一個那麼混亂的時代.
求你給我們更加清晰的一些了解.
主愛的代表.
再次毀到以巴的衝突.
我們禱告.
主愛知道以色列的一方.
慘受.
這些恐怖的襲擊.
人被挾持.
主愛我們不知道他們的生命是怎樣.
主愛我們求你憐憫.
起到為了這個.
彼此之間的衝突.
為了平民百姓.
我們禱告.
他們就是夾在.

$^{721}$這個crossfire(沙灘)裡面.
上主我們求你.
憐憫這個土地.
讓裡面無辜的人.
生命.
受到更大的毀滅.
主愛你是滿有憐憫.
滿有慈愛.
滿有和平的神.
求你憐憫這個地方.
憐憫這個地方上的以色列人.
猶太人.
巴勒斯坦人.
主愛你是那位歷史的主宰.
我們求你繼續憐憫他們.
主愛也幫助我們.
讓我們在基督裡面能夠得到所有認祂的人.
我們的生命是一個祝福.
讓我們在這個生命裡面.
有不和的地方.
我們替做和平.
有需要了解的地方.
主愛我們願意帶著愛.
進入我們的家庭.
進入我們的社區.
使人能夠見到耶穌基督的自己.
主給我們教會.
能夠在地上實行這個使命.
靠著你的恩典.
我們去見證你.
祈禱逢轉美及.
願上帝賜給各位平安.
我愛平安.
\newpage



\section{13:24-30}
\label{sec:Xbkzaf9FOTY}
\textbf{The Power of Analogy: Weed, Wheat, or What? (Matthew 13\_24-30) - Rebecca Wu}
\newline
\newline
連結: \href{https://youtube.com/watch?v=Xbkzaf9FOTY}{\texttt{https://youtube.com/watch?v=Xbkzaf9FOTY}} ~~~~ 語音日期: 2023-10-24
\newline
\newline
\hyperref[sec:RamdTJjsZh8]{\small{< < < PREV SERMON < < <}}
~
\hyperref[sec:index_chronic]{\small{[返順時目]}}
~
\hyperref[sec:index_scriptual]{\small{[返順卷目]}}
~
\hyperref[sec:fRha9TL_FUQ]{\small{> > > NEXT SERMON > > >}}
\newline
\newline
$^{1}$Hi everyone, good to be back..
I was here last week and I think I spoke for about 25 minutes..
So today I'll speak within range again..
Last week I talked about storytelling and for those of you who weren't here, I talk.
about stories because I'm a children's book author and I'm also an artist and an educator..
So stories have a lot of meaning to me personally and I use a lot of stories in my work..
Today we'll be talking about the power of analogy and we're looking at Matthew chapter.
13 verses 24 to 30..
But before we kind of get into it, you'll see the subtitle says "Weed, Wheat or What?".
You're like, "What is that about?".
So before we get into it, let's recap a little about what happened last week..
If you were there, this is a nice reminder..
If you were not, then you got the cliff notes..
So number one, Jesus is the goat..
So he's the greatest of all time storyteller..
And last week we looked at the parable of the lost sheep, both in Luke and in Matthew..
So we looked at the two of them because even though they're the same story about the parable.
of the lost sheep, the way Jesus told it is different, which is why he's the goat..
Because he changes it based on the audience..
Who's watching and who are they?.
Are they the Pharisees and the religious leaders like in Luke?.
Or are they his own disciples like in Matthew?.
A great storyteller also changes the way he or she tells a story based on the context..
So what is happening around them?.
How are they gathering?.
So is it a public setting or is it a private setting?.
And that changes the tone and the voice of what the storyteller says..
[BLANK AUDIO].
(music).
\newpage



\section{19:1-19}
\label{sec:fRha9TL_FUQ}
\textbf{我在這裡作什麼 (列王記上19\_1-19) - 羅懿信牧師}
\newline
\newline
連結: \href{https://youtube.com/watch?v=fRha9TL_FUQ}{\texttt{https://youtube.com/watch?v=fRha9TL\_FUQ}} ~~~~ 語音日期: 2023-10-24
\newline
\newline
\hyperref[sec:Xbkzaf9FOTY]{\small{< < < PREV SERMON < < <}}
~
\hyperref[sec:index_chronic]{\small{[返順時目]}}
~
\hyperref[sec:index_scriptual]{\small{[返順卷目]}}
~
\hyperref[sec:_t1VkAvsN7o]{\small{> > > NEXT SERMON > > >}}
\newline
\newline
$^{1}$大家好.
大家知道我曾經在溫哥華住過.
曾經在加拿大住過.
是一個很好的地方.
我說溫哥華.
不是說其他地方.
那時候在90年代中和90年代末的時候.
我和我太太在溫哥華有幾年時間.
我們在維珍神奴院裝備自己.
回想起來已經20年了.
已經老了20年很多.
那時候我們在神奴院認識了不同的同學.
今天我們說完Anthony是其中一個.
我們很多時候有兩夫婦.
我們一對一.
兩對一.
我們都會出去吃飯.
大家抄功課等等.
大家幫幫大家.
所以在溫哥華有一段短時間都在那裡事奉.
也曾經和Anthony同功過.
很高興這次Anthony能夠經過芝加哥.
他剛剛在這個星期在維珍附近有一個mission conference.
Anthony是OMF的其中一個regional director.
OMF是海外基督使團.
負責亞洲的宣教工作.
所以我們很高興他能夠來到我們中間.
說到加拿大有幾位朋友今天是特別來到的.
Ontario來的.
可不可以舉手.
揮揮手.
我們歡迎他們.
所以我們很高興有Anthony來到我們中間.
來到政道.
我們歡迎Anthony.
各位兄弟姐妹早安.
說到抄功課.
應該Mark沒有抄我的功課.
因為他是我的師兄.
抄功課是我抄他的功課.

$^{41}$很高興來到大家當中.
我想問大家一個問題.
我在這裡做什麼.
糟了.
有一點點思覺失調.
我想問你在這裡做什麼.
這個問題我們廣東人很特別的.
我在這裡做什麼.
或者你在這裡做什麼.
其實可以有兩個意思.
你想清楚一點.
包括「What」.
我在這裡做什麼.
另外也有一點點是在說「Why」.
我在這裡做什麼.
我為什麼在這裡.
我在這裡講道.
為什麼我在這裡講道.
Mark Sir請我.
你在這裡做什麼.
聽道.
為什麼你在這裡聽道.
我空閒時沒有事情做.
不是的.
我只是說這些.
或者牧師迫我的.
這個問題很特別.
我在這裡做什麼.
對於這個可以做「Why」.
又可以做「What」的一個問題.
不知道此時此刻.
大家會在想的是什麼.
我記得大概五年之前.
2018年的夏天.
一個晚上.
我一個人躺在床上.
我問同一個問題.
我在這裡做什麼.
當時我正患上了嚴重的抑鬱症.
我什麼都做不到.

$^{81}$我是一個牧師.
我是教會的主任牧師.
我躺在床上.
我問上帝這個問題.
我在這裡做什麼.
包括你為什麼要帶我來這裡.
你又要我究竟還可以做什麼.
我跟上帝說.
算了吧 上帝.
你拿走我的命吧.
耶和華罷了.
求你取我的姓名.
這個是.
這個是以利亞.
她在跟巴力.
在爭戰的時候.
被皇后耶洗別追殺的時候.
向上帝發出一個呼求.
不知道你這個時候有沒有曾經.
又或者現在.
你也是在向上帝發出.
同樣的呼求.
因為你在爭戰的裡面.
你在危難的裡面.
你在這裡求.
你在問這個問題.
今天選字的經文.
是在《列王記》上.
十九章一到十九節.
這個也是在我抑鬱症的時候.
一段對我很有安慰的經文.
安慰在於我發現原來.
我不是孤單的.
同時也希望透過這段經文.
讓你知道你今天也不是孤單.
經文裡面是這樣說的.
阿哈將以利亞一切所行的.
和她用刀殺眾先知的事.
都告訴耶洗別.
耶洗別就差遣人去見以利亞.

$^{121}$告訴她說.
今日若在這時候.
我若不使你的性命.
像那些人的性命一樣.
願神明重重地降罰於我.
以利亞見這光景.
就起來逃命.
到了猶大的別士巴.
將僕人留在那裡.
自己在曠野走了一日的路程.
來到一顆羅藤樹下.
就坐在那裡求死.
說耶和華罷了.
求你取我的性命.
因為我不勝於我的列祖.
她就躺在羅藤樹下睡著了.
有一個天使拍她說.
起來吃吧.
炭火燒的餅.
她就吃了喝了.
仍然躺下.
耶和華的使者第二次來拍她說.
起來吃吧.
她在當走的路甚遠.
她就起來吃了喝了.
盡著飲食的力.
走了四十晝夜.
到了上帝的山就是荷列山.
她在那裡進了一個洞.
就住在洞裡.
耶和華的說.
臨到她說.
以利亞你在這裡做什麼.
她說.
我為耶和華萬軍的上帝大發熱心.
因為以色列人背棄了你的約.
毀了你的壇.
用刀殺了你的先知.
只剩下我一個人.
他們還要尋索我的命.

$^{161}$耶和華說.
你出來站在山上.
在我面前.
那時耶和華從那裡經過.
在她面前有烈風大作.
崩山碎石.
耶和華卻不在風中.
風後地震.
耶和華卻不在其中.
地震後有火.
耶和華也不在火中.
火後有微小的聲音.
以利亞聽見.
就用愛醫蒙上念.
出來站在洞口.
有聲音向她說.
以利亞你在這裡做什麼.
她說.
我為耶和華萬軍的上帝大發熱心.
因為以色列人背棄了你的約.
毀了你的壇.
用刀殺了你的先知.
只剩下我一個人.
他們還要尋索我的命.
耶和華對她說.
你回去.
從曠野往大馬士革去.
到了那裡就要高哈舌作亞蘭王.
又高靈時的孫子耶護作以色列王.
並高阿伯米荷拉人沙伐的兒子.
以利沙作先知接續你.
將來躲避哈舌之刀的.
必被耶護所殺.
躲避耶護之刀的.
必被以利沙所殺.
但我在以色列人中.
為自己留下七千人.
是未曾向巴力屈膝的.
未曾與巴力親嘴的.
於是以利亞離開那裡走了.

$^{201}$你在這裡作什麼?.
你在這裡做什麼?.
從經文裡我們發現.
上帝兩次問以利亞同一個問題.
分別在聶王記上.
剛才所讀的經文第十九章.
第九節和十三節.
以利亞你在這裡作什麼?.
但兩次所發問的地點是有不同的.
一次是在洞裡.
一次是在山上.
洞裡的以利亞.
事實上又在上帝第一次發問的時候.
他實實在在的在這裡.
他在做什麼呢?.
經文裡說.
以利亞見者光景就起來逃命.
到了猶大的別氏巴.
將僕人留在那裡.
自己在曠野走了一天的路程.
來到一棵羅藤樹下面.
就坐在那裡求死.
說耶和華罷了.
求你取我的性命.
因為我不勝於我的列祖.
他睡在羅藤樹下面.
睡著了.
有一個天使去拍他.
就說起來吃東西.
他看看頭旁有一杯水.
有一個炭火燒的餅.
就吃了.
喝了.
然後再睡下來.
耶和華的師姐第二次來拍他.
說起來吃吧.
因為你要走的路甚遠.
在這一段裡我們看到.
他在做什麼呢?.
他當時的情況是怎麼樣的呢?.

$^{241}$去不到.
他當時的情況.
我在這裡做什麼呢?.
為什麼我在這裡這樣的呢?.
我交給後面這樣吧.
他當時處於一個什麼狀態呢?.
一個三無的狀態.
三無的狀態.
他無什麼呢?.
無處可容.
聖經裡面說.
他做什麼呢?.
逃跑.
他無處可容.
第二個無是什麼呢?.
他無所作為.
他在這裡說.
我還可以做什麼呢?.
他躺平.
是不是?.
他沒有事情做.
什麼都做不到.
那怎麼辦?.
躺平.
我們現在都知道.
躺平的意思是什麼意思呢?.
就是我什麼都做不到.
所以寧願什麼都不做.
第三個呢.
就是求死.
生無可戀.
他說.
罷了.
上帝你取了我的命.
你拿走了我的命.
為什麼他會這樣呢?.
為什麼他會這樣做呢?.
我們看回有一句.
有一個經歷的事蹟.
在美國這個地方裡面.

$^{281}$特別是在1800年代後期的時間.
當時火車的工業.
其實是很興盛的.
在美國的地方.
無論運人或者載貨.
都需要火車去到很多不同的地方.
我們來到芝加哥這個城市的時候.
我發覺還有很多這樣的情況.
不是遺骸.
還在做的.
這些火車.
在這裡有很多的運作.
有一個人.
他寫了一本書.
叫做The Search for Excellence.
這本書裡面.
這位作者Tom Peter.
他就說到.
這個火車的工業.
其實那些巨頭.
他們不明白他們在做什麼.
他們不明白他們究竟在做什麼.
為什麼這樣說呢?.
因為火車一直發展的時候.
很興盛.
到了另一個交通出現的時間.
火車就開始衰落.
那個交通的工具叫做車.
當汽車出現的時候.
火車這個大型的運輸.
就開始衰落了.
這個人這樣說.
其實他們不明白.
他們沒有把握到那時候的機遇.
為什麼呢?.
因為他們在看火車就是火車.
他們在看火車就是工業.
他們看不到原來這個火車.
背後應該是一個交通的工業.
是一個交通的工業.

$^{321}$以致他們沒有把握到機會.
也是當時的時候.
作者就寫了.
他們不知道真正的原因是什麼.
他們不知道為什麼.
不知道原來這個火車所做的.
不只是火車這個物件.
而是背後的交通的行業.
不明白所作為.
所作的為何.
是因為不知道為何而作.
我們很多時候在無處可容裡逃跑.
在無所作為裡躺平.
在生無所戀裡求死.
但我們是否真的知道.
我們為什麼在逃跑呢?.
我們因何在這裡躺平呢?.
很多時候我們求死之前.
我們或者走得太遠了.
我們躺得太久了.
在沒有原因,無因,無故的底下.
我們只求一死.
得以去解脫.
以利亞在這方面比我們更加好.
在他和上帝的對話裡.
他知道自己為什麼逃跑.
何故的在躺.
因何的去尋死.
以利亞說:我為耶和華萬軍的上帝大發熱心.
因為以色列人背棄了你的約.
毀了你的壇.
用刀殺了你的先知.
只剩下我一人.
拋門還要尋索我的命.
他知道為什麼.
因為大發熱心.
所以他逃跑.
因為剩下他一個人.
所以他躺平.
因為有人尋索他的命.

$^{361}$所以他求死.
他知道原因.
我也曾經逃跑.
我曾經躺平.
我曾經尋死.
讓我看到不單止那些觀察.
就是當時的處境.
我在做什麼.
透過那些輔導的過程.
也讓我看到背後的為什麼.
我做傳道人之前.
我是做設計師的.
我大學是讀產品設計的.
這些是我在大學的時間.
一些我的illustration和一些畫作.
當大家看的時候.
你留意到什麼.
有什麼共通點.
當然都是一些產品.
但是很灰沉的.
我用的顏色是很灰沉的.
但是我想告訴你.
我每畫一張illustration的時間.
我開頭所選的顏色.
不是黑黑灰灰的顏色.
很多時候都是比較鮮色的.
但當我在畫的時候.
很多時候我就會不斷在上面再加很多層.
再加很多不同的layer.
又加一個layer.
大家都知道什麼顏色再加.
最後越加得多就會變成黑色.
所以我畫出來的畫作.
當時的教授和同學.
都會經常給一個意見.
就是很灰沉的.
當我在抑鬱症的時候.
我再看這些情況的時間.
我的輔導員讓我看到.
我在做什麼的原因.

$^{401}$為什麼我會不斷地畫.
有幾個片段在我成長的過程中.
讓我再了解我自己的情況.
第一個就是當我弟弟出生的時候.
我弟弟比我小四年半.
當時他出生.
媽媽就在醫院生了我弟弟.
我就和爸爸去醫院接我媽媽和弟弟回家.
我記得當時在拿打數醫院.
不是大埔那一間.
是在香港那邊那一間.
所以你知道是多久之前了.
我不介意告訴你.
我已經52歲了.
在拿打數醫院.
小朋友坐的士回家.
一上的士.
我爸爸和媽媽分別都和我說了一句話.
他說「阿信,你做哥哥了」.
「弟弟將來乖還是不乖,就看你了」.
「你要做個好榜樣」.
當時我不怪我的父母.
他這樣說是沒有任何惡意的.
但原來這句話.
在我整個人生裡.
一直在我裡面.
我就是要乖.
我就是要做個模特兒.
我就是要做好.
後來我做牧師.
我更加加深在我自己的要求上.
而這是其中一個.
我生命裡對自己的要求的「乖」.
而另一個就是.
一直成長,我一直都盡量做好.
我考試,我一直盡量都拿高分.
但是我的成長裡.
在我原生家庭.
我父母的背景裡.
他們當我拿80分.

$^{441}$就會怎樣?.
問我為什麼不拿90分.
我拿90分,為什麼不拿95分?.
我試過拿99分.
很開心地回家.
拿著我的成績給我媽媽看.
我媽媽笑都沒有笑.
也沒有任何的讚賞.
只是問我一個問題.
為什麼你不拿100分?.
下一次我再拿100分.
派成績表回來的時候.
我媽媽看到成績表.
她說為什麼不是全部都是100分?.
我被要求要做榜樣.
但是我卻是在另一個情況裡.
無論我怎樣做都不會做到最好.
在這樣的處境裡.
我只會不斷不斷要求自己做得更好.
而在這個不斷的要求裡.
內心裡又有一種聲音告訴你.
無論你怎樣做都做得不好.
而我甚至將這個代替.
在我地上的父母.
我去預設在天上的父親那裡.
覺得我無論怎樣做.
我的上帝都會告訴我.
我還沒做得夠好.
所以我很喜歡以前的傳道人說.
要為主言進.
我就是不斷不斷不斷不斷不斷.
為主言消.
然後就言進.
我就覺得.
這樣起碼是無功都有祿.
我也言進了.
這個也是我當年.
媽媽不知道知不知道.
我當年去港道的時候.
我練習了20次.

$^{481}$每篇道練習了20次.
我才走上港台.
今天我的師母不在.
你知道當時我讀神學的時候.
她就是三更半夜要聽我讀完3至4次.
我還要說.
不行,你再聽多幾次.
然後她說不行,我要睡了.
我就是這樣.
為什麼呢?.
因為我站上港台的時候.
我覺得我怎樣講都不行.
不過無功都有祿.
起碼我安心一點.
這個就是我.
這個就是那個why.
我為什麼大發熱心.
以至不斷奔跑呢?.
因為我自覺我做極.
我都不夠得不到滿分.
我為什麼孤單一人.
以至我躺平呢?.
因為我做不到帶兒子的好榜樣.
我為什麼在透不過氣裡面.
我寧願去求死呢?.
因為我滿足不了要求.
底下去失去了自我.
今天你知不知道自己可能在逃跑.
但是你知不知道自己為什麼在逃跑呢?.
今天你看見自己躺平.
但是你有沒有想過.
為什麼在這裡睡了呢?.
今天你渴望來到一個終結.
但是你有沒有曾經回望.
在這個去到這個終站的路途.
究竟走過什麼路呢?.
尋恩不是為了找一個理由.
或者找一個藉口.
尋恩是要明白我在做什麼生意.
去看見真正的原因.

$^{521}$聶王記上第十九章.
對於以利亞逃命尋死的過程裡面.
不是要為她去找藉口和找理由.
而是從記載的裡面的何事和為何當中.
讓讀者看到以利亞的生命裡面.
有這樣的反應是為什麼呢?.
以利亞因為大發熱心所以逃跑.
因為剩下一人所以躺平.
因為被追殺所以求死.
但是是不是大發熱心就要逃跑呢?.
是不是剩下一個人就要躺平呢?.
是不是被追殺就要求死呢?.
為什麼當中好像沒有其他的選擇和出路呢?.
因為我們每一個人處於不同的環境.
不同的成長.
不同的際遇裡面.
都會產生出一種的思想行為模式.
而以利亞也有她的一個思想行為模式.
我也有我自己的思想行為模式.
我為什麼知道呢?.
就是當我在病的時候.
我去接受不同的治療的時候.
其中一個叫做認知行為療法.
大家可能都聽過BCT.
Behavioral Cognitive Therapy.
裡面所說的就是基於要去了解我們的想法.
和我們的模式.
我們的觀念.
怎樣影響了我們的情緒和影響我們的行為.
去了解那個病的根本原因是什麼.
背後的理論有兩個假設.
第一個就是我們都受到一些自動化的思維.
Automatic Thoughts 的影響.
這個指的就是我們習慣依賴同一套的思考模式.
而忽略了環境外在的其他的原因.
對我來說就是我做什麼都不夠.
我怎樣都不夠.
我一定要做好榜樣.
而同時第二個假設就是.
我們的身體的反應,行為和情緒.

$^{561}$都是因為由思想而產生的.
這個的確.
為什麼情緒病的人這麼困擾.
就是其他的病你是在外面加諸醫藥的方法.
而情緒病是在腦子裡面.
而你要處理的是腦子裡面的問題.
在我這樣的情況下.
我經常的思想行為模式.
Automatic Thoughts 就是什麼呢.
就是這樣的.
還有還可以怎樣.
你試一下不斷地跟自己說這兩句話.
就是這樣的.
還可以怎樣.
你試一下催眠自己.
你不要回去不斷不斷地說.
我只是做一個假設.
就是這樣的.
還可以怎樣.
就是這樣的.
還可以怎樣.
就是這樣的.
還可以怎樣.
一天是這樣說.
一個星期是這樣說.
一個月是這樣說.
一年是這樣說.
不死都瘋了.
因為大法二心底下去比較.
那個原因.
我想用四個的思去表達.
透過以利亞在我們的生命裡面.
其實很多時候就是這四個的思.
去讓我們變成一種思想行為模式.
第一個的思就是比較.
Compare.
以利亞因為大法二心底下.
以至到他產生逃跑的行動.
他說什麼.
他說我不勝於我的列祖.

$^{601}$他在比較裡面.
落在這種處境當中.
第二個是什麼.
Calculate 計較.
他計較什麼.
他跟上帝說.
我這樣事奉你.
現在怎樣.
我只剩下一個人.
只剩下我一個.
你有沒有試過.
你有沒有跟牧師說.
牧師.
我多辛苦事奉.
我做敬拜.
只有我一個.
我帶團的只有我一個.
我不是說你就是抑鬱.
只剩下一個.
就是覺得我就是最乖的.
我就是最好的.
我最聽話的一個.
第三個可能多一點.
Complain.
Compare Calculate.
Complain.
他跟上帝Complain說.
我為耶和華萬軍的神大法二心.
他們.
他們要拿我的命.
Complain.
然後怎樣.
他就這樣自困.
去到哪裡.
他一直走路.
走到一個他可以停下來的地方.
那個地方叫什麼.
是一個山洞.
他走進了山洞.
將自己圍住.

$^{641}$以利亞在裡面作繭自縛.
或者我們今天都處於類似的狀態.
而人生是不是就是這樣.
故事是不是就是這樣完結.
而以利亞他是不是就老死在那個山洞裡面.
他既然已經來到山洞裡面去作繭自縛.
隨便他吧.
但上帝為什麼又第二次跟以利亞去發問.
再問一個問題.
你在這裡做什麼.
你在這裡作什麼.
問這個問題是要讓他破繭而出.
走出那個山洞.
怎樣破繭而出.
是不是透過認知行為.
療化那個why讓他知道多一點.
還是在watch上讓他看到什麼不同.
我們發覺.
當他再問這個問題的時候.
那個why的時間.
以利亞的反應其實是沒有分別的.
他說的是什麼.
我為耶和華大發熱心.
現在剩下我一個人.
依然的我的命被人尋索.
以前你發覺是完全一樣的.
他跟上帝的回答.
完全是沒有改變.
不過那個what就有些不同.
你在這裡做什麼.
他做了什麼呢.
我們發覺經文裡面說.
他就愛依蒙面.
然後怎樣.
他聽見.
他是不是在山洞裡面.
他走出來.
然後他站在山洞的門口.
就是他站在山洞的門口的時候.
上帝向他顯現.

$^{681}$仍然比較計較自願.
但是不再自困的時候.
願意從山洞裡面走出來.
站在山上的時候.
他終於發現原來他一直在比較當中.
在四個C的裡面.
原來他在當中有四個B.
大家記住了.
四個C四個B.
一會兒還有四個A.
不過今天不是考試.
四個C就是Compare.
Comparing.
我還是要看一看.
Compare Comparing.
Calculate.
Calculate和Constrain.
大家有聽到.
而他站在山洞的時候.
他留意到.
上帝要讓他看見的.
他在比較的裡面.
第一個B.
他break down 崩潰.
他在Compare.
在Calculate的裡面.
計較的裡面.
他的Base Move.
他在自願的裡面的時候.
他Burn out.
他帶來倦怠.
而在他自困的裡面.
他不再Before the Lord.
經歷主的同在.
怎樣看得到呢?.
上帝第二次問以利亞.
你在做什麼的時間.
上帝特別的.
當他站了出來.
在山洞門口.

$^{721}$山上的時候.
藉著不同的景象.
讓他看見.
他仍然自顧.
仍然顧我.
而上帝親自給他答案.
第一個的B.
Break down.
耶和華說.
你站在山上.
在我面前.
那時耶和華從那裡.
經過在他面前.
有裂風大作.
崩山碎石.
耶和華卻不在風中.
風後地震.
耶和華卻不在其中.
地震後有火.
耶和華也不在其中.
火後有微小的聲音.
耶和華在那裡.
向他講說話.
這裡我們看見.
在裂風的裡面.
英文的聖經特別說到.
The strong wind tore the mountain.
and broke into pieces.
崩潰.
崩壞.
Break down.
而地震的處境.
說的是什麼呢?.
地震的時候.
說的是他腳不能夠站穩.
他失去平衡.
而火的形容.
要帶出來的是什麼呢?.
當火出現的時候.
就會被燒盡.

$^{761}$就會帶來的是.
Burn out.
絕桌.
而最後.
上帝要告訴他.
在微聲的裡面.
他在當中.
Before the Lord.
弟兄姊妹.
你在這裡做什麼呢?.
在自動化的思維.
Automatic force 底下.
是這樣的.
還可以怎樣呢?.
在比較裡面.
去Break down 崩潰.
在計較裡面.
去Base move.
失衡.
在自願的裡面.
去Burn out 絕桌.
在自困的裡面.
不再Before the Lord.
經歷主的同在.
如果是這樣的時候.
我們或者應該問的是.
你在這裡做什麼呢?.
上帝在這裡做什麼呢?.
四個C.
Compare.
Calculate.
Complain.
Constrain.
四個B.
Break down.
Base move.
Burn out.
Before the Lord.
這個包括我們個人生命的裡面.
也包括教會的生態裡面.

$^{801}$我們需要四個A.
經文裡面提到的.
當上帝去呼召爾利亞.
跟他說話的時間.
兩次的跟他說.
你在這裡做什麼的時候.
是不是就這樣說.
誰誰.
你在這裡做什麼?.
又或者甚至乎.
就這樣走過去.
你在這裡做什麼?.
不是吧?.
你留意經文.
兩次的.
聖經也很強調.
是將爾利亞的名字擺出來.
爾利亞.
你在這裡做什麼?.
為什麼要叫他的名字呢?.
因為他的名字是有意思的.
爾利亞的意思是什麼?.
爾利亞的意思.
Elijah就是耶和華是神.
Jehovah is Lord.
上帝告訴他.
就是跟他說話.
問這個問題的時候.
已經發出了提示.
不是看你.
不是你是誰.
而是我是誰.
你知不知道.
你的名字的意思就是.
我就是神.
爾利亞是誰.
是基於耶和華是誰.
而當他一直的.
上帝繼續向他顯現的時候.
又有第二個的A出現.

$^{841}$就是Aligned.
在affirm了他的身份之後.
耶和華對他說.
你回去.
從曠野往大馬士革去.
都了那裡就要給哈舌族.
就要告哈舌族亞蘭王.
又告靈士的兒子耶護族以色列王.
並告阿白米荷拉人約沙發的兒子.
約沙發的兒子.
阿拉人沙發的兒子.
爾利沙族.
先知接續你.
這裡面有幾個的名字.
看上去好像.
中執一個西執一個.
但是原來幾個名字是有意思的.
哈舌的意思是什麼呢.
One who sees God.
一個見到神的人.
耶護的意思是.
Jehovah is he.
他就是神.
爾利沙的意思是.
God is salvation.
神是拯救.
你發覺.
首先上帝不斷不斷跟他說.
爾利亞爾利亞.
告訴他我是誰.
然後說完我是誰的時候.
告訴他.
你要見我.
你要在當中的時候.
你要看到我是誰.
而我就是拯救.
然後上帝繼續說.
在以色列人裡面.
我為自己留了七千人.
是未曾向巴力屈膝.

$^{881}$未曾與巴力親嘴的.
這裡說到為自己留了七千人.
這裡不是說.
爾利亞我為你留了七千人.
很多時候我們今天.
事奉到很疲倦的時候.
這句經文很大的鼓勵.
我們就去找那七千人.
那七千人要去擁捉我的.
要去攀附我的.
要去站在我身邊的.
錯.
這裡說的.
上帝提到爾利亞.
我是為自己.
再一次將焦點放在上帝那裡.
聖文是哪一個.
基於的是知道.
為他存留的是主.
而最後的.
於是爾利亞就怎樣呢.
就離開那裡走了.
遇見沙伐的兒子.
爾利沙耕地.
在他前面有十二頭牛.
自己趕著第十二隊.
爾利亞到他那裡去.
將自己的外衣搭在他身上.
聖公為誰而作呢.
是聽候主的差遣.
我在這裡做什麼.
原來關鍵的.
透過剛才那四個A.
你發覺焦點在誰那裡.
不是在我那裡.
原來一路的問題就是.
將焦點放在我那裡.
而我看不到的原來焦點在上帝那裡.
不是在說我.
甚至乎不是在說What和Who.

$^{921}$和那個Why而是在說Who.
在維珍神學院裡.
我和Mark有一位很好的教授.
Eugene Peterson.
他寫了《聖經》的譯本.
在裡面的其中.
在羅馬書第十三章第三節那裡.
中文的翻譯是.
我憑著所賜給我的恩對你們.
說不要看自己過於所多看的.
要照著上帝所分給各人的信心.
但是少看得合乎中度.
這段《聖經》我們好像說.
我們不要做得太好.
我們也不用做得太差.
中規中矩就行了.
但是其實真正的翻譯是什麼呢.
Eugene Peterson有一個很好的翻譯.
他說.
(英文翻譯).
Do not misinterpret yourself.
不要錯誤地去詮釋你自己.
上帝兩次的問以利亞.
你在這裡做什麼.
一次的在山洞裡面.
一次的在山上.
兩次都是問你在這裡做什麼.
有一次在山洞裡面.
但是第二次的時候.
當上帝問的時候.
以利亞就站出來.
離開山洞.
在上帝的面前.
今天上帝依然問你.
你在這裡做什麼.
你又身在何處.
你是在山洞裡面.
住在那裡.
作繭自困.
還是你在山上.

$^{961}$去到主的面前呢.
提醒我們一個救恩的寶貴.
完全的我們在基督裡面.
就被你鑽研的節流.
赦免我們的罪.
以致我們能夠在你的裡面.
知道你在我們生命當中.
是這麼豐盛的.
主我們求你去矚目我們每一個.
當我們可能走過低谷的時候.
我們不知道自己在做什麼.
主讓我們抬頭仰望你.
聽見你的聲音對我們說.
上帝你罪與我們同罪.
在我們心靈當中一主.
我們這樣祈禱奉祝你救命.
再次多謝Anthony給我們的訊息和提醒.
Anthony也寫了一本書.
剛剛出爐沒有多少東西.
叫做《我的牧師還了幽屈症》.
這本書是一本很好的書.
他自己用自己的經歷.
包括聖經和經過自己的治療.
他自己的經歷.
來述說這麼重要的一個課題.
我相信我們華人教會很怕說這件事.
因為我們華人的文化不太能接受.
或者我們的屬靈的傳統.
不太能夠接受這件事.
但這件事很重要.
我相信這本書是一個很大的貢獻.
我們有20個複製.
我們只是在外面有20個複製.
如果你可以出去的時候.
你可以買這本書.
不單止對你自己有幫助.
也對身邊的人有幫助.
願上帝祝福各位.
\newpage



\section{}
\label{sec:_t1VkAvsN7o}
\textbf{聖經地理 EP1 - 地理神學 (陳崇基牧師)}
\newline
\newline
連結: \href{https://youtube.com/watch?v=_t1VkAvsN7o}{\texttt{https://youtube.com/watch?v=\_t1VkAvsN7o}} ~~~~ 語音日期: 2023-10-30
\newline
\newline
\hyperref[sec:fRha9TL_FUQ]{\small{< < < PREV SERMON < < <}}
~
\hyperref[sec:index_chronic]{\small{[返順時目]}}
~
\hyperref[sec:index_scriptual]{\small{[返順卷目]}}
~
\hyperref[sec:H8JKxV9_GHo]{\small{> > > NEXT SERMON > > >}}
\newline
\newline
$^{1}$好,這次我們就講地理.
《聖經》的地理.
我對這方面很有興趣.
因為除了考古之外就是地理.
考古和地理當然是兩者非常有關係.
因為地理就是講到在哪裡.
地點.
然後考古就是背景.
就是那個地方的歷史.
那個歷史的背景.
所以兩者都走在一起.
Geography和Archaeology.
我們首先有些概念來開始.
為什麼要讀《聖經》的地理呢?.
很多時候我們都覺得這件事不是很重要.
因為你知道大家的聖經.
特別是舊的聖經.
紅色那些印出來的.
裡面有地圖.
地圖夾雜在不同的地方.
或者集在最後.
我們很多時候都不會看.
因為我們覺得最重要的是經文.
當然最重要的是經文.
因為上帝的啟示就是在.
在他的話語.
這些text.
就是經文的書.
但是經文有背景.
正如我們之前.
比對的研究也好.
或者背景也好.
我們經常都要解釋的.
有些東西是翻譯不了出來的.
甚至你就算看原文也好.
原文希伯來文.
或者希臘文.
背後都有它的背景.
就算我們看原文也好.
作者寫的經文給讀者.

$^{41}$其實都有一個大前提在那裡.
就是預設了大家知道那個背景.
知道那個地理.
知道裡面提到那個地方的名字是什麼名字.
或者那個地帶是什麼地帶.
或者那個地理的動態是怎麼樣.
這些全部都是有預設在那裡的.
假設了的.
所以我們更加.
作為現代的信徒.
我們拿一本聖經的時候.
翻譯了我們現代的語文.
我們就和那個的距離更加遠了.
但是不要緊的.
我們有一些工具可以幫到我們.
但是我們先拿捏到一些概念.
什麼概念呢.
就是一些神學.
系統神學的一些概念.
其中一個的概念就是.
講到聖經的本質.
我們要知道聖經的本質.
是百分之一百神的話語.
是不是.
上帝啟示他的說話.
聖經每一句的說話.
都是有神聖的權威的.
是神的話語.
但是同時我們不要忘記.
也是百分之一百人的話語.
是不是.
希伯來文不是一些魔術的.
外星的一些特別的語言.
都是人的語言.
希伯來文.
希臘文也是一樣.
都是當時人用的.
一些常用的語言.
所以讀者和作者.
就能夠去溝通了.

$^{81}$正如我們現在能夠溝通了.
我們都懂得這種話.
能夠聽得懂.
裡面的一些概念.
我們用這些說話去表達的時候.
都是人的說話.
所以是百分之一百神的話語.
但是不要忘記.
百分之一百都是人的話語.
在十六世紀的時候.
有一個改教的神學家.
就是宗教改革.
馬丁路德等等.
其中一位很著名的改教家.
叫做嘉爾文.
嘉爾文就是今天的獎勵會.
改革中的始祖.
他所引申出來的傳統.
他就提出了一個概念.
叫做神的遷就的教義.
什麼叫神的遷就呢.
神的遷就的意思就是說.
人是不能夠真正認知到神的本體.
God in himself.
因為神太高了.
太超越我們了.
所以我們是不能夠.
talk about God.
我們不能夠去描述神.
或者講論神.
只有神先開口跟我們說.
我們才能夠認知到神自己.
或者神自己無限本體的其中一部分.
我們自己認知都是非常有限.
所以同時聖經也告訴我們.
神與人 神與我們之間是可以溝通的.
建立關係.
為什麼呢.
嘉爾文就認為.
神一定是將他自己降臂.

$^{121}$他先去跟我們溝通.
而且他溝通的方法.
是遷就了我們人的程度.
就是用我們人類的語言.
希伯來文 希臘文.
這些不是法術.
我們去讀希伯來文 希臘文的時候.
其實都是讀一些完全100\%是人的語言.
但上帝就是遷就他自己.
用一個適合我們聽的程度.
向我們溝通.
等於我們跟一個小朋友去講話.
我們就不可以講得這麼複雜.
我們要遷就我們的概念.
或者用語言.
以至小朋友能夠明白.
上帝是無限的上帝.
跟我們溝通的時候.
他也是一個降臂.
所以這個啟示的教義.
是看到上帝降臂自己.
遷就他自己.
以至我們能夠認識神自己.
這個我們曾經講過.
每一段經文背後.
都有本身經文裡面的上下文.
由段落至一卷書.
甚至整個舊約聖經.
或者新約聖經.
或者整個正典.
也有經背後的世界.
聖經背後的.
有經外文獻 歷史文化 地理 考古等等.
所以為什麼我們需要對聖經地理.
或者聖經裡面印的地圖.
有多少認識.
能夠幫助我們認知聖經裡面的信息.
我們回到這個簡圖.
也告訴我們.
有世界背後的聖經.

$^{161}$聖經背後的世界.
有近東的背景.
有考古的背景.
地理等等.
聖經裡面的世界.
所謂的上下文.
我們看經文一定要看上下文.
文理結構 修辭技巧等等.
聖經前面的世界就是我們.
我們看聖經的時候.
我們都要知道.
我們都接受了某些包袱.
或者我們有一些背景.
有一些傳統.
這些傳統會左右我們解釋聖經.
或者對聖經的認知.
甚至我們現代人的文化價值.
都會左右使我們不能掌握聖經的意思.
所以在這三個角度裡面.
我們需要嘗試去用不同的工具.
讓我們能夠自我批判.
以致我們看聖經的時候.
可以用不同的角度去看.
大家可能玩一個桌面遊戲.
我很喜歡用這個比喻來講聖經地理.
桌面遊戲 例如大富翁等等.
有些桌面遊戲我們不懂得玩.
或者有些人玩一些撲克牌的遊戲.
我們不懂得玩 為什麼呢.
他們玩得那麼開心.
為什麼我們不懂得玩呢.
或者不明白他們在做什麼.
原因就是他們知道遊戲規則.
知道桌面到底在做什麼.
規則 遊戲的規定等等.
我們懂得這些規則之後.
或者桌面應該怎樣運作的時候.
我們才能夠玩遊戲.
其實看聖經有一點像這樣.
認識地理 聖經的地理.

$^{201}$其實就是有這樣的好處.
知道地理 政治上 文化上.
物質上 物理上 怎樣運作的時候.
讓我們知道聖經背後所假設的地理動態.
我們就能夠更加清楚聖經裡面所講的東西.
三一神學院在我們芝加哥附近.
一間很出名的科文派神學院.
有一個學者 一個舊約的老師.
叫做John Monson.
John Monson是做了考古地理很多年.
是北美甚至世界對聖經地理最有研究的.
其中一個最出名的人.
一個地理學家 聖經地理學家.
他經常說.
我們建立自己的門徒生命的時候.
我們是用上而下的方法去講述.
一個人信了主.
我們第一件事就是跟他栽培.
去建立他的生命的時候.
我們通常都是上之下.
告訴他你應該怎樣看聖經.
然後怎樣祈禱.
教會的生活是怎樣.
或者對聖經的認識是怎樣.
我們是上之下.
無他的 很多時候都是這樣.
我們一個人信了主.
教會開始的時候.
我們就認識怎樣做一個好的基督徒.
怎樣去做一個神的子民.
然後我們才開始講.
下面不同的科目.
當然這些就是神學院用的用詞.
但是我們教會.
我們去跟大家查經.
去教導聖經的時候.
其實我們都是涉獵這些不同的科目.
譬如解經學 怎樣看聖經.
大家怎樣看耶穌的比喻.
大家怎樣看保羅的書信.

$^{241}$這些都是解經學的東西.
怎樣查經 怎樣去看聖經.
這些就是做了門徒之後.
我們才慢慢地找一些原則.
來教給他.
以致他能夠成熟一點.
聖經神學就是譬如.
什麼叫聖經神學.
就是聖經由創世記到啟示錄.
有很多的主題.
這些主題就是不斷地在歷史的進程裡面.
在救恩歷史裡面.
不斷地去演變出來.
不斷地有改變.
即是說神對人的救贖計劃.
都是按著不同的時候.
他有不同的啟示.
或者更加多的描述等等.
譬如約 聖約這件事.
上帝呼召了亞伯拉罕.
說要跟他立約.
然後在他身上出以色列國.
這個約就是亞伯拉罕之約開始.
接著就去到摩西的時代.
就是摩西的西乃之約.
上帝因點拯救了他的子民.
然後就出了埃及.
去到西乃的時候就頒布了法給他們.
他們怎樣做一個立約的子民等等.
然後繼續發展下去大衛的約.
到最後就是新約.
就是耶穌所承傳的.
所應驗的.
整個以前的所有重約的應許.
所以就兩樣東西都集中在新約.
即是基督.
所以這個就是聖經神學.
這個主題在聖經由頭到尾怎樣去發展.
這個都是我們做了門徒之後.
慢慢慢慢從教導和學習裡面才認識.

$^{281}$系統神學就有些不同.
系統神學就是整本聖經對某主題的教導是什麼.
譬如整本聖經對天堂的教導.
對新天新地的教導.
這個就不是說某一個主題.
由創世記開始怎樣去發展出來.
然後基督成為一個最高峰等等.
整體的就是總結.
整本聖經對某一些主題怎樣去教導.
譬如罪 人倫 神倫 基督倫等等.
這些都是系統神學裡面的東西.
所以一個人做了門徒.
首先開始認識了一些東西.
然後慢慢慢慢才去接觸.
什麼叫基督.
基督神人異性是什麼意思.
為什麼這麼重要.
然後有實踐神學.
或者有宣教神學.
這些都是比較實踐性.
怎樣去做講道 講道法 教導 基督教教育.
宣教學 怎樣傳福音 宣教等等.
這些都是由開始上面建立下來的.
然後到最後的時候我們才會想一下.
原文到底講什麼.
背後有沒有歷史背景呢.
我們才開始發現多一點.
然後有考古的背景.
當時的人的生活是怎樣等等.
地理 原來在以色列這個地方是怎樣的.
所以為什麼經文會寫這些東西.
有什麼地理獨特的背景.
然後文化 他們近東的價值.
他們的榮辱是怎樣定義.
什麼叫榮耀 什麼叫羞辱.
他們怎樣去定義一個人.
他們不是用個人的方法去定義.
他們用一個群組 家族 宗族去定義.
所以家譜這麼重要等等.
下面黃色這些 橙色這些.

$^{321}$都是我們很多時候最後才會去接觸的.
John Monson 就是說.
其實應該是相反的.
首先我們要認識語言.
起碼一些基本的語言原則.
希伯來文 希臘文.
歷史的背景 聖經有它的歷史背景.
好像一個桌面遊戲.
有它背後的桌面規律.
哪個歷史發生什麼事.
亞述帝國的時候發生什麼事.
巴比倫帝國的時候發生什麼事.
然後才有上帝在中間去救贖.
和對他的子民講話.
考古一些生活上的.
那些人到底他們的物件是用什麼.
以致那些聖經會描述他們當地文化的.
社會的一些背景.
地理就是我們這次的重點.
那些地圖到底有什麼動態.
所謂的Geopolitics 地理的政治.
當時的地理政治.
埃及等等.
和以色列 巴勒斯坦之間的地理互動是怎麼樣的呢.
這些就是最基本的東西.
文化背景當然是.
John Monson 就是說.
應該要由下至上 不是由上至下.
就是說我們做了門徒之後.
信了耶穌之後.
應該起碼對這些東西越來越掌握.
以致我們才能夠建立對聖經的認識.
解經 系統神學等等的東西.
然後才能夠讓我們門徒的生命.
能夠更加站得穩.
這件事他叫做地理正典神學.
Canonical Theology.
就是都是在做神學.
因為每個人都有神學.
你的神學可能是在主要學校.

$^{361}$在港台那裡接觸回來的.
有些是在仁仇院接觸回來的.
但是這些神學應該是怎麼樣做的呢.
John Monson 就認為是Geo-canonical.
為什麼他加了個Geo 地理進去呢.
因為這些是最基本的東西.
所以我們的神學應該是由下面做上去.
不是由上面才做下來.
所以才打好這樣的基礎.
今天講地理就是想講這件事.
所以在我們現在講地理的時候.
我們大致上就會接觸到.
這些最基本最根本的背景.
\newpage



\section{}
\label{sec:H8JKxV9_GHo}
\textbf{Israel-Palestine Conflict - A Reflection (1) - Rev. Mark Chan}
\newline
\newline
連結: \href{https://youtube.com/watch?v=H8JKxV9_GHo}{\texttt{https://youtube.com/watch?v=H8JKxV9\_GHo}} ~~~~ 語音日期: 2023-10-31
\newline
\newline
\hyperref[sec:_t1VkAvsN7o]{\small{< < < PREV SERMON < < <}}
~
\hyperref[sec:index_chronic]{\small{[返順時目]}}
~
\hyperref[sec:index_scriptual]{\small{[返順卷目]}}
~
\hyperref[sec:giSEqlRniIQ]{\small{> > > NEXT SERMON > > >}}
\newline
\newline
$^{1}$- Thank you, James..
Thank you so much for the church members.
who express in their own ways.
to appreciate the pastoral staff..
I want to tell you something,.
even though I'm a lead pastor,.
as of today, I have probably the best team.
of pastoral workers on my team,.
so I really appreciate their,.
they're always hardworking..
I don't have to keep track of their time..
They are always working more than they should,.
so I just want to give a shout out to them.
and really appreciate their hard work..
Well, today, we are going to take a detour..
We're not going to talk about Providence.
for maybe two weeks or so..
I think there are times we need to address.
and think about or reflect upon.
what is going on in the world,.
and we as Christians are being called to be peacemakers.
and to be the bearers of Christ's image in the world,.
so today, we take a short detour.
and think about and reflect upon.
what is going on in Israel-Palestine..
As you know, Israel is probably.
kind of like my second home, all right?.
Kind of like my second home..
Actually, I'm planning to go back,.
leading a tour in May next year..
I hope everything will deescalate,.
at least in the next few weeks or next month or so,.
and so today's sermon is gonna be shorter.
because we have a fire drill..
However, I wanna switch the topic and switch our gear.
to think about what is going on in this piece of land, Gaza..
As you can see on the map,.
Gaza Strip is at the lower left corner.
of the long corridor of Israel,.
and you can see right there, there's a yellow line,.

$^{41}$and Gaza Strip is a small piece of land.
on the Mediterranean Sea along the coast.
and in the southern part of Israel..
Israel-Hamas conflict has captured.
the attention of the world, that's for sure,.
and I'm sure you read many, many news stories.
and many horrible images and videos.
on what is going on and what Hamas has done.
to the Israelis on October 7th.
was just pure horrendous terrorist act,.
especially for the civilians..
Countless, many, many Israelis were terrorized..
Some, many of them are kidnapped,.
and it's horrible to just think about.
what had happened to them and what will happen to them..
And now with a lot of airstrikes.
by the modern state of Israel,.
bombing and ground invasion.
has already taken place in Gaza..
Thousands of people, including the Palestinians in Gaza,.
their lives are being threatened..
Now, there's a danger, there's a danger..
Whenever war breaks out,.
it's particularly in this area,.
and you don't want it to spill out into a regional conflict,.
and there is a possibility of this danger happening.
because in the northern border of Israel.
is the Hezbollah, another terrorist organization.
funded by Iran, and there were different political actors.
coming onto the scene,.
and if this conflict continues to rise,.
it will be a disaster for everyone, everyone,.
not just Israel, but mostly for the Palestinians..
So we have to think about this topic.
because Israel-Gaza conflict seems far, far away..
They are eight hours ahead of us..
With the fallback, it's nine hours away from Chicago,.
but it hit home, even in the United States,.
because 13 days ago, October 16,.
hate crime was committed in Plainfield..

$^{81}$Plainfield is 40 minutes away from Chicago..
What happened was a landlord was overwhelmed.
by his fear for Muslims.
in light of what's going on in Gaza, Israel..
He went into the home of a Palestinian family, a tenant,.
and he went inside and stabbed a six-year-old boy 26 times.
and murdered the boy on scene and wounded his mother..
It's Islamophobia, right, because what is going on.
and what you read on the news stories and all that..
And October 21, October 21,.
the president is a lady rabbi.
of a Jewish synagogue in Detroit..
She was found fatally stabbed outside her home..
Now, whether this was a hate crime, we don't know yet,.
but this tragic story is definitely the tip.
of the iceberg of antisemitism..
It is on the rise..
With my years of listening to stories.
from both sides, Israel and Palestine,.
I observed two things..
I think there is one set of responses.
or one set of viewpoints that are not so helpful,.
and other set of viewpoints or perspectives.
perhaps more helpful..
Whenever a conflict arise in this piece of land,.
people are very quick to go to Bible.
and biblical prophecies being fulfilled and all that stuff..
I do not think this kind of perspective is helpful..
Perhaps we need to study,.
that's why we need to study the Bible, right?.
That's why we need to study the Bible in context, right?.
I don't think it is so helpful.
when we just see what's going on in Israel, Palestine,.
and then we quickly flip the pages of the scripture,.
and then we find this prophecy fulfilled.
or this prophecy fulfilled, and this is going to happen,.
and this may happen..
We have to examine carefully what the Bible claims.
according to its context..
And I hope to address some of the popular viewpoints.

$^{121}$in this short detour..
But let me make this clear..
Wars are ugly and evil..
And we as Christians, especially we believe in Jesus Christ,.
we are to pray for those who suffer,.
and we are being called to do peacemaking, right?.
One of the most quoted Bible passages.
when people talk about the Israel-Palestine debate.
is Genesis 12..
Genesis 12..
Here the Lord, or God, is calling Abraham, or Abraham,.
and he said to him, "Leave your country, your people,.
"your father's household, and go to the land,".
which is the land of Canaan, eventually,.
"and I will show you..
"I will make you into a great nation,.
"and I will bless you..
"I will make your name great, and you will be a blessing,.
"and I will bless those who bless you..
"Whoever curses you, I will curse..
"All the peoples on earth will be blessed through you.".
Now, this passage has often been quoted.
as God's unconditional promise to Abraham,.
the land of Canaan..
The land of Canaan is given to Abraham,.
the father of Israel,.
therefore, it's given to all of his descendants,.
and therefore, as sons and daughters of Abraham,.
you have the deed, you have the title.
to go into the land and take the land..
This is what we call sometimes the Abrahamic covenant..
It means God is making a covenant to Abraham.
and to his generations after him..
And sometimes we quote this passage,.
and it seems to settle the entire debate.
of Israel-Palestine..
The land belongs to Abraham and his generations after him.
being granted to this land unconditionally..
And some would say this is God's land grant.
to Abraham, the father of Israel..

$^{161}$But when you ask more questions,.
it doesn't seem very straightforward..
It doesn't seem as straightforward as we expected..
First, what do we mean by unconditional?.
Now, we use it almost like a cliche for us,.
unconditional, right?.
God's unconditional promise, unconditional blessing..
But sometimes we just don't think enough, right?.
Now, if you think about the background of the Bible,.
all covenants take two parties,.
just like a marriage covenant, right?.
Marriage covenant takes two parties to work together, right?.
Take two people working together,.
committing themselves together..
This is the covenant background of the Bible..
So when we say God is making a covenant with Abraham,.
we have to think very, very careful..
Is it when we say unconditional,.
are we talking about God is making a unilateral commitment?.
Meaning that God is gonna give Abraham this land.
no matter what he does..
If he does bad, he still get the land..
If he does good, he still get the land..
Is that what we mean?.
Now, as you can see in the last line of the slide,.
I wanna do some grammar with you..
I know we hate grammar, right?.
Hate all sorts of grammar, English grammar,.
Chinese grammar, or whatever grammar, right?.
Now, if you take into the original verbs of this verse,.
this passage, you find that some of the key verbs.
are cohortative..
Now, what we mean by cohortative is that.
it's not so much of I promise you something, anything goes..
No stranger attached..
Cohortative can be translated, or more correctly translated,.
let me do this..
Let me bless you..
Let me make your name great..
Now, cohortative seeks to communicate potentiality..

$^{201}$That means there is a possibility,.
there's a potential situation we can reach, we can achieve..
But let me do this, and you being cohortative.
and collaborate with me..
So let me bless you, and you trust in me.
so that you will be a blessing..
You see what I'm saying?.
There are two parties, right?.
It's not a unilateral commitment..
Let me bless you instead of I will bless you..
Let me make your name great.
instead of I will make your name great no matter what..
Now, that's a little bit grammar..
And in fact, if you read this passage this way,.
it makes more sense when you read the entire story,.
the entire narrative of Abraham..
In the book of Genesis, chapter 12 through 22,.
the entire journey of Abraham.
is that God is calling him out of Ur,.
and then he goes to the land of Canaan..
If Abraham respond positively, faithfully,.
he can make this potential situation a reality..
Only when he obey, then this becomes a reality..
If he disobeys, if he detours,.
then there are some barriers..
Something happened and put the covenant in jeopardy..
Now, 12, one, two, three is this passage..
You know what happened in verse four?.
12, verse four..
Let me read this to you..
Abraham went according to the word of God,.
according to whatever God said to him,.
which is verses one to three..
So Abraham responded positively..
And in the entire story, whenever Abraham acts obediently,.
the plan of God advances..
Whenever Abraham does things on his own,.
with his own means, then the plan detours,.
and there's a threat to the promise of God..
So all covenants in the Bible have two sides..

$^{241}$Both sides, both parties, in this case, Abraham and God..
Abraham's descendant, Israel, and God have to work together..
Another theory, or another popular saying,.
is that this land promise was not fulfilled,.
was never fulfilled until 1948, the modern state of Israel..
Now let's examine, or try to trace,.
the theme of land in the entire Bible..
I told you many, many times, the Bible is a grand story,.
from Genesis, the first book, to Revelation, the last book,.
Old Testament and the New Testament..
It's a continuous grand story, grand narrative,.
to tell us what God has done in human history,.
in the midst of us..
From the first book to the last book, it's a grand story..
So there are different themes in the Bible.
that was introduced in the beginning,.
and you'll find that when you trace it,.
it will develop and it will evolve into something.
or becomes a reality..
And as we know, all the themes,.
all the major themes in the Bible introduced.
in the beginning, including the land promise,.
find its ultimate fulfillment in Jesus Christ,.
in Christ our Lord..
So you have to be sensitive when you read the Bible,.
you cannot pick and choose..
You have to read a passage and know where the passage fits.
in the entire narrative of the Bible..
So as the Bible story develops,.
the sons of Abraham, Isaac and Jacob,.
they went down to Egypt and soldiered there for 100 years..
And eventually God was gracious.
and God was faithful to his promise..
And he called Moses and led the people of Israel.
out of Egypt, that was the Exodus, right?.
And then they came out of Egypt and end up in wilderness.
for 40 years and they did not go to the land..
They did not get to enter the land because of their rebellion..
The first generation did not enter the land except Moses..
Even Moses did not enter the land of Kingdom,.

$^{281}$except two guys..
Who are they?.
Caleb and Josh, right?.
They are the only two guy from the first generation.
of Exodus had the blessing to enter the land..
And then with the second generation.
who were born in the wilderness..
So finally, under the leadership of Joshua,.
they went into the land, the new generation,.
the second generation of faithful Israelites,.
they went in and they inherited the blessing of God.
because they put their faith in God..
Again, there's a conditionality there..
There is a criteria there..
And the criteria is trust in God.
and you will inherit God's blessing..
So as we turn to the end of the book of Joshua,.
we see this text..
So the Lord gave Israel all the land he has sworn.
to give their forefathers, Abraham, Isaac, and Jacob..
They took possession of it and settled down there..
The Lord gave them rest on every side,.
just as he has sworn to their forefathers..
Not one of their enemies withstood them..
The Lord hand all their enemies over them..
45, very important, 45..
Not one of all the Lord's good promises.
to the house of Israel failed..
Not one failed..
Everyone was fulfilled..
Doesn't get as clear as this..
Everyone is fulfilled..
Every promise, not even one word,.
not even one word to the house of Israel failed..
Everything is fulfilled..
The land promise was fulfilled..
Now before Joshua passed away,.
he also warned the people of God,.
the people of Israel, 23 verse 14..
I'm about to go the way of all the earth.

$^{321}$and you know all my heart and all my soul.
and not one of all the good promises.
the Lord your God gave you has failed..
Every promise has been fulfilled..
Not one has failed..
You keep on going, 16..
There was a condition there..
So now you settle down in the land..
Now you get the blessing, right?.
However, there was also a condition..
If you obey the commandment of God,.
the covenant of God, then you will stay..
You will enjoy the presence of God, his blessing..
But if you violate, verse 16,.
the covenant of the Lord your God,.
which he commanded you and go and serve other gods.
and bow down and worship them,.
the Lord's anger will burn against you..
You will quickly perish from the good land.
he has given you..
See, the land is not a free giveaway, you know..
It doesn't seem that way, right?.
Even within the Old Testament,.
we are by the way in the sixth book of the Old Testament..
Genesis to Deuteronomy, the Torah, five books..
This is the book following the Torah..
And people of God are expected to live faithfully..
If they violate the covenant,.
they will be driven out of the land..
Remember the original people in the land of Canaan,.
the Canaanites, the Amorites,.
they were driven out of the land because what?.
Because their sin has reached its full measure..
That's why God is driving them out, right?.
Now, as if the Bible is not clear about the fulfillment.
of the land promise to Abraham in the Old Testament,.
the Bible continues to remind us..
Again, we trace that theme beyond the book of Joshua.
and we come to 1 Kings..
This is the time of David and Solomon..

$^{361}$1 Kings 4.20, the people of Judah and Israel,.
they were now 12 tribes united as numerous.
as the sand on the seashore..
They ate and they drank and they were happy..
Solomon, the king, ruled over all the kingdoms.
from the river, Euphrates, to the land of Philistine,.
as far as the border of Egypt..
All these countries brought trouble.
and were Solomon's subjects all his life..
Keep on going..
In the dedication prayer of the first temple,.
Solomon said, "Praise be to the Lord,.
who has given rest to his people Israel.
just as he promised.".
Not one, again, this phrase again,.
"Not one word has failed of all the good promises.
God gave to his servant Moses.".
It seems clear that the land promise.
however you want to describe it,.
conditional or unconditional,.
however you describe it,.
God's promise, the land promise,.
is already fulfilled within the time of Old Testament..
By the time of Solomon,.
all the promises have fulfilled..
Have fulfilled..
That's why when they settle down in the land,.
the Israelites are being warned again and again,.
they need to be faithful to the covenant..
They need to exercise their faith and trust in God.
so that they can enjoy the good land,.
enjoy the presence of God..
If they don't, if they are faithless,.
then they will be kicked out..
And that's why eventually they went into exile, right?.
That's how the Old Testament story goes, right?.
Again and again in the Torah or in other verses,.
God warned the people of Israel,.
when you settle down in the land,.
when you enjoy the presence of God,.

$^{401}$you have to do something and to be faithful to the covenant,.
faithful to the commandment of God..
One of them is to not mistreat an alien or oppress him..
For you were aliens in Egypt..
You had that experience before, right?.
And so you don't oppress other people.
with a similar situation..
It's kind of like us..
Some of us are immigrants, right?.
We come to this land, we come to this place..
Sometimes, you know, sadly, we oppress other people, right?.
Remember when you first came to the United States,.
it was a tough time, right?.
I used to work in a restaurant,.
just do dishes for 10 hours a day,.
cold water, no dish washing machine..
Working in the restaurant 10 hours a day, right?.
It was really, really tough, right?.
Because I, you know, I don't come from a rich family, right?.
I have to work my butt off, right?.
If I turn around, now I'm settled down, right?.
If I turn around and oppress other people,.
this is really sad..
Exodus 22, do not mistreat an alien.
because you used to be an alien in Egypt.
and you were oppressed..
That's why you need to treat them kindly, right?.
That's part of the terms within the covenant..
Leviticus 18, it says, "Even the land was defiled.".
Now, it was talking about the Canaanite,.
the people before Israel took over the land of Canaan..
So God said, "I punish them for a sin..
"And the land vomited out its inhabitants.".
And now you guys, the Israelites, 26,.
"You must keep my decrees and my laws..
"The native born, the aliens living among you.
"must not do any of these detestable things.".
You know, in the context above,.
God has been listing out all the detestable,.
sinful things the Canaanite have done..

$^{441}$The Israelites are being warned not to follow them..
Verse 27, "For all these things were done.
"by the people who lived in the land before you,.
"and the land became defiled.".
And if now, if now, if you defile the land,.
the land is like a person, right?.
Vomiting out, just like it vomit out the nations before you..
So you have to keep the covenant..
Chapter 20, 22, "Keep all my decrees and laws,.
"follow them, so that the land where I'm bringing you.
"to live may not vomit you out..
"You must not live according to the custom of nations.
"I'm going to drive out before you,.
"because they did all these things I abhorred them.".
The land, if you wanna live in there,.
if you wanna enjoy the blessing of God,.
if you disobey, you will be vomited out, right?.
If you don't act kindly to other people,.
if you become the oppressors, you will be kicked out, right?.
There's a conditionality there, there is a condition there..
And why?.
Why is there a condition,.
if the land is given to the Israelites?.
Well, I think 25, 23 from the book of Leviticus.
answer this question..
And the answer is, "The land must not be sold permanently,.
"because the land of Canaan," God said, "is mine,.
"and you are..
"But aliens and my tenants,.
"you don't own the land," God said..
I own the land, God owns the land..
And the people of God are being blessed to enter the land,.
but they are called to be faithful tenants in the land..
So let me sum up a little bit.
and give you a little bit more perspective.
on whatever we hear from the modern conflict..
Number one, the land promise was fulfilled.
in the Old Testament, in the time of David and Solomon..
Not one word failed, remember?.
God has been faithful..

$^{481}$Number two, God owns the land..
The people of God do not own the land..
They are simply tenants..
Number three, living in the land is an opportunity.
for the people of God to testify the goodness of God..
They live in the land of Canaan,.
so that they can tell the world, the nations around them,.
how good God is..
Number four, security in the land depends on your loyalty,.
your faith to God's promise..
Number four, land in the Old Testament.
is a foreshadow of life in Christ..
Put it simply, whenever you read the Old Testament,.
by what we did, you find the people of Israel.
living in the land of Canaan..
You sum it up as life in the land,.
but they are only the first theme being introduced..
The theme keeps on developing..
And when we come to Jesus Christ, our Lord,.
He came and He accomplished the salvation..
He obeyed God to the end, and He died on the cross.
for your sin and for my sin..
He was the perfect Israelite,.
the perfect covenant-keeping Israelite..
And He was faithful to the covenant,.
and He did it on behalf of us..
So when you read something in the Old Testament.
about the land, it's life in the land,.
how they should live in the land..
It points to life in Christ..
In the Old Testament, it was life in the land..
In the New Testament, now, today,.
since the land promise is already fulfilled,.
and it's fulfilled in Christ,.
life in the land points to life in Christ..
See, it's the same way..
If we want to have the blessing in Christ,.
if we want to enjoy the joy and peace and love and harmony.
we have or promise from Scripture in Christ,.
we are also being called to be faithful..

$^{521}$We have to exercise our faith.
in order to enjoy all these things..
We can live in the land, live in Christ,.
enjoy His blessing and His presence.
only if we are faithful, right?.
Therefore, in the New Testament,.
it keeps on encouraging us.
that since Christ already fulfilled His covenant,.
fulfilled the salvation to save all of us,.
that's why in Christ we need to continue to follow Him..
We need to continue to be faithful..
Let me close by citing this verse,.
and I want to take this verse seriously..
2 Corinthians 1:20..
For no matter how many promises God has made,.
they are yes in Christ..
So through Christ, the amen is spoken by us.
to the glory of God..
You realize you speak Hebrew?.
We speak Hebrew because whenever we pray,.
we end with amen, right?.
What is that?.
Amen is a Hebrew word, meaning yes,.
meaning definitely, meaning truthfully..
And what we find the truth,.
what we find the amen is in Christ.
because He did the covenant loyalty for all of us.
so that we, all of us,.
if we begin our faith in Christ,.
we are part of God's people,.
but we are being called to trust in Jesus our Lord..
And only when we trust in Him,.
only when we are truly born again in Christ,.
then we can experience the grace of God,.
the blessing of God..
And church, I want to encourage you today..
Since we have such a huge blessing in Christ,.
therefore we press on..
We need to continue to put our faith in Jesus our Lord..
Put our faith in Him and continue to obey.

$^{561}$and trust in Jesus Christ our Lord..
Amen?.
Let's pray together..
Father, thank you for your son, Jesus Christ..
He has done the covenant of faith.
and covenant all the promises we found on the cross..
Thank you for sending your son to die on the cross.
and rose from the dead..
And he's the true covenant keeper..
And thank you for sending your Messiah,.
your Lord to save us,.
save us from sin, from death..
And I pray that as your church,.
as brother and sister in Christ,.
we not only begin in Christ with faith,.
but we'll continue our faith in you..
And help us as we face different challenges.
and different problems and different issues in our lives..
We turn our eyes to you.
and we continue to follow and obey you..
Thank you for being our faithful God..
May you continue to help us to be faithful,.
to serve you, to worship you,.
to follow our Lord Jesus Christ..
in Christ's name we pray..
\newpage



\section{}
\label{sec:giSEqlRniIQ}
\textbf{以巴衝突的反思 (2) - 陳崇基牧師}
\newline
\newline
連結: \href{https://youtube.com/watch?v=giSEqlRniIQ}{\texttt{https://youtube.com/watch?v=giSEqlRniIQ}} ~~~~ 語音日期: 2023-10-31
\newline
\newline
\hyperref[sec:H8JKxV9_GHo]{\small{< < < PREV SERMON < < <}}
~
\hyperref[sec:index_chronic]{\small{[返順時目]}}
~
\hyperref[sec:index_scriptual]{\small{[返順卷目]}}
~
\hyperref[sec:XjRMb_PqUwc]{\small{> > > NEXT SERMON > > >}}
\newline
\newline
$^{1}$大家好,早安.
這些是甚麼來的.
那些是卡嗎.
那些是卡嗎?謝謝….
我以為是現金.
開玩笑的….
好,今天我們繼續看「以巴衝突」.
希望這個DETOUR可以給大家一些角度.
來審視一下.
現在世界在不同地方發生的事.
特別是在以色列巴勒斯坦地區這個地方.
我們這個世界已經夠緊張了.
再加上一些這樣的情況發生.
確實令我們有很多事情要處理.
所以我覺得我們活在世界裡.
我們真的需要學習如何理解.
學習如何判斷世事裡發生的事.
特別是坊間有很多不同的說法.
對於以巴的衝突有不同的說法.
我們應該如何審視.
不僅是以色列和加沙的戰爭.
好像很遠.
其實也很近.
在爆發的第一天.
是10月7日.
哈馬斯恐怖分子.
他們去了以色列的南面.
然後去 kidnap 或者做一些仇殺等等.
然後以色列就還火.
大家在新聞裡面應該看到很多.
應該是已經開始了.
直接開始了.
開始了這個地面衝突.
雙方一定是非常慘痛的.
這個衝突.
無論結果是怎樣.
過程是怎樣.
雙方一定是有很大的損失.
正如我所說這個事情.
不僅是在遠地發生的事.

$^{41}$也可以很近的.
因為在13天之前.
10月16日.
在我們芝加哥附近.
40分鐘車程.
一個地方叫做 Plainfield.
那裡有一位房東.
進入了一個租客的家.
他用刀刺.
用刀刺一個六歲的小孩.
那個小孩是一個巴西人.
因為這個房東.
他看了一些以巴衝突的新聞太多.
我想不知道什麼原因.
他就進了這個房子.
然後就將這個六歲的小孩去刺殺.
插了他26刀.
他當場死了.
當然也傷了他的媽媽.
是一些很近的事情.
世界已經有很多令我們很煩的事情.
還有我們經過疫情幾年.
社會整個世界已經很分裂.
很分裂.
再加上這些情況.
我們不能夠去理解的時候.
我們就會發現更多這樣的事情發生.
當然不僅是這一宗.
我想說另一宗就是10月21日.
10月21日在德州.
一個女拉比.
女拉比.
在德州一個猶太人的會堂.
那個總管.
她在她的家附近又被刺傷死了.
我們不知道是不是仇恨犯罪.
不過這個其實是一個冰山的一角.
有沒有另外一浪的.
叫做反右的主義會出現.
在這個情況下.

$^{81}$這個世界裡面就是有一種.
很不和氣的一種氣氛.
我們需要了解一下.
到底世界上發生了什麼事情.
在我們上一次開始.
我們就說到這個以巴衝突.
有大概五個命題.
這五個命題.
很多時候我們需要再斟酌一下.
特別是有時候引經據典的時候.
引用聖經的時候.
到底是不是引用得對.
還有我們有沒有深思熟慮.
想清楚這些聖經的意思.
我們上一次停在第三個命題.
就是迦南地是上帝無條件地給以色列人的.
我什麼意思.
我知道網上有很多人有不同的意見.
他說為什麼你將一些無條件的認許.
說成好像有條件一樣.
我想你想清楚.
這一段的聖經.
特別很多時候在以巴的關係裡面.
我們都會使用的.
就是《創世記》十二章一至第三節.
耶和華對亞伯蘭說.
你要離開本地本族夫家.
往我所指示你的地方.
我必使你成為大國.
又必賜福給你.
使你的名為大.
你要使別人得福.
為你祝福的 我賜福給他.
咒助你的 我必咒助他.
地上的萬族 都必因你得福.
我們看這段經文的時候.
我們知道這段經文.
我們首先要看文法.
一定要看文法.
雖然我們看的是翻譯本.

$^{121}$但我已經指出過了.
這裡有好幾個重要的動詞.
是叫做coordative.
就是說讓我的意思.
就是勸告式的意思.
上帝當然有應許.
對亞伯蘭有應許.
但那個應許不是anything goes的.
不是說我這樣跟你說了.
然後你隨隨便便生活.
我都會將那個祝福給你.
我都會將那個基業給你.
不.
在聖經裡面所有的約.
我們這個叫做亞伯拉罕的約.
Abrahamic covenant.
所有的約都是有兩個的parties.
兩個方面來合作一起做的.
當然一定是上帝首先主動.
但他講這段話的時候.
他其實就是一種呼召.
呼召亞伯拉罕.
他說有一個可能會發生.
一個potentiality.
這個就是那個勸告式的意思.
有一個可能會發生.
就是我的應許能夠實現在世界上.
怎樣發生.
你就要相信我 是不是.
你看上文下你都知道.
十二章第一節開始.
就是耶和華對亞伯蘭.
或者亞伯拉罕講的話.
十二章四節.
就是亞伯拉罕朝著上帝的話去.
十二章至二十二章.
整個亞伯拉罕的生平.
就是在講亞伯拉罕.
當他相信上帝的計劃.
用上帝的方式來進行的時候.

$^{161}$神的國 神自己在地上可以做的.
就會推展.
如果亞伯拉罕不願意相信的話.
或者有懷疑 有軟弱的時候.
那個計劃就會有detour.
就會可能搞砸那件事.
所以你要留意在聖經裡面所講的弱.
好像婚約一樣.
Marriage Covenant.
是有兩個的parties.
一個Marriage Covenant.
一個婚約.
是不是應該沒有條件的.
In some sense 是沒有條件的.
Right.
因為大家是對大家毀身.
是不是.
不是說你做30 我就做70.
你做40 我就做60.
In some sense 是沒有條件的.
但in another sense 是有條件的.
是大家需要一起盡力去維繫.
Your Covenant 就是這樣的.
整本的聖經都是一樣的.
這樣告訴我們.
當你看到整本聖經的發展的時候.
你就發現.
原來是需要神的子民去相信 去信服.
他們才能夠得到上帝自己的認許.
另一個出自《聖經》的經文.
我們很多時候說.
這個所謂Land Promise.
這個土地 迦南地的認許.
其實在《聖經》裡面從來沒有應驗過.
在哪裡應驗呢.
我們就說在今時今日.
現代的以色列國.
Modern State of Israel.
那裡應驗出來.
事實是否如此呢.

$^{201}$我們就需要學習看《聖經》.
從舊約看到新約.
你知道整本《聖經》.
我們都不是第一次說.
從頭到尾是一個大的敘事.
大的故事.
一個大故事.
告訴我們整個世界.
上帝在做什麼.
我們叫它舊因的歷史.
或者舊屬的歷史.
Redemptive History.
從頭到尾 從創世記到啟示錄.
是一個大的敘事.
你要知道在這個敘事的時間裡面.
不斷有事情發生.
還有不斷有事情發展.
是有所演變的.
連上帝對世人的約都會有演變.
直到基督來實現所有事情為止.
這個土地的認許.
有什麼值得我們留意呢.
你看完亞伯拉罕.
上帝對他所說的話.
到以撒和雅各.
然後雅各就有十二支派出來.
然後以色列成為一個大的家族.
最後他們成為一個大的國.
他們到埃及避難.
然後摩西被呼召.
去拯救以色列民出埃及.
然後他們出埃及之後.
在曠野走了四十年.
因為他們不願意信服.
所以他們進不了迦林地.
拿不到認許.
他們沒有信仰.
所以他們進不了去.
連摩西最後都進不了去.
到最後是第二代.

$^{241}$在曠野出生的第二代.
因為約書亞帶他們.
他們才能夠過約旦河.
進入認許之地.
所以我們由創出黎文新五經之後.
是哪本書呢.
就是約書亞記.
約書亞就是帶以色列人過紅海.
進入迦林地.
他們打完仗.
然後入主迦林之後.
約書亞記說了什麼.
二十一章四十三節.
他說.
(約書亞記).
因為耶和華將所有的仇敵.
都交在他們手中.
(約書亞記).
有沒有看清楚.
一句一話.
其實原文的意思是.
Not even one word.
Not even one thing.
Not fulfilled.
當他們入主迦林的時候.
沒有一句上帝所認許的話.
是沒有落空的.
全部所有的事都已經應驗了.
然後他怕我們不太明白.
二十三章.
約書亞臨離開這個世界的時候.
他跟以色列人說什麼呢.
(約書亞記).
因為耶和華是認許的.
一切的福氣.
怎樣臨到你們身上.
耶和華亦都.
現在你有福氣了.
你住在迦林地.
但是耶和華也必照樣.

$^{281}$將各樣的災禍臨到你們.
直等到他們從耶和華神所賜你們的.
美地上除滅.
在說什麼.
因為每一個約.
近東的條約.
一定最後的時候.
有祝福的說話.
和有奏奏的說話.
你信服的話.
你就得到這個百姓.
如果你違背上帝的約.
你就會招致滅亡.
這個所有的近東的約.
都會是這樣提出條款的.
你們若違背耶和華.
在耶和華所吩咐你們所守的約.
去事奉別臣的話.
敬拜他們的話.
耶和華的怒氣就必向你們發作.
使你們在他所賜你們的美地上.
迅速滅亡.
你明白嗎.
進到地裡面.
不是說任何東西都可以.
或者給地契.
給古代的以色列人.
以至他們進到樂園.
他們就隨隨便便可以生活.
不是這樣的.
是有祝福和奏奏.
是視乎你信服.
還是你是背弱的.
你記得為什麼.
迦南地.
上帝要趕走本身住迦南地的迦南人.
是因為他們的罪惡滿盈.
上帝是跟他們講條件.
當他們的罪惡滿盈.
亞摩利亞人 迦南人罪惡滿盈的時候.

$^{321}$上帝就將他們趕出去.
現在以色列人進來了.
同樣的條件和條件.
有祝福 也有罪惡.
聖經對土地的認許.
在《約書亞記》已經很清楚.
所以為什麼有些人說.
《約書亞記》其實是第六經.
因為我們有五經.
Pentateuch.
其實應該不要只看五經.
因為五經看完之後.
土地的認許還沒有實現.
所以要看完第六本 第六經.
所以頭六本書可以叫做Hexatew.
就是六經的意思.
看完《約書亞記》你就很清楚.
沒有一句話落空.
上帝已經賜了土地給以色列人.
就看他們能否守住這個地方.
繼續看下去.
王國的時代.
當他們設立了王國的時候.
就是獵王機上第四章.
這裡是所羅門執政的時候.
大衛所羅門已經到了.
設立了他們聯合的王國.
就說猶大人和以色列人.
如同海山那樣多吃喝快樂.
所羅門就統治諸國.
從大河 又花底河.
直到非利士 直到埃及的邊界.
所羅門在世的時候.
這些國都向他進貢.
你繼續看下去.
所羅門的國度.
耶和華怎麼說呢.
所羅門的祈禱.
耶和華說是應當稱頌的.
是照著他一切所應許的.

$^{361}$給了那他的百姓以色列人飯.
藉著他僕人摩西所賜福的話.
又是這句話.
一句沒有落空.
你住在這裡了.
上帝已經給了你一個認許.
一句沒有落空.
但是你怎麼能夠留在這個地方.
這是一個必須的聲明.
是要看你能夠留多久.
所以在摩西的律法裡面.
經常都去告訴以色列人.
說他們進到這個地裡.
他們需要做些什麼呢.
就是要守律法.
守律法的其中一種應用.
就是你怎麼去待其他人.
愛人如己等等.
不要虧待寄居的.
不要去欺壓他們.
因為你以前在埃及裡面試過寄居.
現在你住在迦南地.
就不要這樣做.
利未基十八章.
連迦南地也沾污了.
這裡說的是什麼.
原本的迦南人.
上帝說我懲罰那地的罪孽.
地就套出他的居民來.
二十六節.
但你們就是當代的以色列人.
要遵守我的律例典章.
這一切可憎的事.
之前他說了很多可憎的事.
就是迦南人的風俗.
無論本地的寄居的.
在你中間的外人都不可以做.
在你之前居住那地的人.
做了這一切憎惡的事.
地就玷污了.

$^{401}$不要讓地因你們玷污了他套出來.
將套出你們之前的迦南國民一樣.
你明白嗎.
不要把地玷污.
以免地好像以人法一樣.
土地 迦南地會將他們套出來.
其實更好應該是吐.
吐出來.
因為他們太過罪惡滿盈.
太過為耶和華所憎惡.
所以耶和華要將他們套出來.
照樣以色列人進入迦南地的時候.
他們都要遵守上帝的律例典章.
免得他們被套出來.
你還沒記20章.
22節.
你們要謹守我一切律例典章.
遵行他們.
免得我令你們去住的那地.
把你們嘔吐出來.
我在你們面前所逐出的國民.
你們不可以隨從他們的風俗.
不要行他們一切所行的事.
所以是有條件的.
我們不要當聖經只是看《創世記》那幾節.
然後我們就說這是永恆的地契.
不.
在背景裡面你要明白什麼叫做規範.
什麼叫做約.
所以為什麼在摩西五經裡面說什麼呢.
說到賣地的時候.
耶和華吩咐什麼呢.
耶和華吩咐說地不可以賣團.
就是說地不可以永久地賣.
因為沒有一個人是永久佔有那個地方.
因為真正的地主是耶和華自己.
不是住在裡面的人.
你們在我面前是客裡 是寄居的.
你永遠都是寄居.
永遠都是客裡.

$^{441}$你們不是地主.
誰是地主呢.
耶和華才是地主.
神才是土地的主人.
所有住在這個地方的.
就是在地主的條件下.
因為他是租戶.
租戶忠信的話就可留.
租戶不忠信的話就要走.
在迦林地能夠在上帝的同在裡面.
就是需要守上帝的約.
我這裡很快地說.
舊約的土地.
你繼續去追溯這個主題的時候.
由摩西五經 約書亞等等.
去到新約的時候.
你就發現這個土地的認許.
其實是指向一個現實.
就是我在基督裡面的現實.
說得簡單一點.
舊約裡面我們每次看這些經文.
說到Live in the land的時候.
其實它就是一個路牌指向.
Live in Christ.
在基督裡面.
因為基督是那位真正的以色列人.
他是馴服至底的.
並且是滿足了所有約裡面的條件.
所以上帝就將國賜給他.
所有信奉基督的人.
在耶穌基督裡面的人.
我們就得到這個認許.
不是因為我們厲害.
而是因為上帝的兒子在十字架上.
為我們成就了救恩.
是他馴服至底 他完全馴服.
他遵守了這個約.
試探的時候 他沒有犯罪.
然後他馴服至底至死.
死在十字架上.

$^{481}$去聖咒 實現真正的約.
這個約 我們在基督裡面.
才能夠有恩典來分享.
在舊約裡面Live in the land.
就等於Live in Christ.
最後 今天的經文.
我願意將這段經文帶走.
如果我們要在上帝的同在裡面.
能夠享受上帝給我們的祝福的話.
我們要遵守.
我們必須有信心.
對吧.
我們說要得到這個整個永生的認許.
不是任何東西都可以的.
不是免洗金牌.
而是我們在基督裡面.
我們需要有信心.
需要去篤信耶穌基督.
多倫多後書一章二十節.
神的認許不論有多少.
在基督裡面都是是的.
Amen.
你有沒有realize.
你懂得希伯來文.
因為Amen.
阿們.
是希伯來文.
是Hebrew.
意思就是說是.
Yes.
在基督裡面都是Amen.
我們藉著祂說阿們.
使上帝恩我們.
得永.
今天無論你生命裡面.
經歷任何的事情.
或者你關係裡面.
生活的狀況.
你想有上帝的祝福.
你必須有信心.

$^{521}$不是讓你免洗金牌.
我們需要有信心.
去信耶穌基督.
去遵行祂的心意.
先遵行祂的心意去信服祂.
我們才能夠領受上帝給我們的得勝.
願意我們每一個都是這樣.
我們同心大道基督.
主我們求你.
讓我們在這個世局當中.
我們有不同的混亂.
不同的聲音.
讓我們重返到你自己的華語當中.
我們能夠得利.
主為了教會.
為了弟兄姊妹.
我們在這個世間裡面.
我們需要作分辨的.
給我們智慧.
在我們生命當中.
我們需要去領受.
你給我們的祝福的時候.
主我們求你都給我們信心.
幫助我們的不信.
幫助我們的軟弱.
讓我們篤信我們的主.
讓我們守你的約.
遵行你的約.
以致我們的生命能夠經歷上帝.
你對我們的得勝.
我們求主你去祝福這裡每一位.
祝福我們每一位弟兄姊妹.
祝福我們每一個的家庭.
每一個的生命.
以致我們滿載著上帝你自己的美善.
祈禱奉主名.
\newpage



\section{}
\label{sec:XjRMb_PqUwc}
\textbf{Israel-Palestine Conflict - A Reflection (2) - Rev. Mark Chan}
\newline
\newline
連結: \href{https://youtube.com/watch?v=XjRMb-PqUwc}{\texttt{https://youtube.com/watch?v=XjRMb-PqUwc}} ~~~~ 語音日期: 2023-11-07
\newline
\newline
\hyperref[sec:giSEqlRniIQ]{\small{< < < PREV SERMON < < <}}
~
\hyperref[sec:index_chronic]{\small{[返順時目]}}
~
\hyperref[sec:index_scriptual]{\small{[返順卷目]}}
~
\hyperref[sec:g4cKcCqr1bs]{\small{> > > NEXT SERMON > > >}}
\newline
\newline
$^{1}$Good morning, guys..
Why don't we say hi to each other?.
Just wave your hand..
Good to be here..
Whose bottle is this?.
Oh, it's some living water, yeah..
Okay, today we are going to continue.
with the reflection on what is going on.
in Israel-Palestine..
I'm not going to talk so much about.
what is going on right there on the ground..
It's tragedy to both sides..
No matter what happened, ground invasion.
or spill out into a regional conflict,.
it's going to be bad..
What I would like to do is, in this conflict,.
there were a lot of different arguments.
and a lot of different discussions about the land..
So I would like to walk you through,.
like what I did last week,.
walk you through the entire Bible.
from the first book, Genesis,.
to the last book, Revelation,.
and what is going on with the theme of land..
Definitely in the Old Testament,.
there were a lot of discussions.
and a lot of emphasis on land,.
land of Canaan, which is between the Mediterranean Sea.
and the desert on the right,.
that narrow corridor..
In this conflict currently happening,.
people often come to the debates and the arguments,.
whose land is this and who should stay in this land?.
And it can be very divisive..
But I want to walk you through.
from Old Testament to New Testament,.
what is going on here..
Sometimes we miss the big picture..
The Bible, when the Bible talks about the land of Canaan,.
the promised land, it actually talks about grace..

$^{41}$It actually talks about God's gifting the land.
to the Israelites, to the Old Testament people..
And the land is a sheer, a sheer gift from God..
And the people of God get to enter.
and to enjoy the land all because of God's grace..
Of course, there was a covenant between God and the people,.
and the people are to be faithful..
People are to follow the Torah,.
the principles lay out in the Old Testament,.
especially the first five books, how to manage their land..
So when the Bible talks about God gifting the land.
to the Old Testament people,.
it is not saying, okay, anything goes..
You know, I just give you this piece of property.
and you can just go in and do whatever you want, right?.
And by the way, we talk about the people of God being chosen,.
right, chosen people of God..
The word chosen or the concept of election, chosen,.
does not mean the people being chosen have privileges,.
does not mean the people being chosen are superior..
Chosen or being elected means God is choosing you.
to be his instrument to bless other people, right?.
So it's the land, the land is a gift..
And the Israelites are called to manage it well.
and to give back and to love God..
So there are many different principles in the Torah,.
and I wanna give you a sample of it..
One of the sample that talks about.
how people should behave in the land,.
in the land gifted by God..
Deuteronomy 24, it says,.
"When you beat the olives from your olive trees," okay,.
"don't beat it twice..
"The remaining olives are for the foreigners.
"or the orphans, the widows..
"When you gather the grapes in your vineyard,.
"do not clean the vines after they are picked..
"Leave the remaining grapes for the foreigners,.
"orphans, and widows..
"Remember that you were slaves in the land of Egypt..

$^{81}$"That is why I am giving you this command.".
Now, this is very interesting..
When you beat the olives, what does it mean, right?.
Have you seen the olive tree?.
Olive tree, no?.
Okay, you've seen olives, right?.
In whole foods..
Right, but olive tree, the olive tree,.
it's like an umbrella,.
and when you harvest olive tree,.
the olives are on the tree, on the branches,.
and how you do it?.
You have a blanket under the tree,.
and then you use a wood stick,.
and then you beat the branches..
When you beat the branches, the mature olives will fall,.
and then you will gather it underneath.
with a blanket or with baskets..
So that's what it's talking about..
When you beat the olive tree,.
when you harvest the olive tree,.
don't harvest to the end every single piece of olives..
Leave it there..
Why?.
Because after you harvest,.
then the poor and the needy,.
they will come to your field,.
and they can have their summit field..
You see what it's saying?.
Okay, vineyard as well..
You know, have you seen vineyard?.
Grapes, have you seen grapes?.
Whole food, right?.
Yeah, but vineyard, we're talking about vineyard, right?.
The branches, right?.
When the grapes mature, you're gonna clean them, right?.
You can take them down, cut them down..
But what it says here is that you have to have leftover.
remaining in the vineyard, on the vines..
Please, why?.

$^{121}$Because the needy will come by,.
and they can enjoy the fruit..
That is what land management or farm management.
looks like in the Old Testament..
And usually those laws,.
this kind of law, similar principles,.
are all over the place in the Old Testament..
And usually they are followed by one clause,.
what we call the motive clause..
Motive clause, verse 22..
It gives you the motive..
So the Torah is not just like,.
you know, I give you this principle, this commandment,.
and you have to observe it,.
so you can earn my favor, right?.
No..
It gives you the motive clause..
Remember that you were once slaves in the land of Egypt..
Remember you were saved by grace..
Remember I saved you by grace from Pharaoh, from Egypt..
Now you settled down in the land..
You should not imitate the oppressors, right?.
You should leave behind some resources.
so you can bless other people, right?.
You can be generous..
So the promise of,.
they get to enter the land.
in order the land will be filled with,.
you know, resources, can bless other people..
Again and again, the Bible tells us in the Old Testament,.
the land promise was already fulfilled..
Not one word has failed..
Everyone was fulfilled..
Not one thing has fallen..
Everything God promised to the Israelites were fulfilled..
That's why they can settle down in the land,.
and then they build a kingdom, right?.
But of course, that's not how the story goes, right?.
In the Old Testament, when the people of God,.
they build their own kingdom, then they were divided,.

$^{161}$and they continue to oppress other people..
And until 586 BC,.
God sent the Babylonians to destroy Jerusalem.
and the country of Judah,.
because they violate, they violate the Torah..
So they were in exile..
But then in exile, they came back, right?.
Under Ezra and Nehemiah,.
you can also call that second Exodus..
And other Exodus, they came back from the exile land,.
back to Jerusalem, back to their homeland..
But more importantly, when they came back to the land,.
they did not treat it as we conquest..
They're not conquering the land again..
When they came back to the land,.
they understand it is a regifting..
God is regifting them a new, generous,.
a gracious opportunity to serve God..
So you remember those guys, Ezra and Nehemiah, right?.
Came back, Ezra as a scribe,.
and teaching the people the Torah,.
and lead them into a spiritual revival, right?.
And then there was big, long prayer in Ezra chapter nine..
And then Nehemiah had prayed the same prayer,.
or similar prayer in chapter nine of the book of Nehemiah..
They all talk about the same thing, and with the same tone..
And let's listen to their tone..
In the midst of spiritual revival, back in the land,.
being regifted..
Oh my God, I'm ashamed, utterly ashamed,.
for our sins are power higher than our heads..
But now we have been given a brief moment of grace,.
for the Lord our God has allowed a few of us.
to survive as remnant..
We are the remaining people,.
because we have abandoned your commands..
There's no demand of God, giving them the land again..
There's no demand of God that says,.
no, this land belong to us, because we were here first..
There's no sense of entitlement,.

$^{201}$but only a sense of repentance, right?.
Knowing that God is now gracious enough to regift them,.
so that they can serve God in the land..
And then the Old Testament story continues into,.
between the Old and New Testament,.
there were three to 400 years,.
and the people of God were still yearning.
for the real return..
Because yeah, physically they are returned, right?.
They have returned to the land physically,.
but spiritually, politically,.
they were still under different empires, right?.
Persian empires, Greek empires, and then Roman empire..
They did not have their freedom..
And the people were looking forward,.
someday God will fulfill his promise again,.
fulfill his promise in the Messiah..
So the second theme we have to think about.
is that the entire Old Testament.
is actually pointing forward to a goal..
And that goal is our Lord Jesus Christ..
In Jesus, we find that all the promises of the land.
and other promises are fulfilled..
Yes, the people lost the land because of a sin..
After exile, they were still up in the air..
But not because God failed, God is faithful..
The people failed, the other party of the covenant,.
they failed, they did not keep the covenant..
Now this has to do with Jesus..
Now think about how the New Testament,.
the first book of New Testament introduced Jesus..
What is the first book of the New Testament?.
Gospel of Matthew, right, Matthew..
The genealogies, right, from Abraham.
all the way down to Christ, and then the birth of Christ,.
and then living under an other pharaoh,.
Held the Great, and he was persecuted..
And then the Holy Family, Joseph, Mary, and Joseph,.
they were kind of exiled from the land.
and they went down to Egypt, which is on the left-hand side..

$^{241}$And then after the pharaoh, the new pharaoh,.
Held the Great died, they returned back to the land, right?.
So there's exile and return in the life of Jesus..
So as you can see in the Bible,.
already in the first book of the New Testament,.
the Gospel of Matthew, they look at the life story of Jesus,.
his life, his mission, his ministry,.
his death and resurrection, in light of the story of Israel..
Because in Jesus, no matter what happened to him,.
no matter what he did and what he said,.
they find that Jesus is the true Israelites.
who keeps the covenant perfectly..
He fulfilled all the promise and secure all the promises.
for those who are in Christ..
You are in Christ, I am in Christ..
So in that sense, we share that promise,.
which is included, which the land is included..
So in some sense, as we trace the entire Old Testament.
to New Testament, we know that, so listen up,.
Old Testament, the Old Testament is incomplete.
without the New Testament, right?.
If you read the entire Old Testament,.
which probably you haven't, right?.
But if you read them, all the books, 39 books,.
you come to the end, it's an open-end book..
What is going on, right?.
We're still yearning for another episode, right?.
In some sense, you can imagine the Old Testament.
and New Testament as a unit, as an entire grand story.
that has a lot of episode..
The Old Testament is episode one to episode 98..
After 98, you're still yearning for something, right?.
The hope, the promise is unfulfilled..
And then episode 99 is New Testament, right?.
Jesus came and he fulfilled all the promises..
He secured the fulfillment..
Where's episode number 100?.
It's the new heaven, new earth, right?.
That we're all going in Christ..
In that sense, Old Testament and New Testament are a unit..

$^{281}$So when you trace the land promise, the theme of land,.
you come to Jesus and Jesus fulfilled principally,.
fundamentally, he fulfilled that promise of inheritance,.
the promise of land..
So in the New Testament, we are not looking forward.
to any point or any piece of property on earth..
You cannot find it with a GPS..
You find it in Jesus Christ, the living person.
who loved you and died for you and rose from the dead..
Let me take a look at, fast forward to the book of Hebrews..
Hebrews chapter 11..
Remember this chapter is a list of faithful heroes, right?.
In the Old Testament..
So-and-so is so faithful and they obtained the promise..
And now come to 1110, Abraham, talking about Abraham..
Abraham is the first story, the forefather of Israel..
He was confidently looking forward to a city.
with eternal foundations, a city designed and built by God..
So what was Abraham looking for?.
He was probably in episode three, right?.
Abraham in the Old Testament..
He was looking forward or hoping for not a point.
on the map, not a physical point on the map,.
but looking forward to the heavenly city..
And the New Testament is clear..
It's clear the new city is in episode 100..
Someday in Christ, we will see that new city.
coming down from heaven, the new Jerusalem to earth..
And that is our hope..
The Bible is clear, the fulfillment of land promise..
Indeed, every promise of God is found none other.
than Jesus himself, Jesus Christ himself..
Ephesians chapter one..
It talks about how blessed we are in Christ..
How many promises and how many blessings.
we receive in Christ..
Chapter one, seven..
In Christ or in God, we have redemption through his blood..
The forgiveness of our trespasses,.
according to the riches of his grace,.

$^{321}$which he lavished upon us..
God is so resourceful and he is so willing.
to lavish on us blessing..
I remember when my kids were younger,.
of course, young kids, they like candy, right?.
They like candy..
I used to buy them some candy, right?.
And then they were so happy..
When I buy them a little box of candy.
and they were overjoyed, right?.
Little did they know that my credit card.
can buy millions of candies..
And I can shower upon their heads, right?.
That is what is happening here in Ephesians one..
God is God of blessing..
And in Jesus, we have so much riches of his grace,.
which he lavished upon us in all wisdom and insight..
Verse nine, making known to us the mystery of his will.
according to his purpose, which he set forth in Christ.
as a plan for the fullness of time.
to unite all things in Jesus,.
things in heaven and things on earth..
Everything in the universe, everything in the cosmos.
is united together, found the fulfillment in Jesus Christ..
Verse 11, in him we have obtained an inheritance, the land..
And what is the land?.
Episode 100, new heaven and new earth..
We have obtained the inheritance, the property in heaven,.
in new heaven, new earth,.
having been predestined according to the purposes of him.
who works all things according to the counsel of his will.
so that we who are the first to hope in Christ.
might be to the praise of his glory..
This is episode 99, and telling us someday we will inherit..
The inheritance, the land in episode 100..
I know most of us live in Chicago, right?.
We live in Illinois, right?.
You know the most expensive house in Illinois?.
Where is it?.
The most expensive house in Illinois, in the state of Illinois..

$^{361}$Where is it?.
Let me give you a clue, in Chicago..
It is in Chicago..
The most expensive house in the state of Illinois.
is in Chicago, in the city of Chicago, right?.
Have you seen it?.
Which neighborhood?.
Lincoln Park..
Anyone from Lincoln Park?.
The most expensive house in Illinois.
is located in Chicago, Lincoln Park..
The price for the house is 50 million..
50 million, imagine that, one house, right?.
25,000 square feet, 11 bathrooms..
Praise God, 11 bathrooms, right?.
Eight city lots, eight city lots, right?.
How many houses do you have?.
One house is one lot, right?.
Oh, no, no, no, one lot has so many houses, right?.
Eight city lots, big house, huge house..
Imagine this most expensive house is God's inheritance..
I'm using a very finite analogy, right?.
Imagine this is the inheritance that we have in Christ..
Now, you don't get to live in that house yet, right?.
We're still in episode 99, but the key is given to you, right?.
If the key of the house is given to you,.
how you treat, where you gonna put your key?.
Where you put the house key?.
A 50 million dollar house, the key is already given to you..
Someday you're gonna get into this mansion.
and you can enjoy the inheritance..
Where you gonna keep the key?.
I don't know..
I don't know, somewhere in the safe, right?.
Would you give it to your dog?.
Probably not, right?.
Hopefully not..
That is the new heaven, new of the inheritance.
that Apostle Paul is telling us..
In Christ, we have that gift..

$^{401}$We have that inheritance, secure and fulfilled.
in Jesus Christ..
That's why we are being gifted so much..
And at the same time, if we understand.
that we have experienced grace of God, right?.
The lavish grace of God..
It should motivate us to be generous people, right?.
We are so generous..
We can be so generous because we are so blessed..
I know Thanksgiving is around the corner.
and I've been thinking about grace,.
the goodness of God, the blessing of God,.
and our responses..
I wanna close today with a parable, a parable by Jesus..
And usually when we hear this parable,.
when we study this parable,.
we don't read this parable with the perspective.
of inheritance, the father's house, right?.
The household, the big household..
Usually we read this parable understanding the love of God..
Yeah, definitely..
But I wanna go through this parable with you.
and with this perspective..
Remember the parable of prodigal son?.
The prodigal son?.
The younger son who is rebellious, right?.
And he took his inheritance and go to far country.
and waste them all, right?.
And he was suffering and he just thought on him,.
hey, maybe I should return to the father's house.
and rather be a servant, right?.
That's a parable of prodigal son..
But actually if you read it carefully,.
the parable of prodigal son is parable of two sons..
Two sons, the younger and the older, right?.
The younger and the firstborn..
The younger son, as I said, goes away.
and spend all of his inheritance from the father.
and he suffered so much..
And then the son realized the patience.

$^{441}$of the father is enduring, right?.
So he decided to just go back to the father's house..
He hurts his father deeply and he knows.
he does not deserve the favor of God..
So he made up his mind and say,.
no, I'm not worthy to be called my father's son,.
but I would rather go back to be one of his servants..
So he goes back and he just wanna work for his father.
as a slave working for the master..
But when the son returns home, what did he find?.
He does not see a stern master,.
but a loving and patient father waiting for him..
The loving father is waiting for him..
He saw him and felt compassion and ran.
and embraced and kissed him..
Now, this sort of actions, usually in that kind of culture,.
in biblical culture, means humiliation..
Old men don't run..
They walk slowly, right?.
And respectfully, keeping their honor..
But this father seeing the younger son coming back,.
he ran and embraced him and kissed him.
and threw a party for him and celebrate..
So the son returns home,.
returned to his father's land and inheritance..
Now, at this point, the firstborn,.
the elder son came back, right?.
And how did he react?.
He comes home from the farm.
and he heard about the celebration.
for the long lost younger brother, but he was angry..
He said, "It's not fair.".
Why?.
Because he has been working so hard, right,.
to earn his father's favor..
And he was angry against his father.
and the younger son, his brother,.
because he has been working so hard.
to earn the favor of God..
And the father never threw a party for him..

$^{481}$He, the elder son, is a dutiful son..
He's a very faithful son, but he never realized, sadly,.
he is already living in and enjoying.
all the inheritance of the father,.
every single day since he was born..
He would enjoy it and he would inherit.
all the land and inheritance..
You see here, the two sons,.
one find that God's favor is undeserved..
The other one is trying, working so hard.
in order to earn that, right?.
And the parable ends with, you know,.
you don't have to do that because God is so rich.
and God has blessed you and give you the inheritance..
Now, usually theologians would define grace.
as unmerited favor..
What does it mean, unmerited favor?.
It means that you don't deserve the grace of God,.
the love of God..
Nonetheless, God loves you..
God loves you even you don't deserve it,.
even you feel you don't deserve it..
God still loves you..
That is grace, right?.
No debate about it, unmerited favor..
You don't deserve it, unmerited, but it comes upon you..
But this is only one side of the coin..
I will say grace is also dismerited favor..
Now, what does it mean, dismerited favor?.
I don't know how to translate in Chinese,.
so yeah, don't worry about it..
It's called dismerited favor..
What it means is that you deserve judgment..
We sinners, we deserve condemnation,.
yet condemnation does not come to us..
God forgives..
God forgives us even we do deserve condemnation..
That is dismerited favor..
It's not just unmerited favor, it's dismerited favor..
This ought to be the motivation.

$^{521}$for our entire Christian life, right?.
You don't have to work so hard to earn God's favor,.
but it has given to you unmerited and dismerited..
Once we experience the grace of God,.
our relationship with God should be motivated.
by a grateful response..
Since God loves me so much, and I'm not worthy,.
I'm not dismerited, yet God's love comes upon my life,.
and I want to respond to God with gratitude,.
with worship, with my service..
Yesterday, I read something really,.
I think it's meaningful for this season, Thanksgiving..
That short piece is about gratefulness and gratitude..
It's a reminder for me,.
and I want to close my sermon with this..
There's a fine line between doing things for God.
with a strong sense of responsibility,.
and serving God with gratitude, gratefulness..
So that short piece talks about this..
Let me run it through..
Grateful for, house to clean..
You can take it as, I have to clean my house..
Such a burden for me, right?.
Have you cleaned your house lately?.
Such a burden..
But you can also reflect it on the other side..
Yeah, I can clean my house because I'm grateful..
God has given me a safe shelter to live..
Doing laundry, right?.
It sucks doing laundry..
I don't do it, my wife does..
It sucks..
But consider, we are doing laundry.
because God has been so blessing us with clothes to wear..
Doing dishes, last night I was doing dishes,.
a lot of dishes, because my son came back.
to spend the night..
Yeah, there were so many dishes, right?.
I do dishes, it sucks..
You know, it come from a spirit of just being responsible..

$^{561}$The elder son, by the way, the elder son..
I'm working so hard, right?.
But gratitude says, we have food to eat, right?.
That's why we do dishes..
Grocery shopping, right?.
It's routine, it's tedious..
But God has provided me money, it's proficient..
I have to clean toilets..
But thank God for indoor plumbing..
There were a lot of noise in my life,.
but thank God, we have people around us.
who care about us..
There are endless questions about homework,.
for the kids especially, right?.
Endless assignments they have to do..
Oh, so tired..
But gratefulness says, hey, my kids are growing, right?.
Early wake-ups, just like this morning, right?.
Early wake-ups..
Or Monday, right?.
Anyone has Monday blue?.
I do, right?.
Oh, what a day, we have to drag our body,.
you know, away from the mattress, right?.
To wake up, early wake-up..
But hey, if you're taking care of the young kids,.
children to love, praise God..
Throwing tire in bed, it sucks..
But thank God I'm still alive..
And by the way, ties an offering, right?.
One-tenth of our income goes into offering..
Um, it feels burdensome..
Yeah, I must say that, you know,.
I'm a Christian for many, many years..
I feel like ties and offering, yeah, right?.
Yeah, it would be heaven if there was no ties and offering..
However, gratefulness says,.
I do have the resources to bless other people, right?.
Last week we had the General Assembly.
and usually at the end of the year,.

$^{601}$all the churches and most of the organizations.
are lacking some resources..
We do, our church does, yeah..
It's challenging this year, right?.
Yesterday I was dining out in a restaurant with someone.
and then I paid the bill..
\$30 each person..
\$30 a meal, right?.
That's crazy, I didn't have enough..
I was still hungry, leaving the restaurant, right?.
So I paid \$60, right?.
The biggest struggle I have these days,.
dining out, is the tips..
Can I cut the tips, right?.
Life is hard..
But gratitude says,.
we have the resources to bless other people..
God is too good..
God is too resourceful..
God of blessing..
Let's all pray..
Father, thank you for being our God, our infinite God..
And you are the God of blessings..
You are the God full of resources..
And you are infinite..
And your love never exhausts..
Your patience towards us never exhausts..
Thank you for being our great God..
Thank you for the fact that in Christ,.
we enjoy the unmerited,.
dismerited favor..
That we are in Christ, no more condemnation..
And in Christ, we have new life..
Thank you for all the good things.
that you have showered upon us..
Especially as we end our service with communion,.
we pick up the bread and the cup of grape juice.
to remember how great you are.
and your love is so great..
Beyond our comprehension,.

$^{641}$yet becomes a reality in our life..
Encourage us and give us hope..
Give us this new strength to live for you,.
to serve you..
In Christ, we pray..
\newpage



\section{}
\label{sec:g4cKcCqr1bs}
\textbf{以巴衝突的反思 (3) - 陳崇基牧師}
\newline
\newline
連結: \href{https://youtube.com/watch?v=g4cKcCqr1bs}{\texttt{https://youtube.com/watch?v=g4cKcCqr1bs}} ~~~~ 語音日期: 2023-11-07
\newline
\newline
\hyperref[sec:XjRMb_PqUwc]{\small{< < < PREV SERMON < < <}}
~
\hyperref[sec:index_chronic]{\small{[返順時目]}}
~
\hyperref[sec:index_scriptual]{\small{[返順卷目]}}
~
\hyperref[sec:_21dlkrR_Vw]{\small{> > > NEXT SERMON > > >}}
\newline
\newline
$^{1}$我們今天繼續在以巴衝突的主題裡面.
我們思想到土地的問題.
在舊約裡面說了很多.
Land of Canaan.
Promised Land.
應許之地.
是什麼意思呢.
我們在過去的兩次.
我們都是處理著五個命題.
我們還在第三個命題.
第三個命題就是我們很多時候.
認為迦林地就是在舊約聖經裡面.
上帝寫包丹給以色列人.
只要上帝有這個地契.
他們就可以進入這個地方.
上一次我們解釋了大概一半左右.
其實在舊約的聖經裡面告訴我們.
土地是上帝和他的子民之間一個約.
所以不是寫包丹或者隨隨便便.
而是有恩典賜給土地.
給以色列人.
他們進入這個土地生活的時候.
在上帝的同在下生活.
他們就需要守約.
所以在律法書裡面說了很多不同的律例.
是說到他們得到這個地的時候.
不是說我們霸佔這個地.
而是我們得到這個恩典.
我們去學習回報.
所以第一個主題我們需要想清楚.
就是Land is a gift.
土地 迦林地.
是對上帝的子民是一個恩典.
能夠進入這個土地.
就能夠享受這個土地裡面的土產.
享受了之後.
接受了這個恩典的時候就需要去回報.
當然有很多這樣的律例都告訴以色列人.
他們在土地裡面應該怎樣生活.
或者Live in the land.

$^{41}$應該怎樣去生活.
怎樣去他們的操守行為等等.
其中一段我們可以去看一下.
就是《新民記》二十四章.
這裡說到他們收割的時候.
就是收割橄欖樹.
橄欖樹.
橄欖樹收割大概是每年的九月十月的時候.
當他們打橄欖樹的時候.
為什麼說打橄欖樹呢.
因為橄欖樹是用一支棍子去打.
然後橄欖樹會掉下來.
下面就鋪一塊布或者有個籃.
然後你就會收到那些橄欖.
打橄欖樹的時候.
枝上面所剩下的就不可再打了.
為什麼呢.
因為打剩的當然會有一點點.
所以耶和華就說不要打完.
就是說有錢不要賺完.
有錢不要賺到盡.
剩下一點點在樹枝上面.
要留給那些軟弱有需要的人.
孤兒寡婦寄居等等.
然後二十一節.
你們摘葡萄的時候.
摘葡萄.
我們這裡也有這樣的活動.
每逢秋季的時候.
我們就去摘一下南瓜.
不是摘南瓜.
摘一下蘋果等等.
南瓜太大了.
摘一下不同的東西.
摘一下一些水果等等.
這裡同樣是在說這件事.
同一個季度.
在秋季的時候.
葡萄樹要結果.
然後要去到收割.

$^{81}$掉下來的不可以收取.
就是說那些葡萄熟了.
掉在地上.
你就不要再收拾.
有錢不要賺到盡.
剩下給誰呢.
剩下給寄居孤兒和寡婦.
很多時候律法後面也有一句東西.
就是動機的句子.
Motive class.
就是告訴你為什麼要這樣做.
很多時候我們看律法的時候.
覺得很多事情要做.
很辛苦.
很多工作和責任等等.
但不是的.
他先說律法.
很多時候後面他一定會再提醒你.
提醒你什麼呢.
提醒你恩典.
你要記得.
因為你曾經在埃及地作過奴僕.
所以我就吩咐你進行這事.
你見不見到那個Motive告訴你.
你要記得曾經在埃及.
上帝拯救過你出來.
你已經得到恩典.
能夠進入這個地.
所以你應該.
不要像以前欺凌你埃及的主人一樣.
法老.
你今時今日進入迦牟尼的時候.
就在恩典的土地當中去學習.
感恩去回饋.
剩下的東西給貧苦大眾.
慷慨的來施瑜.
打橄欖打兩下就好了.
你已經有足夠.
其他東西可以奉獻給人.
去施瑜給人.

$^{121}$所以在舊約的律法裡面.
經常提到以色列人.
恩典的土地.
感恩的生活.
這個土地的認許.
已經在舊約裡面.
是實現了.
一句也不落空.
全都應閻了.
上帝對我們多好.
上帝對以色列人.
我曾經承諾過你的東西.
全部都在迦南地.
你們入主迦南的時候.
就已經實現了.
一句也不落空.
當然我們知道.
整個舊約的故事.
我們一直在說.
都是在追溯整個舊約新約的故事.
是一個大的敘事.
一個大故事.
大故事是怎樣去發展和演進.
特別是我們追溯一個土地的主題的時候.
到了他們進入迦南地的時候.
一句也不落空.
之後發生什麼事.
才提醒以色列人.
你們不可以這樣.
因為他們在欺壓其他人.
所以一直到罪惡滿盈的時候.
最後主前六世紀的時候.
就是巴比倫軍打過來.
在耶和華差派巴比倫帝國來.
就是說這一班人.
他們進入了恩典的土地.
居然不懂得去感恩.
去生活.
所以他們就被擄到巴比倫.
土地就沒有了.

$^{161}$然後你記得故事發生什麼事.
在他們被擄的當中.
上帝說我再恩待你.
然後帶你回來.
於是我們就去到索羅巴巴.
以色列和尼希米的時代.
他們少數的人.
可以回到這個應許之地.
但是留意回到應許之地之後.
他們的土地就沒有了.
但是留意回到應許之地的時候.
氣氛是怎麼樣的呢.
這裡是以色列第九章.
或者你可以看尼希米第九章.
差不多.
一個就是聖殿已經重建了.
一個小的聖殿重建了.
然後以色列就作一個文字.
去教導以色列人.
律法書.
然後他們就再一次.
有什麼.
這個屬靈的復興.
然後尼希米記第九章.
就說到城牆已經建好了.
他們再一次有一個聚集.
又是一個很長的禱告.
以色列第九章.
或者尼希米第九章.
那個禱告的.
感覺其實都是一樣.
說什麼呢.
他說我們的上帝我們抱愧蒙羞.
我們的罪孽.
多到滅頂.
多到滅頂是什麼意思.
我們的罪孽多到我們頭上.
整個人都是罪.
多到滅頂.
現在耶和華我們的神.

$^{201}$暫且向我們示恩.
回來這個認祂之地.
為我們留下一些殘存之民.
我們去到.
我們曾經離棄過.
你自己的律法.
你看到那個感覺是什麼.
尼希米.
回到認祂之地的時候.
不是再征服.
不是reconquest.
不是再霸佔那個地方.
而是他們了解.
這個上帝是恩典.
我們不配.
我們暫且可以存留.
re-gifting.
不是re-conquest.
回到這裡.
然後舊約聖經.
埋尾的時候.
你就發現他們好像回去了.
但又沒有回去.
他們還是生存在帝國的欺壓下.
你trace這個story.
一直到兩約中間.
然後到最後就是新約.
新約來到發生什麼事.
你打開新約.
第一本書是什麼事.
馬太福音.
馬太方.
去介紹耶穌的時候.
怎麼介紹.
加普.
十四代.
然後基督的出生.
然後到智利朝拜.
然後應驗.
舊約先知的盼望.

$^{241}$然後耶穌就被.
大希律追殺.
然後他就.
所謂被擄.
下了埃及.
然後大希律死了.
這個好像法老的王.
死了之後他又再出埃及.
然後回到以色列.
被擄又回歸.
你看到他介紹耶穌的Jesus story.
其實就是在說.
以色列的故事.
整個舊約的以色列的故事.
就是在耶穌基督自己身上.
有的故事.
所以去到新約.
耶穌來到的時候.
新約的作者說什麼.
我們見到耶穌的故事的時候.
我們發現原來見到真正的.
以色列人的故事.
並且他是一個守約的.
以色列人的故事.
耶穌出來傳道四十天.
在曠野受試探.
是不是.
正如以色列人四十年在曠野受試探.
他在約旦河受洗.
正如以色列人.
經過紅海.
經過約旦河入應他之地.
整個的Jesus story的輪廓.
就是以色列.
故事舊約故事的輪廓.
所以我們做基督徒.
我們不要只是.
做新約的基督徒.
我們經常看來看去.
都是看新約聖經.

$^{281}$我們要看舊約聖經.
因為舊約.
是不圓滿的.
看完舊約聖經.
總是覺得有一點點.
有點點的不滿.
總是有些事情還沒搞定.
到底誰來幫我們.
守那個約.
以致我們能夠進入真正的應許.
沒有人.
在整個舊約裡面.
亡國被老回歸也是沒有人.
直到新約耶穌來到的時候.
所以舊約裡面.
常常有一個像箭咀一樣.
指著前面.
要等一個來.
那個彌賽亞那個救主.
他來到他就會實現所有的應許.
所以舊約和新約.
是整本聖經.
我其實很擔心.
很多時候我們.
看聖經我們就.
只看三分之一.
有三分之二我們不太理會.
舊約聖經就是三分之二.
舊約和新約是一個.
是一個合一.
你可以想像整本聖經.
由舊約創世記.
到新約最後那本書啟示錄.
其實是一個長篇劇.
舊約就是第一集.
到第九十八集.
你看完第九十八集.
你總是覺得是不夠的.
新約來到了.
就是第九十八集.

$^{321}$新約來到了就是第九十九集.
你說第一百集是什麼?.
第一百集就是.
基督在來的時候.
新天新地要到.
所以第九十九集告訴你.
將來會有第一百集.
將來會完結整個長篇劇.
這個就是舊約和新約.
整個的合一.
所以你看到.
到新約的時候你就會看到.
耶穌的身上就是傳過以色列的故事.
耶穌的身上就看到.
有回歸進入真正的應許之地.
真正的應許之地.
不是在地圖上.
那一格.
那一格長的以色列.
地中海沙漠中間.
那個地方.
真正的土地的應許.
是新天新地第一百集.
告訴你的.
啟示錄告訴你的.
所以我們繼續看下去.
在新約裡面有什麼說呢?.
在希伯來書第十一章.
這裡說到阿巴拉罕.
你知道第十一章.
經常說一些所謂.
List of Heroes.
Heroes of Faith.
信心的偉人.
一個又一個這樣去列出來.
說到阿巴拉罕的時候.
是什麼呢?.
他說阿巴拉罕.
是等候著那有根基的城.
就是上帝所設計和建造的城.

$^{361}$阿巴拉罕所仰望的.
不是這個以色列巴勒斯坦地這麼小的.
是新的城.
天上.
降下來.
新天新地.
這一個城.
我希望你看到舊約和新約是一個的Unity.
而且在基督裡面.
所有的應許.
都Fulfill.
無論你看馬達福音.
很多句說話.
怎樣應許.
先知應許的.
全部在基督身上.
來應驗.
Fulfill.
整本的新約都是這樣說.
在基督裡面我們有的.
就是全部一切的基.
第九十九集.
一百集裡面所說.
所以第三個主題.
我們會看到新約的聖經裡面.
經常提到我們作為基督的人.
我們是多麼的有福.
多麼有多上帝給我們的恩福.
這是已忽所書第一章.
已忽所書第一章一開播.
他就說我們是多麼的多Blessing.
我們在基督裡面有多多的祝福.
他說在基督裡面.
在愛子的血裡面.
得蒙救贖.
過範得以赦免.
這是照著上帝豐富的恩典.
充充足足的賞賜給我們.
在基督裡面有珠般的智慧和聰明.
照自己在基督裡面所訂立的美意.

$^{401}$上帝使我們知道他的旨意的奧妙.
按照所安排的時機成熟的時候.
使天上地上一切所有的.
在基督裡面都同歸於一.
在他裡面能夠找到Unity.
在基督裡面找到Unity.
十一節我們也在基督裡面得到基業.
Inheritance.
Inheritance是什麼意思.
Inheritance就是在說土地.
但是土地不是這一格巴士山地.
是將來新天新地.
這原是那隨己而行萬事的上帝.
照著他的旨意所預定的.
要使我們在基督裡面.
首先把希望託付在基督裡面的人.
頌讚榮耀歸給他.
這個基業就是將來要有的新天新地.
New Heaven and New Earth.
大家在座的我們都住芝加哥.
你知道在Inheritance的這個州裡面.
最貴的那間屋在哪裡.
那些做房地產的.
最貴的那間屋在哪裡.
最貴的那間屋在芝加哥.
在Lincoln Park.
聽說這間Lincoln Park的豪宅.
是值五千萬.
50 million 有機會你開車去看.
二萬五千呎.
只是那間屋二萬五千呎.
11個廁所.
OK吧.
11個廁所.
不用爭了.
8個City Lock.
你住幾個房子.
一個.
他們8個.
8個City Lock.

$^{441}$50個million的大屋.
你可以想像一下.
如果上帝的基業.
用這個所謂有限的比喻.
上帝說在基督裡面有基業給你.
就是這間屋.
50個million.
好不好.
打你地水打死你.
但是他說如果這個基業.
所有地水都給了.
所有的債務都付了.
現在我就給你一條鑰匙.
你會怎樣生活.
你會很想進這間屋.
你會不會將鑰匙到處丟.
掛在狗那裡.
你會保存住這條鑰匙.
因為將來這個基業.
已經是屬於你的.
以法所述第一章最後.
他說這個基業有憑據.
有個銷.
有個deposit.
那個deposit就是聖靈.
聖靈住在每一個信耶穌的人的生命裡.
讓我們大能大力能夠盼望上帝.
這個就是在基督裡面.
我們有的福分.
多大的福分.
恩典有多大.
原來我們就仿似住在上帝自己的基業裡面.
並且我們還在等待第一百集.
這個真正的基業的來臨.
因為十一月是感恩節.
所以我就想起一些和恩典有關的.
我昨天去預備的時候就想起.
耶穌講過一個比喻.
那個比喻我們聽的時候.
很多時候我們不會想起將來盼望基業這件事.

$^{481}$我們很少會想起.
那個比喻是什麼.
那個比喻是浪子的比喻.
大家知道浪子的比喻是什麼.
有一個老爸有兩個兒子.
其實不是浪子的比喻.
是兩個兒子的比喻.
你現在有一個小兒子是一個反斗性的.
敗家子.
他搶了上父家的基業.
然後他就去到外地放蕩.
花盡了他自己所有的基業.
然後他突然間醒悟什麼.
我為什麼要在這裡挨窮.
吃豬鮑.
我不如回到父家那裡.
雖然我已經不配做他的兒子.
不如我回到父家那裡做他的奴僕也好.
做他的傭人也好.
所以他毅然就回到家裡.
然後遠遠看到他家父在等他的時候.
他就跪在他家父面前.
他說我不配做你的兒子.
我做你的奴僕好了.
你記得家父怎樣去回應他嗎.
記得嗎.
馬上幫他換衣服.
戴戒指在他身上.
然後宰一隻肥牛.
然後給他一個大的派對.
因為這個失而復得的兒子終於回來了.
終於回來了.
回到這個父家裡面.
然後第二個兒子出現.
記得嗎.
大兒子出現.
整個比喻的重點是大兒子.
大兒子從哪裡回來.
從田裡回來.
即是父家很多的田.

$^{521}$做公他回來的時候.
他就聽到有一個大的派對.
居然是一個反骨仔.
回來的時候.
居然為他開一個派對.
這個大兒子有什麼投訴.
他說我從小到大都聽你的話.
為你做了這麼多事.
我為你做了這麼多事.
為什麼你沒有給我一個派對.
家父跟他說.
他說你一直住在我的地業裡面.
其實所有的東西.
一切的事.
一切的物業.
都已經是屬於你的.
這個大兒子還沒有搞清楚.
他原來做了這麼多年.
他一直沒有去享受.
一直沒有去享受到原來的家業.
家業是我的.
上帝的恩典原來不是我賺回來的.
是他賦予的.
由他出生那一刻.
他就在這個家業裡面.
所以有神學家說什麼呢.
他說恩典是什麼.
他說恩典當然就是不配得的好處.
Unmerited Favor.
Unmerited.
不是我們能夠賺取回來的.
救恩不是我們賺取回來的.
在基督裡面有的基業有的福分.
都不是我們賺取回來的.
上帝愛你.
縱使你不應得.
縱使你感覺到你不應得.
上帝都給過你.
這個叫恩典.
但是留意這個恩典的定義.

$^{561}$只不過是一個硬幣的其中一邊.
其實恩典不只是Unmerited Favor.
恩典不只是不應配得的好處.
恩典更加是Dismerited Favor.
我中文不知道怎麼翻譯.
什麼叫Dismerited Favor.
意思是什麼呢.
就是說我們應得的是什麼.
我們應得的是審判.
是被停職.
但是我們應得的沒有臨到我們.
因為上帝愛你.
上帝愛我們每一個人.
這個才是真正的恩典.
不只是我們不配得的臨到我們.
硬幣的另一邊就是.
我們配得的卻沒有臨到我們.
Condemnation.
在罪裡面我們是被定案的.
不義的.
沒有臨到我們.
臨到基督的身上.
這個應該是我們基督徒生命裡面.
應該有的Motive.
我們應該有這樣的動機.
我們不是死做爛做.
然後我們想換上帝給我們的恩典.
No,這個是Unmerited.
And Dismerited.
不是靠我們的工作.
而是上帝願意去施恩給我們.
正因為上帝願意施恩給我們.
我們就應該用恩典來回報.
我們會服侍他.
我們會去遵上帝的命.
是樂意的.
因為我們是受了恩的人.
浪子比喻.
大兒子就是.
不知道他是不是真的得到那個恩典.

$^{601}$家父告訴他.
一早已經給了你.
小兒子回歸.
他也得到恩典.
是Unmerited.
And Dismerited.
感恩節的時候.
我心靈很艱苦.
疫情之後.
我知道大家都很艱苦.
昨天我和朋友出去吃飯.
結帳的時候.
嘩,30元一位.
以前多少?.
以前10元都不用.
20元最多已經很好.
嘩,30元怎麼搞的.
我吃得不是很飽.
這樣才合.
我現在最掙扎的.
就是貼屁股.
生活真的艱苦.
物價很高漲.
我們有很多憂慮.
你看看你每一個月的結帳.
真的不同了.
因為不同了的緣故.
所以有很多無形的壓力.
到年尾的時候.
教會.
我真的不只是我們教會.
我也收到很多不同的信.
有關宣教工場.
或者基督教機構.
呼音機構,教會等等.
那個奉獻.
或者運作上.
是遇見有難處.
包括我們UV.
我不想每次都像是奉獻呼籲.

$^{641}$但我相信我們是可以負擔得到的.
如果我們每一個都願意奉獻.
或者應該這樣說.
如果我們每一個都是受恩.
上帝給我們如此多祝福的時候.
我們更加可以慷慨去奉獻.
昨天我看到一個短的文章.
那裡說什麼呢.
他說感恩是什麼意思.
有什麼叫為了責任而做事.
什麼叫因為受恩而回應.
我覺得很有意思.
不過說什麼呢.
他說今天又要洗房間.
又要清潔.
撞鬼而已.
有這麼多責任我要做.
但感恩的心就不同了.
感恩就是我今天可以清潔.
因為我有安全的居住.
要洗衣服.
但感恩我有衣服可穿.
要洗碗.
昨天我就洗碗了.
因為我老婆煮了很多東西.
因為我兒子剛剛回來一晚.
很多碗仔.
撞鬼而已.
但感恩為什麼我可以洗碗.
因為有溫飽.
可以去超市買菜.
很辛苦吧.
但感恩我有錢可以購物.
要洗廁所.
是不是很辛苦.
當然.
但感恩有噴霧.
身邊有很多噪音.
但感恩身邊有很多關心我的人.
每天回去我要幫小孩做家課.

$^{681}$沒完沒了的家課問題.
是出於責任.
但感恩.
因為我的小孩.
他的腦袋開始成長.
一早起床我要去照顧子女.
很辛苦.
但感恩.
值得我愛的孩子.
我當然要照顧.
睡醒的時候身體很累.
很累.
我沒精打采.
我不想面對今天.
出於責任.
但感恩.
我今天還有氣息.
十一奉獻也是一樣.
為什麼我要做這件事.
出乎責任這麼辛苦.
但如果反過來看.
感恩.
感恩我能夠用我的資源.
去祝福其他人.
上帝對我們的祝福是多麼大.
不應有的好處.
祂給了我們.
我們應得的受罰.
祂不給我們.
這才叫做感恩.
土地是感恩.
新約在基督有的生命.
甚至將來新天新地.
都是一場的恩典.
只有恩典.
求主幫助我們.
我今天聖餐結束的時候.
要我去提醒.
不是為了.
我必須要去做.

$^{721}$我已經被賜予了.
我很高興去做.
我誠心帶著祈禱.
主我多謝你.
多謝你每一天給我們.
充充足具的恩典.
雖然在這個世界裡有混亂.
在生命裡有很多挑戰.
有很多難處.
可能我們要節衣足食.
可能我們要節省.
看到生命裡有經濟的壓力等等.
或者看到我們身體有軟弱.
但上主我們來到你面前.
我們知道你是豐豐足足像恩典的神.
你是滿有不成的上帝.
主我們就是實實在在.
在基督裡面有這個特性.
上主我們來到你面前.
我們知道我們是滿有恩典的生命.
也知道你是.
癡下萬有恩典的神.
在你裡面沒有缺點.
主讓我們感恩.
慷慨去服侍你.
以致你的名能夠織住我們.
得到榮耀.
主我們拿起這個聖餐.
這個紙和餅的時候.
為我們成就將來天上新天新地的基業.
感恩 讚美你.
祈禱奉主耶穌的名字.
\newpage



\section{}
\label{sec:_21dlkrR_Vw}
\textbf{Do You Even Lift, Bruh? - Brian Ng}
\newline
\newline
連結: \href{https://youtube.com/watch?v=-21dlkrR_Vw}{\texttt{https://youtube.com/watch?v=-21dlkrR\_Vw}} ~~~~ 語音日期: 2023-11-14
\newline
\newline
\hyperref[sec:g4cKcCqr1bs]{\small{< < < PREV SERMON < < <}}
~
\hyperref[sec:index_chronic]{\small{[返順時目]}}
~
\hyperref[sec:index_scriptual]{\small{[返順卷目]}}
~
\hyperref[sec:z0bvgIPiO7c]{\small{> > > NEXT SERMON > > >}}
\newline
\newline
$^{1}$- Well, good morning, everyone..
It's nice to see everybody up here, from up here again..
And I know that Rebecca, my wife, was here just with you.
all last month for a couple of times, and twice,.
and in an art workshop..
And the first message she gave was--.
(audience applauding).
Was on the topic of the power of storytelling.
and Jesus being the goat..
And her second message was diving into.
one of the parables of the wheat and tare..
And also, she facilitated an art workshop,.
and I really appreciated some of the artwork.
that she had photos of..
And my favorite was one where someone made an ice cream,.
given that's my favorite thing right now..
Time really flies, because the last time I was up here,.
I realized when I was preparing for this message.
was actually last October..
So it's already been a little over a year..
And for that message,.
I actually prepared one called "Lightning Our Yolks.".
If some of you are here, you remember that in Acts 15,.
I talked about how we sometimes put on additional.
individual and corporate yolks on ourselves..
And these are the additional spiritual restrictions.
that we prescribe ourselves,.
that kind of hinder our spiritual faith journeys..
But instead, Jesus prescribes to us.
that he is lightening our yolks,.
and that he will give us greater joy and freedom,.
because we can know that we have these lighter loads.
that we carry, and he's helping us carry them..
At the end of the message, I mentioned that.
even though we're afforded greater freedom and greater joy,.
there's still a yolk that we have to carry..
A lighter one, but we're still expected to maintain.
proper conduct, behavior, and live our lives.
that are, in a way, glorifying to him..
Carrying this yolk, doing this work,.

$^{41}$still requires a degree of focus, intention, and discipline..
I really appreciate coming here,.
because Pastor Mark always gives me freedom to choose.
whatever topic I wanna preach up here..
And it wasn't my intention to actually tie in.
what I talk about today with the message.
that I prepared before, but it kinda ties in together..
So I really kinda like that it worked out this way..
I've been praying about what message to give,.
and the message today is what I got,.
and it happens to still connect with what I did last time..
So anything we do, especially focused work,.
requires a level of preparation, training,.
and right mindset..
Our message today comes from 1 Corinthians 9, 24 to 27..
And the sermon title, as you guys saw,.
is Do You Even Lift, Brah?.
The passage today is from 1 Corinthians 9, 24 to 27,.
and I'll read it..
I know it's a little bit small for you all,.
but I'll read it..
Do you not know that in a race, all the runners run,.
but only one gets the prize?.
Run in such a way as to get the prize..
Everyone who competes in the games.
goes into strict training..
They do it to get a crown that will not last,.
but we do it to get a crown that will last forever..
Therefore, I do not run like someone running aimlessly..
I do not fight like a boxer beating the air..
No, I strike a blow to my body and make it my slave.
so that after I preach to others,.
I myself will not be disqualified for the prize..
In Rebecca's presentation and messages,.
she prepares very beautiful slides..
I'm sorry, mine will disappoint.
because they're not nearly as elegant..
Actually, they're the opposite..
They're very, very crude..
But I lean more kind of your common pop culture.

$^{81}$cliche visual tools..
So hopefully some of you will appreciate..
And I know there's a broad age and generational range here..
Some probably as young as high school or college.
and those that might even be retired..
And so I know that my title requires.
a little bit of explanation for those especially.
that don't know what it actually means..
The phrase, do you even lift brah?.
Is actually a mocking question used for people.
who attempt to look or act capable in something.
but actually aren't..
Case in point is someone like this guy, right?.
I think that he's presenting this persona of a lifter,.
very athletic, but I think some of us question.
whether or not he has the ability..
It's often used against people who want to give,.
this phrase is often used against people.
who give the impression that they're really strong.
and athletic, but in fact, they look very foolish.
and we kind of quickly see through their kind of facade..
Because those who actually engage in something like lifting.
or athletic training know that it requires a lot of focus.
and a lot of preparation and training..
I don't know the closest Mandarin translation.
to this particular phrase, but I know that in Cantonese,.
there used to be something that a lot of people use.
in the movies, in Hong Kong movies..
It's like (speaking in foreign language).
and so that question is, do you have the qualifications?.
And actually it kind of points out to someone like,.
you kind of try to look the part of something,.
but do you actually have the backing, the credentials.
to support the persona that you're presenting?.
When I mention these ideas, I'm sure some of you have.
some memories that come to mind of your own experiences..
In my case, I remember years ago.
when I was still going to the gym..
I don't go to the gym anymore, I have two kids.
and I have a puppy now, so all I do is just take care.

$^{121}$of them and that's sufficient enough exercise for me..
But I remember that there was a time.
when I was going to the gym and this was at Export.
on State Street and I was there for an hour with my friend.
and I was lifting weights and there was a guy that came in.
and he had full track suit, he had a headband,.
he brought his gym bag and he was ready to look like,.
well he looked like he was ready to lift,.
but in fact he sat down at the bench press.
and he pulled out a whole chicken dinner out of his bag.
with a cup of soda and he just proceeded.
to eat dinner at the gym..
For the whole hour I was there,.
he didn't even do one exercise,.
he was just sitting and enjoying his meal..
So that for me is like the perfect example.
of when you would use the question, do you even lift, bruh?.
Because this guy is coming in and showing his athleticism,.
but really there isn't very much of it..
In our passage, Paul uses the visual image of running.
and this is a metaphor, an image that connects the idea.
with our own spiritual journeys, our spiritual walks..
He says, do you not know that in a race all runners run,.
but only one gets the prize?.
Run in such a way as to get the prize..
He's calling us to keep our faith muscles.
well conditioned and activated for the Lord's work.
and for our faith journeys..
Paul highlights that our faith journey.
is demonstrated through action.
and not just a cognitive process that happens in our minds..
Living out our faith is more than the mental belief.
that God exists in identifying Jesus as our savior..
It requires a manifestation of our faith.
through our actions..
We're not just thinking about what we know as truth,.
we're living it out..
It's not a basic matter of belief.
through intellectual insight,.
but it's an actively engaging process.

$^{161}$and Paul, in our passage, uses this metaphor.
to connect this for his original audience..
But it also really kind of ties into what we understand.
in our modern times as well..
Paul's audience is the Corinthian church.
and the believers are in Corinth,.
which was a major Greek commercial city in his time..
The idea of running a race was actually very familiar.
to them because they had a long history.
of holding what were known as the Isthmian Games..
And these were the games that I think.
that our Olympics are actually modeled after..
These games happened every two years.
and athletic as well as musical competitors.
from all of Greece would descend upon Corinth.
and compete in honor of Poseidon..
This happened since like 582 BC,.
so by the time Paul was using this example.
of running a race, people were already familiar.
with this for centuries..
Paul takes an idea that the Corinthians understand.
and uses it to help them visualize.
a right perspective of Christian living..
Paul turns around the idea as well.
of this big competition being in honor of Poseidon.
and instead attributes what we do.
for God's honor and glory and for his service..
Throughout his letter to the Corinthians,.
Paul addresses a laundry list of specific matters.
that range from addressing disorderly church conduct.
and leadership to divisions,.
to sexual immorality in the church..
He addresses things like lawsuits among believers,.
divorce, sexual immorality..
And I mentioned before that Corinth.
was this important city of commerce in Corinth..
And to make a comparison, it's not unlike Hong Kong.
or New York or London..
It was a major port city where there was a lot of trade..
And with trade in these places,.

$^{201}$it wasn't just the trade of physical goods,.
but trade of physical ideas and culture.
and beliefs as well..
And with any big city, we know there's big city problems..
And Corinth was definitely experiencing this as well..
The church was, by comparison to the Jewish churches.
in kind of the Israeli regions,.
not as grounded in moralistic traditions.
and that culture that they had for a very long time..
And so with a younger church like the Corinthian church,.
you have leaders and followers that are.
just kind of understanding the outside world.
and are not as mature..
And they're bringing in the ideas that they understand.
from their usual conditioning into the church..
And so Paul is writing the church to address.
the deficiencies in their theology.
and their practices..
Paul addressed topics that are very much practical..
And he stares away from the idea, again,.
of a passive idea of our faith.
and just understanding it as a concept..
His belief was demonstrated by him..
Well, like looking at Abraham,.
our original kind of forefather of the faith,.
we know that God credited his faith as righteousness to him..
But this wasn't just him believing..
It was backed by action..
And we know that Abraham left everything that he knew..
He left his family, his home,.
and he went to an unknown place.
where God even called him to sacrifice his son..
And he was obedient..
So he didn't just believe..
He demonstrated his belief through action..
And what I'm saying is that genuine faith.
is exhibited through consistent action..
And this is a way that we serve and glorify God..
The race that Paul refers to.
is his race to fulfill the mission of God,.

$^{241}$to proclaim the gospel to as many people as he can.
and to save as many as he can..
And I'm not trying to present the idea.
that, oh, we're trying to earn our way to salvation..
We are justified as soon as we believe..
This is a demonstration of what real faith is.
because it's consistent with how we act..
It's not just what we say..
I think the example of loving someone.
is a very helpful one in showing this..
Authentic love is reinforced by how you express it..
So if I tell my children that I love them,.
but I don't spend time with them to do the things they like,.
I don't support them, I don't encourage them,.
I can't say that my love is really genuine..
In the same way, thinking that just belief in God is enough.
is not enough..
We have to demonstrate it through action.
because even Satan and the demons believe in God.
and they fear him..
If we're looking at the act of running.
as a metaphor for our faith journey,.
there are specific spiritual practices in our lives.
that then represent our spiritual muscles, so to speak..
These are our actions like our prayer life,.
our devotions, our grounding ourselves in God's word..
And like muscles, these areas need to be developed over time.
and conditioned as well so that we can run like an athlete,.
so that we can do the work God calls us to do..
In our physical bodies, when our bodies are not engaged.
in activity and using our muscles, do you know what happens?.
There's something called atrophy that happens..
The muscle starts to weaken and it starts to diminish..
And because the body deems these muscles as non-essential.
'cause we're not using them,.
it just degenerates and fades away..
This is especially visible, I think, in older folks.
where they're just less active and so their muscles.
begin to shrink and they just become smaller..
But this can happen with younger folks too..

$^{281}$If any of you have broken a limb, like an arm or a leg,.
and you put a cast on, and then at the end of your healing.
and you take away the cast, that arm or that leg.
on one side is gonna be smaller than the one.
on the other side because it's atrophied..
And likewise, with our spiritual condition,.
our spiritual practices, if we don't keep engaged.
in things like reading God's word.
and grounding ourselves in prayer and fellowship.
and these type of actions, we also will begin to.
kind of degenerate in our spiritual strength.
and our spiritual muscle..
When we engage in training and keeping ourselves.
sharp and ready, we'll be able to handle times of stress.
and even take advantage of times of opportunity..
There's something called muscle memory..
Some of you might know it if you guys are athletes..
And it happens where when through repeated training.
of something, an action, your body just begins.
to know what to do in specific situations.
where you trained it to do something..
Because your body reverts back to its conditioning.
when you might not be thinking about how to react.
in a certain situation..
This is why athletes and soldiers train.
and train and train..
Because in times of pressure and urgency,.
when we don't have the time to think about a certain action.
we wanna take in a situation, our body naturally,.
intuitively responds in the way that we've trained it..
Likewise, as we engage in prayer and worship.
and grounding ourselves in God's word,.
when there are times of crises,.
we will naturally go to do what we've been training.
ourselves to do..
We don't think maybe, we don't have the ability to think.
'cause we might be so stressed out,.
but we will naturally gravitate toward that type.
of behavior and why training is so important..
And it's not just in times of crises..

$^{321}$It could be training for times of opportunity..
So let's say like you're on the street,.
you're at a restaurant, you meet a friend of a friend.
and they find out you're a Christian..
They want to know more about your faith and what it means..
If you have not been doing the practices that I mentioned,.
you may not have the ability to truly,.
in a well-fashioned way, articulate what the faith.
that you believe in actually means and is..
And so the training is also to take advantage of opportunity..
We need to remember that these spiritual muscles we develop.
serve the purpose of loving our God and loving our neighbors.
for advancing God's mission, his church,.
and to build up our fellow brothers and sisters..
So these muscles that I'm talking about.
aren't for the vanity of our own self-righteousness..
It's not for us to look in the mirror.
and to feel good about ourselves.
or to present ourselves to others..
It's all for a purpose..
And that purpose is to advance God's kingdom.
and to serve him and to build up his church..
We don't engage in action for ourselves.
or for action for action's sake..
So the second part of 1 Corinthians 9, 24,.
Paul says, "Run in such a way as to get the prize.".
There's a particular direction, a trajectory,.
and goal we're supposed to train and aim for.
when we're doing our training and our conditioning.
and our running of the race..
Paul runs with a strong strategic purpose in his mind..
As he says, "I do not run like someone running aimlessly.".
There's an objective to his actions..
The problem is that we can get trapped.
in the cycle of keeping ourselves busy.
and doing very churchy things.
because we feel that that equals making progress..
And like this example of the hamster,.
like he's just running and running and running.
and probably feels like he's going somewhere,.

$^{361}$but he's really going nowhere..
And a lot of organizations today.
and in a lot of work meetings,.
they kind of go on and on.
and you might have meeting after meeting after meeting..
And there are times when I've been in work.
where I feel like it's death by meeting..
Like there's just meeting, meeting, meeting..
And sometimes people feel like just because you're meeting.
and meeting for long periods of time,.
that there's progress..
But that's not true..
I mean, there are probably some of you that have been in,.
I know Pastor Mark has agendaless meetings..
And this could be for work and it could be for church.
and it goes for hours and hours and hours..
At the end of it, you're exhausted and you wonder,.
what did we actually accomplish?.
It doesn't seem like we accomplished anything,.
but some folks feel like by the pure fact.
that we're together, that some progress is made..
And we need to not confuse the flurry of movement.
or busyness as actual progress..
Paul lays out a very specific goal in his running..
His main purpose is to advance the gospel.
to as many as he can..
He discusses this in the section.
right before our focus passage here..
So we're focusing on 1 Corinthians 9, 24 to 27,.
but in verses 19 to 23, I'm not gonna put it up,.
but I'll just read it out loud..
He says, "Though I am free and belong to no one,.
"I have made myself a slave to everyone,.
"to win as many as possible..
"To the Jews, I became like a Jew to win the Jews..
"To those under the law, I became like one under the law,.
"though I myself am not under the law,.
"so as to win those under the law..
"To those not having the law,.
"I became like one not having the law,.

$^{401}$"though I am not free from God's law,.
"but am under Christ's law,.
"so as to win those not having the law..
"To the weak, I became weak to win the weak..
"I have become all things to all people,.
"so that by all possible means, I might save some..
"I do all this for the sake of the gospel.
"that I may share in its blessings.".
Paul is saying that he will do whatever it takes.
to achieve his goal, to proclaim God's word,.
and save as many people as he can..
He's willing to become a slave and serve..
He's willing to become a Jew and take on Jewish traditions..
He's willing to become weak and forgo the privilege.
that he has as a Roman citizen in his particular context..
He's willing to become one apart from the law,.
not having the law, and eat and engage in activities.
with Gentiles, something that the Jews found repulsive..
So he acts in various ways so that he can relate.
to his target audience, and not just with words..
Paul's purpose adheres to what Jesus makes clear.
for all of us as our purpose in our spiritual lives.
in Matthew 28..
And this is, I'm sure, overused,.
but this is the great commission that Jesus leaves us..
He says, "Therefore go and make disciples of all nations,.
"baptizing them in the name of the Father and of the Son.
"and of the Holy Spirit, and teaching them.
"to obey everything I have commanded you..
"And surely I am with you always.
"to the very end of the age.".
So this is a command not just for Paul,.
but actually for all of us..
Paul advocates engaging in action.
and engaging in action with purpose..
Paul advocates this because also, not just with purpose,.
we also need a sense of urgency in our action..
It's not okay to just participate in the race,.
uncaring of whether one actually aims to finish.
or cross the finish line..

$^{441}$He says to run in such a way as to get the prize..
In the ancient races, there was no second or third place..
There can only be one..
And maybe someone can appreciate that movie quote..
But imagine that there's competition.
when you're running the race..
No one waits on you..
You have to push yourself,.
and there's a sense of urgency in all of that..
You have to push your hardest..
And these days we have, and I guess my children aren't here,.
these days we have diminished the need.
for competitive urgency in our society,.
especially for the sake of not hurting children's feelings..
There are awards for everything..
There's participation awards for just someone.
who participates and shows up..
And everybody is crowned a winner..
You might get a medal and a trophy..
And I think that that bleeds into.
how our spiritual activities and lives are..
It seems enough to just show up,.
even when we're not fully involved..
And sadly, I think that breeds spiritual mediocrity..
Regularly attending church and volunteering are empty.
if our daily life does not reflect Christian behavior,.
and there are no efforts to be made.
to advance God's mission..
Without the drive to strive,.
I think people fall into the danger.
of stopping short of completing what their holy duties are..
I think the aspiration of people working hard.
and retiring early is also a very damaging narrative.
for people to follow,.
especially when it bleeds into how we behave..
Even amongst believers,.
there's this aspiration to retire in our 50s.
and maybe in our 60s at the earliest..
And folks dream about time they're gonna spend.
on the golf course or on cruises they're gonna go on,.

$^{481}$vacations they're gonna take..
But this is, in my opinion, a very modern notion..
We have the luxury now to actually consider.
this more leisurely lifestyle in our later stages of life,.
when we couldn't do that before..
But I don't think that's actually a very biblical notion..
And despite modern advances in productivity.
and healthcare and safety in the world,.
there's still many problems where Christians.
are called to answer those issues,.
and where the gospel and Christian work is needed..
Now, I'm not saying that we should maintain.
the same level of intensity throughout our lives.
and how we do God's work,.
but we're still meant to just keep our eyes on the prize,.
to kind of cross that finish line,.
really until the day we die..
I had a friend that used to always say,.
I'll rest when I die, and he was a ministry worker..
And I think that is probably the more biblical thing.
to have in mind when we're actually thinking about.
how we proceed with the later stages of our lives..
So God deserves our first fruit efforts.
in every part of our lives..
And we don't need to officially win to be the winner,.
but we have to be in it to win it..
We have to have a winner's mindset.
where we're aiming for a goal,.
we're pushing our highest, and we actually go for it..
And that only happens with a sense of urgency.
to advance God's kingdom..
Our faith requires action..
That action or these actions require purpose,.
a sense of purpose, frame with purpose,.
and there has to be a level of urgency and immediacy.
to how we approach that work..
Now, this is a picture of the tortoise and the hare,.
and it's a very famous fable that I assume all of you know..
We need to be like the tortoise in the fable here,.
where we're constantly pursuing our goal..

$^{521}$Slow, well, not slow because he's slow,.
but steady wins the race..
Whereas the hare here is actually just relaxing.
because he's relying on his talent and his abilities,.
but he's procrastinating, he's lazy..
We operate in a world now where we have control.
over so many aspects of our lives..
We can control so many things, even from our smartwatches..
And we think that we are the ones in control,.
even of how we handle our kingdom responsibilities.
that the Lord has given us..
Sometimes the problem is that it's not that,.
you know, we're slow in getting to the race.
or slow in the pace of our race..
We haven't even started the race..
And the most blatant experience that I've had.
is when I used to serve as a volunteer.
at a MBA campus ministry..
And I overheard a student speaking to another student,.
and he was telling this other student.
that he's gonna get this very lucrative job,.
and that he's gonna make up to \$10 million,.
and when he does, that's when he's gonna retire early.
and then start serving the Lord..
He thought that it was appropriate to wait.
and allocate time later on in his life to serve the Lord,.
and there was no immediacy,.
and there was no need to start his race anytime soon..
You know, this reminds me of the passage.
in James chapter four, 13 to 15, where James says,.
"Now listen, you who say today or tomorrow.
"we will go to this or that city,.
"spend a year there, carry on business, and make money..
"Why?.
"You don't even know what will happen tomorrow..
"What is your life?.
"You are a mist that appears for a little while.
"and then vanishes..
"Instead, you ought to say, if it is the Lord's will,.
"we will live and do this or that.".

$^{561}$In contrast to the MBA student,.
one of my inspirations is Dr. Stuart Wu..
I don't know if some of you might know him,.
but he's probably one of the most dedicated.
and faithful servants I've ever met in my life..
And I first encountered Dr. Wu.
when I was doing a short-term medical mission.
in Zhongshan, China..
He's a retired surgeon, and he would regularly,.
more than once a year, go out to remote parts of China.
where we would have to track mountainous areas.
to get to people and provide them medical care.
and also to share the gospel..
He did this into his late 80s,.
until he was 87 and then passed away..
And he never made a complaint that I knew of.
when he did this..
And so, I was less than half his age..
I was probably a third of his age when I first met him..
And I didn't have nearly as much stamina.
and as much zeal as he did..
And I wish to be like him,.
because he kind of constantly had this mindset.
of finishing the race and finishing well,.
even until his late 80s..
Nope, that's not it..
I know this last slide says active rest,.
but actually, my focus pivoted toward.
making sure that we have guardrails..
As we run the race and start doing God's work.
to spread his good news and make disciples,.
we need to make sure that we don't take.
an end justifies the means approach..
And that how we behave in the process.
of reaching that goal matters as much.
as actually reaching that goal..
So, even if there are great results.
of people coming to the church, of the church growing,.
it doesn't give you a pass on how you act.
in the way getting there..

$^{601}$The Lord cares as much about that journey.
as he does about the end goal that you're striving for..
You know, the last part of this passage,.
Paul emphasizes that he makes sure.
that he's not disqualified in his attempt.
to obtain the prize by making sure.
that he keeps his body in check..
In verse 27, he says, "No, I strike a blow to my body.
"and make it my slave, so that after I preach to others,.
"I myself will not be disqualified for the prize.".
I believe that he's actually alluding.
to keeping material and physical temptations at bay..
Because earlier in Corinthians,.
he had addressed the audience to flee.
from sexual immorality in chapter six..
And I think that he's alluding to this.
when he's talking about making sure.
that he's not disqualified..
He makes sure that as he's fulfilling the work of God,.
that he's not discredited and doesn't discredit.
the Lord's work by his behavior..
Sadly, in recent years, we've seen a lot of church leaders,.
even in our region, fall to sexual sin..
And this has had severe impact on the churches themselves,.
as well as people coming to the church..
And this is what we want to stay away from..
And it may not be such a blatant sin as sexual immorality..
It could be an arrogant attitude..
It could be poor financial stewardship..
It could be a lack of empathy for the people.
that we're actually trying to target..
In my devotions recently,.
I'm now finishing the book of Deuteronomy..
At the end of Deuteronomy,.
Moses is getting ready to go up to Mount Nebo and to die..
And he will never be able to cross into the promised land..
He gets a view of the promised land,.
but he's not allowed to cross it..
And this happened because he actually.
did not follow God's exact orders..

$^{641}$Even though he is a champion and a hero of the Bible,.
and he saved the Israelites from slavery,.
and took him into, or got him ready.
to go into the promised land,.
the Lord prevents him from going in,.
because instead of whispering to a rock.
to call forth water, he knocks it twice in anger..
And this was enough for the Lord.
to actually prevent him and penalize him.
from going into the promised land..
How we uphold ourselves during our work.
is as important as accomplishing the end result..
It's not only about the end point, but how we get there..
Not just about the finish line,.
but how we run the race itself..
So, I ask all of you, do you even lift, bruh?.
If you have been diligently doing God's work,.
keep at it, and remember that we're in a continuous marathon.
that requires training, diligence, dedication..
But also, the point that I want to highlight,.
and said that I changed,.
was also to make sure that we rest..
We, I think, especially as immigrants,.
or children of immigrants,.
there's this very strong work ethic that we develop..
And sometimes we bring that also into our faith,.
and we don't rest..
And the Lord prescribed the Sabbath for us,.
so that we can make sure that we rest as well,.
and are recharged so we can run even stronger,.
and for God..
Continuously train yourselves,.
keep yourselves sharp and ready,.
and your muscles going, you know,.
reading the word, praying, being in fellowship,.
serving actively in your communities,.
and then keeping yourselves accountable.
while in fellowship with others..
And for those of you who are honest and say,.
no, I don't lift, now's the time to start,.

$^{681}$because there is, it's never too late to get started..
And it's important to begin when you know.
that God's calling you to his great commission..
So, go run the race, and do your part in God's great work,.
and get those muscles going, and well conditioned..
Let me end this in prayer..
Heavenly Father, thank you for this opportunity.
to preach your word, and to understand that our faith.
is one that is demonstrated through action..
And Paul uses the visual of running a race..
I pray that for each one of us,.
you help us to understand the areas in our lives.
where we need to be better conditioned,.
or to start conditioning, and to operate these muscles,.
these spiritual muscles, so that we can be ready.
for your work that you've called us to,.
to advance your kingdom, and to make disciples.
for your kingdom..
In your name we pray, amen..
- Amen..
\newpage



\section{}
\label{sec:z0bvgIPiO7c}
\textbf{認識三一真神 - 陳漢儀姊妹}
\newline
\newline
連結: \href{https://youtube.com/watch?v=z0bvgIPiO7c}{\texttt{https://youtube.com/watch?v=z0bvgIPiO7c}} ~~~~ 語音日期: 2023-11-14
\newline
\newline
\hyperref[sec:_21dlkrR_Vw]{\small{< < < PREV SERMON < < <}}
~
\hyperref[sec:index_chronic]{\small{[返順時目]}}
~
\hyperref[sec:index_scriptual]{\small{[返順卷目]}}
~
\hyperref[sec:__GpGaclvJI]{\small{> > > NEXT SERMON > > >}}
\newline
\newline
$^{1}$各位頂音姐妹 早安.
(早安).
很感恩今天有機會跟大家分享神的話語.
首先我想問大家一個問題.
當我們被問到.
你們基督教所信的神.
和其他宗教的神有什麼不同的時候.
你會怎樣回答?.
我們可能會回答一個很不同的.
就是我們所信的神.
是三位一體的神.
那你怎樣解釋三位一體?.
他問你什麼是三位一體?.
大家知不知道我當時是怎樣回答的?.
我跟他們說了一句很有智慧的話.
就是.
什麼呢?.
這個是神的奧秘.
That's it! OK! 收工!.
所以為了避免以後我又有這麼有智慧的回答.
今天就跟大家分享認識三一真神.
三位一體是基督教裡.
從聖經總結出來的一個教義.
聖經中是沒有這個字眼的.
它是來自拉丁文Trinitas.
意思是三.
三位一體是一位神.
祂有三個位格.
剛才大家應該都答對了.
聖父,聖子,聖靈.
祂是三位一體的神.
聖父是神.
聖子是神.
聖靈是神.
每一個位格都是完全的神.
不混亂.
然而只有一位神.
這個教義是從我們信耶穌的時候就已經接觸到的.
但是在基督教信仰兩千年的歷史當中.
信徒甚至神學家都很難去完全明白和解釋出來.

$^{41}$為什麼呢?.
因為神是零.
是無限的.
祂是超越我們的思想.
超越我們的情感.
我們的理性和邏輯.
而我們人是很有限的.
受著時間和空間的限制.
試問我們又如何可以用我們這麼有限的心思意念和邏輯.
去完全明白這一位無限的神呢?.
所以對於三位一體的教義.
我們的理解是有困難的.
我們可能會嘗試用一些類比.
一些比喻.
但是我們要非常小心.
因為有些類比會將我們引導到一個異端.
比如什麼呢?.
我們會說三位一體會不會像一個人.
他有三個角色.
他又是父親.
又是丈夫.
又是別人的同事.
我們可能會這樣去理解.
但是這個比喻會將我們引導到一個異端.
這個異端叫什麼呢?.
就是形態論.
他們認為父子聖靈是神自我顯現的三個角色.
就好像剛才那個比喻一樣.
他的錯誤在於什麼呢?.
就是在於他否定了每一個位格都是完全的神.
這一條.
這個就是不符合聖經的.
既然我們去理解這位三位真神.
是這麼有困難.
和這些類比又存在一些危險性的話.
那是不是就是沒有辦法去理解呢?.
不是的.
雖然我們不能夠完全明白.
但是神他仍然願意將自己去啟示給我們聽.
在聖經當中去告訴我們什麼是三位一體.

$^{81}$現在我們就去看一看.
在聖經當中神是怎麼去介紹自己是三位一體的神.
首先我們去解釋一下.
什麼是位格.
位格英文是person.
是什麼意思呢?.
一個人.
何謂一個人呢?.
你怎麼知道你是一個人呢?.
你怎麼知道我是一個人呢?.
不是一本書.
不是一隻貓.
一個人他是有理智.
有情感.
有意志的.
是不是?.
我們有的.
我們知道自己存在.
知道自己在做什麼.
貓貓狗狗不會知道自己存在.
他們看到鏡子會嚇到.
有些狗甚至會吠牠.
哎呀 這個是誰啊.
還有一件事就是.
作為人我們是知道.
怎麼去跟其他人建立關係的.
神就有三個位格.
他就是有三個person.
他們各自有自己的理智.
情感.
意志.
他們知道各自的存在.
他們也可以彼此建立關係.
事不宜遲我們來看舊約.
首先我們說舊約.
神介紹自己.
是一個獨一的神.
舊約可以說是為我們奠定了一個很好的一神論.
一神觀.
我們要多謝猶太人.

$^{121}$因為在古近東的世界.
那些人是拜很多神的.
滿天神佛.
但是唯有猶太人.
他是信奉一神.
在《申命記》.
六章四到五節說.
以色列你要聽耶和華.
我們神是獨一的主.
你要盡心盡性盡力.
愛耶和華你的神.
舊約強調的是一位神.
獨一神.
有沒有體現到三位一體呢.
這個教義.
我認為是沒有的.
舊約很明顯是獨一神的教義.
他是反對多神的.
所以如果我跟當時的猶太人說.
神有三位.
他們會怎樣.
用石頭扔死我.
我們現在的基督徒.
之所以將三位一體的教義放入舊約.
是因為我們已經從神得到完全的啟示.
我們有舊約有新約.
在新約裡面已經顯明了神是三位一體的神.
所以當我們回頭讀入舊約的時候.
就會將裡面有暗示的文字.
暗示的經文.
用來作為舊約有三位一體的教義證據.
但是我們要了解一件事.
聖經是關於神的故事.
上星期Mark Sir說.
聖經就好像一個長篇劇.
記不記得.
有多少集.
一百集.
大家有沒有追過劇.
有追過劇吧.

$^{161}$現在我們不叫追劇了.
叫什麼.
對了 煲劇.
為什麼叫煲劇.
是因為我們等齊了.
出了之後我們可以一次過煲完.
十倍速搞定.
以前像我這種年紀.
小時候就追劇.
什麼叫追劇.
就是今晚看完.
要等明晚.
預告都好像不夠.
我要等明晚.
我要追它.
同樣神向人啟示自己.
他也是一步一步進行.
一集一集告訴我們.
他是怎樣的神.
在舊約.
神是介紹自己是獨一神.
來到新約的時候.
神繼續介紹自己.
我們就非常清晰地看到.
這位獨一的神.
祂有三個位格.
聖父 聖子 聖靈.
我們就問一個問題.
聖父 聖子 聖靈都是完全的神嗎.
那是不是有三位神.
我們怎樣去理解這個一體.
怎樣去理解只有一位神.
或者我們被一條數學題難倒了.
我們的邏輯思維裡.
1+1+1=多少.
3.
你們很厲害.
等於3.
懂得計算.
但是這個公式我告訴你們.

$^{201}$在聖經當中是不能用的.
1+1+1不等於3.
是等於1.
是不是.
這個邏輯我們很難去明白.
但是我今天和大家介紹一個方法.
告訴大家我們不用1去計算.
我們用一個數學符號.
這個數學符號相信在座各位讀過數學都應該知道.
就是一個叫做無限的符號.
這個8打橫躺在那裡.
無限.
表示一個常數是去到無限數不盡.
我們想像一下這裡有三個不同顏色的無限.
有紅色 藍色 黃色.
然後我們把它們加在一起.
它仍然是一個無限.
因為無限就是無限.
加在一起三個都無限.
就是一個無限.
但是裡面有三個不同顏色.
有三個無限.
這是一個比喻.
當我們用這樣的比喻的時候.
去理解三位一體.
我們要非常非常清楚一件事.
我們是有限的.
這個用無窮的類比.
可能已經是我們用類比去理解三位一體的極限.
所以我們還是要回到聖經.
看看神怎樣去介紹.
祂自己是一體的.
我們看一段經文.
《馬太福音》28章18到20節.
大家都很熟悉這段.
是耶穌升天之前對門徒頒佈的大使命.
當耶穌說到奉父子聖靈的名的時候.
我們說.
聖父是神.
聖子是神.

$^{241}$聖又是神.
是不是有三個不同的神呢.
如果有三個不同的神.
這裡的名紅色的字.
名.
它應該要用複數.
Names.
但是耶穌在說這個話的時候.
在原文希臘文.
或者一些英文翻譯本裡面.
名是單數.
Name沒有S.
大家可以去查一下.
也就是表明.
神只有一位.
因為祂只有一個名.
在這裡奉父子聖靈的名施洗.
就將這裡三個位格.
區分開來.
就是父,子,聖,靈.
又將三個位格連在一起.
一個名.
是單數.
一位神.
這個就是聖經在比較早的時候.
有關於三位一體的教義.
很明顯的一個地方.
第二段經文我們來看.
約翰福音十章三十節.
又考考大家.
這一節.
下一節是什麼.
三十一節是什麼.
當耶穌說.
我與父原為一的時候.
三十一節是什麼.
三十一節就是.
猶太人又拿起石頭要去打他.
打他.
為什麼要打他.

$^{281}$當時是什麼時候.
大家有沒有聖經可以翻一下.
是修奠節的時候.
修奠節是猶太人的一個節期.
是在十二月冬天的時候.
為什麼耶穌說完這句話.
猶太人要拿起石頭去打他.
我們首先來看修奠節是怎麼來的.
簡單來說.
在主前二世紀的時候.
有一個希臘大帝.
他的名字叫安條克四世.
他竟然自稱為神顯者.
神顯者是什麼意思.
他說他自己是神.
然後他還去攝族聖殿.
這個對猶太人來說是萬萬不能接受的.
你自稱自己是神.
你什麼意思.
於是猶太人就發動了一個革命.
叫做瑪加比革命.
將希臘大國擊退.
然後就拿回聖殿.
重新修整聖殿.
重現聖殿.
所以修奠節就是這樣來的.
所以在修奠節裡.
猶太人就會想起這樣的歷史背景.
在修奠節裡.
最大的罪就是攝族的罪.
這個時候耶穌竟然在聖殿裡.
去說一句.
我與父原為一.
耶穌在大庭廣眾眾目睽睽底下.
去說了一句.
我自己就是神.
猶太人就想起.
那一個很可惡的安條克四世.
哇 耶穌你這麼斗膽.
竟然自稱為神.

$^{321}$就惹怒了這些猶太人.
然後猶太人就在31節裡.
就拿起石頭.
要打他.
就是因為耶穌自稱為神.
我與父原為一.
所以既然聖父 聖子 聖靈是一體.
一位神.
那神的三個位格有沒有區別呢?.
我們來看約翰福音14章16到17節.
這裡就是耶穌快要上十字架.
要離開門徒.
耶穌就鼓勵門徒.
我要求父.
父就另外賜給你們一位保衛師.
叫祂與你們永遠同在.
什麼是保衛師呢?.
保衛師就是中保的意思.
我講起中保 大家想起誰呢?.
耶穌吧.
我想起耶穌就是中保 是保衛師.
保衛師有好幾種解釋.
就是有鼓勵者 有安慰者 有顧問.
或者策士.
策士大家又想起哪裡呢?.
還有一個就是幫助者.
耶穌說我要求父.
父就賜給你們另外一位中保幫助者.
耶穌已經是一個中保.
中保就是耶穌 幫助者.
現在耶穌要求父.
父就賜另外一位幫助者.
而這位另外的幫助者就是真理的聖靈.
這裡大家看到嗎?.
耶穌聖子求聖父.
聖父就賜下另外一位幫助者.
就是真理的聖靈.
很清楚吧?.
三個位格出來了.
聖父 聖子 聖靈.

$^{361}$有沒有區別呢?.
有.
聖父就是聖父.
他不是聖子.
因為只要求父.
他也不是聖靈.
而聖子也不是聖靈.
聖子是保衛師.
而聖靈就是另外一位的保衛師.
聖靈就不是聖子.
也不是聖父.
一點都不混淆.
不要以為我在講急口令.
非常清楚.
所以我們可以在這裡總結一下.
三位一體.
就是神只有一位.
祂有三個位格.
聖父 聖子 聖靈.
我們回到剛才的PPT.
聖父是神.
聖子是神.
聖靈是神.
每個位格都是完全的神.
這句說話很重要.
不混亂.
最後那句更重要.
然而只有一位神.
明白三位一體的教義.
其實對我們基督徒來說是非常重要的.
為什麼這麼重要呢?.
第一個就是我們要分辨異端.
剛才我們講了一個異端.
就是形態論.
異端是什麼呢?.
異端就是將神的自我啟示.
即是神所說的話.
限制在我們人的理性裡.
將所有不能用理性去明白的東西.
改到可以明白為止.

$^{401}$這就是異端.
三位一體的教義.
在他們這些異端來看.
是不能用理性去接受的.
什麼一二三三二一是什麼呢?.
所以他們就將它轉換成另一種東西.
比如大家比較熟悉的.
就是耶和華見證人.
他們其中的一個認訊.
就是將基督說成是被造的.
只不過是最先被造的.
所以就這一點.
他們的聖經翻譯本.
叫做《新世界譯本》.
就篡改了約翰福音一章一節的經文.
約翰福音一章一節的經文.
大家會不會背呢?.
太初有道.
道與神同在.
道就是神.
是吧?.
道就是神.
在他們的聖經翻譯本裡是怎樣的呢?.
道是一個神.
The word is a god.
而且那個god的G是小寫的.
但是我們的聖經翻譯本就是道就是神.
所以我們知道三位一體的時候.
我們就可以用這一點來做一個試金石.
可以幫助我們去分辨異端.
凡是與我們的教義有相反有不同有混淆的.
我們都可以定為異端.
這裡再一次讓我們看到一件事是什麼.
因為我們的理性真的很有限.
我們沒有辦法完美地去解讀三位一體.
剛才的類比幾乎就是道的極限.
所以我們應該怎樣做呢?.
我們應該專心去聽從和信服聖經的教導.
神怎樣去介紹他自己.
我們就要接受.

$^{441}$我們可以存一些疑惑.
為什麼會這樣?.
為什麼是三位一體?我真的不明白.
可以存有這些疑惑.
但是就不應該去強行曲解聖經.
第二個教導.
三位一體帶給我們一個很重要的訊息就是.
三位一體的教義告訴我們.
神的本質不是孤零零的.
不是孤零零一個神.
而是一個群體.
聖父,聖子,聖靈.
他們是有關係的.
他們的關係是什麼?.
父子,聖靈的關係.
其中一個就是.
在《約翰一書》四章七到八節這裡說.
其中一個關係就是.
神就是愛.
約翰這裡沒有說神愛.
我們記得三章十六節的時候說.
神愛世人.
God loves the world.
但是這裡沒有說神愛.
愛是一個動詞.
在這裡.
神就是愛.
God is love.
愛是作為一個本質.
作為一個屬性.
什麼意思呢?.
這裡不單單意味著.
神愛他的創造物.
愛他的子民.
或者不單單只是神.
他是充滿愛的.
這意味著神本身自己就是愛.
其實這句經文已經告訴我們.
三位一體的真義在這裡.
因為神就是愛.

$^{481}$愛不可以沒有對象.
愛.
我愛你.
要有愛者.
有被愛者.
才有愛.
才能構成愛.
在三位一體裡.
我們知道.
聖父愛聖子.
聖子愛聖靈.
聖靈也愛聖父.
聖父,聖子,聖靈之間的關係.
是一個愛的關係.
神是一位有關係的神.
三位一體彼此之間.
他們是建立一個愛的關係.
而神在創造人的時候.
神也是希望我們人和人之間.
建立一個愛的關係.
而且在這裡的第七節說.
親愛的弟兄啊.
我們應當彼此相愛.
因為愛是從神來的.
凡有愛心的.
都是從神生的.
並且認識神.
這裡我們可以看到.
神是希望人和人之間.
建立一個愛的關係.
神是愛的源頭.
而且祂將信徒.
我們描述為愛神.
並且是認識神的人.
因為愛是從神而來的.
所以我們這些認識神的人.
就應該彼此相愛.
由此可見.
我們的神和我們是有關係的.
關係應該是怎樣去建立的呢?.

$^{521}$神是怎樣去建立關係的呢?.
我們又是怎樣去和人建立關係的呢?.
我們通常會做些什麼事去建立關係呢?.
我為你做一些事.
我送一份禮物給你.
我和你一起去做一件事.
做一個事工 做一個工作.
我們有交流.
我們一起去玩 去吃東西.
這些都是可以去建立一個愛的關係的行動.
而在聖經當中.
建立關係的其中一個.
很重要的方式是什麼呢?.
大家猜猜.
猜中有獎.
就是.
記住自己的答案.
就是對話.
神的三個位格之間.
是有對話的.
我要求父.
父就賜給我另外一位聖靈.
賜一位保衛師給你.
是有對話的.
聖經就是神和我們有對話的途徑.
他藉著聖經向我們啟示他自己.
和我們去對話.
在一些人際關係當中.
特別是在教會.
在家庭裡面.
我們怎樣去建立一個愛的關係呢?.
除了剛才說去做一些事.
送一些禮物給他之外.
我們可否從神.
他是一個對話的神裡面.
去看到一件事.
我們亦都可以彼此.
去有一個有品質的對話呢?.
Quality Conversation.
什麼是有品質呢?.

$^{561}$我們看看聖經.
聖經就是最有品質的對話.
怎樣有品質呢?.
就是神不說廢話.
不說廢話.
是就是.
不是就不是.
他的每一句話都是寶貴的.
我們的Quality Conversation.
我不是說我們不可以說笑.
聊天.
開玩笑.
說一些輕鬆的話題.
我不是這個意思.
而是說我們要在對話之後.
我們彼此是心滿意足的.
在靈性上是有得著的.
而且我們的關係是密切了.
Spend some time together.
其實我們在接下來的星期.
就可以這樣做.
找一個人.
或者找一班人.
我們有一次的Quality Conversation.
剛才誰猜到的?.
有沒有人猜到?.
如果猜到的話.
崇拜之後去後台找我.
送一份獎品給你們.
獎品是什麼呢?.
就是可以約我去一起.
有一個Quality Conversation.
OK.
我最近在一次很有品質的對話當中.
學到了一句話.
這句話叫什麼呢?.
Love needs time.
當我們去花時間.
去陪伴.
去聆聽的時候.

$^{601}$這個時候.
就是我們可以彼此相愛.
去做出三位一體真神的形象.
希望今天的分享.
可以幫助我們去認識三一真神.
並且在這個教義當中.
有多一點的看見.
多一點的明白.
和多一點的實踐.
多謝大家.
同心低落祈禱.
各位愛的上帝.
你自己本質就是愛.
在你的本體裡面.
聖父聖靈之間.
彼此的相愛.
從耿古到永遠.
你就是愛.
我們的生命.
是從愛誕生出來.
也以愛,以恩典來拯救我們.
祝我們能夠有著.
父子聖靈的名在我們身上.
因為你對我們愛和救贖.
我們是屬於你的.
有著你自己這個奇妙的名分.
祝我們感恩.
多謝你將生命向我們分享.
祝我們每一個在基督裡面.
我們就是從神的而來.
從愛而來.
主求你幫助我們.
就去彼此相愛.
在我們自己弟兄姊妹之間的關係.
在我們大家的婚姻.
我們的家庭.
主願以你自己三位一體.
這個奇妙的名字.
這個相愛的本質.
能夠在我們生命當中流露出來.

$^{641}$甚至我們能夠帶著你的生命進入這個世界.
進入生命裡面每一個方向,每一個崗位.
能夠將你永恆的愛.
能夠傾倒,能夠分享出來.
主叫人能夠分享到,見到.
神在我們中間.
神在我們生命裡面有奇妙的作為.
我們再一次求你幫助我們.
更加認識你的時候.
更加愛你.
認識你的時候更加認識我們自己.
認識我們教會.
叫我們教會.
能夠著神你自己愛的生命.
我們這樣祈禱.
奉父子聖靈的名.
Amen.
多謝各位.
Phoebe, very good job.
請我喝茶嗎?.
OK, go ahead.
希望上一集讓大家平安.
下個星期再跟大家分享.
四方.
\newpage



\section{5:16-22}
\label{sec:__GpGaclvJI}
\textbf{Daily Reminders from Jesus (1 Thess 5\_16-22)- Rev. Mark Chan}
\newline
\newline
連結: \href{https://youtube.com/watch?v=_-GpGaclvJI}{\texttt{https://youtube.com/watch?v=\_-GpGaclvJI}} ~~~~ 語音日期: 2023-11-21
\newline
\newline
\hyperref[sec:z0bvgIPiO7c]{\small{< < < PREV SERMON < < <}}
~
\hyperref[sec:index_chronic]{\small{[返順時目]}}
~
\hyperref[sec:index_scriptual]{\small{[返順卷目]}}
~
\hyperref[sec:gNNugJCYC2g]{\small{> > > NEXT SERMON > > >}}
\newline
\newline
$^{1}$- All right, thank you, Amir..
And welcome all of you and happy Thanksgiving to you..
I know some of you are out of town.
and you probably have some special plan..
And today we are going to the luncheon..
And so I will keep it really, really short..
So we'll be out and be there by one o'clock.
and we're gonna have a good lunch together.
and play some games and get to know one another,.
connect with one another..
And that's what Thanksgiving is all about.
because it's all about family and taking care.
of one another..
Let's all pray together..
And I got the news that Annie is in active labor right now..
Yeah, so let's pray for her.
and pray for those who are out of town as well..
Let's pray together..
Father, give thanks, give thanks in good times.
and in bad times..
And most of all, we want to give you thanks.
because of what you have done for us in Jesus Christ..
And I pray that we will remember and be reminded.
you have given us the greatest gift of all,.
which is your son, Jesus Christ..
Thank you for the salvation..
Thank you for the identity being your sons and daughters.
and the eternal life you have given to us..
I pray for Annie, I pray for James, I pray for the baby..
I pray that this time you will be with them.
and give strength to Annie as she is in active labor..
And I pray that as this newborn baby comes,.
they will be ready and give them a sense of peace and love.
and to receive this new life in the family..
I pray for those who are out of town.
and I pray that you will grant them safety..
And in this week of Thanksgiving, be reminded.
that we do need to be grateful.
for all that you have given to us in Christ in prayer..
Well, I have understood.

$^{41}$that Thanksgiving is actually a discipline..
Thanksgiving is something that we have to do.
and we learn to do..
And if you ask someone in the Bible.
or the apostles or the disciples,.
pastors in the Bible, in the New Testament,.
probably the one guy who can tell us so much.
about the meaning of true Thanksgiving is the apostle Paul..
He was preaching the gospel and he was imprisoned in Rome.
for the sake of the gospel.
and he wrote this to the church in Ephesus..
He said, "Sing and make music in your heart to the Lord.".
Now again, to sing and make music is not an option..
It's not something you can choose from the choices..
It is a command..
It is an imperative..
It's telling you, you have to do it..
Sing and make music in your heart to the Lord..
Always giving thanks to God the Father for everything.
in the name of our Lord Jesus Christ, Ephesians 5..
Now think about it, always giving thanks for everything.
no matter what circumstances, good times or bad times..
To Paul, the writer of this verse,.
Thanksgiving is not a once a year celebration.
but a daily reality, every single day,.
a daily discipline to cultivate a thankful attitude..
And such an attitude will change our life view.
and make us joyful people in every situation..
The worship leader remind us to reflect on this year,.
this past year, what had happened in your life.
and how would you conclude.
and how would you describe this past year, 2023..
To reflect upon it is to help us to find some compass,.
some direction in what God is doing in our lives..
Thanksgiving reminds us that God is a God of blessing..
God has blessed us and that's why Thanksgiving.
is one of the most important marks of Christian life..
When we follow Jesus, we must not allow.
a spirit of ingratitude because it will harden your heart..
It will chill your relationship with God and with others..

$^{81}$Thank God for the blessing, material blessing,.
spiritual blessing for the people in your life..
Most of all for Jesus in your life..
You see the mind, our thinking is like the belly..
We eat garbage in and garbage out..
If we take on, we take in toxic substance,.
it will ruin our system, ruin our digestive system.
and indeed our whole body, our mind, our spirit..
The same goes with our mind..
If our thinking, if our thought life is filled with myself,.
my own interest, my own voice, my own want,.
your mind will deteriorate..
It will, it will..
'Cause it will turn you into bitter self,.
selfish self, dissatisfied, ungrateful heart.
who do harm to all of us..
The only way to restore it is to regain.
the spirit of contentment..
What we have right now is more than we need..
Yes, there are some needs and there's some wants.
in our lives and we pray for them,.
but the spirit of contentment says,.
"Hey, what I have right now is enough.".
And that kind of spirit will be a good soil.
to yield joy and thankfulness..
Today, I think we are often,.
so full of unthankfulness..
It's too common of a spirit, right?.
Kids forget to thank their parents for all that they do..
Common courtesy is hard to find..
We often take for granted for the things people do to us..
People help us, we take it for granted..
People are gracious to us, love us, we take it for granted..
And above all, if we do that,.
we fail to thank God for His blessing..
Ingratitude in the Bible is sin..
Let's put it this way..
Ingratitude is sin..
As much as lying, stealing, immorality,.
any other sin condemned by the Bible..

$^{121}$Ingratitude is one of them..
So one of the Bible's warnings in the New Testament.
is found in Romans chapter one..
It talks about those who don't know God,.
they knew God, but they neither glorify Him as God.
nor give thanks to Him..
They forget about being thankful..
Those who don't know God, don't sense the existence of God,.
don't worship God, is marked by ungratefulness, ingratitude..
Let me ask you this..
What are you thankful for?.
If your mind goes blank,.
we probably have to cultivate that spirit again,.
the spirit of thankfulness..
From one end to the other end of the Bible,.
we are commanded to be thankful..
I know it's hard..
When things don't go well,.
I speak from my personal experience..
When things go well in my life,.
I am angry, you know..
My family members can tell you..
I'm an angry man..
I complain..
But Thanksgiving reminds us,.
reminds me that this kind of spirit comes from my heart..
I need to be constantly attuned to God..
Psalm 147, saying to the Lord with thanksgiving,.
we are commanded to do that..
Colossians 3, let the peace of Christ rule in your hearts.
since as members of one body, the church,.
you were called to peace..
But be thankful, right?.
Be thankful..
It's not about being polite to God..
It's about worship..
Who is God in your life?.
If Jesus is God, if Jesus is the Lord,.
we give thanks, right?.
It's not about being nice..

$^{161}$Thankfulness is a command.
and it's a mark of our spiritual life..
But we have to learn..
We have to cultivate it..
In bad times, in good times,.
we learn and we discipline our mind..
So what should we be thankful for?.
Let me close with a few things.
to remind us what we need to be thankful for..
Number one, thank God for the material blessings.
that God gives you..
Paycheck, material enjoyment, shelter, a job,.
even though it sucks, right?.
At times it sucks..
But be thankful, be thankful..
Physical health, right?.
Yeah, sometimes we have back pain.
and all these different problems..
But be thankful..
Second, thank God for the people in your life..
Your family members, yeah, they suck sometimes..
They suck sometimes, right?.
They get us angry, they agitate us, you know?.
And I agitate them..
But be thankful for the people in your life..
It's so easy to take people for granted,.
especially people who really care for us and pray for us..
But we need to give thanks for all those around us..
Our spouse, our children, our friends,.
our church members, relatives,.
anyone who come my way, surround me with care..
Number three, give thanks in the midst of trials..
We face different difficulties and sometimes trouble..
And, you know, look at the news..
This world is such a broken place..
It's dangerous, it's painful..
Because of all these difficulties and challenges.
in the world..
Our personal life as well..
But, you know, the Bible tells us,.

$^{201}$the letter of James, consider it pure joy,.
my brothers and sisters..
Whenever you face trials of many kinds,.
be joyful..
Why?.
Because you know that the testing of your faith.
develops something inside of you.
that can never be taken away..
That is your character, your spiritual life..
Perseverance, patience..
How can you have perseverance?.
It has to be developed from trials and difficulties, right?.
Patience and all these virtues and fruit of the spirit..
It has to be cultivated from times of trials..
All the faithful men and women in the Bible.
have their own difficult times..
Think with me..
From the first book of the Bible, Genesis,.
Abraham, Isaac, Jacob, Israel, Moses, David,.
even Jesus Christ our Lord, as the perfect son of God..
He had his own troubles..
He went through trials and difficulties..
But he give thanks..
Well, I don't know about you,.
what kind of trials you are going through,.
or maybe you're having good times of your life, right?.
You think about 2023..
This is such a great year, right?.
How many of you think like that?.
2023 is a great life..
Okay, great..
We got one hand up..
That's nice..
Maybe you're envious..
Think about 2023 to you..
It may be difficult, or indeed it is difficult..
Nevertheless, cultivate a spirit of thankfulness,.
even in the midst of heartaches and trials..
Number four, that's the last one..
Thank God, especially for his salvation in Jesus Christ..

$^{241}$God has given us the greatest gift of all, his son..
He died on the cross and rose again to give us hope,.
to give us eternal life,.
so that we can know God personally,.
and we can know the Father and spend eternity.
with the Father in heaven..
And the Bible tells us we are separated..
We don't deserve this grace of God, because we have sinned..
We are separated from God..
We are enemies to God, but God loves us..
He loves you, and he loves me so much.
that he, God, wants us to be part of his family..
And he did it to save us,.
send his only son into the world to die for us.
as perfect sacrifice..
Give thanks that Christ the Savior has come, right?.
Turn to him, turn to God with a simple prayer of repentance.
to say to God, "I'm sorry..
I take things for granted..
But here I am..
I want to recover that faith, that spirit of thankfulness..
Thank God for what you have done for you and for me.".
Let's all pray together..
Father, once again, we give you thanks..
We're sorry we take too many things for granted,.
for the life that we have,.
for the people we have in our life,.
for school, for work, for family, for friends,.
and brother and sister in Christ, for our church..
We complain too often..
We come before you with a list of needs.
as if we are lacking in our life..
Father, forgive us..
And again, we kneel our mind,.
and we cultivate that spirit of thankfulness..
As we think about this whole year coming to the end,.
no matter if it was a good time.
or there were ups and downs,.
hills and valleys,.
we come before you and give you thanks..

$^{281}$Thank you for being our God..
Thank you for loving us in Jesus Christ..
In your name, O the Son, we pray, amen..
All right, off we go..
We are going to the luncheon, all right?.
So be thankful, all right?.
So say hi to one another before we depart, okay?.
I wanna see you next week, next Sunday..
Let's worship and serve God together..
[BLANK AUDIO].
\newpage



\section{}
\label{sec:gNNugJCYC2g}
\textbf{以巴衝突的反思 (4) - 陳崇基牧師}
\newline
\newline
連結: \href{https://youtube.com/watch?v=gNNugJCYC2g}{\texttt{https://youtube.com/watch?v=gNNugJCYC2g}} ~~~~ 語音日期: 2023-11-21
\newline
\newline
\hyperref[sec:__GpGaclvJI]{\small{< < < PREV SERMON < < <}}
~
\hyperref[sec:index_chronic]{\small{[返順時目]}}
~
\hyperref[sec:index_scriptual]{\small{[返順卷目]}}
~
\hyperref[sec:UIZxTjRVxE4]{\small{> > > NEXT SERMON > > >}}
\newline
\newline
$^{1}$各位好,感恩節快樂.
我們繼續看以巴衝突.
今次是第四講,應該做完了.
我知道在世界當地發生很多不同的事.
我們都需要處理一下.
但亦不能夠經常處理那件事.
但在處理以巴衝突.
想起這些事情,反思的時候.
我們都相信對聖經和信仰的認識可以深一點.
我們一直以來都在處理五個命題.
今次是處理第四個和第五個.
第四個是什麼呢?.
很多時候在我們對聖經的體會和了解.
我們會覺得神的子民有兩個.
一個就是血統上的猶太人.
一個就是外邦人,即非猶太人.
我想華人教會,特別是上一代.
就受一個所謂神學思想的影響.
叫做時代論.
Dispensationalism.
可能你未聽過這東西.
如果未聽過,上帝祝福你.
不需要去想這些.
時代論是什麼意思呢?.
就是將整個舊恩的歷史分成很多不同的時代.
有一個時代,大家可能會聽過.
就是我們不再是律法時代.
現在是恩典時代.
有沒有聽過?.
這個是時代論遺下來的一些思想.
或者一些不太清楚的概念.
無論如何,我們受這個概念影響的時候.
我們就很多時候想.
上帝在地上的人類的計劃是有兩個的.
Plan A就是給血統上的猶太人.
即阿巴罕以色列的後裔.
血統上的.
Plan A就是給他們的計劃.
然後Plan B就是給一些像我們這樣.
外邦人作為教會.

$^{41}$所以我們今時今日就是教會時代.
一個主要外邦人的時代.
所以如果這樣看的時候.
我就發現上帝是有兩個的子民.
福音是猶太人而出的,Plan A.
但是當耶穌基督死而復活,升天之後.
多不多猶太人信主?.
很少,是吧?.
所以上帝有一個Plan B.
所以福音就去到外邦人的人群當中.
教會時代就是今時今日.
直到基督再臨的時候.
猶太人是絕少的.
外邦人是更多更多.
這個就是教會時代.
我們在等什麼呢?.
這個時代論不代表聖經.
一定是這樣說的.
等一下我們處理一下.
我們在等什麼呢?.
等基督再臨.
因為基督再臨的時候.
神秘地再臨.
然後教會就會被提.
我們就會突然消失.
然後七年大災難就會出現.
七年大災難說到那些很轟天動地.
一些災殃會淋到地上.
那些的描述在哪裡找到呢?.
是在啟示錄裡找到的.
然後亂局一輪.
敵基督就出現.
然後統治整個全球.
有七年的大災難.
當七年大災難的時候就非常困難.
剩下來的就是一些不信主的人和猶太人.
上帝對他們有一些恩典.
七年大災難完了之後.
敵基督被消滅.
撒旦被消滅.

$^{81}$然後千禧國度等等.
是不是很煩?.
這個就是我們上一代華人教會留下來的一些概念.
我們需要處理聖經裡面到底是怎麼說的.
首先我們處理這個命題的時候.
神的子民是不是有兩個呢?.
一個是血統上的猶太人.
另一個就是外邦人.
其中一段經文就是羅馬書第九章至十一章.
我們很多時候要鬥賭的這段經文.
一至八章羅馬書是說上帝在基督裡面的救恩.
給我們的兒子在基督裡面成就.
然後到第九章至十一章.
他就處理當時羅馬教會裡面的問題.
你可以想像到當時羅馬教會的背景是怎麼樣的.
外邦人是為多的.
並且我們知道當時第一世紀.
大概主要是四十年代的時候.
四十年代,五十年代.
羅馬經都是一種反猶的意識是非常強的.
Anti-Semitism.
甚至在一個羅馬王叫做Claudius.
克勞蒂.
他在一年的政策中.
趕走了很多在羅馬京城的猶太人.
將他們打散到其他地方.
所以那裡的反猶主義是非常強的.
於是你可以想像當時的羅馬絕多外邦信徒的教會.
他們就帶著這種反猶的心態進入了教會.
然後那些外邦信徒就以為.
所有的猶太人面臨的結局就是永遠的沉淪.
背景的時候我們看羅馬書9至11章的時候.
我們就明白到底保羅想說什麼.
保羅在9章至11章的羅馬書.
他就說原來猶太人是還沒game over的.
上帝仍然在他們的生命裡面.
在基督裡面是有計劃的.
所以在第9章第6節的時候.
保羅就說雖然不是那麼多猶太人信主.
不過不代表上帝的話語落了空.

$^{121}$因為上帝應許過請求萬民.
甚至是他自己本身的選民猶太人.
是不落了空.
什麼意思呢.
因為從以色列生的不都是以色列.
肉身所生的兒女不是上帝的兒女.
唯獨那應許的兒女就是信基督的.
那些先算是後裔.
那你就可以看到有兩個圈.
他說不是從以色列生的不都是以色列.
是什麼意思呢.
就是有兩種sense of Israel.
第一種sense的就是大圈.
就是血統上的猶太人.
第二個就是那些真的以色列人.
就是信耶穌的那些以色列人.
保羅就是其中一個.
所以他說不是上帝的話落了空.
只不過我們要看清楚.
原來在以色列當中.
真正的以色列人不是血統的問題.
真正的以色列人是信耶穌那班人.
信基督.
猶太人要救主.
世人要救主.
信彌賽亞.
這班才叫做true Israel.
或者叫做spiritual Israel.
ethnic Israel不代表什麼.
雖然他們血統上有很多法.
有很多恩典.
有很多認識上帝的道.
有很多不同的渠道.
但是不代表什麼.
ethnic Israel不代表什麼.
spiritual Israel或者true Israel.
才代表他們真正屬於神的子民.
作神的兒女.
所以在第九章說完一大堆的時候.
然後我們來到第十章.

$^{161}$這段聖經我們很熟悉.
口裡宣認.
耶穌為主心裡信神.
叫他從死裡復活就必得救.
因為人心裡信就可以清兒.
口裡宣認就可以得救.
這個很多時候我們信主初信的時候.
我們經常會記起這段經文.
但最重要的十二節說什麼.
為什麼信耶穌就會得救.
口裡相信承認.
心裡認同就可以得救.
只是靠耶穌就能得救.
因為猶太人和希臘人並沒有分別.
這句說話已經很有革命性.
猶太人和希臘人並沒有分別.
不代表猶太人特別好.
也不代表希臘人和外邦人差一點.
因為人人都有一位主.
這位主也口袋叩告他每一個人.
所以外邦人和猶太人在十字架面前.
就是平地.
十字架在十字架的前面.
平地.
不是一個高一點一個低一點.
不是猶太人高一點.
外邦人低一點.
也不是相反不是外邦人高一點.
然後猶太人是避免著.
是平地.
於是保羅就用了一個比喻.
我不知道大家有沒有種花種植物.
種樹.
什麼叫做移花接木.
把一些舊的枝子切掉.
然後把另一棵樹的枝子貼上去.
接住它.
有沒有試過.
有種東西吧.
不知道種了什麼出來.

$^{201}$因為移花接木.
他用了一個比喻.
就是橄欖樹的比喻.
大家知道這裡一棵橄欖樹.
有一部分不結果子的枝子被剪下.
然後剪下之後.
把另一棵樹野橄欖.
黃色那個.
配上原本的橄欖樹上.
保羅想說什麼呢.
他說舊因是從猶太人而出.
猶太人沒錯是一些不信主的.
就被戲絕了.
右下角.
然後上帝就把這個舊因.
向外邦人傳講.
也是按照上帝的計劃.
然後從野橄欖.
信主的外邦人.
就像野橄欖樹.
配上原本的樹幹.
所以整個樹幹.
就有本身信主的猶太人.
和信主的外邦人.
你在哪裡.
我們是野橄欖.
野橄欖.
我們被配上本身舊因的樹.
所以有些不信主.
戲絕耶穌的猶太人.
就像被剪下一樣.
再問清楚一點.
這裡有多少棵樹.
一棵樹.
不是兩棵樹.
舊因只有一棵樹.
神的子民只有一棵樹.
只有一棵.
你和我是外邦人.
我們本身在另一棵樹.

$^{241}$被剪下來.
在野橄欖.
然後配上了主要的樹幹.
這些移花接木的比喻.
想講什麼.
你記得我剛才講的羅馬書的背景.
就是外邦人怎麼想.
他們覺得我們就是祝福那一群.
猶太人就是被戲絕.
你看他們多慘.
所以上帝在他們的生命當中.
在猶太人的計劃當中.
就應該是遊戲結束.
本來就用這個比喻.
來跟外邦的信徒說什麼.
你們不要這麼了不起.
因為你們只是野橄欖.
被配上.
所以你們都要結出悔改的果子.
十一章第十七節.
他就講這個比喻.
若有幾斤的枝子被折下.
這裡講原本的橄欖枝子.
你外邦人.
這野橄欖.
接上去.
同享本身橄欖根的肥汁.
我們能夠生長.
為什麼.
因為本身橄欖的枝幹.
能夠使我們有生命.
你就不可向舊的枝子.
就是那些猶太人.
去誇口.
若是誇口就知道.
該知道不是你托著根.
而是根托著你.
這個很清楚.
我們是野橄欖枝子.
被配上樹幹.

$^{281}$不是枝子承載樹幹.
而是反過來.
樹幹來祝福枝子.
所以他說.
他們.
就是那些不信主的猶太人.
因為不信所以被折下來.
你因為信耶穌.
所以你站立得住.
你就不可自告奮勇.
戰戰兢兢.
所以對著那些外邦人.
信主的.
他們知道.
保羅說這是恩典.
但是猶太人是不是就是game over呢.
又不是.
十一章第25節.
大兄們.
我不願你們不知道這奧秘.
恐怕你們自以為聰明.
你們外邦人自以為聰明.
這奧秘就是有一部分.
以色列人是硬心不信主.
等到外邦人的數目天換了.
以色列全家都要得救.
或者這樣.
以色列全家都要得救.
原來有一個奧秘.
是繼續去unfold.
上帝的plan.
上帝自己救恩的計劃.
特別在猶太人的生命裡面.
就是不斷有一個計劃.
是unfold開展.
現在你未必知道.
但將來你就知道.
什麼時候知道.
等到外邦人信主的數目天換了.
這樣以色列全家就必得救.

$^{321}$這裡有一個很重要的字.
有一個連接詞.
連接詞.
可以有不同的翻譯.
和合本翻譯.
於是以色列全家得救.
新譯本翻譯.
這樣或者同樣地.
以色列家.
都要得救.
第三個翻譯.
ESV.
In this same way.
All Israel will be saved.
這裡先搞清楚.
以色列全家.
All Israel.
這句說話.
以色列全家.
經常在舊約裡面出現.
甚至兩約中間的猶太人著作都出現.
並不代表每一個猶太人.
只不過是說.
很多猶太人.
所以你知道.
All Israel 以色列全家.
不是每一個猶太人.
只不過是一個代表性的數目.
All Israel.
好像我們今天會聚餐.
All Church.
Go to the lunch.
並不代表每一個.
可能你沒空去不到.
但整個教會去了.
同樣的意思.
並且翻譯的問題.
於是這樣.
照樣.
In the same way.

$^{361}$是甚麼意思.
這個連接詞重要的地方.
這個連接詞本身.
不一定代表.
時間上的先後.
不是說某件事發生了.
然後那件事才會發生.
不是A發生了之後.
然後B才會發生.
所以第一條時間線是錯的.
不是說外邦人全數都得救.
然後很多的以色列人就得救.
這個翻譯應該是.
這樣或者同樣的.
In the same manner.
In the same way.
意思是下面的時間線.
外邦人全數會得救.
包括你和我在基督裡面.
我們是計算在這個數目上.
外邦人全數會得救.
包括你和我在基督裡面.
我們是計算在這個數目上.
外邦人仍然不斷會得救.
In the same way.
All is will be saved.
外邦人怎樣得救.
歷史歷代因為耶穌.
相信耶穌受洗歸入基督而得救.
猶太人呢.
全以色列都要得救.
就是In the same manner.
外邦人怎樣得救.
猶太人就是歷史歷代的上帝.
仍然去拯救他們.
我們差不多看完羅馬書9至11章.
請你們聽清楚.
羅馬書9至11章.
沒有講過任何政治上的東西.
沒有講過1948年伊斯蘭福國等等.

$^{401}$全部都是講在基督裡面.
能夠得救.
上帝對世人是不是很多恩典.
對我們很多恩典.
我們先不講猶太人.
我們能夠被接上.
真鑒冷船.
成為上帝的子民.
不是因為血統.
不是因為我們有什麼好處.
而是因為耶穌所做.
我們就是.
真的以色列人.
因為不是在講ethnicity.
是在講your faith in Christ.
我們轉去另一段經文.
加一看書第3章.
這裡講其實你們直著信.
在基督耶穌裡面.
都成為了神的兒女.
你們凡受洗歸入基督的.
都被帶基督了.
不再分什麼.
猶太人和希臘人.
不再分維羅的自主的.
不再分男或女.
為什麼不分呢?.
因為我們都在基督裡面成為一個的子民.
既然你們是屬基督的.
留意這裡.
既然你們是屬基督.
你們就是亞伯拉罕的子孫.
你們不是亞伯拉罕子孫裡面的二等公民.
你們就是亞伯拉罕的子孫.
是照著應許承受產業的.
這裡不是在講replacement theology.
這裡不是在講外邦人.
這裡不是說外邦人我們就取締了猶太人.
我們變成選民.
然後他們就選民不再.

$^{441}$不是replacement.
而是fulfillment.
在哪裡fulfill?.
在Christ 在基督耶穌的裡面被fulfill.
這個就是上帝對我們的恩典.
我們能夠成為上帝的子民.
所以神的子民只有一個.
樹只有一棵.
真橄欖樹.
外邦人被接上是因為恩典.
處理完這個命題.
處理第五個命題.
就是我們很多時候講到中東有什麼衣沃.
或者耳巴有什麼衝突的時候.
我們就說這個末世徵兆越來越多.
就快教會要被提.
然後呢七年大災難就要來到了.
啟示錄裡面說.
七日七晚七日等等.
這些災難就會到.
迪基督就會出現.
迪基督就會統治整個世界.
並且叫以色列和周邊的國家來講和.
立約立和約.
這個就是七年大災難的情況.
所以每逢耳巴有任何張力的時候.
有很多人就會出來用這個時代論的方法來詮釋.
我知道這些東西都很複雜.
什麼叫時代論.
很多時候就會畫出一個這樣的時間表.
開頭的時候我說什麼.
時代論.
如果你沒聽過的話.
上帝祝福你.
因為你不需要搞這些東西.
我問你你現在在哪裡.
按照時代論的時間表.
末世時間表.
教會時代最後的時候.
在等什麼呢.

$^{481}$教會被提.
教會被提之後就會七年大災難.
然後有迪基督出現.
然後搞完一大輪很多災難之後.
基督就親自來到伏林.
千年國度就會展開.
然後有大審判.
新天新地就會到.
這是時代論有的時間表.
是不是真的這麼簡單呢.
我小時候都是在我的教會.
是一些時代論受影響的教會.
我記得1990年代的時候.
我不講那麼多理論.
我講一個故事給你聽.
1990年代的時候.
當時誰是美國總統.
Bill Clinton.
他做了一任之後.
他想再競選.
我記得當時教會有很多末世講座.
加上90年代初的時候.
又有一些波斯灣爭戰.
伊拉克等等.
當時的政局都很混亂.
去那些末世講座.
我記得那個牧者在呼籲我們.
他說你要小心.
因為按照這些末世的徵兆.
敵基督已經在世界上.
加上當時發生什麼事呢.
你記得Clinton其中一個外交政績.
就是去找以色列和巴基斯坦做和約.
叫做Oslo Accords.
1993年的時候.
你記得右邊那個.
就是Isaac Arafat.
是巴解組織的主席.
左邊那個就是以色列的總理.
Rabin.

$^{521}$找了他們去搞和約.
當時那些末世的徵兆的講座等等.
種種的跡象就說.
敵基督已經在人世.
其實那個講完沒有明講.
但是暗示叫我們不要選Bill Clinton.
因為這個人一定是敵基督.
我記得完了講座之後.
我們團契弟兄姊妹.
我們就去喝咖啡.
有個弟兄問得很好.
他說如果這個人真是敵基督的話.
我們一定要選他.
為什麼我們要阻住地球轉呢.
因為如果選了他.
我們就被提出大災難.
多好啊.
你記得我在講哪一年.
90年代中.
如果那時候敵基督真的來到.
7年大災難.
教會被提.
是不是很好啊.
我也消失於人世了.
我們大家還需要搞.
20多年了.
由那時候搞到現在差不多30年.
如果那時候真的有7年大災難.
那就好了.
搞之前幾年有疫情.
搞到世界都連7,8遭.
為什麼.
那個弟兄喝咖啡的時候說.
為什麼我們不選他.
我們要阻住地球轉.
What is going on.
我們混雜了地上一些政治的東西.
進入了全息聖經.
那些不是聖經全息.
那些是現世中東的全息.

$^{561}$末世徵兆的講法.
這些末世.
很多時間表好像畫清楚給我們.
每逢中東以巴有大亂的時候.
就會用不同的版本出現.
以前我們就說.
你記得那時候蘇聯的時候.
誰是所謂的敵基督.
戈爾巴喬夫.
接著蘇共解體之後.
俄羅斯.
現在就說普頓.
其中一個,起碼北方的王是他.
有很多很多用現代的報章.
裡面見到的新聞.
我們就他壓下去.
聖經.
這些絕對不是一些嚴謹的全息.
並且你知道這些講法.
是可以無限後退的.
那些老一點的應該記得.
那時候歐洲搞EU的時候.
是不是有十個國家一起.
現在不止了.
那時候已經說得很厲害.
EU有十個國家.
因為十個國.
說到但以理斯的異象.
有十個國出來.
然後末世就走到了.
結果就不是十個.
繼續多下去.
那你怎麼解釋.
是可以無限後退的.
我想問一個問題.
為什麼這些所謂的末世徵兆等等.
如果你記得我說過.
這些中末論的東西.
你可以回去聽聽.
這裡不處理太多.

$^{601}$被提是不是真的合乎聖經.
不是.
你可以回去聽聽.
或者是一些啟示錄等等.
你可以回去聽聽.
但我想問一個問題.
為什麼這些末世徵兆.
這些說法.
可以在世間上大行其道.
因為合我們口味.
合我們口味.
為什麼合我們口味.
因為世界大亂的時候.
我們總覺得基督的信仰.
就是叫我們離開世界.
脫苦海.
挺好的.
今時今日有痛苦.
我立刻離開.
所以剛才的時間表有一個理論叫.
災前避題論.
七年大災難之前會避題.
因為知道這些理論.
很多時候.
有一個版本.
有些人就說不是災前避題.
應該是災中避題.
三年半之後.
我們就災中避題.
有些是災後避題.
所以災前災中災後.
但大部分人喜歡什麼.
或者選擇什麼.
當然是災前.
災前走人就最好.
這個就合我們口味.
我們覺得信仰就是脫苦海.
因為不需要理這個亂七八糟的世界.
我們基督徒在地上.
行公義好憐憫.

$^{641}$我們最好不要管.
用不著管.
靈魂得救就罷.
你說是不是這些末世的理論.
中末的理論.
就是合我們口味.
還合我們什麼口味呢.
就好像當聖經.
舊約新約的寓言.
就變成.
已經預先寫好的現代的報章.
太好了.
如果是這樣.
為什麼合我們口味.
因為我們好像掌握到世界發生什麼事.
只要我們掌握到時間表.
我們就安心.
但我們卻沒有去信教歷史背後的主宰.
你能否理解我的意思.
如果我們可以神秘地知道什麼時候是基督.
我們是不是會有更多把握.
更多把握選誰呢.
順帶一提,這也是美國的產品.
美國的神學.
因為我們才有這樣的東西去搞.
什麼時候有反應.
然後反應那個候選人.
到底他是不是基督.
這是美國的神學,不是聖經的神學.
在亂局裡面.
我們覺得這樣.
我們有多一點把握就太好了.
你想想,那些人去占卜.
有個水晶球看看.
有紫微斗數看看.
為什麼.
因為我們多一點把握就好了.
我們就有安全,有平安.
這是合我們口味.
但並非聖經裡面所教的.

$^{681}$聖經告訴你.
你不要管世間將來有未知之數是怎樣.
總之你很好奇.
你很想知道.
不一定是世界大的局勢.
可能在你心靈裡面.
你覺得如果我真的掌握到.
下個禮拜發生什麼事.
下一年發生什麼事.
那就好了.
但這個聖經告訴我們.
不是真的信仰.
因為我們所信的就是紫微斗數.
信的不是我們心靈的主.
我們所信的.
只不過是我們想在心靈裡面多一點時間把握.
我們要信的是歷史的主宰上帝.
多亂也好.
祂仍然掌權.
我們可以信靠祂.
道路多艱難也好.
我們知道上帝給我們有力量.
有恩典.
可以走下去.
這真是我們今天感恩節.
因著基督的恩典.
可以成為祂的子民.
成就這個救恩.
我們同心得著基督.
你自己的應許.
在基督裡面.
你所有的應許都是阿們.
因為你是我們生命當中真正的主宰.
在世局裡面有很多令我們很混亂很憂困的事.
但我們知道在你裡面能夠有平安.
耶穌你曾經對我們說過.
在世上可能會有苦難.
但我們可以安心.
因為你已經勝過世界.
主叫我們繼續緊靠,順靠你.

$^{721}$前面的道路我們真的有很多不能把握.
不能夠掌控的.
無論是我們生命裡面任何一些計劃.
可能有一些夢境.
可能有一些道路迂迴曲折.
我們未能夠見到前路.
但我們知道前路在你手.
神求你幫助我們.
繼續順靠你.
在順靠你當中我們繼續在地上.
行公義,好憐憫.
願你的國光臨.
願你的子言行在地上.
從我們個人,家庭,一切的角色裡面.
主你的平安能夠臨到.
祈禱,防止耶穌的名字.
願上帝祝福各位.
好.
\newpage



\section{}
\label{sec:UIZxTjRVxE4}
\textbf{Israel-Palestine Conflict - A Reflection (3) - Rev. Mark Chan}
\newline
\newline
連結: \href{https://youtube.com/watch?v=UIZxTjRVxE4}{\texttt{https://youtube.com/watch?v=UIZxTjRVxE4}} ~~~~ 語音日期: 2023-11-28
\newline
\newline
\hyperref[sec:gNNugJCYC2g]{\small{< < < PREV SERMON < < <}}
~
\hyperref[sec:index_chronic]{\small{[返順時目]}}
~
\hyperref[sec:index_scriptual]{\small{[返順卷目]}}
~
\hyperref[sec:wrvuGGxg0R0]{\small{> > > NEXT SERMON > > >}}
\newline
\newline
$^{1}$- Hey, Chris, it's CEO..
And today, I would like to end this three-part series.
on some reflections on the Israel-Palestine conflict..
If you are following the news stories,.
this conflict is still going on.
with ground invasion from Israel..
And there were a lot of people who lost their lives,.
and many of them, over 50\% of them in Gaza,.
were children and women..
So this is very tragic..
But we'll take this conflict.
in light of this current conflict,.
and we ask the question, what we can reflect upon..
We know that ideas do have consequences..
Theology matters..
Your understanding of God and faith.
in the Bible is critical..
If we have a truncated understanding of faith.
or the Bible or God,.
it will trickle down into our daily life,.
and it will affect us..
Let me give you some example..
If you believed God to be a tiger parent,.
any one of you have a tiger mom or tiger dad?.
If you believe God is like that,.
tiger parent, stern and straight,.
it's probably very difficult for you to translate.
grace and joy in your own life, right?.
If you constantly think that God,.
in that kind of image, stern and straight,.
tiger parent, always a righteous judge,.
it's pretty difficult to live a graceful life..
And/or if you think that the Bible.
is a moral book of legal demands,.
there are hundreds of laws or legalistic demands,.
lists of do's and don'ts in the Bible,.
then you probably have two options..
One option is to be a Pharisee, right?.
A Pharisee would come up with a list of manageable works,.
of good works, that he can attend to..

$^{41}$He can do all of this do's and don'ts.
and feel good about himself, right?.
It's a self-righteousness..
Or the other option is that you live hopelessly,.
and there's no way you can match the standard of God..
You always feel you're not good enough,.
and grace is so far away from you, right?.
So what I'm saying is ideas do have consequences,.
and our understanding of God should be growing, right?.
We get a better understanding of the Bible..
So in this, in the light of this conflict,.
we ask the question, what can we reflect upon?.
And one idea we often encounter in the news story.
is that God seems to have two peoples..
One is the Gentiles..
The other one is the Jews..
So it seems to be that God has two plans, right?.
Plan A is to save the chosen people, so-called,.
ethnic Israel or ethnic Jewish people..
And then plan B is to save the Gentiles, right?.
And we are Gentiles, right?.
We are non-Jews..
Ethnically, we are not Jews..
So there is a popular theory that God has two plans,.
plan A for Jews and plan B for Gentiles..
And so it ends up God has two distinct peoples..
Is that really true?.
Is that really biblical?.
Now, if we think that God has two plans,.
if we think that God has two peoples,.
and one of the two peoples has more privileges, right?.
They are kind of a superior people, right?.
Being chosen, right?.
And the other one, the Gentiles,.
you know, they kind of did plan B, right?.
So plan A was to save one kind of people..
And if that kind of people remain unbelief,.
the gospel spilled out to the rest of the secondary people,.
right?.
If we think that way, this idea has consequences..

$^{81}$We will see one people superior,.
the other people inferior..
So we have to examine that kind of notion,.
one for the ethnic Jews and one for the Gentiles..
Is that really correct?.
So we have to go into a big passage in Romans 9 to 11,.
three chapters long..
So in the letter of Romans,.
Paul is trying to explain to us.
what it means to be in Christ..
We are standing in Christ,.
that means we are fully righteous,.
not because of what we have done,.
but what Jesus has done for us on the cross..
So chapter one to chapter eight,.
he's trying to expand what it means.
to believe in the gospel..
It's all about Jesus Christ..
It's not about us..
And then in 9 to 11,.
he's trying to tackle another tricky question..
The Roman Christians might have been asking,.
and to understand these three chapters,.
we have to understand some sort of background.
of the Book of Romans and the city of Rome.
in that time period, in the first century,.
especially mid first century..
What is going on there?.
What we can do is that we can outline three things..
And the first thing, as you can see on the screen,.
most Roman Christians were Gentiles at that time..
So they lived in the city of Rome,.
they gathered in small groups in the bigger homes,.
and they constitute the early church in the city of Rome..
But most of them are non-Jews, Gentiles..
And then we know that in the city of Rome,.
especially in the first century,.
there was a wave of anti-Semitism..
People just don't like the Jews..
And that kind of feeling,.

$^{121}$that kind of atmosphere was very, very strong..
And in fact, we know that some of the Roman Empire policy.
would kick out some of the Jews out of the city of Rome..
And that kind of public policy will feed.
into that kind of anti-Semitic atmosphere..
With that, you can understand some of the Gentile Christians.
coming to the church, they thought that we are better..
And the Jews are being kicked out, right?.
Being discriminated..
And the Gentile Christian would have thought.
that God's plan for the Jews was over,.
it's game over for them..
Now, they ask questions like,.
we are the people of God, we are the Gentiles..
And there were no more Jews, right?.
They may feel that, or imagine that they are better,.
they feel superior, and they became arrogant..
We are Gentiles, we win..
And the Jews are being kicked out,.
and God's plan for them, God's saving grace for them.
is game over..
Because, as you can tell,.
how many Jews were in the early church?.
Not many, right?.
But the foundational population of Christians.
in the early church were Jewish Christians..
And then, as the church spread to other places.
in the Roman Empire, more and more Gentile Christians.
came into the church, and they became the majority,.
and the minority was the Jews..
So, Gentile Christians having some sense of superiority.
with the culture around them, antisemitism..
So, Paul is trying to write in this three chapters,.
nine to 11, trying to address this problem..
Paul argues that, yes, only a handful of Jewish Christians.
in the early church..
But don't forget, they are the foundational people..
Remember, Apostle Paul, Peter, James, and John,.
all these early disciples, they were all Jewish Christians..
They are the foundation..

$^{161}$But yes, only a handful of them believe in Jesus Christ,.
and the majority of the church people are Gentiles..
But Paul argues that this does not mean.
God's plan has failed..
So, 9-6, he says, "It is not as though God's word.
"had failed, for not all who are descended from Israel.
"are Israel.".
What the heck is he talking about, right?.
Israel in two different senses, right?.
So, we can imagine two circles, right?.
Two circles, the big one is the ethnic Israel..
All the people who are descended from the forefathers.
of Israel, Abraham, Isaac, and Jacob,.
they constitute the ethnic Israel at large..
But there's a smaller circle inside, the true Israel..
They believe Jesus as the Messiah,.
so-called the Messianic Jews today..
They embrace Jesus as the Messiah..
They're not waiting for another Messiah..
They know Jesus on the cross and from the empty tomb,.
He is the Messiah..
So, they believe in Jesus,.
and they call on the name of Christ..
So, they constitute the true Israel, the believing Jews..
So, Paul is saying,.
not everyone descended from Israel, ethnic Israel,.
is the true Israel..
And as I told you in the early church,.
some of the most important people were Jewish..
Paul, Peter, James, and John, et cetera..
So, Paul is trying to argue that God's plan for the Jews.
is not over, it's not over..
Because see, there are true Israel, the remaining Israel,.
the believing Jewish people, they are still with us..
Because what matters is not whether they're ethnically.
related to Abraham, Isaac, and Jacob..
What matters is that they are standing in Christ..
10, nine, familiar verse..
If you confess with your mouth, Jesus is Lord,.
and believe in your heart that God raised Jesus.

$^{201}$from the dead, you will be saved..
Verse 12..
See, there's no difference, there's no difference.
between Jew and Gentile..
The same Lord is Lord of all, and richly blesses.
all who call on His name..
So, God's plan for His original chosen people.
is not over..
Paul is saying, look at the true Israel..
Yes, they are the minority, but they are still there..
Right?.
And then, in these three chapters,.
Paul moves on to an analogy,.
and the analogy is the olive tree..
So, here's my amazing artwork here..
As you can tell, there's one olive tree,.
and God has not rejected His people, the ethnic Jews..
There's still a future for them,.
because God's saving grace is for all, right?.
He will execute His plan..
So, here's just this analogy..
There are a few things I'd like to point out.
from this amazing artwork, is that the root of the tree,.
the stem of the tree, is the Old Testament Israel,.
the Jewish people..
Salvation is from them, and Jesus, by the way, was Jewish..
Right?.
He was 100\% Jewish..
The root of the tree is the Old Testament Israel,.
and then there are natural branches..
On the left, the green ones, are the believing Jews..
Paul, Peter, James, and John, and all that..
But some of the natural branches have been broken off..
Lower right..
They have been broken off because of their unbelief,.
because they hardened their heart..
So, those are the unbelieving Jews..
They are being broken off..
However, from another tree, the Gentiles are being.
crafted in..

$^{241}$They are from another tree, from the wild olive shoots,.
taken from the other tree and crafted into the one tree,.
the olive tree..
Now, this famous olive tree speaks of God's.
undeserved grace..
We don't deserve, if we are Gentiles,.
we don't deserve to be crafted in to the tree of salvation..
Remember in Rome, the Gentiles were bragging about.
the fact that we are the majority, and the Jews are gone..
Right?.
It's game over for them..
That's what the Gentiles was thinking..
But this tree analogy is trying to present to us grace..
It's because of grace we are being crafted in..
11, 17..
Paul says, "If some of the natural branches have been.
"broken off, and you, although you are wild olive shoots,.
"have been crafted in among the others,.
"and now share the nourishing sap from the main olive root,.
"do not boast over those natural branches..
"Do not think we are being crafted in so we are better,.
"so we are in a better place in terms of grace.".
No..
Paul says, "Do not boast over those natural branches..
"If you do consider this, you do not support the root,.
"but vice versa, the root supports you.".
Right?.
We share the salvation..
We are being called the people of God..
We are Gentiles..
Don't boast about the fact that, okay, we are the majority,.
and the minority is over..
The Jews is over..
Now, two questions here..
How many olive tree is there?.
How many?.
One..
There's only one..
So there's only one people of God,.
not two peoples of God, right?.

$^{281}$One tree, two kinds of branches,.
the natural and the wildly crafted, right?.
Second, where are you?.
We are the Gentiles..
We are the yellow, right?.
On the right..
We are not the natural branches,.
but by grace, we are being crafted in..
So the olive tree speaks of God's grace to all of us,.
especially the Gentiles..
Us, all of us, every one of us..
We are being called the sons and daughters of God,.
being crafted into the people of God..
It's all because, not because we are better,.
and it's not because they are worse..
It's because of the grace of God..
We enjoy the blessing, the relationship with God,.
and thank God for his saving grace, right?.
And there should be not boasting, but gratitude..
We can come before God as his people..
We should be thankful,.
thankful for what Jesus has done for us..
At the same time, Paul is arguing.
that God's plan for the Jews.
is yet to come..
It's not a game over for the Jews..
God's promised to save many of those unbelieving Jews,.
and it is yet to come, 1125..
Paul said, "I do not want you to be ignorant.
"of this mystery.".
Now he's like opening the lid of the pot,.
and you can see inside,.
there's a spiritual mystery yet to come..
He said, "Brother and sisters,.
"so that you may not be conceited,.
"Israel has experienced a hardening in part.".
Part of the people are being hard-hearted,.
not believing Christ as, or Jesus as the Messiah..
"Until the full number of Gentiles has come in,.
"and so all Israel will be saved.".

$^{321}$All right, this is a very, very difficult verse,.
and one thing you need to know is that,.
first thing is there's a mystery..
There's a future unseen, not unfold yet,.
that is so all Israel will be saved..
Now, one thing you need to understand,.
first, the key word here..
The word all does not mean each and every single Jew..
Now, how do we know that?.
This is obvious because if you check the Old Testament,.
many times it uses the phrase all Israel,.
and many times from the context,.
you know that it does not mean,.
and cannot mean every, each Israelite..
It simply means many of them,.
representative number of them..
And secondly, there's an adverb here,.
so, NIV translate, so, S-O..
ESV has it in this way,.
or in this manner,.
or in the same manner..
Now, the word translation is very important here..
You know, ESV or NIV are right on,.
especially ESV..
In this way, in a similar manner..
Now, what does it mean?.
This adverb, so, or in this way,.
does not mean a straight sequential order..
It does not mean, okay, A happen,.
and then B happen..
So, cross that..
It does not mean the full number of Gentiles.
come in to the church..
We count one billion, two billion, three billion, and seven..
Done, and then all Israel will be saved..
No, it does not mean that,.
because that adverb is not a straightly sequential adverb..
What it means is that, it is trying to show you.
it is in the like manner..
Now, what does it mean?.

$^{361}$Look at the bottom of the slide..
The yellow line is the full number of Gentiles.
being crafted into the church in the human history..
In the last 2,000 plus years,.
Gentiles from here and there,.
here and there, generation to generation,.
are being crafted..
God's saving grace becomes reality.
to each and every one of them..
And the full number of Gentiles will be crafted, right,.
in the entire human history between the first coming.
and the second coming of Christ..
That's the yellow line..
And Paul said, "In this way, in the same manner,.
"all Israel, all true Israel, will be saved.".
So, the green line..
So, it's not sequential, full number, boom, done,.
and then B happen..
It's actually happening the same time, in the same manner..
So, all Israel, true Israel, will be saved.
because God is choosing and saving.
from generation to generation,.
from different times and space,.
different places here and there,.
and save them to be God's true people, true Israel..
So, that's the amazing thing about God's grace..
We, the Gentiles, we can be God's people..
You can be God's people..
Not because we are naturally related,.
not because of blood relation, not because of ethnicity..
We are saved because faith alone in Jesus Christ..
If we confess with our heart and our lips.
that Jesus is Lord, then we are saved..
It is by grace..
It's by faith alone..
Now, no other verses in the entire New Testament.
is more clear than this one, Galatians 3..
It says about Gentile Christians,.
Paul says, "You are all sons of God.
through faith in Christ Jesus..

$^{401}$All of you who were baptized in Christ.
have clothed yourself with Christ..
The image of Christ is on us,.
and there's neither Jew nor Greek.".
It really doesn't matter..
Your ethnicity, slave or free, economic status,.
male or female, gender, for you are in Christ..
You are all one in Christ..
If you belong to Christ, you are the seed of Abraham..
You are the children of Abraham..
Not because of ethnicity, but because of faith..
You have the same faith with Abraham's faith..
So, you are heir according to the promise of God..
Now, by the way, this is not replacement theology, right?.
If you know that kind of word,.
that means church is replacing the original people of God..
No, it's not replacement..
It is fulfillment, fulfill in Christ..
That's what the apostle says in Romans 9 to 11..
God's plan is still going on..
Yes, there's only a handful of Jews believing in Christ,.
but God's plan will continue..
It's a mystery, but it's active..
And God's plan continues to unfold in humanity,.
in human history..
His saving goal, His saving plan.
will continue to be accomplished..
And that is amazing because God's mysterious plan.
for the future, for His people,.
for you and for me continues to unfold..
If we think about that,.
now we understand why Paul ends these three chapters.
with doxology, with praises..
That's the end of chapter 11..
He says, "Oh, the depth of the riches and the wisdom.
"and knowledge of God, how unsearchable His judgments,.
"His paths, the paths of God beyond tracing out.".
There's no way we can comprehend the future..
There's no way we can advise God what He is doing..
We just cannot comprehend.

$^{441}$because the riches of God's wisdom..
"Who has known the mind of the Lord?.
"Who has been His counselor?" 35..
"Who has ever given to God that God should repay him?.
"For from Him, through Him, to Him are all things to Him..
"Be the glory forever, amen.".
He just end this passage with doxology,.
praising God together..
That's the response song here..
He just cannot keep it to himself..
Paul, when he thinks about the plan of God,.
unsearchable plan and mysterious plan,.
he just wanna praise God and continue to serve Him..
Sometimes we don't see the future..
The future seems unknown to us..
Or even we feel confused about what the heck God is doing.
in the future..
That is what the Roman church was struggling with..
And Paul burst out with worship and praise.
and say that, "Yes, we don't quite know.
"where God is leading..
"We don't quite see the future..
"We cannot just fast forward one year from now,.
"six months from now, six years from now,.
"or 10 years from now, what will be unfold?.
"We just don't know yet..
"It's mysterious to us.".
But that does not mean God is not actively doing His purpose..
Remember this line?.
I'm sure you heard of it..
Never place a period where God has placed a comma..
God is still working in the world,.
saving the Gentiles, saving the Jews, saving us..
And indeed, God's working in our life,.
sometimes it seems mysterious..
But the fact that it's mysterious is unknown to us,.
should not and must not limit the possibilities of God..
God's plan in the world continues..
It continues to unfold..
And that's why after this doxology,.

$^{481}$Paul urges us in chapter 12,.
"Therefore," see the connection here?.
"God's plan is unsearchable yet active..
"God's purpose is going to continue..
"It is not a period, it's a comma..
"God continues to orchestrate things for His glory..
"Therefore, I urge you, brothers and sisters,.
"in view of the mercy of God,.
"to offer your whole life, your whole body,.
"as living sacrifices, holy and pleasing to God..
"This is your spiritual act of worship..
"Do not conform any longer to the pattern of this world..
"Do not think according to the culture,.
"according to the world,.
"but be transformed by renewing our mind.".
Dig into the word, study the word of God,.
and see what God is doing in the past,.
what God is doing in Christ..
And we will be amazed that God's plan is mysterious,.
but yet His plan continues to unfold..
Praise God..
"To Him, through Him, to Him are all things..
"To Him be the glory forever.".
That should encourage us to continue to serve our mighty God..
No matter what, offer our whole life as living sacrifice.
to worship Him, to follow Christ,.
not the world, but our Savior, Messiah..
We're gonna sing the last song..
And I pray that this song will be a prayer for all of us..
Thank God for all the times, good or bad,.
continue to serve Him and follow Him..
Let's pray..
Father, once again,.
we acknowledge that we are so limited in our sight..
We can only see a few things..
A few minutes ahead of us..
We can plan, we can worry about the future,.
but yet your word reminds us that you are the active agent..
You continue to work for your glory,.
saving Jews and Gentiles alike,.

$^{521}$saving all of us in Christ..
Yes, life is full of mystery,.
but we know they are all commerce..
You have not put periods in our life journey..
Your plan, your purpose continues to unfold in our life..
In the midst of uncertainty, in the midst of unknowns,.
pray that you will increase our faith.
to see how amazing you are,.
how rich and unsearchable your plan is..
To you be all the glory in Christ, we pray..
\newpage



\section{18:13-19:37}
\label{sec:wrvuGGxg0R0}
\textbf{最終話事權 (列王記下18\_13-19\_37) - 譚文倩莊傳道}
\newline
\newline
連結: \href{https://youtube.com/watch?v=wrvuGGxg0R0}{\texttt{https://youtube.com/watch?v=wrvuGGxg0R0}} ~~~~ 語音日期: 2023-11-28
\newline
\newline
\hyperref[sec:UIZxTjRVxE4]{\small{< < < PREV SERMON < < <}}
~
\hyperref[sec:index_chronic]{\small{[返順時目]}}
~
\hyperref[sec:index_scriptual]{\small{[返順卷目]}}
~
\hyperref[sec:a0q_gDDhWfE]{\small{> > > NEXT SERMON > > >}}
\newline
\newline
$^{1}$感恩節快樂.
不如我們先和大家互相前前後後.
左左右右.
說一聲感恩節快樂.
寫給我的卡和送給我的禮物.
起初我不太明白.
為什麼Mark Sir這麼隆重期待.
這個教務感恩的事情.
或者我自己覺得.
好像有點不好意思很噁心.
但是想了想又覺得很有意思.
很甜蜜.
我在亞洲大學的時候上網.
我自己也很期待收到這些卡和禮物.
今天的經文是比較長的.
我不會和大家讀完.
但如果你不想只聽我說.
願意在家裡時間去探索一下經文.
今天的經文就在《列王記》下18,19章.
小時候我們三兄弟妹經常打架.
特別是我和哥哥打得最多.
雖然我不小.
但在我大哥哥利亞面前.
注定打不過.
但我的性格不會認輸走老.
所以我只有一個戰略.
唯一的戰略就是死忍.
無論打到多痛.
我都不會認輸.
咬緊牙根作戰到底.
打到大家筋疲力盡.
戰事才會結束.
《列王記》下18章13至19章37節.
處於弱勢的猶大希西加王.
對比這個強大的亞述亞拿基勒王.
打起來誰勝誰負.
勝利的一方又是靠什麼能戰勝對方.
今天講道的主題.
Final Say 最終話事權.
我們先低頭祈禱.

$^{41}$主你的名是應當稱頌的.
在列國中受尊崇.
求你祝福你的話語.
祝福人能夠清心的來到你面前.
你的同在.
你的話語.
充滿等候你的人的心.
這樣祈禱.
奉主耶穌基督的聖名祈求.
Amen.
在舊約和古代的文獻記載.
希西加統治猶大國的年代.
最強大的勢力就是亞述.
大同勢是什麼意思呢.
公元前七百年的亞述.
版圖強大到一個地步.
包括有北至今日部分的土耳其.
南至部分的埃及.
東至以色列.
黎巴嫩.
敘利亞.
約旦.
西至伊拉克.
和部分的伊朗.
而希西加的猶大.
只是剛才提到以色列.
裡面的其中一部分.
公元前七百零一年.
亞述王西拿基納.
入侵希西加統治的猶大.
西拿基納公開嘲笑希西加.
並且一次又一次威脅他使用武力.
要猶太人屈服於亞述.
希西加打仗無數.
戰無不勝.
攻陷了很多個國家.
他只是贏.
沒有輸過.
也因為這樣.
也導致他非常驕傲.

$^{81}$也很殘暴.
被他攻陷的國家.
當中不少人被虐待至死.
先說說希西加是個什麼人.
舊約記載.
北國好的君王是一個都沒有.
而南國只有兩個好的君王.
其中一個是約書亞.
另一個就是希西加.
希西加在二十五歲的時候登基.
在神的眼中.
在上帝的眼中.
他行耶和華漢為政的事.
他也效法他祖先大位一切所行.
他是一個很倚靠神.
是一個好的君王.
在他執政的時候.
他落聖旨拆毀了一些廟宇.
拆毀了一些木造的偶像.
因為當時以色列人為了生孩子.
為了求發達.
求平安.
就去拜這些偶像.
他也打碎了當時摩西的銅蛇.
因為整整以色列人都當他是其中一個偶像去參拜.
列王記下十八章五節.
希西加倚靠耶和華以色列的上帝.
在他之前和他之後的猶大列王中.
沒有一個像他的一樣.
希西加的確是一個很好的君王.
但是他有一個致命傷.
就是他不是一個精通戰術.
善於打仗的君王.
列王記下十八章記載.
希西加執政第十四年.
亞述王西拿基納率領一支龐大的軍隊去攻佔猶大城.
征服了整個地區.
希西加自知打不過.
所以他派人去和西拿基納談判.
希望他不要來打他.

$^{121}$列王記下十八章十四至十六節.
他說:我錯了,求你撤退離開我.
凡你罰我的,我必承當.
於是亞述王罰猶大王希西加三百他連得金紙.
希西加把耶和華殿女和皇宮庫房女所有的銀紙都給了他.
那時猶大王希西加將耶和華殿門上的金紙.
和他自己包在柱子上的金紙都挖下來.
給了亞述王.
訂煲要訂煲費.
講和都不是免費.
為免死傷慘重.
希西加金和銀全然奉上.
最值錢的都給了西拿基納.
傾家蕩產.
可憐到在殿裡柱上的金箔都要挖下來給他.
當然了,是在講紀安和大雄的戰鬥.
如何戰鬥,注定輸定硬.
希西加向亞述王進貢.
想著可以息事寧人.
想著可以保住性命.
誰不知西拿基納王反口不認賬.
結果埃及和希西加伸出援手.
想聯盟能夠強大一些去對抗亞述.
誰不知亞述王西拿基納收到風.
令他大衛震怒.
他派將軍來到猶太國城門口.
鬧爆希西加.
帶著壓倒性的姿態.
要求他打開耶路撒冷的城門.
要求猶太人投降.
於是希西加派三個官員出來迎接.
將軍對他們三人說.
你回去跟君王希西加說.
你不要以為你們可以靠埃及.
以為跟他們聯盟就可以救我打.
你這樣想,我就先打殘埃及.
然後再打你.
你們也不要說什麼信靠上帝.
不是你們之前已經毀壞了電雨.
打爛了偶像嗎.

$^{161}$還拜什麼.
還依靠什麼鬼神.
再說,我給你兩千匹馬.
你們也沒有人有本事能騎.
我的軍隊裡最弱雞的那個也不怕你.
順便說,你有沒有想過.
可能是你所信的,所依靠的上帝派我來毀滅你呢.
《列王記》下十八章二十六,二十七節.
希西加的兒子以利亞徑.
舍伯拿和約亞對將軍說.
求你用雅蘭話對僕人說.
因為我們聽得懂.
不要用猶大話對我們說.
免得傳到城牆上百姓的耳中.
將軍對他們說.
我主差遣我來.
豈是單對你和你的主人說這些話嗎.
不也是對這些坐在城牆上.
要與你們一同吃自己的糞.
喝自己的尿.
人說的嗎.
哇,用不著去到這麼盡.
吃自己的糞.
喝自己的尿.
我看的時候就想起我們廣東話.
一句罵人的句子.
有人問我,聖經是不是真的這樣寫的.
聖經真的是這樣寫的.
英文是說,Eat your own dung.
Drink your own urine.
原文是說吃自己的糞便.
鴨子的糞便和尿.
有興趣可以自己回去看看.
亞述王派來的將軍這樣說.
當然他是要侮辱西西加.
他更加是要告訴所有人.
如果耶路撒冷長期被困的時候.
他們會斷水斷糧.
到時所有人就要吃自己的糞.
喝自己的尿.

$^{201}$而為什麼西西加的人.
要求亞述王的將軍去講亞蘭文呢.
原因是亞蘭文是受過教育的猶大皇家官員.
才懂得聽的.
但是城牆裡一般的居民根本聽不懂.
亞述的將軍堅持用希伯來文跟他們溝通.
就是要所有人都聽得懂.
要他們聽得懂.
你們那個君王.
西西加其實是沒什麼用的.
他們保護不了你們.
你們所依靠的上帝也沒什麼用.
進貢也搞不定.
打也打不夠.
還要被人嘲笑一番.
敵人已經來到城門口了.
還可以怎樣呢.
西西加收到亞述將軍帶來的信息之後.
當然他就坐不安.
吃不下.
帶著生意失敗的心情.
他就撕裂衣服.
披上麻布進入上主的聖殿.
他來到以賽亞先知的面前.
申訴亞述王派來將軍所講的.
特別是藐視上帝的說話.
西西加求以賽亞的上帝施行拯救.
伸張正義懲罰亞述.
以賽亞先知當然是很細心地去聆聽.
去安慰他.
提醒他神與他同在.
答應上帝會親自對付亞述.
在上帝出手之前.
這個亞述王還未知死.
他再一次派人去找西西加.
作出最後的通牒.
威脅他.
提醒他.
以前很多以為救他打的人.
全部都被他消滅了.

$^{241}$你又怎會是例外呢?.
列王幾下十九章十至十三節.
你們要對猶大王西西加如此說.
不要聽你所倚靠的神欺控你說.
耶路撒冷必不交在亞述王的手中.
看那你總聽說亞述諸王.
向列國所行的進行滅絕.
難道你能倖免嗎?.
我祖先所毀滅的.
就是哥散哈那尼西和提拉薩的伊甸人.
這些國的神明何曾拯救他們呢?.
哈瑪的王.
阿爾拔的王.
西伐胡音城的王.
希那和爾戶的王.
都在那裡呢?.
如果你是西西加.
來到這個時候你會怎樣做呢?.
講和又搞不定.
打又打不夠.
傾家蕩產被人侮辱.
敵人處處逼人.
水浸眼眉不單止.
敵人也一次又一次侮辱.
你所信靠的上帝.
今時今日我們身處的這個年代.
我想每一天我們都看到不同的新聞.
說到蘇聯打烏克蘭.
哈瑪斯對抗以色列.
西方國家和亞洲國家.
不停地互相制裁.
最近我接二連三收到.
不同的朋友患病的消息.
又聽到身邊的朋友失業.
但還未供完子女讀大學.
也有未畢業的大學生和父母說.
畢業之後未必找到工作.
因為現在的經濟不太好.
有家長投資529教育基金.
縮水了一半.

$^{281}$一個不幸還買了港股.
我心想,慘了.
怎麼辦.
有什麼可以做.
怎樣作戰到底.
面對接種以來的困難攻擊.
我相信希西加和你和我一樣.
面對這麼多困難的時候.
我們都會心情很煩亂.
會很擔心.
我們會很害怕.
我們會胡思亂想.
我們晚上睡得不好.
血壓高.
想想繼續下去怎麼辦.
獵王幾下十九章.
十五至十六節.
希西加向耶和華禱告說.
坐在基路柏之上.
耶和華以色列的神.
唯有你是地上萬國的神.
你創造了天和地.
耶和華求你則以而聽.
希西加的確似乎走投無路.
他的確是一個不精通戰術.
善於打仗的軍王.
但他有一樣做得很對.
就是他沒有失去重心.
仍然很清楚誰是上帝.
他所倚靠的是誰.
所以在這個絕望的境況裡.
他去尋求神.
因為他知道唯有上帝才有能力施行拯救.
唯有上帝出手能夠叫萬國.
包括亞述.
去知道誰才是真正的軍王.
最近有什麼纏繞著你.
讓你不安.
讓你睡不著.
讓你血眼飄星.

$^{321}$在這些時候.
你尋求什麼幫助.
上帝在這些事情上有沒有分.
有沒有說.
我最近有個朋友的癌症.
第三次復發.
每次回來的時候.
我都會去探望他.
上星期他在家裡坐著的時候.
我很有福氣.
吃到他煮的飯.
我喝了三碗湯.
很久沒吃過.
見完他之後.
這幾天我不停回味.
我們那幾個小時的相聚.
閒話加上很享受.
想起他的時候.
我有很莫名的感恩.
第三次的復發.
我想起他之前第一次有事.
第二次復發.
到這次第三次.
他很平安.
很有趣.
我自己想起初初見他的時候.
第一第二次的時候.
他很不OK.
但這次他很OK.
他還跟我分享神跟他同在.
他說神很奇妙地給了他一個平安.
令他很安然.
他今時今日.
每一天都覺得是神給他的禮物.
每一天都很珍惜地度過.
他很感恩.
因為神給他在這個癌症裡面.
面對這個病裡面.
他能夠有平安.
我想起我們祈禱的時候.

$^{361}$很多時候我們在想我們祈禱.
不只是說你們.
我也試過將我的禱文寫下來.
發覺求的東西真的佔大多數.
只在祈求神應運我的一些祈禱.
也想神盡快答應.
當然快點幫我搞定事情.
但其實有時候我看這段經文的時候.
我自己在想.
有時候我們祈禱沒有再進一步地想.
除了我求神應運我的祈禱之外.
在這個回應裡面.
神的名字是被高舉的.
神在他的回應裡面.
對我的回應裡面是得著榮耀.
就像我這個朋友.
他的祈禱.
當然他的祈禱不會說不用醫好.
他當然會求神.
希望神去醫治他的病.
但他更加想神的名字.
即是被榮耀.
接觸他的人能夠看到.
神在他身上所做的.
神給他的平安.
希西加向上帝求救之餘.
他更加期望上帝出手的時候.
萬國都能夠意識到.
上帝就是掌管萬物的上帝.
只有上帝在他的困難裡面.
有能力.
有這個最後的說話.
有這個最終話事權.
結局是怎樣.
是在神的手裡面.
聶王紀下十九章三十四節.
因我為自己的緣故.
又為我僕人大衛的緣故.
必保護拯救者成.
根據聖經記載.

$^{401}$一到晚上的時候.
上主就派天使進入亞述軍營.
殺了十八萬五千個亞述士兵.
亞述王西拿基納和亞述人.
早上起來.
看到獅王遍野.
當然是嚇得要命.
逃也逃不及.
亞拿基納逃到尼尼尾城.
之後亞述王西拿基納的兩個兒子.
殺了他.
然後他兩個兒子再逃亡.
他的另一個兒子.
以撒哈頓成為新一任的亞述王.
不是像我這樣.
死忍作戰到底.
也不是奉上金和銀講和.
也不是打算輸定了.
算了我投降.
獵王幾下十八十九章.
希西加的故事.
對你有什麼提醒.
詩篇四十六篇第十節.
你們要休息.
要知道我是上帝.
我必在列國中受尊崇.
在全地也受尊崇.
詩篇四十六篇.
鼓勵我們在艱難的時候.
要抱著希望和信任.
唯有神在我們所面對的事情裡.
有主權.
有最後的說話.
有最終的決定權.
願意我們所做的.
我們所祈求的.
我們所面對的困難裡.
神是有份的.
神也是有說話的.
願意神在我們的生命裡得著榮耀.

$^{441}$我很想和大家一起去重讀詩篇第四十六篇.
上帝是我們的避難所.
是我們的力量.
是我們在患難中隨時的幫助.
所以地需改變.
山需搖動到海深.
其中的水需澎湃翻騰.
山水因海漲而顫抖.
我們也不害怕.
有一道河.
這河的分叉使神的城寬起.
這神就是至高者居住的聖所.
上帝在其中.
城必不動搖.
到天一亮.
上帝必幫助這城.
萬邦軒揚.
國度動搖.
上帝出聲.
地就融化.
萬軍之耶和華與我們同在.
雅各的上帝是我們的避難所.
你們來看耶和華的作為.
看他使地怎樣荒涼.
他止息戰爭直到地極.
他折弓斷槍.
把戰車焚燒在火中.
我讀完第十一節.
你們要休息.
要知道我是上帝.
我必在列國中受尊崇.
在全地也受尊崇.
萬軍之耶和華與我們同在.
雅各的上帝是我們的避難所.
這裡有三節經文.
第三節,第七節,十一節.
都有一個字.
「細拉」.
意思是停一停.
Pause.

$^{481}$待會我們會有一首回應的詩歌.
就是想在這個詩歌.
在唱在彈奏的時候.
我們安靜下來.
我想給你們一些時間.
在座位上思想的時候.
有些祈禱的時候.
如果你現在面對一些困難.
面對一些事情.
你需要回到神的面前禱告.
我想給你們這個時間.
你們可以在神面前安靜禱告.
我先跟大家做一個禱告.
統領天使軍隊的神.
萬軍之耶和華.
你必要在列國中受尊崇.
主啊,你是我們的避難所.
你是我們的力量.
是我們患難中隨時的幫助.
奉主耶穌基督的聖名祈求.
Amen.
《天使與神》 詞:陳小霞 曲:陳小霞.
Oh shall I fear.
Shall I fear.
The God who covers.
Angel of peace.
Has sent encampments.
Around me.
Whom shall I fear.
I shall not fear.
The air by day.
Nor shall I fear.
The terror by night.
The God who covers.
Angel of peace.
Has sent encampments.
Around me.
Whom shall I fear.
The God who covers.
Angel of peace.

$^{521}$Has sent encampments.
Around me.
Whom shall I fear.
I shall not fear.
The air by day.
Nor shall I fear.
The terror by night.
The God who covers.
Angel of peace.
Has sent encampments.
Around me.
Whom shall I fear.
The God who covers.
Angel of peace.
Has sent encampments.
Around me.
Whom shall I fear.
我們一起低頭.
主在這個感恩節的節日裡.
再一次提醒我們.
我們要向你感恩.
無論在什麼境況裡.
主我們願意感恩.
願意高舉你的聖名.
因為你是值得稱頌.
你的名是值得被稱頌.
主願意你在我們生命裡.
我們每一個人的生命裡.
去得著當得的榮耀.
困難也好.
疾病也好.
我們知道你在我們的生命裡.
要彰顯你自己的榮耀.
願統領天使軍隊的神.
為你征戰.
主的平安與你同在.
主的名被高舉.
奉主耶穌基督的聖名祈求.
Amen.
\newpage



\section{3:15}
\label{sec:a0q_gDDhWfE}
\textbf{The Advent of Hope (Genesis 3\_15) - George Crawley}
\newline
\newline
連結: \href{https://youtube.com/watch?v=a0q-gDDhWfE}{\texttt{https://youtube.com/watch?v=a0q-gDDhWfE}} ~~~~ 語音日期: 2023-12-05
\newline
\newline
\hyperref[sec:wrvuGGxg0R0]{\small{< < < PREV SERMON < < <}}
~
\hyperref[sec:index_chronic]{\small{[返順時目]}}
~
\hyperref[sec:index_scriptual]{\small{[返順卷目]}}
~
\hyperref[sec:NrY2oymR5Fc]{\small{> > > NEXT SERMON > > >}}
\newline
\newline
$^{1}$Thank you..
Well, it's a blessing and a pleasure.
to be speaking with you guys again today..
So thank you all for allowing me to be here..
Today is the 3rd of December,.
and so if you weren't aware,.
allow me to be the first one to tell you.
that it is officially Christmas season..
Yes, it's the most wonderful time of the year..
And when I was putting this together,.
I thought about singing that line,.
but nobody really wants that..
Nobody wants to hear me sing..
That wouldn't be good for anybody..
But Christmas time, it's the time for Christmas trees,.
Christmas cookies, and Christmas decorations..
It's a time to decorate houses, give gifts,.
wear matching pajamas,.
and one of my favorite parts of the Christmas season.
is that it's the only time of the year.
I can find eggnog at the grocery store..
Are there any eggnog-like lovers in the audience.
by show of hands?.
Oh, only a couple?.
All right, that's fine..
All right, I'll have to scratch the next five minutes.
of what I was gonna talk about..
No, I'm just kidding..
But something that's more important than eggnog,.
pajamas, cookies, is that this is the time of year.
we celebrate the birth of our Savior and Messiah,.
Jesus Christ..
For many Christians, it's a time of Advent.
where we recognize and look back,.
celebrating the first coming of our Lord into this world,.
and we look forward in anticipation.
to His second coming into this world..
So for many reasons, Christmas time.
should be a time of great cheer,.
a time of joy, and a time of hope..

$^{41}$But unfortunately for many people, this isn't the case..
For many people, the holiday season.
is actually a time of stress, a time of sadness..
So I did a little looking into this,.
and in a recent study done by the National Alliance.
on Mental Illness, researchers found.
that three out of five Americans.
feel their mental health is negatively impacted.
during this Christmas season..
Another study done by the American Psychological Association.
finds that 88\% of people in America.
struggle with stress during the holidays..
And there were many other studies done,.
and maybe some of you, people you know,.
have experiences that lead us to this conclusion.
that for some people, the holidays.
are actually a very tough, difficult time..
But I wanna talk a little about why is it..
Why is it that during what is supposed to be.
the most wonderful time of the year,.
people seem to feel anything but wonderful?.
Why during this season, which is a time.
that should be full of good tidings of great joy,.
are there so many people who don't feel that joy,.
who don't feel that cheer, and don't have that hope?.
For some, this could be caused by financial concern,.
worrying about making ends meet.
while you're finding gifts for family and friends,.
while you plan for the year ahead..
That could be a source of stress..
For other people, it's a time of loneliness..
The distance between us and our family and friends,.
it feels a little extra distant..
It feels a little extra heavy.
when we're not together during the holidays..
As you mentioned, the weather also plays a part.
in this kind of dampened mood..
It gets colder, nights are longer, days are shorter..
Here in Chicago, the sun doesn't shine past 4 p.m..
And that's honestly a little bit depressing.

$^{81}$because it's only getting colder and darker from here..
And for some people, the cold weather.
can bring about aches and pains,.
kind of sort of triggers this flare up,.
you know, of maybe arthritis or something in your body..
So the cold and the darkness certainly has an effect.
on people's mental health..
But while this holiday hopelessness.
can be caused by a number of varying factors,.
there's ultimately one reason behind it all,.
one source behind this lack of joy, peace, and hope..
And it's sin, it's the reality that we live.
in a fallen world, a world stained by sin..
Sin is the source of all human suffering,.
sadness, and sorrow..
And this is the unfortunate reality.
that we all face on a daily basis..
We live in a sin-stained world,.
and in this world, even the most beautiful things.
that should bring us the most joy and wonder.
can be corrupted by sin..
That's why it's no surprise to see that.
during a time of year that should be filled.
with joy, peace, and happiness,.
many are suffering..
We see something beautiful corrupted by sin,.
pain instead of peace, strife instead of serenity,.
hopelessness instead of happiness..
Yet in the same way that we see this exchange.
of pain in the place of peace,.
something beautiful corrupted by sin,.
it's not just something that affects us..
We can actually trace this pattern all the way back.
through the history of humankind to the beginning.
when we see that exchange take place..
But in the same way that we can trace the problem,.
sin, back to the beginning,.
we can also find the solution to that problem.
from the very beginning..
And so that's something that I wanna talk about today.

$^{121}$is how God from the very beginning.
has offered his people hope..
And it was this hope, this promise of God.
that the people of God have looked to,.
have held onto for thousands of years,.
and it is this hope that has kept God's people.
even in the midst of great suffering, troubles, and trials..
And so it's this very same hope that we can look to today,.
whether it's Christmas time,.
whether it's after Christmas time, in between,.
we look to this hope as a solution.
for our own sorrows and troubles..
And so I wanna look together at Genesis chapter three..
So we're gonna get into the text.
and I can stop just ranting..
Well, actually, no, I'm gonna talk a little bit more.
then we'll get into the text.
because I wanna set the scene a little bit, right?.
I wanna give you guys some context.
for what is going on when we see here Genesis chapter three..
So allow me to set the scene..
So if you want, this is your thing,.
close your eyes, imagine it..
You don't have to if you don't want to,.
but just picture it..
There's total darkness, nothingness..
The world is empty and blank..
It's void of anything..
And then a voice calls out into this nothingness,.
let there be light..
And there was light..
And then God creates the heavens and the earth..
He puts all the galaxies, the solar systems,.
the stars in their place, the sun and the moon..
He creates the earth and he takes all the water on the earth.
and he divides it and separates it and brings up the land..
Well, then he plants seeds, trees,.
fruit bearing trees on the land..
He fills the land with all kind of animal life, plant life..
He fills the water with all kind of fish and wildlife..

$^{161}$And you see God is building, designing step by step,.
day by day, this beautiful creation..
And all of this creation that he made,.
that he powerfully, wisely, perfectly spoke into existence.
is doing something..
The sun, the moon, the stars, the earth,.
the animals, the trees, the mountains,.
even the smallest grain of sand on the beach.
is declaring the glory of God..
And we see all this happening,.
declaring the infinite glory of God as we go day one,.
day two, day three..
And it's as if it's all leading up to a point, a climax,.
the pinnacle of creation on day six when God creates man..
And so God takes from the earth the dirt and he forms a man..
He forms him and designs him, the scripture says,.
and he breathes into him the breath of life..
And in that we see something that God creates mankind.
in his own image..
Nothing in all of creation can say that they are made.
or it is made in the image of God..
But this one man is made in his image.
and all of his descendants made in the image of God..
And so all of creation, all this perfection,.
all this beauty we just imagined, God gives to the man..
And even amongst the perfect world,.
God sets aside a perfect garden.
and he places the man in the garden to keep it,.
to tend it, to live off of the fruit.
and all the bounty that God has given the man,.
everything he could possibly need..
And then we see as the man is giving names.
to all the creatures, all the animals,.
he has a realization because when he,.
I'm imagining this, when he names the giraffe,.
he sees a daddy giraffe and a mommy giraffe..
And he names a tiger, he sees the boy tiger.
and the girl tiger..
Well, he realizes that all these animals.
have their pair and their mate, but I'm alone..

$^{201}$That as good as these animals are,.
there's none that matches me,.
there's none that is perfect for me..
I don't have a companion..
And so God creates for him, he puts him to sleep.
and takes from his own side, his own flesh and bones.
and creates a perfect, suitable companion for this man..
He creates the first woman and he gives her,.
like a husband giving his daughter in marriage,.
he brings her and gives her to the man..
And the man sings out this praise,.
I think you could say it's the first psalm.
that we see in the scriptures when he thanks God.
and praises him for this woman he's given him,.
this helper, flesh of my flesh, bone of my bone..
And we see that after this happens,.
it's almost as if God stops and he looks at the man.
and the woman and all the animals and all the trees.
and the plants and the waters and the stars.
and the sun and the moon, and he declares.
that it is very good..
Absolute paradise, total perfection,.
no sin, no death, no pain, no toil, no suffering or sadness..
In the garden of Eden, Adam and Eve had everything.
they could ever need and everything they could ever want..
And this is where we begin chapter three..
So I'll be reading from Genesis chapter three,.
verses one through 15, if you wanna follow along..
Now the serpent was more crafty.
than any of the other beasts of the field.
that the Lord God had made..
He said to the woman, did God actually say.
you shall not eat of any tree in the garden?.
And so I wanna stop a few places here.
as I read through this passage to point out.
some of these things that I think we can really.
gain a lot of wisdom from in seeing.
how did this first temptation take place?.
Because I believe that the tactics of the enemy.
have not really changed..

$^{241}$There's nothing new under the sun..
And so this form of temptation that we see.
come against the man and the woman and their response,.
we can experience those same things..
And so we should learn from this..
And so the first thing we see, the very first temptation,.
this first inkling of sin entering the world,.
it is a questioning and a doubting the word of God..
When the serpent says to Eve, did God actually say this?.
What he's doing is he's questioning,.
creating doubt within the woman's mind.
of whether or not God's word is trustworthy,.
whether or not God is credible and reliable or authoritative..
Did he actually say that?.
Do you really have to follow that?.
Are you sure you heard that correctly?.
There's that seed of doubt..
And the woman responds and rightly corrects him..
And she says, we may eat of the fruit.
of the trees in the garden, but God said,.
you shall not eat of the fruit of the tree.
that is in the midst of the garden..
Neither shall you touch it lest you die..
And so two things we see in Eve's response.
that are a little curious is that.
when you look at the previous chapter,.
God never said they couldn't touch the fruit.
or touch the tree..
Somehow Adam and Eve, I don't know if they misheard.
or added it maybe for their own protection,.
but they added something to what God had said..
What's more, we see that Satan,.
and this is, I believe, still one of his tactics today..
He flips what God says..
And so when we see God spoke to Adam that all the trees,.
all the fruit you can eat from, but not this one..
What Satan does is he minimizes what God has given Adam.
and he maximizes that restriction..
It says, did God say you couldn't eat from any tree?.
And so what Satan does is he wants the woman.

$^{281}$and the man to focus on the one thing God said.
wasn't for them that they shouldn't have..
Instead of focusing on all that God had given them,.
I'm somebody who loves fruit..
You can ask my wife, you can ask Rachel..
I love fruit, I'm always looking for new fruit..
There's so many fruit in the world,.
and they had all of it, and I can't believe it..
They allow the devil to tempt them to only focus.
on that one tree they couldn't have..
But as much as I have a hard time understanding.
why they would do that, we do the same thing..
God has given us so much goodness in the world,.
and yet we're often tempted to focus on those things.
that we can't have because there's something in human pride.
that doesn't like God telling us what we can't do,.
what we can't say, what we can't have..
We don't like that..
And so instead of rejoicing and enjoying.
all the good things God has given us when we sin,.
we choose to instead enjoy the few things.
that God says, this isn't good for you..
I don't want this for you, don't have that..
And we rebel against God just as Adam and Eve did..
And so my first thing is that we should learn from this..
The first thing we can gain from this.
is to focus on the good God has given us,.
to overcome the temptation towards the bad.
that is not good for us..
I'll continue from verse four..
But the serpent said to the woman,.
you will not surely die, for God knows.
that when you eat of it, your eyes will be opened,.
and you will be like God, knowing good and evil..
Again, further doubting, now outright denying..
At first it started with questioning,.
did God really say that?.
To outright denying, that's not true..
Okay, God said that, but it's not true..
And I think today we should also be aware of.

$^{321}$and on guard against similar temptations.
to mistrust or outright deny the word of God,.
that he's graciously given us for our good..
So the temptation comes in verse six,.
says that when the woman saw the tree was good for food,.
and that it was a delight to the eyes,.
and that the tree was to be desired to make one wise,.
she took of its fruit and ate,.
and she also gave some to her husband who was with her,.
and he ate..
Then the eyes of both were opened,.
and they knew that they were naked,.
and they sewed fig leaves together.
and made themselves loincloths..
Adam and Eve had been naked from the beginning..
Nothing new happened except this loss of innocence,.
where they realized they were naked..
But the thing with that is that.
there was nothing wrong with nakedness..
God had made them naked, and it was good..
But we see what happens here.
when they took and ate of that fruit,.
their morality, their minds, their conscience was corrupted..
They had lost their sense of innocence..
And we see that entrance of sin,.
guilt, and shame into the world,.
to the point they tried to cover themselves with leaves,.
which you'll see later in the chapter,.
it was not sufficient..
That's another sermon that I could,.
there's so many sermons you can get from this passage,.
but just remember that..
They tried to cover themselves, but it wasn't sufficient..
So, verse eight..
They heard the sound of the Lord God.
walking in the garden in the cool of the day,.
and the man and his wife hid themselves.
from the presence of the Lord God.
among the trees of the garden..
But the Lord God called to the man and said to him,.

$^{361}$"Where are you?".
Now, I wanna make a point here.
that when God asked this question,.
he wasn't asking questions the way we ask questions..
We ask questions because we wanna learn something..
There's some piece of knowledge we don't have,.
and we want to get it..
That's not how it works with God..
God already knows everything..
There was never a time where God learned something,.
because God already knows everything..
So, when he asked Adam, "Where are you?".
It wasn't because he didn't know where Adam was..
It's because he wanted to prick the conscience of Adam,.
so Adam would reflect, "What am I doing?.
"Where am I?.
"I'm hiding from God.".
Which doesn't even, how can you hide from the God.
who knows all things?.
Whose eyes are everywhere, who sees everything?.
That was the point of this question,.
was to prick Adam's conscience..
And Adam says this, "I heard the sound of you in the garden,.
"and I was afraid, because I was naked and I hid myself.".
So Adam hid..
He hid himself from God because of his guilt and shame,.
which we already know..
I mentioned before, there's no logic in hiding from God..
But we tend to do the same thing.
in a slightly different way,.
but it's the same thing at the heart of it,.
that when we sin against God and our conscience is pricked,.
we tend to feel that we ought to distance ourselves from God..
Maybe we feel like, "Whatever I did this weekend,.
"I know it wasn't right..
"I shouldn't go to church..
"I can't fellowship..
"I'm too guilty..
"I'm not right the way others are right.".
Or we feel like, "I can't pray right now..

$^{401}$"I can't go to God..
"He's angry with me, 'cause he knows what I just did,.
"what I just said.".
Or maybe, "How could I pick up the Bible.
"and read it right now?.
"I'm such a hypocrite..
"I just sinned, I can't.".
And so our sin causes us to hide from God in a sense..
We turn from God..
Well, on the flip side, how it should be is that.
when we sin, we ought to go to God,.
because he's the one who forgives us.
and cleanses us and makes us right..
Sin should not separate us from God.
the way it did with Adam and Eve..
We should go to God, confess our sins,.
and receive forgiveness and restoration..
We're not gonna find that by hiding from God.
in many of the ways we try to..
I'm gonna continue..
Verse 11..
God said, "Who told you that you were naked?.
"Have you eaten of the tree.
"of which I commanded you not to eat?".
Now, this is an interesting part.
where we see something enter the world.
that we're also all still guilty of,.
called the blame game..
Verse 12..
"The man said, 'The woman you gave me to be with me..
"She gave me the fruit of the tree, and I ate.'".
So Adam not only puts the blame on the woman,.
well, she gave it to me,.
he also kind of puts the blame towards God,.
because it's the woman that you gave me..
He's already pointing fingers..
He's already passing the blame..
Then what does the woman say?.
"The Lord God said to the woman,.
"What is this that you have done?".

$^{441}$The woman said, "The serpent deceived me, and I ate.".
So we already see this lack of accountability.
and responsibility for sin..
Instead of confessing and admitting guilt,.
it's, well, it's the woman's fault..
Oh, it's the serpent's fault..
Oh, God, it's your fault..
You put us here..
And we still struggle with that..
It's difficult to take accountability when we sin..
It's easier to think that, well, it's this person's fault..
And, you know, no, I wasn't being rude..
They just took it the wrong way,.
and they shouldn't have started with me..
You know, we kind of steal all these things.
that Adam and Eve did and stumbled with here..
They're still stumbling blocks for us today..
We can learn a lot from reading this and from studying it.
and comparing what they did to what we should do..
Okay, but I'll continue..
The serpent deceived me, and I ate..
I do want to note that after Adam and Eve sinned,.
it's important to see this, because in this next passage,.
we see the mercy of God, the justice of God,.
but also the mercy of God,.
because Adam and Eve had just been created by God.
from nothing, given everything..
He gave them everything,.
and they immediately turned and rebelled.
and committed this heinous, high-handed sin.
against Almighty God, against their creator..
And yet, God first turns to the serpent.
and punishes the serpent,.
not the man and the woman who had sinned against him..
So even in that, you see some of God's mercy..
God could have just started over from that second.
and created a new earth and a new Adam and a new Eve,.
but he was merciful..
And the worst of this punishment is on the serpent..
Verse 14, the Lord God said to the serpent,.

$^{481}$because you have done this,.
cursed are you above all livestock.
and above all beasts of the field..
On your belly you shall go,.
and dust you shall eat all the days of your life..
I will put enmity between you and the woman,.
and between your offspring and her offspring..
He shall bruise your head, and you shall bruise his heel..
So surely this was the darkest day in human history..
It was the entrance of sin into the world,.
which at that time was a perfect world,.
an untarnished, unblemished world..
Paradise was lost..
We went from perfection to corruption..
Innocence exchanged for guilt and shame..
Purity exchanged for sin..
Hope exchanged for hopelessness..
Life exchanged for death..
And even as bleak at times,.
depressing the current state of the world may seem,.
when you see all the terrible things happening.
in our city, in our country,.
surely you see it around the globe..
I can guarantee you,.
nothing we have today is as dark as it was on that day,.
when darkness first entered the world..
But as is often said,.
it's always darkest before the dawn..
You see, in the midst of this very darkness,.
we find the advent, the emergence of light,.
a dawn of hope,.
a promise given by God himself,.
in which his people have held onto for thousands of years,.
a promise of salvation, redemption,.
restoration, and justice for mankind..
And we find this promise, ironically, in a curse..
In verse 15, the Lord God said to the serpent,.
"Because you have done this,.
"cursed are you above all livestock.
"and above all beasts of the field..

$^{521}$"On your belly you shall go,.
"and dust you shall eat all the days of your life,".
which is almost a metaphor,.
a picture of utter humiliation, lowliness, defeat..
"I will put enmity," which is conflict, warfare,.
violence, "between you and the woman.
"and between your offspring and her offspring..
"He shall bruise your head,.
"and you shall bruise his heel.".
So this is the curse that God gives to the serpent..
But even in this curse, we find a promise.
when God says, interestingly,.
"I will put enmity between you and the woman,.
"between your offspring and her offspring..
"He shall bruise your head.".
There's a pronoun there of an offspring of the woman..
"He shall bruise your head,.
"and you shall bruise his heel.".
And so I wanna talk a little bit about that last part..
He shall bruise your head, you shall bruise his heel..
It's referring to the offspring or the seed of the woman..
In Hebrew, this is the word zara..
Its meaning is like an offspring, a descendant,.
a child, a literal seed..
And in the Bible, when we see this word used in this way,.
it always refers to a man, the seed of Abraham,.
the seed of David..
But in this one instance,.
of all the places I could find in the Bible,.
the seed describes a woman..
And this is because the seed is in,.
and it comes from the man..
It doesn't come from the woman..
So it's curious that God would talk.
about the seed of a woman..
This wouldn't really make sense normally,.
but I think this is significant..
Because in this promise,.
the seed of the woman is referring to a certain someone..
And if you can think of who it is,.

$^{561}$go ahead and just shout it out..
Exactly, exactly..
'Cause there's one person, one child,.
in all of human history,.
who was born without a human father..
So literally, Jesus, the Messiah,.
born of the Virgin Mary,.
could only be the seed of a woman..
Jesus Christ, the Son of God,.
the seed of the woman,.
promised by God from the very beginning.
that he would be the one to bruise,.
to crush the head of the serpent..
God says, "He will bruise your head,.
"you will bruise his heel.".
And so you can compare the two..
To bruise someone's heel,.
it's like a snake bite to the back..
You don't see it coming..
It's kind of an underhanded, deceptive attack..
But if a serpent were to come on the stage,.
and I don't know if we have snakes in Chicago..
I know we have rats..
Maybe if we had some more snakes, we'd have less rats..
But if a snake came on stage and bit me in the heel,.
it would hurt, it would not be great,.
but it wouldn't kill me..
It's not a fatal blow..
But if I turned around and stomped on that snake's head,.
that's it..
It's a fatal blow..
It's a death strike..
A complete defeat for that snake..
And so in that curse, there's also a promise.
that this one figure, this male offspring of the woman,.
this savior, would be bruised on the heel,.
would suffer some temporary defeat..
But in that defeat, he would ultimately defeat the serpent..
And so this passage right here is known.
as the Proto-Evangelium, which just means the first gospel..

$^{601}$It is the first gospel message in all of the Bible..
The first glimpse of hope, promise of redemption,.
fulfilled in Jesus Christ's death and resurrection..
Because in death, Jesus was rejected and scorned.
by his own people, betrayed by them..
Betrayed even by a disciple who he loved as a brother..
On the way to the cross, he was abandoned.
by the rest of his disciples..
In that crucifixion, in that rejection and abandonment,.
Jesus suffered the serpent's bite to the heel..
Yet in his resurrection, he crushed the head of the serpent.
declaring victory over sin, death, and the grave..
And so this was the hope, this savior figure to come,.
was the hope of Adam and Eve, that one of their descendants.
would one day crush the head of the serpent..
That this figure, this messiah, the savior,.
would defeat Satan and rid humanity.
of his evil schemes once and for all..
It is this savior who is the messiah of Israel,.
the one in which the people of Israel.
hoped in for thousands of years..
The one whose birth we celebrate every year at Christmas..
He is the one upon which all of our hopes rest..
All of God's people rest their hope in this one savior..
He is the solution to all of our problems..
The ones that, the sin problem that has plagued.
all of mankind throughout history and plagues us today..
Our problem is sin, it's devastating effects on our bodies,.
our families, every aspect of our life is marred by sin..
But just as it is written, she will give birth to a son.
and you shall call his name Jesus..
For he will save his people from their sins..
This is the hope that we must cling on to..
The same hope that Adam and Eve, even after the fall,.
even after such a devastating loss of perfection,.
this was the hope they held to..
The same hope we see throughout the scriptures..
So many of God's faithful held on to this hope.
of a savior, of a redeemer..
One example I wanna look at from the book of Job..

$^{641}$A great example, Job was a man who lost nearly everything,.
who by worldly standards had nothing to live for..
His possessions, gone, his children, dead,.
his own friends accused him of evil..
His wife even asked him, you know,.
why don't you just curse God and die already?.
Yet in spite of all his suffering and anguish,.
Job could confidently declare this..
I know that my redeemer lives and that in the end,.
he will stand on the earth..
And after my skin has been destroyed,.
yet in my flesh, I will see God..
I myself will see him with my own eyes,.
I and not another, how my heart yearns within me..
You see, despite Job's sorrow and suffering,.
he was able to hold on to hope in God, his redeemer..
And so I ask, can you, can we say with Job,.
I know my redeemer lives in the midst of our suffering,.
my redeemer lives..
In the midst of frustration, difficulty,.
disappointment with family, with friends,.
my redeemer lives..
In the midst of pain and anguish, physical ailment,.
my redeemer lives..
When you lose a loved one, my redeemer lives..
You see, that is our hope, that Jesus Christ,.
the seed of the woman, he lives, he's our hope..
It's this very same hope,.
which has been held onto for millennia..
The prophets all wrote of it, they spoke of this hope..
Many people have waited for thousands of years.
for this hope..
Now we have him, we have this hope..
Yet unlike all the people of God that came before us,.
we don't have to wait and look forward to it..
We look back to something that already happened..
We know what happened, we have a confidence of it,.
of the time that Jesus entered the world.
to live a perfectly righteous life that we never could..
To suffer and die a painful, humiliating death.

$^{681}$that we never would..
To bear the wrath of God that we deserved..
To defeat sin, to defeat Satan..
And we know that not only was he born in this world.
of the Virgin Mary, not only did he suffer.
and die under Pontius Pilate,.
not only was he buried dead for three days.
after three days, he rose again to life..
And so we can say confidently, as Job said,.
I know my Redeemer lives..
Not only is this a present reality that we look back to,.
that we commemorate yearly,.
we also have an even greater future hope that we hold onto..
The second advent, the second coming of Jesus Christ.
into the world, the Messiah who will once and for all.
defeat the serpent..
And when he comes back, he will raise all his believers.
who have fallen, he will raise us and give us a new body,.
uncorrupted by sin, new minds, unstained by evil..
He will put the enemy under his feet.
and then put an end to sin and death and pain for all time..
And he will give us eternity with God in a new heavens,.
new earth, a new place of perfection..
And this is our joy, this is our peace, this is our hope..
This is our solution to the problem of sin, shame and guilt,.
loneliness, hopelessness, depression, sickness,.
death, frustration, trouble, trial, tribulation..
Jesus Christ is our hope..
So during this Christmas season, we should be encouraged,.
encouraged by this hope of who he is, what he has done.
and what he will ultimately accomplish..
Even in the midst of difficulties and troubles,.
we have a hope, we have something to hold onto.
that will keep us going..
And therefore we should be encouraged.
and we should encourage one another..
Look to a brother or sister,.
if you see someone around you, reach out..
If they're struggling, so many of us are,.
offer them that same hope, remind them,.

$^{721}$encourage them with this hope..
It's not just a hope that we have, we share with one another..
There's a passage I wanna, I'll just mention it,.
I won't read it, but Ephesians chapter two, verse 12,.
Paul is reminding the brothers and sisters in Ephesus.
of the time before they were in Christ,.
when they were alienated from the household of God,.
they were strangers and in the world they had no hope..
They were without God..
And there's people like that today..
And we once were like that with no hope in the world,.
with no ultimate solution or answer to all the troubles.
and pains of this world..
But in Christ, we have that hope..
And it's not just enough for us to hold onto the hope.
or to share it with our brothers and sisters,.
but we have to reach out to those in the world.
who are lost, stumbling in darkness, who don't have hope..
Share that message of hope with them..
I think Christmas time is a perfect time.
because it's already on everyone's mind..
There's gonna be family meals,.
you might be sitting next to somebody in your own family.
who is lost and without hope..
Share that hope with them, a coworker, a friend,.
a stranger on the bus or the train or whatever,.
share that hope with them..
It's the perfect time..
We have the hope that the world needs..
So as we approach our celebration of Christmas,.
of the advent of Messiah, Emmanuel, God with us,.
who as we heard earlier is the wonderful counselor,.
almighty God, everlasting father and prince of peace,.
the seed of the woman and the savior of mankind,.
let us draw close to him, hold onto him, find hope in him,.
and be as willing and as eager to share that hope.
with someone else..
And so my prayer for all of us, for you guys,.
is the same as Paul's prayer for the Romans,.
when he says, "May the God of hope fill you.

$^{761}$"with all joy and peace in believing,.
"so that by the power of the Holy Spirit,.
"you may abound in hope.".
Let's pray..
Heavenly Father, Lord God almighty, I thank you, Lord..
We love you, we praise your precious and holy name..
We thank you that even in the midst of a dark,.
sin-stained world full of evil and trouble and suffering,.
you have given forth a light in the midst of darkness,.
a light that we look to, a light that we have,.
that we hold onto, and God, I thank you for this light..
I thank you for your love that we see in Jesus Christ..
Help us, God, to be encouraged by this hope..
Help us to be strengthened and matured in our faith.
by this hope, and to be so filled with love.
for you and others that we take this hope to the world,.
who needs it more than anything, God..
So that is my prayer today, that we would all be,.
not just hearers of your word, but doers who act,.
who take this message and apply it to our lives, God..
Help us in this, I pray, for the glory.
and in the name of our Lord, Jesus Christ, amen..
\newpage



\section{路加福音}
\label{sec:NrY2oymR5Fc}
\textbf{將臨節的挑戰 - 施洗約翰出生的報喜 (路加福音1-2) \#1 - 陳崇基牧師}
\newline
\newline
連結: \href{https://youtube.com/watch?v=NrY2oymR5Fc}{\texttt{https://youtube.com/watch?v=NrY2oymR5Fc}} ~~~~ 語音日期: 2023-12-05
\newline
\newline
\hyperref[sec:a0q_gDDhWfE]{\small{< < < PREV SERMON < < <}}
~
\hyperref[sec:index_chronic]{\small{[返順時目]}}
~
\hyperref[sec:index_scriptual]{\small{[返順卷目]}}
~
\hyperref[sec:SsOhyCfDa_I]{\small{> > > NEXT SERMON > > >}}
\newline
\newline
$^{1}$各位好 來到我們這個Advent.
即是將臨節我們叫做.
或者迎接一個聖誕的彌臨.
有一件事就是我們幾個星期都會做的.
就是點蠟燭.
點蠟燭是一個很小的傳統.
在教會歷史裡面時常都會做的.
就是聖誕節之前那幾個星期.
每一個星期我們點一支蠟燭.
這個不是應節這麼簡單.
也是告訴我們聖誕節是什麼意思.
我們生活很忙的.
有時候那個節日來到的時候.
我們才預備.
或者甚至忘記了.
將臨節就是幫助我們.
有幾十天的時間 幾個星期.
去想清楚到底聖誕節是什麼意思.
當你去到聖誕節的時候.
或者聖誕節附近的時候.
你說Merry Christmas.
說聖誕節快樂的時候.
到底是什麼意思呢.
有一個問題想問大家.
如果你想去.
如果有人問你.
可能有些人沒看過聖經.
對聖誕節沒什麼認識.
或者他不知道聖經裡面怎麼說聖誕節.
如果用三分鐘跟他解釋.
聖誕節是一回什麼事呢.
什麼叫Christmas.
你會怎麼說.
三分鐘.
聖誕節是一回什麼事.
我們這幾個星期就會去看看.
到底聖經裡面怎麼去描述.
還有怎麼去預備.
救主的降生.
如何去預備所謂聖誕.

$^{41}$我想有一個心態是很重要的.
在這幾個星期裡面.
我都鼓勵大家.
你有聖經的話.
應該每個人都會有聖經.
現在你有電話就有聖經.
你去聖經那裡.
不妨看一看.
這幾個星期.
福音書怎麼去描述聖誕節.
特別是馬太福音和盧家福音.
都會說到一點點.
救主降生的事情.
看的時候我需要有一個心態.
那個心態是什麼呢.
就是一種Wonder.
我們應該有一種Sense of Wonder.
就是有一種驚喜的感覺.
那種驚喜的感覺.
不要因為每年都過聖誕節.
我們就覺得沒什麼所謂.
驚喜的感覺是什麼呢.
就是救主的降臨.
一個生命的降臨.
大家有小孩的.
記不記得第一胎在產房的時候的感覺.
有沒有Wonder.
還是你很害怕.
看到血淋淋很害怕.
抱著小孩.
你不用真正去產房.
你都可以感受得到.
那種Wonder.
一個新生命來到這個世間.
何況是一個神的兒子.
耶穌基督降臨在這個世間上.
那個Sense of Wonder.
我們再回復.
去看聖誕故事的時候.
到底是一回什麼事呢.

$^{81}$所以這個星期我們看看.
章林節或者聖誕節.
聖經是怎麼描述的.
特別是路加福音第一二章.
我們會看多一點經文.
因為我覺得去到聖誕節的時候.
讀多一點經文是應該的.
不過讀之前我們首先了解.
路加福音一二章是說什麼事.
我首先給大家一個地圖.
這裡不是雜亂無章地記載.
一些救主降生之前發生的事.
路加福音第一章開始.
有一個序言.
寫給讀者.
寫給那些初代的基督徒.
告訴他們耶穌基督的事情是怎麼樣的.
然後過了序言之後.
一章五節到二十五節.
看到嗎.
是很工整的.
一對一對這樣.
首先他說不是耶穌降生的事.
是施洗約翰.
那個叫John the Baptist.
施洗約翰出生的一個報喜.
Annunciation.
然後說完這個報喜之後.
就說救主耶穌基督出生的報喜.
一章二十六節三十八節.
然後就說到兩個母親.
就是施洗約翰的母親.
以利沙白.
Elizabeth.
和耶穌的母親是誰.
瑪利亞.
沒錯了.
瑪利亞他們兩個相遇的時候.
母胎怎樣驚喜.
大家快要出生了.

$^{121}$然後呢.
就去到一章五十七節到八十節.
就說到約翰自己的出生.
然後有首讚美詩在那裡.
第二章就說到耶穌的出生.
一對一對的.
然後說到主被奉獻到聖殿裡面.
在耶路撒冷.
然後有兩個老人家.
西冕.
Simeon.
和Ana.
就去迎接這個出生的救主.
然後另外一對.
就去到耶穌已經長大了.
十二歲的時候.
主在聖殿裡面.
就和律法師在聖殿當中.
去到有一個辯論.
人們就很驚奇.
到底這個人是誰呢.
這個就是路加福音第一二章.
給我們的一個地圖.
正如我所說.
我們多看一點經文好嗎.
因為這些我們平日.
我們很少會看的.
聖誕節了.
我們看看聖誕節的經文.
一章第五節.
在希律作猶太王的時候.
亞比亞班裡面有一個祭司.
叫做撒加尼亞.
他妻子是亞倫的後代.
叫做以利沙伯.
他們兩個在上帝面前.
都是二人遵行主的一切誡命和條例.
沒有可指責的.
只是他們沒有孩子.
因為以利沙伯不生育.

$^{161}$兩個人又年紀老邁了.
這裡頭一句.
說希律作猶太王的時候.
這個不是純粹歷史新聞報道這麼簡單.
因為人們聽到希律.
這個是大希律叫做Hildegrave.
聽到這個名字的時候.
不是純粹一個歷史人物這麼簡單.
他會想起當時的政局.
或者當時的情況會是怎樣.
或者今天這樣.
如果我們說在希特拉執政的時候.
發生這樣的事.
你就知道發生什麼事了.
很不妙是吧.
很差的情況.
這裡希律作猶太王的時候.
就是一句這樣的說話.
因為大希律什麼時候做猶太人的王呢.
他在主前40至44年的時間.
大概有幾十年.
他就是鐵拳掌管猶太人的地方.
他在羅馬京都上京的時候.
他就得到羅馬議會的任命.
回去巴勒斯坦地來作王.
然後打了幾年仗之後.
他就將之前的王朝叫做瑪加比王朝.
將末代的王子都殺了.
然後他就執掌猶太人地方的政權.
我們今天提起大希律的時候.
我們就說他是一個最偉大的建築家.
因為留下很多很多的東西給我們看.
馬薩大,耶路撒冷,聖殿.
很多龐大衛大的建築.
我們今天就覺得這些是打卡的地方.
但是當時的人是不會打卡的.
當時是被暴虐的時候.
有誰幫他建這些東西.
就是平民百姓.
誰要立重稅.

$^{201}$平民百姓.
希律作猶太王的時候.
等於希特拉執政的時候.
一個最黑暗的時候.
當時的百姓恨不得將這個大希律剷除.
只不過不能夠將它剷除.
生活是非常之苦.
你這樣看的時候.
你就發現路加想說的聖誕節.
就和我們平日想的聖誕節.
有蠟燭,聖誕樹.
有沒有去downtown看聖誕樹.
怎樣著聖誕樹.
有沒有.
很漂亮.
想想家裡怎樣佈置.
當時不會這樣想.
在黑暗裡面上帝有沒有作供.
這是問題.
所以為什麼在Evan的時候.
我們會點蠟燭.
黑暗與光之間的關係.
不單止大希律作王的時候.
好像希特拉作王一樣.
當時你知道整本聖經的歷史.
當時的百姓正在等上帝和他們講話.
等了很多年.
等了有400年時間.
你想想舊約的歷史.
由創世記到最後那本書.
叫做《馬拉基說話》.
舊約所有的東西都講了.
到馬拉基的時候.
大概是主前300多400年的時候.
然後到新約.
這個路加福音開始的時候.
中間隔了三,四百年的時間.
這段時間我們叫它沉默時期.
Years of silence.
也就是說沒有直接得到上帝的說話.

$^{241}$沒有先知出來.
百姓已經等了很久.
為什麼上帝還沒有說話.
上帝明明說過很多應許.
承諾過很多的東西.
等了幾百年四個世紀都沒有出現.
沒有說話可以講.
到底搞什麼科.
我不知道你和神的關係.
有時候會有這樣的掙扎.
上帝在搞什麼科.
上帝在做什麼.
或者你想想以前.
你心靈裡面經歷過一些低谷.
或者所謂的沉默時期.
不知道上帝在做什麼.
你回首一看的時候.
你才明白原來上帝在做這件事.
但當時你不會這樣想.
這就是當時猶太人百姓的心境.
等了四百年.
現在是2023年.
你倒數四百年.
這個地方是什麼地方.
美國都還沒有在地圖上.
四百年.
當時猶太人或者神的百姓.
就是在等這件事.
就是這個時候.
路加夫說什麼呢.
路加少兩個人出來.
兩個老人家.
撒迦利亞和他的太太以利沙白.
這兩個是什麼人呢.
是祭司家族的人.
說到他們是無可指責.
不代表他們完美.
而是代表他們的承諾.
對上帝的忠誠.
是不計較的.

$^{281}$全然的忠誠來事奉主.
由他們年輕到他們年老.
都是一路忠誠來事奉主的.
無可指責.
還有呢.
接著說了一句話.
就是不生育和年紀老邁.
這個就和當時的社會.
應該會感覺到那種感覺.
你回到當年古代的時候.
不生育年老.
這些是社會的羞辱.
你不生育就沒有盼望.
在社會裡面也沒有地位.
你記得舊約聖經裡面.
經常都有這樣的故事.
不生育的故事.
我們有時候想起的時候就奇怪.
為什麼每個人都生不了孩子呢.
你記得阿巴罕的太太撒拉.
不生育.
以撒的太太尼伯加.
不生育.
亞國的太太拉傑.
不生育.
舊約裡面經常有這樣的故事.
不生育的時候怎麼辦.
在人的角度來看.
不生育就是死路一條.
就是死路一條.
就好像上帝在生命裡面的計劃已經停了.
上帝再不搞你.
你沒有什麼參與.
正是這個時候.
聖經就說.
等一下.
開始有轉變.
在不生育的時候.
上帝的計劃是不是繼續向前呢.
第八節.

$^{321}$撒加利就按班次.
祭司嘛.
按班次在上帝面前要執行祭司的職務.
照祭司的規矩抽籤.
進入主的殿裡面燒香.
燒香的時候眾百姓在外面來禱告.
當時的聖殿有很多事情要做.
就像今天教會.
你進來做了什麼.
吃東西喝咖啡.
聊幾句.
坐在那裡.
聚會結束了.
出去等垃圾然後走.
其實裡面有很多事情要做.
有很多人在預備.
聖殿也是一樣.
百姓去到聖殿敬拜的時候.
不是純粹奉獻一件事.
或者祭生等等.
他們是知道聖殿的運作是很複雜的.
裡面有很多很多事情要預備.
撒加利就是一個祭司.
班次裡面其中一個小祭司.
當時的祭司分24班.
有24個團隊.
每個團隊可以在聖殿裡面做兩次.
並且誰可以進到聖所燒香.
就要抽籤.
如果不是班次不用事奉的時候.
你去做什麼.
你回家耕田.
到你的時候就安可.
離開你的本鄉.
或者在野口山附近.
在耶利哥.
你記得耶穌有一個故事.
說很撒馬利人的比喻.
從耶路撒冷下到耶利哥.
有祭司有你們未人下去.

$^{361}$因為當時有很多祭司住在耶利哥.
所以他們班次來到的時候.
到我們了.
我們就要上到耶路撒冷準備.
誰可以進去燒香.
就要抽籤.
當時的祭司我們知道大概有多少個.
大家聽清楚.
有18000個祭司.
18000個.
加起來有這麼多個班次.
還有那些僕人,利未人等等.
如果你說撒迦利亞可以抽籤.
中了然後可以進去聖所燒香的話.
幾乎是中了透彩票.
是不是.
可能一輩子都沒試過.
現在年老的時候.
突然間就抽中他.
上上籤.
他可以進去.
你想一下那個心境是怎樣.
一輩子我們教會有18000個牧師.
每人按班次講道.
不是輪流的講.
要抽籤.
抽到你講道的時候.
會不會很興奮呢.
還是會覺得很麻煩呢.
不要搞我.
I don't know.
撒迦利亞就是抽到.
幾乎是他死之前唯一一次.
做一個祭司.
如果作為一個祭司進到耶路撒冷的聖殿.
一個這麼堂皇的地方.
能夠在聖殿燒香.
獻上香.
這個可以是不枉此生.
主的殿燒香.

$^{401}$主的殿.
主的殿在猶太人的世界裡面.
主的殿就好像美國總統的白宮.
上帝掌皇權.
在地上作供.
他的situation room.
就是temple.
就是耶路撒冷的聖殿.
你進去見總統.
能夠進白宮.
有誰進過.
我呢.
就沒有進過.
我沒有這個機會.
你想想有白宮的情況.
那種感受會是怎樣.
這個就是撒迦利亞進去.
這個聖殿裡面.
燒香的那種情況.
那種感受.
大希律的聖殿.
大家應該很清楚.
有很多的無論是發掘.
研究都告訴我們.
這個地方是一個多麼堂皇的地方.
這個幾乎是整個羅馬帝國.
當時古代的世界最大的聖殿.
大家有沒有聽清楚.
是最大的聖殿.
有很多其他的廟.
賈臣的廟.
或者奧古士多羅馬王的廟.
都沒有那麼大.
這個是最大最大的廟.
幾十個美國足球的球場.
的size.
還是對猶太人的世界來講.
是唯一一個地方.
唯一一個地方敬拜神.
大希律的聖殿.

$^{441}$在這幅相.
或者在這個重構上.
在橄欖山東面.
猶太山東面.
向西望.
看到的當時所謂的second temple.
大希律的聖殿.
你見到中間坐立的有四條柱.
那個就是聖所的地方.
所謂temple proper的地方.
其他的地方都是.
讓人去聚集.
或者露天的廣場.
會場等等.
然後我們進去裡面的時候.
這個就是以色列人院.
如果我們去耶路撒冷.
敬拜的時候.
你是外邦人的話.
你就進不了這個地方.
如果你是以色列人.
都未必進得到.
如果你是婦女的話.
你就在這個quarter的外面.
只有男士才能夠進去裡面.
你見到左下角那裡有一個平台.
那個就是祭壇.
那裡有一條ramp上去.
那個就是祭壇.
然後進去四條柱那裡.
有一道門進去.
那個就叫做聖所主的殿.
裡面是長方形的一個空間.
然後再裡面就是至聖所.
這個地方就是聖所.
所謂主的殿.
長方形.
裡面有什麼呢.
簡單講一講.
第一樣東西.

$^{481}$就是那個minora.
七燈台.
七燈台是24/7會開的.
對面右邊那裡有一個桌子.
有兩支棍.
那個叫做陳設餅.
又是24/7有餅在那裡.
新鮮的餅.
好像那些西餅店那樣.
每天早上就弄.
弄上去.
然後中間那兩個人.
有一個正方形的壇.
那個就叫做香壇.
香壇是什麼意思.
在香壇燒香.
就代表上帝的子民的祈禱.
香是往什麼方向升的.
向上.
香壇點起的時候.
就代表上帝的子民的祈禱.
去到上帝的面前.
燒香的時候.
眾百姓在外面禱告.
你可以想像一下.
眾百姓在外面等什麼.
當他燒香點著香的時候.
香就向上飄響.
然後藉著門縫和窗縫.
煙就會飄到外面.
眾百姓在外面禱告的時候.
他們就見到一個這樣的情況.
如果你是天主教背景的.
你記得每逢天主教選教宗的時候.
是不是要燒煙.
那些主教就在那裡投票.
投完票之後.
他們就把煙弄出去.
黑色的代表什麼.
還沒行.

$^{521}$那就要再投過.
然後他們又再投過的時候.
終於投到.
他們就加一些化學物料進去.
然後就變成白煙.
然後就說.
OK 總之選了誰做教宗.
當然我們不會參與.
我們不會去參與這些東西.
聖經也沒有說.
但是我想你了解的情況就是這樣.
祭司進去 撒加尼進去.
這些香飄上去的時候.
就代表眾百姓自己有祈禱.
得到上帝的悅耳.
去到上帝的面前.
一個很隆重的儀式.
也是一個很重要的時刻.
一章第十一節.
有主的一個使者.
我們在聖殿裡面.
聖所裡面.
那個長方形的地方.
只有誰在那裡.
撒加尼亞.
有主的使者站在香壇的右邊.
向撒加尼亞顯現.
撒加尼亞看見就驚慌害怕.
害不害怕.
突然間你去洗手間.
突然間旁邊有一個人.
類似這樣的情況.
在聖所裡面.
突然間在香壇的右邊.
出現另外一個人.
看見他就驚怕.
天使對他說.
撒加尼亞不如怕.
因為你的祈禱已經被聽見了.
你的妻子以利沙白.

$^{561}$要生一個兒子.
你要給他起名叫約翰.
你必歡喜快樂.
有許多人因他出世.
必會喜樂.
他在主的面前.
將要為大.
淡酒濃酒都不可喝.
從母腹裡面.
他就被聖靈充滿了.
他要使許多以色列人迴轉.
歸向主他們的上帝.
他將有以利亞的精神和能力.
走在主的面前.
叫父親的心轉向兒女.
叫伴逆的人.
轉向二人的智慧.
又為主預備.
迎接他的百姓.
你的祈禱已經被聽見了.
當然祭司進到聖所的時候.
他帶著一個禱告.
禱告可以是一個national prayer.
一個國民的祈禱.
祈禱什麼?.
剛才我不是說過希特拉專政.
四百年沉默.
你的national prayer是什麼?.
上帝在你的國度來到.
好不好?.
不要再遲延了.
今時今日要來到.
這是national prayer.
他說你的禱告已經聽見了.
似乎不是說national prayer.
我想也不是說當時撒迦利亞.
會為了有個兒子好不好.
這樣祈禱.
因為剛才交代了什麼?.
已經老邁了.

$^{601}$已經無望了.
hopeless.
還生什麼呢?.
你七十歲的時候.
你不會說我希望多生一個兒子.
你不要搞這種事.
你都不夠膽這樣祈禱.
I don't know.
他說你的祈禱已經蒙月立.
我想可能是.
說他年輕的時候.
他曾經苦討過一些東西.
兒子一直沒有來過.
他就在想上帝在搞什麼呢?.
他說你的祈禱應該是他多年前.
為了他的兒子.
能夠得到兒子來祈禱.
多年苦苦的禱告上帝.
就是沉默不言.
就好像silent period四百年一樣.
他已經放棄了.
give up.
他認為上帝的計劃應該沒有他的份.
他繼續會有羞辱.
但過了多年之後.
現在上帝的國臨到的時候.
許多人會得到blessing.
在主面前.
他要為大分別為聖.
被聖靈充滿.
淡酒濃酒都不可喝.
不是代表我們不要喝酒.
不是這個意思.
而是喜樂在哪裡.
他心靈裡面喜樂的傳源.
不是在酒那裡.
是在天上.
喂.
以色列文要迴轉.
他有以利亞的心靈.

$^{641}$這不就是舊約裡面所說.
上帝救恩的日子來到的時候.
聖誕來到的時候.
就是會有以利亞.
因為以利亞就是先鋒.
會為主開路.
為主預備迎接塔利伯式.
幾乎這個描述和這個baby的描述.
施西約翰出生之前.
這個描述這個報喜.
包括整本聖經裡面.
上帝曾經說過的應許.
每字每句一句不漏.
整本舊約聖經就在這個人的身上.
會實現.
喂.
這是一個很大的日子.
剛才我說過我們產房的時候.
BB出生.
我記得我第一胎.
小孩出生的時候.
都沒有人跟我說這句話.
他可能跟我說的就是.
你死了你要開始還債.
你負債在身.
你將不會得到福了.
等等的東西.
這裡反過來.
施西約翰出生.
就是整本舊約聖經所有應許.
能夠得到fulfillment.
這個moment.
其實一個老祭司或者老百姓.
應該很清楚.
你求了這麼久一個national prayer.
就是希望什麼.
就是這些事情會發生.
撒迦利亞怎麼做.
一章十八節.
撒迦利亞對天使說.

$^{681}$我怎麼能夠知道這件事呢.
我已經老了.
我的妻子二十年多老了.
他的意思就是說.
no way.
there is no way.
天使跟他說.
什麼叫no way.
我是站在上帝面前的加百列.
天使有很多個當然.
成千上萬.
但只有幾個在聖經裡面有名字.
Gabriel.
加百列就是其中一個最威風的.
因為舊約裡面每逢有Gabriel出現.
有加百列出現.
就是上帝的計劃要翻開新的一章.
他是揭書那個.
他揭到下一章.
他說他是加百列.
每逢警察已經跟你說這句話.
跟你說這句話.
還有報這個訊息給你.
到了時候這句話必然會應驗的.
take my word.
我說了就算數.
上帝說過的就算數.
到了時候這話必會應驗.
只因你不信我.
你就成為啞巴不能說話.
直到這事情實現的日子.
從上下問下你.
我知道不是啞巴這麼簡單.
他是聾了.
你想像一下你聾了又啞巴了.
會是怎樣.
這個就是撒吉利亞.
我想有兩方面可以想到.
這個不單單是一個懲罰.
正如我所說的一個老祭司.

$^{721}$很熟讀律法.
遵守律法的條例一字不漏.
應該很清楚上帝的心意.
上帝的救恩來到的時候應該會發生什麼事.
但他低估了上帝.
underestimate上帝.
你記得以前有一本書.
很小的書來的英文書.
是一個英文學者寫的.
那本書的title.
你不用看那本書.
我告訴你title你知道title說什麼.
title說什麼.
那本很小的書.
他說your god is too small.
你的上帝太小.
小到可能放在神主位置.
那個就當是上帝.
your god is too small.
有時候低估上帝和叛逆上帝是一線之隔.
這個就是老祭司薩加里亞學到的功課.
一方面是對他一個教訓.
另一方面可能也是一個恩典.
因為老祭司點完香之後.
他出到去聾了啞了.
到他再開口的時候.
他就不是說一些八卦新聞.
他就讚美上帝.
當施洗約翰十月懷胎之後.
出生的時候.
他確實就是衝口而出.
是讚美的聲音.
是知道上帝的權能是不可以低估.
一章第21節.
百姓等候薩加里亞.
就詫異他在聖所裡面.
馳而了那麼久.
你想不想到外面的百姓在想什麼.
煙已經飄了出來.
人都還沒出來.

$^{761}$可能他死了.
因為太老死了.
因為他做錯事.
上帝激殺他死了.
詫異為什麼這麼遲都不出來.
到他終於出來.
推門出來的時候.
卻不能夠與他們說話.
他們就知道在聖所裡面見到異象.
聾了啞了就跟他打手勢.
因為他成為啞巴.
到他公職的日子滿了就回家.
這些日子以後.
他的妻子以利沙白就懷了孕.
隱藏了五個月.
一半的時間.
以利沙白說.
主在眷顧我的日子.
這樣看顧我要除掉我在人間面前的羞恥.
你見不見到這個故事很有趣.
說到聖誕的時候.
他就先說這對老夫婦.
和施洗約翰神跡的出生.
上帝為撒迦利亞成就的.
是多過他所想所求的.
你想一下我剛才說過.
早年的時候他們兩夫婦.
可能苦討了很久很久.
他們想只是一個兒子.
希望上帝給過他們.
但是現在成就的東西.
就不是純粹一個兒子.
是上帝有一個更大的計劃.
要在他們身上實現.
施洗約翰的出生.
不是純粹一個兒子這麼簡單.
是救主的煽風.
是以利亞先知來到.
是上古之年上帝的應許.
全部一句一句要實現的時候.

$^{801}$其實有很多時候我們都是這樣的.
你想一下我們曾經有過的祈禱.
我們和上帝.
或者如果你未信主.
你和上天祈禱.
但當然上帝是掌管一切.
向上帝禱告的時候.
很多時候是出於自己能夠見到的東西.
我們很短視.
在中間有誰是近視的.
遠的東西我們看不到.
這個故事也是這樣說.
曾經有過的祈禱都是為自己祈禱.
但往往上帝是用他自己的時間.
做他的事.
並且做更大的事.
是你未必測得透的事.
你可能遺忘過的祈禱.
但上帝這裡沒有遺忘.
人以為神在答允.
但上帝在他自己的時間裡去答允.
人以為上帝不會成就的.
但上帝在後來你都成就.
或者在你人生以後你都可以成就.
一個更大的計劃.
這個就是有趣的地方.
聖經很有趣.
你看這個故事.
有時候有趣的地方.
連人名都和故事有關.
一個老祭司叫撒加尼亞.
撒加尼亞是什麼意思.
喜伯來文就是耶和華紀念的意思.
撒加尼亞就是神回憶.
撒加尼亞回憶.
路加去到講聖誕的故事.
他用這個故事開始的時候.
就是在說這個老祭司由不信變成相信.
由不信.
怎麼可能會這樣發生.

$^{841}$成為他最終明白他要相信.
因為其實整本路加夫告訴我們的信息.
聖誕就是這個意思.
怎麼會有這個可能.
上帝的兒子會降卑來到我們中間.
出生去拯救我們.
No way.
但我們信了主的.
我們就知道.
由不信變成相信.
原來就是一線之隔.
你們了解上帝所成就的.
不只是單單你所求的.
而是上帝一個更大更大的計劃.
路加第一章講到施洗約翰出世的時候.
你不妨看一下.
撒加尼亞終於開口.
他說什麼.
孩子他是至高者的善子.
走在主的前面要預備他的道路.
這個也是舊約裡面.
耳措書裡面所講的.
叫他的百姓因聚得舍認識救恩.
因我們上帝憐憫的心腸.
叫清晨的日光.
The light.
從高天臨到我們.
照亮我們黑暗裡面.
死任當中的人.
把我們的腳引到和平的路上.
前兩天我再看一本書.
一本舊書.
是Philip Yancey寫的.
Disappointment with God.
他說對上帝很失望.
怎樣處理對上帝的失望.
我想大家都應該有這樣的經歷.
對上帝有些失望.
都不知道他在搞什麼.
其中一個他處理的問題就是.

$^{881}$上帝是否隱藏.
上帝在我生命裡面.
是不是真的隱藏了.
不見了.
我覺得很有趣.
再看回聖經的時候.
讓我們另一個角度去看.
我們希望上帝是做什麼.
我們希望上帝明顯做事.
大軍駕到.
然後有榮耀出現.
有煙花去放.
有紅地毯上帝駕到.
在我們生命裡面.
我求這一樣.
上帝就立刻.
好像變魔術一樣.
在我面前出現.
這個不是一個最大的見證.
但你數一數聖經裡面.
所有的這些故事.
逢真上帝大能駕到.
明顯地作功的時候.
人是反過來叛逆的.
以色列人過紅海.
過曠野.
有回眸給他.
有雲柱火柱.
每一天去引領著以色列人.
他們幾乎不用祈禱.
每一天.
警鐘響了.
就有火柱雲柱帶他們走.
以色列人是怎樣.
在曠野裡面是硬著脖子不順.
以色列人進到迦南地.
有什麼比進迦南美地更好.
可以安居樂業.
等於你買了一個八萬的豪宅.
你能夠安居樂業.

$^{921}$但以色列人進到迦南的時候.
上帝為他們打仗.
安居樂業之後.
他們就敬拜別神.
所羅門王.
第一個聖殿.
有榮耀進入這個聖殿.
連人都進不去.
太重了.
太厲害了.
但所羅門.
第二天就敬拜假神.
上帝很明顯地人卻叛逆.
但每當上帝隱藏的時候.
人卻好像有信心.
創世記約瑟的故事.
上帝不明顯.
但到最後他想清楚的時候.
再回頭看的時候.
他知道上帝的意思是好的.
路德的故事等等.
不同聖經裡面的故事.
全部都是上帝不明顯的時候.
人就開始去相信他.
及至神的兒子上十字架.
為我們罪人死.
他的降生.
死而復活.
都是不明顯的工作.
但他改變世界.
上帝是一個行事出於我們意外的神.
今年2023年.
我不知道你在上帝面前有什麼祈禱.
有什麼祈願.
如果上帝在你的生命裡面.
是有明顯的工作的話.
Praise the Lord.
請你不要忘記上帝.
但如果今年2023年.
兜兜轉轉好像上帝沒有聽你的禱告.

$^{961}$不應運你的祈求的時候.
好像沉默無聲的時候.
因為上帝的計劃.
比你的計劃更大.
真的.
上帝的能力.
上帝在你的生命裡面編寫的故事.
比你能夠看到的.
近視的你能夠看到的更長篇.
重點就是.
你是否願意全年繼續去相信.
我們低頭祈禱.
預備愛心.
領受主的恩典.
主的聖詞.
上帝你是那位行其事的神.
祂在我們面前降服.
去承認我們的不信.
我們真的信心軟弱.
我們看到的只是前面的事.
但聖誕主的兒子.
能夠降生在我們的生命當中.
這一件奇妙的事.
奧秘的事.
主鼓勵我們.
讓我們知道你自己的計劃.
比我們想像中更奇妙更大.
為我們預備的.
過於我們所想所求.
再一次來到你面前.
我們去預備我們的心.
領受你的救恩.
這個餅這個杯.
告訴我們.
十字架的救恩.
為我們所成就的.
確實是過於我們所想所求.
對我們的恩典.
\newpage



\section{}
\label{sec:SsOhyCfDa_I}
\textbf{回應方舟問題續集:詳解局部洪水說的理據,回應抱此看法是不信聖經的質疑,處理聖經無誤,原文釋經,文學釋經及末世時代論等相關課題(2023-12-6 鄺牧道頻.城匯網台聯播)}
\newline
\newline
連結: \href{https://youtube.com/watch?v=SsOhyCfDa-I}{\texttt{https://youtube.com/watch?v=SsOhyCfDa-I}} ~~~~ 語音日期: 2023-12-06
\newline
\newline
\hyperref[sec:NrY2oymR5Fc]{\small{< < < PREV SERMON < < <}}
~
\hyperref[sec:index_chronic]{\small{[返順時目]}}
~
\hyperref[sec:index_scriptual]{\small{[返順卷目]}}
~
\hyperref[sec:aFC7A7v5zzs]{\small{> > > NEXT SERMON > > >}}
\newline
\newline
$^{1}$Hello Mark Sir 你好.
又有機會可以聊天.
長話短說.
上次我們那條片.
我看日期都差不多十月頭.
我們可以回應影音使團.
他們說又再去找方舟.
我們上一條片.
大家可以搜尋成會網台.
或者我國務頻道.
去看回我們已經在回應.
究竟找方舟這件事.
我們覺得有什麼問題.
上次那些我們不重複.
但是Mark Sir都很好.
我們再聊的時候.
因為那兩條片出了.
分別都有一些不同的人.
在留言裡面或者私底下.
都對我們那條片講的論據.
或者一些內容有回應.
或者是挑戰.
我們就覺得值得澄清.
不如我就交時間給Mark Sir.
因為Mark Sir都準備了好幾個澄清.
以致我們可以從那裡回應一些問題.
有些都是很值得回應的.
第一樣我想要回應的就是.
信聖經這件事.
我們說洪水.
從聖經來看或者歷史考古來看.
局部的可能是機會大很多.
大過全球.
上次我們說過全球洪水從地質,考古等等.
甚至經文再深入去看的時候.
我們發現全球洪水.
這個不是經文一定要要求的結論.
有些人可能從小到大都是接受.
以前我小時候都是接受全球洪水.
慢慢認識的時候開始有些改變.

$^{41}$有些人可能聽到一個不同的解釋的時候.
我們就會質疑.
這個不同的解釋是不是不信聖經呢?.
我相信有一樣東西原則上我們需要澄清的就是.
我們大家都是相信聖經無誤和聖經權威.
這個最上.
然後下面就有很多不同的閱讀或者不同的解法.
我們要知道一件事就是.
我們相信聖經無誤.
不代表我們對聖經的解釋是無誤的.
正如我小時候接受的六天創造論.
每天都是24小時的.
硬性的.
聖經就好像一個科學的報告.
告訴我們不相信就好像對聖經不尊重.
小時候是這樣接受的.
慢慢長大之後就發現.
原來創世的第一章第二章等等其他的文獻.
我們就知道創世的故事到底是想說什麼呢?.
是不是在說科學呢?不是在說現代科學.
所以是用一個近東的世界觀宇宙觀去表達出來.
所以解釋上有轉變.
但解釋上的轉變不會影響聖經權威和聖經的無誤.
對嗎?.
我們對聖經的理解會越來越完備.
因為我們越來越多資料和越來越多不同的方法.
這個可能會對我們以往接觸的解說方法.
我們慢慢覺得以前的解釋是比較弱.
現在可能比較完備一些.
加上很多其他的學者的研究等等.
我們應該是有這種上進的.
但一定要搞清楚.
我們大家都贊成聖經無誤.
但不代表我們對聖經的解釋就是無誤.
是不是?.
所以你搞清楚這個時候我們就可以繼續討論.
所以我覺得第一個概念需要澄清的地方.
我想有不少弟兄姊妹.
其實在看這些或者聽一些傳道人講解這些.
就未必有那個盡心去讀聖經的解釋.

$^{81}$或者是原文的解釋接觸.
在聖經無誤這些概念其實都有修正的.
其中有一樣東西就是.
剛才Mark Sir說的.
我們小時候就說聖經摩西五經是摩西寫的.
不止摩西死的那段也是他寫的.
尼泊爾死的那段也是他寫的.
我進了信封然後跟約書亞說.
如果我不學完你就拆開它.
現在差不多在學界或者是在識經界.
摩西五經一定是後面有編修編輯.
甚至那個年期可能有些解釋不少.
尤其是新冥紀根本是很後期著作的.
就算我們不爭辯那個年期或者內容.
成書時間但是他背後有編輯.
那摩西寫但是編輯的手加上去.
他是不是也是入於神的密使的其中一部分.
我們會相信是.
但是問題是摩西五經尤其是最古遠的那幾卷.
其實是不容易去追溯.
抄本也是很後期才找到抄本.
所以我們剛才討論都是新約和舊約.
舊約在這方面是更加困難的.
處理文本上面.
甚至希伯來文本身都是.
如果大家有些認識就知道.
他那個文字本身後來文士都整理很多.
我們現在讀的希伯來文聖經.
有母音和子音.
其實是很多文字做了功夫.
或者是《識經家》做了功夫.
其實本身沒有文字的.
還有一點一劃都很容易會有抄寫的不同.
其實研究那個文字本身.
尤其是講到創世記.
也有很多的文字上.
希伯來人是用很多數字的方式去表達一些信息.
所以我們不是不相信這些.
而是那些聖經的文本本身.
是有很多修辭學的.

$^{121}$或者是他傳遞的東西要用當時的人明白的東西去傳遞.
所以這個就會使我們看著中文聖經.
還要是和合本的很堅持.
來解釋聖經傳遞.
等於全世界.
舉個例子.
Mark Sir會澄清.
如果是這樣.
接著當有人有其他解釋.
就說我們不信聖經.
這個是一個很抱歉.
會請大家多做一些研究.
Mark Sir補充一下.
我覺得大家需要思考得比較仔細一點.
才能夠說這個解釋是不同我的.
或者我接受的解釋不同了.
於是我們就把對方標籤了.
你就是不信聖經自由派.
我覺得這是華人的討論裡面.
經常我們用一些不需要思考的方法.
總之就標籤了你.
你就是不信的.
沒有這樣的事情.
Mark Sir你補充一下.
剛才聊天的時候.
說到芝加哥.
80年代.
對於聖經無誤的一個很出名的.
最保守派的一個聲明.
Statement of Biblical Inherency.
在80年代在芝加哥開了一個.
基要派和保守派的學會.
AIPAC等等他們都在那裡.
他們的聲明都有說.
雖然我們相信聖經無誤.
但我們解經的時候.
我們需要留意聖經的作者.
是會用當時的修辭手法來表達真理.
包括有比喻.
一些metaphor.

$^{161}$和hyperbole.
即是誇式的方法.
誇大的方法.
施辭的方法等等.
我們需要明白.
大家仍然最公認的.
最古老的福音派的宣言.
聖經無誤的宣言.
都有特別說到.
有誇式方法來做修辭.
所以這個和聖經無誤沒有任何抵觸.
我剛才和Mark Sir聊天的時候.
原來我們有common friend.
其中有幾個留言.
有人分享影片給我.
我自己看到.
就是北美一個很出名的牧者.
他是全球大紅水.
其實這樣說.
我不方便開名.
但有些人會以為我們沒有話題.
或者他的解釋和我們的解釋.
就是衝突.
然後他拿那件事來騙我.
你現在不跟這個.
這個應該是你很資深的牧者前輩.
但其實有些人不知道.
其實我和他認識的.
和也談過這個問題的.
十幾年前.
但以我認知.
如果我們談這些.
大家不同的論據.
大家不同的看法.
我們都不會認對方.
是一個不信聖經的人.
倒過來.
我看他那個.
就算我不是贊成他的解釋.
但大家都是牧者傳道人.

$^{201}$信徒.
我們都有我們那個.
研究的背景.
有時候我覺得.
弟兄姊妹就像Mark Sir這樣說.
二分法歸邊.
你不是那個解釋.
就變成你是一個不信聖經的人.
如果這麼簡單.
就不用讀這麼多聖經.
不用去上院讀.
其實不是這麼簡單.
就二分.
我們一會再跟進.
很簡單回應.
剛才說到聖經無誤.
和對聖經的解釋.
我們牧者對聖經的解釋.
其實都可以.
先講清楚.
我們牧者所教的解釋.
不等於是無誤的.
雖然我們看的都是聖經無誤.
是 絕對是.
所以你說.
我們可以拿出權威.
我也想告訴大家.
我們都有權威.
John Wharton.
在Wharton College.
現在剛退休.
Tramp Longman.
這些全部都是相信聖經.
High view of scripture.
的福音派.
很有資格的學者.
這些你有沒有認識到呢.
如果你認識到的時候.
你就發現原來有不同的聲音.
需要繼續去討論.

$^{241}$我想最主要想說的就是.
你怎樣去衡量哪一個解釋比較傳比一點.
首先第一件事.
你退後一步.
你想想方法是什麼.
你想清楚那個方法.
有沒有囊括歷史考古和文本.
上下文.
還有近東的背景.
這些全部都要考慮.
我們才能夠繼續讀經的時候.
好像我說的.
繼續能夠前進.
是嗎.
還有我們剛才也談到.
好像最近.
因為以巴的戰爭帶出來.
其實橫教會最近就講了很多時代論.
就是說時代論對於聖經末世的解釋.
因為很多人都出來講這個末世的問題.
其實時代論是一個很典型的.
我想我和Mark Sir以前也是受時代論的影響.
因為橫教會絕大部分都是受以前那套時代論.
但是其實我們一直看.
我們現在不是討論時代論.
但是問題是很肯定的.
時代論本身以前很出名推動時代論的學院.
他們都一直修正時代論的解釋.
他們自己也在修正.
但是有時候很可惜.
有些橫教會還在修正他們還沒改變的時代論.
停在那裡.
然後就覺得其他的解釋是錯的.
另外就是說.
在那個時代論的解釋裡面.
就算你很舊的新的有不同的時代論.
或者不是接受時代論.
其實我們在末世論的解釋裡面.
一直都有一個接受性.
我們可以有不同的解釋.

$^{281}$這是一個論.
時代論.
對聖經有一個解釋.
但是大家就不會因為你不相信時代論.
你就是一個不相信聖經的人.
所以這件事有一點相似.
但是我也想特別講一下重點.
我們為什麼要提出影音使團.
他現在找方舟的問題.
那些問題和聖經解釋有關.
但是有些問題就是.
他那個解釋是變來變去之外.
他的方法也很有問題.
我覺得大家應該要看清楚這件事.
如果他們是持續論點.
很持續的.
很一致的.
我們也不會有那麼多質疑.
問題是他們的解釋.
然後就做這些事.
然後要人給錢.
這個就是一個問題.
不要把問題放在提出問題的人身上.
阿柏sir不如你講一下一些補充.
你和我說的一些.
我覺得很多都很有意思的.
我會回應一些之前有些人問的.
關於局部性洪水和全球性洪水.
尤其是全球性洪水.
不是全球性洪水嗎.
那個全面性.
你可不可以補充多一點.
我知道你有很多準備.
我們上一次講到.
聖經裡面用一些普世的字眼.
其實可以用普世字眼來形容局部性的事件.
好像我們在創世記後面.
約瑟當時的埃及的饑荒.
講到全天下都是這樣.
包括澳洲都是這樣.

$^{321}$但是我們知道歷史的背景的時候.
我們知道這個絕對不會是全球的一個饑荒.
只不過是局部的中東的一個饑荒.
但是聖經的語言就會用.
Koharis.
All the earth, all the land.
你要明白他為什麼要這樣做.
他其實是想表達嚴重性在那裡.
其實羅亞洪水都是一樣.
不單止是我們這樣看.
Trump, John Wotton.
這些最好的舊約學者.
開始了解到如果我們說Koharis.
不一定是全球性的東西.
包括約瑟的饑荒.
如果你Koharis在洪水的故事.
創世記六至九章.
裡面見到Koharis的時候.
All the earth的時候.
那是不是必然性必要.
一定就是要全球的.
很明顯你這樣去以經解經的話.
是不需要的,不一定要.
當然我們也知道當時的近東的世界觀.
近東的世界觀是一個平面的世界觀.
不是一個全球的概念.
所以不能夠用現代人的字面解說.
我們就把字面的意義套在古代的聖經裡.
應該用當時近東的世界觀.
去講述他們的意義在哪裡.
所以我們看聖經就需要小心一點.
盡量點出這一點.
其中一件事就是有很多論據.
說全球性除了字面之外.
提到世界各地都有很多洪水的記載.
這是他們其中一個覺得頗強的論據.
就是如果不是全球性洪水.
世界各地也不會有洪水的記載.
那裡有甚麼可以回應呢?.
其實回應是很簡單的.

$^{361}$因為世界各地有很多洪水的故事.
無論是北美洲,澳洲,太平洋的島嶼等等.
其實你看到這些洪水的故事表面上很相似.
但其實有些重要的情節是完全不同的.
所以我們不能說世界各地有很多不同的洪水.
就指向一定會有一個全球性的洪水.
順帶一提,這些世界各地不同民族的洪水的故事.
很多時候他們不是在講全球性的洪水.
而是一個局部性的洪水.
當地局部性的大洪水.
或者再問多一個問題.
這些世界各地不同的洪水的故事.
是否出自同一個時間?.
同一段歷史的時間?.
如果不是的話,我們就不能亂將他們.
因為表面上很相似,所以就在講同一件事.
你不能這樣說.
甚至這些洪水的故事.
其實在某些地方是完全沒有的.
譬如非洲,埃及.
亞洲某些地方,歐洲某些地方.
是沒有這些故事的.
有些研究這些故事的起源的學者.
都指出某些很像羅亞方舟的故事.
其實是一些早期的傳教士.
去到那個地方傳了福音之後.
然後對那個地方的人.
如何看他們其他的洪水的故事,災難的故事.
就會有一個羅亞故事的影子在那裡.
美國,美洲的印第安人.
或者中世紀某些歐洲的洪水故事.
都是因為傳教士到了那個地方.
然後那些人接受了羅亞洪水這個故事.
然後他們就用這個故事來詮釋.
其他的洪水故事.
所以不能夠這麼簡單去看.
其實裡面都可以很複雜.
那個傳播裡面可以很複雜.
所以最主要其實就是.
表面上相似但其實很多情節是自然不同.

$^{401}$所以不能夠說周圍都有洪水故事.
所以一定會有一個曾經有一個的普世的洪水.
我印象中如果說是各地.
其實中國有沒有呢?.
好像沒有.
大禹治水.
想來想去都是大禹治水.
但大禹治水很明顯不是一個普世性的.
很明顯是自然不同.
如果這麼粗疏就拉上那個.
OK 中國都有沒有.
我馬上想了我是中國人.
你的帳戶不完整.
好像證實了剛才Mark Sir說的那件事.
那個說法都不同的.
如果硬扭下去就是龍門任搬.
是的 龍門任搬.
你怎麼說都可以.
還有上次說的不是很精細.
可不可以補充一下.
年期的問題.
其實實質.
洪水最有可能性是在哪個時間的範圍發生呢?.
因為其中.
你說.
這個我們需要囊括了不同的資料.
我們才能夠做一些觀察或者一些的立論.
之前說找到山上有木頭.
就說是三千多年前.
就是Early Bronze Age的時候.
然後就改變了.
然後就說有些學者上去探測過.
然後就說那裡的阿拉喇山的遺跡.
就是屬於一萬多年的時候.
Paleolithic的時候.
這樣是有些問題的.
為什麼呢?.
為什麼龍門任你亂搬.
之前就三千多.
之後一推就推到一萬多.

$^{441}$其實你預見的問題更加多.
因為如果我們大家都是接受聖經裡面的歷史性.
你就說回聖經由第一章創世記至到.
起碼十一章.
一至十一章.
裡面所說的話.
你就不能夠說不理這些資料.
因為你說如果羅亞洪水是在一萬多年前的時候.
發生Paleolithic的時候.
那就不對稱不吻合聖經裡面所說的.
因為你記得創世記第四章.
該隱的後代.
發生什麼事?.
該隱首先建造了一座城.
已經有City出現在中東.
我們知道中東的古代歷史.
有City出現有Urbanization出現.
是什麼時候呢?.
是新石器時代才開始有的.
Paleolithic的時候.
那時候是主前大概最早都是八千年的時候.
是嗎?.
然後第四章創世記最後.
就是說到該隱的後代有打銅打鐵的.
是嗎?.
然後有牧羊的.
這些你不能夠不去考慮一些資料.
因為如果你說我已經有牧羊業.
有農業.
這個絕對是去到更加後期一點.
所以你不能夠說搬到一萬多年前.
是嗎?.
因為城市已經說了.
羅亞是什麼時候之後.
是該隱的後代之後.
是嗎?.
所以有銅器鐵器.
所以你可能這樣去看回那個大的歷史.
搞古歷史背後發生了什麼事.
你就可能只能夠推論到.

$^{481}$羅亞的時代雖然我們不能夠偏頗.
但是起碼它不會去到萬多年.
是嗎?.
主前萬多年.
很可能去到大概主前三千年左右的時間.
就是給和拿.
你可以三千多三千多多多少少.
多除少補.
但是你不能夠一推就推到第二.
是嗎?.
因為你需要囊括所有聖經的.
以及考古的這些資料.
是嗎?.
所以有些學者都去指出.
羅亞的時代.
看了一些考古歷史等等.
羅亞洪水這件事情.
是主前大概三千年左右的時間.
所以如果你要在山上找木頭.
你那個化驗的年代都應該要吻合這個時代.
這個文章其實就在一個很保守的神學的期刊裡面出現.
其中一篇文章是Paul Szili.
在《羅亞洪水》的《羅亞學者記事館》.
2004年的時候已經發表了他看羅亞洪水.
它的地域性有多少.
有多大.
以及那個年代.
The date and the age.
這個學者認為.
有很好的理由去相信羅亞的洪水.
以及其他近東所說的洪水.
是在說一個歷史性的洪水.
而發生的時代大概是主前2900年的時候.
我現在分享的畫面都可以看到那篇論文.
Westminster的Theological General.
686期很清楚.
2004年.
我會放一個連結給大家.
我在網上找到的.
你可以訂閱去買的.

$^{521}$多謝馬Sir.
告訴我 我自己也看了一次.
資料都挺硬的.
不過他那個討論的方法.
其實我們所謂很堅定的論據.
那個論據是非常清晰的.
我覺得這個不是說要拋書包.
有人問我們學歷.
其實我讀新約.
馬Sir就讀聖經考古.
大家知道馬Sir的學樣都很好.
我們不是博士.
博士也有很多種.
以及我們學術界裡面.
或者聖經研究裡面.
其實我們都是講論據.
我們去討論一件事.
不然我們全代人都不用講.
我們不是博士怎麼講.
我們受一個神學訓練.
應該都是我們按正義分解真道.
我們都是需要很小心去查考.
就是引用這些所謂權威.
但不是說那個權威.
他是權威所以他的東西都很對.
沒有這件事像聖經解釋.
或者在學界或者聖經研究.
我們都是按著我們能夠去研究到的.
找到的東西.
我們盡量去貼近我們那個解釋.
希望解釋得最好.
所以其實學界三年已經回一回.
不斷有很多新的資料.
給我們知道.
去幫助我們理解聖經.
有時候弟兄姊妹我也想回應.
有時候很兩極.
覺得一定要學者.
講那些才有權威.
但有時候當我們引用那些權威.

$^{561}$又覺得我們比較熟悉一點.
其實是沒有衝突的.
很多學者都是很熟悉的.
但我們一定是.
我們其實是很感恩的.
在聖經研究裡面.
我們閱讀就覺得很多東西不懂.
但是我們去看到.
也有很多人用很多精力去研究.
以致我們學到很多東西.
以致我們去判別.
其實我們也是為了判別.
究竟什麼是對什麼是錯.
還有聖經的解釋有多全面.
好像Mark Sir說的.
我們其實也是很普通的人.
但是我們也是盡量去學習.
和去講解.
Mark Sir你還有什麼補充呢.
我們也不想說.
有一個補充很簡單.
就是那個文學結構.
就是這個洪水故事的文學結構.
在不同的場合.
可能你聽過《牧者》.
或者看過那些書.
就是那個線型結構.
很明顯這個洪水的故事.
是有一個線型的文學的元素.
就是七天等候洪水來臨.
然後四十天洪水氾濫.
水勢浩大一百五十天.
然後就神紀念羅亞.
神紀念羅亞.
然後就將那個數字.
一百五十四十七.
又反過來.
這個很明顯是一個線型結構.
但是很多時候我們聽到.
這些文學元素的時候.

$^{601}$我們就說你一講文學.
就是不相信歷史.
其實絕對沒有這樣的事.
因為你要了解.
經文背後確實是有一個.
歷史的洪水事件.
但是經文本身.
不是要給我們.
很多很多的細節.
讓我們看完經文.
重溫歷史的事件.
其實最簡單就是.
我們今天看的聖經.
經文本身並非.
那個背後事件的本身.
text is not the event itself.
我們要知道.
text 就是文本.
只不過是去到解釋.
那個歷史事件.
那個意義.
所以他解釋的時候.
文本就用不同的文學結構.
元素去表達出一些.
神學的意義.
所以不能夠說.
因為講到文學元素.
就一定說這個背後的事件.
就是虛構.
你要搞清楚.
聖經裡面.
所表達出來的意義.
神學的意義.
是在經文裡面找到.
我們用經文去重構.
歷史的事件裡面.
找到神學的意義.
是不是?.
很簡單很清楚.
譬如耶穌的五餅二魚.

$^{641}$四福音都有記載.
我們查經的時候就在想.
這個五餅二魚到底是怎麼發生的呢?.
但是聖經很明顯.
四次講到神蹟的時候.
他不是集中於.
怎樣發生.
他告訴你發生之後的意義在哪裡.
耶穌就是彌賽亞.
在講話裡面能夠叫人餵飽.
是不是?.
那個理論意義.
在文字裡面找到.
不是在事件裡面找到.
這個值得我們去想一想.
值得我們去思索一下.
聖經怎樣去.
用文學的手法.
來描述.
不會說.
是抵觸於.
他所描述的.
事件的歷史性.
所以這個就是聖經.
還有大家如果有.
讀歷史的背景.
你會知道.
每一個就算寫歷史.
是沒有純歷史的東西.
每一個記載歷史.
尤其是古代的歷史.
皇帝的歷史.
找一些史家幫他寫.
他一定有偏見.
一定有傳息.
聖經也不是為了寫歷史.
寫九因.
寫神的作為.
我們很明顯知道.
所以沒有純歷史.

$^{681}$根本是零.
是沒有可能.
他一定的文本.
寫那些事件.
就算是歷史事件.
一定有神學的意義.
在裡面.
這麼古代的摩西五經.
他也要向當時的人.
講明道理.
他用的東西.
全部都是當時的.
智慧.
或者他們認知的世界.
特別是.
羅亞洪水那裡.
阿不Sir可以補充.
無論創造.
他們的概念.
世界觀.
他寫的.
他寫的人.
要跟這些人講.
其實背後神是怎樣做事.
那個九因是怎樣.
他就用當時的人.
不是用現代的.
那些科學.
去看.
如果這樣就死了.
我們現代科學.
也在變的.
所以.
我們要.
這個方法的問題.
不是用我們現代人的.
眼光.
而是用古代近東人.
來看.
創造的故事都是一樣.

$^{721}$我們不可以用我們現代人.
宇宙觀來看.
用古代近東人的宇宙觀來看.
我們可以很快.
去處理其他問題.
但有一個問題我想處理.
創世的第九章的彩虹之約.
很多時候我們會聽到.
因為上帝說.
有一個彩虹之約.
永遠都不會用洪水滅世.
所以這個洪水一定.
一定是全球性的.
因為不會再有.
由羅亞到現在.
不再有任何全球的洪水.
但有很多局部性的洪水.
又發生.
他的意思就是說.
如果你說羅亞的洪水.
是局部性的話.
即是上帝又毀約.
明明說不再有洪水滅世.
那我就要小心一點.
因為第九章是一個盟約的經文.
說到盟約的時候.
一定有兩個派.
一個派就是上帝.
另外一個派就是人類.
如果看清楚九章.
開頭的時候就是上帝再重新.
創造祝福.
給羅亞代表新的人類.
等等.
然後他就很清楚說到.
大家兩者在約裡是需要守約的.
九章第七節說.
和合本就要你們要生養眾多.
其實那個意思原文是什麼呢.
他應該是至於你們.

$^{761}$As for you.
你們就要生養眾多.
地上要昌盛等等.
然後九章九節就說另外一個派.
就說看啊我要.
與你們你們的後裔.
立我的約.
原文其實就是至於我.
至於我這一邊我這個派.
看啊我要去到.
和你和你的後裔立約.
雙方的結盟雙方的合作.
因為我以前聽過.
有些人說我們.
concern一些.
global warming.
冰川融化.
會帶來水災毀滅地球.
有些人或者牧者會這樣說.
如果你這樣擔心的話.
憂心這些事情會發生的話.
對上帝就沒有信心.
因為上帝說過不會有洪水.
不會有全球洪水.
這個我相信你再看清楚的時候.
你會發現不是這樣的意思.
就是一個.
創世的地球章其實都是說.
雙方的合作不是說.
上帝寫包單單方面.
隨便你怎麼生活都好你隨便毀毀地球都好.
我都不會有洪水.
滅世的其實是說.
上帝我會做我這一邊彩虹.
就是那個立約記號但是你們.
都要如果你們好像以前.
這樣充滿暴力的話我會照樣.
朕盤是不是.
所以如果你這樣看的時候其實那個.
雙方合作都是一樣.

$^{801}$一貫的在聖經裡面.
講到盟約的時候都是雙方的.
連阿巴拉罕的約都是這樣.
第十七章創世記.
神和阿巴蘭來到雙方盟約.
也都說.
至於我要做什麼.
至於你又要做什麼.
阿巴蘭就要守約然後上帝.
一個可能性.
祝福的可能性才能夠.
直至阿巴拉罕能夠.
臨到萬國.
我們一講到盟約的時候.
我們就說上帝寫包彈.
隨隨便便沒有所謂.
沒有這樣的事.
因為講到覆盤的時候就一定是要這樣.
所以這個就.
你不能夠用.
彩虹之約第九章.
就說這個洪水一定是全球性.
正如我們所說.
經文沒有要求我們.
洪水.
一定要全球性.
局部性也都好.
但局部也好.
全球也好.
經文背後的歷史.
到底範圍有多大.
都不影響經文的神學意義.
正如我們所說.
事件和經文.
是可以不同.
經文所說的神學意義.
永恆的意義.
屬靈的意義是沒有改變.
講到羅亞洪水的故事.
和其他古近東的洪水神話.

$^{841}$確實是有些相似.
但你說.
那些是不是聖經抄襲他們.
首先你講抄襲的時候.
我們用現代人的概念.
因為我們有copyright, intellectual property等等.
古代的文獻是沒有這方面.
他們可以照抄.
照搬過來也無所謂.
還有如果你這樣想.
因為聖經比較後期.
然後近東的泥板.
洪水的泥板比較古舊.
所以一定是後者抄前者.
那又不是很對.
忘記了文化是怎樣運作.
因為有很多東西是根深蒂固.
不是說誰抄誰.
曾經有一個歷史.
大洪水死了很多人.
大家其實是在回應不同的角度.
聖經是其中一個.
想去回應說.
這是道德的問題.
為什麼會有這樣的災難.
其他的神話就會說.
神無緣無故不問究竟.
他就弄一些這樣的災難下來.
你要知道文化的運作.
不是說誰抄誰.
也不是說誰引用誰.
而是大家在回應一個.
已經根深蒂固的歷史事件.
在我們中間.
有些東西已經深入了民心.
不是說誰引用誰.
只不過大家做一些echo.
所以我就不需要太過.
覺得.
為什麼要用近東的洪水神話去比較.

$^{881}$因為他們有相似的地方.
也有很多不同的地方.
不同的地方就跟神學意義不同.
所以我們要去理解洪水的故事.
為什麼會寫成這樣.
的原因.
因為他就是在回應當時.
坊間如何去看待洪水的原因.
聖經就說這個世界.
是上帝會審判.
上帝會問責.
有暴力遍地的時候.
上帝就會來審判.
這就是聖經裡面最清楚的.
神學的意義.
我們所回應.
所針對的.
有些回應.
我們局部的洪水的解釋.
不是全球性.
就覺得我們可能背離了.
聖經.
相信聖經那些.
其實我想回應.
那些弟兄姊妹.
你要看回.
當然我們是針對醒.
你再去找方舟.
但那個問題.
不是說退一萬步.
就算解釋成為.
全球性的洪水也好.
那個問題.
是他那個.
去做考古的方法.
那個很大的問題.
你第一個年期.
就很混亂.
第二個地點.
又不肯定.

$^{921}$第三你怎麼處理.
考古的東西.
又不是合一個.
做考古的方法.
然後你的解釋.
又任搬龍門的解釋.
然後最大的問題.
你拿來傳福音.
你拿來叫弟兄姊妹奉獻.
我覺得這件事.
我們覺得.
是很有問題.
所以也都希望.
這個澄清幫到大家.
我想其實.
你說講信心.
其實很多人一讀神學.
信心就很大挑戰.
因為發現原來救約.
很多東西.
不是真的可以.
好像很多考古的證據.
甚至文本都有很多東西.
其實是要探究的.
原來很多東西.
我們真的憑信心.
我們經歷就看見.
在聖經帶出來的神學的信息.
神的救贖.
我們作為一個福音派的傳道人.
我們都會覺得.
我們相信神.
我們相信神的.
在聖經裡面講的.
很多的信息.
包括羅亞松水的信息.
所帶出來的.
但要研究下去.
很仔細的.
就可以容許有不同的解釋.

$^{961}$就不代表這些不同的解釋.
我們就被標籤.
你那個解釋不對.
你那個就等於不相信.
這個絕對不是.
我們都很想澄清.
給大家看一些比較成熟的.
希望大家.
去看這些事.
我們仍然歡迎.
大家在留言裡留言.
這個是成熟的.
一個理性的討論.
但是有時候留言.
請大家都注意一下.
我們都會有些回應.
有些東西如果太過不合理.
我們都會按照情況回應大家.
應該是開放來討論的.
這個是好事來的.
我們今天就講到這裡.
多謝Mark Sir.
我們遲些再談.
拜拜.
\newpage



\section{2:1-12}
\label{sec:aFC7A7v5zzs}
\textbf{The King Is Born  (Matthew 2\_1-12) - Sam Johnson}
\newline
\newline
連結: \href{https://youtube.com/watch?v=aFC7A7v5zzs}{\texttt{https://youtube.com/watch?v=aFC7A7v5zzs}} ~~~~ 語音日期: 2023-12-12
\newline
\newline
\hyperref[sec:SsOhyCfDa_I]{\small{< < < PREV SERMON < < <}}
~
\hyperref[sec:index_chronic]{\small{[返順時目]}}
~
\hyperref[sec:index_scriptual]{\small{[返順卷目]}}
~
\hyperref[sec:iKrjLPh_hzM]{\small{> > > NEXT SERMON > > >}}
\newline
\newline
$^{1}$- Well, good morning..
It's fantastic to be here and enjoying worshiping.
with you all..
Thank you to the worship team for leading us.
in worship this morning..
I'm excited to be here..
I love being at UV..
I love Christmas time..
This is a really fun time of the year..
And to begin with, I'm gonna ask you.
to use your imagination for just a moment.
and imagine one of the announcements that we just heard,.
the Christmas dinner, the UV Christmas dinner..
It's been a few years since I've been to one,.
but I used to be at the UV Christmas dinner every year..
And so just imagine that in your mind,.
or if you haven't been to one,.
you can imagine another feast that you've been at..
Think about all that food that is there.
just spread out on the mini tables,.
and it's so amazing and so good..
One of the problems that I have.
when I get to a feast like that.
is there's so much good stuff and I wanna eat all of it,.
but I know that I physically cannot,.
so I have to kind of choose what I'm gonna eat, right?.
And there's so many options, so many choices..
And the reality is we all have favorites..
We all have kind of our go-tos,.
especially at a big feast like that,.
the things that we want..
And so sometimes, at least what I do,.
is I go to the things I know I really enjoy,.
to my favorites..
And even though there's other really good stuff,.
I never actually get to it..
I never try some of those other things.
just because I always end up eating my old favorites.
and I can't eat it all..
And so there's something off to the side.

$^{41}$that I forget about..
And sometimes that happens to us.
with the themes and the stories of Christmas and Advent..
It's literally like a feast..
There's so much to think about.
in terms of theologically and spiritually.
in terms of Advent and Christmas..
And there's so many beautiful themes here.
of Jesus coming to earth..
We have the marvel of the miraculous virgin birth,.
supernaturally conceived by the Holy Spirit..
We think about the miracle, the incarnation,.
that the word was made flesh,.
God was made human, Emmanuel, God with us..
That's one of my favorite themes, such a powerful theme..
We think about the unveiling of Jesus as Savior.
and the verse that we heard earlier from the worship team.
where the angel spoke and said,.
"This is good news for the whole world.".
All of those themes are so beautiful.
and they're the heart of the gospel..
And so we should enjoy them,.
just like our favorite foods at the feast..
We should go to those places every year.
and be encouraged and blessed by them..
But we can't forget about some of the other themes..
And that's what we're gonna talk about today,.
a theme that is here in the story.
that maybe sometimes we just forget about.
because some of the other themes are so thrilling to us..
And that is that Jesus came as king..
That when Jesus came to the earth,.
the king was born, the king of the kingdom..
And what that means for us,.
because the incarnation and the nativity story is beautiful,.
but it didn't just stop there..
It went on, obviously, to the cross..
But during the ministry of Jesus,.
he talked about the kingdom of God a lot..
He talked about the kingdom that was and is and is coming..

$^{81}$And Jesus gave us this theme..
So that's what we're gonna consider this morning..
We're gonna look at it in Matthew chapter two,.
because interestingly enough,.
even though there's a lot of people in the Christmas story.
who didn't say anything about the king,.
didn't seem to get that,.
this one group of people did..
These people from a whole 'nother land,.
these Gentiles from another land, they were not Jewish,.
but they understood that the coming of Jesus.
was the coming of a king..
And they traveled all this way to come see him.
as the wise men or the magi..
So we're gonna look at the verse 12 verses.
of Matthew chapter two..
It says, "Now after Jesus was born in Bethlehem of Judea.
"in the days of Herod the king,.
"behold, wise men from the east came to Jerusalem,.
"saying, 'Where is he who has been born king of the Jews?.
"'For we saw his star in the east.
"'and have come to worship him.'.
"When Herod the king heard this,.
"he was troubled and all Jerusalem with him..
"And assembling all the chief priests.
"and scribes of the people,.
"he inquired of them where the Christ was to be born..
"They told him, 'In Bethlehem of Judea,.
"'for so it is written by the prophet,.
"'and you, O Bethlehem, in the land of Judah,.
"'are by no means least among the rulers of Judah..
"'For from you shall come a ruler.
"'who will shepherd my people Israel.'.
"Then Herod summoned the wise men secretly.
"and ascertained from them what time the star had appeared..
"And he sent them to Bethlehem,.
"saying, 'Go and search diligently for the child..
"'And when you have found him,.
"'bring me word that I too may come and worship him.'.
"After listening to the king, they went on their way..

$^{121}$"And behold, the star that they had seen when it rose.
"went before them until it came to rest.
"over the place where the child was..
"When they saw the star,.
"they rejoiced exceedingly with great joy..
"And going into the house,.
"they saw the child with Mary's mother..
"And they fell down and worshiped him..
"Then opening their treasures,.
"they offered him gifts, gold and frankincense and myrrh..
"And being warned in a dream not to return to Herod,.
"they departed to their own country by another way.".
Well, these guys, this is an amazing part.
of the Christmas story and supernatural and miraculous,.
like so many other parts of the Christmas story..
And there's a lot of stuff we could talk about here..
We could talk about the star..
That's one of my favorite parts.
of the story of the wise men..
We could talk about just them..
They're very interesting.
'cause there's so little information about them..
Obviously, in tradition,.
people assume there was three of them,.
but that's not told to us in the scripture..
It's just the three gifts..
So people assume there was three people..
But there probably actually were more..
They probably would have traveled in a much larger group..
And even if there were just three,.
they would have had a lot of servants with them..
So this was not just a couple of people..
This was a large group of people,.
which is one of the reasons that it says in verse three,.
"When Herod the king heard this,.
"he was troubled in all Jerusalem with them.".
Whatever was going on was big enough.
that throughout the city of Jerusalem,.
people were taking notice..
So this was a large group of people that came..

$^{161}$But think about the journey that they went on..
These are not Jewish people..
We don't know exactly where they're from..
They're from the east,.
but probably from the region of Persia,.
which there was a large Jewish community in Persia.
that had been there since the time of the exile..
And so that's probably how these men had been exposed.
to the scripture and to the prophecies..
But they traveled an enormous distance at enormous cost,.
gave up months of their life to pursue something.
that they couldn't see,.
that they didn't really know what they were pursuing after..
And they were doing it on the basis.
of their faith and belief.
in the prophecies of the Hebrew scripture..
And so these are actually pretty extraordinary individuals.
in an extraordinary part of the story..
And one of the things that's really interesting.
is that here they come into Jerusalem.
with all these people who are Jewish,.
who know the scriptures,.
who have been anticipating the Messiah,.
and yet they're not getting it..
They're not seeing that the king is being born.
and that this is happening right now..
And the guys who do get it are these people.
from a distant land.
that God has supernaturally drawn and brought..
And one of the things that lets us know.
when we think about the king,.
Jesus as the king and the kingdom that he came to bring,.
is that if we're gonna see the king.
and know what's happening,.
we've got to believe in the kingdom..
That was one of the real differences.
between the wise men, the Magi,.
and the rest, Herod and the rest of the people.
who just didn't know what was going on..
The Magi obviously had been drawn by the spirit of God..

$^{201}$They didn't just come up with these things on their own..
God drew them through the scriptures.
and drew them towards God's truth,.
and they believed it by faith..
They had to have enormous faith to begin this journey.
and to continue on this journey..
When you think about how long their journey was,.
they had to start at some point,.
had to have a lot of faith just to start,.
but then they traveled through the desert.
and traveled and traveled and traveled,.
and they had a lot of faith to keep going, to continue,.
to not say, okay, this was a fun camping trip,.
we enjoyed it for a while, but it's time to go back.
because it's not making sense to me anymore..
Things are getting tough, things are getting challenging..
These men believed in the kingdom..
They believed that it was real..
They believed that they were going to see,.
as they said in verse two,.
where is he who's been born king of the Jews?.
They believed that this king was being born..
That's why they were traveling..
And if they didn't, they wouldn't have made this journey..
But the other people who didn't see it, who didn't get it,.
didn't believe in this kingdom..
They thought the kingdom of the Jews was Herod..
They thought the king of the Jews was right in front of them..
Herod certainly thought that..
They didn't see the kingdom that God was beginning,.
and so they didn't have any idea.
that the king was being born..
They didn't believe it, they didn't see it..
Now, that's not surprising..
Jesus, in John 18, 36, he's talking to Pilate..
This is at the very end of his life..
Jesus answered, "My kingdom is not of this world..
"If my kingdom were of this world,.
"my servants would have been fighting,.
"that I might not be delivered, O, to the Jews,.

$^{241}$"but my kingdom is not from this world.".
So it's not surprising that a lot of people missed it.
because Jesus was coming to set up a spiritual kingdom,.
not a physical kingdom..
But that's something that we need to understand..
That's a challenge for us, too..
Because the kingdom of God is not a physical place,.
it's not a country on the map,.
we can miss out on how significant it is,.
and the world around us certainly doesn't see it,.
certainly often doesn't believe in it..
The kingdom of God is a kingdom that has come,.
and that is coming, and that will come..
It's the rule and reign of God,.
and it's existed since before the world began,.
and Jesus made it real in every way when he came to Earth.
and perfectly lived out the values of the kingdom of God,.
and by his life and death, opened the way for us.
to have relationship with God..
The kingdom has come in that sense,.
but it is coming because it is not yet.
fully manifested here on Earth..
It is partially manifested in us..
Since Jesus came and died and started this new group.
of people, this new kingdom that we call the church,.
as we live under the rule and reign of God in our lives,.
and we live out the values of the kingdom,.
the love, the joy, the peace that God has called us to,.
we manifest the kingdom of God, right?.
We manifest the rule and reign of God in our lives,.
but there's a whole lot of people on our world.
who are not living under the rule and reign of God,.
and so the kingdom of God is not fully manifest,.
and that's why we continue to have a lot of problems,.
a lot of brokenness, a lot of injustice and war.
and all these terrible things,.
because the kingdom of God is not fully manifest,.
so even though it has come in Christ, it is still coming..
It is not fully manifest yet..
It will be fully manifest one day in heaven,.

$^{281}$and we know what the scripture says.
in Revelation chapter 21, these verses that say,.
in verse four, "He will wipe away every tear.
"from their eyes, and death shall be no more..
"Neither shall there be mourning nor crying.
"nor pain anymore, for the former things have passed away..
"And he who was seated on the throne said,.
"Behold, I am making all things new.".
The kingdom will be fully manifested in the future,.
and every brokenness will be healed,.
and there will be no more tears,.
and there will be no more death or sin or injustice,.
but we're not there yet, so the kingdom has come..
It is coming right now through us,.
and it's coming in the future in heaven,.
but because it's not a physical kingdom,.
to believe in it requires faith,.
just like the wise men had to really believe.
that at the end of their journey,.
they were gonna find this king,.
otherwise they would have never stayed on the journey..
We have to believe by faith that God's kingdom is real,.
and we have to believe that what God has said.
in his word is true, that God did lovingly create us.
in his image, that we have fallen away because of our sin,.
that we need a redeemer and a savior,.
that God in his love sent Jesus to be that savior..
We have to believe those things by faith..
If we don't, we're not gonna see the kingdom..
We're gonna miss it..
We're gonna miss what's happening right in front of us,.
just like Herod and all the other people..
It was happening right in front of them,.
and they totally missed it,.
because they weren't looking with eyes of faith.
to see what God was doing,.
and so we have to look with eyes of faith and say,.
yeah, we believe what God is saying..
There is a kingdom here..
It's a spiritual reality that's bigger.

$^{321}$than anything in our physical reality,.
but it's not visibly manifested right in front of us,.
but we believe it by faith,.
and because we believe it by faith, we can know the king..
That's the next thing that we see here,.
is that if we want to know the king,.
then we have to surrender to his reign..
It's one thing to believe in the kingdom and to say,.
okay, I believe that it's real,.
but it's a whole 'nother matter to say,.
not only do I believe the kingdom is real,.
but the king is my king,.
and I'm gonna surrender to the reign of this king in my life,.
and if we think about what it means.
to be a follower of Jesus,.
that's what it means to follow after him..
That's what it means to be a Christian..
To come into the kingdom is what happens by grace.
through faith, where we acknowledge and say,.
I know I've sinned, I need a savior,.
and we trust in Christ to be our savior..
That's how we come into the kingdom,.
but to live day by day as a follower of Jesus,.
we have to consistently, moment by moment,.
live out this process of surrender and saying,.
God, you're the king of my life,.
and I wanna live under your rule and reign,.
and that means I'm gonna live according to your values.
and your ideas, not my own or not anyone else's,.
and if we're gonna know the king, that's what we have to do..
Now, that makes us uncomfortable, right?.
Even thinking about living under someone's reign.
makes us uncomfortable..
There's a lot of reasons for that..
Part of that is cultural, 'cause we live in America.
where we're very uncomfortable with monarchs, right?.
And there's nothing wrong with that..
We're uncomfortable with monarchs for a good reason, right?.
Because we're very familiar with tyrannical forms of rule,.
and we're very uncomfortable with that,.

$^{361}$but this is a different type of kingdom..
The kingdom of God is a perfect reign.
from a perfect kingdom with a perfect king..
There's no tyranny involved at all..
In Matthew 20, Jesus speaks about the type of leadership.
that he gives..
He's talking to the disciples..
He says, "You know that the rulers of the Gentiles.
"lorded over them, and their great ones.
"exercised authority over them..
"Shall not be so among you,.
"whoever would be great among you would be your servant.".
And then he says, "Even as the Son of Man.
"came not to be served, but to serve,.
"and to give his life as a ransom for many.".
Jesus himself said he didn't come to dominate people..
He came to serve..
This is a different type of king..
But just because he's a servant,.
don't think that he's not a king..
We sang that Andrew Peterson song earlier, "Is He Worthy?".
And that song quotes some of the scripture.
from Revelation 4 and 5..
In Revelation 4 and 5, you have this vision..
John has this vision of heaven and what's going on there..
And in chapter four, he talks about the throne.
and the worship of the creator..
I'm not gonna read it right now,.
'cause there's a lot of verses..
But if you have time, you can look at it later..
And he describes these amazing creatures,.
these amazing angelic creatures,.
that are worshiping God every moment of every day..
It says they do not cease to worship God.
in verse eight of chapter four..
And all they say is, "Holy, holy, holy.
is the Lord God Almighty,.
who was and who is and who is to come.".
And there's just this amazing scene.
of power and majesty and worship going on..

$^{401}$And then in chapter five,.
it describes Christ, who's sitting on the right hand of God..
And there's this angel in verse two.
who's proclaiming with a loud voice,.
"Who's worthy to open the book and to break the seals?".
And it says, "No one," verse three,.
"No one in heaven or on the earth or under the earth.
was able to open the book and to look into it.".
No person, no human, no angel, no demon,.
no force, no power, anyone in the entire universe.
has the power to open this book..
And John begins to weep in verse four,.
because no one can open the book..
And one of the elders says to him, verse five,.
"Stop weeping..
Behold, the lion, this is from the tribe of Judah,.
the root of David, has overcome,.
so as to open the book and its seven seals.".
And then we go on in the chapter,.
and the book is opened, and in verse 11,.
it says, "And I looked and I heard the voice.
of many angels around the throne,.
and the living creatures and the elders,.
and the number of them was myriads and myriads.
and thousands of thousands, saying with a loud voice,.
"Worthy is the lamb that was slain.
to receive power and riches and wisdom.
and might and honor and glory and blessing.".
Every being, every power in heaven.
begins to praise Jesus, because of his worthiness,.
the worthiness of his sacrifice, the lamb who was slain..
And then verse 13, "And every created thing,.
which is in heaven and on the earth.
and under the earth and on the sea,.
and all things in them, I heard saying to him.
who sits on the throne unto the lamb,.
be blessing and honor and glory and dominion.
forever and ever.".
Everything else in creation joins this song of worship..
First you have all of the heavenly beings.

$^{441}$and the heavenly powers praising Jesus,.
and then you have everyone on earth and every creature.
and everything on earth begins to join this song of worship..
Verse 14, "And the four living creatures kept saying,.
"Amen, and the elders fell down and worshiped.".
Every thing in the universe worships the king..
So yes, Jesus is a servant..
He loves us, he cares for us,.
but don't think that he's not a king, because he is..
Scripture makes that so clear..
And the world tempts us to choose to surrender.
to other things and tries to get us to serve.
other imposter kings, other false kings..
And the world lies to us, the world says,.
"Well, the choice is between either serving God.
"and obedience to God and freedom.".
It gets old to serve God, it gets old to be a part of church,.
there's stuff that you're supposed to try to live right,.
and that's not fun..
Be free and come over here and do whatever you want, right?.
That's the false choice that the world gives us..
Well, that choice is a lie..
It's a lie because there's no freedom from.
without a freedom to..
And it's not a choice between serving God and being free..
It's a choice between serving God.
and serving some other imposter king, that's the choice..
Are you gonna serve the true king,.
who by the way as we just read is a loving and good king,.
who will lead you into, in a loving and good way,.
into abundant life, or are you gonna serve.
some imposter king?.
Are you gonna serve self, or serve materialism,.
or serve power, or money, or achievement,.
or the influence of other people?.
You're gonna serve somebody, right?.
Bob Dylan wrote a song about it..
It's actually a pretty good song if you've never heard it,.
you can listen to it sometime,.
called "Gotta Serve Somebody,".

$^{481}$and that's the point of the song,.
is saying listen, you're either gonna serve God.
or the devil, but it doesn't matter who you are,.
you can't get around it..
You're gonna serve somebody in the end..
And it's true..
That's the choice..
It's not between surrender to the reign of the king.
and freedom..
It's between surrender to the reign of the king.
and serving somebody else, or someone else..
And what we will find is that every one.
of those imposter kings, no matter what it is,.
will end up being a tyrannical ruler in our lives,.
and will harm us and ultimately destroy us..
The only good king is King Jesus..
He's the only one..
And the New Testament talks about the kingdoms,.
it uses this metaphor of the kingdom of darkness.
and the kingdom of light..
Colossians 1, 12 to 14 says this,.
"Give me thanks to the Father who has qualified you.
"to share in the inheritance of the saints in light..
"He has delivered us from the domain of darkness.
"and transferred us to the kingdom of his beloved Son,.
"in whom we have redemption, the forgiveness of sins.".
And that's the king and the kingdom that we wanna be in..
The kingdom of light, the kingdom of his Son..
Now this passage doesn't tell us.
that the wise men made this surrender..
We know they believed in the kingdom enough.
to come see the king..
We don't know necessarily if they, by faith,.
trusted in Jesus..
Perhaps they did..
I think you could make that argument, certainly..
We don't know for sure..
They certainly were sensitive to God and obedient to God..
But for us, obviously salvation is a one-time event..
We don't have to do that over again..

$^{521}$We come and we put our trust in Christ..
That happens once..
But the surrender to the king is a moment-by-moment process..
That's not a one-time event..
We don't say, "Okay, Jesus, I wanna follow you..
"I'm surrendering to your reign..
"I'm done..
"I'm never gonna be tempted to walk out.
"out of bounds ever again.".
No, it's a moment-by-moment process of saying,.
"God, I need to surrender to you in this area of my life..
"I'm tempted to follow some other imposter king..
"Maybe it's me..
"Maybe it's something else, but I've gotta surrender to you.".
We've gotta constantly be checking where our allegiance is..
And it's a repeated surrender..
And if we live under the reign of the king,.
then we experience the blessings of the kingdom..
There's so many things that we could talk about,.
and the most significant ones are the forgiveness of sins,.
the transformed life that comes through Jesus,.
the eternal life that's promised..
But even if you think about the advent themes.
that we often talk about during advent,.
the candles that represent hope, peace, joy, and love,.
those are things that come to us.
as we live under the reign of the king..
We have hope because the kingdom has come,.
and it is coming..
We have peace because the king who loves us.
is in perfect control of our world, and we are his..
That gives us peace..
We have joy because every need is met.
in our relationship to the king..
We have love because we're freed.
from the transactional nature of broken relationships.
that we experience in this world,.
and we're free to truly love the king and others.
without fear..
We can experience true love..

$^{561}$Now, the world desperately needs the blessings.
of the kingdom..
We want 'em as individuals,.
but the world desperately needs that, right?.
We look at our world, and we see war,.
we see suffering, we see injustice,.
we see all these terrible things,.
and sometimes our heart cries out,.
and we say, "God, why is this happening?.
"How long?.
"When are you going to change things?.
"If only there was some way.
"that we could be changed as human beings..
"If only we could be transformed to be a lot less selfish.
"and stop hating each other and fighting each other..
"If only there was a way for this to happen.".
And God says, "Yeah, I've made a way,.
"but if you wanna experience the blessings of the kingdom,.
"you gotta surrender to the reign of the king.".
And we need to understand that..
We desire the blessings of the kingdom..
We crave them..
We chase after them..
But we all wanna find them.
by building our own tower of Babel to heaven.
instead of just going and surrendering to the king.
who is gonna offer those blessings to us..
So as we close out, I just wanna be practical for us..
How do we continue in this process.
of surrender to the reign of the king?.
If it's a process, moments by moment,.
we say, "Okay, well, if I'm a follower of Jesus,.
"I wanna do that.".
What does that look like in our lives?.
How do we do that?.
In the scripture, in the Old Testament and the New Testament,.
there's this metaphor that's used so much.
about the path or the way.
that describes the lifestyle of following Jesus..
And the early Christians actually were called.

$^{601}$followers of the way..
And so that metaphor is really strong in our faith..
I think that's a good illustration for us.
to think about following Jesus kind of as a map,.
as a pathway, right?.
And scripture gives us turn-by-turn directions..
Like if you think about GPS or your Maps app,.
it tells you how to get somewhere,.
it gives you turn-by-turn directions, right?.
Scripture gives us all kinds of turn-by-turn directions.
about how we're to relate to God,.
how we relate to one another,.
how we're to live our lives.
in all these different areas of our lives, right?.
And God's given us the Holy Spirit to speak to us.
about staying on course.
and staying in fellowship with Him,.
just like the Map app talks to you, right,.
when you miss a turn.
and tells you that you missed the turn, right?.
The problem is a lot of times we don't,.
in the illustration that we're using here,.
we don't consult the Map app.
and we don't listen when it tells us.
we've made a wrong turn..
We just keep doing our own thing, right?.
And that's what it means for us.
to moment-by-moment surrender to the King,.
to continue to soak ourselves in Scripture.
and go back to what God has said and say,.
"God, am I living under your rule and reign,.
"or am I following some other alternate map.
"that's provided by me or by my friends.
"or the world around me or culture or some celebrity?.
"Am I following another map, or am I following your map?".
And we listen to the Holy Spirit,.
and the Holy Spirit is good,.
and He lets us know when we're out of fellowship.
and we feel spiritually distant from God.
or we don't wanna be in church, we don't wanna worship..

$^{641}$God lets us know..
But we have to listen to that.
and not be like I am sometimes when I miss my turn.
and then I drive 10 more miles, right?.
Sometimes in life, I get out of fellowship with God.
and six months later, everything's unhealthy,.
and I finally figure out, oh yeah, something's wrong..
And God's like, "Yeah, I've been trying to let you know.
"for the last six months..
"You made a wrong turn back there, way back there.".
Right?.
And so that's how we surrender.
to the reign of the King moment by moment..
We keep going back to Scripture and saying,.
"God, am I aligning with what your map is for my life?".
We listen to the Holy Spirit..
"God, can you communicate to me in my heart.
"and let me know when I'm making a wrong turn.
"and get me back on course?".
And then we surrender to it, we obey it..
We say, "God, you know what?.
"I don't really like what I'm reading right now..
"I don't really like what the Holy Spirit's telling me..
"I don't want to be unselfish..
"I don't want to care about someone else..
"I wanna make this about me,.
"but I'm gonna surrender to it because you're the King.
"and I wanna live under your rule and your reign.".
Let's pray..
God, we just thank you that you came as Savior..
You came to rescue us, to redeem us..
But we also thank you that you came as King..
Help us to not just receive you as Savior.
and rejoice in the forgiveness of our sins and stop there..
Help us to go on to acknowledge that you're the Lord,.
you're the King, you're the ruler,.
and help us to choose to live under the rule.
and reign of your kingdom,.
to manifest the values of your kingdom in our lives.
to those around us, God..

$^{681}$In Jesus' name, amen..
\newpage



\section{路加福音}
\label{sec:iKrjLPh_hzM}
\textbf{將臨節的挑戰 - 耶穌出生的報喜 (路加福音1-2) \#2 - 陳崇基牧師}
\newline
\newline
連結: \href{https://youtube.com/watch?v=iKrjLPh-hzM}{\texttt{https://youtube.com/watch?v=iKrjLPh-hzM}} ~~~~ 語音日期: 2023-12-12
\newline
\newline
\hyperref[sec:aFC7A7v5zzs]{\small{< < < PREV SERMON < < <}}
~
\hyperref[sec:index_chronic]{\small{[返順時目]}}
~
\hyperref[sec:index_scriptual]{\small{[返順卷目]}}
~
\hyperref[sec:JrlUbBSqsPA]{\small{> > > NEXT SERMON > > >}}
\newline
\newline
$^{1}$我們繼續Evan的系列.
今次第二講.
我們講到.
《盧家福音》第一,二章.
我們上一次講到特別的英孩.
報喜就是施洗若寒.
是一章五至二十五節.
今次我們看到.
另一個英孩的報喜.
他出生的報喜就是耶穌我們的救主.
是一章二十六至三十八節.
我們重申一點.
在聖誕節來到的時候.
我們需要去復原一件事.
記不記得是什麼.
是一個奇蹟.
那個驚喜.
就是一個新的生命的誕生.
還是我們神的兒子.
救贖主耶穌基督誕生.
那種驚喜上帝為我們預備的.
上帝所做的.
是過於人所想所求.
不只是多於你想的.
跟你想的完全不一樣.
這就是上帝在聖誕節給我們的驚喜.
耶穌出生的報喜.
我們一起來看看.
一章二十六節.
到了第六個月.
天使加伯列奉神的差遣.
往加利利的一座城去.
這城名叫那撒勒.
到一個童女那裡.
她已經許配戴維加的一個人名叫約瑟.
童女的名字是瑪麗亞.
天使進去對她說.
蒙大恩的女子你好.
主和你同在.
瑪麗亞因這話就很驚慌.

$^{41}$又反覆思考這些問候是什麼意思.
這裡說到了第六個月.
就是什麼時候呢.
就是說之前的報喜.
就是施洗約翰的報喜.
六個月之後.
這個時候.
以利沙白.
就是她的表姐.
或者表姨媽.
總之是親戚.
已經有新人大概六個月的時間.
六個月是不是要見肚子了.
是不是.
六個月應該見肚子了.
因為我不知道.
六個月應該見肚子.
我們中年發福.
我們就知道很多年會見肚子.
但是那些懷孕的應該六個月就見肚子.
這裡就是天使加巴列.
就去了拿撒亞這個地方.
和瑪麗亞報喜.
說她會有上帝的恩典臨到她.
你想象一下.
以前的人和我們今天的人都一樣.
我們人生裡面總會有計劃的.
雖然我們不是每個人都一樣.
不是每個人都是one size fits all.
但是我們每個人在不同的階段.
我們都會想不同的東西.
我們人生階段到了這裡.
我們就應該要上學.
然後就要畢業等等等等.
不是一定每個人都一樣.
但是每個文化每個時代.
都會有一個人為的.
人生的一些預料計劃.
今天是這樣.
當時也是這樣.

$^{81}$考慮到這裡的時候.
你就可以嘗試去揣測一下.
這裡有什麼驚慌.
有什麼驚訝.
瑪麗亞是出自一個.
敬虔的猶太人家庭.
這個我們知道的.
因為之後她所說的話.
或者她唱的贊尼斯.
都是一些經典的詩篇.
她應該很明白舊約裡面.
一些的應許.
譬如大衛家的應許.
大衛國度.
以色列的君王要嫁到等等.
這些盼望這些的應許.
瑪麗亞應該很清楚.
因為她是一個敬虔的猶太人的家庭.
所以我問一下你.
特別當耶穌出生之後.
她帶耶穌去到耶路撒冷的聖殿.
她所獻的祭物.
是兩隻班戈或者初甲.
按照摩西的律法你就知道了.
如果只獻兩隻初甲.
那就是什麼呢.
就是貧苦大眾.
因為摩西的律法裡面說到.
如果你負擔不住的話.
你就獻上一隻初甲.
另外瑪麗亞的年紀有多大.
這裡.
大概有多大.
我們每個人都比她大.
應該十二三歲.
可能十四五歲.
為什麼因為以前的人.
到了十二三歲的時候.
他們就會談婚論嫁.
然後就建立家庭.

$^{121}$是不是比我們大.
比我們小.
比我們小很多.
所以是一個高中女孩.
可能是專一.
但居然她去到.
接受了這個報喜.
這裡所謂講一章二十七節.
同女已經許配了大衛家的約瑟.
我想這個時候.
嫁妝都已經講了.
準備了.
可能幾個月後.
她就要下嫁給約瑟.
是不是.
當時平凡的男女.
都是預料中事.
每個人都可以做這件事.
搞定這些東西.
應該是順理成章的.
原心計劃都是這樣的.
去到那個階段.
他們就應該很順利.
相看了.
約瑟,瑪麗亞.
搞定了兩家人.
講好了所有的東西.
禮服,婚宴都談好了.
天使就來報喜.
一章三十節.
天使對她說.
瑪麗亞不要怕.
在神面前你已經蒙恩了.
你要懷孕生子.
要給他起名叫耶穌.
他將要為大.
稱為至高者的兒子.
主神要把他祖先大衛的王位給他.
他要作那國家的王.
直到永遠.

$^{161}$他的國沒有窮盡.
瑪麗亞對天使說.
我沒有出嫁.
怎麼會有這事呢.
天使回答她說.
聖靈要臨到你身上.
至高者的能力要庇佑你.
因此那要出生的聖者.
要稱為神的兒子.
況且你的親戚以利沙白.
就算那數來稱為畢生育的.
在年老的時候也懷了男胎.
現在懷孕六個月了.
因為出於神的話.
沒有一句不大能力.
我想像一下.
十幾歲的一個女孩.
去接到一個這樣的消息.
對她來說是非常驚天動地.
我想你們可以想像一下.
到底在一個少女的腦海裡.
她在想什麼.
可能在想為什麼要選中我.
為什麼是我.
首先第一件事.
瑪麗亞這個名字.
或者叫瑪麗亞的女孩.
通街也是.
按照我們當時的名字的研究.
猶太人的社會裡面.
50\%以上的人都是叫瑪麗亞的.
瑪麗亞這麼多個.
為什麼偏偏選中她呢.
跟她年紀差不多的瑪麗亞.
也有很多的.
整班同學裡面.
有半個都是叫Mary.
還有他們預料是談婚論嫁的男女.
我相信也有不少的.
雖然拿撒亞這個地方比較小.

$^{201}$但是你看看整個區域.
其實不止這一對的男女.
是不是.
可能也有很多弱勢瑪麗亞的男女.
是不是.
偏偏就選中她.
為什麼呢.
上帝主權.
三分天下.
上帝就是要選她.
出於上帝的召命.
出於上帝自己的主權.
我想最困擾的問題就是.
他問的問題34節.
我沒有出嫁.
怎麼會有這事呢.
他訂了婚.
但還沒有出嫁.
還有幾個月他才出嫁.
怎麼會有這事呢.
當然這個是神跡的發生.
聖靈的感人.
上帝在瑪麗亞身上做的事.
不單單是一個驚喜.
實際上是一個中斷.
是不方便.
不單止不方便.
是一個很大的打擾.
上帝突然間打擾你的計劃.
突然間中斷你心裡面所想的計劃.
正如我所說.
當時的社會.
大家都是在想和預料.
這個時候應該要做什麼.
瑪麗亞弱勢.
這個時候應該要談婚論嫁等等.
但上帝來到的時候.
就將它打斷.
並且在社會當中.
是很難交代的.

$^{241}$說沒有出嫁就有新人.
是很難交代的.
你不要說以前的社會.
現在的社會都很難交代.
你想像一下我和我太太訂婚之後.
失驚無神.
她說我有了新人.
你信她還是信我.
還是信那支棒.
我不知道.
是很難的.
你知道那個故事.
就是約瑟想在暗地裡將瑪麗亞休了.
但天使也都介入.
然後他遵行了上帝的旨意.
就將她娶過來.
這個行動絕對不簡單.
對當時來說.
是社會裡最大的羞恥.
你將她娶過來.
還是要很多年之後.
你才明白上帝的心意.
在這麼多年裡.
那種中斷的打擾.
真的不簡單.
還是現在開始的時候.
瑪麗亞怎樣回應呢.
38節.
瑪麗亞說我是主的侍女.
願意照你的話實現在我身上.
於是天使就離開.
這個回應和之前路祭司薩格尼亞截然不同.
薩格尼亞雖然懂得所有的律法.
知道所有舊約裡面所有的promises.
但他就說How can it be?.
No way.
所以就由不信之後成為信.
這裡瑪麗亞一個十幾歲的小女孩.
她說什麼呢.
她說願意照你的話實現在我身上.

$^{281}$如果你是有天主教背景的話.
你可能聽過這句話叫Marie's Fiat.
可能在天主教以前讀中學.
或者有天主教傳統的人.
你應該可能聽過這句話.
或者這種神學的概念.
瑪麗亞的兒子聖恆.
你的旨意行在我身上.
Marie's Fiat.
沒有天主教背景的可能你都聽過Fiat是什麼.
是車來的.
意大利的車.
我講道的時候都take了個break.
看看不同的車.
為什麼呢.
因為那時候我曾經買過一輛Mini Cooper.
我買Mini Cooper之前.
我就想買Mini Cooper一定買Fiat.
結果我買了Mini Cooper.
Fiat是意大利的車來的.
剛好由意大利.
天主教.
UV電影節目對名車應該是非常之看慣.
都不是很surprise.
但Marie's Fiat當然不是Marie的車.
什麼叫Marie's Fiat.
原來在主後四世紀的時候.
有一個拉丁文譯本叫做Volga.
中文叫做武加大.
武加大.
武加大武加世.
武加大譯本.
Volga是什麼意思.
就是common language的意思.
就是通用的語言.
當時就將聖經翻譯成當時的拉丁文.
因為當時拉丁文就是通用的語言.
一張三十八字願意照你的話實現在我身上.
第一個字在Volga裡面就是Fiat.
什麼意思.

$^{321}$就是Let it be.
Let there be.
願意.
讓他成就.
讓他去吧.
讓他行出來.
你記得在創世記.
上帝講第一句說話是什麼.
創世記第一章創立天地的時候.
上帝講的說話.
記不記得.
要有光.
中文翻譯要有光.
英文是什麼.
Let there be light.
在Volga裡面就是Fiat.
Lux.
Fiat.
Let there be.
Lux.
Light.
拉丁文.
所以為什麼會有一個這樣的概念.
就是Mary's Fiat.
當瑪利亞被震撼了她世界的時候.
被上帝的主權去打斷了她一切的計劃的時候.
這個小妹.
她就說Fiat.
第一個字.
Fiat.
願者是成就.
你有沒有想過那個Wonder.
一個少女.
沒有任何人生的經驗.
但是她聽見天使向她報喜的時候.
她願意用信心去接受.
救主可以道成肉身.
就是因為她的Fiat.
因為她的願意.
Yes 上帝來打斷我的計劃.

$^{361}$去中斷我有的預料.
去震撼我這個小的世界.
但是Let it be.
Let there be.
是上帝對她的心意.
上帝對她的心意的話語.
當然天主教的教義裡面.
如果是天主教背景的人.
你都知道瑪利亞的地位.
很多時候將她無限放大.
天主教的神學裡面說.
瑪利亞重要到一個地步.
是Co-redemptrix.
就是共同的女救主.
因為她願意.
所以救主能夠降生.
或者Mediatrix.
或者一個女的共同的中保.
或者Perpetual virginity.
終生的童生.
就是說她生了耶穌之後.
就沒有生其他小孩.
或者沒有原罪的理論.
Immaculate conception.
就是說瑪利亞出生的時候.
她是沒有原罪的.
我們每個人都有原罪.
都在罪惡的狀況底下.
Original sin.
但瑪利亞在天主教的教義裡面.
就是無罪之身.
或者有其他升天等等.
瑪利亞沒有死等等.
我們不需要將她無限放大.
看聖經怎麼說.
聖經怎麼說瑪利亞.
或者瑪利亞是一個怎樣的角色.
當然是一個很重要的角色.
但是否重要到這樣呢.
一章46節.

$^{401}$瑪利亞唱出她讚美詩歌的時候.
她說我心尊主為大.
我靈以神我的救主為樂.
她有沒有救主呢.
她都有救主.
耶穌基督就是她的救主.
提摩那前書2章4節更清楚.
上帝願意人人都得救認識真理.
因為只有一位神.
神和人之間.
幾多個宗婆.
一個.
Mediator.
成為人的基督耶穌.
我們不需要全盤去接納這些.
所謂天主教有的.
神學的想像.
有時候需要神學想像.
但有時候神學想像會太過份.
我們始終都要回到聖經裡面.
但同時我覺得我們不需要去貶低瑪利亞.
這個雖然是靚妹.
少女.
但她對上帝救贖的計劃.
縱使是中斷她人生的計劃.
帶來很多的震撼.
但她說Fear.
我肯.
Let it be.
路加福音第一章開始.
整本的路加福音裡面.
信心的人物你可以數.
但你數第一個.
就是這個少女.
這個少女.
有這個Let it be.
願意你的旨意行在我的身上.
到底瑪利亞特別在什麼地方.
如果不是無罪之身.
又不是共同中保等等.

$^{441}$你再看看以下的情節你就知道.
她特別的地方在哪裡.
我們來到一章39節.
在那些日子.
瑪利亞她收到報喜.
Fear之後.
她就急忙前往山區.
如果她有車的話.
應該開Fear.
來到猶大的一座城.
進了撒迦利亞的家.
向伊利沙伯問安.
她的表姐或者她的親戚.
伊利沙伯一聽見瑪利亞問安.
所懷的胎就在腹裡面跳動.
伊利沙伯被聖靈充滿高聲的喊.
說你在腹裡中是有福的.
你所懷的胎也是有福的.
我主的母親到我這裡來.
為何這是臨到我呢.
一件很大的福氣.
主的母親來到探訪她.
因為你問安的聲音一入我腹裡的胎.
就歡喜跳動.
這相信的女子是有福的.
因為主對她所說的話都要應驗.
這裡有兩個母親相遇.
伊利沙伯和瑪利亞.
施洗約翰在哪裡.
在伊利沙伯的肚子裡.
獨個兒懷胎.
應該已經懂得動了.
施洗約翰在母親的肚子裡.
兩個母親相遇的時候.
肚子裡的施洗約翰.
他就歡喜跳動了.
等於有些有新人的姐妹.
有時候.
說回我自己的那個.
我老婆.

$^{481}$她懷胎的時候.
她會動的.
她撐腿等等.
是不是很Wonder.
你可以當她很Creepy.
但她還是很可愛.
是不是.
就是這種情況.
他聽見瑪利亞的聲音的時候.
主母的聲音.
就在胎裡歡喜的跳動.
你知道施洗約翰是什麼角色嗎.
是耶穌的先鋒.
第六個月大的時候.
施洗約翰已經看到.
他要做他先鋒的那位.
救主耶穌.
他已經在做那件事.
Herald is coming.
原來我要介紹那個.
要準備那個.
就是這個救主.
這個先鋒在肚子裡.
已經開始跳動和歡喜.
但這裡最重要的是.
怎樣形容瑪利亞呢.
一張四十五字.
這相信的女子是有福的.
伊利沙伯這樣形容她.
Blessed is she who believes.
這相信的女子.
瑪利亞有什麼特別.
瑪利亞特別的地方.
不是因為她本質上和其他人不同.
不是她出生無怨罪.
不是她的本質比其他人更加.
聖潔更加神聖.
不是因為她本質與眾不同.
是這相信的女子.
這位女子有的信心.

$^{521}$對上帝的信仰.
對上帝的信求.
The fear.
上帝的旨意臨到她身上的時候.
她說Yes 願意.
信心的榜樣.
一個拿撒亞的女子.
居然有這樣的Faith.
這樣的Belief.
但這些Belief不是盲從的Belief.
很多時候我們以為.
聖經所說的信仰 相信.
就是你不要想太多.
你閉上眼跳過去就行了.
聖經所說的信仰不是這樣的.
不是一種非理性或反智的信仰.
聖經的信仰說.
我們所相信的.
就是上帝所說的話.
就是聖經所啟示的一位上帝.
他是信實的.
我們是有理據去相信的.
你可以說一下歷史.
說一下其他東西.
但最重要的是上帝的Character.
上帝既然是Faithful.
他的本性是信實的.
我們就相信他 是不是.
他在華語裡面所啟示的.
我們就相信他.
這個就是Faith.
瑪利亞也是一樣.
她接到這個報喜的時候.
她就說Yes.
這個對我來說是非常麻煩.
要很多年掙扎才知道結果.
有很多Uncertainties.
But Faith.
去過拿撒拉的人.
很喜歡去拿撒拉一個地方.

$^{561}$就是那個大教堂.
Church of Annunciation.
這個大教堂蓋著的地方.
就是傳統上.
紀念就是拿撒拉這個小的村莊.
當然現在很大.
以前這個小的村莊那個地方.
就是天使加伯尼向瑪利亞報喜的地方.
這個教堂是一個天主教教堂.
但很有神秘感.
因為是分兩層.
你看到那個空.
就是那個教堂.
這個就是報喜堂.
分開幾層.
教堂的頂在裡面拍上去的時候.
就有光線從上而下射下來.
光是從天降下的.
上層呢.
右邊那裡.
右上角.
就是一個很華麗的大教堂.
可以容納很多人.
有時候在那裡有彌撒.
有敬拜聚會.
連接上下層的.
左下角.
是一個八角形的天窗.
叫做Oculus.
天窗.
眼的意思.
Oculus.
底層就是入口的地方.
右下角.
是很暗的.
遊人就會從底層開始進去.
是黑漆漆的.
當然拍照拍得很光.
進去那個地方.
有些人在排隊看.

$^{601}$相傳就是瑪利亞的居所.
當然我們不是很清楚是不是.
但拿沙特的村莊不是很大.
不是那裡.
都是附近.
這個就是建築神學.
Architectural Theology.
在教堂的建築裡面.
去到講出神學.
怎樣講出神學.
上層好像上帝.
將光照下來.
下層就是人間.
有一個叫Oculus.
穿過.
然後下去最下最下的地方.
就是瑪利亞的地方.
接通著天地之間的.
就是那個Oculus.
信心是什麼?.
整座教堂就好像講這件事.
Faith.
就是上主啟示他的說話.
上主用他的主權.
在你身上要做事.
地上的人就需要回應.
當你看見上帝的光.
看見上帝的啟示.
他的Revelation下來的時候.
上帝希望和我們同工.
Faith.
這個就是瑪利亞特別的地方.
或者它偉大的地方.
他將45節.
這相信的女子是有福的.
The one who believe.
The girl who believe.
因為其實是路加福音.
由頭到尾的主題.
路加福音裡面.

$^{641}$講了很多耶穌工作的時候.
神蹟奇事發生的時候.
人都有不同的回應.
當然有很多人是reject.
太麻煩.
或者我們不接受你.
發你錯人.
宗教領袖等等.
但路加福音從頭到尾都有不同的.
我們叫信心的故事.
Faith stories.
四個人帶著一個毯子.
去到耶穌面前.
耶穌跟他說什麼?.
耶穌看見他們的信.
Faith.
毯子就行走.
羅馬伯夫長.
他的僕人生病.
封聞耶穌.
派人去求耶穌.
叫他說一句話.
我的僕人.
我所愛的僕人就可以得到醫治.
耶穌說什麼?.
耶穌說我見到一個大的信心.
Great faith.
一個滿身罪惡的女人.
去到一個筵席裡面.
流淚.
然後為耶穌抹腳.
耶穌說什麼?.
對那個女子說.
你的信救了你.
血流十二年的婦人.
她已經無望了.
不潔淨了.
整個人生都不潔淨.
他說只要追趕到這個拉比.
摸一摸她的衣裳的話.

$^{681}$我就可以得到醫治.
耶穌轉身來.
他得到醫治之後.
他說你的信.
你的faith.
十個長大麻瘋的病人.
來到耶穌面前求他去醫治他.
被醫治了.
有一個外族人.
回來見耶穌.
將榮耀歸給上帝.
耶穌說什麼?.
你的信救了你.
你的faith救了你.
聖誕就是一個信心的挑戰.
救主耶穌基督來到我們中間.
最終的目的.
除了出生之外.
就是在十字架上成就救人.
去救我們這一班不該救的人.
不應救的人.
信從頭到尾.
都是接受faith.
瑪利亞.
不只是她才有信心.
信心不是一個monopoly.
對瑪利亞來說.
是每個人都可以有.
faith.
上帝在你身上要做的事.
let it be.
你可不可以figure out all the answers.
No, I don't.
你都不可以.
但你還是活著的時候.
上帝的光.
上帝的啟示臨到你的時候.
上帝的旨意臨到你的時候.
你的回應.
瑪利亞怎麼回應?.

$^{721}$let it be.
天使報喜.
annunciation.
是一個challenge.
是一個很大的challenge.
你記得上個星期我們看第一章.
第一個人物.
魯濟斯·薩加利亞.
上帝對他說什麼?.
他說魯濟斯.
你心裡面的上帝太小.
我還在做事.
請你不要忘記.
他對薩加利亞說.
上帝對他說.
我不但要賜給你一個遺址.
就是多年前你曾經有過的禱告.
我不但要賜你一個遺址解除你社會裡面不育的羞恥.
我還要賜一個你所說的先鋒.
你所生的就揭開上帝要做的計劃最後一章.
It's not only a son.
a prophet.
一個new beginning.
上帝說你的上帝太小.
或許你的祈禱太小.
瑪利亞的挑戰.
就是告訴他.
或許你的計劃太小.
I have a much bigger plan for you.
可能你的計劃只不過是按照人間的定律.
不同的階段我就做現在要做的事.
But your plan may be too small.
神說My plan is much bigger.
當你重拾這種wonder的時候.
你會發現這件故事不是那麼簡單.
我們相信上帝在我們生命當中.
在我們這個世界裡面.
他去到實行的計劃仍然是未完的.
It's to be continued.
怎樣continue?.

$^{761}$At this point 不知道.
但是你是不是願意.
以止成行.
Fiat.
願意照上帝的話語.
讓上帝在你身上來作供.
在未知的裡面我們信.
並且可能發出讚美.
應救發出讚美.
在盧卡福音裡面每一個信的故事.
最後都會有所謂wonder and praises.
不等於他們所有東西都figure out.
不過他們會以信和歌頌去榮耀神.
瑪利亞的歌頌是什麼呢?.
就是他的專主頌.
拉丁文是manifica.
manifica.
大家知道英文有manificent.
芝加哥有manificent mile.
就是這個字.
為什麼會叫做manifica.
Again 又是Vogue.
拉丁的翻譯.
第一句瑪利亞唱出的.
就是manifica.
專主偉大的意思.
專偉大manifica.
我想用這段的專主頌來結束.
瑪利亞說.
我心專主偉大.
我靈異神我的救主為樂.
因為他顧念他使慮的卑微.
從今以後萬代要稱我有福.
因為那有權能的.
為我做了大事.
他的名是聖的.
他憐憫敬畏他的人直到世世代代.
他用傍臂施展大能.
他廣散心裡妄想的狂傲人.
他叫有權畏的失畏.

$^{801}$叫卑賤的性高.
他叫飢餓的飽餐未食.
叫富足的空手回去.
他輔助了他僕人以色列.
不忘記施憐憫.
正如他對我們的列祖說過.
憐憫阿巴罕和他的後裔直到永遠.
fear manifica.
可以同心得來祈禱.
主你在這個章林節.
他都告訴我們.
原來你的計劃比我們的計劃不同.
你的計劃比我們的計劃更大.
主讓我們來到你面前.
我們說依旨承行.
我們心靈裡無論今天去到哪個階段.
哪個境況.
我們仍然深信你自己永美的計劃是繼續.
主我們求你.
讓我們有信心來回應你.
fear let it be.
願上帝你自己的旨意.
行在我們身上.
主幫助我們.
讓我們在我們現在未知的當中.
我們仰望你.
因為你是我們的真光.
我們是你的盼望.
主我們來到你面前.
求你給我們這個信服.
信心.
讓我們相信你.
你那位信實的神.
你所說過的計劃.
你所說過的心意.
絕對會實行.
讓我們繼續的依靠.
仰望你.
祈禱奉祖耶穌的名字.
\newpage



\section{}
\label{sec:JrlUbBSqsPA}
\textbf{Reflection on Advent (Luke 1-2) - Rev. Mark Chan}
\newline
\newline
連結: \href{https://youtube.com/watch?v=JrlUbBSqsPA}{\texttt{https://youtube.com/watch?v=JrlUbBSqsPA}} ~~~~ 語音日期: 2023-12-19
\newline
\newline
\hyperref[sec:iKrjLPh_hzM]{\small{< < < PREV SERMON < < <}}
~
\hyperref[sec:index_chronic]{\small{[返順時目]}}
~
\hyperref[sec:index_scriptual]{\small{[返順卷目]}}
~
\hyperref[sec:vk3Vfk1fF_M]{\small{> > > NEXT SERMON > > >}}
\newline
\newline
$^{1}$- All right, we are in the season of Advent.
and actually to next Sunday,.
we're gonna celebrate Christmas Sunday together..
So I think we should turn our attention to passages.
related to the birth of our Savior, our Lord Jesus Christ..
So today we're gonna turn to the Gospel of Luke 1..
And this opening chapters of the Gospel of Luke.
is to help us to pay attention to stories and records.
or histories related to the birth of our Savior..
As you can see in this summary of Luke 1 and 2,.
the opening chapters of Gospel of Luke,.
you see pairs and pairs of stories.
that Luke is trying to tell us.
to prepare for the birth of our Savior..
As you can see, 1.5-25 is the annunciation,.
that means the good news of giving birth of someone,.
annunciation of John the Baptist, the pioneer of Messiah..
And then it is paired with the annunciation story of Jesus,.
our Savior, 1.26-38..
And then there is a story in between 39-56,.
two mothers meeting together, Mary and Elizabeth..
Elizabeth is the mother of John the Baptist.
and Mary is of course the earthly mother of Jesus, our Lord..
And then there will be a birth stories of John the Baptist.
on the left, 1.57-80,.
and then birth story of Jesus, our Lord, 2.1-21..
And then there are a couple of stories again.
to end this opening two chapters..
Jesus in the temple when he was a baby.
being presented to the temple..
And there were two individuals, Simon, Simeon,.
and Anna greeting and blessing the baby, Jesus..
And then there will be an other,.
the other half of the pair of the story is Jesus.
finally going into the temple, his father's house,.
to debate with the teachers, chapter two,.
at the end of chapter two..
Now Advent happens once a year,.
Christmas is only once a year..
And I think when we read stories about Christmas.

$^{41}$in the Bible, Gospel of Luke particularly,.
we need to recover the sense of wonder,.
the sense of awe and wonder..
Anybody seen a baby born in the delivery room?.
Life, I mean..
No?.
Yeah, right?.
That's the sense of wonder..
You don't have to be there to sense that..
You just think about it.
and you will get that wonder, that sense of,.
it's like magic, right?.
The power of life, right?.
Giving birth to a new baby..
And I believe that we need to recover.
that wonder of Christmas, the baby in the manger..
As you can see, pairs of story gives you a sense of symmetry.
telling you the story of the birth of our savior.
under the sovereign orchestration of God..
And everything's under God's control..
Everything is under God's power.
in the birth story of Jesus..
And Jesus came at the right time, in the right place,.
in the right spot..
So we read the story in chapter one, 126..
In the sixth month, God sent the angel Gabriel.
to Nazareth, a town in Galilee,.
to a virgin pledged to be married.
to a man named Joseph, a descendant of David..
The virgin's name was Mary..
The angel went to her and said,.
"Greetings, you who are highly favored..
The Lord is with you.".
Mary was greatly troubled at his words.
and wondered what kind of greeting this might be..
Now in the sixth month,.
that is six months after the annunciation of Gabriel.
to Elizabeth about the birth of John the Baptist,.
the pioneer..
So Elizabeth was pregnant with John the Baptist,.

$^{81}$and six months later, Gabriel showed up to Mary..
Now at this time, Mary is pledged to be married to Joseph,.
a probably a young man..
And of course, Mary was very young as well..
And Joseph was a descendant of David..
And Mary at this point probably is 12 to 13 year old,.
14 year old..
Yeah, back in the days, you know,.
people lived differently, right?.
Back then people get married very, very early on.
and they build a family.
because most of the people died young..
So from different passages of the Bible,.
we know Joseph and Mary, this couple,.
come from a poor Jewish families..
Like every typical young Jewish couple,.
everything seems to have been arranged..
When you hit 13 year old, 14 year old,.
when you were a teenager, Joseph and Mary,.
they were expected to follow some sort of pattern,.
to build their own small family,.
get married, build a small family, and that's it, right?.
That's your life goal..
You enter that stage,.
you automatically will do something.
as expected by the society..
So Joseph and Mary, I would say,.
were expecting just to get married.
and build their own small family.
and have their own child and have their own children.
and then live in Israel.
to grow their own family,.
to join the extended family to celebrate,.
and that is expected for them..
But comes the divine interruption to their life's plan,.
and that interruption is a big interruption..
The angel Gabriel continues, 1 Dirty..
The angel Gabriel said to Mary,.
"Do not be afraid, Mary..
"You have found favor with God..

$^{121}$"You will be with child and give birth to a son,.
"and you are to give him the name Jesus..
"He will be great and will be called the Son of Most High..
"The Lord God will give him the throne of his father, David,.
"and he will reign over the house of Jacob forever..
"His kingdom will never end.".
What an interesting and astonishing news, right?.
You will be pregnant and the child will be what?.
A king, and his kingdom will be forever, right?.
It's astonishing to just think about it, right?.
When my wife was pregnant with the first child,.
nobody said these words to me..
They probably would say to me,.
"Okay, you better get ready, sleep more,.
"because you will be deprived of sleep,.
"and you better save up some money.
"because you will be in debt," right?.
But here, for the announcement for the child,.
he will be great and will be called the Son of Most High,.
and he will be king, and his kingdom will last forever,.
and there is no end for his reign..
It's astonishing..
134, "How will this be?".
Mary asked the angel, "Since I am a virgin.".
The angel answered, "The Holy Spirit will come upon you..
"The power of the Most High will overshadow you,.
"so the Holy One to be born will be called the Son of God..
"Even Elizabeth, your relative,.
"the mother of John the Baptist,.
"is going to have a child in her old age..
"She who was said to be barren is in her sixth month,.
"for nothing is impossible with God.".
Now, one can imagine the shock.
and the astonishing news, right?.
Gabriel announces God's plan,.
actually it's God's interruption,.
major interruption in Mary's life and in Joseph's life..
It's a major detour, major interruption.
for her original plan..
Remember what I said in the social custom.

$^{161}$of Mary's and Joseph's day?.
They were expected to just get married.
and have their own family, and they just move on..
But one can imagine many whys in Mary's head..
Perhaps Mary would ask, "Why me?".
Because there are so many Marys on the street, right?.
According to some statistics,.
more than 50\% of girls, Jewish girls, were named Mary..
There are tons of Mary out there..
Why me, right?.
And perhaps there are many young couples.
engaged in Israel at that point, right?.
They are not the only one engaged..
Many other are pledged to get married at that point..
And perhaps I think the most perplexing question is 34..
Verse 34, "How will this be?".
The young girl Mary asked the angel Gabriel,.
"Since I am a virgin," right?.
That is difficult to swallow, right?.
A virgin birth..
As for Joseph, I think it will be difficult.
to explain as well..
As you know the story from the Gospel of Matthew,.
Joseph wanted to divorce Mary secretly.
in order to protect her, not to shame her..
But because of divine intervention,.
Joseph obeyed and took Mary as his wife..
In any case, it was a major, major just interruption.
in their life, right?.
When God moves so often, He moves with interruptions..
Is that interesting?.
When God moves in your life,.
sometimes He does so with detour..
Something that is not expected..
Something that you never thought of..
So how did Mary respond to the angel?.
Verse 38, "I am the Lord's servant," Mary answered,.
"May it be to me as you have said.".
Then the angel left her..
Now if you come from a Roman Catholic background,.

$^{201}$you probably heard of the phrase Mary's fiat..
Mary's fiat..
Anyone from the Catholic background?.
Maybe you went to a Catholic school.
or maybe you come from a Catholic,.
Roman Catholic family..
If you have that background,.
you probably have heard the phrase, the fiat of Mary..
It tried to capture Mary's reply to the angel..
If you don't have Catholic background,.
you might have heard the word fiat as well, right?.
Fiat is Italian auto brand, all right?.
So when I prepare this sermon,.
I shop around for fiat as well..
That is fiat..
Of course, when we talk about Mary's fiat,.
we don't mean Mary had a Italian car, right?.
And by the way, before I bought a Mini Cooper,.
I was debating to get a fiat, but I didn't..
I got a Mini Cooper..
Anyways, Mary's fiat..
We don't mean Mary had a car..
The word fiat comes from the Latin translation of the Bible.
some 400 years after Christ..
So it's somewhere around the fourth century after Christ..
Some Bible translator translate the Bible.
into the common language of the day..
And the common language, common,.
is the Latin word voce..
Voce, voce means common language..
Back then, a lot of people were using Latin.
as their common language..
The Bible has to be translated into Latin..
And that translation is very famous,.
very popular at that time, called voce..
And in the voce translation, the words of Mary,.
the first word, may it be, is translated in Latin, fiat..
So you got it, right?.
Fiat is Italian word..
It means let it be, let there be..

$^{241}$You remember Genesis 1?.
When God set up the cosmos, when God set up the universe,.
what was the first word spoken by God?.
What was the first word?.
Let there be, right?.
Let there be light..
And in voce, in the Latin translation, is fiat, looks..
Fiat, let there be, let it be..
Looks, light..
Fiat, looks..
Let there be..
So when Mary says, fiat, let it be,.
it means let your will be done..
Well, I don't understand..
I cannot comprehend everything..
But let it be, fiat, is a word of faith..
Now at this point, I'm sure Mary doesn't even know.
how the birth of Savior can be accomplished.
since she was a virgin..
And what will be the son's destiny?.
What will happen to my son, to this special son?.
What will be her future?.
Nothing is in the answer..
There were tons of uncertainties, unknowns at this point..
For this high school girl, right?.
13, 14 year old Mary, God bless her..
And that is so special about Mary..
Now, of course, if you come from a Roman Catholic background.
you will find doctrines that would amplify,.
that would amplify Mary's status..
If you come from a Roman Catholic background,.
they would describe Mary in highly exalted term,.
like co-redemptrix, co-redemptrix..
So Mary, because of her obedience,.
made salvation possible..
So God, of course, is the first cause,.
but not without Mary's obedience..
So in Roman Catholic theological term,.
it is highly exalted because Mary's obedience.
then make the redemption possible..

$^{281}$So she is called co-redemptrix..
And sometimes called mediatrix, mediatrix,.
the mediator or female mediator between God and humanity..
Or because Mary's high status, exalted status,.
she has, she's being described as perpetually virgin,.
perpetual virginity, or immaculate conception..
It means when she was born, when Mary was born,.
she was born without original sin..
She was not born like us in the sinful condition..
Now this sort of theological imagination.
found no biblical data, no biblical support..
As you can see, Luke 1:46, in one of the Psalms.
that Mary sang to the Lord after receiving the annunciation,.
Mary said, "My soul glorifies the Lord,.
"and my spirit rejoices in God, my Savior.".
Does Mary need a Savior?.
Yes, she does, right?.
1:46, Gospel of Luke, and 1 Timothy 2.5..
It says point blank, there's one God,.
and how many mediator between God and man?.
One, right?.
One mediator between God and man, the man, Jesus Christ..
So it is clear there's no biblical warrant.
to describe Mary in such a high status..
But we still need to wonder,.
we still need to wonder how a teenage girl,.
how a teenage Jewish girl could have such a great faith..
Right?.
You go to a high school fellowship,.
you pick one of those girls, you ask them,.
"Do you believe in Jesus Christ?.
"Do you believe in God?".
They probably will say, "I come to church because it's fun.".
Mary is being set apart here..
But I wanna make sure that you understand.
this is not blind faith..
'Cause so often we morning people would see faith.
as some sort of blind, irrational belief,.
or wishful thinking..
You know, in the Bible, faith is not like that..

$^{321}$In the Bible, faith is based on something solid,.
and that is the character of God, the faithfulness of God..
And God has revealed Himself in the history.
and in Scripture that He is faithful..
So that kind of faith is not blind faith,.
it's reasonable faith..
It's great faith indeed, but it's reasonable..
So no matter what, God will work things out.
for His own glory..
Even though we don't know or we cannot comprehend..
And as for Mary, it's the same thing..
I just don't know how it would be accomplished..
I just don't know how it will unfold, right?.
We are at the end of the story,.
we saw the end of story because Jesus died on the cross.
and on the third day He rose from the dead, right?.
But for Mary at this point,.
she had to wait for another 30 some years.
in order to see the full plan of God revealed, right?.
But at this point, she is the first hero of faith.
in the New Testament..
Let it be, fear it..
And after this, Mary rushes to see her relative, Elizabeth..
And we see here this exchange between two mothers..
We see why Mary was so special..
Luke 1:39, "At that time, Mary got ready.
and hurried to the town in the hill country of Judea,.
where she entered Zechariah's home and greeted Elizabeth..
When Elizabeth heard Mary's greeting,.
the baby lived in her womb.
and Elizabeth was filled with the Holy Spirit.".
1:42, "In a loud voice, she exclaimed,.
'Blessed are you among women.
and blessed is the child you will bear..
But why am I so favored.
that the mother of my Lord should come to me?.
As soon as the sound of your greeting reached my ears,.
the baby in my womb lived for joy..
Blessed is she who has believed.
that what the Lord has said to her will be accomplished.'".

$^{361}$You notice where's John the Baptist here?.
It's in the womb of Elizabeth, right?.
Still in the womb, six months..
And this baby in the womb is already hearing.
the coming of the Savior by greeting the mother of Jesus..
It's jumping inside the womb..
But how does the Bible,.
how does the scripture describe Mary here?.
1:45, "Blessed is she who believe.".
You see here, it's not Mary is essentially different from us..
It's not because she is holier..
It's not because she is set apart from us, essentially..
She is the same human being just as we are..
It is not Mary is essentially different from us..
It is she who believe..
It is her faith, her faith in God,.
in God's plan, in God's bigger plan..
Even her own plan is being disturbed,.
her own plan being interrupted or detoured..
But she puts her faith in God,.
knowing that the plan of God will not fail..
Many years ago, there's a British Bible translator.
or British Bible scholar named J.B. Phillips..
J.B. Phillips translated the New Testament.
into contemporary English at that time..
And he wrote some other books..
And one of the little books that we should all read,.
we should all read..
It's a very little, tiny book..
And sometimes you pick up a book,.
you have to read cover to cover.
in order to find out the content, right?.
But for some books, if you pick it up,.
you read the title of the book,.
you immediately know what the book is all about..
You don't have to read..
You ever encounter any books like that?.
You pick it up, read the title or subtitle,.
you know what the book is all about..
You don't have to read it..

$^{401}$You don't have to buy it..
Some books are like that..
And J.B. Phillips wrote such a book..
And the book title is "Your God is Too Small.".
Your God is too small..
It's a tiny little book..
It would be beneficial for you to pick it up and read it,.
you know, to edify your faith, to build your faith..
But as I said, you pick up this book,.
very thin, it's like a booklet..
And you immediately know what the message is all about..
Your God is too small..
Do we all have the same struggle?.
We do..
We struggle with the fact that our God is too small..
It's way too small.
because we create our God in our own image, right?.
Instead of the other way around,.
when things don't seem the way we expected,.
when we face uncertainties,.
when things seem to be in chaos,.
we reduce God into a little idol..
God is too small..
And sometimes our prayer can be too small..
What have you been praying for?.
What have you been praying for this month?.
Probably some needs or some list of things that you want,.
right?.
And the message seems to challenge all of us.
that sometimes our prayer can be too small.
because our God is too small..
And sometimes our own plan may also be too small..
And we seldom think of the plan of God,.
which is much, much, much bigger..
It takes miracles to see the plan of God..
It takes God himself to accomplish.
something that is worthy of his name..
Sometimes it takes God's power and our obedience,.
our fear in order to see the big plan of God..
Now, the faith of Mary is fear..

$^{441}$Let there be..
Yes, I don't understand..
I don't see all the details..
I don't see the details as yet, but let it be, fear it..
I don't need all the details worked out before I say so..
May this be our prayer, our attitude and our faith..
Let's pray together..
Father, thank you for being our God..
Thank you for being our big, big God,.
much bigger than we think,.
much bigger than our mind and our plan and our thinking..
You are indeed worthy of our worship..
So we come before you in humility,.
knowing that you are the one who makes all things new..
You are the one who control and orchestrate everything.
from the greatest to the least..
All things are in your hand..
So we commit our life into your hand..
And I pray that you'll help us to say fear it,.
to say let it be..
In our life journey, we wanna see your glory..
In our faith, we wanna see your power..
So grant us this obedient, obedient attitude,.
the faith that we're gonna see that you are faithful..
And the plan that you accomplish is much bigger,.
bigger than our own plan..
Father, once again, use us, use our faith,.
increase our faith and help us to be loyal and faithful.
to follow you, to serve you, to continue to worship you.
so that we will see your glory more and more.
as we experience your power..
In this Advent season, come to us.
and we open our hearts to welcome you,.
our Lord Jesus, who come as a baby in the manger.
to bless us, to redeem us and to save us.
so that we can be your sons and daughters..
Thank you and we worship you in Christ and pray..
Well, God bless you and see you next Sunday..
And remember next Sunday is gonna be 5 p.m., right?.
5 p.m., yeah, don't come in the early morning, okay?.

$^{481}$There's no breakfast for you..
Get your ticket and join our Christmas baptism Sunday..
All right, God bless..
\newpage



\section{路加福音}
\label{sec:vk3Vfk1fF_M}
\textbf{將臨節的挑戰 - 神子降臨 (路加福音1-2) \#3 - 陳崇基牧師}
\newline
\newline
連結: \href{https://youtube.com/watch?v=vk3Vfk1fF-M}{\texttt{https://youtube.com/watch?v=vk3Vfk1fF-M}} ~~~~ 語音日期: 2023-12-19
\newline
\newline
\hyperref[sec:JrlUbBSqsPA]{\small{< < < PREV SERMON < < <}}
~
\hyperref[sec:index_chronic]{\small{[返順時目]}}
~
\hyperref[sec:index_scriptual]{\small{[返順卷目]}}
~
\hyperref[sec:SCf1z8t8nOA]{\small{> > > NEXT SERMON > > >}}
\newline
\newline
$^{1}$各位好.
我們繼續去預備聖誕.
我都沒有聖誕帽.
總之我帶頂聖誕帽上來.
Anyways.
聖誕快到了.
我們繼續去看一些經文.
今天可能會看多一點.
因為下個星期已經是聖誕的主日.
所以我們要去完結盧家福音.
一至到第二章.
我們看看張臨哲.
所謂的Advent.
給我們一些什麼提醒.
上一次我們去看盧家福音第一二章的時候.
我們說到一二章.
就是一對一對的經文.
放出來的.
所以一開始的時候.
你會發現有施洗約翰出生的報喜.
Annunciation.
Annunciation就是說.
The news of giving birth.
就是說他快要出生了.
所以有施洗約翰出生的報喜.
一章五至二十五節.
然後第二部分.
一對一對的.
就是耶穌出生的報喜.
二十六至三十八節.
然後兩個母親.
就是施洗約翰的母親以利沙白.
還有耶穌的母親瑪利亞.
他們相遇的一段的記述.
然後就是施洗約翰出生.
然後耶穌的出生.
然後英雄的耶穌就被帶到聖殿.
這裡說了兩個故事主在聖殿裡面出現.
所以今天我們會看比較多的經文.
我們已經看了施洗約翰的報喜和出生了.

$^{41}$這次去到《路加福音》第二章.
在那些日子.
凱撒·奧古斯都降旨.
叫全國人民都登記戶籍.
這第一次登記戶籍.
是在居里樓作敘利亞總督的時候行的.
眾人各歸各城辦理登記.
約瑟也從加利利的拿撒勒城上猶太去.
到了大衛的城民叫伯里行.
因為他是大衛家族的人.
要和他所聘之妻瑪利亞一同登記戶籍.
那時瑪利亞已經懷孕.
他們在哪裡的時候.
瑪利亞的產季到了.
就從胎的孩子.
用布包起來放做馬槽裡.
因為客房裡面沒有地方.
這是一個很熟悉的故事.
聖誕的故事.
熟悉到一個地步.
我們覺得只不過是給小孩子看.
裡面有一些細節.
我們不妨再點一點出來.
在政治和社會或者歷史上.
誰在掌權呢.
凱撒奧古斯都.
羅馬帝國第一個帝王.
第一個的皇帝.
全國人民都要登記戶籍.
或者人口的普查.
我比較喜歡以前舊的和本的翻譯.
就是全天下.
全天下都要去登記戶籍.
所以掌權的人就是奧古斯都.
全國的人或者天下的人.
都要聽他的命令.
因為他降旨下來.
但這裡有一個地方很奇怪.
說到聖誕節.
耶穌出世的事蹟.

$^{81}$或者裡面一些細節的時候.
我不知道你有沒有那種驚訝.
我們看回聖誕節的故事.
我們需要重新回顧那種驚訝的感覺.
什麼驚訝呢.
約瑟和瑪利亞在哪裡呢.
加利利的拿撒勒在北邊.
就是以色列的北邊.
加利利的拿撒勒他們就上猶太去了.
為什麼上猶太去了.
因為猶太在中間山地.
耶路撒冷附近.
所以他們一定會上山.
加利利比較低一點.
所以他們.
Elevation來說.
他們一直爬上去.
上到去猶太地.
就是耶路撒冷,伯利恆等等的地方.
拿撒勒在北邊.
伯利恆就在南邊.
如果步行的話.
距離大概有90英里.
90英里就是150公里.
大概應該要走三四天的路程.
中間要停一些.
可能是Motel或者Airbnb等等.
休息一下.
然後再上路.
每一天早上又再上路.
有沒有覺得奇怪.
為什麼上帝不乾脆在伯利恆那裡找一對新人.
我相信伯利恆.
上一次說過.
伯利恆裡面應該大部分的女孩.
穿大半都是叫瑪麗亞.
是不是.
如果你想找一個瑪麗亞的.
隨處都有.
是不是.

$^{121}$或者不需要找一個叫做瑪麗亞的.
因為你看到耶穌降生的故事.
瑪麗亞這個字.
或者一個名字.
其實沒有什麼特別.
沒有什麼特別的意義在這裡.
可以叫瑪麗亞.
可以叫Susanna.
可以叫其他的女孩的名字.
是不是.
偏偏他就在北面.
拿撒拉.
來找一個的同女.
已經是許配了給約瑟的一個同女.
約瑟和瑪麗亞他們遠道的.
要下去伯利恆.
是不是.
很難搞.
是不是.
為什麼要搞這麼多東西出來.
是不是.
為什麼不方便一點.
我們華人就很方便.
我們很practical.
是不是.
為什麼要這樣做.
繼續看下去.
大衛的城.
名叫伯利恆.
伯利恆是一個小鎮.
大家知道了.
但是在舊約裡面是很有名的.
為什麼呢.
因為伯利恆是誰的祖家.
是大衛自己的家鄉.
這裡說到大衛城.
伯利恆.
City of David.
Bethlehem.
如果熟悉舊約的話.

$^{161}$應該會想到.
在舊約裡面.
有關伯利恆的預言.
或者有關伯利恆的.
應許promise.
其中一段.
是很著名的.
就是尼加書第五章.
那裡說什麼呢.
那裡說伯利恆.
以法他.
伯利恆又叫以法他.
有出名的地方.
通常又多過一個名字.
你在大衛諸城裡面雖小.
將來必有一位.
從你那裡出來.
在以色列.
為我作掌權者.
他的根源是耿古.
從太初就有了.
這個尼加先知.
在舊約裡面其中一個小先知.
在第五章這裡說什麼呢.
其實他在指罵北國的以色列.
或者南國的猶太國.
就說那些國民已經離開了耶和華.
這個國度快要完蛋了.
於是他審判完之後.
他就說將來有一個盼望.
是上帝自己可以成就的.
這個盼望就是在伯利恆這個地方.
會再有一位掌權者.
從你那裡出來.
什麼意思.
如果你想象一下當時.
舊約以色列國.
猶太國.
已經分了南北國.
當時的王子.

$^{201}$或者掌權人.
是從哪裡出來的.
如果是猶太國的話.
首都在哪裡.
猶太國首都在哪裡.
耶路撒冷沒錯.
在耶路撒冷.
OK.
剛好又叫大衛城.
耶路撒冷.
Jerusalem.
所有王室或者王族的後裔.
將來會掌權的.
都會在耶路撒冷出來.
尼加這裡說什麼呢.
尼加就說.
這個新的掌權者.
就會在.
伯利恆出來.
換句話說.
尼加說.
會再在大衛的原生的祖家.
換句話說.
這個伯利恆再出來的掌權者.
就不是耶路撒冷那一批.
耶路撒冷那一批已經.
糾且殘存了.
將來一定是Game Over玩完.
他說伯利恆Bethlehem這個地方.
會有一個新的掌權者.
新的王出來.
換句話說.
就是有一個New Beginning.
不再靠耶路撒冷那一幫.
敗壞了的朝廷的人.
那些大衛家已經玩完了.
重新在.
伯利恆大衛的老祖家.
再會有一個掌權者.
並且他的根源.

$^{241}$他的根.
是Eternal的.
從太初就有.
這裡我們搜集了.
一點點不同的細節.
剛才我們看路加福音第二章開始.
還有現在.
尼加書第五章.
奧古士都降旨.
要人口統查.
每個人要登記戶籍.
其實在京外的一些記載.
我們找到一些.
Census.
是羅馬帝國會做的Census.
譬如高盧.
高盧就是今天的哪裡.
法國.
在那裡有一些人口統查.
在埃及有人口的統查.
路加福音第二章.
這個人口統查.
絕對不是第一次來.
絕對不是整個帝國第一次.
周不時都會做的.
奧古斯都降旨.
要做人口統查.
要抽稅.
這裡告訴我們.
地上的君王可以做很多不同的政策.
降旨.
但是.
上帝的旨意.
是凌駕地上的帝王.
伯利恆的預言.
尼加書五章這裡.
當時尼加做先知的時候.
是主前七百多年的時候.
七百多年的時候.
跨越了七個世紀的歷史.

$^{281}$中間很多改朝換代.
不同的帝國等等.
應許有沒有落空呢.
最後就是聖誕這個時候落空.
跨越所有歷史裡面的.
上上落落.
喜悅服和改朝換代.
仍然是沒有落空.
那撒勒和伯利恆.
九十英里之遠.
是超乎常理的.
並且按當時的習慣來說.
約瑟是不需要帶瑪利亞去的.
因為只需要南丁來做人口統查.
他沒有需要帶著瑪利亞.
那你想一下有什麼原因.
為什麼約瑟要帶著瑪利亞呢.
很明顯應該是新人已經快到了.
應該是約瑟不想留瑪利亞在北面那撒勒.
以免他來回的過程裡面.
他夫君不再.
英孩就出世了.
英孩出世的時候.
媽媽生小孩的時候.
夫君不在很慘.
對夫君來說很慘.
要陪著他.
所以我想多煩都好.
他都要帶著瑪利亞.
可能租一隻路或租一輛Uber.
帶瑪利亞長途不息的.
由那撒勒到伯利恆.
你記得我們曾經說過一個聖經的教義.
叫做神的庇護.
Providence.
這裡其實你可以這樣想.
上帝要去到show off.
上帝要show off.
去到展現出他自己至上的權能.
剛才我們所搜查的細節.

$^{321}$是凌駕地上的帝權.
是超越時間地理的限制.
上帝要做的事.
他怎樣都會做到.
他是show off上帝自己.
show off他最終極的指揮權.
是全部東西.
大小事情少到一個地步.
連時間都是在他自己的掌握.
所以路加福音第二章你可以想像.
每個人都回鄉去登記.
約瑟和瑪利亞到了.
回到祖家.
可能回鄉登記的人太多.
黑房裡面沒有地方.
所以他們就被安排去到馬鵬.
馬鵬就成為產房.
怎樣都好.
黑房裡面沒有地方.
馬鵬誕生.
總是一種謙卑的感覺.
第二章第一節說到奧古斯都.
在上掌權的萬人之上.
這個君王.
神的兒子.
就是謙卑的降生.
馬鵬的救主.
在路加福音第一二章裡面.
我們說過有報喜.
也有出生的故事.
裡面也有很多不同的詩歌.
如果你數一下的話.
應該大概有四首的詩歌.
四首的詩歌.
為什麼會這樣寫.
當時的biographers.
那些傳記.
名人的傳記.
很多時候說到他們出生.
說到他們成長的時候.

$^{361}$都會有一些故事和詩歌.
來讚頌他們.
譬如羅馬王.
奧古斯都.
當然會說一下他的威風史.
還有一些出奇的故事等等.
如果你在當時的話.
有電視.
如果當時的話.
有YouTube的話.
你每一天都會聽到一些宣傳.
就是這個國王是多麼偉大.
多麼威風.
還有會聽到很多的詩歌.
那些詩歌.
有如今天的一些國歌.
National Anthem.
去說羅馬帝國.
要向他忠誠.
羅馬這個帝王.
是多麼的好.
所以路加福音第一,二章.
說到耶穌出生的時候.
說到他報喜的時候.
也有幾首詩歌.
這幾首詩歌我們都看過.
第一首是瑪利亞的Manefica.
說到耶穌的出生.
預言耶穌的出生.
一章68至79節.
和耶穌沒有關係的.
就是約翰的出生.
就是撒迦利亞的父親.
撒迦利亞的父親.
的一首讚美的詩歌.
說到這個孩子.
成為彌賽亞的先鋒.
二章14節.
就是天使的Doxology.
巴黎人的野地.

$^{401}$的牧羊人.
然後第四首.
就是在聖殿裡面.
有一位可能是老人家.
他的頌歌.
說到將來.
這位的孩子.
耶穌基督.
有的榮耀.
我們繼續看第三首.
二章第八節.
在白瀚的野地裡面.
有牧羊人.
夜間值班看守羊群.
在主的一位.
有主的一位使者.
站在他們旁邊.
主的榮光.
四面照著他們.
牧羊人就懼怕了.
那天使對他們說.
不要懼怕.
看啊.
因為我報給你們.
大喜的福音或者訊息.
是關乎萬民的.
因今天在大衛的城裡面.
為你們生了救主.
就是主基督.
你們要看見一個.
嬰孩包著布.
臥在馬槽裡面.
那就是給你們的記號.
忽然有一大隊天兵和.
和那天使讚美上帝說.
在至高之處.
榮耀歸於上帝.
在地上平安歸於.
他所喜悅的人.
這裡有沒有留意一些用詞.

$^{441}$我點給你看.
主Lord.
救主Savior.
天下.
平安或者盛世.
大喜的訊息.
福音Gospel.
這些全部都是.
羅馬帝國的電視官媒.
會說的字眼.
你可以見到.
今天我們說的報街音.
我們今天說教會.
有時候會去報街音.
或者去老人院.
或者去街上報街音等等.
是很不同的.
這是第一次報街音.
是帶著很多.
不同的味道在裡面.
所以為什麼今天有很多極權的國家.
伊斯蘭的回教國家.
或者共產的國家.
是不喜歡人慶祝聖誕.
為什麼.
因為他們明白.
聖誕節是什麼意思.
聖誕節就是嬰孩來到的時候.
他才是主.
其他人變成僕人.
那些在上掌權的.
他以為他是在萬人之上.
無人之下.
這個嬰孩誕生.
真正的主.
Lord and Savior.
真正帶來Peace的.
平安的.
是因為這位在馬槽裡面的救主.
你這樣聽聖誕的故事.

$^{481}$將會完全不同.
牧羊的人在.
白人的野地.
他們看到主的榮耀顯現.
這些也是舊約裡面.
經常有的故事.
當上帝顯現.
我們有個名詞叫Theophany.
就是Manifestation of God.
就是上帝彰顯的時候.
有很多很多這樣的記載.
在西乃山.
上帝去到顯現的時候.
和摩西顯現.
摩西在牧羊的時候.
在曠野的時候.
他就用荊棘裡面的火.
向摩西顯現.
當以色列人過了紅海.
去到西乃山的時候.
在西乃山的山腳聚集的時候.
在山頂上帝的榮耀就顯現了.
Glory.
然後有地震.
有風有雷.
會幕建成之後.
他們在曠野行走.
這個會幕帶領著他們.
這個會幕建成之後.
上帝的榮耀又充滿這個地方.
你知不知道在舊約裡面.
希伯來文.
榮耀這個字.
其實就是重的意思.
Heavy.
大家有沒有舉過重啊?.
有吧?.
有沒有舉過啞鈴啊?.
50磅可以嗎?.
50磅不是很重嗎?.

$^{521}$如果你舉100幾200磅的話.
你就知道有多重.
Glory就是這個意思.
就是上帝的榮耀臨到的時候.
那種威勢.
所羅門聖殿.
被建立的時候.
同樣有上帝的榮耀降臨.
這裡說.
牧羊的人見到.
主的榮光四面照著他們.
這個奇怪的地方就是.
居然不在耶路撒冷.
不在聖殿.
不在當時的京都.
在帕里蘭的野地.
一個他們不預期會有上帝榮耀的地方.
上帝就出現.
四面照著他們.
上帝的榮光.
然後他們就唱出這首歌.
說這一位的救主.
就是至高之處.
歸於上帝.
二章十五節.
眾天使離開他們升天去了.
牧羊的人彼此說我們往帕里蘭去.
看看所成的事.
就是主所告訴我們的.
他們急忙去了找到瑪利亞和約瑟.
還有那英鞋臥在馬槽裡.
他們看見就把天使論這孩子的話傳開了.
聽見的人就詫異牧羊人對他們所說的話.
瑪利亞卻將這一切的事傳在心裡.
反覆思想.
牧羊人回去.
因所聽見所看見的一切.
正如天使向他們所說的.
就歸榮耀於上帝讚美他.
到了白天.

$^{561}$他們就給孩子行國禮.
又給他起名叫耶穌.
這就是還沒有在母胎裡面.
成胎之前天使所起的名.
然後我們就去到.
繼續這個故事的記載.
耶穌開始長大.
英鞋的時候.
約瑟和瑪利亞就帶他去到聖殿.
在盧家福音.
上帝救恩的中心.
或者開始的地方.
就是聖殿.
所以你看到盧家福音從開始到最後.
都是在聖殿.
開始的時候就是聖殿.
第一第二章.
最後的時候二十幾章.
也是回到聖殿.
聖殿就是上帝.
去啟示和啟動.
他自己傳遞救恩計劃的地方.
所以耶穌進入聖殿.
這個絕對不是只是習慣的問題.
不是習俗的問題.
因為將主的榮耀.
這位的英鞋.
帶到聖殿.
就是說這個榮耀終於進入聖殿.
他們帶這個英鞋.
進入聖殿的時候.
就按照摩西律法.
滿了潔淨的日子.
他們就帶著鞋子上耶路撒冷去.
要把他獻給主.
正如主的律法上所記的.
凡投生的孩子必歸主為誠.
又要照著摩西律法所說的.
用一對班戈.
或者兩隻的鞦韆來獻祭.

$^{601}$上次我提過一些背景.
在舊約裡面.
那些比較窮的人.
他們不能夠獻一隻羊羔的話.
因為潔淨的禮應該要獻一隻羊羔.
一隻羊羔大概多少錢?.
你有沒有去Costco買?.
買一隻羊.
好像是八多塊.
三百多塊已經這麼貴了.
是不是?.
三百多塊.
我記得有一次在West Bank.
我問一個人.
我說一隻羊多少錢?.
他就告訴我.
大概一百塊美金.
OK.
因為這裡三百多塊.
因為窮人是不能夠獻一隻羊羔的話.
他就可以獻兩隻的雛鴿.
所以我們知道.
這個社會背景的時候.
約瑟琳瑪利亞應該是窮人.
當然這個是之前.
以前第一世紀的世界.
今時今日不是了.
上星期有個弟兄提過.
今天兩隻的乳鴿.
其實都很貴.
也不是每個人都能吃到.
你這裡能吃到嗎?.
吃不到.
我想說的是.
譬如你今天去西場.
去Western Wall.
Western Wall其實是一個Open Air的聖堂.
很多時候有些人去聚集.
有些什麼慶典等等.
巴默斯福.

$^{641}$或者一些生日會等等.
他們就會大肆慶祝.
有錢的人家.
有錢的猶太人的人家.
他們就會請一些樂隊.
或者有跳舞.
作樂的人.
一起在西場.
他們就會在這裡聚集.
然後大肆慶祝.
似乎瑪利亞和約瑟.
就沒有這些Luxury.
只是欠兩隻乳鴿.
Again.
當時來說是便宜的.
OK.
二十五節.
那時在聖殿耶路撒冷.
有一個人名叫西冕.
這人又公義又虔誠.
素上盼望以色列的安慰者來到.
又有聖靈在他身上.
他得了聖靈的啟示.
知道自己未死之前.
必看見主所立的基督.
受了聖靈的感動進了聖殿.
正遇見耶穌的父母抱著孩子進來.
就照律法的規矩而行.
西冕就把他抱過來.
稱頌上帝說.
這首就是第四首歌.
主啊 如今可以照你的話.
用你的僕人安然去世.
因為我的眼睛已經看見你的救恩.
就是你在萬民前所預備的.
是啟示外邦人的光.
是你聞以色列的榮耀.
西冕之後又有另一個叫Anna.
我相信是代表著整個以色列國.
也代表著以色列上帝的子民.

$^{681}$去到盼望上帝的應許.
終於可以落實.
可以在我們中間成就.
他說他整個人生.
他是盼望著以色列的安慰者.
舊約先知特別爾賽亞書.
經常說到將來.
被擄的流亡的以色列民.
他們會回家的這個promise.
上帝的榮耀將會回歸.
這個就是以色列的安慰者.
他們所盼望.
見到這個嬰孩進到聖殿.
有沒有那個wonder.
剛才我說了很多.
上帝的glory.
進入聖殿的時候應該是怎樣的.
是千軍萬馬.
是很heavy的.
但居然在這裡.
上帝的榮耀不是這樣彰顯.
是一個手無寸鐵的孩子.
被抱到進去上帝的聖殿.
上帝的榮耀駕到.
他有這樣的盼望的信心才能夠看到.
33節.
孩子的父母即約瑟瑪利亞.
因論耶穌這些話就驚訝.
they wonder.
西冕給他們祝福.
又對孩子的母親瑪利亞說.
這孩子避難.
要叫以色列中間很多人跌倒.
很多人興起.
又要成為毀謗的對象.
叫很多人心裡面的意念顯露出來.
你自己的心也要被劍刺頭.
我想對一個母親來說.
這句話就是最拗頭.
最難解的看不透.

$^{721}$因為對瑪利亞來說.
最大的關懷是什麼.
做過母親的就知道.
你最大的關懷是什麼.
你對你的兒女的將來最大關懷.
你很希望他知道他將來是怎樣.
知道將來去哪間學校.
知道將來他的家庭是怎樣.
他的配偶是怎樣.
你最大的ultimate concern.
就應該是你的兒子.
你的兒女的將來.
因為你的兒子的將來.
就是你的glory.
他的榮耀就是你的榮耀.
如果孩子傷心的時候.
你就心如刀割.
這就應該父母.
特別是母親.
對孩子的心腸.
這裡說到耶穌的將來.
說了一點又不說一點.
一方面他使很多人得救.
另一方面.
這位耶穌也使.
有些不信者反對他.
但最後那句才是最難解的.
因為想像一下當時.
孩子只是幾十日大.
你的心將會被劍刺頭.
你會怎樣去明白這句話.
我們當然知道.
因為劇透了.
因為耶穌最後.
是被釘上十字架.
承受他的痛苦.
我們知道在十字架面前.
也有不同的婦女.
瑪利亞也在.
站在十字架面前.

$^{761}$看著她自己的兒子被羞辱.
赤裸裸地被羞辱.
看著她的兒子慢慢死去.
心如刀割.
我們看了這裡不同的故事.
有很多人看到第一個聖誕.
約瑟,瑪利亞.
伯利恆野地的牧羊人.
聖殿裡面的西面.
和之後的阿那.
你可不可以找到一個共通點.
在這麼多人的中間.
有什麼共通點.
每一個他們的共通點.
只是看到救恩上帝計劃的開始.
第一,二章當然.
老家坊頭兩章.
他們看到的只不過是.
上帝的計劃開頭的地方.
他們看不到的.
就是最終上帝會怎麼帶領.
最終上帝的計劃應該怎樣揭曉出來.
Don't know yet.
要等很多年才知道.
我想信心信仰其實就是這樣.
我們作為基督徒.
我們所看到的.
是很有限的東西.
是開頭那一,兩章.
未來或終極是怎樣.
都是模糊不清的.
你想想你現在的生命裡面.
可能面對著不同的十字路口.
可能面對著不同的困難.
你所看到的只是.
這裡現在的一,兩頁.
怎樣揭曉下去.
還是模糊不清的.
有很多東西都解不透.
解不透,猜不透.

$^{801}$但這班見證到第一個聖誕的人.
就有一個這樣的共通點.
雖然他們看到的非常有限.
對將來是模糊不清的.
但他們有另外一個的共通點.
就是他們願意相信.
他們只能夠確信上帝的計劃是不會落空的.
他們的希望,他們的盼望.
是在上帝,信實的上帝身上.
在生命裡面發生很多的事情.
是我們不能夠去估算.
也猜不透.
好像奧秘一樣.
We ask why.
但我們未有answer.
何時有answer.
將來.
但將來是模糊不清的.
最近在網上有一位牧師.
他分享了一些自己的寫作.
你叫他靈修札記.
是他自己對詩篇的一些感受.
他每天看一篇詩篇.
然後寫了自己的心得等等.
他就將它編成書.
然後放上網是免費的.
原因是什麼呢.
這位牧師.
他受了多年的痛苦.
很苦.
現在還患了血癌.
他走過很多痛苦的療程.
每一天他看詩篇.
寫下他的禱告.
如果他精神狀況好的話.
他會寫一些文章.
我就下載了這本書.
我看了好幾篇.
我覺得很感動.
因為這些是有血有淚的經歷.

$^{841}$不只是抽空學術的東西.
我不是說學術不重要.
學術重要.
但這是他生命裡有的經歷.
每一天看一篇詩篇.
能夠爬到起來看詩篇.
能夠寫到一些東西.
我看了一篇.
是講到詩篇第二篇.
當然我們沒有時間看詩篇第二篇.
詩篇第二篇是講什麼呢.
是講到上帝將來會掌權.
現在可能見到很混亂.
外邦的君主來反對.
甚至去攻擊.
但上帝所揀選的受高棍.
將來一定會掌權.
一篇很有寓言味道的詩篇第二篇.
他就這樣寫.
療癒的詩篇的書名.
他說看完這篇詩篇的時候.
教會不可以談政治真是可笑.
但我們永遠都不應將焦點.
依靠放在政權和政客身上.
因為我不知道.
上帝才是我們唯一的主.
除他以外便無拯救.
也沒有讓我們絕對馴服的權柄.
當人還驕傲囂張的時候.
包括我自己.
上帝要向人彰顯那高過我們.
中人腦袋裡面的意念和計劃.
上帝會用你想不到.
甚至你不想也不願意的方法.
可能對我而言就是血病.
去成就它的美意.
同時也可能執行懲罰.
讓原本以為不可能的事成為可能.
誰是真正的掌權者.
一目了然.

$^{881}$不是你.
不是我.
更不是地上任何可憐的政權.
而是萬君之野和他的受告者.
這幾天我再一次確認.
我相信信仰真是這樣.
越來越像長篇劇.
長篇劇大家可能有追劇.
有這樣的經歷.
看首兩集的時候.
有很多事未猜透或未看到.
一定要最後內心看到劇中.
你才知道.
在這個經歷和歷練裡.
只能靠信心.
忍耐.
盼望.
但這個信心和盼望不是盲目的信心.
為甚麼.
因為上帝自己在歷史和話語裡.
已經告訴我們他是一位可信者.
他是信得過的.
他就是告訴你.
你儘管信我.
我會逐步帶領你.
讓我們在聖誕重拾這個信心.
可能跟著主的時候覺得很迷惘.
不知道如何反應.
你要知道上帝仍然是可信的.
求主幫助我們.
幫助我們.
我們同心合力去祈禱.
主我們在人生不同的階段.
不同的歷練當中.
上主我來求你.
我們來仰望你.
在聖誕裡面.
在你自己話語當中.
可以提醒我們.
你是那位可信靠的神.

$^{921}$你是那位信實,信服的上帝.
主我們再一次將我們的人生呈現.
將我們自己人生走過.
不同的起起伏伏.
很多的高山低谷.
特別去到今年的年尾.
主我們回看的時候.
可能有很多迂迴曲折的情節.
是我們解不通,猜不透.
但我們深信上帝.
你仍然是信神.
你會藉著不同的方法.
甚至我們想不到的方法.
達成你自己的計劃.
主讓我們信心繼續的去信靠你.
願意來到我們這個世界.
謙卑的降生在馬槽裡面.
給我們這個憑證.
讓我們知道我們能夠繼續去信靠,跟隨你.
祈禱奉主名.
完善祝福各位.
下星期能夠跟大家一齊去敬拜.
\newpage



\section{}
\label{sec:SCf1z8t8nOA}
\textbf{Atomic Wisdom 1 - Rev. Mark Chan}
\newline
\newline
連結: \href{https://youtube.com/watch?v=SCf1z8t8nOA}{\texttt{https://youtube.com/watch?v=SCf1z8t8nOA}} ~~~~ 語音日期: 2024-01-02
\newline
\newline
\hyperref[sec:vk3Vfk1fF_M]{\small{< < < PREV SERMON < < <}}
~
\hyperref[sec:index_chronic]{\small{[返順時目]}}
~
\hyperref[sec:index_scriptual]{\small{[返順卷目]}}
~
\hyperref[sec:8jU_KxWNams]{\small{> > > NEXT SERMON > > >}}
\newline
\newline
$^{1}$- Good morning, everyone..
Well, I've been sick for a few days,.
so I've been suffering too..
But that's okay, yeah..
It's the last day of 2023..
Today's the last day, as Jinfan said..
This is, I think this is the best way.
we can conclude this year,.
end with the Lord's Day on Sunday..
As we go from one year to the next,.
we will think about New Year resolutions, right?.
- Yes..
- We all do, right?.
- Yes..
- So what do you have in mind?.
Maybe physical exercise, learn something new,.
or maybe financial planning and all that..
Maybe for those who go to school,.
probably you have some new plannings as well..
Maybe for questions,.
sometimes we talk about reading the Bible together..
Have we done the Bible reading last year?.
- Yes..
- We finished it, right?.
No, I didn't finish it, I didn't join it..
- Some of them..
- Some of them, right?.
Well, give them a hand, yeah, give them a hand..
Yeah, very good, very good..
Usually churches do is to have some big sign.
or big banner saying,.
"New Year, we're gonna do this or do that," right?.
"We're gonna expand our church,.
"and we're gonna grow together,".
and something like that..
Today, I want to do something a little bit different..
I want to tell you how life transformation work..
Because as we set the New Year goals,.
as we set the New Year resolutions,.
or think about stuff like that,.

$^{41}$not only we need to set goals,.
but we need to understand how we reach the goal.
and how life change works..
So my sermon title is Atomic Wisdom..
Now, atomic does not mean destructive..
Okay, I know there's atomic bomb,.
but it means tiny and insignificant, atomic particle,.
something really, really small..
And I borrowed this title from a bestseller.
called "Atomic Habits" by James Clear..
And this bestseller is recommended to me this year.
by actually Amanda, recommend me to read this book..
And I find it's very insightful book..
And the thesis of the book is success,.
all kinds of success, okay?.
Whether you're in school, whether you're in career,.
whether you're just doing exercise,.
just trying to be strong, all kinds of success.
is a product of daily, tiny habits,.
not a once in a lifetime action, you know?.
All kinds of changes, all kinds of positive changes.
have to be done in tiny little steps..
And everyone can do that, and that's the good news..
Now, this book is packed with insights.
from behavioral psychology..
You know, behavioral psychologists observe the human nature,.
observe how we work, right, how we behave..
And then they summarize the psychological dynamic..
And as I read this book,.
I instantly, I thought of wisdom literature in the Bible..
You know wisdom literature?.
How do you find the wisdom literature in the Bible?.
It's right in the middle, right?.
Some of the Psalms and Proverbs, Ecclesiastes,.
and the Book of Job, and those are wisdom literatures..
It's always, wisdom literature has a feature,.
it is always short and clever sayings, tiny sayings..
And it seems to instill us with tiny little wisdom.
here and there, and tell you some general truth.
about human nature, general truth about.

$^{81}$living a life of wisdom, living a life that is successful..
Actually, in the Bible,.
wisdom is defined as practical skill..
If you have wisdom, you live skillfully..
You live skillfully to the life God has given you..
So at times in wisdom literature in the Bible,.
observing human nature, and try to summarize it.
and teach us some of the general truths..
So this sermon, I'm going to do some integration..
I'm gonna go for some of the insights.
from behavioral psychology, and then we'll go into.
some of the passages from the Book of Proverbs,.
wisdom literature..
So here, I want you to understand how life change works..
All right, so whatever you have.
for your New Year resolutions, all right?.
If you don't have any New Year resolutions,.
this is time to think about it, okay?.
Time to think about it..
But if you have New Year resolutions, how does it work?.
What are the principles that we can bring.
about changes in our lives?.
So here we go..
First principle is the 1\% power..
1\% power..
That you do something tiny on a daily basis,.
even if this is 1\%, and it's seemingly not obvious,.
the result will be exceptional..
Now, so often in any goal we are pursuing,.
we attempt to set up some of the,.
some high and unrealistic expectations, right?.
Like I'm gonna read the whole Bible in two months..
Okay, good luck to you, all right?.
You're not gonna happen, right?.
It's not gonna happen..
If you set the goal, like next year, 2024,.
you're going to read the entire New Testament in one year,.
that may be more realistic, right?.
But sometimes when we set up a New Year goal,.
we have some unrealistic expectations,.

$^{121}$and it will require exceptional willpower..
That's why we don't set New Year goals anymore,.
because we find ourselves failing all the time, right?.
You don't wanna fail, that's why you don't wanna set goals,.
right?.
That's kind of pessimistic, right?.
It's kind of sad..
But this book says all kinds of changes happen atomically..
Small, tiny changes..
And far too often, tiny improvement.
will create long-term change..
And that's how life change works..
So think about daily habit..
Daily, what can you do with just 1\% effort, right?.
There's a chart in the book, I find it very insightful,.
talking about behavioral psychology..
When you think of growth,.
we often think in a straight line, right?.
Linear progression, that's the red line right there..
But actually, life change happens in the curve..
In the curve..
Notice the beginning of the curve..
The progress is not obvious..
And the author named the first phase,.
you know, it's not obvious in terms of the result..
You don't see a lot of increase,.
you don't see a lot of obvious improvement yet..
That, the yellow part, the author says,.
James Clear says, is the value of disappointment..
Because you have done so much, right?.
And there's no result..
And people usually give up at this point..
So the encouragement is, if we do the 1\%,.
which everyone can do, right?.
And you persist..
And then, exponentially, changes will happen..
So if you do 1\% improvement for 365 days,.
the result at the end is 37.78 times better..
Right?.
Isn't that cool?.

$^{161}$If you do just 1\% every single day, 365 days,.
you will get the improvement result for 37 times..
But if you do 1\% decline, or you do nothing about it,.
or you keep on your bad habits,.
1\% declines for 365 days, your result is close to zero..
There's no improvement at all..
Now, you can talk about any kind of result,.
or any kind of goal..
Relationship, family, career, learning a new skill,.
emotional health, spiritual life, spiritual health..
If you do negative talk all the time, right?.
You will have the graph, have the line at the bottom..
You will decline, or you will see zero changes in your life..
But if you do tiny habit, 1\%, for 365 days,.
you will see a lot of improvements..
Now that makes me think of wisdom literature,.
because all this wisdom sayings in the Bible,.
especially in the book of Proverbs,.
are usually one verse with two sub-verses,.
two parallel statements..
Some of the wisdom sayings are two verses.
with parallel statements, but they are repetitive..
They are tiny little statement, little and often..
That is the spirit right there..
Proverbs 1.7, that's the beginning of book of Proverbs..
The fear of the Lord is the beginning of knowledge,.
but fools despise wisdom and discipline..
The fear of the Lord is the beginning, right?.
The beginning, remember the beginning part.
of the life change curve..
It doesn't see a lot of improvement, right?.
At least in the beginning, but it's so important..
Now Proverbs doesn't do big talk..
He always do small talks..
Talking about beginning here..
Talking about you persist from the beginning,.
then you will see along the line, 1\%,.
remember 365 days you see potential changes, right?.
Three, 30 some times changes..
Now some of the couples who got married in this church.

$^{201}$went through the premarital counseling with me..
And the premarital counseling, I have about 10 sessions,.
and I always end with this one..
It's a book called Love List..
Now this book, we read this book,.
my wife and I read this book in our 10th anniversary,.
some years ago, right?.
So we went on a short vacation and we pick a book.
and we read together, right?.
How romantic is that?.
(audience laughing).
So we read the book together..
It's a tiny, tiny book called Love List..
And it's a very interesting book and it's very insightful..
Because people always think that a great relationship.
requires a lot of energy, requires a lot of expectations..
Right, you gotta set a very, very high standard..
You gotta do tons of things.
in order to have a mature, good relationship..
The book says no..
It doesn't work that way..
Have you think about it?.
Think about it..
It doesn't work that way..
A successful relationship, a good relationship,.
doesn't work that way..
It works because there are tiny, daily habits.
the couples or the people in the relationship.
are committed to do..
And everyone can do it..
And at the end of the book it says,.
once a day you do two things..
Once a week you do two things..
Once in six months you do two things..
Once in a year you do two things..
How many things you have to do?.
Eight, right?.
Only eight things..
And once a day it says take time to touch..
So, like for the couple, you take time to touch..

$^{241}$If only for a minute, right?.
So in my household we express the touch.
by patting on the head..
So I pat on the head of my wife, of my children..
Yeah, you know, don't attack me, okay?.
It's like a dog, right?.
It's kind of cute, right?.
So what I do every single day is that I take the time,.
whether she's doing dishes or whatever,.
I go over there and, you know,.
I give her a touch for one minute..
That's it..
The second thing I do daily is to find something.
that make us both laugh, right?.
Life is stressful..
You want your family, you want your relationship.
to be joyful, then you have to find something.
every single day that make you both laugh, right?.
Now that's easy to do, right?.
Scrolling on the phone and I saw some TikTok video..
It's hilarious..
I show it to her, right?.
She doesn't laugh, I laugh, right?.
That's okay, right?.
And that's what we do every single day..
Once a day, we do two habits..
Now again, the essence is it's the tiny habit.
that can change a lot of different things..
Little and often, that is the key..
Some of you know that I took a PhD seminar.
the middle of this year..
This is the toughest class I've ever had..
Not only me, but other students taking the class with me..
They all think this is the toughest class.
in the entire degree..
The class is advanced Hebrew grammar..
Sounds fun, right?.
Sounds really fun, right?.
Advanced Hebrew grammar, right?.
So in order to pass this class, you have to pass a test,.

$^{281}$which is a biblical Hebrew vocabulary test..
And it tests you for 100 random words,.
any word that appears in the Bible more than 50 times..
All right?.
And the test will pick some of the random words.
and then list them out in Hebrew,.
and then you have to give a translation in English..
And in order to pass the exam,.
you have to get 87 over 100..
87, it's not 50..
I heard it wrong..
No, it's 87..
And half of our class failed..
Half of our class succeeded..
It was kind of scary, right?.
So before the class, I took about two months..
Actually, I delayed my start..
I end up like five weeks before the class, before the exam..
So what did I do?.
I simply used my downtime..
You watch TV, there's commercial, right?.
Three minutes or four minutes..
And then I go into the vocabulary book,.
and then I try to memorize them, right?.
I try to connect them and try to memorize it..
And then there's sometimes downtime, right?.
When we have conversations together..
I'll scroll up, open the book in my phone,.
and then I'll review the vocabularies..
And that's what I do..
And there are some other students.
who try to cram all these words into their brain.
one night or two nights, right?.
It's not going to work, right?.
So I passed the class..
I passed the vocabulary test..
What did I do?.
I did nothing significant..
It's little and often, right?.
Most of you who are students, this is what you do, right?.

$^{321}$This is what you do..
I remember I took Chemistry 101..
There were 10 exams in the entire semester..
I failed the first one because I missed the exam day..
I show up on class, it was an exam day..
I had nothing in my brain..
So I failed the first test..
And there are nine more tests to go..
And what did I do?.
I made up my mind..
I commit myself to be successful in this course.
because I don't want to retake it, right?.
It's gonna cost a lot of money..
So what did I do?.
Every single class, right after we are dismissed,.
I go to the library and I review for one hour, right?.
Now, so often we attempted to think,.
"Okay, maybe we should go home, right?.
"If I go home, I get four hours to study, right?.
"And do you want four hours to study.
"or one hour to study?".
Of course you want four hours, right?.
But what I did was I go to the library.
right after the class, while the material is still fresh,.
I'll review it only for one hour or so..
And then I'll go home..
And then I ace every single test, all nine of them..
And then I got an A at the end..
Again, it's the tiny changes, insignificant habits.
that would pave the path for fruitfulness..
This is what we do for spiritual life as well, right?.
Proverbs 9, 10, "The fear of the Lord.
"is the beginning of wisdom..
"Knowledge of the Holy One is understanding.".
You see the parallels here?.
What is the fear of the Lord?.
The fear of the Lord is that you fear God.
and you relate to God in a proper way..
And that is the real knowledge of the Holy One..
You know God and you fear Him and you are related to Him,.

$^{361}$you have a right relationship with God, your Creator..
That is the foundation of small things..
Small habits will go a long way..
Now I want to test you..
How many of you do push-ups?.
How many of you?.
How many push-ups can you do?.
How many?.
25?.
I do 40 push-ups every single day..
It's not much..
How did I begin?.
I begin with 10, right?.
And I add one the next day, I add five next week, right?.
And now I do 40 push-ups every single day, right?.
It's just my way to strengthen myself..
How about plank?.
How many of you do plank?.
And how long do you do?.
I do four minute plank every single day..
How did I start?.
I start one minute, right?.
And then one minute and 15 second, one minute and 30 second,.
two minute, three minute, and four minute, right?.
And my goal next year is five minutes, right?.
Now how do we do that?.
Small, tiny habit of 1\% power..
This goes the same way for our Bible reading.
and our daily prayer life, right?.
It doesn't have to be huge, significant change..
It's not a massive action that God is calling for..
It's just tiny, little habit that will change you,.
that will effect life changes..
So 1\% power, that's the first principle..
The second principle is incentives.
and this incentive are effective..
The behavioral psychologists find out.
if there's pain involved, if there's suffering involved,.
people are more likely to follow through,.
especially financial pain..

$^{401}$That's why people pay money to go to gym, right?.
That's why you go to exercise class and you pay, right?.
After you pay the amount,.
there's more likelihood for you to go..
The more likelihood you keep on doing physical exercise..
When money is involved, the psychologist tells us.
your chances of success increase by three times..
You pay to go to school, right?.
You better go to school because you pay for it, right?.
There's a website called stick.com.
and you sign up, you set up your goal.
and it will ask you to pay money.
to something that you hate, like an anti-charity, right?.
For those of you who don't support Trump,.
you donate to Trump..
If you fail in your goal,.
if you fail to do your daily exercise.
or fail to meet your goal,.
you're gonna donate money to Trump that you hate, right?.
For those of you who hate Biden, you can choose Biden..
And then when you meet your goal, you report it.
and then you stop the payment for your credit card, right?.
Money is involved..
If you do that or you don't,.
you will kind of lose money.
and you better work on it, right?.
That's how it works..
There's reward and disincentive..
If there's incentive in your life-changing goals.
or new US solutions, the psychology is applied..
The rate of success is better.
if there's incentive or disincentive..
That's the same way in wisdom literature..
It says one man gives freely if someone is generous to give,.
yet gains even more..
Another withholds and comes to poverty..
A generous man will prosper..
He who refreshes others will himself be refreshed..
What is the incentive here?.
What is the reward here?.

$^{441}$The Bible says that if you honor your wealth,.
you honor God with your money, you will increase..
You will benefit others and you will benefit yourself..
That is what Jesus said..
You lay up treasures in heaven, not on earth..
And then your reward will be heavenly, right?.
If you will believe in the gospel,.
you believe in eternal life, right?.
The reward is eternal..
The reward is there..
So once you think about, okay, I need to do maybe New Year..
One of my resolutions is to be faithful in tithes and offering,.
to be generous, right?.
Not to hold tight to the wealth God has given to me,.
but to benefit other people,.
to give it, to support the ministry or support the church.
or support the cause that is heavenly,.
that is gospel related..
You need to understand there's reward.
and the reward will be great..
There is incentive there..
But if you do not want to do that, if you withhold,.
then there's this incentive..
Your life becomes poor, right?.
It's always needy..
If you practice generosity, then God will honor you,.
He will bless you,.
He will continue to make your life rich..
That is the incentive and disincentive there..
So think about your New Year resolutions..
What are the incentives or what are the disincentive, right?.
Yeah, one of the disincentive I can think of,.
if you want to read the Bible next year,.
if you fail your goal, you sign up stick.com.
and you give me \$500,.
that would be a great motivation for you, right?.
Because you don't want the \$500 to go into my pocket, right?.
And I promise you, if you fail in your goal,.
you give me \$500 at the end of the year,.
I'm gonna keep it for you..

$^{481}$I won't use it, I won't spend it,.
and then the following year, if you meet your goal,.
I'll give it back, right?.
And those kind of incentives or disincentives,.
you can think of in your New Year resolutions,.
and that is wisdom..
Number three, specifics are better..
Now the wisdom literature.
use a lot of different illustrations to talk about wisdom,.
usually animals and creatures on Earth..
Sometimes these creatures don't seem to be strong,.
don't seem to be wise, but the observers,.
you know, wisdom teacher are good observers..
They observe the creation,.
they observe how the creatures work, right?.
And then they learn lessons of wisdom..
They learn how to be wise..
Chapter six, it says, "Go to the end,.
"use lucker, consider the ways and be wise.".
They don't have commander, they don't have ruler,.
they don't have king, yet they store provisions in summer,.
and gather the food at harvest,.
at the right time, specific time, right?.
And that is wisdom..
And psychologist tells us the more specific your goals are,.
the likelihood the success will be..
If your goals are general, you know, non-specific,.
you are planning to fail..
And there's a study done in UK on physical exercise..
So they divide the people into two groups..
First group, group A, is that people who are committed.
to do exercise on a daily basis,.
and they do the checkbox, okay?.
I'm gonna commit to do pushup, do crunches,.
do plank, walk a little bit, right?.
They check the checkbox..
That's group A..
And they follow through,.
and the success rate for this group, group one, is 35\%..
Group B, they also do the checkbox,.

$^{521}$but they write specifics to their goals..
Like time of the day, day of the week that they do exercise,.
and how long exactly, how many minutes that they going to do..
20 minutes they gonna check, and then write it down..
And they specify the exact time, place,.
where, when to do the goal..
And the success rate is 91\%..
Three times higher..
So the more specific that you go,.
the more successful you will be, right?.
Your goals will be achieved,.
and the likelihood to have result is higher..
Let's say you wanna commit yourself.
to restructure your worship life, okay?.
Every single week, the Lord's Day,.
you're gonna come here and worship, right?.
And that's a very, very good goal, right?.
That's a very good goal, because as the people of God,.
we need to worship Him, right?.
The day of Sabbath, or day of celebrations,.
day of worship, we set aside time for God, right?.
Every single week..
That's a good goal, but it's very general..
To be more specific, you should say,.
I'm gonna commit myself to worship God.
every single Sunday, and I'm gonna come here.
at 11.16 in the morning..
By 11.16, my butt is here..
And I'm gonna sit at the same place every single week..
The when and where are specified,.
and the rate of success is higher..
That, the same way goes for any kind of goal..
Your physical health, emotional health,.
or even your spiritual health..
That's how it works..
And then another principle, number four,.
and that's the last one I'll continue next week..
Is community influences you..
It has power to impact you..
The people that you hang out with,.

$^{561}$the people that you build friendship with,.
the people who are part of your family.
will have impact on you, and the impact is significant..
Proverbs says the same thing, 13.20..
Whoever walks with the wise becomes wise..
And the friends or the companion of fools will suffer harm..
If you wanna experience destruction,.
you hang out with the fools..
If you want to have wisdom,.
if you wanna have successful life,.
you walk with other who are wise as well..
And there's a study on the effect of people.
on one's behavior..
You know, how people around them behave.
will affect that specific person..
So that study is very interesting..
They send out actors, actresses, to the buffet line..
You know, you go to a buffet restaurant,.
and you pick up some food, right, on your plate,.
and the actors are being taught to go into the buffet line,.
and they pile up all kinds of food on the plate..
You know, the more they get, the better, right?.
So the actors act that way in order to influence.
those who are not acting..
You know, the typical restaurant customers,.
they're trying to line up at the buffet to get their food..
And they wanna see how this will influence the next person,.
and how much food they will get, right?.
It's very interesting..
The result is, the result for men and women are different..
Men has no impact..
For women, because the next person is piling up more food,.
the women are more likely to get more food..
The point is, your community works..
You gather, you study the Bible with other people,.
you join a class for exercise,.
you go to a fellowship group, you study in the library,.
the environment, other people will have impact on you..
So in order to increase the rate of success,.
you have to find the right community..

$^{601}$You cannot do solo..
Let's say you wanna improve your academic study..
What is the best way?.
For me, if I study at home, it would not be helpful for me,.
because my wife doesn't study..
She's listening to music and watching TV show..
Huge distractions, right?.
Huge temptations..
I have to go away from her, right?.
I have to go somewhere that is fruitful.
for my study, like library, right?.
Or be myself, be by myself..
So the whole point is, you have to find friends.
and people who can influence you positively, right?.
So there are more insights I can share,.
or more wisdom I can share, and that will be next week..
But today, I would like to conclude.
with some of the suggestions I have..
It's for the new year, 2024, new year..
I'll suggest three areas you should think.
about new year resolutions, new year goals for you..
First is spiritual health..
How important that is, very, very important..
It's very important you will take time.
to read the word of God, right?.
Not just reading more of the Bible,.
it's learning how to read better,.
how to be a good reader of scripture,.
praying, prayer life, it's part of our spiritual health..
Right, we receive the word of God.
through reading scripture, we talk to God.
in our daily prayer, right?.
And part of the good things about going to fellowship.
is that we can pray together..
And by praying together, somebody is praying,.
people around you are praying,.
and it will influence you, impact you to pray..
Worship is part of it..
Take time that is set apart for worshiping God.
to come in person, right, if possible..

$^{641}$Come in person and worship God..
Spiritual health, you know, you need to set some goals,.
or maybe you should read a spiritual book,.
a Christian book, a good book.
that will help you grow spiritually..
2024, maybe you should pick one of the good books out there,.
right, in the library, or you want some recommendation.
I can give to you..
In the morning service, I talk about this,.
and right after, a couple of people talk to me,.
"Hey, can you recommend one or two.
"for me to read through in 2024?".
Spiritual health, emotional health..
New year, 2024, it just dawned on me this year.
that the emotional life for every single one of us.
is so important..
How do we manage different emotions?.
How do we manage them well?.
I'll have more to talk about next week,.
but this is something for you to think about..
You know, your emotional health impacts.
your spiritual relationship with God,.
and your spiritual relationship with God.
impact your emotional life as well..
And number three is physical health, physical health..
You wanna be stronger and stronger, right?.
And live a healthy lifestyle, right?.
Stay strong and stay healthy, right?.
Set some goals for yourself, right?.
And then remember, 1\% is powerful, right?.
Do tiny little changes, little and often,.
and then at the end of the year,.
you get results 30-some percent of 30-some times results..
And let's do that..
And my final prayer for you,.
and we're gonna pray together,.
is that dear God, as this year comes to an end,.
help me to know what I need to get rid of.
in order to have more room for you,.
for what you have in store for my life..

$^{681}$Let's all pray together..
Father, we wanna give you thanks.
because in this entire year,.
from the beginning to the end to today, you are faithful..
Yes, we go through our life journey with ups and downs,.
and with all these turns and twists,.
but we continue to acknowledge that you are faithful..
You're faithful to us,.
and thank you for your loving kindness..
In the next year, in the new year to come,.
we know we are not in control,.
but we do know that we need to make room for you,.
for your blessing, for your grace,.
for your work in our lives..
And I pray that you will continue to help us.
to be obedient, to be open,.
to listen to your voice,.
so that in the next year, 2024,.
we look forward to walking with you,.
to experience you, to go deeper with you,.
to get to know you better, get to know you more..
Help us to rebuild, reestablish.
some of the tiny daily habits,.
habits that are life-changing,.
habits that are seemingly tiny and insignificant,.
but will make our life fruitful..
And I pray that we'll read your scripture more,.
pray to you more, and worship you more, and give more..
And I pray that you'll grant us wisdom.
to manage our emotions..
Help us to go to you when we are in need..
And Father, we commit our whole body into your hand,.
knowing that our physical health.
will impact every single area,.
our emotional life, our spiritual health, and vice versa..
And I pray that we'll grow holistically..
So we put our hope in you, knowing that in year 24,.
you have so much in store for us,.
because you are our good God..
I pray that we'll experience your love and your grace..

$^{721}$And help our church to grow together,.
so that we will honor you,.
and proclaim your name in Christ, in praise..
Well, God bless you, and see you next year..
(claps).
\newpage



\section{}
\label{sec:8jU_KxWNams}
\textbf{原子智慧 1 - 陳崇基牧師}
\newline
\newline
連結: \href{https://youtube.com/watch?v=8jU-KxWNams}{\texttt{https://youtube.com/watch?v=8jU-KxWNams}} ~~~~ 語音日期: 2024-01-02
\newline
\newline
\hyperref[sec:SCf1z8t8nOA]{\small{< < < PREV SERMON < < <}}
~
\hyperref[sec:index_chronic]{\small{[返順時目]}}
~
\hyperref[sec:index_scriptual]{\small{[返順卷目]}}
~
\hyperref[sec:YZjpAlGepXw]{\small{> > > NEXT SERMON > > >}}
\newline
\newline
$^{1}$Harry Tai的節目好.
過去兩天我都病了.
所以真的不是很好.
年尾又病了.
因為我家有幾個人都病了.
所以沒辦法.
我們重災區.
我都想著可以過年.
誰知道都過不了年.
不過不要緊.
我們繼續去事奉.
我們快到新的一年了.
明天就是了.
是不是.
雖然都是這樣過一天.
也是這樣過.
二十四時也是這樣過.
這些日曆或者年曆的事情都是人為的.
我們有一些機會.
我們能夠去回望一下.
然後去做一些計劃.
2024年怎樣也好.
也是一個新的一年.
大家可能想到一些不同的新年目標.
有沒有.
財政上可能又要想過.
今年新的一年的時候.
財政上財務上應該怎樣去做一個好的管家.
或者減肥.
是不是.
也是時候了.
想了幾年.
今年終於要了.
一直說了幾年.
或者想學一些新的東西.
也可能是.
男女老幼也好.
我們都可能想學一些新的東西.
或者我們作為基督徒.
我們也會想一下.

$^{41}$也是時候了.
要看看聖經.
說了五六年了.
終於到新的一年了.
又想一想.
是不是應該要看聖經看一次.
或者看某部分等等.
這些不同的考慮.
我們都會想的.
而且也是一個好的時間去想一想.
過年了.
去了一年.
然後再迎接新的一年.
應該怎樣做呢.
很多時候教會我們都會定一些新的目標.
不過我們教會很少會做這些.
因為很多教會.
到了新的一年的時候.
我們就會有一個橫額.
然後有一個Banner.
有一個很大的口號等等.
或者譬如我們教會要Expand.
或者教會要成長等等.
說一些很大的東西.
但是我這次想跟大家分享.
不只是說我們有一些很遠大的目標.
而是到底一個生命的改變是怎樣發生的.
怎樣有生命的改變.
怎樣在我們生命裡面可以達到一個目標.
到底當中有些什麼事情會發生呢.
以致我們能夠有所謂的Transformation.
我們生活裡面.
生命裡面可能有一些改變.
有一些向前.
我這個講道的標題叫做原子智慧.
Atomic Wisdom.
Atomic不是原子彈那種.
原子意思是很小很小的意思.
一些好像很不顯眼的東西.
做下去好像很少的東西.

$^{81}$這些東西應該怎樣去到.
讓我們了解我們生命裡面向前.
生命裡面去到成長的時候.
到底應該怎樣留意呢.
這個標題是來自一本書的.
叫做Atomic Habits.
James Clear所寫的.
翻譯成中文.
其中有一個人推薦我看這本書.
我看這本書的時候發現裡面有很多的智慧.
裡面說什麼呢.
Atomic Habits.
原子習慣.
他說我們生命裡面要做的事情.
或者我們想達到的一個目標.
很多時候不是一些一次過做很大的事情.
我們就會達到的.
相反是一些很日常.
每一天無時無刻做一些很小很小的事情.
我們就會達到所謂的成功.
達到那個目標.
或者我們追趕著一個目的.
怎樣能夠達到呢.
無論是一些學業,事業.
或者我們人生等等.
是一些很細微的東西.
以致我們能夠去達到的.
所以我看這本書的時候.
我就想起裡面有很多的靈感.
都是來自Behavioral Psychology.
一些行為心理學.
行為心理學是做什麼的呢.
就是你去觀察和去明白人性是怎樣的.
然後你將這些的Insight放進去一個Life Change的Process.
一個生命改變的過程裡面.
這本書就結集了很多這樣的智慧.
我就想起看這本書的時候.
聖經裡面的智慧文學.
Wisdom Literature.
大家記得了.

$^{121}$Wisdom Literature在哪裡.
在聖經中間.
是不是.
智慧文學.
有幾本書是屬於智慧文學.
詩篇裡面有不同的詩歌.
《約伯記》.
《箴言》.
《傳道書》等等.
《箴言》特別是Proverbs.
就好像這些所謂行為心理學裡面所觀察到的一些結果.
你知道《箴言》Proverbs.
你看的時候.
你想到一個口吻.
你get到一個口吻.
那個用詞.
就是一個爸爸去教他的子女.
要做一個有智慧的人.
Live a life of wisdom.
就是教一個他初出茅廬.
沒有什麼經驗.
未成熟的一個人.
讓他的生命能夠成功.
能夠Successful.
所以在《箴言》裡面說很多一般性的真理.
你看《箴言》的時候.
或者看《智慧文學》.
不是在說一些Guarantees.
不是在說一些保證.
他不是說你這樣做的話.
你一定會成功.
而是在說人性裡面一些一般性General的principles.
大概一般性是這樣的.
勤力的話.
你就會得到有成績.
你懶的話.
你就會自己受苦.
等等是一些General的真理.
General的principles.
並且在這些《智慧文學》裡面.

$^{161}$是結集了外面.
聖經以外.
一些的觀察和一些的真理.
所以就將它收集了.
然後放在上帝的啟示底下.
All truth are God's truth.
這些東西放在一起.
然後就教導我們怎樣去過一個智慧的生活.
所以在這個講道裡面.
我想做一些integration.
做一些的結合.
怎樣將這些所謂世界裡面.
行為心理學裡面所說的話.
或者這本書裡面所說的話.
其實不是單單這本.
有另一本我也在看.
叫做《Much》.
也是說類似這些東西.
說一個人的本性.
或者人的心理等等.
明白人的本性.
然後怎樣去明白呢.
Life change.
Atomic wisdom.
所以我們就在這兩個星期.
我們試一下去看一看.
還有留意一下.
到底我們所謂新的一年.
我們有一些的New Year resolution.
一些新年可能要有的目標等等.
這些目標定下來之後.
到底應該怎樣去想.
怎樣達到呢.
由A到B.
中間有些什麼我們需要留意的.
我想在這裡做一些的總結.
第一件事.
他講到1\%的威力.
我們想到達成一個目標.
例如你新的一年.

$^{201}$你想減肥.
你想減多少磅.
10磅好嗎.
10磅好像太多了.
好 3磅吧.
平時吃叉燒都吃3磅.
3磅.
計算你3磅.
我們以為很多時候.
成功地達到那個目標的時候.
我們就要用很大很大的氣力.
很多的行動我們才能做到.
任何目標我們都是這樣想.
但是這裡的智慧告訴我們.
是不需要的.
我們以為定一個目標.
然後定一些預期.
並且那些預期是.
一些很高的預期.
一些很難達到的期望.
我們以為這樣才能達到那個目的.
這些行為心理學裡面講的.
不是這樣.
他說真正的改變來自什麼.
是一些很小很小很小的改變.
是每一天都會做的.
是每個人都能做到的.
裡面有一個表.
你就會看到.
一個改變應該是怎樣.
我們經常想著一個成長或改變.
達到一個目標.
就好像黃色這條線一樣.
一定是直線上去的.
但是真正來說.
一個改變是怎樣發生的.
就是曲線上去.
你留意開頭那裡.
開頭那條線.
你一直去做一些事情.

$^{241}$在時間裡面一天一天.
是不會有明顯的進度的.
特別是開頭的時候.
這位作者就說.
這個就是黃色那個地方.
就叫做失望谷.
就是Valley of Disappointment.
如果你想減肥.
想每天做運動.
你做著做著.
做了一個星期有沒有效果.
沒有效果的.
沒有效果.
你就想著有效果.
你就很容易會放棄.
但是他說如果你一直做.
如果你過了這個所謂失望谷.
你才會看到有很大很大的改變.
這個給我們一些益處.
讓我們想到.
我這樣熬過之後.
小小的事情做了之後.
然後我們的改進就會突飛猛進.
所以他說如果你每一天.
你做一個百分比的改進.
你計算一下365天.
2024年.
你所得到的改進的程度.
就是37倍.
37倍就想一下.
剛才你想減肥減多少磅.
3磅是吧.
如果要拿到37倍的效益.
你會減到多少磅.
30磅一次.
開心.
但是如果你每一天.
你不做任何事情.
或者你仍然活在壞的習慣裡面.
你1\%的改進.

$^{281}$計算一下365天.
你所收益到的改變.
就是幾乎零.
所以這裡很能鼓勵你.
只要1\%的改變.
你365天之後.
你所收益到的效果.
就是37倍.
很好是吧.
所以如果你想達到什麼目標.
比如關係上的目標.
婚姻關係 家庭關係 朋友關係等等.
事業上.
或者一些生命裡面的技巧.
一些技巧你要學的.
或者情緒的健康.
屬靈的健康等等.
你想一下如果1\%每一天.
你所最後得到的收益.
其實就是37倍.
所以不要小看那個1\%.
特別開始那段時間.
其實《箴言》也是這樣說.
《箴言》一章第七節.
那裡說.
敬威爺說是知識的開端.
慾望人就是藐視智慧和憤悔.
敬威爺說.
The fear of the Lord.
就是知識的beginning.
是開始.
起始.
那個好像一些很細微的東西.
或者很基本的東西.
好像是看不到的.
但是你在那個起始的時候.
敬威爺說你和上帝有一個正確的關係.
由這裡起始的時候.
你又忍耐.
過了那個valley的時候.

$^{321}$你就發現.
咦 我的生命裡面的智慧.
就越來越累積.
生命的改變是應該這樣的.
任何一個關係的維繫.
或者生活裡面的經營.
其實都是小小開始.
還有在那個起始裡面.
你去留心的.
在我中間有一些couple 夫妻.
跟我做過一些婚前輔導的時候.
你記得我婚前輔導嗎.
我有十堂.
最後一堂我們說什麼呢.
最後一堂我講一本書叫做《Love List》.
《Love List》這本書是我和太太十週年紀念的時候.
我們度假了.
我們那幾天度假的時候.
我們就看書.
是不是很浪漫.
我們看書很薄的.
我們看了那本書 很薄.
叫做《Love List》.
《Love List》講什麼呢.
它說怎樣去維繫到.
或者經營到一個美滿的關係.
我們就以為一個美滿的關係.
可能和我沒有什麼關係.
因為太遙遠的事.
和太高的預期.
或者我們覺得這個美滿的關係.
一定要很多很多的能量放進去.
input進去才可以做得到.
但是那本書很奇怪.
它不是這樣講的.
它說不是的.
它說一個美滿的關係.
只需要做些什麼呢.
它就說有一個love list.
每一天你做兩件事.

$^{361}$然後它說每個禮拜你做兩件事.
每半年你做兩件事.
每一年你做兩件事.
其實加起來只有八件事.
很容易吧.
它說每一天要做些什麼呢.
它說take time to touch.
觸摸一下大家.
夫婦之間.
我說搞錯了.
這麼簡單.
就是這麼簡單.
就算是一分鐘也好.
你試一下去碰一碰對方.
我家裡喜歡拍拍他的頭.
所以我每一天定時定候.
我就像小狗一樣拍拍他.
他在洗碗也好.
他在摺衣服也好.
我去拍拍他.
這樣就搞定了.
我做完了.
take time to touch.
然後就是find something that.
makes you both laugh.
是不是.
在電話裡看到什麼好笑.
讓他笑一下.
他不笑也好.
我也笑一下.
容易嗎.
其實兩分鐘就做完了.
但它說你每一天這樣做.
這就成為你關係維繫的基本.
絕對不是難的.
但我們肯去做一些tiny changes.
1\%就已經足夠.
我記得跟弟兄姊妹說過.
在今年年中的時候.
因為我現在正在進修.

$^{401}$有一堂最難的.
就是advanced Hebrew grammar.
即是進階的聖經希伯來文的文法.
那堂是非常難的.
我覺得每個人都很辛苦.
每個同學都覺得很辛苦.
到底一堂這麼難的科目.
怎樣去讀呢.
因為我們做牧師做了這麼多年.
我們都是生手的.
那些希伯來文.
我們要pass這堂的時候.
我們就需要整本舊約聖經.
裡面凡是出現過50次以上的字.
測驗就是測你.
你懂不懂那個字.
每逢一個字出現過50次或以上.
譬如哈維這個字.
Lord這個字.
God這個字.
Elohim.
那你要知道測驗的時候.
你就需要寫那個意思出來.
那這堂不是50\%pass.
87\%.
即是起碼要拿B+.
很難.
那我就決定怎樣做呢.
我決定那堂未來之前.
大概一個多月.
其實我都遲了.
我其實想著三個月.
怎知道拖著拖著變成四,五個星期.
在那四,五個星期裡面我做什麼呢.
我就不是天天熬夜讀書.
我只是每天有空閒的時間.
我在電話裡看那本書.
那個字典.
記一下.
或者和家人聊天.

$^{441}$當然有些時間是空餘的.
譬如你看電視.
它會有廣告.
那我就會看那字.
記一下那字.
那我就pass了.
但有些同學和我一起讀書的時候.
剛剛去到校園.
第一天上課要測驗.
一起讀書的時候.
他們很辛苦.
他們就想cramp.
想開通宵.
然後記那些字.
這就是如何的.
就是你慢慢慢慢一天一天.
很tiny的一些change.
你就可以達到那個所謂生命的改變.
你學一件事也好.
你想改進也好.
都是一些很細微很細微的東西.
你養成了這個習慣.
你就可以做到.
並且那個習慣並非是一些遙不可及的事.
是一些很tiny的一些change.
九章第十節.
敬畏耶和華是智慧的開端.
認識至聖者便是聰明.
你留意到這裡.
箴言很多時候是有A句和B句的.
有時候是一節.
有時候是兩節.
總之是有一個對聯.
敬畏耶和華是什麼意思.
認識至聖者就是.
認識我們的上帝.
和上帝有一個正確的關係.
這一個就是智慧人生的tiny change.
小小的開始.
所以不要忘記.

$^{481}$你每一天做一點的事.
只要那些事是一些好的習慣.
只是1\%而已.
365天你所收益到的是37倍.
是吧.
大家都有每天做運動的.
有嗎.
很少人說.
你做運動.
大家有沒有做push up.
想像一下.
你做了多少個.
35.
35啊.
很厲害.
我每天是做40個的.
當然我不想繼續多做.
因為多做會腦子.
我試過多做會痛.
有沒有做plank.
那些直身的.
叫什麼.
就是撐起自己.
看你撐多久.
你撐多久.
1分鐘.
我開始的時候也是1分鐘.
我現在4分鐘.
我4分鐘.
怎樣做.
30秒開始.
1分鐘.
2分鐘.
小小的tiny change.
你就會發現.
原來你所收益到的是30多倍.
這是第一個原則.
第二個原則就是incentive.
或者disincentive的power.
賞罰的效果是怎樣.

$^{521}$那些行為心理學家說.
如果你想有一個life change.
最好就是有賞罰.
你對你自己有賞罰.
特別有錢的賞罰.
如果跟錢有關的話.
或者你會不見錢的.
你的生命的改變.
那個動機.
就會更加大.
是吧.
譬如你想每天做運動.
如果你不做的話.
你就要給我50元.
好不好.
ok.
好啊.
你等一下可以sign up.
如果一個月裡面.
不可以堅持做運動的話.
你就給50元給我.
是吧.
在網上有一個網頁.
或者一個網站.
叫做stick.com.
stick.com是什麼.
就是你可以sign up.
然後set up你自己的目標.
然後你sign up的時候.
你就可以用你的credit card去付款.
如果你達到目標的話.
你去報告說我已經達到目標.
這樣他就不會charge你.
如果你達不到的話.
你就會虧錢.
並且你虧錢.
就是選一些你最討厭錢去的地方.
所以有所謂什麼.
anti charity.
你最討厭的.

$^{561}$譬如你不喜歡川普.
你做不夠運動.
你就給錢川普.
那你做不做.
你當然做.
如果你不喜歡拜登.
你就sign up.
那個anti charity就是拜登.
做不夠運動.
我就給錢他.
你有沒有搞錯.
怎麼會給他.
所以你有incentive.
有disincentive的時候.
你就會有更大的motivation.
是吧.
譬如你新的一年.
你sign up我網上的課程.
你先給錢.
一鋪過你都給他.
yeah right.
exactly.
你給他.
你會不會上課.
你多大都上.
可能你會miss.
不過你都覺得不要給他好處.
最重要不要給人好處.
你有disincentive的時候.
你就更多的會去改進.
曾賢都是這樣說.
有施捨的人.
錢財就會增添.
孤寒過度的.
反而會貧乏.
慷慨待人的.
必然會豐原.
滋潤人的.
自己的心也會被滋潤.
這裡說什麼.

$^{601}$奉獻.
Amen.
當你奉獻的時候.
當你practice這件事的時候.
每個人都想受奉獻.
有誰不想受奉獻.
每個人都想學國粹.
但你開始做的時候.
tiny changes.
5\%開始.
你要學習慷慨過你的生活.
上帝給你的資源.
給你的資財.
將財寶積聚在天上.
這些全部都是incentive.
是有reward.
是給你有reward.
你想起這件事的時候.
你就發現就是.
《箴言智慧文學》所說的.
如果你肯去施與.
自然你都會得到祝福.
如果你孤寒的話.
你反而會缺乏.
經常覺得你自己不夠.
慷慨的去到生活.
你自己的生命.
新生命都會得到豐原.
滋潤到人.
益到人.
同時益到自己.
就是reward的psychology.
對嗎.
如果你經常想著.
奉獻是一個很慘的東西.
《箴言智慧文學》說.
你想清楚.
那個incentive.
你真的相信一個永恆的福音是永生的.
你就會想清楚.

$^{641}$咦 我不是想明白.
這個奉獻是什麼意思.
可以祝福到上帝的國.
祝福到人家.
自己也能夠受益.
就是tiny changes.
365日之後.
你的生命會是多麼的風俗.
第二個原則說完了.
第三個.
他說越具體的目標是越好的.
越具體的目標.
譬如你下一年.
你打算讀經才算.
讀聖經.
讀新約好嗎.
讀新約27卷算.
一年讀完.
你這樣定目標是不行的.
一定要有具體的細節在裡面.
你才能夠有更大機會.
可以成功.
在智慧文學裡面.
很多時候都會說到一些.
受造之物.
他們會觀察一些東西.
智慧文學就是這樣.
不是直接啟示下來的.
不是thou shall.
或者thus says the Lord.
而是說如此說.
好像先知書一樣.
智慧文學是知者去observe.
去觀察人的人性.
大地裡面的本性.
然後領略到智慧.
這個就叫做wisdom literature.
很多時候會說到一些受造之物.
或者一些動物等等.
這裡說到螞蟻.

$^{681}$螞蟻做什麼.
它說懶惰的人.
你們觀察observe螞蟻的動作.
就可以得到wisdom.
螞蟻沒有君王.
沒有領袖.
尚且在夏天時間.
正確的時間.
它就預備食物.
收割的時候.
它就儲存糧食.
是很specific.
還有很具體的.
這個就是生命改變的智慧.
你需要明白.
在裡面原來要有一些目標.
不只是目標這麼簡單.
一些具體的細節.
你成功的機會就更加大.
在英國做了一個研究.
研究一些人去做運動.
適合我們新年.
他就分了兩組人.
第一組人.
他們說新的一年.
我們決定每一天做運動.
於是他們給一個表格給他.
他就在表格上check.
我每一天要做1,2,3,4的運動等等.
然後跟他follow up.
看看他們的成功率有多少.
這是第一組.
第二組.
不只是那些表格.
他還要親自寫.
specifically.
具體來說.
哪一天.
哪個時辰.
會做些什麼.

$^{721}$還要做多久.
那些人就寫.
哪一天哪個時間.
哪一天我們做20分鐘.
這些大量的運動.
所以他將時間,地方都放在裡面的時候.
他的success rate.
成功率是大很多的.
第一組成功率有多少呢.
35\%.
100人裡面只有35人真的可以堅持.
但是第二組.
是有具體的時間地點.
就有91\%.
91\%.
譬如你下一年.
你說我敬拜的生活做得好.
每個禮拜日我就要來敬拜神.
這是第一組.
第二組是什麼呢.
不只是我每個禮拜日來到敬拜神.
我還是9點31分就會到.
9點31分我會到.
並且我進來的時候.
我坐在同一個位置.
坐在那裡去敬拜神.
你的成功率就會更加高,更加大.
所以當我們看.
好,我們下年.
我們看新約聖經.
對嗎.
不可以就這樣看.
不可以就這樣定目標.
具體的細節是什麼.
每天早上我在幾點鐘.
或者每晚我幾點鐘.
看十分鐘.
你這樣定了的時候.
你的成功的機會率就會大很多.
這個就是life change.

$^{761}$生命的改變.
無論你想做exercise也好.
或者定任何新年的目標也好.
都需要有這樣的智慧.
第四方面.
然後有其他的我們下個星期再說.
第四方面就是.
the community will impact you.
你的社群.
是對你有影響力的.
你的朋友或者你身邊的人.
是對你有很大的影響力.
特別當你要去到生命改變的時候.
十三章第二十節.
他說與智慧人同行的.
必得此為.
與愚昧人作伴的.
就必受虧損.
你想得益的話.
你就跟一些有智慧的人.
一起見面.
跟他聚會.
請他吃飯.
跟他hang out.
但如果你想虧損的話.
就自然會選一些愚昧的人.
愚昧的人就是.
我會每天做exercise.
我們一起吃飯.
愚昧人.
思維的人是什麼.
我會做exercise.
並且每時每刻.
哪個時間我會做多久.
這些就叫做wisdom.
你跟這些人hang out.
你自己就會有不同的.
得到不同的好處和影響.
有另一個行為的心理學的研究.
發現什麼呢.

$^{801}$這個研究很有趣.
他找一些演員.
臨時演員.
他就去那些自助餐.
拿東西吃的地方.
派他們去.
那些演員就會.
因為他們是臨時演員.
他們就在排隊.
拿那些自助餐吃東西的過程裡面.
他就會拿很多很多.
疊起來疊起來.
研究就是想做什麼呢.
就是想看這些人.
拿很多東西.
對身邊的人有什麼影響.
發現對男對女.
有不同的影響.
隔壁那個不停拿很多東西的時候.
你應該會怎樣呢.
這個行為心理學.
男女的效果是不同的.
男是沒有效果的.
女的很大效果.
因為見到別人拿.
你又拿很多.
吃不完也好.
我都會去拿.
但是重點是.
重點是你和誰一起.
或者你的社群.
對我們有很大的影響.
如果你想生命去到改變.
能夠成功率高一點的話.
你不可以獨立的.
你不可以獨就.
所以為什麼有些人.
給錢去參加一些訓練班.
或者做運動的班等等.
那些賺很多錢.

$^{841}$是一個行業來的.
為什麼.
因為有其他人和你一起做.
對你有影響.
或者參加一些查考聖經的小組.
是不是.
別人查經.
你又去查經.
就對你有影響.
所以為什麼有些人.
要去圖書館讀書.
自己家裡讀書不好嗎.
當你去到圖書館的時候.
就有人影響你.
你就自然而然會努力一點.
所以不可以獨立.
到了最後的時候.
我想去到建議大家.
2024年你考慮一下.
一些應有的目標.
我可以分三類.
Spiritual.
Emotional.
和Physical.
Spiritual的有什麼.
學習一下去看聖經.
你留意我說的.
不是看多聖經,少聖經.
這麼簡單.
學習一下怎樣去看聖經.
How to read.
掌握一下怎樣去看好半聖經.
祈禱會或者去學習.
建立我們祈禱的生活.
很多時候我們想有.
豐富的祈禱生活.
是一個人做不到的.
所以為什麼會有一班人一起做.
別人影響你.
你也去影響別人.

$^{881}$敬拜的生活.
是嗎.
都是需要一同一起做.
我想幾年疫情.
已經打散了我們對這些的體會.
亂了.
敬拜一定是一切都可以做.
奉獻的生活.
慷慨的生活.
剛才我說了.
如果你還沒有實行奉獻的生活.
你不妨開始小小的改變.
你會發現到年底的時候.
你有三十多倍的修營.
但如果你一點都不開始.
並且你只是定一個很高的目標.
你就一定不會成功.
Tiny changes.
一些很小很小的改變.
屬靈的健康做什麼.
可能新的一年你看一本書.
看一本屬靈的書.
今時今日的人都不看書了.
這個真的很可惜.
去買一本好的屬靈的書去看.
如果你需要推薦的話.
我可以有一個清單給你.
不同程度的.
這個是屬靈的健康.
第二個是情緒的健康.
我漸漸今年開始了解情緒健康.
真的很重要.
今年為你情緒的健康.
去定下一些目標.
怎樣去管理你的情緒.
你的脾氣.
你的心境.
這些都需要一些tiny changes.
第三就是身體的健康.
我不管你減肥.

$^{921}$做運動也好.
什麼都好.
定一些特定的目標.
讓你的身體能夠壯健起來.
你有沒有留意這三件事.
其實是大家有連結的.
你身體不好.
你很難屬靈上有健康.
日日這樣污糟.
睡眠質素不好.
身體又差.
你的靈性怎可能提升起來.
這三件事其實都是有關的.
求主幫助我們.
我們希望最後一個祈禱.
今年最後一個祈禱.
就是我們來到上帝的面前.
今年要結束.
求主幫助我擺脫一些壞習慣.
一些霸佔我生命空間的東西.
讓我的生命騰出更多空間給上帝.
我同心協力地祈禱.
最後一分鐘請大家花一點時間來到上帝的面前.
我給你一點時間去禱告.
上帝你是從歲守到年終.
都去帶領我們的神.
雖然我們人生裡面有起起伏伏.
生命裡面可能有迂迴曲折的事.
但我們知道你是掌權.
在我們生命裡面帶領著.
愛護著.
去供應著我們的上帝.
來到這個2023年最後的一天.
主我們仰望你在新的一年裡面.
求你幫助我們.
去擺脫一些霸佔我們生命空間太多的東西.
一些壞的習慣.
一些壞的思維.
主我們求你將這些東西洗淨除去.
幫助我們生命裡面更多的空間.

$^{961}$是預備為你的.
願上帝我們的主.
在我們生命裡面.
生活裡面.
繼續地去帶領我們.
在靈性的健康.
我們和你之間的關係.
上主我們求你幫助我們.
和你建立一個更活潑的關係.
無論我們是讀聖經.
看你自己的話語.
聽你自己的說話.
去向你禱告.
去敬拜你.
去思考我們的信仰.
主我們求你給我們靈性上健康.
更健康.
在我們情緒的管理上.
上帝我們都來到你面前.
求你幫助我們.
幫助我們去在我們眾多反反覆覆.
喜怒哀樂的情緒裡面.
來到你面前.
去禱告.
去整理.
主我們的身體.
我們都教導.
求你保守.
讓我們能夠有壯健.
有力量.
繼續地為你去事奉.
願上帝在新的一年裡面.
讓我們見到更多.
三十倍.
五十倍.
一百倍.
我們見到修成.
讓我們新心靈.
能夠成為一個土地喜悅的人.
我們祈禱奉著名.

$^{1001}$願上帝祝福各位,下一年見!.
(觀眾笑聲).
\newpage



\section{}
\label{sec:YZjpAlGepXw}
\textbf{Atomic Wisdom 2 - Rev. Mark Chan}
\newline
\newline
連結: \href{https://youtube.com/watch?v=YZjpAlGepXw}{\texttt{https://youtube.com/watch?v=YZjpAlGepXw}} ~~~~ 語音日期: 2024-01-09
\newline
\newline
\hyperref[sec:8jU_KxWNams]{\small{< < < PREV SERMON < < <}}
~
\hyperref[sec:index_chronic]{\small{[返順時目]}}
~
\hyperref[sec:index_scriptual]{\small{[返順卷目]}}
~
\hyperref[sec:C09znJHrIGU]{\small{> > > NEXT SERMON > > >}}
\newline
\newline
$^{1}$Hey guys, I know a lot of people are sick in the last week, right?.
You sick too? Yeah. And I've been sick..
And since last Sunday evening,.
after all this talking and meeting and I lost my voice for three straight.
days, three days. Yeah. I completely lost my voice..
This is the first time happened in my life. I could not speak, you know,.
you know, this is how,.
how devastating it is for me to not to be able to speak a word,.
right? Especially because we know we are pastors. We, uh,.
we live by what we say, right? And we have to say,.
we have to say things that we have to speak. Right..
And what is most frustrating for me is actually at home when I'm trying to.
command my children to do something good for the family. Right..
I said, Nathan, can you do the dishes? Because of my crack voice,.
you know, I couldn't speak up. Right. And then how they deal with me,.
it was that, what can you say that again?.
And it was very, very frustrating. And, uh, on Wednesday, last Wednesday,.
I slowly get my voice back. Um, so, you know,.
as you can tell my voice has not fully recovered. Um, but, um,.
because, because of my voice, it's not working so well,.
I will make today's sermon a little bit shorter. Praise God, right?.
Happy new year. So anyways, uh,.
last Sunday we began the series on atomic wisdom..
And this is part two. As I said, atomic does not mean destruction..
Atomic means tiny, insignificant, but as we, uh,.
do life change,.
we need to understand the dynamic of making life change successful,.
making life change sustainable. We want to achieve our goals in the new year,.
whatever new year resolutions you may have..
We need to understand on what conditions that we can improve ourselves..
Whether you are doing your schoolwork or you're doing your career or your.
family plan or financial plan or spiritual plan, whatever that is,.
we have to understand what conditions are necessary to make life change,.
uh, sustainable. So that is what, uh, this is all about. Uh,.
we trying to integrate behavioral psychology insights from a book like.
atomic habits by James clear on the left, or, uh,.
we integrate it with, uh, the book of prophets,.
or especially a chunk of writings in the Bible called wisdom literature..
And as you can tell, if you find wisdom literature in the Bible,.
usually is in the middle of the Bible,.

$^{41}$some of the Psalms and the book of prophets,.
they are full of wisdom and insights of how to live a skillful living,.
how to live a life that will honor to God..
And there's one characteristics about, uh, wisdom literature. Um,.
it is bottom up literature, not top down..
As I told you, the prophetic literature,.
what the prophet says or the Torah, the law, this all top down, right?.
God reveal himself..
God tell you what his demands are and what his principles are from top down.
revelation from top, top down and wisdom literature is bottom up..
So the wise teacher will observe the creation, the creatures, the human,.
humanity,.
how human nature is and trying to draw lessons and revelation of God from.
bottom up from the creation..
So I think this kind of literature will help us or inspire us to open our eyes.
to observe how the order of the cosmos,.
how the order of God's creation is and learn wisdom here and there..
Like this one, 25 from the book of prophets, it says the perfect,.
the purposes of a man's heart are deep waters,.
deep waters. So you look into yourself,.
you find you're doing some thinking all the time and your thought life,.
whatever you're thinking of the purposes of the thinking of a man's heart,.
unlike deep water, deep water..
How many of you do fishing from time to time? Right? Like deep water fishing..
No shaking your head. Boring, right? Like deep water fishing..
You know,.
you gotta have your bait down deep in the water and you can wait and once you.
get your fish, you gotta pull it out, right? And patiently pull it out..
This is what wisdom is all about,.
especially wisdom and insight into human nature..
It's like man's heart are like deep water..
Even we sometimes don't understand what we are thinking about,.
but a man of understanding, a wise teacher will draw them out..
Use patience, use principles to observe and to understand human nature..
Then you understand how we work, understand the general human nature..
That is wisdom literature..
And European journal did a study of 96 people some years ago and find that if a.
person wants to form a new habit, say you have new year resolution,.
say you don't read your Bible every day,.
you want to have a new habit to read the Bible every single day..

$^{81}$How many days would it take for a person to form that new habit?.
It finds that the shortest time for someone to develop a new habit is 18 days..
18 days doable, right? 18 days..
The longest time for some people,.
the lazy people is 254 days..
That's still doable, right? Less than a year. And you can form a new habit..
How many of you floss every single day? Floss your teeth..
Only a handful. How disappointing, right?.
But good news to the dentist, right?.
How many of you go to the dentist like once every six months, once every,.
every year? You should go. You should go. Okay..
Flossing is a new habit, right?.
And how many days would you take to form this new habit?.
Shortest is 18 days. Longest is 254 days..
And in average is 66 days, 66 days..
So those of you don't floss your teeth, start flossing today..
For 66 days in average, you will have this new habit. Hallelujah..
Right? How cool is that? Right?.
So you have to plan ahead and you have to understand how we work, right?.
So let's reveal what we have talked about last week is the power of 1\%..
If you make 1\% change every single day for 365 days,.
you will be like that curve..
You will accumulate the result and the effect or the impact of your life change.
for 37 times or even better..
Once you have the patience to, to,.
to grind through the valley of disappointment, right?.
However, if you don't do any improvement,.
if you are 1\% decline for 365 days, bad news for you,.
you have no result at all. Right?.
So no matter what solutions you have in mind, keep that in mind..
1\% only for every single day, 365 days,.
you will see impactful result..
That's principle number one, power of 1\%..
And then the power of incentives or disincentives..
Once you have incentives in your goal, in your plan,.
then the possibility of achieving your goal will be higher..
And the specifics are better. Remember I said,.
if you have a new goal in mind, a new year resolution,.
make sure you specify the when and where and how..
The more specific you are, the better you will achieve your goal..

$^{121}$Number four, community, the people you hang out with,.
will make influence on you, will make influence on you..
So you have to understand how it works..
A lot of this are actually common sense, right?.
They're all common sense..
And of course behavioral psychologists will solidify our understanding how we.
work. So we'll continue today with principle number five,.
is the power of default, power of default..
Now what does it mean? Default..
Default means something preselected for you,.
preselected for you. Now a lot of you use computer, right?.
Right. How many of you will browse the internet? When you browse the internet,.
you have to choose which browser to choose. No, right?.
How many of you use Chrome? Chrome browser? Edge?.
Fox? Firefox? No? I mean,.
you browse the internet on your phone or your computer..
There's always default browser, right? I use edge..
And every time I open edge, sometimes the right-hand corner will ask me,.
do you want to make Chrome the default? I said, no,.
don't bother me. Right?.
That is the power of default because you don't have to choose because it's.
already preselected for you. That's how live change works, right?.
Recently, my wife just got a new phone..
Actually it's an old phone from Rachel because Rachel got a new phone and I give.
the used phone to mom as new as iPhone..
She has been using Android for years and she just.
got switched to iPhone..
So my whole family, everyone is using iPhone except me now..
I'm the only Android user and it took her days to get used to the iPhone..
Yesterday, last night,.
she was still struggling with the iPhone operation.
system. Right? I said, good luck to you. Right?.
And it takes time to make the default happen. Right?.
For me, Android is the default. For some of you, iPhone is the default,.
but you have to understand the power of change for default. Right?.
So psychologist tells us usually when, when,.
when there's something new happen in the calendar and you make some new year.
resolution, some new goals,.
the chances of succeeding is better..
Say you have some new life change,.

$^{161}$a fresh start strategy, like new year..
Usually people make new goals or plan for new.
endeavors when,.
when the new year hits,.
because this is first strategy and the behavioral scientists are telling us.
people would like to tie new goal at some turning point in their lives..
So this month is the best day or best time to make new year resolution. Right?.
Right. You don't want to make some new goals on September 19th. Right?.
I mean, what does that have to do with you? Right? September 19th,.
by the way, is our birthday. Right? It has nothing to do with you. Right?.
But new year, everyone's new year..
Everyone is going through new year and first start strategy usually work..
22 six from the book of Proverbs..
It says to train a child in the way he should go..
When he's old, he will not turn from it..
Train, discipline or dedicate,.
begin or inaugurating the life of a child when he was born or she was born..
You want to set the default for her..
She'll come to Sunday school. She will come to learn the faith. Right?.
By the way, this is not brainwashing. Right? It's teaching..
We all want to teach the best thing for our child and our faith and our values.
are the most important things that we want to teach our child..
And the Proverbs sets train your child or dedicate him in the way he should go.
the rest of his life. Then in general, when he is getting older,.
he will not depart from it because the beginnings are critical..
The default is critical..
I'm talking about the power of default. One of the scholars, um, you can,.
you know, do a Google search is James Choi. James Choi..
James Choi is a professor at Yale school of business..
It's very interesting, right? Uh, Choi is a professor as Yale school..
Yale is one of the best schools in the world..
And, uh, Choi is a solid question..
When you Google on YouTube talking about James Choi, uh,.
he will be so point blank telling people about his faith,.
right? And news channel interview him. He always say, I am a Christian..
So choice is solid question. And he has done much research on 401k,.
401k. How many of you have 401k or some sort of retirement plan? Right?.
You do, right?.
And James Choi is one of the scholars who has done so much research on a 401k..
And he found that, uh,.

$^{201}$the ideal saving rate for today's Americans is.
14\% of your income, 14\% of your income..
Um, you invest 40\% of your income into 401k..
Then, uh, you have a really solid retirement, uh,.
arrangement for yourself and his study. Uh,.
one of his studies try to follow up on people, uh,.
employees signing up for 401k..
Now the first option of the study is that the first group of people to have the.
employees or new employees, when they, uh,.
go to work first day to give them the choice to opt in,.
to sign up for 401k and give them 30 days to do.
so. So when you're new to the company, when you're a new employee,.
the boss will tell you, the HR will tell you, okay,.
you have 30 days to sign the full 401k plan. All right..
And this assess rate for people opting in the.
401k plan in 30 days is 75\%..
Not bad, right? Not bad. However, if you reverse it,.
but the second group,.
the second group of employees is that when they come to work the first day,.
you already signed them up for 401k..
So it's automatic enrollment and then you give them 30 days to.
opt out. All right. So first day of work,.
HR telling you, you are signed up for 401k already. All right..
You have 30 days to opt out. You can sign off from the plan..
And the end,.
we saw the success rate for people signing up for 401k with.
automatic enrollment is 100\%..
100\%. Why?.
Because the power of default, right?.
You presuppose something in your life,.
some sort of habit like this day of the week,.
I'm going to come to worship. This time of day,.
I'm going to read the Bible. That is the default position..
And you have the default presuppose and it will have 1\% chance,.
1\% effect on you and 365 days you have 30 plus life effect change in your life..
Every single morning after I had my,.
when I have my breakfast during my breakfast time,.
I would turn to the phone..
I'll go to YouTube because I subscribe to a channel..
It's a devotional channel and it will take me a reading one verse from the Old.

$^{241}$Testament and one verse from the New Testament and each video take one verse is.
only three to five minute long. So during the breakfast,.
I will turn to the YouTube channel and I will watch that video and the teacher.
will explain to me what this version is all about. Right?.
Is that easy? It's easy, right? But that's my default position..
So I've been doing it for many, many, many, many,.
many times and for a very, very long time..
And that is the power of default, right?.
So think about it when you plan your new year,.
when you plan your life goals are important,.
but you have to understand the power and the effect of default presupposing in.
your life, then it will have life, permanent life change on you..
Number six, make it personal..
The more personalized your goals are,.
the interventions are, the more successful the outcome..
The more personal it is, you make it more personal,.
then it will make it more sustainable..
And that's why people hire personal trainer in the gym, right?.
How many of you have personal coach to help you do exercise?.
No, huh? I know, I know, I know you guys don't do that. Yeah..
But you know, because you have, people have personal trainers,.
so they've, they will feel accountable, right?.
It's a person to person relationship..
Proverbs 27 says,.
like a bird that strays from his nest is a person who strays from his home..
So talking about people who left their home, right?.
Like many of you left your homeland and you're trying to take root here in the.
foreign land. You are like a bird, right? You're like a bird..
You're flying here and there. You left your nest..
Then we'll skip to verse 10. It says,.
but do not forsake your friend and the friend of your father and do not go to.
your brother's house when disaster strikes you. Why?.
Why don't you go to your brother's house when you need help?.
Because here's the wisdom,.
better a neighbor nearby than a brother far away because you already move away.
from your homeland, right? Whoever have blood relationship with you,.
your dad or your mom, your siblings. Yeah, they, they, they,.
they are close right in terms of blood relationship..
But since you move away from them, you are far away from them..
And who is better when this, when you need help?.

$^{281}$A neighbor without blood relationship.
nearby would be more valuable, right?.
That's what I find church family valuable, right?.
Sometimes we bump into each other and we learn each other and Oh,.
you're a Christian too. I'm a Christian too..
It's for some unexplainable reason, some mysterious connection, right?.
Because our faith, our core values are the same,.
even though we are not blood related, right? And we feel close, right?.
That's why we have fellowship and that's why we foster the relationship and our.
faith community so that we get this wisdom..
What is it? Better is a neighbor nearby than a brother,.
a blood related brother far away, right?.
That's what church is for, right? When you need help,.
people around you and nearby who can help you more so than your siblings,.
right? Amen to that. Sometimes it's like that, right?.
I heard a story about a PhD student. The PhD student, you know,.
you need to graduate from a PhD program. You have to write dissertation, right?.
And this station runs around 200 to 300 pages long and you have to write a lot.
of chapters and you have to write those chapters by a certain date, right?.
So the story I heard is that this PhD student,.
because he has to write a certain chapter, like chapter one,.
by a certain date, chapter two, and then by a further day,.
and then chapter three and so on, right? In order to motivate himself,.
he writes a series of checks, post dated,.
date them accordingly and give the stack of checks to his friends and saying,.
you know, if I don't write a chapter by a certain day,.
then you can cast a check. Right? Sounds like a good strategy, right?.
And make it really, really personal, right?.
Because not only the checks are with the friends,.
but the checks are your money, right? If you don't,.
if he doesn't write a certain chapter by a certain deadline,.
then he will lose the money, right? But if he does so,.
then he will get the check back, right?.
Make it personal. As you know, I'm in a PhD program..
I don't think I would do that. Yeah..
But think about it, you know,.
spiritual goals or some sort of new year resolution,.
Bible reading plan. You ask a friend to keep you accountable, right?.
You ask a neighbor, "Hey, check on me." Right?.
Did I keep the reading going or not?.

$^{321}$Well, as I said, my sermon will be shorter..
So I would like to do a short conclusion today..
And coming back to here,.
what I shared to you last week is that there are three big areas we have to.
think about our life. When we manage our life,.
God has given us, think of this three areas of growth,.
spiritual growth, spiritual health, and emotional health,.
and physical health. Set some goals or have some ideas, specific ideas,.
how you can live a life of wisdom,.
manage your life better..
And today I would like to emphasize on emotional health..
And I know to many of you, perhaps emotional health is a new language..
A new aspect. Right?.
But emotional health is very, very, very important. Very important..
And you know, my generation or my previous generation,.
would like to sweep things under the rug. You know,.
that's how we deal with emotions. Right? You feel sad..
You tear up, you know, that's not a man's thing. Right?.
That's what my dad taught me. Right? Sweep it under the rug. Right? Yeah..
You grind your teeth and then you move on. Right?.
And who cares about your emotional health? Right? But you know,.
in this generation, I think we need to learn this language..
We need to learn how to live with emotional health..
If you read some of the Psalms in the Bible,.
the Psalms are prayers of the people of God,.
to God, people of God, to the Lord..
And they write Psalms and their prayers and that become the book of.
Psalms. Right?.
Many of these prayers are not rosy and happy.
prayers. Right? Yeah..
There are praise Psalms and Thanksgiving Psalms and all that. Right?.
But many of the Psalms are written.
from some sort of negative emotions..
The Psalms you find,.
they are in fear and they are struggling with anxiety..
They are struggling with fear and depression and.
sadness. Sometimes they cry out to the Lord. Right?.
When they wake up and prosper, their hearts are frustrated. Right?.
But they write the prayers out and they go to God. Right?.
So the Psalms are full of laments,.

$^{361}$will and honest prayers and they grieve before God..
And the people of God did not hold back their true feelings toward God..
And bow means because God is all knowing. Right?.
God knows how you feel. Right? You can't hide it..
You can't sweep under the rug, but it's there. Right?.
But you need to understand God is big enough to handle your complaints,.
your griefs, and your laments..
And we can draw close to God with deeper intimacy..
I find that sometimes Christians always like to over-spiritualize things..
There is a certain tendency,.
especially conservative Christians. We like to spiritualize things..
You know, when you feel sad and we say, "Oh, we need to give thanks." Right?.
Because give thanks every day and be joyful every day. Right?.
Are those true? Yeah, those are true. Right? In a certain sense..
But you know, when we feel sad, sometimes we sweep it under the rug and say,.
you know, if we are living faithfully, we have faith in God,.
we should not feel sad. Right?.
And that is tragic because that will suppress our emotions and sweep it under.
the rug and it will do damage to ourselves. You know,.
recently you may have heard some news stories about, you know,.
some celebrities, right, commit suicide..
I heard one from someone in the morning service. He said,.
there's a Korean actor or actress, something like that,.
celebrity who committed suicide. And days before,.
nobody knew that he was going through something negatively..
There was no signs. Right?.
Because you think about physical symptoms are visible. Right? Like my voice is,.
you know, you can hear it, right?.
It's visible and you know I'm sick. Right?.
You have runny nose, you're sick. There's a physical, visible,.
you know, bodily symptoms. But emotional,.
usually has no visible signs. Right?.
People may sound positive, but you know,.
in the middle of the night you're struggling with some griefs or some dark.
emotions. And what do you do about it? Right? Instead of,.
instead of suppressing it, you know,.
the Psalm gives us the channel to go to God in prayers..
And God can certainly beg enough to hear your cries..
And it's okay to go to heavenly daddy to say, you know, I,.
I feel sad and I struggle. Right? I complain..

$^{401}$God can certainly handle that..
Many years ago, I think almost 20 years ago,.
I was in a state of some sort of fatigue, you know,.
very, very exhausted, you know, emotionally and spiritually..
And I was trying to figure out what is going on in my life. You know,.
from the outside, nobody can, nobody can tell. Right..
Even my wife cannot tell. Right..
My wife cannot tell because I take care of her. I do my best. Right..
Nobody knows. You know, I feel really,.
really restless and exhausted. So I got some books and, you know,.
I know where the resources are. So I got this book..
This book is called, um,.
"Unmasking Male Depression" written by a Christian psychologist,.
Archibald Hart..
Now Hart has been a professor of Christian psychology in Fuller Seminary in.
California for many, many years from 1970,.
early 1970s to, to the 2000s. Um, he ended his career,.
um, uh, teaching at Fuller until 2016..
And he just passed away at Be With The Lord in 2021 during the.
pandemic. And he wrote this book. And by the way, this is not,.
this may not be one of the best books on market,.
but this is a very, very helpful book. So I picked up this book..
I didn't know what this book is all about. Right..
It talks about male depression. Hey, I'm male..
I feel depressed. Right. And by the way, there's,.
there's also another book he has written with another psychologist,.
um, on female depression. So I picked up this book after the first chapters,.
I could not put it down..
I could not put it down because it's like a mirror and.
telling me what kind of depression I am going through,.
what kind of emotional state I'm going through and the signs and symptoms and all.
that, and how to deal with that,.
how to allow me to have room to deal with the dark emotions.
inside my mind, right inside my feeling..
So I find this book is a life-saving book for me. Now I struggle for it..
I struggled for probably two plus years. Yeah. Two plus years serving in the church..
And, uh, you know, my plate is full and I'm full of activities. Right..
But in the middle of the night, I'm just find myself, you know,.
struggling and exhausted..
And this book helped me to clarify a lot of issues I was going through..

$^{441}$Now I give you this example because I want to tell you, we have to,.
um, have the wisdom to manage our emotions..
We do. We do. Okay. Don't learn from the previous generation..
Right. They're doing the wrong way. You know, they deny it,.
but what is the way out? Right..
How do we manage emotions correctly, properly and wisely?.
Right. So these are some of the resources that you should have..
And in the morning service, I also mentioned, uh, probably there was a necessity.
for us to go to see a counselor. You know, in Chinatown,.
we have a wonderful growing ministry of counseling..
Um, I don't know what the English is, but, um, they are trilingual..
So English, Cantonese, and Mandarin,.
and you can call them and set up an appointment and you can see.
a counselor, a professional counselor, uh,.
well-equipped to help you process, you know, what is going on in your life..
And usually it's about what four to six sessions..
Each session is about 50 minutes..
You talk to a skilled counselor and she or he will be able to draw.
things out of you to help you manage your emotions better..
I believe there's no healthy question unless we are emotionally healthy,.
unless we're physically healthy, unless we're spiritually healthy..
And those three things, or even more mental and all the,.
all the other aspects of our lives are intimately integrated together..
And we need to make sure that we are healthy. Right. So, um,.
if, if that rings a bell, you know, you can come to me and, um, you know,.
I can give you a brochure and give you the information. You know,.
one of my good friends, uh, Dr. Jane Mann,.
Jane Mann is one of the skilled counselors, uh,.
working there and serving the community. And definitely you should try it. Right..
Um, actually they charge, but it doesn't charge that much. You know,.
I talked to Jane, I said, how do you earn money?.
How do you make it sustainable? Right..
But the goal of the ministry is to serve the community,.
to serve people like you and me. Right..
And it's necessary. It's necessary. Right. Um,.
I don't know what happened to you in the past year..
Maybe the past year has been a really rough year for you. Uh,.
you really struggling. Sometimes we just strip things under the rug..
Maybe you need to see some of the resources that's out there..
Maybe you need to see a professional counselor. And by the way,.

$^{481}$don't have that misconception that, Oh,.
I'm going to see a counselor. That means I'm insane..
No, you can see a counselor and enrich yourself. Right..
You can enrich yourself by talking to a well-trained.
counselor. Right. Everybody should see a counselor. Right..
We all need another wise person in front of us to.
draw things out of us so that we can live, um,.
healthily. Right..
We can live a life that would bring honor to God..
All right. So we understand all this principles and let's go to God.
and pray that we will continue to live a.
faithful living and, um, uh, live a life of wisdom,.
honoring to God. And then we're going to partake the communion together..
Let's all pray. Father, thank you..
Thank you for all the bad and good experiences that we have..
Curses and blessing you give and take away..
And you are the Lord who is sovereign over our lives and you're in control of.
all things. Father, once again, we come before you as the first Sunday of.
2024, we commit ourselves into your hand..
And knowing that you are the God who is kind,.
who is gracious, who loves us so much..
Thank you for all that you have done through your son,.
Jesus Christ on the cross. You redeem us holistically..
You'd only not only save our soul, but you save our body..
You grant us holistic redemption..
And I pray that our lives will rest in your hand and we'll find the true.
peace..
Help us God to live in a way that will honor you..
As we launch into a new year,.
we don't know what you have in store for us, but we know your character..
We know you are a good God. So once again,.
we pray that you will hear our cries, hear our struggle,.
help us to grow spiritually, physically, and emotionally, mentally,.
healthy in such a way that we will draw closer and closer to you..
Thank you for being there and walking alongside with us,.
even in our dark times and also in our good times..
Thank you for being our faithful God..
As we enjoy the communion together, take the bread and the cup..
Help us be reminded how much you love us. In Christ name we pray..
\newpage



\section{}
\label{sec:C09znJHrIGU}
\textbf{原子智慧 2 - 陳崇基牧師}
\newline
\newline
連結: \href{https://youtube.com/watch?v=C09znJHrIGU}{\texttt{https://youtube.com/watch?v=C09znJHrIGU}} ~~~~ 語音日期: 2024-01-09
\newline
\newline
\hyperref[sec:YZjpAlGepXw]{\small{< < < PREV SERMON < < <}}
~
\hyperref[sec:index_chronic]{\small{[返順時目]}}
~
\hyperref[sec:index_scriptual]{\small{[返順卷目]}}
~
\hyperref[sec:RlCj0uSBVME]{\small{> > > NEXT SERMON > > >}}
\newline
\newline
$^{1}$看完節目.
我生病了.
上星期講完道之後.
三天都不能發聲.
我覺得最激氣的是.
我家中的小朋友.
我叫他們幫忙收拾東西.
他們就問我「你說甚麼」.
真是豈有此理.
結果三天都不能發聲.
你聽到我的聲音都還未回來.
所以剛才說開課.
我都未夠膽說.
因為我從來未試過三天都不能發聲.
發不了聲很慘.
我們靠嘴巴吃飯.
三天我靜修.
所以我都不知道會否好.
不知道何時好.
所以我想給自己一點空間.
聲音可以恢復.
所以開課日期和時間和資料.
我們都未決定.
希望今天講完道之後.
叫團會後回去休息一下.
希望小朋友乖乖地.
等我恢復之後.
資料就可以出來.
所以我打算今天講道短一點.
新年快樂.
好 短一點.
盡量吧.
我們一直說原子智慧.
Atomic Wisdom.
這是一種結合.
我們從一些行為心理學.
Behavioral Psychology.
給我們一些人性的Insight.
我們能夠結合聖經裡面.
Proverbs或者智慧文學裡面所講的東西.

$^{41}$幫我們搞清楚.
首先我們新一年的時候.
我們當然會有一點點目標.
有一點點新年的立志等等.
但未立志之前.
或者立志了.
我們要明白一個生命的改變.
一個生命的向前.
到底有什麼動態在裡面.
這本書叫做Atomic Habits.
是James Clear寫的.
裡面講了很多人性心理學.
有關的Insight.
一些觀察.
再結合聖經裡面所講的.
智慧文學 箴言等等.
其實箴言或者智慧文學.
Wisdom Literature.
都是在做這件事.
它是由下至上去領略上帝的啟示.
什麼叫做下至上呢.
就是說觀察人性.
觀察一下自己.
觀察這個世界.
以致我們能夠知道.
人生是怎樣Work out的.
在裡面找到一些屬靈的真理和啟示.
譬如在20章第5節.
箴言在聖經裡面中間.
其中一本書叫做箴言Proverbs.
這裡講到人心裡的籌算如同深水.
唯聽聰明人才能將它吸收出來.
什麼意思呢.
其實是一種Observation.
所以看箴言.
聖經裡面一些智慧文學.
其實是在講智者.
留意到人性是怎樣的.
人心裡面所想的東西.
好像深水一樣.

$^{81}$深水,你們有沒有去釣魚.
釣過深水魚.
裡面不知道是什麼.
下面不知道是什麼.
人心就好像這樣.
連我們自己想的東西都不是很清楚.
唯獨是有智慧的人.
有上帝智慧的人.
慢慢可以將這些東西.
Draw them out.
好像井水打水一樣.
慢慢將它拉上來.
以致我們能夠明白人性是怎樣的.
所以新的一年納節的時候.
請你未納節之前.
想清楚.
有什麼我們需要去建立.
一些所謂原子習慣.
一些很細的習慣.
這些習慣能夠幫助我們.
有更加長遠的生命的改變.
在歐洲有一個研究說到.
一個人養成一個新的習慣.
最短.
調查裡面最短.
一個人能夠去建立一個新的習慣.
就是需要18天.
18天.
譬如你不用牙線.
有誰用牙線的.
很乖.
對牙醫來說可能不是很好.
因為牙醫想見你.
你養成一個習慣.
牙線才算.
最短是要18天.
即是你要做18天這個新的習慣.
你才能夠養成.
最長的.
那些不是很堅持的人.

$^{121}$最長的就需要254天.
都不是很長.
一年都不用.
平均來說是66天.
你要做那件事66天.
你才能夠去養成一個新的習慣.
一個長遠的習慣.
所以想一想.
我們生命裡面想改變的時候.
有些事我們需要作一些改善.
養成一些新的習慣.
我們需要計劃得好一點.
起碼66天.
你堅持到66天.
你就能夠有大一點的機會.
可以養成一個新的習慣.
看回我們上個星期的一些評價.
就是1\%的力量.
什麼意思呢.
你看看這個改變的表.
一直向上曲線.
那條黑色的線.
就是我們一般人.
我們能夠做到生命改變的情況.
那個過程是怎樣呢.
由開始的時候.
是沒有什麼改變.
沒有什麼效果.
開頭的時候.
結果不是那麼明顯.
你過了那個所謂失望谷.
然後那個回收.
就會繼續上去.
但是你只需要做1\%的事.
很少很少原子的新的習慣.
以致你慢慢慢慢.
你生命裡面就會不斷有影響.
直到你365天.
你所收到的效果.
就是37倍或以上.

$^{161}$這個是1\%的力量.
這是第一個原則.
第二個原則就是賞罰的效果.
如果我們去到做一個新的習慣的時候.
我們給我們自己有勵忍.
或者有不勸忍.
我們可以養成一個新的習慣.
機會就會大很多.
第三個原則.
剛才所講的.
上個星期所講的.
就是具體性.
越具體性就越好.
我們訂立一個新的習慣的時候.
要有具體.
時間,地點.
when and where.
想清楚.
譬如你想讀經才算.
讀聖經.
每天讀聖經.
想清楚時間和地點.
是不是要多一杯咖啡.
是不是.
那你養成一個新的習慣.
機會率就會很大.
還有你的生命能夠有更多的效果.
第四.
我們上一次講到的.
是社群的影響力.
身邊的人給我們的影響.
或者團契.
身邊的弟兄姊妹.
朋友給我們的影響是很大的.
所以一個人獨自是很難做到的.
有一個新的習慣.
是需要有其他人來支援我們.
這次我們繼續看下去.
其他的原則.
再看第五個原則.

$^{201}$第五個原則是什麼呢.
The power of default.
預設的力量.
什麼叫做default.
大家用電腦的就知道.
那些browser經常彈出來.
用不用make it default.
chrome經常問你.
做default好不好.
或者microsoft edge的瀏覽器.
就想你將我變成default.
就是什麼.
每一次你上網看東西的時候.
你就用那個browser.
不用其他browser.
那個就是預設的選項.
就是說你不需要再選擇.
我們很少會上網看東西的時候.
我們就會先選擇你用哪個browser.
你用firefox.
microsoft edge.
還是google chrome.
因為你已經選擇了default.
預設的能力.
在我們生命的改變裡面都是一樣.
什麼叫做預設能力.
就是說人的行為心理的時候.
凡是真的有fresh start的時候.
時間裡面有fresh start.
或者日曆裡面有fresh start.
那個就成為一個改變的契機.
譬如新的一年.
就是新的立字好很多.
比你在9月19號去訂立.
我要新的一年有什麼立字.
你會不會.
當然不會.
9月19號派到威嚇我生日.
你通常新的一年1月1號.
或者新一年的時候.

$^{241}$我們就會想.
有什麼新的計劃我們要進行呢.
所以在那個fresh start裡面.
birthday或者一些new的life changer.
通常是一些turning point.
在我們時間生命的旅程裡面.
一些的turning point.
在那裡做一些改變.
或者做一些開始.
一些新的原子的習慣.
那個機會率是大很多的.
《聖經》22章第6節.
那裡說.
教養孩童當行的路.
就是道路也不偏離.
教養孩子當行的路.
train這個字.
其實就是dedicate的意思.
奉獻的意思.
聖殿會奉獻的.
做了聖殿.
新的聖殿的時候.
首先會有一個奉獻禮.
或者inaugurate.
去到啟始.
去到將他啟動開創.
就是說他帶一個小孩的時候.
這個智慧就是說.
從開始.
fresh start.
你就去教他走當行的路.
一般來說.
到他老的時候.
都不會改變.
就是說那個beginning.
啟始.
fresh start那裡.
是給一個人的生命的改變.
是非常重要的.
在耶魯大學有一個很出色的教授.

$^{281}$叫做James Choi.
James Choi.
你去google一下.
或者做finance的.
應該很清楚這個人.
James Choi是耶魯大學.
一個business school的教授.
James Choi.
是一個很solid的question.
很多時候新聞電台去訪問他.
他第一句就會說.
我是一個基督徒.
所以我說的話.
是從一個基督教的立場來說.
在耶魯大學.
business school的教授.
他專門研究那些401k.
就是那些.
就是那些退休金.
養老金等等.
他就做了很多這樣的研究.
他的研究發現是什麼呢.
他發現.
他曾經做過一個研究.
就是一間公司裡面的員工.
年頭的時候.
你叫他sign up.
或者他剛剛上任的時候.
剛剛開始他的職事的時候.
他就去sign up 401k.
有一組人.
他就叫這班員工.
給你30天.
你去sign up 401k.
sign up退休金的計劃.
給你30天.
結果他的成功率是75\%.
就是說新的員工.
在30天之內.
他會up in.

$^{321}$去sign up一個401k.
但是第二組人就不是了.
第二組人就是他是員工.
他一上任的時候.
一請他的時候.
就已經幫他sign up了.
那就給他30天.
你去up out.
如果你不要401k的話.
你就儘管up out.
給你30天的機會.
結果他發現成功率是100\%.
為什麼.
因為default的power.
你開始已經設定了.
你要做那件事的時候.
你就會容易些改變.
如果開始設定是沒有的時候.
你再加上去做一個抉擇.
那就會難一些.
所以James Choi就做了很多研究.
然後Porsche很多公司.
今日公司都是這樣.
automatic enrollment.
自動的你就是sign up了.
然後那個participation rate.
就會更加高.
參與的效率就會更加高.
這些都是behavioral psychology.
如果你想清楚.
到時候這個星期.
這一天我是做不了一件事.
我是不會改的.
這個是default.
你的生命就會有改變.
好像我這樣.
我每一天都會打開youtube.
去看一個channel.
那個channel就是有一條片.
3至5分鐘的一天.

$^{361}$他就講一節《成競》.
一節舊約.
一節新約.
我每一天.
我就是吃早餐的時候.
我就會看3至5分鐘.
去幫我靈修.
吃東西的時候.
總會祈禱.
感恩祈禱.
短也好長也好.
the power of default.
你想想在生命裡面.
要有改變的時候.
特別你和上帝之間的關係.
如果有一些default已經in place的時候.
那就會容易很多.
另一個principle.
第六個.
就是越個人就越有效.
如果你訂立的目標.
訂立你達成目標的方法.
越個人就是越好的.
什麼意思呢.
就是有人去check你.
身邊有人去keep you accountable.
20世紀第八節.
他說人離開故鄉漂泊.
就好像鳥鳥離窩四處飛翔.
這些就是我們講的沒有腳的鳥.
到處飛.
這個就是知者去觀察鳥鳥.
到處飛.
然後他就領略什麼呢.
可以跳到第十節.
他就說你的朋友.
和你父親的朋友.
你不可以離棄.
就在遭難的日子.
在遭難的時候.

$^{401}$你有困難的時候.
不要上你兄弟的家去.
為什麼.
因為離開了故鄉.
在遭難的日子要做什麼.
相近的鄰舍.
強於遠方的兄弟.
遠方的兄弟可能和你有血緣關係.
但始終都是遠方.
但相近的呢.
縱使他沒有血緣關係.
他都是你一個更親的人.
見著你 陪著你 和你同行.
所以如果你離開了故鄉.
你在這裡沒有腳的鳥.
你就是需要建立一些關係.
那些關係很多時候不是和血統有關的.
和你的血緣無關.
那些人就是你更加相近的鄰舍.
比遠方的兄弟更親近.
這裡講到我們人的承擔.
我們人總是需要一些個人的關係.
其實在教會 團契.
在我們弟兄姊妹 基督徒之間的關係都是這樣.
我覺得很奇妙.
基督徒和基督徒之間.
就是有一種非血統.
但又可以很親的關係.
所以我們核心的價值是一樣的.
是拿不走的.
我們都是相信耶穌基督.
都是篤信主的.
都是行在道裡面的.
我們這班人就是相近的鄰舍.
有很多時候就是強如.
比遠方的兄弟姊妹更加親近.
在南美洲哥倫比亞.
哥倫比亞這個國家.
他們發現很多人.
哥倫比亞這個國家裡面.

$^{441}$很多人民是不喜歡交稅.
這是永恆的真理.
有誰想交稅.
如何解決這個問題呢.
社會學的研究發現.
如何逼那些人交稅.
或者提那些人去交稅.
以免他們遲了交稅.
然後他們要再追他們等等.
於是他們就找了一百個人.
寫信寄信給這一百個人.
寄信給他們.
要交稅了.
結果收效就是.
一百個人裡面八個人.
肯乖乖地去交稅.
然後他們用另一個方法.
就是打電話給那一百個人.
替他們去交稅.
恐嚇他們也好.
罵他們也好.
總之提醒他們.
這個時候一定要交稅.
結果收效就是三十人.
即是一百個人裡面.
本來是由八個人.
接到一封信才會交稅.
變成三十個人會交稅.
這樣有沒有結果.
有一點吧.
但他們用第三個方法.
第三個方法就是敲門.
派那些IRS那些稅務員.
上門敲門.
就叫他們是時候交稅了.
面對面跟他們說.
收效是百分之九十.
我們的tax day都快到了.
如果IRS的staff上門敲門.
這樣就腳都會震.

$^{481}$意思就是說.
如果make it personal的時候.
那個果效就會更加大.
有一個PhD student.
我聽過一個這樣的故事.
一個讀博士學位的人.
他最後要寫一個博士論文.
博士論文要寫二三百頁.
你寫的時候.
其實是一個馬拉松.
很辛苦的.
基本上是在寫一本學術的書.
博士學生要寫論文的時候.
他知道這個挑戰太大.
所以他就決定.
每一個chapter.
他要寫一個chapter.
他要在某一個時間之前寫好.
怎樣逼他自己.
能夠在那個時間之前寫好.
他就拿他個人支票簿.
他就寫一些check.
寫給他朋友.
但你知道check是可以post date的.
譬如現在是一月.
三月三十一日.
他一定要寫好chapter one.
他就寫一張check.
三月三十一日才可以cash.
他就send給他朋友.
就告訴他朋友.
如果三月三十一日.
我都未寫好chapter one的話.
你就可以cash出這張check.
結果他就是在每一個chapter的deadline之前.
他就真的每一個chapter去寫.
為什麼.
你哪裡想便宜你朋友.
當然如果我寫我的博士論文的時候.
我就不會這樣做.

$^{521}$如果你要讀經祈禱.
或者一些屬靈生命裡面的雜運.
如果你make it personal的話.
告訴你的朋友.
你check我.
你check住我.
事半功倍.
所以今天我們結束的時候.
去講一講2024年的目標.
我上一次提議到.
在你的人生裡面.
應該要想起碼這三方面.
你的屬靈的健康.
作為一個基督徒.
我是不是過著一個正常的.
起碼的一個屬靈的生命呢.
第二 想想我情緒的健康.
第三 想想我身體的健康.
今天我想特別強調.
就是情緒的健康.
我們這一代的人.
包括在教會.
我們已經不能再避不去提及這件事.
或者不去學這件事.
情緒這件事.
對大家來說可能是一個新的語言.
特別是上一代.
我們是不講情緒的.
哭什麼哭.
咬緊牙根就照做.
起碼我那一代.
和我最上一代都是這樣.
這一代我們慢慢意識到.
任何的年紀 男女老少.
我們的情緒健康都是很重要的.
你留意一下詩篇裡面所講的情緒.
有很多不同的情緒.
留意這個表我會放在社交群組.
就是講到焦慮的時候.
你看什麼詩篇.

$^{561}$如果將詩篇分類.
按照情緒.
emotion, mood來分類.
你會發現有很多不同的詩篇.
講到恐懼的時候.
你應該要看哪一篇詩篇.
困惑的時候應該要看哪一篇詩篇.
當你有羞愧的時候.
無助的時候.
當你抑鬱的時候.
或者想尋求公義的時候.
感恩讚美也好.
所有情緒的範圍裡面.
你都可以找到一篇詩篇.
所以其實詩人都是這樣.
詩人是很誠實.
和上帝來表達他的情緒.
是不會藏起來的.
反正上帝都是全知的.
上帝都知道你自己.
心裡面在想什麼.
和在過渡什麼.
你也要在上帝面前祈禱.
Just be real and honest.
或者去看這些詩篇.
幫我們怎樣去將我們心裡面的情緒表達出來.
在很多年前.
大概都幾乎是二十年前.
我記得我自己在過一些很困難的時間.
那時候我是很疲勞的.
我是很無力的.
在宿命上或者情緒上.
是很累的.
很辛苦.
我自己都意識到自己有些問題.
當然別人外表看是看不到的.
每個人都說Mark Sir很堅強.
其實我不是很堅強.
我都有情緒.
大概二十年前的時候.

$^{601}$我就是在過一個低谷.
我也不知道怎樣去表達.
總之就是每件事都好像沒有趣味.
很depressed.
外表去work out一些東西是沒有問題.
但總是裡面很不妥.
於是我就買了一些書看.
其中一本就是這本.
這本叫《Unmasking Male Depression》.
這本作者是Aperture Heart.
Aperture Heart在加州的呼拉神學院.
教了很多年基督教心理學.
1973年教到201幾年.
然後他在2021年的時候pass away.
2021年剛剛pass away.
沒有多久.
他寫了一本書叫《Unmasking Male Depression》.
我打開這本書都好好奇.
到底要講些什麼呢.
於是我打開看了兩個chapter.
我就發現原來這些我沒有留意.
什麼叫做male depression.
他也有另一本書是講female depression.
這兩本不一定是最好的書.
但我想告訴你.
我就是高呼低谷的時候.
我就看這本書去了解自己.
和可以找人去傾.
我們每個人都有情緒.
情緒是需要誠實去面對的.
我們需要去誠實面對.
我們不能夠sweep it under the rug.
掃到地上才抵.
我們就不需要理會它.
或者我們很多時候.
算了自己搞定.
其實是不行的.
有很多時候你需要找一些專業輔導.
拜登說找專業輔導.
不等於你是弱者.

$^{641}$我們經常有那種意識.
或者一種印象.
你去找輔導.
就是說你差一點.
不是的.
你找一個專業的輔導員.
去幫你處理一下你裡面的東西.
在基督徒裡面.
我想告訴你.
很多時候我們就是sweep it under the rug.
我們就說什麼.
我們要上上喜樂.
我們要上上感恩.
我們要讚美.
為什麼呢.
因為信心應該是這樣.
當然這個沒有錯.
但不代表你心靈裡面的東西.
裡面一些譬如抑鬱.
或者一些負面的情緒.
就會解決的.
我們把它overly spiritualize.
我們把它過分地去屬靈化的時候.
反而是誤了一件事.
反而我們就是不理這些東西.
然後就等它去爆發.
或者去損毀我們的人生.
最近香港有一位姓劉的新聞工作者.
大家有聽到這個新聞.
本身是一個很正向的人.
我看回他前兩天.
他在他的Instagram IG.
或者網上他的網台.
他都在說做運動.
跑馬拉松.
然後叫人正面一點去思考等等.
誰會想到.
他就自殺.
情緒是需要面對.
我們病就自然有外在的病徵.

$^{681}$沙聲或者咳嗽等等.
但情緒很多時候不是visible.
但它是很真的.
所以去了解一下.
去找一些好的資源.
去幫助你情緒上能夠更加健康.
或者在我們芝加哥附近.
有一個剛剛設立的心理輔導服務.
在葡萄中心.
Dr. Jane Man在那裡.
去看看這些專業的父老.
以至有個人可以幫你處理一下自己裡面的事.
2023年過去了.
可能2023年對你來說都很tough.
可能有很多掙扎,失望.
很多我們需要處理的.
一些所謂的負面情緒.
怎樣都好.
上帝都是要我們有一個情緒的健康.
如果你需要輔導的話.
你儘管找我們教務同工.
我可以給你一些資料.
見一見四個session.
每個session 頗便宜.
便宜到我跟Jane說你怎樣賺錢.
但需要的.
正如我們需要找醫生一樣.
情緒健康的話.
屬靈生命就會健康.
情緒不健康.
你講多少屬靈的話語都好.
都是搞不定.
身體當然需要健康.
情緒也是.
屬靈生命.
這件事都是integrate.
求主幫助我們.
我們同心低流祈禱.
然後領受聖災.
因為你創造我們.

$^{721}$並且我們歲守到年終.
我們的人生都在你自己的大靈裡面.
去照顧我們每一個.
你對我們的慈愛.
對我們的恩典.
無微不至.
主我們就是來到你自己的外宗.
我們能夠享受.
在基督裡面的愛的關係.
主我們為了我們每一位.
為你禱告.
我們生命裡面.
我們身體裡面.
可能有不同的病痛.
上帝求你.
讓我們的身體.
能夠得到你的醫治.
讓我們安息.
主我們對你屬靈的生命.
去聽你自己的話語.
在你面前我們能夠.
維持住一個聽道.
對你性靈敏感的一個態度.
主我們求你賜給我們.
我們自己情緒的健康.
喜怒哀樂.
主這些都是你所創造給我們的.
求你幫助我們.
好好地去管理.
你聽見我們心裡面有很多的苦澀.
有很多的灰心,失望,抑鬱等等.
甚至一些我們未能察覺的聲音.
主求你帶領我們.
去到一個情緒健康的境地.
因為你就是我們的醫治者.
帶給我們全人的醫治.
全人的平安.
對我們的厚愛.
主我們今天來到你面前.
繼續去敬拜你.

$^{761}$和你親近.
去拿起聖餐的餅和杯.
我們更加知道.
裡面所代表的.
具體所代表的.
就是你對我們無微不至的恩典.
感恩贊給你.
祈禱給他.
\newpage



\section{1:15-23}
\label{sec:RlCj0uSBVME}
\textbf{Know Your Hope (Ephesians 1\_15-23) - Sam Johnson}
\newline
\newline
連結: \href{https://youtube.com/watch?v=RlCj0uSBVME}{\texttt{https://youtube.com/watch?v=RlCj0uSBVME}} ~~~~ 語音日期: 2024-01-16
\newline
\newline
\hyperref[sec:C09znJHrIGU]{\small{< < < PREV SERMON < < <}}
~
\hyperref[sec:index_chronic]{\small{[返順時目]}}
~
\hyperref[sec:index_scriptual]{\small{[返順卷目]}}
~
\hyperref[sec:__gIiUwiPLI]{\small{> > > NEXT SERMON > > >}}
\newline
\newline
$^{1}$- Well, good morning..
It's fantastic to be here as always..
Really grateful to be with you all.
and to worship with you this morning..
And appreciated the worship..
Thank you, Jason..
And I'm glad to see all of you..
I know it's not easy to get here on a morning like this.
when it's this cold..
I think, there we go..
I know it's not easy to get here on a morning like this..
When I woke up at my house, it was about negative eight.
or negative nine..
So I know it's an effort just to get out the door.
and to get to church today..
So thank you for being here..
And for those that are watching online,.
hope that everyone is warm and safe..
That's important today..
Make sure that you're warm and safe.
and check on the people in your family,.
in your neighborhood,.
make sure everybody else is warm and safe as well..
Since this is the beginning of the year,.
we're gonna be talking about the hope.
that we have in Christ..
And we're gonna be looking in the book of Ephesians today..
So if you wanna turn there in your Bible.
or in your Bible app,.
we're gonna be in Ephesians chapter one,.
verses 15 to 23,.
talking about the hope that we have in Christ..
And in chapter one of Ephesians,.
this is a letter written by Paul to this church at Ephesus,.
this group of believers..
And he starts talking in the chapter one.
about what God has done on their behalf..
And he ends the chapter with this beautiful prayer for them..
And I think this is really a powerful prayer.
for us to consider at any time,.

$^{41}$but especially as we begin the new year.
and we're thinking about some of the things we might face..
So let's read these verses together..
Ephesians one, beginning in verse 15,.
it says, "For this reason,.
"because I have heard of your faith in the Lord Jesus.
"and your love toward all the saints,.
"I do not cease to give thanks for you,.
"remembering you in my prayers,.
"that the God of our Lord Jesus Christ,.
"the Father of glory, may give you the spirit of wisdom.
"and of revelation in the knowledge of him,.
"having the eyes of your hearts enlightened,.
"that you may know what is the hope.
"to which he has called you..
"What are the riches of his glorious inheritance.
"in the saints?.
"And what is the immeasurable greatness of his power.
"toward us who believe,.
"according to the working of his great might,.
"that he worked in Christ when he raised him from the dead.
"and seated him at his right hand in the heavenly places,.
"far above all rule and authority and power and dominion,.
"and above every name that is named,.
"not only in this age, but also in the one to come..
"And he put all things under his feet.
"and gave him as head over all things to the church,.
"which is his body, the fullness of him.
"who fills all in all.".
God, thank you for your word, thank you for your truth,.
thank you for the hope that we have in you,.
and I pray that our hearts would be.
filled with that this morning as we look at your word,.
and I pray that as we consider this year that is to come.
and what you might have for us,.
the joys as well as the adversity,.
that we would be filled with hope.
knowing that you are at work in our lives,.
and that you have grace that's sufficient.
for anything that we're gonna face this year, God..

$^{81}$In Jesus' name, amen..
I wanna remind us of a story, this is in Mark chapter two,.
the beginning of Mark chapter two,.
there's a man who is paralyzed, he's unable to walk,.
and this man has four friends.
who want him to be healed, want him to be made whole..
And apparently one of them, or maybe all of them,.
we don't know, it's not clear from the story,.
but somebody in the story, one of his friends,.
had heard about this itinerant rabbi, this teacher,.
who was teaching and who was also healing people..
And so these four friends get together,.
and they take this man who's unable to walk,.
and they say, "We're gonna take our friend to see.
"this teacher, this healer.".
And they go and they find out where Jesus is at that time..
They go and he's teaching in this house,.
probably not a real big house..
They get there, they've carried their friend all the way.
to bring him to Jesus..
They get there and the house is so full.
that there's no way they can get in..
And scripture tells us that even the doorway is blocked..
They can't get in at all..
But they're not deterred by that..
They say, "We really want our friend.
"to be able to encounter Jesus and to be able to have.
"this time with Jesus because we believe, we hope,.
"we have faith that Jesus can heal our friend.".
And so they find a way up onto the roof of the house..
They somehow make a hole in the roof..
We don't know if it was a mud house.
and they actually tore tiles off of the roof.
or it was more of a thatched roof and they just opened it up.
but that's not real clear..
But they make a hole in the roof.
and lower their friend down through the roof..
And then Jesus ends up healing their friend..
We're probably familiar with that story..
Maybe you've heard it in church before..

$^{121}$Maybe you heard it in Sunday school.
when you were a little kid..
'Cause it's an amazing story..
But what I want us to think about in that story today.
is those friends and the motivation,.
the hope that they had in order to take their friend.
who couldn't walk and say, "We're gonna get up this morning.
"and we're gonna take this person..
"We're gonna go find this teacher that we've heard about.
"and we're not gonna be deterred by any obstacle.".
And when they got there and there was too many people.
and people said, "There's no way you can't even get in..
"The door's blocked.".
They didn't say, "Well, I guess we'll just go back home.".
They said, "No, we're gonna find a way.".
And however it happened, they end up on the roof.
and they tear this hole in the roof.
and lower their friend down..
And then what happens?.
Their friend is healed and he's able to walk..
He's healed by Jesus..
And that story reminds us of how important hope is.
in our lives, how it propels us forward..
Without hope, we struggle to move forward..
We say, "Well, I don't see why I should go forward..
"What's the point?.
"It's not gonna work, nothing's gonna happen.".
But when we have hope in the work that God is doing,.
it propels us forward..
And that's a great consideration for us.
at the start of a new year..
All of us are probably thinking about some things.
that are in our life during 2024..
Maybe we're reflecting on some things.
that were in our life during 2023..
Hopefully we have some prayers, some desires.
that we're asking God for and we're considering..
And all of us know that this year is gonna hold some things.
that we are fully unaware of.
and unprepared for at this point, right?.

$^{161}$I remember January of 2020, right here in this room,.
having conversations with Pastor Mark.
about what was gonna happen..
At that point, there were no,.
the lockdowns hadn't started in Chicago yet..
It was just, things were just starting here..
Obviously, we had no idea what 2020 was gonna hold.
and the adversity and the challenges.
that it was gonna hold for so many people..
Hopefully this year doesn't have anything like that for us..
But still, there are things ahead that we don't know about..
Some of them are joyful things..
Some of them are really, could be challenging things.
in this year for us..
And we need to have hope in God's work in our lives..
And one of the things that I see in this prayer of Paul.
in Ephesians 1 that's really beautiful.
is he's letting these believers know,.
he's saying, "Hey, the climb ahead of you is steep,.
"but have hope because God has given you.
"all that you need in Christ.".
Now, I wanna just do a quick flyover.
of the book of Ephesians and the things.
that he's talking about..
There's a lot of big issues in this letter..
So in chapter two, he starts talking.
about unity in the church..
That's a big issue..
That's a challenging thing, to stay unified.
as the people of God..
He talks about how God has brought us.
from this spiritual state of alienation.
and being separated from God into reconciliation..
That's huge, that's cosmic, right?.
To go from separation from God into reconciliation..
In chapter three, he talks about the mystery.
of the gospel, how God's brought us.
from spiritual darkness into understanding..
In chapter four, he talks about the new life.
we have in Christ, how we're to put off.

$^{201}$the old self and put on the new self..
Well, again, how do we even comprehend that,.
much less live it out?.
These are big things that he's gonna be putting.
in front of this group of believers..
In chapter five, he talks about how we're supposed.
to walk in love in our family relationships,.
in the practical relationships of our lives..
In chapter six, he talks more about relationships.
and spiritual warfare and the armor of God..
The point I'm making here is that this letter.
is full of really heavy topics and really big things.
and things that, as a Christian, we say,.
wow, how is this possible?.
God, how could I ever live this out?.
How could I ever see these things happen in my life?.
How could I be transformed in this way?.
How could I be brought from darkness to light?.
These are really big things..
The climb is steep, and that's what Paul is letting us know..
That's why this prayer is in the beginning of,.
it's in the first chapter, it's in the beginning.
of the letter, he's saying, listen,.
there's some tough stuff ahead of you..
It's not gonna be easy..
God's called you to something big and something important.
and something significant, and it's not gonna be simple.
or easy, but you can have hope because God's given you.
everything you need in Christ..
And that's a really beautiful message for us..
We need to know that..
We need to know that as we look forward in our own lives.
into this year and into our future, and we say,.
hey, there's probably gonna be some challenging things..
Some of us may be sitting here already knowing.
what some of those are..
We may say, okay, I already know that there's some things.
coming up in my life, in relationships that I have.
that are gonna be challenging, or in my workplace,.
and it's gonna be really tough, or in my school,.

$^{241}$or in some decision I have to make,.
and it's not gonna be easy..
Maybe we already know what some of those things are..
Maybe we don't know..
Maybe they're gonna hit us in a few months..
Or maybe for some of us, those are these things.
that we see in the book of Ephesians..
We're saying, I want to be transformed..
I wanna be more like Christ..
I want God to use me more in the lives of others..
I wanna grow as a disciple maker..
I wanna grow in my service to God in the church,.
and I know myself, and I know that that's not.
gonna be easy..
And as we think about those challenges,.
and the fact that the climb is steep,.
we can have hope because we have the power of Christ..
This prayer is prayed over people who are facing.
a lot of challenges, and they need to know.
the hope that they have..
And this prayer isn't just, it's a pep talk,.
but it's more than a pep talk,.
because it has spiritual power..
He's not just saying, hey guys, let's think positive,.
and I'm gonna tell you some stuff to make you feel good..
It's more than that..
He is encouraging them, but he's also saying,.
God is at work through the power of the Holy Spirit.
in your lives in a supernatural way,.
and God will carry you through these things..
Look at what he says..
He says that you may know what is the hope.
to which he has called you..
What are the riches of his glorious inheritance.
in the saints, and what is the immeasurable greatness.
of his power toward us who believe.
according to the working of his great might?.
That's the hope that we have..
Our hope is in the gospel, the good news of transformation.
by grace through faith in Christ..

$^{281}$Our hope is in the power of God..
And this is, we need to understand.
that this power is God's power..
It's given to us as his people through Christ.
and by the Holy Spirit, but it's God's power..
It's power, as he says here in verse 19,.
beyond measurement, it's power beyond comprehension..
This is the power of God the Father,.
the King of the universe, the creator who made us.
and everything that we see in the entire universe..
And sometimes we need to be reminded of that..
As we face challenges in our lives,.
and we say, these are big challenges,.
this is a big deal, I don't know how I'm gonna.
get through this, and it seems like there's no way out,.
and this situation seems impossible,.
we need to be reminded that we're not just trusting.
in a God who can kinda work things out,.
or who is a good negotiator..
That's not the God that we're trusting in..
We're trusting in the God who has absolute power.
over everything in our life and in the universe,.
the God who created it all, and the God who loves us..
And so we have a sure hope, as the scripture says,.
we have something that is not gonna change,.
it's not gonna fade away, it's not gonna evaporate,.
because we're trusting in the power of God..
And Paul reminds us of this in the next verse..
In verse 20, he goes on to say that he worked in Christ.
when he raised him from the dead and seated him.
at his right hand in the heavenly places,.
far above all rule and authority and power and dominion,.
and above every name that is named,.
not only in this age, but also in the one to come..
Paul's reminding us, hey, you have this power,.
and these immeasurable riches of grace.
that are given to you by God..
What God?.
The God of the universe, the creator God,.
the God who raised Jesus from the dead,.

$^{321}$the resurrection power of Christ,.
is at work in our lives as believers,.
and we need to understand that..
This is not a small thing..
This is the same God, the same power.
that raised Jesus from the dead.
that is at work in our lives..
It's an eternal power..
And this power is fully vested in Jesus Christ.
as the son of God, but through the Holy Spirit,.
it's at work in us as well..
It's not something that we take hold of.
and we use for our own purposes..
It is God using it for his purposes through us,.
but we can have confidence in that..
Not a confidence where we believe in ourselves.
and say, well, I know my year's gonna be great..
I know things are gonna go well for me.
because I'm really confident in myself,.
and I'm gonna do such a great job.
of navigating all these challenges..
No, it's a confidence in Christ where we say,.
you know what?.
I know my year's gonna be challenging,.
and I know I'll probably make some mistakes,.
but God is at work in my life through the power of Christ..
His grace is gonna be at work even when I mess up,.
even in my relationships or my workplace.
or whatever other area of challenge I have..
Even when I make a mistake, God's grace is at work..
God has the power to heal..
God has the power to redeem even my mistakes..
God has the power to change people's minds..
God has the power to work in situations,.
and so I can have confidence in God's work.
in my life this year,.
and peace and comfort that comes along with that confidence,.
not because of me, but because of Christ,.
because of my relationship to him..
We can believe in him and believe in the power of God.

$^{361}$at work in our lives..
One of the things we see here is that this power.
is for God's glory, it's not for ours..
This prayer is so beautiful,.
and it's there to encourage us as believers,.
but it's not centered on us..
It's centered on God..
God doesn't just empower us for our own purposes.
so that we can say, okay, now I can go out.
and achieve my objectives in life,.
and that's what it's about..
God gave me power so I can do what I want..
That's not what the scripture teaches us..
That's not the gospel..
The gospel is that God has his objectives..
He's at work in our world..
He's at work in our life,.
and as we align with him in obedience,.
his power flows through us,.
but it's to achieve God's objectives, God's purposes,.
not just our own..
Now, the beautiful thing about that is that.
our God is a God of love..
The scripture's clear about that as well,.
and so as we are walking with God.
and we're living in obedience to him,.
he's going to bring goodness and blessing.
and joy into our lives because he's a God of love,.
but we need to understand that ultimately,.
the glory is his and the power is his..
It's not mine, and we see that in these verses.
when he talks about what the power of God.
at work through Christ, it's all centered on Jesus..
It's all centered on him..
Far above all rule, verse 21,.
in authority and power and dominion,.
and above every name that is named,.
not only in this age, but also in the one to come..
Whose name is above every name that is named?.
Jesus..

$^{401}$That's, it's centered on him..
22, he put all things under his feet.
and gave him as head over all things to the church,.
which is his body, the fullness of him.
who fills all in all..
We see Christ at work in us, the body, the church..
He's filling all of us with his fullness,.
but it's all centered on him..
It's not about drawing attention to ourselves..
It's not about exalting ourselves..
And the mission that we have in life.
is for us to know him and experience his shalom,.
his peace, his wholeness, his health, his blessing.
as we live in right relationship with him and one another..
And for us to make disciples..
That's the mission that we have in the church..
Well, that's not an easy thing..
How does that happen?.
It happens as we align with what God is doing in our lives,.
as we live in obedience to him.
and his power flows through us..
And then we see those things begin to take place..
We begin to live a life that's defined by love for others.
and compassion and defined by holiness and righteousness.
and defined by service to others.
in the community of the church..
And all these things that we want to see take place,.
those things take place as God's power.
begins to work through us..
And it's for his glory..
The other thing that we really see in this passage.
is that Paul is trying to help us understand our identity.
so that we can live out the mission God has for us..
If we're gonna have hope for the future,.
a part of that is understanding who we are..
And I love that in this passage,.
Paul is really focused on our identity..
He's not focused on the forms of what we're doing..
Now, here's what I mean by that..
Paul is praying for this church,.

$^{441}$this small group of believers.
who are facing a lot of challenges..
And if we went back into the background of Ephesus,.
we could explore that further, right?.
They're in this very secular city.
that is known throughout the known world.
as one of the largest and most important centers.
of the kind of European Roman Empire there at the time..
And it's known throughout that area for its idol worship..
It's known throughout that area for its secularism,.
we would say, and its paganism, right?.
So this is the background of where they live..
So obviously, they're facing all kinds of challenges,.
both in their personal lives,.
but also in their life as a church.
in order to be disciples together and make disciples.
and help people know Jesus..
They've got this really steep climb ahead of them.
as we talked about earlier..
And what does Paul pray for them?.
Look at what he prays..
He doesn't pray for the forms of ministry or lifestyle..
He doesn't say, "Church of Ephesus,.
"I'm praying that you guys.
"would really have an effective ministry,.
"and you'd learn how to start a great men's ministry.
"or a great women's ministry,.
"or you'd learn to do great work with kids,.
"or you'd learn to really have strong marriages,.
"or that you would learn how to.
"make great financial decisions,.
"or that you would all be successful in your jobs.".
That's what's important, right?.
All those practical things..
But he doesn't pray any of that..
He's pointing them towards their identity in Christ..
He's praying that they'll be full of hope.
because they know who they are in Christ.
and the power that they have..
Not that they'll have great programs or lots of money,.

$^{481}$or even be really effective.
in what they're trying to do..
I think that's really huge for us..
When we think about God's work in our life,.
and we look forward to this year,.
and we say, "Okay, here's some of the things.
"that I'm praying for.".
That's great..
All those practical things are important,.
and God loves you,.
and he cares about those things in your life,.
and he's engaged in those things,.
in your job, and in the needs that you have,.
and in those relationships..
God's engaged in all those things..
They're important to God..
But the most primary thing for you.
to be healthy in all the other stuff.
is to know who you are in Christ,.
and to walk in that..
And that's why Paul prays this,.
because he knows for this little church,.
this group of believers,.
maybe it was big..
I don't really know how many people were there..
But for this group of exiles,.
we'll say it that way,.
I want us to understand that they were facing challenges,.
whether there was a lot of them or a few of them..
They were in a difficult place..
And for this group of exiles in a very difficult place,.
what does he pray?.
He prays that they'll know who they are in Christ,.
and know what they have in Christ..
Because he believes,.
I think the reason he prays this.
is because he believes.
that if they know who they are in Christ,.
and what they have in Christ,.
then everything else gets fixed, right?.

$^{521}$Then they begin to have healthy marriages,.
and they begin to be effective.
in the ministry they're doing,.
and they begin to effectively make disciples.
with their neighbors,.
and they begin to truly love others,.
and serve in their community,.
and all that other stuff,.
the practical stuff that he could've prayed for them,.
he doesn't pray that,.
because he knows if they get this,.
then the rest of it falls into place..
Like Jesus, who told us to seek first the kingdom of God,.
and his righteousness,.
and everything else will be added to us in Matthew 6:33..
I think that this prayer has a very similar theme to it..
Paul's saying if you will know who you are in Christ,.
and know the hope that you have in him,.
then the rest of this stuff will sort out,.
and God will be at work in those areas of your life..
And so I think that's a powerful consideration for us.
at the beginning of this year..
As we think about all those specific.
and practical challenges,.
again, nothing wrong with thinking about those things,.
and nothing wrong with asking God.
to be active in those things..
That's important, God loves you,.
he wants to be at work in those parts of your life..
But the thing that we need most.
is to know who we are in Christ,.
and to be really secure in that identity,.
and to be really focused on living into that..
And if we're doing that,.
then those other areas,.
God is gonna be at work,.
and they're gonna get sorted out..
And that's what's happening here..
Paul is praying that these believers.
would walk in the hope of God's work,.

$^{561}$and God's power in and through them,.
as they seek to know and experience him..
Walk in that hope as they seek to share the good news,.
and make disciples in their very broken and hostile world..
And that's true for us as well..
We need to walk in that hope.
of God's work and power in and through us..
Walk in the hope that we belong to Christ,.
that we're forgiven, that his grace covers us,.
that we are being transformed by him,.
that during 2024 we're gonna be able to seek him,.
and know him more, and obey him more,.
and become more like Christ..
Walk in that hope..
And as we seek to share the love of Christ with others,.
and share the good news,.
and make disciples in our broken and hostile world,.
God's power is gonna be at work through us..
Walk in that hope..
We need to know the hope that we have..
Verse 18, "Having the eyes of your hearts enlightened,.
"that you may know what is the hope.
"to which he has called you..
"What are the riches of his glorious inheritance.
"in the saints?.
"And what is the immeasurable greatness.
"of his power toward us,.
"who believe, according to the working of his great might,.
"that he worked in Christ when he raised him from the dead,.
"and seated him at his right hand in the heavenly places?".
That's the hope that we are called to walk in as well..
And I would just encourage us,.
as we think about that,.
that walking in that hope.
allows us to have peace,.
it allows us to have security,.
and it also allows us to be prepared.
for whatever's in our future..
All those things that we don't know about,.
that will come in 2024 and beyond..

$^{601}$I remember, today's the 14th,.
so tomorrow, the 15th,.
is kind of a special anniversary for me,.
'cause we moved to Chicago on January 15th,.
as of tomorrow, 10 years ago,.
so when we moved here..
And our family came,.
we just had these two at the time,.
our other three were born here,.
and we didn't grow up here,.
we were coming here from a state pretty far away,.
that's further south than the United States,.
so we were not used to these winters..
And when we got here, it was very cold..
We moved into an apartment,.
sight unseen, we had just found it on Craigslist,.
which was a thing back then,.
I know no one uses Craigslist anymore,.
but back then, people used to use Craigslist..
And we found an apartment,.
and signed a lease,.
and pulled up in a big van with all of our stuff,.
and that was the first time we had actually seen it,.
other than the pictures,.
and went up the stairs..
And we moved here,.
'cause we felt like God called us here to do ministry,.
so we were just trying to respond in obedience.
to what God was doing in our lives..
And we did, all those things hit us right in the face,.
just the reality of being in the city,.
we hadn't lived in a city this big,.
but also the place we were living,.
and all the challenges of it..
I mean, from the very first day,.
it was just all right there, and the cold..
There was like a record-breaking,.
not record-breaking,.
but one of the top five snowstorms in Chicago history,.
I think, was about a week after that,.

$^{641}$snowed like 20 inches..
And I remember just being like, what's happening?.
(laughing).
And it was also a very cold winter,.
it was one of the winters when like 97\% of the lake.
was frozen, and nothing thawed out.
until the third week of April..
So it was just a shock to our system in so many ways..
But like I said, we were just trying to obey the Lord..
And I can honestly say that my wife and I had no idea.
(laughing).
the things that we would encounter.
over the next 10 years in our life..
A lot of joyful and beautiful things.
that we've gotten to experience,.
but a lot of really challenging things as well..
Challenging things for my wife and I,.
challenging things in the life of our family,.
but very challenging things in ministry.
that were ahead for us..
The thing that propelled us forward 10 years ago.
and still does today is this, the hope that we have in Christ..
If it wasn't for that, there's no way I would have made it..
No way I would have made it..
A lot of people that I met at the time.
thought my wife and I were crazy for even trying,.
and we were..
From an earthly perspective, it didn't make sense.
what God had called us to do..
It really didn't make sense..
And there's still plenty of days now.
when it doesn't make sense..
But we have hope because of the God that we serve,.
the God who loves us, and because of the power of God.
that's in work in our lives through Christ..
And that's true for each one of us in this room,.
for each one of us who's watching today..
If you know Jesus and you have relationship with him,.
as you look forward in your life, you say,.
"Hey, I wanna obey..

$^{681}$"And if I'm honest, I have no idea what that means.
"over the next two years or five years or 10 years,.
"or even over the next one year..
"I have no idea what it means for me to walk with Jesus.
"and try to surrender to him and do what he has for me..
"And there's probably stuff that I'm not prepared for,.
"and that if I knew about it now,.
"I would not wanna go forward.".
(laughs).
Definitely some of those things in our life,.
that if we had known 10 years ago.
what we were gonna experience,.
we would've had some serious second thoughts..
But we have this hope, and we need to know that hope,.
this hope that Paul prays over the Ephesian believers.
of the immeasurable riches of the grace of Christ..
It work in us, and it pushes us forward,.
it propels us forward..
Even when we encounter adversity and challenges,.
to know that God's power, the resurrection power,.
is at work in our lives through Christ,.
and that he has blessing and goodness.
and protection and peace even in those difficult times,.
and that he's at work, and that he's gonna use us,.
if we'll be obedient to him, he's gonna use us.
for his glory, and it will also be for our good.
because of his great love for us..
That's the hope that we have, and that's the hope.
that as we start this year, I pray that we know.
and rejoice in that hope..
Let's pray..
God, I thank you for the hope that you give us,.
and God, I just pray over everyone in this room right now,.
and anyone who's watching online, God,.
I know that there are a lot of challenges.
and a lot of questions in this room,.
and a lot of things that we look at in our lives.
that maybe are happening right now,.
or maybe they're happening in the weeks and months to come,.
things that we look at and we say,.

$^{721}$God, we don't understand, we don't know how to navigate this..
We see the brokenness of our world,.
we see the struggles in our own lives,.
and yet, God, there is such great hope in you..
There is such great hope in the immeasurable riches.
of Christ that have been given to us by grace..
So God, I pray that we would know that..
We would know the hope that we have,.
and that it would propel us forward with strength.
and with confidence in you and with joy.
into whatever you have for us this year,.
into walking more deeply with you,.
into serving you in a deeper way than we ever have,.
into loving you and seeking you and knowing you.
more than we ever have in the past,.
into loving the people around us.
deeper than we ever have in the past..
God, I pray that that hope would really propel us forward.
into what you have in our lives this year..
In Jesus' name, amen..
[BLANK AUDIO].
\newpage



\section{}
\label{sec:__gIiUwiPLI}
\textbf{原子智慧 3 - 陳崇基牧師}
\newline
\newline
連結: \href{https://youtube.com/watch?v=--gIiUwiPLI}{\texttt{https://youtube.com/watch?v=--gIiUwiPLI}} ~~~~ 語音日期: 2024-01-16
\newline
\newline
\hyperref[sec:RlCj0uSBVME]{\small{< < < PREV SERMON < < <}}
~
\hyperref[sec:index_chronic]{\small{[返順時目]}}
~
\hyperref[sec:index_scriptual]{\small{[返順卷目]}}
~
\hyperref[sec:G8td0HdJn8w]{\small{> > > NEXT SERMON > > >}}
\newline
\newline
$^{1}$大家好.
很高興見到大家.
雖然今天非常冷.
聽到我的聲音已經好了85\%左右.
還有15\%.
還要休息一下.
我也不推動自己那麼多了.
無論如何我們繼續這個系列.
這次是第三講.
這個系列是一些integration.
我們嘗試去結合一下.
所謂一些外面所說的行為心理學的insight.
來看看聖經裡面的一些智慧文學.
Wisdom literature裡面是說什麼.
所以這個三部分的講道系列今天完結了.
這個講道系列就叫做原子智慧.
出自這本書叫做Atomic Habit.
是一本很流行的書.
在去年到近幾年.
是非常多人看的.
其實我也沒有留意.
是Amanda叫我.
不如你看看這本書.
我說看一下.
誰知道就對了.
在這個年頭的時候.
我們可以思想一下.
這些Atomic Habit.
一些很小的習慣.
如何幫助我們心靈去改變.
去結合一下聖經裡面智慧文學.
特別是Proverbs.
《箴言》裡面所說的話.
我們review一下在過去兩次.
我們曾經說過的一些原則.
大家都記得了.
這六個原則.
1\%的power.
每一天你做1\%的改變.
365天之後.

$^{41}$你有的效應是30多倍的效應.
然後有賞罰的效應.
就是每一次我們定立目標.
或者去做一些習慣的時候.
我們定立一些incentives.
或者disincentives.
也可以幫助我們的.
定立一個目標的時候.
我們更加需要去明白.
where and when.
記得嗎.
specific.
一些很具體的東西.
你要去定立的時候.
就會更加容易去養成一些良好的習慣.
第四就是我們身邊的朋友.
這個社群對我們的力量.
第五就是default的power.
預設的一種影響力.
我們一些習慣預設了.
然後生命裡面才會養成這些習慣.
然後我們生命裡面有一些效應.
第六就是make it personal.
越個人越有效.
這次我們看第七個principle.
就是environment的principle.
就是環境對於我們生命的改變是重要的.
我記得很多年前.
那時候的電視都沒有那麼扁的.
CRT的電視.
大家有沒有.
家裡有幾部放在這裡.
不知道有沒有封塵.
我記得我們睡房裡面是有一部這樣的電視.
我是拉上牆上.
你覺得這個習慣好不好.
不行.
很不妥.
不過我拉上去的時候是有趣的.
因為多了一部.

$^{81}$不如拉在那裡.
但是那個環境根本就是幫助不了你的.
所以一段時間之後我就拆了下來.
然後放在後面.
封塵.
我們的環境對我們影響很大.
有一些study告訴我們.
譬如.
你說你想on diet.
就是不要吃那麼多東西.
你知道我們美國人吃很多東西.
太多東西吃.
study告訴我們.
如果你真的on diet的話.
有什麼habit可以幫到你的.
就是每次吃飯之前.
你喝兩杯水.
那你就自然無謂吃那麼多東西.
所以你的環境是幫助到你的.
很多年前我們小朋友小的時候.
我們芝加哥的人有時候上去人多.
小朋友小的時候他們不懂吃什麼好東西.
他們懂得吃什麼嗎.
我們就想了一個方案.
就是先帶他們去吃一些麵包.
先把他們包起來.
兩塊錢一個麵包.
包起來.
好吃啊 蛋撻好吃啊.
然後我們才去吃海鮮.
結果他們不用吃東西了.
這個就是環境對我們的生命有很大的影響.
反過來想.
養成一個良好的習慣.
應該有什麼環境.
以致我們能夠更加容易做到.
有些study去研究.
就是其中一個.
在波士頓有一間醫院.
那個醫院的方案就是.

$^{121}$希望能夠鼓勵那些員工.
去cafeteria的時候.
他們就要吃一些健康的東西.
沙拉 菜等等.
他們結果做了什麼.
就只是把cafeteria的layout改一改.
那些水 一瓶一瓶的水.
就放在門口 最容易拿到的地方.
然後那些fountain drink 即是汽水.
就放在最後最後 最難拿到的地方.
對不對.
三個月之後他們發現.
員工喝水的比率.
是多了25.8\%.
三個月之後.
就因為這樣的layout的緣故.
然後那些soda或者汽水.
就少了11.4\%.
就是那個環境上面的layout.
就可以改變了人的習慣.
同樣的我們辦公室.
辦公室.
那些cafeteria都是一樣.
你公司的cafeteria.
應該有jar of candy在那裡.
candy junk food.
如果將那些junk food.
放在break out room那裡.
rather than.
方便你在辦公室的桌頭.
人的習慣就自然會改變.
壞的地方就是.
過去幾年pandemic.
每個人都work from home.
你以為work from home很好.
誰知道work from home就最多junk food.
你沒事做 口癢的時候.
你又拿一包薯片去吃.
然後又拿一些不利的東西去吃.
但是point就是環境.

$^{161}$surrounding是會影響我們的.
在浸賢27章第17節.
這裡說.
用鐵磨鐵.
磨印月離.
越磨越離.
朋友當面剁磨.
也是如此.
什麼意思.
如果看27章第17節之前.
幾節.
是在說關係.
一些不良的關係.
浸賢經常都告訴我.
人生裡面有很多.
眾多的關係要處理.
也有一些不良的關係.
不良的關係會帶來.
生命裡面有危機和有試探.
譬如一些人.
他們溝通是沒有智慧的.
或者開罪人多.
稱呼人少.
好爭辯的那些人.
他說這些人要遠離.
你畫清楚boundary.
然後27章第17節.
他就說什麼.
他說鐵磨鐵就越磨越離.
你想一塊鐵.
它本身是有價值的.
但是這塊鐵和另一塊鐵.
一起相磨的時候.
就會變成更加有價值.
你要如何increase你的value.
就是兩塊鐵一起打磨的時候.
磨刀有沒有.
有沒有磨刀.
越磨就會越離.
但是一定硬物和硬物一起.

$^{201}$去打磨的時候.
就會使那件事更加有價值.
其實我們的生命都是一樣.
你的生命越來越和其他人.
去到打交道.
其他生命裡面是有智慧的人.
你的生命自然就會多智慧.
希伯來文就很妙.
刀鋒.
是什麼字.
Blade of the sword.
Blade of the knife.
其實就是face of the knife.
The face of the knife.
就是說那把刀的面.
就是那個刀鋒.
所以他下一句說什麼.
他說朋友應該當面去琢磨.
One sharpens the face of the other.
就是大家都是有刀鋒的.
我們就一定要彼此去磨擦.
Face to face.
Person to person.
以致我們能夠成為一個更加成熟.
和有智慧的人.
這個就是一個大的principle.
就是我們身邊的環境.
是會造成我們的影響.
所以就說你應該去適應.
或者包圍你自己.
更多一些有熟齡智慧的人.
第七個principle.
第八就是Habit stacking.
或者Habit bundling.
這個concept都是一些很常識的原則.
就是說你本身已經有一些的習慣.
是很成立的習慣.
你是不需要想的.
大家都有的.
今早起床就一定要刷牙.

$^{241}$You better.
刷了牙.
然後你有一些已經固定了的習慣.
你想一想你生命裡面.
其實一日二十四小時.
經常都有這樣的習慣.
我今天上車的時候.
很多時候我左腳都想找個按鈕踩.
原因就是因為我開過人家的摩托車.
所以我上了一輛自動車的時候.
你那些不開過人家的摩托車.
你不知道你有第三隻腳.
第三個按鈕要踩.
所以很多時候我上了車的時候.
踩都沒有東西踩.
原來這個是自動波.
這個就是我們的second nature.
就是說我們已經有一些establish了的habit.
一些已經成立了的習慣.
什麼叫做習慣的堆疊.
就是你要養成一個新的習慣的時候.
將一個新的習慣扣上一個已經成立了的習慣.
那你就會更加容易一點.
我記得上一次問過大家有沒有用牙線.
有沒有.
你怎樣才會養成一個習慣用牙線.
你當然不會說臨睡前想起就用牙線.
刷完牙就用牙線.
就這麼簡單.
牙線放在哪裡.
就是放在你牙刷的地方.
然後你就會養成一個新的習慣.
這些都是common sense.
在proverbs裡面.
在箴言裡面都有這樣說.
堆疊的習慣.
四章二十節.
父親就教那個兒子.
他說我兒要留心聽我的話.
則以聽我的言語.

$^{281}$不要把它偏離你的眼目.
要存在你的心裡面.
因為找到他們的.
找到智慧或者智慧的言語的話.
就找到生命了.
你就得到全新的意志.
你要保守你的心.
勝過保守一切.
因為生命的傳言由心發出.
箴言頭九章一至九.
都是introduction.
說了很多很多的東西.
其實就是父親去叮嚀他的兒女.
就是勸導他.
你初出茅廬.
你沒有智慧.
不夠智慧的.
你就應該要學智慧.
所以頭九章經常有金句.
譬如這句四章二十三節.
你要保守你的心.
勝過保守一切.
歌也有一場.
因為生命的源頭是由心發出的.
等等等等.
你會發現這些這樣的箴言.
或者智慧文學.
經常都是一對一對的.
有A句有B句.
每一節有A句有B句.
有時候有三句ABC.
所以一句疊一句.
一句疊一句.
在這裡說留心聽.
你要留心聽我的話.
留心聽.
其實意思就是說守望.
如果你是士兵.
你去到守城的話.
在城樓上面你做什麼.

$^{321}$就是留心.
留心聽.
留心聽什麼.
就是有沒有馬匹的聲音.
有沒有軍隊來的聲音.
有沒有人來攻擊你的城市.
好像士兵那樣看守.
然後他再疊上去的就是.
則以而聽.
然後第二節.
他說不要偏離你的眼目.
整天要在你的眼前.
就是說將上帝的話語放在你的眼前.
你才會養成這個習慣.
不單止這樣.
再疊上去的就是.
傳記在心裡面.
保存在心裡面.
儲存在你的心裡面.
最近電影節目經常說.
那些CD那些百分比的利率是高了.
定期存款.
存在心裡就是這樣的意思.
將它放在那裡.
你不要動它.
留在你的心裡面.
以至你心靈裡面有更多肅靈的習慣.
找到智慧就是找到心靈.
並且你要去保守你的心.
去到保守你的心.
都是一疊一疊再一疊加上去.
這個就是行為心理學裡面說的原則.
就是習慣是需要堆疊上去的.
已成立的習慣再加一些新的習慣上去.
那你能夠達到你的目標的機會率就會高很多.
以前讀大學的時候.
讀大學的時候.
有一科我一定要讀的.
就是chemistry.
還記得是chemistry 1B.

$^{361}$整個學期裡面有十個考試.
十個測驗.
誰知道第一次測驗的時候.
我去到上課的時候.
我忘記了那天是測驗.
所以第一次測驗我就空窗上陣.
結果就拿F.
那我就覺得很糟糕.
這個一定要pass的.
不pass就很糟糕.
還有九個測驗.
我就決定不如去到盡力去做.
於是我每一次上課之後.
我就馬上上完課後就會去圖書館.
我就用一個多小時的時間去review一下.
剛剛lecture裡面講過的是什麼.
或者教授講了的是什麼.
一個半小時.
有很多同班同學.
他們就去coffee shop或者回家.
這個習慣很自然.
因為你去coffee shop那裡舒服一點.
可能那裡你會覺得有更多時間可以溫習.
或者你回家更多時間.
三四個小時.
多好.
我就立刻收拾東西.
馬上去圖書館.
然後我就review剛剛那堂所講的東西.
我就一天一天這樣做.
結果之後那九個測驗.
我都是拿A的.
結果全班最高分的是我.
其他人他們開始拿A的.
反而最後是拿不到A.
很簡單的一些strategy.
就是habit stacking.
habit bundling.
結果就可以很大.
智慧的文學都是這樣教我們.

$^{401}$將你的習慣堆上去.
如果你說你沒有讀經的習慣.
如果你說沒有祈禱的習慣.
你將這些習慣放在一些已經成立的習慣.
將它加上去.
以至你堆疊上去的時候.
你每一天做的改變只是很小的1\%.
但是你end of the year.
三百多天之後.
你所受的效益可能是很多很多倍.
第九個principle.
就是讓改變使它吸引.
去到描述改變.
或者去到想起改變.
用一些attractive的wording.
或者一些吸引力的描述.
以至你覺得改變是更加容易的.
你知道那些酒店的房間.
很多時候要浪費很多資源就是洗那些毛巾.
因為你去完酒店.
你抹完臉.
毛巾放在哪裡.
很多時候放在地上.
放在地上收房.
你拿毛巾去洗.
但是你抹完臉.
臉都很乾淨.
那麼大塊毛巾.
你只是用一個很小的地方.
但是你又要洗.
所以花很多水和人力物力.
去洗這些毛巾.
於是很多酒店.
特別是美國.
就希望能夠節省這些資源.
所以在酒店裡面.
他們發現有些study.
可以幫助人再用回那些毛巾.
結果在酒店裡面.
很多時候有這樣的sign.

$^{441}$那些sign就是說.
75\%的旅客都會reuse他們的towel.
他們會reuse他們的towel.
然後你看到這些wording的時候.
是喔.
我是不是那75\%的人呢.
有一個常規給你.
然後你就會confirm.
你就會去跟.
就告訴你這是保護環境.
甚至是save the planet.
是不是很偉大.
你多用一次毛巾就可以save the planet.
多厲害.
他就用這樣的wording幫助人.
去再reuse再掛那條毛巾.
不是這樣丟在地上.
這些是很常規的psychology.
behavioral psychology.
在浸泳裡面.
同樣是有這樣的東西.
說到我們生命改變的時候.
得到智慧的時候.
是怎樣去描述的呢.
他說智慧為首.
所以要獲得智慧.
要用你一切.
生命裡面的一切.
來去換取肅靈的智慧.
跟上帝之間的關係.
在上帝的生命裡面.
去到活出來.
用一切來去換取這些的智慧.
高舉智慧.
他就使你勝過.
擁抱智慧的他.
就使你得到尊榮.
他必將你.
一個欣慰的華官.
加在你的頭上.

$^{481}$將榮免賜給你.
叫你升高.
叫你尊榮.
叫你有欣典.
欣慰的華官.
好像做皇帝一樣.
有智慧的就好像做皇者一樣.
Crown of Splendor.
是一些很好的描述.
不單止是一些純粹很好的描述.
亦都是一些真實的情況.
一個有智慧的人.
在上帝面前就得到上帝給我們永遠的尊榮.
這起碼就是make it attractive for you.
我剛才說過.
箴言頭九章是什麼.
是序言.
是introduction.
introduction想做些什麼.
就是希望推動你追求智慧.
追求上帝給我們的spiritual wisdom.
去稱讚你.
去說一些好聽的話.
告訴你真正的價值在哪裡.
價值鍊成Crown of Splendor.
你就可以得到這件事.
但你需要追求.
在英國以前有一個很著名的speech therapist.
叫做Lena Rustin.
如果想搜尋一下這個人.
Lena Rustin.
是英國的一個speech therapist.
她是專門去醫治那些兒童.
幾歲四五歲有口窒的問題.
有些孩童說話可能說不到.
或者發音發得不清楚等等.
或者沒信心去說話.
很多治療師都是去做一些個人的事情.
幫助小孩怎樣發音.
怎樣說pronunciation做得好一點.

$^{521}$怎樣呼吸等等.
都是一些individualistic therapy.
Lena Rustin除了做這件事之外.
她就加多了一個規矩.
她就說見了一家人.
小孩的父母等等.
全部人見了他們.
她就說由今天開始.
你家裡就需要做一件事.
每天都要做的.
每個人都要做的.
她說每一個家人.
必須每一天都去稱讚另外一個家人.
找一些東西稱讚他.
並且被稱讚的人.
就要樂意接納讚揚.
要樂意接納.
不是我們華人這樣.
不要稱讚我.
然後我們就踩低自己.
我們最興奮是這樣.
Lena Rustin就說所有的家人都要一起做.
開始可能有些困難.
可能覺得沒什麼動力去做.
但你要養成這個習慣.
就是說一些讚賞的話.
然後聽到讚賞的話的時候.
必須每一個人每一天都要做的.
譬如媽媽.
媽媽就說John.
I love how you share with your brother.
Good job John.
John就不可以說.
我們說No No No 不用稱讚我.
John就要說什麼呢.
Thanks.
I appreciate the comment.
It makes me happy.
Thanks mom.
一天一天就是做這件事.

$^{561}$純粹在說話上做.
Lena Rustin加了這件事的時候.
她就發現.
這可能是一些很簡單的習慣.
但卻是比有口吃或有口吃的孩子.
能夠改變得更快.
很奇怪吧.
因為Lena Rustin認為.
一個人能夠得到改變.
一個口吃的兒童能夠改變到他.
以至他有信心來溝通.
這個不是他一個人的事.
是一家人的事.
也不是單單家人這麼簡單.
是有環境使他的改變能夠更加有效.
讚賞和接納讚賞.
不單單是改善了口吃的問題.
Lena Rustin甚至有很多的客戶.
回來告訴他.
你不單止改變了我們的孩子說話.
還改變了我們夫妻之間的關係.
改變了我們整個家庭的關係.
這個很驚人.
你知道我們教會.
我們跟另一個的教會合作.
我叫UCA Urban Christian Academy.
我們經常接觸一些社區裡面的家庭.
特別是華人的家庭.
有些小朋友不敢說話.
沒有信心說話.
並且說話是有口窒.
因為我老婆在那裡工作.
所以經常在家裡聽到這樣的現象.
知道這樣的現象.
華人的家庭就是.
你做得不好.
就罵你.
你做得好.
就不說任何話.
你有口吃或口窒.

$^{601}$就罵你.
罵到你沒有口窒為止.
我們發現一些小朋友和學生.
口窒變得更加厲害.
溝通更加難.
你知道這個世代多煩嗎.
每天都看著手機.
更加不懂得跟別人溝通.
你心明意有智慧.
將來初出茅廬的時候.
你都要靠口賺錢.
不只是牧師靠口賺錢.
你都要跟別人溝通.
跟老闆溝通.
這就是環境問題.
所以我們為什麼要跟UCA合作.
因為他們的哲學.
不只是學生來到.
我們幫他們擊敗所有功課.
完聖所有功課.
然後考試.
考到最好的學校.
最好的學校.
我們超級市都有這樣的壓力.
雖然不及香港那麼厲害.
這裡也是.
入到哪間很好的名校.
入不到的時候.
就要摒他.
希望摒到他入到為止.
越摒他越難入到.
目標教錯了.
我兒子今天經常說這句話.
因為他在U of I.
他入了U of I.
那時候他中學的時候.
入不到他每個同學都想入的那間.
好一點的學校.
他就一個人去另一間學校.
他說現在不是end up一樣嗎.

$^{641}$那時候入了好一點的學校.
反而他入了一間差的大學.
為什麼是一樣的.
為什麼他們的哲學.
不只是學生的問題.
也不只是家長要幫忙.
告訴家長.
勸告家長.
要除去華人家庭的壞習慣.
養成一些新的好一點的習慣.
是家庭和公立學校.
或者外面學校的老師.
和他們補習的老師.
是全國的環境一起來做的.
不過你不會使得那個人的生命.
能夠得到這裡所說.
Crown of Splendor.
榮美的冠冕.
有了智慧.
成熟的一個生命.
在真言裡面.
何況何尚不是這樣說.
當我們追求到智慧.
擁抱到智慧的時候.
我們就能夠得到一個有尊榮的生命.
第十個原則.
就是要確立一個系統.
而非確立一個目標.
這個聽起來有一點反常理.
因為我們經常說.
新年有什麼目標.
減肥要減多少磅等等.
這不是重點.
重點不是要訂立一個目標這麼簡單.
目標好不好.
是否需要.
當然是需要的.
但最重要的是確立達到目標的系統.
那個微小的習慣.
很多不同微小的習慣.

$^{681}$Tiny Habits.
能夠使你去到那個目標.
如果你沒有那些習慣.
沒有那些系統的時候.
你是絕對去不到的.
我們看看這本書.
我想引用一段.
是James Clear在Atomic Habits裡面所說的.
他說成果和目標沒有太大關係.
他與我遵從的系統息息相關.
目標有助於確立方向.
但是系統是進步的最佳解答.
如果花太多時間思考目標.
他花太少時間來設計那個系統.
就會浮現很多問題.
每一個奧運的選手都想奪金.
有誰不想奪金.
每個目標都一樣.
如果成功者和失敗者.
有著同樣的目標.
贏家和輸家的差別就不在於目標.
因為每個人都想奪金.
為什麼有些人奪金.
有些人奪不到金.
目標一直在那裡.
直到他們執行了持續追求微小改善的系統.
才產生生命裡面的成果.
Does it make sense?.
不只是訂立目標這麼簡單.
目標可能幫到你.
但是如何去到目標中間的那一句.
執行持續追求微小改善的系統.
這個更加重要.
一步一步加上一些良好的習慣.
有沒有留意那些音樂家.
他們練習的時候.
當然有目標.
但他們總是有一些習慣.
練來練去都是練那一句.
然後一個大的樂章.

$^{721}$他會分析哪裡最難.
哪裡最有挑戰性.
要如何去練習.
每一次練習的時候會有熱身等等.
都是一些很小的習慣.
這些就是那個系統.
那個系統如何去推動他.
能夠在他的藝術上再進一步.
我們的職業.
甚至我們屬靈的生命.
其實都是一樣的.
屬靈的生命都是一樣.
十個的原則.
也都扣上去聖經裡面.
哲偉文學裡面所說的.
在《箴言》我講過頭九章.
就是那個introduction.
這個是當時希伯來文化.
有特色的地方.
因為智慧這個字.
wisdom這個字.
是陰性的.
是女性的名詞.
我們中文就沒有這個東西.
有很多語文有分男性,中性和女性.
陽性,陰性.
智慧這個字.
是陰性的.
是女性的.
所以在頭九章.
用擬人化的方法去形容智慧.
什麼是智慧呢.
就是一個女人.
那個女人我們可以叫她做Lady Wisdom.
或者是Woman Wisdom.
但同樣是另外一個女人.
那個女人就是和智慧相反的.
就是使你生命裡面敗亡的.
那個就是愚昧.
Folly.

$^{761}$或者是Woman Folly.
Lady Folly.
所以在《箴言》裡面.
很奇怪.
經常提到一個女人.
但其實和背景很有關.
因為在說什麼.
父親和初速茅廬的年輕人.
跟他說你生命裡面需要處世有智慧.
有什麼可以吸引到他.
就是找一個女孩子.
他希望找到一個什麼女孩子.
就是有智慧的女孩子.
所以智慧就是擬人化的.
個人化的.
你都可以描述她.
什麼叫做Wisdom.
就好像娶一個好的女子.
你娶到一個好的女子.
就使你生命裡面.
能夠有很多很多的益處.
所以在第九章.
這個Introduction.
最後那一章.
他說什麼呢.
就有兩個女子在那裡說話.
九章第一節.
智慧建造自己的房屋.
作成七根柱子.
她宰殺了牲口.
然後有美酒有筵席.
她都差派幾個侍女出去.
自己在城裡面高處呼喊.
呼喊什麼.
就是那些想有智慧.
沒有智慧的年輕人.
希望她能夠被吸引.
追求她.
同樣有另外一個Lady.
就是Lady Farley.

$^{801}$Lady Farley都是在說同樣一番話.
這裡怎麼描述她呢.
十三十四節.
愚昧的女子喧揚.
她是愚夢.
一無所知.
坐在自己的家門口.
在城裡面高處的座位之上.
她所說的話.
和Lady Woodson說的話.
差不多是一樣.
其實是兩個的妖虐.
兩個的Introduction.
一個就是智慧婦人的邀請.
另外一個就是愚昧的婦人的邀請.
留意兩個人.
都是在城的高處坐立.
是什麼意思.
城裡面的高處.
山上面.
高的地方.
所有以前古代的城市.
最高的地方是什麼地方.
就是聖殿.
廟.
無論是假城的廟.
或者是真神的廟.
都是在山的高處.
所以這裡說的是什麼呢.
智慧其實就是上帝自己.
上帝在呼喚我們.
要過一個有智慧的生活.
另外一個Lady Folly.
右邊那個.
女士Folly.
如絕的生命.
同樣做同樣的.
有力的呼喊.
就叫你去到她的山上.
城的高處.

$^{841}$這個代表什麼.
就代表假神.
代表假神.
兩個的妖虐.
就叫你選兩個的生命.
到底你選哪一個.
你選智慧.
上帝.
敬拜上帝.
還是你選的是愚昧.
假神.
所以你這樣說智慧文學.
其實不只是單單說一些常識.
說生活.
過得好.
其實是說敬拜的生命.
說屬靈的生命.
到底生命真正效忠於的是哪一個.
你的生命就應該要先去敬拜耶和華.
祂在高處去呼喊.
我們去回應.
但願我們的生命都是這樣.
無論你新的一年.
你有什麼立志也好.
沒有立志也好.
都是需要做一些什麼.
一些微細的屬靈的習慣.
將上帝的道放在你面前.
將上帝要你過的生活.
聖經裡面有個價值.
事實上放在你的面前.
以至你有一些屬靈的習慣.
是習慣去練習那件事.
習慣去實行我們的信仰.
這個就是智慧的生命.
這個就是敬拜神.
我們動心低頭祈禱.
希望我們的生命是照向上帝.
主我們多謝你.
因為你是智慧的源頭.

$^{881}$神你就是智慧.
你都是愛和平安.
你都是喜樂.
主我們求你.
讓我們的生命被你自己的恩典吸引.
你坐在高處.
主讓我們被你吸引去敬拜你.
以至我們的生命是真的屬於你.
主我求你來檢查.
我們自己的生命裡面任何一個方位.
是不是有些地方我們藏起來.
不願意來掌管.
主我們將我們的罪.
我們的軟弱.
我們再一次呈現在你的面前.
我們知道你知道所有東西.
求你來洗淨 潔淨.
主給我們一個純潔的動機.
來追求智慧.
追求生命裡面和你的關係.
主讓我們想起我們生命裡面有很多的未逝的習慣.
讓這些習慣能夠帶領我們.
使我們的生命更加成熟.
更加像你的兒子耶穌基督.
因為在耶穌基督裡面我們能夠有十字架的九印.
我們能夠成為神你的兒女.
主我們在基督裡面就是一個新造的人.
主將這個新造的心賜給我們.
以致我們每一天的習慣.
每一天我們的生活.
我們的方式.
是去擁抱你.
去追求你.
願以我們教會.
能夠成為一個更有智慧的教會.
願以我們教會成為一個屬於生命更加壯健的群體.
去榮耀神靈.
我們這樣祈禱 奉主名禱.
願上帝賜福給各位.
希望下個星期能夠見到大家.

\newpage



\section{18:1-4}
\label{sec:G8td0HdJn8w}
\textbf{Are You Limping Through Your Faith? (1 King 18\_1-4, 17-40) - Tommy Lee}
\newline
\newline
連結: \href{https://youtube.com/watch?v=G8td0HdJn8w}{\texttt{https://youtube.com/watch?v=G8td0HdJn8w}} ~~~~ 語音日期: 2024-01-24
\newline
\newline
\hyperref[sec:__gIiUwiPLI]{\small{< < < PREV SERMON < < <}}
~
\hyperref[sec:index_chronic]{\small{[返順時目]}}
~
\hyperref[sec:index_scriptual]{\small{[返順卷目]}}
~
\hyperref[sec:9LbyKyl1jKw]{\small{> > > NEXT SERMON > > >}}
\newline
\newline
$^{1}$Here we go..
Mark, it is good to see you..
Happy New Year to everyone..
For those who are here, it is great to be back..
For those who are online, it is great to be back..
For those who are not familiar, my name is Tommy Lee..
And I get a chance during the week, actually,.
when it's quiet here, my office is right here.
because we get a chance to share some space.
with Urban Voice and they have graciously,.
continuously allowed us to work here during the week..
Born and raised here in the Chicago area,.
born in Bridgeport, born in Chinatown,.
and now I live right over on 35th in Indiana..
And so a lot of times, I've been working.
on this series on Elijah..
And so I wanna go through a passage.
that we looked at the last time I was here..
I wanna at least close the loop and talk through.
some of the lessons that we learned..
Then I wanna do one more at 1 Kings 19,.
and then we're done with the story of Elijah..
And Elijah is this one guy who's really unimpressive..
He's not someone who kinda, someone that you sit there,.
he's not David, he's not Solomon, he's not even Saul,.
he's not Peter, he's kinda rough around the edges..
He comes from a rural area..
Scripture calls him, he's this Tishbit from Tishbe.
in Gilead, as 1 Kings 17 tells us..
He's this guy who just comes out of nowhere..
And look, I mean, the author even tells us he's from Gilead.
because we don't even know what a Tishbit is..
We don't even know what Tishbit is..
We know it's somewhere in Gilead,.
and that's what 1 Kings 17 tells us..
And so he comes out of nowhere in 1 Kings 17,.
and he gets a meeting with the king,.
King Ahab of the Northern Kingdom..
As we talked about the last time I was here,.
and we looked at 1 Kings 17, post-Solomon,.

$^{41}$post-David, post-Saul, the kingdom of Israel.
was split into two parts..
You had a northern part, northern kingdom,.
which was called Judah, and the southern part,.
which was called Israel..
Israel was the site of the Babylonian invasions..
That was where, when you look at Nehemiah, Ezra,.
that's where the story takes place..
Our story with Elijah takes place actually.
right in the northern kingdom..
And out of the northern kingdom,.
there's 19 kings who ruled for a long time post-Solomon..
And scripture tells us all 19 kings.
were evil in the eyes of the Lord..
None of them were good..
All of them were bad..
But the king that we're looking at today.
during the time of Elijah was King Ahab..
He was the worst of the worst of all of the kings..
But ironically, he actually was not the worst..
His wife, Queen Jezebel, actually ends up.
being the villain of the story, as you will see,.
throughout 17, 18, and 19, the first kings..
And so this guy, Elijah, this Tishbit the Tishbi in Gilead.
gets a meeting with King Ahab..
How this real guy gets a meeting with the king,.
we don't know..
But he shows up one day and he gets a meeting with the king.
and he tells them and the king this,.
"Neither dew nor rain except by my word.".
He predicts a drought..
He says, "You will not have dew, you will not have rain,.
"and it'll last for three and a half years.".
And he tells them that my motivation is.
the God before whom I stand..
That's why I'm doing all these things..
I'm doing it on behalf of the God before whom I stand..
And so from that point on, later on,.
once he says that to the king,.
he then ushers him over to Cherith..

$^{81}$God says to him, "Go to Cherith.".
And he walks over to Cherith..
And therefore, now he just doesn't,.
he lives there for at least a year.
or maybe a long number of months,.
not definitely, not a couple weeks,.
where scripture tells us he's fed by the ravens.
and he is drinking from the brook over in Cherith..
Soon after, the brook dries up in Cherith,.
and God tells him, "I want you to go to Zarephath.".
And it's a hundred miles walking..
By the way, it's the birthplace of Queen Jezebel,.
who is the villain and the queen of this entire story,.
married to Ahab..
"Go, walk over to Zarephath,.
"and there you will find a widow.".
And he comes across this widow and she's picking up sticks..
And this widow is picking up sticks and he approaches her,.
and she says to him, "Look, whatever you need,.
"I don't have anything..
"I'm picking up sticks..
"I have a little flour, a little oil..
"I'm gonna mix it with water..
"I'm gonna make some bread,.
"and then my son and I are gonna die..
"And we're just gonna eat and die..
"That's it, we have nothing..
"So Lord, you made me walk a hundred miles.
"to Zarephath for this, yes?".
The Lord knows where you are..
The Lord knows what you need..
The Lord knows where you're going..
That's the key that Elijah had to learn..
The God before whom I stand, that's my motivation..
Go to Zarephath, go to Cherith, all of those things..
We follow the Lord because one, he knows where we are,.
he knows what we need, and he knows where we're going..
It doesn't go according to plan..
It may not be where we wanna go..
It may not be what we think that we need,.

$^{121}$but the Lord knows what we need,.
where we're going, and where we are..
Later on, stories, his son dies and all that other stuff,.
and now there's a miracle that performs,.
and now today in 1 Kings 7, 18,.
which some of you guys who have been here have heard,.
and so we're gonna go through that story one more time,.
is the showdown and the battle at Mount Carmel..
And so now it's been three and a half years..
The drought is finally over..
And 1 Kings 18 tells us,.
"After many days, the word of the Lord came to Elijah.
"in the third year, saying, 'Go, show yourself to Ahab,.
"'and I will send rain upon the earth.'".
After three and a half years, go to Ahab, the drought is over..
So Elijah went to show himself to Ahab..
Now, the famine was very severe in Samaria..
Samaria is the capital city of the northern kingdom.
where Ahab has his kingdom..
The famine is very severe, the drought is very severe..
It's been three and a half years of no water..
And there later on, there's Obadiah,.
we don't have time to go through that,.
that's the chief of staff, his conversation, Obadiah..
And so now he gets to Ahab and he proposes a challenge..
Ahab, we're gonna meet at Mount Carmel..
You bring all your prophets, Abau and Asherah..
I will bring my God and we're gonna meet here.
and we're gonna call upon the name of the Lord.
and whoever, whatever happens, if the Lord brings fire,.
whichever God, whether it's Abau or my God brings fire,.
that is the true God, what do you think?.
Ahab says, I like this idea..
You have one God, I have 450, let's do it..
So they came together and they meet at Mount Carmel..
450 prophets, Abau and Asherah shows up..
Elijah shows up and all of Israel, because they're curious,.
they wanna know what happens..
They all show up and here's what he says..
So Ahab sent to all the people of Israel.

$^{161}$and gathered the prophets together at Mount Carmel..
This is verse 20, chapter 18..
And Elijah came near to all the people and said,.
"How long will you go limping.
"between two different opinions?.
"If the Lord is God, follow him..
"But if Baal, then follow him.".
And the people did not answer him a word..
The Israelites sat there and stood there..
They had no response..
They did not answer him a word..
One of the things that we talked about,.
and there's two things that I wanted to touch base.
on the second part, is are we limping.
between two different opinions?.
As Christ followers, as Christians,.
are we limping through two different opinions?.
Let me give you two illustrations of what that means.
for us in our context..
There's the first one, there's more the black.
and white version of it, right?.
In terms of limping between two different opinions..
The first one is that we limp between following Christ.
and we limp between the world..
It's black and white..
We have one foot following Christ, we're all in..
We're all in in terms of following Christ..
But there's another part of our life.
where we're attracted by what the world offers..
We're attracted by the power, the influence,.
the money, the wealth, the promises..
We limp between following Christ.
and following what the world has to offer..
Because that was Ahab and Jezebel's problem..
They wanted the prosperity and fertility.
which they believed that Baal could offer them..
For us as Christians, we want what the world has to offer.
because it's the money, the relationship,.
the job, the status that we believe.
is the answer to all of our problems..

$^{201}$That's what Judas believed..
That's why he turned on Jesus..
That's what David believed.
and that's why he made the decisions.
when he saw Bathsheba, he wanted Bathsheba..
That's what Solomon believed and all of his wisdom,.
the wisest of all the kings,.
fell into a lust of women.
because for him, that gave him fulfillment..
But ultimately, did Judas find fulfillment?.
Ultimately, did David find fulfillment?.
Ultimately, did Solomon find fulfillment?.
The rich, young ruler, when Jesus would come in front of him,.
the rich, young ruler wanted to follow Jesus.
and Jesus says to him, "Sell everything and follow me.".
The rich, young ruler had a hard time..
The black and white issue is that as Christ followers,.
we limp between two different opinions..
We limp between the opinions of following Christ.
versus the opinion of what the world has to offer..
Now, here's the one thing..
And some of you guys are sitting here right now..
So it's like my friend who I was texting.
and when I was originally did this sermon,.
this one friend texted me and she says,.
"You know, I have a hard time following.
"and giving up my life because I actually,.
"do I really wanna sell everything.
"and be poor all my life just to follow Christ?".
I actually think sometimes people believe that..
I believe that a lot of times with Christians,.
we sit here and we say, "Okay, following Christ.
"means I have to sell everything,.
"I have to be poor and give up everything.".
Is that what, no, no, no, no, no..
I don't think that's really what God is trying to say..
Being a Christian does not mean that you sit there.
and you sell everything and you become poor.
and you give up everything..
For some of us who are sitting here,.

$^{241}$God has placed us in a position,.
whether it's your job, your gifts, or your creativity,.
all of those things, he has given you certain things..
Some of you guys who are sitting here right now.
have the capability of making money..
If God is giving you the skills and the gifts.
to make money, make money..
If God has given you the skills to be a pharmacist,.
be the best pharmacist that you can..
If God has given you the skills to be an engineer,.
be the best engineer that you possibly can..
Whatever you do, thrive and be successful at it..
That doesn't mean that you can't have vacations.
or have nice stuff and all that stuff..
You know what it means to follow Christ?.
It's whose life do you belong to?.
That when you sit there and look at all your skills.
and all your gifts and all your money and all your time,.
whose life does it belong to?.
Do you sit there and just say, "A part of it is mine,".
or, "Lord, here's my entire life..
"Here's all my skills that you've given me..
"Here's all my talents..
"Here's everything I'm able to achieve.
"because of your name, and I lift it up to you.".
Whose life do you belong to?.
That's one aspect of wrestling.
between two different opinions..
What's the other one?.
Here's another one that's kind of new,.
that's emerging as well, too,.
that I really wanna talk with,.
because I'm not sure, I bet you some of you guys here,.
especially living in the US and being in America,.
you probably see, it's this..
We are disenchanted by the local church..
We get frustrated by Christians..
We get frustrated by what we see.
on behalf of the church in America and over the world..
That has created a new phenomenon.

$^{281}$that we've started seeing in the last 10 years..
Let me give you some examples of what I mean by that..
And I wanna land on that..
That's one of the reasons why I wanna spend.
some time on this..
There's this whole aspect of limping.
between two different opinions, black and white..
One way in the world, one way of following Christ,.
but there's another one where we're limped.
between two worlds, and it's a little more gray..
What do I mean by that?.
Let me talk about this..
It's this whole idea..
I don't exactly like this word,.
but people have been using it,.
calling it the deconstruction of the Christian faith..
Here's what we're beginning to see a lot..
Church attendance over the last 10 years.
have continued to decline at a rapid rate..
Here's some reasons..
We're dissatisfied with many of our church options..
I'm not, and this is not Urban Voice, please,.
but this is just a general church in America..
People in America are frustrated.
by the behavior of church members..
We don't know how to disagree anymore..
We fight too much..
In too many churches, the whole idea of proclamation.
refers to justice, politics play a role in it,.
everything like that..
People are frustrated..
People are frustrated with what they see.
and the behavior of different church members..
People are frustrated because they feel like.
when they come to church, it's just a show..
We sit there, we listen to music,.
and some of these guys just have all these different pizzazz.
and lights and smoke shows and all these different things..
We feel like it's a show,.
and we're frustrated by what we see..

$^{321}$More and more, you're hearing about.
the controversy of pastors,.
the things that they do,.
the things, times when they get in trouble..
We're frustrated by how much that we read in the news.
about how bad pastors are..
The church is intolerant and exclusionary to people.
who do not believe like them or act like them..
The church in America is unwilling to engage.
in contemporary issues..
We don't talk about the role of the church.
as dealing with migrants..
We don't deal with church issues in dealing with racism..
We do talk a lot about black, white..
Sometimes we have a hard time with that..
We don't talk about Asian hate.
a lot of times in many churches..
Churches struggle with talking about LGBTQ..
We don't talk about cultural issues in churches..
There is an overemphasis on political and social agendas..
There's many belief that if you're a Christian,.
you have to be a Republican..
You can't be a Democrat..
Democrats disagree with all of these things..
And people may say,.
"Hey man, it's been really developing.
and building up since COVID.".
No, no, it's been building up for over 10 years..
Back in the early 80s and 90s,.
a few churches, a few pastors in the US.
began to experiment with a new way of doing church,.
according to different people who are out there.
in terms of Barna studies and also Gallup studies..
These ministries began to design a church experience.
that will be available to everyone.
in which it reaches out to everyone..
It would draw and make church attractive..
One author writes this,.
"The logic was simple..
If the baby boomers in the US.

$^{361}$did not feel the need to connect with God,.
then perhaps another felt need.
would draw them into the church,.
the need for community, entertainment,.
help with their children and their marriage..
They would include music, improve the music,.
support groups, drama, sermons.
that reached out and ministered to you..
And hopefully you will find God as well.".
And that was a strategy..
And that strategy worked very well in the 80s and 90s..
Evangelical churches began experimenting.
on Sunday morning services with worship experience.
with the emphasis on production..
In the 1990s, the emphasis starts shifting.
to program, art, and services.
that really excited and drew in different people..
New churches strip religious icons from their exteriors.
and interiors and gravitated toward an architectural style.
that blended the feeling of a mall,.
a university campus, movie theater versus a church..
Church could be anything.
as long as it wasn't boring..
But the problem is,.
the problem with creating a church that attracted people.
that we saw in the US in the 80s and 90s,.
created another problem..
The growth and direction of the modern church.
was fueled and influenced by money.
donated by all of these people.
who were attracted by this model.
because the plan worked..
People started coming to church.
and there started being a lot of money.
flowing into churches..
And now churches began to really fall.
in the temptation of consumerism..
And soon, many other churches,.
not only in the US but around the world,.
began to copy this model and it worked..

$^{401}$But there was a cost..
One author writes this,.
"Instead of the church becoming a place.
"where the people of God are equipped for work and service,.
"showing up to church began to be an end in itself.
"and instead of a launch pad into the world,.
"Christianity became a belief system.
"that insulated us from the world..
"Through consumer Christianity,.
"the value of God in our lives became predicated.
"on how well he fulfilled our needs.".
And now you started creating a church.
full of different consumers..
In a church as entertainment culture,.
instead of seeking to be equipped as disciples of Jesus,.
we are slowly formed into consumers and others.
who give ratings and reviews on local church performance..
But when we expect the church to entertain us,.
it limits the church's ability to challenge us..
And now here we are for what we have..
And we've seen it now,.
church attendance becoming a drop..
And in some sense, actually,.
you would think that as a drop of attendance,.
people actually are walking away from religion..
They're not walking away from religion..
They're still saying they're Christ followers..
They're still saying they pray..
They're still saying they read the Bible.
and all that stuff..
They're still religious..
They just don't go to church anymore.
because they just wanna be consumed and entertained.
and go from one place to another to another..
And when it doesn't fulfill them, they just don't go..
But you know what's going on instead of that?.
Even though they're no longer going to church.
and some of you guys may feel that way,.
this is what's happening in the American church..
Even though they don't go to church,.

$^{441}$you know what's still happening?.
They're still watching cable news..
They're still paying attention to their social media..
They're still watching YouTube.
and they're still watching and listening to their podcasts..
But the accountability, the community,.
the fellowship, the guardrails are no longer there..
And so they're listening to what they want on cable news,.
on podcast, on YouTube,.
and they now have become a different person..
Theologically, they're changed,.
but they have no accountability anymore..
That's what's going on..
There's one aspect of living.
between two different opinions,.
one foot in the world, one foot following Christ,.
but there's another part of us.
in which we're frustrated by the church experience..
We don't want anyone to tell us what to do anymore..
And if we don't like what we hear,.
then we'll move on to something else..
But yet at the same time,.
we're still consuming things that reinforces our view..
And now we're limping between two different opinions..
And that's where the current state of the church is at..
Are you limping between two different opinions?.
For me, that the second option is where I struggle with..
I have grown as a Christ follower,.
as I sit, stand here with you,.
disenchanted with the church..
I get frustrated with the church..
I've been in ministry doing this for about 20 years now..
I get frustrated by the pastors that I meet,.
not Mark, I love Mark,.
but there's so many pastors who I read when I was a child,.
who I sat there and listened to..
And when I meet them, I'm sitting there,.
seriously, you're all consumed about money and power..
The people that I, all the things that they're,.
I get frustrated when I sit in a church,.

$^{481}$I get frustrated by the politics that they talk about..
I get frustrated by Republican party..
I get frustrated by the Democrat party..
I get frustrated when we ignore.
all of these different things..
You know why I also get frustrated at?.
The conversations, even when we, as Christians,.
we dwell on the race..
I sat there with Latino pastors, African-American pastors,.
all of us are Asians..
Well, most of us are Asians,.
but I get frustrated sometimes even as Asians,.
because as Asian Christians,.
we're expected to be the quiet, obedient culture..
We're expected not to cause trouble..
I had a pastor who says to me,.
"Tommy, racism here in America is a black brown thing..
"You Asians, you have it easy..
"You guys are the ideal culture..
"You guys get everything.".
And I had to explain,.
"Hey, look, when my parents immigrated here in 1972,.
"they worked hard to get to where we are.".
And so a lot of times, you can't expect us,.
as Asian Christians, to take a back seat.
and be the stepping stone..
I get frustrated by all of these conversations..
Are you limping between two different opinions?.
Are you seeing the problems that you have?.
And for the Israelites, when Elijah challenged them,.
they said nothing..
Do you even wanna do anything about it?.
And the answer is yes..
What are you gonna do about it?.
In light of how God has gifted you and all of that,.
what are you gonna do about it?.
Or you'd be like the Israelites,.
who had nothing to say..
And you were quiet..
And maybe some of us who are sitting here,.

$^{521}$it says, "I'm good..
"I'm good, I see the problems of the church,.
"I see the problems of my faith,.
"but you know what, I'm happy with coming in on Sundays..
"I'm happy just to be sitting here,.
"and then once I walk out of here on Sunday,.
"I'm gonna do my own thing..
"I'm good, I'm good, I'm good.".
Because if that's your answer,.
then there's nothing for us to say..
Then we're done..
Go do your own thing..
But if you're sitting there,.
I wanna do something about it..
What are you gonna do about it?.
Elijah goes in, in verse 30,.
and he said to all the people,.
"Come near to me.".
And all the people came near to him..
And he repaired the altar of the Lord.
that had been thrown down..
He took 12 stones according to the number.
of the tribes of the son of Jacob,.
and with the stones he built an altar.
in the name of the Lord..
If you wanna do something about it,.
you're wrestling between,.
you're limping between two different opinions..
If you wanna do something about it,.
how do you rebuild and repair your altar?.
How do you rebuild and repair your altar?.
I wanna find one way that you could do it..
Do you have people in your life.
that speaks truth to you?.
Do you have people in your life.
that says no to you?.
Do you have people in your life.
that stretch you to become the version.
that Christ Jesus wants you to be?.
Do you have people in your life.

$^{561}$that pushes you?.
Or do you surround yourself.
with people in your life.
that tells you what you want to hear?.
Let me tell you one story..
Back in late November,.
right after Christmas,.
my friend Sharon was a paralegal.
who was diagnosed with stage four cancer.
right on her last day of work..
Her goal was to retire early..
She retired early at 48..
On her last day of work,.
she was diagnosed with cancer..
And there she had to battle for two years.
through chemo, radiation, and all of that..
On a Thursday night,.
the first week of December,.
she was slowly struggling,.
slowly just not doing well..
The doctors had told her.
she did not have time to live..
She wanted all her friends,.
different friends,.
to stay with her at night,.
and it happened that on that first week,.
right after Thanksgiving,.
on that Thursday night,.
I was assigned to stay with her..
And I walked in,.
and I actually did not want to be there..
I had just gotten back from out of town,.
did not want to be there,.
but my wife says to me,.
"You don't know what's gonna happen..
"She's not doing well..
"You need to go see her.".
So I got all my stuff,.
was gonna stay there with her that night,.
and I got to the hospital,.

$^{601}$and when I walked in,.
her face was white,.
and she was just staring out.
into the darkness..
And the nurse and doctor says,.
"She has not been doing well.".
And I said, "Okay.".
And her oxygen was dropping,.
she stopped breathing,.
she was having a hard time breathing,.
all of that stuff,.
and she got stabilized that very night..
And it was about midnight,.
and she had fallen asleep,.
and she was snoring, actually..
And the nurse comes in and says,.
"Wow, that's pretty good..
"She's snoring.".
And at 12 midnight,.
her oxygen started dropping,.
and it was about 90-something,.
and it kept dropping..
I was the only one in the room..
It's midnight..
I was the only one in the room..
And it kept dropping from 90,.
to 80,.
to 70,.
to 60..
Suddenly, all these nurses,.
all of these guys,.
there was about 20 people in the room,.
and as it was dropping to 59,.
the nurse says to me,.
"Whoever you have to wake up,.
"call them.".
Because when I lift this ventilator,.
this mask off her,.
she has three minutes before she dies..
And now I had to call everyone..

$^{641}$And there,.
she never,.
her family came,.
her friends came,.
she never made it,.
and there she died in front of me..
But the week before,.
when I saw her,.
on that Sunday,.
Sharon knew that her time was about to be gone,.
and she grabbed my hand,.
and she says,.
"I am angry at God..
"How did God allow this to happen to me?.
"How did God allow me to suffer so bad?.
"Help me understand.".
And that's where we sat there,.
and we talked,.
that sometimes we don't understand.
how God does things..
God never promises that he'll save us..
He promises that he'll be with us,.
but he doesn't promise us pain-free..
He doesn't promise that we're gonna be free of trials..
He promises that he'll give us strength,.
he'll be with us..
How do we rebuild and repair the alters in our life?.
Do we surround ourselves with people.
who are willing to tell us things.
that we don't wanna hear?.
Are we willing to surround ourselves.
with people who tells us no,.
who guides us,.
who tells us and help us see God clearly.
in the midst of all these things?.
If we're limby between two different opinions,.
whether it's with the world or with our faith.
or with what we're seeing in the church,.
that's great,.
I totally understand..

$^{681}$What are you gonna do about it?.
'Cause if you don't wanna do anything about it,.
then we don't have to do it..
There's nothing for us to talk about..
But if you wanna do something about it,.
what are you gonna do?.
How do we repair and rebuild the alter?.
And lastly, as we look at things,.
once they rebuild the alter,.
he praises one thing in verse 36..
"At the time of the offering,".
and so the bowels and all prophets bowed,.
they couldn't call their God, no one came, all right?.
They doused it with water, everything, nothing came..
And at the time of the offering of the oblation,.
Elijah, the prophet came near and says,.
"O Lord, God of Abraham, Isaac, and Israel,.
"let it be known this day that you are God in Israel.
"and that I am your servant.
"and that I have done all these things at your word..
"Answer me, O Lord, answer me,.
"that this people may know that you, O Lord,.
"are God and that you have turned their hearts back..
"Then the fire of the Lord fell.
"and consumed the burnt offering.
"and the wood and the stones and the dust.
"and licked up the water that was in the trench..
"Are you praying that God will show up in your life?".
Elijah prayed and asked the Lord,.
"Let it be known this day that you are God in Israel.
"and that I am your servant..
"I need you to show up..
"Are you praying that God will show up in your life?".
Now, here's one thing, God is always at work in your life..
God is always present in your life..
But sometimes what you and I need.
is that we're praying that God show up in our life..
Make it known..
Show us that you're real..
Show us that you're,.

$^{721}$because you know what, we need to be reminded..
Because life is hard..
It's not easy..
We get discouraged..
The things of the world drives us crazy,.
the discouragement, the things that we face.
each and every single day..
And there are times in the midst of following the Lord,.
we're gonna need strength..
For those who are sitting online.
and for those who are sitting there,.
we're gonna need some encouragement along the way..
Are we praying that God will show up?.
Because if we don't,.
we're gonna struggle along the way..
Can I read you one more story?.
Or one more email that I got?.
And I apologize that this,.
all my illustrations have been on death and struggles,.
but here's one more of an email.
that I actually received this week that actually struck me..
And it's this guy who is an inter-varsity staff worker,.
and he has been working inter-varsity.
for a while, he says,.
"Some of you have prayed for me for my whole life..
"Ever since I joined inter-varsity as a freshman.
"at the University of Illinois in 1992..
"Thank you to your prayers, love, and support..
"God's ministry through me has touched.
"over 100,000 people because of through you..
"The true extent won't be known until we reach heaven..
"And now heaven is where we will meet next.".
As I've lost my 27-year-old health battle.
with migraines, with Crohn's disease, with all of that..
He talks about all the wonderful things.
that God has allowed him to do through their support..
But my death will have been several days ago.
by the time you get this..
God will judge me by the 50 years of my life,.
not the moment of my death..

$^{761}$When I read this email,.
I actually had to read this a couple of times..
I said, "Did he just say what he just said?".
He did..
This guy who has been serving for 27 years of his life.
needed some encouragement from God..
And he didn't, he couldn't find it..
And he ended his life earlier this week..
And now he's sending an email to all his supporters saying,.
"It has been a couple of days and you will soon,.
"by the time you read this, I will be dead.".
Yes, I will have passed..
Our life, and the thing is, for a guy who struggled cancer,.
and some of you guys know I've struggled cancer,.
I know what that means..
I know what he's dealing with..
Every day I wake up, my body aches..
Every day I wake up, I have chills throughout my body..
I probably don't have what he had..
Every day I struggle migraines..
Every day I have no saliva, no tears, all that..
My head is just pounding..
I just dealt with three weeks of a sinus infection,.
and since I've weakened the immune system,.
my sinus and my headaches were just attacking me..
It just kept painful, pounding and pounding and pounding..
This guy had to deal with 27 years of it..
He lost sight, and he ended his life because of it..
We have to ask God to show up in our life..
We need God to give us strength..
The reason why we're not, we're wrestling between two,.
we're limping between two different opinions..
The reason why we don't have the strength to follow God..
The reason why we're so attracted by the world..
The reason why we get attracted.
and disenchanted by the church.
is because we don't know how to repair or rebuild our altar,.
and one of those ways to rebuild and repair the altar.
is the community of believers,.
but also praying that God will show up.

$^{801}$and that we will see God every aspect of our lives,.
because you and I are in a journey of following Christ,.
and that journey is not easy..
If anything that the story of Mount Carmel teaches us.
as we wrap up is that what it means to follow in Christ.
is that you're gonna go through tough times..
You're gonna have to have choices that you have to make..
You're gonna have to make decisions.
about what you're gonna do with my faith,.
because you can sit there and say,.
"I'm a Christ follower," and go through the routines,.
but if you truly wanna make a difference with your faith,.
you're gonna have to surround yourself.
with believers who pushes you,.
but you're gonna have to ask God to make himself known,.
to show up in your life and remind you.
of what it means to follow after him,.
and before I pray, I wanna say one thing,.
and later on as we look at that story,.
the fire comes down, and Elijah, as he prays,.
the fire comes down and burns up the burnt offering.
because it's a simple reminder.
of what happens when Jesus comes,.
because the fire could have destroyed.
and killed all the Israelites that were gathered there.
for their disobedience..
It could have killed all the folks.
that bowed for their disobedience,.
but no, the fire destroyed the burnt offering..
When Jesus came, the judgment of God.
could have been placed on all of us,.
but no, the judgment of God was placed on the sacrifice.
in the form of Jesus Christ that was found on the cross..
That's the gospel message..
That's what Jesus reminds us of..
That's why we follow after him..
Let's pray..
Father, thank you for the message of Elijah..
Thank you for how this Tishbite of Tishbe and Gilead.
continues to teach us all these wonderful lessons,.

$^{841}$and the lesson today of what it means.
to live between two worlds,.
what it means to repair and rebuild our altar,.
and what it means for you to show up.
and for us to pray that you will show up..
Give us a strength..
Let us draw that strength..
Let us see you..
In Christ we pray, amen..
(piano music).
\newpage



\section{創世記}
\label{sec:9LbyKyl1jKw}
\textbf{上主的護理 (4) 細想約瑟的故事 4 (創世記 42-43) - 陳崇基牧師}
\newline
\newline
連結: \href{https://youtube.com/watch?v=9LbyKyl1jKw}{\texttt{https://youtube.com/watch?v=9LbyKyl1jKw}} ~~~~ 語音日期: 2024-01-24
\newline
\newline
\hyperref[sec:G8td0HdJn8w]{\small{< < < PREV SERMON < < <}}
~
\hyperref[sec:index_chronic]{\small{[返順時目]}}
~
\hyperref[sec:index_scriptual]{\small{[返順卷目]}}
~
\hyperref[sec:TqQFqNVa814]{\small{> > > NEXT SERMON > > >}}
\newline
\newline
$^{1}$各位電視節目好.
在2024年我終於回到這個系列.
God's Providence.
因為去年我們講了三堂.
這個God's Providence發生什麼事.
記不記得.
就是以巴衝突.
以巴衝突又衝突了幾個星期.
現在還沒完.
越搞越大.
我原本想著5月去以色列當然是完好.
可能還是去不了.
因為如果你看到一些新聞的話.
你知道戰火已經蔓延到以色列北面的哈斯波拉.
南面你可能收到一些消息.
在紅海附近的那些叫什麼.
Hoovis.
Hoovis在也門那裡有不同的襲擊,攻擊等等.
在以色列等等不同的地方.
所以我都沒有打過很大的問號.
5月底我能不能去到呢.
我希望這個戰火能夠快點完結.
因為最主要不是我能不能去.
最主要是那裡的一些平民百姓.
以及一些婦女,兒女.
很多小朋友因為戰火的緣故.
是受到很大的影響.
願上帝去憐憫那個地方.
因為以巴的緣故.
所以就停了上帝庇護.
上帝的Providence系列.
然後就聖誕節,過年等等.
然後開年等等.
我們回來這裡就講到上帝的Providence.
第四堂了.
我打算講完這一堂再講多一堂.
就完結一個約瑟的故事.
我都不是很想中斷了.
還沒完結就由他.
我希望能夠跟進.

$^{41}$這裡講到Providence是什麼呢.
是聖經裡面一個很清楚的教導.
教導我們上帝自己的本性是怎麼樣.
並且這個教導不是一些抽空的教導.
是跟我們生命每一步.
我們人生每一步都是實用的.
因為人生裡面你想一想.
我們總是會覺得.
就算你不是信耶穌.
你都會覺得世事的發生.
或者生命裡面遇見的人事物.
有時候會令到我們感到無奈.
是我們不能夠去控制得到的.
即使你的人格是比較樂觀的.
有誰是比較樂觀的.
天跌下來當被掩蓋的人.
即使你是樂觀的.
你都會覺得生命裡面很多時候是失控的.
無論是廣大的政局也好.
或者我們自己人生的經歷也好.
所以聖經裡面告訴我們.
我們知道這個教導的時候.
可不可以給我們多一點的平安.
多一點的盼望.
讓我們能夠有目標去走.
所以要明白上帝庇護的教導的時候.
我們知道在你人生裡面.
背後是有一個總指揮.
有一個導師.
他是在幫你.
以及他安排了一切.
縱使你現在覺得失控.
但一切都在主自己的帶領.
以及在向他自己的計劃.
這個就是約瑟的故事的中心.
就是講到Providence.
我們review一下在37章開始.
約瑟的故事.
當時他是17歲.
你記得他穿彩衣那個.

$^{81}$那個是被雅各加夫.
去到偏愛的一個兒子.
排第十一.
他就因為兄弟對他不滿.
或者對他嫉妒.
就將他賣到埃及.
死不去就將他賣到埃及.
然後38章有一個小小的插曲.
講到猶大和他馬的故事.
有一點伏筆.
將來猶大會不會成為主角.
無論如何39章就講到.
約瑟在埃及波提法的家.
是一個官的家.
那裡他就被誣告.
波提法的妻子去引誘他.
然後他就保守著他的誠信.
但是因為這個緣故他就受苦.
然後到第40章.
他被監禁.
不單止是這樣.
再差一點.
到第40章就在監獄裡坐牢.
然後他就解了兩個夢.
一個是施酒長和施膳長的夢.
誰知道施酒長就忘記了他.
沒有報恩給他.
所以他也留在監獄裡.
結果過了兩年.
法老就做夢.
然後法老的內閣就不能夠將他解夢.
不知道發生什麼事.
誰知道就記起約瑟.
然後將他在監獄裡.
終於去reversal 提起上來.
讓他出監.
然後他解了法老的夢.
就是兩個七年的夢.
你記得嗎.
兩個七年的夢.

$^{121}$就是講到七年的豐年會到.
豐年之後.
就是七年的荒年.
當時約瑟是30歲.
所以你算一下17歲到30歲.
生命裡面一路走下坡.
一坡未平一坡又起.
然後到大概30歲的時候.
他才被提拔.
然後成為埃及的大官.
七年荒年之前.
就是豐年的時候.
約瑟發生什麼事呢.
41章45節.
法老將他提升之後.
將他成為大官.
法老就給約瑟起名叫做.
薩法那達巴萊雅.
哇 真是難.
給他這個名字.
約瑟不就好了.
又將安城的祭司波提.
菲拉的女兒.
阿西娜.
賜給他為妻.
約瑟就在埃及王法老面前.
是時年30歲.
接著50節.
荒年未到已先.
安城的祭司.
不說他的名字了.
他的女兒叫做阿西娜.
就為約瑟生了兩個兒子.
約瑟就為他長子起名叫做嗎哪西.
因為他說上帝使我忘記了.
一切受苦和我負的全家.
52節.
為他第二個兒子起名叫做爾法蓮.
因為他說上帝.
使我在受苦的地方昌盛.

$^{161}$我這裡想說的是.
在我們第一講的時候.
我們講上帝的聖經.
很多時候去到講一些教義.
一些教導的時候.
不是好像我們西方的那種.
列出了點的形式.
不是說抽空的說.
而希伯來的聖經.
很多時候是用故事的方法.
用故事的方法.
好像劇集一樣.
來告訴你一些背後的教義.
一些上帝的本性.
用故事來表達.
所以我們要懂得去讀.
希伯來聖經的故事.
怎樣讀呢.
因為很多時候我們打開聖經.
無論我們小組一起研讀聖經.
或者我們自己去看聖經的時候.
想起聖經的時候.
很多時候我們看到聖經的故事.
一些敘述的時候.
我們就會很自然地說.
故事裡面的角色.
裡面所作所行.
有什麼榜樣.
有什麼反面教材等等.
這些都可能是重要的.
但最重要的.
我想告訴你.
要懂得看聖經的話.
就要明白.
最主要其實不是那個角色所言所行.
不是那個人物做了什麼.
而是那個旁白.
想藉著那個故事想講什麼.
你留意到我說什麼嗎.
我們芝加哥經常有些show看.

$^{201}$有沒有去看過一些drama.
或者音樂劇等等.
舞台劇那些.
舞台劇就是那個角色.
在台上.
來跟你表達整個的訊息.
你去看舞台劇的時候.
要留意那些角色在做什麼.
但那個導演呢.
是不在的.
也沒有旁白來做媒介.
但希伯來聖經.
Hebrew Bible就不同.
Hebrew Bible很多時候.
主角是我們沒有留意的.
就是那個narrator.
那個旁白.
就好像你看這段聖經一樣.
那個旁白跟你在說話.
不是約在那裡生小孩.
在產房那裡做了什麼.
他沒有任何的對話.
不過那個旁白現在告訴你.
發生了什麼事.
所以你留意那個旁白想做什麼.
怎樣刻劃那個角色和他的際遇的時候.
去藉著這個故事.
那個旁白是主角.
去告訴你上帝是怎樣的.
去告訴你上帝的說話.
上帝的本性.
就好像看這段聖經一樣.
你看我highlight所有的東西.
什麼意思.
你看約瑟是一個怎樣的人.
當他受苦了很多次之後.
到他被提升成為埃及的大官.
這段旁白告訴我們.
法老將他的名字改了.
那個很長很長的名字.

$^{241}$薩法那特巴萊雅.
什麼意思.
在埃及文.
大概意思就是說.
上帝說了話.
上帝給他傳喚.
或者上帝祝福他.
這樣的意思.
他有一個完全埃及的名字.
安城的祭司.
或者那個Greek的city name.
就叫做Helopolis.
太陽神的城市.
就是說一個外邦的.
埃及的宗教的中心.
那裡有個祭司.
將他的女兒許配給他.
這個就是一個正正式式.
百分之一百的埃及人.
所以想一想.
描述這個約瑟.
已經不像一個希伯來人.
是一個埃及人.
他的妻子是那個祭司的女兒.
阿西勒.
Esneth.
翻譯出來.
意思就是說.
She who belongs to Nath.
Nath就是一個女的神明.
埃及的女神明.
就是說一個完完全全的.
外邦的女子.
然後她生了兩個兒子.
第一個兒子叫什麼.
Manassar.
Manassar意思就是說忘記.
遺忘這樣的意思.
Forget it.
我們年紀大了.

$^{281}$經常都沒有記性.
我老婆罵我.
怎麼罵我.
馬來西亞人.
你很馬來西亞人.
就是這個意思.
Forget.
可以是forget一些好的東西.
也可以是forget一些不好的東西.
他這裡說什麼.
上帝想我忘記我的父家.
和我一切的困苦.
Ephraim.
以法連.
就是昌盛的意思.
就是在外地的地方.
一個異地的地方.
我落地生根了.
能夠昌盛.
約瑟是一個怎樣的人.
是完完全全的一個Egyptian.
一個埃及人.
已經洗底了.
亞國家.
以色列人的身份.
希伯人的身份.
全部已經洗清了.
這個時候.
就是說以往上帝有的夢.
上帝有的計劃的彰顯.
似乎這個情況.
這個時間.
約瑟已經forget it.
已經沒有了.
但是故事當然不會是這樣.
希伯來聖經的narrator.
就告訴我發生什麼事.
就到42章.
42章你看到嗎.
就怎知道有饑荒發生.

$^{321}$饑荒發生.
有什麼.
去繼續推進上帝的計劃.
因為饑荒太大了.
影響不單止是埃及.
影響迦南地.
迦南地有誰.
就是約瑟的本家.
亞國和他那11個兒子.
還在那裡.
所以42章開始就交代.
他們要下去埃及.
去扔糧.
去拿糧.
亞國就說.
你10兄弟不如你們下去吧.
但是就減了便衙門.
就不要給便衙門去.
於是發生了很多事情之後.
在埃及一次扔糧的journey裡面.
結束的時候.
第二兒子西勉就被扣留在埃及.
然後43章就講到第二次.
他們下去埃及取糧.
我們試一試看第42章.
42章這裡我們只是outline地告訴你.
頭五節就講到亞國.
亞國就說.
大哥我們現在沒有東西吃.
我們快吃光我們的東西.
你們10兄弟可不可以下去埃及去扔糧呢.
不過就不要帶便衙門下去.
為什麼呢.
這裡經文有所交代.
42章第一節.
他見到埃及有糧.
加附亞國就說.
你們這班人在做什麼.
在彼此眼望.
我聽見埃及有糧.

$^{361}$你們不如下去吧.
於是第三節.
約瑟的10個哥哥.
10兄弟他們就下去埃及那裡買糧食.
至於約瑟的弟弟便衙門Benjamin.
亞國沒有派他和哥哥同去.
因為亞國心裡面說什麼.
恐怕他被遭害.
為什麼呢.
首先約瑟和Benjamin是同一個女人生的.
那個是誰.
是拉傑.
拉傑是誰.
是加附亞國最愛的女人.
這些大家族很多時候都很煩的.
這樣的關係.
但要知道就是.
約瑟是第十一兆.
拉傑是Benjamin最小的那個.
這兩個是他心肝寶貝.
所以他就說.
你們10個兄弟下去吧.
不要叫我最小的那個.
因為我最疼愛那個.
如果他再下去的話會變成怎樣.
他又怕好像.
開頭的時候.
約瑟就猜他去.
然後他遭害.
他以為他死歸了.
所以這次他保住這個拉傑.
然後叫這10個兄弟.
你們自己去吧.
這個是頭五節所交代的.
然後六節到十七節.
就說到這10個兄弟下去埃及的時候.
居然發生什麼事.
最詫異的地方.
就是他們和約瑟又重覆.
10個兄弟去到約瑟面前.

$^{401}$就向他下拜.
是約瑟.
因為約瑟已經是一個埃及人.
說的又是埃及話.
跟他們之間說話的.
又有一個傳譯官.
有一個翻譯的.
他又高高在上.
在寶座上面.
所以是約瑟認得那10個兄弟.
但10個兄弟就不認得他.
當這10個兄弟來到約瑟面前.
去下拜的時候.
約瑟在想起什麼.
他說忘記所有的東西.
但當10個兄弟向他下拜的時候.
他就想起開頭那兩個夢.
這個上帝一早啟示了.
說10個兄弟或亞國家.
會來求他.
叫他去救他.
但約瑟沒有表露他身份.
所以他就嚴格地去斥責他們.
然後指控這10個兄弟.
其實是迦南派來埃及的奸細.
是spies.
所以就需要求證.
他們到底是不是真的人.
結果說著說著就說了什麼.
就說問他們查家宅.
問這10個兄弟.
約瑟就問他們.
到底是什麼人.
然後你們是不是還有一個.
另外一個的弟弟.
這個就是第三段的情節.
約瑟就差派這些兄弟.
就說你們回去.
然後在對話裡面.
這10個兄弟自己用希伯來話說.

$^{441}$當年我們錯了.
因為約瑟去求我們去救他的時候.
我們沒有救他.
誰聽到呢.
約瑟聽到.
因為約瑟是雙語的.
他會埃及話,希伯來話,迦南語.
所以這裡有一個很有趣的情節.
42章18節.
這裡說的是什麼呢.
10個兄弟被約瑟還押在監牢的時候.
第三天約瑟就跟他們說好一點的話.
他說我是敬畏上帝的.
你們這樣做就可以存活了.
如果你們真的誠實的話.
如果你們真的靠得住的話.
留你們兄弟裡面一個還在監裡面做人質.
你們就先帶糧食回去.
然後先救了你們的饑荒.
再把你們最小的弟弟.
就是誰.
卞雅文.
帶到我這裡來去證實一下.
你們說的話是不是真的.
你們就不至於死.
於是他們就說好.
這樣做吧.
這樣就沒辦法了.
因為你們試過過海關嗎.
過海關是特別乖的.
當時也有海關的.
古代的社會也有海關的.
他們進入埃及的時候.
我們有一些文件告訴我們.
是有一張簽證的.
要給人的.
以前已經有了.
你現在拿著一本護照.
你的護照用了幾千年了.
都有這樣的過海關.

$^{481}$你過海關特別乖.
為什麼.
因為玉隨針板常.
他是最大的.
你有什麼事.
搞不定就死了.
這裡說什麼.
玉瑟是想試驗一下這十兄弟.
他的目的是什麼.
可能是我們閱讀這個故事的時候.
可能是報仇.
你以前耍我.
現在耍你.
死了沒有.
或者可能是他記起.
他們十兄弟.
是那些不擇手段.
想搶家產.
對一些被偏愛的人.
或者家父喜歡的人.
他們就生疾度.
所以玉瑟想試一試.
這麼多年了.
有幾年了.
起碼二十年了.
二十年之後再見這個人.
或者再見這班人的時候.
他覺得這班人的品性有沒有改變.
還是好像以前那樣.
所以他遭殃.
是否依然顧我.
你知道江山易改.
品性難移.
所以你留意龐百強在強調什麼.
誠實的人.
你們的話是否真的.
這個就是玉瑟想試驗十兄弟的目的.
再把你最小的弟弟.
帶到我這裡來.
對玉瑟來說.

$^{521}$他就想起來.
因為他最掛念的是誰.
不能忘記的是誰.
卞雅文.
他的親弟弟.
也是雅各家父最愛的.
這次是二十年之後.
首次的重逢.
見到十個人.
欠誰.
卞雅文.
卞雅文在哪裡.
不知道.
是不是他們又埋在別處.
或者把他殺了.
所以這個可能是玉瑟最關心的問題.
卞雅文現在是怎樣.
我最愛的那個現在是怎樣.
二十年之後他想找出.
但最主要他想知道這十兄弟到底有沒有改變.
四十二章二十一節.
於是這樣說的時候.
這十兄弟彼此去到傾密計.
他說我們在弟弟身上實在犯了罪.
這個弟弟就是弱失.
他哀求我們的時候.
我們見到他痛苦.
我們卻不肯聽.
所以這場苦難現在臨到我們.
現在我們有收場.
裡面是誰.
最大那個就是長子.
裡面是Rubin.
他就回答其他九個兄弟.
他說我當時不是跟你們說過.
不要害這個孩童.
就是弱失.
只是你們不肯聽.
看他的血現在向我們追討.
他一說不說尤志可.

$^{561}$說起就很大件事.
為什麼.
因為裡面這樣說的時候.
那時候我已經告訴你不要害他.
當時你又不聽.
現在血或者罪債追討我們.
他們不知道弱失在聽.
因為弱失是白lingo.
因為他們之間有傳譯官.
有個translator.
弱失轉身就離開他們.
哭了一場.
又會去跟他們說話.
就從他們中間抓住西眠.
在他們面前將他綑綁.
搞什麼鬼.
他想起.
可能他二十年裡面都是在怪長子.
在裡面出計和去到施計.
去到害他.
因為近乎文化裡面.
通常話事高是誰.
一定是長子最大那個.
不是那些弟弟.
一定是長子.
原來他二十年裡面都誤會了.
原來不是裡面.
其實裡面是想救他.
於是他就離開他們.
然後哭了一場.
抹完眼淚然後再回來.
結果就抓了西眠.
將西眠扣留作人質.
西眠是誰.
是第二仔.
不是大仔所以是誰.
近乎文化就是這樣數下來.
抄家都是這樣.
最大那個抄去最小的那個.
所以西眠就要留在埃及.

$^{601}$然後其他的兄弟就回去.
拿糧回去亞國家.
很精彩.
所以這裡二十年之後.
可能藥室就聽見他們.
大家講的東西.
和裡面所講的東西.
就想起當時那一幕.
想起當時他受創傷的時候.
有沒有想起一些往事.
令你大哭一場.
重播以往那種畫面.
當時不知道誰錯誰對.
到底始作俑者是誰.
怎樣也好.
是一個很有人性的黑話.
是一個旁白告訴我們.
扣留西眠就是第二仔.
然後九兄弟就回去.
向亞國回報.
之後四十二章結尾的時候.
就告訴我們.
九兄弟終於回到迦南.
他們就說如果我們要救西眠.
或者再拿多點糧的話.
就需要做什麼.
就要帶這個拉仔走.
如果你是亞國你信不信.
這班人是瘋子.
被人查家宅的時候.
你就什麼都說出來.
現在連我的拉仔都保不住.
所以亞國就反對.
但沒得反對.
因為是海關.
海關沒得跟他爭辯.
海關叫你拿文件出來你就要拿.
無論你feel good也好.
不feel good也好.
這個就是四十三章的場面.

$^{641}$亞國一開始是不想的.
但猶大出來.
猶大終於勸服了家主.
他說我保他.
猶大第四個兒子.
我們之前說過.
十二個兒子裡面.
頭三個已經DQ了自己.
他們不能夠成為長子名分的產業繼承人.
第四個兒子猶大.
最終他勸服了家主.
肯帶那個拉仔去.
於是他就將這個拉仔帶走.
然後第二部分就是十五至三十四節.
這個大的部分.
就是說到這班十一個兄弟.
帶著這個綿阿敏去埃及的時候.
約瑟就很奇怪.
就在家裡面去款待眾兄弟.
這個就說回二十幾年之後.
他就重遇他最親的弟弟Benjamin.
當十一個兄弟回到埃及的時候.
他就進了這個大官的家.
他們就很害怕.
也都覺得很奇怪.
很混亂.
為什麼.
因為之前那個海關很惡.
這次他再過關的時候.
居然去海關的家裡吃東西.
他請你喝咖啡.
給你吃東西.
非常之迷惘.
聖經裡面旁白就說什麼.
二十七節.
約瑟居然問他們安和.
你試過海關的海關官員問你.
你家裡有沒有病.
居然是這樣.
他說又說你們父親.

$^{681}$就是你們所說的那位老人家平安嗎.
他還在嗎.
這個嚇死人.
你過海關入境的時候.
你爸爸怎麼樣.
爸爸在家裡.
他們說.
李僕人你的父親平安.
他還在.
其實這裡已經在顫抖.
於是他們低頭下拜.
約瑟舉目看見.
同母的弟弟辯雅文.
就對他說.
你們向我所說的最小的弟弟是這位嗎.
順便說這裡是經過翻譯的.
他用埃及文講.
然後用翻譯再講.
又說我兒願上帝賜恩給你.
第三十節.
約瑟愛弟之情激動.
就急忙找個地方又再去行.
進去自己的房間哭了一場.
洗了臉之後.
勉強的忍住.
就說開飯.
開飯.
沒有見這個親弟弟辯雅文.
已經有二十多年的時候.
以為他已經遭殃.
以為他已經被埋.
誰知道親眼見到他.
接著這個旁白就告訴我們.
在海關的家裡有個筵席.
三十二節.
他們為約瑟單獨的去到擺了一席.
為那些人.
即是那些埃及人又擺了一席.
又為和約瑟同吃飯的.
那些人是希伯來人.

$^{721}$又為和約瑟同吃飯的埃及人另外擺一席.
因為埃及人和希伯來人是不可以吃飯的.
因為加盟人是不喜歡和他們吃飯的.
三十三節.
兄弟被安排在約瑟面前坐席.
從最大的去到最小.
長幼的次序.
這些人彼此的感到詫異.
是不是感到詫異.
誰知道他們是誰大誰小.
就是埃及的海關官.
就是他知道.
但又不知道他是誰.
約瑟就在他面前的食物分給他的弟兄.
但是辯雅文所得比別人是多五倍.
他們就一起喝酒.
然後一同唱音.
約瑟將他面前的食物分給他的重弟兄.
這個似乎第一次再有伏筆告訴我們.
約瑟就是成為埃及國家的救主.
是他的食物分給他們.
以至他們能夠有東西吃.
救到埃及國家.
辯雅文就多得五倍.
這裡又是另一個測試.
最小的被偏愛的.
這10兄弟會不會窒到.
會不會好像以往家父被約瑟的採依.
他們窒到和他起降.
你想想一起吃東西.
最小的那個.
大家一塊牛扒.
他有五塊.
很明顯這裡又是另一個測試.
現值這裡有多少位呢.
旁白告訴我們有多少位.
希伯來人一位.
11個兄弟.
埃及人另外一位.
然後約瑟一個人一位.

$^{761}$為什麼呢.
不是因為他大官不能夠和這些下人一起吃飯.
如果是的話.
旁白已經交代了.
他交代了希伯來人和埃及人是不能夠同台的.
明白了.
你記得我們之前說什麼.
約瑟在埃及.
是一個完完全全的埃及人.
完全埃及化了的希伯來人.
為什麼他自己一位.
這些似乎就是旁白想說的一些東西.
約瑟的位置是怎樣.
可能他還在糾結.
在這個時候.
在糾結到底我是一個埃及人.
還是一個希伯來人.
我們可以和他們復和.
還是應該和他們劃清界線.
所以他自己一台.
可能他正在處理以往的創傷.
我們人就是這樣.
我們要處理我們以往的創傷.
總是需要時間.
解怨不是按個按鈕.
是要慢慢來的.
就算解怨.
就算要去原諒兄弟.
都是二十年多的恩怨.
所以要調整一下兄弟之間的關係.
這個就是我之前看過.
就是你看聖經的時候.
你看的不是那些角色那麼簡單.
而是一個旁白用角色想和你說什麼.
你明白嗎.
不是說我們有什麼榜樣.
有什麼反面教材.
而是一個旁白的說話者.
想說什麼.
他有什麼目的.

$^{801}$他就說上帝一步一步帶領著十二個兄弟.
要整個的亞國家能夠復合.
用了很長很長的時間.
用了很多不同的際遇.
很多很壞的際遇.
很多低谷的時候.
饑荒等等.
上帝是一步一步推展著.
雖然人已經覺得他已經忘記了.
馬來西亞法.
人覺得這些往事已經要去除掉.
Erase.
Delete.
但是上帝沒有忘記.
這個就是我們很多時候看約瑟的故事.
我們有一點掙扎.
就是覺得人之常情.
就是上帝如果有能力的話.
應該要叫約瑟在他的苦難裡面.
能夠馬上得解救.
他坐牢將他提拔出來.
他在坑裡面將他提拔出來.
他在埃及將他帶出埃及.
上帝絕對有這樣的大能.
上帝絕對有這樣的能力.
可以立刻馬上將約瑟解救過來.
但是如果我們想馬上要這樣做的時候.
我們就不是在想Providence.
也不是在想我們人生的常規.
因為人生常規上帝帶領的步伐.
往往不是馬上問題得到解答.
苦難得到拯救.
因為如果是馬上將約瑟解救的話.
早年將約瑟已經解救出來了.
你就不會看到什麼呢.
你就不會看到十兄弟的改變.
因為以往他們作惡多端.
他們是殺害親弟的.
出賣他的.
這班人怎麼辦.

$^{841}$如果是立時解救約瑟的話.
你只會看到約瑟脫離苦難.
That's it.
但其他的十兄弟他們的生命有沒有改變.
如果生命改變的話.
就一定要逐步.
要有時間.
或者這麼說.
如果將約瑟當時馬上去解救的話.
你就不會看到這個局面.
你只會看到一個人得拯救.
而見不到上帝更大的計劃.
上帝更大的計劃就是要亞國全家能夠復合.
而且藉著約瑟來救萬國.
這個生命裡面一個很大的謎.
你可能正在經歷一些事情.
經歷過一些事情.
你會覺得很奇怪.
為什麼上帝不去立時救我.
生命裡面很多時候要等.
要時間.
早前我看到植物.
不是因為我喜歡養植物.
我每次養植都會死.
但AI突然給我一些東西.
我就看看.
講到竹.
這個竹.
竹是怎樣的呢.
竹很高.
講到竹的生長是怎樣.
他說頭三年竹是不會長高的.
用了三年時間生根.
你要等三年.
然後等了三年之後.
竹就開始紮好根.
它就每天可以高三呎.
每天高三呎.
即是每小時高一至兩吋.
可以一直這樣飆上去.

$^{881}$但頭三年你好像是白等.
其實就不是白等.
You are not wasting time.
Time is needed.
時間是需要的.
你想想我們生命裡面是不是有很多東西.
我們覺得好像白等.
無論你是過著逆境.
一些很不順利的事情.
一些很苦難來說.
好像在白等.
但在上帝的prophets裡面.
是不會白等的.
Waiting is not wasting.
等待不是白白花時間.
兩個星期前我病了.
沒有聲音.
三天沒有聲音.
我都未完全recover.
現在都做好肝.
我學到什麼叫等候.
有什麼不滿.
對我太太有什麼不滿.
先不說.
多等你幾天.
生命裡面有很多的waiting.
But it's not wasting.
約瑟等了二十年.
不是白等.
原來這個時候才見到.
上帝有一個更美的傳盤計劃.
這些不是叫做事後孔明.
這些不是叫做自圓其說.
這些叫做Hindsight.
Insight.
很多時候回看.
後見之明才是智慧.
明白我們生命裡面的取態是怎樣.
因為上主是控制一切.
我們結束的時候.

$^{921}$我們回看衛斯敏特信經.
這個十七世紀.
一個很好的很傳避的信經.
用一個西方的方法來告訴我.
不是說故事.
他spell out到底聖經裡面怎樣說providence.
神的庇佑是怎樣.
是彰顯上帝傳避的能力.
莫測的智慧.
無限的善良.
彰顯到一個地步.
涵蓋人的罪.
包括第一次的墮落.
天使的犯罪.
人的犯罪.
但是神不是僅僅允許罪發生而已.
而是藉著祂的全知和全能.
來限制支配和管轄罪.
這一切的安排.
都是達成祂最至聖的目的.
雖然如此.
罪是出於受到者.
不是出於上帝.
上帝是至聖結至公.
而絕對不是也不會是.
罪的創始者和贊同者.
但這些遺思問答.
最精彩的地方是什麼.
是那些形容詞.
選得很仔細的形容詞.
莫測的智慧.
Unsearchable wisdom.
上帝管理萬有.
不一定是按照我們這樣出牌.
按照祂.
莫測的智慧.
在我們心靈裡面都是一樣.
當我們了解這件事的時候.
不代表我們能夠知道將來會發生什麼事.
每一步我們都知道.

$^{961}$這是上帝隱密的智慧.
但當我們明白這件事的時候.
起碼讓我們心靈裡面的取態.
我們能夠有安穩.
因為正如這裡所說.
上帝的庇佑.
是要我們做什麼呢.
是鼓勵我們.
愛我們的上帝.
以致我們能夠得到生命裡面的平安和安慰.
最近我看了一個網上的祈禱.
我覺得這個祈禱可以作一個結束.
你問一問你自己.
可能你過去有一些很不順利的事.
或者現在仍然覺得一些很不順利的事.
你正在過渡.
上帝好像聽禱告又不聽禱告.
不按照你的答案.
這個禱告可以作一個結束.
既然這是一個莫測智慧的上帝.
帶領著安排萬有.
天父對你說什麼.
他說孩子我沒有馬上成就的事.
不是因為我忘記.
而是我正在跳渡萬有.
我沒有按照你所求的成就.
不是因為我做不到.
是我要給你更好的.
在你還沒有開始禱告的時候.
我的工作已經再進行了.
你的禱告我從來不會遺留.
你的需要從來不會忽略.
我們同心低頭祈禱.
在祈禱當中我們交託.
我們的人生.
再一次來到上帝面前.
我們知道他用莫測的智慧.
傳遞的良善來帶領我們.
多謝你不是按照我們的智慧來帶領我們.
也不是按照我們的時間來進行.

$^{1001}$因為你的智慧高過我們的智慧.
主要我們今天來到你面前.
循心的信求.
讓我們知道你是那位有原備能力.
莫測智慧.
知性知潔的上帝.
你有知性的目的.
最好的是你會給我們.
甚至你自己的計劃.
你自己的安排.
你的庇佑.
是高過一切我們所作所言.
甚至人類有的挫犯.
你都可以去限制支配和管轄.
一切的安排都是要達到你自己.
在我們身上知性的目的.
主要我們來到你面前.
信福信教.
在我們心裡面我們有安穩.
在我們生命裡面有不測.
在我們人生裡面有很多混沌.
但我們知道在你裡面能夠有平安.
我們生命裡面人生裡面.
大大小小作總指揮的神.
並且你是美善的主.
祈禱奉祝明教.
願上帝賜福各位.
下星期跟大家一起事奉上帝.
拜拜.
\newpage



\section{}
\label{sec:TqQFqNVa814}
\textbf{Atomic Wisdom 3 - Rev. Mark Chan}
\newline
\newline
連結: \href{https://youtube.com/watch?v=TqQFqNVa814}{\texttt{https://youtube.com/watch?v=TqQFqNVa814}} ~~~~ 語音日期: 2024-01-30
\newline
\newline
\hyperref[sec:9LbyKyl1jKw]{\small{< < < PREV SERMON < < <}}
~
\hyperref[sec:index_chronic]{\small{[返順時目]}}
~
\hyperref[sec:index_scriptual]{\small{[返順卷目]}}
~
\hyperref[sec:WYvLrPDJi5g]{\small{> > > NEXT SERMON > > >}}
\newline
\newline
$^{1}$- Thanks, Kelly..
Well, why don't we say hi to your neighbors?.
I know some of us are out of town.
and cannot be here this week,.
and there are people who have gotten.
to some minor or major car accident.
in the past week or two,.
and our prayers go out to them.
and to continue to pray for them,.
especially Phoebe..
Phoebe's still on break.
because of the serious car accident..
She's fine, but she does need more time.
to recuperate, all right?.
So let's keep praying for them.
and praying for our church ministry..
Today, we are going to finish our series.
on atomic wisdom, okay?.
Atomic wisdom..
This is part three, okay?.
Let me move this around..
Part three..
What we are doing is that we take.
one of the best seller out there,.
and the book is called "Atomic Habits," all right?.
Some of you might have heard of this book,.
"Atomic Habits.".
It's a very popular book..
It talks about behavioral psychology..
It talks about how we build good habits,.
and we ask the question,.
what kind of conditions are necessary for life change?.
What kind of conditions are necessary.
for the probability of achieving your goal?.
In the new year, some of us might have some new goals, right?.
You may have new goals to be on diet.
or do more exercise or be more healthy,.
something like that..
We try to draw insight from this book.
called "Atomic Habits,".

$^{41}$and it's already translated into different language,.
Chinese included,.
and it draws so much insight from behavioral psychology.
and trying to help us to set up a system,.
how to set up our life in general.
to have good habit, to have healthy habit,.
and we wanna combine this so-called secular insights.
into biblical wisdom,.
and that is on the right, the book of Proverbs,.
and if you wanna find the book of Proverbs in the Bible,.
it's right in the middle after the book of Psalms..
The book of Proverbs is the part of the wisdom literature,.
and this is part three,.
so we're going to reveal what are the principle.
that are necessary for life change,.
and number one, we talk about 1\% power, right?.
Remember, if you do a tiny change every day.
after a whole year of 360-some days,.
you will receive or you will achieve your goal.
in 30-some fold multiplication,.
and do not estimate the tiny power.
of every tiny change that you do every single day,.
and principle number two is incentives.
or disincentives are effective,.
incentives and disincentives..
If you set up some goal in your life,.
for our church, for our fellowship,.
probably we're gonna read the Bible a little bit more,.
or we wanna pray more,.
we wanna build up some of the good habits,.
and we do need to set up some incentives and disincentives,.
but tell you the truth..
I finished this series in the 10 o'clock service,.
and after the part three, after this one,.
a brother came up to me, and he said,.
thank you for your message,.
thank you for encouraging us to set up a good habit,.
good life, and he pulled out an envelope from his pocket.
and gave me \$50, because he heard the first part of message,.
and he wanna build up some good habit,.

$^{81}$and he wants to have some incentive.
or disincentive in his life,.
so he was encouraged by the message, by the insight,.
and he built up some good habit,.
and after part three of the series,.
he came up to me and gave me \$50, yeah,.
praise God, right, praise God, \$50..
Anyways, so you have to think about life changing,.
it's necessary to have some benefits, some incentives,.
right, you go to the gym, you get some incentive,.
and you're gonna praise yourself,.
or you wanna uplift yourself,.
so you have system with incentive or disincentives..
Number three is specifics are better than generals, right,.
set up a goal, set up a system,.
which is as specific as possible..
Number four is that our community is important,.
because the people that we hang out with,.
the people that we see a lot,.
will have influence on us, so the community is important..
Number five is the power of default, default, all right,.
so you set up your habit, and you say,.
this is the priority, and this is the default of my life..
Every single morning, what are your defaults?.
You wake up, and you brush your teeth,.
wash your face, right, these are your default, right,.
so in order to have good habit in our life,.
you have to have, you have to understand.
the power of default, the effect of default, all right,.
so the assumption to make this habit possible..
Number six is making it personal, all right,.
making it personal, either you find a friend.
who is keeping you accountable,.
making it personal will make it even more possible.
to have your new habit set up..
So number seven now, number seven.
is the power of environment, all right,.
so you need to understand, if you want your life to change,.
or you want a kid's life change,.
we need to set up a good environment, good environment,.

$^{121}$the physical environment have so much impact on us,.
all the people around us, all right..
Many, many years ago, many years ago,.
many, many years ago, I used to have a TV in my bedroom,.
a TV in my bedroom, somebody gave me a TV, right,.
so I said, where should I put it?.
I put it in my bedroom, right,.
you think it's a good environment?.
No, right, no, it was so bad, it was so bad,.
so I took it off, right, because when we had the TV.
in the bedroom, you know, it disrupt my sleep, right,.
and my kids, when they were so much younger,.
they come into my bedroom all the time,.
because they wanna watch TV, or lying on the bed, you know,.
so I remove it, because it was a bad environmental setup..
So behavioral psychology here, study finds that.
how people maintain a diet,.
how people maintain a good diet..
As you know, in the United States, we eat too much, right,.
we have so much to eat, we have so much to eat,.
we always eat too much..
Well, study finds that it is simple to maintain a diet,.
to eat a little bit less, and the method is very simple..
You make a habit before you have a meal,.
you drink two glasses of water, right,.
you drink two glasses of water,.
then you will eat a little bit less, right,.
that makes sense, right?.
Many, many years ago, when my kids were still so much younger,.
five year old, six year old, you know,.
sometimes we, in Chicago, we drive all the way to Toronto,.
right, to have some good Chinese food, right, yeah..
Here, the Chinese food is fine too,.
but you know, you go in Toronto,.
you got so much better, right?.
So many years ago, we go to Toronto for family trip,.
because the kids were so young,.
they did not know how to appreciate good food, right?.
The expensive food, they don't know how to appreciate it..
So we, as good parents, we bring the kids first.

$^{161}$to the bakery, so we buy them a bun or so,.
and fill up their stomach,.
and then we go to a good restaurant, right?.
So we adults can enjoy it, right?.
That sounds like a very, very good insightful wisdom, right?.
We save a lot of money, right?.
Because we only buy good food for ourselves, right?.
Because they don't know how to appreciate it,.
why spend the money on them, right?.
So, the study finds that if your stomach.
is a little bit full, then you go to the meal,.
then there's a bigger chance for you to maintain a good diet..
And the flip side is telling us that new habit,.
you have to make it easy and natural..
And there are some behavioral psychological study.
done for cafeteria of a hospital in Boston..
And the study is trying to find out.
how to encourage the staff to get healthy food.
in the cafeteria..
And the plan is very simple,.
is to change the layout of the cafeteria..
Change the layout of the cafeteria..
Simply by changing the layout, you get so much effect..
Now, the layout change has to do with bottled water.
being put at the entrance of the cafeteria..
And fountain drink, where do you put it?.
You put it all the way back, right?.
Make it difficult for you to get fountain drink, soda..
So, after three months, the study finds.
that water consumption increases 25.8\%,.
and soda consumption drops by 11.4\%..
Isn't that amazing?.
It's the same way..
I mean, many of you go to office and work in the office..
Sometimes you find office cafeteria layout intentionally.
to discourage you to get snack food..
Snack bowl or candy bowl or junk food.
is being placed in the breakout room..
And then the healthy food or water is being put.
on some place where it's convenient..

$^{201}$What goes wrong in the last few years.
is that people working from home now, right?.
And that is bad, right?.
That is bad..
Because in the pandemics, everybody works from home,.
works at home, and there's a huge temptation, right?.
You work from home, right?.
With your computer, and kitchen is right next to you,.
and you get hungry, right?.
You get boring, and then you go to find some potato chips.
or some snack food, right?.
And that is unhealthy..
And actually, the biblical wisdom has so much to say.
about our environment, especially about the people.
that we hang out with..
They have effect on us..
Proverbs 27, 17..
27, 17..
It says, "As iron sharpens iron,.
"as one man sharpens the face of his friend.".
It says, "If you surround yourself with good friends.".
Now, in the immediate context of chapter 27,.
it talks about bad relationships..
So 15 or 14, verse 13, talk about bad relationship,.
bad community..
It will have a negative effect on your life..
But here in verse 17, the wise teacher is telling you,.
if you surround yourself with good and spiritual people,.
you will grow, you will grow..
If not, your life will decline..
You will slowly drift away from God..
You will slowly drift away from the fear of the Lord..
And it tells you, nobody can be a Christian alone, right?.
No one's an island..
Everyone is connected, and that's why you need to belong.
to a spiritual community, a home church..
And that is important..
You cannot do Christian from home, all by yourself, right?.
And the biblical language for Hebrew is very interesting,.
because in Hebrew, here it tells you a piece of iron.

$^{241}$sharpening another piece of iron, right?.
And slowly, you will increase the value.
of both pieces of iron, right?.
The iron, one piece by itself, it has value..
However, if you rub it with another piece of iron,.
both of them will get sharpened and becomes a knife, right?.
It becomes a knife or some tool,.
and the value of the piece of iron will increase..
And what it said about, the funny thing about Hebrew.
is that in Hebrew, the word faith, faith,.
is actually used to describe the blade of a knife..
The blade of the knife is actually the face of the knife..
So the wise teacher is saying, as iron sharpening iron,.
as one friend sharpens the blade of the other friend,.
the face of the other friend..
So no one can be a Christian alone..
You need to plug in to some sort of community.
or small group or some local church.
in order for you to grow spiritually,.
to grow your relationship with God..
So that is the power of environment..
Principle number eight is habit bundling or habit stacking..
Now what the psychologist is telling us.
is that if you wanna build some new habit,.
you need to tie the new habit to some established habits..
Some habits that you are doing like your second nature..
How many of you drive a stick shift?.
There you go, all right, here's the man..
I drive a stick shift for so many years..
So even today, I get into a car,.
my left foot will do something..
I will try to find the clutch..
And then there are times,.
because I have automatic transmission car as well,.
so sometimes I jump into the car.
and then try to find the clutch..
There's no clutch there, right?.
Because it becomes a second nature for me, right?.
So that's established habit..
So the psychologist is telling us.

$^{281}$if you wanna build new habit,.
you have to attach it or tie it with some established habit..
And the easiest way to develop this.
is to think about, okay, established habit,.
and then I wanna add something to my established habit..
Every morning or every night, you brush your teeth..
When is the best time to floss your teeth?.
It's after brushing your teeth, right?.
Brushing your teeth is already an established habit, right?.
I hope that's true to you..
It's established habit..
It's kind of, you don't have to think about it..
But if you are not flossing your teeth regularly,.
the best way, the behavioral psychologist is telling us,.
is to stack it onto some established habit..
If you check the wisdom literature in the Bible,.
it actually talks about stacking,.
something adding to some of the habits.
that you may already have..
Here, chapter four, verses 20 and following..
The father is teaching the son to have wisdom..
So he says, "My son, pay attention to what I say..
"Listen closely to my words..
"Do not let them out of your sight..
"Keep them within your heart,.
"for they are life to those who find them,.
"and health to a man's whole body..
"Above all else, guard your heart,.
"for this is the wellspring of life.".
You see the stacking here?.
One of the way you appreciate the book of Proverbs.
or wisdom literature is that there are short sayings,.
and sometimes one short sayings has two short verses..
Verse B is stacking on verse A, all right?.
It's like staircase, A and then B, A and then B..
Sometimes A, B, and C..
Here, it says, "My son, pay attention.".
Pay attention..
It's actually a military word..
It describes soldiers, you know,.

$^{321}$standing, guarding on the top of watchtower..
And what the soldier is doing is actually paying attention,.
paying attention to noise from the enemies..
Maybe there's approaching army,.
approaching horses to the city,.
and you pay attention to it..
So he says, "My son, pay attention to what I say..
"Listen closely to my words.".
Pay attention, listen closely..
Do not let them out of your sight..
Put it right in front of you..
Keep them within your heart, right?.
Stacking one on the other..
Because they are life, right?.
They are the wellspring of life..
I remember I told you when I was in college,.
I failed the first test of my chemistry class,.
and then the following nine tests, I built a new habit..
After the lecture, after the class,.
I don't go with my classmates to go out to Starbucks, right,.
to have a good time there..
No, I go immediately to the library.
and to review the materials..
And I actually use less time to ace the test..
All my classmates, they have fun, they have good time, right?.
But at the end, I ace all nine tests..
And that is habit stacking..
You stack one new habit to your established habit..
And that is habit bundling..
And principle nine is to make your goal attractive..
Make it as attractive as possible..
Use different wordings to describe your goal,.
to describe your system, to describe your new habit..
If you use attractive words, you will tend to change..
You go to the gym, you have to do some exercise..
If in your mind, if in your mind,.
you use some negative description.
to describe that new habit, you will bound to fail, right?.
Oh man, I have to go, I have to drag my body to go, right?.
But if you describe your new habit in a nice way,.

$^{361}$there is a bigger chance for you to achieve it..
You know, in the United States hotel rooms,.
they are trying to find ways to encourage the customers.
to reuse the towel..
That will save a lot of money, right?.
Because sometimes you go to a hotel, you stay there,.
and then you use the towel only once or just part of it,.
right, and then you throw it on the ground, right?.
And then you use, and after that, you use a new towel, right?.
And by the end of the day,.
you probably use three or four towels, right?.
Which are unnecessary..
And the hotel manager is trying to do a study..
It says that if you wanna encourage guests.
to reuse their towel, you describe reusing towel.
with some good words..
As you can tell here, it says,.
dear guests, every day, tons of water, gallons of water,.
are used to wash towels that have been only used once..
And this hotel is committed to conserve natural resources,.
and this is how you can help, right?.
If you put it back on the rack, I'll use it again..
If you throw it on the floor, please exchange it, right?.
And the study finds out, because of this wording,.
even at the wording, save the planet, right?.
You don't have to do anything to save your planet..
You just reuse your towel, right?.
It makes you a hero, right?.
The study finds that this means it's effective..
And then the end, the result is 75\% of the guests.
in the hotel, reuse their towel more than once..
Because you describe in such a way to save the planet,.
right, we do this together, you establish some social norm.
for people to conform to it, right?.
And the wisdom literature has so much to say about this..
It describes wisdom with beautiful, glorious terms..
Chapter four, wisdom is supreme, therefore get wisdom..
Though it costs all of you, get understanding..
Esteem wisdom is being personalized as a woman..
Esteem her and she will exalt you..

$^{401}$Embrace her, she will honor you..
It is a garland of grace, it's a crown of splendor..
If you get wisdom, you get crown of splendor, right?.
Supremely, you know, super positive terms to describe,.
you know, living a life of wisdom in front of God..
Developing our spiritual life,.
developing our close relationship with our God, our creator..
Positive terms, you know, positive effect.
of the descriptors that we use..
Some years ago, there was a famous,.
a very famous speech therapist..
Her name is Nina Rustin, Nina Rustin..
She is a British therapist to help children.
how to correct their pronunciation,.
how to correct their speech..
And what Nina Rustin does is more than.
the typical therapist will do..
A typical therapist will correct the pronunciation,.
the saying of how they pronounce terms.
and how they make sound for the children..
And that is necessary, but Rustin does more than that..
Rustin call on the whole family of the child.
into her office and she assign assignment.
to each of the family members..
And every family members, she said,.
"You have to do something in your family..
"And what you do is you have to go home.
"and every single day make it a daily habit.
"to praise one of the family members.".
To find something specific,.
to praise one of the family members,.
maybe your sibling, maybe your parents..
But you find something specific to appreciate,.
to thank, to praise..
And then the one who is receiving the appreciation.
has to say something to receive it..
You can't say a typical Asian way, right?.
No, no, no, no, I'm not good, right?.
We always pretend to be humble, right?.
This is our Asian culture..

$^{441}$The one who is receiving has to say something.
in positive term..
So for example, mom can say,.
"John, I love how you share with Matthew.".
And John cannot refuse it..
John has to make it happen as well..
John has to say something like,.
"Thanks, I appreciate the compliment..
"It made me happier, thanks mom.".
And this dialogue has to be a daily habit.
in the whole family..
And Linda Woodson finds that.
if you wanna improve a child,.
improve her speech,.
you cannot do it alone, right?.
You cannot just fix the child..
You have to have the entire family get involved..
And by doing so, the entire family culture has changed..
We appreciate each other and encourage each other.
and build up the confidence of one another..
And especially, it will help the child.
who is struggling with speech..
And what's funny is that right after the therapy,.
not only the child, speech has been corrected..
The husband and wife relationship has improved..
And they come back to Linda Woodson and say,.
"Thank you, thank you for improving our marriage.
"and our family.".
So that is the principle..
It is common sense, right?.
But to make our goal, to make the system attractive,.
make it positive, will increase the possibility.
of achieving our goal..
The last principle, number 10..
It is to establish system, not goal..
Think about not goal, but the system achieving the goal..
It's not the end that you most concerned about..
It's the means to the end that is most important..
So here I want to quote a long quote from Atomic Habits..
And I only have the Chinese, but what it says is that,.

$^{481}$goal and results are not so much important..
Why?.
Because he gives the example of Olympic athletes..
All the Olympic athletes are trying to get the goal, right?.
Every one of them has the same goal, to win,.
to get the goal..
So what is the difference between a winner and loser?.
The winner has set up a system with tiny habits.
of training himself or herself.
in order to get the gold medal, right?.
That is the only difference..
The difference is not at the goals that they set..
It's the system, it's the habits that they have..
So all of this interesting, insightful psychology.
is telling us how to build a good life..
And actually if you find, if you read the book of Proverbs,.
it has so much to say about how we should set up.
a good system in our life, right?.
Now in the book of Proverbs, the first nine chapters,.
first nine long chapters,.
are the introduction of the whole book..
The introduction of the whole book..
And this introduction is very interesting..
It introduced two characters..
The first character we can call lady wisdom,.
or woman wisdom..
So wisdom is being personified,.
being described as a person, as a lady, as a woman..
And the other character, the second character,.
is lady folly, or woman folly..
That's very interesting because you find wisdom,.
it's like you look for a good wife, right?.
If you reject wisdom, if you love folly or foolishness,.
you will find yourself marrying to a wrong woman, right?.
That's very interesting, but if you consider the context,.
you understand why it is being described in such a way..
Because the book of Proverbs is teaching of wisdom.
from the father to the son..
When the son is growing up,.
one of the main goals in his life.

$^{521}$is to establish a good family, right?.
To find a good wife..
So the father is using this means, this metaphor,.
to teach the son what is most important,.
is to find wisdom in God..
In chapter nine, verses one to three,.
it says wisdom has built her house,.
he has held out his seven pillars..
The wisdom, or lady wisdom, has prepared her meat.
and mixed with her wine..
She has also set her big table, big table of good food,.
sent out her servant, her mates,.
and she calls from the highest point of the city..
Is calling out to the young men, the young people,.
to come to her table, come to her house..
She's calling out and yelling out.
from the highest point of the city..
That is what the lady wisdom is trying to say..
And then also, the lady folly has so much to say,.
and what she said is actually similar.
to what wisdom wants to say..
Chapter nine, verses 13 to 14..
The woman folly is loud,.
she's undisciplined, without knowledge..
She sits at the door of her house,.
again prepare a table,.
on the seat at the highest point of the city..
So the lady wisdom is calling out.
from the highest point of the city,.
and then the lady folly is calling out.
from the highest point of the city..
And what's interesting about this phrase,.
at the highest point of the city,.
is at the top of the mountain..
Now in the biblical culture,.
when you see a mountain,.
mountain top,.
you usually think of the temple,.
because all the temples are built on the Acropolis,.
the highest point of the city..

$^{561}$So here lady wisdom at the highest point of the city,.
calling out to those who need wisdom,.
and the lady folly also calling out from the same place,.
from the highest point of the city..
What does it mean?.
It means wisdom is not just secular,.
self-improvement pieces of advice..
Wisdom is actually God himself, right?.
God is calling out from the highest point of the city.
to his temple, his house,.
to invite us to build relationship with him, right?.
But at the same time, the false gods,.
the gods of the culture,.
the gods out there, the false gods,.
are also calling out from the same place,.
from their temples,.
to entice people or to tempt people to follow them..
So if you come to this point,.
you understand the book of Proverbs.
is not just talking about self-improvement strategies..
It's not just being better..
It's not just becoming a better man or woman..
It's actually about worship,.
about your relationship with God..
Remember the fear of the Lord is the beginning of wisdom..
The fear of the Lord is the proper relationship.
between you and God is the beginning,.
is the foundation of wisdom..
So it's not just talking about self-improvement..
It's talking about us worshiping the right God,.
the true God..
So I hope that we will end this series.
by this appeal,.
that God is calling you from his temple,.
from his dwelling place,.
to invite you to build a right relationship with him..
And the good news we have from the Bible.
is that we, as sinners,.
we are being prepared for the salvation of Christ..
And Jesus came in our place,.

$^{601}$in the place of sinner,.
to die for us, to sacrifice himself for us on the cross,.
in order to bring us back and redeem us.
so that we can be the children of God..
And I hope that this will be true to each and every one of us.
that we will build on the beginning of wisdom.
to get to know God in Christ..
Everyone who has accepted Jesus Christ.
as his personal savior and the Lord of his life.
will be related to God,.
totally forgiven,.
and stand righteous in front of the Holy God..
And that is the beginning of wisdom..
And I hope that this will be your decision today.
to get to know God in Jesus Christ,.
and you will begin a life of wisdom..
Let's pray together..
Father, we love you because you came to save us..
You came to redeem us in Christ Jesus..
Thank you for the way of wisdom.
that is laid before us..
Thank you for bringing us back to you..
And our Lord Jesus is the way, the truth, and life..
So we come and approach you without fear,.
approach you afloat,.
fearlessly,.
all because Jesus Christ has done so much for us,.
forgiving our sins and accepting us just as we are..
Thank you for what Jesus has done for us.
to die on the cross, die for our sin,.
so that we can live..
And thank you for inviting us into your relationship..
And I pray that each and every one of us.
will begin with our relationship with you,.
to invite Jesus to be our savior,.
that we can call you as our father..
Thank you, and make us new..
Help us to build new habit of wisdom,.
to live wisely,.
to live in the fear of you,.

$^{641}$and to live wisely, to walk wisely,.
to bring glory to you..
In Christ's name, we pray..
Well, God bless you all,.
and I hope to see you next week,.
and we continue to worship our God..
Stay behind for our lunch,.
and after our lunch, around one o'clock,.
we will split into different groups for Bible study,.
and we'll continue to study the Gospel of Mark..
All right, God bless..
\newpage



\section{}
\label{sec:WYvLrPDJi5g}
\textbf{The Chosen (Sermon Series) \#1 - Rev. Mark Chan}
\newline
\newline
連結: \href{https://youtube.com/watch?v=WYvLrPDJi5g}{\texttt{https://youtube.com/watch?v=WYvLrPDJi5g}} ~~~~ 語音日期: 2024-02-06
\newline
\newline
\hyperref[sec:TqQFqNVa814]{\small{< < < PREV SERMON < < <}}
~
\hyperref[sec:index_chronic]{\small{[返順時目]}}
~
\hyperref[sec:index_scriptual]{\small{[返順卷目]}}
~
\hyperref[sec:QExp8HlGgEo]{\small{> > > NEXT SERMON > > >}}
\newline
\newline
$^{1}$- All right, thanks Kelly..
I know some of you cannot be here.
and some of you are online being with us.
because they got sick and we wanna pray for them..
And some of those in our community also out of town,.
they went back home and visit their own country.
and they will be back, all right?.
So let's pray..
Father, give thanks for giving us this community,.
this spiritual family..
We can study your word and we can worship you together..
And we pray for those who are not with us today..
Maybe they are sick,.
maybe they're going to their home country for visitation.
and I pray for safety, I pray for healing..
And especially those of us who are struggling.
with our own personal challenges,.
remind us God, remind us that you are our great God.
and you are greater than all the issues.
and all the problems and challenges that we are facing..
So we come before you and to get to know you,.
come before you to worship you and exalt your name.
because we know that you are a great God.
and we can trust in you..
And increase our faith, increase our love for the community.
that we can grow your church together..
We can glorify your name together in Christ name pray..
Well, as you can see on the screen,.
I'm going to start a new sermon series..
This series is on the TV series,.
very successful, very popular these days..
It's been popular for a while, "The Chosen.".
How many of you have been watching "Chosen"?.
None of you?.
Only two, oh my God..
Well, you have to watch it because for a long while,.
I resist to watch it because I just didn't have time, right?.
But in the past few weeks,.
my daughter finally convinced me to watch episode two..
And I was kind of trying to entertain her desire.

$^{41}$and sit down and watch with her..
And watching like a few episodes and I just love it.
because the scenes that they take.
and the storylines that they create,.
although it's not straightly biblical storyline,.
but of course they add some poetic license.
and some creativity there and here..
But I'm just amazed how well it is done..
And I have watched quite a few episodes already..
And as you know, this is a historical drama.
based on the true story of Jesus..
And we know that we are followers of Jesus Christ.
and Jesus is a historical figure..
Now you can find the record of Jesus' life.
in sources outside of the Bible..
There are some Roman sources, there's some Jewish sources..
We can find a little bit the story of Jesus,.
but mostly we find in the gospels, in the New Testament..
And as you know, the first four books of the New Testament,.
Matthew, Mark, Luke, and John, the gospels,.
the four gospels have recorded quite a bit details.
about the life of Jesus..
And I know the church is always composed of people.
of different walks of life..
Some of you may be just beginning in your faith journey.
trying to know Jesus,.
trying to understand who this guy is, right?.
And some of us have been a Christians for a long, long time.
and your experience for church.
and your experience of what the Bible says..
I think as we watch this series,.
and I do encourage you to watch episode two,.
this series will help you picture.
the first century life of Jesus..
It will help you to understand.
what the Bible is talking about..
And for seekers, it is helpful..
For new Christians, it is helpful..
Then for experienced long-time Christians,.
it is also helpful..

$^{81}$Even I find it helpful..
And I'm going to do the series by referring to the episode,.
and then I'll come back to the Bible and say something.
and pinpoint something.
that is highlighted in different episodes..
I hope you understand, this is not the Bible,.
but it's a recreation of the biblical stories about Jesus..
So do watch "The Chosen.".
I don't get commission, right?.
We don't get commission..
I don't get any money,.
but do watch it because it is so well done..
And watch "The Chosen" and think the Bible and your faith..
That would be the slogan here, all right?.
As I told you, the gospels, we call it the life of Jesus,.
and Jesus is a historical figure..
And you can find sources outside of the Bible..
You find something about Jesus and his context.
and his background..
Now, Jesus was a Jew living in Israel, Palestine..
And back in Jesus' day, the Jews, or the Jewish people,.
were under the military occupation of the Roman Empire..
They had a pretty tough life..
And they were under heavy weight of taxes,.
all kinds of taxes, income taxes, property taxes,.
all kinds of taxes pressing on the Jewish people..
And as we all know, we also have our tax day, right?.
And we don't like the IRS, right?.
We don't like them..
I'm sorry if you work for IRS, good job..
But we generally refuse to think about taxes, right?.
And you can just imagine what life would have been.
for the Jewish people under the Roman occupation..
They had a very, very tough life..
And because the Jewish people had a long traditions.
for yearning for freedom, yearning for liberty,.
yearning for having their own nation,.
and every single year in the Passover meal,.
they will recite one story,.
and the foundational story in the Old Testament..

$^{121}$And what was that?.
The Passover, right?.
The Passover..
Moses came and liberate them from the land of slavery,.
defeat the oppressor, the Pharaoh, the king of Egypt,.
and then free the people up,.
and then they can really be the people of God, right?.
So they were yearning for freedom,.
they were yearning for liberty,.
and that's why in the Jewish tradition,.
there's always been people who was trying to fight back,.
right?.
That is the entire Old Testament,.
trying, saying something about the hope.
that God will reestablish his long lost kingdom on earth..
Now the Jewish people were under Rome..
And that is the context of the life of Jesus.
and the first four books of the New Testament,.
Matthew, Mark, Luke, and John..
Now if you compare the Gospels, the four Gospels,.
and other ancient biographies of the day,.
you'll find that there were similarities.
because the ancient biographies,.
they were not interested in the personal life.
of the public figure..
You know, we, modern people, we are really interested.
in the private life of the public figure, right?.
You know, the autobiography of Obama, right?.
Trump, Hillary Clinton, right?.
You check all those books, you know,.
they're trying to dig out, you know,.
their romances, right?.
Their history of personal life, right?.
Did they ace the test?.
And did they fight with their classmates, right?.
What happened when they grow up?.
And all that stuff..
This is modern biographies..
We're interested in the private life of the individuals..
But ancient biographies, like the Gospels,.

$^{161}$they're not interested in the private life.
of the individuals..
They are interested only in the public record.
of the individual..
Something you can find in the official archive,.
you know, from the database..
Everybody knows about it, right?.
And then they write autobiographies of the famous people..
People like Plato, Aristotle, philosophers,.
Augustus, Roman Empire, Pilate,.
the Roman governor of Palestine at the time..
They are only interested in their public life..
That is ancient biographies, and that is Gospels..
So you flip the pages of the first four of Gospels,.
Matthew, Mark, Luke, and John,.
it doesn't talk about, you know, Jesus..
Any girl had a crush on Jesus when he was in high school..
Nobody cares, right?.
They only care about his public ministry, public life..
So let's say we are in the fellowship,.
the Paneo fellowship, we are doing Bible study,.
and on chapter one, Gospel of Mark,.
all of a sudden Jesus came, right?.
And he started preaching the message of the kingdom..
The kingdom is here, right?.
So you repent, and we turn back to God..
Jesus was already an adult..
So they only are interested.
in the public life of the individual..
So we have to keep that in mind when we read the Gospels,.
Matthew, Mark, Luke, and John..
So here we compare the Gospels and the Chosen, all right?.
So this is episode one, all right?.
So go back home, and I think it's on Netflix, right?.
It's on Netflix, or Amazon Prime, or even on the website..
You can watch the Chosen..
Let me highlight episode one here..
Of course, it's based on the true stories of Gospel,.
and the movie makers make their own creativity.
and add their poetic license in there..

$^{201}$This is one of the screen cap in the beginning..
But episode one, let me summarize it,.
and I translate for those who are reading Chinese..
But on the left column, episode one is called.
I Have Called You By Name..
Now, let me give you a summary, only four points..
Now, the story start with 2 BC, all right?.
Two years before the common era..
There was Mary Madeline..
She was a little girl, and she was with her father,.
and she wakes up all the time in the middle of the night,.
and don't know why, she was scared of something..
And her father reminds her, Isaiah 43, one..
And that is a scripture we're gonna read later..
But she was afraid, but because of this word of God in her,.
and she could go through all these dark times..
And then, the second point, the story is fast forward,.
28 years later, now Mary Madeline was an adult..
And she goes by the name Lilith..
Lilith, because she was possessed by demons, many demons..
And she was suffering so much..
Perhaps in the video, in the movie,.
she is involved in prostitution,.
or she was attacked by people,.
but she fought back and she hurt someone,.
and so the community were very, very concerned.
about her being heavily possessed by the demons..
Then number four, no, number three..
There was a rabbi, well respected in Jerusalem,.
and also in Galilee, but his name is Nicodemus..
He travels to Capernaum, which is a fishing port.
by the Sea of Galilee, and some fishermen were catching fish.
on the Shabbat, which was illegal..
And they did so because they wanna avoid taxes,.
the Roman taxes, 'cause whatever you catch on the Shabbat,.
it's not counted as your income, I guess..
So they were violating the law of Shabbat,.
and they don't wanna pay extra income tax..
But the Romans asked Nicodemus,.
the well-respected leader, the rabbi,.

$^{241}$to help catch or arrest some of the Jewish fishermen..
And then later on, Nicodemus found himself.
facing the demon-possessed Mary Magdalene, or Lilith,.
trying to cast the demons out of her, but he failed..
He failed miserably, and finally,.
Jesus comes on scene, unexpectedly,.
and he cast out the demons and free Mary.
from demon possession..
So Jesus call her not Lilith,.
call her by name, Mary..
So that's episode one..
We do some biblical reflection here..
Another main storyline in episode one.
is definitely Mary Magdalene..
She goes by the name Lilith..
Now, you have to understand why Lilith,.
why people call her Lilith,.
because when she was possessed by demons..
Because in later Jewish traditions,.
or later legends in Jewish mythology,.
Lilith was the name for the female demon..
So the entire community, according to this episode,.
call her Lilith, because she was heavily possessed.
by demons, and she suffered quite a bit..
And Jesus came to redeem her..
And you have to understand,.
this is actually mentioned in the gospels..
Let me turn to the gospel of Luke, chapter eight..
In chapter eight, we find the name Mary Magdalene..
In fact, we find a list of names,.
people, male, and especially female, following Jesus..
Here Luke says in chapter eight, verse one,.
"Soon after, Jesus began a tour.
"of the nearby towns and villages in Galilee,.
"preaching and announcing the good news.
"of the kingdom of God..
"He took with him 12 disciples," the 12, right?.
"And along with some women,.
"who had been killed of evil spirits and diseases..
"Among them were Mary Magdalene, Mary Magdalene,.

$^{281}$"from whom Jesus had cast out seven demons.".
So again, is it really seven?.
Count one, two, three, four, five, six, seven..
Or maybe this is a complete number..
That means she was being possessed by many demons,.
completely controlled by the demon..
Mary Magdalene, and then Joanna, and Susanna,.
and some other women, they were following Jesus,.
contribute to his ministry and support Jesus.
and his disciples with their own financial resources..
And that is the first time Mary Magdalene.
show up in the Bible..
If you chase this name, this character in the New Testament,.
you find an amazing character being transformed.
by the grace of God, Mary Magdalene..
So here I summarize Mary Magdalene in the Bible..
Now she's called Mary Magdalene.
because Magdalene means the person from Matalah..
And Matalah is a fishing port just south of Capernaum..
Now some people would follow me to tour Israel.
from time to time, right?.
I lead a tourist group to Israel almost every single year..
You'll find if you drive today, if you drive,.
usually we drive, I drive from Capernaum to Matalah,.
it only take about 15 minutes, right?.
If I drive faster, probably 10 minutes, I'll be there..
Right, so it's very close by fishing ports..
And both these fishing ports are on the main.
international highway of the ancient day..
So they're small towns, but they're very, very important.
because they are by the seashore of the Sea of Galilee,.
and they're selling fish and catching fish..
And Jesus began his public ministry in Capernaum..
That is the reason, because he could preach the good news.
of the kingdom of God on the international highway, right?.
On the 1994, right, Kennedy Highway, right?.
And who would not catch that, right?.
And he was announcing and preaching the good news.
in Capernaum..
And just some of it is this fishing town,.

$^{321}$fishing port called Matalah..
So Mary Madeline simply means Mary from Matalah..
And then you chase the story, we read Luke 8 too..
She was one of the female followers of Jesus..
Now in this list, what is surprising is, Luke 8,.
we just read there are many names of women..
What is surprising, and you have to understand the Bible.
in its own context..
Because in Jesus' day, Jewish rabbis, there were many.
of them, they usually will find, will have male students..
You know, male disciples would follow a certain famous.
rabbi, right?.
Just like us, we wanna, when we go to university,.
we wanna go to the best university, right?.
We wanna go to Hartford, right?.
Go to Stanford..
In those days, you know, male students, you know,.
like the 12, they really, really wanna get into.
the best school of their day, right?.
And Jesus was one of the rabbis..
And of course, male students would want to follow him..
But here's the day, this is the norm of the day..
Rabbis only have male students, male disciples..
Hardly you find any rabbi would have females following him..
And Jesus broke the norm..
He was calling and allowing females to follow him.
from town to town, village to village, city to city,.
to do the ministry and to preach the gospel..
And Mary Madeline was one of them..
And how did the Bible introduce Mary Madeline?.
Not in a very good way, right?.
Because she had been cast out, seven demons out of her..
She was heavily suffering, oppressed by the demons..
But you can see how the Bible introduced her..
Her life has been changed..
Her life has been changed by Jesus..
Now, you know a thing or two about life change..
When we make the commitment to Jesus Christ,.
our lives should have been changed, at least fundamentally..
Now, it doesn't convert me from a sinner to a saint..

$^{361}$It doesn't do that..
Because no matter how many years I'm in Christ,.
I'm still 100\% a sinner, right?.
I've known Jesus for, for how many years?.
Oh boy, 30 some years..
Yeah, 30 some years, right?.
Since I was a youth..
No, my life has been fundamentally changed..
You know, my value system, my perspective on the world,.
my perspective on myself, my self image,.
how I understand other people..
These are all changed because of Jesus,.
because of the gospel..
How I see, you know, how the world.
has been changing around me..
And how do I see that from the kingdom of God.
perspective?.
Those are the things that should happen in your life.
when you make a commitment to make Jesus Christ.
your Lord and Savior..
That is the life change that should be a reality.
for all Jesus followers..
And that was the reality for Mary, Marylyn..
Now, I think the best analogy for talking about life change.
is immigration, immigration..
Now, some of you, some of you know what immigration.
is all about, right?.
Because you know it firsthand..
Immigration means you leave your own country, right?.
And you come to a new country,.
and you try to settle down here..
When you come to this new country,.
you learn new culture, right?.
New customs, other language..
And you still remember your old ways of life, right?.
As an immigrant, you still remember your old ways of life,.
but now you're in a new land, right?.
You probably still value your cultural root, right?.
Like Chinese New Year, right?.
I mean, who cares?.

$^{401}$Oh, you care, right?.
I've been to America for many, many years,.
and I lost track, right?.
I don't do all of this..
I don't give money to my kids..
My kids are born here..
And once they find out, oh, they have the Chinese heritage,.
sometimes when they were still younger,.
they were still hopeful that I'll give them money.
in Chinese New Year..
I said, no, this is America..
We don't do that..
Now they don't have that hope anymore, right?.
Because we are in a new land, new custom, right?.
And you understand you still value your cultural tradition,.
but there's a huge difference..
The huge difference for you is the previous kingdom,.
the previous authority, the political power,.
or rules no longer have any holds on you, right?.
You drive here, it's not the same way you drive.
in Hong Kong or in China, right?.
Or in Japan..
You drive differently in America..
In Chicago, you drive differently, right?.
You drive totally different,.
because you belong to the new country, not the old one..
Yes, you still have some memories,.
but now you're in the new land..
You're under a different kingdom, right?.
The previous political power had no holds on you..
And the Bible tells us the same sense..
We were under the kingdom of Satan,.
darkness and sin and shame before Christ..
And now we came to Christ, we came to know Jesus..
The power of sin and shame will still have influence on us,.
but it shouldn't have any authority over you..
It may influence you..
It may affect you, still, it may..
But you don't answer to those authorities..
You don't answer to sin..

$^{441}$You don't answer to darkness..
You don't answer to Satan..
You only answer to the new authority,.
which is the kingdom of God..
God is the only king in this new order of things, right?.
You answer to Jesus only..
Whenever you make decisions, especially important decisions,.
you should think in terms of what the true king,.
what the true king desires..
You should always orient your life.
according to this new kingdom of God..
And that is the life change that should be happening.
in all the followers of Jesus..
And that is what happened to Mary Magdalene..
Her life was changed..
And that's why in this episode, when Jesus came,.
cast out all the demons,.
the demons had no holds on Mary Magdalene..
And that's the same story for all of us..
Jesus came to redeem us, to save us.
from the old order of things,.
and brought us into a new order, a new kingdom,.
the kingdom of God..
And one of the passages that captures this message.
is found in episode one of "The Chosen.".
And that is Isaiah 43..
Isaiah 43, verses one and two..
It says, "But now, O Jacob,".
and Jacob is the nickname for the people of God,.
the nation of Israel..
"Now, O Jacob, listen to the Lord who created you..
O Israel, the one who formed you says,.
do not be afraid, for I have ransomed you..
I have called you by name, and you are mine.".
Isaiah 43..
Now, one more thing you need to understand,.
because when the Bible quotes a verse.
from the Old Testament,.
people don't take verse out of context..
They quote the entire chunk, the entire block of materials..

$^{481}$And that material is the book of Isaiah, right?.
Have you read the book of Isaiah before?.
It's a huge, long book, right?.
So whenever the Bible refers to this verse, 43, one,.
it's actually referring to the huge chunk of materials.
found in the book of Isaiah..
And you need to understand why the book of Isaiah.
was written..
It was written for those who were under oppression..
The Jewish people at that time, in Isaiah's time,.
100 years before Jesus, were under the Babylonian Empire..
They were under the military oppression of Babylon..
And these chapters, including 43,.
summarizing all the Jewish hopes that we have in God..
The people of God were waiting..
They were yearning for a new king to set them free,.
to rule over them..
The new order of things they were yearning for..
A new son of David, a new descendant of David,.
like Solomon, who established the first kingdom..
Someone like the son of David will come.
and restore us back to God, from oppression to liberty,.
from oppression to freedom..
And God will give us power to witness His greatness.
to the ends of the earth..
That's what Isaiah's all about..
I just summarized the whole book for you.
in the last 30 seconds, all right?.
And God tells the people,.
no matter how difficult your situation is,.
fear, unknowns, or challenges, I am with you..
Because I have called you by name..
I have called you by name..
In the biblical culture,.
whenever you say, "I call you by name,".
whenever they talk about naming,.
naming someone, or renaming someone,.
like Abram to Abraham, right?.
Sarai to Sarah..
Whenever you have name changing,.

$^{521}$you know, naming in the biblical culture.
is not simply a label..
It's not like among the crowd,.
"Hey, Mark, can you stand up?".
That's just a label, right?.
That's how we modern people understand naming..
But in biblical culture,.
naming is a totally different booking..
Because naming in the ancient culture.
has to do with authority and belonging..
You are being named because you belong to someone.
who named you, right?.
You remember there's a,.
you understand there's a very, very important ritual.
in the Christian journey,.
which is baptism, right?.
Baptism, you dunk someone in the water,.
and then you lift him or lift her back up, right,.
out of the water..
And what did we say when we do that?.
We say we baptize you in the name of the Father,.
the Son, and the Holy Spirit, right?.
And what does it mean?.
Does it mean getting a license?.
Do pastors get a license to do that,.
to dunk people in the water?.
No, the formula says,.
in the name of the Father, the Son, and the Holy Spirit,.
I baptize you..
That means you have the name of the triune God..
You have the name of God of Trinity..
The name, singular,.
the Father, the Son, and the Holy Spirit,.
and that means you belong to God himself..
When you get baptized, God's name is on your life,.
and he has authority because he is the new king..
And according to him, there will be life change,.
and new order of things..
So call you by name, it says Isaiah..
Do not be afraid, I have ransom you..

$^{561}$I have called you by name, and you are mine..
You belong to me..
No longer you belong to the oppressor,.
or your previous political power,.
or anyone that has a hold on you,.
or anything that has hold on you..
Whether it's the shame, guilt, depression,.
all those things no longer have any hold on you.
because you don't belong to them..
You belong to the God who ransomed you, who saved you..
God knows you by name and call you by name.
because you are his..
And he said, "You are mine.".
God know us, thank God, right?.
We belong to the new king, Jesus Christ,.
who died on the cross for our sin,.
to set us free from any oppressive powers,.
ultimately sin and darkness and death..
And Jesus free us and give us this new life, right?.
And Mary Magdalene experienced it,.
and we should too, we should too..
Now, before we close in prayer,.
we would close today's service with a communion..
I think we have the communion elements ready..
And by doing that,.
we proclaim to the world,.
as we protect the bread and the cup..
Remember, this is the Passover meal.
Jesus set up for generations of Christians, right?.
The Passover story, find fulfillment in Jesus.
because Jesus is the lamb of God..
He set us free, right?.
He set us free and we belong to him..
And by protecting the cup and the bread,.
we know we belong to Jesus and to him alone..
Let's do that because of grace of God,.
then we can be called the children.
of the kingdom of the light..
Let's all pray..
Father, as we bow our heads to pray to you,.

$^{601}$we know we belong to you..
And thank you for being our king,.
being our Lord..
You have called us by name and we are yours..
Thank you for loving us so amazingly and so unconditionally..
As we see the cross,.
as we see Jesus hung on the cross because of our sin,.
in order to love us and to redeem us and ransom us.
from all kinds of powers on the earth,.
sin and darkness and death,.
they no longer have any holes in our lives..
Thank you for the change that you brought about.
in your grace..
It's not because of what we have done,.
but because of what your son has done for all of us..
So as we observe the communion together,.
taking the bread and the cup,.
bring us back to your table..
Bring us back to your table of grace,.
knowing that we can partake in eternity.
because of your sheer grace..
Thank you, Jesus..
In Christ's name, we pray..
\newpage



\section{創世記}
\label{sec:QExp8HlGgEo}
\textbf{上主的護理 (5) 細想約瑟的故事 5 (創世記 44-50) - 陳崇基牧師}
\newline
\newline
連結: \href{https://youtube.com/watch?v=QExp8HlGgEo}{\texttt{https://youtube.com/watch?v=QExp8HlGgEo}} ~~~~ 語音日期: 2024-02-06
\newline
\newline
\hyperref[sec:WYvLrPDJi5g]{\small{< < < PREV SERMON < < <}}
~
\hyperref[sec:index_chronic]{\small{[返順時目]}}
~
\hyperref[sec:index_scriptual]{\small{[返順卷目]}}
~
\hyperref[sec:eEU2FEv4C5I]{\small{> > > NEXT SERMON > > >}}
\newline
\newline
$^{1}$看完節目就好.
這次我們繼續看約瑟的故事.
我們連續有五講,這次是第五講.
我們說到約瑟的故事.
我們說是今場的一個故事在創世記.
其實是在說一個聖經裡面很寶貴的教導.
就是Providence.
即是上帝對我們生命裡面的護理或者庇護.
上一次我們說到42和43章.
這裡說到約瑟和他十個哥哥重遇.
因為饑荒緣故所以他們去了埃及去吃糧.
所以42和43章我們上一次說到.
在這兩章裡面約瑟就去掙扎一個問題.
就是怎樣去測試他的兄弟.
這些兄弟有沒有改變.
還是好像以前一樣.
因為要分家產的緣故他們就弱肉強食.
然後就將他們家父所最愛的弟弟就將他剷除.
約瑟就是其中一個.
現在剩下一個就是約瑟的親弟弟是誰.
Benjamin沒錯,就是卞雅敏.
所以就希望能夠帶卞雅敏下來.
能夠測試到到底卞雅敏是否仍然還在.
如果還在的話.
我想約瑟希望的是保護到他.
因為他怕那十兄弟會殘害他.
所以我們上一次去到43章最後的時候.
就是一個筵席.
你記得約瑟到底他是埃及人還是希伯來人.
剛好在中間.
44章第一節.
在這裡筵席之後就打發這些兄弟回到迦南.
在這裡約瑟就出了一個主意.
他就吩咐他的管家說什麼.
他說按照他們奴隸所帶來的糧食.
將個人的銀子還給他們口袋裡面.
第二節.
我的杯.
約瑟說我的銀杯.
就要放在年輕人的袋口裡面.

$^{41}$最小的Benjamin的袋口裡面.
那個銀杯是搞什麼的.
在近東的世界裡面.
一些有權有勢的人.
很多時候都有銀杯是做占卜所用的.
就是說他料事如神的意思.
當然這裡不是說約瑟一定會做這些通靈的事情.
但是在這裡.
你會看到約瑟和他的哥哥.
這麼多哥哥之間的關係.
其實約瑟在故事裡面.
好像比他的哥哥更高.
知道的東西更多.
好像在他的哥哥面前.
好像一個神一樣.
因為知道他的哥哥的過去,現在.
還有知道他們的去留.
掌握到他們的行動等等.
所以這些兄弟就會覺得很奇怪.
為什麼約瑟會知道我們誰大誰小.
然後知道我們家父是怎樣的.
問我們的家宅等等.
無論如何約瑟就有一個策劃.
就將銀杯放在最小的Benjamin的口袋裡面.
然後就將他怪罪於Benjamin.
說他偷了他的銀杯.
然後就將所有的兄弟再押回來埃及.
所以就再去搜他們的口袋.
搜查一下誰偷了銀杯.
結果是誰呢.
就是Benjamin.
一搜出來的時候.
由大到小一直搜到最後Benjamin.
就搜到銀杯在他的口袋裡面.
這就叫做斷帳.
所以約瑟就想去扣留Benjamin.
情形就好像你過TSA的時候.
搜你的行李.
搜你的白粉.
是不是很嚴重.

$^{81}$這就是約瑟所設的計劃.
就是希望Benjamin怪罪於他.
讓他斷帳.
以至去試一試那些兄弟怎樣去對待Benjamin.
如果他們死性不改的話.
那些兄弟會怎樣.
我們走了.
我們有東西吃就可以了.
但就是這個測試去看一下.
他們的本性有沒有改變.
我們就去到44章第16節.
由大第四公子.
由大就發現Benjamin是斷帳的時候.
他就說了什麼呢.
他說你僕人的罪孽.
你看一下我們上帝已經查出來了.
我們看到那個杯子.
是在Benjamin的口袋裡面.
他說我們和在手裡找到杯子的人.
就是Benjamin.
都是我主的奴僕.
他說不如我們一齊去做奴僕.
約瑟怎麼說.
絕不能做這樣的事.
誰的手中找到那個銀杯.
誰就做我的奴僕.
至於你們.
可以平平安安上到家父那裡.
就是說想扣留Benjamin.
這個就是make or break的moment.
到底這些兄弟怎樣反應.
而他們的反應是不是表達出.
他們的生命終於有改變了這麼多年.
以下就是這個所謂由大的講話.
44章第18至34節.
這個在聖經的文學裡面.
是一個最著名的一段講話.
你沒有聽過吧.
那些研究古代文學.
研究聖經文學的.

$^{121}$整本聖經你找遍.
所有的literature.
所有的文學.
最好最好的一段就是這一段.
居然是由大講的說話.
就是由這段由大的講話.
和約瑟的求情的講話.
好像聖經裡面的莎士比亞.
有沒有讀過莎士比亞.
我們都是沒有什麼文化的人.
但是我想告訴你.
這個是聖經裡面最有文化.
和最有文學味道的一篇講話.
正因為這樣的緣故.
所以我們應該去看一看.
18至34節.
讓我讀出來或者演繹一下.
18節.
由大靠近約瑟的時候.
他說我主求求你.
讓僕人說一句話給我主聽.
不要向僕人發裂腦.
因為你如同發腦一樣.
我主曾問僕人們說.
你們有父親兄弟沒有.
我們對我主說.
我們有父親.
他已經年老.
還有他老年所生的一個小兒子.
就是卞雅文.
他哥哥死了.
就是說約瑟.
他的母親拉傑.
又是死了.
只剩下他一個的孩子.
父親也痛哀他.
你對僕人說.
把他帶下來.
我姐女來.
讓我親眼看著他.

$^{161}$我們對我主說.
這年輕人不能離開他老父親.
如果離開的話.
父親就會死去.
你對僕人說.
你們最小的弟弟.
若不和你們一同下來.
你們就不要來見我的面.
我們上到去.
你僕人我們父親那裡.
將我主的話曾經告訴過他.
25節.
後來我們的父親跟我們說.
你們再去買奢良.
我們說我們不能下去.
最小的弟弟就是卞雅文.
如果與我們同去.
我們就可以下去.
因為最小的弟弟卞雅文.
若不和我們同去.
我們不能夠見那個人.
就是約瑟.
你僕人我父親對我們這樣說.
你們知道我的妻子拉傑.
給我生了兩個兒子.
一個已經離開了.
我以為他是被野獸撕碎了.
直到如今我再沒有見過他.
現在你們又要把這個卞雅文.
從我面前帶走.
若卞雅文遭難的話.
那麼你們就害我白髮蒼蒼.
悲悲慘慘下到陰間去.
如今我回到你僕人我父親那裡.
若沒有這年輕人卞雅文與我們同去.
我們父親的命.
是與這年輕人卞雅文的命是相連的.
當我們的父親看見.
沒有這年輕人的話.
他就會死硬.

$^{201}$這樣我們就會害你的僕人.
我們的父親白髮蒼蒼.
悲悲慘慘下到陰間去.
是不是很傷心.
真的很慘.
這個時候你有沒有看過《月月殘片》.
那些男女相愛.
相愛了很多年.
為什麼我們這麼合得來.
然後結婚的時候.
洞房花燭夜的時候發現.
原來我們是失散多年的兄妹.
然後你們聽到小提琴就在那首樂曲.
然後就晴天霹靂.
就是這個時候悲悲慘慘下到陰間去.
怎麼辦.
猶大繼續說.
32節.
僕人曾向我老父為這年輕人卞雅文作擔保說.
我若不帶他回去交給父親.
我就在父親面前永遠擔當這罪.
現在.
這是猶大.
跟約瑟說的.
現在求你把僕人.
就是把我留下來.
代替這年輕人作我主的奴僕.
讓這年輕人和他哥哥一同上去.
若這年輕人不和我們一起.
我怎能上去見我的老父.
恐怕我要看到災禍臨到我的父親.
這是聖經最偉大的文學之一.
剛才我們說過.
為什麼這是一個這麼重要的講話.
或者一個文學氣味很重的一段講話.
因為猶大這一段的講話.
其實幾乎每一點.
都是撤銷了多年的恩怨.
多年有什麼恩怨.
感情上的恩怨.

$^{241}$你留意整段的說話裡面.
十四次看到父親.
我們的父親.
Our father.
如果你是約瑟.
你還沒有表露你的身分的時候.
他說我們的父親.
約瑟聽到我們的父親.
Our father.
他會想到什麼.
感情上他就放不下.
第二點.
孝道.
父親的命根.
就是最小的弟弟.
辯雅文.
他不像其他的兒子.
因為父親的命.
是和辯雅文小兒子的命連在一起.
有小孩的你明白吧.
有時候拉兒子拉心肝.
就是這個意思.
孝道.
為什麼要保護辯雅文.
不可以讓他做奴隸.
不可以讓他挾持.
第三.
情理上.
猶大多次提到.
家父最愛的是誰.
拉兒子.
辯雅文.
幾乎是唯一的兒子.
奇怪 明明有十幾個兒子.
只是愛那個.
只是偏愛那個.
我們不說偏愛是對還是不對.
這不是聖經處理的問題.
聖經處理的問題就是.
他確實是偏愛她.

$^{281}$你們這班兄弟.
是否願意接受.
你們肯不肯接受這個事實.
接受這個事實.
就是承認家父對弟弟的愛.
是寧死不同的.
你們肯接受的時候.
你就不會發生一些狗咬狗骨.
彼此傷殘的事.
這是事實.
家父雅各.
真的拉兒子拉心肝.
我們不是.
這麼多年為何有這麼多的恩怨.
就是他們不能夠去接納.
這個家庭的實況.
家父是真的偏愛她.
不說對錯.
是說事實.
有時候家家有本難唸的經.
你就是接受不到那個事實.
所以心裡面為何這麼難過.
這麼難受 這麼不安.
還有很多事情.
為何心裡面這麼難過.
這麼難受 這麼不安.
還有很多口角.
接受了這個事實.
你才能夠幫助你.
Move on.
雅各就是愛.
變雅文多過愛我們任何一個人.
Let's move on.
第四.
就是說誰做奴僕.
這是猶大做奴僕.
劉大願意.
金元來做奴僕.
這是最大的改變.
因為你看猶大的故事的時候.

$^{321}$你就會發現.
這個人本身不是這樣.
所以他說願意做奴僕.
代替這個小子.
他留下來.
去取代他.
所以你看完這段說話.
聽完這段說話的時候.
應該發生什麼事.
就真是復和了.
四十五章.
約瑟對所有士立的人.
他在他們面前情不自禁放聲大哭.
埃及人都聽到.
法老全家都聽到.
約瑟就跟他的兄弟說.
我就是約瑟終於表露他的身分.
對嗎.
我的父親還在嗎.
我兄弟不敢回答他.
都市驚惶.
約瑟就說靠近我一點.
你過來.
看一看我.
我真是你們多年前出賣的約瑟.
我就是被你賣到埃及的兄弟約瑟.
現在不要因為我賣到這裡而憂傷.
為自己生氣.
因為上帝早就差我來.
為要保存大家的性命.
這個很感人的一幕.
放聲大哭.
十二兄弟終於相認.
剛才我說猶大有改變.
這是最詫異的地方.
猶大由開始到現在.
你會發現開頭的時候.
這個人不是這樣的人.
猶大有什麼改變.
三十七章開始的時候.

$^{361}$誰出主意出賣約瑟.
四公子猶大.
三十八章我們插了一個故事在這裡.
我們沒有說.
但這個故事跟傳宗接代有關.
就是他寧願家族的沒落.
都不願意插手做任何事.
他媽就不是.
他媽就是他媳婦,寡婦.
他就不擇手段為了保留家族.
整個故事就是一個傳宗接代的故事.
他是不在乎.
弟兄不在乎下一代的人.
他只在乎自己.
四十三章最後的時候.
竟然他就說.
他願意去擔保他的弟弟.
然後四十四章我們剛才看的最偉大的講橫.
就是他願意作奴僕.
代替彼雅文.
代替他的弟弟.
四十四章三十三節.
現在求你北僕人留下來.
代替這年輕人作我主的奴僕.
讓這個年輕人和他哥哥回去.
一個徹底的改變.
以往就是出賣弟兄.
去使用弟兄,運用弟兄.
用手段來爭取自己的利益.
現在他就放下自己.
整個人不同了.
願意作奴僕.
你知道為什麼君王是猶大之派出來的嗎?.
為什麼猶大會最後出盛了.
能夠成為十二之派裡面的王.
就因為聖經定義王者是servant king.
是一個願意等下自己.
serve其他人的.
這個才叫做king.
和當時的社會和現在的社會.

$^{401}$做king是兩回事.
所以為什麼猶大這個講話這麼重要.
就是他的人生的改變.
整個人的改變.
以前就是必人做他的奴僕.
現在就是他作人的奴僕.
甘願作人的奴僕.
事隔多少年了?.
我相信到現在.
二十多年.
很多時候我們覺得.
上帝在一個人生命裡面的工作.
我們都希望或者我們祈禱.
我們希望上帝在生命裡面的工作有改變.
最好是突然來的.
馬上有就好了.
幾個星期他就可以改變過來.
我們很想他改善的時候.
如果上帝可以做一個transplant.
改了他的心.
立刻他的搖身一變.
好像超人換衣服一樣.
變了另外一個人.
但很多時候上帝的providence裡面不是這樣.
上帝去改變一個人很多時候.
是按照時間.
因為你知道神給我們每個人.
最大的禮物是什麼?.
神給人類最大的禮物.
是free will.
是自由的意志.
我們不是robot.
你記得我們在電腦或者電話.
sign up或者authorize的時候.
我們不是要click那個box嗎?.
那個box很煩.
not a robot.
我click他就是我不是機械人.
我們就不是機械人.
因為上帝不會強制.

$^{441}$強迫我們.
去做一些我們不願意做的事.
你問了一個問題.
因為你說的providence.
上帝是掌管萬有.
最大的事和最小的事.
上帝都可以掌握.
你怎麼可以說人有自由.
但是神又掌管一切.
這兩方面是一個奧秘.
但這裡告訴我們.
神怎樣去掌管.
甚至人的自由都不會令他有surprise.
因為上帝可以control.
你的外在環境.
用你的外在環境.
用你某一某的情況.
以至你會作一些選擇.
因為上帝知道你的心思意念.
可能會安排一些機遇.
一些外在的環境.
不是去逼你.
而是在那些際遇當中.
祂知道你會怎樣去作任何的決定.
舉個很簡單的例子.
講越深就越好.
快要講得更加深.
譬如我們不要講猶太.
我們講新約賣耶穌的猶太.
我們去到假想.
猶太是用三十兩的銀子出賣耶穌.
即是猶太人去到和猶太人講.
給三十兩銀子給你.
猶太變出賣耶穌.
但我們可以假想假設.
如果猶太人是cheap一點的話.
他們出的只是三十兩銀子.
你猜猶太會不會出賣耶穌.
猶太可能會說我哪有這麼cheap.
三十兩銀子就想出賣我的拉比.

$^{481}$猶太就不會出賣耶穌.
有兩個的可能性.
但這個就是上帝.
去到控制萬有甚至人的自由意志.
都不會出於他的預料之外.
這個理論叫做Middle Knowledge.
越講越深.
不講了 我停在這裡.
你回去google一下.
中間知識的理論.
即是說上帝對於人.
在某個處境裡面.
他很清楚那個人會作什麼的抉擇.
你如果這樣看的時候.
上帝的預知.
就和我們的自由.
雖然是一個奧秘.
但可以澄清到一些東西.
澄清到一些東西.
我們做父母的.
我們都有些預知.
雖然是人間的預知.
我們對我們的小朋友有沒有預知.
都有的.
因為知道某個情況下.
小朋友必定會作某種選擇.
今晚吃飯的時候.
你在小朋友面前放晚飯.
和放糖果.
你都差不多知道他會怎樣選擇.
所以你不會放糖在他面前.
不會讓他看到.
不會讓他想到.
你就在環境裡面.
塑造一個環境.
以至你知道他作什麼選擇.
這個就是上帝的庇護.
上帝的providence.
providence裡面就是上帝.
去控制我們面對的環境.

$^{521}$一個人面對的環境.
上帝是完全清楚你.
清楚你在某個環境裡面.
你會作什麼決定.
這個就是上帝對我們的providence.
正因為這樣的緣故.
我們就有平安.
在環境裡面.
上帝可以安排.
orchestrate.
以至我們去尋找祂.
你想想我們信主的時候是怎樣.
都是某種環境.
以至你去抉擇去信主.
之前你沒有信主 為什麼.
因為在那種環境裡面.
你就是不會選擇.
當然最終是一個奧秘.
但給我們卻是有一個安穩.
原來上帝是控制所有的環境.
我們自己的自由.
我們的決定.
沒錯 是真真實實的.
我們是有free will.
not a robot.
但對上帝的預知來講.
他不會有surprise.
我預備的時候想起很多年前.
有一個青年團契的學生.
以前我們青年團契的時候.
一班的年青人.
這個女孩.
是party girl來的.
令我們很頭痛.
我們兩夫婦都很頭痛.
因為教極她都不聽.
她傷害自己.
然後我們真的沒有她扶.
半夜兩三點.
喝醉酒.

$^{561}$然後我們又出去救她等等.
好幾年後.
她離開了.
然後再聯絡.
發現她搖身一變.
變成另一個人.
認識回耶穌.
你可以見到上帝就是用不同的環境.
不同的際遇.
當中去改變了她.
這個就是上帝對我們有的privilege.
在整本創世記裡面.
其實就是想convince我們.
去告訴我們.
上帝是一個怎樣的上帝.
你記得在創世記裡面.
開始的時候.
上帝呼召亞伯拉罕.
跟他說什麼.
要藉著你.
使萬國能夠得福.
使萬國能夠認識你.
你能夠去見證上帝.
是一個獨一的神.
你就發現每一個的族長.
亞伯拉罕 以撒雅各.
都是有他們面對的問題或者威脅.
威脅到幾乎是上帝的計劃不能夠進行.
是game over.
亞伯拉罕有什麼威脅.
他所遭到的威脅就是不辱.
沒有兒子.
所以就成為一個危機.
就是一個繼承的危機.
誰去繼承他的遺產.
誰去繼承他的弱.
最後上帝神跡地.
使以撒出生.
Isaac出生.
然後到以撒.

$^{601}$以撒又是同一個問題.
又是不辱.
又是同一個繼承的危機.
然後上帝神跡地.
使一對雙胞胎出生.
以素和雅各.
然後揀選雅各來做繼承.
雅各晚年的時候.
就是我們弱勢的故事.
你發現那個威脅.
他不是不辱.
是什麼.
是他內部的威脅.
是兄弟之間彼此相殘.
這個威脅更加要命.
因為弟兄相殘.
將家拆散的時候.
這個家怎麼能夠繼承上帝的計劃.
使萬國能夠得福.
上帝怎麼能夠使萬國得福.
你就會看到整篇的故事.
由亞伯拉罕以撒雅各.
去告訴你.
好像一切都是遊戲結束的時候.
弟兄之間躍躍強盛.
但是上帝仍然在進行他的計劃.
仍然能夠克服這些所有的威脅.
聖經就不斷要去.
convince你.
你所信的上帝就是一個這樣的上帝.
弱勢被賣.
一波又一波跌落低谷.
卻被高升.
然後弟兄和弱勢之間是拆散了.
上帝就在那個環境裡面.
讓什麼發生.
饑荒發生.
以致他們就落埃及.
然後最後來到重逢.
你讀這個故事的時候.

$^{641}$所有角色其實很confuse的.
大家不是很清楚發生什麼事.
因為上帝好像顯明.
但是又不顯明.
上帝好像在那裡.
但是又隱藏.
Is he visible or hidden?.
整個弱勢的故事.
不是很清楚.
弟兄有少少confuse.
弱勢最後也有少少confuse.
但是上帝就用大的小的環境.
七年的豐年.
七年的荒年.
來解除人為所造成的威脅.
十二支派最後就復和.
十二個兒子整家能夠復和.
因為弱勢一早被拆散.
就是要救活整個亞國家.
甚至萬國的人.
當時饑荒很大.
萬國的人都能夠得逢.
所以看回最後的時候.
我們說過.
Hindsight is insight.
看回我們過去經歷過上帝的事.
我們就有insight.
弱勢也是一樣.
弱勢的其中最後的講話說什麼.
五十章二十節.
他跟哥哥說.
你們以前的意思.
從前你們是要害我.
但是上帝的意思.
願是好的.
要使眾百姓能夠得到存活.
成就今天的光景.
所以你們今天不需要怕.
我必養活你們.
於是乎就安慰他們.

$^{681}$Again你看聖經.
我們不是在說道德的榜樣.
我們在查經的時候.
我們看到弱勢.
弱勢在饑荒的時候.
他做財政司.
他幫很多人.
我們將來的願望都是做財政司.
我們能夠祝福他們.
不是說這些moral examples.
故事是想告訴你.
神學.
上帝是怎樣.
上帝是一個有權能.
並且操縱所有東西的.
最後他的美意能夠成功.
萬國能夠得福.
你看到一個字.
一個好字.
好字good.
創世記何時出現一個好字.
第一章.
上帝創立萬有的時候.
整個宇宙都捍衛好.
good.
甚好.
到最後.
最後一個的威脅.
上帝去克服了之後.
他說上帝的意思.
完是好的.
你有平安.
你願意去信祂的話.
你願意知道上帝.
是一個更大是美好的計劃.
你的生命是不會浪費.
雖然你覺得很confuse.
雖然你覺得很浪費時間.
但你看到上帝.
是一個怎樣的上帝.

$^{721}$他的美意一定會行得通.
一定會克服所有的困難和威脅.
這個就是聖經嘗試去告訴你.
convince你所信的上帝是一個怎樣的上帝.
我記得我小時候.
幾歲大的時候.
大家應該有這樣的經驗.
有時候我們在美國.
我們有road trip.
我們一時可能開車幾個小時.
甚至九個小時.
十個小時.
小時候他們一上車.
開車不到三十秒.
他就問.
什麼時候到.
還有多久等等.
因為小朋友沒有sense of time.
他不明白開十個小時多不多.
是多麼美好的事.
所以我們做父母的.
有時候我們會跟他說什麼.
解釋不到.
總之還沒到.
就去occupy他.
有時候可能我們就將整個行程分段.
我們首先去麥當勞.
他覺得真的美好.
特別是有playground的地方.
playland的地方.
然後我們下一次就去那裡.
可能另外一間fast food.
或者一個mall等等.
我們就在中間有很多的終站.
一段一段的.
到達真正目的地的時候.
他就發現.
這個真的很美好.
整個的行程是需要時間.
是需要等.

$^{761}$在不知道裡面等.
在不知道裡面去信.
曾經有一個神學家.
他說什麼呢.
他說God's will does not come to us in the whole.
but in fragments.
in fragments.
and often in small tiny fragments.
我們看到上帝在我們生命裡面.
所揭示的美意.
不是一次過就讓我們看到大結局.
而是逐少逐少去給我們.
但每一片都是上帝的美意.
所以我們可以做什麼呢.
我們其實可以去回頭.
去看一看.
在生命裡面.
試一下發掘一下.
上帝有沒有一片一片的美意.
而一片一片會組成最後整幅圖畫.
整幅圖畫你現在看不看到.
看不到.
但我們深信上帝.
是對我們人生.
對他自己的榮耀.
對他自己的計劃.
是繼續進行.
所以有人說得很對.
Experience is not the best teacher.
聽到嗎.
經驗不是最好的教師.
Evaluate experience is the best teacher.
你要去評估一下你曾經經歷過的東西.
以至你才能夠去學到東西.
不是說我吃鹽比你吃米多.
問題是你以前吃鹽的時候.
你都沒有去evaluate過.
你去評估一下你以往有的經驗.
你就明白上帝是一個怎樣的上帝.
因為這樣的緣故.

$^{801}$在Prophetess裡面我們就有安穩.
我們低頭祈禱的時候.
可以你帶到上帝面前.
可能你今天在上帝面前有confusion.
或者有些不知的東西.
不要緊.
上帝說我就是一個實行美意的神.
實行美意的主.
你在不知裡面繼續信他.
你就明白上帝的計劃.
是一個美好的計劃.
可以同心低頭祈禱.
神可以多謝你.
因為你是那位美好的神.
縱使在我們生命的際遇裡面.
我們眼前未必能夠見到.
很美善的光景.
但我知道你是美善.
我知道你是good God.
你自己的計劃繼續推進.
你自己在我們身上有的旨意.
會繼續推進.
主讓我們在不知的當中.
不穩定的裡面.
在十字路口裡面.
我們去信靠你.
知道你是我們美善的主.
在你自己的護理.
在你的庇護底下.
我們無驚怕.
主我們今天來到你面前.
聖餐桌的前.
我們都想起你自己對我們美意.
在印刊內是一個羞辱的十字架.
但它在十字架裡面.
你彰顯你自己救贖復活的大能.
去拯救我們一班不配的人.
主我們拿起這個餅這個杯.
當我們看回我們的人生.
經歷過起起伏伏.

$^{841}$有多少跌碰.
主我們在評估的時候.
我們會發現你自己是一位美善的神.
一直帶領我們前面的人生.
讓我們信服.
走在你自己的美意當中.
讓你自己名得榮耀.
祈禱奉傳明教.
\newpage



\section{2:1-3}
\label{sec:eEU2FEv4C5I}
\textbf{The Chosen - Sermon Series \#2 (Genesis 2\_1-3) - Rev. Mark Chan}
\newline
\newline
連結: \href{https://youtube.com/watch?v=eEU2FEv4C5I}{\texttt{https://youtube.com/watch?v=eEU2FEv4C5I}} ~~~~ 語音日期: 2024-02-13
\newline
\newline
\hyperref[sec:QExp8HlGgEo]{\small{< < < PREV SERMON < < <}}
~
\hyperref[sec:index_chronic]{\small{[返順時目]}}
~
\hyperref[sec:index_scriptual]{\small{[返順卷目]}}
~
\hyperref[sec:aIlHRCRsv_k]{\small{> > > NEXT SERMON > > >}}
\newline
\newline
$^{1}$Well, why don't we say a happy new year.
to the neighbors next to you, in your own language..
All right?.
And then we'll begin our sermon series..
All right?.
Great, all right, all right..
That's cool..
We do have some refreshments out there, right?.
Yeah, I hope those guys won't devour everything..
(laughs).
And save some for us, all right?.
So it's very, very happy you come in..
There's a photo booth you can use..
Right after service, you can take some pictures.
or some family pictures with your friends, okay?.
So happy to celebrate China's New Year with some of you..
We will continue our sermon series on The Chosen,.
and this is part two..
And I have picked episode two of season one..
How many of you checked out The Chosen since last week?.
Please do, please do, all right?.
We will send out the link in the social media group,.
and you can watch it for free..
And Josh remind me, not only all the episode.
have the captions for different languages,.
including Chinese or different languages,.
but also there's audio, right?.
There's audio for different languages as well..
I'm not sure how I feel if Jesus speaks in Mandarin..
Not quite sure, yeah..
But anyways, but in Chosen, Jesus speaks English,.
which is kind of strange, right?.
Because the original language he spoke.
was Aramaic or Hebrew, right?.
In any case, it is a well-crafted TV series,.
and it's become so popular..
No matter you like the episodes.
or you like how they craft the story or not,.
I think watching it will help you picture.
what life was like in times of Jesus,.

$^{41}$because all the scenes that they pick,.
all the, you know, the houses.
or the background they pick are pretty realistic..
It's very realistic..
If you understand history,.
you understand the historical background..
They actually did a good job..
So watch The Chosen and think about the Bible.
and your faith..
Let's see, last time, we were at.
episode two, episode one..
Now, episode two, let me give you a summary,.
just highlight..
Now, remember episode one,.
it was Mary Madeline was tortured by many demons,.
and Jesus finally cast out all the demons.
and free her from the tyranny of Satan, of the demons..
And then she is now free,.
and now Nicodemus, the respected rabbi,.
is trying to find out what happened to Mary Madeline,.
because remember, Nicodemus, the rabbi,.
was being called upon to minister to Mary Madeline.
when she was tortured and possessed by the demons,.
and Nicodemus failed to cast out all the demons..
But that's how the story goes..
It's not in the Bible, right?.
But it's kind of interesting..
You try to creatively, you know,.
woven all the stories together..
But in any case, Nicodemus failed,.
and he wanted to find out what happened to Mary Madeline,.
and who cast out the demons, which he could not..
And now, back to Mary, in episode two,.
she prepares the first Shabbat meal,.
and she has invited some of her friends.
and those who are socially outcast,.
the blind, the disabled, and two of the disciples of Jesus,.
being invited to the home of Mary Madeline..
And of course, at the end, the surprise guests show up.
at the first Shabbat of Mary Madeline, right?.

$^{81}$And that surprise guest is Jesus, right?.
Jesus show up at the end of the episode..
Now, the key issue here I want you to think about.
is the Bible background..
Why they were so concerned about casting out the demons,.
like miracles that is done by Jesus..
Why did the Pharisees was so upset.
or so concerned about who is doing what,.
and why were there miracles?.
They want to authenticate all the miracles..
No, you get that kind of sense.
in the opening chapters of all gospels..
When Jesus came out publicly,.
he did a lot of miracles, signs and wonders,.
including casting out the demons..
And then the Pharisees from Jerusalem.
came out to Jesus and challenged him..
They are trying to authenticate in whose authority.
Jesus did all those things, right?.
Why did they so concerned about all of this miracles.
and they want to authenticate.
whether someone is using the authority of God,.
the name of God properly?.
Because the Jews in those days,.
they thought if true miracle happened in our time,.
that means the kingdom of God is finally here..
It's finally here..
Remember the Bible background before the New Testament,.
there was an Old Testament, right?.
Between the Old Testament and the New Testament,.
there were several hundred years.
that God seems to be silent..
What we call the silence period..
It's about three to 400 years..
We call it intertestamental period between the old and new..
And the Jewish people, they were waiting and yearning.
for God to come and to arrive..
The king will arrive, finally will restore his kingdom,.
the long lost kingdom of David..
And when the new Messiah,.

$^{121}$when the new king arrive at the scene,.
he will arrive with signs and wonders.
and casting out the demons..
And that's why the New Testament opens.
with signs and wonders done by the Messiah, Jesus Christ..
Now we were all in the Bible study, right?.
The gospel of Mark..
We just finished chapter nine, right?.
We just finished chapter nine..
And the first eight chapters were all about the signs.
and miracles that done by Jesus..
And, you know, reading between the line,.
you see the Pharisees and the authorities from Jerusalem.
is trying to check them out, right?.
To challenge Jesus and to ask Jesus by who's the authority,.
you are doing all of this, right?.
And of course, for those who are unbelieving, they resist..
They resist to acknowledge the divine authority Jesus had..
And they try to blame him for casting out demons.
with some other power, maybe demonic power, right?.
In fact, in chapter 11 of the gospel of Luke,.
Jesus said, after he cast out demons,.
and the authorities, the Pharisees,.
the religious teachers from Jerusalem challenged him,.
by whose power you're doing this, Jesus said,.
but if I am casting out demons by the power of God,.
then the kingdom of God has arrived among you..
Jesus speaks here, he says, you know,.
because I'm casting out demons with the authority of God,.
you know, the kingship of God has already arrived..
And because of me, when you see me work,.
when you see me saying things and teaching and preaching,.
you know the power of God has already arrived in our midst..
And that's what he said here..
The kingdom of God has already arrived..
Pay attention to the tense here..
It's not a past tense, it's not future tense,.
it's a present perfect tense, right?.
Has already arrived among you..
And you better acknowledge that because of power of Jesus..

$^{161}$Now this episode, try to, you know, intervere,.
interconnect all the stories together.
and talking about the first Shabbat.
prepared by Mary Madeline, all right?.
Of course that is not in the Bible..
The Bible never recorded, you know,.
Mary Madeline preparing for Shabbat,.
preparing for a meal on the seventh day of the week..
But we know that is the background.
because it's Jewish people..
For them, the most important thing is the Shabbat meal..
The most important thing..
Even today, modern Israelis,.
they will celebrate Shabbat.
and they will take it very, very seriously, right?.
It's kind of like those of you.
who have Chinese heritage, right?.
You have Chinese heritage.
and you have that cultural background.
and you want to celebrate Chinese New Year, right?.
Some of you do, right?.
We don't..
I don't..
Because we just too lazy..
So my son is back for the weekend.
and all we did is that we went out for lunch yesterday.
and we found a restaurant that is less crowded..
So we had lunch, right?.
And that's all, right?.
I don't give them money..
I wish they can give me money..
But we lost kind of the Chinese heritage..
But for many of you, because of the cultural background,.
you still take it seriously, right?.
And that's why you've got some refreshment out there..
That's why you've got some, you know,.
decoration or the red color, right?.
That you have..
This is only red shirt, kind of red shirt that I have..
So I save it for Chinese New Year..

$^{201}$But you know, for the Jewish people,.
Shabbat meal or the Shabbat.
is a super important sign and ritual..
And where does it come from?.
It comes from the second book of the Old Testament..
Remember the first book is Genesis..
Second book is Exodus..
Exodus chapter 20, here this entire passage.
is the 10 commandments..
Now, why there were 10 commandments?.
Because when God saved the Israelites out of Egypt,.
the land of slavery,.
he redeemed them and saved them by grace.
and saved his people and out of Egypt,.
they passed through wilderness and finally they came.
to the first most important mountain called Mount Sinai..
And they were there and God show up in clouds and thunders,.
clouds and thunders and lightning and fire.
at the top of the mountain..
And at the foot of the mountain,.
the people of God just redeemed..
And God spoke to them because I save you by grace..
I led you out of the land of slavery, land of sin..
Now you become my people and I give you the principles.
to live by because you have a new identity as my people,.
as the new people, that's why you live differently..
And God spoke the 10 words, the 10 principles,.
what we call the 10 commandments..
And those commandments are the basic constitutional.
principles lived by the new people of God..
And in the fourth commandment,.
it speaks of the Shabbat or the Sabbath..
Now this is super important..
Chapter 20, verse eight and following..
It says, God said to the new people of God,.
remember to observe the Sabbath day by keeping it holy..
You have six days each week for your ordinary work,.
but the seventh day is a Shabbat day or Sabbath day of rest,.
dedicated to the Lord your God..
On that day, no one in your household may do any work..

$^{241}$This includes you, your sons and daughters,.
your male and female servants, your livestock,.
any foreigners living among you..
For in six days, the Lord God made the heavens,.
the earth, the sea and everything in them..
But on the seventh day, he rested..
That is why the Lord God plus the Sabbath day,.
the Shabbat day and set it apart as holy..
Now this is very interesting..
Let me point out something here..
The Shabbat, what is interesting, the Sabbath is time..
It's time..
All right..
From a particular time to another particular time..
You know, time is being set apart here, being set aside.
and being seen as something more important.
than all those times, the ordinary time..
That's very interesting..
Right, the first thing God told the people of God,.
God said to them, God said to us,.
there was certain time period of the week.
that you have to see it differently.
from all those other times, right?.
Monday night, Tuesday night, Wednesday afternoon,.
those are ordinary time..
And there's a special time during the week.
that God sees it as important..
And he said, remember to keep that holy..
Now people of God are to refrain from working,.
the normal, ordinary work, right?.
And then on the seventh day, what we call Shabbat, right?.
They stop working and set it aside.
and keep that apart from all the other time..
Now what does the word Shabbat or Sabbath mean?.
What does it mean?.
It means stop, right?.
Cease, ceasing, ceasing to do something.
you have already done..
Just stop doing this and then you do something else..
Well, as I said, modern Israelis.

$^{281}$would take Shabbat seriously,.
especially those who are on the orthodox side,.
and they will stay home and they would rest..
They would go to synagogue to worship,.
and they will have, of course, the Shabbat meal.
as their family time, they enjoy it..
You go to modern Israel on the Shabbat day,.
which begins on Friday at sundown,.
the streets are suddenly quiet,.
no cars traveling, no one driving,.
no one on the street..
Those people you usually find on the street,.
they are on the way to somewhere else..
They may go to a synagogue,.
they may go to a Western wall to pray, to worship, right?.
So what an interesting social phenomenon, right?.
On Shabbat, they just cease everything.
and go to synagogue and stay home for Shabbat meal..
Now the Sabbath is important because it is the sign,.
the indication of who you are..
And for the people of God in Israel,.
for the Old Testament people,.
as the people of God, Shabbat, Sabbath,.
is a special sign of their relationship with God.
because God has entered a covenant,.
sort of a serious agreement with them..
God, Yahweh, and the people of God,.
and between them, there's a covenant,.
there's a relationship..
And God said, "One of the signs that you show the world,.
"you show yourself, is to stop working.
"and you come to worship me..
"You rest so that I can work," right?.
To acknowledge my kingship,.
to acknowledge that you are in control.
and I am not in control..
That is the most important sign in the Bible..
And according to Jewish tradition,.
Shabbat begins a few minutes before sundown..
There was some rabbis in Jesus' day,.

$^{321}$they said, "How do you know how to tell time?".
Today we have the Apple Watch, we have the clock, right?.
We have the internet to tell us the world clock, right?.
What time it is right now..
But for those people in Jesus' day,.
they didn't have a watch,.
so those rabbis were discussing.
when exactly does the Shabbat begin?.
And some of the literature tells us.
the exact time to begin Shabbat.
is when three stars appear in the night sky..
So they observe the stars and the moon, right?.
They find three stars appearing in the sky..
That's when Shabbat begins..
And you see it's working..
Remember, Shabbat or Sabbath means stop, right?.
And they just drop everything, right?.
They don't even cook..
They don't even cook, right?.
One of the most interesting experience.
that we have when we go to Israel,.
when I lead church groups to Israel,.
is that sometimes we'd stay in the hotel, right?.
And Friday night, when Friday night comes,.
and all through Saturday for 24 hours,.
the hotel kitchen is closed, right?.
So what are we gonna eat, right?.
What are we gonna eat?.
We eat the leftover..
Because what they do is that they leave the steamer on,.
and all those trays of food are being heat up, right?.
So they don't do any work in Shabbat, right?.
Because starting a fire or electricity is considered work..
So they turn it on on Friday and leave it on for 24 hours,.
and that's all you eat, right?.
So Thursday night, whatever you eat,.
you probably will eat the same thing Friday night, right?.
And then Friday breakfast becomes Saturday breakfast,.
because they take this seriously..
And in the entire Bible, anyone in the nation of Israel,.

$^{361}$if they fail to observe Shabbat,.
that will show that this thing or this respect.
for the special relationship between God and Israel..
Like for married couples, right?.
Some of us are married couples..
If one side, with one person,.
trash some of the symbols of the relationship,.
trash it, not losing it, not unintentionally,.
but if you trash it, some of the symbols,.
like wedding ring, certificate, right?.
Or something that is important between the two people,.
if you intentionally trash it,.
that means you show disdain toward that relationship.
and toward the person in that relationship..
And that is what Shabbat becomes..
It becomes a sign to tell you,.
that tell ourselves and tell the world.
that God is important, right?.
God is important..
So the Jewish people have already respected this tradition.
throughout the history,.
because they understand it's not just a commandment..
And this commandment rests on.
how God order the entire universe..
Now we read this passage, Exodus chapter 20..
It says you have to remember this day,.
to observe this day, and then guard it, to keep it holy..
You know, I like that word, guarding..
You know, guarding is actually used for military context..
It's like a soldier trying to watch out.
for enemies approaching, standing on the watchtower..
You know, you're always ready and self-disciplined, right?.
It's not a come and go..
You guard it and you keep it,.
because it's so important, right?.
It's not by accident..
To guard the day..
Now this passage is based on how God work, right?.
Verse 10, seven day, right?.
Seven day, or six days each week for your ordinary work,.

$^{401}$but seven day is the Shabbat day of rest,.
dedicated to Yahweh, the Lord your God..
On that day, no one in your household may do any work..
Now it's actually referring to the book of Genesis,.
the first chapter..
The first chapter of Genesis from verse,.
chapter one, verse one to chapter two, verse three..
We recorded how God work in the entire week..
Let me read this passage, chapter two, verse one..
So the creation of the heavens and the earth.
and everything in them was completed..
On the seventh day, God had finished his work of creation,.
and so he rested from all his work..
And God blessed the seventh day and declare it holy,.
because it was on this day, on the day when he rested.
from all his work of creation..
Now this entire creation week being recorded.
and being narrated from Genesis 1, 1,.
all the way up to chapter two, verse three..
Now, of course the Bible is not interested.
in scientific stuff, right?.
Moses or people in those days,.
they didn't have iPhone, right?.
They didn't have iPad,.
they didn't have all these modern technology..
They didn't even have electricity back then..
So the creation week is not a scientific account..
It is a spiritual account..
It is a theological writing..
It's more interested in theology..
It's more interested to tell you who God is.
and who we are if we relate to God..
So God's act of creation is being presented as one week,.
the work of one week, what we call six-day creation..
Now, if you remember chapter one,.
day after day, day after day,.
talking about how God set up the cosmos..
He set up the entire universe,.
and after six day of work, God rested from all his work..
He rested on the seventh day..

$^{441}$Now, we modern readers, we come before the Bible and say,.
"God rested.".
That's interesting, right?.
The eternal God would rest?.
Does it mean he is getting tired, right?.
Just like we get tired, right?.
After all day of studying, writing paper,.
or going to work, geez, we are so tired,.
and we just sit on our lazy boy, right?.
And sip a cup of wine or something,.
and then we want to rest, right?.
Now, in the ancient time,.
when people talked about God resting,.
it does not mean God is exhausted..
It does not mean God is tired like we get tired..
No, when the ancient people, the Bible people,.
when they talk about resting, God resting,.
they said, what they meant was,.
God enter into his temple, and he started to rule..
It's actually going up and be seated on his throne..
That is the meaning of God resting..
So it's not because God is tired..
God is going into his house, the temple,.
and inside the temple, there is one part in the temple,.
the very center of the temple, called Holy of Holies..
And there you find the Ark of the Covenant..
And the Ark of the Covenant is actually God's throne,.
God's own seat, his kingly seat..
And he was sitting there, and that is what it meant,.
what it meant, God rested..
So God rested means God enter into Holy of Holies,.
and God is in charge..
God is now the king..
He runs the entire universe..
So God is in control..
A modern analogy is that,.
how many of you have recently purchased a new computer?.
A new computer?.
There are some on sale in Costco, gaming PC, right?.
You can check it out..

$^{481}$I mean, you buy a PC, or you buy a computer,.
and you move it to your house, right?.
And you open the box, right?.
You get the tower, and you get the monitor,.
you get the stand, and keyboard, and all that stuff, right?.
And then you plug in all the cables,.
and all the connections,.
and whatever would make the computer work..
So you work a whole day in order to set up your workspace,.
right, to set up your work computer..
And after you set them all up,.
plug in all the cables with electricity,.
and all that stuff, what do you do?.
Do you go to bed?.
No, you don't..
You turn it on, right?.
And that's when you start to rule, right?.
You control it..
And that's what it means when God enters.
His holy of holies..
God rested..
It means He controls everything..
The entire universe, your life, your family,.
your career, your work, everything..
So that's what this Shabbat commandment mean..
It means honor the day..
Honor this 24 hours..
Keep that holy to worship God..
What it means is that God is the one who is in control,.
not me..
God is operating my life..
God is orchestrating everything around me..
God is in control every single thing.
from the greatest to the least..
God is the king, God is the Lord..
And after six days of labor, I've done enough..
The Shabbat is telling you enough is enough..
This is time for you to stop.
and come before the presence of God to worship Him..
Because you know, ultimately,.

$^{521}$God is the one who controls everything,.
not you and not me..
Do you ever find it interesting.
in the entire week of creation,.
God never set apart some things as holy..
The entire creation of Genesis,.
God did not sanctify or set aside a tool.
or a vessel or a building or a piece of land..
God never did that..
In the entire creation week, God set aside time..
God set aside time to be holy, to sanctify it..
What is it?.
It's a period of time during the week..
Now what it reminds us is that how we use our time.
tell us who we are..
How we use our time, how we manage our time,.
tell ourselves and other people who our true God is..
You know, we human beings, we live in time and space,.
concrete time and space, right?.
We watch the clock, right?.
We follow the clock..
Sometimes we have to clock in and clock out for our work..
Right?.
We live in time..
And sometimes we are busy, right?.
Probably most of the time we are busy, right?.
And we don't have time to do this or to do that..
But if we say we don't have time to worship God,.
to serve God, what it really means is that.
it's not that we don't have time..
Because God is so very fair..
Each one of us gets 24 hours a day, right?.
Nobody gets 26 hours, right?.
Nobody get more, nobody get less..
The matter is whether we have time or not,.
it's not the matter..
It's whether we have the priority right, right?.
What it really means when we say we don't have time.
is that we don't have the priority, right?.
Such and such does not have my priority, right?.

$^{561}$It's not number one..
That's why we say we don't have time..
Actually, we do have time..
What goes wrong is our life priority..
And Genesis 2 tells us, with Exodus 20,.
is that God is most important..
God is most important..
We rest and we come before God to worship Him..
Now when we came to Christ, our view of time,.
or our perspective on time should change,.
should be modified, should be reconfigured.
in such a way that we will honor and keep that holy.
one day a week..
And realize, because we realize that time.
is a gift of God, right?.
Don't you realize that?.
Some weeks ago, one of the guest speakers talks about.
young people, or relatively young people passed away,.
right, because of cancer, because of some terminal disease..
This is tragic, right?.
And life is not in our control..
Time is not in our control, right?.
And that's why the Bible tells us we have to value time.
because time is gifted by God..
It's a sheer gift..
And time is entrusted to us as treasure, and as money..
Just like money, just like our investment,.
that we need to take care of it,.
and we need to manage it well..
And the Bible, entire Bible, from the first page.
to the last page, it says, if God redeem us,.
if we know Christ is our Lord and Savior,.
then we are called to faithful management.
of all those good things God has given to us,.
and that includes time..
How we manage our time, right?.
Now, with the resurrection of Jesus,.
of course, the calendar for Christians has changed..
It's shifted from the seventh day to the first day..
And why?.

$^{601}$Gospel of John, chapter 20, it says,.
early on the first day of the week,.
Mary Mellon found the empty tomb..
On the first day of the week, Jesus rose from the dead..
And the New Testament tells us,.
because Christ came and Christ is true Shabbat..
Christ is a Shabbat..
In Him, we find rest..
In Him, we find the kinship of God.
because Jesus died on the cross,.
and He redeemed us from sin..
And He was buried in the tomb for three days,.
and on the third day, He rose from the dead, right?.
He rose from the dead to give us new life and new hope..
And that's what we have in Jesus..
So the day of Shabbat gradually.
shift from the seventh day to the first day..
And that's why we gather here, right?.
Right, whether it rains or it shines,.
we gather here because we acknowledge.
that only Jesus is our King..
Shabbat, faithful management of time, right?.
Now, my questions toward the conclusion is,.
are these, there are some evaluation questions.
we need to ask ourselves..
How often do I use Sunday, today,.
the day of worship, as an extension of six days of work?.
Right, from Monday to Saturday,.
oh, I got so much to do, right?.
That's okay, I'll use Sunday as the extension, right?.
I'll keep on doing stuff, right,.
to do my own thing on the day of rest,.
or on the day, the worst day of worship..
How often do we do that?.
Yeah, we do that all the time, you know,.
including myself, right?.
Now, that does not mean we just don't do anything.
and sit on the sofa for 24 hours, right?.
But today, the day, the time of Shabbat.
is to be set aside, and the first priority.

$^{641}$is to worship God and to serve Him, right?.
To serve Him..
To rest and to read scripture, to pray,.
to reflect on the past week,.
and to chart the course for the next week,.
for the coming week..
And that's what Shabbat should be like..
But how often do we use it as an extension.
of the six days of work, right?.
Let us repent and let us change..
You pick up the Bible today,.
you know, think about God in a more deep way,.
to reflect on His word,.
to reflect on what God has done in our lives..
And how do I honor the lost day in action?.
Question number two..
How do I honor it, right?.
I show up to worship God, no matter what..
To serve Him, no matter what..
And how should I plan ahead the lost day.
so I can truly worship God, right?.
Remember what I said, the Jewish people,.
they would take time to prepare for Shabbat, right?.
It's not like, okay, I wake up this morning,.
I find out, oh, it's time to go to worship..
Oh man, I didn't plan for it..
But okay, let's go, right?.
Or I don't go..
It just happened incidentally to us, right?.
No, Shabbat has to be prepared..
Observe and keep it holy..
That means we take action to guard it,.
to plan ahead, so that my time worshiping God.
will be pleasing to Him..
Let us all do that..
Why don't we close in prayer?.
And just remember that God is the King of the universe.
and He controls my life..
And how should I honor Him with my time?.
Let's pray..

$^{681}$Father, we give you thanks.
because you are the God of creation..
Everything is in your hand,.
including all the details in our lives..
We can rest in you because you rested..
You rested in your holy of holies..
You are ruling the entire universe..
You are managing every details of our lives..
So we rest in you..
We don't have to over worry about things in our lives..
We don't have to overwork.
because we know you are in control.
when things don't seem to be in control..
We know that you can calm the storm,.
even though our life can be going through.
some stormy period of time..
We know that you are our King..
And thank you, Jesus, for loving us..
And today we come before you and honor your kingship.
to worship and to serve you..
And I pray that you will bless us the rest of the week,.
that we'll continue with our labor of six days,.
but we'll worship you in our attitude,.
worship you in our heart..
And when we honor and observe the Shabbat,.
we know that you are always in control.
and you are indeed worthy of our praises and our worship..
Thank you for calling us to be your people.
and help us to have wisdom how to manage our time well,.
because you are a good God..
In Christ's name we pray..
Well, God bless you all and happy new year to all of you..
And hopefully I'll see you next week..
All right, make sure if you have time,.
join our Super Bowl party tonight..
You can just bring some stuff.
and we can enjoy our time together..
\newpage



\section{路加福音 8:1-3}
\label{sec:aIlHRCRsv_k}
\textbf{The Chosen - 講道系列 \#1 (路加福音 8\_1-3) - 陳崇基牧師}
\newline
\newline
連結: \href{https://youtube.com/watch?v=aIlHRCRsv_k}{\texttt{https://youtube.com/watch?v=aIlHRCRsv\_k}} ~~~~ 語音日期: 2024-02-13
\newline
\newline
\hyperref[sec:eEU2FEv4C5I]{\small{< < < PREV SERMON < < <}}
~
\hyperref[sec:index_chronic]{\small{[返順時目]}}
~
\hyperref[sec:index_scriptual]{\small{[返順卷目]}}
~
\hyperref[sec:LTR4cqb1AYI]{\small{> > > NEXT SERMON > > >}}
\newline
\newline
路加福音 8:1-3
\newline
\begin{longtable}{cl}
\hline
\hline
章節 & 經文 (和合本修訂版)\\
\hline
8:1 & \begin{tabularx}{0.7\textwidth}{X} 過了不久，耶穌周遊各城各鄉傳道，宣講神國的福音。和他同去的有十二個使徒， \end{tabularx} \\ \\ \relax
8:2 & \begin{tabularx}{0.7\textwidth}{X} 還有曾被邪靈所附，被疾病所纏，而已經治好的幾個婦女，其中有稱為抹大拉的馬利亞，曾有七個鬼從她身上趕出來， \end{tabularx} \\ \\ \relax
8:3 & \begin{tabularx}{0.7\textwidth}{X} 又有希律的管家苦撒的妻子約亞拿，和蘇撒拿以及好些別的婦女，她們都是用自己的財物供給耶穌和使徒。 \end{tabularx} \\ \\
[1ex]
\hline
\hline
\end{longtable}
$^{1}$各位好 新年快樂.
跟大家拜個年.
願意每一位都有上帝的福分.
剛才我們唱詩的時候也被提醒了.
真正的福分從何而來.
耶穌基督.
耶穌基督給我們的永生.
給我們的盼望.
這是最大的福分.
你不要忘記這件事.
也跟網上的大英姐妹.
先去拜年.
農年.
希望大家農晚精神.
好嗎.
我們有16000多位的訂閱者.
我們希望你能夠繼續支持我們使公.
所以無論如何.
新年快樂.
從這個主日開始.
我想開始一個系列叫做「Chosen」.
我知道這個電視系列.
是非常…至少在美國.
很受歡迎的.
很受歡迎.
我不知道大家有沒有看過.
或者有沒有聽聞過.
有誰聽聞過.
也有吧.
是吧.
如果你有聽聞過.
可能你也有看過.
是吧.
那些比較老練的.
或者我們比較老一點的信徒.
我們以前看過耶穌的電影.
那些將耶穌的生平拍成一個電影.
那些不同的電影.
我們會….
你有沒有看過.

$^{41}$很悶的.
我記得70年代有一套叫做耶穌傳.
就是按照路加福音拍的.
就是一字不漏.
是根據路加福音.
就是耶穌說什麼路加福音有說的.
祂也說.
路加福音沒有說的.
耶穌也不說.
早晨,晚安等等.
耶穌全部不說.
那時候年少的時候看這套電影.
因為那時候沒有選擇.
唯有看那一套.
但是進度非常慢.
而且你知道路加福音寫出來的時候.
不是為了拍電影的.
所以你不可以照樣搬過去.
在銀幕上.
耶穌就是這樣.
不是的.
那些是故事的故事.
所以是不行的.
不久之前.
就開始了這個Chosen的電視系列.
我本身也沒有什麼興趣.
因為我沒有什麼時間.
但是我女兒就很有興趣.
她就經常看.
經常看.
她就經常鼓勵我.
你一定要看.
我沒有什麼時間.
我又沒有什麼興趣.
所以我就推了她好幾次.
結果最近.
她就成功地說服了我.
我就說好吧.
我跟你一起看.
她又想重看.

$^{81}$她說竟然要我重看.
以前的那些看過耶穌傳有沒有重看的.
不會的.
她說以前我們看.
不能說兩倍的.
是不是.
快一點.
耶穌走快一點.
說話說快一點.
沒有這招的.
她說竟然重看.
這麼好看.
可以讓她重看.
結果我.
應酬一下她.
我就看了一下.
結果我看了一點.
我覺得拍得很好.
拍得很好.
我所著眼的地方是什麼呢.
就是那個外景.
她一定不會在以色列拍的.
因為那個號支很龐大.
現在以色列搞成這樣.
又不是很安全.
我看了一下那些外景.
那些外景都挺像的.
就是加利利湖.
雖然它有些地方在德州.
或者在不同的地方拿外景.
不過看下去的時候.
這個場景都挺嚴肅的.
還有.
那裡的.
就是TV series裡面有的故事.
我覺得穿插得不錯.
我剛剛前兩天.
我看完season 1.
就是有八集.
season 1.

$^{121}$現在已經是.
好像是season 4.
season 4.
所以第一季我已經看完了.
我覺得真的不錯.
我相信在教會.
我們很多時候都有不同的朋友來到.
有一些可能是比較初初去接觸聖經.
接觸信仰.
就是聽到耶穌這個名字.
有一些是像我們這些信了幾十年的.
我們對聖經好像很熟悉.
是不是?.
對教會好像很熟悉.
對耶穌的生平或者記載.
我們好像比較熟悉一點.
我們有不同的人.
different walks of life.
但我相信.
我們如果去看回耶穌的故事的時候.
其實每一個人都是有好處的.
即使你是剛剛去認識.
或者你認識了很多很多年.
所以在這個講道的系列裡面.
我就是想按照.
找幾集.
Chosen.
你都可以講一講.
當然了.
這個Chosen的TV series.
就是他開頭講的時候.
就是按照一個true story.
是嗎?.
當然這個true story就是耶穌的.
生平的故事.
是嗎?.
四福音裡面所講的.
但是為了增添真實性等等.
所以就加了一些經外的內容.
或者一些所謂poetic license.

$^{161}$加了一些藝術的想像在裡面.
但是最重要的.
我們回到福音書裡面.
怎麼去介紹.
怎麼講耶穌的故事.
我相信這個系列.
大家回去.
儘管回去看看.
OK.
儘管回去看看.
在網上是免費可以看的.
還有字幕的.
有人告訴我.
甚至是有普通話的配音.
我不知道.
我不知道耶穌講普通話會怎樣.
如果耶穌講普通話很搞笑.
Anyway.
不過你Chosen.
耶穌都是講英文的.
是嗎?.
本身是講亞文的.
Anyway.
儘管回去看看.
我鼓勵你回去看看.
你可以從Episode 1開始看.
但是看的時候想一下.
聖經裡面怎麼講.
看的時候想一下.
福音書裡面怎麼講.
還有我們信仰的生命是怎麼連結在一起的.
每一集.
大家要知道.
每一集.
都有耶穌出現的.
他不會藏起來的.
我也知道香港的朋友最近很失望.
看《求王》,《美斯》.
誰知道陪他一起看球賽.
是嗎?.

$^{201}$但是Chosen不會的.
Chosen是每一集都有的.
每一集都有耶穌出現一點點的.
起碼一點點.
所以是一定不會缺席的.
你不會跟耶穌一起看戲的.
我就想從Episode 1開始講一下.
但是我想講一下.
耶穌確實是一個歷史的人物.
除了聖經裡面有記載之外.
聖經以外的一些文件.
譬如羅馬人的一些文件.
或者猶太人的一些文件.
都有提過耶穌的名字.
或者耶穌這個歷史的人物的.
他確實是一個以色列巴斯坦第一世紀.
出現過的一個人.
在耶穌的時代.
當時的猶太人.
是被羅馬帝國軍事佔領的.
所以那些猶太人都是活得很苦.
稅項非常重.
是嗎?.
我們美國快要交稅了.
你覺得交稅要交很多人嗎?.
那時候他們更多.
那時候更多.
如果你是漁夫.
或者你是農民的話.
你交稅是非常非常多.
是非常重的.
就是因為這樣的背景裡面.
耶穌是出現的.
舊約裡面經常看到.
猶太人的聖經經常看到.
就是上帝一定會來到.
重建他自己的國度.
就是說上帝的子民能夠復興起來.
能夠有自由.
你知道每年一次.

$^{241}$每年一次.
他們一起有團年飯的時候.
我們華人就會在初一十五.
正月的時候有團年飯.
但是猶太人的團年飯就是.
吃什麼?.
什麼時候吃?.
御月節是嗎?.
御月節就是他們每年一次的團年飯.
那個團年飯是做什麼的?.
他們會講故事.
講什麼故事?.
出埃及.
就是摩西會帶領我們.
離開埃及遺老之地.
然後呢.
上帝會大能的光能.
然後會打敗欺負我們的那個.
那個壓迫者.
然後將他們自由起來.
讓他們能夠真真正正成為神的子民.
每年一次.
每年一次過節的時候.
吃團年飯的時候.
他們就會去講這個故事.
就是出埃及記的故事.
所以在羅馬帝國的時候.
第一世紀的時候.
猶太人都是照樣這樣做.
猶太人都是一想到.
以前上帝給我們的認許.
上帝給我們的承諾.
是一定會兌現的.
他們的盼望就在上帝那裡.
所以上帝一定會有救主.
派到我們中間.
然後重整這個已經.
失去了很久的一個國度.
所以因為這樣的緣故.
在這個背景底下.

$^{281}$耶穌出現.
所以你要記住這個背景.
當你看耶穌的故事的時候.
你就明白.
到底當時的人掙扎著什麼.
我們首先呢.
由第一集開始.
第一集.
第一集講什麼呢.
那我想問你.
當我講到第一集裡面的情節的時候.
有什麼在聖經裡面有提過.
有什麼是製片商.
或者作者.
電影的作者.
加上去一些想像.
第一集講什麼呢.
大概主前二年的時候.
有一個鬼的少女叫做莫大拉瑪利亞.
莫大拉瑪利亞當她小時候.
她整天晚上就會醒來.
她會很害怕.
為什麼這麼害怕呢.
因為怕不知道有什麼騷擾她等等.
她爸爸就經常提醒她.
一句經文.
以賽亞書四十三章.
就告訴她四十三章裡面講什麼.
以致她不會懼怕.
這樣開始之後.
轉眼間就跳到很多年之後.
就是二十八年之後.
就是莫大拉瑪利亞已經長大了.
這個莫大拉瑪利亞在群眾的眼中.
其他人怎麼看她.
或者叫她什麼名字呢.
就不是叫她Mary.
是叫她Lilith.
叫她Lilith.
這個莫大拉瑪利亞.

$^{321}$就是被很多的惡鬼所附的.
她是有鬼附在身.
因為這樣的緣故.
所以她就受了很多苦.
她可能襲擊了某人.
然後人們就會懷疑她真的被鬼附.
然後就找了一個人去驅鬼.
那個人的名字叫做尼哥底姆.
Nicodemus.
尼哥底姆是當時一個非常出名的拉比.
適逢他就去到莫大拉瑪利亞住的地方附近.
就是去到加伯隆.
尼哥底姆去到加伯隆.
加伯隆是一個漁港.
在加利利海附近的一個漁港.
那裡就說到一些漁夫.
猶太人的漁夫.
就是因為要納稅.
因為有很重稅.
所以他們很多時候就違例地.
在安息日的時候都做事.
因為猶太人安息日是不做事的.
所以安息日他照打魚.
因為打了那些魚就不當是收入了.
不用上報.
所以那裡就可以逃避一些的收入稅.
我們很熟悉.
我們經常都做這些事.
就是想清除那些收入稅.
猶太人就想做這些事.
安息日的時候打魚就不算了.
但是那些羅馬人當然不妥.
所以就找了一個大的拉比叫做尼哥底姆.
去到想拘捕這些猶太人的漁夫等等.
因為尼哥底姆去到加利利附近.
所以有些人就找了尼哥底姆.
去幫Lilith或者莫大拉瑪尼亞.
希望能夠去拘捕.
但是尼哥底姆的能力不夠.
就不能夠將他拘捕.

$^{361}$他就失敗.
結果呢.
誰出場呢.
就是耶穌.
耶穌就去拘捕.
我想問你整個情節裡面.
有些什麼在聖經裡面.
有些什麼不是在聖經裡面.
有很多都不是在聖經裡面.
有很多都不是.
但是有一些是聖經裡面確實有提過的.
你知道看這些電影的時候.
你就留意有些情節是別人加上去的.
以至那個故事能夠吸引多一點.
那些聽了看了是OK就行了.
但是最重要是什麼.
我們回到聖經裡面.
聖經裡面怎麼說.
有一樣東西想提一提.
那個莫大拉瑪尼亞.
他是被鬼附身.
受了很多苦.
那些人怎麼叫他.
叫他什麼名字.
不是叫他Mary.
是叫他Lilith.
為什麼會叫他Lilith.
因為猶太人後期的傳統.
主後很多世紀的時候.
他們的傳統就是說.
女鬼叫什麼名字.
叫做Lilith.
就是叫Lilith.
我知道新年的時候說這些不是很利事.
但是我們百無禁忌的.
因為耶穌最大.
所以就叫他Lilith.
耶穌最後就去幫助莫大拉瑪尼亞.
以至這些鬼去纏身的就被驅.
然後被趕走.

$^{401}$聖經裡面怎麼說.
這些人物或者有什麼情節裡面.
當然不多.
但是這裡說到什麼.
路加方第八章.
路加方第八章由第一章到第七章.
都是在說耶穌公開的事公.
他公開的事公.
他傳道之後也做了很多趕鬼驅鬼的行動.
以至將神的國帶來.
說了一輪.
然後到第八章第一節.
他說過了不久.
耶穌周遊各城各鄉傳道.
宣講上帝國的福音.
和他同去的有十二個使徒.
當時已經召集了十二個使徒.
還有曾被邪靈所附.
被疾病所纏.
而已經治好的幾個婦女.
原來不只是一個.
有幾個.
其中有一位叫做莫大拉瑪利亞.
Mary Madeline.
曾有七個鬼從她身上趕出來.
又有希律管家撫薩的妻子約阿娜.
和蘇薩娜和好些別的婦女.
她們都用自己的財物.
去供給耶穌和他的十二使徒.
這裡提到莫大拉瑪利亞Mary Madeline.
什麼叫做莫大拉瑪利亞.
莫大拉瑪利亞其實很簡單.
就是一位叫做瑪利亞的女孩.
是從莫大拉這個地方出來的.
Metala這個地方.
Metala在哪裡呢.
今天去以色列也可以去看的.
因為那裡找回一些第一世紀的房屋.
和一個很漂亮的猶太人的會堂.
莫大拉其實就在加巴隆附近.

$^{441}$如果開車的話.
我開車.
應該十分鐘就到了.
十分鐘左右.
你開車應該也是十五二十分鐘就會到.
很近的.
就在加利利沿海.
一直南下.
然後就去到加利利的西岸一個地方.
一個漁港.
那裡就叫做Metala.
現在是屬於一個天主教的物業.
在那裡有很多的發掘.
很多的發言等等.
在一個漁港.
所以莫大拉瑪利亞就是Mary Magdalene.
從莫大拉這個漁港.
長大或者是一個出生的瑪利亞.
路加福音第八章這裡這樣介紹.
他說她是其中一個.
很多婦女裡面其中一個.
去跟隨耶穌和支持耶穌的傳道.
你知道當時耶穌是各城各鄉來傳道.
他起碼也要住Airbnb.
起碼也要找個地方睡.
就算不住屋.
他也要紮營.
也要伙食等等.
所以是一個很實際的需要.
所以莫大拉瑪利亞和等等的婦女.
她們經歷過耶穌所給她們的工作之後.
她們就跟隨耶穌和十二個滿徒來支持她們.
八章第二節這裡提到.
莫大拉瑪利亞.
和很悲傷的婦女.
這裡就是.
我們看聖經要懂得看.
逢聖經提到女人的時候.
我們就要豎起對耳朵.
我們現代人不會這樣看.

$^{481}$現代人看一本小說或者看一個故事.
提到男提到女很簡單.
因為世界上有男有女.
但以前的一些文獻.
很少會提到女人.
特別是寫男人的故事的時候.
因為不需要提到女人.
是不是.
當時的第一代讀者.
第一批的讀者.
都是男士.
因為父權社會就是這樣.
男人大了.
是不是.
男人聽到故事.
回家煮飯的時候跟老婆說.
他不需要做直接的讀者.
所以你看到這些聖經.
等等這些古代文獻.
一提到女人的時候.
我們就要豎起對耳朵.
為什麼呢.
因為提到女人的時候.
就代表是非常重要.
所以他就特意提.
是不是.
這裡提到什麼.
婦女.
幾個的婦女.
已經自豪的.
當中有稱為梅達拉·姆瑪利亞.
曾經怎樣.
然後有一個怎樣介紹的約阿娜.
有一個叫做蘇薩娜.
和好些別的婦女.
這個就是奇怪的地方.
是不是.
大家知道第一世紀.
猶太人的拉比.
猶太人的律法老師.

$^{521}$通常他們都會有門徒的.
有一群門生跟著他們.
但通常是學生找拉比的.
學生找拉比.
學生去申請.
或者想找拉比.
我想跟你等等.
這個是當時拉比的文化.
等於今天.
你想不想入哈佛大學.
想吧.
你不想嗎.
知道你能進去.
但讀不了.
你想入名門的學府.
就是你去申請的.
有什麼時候哈佛大學會去拍門說.
我想你進我的大學.
對不對.
嚇得你腳都發抖.
對不對.
當時就是這樣.
是門生去找老師的.
很少會去老師找門生.
首先這是第一樣.
耶穌是去呼召門徒的.
不是門徒來找他.
是他去找門徒.
並且.
很少會有女門徒.
很少有女人跟著他.
耶穌做了兩件.
當時革命性的事.
他呼召門徒.
並且呼召女的門徒跟隨他.
也讓女婦女跟隨他.
支持他的工作.
莫大拉在聖經裡面.
是一個跟隨耶穌的其中一個門生.
雖然他不是十二使徒之列.

$^{561}$但你看到整個聖經裡面.
去介紹莫大拉在馬尼亞的時候.
由他開始.
他是由莫大拉這個漁港出身的一個人.
跟隨耶穌去支持他.
到了漢福音第十九章.
他就在十字架之前.
就是耶穌被釘死的時候.
他是親眼目擊的.
並且第三天之後.
耶穌復活.
七天的第一天.
耶穌復活的時候胸墳無.
他是第一個去見證這件事.
這個也是非常詫異的.
因為當時的人是不會相信女人說的話的.
今天也未必相信.
不相信女人說的話.
特別是法庭上.
女人的證供是有嫌疑的.
男人是歧視女人的.
這個就很奇怪了.
福音書裡面說.
第一個見證胸墳無的人.
是一個女人.
這個女人就是莫大拉在馬尼亞.
並且福音書裡面提到這些名字.
無論是瑪利亞.
約阿娜.
蘇薩娜.
這些一定是當時的讀者熟悉的名字.
對不對.
寫一個當時初代教會所熟悉的故事.
耶穌的故事.
耶穌傳道的故事.
工作的故事.
提到一些名字.
特意將他們的名字說出來.
一定是當時的讀者知道.
莫大拉我們聽過.

$^{601}$甚至我們認識的.
所以在初代教會裡面.
莫大拉瑪利亞.
很肯定.
是一個很重要的人物.
無論是因為他用他的方法.
或者他的資源.
來幫助耶穌的運動也好.
什麼原因也好.
或者他開放自己的家裡.
以至初期的教會.
他能夠去聚會等等.
總之用了他的資源.
來支持耶穌的事件.
所以當時初代的基督徒.
他們拿起福音書的時候.
他們就知道.
這個人就是經歷過.
耶穌.
給他的工作.
然後他就去奮身的.
來事奉耶穌.
八章第二節.
提到莫大拉瑪利亞.
曾有七個鬼從他身上出來.
被趕出來.
七個鬼.
按道理就應該是數著七個.
第一ABCDE.
這樣數下來.
有很多隻.
你可以想像他的生命.
本身是一個.
多麼悽慘的生命.
但是耶穌去.
使他的生命能夠得到改變.
但這個其實應該是.
每一個信徒.
有的改變.
我們未必好像他那樣.

$^{641}$但我們應該有一點點.
Life Change.
每一個信耶穌的人.
應該基本上.
是有一個生命的改變.
因為聖經裡面告訴我們.
救恩是什麼.
救恩就是從.
以前的國度.
我們所屬的國度.
然後將我們救出來.
以至我們進入一個新的國度.
一個新的Kingdom.
以前那個國度.
裡面有的原則.
對我們有的權能.
因為信耶穌的緣故.
已經被破解.
現在我們進入一個新的國度.
那個新的國度.
是屬於神.
屬於光明的國度.
那你說什麼呢.
什麼意思.
最好的比喻就是移民.
如果你想一下移民你就知道.
移民就好像救恩一樣.
在耶穌裡面有的恩典.
就好像移民一樣.
在我們中間有很多移民.
很多移民.
移民是什麼一回事.
移民就是離開你以前的國度.
來到一個新的國度.
你可能還記得以前的國度.
以前那個國家.
裡面有的法例也好.
有的權能也好.
你還會憶想到記得的.
但是你已經離開了.

$^{681}$你來到這個地方.
你就跟一個新的國家.
新的政權.
在這裡的法律底下.
在這裡的管治底下.
在這個新的權能底下.
你也要生活.
你聽的權能.
就是這裡的權能.
不是以往的權能.
是 可能以往的權能.
對你來說還是有一些陰影也好.
有影響也好.
對你一些的掣肘也好.
但是它沒有影響到你.
對你來說已經不再有任何的箝制.
來到一個新的地方.
你就是在一個新的規律下.
以前你怎麼開車.
今時今日你在這裡怎麼開車.
是兩回事.
你看的法律是看這裡的法律.
不是看以前那個國度的法律.
這個就是最大的分別.
信耶穌其實也是一樣.
當我們說耶穌是我們的主.
我們認識耶穌的時候.
我們就離開以往的國度.
聖經告訴我們.
以往的國度.
就是撒旦.
開車的.
是撒旦的權勢.
罪惡的權勢.
死亡的權勢.
黑暗的權勢.
現在我們來到耶穌的國度裡面.
這些所有的權能.
罪惡的權能.
撒旦死亡對我們的權能.

$^{721}$罪惡.
羞辱等等這些權能.
可能對我們仍然是有一點點影響.
可能我們仍然會在陰影底下.
不過我們不回答那些權力.
我們不再被祂問責.
對嗎?.
真正問責我們的.
就是這位愛我們的上帝.
所以在剛才說的情節裡面.
有一段聖經.
就是《以賽亞書》第43章.
你記得這個故事裡面.
就是穆達拉瑪雅小時候.
爸爸跟她說.
《以賽亞書》第43章.
第一節.
這裡說什麼呢?.
他說:亞國創造你的耶和華以色列.
造你的那位現在如此說.
你不要害怕.
因為我救贖了你.
我曾提你的名.
召你.
你是屬我的.
這句話有什麼特別呢?.
除了在情節上你覺得很感人之外.
I call you by name.
我用你的名.
你都可以召你.
這段聖經《以賽亞書》第43章.
大家知道以前的人用聖經.
他們不只是引用那一節.
他們引用的時候引用一整段.
上去問問你.
全部都知道.
所以引用一段.
然後有很多其他東西.
這一段就是《以賽亞書》第二部分.
《以賽亞書》第43章到第2部分.

$^{761}$那裡說什麼呢?.
《以賽亞書》是寫給那些被擄巴比倫.
他們在帝國裡面被鎮壓的那些賤民.
猶太人被人欺負.
《以賽亞書》就是寫給他們.
不要緊上帝要來幫你.
上帝會給你救贖.
會用祂的名來呼召你出來.
救贖你離開那個國度.
以至你能夠享受上帝給你的平安.
這就是《以賽亞書》下一半部.
裡面所說的話.
上帝的救贖會臨到你們這群被欺負的人.
你可能有很多恐懼.
有很多未知之數.
不過上帝說我與你同在.
因為我是創造你.
以及去救贖你的主宰你.
其中一節就是43章第一節.
其實耶穌的時候.
當時很多人.
去看聖經的時候.
其中一個最多人去引用的金句.
最多人引用的那段經文.
其實就是這一段.
《以賽亞書》43章.
40至55章裡面說到.
上帝要來拯救你.
以前你未認識上帝的時候.
上帝已經認識你.
因為上帝愛你.
這段聖經是多麼大福氣.
我們未認識上帝的時候.
上帝說你不要懼怕.
因為我救贖了你.
提你的名去召你出來.
因為你是屬我的.
我們中間有很多做父母的.
我還記得第一胎出生的時候.
記不記得那個情況.

$^{801}$我還記得我女兒第一胎出來的時候.
10月懷胎.
她還未出生.
在人世上未出現的時候.
我已經認識她.
已經有名字給她.
已經等她出來.
已經知道她的存在.
縱使她未知道我的存在.
當她出生的時候.
我就是用我給她的名字.
去叫她.
你記不記得產房的時候.
就算我未經歷過.
你想像都想像得到.
在產房的時候.
護士將我的寶寶放在我的手上.
我去叫她的名字.
她還未懂得說話的時候.
這個就是那個景象.
上帝說.
你不要害怕.
我救贖了你.
用你的名來召你.
你是屬於我的.
你要知道聖經裡面說到名字的時候.
不是我們今天說的名字.
今天名字純粹是一個標籤.
我們叫John就叫你John.
聖經裡面說到名字.
是說到權能.
不是純粹一個標籤.
聖經說到名字的時候.
是說到所屬.
你的所屬權在哪裡.
等於我們小時候.
有那些圖章.
有沒有圖章.
那本書是屬於你的.
你就用圖章去印上去.

$^{841}$那本書就是屬於你的.
在聖經裡面的文化都是一樣.
所以那些大的王.
佔據了一個地方的時候.
那裡的小王就把他的名字改了.
等於列王記裡面有很多這樣的君王.
都是這樣的.
巴比倫王來到的時候.
就把他的名字改了.
以前他的名字是甚麼.
不要緊.
現在改了新的名字.
西底加拉.
代表甚麼.
不只是給他標籤.
而是說你是屬於我的.
你在我的權能底下.
這裡耶和華說甚麼.
耶和華說我已經救贖了你.
Call you by name.
你是屬於他.
我知道我們在生命裡面.
經常會遇見很多.
Unknown.
或者很多混沌的時候.
今天這段聖經.
希望給你一點安慰.
給你提醒.
你在耶穌裡面.
神說我救贖了你.
已經用我的名字照料你.
你是屬於我的.
我們同心協力地祈禱.
上主我們多謝你.
因為你是我們的救贖主.
主我們在這個國度裡面.
我們不再有羞辱.
不再被黑暗的權勢所勝過.
因為我們的生命是屬於你的.
神愛多謝你對我們的愛.

$^{881}$對我們的恩典.
你的恩典臨到的時候.
我們在生命裡面.
就有開始這種改變.
我們的價值.
我們對世事.
對我們人生.
對我們自己.
對你的看法.
都要基本上有一個這樣的改變.
主多謝你因為你的福音.
你自己的信息.
你自己的救恩.
有一個生命改變的救恩.
主今天我們來到你面前.
我們知道你是那位創造我們.
救贖我們的神.
二主你自己.
這個生命的改變.
繼續在我們生命裡面去產生出來.
讓我們知道我們屬於光的國度.
屬於你自己的國度.
主讓我們的家庭.
讓我們的關係.
讓我們的生命.
我們的職業等等任何一件事.
慢慢繼續的.
去到兌現出.
上帝你自己的權能.
我們屬於你自己國度的國民.
求你去使用我們.
祈禱奉主名.
\newpage



\section{19:1-8}
\label{sec:LTR4cqb1AYI}
\textbf{What Are You Doing Here? (1 Kings 19\_1-8) - Tommy Lee}
\newline
\newline
連結: \href{https://youtube.com/watch?v=LTR4cqb1AYI}{\texttt{https://youtube.com/watch?v=LTR4cqb1AYI}} ~~~~ 語音日期: 2024-02-20
\newline
\newline
\hyperref[sec:aIlHRCRsv_k]{\small{< < < PREV SERMON < < <}}
~
\hyperref[sec:index_chronic]{\small{[返順時目]}}
~
\hyperref[sec:index_scriptual]{\small{[返順卷目]}}
~
\hyperref[sec:Lczwsjcf0tU]{\small{> > > NEXT SERMON > > >}}
\newline
\newline
$^{1}$Here we go. Let me turn this on..
We'll have a point with that one, Mark. I will promise I am not making a mess today, Mark..
Good morning. Today we finish up our series on Elijah..
And so we've been looking at chapters 17, 18, and 19..
They span about three years of Elijah..
And it is interesting how, for some odd reason, God works in mysterious ways..
I think all of us who are sitting here and when I travel and talk with different leaders around the world,.
they all want to know what God is calling them to do..
But you look at the life of Elijah, and it was in that three-year period that he says,.
"Okay, Elijah, this is what I want you to do. That's my calling.".
Whatever he did had such a ripple effect that we're still reading about him today..
We're still studying what he did in that three-year period..
Very similar to Jesus, very similar to Joseph or Moses, a certain period in time..
And after that, Elijah is whisked away, and we don't see Elijah until the Transfiguration,.
when he's there with Jesus and Moses..
And so, but at the same time, he calls Elijah to be faithful and continue to call us to continue to live your life,.
to be faithful each and every single day..
And for all of us, if there is anything we remember about Elijah, is that there are defining moments in your life..
You don't know when those defining moments are coming..
You don't know when God calls you to do something..
And when he does call us, it'll come out of nowhere, but we have to be faithful during that whole period of time.
in terms of waiting and anticipating for what God calls us to do..
But for all of us, he continues to call us each and every single day to live as Christ would be,.
to be faithful in how you act, how you think, how you treat other people,.
because the gospel motivates all of us to do that..
You look at Elijah, we first start in chapter 17, right?.
And he just comes out of nowhere. There's nothing about Elijah beforehand..
And you guys have heard me summarize this over and over again..
This guy is the Tishbit of Tishbe in Gilead..
That is the first description of Elijah..
He comes out of nowhere. He comes out of this rural area..
He's a Tishbe, and the author doesn't really know where he's at, but he says it's somewhere in Gilead..
And this guy, this rural guy, who's not very presentable, rough around the edges..
He's not tall and strong like Saul. Somehow he gets a meeting with King Ahab of the Northern Kingdom..
There's a Northern Kingdom. There's a Southern Kingdom..
And Ahab is the king of the Northern Kingdom. He gets a meeting with the king..
And he tells the king that there will be neither dew nor rain until my word..
And we will see that there is a huge drought that takes place for the next three years..
And he reminds him, "Why am I doing this? I represent the God before whom I stand.".
That's it. That's my motivation. That becomes a theme of Elijah's life..

$^{41}$The God before whom I stand..
And after that, he's whisked away to Shareth..
And there he lives in solitude. He's fed by the ravens..
He's fed drinks from the brook until it dries up..
And he stays there. And after that, God says, "You know where I need you to go?.
I need you to go a hundred miles. I need you to walk a hundred miles to Zarephath..
I see where you are. I know where you are. I see you..
I know where you're going. You want me to walk a hundred miles to Zarephath,.
the hometown of Queen Jezebel, the wife of King Ahab, who is the worst of all the kings? Yes..
And there you're going to encounter a widow. And that widow will provide you for all your needs.".
So he shows up, and this widow is sitting there. She's not of the faith and says to him,.
"Hey, I'm not sure if you are here because you need food. I have no food..
I have a little flour. I have a little oil. I'm collecting sticks so I can make a fire..
I'm going to make a little bread so my son can eat it and die.".
"Really, Lord? You sent me this?" "Yes. I know where you're going..
I know what you need. But you're going to have to overcome first impressions..
You're going to have to overcome physical impossibilities.".
Later on, as the story goes, the widow's son dies, and she blames him..
And she says, "Your God did this to us. He is punishing me for one of my sins.".
And he lives with that fear. Elijah takes the boy up, brings him to his room..
And you know what? He doesn't know the answer..
He comes and lays the boy before the Lord and says, "Lord, whatever, can you please heal him?.
Because I don't know what's going on. I don't know what you're doing.".
There's a genuine relationship that Elijah has with the Lord..
And somehow the Lord heals the boy. And the widow says,.
"If you guys look at the end of chapter 17, this woman who didn't know who the Lord was,.
was not part of the religion, says to him, 'Clearly you are a man of the Lord..
Clearly the Lord is present in your life.'".
She doesn't say, "Look at you. You're such a powerful man. You're such a strong man.".
No, no. "Clearly, clearly the Lord is present. You are a man of the Lord.".
18 comes in. It's the last time I saw you guys, Mount Carmel..
There, a lot of times after three years, God says, "Go confront Ahab.".
Yes. Time, drought is over. And then they meet and Elijah says to him, "Let's have a showdown..
You bring all your prophets to Baal. Bring all your prophets to the Asherah..
And I will bring my God. And we're going to ask those gods to come and bring fire..
And whoever calls on the Lord and they bring fire, that is the true God..
Yes, let's do it." And everyone comes. And all the Israelites of that land comes..
And if you remember, Elijah looks upon the Israelites..
And he says to them, "Are you limping between two different opinions?".
And they said nothing..

$^{81}$And it's an illustration for us, is for all of us, are we limping between two different opinions?.
And also, do you even care to do something about it?.
Story goes, the prophets of Baal call upon their God. Nothing happens..
And they keep calling and calling and calling. Nothing happens..
Now it's Elijah's turn. And the word over in 1st Kings tells us he repaired and rebuilt the altar..
Are you limping between two different opinions?.
Do you care to do anything about it? And so how do you repair and rebuild your altar?.
And also, are you asking God to show up?.
And boy, does God show up..
And the fire comes down from heaven, and it consumes all of the burnt offering..
It doesn't bring judgment on the people for their disobedience..
It consumes the burnt offering..
And this is where we are in chapter 19, as we wrap up our series of Elijah..
And so the story goes, "And in a little while, the heavens grew black with clouds, and wind, and there was a great rain..
The drought is over. After three years, the rain is finally there..
And Ahab rode and went to Jezreel..
And the hand of the Lord was on Elijah, and he gathered up the garment and ran before Ahab to the entrance of Jezreel.".
Why does Elijah show up to Jezreel?.
Now, the capital city is in Samaria. We talked about that, right?.
Capital city is in Jezreel. But there are little remnants of a small little palace in Jezreel..
People believe that Ahab and Jezebel probably had a summer home, a small little home over Jezreel..
And it's about 20 miles away from Mount Carmel..
And so a lot of times Elijah goes there. Why?.
Look at how the Lord works. He brought fire to Mount Carmel..
He destroyed all the prophets of Baal..
The people of Israel are going to turn back to God..
Ahab, Jezebel, they're all going to turn back to God..
And I'm going to be there to see it. This is the triumphant return of our Lord after this three-year drought.
because of everything that happened at Mount Carmel..
And so he gathers up his garment, and he gets to the entrance of Jezreel..
But that's not what happens..
The people of Israel doesn't come back..
Ahab and Jezebel don't turn their lives back to God..
But rather Ahab told Jezebel all that Elijah had done and how he had killed all the prophets with the sword..
It's funny. We talked about how all the 19 kings, Ahab was the worst..
But we truly see who the true villain of the story was..
Jezebel had all the power, and Ahab goes crying to her..
Then Jezebel sent a messenger to Elijah saying,.
"So may the God do to me more and more also, if I do not make your life as the life of one of them by this day tomorrow.".
Then he was afraid, and he arose and ran for his life and came to Beersheba, which belongs to Judah,.

$^{121}$and left his servant there..
Point number one, there are times when things do not go according to plan..
There are times when things don't go according to plan..
Elijah had it mapped out perfectly..
Elijah thought, "After three years, I make my triumphant return..
Lord, I've obeyed you. You've followed me. I mean, you've listened to me..
We've done all these things..
I've gone to Jareth. I've gone to Zarephath. I've come back..
Mount Carmel, I've called upon the name..
Clearly, this is the time. Clearly, this is what you're going to do.".
And he goes over to Jezreel, and he is expecting everyone to turn to the Lord..
But that's not what happened..
Things don't go according to the plan that we want at times..
And when that happens, are you okay with that?.
When we expect God to work in a certain way,.
when we expect God to do what we want him to do,.
and it doesn't go according to plan, are you going to be okay with it?.
How are you going to react to it?.
It's like the story of Dave and Goliath, isn't it?.
We've talked about this before..
The story of Dave and Goliath, because here's this young shepherd boy.
who is the one person that no one ever expected would battle Goliath..
He's the underdog in the story, and he kills Goliath..
And suddenly, we think he is going to be--.
songs are sang about him, stories are written about him,.
and this is his time to shine..
He's going to be the best warrior..
Saul is going to cede him the kingdom..
Saul is going to retire, and David now truly becomes king..
That's not what the story happens..
Saul becomes jealous of him, and eventually Saul takes away his job,.
and Saul goes after him, wanting to kill him..
David runs to his wife. Saul goes after him..
David then runs to his mentor Samuel. Saul's after him..
He runs to his son Jonathan, who runs to his best friend Jonathan,.
who is the son of Saul. Saul finds him there..
He goes to the city of Gath, where Goliath is--that's his birthplace, his hometown..
He stays there for a period of time, pretending to be a crazy man,.
and he ends up in the cave of Adullam by himself,.
where he writes Psalm 57 and Psalm 132..

$^{161}$Things don't go according to plan..
And when things don't go according to plan, how are you going to react to it?.
I remember in my own life as well, too, when I first started Resource Global..
This was about 10 years ago..
The premise of Resource Global was that young Christian marketplace leaders,.
young Christian business leaders--not pastors, not seminary students,.
not ministry leaders--young Christian marketplace leaders.
would be the catalyst for gospel growth all around the world..
That was the catalyst. That was my thought..
And the Foundation says, "Tommy, I love this idea..
Here's some funding." And we started in Jakarta, Indonesia..
And the plan is once we start with the Chinese Indonesians in Jakarta, Indonesia,.
and then we started one in Chicago, and then we're going to do something in Malaysia..
We're going to do something in Hong Kong. We're going to do something in Nairobi..
We're going to do something in LA. That was the plan..
And I had it all laid out..
And I remember I brought in my first team of people..
I hired a director of communications. I hired a program manager..
That was the plan. That's what I believed that God was calling me to do..
But that plan didn't work out..
I didn't really get along with those individuals..
We were on the same page. We were constantly disagreeing..
And I thought, "Lord, you told me to do this. I really believe you told me to do this.".
I hired a team. Why isn't this happening?.
So we got rid of the team. And I said, "Okay, okay..
I probably misinterpreted you, right? I'm going to bring in a second team..
And I'm going to bring in a team who's completely different to me in terms of personality..
That will work." It didn't work. It didn't work..
"Lord, I'm so sure that you called me to do this.".
Next thing you know, as the story goes, that was about five, six years ago,.
I was diagnosed with cancer, and then I had to leave for eight months..
And I thought we were going to shut it down..
But that's when the ministry started growing..
Jakarta started growing when I wasn't around..
Then from that point on, Malaysia started. Singapore, when I was just starting..
Nairobi started when I wasn't around..
Our young leaders in Jakarta became the leaders in their city..
They became the mentor. They became the donors..
They became the board members..
All because they started rising up to the challenge..

$^{201}$Now here we are, 13 cities later..
We just started Cluj. We're working on Tokyo at this point right now..
I'm off to Cluj in a couple weeks..
And I look back and I said, "I would have never scripted this.".
It doesn't go according to plan..
That doesn't work in the way that you think it'll work..
Elijah had to understand this..
Elijah gets up there and he thinks that Mount Carmel, everyone's going to go according to plan..
You're going to respond. Everyone's going to turn back to the Lord..
It didn't happen that way..
And when it doesn't go according to plan, are you okay with it?.
Another thing also as well too is, and so Elijah then in verse 4,.
"But he himself went a day's journey into the wilderness..
And he came and sat down under a broom tree..
And Elijah asked that he might die, saying, 'It is enough now, O Lord..
Take away my life, for I am no better than my father's.'".
Here's something, a side note as well too..
Different people respond in different ways, right?.
Different people respond in different ways..
And that's something we have to honor and respect..
Sometimes with different people, when we come across a situation,.
we just always have to understand that people respond in different ways..
One of the things that we have seen, and I think a lot of times you're going to see as we map out the story,.
Elijah gets discouraged easily..
And that's okay. It's understanding that personality..
But here's another personality I want to show you..
Isaiah chapter 6. You don't need to turn into it..
Isaiah chapter 6 talks about this story..
And here's where it begins with. Isaiah 6, and right in verse 10..
"And I heard the voice of the Lord saying, 'Whom shall I send, and who will go forth for us?'".
That's what the voice of the Lord says..
Isaiah says, "Here am I. Send me.".
And here's what the Lord says to him. Isaiah 6, verse 10..
"Keep on hearing, but do not understand. Keep on seeing, but do not perceive..
Make the hearts of the people dull, and their hearts heavy..
And blind their eyes, lest they see with their eyes, and hear with their ears,.
and understand with their hearts, and turn, and be healed.".
Then Isaiah says, "How long, O Lord?".
And he says, "Until cities lie waste without inhabitants, and houses without people,.
and the land is a desolate waste, and the Lord removes people far away,.

$^{241}$and the Israelite places are many in the midst of the land..
And though a tent remain in it, it will be burned again.".
Do you understand what Isaiah is saying as well too?.
Here, Elijah gets mad because he sees everything that's happening..
He's all discouraged, right?.
Isaiah is literally being told in Isaiah 6, "I'm sending you,.
and when you preach, you will see no results..
You will keep putting all this effort, and you will keep preaching the Word of God,.
and you'll keep serving, and serving, and serving,.
and you will never see a lack of results.".
And Isaiah is saying, "So how long am I doing this?".
"Until nothing remains..
That's when you'll be done.".
But Isaiah is strong..
Different people respond in different ways..
And we see here with Elijah, a lot of times, it's nothing wrong,.
but he just needs a little encouragement along the way,.
which leads us to the next point..
Because when things--when we get discouraged--.
I love what Tim Keller says on this one..
He says, "If you look at this passage in Isaiah--in Elijah in 1 Kings 19,.
God uses the story in Elijah where he encounters both the angel--.
an angel--and the Lord..
Both of those guys are showing up in this particular situation,.
and we see God ministering to Elijah in three ways..
He ministers to him physically..
He ministers to him mentally..
And he ministers to him spiritually--.
physically, mentally, and spiritually..
Let's go through that and see what this means..
Verse 5, "And Elijah lay down, and he slept under a broom tree..
And behold, an angel touched him..
And he said to him, 'Arise and eat.'.
And he looked, and behold, there was at his head a cake.
baked on hot stones in a jar of water..
And Elijah ate and drank and lay down again..
And the angel of the Lord came again a second time.
and touched him and says, 'Arise and eat,.
for the journey is too great for you.'.
And he arose and ate and drank.

$^{281}$and went in the strength of that food forty days and forty nights.
to Horeb, the Mount of God.".
So I know Horeb is a very special place..
It is the place where Moses received the Ten Commandments..
That's where Elijah goes..
So how does God minister in the midst when we get discouraged,.
when we get frustrated for some reason or another?.
And here in this situation, there are unmet expectations.
that Elijah is dealing with..
And after all the things that's happening, he is frustrated..
But God, in the midst of our frustration,.
ministers to us in three different ways--.
physically, mentally, and spiritually..
For the physical part, he comes to Elijah,.
and he allows him to sleep..
That's number one..
He allows him to sleep..
Sometimes we just need to--.
when we're in the midst of situations, when we are discouraged,.
we just need to walk away for a little bit..
We have to do something different..
We have to regroup ourselves..
He allows him to sleep..
He also feeds him, not once but twice,.
"Arise, get up, and eat.".
That cake sustains him for forty days.
while he travels to Mount Horeb..
And then from that point on, he touches him..
I love what one author says to him,.
is that we are physical beings.
who live in the physical world..
The things of words of affirmation,.
touch, all of those things,.
mean something to us because we are physical beings.
in the physical world..
When I worked at Moody Bible Institute,.
I remember I got a report about the number of books.
that Moody-- in terms of their revenue sales, right?.
And so did you know under Moody Publishers,.
there is one book that makes up 90\% of their revenue.

$^{321}$every single year?.
If you guys could guess, that book is.
The Five Love Language by Gary Chapman..
After all these decades,.
The Five Love Language of Gary Chapman,.
which is, I think, in its eighth revision,.
in eighth edition or something like that,.
makes up 90\% of the revenue..
And that man is about 80 years old,.
and he is still out there,.
and that message never gets old..
Why? Because we all want to know how we're wired..
We all want to know how to love..
We all want to know how to show love and to receive love.
because we are physical beings in a physical world..
God ministers to us physically..
Secondly, another thing is,.
he ministers to us mentally, emotionally..
What does he say to us?.
And a lot of times this goes back to Elijah..
Remember, he gets discouraged easily..
He says, "Arise and eat," in verse 7..
For what?.
The journey is too great for you..
For some of us who are parents,.
for some of us who are counselors,.
for some of us who are teachers,.
for some of us who are leaders,.
have to manage staff,.
you'll understand what I'm trying to say..
We manage people differently..
Different people are managed differently.
according to their persona, personality, all of that..
A good parent, a good leader, a good counselor, a good boss,.
understands how each person reacts to different situations..
You'll understand..
What do I mean by that?.
There are some people,.
I have my two daughters here with me this morning,.
and so there are some people where you just have to encourage them..

$^{361}$They just need a lot of encouragement, right?.
There are some people we manage,.
there are some people, some of our kids,.
or there are some people that we work with, youth and all that stuff,.
where you have to be very direct with them..
You've got to be very hard on them,.
and a lot of times what helps them.
is when you're sitting there and telling them the truth..
Some people, they just need a lot of encouragement,.
but there are some people who you've managed,.
you kind of let them fail..
You don't say anything to them,.
because you know the best way for them to learn.
is when they fail and they're able to work through their situation..
They know what to do,.
but they just need you to support them..
Each person is managed differently..
The angel, the Lord, God knows us so uniquely sometimes,.
where not only does he heal us and minister to us physically,.
we see here a concrete example.
where he ministers to us mentally and emotionally..
He knows how Elijah is wired..
He needs encouragement..
Elijah doesn't need someone to tell him.
and lecture him of what he needs to do..
What he does need someone is to tell him, reaffirm him,.
I know the journey is too great for you..
Rest, rest, because the journey is too great for you..
Exit angel, get to Mount Horeb..
He comes before her in the cave, and now the Lord comes to him..
And the Lord says to him in verse 9,.
"What are you doing here, Elijah?".
It's very similar to Genesis chapter 3,.
where God says to Adam and Eve, "Where are you?".
Did God not know where Adam and Eve was?.
Did God not know what Elijah's doing?.
No, no, no, no..
He's asking questions, which forces Elijah to reflect..
All of you who have counseled people.
and come alongside people, you realize.

$^{401}$the worst thing when someone is dealing with a problem,.
they are not expecting you to sit down.
and tell them all your answers..
When my wife sits there and tells me about her day.
and tells me a problem with her kids in her first grade class,.
she is not expecting me to solve her problem..
She is expecting me to listen..
And sometimes in the midst of doing that,.
you ask clarifying questions to lead them to the point..
And sometimes the best counselors are the ones.
who allow people to talk 80\% of the time,.
and now you're talking 20\% of the time..
You're asking them, you're leading on a journey..
That's what's happening here..
"What are you doing here, Elijah?".
It forces him to think..
It forces him to reflect..
It forces him to come upon his own conclusion..
God ministers to us physically..
He ministers to us emotionally..
And he ministers to us spiritually..
Here the story comes, and one of the best parts.
of Scripture that I love..
If you have never, ever read this part of Scripture,.
it is a story that I love where God shows up..
Verse 11, God says to him, "Go out and stand on the mount.
before the Lord.".
"And before and behold, the Lord passed by.".
He is now Elijah's Mount Horeb..
That is the place where Moses got the Ten Commandments..
"And a great and strong wind tore the mountains.
and broke the mountain in pieces,.
the rocks before the Lord..
But the Lord was not in the wind..
And after the wind, an earthquake..
But the Lord was not in the earthquake..
And after the earthquake, fire..
But the Lord was not in the fire..
And after the fire, the sound of a low whisper..
And when Elijah heard it, he wrapped his face in his cloak.

$^{441}$and went out and stood at the entrance of the cave..
And behold, there came a voice to him that says,.
"What are you doing here, Elijah?".
God ministers to us spiritually..
He reminds Elijah, and he reminds us,.
"You expect that I'm going to come in power..
You perceive power in the form of a wind,.
an earthquake, and fire..
That's not how I'm going to show up..
You expect me to come in the form of judgment..
That's not how I'm going to show up..
You got to hear my voice.".
And that voice is surrounded by grace..
And that voice is soft..
And that voice comforts..
And through that voice, you experience my presence..
It's sort of like the gospel, isn't it?.
In Matthew chapter 11, when Jesus comes and starts ministering,.
John the Baptist, if you guys remember the story of John the Baptist,.
what was he saying to the Israelites?.
"Behold, someone greater than me is coming..
The ax is at the root of the tree..
Repent, repent, repent..
The ax is at the root of the tree.".
John gets imprisoned..
In verse 11, he hears about all the things that Jesus is doing..
He sends his disciples over to Jesus..
And he relays a message, and they say, "John wants us to ask you,.
are you the one who is to come,.
or should we expect someone else?".
Because Jesus' message was about healing and grace and forgiveness..
John was expecting someone who was going to be sitting there with fire and brimstone.
and was going to preach judgment..
And that's why he was telling everyone to repent..
But Jesus, you're coming to heal your--.
wait, are you the one who's supposed to come,.
or should I be expecting someone else?.
Elijah, when I come in the form of you think I'm going to be in the wind,.
you think I'm going to be in the earthquake,.
you think I'm going to be in the fire,.

$^{481}$these are all forms of judgment..
You've got to hear my voice,.
and that voice is surrounded by grace..
And as we end here, he wraps up,.
and he tells them, "Here's what I need you to do,.
because I have a plan..
Not only am I going to minister to you,.
but here's our plan..
After all, in verse 15, "And the Lord said to him,.
'Go, return on your way to the wilderness of Damascus,.
and when you arrive you shall anoint Hazel to be king over Syria,.
Jehu the son of Nimshi you shall anoint to be king over--.
Jehu the son of Nimshi you shall anoint to be king over Israel,.
and Elisha the son of Shabbat of Ebo-Meholah.
you shall anoint to be prophet in your place..
And the one who escapes from the sword of Hazel shall Jehu put to death,.
and the one who escapes from the sword of Jehu shall put Elisha put to death..
Yet I will leave 7,000 in Israel,.
all the knees that have not bowed to bow,.
and every mouth that has not kissed them..
Elijah, you think I had no plan?.
You think you are the only one left?.
I have 7,000 people waiting for me to be used to them,.
and I have a plan to put all these kings in place,.
and I've also found your replacement in Israel too,.
so you're out of a job..
That's my plan..
You think you're the only one left..
You think after Mount Carmel all of these things were going to happen..
That's not what's going to happen..
Because things don't go as we expect it to be..
We think that we could put God in a box,.
and all these things that will happen will go according to plan,.
but it doesn't happen the way we expect it to go..
That's the stories of Elijah's life, isn't it?.
That's the way--it's interesting..
Three years of this man, right, comes out of nowhere..
He kind of shows up, and he goes from place to place to place,.
all of these places, and at the end of the day,.
God ultimately has a plan..

$^{521}$He knows what's going to happen..
Don't put God in a box..
Sometimes your plan of what you expect God to do,.
it may not go according to plan..
And when that happens, are you okay with it?.
That's the story of Elijah,.
the Tishbit of Tishbe in Gilead,.
who comes out of nowhere,.
who we didn't expect him to do anything,.
shows us what it means to follow the Lord,.
to be obedient to the Lord..
In summary, you're wondering,.
"What in the world is he doing with this,.
and why is he putting water in this?".
That when we look at Elijah, we look at who he is,.
don't put a person in a box..
Just because he's some real guy who comes out of nowhere,.
and he's kind of not fancy, and kind of rough around the edges,.
God could use that person..
Don't put a person in a box..
Sometimes when we do things, and we see God work,.
or we try to perform God, we sit there and do all these things,.
and we come before the Lord expecting certain things to be done,.
but things may not go according to plan..
And that's okay, because sometimes,.
no matter how much water we put in,.
it may not even work..
God works in mysterious ways..
Oops, we're going to have to work on the illustration..
That's the story of Elijah..
That's the story of what God calls us to do..
The tishbit of tishbi in Gilead..
Let's pray..
Father, thank you for this time..
Thank you for the story of Elijah..
Thank you for what you taught us..
Thank you for how you ministered to us..
Thank you for the example that you gave,.
and how you showed how one person worked,.
and for you truly followed after you..

$^{561}$May our theme of constantly, no matter what,.
even when we don't understand,.
allow us to seek after you..
In Christ we pray, amen..
\newpage



\section{箴言 1:1-7}
\label{sec:Lczwsjcf0tU}
\textbf{用盡一生的腦袋 (箴言 1\_1-7) - 譚文倩莊傳道}
\newline
\newline
連結: \href{https://youtube.com/watch?v=Lczwsjcf0tU}{\texttt{https://youtube.com/watch?v=Lczwsjcf0tU}} ~~~~ 語音日期: 2024-02-20
\newline
\newline
\hyperref[sec:LTR4cqb1AYI]{\small{< < < PREV SERMON < < <}}
~
\hyperref[sec:index_chronic]{\small{[返順時目]}}
~
\hyperref[sec:index_scriptual]{\small{[返順卷目]}}
~
\hyperref[sec:ON4fMvKp6vs]{\small{> > > NEXT SERMON > > >}}
\newline
\newline
箴言 1:1-7
\newline
\begin{longtable}{cl}
\hline
\hline
章節 & 經文 (和合本修訂版)\\
\hline
1:1 & \begin{tabularx}{0.7\textwidth}{X} 大衛的兒子，以色列王所羅門的箴言： \end{tabularx} \\ \\ \relax
1:2 & \begin{tabularx}{0.7\textwidth}{X} 要使人懂得智慧和訓誨，明白通達的言語， \end{tabularx} \\ \\ \relax
1:3 & \begin{tabularx}{0.7\textwidth}{X} 使人領受明智的訓誨，就是公義、公平和正直， \end{tabularx} \\ \\ \relax
1:4 & \begin{tabularx}{0.7\textwidth}{X} 使愚蒙人靈巧，使年輕人有知識，有智謀。 \end{tabularx} \\ \\ \relax
1:5 & \begin{tabularx}{0.7\textwidth}{X} 智慧人聽見，增長學問，聰明人得著智謀， \end{tabularx} \\ \\ \relax
1:6 & \begin{tabularx}{0.7\textwidth}{X} 明白箴言和譬喻，懂得智慧人的言詞和謎語。 \end{tabularx} \\ \\ \relax
1:7 & \begin{tabularx}{0.7\textwidth}{X} 敬畏耶和華是知識的開端；愚妄人藐視智慧和訓誨。 \end{tabularx} \\ \\
[1ex]
\hline
\hline
\end{longtable}
$^{1}$早晨各位,這兩個星期我都在Chicago,所以今天我們暫停Mark Sir的港島系列Chosen,不用失望,下個星期會繼續,到時記得準時出席,今天大家稍安勿躁,先聽我分享..
今年我用Timothy Keller和他太太選寫的每日心言去靈修,每日只有一版,不是很長,.
久不久有一兩版或一兩節經文對我特別深刻或很有趣,那幾天我會想起靈修篇章的心言.話說有日我和我老公出去食飯,.
最近一次,我經常笑我老公有選擇困難症,因為無論每次我們出去食飯,去那間餐館去過多少次都好,.
他會望著餐牌,要用很長時間去研究餐單,無論裡面有多少選擇也好,餐館內去過多少次也好,他都要用很多時間,很多掙扎才可以下單..
不單止這樣,還要食一食,通常他會說早知不要點這個,其實我應該點那個好些,因為那個似乎多些食物,.
抵食些,甚至有剩餘的可以帶回家,但價錢差不多.今次和他食飯,這件事當然是歷史重演,.
我突然間腦海裡面出現一句最近我靈修時的經文,在《箴言》26章11節,.
他說:如未人一再做蠢事,如同狗轉身吃自己吐出來的東西.我和我老公說:Daddy,《箴言》裡面有這句經文,.
在你這個情況下,是否蠢人一再做蠢事,如同狗轉身吃自己吐出來的東西..
我老公看著我,扁嘴,低頭,《箴言》真是句句精金.今天講到的主題,用盡一生的腦袋,經文出自《箴言》1章1至7節,.
我們先低頭祈禱.親愛的主,我們一生,願也都尊你的名為聖,願你的國降臨,.
願你的旨意行在地上,如同行在天上,你祝福我們用心靈和誠實去敬拜你,虛心聆聽你話語的人..
我先跟大家讀一讀這個《箴言》1章1至7節..
大衛的兒子以色列王所羅門的箴言,要使人懂得智慧和憤悔,明白通達的言語,使人領受明智的憤悔..
就是公義,公平和正直,使愚民靈巧,使年輕人有知識,有智慧,智謀..
智慧人聽見,增長學問,聰明人得著智謀,明白箴言和曉悅,懂得智慧人的言辭和謎語,敬威耶和華是智識的開端..
如往人藐視智慧和憤悔..
箴言1章1開始,所羅門就解釋箴言的目的,是要讓讀者去研讀,去思考箴言,從中得著智慧去生活..
人的一生有很多起起跌跌,充滿了難以預料的事情..
你有錢或沒有錢也好,煩惱自然會找上門..
面對每天的挑戰,箴言所提供的智慧,幫助人有洞察的能力和理解能力去辨別出,作出有智慧的選擇..
有知識不代表有智慧,因為知識只是在說對事實的認知,智慧是將所知道的,所認識的運用出來..
例如人可以明白金融系統的運作方式,智慧是在說你能夠利用這個金融系統的運作方式去管理金錢..
使人領受明智的憤悔,就是公義,公平和正直..
能夠對事物有正確的概念,也有清澈明確的想法,讓我們有能力去辨別有智慧的教導..
我們聽了很多東西,但究竟哪些才是有智慧的教導呢?.
區別真偽,善惡,理解一些有智慧的說話..
有領悟,有判斷能力防範錯誤,並且將所學的應用,甚至幫助我們身邊的人..
菲勒比書一章十節,使你們能分辨是非..
使愚民靈巧,使年輕人有知識有智慧..
《箴言》十四章十五節,無知的人什麼話都信,通達人謹慎自己的腳步..
當然,智慧對我們所有人都很有用,我們都想擁有..
但特別對一些思想比較單純的人,就更加需要聽到多一些智慧的建議,能夠重組獲得益處..
幫助那個人思考,釐清真偽,避免做錯事..
當然更重要的是逃避誘惑,或者中了壞人的詭計..
像所羅門所說,壞人的外表可以很善良,但其實他們就像蛇一樣假惡..
青年時期是接收各種訊息的年紀,我們經常聽到年輕人就像一個海綿一樣,.
很容易吸收很多東西,他們心智未成熟,容易衝動,性急,很多時候想做就去做,容易受人唆擺,不考慮後果..
我想我們在座每一個人,要去想一想的時候,都會想起我們小時候做了多少蠢事..

$^{41}$現在的確是一個資訊的年代,不同的資訊其實是誰手可得,你只要有個電話,基本上你想知道什麼,你打進去,你就可以找到答案..
只不過你不知道這些資訊是從哪裡來的,對你有益還是有害,是真還是假,或者製造這些資訊的人,背後到底有什麼動機..
我讀神學院的時候學了一個很重要的功課,就是每一次當我們要引用一些文章,或者一些資料的時候,.
我們都需要探究清楚知道我們所引用,接收的訊息是從哪裡來的,是誰寫的呢,是真還是假的呢,.
不是每本書都是好書,不是所有資訊都是正確的..
今時今日我們看的新聞,其實都是要千挑萬選,因為有些新聞已經不是再用公正無私的態度,客觀的去報導新聞,.
只是淪落成為了國家的宣傳工具.還記不記得在2000年,香港有不少騙徒,宣稱有千年蟲的蛀蟲藥,.
有騙徒對千年蟲問題一知半解的市民去購買,以訛傳訛,你怕我又怕,.
有部分不熟悉電腦的市民就以為千年蟲好像一個瘟疫一樣,一個傳染病一樣的病毒,.
花了很多錢來買這些解藥,被人騙錢.上星期我才有個朋友跟我說了一個笑話,.
他說他在一套劇集裡面,其中一幕講到一條村裡面的村民衝著去教會信耶穌,.
因為教會門口有一塊牌寫著信耶穌得水牛,.
個個信以為真,走了去教會等牧師派水牛,.
但其實教會門口那塊牌是日久失修,永生這兩個字少了一點,同一畫,.
雖然聽起來很好笑,是一個笑話,但想想如果真的發生好像不太奇怪,.
智慧人聽見增長學問,聰明人得著自謀,明白真言和巧語,懂得智慧人的言辭和謎語..
說開又說,你身邊有沒有一些你信得過的弟兄姊妹,.
能夠在你的生命裡面,你的生活裡面面對一些困難,或者一些你掌握不到的事上,.
能夠跟你聊天,能夠跟你一起分析,給你意見,跟你一起祈禱呢?.
忘記多少歲開始,我發覺自己想事情不是很全面,我是一個過分樂觀的人,.
由那時候開始,我很喜歡找一些認識我,跟我很親密的弟兄姊妹一起分享,去問他們意見,.
慢慢這個動作就成為了我的習慣,多年來,加加減減,今時今日我存了一班弟兄姊妹在我身邊,.
在不同的事情上,我都可以跟他們一起分享,有時這件事我去問A,那件事B比較厲害,我會找B去分享..
這班跟我一起追求成長的智囊團,陪伴了我很多年,跟我一起成長,看著我掙扎,支持我,.
在需要的時候跟我一起祈禱,跟我一起仰望神,同樣我都是他們的智囊團..
《箴言一章五節》這裡,正正就是在說我們人願意謙卑聆聽,尋求意見是多麼重要,不是靠我們自己..
更加重要當然是我們找一些比較淑齡的弟兄姊妹,有追求心的弟兄姊妹一起去分享,一起去成長..
願意聽意見受教的人,在《箴言》這裡,就是在說他們是一些有智慧的人,.
懂得找適當的人去尋求意見,可以對事情有更透徹的了解..
說起問問題,說起來我們基督徒很喜歡問神的心意,神的旨意,.
我們很希望當我們問的時候,神有一個很清澈的答案給我們,最好就幫我選好,不用我煩..
其實什麼是神的心意,什麼是神的旨意?.
如果你經常問這個問題的時候,特別是在你遇到困難,想神幫你選的時候,其實這個問題是問錯了的..
因為神的心意,神的旨意不是一個功能,不是用來幫人解決問題,化解疑惑..
如果你是真心想知道什麼是神的心意,神的旨意的時候,其實答案就在聖經裡面..
在聖經裡面,神已經很清清楚楚地告訴我們他的心意,他的旨意..
例如約翰福音13章34節,這裡說你們要彼此相愛..
亞馬太福音6章33節,要先求他的國和他的兒..
約翰一書2章15節,不要愛世界和世界上的東西,若有人愛世界,愛父的心就不在他裡面了.等等..
這些就是神的心意,這些就是神的旨意..

$^{81}$在聖經裡面清楚列明,可能不是人想找到的答案,甚至可能聽過很多次,你沒有什麼特別感覺..
要先求他的國和他的兒,這些東西都要加給你們..
是一個命令,也是一個應許..
如果你花點時間去消化這節經文,其實我們做基督徒是很有福氣的..
我們先追求神的兒和神的國,其他東西是自然神就給你的,你不需要擔心,你不需要去求..
經威耶和華是知識的開端,如往人藐視智慧和憤悔..
真言裡面不單止一章七節提到經威耶和華是知識的開端,.
真言九章十節,經威耶和華是智慧的開端,認識自聖者便是聰明..
上年我去伊拉克短孫的時候,我認識了一位姐妹,她在香港美國的英特爾工作了27年,她越做越高級..
經常會坐公司的飛機到處去工作,逢星期一就飛,星期四晚上就會回到家..
每個星期都如是,因為她經常不在家,她的女兒長期不開心,因為一會兒要去公公婆婆家住,一會兒又奶奶爺爺,一會兒可能會去另外一些人家..
她的工作很繁重,壓力很大,她說當她的同輩,她的同事升職快過她,那個壓力會更加大,會令她懷疑自己的能力..
單親媽媽做公養家,照顧女兒,處理家裡的事,她結果捱到身體出現毛病..
當她病得快要死的時候,她突然開始在想,我是不是快要死了?.
如果我突然死了怎麼辦?她在想,其實我不認識神,我雖然是一個基督徒,但我不太認識神,我不太認識耶穌..
她一想到如果我死了之後上天堂,見到神,三不識七,還要24/7天天對著祂..
她說再想想,還要永遠,她說想起我會很害怕,我不想和一個陌生人長期對著,她說實在很難想像..
因為那場大病,她重新檢視生命,後來她康復了,她決定要滴起心肝去認識給她生命的主..
沒多久公司有個不錯的機會,她乘機把握機會提早退休..
今時今日她已經退休了八年,她過著一些很簡樸,很簡單的生活..
一個星期有兩天去參加Bible Study Fellowship,查經團契裡面去查經和服侍,教會有需要的時候去教會幫忙..
和她聊天的時候,我覺得有很大的得著,我很欣賞她,覺得她很有智慧..
我也想起她是一個很謙卑的人,我和她去短孫的時候,她做的崗位其實是很卑微的,其實她是全部都懂的..
做她的頭的那個,反而沒有她那麼懂事,但是她從來沒有說過一句話,令你覺得她很厲害..
她也沒有挑戰給她工作的那位姐妹,有些事情可能可以做得更加好..
她跟我說她很感恩有這場大病,令她醒過來,她多了很多時間可以在家陪她女兒,也可以參加查經班去認識神..
在人的眼中,她當然有一份很令人羨慕的工作,賺很多錢,她住在一個很貴的社區裡面..
但是她沒有留戀這些東西,她也沒有留戀以前花錢買東西,她從來不需要想,不用看價錢,不需要再去找更豐厚的人工升職..
生命裡面的人和事,好像重新排位過,因為她現在很清楚知道在她生命裡面,什麼是最重要..
《箴言》二十四章三至四節,房屋因智慧建造,因聰明立穩,又因知識屋內充滿各樣美好寶貴的財物..
詩篇一百二十七篇一至二節,若不是耶和華建造房屋,建造的人就枉然勞力.若不是耶和華看守城池,看守的人就枉然警醒..
你們清晨早起,夜晚安歇,吃奴奴得來的飯本是枉然,唯有耶和華所親愛的必叫他安然睡覺..
若不是耶和華建造,建造的人枉然勞力.早上起床,夜晚休息,吃奴奴得來的飯本是枉然..
願望人藐視智慧和憤悔,《箴言》一章七節提醒,基本上只有傻瓜才會用我們一生的腦袋去建立我們眼前或今世以為能滿足,能給我們幸福安穩的生活..
卻輕視從神而來的智慧和憤悔,有屬靈遠大的眼光去鋪排永恆的生命..
2023年接二連三,我們家裡也收到很多沉重的消息,其中六月的時候,我小時候的偶像和我很親近的小表姐離開人世..
由她覺得不舒服入醫院,兩個星期後確診肺癌末期,再不夠兩個星期她就離開了世界..
她入院之後,我很多時候都會跟她互傳短訊,因為我知道她每天晚上都睡在病床上,特別是在日本凌晨的時分,.
我知道她白天眼睛都不見,所以我特別多跟她聊天..
如果那天我出門,我會拍一些照片給她看,因為有些地方是我們小時候一起去過的,我們一起懷緬過去..

$^{121}$她離開了芝加哥差不多三十年,當中只回來過幾次,所以有很多地方都不同了,有些新的東西,我又會拍照給她看,跟她介紹一下..
她住在醫院的時候,她跟我說,很諷刺,就像Katrina那樣說,她說我一直在減肥,有一個目標的磅數,她說進來之後瘦了幾磅,.
我現在的磅數就是我一直夢寐以求的.我們又會說回我們小時候,以前我去找她,在華富村一起玩,一起出街的時候,.
剛開始移民的時候,我們兩位表姊妹住在同一個房間,很想念香港,我們兩個一起哭了一個月,甚至很多年後,我去日本探望她,跟她傳福音,我們談到死亡的問題..
每天如是,傳著傳著的短訊,有一天她突然開始沒有回覆,一天沒有回覆,兩天沒有回覆,第三天我的大表姐告訴我,她已經離開了這個世界..
她入院之前一點舒服,一點不舒服都沒有,還安排了一家人準備去旅行.12月初回香港之前,我們收到一個消息,我們的伴郎在香港入了醫院..
回到香港之前,我們有點擔心,不單單是擔心他,其實我們在想,不知道香港的醫院情況如何,我們回到香港之後,有沒有可能去探望他呢?.
他12月底生日,我丈夫當時希望回到香港的時候,沒有多久出院,我們可以跟他一起慶祝生日.結果我們還沒有上機,他已經離開了..
消息來得很突然,我們基本上好像完全消化不了,回到香港,結果我們在他生日那一天,去參加他家人為他舉辦的追思會..
2023年終於過去,回到香港之後,回到美國,其實我一直都很忙.2月稍為回復正常的時候,有一天坐在客廳,我丈夫問我,.
媽媽,去年發生了那麼多事,又有幾個跟我們那麼親密的人離開了,其實對你有什麼影響?你猜我怎麼回答他?.
這一分鐘好好的,還在計劃去旅行,突然之間不舒服進了醫院,不到一個月就離開了,大家都還沒有消化到他病重,轉頭已經離開..
剛才說的那兩個人跟我年紀差不多,你問我對我有什麼影響?我在想的是,怎麼會沒有影響?.
有人可以每天做4分鐘的plank,40個掌上壓,少糖少鹽,少澱粉質,只是吃有機食物,早睡早起,不過分操勞,操練有好的EQ..
但你知不知道這樣做也好,其實也不保證你健健康康,長命百歲..
我有個朋友熬到五個兒子都終於出生了,他獎勵自己買了一輛Jaguar,搬到downtown,想著可以走路上班,.
結果車子落地不久就中風,左邊身體沒有活動能力,坐輪椅..
他太太結果把Jaguar換成了minivan,方便推他進去..
想想人生是不是真的就是這樣,好像傳道書那樣形容,一章二至三節說,虛空的虛空,虛空的虛空,全是虛空..
人一切的牢獄,就是他在日光之下的牢獄,有什麼益處呢?.
人的一生是不是就是這樣,是不是就是這麼恐怖,就是這麼虛空?.
我們窮一生的精力去計劃,去籌算我們的人生,去到這個年紀,可能有時還要想著,.
怎樣也不敢吃,那樣也不敢吃.我昨天才跟一個朋友說,他舅父年輕讀書的時候很窮,.
天天吃雞,只可以吃雞,到年紀大的時候身體不好,醫生叫他最好多吃點雞胸,.
因為這個人是醫生,到我已經有錢了,我什麼都可以吃的時候,結果還是天天吃雞..
我們看N這麼多篇的文章,聽N這麼多個朋友的意見,我們都是想確保我們的生活能夠有保障,.
我能夠安心,我能夠得到最好的資訊去幫助我去過一個平平穩穩,很安舒的人生..
但是我想你和我都知道,這個世界不是這樣的,當你的投資持續虧損,.
努力經營的得不到回報,突然病起來,生命倒數,死亡臨近的時候,.
所有你平時很緊張的東西,你看得很寶貴的東西,突然全部都沒有意思..
睡在床上的時候,無論你的戶口裡有多少錢,你衣櫃裡有多少漂亮的衣服,.
房子有多大多漂亮,我想起我自己,廚房裡有多少我看為很寶貴的沙煲冷炒,.
或者你在人面前有多麼的德高望重,誰在乎?.
用一生的腦袋去追求今世的款足,但有時我們甚至沒有好好經營我們和我們家人的關係,.
和教會弟兄姊妹的關係,有沒有想過病起來沒有人探你,很慘..
武漢重掌管生命創造你,給你生命的神,去問他,.
知取智慧,其實我應該怎樣去過我的一生,是先後次序的問題..
弟兄姊妹當你看到生命的無常,或者像吳牧師禱告提起有兩位姊妹昨天離開,.
像我這樣去年也有些和我們很親密的至親朋友離開,.
你看到人多麼無能為力失控的時候,你是否只懂得擔心,.

$^{161}$發財賣財去找一些方法去確保自己長命一點,不要生病,.
你是否很怕死,所以經常都憂憂愁愁,終日唉聲嘆氣..
其實死,我們每個人都有一死,走不掉的,.
Life is uncertain,但是Death is certain,問題不是我們想想要做些什麼不用死,.
因為死是必然的,但是會不會是需要多花時間去想想我們怎樣去過我們這一生,.
怎樣去預備我們自己去創造我們的做物主,.
我們一起低頭祈禱..
主啊求你引導我們的心思意念,.
教準我們的眼光,我們的心思,.
是在你裡面去尋求得著智慧,.
尋求你的引領,我們去怎樣面對生命,.
特別在我們生命遇到困難,病痛,挑戰的時候,.
拿開人的心思意念,就算我們突然想的時候也好,.
轉頭幾個second,你將我們的心思意念轉化到你裡面去想,去思考,.
在你的話語裡面去尋求,怎樣繼續前進,.
繼續怎樣去過我們的生活,.
看我們的生命,幫助我們每個人,.
怎樣去預備好我們自己,有一天我們要見你的面,.
都是無愧的,都是嚮往你與我們的同在,.
這樣去祈禱,奉耶穌得勝命求,.
如果你想的時候,可以在你的座位上,.
花一點時間祈禱,然後我們散會,.
祝福大家.
\newpage



\section{18:21-35}
\label{sec:ON4fMvKp6vs}
\textbf{Forfeiting Forgiveness (Matthew 18\_21-35) - George Crawley}
\newline
\newline
連結: \href{https://youtube.com/watch?v=ON4fMvKp6vs}{\texttt{https://youtube.com/watch?v=ON4fMvKp6vs}} ~~~~ 語音日期: 2024-02-28
\newline
\newline
\hyperref[sec:Lczwsjcf0tU]{\small{< < < PREV SERMON < < <}}
~
\hyperref[sec:index_chronic]{\small{[返順時目]}}
~
\hyperref[sec:index_scriptual]{\small{[返順卷目]}}
~
\hyperref[sec:LB9q6sOwuKc]{\small{> > > NEXT SERMON > > >}}
\newline
\newline
$^{1}$All right, good morning, everybody..
Let's quickly pray with me..
Heavenly Father, I just pray now that you would use this time, use this word..
Let it not be a man's words, but let it be your word for your people and for your glory..
So I pray that you would speak through me, speak to each of our hearts, Father God, to.
make us more like the Lord Jesus Christ..
And I pray this in his name, amen..
So in December of 2003, a man named Gary Ridgway was sentenced by a Seattle judge to 48 life.
sentences..
In addition to this, he was sentenced 10 years for every one of those life sentences..
So he was facing 480-year sentence on top of 48 consecutive life sentences..
And it may seem extreme, but legally speaking, the sentence, this punishment was well-deserved.
because Gary Ridgway is more commonly known as the Green River Killer..
And he is the second most prolific serial killer in the history of the United States.
of America..
And while he was standing before the judge on trial for these 48 murders, he was stone-faced..
He had no emotion..
And he had actually confessed to killing over 70 women throughout Washington state, although.
he was only tried for 48 of them..
And through the entire sentencing process, he was stone-faced, he was emotionless..
And in the court proceedings, the judge had invited family members of these victims to.
come and speak before the court..
And so these families came up, and they bared their soul, they poured out their grief, their.
heartbreak, their rage towards this man who had taken everything from them—mothers,.
daughters, wives, sisters—all lost to this man who sat in the very same courtroom as.
them, absolutely remorseless..
And I got a few quotes here of some of the things that these families said, some mocking.
Gary Ridgway, some condemning this man for what he did to their families..
And so a few of the things that they said I have here..
"He's an animal..
I wish for him to have a long, suffering, cruel death.".
One woman said, "He's going to go to hell, and that's exactly where he belongs.".
Through all these statements and many more just like them, Ridgway remained emotionless..
But that is until one man took the podium, a father whose daughter had her life snatched.
away by this man..
And the father took the podium, and to everyone's surprise, this is what he said..
"Mr. Ridgway, there are people here that hate you..
I am not one of them..
You've made it difficult to live up to what I believe, and that is what God says to do,.
which is to forgive.".

$^{41}$And at that moment, the entire courtroom was stunned in silence as he said this next sentence..
"You are forgiven, sir.".
And at that, that shocking display of mercy and forgiveness, for the first time in months.
of court proceedings, Gary Ridgway broke down in tears, and he cried..
And there's a clip of this on YouTube..
If you look it up, you can search "The Power of Forgiveness, Gary Ridgway.".
And every time I watch that clip, I find it difficult not to myself tear up because it's.
just so powerful..
It's so full of emotion on display in the courtroom..
The heartbreak, the pain, the anger of those families who've just lost a piece of themselves,.
their desire for vengeance..
And then one man among many stands up and offers, extends forgiveness and mercy to a.
guilty murderer, a serial killer..
And when I see that, I feel a tug at my own heart, a tug of conviction, a conviction that.
asks me, "What would I do in that same position if it was my mother or my sister, my wife.
or my child?.
Would I be able to forgive like that?.
Could I forgive like that?.
Should I even forgive something so egregious, something so terrible?.
Does a murderer like that, does he deserve that forgiveness?.
How could that father have forgiven him?.
How could he do that?.
Why would he do that?".
And so today we'll be reading a portion of teachings from the Lord Jesus Christ that.
answers those very questions and it answers our questions when we find ourselves struggling.
to forgive, not wanting to forgive..
And so my passage I'll be reading from today is Matthew chapter 18, verses 21 through 35..
And before I get into the passage, I kind of want to explain and set up a little bit.
what is going on in the context of this teaching Jesus gives..
And so prior to this, Jesus had just taught his disciples his way of forgiveness, how.
when a brother or sister sins against them, how are they to respond?.
And what he tells them is when a brother or sister sins against you, you go to them and.
you tell them what they've done in hopes that they will repent and you will forgive them.
and you'll be restored, reconciled together..
But if they refuse to repent, you go and take two or three witnesses with you and together.
you'll go to this brother or sister who is in sin against you with the hope that you'll.
be restored and reconciled to them, that they will repent and you will forgive them..
But if they still refuse to repent, then you take the matter before the entire church with.
the hopes that the weight of the whole body of believers will cause this person to repent.
and then you forgive them and be reconciled and restored in brotherhood..

$^{81}$But if even then they repent, Jesus says to treat them, deal with them as a pagan, as.
a tax collector, as an outsider..
And so this is what Jesus is now responding to in this passage..
And after teaching these steps to his disciples, how they are to forgive a brother who sins.
against them, Peter comes up, he comes to the Lord and he asks him this question, "Lord,.
how often will my brother sin against me and I forgive him?.
As many as seven times?".
Which is a pretty good question Peter asks here, because if we're honest with ourselves.
we would probably have a similar question..
"Well, how many times Lord should I forgive?.
Yes, I know that you want me to forgive.".
And we generally agree that extending forgiveness is a good thing..
And we certainly know how good it feels to be forgiven..
But I think we all kind of think to ourselves, "There must be a limit, right?.
You don't expect me to forgive unconditionally if they keep hurting me or sinning against.
me..
My mercy, my patience, my forgiveness, it can only go so far..
I don't have to keep forgiving, right?".
And according to the Jewish rabbinical tradition of that day, of that time period, you were.
obligated to forgive three times..
So if a brother sinned against you, you forgave him..
And after the third time of extending that forgiveness, you were to show them no more.
forgiveness..
There was no fourth chance..
And it's kind of similar even 2,000 years later to our culture, to our day and age,.
and the world's approach on forgiveness is pretty similar..
And so I have a few sayings here that maybe you've heard of..
Maybe you've said it before or had it said to you..
"I'll only give you three chances..
That's one..
That's two..
Don't let me get to three or else.".
Life is like baseball..
Three strikes and you're out..
Or maybe this one..
"Fool me once, shame on me..
Fool me twice, shame on you.".
And you see the idea behind that is you're not going to fool me..
You're not going to cross me or offend me a third time..
I won't give you the chance..

$^{121}$And that's exactly how it was for the Jews in Jesus' day and age..
And Peter knew this..
And so he asked Jesus, "Should he forgive up to seven times a brother who sins against.
him?".
And so you can imagine the reaction that may have come from those who were around him..
Hearing Peter say, "Should I forgive seven times?.
Seven times?.
Are you kidding me?".
This guy, Peter, he's so forgiving..
He's so merciful..
No one else is forgiving seven times..
Not even the rabbis are forgiving seven times..
This guy, Peter, is something special..
But surely Jesus doesn't expect me to forgive seven times, right?.
No one else is doing that..
And so if Peter's question was surprising to those who were around at the time, Jesus'.
response to Peter would have been even more surprising..
And I think it's, to an extent, even surprising to us when we think about the implication..
And so Jesus answers Peter and says to him, "I do not say to you seven times, but seventy-seven.
times.".
Seventy-seven times..
Remember, the Jewish rabbis, the holy men of the day, they taught three times is enough..
And Jesus says, "No, not seven, but seventy-seven.".
And actually, depending on how you take the Greek that's behind this, how you translate.
it, it could be translated as seventy times seven, which is even a greater number of forgiveness..
But whether it's seventy-seven or seventy times seven, that's not exactly the point,.
because Jesus is not giving a specific number to teach the people, okay, you keep track..
And once they've sinned against you seventy-eight times, that's it..
Once they've crossed this certain number of times sinning against you, offending you,.
then you cut them off..
You offer them no forgiveness and you hold that against them..
That's not what He's saying..
He's not literally saying, "Only to forgive seventy-seven times," and that's it..
That defeats the purpose of what He's saying..
What He's really teaching them is that for the followers of Christ, there is no count..
There is no number of times you forgive and then you're off the hook..
But He wants you to go out there and continue to forgive..
Don't keep track, don't keep score, you just forgive, and they sin against you, and you.
forgive..
They sin against you again, and you forgive again..

$^{161}$It's an innumerable forgiveness..
That's what Christ is teaching His followers to do..
Not to keep score, not to get even..
And it's likely that the answer Jesus gives is actually a reference to a passage of the.
Old Testament Scriptures from the book of Genesis chapter four, where we see a descendant.
of Cain by the name Lamech, who is singing this song, so to speak, to his two wives..
And it's a song of revenge that demonstrates how great his vengeance is against those who.
sin against him..
And so this is what he says, "Ada and Zillah, listen to me, wives of Lamech, hear my words..
I have killed a man for wounding me, a young man for injuring me..
If Cain is avenged seven times, then Lamech seventy-seven times..
And sometimes, if we're being honest, this is how we feel..
Someone hurts us, they're in for it..
They're going to get it..
Whatever you do to me, I'm going to get you back with everything you gave me and then.
some, because I want revenge..
I want to get even..
I want to make you feel how bad you made me feel..
I'll teach you a lesson..
But for the followers of Jesus Christ, this is, it shouldn't be this way..
As a matter of fact, it cannot be this way for those of us who follow Jesus, because.
Jesus flips this entirely on its head..
So it's not that when you are sinned against, offended, or hurt, that's an opportunity to.
display vengeance and to get revenge..
For us, when we're sinned against, it's an opportunity to extend mercy and to forgive..
And see, that is the standard of forgiveness, which Christ holds all those who call on His.
name..
And this level of forgiveness, for some, for many, it is shocking..
And for some people in our modern sensibilities, it's even offensive..
Think about the story I opened with, this serial killer, this murder, just talking about.
that..
Some of us just have a reaction where we're repulsed..
And if you put yourself in the place of those families, it's very difficult to say that.
you would stand up and forgive someone who took the life of a loved one from you..
It's offensive to suggest that you should forgive him, that you must forgive him..
And so for many people in our world today, this level of forgiveness is extremely offensive..
And I'm speaking from personal experience..
I recall a time around 2018, 2019, there was a case in the news of a black man, his name.
was Botham Jean..
And he was sitting in his apartment on the couch eating ice cream, and a white police.

$^{201}$officer opened the door and came in and shot him and killed him..
She claims she lived in the same building, that she thought she was entering her apartment.
that he was an intruder..
That was her defense..
At the time, I didn't care..
I didn't care what she said..
She, in my brain, was a murderer and needed justice, because what took place was terrible..
So, during the court case, when the judge gave the family of this slain man a chance.
to speak to his killer, what happened surprised me..
His only brother stood up to the judge and when allowed to speak, he looked at this killer.
and he told the police officer, "I forgive you.".
And he went so far as to ask the judge, "Can I come from behind the stand and can I give.
her a hug?".
And when I saw that, I was angry..
I could not understand how you could forgive and embrace your brother's killer..
And when I went to my dad about this, and we saw it on the news, he told me that what.
had taken place was actually beautiful, and that what the family did was the right thing..
And at the time, I could not come to grips with that, that my dad, a black man, did not.
have that same anger I had at what took place..
That forgiveness made no sense to me..
And the truth is that at that time, I was not a Christian..
I did not understand, I could not understand the forgiveness of Jesus Christ..
It was offensive to me..
It was infuriating to me..
It baffled me..
And this radical call to forgive still baffles many people today..
And so, we want to know why..
Why Jesus?.
Why such a high standard?.
Why is it so important that we forgive, that followers of Christ have to forgive like this?.
And this is a question that would have been on the minds of all those around that day.
who heard Jesus' answer given to Peter..
Not seven, but 77..
"Why, Jesus?.
Why is that so important?.
How could I possibly forgive that much?".
And so, it's Jesus follows us up with a parable, explaining, giving an answer to those questions.
that we sometimes have..
This is what he says, "Therefore, the kingdom of heaven may be compared to a king who wished.
to settle accounts with his servants..

$^{241}$When he began to settle, one was brought to him who owed him 10,000 talents..
And since he could not pay, his master ordered him to be sold with his wife and children.
and all that he had, and payment to be made..
So the servant fell on his knees imploring him, 'Have patience with me..
I will pay you everything.'.
And out of pity for him, the master of that servant released him and forgave the debt.".
And so, what just happened here in this parable is unthinkable, especially to the people of.
that day who Jesus was speaking to..
Utterly unthinkable..
This large amount of debt that the servant owed to his master was uncountable..
It was 10,000 talents..
In that day and age in the Roman world, the talent was the largest denomination of money..
You couldn't go any higher than that..
Talents was, you can almost think of it in this way, talent, that amount of money was.
how you paid kings, how you dealt with prices of entire kingdoms..
And the number, 10,000, was the largest number employed in the day..
You couldn't think of, you couldn't write a higher number than 10,000..
The Greek word for 10,000 is myrias, which I say that because that's where we get the.
word myriad, which is an incountable, innumerable amount..
And so that's how much this servant owed his king, owed his master..
And so to put it in kind of common terms, modern terms for us to understand, imagine.
a low-level worker who owed the CEO and founder of a Fortune 500 company, Google, Microsoft,.
Tesla..
He owed him a billion dollars..
We might use the word gajillion or bazillion, right?.
It's nonsense because it's so much money, there's no way in the world you could pay.
it back..
And that's the debt that this servant had to his king..
And because he owed him such an incredibly unpayable amount, the king's response was,.
we'll sell all you have and we'll sell you and your family into servitude..
And at least I can regain a small fraction of a fraction of a fraction of a fraction.
of what you owe me..
But in a move of total and utter desperation, the servant throws himself before the king.
and begs and pleads for mercy, claiming that, give me time, I'll pay everything back, which.
both he and the king knew was impossible..
There's no way in the world that man could ever have enough time and do enough labor.
to pay back a debt of 10,000 talents..
The king knew this..
Yet for some reason, he was moved to pity this man and the king released the servant..
And not only did he release the servant, but he took that enormous incalculable debt and.

$^{281}$erased it completely..
This servant, this lowly servant, forgiven by a king, the servant who did nothing to.
earn this mercy, this incomprehensible debt, it's unheard of for that kind of thing to.
happen..
And so for Jesus' audience, hearing that from Jesus, it would have been shocking..
And then what happens next in the parable would have been just as shocking..
So Jesus continues in the parable..
But when that same servant went out, he found one of his fellow servants who owed him a.
hundred denarii and seizing him, he began to choke him saying, "Pay what you owe.".
So his fellow servant fell down and pleaded with him, "Have patience with me and I will.
pay you.".
He refused and went and put him in prison until he should pay the debt..
Despite the complete mercy, the grace, the forgiveness that this servant had just received,.
he refused to give that same mercy to someone else who was in the same position that he.
was just in..
Now it's worth noting that this hundred denarii is equivalent to about a hundred days' worth.
of wages..
So the debt owed to the servant, it was not insignificant..
That's a significant amount of money..
But it is payable..
And in comparison to ten thousand talents that were owed to the king, this one hundred.
denarii debt was, is not even worth comparing..
Right?.
Ten thousand talents pales in comparison to a hundred denarii..
But despite this, despite the fact that it was repayable, the mercy who had just received.
immense mercy and forgiveness, the servant punished the man, refused to pay him..
And that reaction from the servant, it demonstrated that his heart was full of unforgiveness and.
full of wickedness..
And so the scene actually plays out almost exactly parallel to the first..
The man, the servant is brought before the one he owed the debt to..
He pleads for mercy..
But the difference is, he receives none..
The servant is merciless to him..
Even though the debt was way less, even though the debt could have been paid off, it didn't.
matter to the servant, who drags the man, throws him into prison..
Jesus continues, "When his fellow servants saw what had taken place, they were greatly.
distressed and they went and reported to their master all that had taken place.".
And so the servant's mistreatment and lack of mercy to his fellow servant, it was appalling.
and it was offensive..
So much so that the other servants knew that what happened was wrong..

$^{321}$And so their first response was to go straight to the king and tell the king what had just.
taken place..
And we're going to see the king's response shows that he was just as deeply disturbed,.
even angered, by what took place..
And so here's what happens next..
"Then his master summoned him and said to him, 'You wicked servant, I forgave you.
all that debt because you pleaded with me..
And should not you have had mercy on your fellow servant as I had mercy on you?'.
And in anger, his master delivered him to the jailers until he should pay all his debt..
So the master rebukes him..
He calls him wicked and delivers him over to the jailers where he will be tortured until.
his debt is paid off..
Although he was at first abundantly merciful and forgiving the servant, now we see he has.
nothing left for the unforgiving servant but condemnation, contempt, justice..
And the reason the king gives is that this servant did not demonstrate the same mercy.
he had just refused, received..
And in doing that, in refusing to show that mercy, it's almost as if the servant holds.
in contempt the mercy he had just received..
As if he is too good to show mercy..
Better than the king was because the king showed mercy..
And so Jesus warns his followers that just like this servant was brought before the king,.
rebuked, sent off to the jailers, Jesus warns his followers that it will be no different.
for them if they act as the unforgiving servant acted..
And Jesus gives this very stark, very serious warning and it's a warning that extends to.
all of us today..
He ends this passage by saying this, "So also my heavenly father will do to every one of.
you if you do not forgive your brother from your heart.".
Exactly what happens in the parable to the unforgiving servant, Jesus says that is exactly.
what will happen to you from God if you refuse to forgive your brothers and sisters from.
your heart..
And so now Peter has his answer and now we also have ours..
Why should we forgive?.
Why should we forgive those who hurt us, who offend us, who abuse or mistreat or misspeak.
about us?.
Is it not because we, like the servant, have been greatly forgiven by God?.
That is the entire point of Jesus' parable..
And so you see, when you look at the parable, what does it all mean?.
What does it represent?.
We are the servant..
We owe an unimaginable debt to God..

$^{361}$Scripture teaches that all have sinned and fallen short of the glory of God..
And that the wages of sin is death..
So there's nothing we could ever hope to do, given all of time and eternity, to pay off.
our debt of sin to God, a righteous, holy, infinitely holy and good God..
There's nothing we could do, no amount of good deeds, good works we could ever do to.
pay off that debt..
Scripture teaches there's no one righteous, no not one..
No one understands..
No one seeks for God..
All have turned aside..
Together they have become worthless..
No one does good, not even one..
Elsewhere we see the prophet Isaiah writes that all our righteous deeds are like filthy.
rags..
Nothing we could ever do could pay off our debt to God..
But God is rich in mercy, abounding in steadfast love..
He has forgiven us completely the debt of our sin at the high price, the high cost,.
not of denarii, not of talents, not of silver, not of gold, but of the precious blood of.
Jesus Christ, the Lamb of God who takes away the sin of the world..
So therefore, though our sins be as scarlet, they shall be white as snow..
That is the forgiveness we receive from God in Jesus Christ..
Our debt to God entirely wiped out..
How then could we who know this mercy, we who have received this forgiveness, we who.
have believed in this Son, and we who profess this truth, how then could we turn around.
and withhold forgiveness from someone else, especially another brother or sister in the.
Lord?.
Remember that the debt that was owed towards the servant from his fellow servant, it wasn't.
insignificant..
It was a sizable debt..
And likewise, there are times when we are sinned against, when we are offended, and.
it is not insignificant..
It does hurt us..
It does cause issues..
It's not a light thing, but God understands that..
God understands our pain, our hurt, our feelings of betrayal when we are hurt like this..
But you see, God in the person of Jesus Christ has experienced the same pain, the same hurt,.
the same betrayal..
And so He can empathize with us, and He can give us the grace and the strength to show.
mercy to those people who hurt us..
An important thing to note here too, and this is a very, can be a hard pill to swallow..

$^{401}$It could be a difficult truth, but it is true nonetheless..
No matter what anyone on this earth does to you, no matter how terrible their sin against.
you may be, it will never come close to how offensive and egregious your sin, my sin,.
our sin is to God..
Because God is infinitely holy and pure and just and righteous..
And so our sin against Him, way worse, incomparably..
To try to compare the sin that we have towards God with the sin that someone has against.
us, it's pointless..
Our sin against a holy God is both incalculable and incomparable..
Yet in Jesus Christ, we have complete forgiveness..
That is why we must forgive..
It's because we have been greatly forgiven..
It has been said that there is never a time you are more like God than when you forgive..
Well, then the negative of that must also be true..
There is never a time you are less like God than when you refuse to forgive, especially.
if you know that you have been forgiven by God..
How many times have we all said this while reciting the Lord's Prayer, "Forgive us.
our sins as we forgive those who sin against us.".
That should be more than a few words we recite from memory here and there..
It should be, for us, a way of life..
It's more than a prayer..
According to Jesus, it's a command..
You know, we all know how good it feels to be forgiven when we have wronged someone..
Even more than that, we know how great it is to be forgiven by God..
So why would we withhold that forgiveness from a brother or sister who sins against.
us?.
Why would we hold a grudge?.
When you ask yourself, when you're honest, think about it this way..
What has holding a grudge ever done for you?.
Has anyone ever held a grudge and it solved the problem?.
Has anyone ever held a grudge and it went back in time and it righted the wrong?.
It never does..
It never helps..
And although it might initially feel good, if we're honest, to be angry and to stir.
in that anger and think thoughts of revenge, of how I'm going to get back and how I'm.
going to hurt you because you hurted me, in the end, it only hurts us..
It only makes us feel worse..
I was here at Summit on Friday speaking with the middle school and high school students.
and I told them that holding a grudge against somebody, withholding that forgiveness, harboring.
that anger towards them, is like drinking poison and expecting the other person to die..

$^{441}$It's ridiculous..
It's totally ridiculous..
How much more ridiculous is it then for us, those who have been greatly forgiven by God,.
forgiven more than anyone else in the world, to refuse to forgive someone else, especially.
a brother or sister in the Lord?.
Not only is it ridiculous, but according to Jesus, it is wicked..
And what we see the King in this parable do is exactly what we know that God will do to.
those who refuse to forgive their brother or sister..
It's true, God is a God of great mercy, forgiveness, patience, love..
He is also a God of great justice..
Hebrews chapter 10 verse 13 tells us this, "It is a fearful thing to fall into the hands.
of the living God..
Yet those who have received the forgiveness found only in Jesus Christ have no need to.
fear..
We should never forget, as Jesus teaches in Matthew 6, 15, "If you do not forgive men.
their sins, neither will your Father forgive yours.".
Brothers and sisters, if you leave here today and you only remember one thing that I say.
this entire morning, please let it be this..
If we refuse to extend forgiveness to others, we should not expect God to extend forgiveness.
to us..
Jesus speaks very clearly and forcefully on this matter..
There's no ifs, ands, or buts about it..
For the followers of Jesus Christ, forgiveness is not optional..
It can't be..
Now that's not to say, and don't take this to mean that our forgiveness is earning us.
God's forgiveness and that if we forgive, we earn God's forgiveness..
That's not how it works..
That's a misunderstanding and a misordering of forgiveness..
But what it is is that if we truly have been forgiven by God, if we understand and believe.
all that Jesus Christ has done for us to earn our forgiveness, we would then, knowing that,.
be moved to forgive those who sin against us..
We would have a new heart transformed by the Holy Spirit, which understands how much we've.
been forgiven and is quick to forgive those who sin against us..
Yet if we refuse to forgive others, that may be a sign that we have not truly understood.
and believed the forgiveness of God in Jesus Christ..
And that is a very dangerous place to be..
So brothers and sisters, I plea with you, search your heart..
Be honest with yourself..
Be honest before God..
Is there someone in your life who you cannot bring yourself to forgive, that you are not.

$^{481}$willing to forgive for some wrong they've done against you, some way they've hurt you,.
whether recently or in the past, whether it's a family member or a friend, a co-worker,.
a neighbor, a stranger?.
Is there someone that you cannot forgive?.
Now is the time to forgive them..
And the chances are, for some of us, if not all of us in this room, there is some real.
serious pain and hurt that we've experienced, trauma we've been through..
Give that pain to God..
He understands it..
He can bear that pain for you..
And maybe you're like me and you think, "Well, I'm not the type to hold grudges..
I don't do grudges..
I don't hold grudges..
I forgive people.".
Okay, but I challenge you to ask yourself, are you the kind of person who's quick to.
get angry?.
Yeah, maybe you're quick to forgive, but are you quick to get angry?.
Are you quick to have a bad mood because someone offended you at work, at home?.
Are you quick to become upset with a spouse for something they said or did or didn't do.
that you didn't like?.
Are you quick to have an attitude against them and think how you'll get back at them?.
Maybe it could be a stranger in traffic that cut you off and your first thought is, "Okay,.
I'm going to get around you and cut you off..
I'm going to get in front of you now and see how you like it.".
Is that your disposition?.
Sometimes that's me..
I'll be honest..
That's a heart of unforgiveness..
I don't hold grudges, but I ought to be quick to forgive because you see, God does not do.
that with us..
He doesn't hold those sins against us and think about how to get back..
So we should be quick to forgive because God is quick to forgive..
When God forgives us, He forgives quickly..
He forgives completely..
He doesn't keep a record..
"Okay, I forgive you this time, but next time I'll remember it because next time maybe there.
won't be a next time..
We should not be like that.".
Because God forgives quickly and completely and without keeping a record of wrongdoing,.
so should we..

$^{521}$So remember, all that Jesus Christ has done to forgive you your sin, remember how great.
your debt of sin is in the eyes of God and praise God that He has forgiven you..
Praise the Lord Jesus Christ that He has shed His life's blood to purchase your forgiveness.
and in giving Him the glory for that forgiveness, for that changed life, that changed heart,.
forgive others..
Please pray with me..
Father, forgive us our sins as we forgive those who have sinned against us..
Amen..
\newpage



\section{創世記 2:1-3}
\label{sec:LB9q6sOwuKc}
\textbf{The Chosen - 講道系列 \#2 (創世記 2\_1-3) - 陳崇基牧師}
\newline
\newline
連結: \href{https://youtube.com/watch?v=LB9q6sOwuKc}{\texttt{https://youtube.com/watch?v=LB9q6sOwuKc}} ~~~~ 語音日期: 2024-02-28
\newline
\newline
\hyperref[sec:ON4fMvKp6vs]{\small{< < < PREV SERMON < < <}}
~
\hyperref[sec:index_chronic]{\small{[返順時目]}}
~
\hyperref[sec:index_scriptual]{\small{[返順卷目]}}
~
\hyperref[sec:_33GO1favQo]{\small{> > > NEXT SERMON > > >}}
\newline
\newline
創世記 2:1-3
\newline
\begin{longtable}{cl}
\hline
\hline
章節 & 經文 (和合本修訂版)\\
\hline
2:1 & \begin{tabularx}{0.7\textwidth}{X} 天和地，以及萬象都完成了。 \end{tabularx} \\ \\ \relax
2:2 & \begin{tabularx}{0.7\textwidth}{X} 到第七日，神已經完成了造物之工，就在第七日安息了，歇了他所做一切的工。 \end{tabularx} \\ \\ \relax
2:3 & \begin{tabularx}{0.7\textwidth}{X} 神賜福給第七日，將它分別為聖，因為在這日，神安息了，歇了他所做一切創造的工。 \end{tabularx} \\ \\
[1ex]
\hline
\hline
\end{longtable}
$^{1}$好, 請主人.
我們今天繼續這個Chosen的系列.
不知道大家有沒有八卦在看.
看看這個Chosen八成怎樣.
有沒有興趣?.
沒有興趣是吧?.
你喜歡看舊的那套嗎?.
舊的那套耶穌傳很悶的.
悶到瘋了.
這套好很多, OK?.
所以我鼓勵大家, 真的.
如果你想重溫一下耶穌的故事.
其實看一看他的場景.
看一看他的故事的鋪排.
都可以給你一種具體的.
去到知道耶穌當時是怎樣.
他的生平是怎樣.
特別是我們壽苦節和復節快到了.
這個時候還有幾個星期.
我們儘管去到重溫一下.
阿科林書裡面的一些事蹟.
這個Chosen已經拍攝到.
好像是第四個季節.
四個季節, 很厲害.
用耶穌做主意.
可以拍攝到四個季節.
已經不錯了.
希望大家可以去看看.
當然裡面的故事內容.
是作者或製片商自己加上去.
我們需要去設想一下.
或判斷一下到底有什麼在聖經裡面.
有什麼是聖經以外的.
或者有什麼是聖經以外的.
加上去的一些情節和插曲.
是否合理呢.
我看了一個季節.
我覺得是合理的.
不是亂作的.
我們來到Episode 2.

$^{41}$講到Episode 2我就停了.
講完這次的Chosen我就不再講了.
因為大家沒有什麼興趣.
Anyway, Episode 2.
上次我們講Episode 1.
就是莫達拉瑪利亞被鬼附.
然後耶穌來釋放她.
裡面用了一些經文.
我們有解釋過.
Episode 2就是莫達拉瑪利亞.
她不再被鬼附的時候.
她已經康復了.
她開始有自己正常的生活.
上次我們講過什麼事呢.
就是尼哥底姆.
尼哥底姆在約翰福音第三章.
提到猶太的拉比.
曾經嘗試過為莫達拉瑪利亞驅鬼.
但是她失敗.
拉比失敗.
她不明白為什麼她會失敗.
知道莫達拉瑪利亞身上的惡鬼被驅走之後.
她就想找出.
想去發現到底發生什麼事.
為什麼她不能驅鬼.
而有另外一個人幫她驅鬼呢.
你知道當時的猶太人.
特別是兩約之間.
他們開始有一種概念.
就是魔鬼的概念.
他們明白魔鬼是一回事.
所以他們在兩約之間的著作裡.
經常提到.
如果妮塞亞來到.
如果上帝的國度來到的話.
他就會將人從撒旦的權勢釋放出來.
所以驅鬼的故事.
或者驅鬼的方法.
在兩約之間特別多提及.
當然聖經裡面沒有.

$^{81}$但是聖經以外的一些寫作.
尤其是猶太人的寫作.
例如《聖經》.
等等這些著作.
經常提到如何驅鬼.
你可以想像當時的人.
他們受的壓迫.
受的騷擾.
他們會覺得黑暗的權勢來到.
去籠罩著他們的時候.
他們在等一個王者來到.
在舊約裡面經常提到.
將會有一個王者.
大衛的君王.
或者大衛之子.
所羅門.
新的所羅門.
新的大衛.
新的摩西.
就會來釋放他們來幫助他們.
所以他們就是在等一個人可以來驅鬼.
所以為什麼他們這麼緊張.
為什麼會有人可以驅鬼.
你打開科堪書開頭那幾章.
經常提到什麼事.
耶穌出來傳道的時候.
就是他去驅鬼.
他驅鬼的時候.
就代表上帝的國度已經臨到.
而那些犯人就很不驅他.
為什麼.
因為到底你是按照什麼權柄來驅鬼.
經常問他.
你是鬼王的權柄來驅鬼.
還是真的有上帝的權柄來驅鬼.
所以經常有這樣的爭拗.
經常有這樣的議論.
所以如果背景的時候.
你明白科堪書為什麼會這樣提.
這個Episode 2說到尼哥底姆.

$^{121}$很惆悵.
到底為什麼他可以.
有人可以幫莫德拉瑪利亞來驅鬼.
到底是誰.
這些是後面的Episode.
繼續去交代.
無論如何.
瑪利亞就回復正常的生活.
開始他拍攝的.
所說的故事.
就是他去到預備一個安息日的晚餐.
一個Shabbat dinner.
他就邀請了一些朋友.
特別是一些社會裡面.
沒有人接受的朋友.
去到他家裡.
包括耶穌當時已經有兩個門徒.
這個Episode 2純粹就是講到Shabbat dinner.
就是安息日的晚餐.
當然每一集都有耶穌出現.
這一集最後耶穌又出現了.
所以不需要失望.
一定有耶穌在.
我們美斯不踢球.
OK.
耶穌一定在.
無論如何.
這裡看到他去到主持.
第一次的Shabbat dinner.
我覺得裡面所刻畫的.
一個Shabbat dinner.
一個Jewish的Shabbat dinner.
猶太人的安息日晚餐.
我覺得是很真實的.
在猶太人的安息日晚餐.
很多時候他們會做一些儀式.
那些儀式是一些講到聖經的故事.
還有會去到一些祈禱.
其中一個的祈禱.
就是這個.

$^{161}$Blessed are you O Lord our God.
King of the universe.
Who creates the fruit of the vine.
這個其實就是當時猶太人.
守禦節晚餐的時候.
他們會有四隻杯.
四隻杯.
你記得我們聖餐.
我們都有守這個杯.
領受這個杯.
但當時猶太人的晚餐.
禦節的晚餐.
是有四個杯.
四個杯代表四個上帝對他們的認許.
你要喝完四杯酒.
你才完結這個整個的禦節的晚餐.
Passover的晚餐.
講這個禱告.
Blessed are you O Lord our God.
King of the universe.
Who creates the fruit of the vine.
這個禱告是第三個杯.
喝的時候作的禱告.
所以這些是有根有據的.
這個Chosen是有根有據.
他們就會講這個禱告.
然後大家就會說Amen.
然後晚餐會繼續.
然後就會喝完第四隻杯.
就是第四個最後的認許.
在中間那邊.
然後他們又會有一個祈禱.
然後要結束這個晚餐.
剛才我提過.
我們每個月一次都會守聖餐.
你記得守聖餐的時候.
福音書怎麼說.
耶穌說我不再喝這葡萄汁.
直到神國的日子.
記得嗎.

$^{201}$所以其實就是跟門徒喝完一隻.
喝完另外一隻杯.
第二隻.
再喝完第三隻.
就是這個祈禱.
然後他就說什麼.
耶穌說我不再喝這葡萄汁.
直到上帝的國度來到.
你覺得很奇怪.
如果跟他一起去吃過聖餐.
吃一吃就走人.
沒有結束.
沒有去第四杯.
原來那個晚餐.
就是期待著.
將來上帝的國度真正到來的時候.
所以耶穌的晚餐.
舉起杯子來.
他說這個杯子就是代表我的血.
其實他們在說這個杯子.
就是門徒跟他有一個晚餐.
就是朝央行.
其實就是在說我們每次.
教會的時候.
我們去到首聖餐.
其實就在第三個杯子.
等著耶穌再回來.
這個是索伯克的背景.
猶太人的背景.
索伯克安息日.
跟我們有什麼關係.
安息日是從哪裡出來的.
第一段的經文.
我想和你處理的.
就是創世的第二章.
他說天和地以及萬象都完成了.
到第七日.
上帝已經完成了造物之功.
就在第七日安息了.
揭了他所做的一切的功.

$^{241}$上帝賜福給第七日.
將他分別為聖.
因為在這日.
上帝安息了.
揭了他所做一切的創造之功.
這個很熟悉了.
因為第一章說什麼.
六日創造所為.
六日上帝創立這個世界.
這個想再提一次.
如果第一次看聖經的話.
聖經是對現代科學的細節.
沒有什麼興趣的.
所以你不要拿創世的第一章.
想著這裡解答我們一些科學的問題.
最近這兩個星期.
我和子女聊天的時候.
經常都回到創世的第一章.
經常回到那裡.
他們問到底為什麼沒有提那些東西.
我們科學宇宙的起源.
我說聖經沒有興趣要回答這個問題.
聖經沒有興趣回答.
How it came to be.
這個宇宙到底是怎樣來的.
聖經沒有興趣回答這個問題.
聖經想解答的一個更重要的問題.
就是我們為什麼在這裡.
為什麼我們存在.
為什麼會有宇宙.
為什麼會有我們生命的存在.
Why 不是How.
為什麼我們在這裡.
意義何在.
人生的意義何在等等.
我們和上帝的關係.
應該是怎樣.
這是更加重要的問題.
聖經想去解答的.
所以聖經說到六日的創立世界.

$^{281}$是在說上帝為什麼要創立這個世界.
創立這個世界要做些什麼.
當然最後第六日就是創造的高潮.
人就是按照神的形象.
我們就好像上帝一樣.
我們原本的形象就好像上帝的尾膳.
我們的本性有上帝的神性.
有上帝的聖潔.
到第七日.
這裡就說天和地以及萬象都做完了.
這個就是所謂的第七日.
不是在說一些現代科學的東西.
是在說什麼呢.
是在說屬靈的事.
在說神學的事.
在說上帝.
在介紹上帝是怎樣的.
所以用了六日來創造.
然後第七日就說.
上帝就揭了他一切的功.
揭了是什麼意思呢.
停止了.
做了六日的事.
然後第七日他就停止了.
我們現代中文或者現代英文.
我們就會說.
God rested.
上帝休息了.
或者上帝揭了.
安揭了.
安息了.
我們的印象就是.
好像我們每天上班.
做了九個小時.
十個小時.
回到家.
我們很累.
然後我們坐在Lazy boy那裡.
我們休息一下.
那些Work from home更慘.

$^{321}$你休息不了.
因為你整天在家.
就等於工作.
總之我們的想法就是.
做了很多事.
然後我們休息.
或者我們累了.
這個就是我們現代人有的印象.
但是聖經裡面.
在以前的背景.
說到上帝安息.
一個神.
一個神明.
去到安息的時候.
是什麼意思呢.
就是說到他.
在他的聖殿的寶座.
坐在那裡.
坐在那裡的是寶座.
是掌管萬有的龍椅.
不是Lazy boy.
不是沙發.
所以說到六日創造完.
工作完之後.
上帝進入他的聖殿.
在聖殿裡面最重要的地方.
你知道聖殿裡面最重要的地方是什麼嗎.
致勝所.
Holy of Holies.
進去的時候已經越來越聖潔了.
但是進去裡面正方形的那一塊.
就是最神聖的地方.
那裡就是上帝的約櫃.
上帝的約櫃就是他的座駕.
安息了.
就是他進到致勝所.
他的寶座.
他坐下.
坐下做什麼.
就是去統治這個宇宙.

$^{361}$這個宇宙就好像他的聖殿一樣.
他去統治一樣.
所以不是說他做完事很累.
然後休息一下.
其實就是代表著上帝.
現在已經統治萬有.
這樣的意思.
這種描述.
說到上帝的停止.
上帝的安息.
上帝進到致勝所.
就是說他去in charge.
他去運作整個宇宙.
他是in control.
這個就是第一次介紹.
所謂安息這個概念.
所以為什麼在舊約裡面.
說到要守安息日.
這回事.
到底是什麼意思.
守安息日.
當他們進入聖殿來敬拜.
或者進入會堂來敬拜的時候.
他們祈禱.
去到上帝reconnect的時候.
他們知道.
我們現在放下.
我們六日守所所做的事.
去到聖殿敬拜.
我們知道誰是in control.
是誰in charge.
原來不是我們in charge.
原來是上帝in charge我們世界.
管理著我們整個世界.
這個值得留意的.
就是整個六日的創立世界的技術.
你記得由第一章第一節.
起初上帝創造天地.
地是空氣混沌.
淵源下流.

$^{401}$然後有光等等.
二三四五六日.
經常有一句話出現.
就是神看著是好的.
God saw it was good.
到最後God saw it was very good.
有多少次?.
有七次.
一次又一次.
上帝說看著是好的.
光他認為是好的.
他看著是好的.
等等等等不同的東西.
上帝在萬有裡面去分別.
什麼是好.
什麼是壞.
什麼是好.
什麼是美善.
是由他去分別的.
聖經第一章第一節.
一路講.
上帝就認為.
這件事是美好.
好像藝術家.
他做了一些藝術作品.
在中途裡面.
他說這顏色多好.
我就塗上去.
這樣畫多好.
塗上去.
然後整幅畫塗上去之後.
他就說.
It is very good.
有一個這樣的動作經常出現.
但到這裡第七天.
他說天地萬物已經做齊了.
到第七天的時候.
他安息了.
上帝就賜福給第七天.
這個很奇怪.

$^{441}$因為由他創造開始.
一路做到有乾地.
有海洋.
有天空上的飛鳥.
海裡面的魚.
地上的走獸.
然後地上有人類.
是按照他的形象所做的.
他一路都沒有說.
上帝賜福.
或者分別為聖.
這個動作.
你知道整個舊約裡面.
如果看回後面.
特別整聖所的時候.
整聖殿的時候.
神很多時候將人,事或者物.
分別為聖.
人會分別為聖.
為什麼?.
因為那些利未人.
祭祀神職人員所為.
有些物件是會分別為聖的.
聖所裡面的器具.
金燈台.
聖所本身的地方.
那個地方.
那個空間.
是會將它分別為聖的.
但在聖經裡面.
第一次上帝將那件事.
分別為聖的是什麼呢?.
是時間.
是一段時間.
一段二十四小時的時間.
他賜福給第七天.
將這段時間分別為聖.
Sanctify.
這個非常有趣.
不是一種物件.

$^{481}$不是一個地方.
而是一段時間.
一天第七天.
他將它分別為聖.
是和其他東西不同的.
這段時間和其他時間不同.
所以猶太人的傳統裡面.
經常提到說安息日.
是上帝賜給人類.
賜給你和我.
一個最大的禮物.
The greatest gift.
就是賜給我們時間居然.
分別為聖的時間.
安息日這段時間.
六天你要工作.
六天不是休息.
六天你肯定要工作.
但第七天.
一段時間.
他就說這是分別為聖.
猶太教的傳統和神學.
一直告訴我們.
這個安息日的日子.
安息日的時間.
是上帝給人最大最大的禮物.
這裡我不想引用一些.
現代的調查.
一些臨床的調查.
我昨晚看了幾個.
找了一些Journal.
就說守安息日.
對心理和精神的健康.
是有關聯.
Correlation.
不是Corsation 是Correlation.
不是因果的關係.
而是一種相關性.
這裡不引用.
因為大家都知道.

$^{521}$六天做工.
然後第七天那段時間.
他分別出來.
這個我覺得非常有趣.
上帝將一段時間.
首先將一段時間.
來到分別為聖.
我們每個人都是時間的活物.
對不對.
我們都是活在時間裡面.
我們沒有選擇.
你今年比去年老了.
你今年比去年殘了.
肯定這個我們沒有選擇.
對不對.
總會是這樣.
我們常說日月如梭.
因為我們活在時間裡面.
我們沒有選擇.
上帝創造我們就是.
活在一個時間裡面.
所以當我們認識上帝的時候.
我們和時間之間的關係.
是應該有些改變的.
對不對.
當我們信了耶穌.
承認這個是真神.
這個是真上主.
你和時間的關係應該有改變.
當然我們很多東西.
其他都會有改變.
我們和我們的財務.
都會有改變.
我們認識了耶穌之後.
我們和我們的金錢之間的關係.
會有改變.
我們和我們食物之間的關係.
都會有改變.
我們和我們兒女之間的關係.
和我們身邊的人.

$^{561}$無論配偶或者朋友.
所有東西都會有發生改變.
基本上會發生改變.
但在這裡說的是.
我們認識主的時候.
我們和時間本身.
看中做人.
是有關係的.
信了耶穌之後.
我們看時間.
或者看時間的管理.
應該是有不同.
這個就是安息日的意思.
提到安息日的時候.
第二段我們處理經文.
就是《出埃及記》第二十章.
我們先讀.
當紀念安息日首為聖日.
六日要勞碌做你一切的工.
但第七日是向耶穌說.
你的上帝當守的安息日.
這一日你和你的兒女.
老婆,婢女.
生畜以及城裡面寄居的客人.
都不可作任何的供.
因為六日之內.
耶穌說做天,地,海和其中的萬物.
這第七日就安息了.
所以耶穌說賜福與安息日.
定為聖日.
分別為聖的日子.
這個應該大家清楚.
就是十屆裡面其中一屆.
第幾屆?.
第四屆.
十屆裡面的第四屆.
十屆就是以色列人出埃及之後.
被恩典拯救了.
然後他們終於去到一個山腳.
就是西大山.

$^{601}$上帝就在密雲和有火出現.
就向這些以色列人納藥.
然後摩西就拿著法板下山.
那些情節我們不說太多.
但那個法板就寫著.
上帝和神的子民之間的約書.
上帝就說我已因恩典拯救了你們.
你們是一班新的人.
新的關係.
你們是一個新的人.
你們應該怎樣去生活呢.
然後就告訴他們.
第一屆,第二屆,第三屆.
直到去第十屆.
這是第四屆.
每一段的屆命.
是有長有短的.
我告訴你第一屆是什麼.
除了我以外不可有別的神.
如果你看文字的話.
第一屆是用了七個字.
第二屆不可為自己雕刻偶像.
就長一點.
有43個字.
然後你繼續數下去.
譬如第十屆也算很長.
有多少字呢.
大概有15個字.
這裡有多少字.
第四屆是有55個字.
在十屆裡面很奇怪.
第四屆就是有最長的描述.
最長的描述居然不是說.
不要謀殺,不要犯奸淫.
不要貪戀別人的東西.
不要偷別人的東西.
不要製造偶像等等.
居然最重要或最長篇的.
介紹一個的屆命.
第四屆就是安息的屆命.

$^{641}$用了55個字.
最長.
我們要注意到為什麼.
安息是勞碌做一切的工作.
六天裡面.
然後第七天就要停止.
其實這個就是Sabbath的意思.
Sabbath或Sabbath.
希伯來文.
其實就是說停止的意思.
停止的意思.
就代表上帝指明他們在第七天.
這段時間不可以做事.
停止下來.
跟之前的六天是不同的.
有個六天一天.
6161這樣的圖案.
今天的以色列國.
比較保守的地方.
當然有很多地方比較.
比較傳統的地方.
有些人不理會.
照樣繼續做事.
但是比較保守的地方.
比如耶路撒冷.
你就很清楚看到.
星期五的Sundown.
太陽下山的時候.
人們就不做事.
回家.
或者去會堂.
去敬拜 去祈禱.
巴士也不走.
然後食肆也不開.
如果你去到那裡.
就餓死你了.
不過你如果知道.
認識一個好的導遊.
他就帶你去基督徒那區.
基督徒那區是照樣辦事的.

$^{681}$或者阿拉伯人那區.
就照樣開的.
你去那裡就有東西吃.
你不認識的話.
你去猶太人區.
鬼影都沒有.
沒有人理你.
這是一個很有趣的現象.
因為Shabbat就代表Stop.
去停止.
因為是一段時間.
是分別為勝的.
亦都是上帝和神的子民之間.
一個很重要的記號.
是一個特別的記號.
在舊約裡面.
當然有很多不同的記號.
割禮等等.
聖殿等等.
但是Shabbat有一個特別的標誌.
一個外在的記號.
來記錄.
來標記.
這一班人和其他人.
全世界的人是不同的.
所以有一個這麼奇怪的現象.
Shabbat停止的時候.
開始的時候.
車也不走.
食肆也會收.
但是一到Shabbat完結了.
星期五晚開始.
按照阿拉比的文學.
就是你望天.
有三顆星出現的時候.
Shabbat就開始了.
然後他們回家.
他們會點蠟燭.
開始他們的Shabbat.
但是到星期六的聖誕.

$^{721}$完結了Shabbat.
24小時之後.
車就全部走出來.
夜市又會再開.
真的奇怪.
但無論如何都是一個停止.
在舊約裡面非常看重.
所謂安息日的時間.
每逢有人來干犯安息日.
就是一件很大的事.
當然我們要搞清楚.
施子舒經常說.
神的子民.
你們守安息日禮儀是沒有用的.
他們不是說完全廢去禮儀.
他們只不過說他們的生命還沒搞定.
然後他們走進聖殿裡面.
去守安息日.
就以為這樣就OK了.
這種虔誠的假裝是不行的.
這是施子舒經常斥責的事.
但主要他說的就是.
這個Shabbat的時間.
Shabbat的停止.
是說他們和世界各地的人是不同的.
就等於我們今天.
很多時候有些關係.
我們都有些外在的記號.
譬如婚姻.
我們有戒指.
戒指你會戴上.
不見了就很嚴重.
不過我不見了就OK.
為什麼呢.
因為我是那些蛛絲繩.
我結婚的時候沒有錢.
隨便買一隻.
然後我又肥又瘦.
有時候又跌了.
我就買了幾隻.

$^{761}$幾隻這些蛛絲繩.
但我都戴上.
或者兩個人的婚約的證書.
如果你燒了它.
很嚴重的事.
這些外在的標記.
就是代表著.
他和上帝之間的約.
這個就是Shabbat.
正如我所說.
我們來到認識主的時候.
我們的時間.
和我們之間的關係.
沒有改變.
因為真正的神.
是耶穌基督.
在舊約裡面就是第七天.
繼續發展下去的時候.
你就發現慢慢地就變了.
變成了第一天.
為什麼呢.
我們看兩節經文.
約翰福第20章第一節.
七天的第一天清早.
天還黑的時候.
末大拉瑪尼又出現了.
來到墳墓看見.
石頭封住墳墓的石頭.
拿開了.
這裡.
First day of the week.
是耶穌復活的.
那個清晨.
First day of the week.
然後我們再看.
譬如.
啟示錄第一章第十節.
這裡說到約翰他被流亡.
到一個荒島拔摩海島.
因為福音被迫害.

$^{801}$他在主人.
被聖靈感動.
然後他聽見大的聲音.
如同皓角跟他說.
你現在要寫信給七個教會.
以弗所至.
到老弟家.
他說有一主日.
The lost day.
應該就是.
First day of the week.
所以舊約就是第七天.
去到新約的時候.
因為耶穌基督所做的.
慢慢地將這個七天.
初期的教會就變成了第一天.
為什麼?因為沒有一天比一天更加重要.
就是耶穌的空墳墓出現.
耶穌從死裡面復活.
七天的第一天清早.
主日.
The lost day.
因為基督才是我們真正的安息.
基督真正的作王.
死十家上面.
三天之後他復活.
他坐在神的右邊.
他去掌管一切.
掌管你的生命,掌管我的生命.
所以十八天由第七天.
就變成了第一天.
就變成了所謂新地出現.
所以我們就不再守第七天了.
所謂.
這個就是福林安息日會.
或者Adventist他們會做的.
他們就會說一定是第七天.
其實也不是說太拘泥哪一天.
但肯定是有一個六天.
一天,六天,一天.

$^{841}$這樣的一個模式.
認識了主.
我們對時間.
是會有改變的.
我想結束的時候可以提一提.
一個很簡單的提醒.
我們時間上面的紀律.
時間上的紀律.
正如我所說.
你認識了耶穌之後.
你生命裡面任何的東西.
都會再整理過位的.
那個關係.
或者那個架構會不同.
但這裡說的.
一個更重要的.
就是你和時間的關係.
會有改變.
你怎樣去看你的時間.
看每一個星期都會有分別的.
有幾條問題我想問你.
我們應該怎樣去紀念.
或者去尊重.
這個所謂的主日.
我們是不是經常.
將主日.
當作我們六天工作的延伸.
六天我們又不勞碌工作.
然後我們說第七天.
我們就拿來做延伸.
不要緊,昨天做的東西做不完.
今天我們補回吧.
反正我有一天.
上帝給我的.
不是.
是紀念安息日的意思.
我在行動上應該怎樣去.
紀念這個主日.
記起這個主日.
怎樣去預備這個主日的來到.

$^{881}$以致我真的能夠在主日裡面去敬拜上帝.
為什麼要敬拜他.
因為他才是那位.
掌管我一切.
的上帝.
在這裡有一些很實際的形容.
我希望能夠提醒大家.
第一,講到安息日的時間.
七天裡面的一天.
基督徒的教會.
由第一世紀到現在.
就轉成一個星期的第一天.
主日.
怎樣去紀念,怎樣去記憶.
你想起紀念.
想起記起.
一些重要的東西.
譬如某某人.
你很關心的人,他生日.
你就會怎樣呢.
你會記憶他的生日.
你會做什麼呢.
你會預備他來到.
快要誰誰生日了.
先幫他擺十圍酒.
你就要訂好酒席.
或者買好蛋糕.
或者想起這一天的來到.
起碼心靈上你都想起.
心靈上.
你生日了.
這個就是紀念的意思.
記憶.
主日.
我們每一個禮拜的計劃裡面.
其實都應該要將這件事放在裡面.
這一天快要來到了.
我要怎樣預備好這一天.
預備好這一天的來到.
猶太人很多時候花了幾天時間.

$^{921}$來到去預備他們的安息日.
無論是星期六早餐也好.
或者家裡的.
要去主持.
然後有些款待.
他們會預備這一天的來到.
渴望這一天的來到.
就不是說禮拜天早上起床.
哎呀.
今天是主日.
死貨.
我們飛車去教會.
紀念.
真的在計劃裡面.
每個星期的計劃裡面.
我們要看前面.
還有去安排預備這一天.
第二樣東西.
是安息.
神的安息就是他統治一切.
但我們的安息.
就是他統治一切.
所以我們可以有安穩.
現代人是很忙的.
有千萬樣的事.
在你的心思意念裡面.
你今天走進來.
直到你走出去.
你有很多很多的事情要煩.
小孩子的事情.
大人的事情.
要照顧不同的事情等等.
是要安息的.
在24小時裡面.
我們需要take some moments.
一些時刻.
來停下來.
好像你的指南針.
有沒有玩過指南針.
以前我們小時候.

$^{961}$現在出了電話很少.
有支針發來發去.
有時候指南針會亂.
亂指.
然後你要calibrate它.
你要去校正它.
以致你上路的時候.
你知道方向.
其實做人都是一樣.
上帝說將這一天分別為勝.
其實對我們是一個最大的禮物.
就是我們需要reset.
重新再來過.
校正接下來這個禮拜.
我的優先次序.
是不是對呢.
還是我是被這個忙碌.
來騎劫我的生命.
校正.
以致我們能夠在上帝面前有安穩.
第三.
安息是做什麼的.
端午就是停下來.
我們就在這裡睡覺.
在舊約甚至新約.
安息是用來敬拜的.
敬拜就是上帝的子民.
走在一起.
告訴這個世界.
我們宣告耶穌才是主.
這個是最重要的一個服事.
就是我們整體走在一起.
去告訴人.
耶穌才是主.
上帝才是作王者.
不是我們.
不是任何人.
不是政權.
不是身邊任何勢力.
是上帝是耶穌基督作主.

$^{1001}$所以我們來到安息日的時候.
我們就去敬拜.
我們會spend time with family.
和家庭和家人.
有些親近的時間.
亦都是一種reset.
第四就是默想.
今天不妨去想一想.
停下來一點.
慢下來.
去思想上帝在你生命裡面的恩典.
剛才我們提到十誡.
十誡我們經常說第一誡是什麼.
除了我以外不可有別的神.
但其實不是第一句說話.
十誡你回去看看.
第一句說話.
是我師爺說.
曾經從埃及遺老之地.
拯救你出來.
除了我以外不可有別的神.
第一句說話是恩典的說話.
是恩典的救贖.
上帝無條件去拯救你出來.
有恩典給你.
以至你去到守戒名的時候.
就是跟上帝說.
多謝你來愛我.
用恩典拯救我.
以至我整個的生命.
所有的關係完全都改變了.
七天裡面的一天.
主日.
我就來到敬拜你.
我做好這個功課.
以至教會能夠在這個世界裡面.
繼續去見證神.
我跟心大留祈禱.
主日在我們忙碌的當中.
我們停下來.

$^{1041}$今天是主日.
是你的聖日.
是你自己用你的恩手.
去到分別為勝的時間.
主日求你幫助我們.
讓我們不要當這一天.
好像以往的六天那麼普通.
因為是你所審證的.
是你所祝福的一段時間.
主日叫我們敬拜.
我們來到你面前.
知道你是那位真正作王.
當我們在你面前安息.
我們知道我們有安穩.
有平安.
無論生命裡面有多少混亂.
有多少忙碌也好.
有多少挑戰也好.
我們知道你比任何一切的事情.
更大更重要.
主日我們今天.
將這一天嘗試去分別為勝.
主教正我們的焦點.
我們生命裡面的優先次序.
要我們尊你自己先為大.
接下來這個星期.
求你教正我們優先的次序.
在工作裡面祝福我們.
在我們的家庭.
在我們的關係裡面.
更多可能錯綜複雜的關係裡面.
教我們怎樣去活得像你.
神我們求你來祝福.
讓這一天成為我們每一個星期新的開始.
讓我們敬拜夢裡的喜悅.
祈禱奉祖耶穌賜給我們.
願神祝福各位.
\newpage



\section{6:5-8}
\label{sec:_33GO1favQo}
\textbf{Remember: You are speaking with your father (Matt 6\_5-8) - John Yi}
\newline
\newline
連結: \href{https://youtube.com/watch?v=-33GO1favQo}{\texttt{https://youtube.com/watch?v=-33GO1favQo}} ~~~~ 語音日期: 2024-02-29
\newline
\newline
\hyperref[sec:LB9q6sOwuKc]{\small{< < < PREV SERMON < < <}}
~
\hyperref[sec:index_chronic]{\small{[返順時目]}}
~
\hyperref[sec:index_scriptual]{\small{[返順卷目]}}
~
\hyperref[sec:_2c5tRWbVyU]{\small{> > > NEXT SERMON > > >}}
\newline
\newline
$^{1}$(audience laughing).
There we go..
Sorry about that..
For all you folks online,.
sorry you're gonna have to ask somebody else what I said..
(audience laughing).
So if you will turn in your Bibles.
or you can take a look at the screen.
and we'll have the scripture text available for you..
I'm really glad to be back.
and it's good to see familiar faces.
and some new faces too..
If I remember right,.
I think it might have been the very last sermon.
that one of the kind of things that I had said.
that I was thinking like,.
no, this might be controversial,.
but I said, okay, enough's enough..
Let's all come back to church, right?.
And so I know that probably may have upset some people,.
but I still stick by those words..
I'd say, all right, enough's enough..
Let's come back to church..
Let's enjoy the fellowship that we have.
with our brothers and sisters..
And with that, maybe there's all these accounts.
is clicking off right now..
They're like, okay, I don't wanna hear this guy anymore.
after he said that,.
but I hope you will stick around..
And I believe God's got something for us.
that is just so timely and it's timeless too,.
because it's something that Jesus said 2,000 years ago.
and it still meets us right where we are today, right now..
So let's just get right into it..
So over the last couple of years, it's nothing new..
I'm not saying anything revolutionary here.
that we've all gone through stuff together.
that is probably more difficult,.
more challenging than anything,.

$^{41}$any season that we have experienced in our lifetime..
So while only one or two things.
of these things that we've experienced.
over the last couple of years,.
one or two things would just shake the foundation.
of our lives, of our daily lives,.
we've experienced collectively a whole bunch.
of what I'll call fear factors..
And the challenges of a global pandemic,.
the widespread racial tensions,.
and there's systematic supply chain issues.
we keep hearing about,.
and wave after wave of geopolitical intrigue..
And then it's compounded by, you all know it,.
go to the gas station, go to the grocery store,.
economic instability, maybe it's in the workplace..
And then that doesn't even include.
like what are we experiencing when we're all by ourselves?.
Maybe it's a sadness or loneliness or depression,.
feelings of inadequacy, feeling like we're failures,.
all these things that are just weighing on us..
And while worry and fear are going to try.
and seize this opportunity to lead and control our lives,.
Jesus invites us to authentic faith in God,.
to not be ruled by fear and worry,.
but that our lives will be ruled by faith in God..
And one of the principle ways that we express our faith.
on a day-to-day level is in the way that we pray..
And so today's passage,.
Jesus has something to say to us about prayer..
So why don't we begin with prayer?.
Father God, thank you so much.
that you invite us to come and speak with you,.
to meet with you, to know you,.
have relationship fellowship with you..
And sometimes God, the reality is we don't know what to do.
once we are with you..
And so would you today speak to us and teach us.
how do we relate to you?.
Who are you?.

$^{81}$And teach us Lord how to pray..
Just as the disciples asked Jesus, your son,.
the very first time we ask joining the generation.
after generation of disciples who have said,.
teach us how to pray..
And it's in Jesus name we ask it, amen..
So God, our father in Jesus Christ.
is going to reveal himself today, I believe that..
So let's be ready to receive what he has for us..
So we go to Matthew chapter six, verses five through eight..
We're in the Sermon on the Mount,.
you know, Jesus' famous sermon..
And this is what he says in chapter six of Matthew,.
verses five through eight..
"And when you pray, do not be like the hypocrites.".
I should give you the passage, there it is..
"Do not be like the hypocrites,.
for they love to pray standing in the synagogues.
and on the street corners to be seen by men..
I tell you the truth,.
they have received their reward in full..
But when you pray, go into your room,.
close the door and pray to your father who is unseen..
Then your father who sees what is done in secret.
will reward you..
And when you pray, do not keep on babbling like pagans,.
for they think they'll be heard.
because of their many words..
Do not be like them,.
for your father knows what you need before you ask him.".
So while this passage from Jesus' Sermon on the Mount.
might look like, okay, it's his instructions on prayer,.
and it is, it is his instructions on prayer,.
but it's more..
The whole Sermon on the Mount I submit to you.
is really Jesus' Sermon on what does it mean.
to have a real relationship with God..
Because how we pray reveals the one to whom we pray,.
how we pray reveals our relationship to the one.
to whom we pray,.

$^{121}$and third, how we pray reveals what we believe.
about the one to whom we pray..
So here, Jesus first approaches the subject.
from the negative,.
and then he's gonna go back and forth.
between positive, negative, positive..
So he starts and he says in verse five,.
"And when you pray," here's the negative,.
"Do not be like the hypocrites,.
for they love to pray standing in the synagogues.
and on the street corners to be seen by men.".
Simply don't pray like hypocrites..
Don't pray like those who have all the wrong reasons,.
all the wrong motives for praying..
Don't do it for show..
Don't do it to keep up appearances,.
not even for the appearance of spirituality..
Don't pray so that people think you're godly..
Don't pray just so that people can say,.
"Wow, look at that relationship they have with God..
Look how close they are with God.".
Because that's how the hypocrites pray..
They are people who have little.
and maybe no relationship to God..
Yet they wanna suggest,.
and you could say they're all about this hustle..
I want everybody to think I'm something I'm not..
Remember, we are speaking to our Father in heaven..
And what does Jesus say next?.
And this ought to be a shock, a surprise.
for those who have settled for that kind of communication.
and relationship with God..
Jesus says, "I tell you the truth,.
they have received their reward in full.".
Why is that a shock?.
Why is it like a punch in the stomach.
for those who would pretend,.
act like they are spiritual,.
who are acting like they're talking with God?.
Because let's unpack something here..

$^{161}$You know, Jesus often points to an absurdity.
so that he can show us what's so obviously true..
And what's the absurdity here?.
I mean, think about it..
What is the purpose of prayer?.
If we go back to Sunday school,.
and I don't know if you guys learned this,.
I learned this in VBS in Sunday school..
Hey, you know, if you don't know how to pray,.
consider using this acronym, the Acts of Prayer..
A stands for adoration..
C stands for confession of your sins..
T for thanksgiving and gratitude to God.
for what he has done for us..
And then S, supplication,.
which means asking God what we need, our supply..
So think about it..
If you're a hypocrite and you pray great words,.
beautiful words of worship and adoration to God,.
sing wonderful songs to him,.
but you're doing it for show..
Or you confess your sins,.
but you're not doing it.
because you want God to know that you're repenting,.
you're admitting to your flaws,.
but you just want people to think like,.
wow, look how godly he is, how humble he is..
He's confessing his sins..
And if you're thanking God for all that he's done for you,.
but you're not really talking to God,.
you're talking to the people..
Or if you are presenting this long list of all your needs,.
and we all have long lists of needs,.
but we don't bring it to the one who can actually supply us..
How ridiculous is that?.
When the opportunity is right there,.
the invitation is there, come, come, let's hang out..
Let's fellowship together..
Let's have a personal one-on-one appointment together..
And then instead of taking advantage.

$^{201}$of such an amazing opportunity, an invaluable,.
such a priceless opportunity to be spending time.
with the God of the universe, the creator,.
our heavenly father,.
instead we're spending time with the people.
so we can show off..
We have failed to do the very thing,.
the most important thing, the very purpose of prayer..
We have not communicated with God..
We have not spent time with God..
Instead, we receive our reward in full..
And what is that?.
That somebody might pat you on the back and say like,.
boy, aren't you a good Christian?.
Aren't you spiritual?.
And what good is it to receive that reward?.
What is that compared to the reward.
of spending real time with God?.
Then Jesus declares,.
truly they have received their reward in full..
People have seen them..
People have congratulated them, approved them,.
thought highly of them, respected them..
If we're gonna do that, why would we do that?.
But if we're gonna do that, why do we go to church.
if it's not just to be able to put on our calendar,.
like check, we are at church..
Look at our attendance record..
Or why would we carry a Bible just so people would say like,.
wow, look at that, that guy still carries a Bible..
All I carry is my phone..
Or why would we teach or preach or serve or lead.
just so that people can say great job?.
When the opportunity is there to receive a greater reward.
from our heavenly father..
So then he pivots to the positive in verse six,.
and he says, instead pray like this..
But when you pray in verse six,.
go into your room, close the door..
I'm gonna go back here so you can have the passage again..

$^{241}$Sorry, I think the app froze..
Oh, there it is, okay..
But when you pray, go into your room, close the door.
and pray to your father who is unseen,.
then your father who sees what is done in secret.
will reward you..
You know, I had a brother that if you would ask him.
to pray during worship service or in a Bible study.
or in a meeting or before a meal, he'd say, no, no, no, John,.
I don't do that because citing this very passage,.
Jesus says we should pray in secret..
Well, friends, I hope we can all agree that Jesus.
in saying what he is saying here.
is not prohibiting public prayer or vocalized prayer..
He's not saying the only way to pray is in a closet.
or in secret, but that he is emphasizing.
the opportunity we have that we can have.
an intimate relationship with God, a one-to-one,.
almost like a secret relationship with God..
And you know what we do, right?.
When we know that we have, that somebody is gonna be,.
you know, they're gonna value a secret that we share,.
they're gonna guard the secret,.
we know we can say anything..
When we're with somebody we trust, somebody we love,.
and we're close right next to them,.
and we can whisper into each other's ear,.
we know we can say anything and it's safe..
We don't have to hide, we don't have to pretend,.
we don't have to put up some kind of air of,.
oh, you know, I'm actually better than I really am..
No, we can be completely honest, revealing everything..
And Jesus says, you can be that way with God..
You can be that way with God, going into a secret place..
You see, God is not a disinterested God..
He is not aloof and far away..
He is our Father who invites us to come near to Him,.
who invites us to seek Him out,.
and then promises us that we're gonna find Him,.
that He's not like trying to be evasive..

$^{281}$You know, you ever get ghosted by somebody?.
I get, that happens to me..
I know, you'd think like,.
yeah, how could that happen to this guy?.
He seems like such a likable guy..
No, I get ghosted all the time, and I'm like,.
man, I know how text messages work..
He just got that message..
Why isn't he answering, right?.
So you just hope they're probably in a meeting,.
they can't answer right now..
God's not gonna ghost us..
He doesn't get tired of His children and say,.
you know, oh, gosh, this guy, he won't stop talking..
He always wants to hang out with me..
God's not that kind of a Father..
He loves it when we spend time with Him..
He wants us to come closer and closer to know Him better..
'Cause remember, we're speaking with our Heavenly Father..
I hope we all believe this,.
that this is the God that invites us to come near to pray..
He's intimately close, He loves us,.
and in fact, the Bible says He loves us so much.
that His own Holy Spirit has made our bodies His temple,.
His home..
But He says that we are to abide, have unity with His Son..
And then He tells us, we're so close,.
that nothing can separate us from His love.
in Christ Jesus, our Lord..
And that's close..
And this is the Father to whom we pray..
And as in a secret place, we can enjoy unparalleled,.
uninhibited access to the Father.
so that this reward of knowing Him.
is so much greater than that meager reward you get.
when you pray for the sake of what people might think..
Sure, there's stuff that comes with this reward..
You know, we bring our lists and we say,.
God, please help me, please provide for me..
I have this need..

$^{321}$And sure, God takes care of material needs too..
But that's like having a steak and saying like,.
man, this is the most awesome salt I've ever had..
(congregation laughing).
There's something even better.
than the stuff that God provides..
It's Himself..
So let's not let the stuff that is good.
distract us from what's really the best..
It's God Himself, the Father..
So in verse seven, then He said,.
well, let's say, you know,.
the authentic relationship with the Father.
requires us to believe who He is..
And actually it requires us to believe who we are too..
That when we pray, we're not praying as strangers.
or the enemies of God, but in Christ Jesus,.
we are His children..
We are talking to Daddy God, Abba Father..
So we had to know who He is..
We got to know who we are..
And then when we pray,.
it also includes knowing what kind of Father is He..
So in verse seven, Jesus says,.
"And when you pray," going back to the negative,.
"do not keep on babbling like pagans,.
for they think they'll be heard.
because of their many words.".
So what kind of Father is He?.
The first thing we see is He's not the kind of Father.
who's really impressed by lots of words..
Words by themselves don't do it for Him..
You know, in many religions,.
including many traditions of Christianity,.
you know, writing down prayers, reading prayers,.
memorizing and reciting prayers,.
and even praying a long time..
These are valuable things, highly valued things..
And don't, just like we're not gonna do.
Bible studies anymore,.

$^{361}$doesn't mean there's no more prayer meetings either, okay?.
By the way, we are gonna have Bible studies,.
right, Pastor Mark?.
And yes, we should pray..
So these are all still valuable..
Yes, write down your prayers, journal your prayers,.
because you know, you're gonna come back one day.
when you're really discouraged,.
and you're gonna go back and read your journal,.
go, "Man, I totally forgot about this..
Six months ago, I was so depressed,.
and I asked God, 'Lord, would you please give me joy.
and peace?.
Would you give me contentment?'.
And I have it, and I totally forgot I asked Him.".
And then you can write next to your little prayer.
that you wrote down, "God did it..
He came through..
How faithful He is.".
And then we'll worship Him..
So write them down, recite prayers, memorize prayers,.
'cause sometimes, you know this experience,.
you get to your quiet time, your prayer closet,.
and you're just like,.
"Man, I don't even know what to say to God today.".
So what do you do?.
You might go back and read the prayers of others..
That's valuable..
And then sometimes, you wanna discipline yourself and say,.
"You know what?.
I don't have much to say, but I am not leaving this room.
until I have prayed 10 minutes, or 15 minutes,.
until I've talked to God for half an hour,.
or prayed for an hour, or more.".
Those are great disciplines..
But even those valuable things can become empty and void.
of value and meaning if we're not careful..
Let me give you an example..
You know, when I was a small boy,.
you know, I probably learned from my parents.

$^{401}$more than anybody else how to pray..
And one day, when my dad would hear me saying the same things.
over and over again in prayer, like at the dinner table,.
he said, "Hey, why don't you kinda mix it up a little bit.
and be more creative?.
'Cause it doesn't sound like you mean it when you say it.".
So he said, "For example, maybe you might wanna start.
with not just, you know, 'Thanks for the food, God,'.
but you might say, 'Thank you, God, for a beautiful day..
Thank you for taking care of me.'".
And, you know, so he would provide these suggestions..
And I remember thinking, "Yeah, that's a great idea.".
So I added in my prayers,.
"God, thank you for this beautiful day.".
As a child, that just like, "Hey, that's simple..
That makes sense.".
And, but I remember, I was in elementary school,.
and I was praying over my cereal and milk..
And I remember praying,.
"God, thank you for this beautiful day.".
And immediately, I looked out the window.
and noticed it was dark, it was raining, it looked cold..
And here I was thanking God for a beautiful day..
And immediately I realized, like,.
"Man, I just said something I don't even mean.".
Unless, you know, even as a kid, I was super spiritual,.
like, "I'm praying for a beautiful day..
I'm believing God's gonna give it to me.".
Or, "I'm so spiritual, I'm gonna be content.
as I'm getting soaking wet and freezing in the rain.".
No..
Even something that was meant to be good.
had become meaningless 'cause it was just words..
And you know what I'm talking about.
'cause I'm sure you've experienced that too..
We gotta check ourselves..
And sometimes it's even the very name of God,.
the calling God Lord..
You know, have you ever had,.
sorry, this isn't in the notes,.

$^{441}$have you ever had one of those conversations with somebody?.
Wouldn't it be annoying if I was talking to Mark.
and I said, "Hey, Mark, how's your day, Mark?.
Mark, you wanna go out for lunch, Mark?.
Is Wendy doing okay, Mark?.
Are your kids doing okay, Mark?".
And if I kept saying his name over and over again,.
and you know what?.
If you think about it, that's how we often pray, right?.
Like, "Lord, thank you so much for this beautiful day, Lord.".
And then, "Lord," and then now you're all gonna get.
really self-conscious about public prayer or something..
But like, even Lord,.
that He is the master of all things, King,.
and it becomes meaningless..
It's just another word..
Because we don't mean it when we say it..
Or Father, or maybe even the name of Jesus..
It just becomes something that we automatically insert.
into every sentence at the beginning,.
in the middle, at the end, you know?.
As if somehow, like the pagans think,.
that their many words are gonna get God to hear them..
God is not a respecter of lots of words..
So when we pray, let's pray,.
meaning what we say and saying what we mean..
And here's another thing..
You know, God is not like other gods, and He's not,.
you know, though He is our Father,.
He is not like our earthly father..
And here's what I mean..
How many parents do we have in this room?.
Well, we got a young room, not a lot of parents..
Two parents..
Well, when you're parents, you're gonna experience this,.
and this is where your children will,.
well, we've all been children,.
so maybe we remember when we've done this,.
that we negotiate with our parents..
That we might even say things like adore them,.

$^{481}$love them, flatter them, like worship and praise them,.
just so we can get what we want..
God's not like that..
We don't adore Him, confess our sins, and thank Him,.
just so we can hurry up and get to the list of things.
that we need, thinking that if we butter Him up,.
if we get Him in a good mood,.
He's gonna give us what we need and what I want..
Or we might bargain..
God, if you'll do this for me, I'll do this for you..
You know what my favorite one is?.
God, if you make me rich, I promise I will start to tithe..
Give me a million dollars,.
and I promise I'll start to tithe..
And you know how ridiculous that is?.
It's like, God, give me a million dollars,.
and I promise I'll give you 100,000 back..
What a great deal for God, right?.
But that's how we pray sometimes..
We're thinking we're doing Him a favor in these deals..
Let's not do that..
That's not the kind of God we serve..
That's not the Heavenly Father..
Are we using many words, strategic words, insincere words,.
maybe even empty words because we're forgetting.
who we're talking with, who we're dealing with,.
that He's not impersonal?.
He's not an absentee father,.
but rather we are speaking to our Heavenly Father.
who loves us..
Yes, we believe our God is our Father,.
but do we really believe what kind of God He is?.
And then Jesus reminds us again,.
He doesn't respect words.
because that's 'cause He's not an empty God..
He's not a dead God who receives empty prayers.
from empty pray-ers..
So He says in verse seven again,.
"When you pray, do not keep on babbling like pagans,.
"for they think they will be heard.

$^{521}$"because of their many words.".
Second, though, we have to remember that our Father.
is one who already knows what we need..
In other words, when we give our list of needs,.
God's not like, "Wow,.
"here's information I did not have..
"I'm so glad He gave me a detailed list.".
Or it's not like time goes by.
and we forget to tell God we need something,.
and God says, "Oh man, why didn't you tell me.
"you needed that?.
"I would've taken care of it,.
"but you never told me..
"I had no idea that was going on in your life..
"I didn't know.".
God is a Father who already knows what we need..
You know, in a season when our prayer life.
can become very imbalanced,.
you know, when it seems like there's so many needs,.
that the list at the end.
becomes the large part of our prayer.
that we kind of rush through praising and worshiping,.
and we rush through.
sincerely confessing our sins.
and asking for forgiveness,.
where we rush through and try, you know,.
have a hard time remembering,.
"What has God done for me lately?.
"I gotta thank Him for something..
"I gotta come up with something.
"just so I can get to the list.".
Let's not forget.
what kind of God, what kind of Father is He?.
That He is one who knows what we need,.
and in knowing what we need,.
we have a God who is able,.
He is capable of providing what we need..
Do we believe that?.
Amen?.
This side doesn't sound convinced..

$^{561}$Amen?.
Can He take care of us?.
Yes..
But here's something that's equally important..
Yes, He is able to take care of us..
The Bible tells us He is a king.
that owns cattle on a thousand hills..
I mean, maybe like in today's vernacular,.
that might be like,.
He's a car dealer who owns a thousand lots of cars,.
you know?.
If you need a car, He's got one for you, okay?.
By the way, I don't know how many times.
God has provided a car when I've desperately needed a car,.
and others have said the same thing to me..
But not only can He do it,.
what's equally important is that He is gracious..
He cares, He loves us, so He wants to do it..
He wants to take care of us..
He wants to bless us..
He doesn't want us to go without..
He wants to bless us..
He wants to make sure that He is supplying us..
These two things, it might seem obvious,.
but these two things in combination.
are so incredibly important..
You see, if He was a powerful and wealthy God,.
it would be like having a father who is a billionaire,.
but he never asks, "Hey, did you have lunch today?.
Hey, there's holes in your clothes..
Do we need to get you a new shirt?.
Hey, you look sad today..
How can I help you today?".
It would be having a God who is a billionaire,.
who has everything that a person could want or have,.
and then, but they don't care about us..
They won't lift a finger to help us..
But what if God was unable to help us,.
but He really cared about us, and that was all He,.
you know what?.

$^{601}$It would be like the impotent father.
who desperately sees his desperate children.
starving to death, freezing in the cold,.
having no roof over their heads, no bed to sleep in,.
and because He's powerless, He's helpless,.
He can do nothing..
But God is able, and He so cares about us..
There's no place I see it better than in the words of Paul..
I gotta ask, sister, I don't know if you read.
the sermon manuscript before I preach today,.
but I hear she quoted the very verses..
In Romans 5, six through eight, Paul says,.
"You see, at just the right time,.
when we were still powerless, you know, helpless,.
Christ died for the ungodly..
Very rarely will anyone die for a righteous man,.
though for a good man, someone might possibly dare to die.".
And here it is in verse eight..
"But God demonstrates His own love for us in this..
While we were still sinners, Christ died for us.".
Our father knew exactly what we needed.
more than anything else..
We had a problem that was the elephant in the room..
We could not get past it, we could not ignore it..
It is a sin problem that we have offended.
our heavenly father, our creator..
And God had the solution because we were helpless,.
powerless to do anything about it..
And the solution was in His son..
And God could say, "Oh, no, no, wait,.
the price is way too high..
I cannot give up my son for these losers..
They're my enemies..
Why would I sacrifice my son for them?".
But He didn't do that because He demonstrated.
His love for us like this..
And while we were still sinners, still His enemies,.
He gave His son..
Jesus died for us..
Amazing..

$^{641}$This is amazing love, the song that we sang..
This is amazing..
We needed a savior..
The father was able to provide a savior..
And because He loves us and cares for us,.
He was willing to pay the price,.
that high price so that we could be saved from our sin..
Jesus' words on prayer reinforced for us.
that God desires an authentic relationship with us..
Not just religiosity, not just piety,.
not just ritualism, He wants a relationship..
Come close to me..
Come close, come even closer..
You're still too far away, come closer..
And He gives us access through faith in Jesus Christ,.
His son..
If you don't know Jesus as your Lord and savior,.
why wait a moment longer?.
Trust in Him today..
He said Himself, no one comes to the Father.
except through Him..
I'm gonna ask you to just bow your heads for a moment..
And I'm gonna ask you to,.
we'll just go through the acts, the acts of prayer..
God, thank you that you are indeed a loving Father..
Thank you that you are a God who is powerful.
and you are able..
We don't have a weak God, we don't have a distant God,.
we have a God who is intimately close..
A God who doesn't have to hang out with us,.
but God who wants to..
We praise you for being a God like no other..
For being our heavenly Father..
And we confess our sin to you,.
we admit that we have sinned against you.
when we have tried to live life our own way..
God, we have said no to you..
We have said your way is not good enough..
Your way isn't good enough for me, it's not my way..
And we have chosen to rebel and we ask you for forgiveness..

$^{681}$For every choice that we have made,.
which only confirmed that we are in rebellion against you..
But Lord, thank you that you have sent your son, Jesus..
Knowing exactly what we need, a savior,.
you have provided a savior,.
the perfect Lamb of God in your son, Jesus Christ..
And you have said to us,.
if we will just believe in your remedy.
for our sin problem, the solution to our sin problem,.
if we would just receive in faith that remedy.
that you have said, you will forgive us..
You'll give us a new life, a new identity,.
no longer an enemy of God,.
but now a child of God, a friend of God..
That you promise us eternal life and not death..
Thank you for the gift of salvation in Jesus Christ..
And Lord, we do lay before you all our needs..
And we know that we're going to keep being tempted..
We're gonna keep being distracted..
We're gonna think that things that are fool's gold.
is real gold..
We're gonna think that what isn't valuable is valuable..
Our priorities are gonna get messed up.
and we need your help, Lord..
Teach us, teach us how to get it right..
We can't do it without your help..
And there are other needs in this room,.
financial needs, concerns about career or education,.
relational issues, broken families, broken friendships,.
broken relationships with, in our workplaces,.
in our schools, Lord, with this world,.
we pray you will heal those relationships..
And God, we also ask, again,.
that you would teach us how to pray..
Because we want to have that relationship with you,.
that communication with you,.
that Jesus has now, he's wet our appetite.
and has shown us how much better.
how much more there is for us..
Lord, teach us how to pray..

$^{721}$I'm gonna ask you to open your eyes for a moment..
And we're gonna finish prayer this way..
'Cause Jesus did give us a specific prayer..
And it's probably one that everybody in the room.
has memorized..
But how many times when somebody says,.
"Okay, let's pray the Lord's Prayer together,".
we're able to just, you know,.
we could be like texting somebody.
and we could still do the Lord's Prayer, right?.
You're watching a video, YouTube,.
and we can recite the Lord's Prayer..
But how about right now, we'll just focus,.
pull away from all the distractions,.
and let's look at the Lord's Prayer,.
the prayer that Jesus taught us afresh..
'Cause he said, "This then is how you should pray.".
Let's do it together..
Our Father in heaven,.
hallowed be your name..
Your kingdom come,.
your will be done on earth as it is in heaven..
Give us today our daily bread..
Forgive us our debts,.
as we also have forgiven our debtors..
And lead us not into temptation,.
but deliver us from the evil one..
And Jesus invites us to an authentic relationship.
with his heavenly Father..
And I ask you,.
do you accept that invitation?.
Are you all in on this relationship?.
I hope your answer is a resounding yes,.
absolutely, that's what I want..
That's in Jesus' name we have prayed, amen..
[BLANK AUDIO].
\newpage



\section{}
\label{sec:_2c5tRWbVyU}
\textbf{The Road Unknown - Joshua 3 (Speaker: Tommy Lee)}
\newline
\newline
連結: \href{https://youtube.com/watch?v=-2c5tRWbVyU}{\texttt{https://youtube.com/watch?v=-2c5tRWbVyU}} ~~~~ 語音日期: 2024-02-29
\newline
\newline
\hyperref[sec:_33GO1favQo]{\small{< < < PREV SERMON < < <}}
~
\hyperref[sec:index_chronic]{\small{[返順時目]}}
~
\hyperref[sec:index_scriptual]{\small{[返順卷目]}}
~
\hyperref[sec:q3J7I2JVwVE]{\small{> > > NEXT SERMON > > >}}
\newline
\newline
$^{1}$right? Tommy, I don't know how you can do so much. We want to welcome Tommy. Let's welcome him..
Well, thank you very much everyone. And Pastor Mark, it is great to be here with all of you guys.
this morning. And you know, I never thought I would be on this journey that God has called me.
on. And when I graduated from high school, my parents, many of you guys who are listening are.
Asian and Chinese. So my mom says to me, "Pick a goal. Doctor, engineer, or lawyer? Pick one.".
So I said, "Okay, I guess I really enjoy chemistry. I guess I'm going to be a doctor. I want to apply.
for Washington University in St. Louis." And got into their medical program, a seven-year medical.
program, right? And from that point on, after the first semester, I realized, man, God did not call.
me to be a doctor. I'm not sure if I enjoy doing that. And so as a result, I went to the business.
school at Washington University in St. Louis. Graduated from that. And then from that point,.
I built cell towers, telecommunications towers for about 12 years. But during that period of time,.
I'm at my church in Chinatown, where some of you guys were. And so I'm working with our high school.
and college students. And the big question I asked myself is, "Lord, are you calling me to go into.
full-time ministry?" And my thought was full-time ministry at that time was pastoral ministry,.
right? So now I'm taking classes at night at Moody Seminary here in Chicago. And so after I.
graduated, I went and worked in a church, a Korean church. I wanted something new of all places. And.
for about three years, I was a pastor. Realized, man, God did not call me to be a pastor. I'm not.
sure if I enjoy doing that. I just was not wired to do that. So then from that point, I went to.
work at a Bible school, Moody Bible Institute, where my title was special assistant to the.
president. I did major donors, Founders Week, and Pastors Conference. I realized the non-profit.
ministry is too slow for me. It's too hierarchy. It's too top-down. I'm not sure if I enjoy doing.
that. So now for the last 15 years, I get a chance, as Pastor Mark says, we started Resource.
Global that works with young Christian marketplace leaders in different global cities. We started in.
Jakarta, Indonesia. Then we went to Singapore and Chicago and Kuala Lumpur, Nairobi, Austin..
From that point on, we're starting one in Los Angeles, in South Africa, here at Houston in.
Sydney in the next couple of years. And then from that point on, Northern Seminary called and says,.
"Can you do the same thing that you did at Resource Global with ministry leaders and pastors?".
Because many of these pastors from around the world, ministry leaders, Christian lay leaders,.
cannot afford your full-time seminary that many American seminaries offer. So then we started.
doing that. I get a chance to now operate as the executive director of Together LA. We opened up a.
business in Indonesia called City Leadership Center. And so I get a chance to work with different.
other groups as well too. But the whole idea is, how has God trained and created you? And we're.
going to talk a little bit about that because all of us individually are designed uniquely by God..
Each one of us has a certain story that's unique to all of our lives. As you're sitting here.
listening to me right now, your story is unique. Your story is different. And God has made that.
story unique with all of the adventure, with the different chapters in your life..
Now, where are we? Why are we calling us to go? What do we need to do? That becomes the question.
that all of us will have to ask. And if we could start off our sermon title today is "The Road.
Unknown." And we're going to take a look at Joshua chapter 3, verses 1 through 8. Joshua chapter 3,.
verses 1 through 8. And let me read that as well too..

$^{41}$"Then Joshua rose early in the morning," I'm reading from the ESV, "and they set out from.
Shittim. And they being the Israelites came out to the Jordan, he and all of the people of Israel,.
and lodged there before they passed over. At the end of three days, the officers went through the.
camp and commanded the people, 'As soon as you see the Ark of the Covenant of the Lord your God.
being carried by the Levitical priests, then you shall set out from your place and follow it..
Yet there shall be a distance between you and it, about 2,000 cubits in length. Do not come near it,.
in order that you may know the way you shall go, for you have not passed this way before.'.
Then Joshua said to the people, 'Consecrate yourselves, for tomorrow the Lord will do.
wonders among you.' And Joshua said to the priests, 'Take up the Ark of the Covenant and pass out.
on before the people.' So they took up the Ark of the Covenant and went before the Lord..
The Lord said to Joshua, 'Today I will begin to exalt you in the sight of all Israel,.
that they may know that as I was with Moses, so I will be with you. And as for you, command the.
priests who bear the Ark of the Covenant, when you come to the brink of the waters of the Jordan,.
you shall stand still in the Jordan.'" We're going to look at five questions this morning.
that emerges out of these eight verses, questions that emerges that guides us as we look at these.
different verses. And we're going to see exactly what God is teaching us in this theme of the road.
unknown, because all of us who are sitting here will have to cross that road unknown. All of us.
sitting here will get to a point where you don't know what to do, where you're trying to figure.
out, "Lord, what is it that you're calling me to do? What is it that you want me to do to continue.
to follow you?" Five questions. Number one, first question is this, and I hope a lot of times, and.
Mark, a lot of times I am trying to figure out, here we go. Question number one, what happens.
to Joshua in chapter one that sets us up for what happens in chapter three? Chapter three,.
verse one begins with this word, "Then Joshua rose early in the morning, and they set out from.
Shittim." They rose early in the morning, and he knows that this day is there. After 40 years.
of living in the promised land, after their time in Egypt and slavery,.
it is now time to pack up. It is now time to go into this promised land..
And so the question begins, what happens to Joshua in chapter one and two that sets us up for.
what happens in chapter three? If you remember what happens in chapter one, some of you guys.
have read this. He tells him at the beginning, verse two, "Moses, my servant is dead. Now therefore,.
arise, go over to Jordan, you and all of this people into the land that I am giving to them,.
to the people of Israel." God tells Joshua, "Joshua, your leader Moses, the guy who has been.
your mentor for so many years, the guy who led you guys out of slavery in Egypt, the guy who took.
you guys through the Red Sea and parted the Red Sea, you know that person? The one who brought.
you the Ten Commandments, who you thought was going to lead you to the promised land? Joshua,.
he's not going to be there. Moses, your leader is dead. It's on you now, Joshua. You're going to be.
the one that leads everyone there. You're going to be the one that will have to bring them to the.
promised land." And he says to this, "Only be strong and very courageous. Be strong because.
you're going to have to lead this group after 40 years in the wilderness to a place where they.
don't know. You have to be strong and very courageous. And how you're going to do that,.
you have to be careful to do according to all the law that Moses, my servant, commanded you..

$^{81}$And you're going to be strong and successful when you do not turn from it to the right hand.
or to the left, that you may have good success wherever you go. Do not let the book of the law.
depart from your mouth, but you shall meditate on it day and night so that you may be careful to do.
everything according to all that is written to it." Chapter one begins with the whole fact that,.
let's put the whole situation in reality. Moses, your leader. Moses, the guy you followed. Moses,.
the person you're looking to. He's dead. Let's just get that out of the way. Joshua, you're.
going to be the one that leads them to it. And Joshua, it is not going to be easy because later.
on he finds out in chapter two when he sends the spies over to Jericho, they're not welcoming them.
there. Those guys in Jericho, those guys in the promised land, they heard that the Israelites are.
coming. There's this God who brought them out of Egypt through the Red Sea. They've been living in.
the promised land, this wilderness for 40 years. And now they're about to cross the Jordan River.
and they're coming to the promised land. They're not just going to come over and say, "Here you go,.
everyone." They're going to try to get rid of you. You're going to take their home. And Joshua,.
it will not be easy. People are going to be complaining. People are going to be critical..
You're not going to know what to do. You're going to be scared. So Joshua, my words to you is to be.
strong and very courageous. Very, very courageous. And how you're going to do that, we're going to.
talk more about that. Do not turn from the words of this mouth. To meditate on it day and night..
And those are the words that Joshua hits. And we head into verse one. He rose early in the morning..
He knows what is at stake. And he sets out from Shittim. And he tells them,.
"And they came to the Jordan, he and all of the people at Israel.".
Today is a day which leads us to the second question..
As the Israelites are going into packing up for the last time, ready to go to the promised land,.
what do you think was going through the minds of the Israelites as they were packing up for the.
final time in a home they lived with for 40 years? For 40 years they lived in this wilderness..
For 40 years some of the kids were born in that wilderness. For 40 years they lived on this thing.
called manna, a rice little cake. For 40 years they lived on it. It's not the best, but they.
lived on it for 40 years. At night for 40 years a flame came in and gave them warmth and gave.
them light. For 40 years in the wilderness this is what happened. And after 40 years the day comes.
and now they are ready to pack up. What do you think went through their mind as they were packing.
up knowing that after 40 years in the home they lived in 40 years? It is now time to go. For some.
of you guys you've lived in the home for 10 years. You've lived in your parents' home for 20 years..
Your parents, look my parents lived in Chinatown, Chicago. They've lived there for over 20 years..
There's a lot of stuff if they have ever have to pack up and leave. I just moved in the house that.
I lived to within four years. There's a lot of stuff to pack. They now live, they are ready to.
leave a home they live for 40 years. It depends on who you ask, right? There are some, the older.
ones who came out of Egypt into the province land, they probably passed. Some of the younger guys.
like Joshua who are now older, they were probably in the prime of their life when they were in.
slavery in Egypt. They remember the slavery in Egypt. They remember packing up and leaving through.
the Red Sea. They remember all that. They remember the early days when Moses brought them there in.
the Ten Commandments. These guys were the prime in their life. They were the leaders. They were.

$^{121}$the warriors back in those days. They now are the elders of this group of the Israelites. A day.
they have been thinking about for 40 years. The day where God has promised them for 40 years.
is finally there. There is a little worry, but there is excitement. That's that group..
You have some of those who were babies when they were just in Egypt. They don't remember anything.
about what happened in Egypt. They remember probably sitting on the cart and donkey going.
through the Red Sea, but for the majority of their homes, their type, they've lived in the wilderness..
That has been their home. They are now your leaders. They're now your warriors. They're now.
in the prime of their life. They're worried because if there's battle, those guys are the.
ones who will have to take care of everything. And so yes, they do remember a little bit..
Moses is kind of, they kind of remember Moses, but we've already sent out two spies..
They're coming to get us. So there's a little hesitation. Maybe the youngest ones, the ones.
who were born in the wilderness, who are just growing up right now, they're excited, but similar.
to their parents. Man, why can't we stay in the wilderness? Why can't we just stay where we are?.
You know, when we go to the promised land, they're going to try to kill us. They're not going to just.
turn over their home to us, their lands. Can we just stay in the wilderness? We're fine with.
there. We've been living there for 40 years. Depends on who you ask. There's so many different.
feelings that go, that are going through their mind, but you're not going out. It's a community,.
isn't it? And we find this even in church as well too. That as Christ followers, you are not alone..
That even though I am speaking here and I'm just getting to know you, I don't know where some of.
you guys are watching from. It could be in Chicago. It could be other parts of the US. It could be.
parts of the other world. We don't know. But it's a community. With all of our differences in.
opinions, with all of our feelings, with all of our stories, our unique stories, we go together.
with our feelings, with our thoughts. And so it depends on who you ask with number two..
There's anticipation. There's fear. There's wondering. There's excitement. There's,.
are you sure we should be doing that? All of this stuff. But we go forth as a community..
And they came to the Jordan, he and all of the people in Israel, and lodged there.
for three days as they pass over. At the end of the three days, the officers went sent.
throughout the camp and commanded the people, which leads us to question number three..
Why do you think Joshua had them camp out at the base of the Jordan River for three days?.
They packed up. Now they get to the Jordan River..
Joshua, go do your thing and park the Jordan River. What are you going to do?.
I want you guys to sit there for three days. And I want you to stare at the Jordan River for three.
days. Do you know what's happening at the Jordan River during that time? Chapter three doesn't,.
you don't need to turn to it, but here's what happens in chapter three, verse 17..
And it tells us about the Jordan River. Now the Jordan is overflowing all its banks.
during this time of harvest is what happens, it says in chapter verse 17..
So Joshua, you are actually asking us to pack up, go to the Jordan River during this time of harvest.
in the middle when the waters are raging. And it's crazy. And you want us to cross the Jordan River.
during this time. And now for the next three days, we're staring at the Jordan River so we can see.
how high the rivers are, how raging the waters are and what we're going to do. Really?.

$^{161}$Why do you think Joshua had those guys camp out at the base of the Jordan River for three days?.
Here's what I'm going to say about this. Don't ever underestimate how God uses experiences.
in your life to show his faithfulness. Don't ever underestimate how God uses experience.
in your life to show his faithfulness. They sit there at the Jordan River,.
they're staring at the Jordan River for three days as the waters are raging..
In order to show I'm going to do something work and I need you guys to see what you're going to.
deal with and what I'm going to do in your life. Because when you go to the Jordan River into this.
land that is unknown, you're going to have to grab upon these memories of how God has been faithful,.
how I have been faithful to you in order to guide that every step of the way. Joshua teach them.
how to use the words of God, help them to meditate. But now I'm going to show them memories.
in order for me to show I am faithful. Do not underestimate how God uses memories in your life,.
situations in your life to prove how faithful he is to you. Let me tell you a quick story as well.
too, okay? And then we're going to go on to the next point. I stand here right now, if you notice.
to me I have a lot of different water. Three years ago the doctors told me that I was diagnosed with.
cancer in my sinus. That cancer was going to destroy everything inside of my nose and all of.
it inside my mouth. I stand here today with no saliva, little bit of taste buds, my hearing is.
down to 30\% in my left ear, my ear is going bad, I'm down to about, eventually it's going to be.
down to 20-30\%. I have tears in my eyes, I have no tears in my eyes, I live with headaches and.
pain throughout my body every single day. Two and a half years ago when I was going through my chemo.
radiation, the doctor says you're going to be in pain and I never understood that pain. And I lived.
in that pain by week five of the radiation. I lived with pain every single day. That pain didn't stop.
everyone. That pain affected the inside of my nose, the inside of my mouth, the outside of my.
skin on my neck and I lived with intense pain and that led on for four weeks. I lost 80 pounds.
during those weeks because I couldn't eat. Every day I sat in my home, every day for those four.
weeks I cried to the Lord, take away that pain. I can't deal with that pain anymore, take away that.
pain. And you know what happened? The pain never left, went away. God was quiet during that period.
of time. And you know what I had to depend on every single day? I had to always remember God's.
faithfulness to me during those years. All of the memories of how He was faithful, I had to remember.
and I had to remember the truth of the scriptures here in this world to guide me those four or five.
weeks. Because folks, there are times in your life you're going to go through trials, you're going to.
go through hard times, you're not going to know what to do. And you're going to have to cling upon.
the words because you can't trust your heart. Your heart's going to tell you, be angry. Why is God.
doing this to you? But you're going to have to trust upon those words and cling upon those words.
as truth. And you're going to have to cling upon those memories to guide you every step of those.
days. That's what you're going to do. And do not underestimate that God will use memories in your.
life to guide you every step of the way. And He says this in this wonderful things, as soon as you.
see the Ark of the Covenant of the Lord your God being carried out by the Levitical priests,.
then you shall set out from your place and follow it. And watch this, underline this,.
yet there shall be a distance between you and it about 2,000 cubits in length. Do not come near it..
If you have an underlined pencil, highlight or highlight this. Do not come near it in order that.

$^{201}$you may know the way you shall go. For you have what? Not passed this way before. Underline it..
Underline it. You've never passed this way before. You don't know. I know this..
I've wrought, I gave you the promise of life. I've determined your path. I've shown you what.
you need to do. Follow the Ark of the Covenant. In our case, follow the Holy Spirit. When I came.
and died on the cross, I gave you a new life. I've given you the Holy Spirit. Follow the Holy Spirit.
in your life. Let the words of Scripture guide you. Let the memories of those guide you. But.
follow the Holy Spirit and allow you to really be able to give you direction for you don't know.
where you are going because you've never been in this way before..
This morning I get out of bed, I see this text. It is a pastor friend of mine..
He's a young adult pastor, youth pastor with two little kids, two girls. One just born,.
about a year old. Another one about five years old. He tells me, "Tommy, my wife left us..
She didn't want to live the pastor's wife life anymore. She picked up, left me with my two girls.
and moved to New York to do what she always dreamed about. I don't know what to do..
I have one of our other co-workers who recently said to me, "Tommy, my dad has a spinal injury.
and he needs spine surgery. My mom has dementia and they both live with us. I also just want to.
let you know my husband and I have been trying to go to marriage counseling for the last six months..
It's not working. We're thinking about separating but I can't separate because I've always been a.
stay-at-home mom until I started working part-time with you. I don't know if he divorces me and we.
separate. What am I going to do with money? How am I going to handle my parents? We have three kids.
in college. In this COVID season, some of you guys are dealing with loss of loved ones wherever you.
are. You're dealing in situations in many countries, maybe where you are, where it's not safe to be a.
Christ follower. You're dealing with loss of jobs or loss of incomes because of COVID. Your business.
is shutting down. You may be on the brink of being laid off. You're looking at your savings.
account. You're trying to figure out what to do. You're dealing with kids who are acting out. You're.
dealing. You just had a baby. You don't know what to do. You're arguing with your spouse. There are.
situations in your life where you don't know what to do because you've never gone this way before..
And God is saying, "I understand. I've written your story. Your story is unique.".
Know I am part of that story. And as I led you to the promised land, Joshua, the same words that.
ring to Joshua, rings to us. Meditate on your words day and night. Do not allow the words that.
depart from your mouth. Remember the memories of you, but follow the ark of the covenant. Follow.
the Holy Spirit because you have not been this way before..
Sorry about that. And next one, question five..
But consecrate. Oops, sorry about that. Consecrate yourselves..
I am so, I'm getting used to, Mark, sorry about this. I'm getting used to number five here..
And then from that point on, it goes to this. Consecrate yourselves for tomorrow the Lord.
will do wonders among you. Consecrate yourselves for the Lord will do wonders among you. I'm going.
to send you on a journey. I'm going to send you through the promised land. After 40 years,.
I have a plan for you, says the Lord. Follow the ark of the covenant for you have never been this.
way before. But before you do, consecrate yourselves for tomorrow the Lord will do what?.
Wonders among you. And that will be my last point. And what does it mean by wonders? Let's take a.

$^{241}$look at chapter seven, right? Sometimes here's the thing I want to encourage you. You're sitting.
there, consecrate. What in the world does it mean, consecrate? If you have Google, just do Google,.
go to BibleGateway.com. And in the book of Joshua, just do the word consecrate and see.
exactly how it appears. So here's what it says in verse 13 of chapter seven, the word consecrate.
appears twice in Joshua. It appears in chapter three and now it appears in chapter seven..
Here's what it says. Get up, consecrate the people and say, consecrate yourselves for tomorrow..
For thus says the Lord God of Israel, there are devoted things in your midst, O Israel,.
that you cannot stand before your enemies until you take away the devoted things among you..
Israelites, I'm going to do something powerful again for you. I'm going to lead you to the.
promised land, I'm going to write you a new story and there you will become one of the greatest.
nations. But it's not going to be easy. You're going to go into a period where it is the unknown.
and you're not going to know what to do. But I will be with you because I will give you the ark.
of the covenant. I will give you the Holy Spirit. But before you do, and I'm going to do wonderful.
things in your life, but before you do, you know what I'm going to ask you? I want you to.
consecrate yourselves. I want you to take the devoted things away in your life..
Where's your heart? What are the things in your life that God says to you when he says to.
consecrate? As you are sitting here right now, what does he have you to consecrate? What needs,.
what are some of the devoted things in your life that you need to come.
and deal with the Lord? Because he's going to want to do wonderful things about you. He has a story.
for you. He has wonderful things. He has a life planned out to you. But he's asking you to.
consecrate your life. To take the things in your heart, the devoted things in your heart,.
and get rid of it. And for some of us, it may be our plans. It may be the things that we want in.
our life. The things that we inherit, our dreams, our goals. It may be that. For some of us, it may.
be guilt that you're dealing with. You're sitting here listening right now and you're just dealing.
with guilt. Some of you may be anger or jealousy. All of that. When God says to you, "Take out,.
consecrate yourselves." What is he saying to you? What is he asking you to give up and come before.
the Lord, the devoted things in your life, and bring out to him? Because he wants to do a wonderful.
thing in your life. He wants to do something that you can't even imagine. He wants to write a story.
in your life that's unique and do things and lead you on that journey. But before you do, what I.
need you to consecrate yourselves to take out the things in your life that will hinder my relationship.
with you. That becomes the whole gospel message, isn't it? That Jesus Christ, upon the death of.
his cross, dies for our sins. We live a matter of grace. It is not a matter by works. He promises you.
that you have to come before the cross. Come before him and say, "Lord, I accept you into my life.".
And through that, he opens the door to a new journey. And he says to you, "Continue to follow.
the words I give you. Continue to follow the memories. Continue to follow the Holy Spirit I.
give you. But every single day, as I lead you on this unique journey to consecrate yourself, to.
sanctify yourself, to continue to cleanse yourself of your heart and the things that will bring.
a divide between me and you." And so that becomes the question for you..
What do you need to do to do business with God to consecrate yourselves with him?.
Let's actually stick with four questions because we're out of time as well too. The role to know..

$^{281}$What happens in Joshua chapter one that sets up for chapter three?.
You are my leader. Do not depart from the word that I meditate on day and night. That's number.
one. Number two, I'm going to give you memories to show you that I'm proving to be faithful. Do not.
ever underestimate how God is faithful in your life. Those memories, not only the word of God,.
but the memories. And he said, "Give you the Holy Spirit to guide you every step of the way.".
But we do it in community together. And that's what our calling is as Christ's fathers. And along the.
way, consecrate yourselves. For there are devoted things in your heart that will affect that I want.
you to do business with. Jesus is there always in all parts of your life, but he calls us to work.
and follow him in this journey, on this wonderful adventure that he is ready to leave you..
And he calls you on that path. Let's pray. Father, thank you for your words..
Thank you for your guidance. Thank you for how you speak to us each and every single day..
Thank you for the journey you give us. Thank you for how you give us your word. You give us memories..
You give us the Holy Spirit. But you challenge us to consecrate, to come and continue to do.
business with you. And so, Lord, I lift those things up to you, dear Father. And to pray for.
all of us who are listening, Lord, that we continue to hear your voice and follow you in this journey,.
that we come before you and lay it all at your feet. In Christ we pray..
\newpage



\section{2:1-11}
\label{sec:q3J7I2JVwVE}
\textbf{The Chosen - Sermon Series \#3 (John 2\_1-11) - Rev. Mark Chan}
\newline
\newline
連結: \href{https://youtube.com/watch?v=q3J7I2JVwVE}{\texttt{https://youtube.com/watch?v=q3J7I2JVwVE}} ~~~~ 語音日期: 2024-03-05
\newline
\newline
\hyperref[sec:_2c5tRWbVyU]{\small{< < < PREV SERMON < < <}}
~
\hyperref[sec:index_chronic]{\small{[返順時目]}}
~
\hyperref[sec:index_scriptual]{\small{[返順卷目]}}
~
\hyperref[sec:LtoT2yWzuEc]{\small{> > > NEXT SERMON > > >}}
\newline
\newline
$^{1}$Alright guys, welcome, and I want you to,.
let's see, connect with your neighbor,.
one of your neighbors, alright?.
I have a question for you, I have a question for you..
You have to give a quick answer to your neighbor, okay?.
What is the most, think of the most embarrassing moment.
in your life, embarrassing moment in your life, alright?.
Can you talk to your friend about it?.
What is that?.
What happened?.
The most embarrassing, most shameful moment.
that you can think of, alright?.
Can you think of any?.
There has to be an embarrassing moment in your life, right?.
So I want you to think of the most embarrassing moment..
Most embarrassing moment..
Now you don't have to go into details,.
but you can just, you know, tell them the basics,.
the basic information..
Most embarrassing moment..
(audience chattering).
Alright, so let's see..
Let me tell you one of the most embarrassing moments.
I ever had, is that many, many years ago,.
I officiate a funeral for someone who passed away..
Alright, it's someone's father, alright, someone's father..
So I officiate the funeral,.
and I actually know the deceased personally..
So he passed away, and I officiate the funeral,.
and a couple months later,.
I met the daughter of the deceased,.
daughter of the father who passed away..
And then we were chit-chatting,.
and I asked the daughter, I asked,.
"Hey, how's your father?".
(laughing).
And then she replied,.
"Well, you officiate the funeral of my father.".
I was so embarrassed, you know,.
because I officiate a lot of funerals, right?.

$^{41}$'Cause I can't keep track, I can't keep track..
So that's one of the most embarrassing moments in my life..
So maybe you find your own moments, right?.
But today's passage is about a shameful moment,.
a moment of embarrassment, a moment of shame, alright?.
So we are going through this series on Chosen,.
and one of the episode, episode five of season one,.
is actually called Wedding Gift, alright?.
Wedding Gift..
And the story is very simple,.
because it follows chapter two of Gospel of John, alright?.
So Jesus and his disciples were invited to a wedding feast.
at a place called Cana in Galilee..
It's a little town in Galilee,.
in the northern part of the country..
So they went, and then the hosting family,.
the Bridegroom family, prepare a big banquet, right?.
And at the end, all the wine has been used up..
And that means a major disaster,.
or major social shame of, you know, Jesus' day..
And then Mary, the mother of Jesus,.
comes and asks Jesus to help out, alright?.
And you know the story, Jesus turn water into wine, right?.
And then saving the honor of the family,.
turning shame into honor..
So that's how the story goes..
It's a very familiar story,.
and perhaps it's one of the beloved stories from the Gospel,.
one of the miracles that Jesus has done..
And it talks about Jesus and his disciples.
going into, joining a wedding banquet,.
and let's dive in and see what this story is all about,.
and how does it relate to you and to me..
John chapter two, verse one..
"The next day, there was a wedding celebration.
"in the village of Cana in Galilee..
"And the mother of Jesus was there,.
"and Jesus and the disciples were invited.
"to the celebration, to the wedding feast.".
Now, Jewish weddings in those days, or even in today,.

$^{81}$were very important occasions, right?.
Well, actually, that applies to all culture,.
across the cultures..
If you've been to a wedding banquet, you dress up, right?.
Hopefully you don't wear flip-flops, right?.
You dress up, and then one of the most important thing.
you bring with you to a wedding feast,.
or a wedding banquet, is what?.
Money, yes, yes..
You gotta give some money, right?.
You gotta give some gifts..
As a token of honor to the hosting family..
And it was a very important social occasion..
And often in those days, in Jesus' day,.
the whole community and the whole village is being invited..
And of course, the extended family,.
or friends of extended family, were invited..
So it's not surprising that Jesus.
and the mother of Jesus, Mary, and his disciples.
were also invited to celebrate..
So they went to Cana..
And then, if you go to Israel,.
there's a little town called Kefa Cana..
Kefa Cana is a small village..
It's still keeping the name in biblical time..
It's called Cana, right?.
And there's an ancient church there,.
and some people would go to the church.
and read the story, chapter two of 1 John..
But in order to understand the magnitude of this story,.
we need to understand the biblical culture..
In Bible times, culture is often called.
honor and shame culture..
Honor and shame culture..
Now, what is this?.
On the right, you see some words like glory,.
honor, good name..
Now, these are the things that people are striving for..
These are the things that people are trying to protect.
and to defend..

$^{121}$Defend your own glory, defend your own good name,.
your own reputation, and that is very important..
And for many of us who are Asians, who are Chinese,.
we do understand this kind of culture.
because we have this faith culture, right?.
And we wanna save faith all the time, right?.
To people, Caucasian, Caucasian culture,.
or the American culture,.
sometimes they go to,.
we enjoy a meal together..
At the end of the meal,.
they don't understand there is something called.
fighting to pay the bill..
We will do our best to pay the bill, right?.
And why is that?.
You know, other cultures would not understand.
what is going on here, right?.
Why do you have to fight for paying the bill, right?.
Because paying the bill in the Asian culture,.
in the Chinese tradition, means honor and faith, right?.
You earn your respect and you save your faith.
by I pay the bill, right?.
So you owe me, right?.
I know a lot of people don't do this these days..
Yeah, which is, it can be good, right?.
But we are fighting for social honor..
And that is what this culture is all about..
It's very, very interesting.
because you are fighting for something.
that is not tangible..
You cannot touch it, you cannot grab it,.
but it's so real..
It's honor and shame..
So honor and shame is kind of like a social credit rating..
Here, the diagram I show you..
When you get more likes on your TikTok,.
on your Instagram, on your social media,.
you have more social currency..
So honor can be thought of as social credit rating..
If you get more likes, if you get more praises,.

$^{161}$then you get better ratings..
And you get more currency, social currency,.
in your society, in your group..
And you are better..
And what is more important to observe in Biblical culture.
is that social honor is a group thing..
It's a group thing..
It's not individualistic value..
Now, in American culture, in the Western culture,.
or Westernized areas,.
the culture is more individualistic..
So we talk about honor, good name,.
and reputation is just about me, right?.
Mark Chan has a good reputation..
And the reputation is pretty much individualistic, personal..
It's a personal currency..
People don't care whether my son or my daughter.
has the same reputation, right?.
Because we are very individualized..
But in Biblical times, in Biblical culture,.
honor and shame is a group thing.
because you don't represent yourself only..
You represent a group,.
a big group, your family, your extended family,.
your uncles, your aunties,.
your cousins, your niece, your close friends..
So when you have honor, when you have a good name,.
when you have good reputation,.
your whole tribe will have honor..
But if you lose it, when you are being shamed,.
your whole family, including yourself,.
will face social stigma..
Now that's how it goes..
So you have to understand the magnitude of this story.
in Gospel of John chapter two..
You have to understand this cultural background..
What is going on?.
What is the big deal?.
In Jewish culture, banquet is the most important occasion.
you can show off and you can defend the honor.

$^{201}$of your family, of your tribe..
Some of the families in Jesus' time,.
they would actually invest all of their energy.
and all of their money, all of their resources.
in order to make the banquets look good.
and all the guests are being entertained..
And of course, wine is indispensable in those days..
Even today, we go to a banquet..
We want to open bar..
Hallelujah, we will have open bar..
People can drink all they want, right?.
And that's honor, right?.
The hosting family would like to offer.
something like that, right?.
The more you offer, the more lavish your banquet is,.
you know, you get more social currency..
So some of the families, even small families,.
they would use up or they even take on a loan.
in order to set up a very, very nice and lavish banquet..
And sometimes the feast can go on for seven days, right?.
Seven days..
Today, when you go to a wedding banquet,.
it's probably just one evening, right?.
But there is some Asian culture, like Vietnamese culture,.
they still do this, you know?.
They have a whole week of feast and invite all the friends.
and all their related to come to enjoy the meal, right?.
So this is the Jewish culture..
So in the story, when the wine supply went out,.
it means a social disaster..
Now I had you in the beginning to think about.
a moment of embarrassment, right?.
Now those moment, I hope, were not disasters, right?.
We're still alive, right?.
So it was not disasters..
But in those days, if you fail in your wedding banquet,.
it means social disaster..
And the worst form of social shame will happen to you.
because your name will be smeared for decades to come,.
for decades to come..

$^{241}$So Mary, the mother of Jesus, came in.
and asked Jesus to help out..
And ask him, you have to help them..
You have to help them, not only we solve the issue.
of the lack of wine, but to save the honor of the family,.
to defend the honor of the family..
It was a very, very bad situation..
And Mary, as the mother of Jesus,.
and you know Jesus, how old was he?.
He's already over 30 years old, right?.
So as you know, if you're a mother.
or you're a sibling of someone,.
you observe someone for decades,.
and you know this guy is special, right?.
Maybe he could do magic or something like that..
So mother of Jesus, Mary came,.
and I think she knows that her son is special,.
and he can do something about it..
How to do it, what to do, she doesn't know..
But she just trusts that her son, Jesus,.
has the ability to save the family in this situation..
And Jesus respond by saying,.
this is very interesting reply..
He said in verse four, "Dear woman, that is not our problem.".
Jesus replied, "My time has not yet come.".
Now, dear woman is actually a very polite way.
of addressing a woman or a mother, all right?.
So there's nothing wrong with that, all right?.
So we don't do that, right?.
I don't call my mom dear woman, what's up, right?.
We don't do that..
But in those days, you know, it's just a normal title, right?.
But Jesus' reply is interesting.
in that it's being highlighted here..
"My time has not yet come.".
What does it mean, right?.
"My time has not yet come.".
Now, this phrase, Jesus said, "My time is not here yet.".
"My time has not arrived yet.".
"My time has not yet come.".

$^{281}$Something like that, similar phrases,.
occur many, many times in the Gospel of John..
And it shows up frequently in the Gospel..
And Jesus often say, "My time is not here yet.".
And what does he mean?.
He said, "My time, the time of mine, is the hour,.
"is the time that Jesus will display all the glory of God.".
The whole glory of God, he will display,.
he will show the whole glory of God..
And that is the time he's referring to..
And here is the, in the early years of his ministry,.
and Jesus, of course, would say,.
"Hey, my time is not here yet..
"I'm not ready to show all the glory of God at this point.".
Because if Jesus displays the whole glory of God.
in a perfect way, in an immature time,.
people would have misunderstood why Jesus came, right?.
People would be expecting someone who can be their leader,.
right, to solve their political problem,.
or solve their personal problem..
But that's not why Jesus came, right?.
That's not the real glory of God..
So it is in this context that Jesus said,.
"My time has not yet come.".
But I think when he said that,.
Jesus already knows what to do in this situation in Cana,.
trying to save the honor, or defend the honor of the family..
He already knows what to do,.
because he understand the level, or the magnitude,.
of social disaster for the family..
And that's why the mother said,.
"Hey, whatever Jesus tells you,.
"tells his servant, you do it, okay?".
So we come to verse six..
"Standing nearby were six stone water jars.
"used for Jewish ceremonial washing..
"Each could hold 20 to 30 gallons,.
"and Jesus told the servants, 'Fill the jars with water.'.
"When the jars had been filled, Jesus said,.
"'Now dip some out and take it.

$^{321}$"'to the master of the ceremonies.'.
"So the servants followed his instructions.".
So Jesus knew what to do..
Now what is going on here is that.
in many of the Jewish households,.
you will find jars like that..
It's about this size, this high..
This tall..
And they're stone jars, and they would hold water..
Hold water for what?.
Because water is a very, very valuable thing in Israel,.
because they all depend on rainwater..
So they gotta have a lot of containers.
to contain usable water, right?.
So these jars were commonly used.
by Jewish families in Jesus' day to hold water.
and for purpose of purification..
Now you understand in Jesus' day,.
Orthodox Jews were cleaning freaks..
They loved to clean..
They loved to purify themselves..
And it's not a matter of hygiene..
You know, it's not because of germs,.
those that we have to wash our hands..
It's ritual cleansing that is their most concerned..
They're most concerned about..
Ritual purity is a very, very important value,.
because you have to make sure you are ritually pure,.
you're ritually cleansed,.
when you partake important occasions,.
like having a dinner..
You gotta wash your hands before you're having food..
Now again, it's not hygiene..
It's about ritual, right?.
It's a religious thing..
Or you go to the temple, you go to a synagogue.
to pray or to worship..
You gotta make sure you are cleansed..
So they have all sorts of ritual cleansing mechanism.
to make sure that they can stand in the presence of God.

$^{361}$and they will not lose their place..
So here is the background of the story..
There are six jars of this stone vessel, right?.
Now why six?.
In the Bible, perfect number is seven, right?.
I know in your culture, it may be eight, right?.
In your license plate, it may be whatever, eight, eight,.
eight, right?.
More eight, better, right?.
The more, the better..
But in Hebrew culture, the perfect number is seven, right?.
Perfect number is seven..
So only six jars here and that implies imperfection..
And what kind of imperfection is trying to imply here?.
In the Gospel of John, it's always trying to tell us.
the traditional value or traditional way of life, Judaism,.
the Jewish way of life is not enough..
It's not perfect..
Now in this situation, there is a social disaster, right?.
There's social shame being involved here..
And you can see six jars of, six stone jars filled with water.
cannot save this family from social shame..
Cannot turn shame into glory..
Only when Jesus comes and he can do something about it..
So Jesus had the servants fill the jars with water.
and then he turned water into the best wine, right?.
Best wine..
And how did Jesus do that?.
Well, the Bible doesn't say..
What mechanism did he use, right?.
We don't know because the Bible is not so much interest.
in the mechanism of miracle..
The Bible is more concerned about the effect of the miracle..
The effect of the action of Jesus..
What does it accomplish?.
What did Jesus accomplish through the miracle?.
As you can tell from the story,.
from verses nine and following, as we close this story,.
it says, "When the master of ceremonies tasted the water.
"that is now best wine," indeed, best wine,.

$^{401}$"not knowing where it had come from.".
Though of course the servants knew..
He called the bridegroom over and he said,.
"Posts always serve the best wine first.".
You know, because the banquet is so long, right?.
So many days..
So they usually serve the best wine in the beginning..
And then when people get drunk, right,.
they can't tell whether it's beer or wine.
or bad wine or good wine..
At the end of the banquet, they cannot tell..
When everyone has had a lot to drink,.
he brings out the less expensive wine..
However, you have kept the best wine until now..
Verse 11, "The miraculous sign at Cana in Galilee.
"was the first time Jesus revealed his glory.
"and his disciples believed in him.".
Now this interesting miracle shows the glory of Jesus..
Part of his glory, not the full glory of God.
being displayed here, but his miracle.
saves the family from shame, disaster,.
to honor and glory..
And that is the first miracle recorded by John..
And the miracle tells what Jesus does.
in the lives of the people..
And often we find ourselves in situation of embarrassment..
We find situation of shame..
But Jesus is able to save us from shame..
Now you think about it, this story is being recorded.
and being included in the gospel..
It's very meaningful..
Because that is what the gospel of Jesus is all about..
To save us from shame and embarrassment..
Save us from shame into glory..
Now let me close by reminding you.
what the gospel is all about..
If you read the Bible from the first page,.
almost to the end, read the whole Bible through and through,.
you will see a lot of examples,.
a lot of lives of people that have been saved by God,.

$^{441}$saved from shame to glory..
Let me give you some examples..
Adam, the first man, when he sinned against God,.
when he rebelled against God,.
he was hiding from God, right?.
And he discovered that he was naked..
It's a metaphorical way to describe shame, right?.
Guilt..
And Adam, a sinful Adam, was shame covered, right?.
But God made a garment to cover his shame, his guilt..
The Israelites in Egypt, they were in a shameful situation..
They were being enslaved..
But God saved them by grace out of Egypt, the Exodus..
And Rahab, you remember the Israelites,.
after 40 years of wandering in the desert,.
and they came into the land of promise,.
and they came to Jericho,.
and they found this woman called Rahab..
And Rahab was a prostitute,.
and he was doomed to destruction..
But she turned to Yahweh, the true God,.
and she was saved from shame to honor..
And Ruth, you know, the great-grandmother of King David,.
she was originally a foreigner, a Moabite,.
someone despised by the Israelites' society.
because they are just foreigner, Moabites..
She was a widow, so, so shame..
But she was redeemed, being included.
into genealogy of King David,.
and indeed, the genealogy of Jesus, from shame to honor..
The Jews who went to exile in the Babylonian period,.
they returned from exile, and God saved them,.
again, from shame to glory..
Now to the New Testament, Peter, the head of the disciples,.
the leader of the disciples,.
three times he denied Jesus, right?.
He was being shamed, failed, but Jesus came back..
He rose from the dead,.
and he came to Peter and restored him..
Even our Lord Jesus, he suffered shame,.

$^{481}$but honored by God's glory..
Now, as I remind you, Good Friday is approaching..
Easter is approaching..
So often when we think of Good Friday,.
Jesus being crucified on the cross,.
we think of something physically painful..
Of course, it was painful, extremely painful..
But the cross, crucifixion, is not just about pain..
It's about shame..
That's why victims of crucifixion.
is always being crucified, being displayed.
in a public way, right?.
On the street..
But our Lord Jesus suffered shame,.
but honored by God's glory..
He died in shame, but he rose in glory.
and seated at the right hand of the Father..
And that is the essence of the gospel,.
turning our shame into glory..
And that is only available in Jesus Christ..
We were sinners, but we are now forgiven.
because of what Jesus has done for us..
We may have loads of guilt in our life..
We feel ourself unworthy..
We question our own self-worth..
We look in the mirror and we always think.
we are not worthy or we're worth less than other people..
Right?.
Physically or even our own life..
We just cannot compare with other people.
and feel kind of shameful..
But lo and behold, the gospel tells us.
Jesus raised us with Jesus..
Jesus raised us from death, from shame..
And we stand totally righteous before God..
And that is the gospel..
We may feel unworthy..
We may feel sin-loaded..
But Jesus came to forgive all of our sins..
And that is the freedom of the gospel..

$^{521}$Right?.
Praise God..
Let's all pray together..
And come before God..
Yeah, we know we have been shamed..
And sometimes we shame ourselves..
Right?.
We feel unworthy..
Yet, at the foot of the cross,.
we are accepted, redeemed, and totally forgiven..
Let's all pray together..
Father, we want to give you thanks..
For Jesus, your son, has died on the cross.
to pay the price.
in forgiving all of our sins..
Father, thank you..
Thank you for saving us from sin, from guilt,.
from shame, from unworthiness..
And now we are your glorious sons and daughters.
standing righteous before you..
Only because Jesus has paid the price,.
paid for our debt,.
and suffered our shame..
You turn us from shame to glory,.
from dust to glory..
And thank you for your love..
May our lives be encouraged and transformed.
by your gospel,.
by what Jesus has done for us on the cross..
And thank you..
In Christ's name, we pray..
All right, thank God for seeing all of you..
And next week, we're gonna worship together..
And I know there's some freshmen out there, right?.
Yeah..
And enjoy our time of worship, fellowship together..
All right, God bless..
\newpage



\section{11:13-16}
\label{sec:LtoT2yWzuEc}
\textbf{客旅人生 (來 11\_13-16, 彼前 2\_10-11) -  黃茂林傳道}
\newline
\newline
連結: \href{https://youtube.com/watch?v=LtoT2yWzuEc}{\texttt{https://youtube.com/watch?v=LtoT2yWzuEc}} ~~~~ 語音日期: 2024-03-05
\newline
\newline
\hyperref[sec:q3J7I2JVwVE]{\small{< < < PREV SERMON < < <}}
~
\hyperref[sec:index_chronic]{\small{[返順時目]}}
~
\hyperref[sec:index_scriptual]{\small{[返順卷目]}}
~
\hyperref[sec:01fKINSJbMU]{\small{> > > NEXT SERMON > > >}}
\newline
\newline
$^{1}$大家早晨 聽到嗎?.
我知道你們的音響都很好.
試過都很喜歡.
我們先讀指定的經文.
希伯來書11章13到16節.
如果你有聖經的.
你可以跟我一起讀.
希伯來書11章13到16節.
這些人都是存著信心死的.
並沒有得著所應許的.
他從遠處望見且歡喜迎接.
又承認自己在世上是客裡是寄居的.
說這樣話的人.
是表明自己要找一個家鄉.
他們若想念所離開的家鄉.
還有可以回去的機會.
他們卻視每一個更美的家鄉.
就是在天上的.
所以神被稱為他們的神.
並不以為恥.
因為他已經給他們預備了一座城.
還有一處的經文.
在彼得前書的4章10和11節.
我想是…….
好.
應該可能是寫錯了給大家.
我想我不讀這裡了.
在這裡再一次多謝成會教會.
陳牧師和大家的同工.
你們的邀請來到當中來分享.
今天很想和大家講一個主題.
就是「客裡人生」.
我不知道大家喜不喜歡去旅遊.
科技發達 互聯網.
互享經濟 建立地球村.
什麼都很近 很方便的時候.
自助旅遊各種形式的旅行.
是以前20年30年前都是不方便的.
不可以有的.
現在有人這樣說.

$^{41}$是前所未有從海陸空.
包括深海潛水的方式.
去體會文化和不同地方旅行.
有一個很重要的.
武長的問答遊戲.
就是幾條.
你在心中 你知道的.
我答得到 挺好.
我認識這個世界.
我知道什麼是旅行.
第一條就問.
2023年全球統計.
究竟有多少人次的國際入境旅客.
是指飛的.
如果你陸路走路就不能統計.
心裡面你想一個答案好不好.
多少億呢.
答案就是13億.
第二條.
看你最後答到多少條.
國際入境旅客人次最多的是哪個國家呢.
B.
我忘記了.
亞洲人的答案多數是B.
其實答案是C.
是C.
最後.
這個關係我們很重大.
當年環球和平指數.
美國在安全旅行指數上的排名.
如果用168個國家.
因為他有某個統計.
168個國家裡面排行多少呢.
A 1到20.
B就是20到50.
這是安全的意思.
C就是51到84.
D是85以上.
就是一半以上.
都是嗎.

$^{81}$D最差.
多是嗎.
是呀.
是真的很差.
原來是百多個.
所以不好意思.
大家在美國生活.
原來差不多可以這麼說.
我們在美國出發的話.
去大部分國家都是安全的.
反過來.
回程那一下就大件事了.
回程那一下就危險了.
剛才讀的聖經.
其實是在說一些經文.
就是寫給去了別的地區的朋友.
別的地區的.
主要是一些猶太的信徒.
那些地區包括今天.
其實都很喜歡去的地方.
土耳其.
亞細亞.
加拿大那些地方.
無論是希伯來書那段經文.
或者是彼得前書那段經文.
都是給這些這樣的朋友.
給這些都是猶太背景的文化.
但是去了一個很不同的地區的朋友.
都是信了耶穌的.
但是大家留意一下.
這些移居去那些地方.
加帕多加加拉泰.
亞細亞地方的人.
他們不是去旅行.
你知道有什麼分別了.
我們過境的時候.
你去旅行的話.
很多時候會問.
什麼時候回來.
如果你說.

$^{121}$我不回來了.
很久之前有一位女士.
她想去英國.
去英國的時候.
過境的時候.
她說我還沒訂回程機票.
然後就問她.
那你怎麼生活.
她說我有錢.
我說多少錢.
她大概有2,300鎊左右.
你玩多久.
三個星期.
你有沒有信用卡.
沒有.
那你猜最後她能不能過境.
因為那個官都不覺得.
我看你都不可能回到.
再來這個機場走了.
又沒有機票.
又沒有信用卡.
不知道為什麼.
所以弟兄姊妹.
那個是旅行的.
你覺得你旅行.
應該是旅行的表現.
這些人去到這些地區.
他們不是旅行.
他們是移居那個地方.
彼得和寫希伯來書的作者.
就是告訴他們很重要的東西.
就是原來在這一幫.
各個地區信了耶穌.
又在那裡生活的.
但是猶太背景的人.
他們很需要知道一件事.
就是剛才讀的.
包括希伯來書十一章的這段.
信心.
這個很重要的主題.

$^{161}$就是這一幫人.
他們有一個難題.
如果他們在.
不是自己原居地文化的地方生活.
他們會有逼迫.
為什麼呢.
因為信仰的緣故.
可能我們在北美洲.
覺得被逼迫的機會很微的.
很少的.
我前幾天就看回.
有位牧師叫舒密.
我看回這個舒密.
他那時候就在街頭講福音.
他有一個慣例.
他很喜歡崇拜早半個鐘.
他就在對開公園的地方.
講幾個字福音.
講完就回去崇拜.
那天他講完之後.
很奇怪.
好像有些聲音.
就是左邊太陽穴.
好像被一些東西撞到.
他沒事.
他又去收拾東西.
收拾完東西.
然後回到教會那裡.
準備的時候.
那些人看著他.
牧師.
太陽穴有血.
原來他中了槍.
那些子彈卡在他腦袋裡.
馬上送去醫院.
我看新聞的時候.
幾個月.
昨晚想起他的時候.
就看他的更新.
本來醫生就想著他不能生存.

$^{201}$現在他竟然可以起來.
還可以生活.
但是暫時說不了話.
他為了這個牧師.
在美國祈禱.
原來在美國.
他在那裡的蓬蓉位置.
已經傳到了.
這樣街頭禱告.
一年多了.
都沒事.
不知道那天又會有這樣的人攻擊他.
還加上那顆子彈.
因為在一個很特別的位置.
是沒有辦法用手術拿走.
真的要求神憐憫.
所以你說亞特蘭提姐妹.
我們身處一個環境.
會面對壁壩的.
我們應該要怎麼做.
有人說不如我練習好一些槍法.
有些智慧的方法.
是不是有一個途徑呢.
有些人認為學智慧.
又或者什麼地方都不要去.
彼得和希伯來書的作者.
我相信他們都有一個看見.
其實我們要用一個方式.
怎麼去面對呢.
改變我們這個人的價值觀.
改變我們這個人的思想方式.
在我們內心裡面來有改變.
以致內心強大.
面對不同的事情都有力量.
好了.
我想說幾樣東西.
主要是四大點.
就是關於今天的主題.
嚇女人生.
不是旅遊人生.

$^{241}$是嚇女人生.
不是嚇壞女生.
而是我們看自己.
用幾個角度去看.
第一樣.
經文告訴我們.
憑著信來跟隨.
雖然我們今天所讀.
所領受的應許.
還沒有成就.
還沒有得著.
經文裡面就引一個.
很出名的猶太的教父.
猶太的先祖.
亞伯拉罕做例子.
亞伯拉罕就是一個怎樣的例子呢.
75歲才回應神的呼召.
離開迦勒底的梧椅.
其實75歲我們應該.
離開了工作崗位而已.
他75歲才領受這個機會.
然後去到迦南地.
住了一大段日子.
神給了他應許.
如果你看聖經裡面.
神應許什麼呢.
應許他有幾樣東西.
第一就是神會賜福給你的.
人生.
其中裡面包括.
你有後裔.
你會有地.
他們會成為大國.
所以這個應許真的很厲害.
很多福.
有地.
有國.
但是很重要.
誰承受呢.
所以神答應他有後裔.

$^{281}$當時他還沒有子女.
很振奮的應許.
很令他激勵.
但是等來等去都等不到.
大神最後在他100歲的時候.
給他一個兒子.
就是他應許所生.
他的元配.
撒拉所生的以撒.
但是地和國還沒有實現.
到他離開世界.
175歲的時候.
有什麼呢.
姑且這麼說.
三大類.
他起碼有兩大類的東西都沒有.
他還沒有真正有地.
他還沒有真正有一個國.
他是怎樣離開世界的呢.
就好像這個經文所記載.
他都是憑著信心這樣死的.
沒有得到這些應許.
但是他仍然是一個什麼人.
跟從神的人.
弟兄姊妹.
這個信心就是很寶貴的.
神學方面有一個說法叫做.
Already but not yet.
已經大患沒有.
有些東西.
神是應許給我們.
我們得著了.
大部分.
或者一部分.
但是還有沒有完全應驗的.
今天弟兄姊妹.
我知道大家遲些會浸禮.
其實在浸禮.
你能夠去受浸之前.
應該發生什麼事呢.

$^{321}$就是你真心跟隨.
接受主耶穌基督.
做你的救主.
跟從祂.
但是這個生命開始有一個改變.
但是你要發覺這個生命的改變.
按聖經說.
我們本應就好像主一樣.
不受這個環境限制.
穿牆入室.
可以吃又可以不吃.
又不會餓著.
還沒到.
我們形容這個就是.
完全的救贖的狀態.
那是什麼呢.
Not yet的.
但Already是什麼呢.
已經發生的.
就是已經有一個改變的生命.
裡面.
想事情很不同.
愛慕的事情也不同.
這個是一個改變.
在這個狀態很需要什麼呢.
雖然還沒領受應許.
仍然去跟隨.
這個就是這群人.
很特別的地方.
第二樣我留意到他們.
認定自己的身份.
是什麼身份呢.
客女和寄居的身份.
客女和寄居是兩個字.
我分開來講講.
客女有不同的譯本.
翻譯成外國人.
或者異鄉人.
Foreigner, Stranger.
那個字.

$^{361}$Xenos這個字.
其實可以解作外人.
或者局外人.
就是沒有你的份.
一個人不關你的事.
可能這個人是在這個地方.
到時時出現的.
但是他是局外人.
他是外人.
他跟那些人是沒有這種聯繫的.
有些什麼人會是.
完全沒有聯繫的人.
在關係上都沒有侵略你.
那你就是那些人.
好了.
客女就講完了.
寄居.
寄居也是翻譯成Foreigner.
但是有一個翻譯是Sojourner.
就是住來講.
就是說在外地來到本城.
本鄉來居住的.
在我們本地人中間.
我也記得.
其實拿了綠卡.
你知道你的身份是什麼嗎.
有些人拿綠卡就應該很什麼.
他不知道.
美國喜歡用什麼.
叫Resident Alien.
以前不知道.
我什麼外星人.
外星人.
不是.
他是有別於你是本地.
這個國家的人.
你還是外面的人.
從外地來居住.
那這個人包括什麼呢.
那個人當然是不同這個國家.

$^{401}$不同這個文化 語言的.
他就專門來到這裡居住.
這些就是寄居的.
簡單來講就是.
這一類人.
不是慣常是我們中間這樣的人.
弟兄姊妹.
那是什麼意思呢.
我們明明是芝加哥人.
或者我明明是在德薩斯州住.
或者我在加州住.
我是美國人.
我這個是什麼身份呢.
聖經就告訴你.
這個身份.
就是你不以這個世界.
你所身處的這個世界.
為你最終的家鄉.
為你的原居地.
原來你所屬的地方.
也就是說.
你效忠的地區地方是什麼呢.
不是這個世界.
就好像主耶穌說的.
他的國不屬於這個世界.
就是他這樣講法.
很權威.
就好像他有個國.
我的國不是這個世界.
就令到周圍的人都很驚訝.
有些人抓他的時候.
他說他大權柄到什麼地步.
如果他不想被人抓.
他可以求父.
叫十二營多個天使.
猶太人的編制.
十二營可以是過千人的編制.
有些人的講法.
可能是有五六千人.
就是一個營.

$^{441}$十二營可以有很多很多天使.
就是說他有很大的權柄.
因為他的主是屬於天上的.
那我們屬於他的人.
在這個世界來講.
我們看自己是客類和妓女.
明明我們是地球人.
現在就因為這樣.
我們相信了主.
原來我們身份.
對這個世界來講是局外人.
我們對這個地球來講.
我們是這個地球以外的屬神.
天上的人.
關鍵在哪裡呢.
我們雖然身處在這個地方.
我們認定自己是什麼人.
但是我可能和你很不同.
我雖然在美國事奉生活.
但是我以公民的角度來說.
我是加拿大人.
這個地區對我來講.
都不是完全自己的地方.
雖然在這裡生活.
同樣應用在我們基督徒.
在這個地上.
怎麼看自己.
就是說我雖然在這裡生活.
我不要將我人生希望.
所有的盼望.
包括我個人的理想力量.
全部放在這裡.
不要以為這個地上的所有東西.
包括你的工作,事業,家庭.
或者你所擁有的東西.
就代表了你自己一切.
怎麼看自己的身份.
身份來講.
有時候是有兩個角度的.
一個角度就是.

$^{481}$你要看自己主觀來講.
你覺得自己是什麼身份.
有些人在美國生活很多年.
都不覺得自己是美國人.
仍然是喝奶茶.
港式奶茶那種.
很認同這個文化.
你們認同嗎.
當然了.
我是香港人.
喝著那些熱樽裝.
我那個年代.
熱樽裝的玻璃樽的維他奶.
那些才好喝.
客觀上的身份就很不同了.
有一天我去看房子.
有個人看上去.
好像印度人一樣.
講著講著才發覺不是.
他是柬埔寨人.
是喔.
看起來差不多.
但其實真正的身份是會不同的.
這裡有個年輕人.
你猜.
你猜他是墨西哥人還是美國人呢.
這位年輕人從小就在美國長大.
他在2020年.
他在喬治亞州.
他17歲的時候.
他貪玩.
拿大麻 又玩又搗亂.
然後就被警察抓了他.
然後他因為有小噓大罵.
他又說自己的歲數多大.
最後就被抓了.
然後就發覺一件事.
其實他是美國人.
他一直都看自己是美國人.
去到警察局才發現.

$^{521}$他是undocumented的人.
即無證的.
原來他媽媽一歲不夠的時候.
就帶他從墨西哥來美國.
死了.
根本他連西班牙話都不會說.
只會說英文.
只會想美國的東西.
看美國電視.
結果很災難性.
長話短說.
求情.
又說不懂得找律師.
他就竟然自辯.
於是這個叫Brandon的年輕人.
就真是被在2020年遞解出境.
按他自己寫給報章的記者.
他說原來我十年以來都不能再回美國.
在他的印象中自己是什麼人呢.
美國人.
在實際來說.
他覺得他是墨西哥人.
因為他墨西哥出生.
反過來說.
會不會有些自己是美國人身份.
他忘記了自己是美國人呢.
有些孫教士去了外地生活.
他回到美國很不習慣.
他幾個星期就說不行.
我要回到我的工場.
他不覺得自己是美國人.
在曠野的以色列人.
他們看自己的生活是怎樣.
住在曠野.
是一個客旅.
是暫居.
住張牌隨時會動.
他們不會永遠將自己放在這個地方生活.
真正的基督徒會不會有一個毛病.
我在這個世界生活了很久.

$^{561}$這個世界好像是我的一切.
這個世界就是我最重要要建立的地方.
親愛的朋友弟兄姊妹.
我們認識了救主耶穌基督.
我們看自己是什麼身份.
有沒有人去旅行.
會將整間屋子背在他身上離開.
搬屋的時候可能有一個車這麼多東西.
你旅行的時候會不會.
我也認識一些人旅行很厲害.
我說你有多少個旅行行李箱.
我這次也算少了.
我帶十個.
當然他不止特別的會員.
帶無限的行李箱也可以.
我們覺得很誇張.
你旅行當然很輕便.
我們在地上作客旅.
我們用什麼心態.
我們會不會在這裡有幾百個箱子的東西.
想帶回天國.
帶不走的.
所以希伯留斯所說的這班人.
他有客旅的身份.
他們不只是很無奈地.
什麼時候回到自己真正天上的家.
他們有一個積極的態度.
他們羨慕天上的家.
經文中說他們羨慕天上的家.
我說他們是一家嚮往天家的一班人.
什麼叫做羨慕.
羨慕這個字很有趣.
羨慕這個字用過幾次.
這個字原來的意思就是指.
因為黃泉道不是很高大的人.
有時候有些地方 高一點的地方.
我們怎樣 伸到盡去拿.
就是這個字.
撐到最盡.
我們看人生.

$^{601}$我們伸到盡去拿的是什麼.
當然是最高一點的職位.
最多一點的回報.
最大的利益.
但是朋友 弟兄姊妹.
我們羨慕什麼.
就是反映我們以什麼為我們的家鄉.
以什麼是重要.
這班人我們看到他們很特殊.
他們羨慕什麼.
一個更美的家鄉.
有個翻譯 哪個更美.
其實原來的文字.
因為在希臘文.
它就沒有個最 沒有superlative.
所以其實那個更美.
有時候就是指最美的意思.
更美.
那個就好了.
現在這些東西比我們好.
親愛的弟兄姊妹.
什麼是我們羨慕呢.
就是決定了我們究竟.
是不是看我們這個客裡的人生.
客裡的人生對我們有什麼意義呢.
我相信其中一樣東西.
我們嚮往的一樣東西.
影響我們兩方面.
第一 歸屬感.
有些人對於自己的國家.
很有歸屬感 是吧.
我在多倫多生活.
因為它是一個多元文化.
大概有七八十個country的人在那裡移居.
世界盃 什麼人射入個球.
都有不同的人起來cheer的.
所以一定有人很開心的.
一定有那個國家的人會中的.
所以有時候在街上.
那天下午世界盃決賽.

$^{641}$我就不看了 是吧.
開完會才回去看.
都沒有用的.
因為開車回家的時候.
哎呀 糟了 那些意大利旗到處飄揚.
知道誰拿了.
因為那個環境是很歸屬這件事.
它不掉下來的.
弟兄姊妹.
會不會我們慢慢在地上生活.
到那個地步就是覺得.
我忘記了天國.
我忘記了我要去面對永恆.
面對主耶穌.
所以我奉勸我們當中的弟兄姊妹.
怎樣可以幫我們人生看.
下一代的人生.
更加積極 我認為.
我們歸屬於哪裡.
第二 變成什麼呢.
優先次序的決定.
你怎樣認為什麼是重要呢.
就表示你以哪裡是你效忠.
我記得那時候金融風暴.
2000年2008年亞洲.
有些國家很有趣的.
國家不夠外匯 不夠錢.
他們排隊去那個國家.
去自己國家銀行那裡給錢.
其中一個我記得就是韓國.
那些人都艱難的.
但是韓國人他們會.
效忠他的優先次序.
將他的錢都給國家覺得有需要.
弟兄姊妹.
我們有時候沒時間參加教會聚會.
其實是歸屬感和優先次序.
有時候我們覺得沒時間靈修.
其實都是歸屬感和優先次序.
我是主耶穌的弟女 天父的兒女.

$^{681}$我當然要認識祂.
我應該要有個優先次序.
甚至我們退休.
我們不是打算什麼都不做.
其實我們仍然都有個優先次序.
去事奉 去將我們的時間.
放在有意義的東西上.
同一道理.
你認同這個是你的天國的話.
主耶穌的吩咐.
我們要見證祂的傳揚福音.
如果我們歸屬感強的話.
我們就覺得自己是天國的子民.
我去到哪裡.
我就要講出那個意義出來.
並且我會優先次序安排時間.
所以鼓勵大家.
你嚮往這個天家.
你對這個國有歸屬感.
優先次序.
所以我勉勵大家.
最後我們要去思想一些實際的東西.
好了.
當我們這樣做.
神又給我們一個回應.
聖經說神不以我們為恥.
我覺得很奇怪.
神以我們為榮正面一點.
不以我們為恥.
這麼負面的說.
因為人本來是離開神.
很難會做到這件事.
而竟然我們有這種改變.
所以神真的以我們.
看到我們不以我們為恥.
真的完全是一種180度的改變.
好.
最後這個關於這個客女人生.
在她的生命裡面.
我覺得最寶貴的是什麼呢.

$^{721}$她是有選擇權.
選擇回去的機會.
而她選擇什麼呢.
仍然選擇神.
在舊約裡面.
路德和俄爾巴兩個都是做媳婦的.
當她的婆婆拿俄米要回到原居地的時候.
他們有一個選擇.
他們選擇可以留在原居地.
又或者回到他們夫家的鄉下.
就是巴黎行.
在那個過程出現什麼呢.
就是她的愛.
她的感情.
她的效忠在哪裡.
我們跟從主耶穌基督.
或者我們真的過一個客女人生.
離不開我們對這位神的那種愛.
那種的揀選.
我記得現在算起來是幾十年前.
有一對夫婦.
就是一對男女朋友.
就拍拖了一段時間.
他們都是想著前面會一起.
會組織家庭等等.
有一天那個女的就跟那個男朋友說.
我可能會在短時間之內.
我要過去美國.
當時在香港.
我要過去美國讀書了.
這個簡直是一個炸彈.
為什麼呢.
他們只是想著一起有自己的家庭.
有他們的一個夢想.
讀書而已.
讀書就是異地之情.
連繫一段時間之後.
再看看方向怎麼樣.
不是的 弟兄姊妹.
在幾十年前還沒有那麼多科技的時候.

$^{761}$其實一個飛去外國的人.
基本上是沒有太多機會連繫的.
所以大家都很現實去想這件事是什麼呢.
就是真的要分手了.
在那個時候.
攤牌那樣.
真的要想清楚了.
是不是這段感情是會繼續.
又或者各自去追尋.
自己另外一個夢想計劃呢.
就不知道要怎麼抉擇.
希伯來書作者可能也是.
一個人的人生就是.
我是選神.
還是我真的以這裡為我永遠的家鄉.
還是這個只是一個暫時.
我朝向永恆的生活呢.
你會怎麼選呢.
最後過了一段時間.
他們又祈禱去談的時候.
女兒最後就說.
我決定不過去 過去讀書了.
我決定留在這個地方繼續.
而男兒就答應他.
你放心吧.
什麼時候有機會我們去外國呢.
我就會帶你去外國.
去外國生活 好不好 好.
這個故事最後就很開心.
幾年之後他們就結婚.
有他們的家庭 兒女.
接著他們就在94年的時候.
就移民到多倫多.
那個就是我和我太太了.
最奇妙的一件事就是.
我心裡面說.
我主啊 我還是沒辦法帶我太太來美國生活.
誰不知到了某個年紀的時候.
因為來芝加哥事奉.
2019年我就帶我太太來事奉.

$^{801}$但是我也告訴自己.
主啊 我明白了.
我不以這個地上的地方為我的家.
不過神你要帶我去哪裡.
我就去哪裡.
所以我發覺我一生裡面去過很多地方.
去澳洲 去加拿大 英國等等.
多美的地方我不會留戀.
因為我們有天上永遠的家鄉.
最重要的是在這個過程裡面.
弟兄姊妹.
你隨著神的帶領.
你揀選最寶貴 最重要的東西.
我提醒大家.
問一下大家.
你甘心情願地去屬於什麼地方呢.
是天國永恆的家鄉.
還是你覺得手拿著的.
你所擁有的一切 地上的東西呢.
你怎樣去實踐你這個決定呢.
我挑戰你.
今年耶穌愛之城早些舉行.
8月21日到24日.
想一想它.
第二天 星期四還在UV那裡做場地.
第二個挑戰大家.
3月16日早上10點半.
芝加哥斷宣中心會借用.
CCUC的地方有個聯合祈禱會.
是為了整個芝加哥的.
科研工作圖告.
你是不是真的能夠隨時上路的一個旅客.
你是不是一個真正的旅客.
客旅人生的一個基督徒呢.
有人這樣說.
The biggest assumption we make in life.
is that we will have tomorrow.
人生命中最大的假設就是.
我們一定有明天 是不是呢.
這個威爾遜醫生呢.

$^{841}$他是一個很夢幻的故事.
他不夠30歲就在做他夢想的醫生的工作.
他是一個麻醉師醫生.
他已經有他第一個兒子 Lincoln.
他覺得有一次他自己野餐.
都說自己過著夢幻的生活.
很enjoy.
一年多之前.
他的太太懷孕生第二個小孩.
當他生第二個小孩的時候.
太太的嬰兒的位置有問題.
臨時要emergency C-section.
在那個過程中很複雜.
這個醫生因為他是那個麻醉科醫生.
所以他在幾份雜誌裡面.
很詳細地講到那種複雜的complication.
竟然是幾種情況同時發生.
最後他太太在很嚴重的病中.
生了他第二個兒子.
生了這個兒子 Gabriel.
然後就是要等待治療.
在這個時候就更加出現感染.
和不同organ的衰竭.
最後他是離開了世界.
他是在聖誕節前的幾天.
12月22號離開世界.
最後右手邊的lockbook上.
希望你看到.
就是這個爸爸在第一個聖誕節.
就是三天之後抱著大兒子.
抱著他小兒子.
小兒子望著他爸爸.
大兒子又不知道做什麼事.
望著相機.
應該是拿相機的.
是他的爺爺 阿媽.
他們正在過他們第一個.
沒有母親的聖誕節.
這位醫生說回他.
原來我假設我的明天是擁有一切.

$^{881}$或者我掌控這一切.
這個是假設.
弟兄姊妹 什麼叫做客裡人生.
這個世界的東西我們拿不走的.
為什麼希臘的亞歷山大大帝的時候.
他32歲病到要死的時候.
他就做了一個很諷刺的行動.
他雖然擁有歐亞非三大洲的領域.
但是他就是32歲.
有人說他那天不知道是喝飽了.
就有事.
有什麼病.
於是急病死.
他在棺材伸出兩個洞.
他做了一個諷刺的動作.
他以為自己有歐亞非三洲.
其實他只是掌控了nothing.
他不知道什麼叫客裡人生.
他想著在這裡建立他的一切.
但是失去.
弟兄姊妹.
我們將會怎樣過我們的人生呢.
我很盼望我們能夠有一個體見.
主啊 我雖然在這個地上用世物.
但我不是被這些東西控制我.
我在這個世界生活.
我最終目的是為了永恆的國度.
歡喜快樂 勇敢 積極的去面對.
我盼望我們互相勉勵好不好.
誠惠的弟兄姊妹 或者來賓朋友.
客裡的人生才是真正的人生.
因為我們有一個更美 更好的家鄉.
我們不是無奈 記住.
有這個很積極的天國.
我們要奮鬥得著.
我們一起祈禱.
原恩主耶穌基督激勵我們.
幫助我們.
主啊 我們真的體見.
當很多世上的人.

$^{921}$圍到地上這麼短暫的形形形形.
擺上一切 犧牲一切.
但是可以一剎那 幾天.
很短時間 歸於無憂.
恩主啊 我們實在空空的來.
赤裸裸的來到地上.
我們也同樣赤裸裸的來歸活.
主啊 願意我們的生命豐盛.
我們認定我們是客裡的人生.
我們來這樣跟從神.
不以這個地上短暫的.
牢籠了我們這個更積極.
永恆的天家的盼望.
恩代我們每一個.
聽我們禱告.
奉靠主耶穌基督的聖名.
阿們.
\newpage



\section{}
\label{sec:01fKINSJbMU}
\textbf{The Last Days According to Jesus 1 (Matthew 24) - Rev. Mark Chan}
\newline
\newline
連結: \href{https://youtube.com/watch?v=01fKINSJbMU}{\texttt{https://youtube.com/watch?v=01fKINSJbMU}} ~~~~ 語音日期: 2024-03-09
\newline
\newline
\hyperref[sec:LtoT2yWzuEc]{\small{< < < PREV SERMON < < <}}
~
\hyperref[sec:index_chronic]{\small{[返順時目]}}
~
\hyperref[sec:index_scriptual]{\small{[返順卷目]}}
~
\hyperref[sec:o3vdukdwyRI]{\small{> > > NEXT SERMON > > >}}
\newline
\newline
$^{1}$- Welcome everyone, and then Mark will take us further..
- Alright, thank you Kelly and Pastor J..
Welcome, good to see you..
And we're gonna pray for you at the end..
Make sure we pray for Pastor J until he,.
I guess, you know, settle on a.
stable ministry in Korea, alright?.
So we're gonna pray for you and pray for your family..
Alright, so why don't we take some time.
and say hi to your neighbor..
You know, just wave your hands, you know, front and back..
You know..
We have been talking about the book of Revelation,.
the last book of the Bible, about end times,.
about the second coming of Jesus Christ..
And before we go back to the book of Revelation,.
I think it would be a great idea to digress a little bit,.
you know, to study something outside of the book.
of Revelation to get some idea, some basic concept.
about what the end time is gonna be like..
So today, we're going to, not going to go to the book.
of Revelation today, but we're going to read.
and study Matthew 24..
It is a famous passage, a famous well-known sermon.
given by Jesus called Oliver Discourse, Oliver's Sermon..
He gave this sermon on the top of Mount Olives.
facing Jerusalem..
So we're gonna study a little bit, but as I said before,.
as I told you before, remember when you read the Bible,.
whatever passage you're reading, you need to understand.
the context, right?.
What goes before and what goes after..
Matthew 24, if you go back a few chapter before Matthew 24.
to Matthew 21, it is the triumphant entry of Jesus Christ..
He rise the donkey and enter the capital city.
of the Jewish nation..
That was the Palm Sunday, so-called,.
the triumphant entry into Jerusalem..
You know, Palm Sunday is coming up, right, next month,.
right, this Sunday before Easter Sunday is Palm Sunday..

$^{41}$And then after Jesus enters Jerusalem,.
he cleansed the temple, right?.
It was a pretty, you know, angry Jesus, right?.
Not the typical gentle, soft-spoken Jesus.
that the culture has us to believe..
But he cleansed the Jewish temple,.
and because of the corruption there,.
he turned the table and all that..
And then he enter, Jesus enters into a series.
of controversies, a series of arguments,.
very heated debate with the religious authorities.
in Jerusalem, with the Pharisees and with the scribes.
and, you know, St. Joseph and all those people..
So it's a very heated debate if you read Matthew 22 and 23..
And then come to our passage today,.
Matthew 24, all of the discourse..
So I want to remind you by going back.
a few chapters before 24,.
this is the last week of Jesus' life on earth,.
the last week..
Remember Palm Sunday, and then cleansing the temple,.
couple of days debating and having controversies.
with the religious authorities in Jerusalem,.
and then come to Matthew 24, the all of the discourse,.
and Jesus is talking about his second coming..
This is the last week of Jesus' life on earth..
So let's look at this sort of a last hour sermon.
given by Jesus to his disciples..
That is the context here..
That is whatever being recorded in these chapters.
and following happened in the very last week,.
the last few days of Jesus' life on earth..
And then he was crucified on the cross,.
died for our sins, and then on the third day,.
the next Sunday, he rose from the dead..
All right?.
24, one..
Jesus left the temple and was walking away.
when his disciples came up to him.
to call his attention to his buildings..

$^{81}$Now this picture is taken on top of Mount of Olives..
We are looking westward, as you can see,.
the old city of Jerusalem, and then beyond that,.
the high rise and the modern building,.
the modern Jerusalem, West Jerusalem, so-called..
So this verse says,.
when they left the temple,.
the disciples were astonished by the beautiful building,.
the magnificent building of the Jerusalem temple.
in Jesus' time..
Of course, you cannot see it right now,.
but in this picture, you can see there's a rectangular.
platform, stone-built platform, it's still standing there..
And this is what we call the Temple Mount..
The Temple Mount equals to 25 American football fields,.
man-made, 2,000 years ago..
So the foundation stone is still standing there,.
19 soccer fields..
Can you see the corner on the left?.
All the way extend beyond this picture..
25 football fields..
I mean, it was a magnificent building..
If you visit Israel with me,.
I'll bring you to the Western Wall Plaza,.
all right, as your best tour guide..
I'll bring you to the so-called Waiting Wall..
Now we usually call Western Wall Plaza..
As you enter the plaza, it's an open-air synagogue..
From the floor of the plaza,.
you count seven rows of stones..
Seven rows of stone belong to 1st century,.
in times of Jesus..
And above that, the blue arrow, later Middle Ages.
and beyond..
So there are at least seven levels of stones.
are Herodian stones..
This is kind of like a picture frame being dug out.
from natural stone, very beautiful..
1st century still standing there after 2,000 years..
In fact, there is a Western Wall tunnel.

$^{121}$beneath the Western Wall Plaza..
You go down about 100 feet below the plaza,.
you can walk alongside the 1st century Western Wall,.
about 500 meter long..
1600 feet, five city blocks..
The biggest stone we find is 40 feet long.
and 13 feet wide and weight about 550 metric tons..
That one stone in the yellow box there,.
I hope you can see it, there's the edge,.
and all the way to the lady taking picture there..
One stone, one stone..
500 to 600 metric tons..
That is one stone..
A male African elephant weighs about seven metric tons..
So that one stone equals 78 elephants..
How did the ancient people do that?.
How would they move the stone right there, right?.
It's a remarkable engineering work..
That's why the disciples, they came from,.
if you remember, from the village of Upper Galilee,.
from simple villages..
They came to the urban Jerusalem.
and they saw the Herodian Temple and they were astonished..
Of course, Jesus had different reaction..
Jesus said, "Do you see all these things,.
"the temple buildings and the retaining walls.
"and all the building and all the structure?.
"I tell you the truth..
"Not one stone here will be left on another..
"Everyone will be thrown down.".
Again, this picture is taken from the top of Mount Olives.
to our west, toward the old city of Jerusalem..
The huge rectangular platform you can still see today,.
the Temple Mount, is an extraordinary claim that Jesus made..
Jesus is talking about the destruction of Jerusalem.
and his temple in AD 70..
About 40 years later, it did happen.
because the Roman armies came to the capital city.
and they stormed the city and burned the city.
and teared down almost every single one of the stones..

$^{161}$Now, Matthew 24 is, with this background,.
is one of the passages that actually support.
the historical reliability of the words of Jesus..
He spoke those words, Matthew 24, 40 years before AD 70,.
40 years before the historic destruction of Jerusalem..
So, Matthew recorded the words of Jesus.
spoken before AD 70..
It was a prophecy fulfilled..
Now, some people will say, "This is not true.
"because it can be prophecy written after the event,.
"after the fact, after AD 70,.
"and the disciples, because they want to exalt Jesus.
"or they want to protect or defend Jesus,.
"so they made up a story.".
But this is not likely..
Let me tell you why..
This is historically not likely..
Why?.
Because if Matthew 24 is prophecy written after the fact,.
then Matthew 24, as we read the passages,.
should be more precise..
If it is prophecy written after the fact, after AD 70,.
it should be more precise because,.
as you can see, many of the stones.
were still left standing on others, right?.
If it's prophecy written after the fact,.
then if the disciples are trying to make up a story,.
then they probably won't put some of the words.
in the mouth of Jesus, right?.
If it's prophecy after the event, trying to prove Jesus,.
they would have mentioned that the temple of Jerusalem.
was burned by fire, destroyed by fire,.
but none of those were in Matthew 24..
So, as it stands, Matthew 24 is very likely historical..
Words of Jesus spoken 40 years before AD 70,.
before the destruction of the temple..
That is to say, Jesus did predict.
the destruction of the temple, and it did happen..
Matthew, parallel passages of Matthew 24, Mark and Luke,.
faithfully reports the words of Jesus..

$^{201}$And it came fulfilled..
Now we're back to 24, we read verse three..
As we go on, what is recorded here is that,.
as Jesus was sitting on the Mount of Olives,.
they exited Jerusalem and they climbed..
It's about 15 minutes walk from the Golden Dome..
That marks the holy of holies.
of first and second temple, right there..
You go down the valley of Kidron,.
and then you come back to where we took the picture,.
top of Mount of Olives..
Sitting on the Mount of Olives,.
the disciples came to him privately,.
tell us, when would this happen?.
That is, the destruction of the temple..
What would be the sign of your coming?.
What would be the sign of the end of the age?.
So the disciples are very curious..
They came forward to Jesus and say,.
they ask two questions..
When will the destruction of the temple happen?.
What would be the signs?.
And what would be the signs of the end time?.
What would happen before you come back?.
Now, Matthew 24, 25 and following,.
is Jesus' answer to those two questions..
Signs of times..
What will you see before 80, 70?.
Before the end of time..
Now, as we talk about the last days, according to Jesus,.
we enter into a subject called eschatology..
Now, eschatology sounds like a big word, right?.
But it's actually very simple in terms of what it says..
It came from a Greek word, eschaton,.
meaning ending, end time..
The destination, the end..
So eschatology is the study of what the Bible says.
about end times, about the last days..
Now, in New Testament terms,.
the last days is the final chapter of human history..

$^{241}$Our history..
Our timeline is the final period..
That final period in New Testament terms.
is between the first coming of Jesus.
and the second coming of Jesus..
And the Bible has a lot to say about the second coming,.
besides the book of Revelation..
It's right here, 24, verse four and following in Matthew..
Jesus answer, "Watch out, no one deceives you..
"For many will come in my name, claiming I am the Christ,.
"and will deceive many..
"You will hear wars and rumors of wars,.
"but see to it that you are not alone..
"Such things must happen..
"But the end, the telos or the eschaton, is still to come..
"Nation will rise against nation, kingdom against kingdom..
"There will be famines, there will be earthquakes.
"in various places..
"All these are the beginning of birth pains.".
So Jesus goes on to talk about the signs of end times,.
signs of time before the temple was destroyed..
Before the end of age..
We're not going to read it word by word,.
but let me make a list here..
These are the things that Jesus talk about..
The signs of times, before the destruction.
of the Jerusalem temple, before the end of our history..
False prophets, false messiahs..
People would go on and say, "I'm the Christ,.
"I am the prophet.".
But they are false prophets..
These guys will come out and say words.
that will deceive many, many people..
Number two, wars and rumors of wars..
Civil wars, military invasions..
Three, earthquakes, disasters..
With natural disaster, earthquakes, comes famines..
Number five, persecutions of the church..
Being a believer will be in a tougher and tougher time..
Number six, the gospel will go to the end of the earth..

$^{281}$Now intuitively, intuitively, when we read those words,.
when we read the signs of times,.
we naturally think that these things are happening.
in our time, right?.
Is that true?.
Famines, earthquakes, wars, and rumors of wars..
Kind of like a checklist, right?.
Okay, earthquake, check..
Famines, check..
Wars, check, check, check, right?.
We have a signs of times checklist..
So Jesus is coming any time..
False prophets, false messiahs,.
false prophets, right?.
Check, right?.
Can you check that?.
Couple days ago, I received an email.
from one of my online students..
So she sent this email to me and said,.
one of her friends sent her some YouTube video..
It's actually one in English, one in Mandarin..
So the YouTube videos are saying.
that Russia invading Ukraine right now.
is that a fulfillment of Bible prophecy..
Is that true or not?.
And those videos claim that to be..
What is going on right now in Ukraine.
is a fulfillment of biblical prophecy..
Now, those videos are referring to a passage.
in Ezekiel 38 to 39..
38 to 39..
It's talking about the end time again.
in different context..
And talking about Israel final enemy..
And that final enemies, one of them,.
or the head of them, is called the Prince of Rosh Meshach..
Rosh Meshach, okay?.
So some of those prophecy teachers out there.
have a lot of likes on YouTube and social media..
They will interpret this phrase..

$^{321}$The final enemy of Israel is called Prince of Rosh Meshach..
Does that sound familiar?.
They will interpret Rosh as Rus,.
changing the vowels, Russia, Meshach,.
is changing the vowels, Moscow, okay?.
All right, you got Russia and Moscow..
What is so interesting about this passage.
is that it's talking about the final enemy of,.
the end time enemy of Israel..
They come from the far north of Israel..
Now, if you know the map a little bit more,.
you understand if you draw a straight line from Israel,.
according to the YouTube prophecy teachers,.
if you draw a straight line from Jerusalem northward,.
you will hit Moscow..
That depends on how you spin the sphere, right?.
How you spin the earth..
But if you know the background a little bit,.
Meshach is a place name..
Refer to Asia Minor..
Today's Turkey..
Rosh is not a place name, never a place name..
In the Bible, outside of the Bible,.
Assyrian records, Near Eastern records, all that..
Rosh, never a place name..
In fact, the term, the modern term, Russia,.
did not appear until 11th century AD, after Christ..
More than a thousand years after the first century..
So this passage, no matter how you interpret it,.
has nothing to do with Russia, all right?.
So one more on our checklist, false teachers,.
false prophecies, false messiahs, false prophets..
Exactly what Jesus said, right?.
That's why Jesus start his Olivet Discourse, his sermon,.
asking his disciples to watch out..
No one deceives you, because there are many.
false prophets out there, nowadays,.
and back in the days of Jesus as well..
Now if you want to read the Bible,.
I want you to read the Bible well, okay?.

$^{361}$I want you to read the Bible well..
Back to this list, all right?.
So forget about Russia, okay?.
Forget about Russia, it's a tragedy, all right?.
We should pray for it, we should support it,.
support the Ukrainians, you know,.
you can send money to relief funds and all that..
Nothing to do with biblical prophecies, okay?.
However, if you want to read the Bible well,.
we must begin with the first audience in mind..
So as we read this list of signs,.
remember this list of signs are words.
spoken to the first disciples, not to us, right?.
The first audience is not us, okay?.
So we need to read the word of Jesus.
according to their context, not our own context..
So don't do the newspaper theology, all right?.
Okay, you know what newspaper theology is?.
Okay, we look at the headline news,.
and then we try to match what the Bible says..
Go to the book of Revelation, go to Matthew 24,.
and try to match, okay, what event looks like this,.
and what event looks like that, right?.
No..
If you know a little bit more.
about the first century context,.
you would know how to understand.
how the disciples understood the words of Jesus..
Now back to this list of signs,.
Jesus is talking about the time before AD 70..
I was in seminary for some years..
I always wonder, you know, one seminary class,.
like New Testament studies,.
you are required to read about 3,000 to 4,000 pages.
of reading, academic reading..
You read tons of things, right?.
You read tons of things, just for one class..
When I was in seminary about 15, 20 years ago,.
I always wondered why we didn't spend much time.
on the literature outside of the New Testament, right?.

$^{401}$New Testament, of course, is a lot to read, right?.
27 books, but there are tons of things.
outside of the New Testament, extra-biblical writings,.
that should inform us about the immediate context,.
the first century context of New Testament..
If you read the Jewish Wars, chapter five and six,.
by Josephus, a Jewish historian in Jesus' time..
He was not a Christian,.
fairly trustworthy in general, a historian..
He wrote tons of things,.
and Josephus talks so much about the time before AD 70,.
before the destruction of the Jerusalem temple.
by the Romans..
He mentioned false prophets..
Many people claimed to be the prophets..
Josephus talks about wars, and rumors of wars..
Earthquakes happening in Italy, in the land of Palestine,.
in Turkey, in today's Turkey, famines..
So if you go with this list,.
you go to Josephus, Jewish War, chapter five and six,.
you actually can check off quite a bit from this list..
Jesus was referring to his generation,.
AD 30 to AD 70, 40 years..
What happened in their world?.
It was such a chaos, right?.
Of course, persecution of the church,.
Josephus did not mention much,.
but you can tell from just New Testament writings alone..
Persecution of the church..
So many of the things mentioned by Jesus.
happened in the first century..
Josephus was one of the eyewitnesses..
The first generations of Jewish Christians,.
they saw all of this..
But how about gospel to the end of the earth?.
You ask the question, no way that happened.
in the first generation, right?.
We have internet, we have IG, we have social media,.
now we can do that, right?.
Back then, no way..

$^{441}$Well, you check some other New Testament passages.
like Colossians 1, 5 to 6..
Paul wrote the gospel that has come to you, Colossians..
All over the world, this gospel is bearing fruit.
and growing..
1, 23 and other passages in Romans also talk about.
this gospel has been proclaimed.
to every creature under heaven..
So in some sense, in real sense,.
the gospel has already reached.
over, all over the known world of the apostles..
In some sense, in some real sense,.
because all over the world can mean all sorts of people..
When the Jewish Christians broke out of Jerusalem.
and spread the gospel everywhere they went,.
they bring the gospel with them..
So all sorts of people are hearing the gospel..
The Greco-Roman world, the non-Jewish world,.
they heard the gospel..
So they have seen all the signs.
and they have the reasons to expect.
the second coming of Jesus happening anytime,.
especially when they saw the destruction.
of Jerusalem temple in 1870..
Now, of course, there were many things in Matthew 24..
Some of them are fulfilled in the first generation.
and some of them are pointing beyond 1870.
to the end of the age..
We can go through more details next time,.
but in any case, I want to wrap up today's message.
by this timeline, this chart..
If you want to understand what Jesus,.
what the New Testament talks about end times,.
this is the chart..
The bottom line is our age..
It will come to the end and Jesus has come..
The first coming of Jesus happened already..
But when Jesus came, he came with power..
He brings in the future age,.
the new heaven, new earth..

$^{481}$In some sense, in a very real sense,.
because of first coming of Jesus,.
the next age, what we call heaven,.
but actually it's heaven and earth,.
new heaven and earth,.
has already been in breaking into our age..
That is the top arrow right there..
So right now, we are living in kind of in between,.
between first coming and second coming of Christ..
The future age, in some sense, fundamentally,.
spiritually, in a very real sense,.
has already in breaking into our age..
So we are in this in between time,.
we can call it already, but not yet..
Already, meaning that the future age.
is no longer just in the future..
It is now, the kingdom of God is now in our hearts,.
in our midst..
The spirit is working in our hearts and our lives..
The future age has come here,.
already been inaugurated,.
already been in breaking into our present age..
Now, most of the signs have happened.
in the first coming of Jesus, as I told you, AD 70..
But the end of our age,.
with the second coming of Christ, is not yet..
Therefore, when you think of our time,.
we're living in this time in between..
It's a already, but not yet..
However, we are definitely getting closer and closer.
to the eschaton, to the second coming of Christ..
That's why Jesus said,.
"You do not know when exactly I will come,.
"but I will come any time.".
Right?.
I will come any time..
Now, this realization should affect how we live today..
Already, meaning the future kingdom has made an in-breaking.
into your world, into your life..
And you are no longer what you used to be,.

$^{521}$because the life of the kingdom of God.
has already been a reality in your life..
However, not yet..
On this side of heaven, we're not in heaven yet..
We're not in heaven in new earth yet..
But it will happen any time..
Now, we modern people are very impatient..
Right?.
Where are we?.
If you say, in Matthew 24, go through the checklist,.
the first generation of disciples,.
they understood all the signs have been fulfilled..
In their generation, AD 30 to AD 70..
That means, in hindsight, from our perspective now,.
the first generation of disciples,.
they were living in the beginning of the end time..
How about us?.
We are living at the end of the end time..
They were expecting the second coming of Jesus..
We are also expecting the coming of Jesus..
Yeah, we say, why the delay, right?.
2,000 years, come on, man..
If you come, please come, right?.
When are you coming?.
Well, let me remind you, we have one verse.
from 2 Peter 3.8..
But do not forget this one thing, dear friends..
With the Lord, a day is like 1,000 years..
And 1,000 years are like a day..
The Lord is not slow in keeping his promise,.
as some understand slowness..
God is patient with you, not wanting anyone to perish,.
but everyone to come to repentance..
One day is like 1,000 years in God's sight..
So if you remember this chart,.
so far in God's timing, it's only two days..
Right?.
Remember what you did two days ago?.
Remember what you did yesterday?.
You may have forgotten, right?.

$^{561}$It's gone like that..
In God's timing, two days only..
However, Jesus keeps reminding us to wait patiently,.
and to wait and be alert, and be prepared, right?.
Be prepared for the coming of our Lord..
Because we need to give an account of our lives..
How do we use our time, our treasure, our talents?.
Right?.
The three Ts..
How do we use our time and treasure,.
and our resources?.
Do we use it to extend and advance the gospel?.
To extend the kingdom?.
Or do we just bury it underground, right?.
We need to give an account, and we need to wait and prepare..
Prepare for the coming of our Lord anytime, anytime soon..
When I was a really young kid, all right?.
I was in Hong Kong back then..
Sometimes my mom, in the afternoon,.
she would say to me, with some tone of warning,.
she would say to me, "Wait 'til your dad comes back.".
It's a very meaningful sentence..
Now, you could have two different reactions, right?.
As a small kid, right?.
Wait 'til your dad comes back..
The first reaction, of course, is you're afraid, right?.
You have done something wrong..
And you should repent..
And then you can wait in fear..
That's your first reaction..
The second reaction, as a small kid,.
wait 'til your dad comes back..
That you're excited for the coming of your dad,.
because you have done well..
Well, as Asian, you have done well in school, right?.
And then you actually anticipate.
the coming of your father, right?.
Either or, either you're afraid,.
or you're anticipating the coming of the Lord..
There's no time to waste, guys..

$^{601}$We are living in the end of the end times, right?.
Now, I'm not suggesting, okay, please understand,.
I'm not suggesting, okay, let's quit our day job,.
and then we get together, we worship here, okay?.
Sell off all of your properties,.
you bought a new house, sell it, all right?.
If you have any stock, you can sell it, okay?.
We're gonna gather here, and then we'll wait.
for the alien to come, and something like that..
I'm not suggesting that..
And in fact, the first century Christians,.
they were eagerly expecting the second coming.
of Jesus Christ..
They never done something like that, right?.
Only the occult, or the cults would do that, right?.
What did they do?.
First generation of Christians..
They expect, they saw the signs..
What did they do?.
They set the orientation of their life.
toward the future kingdom..
No matter what they do, whether they are farmers,.
they are slaves, they are masters, they are businessmen,.
they orient their direction toward the kingdom,.
toward the second coming of Christ..
We are living at the end of the age..
Set your orientation, set your direction..
Toward the coming kingdom, toward the coming Christ..
If you are going to school, pray and be alert..
Do your, go through your school with the wrong direction,.
with the general direction to glorify God, right?.
If you're working adults, do so to honor God.
with your three Ts, your time, your talent,.
and your treasure, to advance the kingdom of God, all right?.
So if you understand eschatology in the Bible,.
it doesn't do any sort of escape reality..
It's not escape game, all right?.
It's not like we get together and then we wait.
for Jesus to come and then, okay,.
we'll be raptured out of this world, right?.

$^{641}$It's not escape, it's to engage..
The New Testament eschatology tells you and me.
that we are future people..
We are in some sense from the future.
because our home is in the future..
Whatever characterize our future home,.
love, peace, righteousness, justice, we bring it over here..
We make it the reality..
Now this is the life transforming idea..
If we understand we are living at the end of the end time.
and Jesus is coming back anytime soon,.
we are called to change our orientation of our lives.
toward the new creation..
That is the gospel..
I hope you and I are ready,.
ready for the second coming of Christ,.
ready to give an account to our master, to our Lord, right?.
You know, give what you have, right?.
Give what you have toward the kingdom..
Be a good steward, biblical stewardship..
How you manage your time, are you wasting your time?.
How you manage your resource, are you wasting your resource.
for something unrelated to the gospel,.
unrelated to the kingdom?.
We repent and we change..
And we yearn for the coming back of Jesus Christ..
Let's all pray together..
Dear God, we give you thanks.
for you are the master of history..
You created us all..
You are the creator and you are also the redeemer..
You are the Lord of our history..
Help us understand we are living in the end of the end times.
and we are expecting your coming..
And we pray that we will anticipate your coming,.
not in fear, but in faith..
Help us to use our resources to extend your kingdom,.
to worship you, to serve you..
Use our time, our talents, our treasure to honor your name,.
to extend your gospel,.

$^{681}$to let your gospel reach to the end of the earth..
There are so many people.
haven't heard of the name of Jesus..
There are so many people who have not received the gospel..
And we're gonna send our missionaries.
and we're gonna support the movement of the gospel..
We're gonna build up your church.
so that your name will be glorified..
And we pray that you will help us to be faithful stewards.
of the things that you have given to us..
No matter we are in school, in work, in our family,.
in our marriage, we wanna do it to honor your name,.
to make your name famous..
In Christ's name we pray..
Amen, right?.
\newpage



\section{2:1-10}
\label{sec:o3vdukdwyRI}
\textbf{Know Your Identity (Ephesians 2\_1-10) - Sam Johnson}
\newline
\newline
連結: \href{https://youtube.com/watch?v=o3vdukdwyRI}{\texttt{https://youtube.com/watch?v=o3vdukdwyRI}} ~~~~ 語音日期: 2024-03-12
\newline
\newline
\hyperref[sec:01fKINSJbMU]{\small{< < < PREV SERMON < < <}}
~
\hyperref[sec:index_chronic]{\small{[返順時目]}}
~
\hyperref[sec:index_scriptual]{\small{[返順卷目]}}
~
\hyperref[sec:Ep4OFfIa5EM]{\small{> > > NEXT SERMON > > >}}
\newline
\newline
$^{1}$Good morning, UV, it's good to be here..
Excited to worship with you today.
and getting the word together..
And thank you, Jason, for leading us this morning..
We're gonna be looking in Ephesians chapter two.
this morning, Ephesians chapter two..
If you wanna look there, turn there,.
look there in your Bible app..
We're gonna be looking at the first 10 verses.
of Ephesians chapter two..
And we're gonna be talking about our identity.
and knowing our identity, these verses give us.
this powerful picture of who we are in Christ..
So we're gonna read these verses together.
and then apply them to our lives.
and specifically what they mean for us in the church.
as a group of people who share this identity in Christ..
So let's read this..
We're gonna begin in verse one of Ephesians chapter two..
And you were dead in the trespasses and sins.
in which you once walked, following the course.
of this world, following the prince of the power of the air,.
the spirit that is now at work in the sons of disobedience,.
among whom we all once lived in the passions of our flesh,.
carrying out the desires of the body and the mind,.
and were by nature children of wrath.
like the rest of mankind..
But God, being rich in mercy because of the great love.
with which he loved us, even when we were dead.
in our trespasses, made us alive together with Christ..
By grace, you have been saved and raised us up with him.
and seated us with him in the heavenly places.
in Christ Jesus, so that in the coming ages,.
he might show the immeasurable riches of his grace.
in kindness toward us in Christ Jesus..
For by grace, you have been saved through faith.
and this is not your own doing, it is the gift of God,.
not a result of works, so that no one may boast..
For we are his workmanship, created in Christ Jesus.
for good works, which God prepared beforehand.

$^{41}$that we should walk in them..
Let's pray real briefly..
God, thank you for your word and for the truth.
that it has for us..
I pray that you would speak to our hearts this morning.
and make your truth plain and clear.
and that we would be changed by it..
I pray that your truth would change our hearts this morning.
and through that change, it would change our actions..
I pray that we would go from this place different.
than when we came in, in Jesus' name, amen..
All right, so this is probably kind of a familiar passage..
Maybe you've heard it before, read it before,.
maybe you've read it many times, I hope so..
This passage is such a powerful statement.
of God's grace and of salvation, redemption.
that we have in Christ and where it comes from..
And we could spend a lot of time here.
just talking about that, about the grace of God.
and what this passage tells us about the grace of God..
But this morning, we're gonna talk specifically.
about what this passage on the grace of God.
tells us about who we are.
and specifically about who we are in the church..
And that's important because this letter.
was written to believers in churches,.
obviously in a very different time.
and a different culture than where we are right now,.
but in a lot of other ways, people just like us,.
people who were probably in large cities.
surrounded by what we might call secularism,.
it was really paganism..
They were surrounded by idol worship,.
they were surrounded by people chasing.
after destructive things in their life,.
they were surrounded by people who rejected.
the God of the Bible and sought their own happiness.
and pleasure and fulfillment in all these other things..
They were surrounded by people who disagreed with them.
and who mocked their beliefs..

$^{81}$So in a lot of ways, the experience of these believers.
was very similar to ours..
And Paul's writing to them and he's telling them things.
that are important for them to know.
and then he's telling them things.
that are important for them to live out..
And in the rest of chapter two,.
he's gonna talk about how important it is.
for them to live out unity in the church.
and how Christ has broken down these walls.
that divide us as human beings.
and made us into this new family..
He's gonna go on and talk about that..
But this is where he starts..
He doesn't start with unity and with breaking down the walls..
Before he gets to that, he starts with identity..
And he starts with identity by talking about grace..
And that's the most significant thing.
for us to understand about our identity.
is that in Christ, we are given this grace-based identity..
That's who we are..
We see it in this passage..
There's this exchange from death to life.
that Paul makes really clear..
And notice there are no qualifiers in these verses..
Every believer who comes into relationship with God.
has been given this exchange..
Every one of us has been brought.
from a state of spiritual death.
into a state of spiritual life..
If you look at the first five verses here of the chapter.
and summarize them kind of in one sentence,.
it would be you were, but God, you are..
That's the theme of what he's saying here..
He's saying this is who you were..
Then God intervened in your life..
Now this is who you are..
You have been changed..
Your identity has been changed..
And this is not an exchange of function or vocation..

$^{121}$He's not saying you guys were peasants.
and now you're farmers..
You guys were servants and now you're royals..
He's not talking about a vocational change..
And that's important for us to understand.
because in 21st century, especially in North America,.
we tend to identify ourselves by our vocation..
We think who we are is what we do..
So we think our job is our identity.
or our relationships are our identity.
or some choices that we make in our life are our identity..
That's not what we're talking about here..
It's not a change of our vocation or our function..
It's a change of our essence..
Death to life is a change of essence..
It's a change of a state of being..
It's not a change in your job..
And that's what he's talking about..
He's saying you were something before.
and God intervened and now you have.
a completely different essence,.
a completely different identity..
And this is the gospel..
This is the doctrine of regeneration..
This idea that there's this supernatural work of God.
through the Holy Spirit whereby God gives life.
to sinners who had rejected God.
and were in a state of spiritual death..
Our nature is fundamentally changed.
out of salvation by the grace of God..
That's what we're being shown here..
And we're further being shown that this exchange.
is made solely by God on the basis of his grace..
We don't do anything to deserve it,.
to earn it, or to make it happen..
That's really clear in these verses..
Now, like I said, hopefully we all.
at least have heard that before.
and thought about it a little bit..
And I'll be perfectly honest,.

$^{161}$all of us should meditate on that.
and really reflect on it a lot more,.
myself certainly included,.
because this is really the basis of our relationship with God.
is the basis of our relationship with one another..
And so even though we've heard it before,.
we're like, yeah, I get it..
For by grace are you saved, it's not of works,.
I know the verse..
Still, we need to really reflect on that,.
consider it, and let it sink into our hearts.
and begin to change who we are..
And that's, I think, what he's trying to communicate,.
why he's starting with our identity..
He's telling us our identity.
is completely based in God's grace..
It's founded on God's grace..
That's where it comes from, that's the source of it..
It's founded in God's unmerited favor towards us..
We're forgiven because of God's grace,.
we're cleansed because of God's grace..
What Jesus did on the cross,.
the sacrifice on the cross for us..
We're accepted in the family of God because of God's grace..
We're made like Christ, transformed into his image,.
sanctified as we live out the Christian life.
and as we follow Jesus,.
that process is also because of God's grace..
And that's the essence of what it means.
to be a follower of Jesus.
and what it means to be in the church..
I'll also just say this,.
it's not really where we're headed today,.
but it's just something important.
to understand about identity,.
is that the security of your identity.
is only as secure as the source of your identity..
And that's important for us to understand..
Even as believers, we can place our identity.
in other things..

$^{201}$We can place it in our vocation,.
we can place it in our relationships,.
we can place it in our achievement,.
our education, our status..
We can place it in our financial position,.
we can place it in our class,.
we can put our identity in all types of things..
And the problem with that is,.
those things are all earthly, temporal things.
that can change and pass away..
Our relationships can change,.
people that we love can move or they can age..
And if our identity is in a relationship,.
even if it's in our own family,.
that can be taken away, that can be shaken..
If our identity is in our job, we can lose our job..
If our identity is in our possessions and our wealth,.
that can all disappear..
If our identity is in our ability.
and the things that we're good at,.
that can be changed..
We may have great beauty.
or we may have great athletic prowess.
or we may have great intelligence,.
but those things can go away, right?.
Our beauty fades with time.
and so does our physical strength and vigor..
And our intelligence,.
we can still be outsmarted,.
we can still make a bad decision..
All those things that we place our identity in.
ultimately end up being insecure..
And that's why it's so important for us as believers.
to place our identity in Christ,.
because Christ is the only source of our identity.
that's not gonna change.
and it can't be taken away from us.
and won't be corrupted or won't disintegrate or diminish..
And so if we have our identity in Christ.
and we lose our job, our identity is still secure..

$^{241}$If we have our identity in Christ.
and someone mistreats us in a relationship,.
our identity is still secure..
If we have our identity in Christ and we get old,.
our identity is still secure..
And so it's really important for us as believers.
to put our identity in Christ,.
because the source of our identity.
determines the security of our identity..
And that's the only place for us to be..
But if we're placing our identity in Christ.
and understanding his grace,.
it has real significance for the way.
that we relate to one another, especially in the church..
And that's what I want us to consider today.
because I think that's really why Paul starts here.
with identity and with grace.
and then moves on towards unity..
So let's think about this..
First, in Christ we're given this equality of identity..
We're all equal in the grace that God has given us..
But not necessarily uniformity of expression,.
and that's important for us to understand..
In the first few verses here between verse one and three,.
there's a shift in speech from you to we..
In verse one, Paul is speaking to the Gentile believers.
in Ephesus and speaking to them as a Jewish person..
That's why he's saying you..
You were in this condition..
In verse three, he's being careful to make it clear.
that Jews were also in the same condition.
of sinfulness, fallenness, and need for Jesus..
As he does this, what Paul's doing.
is he's both acknowledging the distinction.
between Jew and Gentile and also erasing.
the inferiority, superiority dynamic of the distinction..
That's important for us to understand..
We're all given this equal identity of grace in Christ..
We have equal worth in the essence of our redemption,.
but we're still different in different ways..

$^{281}$We still have different expressions.
of God's work in our life,.
and that's okay..
We still have different cultures..
We have different ways of doing things..
We have different ways of expressing the gifts.
that God's given us..
We have different gifts,.
and those are all beautiful things..
The scripture's really clear about this..
First Corinthians 12, verses four to seven says this..
Now there are varieties of gifts, but the same spirit..
There are varieties of service, but the same Lord..
There are varieties of activities,.
but it's the same God who empowers them all and everyone..
To each is given the manifestation of the spirit.
for the common good..
That means even though we're equal in Christ,.
we have this shared identity,.
our God-given distinctions and diversity are preserved,.
and that's a good thing..
Our physical differences, our personality, our skills,.
our talents, our abilities, our spiritual gifts,.
even our cultural background and expressions,.
all those things are preserved and valued in the church,.
in the community of God's people,.
and that affects our relationships, right?.
When we understand, hey, I've been saved,.
I've been brought into this family by grace..
My brother or sister,.
they do some things differently than me..
They have different gifts..
They have different things that are important to them..
They have different ideas about how we should connect.
with the community and serve,.
or what the worship service should look like,.
or whatever it is, but I can value them.
because they were also brought into this community by grace,.
and they were saved with the same grace that I was,.
and so even though there's a different expression,.

$^{321}$we share that identity of grace..
There's no inferiority or superiority here,.
and that's clear throughout the New Testament..
Galatians 3, Paul talks about that..
He says, in verse 27, this is Galatians 3, 27, 28..
For as many as you were baptized into Christ.
have put on Christ, there's neither Jew nor Greek,.
there's neither slave nor free,.
there's no male nor female,.
for you are all one in Christ Jesus..
That's what he's saying..
He's not saying that once you come into the church,.
those things don't exist,.
that there's no longer Jewish people and Gentile people..
Obviously, they still existed as Jews and Gentiles.
in the first century church..
He's not saying that there were no men and women.
in the church..
Obviously, they still existed as men and women in the church,.
but he's saying there's no dynamic of inferiority.
or superiority, that each one of you has equal worth.
and value and redemption and grace through Christ,.
and that when we come into the church,.
yes, we're still different,.
and we express that in different ways,.
we serve God in different ways,.
but each one of us has the same identity.
of grace and forgiveness..
Now, obviously, we have to submit our differences.
to God's word..
Sometimes we have differences of preference.
and things like that,.
that may be unbiblical,.
some of them may be unhealthy,.
so we have to submit those things to God's word..
We have to say, okay, what are the things.
that are a part of my preferences or my background.
that maybe I need to let go of,.
maybe I need to release?.
We need to understand that all of us, as believers,.

$^{361}$are beautiful and broken..
We're redeemed, but we still have sin,.
we still have brokenness,.
and so we have to submit those things to God's word.
and lift up the beautiful God-honoring things.
in our community of the church.
and be willing to lay down the broken things.
in the community of the church..
There are things that we have to say,.
okay, this is a part of who I am,.
but it's actually not honoring to God,.
and so I'm gonna lay this down.
for my brothers and sisters in the church,.
but maybe there's something else.
that is not a part of who I am.
that's a part of my brothers and sisters,.
but there's nothing wrong with it,.
it's actually a beautiful thing, I can lift that up,.
even though it's unfamiliar to me,.
it's still something that I can celebrate.
because we have the same identity of grace,.
we have the same essence of forgiveness,.
and those things ultimately, those different expressions,.
even if they're unfamiliar to us,.
like I said, as long as they're not unbiblical,.
as long as they're not in a violation of God's word,.
those things become adornments to the gospel,.
and the New Testament's clear about that,.
that God desires a host of people.
from every nation, every background,.
with all different types of skills and talents and abilities.
and all different types of people.
to be in his family and to worship him,.
and that's really clear in the New Testament,.
and so if God wants that, then we can be okay with that too..
We can say, you know what,.
my brother and sister,.
they do it a little bit differently than me,.
have some different ideas than me,.
but they have the same redemption that I have,.

$^{401}$they've been brought into the same family of God.
that I have been,.
and so I can value and love and respect them,.
and even experience fellowship with them.
because we have this same redemption in Christ..
Our grace-filled identity's really important..
To our unity..
We see in verses four to seven.
the equality of that redemption in Christ..
That God, being rich in mercy because of the great love.
with which he loved us even when we were dead.
and our trespasses, made us alive together with Christ..
And by grace you have been saved and raised us up with him.
and seated us with him in the heavenly places.
in Christ Jesus, so that in the coming ages.
he might show the immeasurable riches of his grace.
in kindness toward us in Christ Jesus..
Everyone in the church who truly knows God.
and has relationship with him experiences the same grace..
There's no one who gets like class one salvation,.
there's no first class salvation,.
an economy class salvation..
We all get the same grace,.
we all get the same forgiveness,.
we are all accepted and embraced by Jesus in the same way..
We're all welcomed and valued into the family of God.
in the same way..
We all have equal worth and equal value before Christ..
And every person in the church matters equally..
Sometimes that's hard for us to understand.
because we say, well, there's different roles.
and there's different functions in the church..
And so it seems like one role or function.
is more important than the other..
There are different roles and functions in the church,.
that's biblical..
The Bible lays that out for us.
and it gives us those things so that the church.
can function well in pursuit of its mission.
of glorifying God, knowing him,.

$^{441}$and reaching the world, making disciples..
So that's something that we should have..
But we shouldn't misunderstand and think that.
because there are different roles and different functions,.
that that means people have different worth.
in the body of Christ..
It's not true at all..
Just because someone is a part of serving in the church.
or because they stand on the stage,.
doesn't mean they have more worth..
They don't have more worth in God's eyes,.
they don't have, like I said,.
they don't have a higher level of salvation..
All of us are given the same grace.
and all of us are valued equally before our savior..
And each one of us matter in the same way..
And everyone is essential to the body..
Again, in 1 Corinthians 12, talking about the unity.
and the working of the church,.
in verse 20 and 21 it says this,.
"As it is, there are many parts, yet one body,.
"that I cannot say to the hand, I have no need of you,.
"nor again the head to the feet, I have no need of you.".
Everyone who is a part of the church.
is essential to the body..
Even though the roles are different..
And I love that illustration that Paul uses.
of the physical body,.
because there are parts of our physical body.
that we think are more important, right?.
Like we think our hand is more important.
'cause it does so much stuff..
But there are other parts of our body.
that we may notice much less,.
but they're essential to the working of our body..
And that's this illustration in the church..
We can't look at different roles and say,.
this is worth more, this person has more value.
in God's sight, that's not true..
Every one of us is essential to the body..

$^{481}$And in verses four to seven,.
we also see as he's talking about what God has done for us,.
he's shown us the immeasurable riches of his grace,.
that that is again given to all of us..
Think about that..
In the church, all of us are given the same Holy Spirit,.
all of us have the same access to God relationally..
Now that doesn't mean we all walk with God in the same way,.
and that doesn't mean that all of us are at the same place.
in our walk with God, right?.
There is such a thing as spiritual maturity,.
and there are people who are walking with God.
in a deeper way than others..
And again, that's biblical..
But we're all given the same access..
We all have the same access to God..
We all have the same opportunities to know God..
We all have the same privileges and spiritual riches.
as children of God..
God doesn't say, well, I'm gonna give some blessings.
to this person 'cause they're really special,.
but the other believers don't get to have that..
No, we all have those same privileges as a child of God..
We all have those same spiritual riches.
that we have access to..
And we all have the same destiny in heaven..
Bible's clear about that..
Now, that doesn't mean everything is the same.
in our lives, right?.
And again, that's okay..
The Bible speaks to that..
We're gonna have different experiences in our life..
We're gonna go through different things..
Some of us are gonna grow in spiritual maturity.
at a different rate, and that's okay..
But what we need to understand is that our physical resources.
and opportunities and the things that we go through.
on this earth do not mean that we matter more.
or have more worth in God's eyes..
And that's been a problem in the church.

$^{521}$since the beginning, right?.
James talks about it, how we're not supposed.
to show partiality when someone comes into the church.
and we say, oh, this person's wealthy or influential,.
so we're gonna make sure this person.
has the first-class experience..
And he says, don't do that, that's wrong..
Because that's always been a part of who we are as humans..
We're tempted to say, well, this person,.
their physical experience, their life experience.
here on earth, they have more wealth.
or they have more influence or there's someone with power.
or there's someone, so maybe we should treat them.
like they have more worth in the church as well..
But that's not true..
And that's what we have to remember..
It's fine that we have different experiences,.
that we have different resources,.
that we're in different places in life..
That's fine..
But don't get confused and think that that makes people.
worth more in the eyes of God, because it's not true..
What we're seeing in this passage is that in Christ,.
the foundation and the ground of our equality,.
which ultimately results in unity, is one of grace..
Verses eight to 10 make that so clear..
For by grace you've been saved through faith,.
and what you're doing is the gift of God,.
not a result of work, so that no one may boast,.
for we are his workmanship, created in Christ Jesus.
for good works, which God prepared beforehand.
that we should walk in them..
None of us have room for pride or to think.
that we're better than someone else..
That's what he's saying there..
In the church, this equality of worth that we have.
before Christ comes from grace..
And that's really important for us to understand..
We live in a world that really wants equality, right?.
It's an important value in our world,.

$^{561}$and people are seeking it..
Unfortunately, people usually seek it in the wrong ways..
And one of the really backward conceptions.
that the world has of equality that just doesn't work.
is we're trying to seek after an equality of exaltation..
We're saying, who is the person that has the most,.
and let's try to raise everybody to that level, right?.
We think that would be fair..
That's not what the Bible does..
The Bible gives us a true equality,.
but it's totally different..
The Bible says, hey, everyone's equal,.
but not in exaltation..
Everyone's equal in need..
Everyone's equal in brokenness..
The most equal place in the universe.
is the ground at the foot of the cross,.
where we come before Christ, and every one of us,.
regardless of our background or our family heritage.
or our culture or our possessions, our wealth,.
regardless of all of that,.
every one of us is at the foot of the cross saying,.
I have no hope..
I'm spiritually dead..
I'm separated from God..
I'm broken..
I'm in need..
And everything else about me, my money, my background,.
whatever, can't help me, and I need a savior..
And that's the place where we find really true equality,.
because that puts us in a place.
that is really an antiseptic to our pride..
It destroys that pride that we have,.
thinking like, well, I'm better than someone else.
because I did this or I did that..
No, at the foot of the cross, we did nothing..
We earned nothing..
There's nothing for us to boast in,.
as this passage tells us..
At the foot of the cross, we can't look at someone else.

$^{601}$and think, well, there's kind of a moral hierarchy here..
I'm better than them in some way, because we're not..
Even if we're a person who knows Jesus.
and we're following him and God's at work in our life.
and there's good works coming out of that,.
because that's what verse 10 says is gonna happen,.
even if we're that person, we have to say, wait a second..
Even any good work coming out of my life.
is because of Christ..
So I don't get to look at anyone else,.
think I'm better than them..
There's no room for pride, there's no room for boasting..
Those types of things are poisonous.
to a spirit of equality..
They're poisonous to a spirit of unity..
And that's why Paul is starting with this..
He's saying, look, if you guys are gonna be unified.
in living out the mission God has for you,.
you gotta understand who you are in Christ.
and where that comes from..
It comes from grace..
And because of that, you can treat one another.
like everybody matters just as much as you do,.
because that's the truth..
And you can treat one another like the only way.
you got in the door was because of Jesus..
That's the truth..
No one got in the door because they were wealthy..
No one got in the door because they were well-connected..
No one got in the door.
because their family knew somebody..
No one got in the door because they were smart..
No one got in the door because they were gifted.
or a celebrity or they were beautiful or none of that..
In the church, we only come in.
because of Jesus and what he did for us.
and because of his grace and the forgiveness.
and acceptance that he offers us..
And that's so important for us..
And so in closing, I want us to just think about that.

$^{641}$for a minute..
So much of identity is enculturated..
It's as we grow up, it's the water we swim in.
and we're taught that, right?.
You can, I'm sure, remember instances like that.
from when you were a kid when maybe you had a question.
about why do we do this?.
And your parents said, "That's what we do, right?.
"This is who we are..
"This is why we eat what we eat..
"This is why we go where we go..
"This is why our house looks like this.".
So much about our identity and who we are.
is just taught to us as we grow up.
and then we live it out, right?.
But think about that in the church.
and think about that in light of these 10 verses..
If we start to live out this culture.
that Paul is pointing us towards here,.
where our identity is grace-based,.
where it has, it's nothing to do with our performance.
or our possessions or anything else..
It's just grace-based..
What kind of culture does that build, right?.
'Cause in our human cultures,.
we can build identities that have beautiful things.
but also very broken things in our human cultures, right?.
We see that in our human cultures..
But Paul is putting in front of us.
something that's, there's no brokenness in it..
This is the grace of God..
And if we built a culture based on this in the church,.
that would be something, that would be truly beautiful.
and would be truly life-giving to people who are in it.
and people who are outside it and would see it.
and would wanna be part of it..
And it certainly would lead to the type of unity.
that is discussed later on in this chapter..
And this is where it starts..
It starts with this identity of knowing.

$^{681}$that we are here because of God's grace..
And because of that, we can have that same grace.
and humility and love and care for everyone else.
that God has put in our church, in our local body.
as we live life together..
We can treat them with that grace.
because that's why we're all here..
And each one of us has received that from Christ..
None of us earned it, none of us got it..
We all were given it by Jesus..
And each one of us got it..
No one got it halfway..
That's who we are, and that's how we can live.
into our relationships with one another in the church..
Let's pray..
God, I thank you for the truth of your grace,.
and I thank you for the way that it levels our pride.
and it turns us towards one another with a heart of love..
And I pray that would be true in our lives, God,.
as we live out what it means to be your followers,.
your church, both in this setting where we come together,.
we gather with one another and we're facing one another,.
but also then tomorrow and through the week.
as we go out into the world and we face the world..
God, may we, in both of those places,.
may our identity be grounded in your grace..
May we see that we're only here.
because of what you did for us..
And that's true of everyone else in the room,.
and may we really love one another.
and be unified and care for one another.
because we see the grace that you've given.
to each one of us..
God, we just thank you and praise you for what you've done..
In Jesus' name, amen..
\newpage



\section{馬可福音}
\label{sec:Ep4OFfIa5EM}
\textbf{耶穌最後的一週 \#1 (馬可福音 11) - 陳崇基牧師}
\newline
\newline
連結: \href{https://youtube.com/watch?v=Ep4OFfIa5EM}{\texttt{https://youtube.com/watch?v=Ep4OFfIa5EM}} ~~~~ 語音日期: 2024-03-12
\newline
\newline
\hyperref[sec:o3vdukdwyRI]{\small{< < < PREV SERMON < < <}}
~
\hyperref[sec:index_chronic]{\small{[返順時目]}}
~
\hyperref[sec:index_scriptual]{\small{[返順卷目]}}
~
\hyperref[sec:1C8pDQHXwCU]{\small{> > > NEXT SERMON > > >}}
\newline
\newline
$^{1}$弟兄姊妹早晨.
在三月底開始就是聖週.
Holy Week.
在我們基督徒的日曆裡面.
這個禮拜是非常重要的.
因為這個禮拜是耶穌最後的一個禮拜.
我們叫它Holy Week.
或者叫它Passion Week.
Passion意思是受苦.
就是受苦的意思.
耶穌在地上的生命.
最後的一個星期.
最後的七天.
聖經裡面居然是用最多的篇幅.
最多的份量來記載的.
你想一想.
你人生裡面如果是.
說得不好聽的話.
你剩下一個禮拜.
OK.
剩下這個禮拜.
你會去哪裡?.
喝茶,吃東西.
去多倫多.
是不是?.
最後一個禮拜你應該怎麼過?.
這個就是科恩書想聚焦.
告訴我們知道.
耶穌在他生平裡面.
最後一個禮拜.
最重要的一個禮拜.
到底他做了什麼.
在他開始傳道的時候.
他不斷提醒門徒.
提醒教會做什麼.
他說.
人子一定會上耶路撒冷.
然後受辱.
受苦.
受傷.

$^{41}$受死.
然後才得榮耀.
一次,兩次,三次的.
來去吩咐門徒.
人子.
你們的主就是這樣.
我就是一個背十字架.
願意上十字架的一個生命.
所以你們跟隨我的.
都是一樣.
所以他告訴所有的弟子.
我是這樣.
你們都是這樣.
我是先背十字架.
先受苦.
先受死.
然後才得榮耀.
有受苦.
然後才有榮耀.
我的生命的形式是這樣的.
所以每一個稱耶穌為主的人.
都依然是一樣.
我們作為門徒.
我們作為基督徒.
我們期望.
不用驚喜.
我們是一定會.
經歷一些十字架的痛苦.
拜道威我們想起十字架.
通常都是想起痛的.
釘下去那一刻.
有些電影去拍.
耶穌被釘的時候.
都是痛苦.
說到痛苦.
上個星期我的兒子.
他在讀大學.
他家裡有塊刀片.
他去清理有塊刀片.
那塊刀片掉在地上.

$^{81}$剛好那塊刀片掉在地上的時候.
就是地板與地板之間的縫隙.
掉下去是向上的.
誰知道他清理完之後.
他就一腳踩下去.
痛吧.
百多磅的踩下去.
然後晚上兩點鐘.
他就打電話給我們.
我們怎麼救你.
幸好他有朋友在身邊.
還有朋友做救護車.
就載他去醫院.
他被連了三針.
昨天載他回來.
因為這個星期是春假.
我看一看他的腳.
哇,真慘.
打了四針才連了三針.
如果你說.
如果你說他要走得了之前.
要受這樣的痛苦.
很多人都覺得很心痛.
我們想起十字架.
我們都想起痛苦拼的.
但我想告訴你.
十字架不只是拼這麼簡單.
十字架是羞辱.
是羞辱.
一個赤裸裸的生命.
犧牲在上面.
來做代價.
這是一個很大的羞辱.
一個很大的羞辱.
不只是痛苦這麼簡單.
但耶穌說.
我的生命是這樣.
門徒的生命.
朝陽是這樣.
所以最後一個星期.

$^{121}$他就去表明.
表明他自己的生命.
是一個怎樣的生命.
在這個最後的星期.
聖誕週裡面.
所有科幻書裡面記載的情節.
都是告訴我們.
耶穌的死.
耶穌的受辱.
不是被動的.
不是意外發生的事情.
而是刻意的.
他刻意的來到.
是主動的上十字架.
主動的去安排.
主動的去上十字架.
成就上帝的救恩.
在最後一個星期裡面.
就呈現到.
或者表達到.
清清楚楚表達.
耶穌自己的生命.
是成為我們每一個門徒.
有的生命.
所以在這三個星期裡面.
我們會去看.
聖誕週裡面.
四幅暗書.
裡面一些的記載.
所以我們看的經文.
可能比較多一點.
但是我想給大家一個重溫.
整個星期發生什麼事情.
字都寫好了.
不用看字的.
你看左邊就行了.
左邊的圖案是什麼.
一個星期.
我們就從這個星期日開始.
或者前兩天開始.

$^{161}$首先在星期五.
他在百大尼.
就受膏於瑪利亞.
那個拿達香高.
為他安葬的事.
來到去高尚的.
然後有安息日.
然後到禮拜日.
你看到那個.
宗之主日.
那個就是我們所謂叫做.
榮進耶路撒冷.
或者Palm Sunday.
因為那些人歡迎這個王者進入耶路撒冷.
然後他經過耶路撒冷.
進入.
經過鑒衙山.
快要進入耶路撒冷的時候.
他就為這個快要被滅亡的耶路撒冷哀號.
然後他又回到百大尼.
第二天星期一.
看到嗎.
星期一.
他就就坐無花果樹.
然後就進入潔淨聖殿.
然後同樣回到百大尼.
耶穌很有趣.
他不住在耶路撒冷.
不住在Urban地的地方.
住在濕部的.
有沒有人住濕部的.
他住濕部的.
耶穌住濕部的.
在百大尼.
耶路撒冷的濕部.
然後他星期二又回到耶路撒冷.
然後和那些宗教的權貴有很多的爭論.
很多的爭論和公開的教導.
然後他回到.
耶路撒冷對面的鑒衙山.

$^{201}$去到.
禹賢聖殿被毀.
然後又回到耶路撒冷.
去到教訓人和那些宗教的權貴.
有很多的爭論.
那些權貴.
結果是想.
去到設計來殺害他.
然後到禮拜四.
就是耶穌和門徒.
去到守最後的晚餐.
你記得有個farewell discourse.
在約翰福音.
一段很長的道別的講話.
然後晚上的時候.
他就去到赫西瑪利園.
被出賣.
被捕.
然後帶到.
大祭司和猶太人的權貴面前.
他就被審查.
然後被定罪.
彼得否認主.
這裡是禮拜五.
被彼拉多.
希律王.
來審查.
最後被判釘十字架.
死十字架之後.
聖殿的萬子.
就由上至下分為兩半.
然後耶穌在十字架上死了.
被埋葬.
另一個安息日又到了.
然後禮拜日.
就是復活的主日.
然後門徒有.
耶穌從死裡復活的見證.
這個大概是我們再重新建立.
類似這樣的一個時間表.

$^{241}$這個就叫做聖週.
你不需要現在很留意.
不過我會上傳到社交媒體上.
這三個禮拜不妨.
你去看看.
默想一下.
福音書裡面.
如何記載耶穌最後最後的一個週.
最後一個禮拜.
今天我們看.
馬可福音第十一章.
就是說這個宗之主日.
那個的.
技術.
宗之主日.
Palm Sunday.
發生了什麼事.
十一章.
馬可福音.
耶穌和門徒快到耶路撒冷.
來到伯伐奇和伯大利.
在鑒欖山那裡.
耶穌打發兩個門徒對他們說.
你們往對面的村莊去.
一進去的時候.
就會看到一匹的驢.
綁在那裡.
從來沒有人騎過的.
將他解離牽來.
若有人對你說.
為什麼做這件事.
你們就跟他說.
有個密碼.
有個password.
主要用他.
便立刻.
牽回到這裡.
用完之後就會還給你.
第四節.
他們去了.

$^{281}$看見.
果然是一匹驢.
在街上.
然後他們將他解去.
然後幾個人問他.
你解這匹驢做什麼.
門徒就照耶穌的說話說.
用那個密碼.
然後那些人就.
任憑他.
牽來.
他們把驢驅.
牽到耶穌那裡.
把自己的衣服搭上去.
耶穌就騎上了.
有很多人把衣服.
鋪在路上.
有人將田間的樹枝.
撞下來鋪上.
前呼後擁的都喊著說.
和撒拿.
意思就是說.
save.
salvation.
奉主明來的.
是應當稱頌的.
那將來.
的我祖大位之國.
是應當稱頌的.
至高無上的和撒拿.
save us save us.
就是歡迎王者進入京都.
耶穌進了.
耶路撒冷進入聖殿.
看了周圍的一切.
天色已晚了.
他就和十二個門徒出城.
返回塞浦爾.
這個伯大利的地方.
這是頭一段.

$^{321}$說到這個Palm Sunday.
Right?.
我們看地圖.
大家一定要看地圖.
因為在右邊.
就是東面.
是那個峽谷的地方.
死海峽谷的地方.
耶穌就是從加里尼北面一直下來.
經過這個峽谷.
然後在右邊東面一直上山.
因為耶路撒冷在高山上面.
他一直爬上山.
一直上到一個村莊叫什麼名字.
伯大利.
Bethlehem.
然後呢.
去到伯大利.
再上一點.
就是這個Bethphage這個地方.
就是百發奇這個地方.
在那裡做了什麼.
就是安排那條路過來.
突然叫了一個Uber過來.
他說你帶一條路過來.
我預備好了.
有埋密碼.
將這條路帶到來的時候.
他就騎上.
然後跨過這個橄欖山.
看到嗎.
橄欖山山頂.
上到去的時候.
向西看或者向左邊看.
就看到聖城耶路撒冷.
他就下山.
經過紅色那條線.
吸倫谷.
去到吸倫谷再上一點.
就是進入耶路撒冷.

$^{361}$這個就是那個路線.
那個距離其實只是.
一英里半而已.
一英里半.
他在這個地方就騎上路.
借用一條路.
在這裡看到他很intent.
是想做一件事.
在經文裡面有一些交代.
譬如百發奇.
為什麼會提到Bethphage這個地方.
馬太福音又提到.
馬可福音又提到.
這個小村莊在新約裡面出現過三次.
三次都是說.
耶穌詠進耶路撒冷.
Bethphage.
但是在經外.
在新約意外.
特別是拉比的文獻.
就有時候會提到百發奇Bethphage這個地方.
你說有什麼特別.
你知道那些猶太人拉比很緊張.
就是你帶到耶路撒冷聖殿裡面奉獻的東西.
一定要在耶路撒冷的城市限制裡面做.
譬如你奉獻那些陳設餅.
你燒過餅的.
你一定要在耶路撒冷的城市限制裡面去燒那些餅才算數.
你出去了就不算數.
所以你奉獻進去聖殿的東西.
都是要在城市限制裡面做的.
包括你搓的麵粉和做的那些餅.
拉比的文獻就很緊張這件事.
他們就說.
先定耶路撒冷的界限.
耶路撒冷都要延長到什麼時候才算是耶路撒冷.
就像我們住芝加哥那樣.
有很多住Support那些人其實他們不是住芝加哥的.
他們說他們住芝加哥的.
如果你在Frankford.

$^{401}$有時候那些Support那些人.
我們有時候在區會裡面開會的.
有時候來這裡開會.
大家介紹一下.
他們住哪裡.
住Rockford.
他們說我們是芝加哥來的.
我說什麼芝加哥來的.
Rockford和芝加哥差多遠.
但他們用了大城市來做一個座標.
你知道City Limit of Chicago.
沒有去到Rockford那麼遠.
沒有去到那麼遠.
所以拉比就很緊張.
耶路撒冷的City Limit去哪裡.
所以拉比的文學裡面就說.
八法旗.
就是東邊的界限.
所以八法旗那裡做的那些餅.
就可以獻在聖殿.
是合格的.
所以耶穌就在那裡.
帶一條路來.
在八法旗.
騎路進入耶路撒冷.
大家可以看到是刻意的.
十字架的道路.
他獻上他自己是刻意的.
為罪人獻上的祭.
他就是完美的祭.
他在那裡騎上路.
當然這段聖經也提到.
舊約一個所謂預言.
或者Promise.
就是撒迦利亞書.
撒迦利亞書第九章那裡說到.
將來會來到的王者.
帶來平安.
這位王者.
是騎著路進來的.

$^{441}$謙謙和和的騎著路進來.
所以耶穌是知道這些的應許.
知道這些的習慣.
他就預備好.
將路預備好.
有密碼.
將路預備好.
騎上去.
知道舊約裡面.
一直有對王者的應許.
他就騎上去.
他就是這樣進入耶路撒冷.
是刻意的.
不是意外的事.
然後故事繼續.
進入耶路撒冷之後.
他說到了第二天.
他們從伯大利.
Right outside the city limit.
Bethany.
出來.
耶穌餓了.
他遠遠看見一棵無花果樹.
樹上有葉子.
就過去了.
在樹上他看看.
樹上可以找到什麼果子呢.
到了樹下的時候.
竟然找不到.
只有葉子.
因為不是無花果的季節.
我想問你.
大概Holy Week的時候是什麼季節.
現在是什麼季節.
春天.
雖然你不太覺得.
馬上說.
在巴勒斯坦地區是春天的時候.
所以春天的時候.
不是收無花果的時候.

$^{481}$收無花果是秋末的時候.
大概十月.
西曆十月九月的時候.
那時候才會預期.
你會預期無花果樹.
是會有無花果給你的.
是不是.
現在當然不是了.
所以因為不是無花果的季節.
所以在樹上.
你就找不到任何果子.
耶穌就對樹說.
從今以後永遠沒有人吃你的果子.
他的門徒都聽見了.
什麼事啊.
明明不是收割的時候.
從今以後就不會有人再吃你的果子了.
就是以後你都.
不用再生果子了.
然後第十五節.
他們終於來到耶路撒冷了.
又進去.
進了聖殿.
就趕出聖殿裡面做買賣的人.
和兌換銀錢的人.
夾子的桌椅等等.
也不許人拿器皿.
從聖殿經過.
他就教導眾人了.
經上不時記著說.
我的殿要稱為萬國土地的殿嗎.
你們倒洗它成為賊窩.
就是藏在一個宗教的盒子裡面.
宗教的外殼裡面.
居然你們就是一幫強盜.
你們做了很多不好的事.
這樣去斥責他們的時候.
十八節.
祭司長和文士聽見這話.
他們就想發旨去除掉耶穌.

$^{521}$但是又怕他.
因為他是非常受歡迎的.
很多人喜歡他.
十九節.
每天晚上.
耶穌和門徒就出城外.
回到伯達利.
繼續.
第二十節.
早晨又來了.
他們從那裡經過的時候.
見到那棵無花果樹.
就見到他連樹根都枯乾了.
整棵樹就枯乾了.
彼得就想起耶穌的話.
他就對他說.
拉比.
這個就是你就坐的無花果樹.
已經枯乾到根為止.
耶穌就對他說.
你們對上帝要有信心.
我實在高數.
你們何人對這座山樹.
離開此地投在海裡.
他心裡若不疑惑.
只信所說的成就.
就為他實言.
所以我告訴你們.
凡你們禱告祈求的時候.
無論是什麼.
只要信你們已經得罪了.
就為你們實言.
你們站在禱告的時候.
若想起有人得罪你.
就該饒恕他.
好讓你們在天上的符也饒恕你們的過犯.
這段經文很難搞.
是不是.
很難搞.
首先第一件事.

$^{561}$我和第一節茶經一樣.
首先你們搞不清楚這段經文的時候.
有一個方法的.
就是你把鏡頭zoom out一點.
到底前文後理整幾個東西到底發生了什麼事.
我在這裡就zoom out了.
你就發現你在做什麼.
第十一章.
說到耶穌泳進耶路撒冷的時候.
騎路進入去.
一至六十一節.
然後有三個段落.
有三個東西.
是對稱的.
首先他經過一個無結果的無花果樹.
然後就詛咒他.
然後就進聖殿去潔淨他.
其實就是審判聖殿.
不是潔淨了之後再用的意思.
是審判這個聖殿.
condemn the temple.
這個是中間那一塊.
十五節,十九節.
然後又回頭.
第二天他們經過那棵樹的時候.
發現那棵樹被詛咒了.
到根.
到枯乾了.
你看到是一個對稱.
所以如果你這樣看的時候.
那個無花果樹.
死得很慘.
其實那個無花果樹是一個比喻.
耶穌去詛咒那個無花果樹.
其實是詛咒.
聖殿.
或者聖殿所代表的宗教的權貴.
當時那種tradition.
因為你去聖殿.
才可以見到天父.

$^{601}$你要去聖殿祈禱.
才靈的.
你要去聖殿獻祭.
才有效的.
耶穌就咒詛並且去審判這個聖殿.
就好像審判無花果樹一樣.
無花果樹發生什麼事?.
真的是無花果.
是不是?.
因為他們不結果.
是不是?.
不結果.
所以你看到咒詛無花果樹.
之後.
說到這棵無花果樹.
中間夾著的.
就是結證和審判聖殿.
馬可呢.
你要知道馬可福音很出名的一件事.
馬可福音.
最出名的就是三文字.
因為他寫作的手法很多時候是ABA那樣.
就像你吃蛋餅.
中間夾著一樣東西.
最重要的那樣東西.
所以馬可福音有十幾次.
是用這樣的方法.
來寫出.
到底重點在哪裡.
重點就是這個聖殿.
要過去.
所以如果你這樣看的時候.
你就明白下面到底在說什麼.
11章22節.
耶穌回答說.
你對上帝要有信心.
什麼事?.
「我實在告訴你們無論何人對這座山說」.
這座山說.
「離開此地」.

$^{641}$「投在海裡」.
心裡面不疑惑的話.
所說的就會成就.
這裡不是在說.
寫包單.
你求什麼都可以有.
不是心想事成.
求什麼?.
這座山.
離開此地.
投在海裡.
耶路撒冷是很多山的地方.
他說這座山.
This mountain.
是說到specifically.
具體的一座山.
哪座山?.
就是石安山.
就是聖殿山.
Temple Mount.
就是聖殿坐立的地方.
就是猶太人最愛或者最崇尚的地方.
那個地方就是聖殿.
就是天和地之間的疊通的地方.
就是那一點.
知聖所.
上帝就在那裡.
他說你對這座山說.
This mountain.
離開這個地方.
投在海裡.
我們跟大家 姐妹 爺爺.
有時候春天的時候.
去以色列.
當然今年去不了.
本來五月去.
又取消了.
真的麻煩.
經常打仗.
搞到民不聊生.

$^{681}$你那裡.
搞多久.
真的很傷心.
我們去橄欖山的時候.
如果春天去.
特別是.
下雨.
洗清空氣之後.
你可以看到什麼海.
你可以看到死海.
在橄欖山山頂.
你可以看到死海.
你看到另一個方向.
另一個方向就是聖城.
另一個方向就是Dead Sea.
拉比文河.
說那些垃圾.
不要的東西.
就是丟到死海.
是嗎.
我們今天去死海.
就去浮浮玩.
其實是垃圾缸.
不要的東西就丟下去.
死海用來做什麼.
特別是以前.
不會有護膚品.
現在就說會提煉護膚品.
死海是垃圾的地方.
拉比文河說不要的東西.
垃圾的東西就丟下去.
所以這句話.
不單地理是合理的.
耶穌所做的事情.
屬靈的意義.
都是合理的.
你對這座山說.
頭在海裡.
因為這座山已經被審判.
是腐敗的.

$^{721}$是不再需要的.
所以說你們要信上帝.
不是信那座山.
Mount Zion.
是頭在海裡.
已經要過去了.
因為什麼緣故呢.
這裡就是.
耶穌所做的事情.
我們zoom out一下.
到底耶穌來做些什麼.
在聖州裡面做些什麼.
他說你要信上帝.
因為這個聖殿將要告終.
基督已經到.
真正的王者已經到.
王者才是.
神人之間的橋樑.
上帝的兒子死在十字架上.
這個救恩才能夠帶你.
去到上帝面前.
你才有真正的禱告.
你有真正的禱告.
你能夠去接受到上帝.
對你的罪德赦免.
這個群體.
是一個十字架的群體.
你就能夠彼此赦免.
就是這樣的意思.
你再不需要這個聖殿.
它多宏偉也好.
它多受人愛戴也好.
It's about to end.
我們又zoom out一下.
這個聖州到底發生了什麼事.
其實是很多的結束.
第十一章我們看到.
就是結束聖殿.
每一個猶太人都很喜歡的聖殿.
甚至希望它屹立不倒.

$^{761}$耶穌說這個是賊窩的地方.
不再需要.
The end of the temple is coming.
然後他就和那些宗教權威.
是不是有很多的爭議.
其實就是說.
The end of the leadership.
他們的領導.
他們敗壞.
要告終.
真正有聖靈內住的群體.
他們才能夠領導上帝的群體.
然後第三段.
就是結中聖城耶路撒冷.
這裡你記得.
就是那個橄欖山的廣道.
耶穌說你們看到這些石頭很美.
主後七十年的時候.
他們就會還原.
一切都會結中.
三樣東西結中.
聖殿的祭祀.
在那裡祈禱會結中.
現在你在基督裡面祈禱.
你在基督裡面去到上帝面前.
你都敬拜.
第二樣東西.
就是那些宗教的權貴.
大祭司等等.
他們會結中.
然後聖城耶路撒冷的地位.
也都會結中.
為什麼呢.
因為基督.
來到受死.
受難.
死而復活.
十四至十六章結束的時候.
馬可科.
就是說這件事.

$^{801}$基督的救恩.
來到的時候.
這些舊事就會過去.
新的事.
就會來到.
這個就是耶穌最後一個禮拜.
並且是一個改變世界的.
最後一個聖祭.
他是計劃好的.
不是意外的.
是刻意的.
是掌握所有的細節的.
你看到耶穌.
將耶穌整個的生命三十多年.
濃縮成為一個星期的時候.
你就發現.
原來祂的生活.
祂的生命.
只有一個目標.
目標不是.
祂的生命的導向.
不是亂打亂撞的.
有一個很清楚的目標.
就是要完成上帝在祂身上的心意.
完成上帝在祂身上.
要成就的.
救贖的計劃.
去救你.
和救我.
以致我們能夠.
因為祂的十字架.
我們的罪得赦免.
我們能夠來到上帝面前.
控制所有的東西.
為了要行上帝的心意.
上個禮拜我去一個.
老人家的安息禮拜.
那個老人家我認識了很多年.
一個很忠心的.
一位婆婆.

$^{841}$很愛住的一位婆婆.
我去到安息禮拜的時候.
其中一個兒子.
他就來跟我說.
牧師 這個姐妹認識了你很久.
我說 耶 耶 耶.
他就在她的袋口.
拿了兩封信封出來.
他說.
這兩封信封是給你教會的.
我說.
為什麼要這樣做.
他說.
他說很奇怪.
他兒子說.
他說好像.
好像那些人.
他們在離開這個世界之前.
快要回天家之前.
他們有很多事情.
他們都做好了.
好像他們自己知道.
就要離開一樣.
所以他們就將這些事情.
全部都安排好了.
所以他們定時定候的奉獻.
他都放好在兩個信封那裡.
好像他知道自己快要離開了.
準備好這些東西.
對這位姐妹來說.
對生命來說.
最重要的東西是什麼.
最後一個禮拜.
臨離開這個世界.
刻意地.
揀選.
去事奉上帝.
刻意地將他優先的次序.
搞清楚.
人生最後一個的目的.

$^{881}$就是去敬拜神.
終之主日.
三個禮拜之後就到了.
最後一個禮拜.
耶穌的生命.
我們很多時候去完這個Palm Sunday.
然後第二個Sunday.
我們就去復活節.
很多時候.
或者有些傳統的人.
都會這樣做的.
去了Palm Sunday.
然後下一次就是Easter Sunday.
但Good Friday他不太理會.
很奇怪.
我認識一個人.
每逢去到復活節的.
我們有Sunrise Service.
在Chicago.
我也會見到他.
打了西裝.
打了領帶.
穿好了.
他就來Show Up.
去敬拜.
但平時不會見他.
他好像定時定候來喝茶.
這個Sunrise Service就到了.
我就會在這裡.
我不知道.
就是一些宗教的儀式.
但很多時候我們由Palm Sunday.
跳到Easter Sunday.
那個禮拜發生什麼事.
耶穌最後一個禮拜.
就是告訴你.
我背十字架的心靈就是這樣的.
一次又一次.
耶穌說.
我來到.

$^{921}$是要完成上帝給我的心意.
我是受苦,受難,受凌辱.
被唾罵.
被釘死.
然後才得榮耀.
這個是我的Pattern.
你們作門徒的.
同樣都是一樣.
背起十字架來跟隨我.
寫記.
聽天父的旨意.
尋求天父在你身上的旨意.
就算是十字架.
我也肯背.
我好幾年前看了一篇靈修的作品.
我送給你.
這個靈修的作品來做結束.
這個靈修作品叫做《我的十字架》.
我希望給大家一個提醒.
到底我們跟隨這位主.
這位愛我們的主.
願意刻意的.
主動的.
來犧牲自己.
來買贖我們這位主.
我們作為基督徒.
我們應該怎樣背十字架來跟隨他.
靈修小賓這樣說.
他說我主所吩咐我的十字架.
形狀大小都是他親自預定的.
也許我的十字架.
是當我覺得有能力做大事的時候.
卻要懷著知足的心.
在簡陋的工場.
忠心的工作.
或許我的十字架.
是在我一個似乎不結果的田地裡面.
一年又一年繼續的耕耘.
卻信上帝帶領之下.
忍耐不會白費.

$^{961}$或許我的十字架.
是要用我恩慈.
仁慈.
來對待曾經得罪我的人.
用溫柔的態度與仇敵來相處.
或許我的十字架.
是我先站在對方的一邊來處理事情.
先盡力的了解對方的角度.
不妄下判斷.
或許我的十字架.
是要在那些不願意聽見耶穌名字的人中間.
去表彰基督的品性.
十字架大小種類各有不同.
但是每一個都是痛而且重的.
每一個是合乎自己理想.
但是當你背起十字架.
馴服的.
忍耐的.
歡迎上帝給你的十字架.
那時候就是主對我.
和我最近的時候.
用你的十字架來幫助你前行的拐杖.
不要用你的十字架.
做叛徒你的絆腳石.
你的十字架.
我的十字架.
都是跟隨耶穌的生命.
想一想你的十字架是什麼.
求主幫助我們.
我們同心得來祈禱.
上主我們紀念到思想到.
你是那位背十字架.
受凌辱.
被出賣.
孤單的.
你都可以負上代價.
去改變萬民.
改變我們生命的主.
主是你這個生活的方式.
你這個生命的導向.

$^{1001}$主求你再一次提醒我們.
我們作為.
跟隨耶穌基督的人.
我們都願意背上我們的十字架.
我們生命裡面種種不同的難處.
不同的挑戰.
主你與我們親近.
讓我們知道.
這個就是基督的生命.
我們今天的禱告就在你的面前.
坦然無懼.
讓我們的生命能夠得到一個.
得到一個不同的生命的導向.
求聖靈幫助我們.
背起我們的十字架.
去跟隨你.
討你喜悅.
祈禱奉祭耶穌基督的名字.
願上帝祝福各位.
下星期我們繼續去紀念聖舟.
好.
\newpage



\section{}
\label{sec:1C8pDQHXwCU}
\textbf{Family Sunday School - EP 8 (Yes, He Is Risen!)}
\newline
\newline
連結: \href{https://youtube.com/watch?v=1C8pDQHXwCU}{\texttt{https://youtube.com/watch?v=1C8pDQHXwCU}} ~~~~ 語音日期: 2024-03-13
\newline
\newline
\hyperref[sec:Ep4OFfIa5EM]{\small{< < < PREV SERMON < < <}}
~
\hyperref[sec:index_chronic]{\small{[返順時目]}}
~
\hyperref[sec:index_scriptual]{\small{[返順卷目]}}
~
\hyperref[sec:2mbeDPTg12s]{\small{> > > NEXT SERMON > > >}}
\newline
\newline
$^{1}$[music].
Hello, everyone! Welcome back to Family Sunday School!.
Today is the last episode, the eighth episode..
Last time, we talked about Good Friday,.
which is when Jesus died on the cross and was buried..
Today, we're talking about something else..
Have you noticed what Mr. Becker is wearing?.
What color did he wear last time?.
Yes, black!.
What color is he wearing this time?.
A shiny silver color!.
Huh, why?.
Last time, he wore black like a bride..
But today, he's wearing a shiny color..
Huh, why?.
Today, we're going to talk about a story..
But before the story, let's play a game..
Come on!.
Let's play Guess the Surprise!.
Let's go!.
Behind this rock, there's a different surprise..
Take out each brick, each stone, and see what you can guess..
What's the surprise?.
[heart beating].
Can you guess?.
[heart beating].
Cotton candy!.
Surprise!.
Okay, next one..
What's the surprise behind it?.
Guess the surprise!.
[heart beating].
Take out each brick..
[heart beating].
What's this?.
[laughing].
Yay! Ice cream!.
What flavor is it?.
Strawberry, chocolate, mint..
Ice cream! Surprise!.

$^{41}$Next one..
[heart beating].
Do you know what this means?.
Hmm....
Do you know?.
Popcorn! Surprise!.
Next one..
What's the surprise behind it?.
Take out each stone..
Take out each brick..
Hmm....
Can you guess?.
[heart beating].
What's this?.
[heart beating].
Puppy! Surprise!.
Okay, next one..
[heart beating].
Hmm....
Take out each stone..
What's the surprise behind it?.
Hmm....
Oh!.
I don't know who likes it..
Who likes to eat sushi?.
Surprise!.
Next one..
What's that?.
[heart beating].
What's the surprise behind it?.
Hmm....
Oh!.
Strawberry! Surprise!.
Okay, last one..
Do you know?.
What's the surprise behind it?.
Take out each stone..
What's the surprise behind it?.
Oh!.
What's that?.

$^{81}$Easter egg! Surprise!.
[HK].
Alright, Bible story time!.
[HK].
Tough!.
[HK].
\newpage



\section{創世記}
\label{sec:2mbeDPTg12s}
\textbf{聖經考古學 - 創世記,洪水,方舟  (1) -  陳崇基牧師}
\newline
\newline
連結: \href{https://youtube.com/watch?v=2mbeDPTg12s}{\texttt{https://youtube.com/watch?v=2mbeDPTg12s}} ~~~~ 語音日期: 2024-03-18
\newline
\newline
\hyperref[sec:1C8pDQHXwCU]{\small{< < < PREV SERMON < < <}}
~
\hyperref[sec:index_chronic]{\small{[返順時目]}}
~
\hyperref[sec:index_scriptual]{\small{[返順卷目]}}
~
\hyperref[sec:3W39B_RRLi0]{\small{> > > NEXT SERMON > > >}}
\newline
\newline
$^{1}$我們一向都在講這個創世記.
我們終於來到第六章了.
第六章就講到羅亞洪水的故事.
這裡不如我們暫時不去處理那些經文.
當然有些經文都要處理.
並且六至八九章這麼長的段落.
講到一個這麼有名的故事叫做洪水的故事.
我們需要去看看聖經裡面怎麼說.
之前有提過一點.
但在這裡不如處理一些問題跟《釋經》是有關的.
但不是直接《釋經》.
這個就是聖經考古學的問題.
看到洪水,方舟.
很多時候我們在傳統的教會.
或者在華人教會的文化裡面.
我們都很著迷一件事.
有時候我們比較保守的信徒.
我們就說要找一些考古學的證據.
去證明某些聖經的事跡.
所以這裡不妨探討一下.
特別著迷或者很obsessed(附和).
在傳統教會的圈子裡面.
華人都是想找方舟.
近年來都是偶爾會聽一下這些故事.
誰誰找到方舟.
誰誰可以證明有大洪水等等.
到底有沒有這回事.
或者應該怎麼去入手去想這些問題呢.
我覺得我們作為教務,作為童工.
我們起碼有一些基本需要去認識一下.
所以在這裡我們先講聖經考古學.
我是讀聖經考古學的.
我認為聖經考古學有三方面可以做到的.
三方面功能,三方面的目的.
第一方面目的就是印證.
有一些聖經的事跡是可以在考古學上.
可以找到證據,找到一些佐證.
去證明它的歷史性.
舉個很簡單的例子.
譬如有個石碑在以色列一個北面的地方叫做丹城.

$^{41}$那裡是1993年出土出來的.
是一塊石碑提到有House of David.
大衛之家,大衛王朝這一句話.
這個石碑我們不詳談.
但是一找到出來的時候.
1993年就可以告訴世界.
大衛王朝的歷史性.
它有Historicity.
因為這塊石碑不是以色列的人自己雕出來的.
也不是大衛王朝的人雕出來的.
是大衛王朝或者以色列人的敵人.
亞蘭人的一塊石碑.
提到他和大衛家House of David去打仗.
這個戰碑當然是說他贏了他們.
所以就雕了這個戰碑出來.
所以有大衛的歷史性.
所以某一些的人事物.
聖經考古是可以印證出來的.
第二方面就是修正.
修正什麼呢?.
修正我們以往對聖經的理解.
以往聖經的詮釋.
考古學也算是一種科學.
那種科學提出來的理據.
一些資料或者已經確認了一些現代科學的看法.
可以修正一下我們現在對聖經的理解.
其實這個絕對不是新的.
因為我們創世記第一章的時候已經說了這件事.
現代的科學可不可以吻合.
創世記第一章六日七日的創造次序呢?.
是不可以的.
所以可能以前教的那種字面式經.
就是覺得我們只要照著那七日的順序.
我們就說宇宙的起源的過程就會是這樣的.
這種是現代人字面式經.
當我們現在知道科學或者太空科學等等.
他們所表述出來的物質形成的過程.
宇宙的來源的歷史等等.
我們發現是很難能夠合得上這六日和七日.
事實上聖經六日的創造次序.

$^{81}$本身已經不可以科學性地去理解.
我們之前已經提過了.
很簡單,一直以來在過去這麼多年的誡經歷史裡面.
都已經有人提出.
怎麼會頭三日有晚上有早晨呢?.
都還沒有太陽,月亮,星星.
到第四日才有.
那三日是什麼呢?.
我們很多時候就按照現代字面式經.
嘗試去解通它.
可能第一日已經有太陽,月亮,星星.
只不過第四日才呈現出來.
這個在原文的文法語言裡面是不能夠通過的.
我們只能夠在第一章的時候說.
根本古代的聖經古近東的寫作.
講到創造的時候.
其實就不是回答科學的問題.
是回答比科學更重要的問題.
身份的問題.
還有誰是創造主的問題.
不是who,不是how.
也不是when.
所以有一些修正.
correction是可以幫助到我們.
考古學裡面找到的東西.
可以告訴我們.
矯正到我們.
以往我們對聖經的詮釋.
可能需要一些修改.
但當我們修改我們詮釋的時候.
我們不需要怕.
因為不會影響聖經的權威.
因為只會影響.
我本身對聖經的詮釋.
那個解法.
所以解法是可以變.
聖經的最高權威.
沒有變.
只不過現在越來越修正.
越來越correct的時候.

$^{121}$我更加預備了解.
第三方面最重要的.
這個也是.
在華人的信徒裡面.
認識.
所謂考古學的時候.
很多時候都沒有留意.
其實考古學最重要的是什麼.
就是讓我們知道聖經的背景.
不是confirm什麼.
如果我剛才說找到.
大衛的歷史城.
大衛皇朝的歷史城.
你只能夠說.
真的有一個人存在過.
但你confirm不到.
很多聖經裡面說的.
大衛的信心.
大衛的豐功偉績.
大衛的屬靈心靈.
那些根本是不可以.
在考古的東西裡面.
能夠證明出來.
是吧.
所以.
可以confirm到的東西.
只不過是歷史城.
但很多時候神學的真理.
是不能夠歷史的方法去confirm的.
對嗎.
在這裡順帶一提.
快到復活節了.
復活節的時候.
我很喜歡一本書.
是一個以色列的考古學家寫的.
我這裡沒有列出來.
叫做Final Days of Jesus.
這個以色列的考古學家是猶太人.
他是不信耶穌的.
他就將四福音.

$^{161}$當成一些.
是.
有歷史.
真實性的文件.
來去查考.
他就研究.
他在耶路撒冷找到的東西.
然後配合.
這個福音書裡面.
所講的一些細節.
等等.
他就發現原來是match的.
原來是match.
但他這本書.
在最後的conclusion.
他就說什麼呢.
他說.
耶穌的墳墓.
是空的.
他說有很多人.
以往曾經用不同的方法.
用不同的歷史方法.
來去否定他.
又或者去到.
自言其說等等.
他就認為這些是很奇怪的理論.
根本是不合理性的.
他說最終.
這個墳墓是空的.
但為什麼是空的呢.
他說.
不需要用一些.
旁門左道.
就是說可能耶穌沒有死.
等等不同的想法.
去解釋耶穌的墳墓.
是空的.
他說什麼呢.
他說除非你去到.
接受聖經的解釋.

$^{201}$聖經如何詮釋耶穌的空墳墓呢.
就是神使基督復活.
這是神學性的真理.
他就說.
這是超過了.
超過了.
這個歷史.
和考古的方法.
能夠溝通的範圍.
所以他就說.
留待讀者你自己決定.
換句話說.
這個所謂信仰神學的真理.
不是考古學可以觸及的.
所以.
最重要.
我們不是.
確定一些東西.
正確是可以有.
但最重要是context.
讓我們知道聖經的世界.
聖經背後的世界.
你記得這個.
圖表.
我們看聖經的時候有三個視野.
其中一個視野我們往往很多時候.
我們之前說過的.
都是說背後的世界.
就是說context.
找到的東西我們能夠重建.
當時那些人的生活.
或者他們的.
留下的文獻.
等等.
這是考古學的部分.
所謂找到的東西可以讓我們知道.
當時人怎麼想創造.
怎麼想歷史.
怎麼寫他們的東西等等.
這些都是聖經背後的世界.

$^{241}$所以我們說.
三個horizon.
三個視野.
the world behind the bible.
然後.
聖經裡面的世界.
然後我們.
或者我們以往歷史裡面.
一群人一起看聖經.
就是world in front of the bible.
三方面都需要.
考古學做什麼呢?就是背後.
context或者background.
在這裡先看看書.
有些書我覺得.
是很好的.
是近幾年寫的.
現在是2022年.
這些書.
叫做.
大部分的作者就是John Wharton.
John Wharton現在在Whitman College.
他教近東教.
他由.
第一本叫做Lost World Series.
就是失去的世界.
意思是說.
聖經背後的世界.
很多時候我們現代人已經失傳了.
我們已經忘記了.
藉著很多不同的方法.
譬如近東的文獻,考古等等.
這些就呈現了.
Lost World.
他由第一本開始出.
就是講Lost World of Genesis 1.
然後第二本.
就是Lost World of Scripture.
聖經文學背後的世界.
第三本就是.

$^{281}$Lost World of Adam and Eve.
就是亞當夏娃遺失了的世界.
是怎樣呢?.
第四本就是.
Lost World of Israelite Conquest.
就是講到約書亞進入.
迦南.
這本書也是很有趣的書.
也是考古裡面幫到我們.
我們經常說.
滅迦南那麼多族.
死了很多人.
什麼人都要殺.
到底怎樣去解釋.
戰技裡面的技術呢?.
就是.
每個人都死了.
近東有很多同類的戰技.
我們需要去比對.
去看.
你就知道實況未必是我們.
現代的字面.
這樣去解釋.
第五本就是Lost World of the Flood.
就是講到羅亞方舟.
這本就是兩個.
美國福音派的神學的.
學者巨人.
所寫的,就是John Wotton和.
Trevor Norman.
兩個都是很有資格的.
舊約的識經者.
然後第六本就是.
Lost World of Torah.
就是講到解這些律法.
就是五經裡面.
不同的律法,到底應該怎樣解釋呢?.
Again,這些也是.
用一些比對的研究的方法.
讓我們知道當時的背景.

$^{321}$是怎樣,然後.
告訴我們聖經直接寫給當時的.
古利史人的時候.
他們怎樣理解那個信息.
然後那個信息才可以.
用到在我們現在.
現代的社會裡面.
所以第一.
這裡我想提的就是Lost World of the Flood.
之前我們有很多課.
都是有很多的.
原則和觀察.
都是從這些.
這個.
Lost World of Genesis 1.
World of Adam and Eve等等.
那裡拿出來.
所以我們今天可以看一下.
這本Lost World of the Flood.
逐字去看.
大概上他們的立論是什麼.
還有一些什麼.
傳息的原則.
能夠讓我們學到一些東西.
我引用這些書的原因.
就是這樣.
就是在.
我現在針對的教會.
這個世界,很多時候我以為.
哇,Mark Sir,你說的東西.
很不同我們平時.
從小到大在教會.
學那些東西,傳統學那些東西.
你現在是新派了.
我想告訴你,我引用這些書的原因.
就是你要知道.
我們都要update(更新).
因為這些書.
都是誰寫的?.
都是一些很穩打穩紮的.

$^{361}$保守的.
福音派的學者寫的.
我們不知道那些.
過去我們講過的.
所謂新的解釋.
原因是因為我們沒有upgrade(更新).
其實這些.
知識.
這些發現的.
一些資料,藉著考古.
等等,其實已經.
存在了幾十年.
只不過很多時候我們做教務太忙了.
我們離開了.
神樂園太久.
我們就沒有upgrade(更新).
我們沒有.
所謂跟隨.
學界裡面.
有的新的發現.
所以我們脫節了.
我們以為以往.
我們有西都大學那套就是對的.
Again.
我們在詮釋上需要修改.
因為我們越來越知道.
背景的時候,我們越來越容易.
去掌握.
聖經裡面的意思.
這個Laws of the Flood.
我們先講.
Laws of the Flood.
如果你能夠看到英文的話.
我想這本書是.
美中文翻譯的.
如果你看得到的話.
我非常推薦.
你去看,Again是Trumper Lohmann和John Wharton.
兩個神學界.
舊約的巨人.

$^{401}$書的內容.
不是那麼難看的.
因為它是分開不同的propositions.
就是每一個chapter.
每一個章.
就有一個的命題.
所以它有十幾個的命題.
每一個命題.
我們討論一個中心的思想.
譬如第一個命題.
就是.
創世記是一個.
古代近東的文獻.
Ancient Document.
你不要當它是現代科學的報告.
這樣就不會那麼笨.
這是第一個命題.
第二個命題就是.
創世記1-11章.
宣稱是一些真實過去.
發生過的事.
真實的事件.
我們從頭到尾一直說.
如果你聽得清楚的話.
聽得明白的話.
仔細的話.
你就知道我從來沒有.
去質疑過.
有歷史的亞當夏娃.
也沒有質疑過有歷史的創造.
創造真的發生過.
It's a fact.
所以聖經1-11章.
是在宣稱一些.
Historical events.
第三個命題.
就是.
頭十一章.
運用了一些修辭手法.
Rhetorical devices.

$^{441}$這些都是文學上的修辭手法.
說到文學的時候.
傳統的教目.
或者比我更老的教目.
一聽見文學式經.
他就很怕.
他就以為文學式經和歷史是對立的.
一直學界以來.
都沒有這節過長.
這些保守的學者.
John Ward and John Berman.
都是用文學式經.
來解釋聖經.
所以你不要想著.
文學式經.
Literary approach.
和歷史性.
Historical history.
就是對立的.
兩者不是對立的.
但是我們要留意.
這些古文獻是有寫東西的手法.
是跟我們寫東西是不同的.
正如我們講到.
創世的第一章.
123456.
兩個是平衡的.
1和4 2和5 3和6.
這些古代近東的.
修辭手法.
不是我們今天寫東西的手法.
我們掌握了的時候.
就明白多一點.
經文想表達什麼.
第四方面.
聖經是用誇式法來描述歷史事件.
Hyperbole.
即是說.
一件事件.
可以將它誇大來形容.

$^{481}$講到這裡我們就很緊張.
誇大什麼?即是說謊?.
不是的.
因為這些古代文獻.
很多時候是用一個這樣的.
Rhetorical device.
Hyperbole 來表述一些東西.
等於.
耶穌在比喻裡面說什麼.
他說.
界賽種就是百種裡面最少的.
你真的見過界賽種你就知道.
不是百種裡面最少的.
你可不可以說耶穌說錯話.
當然不可以.
Hyperbole 是一種誇式的手法.
讓聽眾.
讀者.
他們能夠.
更加容易理解到這個點.
所以我們都要留意.
如何去解釋.
聖經裡面的一些資料.
一些經文.
往往很多時候.
都會記載一個歷史的事件.
是用一些誇式的手法.
告訴你那個嚴重性在哪裡.
其實.
在這裡提一提.
如果你還沒聽過的話.
有一個叫做聖經的.
教義的宣言.
一個Statement of Faith.
聖經教義的宣言.
或者聖經無誤的宣言.
是在70年代的時候.
產生的.
在哪裡產生的呢.
就是我住的這個地方.

$^{521}$芝加哥.
當時是搜集了所有.
在北美最出名的福音派的.
學者.
就來到芝加哥這個地方開會.
去討論什麼叫做聖經無誤.
在這個Statement of Faith.
講到.
聖經的教義的時候.
聖經無誤的時候.
也提過.
聖經的作者.
是會用誇式的手法.
來形容表述一件事.
所以一講到誇式.
哇.
你不要以為這個人是新派.
Marx是Liberal.
Chicago Statement of the Bible.
也是這樣說.
是不是.
所以見到這些東西我們不需要說.
哇這麼緊張.
因為確實聖經裡面有很多東西.
明明不是講到全世界的東西.
他描述的那件事.
是不一樣的.
包括洪水.
第五個命題.
聖經很適當地用一些誇式的修辭.
來表述洪水的事件.
Hyperbolic account of the biblical flood.
用一些Hyperbole.
全地.
整個世界.
這些字眼.
來形容洪水.
所以這個是適當的.
因為其他的.
例如近東.

$^{561}$講到洪水的文獻.
都是用一些這樣的手法.
這個後話.
第六個.
就是講到.
洪水是一個.
Global event.
創世記將洪水.
表達成為.
描述成為一個全球的事件.
第六個命題.
第七個.
講到背景.
就是古美索不達米亞.
也有其他洪水的故事.
這些其他洪水的故事.
都是講到World Wide Flood.
世界性的洪水的故事.
第八個命題.
就是聖經.
洪水的技術.
與古近東.
洪水技術有相似.
也有很多差異的地方.
這個第八個命題.
第九個命題.
就是一個.
局部的.
地區性的災難性的洪水.
被刻劃成.
全球的洪水.
以達致什麼呢.
就是一些修辭.
或者神學的原因.
以達致一些這樣的.
解釋一些這樣的理由.
所以雖然那個事件.
本身.
可能是局部性的.
但是將它去到刻意描述.

$^{601}$成為全球性.
其實是想做到.
一個神學的教導.
一個神學的.
元素表達出來.
這個第九個命題.
作者說.
然後第十個命題就是洪水的.
技術是罪惡審判的一部分.
這個就作為.
所謂的covenant.
盟約背景的故事.
因為在.
第一至到第十一章都是在講一個covenant.
然後我們曾經說過.
作者在.
頭十一章就用.
不同的事件.
去到講出什麼.
講出這個罪與審判.
的故事.
第十一.
的命題.
就是神學史.
的重點.
就是神的同在.
神要做什麼呢.
就是去確立那個秩序.
還有去講出.
在第十一章裡面.
創世記.
就是神所訂立的.
秩序.
是怎樣被罪所破壞.
然後第十二個.
就是我們講到神的眾子.
這個段落.
不是洪水的前奏.
這麼簡單.
而是King and Abel.

$^{641}$這個.
該隱與阿伯的.
敘事這個故事的.
一連續下去.
第四章一直下去.
去到第六章就是洪水.
然後第十三個命題.
就是最後一個故事.
第十一章巴別塔的故事.
就是這個所謂.
起源敘事.
這十一章.
的起源史.
一個最恰當的結論.
因為到第十一章的時候.
我們假意事又去到第十一章.
就會告訴我們.
這個世界是需要救贖.
這個世界被罪破壞得太厲害.
上帝呼召亞伯拉罕.
和他的子民.
這個屬靈的子民.
來拯救這個世界.
我們繼續下去.
第十四個的命題.
就是洪水故事背後.
是一個真實的故事.
真實的事件.
另外.
地質學並不支持.
有全球洪水.
這個第十五個的段落.
就是一位.
如果比較細一點.
看不到的叫做斯蒂夫·摩西.
斯蒂夫·摩西.
也是在威廷大學教地質學的.
一個基督徒的地質學家.
他就寫.
第十五段.

$^{681}$他就說Geology.
沒有任何證據.
去證明曾經有一個全球性的洪水.
第十六.
世界各地有洪水的故事.
但不能證明.
曾經有一個全球的洪水.
這些可能.
你和我.
從小到大都是這樣聽回來的.
在主要學裡面.
主要學老師就說.
什麼民族都有洪水的故事.
所以曾經有一個全球的洪水.
如果你想清楚.
邏輯上.
是不carry的.
你可以同樣說.
全球各地.
有很多洪水.
以往古代流下來的故事.
其實就證明不是全球洪水.
是證明一個很大的地區性的洪水.
一樣可以.
因為人擴散.
繼續搬去別的地方.
失傳了或者遺留了.
就說回有些洪水的故事.
所以不一定.
去證明到全球洪水這件事.
第十七.
科學能夠淨化我們的信仰.
其實這裡說得清楚.
能夠修正.
修正和淨化我們對聖經的傳承.
信仰.
反過來.
就能夠去purify.
能夠淨化.
科學.

$^{721}$現代科學那種偶像.
或者虛假的絕對.
以為科學就是陳述整個所有真理.
的唯一方法.
但信仰.
神學就是告訴我們.
這個神學.
告訴我們.
不一定是.
全部真理.
不一定是真理的全部.
科學不是真理的全部.
只是一部份.
另外一部份就是.
信仰,宗教.
聖經能夠給予我們.
所以有大概17個.
的命題.
所以這裡每一章.
都是這樣.
組織出來.
如果你喜歡我們的影片.
就記得要讚好,留言和分享吧!.
謝謝啦.
(音樂播放).
\newpage



\section{}
\label{sec:3W39B_RRLi0}
\textbf{Apostles' Creed \#6 (Death \ and  Resurrection) - Rev. Mark Chan}
\newline
\newline
連結: \href{https://youtube.com/watch?v=3W39B-RRLi0}{\texttt{https://youtube.com/watch?v=3W39B-RRLi0}} ~~~~ 語音日期: 2024-03-19
\newline
\newline
\hyperref[sec:2mbeDPTg12s]{\small{< < < PREV SERMON < < <}}
~
\hyperref[sec:index_chronic]{\small{[返順時目]}}
~
\hyperref[sec:index_scriptual]{\small{[返順卷目]}}
~
\hyperref[sec:EpFmOvtCyQg]{\small{> > > NEXT SERMON > > >}}
\newline
\newline
$^{1}$Thank you, James, for giving us a short announcement..
Today, we're going to continue with our series,.
and there are some questions I want you to do together..
There will be some riddle to solve, okay?.
All right, so it's coming up, all right?.
So, at the Apostles' Creek,.
we are talking about this ancient statement of faith..
It's been around for more than 1,500 years,.
almost 2,000 years, to give us a world map.
of what it means to follow Jesus Christ..
We are coming to, I don't know,.
statement number one, two, three, four,.
maybe four or five..
It talks about Jesus..
Remember, the center of the entire Apostles' Creek.
is on Jesus Christ himself,.
what he has done and who he is..
So, right here, we have a statement that says.
he was crucified, he died,.
and was buried, and he descended to the dead..
That means he really died, right?.
Was crucified..
All right, that is a technical phrase in the first century..
Now, who is being crucified in the first century?.
Political criminals..
They will be crucified..
In fact, Jesus was not the only one.
being crucified in the first century..
I think, as I prepared this message,.
I remember it was back in the 1970s..
Yeah, '70s..
Before you were born, right?.
Wait, I was born, okay..
I was born..
There in Jerusalem, they dug out a tomb,.
and in the tomb, they found some human bones..
One of the bone pieces is.
a heel bone of a man, all right?.
Heel bone of a man..
You can see on the left..

$^{41}$There is still a Roman nail going through it..
We guess that maybe in,.
somebody being crucified,.
that man being crucified, the victim of crucifixion,.
and he was buried, and then as they try to bury him,.
because there was a big nail going through his heel bone..
You know where the heel bone is?.
Yeah?.
Going through it, professionally,.
going through the entire heel bone without breaking it..
And then when they buried the guy,.
they could not take the nail out,.
probably because the nail was bent or something like that..
So as you can see from this heel bone,.
you can reconstruct kind of what it would look like.
for a crucifix victim..
So,.
both the feet being nailed to the vertical wood beam,.
and then.
hands and arms extended..
Usually, the nail would go through your wrist..
So this reconstruction does not show it..
However, sometimes, a crucifix victim.
would be nailed to the cross by four nails,.
just like what the gospel have described..
All right?.
The nailed hand, right?.
Both hands..
And then, nailed feet, right?.
So this is something historical..
We are looking at what's crucified, this face..
Only the political criminals,.
political rebels, they were being hunt down.
by the Roman authorities, the Roman Empire police..
And then they would crucify these people..
Now, as mentioned, crucifixion is a political penalty.
for all those who have committed treason.
against the Roman Empire..
Now, with this picture, we can see that, actually,.
when Jesus died on the cross,.

$^{81}$he let us see.
there was a political side to the kingdom of God..
By dying on the cross, by suffering from crucifixion,.
he's actually showing us there was a different reality,.
different kingdom, different political reality..
And we, as Christ followers, we have to see that..
Just happened today is the happy independence day,.
well, tomorrow, right?.
The happy July 4th weekend..
Many Americans are celebrating nationalism, right?.
Patriotism..
There's nothing wrong with loving our country,.
but that can never trump our allegiance.
to the kingdom of God..
Because followers of Jesus Christ,.
we are proclaiming a totally different reality..
And Jesus, by dying on the cross,.
he proclaims a totally different world order,.
a totally different order of power,.
reconfiguring our whole life..
When Jesus came to your life, when Jesus came to my life,.
there should be some reconfiguration of powers,.
reconfiguration of functions in our lives..
Jesus died on the cross, he sacrificed himself,.
he laid down his own right..
That is a different reality, isn't it?.
Right, crucifixion, Jesus died on the cross,.
it's other center, he died on the cross.
so that we can be saved..
He did not grasp his own right in heaven,.
but he laid down his own right..
That's why the Bible consistently emphasizes one fact,.
is that Jesus is the servant king..
And that is the status for all church communities..
We follow Jesus, we profess his faith,.
we believe in him, as the apostle of Crete says,.
we believe that he was crucified, he died..
And that is the example of Jesus..
Let me read a passage from Philippians 2..
It says, "Your attitude, your mindset,.

$^{121}$"should be the same as that of Christ Jesus.".
And then he goes on, "Jesus, being in the nature of God,.
"very nature of God, did not consider equality with God.
"something to be grasped, but made himself nothing.".
Those of us who are versed in English grammar,.
you should know something called reflexive pronoun..
You know what reflexive pronoun is?.
No?.
Those who are teachers, students,.
reflexive pronoun means something that you do to yourself..
Right?.
Something that you do to yourself..
That is reflexive pronoun..
Can you find a reflexive pronoun here?.
Made himself nothing..
Nobody forced him to do that, but he chose to do that..
He made himself nothing..
This is the example of Jesus..
And we have to ask one more question..
Why did Paul write 2-5, right, Philippians?.
There was something before 2-5, which is 2-1-4, right?.
Actually, Paul is actually teaching the Christians.
in Philippi, the small group churches,.
maybe five, 10 people a group..
10, 15 people is called a church community.
in the city of Philippi..
And he is trying to teach them.
what it means to live as a church..
What it means to be the church..
And church is a community where we lay aside.
our own interests..
We come into the church, we gather with our,.
with the Bible study group or college Bible study group.
or high school group, or coming to worship..
No matter where we gather,.
if we gather in the name of Christ,.
if we gather in the name of Jesus,.
first presupposition is to lay down your own rights,.
your own priority, your own interest..
Working for the interest of others..

$^{161}$That is the community of church..
That is what Christ has called us to be..
Now many people today, when we come to church,.
sometimes we forgot to check our attitude..
We come here to be entertained, right?.
Have you ever noticed the worship team.
played some wrong notes?.
Probably they did, right?.
You walk out the church and would you give five star.
to the worship experience?.
Probably not, right?.
Probably not..
But if we do that, if we do that,.
we miss the whole point..
We are not experiencing the presence of God..
We are not experiencing the presence of Jesus..
We come here not as a customer..
We come here as a community of Jesus Christ himself..
And what did he do?.
He laid down his own rights..
He made himself nothing..
That is the church..
That is what the church is all about..
One more example..
We can quote other passages from the Bible.
talking about the church being a different order.
from the world, totally different from our culture.
because the culture, there is a culture script out there..
The American culture, happy July 4th,.
we celebrate individualism..
That is the meta-narrative of our culture, of our world..
And it's constantly brainwashing us to tell you.
that if you follow this meta-narrative,.
this big narrative, this cultural script,.
you'll be great..
You define yourself by individualistic choices, right?.
I make my own choices..
I want to define myself the way I want to..
That is the meta-narrative out there..
But when we come before Christ,.

$^{201}$it's a totally different ballgame, isn't it?.
We lay down our own rights..
It's no longer individualism..
It is other-centered community..
We lay down our own rights..
We lay down our own status..
I think of some other passages from the Bible.
that also proclaims this kind of different order of power..
Remember that there were some passages in the Bible,.
in the New Testament..
It sounds funny to us as modern readers..
I'll give you an example..
One of the example is 1 Timothy 2..
Paul there is teaching the women in the church.
not to dress up to go to church, but to dress down..
Have you ever noticed that?.
No jewelry, Paul is saying,.
"Don't come here with nice clothing..
"I see that you are not coming with very nice clothing..
"That's good.".
All right?.
And why did Paul do that?.
Why did Paul do that?.
You think your clothing is nice?.
Well, anyways..
(audience laughing).
Well, you know, you picked that this morning?.
Cool..
(audience laughing).
Anyways, 1 Timothy 2..
If you have time, go back home and read that passage..
Paul is actually doing something very funny to us..
He said, "Don't dress up when you come to church..
"You dress down.".
Now the question is why?.
Because in the Roman Empire,.
the social order so often is defined by.
how you wear, right?.
Now it's funny to us because we no longer do that, right?.
We no longer emphasize so much on this, right?.

$^{241}$I'm a CEO of a big company..
They can give a presentation,.
a formal presentation with jeans, right?.
I mean, we don't do that anymore, do we?.
We don't..
I remember in the old days, in the '80s, in the '90s,.
when we go to church,.
sometimes we are required to put on suit and tie..
Well, unspoken rule..
I remember one time we were still in college..
We were going with some friends to Los Angeles..
So we were driving down from Northern California..
I was in Sacramento at that point..
And then we drove down to Los Angeles for vacation..
So some of our friends, we went to Hollywood,.
we went to.
Universal Studio and all those stuff, right?.
The fun stuff..
So on Sunday, we found a church..
See, we have good Christians, right?.
We still find a church to go to..
So Sunday morning, we drove to one of the churches.
in the area..
So we walk into the church, right?.
So we were young and fun people..
I remember one of my friends was wearing shorts..
As a group, we walk into the church..
And then we were stopped by a deacon..
And the deacon stopped us,.
but not stopped me because I was wearing pants, right?.
Long pants..
Stopped the guy wearing shorts..
And started yelling at him.
that you don't come to the sanctuary with shorts..
I mean, is that crazy?.
That's kind of crazy, right?.
That's kind of crazy..
I still remember some of the churches.
back in my generation,.
they have ties.

$^{281}$at the door..
And whenever they pick,.
they find a guy without wearing a tie,.
you can borrow a tie, don't worry..
You can borrow a tie and tie it around your neck.
and you come to the holy sanctuary to worship God, right?.
Do you want to do that?.
No, right?.
I find it funny..
Now I've grown.
after all these years..
And I find it funny because.
the New Testament actually tells us.
to dress down..
In the Roman order of things,.
you dress up.
to focus on your social status, right?.
That's why the wealthy woman.
would dress up when they go to church..
And when they dress up to go to church,.
they will emphasize their own high status in the society..
So that other poor people also come to the church,.
then there is a contrast, right?.
Now, poor is,.
poor is not trying to censor their clothing.
or what they wear..
Poor is trying to say in the community of the church,.
we don't emphasize our social economic differences.
because the ground around the cross is level..
At the foot of the cross, we are all sinners.
and we only justify by grace..
And that's the spirit, right?.
That is what the church means..
We come as a community to focus on others, not ourselves..
Great example, right?.
If you take a look at the previous context of 2-5,.
you know, the hymn for Christ,.
how he made himself nothing,.
how he died on the cross as a slave..
Paul was actually teaching the Christians.

$^{321}$how to be together, how to form a community together..
He said, if you have any encouragement.
from being united with Christ,.
with any comfort from his love,.
any fellowship with the spirit,.
any tenderness and compassion,.
then make my joys, Paul says, complete.
by being like-minded, having the same love,.
being one in spirit and purpose..
Do nothing out of selfish ambition..
In humility, consider others better than yourselves..
Each of you should look not only to your own interests,.
but also to the interests of others..
Jesus was crucified and died as a political victim..
That is the example of servanthood..
We follow Jesus, we do the same..
Now, the apostles continue with the phrase.
"under Pontius Pilate.".
You know who that guy is?.
Pontius Pilate?.
Pontius Pilate..
Now, under the Roman occupation at the time,.
Rome will send a local governor, a Roman,.
to rule over conquered cities,.
like Palestine, the land of Israel.
at that time in the first century..
So, the Roman governor would rule over the Jewish people,.
and of course, the Jewish people would hate the governor..
So, there's a huge political disconnection there..
There's a huge political tension between Pontius Pilate,.
who ruled Palestine at the time, AD 26 to 36,.
against the Roman population..
There was story after story outside the Bible.
talking about people rebelling against the authorities,.
protesting, and all that stuff..
And Pilate was a crafty politician..
Remember, the Jewish authorities trying to persecute Jesus.
and crucify him and crucify him, they yell,.
and then they bring Jesus to Pilate, the governor,.
because as local authorities, the Jerusalemites,.

$^{361}$they could not impose a death penalty..
So, they bring him to Pilate..
As the Roman governor at the time,.
Pilate wanted to keep the peace, to please the people..
That's why he had Jesus crucified..
So, Pilate was a crafty politician..
It's like today, right?.
You hear about who you should vote for,.
someone's going to try to win your vote, right?.
Don't get into the trap,.
because they are all politicians, right?.
They are all politicians..
No political party can represent the kingdom of God..
The kingdom of God is so big,.
bigger than left or right or center..
A different reality, different order of things..
Now, here's the riddle..
As we talk about the Apostles' Creed,.
under Pontius Pilate, that's very interesting,.
because there's a bad guy in this Apostles' Creed, right?.
It's supposed to be something really pure, right?.
Document, talking about the Christian faith..
Why did they put Pontius Pilate there, right?.
A bad politician there..
Well, this is a historical note, right?.
Christianity is historic..
Something that happened in real time and space, right?.
It's not some mythic stories..
It's not some fictional stuff that we believe in,.
but something that actually happened.
in real time and space..
Pontius Pilate, governor of Palestine, AD 26 to 36, right?.
If you check the Roman history,.
or the history of the eastern Mediterranean area,.
you would find this name show up.
on other different historical documents..
Many years ago, there was a city in the land of Palestine.
called Caesarea Militima..
It's a harbor..
It's a huge man-made harbor..

$^{401}$People were digging there, and they found a piece of stone.
being reused later time in a Roman theater..
And on that stone, there were some inscriptions..
Now, the riddle is, can you find the name Pilate.
on the inscriptions?.
Can you find it?.
Can you make it out?.
Where is it?.
Can you find P-I-L-A-T-U-S, or V-S?.
Can you find it?.
Where is it?.
The first line says Tiberium, right?.
Tiberius was also the Roman emperor at the time, Tiberius..
So there's a Tiberium, some sort of temple,.
dedicated to the boss, right?.
The CEO of the entire Roman Empire, right?.
Tiberium, so this kind of a dedication inscription.
for a temple in Caesarea Militima..
Can you find Pilate?.
Can you see that?.
Yeah, Pontus Pilate, right?.
Pontus Pilate, and the prefect of Judea, right?.
So Pilate is this historical thing,.
he's a historical figure, right?.
When they unearthed this thing,.
they finally find something in the stone.
that says Pilate was a real guy..
Of course, Jesus was a real guy, right?.
As the Creed says, he was crucified on the third day,.
and then he rose again, ascended into heaven,.
and is now seated at the right hand of the Father..
He will come again to judge the living and the dead..
The empty tomb is being affirmed.
as an other historical fact..
The tomb of Jesus is empty..
Point blank, telling you, the tomb of Christ is empty..
But how do you explain the empty tomb?.
I want to close today's message with a pretty long quote..
I find this book is one of the best books.
on the reality of the empty tomb..

$^{441}$This guy is an Israeli archaeologist,.
Israeli historian, Shimon Gibson..
Now I find this book pretty interesting.
because he's, you know, Gibson is not,.
he's not a Christian, obviously..
He's a Israeli Jew..
He's a historian, and he has written this pretty good book,.
'cause this book begins with some sort of presuppositions..
The presupposition is that the four gospel,.
Matthew, Mark, Luke, and John,.
are quite reliable documents..
So he started to investigate what is being told.
by the four gospels..
Are they really eyewitnesses?.
Are they really historical accounts, reliable?.
We can trace back to the time of Jesus,.
and then he's tried to unearth other stuff.
and to match the evidence..
Toward the end, he says,.
"There were different strange theories.
"having been put forward to explain away the empty tomb.".
But Gibson said, "They are all based on nonsense..
"Archaeology and history cannot prove or disprove.
"the possibility that Jesus' body was smuggled out.
"by persons unknown, but this is a matter of speculation.".
Now here's the important point..
The reality is there is no historical explanation.
for the empty tomb..
Other than, now if you try to come up with different theories.
different ways to explain away the empty tomb,.
Gibson is saying, there's no way you can do that..
Because there's no historical scientific explanation.
for the empty tomb..
Unless you adopt the biblical reason..
The biblical reason..
The theological explanation..
Which says, the resurrection..
And he says, "I leave it up to the reader.
"to make up your mind.".
I find this very interesting conclusion, right?.

$^{481}$You can do historical inquiry.
of why the tomb is empty..
Somebody stole the body,.
or maybe the disciple went to a wrong tomb,.
or something like that, right?.
You can explain away the empty tomb,.
but ultimately, there's no historical scientific ways.
to explain why the tomb was empty..
There's no way..
Because that goes beyond historical inquiry..
That goes beyond all scientific knowledge can tell you..
At the end of the day, you can only conclude.
that the tomb is actually empty..
But why?.
When you talk about the reason,.
there's only one reason given by the Bible..
Whether you accept it or not..
It's a theological reason..
That God raised Jesus from the dead, right?.
God raised Jesus from the dead..
It's a miracle..
It's something only God can do..
But to accept that reason,.
you cannot stay in your own order of things..
You cannot maintain your own worldview, so to speak..
All right, you have to change your own worldview radically.
in order to accept that biblical theological reason.
that God did something to the body of Jesus,.
raised him from the dead..
And that worldview transformation.
is to say that, okay, I don't maintain my own.
closed universe, closed system kind of worldview..
Reality is not just here, now, and something I can touch..
Right, I can sense..
It's not just that..
There's a reality higher than the order of things.
that I can see..
And that is a supernatural metaphysical reality..
I think Gibson, I have to agree with Simone Gibson..
You have to make up your mind.

$^{521}$whether you accept the theological reason or not..
Now, I want to close today's message by quoting this verse..
Galatians chapter two, verse 20..
Paul says, "I have been crucified, the old self.".
Paul is talking about his old self..
"I have been crucified with Christ,.
"and I no longer live, but Christ lives in me.".
The new self..
When you accept Jesus, there should be.
a fundamental change in our lives..
Our attitude, our worldview,.
our perspective on order of things..
Christ lives in me..
The life I live, the new life I live in the body,.
I live by faith in the Son of God..
And He loved me and gave Himself for me..
This is no longer something out there,.
some objective fact or something that happened.
in the world, in time and space..
Yes, that's true, historical fact,.
and Pontius Pilate, Jesus, the real guy,.
and all that stuff, empty tomb..
It becomes a subjective experience..
It becomes your own personal experience..
Christianity, biblical faith, it's all about relationship..
It is your relationship with Christ..
But please understand, this sort of relationship.
is not just, "Oh, you know, patting my back,.
"and you'll be okay.".
It's a radical transformation, how we see the world,.
and how we see ourselves..
Well, I want to recite the Apostles' Creed together,.
and then we'll pray together..
That's all we recite together..
You all stand..
Apostles' Creed..
Let's recite together..
I believe in God, the God Almighty,.
creator of heaven and earth..
I believe in Jesus Christ, His only Son, our Lord..

$^{561}$He was conceived by the power of the Holy Spirit,.
and born of the Virgin Mary..
He suffered under Pontius Pilate,.
was crucified, died, and was buried..
He descended to the dead..
On the third day, He rose again..
He ascended into heaven,.
and is now seated at the right hand of the Father..
He will come again to judge the living and the dead..
I believe in the Holy Spirit, the holy universal church,.
the communion of saints, the covenants of sin..
Let's all pray..
Father, I give you thanks for breaking into our reality,.
breaking into our world,.
to tell us that you are real,.
to tell us that you have done something mighty,.
something that is beyond our explanation.
in real time and real space..
In history, you have come,.
and you die on the cross for our sins,.
and you rose from the dead.
to give us new hope and eternal life..
Father, this new reality, may you help us to be reminded.
and to stand on it,.
to base our whole life on this foundation,.
knowing that your kingdom is eternal,.
and your kingdom will come..
And we pray that your kingdom will come into our lives,.
change us gradually,.
radically transform us into the image of your Son..
And thank you for the relationship that we have.
at the foot of the cross..
Thank you for saving us and changing our identity..
And we are your children..
We are your sons and daughters..
And may you use us, may you continue to change us.
into your Son's image..
In Christ's name we pray, amen..
Well, you're dismissed..
Thank you very much, and I will see you next Sunday,.

$^{601}$and worship God together..
[BLANK AUDIO].
\newpage



\section{}
\label{sec:EpFmOvtCyQg}
\textbf{God's Design and Calling (Joshua 6) - Tommy Lee}
\newline
\newline
連結: \href{https://youtube.com/watch?v=EpFmOvtCyQg}{\texttt{https://youtube.com/watch?v=EpFmOvtCyQg}} ~~~~ 語音日期: 2024-03-19
\newline
\newline
\hyperref[sec:3W39B_RRLi0]{\small{< < < PREV SERMON < < <}}
~
\hyperref[sec:index_chronic]{\small{[返順時目]}}
~
\hyperref[sec:index_scriptual]{\small{[返順卷目]}}
~
\hyperref[sec:3t6kNRmRV_s]{\small{> > > NEXT SERMON > > >}}
\newline
\newline
$^{1}$Right here?.
All right..
Okay..
Today we're going to look at a passage that I had asked Kelly..
I said, "Hey, Kelly, when was the last time you preached this passage?".
Because I wanted to make sure that it's been two years since we looked at Judges chapter 6..
And when I first looked at the story of Judges 6 with Gideon,.
we looked at the idea of tearing down the altars and putting the proper altar in its place..
But today we're going to talk about a topic that I think applies to every single one of us in the room,.
is the idea of work..
60\% of all of our time throughout the entire week is spent at a job..
And a lot of times within that job, some of us may love what we do..
Some of us may just sit in there and say, "Man, is this all there is to my life? I hate my job.".
And we're left wondering, "Lord, what is it that you're calling us to do in the idea of work?".
Faith, work, integration is so important for us as Christ followers.
because we do that all the time for a good chunk of the rest of our lives..
We're at work until you decide to retire..
Whether it's at 60 or whether it's 75 or 80 years old, you will be working..
And I wish that some of us were probably sitting there, could be like David,.
where Samuel calls out and is trying to select the next king..
And David shows up and he whispers to David, "You're going to be the next king.".
But it doesn't always happen, does it?.
Some of us are like Joseph, where we kind of land in exactly at a time when God uses us..
Joseph is a guy that for some odd reason, his brother sold him into slavery..
And now he is a slave..
And somehow or another, God has given him the skills to rise up in the ranks in the palace.
of where he sold in slavery..
But then his boss's wife accuses him of sexual malpractice and that he's thrown into prison..
And later on, because of all the things that are happening,.
he rises up to be the chief of staff in the political system..
Some of us are like Nehemiah, who is a cupbearer to the king..
He tastes food so that the king, his boss, doesn't get poisoned..
And somehow or another, God uses the cupbearer to the king.
to then rebuild the walls of Jerusalem because he has given him that burden..
Or then some of us look at the two tax collectors in the Gospels..
There's the tax collector of Matthew, where Jesus says to Matthew,.
"Pack up everything and follow me.".
But yet at the same time, he says to Zacchaeus, who climbs up a tree to see him,.
"I want you to follow me by staying as a tax collector.".
So how do you know when to be a tax collector?.

$^{41}$How do you know when to drop everything and follow Jesus?.
How do you know when to be a king?.
How do you know when to do all these different things?.
All of us are on the journey, aren't we?.
Some of us are called to be like Pastor Mark is a full-time vocational ministry..
Some people are called to be a stay-at-home mom..
Some people are called to be a teacher or a nurse..
Some people are called to be an engineer, an entrepreneur..
How do we begin to know what God has gifted us to be?.
And it becomes a journey that he leads all of us on..
And sometimes within the journey, there's different chapters in our life within that journey..
And within those chapters, is that where you want me to be?.
Maybe, but maybe not..
And so sometimes he'll ask us to pivot, and 20 years from now, we'll look back at our life and say,.
"Wow, that's what you've been preparing me for.".
I think I've told all of you guys I spent the first 10 years of my life building cell towers..
Then from that point on, I spent about three years working in a Korean church,.
while at the same time, also at that point in time before that,.
serving in a church over in Chinatown at CCUC..
And then from that point, I worked at Moody Bible Institute for a number of three years..
And then now for the last 10 years, I work at Resource Global,.
and Wheaton College, Northern Seminary for a period of time..
Lord, is that what you're calling me to do?.
I thought it was building cell towers, or rather that..
I went to college thinking I was going to medical school..
God calls us to do different things..
So today, he calls this particular person,.
and we're going to look at that person,.
and we're going to glean some stories that he talks in..
What the whole idea is, what is God calling you to do right now?.
And can you and I be faithful to where God has placed us in?.
Because God, for a period of time, has called you at that point in time,.
at that point in place, to make a difference for the gospel..
We start in Judges 6..
I studied this for a number of years..
And it is such a powerful--it begins with this whole idea,.
"The people of Israel did what was evil in the sight of the Lord..
And the Lord gave them into the hand of Midian for seven years..
And the hand of Midian overpowered Israel..
And because of Midian, the people of Israel made for themselves.

$^{81}$the debris that are in the mountains and the caves and the strongholds..
For whenever the Israelites planted crops to Midianites and the Amalekites.
and the people of the east would come up against them..
They would encamp against them and devour the produce of the land as far as Gaza.
and leave no sustenance in Israel and no sheep or ox or donkey..
For they would come up with their livestock in their tents..
They would come like locusts in numbers..
Both they and their camels could not be counted,.
so that they laid waste to land as they came in..
And Israel was brought very low because of Midian..
And the people of Israel cried out for help in the Lord..
Don't ever misinterpret how God uses events in our life to turn us back to him..
And too often with the Israelites, time and time and time again,.
the Israelites go and do whatever they want to do..
But yet God uses situations and circumstances in our life.
to call and draw us back to them..
And some of those circumstances may be trials in hard times..
And here we see an example of that in the hands of the Midian..
When the people of Israel in John chapter, Judges chapter 6, verse 7,.
when the people of Israel cried out to the Lord on count of the Midianites,.
the Lord sent a prophet to the people of Israel, and he said to them,.
"Thus says the Lord, the God of Israel,.
'I led you up from Egypt and brought you out of the house of bondage,.
and I delivered you from the hands of the Egyptians.
and from the hand of all who oppressed you,.
and drove them out before you and gave you their land..
And I said to you, 'I am the Lord your God,.
and you shall not fear the gods of the Amorites in whose land you dwell,.
but you have not obeyed my voice.'".
And now the angel of the Lord, when the angel of the Lord appears,.
you listen to him, "came and sat down on the terebth at Oprah," in verse 11,.
"which belonged to Joash the Abyssalite,.
while his son Gideon was beating out wheat in the winepress.
to hide it from the Midianites.".
And here's a key verse, and we've looked at it before, verse 12,.
"And the angel of the Lord appeared to him and says to him,.
'The Lord is with you, O mighty man of valor.'".
In verse 12, "The Lord says to him, 'The Lord is with you, O mighty man of valor.'".
If you were to look at the NIV, I'm reading from the ESV,.
if you were to look at the NIV, it says, "O mighty, mighty warrior.".

$^{121}$"The Lord is with you, O mighty warrior.".
And when you're looking at that, your first thought is,.
"Wait a minute, Lord, why are you calling him a mighty warrior?".
That is the verse that sticks out..
Why is the angel of the Lord calling him a mighty warrior?.
Why is that?.
He's threshing wheat in a winepress..
That and being a mighty warrior don't go hand in hand..
And you're wondering why..
Because if you go into it, it's harvest time, right?.
It is harvest time, the wheat is ready to be harvested,.
and you're pulling out all of the wheat, and you pull out all the stalks.
and everything like that, and you usually grab all your wheat,.
and you lay out on top of a threshing floor..
And a threshing floor is usually an inclined space out in the open,.
on a hill and everything like that, and you'll understand exactly what it is,.
and you'll lay out all your wheat on the threshing floor..
And from that point on, if you have a lot of wheat, you'll get your oxen,.
you'll get your horses, you'll get all that stuff,.
and you're going to stomp on it, stomp on it, stomp on it..
If you have a tiny little harvest, you'll take a little flail or hammer,.
and you sit there and you'll hit the wheat over and over and over again,.
and it separates the stalks and it separates the head..
And from that point on, once you do that, remember it's out in the open,.
you'll take a little rake or a winnowing fork, and you'll sit there.
and you'll lift it high and you'll throw it high up in the air..
And it separates the two areas..
It separates the heavier grain on one side and the straw on the other side..
And then there's seed covering the chaff, and as you're throwing it up.
with this winnowing fork, it's blowing all over..
And that's why it has to be up on an inclined space out in the open.
where there's a lot of breeze that blows everything away..
And that is threshing wheat on the threshing floor..
But where is this guy at? He's in--what? He's threshing wheat where?.
In a wine press. What's a wine press?.
It's literally a hole, a rectangular hole, in which when I stand in there,.
it goes up to my waist, and you lay out all your grapes right there..
And you sit there and you stomp on it..
And there's a little hole which leads to a set of cisterns or jars.
and everything like that..

$^{161}$And as soon as you sit there and you stomp on it and you stomp on it.
and you stomp on it, all the juice rolls into a cistern..
But yet it's not an inclined space..
And yet he is threshing wheat. In where?.
In a wine press..
And the angel of the Lord, the Lord, the angel of the Lord calls him,.
"The Lord is with you, mighty warrior.".
And when you're sitting there in verse 12, looking at it,.
you're sitting there, "Okay, Lord, what are you seeing?.
He's not a mighty warrior..
He's afraid. He's a coward..
So why are you calling him a mighty warrior?".
Two things, two things, two very important things..
Number one, the Lord sees you differently..
The how you see yourself and not how others see you..
The Lord sees you differently to how you see yourself.
and how others see you..
The Lord sees in him Gideon, a mighty warrior..
We see in him a guy who's afraid..
That's why he's threshing wheat on the wine press..
We see a guy who's scared..
We see a guy who lacks courage..
And we don't understand what God is trying to say to him..
"The Lord is with you, almighty man of valor..
The Lord is with you, mighty warrior." Yes!.
Because God sees you differently.
to how you see yourselves and how others see you..
That's number one..
Number two, in terms of this, God has a name for you..
God has a purpose for you..
God has a vision for you..
For Gideon was a mighty man of valor, mighty warrior..
What's the name that he has for you?.
What's the name that you believe that he's calling you to be?.
If we all sat here in a group and had an exercise.
and we got a chance to brag each other,.
and I know a lot of times primarily with us Asians.
and other stuff is we don't like to talk,.
we don't like to brag about ourselves,.
but we are familiar with what skills and gifts we do have..

$^{201}$What's the name that he's given to you?.
What vision does he have for your life?.
For Gideon is the mighty warrior..
Maybe for Mark, a lot of times, it's leadership.
or it's shepherding or something like that, administrative..
Maybe it's the ability to convene people..
Maybe it's the ability to inspire people..
Maybe it's the ability to build together a team.
or whatever it is..
Maybe a lot of times your intellect is so brilliant.
because that's how God is going to--.
he's using you to do things and change education..
Maybe it's because your skill, you're going to change healthcare..
Whatever it is, what is that vision?.
Maybe it's being an engineer to create innovative tools..
What is that vision in your life?.
As we continue to seek the Lord,.
what is he calling us in our life?.
There's this one lady that I got a chance to meet..
She's Filipino, and if you came, she's very quiet,.
and she would just sit there, and you probably wouldn't even know her..
Her name is Nikki..
But what you don't know about Nikki is she is a brilliant lady..
She is extremely talented, teaches part-time at Stanford University..
Then she's an entrepreneur, started working in education,.
worked in colleges and universities and the executive level,.
but then she started realizing she doesn't have a medical degree.
and all that stuff..
She was great at entrepreneurial things..
She was great at starting up and launching different ideas..
Now she teaches part-time at Stanford..
She's the founding CEO of a company that is working with.
artificial intelligence and nanobots to replace--.
Kelly, what is that thing when you put oxygen in your--.
- [Kelly] The nasal canthopathy?.
- Something like that, where it puts oxygen in your--.
You can breathe through oxygen, right?.
I completely blanked on it..
That's what she's working on right now..
And why in the world was she doing that?.

$^{241}$Because during COVID, her mother died because she wasn't getting.
enough oxygen in her body..
And she's a single older lady, and she started searching.
while she was in San Francisco with her mother passing away..
How do we address this problem?.
Because her mother was at older age, and she wasn't able to get.
enough oxygen..
They weren't able to put enough oxygen in her body, and she.
passed because of COVID..
And God used that experience as hard as it was as a Christ follower.
to look and address the problem..
And she started finding out these different doctors and what they're.
doing, and through her entrepreneurial skill, launched this entire.
company that she's doing to change how she integrates her faith.
in her work..
What is the name that God has given us?.
Gideon, it was "Mighty Warrior.".
What is the name that he has given you?.
Because how God looks at you, how God sees you, is different than.
the picture that we see of ourselves and that others see us..
Go on, and Gideon says to her, "Please, sir, if the Lord is with.
us, then why then has all this happened to us?.
And were all his wonderful deeds that our fathers recounted to us.
saying, 'Did not the Lord bring us up out of Egypt?'.
But now the Lord has forsaken us and given us into the hand of.
Midian.'.
And the Lord turned to him and says, 'Go in this might of yours.
and save Israel from the hand of Midian.'.
Do not I send you?'.
And verse 15, "And he said to him, 'Please, Lord, how can I save.
Israel?.
Because my clan is the weakest in Manasseh, and I am the least in.
my father's home.'".
Take a look at that..
You know how Gideon saw himself?.
The angel of the Lord, the Lord sees him as mighty warrior..
Gideon himself sees him as the weakest clan and him as the.
weakest in the weakest clan..
"Behold, my clan is the weakest in Manasseh.".
Of all these clans that we have here, my clan is the weakest, and.

$^{281}$I am the least in my father's home..
You don't understand this..
My clan of all the clans is the weakest, and I am the weakest in.
the weakest clan..
You know where that put me?.
Dead last, dead last..
That's how Gideon sees himself..
That was the name that he has given himself..
He has defined himself as someone who couldn't accomplish much.
because I am the weakest in the weakest clan of everyone, and I.
am last, and what in the world can you do for me?.
And now here I am threshing wheat in a wine press because I'm.
scared of the Midianites, and even though this threshing wheat.
on the threshing floor is not gonna cut it in the threshing,.
I'll take whatever I can get because my name is defined by.
being in the weakest clan and me being the weakest..
What defines you?.
What are the insecurities that you have?.
How do you see yourself?.
How do you look at yourself in the mirror?.
Are you not the smartest?.
Are you not the brightest?.
And you may have all of these ideas, and if you look at the.
story in verse 16, the Lord said to him, "But I will be with you,.
and you shall strike the Midianites as one man.".
Do you not see what's happening here?.
When you look at the verses, you know what the Lord never says.
to Gideon?.
He never disagrees with him..
When he says to him, "My clan is the weakest, and menace I am,.
and the weakest in my weakest clan," the Lord never disagrees.
with him..
The Lord doesn't say to Gideon, "Hey, hey, hey, hey, your clan's.
not the weakest..
So-and-so, I think, is the weaker one..
Look, you're not the weakest of the weakest clan, all right?.
Maybe you're not the weakest, but you're not the weakest..
So-and-so probably is weak.".
No, he doesn't disagree with him..
He doesn't say otherwise..

$^{321}$He agrees with him..
He says, "Go in the strength that you have.".
That when you and I sit here, "Lord, I'm not that smart.".
You know what the Lord says?.
"I know..
Lord, I'm not that talented..
I know..
Lord, do you not know the history that I came from?.
Do you not know how bad I was, my family, the path I came in?.
Yes, I know..
Lord, I came through a poor family..
My parents immigrated here..
My parents immigrated in '72, and a lot of times, I don't know.
how to operate in these worlds..
I know..
You don't have that experience..
I know, I know, I know..
Lord, I'm a little socially awkward..
I get a little shy..
I don't know what to say..
I know..
You know why I know?.
Because I made you..
I created you..
I painted your life..
I designed your entire path in your life..
I know what's going to happen in your life..
I know these things..
I gave you your strengths..
I gave you your weaknesses..
You're stuck, Gideon, that you just see what you have..
You come from the weakest clan, and you're the weakest in the.
weakest clan..
But despite all that, Gideon, you're a mighty warrior, because.
that's the vision I have for you..
What's the vision that he has for your life?.
And how is that manipulated by sometimes the name that we have.
given to ourselves?.
My biggest encouragement to each one of you, and those who are.
watching, those who are sitting here, don't you ever, ever let.

$^{361}$anyone say anything, because your age, your sex, the color of.
your skin, your education, how wealthy you are, how bright you.
are, the quality of your work, your job title, all of that,.
don't ever let people look down on you..
You come before the Lord each and every single day..
You say, "Lord, what is it that you're calling me to do?.
And whatever path you ask me to do, I will be faithful to that,.
and I will give you the best I can to glorify and honor you,.
because you have called me to do it with excellence and to.
continue to pursue the calling that you have for your life.".
And we may not know what that calling is..
God may not reveal to you what that calling is..
It may take 20 years..
It may take five years..
But we're pursuing every single day, and we are faithful every.
single day, and we are coming before the Lord to seek what.
that means..
Let me give you an example of that..
Here's an example..
And I hope as I tell you this story, you're going to sit.
there, "Tommy, what are you trying to do?".
It's not me bragging about myself, but this is what I.
learned about faith, work, integration, right?.
When I graduated from Washington University in St..
Louis, I actually didn't know what I was going to do..
I had no job..
And so I got introduced to a guy who introduced me to a.
company that built cell towers for a living..
And so a lot of times I got a job..
I had no engineering degree..
I was a general management major..
And a lot of times, where am I going to go to work?.
I don't know..
So I get an introduction..
I get an interview..
I sat with this guy named Mike Mitchell, and he's a director.
at Tower Operations..
Later on, I reported this guy named Dan..
But Mike sits there, and he looks at my resume, and he.
says, "Did you have an internship?".

$^{401}$"Nope.".
"Do you have an engineering degree?".
"Nope.".
"Do you have any experience in the cell phone industry?".
"Nope.".
And he puts it down..
He says, "Tell me why I need to hire you..
Just give me a reason why.".
And I said, "There's no reason why you should hire me.".
And he looks at me, and he says, "Really?".
"Yes, there's no reason why you should hire me, but give.
me one year..
I will work as hard as I can, and I will give you everything.
I have, and I will learn this business, and in one year,.
I will make a difference in this company.".
He says, "Oh, yeah..
I like your answer..
You're hired.".
And so the first thing that he says, I get to work, and he.
says, "We just acquired 200 towers from Verizon Wireless,.
and those cell towers and all the file folders and all the.
construction drawings all came in, but here's my problem..
They didn't put any order to it, so they're all in different--.
they're all in this room..
Your first job for the next five months, go organize.
the file room.".
I'm a college grad from Washington University in St..
Louis..
I was a pre-med major..
You want me to be a file room, yes?.
And the Lord said to me, "Tommy, wherever I put you, may.
you be faithful to where I place you.".
And so I went, and I had the best attitude because, "Lord,.
I don't know why you put me here, but whatever it is, let.
me be faithful to where you have put me.".
And I learned that file room really well, and I learned to.
tell the difference between FAA, FCC, environmental, structural.
analysis, everything..
But what I didn't realize during all that time was God was in.
the midst of working and creating me the habits, the.

$^{441}$discipline of looking at details and being organized and seeing.
structure every single day for five to six months..
I had to look at these details..
After that five months was over, a lot of times they promoted me.
over to scanner guy..
I had to scan all these documents..
I said, "Oh, my gosh, I've got to scan everything, create a.
system to scan everything.".
So I did that for another four months..
During that period of time, I had this one guy named Dan..
He was our project director, and he says, "Tommy, you've done.
really well..
Get some files and meet me every single day," and we did this.
for three months at 5 o'clock, and I will show you what it.
means to be a project manager..
We did that for another three months..
So all together, I'm now in my company for nine months, and.
after that, they promoted me as project manager with no.
engineering skills or anything like that..
Here's a couple of other things as well, too..
Now I had a rough boss, but, Lord, are you helping me?.
Are you working with me?.
Are we faithful to where our -- the one thing I encourage you.
wherever you are, are you faithful to where God has.
placed you to be?.
Here I am as a project manager, and now I am working on the.
Sprint project, and usually a lot of times you're building.
20 towers..
This year, I had to build 40 towers, and then one day, I'm in.
a deployment meeting at Sprint, and there's this guy named.
Sonny Peterman..
He's the vice president of our region..
Sonny, you should never work for a boss like that..
The guy was a crook..
I mean, the guy was just terrorizing..
So he walks in, and he sits down unannounced, and the guy.
from Sprint -- now, back in the days, we had audio.
conferencing and he says to me, "Tommy, can you get this done?".
And I said to him, "I will try my best.".
Sonny then says to the guy at Sprint, "Let me mute you, and.

$^{481}$I'll get right back to you.".
And he mutes him, and he says, "Tommy, can you get this job.
done?".
And I said to Sonny, "I will try.".
And you know what he says to me?.
"I don't care if you try..
I didn't hire you to try..
I hired you to get this job done..
Now, if you can't get this done, I will fire you right now and.
replace you with someone who can..
Can you get the job done?".
I says, "I will.".
And he unmutes me and says, "We'll get the job done.".
He walks out of there and says, "Always remember, don't try.
hard..
Get the job done.".
And he walks out..
A year later, and this is the biggest lesson I learned, and it.
ties in with everything that we're talking about a year.
later..
Sonny comes in, and we have an employee gathering, and he says,.
"This year's Employee of the Year is Tommy Lee.".
And he says, "Tommy Lee, wherever I go, in all of the.
travels, and all the -- when I talk with clients like AT and T,.
Sprint, and Verizon, they all say, 'If Tommy Lee's on the job,.
we'll hire you and give you the business.'.
All of the sales development, all of our -- and when I talk.
with the janitors, they all know who this guy is..
And then I talk with him, and all of us go out to the bars on.
Fridays, but this guy goes to his church and works with his.
youth..
That's why I hear him..
And he says to this, he says, "I don't follow Jesus, but if this.
is what it means to follow Jesus, then sign me up.".
And he says that..
And you know what happened after that?.
I had a girl who comes over to me, and she was my admin..
And she says, "Tommy, how are you so joyful?.
I've had three miscarriages..
How do you see God in the midst of all these things?".

$^{521}$I had another guy, his name was Patrick..
Patrick would go out to the bars all the time..
Patrick was always getting drunk..
And he says to me, "Tommy, what keeps you going?.
What's your purpose in life?.
Because I haven't found it yet.".
My boss, Dan, who promoted me as project manager, calls me up one.
day and says, "Tommy, I need you to stay after work..
I need to talk with you about something.".
I walk in after work, and he says, "Hey, Tommy, my daughter.
has been diagnosed with a heart issue..
Can you pray for me?.
Can we spend some time praying?".
And during Christmas, I would give their names,.
give them all Bibles..
And that was what I did for Christmas, and I would put their.
name on it and gave it to each one of them..
And during my cancer journey when I was gone,.
I got these Facebook messages..
And this is now 10 years later..
And I get these Facebook messages..
One of them was from Erin, the girl who had three miscarriages,.
and says to me, "Tommy, I have three beautiful daughters now..
I lead the VBS in their church..
I came to know the Lord, and I read your Bible..
Your Bible's still with me.".
The other guy, Patrick, Patrick says to me, "Tommy, I'm still,.
like, going out all the time and drinking, but I got married..
But I did start going to church, and your Bible's the one that I.
bring to church every single day.".
Dan, these are all 10 years later, his wife calls me up.
and says to me, "I saw on Facebook that you were.
struggling with cancer..
Your former boss, Dan, is also struggling with cancer..
Can you call him because your Bible is still sitting right.
there by his bedside?".
You never know how God uses you..
You never know what God is calling you to do,.
all because you do your job well..
And you are faithful to wherever God has called you to be..

$^{561}$Whether it's a job that you can't stand,.
whether it's a job that you hate,.
whether it's a job where you just don't like the industry,.
where you don't like your characters, all of those things,.
are you faithful to where God has placed you at?.
And are you always praying and asking,.
"What is that vision that you have for me?.
What are you preparing me to do?".
For Gideon, he leads him on this journey..
This guy who is at the weakest of all of the clans,.
that he is the weakest, he is the lowest of all the clans,.
God is preparing him to be the leader of--.
to fight against the minions..
He is going to be their victor..
What is God preparing you to be?.
And how is God using all the things that you're doing right.
now in your job to prepare you for what you're doing right now?.
Because you don't know..
He may call you now..
He may be calling you 10 years from now,.
do something unique or--.
what he's calling you to do is be faithful..
If you look at Nehemiah, he didn't call him until the late.
'70s, but he was faithful as compared to the king..
And when it was time to rebuild the walls,.
do you know what the king did?.
His boss, the king, gave him time off,.
gave him all the money that he need,.
gave him all the papers that he need to get all of those.
different things, and gave him all the resources and materials..
Why?.
Because he liked Nehemiah and he had nothing better to do?.
No..
He was faithful as compared to the king..
Do your job well..
Be faithful day in and day out in your job..
Make such a difference and impact in your job that you leave.
a lasting impact on the lives of people around you,.
that when you walk away from that job,.
they're going to sit there and say, "Wow, look at that person..

$^{601}$I'm able to see Jesus through that person.".
That's what it means to be a Christ follower..
That means you treat the janitors well..
That means you treat the admins well..
That means you treat your boss well..
That means that every single person you treat well,.
that you do a good job, that you are productive,.
that because God has put me in that situation,.
I will honor my boss by doing my work with such excellence..
It glorifies God in all that we do..
And during that time, you are always praying,.
"Lord, what is that vision that you have given me?.
What are you calling me to do?".
Because God may choose and show you different things.
that he wants you to do..
That's so important about faith and work..
Sixty percent of our day is spent at work..
God calls us to work..
God calls us to flourish and redeem..
God has given you skills and talents to change your city,.
to change your community, to change your work,.
to make this world a different place..
What is God calling you to do?.
What is the picture and name that he has given you?.
For Gideon, it was "Mighty Warrior.".
But too often, our insecurities causes us to look past.
what God is calling us to do..
And for Gideon, he had to get past that insecurity.
of being from the weakest clan.
and being the weakest in the weakest clan..
What is God calling you to do?.
We are motivated by grace..
We are saved by grace..
The cross causes us, motivates us to be where we are..
Our job now is to take that grace that we have in our life.
and to make a difference in all that we do..
In Christ we pray. Amen..
Father, thank you for how you continue to serve and use us,.
how you've created us to become the person you want us to be..
And so, Lord, may you continue to paint that journey over to us.

$^{641}$and allow us to see what that name that you have given us is, Lord..
What is that name you have given us?.
How have you uniquely designed us to flourish.
and change this world and redeem this world for you?.
I pray all of us will be in that journey to see you..
In Christ we pray. Amen..
Thank you. Thanks, Mark..
\newpage



\section{4:13-17}
\label{sec:3t6kNRmRV_s}
\textbf{Planning For Tomorrow (James 4\_13-17) - Rev. Mark Chan}
\newline
\newline
連結: \href{https://youtube.com/watch?v=3t6kNRmRV-s}{\texttt{https://youtube.com/watch?v=3t6kNRmRV-s}} ~~~~ 語音日期: 2024-03-19
\newline
\newline
\hyperref[sec:EpFmOvtCyQg]{\small{< < < PREV SERMON < < <}}
~
\hyperref[sec:index_chronic]{\small{[返順時目]}}
~
\hyperref[sec:index_scriptual]{\small{[返順卷目]}}
~
\hyperref[sec:3kSkaNorqlI]{\small{> > > NEXT SERMON > > >}}
\newline
\newline
$^{1}$Yeah, that will be the announcement for all..
And let's welcome Pastor Mark..
- Hey, good morning, everyone..
Good to be here..
And I'm sure you are very, very happy that I'm back..
And can you check the mic for me?.
Yeah, there you go..
All right..
Well, anyways, so today I have two parts..
The first part is I will talk about my experience.
going to California..
It's a heck of a trip, you know?.
So we spent two weeks..
I take a long break from the church, two weeks..
And then we travel to California..
And it was a very exhausting trip, very exhausting..
And I'll tell you why..
This is the first part..
And then the second part is I'll tell you.
what my spiritual reflection is.
during this crazy trip, all right?.
So you remember December 25th, Christmas Sunday,.
right, remember that, right?.
We had a big celebration here..
And then we had baptism and we had the luncheon..
And then we played the white elephant game, right?.
So it was such a fun, right?.
And then we finished it in about one-ish.
or almost two o'clock and we depart from this place..
And because our flight flying from Chicago to Phoenix.
is gonna take off around 5.45 p.m..
So I still had some time..
So my whole family, actually they all left already.
after the Christmas program..
And then I was the only one at the church..
And I was rushing myself and just wrap up some stuff..
And then I went home..
I didn't pack at all..
So all my three family members, they have already packed..
And then I went home in such a hurry..

$^{41}$I try to pack my clothing, pack all my essentials.
into my suitcase..
So December 25th, right?.
And then around two o'clock, 2.30,.
I keep checking my Southwest app,.
finding out my first flight from Chicago to Phoenix,.
what time is departing..
Because my friends and my family members in California.
reminding me, "Are you sure your flight.
"is departing on time?.
"Please check your flight.".
So I keep checking and it was okay at 5.45..
And then a little bit later, got delayed for 15 minutes..
And then later on, got delayed for 30 minutes..
So it was good..
It was all good because I was in such a rush, right?.
I don't wanna forget something, right?.
So it give me some time and almost an hour delay..
And then I sense that it's gonna be a trouble, right?.
It's gonna be trouble..
Because I checked the phone app.
and then I go to the website Southwest.
to check my flight departure schedule.
and they don't match..
So amazingly, they don't match..
And then I was talking to some of my friends.
and relatives in California..
Hey, maybe my flight is gonna be delayed..
Now that's a problem because I'm flying to Phoenix.
and then there's a layover less than an hour.
and then connecting flight from Phoenix to Sacramento, right?.
So if the first leg got delayed,.
then I would not be able to catch the second leg, right?.
So I talked to Nathan, Rachel and Wendy..
I said, we gotta go to the airport physically.
to go to the counter and check the ticket, right?.
Or we may have to rebook the ticket..
So we have Ivan, one of our neighbors,.
have us drop off at the airport.
about 3.30, right?.

$^{81}$Remember my first leg was departing,.
originally departing 5.45 p.m..
So we were there two hours before.
and then we got dropped off at the sidewalk..
And then I walk into the departure hall,.
going to the Southwest counter..
Go inside, there was a long line.
of almost, I think, almost 300, 400 people..
They're waiting to go to the Southwest counter.
and have their tickets rebooked, right?.
So we got in line, 3.30,.
and then wait and wait and wait..
The line's moving so slow.
because the counter has only three or four people working..
And then we waited, I tell you,.
we waited for almost four, five hours..
And almost the counter is in sight, right?.
Almost the counter is in sight..
And then one of the managers step up.
and shout out loud and said,.
"Sorry, computer crash..
"And you have to wait for a little bit more.".
And I said to myself,.
"It's gonna be a big trouble for me, right?.
"There's no way I can catch the second leg, right?".
And then, there's a huge line in front of me..
And then I keep checking all the Southwest flight,.
almost all of them got delayed..
Or many of them got canceled..
And my own flight keep delaying, right?.
5.45 to six o'clock, and then eight o'clock,.
and then 9 p.m., right?.
So by the time we got to the counter,.
it was almost 9 p.m., right?.
Right, almost 9 p.m., something like that..
It was seven, seven, eight, 8 p.m., something like that..
It was a long wait..
Now during the wait, I was like,.
there was no way I can get the second leg, right?.
So my decision was whether I would catch the first leg.

$^{121}$and go to Phoenix, right?.
And going to Phoenix for what, right?.
So I was doing my research on my phone.
and see if we can just fly to Phoenix.
and then just figure out what to do next, right?.
So I was trying to figure out what to do..
And I text Joyce, one of our pastoral staff in Phoenix,.
and then I was doing some research on the distance..
How long does it take from driving from Phoenix.
all the way to Sacramento?.
So I was doing some research, and I had done my research,.
and I talked to my family members,.
my daughter, my son, my wife,.
and I called them out and said, "I have a crazy idea.".
And the crazy idea is we're gonna fly to Phoenix,.
and then we're gonna rent a car,.
we're gonna drive to Las Vegas,.
and then we're gonna stay over there,.
and then we're gonna drive from Las Vegas.
all the way up north to Sacramento..
It was a long drive..
So I did my research, it's about 15 hours of driving,.
and then there's a midpoint we can stop at Las Vegas..
So I asked them, "This is my crazy idea..
"Are you up for it?".
And you know what they said?.
They said, "Let's go..
"Let's go.".
So we went to the counter and talked to the lady and said,.
"Hey, just let us fly to Phoenix.".
And then she was questioning us,.
"What are you gonna do in Phoenix?".
"Hey, no problem, we have a good friend there,.
"and we're gonna stay overnight there,.
"and we'll figure out what to do.".
So we fly to Phoenix, and then finally,.
by the time we board the plane,.
it was, I think, almost 11 p.m., right?.
Almost midnight..
And we were in the airport for nine hours already,.

$^{161}$almost nine hours..
And I went to the gate, and it was chaotic, right?.
Yeah, James was there too, right?.
It was crazy..
I was happy to see James, but I'm sad.
that we got stuck right there, right?.
And then there were a lot of people I overheard..
Their flight got delayed and delayed,.
and then finally canceled, and they gotta go home..
So I was praying that my flight would be okay,.
and then we will fly to Phoenix..
So by the time we board the plane,.
it was almost midnight, and we were ready.
for a three-and-a-half-hour flight.
from Chicago to Phoenix..
So we arrived at Phoenix Airport..
It was about 3 a.m., I think, 3 a.m..
So we got off the plane and went to the baggage claim,.
and with a high expectation.
that we will get our suitcases, right?.
And then the bell started rolling..
There's only a few bags there,.
and our bags were not found..
They were not being sent to Phoenix..
They didn't follow us..
So another long line, we go to the Southwest baggage claim.
and try to find out where our bags are,.
and it turns out our bags did not follow us..
They stayed in Chicago..
So we were separated from our suitcases..
And 3 a.m., 4 a.m. at Phoenix,.
luckily, fortunately, we had our hand carry.
with one set of clothing, our essentials,.
in our hand carry so we can live by it.
with just the minimum..
So we were pretty, pretty sad, pretty sad..
Went to the car rental and went to the parking lot..
It was kind of scary because 3 a.m., 4 a.m.,.
there's nobody there, and there are not many cars..
I guess there were so many stranded travelers,.

$^{201}$they were trying to rent their own cars, right?.
So the car rental guy said to me,.
"You know, I'm sorry, we don't have.
"the standard size you requested,.
"but we have a free upgrade for you.".
Yeah, free upgrade for you..
Okay, okay, that's great..
And then he said a number of brand,.
number of different models of cars,.
and I heard Camry, I got out of Camry..
(audience laughing).
It turns out it's a new Camry, all right?.
So we hopped into the car and drove to Joyce,.
and we spent a day in Phoenix..
You know, with all the trouble,.
I think it was a blessing to see some friends.
and to visit Abby and had a good time there..
So it was great..
And then after a day in Phoenix,.
we decided to drive all the way from Phoenix to Las Vegas..
Now this is a crazy idea,.
but I think it's a very good idea, right?.
Because Rachel and Nathan,.
they'd never been to Las Vegas..
Yeah, why not, right?.
We go to Sin City at a stop in a trip,.
and so much fun, right?.
So we got to the hotel, it's a big hotel,.
and again, it was a long line,.
because I guess all the stranded travelers.
were looking for a place to stay..
So we waited in line for an hour,.
and finally we got our room keys..
And it was pretty tiring, pretty exhausting..
But hey, we're in Las Vegas, we could have some fun..
And so finally we were there..
You know, I did not gamble,.
because the whole trip is a gamble..
(audience laughing).
Because the whole trip is already a gamble..

$^{241}$It was crazy..
We got to Las Vegas, and then took some pictures,.
but we were really, really tired..
We were really tired..
But we had to go on, right?.
So I think Tuesday afternoon,.
we decided to drive all the way.
from Las Vegas to Sacramento, right?.
So we drove to Bakersfield, north of Los Angeles,.
and then all the way up I-5.
to Bay Area and Northern California..
That's how the trip was,.
and it was a really tough drive,.
because it just happened during that week,.
California was hit by huge rainstorms,.
and we were almost in every hour driving in heavy rain..
And it was crazy..
I remember there was a hilly portion of the highway.
near Bakersfield in Southern California..
It was a very mountainous region,.
and I had to turn, twist up and down,.
and it was at night,.
and I could not see the lane line, you know?.
Now my kids can tell you,.
if I drive below the speed limit,.
you know, it was pretty scary drive, right?.
When I drive slow, it was scary..
When I drive faster, it was okay, right?.
I couldn't see a line because I'm getting old as well, right?.
And then I turned on the lane detection on the Camry..
Hopefully it will help me to steer the turn, right?.
And the lane detection did not work,.
because it was so dark..
Even technology cannot see it, right?.
Oh, it was crazy..
But finally we arrive at Sacramento,.
and it was nice..
Sorry for the picture, not the best picture,.
but you know, as I told you,.
I usually don't take long breaks from my work, from church..

$^{281}$You know, some brothers and sisters,.
they are really caring about my wellbeing..
They always ask me and tell me,.
"Hey, Pastor, why don't you take a long break?".
Right? Yeah..
I know why I don't take long breaks..
(laughing).
Now I know, right?.
Because once I take a long break,.
this is a crazy break, right?.
But in that two weeks, I think it was still a blessing..
It was still a blessing,.
because I see all the friends and friends and relatives..
We had a good time,.
and there were also some people I have met online before,.
and I have never seen them..
I was able to have a meal with them and to talk to them..
So it was such a blessing to see.
all the friends and family together..
I think every gathering,.
every moment that we can see each other,.
see some of the friends and relatives, was a blessing..
Every single moment, I enjoyed it,.
even though it was exhausting..
During the whole week,.
I was constantly reminded of one single passage in the Bible..
It is in the letter of James, chapter four..
Chapter four..
Chapter four, 13..
"Look here, you who say,.
"today or tomorrow we are going to a certain town,.
"and we'll stay there a year..
"We will do business there, and we will make a profit..
"How do you know what your life will be like tomorrow?.
"Your life is like morning mist..
"It is here for a little while, and it will be gone.".
Now this is the passage I keep reminding myself.
during the entire journey..
James, as a pastor, he was writing this letter.
to the brothers and sisters who were scattered.

$^{321}$because of persecution in the church in the first century..
And reminding them, or even confronting.
some of the people in the congregation.
who had a wrong attitude toward life and toward God..
You notice here in this short verse,.
James is quoting some of the rich guys in the congregation.
who were so resourceful..
They were merchants, they probably are traveling merchants.
in those days..
And notice how he quote them, saying,.
"We will go, we will stay, we will do business,.
"and we will make a profit.".
He was confronting these people who thought.
they control, that they are in charge of their lives..
And with this confrontation, James is saying,.
"To the people, you are not in charge, so be humble.".
We are not in charge of our own life..
Every single moment is all by the grace of God..
Right, this whole trip, I really, really learned this lesson..
It didn't mean I didn't know before,.
but it was just a personal experience that is telling me.
God is the one who is in charge..
God is sovereign, and we are simply stewards..
We are simply servants..
So maintain a constant spirit of humility.
to fear God, to know that He is the one in charge..
We were in Las Vegas, and walking on the.
Hollywood Walk of Fame, right?.
Have you ever been to Las Vegas?.
No?.
Maybe you should go, yeah..
- Wow..
- Maybe you should go..
Go to Las Vegas, you know, not to gamble,.
but to just walk around..
So we were in Hollywood Boulevard,.
and walking the Walk of Fame,.
and you see all the stars on the street, right?.
And with the famous people,.
their name engraved on the walk, right?.

$^{361}$I saw this one, Elvis Presley..
Ring any bell?.
Huh?.
Presley..
His daughter just died, just passed away last week..
All of a sudden, 54 years old..
You think about your life,.
think about how our life is like what James says..
It's like a morning mist..
We were in San Francisco for a trip,.
and it was a rainy day, it was cool and foggy..
San Francisco is always like that, you know?.
You go to the beautiful Golden Gate Bridge,.
just mist and fog will just come by,.
and then it will be gone..
This is exactly what James is talking about..
What did he say?.
Your life is like morning mist..
You cannot hold onto it..
You cannot control it,.
because it comes, and it will be gone..
We cannot control many things in our lives..
We can't even control the weather, right?.
People in Chicago should know.
that our weather is unpredictable, right?.
Even our economy, especially after COVID,.
is kind of unpredictable..
There's no way you can be in charge.
of every single thing, right?.
I saw this picture a couple days ago..
The best Valentine gift is to give eggs, right?.
Because the price of egg is going up and down.
in an unpredictable way..
Now, James is not against our human responsibility..
No, we are stewards of our time..
You are a steward of opportunities,.
resources God has given you..
But what he's confronting is the wrong attitude.
toward life, toward God..
What James is confronting is the wrong,.

$^{401}$misplaced attitude toward life,.
thinking that we can be in control..
And the Bible says, "Hold on, hold on..
"Life is unpredictable, and life is full of mystery..
"So you better recognize who is in charge,.
"and God himself is in charge.".
4.15, "What you ought to say is if the Lord wants us to,.
"we will live and do this or that..
"Otherwise, you're boasting about your own plans,.
"and such boasting is evil.".
And remember, it is sin to know what you ought to do,.
and yet not do it..
The right attitude toward life is to fear God,.
to acknowledge that God is in charge,.
he's directing your path, not you..
You simply follow..
Now, let alone James, it's like a lot of wisdom literature.
in the Old Testament, like Book of Proverbs.
and Ecclesiastes, wisdom literature..
Now, what is wisdom literature?.
Wisdom literature is telling us how to live our life.
under the will of God..
To fear God is the foundation of life wisdom..
To fear God is to relate to God properly..
And we relate to God properly, and in Christ..
Through Christ, we relate to God as our Father,.
and we are his children..
And one of the main points in the letter of James.
is that you have to figure out, you have to make sure.
you have the right relationship with your creator,.
your Lord, and you are the follower,.
and you follow the sovereign will of God..
Last Monday night, we had our English team,.
leadership team meeting together,.
and one of the conversations that we had was,.
we were all drivers, and we used to teach someone else.
how to drive, right?.
So we have kind of a sharing our own,.
you know, frightening experience, right?.
Teaching a new driver how to drive, right?.

$^{441}$Have you ever done anything like that?.
You know, you're an experienced driver,.
and you teach someone else how to drive,.
and you are in the passenger seat,.
and the new person is in the driver's seat, right?.
How do you feel?.
You feel out of control, right?.
You feel out of control..
Because, you know, you can only control the car.
with your voice command, right?.
And sometimes the driver, the new driver,.
doesn't listen to you..
So it was kind of scary..
But think about it..
If the driver is an experienced driver,.
someone you can trust,.
then your experience will be totally different, right?.
Have you ever seen, have you ever, you know,.
been in my car, right?.
Do you feel safe in my car?.
Oh, you feel dangerous, right?.
You feel dangerous, right?.
But a different kind of danger, right?.
Different kind of danger..
It's not because of inexperience, right?.
Because of experience, right?.
(audience laughing).
But think about it..
If the driver is totally trustworthy,.
then your experience will be totally pleasant, right?.
God is in the driver's seat..
God is the driver..
He is the creator, He is our loving Lord, right?.
He is taking control..
And we are in the passenger seat..
And when you can trust God totally,.
and He is indeed trustworthy,.
then we in the passenger seat,.
we will say and we will pray,.
take me according to your will..

$^{481}$And we will go to the destination, right?.
Yes, there were different turns,.
there were different corners,.
ups and downs during the journey..
However, the driver who is trustworthy,.
our God Himself will bring you to the destination..
That is what James is talking about, right?.
The sovereignty of God..
Again, it's not saying we do not need.
to take responsibilities of our own affairs..
You're a business owner, you have to be faithful.
because you're a steward..
You are a worker in your workplace,.
you are to be faithful..
You're a family member, you have to be faithful..
This is all about spiritual stewardship,.
human responsibilities..
However, we have to recognize that God is in control..
God is in control..
Let me close with a quote from Charles Spurgeon..
He's one of the well-known preachers 100 years ago..
He talks about the sovereignty of God..
He said, "When you go through a trial in your life,.
"the sovereignty of God is the pillow.
"upon which you can lay your head.".
It's comforting to know that the driver of our life.
is God Himself, not us..
But as we figure out the right relationship with God,.
we are in the passenger seat..
We follow His will and be humble always before Him.
and follow and seek the will of God..
Now, I hope this lesson will be your lesson as well..
As you go through the journey of life,.
probably not as crazy as the two weeks.
that I experienced with Southwest Airlines..
So thank you, Southwest Airlines..
But we all go through up and down, right?.
And twists and turns in our own life..
But do remember, do remember the sovereignty of God..
He is guiding..

$^{521}$And your lesson and my lesson is to obey.
and to seek and follow His will..
Well, let's all pray together..
And I want you to have a moment.
to bring your concern before our God, before your God..
What is it that you are thinking?.
What kind of decisions you are facing?.
Or maybe some uncertainties in your life..
Maybe you're praying for something.
or seeking something in your life..
Bring it, bring your concern before our loving God..
And then we'll close in prayer..
Dear God, we thank you because you are sovereign..
You are the one who is in charge..
All of our concerns, all our cares,.
we bring before you,.
knowing that your will is always the best will..
In our journey of life,.
we have ups and downs and twists and turns..
Yet we come before you knowing that your will is sovereign..
And you are not only sovereign, you are our loving God..
And you are the one who is guiding our path,.
giving us guidance and wisdom,.
and help us to be humble, to follow your will..
And we praise you because this is who you are..
And you are a great and sovereign God..
And help us serve you and worship you and follow your will..
In Christ's name, we pray..
Well, God bless you and I wanna see you next Sunday..
For the bag of blessing, you need to come early.
and have some time of fellowship together..
\newpage



\section{4:13-18}
\label{sec:3kSkaNorqlI}
\textbf{Ready for Christ's Return (1 Thess 4\_13-18) - Rev. Mark Chan}
\newline
\newline
連結: \href{https://youtube.com/watch?v=3kSkaNorqlI}{\texttt{https://youtube.com/watch?v=3kSkaNorqlI}} ~~~~ 語音日期: 2024-03-19
\newline
\newline
\hyperref[sec:3t6kNRmRV_s]{\small{< < < PREV SERMON < < <}}
~
\hyperref[sec:index_chronic]{\small{[返順時目]}}
~
\hyperref[sec:index_scriptual]{\small{[返順卷目]}}
~
\hyperref[sec:dON_DJyLmhA]{\small{> > > NEXT SERMON > > >}}
\newline
\newline
$^{1}$>> Thanks, it's always good to see you..
Why don't we say hi to your neighbors?.
Say hi and give them a fist bump or something..
Greet them in the Lord..
It's always good to worship together as a church..
This is what we are being called to..
One thing is very important that Muir announced is that the new devotion series on Apostles.
Creek..
All right, so if you don't know anything about Apostles Creek, I hope you will enjoy the.
devotional..
Every single day we're going to send it out to your chat group, to your social media group..
So please read that because in the following weeks I'm going to go through sermon series.
on the Apostles Creek, which is very, very important..
Apostles Creek is a summary of our faith..
How many of you have read through the Bible, the whole Bible?.
From cover to cover, from Genesis to, wow, that's great, that's great..
Two, all right, three, four, all right, are you sure?.
Four, five, okay, more hands, okay, six..
All right, how many of you have read through the entire Bible but you don't understand.
much of it?.
Yeah, I mean, right?.
We all struggle to get to know the Word..
We believe that the Bible is written for us..
It's for all generations of Christians, all right, even though it's not directly to us,.
but it's for us, it's written for us..
However, when you read the Bible, there are tons of things that we may find confusing.
or we just don't have enough background information to know about it..
However, the Apostles Creek is a summary to help us, it's like a big scale map to understand.
what the Christian faith is all about, all right?.
So it's very important, so make sure you go through the daily devotional on Apostles Creek,.
all right?.
So today, we're going to continue our series on the last days, the eschaton, what we call.
eschatology, right?.
So we're going to read the passage from 1 Thessalonians 4, but before we read that,.
let's look at this chart..
This is a very popular understanding of what is going on in the last days, what is going.
on before, right before the yellow arrow, the second coming of Christ..
This is very popular in our American culture especially, even in Chinese churches, even.
in our culture..
There are many, many people who promote this kind of understanding..

$^{41}$Now as you can see, the kind of the light blue line is the timeline from left to right..
So we're going ahead into this timeline, going into some sort of future, right?.
And the church is right there, we are right in between..
I have talked about the already and not yet concept..
Christ has already come..
He died on the cross and he rose again..
And then there are generations of churches, right?.
And before that yellow arrow, the really glorious second coming of Christ, most people will.
find that this kind of popular understanding is what you hear a lot of time..
Is that before the second coming of Christ, before the yellow arrow, before the glorious.
second coming of Christ, there will be seven years of tribulation..
There will be unprecedented suffering on earth..
And then, right before the seven years of tribulations, there will be something called.
rapture..
Something called secret rapture..
So Christ is coming in this scheme, in this kind of understanding, Christ is going to.
come two times..
The first time is in secret, alright?.
Suddenly in secret, nobody knows..
And then rapture the believers, rapture the church, disappear out of the world, okay?.
Out of the world and taken us to heaven, alright?.
And then the seven years of tribulation will kick in on earth..
And there will be wars and there will be tons of, you know, crises happening on earth for.
seven years..
And then after the seven years, after this crisis, Christ will come again..
That's the yellow arrow..
And then the new heaven and new earth will be launched, will be escort on earth..
Now this is, as I said, a very popular understanding of what is going on in our days, alright?.
Especially in American culture..
There are many people who talk about rapture..
You can see on the upper left is the left behind series..
This is a very popular novel series trying to put together a story of some kind to talk.
about the secret rapture and then seven years of tribulations..
Very popular and millions and millions of copies are being sold..
And even being translated into different languages, including Korean, including Chinese, and all.
sorts of Asian languages..
Very popular..
And also there are on the lower left is some sort of bumper stickers..
You can buy and talk about rapture..
In case of rapture, can you see that, the white one?.

$^{81}$This car will be unmanned because you will disappear according to the scheme..
And then at the lower center, warning in case of rapture, this Facebook page will no longer.
be updated because you will disappear, right?.
As the church, as believers, you will disappear..
And then on the right, on the background, there is a rapture ready index webpage..
You can go check it out..
But I advise you not to because this is a popular cultural understanding of second coming..
It may not be the right, correct understanding because when we read the Bible, we need to.
understand there is a background to every single passage, including the passages we.
are going to read, the so-called rapture passage..
So we're going to read this passage line by line together and see how it applies to us..
So I'm going to set it right out front..
I'm going to debunk this popular understanding..
There's no secret rapture of this kind..
So what this passage is all about..
Let's do this..
First Thessalonians chapter 4 verse 13..
We're going to start from there..
"Brothers, we do not want you to be ignorant about those who fall asleep or to grieve like.
the rest of men who have no hope.".
Now you need to understand when we read the letters to the churches, Paul's letters to.
churches or other letters to the church in the New Testament is always a reply email..
What do I mean by that?.
You do all, you all do emails, right?.
Let's say you get the reply emails only..
You don't get the original email..
But from the reply email, you kind of figure out what the original email is all about,.
right?.
You know what I'm saying?.
That's how you read the letters..
That's how you read First Thessalonians..
Paul wrote this stuff to the church in Thessalonica..
It's because he is trying to answer some of the questions that they have asked..
So the churches in Thessalonica, they have some problems..
They have some issues..
So they wrote to Paul or they have a messenger sent to Paul..
And Paul is writing First Thessalonians chapter 4, this passage as a reply to tell them what.
the answers are to address their questions..
So what he's saying is that there seems to be some vital concern for the early Christians.
in Thessalonica..

$^{121}$The early Christian church, they, of course, they had hope for the future coming of Christ,.
the future resurrection, right?.
Because Christ has already rose from the dead, right?.
Someday he will come according to his word..
And the return of Christ in glory, everyone in the early church, in all churches would.
believe that..
So before these things happen, some of the people in the Thessalonian church have passed.
away..
Some of those in Christ, brothers and sisters in Christ, they have passed away or some of.
those even martyr, died for their faith..
A few of them..
So obviously the Christians in Thessalonica were doubting because they have no hope..
What about these people?.
Would they miss, would they miss out the glorious coming of Christ?.
If he comes anytime, right?.
If human history will conclude with the end, the eschaton, right?.
God, God's son will come in glory..
Those of us who receive Christ, right?.
Will welcome Christ in glory..
What about those who have passed away in Jesus?.
What about those Christians, brothers and sisters who have passed away?.
Will they miss the glorious coming of Christ promised to us?.
Will they?.
And Paul is saying, absolutely no way..
Paul is saying, we do not want you to be ignorant about those who fall asleep..
Those who have passed away..
Now falling asleep is a metaphor, meaning dying in Christ..
So Paul is saying, absolutely no way..
Let's see how he answer the questions..
Verse 14, we believe that Jesus died and rose again..
So we believe that God will bring with Jesus those who have fallen asleep in him..
According to the Lord's own word, we tell you that we who are still alive, who are left.
to the coming of the Lord, will certainly not precede those who have fallen asleep..
Now what is he talking about?.
He said those who have fallen asleep, again it's a metaphor for dying in Christ..
They will not miss out the glorious second coming of Christ..
They will actually be the first in line to welcome Christ, to see Christ, to see his.
glory..
So the apostle is saying, the dead in Christ, they will rise from the dead first, followed.
by us, right?.

$^{161}$They will rise first..
They will be taken up into the air..
And we who are alive as coming, will also be taken up in the air to join the Lord, so.
called to meet the Lord in the air as Jesus is coming down in trumpet sound and glory..
Verse 16, for the Lord himself, he's talking about second coming, will come down from heaven.
with a loud command, with the voice of archangel, with the trumpet call of God, the dead in.
Christ will rise first..
After that, we who are still alive and the left, will be caught up together with them.
in the clouds to meet the Lord in the air..
So you got a picture, at the end of human history, when Christ comes again, it will.
be the appearance of a public appearance, a loud command, the voice of archangels, the.
trumpet call of God, the dead in Christ will rise first..
They will welcome Christ first..
Now it's not that Christ will come so far in the mid air, and then some of those will.
rise from the dead and meet him in the air, and then we turn back to heaven..
No as you see the context here, Christ is on the way from heaven to earth..
It's a one way street..
He is on the way coming, and then those who are on earth, those who are dead in Christ,.
they will rise first and will meet him in the air as he coming down..
So it's kind of being sent to rise and then to meet Christ in the air and then come back.
to earth..
So the imagery is this, we who are still alive, when the second coming happens, we will be.
with those who have died in the Lord..
We will be meeting the coming Christ, descending Christ as he is returning in glory..
Actually this is the technical term here, to meet the Lord in the air, to meet, to meet..
What does it mean?.
It is a technical term frequently used in the Roman Empire..
For a city, they will have a welcome committee, get ready and prepare to welcome a governor.
or a victorious general or the Caesar, the king, the emperor to come to town..
When emperor comes to town, when an important government official comes to town, the city.
will organize a welcoming party, a welcoming team..
They will go out in the midway to meet the emperor first, okay, to worship him, to meet.
him and then escort him back to the city, back to the town..
That is what it means to meet the Lord in the air..
So Paul is using a first century political scene to talk about the coming of the true.
king..
The true king is Jesus, is Christ himself..
Now when the Roman armies come back from a campaign, when the Roman armies go out to.
the east or west or north or south, when they conquer their enemies, they have taken some.
captives and they will return to the city of Rome and before they will enter the city.

$^{201}$of Rome, they will camp outside of the city, about a mile or so away..
And then they will send a messenger into the city and talk to the senate of Rome to announce.
the immediate arrival of the victorious king, of the general, of the victorious army..
Then they will give the city some time..
The city planners will put together some sort of people to decorate the city and to prepare.
for a parade for the conquering king, the hero for homecoming..
So this is the sort of the political scene that often happens in the city of Rome and.
many different cities in the Roman Empire..
And then after the planning, there will be a trumpet blast being made..
And then the signal will be for the armies to march in triumphant into the city..
That's the meeting committee, the welcoming committee, to welcome the king into the city..
So the citizens of Rome sometimes are being invited to join the welcoming committee, to.
be invited to go outside of the city, outside of the city wall, to join the parade and march.
back through the arch of triumph with the victorious army..
Some of the officials, some of the high status citizens will participate in this kind of.
victory march..
Now this sort of things actually also happen in our days..
When I first came to Chicago, it was about 20 some years ago, some of the locals who.
know all the local history here, they told me some interesting stories..
And I remember this part, we were driving in South Loop..
He told me this story, a local friend..
He said in the mid '90s, in the mid '90s, those of you who are old enough to remember.
this, right, in Chicago, mid '90s, when President, ex-President Clinton came to Chicago, he landed.
in Grand Park with a military helicopter in the mid '90s..
And then there were a party of people welcoming him..
In order to prepare for his arrival, the city of Chicago actually renovate part of Roosevelt.
Road..
Have you been there?.
There was a stretch of Roosevelt Road nicer than the rest of it..
You know why?.
My local friend told me it was because the city, in order to prepare for the arrival.
of the ex-President, arrival of Clinton, they renovate this entire stretch of Roosevelt.
Road so that when he comes for the National Convention, there will be honor, there will.
be a glory, right?.
It's kind of a celebration..
Do you know that?.
You didn't know that, right?.
You pay for it..
(Laughter) Now this is a sort of the political scene.
Paul was using..

$^{241}$And he was obviously using the first century political scene from the Roman Empire and.
trying to tell you who is the true king and how we should prepare ourselves for his second.
coming..
Now you can see throughout Paul's writings that he often borrows or subverts the Roman.
political images and apply them, subvert them, apply them to Christ..
Because Christ is the only true king, not Caesar..
Now you can imagine Paul is trying to use that scene to talk about the arrival, second.
coming of Christ in a horizontal way..
Like Christ is like the victorious general, the victorious king marching toward the city..
And then we as a city, we will prepare ourselves for the arrival of the true king..
We will go out midway to meet him, right?.
And then the first in line to meet him would be those who are already dead in Christ..
Now in the ancient days, when you die, your tomb is often outside the city wall..
You got it?.
Because it's unclean, right?.
It's unclean..
It's always outside the city wall..
So you go out of the city, out of the city gate, you look around, there are tombs, right?.
That is the picture in Jerusalem, there is a picture in Rome, there is a picture in a.
lot of ancient Mediterranean cities..
The tombs are outside the city wall..
So you can see what Paul is doing..
When the king is coming toward the town, the first to meet the victorious king of course.
are those who are dead already..
They will rise first..
They will be the first in line to meet the true king..
And then we who are still alive within the city, we will join the welcoming committee..
And we go out midway to meet Christ..
Now that's the horizontal picture..
Now I want you to turn it 90 degree, vertical..
That is the picture of the second coming of Christ..
Of course Christ is not coming horizontally..
Christ is coming vertically..
He is from heaven..
He is from above..
So when Christ comes in the time of second coming, midway, those who are dead in Christ,.
they will rise first, be the first to meet Him in midair..
Not to go away and disappear..
But as He is descending, as He is coming on earth with glory, we rise with those who are.
dead in Christ, will come in midway to meet Him and come down..

$^{281}$You see what I'm saying, right?.
We will be joining the welcoming committee to welcome the true king..
Now as you can imagine, Christ coming in glory from heaven..
Those who die in the Lord will go up and meet Him in the air..
That is the second coming Paul is talking about..
Now why does he talk about these things?.
As I told you again and again, in this series when we talk about the book of Revelation.
or the second coming of Christ, eschatology, eschaton, the end of time, the last days..
You need to understand in the New Testament when it talks about the second coming of Christ,.
it's not some abstract idea..
It's not for us to, okay let's figure out, are there any things happening, what kind.
of terrible things will happen, is it a weather forecast or something like that..
It's not abstract idea..
It is not newspaper written ahead of time..
It's actually a very practical teaching..
Eschatology is always practical..
You know the end is always practical..
Now why do I say that?.
418, remember Paul is trying to answer a question from the Thessalonians questions..
What kind of questions did they have?.
Some of them are martyrs..
They passed away in the Lord, right?.
And they're really really sad..
They're really heartbroken..
Now Paul is trying to say to them, they will not miss out the glorious coming of Christ..
In fact they will be the first..
They will be the first to see the glorious kingdom, right?.
Those who are dead in Christ, they will rise first..
They will welcome Him first and followed by us..
418, he says, "Therefore, encourage each other with these words.".
How can these words, the things that we talk about, the second coming of Christ, the coming.
to meet error and welcome Him back..
How can we encourage one another with these words?.
You remember at the beginning he says, "We grieve for those who passed away in the Lord.".
We grieve..
Our brothers and sisters, our family members, those who are in Christ, those who are Christians..
They passed away before we do..
Of course we will grieve, right?.
Tomorrow I'm going to a funeral, a Christian funeral..
Of course we will grieve..

$^{321}$We will weep..
Because such a loss to all of us, right?.
The death of your loved one is the death of part of you..
It's taking away part of you, a big part of you..
We will grieve..
But Paul is saying, as we grieve, we don't grieve with those who are without hope..
We have a hope that transcends death..
That transcends all of our tears..
Because someday Christ will come and make things right..
When the new heaven and new earth come down from heaven, Revelation, the last two chapters,.
there will be no pain and no death and no tears..
There will be no evil..
That is the hope that transcends everything else, even our final enemy, death itself..
I always say when I officiate a funeral, I always say in a Christian funeral, we don't.
say goodbye to our loved one..
What do we say?.
We say, "See you later.".
That is the glorious hope we have in Christ..
So don't forget that..
Don't forget that..
Secondly, eschatology is practical because it help us to put things in perspective..
Put our lives and our lifestyle, today's work, today's how you live, put things in.
perspective..
Why?.
Because we know there's a final destination we're going to..
And Christ will come again..
The new heaven and new earth, the new kingdom, the future kingdom, is the country that we.
are going to..
I always tell Christians, we have two passports..
The first passport is our earthly passport..
Maybe China, maybe UK, maybe American passport..
That is one of our passports..
The other passport is the kingdom of God passport..
Now if you compare it side by side, the earthly passport will expire..
You have to renew it..
And sometimes you cannot renew it..
You got to change it..
That is your first passport..
Your national identity..
Where you belong on earth..

$^{361}$It will expire..
It has expiration date..
Did you check that?.
Your passport?.
It has expiration date..
Your early citizenship will expire..
But our kingdom of God passport will not expire..
It's eternal..
Okay, you don't have to check it..
But what you do need to check is that where you are living up to, that passport, that.
identity..
We are not so much living for our earthly identity..
No matter what nationality you belong to, it will pass away..
You know that?.
It will pass away..
It really doesn't matter at the time of funeral..
What will last forever is your heavenly citizenship..
That you belong to God..
You belong to true king, Jesus himself..
So it's practical in that sense we need to remind ourselves..
Are we really the heavenly citizen?.
So as you can tell, there were a lot of passages we have worked through in the past many sermons.
of this series..
Be it book of Revelation or book of Matthew or other passages..
When Jesus talks about the last days, when Paul talks about the last days, he always.
talks about how we live today..
How we live in light of the future..
How we live in light of the coming kingdom..
How we live in light of that future..
Remember the gospel of Matthew?.
He talks about, Jesus talks about his coming..
It's a list of signs, right?.
There are many things that will tell you his coming is imminent..
His coming anytime..
And that's always followed by some sort of practical teaching..
The parables..
Remember the parables of ten bridesmaids, right?.
Five of them prepared the oil..
Five of them did not prepare for the oil..
So when the groom comes, finally they have to go, again the same picture..

$^{401}$They go halfway out to meet the groom, right?.
To join him and escort him back to the banquet..
Five of those who are prepared, they can enter the banquet..
The five of those who thought they have prepared, but they did not..
So they could not go into the banquet..
The master said, "I don't even know you.".
That's the first parable..
The second parable is the talent parable..
With the master going away, and give his talents to his three servants..
Five thousand, two thousand, and one thousand..
Those who have five thousand and two thousand, they took the risk and invest and increase.
the asset of their master..
And they are prepared when after a long time, after a time of delay, finally the master.
comes back..
Those who had, he who had the five thousand and he who had the two thousand, they are.
prepared to give an account..
Because they did serve the master, increase his asset..
The last one, with the one thousand, didn't even bother to take the smallest risk to increase.
the master's asset..
So always, when Jesus talks about the last days, always reminding us how you should live.
today..
Are you living in light of that future?.
Are you living worthy as a citizen of heavenly kingdom?.
You know, whatever the kingdom characterize should be our characteristics..
If the future kingdom is full of love, then our life should live in light of that, right?.
We will practice love..
If the future kingdom is nothing but justice, but righteousness, then our lives, our decisions,.
our choices should reflect that set of values..
Right?.
Live according to the future..
Some people said, and I'll close it with this thought..
Some people said questions are actually aliens, right?.
Aliens, not the monsters, okay?.
Foreigners, okay?.
Where are we from?.
We are from the future..
Because we are heading to the future, we belong to the future..
And we bring stuff from the future to the present..
We are aliens in that sense..
We live differently..

$^{441}$Why?.
When people look at your lifestyle, look at your marriage, look at your family, look.
at how you work, look at how you make decisions, they will wonder why you are different..
Well, you are different because you belong to the future kingdom..
You don't belong to the earthly kingdom..
You make decisions differently..
You set a priority differently..
Why you come to worship today?.
Instead of, you know, just get a couple more hours of sleep, right?.
Because your decisions is affected by your hope, by your future..
Let's do that..
Share the gospel..
You see the gospel is a huge thing..
Because it talks about, it affects your eternal destiny..
Live according to the future kingdom..
Prepare yourselves to meet the Lord..
Alright, one question, final question you need to ask yourself..
If Jesus comes today, or if you are being called back home to see, one cosmic question.
you need to ask yourself is that will God see you and you can give an account?.
And Jesus will say, "You are my faithful and wise servant..
You can enter the kingdom.".
Right, that is a cosmic question for you this week..
Live according to the future..
Let's all pray together..
Father, we give you thanks because of your grace you have called us into your kingdom..
It is your grace we have been called your children..
And I pray that each and every one of us will prepare ourselves knowing that our identity.
belongs to the future kingdom..
That we need to prepare ourselves to give an account to you, Master..
To serve you, to increase the value of the kingdom, to build up your church, to build.
our lives so that you get to grow..
And I pray that when you come, you ask us whether we are being faithful to what you.
have entrusted to us..
And we'll say, "Yes, we have been faithful, Master.".
So pray that we'll grow spiritually in this way, to live according to our future hope..
To live according to the future eschaton..
Help us to live as your citizens, as your heavenly citizens..
That when Christ comes back, we'll be able to give an account..
We praise you week after week..
We strive to be a better citizen of your heavenly kingdom..

$^{481}$And pray that you will continue to transform us with your grace..
That we'll be more like your son..
In Christ's name we pray..
Well, thank you..
\newpage



\section{馬可福音 12:38-44}
\label{sec:dON_DJyLmhA}
\textbf{耶穌最後的一週 \#2 - (馬可福音 12\_38-44) - 陳崇基牧師}
\newline
\newline
連結: \href{https://youtube.com/watch?v=dON_DJyLmhA}{\texttt{https://youtube.com/watch?v=dON\_DJyLmhA}} ~~~~ 語音日期: 2024-03-19
\newline
\newline
\hyperref[sec:3kSkaNorqlI]{\small{< < < PREV SERMON < < <}}
~
\hyperref[sec:index_chronic]{\small{[返順時目]}}
~
\hyperref[sec:index_scriptual]{\small{[返順卷目]}}
~
\hyperref[sec:55R6ROPsg74]{\small{> > > NEXT SERMON > > >}}
\newline
\newline
馬可福音 12:38-44
\newline
\begin{longtable}{cl}
\hline
\hline
章節 & 經文 (和合本修訂版)\\
\hline
12:38 & \begin{tabularx}{0.7\textwidth}{X} 他在教導的時候，說：「你們要防備文士。他們好穿長袍走來走去，喜歡人們在街市上向他們問安， \end{tabularx} \\ \\ \relax
12:39 & \begin{tabularx}{0.7\textwidth}{X} 又喜愛會堂裡的高位，宴席上的首座。 \end{tabularx} \\ \\ \relax
12:40 & \begin{tabularx}{0.7\textwidth}{X} 他們侵吞寡婦的家產，假意作很長的禱告。這些人要受更重的懲罰！」 \end{tabularx} \\ \\ \relax
12:41 & \begin{tabularx}{0.7\textwidth}{X} 耶穌面向聖殿銀庫坐著，看眾人怎樣把錢投入銀庫。有好些財主投了許多錢。 \end{tabularx} \\ \\ \relax
12:42 & \begin{tabularx}{0.7\textwidth}{X} 有一個窮寡婦來，投了兩個小文錢，就是一個大文錢。 \end{tabularx} \\ \\ \relax
12:43 & \begin{tabularx}{0.7\textwidth}{X} 耶穌叫門徒來，對他們說：「我實在告訴你們，這窮寡婦投入銀庫裡的比眾人所投的更多。 \end{tabularx} \\ \\ \relax
12:44 & \begin{tabularx}{0.7\textwidth}{X} 因為，眾人都是拿有餘的捐獻，但這寡婦，雖然自己不足，卻把她一生所有的全都投進去了。」 \end{tabularx} \\ \\
[1ex]
\hline
\hline
\end{longtable}
$^{1}$大家午安.
這次我們繼續聖周的系列.
我們看到耶穌最後一個星期.
上一節我們說.
由榮晉耶路撒冷開始.
我們曾經給大家看過最後一週的時間表.
這次第二節我們說到大概是星期二的時間.
星期二的時間.
我們上次說過榮晉耶路撒冷之後.
耶穌進入聖殿去結證聖殿審判聖殿.
然後再進入耶路撒冷.
大概是星期二的時候.
他就和宗教的權貴有很多的爭辯.
我們要處理的經文是很多的.
有整大段在中間.
所以我們不能夠很仔細地讀.
但是正如我所說.
有時候我們zoom out一點.
用一個宏觀一點的角度.
來看一下整段的聖經在說什麼.
和耶穌與宗教的權貴爭辯之間.
到底是有什麼在談或爭執.
所以我們zoom out一點.
我們就發現.
11章27節耶穌的位置在哪裡.
進入聖殿.
在聖殿裡面行走的時候.
然後祭司長,法利賽人和民事長老.
就去和他辯論和爭辯.
然後有六次耶穌去搏鬥.
當時耶路撒冷的宗教權貴.
大家有六次的對峙.
大家彼此去問答.
彼此去挑戰大家.
或者挑釁大家.
是一些很hardly debated的episode.
第一次就是耶穌去搏鬥祭司長和民事長老.
他們說的內容大概就是權兵的問題.
到底你憑什麼權兵來審判聖殿.
去挑剔我們或者批評我們呢.

$^{41}$耶穌的對答就有些很有趣味.
然後他答完之後.
搏鬥他們之後就說出一個比喻.
那個比喻可以叫做邪惡原獨的比喻.
這個比喻也是挑釁他們.
去挑戰他們.
他們就越發地去憎恨耶穌.
然後第二次的辯論.
第二回合就是說到納粹的問題.
你記得那個錢幣.
耶穌說有該殺的形象的.
就歸給該殺.
等等.
這個整個的辯論就是耶穌搏鬥.
法利賽人和希律黨的人.
這個很有趣.
因為法利賽人就不喜歡納粹.
因為法利賽人是在求上帝的國來到.
我們有自己的獨立.
為什麼要給稅給羅馬帝國呢.
希律黨就不是了.
希律黨就是希律的家族.
或者王族剩下來的人.
他們是親羅馬的.
兩班人居然是聯合去試探耶穌.
就是問他我們向羅馬帝國納粹.
可不可以呢.
結果耶穌也搏倒他們.
第三次的辯論就是耶穌和薩都哥人.
薩都哥人當時只是信五經.
創出黎文新.
他們不信復活的.
所以說到復活的問題.
耶穌就搏倒他們說你們看錯了.
因為上帝是一個活人的上帝.
亞伯拉罕以撒瓦戈是一個的神.
這個神是一個活人的神.
所以你們不要忘記上帝的工作.
是超過你們人所想的.
或者人所能夠觀察到.

$^{81}$我們理性能夠去追溯到的領域.
上帝是可以使人復活的.
因為他是活著的神.
第四次的辯論就是他搏倒一個民事.
說到什麼是最大的解密.
第二解密.
然後十二章三十四節.
就從此以後沒有人再敢問耶穌什麼.
所以他四次的辯論裡面.
他就搏倒他的敵人.
然後第五次所謂辯論.
他就去解釋基督是大衛之子.
你記得那裡引用了一段經文.
就是詩篇第110篇.
主對我主說.
你坐在我的右邊.
那裡說到大衛之子君王要嫁到的時候.
這個以前的君王.
那個最好的君王大衛.
他在上也有一個主.
這裡就解釋基督是大衛的子孫.
大衛的後裔.
第六次就是去譴責民事的偽宗教.
所以他一次又一次告訴他們.
他們需要去到悔改.
他們要被審判.
因為他們是愛權兵等等.
然後最後他就說出一個故事.
就是寡婦獻上民前的奉獻.
其他人有錢有勢的人.
都是正如下來他們才奉獻的.
這個整段的聖經.
六次的爭辯.
六次的debate或者controversies.
去到駁倒當時的宗教的權貴.
然後十三章第一節.
他就離開聖殿了.
然後去到聖殿對面的橄欖山.
他就宣佈聖殿.
在一代之內.

$^{121}$就是主後七十年.
就會被毀的.
所以整段六個的爭辯裡面.
和宗教權貴的爭議.
都是發生在聖殿當中.
其實都make sense的.
因為第十一章.
他進入耶路撒冷的時候.
永進耶路撒冷.
他就進入聖殿了.
然後就是結正聖殿了.
就結束聖殿的祭祀敬拜.
在基督裡面敬拜.
不再是一個地方來敬拜.
然後這六次的爭辯十二章.
就是講到結束耶路撒冷.
宗教的權貴.
結束他們的leadership.
不需要有他們在這裡.
因為聖殿已經要過去.
聖殿的權柄.
聖殿的權貴.
也都會過去.
我們會選一幾段的經文來看一看.
在這一大段裡面.
但我有兩個重要的point我想講.
就是寫這樣的故事留給初代教會.
為什麼呢?.
你可以想像馬可.
或者其他福音書的作者.
都是大同小異在講這些東西.
審判聖殿結成聖殿之後.
就是和宗教的權貴在耶路撒冷有很多的爭辯.
然後結束他們的領導.
寫給讀者.
第一代讀者.
當他拿起四福音的時候.
就是初代教會的人.
差不多他們已經有一代的時間了.
為什麼要寫這樣的故事呢?.

$^{161}$為什麼要和他們憶想.
講這些宗教的權貴.
和耶穌之間的爭議呢?.
有什麼警告呢?.
初期的教會.
會不會可能都有這樣的危機呢?.
所以在這一大段裡面.
我們抽一些經文的時候.
我有兩大點.
第一大點.
就是宗教政治的本相.
在你的經文裡面.
告訴我們什麼叫宗教政治的本相.
第二大點就是.
如何防備宗教政治.
或者教會政治.
第一點.
宗教政治的本相.
什麼意思呢?.
我們看一段.
就是開頭的爭議.
第一段.
十一章第二十七節.
我們先讀經文.
他們就是耶穌和門徒.
去到耶路撒冷.
耶穌在聖殿裡面行走的時候.
祭司長問事.
和長老進前來問他.
你將著什麼權柄.
做這些事.
給你權柄.
做這些事的是誰呢?.
可能是說之前.
一大堆他所做的事.
或者說的事.
耶穌說過什麼?.
做過什麼?.
就是在聖殿的商場的地方.
他去推翻那些兌換銀錢的桌子.

$^{201}$或者桌子等等.
去審判聖殿.
你將著什麼權柄.
做這些事?.
二十九節.
耶穌對他們說.
我要問你們一句話.
你們回答我.
我就告訴你.
我將著什麼權柄.
做這些事?.
若翰的洗禮.
是從天上來的.
還是從人間來的?.
你們先回答我.
他們就彼此相依.
如果你說是從天上來的話.
耶穌一定會說.
為什麼你不相信他呢?.
但是如果我們說.
若翰的洗禮.
是從人間而來的.
又怕眾人.
因為眾人.
所有民間的人.
都認為若翰確是先知.
於是他們就回答耶穌.
我們不知道.
耶穌就說.
你不知道嗎?.
我也不告訴你.
我不告訴你.
我將著什麼權柄.
來做這些事.
這裡不是我們平時.
很多時候查經我們有個錯覺.
我們覺得耶穌迴避問題.
我們查經問大家.
有什麼意見的時候.
他說你覺得呢?.

$^{241}$你又問你呢?.
你有什麼意見?.
不是這個意思.
而是這個對答.
其實是當時猶太人的拉比.
和他們的學生.
那種對答的方法.
是用問題來回應問題.
或者用問題來作回答.
所以其實耶穌已經回答了.
回答了他的權柄在哪裡來.
就好像施首約翰一樣.
只不過你們不承認而已.
所以他將問題問出的時候.
其實就是在回應一件事.
就是他們最終不是權柄在哪裡來.
而是他們信與不信的問題.
他們肯不肯服的問題.
既然施首約翰的權柄.
是從天上來的.
我的權柄當然就是從天上來.
但如果我告訴你.
我的權柄從天上來的話.
你不信服的話.
告訴你也沒有用.
所以是一種拉比和學生.
經常有一種這樣的回答.
或者對答的風格.
但在這裡經文講清楚.
他們有什麼考慮.
或者什麼政治考慮.
法理賽人.
祭司蔣文士.
為什麼不信服施首約翰.
或者耶穌從天上來的權柄.
31至32節.
講什麼.
卻又怕眾人.
就是怕自己民間的地位會不保.
會失去民間的地位.

$^{281}$你想想祭司蔣文士,法理賽人.
他們這班人就是掌握猶太人整個的社會.
整個民族最重要的命脈.
是他們去掌管的.
是他們去掌握的.
那個命脈是什麼.
就是聖殿.
整個聖殿的體制.
聖殿的機制.
都是在他們的手中.
如果民間不能得到他們的尊重的話.
他們地位不保的時候.
他們就會坐不穩.
對於我們現代人來說.
我們看這些經文的時候.
因為我們有一種錯覺.
我們覺得宗教信仰.
是和我們公共的生活是分開的.
信仰是一些很私人的東西.
我們就說個人救主.
只是我.
耶穌只是我.
和我的公共生活.
外面的生活是沒有什麼關係.
信仰很多時候我們淪為.
只是一種個人修養.
敬拜禱告.
純粹是一些信仰宗教.
私人的空間或領域.
政治,經濟,社會,公共的生活.
是那些世俗的東西.
Holy,Sacred和Secular.
Sacred和Secular我們分得很開.
有什麼是屬靈的東西.
有什麼是屬世的東西.
我們全部都分開的.
工作和信仰可以完全沒有關係.
我們有一種錯覺.
所以我們看經文的時候.
我們就會說.

$^{321}$聖殿只不過是宗教的建築物.
所以河水不犯井水.
聖與俗.
大家兩樣東西可以分得開.
我們有一種錯覺.
其實這種錯覺是不對的.
亦都是不可能的.
不可以將我們聖與俗的東西.
分為二元的.
大家河水不犯井水.
因為聖經裡面的信仰.
信耶穌的.
這個信仰是涉及我們人生每一大小部分.
沒有一樣東西是信仰救不了的.
當你信耶穌的時候.
我們所有的東西都會重新再結構過.
是耶穌作主.
祂在中間.
然後其他東西.
無論我們私人的生活.
關係.
家庭.
工作.
我們的政治生活.
我們的社會的生活.
公共的生活等等.
都是因為基督已經作主.
天國要從你的生命裡面流露出來.
所以全部的生命.
每一部分都會被觸及被牽涉.
所以不可以這樣分聖族.
這個只不過是我們現代人的錯覺.
回到耶穌的時代.
其實都是一樣.
宗教.
經濟.
政治.
是彼此之間分不開的.
聖殿不是一個純粹.
宗教性的建築物.

$^{361}$而是宗教.
政治.
經濟.
的合一.
所以你看聖殿的時候.
就不是純粹一個寺廟那麼簡單.
不是一個你去那個地方.
然後有祭祀獻上.
不只是一個宗教性的建築物.
也是一個市場.
叫做Start Exchange.
就好像今天去Chicago Start Exchange.
New York Start Exchange.
或者.
聖殿是一個政權的地方.
好像City Hall那樣.
市政府.
或者是一個中央銀行.
一個銀行.
是全世界.
整個猶太人的世界.
所有的奉獻的資金.
都會流向聖殿在耶路撒冷.
所以你說聖殿是多麼複雜的東西.
多麼大的一件事.
對不對.
我記得好幾年前.
自己在那裡閒讀一些書的時候.
我讀到一個.
就是兩約中間.
舊約之後.
然後新約之前.
那幾百年.
有一些猶太人的著作.
我就純粹因為.
想休息一下.
我就想看看這些東西.
我看.
其實很好看.
我告訴你,報紙比上網的好看.

$^{401}$我打開這個瑪嘉比爾書.
有一本書叫做Second Maccabee.
裡面提到.
他說聖殿裡面的庫房.
那些資金流到聖殿.
或者奉獻到聖殿的時候.
無論是一些能吃的東西.
還是一些金,銀等等.
這些的奉獻.
就會去到聖殿裡面的chamber.
一個庫房的地方.
那裡就存放人民所有的奉獻.
目的是做什麼.
目的是要去救濟孤兒寡婦.
社會裡面真的有需要的人.
就是用聖殿庫房裡面的錢.
來給他們救濟.
其實是一個社會福利的系統.
當時是瑪嘉比爾的時期.
瑪嘉比爾就是耶穌之前大概150年的時間.
按照那個估算.
那段經文說.
瑪嘉比爾書有第三章.
那個講庫房裡面.
是有400個他連德的銀子.
200個他連德的金.
你說多少.
換算一下的話.
就是300多萬美金.
就是在耶路撒冷的聖殿.
300多萬美金.
你說不是很多你也有.
300多萬美金.
你查一下你的錢包有沒有.
現在說百萬富翁用什麼.
看你是不是億萬富翁.
300多萬美金.
在耶穌之前150年.
很多錢你知道嗎.
你今天一公斤有4000元.

$^{441}$你記得我剛來美國的時候90年代.
9號9日.
你這樣換算一下.
你想像不到那種錢銀是有多少.
可想而知如果大祭司.
他們管經濟的命脈的話.
如果他們有誠信.
百姓應該是無憂.
但如果你敲詐的話.
那些無依無靠的人就苦不堪言.
第二段文獻或者第二段文字.
就是約瑟夫寫的.
約瑟夫講到在主後大概60年代的時候.
那時候有一個反羅馬的第一次的暴動.
革命.
那些猶太武裝的分子.
他們去到反抗那時候的貴族.
或者親羅馬的人.
耶路撒冷的權貴.
講到那些革命分子衝入耶路撒冷和聖殿的時候.
他們就去到將那些大屋.
或者那些黃色的物業來到墳位.
然後有一句話約瑟夫就講.
他說之後這些革命分子就去到儲藏檔案的地方放火.
然後急忙的去銷毀所有債權人的合約.
從此他們就解決了這個負債的問題.
而這樣做也可以獲得曾有債務的百姓支持.
又能夠說服那些貧民加入他們的革命.
不再怕反抗這些富戶.
這樣守衛黨庫的人就逃離了.
他們就將這些紀錄去焚燒.
就是走進去中央銀行燒掉所有的資料庫.
燒掉所有的伺服器.
全部紀錄就沒有了.
這些革命分子就是想做這些事情.
但我想知道這段經文背後是在講什麼.
那些百姓被敲詐.
那些百姓有很多的債務在身.
是不公的債務在身.
這些資料和聖經有些不謀而合.

$^{481}$因為耶穌結證聖殿的時候.
不是純粹結證一個宗教的建築物.
他說聖殿是賊窩.
賊窩是什麼呢?.
你偷了錢就放在那裡.
然後用聖殿的外殼包著.
我們很剛鮮的.
聖殿是賊窩.
另一段文字也是我閒讀的時候看到的.
也是約瑟夫說的.
這裡在講什麼呢?.
他說聖殿北邊有一個駐軍的地方.
叫做安東尼亞堡.
裡面是存放大祭司的禮服.
這件禮服只是在大祭司獻祭的時候才穿的.
希律王就是大希律.
將這件制服放在安東尼亞堡.
到他死了之後.
這件禮服就交由羅馬人來管理.
就是說每逢大時大節的時候.
祭司就要穿這件禮服.
特別是雍歇波的時候.
贖罪日的時候.
他進去至聖所去打理一下.
他就需要穿這件禮服.
這件禮服就不在他手裡面.
是在羅馬人的手裡面.
是在皇的手裡面.
所以有很多政治的考慮.
有很多政治的張力.
如果得罪羅馬人會怎樣.
禮服你都不能穿.
禮服不能穿.
你怎樣做民間的禮儀.
怎樣可以說服民間.
所以一定要討好他和要親他.
這裡講到耶穌最後的一周.
耶穌進入耶路撒冷之後.
其實有很多的經文.
四福音裡面都有提到.

$^{521}$和我們剛才所看的.
不在聖經裡面的經文.
都是不謀而合的.
約翰福音第十一章.
耶穌進了耶路撒冷之後.
約翰福音就很有趣.
他說於是.
來到內漢瑪利亞的猶太人中間.
有很多人就看到耶穌所行的奇蹟.
他就信了他.
剛剛拉撒路.
從死裡復活.
然後他進入耶路撒冷.
然後有一些猶太人就信了耶穌.
但其中也有人去見法利賽人.
去告訴法利賽人.
耶穌所做的事.
他告訴法利賽人.
祭司長和法利賽人就召開議會.
他說這人耶穌.
行了好些神蹟.
我們應該怎樣處理他.
若讓他這樣做的話.
人人都信他.
整個平民百姓都信他的時候.
羅馬人就會來毀滅我們的聖殿和我們的民族.
這個大不.
這個真的很嚴重.
因為我們的地位就會不保.
其中有一個叫做戈亞法就是當年的大祭司.
他說你們什麼都不知道.
你們真是傻的.
你想想.
有沒有想一想.
一個人如果替百姓死.
免得我們整個民族滅亡.
是對你們是有利的.
這些是政治辦公室裡面說的東西.
是不是.
如果將他殲滅的話.

$^{561}$將耶穌殲滅.
就一定對我們是有利的.
如果不將他殲滅的話.
我們就會很嚴重.
因為我們沒有了聖殿.
沒有了我們的地位.
沒有了我們的權勢.
我們整個民族就會滅亡.
其實不是整個民族滅亡.
是他們滅亡.
我想看大點是說什麼.
就是宗教政治的本相.
人生一定有政治.
在生活裡面你想想.
大把政治.
說到權宜的張力.
或者關係上的矛盾.
或者問題等等.
你和我都遇見很多不同的政治.
大大小小.
美國兩黨的政治.
市政府的政治.
辦公室的政治.
office politics.
家庭都有政治.
因為我們總是會有不同的張力.
但我想告訴你.
最差的一種政治.
是教會政治.
是宗教政治.
最差最壞的一種政治.
是宗教政治.
religious politics.
因為這一種的政治.
是in the name of God.
贊成嗎.
是有神聖的明字牽涉的.
in the name of God.
我要這樣做.
in the name of God.

$^{601}$我要這樣保住自己.
所有的政治.
人間裡面.
最差的一種政治.
我不敢說是影響最大的一種政治.
但最差的一種政治.
就是in the name of God.
這是宗教政治的本相.
當然有錢.
本身不一定有問題.
但如果有錢有權有勢.
而沒有平衡的話.
沒有check and balance.
就會成為很大的問題.
因為這種是宗教政治.
我相信.
去到列舉這些爭辯的時候.
耶穌和聖殿和聖殿裡面的權貴.
寫給初代的教會.
不是單單要憶想.
上一代他們發生什麼事.
而是警告他們現代.
他們需要防備什麼.
所以福音書裡面.
你記得耶穌曾經說過一句話.
在同門.
他說你要防備什麼.
法利賽人的孝.
East of the Pharisees.
法利賽人他們那種.
愛權勢的那種孝素.
馬克方第十二章.
三十八字.
說到如何防備.
應該防備法利賽人的孝.
耶穌在教導人的時候.
這個by the way是六次的爭辯裡面.
其中最後一段.
他在教導的時候.
你們要防備民事.

$^{641}$他們愛穿長袍走來走去.
喜歡人們在街市上.
向他們問安.
又喜歡會堂裡面的高圍.
宴席裡面的手鐲.
他們侵吞寡婦的家產.
假以作狠詳的禱告.
這些人要受更重的懲罰.
Worst kind of politics.
In the name of God.
為了自己的地位.
為了自己的權勢.
我們將他用聖殿包裝.
用敬虔包裝.
這不是最差的一種politics.
耶穌說你們要防備.
你記得我剛才在大綱的時候說到.
這裡之後就有一個故事.
耶穌見到寡婦將他一生的奉獻.
放在聖殿.
其實那個奉獻不是鼓勵我們.
我們要奉獻這麼簡單.
寡婦的奉獻就是審判聖殿.
那種腐敗.
怎樣防備.
我記得好幾年前看了一篇文章.
是一個心理學家寫的.
他就講到宗教心理學.
宗教信仰的心理學.
人很自然和一些人事物物.
是會建立一些情意結的.
很容易.
特別是一些神聖的東西.
所以那個心理學家就寫了這件事.
他說為何聖殿這麼吸引.
因為聖殿裡面曾經發生過一些事情.
或者你獻祭的時候做過一些事情.
感受過一些事情的時候.
你就很容易有一種attachment.
如果聖殿要過去.

$^{681}$就好像人生要過去一樣.
很大件事.
原因是因為有一種情緒的.
感情的情意結.
連結在那種人事物物之上.
這個我覺得是值得我們去想的.
會不會有一些情意結.
我們是建立了.
而我們不知道的時候.
我們就假作是以為是神聖的東西.
或者是一些不能夠取締的東西.
在美國南邊.
Deep South.
有一些白人的教會.
每間都是過百年的老店.
歷史悠久.
他們今天還是在用King James Version.
英文的聖經翻譯版.
King James Version.
我有一個教授是教舊約的.
有時候去這些地方講道.
那些執事會就問他.
你用哪個版本.
是不是用King James.
那個教授就說.
我什麼版本都不用.
我是希伯來文的.
有些東西我們習慣了.
然後有一種情意結建立了的時候.
我們就覺得是不能不變.
或者年輕的時候.
在一些夏令會唱過的詩歌.
覺得很感動.
或者改變了一生.
成為了人生裡面的金曲.
勁歌金曲.
到你長大了.
老了的時候.
你又不會學新歌.
你永遠都是憶上那種敬拜的詩歌.

$^{721}$你就覺得只能夠接受那一套的詩歌.
這些就是那種情意結建立.
或者某些教會因為架構的緣故.
或者傳統的緣故.
可能有些牧者.
或者有些長輩.
可能會退了.
但是有榮譽的下午頭.
他仍然可以在這裡.
我其中一個Mentor.
有時候跟我談電話的時候.
他就說.
你教會裡面有沒有硬硬的那些人.
大家明白我的意思嗎.
宗教政治如何防備.
或者我們結束的時候.
可以看其中一段耶穌的講話.
就是我們剛才大綱的時候說到.
天上權柄的問題.
耶穌講完之後.
反駁了之後.
他就講了一個比喻叫做.
邪惡玄護的比喻.
邪惡玄獄的比喻.
十二章第一節.
耶穌就用比喻對他們說.
有人開墾一個葡萄園四圍.
捉上泥巴.
挖了一個榨酒池.
蓋上一個流氓池.
這些東西全部都是Standard.
當時Renovation一定要做的.
就是你要建一個葡萄園.
山那裡要作一些梯田出來.
然後要捉起一些石牆.
泥巴.
去留著好的土壤.
然後在梯田上面就可以去種果樹.
然後挖一個榨酒池.
那些葡萄就可以放在榨酒池那裡.

$^{761}$然後你可以把它釀成葡萄汁.
或者釀成葡萄酒.
當你收割的時候.
自然就會有一個守望樓.
因為你在田地那裡過夜.
看著你的收成.
你就很忙的去收割.
他做完所有事情.
這個園主就遭給了玄戶.
出外遠行.
到了那個時候.
他就打發一個僕人到玄戶那裡.
向他們收果子.
他們就拿著他.
打了他.
叫他空手而去.
園主又再打發一個僕人到他們那裡.
他們就打傷他的頭.
然後侮辱他.
園主又打發一個僕人.
又再去.
他們又將他殺死.
之後又打發一個很厲害的僕人.
又被他們打了.
然後殺了.
園主還有一位就是他的愛子.
最後又打發他的愛子去.
他心想應該會尊重我的兒子吧.
那些園戶.
他們彼此來說.
這個就是承受產業的.
等我們來到去殺了他.
產業就歸給我們了.
於是他們就拿著他殺了他.
將他扔出葡萄園.
這樣葡萄園主要怎麼做呢.
他就必定會來審判.
將葡萄園交給別人.
這個很短的比喻其實是說整個舊約的歷史.
上帝對以色列人.

$^{801}$給了一個葡萄園給他們.
然後租給他們來看守.
讓他們好好管理.
誰知道他們一個才知道.
兩個才知道.
三個才知道.
全部將他們趕走和殺害.
最後耶穌愛子主基督來了.
他們都會將他殺了.
上帝要做什麼呢.
就是要除滅這一幫惡僕.
然後將這個基業交給別人.
交給別人.
這個比喻有什麼可以幫助我們來防備呢.
看完這次舊約的歷史.
是一個這麼慘痛的歷史.
這裡說出幾件事.
第一件事.
我們就是管家.
或者應該說我們只是管家.
上帝按照他們的心意.
按照他們的主權.
交託東西給我們.
對管家的要求是什麼呢.
管家是沒有永久擁有權的.
沒有永久的佔有權的.
請你不要將管家的權.
當作是忠心的擁有權.
交給你管一段時間.
不代表你會永恆擁有.
押在這裡.
你就當是忠心.
你要記得.
耶穌警戒的門徒就是防備法利錯人的孝.
最後三點第一點.
第一件事我們要認清楚.
宗教政治的本相.
那種的黑暗.
信仰是可以賜給我們生命.
但是我們每一個人有的罪性.

$^{841}$我們都可以利用信仰.
然後愛權給.
宗教政治的本相我們認清.
第二件事我們要知道的.
是上帝託管我們的東西.
是有限期的.
十二章第二節.
到了時候.
要向他們收葡萄園的果子.
When the time comes.
要有exploration.
交託給你管的東西.
無論是恩賜財幹權柄.
你的信仰的角色.
全部東西都好像定期存款.
是有時限.
他會mature.
上帝就要求你成熟的時候有收成.
這樣警覺我相信.
起碼給我們一種防備的心態.
第三.
最後一樣.
就是耶穌的生命.
最後一周他做什麼.
就是十字架的道路.
十字架的生命.
他以捨己.
放下他自己的權柄.
來拯救萬民.
叛徒就要學習耶穌的生命.
一方面.
交託給我們的.
我們就忠心去事奉.
但我們知道有限期.
限期來到的時候.
我就需要交出果子.
我會停下.
我現在得起也放下.
我努力.
為了下一代去鋪路.

$^{881}$為了將來的人去鋪路.
不是我自己.
但願我們有一個謙卑.
十字架的生命.
來到耶穌最後一周.
不只是一些儀式上我們紀念耶穌為我們所做.
耶穌的生命在我們生命當中.
展開的時候.
同樣要求我們捨己.
背十架.
好像我們主一樣.
我們同心低頭祈禱.
上主我們求你憐憫我們.
因為我們每一個都是頗衰有罪的.
主你交託給我們一個美好的生命.
主你我們求你給我們效法.
耶穌基督自己的榜樣.
讓我們的生命給你去用盡.
並且讓我們被警惕被警戒.
到了時候我們需要交出果子.
上主我們求你給我們謙卑.
謙卑,忠心來服侍你.
你為我們所做的一切.
是歸榮耀給你的.
讓我們有限的生命.
能夠被你用盡.
求你自己親自預留.
預留我們的事奉.
也幫助我們謙卑.
走十架的路.
學習我們耶穌的榜樣.
我們的主的榜樣.
捨己背十架.
去事奉你,祈禱奉主名.
好,元常大利,祝福各位.
\newpage



\section{}
\label{sec:55R6ROPsg74}
\textbf{Kingdom Values XXV: Matthew Chapter 6 - Pastor Sam Johnson}
\newline
\newline
連結: \href{https://youtube.com/watch?v=55R6ROPsg74}{\texttt{https://youtube.com/watch?v=55R6ROPsg74}} ~~~~ 語音日期: 2024-03-22
\newline
\newline
\hyperref[sec:dON_DJyLmhA]{\small{< < < PREV SERMON < < <}}
~
\hyperref[sec:index_chronic]{\small{[返順時目]}}
~
\hyperref[sec:index_scriptual]{\small{[返順卷目]}}
~
\hyperref[sec:5fl01IlnFNQ]{\small{> > > NEXT SERMON > > >}}
\newline
\newline
$^{1}$Good morning family. It is great to be together again. Last week we finished up.
Matthew chapter 5 and so this week we're going to be beginning Matthew chapter 6..
And Jesus opens Matthew chapter 6 talking about hypocrisy and he actually.
talks about it for quite a few verses. He gives us three really powerful.
illustrations of how damaging hypocrisy is to our lives as subjects of the.
kingdom. And so we'll be talking about this topic for the next few weeks.
because we're going to go through these these verses. It's a pretty big section. But.
today we're just going to kind of introduce it and scratch the surface..
We're going to look at the first verse of Matthew 6. As usual I have some.
footage for you to look at today. It's hiking footage but I'm going to let you.
know it's not from Cook County. All the other hikes that you've seen in the.
past few weeks have been from Cook County, mostly the forest preserves right.
here around the city of Chicago. And most of them are in southwest Cook County.
because that's close to where I live. But there are places you could get to in 30.
minutes or so from the city. But this week we're going to have something.
different. On the last week of September, Jessamyn and I took the Z's to Great.
Smoky Mountain National Park and we really enjoyed it. We stayed in the North.
Carolina side. It was really quiet there, not many people, easy to social distance.
and it was very peaceful. And we did a lot of hiking with the kids in the park..
So I'm going to share some footage of that with you. But as you see it you'll.
probably see some things you'll say, "Wait a second, I don't think that's anywhere.
around Chicago." It is really beautiful and it was a beautiful time of year to.
be there. And so I just want to encourage you, as I often do, to get outside and.
enjoy God's creation. It's something that's really healthy for us, especially.
during times of uncertainty. And it helps us to connect with God and gives.
us a sense of peace. And this is a great time to do it. Autumn is coming and so.
the trees are changing colors. Everything is really beautiful. So I encourage you.
to do that right here close to home. Find a place with not many people and stay.
safe and wear a mask if you need to. Do whatever you need to do. But get out and.
enjoy God's creation. It's an encouraging and uplifting thing..
Back to Matthew chapter 6. Let's read verse 1 of chapter 6 together. In Matthew.
6 1 Jesus says this, "Beware of practicing your righteousness before.
other people in order to be seen by them. For then you will have no reward from.
your Father who is in heaven." Now as I said, Jesus goes on in this passage to.
give some illustrations of hypocrisy and of what this looks like. But in this.
first verse, Jesus is kind of laying out a general truth about hypocrisy. That we.
should be careful and avoid doing things for people rather than God. Caring about.
what people think of us rather than God. Caring about what the appearance, what.
the image of what we're doing is rather than the reality of our hearts and our.
lives. And that's what Jesus is pointing us to. And I want to call us back real.

$^{41}$quickly as we start to talk about this to the last verse of Matthew chapter 5.
Matthew 5 48. Remember we talked about that verse last week where Jesus says.
you therefore must be perfect as your Heavenly Father is perfect. And we kind.
of scratched the surface of talking about some of our responses to that.
verse. And we spent most of our time last week talking about the right response to.
that verse. The correct response that is to surrender to the fact that we can't.
live up to Matthew 5 48. We can't be perfect in the way that our Heavenly.
Father is perfect. We can't live up to God's standard. And so we need to.
surrender to that and we need to realize that Jesus wants to live through us. And.
that's the way that we begin to be who God wants us to be. And we begin to be.
who we want to be. But what I want us to see this week as we start to think about.
what it looks like to avoid hypocrisy in our lives. I want us to see that when we.
are confronted with the fact that we can't live up to God's standard. That we.
are imperfect as humans. That we make mistakes. That we sin. That we fall short..
There are a few different responses we can have. And we mentioned this last week.
but we didn't dig into it as much. And the first response that we could have.
is just kind of a response of rebellion. Basically of saying, "Okay God I don't.
measure up to your standard and I don't want to. And I'm not gonna try and I'm.
gonna turn away and just do my own thing and live my life." And that's honestly the.
response that most of the world around us has to God's standard and to realizing.
that they fall short. Or maybe they don't even believe that they fall short and.
they say, "You know what I don't think I fall short." But their response is saying,.
"I want to rule my own life. I want to do what I want to do. And I'm not interested.
in listening to God. I'm gonna rebel against whatever God has said in this.
world and I don't care." That's one response that many people have. And that.
we have probably had at different times in our lives. Where we said, "God I don't.
measure up to the standard and I'm not gonna try. I want to just go my own way.".
Another response that we did mention last week is the response of despair..
Just saying, "Well I can't measure up." And really allowing that to call into.
question our self-worth and attacking ourselves and getting to a point of.
condemnation and shame and saying, "I'm just a bad person and I can't do it and.
and I'm just gonna fall into despair and hopelessness." And obviously that's not.
the response that God wants us to have. God desires for us to grow close to Him.
and to change, to be transformed, to be made into His image. So God doesn't want.
us to get stuck in a place of despair or hopelessness. But that is something that.
can happen to us. And we've probably all been there as well in our journey as a.
believer. Where we said, "Okay I'm just struggling and I'm not who I want to be..
And I feel hopeless and I want to give up that place of despair." Another.
unhealthy response is this response that Jesus is going to focus on in Matthew.
chapter 6. And that's the response of dishonesty or hypocrisy. And that's when.

$^{81}$we basically say, "You know what? I am going to pretend that I measure up and.
I'm going to tell a lie about the fact that I fall short." And we can tell that.
lie to different people. Sometimes we tell that lie to ourselves and we say,.
"Yeah I don't fall short. I do measure up. I am everything that God wants me to be.
and everything that I want to be." Sometimes we tell that lie to ourselves..
Sometimes we tell it to God. Sometimes we say, "God I don't need you in my life.
because I'm good enough on my own." But sometimes we tell it to other people and.
we try to put forward this image that we are something that we're not. We try to.
put forward this perception, this illusion that we are keeping God's law.
and we're living right and we're doing all the right things in life even when.
that's not the reality of our heart. And that's a pretty good working definition.
of what hypocrisy is. Hypocrisy is when we are pretending, when we are acting,.
when we are giving off an image or creating some illusion that we are.
something that we're not. And when we're pretending to be something we're not..
When we're pretending to be someone that we're not. That's when we fall into.
hypocrisy. And we know that hypocrisy is so hurtful, it's so damaging. And.
hypocrisy in relationship to our faith and our following of Jesus is probably.
some of the most damaging of all because it causes people to just discard God's.
truth. And people look at someone's life who's hypocritical and they say, "Well you.
know there's some things there, I've heard some truth about Jesus that.
maybe I might be interested in but then I see hypocrisy in someone's life and it.
just turns me away. I just want to discard, even the good things, I.
don't want to listen to them anymore because of hypocrisy." And so we know it.
can be really unhealthy and we've probably all been hurt by it in our.
lives. We've probably been hurt by someone who held us to some standard.
or told us that we were supposed to be a certain way but then they didn't live.
that out themselves. And their words and their actions didn't agree and that.
hypocrisy, that dissonance bothered us and hurt us and made us upset. We've.
probably all experienced that. And the reality is all of us as humans are.
hypocrites. We've all acted out hypocrisy. All of us have. It's something that's.
easy for me to talk about because I'm an expert in it. We all are experts in.
hypocrisy. We're all really good at pretending to be something we're not..
We're all really good at trying to make other people think that we are someone.
that we are not in our heart and on the inside. It's something that all of us do..
And last week we talked about the one right response to Matthew 5:48 and that.
is that response to surrender. We really covered that last week so I'm not going.
to get into it again but I just want to draw one thing as we move on this week.
and that is the fact that surrendering to God and allowing Jesus to live.
through us is really a journey and not a destination and that's so important for.
us to understand. And it's important that we understand that in order to avoid.

$^{121}$hypocrisy in our life because sometimes as followers of Jesus, as Christians, as.
subjects of the kingdom, we can think that surrender is a destination. That we.
get there once and once we've arrived that we're there for the rest of our.
life. That we get there and we say, "God, I'm gonna surrender. I'm gonna give it all.
up to you." And then after that it's done. And maybe we think that happens when we.
begin a relationship with Jesus when we're saved or maybe we think that.
happens when we have some type of experience in church where God really.
moves on our heart and we say, "Okay God, I'm gonna surrender everything to you.".
And those are beautiful moments and yes when that surrender happens it's real.
and it's complete but we still have this natural part of us, what the.
scripture calls the old man or the flesh, this part of us that's still.
really selfish and really sinful and that part of us fights against the.
spiritual part of us. The New Testament talks about this and that part of us.
wants to grab control back that we surrender to God. And so we don't just.
surrender once. It's not something that we get to that point and we surrender.
and then we never have to surrender again. No, we have to surrender and we.
have to surrender over and over again. It's a journey. It's a constant.
journey of surrender. We have to daily surrender to God's control and moment by.
moment when that part of us rises up and says, "No, I want to be selfish and I want.
to focus on me and I want to do what's right for me instead of trying to obey.
God." We have to surrender again and again and again throughout life. And that.
process, that journey in church, a lot of times people call it sanctification. It's.
just a big church word that means that we're becoming more like Christ..
And as we continue to surrender, that's what happens. As we surrender more and.
more of our life and we surrender in those areas that we struggle in again.
and again, we begin to become more like Christ. We begin to be more of who God.
wants us to be and more of who we want to be. And our behavior changes, our.
attitudes change, the way we treat other people changes. Those things happen, but.
it's a process. It's not just a destination. And so if you find yourself.
saying, "Okay, I see that I fall short. I do want to surrender to Jesus, but then.
what about next week when I struggle again? What about a month from now when.
I'm still having problems with the same issues?" That's okay. It's a journey.
of surrender whereby we continue to walk in those things. We continue to repent of.
our sin. We continue to give control to Jesus and we continue to release control.
from our own hands. We continue to try to abide in Him like we talked about last.
week in John 15. It's a journey. It's a process. And that's important for us to.
know because sometimes in the church we say, "Well, I surrendered. I think I'm done..
I've got it all figured out." And we start to walk in our own strength again. And.
once we do that, we start to fall into hypocrisy. We start to become someone who.
is acting like we're something that we're not. So today as we kind of.

$^{161}$introduce and think about hypocrisy, I want us to think about a few kind of big.
ideas about hypocrisy and maybe they'll help us to be aware of it in our.
life and to avoid it. And the first thing I want us to think about is why we're.
tempted towards hypocrisy. Why as humans there's such a strong temptation towards.
it. And I think the biggest reason that we're tempted towards hypocrisy, if we're.
honest with ourselves, is pride. We don't want to believe that we fall short of.
God's standard. We don't want to believe that we're something less than who we.
want to be. We don't want to believe that we're something less than perfect. And we.
don't... Our pride tells us, "No, don't let anybody know that you make mistakes. Don't.
let anybody know that you're less than perfect. Don't let anybody know that you.
could have a weakness or that you could struggle in any way." And so our pride.
puts us in this position where even though in our heart often we know that.
we we're broken and we struggle and we're not perfect, we don't want others.
to know that. And so we try to hide that. We put up a facade. We begin to build.
that image, that illusion that we're something that we're not. And our pride.
really tempts us towards hypocrisy. And there are other things that can tempt us.
towards that as well. Sometimes it's fear. It's fear of what other people will.
think or how other people will treat us. If they know that we make a mistake or.
if they know that we've failed in some way or if they know that we're imperfect,.
sometimes it's fear of the consequences. If someone finds out that we've made.
some mistake, what will happen and what will happen to us? And we might be afraid.
of that. Sometimes it's... Maybe it's just deception. Maybe sometimes we.
really are convinced that putting forth an image that is not really who we are.
is the best way to to live life or that maybe we actually are convinced that.
that's gonna help us or it's gonna protect someone else. There's a lot of.
reasons that we can do that. But I think a lot of them go back to that sense of.
pride. The fact that we don't want other people to know that we need any help. We.
don't want other people to know that we could struggle. We don't want other.
people to know that we fall short of God's standard. And the reality is the.
Bible tells us all of us fall short and that's why all of us end up being.
hypocrites because we all try to hide the fact that we fall short. We all try.
to hide it from one another. Now what is hypocrisy? When we talk about it in a.
biblical sense, I think we understand that what hypocrisy is, it's when we.
pretend to meet God's standard and we don't. None of us can meet God's standard.
on our own. Only through Christ can we meet that standard. We talked about that.
last week. And so anytime we start to pretend, whether it's just in our own.
mind or in the words or the actions that we display before others or the.
impression that we give other people, when we start to pretend that we're.
meeting God's standard and we're doing it, I'm doing it, when I pretend that I'm.
doing it, that's hypocrisy. I can't meet God's standard. Jesus can. And because of.

$^{201}$His love for me and His forgiveness and His grace, He can live through me and I.
can meet His standard with the power of Christ in my life. But on my own, I cannot..
But oftentimes we begin to act like we can and do meet God's standard on our.
own, apart from Christ. We begin to act like we are making all the good choices.
and we're doing everything right and we are figuring everything out in life and.
that something is happening in our life because of our own goodness and that.
we've done all these things. That's just not true. It's when we pretend to.
meet God's standard, even though we don't and we can't, that's hypocrisy. And Jesus.
gets after the hypocrisy that we sometimes live out in our faith and He.
gives us these illustrations that we'll talk about in the next few weeks of.
people who do things to try to impress others. Sometimes we can get the idea.
that, well, if that's wrong, then maybe the way to avoid hypocrisy is just to.
hide everything I do. But that's not what Jesus is trying to say. Remember.
Matthew 5:16, the verse that we talked about a few months ago, it says.
we're supposed to let our light shine before men that they can see our good.
works and glorify our Father in heaven. Now the distinction here is that in.
Matthew 6:1 Jesus tells us, "Don't do things to please men," and in Matthew 5:16.
Jesus tells us, "Do things to please God." It's not really about hiding your.
actions from other people. It's about having the right heart motivation as you.
do them. And so that's another thing we need to understand about what hypocrisy.
is. Hypocrisy is when we undertake actions, even good actions, in order to.
give ourselves an image in front of other people rather than in order to.
please God. So how can we avoid hypocrisy? How can we avoid doing this? How can we.
avoid doing things for the wrong reasons? How can we avoid doing them to please.
other people or because other people are watching instead of doing them for God?.
How can we avoid giving other people an image of ourselves that doesn't agree.
with who we actually are in our heart? Well, Jesus has already given us a pretty.
good formula in the Beatitudes. Remember the first Beatitude, Jesus talks about.
poverty of spirit, and we talked a lot about what that is, that it's.
understanding your own spiritual bankruptcy and your own spiritual.
brokenness before God. And I think that's a foundational value that will help us.
to avoid hypocrisy. If we're operating from a place of humility where we.
understand our own need for God and we understand how broken we are, then we.
won't be trying to show other people that we're better than them. We won't be.
trying to give someone else an image of our own greatness or our own goodness or.
our own righteousness, because we'll understand who we really are. And I think.
that's a really important place for us to start as we think about how to avoid.
hypocrisy. But I think also, as I said, all of us as humans are hypocrites. All of us.
don't want other people to see that we might make a mistake or that we.
might be less than what we want to be. And so another important part of.

$^{241}$avoiding hypocrisy is being able to repent, being able to look at our life.
and say, "You know what? I wasn't fully honest with the people around me about.
who I am. I tried to give them a perception or an image that I'm.
something that I'm not, and I need to repent of that." Be honest and repent.
before God, first of all. Take that to God and say, "God, that wasn't right. I know.
it was wrong. Will you forgive me? Will you help me not to do that in the future?".
But also, if you need to repent before someone else, and what I mean by that is.
you can ask their forgiveness. If you've done something that was hypocritical.
and it's been hurtful to someone in your life and you know that it's been.
hurtful to them, then you should go to them and ask forgiveness and say, "Hey, I.
know I represented myself as this and that's not who I am," or "I know I asked.
you to live to some standard that I didn't live up to. I know I asked you to.
to do things and respond in a certain way, and then I didn't respond in that.
way. I know I got frustrated at you for getting upset or for struggling in.
this area, and that's the area where I struggle, and I get upset about those.
same things all the time." If your hypocrisy has harmed someone else, then.
repent to them. Go say, "Hey, I'm sorry, and that wasn't right, and I want to.
be more transparent in the future, and I want to be more honest about my.
struggles, and I don't want to present myself as something that I'm not.".
Obviously, we can always ask God, and we can say, "God, can you really help me?".
This is an area where you struggle, where you find yourself a lot of times.
presenting yourself to other people as someone that you're not, and trying to.
show yourself as better than others, or trying to show yourself as more.
righteous than others, or whatever it is. If you struggle in that area, ask for.
God's help, and say, "God, can you help me to stay humble? Can you help me to be.
honest and transparent about my struggles in my life, and not to hide.
that from other people? Can you help me to not try to put some image out there,.
to not act like I don't make mistakes?" Then, obviously, the truth.
that we talked about last week, that surrender, whereby you say, "God, I know.
the reality is I'm not going to be perfect as my Heavenly Father is.
perfect. I'm not going to live up to your standard, and I'm going to surrender to.
you, and let Jesus live through me, so that when people look at me, they will.
see the reality of the situation. They will see the reality of what it means to.
be a Jesus follower, to be a subject of the kingdom. Here's the reality that.
people should see in our lives. Here's the reality that I pray people see in my.
life. When someone looks at my life, I want them to see me as a broken person..
I don't want them to see me as some type of great spiritual giant. I want them to.
see me as a broken person, but I want them to see the love of Jesus and the.
attributes of Jesus shining through me. That's what we should pray for, and.
that's what we should surrender to. We should say, "God, I want Jesus to live.

$^{281}$through me, so that when people look at me, they won't see my perfection. They'll.
see Jesus' perfection." Because when we are hypocritical and we act like we're.
something we're not, we put the focus on us. It's hard for people to see Jesus..
Even if they buy into our hypocrisy, they're just going to think we're the most.
amazing person that ever lived, and they'll never see Jesus. We don't.
want that. We don't want that. It's unhealthy, and it puts us on a pedestal.
we were never meant to be on. It puts us on a pedestal that, once we fall.
off, it disappoints and discourages people. We don't want to live that way. We.
should want people to see our brokenness and to see the perfection of Jesus..
As John the Baptist said, we should be decreasing so that Jesus can be.
increasing. That point of surrender will help us to avoid hypocrisy in our.
life, to say, "God, I want to surrender to you and let you live through me, so that.
people are seeing you and not me." Let's pray. God, thank you for the reality of.
your truth. Thank you for the way you don't let us get away with our mess..
God, our hypocrisy is mess. It is a problem. It is so toxic and so hurtful..
It's so hurtful to us to pretend that we're something we're not. It's so.
hurtful to the people around us, and it's definitely hurtful to our witness in the.
world as believers, as we try to share the gospel and share your love. If we are.
being hypocrites, then it destroys our witness. God, I pray that you would.
really work in our hearts. Help us to see the reasons that we act that way. Help us.
to see our pride or our fear or ways that we've been deceived and to turn away.
from those things. God, help us to see that we can't meet your standard and we.
shouldn't act like we do. God, help us to see that you desire for us to live.
in a way that's pleasing to you, not just trying to please other humans. God,.
help us to live in a place of poverty of spirit, humility, repentance. God, help us.
to live from the place of surrender so that Jesus is able to live through us..
God, help us to avoid hypocrisy. Help us to avoid putting out an image that we're.
perfect or that we are doing everything right or that we're something that we're.
not. Help us to live in an honest and transparent way so that our witness is.
amplified and so that when people look at our lives, they see less and less of.
us and more and more of Jesus. God, we just ask these things in your name. Amen..
\newpage



\section{}
\label{sec:5fl01IlnFNQ}
\textbf{聖經地理 EP10 - 地理與釋經 - 地景/氣候 (陳崇基牧師)}
\newline
\newline
連結: \href{https://youtube.com/watch?v=5fl01IlnFNQ}{\texttt{https://youtube.com/watch?v=5fl01IlnFNQ}} ~~~~ 語音日期: 2024-03-23
\newline
\newline
\hyperref[sec:55R6ROPsg74]{\small{< < < PREV SERMON < < <}}
~
\hyperref[sec:index_chronic]{\small{[返順時目]}}
~
\hyperref[sec:index_scriptual]{\small{[返順卷目]}}
~
\hyperref[sec:6DyYwyvwR5E]{\small{> > > NEXT SERMON > > >}}
\newline
\newline
$^{1}$我們繼續看下去.
講到這個地景的時候.
就可以看一下氣候.
整個迦南地的氣候到底是怎樣的呢?.
就是那個氣候是怎樣的呢?.
就是長年累月的循環是怎樣的呢?.
在這裡我用了一個圖來表述出來.
藍色這個地方.
藍色的1,2,3,4,5,6,7,8,9,10,11,12.
就是今天的西力.
今天的西力.
中文紅色這些數字就是農曆.
猶太人的農曆.
或者好像我們華人的農曆一樣.
都是跟月份的.
跟月亮的.
月曆是不是?.
所以由1月就是今天的西力3-4月左右的時間.
直到12月春節的時候.
1月就到春節了.
在這個12月的循環裡面.
你就看到有兩個季節.
第一個季節就是雨季.
雨季是9月左右開始.
有第一陣雨在希伯來聖經裡面叫做初雨.
有一個特別的字給它.
然後過了這個初雨之後.
那些泥土開始鬆了.
因為已經過了夏季了.
泥土開始鬆了.
灑了一陣水.
然後就開始農夫可以耕地了.
乾不可以.
因為乾到裂了.
所以浸了第一陣的初雨之後.
特別是鑄棚節之後.
這些土地開始鬆的時候.
他們就可以耕地.
然後就去撒種.
播種.

$^{41}$撒種的比喻.
就是這段時間.
就是這個背景.
然後第二個的雨季就是冬雨.
Winter Rain.
就是12月至2月.
撒了種之後有Winter Rain.
那些種子就開始可以發芽.
可以成長.
然後到最後一陣的雨.
就是3月至5月.
西曆的時間.
下究的時候.
夠8吋以上.
大麥就可以成長了.
就可以做收割.
夠12吋就是小麥.
就可以收割了.
這段時間剛好就是Passover.
就是雨節的時間.
就是我們的Good Friday.
秀蘭節和復節的時間.
過了這些雨季.
4-5月左右.
就開始乾旱的季節就到了.
6-7月正式沒有雨水.
黃色這個地方就是所謂的夏季.
是沒有雨水的.
整個天很藍的.
這個地方.
這個時間唯有靠金露.
記得黑門山的金露.
早上很薄的一層霧水.
就可以使果樹能夠熬過最乾旱的夏天.
過了這個夏天.
8月至9月左右.
就開始有收成.
一個Good Year的時候.
一個夠雨水的年度.
果樹吸收了足夠的水份之後.

$^{81}$首先有葡萄的收割.
然後有無花果樹,石榴,綠橄欖,蜜棗.
不是蜂蜜.
雖然也有蜂蜜.
但主要是那些Day Palm,蜜棗.
最後有黑橄欖的收成.
說了這些農耕的循環.
我們知道這些紅色的地方.
是大節日.
其實是跟著農耕的循環.
雨節,大小麥的收割中間.
然後到收割節.
收割節就是七七節.
就是五旬節的時候.
夏季來到的時候.
可以上路.
因為乾了.
不再有雨水.
走路就方便一點.
去別的地方就方便一點.
行軍行商就方便一點.
五旬節起不就是說這件事.
從耶魯省出發.
出發到其他地方.
宣教,福音傳到地極.
都是這個季節最適合去說的.
所以過了夏季.
有足夠的金露.
就有另外一個循環.
就是果樹就可以收割.
所以看看新面記第八章第八節.
我們就了解了.
裡面說到一些七種本地的植物.
果實.
有小麥,大麥,葡萄樹,無花果樹,石榴樹.
也有橄欖樹或者橄欖油.
還有一個棗蜜,蜜棗等等.
七樣.
這七樣就是整年裡面.
有的最基本本地的出產.

$^{121}$OK.
所以就想起有很多不同的聖經.
都可以放在這個.
迦林地氣候循環,龍更循環的時間裡去理解.
譬如葡萄園的比喻.
那些工人去等進去葡萄園裡面去工作.
家主就給他們多少銀子.
這個就是去到收割的時候.
就是葡萄收割的時候.
這些就是其中一些的例子.
另外如果你看迦林地.
你就知道整個迦林地本身對雨水來說.
都會有偏差,有相距的.
比較北一點的地方.
就自然雨水比較多.
譬如嘉里尼湖附近.
黑門山等等.
這裡你記得一隻角就是哪裡.
就是加密山,海發市這個地方.
北邊的雨水是比較多的.
並且西邊的雨水也比較多.
為什麼呢?因為雨雲是由地中海吹過來的.
還有比較高的地方都會有比較多的雨水.
因為你越高的話.
就像耶魯山一樣.
水氣就會比較多.
乾旱的地方就自然是南邊.
南邊的沙發附近.
你看到這條黑色的線.
其實就是說剛剛足夠有雨水.
你可以在這裡做耕種.
這些黑色的地方都是剛剛足夠有雨水.
每年12吋的雨量.
東邊就自然是越來越少雨水.
比較乾一點.
低地的地方.
譬如這些死海峽谷.
約旦河流經的地方.
都不是很適合耕種.
因為太低了.

$^{161}$雖然是有河流經過.
所以你看到.
在十二支派的分佈裡面.
都有各自的不同的.
在北邊的支派.
拿忽塔利.
或者馬來西亞.
以法連.
這些都是比較好的地方.
北國以色列.
分了南北國的時候.
北國是強勝過南國的.
原因很簡單.
因為其實雨水都比較多.
還有它能夠拿到這些國際路線.
都是比較多一點.
南國 猶大國.
是比較弱一點的.
但上帝卻揀選南國 猶大國.
揀選這個小的國度.
揀選耶路撒冷為他命所居住的地方.
所以這裡地理.
不單單說到地景.
從地理我們看到.
有神河.
就是選舉.
神的恩典揀選的神河.
《申命記》第十一章.
因為你們要進入的那地.
並不像埃及地.
在埃及你們殺了種.
要用腳來灌溉.
像菜園一樣.
所以他們說你們進去的地方.
摩西話.
就是迦南地.
那個地方絕對不像埃及地.
埃及地你殺了種.
用腳去灌溉.
好像菜園一樣.

$^{201}$什麼意思.
你看看埃及地.
利羅河就在這裡.
利羅河附近兩邊有很多平地.
這些地方都是一些大的菜園.
農田等等.
你看到那些河道.
人工河道.
河道再進去的時候.
你幾乎用腳去刮一刮.
水就會進去農田.
你看到多綠.
迦南地不是這樣的.
迦南地是什麼地.
他說是一個有山有谷之地.
不同埃及是一個平原.
有天上雨水滋潤之地.
就是耶和華利神去到眷顧之地.
從歲守直到年終.
耶和華利神的眼目常常去到眷顧之地.
如果你留心聽我的誡明.
愛耶和華利神一心一意.
事奉他.
不敬拜別神.
他做什麼.
必定會按時降雨在地上.
就是降秋雨.
Early Rain.
初雨.
和春雨.
使你們可以豐收五穀.
新酒新油.
也使你們田野長出青草.
你牧羊羊群畜牲就有草食.
你見到流奶與蜜之地嗎.
蜜就是蜜棗.
新命紀八章第八節.
七種本地的蔬果.
最後最甜的就是蜜棗.
流奶與蜜.

$^{241}$奶在哪裡.
田野有雨水.
為你們的畜牲長出青草來.
他們得吃飽足.
流奶與蜜講到兩種生存經濟的動力.
畜牧和農業.
最主要是有山有谷.
他們靠天上的雨水生存.
因為在迦南地不會有大河.
再講多一次.
埃及就有大河.
平地一塊在埃及三角洲.
你就在那裡用腳耕種就行.
這邊也是一樣.
米索不達米亞.
兩條大河.
迦南地是怎樣.
沒有河流.
沒有自然天然的上游清水的河流.
唯有靠雨水.
這個就是Mediterranean Climate.
雨雲由地中海捲過來.
先著陸的地方就是加密山.
所以以利亞和巴力的先知鬥法.
就是這個地方.
這個也是中間之地的中間之地.
遲些有機會再講.
所以這裡是靠Rain from Heaven.
天上降下來的雨水.
你才可以有耕種.
能夠有飽種.
但是是一個試煉的地方.
試煉他們是否一心一意跟隨耶和華.
如果一心一意跟隨.
耶和華必定會眷顧稅收到年終.
眼目上去到眷顧這個地方.
有山有谷的地方是怎樣.
我們給大家一些樣板.
這裡就是沿海平原.
地中海.

$^{281}$廣闊的地中海.
然後陸路都是比較平坦.
沒有什麼大的山谷.
我們在海邊走過.
向南走向北走.
行商行軍很容易.
這個就是沿海平原.
但這只是加密山的其中一部分.
我們看的這個地方是沿海平原.
有白港加密山附近的地方.
但過了這個地方就去哪.
第二個地帶.
由左邊去右邊.
西邊去東邊.
低山地.
低山地和沿海平原不同.
這裡就是以拉谷.
剛才我展示給大家看.
這個山後面這個地方就是以拉谷.
我們在阿西加土丘上拍照.
後面再往東看就是中央山脈.
中央山脈後面就是死海.
在中央山脈的以色列人和這邊的菲律賓人.
在沿海平原中間的地方就是一個爭執.
一個門口.
經常在門口這裡有衝突.
大衛和哥利亞就在以拉谷這個地方.
一邊和另一邊.
大家在罵陣.
然後兩個人去決鬥.
再看看以拉谷在哪.
以拉谷就在這裡.
我們在阿西加這個地方.
往過去.
有一個山.
兩個山中間這個地方.
如果你今天選擇路線.
展示路線的時候.
路線都是這樣走.
這條是38號公路.

$^{321}$轉左375.
一直這樣上.
就會上到去.
375.
就會上到去耶路撒冷.
耶路撒冷就在這裡.
他都是跟著這些山脊.
這些山谷.
來建造今天的公路.
雖然有現代的公路.
有科技.
都是一樣跟著這個地景.
以拉谷.
大衛和哥利亞的關口.
低山地.
然後再上去.
就是到中央山地附近的地方.
你看到.
今天的公路.
都是非常狹窄.
只有兩條線.
還要掘水道.
所以這裡就是中央山地.
近百里行的地方.
連這些.
現代的公路.
我們都看到是幾崎嶇.
有山有谷的地方.
正如聖經裡面所說.
今天以色列的公路都是這樣.
這裡就是梳列谷.
聖經裡面都有提過.
就是中央山地裡面.
耶路撒冷附近的一個.
就是.
山谷.
叫梳列谷.
你看到這裡沒有什麼平地.
去種一些大的麥田.
小麥大麥.

$^{361}$通常都是種果樹.
種果樹的方法.
用梯田的方法.
我們住在.
起碼我的背景都知道.
我在香港長大.
廣東省南面的地方.
很多山.
經過山旁邊就會有梯田.
農夫就在這裡.
一級篤了山.
將石頭搬到下面.
然後攔住它.
然後土壤沖下來的時候.
就能夠留下來.
然後他們就在那裡種東西.
這個中央山地的地景都是一樣.
你看到一級一級.
有不同的小的農田.
還有很多的果樹.
這些橄欖樹等等.
只能夠藉著這些方法.
來繼續生存下去.
所以就靠這些梯田.
這個是中央山地.
再去中央山地.
過了耶路撒冷.
過了橄欖山.
之後就是什麼地方呢?.
耶路撒冷在這裡.
今天白色的就是這個OCP.
過了耶路撒冷.
過了這個.
這裡就是橄欖山.
一過橄欖山就是伯特利.
窮人就住在這些地方.
因為他背著.
伯特利之後一直下去.
今天的衛星圖片都是橙色的.
沒有綠色的.

$^{401}$就是猶太曠野的地方.
所以就去到猶太曠野.
我們在這裡看猶太曠野的地景.
猶太曠野背後.
就是耶路撒冷.
橄欖山等等的地方.
是綠的地方.
一過了那個地方.
就去到這個猶太曠野.
所以是一個存在不生的地方.
有山有谷.
要靠天上的雨水.
Rain from heaven 滋潤的土地.
所謂迦南美地是什麼意思呢?.
Good land 有多good.
都沒有埃及那麼好.
沒得比.
但good是什麼?.
就是有上帝的看顧.
上帝的恩典.
上帝的同在.
這個才是good land.
有上帝的prayer.
有上帝的祈禱.
有上帝的good land.
有上帝的presence.
但這個有條件.
以色列人甚至是知民.
就是要見到上帝公認.
你就要一心一意去事奉他.
不要離棄他.
不要敬拜別神.
不要敬拜巴黎.
只有耶和華才是歲守到年終.
眼目相照看顧這個地方.
\newpage



\section{}
\label{sec:6DyYwyvwR5E}
\textbf{Faith and time: Sabbath - Rev. Mark Chan}
\newline
\newline
連結: \href{https://youtube.com/watch?v=6DyYwyvwR5E}{\texttt{https://youtube.com/watch?v=6DyYwyvwR5E}} ~~~~ 語音日期: 2024-03-25
\newline
\newline
\hyperref[sec:5fl01IlnFNQ]{\small{< < < PREV SERMON < < <}}
~
\hyperref[sec:index_chronic]{\small{[返順時目]}}
~
\hyperref[sec:index_scriptual]{\small{[返順卷目]}}
~
\hyperref[sec:6uSoGUnudqk]{\small{> > > NEXT SERMON > > >}}
\newline
\newline
$^{1}$Can you say hi, like high five or whatever?.
Just wave your hand, okay?.
Good to see you all..
And it's good to be back to think about the Lord's,.
the Lord's word, what it's gonna say to us..
We're gonna go through our series,.
continue with our series, Faith and Time..
Think about how we use our time..
Think about how God will put different markers.
in our time, in our journey,.
to remind us of his work and his grace..
Today, we're going to talk about the concept.
or the practice of Sabbath..
Sabbath..
Sabbath is what day of the week?.
First day, third day, fifth, sixth?.
What's your guess?.
The seventh, right?.
Yeah, exactly, seventh, yeah..
The seventh, okay..
Well, according to Jewish calendar,.
it's Friday, sundown, you're gonna launch the Shabbat..
And then for 24 hours until Saturday,.
which is the seventh day, right?.
Seventh day of the week,.
until the sundown of the next day on Saturday..
So 24 hours is the Sabbath or Shabbat, all right?.
It is a regular weekly practice,.
even for many modern Jewish people,.
Israelis, or even in Jewish communities today,.
like community in Skokie or in other areas,.
Brooklyn or anywhere you get a concentration.
of a lot of Jewish people,.
they will continue the practice of Shabbat..
Friday night to Saturday night,.
for 24 hours, they don't do any work..
They stop anything that they do in the past six days..
So they cease all the work,.
and then they don't do a lot of things,.
like they don't cook..

$^{41}$They don't start any fire..
So in modern day technology, they put their phone away..
They don't push any button, you know what I'm saying?.
'Cause you push a button, the light will go up, right?.
And that is considered a violation of the Shabbat,.
because you're lighting things up, right?.
Starting a fire..
Strange, right?.
Do you wanna live like that?.
No, right?.
(audience laughing).
But if you ask them, many of them will be grateful.
for such a practice..
Interesting, right?.
To stop anything that you're doing,.
just to take a break,.
especially in our technology tyranny..
We are living in a technology tyranny..
We are being massacred actually by our phone,.
by technology, by a lot of different things..
So it's a very interesting practice..
And actually the Bible emphasizes once and again.
that the people of God should do something like that,.
should observe Sabbath..
As I said, many modern Jewish people,.
modern Israelis, they are practicing this sort of Shabbat..
So let me give you some example..
If you go to Israel or any community.
that has a lot of Jewish people,.
you'll find a Shabbat elevator..
A Shabbat elevator is a Jewish elevator..
So whenever there's a weekly Shabbat or high holidays,.
they already program the elevator.
so that you don't have to go into the elevator.
and push a button, right?.
Again, that's why I said, if you push a button,.
something light up, it consider violation.
because you are lighting things up, right?.
You're starting a fire..
So here's one of them..

$^{81}$I think this picture was found in Brooklyn, in New York..
The second one is another form of Shabbat elevator..
So you go to a high rise, right?.
And this elevator, the one on the right, elevator A,.
and which floor do you go to?.
This Shabbat elevator, by the way,.
we will stop at every single floor and then come back..
So you go into the Shabbat elevator,.
you don't have to press the button..
It will close automatically..
It will bring you to the floor that you want.
and then you will exit, right?.
And then the elevator A on the right is that,.
from basement to ground floor is zero,.
to one, two, three, four, and then come back down.
to zero and then minus two or basement,.
and then go back up..
Elevator D is a little bit different, right?.
If you want to go to floor eight,.
then you take A or you take D..
That's right, if you take A,.
then it will take a long time for you.
because you have to come back..
And then you find out I was in the wrong elevator..
And then you have to go to D, right?.
So it will stop every single floor.
in order for you to cease doing any work to observe Shabbat..
Last one I want to show you is this one..
It's a warning sign in case of fire,.
do not use the elevator..
Do you know that?.
We shouldn't use elevator in case there's fire, right?.
And even worse, if you go into a Shabbat elevator, right?.
'Cause it will bring you to the floor of fire, right?.
(laughs).
And then you have to survive the floor.
and then you come back, right?.
It will be a big trouble for you..
This sort of thing is very interesting..
But as I said, the Bible keeps on talking about.

$^{121}$the marker in the lives of the people of God.
that they need to do something like that, all right?.
In their time, in their schedule,.
there has to be something that will show,.
that will mark themselves away from.
or different from other people in the world..
Shabbat is one of the killer, one of the sure marker..
So what does it mean?.
Shabbat means stop..
You stop doing anything..
You cease to work..
One day a week..
You stop doing what you have been doing for six days..
And then you stop for one day..
That is what Sabbath means..
The first time the Bible mentions Shabbat is.
Shabbat is the creation story, Genesis 2..
The creation of heaven and earth,.
everything is done, was complete..
On the seventh day, God had finished his work of creation..
So he rest Shabbat from all his work..
And God blessed and sanctified, declared holy the seventh day..
Because it was the day when he rested Shabbat.
from all his work of creation..
So right off the bat, the very first page of the Bible,.
God works like that..
Six days he created or he organized the entire cosmos,.
arranged it in such a way that is in order..
And then the seventh day, God blessed and sanctified.
and set it apart, declared holy, or the seventh day..
And then there is the fourth commandment.
in the 10th commandment passage..
One of them is Exodus 20..
Remember to observe the seventh day by keeping it holy..
And by the way, this is the fourth commandment..
This is the longest explained commandment of all 10..
You know, other commandments are just one or two sentences..
Here, talking about Shabbat, the fourth commandment..
Observe the seventh day by keeping holy..
You have six days each week for your ordinary work..

$^{161}$But the seventh day is special..
It is Sabbath, day of rest, dedicated to the Lord, your God..
On that day, no one in your household may do any work,.
including your sons and daughters, your male servant.
or your female servant, your livestock, everything else,.
even your sheep, even your pets, your dog, your cat..
They stop doing any work..
I know your dogs and your cats don't do any work..
You do works for them, right?.
In six days, the Lord made the heavens and the earth,.
the sea and everything in them..
On the Shabbat, on the seventh day, He rested..
And that's why the Lord Yahweh blessed the seventh day.
and set it apart, sanctified it as holy..
This is how God works..
Six days of work, and then one day, He set it apart,.
sanctified it, declared it special and holy..
Holy and special, distinct from all the other six days..
A six plus one, parallel, right?.
Six days, and then one day, six day and one day..
It's like you read the music scores, right?.
You know how to read music scores?.
Kind of, right?.
There's always rest at some places, right?.
Good music does not mean you jam everything together.
and there's no space, right?.
You hear someone play piano, if he's good,.
play some sort of instrument,.
there has to be intentional rest in between, right?.
Do you get it?.
Right?.
It's like you writing a novel or something..
There's always commerce in between..
And then period, here and there..
Punctuations..
Sabbath is our life punctuation..
It's our rhythm..
God set the rhythm of our life..
Six plus one, and then six plus one..
This is how God works..

$^{201}$God bless the seven day and declare it holy..
Can you see what God does here?.
It's very interesting..
I read, you read it again..
I mean, there's so many familiar passage in the Bible.
that we thought we know what it means.
because it's so familiar..
But if you slow down and read the key words,.
you pick up the meaning..
Genesis one and two is one of the examples..
We all know the creation account, right?.
We read it and we feel it's so familiar, right?.
There's nothing surprising..
But if you come to the seventh day, chapter two,.
it's very interesting..
God bless, incentivize, declare holy the seventh day..
Now what's so special about it?.
God created everything, right?.
Day one to day six..
Then found the climax is creating humanity, right?.
There were a lot of things..
There are a lot of stuff that God put in order.
in the universe..
Now those are the things that can be set apart, right?.
Just like the temple, you go into a holy place,.
there were something more sacred than others, right?.
There's always distinction between what is sacred.
and what is not sacred, what is holy and what is profane..
Things can be set apart, right?.
People can be set apart..
Priests serving in the temple services,.
they can be set apart, they can be sanctified..
They can be declared holy, different from other people..
So things can be set apart and personnel can be set apart..
However, the first thing God set apart.
is not a person or a people..
It's not a thing or things..
The first thing God sanctified is time..
24 hours out of one week..
He bless it, seven days, and declare it holy..

$^{241}$Find it interesting that in the creation rhythm,.
there is a time pattern..
There is a slot of time in the entire week.
that is called sacred time..
Not anything, not any person, not any object,.
but a period of time, right?.
There's one very famous rabbi.
who wrote a very well-known work on Shabbat..
His name is called Abraham Herschel..
Now let me give you a quote..
I think this quote worth your reflection..
He talks about Shabbat, a time that is set apart..
Sacred time, special time..
So Herschel said, "There is a realm of time.
"where the goal is not to have, but to be.".
Not to own, but to give..
Not to control, but to share..
Not to subdue, but to be in accord..
Now think about this thing..
Consider all things in the universe,.
all things that God created, all the small, big things,.
you know, down to the molecules, down to the atoms,.
up to the planets, right?.
The galaxy, the universe..
But there's only one realm that is different..
It's the realm of time..
The realm of time is purpose..
It's not to subdue..
It's not for us to control..
That's interesting because, you know,.
like objects in the universe, you can subdue.
and you can control, you can take possession of it..
Like this table..
You know, I can take this table.
and you don't have this table..
If I take it..
Like the chair..
You're sitting on your chair..
Nobody's gonna sit on your chair.
because you're sitting on it, right?.

$^{281}$Unless, you know, the other person's sitting on your lap..
You take possession of it..
You control it..
You subdue it..
Now this is our desire, this is our purpose all the time..
We take control of something.
and then nobody else can control it..
It's kind of taking exclusive possession of something..
Some property, right?.
Your house, right?.
You just bought a house, right?.
Somebody?.
You bought a house, it's your house,.
it's not somebody else's house anymore, right?.
Nobody can knock the door and say, "I wanna come in," right?.
You can only come in by invitation..
It's your house..
You take possession of it..
But Herschel said, "Time is different.".
Think about it..
You have 24 hours, I have 24 hours, right?.
You can take possession of 24 hours.
and take my hours out of mine, right?.
I have 24 hours, you have 24 hours,.
you have seven days, I have seven days..
It is a realm not to own but to give..
God created us as beings in time..
You realize that?.
We live in time..
We watch the clock..
What time we go to work?.
We don't wanna be late..
What time we go off work and then we have our time, right?.
What time do we go for dinner?.
After the service, what time are we getting there.
to have lunch, right?.
We don't wanna be late, right?.
What time is service?.
We wanna be there, right?.
But time is a realm in the entire universe.

$^{321}$not to control but to share..
Think about it..
When we have fellowship together,.
when we worship together, it's a time to share..
That's why worship in the Bible has always been,.
has always been corporate, a group, a people together,.
not alone, right?.
Not individual enjoyment..
We come together to share..
It's very interesting..
God blessed the seventh day..
The first thing that God sanctified.
is the period of time in our lives..
In your week of work, 24 hours, God set it apart..
That's why when the Bible talks about Shabbat,.
it always relate to other things..
Leviticus 23, six days each week for your ordinary work..
But the seventh day is the Sabbath day,.
the Shabbat of complete rest,.
an official day for holy assembly..
It is the Lord's Shabbat day..
It must be observed whenever you live, wherever you live..
In addition to Shabbat,.
there are the Lord's appointed festivals,.
official days for holy assembly.
that are to be celebrated at their proper times each year..
Now, what does it mean?.
Shabbat for some other purpose..
So in some sense, you start working,.
but in another sense, you do something together,.
which is holy assembly, to get together to worship,.
to get together to celebrate,.
to acknowledge who is in charge..
And who is in charge?.
God is in charge..
You go back to Genesis 2,.
when the Bible describes the work of God,.
the six days of creation,.
and then He rests on the seventh day..
Now, it does not mean that God is tired..

$^{361}$It does not mean that God is not human..
It's not like us, right?.
We work for a whole day,.
and then we go back home and we sit on the lazy boy,.
and then we say, "We are tired.".
We want to get refreshed..
No, God is not like human, but He rested..
Now, what does it mean?.
Now, in the biblical world, in the ancient world,.
when they think about God resting,.
they are thinking of God being enthroned..
Being the king..
So whenever God rests, God rests in His temple..
His temple is where He operates the entire world..
So when Genesis 2 talks about God resting,.
it's talking about God being the king..
He is seated on the throne,.
and He controls and organizes the entire cosmos..
He operates the entire cosmos..
That's what it means..
Doesn't mean six days, He works so hard,.
and then the seventh day, "Oh, He's so tired..
"He's exhausted..
"He needs some rest.".
No, sixth day, He put together the entire cosmos,.
arranged it in such a way that it's in order..
And then the seventh day,.
He sits in the control room in the temple.
to operate the entire universe..
It's kind of like, it gives you a modern analogy..
It's like, remember when the pandemic started.
about two, three years ago?.
Remember we have to switch everything back.
from in-person to online?.
So I bought a new computer..
So I went to my basement, right?.
And then I got all the computer parts,.
my standing table, my monitors, my keyboard,.
and all the parts..
I got everything in the basement..

$^{401}$And then I spent almost like two to three days.
to put the entire computer system together,.
camera and all this stuff..
Now, you got the wires,.
you got the different big or small parts..
You put everything together..
Now, if you imagine,.
we put the entire computer system together,.
what do you do after you're done?.
- Rest..
- You rest..
What do you do?.
The first button I push is the power on button..
'Cause that is when the time you control, you operate..
Same way, God put everything in order,.
like putting together a computer system..
And then He push not the off button,.
He push the on button..
So He's the King over the entire cosmos..
He control everything..
Now, that's what it means.
when the Bible talks about Shabbat..
It doesn't really,.
it means more than we imitating God,.
'cause God works six days in one day, right?.
Six in one, six in one..
It's more than imitation,.
but it's more on participation..
We can participate, we can join..
Or live under God's control..
So we set apart one day a week to worship,.
because we know the universe, your universe, your world,.
is not controlled by you..
It's by the one who created all things..
That's what worship is all about, right?.
That's what worship is all about..
Worship meaning we know we are not in control..
We are not the ultimate operator of our lives..
There's one who is Lord of all, right?.
Who is the God of all..

$^{441}$He is involved in everything that you do..
So one day in seven days,.
to be set apart in order to give Him honor,.
to tell Him, "Yes, you are the Lord who control my lives..
You are the Lord who controls my career..
You are the Lord who controls my future..
I participate in your kinship..
I participate in your worship..
You are everything, and I want to do your will.".
Right?.
That's what worship is all about..
You come in here to join our worship..
It's not just going through some motions,.
you know, joining a program..
It's to look upon God Himself..
He is the one who set up everything,.
and He's the one who is in the ultimate control, right?.
Let me close by giving you some suggestions.
how to practice Shabbat..
How do you practice Shabbat?.
You know, I tell you, I do that all the time too..
Most of the people assume.
that I go to church because of work, right?.
Yeah, yeah, I work, right?.
But you know, some people would come to me,.
"Oh, pastor, do you really have your Shabbat?.
Because you come here to work, right?".
And I want to tell you, it's not really work for me..
Seriously, it's worship..
I relate to God..
Yes, I do some work, right?.
I do some preparation..
You too do..
You need to do some preparation too, right?.
So I come here, I really enjoy the worship service..
So I come with a mentality, with an attitude.
that I'm not here to just output..
I'm getting some input too..
When I worship with the people of God together,.
you know, my spirit is refreshed.

$^{481}$because I draw closer to God..
With that kind of attitude,.
this is how we practice Shabbat, right?.
Well, you know, let me explain..
In the Old Testament, it's the seventh day of the week..
In the New Testament, it's the first day of the week,.
Sunday..
Why?.
Because the real work, the real creation is in Jesus Christ..
He rose again from the dead on the third day..
He died on Friday, on the third day, on Sunday,.
the first day of the week..
That's why in the New Testament,.
you always see the phrase, first day of the week..
Jesus rose from the dead..
First day of the week, they collect offering..
They worship together..
First day of the week..
This is the weekly example..
Now again, I come to church..
I'm not here just working..
I'm here to enjoy and I am here to get closer to God.
because I know this is the time that is sacred..
This is a time that is set apart.
from all these other things that I did in the past six days..
It's a sacred time..
It's a holy time..
Sanctified..
I spend time with God in the worship service..
So let me suggest to you,.
one thing is to adjust your mentality.
when you walk into a worship service..
To stop your work..
Stop checking your phone..
Stop thinking about the stock market..
Stop thinking about your career..
Yeah, I know it's been stressful..
I know it's been stressful, you know,.
being more than men and women..
There were a lot of things that we wanna do..

$^{521}$And sometimes we treat the Shabbat,.
the one day of the week, as our extension..
Oh, for six days I have so much work..
And then I spill over to the seventh day, right?.
Oh, I'd rather not go to observe it..
I want another day to finish my work, right?.
I want another day to earn some more money, right?.
I want another day to increase my income, right?.
I mean, what other time do I have?.
I have six days of work, right?.
And then the seventh day,.
I better use it for my own purposes, right?.
I wanna earn some money..
Now, Shabbat is a spiritual discipline, right?.
Well, ask yourself, when do you have enough money?.
When will you have enough money?.
When will you have enough salary?.
When?.
How much is enough?.
(laughs).
When is enough is enough?.
It's never enough, right?.
You never stop..
So Shabbat is to stop, to cease, cease our work,.
and know that God is the God of all..
He's the God of my finance..
He controls my finance, He controls my money..
He controls my time, therefore..
And I wanna set this time apart as a sacred time..
That's a weekly Shabbat..
So let me challenge you, let me invite you..
Come to the worship service with your heart fully prepared,.
fully prepared for worship, right?.
Well, learn from our worship team..
They always come early, so much early,.
as early as they can in order to prepare, right?.
In order to prepare the worship service.
for other people, for all of us..
We come early because we wanna set apart this time as sacred..
Then secondly, I wanna invite you.

$^{561}$to do a little Shabbat every day, a little Shabbat..
Many, many years ago, probably 500 years ago,.
there was a guy called Brother Lawrence..
Now, Brother Lawrence was in Europe at that time,.
and he was a monk, a little monk in the monastery..
The monastery has a lot of people..
Has a lot of people who wanna be priests.
or wanna be pastors in the monasteries..
And Brother Lawrence was working.
in the kitchen of the monastery..
And he wrote a little book, "In the Presence of God,".
or "Practice the Presence of God.".
Now, that little book, he shared his secret.
of how to live in the presence of God every day,.
no matter what I do..
And his secret is, when he was struggling.
with his relationship with God,.
he find himself sometimes distant from God himself..
So he figured out some way to keep himself.
in the presence of God..
So this is what he did..
Every day he was in the kitchen doing dishes,.
very mundane work..
But he find the secret is, whenever he has time,.
he would do a little prayer..
And that little prayer is very simple..
It's like, "God, I wanna do this for your sake..
"And God, I wanna do it for your honor, and help me.".
It's just that simple prayer..
So he start his work in the morning,.
pray that prayer..
And then he find that it's helpful for him.
to be aware that he is doing everything.
and living in the presence of God..
And then he increase it by doing that prayer at noon,.
and then in the afternoon..
And then pretty much it becomes his second nature,.
that whenever he does the dishes,.
whenever he mopping the floor,.
he will pray to God silently, quietly,.

$^{601}$and pray that, "God, I wanna live in your presence..
"I wanna do this for your honor.".
That book, that secret,.
became one of the Christian classics..
Everyone finds it helpful..
Find rest in every single moment that you can in your day..
Even when you're driving,.
even when you're doing your work,.
pray that little prayer,.
and pray that, "God, I wanna live in your presence..
"Help me to be faithful..
"Help me to love you more.".
No matter what you do, changing diapers,.
mopping the floor, doing some mundane work,.
you can live in the presence of God..
So I wanna invite you to be more faithful,.
to really enjoy this sacred time,.
one day a week, first day of week..
It's not just coming here to follow the motions.
and follow the program,.
but you really come to God himself..
And secondly, on your daily basis, around the clock,.
starting tomorrow, Monday,.
do the little prayers..
Live in the presence of God, right?.
Knowing that God is always there with you..
Let's all pray together..
Father, we wanna give you thanks.
for the fact that you are ever-present,.
the fact that you are everywhere..
You are in the world..
You are especially in our heart..
And thank you, Jesus, for your cross,.
for the empty tomb that we can be called.
your sons and daughters..
Now pray that we will learn how to do Shabbat..
We'll learn every single day how to stop.
and to think about you, stop and pray to you..
And help us to learn how to worship you,.
one day a week, to cease our works of the six days.

$^{641}$and to draw closer to you..
And I pray that we wanna give you thanks.
for such a rhythm that you have set up in our life.
so that we will run in your order,.
we'll follow your order..
And God, be with us and be present with us.
so that we'll grow in our relationship with you..
In Christ's name we pray, amen..
Well, I'll see you next week, the next Shabbat..
\newpage



\section{}
\label{sec:6uSoGUnudqk}
\textbf{神的名字5\_公義的神 El Tsaddik (靈思牧禱)}
\newline
\newline
連結: \href{https://youtube.com/watch?v=6uSoGUnudqk}{\texttt{https://youtube.com/watch?v=6uSoGUnudqk}} ~~~~ 語音日期: 2024-03-26
\newline
\newline
\hyperref[sec:6DyYwyvwR5E]{\small{< < < PREV SERMON < < <}}
~
\hyperref[sec:index_chronic]{\small{[返順時目]}}
~
\hyperref[sec:index_scriptual]{\small{[返順卷目]}}
~
\hyperref[sec:lvZuhPuvg84]{\small{> > > NEXT SERMON > > >}}
\newline
\newline
$^{1}$[音樂].
今天在香港 政治動盪.
而在美國 也是嚴峻地政治動盪.
而且不單動盪 更加撕裂.
為什麼呢?.
其中一個原因是因為大家見解強烈分歧.
我們聊天的時候 很多時候都客氣尊重.
但是說到政治 我們就會很激動氣憤.
我個人認為 因為政治切中人心裡面一份公義感.
很多人認為 我覺得某某政客是獨裁欺壓的.
而某些也會覺得那個政客是貪污腐敗的.
誰是誰非 大家都認為自己是.
而當人覺得自己站在公義一方的時候.
傾向聽自己喜歡聽的 信自己喜歡信的.
恐怕就更加難尋求公義了.
而在這個短短的靈修時間.
我不是 也不能夠和大家去斷定.
哪一方是公義的.
我只是覺得 上帝的另一個名字.
很值得我們在這個時候思想.
就是El Shaddik.
中文譯為「公義的神」.
公義 Shaddik 希伯來文的原意是.
公平 公正的意思.
沒錯 聖經描繪的上帝是本性公義.
喜愛公義 行事公義 執行公義的神.
事實上 整本聖經都指出.
上帝是非常重視公義.
就如他頒給他的子民的律法裡.
都是貫徹地強調公義.
在《舊約聖經》記載.
上帝有無數次派遣先知.
斥責社會的不公.
這些先知他們就在公共場合.
好像聖殿 街市 皇宮那樣.
為神發聲 斥責權貴種種不公義行為.
例如 阿摩斯先知 清楚指出.
當時在耶羅安統治下.
安定繁榮的以色列.
竟然肆無忌憚地欺負貧寒人.

$^{41}$壓歲貧窮人.
而對於這些不公義的子民.
上帝就嚴厲指責他們說.
「我厭惡你們的節旗.
我不喜歡你們的嚴肅會.
你們雖然向我獻繁祭和掃祭.
我誤越了.
我要使你們歌唱的聲音遠離我.
因為我不聽你們的禽塞樂曲.
為願公平與帶水滾滾.
公義與江河滔滔」.
而在《新日聖經》也有類似的事.
《馬太福音》有這樣的詳盡記載.
當時耶穌在聖殿裡.
痛罵文士和法利賽人不公不義的行為.
他們在摩西的律法上自行加添了很多要求.
其中包括了繁瑣的規條和財產的賦賞.
但這些只是要求百姓遵守.
他們卻豁免了.
因此耶穌對於這些不公義的人.
說了一連串的「你們有何了」.
意味著他們不公義的行為.
必會受到神的審判.
沒錯,公義的上帝要求他的子民.
在生活裡行公義.
這個心意,即美加先知說得非常清楚.
當時神的子民在想.
我們應該獻上什麼才算敬畏神呢?.
是否獻上最上選的牛羊.
或者豐盛的橄欖油呢?.
還是我們要獻上我們親生的獎枝來贖罪呢?.
而美加先知就指責他們說.
「神的子民,難道你們還不明白嗎?」.
「在耶和華的眼裡,什麼是最寶貴的?」.
「就是你們要行公義,好憐憫」.
「全謙卑的心,與你的上帝同行」.
說到公義的神,有我們記住兩個重要特點.
第一,神的公義和誠信是分不開的.
這相當重要.
如果沒有誠信的本性.

$^{81}$公義其實是沒有意思的.
就像今天執法的人.
他們說一套做一套.
你怎麼期待他們是公義的呢?.
但是聖經記載.
每一件神所做的事.
都表現出他是100\%的誠信.
他說過的,他一定認數.
他的公義是令人信服和信靠的.
第二,上帝的公義和慈愛是分不開的.
這一點對我們更加有盜慾.
當我們說到公義的時候.
我們不要自欺.
其實我們每一個都不是二人.
縱然你和我未必是大奸大惡的人.
但是在犯罪這件事上.
也不過是50步少100步.
正如羅馬書第三章23節說.
世人都犯了罪,虧缺了上帝的榮耀.
不單止人類的始祖亞當夏娃是這樣.
歷世歷代萬國萬邦.
男女老少都是這樣.
也因為這樣.
我們必要受審判和刑罰.
這是神的公義.
但是美妙的是.
聖經又告訴我們.
唯有基督在我們還做罪人的時候.
為我們死.
神的愛就在此向我們顯明.
沒錯.
公義的上帝也是慈愛的.
於是祂會犧牲祂的兒子.
為人承擔刑罰.
為罪人開一條生路.
在這個因罪而殘缺敗壞的世界.
到處都會有不公義的事發生.
無論是國家,政治,工作場所.
甚至是在家裡.
可能都有強權肆虐,剝削,欺壓,欺凌.

$^{121}$而受害的人.
痛失去財產,尊嚴.
甚至是自由.
但是因為上帝是公義的神.
祂給我們帶來安慰和盼望.
我們相信祂必會按照祂的時間和旨意.
懲治惡人.
為受害的人討回公道.
即使今天不公義的事.
囂張不惑,張牙舞爪.
但是讓我們相信公義的上帝.
祂必不會讓邪惡永遠得勢.
公義的上帝必會懲治惡人.
替受害的討回公道.
公義的上帝也已經附上了.
祂的兒子耶穌基督.
拯救了一個因罪而殘缺敗壞的世界.
為罪人帶來一條新路.
而只要你謙卑真誠地相信這位兒子.
這個拯救也會令到你.
這是公義的上帝說的.
現在讓我們一起來祈禱.
上帝啊.
今天我們很多人.
在這個充滿不公義的世界裡面.
難免受到不公義的打擊.
無論是國家,政治,工作場所.
甚至家庭.
我們所受的或大或小.
你都一一看到.
我們要向你呼求.
因為你在歷史上.
不斷地向我們彰顯.
你是公義而誠信的神.
這個令今天受到打擊的我們.
得著莫大的安慰和盼望.
因為知道你必會施行公義.
為人取回應有的公道.
願今日在忍受不公義的事.
感到沮喪甚至痛苦的我們.

$^{161}$對你懷著堅定的信心.
忍耐守望.
我們更加多謝你.
因為你是公義而慈愛的神.
因為雖然你也會審判.
這個因罪而殘缺敗壞的世界.
你卻首先帶來拯救.
就是賜下你聖潔無罪的愛子.
耶穌基督.
為人的罪死在十字架上.
洗到墳為自己的罪.
謙卑懺悔.
轉向耶穌的每一個人.
都得到這個拯救.
你實在為了絕望的世界.
帶來了盼望.
榮耀歸給你公義的上帝.
我們禱告奉救主耶穌基督的命求.
阿們.
\newpage



\section{}
\label{sec:lvZuhPuvg84}
\textbf{Final Week of Jesus (Mark 14-15) - Rev. Mark Chan}
\newline
\newline
連結: \href{https://youtube.com/watch?v=lvZuhPuvg84}{\texttt{https://youtube.com/watch?v=lvZuhPuvg84}} ~~~~ 語音日期: 2024-03-27
\newline
\newline
\hyperref[sec:6uSoGUnudqk]{\small{< < < PREV SERMON < < <}}
~
\hyperref[sec:index_chronic]{\small{[返順時目]}}
~
\hyperref[sec:index_scriptual]{\small{[返順卷目]}}
~
\hyperref[sec:7tBg4MSn2jE]{\small{> > > NEXT SERMON > > >}}
\newline
\newline
$^{1}$Just a word of reminder, I'm looking forward to next week,.
Dr. Ko-Hee Sun will be with us,.
and maybe a few more people will join us as well..
So please come here early, all right, early,.
which is 11, and we're gonna share a portion.
of the whatever you call it, roasted pig, right?.
It's gonna be great, yeah, it's gonna be great..
I talked to the cafe team,.
and we originally planned to have a birthday cake.
so we can cut it together,.
but Ada, one of the heads of the cafe team,.
called me and said, "Oh, I compared the prices,.
"and roasted pig is actually the same price as the cake,.
"right, so why don't we just go for roasted pig?".
And okay, I said, "Yeah, why not?".
Since it will be more delicious, I guess, right?.
So I come here at 11, so we're gonna join together.
in the cafe and have a good meal together, all right?.
So today, we are going to have a short message.
on the importance of this week, which is the holy week..
Today is the Palm Sunday,.
that's why we have a brother getting baptized today,.
and seriously, you need to read his testimony..
He came from China, from a pagan worshiping background,.
and he was led into a government church,.
and he found it uncomfortable,.
and eventually, he went to a family church, or house church,.
and he was introduced to Jesus,.
and then he came over to the United States not too long ago,.
started to come to our church actually months ago,.
and then I was talking to him.
about his decision to get baptized,.
and I actually asked him if he minds to wait a little bit,.
you know, wait till Christmas,.
we may have a few more people getting baptized, right?.
Because we're going to prepare one water tank anyway,.
whether you have 10 people getting baptized.
or one person getting baptized,.
it's gonna be one big water tank, right?.
So I asked him whether he's okay to wait until Christmas,.

$^{41}$we're gonna have baptism in Christmas,.
and he said, "No, I wanna get baptized.".
I said, "Whoa, okay, all right, then how will you answer?.
"How will you answer such a desire.
"to get baptized just himself, right?".
So I talked to him, "You'll be just you,".
and I didn't tell him, I said,.
"In my mind," I said,.
"We're gonna prepare one water tank anyway,.
"so may as well wait for other people to get baptized.".
And he said, "No, I wanna get baptized.
"because God has washed away my sin,.
"and I wanna tell the world that Jesus is my Lord,.
"and dude, amen to that, right?".
So I had no answer, I said, "Yeah, let's do it, let's do it.".
So our team prepared the water tank for just one person,.
you know, one person getting baptized,.
and that is the most meaningful thing, right?.
We believe in Jesus Christ, we are in the name of Christ,.
and we receive the name of Triune God,.
Father and Son of the Holy Spirit,.
and the name being upon the person,.
and what a celebration, right?.
So read his testimony, it's very, very interesting..
So this week is all about this,.
it's all about the salvation Jesus has done.
for all of us on the cross,.
and that is the last week, the final week of Jesus..
Just very, very quickly, I wanna give you some survey..
How important this week is,.
the last seven days of Jesus Christ.
on his earthly ministry, his earthly life..
The last seven days is the most important thing.
recorded in the Bible..
All four gospels, Matthew, Mark, Luke, and John,.
spend over a quarter of its content,.
its pages, to record and to tell us the story,.
the crucified Messiah, Jesus Christ..
If you turn to the fourth gospel, the gospel of John,.
actually half of the book is about Jesus Christ.

$^{81}$on the cross..
He prepared himself to be nailed on the cross,.
and that is how important it is..
If you add those up, if you add those up,.
if you're good at math, you add those up,.
you've got 30-some chapters.
talking about just this final week of Jesus,.
almost one-third to one-half of the four gospels combined..
So this is the very, very important week.
that changed the world,.
changed the lives of the people who believe,.
all of us, each and every one of us..
So it's a good thing to remind us.
to reveal what is going on in this final week..
So we backtrack two days on Friday.
before the Palm Sunday,.
Jesus was annoyed by Mary at Bethany,.
and then it was Shabbat, resting day,.
and then Sunday is the Palm Sunday, which is today,.
and Jesus riding on a donkey,.
he enters Jerusalem as the king,.
and then he weeps for the city.
because the city of Jerusalem does not repent..
And in one generation in AD 70,.
in that year, the Romans would come.
and destroy the entire city and the temple..
And he weeps about the end, the judgment on the city,.
and then he returns to Bethany..
And then on Monday and Tuesday,.
he returns to Jerusalem and have debates.
with all these religious leaders who would not repent,.
who would not recognize him as a Messiah..
He also cursed a fig tree,.
which is a symbol of the Israel, the nation Israel,.
not willing to repent and bear fruit,.
which symbolizes the temple and the city of Jerusalem..
And the Jewish leaders plot to kill him..
So we come to Thursday and Friday,.
Jesus and his disciples share the last Passover meal,.
what we call the chameleon in later generations..

$^{121}$And he talks so much about things for the disciples,.
which we call farewell discourse..
So he's gonna say goodbye,.
but he's gonna have tons of things.
reminding the disciples, right?.
He's giving them the Holy Spirit to be with them,.
to walk with them so that they can grow in discipleship..
And finally, he prays at Gethsemane,.
which is a garden at the foothill of Mount Olive..
And he there was betrayed and arrested..
And it was in the middle of the night,.
he was on trial and condemned by the Jewish court,.
by the Jewish leaders..
And then Peter, as you remember, denies Jesus three times..
And then it was sunrise..
And Jesus was tried by the Roman governor Pilate,.
and was mocked by King Herod..
And then he eventually was crucified on God Gotham..
And at the time of his death,.
the veil of the temple,.
which separate the Holy of Holies and the Holy place,.
there was a long veil hanging down.
from the roof of the temple all the way down to the ground..
When Jesus died on the cross,.
that veil, according to the gospel of Mark,.
split into two from top to bottom,.
not bottom to top,.
from top to bottom, which is a miracle..
And the veil was open,.
because Jesus now is the center.
of people meeting God himself,.
not the stone building, not the temple,.
not in Jerusalem, but in Jesus Christ,.
no longer in the city of Jerusalem,.
no longer in the Holy temple,.
but in the person and work of Jesus Christ,.
so we can be saved and come to the Father.
and be restored to him in relationship..
And then Shabbat came,.
because he was buried in such a hurry manner,.

$^{161}$and the disciples or people who asked for Jesus' body.
from the cross, they just didn't have time.
to prepare the whole body,.
because Shabbat, there's no work,.
so they have to call it off..
And then they left the tomb,.
you know, happily prepare his body,.
and then it was Shabbat,.
for 24 hours, the Jews were not allowed to work..
And then Sunday morning,.
when the sun rises,.
the women go back to the tomb, why?.
Because they want to finish up the preparation.
for the body of Jesus,.
but they discover an empty tomb..
And the empty tomb can be explained in a lot of ways..
There are a lot of strange theories,.
you can explain away the empty tomb..
But the Bible give us the biblical reason.
why Jesus' tomb is empty,.
why only his tomb is empty,.
because by the power of the Holy Spirit,.
Jesus is risen from the dead..
Now, I've talked through the entire week.
in two and a half minutes, right?.
Three minutes, right?.
And this is the most important week of all in human history..
What Jesus has done in this seven days.
changed your life and changed my life,.
and changed the lives of many, many people.
from generation to generation.
in the past 2,000 plus years..
Jesus has changed the course of our lives..
Just as I told you,.
the person who received baptism today.
in the 10 o'clock service,.
he gave the testimony how God has touched his life,.
how Jesus has touched his life,.
and he wants to follow Jesus,.
and change the course of his life journey.

$^{201}$toward the kingdom..
And that is why the gospels, all four gospels,.
and even the gospel of Mark,.
spend so much energy,.
and so much verses and chapters.
on this past,.
on this holy week,.
especially on Thursday and on Friday..
So I want to take you a very, very quick survey.
of what is going on at the end,.
the ending two days,.
the last two days of Jesus' life..
If you turn with me to the gospel of Mark chapter 14,.
here I give you some sort of a map, right?.
I want to encourage you in the holy week,.
go back home in your devotional,.
spend five minutes a day,.
pick the best time of your day..
If you're an early riser,.
how many of you are early riser?.
You wake up, you feel refreshed all the time?.
None of you..
No wonder you guys are late..
(audience laughing).
No wonder you guys are always late, right?.
All right, how many of you,.
your best hour is at night?.
Best hour is at night..
Yeah, yeah, yeah..
Nothing wrong with that, okay?.
Nothing wrong with that..
But use it..
Use three minutes, five minutes,.
turn to the gospel of Mark chapter 14 and 15..
That is all about the Thursday and Friday,.
the last two days of Jesus' earthly life..
And here I give you a map..
Imagine yourself to be a movie director..
You're gonna shoot the scenes of Jesus' last 48 hours..
Here in chapter 14 of gospel of Mark,.

$^{241}$you see it's switched from A to B, A to B, A to B, right?.
All the A's are characters other than Jesus..
All the B's are pretty much the focus of the camera.
is on Jesus himself, of what he said and what he did..
A, you would start with 14, verses one and two..
The Jewish leaders trying to kill Jesus.
and they have a plot to do so..
So the focus is on the Jewish leaders,.
those who betrayed Jesus out of self-interest..
They're trying to protect their own authority..
The religious leaders in Jerusalem.
trying to protect their own power.
because Jesus is getting the popular movement going on..
So that's why they hate Jesus..
They don't want Jesus to take their place..
So they try to take away Jesus, trying to kill him..
And then B, the second part is Jesus is anointed.
at Bethany by Mary..
And Mary did that in order to prepare Jesus' death, right?.
You anoint the body, the dead body,.
in order to prepare the body for burial..
And Mary probably understood the Messiah has to die..
That's why she came with expensive perfume.
and trying to anoint the body of Jesus..
And then it's switched to A again,.
which is the third paragraph, is Judah, one of the 12,.
joins the plot to kill Jesus..
And then the camera shifts back to Jesus.
preparing for the Passover meal..
And then one of the 12, again, Judah,.
will betray Jesus..
And then switch to number six,.
which is the sixth section, verses 22 to 25..
Jesus shares the Passover meal with his 11 disciples..
And all the disciples will betray Jesus,.
and Jesus foretells that..
And then switch back to Jesus praying.
in the garden of Gethsemane..
And then the camera switch back to Judah, one of the 12,.
betrays Jesus, and all the disciples flee in the night..

$^{281}$They all escape, trying to save their own lives..
And then he was being brought in front of the Jewish court..
Jesus revealed himself to be the Messiah, to be the Christ..
And then lastly, we all remember this story..
Peter's trying to get warm.
and get into the court of the high priest..
And then he was being interrogated, being interrogated..
And he was afraid..
And he denied Jesus as his Lord,.
first time, second time, and third time..
And then everybody fled..
Jesus was alone..
That is chapter 14 of the Gospel of Mark..
And you see from top down, from top to bottom,.
all the people around Jesus trying to betray him,.
flee from him, to leave him alone, right?.
To abandon him..
From the Jewish leaders trying to kill Jesus,.
one of the 12 Jews trying to plot against Jesus,.
and all the so-called faithful disciples,.
they thought they would be strong enough.
to follow Jesus all the way to the end, but they did not..
They were fragile..
Their faith is shaken..
And they left..
From top to bottom, all the people around Jesus left..
But from top to bottom, we find that Jesus is determined.
to go to the cross..
Step by step, he prepared himself to go to the cross,.
to be crucified, shamefully, on the cross..
Because he understands this is the will of the Father.
to save humankind, to save each and every one of us..
Right at the center of the story,.
Jesus shared the last meal with his disciples,.
encouraging them and saying to them,.
"This is the covenant..
"This is the promise I give to you..
"As I die, as I shed my blood, and we drink this cup,.
"you know you belong to me, no matter what.".
And he foretold the disciples would betray him,.

$^{321}$but he also promised, "I will go to Galilee.
"after I resurrected from death..
"So I will restore you guys.".
Right?.
What a promise, right?.
In the entire story, in the entire chapter,.
no one can boast..
No one can boast his faith or her faith is strong enough.
to follow Jesus to the end..
Nobody..
Everyone around Jesus is powerless..
Everyone is going after his own self-interest..
Chapter 14..
In chapter 15, also the camera is switching back and forth.
from people around Jesus and Jesus on the cross..
First of all, B, all of B is about Jesus..
He revealed himself to Pilate, the Roman governor..
Now at this time, no longer the Jewish court,.
but the Roman court..
And then the crowd wanting to crucify Jesus,.
they asked for the release of Barabbas..
Barabbas was a criminal convicted in Jesus' day..
And the crowd want to release this guy, a rebel,.
but they want to crucify Jesus..
What injustice, you know, the greatest injustice.
is being done to an innocent man, Jesus himself..
But Jesus was silent because he accepted.
as the will of the father to die on behalf of all of us..
And then Jesus is proven to be innocent.
by the Roman governor, Pilate..
And mockingly, Pilate called Jesus king..
"Are you the king of the Jews?".
And Jesus said, "I am.".
Remember we did the Bible study of Gospel of Mark?.
Whenever someone says, "I am,".
it's not just a simple statement, right?.
It's the declaration of his divinity..
He is God, he's the son of God..
So he said to Pilate, the Roman governor,.
"I am the king of the Jews.".

$^{361}$But mockingly, Pilate, we got him as king..
And then the soldiers mock him as king..
He put on a purple robe, purple robe on Jesus and mock him..
And then at the center of the story, Jesus is crucified..
And then as we go on in this chapter, chapter 15,.
all these people around Jesus mock him as king..
And then finally Jesus dies..
A Roman centurion, a commander,.
proclaims at the foot of the cross,.
"Truly, this is the son of God..
"Truly, this is the king.".
And he was looking at Jesus on the cross,.
who shamefully died and sacrificed himself.
on behalf of sinful mankind..
This is kingship, this is leadership..
He died so that others can live..
He sacrificed himself so that we can be released,.
we can have freedom..
And truly, this is the son of God..
This is the true king, the servant king..
And then finally, the women watching from a distance..
And then at the end of this chapter, Jesus is buried..
Now, as you can see, very quick,.
from 14 to 15 chapter of the Gospel of Mark,.
Jesus becomes more passive and more passive..
He's taking on the actions of others..
People shame him, people mock him,.
and crucify him and kill him..
Jesus is no longer the active agent..
He's becoming more the object of the actions of other people..
Some Bible scholar has counted how many verbs.
are there in these two chapters..
And there are 50 some verbs,.
50 some actions done to Jesus..
He's being passive, the victim,.
until he's crucified and died and buried in the tomb..
Now, at the center of it all is Jesus on the cross..
Truly, he is the son of God..
That is the climax of the entire Gospel of Mark..
And it's very ironic, right?.

$^{401}$Now, we just had our primary election week, right?.
Did you vote?.
All those of you who could vote, you should have voted..
But we're looking for leaders, or political leaders,.
to bring us hope, right?.
To make some positive changes in our neighborhood,.
in our city, in our government..
We're looking for strong men and women..
We're looking for someone who is qualified to lead us, right?.
Who has the power, military power, political power,.
or willpower to make changes, right?.
To make people follow him..
Now, this is totally different..
The leadership of Jesus is the other way around..
It's the upside down leadership..
He died, he sacrificed himself so that we can be born again..
Here, I want to bring out two points..
The first point is,.
Jesus sees through the logic of the cross..
We kind of read through chapter 14 and 15.
of the Gospel of Mark..
And you see, what Jesus was doing.
is that he lays down his life..
He seems to be absolutely powerless..
He has no power in his hand..
He was nailed on the cross, he was totally passive..
But he alone fulfills the will of God..
He alone can save us from sin..
A sinless one became sinful.
so that we, the sinners, can be forgiven..
And that is why we call the Friday of the Holy Week,.
what do we call it?.
Good Friday..
Why is it Good Friday?.
Jesus died on the cross on Friday, right?.
There's nothing good about him..
But it was nonetheless called Good Friday.
because he has done the greatest great things for all of us..
The salvation on the cross..
That is the logic of the cross..

$^{441}$And Jesus is condemned according to the will of the Father.
so that we are accepted..
Jesus is punished so that our sins can be forgiven..
We can be redeemed..
He is cursed so that we can be blessed.
and we can be restored to God himself..
He suffers so that we have peace with God..
He dies so that we can have eternal life..
And that's the most amazing logic..
Nobody would have understood it.
until Jesus died on the cross..
And I think the life of Jesus being summarized.
in four gospels is to challenge us and to remind us.
if we say Jesus is our Lord and our Savior,.
then our lives should reflect the value of the cross..
What is the value of the cross?.
The logic of the cross is that we lay down our own will.
and we say to God, God our Father, your will be done..
Every day you pray, you do devotion..
The ultimate attitude we have is in front of God,.
we say, hey God, your will be done..
Ultimately, it's not my will..
Ultimately, it's not my interest..
But your interest, your kingdom's interest.
is the utmost important to me, right?.
Your kingdom come and your will be done..
That is the logic of the cross..
That is the logic of Christian life..
We come here to worship God,.
not because out of self-interest,.
because we know that Jesus should be praised..
God should be worshiped today..
We lay down our own interests.
so that we can exalt the name of God, right?.
That is the life pattern of all Christians..
All normal Christian life should be like this..
I asked you what's your best hour of the day, right?.
Probably it's the middle of the night, right?.
You do your own quiet time..
Now, quiet time is very important.

$^{481}$because we kind of calibrate our own heart to the kingdom..
We spend time with God so that God can speak to us.
when we read the word of God from the Bible..
And we also pray to him and we speak to him our desire.
and ask him to change us, ask him to lead us..
And that quiet time should be a calibration time.
so that we can calibrate our heart toward Jesus our Lord,.
so that our life can be constantly aligned.
with the discipleship pattern, right?.
The cross, which means thy kingdom come,.
thy will be done, not my will, but the will of the Father..
Now, I wanna point you to the logic of the cross.
by quoting from chapter 10, verse 45 of the gospel of Mark..
Here's one verse, and some of you already memorize it, right?.
"For even the Son of Man came not to be served but to serve,.
"but to give his life as a ransom for many.".
Some of you may have heard about this verse.
or even memorize it..
Here, Jesus summarized the entire gospel.
or his entire life and our life as Christian.
in just one verse..
Here he states two things..
He makes two claims..
He came for the purpose of serving, to serve other people..
And he came to die..
He came to die..
Now, when we do Bible study,.
sometimes we pay attention to some little words.
that usually we just don't focus on.
because the word is so little..
And I have a highlight for you..
And the word is actually,.
you can translate the word as "and," A-N-D..
It's a conjunction in Greek..
It's a very simple word, right?.
And he does that, and they do that,.
and something happened, and something happened..
The conjunction, right?.
You can translate it as A-N-D, "and.".
Simply, simple translation, literal translation..

$^{521}$But you can also translate.
because of the context requires you.
to translate it differently..
So here, all the English versions,.
English translations, render it as "even.".
You see that?.
Even..
Now, I was trying to dig through all the Chinese translations..
They all missed this, except one,.
the contemporary Chinese version..
It says, "Even the Son of Man.
"came not to be served, but to serve.".
Now, what's the effect?.
What's the difference?.
There's a huge difference.
if you translate this little word differently..
If you translate "and," you still get the sense..
But if you translate "even," it becomes more important.
because it tells you, when you translate as "even,".
even for this, for "even the Son of Man came to be served,".
or came not to be served, but to serve..
If you translate "even," the effect would be,.
it tells you,.
if anyone could reasonably be expected not to serve,.
but to be served by others..
It is the Son of Man..
It's Jesus himself..
Jesus reasonably expect all the other people around him.
to serve him, right?.
He's not here to serve other people, right?.
Think of yourself as king and queen,.
princess or prince, right?.
Think of yourself on a high status..
You expect other people to serve you, right?.
I have some friends in Hong Kong,.
they hire servants to do the housework, right?.
Do you have a servant in your household.
to manage the do-the-dishes?.
Maybe your husband is, right?.
Somebody do the work for you..

$^{561}$You expect someone to serve you..
Here Jesus said, "Even the Son of Man came not to be served.".
And the effect is, even he himself.
can reasonably, legitimately,.
require everyone around him to serve him..
But even the Son of Man came not to be served,.
but to serve other people, right?.
Think of the 24 hours each day..
You know, what is your attitude, right?.
Our normal, typical attitude as a human being.
is being occupied by our own self-interest, right?.
Our own self-desire usually, or most of the time,.
is the top priority, right?.
Even for us who come to church,.
our desire mostly are out of self-interest..
I come because I want to get something out of it, right?.
Oh, you got good food, okay?.
So I come, right?.
Oh, your food is crappy..
I don't come, right?.
I'd rather spend my time somewhere else, right?.
I come to church because I expect to be served, right?.
I'm preoccupied by my own self-interest..
People come and pray for me,.
or care about me, right?.
Or treat me well, right?.
That is our normal, typical expectation..
But here, Jesus said, "Even the Son of Man,".
even the King himself,.
"came not to be served,.
"but the other way around, to serve other people.".
So I want to remind you in this Holy Week,.
the life of Jesus, as all four Gospels emphasize.
how important these seven days,.
these seven last days for Jesus..
He came to serve us,.
not to be served by us..
And you would say Jesus is your Lord..
And if you say, "I am a Christian, I follow Christ,".
shouldn't our pattern,.

$^{601}$shouldn't our expectations be different?.
Our life pattern should be different..
To reflect a little bit more what Jesus is,.
to serve other people, and not to be served..
Think of one thing you can do for the church,.
for the church family,.
that you are not being served, but to serve other people..
One thing they all suggest, to pray for the church, right?.
Spend some time and pray for the church..
Pray for the church ministry..
People devote their time and energy.
in order to make it work, right?.
In order to run the church, in order to grow the church..
Pray for them..
And secondly, I would suggest,.
maybe find something you can plug in, you know?.
You can try to help out, you know?.
Maybe just talk to new people, right?.
Come in with a smile, as simple as that..
The pattern of discipleship..
Not to be served, but to serve, right?.
Pray for those who are preparing the service next week..
You know, Josh was leading worship today,.
and we spent a couple hours practicing, right?.
That's why he didn't sing our tune, right?.
He practiced..
He came to serve, not to be served, right?.
So that we can enjoy worship together, right?.
There are many, many things I would suggest.
that you have to think just upside down, right?.
Not to be served, but to serve other people, right?.
To do good to other people, right?.
Or even after the service, you know,.
go to the people in the backstage,.
encourage them, appreciate them..
'Cause they took their time, they took their energy,.
not out of self-interest, right?.
There's no benefit, right?.
But they serve, nonetheless..
Do something so that we can build up the church together..

$^{641}$Do something so that we can walk.
a little bit more like Christ himself, right?.
Let's do that, amen?.
Let's pray..
Father, we come before you in this holy week,.
on this Palm Sunday, to remember what you have done for us..
And thank you for being our sermon king..
Thank you for dying on the cross so that we can have life..
Thank you for experiencing the shame of the cross.
so that we can be honored..
And thank you for your redemption..
Thank you for granting us this new life.
so that we can reflect the value of the kingdom..
And Jesus, we come before you..
We ask that you change our heart.
so that we are not to be served, but to serve others,.
just like Jesus did for all of us..
Thank you for your love,.
and thank you for the new life we can experience,.
and help us to know how to follow you,.
to build up your body, to build up your kingdom,.
to extend your kingdom, to preach the gospel,.
to invite people into your kingdom..
There are so many different things we can do.
in the pattern of Jesus..
And I pray that you will lead us,.
help our church to grow spiritually,.
to grow in the image of your son, more and more..
In Christ's name we pray, amen..
Well, God bless you,.
and make sure you come a little bit early.
so that we can enjoy the eighth anniversary,.
I don't know what you call it, what, refreshment?.
special refreshment..
OK?.
All right..
\newpage



\section{創世記}
\label{sec:7tBg4MSn2jE}
\textbf{創世記第6章 (6.7) - 背景釋經 [方舟,再創造 2] (陳崇基牧師)}
\newline
\newline
連結: \href{https://youtube.com/watch?v=7tBg4MSn2jE}{\texttt{https://youtube.com/watch?v=7tBg4MSn2jE}} ~~~~ 語音日期: 2024-03-27
\newline
\newline
\hyperref[sec:lvZuhPuvg84]{\small{< < < PREV SERMON < < <}}
~
\hyperref[sec:index_chronic]{\small{[返順時目]}}
~
\hyperref[sec:index_scriptual]{\small{[返順卷目]}}
~
\hyperref[sec:IxOjYdkRpUg]{\small{> > > NEXT SERMON > > >}}
\newline
\newline
$^{1}$我們說到羅亞方舟的尺碼 尺碼不是亂來的數字 30乘50乘30乘50乘30 長方的box 一個盒 一個箱 不是一艘船 古代作者不是想回答我們的問題.
我們再繼續比較一下 希伯來聖經裡面.
有些東西好像是Mishapic裡面 近東的洪水故事會想做的事.
而希伯來聖經裡面有種所謂的模仿 我說模仿當然要聽清楚.
我說在近東文學的世界裡面 模仿是沒有問題的.
抄襲是沒有問題的 抄襲是我們現代人的概念.
所以如果你說羅亞方舟的故事是一個以正視聽.
正如我讀911的比喻一樣 以正視聽 錯誤的說法是錯的.
這些陰謀論是錯的 寫一本正式的911發生了什麼的書 記載.
以正視聽 所以從這個角度看羅亞方舟可能是在做這些事.
都是用近東的背景 但回應近東洪水的記載.
這個不足為奇 因為我們講了這麼久 比對研究等等.
有它自己的context 聖經有它的背景.
只不過是上帝藉著這個背景來啟示.
我們可以比較一下 也有不同的學者將這些數字.
我們講過這些數字不是隨機的數字.
有不同的學者指出 在摩西五經裡面.
只有兩個所謂神叫人造的建築物.
第一個就是方舟 第二個就是回望.
有沒有一些東西可以像近東的洪水故事.
講到方舟的時候 基金會的Shepard就說像廟塔.
他們的ideal 這些數字都是羅亞方舟的赤碼數字.
都是用一些Base 60 60進制.
再比較一下 回望的所謂外院.
大家知道這個回望有三個.
第一個就是最外面有蓮蓬的 有些簾遮著的外院.
裡面有祭壇 有水等等.
裡面就是聖所 也是祥方圓的.
再裡面就是正方形的致聖所.
所以回望有三個 長方形的外院.
聖所 最後最入的致聖所 藥櫃的地方.
就是正方形的.
很巧妙地 廣場一的赤碼就是長100 闊50.
再比較一下 方舟的大小.
其實它的平面面積就是匯膜的九倍.
你自己回去計算一下.
當然不會有高 因為匯膜只有一層.
在礦業裡面四處走 一個可接收的聖所.
不會有層數 但是到了第一聖殿的時候.
就有三層 方舟也有三層.

$^{41}$剛才看了聖經也知道.
所以你會看到作者可能用不同的手法.
來表達這個長方體是甚麼.
我不是說這個方舟就像匯膜 我不是這個意思.
我相信沒有很清楚的.
這些所謂經來的回響.
這些inner biblical echoes.
沒有很清楚 原因是甚麼呢.
這些數字不是亂來的 60進制.
有點像以色列人匯膜的形狀.
這裡告訴我們這是一個希伯來人所想的ideal.
這個ideal到底是甚麼呢 不是一棵樹.
到底是甚麼.
我們繼續看經文的時候.
或者可以幫助我們去理解.
到底洪水的故事是甚麼.
第十七節.
「看!我要洗洪水泛濫在地上」.
神說的.
「毀滅天下凡有生命氣息的血肉之軀」.
「地上一切都要滅亡」.
「但我與你 羅亞 立約」.
「你與你的兒女」.
即是你的兒子 妻子.
你息父 八個人進入方舟.
「凡有血肉的動物」.
「每養一對 一公一母」.
要帶進方舟裡面.
「好跟你 一同來到去保存生命」.
「飛鳥各種其類 牲畜各種其類」.
「地上爬行的動物又各種其類」.
「每養一對 都要到你那裡去保存生命」.
「你要拿各種可吃的食物」.
「儲存在你那裡」.
「作你和他們的糧食」.
「羅亞就去造了」.
「凡上帝吩咐他的 都照樣造了」.
這裡 摩西五經講到會幕時.
亦有這句所謂的「Fulfillment Formula」.
上帝吩咐摩西的 都照樣造了.

$^{81}$所以 這裡要想一想.
到底這些字眼 讓我們想起什麼呢.
當然 洪水滅世.
大水又回到岸邊.
我們自然一定會想到.
我們從第一章第二節已經看到.
大水又回到岸邊.
即是整個世界又重新來過.
所以我們神學上 叫它做「De-creation」.
然後 方舟造了之後.
亦都要帶一些東西進去.
你比較剛才說到方舟整的時候.
有不同的空間等等.
這個像什麼呢.
像創世記第一至第三章.
先做了一些forms.
一些空間.
第一天 做光與暗 時間.
第二天 創世記第一章.
第二天 水分開.
水分上和下 有空間出來.
第三天 海地露出來.
第一 第二 第三天 都是空間有關.
第四 五 六天.
就是與空間有一個平衡.
將一些光體 第四天就創造了.
第五天 天空上的飛鳥和水中的活物.
第六天 地上的動物.
最後 高潮的時候.
創造人類 按照神經現象.
所以 四五六天就是製造這些活物.
來充滿頭三天裡面有的空間.
方舟的故事 都有這樣的脈絡.
方舟造了 三層.
希伯來人近東的宇宙觀.
有三層 上 中 下.
不是只有聖所才是這樣.
所以 造方舟的時候就要造空間.
造完之後 水回頭.
大水再回來的時候.

$^{121}$就要將什麼帶入方舟.
就是將這些動物 人的代表.
地上的走獸等等.
去到保存他們.
所以 方舟不是走難這麼簡單.
是和世界有關.
方舟裡面有什麼.
有新的人類 羅亞和他一家.
地上的活物和飛鳥.
還有一個字 經常出現.
就是 各從其類.
各從其類這個字.
Min.
在舊的聖經裡面.
只有三處集中出現.
第一處 就是創世的第一章.
你記得了 這些我不需要重覆了.
什麼各從其類 什麼各從其類.
然後就是方舟這個故事.
經常提到各從其類.
用了Min這個字.
第三處 聚集 集中用這個字的地方.
就是食物的調例.
你未記 或者新年記.
裡面又是說到 潔淨與不潔淨的活物.
各從其類.
每次提到各從其類的時候.
一定是回到創造的規律.
上帝告訴希伯來人 以色人.
創造的規律是怎樣.
所以方舟的故事.
洪水的故事.
一提到這些字眼的時候.
一個熟知舊約 聖經的人.
他一定會回想 創造.
與創造有關.
剛才我們已經說了好幾樣東西.
與創造有關.
水 大水再淹沒回來.
然後有各從其類.

$^{161}$按道理 聖經的經文裡面提出的神學理念.
就是這個方舟是一個新的宇宙.
一個重組的宇宙.
再來過的一個Cosmos.
而這個方舟就是一個Mini Cosmos.
就是宇宙的縮影.
有三層 各從其類.
在大水裡面 有各種動物.
亦有新的人類再來過.
所以不是單單去避水這麼簡單.
不是單單去不要淹沒這麼簡單.
不是單單去救命這麼簡單.
所以不是救生艇.
是一個重新再組的Cosmos.
重新再組的宇宙.
所以上帝在做的一件事.
在洪水故事裡面 我們不可以錯過.
我們大茶經也好 港道也好.
講到這裡的時候.
我們不可以單單說救命.
很多水在來 有災難在來.
我們就救命 不用死.
絕對不是.
以保存生命 經常說這句話.
羅亞就好像一個新的亞當.
他去管理萬物 各從其類.
然後在上帝的宇宙裡面.
這個方舟 去管理.
上帝亦在洪水之後.
將新的世界帶來.
這是一個Mini Cosmos.
所以我們將創世記的記載.
和洪水的記載.
我們可以留意一下.
有沒有一些東西是平衡的.
這些平衡很明顯就是刻意的.
刻意的原因.
就是想表達一個神學的訊息.
就是上帝在作甚麼.
在近東的洪水記載.

$^{201}$其他幾家美紙等等.
就可以所謂以證視聽.
正如我所說.
大家在近東裡面.
都知道有洪水大災難.
問題就是這樣.
為甚麼會發生.
和所謂神在裡面有甚麼角色.
你將創世記第一章.
裡面的一些特別字眼.
你將它留意 強調出來.
然後你再將洪水的故事.
再將它比對的時候.
你會發現更加清楚.
創世記第一章.
我們在左邊就將創造的日數.
1 2 3 4 5 6 日.
就將它列出來.
然後經文的參考.
哪一節.
然後裡面最主要的元素.
內容是甚麼.
一章第二節 你記得.
就是有零這個字 有風這個字.
我們不妨去看看這裡.
創世記第一章第二節.
就在我上面這裡.
第二節.
第一節就是講起初神創造天地.
然後他這裡講到甚麼.
後面這個最後的分句.
就是Verruha Elohim.
Verruha意思就是風或是零.
我們通常翻譯就是神的零 上帝的零.
Verruha Elohim 神的零.
Merhevet Alpanei Ma'ayim.
就是他在拍翼 旋轉或是運轉.
Alpanei Ma'ayim.
Alpanei就是on the faces of.
即是在表面上 甚麼是表面上.

$^{241}$最尾就是Ma'ayim.
The waters.
所以這裡最後一句就是Verruha Elohim.
即是神的風.
就出現了.
然後第一章六至八節.
第二天就把水分開.
到第五天就有天空上的飛鳥.
然後第六天就有動物和人類被創造.
然後叫他們生養眾多等等.
你這樣看的時候.
你就發現洪水的故事是按照這樣脈絡去寫的.
第八章 洪水氾濫.
到最後第八章第一節.
神使風吹在水上.
神使風吹在水上 八章第一節.
首先第一件事.
我們之前說過的一個chaostic structure.
那個交叉的結構.
八章第一節就是轉折的地方.
上帝紀念諾亞.
Phaedrus Call Elohim at Noach.
Phaedrus Call就是Remember.
Remember不是說他沒有記性所以才會記.
你要知道Remember 紀念.
不一定是沒有記性才會記.
紀念是甚麼.
週年紀念 或者國殤日.
世界上一些很重要的事情發生.
很重要的劃時代的日子.
我們紀念.
所以不要以為人們很片面地看.
一定是上帝沒有記性.
很forgetful.
然後現在就記得.
哎呀 死了洛亞在船上.
No No No.
紀念很多時候就是上帝記起.
紀念他曾經說過的話.
和他的承諾 特別是他的約.

$^{281}$紀念他的covenant的時候.
他就開始行動.
所以紀念這個字是有神學意思的.
不是說他老懵懂忘記了.
現在記得了.
怎樣也好 這裡轉折的地方就是.
神紀念諾亞和方舟裡面一切的走獸和生俗.
中文就是 神使風吹地 水勢漸落.
神使風吹地.
風這個字又出來了.
Rock.
Arist.
神 Al-Him 這裡開頭第一個動詞.
Fire of Veil.
Fire of Veil 就是Cross the Pass的意思.
就是使他經過.
神使一些東西經過.
那是什麼呢.
就是Rock 風 或者零 吹在地上.
水就開始漸退.
這個是一個創造的動作.
你記得第一章第二節 都是一樣.
大水在那裡.
然後上帝又再使風或者零 吹過.
然後又說到兩方面的水.
天上的水和下面的深淵的水.
關閉 即是水停了.
然後分他出來.
這個就是八章第二節.
很明顯就是回應第二天的創造.
第三第四天就沒有什麼平衡.
到第八章十二節.
放鳥去偵察水有沒有乾.
為什麼提鳥呢.
故事是這樣發生就這樣提.
很明顯作者是有回應的.
有呼應的.
呼應第五天發生的事.
然後到第六天.
所做的事 原本創造的事.

$^{321}$就說羅亞的故事.
他就在方舟出來.
然後大地又再生養眾多.
遍滿這地.
這個很明顯又是一個呼應.
所以創世記第一章就說到一個創造的故事.
羅亞的洪水.
很明顯就是用方舟的故事.
說的是recreation 新的創造.
上帝看到本身的創造被罪所敗壞.
上帝的審判就來到了.
上帝公義的審判施行之後.
他又再造這個世界出來.
這個方舟就好像一個小宇宙一樣.
所以在說的是.
上帝在一個災難裡面.
他要做什麼 更新的工作.
重整的工作.
重整這個世界的秩序.
使它重新回到本身的美好.
當然我們繼續看創世記的故事的時候.
我們就知道原來新的亞當羅亞.
都是和本身的亞當一樣.
都是被罪闖入這個世界.
也是繼續敗壞下去.
正如我開頭所說.
創世記一至十一章是一個單元.
就是在說這個美好的世界.
因為罪 因為邪惡.
傳言敗壞了.
到第十一章最後的時候.
上帝到底怎樣去.
不放棄這個被破壞的創造界.
繼續用恩典來拯救.
你記得第十一章發生什麼事.
最後的時候.
就是有哈拉的後代.
阿不拉被呼召出來.
上帝就跟他說.
你離開你本族本地本家.

$^{361}$我要使萬國能夠得福.
救贖的計劃會繼續彰顯.
繼續去行動.
所以亞伯拉罕就被呼召出來.
就變成一個任務的故事.
一個使命.
一個抱天下.
叫人歸向耶和華的使命.
就從亞伯拉罕開始.
然後我們知道整個舊約.
其實就是追溯這個故事.
到最後當然就是耶穌基督.
所以一定要掌握到.
這個theology是什麼.
這個theology不是我們現代人所想.
整隻船出來救命.
是上帝正在做一些重新的事.
重新這個世界.
如果你這樣看的時候.
或者可以給我們更多現代的神學訊息.
怎樣看一個disaster.
怎樣看一個災難.
很多時候我們說一些災難包括疫情.
好像一個洪水.
在一個災難裡面上帝想做什麼.
這裡起碼有一個神學的角度.
讓我們看到原來上帝會藉著洪水.
重新創造世界.
重整創造世界.
問題是我們是否有參與他的工作.
我們是否終於去到obey.
去到順從.
去到參與他的工作.
參與他重整世界的工作.
暫時我們先說到這裡.
如果你喜歡我們的YouTube頻道.
就記得Like,Comment和Subscribe.
\newpage



\section{馬可福音}
\label{sec:IxOjYdkRpUg}
\textbf{耶穌最後的一週  \#3 (馬可福音 14-15 ) - 陳崇基牧師}
\newline
\newline
連結: \href{https://youtube.com/watch?v=IxOjYdkRpUg}{\texttt{https://youtube.com/watch?v=IxOjYdkRpUg}} ~~~~ 語音日期: 2024-03-27
\newline
\newline
\hyperref[sec:7tBg4MSn2jE]{\small{< < < PREV SERMON < < <}}
~
\hyperref[sec:index_chronic]{\small{[返順時目]}}
~
\hyperref[sec:index_scriptual]{\small{[返順卷目]}}
~
\hyperref[sec:_seeVdxX33U]{\small{> > > NEXT SERMON > > >}}
\newline
\newline
$^{1}$我們繼續聖周的講道.
Holy Week.
今天剛才敬拜隊已經提醒我們.
今天是Palm Sunday.
中樹節的主日.
所以是一個很重要的.
在基督教的日曆裡面.
一個最重要的聖祈.
就是耶穌在地上的最後一個星期.
The Final Week.
我們有兩次講到這個主題的時候.
我們一直都在跟進.
就是在這個Final Week裡面.
發生什麼事情.
但是我這裡想大家留意.
在我們這個聖周Holy Week.
或者Passion Week的時候.
我們希望大家回家去看看.
聖經裡面四福音怎麼去記載.
耶穌基督最後最後一個星期.
但是這裡很簡單.
我只是將一些章節總結出來.
就是每一個的科林書.
馬太,馬可,路加和約翰.
都花了很多篇幅.
來記載最後這七天.
耶穌發生了什麼事情.
The Final Week of Jesus.
馬太21章至28章有8章的聖經.
就是整本馬太福音裡面的四分之一以上的篇幅.
就是在講最後這七天.
馬可福音11至16章.
6章也是多於四分之一的.
路加也是多於四分之一.
約翰更加不需要講了.
約翰你將他斬一半.
下面那半全部都是講他最後最後這七天.
所以內容的比重是很重要的.
為什麼呢?.
保羅不只是他才說.

$^{41}$說我們傳的是釘十字架的耶穌.
福音書也是這麼說的.
有很多篇幅.
最後這七天就是最重要的.
如果你加起來的話.
整本新約的聖經只是計算四福音的話.
有三分之一至一半的內容.
都是講到最後這七天的.
所以我們不妨看看這七天到底發生了什麼事情.
這個表我們三個星期都有看.
不用很留心地看.
因為字很小.
你只需要留意左邊.
有一個星期.
大概是這樣.
我們講到中樹的主日就是第一個星期日.
Palm Sunday.
耶穌永進耶路撒冷.
然後第二第三.
Monday to Wednesday.
就是耶穌去到耶路撒冷.
和這些文士,祭司長.
那些權貴去爭辯.
然後今天我們去到大概星期四至星期五的時間.
就是耶穌去預備最後的晚餐.
道別的廣橫.
然後他去到鑒山的山腳.
一個地方叫做喀西瑪利園.
他去祈禱.
不按照我的意思.
按照天父的意思.
然後他被出賣.
被捕.
然後就被審查.
門徒全部都雞飛狗走.
然後.
耶穌接受不單是猶太人議會的審議.
也是接受羅馬人議會的審議.
就是彼拉多.
於是他就被判釘十字架.

$^{81}$死在十字架上.
然後聖殿的萬子上下分為兩半.
耶穌就被埋葬.
然後就來到了.
然後就到了下一個星期日.
就是Easter Sunday.
就是復活節的主日.
門徒.
特別是女的門徒先發現這個空墳墓.
然後他們就將這個消息告訴門徒.
這個就是很簡單.
這個星期.
聖周Holy Week裡面發生過的事.
有什麼章節可以自己回去看一下.
我們在三大平台裡面都有發放出來.
但我想你知道的就是.
那個重要性.
就是耶穌最後這個星期的重要性.
我們不妨花一點時間.
只是看馬太福.
後面那兩張.
第十四張和第十五張.
我知道有很多的經節.
大家覺得很辛苦.
在我們中間有誰是讀書人.
沒有吧.
每個人都文盲.
我們不喜歡讀書的.
不要緊的.
我告訴你.
我拋磚引玉.
希望你回去試試讀書.
回去試試看.
只是馬太福第十四張和第十五張.
我們不需要很仔細地看.
我們三個星期都是這樣.
我們說不如放大一點.
將鏡頭調遠一點.
我們先看一下整個森林.
整個圖案.

$^{121}$整個分佈.
整個故事的結構是怎樣.
然後從這個結構可以看出一些訊息.
到底他怎樣設計.
寫出最後七天.
特別是最後兩天.
耶穌的故事.
馬可福音第十四章.
這裡你這樣看一看.
有中有英.
左邊你看到.
好像梅花建築一樣.
有A有B有A有B.
由A開始就是猶太人的領袖.
來設計殺害耶穌.
十四章一到第二節.
然後耶穌再到伯大尼.
這個索比爾的地方.
索比爾的地方被告.
然後就到第三段.
十二使徒之一.
猶大.
來設計出賣耶穌.
他就加入那些祭司長.
猶太人的領袖.
來設計殺害耶穌.
然後回到鏡頭.
又放大.
又將鏡頭轉回耶穌.
耶穌就預備最後的晚餐.
跟門徒一起吃.
然後鏡頭又放在猶大.
耶穌以外的人.
猶大十二使徒之一.
將會出賣耶穌的預言.
然後第六段.
耶穌就跟十一個門徒.
來分享最後的晚餐.
然後鏡頭又轉回.
又不是耶穌.

$^{161}$是耶穌身邊的.
眾門徒.
就被預言.
會否認耶穌.
會出賣耶穌.
然後鏡頭又回到耶穌那裡.
耶穌在黑西瑪利園祈禱.
然後十二使徒之一.
猶大出現了.
在黑西瑪利園夜間的時候.
就出賣耶穌.
然後眾門徒就逃亡.
然後又回到B.
就是耶穌在猶太人的審訊裡面.
去啟示他自己.
他就是人子.
就是彌賽亞.
就是基督.
然後最後.
就是彼得.
走剩遠遠的跟隨耶穌.
連他都三次不認主.
你看到這裡梅花間中.
好像拍戲一樣.
鏡頭.
你是導演.
就將鏡頭轉來轉去.
每逢是A的地方.
都是放在其他人身上.
出賣主的.
祭司在民事化來創人.
門徒.
然後最後是彼得.
B就永遠看完其他人的時候.
又回去看耶穌.
耶穌無論是有話說也好.
或者有些行動也好.
那個焦點都是放在耶穌身上.
你會看到一路由上至下.
整個十四章的故事脈絡.

$^{201}$好像講故事一樣.
門徒開始的時候就很主動.
幫耶穌預備晚餐.
在上樓房預備好筵席.
越來越到下面.
就是被動.
或者門徒開始的時候.
他們是很雀躍.
或者很主動的跟隨耶穌.
但最後的時候.
他們就越變得脆弱和無能.
全部人.
開始是很有力量.
但到最後的時候.
他們都雞飛狗走.
全部走了.
全部都是脆弱無能的.
就在中間這個地方.
耶穌就跟他們設立晚餐.
而晚餐不是吃這麼簡單.
晚餐是去到應許承諾.
承諾門徒.
承諾我們這一群.
是會脆弱無能的人.
耶穌說我用我的血去立新約.
你就是我的子民.
而且從上至下你見到耶穌.
是不斷堅定的.
去步向一個目的地.
那個目的地就是被出賣.
走向十字架.
門徒.
就是脆弱.
耶穌是一步一步走向十字架.
去進行上帝在他身上的旨意.
這個就是十字架的生命.
十四章講完了.
是不是很容易.
等你回去看聖經.
你跟著看這個很過癮.

$^{241}$第十五章.
第十五章很奇怪.
又是這樣.
又是梅花間褶.
B的全部都是耶穌.
A全部都是耶穌以外的人.
你見到攝影機在轉動.
開頭耶穌就在羅馬的審訊.
就是彼拉多面前.
去到自我的啟示.
我就是猶太人的王.
然後就把鏡頭轉去巴拉巴.
那些群眾就說要釋放巴拉巴.
不釋放耶穌.
巴拉巴是一個強盜.
是一個犯罪犯.
他們寧願有一個強盜釋放.
就不要無罪的耶穌被釋放.
然後鏡頭又回到B.
耶穌就被彼拉多來審核.
發現他無罪.
但就稱他為王.
很ironic的.
是諷刺地說你就是王.
然後羅馬的兵丁就戲弄耶穌.
把他掛上紙袍.
稱他為王.
然後耶穌就被釘十字架.
在十字架下面發生什麼事呢.
龐人和猶太的領袖就戲弄耶穌.
稱他為所謂王.
然後耶穌就在十字架上面棄絕身亡.
羅馬的伯夫長.
竟然宣告在馬可福音中最重要的一句話.
他說:這真是上帝的兒子.
這就是上帝所設立的基督.
這就是上帝的君王.
死在羞辱的十字架上.
這就是The Son of God.
然後婦女遠遠地望著十字架.

$^{281}$然後最後第九段.
就是耶穌被埋葬.
我們看到十四十五章.
耶穌一路走向十字架的時候.
他就越來越做一個Victim.
一個無能者可以這麼說.
一開始他設計了所有的東西.
還有引出所有的行動.
但是到最後的時候.
他就接受身邊的人.
那些釘他十字架,凌辱他的人.
所有的行動.
如果你查一下這裡大概的動詞.
有人查過.
耶穌接受的動詞是五十多個動詞.
他做的事或者他有的動詞.
只有九個.
就是說他越來越變得被動.
因為他要上十字架.
他是被Crucified.
然後被埋葬.
在十五章的中間.
就正正是耶穌在十字架上死.
在十字架上他受苦.
受苦之後.
死之後.
他就被稱為神的兒子.
這個才是上帝所派來的王者來拯救我們.
在聖咒裡面特別是最後這兩天.
你會發現有很多不同的角色出現.
這些角色包括.
有一些猶太人的領袖.
一些Jewish leaders.
祭司長,文士,大祭司等等.
這些都是以自身的利益為上的人.
他們見到耶穌很受歡迎.
然後他們自己的地位受到威脅.
於是他們就想出這個方法.
希望用這個方法來剷除耶穌.
來跟隨耶穌的群眾.

$^{321}$所以他們要釋放巴拉巴.
這是第一班人.
第二班人就是那些我們可以叫做.
morally indifferent的人.
或者道德上,道義上.
是站在旁邊的.
隔岸觀火的人.
這個大概是彼拉多的描述.
你不要以為彼拉多是好人.
有個人在第一世紀寫了很多東西.
叫做Josephus.
他寫了很多有關彼拉多的東西.
這個人是很賤格的.
但為什麼會在這裡.
他就好像裝作好人.
其實不是的.
他只不過假意.
其實都是為了自己的self-interest.
因為巡撫或者一個羅馬的總督.
去守一個縣,一個省.
就好像特首一樣.
特首最主要要做的是什麼.
不是要在那個地方.
使他興旺繁榮.
總之不亂就行了.
總之大家平安.
我有茶喝,你有茶喝.
大家不要出聲.
中央就不會下來打壓我們.
這個就是彼拉多要做的工作.
所以他見到耶穌的時候.
他都不是很care的.
為什麼.
因為他都是自我利益為上.
都滿途了.
這班跟隨耶穌的門徒說.
誓死都會跟他.
誓死的都不會否認他的.
到最後的時候.
他們又是各散東西.

$^{361}$因為要保住自己的命.
唯有耶穌自己一個.
我們的主.
他就走上十字架.
被釘十字架.
好像一個無能的人.
但卻是因為他.
去到跟隨天父的旨意.
以至他能夠改變世界.
一個無能的人能夠改變世界.
一個無能的人能夠去.
改變許多人的生命.
改變許多人生命的軌跡.
耶穌願意受刑罰.
讓我們在上帝面前能夠被治下.
他願意受審判.
以至我們能夠聚得赦免.
他願意受咒詛.
以至我們能夠有祝福.
我們能夠與上帝復和.
他願意受苦.
我們能夠有平安.
他願意受死.
我們能夠有永生.
這就是耶穌將無能的生命.
帶在十字架上為我們成就.
改變世人的命運.
改變我們的命運.
我們今天這個中樹主人的浸禮.
我們本身都想將它押後.
因為我們兩堂裡面.
都好像沒有人想浸禮.
還沒有預備好.
我們就去問了幾個.
有幾個有意思的.
但個個都是等下再說.
唯獨小優就決定.
她說一定要受浸禮.
我們同工就跟她說.
我們弄一盆水.

$^{401}$等十個就可以了.
因為Ronald弄了一整天.
又要煲水煲一天.
水在白鷺鷥那邊亂的.
我們問了小優.
她說你想不想等到聖誕節的時候再浸.
她居然說不肯.
她說不行.
我一定要接受水禮.
她才跟我說.
不行 如果她真的決定要接受水禮.
為什麼要攔阻她.
怎樣辛苦都好.
待會白鷺鷥幫幫手.
怎樣辛苦都好.
一個罪人悔改.
宣認耶穌在她身上.
是她的王 是她的主.
這個是不是最美的聲音.
這個真是最好聽的見證.
所以我不知道Ronald弄一盆水的時候怎樣.
是不是覺得怎麼搞的一個人弄這麼多東西.
但你知道這個是無價的奉獻.
一個罪人願意去接受水禮.
去到說耶穌在我身上就是作主.
因為救恩是無價的.
因為耶穌十字架的救恩是沒有價錢的.
放不到價錢 是無價的.
在這最後我想點出兩點.
就是主為我們所做的.
就是他們看清楚認清楚.
看透這個十字架.
羞辱十字架的邏輯.
邏輯是什麼.
他受羞辱以至我們能夠有尊榮.
他受苦以至我們能夠得到生命.
這個就是十架的非世人的邏輯.
世人看見十字架的時候.
就說這個是絆腳石.
或者世人最愚蠢的東西.

$^{441}$一個流血的十字架.
一個死刑.
一個犧牲.
一個羞辱.
但是在上帝眼中.
這個就是上帝的智慧.
就是拯救我們這班貪生怕死的人.
從罪裡面能夠得以釋放.
這個就是十架的邏輯.
就是耶穌看透這個邏輯.
以至他願意馴服天父的旨意.
他就上十字架.
去改變我們的生命.
第二方面.
就是耶穌的生命去到提醒我們.
怎樣活出十字架的邏輯.
十字架就是說.
不照我的意思.
最終要照上帝的意思.
我們生命裡面有不同的際遇.
有不同的問題.
但我們問一個終極的問題.
就是我們不是按照我的意思.
上帝要按照你的意思.
是我會掙扎.
但我最終會問.
到底上帝要我在這個情況裡面.
我怎樣去尊崇他.
不照我的意思.
照上帝的意思.
在馬科宏第十章四十五節.
這一節的聖經.
一節而已.
可以這麼說.
將整個福音.
整個基督徒的生命.
耶穌自己的生命.
做一個總結.
他說因為人子來.
並不是要受人的服侍.

$^{481}$乃是要服侍人.
並且要寫明.
作多人的屬格.
我們剛才說到十字架底下.
有很多不同的角色.
每一個都是以自己利益.
自己的意思先行.
自己的自我興趣先行.
我們現代人都會這樣.
有一種obsession.
或者一種preoccupation.
就是我自己的東西先行.
我現在的處境.
我現在的感覺.
我現在的聲音.
但耶穌說什麼呢.
他說人子來不是要受人的服侍.
乃是要服侍人.
並且寫明作多人的屬格.
在這裡你會留意到有不同的聖經翻譯.
我們曾經有一些查經班.
查經小組.
我們經常說我們對照一下.
其他的翻譯本有什麼東西看.
但其實很多時候看不同的翻譯本.
是沒有用的.
因為是同一回事.
但有時候翻譯本是有意思的.
就好像這裡.
你見到英文全部都翻譯成.
For even the son of man.
For even the son of man came not to be served but to serve.
中文我找遍了.
很多都沒有翻譯出even.
接著我就找一下當代翻譯本有沒有.
因為當代翻譯本是跟NLT的.
英文的NLT.
它是怎麼翻譯的呢.
我覺得是翻譯到那個的意思.
因為就連人子也不是來受人服侍.

$^{521}$那個字是很小很小的字.
希臘文叫Kai.
其實可以直譯就是and.
就是連接詞.
接著發生什麼事.
還有發生什麼事.
你可以這樣直譯.
但上文你也要了解.
為什麼要做就連呢.
even.
因為耶穌說這句話之前.
就是有門徒去到他面前來求他.
他說主你德國的時候.
可不可以我們一個坐在你左邊.
一個坐在你右邊.
約翰和雅各.
他們十二門徒裡面.
他們想爭最大.
這個是我們的意思.
你可不可以照我們的意思做.
當你德國進入Glory的時候.
你有榮耀的時候.
你有威的時候.
讓我坐在你左邊和右邊.
耶穌就說這句話.
for even the son of man.
就連人子也不是來受人服侍.
乃是要服侍人.
就是這麼簡單的意思.
那效果是什麼.
效果就是告訴我們.
人子是應該受人服侍.
人子是大條道理.
應當受人服侍.
最reasonable他expect的東西.
就是to be served.
and not to serve.
最近不是有個新聞.
哪個公主.
英國.

$^{561}$是不是有癌症.
她真是好.
因為她的醫藥保險一定很好.
她是受人服侍.
她可以受人服侍.
應該受人服侍.
因為她有皇室的尊榮.
但耶穌說就算就連人子.
她都不是要受人服侍.
朕在開始的時候提到.
人子是什麼意思.
在舊約裡面.
son of man就是.
那個戰勝帝國.
那位皇者.
有天上地下權柄的皇者.
就連這位皇者.
人子.
來到都不是要受人服侍.
而是服侍人.
在聖咒裡面.
求上帝一巴掌提醒我們.
我們接受了耶穌作為救主的人.
接受了這位人子來服侍的人.
我們應該怎樣去活出十字架的生命.
告訴世人聽.
這位主值得我們敬拜.
值得我們事奉.
祂救我.
我就好像祂一樣.
不是受人服侍.
而是服侍人.
我們請敬拜隊上來.
我們唱最後回應的詩歌.
我們唱的時候.
希望你去思想.
上帝去請教我們.
我們這個生命.
屬父他.
不是要受人服侍.

$^{601}$而是服侍人.
當你有這樣的邏輯的時候.
我相信你的個人生命.
和你自己家庭的生命.
任何一個角色.
都會有基本上的改變.
我們的心態.
我們的取態上.
都會有不同.
用我們的生命去服侍人.
不是要受人服侍.
(前透).
我們同心低落祈禱.
主我們為了我們的生命.
一個罪人的生命.
能夠得著你的救恩.
你愛我們的聲音.
十字架上都有我們的恩典.
今天我們站在你面前.
就如站在十字架的面前一樣.
我們的主耶穌基督.
真正的彌賽亞.
真正的皇者.
願意以愛以犧牲.
來救贖我們.
主我們的生命.
是有著你的名字.
求你讓我們的言談舉止.
我們的生命.
我們的行為.
我們的取態.
都猶如我們所跟隨的耶穌一樣.
帶領我們.
差點我們離開這個地方.
去到我們自己生命裡面.
每一個的崗位.
大大小小的角色.
我們能夠去演繹出十字架的道理.
演繹出服侍人生命的道理.
主求你賜力給我們.

$^{641}$讓教會能夠成為你美好的見證.
我們這樣祈禱.
奉主的名.
各位請坐.
\newpage



\section{6:25-34}
\label{sec:_seeVdxX33U}
\textbf{Seek Him (Matthew 6\_25-34)- Rev. Mark Chan}
\newline
\newline
連結: \href{https://youtube.com/watch?v=-seeVdxX33U}{\texttt{https://youtube.com/watch?v=-seeVdxX33U}} ~~~~ 語音日期: 2024-03-29
\newline
\newline
\hyperref[sec:IxOjYdkRpUg]{\small{< < < PREV SERMON < < <}}
~
\hyperref[sec:index_chronic]{\small{[返順時目]}}
~
\hyperref[sec:index_scriptual]{\small{[返順卷目]}}
~
\hyperref[sec:87VD_Eud_tU]{\small{> > > NEXT SERMON > > >}}
\newline
\newline
$^{1}$Good morning, welcome to our worship service..
This is the first sermon of the year 2023..
As we say goodbye to 2022,.
welcome to the new year 2023..
We are probably thinking about a lot of different things..
The Christmas season was just over..
Probably your life or your house was in chaos or in a mess.
and you're trying to fix it.
and trying to put your house or your life back in order..
The vacation mode is over.
and we're back to the so-called reality..
And we should thank God in our calendar,.
in our daily life pattern,.
there's always time and always opportunity.
for us to say goodbye to the old and welcome the new..
And for us to transition from one period to another period..
I think this is a God sent moment.
for us to think about what our lives should be.
and where our life is heading to..
In the new year, many of the people, including ourselves,.
are probably thinking about some resolutions, right?.
Some new year's resolution, probably losing weight..
Any of you are thinking about this?.
Or quit some of the bad habit, going back to gym,.
paying off some debts or saving more money,.
building up a better or good habits..
This kind of new year's resolutions..
And I keep thinking about this question..
I'm sure you do as well..
What, as questions, should be thinking about.
in terms of new year's resolutions?.
As we transition from one period to another period,.
as we say goodbye to the old and welcoming the new year,.
what kind of new year's resolution.
should we be thinking about as the people of God?.
So I have chosen a fairly familiar passage.
to share with you, to reflect upon,.
so that we know what Jesus wants us to do..
What kind of new year's resolution.
should we be keeping in mind at this time of the year?.

$^{41}$So as we welcome the new year,.
why don't we just read through this passage.
and see what we can take out of this passage?.
What kind of lessons we can learn?.
Here is Matthew chapter six..
Now, Matthew chapter six is a familiar passage..
It's actually the, sort of the conclusion.
of Sermon on the Mount,.
the famous, well-known sermon of Jesus.
giving to his disciples while on earth,.
teaching us what kingdom living is all about..
So right at the end, almost the end of this long sermon.
from Matthew chapter five, all the way to chapter seven,.
we come to the end of chapter six..
Six, 25, it says, "Therefore I tell you," Jesus says,.
"do not be anxious about your life,.
"what you will eat or what you will drink,.
"nor about your body, what you will put on..
"It's not your life more than food.
"and the body more than clothing..
"Look at the birds of the air..
"They neither sow nor reap.
"nor gather into bonds,.
"yet your heavenly Father feeds them..
"Are you not more value than they?.
"And which of you, by being anxious,.
"can add one single hour or cupid to his span of life?".
So Jesus said, "Don't you worry about your essential needs,.
"the money, your bank account, your investment,.
"all this earthly concern that we have.".
Now, Jesus is not saying that we should be irresponsible..
That's not the teaching of the Bible..
Jesus is saying our mind and our heart.
should not be so fixed upon things like these..
He says, "Look at the birds of the air..
"The Father is feeding them..
"They don't work as hard as we do..
"They don't gather into bonds, they don't harvest,.
"but the heavenly Father is feeding them.
"and keeping them alive.".

$^{81}$And he compares that to our lives..
"And which of you, by being anxious, just worrying,.
"working so hard, right, or overworking,.
"how many of us can do this, being anxious,.
"can add one single hour to the span of our life?".
Now, the original word behind the word hour is cupid..
Cupid is the arm length of the king,.
most probably Egyptian king, Pharaoh,.
using his arm as the standard of measurement.
in the ancient days..
Cupid is arm length, is the smallest unit of measurement,.
smallest unit of length..
So Jesus is saying, by having anxiety,.
cannot add a single cupid, cannot add the most basic,.
the smallest unit to the span of life of hours..
So he keeps on saying,.
"Why are you anxious about clothing?".
Right, what we wear, the outward appearance,.
the outward things..
He said, "Look at the lilies of the field, how they grow..
"They neither toil or spin.".
But I tell you, even King Solomon in all his glory.
was not to be compared to one of these,.
to the bird or the wild flowers, the lilies of the field..
"But if God so clothes the grass of the field,.
"which today is alive, tomorrow will be gone.
"and thrown into an oven,.
"will he not much more clothe you,.
"O you of little faith?.
"Therefore, do not be anxious, saying,.
"What shall we eat?.
"What shall we drink?.
"What shall we wear?".
For the Gentiles, those people who do not follow Jesus,.
do not know God, they seek after all these things,.
all this eating, drinking, clothing,.
but your heavenly Father knows what you need,.
that you need them all..
But seek the kingdom of God and his righteousness.
and all these things will be added to you..

$^{121}$Now, therefore, do not be anxious about tomorrow,.
for tomorrow will be anxious for itself..
Enough for the day, sufficient for the day.
is its own trouble..
Now, this is the one unit, one of the units.
of the Sermon on the Mount,.
talking about how the disciples should be living..
What is the disciple of Christ should be seeking after all..
Jesus is comparing some of the things.
that those people who don't know Jesus,.
what they're seeking, what they're going after,.
and compared to what a disciple of Christ.
should be seeking and pursuing..
I remember years ago,.
I left Israel after my trip..
It was about more than 10 years ago..
I left Israel and I was waiting for the flight.
in the Tel Aviv airport..
So I was walking around in the hallway of the airport.
and then I bum into a map.
outside of the Elor Airlines..
Now, the Elor Airlines, as you can see here,.
is the national Israeli airline.
and it's one of the most well-known airlines in the world..
So I went by the counter.
and then I saw this map..
You know, this map, I was pretty curious,.
why is there a map like this by the counter of Elor?.
Can you guess what this map is all about?.
What is this map?.
Now, I have some options here..
A, intercontinental highway system..
So this Elor is intercontinental Africa to Asia or Europe,.
probably a highway system road map.
or the airline route map.
or none of the above..
Which one do you think this map is all about?.
Is it the highway system?.
Is it the route map?.
Or is it none of the above?.

$^{161}$Now, the answer is, it's none of the above..
It's not the highway system..
This is not a map of the Elor airline route map..
And now what this is all about.
is actually routes or map of routes.
of the migratory birds in Israel..
Migratory birds, as we know,.
they fly from south to north and north to south,.
depending on the season..
When it is in the winter, it's getting so cold,.
the birds will migrate down to the warmer climate, right?.
In the continent of Africa..
So this map is very interesting,.
just right outside of Elor airline back then 10 years ago..
It's not there anymore..
So I took a picture..
So I took this picture because as I saw this map,.
I immediately thought of Matthew 6,.
when Jesus said, "Do not worry about what you eat.
"and what you drink.".
Look at the birds of the air, check them out,.
see how they fly..
And I immediately thought of this passage.
of the Sermon on the Mount,.
because Jesus was saying,.
referring to the birds of the air..
And this is not by accident that Jesus knew something.
about the geography of Israel,.
because the birds are traveling north to south.
and south to north..
They have to travel through the narrow strip of land.
of Israel Palestine..
As you can see on the left or the west side.
of this long narrow strip is the Mediterranean Sea..
And on the right is the desert of the Syrian,.
Saudi Arabia desert on the right or in the east..
Now let me show you a map here..
As you can see, this is the Google map..
And you can see Israel is right here..
And this long narrow corridor is the land bridge.

$^{201}$of all the continents, Africa, Asia minor and Europe,.
and Asia..
This is the land bridge..
And on the left is the huge Eastern Mediterranean Sea..
And then on the right is the Saudi Arabia or Syrian desert..
So the birds have to travel through from north to south.
or south to north, they have to fly through.
the land of Israel..
In fact, in the land of Israel,.
just north of the lake of Galilee or Sea of Galilee,.
here there's a small strip of flat land called Hula Valley..
This Hula Valley is famous for bird watching..
In the end of the year, around the month of November.
or December, bird watching people would gather.
in this Hula Valley north of Sea of Galilee.
in order to just look at the birds,.
look at all this beautiful migratory birds.
flying through Israel..
As you can see this beautiful pictures of birds.
and migratory birds are flying through the land of Israel,.
particularly through that corridor of Galilee..
Migratory bird flying through Mount Hermon right here..
So Jesus was saying, pointing to the birds.
and the disciples knew about this..
It's about time of the year.
when the migratory birds will fly through.
this strip of land and they can see all sorts.
of different birds and different big and small.
flying creatures going through the land..
So they are not unfamiliar with all of this.
beautiful birds and different migratory pattern..
Jesus was pointing to something that they could see.
and they could observe and just ask them to open your eyes.
and see how the creatures of God, how the animals,.
how the birds are being taken care of by the heavenly father..
So compared to our own lives, we are so much more valuable.
than one of these sparrow or birds..
And then he take another example from the field..
So the first example was to direct our eyes,.
direct our perspective toward the sky,.

$^{241}$see the migratory birds, right?.
They are being taken care of by the heavenly father,.
just like God will take care of you..
And also he direct the perspective on the field..
On the land..
So we look up and then we look down.
and we see all kinds of wild flower,.
the lilies of the field..
The wild flower growing in every single season,.
every single month, they have different colour.
in every month..
Sometimes it's red, sometimes it's white,.
sometimes it's yellow, sometimes it's like.
this beautiful purple..
So he was saying to these disciples, saying to us,.
look at the wild flower growing almost by themselves.
because the heavenly father is taking care of them..
Look at the flower, how glorious.
and how beautiful they are..
He said, compare the beauty of the wild flower in gallery..
And you will notice how beautiful they are.
compared to the time of Solomon..
Remember Jesus said,.
you see the wild flower, the lilies of the field,.
how glorious they are..
They are more beautiful, more glorious,.
more majestic than the time of Solomon..
In times of Solomon,.
the royal purple or the royal building.
were the glory..
Now, if you know about something,.
if you know something about purple dye,.
the colour purple is really hard to obtain,.
even these days..
Purple is the hardest dye,.
the hardest dye to obtain..
Now, back in Jesus' day,.
when people need to have the purple dye,.
they have to go to the Mediterranean seacoast.
to gather the myrrh shells..

$^{281}$And one of the, take the shells and open it up.
and you harvest the gland,.
a very small part of the myrrh..
And then you gather all the glands.
and you can get some purple dye..
Now, as you can see in this picture,.
you have to do that one by one.
in order to have purple,.
in order to have the most valuable,.
most royal colour in Jesus' day..
How many do we need in order to dye purple?.
How many do we need in order to dye clothing into purple?.
Now, as you can see in this formula,.
you need 54 kilogram of myrrh shell.
that is equal to 119 pounds.
of different, many myrrh shells.
in order to produce only one gram of purple..
One gram of the colour..
It's so hard to get, right?.
It's so valuable..
But look at the lilies of the field..
Look at the pictures..
Look at pictures how the purple flowers.
are just growing on the field..
And God is taking care of them..
God is clothing them with the best colour on earth..
You cannot compare the beauty of this wildflower,.
this purple flower, red flower,.
with all these different kinds of beautiful colours..
Compare them to the time of King Solomon,.
what he was wearing, right?.
These flowers is more beautiful.
than what Solomon has on his body..
So don't you worry..
Don't you overwork so hard.
in order to obtain one of these colours, right?.
One of these glories..
Because you know that the heavenly father.
is taking care of you..
Now, not only the purple colour.

$^{321}$was such a valuable in Jesus' day,.
even the buildings, the royal buildings, the temple,.
they were the most magnificent buildings.
on earth in those days..
Here I show you a picture of a model of Jerusalem..
The ratio is one to 50..
You can see that in Israel Museum..
As you can see this man here,.
you know, this model of Jerusalem is already,.
is very, very, very big..
And the scale is only one to 50..
Now, as you can see this model,.
it is gigantic temple mount..
It is a glorious golden building of the temple..
Now, when people think of Solomon,.
they think of the most glorious period of Israel's history..
I remember I had a dinner with some of the classmates.
when I was studying in Jerusalem..
And we had the opportunity to sit down.
with some of the most,.
the best Israeli Jewish teachers in Jerusalem..
So I remember one time we were sitting down around the table.
with one of the teachers that I really respect..
He is Israeli Jewish teacher..
And we were, you know, chit-chatting.
and talking about different things..
And then one of the classmates asked the teacher,.
the Jewish teacher,.
and if you had the opportunity,.
if you can travel back to the Old Testament time.
or the history of Israel,.
which period would you choose in order to witness.
the glory or what happened in those days?.
And without even thinking,.
he respond by saying,.
I would want to travel back to the time of King Solomon.
when he dedicate the temple,.
when the glory of God fills the temple.
and people cannot even enter..
You know, that sort of glorious day,.

$^{361}$the good old days,.
was the time of King Solomon..
But what was Jesus saying?.
Jesus was saying,.
even when you consider the most glorious,.
magnificent building,.
magnificent period of your history,.
this cannot be compared to the lilies of the field..
How the father has clothed the field.
with all sorts of different valuable colors..
You cannot compare them to one of these..
One of these purple flowers or different kinds..
They grow by themselves.
because the heavenly father is taking care of them..
Now, the take home for us, as Jesus said,.
is that do not be overly anxious about.
what we eat and what we drink..
Or housing, right?.
What kind of house do we have?.
What kind of house do we buy?.
What kind of car do we drive?.
Again, it's not suggesting.
in any way that we should live irresponsibly..
We should be good stewards.
of all the resources God has given to us..
But Jesus is saying, do not be overly anxious..
Do not be just overly worrying..
Or work just so hard on where you live,.
what you wear, and how you look..
God is taking care of you..
But seek first the kingdom of God and his righteousness..
And all these things, what we eat and what we drink,.
where we stay, will be added to you..
Seek first the kingdom of God..
I think a better translation would be.
pursue the kingdom of God..
Run after first the kingdom of God..
Be concerned about the name of God..
Making God's name famous,.
it should be our first and foremost concern..

$^{401}$And his righteousness, all the things.
that the word of God is teaching us to do..
Be concerned and pursue above all.
the righteousness of God, what God requires us to do..
And then all these things will be provided for you..
The priority of our lives,.
as we say goodbye to the old time, 2022,.
and welcome a new year, 2023,.
let us be reminded, seek first, run after,.
be the first and foremost concern.
in our desire to glorify God, to serve him..
So talking about New Year's resolutions,.
how do you fill in the blanks?.
How do you fill in the blanks?.
What are those things that we should first seek after.
and desire most?.
Let me suggest to you five of them,.
before considering any other things..
First, live in the word of God..
Read the word of God, meditate on it,.
reflect upon it, so that the word of God.
can fill our lives and inform our living..
That is the first and foremost thing we can do.
and we have to do as disciples of Christ..
When the word of God is filling your lives,.
your priority is correct, and your priority of lives,.
you know what is your foremost concern and desire..
Then everything else will be configured correctly..
Live in the word of God, read the word of God daily,.
every single day, meditate upon it..
I'm not saying you have to read the word of God 24/7..
I'm saying that should be our consistent pattern of thought..
Think about what the Bible says,.
how we should make any decision,.
how we can please God..
These are all coming from scripture..
So we have to live in the word..
And then secondly, I suggest that we have to pray daily..
Pray for ourselves and most formally,.
and pray for others..

$^{441}$Not just be concerned about our own needs.
and most of the time our own wants..
You know, those are really insignificant.
compared to the eternal kingdom of God..
Pray daily for other people,.
those who are in need around us..
Pray for them and care for them..
And third, worship wholeheartedly..
Worship God with all of our heart.
and all of our soul and all of our strength..
You know, we worship God because whenever we worship,.
we are reminded who is really the king,.
who is really the leader of our lives,.
not ourselves, God himself..
And we can reset our priority on a weekly basis..
That's why we insist on church not closing, right?.
Not be so concerned about, you know,.
what kind of pandemic out there and disease out there..
A church should be operating every single week..
We worship God every single week, right?.
Wholeheartedly..
Fourth, build up the family of a church,.
the family of God..
The New Testament and the Old Testament.
regularly, consistently reminds us.
our number one concern to live socially.
is to build up the family of God..
Use our gifts, use our talents to build up the church,.
to plan a church, to strengthen our church..
Go to a fellowship and pray with other believers..
Let other believers encourage you.
and let your presence be an encouragement to other people..
Build up the family of the church..
Fifth, practice stewardship, faithful stewardship..
The Bible tells us use our time faithfully to serve God..
Use our talent to serve God..
Use our treasure, our money, our offering.
to build up the ministry for the glory of God..
So seek first the kingdom of God and his righteousness..
Let this be our priority..

$^{481}$Let us reset our list of things that we wanna do,.
our list of resolutions so that we can glorify God together..
Let us all pray together..
Father, we wanna give you thanks for a time of renewal..
As we say goodbye to year 2023 and welcome a new year,.
let this be a reminder to all of us to reset our priority,.
to seek first your kingdom and your righteousness..
And let us not be too concerned and overly anxious.
about the essential things,.
because we know that you will take care of us.
and help us to serve you, to worship you wholeheartedly..
Help us to live faithfully according to your word..
In Christ, we pray..
[BLANK AUDIO].
\newpage



\section{馬可福音 4:35-41}
\label{sec:87VD_Eud_tU}
\textbf{漫步風雨中 (馬可福音4\_35-41) - 張文龍牧師}
\newline
\newline
連結: \href{https://youtube.com/watch?v=87VD-Eud_tU}{\texttt{https://youtube.com/watch?v=87VD-Eud\_tU}} ~~~~ 語音日期: 2024-03-29
\newline
\newline
\hyperref[sec:_seeVdxX33U]{\small{< < < PREV SERMON < < <}}
~
\hyperref[sec:index_chronic]{\small{[返順時目]}}
~
\hyperref[sec:index_scriptual]{\small{[返順卷目]}}
~
\hyperref[sec:9_6WJzTToTc]{\small{> > > NEXT SERMON > > >}}
\newline
\newline
馬可福音 4:35-41
\newline
\begin{longtable}{cl}
\hline
\hline
章節 & 經文 (和合本修訂版)\\
\hline
4:35 & \begin{tabularx}{0.7\textwidth}{X} 那天晚上，耶穌對門徒說：「我們渡到對岸去吧。」 \end{tabularx} \\ \\ \relax
4:36 & \begin{tabularx}{0.7\textwidth}{X} 門徒離開眾人，耶穌已在船上，他們就請他一同去；也有別的船和他同行。 \end{tabularx} \\ \\ \relax
4:37 & \begin{tabularx}{0.7\textwidth}{X} 忽然狂風大作，波浪打入船內，以致船灌滿了水。 \end{tabularx} \\ \\ \relax
4:38 & \begin{tabularx}{0.7\textwidth}{X} 耶穌在船尾上，枕著枕頭睡覺。門徒叫醒他，說：「老師！我們快沒命了，你不管嗎？」 \end{tabularx} \\ \\ \relax
4:39 & \begin{tabularx}{0.7\textwidth}{X} 耶穌醒了，斥責那風，向海說：「住了吧！靜了吧！」風就止住，大大平靜了。 \end{tabularx} \\ \\ \relax
4:40 & \begin{tabularx}{0.7\textwidth}{X} 耶穌對他們說：「為甚麼膽怯？你們還沒有信心嗎？」 \end{tabularx} \\ \\ \relax
4:41 & \begin{tabularx}{0.7\textwidth}{X} 他們就非常懼怕，彼此說：「這到底是誰？連風和海都聽從他。」 \end{tabularx} \\ \\
[1ex]
\hline
\hline
\end{longtable}
$^{1}$最後Mark Sir會帶我們一起搜聖餐.
我現在將時間交給Mark Sir.
大家好.
今早如果是起床很難的.
應該最難的是我.
因為昨晚11點晚上才回來.
這兩個星期我們去了加州探親等等.
下個星期再給大家多一點分享.
我總結三個字就是死不去.
無論如何今天是Pastor William和我們正道.
很有趣認識到William.
因為你知道小弟在神童院進修中.
有一次在下課的時候.
出去拍攝拍攝.
碰見William.
大家聊天的時候才知道原來大家都是香港人.
廣東人.
我讀舊約的.
他讀新約的.
所以很有趣.
我們大家一聊天就聊到別人打樣.
我們都還沒走.
所以William來到芝加哥的時候.
一定要請他來分享廣東.
他太太和他的子女都在我們中間.
我們歡迎他們.
我們把時間交給William.
各位誠惠的弟兄姊妹你們好.
多謝陳牧師多謝Mark Sir的邀請.
那天在神學院碰到他.
是我第一次面對面見到他.
我之前不認識他.
是我教他的弟兄姊妹.
你有沒有看過Mark Sir的YouTube.
香港的弟兄姊妹都和我這樣說.
在加州的弟兄姊妹都和我這樣說.
我見到他.
我離遠見到他.
我就很開心.
Mark Sir Mark Sir.

$^{41}$我差點叫他幫我簽個名.
可以嗎.
多謝多謝陳牧師.
我今天想和大家去.
透過馬可福音的一段經文.
大家都熟悉的.
其實是耶穌平靜風浪.
和大家去看一個人生的題目.
大家有沒有聽過.
一個古詩人叫做河馬.
不是那隻動物.
是一個在公元前八百年的.
一個古希臘雙眼失明的詩人.
叫做河馬.
相傳他寫了兩部很長的著作.
我們叫做史詩.
是用詩歌題材來寫一些人物歷史的.
一些古典的著作.
其中一部的著作.
大家認識的名字叫Odyssey.
但不是那輛車.
你知道牌子的車的model叫Odyssey.
不是那輛車.
兩卷著作我們統稱叫河馬史詩.
這兩卷的河馬史詩的著作.
其實是成為了日後古希臘的文獻.
一些文學著作的基石.
其實到今天有很多不同的文學作品.
一些音樂.
藝術甚至乎是一些電影的題目.
都是受到河馬史詩的影響.
我為什麼要提河馬史詩.
因為很多學者對河馬史詩的研究的分析.
他發現到這部著作最吸引的地方在哪裡呢.
就是貫穿著這兩卷著作的一股.
我們叫做秩序.
有一股秩序或者我們叫做一個對稱的意識.
而有趣的地方就是這股秩序並非不可動搖的.
是經過一些時間變遷之後.
這股秩序會失序的.

$^{81}$會失去這個秩序.
但很奇怪的就是這個失序之後.
它會慢慢又會安然的來到留下.
來到繼續下去.
當我們這樣看的時候.
其實和我們人生有什麼關係呢.
我們人生的歷程.
其實就是不斷的在一個秩序和失序的意識形態底下.
一路繼續走下去.
我們時常都覺得.
我們時常都想.
每一個人都很想能夠找到一個安身立命的地方.
但很可惜的就是.
我們想有秩序地生活.
但很可惜的就是.
不時受到一些很無情的混亂的威脅.
使得我們想秩序生活.
但變得失序.
就正如我們都聽過.
變幻才是永恆.
究竟怎樣才能夠可以很安然的繼續生活下去呢.
我們不需要.
其實對人生的描寫.
都不需要只有古人才會讀.
大家有沒有聽過這首流行曲.
當然流行就是以前.
叫《幾許風雨》.
你有沒有聽過.
你聽過大概都知道.
我都猜到你幾多歲.
是二十世紀八十年代.
在香港一首很流行.
記不記得是誰寫的.
是誰唱的.
羅文唱的.
大家懂不懂得唱.
你不用現在唱.
一首很潰滯人口的流行歌.
叫《幾許風雨》.
這首歌一開始的時候就說.

$^{121}$無言輕衣.
唱遍凝望與世急也亂.
這句說話已經告訴我們.
人生是有很多痛苦的.
他用風雨來描寫痛苦.
或者人生裡面很多的混亂.
但這首詩歌去到下面.
他說《幾許風雨》.
我也經過屹立到木前.
意思就是縱然有很多混亂.
很多痛苦.
很多挫敗.
但是經歷了之後回望.
我們都可以站得穩.
這首流行歌就是想說這件事.
究竟怎樣能夠可以站得穩.
究竟我們怎樣能夠從秩序.
變成失序.
然後沉回秩序.
很多很多的智者都去嘗試回答這個問題.
我今天想和大家看這段經文.
就是馬學福音.
耶穌平靜風浪.
同樣用風浪來描寫.
好像一個人生的旅程.
究竟耶穌和門徒.
兩班人兩組人的對比.
給我們一個甚麼的功課.
可以令我們在風雨中都能夠站得穩.
我們在風雨中都好像可以漫步.
我猜大家應該比較熟悉這段經文.
這段經文的開始的時候.
就在馬學福音第四章第35節.
當天黃昏耶穌對門徒說.
我們到海那邊去吧.
門徒離開群眾.
耶穌已經在船上.
他們就載他過去.
也有別的船和他同去.
第37節.

$^{161}$突然間起了狂風.
波浪不斷地打進船內.
船艙裡面都積滿了水.
故事一開始的時候就是晚上.
晚上耶穌出了一個主意.
不是門徒.
是耶穌出主意.
不如我們去加利利海或加利利湖吧.
這個不是門徒自己提議.
不如晚上去賞月.
在船上賞月好看一點.
不是.
是耶穌叫的.
但是出去到海或湖的時候.
突然間起了風暴.
為甚麼會突然間起風暴.
我想你問Marshall會更加清楚.
更加認識當地在加利利湖的地理環境.
可能Marshall會解釋得更好.
簡單來說就是加利利湖是在海平面以下的七八尺.
但是加利利湖是被很多高山圍住.
特別是在東面一個叫做黑門山.
高達九千二百尺.
你看看七八尺水平面以下.
九千二百尺海平面以上.
這個這麼大的差距就造成了在山上的凍空氣和在湖面的冷空氣.
暖空氣上升的時候就產生一個很對價的風暴.
是突然間會發生的.
如果你去過以色列.
跟Marshall去以色列聖地行的時候.
你都可能會經歷類似的情況.
在這樣的情況下.
門徒和耶穌經歷了的是一個突然間來的風暴.
是沒有預計到的.
但總之這個風暴出現了.
要留意的就是在風暴出現的時候.
耶穌不是站在岸上看著你們門徒有什麼反應.
是耶穌在船裡面跟門徒一起經歷這個風暴.
我們要問的為什麼耶穌要帶門徒出去經歷風暴呢.
其實風暴對人有什麼意義呢.

$^{201}$風暴就是叫人看到自己人生的限制.
人是很脆弱的.
風暴叫人失控.
好像失序一樣.
風暴叫我們無力抵抗.
在現代的社會.
在現代的科技的社會裡面.
有一個迷思就是我們能夠掌握一切.
我們有很高的醫學技術.
我們可以上太空.
我們會被一個科技的迷思來騙了我們.
以為人可以掌握一切.
以為人好像上帝一樣可以創造秩序.
可以穩定的活下去.
但風暴其實是壓倒我們人的能力的.
這個就是耶穌想門徒跟他一起經歷的人生的一課.
在這個風暴的發生的底下.
耶穌跟門徒的反應是截然不同的.
大家知道門徒的反應.
頌賢裡面有很多個漁夫.
他們打魚為生.
但他們竟然驚惶失措.
就好像經文那裡在第39節說.
老師我們要死了你都不理嗎.
其實帶有一點責備的.
其實很自然很明白.
當然是責備.
我們明明在岸上好好的.
你為什麼帶我們出海.
你帶我們出海要我們經歷風暴.
我想我們每一個人都會將責任來賴耶穌.
但耶穌呢.
耶穌在船尾那裡睡覺.
或者還可能拿著他媽媽送給他的毛公仔在那裡睡覺.
很安然的在睡覺.
耶穌不是累到連水進來船裡面都不知道.
我想沒有可能的事.
耶穌只是沒有打算起身.
耶穌沒有打算要平靜這個風浪.
在第39節那裡說耶穌起來.

$^{241}$其實是在說耶穌被叫醒.
耶穌沒有想起身.
耶穌起身的原因是因為門徒驚惶失措.
所以耶穌起身為了門徒平靜風浪.
讓他們可以安心.
真正的風暴.
當我們見到耶穌的反應.
和門徒那種亂章反應的時候.
我們知道真正的風暴其實不是在外在的因素裡面.
今天我們見到的風暴不是因為政治的不穩.
不是因為經濟走下坡.
不是因為社會上有什麼特別的變遷.
真正的風暴正如我們見到耶穌和門徒的分別.
其實在人的心裡面.
耶穌在經文的第40節.
耶穌說為什麼這樣膽怯.
為什麼沒有信心.
膽怯信心是一個內在的質素.
耶穌對門徒的責備也好.
或者提醒也好.
為什麼你們這麼擔心.
為什麼你們裡面這麼不穩定.
當然什麼叫信心.
我們需要慢慢來認識.
這裡的信心不是說保羅恩信稱義的意思.
這裡所說的信心.
耶穌說不要膽怯.
是並非我們相信.
信心並非相信上帝一定會除去風暴.
這個不是信心.
這個不是耶穌想說的信心.
信心是好像耶穌一樣.
能夠在狂風大雨底下繼續睡覺.
不是除去外在的威脅.
而是內在很穩定繼續可以睡覺.
不被風暴壓倒.
不會變得失序.
能夠漫步風雨中.
這個就是信心.
耶穌讓我們看到.

$^{281}$我們會問一個問題.
我們怎樣才能得到這個信心呢.
這段經文去到後面.
由第37節一直去到第41節.
我們留意到作者描寫的時候.
有一個漸進的描述.
第37節他說那時候起了大風.
然後去到第39節.
當耶穌斥責風的時候.
然後風就大大平靜了.
然後去到最後第41節.
經文是這樣翻譯的.
他說門徒大大大大大地懼怕.
要留意這個漸進的描寫.
風暴平靜了之後.
門徒的反應竟然更加戲劇化.
門徒的反應就是當他見到耶穌平靜了風暴的時候.
他對耶穌所產生那份我們叫做懼怕.
應該反應得更好.
就是剛才詩歌所唱的敬畏.
門徒見到耶穌平靜了風暴.
門徒對耶穌所產生的那份敬畏.
比起剛才見到福天蓋地的風暴.
更加提升了一層.
而這種敬畏.
在門徒心裡面所產生對耶穌的敬畏.
和當時突然間水面平靜了.
形成一個很強烈的對比.
當我們能夠被耶穌壓倒了我們的時候.
我們就不會被任何東西壓倒.
當我們能夠經歷一個比風暴更大的風暴.
我們就能夠得到這份所謂耶穌所講的.
不要膽怯 要有信心.
可以安然入睡.
所以問題是誰才是最可畏的呢?.
誰才是我們要敬畏的對象呢?.
是經濟的崩潰?.
是政治的不穩?.
還是敬畏的是耶穌基督?.
面對其他一切的突如其來的風暴.

$^{321}$我們都可以不需要害怕呢?.
在上一年.
Colin這本英文字典.
選了一個字來總結2022年.
這個字叫做Permanent Crisis.
是一個造出來的字.
Permanent Crisis.
就是一個永久的危機.
總結了2022年的整個世界的狀況.
其實這個字Colin選出來.
是真的要來描寫.
特別是在英國.
因為Colin是英國出版的字典.
特別是在英國和世界各地.
2022年面對的.
或者將來要面對的一個危機.
包括什麼呢?.
剛才講了政治的危機.
俄烏戰爭所帶來全世界.
無論是經濟還是政治上的危機.
等等的危機.
我們真的生活在一個很多轉變化.
很多突如其來的變化的世界裡面.
在這個突如其來的變化.
在這個危機的底下.
究竟怎樣能夠站得穩呢?.
倫敦大學有一個我們叫做組織行為科學.
叫Organizational Behavior的一個老師.
我都不懂得讀了.
應該是法文來的.
他說原來在危機的底下.
我們是怎樣能夠尋回秩序呢?.
原來是靠的不是一句說話.
沒事發生 沒事發生.
大家安好就好.
他說這樣是不行的.
所以我們見到弟兄姊妹有困難的時候.
你不要跟他說上帝與你同在.
這樣就搞定了.
一兩句是不行的.

$^{361}$他說原來能夠尋回秩序.
能夠站得穩在危機底下.
最大的動力 最大的支撐點.
是有人跟你一起.
我們的信仰是賦予給我們的是.
道成肉身的耶穌基督.
是與我們同行 同住那位上帝.
教會群體弟兄姊妹大家的相交.
就能夠給大家有力量去面對一個.
突如其來你經歷的人生風暴.
我們不要將我們的專注力.
放在無論是科學 人的理性上.
想著透過這些就可以得到.
我們能夠安身立命的方向.
這個不行的.
基督信仰 耶穌的同在.
弟兄姊妹信徒群體的相交.
就能夠讓我們在一個風雨當中.
能夠站得穩.
弟兄姊妹 接下來2023年.
我們不知道究竟會再有多少風雨.
多少突發的事件會發生.
但我們肯定知道一件事.
你身邊有耶穌基督與你們同在.
願上帝常與你們同在.
謝謝.
我們低頭祈禱.
主啊 多謝你一個這麼切身的訊息.
去提醒我們在人生的風暴當中.
有你自己在船上.
主啊 求你提醒我們.
見到外面很大風暴的時候.
我們敬畏的是你自己.
我們倚靠的是你自己.
有你同在的時候.
我們就知道你自己是那位.
在我們生命中能夠重尋方向.
重新穩定的一位主.
主啊 叫我們今天來順靠去敬拜你.
讓牧師的提醒.

$^{401}$一個經文的訊息.
能夠成為我們今天的力量.
在這個星期當中.
無論我們弟兄姊妹遇見任何的難處.
有任何大的問題.
我們知道我們需要來敬畏去懼怕的.
是你自己.
主啊 為我們中間.
我們的弟兄姊妹.
我們的教會.
我們知道有很多的肢體.
都是遇見很多不同的問題.
無論是身體上的不適.
或者是生命的方向裡面.
有很多的未知之數.
主啊 我們知道你自己在船上.
在我們人生當中.
你是道成的新的主.
與我們一起同行 同在.
我們這樣祈禱 奉主名.
\newpage



\section{}
\label{sec:9_6WJzTToTc}
\textbf{聖地網上行 - 步主腳蹤 [第十集] 潔淨聖殿}
\newline
\newline
連結: \href{https://youtube.com/watch?v=9-6WJzTToTc}{\texttt{https://youtube.com/watch?v=9-6WJzTToTc}} ~~~~ 語音日期: 2024-03-31
\newline
\newline
\hyperref[sec:87VD_Eud_tU]{\small{< < < PREV SERMON < < <}}
~
\hyperref[sec:index_chronic]{\small{[返順時目]}}
~
\hyperref[sec:index_scriptual]{\small{[返順卷目]}}
~
\hyperref[sec:AndEm9fTjLE]{\small{> > > NEXT SERMON > > >}}
\newline
\newline
$^{1}$上一集我們講到在鑒嵐山.
耶穌騎上路.
祂就是所謂的榮晉耶路撒冷.
現在我們看到的位置大概就是這樣.
我們站在鑒嵐山的山頂.
前面就是看到耶路撒冷.
現在看到的當然是耶路撒冷的舊城.
眼前所見的.
這裡就是全世界.
政治,民族或者宗教的.
是全世界政治,民族或者宗教上最敏感的地方.
我們今天看到的耶路撒冷.
最矚目的地方當然就是這個金頂清真寺.
或者盆石清真寺.
Doom of the Rock.
事實上這個不是一個墓地.
不是清真寺.
因為清真寺是以色列教徒去聚集.
去祈禱的地方.
實際上很少會在這裡聚集.
事實上這裡是一個宗教的座標.
一個所謂的神社.
去紀念以色列教第三個大的聖城.
耶路撒冷就是第三個聖城.
除了在沙特阿拉伯的麥加和麥丁娜.
在這個金頂下就是耶路撒冷東山最高的一個盆石的地方.
也是猶太人的第一第二聖殿.
知聖所在的地方.
先講一下歷史背景.
大概公元後70年的時候.
大家知道就是耶路撒冷的聖殿.
猶太人的聖殿被毀.
當時有很多猶太人被驅逐.
他們離開了耶路撒冷.
所以這個地方多年來是被荒廢的.
猶太人只准每年一次回來這個地方.
在他們的聖殿被毀紀念日.
Tishab Av.
大約是我們西歷7月8月左右的時間.
來到這個地方他們在西面的牆上祈禱.

$^{41}$因為那裡是知聖所最近的.
歷史上的猶太人來到西牆上祈禱的時候.
因為他們紀念到聖殿被毀的災難.
所以他們就在那裡痛哭.
所以叫做哭牆.
但是現在很少人叫做哭牆.
因為特別到1967年之後.
以色列軍重奪耶路撒冷.
猶太人自由地來到西牆祈禱.
所以較多人叫做西牆.
或者希伯來文叫做Catel.
無論如何在一世紀之後.
被毀之後荒廢了好幾百年.
直到638年的時候.
伊斯蘭軍隊征服了中東和耶路撒冷.
他們就在這個地方建立了一個金頂的盆石清真寺.
就是我們現在看到的這個.
之後在715年的時候.
在南邊建造了另一個大的清真寺.
叫做阿克薩.
Al-Aqsa的意思是遙遠.
在伊斯蘭的傳統裡面.
穆罕默德在這裡夜間升天.
另外你要留意的就是一個很大的石平台.
這個叫做聖殿山.
Temple Mount.
或者伊斯蘭教徒會叫做Al-Aqsa.
西面是488米長.
東面也是有400多米長.
南北面分別有200-300米長.
總共有144,000平方米.
等於20個足球場.
或者27個美式足球場.
這是一個很大的平台.
多德在2000年之前是大希律王.
他執政的時候.
作為羅馬帝國授權的猶太人王.
他大興土木.
然後擴建聖殿.
在這個平台下面加造了很多羅馬人很出名的拱門.

$^{81}$去頂著這個平台.
這個平台是一個很大的平台.
容納了很多人.
聖殿的歷史是怎樣的呢?.
當時第一世紀耶穌的時代.
人對聖殿是讚口不絕的.
其實在羅馬帝國裡面.
是舉世聞名的.
猶太人甚至說.
你沒有見過這個聖殿.
你是從來沒有見過美妙的建築.
這裡被形容為人間裡面最偉大的工程.
事實上在現代之前.
這個真的是數一數二.
人造的一個大的建築物.
其實猶太教跟其他外邦的宗教有很大分別.
就是猶太教只有一個聖所.
所以很多資源都集中在這個聖殿裡面.
甚至猶太人會認為這是萬夫莫敵.
就是永遠不會滅亡的.
神一定會捍衛的一個地方.
因為上帝的名字就在這裡.
到了耶穌的時候.
聖殿已經有幾十年的歷史了.
難怪我們看福音書的時候.
門徒從加利利來到地老山首節.
他們就指著石頭讚不絕口.
好幾年前 UCLA 有一個團隊做了一個電腦的模型.
我認為這個電腦模型是非常準確的.
是由南面上去耶路撒冷.
看到的景象.
這個立體動畫是十年前製作的.
所以解像度是比較低一點.
我們已經盡量提升了.
你可以想像一個朝聖者從南面大衛城.
一直攀上去耶路撒冷的聖殿.
我們看到前面這個聖殿山.
是一個多麼的震撼和偉大的建築.
前面我們看到一個很大的樓梯.
就是上到聖殿山的上面.

$^{121}$我們前面下面的石板街.
兩邊都是商舖.
可以兌換銀錢和買祭物.
接著樓梯上去就真的上到聖殿山.
現在這個地方叫做外邦人園.
外邦人只能夠去到這個地方.
不可以再繼續親近神.
再轉過來南面這個地方.
聖殿山很有名的.
叫做 Royal Stuart.
是一個有瓦遮頭的地方.
很多的雲石的石柱頂著.
整個聖殿山的平台都是露天的會堂.
有不少猶太人的拉比會早早的去祈禱.
去讚美,敬拜.
甚至去聽道.
有很多福音書的情節都是在這種環境發生.
包括耶穌經常去聖殿教導人.
現在我們上到聖殿的西南角.
這裡有一個祭司.
有一個地方刻著說這裡是摧國的地方.
去通知城裡面的人.
聽到報國聲的時候就知道.
Shabbat 到了.
看完剛才那麼偉大的聖殿的時候.
每一個朝聖者從橄欖山一直下山.
要進入耶路撒冷的時候.
那種興奮和期待.
但是耶穌在橄欖山騎驢下山.
進入聖殿面對聖殿的時候.
祂就說了一些很嚇人的話.
路加方第十九章四十一節.
耶穌快到耶路撒冷看見城.
就為他哀哭說.
巴不得你在這日子知道關係你平安的事.
無奈這是現在是隱藏的.
叫你的眼看不出來.
因為日子將到你的仇敵必捉起土裡.
周圍環繞你四面要困住你.
並且要去掃滅你和你裡頭的兒女.

$^{161}$連一塊石頭也不留在石頭上.
因你不知道眷顧你的時候.
路加記載這一段.
主耶穌騎著驢.
在橄欖山的半山面對耶路撒冷的時候.
說出了一個哀嘆的訊息.
聖殿將會在一代之內被審判.
預言羅馬軍隊在主后七十年將聖殿毀滅.
這裡有一所教堂叫做Dominus Flavit.
中文翻譯叫做主哭耶經堂.
就是主入堂.
教堂的外形你覺得像什麼呢?.
就像眼淚的形狀.
一滴的眼淚.
當我們進入城的時候.
福音書就這樣記載了.
馬科第十一章.
他們來到耶路撒冷.
耶穌進了聖殿.
就趕出殿裡面做買賣的人.
推倒兌換銀錢的人的桌子.
和賣甲子之人的凳子.
也不許人拿器具.
送電鋁經過.
便教訓他們說.
經上不事記著.
我的電必稱為萬國禱告的電.
你們都使它成為賊窩了.
祭司長和民事聽見者說.
想法是想去除滅耶穌.
但又怕他.
因為太多人跟隨他.
他有民眾的力量.
因為民眾都稀奇他的教訓.
每天晚上耶穌出城去.
這個所謂耶穌潔淨聖殿的故事.
四福音都有記載.
聖殿裡面有買賣.
那你就問.
有什麼大不了.

$^{201}$然後我們可以再問.
到底出了什麼問題.
然後耶穌做的事情有什麼意思.
到底買賣有什麼大不了.
其實沒有什麼大不了.
其實有實際的需要.
因為你想想.
如果我是外地來的朝聖者.
在加利利或者更遠地方.
來到野口山朝聖的時候.
我起碼要走幾天的路程.
至少.
如果我帶著獻祭的羊,牛等等.
如果中途很容易會有損傷.
來到野口山我就不能夠獻祭.
所以呢.
律法書也有一些實際的安排.
申命記第十四章.
當然會說你的神賜福給你們的時候.
你進到迦南地.
耶和華神就會選擇立他名的地方.
如果離你們太遠的話.
那個路太長.
就使你不能夠將祭物帶到那裡.
第二十五節.
他說你就可以換成銀紙.
將銀紙包起來.
拿在手上.
然後做去耶和華神所選擇的地方.
耶路撒冷聖殿.
你就用你所帶來的銀紙.
隨心所欲地買祭物牛羊.
或者買清酒,龍酒.
隨心所願.
然後你跟你的家屬在耶和華面前.
就可以在愛快樂的境拜.
這裡講到遠地來到朝聖的人.
可以帶著銀紙.
或者當時古代所謂的銀錠.
Silver Ingots.

$^{241}$來到耶路撒冷.
就將這些銀紙換成祭物.
他們能夠去境拜.
這些都是銀錠.
在不同的位置可以找到的.
到了巴比倫時期之後.
中東就開始通用一些錢幣了.
在波斯時期.
這些錢幣叫做猶大幣.
Yehu Coin.
你看到這些小錢幣上面刻了亞蘭文.
代表猶大省.
亞蘭文叫做Yehu.
Y-H-D.
幣上面有時候是一些動物.
或者有時候是一些花.
這裡所看到的就是百合花.
到了一世紀的時候.
當時耶路撒冷的聖殿當局規定.
要用當時通行的所謂推羅銀幣.
Tyrant Shackle.
推羅就是以色列北面.
推羅西頓等等這些大的海港.
當時是紅霸著地中海.
所謂海運貿易的大海港.
就是今天黎巴嫩一帶的地方.
所以當時推羅的銀幣.
純度是最高的.
也是風行世界的.
所以兌換銀錢是一些很實際的需要.
其實律法書也是這樣允許的.
因為你想想.
如果朝聖者從不同的國家.
不同地方來到的時候.
他就拿著不同的錢幣.
不同的貨幣.
你這樣奉獻是不行的.
相等於今天.
如果你在教會裡面.
有一些外地來的人.

$^{281}$他又用台幣,人民幣,美金,加幣.
又拿他的禮物卡.
丟到奉獻箱裡面.
那麼那個奉獻就很難去處理.
或者家裡有些爛銅爛鐵又拿來.
所以是處理獻金上.
一種很實際的安排.
所以猶太人的典籍.
頭五也有提過.
在聖殿山上有兌換銀錢的桌子.
這裡說的是亞達月十五日.
兌換銀錢的會在各省.
來擺設這些桌子.
在聖殿山擺設這些兌換銀錢的桌子.
亞達月就是農曆年的年尾.
下一個月就是彌撒月.
就是西曆的三四月.
就是 Passover 預節的時候.
聖殿有專人來提供這些兌換.
以及買賣祭物的服務.
也有猶太人的律例.
來防止濫用或者防止高利貸的出現.
兌換銀錢有需要.
祭物也有需要.
所以猶太人的法典也有提過.
在聖殿山南邊有一個地方叫做.
Hanunyot 或者 Hanut.
就是商鋪的意思.
Shops 的意思.
大概就是剛才我們經過南邊這個地方.
出入口這個地方.
或者在聖殿山上面.
愛邦人院那裡.
也有一些商鋪.
另外在南邊就找到不同的.
關於憲制聖殿敬拜的文物.
其中一塊就是一個很小的石碼.
上面刻了希伯來文.
Koban.
就是祭物的意思.

$^{321}$上面也有刻了一個圖案.
兩隻反轉的甲子.
似乎這個石碼是屬於其中一間買賣祭物的商鋪.
所以兌換銀錢和祭物.
有買賣是很正常的事情.
另外如果你看看舊約的經文.
你就知道聖殿不只是一個宗教的中心.
也像銀行一樣.
聶王記下十六章第八節.
這裡說到阿哈斯王將耶和華殿裡面的皇宮庫房.
庫房裡面所有的金銀.
拿去給亞述王為禮物.
因為他們攻打你.
李遷米記十三章.
就是他們猶太人被擄.
然後再回歸回來的時候.
這裡說到猶太眾人.
將五穀新酒,油等十分之一的奉獻放入庫房.
兩個中間的瑪加比二書也提到.
大祭司聲明說庫房裡面的存款是給孤兒寡婦的.
有多少我們不知道.
但是按照一個學者的估計.
瑪加比二書看到的.
相等於今天的三百萬美金.
可想而知.
如果管理聖殿的權貴大祭司家族.
他們有誠信的話.
百姓都可以無憂.
所以有買賣有兌換銀錢是沒有問題的.
是實際需要.
但是買賣中間有沒有出現問題呢?.
這裡就需要參考一些經外文獻.
這些經文出自什麼地方呢?.
就是猶太人自己的拉比文學.
昆蘭古卷或者約瑟夫所寫的等等.
之前我們提過.
說去曠野死海北邊.
有一個地方叫做昆蘭昆蘭.
那裡有一群祭司的人士.
他們不滿耶路撒冷的權貴.

$^{361}$他們就退到東邊的死海上面昆蘭.
他們就在曠野那裡聚居.
在那裡抄很多聖經.
做很多解經.
然後自己組成一個社群.
他們說這個社群才是真正的聖殿.
真正的聖城.
有不少的解經書都提到.
耶路撒冷山上的那些大祭司.
是邪惡的祭司.
Wicked Priest.
說他們貪財.
說他們奢華燕樂.
說他們積聚財寶.
囤積很多百姓的資源.
敲詐窮人.
掠奪人民等等.
我們看一段哈巴谷書的解經.
他們怎麼解釋這段經文呢?.
哈巴谷書二章.
他們說這段經文是關於邪惡的祭司.
就是耶路撒冷.
在他任職之初的時候是忠誠的.
但是當他管治以色列的時候.
內心自負了.
為了財富.
棄絕了上帝.
以及出賣律法.
就是跟那些壓迫他們的外邦人.
反對上帝的野蠻人.
一起去同流合污.
積聚很多財寶.
並且沒收公款.
就說他招致更多嚴重的罪惡.
也行了很多淫蕩骯髒的事情.
約瑟夫是第一世紀的猶太人學者.
他寫猶太人戰史的時候.
他就說到.
在第一世紀.
主後66年左右.

$^{401}$第一次猶太人的暴亂爆發的時候.
當時耶路撒冷的羅馬兵力是不夠的.
所以那些革命分子就攻進了耶路撒冷.
然後燒毀了大祭司的房子和物業.
之後他有一段這麼說.
他說這些革命分子到了檔案收藏的地方放火.
他們急忙地去燒毀屬於債權人的合約.
去解除負債之責.
這麼做就可以贏得百姓的支持.
又可以說服更多窮人加入革命.
不怕去反抗這些富戶.
這裡就表示一世紀的聖殿權貴是敲詐百姓.
貪污欺壓他們.
所以叛亂分子做出了這些事情.
等於今天那些人走進去中央銀行.
燒毀了那些負債的記錄.
燒毀了那些伺服器和電腦.
有了這些背景的時候.
就知道耶穌首先榮進耶路撒冷.
是一個政治的行動.
潔淨聖殿更加政治性.
因為王者來到的時候.
他震動了整個國民每一方面的生命.
對上帝的敬拜和他們社會的責任.
兩者是不能夠分開的.
耶穌怎麼說都是大怒.
意憤填胸.
他就去清場.
在聖殿清場之後他就說什麼呢?.
他就教訓他們.
經常不是記著說.
我的殿必稱為萬國禱告的殿嗎?.
你們倒是他成為賊窩了.
有兩方面的指控.
第一方面就是萬國禱告的殿.
舊約第一個所羅門建的聖殿.
以色列人和俄國人在聖殿山中間沒有分開.
事實上所羅門獻殿的時候禱告.
他說歡迎萬國的人來認識耶和華.
來敬拜神.

$^{441}$但是到了第二聖殿的時候.
出於那些權貴發展了很多潔淨的規條.
大希律的聖殿就多了一個外邦人院.
有圍牆隔住.
今天也找到圍牆的石碑.
叫做 Temple Warning Inscription.
聖殿警告石碑.
翻譯出來大概是這個意思.
外邦人嚴禁內進.
是圍著聖所的.
被捉拿者歸咎自己殺生之禍.
警告外邦人千萬不要進來.
不要讓聖殿變得骯髒.
其實耶穌是在引用《以賽亞書》第56章.
上下文說什麼呢?.
他說神會使外邦人成為他的子民.
甚至太監.
如果守約歸向神的話.
他也成為神的子民的一部分.
能夠去敬拜他.
身體有殘疾.
不能夠去就近神的那群人.
那些人上帝都歡迎他來到他的面前.
《馬太福音》第21章說到結證聖殿的故事.
加上了這一句.
在殿裡面有黑眼的有瘸腿的.
來到耶穌面前.
他就治好他們.
然後有小孩子唱讚美的聲音.
在殿裡面當然不是在聖所裡面.
在聖殿山的平台上.
有殘疾的人來到主的面前.
他就治好他們.
歡迎他們來到神的面前.
第二個指控就是拆窩的指控.
剛才那些經濟的背景.
你就知道為什麼耶穌指控這個聖殿是一個拆窩.
我們到目前為止認識了一點點地理.
我們就知道強盜躲在哪裡了.
強盜躲在耶路撒冷東邊以外.

$^{481}$加拿大以外的猶太曠野的山洞裡面.
他們就拿了不義之財收藏了.
囤積在自己的山洞裡.
耶穌就反過來說.
聖殿那些人.
那些權貴才是賊.
他們就匿藏在石牆後面.
在聖殿裡面.
以為自己很安全.
其實他們在引用《先知耶利米書》第七章.
那裡說什麼呢?.
第一個聖殿.
所羅門的聖殿.
那些人就自我安慰.
他說總之耶和華的殿在我們中間的時候.
我們就一定會平安的了.
不會有難的.
耶利米就說.
你們不行公義.
你們倚靠謊言.
你們欺壓寄居.
孤兒寡婦.
不理被欺負的人的時候.
如果你們不改變.
不行公平的話.
聖殿都會滅亡.
結果當然就是在公元前586年.
巴比倫君來遷滅了這個聖殿.
銷毀了這個聖殿.
而耶穌現在結成聖殿.
其實也是審判第二聖殿.
他就告訴他們.
挑戰他們的權力.
如果他們不肯行公義的話.
聖殿必定會被毀.
結果就是一代之內.
在主後七十年就是被毀.
我們看這個結成聖殿的故事.
其實叫做結成聖殿也不是很對.
因為不是結成了再用.

$^{521}$其實是對聖殿的一個審判.
對聖殿的一個抗爭,抗議.
Temple Protest.
你讀真一點這個故事的時候.
你想一下真的很激烈.
第一方面我想可以想的.
就是你對耶穌基督的印象是什麼.
你怎樣認識耶穌.
正如之前所說.
就影響我們的生活.
我們作為基督徒.
很多時候心目中的耶穌.
是很溫柔,很善良的.
這個沒有錯.
但是這段聖經顯示出.
耶穌的靈裡面.
就是祂的義憤,祂的怒火.
祂的公義憤怒.
如果耶穌是這樣的時候.
我就問一下.
我們自己在社會當中看到的東西.
我們有沒有同樣的義憤.
不是要去破壞其他東西.
而是我們有感受到.
上帝對不公不義的感受.
然後我們去捍衛上帝的公義.
上帝的真理.
所以我們想到.
上帝看敬拜的時候.
不是純粹是我們個人.
在上帝面前的感受那麼簡單.
也說到我們心靈的時候.
不是純粹我們個人做善事得行的問題.
而是說到社會的結構.
社會裡面我們有的責任.
我們現代人很喜歡將.
宗教信仰和社會的生活拆開.
宗教信仰就是個人修養.
敬拜禱告.
是我自己個人的.

$^{561}$私人的東西.
而外面的社會.
經濟,政治等等.
這些是世俗的事.
我們分得很開.
喝水不犯井水.
但聖經信仰不是這樣的.
並且你知道實際上根本分不開.
這樣看來.
宗教,經濟,政治都是不能夠分割的.
正如當時的聖殿.
不是純粹宗教性的建築物.
是宗教,政治,經濟的中心.
我們觀察一下這個世界.
我們看到很多東西.
我們不能夠單單說.
這些政治的東西我們不要理.
其實有些什麼真是純粹政治的.
最近在黎巴嫩的大爆炸.
很多平民很慘.
但背後是什麼問題.
是政府官僚的問題.
有些爆炸品放在樹上有好幾年的時間沒有人理會.
所以令平民百姓受苦.
你想想有些國家有些政權.
將異見者打壓.
將他們收監.
這個不是政治問題.
這是上帝看到的不公義的事情.
傷上帝的心.
並且令上帝憤怒.
並且關心的不單是某某人做錯事.
而關心的有時候是社會制度的結構出現問題.
正如第一世紀的聖殿.
制度做到一個地步.
他們是敲詐,坐享其成.
然後欺壓一些無權無勢的人.
耶穌所見到的是不義的事情.
是需要去抗衡.
所以救恩不是純粹救我個人的事.

$^{601}$而是神在這個世界作王.
他是王者來到.
在我生命每一個角落.
在我個人,社會裡面.
他都需要作王,去做統治.
你留意耶穌的潔淨聖殿不是潔淨了之後再用.
四福音都有講潔淨聖殿.
但約翰福音講到耶穌挑戰聖殿的時候.
就特別強調耶穌的死和復活.
約翰福音第二章十八節.
因此猶太人就問耶穌.
你既作了這事還顯什麼神蹟給我們看.
耶穌回答.
你們拆毀這殿.
我三天內再建立起來.
猶太人的貴族,權貴就問他.
這個殿石造的殿四十六年才造成.
你三天內就可以建成.
這裡就加了一句.
事後門徒再想起這件事的時候.
他們想起什麼.
但耶穌說是指他的身體為殿.
所以他從死裡面復活的時候.
門徒就想起他說過的話.
信了聖經和耶穌所說的.
耶穌潔淨聖殿是帶來一個新的秩序.
一個新的聖殿.
承傳了舊約裡面的聖殿.
真正的聖殿是什麼.
是耶穌的身體.
是那些門徒.
是應信耶穌的社群.
所以舊約說到聖殿.
就是一個地方.
現在說到聖殿.
就是基督裡面的一群人.
在新約裡面其他的經文都是這樣告訴我們.
比特前書舉個例.
二章說主耶穌基督是活石.
固然是被人所棄的.

$^{641}$卻是上帝所揀選的所寶貴的.
你們來到主的面前有一塊活石.
我們就是活石.
是建造一個聖體的寺廟.
我們作為祭司.
藉著耶穌基督獻上上帝所閱立的靈祭.
唯有我們就是被揀選為族裔軍專的祭司.
聖經的國度.
肅神的子民.
我們就是那群人.
要叫你們宣揚你們出黑暗.
入光明者的美德.
所以教會就是基督的身體.
就是聖殿.
我們在生活當中就是在打人家來到主的面前.
宣告耶穌已經是作王.
結束的時候.
我覺得有一個早年寫的.
一個世界夫君派的信仰宣言.
我揭露其中一段.
他說得很對.
這個叫做落喪信約.
是一個世界無論什麼話.
夫君派的領袖都走在一起.
來到有一個這樣的綜合宣言.
其中一段是講基督徒的社會責任.
這裡說.
我們確信.
上帝是全人類的創造者和審判者.
所以我們應當共同擔負起.
他對人類社會的公義和和好的關注.
以及對那些受各種壓迫的人的自由的關注.
救恩的信息也都包含對各種形式的疏離.
壓迫,歧視的審判.
無論何處有罪惡與不公正的事.
我們都要勇敢地斥責.
不僅必須努力在這不義世界中彰顯上帝的公義.
還要傳揚他的公義.
我們所宣告的救恩.
應當在個人生命和社會生活各方面都改變我們.

$^{681}$救主憐憫我們.
讓我們教會成為真正主的聖人.
宣揚他的真理,他的救恩,他的公義.
我們同心低流祈禱.
主我們求你來憐憫我們.
很多時候我們把救恩當作私有的東西.
只是為了我們自己的好處.
只是為了我們有平安.
主在這個不義的世界中.
一個這麼破損的世界中.
讓我們宣告自己的真理,自己的王權.
主讓我們的生命,我們的社會.
每一方面都能彰顯你的榮耀.
讓我們的福音能走進社群中.
走進這個世界中.
我們去表揚主你自己的美德.
我們這樣祈禱.
奉主耶穌基督的名字.
\newpage



\section{}
\label{sec:AndEm9fTjLE}
\textbf{UV 匯客室 - 鮑維均博士}
\newline
\newline
連結: \href{https://youtube.com/watch?v=AndEm9fTjLE}{\texttt{https://youtube.com/watch?v=AndEm9fTjLE}} ~~~~ 語音日期: 2024-04-01
\newline
\newline
\hyperref[sec:9_6WJzTToTc]{\small{< < < PREV SERMON < < <}}
~
\hyperref[sec:index_chronic]{\small{[返順時目]}}
~
\hyperref[sec:index_scriptual]{\small{[返順卷目]}}
~
\hyperref[sec:BAxyx20fKwY]{\small{> > > NEXT SERMON > > >}}
\newline
\newline
$^{1}$各位城會的網友 每一位觀眾.
我們今天有一個很特別的節目.
就是訪問一些世界各地不同的弟兄姊妹.
今次很榮幸找到Dr. Pao 包惠軍博士.
和我們聊聊天 特別是講到他的學術旅程.
他在聖經研究上的學術進程.
在裡面有甚麼起起伏伏的經驗可以和我們分享.
我們大家都知道包惠軍博士剛剛成為了.
三藩市神學院的學術院院長.
是一個很難得的機會.
我們覺得一個華人的學者.
能夠晉身在英文當道的世界.
在聖經研究裡面有很多貢獻.
Dr. Pao 我們聽你講道 聽了很多次.
我想世界各地很多的觀眾 很多的信徒.
都得到很大的益處.
特別是你的著書 不單止講道.
你的講學等等 我們都非常得到助手.
但我們很少聽你自己的學術生涯是怎樣的.
所以今次希望能夠在對話裡面.
能夠聽一下你的心聲.
你這麼多年來的學術旅程裡面.
有甚麼經歷.
第一個問題我想問Dr. Pao的就是.
你是怎樣進入聖經研究學術的範圍.
這個領域.
首先多謝Mark Sir的問題.
亦多謝成會等姐妹.
給我機會分享我自己的心路歷程.
其實今年對我來說是一個特別的一年.
因為我在中學年代奉獻.
是1981年 所以如果數起來.
這就是剛剛40年的時候.
所以在1981年奉獻的時候.
我還是一個中學生.
那時候是不太知道甚麼叫做奉獻的.
只知道神的愛呼召我.
我就應該有合適的回應.
大家知道那時候的香港中學教育.
是有分文科,理科.

$^{41}$所以我那時候還是理科的學生.
但是那時候因為這樣的奉獻.
在會考過後.
我就在中六變成文科的文科生.
就是為了以後希望能夠在神學訓練那裡.
能夠有份參與.
所以那個是一些背景.
1980年中期的時候就來了美國.
我大學的時候.
就是讀希臘文和哲學的.
都是為了以後讀神學的路去鋪排.
那時候你奉獻的時候.
就已經想做學術的教學.
其實也不太知道是怎樣的.
因為我自己爸爸也是一個牧師.
在仇人院裡面教書.
可能也有受他影響.
但是其實那時候有一個很深刻的印象.
最少在當時的華語世界裡面.
我感覺就是很多的學者能夠對文化作出回應.
但是那些學者都是神學科的學者.
而不是聖經科的學者.
所以在中學年代有這樣的念頭.
就是說如果我是讀聖經科.
能不能夠對時代有回應呢?.
所以有這樣的念頭.
就是希望讀新九約.
而不是讀神學科.
或者系統神學科.
所以在中學年代有這樣的念頭.
希望是這樣.
還有另外一個背景就是.
身為牧師仔.
我媽媽那邊其實我是第五代的傳道人.
所以爸爸是第一個基督徒在她家庭裡面.
但媽媽那一堆在我身上也是傳道人.
所以我想牧師兒子在家庭長大.
也有一個感覺就是.
究竟這個信仰是遺傳還是自己.
或者是他們的信仰.

$^{81}$是否能夠回應今天的挑戰.
所以我從小已經對一些傳統看聖經的.
standard answer(標準答案).
已經有些懷疑或者不滿.
很多時候有些比較難的經文.
我們說為什麼會這樣?.
很多時候是一些quick answer(快答).
叫我們不要問.
所以這類的不滿.
可能令我對聖經有興趣.
究竟聖經是否真的這樣?.
是否這麼容易回答一些很膚淺的問題.
有很膚淺的答案.
所以可能有這樣的背景.
如果在一個傳統的環教會.
像你這樣說.
問一些難的問題就不想回答或者不讓問.
會不會是年輕的時候已經有一種問題少年的形象?.
可能是的.
我也跟朋友聊過天.
他說你有沒有反叛期?.
我說當然沒有.
他就說不是.
你的反叛期不是吸毒不是吸煙.
他的反叛期是對聖經解釋有些不滿.
是你那種反叛.
我回想也可能是的.
我的反叛期不是吸毒吸煙的反叛.
而是對一些傳統的答案有些不滿的反叛.
另外在華人的文化裡面.
我不知道你第五代的基督徒會不會面對這個問題.
讀理科是比較好的.
讀文科是看低一點的.
當時有沒有這樣的經歷?.
也有的.
在中學年代.
我學校裡面.
有五班是理科一班是文科.
特別是我男校.
所以文科基本上是沒有人看重的科目.

$^{121}$也是賺不到錢的科目.
所以我想在香港就很明顯有這樣的感覺.
一來到美國就沒有了.
因為來到美國之後就覺得什麼科目都是自己的興趣.
那時候你奉獻的時候.
有沒有這樣的心智呢?.
父母或者家庭是不是也有支持和鼓勵呢?.
我想我最初奉獻的時候.
我還記得是1981年的8月培靈會晚上的奉獻.
我回來的時候吵醒了爸爸.
我說奉獻了是什麼意思呢?.
他說有很多不同的方法神父像一個人.
其實那時候對我來說是很混亂的.
對我父母來說他們差不多肯定這個不是一個奉獻.
就是覺得你小朋友可能是一時衝動.
所以教會我還記得有一些奉獻的團體邀請我參加.
父母都叫我不要去.
因為覺得我太小了.
或者是不會想清楚這個問題.
但是這個呼籽在我心裡面.
所以其實是真的清楚呼籽是兩次.
一個就是8月的培靈會那裡.
一個就是12月的東靈會.
所以給我大概半年時間去想.
到12月我再一次走出來.
我就大概清楚自己是做了什麼.
知道是要全時間在全職裡面去獻上.
當然慢慢觀念就清楚一點.
就發覺原來全時間奉獻不需要全職奉獻.
平信徒也可以全時間奉獻.
但是我的應該是全職奉獻.
但是在我全職裡面.
我怎樣在哪個崗位裡面做得更加好呢?.
這個就慢慢摸索.
剛才說到你來美國.
有什麼考慮以至你決定來美國?.
還有為什麼讀希臘文和馬來西亞的學?.
其實我是在香港想讀歷史的.
在香港一間大學裡面讀歷史學.
歷史是這樣的.

$^{161}$但是如果是我這樣上一年級的人.
都會記得香港一些暫取生的制度.
當時的意思就是說.
你中學會考完之後.
中文大學就會收你.
可以早一年進去.
或者是收了你之後就叫暫取生.
早一點就能夠進到這間大學.
但是當時我的中學.
只可以讓我們選中史或者西史中文.
原來中文大學的歷史系.
是要求中史和西史都有.
所以我就不能進入歷史系裡面.
所以其實是這樣令我再想.
如果真的發展的時候會怎樣呢?.
當時我姐姐和我爸爸都是在美國讀書.
媽媽也是.
所以當時就想到不如去美國讀書.
是這樣開始的.
為什麼會選擇Whitman College?.
其實沒有選擇的.
其實我奉獻的經歷.
我以為自己很特別.
是12歲信主.
中學年代再一次決志.
大學年代有懷疑.
然後慢慢再找回自己的信仰.
我後來去到Whitman的時候.
就發覺原來很多牧師都是這樣.
都是12歲,中學,大學.
完全的旅程是一樣.
我自己去Whitman.
因為我爸爸是Whitman的畢業生.
所以感覺就是沒有選擇.
我去大學去那裡呢?.
爸爸說去Whitman.
我姐姐就去了.
所以當時是沒有選擇的時候.
就去了Whitman.
去到後也是沒有後悔.

$^{201}$因為讓我能夠在哲學和希臘文裡面.
有些紮根的機會.
你選擇哲學和希臘文.
當時已經想了可能是在新約世界這方面.
最初是想去新約的.
但也對教會歷史有興趣.
所以中間在碩士課裡面.
在短暫的時候是轉入教會歷史.
但最後還是回到新約那裡.
所以最早也是對新約研究有興趣.
所以你在Whitman就完成了你的碩士.
讀到碩士.
然後你就去了哈佛.
那時去的原因是因為.
可能在福音派圈子裡面也有四,五年的時候.
想嘗試一些不同的學術體驗.
所以那時在哈佛有一個很靈活的程式.
就是能夠很容易讓我自己選擇科目.
那時我的強調或想讀的就是兩個中間歷史.
所以那個程式讓我能夠在兩個中間歷史裡面.
碰多一點.
所以就這樣進入了哈佛.
這個可能是大家可以好奇的問題.
就是為什麼你會選擇兩個中間新約這段時間.
是你以往年青的時候一直問的問題都是這方面.
還是你自己的興趣.
是怎樣去縮減到有這麼多領域去選擇這個.
其實真的怎樣選擇我也不太記得.
但如果真的想起來的時候.
有一個可能性就是覺得.
當時華人比較多的是讀舊約.
所以我就想去讀新約.
其實每一波的,每十年有些不同的.
有段時間有些人讀新約,有些讀舊約.
在我年代似乎是讀舊約的人比較多.
所以我是這樣選擇了新約.
那你很順利就進了哈佛大學.
在那裡的經歷是怎樣.
我想是第一個星期的感覺是最特別.
因為當然我知道神學研究的領域是很大的.

$^{241}$有不同的取向等等.
但沒有想過一個無神論者可以教新約.
也沒有想過可以由很不同的眼光去看新約聖經.
所以第一個星期的那種心靈的掙扎.
我還記得的.
覺得這樣也可以這樣去看.
所以最初那兩年是覺得.
很大的好奇心.
想迫自己擴闊自己的思維.
覺得應該學術誠實.
不要太執著.
所以讓自己嘗試不同的立場.
所以那兩年有這樣的感覺.
掙扎會不會是一些.
好像我們經常聽到.
去到一個聖教環境的時候去做研究.
對信仰有些衝擊.
是不是這一類的掙扎.
是的 但其實我最初去的時候.
我的心態也不是一個護道者的心態.
我的心態不是去打一頓死.
我的心態就是我想擴闊自己.
給自己一些思考的空間.
究竟我的信仰是否真的這樣.
我真的想有些開放.
我還記得當時很出名的新儒家.
一個大師帶一個研討會.
我也很坦白地分享.
因為我最後的程式.
其實是一個比較宗教的程式.
所以我八個同學裡面.
有讀猶太教.
有人讀儒家.
我是讀基督教之類的.
不同的宗教.
這次我對同學分享的是.
其實我在這裡想尋求.
什麼是真,什麼是對.
所以我當時的心態.
其實是一個開放的心態.

$^{281}$想給自己一些空間.
可能因為自己是一個牧師仔.
所以覺得應該給自己一些空間.
去尋找自己的信仰.
所以有這樣的心態.
所以是有掙扎.
但不是一個護道者的掙扎.
是一個可以說是尋道者的掙扎.
就是什麼是真.
什麼是可以成為我們信仰的根基等等.
會不會在掙扎裡面.
有拆毀了一些傳統對信仰.
或者對聖經的了解呢?.
我想有很多都是.
雖然不是很清楚能夠說出.
拆毀一件什麼.
因為可能當時我架構的思維.
都不是那麼清晰.
沒有一件東西是建立了要拆毀的.
都是一些零碎回來的思維.
所以是有拆毀.
而且過程是有一點痛苦的.
因為何去何從.
不知道應該跟哪一範的感覺.
同時也會在學術環境裡面.
有些壓力應該怎樣去看文本等等.
所以都有不同方面的掙扎.
如果有那麼多年在哈佛大學的經驗.
你如果問你最深刻的一個掙扎的經歷.
你會不會有一些印象呢?.
其實感恩神在那段時候.
有幾方面是對我有幫助.
一方面就是一班朋友.
那時候我們有大概四個都是.
同一間學校讀書的.
不同系裡面的.
我們經常都會走在一起討論.
一個讀心理學.
一個讀Genetics.
我太太讀的.

$^{321}$一個就是舊有聖經.
所以我們四個人不停地討論.
大家的衝擊,學術裡面面對的困難.
不同的課題.
所以在這個群組裡面.
我覺得是對我有幫助的.
另外一個對我有幫助的就是教會.
在所謂的Local Church.
我慢慢覺得讀神學一定要plug into一個Local Church.
在那個地方要有成長.
以及要accountable.
因為我發覺讀神學的人可以越讀越囂張.
覺得自己懂很多東西.
可以超越,用理論去概括我們熟悉的經歷等等.
但是一個Local Church逼我們去做一些不是離地的基督徒.
所以這是第二個對我的影響.
第三個影響就是.
多數的聽眾都知道我太太在那個時候生病的.
所以我在90年代.
90年代去到哈佛.
她在93年患上癌症.
所以那個是我讀PhD的第一年的時候.
所以那個的經歷就很深刻了.
因為那次是第一次.
我自己覺得是真的要跟神摔跤的感覺.
就是這個信仰是否真的可以持續下去.
還是這個時候應該全放棄.
所以那一刻可以說是我八年教育中最重要的時刻.
亦想起跟教授的交往.
其中一個教授是Brundtmann的最後一個弟子.
他在哈佛很有地位.
那時候我跟他做助教.
也聽了他很多時候的故事.
他參與納粹黨為希特勒打仗.
然後回到大學讀書之類.
所以聽了他的故事.
但他是一個相當極端的人.
其中一個極端就是他堅持復活是不可能發生的.
我記得那時候看了一本書.
是James Dunn的一本書.

$^{361}$是講耶穌的證據.
那本書裡面還有一個部分是專門寫明.
是反映這個教授的復活立場等等.
所以他是有名的.
但我覺得特別的是.
我太太病的時候.
最後離開世界.
她竟然寫了一封信給我.
到今天還留著.
那封信裡面說如果人不信復活的話.
就比任何的動物更加悲慘.
這當然是保羅的說話.
但我心想為何你可以在學術研究裡面堅持復活是不可能的.
但到了人生經歷某些事情的時候.
或者是一些存在的時刻.
Existential moments的時候.
為何又會提起復活這回事呢?.
所以就逼我想究竟我現在讀的東西.
和我的生命有沒有掛勾.
所以那個經歷對我來說是一個很重要的經歷.
這個學術的期間有多少年?.
你講到90年初.
90年夏季我去到波士頓.
然後是98年離開的.
所以有八年時間.
八年時間.
我覺得很好.
你講到我們生命裡面有苦難和有群體.
學術的研究裡面是不會離地.
是能夠貼地.
是能夠幫到教會.
是能夠為教會服侍的.
你什麼時候入現在的Trinity?.
在大概95,96年的時候.
那時候我離開世界還沒有結婚的時候.
還沒有再結婚的時候.
我和一個弟兄一起住的.
大家都是同學.
他是心理學.
在我們交談當中.

$^{401}$我突然間說了一句話.
如果來Trinity教書就好了.
但我還沒有在三藏院讀過書.
但我有個印象就是.
三藏院能夠堅持聖經的權威.
但亦能夠在學術裡面開放和有研究.
所以對我來說這是一個很好的平衡.
那時候我講了他這樣的說話.
但我亦想起我讀書的時候.
其實是答應神一件事.
就是不會報美國的學校去教書和教學.
因為那時候我答應神會回亞洲事奉.
我和神說如果讓我在美國留下.
你要找學校接近我.
我不會接受.
這一點是沒有可能的.
在美國誰會接近你.
但我和同房說完這句話的時候.
三藏院請我真的好.
真的Trinity有人在那時候接近我.
是當時的新約系.
所以就是這樣開始一個面試的過程.
這個是你已經認識的教授.
認識的朋友.
其實是我學院的教授.
但我都不知道他去了Trinity教書.
學院認識的教授已經離開了Whitman.
去了另一間神院教.
後來再來了Trinity教.
但我都不知道.
突然間收到一個電郵.
叫我考慮.
好像很順利.
你好奇問一句.
我也覺得很奇怪.
上來就帶領你.
為什麼會這樣發生.
對我來說是一個神的引證.
叫我留在這裡.
所以都覺得很奇妙.

$^{441}$如果那時候亞洲的神學院去接近你.
他可能就會接觸仙燈.
其實已經在接近.
我收到這些消息的時候.
我在香港教短期課程.
我還未承諾.
還在想應該怎樣做的時候.
如果我想回時間線.
我是教短期課程之前.
Trinity接觸我.
但接觸的過程很長.
不是說接觸就請.
所以那個過程大概有一年.
所以那一年之後我就決定留在Trinity.
所以你完了哈佛.
幾乎立刻就去了Trinity.
是.
其實是如果說時間線的時候.
Trinity請我.
我還未寫論文.
還未在哈佛寫論文.
那是很普遍.
因為現在我們不會請一個未有PG的人.
但當時我只是考了Holmes.
未考試但未寫論文.
所以我回去寫的時候.
我剩下大概少於一年可以寫論文畢業.
我就答應自己.
我要半年做好論文.
教授就說三年.
我說不行.
我只剩下半年.
我不想再寫.
所以那半年我真的沒有動過.
我連下街查信都沒有.
所以是一個很有趣的故事.
我坐在那裡坐了半年.
之後我要針灸了一年.
我失落過才沒有痛.
因為那半年基本上沒有動過.

$^{481}$即是全程24/7寫論文.
所以你入了Trinity.
又是你自己心願裡面想入的散落院.
入到來的時候有沒有一些令你驚訝的地方.
令到你覺得我本身不是預期這樣.
我想兩方面的驚訝.
由一個新學位入到一間受院教書.
第一個驚訝就是.
原來我們要在教學裡面很看重應用.
因為多數學生最後都是一個牧者型的學生.
所以那種應用是驚訝我.
因為我們任何一個新學位入到來.
都以為在教書的時候要展示我們學了多少東西.
學了多少東西.
原來真的不是學了多少東西.
而是怎樣應用.
所以這是驚訝.
第二個驚訝就是.
在一個白人社會裡面.
我想那個是有少許預料到會有掙扎.
但是始終來到都是一個驚訝.
因為特別是福音派的受院.
如果不是福音派的受院.
其實是很大規模的.
但是福音派的受院基本上是很白人的文化支配.
所以我還記得來到第一個星期.
去教書的時候遇到一個學生教授.
他問我在做什麼.
我就說我在教書.
他覺得沒有可能.
所以第二天再遇到我的時候.
他就說你在這裡教書嗎.
你一定是助教.
覺得你沒有可能是一個全職教書.
所以頭幾年其實不停有這樣的掙扎.
這些才是最難答的問題.
很多年都有人問我何時畢業.
因為我一直以來都在讀書.
你在三藝神學已經教了二十年左右.
二十三年了.

$^{521}$你講了兩方面的掙扎.
第一方面就是怎樣落地.
學術怎樣幫助牧者形的進度.
這個掙扎怎樣調整.
第二個就是在一個白人的宗教圈子裡面.
你怎樣去不能夠說脫穎而出.
怎樣去中心繼續去做你應該做的事.
神父照你的事.
先講第一個.
一個學者去到神學院教書.
我們那些牧者讀過神學的.
每個人都牙痛錦聲.
他們覺得讀一些東西跟他講台或是教導.
是沒有關係的.
他們最大的抱怨就是這樣.
看的書對他預備講章沒有什麼好處.
作為一個教授你怎樣看.
怎樣去調整這方面.
或者怎樣去隔離兩方面.
其實我在教神學的時候.
我印象當中的模式都是受我長大.
在香港環境的模式影響的.
因為在中國人社會裡面.
一個神學教授其實不是一個象牙塔式的教授.
是一個跟教會有很多交往.
和很多服事的教授.
但我覺得很驚訝的.
很多在美國神學院的教授.
是從來不會講道.
他們是一個研究式的教授.
所以我最初開始的時候.
是想多一點在講台裡面事奉.
因為有兩個原因.
一方面是想不離地.
另一方面是發覺自己這麼多年的研究.
其實很狹窄.
所懂的東西可以很深但很窄.
所以很多東西真的不懂.
所以就想用講台去豐富自己.
特別是我喜歡的就是講多一點舊約聖經.

$^{561}$因為是補足我自己的不足.
所以在那些經文裡面的掙扎.
其實對我來說是一個教育.
Mark Sir也知道.
我和你差不多是時候出道.
所以我講台裡面的事奉.
也是逼我自己去嘗試.
不要太離地.
去看看聖經和信仰的掙扎的銜接是怎樣.
可以講一下你當時三十年代.
或者現在三十年代的教授事奉的環境.
你覺得當時你這樣做是否比較獨特.
逼自己去講道.
你剛才講得很對.
有些神樂院的教授講得不好聽.
其實他講道是不行的.
但在課堂學術上當然是非常深厚.
功力很深厚.
你是否比較特別一點.
逼自己去教會去事奉.
其實是,但我想也是一個過程.
最初幾年我都不是想在這方面發展.
我只是覺得講道是一個平時偶爾做的事情.
不是我自己的事奉中心.
但在川梯教五年之後.
有一間神樂院接觸我在亞洲.
當時他們的院長退休.
問我有沒有感動能夠回去做這個崗位.
當時其實是很大掙扎的.
因為我還覺得亞洲是我的家.
或者香港是自己的家.
第二個感覺就是覺得自己應該事奉自己群體.
自己膚色的人.
第三個感覺就是川梯教五年對我來說.
都不是自己的家.
還是一個客女的感覺.
所以這三個不同的感覺.
令到我十八年前很想離開川梯教.
當我在那段時間和部門的人分享的時候.
他們就說過,不如你自己留在這裡.

$^{601}$同時間可以服侍中國人教會.
做你想做的事情.
但也有空間讓你寫作或是預備講章等等.
所以當時我就開始去決定.
不是想一個典型的白人仇家會做什麼.
而是想想自己想做什麼.
所以在那時候我開始講道多.
或者是服侍不同的群體多一點.
所以那個轉變是十八年前.
所以那次之後我就覺得.
我自己做的事情是我想做的.
就這樣算了.
不只是想著我在學術的地位是怎樣.
所以我嘗試是這樣開始的.
另外一個掙扎就是在白人社會.
這個我想很多觀眾都感興趣.
怎樣能夠在這個過程裡面.
由一個亞洲人在一個幾乎所有白人的圈子裡面.
怎樣在學術上能夠有所成就或是有所進步.
裡面的掙扎就是你剛才說的.
別人以為你不是在教書.
一直以為不是在教書.
你覺得有那種的認受.
起碼在你的圈子裡面.
我想我們在這裡做的一定要出版.
甚至在英語世界裡面都要出版.
無論是期刊和書都要出版.
所以那裡差不多成為一個教會的入口條件.
Publish or perish.
所以就要進入這樣的圈子.
但其實在過去二十年.
我覺得在福音派的神學院的文化都開始轉變.
是二十五年前很難幻想的.
現在的教育學院是看重多元化多於以前.
因為明白原來每個人都戴著自己的眼鏡去看聖經.
所以比較年輕的教授就容易接納這種觀念.
所以我想這是二十多年的發展.
但對我個人來說最大的轉變.
就是在大概十年前開始做新約系的department chair的時候.
那時候我終於開始覺得這是我的家.

$^{641}$因為能夠建立一個faculty是我自己想建立的.
所以在我做的時候都請了幾個教授.
所以現在department裡面差不多全都是我自己請的.
所以在過程當中我們開始建立了一種大家共識.
亦都覺得這是一個我們team可以一起work with的.
所以最後我們建立了一個相當元化的faculty在我們新約系裡面.
當時我們請的有澳洲人,有新加坡人,有男有女等等.
有不同的背景的人來到.
所以感覺跟之前的感覺很不同.
是大概十年前開始的.
你剛剛在這裡做了一個bridge.
可能舊的一代剛剛要飛樓.
現在新的一代又可以去建立一些新的東西.
建立一些新的環境等等.
是的,最初我做的時候都覺得.
如果我去外面事奉或是港島.
二十年前我覺得是一些科娛活動.
或者是一些我事奉的關係不大.
現在反過來看原來發覺那些東西很重要.
慢慢發覺原來在方牌圈子裡面.
他都開始value這種經驗或經歷.
所以我覺得是相其妙的.
我最初想著一些科娛活動的經歷.
成為了一些最後跟我主流事奉的關係的背景.
不如說一下你過去這麼多年學術著作.
有沒有一些你覺得很滿意.
很有sense of fulfillment.
很滿足的.
會是哪一方面或者哪本書,哪一本著作.
我想最滿足的就是一本未出版的書.
寫完了.
如果學術研究,大家知道Brio出版社出版一些比較學術的書.
他們剛剛committed出一套色經書.
我負責的是提摩太前後書和提多書.
主編是Stanley Porter,大家讀新約的人認識.
所以這套書我寫的時候有很大掙扎.
因為我最初研究都是科文書.
但被邀請寫提前,提後,提多書.
其實不是我自己的範疇.
但既然被邀請也不會死路.

$^{681}$因為這套著作.
所以用了十年.
十年都是我覺得挺痛苦的時候.
因為要處理一些我未處理過的問題.
去分析一些沒有一套色經書會用這種方法去分析的.
在希臘文學裡面的研究.
所以每一步都是一個掙扎.
但到最後完的時候是成因點.
但過程是十年.
現在過了就是遲些會出版.
不如說說將來.
你現在做了academic dean,院長.
首先第一,院長你主要的job description是做哪一瓣呢?.
其實比較特別的是.
視乎你的神學院的架構是怎樣.
在一間獨立的神學院.
其實校長就是校長.
但在Trinity的大學制度裡面.
每一個學院的dean就變成院長.
所以其實可能很多人知道Trinity是一個law school在加州.
Trinity在佛羅里達州是一個graduate school.
所以我們有幾個不同的學院.
有undergrad.
所以神學院院長的職責主要是.
去想想神學院自己長期的發展.
特別是教授團裡面的.
怎樣組成,怎樣能夠繼續有效去延續這個時風.
所以處理的多數是跟學術有關係.
跟教授有關係.
跟學生的program有關係.
或者是以後怎樣去develop新的program.
是跟這些有關係.
真的需要fundraising.
我們就不需要太多去想這方面.
雖然是這樣說.
但其實我們現在每天都差不多在見一些人.
在東德薩等等的人.
但這個不是我們最主要做的工作.
當然一些PR一定要做.
但要有一種學術上在這個領域有什麼觸覺.

$^{721}$將來發展的trend形勢會是怎樣.
然後請不同的人.
看看他們的著作會fit到你的vision.
你可不可以很簡單說說.
三個神學院現在做academic dean做院長.
你將來有什麼意象.
或者什麼重點.
你覺得應該向那個方向發展.
我想Trinity給人的印象是一個保守的神學院.
但也是願意跟學術有交接的神學院.
我覺得在這個年代世代.
我們需要繼續這兩方面的平衡.
因為在美國特別是南方.
很多的神學院會強調聖經權威.
但完全不會接觸到社會公義的問題.
或者是性倫理的問題.
或者是教會裡面女性文化的衝突之類.
姐妹教會裡面事奉的角色等等.
所以Trinity希望做到堅持聖經為本.
但也能夠跟文化對話.
這是我們希望做到的特色.
但其實一句話就引來很多不同的爭議.
但我們希望能夠在爭議裡面保持清晰的立場.
Trinity的另一個特色是一個堅守聖經權威的仇院.
但在這個權威上我們也堅持有不同的意見.
例如我們有Calvinist,有Arminian.
我們有女性角色的Communitarian.
所以我們希望每一方面都能夠有一個平衡.
讓人進來的時候能夠聽到不同的聲音.
也能夠有合適的對話.
所以這是我們希望能夠做到的工作.
這個當然是Trinity多年來建立的學術傳統.
一個文化.
有不同的聲音.
大家一起對話一起去爭論.
好像你提到南方的神學院.
美國南方的神學院很多時候有自己的泡沫在裡面.
繼續鞏固自己的方向.
你現在作為院長.
如果有不同的教授不同的學者.

$^{761}$他們有不同的意見.
是不是需要很多時候去調和.
怎樣能夠拉著大家.
大家可以繼續去談.
我想是.
特別是在社交媒體的世代.
是很難去將一個多元的教授.
能夠繼續有這個空間去對話.
因為很多時候某個人說的話.
會傳播到外界.
可以極端化了.
另一邊就很難去接受.
但我們教授其中一個特色就是.
我們有一個群組.
一堆教授一年四次在一個教授家裡.
輪流去發表文章.
輪流去回應文章.
所以我們教授之間的關係維繫.
雖然好像很私人.
一年四次在某個教授家裡吃飯聊天.
但其實這個成為了過去十年很重要的橋樑.
容許我們不同的意見.
可以在一些共同的屬靈.
或者是一些友誼的基礎上.
能夠有對話的空間.
所以這是我們珍惜的.
希望繼續延續.
這個是要求嗎?.
你做院長你會要求所有的教授都要參與?.
這個群組不是要求.
也不是每個人都會參加.
特別是一些比較多著作的教授就很參加.
一參加就要輪流去發表文章.
回應文章.
但最少是不同部門有代表去參加.
所以這個就成為了一個橋樑.
有沒有一些新的東西.
你想到作為一個院長.
現在或者將來你會應用的?.
我想首先是MDiv一定要是flexible.

$^{801}$因為全美國現在MDiv讀的人是越來越少.
有兩個原因.
第一個原因就是很多大教會已經不需要全團有正規的受訓練.
第二個原因就是越來越少的神學生.
願意半屋子來到一個地方住三年.
特別是越來越多second career的神學生.
就不能帶整個family來到.
所以我們希望program要是flexible.
但仍然能夠維繫一個群體屬靈成長的氣氛和裝備.
所以Trinity有一個特色.
我們有大學補的.
所以我們有幾種MDiv.
有一個叫5 year MDiv.
是可以五年之內讀完你的undergrad和MDiv.
可以五年完成.
另外一種MDiv叫APMDiv.
如果你是聖經學校出來.
可以在短時間內完成MDiv.
當然有傳統的MDiv也有Online的MDiv.
Online就是全靠covid.
迫我們faculty開始用到科技.
feel comfortable可以開始這件事.
但挑戰就是如何能夠在Online裡面建立一個群體.
我們最近也有一個grant收到.
專門研究如何在Online裡面建立一個屬靈群體.
以致這個program不只是收到學術的data.
而是能夠一起成長.
所以我覺得最大的挑戰就是如何能夠做一個responsible的Online program.
過去一年多covid迫了全部神童院.
三個神童院都要Online.
可不可以分享一下你在神童院裡面學到甚麼功課.
有甚麼pros and cons 有甚麼barrier.
有甚麼阻力需要去克服.
以致能夠成功地去實行這個模式.
我想以前的Online program.
所謂Online多數是一些銜售課程之類的Online.
就是你收到一些video tapes 隨時可以上網聽.
因為covid,zoom和不同的緣故.
我們懂得不只是有那類的Online.
是可以有同時間的同步Online synchronous.

$^{841}$所以我覺得synchronous的Online program可以做得相當成功.
因為我發覺甚至教希烈文一班.
有些學生會躲在後面.
過了頭不作聲.
但上zoom的時候.
大家距離都一樣短.
所以你不能躲起來.
所以我覺得在某些課堂裡面.
那種Online classroom甚至比現場更加有效.
但我發覺如果真的做得好的時候.
要做到一個hybrid.
意思是有學生在課堂裡面.
亦有學生在Online一起參與同一課.
做得好的時候要有很多投資在科技上.
我們都收到一些donation.
我們可以在過去一年做很多upgrade.
但仍然要繼續做下去.
希望能夠在一些hybrid的情況下.
令Online和現場的學生都能夠有同一個群體的經歷.
所以這是一個科技的挑戰.
但對我們教授來說.
年紀大的教授難以接受和適應.
但年輕的就很容易適應.
大家都要嘗試在一個新的環境學習.
現在Trinity的master program MDiv都可以全線上學.
其實很特別.
因為在美國神樂協會.
只不過是2020年7月才通過了一個法案.
一個policy.
MDiv可以全線上學.
有時候都可以.
但有時候要申請exception.
現在終於可以這樣做.
但2020年7月不是因為covid.
我們已經計劃了這樣.
所以剛剛在covid的時候.
以致我們能夠容易一些去考慮這種的可能性.
但其實我們都希望能夠藉著這樣.
令國際學生能夠多些機會去接觸.
但我們想做的fundraising.

$^{881}$就是希望能夠在一個剛剛進入program的學生.
能夠有funding來到Trinity或某個地方.
開始一個星期先建立一個群體.
所以用三年時間都會跟一個群體能夠一起相處在online裡面.
但online之前希望有些實體的對話.
所以我自己最想能夠做到的就是.
有些grant能夠讓我們這些online學生能夠travel來到這裡.
如果亞洲可能有一個location.
澳洲可能有一個location.
北美有一個location.
這樣就能夠容易做一些responsible的online program.
都說起Trinity.
時間都差不多了.
我們阻了Dr. Fowl很多時間.
我們想幫Trinity買廣告.
我們都不收錢.
純粹因為時報.
大家都很欣賞Dr. Fowl的時報.
以及Trinity所做的.
造就很多的工人.
Trinity的連結就在下面.
如果你是在考慮讀神學.
或者考慮想支持神學院.
Trinity神學院的發展.
請你去那條連結.
你都可以考慮一下怎樣去support.
其實我覺得這個新的世紀.
好像是Dr. Fowl提醒我們.
有很多的挑戰.
有很多的地方我們需要去神的話語.
能夠落到地.
能夠去回應現在的時代的挑戰.
再一次多謝Dr. Fowl.
和我們做這個訪問.
\newpage



\section{7:21-23}
\label{sec:BAxyx20fKwY}
\textbf{Kingdom Values Matthew 7\_21-23 Sermon on the Mount}
\newline
\newline
連結: \href{https://youtube.com/watch?v=BAxyx20fKwY}{\texttt{https://youtube.com/watch?v=BAxyx20fKwY}} ~~~~ 語音日期: 2024-04-02
\newline
\newline
\hyperref[sec:AndEm9fTjLE]{\small{< < < PREV SERMON < < <}}
~
\hyperref[sec:index_chronic]{\small{[返順時目]}}
~
\hyperref[sec:index_scriptual]{\small{[返順卷目]}}
~
\hyperref[sec:BG6FWzPxpgI]{\small{> > > NEXT SERMON > > >}}
\newline
\newline
$^{1}$Good morning family. I'm excited that we get to be together again and open God's Word again today..
We're continuing on in Matthew chapter 7, the end part of the Sermon on the Mount here..
And today we're going to be looking at verses 21 to 23. So let's read these together..
"Not everyone who says to me, 'Lord, Lord,' will enter the kingdom of heaven,.
but the one who does the will of my Father who is in heaven..
On that day many will say to me, 'Lord, Lord, did we not prophesy in your name and cast out demons in your name.
and do many mighty works in your name?' And then I will declare to them,.
'I never knew you; depart from me, you workers of lawlessness.'".
These are some really good verses but really challenging verses, really scary verses..
Jesus is showing us that not everyone who thinks they know God actually does know God..
And that's what we're going to talk about today, how we can know that we know Jesus.
and know that we have a genuine and authentic relationship with Him..
And I'm going to remind us that in this part of the Sermon on the Mount,.
Jesus is wrapping up the Sermon on the Mount and He's giving us words of application..
We've been talking about that the past few weeks..
And so this is an important application to our own hearts..
And in this passage, Jesus is shifting the focus from evaluating others to evaluating ourselves..
If you remember the verses that we read last week and the message that we talked about last week,.
Jesus was talking about evaluating false prophets or false teachers.
and how we could evaluate their lives and their teaching to discern.
if it's something that we should allow to influence us..
And these verses are the verses that immediately follow the passage we read last week..
It's definitely related, but Jesus is shifting the focus..
He's encouraging us now to evaluate ourselves, to evaluate our own hearts.
and think about our own relationship to God and our own standing in the kingdom of God..
And Jesus is teaching us that kingdom values are not simply about good works..
We've been talking about kingdom values for almost a year now..
We've been in this study and it's been really good..
But one of the things I want us to see and see very clearly as we discuss this passage today.
is that kingdom values are not just a list of do's and don'ts..
That to be a follower of Jesus, we just have to check the boxes and do the right stuff.
and stay away from the wrong stuff. That's not what kingdom values are about..
We can get that idea sometimes because as humans, we want to make our faith into a ritualistic religion..
It's just a part of human nature. All of us want to do that..
We gain security and a feeling of safety if we feel like we know exactly what we're supposed to do.
and we're doing it. And so we always try to do that..
We see that everywhere in the world in religion..
We see people trying to make it about works and doing good things and checking boxes and keeping the list..
And we just want to do that..
But Jesus is showing us here that that's not what kingdom values are about.

$^{41}$and that's not what a relationship with God is about..
And I want to remind us here that Jesus, as he's speaking in the Sermon on the Mount,.
he's speaking to a very diverse crowd like we talked about when we began this study..
But he's speaking to a lot of people that are very devout..
So Jesus is speaking to religious people, people who are serious about their faith.
and serious about following God..
And many of these people would have been Jewish people who were trying to please God.
through the system of the Old Testament law..
And we know that some of the Pharisees would have been in this crowd.
and the Pharisees had created this works-based religion..
But it was not true to God's faith that was revealed in the Old Testament..
That's one of the interesting things..
If you go back and look at the Old Testament, you see that God, even in the Old Testament,.
was calling people into relationship with him..
The men and women who were heroes of the faith in the Old Testament.
are men and women who walked with God and who had a relationship with him and who trusted him..
You can see that..
Look at their lives and read about their stories..
They are pointed to in the Old Testament because they knew God and had a relationship with him,.
not because they checked all the boxes and kept a list of good works and bad works..
They were obviously righteous people that honored God and loved him and followed him.
and they did good things..
But that was not the reason that Scripture points to their lives in the Old Testament..
So even in the Old Testament, God is calling us into relationship to him..
But by the time Jesus was on the earth, many people misunderstood that..
The Pharisees had twisted the Old Testament law into this works-based religion..
So a lot of Jesus' listeners in the Sermon on the Mount were coming from that paradigm.
of a really works-based religion..
They were very serious, they were very devout, but they were very focused on good works..
Jesus is pointing out to them that being his follower, living in the kingdom of God,.
being a Christian, is not about a works-based religion..
It's about a relationship with God and ultimately with Jesus..
And it comes through the work that Jesus was here on earth to do..
And we're going to talk about that today..
But I want to point us to this truth first..
I want us to consider this, that following Jesus is not simply a matter of correct understanding.
about truth or correct doctrine or correct thinking..
It's not simply a matter of correct understanding..
Now that's important for us to get..
Understanding is important..

$^{81}$We need to have a correct understanding of truth..
In fact, you can't know Jesus without a correct understanding of truth,.
of the truth of who he is and why he came to the earth..
You can't know him..
But here's the thing that we need to understand..
You can have a correct understanding of who Jesus was and of the truth about Jesus.
and still not know Jesus..
And that's really significant..
We need to see that here..
We don't know Jesus simply because we're orthodox in understanding of truth or doctrine..
That's not what brings us into salvation..
That's not what brings us into relationship with God..
And Jesus is pointing us to that here in verse 21..
He says, "Not everyone who says to me, 'Lord, Lord,' will enter the kingdom of heaven,.
but the one who does the will of my Father who is in heaven.".
Jesus is referring to people who are acknowledging that he is Lord..
These are not people who don't believe in Jesus..
These are not people who don't believe that he existed or that he was the Son of God..
These are people who have a correct understanding..
And Jesus is saying there are people who will have a correct understanding..
They will understand that he is Lord..
They will understand that he is God, but they still will not fully believe.
and will not actually know God..
They will understand things correctly, but not know God..
And Jesus is pointing us to that..
Romans 10, 9, which is a really important verse about salvation,.
tells us how we enter the kingdom of God, how we enter into relationship with God..
It says this, "If you confess with your mouth that Jesus is Lord.
and believe in your heart that God raised him from the dead, you will be saved.".
And there are two really important factors in that verse..
We have to understand it..
We have to understand the truth about Jesus, that he's the Lord, that he's the Savior,.
that his death on the cross paid the penalty for our sins..
We have to understand that..
But we have to go beyond simply understanding it,.
simply giving it mental assent and saying, "Yeah, I can assent to the fact that that is true.".
And we have to believe it..
We have to put our trust in that truth and in that reality..
And that's what Jesus is pointing us to here in verse 21..
He's saying there are people who will understand that I'm the Lord..

$^{121}$They will understand who they are before me, but they are not believing it..
They are not putting their trust in it..
And because of that, they don't actually know me..
There's a verse in James 2 where James is talking about our works.
and the relationship between our faith and our works..
And he says this really interesting thing..
He says in verse 19, "You believe that God is one, you do well, even the demons believe and shudder.".
And what's James saying in that verse?.
He's saying, "Hey, if you believe correctly about who God is, that's great..
But that doesn't mean you have a relationship with God just because you understand the correct nature of who God is.".
He says, "Even the demons know who God is..
They tremble. They are afraid of God because they know exactly who He is..
They know exactly how powerful God is..
They have a correct understanding of who God is and who Jesus is as the Son of God.
and of the power of God and of the sovereignty of God in the universe..
They understand that, but they don't have a relationship with God..
Obviously, they are full of evil works and full of terrible things..
And that's a really important thing for us to see..
Because sometimes in the church, we think, "Well, I'm pretty sure I know Jesus..
I know God because I understand some things about God..
I know some of the truth. I know some of the right doctrine..
I have an orthodox understanding of who God is.".
But Jesus is telling us here that it's not just enough to have a right and correct understanding..
That we need to have belief and trust in Jesus, and we need to have relationship with God..
Another way that helps me think about this is if you think about the difference between knowing about someone and actually knowing them..
You can think about someone that you admire, maybe a famous person that you admire in our world..
Maybe it's a celebrity, maybe it's an athlete or a musician or someone who's in the movies..
Or perhaps it's someone who is some famous intellectual or an author that's written books that you love..
But think about someone that you admire in our world..
You may know a lot about that person..
You may know when they were born..
You may know something about their personal life..
Maybe you've read on Wikipedia about them or from other sources..
And so you know something about their marriage or their family or when their birthday is..
You may know all kinds of things about that person..
And if you're in a conversation with someone else, you may be able to tell those things..
Say, "Yeah, this person, they were born here and they grew up here and they spent five years in this place..
This is where they went to college.".
You may even give the impression to someone else in a conversation that you know that person because you have so much information about them..
But that doesn't mean you know them..

$^{161}$That doesn't mean that if they were walking down the sidewalk and you saw them, that you could go up to them and say,.
"Hey, it's great to see you," and that they would have a conversation with you like a friend..
Because you can know all kinds of things about someone but not actually know them and have relationship with them..
And I think that's a helpful illustration of what Jesus is saying to us in this verse..
He's saying it's not enough just to know about God..
Even if you know the right stuff, even if you have a correct understanding of who God is and who Jesus is, that's not enough..
He needs to go beyond that to belief and to trust..
And if you want to be sure that you know Jesus, then you have to do more than just know about God and have a correct understanding..
You need to have your trust in Him..
You need to be believing in Him for salvation..
The next thing I want us to see is that following Jesus is not simply a matter of doing good works..
And this is so important. This is hugely important..
In verse 22, Jesus says, "On that day many will say to me, 'Lord, Lord, did we not prophesy in your name and cast out demons in your name and do many, many works in your name?'.
Then I will declare to them, 'I never knew you; depart from me.'".
The people in verse 22 are people who had done some really good works..
And they said that they'd even done those works in the name of Jesus..
But Jesus says that they don't know Him..
And we have to understand that. We have to wrap our minds around that..
And the church, as followers of Jesus, as people who are trying to live out the values of the kingdom, we have to understand.
that following Jesus is not simply a matter of doing good works..
We don't know Jesus because of the good things that we do..
And that's really important for us to understand..
As I said earlier, it's so easy for us as humans, as a part of our nature, to try to make our faith about some type of system of good works..
And about a checklist and about things that we can see and tangibly do..
And say, "I know I'm right with God. I know it's okay because I can see it..
And I tangibly did this. And I tangibly did that. And that makes me feel good.".
And that's not the basis of our salvation in Christ..
It is not through our works or through our efforts or through our performance or through anything that we do..
That is not how we know God..
And Jesus is pointing us so clearly to that right here in this passage..
And this is important for us to understand, especially in our modern day..
Especially for young people in the world that we live in..
There is so much focus on doing the right thing on morality in our world..
And that's a good thing because we want to be right before God..
We want to do what's right..
As followers of Jesus, that's an important focus for us to have..
But often, the emphasis that our world gives us, especially for younger people,.
is on doing the right thing, advocating for the right causes, volunteering in the right places, being a part of these things..
And that if we do that, then that makes us a good person..
And we can feel like everything is okay..

$^{201}$Now those things are morally therapeutic..
And what I mean by that is they make us feel good morally..
They make us feel like, "Yeah, I am a good person. I volunteered and I helped someone else..
I am a good person. I advocated for this cause.".
Or, "I am a good person because I went to this protest or I was a part of this event.".
Or, "I tried to help in this way.".
Those things are morally therapeutic. They make us feel good..
But they do not mean that we know Jesus..
And they do not get us into heaven..
And we need to understand that..
Jesus is describing people in this verse that have done not just small works, but amazing good works..
And said they were doing them in His name..
These people said they had prophesied..
They had spoken words from God about the world around them in the name of Jesus..
They said they had cast out demons..
They had done supernatural miracles in the name of Jesus..
And many mighty works..
They said, "We have done this big, huge, difficult, hard, supernatural, amazing things in your name..
That is why you should let us into heaven..
That is why you should believe that we know you..
That is why we should be okay with you, Jesus..
Because we have done all this mighty stuff in your name.".
And Jesus says, "I do not know you.".
That is one of the powerful and deeply challenging things about being a follower of Jesus..
And about the ethics of being a Christian..
The ethics of the values of the kingdom..
Is that in Christian ethics, God calls us to do the right thing..
And to do it for the right reason..
And that stands apart from other moral systems in our world..
Other religions and other moral systems say, "As long as you do the right thing, you are okay.".
If you give to the poor, then you did a good work..
You gave to the poor and that is all you need to do..
It is fine if you give to the poor out of a selfish motivation because you are trying to look good in front of other people..
Or if you give to the poor out of a selfish motivation because you want to gain something from it..
That is okay. You did the good work and that is all you need to do..
But our faith, the values of the kingdom say, "No, you have to do the right thing and do it for the right reasons.".
In 1 Corinthians 13, Paul talks about doing good works..
And he talks about the fact that if we are not motivated by the love of Christ and the love of God,.
then those good works are ultimately in vain..
That they do not have the meaning that we think they do..

$^{241}$They do not have the significance that we think they do..
He says that they are just like banging cymbals together..
They are just like making a bunch of noise..
But they are not actually having any deeper significance..
And that is because God is calling us to do the right thing with the right motivation..
God is calling us to do something that is coming out of our relationship to Him..
Not because we are trying to prove our relationship to Him..
Not because we are trying to gain a relationship to Him..
Not because we are trying to please God in some way..
But God is calling us to do good works that are an overflow of our relationship to Him..
Now good works are really important..
Obviously, the values of the kingdom matter..
We have been talking about them for a long time..
So they are very important..
And we want to walk in them and we want to live them..
But we want to walk in them and live them for the right reason..
We do not want to walk in them and live them because we are trying to prove that we are right with God..
And we are trying to earn salvation..
We want to walk in them and live with them because they are an evidence of God's salvation..
They are an evidence of what God is already doing in our life..
And they are an overflow of that change that God is bringing about..
We do good works because we are saved..
Not in order to be saved..
And we need to understand that as followers of Jesus..
That is clear in this verse..
If we think we are doing good things in order to be right with God,.
then we are missing the point and we are not right with God..
But if we are doing good things because we are right with God,.
because Jesus saved us and redeemed us,.
and those good things are an overflow of that relationship,.
that is the place God wants us to be..
And I want to say this..
You cannot know Jesus without doing good works..
The verses we read last week make that clear..
That if you know Jesus, if you are a genuine follower,.
and you genuinely love Him,.
then good works are going to come..
You cannot know Jesus without doing good works..
But you can do good works without knowing Jesus..
And that is really important for us to understand..

$^{281}$Just like the matter of correct understanding..
You cannot know Jesus without a correct understanding of who He is..
But you can have a correct understanding and not know Him..
The same is true in our actions..
If we know Jesus, we are going to please Him with our lives..
It is going to happen..
He is going to work in our hearts and change us,.
and we are going to do things that are pleasing to Him..
But we can do things that look good and not know Him..
And that is what verse 22 is telling us..
So if following Jesus is not just about a correct understanding and doing good works,.
then what is following Jesus about?.
How do we know that we know Jesus?.
Well, the Bible is clear..
It is much more simple than we can imagine..
Following Jesus is simply a matter of trust in the forgiveness and grace of the cross..
Salvation is by grace alone through faith alone..
What I mean by that is it is by grace alone..
We do not deserve it..
We do not earn it in any way..
It is an unmerited gift that God gives to us in His grace and love for us..
And it is by faith alone..
Because we do not grab hold of it by anything that we do..
We simply trust in the sacrifice that Jesus gave on the cross..
We trust in the fact that God loved us..
That Jesus came into our world because of that love for us..
That God wanted a relationship with us..
That Jesus came into our world because of God's desire for a relationship with us..
And that Jesus lived a perfect life, did not sin, and because of His perfection,.
He was able to give His life in our place as our substitute on the cross as a perfect sacrifice..
And that because of what He did in our place as our substitute,.
that all of our sins can be forgiven by God..
That God placed His justice and His wrath on Jesus on the cross..
The stuff that we deserve, Jesus took it in our place..
And because of that, we can experience complete forgiveness,.
complete healing and freedom from sin..
That is the good news of the gospel..
And we believe that by faith..
We do not grab hold of it by our good works or by anything else we do..
We believe it by faith..

$^{321}$We say, "God, I trust You that that is why I have a relationship with You.".
Because of what Jesus did, not because of anything that I do..
We know Ephesians 2, 8 and 9, it talks about it's not our works..
It's not anything that any of us have done..
It's not anything that any of us can brag about or boast about..
It's because of Jesus and His gift of giving His life for us on the cross..
Following Jesus is simply a matter of trust in the forgiveness and grace of the cross..
That's how we come to salvation..
That's how we begin relationship with God..
And if you've never done that, then that is open to you at any time, at any moment..
Even right now, you can begin that relationship with God.
through believing that Jesus is who He says He is..
Having the correct understanding, but then putting your trust in it and saying,.
"God, I believe it's true. I believe that Your forgiveness covers me..
I believe that Your grace covers me.".
It's as simple as that..
But what about our life in Christ?.
After we begin that relationship and we're saying,.
"Okay, I'm a Christian and I'm trying to live out the values of the kingdom..
I'm trying to follow Jesus.".
Well, even then, following Jesus is simply a matter of a relationship to God..
We see that in Galatians 2, 20..
It says, "I have been crucified with Christ..
It is no longer I who live, but Christ who lives in me..
And the life I now live in the flesh, I live by faith in the Son of God.
who loved me and gave Himself for me.".
That verse is clear that even after we know Jesus and we walk in a relationship with Him,.
the life that we live is Him living through us..
It's that relationship with Him..
John 15 talks about abiding in Christ..
Verse 4 says that we abide in Him and as we do, we bear fruit..
So those things that we talked about last week, the fruit, the good works that come out of our life,.
they are important..
But they come from Jesus working through us..
They come from us living in relationship with Him..
Not us trying harder or us figuring out more stuff.
or us coming up with a way to be a better person or us rehabbing ourselves..
They come from God doing work in our life to transform us and to live His life through us..
And the New Testament is clear..
God gives the Holy Spirit to every believer..

$^{361}$When we come into relationship with God, the Holy Spirit comes into our life and is present with us..
And the Holy Spirit teaches us and empowers us to begin to become like Christ.
and to begin to live in the way that Jesus wants us to..
Through that relationship with Him..
Unlike any other relationship in our life, we can't neglect it..
We have to give it time and attention and focus..
That's why I always talk about giving our time to God and time in prayer,.
time in God's Word, time in worship and doing that every day..
That's something that we have to give time to for that relationship to be healthy and flourish..
But as we give ourselves to that and as God begins to work in our life.
and that relationship is healthy, that's how we begin to live out the values of the Kingdom.
and to be obedient to Christ..
Ultimately, following Jesus is simply a matter of surrender to God's work in me and through me..
It's simply a matter of surrendering to God's work in me and through me..
2 Corinthians 5, 17 and 18 says this,.
"Therefore, if anyone is in Christ, he is a new creation..
The old has passed away; behold, the new has come..
All this is from God, who through Christ reconciled us to himself.
and gave us the ministry of reconciliation.".
The work of Jesus after we come to salvation, after we begin relationship with God,.
the work of Jesus is to make a new creation in us.
for the old selfish desires and parts of us to pass away.
and to be replaced with love and service to others..
And that's what God is doing..
So yes, good works, Kingdom values, they're very, very important..
But they come as a result of the work that Jesus does in us..
And we have to surrender to that work and say, "Okay, God, I want you to be at work in my heart.".
And at times it may be difficult..
At times it may be hard to let go of those selfish desires..
At times it may be hard to not just seek what I feel like is best for me..
It may be hard to obey you..
It may be hard to be counter-cultural and go against the grain of what everybody else is doing..
Like we talked about a few weeks ago when we talked about the narrow path..
Those things are not easy, but as we surrender to them,.
God works in our life and makes something new and something beautiful.
and something that is conformed to the image of Christ, we become more like Jesus..
And those good works and those Kingdom values start to flow out of our life..
Galatians chapter 3, the beginning of the chapter, Paul says this,.
"O foolish Galatians, who has bewitched you?.
It was before your eyes that Jesus Christ was publicly portrayed as crucified..

$^{401}$Let me ask you only this, did you receive the Spirit by works of the law or by hearing with faith?.
Are you so foolish, having begun by the Spirit, are you now being perfected by the flesh?".
What Paul is saying in that chapter is saying, "Listen, you were saved through your faith,.
by the grace of God, and now that you're in relationship with God,.
do you think suddenly it's about your works and that you're going to become pleasing to God.
through learning it and through trying hard and through keeping the law?".
Apparently the Galatian church had gotten confused about that,.
and Paul is calling them back and saying, "Look, you got saved by the grace of God.
and through faith in that grace..
And you continue in your relationship with God by the grace of God.
and through surrendering to the work of God in your life, not through your own effort.".
And we need to understand that as well, that we know Jesus through grace.
and we live as followers of Jesus and we live out the kingdom of values through grace.
and through surrendering to the work of God in our life, not through our own effort..
So how do we know that we know Jesus? How do we know those things?.
I would point us to a couple questions here as we close.
and call us to evaluate ourselves like Jesus does in these verses..
It's scary to think about the fact that there are going to be people who think they know God,.
but they don't, and they don't find out until this day that Jesus speaks of..
And we obviously don't want to be in that position..
And so I would encourage us to ask ourselves, "What is our faith or our confidence based on?".
If someone were to ask you, if I were to ask you right now,.
"How do you know that you're going to be in heaven?.
How do you know that you are in relationship with God?.
How do you know that you're in the family of God? How do you know that you belong to Jesus?".
What would your answer be?.
If your answer is something like, "Well, I'm a good person. I've done good things..
I've gone to church a lot. I know a lot about the Bible. I know a lot about truth..
I know a lot about church and I know a lot of Christian songs and I do good stuff.".
If that's what you would answer, then your confidence is in the wrong place..
Your faith is ultimately in yourself and in your understanding and in your works..
And that's a misplaced confidence. That's a misplaced faith..
The only answer we have to how we know that we're in relationship with God.
and how we know that we'll be in heaven someday is because of the grace of Jesus.
and what He's done for us..
And the same question goes for our growth in Christ..
How do we know that we're growing in Christ?.
And how do we know that God is at work in our life?.
Is it based upon, "Well, I'm doing all this good stuff and I'm doing better than other people.
and I feel like I'm a pretty good person and here's the things that I do that are good.".

$^{441}$Hey, all those good works are great..
But if that growth is not based upon the work that God is doing in your heart,.
then it's ultimately in vain like Paul talks about in 1 Corinthians 13..
It's based on the fact that you are trusting in the work of Christ in your life..
You're trusting in God's work to renew you, to change you..
You're being empowered by the Holy Spirit to live out the values of the kingdom.
in a way that you couldn't on your own..
Then that's a solid foundation..
But if your confidence is in yourself for your growth in Christ, then that's misplaced..
If you're just trying to check boxes and keep lists and do all the good things.
so that you can be a good Christian, if you're just trying to keep the law in some way,.
then that's misplaced..
Our salvation is by grace alone, through faith alone, and our growth and our relationship to God.
is through surrender to His work in our life..
And we need to know that and we need to continue to point ourselves to it..
Even if you've known Jesus for a long time,.
you can begin to slowly turn back towards doing good things.
and towards walking in the law and towards trying to keep the checklists..
And we need to constantly be pointed towards the fact that it is only through Christ.
that we know God, it's only through Christ that we are changed and that we grow in knowing Him..
And that's how we know, that we know Jesus..
If we're confident and trusting in the work of Jesus in our heart,.
we're going to have a relationship to Him and for growth.
and for living out the values of the kingdom.
and for the good works that God wants to do in our life..
Let's pray..
God, thank You for Your truth..
Thank You for the way it challenges us and points out our weaknesses.
and points out the wrong things that we're tempted to believe..
And God, I pray that You would help us to be careful that we're not trusting in ourselves.
and that we're not trusting in our relationship to You,.
that we're not trusting in just what we know about You or our understanding of the truth,.
that we're not trusting in the good works that we feel like we've done..
God, help us to make sure that we are trusting in You and You alone.
and that we understand it's nothing that we could do to gain salvation..
It's by Your grace alone and it's through our faith and trust in what You've done alone..
God, we live in relationship with You and as we have a desire to live out the values of the kingdom,.
which is a good and beautiful desire,.
help us not to try to make that happen in our own strength.
and help us not to try to live out good works through our own effort,.

$^{481}$but help us to try to be surrendered to You and to allow You to work in our hearts.
so that You can bring to pass the works that are truly honoring to You.
so that You can bring kingdom values out of our lives,.
so that they literally overflow from our lives into the lives of the people around us, God..
God, may we know that we know You and may we live for You in a way that glorifies You.
and in a way that loves others..
In Jesus' name, Amen..
[Music].
(gentle music).
[BLANK AUDIO].
\newpage



\section{}
\label{sec:BG6FWzPxpgI}
\textbf{神的名字3\_Yahweh Rapha 耶和華醫治(靈思牧禱)}
\newline
\newline
連結: \href{https://youtube.com/watch?v=BG6FWzPxpgI}{\texttt{https://youtube.com/watch?v=BG6FWzPxpgI}} ~~~~ 語音日期: 2024-04-02
\newline
\newline
\hyperref[sec:BAxyx20fKwY]{\small{< < < PREV SERMON < < <}}
~
\hyperref[sec:index_chronic]{\small{[返順時目]}}
~
\hyperref[sec:index_scriptual]{\small{[返順卷目]}}
~
\hyperref[sec:D8rFOM5_n_k]{\small{> > > NEXT SERMON > > >}}
\newline
\newline
$^{1}$喜不喜歡也好.
人生裡面有些事真是.
遇影隨形 揮之不去的.
其中一個就是疾病.
無論是小病.
好像被蚊蟲咬傷.
大病好像末期癌症.
無論是身體的疾病.
或者是心靈的疾病.
更加不要說這個2020年的.
世紀大疫症.
疾病使人辛苦.
以至痛苦 入心入肺.
我們都很希望能夠找到.
大能的醫生將我們醫好.
在這一集.
我就想和你們看一位.
大能的醫生.
不用排期 不用掛號.
我們可以從聖經裡看到.
就是Yahweh Ravah.
中文譯作Yahweh Yizhi.
是神的另外一個名字.
也是祂要向人顯明的.
自己另一面.
要明白神的這個名字.
我們要從聖經出埃及的.
一個小小的故事說起.
那裡記載當時上帝的子民.
剛剛逃離埃及.
遺老的地方 走到馬拉的曠野.
在那裡艱險荒涼.
而他們的水肚裡面.
最後一滴水都喝光.
那個時候他們來到一個水池.
本來是很開心的.
但那裡的水是苦的 喝不到的.
在乾涸的折騰裡面.
他們開始發怨言.
雖然他們表面是埋怨摩西.

$^{41}$但心裡面其實是.
埋怨上帝.
但在極口渴裡面.
埋怨是人之常情.
但這個人之常情.
亦都反映他們對耶和華的恩典.
是多麼善忘.
他們對他的大能.
是多麼沒有信心.
為甚麼這樣說.
因為在出埃及之前.
他們看到上帝怎樣降災.
使出埃及由不可能的東西.
變成有可能.
亦都就在幾天之前.
耶和華將紅海的水分開.
拯救他們逃離法老軍隊的追擊.
現在雖然口渴.
痛苦.
為甚麼不對上帝多一點信心.
等候他的拯救來到.
而竟然這樣埋怨法老.
再讀下文.
我們看到.
耶和華沒有因此撇棄他們.
他叫摩西拿一片樹木.
放在水裡面.
苦水變成可以飲用的甜水.
耶和華的子民.
又一次經歷他拯救的大能.
微妙的就是.
在這個時候.
上帝向他們曉喻自己的一個名字.
就是.
Yawe Rava.
耶和華伊哲.
大聖聽相當是這樣說的.
他說.
若留意聽耶和華.
你神的話有行我眼中為證的事.

$^{81}$留心聽我的誡命.
守我一切的律例.
我就不將所加與.
埃及人的疾病.
加在你們身上.
我耶和華是.
醫治你的.
這句話是甚麼意思呢?.
我們可以看到.
兩個重點.
第一.
耶和華既願意.
亦有大能.
去醫治.
第二.
耶和華非常重視.
他和他的子民立約的關係.
請你注意.
上帝在這裡不是說我是醫治的.
而是說.
我耶和華是醫治的.
耶和華正是上帝和他的子民立約的名字.
他這樣說.
是要他的子民不要忘記.
要求他們忠於立約的關係.
就是忠誠.
遵守他的誡命律例.
行他眼中為證的事.
事實上.
在整本聖經無論是舊約或新約.
都記載了.
無數耶和華醫治的例子.
包括了新心靈的.
傷痛和疾病.
包括了個人的.
亦包括了整個國家民族的.
然而.
耶和華的醫治經常都是強調.
立約的關係和要求.
事實亦都是.

$^{121}$耶和華的醫治.
往往是出人意外.
超過人所求所想.
有時我們身體的疾病.
祈求醫治.
祈求醫治我們的心靈.
有時我們為個人的疾病求醫治.
或是整個家庭或群體.
但無論如何.
只要我們忠誠於立約的關係.
我們每次呼求.
都會得到他大能的醫治.
雖然他的醫治方法未必是.
依照我們所想的.
「也為拉法」.
「忠誠」這個名字.
他告訴我們.
他明白我們在疾病裡的痛苦和掙扎.
無論是在身.
心或靈上.
「也為拉法」.
「忠誠」這個名字.
他告訴我們.
他有能力去醫治.
「也為拉法」.
「忠誠」這個名字.
他提醒我們.
他重視和我們立約的關係.
更震撼的就是.
為了這個因罪而受重病的世界.
神親自來到.
給我們一個永遠的醫治.
「也為拉法」.
「忠誠」這個名字.
他提醒我們.
他有能力去醫治.
「也為拉法」.
「忠誠」這個名字.
他提醒我們.
他有能力去醫治.

$^{161}$「也為拉法」.
「忠誠」這個名字.
他提醒我們.
他有能力去醫治.
「也為拉法」.
「忠誠」這個名字.
就是神要拆去他們一切的眼淚.
不再有死亡.
也不再有悲哀.
哭號,疼痛.
因為以前的事都過去了.
現在讓我們一起.
在「也為拉法」面前祈禱.
「也為拉法」.
我們要讚美你.
因為你有絕對的能力.
和無邊的大愛.
醫治好一切.
無論是身心靈上的疾病.
整個世界的殘缺.
而你的醫治.
不單是短暫.
更加是永遠的.
這對於今天在疫情和病患中.
因著世界各種敗壞而受傷的我們.
是莫大的安慰和激勵.
「也為拉法」.
我們求你赦免.
因為作為你立約的子民.
我們經常都忘記.
甚至背棄這個神聖而寶貴的約.
我們沒有遵守你的律例.
無論是在安息或危難的時候.
我們都會遠離你.
甚至建立很多偶像.
羞辱你的命.
對於這些過犯我們沒有藉口掩蓋.
只有因著耶穌基督的命.
我們求你赦免.
求你醫治.

$^{201}$讓我們在這個時候誠心對你說.
我們要立志迴轉.
求聖靈幫助我們.
兼顧我們對你的忠誠.
我們的禱告是奉耶穌基督的聖名.
阿們.
拜拜!.
\newpage



\section{}
\label{sec:D8rFOM5_n_k}
\textbf{聖地網上行 - 步主腳蹤 [第十二集] 出賣與受審}
\newline
\newline
連結: \href{https://youtube.com/watch?v=D8rFOM5_n-k}{\texttt{https://youtube.com/watch?v=D8rFOM5\_n-k}} ~~~~ 語音日期: 2024-04-05
\newline
\newline
\hyperref[sec:BG6FWzPxpgI]{\small{< < < PREV SERMON < < <}}
~
\hyperref[sec:index_chronic]{\small{[返順時目]}}
~
\hyperref[sec:index_scriptual]{\small{[返順卷目]}}
~
\hyperref[sec:DKQrjJlzvbM]{\small{> > > NEXT SERMON > > >}}
\newline
\newline
$^{1}$耶穌和他的門徒吃完最後晚餐後.
來到橄欖山的山腳.
一個地方叫做赫西瑪利縣.
這是亞蘭文.
亞蘭文有兩個字組成.
赫西瑪利是Gat-Chamonix.
Gat是油炸的意思.
西瑪利後面那一段是橄欖油的意思.
這裡是橄欖山山腳將橄欖壓碎成為油的地方.
似乎象徵耶穌要面對死亡時.
他要面對的壓力和掙扎.
在經文中告訴我們.
他在赫西瑪利縣的禱告.
三本福音書都有記載.
路加福音甚至記載什麼呢.
他說耶穌極其相同.
禱告更加懇切.
看諸如大血點滴在地上.
其實這是有醫學解釋的.
有些醫學解釋說.
面對很大壓力的時候.
一個人可能會極度的缺水.
然後表面的表皮的血管就會爆裂.
然後流出來的汗就像流血一樣.
為什麼會在這個情況下.
他突然間好像有很大的掙扎.
突然間很驚慌.
第一我們可以想起的就是這裡的地理.
因為鑒嵐山較北的山腳.
那裡比較肥沃一點.
那裡可以種鑒嵐樹.
但是鑒嵐山的山脈.
比較藍的山頭.
還有山腳底下的吸塵谷.
地質一直都是這麼差的.
不是很肥沃的.
所以歷世歷代都是墳墓的地方.
等於有人說我們不怕死.
但如果你去到墳場的時候.
你就未必會這樣說.

$^{41}$另外我們可以看看地理.
就知道那個掙扎有多大.
因為在地理上鑒嵐山其實就是耶路撒冷.
東面的出入口.
要逃亡的話.
上山十分鐘就可以跨過鑒嵐山.
然後去到猶大講義.
耶穌可以一走了之.
但祂願意去馴服天父.
我們來到鑒嵐山的山腳.
赫西瑪利這個地方.
按照當時猶太人的預節慣例.
預節晚餐完結後.
他們會留在耶路撒冷城的範圍內.
耶穌和門徒可能來到這個山洞.
門徒就在山洞內睡覺.
然後耶穌就帶著他們三個門徒.
去到不太遠的地方.
在那裡守望祂的禱告.
在另一邊我們看到一間大教堂.
天主教教堂叫萬國教堂.
Church of All Nations.
也有另一個名叫壽府大教堂.
Basilica of Agony.
外面是一個很漂亮的花園.
有橄欖樹的樹園.
這些橄欖樹經過化驗其實歷史悠久.
有些是十字軍時期種植.
我們走進萬國教堂.
教堂的設計是很幽暗的.
其實是想表達赫西瑪利園那種黑夜的掙扎.
在赫西瑪利園當中我們可能會想起.
第一件事就是神未必一定會按照我們祈禱的心願去答應.
因為天父沒有按照祂的兒子起初祈禱的心願.
沒有把祂的墓碑拿走.
相反的,在祂掙扎的祈禱當中.
賜給祂有勇氣,有力量去面對十字架.
順福臣至死.
另外耶穌被捕.
到祂死而復活的期間.

$^{81}$有很多的事跡,很多的細節.
反映出聖經是想教會來聽清楚.
這些故事.
中間也有很多人物.
第一個我們可以提一下的人就是加勒人猶太.
這個賣耶穌的猶太.
我們知道在晚餐結束之前已經離開了.
他找大祭司.
出賣耶穌的價錢是三十塊錢.
如果計算的話大約是今天美金五千元左右.
是當時的四五個月工資.
或者是當代的墓碑的價錢.
當然我們不知道真正的動機是怎樣.
但可能就是警惕我們.
會不會就是為了眼前的利益.
我們出賣信仰,背棄神.
聖經的作者特別將猶太.
和巴大利的瑪利亞作一個對立.
瑪利亞大家可能聽過這個故事.
就是耶穌在進耶路撒冷之前.
他拿了一斤極貴的珍拿達香膏來抹耶穌的腳.
這個珍拿達香膏是當時很名貴的香水.
可能是瑪利亞一生的嫁妝.
或者是她一生的經濟倚靠.
她可以傾倒在耶穌面前.
因為她看到耶穌的死.
是帶來一些比她現世一切更加寶貴.
是一種無價之寶.
這種角色的對立.
告訴我們究竟救恩在我們心中.
是可以有價還是無價.
我們追溯一下耶穌被捕之後的行蹤.
經書告訴我們.
人們把耶穌帶到大祭司的宅裡.
這裡說到當時的大祭司.
哥爾法和他的岳父亞拿的家.
他在那裡被拷問.
彼得當然跟著想知道發生什麼事.
之前的基督教傳統會說.
這個大祭司哥爾法的家.

$^{121}$就是現在的西山.
這裡有一個一世紀的古樓房.
地庫裡有很多山洞甚至密室.
相信耶穌在這裡被羈留和被拷問.
這裡有一座天主教教堂.
叫做雞鳴堂.
Church of St. Peter in Calicutto.
意思是雞啼.
這裡當然是說到彼得三子不認主的事蹟.
你記得耶穌對彼得說.
他警醒他.
雞鳴雞啼之前你要三次不認.
耶穌對彼得說這句話的時候.
是黑夜.
意思是單憑血肉之軀.
你的信心不會捱過一個晚上.
當然我們知道彼得最後是三次不認主.
但經過一些考古的發現.
我們發現原來大祭司等等的家族.
他們應該在聖殿對面的上城.
上一集我們去過一個地庫的博物館.
叫做World Archaeological Museum.
那裡看到一個五千呎的豪宅.
裡面有很多的潔淨寺.
其實牆上也找到一個雕刻.
這個雕刻是一個聖殿的金燈台.
那些雕刻的圖像很仔細地刻畫到金燈台的一些細節.
換句話說,刻這個金燈台的人.
很清楚和親眼見過聖所裡面的金燈台.
這個豪宅裡面除了很多潔淨寺.
是祭司專用,有金燈台.
也很吻合當時文獻所說的大祭司的家族.
貪財,積財,暴力.
甚至控制聖殿的經濟等等.
這裡也有一個很大的羅馬式的大庭園被發掘出來.
面對著對面的聖殿.
這裡較有可能是大祭司家族的庭園.
耶穌被審問的地方.
耶穌的罪狀之一是什麼呢?.
就是這個人曾經說.

$^{161}$我能夠拆毀聖殿,三天之後建造起來.
對於大祭司來說,更加是挑戰聖殿的權威.
然後耶穌被帶到公會受審判.
路加福音第22章.
天一亮,民間的眾長老和祭司長.
帶著民事都聚會,開公會.
他問祂是不是基督.
耶穌回答他,從今以後,人只要在上帝權能的右邊.
他們說,耶穌不需要再作見證.
他犯了褻瀆,自稱為神.
於是他們把耶穌帶到羅馬的總督彼拉多面前.
這裡是路加福音第23章.
眾人都起來把耶穌帶到彼拉多面前.
告他什麼呢?.
告他誘惑國民,禁止納粹給亥薩.
又說自己是基督,是王.
所以之前告他的是宗教的罪,褻瀆.
現在告他是政治上的罪,顛覆羅馬帝國的罪.
因為說他自己是王.
彼拉多本身駐守在海薩利亞的海港.
但每逢大時大節就會來到耶路撒冷.
看守人民,還會做巡迴的法官.
在不同的地方聽證.
彼拉多剛剛在雨節的時候來到耶路撒冷.
他住的地方一定是大希律之前所建造的皇宮.
這是在城西的一個皇宮.
大約2000年的時候.
考官人員在已經在博物館的地底下發掘的時候.
發現了希律皇宮的遺跡.
這不是驚喜,因為我們文獻上一早知道.
皇宮就在這裡.
這裡就是彼拉多去聽證,去審問耶穌的地方.
之後老加方案23章有說.
彼拉多聽見耶穌原來是加利利人.
又知道耶穌是屬於希律王所管的.
所以就把他解去希律王的那裡.
希律是其中一個大希律的兒子.
他是守加利利省的.
你記得他是斬了斯大約翰.
守則的時候當然他就是來到耶路撒冷.

$^{201}$而他住的皇宮應該是按照文獻.
就是之前馬加比時代所建的皇宮.
都是在上城的樹下.
希律王氣憤耶穌一番的時候.
又把他解回彼拉多的樹.
老加方案第23章繼續.
彼拉多願意去釋放耶穌.
又勸解他們,即是猶太人.
慕內地猶太的貴族就喊著說.
釘他十字架.
彼拉多第三次對他們說.
為什麼呢?這個人做了什麼惡事.
我查不到他有什麼該死的罪.
所以我要責打他,希望釋放他.
但是這些猶太的貴族.
大祭司等等.
他就催迫彼拉多.
求他把耶穌釘在十字架上.
於是最後他們的聲音就得了聖.
剛才我們提到耶穌被捕到死而復活.
中間的過程有很多事跡和很多人物.
最重要的是有兩個人物一定要看.
第一個就是彼拉多.
彼拉多是猶太地的總督.
他是替羅馬帝國鎮守這個地方.
也是直接授權於羅馬帝王.
我們知道彼拉多在這裡做了總督有十年.
公元後26至36年.
在猶太人的世界裡.
他就是最大權力.
並且我們知道彼拉多的生性是很凶殘的.
他是很仇視猶太人的.
貪污不義,未審先殺.
很多次是和猶太人特意作對.
把偶像帶入耶路撒冷的聖殿.
激怒猶太人.
甚至濫用聖殿裡的錢去做引水道.
但是你看福音書審耶穌的過程.
彼拉多不像是這麼仇視猶太人.
馬太方第17章.

$^{241}$彼拉多見無濟於事,反要生亂.
他就拿盆水出來在猶太人面前洗手.
這個洗手一幕就很有名.
彼拉多有些猶豫不決.
不想處死耶穌.
他就想洗脫自己的罪.
為什麼彼拉多這麼猶豫呢?.
和他一向的作風很不同.
如何解釋他猶豫不決的行動呢?.
在這個時候我們可以知道背後發生什麼事.
在羅馬朝廷有一個高官.
他是羅馬帝王提比留的第二把交椅.
這個人叫做聖詹姆斯.
他和彼拉多有很密切的關係.
誰知道他被懷疑謀朝篡位.
提比留決定立刻處決這個塞雅路斯.
並且我們知道他被處決的那一天是哪一天.
就是公元後31年10月19日.
他被處決後,他的同黨也全部被處決.
這就解釋了為什麼彼拉多在這個時候這麼害怕.
突然間他對猶太人的仇視和作風改變.
不是因為他同情耶穌.
而是他在這個時候想自保.
約翰福音19章12節這樣記載.
有了這個背景,這裡的記載就很寫實.
他說從此彼拉多想要釋放耶穌.
無奈猶太人哭著說.
你若釋放這個人,就不是凱撒的中神.
凡以自己為王的,就是背叛凱撒了.
似乎猶太人貴族大祭司等等都知道.
到底羅馬帝國在羅馬京城發生什麼事.
似乎捉到彼拉多有些痛腳.
萬一民間生亂,他自己就人頭不保.
鮑爾根博士說得很對,他說這是人在江湖.
在政界裡面的遊戲,彼拉多要討好他的老闆提比林.
結果就是不問真理,不問公義.
提到彼拉多的時候,我們要提一提一個很出名的考古石碑.
後人叫他做彼拉多石碑,Pilot inscription,Pilot stone.
發現的地方在凱撒尼亞海港,時間是1961年.
當時發掘的時候,人們在羅馬劇場的石級其中一塊.

$^{281}$找到這個石碑,上面找到一些刻文.
換句話說,在後人去裝修羅馬劇場的時候.
拿了這塊紀念碑,當成石級其中一塊.
雖然有一些文字已經拆去了.
但是大概知道說什麼,中文翻譯大致上是.
獻給神聖阿古斯督,此提比林大殿Tiberium.
本屌彼拉多,猶大總督,Prefect of Judea.
獻此殿,這個建築物.
可能是一個神廟的大殿來Honor彼拉多的老闆.
提比林大殿Tiberius.
發現的時候很轟動,因為這塊石是唯一一塊.
在文獻以外找到的考古實物,證明彼拉多的歷史性.
但如果看清楚內容,其實不單單是說.
彼拉多存在的證明這麼簡單.
這裡說得很清楚,是彼拉多和羅馬帝王之間的政治關係.
Patron-Client的關係.
Patron的主堡當然是羅馬皇帝比林.
Client,即能夠拿到一官半職的,就是彼拉多.
所以彼拉多就討好他的主人.
建這個廟或建築物來尊榮他.
你看到這是一個政治的關係.
你是一個政客,你就要討好你的主堡.
你心裡面有什麼原則,這不是最重要.
最重要是主人說了算,誰是他的主,提比林.
如果你這樣看石碑,這樣看Patron-Client的關係.
看經文可能會有不同的味道.
又是第十八章,這裡提到彼拉多審問耶穌的時候.
彼拉多對耶穌說,這樣你是皇嗎?.
耶穌回答說,你說我是皇,我為此而生,也為此來到世間.
特為給真理作見證.
凡屬真理的人就聽我的話.
彼拉多回答,真理是甚麼?.
當然他的語氣,肯定不是說我現在想尋道.
想知道真理是甚麼,耶穌你可不可以告訴我.
這裡的語氣,大概彼拉多所說的是.
真理,公理可以當飯吃,是否可以保住我的位置.
我就是人在江湖,身不由己.
你想起歷世歷代有很多這個故事.
重症的人出賣自己,出賣公理.
為了眼前的利益,說甚麼真理,說甚麼公理.

$^{321}$這裡提醒我們,到底我們是否屬於耶穌的人.
耶穌是否我們唯一的主,唯一討好的主宰.
另一個與彼拉多相關的人物,就是大祭司戈爾法.
戈爾法是一個家族名.
這個大祭司的全名是約瑟戈爾法,Joseph Kyrvis.
我們對他的出身所知不多.
只知道後期他娶了一個大祭司叫做阿娜.
她的女兒,然後成為他的女婿新年郎.
進入了大祭司的家族.
在當時的處境,大祭司是由羅馬總督,好像彼拉多那些直接委任.
而且大祭司的店舖也是由羅馬總督管理.
放在總督府裡面,大時大節才讓他穿著.
大祭司的工作,就像戈爾法所說.
就是在討好羅馬總督的大前提下.
維持聖殿和維持民間安定.
一世紀的大祭司,直至公元後70年聖殿被毀.
總共有19個,你想像一下,換了19次.
很多大祭司都是早一兩年之後就被革職.
戈爾法在公元後18年至36年,總共做了18年大祭司.
他是第一世紀任期最長的大祭司.
所以從這個角度看,他是最成功的.
但是不是因為他廉潔?.
很可能是因為他只懂得玩政治遊戲.
剛才我們提過彼拉多濫用聖殿的公款去建引水道.
約瑟夫的記載裡面,都沒有提過大祭司有抗議.
所以你想到羅馬總督和猶太人的大祭司中間有多少利益.
中間有多少不見光的東西,大家可以想像.
1990年11月,當時工人在耶路撒冷那邊建公園.
突然間碰到一世紀的墳墓.
一打開墳墓,叫考古學家到場時.
發現裡面有12個骨盒,叫做Ossuary.
在一世紀的時候,猶太人特別的葬禮形式.
人死後將屍身放在墳墓裡面.
家人一年後會回到墳墓撿骨.
將骨放在骨盒裡面.
這12個骨盒裡面有兩個刻上文字.
其中一個骨盒裡面找到的骨.
有一個女人的骨頭,兩個小孩的骨頭.
兩個嬰孩的骨頭,還有一個大約60歲的老人家的骨頭.
如果是的話,很有可能是戈亞發的骸骨.

$^{361}$骨盒外刻上的銘文是Joseph Bar-Kyava.
即約瑟戈亞發之子.
在福音書記載,大祭司戈亞發決心剷除耶穌.
很明顯不是為了公義.
約翰福音第十一章.
祭司長和法利賽人聚集在公會時.
他說這個人耶穌,行了很多神蹟,我們怎樣辦呢?.
如果由他的話,人人都信耶穌的話.
羅馬人就會來奪我們的土地和百姓.
然後內中有一個人,就是大祭司戈亞發.
他說你們不知道,獨不想一個人替百姓死.
免得通國滅亡,就是你們的益處.
他一個人死,比我們這麼多人一起被羅馬人滅絕好.
你記得我剛才提過在聖堂中找到大祭司家族的豪宅.
上面有七燈台,畫在牆上.
如果你這樣看,你就更加知道.
大祭司生存的目的就是要維穩.
要保住聖殿,其實是為己.
在以色列博物館,有張戈亞發的骨盒.
很多時候放在彼拉多石碑上.
一個彼拉多,就是用政權.
忘卻真理,不理公理,出賣自己,出賣良知.
另一個戈亞發,就是以宗教為名.
去鞏固自己的勢力,為了自己有飯吃.
兩個人很巧合,也是同命相憐.
在公元後三十六年,敘利亞省長把這兩個人都革職.
這就是約翰福音所說的世界.
世界不是物質碰到的世界.
在約翰福音中,世界是撒旦掌握的秩序.
這就是世間的遊戲.
猶大衛了一些利益,出賣自己,出賣主.
彼拉多,戈亞發為了自己的利益.
出賣良心,出賣公理,出賣宗教.
人不為己,天主地滅.
人為了自保,沒有辦法.
在這樣的網絡中,你能去到多盡?.
為了自己,可以埋沒良知,埋沒公理.
甚至出賣真正的宗教,可以去到多盡?.
在今天的新聞中,在今天的政治新聞,宗教新聞中.
你可以想像到有多少人受到世界秩序,撒旦所控制.

$^{401}$耶穌基督卻來到這個世界拯救我們.
而且耶穌的生命是完全不同的.
祂甘願信服真正的主宰,父神.
甘願受苦受死.
我們看完這些所謂奸人的故事,我們有什麼反應呢?.
第一個反應可能你會想,是一個很大的警惕.
這是一個很好的反應.
問清楚我們真正的主是誰.
是誰在操縱我們生活,行為取態.
是世界的王,是世界的秩序,背後的撒旦,還是上帝.
第二個反應可能我們會想,我們真的很無助.
因為苦心自問,如果我們在那種場景下.
我們可能每個人都是猶大,彼拉多,戈亞凡.
這可能就是聖經福音書的目的.
就是要我們知道在罪海裡面我們真是無能.
我們是一班無能的罪人.
這樣我們才會明白救恩的寶貴.
羅馬書第五章,這裡說到福音的時候.
保羅說原來我們還軟弱的時候.
基督按照所定的時候為不敬虔的人死.
為異人死有少的,為好人死的或者有人夠膽做.
但是基督卻是在我們還作罪人的時候為我們死.
神應該就是這樣去顯明.
這裡說到軟弱,powerless,我們確實無能自救.
好像我們患上了絕症,疫病的時候.
沒有人救我們,唯有他人能夠救我們.
唯獨是耶穌基督才能夠救我們.
幫助我們,讓我們能夠脫離世界的王的擺佈.
讓我們以主為我們唯一的主.
求主幫助我們,讓我們生活為人的取態.
不要忘記主,不要忘記真理,公理,公義.
讓我們在我們生活中,忠心地去跟隨耶穌基督.
我們一起低頭祈禱.
每一個我們誤會的罪人.
我們無能的底下,你都去拯救我們.
你顯出你的大能,顯出你的恩典.
求主幫助我們,讓我們知道你在十字架上為我們所做的.
是多麼寶貴.
讓我們知道我們單單是屬於你.
單單聽命於你.

$^{441}$幫助我們,忠心地去跟隨.
祈禱逢轉命給我們,Amen.
\newpage



\section{}
\label{sec:DKQrjJlzvbM}
\textbf{舊約網上行 - 上帝的足跡: EP18 揀選}
\newline
\newline
連結: \href{https://youtube.com/watch?v=DKQrjJlzvbM}{\texttt{https://youtube.com/watch?v=DKQrjJlzvbM}} ~~~~ 語音日期: 2024-04-05
\newline
\newline
\hyperref[sec:D8rFOM5_n_k]{\small{< < < PREV SERMON < < <}}
~
\hyperref[sec:index_chronic]{\small{[返順時目]}}
~
\hyperref[sec:index_scriptual]{\small{[返順卷目]}}
~
\hyperref[sec:DbCFhByjju0]{\small{> > > NEXT SERMON > > >}}
\newline
\newline
$^{1}$來到九日網上第十八講.
創世記二十五章.
是亞伯拉罕故事最後的一段結語.
去交代亞伯拉罕的眾子孫及他們的後人.
亦交代亞伯拉罕將財產分給二子.
二十五章頭三段.
第一至第四節.
就是基土拉所生的東方國族.
基土拉是亞伯拉罕另一個妻子.
為了她生不同的兒子.
然後交代了她每一族的族人及名字.
五至十一節.
就是講到撒拉為亞伯拉罕所生的二子.
這個當然就是將來以色列國的族長.
亦是神跡出生的.
是特夢神揀選.
第三段就是十二至十八節.
就是夏格.
這個很熟悉的名字.
亞伯拉罕另一個妻子.
為亞伯拉罕所生的二十瑪利.
然後講到她的後人二十瑪利人.
我們先看第三段.
記得這個很熟悉的名字二十瑪利.
亞伯拉罕和這位妻子夏格所生.
二十五章第十二節.
撒拉的侍女埃及人夏格.
及亞伯拉罕所生的二子是二十瑪利.
二十瑪利的二子們的名字.
按照他的家譜記載在下面.
然後講出這些名字之後.
他說照著他們的村莊營寨.
就作十二族的族長.
二十瑪利享壽一百三十七歲.
氣絕而死歸到他的列祖那裡.
這裡交代二十瑪利人.
住在以色列地的東面約旦河東的國族.
第一段講到亞伯拉罕另一個妻子.
叫做基圖拉.
二十五章第一節這樣交代.

$^{41}$亞伯拉罕又娶了一妻.
其實是原意.
其他的經文告訴我們.
這不是正室的妻子.
是一個妾侍.
這個人叫做基圖拉.
基圖拉為亞伯拉罕生了不同的兒子.
講了六位.
這個基圖拉另一個妻妾.
你記得之前雖然沒有講過.
但是情節你可以明白.
為什麼會是另一個妾侍.
你記得之前提過撒拉不能夠生育的時候.
他就把埃及人的侍女夏格給了亞伯拉罕.
以至他能夠生一個兒子.
基圖拉是另一個妾侍.
應該不是撒拉死了之後才娶的妾.
基圖拉可以跟之前提過的夏格地位相似.
只不過之前講到亞伯拉罕的故事的時候.
沒有交代基圖拉這個人.
無論如何這裡講到基圖拉的後代.
也講到夏格所生的伊什瑪利的後代.
他們這些族人.
按照這些名字,地理等等.
都是分布在約旦河東的一帶.
沙漠和曠野的一些游牧民族.
例如這裡提到一個名字.
或者提到一個地方.
一個族名大家都很熟悉.
就是米殿.
米殿就在沙地阿拉伯半島的北面一帶.
為什麼提到這些別國的族類呢.
25章第五節.
亞伯拉罕將一切所有的都給了伊什.
就是他真正的正室.
撒拉所生的伊什.
然後第六節.
亞伯拉罕就將財物分給輸出的種子.
就是非正妻.
妾侍所生的種子.

$^{81}$趁著自己還在世的時候.
就打發他們離開他的兒子伊什.
往東方去.
亞伯拉罕一生的年齡是175歲.
亞伯拉罕受膏年邁.
氣絕已死.
歸到他的列祖那裡.
近東的社會裡面.
每逢一個家父死之前.
他都會分佈他的家產.
提到亞伯拉罕其他的後代.
他們是什麼族名.
以及將來都會成為邦國.
成為不同的民族.
為什麼要交代這件事呢.
就是完全因為上帝曾經對亞伯拉罕.
應許過一些的承諾.
在十七章第四至第六節曾經提過.
亞伯拉罕得到上帝的應許.
他會成為多國之父.
Father of Nations.
Father of Many Nations.
其實亞伯拉罕這個字.
意思就是Father of Many Nations.
在他身上會出現很多不同的邦國.
不同的民族.
來到亞伯拉罕故事終結的時候.
作者就要交代.
神承諾的每一句話.
每一句都沒有落空.
所以說到這些別國的時候.
說到他們都得到上帝的祝福.
這個其實都是創世記.
由始至終都有的特別的神學.
這種神學就是告訴我們.
很多時候以色列周邊的邦國.
和族長亞伯拉罕有所謂的遠親的關係.
有些親戚的連帶關係.
所以不單止以色列得到上帝的祝福.
周邊的邦國.

$^{121}$特別是和亞伯拉罕有關係的邦國.
都得到上帝自己的承諾.
創世記你記得在第十章的時候.
有七十個列國的列表.
這些邦國都得到上帝自己的祝福.
能夠分散遍滿全地.
你也記得阿們人和摩訶人.
這兩個民族和邦國.
也是和亞伯拉罕藉著諾德有所關係.
所以他們一般來說.
都得到上帝給他們的一般的祝福.
現在說到以實瑪利.
說到基督拉的後代.
東方的國族.
第一段和第三段所交代的民族.
都是說到他們得到神的普遍祝福.
你也記得創世記之後.
一些不同的民族.
也是和亞伯拉罕以實瑪國有關聯的.
譬如以素,以東.
也是和亞國有所關聯.
是兄弟.
這些別國得到祝福.
傳恩合位是神對萬國.
有一個普遍的恩典.
一個common grace.
以色列雖然被揀選.
在萬國當中.
在他們中間成為chosen people.
他們是不能夠壟斷上帝的普遍祝福.
因為這位上帝是創天造地.
是萬國的主宰.
他對每一個邦國都有一個普遍的恩典.
我們繼續問.
神揀選以色列民.
有什麼特別的地方呢?.
我們再看這25章的分段.
第一段一至第四節.
是說到基督拉所生的東方國族.
五至第十一節.

$^{161}$是說到撒拉所生的以撒.
將來成為以色列民.
以色列國.
然後十二至十八節.
第三段.
就是夏革所生的伊什瑪利人.
你看到這個分佈.
三個民族表面上都得到上帝的祝福.
但唯獨中間有分別.
有什麼不同的地方呢?.
只有以撒才是神跡出生.
只有以撒才是上帝到選選.
是他超越人類的不能.
不能生育的撒拉.
居然神跡地能夠生到一個兒子.
就是以撒.
只有以撒才有上帝親自的創造.
親自的恩典.
獨特的恩惠.
能夠成為一個邦國.
將來成為十二支派.
成為大國.
有什麼特殊的目的呢?.
我們又回到創世記第十二章.
亞伯拉罕的起始.
那裡說到.
選擇亞伯拉罕.
上帝說你要成為祝福.
以至萬國能夠成為祝福.
受到上帝的選擇.
亞伯拉罕,以撒,雅各,以色列.
到底是特權多一些?.
還是使命多一些?.
到底是一個優待.
還是一個任務?.
你們想一想這個問題.
其他的民族.
這段聖經告訴我們.
都得到上帝自己的普遍恩典.
Common Grace.

$^{201}$但是只有以色列民.
只有以撒,雅各,十二支派.
以色列國才能夠得到上帝的揀選.
藉著他們.
來播放出上帝的Saving Grace.
能夠傳出.
能夠成為流通的管子.
流出上帝的救贖的恩典.
這個特殊的恩典.
只有以撒.
他的後代.
才是神特定揀選的.
所以被上帝揀選.
不是一個優待.
而是一個任務.
上帝揀選神的子民.
往往就是一個使命.
要使萬國能夠得福.
使萬國能夠認識耶和華.
在萬國當中.
只有神的子民.
只有以色列國.
只有教會.
藉著上帝的揀選.
能夠將祝福帶給其他人.
到底我們有沒有這樣做?.
被揀選當然就是為了一個使命.
使萬國能夠得福.
能夠認識耶和華真神.
但是被揀選有什麼原因呢?.
這個家庭.
或者這個家族.
這個邦國.
到底是不是有特別的原因.
使上帝揀選他們呢?.
或者我們問一問.
到底在作者的筆下.
如何刻畫亞伯拉罕的家庭.
如果你留意上文.
留意以往亞伯拉罕的故事.

$^{241}$表述的時候.
你其實都可以回答這個問題.
就是亞伯拉罕這個家庭.
其實是一個破碎.
不完美的家庭.
至少在《創世記》第二章尾.
你要記住.
人要離開父母.
一夫一妻的理想地下.
亞伯拉罕的家庭絕對不是完美.
怎麼說都談不上.
是一個達到上帝標準的家族.
一夫又多妻.
而《創世記》的敘述裡面.
很明顯地告訴我們.
多妻是墮落了的人類的實況.
神容許容忍.
但不代表它是完美.
之後我們也看到族長的故事.
不單止亞伯拉罕的故事當中.
讓我們看到這個家庭.
很多時候有糾紛.
有很多破碎的一面.
之後以撒和利伯加作為夫婦.
在家庭裡裡製造了很多.
想拆散家庭的張力和誤會.
之後生了雙胞胎.
就是以素和雅各.
當然以素是短視的.
只是為了眼前的慾望而出賣長子的名分.
雅各呢?.
神選雅各是否代表雅各特別好呢?.
當然不是.
因為他前半生都是行鬼詐.
去搶奪上帝的祝福.
或者用欺騙的手法來得到上帝的福氣.
再繼續讀下去的時候.
雅各的十二眾子大家更加清楚.
彼此出賣彼此的嫉妒.
差點連親兄弟的弱勢.

$^{281}$他的命都保不住.
如果不是上帝的恩典的緣故.
這個家庭早就亡了.
但是人呢?.
你會發現三番四次的威脅.
去刺然了上帝的計劃.
刺然了上帝的工作.
如果這樣看的話.
創世記從來都沒有掩飾.
這些所謂族長.
以色列國的英雄.
這些創始人的醜惡的一面.
或者他們的醜聞.
很實在地將這些人性真實.
醜惡破碎的一面呈現出來.
這個其實就是整個阿巴罕故事.
結中的時候給了我們一些線索.
告訴我們作者到底想表達一種什麼樣的神話.
最近橫掃全球的一個韓國的影集.
叫做《魷魚遊戲》.
不知道你們有沒有看過.
我已經看完了.
你不妨去看看.
到底這個劇情想表達的是什麼.
用一個遊戲的方式.
看到這些參賽者為了贏.
為了利益.
為了生存.
他們可以不擇手段.
很多方面這些情節.
想去反映人性黑暗的一面.
人性的真實和醜陋表露無遺.
其實創世記也有這種意義.
地上的蠻族.
你問問哪一個國家.
自古到今.
永遠都是人不為己 天誅地滅.
每一個邦國 每一個民族.
都是為了自己的生存 為了自己的利益.
而去建國.

$^{321}$而去將自己的勢力擴張.
無論你說巴比倫帝國 埃及帝國.
赫人 以東 摩押 阿們.
以實 馬利 東方國族等等.
萬國萬族在地上.
沒錯得到上帝的普遍的恩典.
但他們存在的目的.
每一個邦國都是為了自己.
所謂人不為己.
但在神的計劃當中.
神卻要將以色列國.
這個破碎的民族.
這個有罪的民族.
將它救贖出來.
放在萬族萬邦的中心.
為了不是自己的存在.
而是為了其他人的特殊恩典而存在.
再回到十二章第一節.
亞伯拉罕故事的起初.
這裡有少許回響.
你看回這段經文十二章第一節.
對亞伯蘭說話第一次呼召.
創建以色列民族出來的時候.
他說你要離開本地本族富家.
往我指示你的地方去.
我必叫你成為大國.
我必致福給你.
叫你的名為大.
你也要叫別人得福.
萬族要因你得福.
地上的萬邦要因你得到祝福.
You will be a blessing to others.
只有以色列國不是為了自己存在.
只有神的子民存在的目的.
不是為了自己的益處.
而是為了萬國周邊的人的好處而存在.
說回魷魚遊戲.
沒錯是一個風靡全球的影集.
我看了一個年輕的台灣人.
他去評論這套影集.

$^{361}$我覺得評論得非常到位.
他指出在劇集裡.
有很多地方諷刺基督徒.
諷刺他們虛偽,偽善.
這個影評很簡單的評論.
我得到的地方是.
他說要給我們做一些謙卑的反省.
反省一下我們的信仰.
是不是搞錯了一些地方.
我們在上帝面前悔改.
我們也可以改變.
其中一點我覺得他說得很對.
很多時候被揀選的人.
新約是基督徒教會.
舊約就是以色列國.
被揀選的人往往有一種錯誤的思想.
就是一種選民的思維.
一種選擇人的心態.
什麼意思呢.
就是說覺得自己被揀選.
覺得自己被上帝寵愛.
我得救了,別人不得救.
不自覺地有一種優越感.
覺得自己高高在上.
這是一個很大的危機.
另一個危機就是.
以為自己比別人好了.
我就配受上帝的恩典.
創世記最後這一章.
都是說原來不是神的恩典.
來揀選亞伯拉罕和以色列的話.
他們不會是神的選民.
也沒有這樣的福氣.
能夠成為神的祝福.
以至流向萬國.
因為上帝的恩典揀選.
不是因為亞伯拉罕,以色列.
甚至新約的教會.
有什麼好處.
完全不是因為我們有好處.

$^{401}$我們都是百分之一百有罪.
是破碎的.
既然上帝的恩典能夠淋到我們的生命.
淋到我們的人生的時候.
我們應該怎樣去回應.
我的子女已經在讀大學.
在大學的城市裡面.
本身有一間很好的教會.
一間亞洲人的教會.
很多時候人的眼中都是一間相當成功的教會.
但偏偏在上一年.
在疫情的一年.
就爆出了教會裡面有很多的醜聞.
包括一些性騷擾.
一些cover up.
做錯的事不願意去承認.
不單止這樣.
這個醜聞爆出之後.
不但教會沒有直接去承認.
連區會中派都視而不見.
處理手法令人有一種感覺.
就是一種優越感.
高高在上.
不需要理會其他人的看法.
原本子女覺得他們這間教會是很好的.
但因為爆出了一連串的醜聞.
我的子女和他們本身有的朋友.
都不願意去這個教會.
一種高高在上 優越感的危機.
這些是上帝要去管教我們.
上帝要去提醒.
當頭棒喝.
告訴我們我們緣末有可跨.
你有沒有想過.
人海當中.
為什麼上帝的恩典.
偏偏臨到我和你.
不是因為我們完美光鮮.
是因為上帝的恩典.
以色列的族長.

$^{441}$他的故事 他的刻畫.
都說到他們脆弱 失敗.
偶爾會有信心的一刻.
但大致上 從頭到尾.
都說到上帝揀選的這一班.
是有血有肉的人.
神居然去揀選一班破碎的人.
來和他同工.
成就他自己的工作.
這才是最震驚的地方.
不是叫我們更謙卑地回應.
去事奉我們的主.
而是更謙卑地鼓勵我們.
以致我們更加願意被上帝使用.
如果神能夠藉著阿布拉罕.
以撒雅各這些罪人.
不要想到這些父親.
任何英雄.
他們都是破碎 和你和我一樣.
神能夠去用他們.
豈不是能夠用你和我.
在恩典揀選出來.
叫你和我沒有任何東西可以誇口.
誇的不是自己.
而是誇十字架 耶穌基督的恩典.
要誇口的 誇耀的.
是耶穌基督十字架的救恩.
是祂自己的名字.
在萬國當中.
能夠被上帝恩典揀選.
被上帝使用.
有機會認識上帝.
有機會去作上帝的工作.
這是何等大的恩惠.
你和我願不願意在上帝的恩典中.
向上帝說話.
求你謙卑我們 使用我們.
我們偶爾會跌碰.
但是我們真的願意用這樣的心.
去回應上帝.

$^{481}$說回魷魚遊戲的評論員.
他是一個年青的台灣人.
這個人原來本身也是基督徒.
他在一個基督徒家庭中長大.
他接受一個報章訪問他的時候.
他亦提出.
沒錯 有很多這些醜聞.
但他不能夠去否認.
就是在他成長的過程中.
有很多真心來事奉神.
真的能夠活出基督生命.
在背後默默服侍.
不會講求名氣.
關注到成就的這班人.
是最影響到他.
影響到他在信仰當中.
縱使看到有很多破碎.
有很多令他跌眼鏡的事情.
但鼓勵他繼續向前.
繼續事奉主.
你猜一下這個評論員.
這位年青人.
在他身邊對他有正面影響的基督徒.
第一個他會說誰呢.
就是說他父母.
原來在他身上.
父母 基督徒的父母.
對他最大的影響.
因為他的父母是宣教士.
默默在背後工作.
憑信心跟隨神.
縱使有很多難處.
很多困境.
但他願意依靠上帝的恩典.
活出基督的榜樣.
活出基督的生命.
就是這樣對他從小到大.
有一個很大的正面影響.
你和我應該是怎樣.
被上帝揀選.

$^{521}$畏著實行上帝的使命.
被上帝揀選.
是完全因為他的慈愛.
求主幫助我們.
以感恩 謙卑.
來回應神自己的揀選.
被上帝所用.
我們低頭祈禱.
主我們求你提醒我們.
我們實在是完完全全破碎的.
我們沒有可掛.
我們是一群完完全全墮落的罪人.
居然蒙你自己的恩典揀選.
求你謙卑我們.
讓我們能夠被你所使用.
去到和你同工.
去到履行你給我們的使命.
叫萬國能夠得福.
讓我們的生命默默地來事奉.
去到跟隨你的時候.
活出基督的生命.
以致更多人欣賞我們的生命.
所垮口的不是我們.
有什麼好處.
而是神你自己救恩的偉大.
祈禱奉主名禱.
\newpage



\section{}
\label{sec:DbCFhByjju0}
\textbf{Sing Hallelujah to Lord 詩歌意義 (黃嘉漢牧師)}
\newline
\newline
連結: \href{https://youtube.com/watch?v=DbCFhByjju0}{\texttt{https://youtube.com/watch?v=DbCFhByjju0}} ~~~~ 語音日期: 2024-04-06
\newline
\newline
\hyperref[sec:DKQrjJlzvbM]{\small{< < < PREV SERMON < < <}}
~
\hyperref[sec:index_chronic]{\small{[返順時目]}}
~
\hyperref[sec:index_scriptual]{\small{[返順卷目]}}
~
\hyperref[sec:Dw3oWngj5UA]{\small{> > > NEXT SERMON > > >}}
\newline
\newline
$^{1}$I believe that in 1974, when Linda Stratton wrote "Sing Hallelujah to the Lord",.
she never imagined that this song would become a song that has a political connection to the present..
In Hong Kong, when there are police officers and government officials,.
the public surround them and sing "Sing Hallelujah to the Lord"..
Indeed, when we listen to this song,.
the tune of this song is very calming and calming..
When we sing this song, we can feel that there is a greater truth and peace in this world..
In this turbulent and unstable world, everyone wants to be safe..
This song has five verses..
Many people sing this song, including those who do not know Jesus Christ..
What message does this song want to convey?.
Let me explain the meaning of the five verses..
The first verse is "Sing Hallelujah to the Lord"..
"Hallelujah" means "Praise the Lord"..
It is a praise from the bottom of the heart..
How beautiful is God?.
The four verses of this song give us some inspiration..
The second verse is "Jesus is risen from the dead"..
In the Christian faith, Jesus Christ is the Son of God..
He died on the cross to save people..
He saved people from being released from sin..
When he died, he conquered death and the evil that controls people's lives..
Why do we praise the Lord?.
Because God saved us from death and saved us from being released..
So that people can have a new beginning..
The third verse is "Jesus is Lord of heaven and earth"..
This is a very important message..
In the first chapter of Genesis, God created the world..
When we understand the word "creation",.
we understand that it is not about evolution or creation..
For the Jews, it is very practical..
If this world is created by our Lord,.
it means that this world is owned by this Lord..
If this world is owned by this Lord,.
we must obey him because he is the master of this world..
He is the Lord of heaven and earth..
That is, no sovereignty in the world can be greater than Jesus Christ..
No sovereignty is greater than the Lord of heaven and earth..
The Lord of heaven and earth reminds those who believe in the Lord that there is only one Lord in the world..
The fourth verse is "Jesus is living in his church"..

$^{41}$It is true that the Bible always says that the head of the church is the Christian,.
but the body is the church..
It means that people should see Jesus,.
Jesus' goodness, Jesus' kindness, Jesus' kindness, Jesus' justice,.
through the true believers..
Some people say, "Ask Jesus to come and see me.".
Jesus has already come to see you and me..
He tells people through the church what is justice and what is love..
The fifth verse is "Jesus is coming for his own"..
When Jesus Christ created this world,.
he saw that the world was harmed by sin,.
and that the relationship between man and God was completely destroyed,.
and that when people lived in sin and tears,.
God came to this world to conquer evil and death..
He came to this world again..
In the 21st chapter of the Book of Ecclesiastes,.
the third and fourth verses say that God is in the world,.
and that God will be in this world again..
There will be no more death, sorrow, and pain..
These teachings are included in the verses..
But what do they mean to us?.
They give us a few meanings..
First, when we sing this song,.
we must remember that there is a God in heaven who loves us..
He is not only coming to get rid of the unjust world government.
and the unfair social system..
Those are important..
But most importantly, he wants you to know him,.
so that our lives can be changed,.
so that the evil and death in our lives can be solved..
When I was young, there was a saying that came out of my mind..
"Every day, five or four young people visit each other.".
Every person, every regime can turn from love and hatred to evil..
Only when Jesus Christ returns to this world,.
can we turn the world around..
One of the ways to change our lives is to save us,.
so that we can be set free from sin,.
so that people in this world can see what love and justice are..
So, first, he hopes that when you sing this song,.
you will not only see the problems in society,.

$^{81}$but also your own problems, and be willing to return to him..
He also wants people to know that he is the true master of the world..
All the regimes in the world,.
whether you are a democratic or an extreme right wing society,.
all have the tendency to tell us that he is the true master..
Sorry, the true master of this world is Jesus Christ..
So, the church in this world is telling everyone in this world.
that there is only one master in this world..
And this master invites us to live a new life of another kind of nation,.
a new life, a life of love, a life of justice,.
a life of justice, peace and love..
The regime is not our master..
We are not our own masters..
The rich are not our masters..
Only Jesus Christ is our master..
So, here is the second reminder for us..
When we sing this song,.
we are welcoming the fact that there is a master who loves us..
He is waiting for us..
Finally, the church is telling people to know Jesus,.
to believe in the Christian community of Jesus Christ..
The church on earth is telling people that there is love,.
there is light, and it comes from Jesus Christ..
He will ultimately rule this world,.
and bring people back to his side once again,.
and give them true freedom..
So, when we sing this song today,.
whether you are a Christian or not,.
I hope you will sing this song,.
and you will see the God who commands justice and peace..
I do not deny that God really wants us to.
carry out justice and peace in society..
This is important..
But at the same time,.
he also invites you and me to walk with him..
Our life is to be returned to him,.
and our life is to be led by Jesus Christ,.
not by any regime..
We are to be led by Jesus Christ,.
to reveal the love of justice and peace.

$^{121}$in our society,.
and to wait for him to return,.
to show the love of justice and peace.
to the whole world..
You and I live every day.
to reveal the love of justice and peace of Jesus Christ..
I hope that at this moment,.
we can experience God's peace,.
peace, and blessings to everyone..
[END].
[BLANK AUDIO].
\newpage



\section{}
\label{sec:Dw3oWngj5UA}
\textbf{Faith and time: Feast of Tabernacle 2 (Rev. Mark Chan)}
\newline
\newline
連結: \href{https://youtube.com/watch?v=Dw3oWngj5UA}{\texttt{https://youtube.com/watch?v=Dw3oWngj5UA}} ~~~~ 語音日期: 2024-04-07
\newline
\newline
\hyperref[sec:DbCFhByjju0]{\small{< < < PREV SERMON < < <}}
~
\hyperref[sec:index_chronic]{\small{[返順時目]}}
~
\hyperref[sec:index_scriptual]{\small{[返順卷目]}}
~
\hyperref[sec:DrODNpzG4PQ]{\small{> > > NEXT SERMON > > >}}
\newline
\newline
$^{1}$Okay, our colleges are closed, okay?.
All right, let's see, how about next month?.
Any important dates next month, November?.
Thanksgiving, yes, Thanksgiving..
So we're coming to the end of the year..
In fact, as we are talking about this series.
on faith and time, think about how you use your time..
What are the important reminders in our time,.
in our calendar?.
And we're talking about the Feast of Tabernacle, right?.
The Feast of Tabernacle, we talk about in biblical times,.
in the Old Testament or in the New Testament..
It's a very important date..
In fact,.
today,.
after sundown, Sukkot begins..
Sukkot, I know, we in United States,.
so we are not really, we don't really care, right?.
So I flew to, last week I flew to.
this Midwestern seminary,.
and I was going through one of my seminars..
So I flew back last Friday,.
or Thursday, actually, Thursday evening..
So I got off the plane, in Midway Airport, if you remember,.
from the gate all the way to the baggage claim,.
there's quite a walk, right?.
Five to seven minutes or 10 minutes,.
depends on how fast you walk..
But I met, I came across a couple,.
they were in black and white,.
and they were holding something, some plants in their hands..
And then I said, "Happy Sukkot,".
because they're a Jewish couple..
They are heading to probably Skokie.
or some other Jewish community in the city of Chicago..
So, hey, there you go, I saw them,.
"Hey, you know, Sukkot," right?.
So, happy Sukkot..
So sundown today begins the Sukkot for seven days..
So tomorrow is the first day of the Feast of Tap Neckles..

$^{41}$So we will take some time to examine.
some of the passages in the Bible..
Remember I said in the Old Testament,.
as well as in Jesus' time,.
there are three main festivals.
to remind the people of God, the grace of God..
The Passover, around what month?.
In our calendar..
Yeah, April, yeah, you're right..
April, the spring..
And then there's Shavuot, or Pentecost,.
which is seven weeks after Passover..
So it's another half a day for the wheat and barley..
And then the third one, the biggest festival is Sukkot..
All right, so it's this week, all right?.
So let's take a look what the people of God.
in the Old Testament, how they celebrate Sukkot..
So here I have a couple of verses from the book of Nehemiah..
So when the Jewish people,.
they were in Babylonian exile for 70 plus years..
So they returned to the land, to Jerusalem,.
and this is what they did..
After building temple, we building a temple,.
we building the city wall, everything is done..
And one more thing they do is to celebrate the festival..
So chapter nine, verse 15,.
the leaders are telling the people to go to the hills.
to get branches from olive, wild olive, myrtle, palm,.
and all these kind of trees..
And they use these branches to make sukkah, shelters,.
in which they will live during the seven day festival.
as prescribed in the law of Moses..
So the people went out and cut branches.
and used them to build shelters, sukkah,.
on the roofs of their houses,.
because most of the Palestinian houses,.
they have flat roofs, unlike Chicago..
In the courtyards, in the courtyards of God's temple,.
in the squares just inside the water gate and Ephraim Gate,.
so everyone who had returned from captivity.

$^{81}$lived in these sukkahs during the seven day festival..
They were all filled with great joy..
And the Israelites had not celebrated sukkot.
since days of Joshua..
Somehow they have forgotten,.
or they have not had the full scale of celebration.
for so long, since the time of Joshua,.
since the time they entered the land..
So this is what they do..
When on sukkot, they will prepare, as I said,.
in Midway Airport, I met a Jewish couple..
They were holding in their hands,.
you see the Jewish man in front of Western Gate,.
or the Western Wall,.
he was holding a bundle of plants..
And just like the book of Leviticus talks about,.
on the first day, gather branches from magnificent trees,.
palm trees, leafy trees, willow, myrtle, or whatever,.
and celebrate with joy before the Lord your God.
for seven days..
So what they have today, the modern Jews,.
they would put palm branches, myrtle,.
and willow together, three kinds of green leaves together..
And then there's one citrus fruit called itro,.
and then they put together, and put in the bundle,.
and then they would shake it, front and back,.
left and right, up and down..
And this is how they pray..
So it's a time of celebration,.
it's a time to pray together..
In fact, in Jesus' time, the Jews in Jerusalem,.
or in the land of Israel, this is how they celebrate,.
sukkot, they would gather some branches,.
maybe different, but something similar..
They would pray together because in sukkot,.
it's actually a harvest day..
It's the last harvest of the year..
Okay, it's like you getting your paycheck, right?.
At the end of the year, it's my last paycheck, right?.
I'm waiting for my last paycheck or bonus, right?.

$^{121}$Yeah, there you go..
And then you get your check or whatever, your deposit,.
and then you hold it in your hand,.
and you kind of shake it, you know, in all directions,.
and then you pray..
Now, what do they do?.
In fact, they are thanking God.
for what God has done in their lives,.
something that doesn't seem very important.
because year after year, you get your harvest, right?.
Month after month, year after year, you get your bonus,.
right, you get your paycheck, you get your salary,.
something that occurs over and over again, right, routinely..
But sukkot is to remind the people,.
you know, whatever happened in your lives,.
God's provision, you should not forget,.
to remember, to pray with thanksgiving..
And then at the same time, sukkot falls in the days.
or the time that the summer is fading away,.
and the autumn and winter is coming..
You know, you feel like it, right, in Chicago, right?.
Summer is gone, too bad,.
and then autumn or winter is coming..
It's the changing of season..
Now, in the Bible, then, the changing of season.
is even more contrasting..
Now, in Chicago, summer is about, what, 80s, 90s,.
good days, right?.
80, 90s, Fahrenheit..
That's pretty hot and a little bit humid..
And then autumn and winter, what is the temperature now?.
50s, right, 50s, 50s, 40s, something like that..
And then it will go down, right?.
It will go down, all the way down to negative, right?.
So you feel the temperature change, right?.
In the land of Bible, you also feel that contrast.
of summer and winter,.
because there's actually two seasons only.
in the land of Bible,.
summer, which is hot and dry, right?.

$^{161}$90 degree, 100 degree, very sunny, scorching heat..
And then after Sukkot, the winter will come,.
and winter in the land of Bible is wet and colder,.
cold and wet..
Summer is hot and dry, and winter is cold and wet,.
because the winter season kicks in the first rain,.
the first season of rain..
Now, when I say rainy season,.
I don't mean one month has 30 days,.
and then 30 days is gonna get rain, no..
A rainy season means in a month,.
you get like nine days or 10 days of rainy days..
So when Sukkot comes, rain is expected..
Rain is gonna come..
So the Jewish man here in front of Western Wall,.
he and with other people are actually praying for rain,.
praying for another year that God will bless them,.
God will provide for them,.
because in the agricultural society,.
you need water and you need rain in order to get harvest..
Right at this time in Sukkot,.
they will keep praying and praying.
that this coming year, God will bless them,.
God will bless us, so we can have enough water, enough rain,.
so that would guarantee our harvest coming, all right?.
It's kind of like you starting a new fiscal year,.
what do you expect?.
You anticipate growth, right?.
Hope that there will be a financial gain,.
or the report, the first quarter report,.
will match your whatever expectation.
or plan that you have, right?.
It's kind of like that in Sukkot..
So praying for rain..
In fact, for seven days in Sukkot,.
the last day, the seventh day,.
is the highest holiday of the entire week..
And the priest in Jesus' day,.
they would go down from the temple.
to the downhill area of Jerusalem,.

$^{201}$and there's a big water pool there,.
it's called Saduam Pool..
They would use a golden flask, golden flask,.
and pick up some water from the pool,.
and then bring it all the way back.
with processions and celebrations,.
go back to temple and the altar,.
and they would do a libation..
You know, they would sprinkle the water.
in front of the temple,.
praying that God will give them a new year.
with provision and blessing..
Now, this kind of festival or celebration.
reminds the people of God,.
everything comes from God,.
as basic as the rain from heaven..
It's only because God provides, right?.
What a reminder for you, too..
Whenever you pick up your paycheck,.
whenever you get your bonus, right?.
Have you ever thought about,.
we should be thankful?.
God provides, right?.
There's no guarantee, right?.
I was just talking to a friend,.
you know, we're getting old, right?.
We're getting old, we're like old cars..
Some of our parts are not working, right?.
(congregation laughing).
And we don't expect that..
We don't expect that part will recover, right?.
Once you get some permanent damage.
or permanent deterioration, right?.
It's gone, right?.
It's kind of like your old car, right?.
The water pump is not working..
You can replace it,.
but once in a while, it will get stuck again,.
and then you have to replace it again..
And then come to a point that.

$^{241}$you realize the old car is not worth to fix it anymore..
Right?.
It's kind of how life is, right?.
So I was talking to some of my friends.
because we are over 50s,.
and we are getting to the point.
that we know that some water pumps are not working,.
and some of our parts are not working..
And we need to go to a physical therapist,.
we need to go to see doctors,.
we need to go to see dentists,.
and there's some deterioration..
When it happens, it will not be renewed again, right?.
Now last week, I remind you,.
life is short, life is fragile..
It is..
We are broken, and we will be broken..
And we better think about how we invest.
our lives and our times and our effort.
into something that will last..
Sukkot is one of those festivals,.
year after year, to remind the people of God.
what you're living for, right?.
Who is providing for you, and for what purpose, right?.
I know we all have our own anxiety and our own concerns,.
right?.
Some of you may be worrying about your jobs,.
your career, your education,.
your family, your visa..
Those are the practical stuff..
That causes anxiety, right?.
Sukkot, one of those time reminders,.
is to remind us every single day,.
it's by God's provision, right?.
Now if I ask you, what are you thankful for these days?.
In the past few weeks, what are you thankful for?.
Can you think of something that you are thankful for,.
for the person next to you?.
That's great..
So nice..

$^{281}$(audience laughing).
What are you thankful for?.
Right, think about it..
Sukkot is one of those days..
In fact, I told you, the last day of Sukkot,.
the seven day festival, is the day of greatest expectation..
The people of God will pray for rain.
for the next season, for the next year..
That God will give them blessing, right?.
And here, in the book of John, the gospel of John,.
we see what happened in the last day of Sukkot,.
the seventh day..
Here is the story of Jesus..
Now what happened before was that.
some of the brothers of Jesus,.
they went from the north, Galilee,.
going to Jerusalem to celebrate Sukkot..
And those brothers did not believe in Jesus at that time..
So they were telling Jesus,.
hey, if you are really the Messiah,.
why don't you go to the Sukkot with us?.
Go to the temple, go to Jerusalem,.
and announce yourself to be the Messiah, to be the king..
And then everyone will believe in you..
Of course, Jesus doesn't do any,.
that kind of show and tell, right?.
So Jesus said, I'm not going to Jerusalem,.
I'm not going to Sukkot..
In fact, he went secretly..
He went secretly in Jerusalem..
He didn't show up on stage..
He didn't talk to the people until the last day of Sukkot..
So 737, Gospel of John..
On the last day of Sukkot, the climax,.
the highest day of the festival of Sukkot,.
the tabernacles, Jesus stood and shouted to the crowd,.
"Anyone who is thirsty may come to me..
"Anyone who believes in me, come and drink..
"For the scriptures declare,.
"rivers of living water will flow from his heart.".

$^{321}$Now, when Jesus said living water,.
he was speaking of the spirit,.
who would be given to everyone believing in him..
But the spirit had not yet been given.
because Jesus had not died and rose again.
and entered into his glory..
So at that point in chapter seven of Gospel of John,.
on the last day of Sukkot, when people are so thirsty,.
so thirsty because after four or five months.
of dry and heat, dry and hot summer,.
they were expecting, they were wanting,.
yearning for water, for rain..
The last day, the best time to announce this gospel.
is that anyone, if you are thirsty,.
everyone you can imagine,.
you know, when their water tank is lowest,.
when their bank account is lowest,.
they were expecting a new year, right?.
Kind of like you spend all your money, right?.
At the end of the month, at the end of the year,.
you're expecting God will provide for you, right?.
When you are poorest, when you're a lack of money,.
when you're a lack of water,.
you're expecting God will bless you..
And Jesus said and announced on the climax of the day,.
if you are thirsty, come to me,.
because I give you more than literal water,.
living water, eternal life, eternity..
This is what you want..
What you really want is not just a paycheck..
What you really want is me giving you the life of God.
that will last forever lasting..
It will not end..
And sometimes we think what we want.
is the ultimate need that we have..
Jesus said, no, it's only for temporary relief..
Next year, you're gonna need water again..
Next month, you need another provision again..
You need another paycheck..
You need another level of happiness.

$^{361}$or whatever you're seeking for..
The true need, the ultimate need that we need.
is eternal life..
And that's what Sukkot is all about..
On the climax of the festival, Jesus said,.
if you're thirsty year after year, month after month,.
you find yourself lacking week after week,.
I announce to you, I give you river of living water,.
the best water, the best water..
The everlasting supply of water is from me,.
from the Spirit, right?.
Now we're in Christ..
We enjoy this kind of living water.
because Spirit is living in you..
Now we often struggle after we follow Christ,.
is the Spirit speaking to me, right?.
Have you ever wondered, the Spirit, God is speaking to you?.
I think God is constantly speaking to us..
The problem is our world is too noisy.
and we cannot hear the voice of the Spirit, right?.
We're too busy with our own things,.
too noisy in our environment,.
and we hardly hear the word of God, right?.
The problem is not the speaker..
The problem is the audience, it's us..
Sometimes we don't slow down..
We don't read the word of God, right?.
And we treat our relationship with God as transactional..
Whenever I have a need, right?.
Okay, I am building my family, I need some more money..
You know?.
God, I pray, and out of the blue, you drop some money, right?.
Well, you don't do that..
So, you know, I'm gonna forsake you..
You see that kind of relationship?.
We have a list of needs we present to God.
and hope that God will fulfill those needs, right?.
If God fulfill those needs,.
then I'll come draw closer to God, right?.
For what?.

$^{401}$For more needs, right?.
We need, we have more wants..
And then if God satisfy those wants,.
we are satisfied and then we come back to God.
with even more needs, right?.
So, the cycle go back and forth, back and forth,.
you know, just continuously..
And you see that kind of relationship with God.
is transactional..
You know what transactional is?.
It's like husband and wife, right?.
Well, how do you tell a husband and wife.
really love each other?.
How do you tell they don't really love each other?.
It's when, you know, I do the dishes.
so that you do other things, right?.
I mop the floor and you're gonna bring food on the table,.
right?.
If I don't mop the floor, right?.
Okay, then you don't bring food on the table..
In order for you to bring food on the table,.
then I'll have to do something, right?.
It's transactional, it's an exchange..
I do something and you do something, right?.
I do something to earn you something..
Now, that kind of relationship is actually idolatry..
It's worshiping idol, right?.
How do you know two people are really loving each other?.
Is that the husband is loving the wife 100\%,.
unconditional, and the other person is loving back,.
unconditionally..
That is a real relationship..
And our relationship with God is the same thing..
It should not be transactional..
It should be a relationship we can enjoy, right?.
I serve God, it's not because I wanna get good stuff.
from God..
I'm faithful to God, it's not because God is good to me..
God is sprinkling some money to me,.
and then I'll do something back to earn the blessing..

$^{441}$No, that is not a real relationship.
because it becomes transactional, it becomes an exchange,.
becomes a financial transaction, right?.
In order for us to really relate to God,.
I mean, eternal relationship,.
is for us to understand God loves us unconditionally,.
fully, 100\%..
So we love Him back, not because we wanna gain something.
or earn something from Him..
It's because we respond thankfully, right?.
All the things that we do for God, in ministry, in service,.
by living closely to the Word of God,.
what we do for God is a response,.
is to tell God, "Thank you..
"Thank you for your grace, right?.
"I don't earn it, I don't deserve it.".
And that kind of relationship.
is the biblical relationship with God..
Now I hope this analogy will remind you.
what does it mean to worship God Himself?.
What does it mean to serve Him?.
What does it mean to follow Christ?.
Now I'll close with a short story about Sukkot, okay?.
About building the tabernacle..
I think years ago, there was a chief rabbi in England,.
in United Kingdom..
So it was in UK, his name is Jonathan Sachs..
He was a chief rabbi in England, everyone respects him..
And he was telling a really good story about Sukkot..
When he was younger, so much younger, he just got married..
So as a Jewish man, he got married.
and building his own family..
Sukkot comes and as a new family, husband and wife,.
they are supposed to build Sukkot, right?.
So he was not quite handy..
Anyone of you are handy here?.
Nails and hammer, putting things together, you are?.
Yeah, okay, you're not..
Well, Rabbi Sachs was not very handy..
So he didn't know what to do, right?.

$^{481}$How do you put things together like a tent or a Sukkot?.
How do you do that?.
So he doesn't know..
So he asked his friend, how to build a Sukkot?.
And his friend happens to be very, very handy..
So the friend invite him, hey, I'm going to like Home Depot,.
lumberjack or something..
And then we're gonna buy something for building my Sukkot..
You wanna come along?.
And Rabbi said, of course, yeah, I wanna come along, right?.
And this friend of Rabbi Sachs was very handy..
And he went into the Home Depot store.
and with a big list of materials and all the detail plan,.
right, how to nail the tent together.
with lights and all that stuff..
And it's kind of like that..
Kind of like that..
There's some Sukkot are really, really elaborate..
So his friend had a really elaborate plan,.
a blueprint, how to build a Sukkot..
So he was going with his friend, he had no plan at all..
So he ended up buying just three, four by eight,.
some nails, brought them all home..
And then his friend got all these plans, right,.
with all the lights and all the decoration on it,.
very beautiful..
So his friend built his Sukkot.
and he built his Sukkot outside of his home..
Each one has a Sukkot..
And his friend's Sukkot, of course, it's very elaborate,.
but his Sukkot is like a cardboard,.
four by eight nailed together only..
And then that night, a powerful storm came by,.
strong wind and strong wind..
And then the next morning, both of them got up.
and checked on their Sukkot..
His friend's Sukkot fell down, collapsed altogether..
Everything's gone..
And then he called Rabbi Sacks and said,.
"How about your Sukkot?.

$^{521}$"Is it still standing?".
And Rabbi Sacks went out to the patio and checked the Sukkot..
Yeah, my Sukkot is still standing..
And guess what, it's still here..
And his friend, in disbelief,.
he drove all the way to the Rabbi's house.
and checked on his tabernacle, his Sukkot..
And he was checking why,.
why my Sukkot, so elaborate, fell on the strong wind..
But his Sukkot is so poor and still standing..
And then he's checking his architecture,.
how he built things together and put things together..
He found there was one nail.
of the Rabbi's Sukkot was nailed to the wall of the house..
And that's why the Sukkot did not fall..
His Sukkot, however elaborate it is,.
fell after the wind and the storm..
And his friend, all of a sudden, just dawned on him..
"I know what that nail is..
"That nail is the nail of faith..
"It's the nail of faith.".
See, how elaborate you build your life,.
Sukkot, as your life,.
how many plans that you have,.
how many investments you have,.
it really doesn't matter at the end..
Because when the test of life comes,.
when storm hits you,.
your life will collapse,.
unless you have the nail of faith..
And that nail of faith, his friend realized,.
is our faith..
It's our faith in God..
It's our faithfulness in God,.
to walk in God faithfully..
I hope this lesson, this story,.
is not just interesting to you..
But it's insightful and transforming to us..
To realize what are we planning for,.
our 401k, our retirement, our family, our whatever..

$^{561}$If you don't have the nail of faith,.
when storm comes,.
when wind and rain hit you,.
you will collapse..
Only with the nail of faith,.
then you can stand..
C.S. Lewis, I mentioned to you last week,.
he said, "Do not let your joy.
"depend on something you may lose.".
There are many things that don't last into eternity..
There are many things like that..
And you will lose them..
You will lose all of them..
But there is something that will last,.
eternal life..
Right?.
Our service to God,.
our investment into His kingdom,.
they will all last..
And they will not collapse..
You know, in times of tabernacle,.
in times of Sukkot,.
if you go into a Jewish synagogue today,.
one of the scriptures that they will read.
during Sukkot, during this week,.
is the book of Ecclesiastes..
Now, the book of Ecclesiastes is a very long book..
It's 12 chapters..
Not gonna go through all the details.
'cause we're gonna end here..
But the central idea, the big idea.
of the book of Ecclesiastes.
is that everything is meaningless under the sun..
Now, under the sun, everything will repeat again and again..
It will be occurring all the time, right?.
No matter you live in Stone Age.
or live in today,.
everything will continue to repeat..
So everything is meaningless under the sun.
because someday you're gonna die, right?.

$^{601}$And we're gonna die..
So what makes life meaningful at the end,.
the 12th chapter,.
is that you don't live under the sun..
You should have a perspective to live above the sun,.
above this world..
You get an eternal anchor, the nail of faith,.
the eternal perspective to see what you're doing,.
big or small..
You got another anchor, eternal anchor,.
to evaluate what you're doing,.
whether it's worth it or not..
And that is life above the sun..
And Ecclesiastes says that is a meaningful life.
because whatever you do,.
if you have that perspective to fear God,.
it will last..
It will work eternally..
12.13, I'm gonna end with this last two verses..
Now, this is the whole story..
So he was arguing for 12 chapters..
Life under the sun is meaningless.
because we're gonna die, right?.
Here now is my final conclusion..
Fear God and obey His commands.
for this is everyone's duty..
God will judge us for everything we do,.
including every secret thing,.
whether good or bad..
Life of wisdom, life of worth,.
is to invest into something that is eternal, right?.
Let's all pray together..
Father, wherever we are in our life journey,.
we pray that you will remind us time and time again.
that our life and our time needs to be invested.
in eternal course,.
into your gospel, into your people, into your kingdom..
And I pray that you will help us to return back to you..
We return to you with a heart to serve you 100\%.
because you love us so much..

$^{641}$You love us 100\%, 200\%, beyond measure..
That we enjoy the eternal life.
because of the cross of Christ..
And thank you we can worship you and listen to your word..
And I pray that this coming week,.
no matter what we do,.
any things that seems insignificant,.
we do this to honor your name..
We do all of our works.
and everything that we do,.
pray that we will do something worthy for your kingdom,.
to invest our life and our time in you,.
to listen to your word, to pray to you, to worship you..
And I pray that you will grant us strength..
And thank you for being our God.
and being our loving God in Christ, we pray..
Well, God bless you.
and pray that you will continue to worship with us.
\newpage



\section{10:11-18}
\label{sec:DrODNpzG4PQ}
\textbf{I am a good shepherd (John 10\_11-18; Psalm 23) - Samantha Lockenour}
\newline
\newline
連結: \href{https://youtube.com/watch?v=DrODNpzG4PQ}{\texttt{https://youtube.com/watch?v=DrODNpzG4PQ}} ~~~~ 語音日期: 2024-04-07
\newline
\newline
\hyperref[sec:Dw3oWngj5UA]{\small{< < < PREV SERMON < < <}}
~
\hyperref[sec:index_chronic]{\small{[返順時目]}}
~
\hyperref[sec:index_scriptual]{\small{[返順卷目]}}
~
\hyperref[sec:_tKyeazYj4I]{\small{> > > NEXT SERMON > > >}}
\newline
\newline
$^{1}$Well, today we have our guest speaker, our youth minister,.
Sammy, to give us the word..
We want to welcome her..
All right..
OK..
That was such a sweet. Nathan, wow..
Yeah. What a crew..
I was blessed by that..
So last time I was here preaching at Urban Voice,.
I talked about how Jesus says I am the bread of life..
And I mentioned how there are seven I am statements in the book of John..
And I'm going to continue on that train today..
And I want to talk to us today about how Jesus calls himself.
the good shepherd later in John..
I'm really excited to be with you all today..
I've been thinking a lot about how in the book of Haggai,.
which is this really small.
prophetic book in the Old Testament,.
God asked the prophet Haggai and the people of God.
at the time to consider their ways together..
I've been reflecting on what a privilege it is to be here in this community.
with each one of you.
to consider our ways under the authority of God's word together..
So I'm really thrilled to be here with you..
And I've also been thinking a lot about why it's important.
to talk about who God says he is..
It can be said that all sin in our lives is rooted in a lie.
that either God is not who he says he is or we are not who he says we are..
Scripture says that our hearts are deceitful..
Above all else, they lie to us. They suppress the truth..
And so when we spend time meditating on who God says he is,.
we can destroy those lies in our hearts that distort our view.
of who God is. And we can only really do that by spending time in his word..
That's our safeguard for our deceitful hearts..
Yeah, so I am going to take some time to pray for us this morning.
before we jump in to talking about Jesus as our good shepherd..
Heavenly Father, you are so good to us..
You are mighty in all of your ways..
It's with your rod and your staff that you lead and guide us.
to life and life in you that is abundant and eternal..

$^{41}$Would you help us today.
to hear truth from your word that is bigger and greater than the lies.
that our hearts can tell us about who you are and how that you see us,.
your sheep, your people..
We are so thankful we belong to you, and we ask that your word be mighty.
in our hearts today, that it would speak to us, that it would convict us.
of what it means to belong to you and to worship you..
It's in your name, I pray..
Jesus, the good shepherd. Amen..
So two summers ago, I was working for a place.
that a lot of you are familiar with..
It's called UCA, and my husband and I were trying..
We had this really fun idea, along with the rest of the organization,.
to take 30 middle school students on the train to play mini golf..
Just the two of us. Super fun..
And to corral 30 students who aren't really familiar with the CTA.
and to get them downtown to Millennial Millennium Park.
on the hottest day of the summer was quite an adventure..
Not only did we have to kind of like shield the children.
from the interesting people on the train, but also we had to pass around.
all these water bottles..
We were basically throwing water at all these middle schoolers..
And I think the only reason we made it, as I reflect on that experience,.
is because my husband is so tall..
And so he stood out among the crowd of people..
And I remember vividly just shouting, "Look for Mr. Brock!".
And so when I think about a shepherd experience, I want you to consider.
that's one example of leading sheep..
But Jesus is a much more competent and capable shepherd.
than this example of just someone being tall leading a crowd of kids..
So I'm going to go ahead and I'm going to start by reading our passage today..
I've actually chose two passages again, because I think it's so rich.
how the Old and New Testament both reflect.
who Jesus says he is in the New Testament..
So and I keep finding that these I am statements.
have parallel texts in the Old Testament, which is so cool..
So today I'm going to read our passage over us,.
and I want you to listen and hear Jesus's deep love for you,.
his sheep, in each of these scriptures..
OK, so I'm going to start with John 10, 10 through 18..

$^{81}$And if you want to get there in your Bible, you can also keep a thumb in Psalm 23..
So John 10, 10..
The thief comes only to steal and kill and destroy..
But I have come that they may have life and have it to the full..
I am the good shepherd..
The good shepherd lays down his life for the sheep..
The hired hand is not the shepherd and does not own the sheep..
So when he sees the wolf coming, he abandons the sheep and he runs away..
Then the wolf attacks the flock and scatters it..
The man runs away because he's a hired hand and cares nothing for the sheep..
But I am the good shepherd..
I know my sheep and my sheep know me..
And just as the father knows me and I know the father,.
I lay down my life for the sheep..
I have other sheep that are not of the sheep pen, and I must bring them in also..
They too will listen to my voice and there shall be one flock and one shepherd..
The reason my father loves me is that I lay down my life only to take it up again..
No one takes it from me, but I lay it down of my own accord..
I have authority to lay it down and authority to take it up again..
This command I received from my father..
And then moving on to Psalm 23, the Lord is my shepherd..
I lack nothing..
He makes me lie down in green pastures and he leads me beside quiet waters..
He refreshes my soul..
He guides me along the right paths for his name's sake..
And even though I walk through the darkest valley, I won't fear any evil.
for you are with me and your rod and your staff, they comfort me..
You prepare a table before me in the presence of my enemies..
You anoint my head with oil and my cup overflows..
Surely, surely your goodness and love will follow me all the days of my life.
and I will dwell in the house of the Lord forever..
This is the word of the Lord..
Wow, I just even just need to take a moment because those promises are for us..
That is insane..
That is how the good shepherd sees his sheep..
And we're going to talk about three truths that we can glean.
from these scriptures that help us to know more of who God is.
and how he sees us..
The first section today that I want to talk about.
is the good shepherd leads us toward abundant life..

$^{121}$The good shepherd leads us toward abundant life..
And zeroing in on those first couple of verses here, John 10, 10,.
the thief comes only to steal, kill and destroy..
But I have come so that you may have life and have it in abundance or to the full..
And then the first couple of verses of Psalm 23,.
the Lord is my shepherd, I lack nothing..
Or other translations say no good thing..
He makes me lie down in green pastures..
He leads me beside still waters..
He refreshes my soul and he leads me in paths of righteousness.
for his name's sake..
The good shepherd leads us toward abundant life..
So taking a look at this first verse,.
Jesus says that he is all about abundant life..
In the first few verses of Psalms, just to repeat again,.
it says that the shepherd leads by green pastures and quiet waters..
When you hear these terms, what kind of feelings does that evoke for you?.
What pictures in this life?.
OK, maybe in the city we don't have that many still waters or green pastures,.
but what what is that.
abundance look like?.
What do you think of when you consider what abundance.
in a person's life really means?.
And a deeper question, does following Jesus.
seem like it holds abundance for you?.
Jesus, as the good shepherd, promises that he leads us.
toward abundant life and not away from it..
And this might feel really contradictory for some of us.
when we weigh the cost of following Jesus..
It can be so easy to think about how following Jesus is very, very difficult..
And let's be real, it can be and it is, especially in our society and culture..
This isn't unique to just us here in this church or in this time..
But Paul, the apostle, if you think about him in the book of Acts,.
he has a laundry list of ways that he is beaten, scorned, abused, mocked..
And that's in the book of Acts laid out for us..
And so sometimes the work of choosing to trust in God.
and that belief he will come through for his people is even exhausting to consider.
when we think about what our present circumstances are..
And yet scripture says that he gives us everything we need for life and godliness..
So sometimes we don't trust that God, in God, is all life in abundance..

$^{161}$And we try to manufacture our own happiness or chase after things.
that make us feel happy and comfortable in the moment..
We strive after that picture of abundance rather than asking God,.
what does it really mean to have abundance in our lives?.
Scripture promises that the way that seems right to man actually leads to death..
Therefore, the opposite is true in scripture,.
that the way that doesn't seem right to man,.
the way that seems really hard and really difficult.
is actually the way that leads to life and abundant life..
Jesus is talking about a spiritual reality here.
that is greater and better than our physical reality,.
and it actually fuels the whole rest of our lives.
and our journeys here on this physical earth..
Jesus says in another passage to his followers.
who are worried about what they will eat or what they will drink or what they will wear,.
if they follow Jesus, they choose this cost,.
this costly journey of following Jesus where he leads..
He says to those people, seek first the kingdom of God.
and then all these things that you're longing for will actually be added to you..
I give you everything you need and just like Psalm 23 says,.
you will lack no good thing..
Jesus isn't saying, OK, if you follow God,.
you automatically get to date that guy you were really hoping to date.
or get that meal you were really hoping that you would be able to afford,.
but rather he is the one who gives us ultimate and full peace..
He gives us deep peace and he provides that stillness.
and security and life in himself and not in our present circumstances..
Everything else God provides for us on top of us, on top of that, well, that's a bonus..
Following this good shepherd, the one whose ways lead to abundant life,.
revives our souls when we start becoming weary of doing good..
When the weight of following Jesus feels never ending and exhausting,.
trusting the good shepherd means leaning into the rest that knowing Jesus provides..
Can you think of a time recently when you felt at the end of your rope,.
maybe you felt burnt out?.
Jesus sees you in the midst of this exhaustion..
Of course, Jesus can provide us seasons of physical rest,.
but he also provides us a deep spiritual rest in himself.
in the context of a relationship with him..
The good shepherd leads us toward an abundant life..
So we are fine when he leads us by still pastures or leads us by quiet waters,.

$^{201}$but what happens when we are walking through the valley of the shadow of death?.
Yikes..
Does our confidence in God begin to waver?.
The good shepherd is our defender in the darkness..
Let's go back to our passage here..
Next few verses of John 10..
I'm having you flip back and forth a lot today, so sorry about that..
John 10 verses 11 through 13..
"I am the good shepherd..
The good shepherd lays down his life for the sheep..
The hired hand is not the shepherd and does not own the sheep..
So when he sees the wolf coming, he abandons the sheep and runs away..
Then the wolf attacks the flock and scatters it..
The man runs away because he's a hired hand and cares nothing for the sheep.".
And then the next couple of verses in Psalm 23..
"Even though I walk through the darkest valley or valley of the shadow of death,.
I will fear no evil for your rod and your staff, they comfort me..
You prepare a table before me in the presence of my enemies..
You anoint my head with oil and my cup overflows..
The good shepherd is our defender in darkness.".
So I ask you again,.
what happens when you are no longer experiencing the highs of life.
or the peacefulness of a calm season?.
What happens when the what ifs you are most afraid of actually happen.
and your fears are realized before you.
in some of the most painful and traumatic moments of your life?.
How is the good shepherd both good and a shepherd in those seasons?.
Here in this passage,.
we see a mighty God who takes on the evil of the wolf fearlessly,.
who carries a rod and a staff..
And takes the wolf who's trying to attack the sheep in his sheep pen.
by the teeth and throws him out of the pen..
This is a depiction of mightiness and authority..
This is the God who has fearlessly defeated any evil.
that can possibly triumph over our lives, whether that's our own sin.
or the sin of others that affects us..
For some of you, maybe this mightiness in God has freaked you out,.
either now or in the past..
Maybe you don't want to believe in a God who has complete sovereignty.
over your life because you're fearful about how God might use.

$^{241}$that rod and staff against you..
But if you belong to God, if you're one of his sheep,.
then you must know that that rod and staff functions to protect you..
If you've been walking with God for a while now,.
perhaps what used to be a fear of God's mightiness.
is something you now find comfort in, knowing that God has the final say..
Maybe you feel safe and secure and comforting to know..
But for all of us, the truth of this scripture can become comforting.
when we consider the good shepherd in the image of us.
walking through the valley of the shadow of death..
In Israel, this is an actual place..
And this might be wrong, and probably Pastor Mark knows the right answer for this,.
but I believe it's called the Valley of Hinnom..
And I think that I've been there..
And from what I remember of my memories, it's this rocky crevasse.
that has steep rocks on every single side..
And historically, this would be a place where wild animals would run..
Where outside of the old city of Jerusalem,.
they used to dump trash into this valley..
So kind of a disgusting and dark place, and it makes sense.
why all these wild animals were roaming there..
Pretty that, would you have complete confidence that you would be safe.
walking through that valley?.
What does it mean for us as believers to believe in a God.
that makes us fearless in the face of the valley of the shadow of death?.
How do we move from fear to confidence in God?.
And how do we move from what if to even if this happens,.
my God will come through for me again and again?.
Here's how..
We can trust that the Good Shepherd always destroys the wolf..
The wolf of sin and death can make us afraid..
It can come to make us intimidated or fearful that our own sin.
or the sin of others will be ultimately victorious..
But the Good Shepherd, he's there to protect us..
In Isaiah 118, it says, "Come now, let us reason together..
Though your sins be scarlet, they will be white as snow,.
and though they be red as crimson, they will be as wool.".
This is a promise of God in our lives for how he deals with our own sin..
His wrath was placed on Jesus so that his wrath is directed.
towards protecting us from our own sin..

$^{281}$I don't know if any of you have ever seen.
like those people who play flutes for boa constrictors before,.
and they make, I'm not describing it well,.
but they play flutes and they make the boa constrictor dance..
And they've always been amazed, like, how terrifying,.
in the face of a boa constrictor, you're just playing this cute little song..
How can you feel safe and secure that the boa constrictor.
won't one day reach out and all of a sudden,.
how can you be sure that you're safe?.
Well, I learned recently that actually,.
these people who play the flutes for the boa constrictors.
glue the mouths of the boa constrictors shut..
Now, that's kind of a disturbing thing to realize about that..
But I think that picture can translate to how we can view.
the sin in our own lives or in the lives of others.
and the ways that it affects us, that God has complete control.
and ultimate victory over the sin and darkness in the world.
and the evil in the world..
And because of this, we can be bold and courageous.
because we have the right man on our side..
When Brock and I were seniors in high school.
and recent Christians, we had the opportunity to interact with,.
we had the opportunity to visit North Africa and sit with missionaries.
who spent their entire lives working with shepherds.
in the hills of this North African country..
And one of the missionaries was telling us a story about his interactions.
when he first got to the country with one of the shepherds..
He said he he had stored up all these stories in his mind.
that he'd heard from preachers over the years about how dumb sheep are..
And he brought those to the shepherd.
and started telling him, I've heard that sheep are so dumb.
that they just kind of wander in the same patch of grass.
and they eat all the grass there..
And instead of moving on to another patch, they just don't realize it..
They just start eating each other's manure in a circle like,.
they are so dumb, aren't they so dumb?.
And I will never forget the missionary turning to Brock and I and saying,.
you know, that shepherd looked me right in the face..
And he said, don't call sheep dumb..
These are my sheep. They matter to me..

$^{321}$The shepherd had incredible love for these animals,.
and it's the shepherd's love and care for these animals.
that helped me begin to see this passage differently..
And I think one of the most beautiful truths from this passage.
is that the good shepherd is not a hired hand..
He's not going to abandon us..
He chooses to be with us. It's his choice..
That's not something the father makes him do..
As it says in the passage, he has his own authority.
to lay down his life and pick it up again..
It's his choice, Jesus's choice to sacrifice everything for his sheep,.
for you, for me..
And in the face of great evil, because he's not a hired hand,.
he will defend us again and again..
This isn't just his job..
It's his never ending love for us that compels his actions toward us..
The good shepherd is our defender in the darkness..
So God's word promises that his voice is powerful enough.
not to just defend us, but also to lead us..
Let's take a look at the next few verses here as we wrap up our time..
In John 10, verses 14 through 18 say, "I am the good shepherd..
I know my sheep and my sheep know me..
Just as the father knows me and I know the father,.
I lay down my life for the sheep..
I have other sheep that are not of the sheep pen, and I must bring them in also.".
I'm just going to pause and say that means us..
If you are not an ethnic Jew who was born in Jerusalem,.
this promise is talking about Jesus is saying,.
"I'm going to bring others into my kingdom.".
So this is really sweet to realize Jesus is talking about you and me in this passage..
So anyway, "I must bring them in also..
They too will listen to my voice, and there will be one flock and one shepherd..
The reason my father loves me is that I lay down my life only to take it up again..
No one takes it from me, but I lay it down in my own accord..
I have authority to lay it down and the authority to take it up again..
This is the command I received from my father.".
And the last verse of Psalm 23, "Surely, surely, your goodness and love.
will follow me all the days of my life, and I will dwell in the house of the Lord forever.".
The good shepherd leads us to everlasting life..
For many of us, we have asked at various points in our lives,.

$^{361}$"What is the purpose of my life? What is God's will for my life?".
I'm not saying these are in and of themselves bad questions..
These are questions in which to sit with before God and faithfully ponder..
But at one point in the book of Isaiah, the prophet Isaiah hears—.
and this is why it's good to sit and ponder with these questions—.
in the book of Isaiah, the prophet hears from God that if he walks down one possible path,.
and then another, he will hear a voice from the Lord saying, "This is the way. Walk in it.".
This is an illustration of what it looks like to faithfully discern the paths God might be leading you towards..
However, sometimes it can be so easy to become so caught up in God's will or purpose for our lives,.
like, "What good task does God want me to do? What college does he want me to attend?.
Does he want me to take this promotion or stay in my current position?.
What do the next five, ten years of my life look like?".
That we can actually miss the true life God has called us into,.
where we are deeply attuned to every wavelength of his voice in our lives..
God desires our dependence on his voice..
In this passage, we are promised that if we belong to God, we can know his voice..
That's a really sweet promise to rest in..
And we can trust his leading voice in our lives, because as the passage shares,.
at the end of this part in John 10, he has eternal authority..
And at the end of Psalm 23, that this kind of authoritative leading goes beyond this life,.
but actually continues in the life to come, where if we belong to the shepherd's flock,.
we experience the goodness and love that will continue in our lives as we dwell in God's house for eternity..
I've also heard stories about sheep before, that a shepherd will actually take a newborn sheep.
that can easily run away from him, because maybe the sheep isn't yet familiar with the shepherd's voice,.
and that the shepherd will actually break the legs of the sheep,.
and put the sheep around the shepherd's neck for a few years, until the sheep's legs heal..
Now this might seem really cruel and unusual, and it is kind of scary,.
but the reason the shepherd does this is actually for the safety of the sheep..
That while the sheep is around the neck of the shepherd,.
the sheep is learning to hear the voice of the shepherd..
The sheep is learning the voice so that one day, when the shepherd leads the flock,.
that the sheep can't go astray, that he knows the voice of the shepherd, and is able to follow after him..
God's love in our lives compels us to hear his voice,.
and it's because of his love that he wants us to know his voice, his leading voice in our lives..
A butcher would drive the sheep towards something,.
would try to get them to go into pen to be slaughtered,.
but a shepherd compels them to follow him..
The voice of the Lord leading in your life is one of deep and everlasting peace and joy..
It's from a place of deep and personal love for you, that he leads you and me..
The good shepherd leads us to everlasting life..

$^{401}$Jesus says in both of these passages that he is the good shepherd,.
and I believe that he longs for us to know him as the good shepherd,.
as one who leads us towards abundant life, those still waters and green pastures,.
as one who is our defender in the darkness, who destroys and defeats the wolf..
He's already done that and will continue to defend us in the valley of the shadow of death,.
and as one who leads us, not just leads us in our day-to-day decisions,.
but toward everlasting life, which is ultimately found in him..
Thanks be to God for that. Let's pray together..
God, you are so merciful that we would be counted as yours..
God, thank you for the opportunity we have to know you,.
and to be able to have confidence that if we know you and we belong to you,.
we can hear your voice..
God, thank you that in the face of our worst,.
\newpage



\section{詩篇}
\label{sec:_tKyeazYj4I}
\textbf{聖經地理 EP3 - 地理與釋經 - 詩篇125-2 (陳崇基牧師)}
\newline
\newline
連結: \href{https://youtube.com/watch?v=-tKyeazYj4I}{\texttt{https://youtube.com/watch?v=-tKyeazYj4I}} ~~~~ 語音日期: 2024-04-08
\newline
\newline
\hyperref[sec:DrODNpzG4PQ]{\small{< < < PREV SERMON < < <}}
~
\hyperref[sec:index_chronic]{\small{[返順時目]}}
~
\hyperref[sec:index_scriptual]{\small{[返順卷目]}}
~
\hyperref[sec:vWnZGyYbGoI]{\small{> > > NEXT SERMON > > >}}
\newline
\newline
$^{1}$我們說大衛城 石安 耶路撒冷 原本的耶路撒冷.
是四面環山 所有的山都比它更高.
這裡有一個影片是我在大衛城上拍攝的.
做一個360度的循環的時候 你就可以看到地景是怎樣.
現在我們在大衛城站著 向著北面看.
北面就是聖殿山 金頂清真寺就在這裡後面.
金頂清真寺前面有一個叫Alakzmar 另外一個清真寺.
你看到地理上是比較高一點 比我站的地方高一點.
現在向北看 如果我繼續向東移動的時候.
你會發現我這裡看到的是鑒欖山.
鑒欖山上面有畫圖 說到耶路撒冷被毀.
這個很合理的 因為它對著耶路撒冷說.
馬卡科第24章 馬卡科第13章.
鑒欖山比我現在站的地方 拍攝影片的地方更高.
然後我再向東看的時候 你會看到山的地景比較高.
因為你不會看到山背後的東西 山背後的就是曠野到死海.
但是在大衛城 耶路撒冷是看不到曠野的.
因為有個山在擋住 然後到南面.
大衛城的南面 下面就有一些考古的法國等等.
是屬於耶布斯人的時期.
這裡就是我所說的公園 還有聯合國的總部.
聯合國在以色列的總部 就在這裡上面.
這裡看過來就是往下.
我們現在錄影的地方就是往上.
所以南面也是有山峰阻礙著.
然後我們再往西去看 是不是又高一點.
這裡有個天主教教堂 是紀念聖母瑪利亞的.
也有個教堂的鐘塔.
這裡就是所謂的西山.
我們站的這個地方就是西山.
我們再回到360度 回頭的時候.
你就會看到 回到北面.
又是回到青春節.
所以東南西北都是有山環繞著.
這個正是詩篇125篇裡面所說的.
說什麼呢 以靠耶和華的好將石安山.
坐穩坐定永不動搖.
眾山怎樣圍繞耶路撒冷.
他就是在說本地的地理.
本地的地理就是要有東西圍著它.

$^{41}$你才能夠說到這件事.
所以無論誰寫詩篇125篇.
一定是很熟悉當地的地理.
如果不是的話就寫不出來.
眾山怎樣圍繞耶路撒冷.
他就看到地理是這樣.
然後又從地理引伸出神學.
眾山怎樣圍繞耶路撒冷.
他就不是停在那裡.
耶和華也好像這樣圍繞他的百姓.
從今時時到永遠.
山是動不了的 大家都知道.
甚至現代科技可以把山爆破.
或者建造隧道.
我們都知道山是動搖不了的.
所以他說在耶路撒冷.
在石安山.
看到眾山環繞的時候.
他就發現如果我們去依靠耶和華的時候.
我們都可以安穩坐定.
這個就是地理可以幫助我們.
他就是在說這一點.
眾山環繞耶路撒冷.
東南西北都是環繞著大衛城.
原本的石安.
不單止是這樣.
他說依靠耶和華的好象石安山.
必不動搖.
永久安居.
眾山怎樣環繞耶和華.
也照樣環繞他的百姓.
從今時時到永遠.
必不動搖.
山當然不動搖.
大家都知道.
但這個永不動搖.
在舊約很有趣.
因為在不同的地方都會形容一些東西會動搖.
有些什麼東西會動搖呢.
我們看一些例子.

$^{81}$以賽亞書四十章十八節.
耶和華對那些被露的神的子民.
說要帶他們回歸.
以賽亞書四十章開始.
就是以賽亞書的下一半.
安慰耶路撒冷和神的子民.
十八節的時候.
他就說神是多麼偉大.
神是不可以和偶像假神相比.
究竟你們將誰給神.
用什麼形象與他相較呢.
至於偶像.
像人去鑄造他.
銀像用金紙包裹他.
又為他鑄造銀鏈.
沒有能力捐獻的人.
就挑選不易扭壞的木頭.
為自己尋找巧像.
樹立不會倒的.
No Emote 偶像.
會倒的偶像.
No Emote.
但誰使他不倒呢.
就是人使他不倒.
其實以賽亞取笑.
這是文化背景.
或者考古背景.
取笑什麼呢.
偶像 假神.
要放偶像在他面前參拜.
但偶像會如何不倒呢.
就是人使他不倒.
使他不動搖.
偶像才能站得住.
這是取笑偶像的方法.
是一種羞恥.
你看他們多容易動搖.
因為人使他能站立.
如果不是人做的事.
偶像就會跌.

$^{121}$為什麼要敬拜他呢.
絕對不能與他.
與真正的創天造地相比.
然後41章又繼續說.
木象鼓勵銀像.
用錘子打光.
鼓勵家賤.
對看工說.
看得好.
因為將銀,金.
看上去偶像身上.
用釘子釘穩.
免得他倒下.
所以以賽亞就是.
用拜假神.
做偶像的工匠.
去取笑他們.
要很小心地看好東西.
使偶像不會跌倒.
這就是巴力的神像.
很小的.
幾吋高.
你看到巴力神像.
是看了金上去.
在巴勒斯坦地.
是很重要的神明.
因為巴力神.
是雷電風雨神明.
巴力是帶來雨水.
所以你看到巴力神像.
兩隻手拿著的東西.
已經不見了.
但右手在很多巴力神話.
巴力圖像等等.
他拿著什麼呢.
就是拿著雷電.
拿著一把刀.
是曲的刀.
代表是雷電.
即是說他伸手.

$^{161}$能夠將雷電打下來.
有雷電就會有雨水.
所以巴力是風雨神明.
但你看到下面是怎樣.
這個出土的巴力像.
腳下面有一個釘的地方.
釘住它.
讓它不會動.
這就是以賽亞所說.
那些假神.
他們為何會去拜他們.
為何會以他們為安身立命之道.
根本沒有可能.
因為他們自己都站不穩.
是人使他們站穩.
另外一些猶大處像.
Judean Pillar Figurine.
可能是一些女神.
生育的神.
學者都在猜測.
她本身的功能是什麼.
很小的 只有幾吋高.
一個用泥或瓦去造的.
這些假神 偶像.
這些偶像可能是亞瑟拉.
可能是一些婦女.
她們要生育的時候.
她們有一些辟邪的偶像等等.
可能有不同的功用.
怎樣都好.
這裡想告訴你.
做法也是要讓它站穩.
所以叫做Pillar Figurine.
叫做處狀偶像.
為什麼呢.
因為這些偶像在以色列的地方.
猶大國的地方找到很多.
被擄之前 猶大國亡國之前.
特別多在耶路撒冷找到.
幾百個在耶路撒冷不同家居的地方找到.

$^{201}$可以看到那些人去拜這些假神.
依靠虛妄的東西.
你會看到舊約的先知.
去斥責這些以色列人去背道.
去敬拜這些假神.
巴力 亞瑟拉等等.
不去敬拜耶和華.
次生育的是誰 是耶和華.
是創造力的主 不是這些偶像.
但你會看到這些偶像的做法.
就像以賽亞所說.
No Emotions 不動搖.
為什麼呢.
因為人做了一些事 使他不動搖.
這裡有另外一些.
也是迦林地出土的.
這些偶像.
這裡就讓我們看到一個模型.
讓他大量生產.
人們就買這些偶像.
去拜也好 拿著祈禱也好.
詩篇125篇就是在說.
真正不倒的東西是什麼.
依靠耶和華的人.
是永不動搖.
永久可以長居.
為什麼呢.
因為好像種生去到.
環繞耶和華.
好像山崗一樣.
是永遠都不會動搖的.
依靠耶和華的人.
就會是這種人.
這首詩篇就是挑戰我們.
到底我們是否真心真意去依靠耶和華.
就是用一個地理的方法.
來表述出來.
讓我們能夠更加清楚.
知道這篇詩篇所說的神話訊息.
到底是什麼.

\newpage



\section{39:1-13}
\label{sec:vWnZGyYbGoI}
\textbf{How Do We Suffer Differently as God's People (Psalms 39\_1-13) - Chee Yie Jong}
\newline
\newline
連結: \href{https://youtube.com/watch?v=vWnZGyYbGoI}{\texttt{https://youtube.com/watch?v=vWnZGyYbGoI}} ~~~~ 語音日期: 2024-04-09
\newline
\newline
\hyperref[sec:_tKyeazYj4I]{\small{< < < PREV SERMON < < <}}
~
\hyperref[sec:index_chronic]{\small{[返順時目]}}
~
\hyperref[sec:index_scriptual]{\small{[返順卷目]}}
~
\hyperref[sec:u_YoyAgPSkE]{\small{> > > NEXT SERMON > > >}}
\newline
\newline
$^{1}$(audience clapping).
- Hello, hello..
Okay, good morning, UV..
Thank you for welcoming me to your church and your community..
It was a pleasure to spend some time.
with the kids this morning and to be able to,.
yeah, speak from the word of God with you all today..
So if you have some non-believer friends or family.
and you have tried to share the gospel with them,.
I'm sure at some point you will run into some questions.
along the line of, if there is a God,.
or if God is really good, why is there so much suffering.
and evil in our world?.
Well, for some people, this are maybe just a way.
for them to discredit the Christian God,.
but I feel like for the rest of the people,.
this is a legitimate and reasonable questions.
just coming from either their personal experience.
or just the things that they have seen around the world..
And if I'm just being honest,.
I feel like even a lot of times as Christians,.
we don't really have any adequate answers.
to this question as well..
And living in a big city like Chicago,.
if we just pay attention or open our eyes.
to our surrounding, it is not hard for us.
to see a lot of brokenness, injustice, and suffering.
that is in such close proximity with us, right?.
And at times like this, these questions of evil.
and suffering become even more relevant..
So my questions for today is, as the people of God,.
how should we face the suffering and pain in our lives?.
As believers, how do we suffer differently.
than the rest of the world?.
So today, we're going to take a look at Psalm 39.
and see how it approaches these questions..
I will be reading from the NIV version,.
so if you have a different version,.
you can feel free to follow along on the screen behind me..
So Psalm 39 is attributed to David,.

$^{41}$but the actual circumstances behind these Psalms.
is not clear..
But what we can know is this,.
this is an individual lament to the Lord,.
and this is a Psalm that is basically lament,.
pouring out a complaint to the Lord..
So let's look at Psalm 39 together..
(pages rustling).
Verse one..
I say, I will watch my ways and keep my tongue from sin..
I will put a muzzle on my mouth.
while in the presence of the wicked..
So I remain utterly silent, not even saying anything good..
But my anguish increases..
My heart grew hot within me..
While I meditated, the fire burned..
Then I spoke with my tongue..
So before we can answer the questions.
of how, as the people of God,.
we can suffer differently than the rest of the world,.
we first need to come to terms that, as believers,.
we are actually not exempt from suffering in this life..
Again, as believers,.
we are not exempt from suffering in this life..
In verse one, we see that the Psalmist,.
he chooses to remain silent and to restrain himself,.
and this is for two reasons..
First, so that he does not sin against the Lord..
And secondly, because the evil and the wicked are present..
Again, the actual circumstances behind the Psalmist.
is not clear, but it is very possible.
that these wicked people have somewhat contributed.
to his suffering as well..
And one of the realities of living in a broken.
and fallen world is that even sometimes.
our best attempt to do the right thing.
might actually bring us undesirable results..
Maybe some of you here,.
you might lose a job promotion at your work.
because you choose to hold on to your Christian ethics.

$^{81}$or Christian morals..
Or maybe for some of you,.
you might become a target or outcast.
in different social settings because, once again,.
you choose to remain outspoken about your faith..
And maybe for some of you, you might lose a friendship.
because you choose to speak the truth and remain truthful..
And this is exactly what the Psalmist is experiencing here..
Despite his best effort to do the right thing,.
which in this case is to restrain himself.
and to not speak up,.
we see that his situation actually become worse.
and that the anguish increase within him,.
his anguish intensify..
So, can you sense the frustration.
that the Psalmist is experiencing right here?.
And can you think of similar situation.
in your personal life as well?.
Maybe you know that you are doing the right thing..
You know that you are doing the way.
that God would want you to do,.
but the outcome of it might just not be favorable.
towards you..
And as a result of that, we see that the Psalmist,.
he could no longer hold back his silence.
and he finally decide to speak up..
So, listen to what he has to say.
when he finally decide to speak up..
So, verse four..
Show me, Lord, my life's end and the number of my days..
Let me know how fleeting my life is..
You have made my days a mere hand breath..
The span of my years is as nothing before you..
Everyone is but a breath, even those who seem secure..
Surely everyone goes around like a mere phantom..
In vain they rush about, heaping up wealth.
without knowing whose it will finally be..
And so, let us just pause here for a moment..
This three verse that we have just read through,.
the imagery that if you just try to like.

$^{121}$paint a picture of it, it is actually kind of bleak.
and in some ways hopeless, right?.
And so, again, before we can answer the questions.
of how as the people of God we can suffer differently,.
we need to acknowledge that suffering and death.
is part of our human experience,.
living in a broken and fallen world..
So, what is the Psalmist's request here.
to the Lord in verse four?.
He asks, "Show me, Lord, my life's end.
"and the number of my days..
"Let me know how fleeting my life is.".
In some of the more recent scholarship,.
verse four is interpreted not as a request.
to know how brief his human life is,.
because you already know and acknowledge that.
in verse five and six..
But right here in verse four,.
this is actually a request to know.
how long his suffering would last..
Again, because he understands that suffering and death.
is just part of human experience, living in this world..
And the urgency of this request.
to know how long his suffering would last.
is then further reinforced by his own acknowledgement.
in verse five and six that,.
"Yes, surely, life is indeed brief,.
"and even our human effort can be meaningless as well.".
So, our English translation might not capture this well,.
but in these three verses,.
there's a Hebrew word that is being repeated twice.
that is the Hebrew word hevel..
It is being translated as mere hand breath in verse five,.
and in verse six, it is being translated as in vain..
So, this word hevel is actually a major theme.
in the wisdom literature of Ecclesiastes..
And in a more literal sense,.
it just means something like a breath or a vapor..
And in a more metaphorical sense,.
it can refer to something that is fleeting or elusive..

$^{161}$So, in the book of Ecclesiastes,.
with everything under the sun being deemed as hevel.
or being deemed as meaningless,.
the author actually concludes at the end of the book.
that, "The end of the matters, all has been heard..
"Fear God and keep his commandment,.
"for this is the whole duty of man.".
In other words, the author of Ecclesiastes.
is trying to say is that,.
"The only things that truly give meaning and purpose.
"to this life is a life that is live with God,.
"a life that is live with God,.
"that is fear God and keep his commandment.".
So, similarly here in Psalm 39,.
we see that by acknowledging.
how short and transient human life is,.
the psalmist is actually agreeing.
with the book of Ecclesiastes that,.
"Yes, life apart from God is indeed meaningless or hevel.".
And we will actually see this being more clearly.
pan out in the second half of the Psalms..
So, now knowing that sin, suffering, and death.
are just inevitable part of our human life,.
where does this leave us as the people of God?.
Where does this leave us as believers?.
And what is the main thing that actually differentiate us.
than the rest of the world in the way that we suffer?.
If we continue to look in verse four,.
I think the psalmist actually provide us with some clues..
So, look with me in verse seven..
So, verse seven,.
"But now, Lord, what do I look for?.
"My hope is in you.".
So, yes, we know that sin and suffering.
are just inevitable part of our human life.
because we live in a fallen world..
However, as the people of God,.
we do not have to face these realities.
with despair or hopelessness..
As the people of God,.

$^{201}$we have a greater hope in the eternal God himself..
So, today, what I really want to press in on our heart.
is this, that even in the midst of our sin and suffering,.
we can hope in the Lord..
Even in the midst of our sin and suffering,.
we can hope in the Lord..
And this is a difference, brothers and sisters..
This is the thing that actually set us as believer,.
as people that believe in the Lord,.
that believe in Jesus..
This is the thing that sets us apart.
from the rest of the world in the way that we suffer.
because with God as our eternal hope,.
we can put our hope in him,.
even in the midst of our sin and suffering..
And I want to be clear right here..
The way that the psalmist used the word hope in verse seven.
is probably very different than the way.
that we use the word hope in our current day and age..
Maybe some of you here are students.
and or think back to your time as a student,.
you might completely butcher your final exams.
and when it is time to get your grade back,.
you might use the word hope in this sense of like,.
"Oh man, I just hope that I can pass this test.
"so that I don't have to retake the entire class.".
Or maybe some of you are new parents.
that have some younger kids at home.
and you just haven't been able to sleep.
for the past, I don't know, week or two weeks,.
and you might use the word hope in this sense of,.
"Oh, I just hope that I can get.
"maybe just four hours of sleep tonight.".
Again, this is not what the word hope means here.
and this is not how the psalmist used the word hope.
in verse seven here..
In contrast to the meaningless or the heavel of life.
that he talks about in the previous verse,.
right here in verse seven,.
the object of the hope is actually anchored.

$^{241}$in the one that is from everlasting to everlasting..
The object of our hope is anchored.
in the one who never fell on his promises.
and the one who will never change..
In other words, the object of our hope.
is anchored in the eternal God himself,.
just as the psalmist confidently proclaimed,.
"My hope is in you, Lord.".
So, what are some of the grounds.
of our confidence in the Lord.
and why is he our only hope as the people of God?.
So, let's follow along the psalmist.
as he pour out his request to the Lord..
Verse eight..
Save me from all my transgression..
Do not make me the scorn of fools..
I was silent, I would not open my mouth,.
for you are the one who has done this..
Remove this scourge from me..
I am overcome by the blow of your hand..
When you rebuke and discipline anyone for their sin,.
you consume their wealth like a moth..
Surely, everyone is but a breath..
Hear my prayer, Lord..
Listen to my cry for help..
Do not be deaf to my weeping,.
for I dwell with you as a foreigner,.
a stranger, as all my ancestors were..
Look away from me, that I may enjoy life again.
before I depart and am no more..
So, the biggest reason that we can put our hope in the Lord.
and that he is our only hope.
is because ultimately, he is the only one.
who can deliver us from our sins..
So, more than the threats of the wicked.
that we have seen in verse one.
and the threat that comes from the outside,.
the psalmist understand that ultimately,.
what he needs to be delivered from.
is something that is within him, that is his own sins..

$^{281}$And again, I feel like a lot of times,.
speaking for myself, when things doesn't go right,.
it is easy, again, for me to point my fingers.
at things that is outside of me,.
maybe my circumstances, different people,.
different relationships..
But how often do we actually take a look within ourself.
and to examine our own heart?.
I think this is a good practice that we should all,.
yeah, be aware of..
A lot of times, instead of blaming the things.
that is outside of us, that is not out of control,.
we also have to take a look within ourself as well..
So, yes, we can hope in the Lord.
even in the midst of our sin and suffering.
because he is the only one that can deliver us..
But I also want you all to hear this..
God's deliverance does not necessarily.
preclude divine discipline..
Or sometimes, the way that God delivered us.
is actually through his divine discipline.
or his divine wrath..
This is because the divine wrath and discipline of the Lord.
is actually a means of grace that is with the purpose.
to bring us into an awareness of our own sins,.
of our own shortcoming, so that we can come before the Lord.
in repentance, in confession, asking him for forgiveness..
In other words, it is actually for our own good.
that the Lord has disciplined us.
and to deliver us from our own sins..
And again, this is what the psalmist is experiencing.
right here from verse nine to 11..
So, having the awareness of the futility of life.
and also the awareness of his own sins,.
it prompts the psalmist to cry out one last time.
to the Lord..
He say, "Hear my prayer, Lord..
"Listen to my cry for help..
"Do not be deaf to my weeping,.
"for I dwell with you as a foreigner,.

$^{321}$"a stranger, as all my ancestors were.".
So right here, again, he cried out to the Lord,.
asking the Lord to hear his prayer,.
do not be deaf to his cry,.
and to not be silent to his tears..
And what is the reason of that?.
Because he understands that, he say,.
"For I dwell with you as a foreigner,.
"a stranger, as all my ancestors were.".
And by comparing himself to a foreigner or a stranger,.
the psalmist understand that,.
ultimately, his final destination is not this short.
and transient life on this earth,.
but he's actually on his way to a life.
that will last for eternity..
He understands that the power of sins over him.
is only temporary,.
because one day he will not only be free.
from the power and the penalty of sins,.
but when he spend eternity with the Lord,.
he will be free from the presence of sins forever as well..
In other words, he knows that he has a place.
in the family of God..
This is why he emphasized to the Lord that,.
"I dwell as a foreigner with you..
"I dwell with the Lord as a stranger.".
And if we just continue, in verse 13,.
we see that the psalm just sort of ends abruptly,.
with the psalmist asking the Lord.
to turn away his angry gaze from him.
so that he is able to enjoy the remainder of his life..
And we actually find no resolution.
or no response from the Lord in the psalm itself..
So, we've come to the end of these psalms,.
but I want to turn our attention to this..
The silence that the psalmist is facing from the Lord,.
pouring out his request before the Lord,.
but hearing no response, at least within the psalmist..
I think the psalmist is actually not alone.
in facing this deafening silence from the Lord..

$^{361}$I know last week we observed Good Friday.
and celebrate Easter Sunday..
So, we're probably very familiar with this..
Right before Jesus' death at the cross,.
he pours out the cry quoted from Psalm 22..
"My God, my God, why have you forsaken me?".
And I feel like a lot of time we can just.
kind of like brush off the magnitude.
of this cry from Jesus..
But this person that make this cry at the cross,.
this is the outcry of Jesus, the Son of God,.
bearing the sins of the world.
and facing kind of like the wrath.
from his own Father as well..
But we know that as the obedient Son of God,.
Jesus understands that the only way.
to resurrection and glorification.
is actually through his suffering at the cross..
So, speaking about this, the book of Hebrews.
described that Jesus, as the author.
and perfecter of our life,.
"For the joy sets before him,.
"he endured the cross, despising its shame.".
And so Jesus' example actually becomes a paradigm.
for all of us as believers on how we are.
to make our journey on our life on this earth.
with all his suffering, with all his pain..
And we are able to endure our journey through this life.
because Jesus not only provide us with the way.
through his death and resurrection at the cross,.
but Jesus actually also sets an example of endurance.
for us to follow..
And this is why he is our only hope.
in the midst of our sin and suffering..
The reason that Jesus is our only hope.
in the midst of our sin and suffering.
is because he is the only one not just providing.
the means for us to be reconciled to God.
through his death and resurrection,.
but he's also the only one that lives a perfect life.

$^{401}$for us to follow..
So from the time that we put our faith in Jesus.
to the time that Jesus comes again,.
Jesus basically model for us on how we are.
to live our life..
And he is the only perfect example for us..
So maybe some of you here,.
you're in a season where you're just standing face to face.
with suffering or just with a lot of brokenness,.
whether it is physically, relationally,.
emotionally, or mentally..
Or maybe are you in a season.
where the realities and pressure of life.
just feels so overwhelming..
Or maybe like the psalmist here in Psalm 39,.
you are currently just facing the deafening silence.
from the Lord, despite countless days and nights.
of prayers and fasting..
And if I'm being honest, I am familiar with this currently..
Or perhaps there's just a particular sin.
or addiction in your life that you have prayed.
and asked the Lord to kind of like remove you from,.
but you're just struggling to overcome..
Of course, I want to acknowledge.
that all of these experience are true and real,.
and I don't view them lightly..
But maybe, just maybe at the same time,.
this might be an opportunity for us.
to learn how to hope in the Lord..
How to put our faith and trust in Jesus..
Again, because our faith does not just get perfected.
in an instance..
To be able to hope in the Lord.
actually requires a lot of practice..
It requires constantly coming before the Lord,.
asking Him, and in some ways, when He responds to us,.
that are the basis of our future hope in Him as well..
So once again, before I end,.
I want to encourage you all with the truth.
that the psalmist here in Psalm 39.

$^{441}$that even in the midst of our sin and suffering,.
we can put our hope in the Lord..
So I'm gonna pray us, so let's bow our head in prayers..
Heavenly Father, we just want to give our praise.
and give our thanks to you..
For you are our hope,.
and we are grateful to you..
For you are the only one that we can put our hope in,.
and the only one that would never fail us..
Jesus, we praise you,.
for you are the only one that provide the means.
for us to be able to reconcile to the Father again..
And you're not only a God, but you are fully human as well..
A great high priest who is able to empathize.
with our weakness,.
because you have been tended in every way,.
just as we are, yet without sin..
So even when we face just a different suffering.
and challenges that life thrown at us,.
we know that Lord, we know that Jesus,.
you know how he feels,.
and you are with us in the midst of our sin and suffering..
Holy Spirit, we thank you for,.
you are the living and active spirit living in us..
You intercede for us when we do not know.
what we ought to say or what we ought to pray for..
And yet, you know our hearts,.
and you pour out your prayers on our behalf before the Father..
Lord, would you just teach us.
to learn how to put our hope in you.
in a way that honors you,.
in a way that brings glory to you,.
in a way that show the broken world.
that our God is a God that is worthy of all the praise.
and all the glory and all the honor..
So we pray this in your son's name, amen..
(piano music).
\newpage



\section{約翰福音 4:1-42}
\label{sec:u_YoyAgPSkE}
\textbf{打破性別定型的女人 (約翰福音 4\_1-42) - 何建邦傳道}
\newline
\newline
連結: \href{https://youtube.com/watch?v=u_YoyAgPSkE}{\texttt{https://youtube.com/watch?v=u\_YoyAgPSkE}} ~~~~ 語音日期: 2024-04-09
\newline
\newline
\hyperref[sec:vWnZGyYbGoI]{\small{< < < PREV SERMON < < <}}
~
\hyperref[sec:index_chronic]{\small{[返順時目]}}
~
\hyperref[sec:index_scriptual]{\small{[返順卷目]}}
~
\hyperref[sec:HYnhtr6gM20]{\small{> > > NEXT SERMON > > >}}
\newline
\newline
約翰福音 4:1-42
\newline
\begin{longtable}{cl}
\hline
\hline
章節 & 經文 (和合本修訂版)\\
\hline
4:1 & \begin{tabularx}{0.7\textwidth}{X} 耶穌知道法利賽人聽見他收門徒和施洗比約翰還多， （ \end{tabularx} \\ \\ \relax
4:2 & \begin{tabularx}{0.7\textwidth}{X} 其實不是耶穌親自施洗，而是他的門徒施洗，） \end{tabularx} \\ \\ \relax
4:3 & \begin{tabularx}{0.7\textwidth}{X} 他就離開猶太，又回加利利去。 \end{tabularx} \\ \\ \relax
4:4 & \begin{tabularx}{0.7\textwidth}{X} 他必須經過撒瑪利亞， \end{tabularx} \\ \\ \relax
4:5 & \begin{tabularx}{0.7\textwidth}{X} 於是到了撒瑪利亞的一座城，名叫敘加，靠近雅各給他兒子約瑟的那塊地。 \end{tabularx} \\ \\ \relax
4:6 & \begin{tabularx}{0.7\textwidth}{X} 雅各井就在那裡；耶穌因旅途疲乏，坐在井旁。那時約是正午。 \end{tabularx} \\ \\ \relax
4:7 & \begin{tabularx}{0.7\textwidth}{X} 有一個撒瑪利亞婦人來打水。耶穌對她說：「請給我水喝。」 \end{tabularx} \\ \\ \relax
4:8 & \begin{tabularx}{0.7\textwidth}{X} 因為那時門徒進城買食物去了。 \end{tabularx} \\ \\ \relax
4:9 & \begin{tabularx}{0.7\textwidth}{X} 撒瑪利亞婦人對他說：「你是猶太人，怎麼向我一個撒瑪利亞女人要水喝呢？」因為猶太人和撒瑪利亞人沒有來往。 \end{tabularx} \\ \\ \relax
4:10 & \begin{tabularx}{0.7\textwidth}{X} 耶穌回答她說：「你若知道神的恩賜，和對你說『請給我水喝』的是誰，你早就會求他，他也早就會給了你活水。」 \end{tabularx} \\ \\ \relax
4:11 & \begin{tabularx}{0.7\textwidth}{X} 婦人對耶穌說：「先生，你沒有打水的器具，井又深，哪裡去取活水呢？ \end{tabularx} \\ \\ \relax
4:12 & \begin{tabularx}{0.7\textwidth}{X} 我們的祖宗雅各把這井留給我們，他自己和兒女以及牲畜都喝這井裡的水，難道你比他還大嗎？」 \end{tabularx} \\ \\ \relax
4:13 & \begin{tabularx}{0.7\textwidth}{X} 耶穌回答，對她說：「凡喝這水的，還要再渴； \end{tabularx} \\ \\ \relax
4:14 & \begin{tabularx}{0.7\textwidth}{X} 誰喝我所賜的水，就永遠不渴。我所賜的水要在他裡面成為泉源，直湧到永生。」 \end{tabularx} \\ \\ \relax
4:15 & \begin{tabularx}{0.7\textwidth}{X} 婦人對他說：「先生，請把這水賜給我，使我不渴，也不用到這裡來打水。」 \end{tabularx} \\ \\ \relax
4:16 & \begin{tabularx}{0.7\textwidth}{X} 耶穌對她說：「你去，叫你的丈夫，再到這裡來。」 \end{tabularx} \\ \\ \relax
4:17 & \begin{tabularx}{0.7\textwidth}{X} 婦人回答，對耶穌說：「我沒有丈夫。」耶穌說：「你說沒有丈夫是對的。 \end{tabularx} \\ \\ \relax
4:18 & \begin{tabularx}{0.7\textwidth}{X} 你已經有過五個丈夫，你現在有的並不是你的丈夫。你這話是真的。」 \end{tabularx} \\ \\ \relax
4:19 & \begin{tabularx}{0.7\textwidth}{X} 婦人對他說：「先生，我看你是一位先知。 \end{tabularx} \\ \\ \relax
4:20 & \begin{tabularx}{0.7\textwidth}{X} 我們的祖宗在這山上敬拜神，你們倒說，應當敬拜的地方是在耶路撒冷。」 \end{tabularx} \\ \\ \relax
4:21 & \begin{tabularx}{0.7\textwidth}{X} 耶穌對她說：「婦人，你要信我。時候將到，你們敬拜父，既不在這山上，也不在耶路撒冷。 \end{tabularx} \\ \\ \relax
4:22 & \begin{tabularx}{0.7\textwidth}{X} 你們所敬拜的，你們不知道；我們所敬拜的，我們知道，因為救恩是從猶太人出來的。 \end{tabularx} \\ \\ \relax
4:23 & \begin{tabularx}{0.7\textwidth}{X} 時候將到，現在就是了，那真正敬拜父的，要用心靈和誠實敬拜他，因為父要這樣的人敬拜他。 \end{tabularx} \\ \\ \relax
4:24 & \begin{tabularx}{0.7\textwidth}{X} 神是靈，所以敬拜他的必須用心靈和誠實敬拜他。」 \end{tabularx} \\ \\ \relax
4:25 & \begin{tabularx}{0.7\textwidth}{X} 婦人對他說：「我知道彌賽亞—就是那稱為基督的—要來；他來了，會把一切的事都告訴我們。」 \end{tabularx} \\ \\ \relax
4:26 & \begin{tabularx}{0.7\textwidth}{X} 耶穌對她說：「我就是，正在跟你說話呢！」 \end{tabularx} \\ \\ \relax
4:27 & \begin{tabularx}{0.7\textwidth}{X} 正在這時，門徒回來了。他們對耶穌正在和一個婦人說話感到驚訝，可是沒有人說：「你要甚麼？」或說：「你為甚麼和她說話？」 \end{tabularx} \\ \\ \relax
4:28 & \begin{tabularx}{0.7\textwidth}{X} 那婦人留下水罐，往城裡去，對眾人說： \end{tabularx} \\ \\ \relax
4:29 & \begin{tabularx}{0.7\textwidth}{X} 「你們來看！有一個人把我素來所做的一切事都說了出來，難道這個人就是基督嗎？」 \end{tabularx} \\ \\ \relax
4:30 & \begin{tabularx}{0.7\textwidth}{X} 他們就出城，來到耶穌那裡。 \end{tabularx} \\ \\ \relax
4:31 & \begin{tabularx}{0.7\textwidth}{X} 就在這個時候，門徒求耶穌說：「拉比，請吃吧。」 \end{tabularx} \\ \\ \relax
4:32 & \begin{tabularx}{0.7\textwidth}{X} 耶穌對他們說：「我有食物吃，是你們不知道的。」 \end{tabularx} \\ \\ \relax
4:33 & \begin{tabularx}{0.7\textwidth}{X} 門徒就彼此說：「難道有人拿甚麼給他吃了嗎？」 \end{tabularx} \\ \\ \relax
4:34 & \begin{tabularx}{0.7\textwidth}{X} 耶穌對他們說：「我的食物就是要遵行差我來那位的旨意，完成他的工作。 \end{tabularx} \\ \\ \relax
4:35 & \begin{tabularx}{0.7\textwidth}{X} 你們不是說『到收割的時候還有四個月』嗎？我告訴你們，舉目向田觀看，莊稼熟了，可以收割了。 \end{tabularx} \\ \\ \relax
4:36 & \begin{tabularx}{0.7\textwidth}{X} 收割的人已經得工錢，為永生儲存五穀，使撒種的和收割的一同快樂。 \end{tabularx} \\ \\ \relax
4:37 & \begin{tabularx}{0.7\textwidth}{X} 『那人撒種，這人收割』，這話可見是真的。 \end{tabularx} \\ \\ \relax
4:38 & \begin{tabularx}{0.7\textwidth}{X} 我差你們去收你們所沒有辛勞的；別人辛勞，你們享受他們辛勞的成果。」 \end{tabularx} \\ \\ \relax
4:39 & \begin{tabularx}{0.7\textwidth}{X} 那城裡有好些撒瑪利亞人信了耶穌，因為那婦人作見證，說：「他把我素來所做的一切事都說了出來。」 \end{tabularx} \\ \\ \relax
4:40 & \begin{tabularx}{0.7\textwidth}{X} 於是撒瑪利亞人來見耶穌，求他在他們那裡住下，他就在那裡住了兩天。 \end{tabularx} \\ \\ \relax
4:41 & \begin{tabularx}{0.7\textwidth}{X} 因為耶穌的話，信的人就更多了。 \end{tabularx} \\ \\ \relax
4:42 & \begin{tabularx}{0.7\textwidth}{X} 他們對那婦人說：「現在我們信，不再是因為你的話，而是我們親自聽見了，知道這人真是世界的救主。」 \end{tabularx} \\ \\
[1ex]
\hline
\hline
\end{longtable}
$^{1}$.
各位弟兄姊妹大家好.
聽到我說話了嗎?.
不知道大家還記不記得我.
其實我在大概兩年還是三年前.
來過UV講過一次到的.
那次我是講馬太福音的第一章.
好像是講耶穌基督的奇利加普.
我的題目是.
今天我就不講馬太福音了.
今天我就想講約翰福音.
今天我們看的經文會比較長.
但不需要擔心.
我是不會看完所有經文.
我今天要看的經文是約翰福音第四章一至四十二節.
大家記不記得約翰福音四章一至四十二節講什麼?.
就是講一個女人.
所以今天我的題目就是.
打破性別定型的女人.
就是講約翰福音四章的這個女人.
我想講的是什麼呢?.
通常華人教會去講約翰福音四章的時候.
我們的關注焦點在哪裡?.
通常華人教會的關注焦點.
就不在這個女人那裡.
是在一些屬靈的議題那裡.
例如是什麼?.
用心靈和誠實去敬拜上帝.
究竟哪裡才是敬拜的地方?.
是聖殿?是基里森山?.
還是哪裡呢?.
耶穌就說要用心靈誠實去敬拜上帝.
我們華人教會去看的時候.
就會看很多這些屬靈的東西.
但我想講的是什麼呢?.
約翰福音四章是關於一個很重要的女人.
一個撒瑪利亞的婦人.
今天我想跟大家.
不是用這些所謂的屬靈的角度去看這段經文.
我不是說這些屬靈的議題是錯的.

$^{41}$不過我想試一下跟大家.
用另一個角度,另一個觀點去看這段經文.
我想試一下用一個性別的角度.
性別議題的角度去看這段經文.
看一下究竟約翰是怎樣描述這個女人的.
華人教會是比較少用性別議題的角度去看這段經文.
或者去看任何經文.
但其實這段經文對於性別的強調是非常明顯的.
因為如果大家記得這段經文有個背景是有講的.
就是耶穌和這個女人去聊天的時候.
門徒去了哪裡?去了買東西吃,是不是?.
當門徒買完東西吃回來之後.
見到耶穌的時候,他很詫異.
為什麼詫異?.
是因為他們詫異.
耶穌竟然會和一個女人去聊天.
所以在四章二十七節那裡.
經文是很強調性別這個話題.
是很強調門徒詫異耶穌和一個女人去談話.
所以作者似乎是透過門徒的詫異.
讓讀者去注意到.
撒瑪利亞婦人不單止作為一個外邦人.
不單止作為一個撒瑪利亞人的身份.
而是作為一個女性的身份.
在這個故事裡面去存在.
所以我們就要去看看.
今天如果我們從一個性別的角度去看這段經文.
究竟會有些什麼?.
看到些什麼?.
好,首先要去解釋一下.
如果我們要從一個性別的角度去看這段經文的時候.
我們一定要從一個舊約的切入點去看.
什麼切入點呢?.
今天很想和大家介紹.
就是一個叫舊約的井旁或者井邊的場景.
因為這段經文是呼應舊約的井旁或者井邊的場景.
什麼是井邊的場景呢?.
有三段經文是很經典的井邊場景的經文.
第一段就是《創世記》的24章1到17節.
就是在說誰呢?.

$^{81}$說亞伯拉罕的僕人為到亞伯拉罕的兒子以殺去找老婆.
最後找到李伯嘉.
第二段經文就是《創世記》29章1到12節.
說的是誰呢?.
說的是亞伯拉罕的僕人去扣府.
然後如何打稅.
就遇見他將來的一個老婆.
拉潔.
然後第三段經文就是《出埃及記》.
就是在說摩西和他老婆的故事.
我想說的是什麼呢?.
什麼叫做井邊場景呢?.
井邊場景有什麼獨特的地方呢?.
井邊場景通常是這樣的.
就是一個男性的角色.
或者男主角.
在井邊那裡.
遇見一個剛出來打稅的女人.
而這個女人就是男主角將來的老婆.
這個就叫做井邊場景.
明白了嗎?.
這是一個很經典的場景.
在舊的三段經文.
就是經常會出現的一個這樣的場景.
所以如果擺現代的話.
一個愛情的故事就是.
當你在水機那裡倒水的時候.
你就可能會遇見你未來的老婆.
所以大家單身的弟兄.
待會在外面倒水的時候要小心點留意.
誰來倒水.
那個可能是你將來的老婆.
就好像這個井邊場景的現代版一樣.
所以這個就是舊約的井邊場景.
其實是和愛情有關.
和結婚有關.
和男主角和女主角有關.
我想說的是.
其實這一段經文是在呼應舊約的井邊場景.
為什麼呢?.

$^{121}$就是經文的一開始.
交代背景的時候說得很清楚.
就是怎麼說呢?.
耶穌因為旅途疲勞.
就坐在井邊那裡.
那時候大概是第六時.
第六時就是正午的時間.
就是太陽最猛烈的時間.
就坐在井邊那裡.
就有一個撒瑪利亞婦人來倒水.
很明顯.
井邊這裡是很明顯出現了.
耶穌坐在那裡.
然後有一個女人出來倒水.
然後如果我們再看上文的話.
我們就更加會發覺.
作者似乎刻意想我們聯想到井邊場景.
為什麼呢?.
就是因為在上文是說什麼呢?.
上文的約翰福音三章27到30節那裡.
是在說施洗約翰.
在說關於耶穌的說話.
他那裡形容耶穌做什麼呢?.
施洗約翰說.
我不是那個新郎.
我只不過是新郎的朋友.
耶穌.
耶穌才是那個要娶新郎的新娘.
耶穌才是那個要娶新娘的新郎.
差點說了新郎.
新郎的新郎.
他要娶新娘的新郎.
耶穌才是那個.
所以耶穌才是那個將要娶妻的人.
所以如果我們這樣看上文的話.
其實當你作為一個讀者.
你看了上面.
施洗約翰形容耶穌.
是一個將要娶妻的新郎的時候.
然後你看到第四章一開始.

$^{161}$耶穌就坐在井邊那裡.
然後就有一個女人.
撒瑪利亞的女人出來打水.
你作為當時的猶太讀者.
或者熟悉舊約的讀者.
你很容易很自然就會聯想.
啊!這是什麼?.
井邊場景了.
你明白我的意思吧?.
所以這裡就很刻意地去呼應井邊的場景.
我們就去看看.
究竟舊約的井邊的場景.
是如何形容女角色或者女主角.
而在約翰福音又如何形容女主角撒瑪利亞的婦人呢?.
我們只針對兩段經文去看.
就是《創世記》的24章和29章.
只看這兩段井邊場景.
在《創世記》的24章16節那裡是這麼說的.
那女子容貌極其盡美.
那女子就是在說什麼?.
尼伯加.
就是那個出來打水的尼伯加.
是以撒將來的老婆.
那女子容貌極其盡美.
在《創世記》29章的17節.
那裡就在說利亞.
亞國將來的老婆.
就在說什麼?.
不是利亞,對不起.
拉潔.
拉潔就是將來亞國最喜歡的那個女人.
這裡就說利亞.
就是拉潔的姐姐.
利亞的眼睛是沒有神器的.
拉潔她是長得美貌盡秀.
那有什麼共同點呢?.
就是漂亮.
第一個特點.
在井邊場景出現的女人.
第一個特點是漂亮.

$^{201}$是很漂亮.
無論是尼伯加還是拉潔.
第二個特點.
就是在《創世記》24章16節那裡這麼說.
那女子容貌極其盡美.
還有什麼?.
還是處女.
也未曾有人親近她.
這個是在形容尼伯加的.
然後在《創世記》的29章26節那裡.
是形容拉潔的.
那時候大家熟悉的故事就是.
拉班騙了亞國.
就是先讓大姐姐娶了.
然後要她再服侍七年才讓拉潔.
記得這個故事.
那時候拉班就這麼說.
大女兒還沒有給人.
就將小女兒給人.
小女兒就是拉潔.
在我們這個地方沒有這個規例.
那這裡在暗示什麼呢?.
暗示拉潔是處女.
暗示拉潔是未曾給人.
拉潔是沒有出嫁過的.
拉潔還是貞潔的.
所以在井邊場景的女人.
第二個特點是什麼呢?.
是什麼?.
處女.
就是貞潔.
第二個特點是貞潔.
我們看看第三個特點.
在《創世記》的24章那裡.
這裡又是形容尼伯加.
亞伯拉罕的僕人跑上前去迎著尼伯加的船.
求尼將瓶裡的水給我喝一點.
女子就說.
我主請喝.
就急忙拿下瓶來.

$^{241}$托在手上給他喝.
然後就看看《創世記》的29章12節.
這裡又是形容拉潔.
雅各就跟拉潔說.
自己是他父親的外甥.
是尼伯加的兒子.
於是拉潔就跑去告訴他父親.
這裡要解釋一點.
這裡雅各其實是跟拉潔交代自己的身世.
跟拉潔交代為什麼我要來找你爸爸拉班.
於是拉潔就很乖地跑去跟他爸爸拉班說.
明白我說什麼了嗎?.
大概的邏輯是這樣.
我想問.
在這兩段經文裡面的第三個特點是什麼呢?.
有沒有人回答到?.
第三個特點.
是什麼?.
是對男性的信服.
留意到對男性的信服嗎?.
亞伯拉罕的僕人說要點水.
尼伯加就這樣說.
我主請喝.
還要好好服侍.
拿著瓶子.
用手托著給他喝.
是對男性的信服.
而在下面.
雅各一說給拉潔聽之後.
拉潔就二話不說.
馬上回家找他父親.
將雅各跟他說的話.
要告訴他父親.
以致之後拉潔就出來迎接雅各.
所以是很乖很聽話.
是對男角色的信服.
對男性的信服.
所以第三個特點是什麼?.
第三個特點就是信服男性.
我想說的是什麼呢?.

$^{281}$我想說的是.
在井邊場景所出現的女人.
是一個好女人的形象.
是完美女性的形象.
為什麼呢?.
因為男主角將來要娶這個女人.
是一個女主角.
是一個完美女性的形象.
包括漂亮,貞潔和信服男性.
我想說的這三點.
在井邊場景.
一個完美女性.
好女人的形象.
其實相當符合.
當時古代社會對於女性的要求.
或者想像.
古代社會要求女性.
好女人定義是什麼?.
就是漂亮.
漂亮是很重要.
我們每個時代都要愛貌協會.
要漂亮.
還有什麼?.
要貞潔.
貞潔很重要.
在古代.
貞潔很重要.
第三是什麼?.
在古代是什麼?.
是一個男權主義的社會.
一個老婆最重要是什麼?.
聽話.
老婆最重要是聽話.
要信服男性.
所以在井邊場景所出現的女人.
其實是古代的人在想的.
古代社會所需要的一個好女人的典範.
好女人的形象.
我們就去看看.
究竟約翰是如何描述.

$^{321}$他自己的井旁的女人.
這個角色.
撒瑪麗亞婦人.
我想說的是三點.
對應剛才的三點.
第一.
約翰有沒有描述撒瑪麗亞婦人的樣子?.
是沒有.
讀者不知道撒瑪麗亞婦人的樣子是漂亮還是不漂亮.
對於約翰來說.
她的樣子是不重要的.
約翰沒有提及過撒瑪麗亞婦人的外貌.
第二.
這個撒瑪麗亞婦人是否貞潔的女人?.
很明顯不是.
約翰還要告訴你她不貞潔.
他還要很強調她是不貞潔的.
為什麼?.
因為他那裡說到.
他在四章十八節強調撒瑪麗亞婦人曾經有五個丈夫.
而現在同居的男人不是她的丈夫.
她不僅不是處女.
而且曾經有五個丈夫.
現在還要同居.
不要說古代.
你放在現代.
你也接受不了.
你接橫教會.
你同居也不行.
真的不行.
同居也不行.
你不要說同居.
一起去旅行也不行.
在橫教會.
死定了.
我想說.
就是這樣.
教會就是這樣.
我想說.
不過我先不說教會.

$^{361}$一會兒反思再說.
回到這裡.
我回到這裡.
這裡是說.
約翰要特意強調.
她不貞潔.
她曾經有五個男人.
現在還要同居.
那個不是她的丈夫.
你明白嗎?.
四章十八節這樣強調.
還要是.
她絕對不是井旁的女人.
去馴服男性.
因為這裡記載了兩次.
撒瑪利亞婦人拒絕耶穌的請求.
第一次是在四章九節.
耶穌向她求水喝.
然後怎麼樣?.
撒瑪利亞婦人有沒有給她?.
沒有給她.
撒瑪利亞婦人沒有給水耶穌.
而是反問耶穌什麼?.
反問耶穌說.
你作為一個猶太人.
為什麼問撒瑪利亞人要水?.
很囂張吧?.
沒有給水不但囂張.
還要囂張耶穌.
你這個猶太人.
為什麼問我這個撒瑪利亞人要水?.
然後第二次是什麼?.
第二次就是.
對話去到一個段落的時候.
耶穌就說.
你回去叫你的丈夫來吧.
就是剛才說的那裡.
她曾經有五個丈夫.
現在那個不是她的丈夫.
那一段.

$^{401}$耶穌就叫她.
你回去叫你的丈夫來吧.
那撒瑪利亞婦人有沒有照做?.
沒有照做.
而是說什麼?.
我現在沒有丈夫.
現在那個跟我同居的不是我丈夫.
是不是?.
就是和拉結是很不同的.
拉結是雅各叫她回去說.
她就馬上去說.
這裡耶穌叫撒瑪利亞婦人.
回去找你的丈夫來.
撒瑪利亞婦人沒有去.
還要質問.
就是還要反駁耶穌.
現在那個不是我丈夫.
是我同居的男人.
OK?.
你明白我的意思吧?.
所以這裡我想說什麼呢?.
在這裡撒瑪利亞婦人.
在約翰的筆下.
其實是在打破當時古代對於好女人的想像.
是在打破舊約的井邊場景.
對於好女人的形象的描述.
撒瑪利亞婦人沒有一點.
是符合好女人的形象.
OK?.
你明白吧?.
她的樣子是怎樣?.
沒有人知道.
她不貞潔.
還要跟別人同居.
她不信服男性.
兩次拒絕了耶穌的請求.
那我想說的是什麼呢?.
約翰在這裡究竟想怎樣去描述這個女人呢?.
我想說的是.
約翰是特意.

$^{441}$在約翰的筆下.
他特意要打破好女人的形象.
他要重構.
重新去定義.
什麼才是好女人.
他要透過撒瑪利亞婦人去重新定義.
什麼才是好女人.
不是古代文化的定義.
而是基督教價值的定義.
什麼才是好女人.
我們再去看看.
究竟什麼才是好女人呢?.
撒瑪利亞婦人在說什麼呢?.
在約翰眼中.
撒瑪利亞婦人才是真正的女性的典範.
第一.
撒瑪利亞婦人是.
在約翰福音中其中一個.
是很稀少的角色.
能夠在約翰福音中.
成功地認出耶穌身份.
為什麼我說很稀少呢?.
如果你看回上文.
上文耶穌是跟一個角色有對話.
那個角色是一個男人.
知道是誰嗎?.
如果熟悉約翰福音.
是尼哥德姆.
在上文.
耶穌是跟尼哥德姆有一番很長的對話.
一個男性角色.
尼哥德姆有什麼特別呢?.
他是一個男性的中產階級.
因為他是一個拉比.
他是猶太人的老師.
但在上文是很明確地說到.
尼哥德姆跟耶穌的一番對話.
完了之後.
尼哥德姆是完全不知道耶穌在說什麼.
是完全認不出耶穌是誰.

$^{481}$完全是九不搭八的整個對話.
但這裡特別的地方是.
一個女人.
一個外號.
一個被邊緣化的撒瑪利亞婦人.
她能夠透過跟耶穌的對話.
逐漸明白到耶穌的身份.
我們看看為什麼這樣說.
在四章九節那裡.
撒瑪利亞婦人稱呼耶穌為猶太人.
然後四章九節就稱呼他為先生.
四章十一,十五和十九節.
就開始稱呼耶穌為先知.
到了四章二十九節.
就稱呼耶穌為基督.
彌賽亞.
你見不見到有一個演變在這裡.
撒瑪利亞婦人透過跟耶穌的對話.
她慢慢認清.
原來耶穌不單只是一個猶太人.
不單只是一個先生.
不單只是一個先知.
而是彌賽亞基督.
這是之前的很厲害的男性角色.
一個猶太的拉比都做不到的事.
但撒瑪利亞婦人能夠做到.
在《約翰福音》裡面.
很多人都很傻.
是認不出耶穌是誰.
但是撒瑪利亞婦人.
能夠透過對話認出耶穌是誰.
所以在這裡透過這個描述.
就見到撒瑪利亞婦人作為一個女性.
將尼哥德姆猶太的拉比.
和耶穌的其他門徒.
男性的門徒.
比下去.
撒瑪利亞婦人是真正透過對話.
能夠認出耶穌的人.
更加重要的是.

$^{521}$第二點更加重要的是.
撒瑪利亞婦人在整個故事裡面.
其實她最終說完這番對話之後.
她做的事.
正正就是耶穌做的事.
在這裡.
約翰是很隱晦地描述撒瑪利亞婦人.
他告訴我們撒瑪利亞婦人.
是在模仿基督.
為什麼這樣說呢.
第一點.
在四章三節那裡.
耶穌做了兩個行動.
是什麼呢.
就是耶穌離開一個地方.
離開猶太地.
然後進入撒瑪利亞婦人的地方.
離開和進入有兩個動詞.
就是離開和進入在四章三節.
我們會看到在四章二十八節.
這兩個動詞同樣都是用在撒瑪利亞婦人那裡.
撒瑪利亞婦人在四章二十八節那裡.
同樣地都是撒瑪利亞婦人離開和進入.
這裡要說一點原文.
我不會說原文出來.
但是二十八節的放下水瓶.
或者放下水罐.
放下這個字其實就是第三節.
耶穌離開的離開的動詞.
所以你看到這裡是echo.
撒瑪利亞婦人做的行動.
離開和進入.
正正就是耶穌一開始的時候.
所做的行動.
離開和進入.
第一點.
第二點.
《約翰福音》的四章四十一節.
在第四十一節那裡強調.
撒瑪利亞人是因為耶穌的話.

$^{561}$就相信耶穌.
四章四十一節強調.
撒瑪利亞人因為耶穌的話.
相信耶穌.
同樣地我們看到.
在四章的三十九節那裡強調.
是因為這個婦人的話.
撒瑪利亞人就相信耶穌.
你看到是怎樣.
又是一個對應的關係.
就是撒瑪利亞人因為耶穌的話.
同樣地也是因為撒瑪利亞婦人的話.
而相信耶穌.
明白了嗎.
第二點.
第三點就是.
在第四十四節那裡強調.
耶穌作見證.
而在第三十九節那裡.
也強調撒瑪利亞婦人作見證.
大家都作見證.
所以你看到在這裡.
透過一個這樣的描述.
我們叫做characterization.
就是角色的描繪.
透過這些這麼implicit的角色描繪.
其實約翰就將撒瑪利亞婦人描述到.
她真的是在做耶穌所做的事的人.
真正imitate Christ.
模仿耶穌學校基督的人.
並不是那些出去買東西吃的門徒.
或者回來見到耶穌很驚訝.
你為什麼跟一個女人說話的那些門徒.
真正追隨著基督.
真正仿效著基督的人.
是這個撒瑪利亞婦人.
這個女人.
所以透過一個這麼厲害的描繪.
約翰將撒瑪利亞婦人描繪成為.
仿效耶穌的典範.

$^{601}$將男性的角色尼哥德姆和門徒.
都給了進去.
在這裡對約翰來說.
女性的典範不在於漂亮.
不在於貞潔.
也不在於對男性的信服.
而是在於什麼.
在於對耶穌身份的悟性.
和仿效耶穌的行動.
這兩點就是撒瑪利亞婦人做得到的事.
今天我想說的是什麼呢.
去到反思的時間.
《約翰福音》是一個早期的基督教文獻.
如果我們看《約翰福音》四章的故事.
我們就留意到.
耶穌是沒有批判撒瑪利亞婦人.
雖然她不貞潔.
但耶穌沒有說你悔改.
不要再同居.
沒有這樣說.
你離開男人.
正正經經結婚.
耶穌不是說這些.
我們華人教會說這些.
但耶穌在這個經文裡面.
不是說這些.
在這段經文裡面.
是強調這個女人.
怎樣去認出耶穌的身份.
和怎樣去學習耶穌的行動.
我想說的是什麼呢.
有一個社會學家叫做Watney Stark.
他寫了一本書叫做.
《The Rise of Christianity》.
他有一個論點.
他的論點就是說.
早期基督教之所以能夠成功.
是因為它有強大的社群凝聚力.
和進步的人民精神.
包括對女性的尊重.

$^{641}$我這裡說進步的人民精神.
不是和現在這個時代比.
是和同期的古代社會比.
是有進步的人民精神.
基督教是比較尊重女人的.
所以這個是真的.
因為很多在早期教會裡面.
很多領袖都是什麼呢.
是女人.
很多教會領袖.
女先知,女牧者,女傳道.
在早期的教會裡面.
基本上是女性的參與.
在領袖層是很多的.
所以的確是比較進步的人民精神.
無論是對奴隸還是對女人.
在同期的社會文化裡面.
如果Watney Stark.
她說的話是對的話.
值得我們反思的一個問題是.
這個Watney Stark是說.
早期基督教成功的優勢.
是因為他們有很強的凝聚力.
和比較進步的人民精神.
我們要問的一個問題是.
今天在這個時代的基督教.
還有沒有這些呢.
我們還有沒有以前早期基督教成功的優勢.
我們是不是還有很強的凝聚力.
和比較進步的人民精神呢.
和社會比較.
嗯.
凝聚力可能還有.
不過我想說.
我不是特別批判任何教會.
不過我純粹想拿這個反思出來.
大家一起想想.
今天.
以前的基督教.
是走在主流社會的一個解放的那邊.

$^{681}$他們是要對奴隸.
不要壓制奴隸.
你知道基督教是對奴隸好一點.
對女性好一點.
尊重一點.
是走在社會解放那邊的.
在古代.
今天基督教是站在哪邊.
我們的基督教是站在什麼立場呢.
我們還有沒有以前的優勢呢.
我不說那麼多教會.
就是說美男浸信會.
大家知不知道有沒有看基督教的新聞.
美男浸信會經常都有一個新聞.
就是除了性侵之外.
除了性侵這件事之外.
就是有些堂妹.
安納了女性牧師.
然後怎樣.
就被人逐出會.
是不是.
喂.
古代的基督教是怎樣的.
古代的基督教是女性.
做領袖的角色.
我只是純粹說美國.
我不是不說香港.
美國的福音派教會.
其實是怎樣的基督教呢.
回到我們自己.
其實回到華人教會.
其實我們華人教會.
我們很多時候以為我們自己的教導是出於聖經.
但其實我們都是.
受到我們自己華人的文化去影響.
例如是甚麼呢.
男孩子就應該喜歡藍色.
女孩子就應該喜歡粉紅色.
這個不是作出來的.
我只是聽過.

$^{721}$我以前是做青少年工作的.
就是牧養一些青少年.
有些青少年人是會.
就是男的青少年人.
男孩子是喜歡粉紅色的.
但是他們小時候就會被爸爸媽媽說.
為甚麼你要喜歡粉紅色.
你不喜歡藍色.
我講的是真實的例子.
因為我們在華人社會裡面.
我們有一個想像.
男孩子應該是這樣的.
女孩子是好的男人.
這樣就是好的女人.
我們同樣地.
都是被我們的社會文化框住.
以致我們對於好的男人.
好的女人.
是一個定義.
問題就在這裡.
今天.
楊寬四尊在告訴我們.
不是再用這些文化價值.
用這些典範去規限.
甚麼是好男人和好女人.
耶穌做的事不是這樣.
耶穌沒有說.
這個撒瑪利亞婦人.
你不貞潔.
你要悔改.
不是這樣.
但我們橫教會就是在講這些.
我們一起去旅行都不行.
不要說一起去旅行.
離婚的人是做不到執事的.
在橫教會你試試離婚.
不要說做牧者.
你做執事都不行.
我講的是真事.
在我們的現狀.

$^{761}$在我們現況裡面做的事.
今天我想講的是甚麼呢.
今天耶穌基督.
要帶給我們生命.
真正的拯救是甚麼.
我們不要再混合.
我們的文化價值觀在裡面.
我們要問聖經.
基督教的價值是甚麼.
我們很受我們社會的包袱影響.
我們覺得好的女人.
其實每一個時代都是愛馬協會.
靚 瘦就是一個好女人的形象.
今天如果回到我們的個人層面.
其實我們是否同樣地.
都被很多文化的包袱.
影響我們怎樣看自己.
如果我們達不到這些標準.
如果我們不是這樣的話.
我們就不是一個好的男人.
我們就不是一個好的女人.
我們同樣地被我們身處的社會.
的文化價值去標籤.
今天耶穌來.
是要拆除這些標籤.
今天我們怎樣重新去看.
甚麼才是真正的好男人.
甚麼才是真正的好女人呢.
今天是聖餐的主日.
一會兒Mark Sir會帶領我們一起去聖餐.
聖餐是代表著耶穌基督.
臨在我們當中.
有一個神套的說法.
聖餐就是基督臨在的象徵.
當我們一會兒領受聖餐的時候.
我們要去想.
今天耶穌臨在的時候.
耶穌與人同在的時候.
祂帶來的是甚麼改變.
今天祂要重新去定義我們的價值.

$^{801}$我們生命的價值.
我們作為一個男人.
作為一個女人的價值.
不再被這些文化的框架所定義.
祂要幫助我們重新去活出.
甚麼才是真正的男人.
甚麼才是真正的女人.
我們正在經歷耶穌的一份愛.
這份擁抱正如昔日.
耶穌與撒瑪利亞的婦人.
面對面的相遇一樣.
今天在聖餐當中.
當基督臨在我們當中的時候.
我們能不能夠.
就好像昔日的撒瑪利亞的婦人一樣.
能夠與耶穌面對面.
讓祂重新定義我們的價值.
重新定義我們的身份.
這就是我今天想說的東西.
讓我們一起去禱告.
親愛的天父.
今天我們來到你的面前.
主佑我們.
無論是華人的社會還是美國的社會.
對於男性和女性都有很多框架.
有很多定義.
主今日求你幫助我們.
讓我們能夠經歷你這份解放的愛.
讓我們的生命被你重新肯定我們的價值.
我們不是被這些社會的標籤.
去標籤我們是誰.
我們不再被這些社會的標籤.
去標籤我們是否一個好男人.
是否一個好女人.
主今日願意讓你的愛淋在我們的生命當中.
在一會領受聖餐的時候.
你就跟我們說話.
正如你昔日與撒瑪利亞的婦人面對面對一樣.
今天你將我們從這些標準.
從這些社會的包袱.

$^{841}$文化的包袱當中去解放出來.
讓我們可以在你裡面找回我們的價值.
讓我們可以在你裡面.
去活出什麼是真正的好男人.
什麼是真正的好女人.
主今日願意淋在我們當中.
去跟我們說話.
願意釋放我們的生命.
願意改變我們的教會.
改變我們對於性別的想像.
是奉主耶穌基督的聖名祈求.
Amen.
\newpage



\section{}
\label{sec:HYnhtr6gM20}
\textbf{Everlasting Praise (Psalms 8) - Rev. Mark Chan}
\newline
\newline
連結: \href{https://youtube.com/watch?v=HYnhtr6gM20}{\texttt{https://youtube.com/watch?v=HYnhtr6gM20}} ~~~~ 語音日期: 2024-04-15
\newline
\newline
\hyperref[sec:u_YoyAgPSkE]{\small{< < < PREV SERMON < < <}}
~
\hyperref[sec:index_chronic]{\small{[返順時目]}}
~
\hyperref[sec:index_scriptual]{\small{[返順卷目]}}
~
\hyperref[sec:DfbxT6XIBuk]{\small{> > > NEXT SERMON > > >}}
\newline
\newline
$^{1}$Well, thanks Kelly. And we do have some staff transition this month and this year.
uh, Reverend Ng, our senior adult minister is going to retire by the end of this month and.
actually some of the senior adults are having a.
farewell with him today.
So we're going to pray for him because he's getting weak and he cannot serve us.
Any anymore because he's physically and mentally being challenged.
um, so continue to pray for Reverend Ng and he's going to retire and the staff will.
um.
Be making some transition in our staffing and in our ministries as well.
Uh, so more details to you and Kelly. Thanks for announcing the general assembly and please.
Make time to be here.
because we're going to go through our finance and our needs in the church and.
We invite you to participate and to contribute.
and once again, uh, Chi is going to be our assistant pastor and.
We can only afford him for part-time position.
Um, so, uh, please continue to give and please continue to support the church ministry so that we can grow together. All right.
um, so, um.
Do come to the general assembly and we'll tell you all the details and uh things that we forecast and we do need.
every sister and brother in Christ to.
Participate and to build this church together because this church is the body of Christ. It's our body is our family.
So we are all part of it and we are all members and we need to.
Work together and build this church. All right.
Um today I would like to start a mini.
sermon series on worship.
uh worship.
Uh, you probably come to church for some time.
Uh, probably you're the first time here. But in any case sometimes we gather.
uh together to sing.
to pray.
and to read scripture and reflect on the word of God and we're going to.
Care for one another and this is what worship service.
Uh looks like.
and I think years of.
COVID and pandemics have deteriorated.
our sense of worship.
Don't you agree?.
Right. I mean remember those three plus years we had the online and live streaming, right?.
Where were you right? You were probably on your bed, right?.
You're sleeping in and then oh, it's about time to uh, join the online worship service and you turn on your phone, right?.
And then you kind of like you didn't brush your teeth and you know, you just watch your phone and oh I worship.

$^{41}$You know, you know, we are deteriorate into that kind of uh,.
You know low state of worship. So, uh, we have to pick it up.
We have to understand what worship is all about and we have to understand why we are here.
Right, we are here not only because of some of the benefits that we get, you know people and friends and uh,.
You know food and all that stuff, but we need to understand why we are here to worship and the word worship.
Is actually coming from uh an old anglish word of worship, right?.
uh, what is what what it does, uh mean is that um.
something of value.
and something of honor.
Either a person or an object that we assign or we ascribe to someone or something, right?.
Someone or something.
Um, you go buy a car, right?.
These days cars are very very expensive, right?.
Uh, someone just bought a new car.
uh lately.
Uh, and she paid for over 40k.
For japanese car, right? Of course. It's a great car, right?.
But is it worth right?.
That's the word is all about.
you.
Think of value and honor of a person on object and you ask the question whether he or she or it worth.
The uh the value that we assign or we ascribe to it and in some sense worship is all about value.
And honor we ascribe to god.
God is great and god is to be honored and he deserves the honor and that's what worship is all about.
And you flip the pages of scripture.
from the book of psalms from the old testament from the new testament and.
It talks about the.
worthiness of god.
Worth our time and our effort to bring him the glory and to praise him.
sample here.
Revelation chapter 4 11 you are worthy.
Oh lord our god to receive.
glory honor and power.
both revelation talking about having you.
Worship what we do eternally.
It's all about worship. It's all about the worth of god that we bring to him and we praise it.
And that's what worship is all about. We have to see.
the value of god.
in your life.
That god is worthy of your everything.

$^{81}$Your praises and your prayers and your worship.
So we're going to turn to psalm 8.
And to think about this topic of worship and you have your bible, please open it up to psalm 8.
It says it's a very very short, uh, nine verses long psalm.
That talks about worship and there are three points I want to go through with you.
First two verses it talks about the worth of god.
And then the second part the middle part talking about our worth as human beings as creatures of god.
And then the last part talking about our mission our work.
Uh, we can.
Collaborate or we can serve god and work with god together in this world of creation.
So let's go through this psalm together psalm 8 a psalm of david. So it's written by david the psalmist.
It says oh lord our lord.
How majestic is your name in all the earth you have set your glory above the heavens.
You have set your glory and honor.
above heavens.
So you can see what david is trying to do.
He observes the heavens the sky.
And then the earth the creation the cosmos the vast universe.
And then he exclaimed that god is to be worshipped.
How majestic your name in all the earth I can see your name. I can see your power.
I can see your authority over everything that you have created and that's what david is doing.
He looks at the universe and then he bursts out with.
Uh.
awe and wonder.
Now I think I think if we have to rebuild our life of worship.
If we need to know how meaningful it is to gather as church.
No matter how many people we have to worship god together. We have to recover a sense of awe and wonder.
Don't you agree?.
the awe and wonder.
when we observe the vast universe how amazing it is.
And then we think of the the designer or the the god who created everything.
And we think of him as someone.
and that.
We cover that sense of uh, uh, awe and wonder and we praise him and we worship him. That's what david is doing here.
Now last week millions of people.
Woke up early or even be there, you know a day or two before and they went to.
Uh the the path of totality, right?.
and to see the.
100.
solar eclipse.

$^{121}$How many of you did see it?.
Did you see it?.
No. Yeah, right. So you went out to the street and it's all there, right?.
And you look up into the sky and then you see the wonder of creation.
Many of us or I should say well, probably many of us.
We went.
uh monday morning.
and drove many hours.
To reach that path of totality and in order to observe the the amazing universe.
right the vast, uh, uh universe the uh,.
awe and wonder.
Some people even tear up when they see.
The solar eclipse right? Isn't that amazing? I saw it too, right?.
And you know according to some survey there were about four millions of people.
four million people.
They don't live along the path of totality.
but they.
They they drove there.
Right, they they they drive there or they fly there in order to see that right?.
um.
It's just amazing number, right?.
And I did a short video and um, I I thought of uh this book.
called uh, the privilege planet.
Uh, it was written some years ago by two.
scientists.
and.
What they are trying to bring is the idea this the concept of you know, our earth.
Our planet earth is a privileged position to do all of this.
um.
observation.
and we.
they argue.
On this planet earth. We are in this swiss spot, you know.
the optimized place in the entire universe.
So that we can discover the laws of physics so that we can do a lot of different scientific discoveries.
And because we are being put in this.
Situation right the atmosphere is protecting us and we're not too far from the sun. We're not too close to the sun.
and then the moon and the stars and everything is being designed in such a way that.
makes scientific discoveries possible.
And I think that's a very very interesting idea. I know some people would disagree with it. But um.

$^{161}$If you look at psalm 8, that's what david is trying to say.
We are in the swiss spot on this planet earth and we can observe the wonder of creation the greatness of creation.
So we can appreciate how great our god is, right?.
verse 2.
from the lips of children and infants.
You have ordained praise because of your enemies to silence the foe and the avenger.
Now, what is what is he talking about? This is.
kind of a difficult verse.
And there are different translations.
Of this verse but the idea you can grab is pretty clear.
Uh, first of all is it talks about children.
and infants.
the babies.
Right the babies, you know our church has.
Many many babies, right? We have uh, some twins, you know, we have many cute babies.
you know.
We modern people we hold a baby in our arm, right hold an infant. Um, what what is your what's your feeling?.
Right.
It's awesome, right?.
It's cute baby.
Right. This is how modern people see babies.
But I probably explained this to you before.
for the ancient people in.
biblical culture.
Babies are not cute, right? Even we understand it, right? Well, I have two.
children, right?.
Uh, they are not cute anymore.
Um, they are older now and I continue to spend money on them, right?.
And they give me suffering.
So we don't we don't optimally think of they are cute, right?.
Now the ancient people they understood this because it's liability.
And babies are powerless people.
We understand it and they understood it because you know, we look at babies, uh, you know, we we have to feed them.
We have to change diapers for them.
And we have to protect them, right?.
my wife.
wakes up every saturday morning.
And she always has the habit to catch a tv show and the tv show is it's all about nature.
We're talking about some environmentalist going out to the nature and try to protect some of the extinct animals and all that stuff.
and she's always.

$^{201}$um, you know.
uh amazed by you know, how how.
How how interesting the universe is how interesting the creation is and you know, if you observe the creatures.
the animals.
most most of the animals.
when they.
give birth to.
The babies so-called the babies usually the babies can take care of themselves, right?.
either deer or birds.
uh.
Eagles, right? Usually they can take care of themselves, right?.
And the effort to feed them the effort to protect them are usually uh, pretty minimal.
but for humans.
We are not.
We're not you you think about it when we when we were born.
We were helpless.
We were powerless and when it's someone our parents.
you know to.
Love us to protect us to feed us.
And this is what verse two is all about.
It says god.
has ordained.
praise.
In the babies.
Right, I think a better translation is not the niv translation here. I think better translation is the word praise is actually strength.
Fortress or protection that means god is to protect.
the helpless babies.
Those people who find them helpless.
Those people who find them powerless and weak and they humble themselves before god.
And they're not proud before god. They will receive the protection and strength from god.
That's that's what psalm 8 is all about.
David looked at the sky and the universe and he.
Recovered the sense of wonder god is so great because the creation is so so great and we are so helpless.
And we are so powerless and we have to look up to god for help and strength and protection.
That's what psalm 8 2 is all about.
Now, um, we have to understand the psalm we have to understand the.
Design of the psalms, you know, the book of psalm has 150.
Psalms in it, right?.
And uh here I list out the first 14.
14 psalms and we know the the entire book of psalms the psalter.

$^{241}$has.
Two two psalms in the beginning psalm one and two as introduction.
Is a dual introduction to the entire book of worship book of psalms.
And then there's a group of psalms psalm three to seven we call them psalms of david is written by david.
um, and all of this psalms from three to seven has has one.
single central theme.
Is that david?.
was helpless.
he was being anointed king, but.
Uh, he was pursued by his enemies and he was hopeless to defend himself.
And he prayed to god and he's trusting god and god becomes his protection.
And he puts his faith in god and god defends him god defends the humble and overcome his enemies.
That's what psalm three to seven is all about and then there's another group of psalms.
nine to fourteen.
Uh runs along the same idea.
They are the individual laments.
the psalmist is crying out to god because he.
Is powerless against his enemies all these persecutions and all these sufferings.
Uh difficulties in life and he look up to god psalm nine to fourteen.
It's all about humble people finding protection in god.
And psalm eight is right.
Between right at the center and it celebrates the power of god.
power of god over those.
Who are like babies?.
Who are defenseless?.
Who are powerless?.
Now should we get that feeling when we read the psalm along with the psalmist david.
And we are all struck by the power of god in creation.
And at the same time the power of creation the power of god humbles us.
And in light of god's greatness, we are like babies.
We exist we have our beings all depend on.
the god.
Who create everything who created you and who created me who sustained you?.
And who also sustained me?.
Without god we cannot.
be.
We cannot exist.
A couple weeks ago. Um, I have some serious, uh lower back pain.
You know, I haven't had back pain for a long time, but in a couple weeks ago, I just had this, you know.
Very quite serious, uh lower back pain and uh, I use a lot of uh, salonpas, you know, I put some all around me.

$^{281}$All around me, but you know, it didn't really help right and then my my skin got rashes in.
It's just terrible right?.
And I and I don't know what to do and I didn't get hurt, you know, I didn't do anything to get hurt.
And um, I don't know. Yeah, i'm over 50 years old, right?.
You look at me.
I'm in pretty good shape, right?.
Compared to other over 50 years old man, right?.
If you grade a man who were 50 years old, you know, great a b c d f right? I'm probably on the uh,.
You know the top level, right?.
But you know, I just find it difficult to to recognize that I I have this back pain, you know for days.
And the only conclusion I had is that i'm getting old.
I'm getting old.
And don't we all get get old, right?.
um, we we we get weak and we realize that only god.
Is our strength and that's why god is worthy of our praise when we see and sense his power in the creation.
Uh, psalm 8 continues to to tell us our worth.
Right. So david look at the universe and he think of how majestic the power of god is and he also.
Like looking in into the mirror.
He understand his own position in the universe.
Verse three he says when I consider your heavens the work of your fingers.
The moon and the stars which you have set in place.
What is man that you are mindful of him the son of man that you care for him?.
So he exclaims.
When I consider the heavens when I see that how big the universe is it tells us.
Hey, we are not worthless.
Even though we are so small.
And we are so insignificant.
But god cares for you and for me.
Sometimes don't we feel that way, you know go for our daily life, you know, we we.
We have to tackle different challenges and you know life is tough, right?.
And sometimes we do feel that way. We are insignificant. Nobody cares about us, right?.
Psalmic says no you are worthy in god's sight.
Yes, we are not worthy of the grace of god, right no one can earn the grace of god no one can.
Right. We're all sinners. We're all broken. However.
in the other hand.
The bible reminds us and psalm 8 reminds us.
We are worthy.
of god's sight.
And that's why some people read the psalm.
as.

$^{321}$Telling us the paradox.
of human being.
You know what paradox is.
It's not quite contradiction. It's not logical contradiction.
Paradox, it's like it's like two sides of the same coin.
Right, it tells different.
Aspect of the same entity.
Right, but it's a different angle.
Here psalm tells us the paradox of human being and we find both sides of being human.
Yes on the one hand.
In the vast universe. We are just like a speck of dust. We're so insignificant.
Yet on the other hand.
In god's sight you are worthy.
It says what is man?.
man in hebrew is.
And knows it's it's a it's a common word in hebrew to describe mankind. I think the nlt gets it.
Right nlt new living translation gets it right it translated as mere mortals.
Mere mortals and we are all mere mortals, right?.
And we'll die someday. We'll pass away someday, right?.
The song leader the worship leader reminds us right?.
Our life is broken, right? We are fragile.
We pass away quickly, right time flies. I was talking to one of the brothers.
Uh, you know, he said oh my kid is eight years old ten years old now, right? I said time flies.
Right, we pass away. I pass away quickly and we're mortals.
and then.
The son of man. What is the son of man that you care for him?.
right.
Here the original language is son of adam.
Where we get the english name adam, right? So we all know some hebrew, right?.
adam.
son of adam.
And what does it mean son of adam, you know hebrew hebrew is a pretty interesting, uh language sometimes.
Two words sound alike and the meaning are related, you know and adam or adam is one of the words.
Adam means you know man.
Person or a human being right and then in hebrew.
ground.
the earth.
You know the dust.
The earth what we're stepping on is adamah.
Right. So adam and adamah.

$^{361}$They're related right with dust.
We're turning to dust right? We are so close to the ground, right? So when when you die, you know six feet under right?.
We are mere mortals.
And the psalmist david is using those terms to remind us.
That we are limited. We are powerless and we are insignificant.
So yeah, there you go. If you feel that you are insignificant.
Nobody care about you and you sometimes feel.
That you are not good enough or you don't measure up when you compare to other people.
You're exactly right.
You should feel that way I should feel that way but the song reminds us the other side of the same coin.
That we marvel how valuable we are in the sight of god.
Again, two terms tell you.
How much worth you have in the sight of god?.
That you are mindful of him. What is meant that you are mindful of him?.
Mindful of him can be translated. God remembers you right?.
Uh when you're happy.
God remembers you when you are crying when you are weeping when you're sad when you're depressed god think about you.
think of you.
That's what mindful means.
and then.
What is the son of man that you god you care for him?.
and the word care.
Uh, it's a very interesting term again in hebrew. It can.
It basically means.
Role calling role calling, you know when we go to primary school when you go to high school the teacher walks into the classroom.
and she does.
the.
Attendance sheet, right?.
And then we'll call every single one of you.
Call you by name. Are you here?.
That's what care for means.
You are on the list of god's.
attendance.
You are here and god is mindful of you that god cares for you.
So.
Here it reminds us that we are valuable in the sight of god.
and then at the same time he continues in verses five to nine to tell us how.
You know how how um, how glorious we are.
Versus five and following god you made him man.
Every one of us a little lower than the heavenly beings or angels.

$^{401}$And crown him with glory and honor.
You see that in the sight of god.
We are crowned.
with glory and honor.
We are so much better. We we this phrase is with glory and honor crown with crown him with glory and honor.
Basically means you have the highest honor in the entire universe.
No, you're different from all the other creatures.
Because we human beings are being created in the image of god and we have so much worth in the sight of god.
Yes, we are limited.
but we are.
Valuable in the sight of god.
Let me continue verse.
verse six and following.
You made him ruler over the works of your hand and you put everything under his feet.
all flocks and herds and beasts of the field.
The birds of the air the fish of the sea all that swim the paths of the seas.
Everything is under.
the feet of men.
Under the feet of humankind.
Um god made us god made human beings.
This is not an accident this is not an afterthought.
The psalmist is reminding you that god made you.
God made you however insignificant you feel god made you god created you.
Now you remember some some years ago. I talked about this picture pale blue dot remember this picture.
Um, it's it's taken by uh, nasa's uh voyage one.
Uh many many years ago. It was taken in the year of 1990.
February 14 valentine's day.
um, it's the voyage one is is a is a.
photographing.
Machine and just being launched by nasa into the deep space.
and on that day.
1990 february 14.
The cameras switched back before, you know, uh, the earth disappeared.
And.
It took the picture and you can find the very very small dot.
the pale blue dot so-called.
still visible in the milky way.
Now that picture reminds us.
That we are a speck of dust in the universe. We're right there, right? Can you see that?.
We're right there.

$^{441}$and sometimes our secular culture.
Will have us believe that we are.
Not significant.
Or we have no meaning in our existence, right matters are eternal, right? All that matter.
Is the matter right?.
We're just a bunch of people, you know with our chemical makeup or physical makeup, right? And that's it.
Where is the?.
Where where is the intrinsic value of human being?.
We have no.
Intrinsic value. We're just accidents.
Um and nature just take its course.
And there's no inherent value or dignity in our existence.
Um, that's what the uh secular culture or naturalistic culture is trying to have us believe.
In the same way in the biblical culture.
in bible bible time.
You know all these neighbors of israel.
babylon egypt.
Or or whatever cultures are surrounding israel. They believe in some of the myths.
And some of the myths were saying that humans were being made to be slaves to other gods.
Humans exist because humans are to be slaves.
to the gods the lesser gods or the uh, the higher gods and all the gods are in in.
up there having vacation.
And you know from time to time we have to you know, bring in uh our offerings to please the gods, right?.
So the gods can have a good time up there.
And that is the cultural myth of the day.
if you if you if you hear the uh similarities.
The ancient stories.
the ancient myth of humans being slaves to the gods.
Right no values and then our western.
secular culture.
They're actually pretty similar.
We are here with no reason.
but the bible.
Says something totally different the bible says we human beings are no accident.
And you are not.
an afterthought to god.
The psalmist says god made us.
a little lower than the angels.
The highest honor we can have in the entire universe and secondly, it says god made human beings.
Rulers.

$^{481}$So human beings are unique. We have a royal.
status.
We're like kings and queens.
Only human beings are made in the image of god.
You know you we call genesis 1.
In fact, psalm 8 is a summary of genesis 1.
And it tells us god set us apart from all other things in the universe and all other creatures in the universe.
Because we are the image of god and we are the representatives.
Of god the king god the creator on earth.
So you read psalm 8 only nine verses.
Is actually a brief summary a condensed summary of genesis 1 god made everything in the universe for humans to flourish.
God made you and you are unique. There's no one like you.
There's no one like me and god gives.
Each and every one of us the honor to rule with him to govern.
to.
operate.
to.
take care of.
uh.
The creation and we have this honor to be part of the work of god. All right.
As we close.
I want you to take three things home.
first.
You're worthy of god's care.
Your life your existence.
It's not talking about your work.
It's not talking about your your achievement.
you.
as a human being.
Man or woman.
You are worthy of the care of the creator god.
right.
Your existence your being.
It's not talking about your work right the asian culture. We are all performance based right? Oh, we got a four point also, you know.
My parents would praise me right so I get more values or how much I earn how many figures on my paycheck.
Right. Those are our achievement that bring us values.
But some is telling you your being your existence alone.
is.
worthy.
of god's care.

$^{521}$And then secondly we are being called to care for the kingdom of god the creation work of god.
You go back tomorrow monday.
your accountant your teacher your.
Whatever industry you are in.
God see values in your work because you are making the you uh, the the creation.
working.
operating.
God can use the angels to do all your works.
Right to make the society running.
God can use other means but he has chosen you.
To be one of the workers to work alongside with him.
And that is the value of psalm 8. It's telling you god may you and me ruler.
We can take care of god's creation together.
And that's why we are different. That's why we're unique.
number three.
God is worthy of our worship.
Uh, we have to recover the sense of wonder when we see the universe when we see the god behind the universe.
When we sing songs to god.
When we sing praises to god we just wonder man god is so great.
And yet he takes care of me, right? He can have relationship with me. He's like friend to me.
and I can have a personal relationship with.
with this creator.
Um god is worthy of our worship and we have to recover that sense.
You know, it's a sense of wonder all the time when we walk into the church when we worship together.
You know, I did not go to.
Uh the path of totality to see the full 100 solar eclipse. I should have gone. I should have gone now.
I regret it.
I should have gone. However, you know, my son went, you know, my son and his friends, you know about 10 of them.
uh, they drove a couple of cars and uh.
You know, my son nathan is in uh, uh champagne, right? He's in june. It's junior in in champagne.
So, uh, he he told us.
Hey dad, we we want to go to see uh, the uh, the solar eclipse.
I mean it would take us two hours of driving from champagne to somewhere called paris.
illinois.
Paris, you know not not france paris, illinois.
So, uh, so he told me about it, but we have to skip class.
I said sure. Yeah, right. Well you skip class anyway, so.
uh, so they drove they drove two hours to paris in the open field and then they.
They bring the cameras and they they bring the devices and uh, they do the time lapse and stuff and it's just amazing.
And he sent back some of the video files to me.

$^{561}$and um, I I watched those videos and there's nothing to to really see except the the eclipse because the.
The camera is is facing upward right? You can only see the sun, but I can hear what they.
say.
In the entire time the full 100 solar eclipse takes about 3.5 minutes.
And I watched those videos looking at the sun.
And hearing what they say, you know three and a half four minutes long video.
And all I hear is wow.
Whoa.
Ah.
And all that's it's meaningless, you know words, right?.
But it's their the expression of sense of wonder.
When you see the eclipse and you understand how vast the universe is and how how wonderful the universe you are all struck.
right.
And all they said in the in the entire four minutes is just that right? Wow.
This is great, this is wild.
Psalm 8 reminds us when we worship god we should have that sense of wow.
God is so great and yet he loves us.
Right you think about salvation we have in jesus christ.
God came to the speck of dust.
In this vast universe.
Jesus came and saved us.
Jesus came and died on the cross so that we can restore our relationship with him.
Wow.
We should be all struck.
By the greatness of god and by the love of god, right? Amen.
Let's all do that every single week.
When we come here.
Now we have to come on time. Okay.
But come on time because it's just like solar eclipse, right? You missed out you missed out, right?.
The entire worship service is the offering to god and to see that god is worthy of our praise.
Let's all pray together.
Father we give you thanks because.
You are a great god.
And you are great and awesome.
And god we see this universe we see this creation that you have made and we see how wonderful it is.
How great this universe is and yet we sense that we are so insignificant.
We're so small.
We're so tiny.
But you are mindful of each and every one of us.
So we thank you for your love.

$^{601}$Thank you for your grace.
And only in jesus christ we can be your people and we can be called your children.
So we come before you we just want to recover the sense of wonder.
To worship you to praise you and to see how great you are.
And how worthy we are in your sight.
So thank you for your love and send us back into our daily life this week so that we can be your representatives.
We can work for you.
We can relate to other people.
And show them your glory your reality.
So we pray that you will continue to grow us into this spiritual awe and wonder.
That we can truly worship you.
in christ name pray.
So god bless you and we want to see you next week.
Uh, we can we're gonna recover and we may and maintain our sense of wonder as we worship and serve god together.
All right, god bless.
[BLANK AUDIO].
\newpage



\section{詩篇}
\label{sec:DfbxT6XIBuk}
\textbf{永恆的敬拜 1 (詩篇 8) - 陳崇基牧師}
\newline
\newline
連結: \href{https://youtube.com/watch?v=DfbxT6XIBuk}{\texttt{https://youtube.com/watch?v=DfbxT6XIBuk}} ~~~~ 語音日期: 2024-04-15
\newline
\newline
\hyperref[sec:HYnhtr6gM20]{\small{< < < PREV SERMON < < <}}
~
\hyperref[sec:index_chronic]{\small{[返順時目]}}
~
\hyperref[sec:index_scriptual]{\small{[返順卷目]}}
~
\hyperref[sec:yghFb3uoZ_c]{\small{> > > NEXT SERMON > > >}}
\newline
\newline
$^{1}$我們在這個月和這一年都會有些轉折和過渡.
由我們UV開始第一天.
我們還沒開始之前已經喝了茶.
我們就一起去同工.
吳牧師是一個很好的榜樣.
在我們中間很忠心事奉.
無論是年長的或年少的.
其實在吳牧師身上都有很多學習.
什麼叫忠心和堅持到底.
因為他的身體的緣故.
所以我們上一個主日.
我們在Call Team有一些討論.
吳牧師提出會在四月底離職.
所以有一些改變.
教育同工和Call Team都會有不同的安排.
所以其他的細節我們會容後公佈.
另外我們知道兒童侍工.
青年侍工.
英文侍工.
其實我們都需要一位同工.
我們一直都有禱告.
最近我們開始聘請一位大兄.
叫楊之逸.
楊之逸傳道.
叫Chee.
他是馬來西亞華人.
我們很巧合地遇上了.
我們討論過後.
面試過後.
我們覺得楊傳道很適合我們.
所以我們就開始跟他合作.
因為他是馬來西亞華人.
我們跟他合作H1B.
不同的工作簽證.
希望大兄姐妹為了這方面去禱告.
我們暫時只是聘請他兼職.
我們只能夠支付他兼職.
希望他能夠很快上到軌道.
因為我們實在很需要一些同工.
這些細節我們在嘉士大會.

$^{41}$也會跟大家多說一點.
如果你想知道誰是Chee.
誰是楊傳道的話.
做父母的.
今天已經取消了.
你隨便過去.
那裡看到有一個年輕人.
就是他.
比我年輕.
怎樣也好.
希望大家大兄姐妹.
為教會的歌道來禱告.
好嗎?.
今天我們開始一個系列.
是講到敬拜的.
講到敬拜.
可能你第一次來教會.
或者你來教會已經很多年了.
你是那些老差骨.
已經習慣了教會發生什麼事.
你療癒之障.
那些老差骨.
都是需要知道什麼叫敬拜的.
我們進來這裡.
吃完東西.
喝完咖啡.
然後我們唱幾首歌.
祈禱.
讀一下聖經.
聽道等等.
我們就說這個是敬拜.
敬拜是一回什麼事呢?.
我希望在這個系列裡面.
跟大家思想一下.
因為幾年的.
之前幾年的疫情.
已經令到我們的敬拜生活.
大受打擊.
我們墮落到一個地步.
睡在床上.

$^{81}$然後拿著手機.
我告訴你.
每個人都是這樣的.
拿著手機.
眼睛蒙著.
還沒刷牙.
然後我們敬拜了.
我們墮落到這樣的程度.
真的搞錯.
你不明白實體現場的重要位置.
昨晚我和我太太去了芝加哥戲院.
我們看一個競技節目.
就是一個競技節目.
幾千人在那裡.
和你YouTube看相差很遠.
我完了之後.
我和我太太說.
這張票值得買.
因為是我買的.
所以我們兩個拍拖.
看完之後.
現場的實體氣氛.
是多麼的不同.
你在教會的崇拜.
其實都是一樣.
剛才你唱詩歌.
去讚美上帝的時候.
你有沒有留意聽到其他人唱.
有吧.
那些YouTube是聽不到的.
當然我明白.
我們有一些直播.
和看YouTube的大英姐妹.
可能在芝加哥以外.
可能找不到教會等等.
我們非常歡迎你.
但實體是不能夠取代.
實體不能夠取代.
什麼叫做敬拜呢.
敬拜Worship這個字.

$^{121}$其實是一個很老的.
很舊的一個英文字.
Worship.
大概是主後十幾世紀.
才開始有一個這樣的字.
其實是什麼意思呢.
就是說值得的意思.
Worth.
我們說你是不是Worth.
是不是Worth my time.
是不是值得我的時間.
有沒有這樣的尊貴.
有沒有這樣的配得.
我們是說一種價值.
無論是一個物件的價值.
Airview手袋為什麼要二千元.
哪值得呢.
我們就會這樣.
或者一個人.
我們值不值得花時間給你.
一個這樣的人.
值不值得花時間給你.
說價值.
Worth.
所以聖經裡面.
無論舊約或者新約.
這裡其中一個例子.
啟示錄最後那本書.
他說我們的主.
我們的上帝.
你是配得.
You have the worth.
榮耀尊貴權柄.
因為你創造了萬物.
萬物是因你的旨意而被創造的.
你是Worth.
Worth my time.
值得我的時間.
值得我的精神.
來將榮耀尊貴權柄歸給上帝.

$^{161}$所以一個敬拜.
真正的敬拜應該問什麼.
不是問今天音樂做得如何.
不是問氣氛做得如何.
別人給了我什麼.
真正的敬拜是你給的Worth.
你看到上帝有這樣的價值.
所以你發出你的祈禱.
發出你的詩歌.
你的讚頌.
去到讚美上帝.
這個就是敬拜.
詩篇裡面有很多很多.
上帝是配得.
這種說話就是在說.
Worship 敬拜.
這次我們去看詩篇第八篇.
我有三點.
詩篇第八篇說到上帝的Worth.
他的價值.
然後就說到我們作為人的價值.
然後告訴他.
我們到底應該做些什麼.
和上帝有些什麼合作.
和上帝有些什麼.
在他的計劃工作裡面.
我們能夠參與的.
這個就是詩篇第八篇.
詩篇第八篇第一節.
耶和華,我唔的主,你的名在全地,何其美.
你的榮耀彰顯於天.
這是大衛所寫的詩篇.
就是說他見到大自然.
見到天空大地的時候.
他就忍不住.
他就發出了一個讚美的聲音.
他說耶和華,我唔的主,你是最高之主.
你的名字,或者你的權能在全地.
不是一隻大.
是美,是Beautiful.

$^{201}$漂亮.
你將你的榮耀彰顯於天.
見到天,見到地的時候.
就不禁說出How majestic.
多麼的尊貴.
是很美好的東西.
大衛有這樣的反應.
我相信我們擔高頭.
不要只看電話.
我們都應該有這樣的反應.
上個星期,大家有誰看日蝕.
有呀,各位.
每個人都去看日蝕.
很多人看完日蝕.
心裡面想到的東西.
或者感覺到的東西.
就是那種震撼.
那種影響.
見到宇宙是多麼偉大.
一個日蝕,一個全日蝕.
我們能夠見到.
是一個很精心的創造設計.
fine tune了.
我們才能夠看到.
如果不是,我們就不能夠欣賞到.
見到宇宙多麼浩瀚的時候.
我們應該有一種All in wonder.
其實就是敬拜應該有的spirit.
今天你走進來.
敬拜我們的主,我們的上帝.
對象是他.
敬拜的對象是上帝.
一位創天造地的主.
我們應該要重新找回那種.
All in wonder.
上個星期,在全日蝕的軌跡上.
軌道上,大概從德州開始.
一直到美國東北的角頭.
整個軌道裡面,估計大概有400萬人.
他們不是住在軌道上.

$^{241}$住在軌道上已經有31萬人.
3100萬人住在軌道上.
可以從角度上看到.
有400萬人.
他們不是住在軌道上.
他們開車走到軌道上.
就是為了看什麼.
就是看全日蝕.
中間我們也有不少.
故意這樣路遠開車.
開幾個小時走到那裡.
就是看那個日蝕.
你想想,香港有多少人口?.
700萬?800萬?.
香港人口突然間沒了一半.
去了哪裡?.
走了去那個軌道.
Path of Totality.
就是看那個日蝕.
大家看過我一個短片.
當我看日蝕的時候.
我在天台看.
我就想起這本書.
叫做Privileged Planet.
就是一個獨一無二.
一個恩寵的地球.
上帝所恩寵的地球.
一個得天獨厚的地球.
一位科學家.
兩位科學家.
寫了一本書.
他就爭論一個idea.
你可以不贊成那個idea.
你覺得可能是一個圓形的idea.
Any case.
他就說我們的地球.
其實是一個最好和最optimized.
讓我們有科學的發現.
能夠觀察到宇宙的近和遠的地方.
我們現在站立坐在這裡的地方.

$^{281}$就是上帝設計到一個地步.
以至我們能夠在這裡有科技的定律的發現.
然後可以使用這些科技的定律.
去應用在我們生活裡面.
你今天拿著手機.
今天我們有享受的科技.
是怎樣來的.
讓世代的人能夠去discover.
能夠去發現這個世界的奇妙.
在哪裡發現.
就是在這裡.
就是在這個地球裡面.
所以我們看見日蝕的時候.
你知道是多麼精密的一個宇宙.
有沒有玩過高爾夫球.
有誰玩過.
我曾經試過玩過一次.
那個人教我.
他們說sweet spot.
你甩那支東西的時候.
有個sweet spot.
他就去了那個位置.
我甩來甩去都甩不到.
所以我還是做回教會好一點.
就是那個sweet spot.
我們就在那個sweet spot.
我們能夠見到宇宙的奇妙.
你的名在全地.
是何其美.
How majestic is your name.
in all the earth.
然後大衛說什麼.
你因仇敵的緣故.
在孩童和吃奶的口中建立能力.
使仇敵和報仇的閉口無言.
這節比較難一點.
什麼意思呢.
我們說的是孩童.
說的是小孩吃奶的.
在聖經文化裡面不同我們的文化.

$^{321}$我們今天聽見小孩.
聽見小孩哭我們都覺得很煩.
不是說他.
當我們見到小孩的時候.
我們覺得很可愛.
很可愛.
但其實古代的文化不是這樣的.
古代文化見到baby的時候.
見到小孩的時候.
就是見到他們無能.
他們確實是無能之輩.
powerless.
人真的很無能.
特別是小孩.
就算今天也是.
由他出生到他長大.
整個過程.
你都要給他很多能量.
給他很多照顧.
還有很多錢.
我的兒子已經很大了.
二十歲又多了.
不斷給錢.
是多麼無能.
他很無能.
我幫他找錢等等.
特別是古代聖經文化就是這樣.
因為你要餵他.
幫他換屎片換尿片.
他根本是沒有能力自己去生存.
人類真的很奇怪.
我太太經常.
有時候星期六早上.
她就在電視上看一個show.
就是講那些自然界的東西.
有些人去保育.
或者去護理大自然.
救回動物等等.
很奇怪.
所有動物一出生之後.

$^{361}$他們自己可以處理.
通常都自己可以處理.
人就不是.
人就處理不了.
很老都可以處理不了.
這個就是大衛所講.
他說孩童和吃奶的.
是沒有能力的.
但是誰給能力.
或者應該怎麼說.
更加好的反應是建立了保障.
就是那些無能的人.
覺得自己生命裡面的能力.
是非常有限的.
我們是有限的人.
我們的能力是非常有限制的.
所以我們要依靠上帝.
上帝就成為我們的保障.
這裡講到上帝很偉大.
講到我們多麼無能.
但是同樣我們經歷到上帝.
給我們這些願意謙卑下來的人.
有保障,有能力.
這一班就是那些能夠親近上帝.
能夠了解上帝的權能的人.
詩篇第八篇.
我們兩個星期前.
聽謝霆博士的研究.
我們就說有一種進路.
現在叫新的進路.
其實都不新.
二十幾三十年了.
新的進路.
看詩篇的時候.
不可以看一篇.
要看旁邊附近左右.
第七篇,第九篇等等的詩篇.
看他有什麼要講.
詩篇第八篇.
大家知道詩篇.

$^{401}$一百五十篇.
由兩篇開始.
第一篇,第二篇是悼論.
然後三至七篇就是大衛之事.
三至七篇都是大衛的哀歌.
因為他是一個無能之輩.
他就祈求上帝幫助他.
因為他沒有能力去打贏他的仇敵.
仇敵就追殺他等等.
三至七篇都是講這件事.
大衛這位君主.
受高之君.
他仍然要依靠上帝.
他知道上帝才能給他權能.
開路給他.
他就去祈禱去祈求他.
然後他就經歷上帝.
給他生命裡面的逆轉.
然後九至十四篇.
又回到講一些的Lament.
一些的哀歌.
就是講到詩人,講到以色列人.
那些信靠上帝的人.
他們願意去謙卑他們自己.
知道自己能力不是從自己來的.
他們依靠上帝.
以至他們能夠經歷上帝的大能.
第八篇就夾在中間.
第八篇告訴我們.
當我們見到上帝的大能.
見到大自然的時候.
天地的浩瀚.
我們就是awestruck,驚訝.
我們會震撼.
震撼於上帝自己的權能裡面.
同時我們也謙卑我們自己.
你看見日蝕也好.
看見宇宙的偉大也好.
應該是看見我們自己的渺小.
看見我們自己太多軟弱的地方.

$^{441}$我們是有很多不能的.
我五十有多.
三個星期前.
我突然背痛很嚴重.
我用背鏡.
我也不知道發生什麼事.
我也不是很勤力工作.
很痛 痛了一個多星期.
我貼了一些貼紙.
貼到我的皮都癢了.
出了疹子.
我五十有多.
我自己心想.
五十有多的人.
我的身材也算是這樣吧.
我覺得已經很fit了.
五十有多的人.
很驚訝吧.
很驚訝.
我心想為什麼會背痛呢.
令我很煩.
因為我做不到事.
唯一的結論是什麼呢.
唯一的結論就是.
老了.
還沒那麼快吧.
比你還要快.
但你知道.
我了解的是什麼.
上帝給我們的健康.
上帝給我們的能力.
不是從我自己來的.
有這樣的認同的時候.
你就發現.
上帝才是值得我們讚美.
上帝才是值得我們將權能.
榮耀 尊貴 歸給他.
如果不是他給我能力.
不是他給你能力.
我們怎麼賺錢.

$^{481}$我們怎麼能夠生活.
健康是從他而來的.
是吧.
你多麼著重健康.
我告訴你.
我很著重健康.
有眼力你看吧.
但是最終給力量給我的.
不是我自己.
我們本身就是脆弱的.
上帝才是我們保障和權利.
所以這篇詩篇一開頭.
講到上帝的權能.
上帝的價值的時候.
然後他就反映到我們自己有什麼價值.
八章第三至第四節.
他說當我看見.
或者你手指點出來的日月星辰.
並你所設立的月亮星辰.
他就反映出一件事.
領悟到自己.
他就看上天的時候.
其實就好像看一塊鏡子.
你就看見你自己人算什麼.
你竟眷顧他.
世人算什麼.
你竟眷顧他.
這裡用了兩個字.
第一個就是人字.
你說有什麼特別.
人字.
Anoche.
希伯來文.
其實那個字的意思就是說.
是脆弱的人.
就是Mortals.
NLT翻譯出來.
Mere mortals.
意思就是什麼.
終會死去的凡人.

$^{521}$就是這樣翻譯.
Man這個字.
What is man.
人算什麼.
其實就是Mere mortals.
我們算什麼.
廣東話翻譯.
我想廣東話才能翻譯出來.
就是死硬的人.
死定了.
我們是非常脆弱的人.
然後他說世人算什麼.
Son of man.
原文已經給你了.
其實我們只懂希伯來文.
Adam這個字.
Son of Adam.
Man Adam.
Son of Adam.
其實就是世人這個字.
Adam之子.
意思就是說什麼.
就是凡人的意思.
是普通人一個的意思.
希伯來文很妙.
有些音很像.
但是又有意思.
Adam這個字.
其實就是和哪個字相似.
就是和地下.
Adam.
我是Adam.
下面就是Adam.
當你死的時候.
你再進去Adam.
哇 Adam.
我們和地上的塵土很有關係.
塵歸塵 土歸土.
這就是詩篇裡面所說的.
Paradox of human being.

$^{561}$就是人生的弔詭.
不是矛盾而是弔詭.
就好像一個硬幣有兩面.
一方面我們是無能之輩.
但另一方面就是詩篇所說的.
居然我們在上帝眼中是有價值的.
我們多麼脆弱也好.
多麼覺得自己不重要也好.
或者我們整天覺得我們不適合.
我們沒有別人那麼厲害.
我們可能會問我們自己個人的價值.
懷疑我們自己的個人價值.
但這個詩人說什麼.
大衛說見到天上的上帝大能的彰顯的時候.
他發現原來上帝去眷顧和顧念我們.
我們算得是什麼.
顧念他.
你竟顧念他.
其實就是Remember.
Mindful of.
記起那件事.
整天想著他.
我們中間有很多做家長做父母的.
小朋友還小的時候.
為什麼你睡得不好.
不是因為他吵著你那麼簡單.
因為你Mindful.
Mindful of the baby.
整天想著他.
他有沒有事呢.
他會不會有病呢.
他哭什麼.
You are mindful of human being.
上帝整天想起我們.
當你不開心.
當你有憂困的時候.
上帝是Mindful of you.
當你開心的時候.
上帝又和你一起喜樂.
第二個字.

$^{601}$眷顧.
是另一個字.
眷顧是什麼意思.
有很多翻譯.
其實最主要的用法是什麼.
大家以前有讀過小學中學.
一進到班房的時候.
就是Roll call.
叫你名字.
跟著學生的名字.
我很清楚我姓陳.
經常都叫我先.
C嘛.
叫你名字.
告訴你你在不在.
這個就是眷顧的意思.
就是他Roll call你.
Take attendance.
就是說你有沒有到.
你有到就真是好了.
這個就是上帝對你的意念.
就是我們有這麼高的價值.
居然我們有這麼高的價值.
我們繼續看下去.
有什麼價值.
他說你是他人類.
人啊.
我和你.
比上帝或者天使微小一點.
似他榮耀尊貴為冠冕.
你派他管理你所做的.
使萬物都服在他腳下.
就一切的牛羊田地裡面的畜生.
還有空中的飛鳥.
海裡面的魚.
飯又在水裡面的.
你都給他們做什麼.
派他管理.
去管理上帝所做的一切.
這個paradox(幻想)多大.

$^{641}$我們雖然是死硬的人.
雖然我們是脆弱的人.
我們的能力不是我們自己來的.
但這裡說出什麼呢.
就是上帝做人的時候.
居然給他一種榮耀尊貴為冠冕.
就是在上帝心目中.
你的出現.
你的生存.
你為什麼會在這裡.
your being.
不是一個意外.
不是突然間想了想.
不如加個輔助下去.
多你一個,少你一個,不少.
不是的.
這裡說上帝去做我們.
比天使少一點.
你怎麼翻譯都好.
Elohim這個原文.
你翻譯做神也好.
翻譯做天使也好.
這裡怎麼都在說.
在上帝眼中人有什麼價值.
就是在這個創造界裡面.
最高的價值.
沒有一樣東西比人類更高的價值.
世界上有任何不同的活物.
可能是很奇妙的.
鯨魚也好,雀鳥也好.
微生物也好.
是多麼精密的結構.
但都不夠你在上帝眼中.
有Glory and Honor.
你知不知道在聖經古代的時代.
那些人都會有些故事.
告訴自己.
我們自己的價值是什麼.
比如巴比倫,埃及等等.
他們自己創造了神話.

$^{681}$神話說什麼呢.
大致上就是在說.
人類為什麼會在這裡.
人類在這裡.
原來天庭上的眾神.
談不攏.
沒有人服侍他們.
小的反叛.
小的神明.
然後大家反叛的時候.
大家不要吵.
大臣就說不如這樣.
我們做些人.
當他是奴隸.
有些飛蛹等等.
弄他出來.
那些就是人類.
就是幫我們打工的.
初一,十五就給我們吃東西.
弄個白切雞給我們.
我們就大家在上面有Vacation.
Essentially就是近東的神話所說的.
你想想這些神話,這些故事.
人類有沒有價值.
人類是沒有價值的.
人類沒有價值.
所以我們廣東話是人渣.
其實每個人都是人渣.
你人渣,我人渣.
因為我們每個人都是沒有價值的.
但是在上帝眼中.
在聖經裡面的這位上帝就不是.
而是在上帝眼中.
他是精心去設計.
我們的存有.
是在於這位愛我們的上帝.
A little bit less than the angel.
這個就是你的worth.
和我的worth.
我不知道你有沒有check過.

$^{721}$你自己的net worth有多少.
有誰check過.
自己的net worth.
Elon Musk有多少net worth.
百多個billion.
他一滴水都淹死你.
check一下自己的net worth.
但是聖經裡面告訴我們net worth是什麼.
Crown with glory and honor.
有榮耀尊貴為冠冕.
好像做皇帝,做皇后一樣.
並且上帝派我們去管理.
他所做的一切.
人類就是這麼獨特.
我們和其他活物是完全不同的.
因為我們是按照上帝的形象所做.
所以我們在這個地上.
我們是代表天上的上帝.
是representative.
是一個代表.
是一個使者.
這裡詩篇第八篇.
就是九節我們看完.
其實就是創世記第一章的總結.
用了九節的聖經.
就去講出創世記第一章.
裡面某一些的真理.
就是說原來上帝是很佩得我們敬拜.
而我們的價值雖然是脆弱.
但是上帝給我們榮耀冠冕.
然後他就告訴我們.
原來我們在這個地上.
是有上帝給我們的使命.
我們可以有這樣的福分.
有這樣的機會.
參與上帝去管理萬物的工作.
這個就是人有的特別的地方.
所以結束的時候.
我想告訴大家.
我們是多麼的獨特.

$^{761}$大家記得我曾經講過的這幅相.
1990年2月14日情人節拍的.
那個Voyage 1.
射到深大空.
然後回頭拍了一張相.
地球就是那一點.
淺藍色的那一點.
原來人的價值是這麼微小.
你就在那裡 我就在那裡.
我們算得了什麼.
但是詩篇第八篇就說.
原來我們是完全不同.
我們在上帝眼中是有榮耀,尊貴為冠冕.
詩篇第八篇.
告訴我們你是尊貴的.
你回去照照鏡子.
你想想你自己.
我可能不夠別人厲害.
或者我網絡不夠別人多.
但是我在上帝眼中.
是有獨特的價值.
獨特的價值.
第二方面.
告訴我們我們有照顧神的國度.
大地就是上帝的國度.
我們有這樣的照明.
當你今天上班.
你在人類的社會裡面.
去令到他運作.
這個本身已經是有價值的.
就是上帝的創造裡面.
應該有人去參與.
一起使他運作.
使他能夠增加他的價值.
第三方面.
就是上帝是值得我們去敬拜.
包括我們宿命上的使命.
其實我們都是參與上帝在地上有的工作.
你說有多榮耀,尊貴.
上帝不需要你和我將福音傳給別人.

$^{801}$他可以自己搞定.
上帝不需要你和我去建立教會.
他可以找天使去做.
但是他卻將這個照明放在我們身上.
叫我們能夠去事奉.
我們有的恩賜.
好少也好.
多也好.
我們都有這樣的尊貴.
可以參與在上帝的工作裡面.
這個就是我們有的榮耀和榮耀.
只不過我們有時.
只是留意著自己的工作.
只是留意著自己眼前的事.
我們不會關注其他事.
《詩篇》第八篇就是喚醒我們.
原來我們是參與上帝在地上的工作.
包括將福音傳給別人.
幾個星期前.
有三個漁民在南太平洋捕魚.
你有沒有聽過這宗新聞.
然後他們的船壞了.
然後不知道怎麼漂浮.
他們就漂到一個很小的島上.
在哪裡附近.
你知道有個地方叫做米古尼西亞嗎.
你知道在哪裡嗎.
米古尼西亞.
屬於美國領土.
在南太平洋.
我跟我老婆說你知道米古尼西亞在哪裡嗎.
我用Google Earth告訴她在哪裡.
她說你怎知道.
你怎知道有這樣的東西.
我說我經常看地圖.
所以我知道在哪裡.
在一個很小的島上他們漂浮了.
沒有人理會他們.
只有三個人.
如果你在羅洛荒島的時候.

$^{841}$你要做什麼.
吃椰子.
一味吃椰子.
他們就是這樣.
在一個多星期就去吃椰子.
然後在沙灘上他們就找椰子樹.
弄了一個help這個字.
希望有人去救他們.
誰知道就是美國的Coast Guard.
那些軍人等等.
他們就在上空經過.
然後見到他們.
help這個字.
他們pay attention to他們.
然後他們就派人去救他們.
一個這麼渺小的人.
派出救援的隊伍來救他們.
這個不是像上帝看我們一樣.
對不對.
我們以為在大地裡面.
我們只不過是一粒塵.
但上帝祂眷顧我們.
但當我們參與在上帝.
在地上做的工作的時候.
譬如我們將福音帶給別人.
將我們的見證帶給別人的時候.
我們都要pay attention.
其他人.
是需要我們的幫助.
誰去代表上帝.
告訴他們上帝是真實的.
你和我.
有這樣的尊貴和榮耀.
對不對.
是不是上個星期.
台灣花蓮大地震.
你知道台灣.
大家應該很熟悉.
台灣是佛教或者道教.
很旺盛的地方.

$^{881}$當然他們的救援是大把錢.
突然有些教會.
雖然是很少.
他們都在聯合.
去做一些救災的工作.
去take care那些小朋友.
婦孺等等.
這個就是基督徒在地上.
我們能夠去參與上帝.
工作的一幅圖畫.
上帝是值得我們讚美和敬拜.
但我們的崇拜.
只是說我們實體的罪惡.
是真的將尊貴榮耀歸給.
配得我們讚美的上帝.
離開這裡的時候.
我們知道我們是多麼的尊貴.
我們就是代表上帝.
走進我們的生活.
每一個角落.
你的家庭,你的職場等等.
不同的地方.
上帝是值得我們去敬拜.
我們應該再重新回憶.
什麼呢.
那種All and wonder.
上帝那種偉大.
以致我們的敬拜.
是真的能夠發出那種奇妙.
上帝是多麼偉大.
和愛我們多麼深.
耶穌基督來到拯救我們.
我沒有去全日食的餐廳.
我後悔.
我應該去的.
誰知道我沒有去.
Anyway.
我兒子去了.
我兒子在香檳.
他約了幾個朋友.

$^{921}$好像開了兩個多小時.
去了一個地方叫Paris.
不是法國的Paris.
Paris, Illinois.
不知道是Backwater Town.
他們開了兩個多小時.
他問我去不去.
我說好啊.
我們是酒堂.
我說好啊.
酒堂啊.
值得去嗎.
我說儘管去.
他們大調導你去.
然後他就拍了不同的東西.
然後發了一些影片給我.
我看了他的影片.
我做的短片.
其中最後一段就是他們拍的.
我聽了他們的影片.
我覺得很有趣.
因為聽見他們說什麼.
因為拍著太陽.
聽見他們在說話.
他們三分多鐘看的日食.
全部都是.
哇.
哇,好酷.
三分多鐘就聽.
哇,怎麼說這麼多次.
我變成那種奇妙.
對上帝應該要有.
我們要重新回應那種奇妙.
哇,對不對.
每一個敬拜應該是這樣的.
整個鐘就是.
哇.
好酷,對不對.
好酷.
上帝滿有尊貴.

$^{961}$在他眼中.
我們也有尊貴.
他值得我們去敬拜.
我們同心大同祈禱.
上主.
我們今天來到你面前.
我們要敬拜.
作為你的教會.
作為你自己的子民.
我們能夠體現到你的愛.
你愛的偉大.
你恩典的廣闊.
你在宇宙間所顯出的大能.
主,我們感到奇妙.
我們應當深感到震撼.
震撼於你自己的權能.
震撼於我們幾個微小.
但你卻看顧.
甚至你自己的兒子來到這個地上.
為我們捨身.
以致我們心靈能夠再不同.
給我們永生.
給我們一切屬靈的好處.
上帝我們在你面前.
感歎.
主讓我們的心靈能夠顯出你多一點的榮耀.
主,差遣我們出去.
以致我們和我們的兒女和我們配偶和我們家庭和我們身邊所有的人.
主,我們帶著這個榮耀冠冕.
去代表你.
去事奉你.
願上帝祝福各位.
拜拜.
\newpage



\section{}
\label{sec:yghFb3uoZ_c}
\textbf{11/11/18 主日講道 (Pastor Tony Wong)}
\newline
\newline
連結: \href{https://youtube.com/watch?v=yghFb3uoZ_c}{\texttt{https://youtube.com/watch?v=yghFb3uoZ\_c}} ~~~~ 語音日期: 2024-04-19
\newline
\newline
\hyperref[sec:DfbxT6XIBuk]{\small{< < < PREV SERMON < < <}}
~
\hyperref[sec:index_chronic]{\small{[返順時目]}}
~
\hyperref[sec:index_scriptual]{\small{[返順卷目]}}
~
\hyperref[sec:GRZOoQIm4Fg]{\small{> > > NEXT SERMON > > >}}
\newline
\newline
$^{1}$當我和Mark Sir聊天的時候.
我說你們教會挺好的.
挺多年輕人.
他說不是吧.
不過我和Mark Sir說一件事.
我在現在的教會.
已經大約在兩年.
在這兩年.
我帶了兩個家庭信耶穌.
是因為這兩個家庭的人信耶穌.
我就成為了教會的中文堂.
不再是教會的中文堂.
我帶了兩個家庭信耶穌.
不管怎樣.
帶了人信耶穌.
是很開心的.
不管他是三十多歲.
或者七十多歲.
我大約是.
應該是今年的.
是今年的復活節.
幫一對夫婦受洗.
七十歲.
很開心.
幫他們受洗的時候.
我想起我自己的生命.
我在1988年的時候.
去接受耶穌基督.
成為我個人的救主.
那時候的感受很深刻.
因為經歷到耶穌基督.
在十字架上的犧牲愛.
我就接受了這個救恩.
我還記得在哪裡接受救恩.
如果你在香港長大.
你應該聽過一個地方.
叫做烏溪沙.
烏溪沙是一個青年營.
我幾年前再回到烏溪沙的時候.
真的完全不同了.

$^{41}$三十年前巴士都不能進營.
我要走進去.
三十年後我走進去.
已經被高樓大廈圍著整個營會.
景物變了.
但是心情沒有變.
因為想起了耶穌基督的愛.
但是這三十年信主的時候.
我又想起了很多不同的故事.
我想起了一些和我一起信耶穌的人.
信仰的路走著走著.
突然間消失了.
不見了.
我記得有一位弟兄.
他和我一起受洗.
我們一起上受洗班.
洗完之後我還在.
他就不知道去了哪裡.
我有時候都說你是不是一淹死了.
我有點害怕.
不是吧.
就算我也認識一些人.
就算他仍然在教會裡.
但是他的信仰生命.
屬靈生命.
一落千丈.
但是與此同時我又認識一些弟兄.
姐妹.
他們的信仰不單止沒有停頓.
而且不斷成長.
成為別人的榜樣.
大約一個月之前.
Mark Sir說我剛剛.
來了三德神學院讀書.
然後飛回多倫多.
用了48小時探望父母.
在探望父母當中.
有位弟兄在Facebook上.
看到我,放了一張相片.
然後馬上把我拉出來.

$^{81}$聊天的時候我很開心.
因為這位弟兄是我們以前在教會.
籃球隊一起打球的.
他做中鋒.
你猜猜我做什麼位置呢.
我做Power 4.
我經常說我們打球還是打人呢.
但是回望今天.
他剛剛神學院畢業.
做傳道人.
我和他聊天的時候.
我們在戰場上一起並肩作戰.
今天就不在球場.
在淑玲的戰場上繼續並肩作戰.
我問一個問題.
為什麼有些人.
同樣都是信耶穌.
有些人.
繼續不斷地成長.
但有些人就退場了.
甚至消失在教會.
究竟發生了什麼事呢.
說真的,是不是我完美呢.
也不是,回想自己的基督徒生命.
我也是有起有跌.
我不斷地問自己一個問題.
什麼使我繼續信靠耶穌.
在信仰的路上不放棄.
在生命上繼續為主耶穌.
勇往直前.
如果我今天.
問你.
在座當中的每一位.
你是因著什麼原因.
繼續相信耶穌.
也不放棄這個信仰.
繼續在你的生命上.
為著耶穌勇往直前.
為什麼.
不知道.

$^{121}$如果真的有一個人在你面前.
這樣問你的話.
你會怎樣回應他呢.
請我們一起打開.
今天的聖經.
約翰一書.
第一章的一至二字.
約翰一書是一本很有趣的書信.
有趣在於.
這本書信.
是試圖約翰.
回應的書信.
他在約翰一書.
回應的書信.
回應什麼呢.
當時有學者指出試圖約翰.
大約在以弗所.
一大地方的牧羊教會.
但他發現牧羊教會在小亞細亞的地方.
當時的教會.
正受到一些人.
傳一些不正確的訊息.
傳入教會當中.
我們稱為疑端的訊息.
傳入教會.
使教會的團契生活.
信徒生活都受到破壞.
約翰.
作為一個牧者.
見到這樣的情況.
他就執起筆.
守著聖靈的帶領下.
寫下約翰一書.
教導信徒在信仰上.
在見證上.
和團契生活上.
如何回到基督信仰的根基.
根本.
而約翰一書.
今天的經文一章一至二字.

$^{161}$其實是約翰一書一章一至四段.
整段經文的其中一半.
一章一至四字.
是整段約翰一書的序言.
如果你能夠.
讀希臘文的時候.
你會更加發現.
一章一至四字是一句說話.
不是幾句.
是一句說話.
我們今天只看這句說話的一部分.
只是兩節.
但我們看這兩節聖經的時候.
我們會發現.
這段經文提醒你.
提醒我.
如果我們合情合理.
去接受基督信仰的時候.
我們的生命.
是能夠可以.
被上帝更生.
我再說.
如果一個人願意合情合理.
去感恩.
接受耶穌基督的福音的時候.
我們的生命.
是可以能夠.
被主耶穌基督更生.
首先如果我們要成為一個.
經歷上帝更生的.
真信之人.
我們需要的是.
合理地去接受.
這個福音.
經文的第一節.
我們看到.
論道起初原有的生命之道.
就是我們所聽見的.
所看見的.
親眼看見的.

$^{201}$親手摸過的.
如果你細心讀這段經文的時候.
我們會發現.
我們所接受的福音.
其實是合理的福音.
這個福音為什麼合理呢?.
就是因為這個福音.
基督的教導.
和他所成就的救恩.
是真真實實.
在人間裡面.
發生過.
你留意.
在經文裡面約翰說.
生命之道是甚麼?.
聽過,看過,摸過.
這些都是表明.
這一班門徒.
和約翰自己.
是親身接觸過耶穌.
有些學者說.
這個生命之道就是耶穌基督.
本身和他所傳的一言一行的教導.
但同時.
約翰很有趣.
為什麼不寫耶穌.
而寫起初原有的生命之道呢?.
很有趣.
為什麼要加起初原有呢?.
這是想表達.
耶穌基督在世界未創造以前.
已經存在.
為什麼有這個需要要說呢?.
為什麼約翰要加幾個字呢?.
是不是約翰手癢呢?.
相信不是.
其實約翰那時候.
特意要加這句說話.
是有意思的.
因為當時的教會.

$^{241}$是受著一個異端的思想.
洛斯底主義影響.
基督教的信仰.
他們不相信耶穌基督.
真的造成了肉身.
因為他們覺得.
只有靈魂是最寶貴.
肉身是邪惡.
所以當耶穌.
在地上的時候.
他們不覺得耶穌是一個人.
就像中國人說的.
鬼上身.
聖靈上了一個肉體的身.
所以肉體所造的不是耶穌.
只是靈魂.
無論怎樣說.
是當時的一個想法.
約翰知道這個想法有問題.
所以約翰就很清楚.
告訴當時的信徒.
耶穌在創世的時候.
已經存在.
而且他還要是道.
成了肉身.
成為了約翰和當時.
其他的門徒.
可以活生生接觸到的.
一起生活過.
經歷過.
見過他的復活和死的耶穌.
耶穌所傳的福音.
在耶穌自己的身上.
都成就了.
耶穌自己.
傳揚天國的訊息.
叫人要從罪當中.
回轉過來.
還有當耶穌.
成為一個人的時候.

$^{281}$這個耶穌.
他將他的主權歸給.
天父上帝.
還有就是.
他自己死在十字架上.
和復活.
他使得接受救的人.
能夠與神和好.
如果耶穌基督.
不是道成肉身.
如果耶穌基督不是神成為人的話.
他根本救不了我們.
他根本就不是救贖主.
因此.
約翰在這裡.
不只是為了一個很正統的神學組的辯護.
約翰在這裡.
為一個他見過.
經歷過的救主耶穌.
親手摸過.
親手見過的耶穌.
去做見證.
這個信仰合理.
是因為耶穌.
是因為這群門徒.
經歷過,見證過.
對我們來說.
真實的信仰是需要合理.
是一定要合理.
因為合理的意思就是.
我們要按照事實.
來接受一個信仰.
假若我們的信仰是基於錯誤的內容.
我們的信.
都是一個錯誤的信.
舉個例子很簡單.
我經常說我這陣子坐飛機.
坐到我有點害怕.
一年坐六次不是開玩笑.
但是.

$^{321}$所以如果我拿著我的護照.
去到加拿大的美國海關.
遞給他看.
看到我的學生簽證.
看到我的護照.
有我這樣的樣子.
他讓我進來是一件正常的事.
很正常的.
但如果有一天我不知道發生了什麼事.
我拿著Mark Sir的護照.
走上去.
美國海關就說.
我是陳崇基.
你讓我入境吧.
我想海關就會想.
有一個問題就是.
第一.
我沒有那麼囂張.
我也認同.
第二.
你好像姓王.
那個護照姓陳.
是不是.
但如果我是一個.
錯誤的思想.
就引起了我一個很錯誤的決定.
拿著Mark Sir的護照.
過關.
所以.
任何的信.
是必須基於真實的事實的根據.
我們從聖經的自證.
到歷史和考古學.
所得出來的資料.
我們都能夠確認.
耶穌是死在十字架上.
也復活過來.
正如聖經所宣告.
耶穌基督是神.
成為了人.

$^{361}$祂為我們的罪死在十字架上.
祂所流出來的血.
為你也為我.
付上罪加.
使得我們接受這份救恩的人.
能夠與神和好.
祂從死裡復活.
叫我們信主的人.
從信主那天開始.
我們就能得到新的生命.
舊有的生命.
就是那個叫我們遠離神.
抵擋神的意念的舊生命.
都已經被釘在十字架上.
我們舊的生命.
因著被釘在十字架上.
就死去.
有時死得慢一點.
但是就死去.
但是那個新的生命.
就是我們願意信服神.
願意跟隨神指引的生命.
慢慢在我們生命中生出來.
我們靠著信服聖靈.
在我們生命中的引導.
叫我們活出一個新我.
自死救我的生命.
我們因著福音的緣故.
被召成為教會.
一個屬命的家.
一起經歷學習一個新的生命.
的群體.
使這個群體能夠彼此相愛.
使這個群體能夠有能力.
去傳福音去關懷社會.
不代表我們的基督徒生命會一帆風順.
但是我們在等候主再回來.
是積極地等候主再回來.
等到與主同在.
與主一起作王的時候.

$^{401}$我們不再有罪.
被痛苦,被罪惡,被死亡.
去傷害我們.
神在大地作王.
祂的愛,祂的公義,祂的榮耀.
都在全地當中.
在那一刻我們真的經歷了.
神與人的和好.
完全的和好.
神與人之間的和好.
人與人的和好達成.
人與自己的和好也達成.
最後人與世界的關係也和好.
我們將來的新天新地所經歷的生命.
是真善美的生命.
我剛才說了幾分鐘的東西是什麼?.
就是福音.
就是一個美好的福音.
這個就是從歷世歷代一代傳一代.
傳給我們的福音.
我今天想問大家.
究竟我們的信.
是不是基於.
剛才所說的真理.
而去相信這個信仰呢?.
我們的信仰是不是基於聖經的啟示呢?.
我們的信仰是不是真的經歷.
是不是因著耶穌基督的救恩組合的核心呢?.
我再說一個錯誤的信仰.
會使我們有一個錯誤的生命觀.
我舉一個很簡單的例子.
大約幾個星期前.
我曾經去跟一位朋友聊天.
這位朋友在報道會中缺了字.
在一些名人所說的報道會中缺了字.
我很開心.
他不是回我教會,只是有人介紹我認識.
我跟他聊天.
我問他:你為什麼信耶穌?.
他說:多好,不錯,多舒服.

$^{441}$跟信觀音沒有分別.
我已經呆了.
我跟探訪隊的弟兄.
打了眼色,然後不肚子餓了.
我問他:弟兄,這位先生.
在聚會中,你有沒有聽過什麼聖經?.
什麼聖經?.
我已經呆了.
我問他:這位先生.
你有聽過什麼耶穌的十字架嗎?.
有嗎?.
我心裡只是故作鎮定.
我告訴你,我那一刻是故作鎮定的.
我心裡流了冷汗.
我心裡想:是誰說的?.
為什麼會這樣?.
是不是他喝醉了,眼睛聽不到?.
我不管是說了問題還是聽了問題.
我立刻跟他傳了一次福音.
真的很簡單,將福音再說一次.
說完之後,我問他:這位先生.
你的決定是怎樣?.
他就說:我要先想一想.
但我接著說了一句話.
你想一想也好.
我起碼今天肯定幫了你一件事.
不會覺得自己買了一張機票.
以為上硬機.
誰知道原來去到機場的時候.
那張機票是假的.
大家聽不明白我的意思?.
真信是基於事實的.
我們的信是不是基於事實呢?.
我們需要合理地去信耶穌.
的確,如果我們要成為一個真信的人.
經歷生命根深的一個基督徒.
我們真的需要合理地去接受這個救恩.
但同時我們也需要一個.
怎樣呢?合情的態度.
合情的經歷.

$^{481}$去體驗耶穌基督的福音.
如果剛才我提及的經文第一,第二節.
我已經提及過,在經文裡面提及過.
這些門徒跟耶穌的關係.
是他們親耳.
對不起,是親耳.
我們是聽過,看過,親眼見過,親手摸過.
這是講述他們第一生的經歷.
是一個很親身的接觸.
而在第二節的講述是.
「這生命已經顯出來,我們也看見過.
現在又作見證.
將原來與父同在.
且顯現與我們那永遠的生命傳給你們」.
你可以這樣說,你細心留意的時候.
這段經文不僅僅說明.
這班門徒是理性上知道信仰.
亦在他們的經歷上體會過.
約翰的見證這麼有能力.
是因為他的信仰從來不是二手信仰.
不是借回來的.
而是他親身體驗和耶穌的關係.
而且約翰和他的門徒.
不單單在耶穌在世的時候經歷.
當耶穌死於福音,活在聖天的時候.
如果你讀聖經就知道發生什麼事.
聖靈降臨與他們同在.
基督的靈降臨在他們身上.
使他們每天都可以經歷神的同在.
我們可以說約翰的信.
不僅僅是一個合理的信.
亦是一個合情的信.
他的信仰合情.
是因為他不單單靠他的頭腦對福音的認知.
而且是他信靠這位救主.
他和這位救主經歷人生當中的一切.
是起起跌跌得得失失.
而且在這個信仰上成為一個有關係的信仰.
有經歷的信仰.
有感情的信仰.

$^{521}$有體會的信仰.
從而經歷過神的愛.
去愛主更深.
願意信服這位救主的教導.
對我們來說.
我們的信不僅僅是要合理.
亦要合情.
去經歷神的靈在我們生命當中.
每一天共在.
可能我們再問.
可能對我們自己再問一個問題.
其實我們可以對一個福音信仰有所理性上.
在理性上可以有所認同.
我們可以認同它.
但是不等於我們和這位救主建立了一個很親密的關係.
更不要說我們順從聖靈.
在我們生命當中的帶領.
我只想到一個例子.
做牧師其中一件事.
就是幫人證婚.
你試想想有一對男女在我面前.
我要和他們證婚.
通常牧師會問那對男女什麼問題.
一定有一個問題一定要問.
某某某你今天是否願意娶這個人.
或者嫁這個人.
通常答案就是.
你要懂得.
你可能聽過.
我願意.
假設我在一個婚禮中.
我這樣證婚的時候.
突然間那個男的不是說I do.
而是他冷冷地說一句.
我認同的.
你猜猜新人會怎樣.
我在想可能新人已經呆了.
不懂得反應.
可能說我為什麼會選擇這個男人.
我做牧師會怎樣.

$^{561}$我可能會問那個男士.
你是不是一時口快說錯話.
I do和I agree.
是很不同的.
我願意和我認同是很不同的.
不同在於什麼.
I do我願意是包含一份毀心.
包含一份信心的投放.
I agree我認同不一定有信心.
不一定有信心的投放.
我講加拿大政治比較好.
不要講美國政治.
我認同.
我一個朋友.
我認同我們現在總理所做的事.
認同.
他認同而已.
但我問他你投票是不是投他.
不是我打算投給反對黨.
是不是.
認同和你真實去做的事.
可以是兩回事.
但我們的信仰.
他需要的不是認同.
他需要的是我們對主耶穌說.
I do我願意.
各位.
我相信大家相處已經一段時間.
或者接觸信仰.
但我想問一個問題.
究竟你的信仰或我的信仰.
是否貼身的信仰.
顧名思義.
我們的信仰不是一個經歷神的信仰.
我們現在信主的光景.
是否只是停留在理性上的認知.
而沒有打開自己的心去經歷神.
上帝耶穌基督.
如果他呼召你.
要為他做一些事.

$^{601}$你不可以只認同.
你可以認同的.
但你可以什麼都不做.
但耶穌要求什麼.
實踐和實行.
各位.
你知不知道理性上認知穿上救生衣.
和真實地穿上救生衣的最大分別是什麼.
我想問在座的有沒有人是不懂游泳的.
我懂游泳的.
但不懂游泳的.
這個例子你會很深刻.
如果你不懂游泳.
你又一個不經意滾下水.
假設水溫夠暖.
你不會冷死的.
假設水溫夠暖的話.
有救生衣和沒有救生衣.
是否差很遠.
差很遠.
有救生衣水夠暖.
穿著也會浮浮浮.
最多等人將你救起.
沒有救生衣水溫夠暖.
但你不懂游泳.
就像石頭一樣浮浮浮沉下去.
各位.
真正的信是一份經歷.
真正的信是一個體會.
在我們這個年代信耶穌.
我們可能對很多事情.
我認同是這樣.
但我不得不承認我們很多事情.
是有保留的.
有沒有?.
有.
你可能問我牧師你有沒有?.
你猜呢?.
當然有.
我也曾經經歷過.

$^{641}$有抓得很緊的時候.
但問題在於.
原來當我們抓得太緊.
不肯放手的時候.
可能你真的未經歷過.
上個救恩的戰利.
我理性上或理理性上.
知道他是誰.
但你根本與他無關.
我很喜歡.
在約伯記有句說話.
我從前風聞有你.
我今日親眼見你.
各位.
我們的信.
是否只是停留了風聞有你.
而沒有親眼見你的部分.
這不只是你決志的那一刻.
而是你基督人生中.
不斷不斷的一個基督徒的生命路.
可能你信了很多年.
可能到了某一刻.
你跟耶穌說.
耶穌不如我現在跟你請假.
我跟你請假.
你讓我做我自己的事.
我講這件事.
其實我引用一個人的例子.
是一個黑社會大哥.
信耶穌的例子.
他跟耶穌說.
你讓我賺夠錢.
我跟你請假.
但最後結果是耶穌不放過他.
要他降服.
最後放下一切.
可能我們的生命沒有他那麼極端.
但每一個基督徒一生當中.
很多時候.
最大的挑戰就是.

$^{681}$放手.
如果信耶穌決志是第一步.
是生命開始放手的時候.
我們基督徒的生命是不斷放手.
是不斷放給耶穌.
但我們人的感覺就是.
很多時候要自己拿著.
今天你的手在哪裡.
是牽著耶穌.
還是只是抱緊.
抓緊一些東西.
死都不肯放手給耶穌.
去處理你的生命.
我們每個人都有生命的優點.
有生命的破碎點.
可能你的成長當中.
你的家庭是你抓緊的地方.
可能是你的婚姻問題.
令你抓緊的地方.
可能是你的事業和錢.
是你抓緊的地方.
可能是你的野心.
你抓緊的地方.
但耶穌說真信的人.
就是一個放手的人.
約翰親眼看見.
親眼看見.
你猜約翰有沒有被上帝雕琢過.
有.
約翰和他哥哥的花名是什麼.
雷子.
Son of Thunder.
很暴躁.
而且他媽媽曾經求過耶穌.
可不可以讓他兩兄弟.
一個服侍在耶穌左邊.
一個服侍在耶穌右邊.
權力抓得緊.
但由信耶穌開始.
一路放手.

$^{721}$今天的約翰是一個什麼約翰.
一個可以寫一本書信.
你想想.
雷子寫一封講愛的書信.
是一個很奇妙的事情.
一個這麼著重權力的人.
最終被人放到拔摩海島.
圍著信仰.
很多個放手.
信耶穌是一生的凝聚.
一生求主的改變.
真信的人.
決志的刻子是開始.
信一路帶你回到天下.
一路經歷放手的階段.
各位.
我們是否願意做一個.
願意經歷神.
信服神.
與主建立關係的一個信徒呢.
我現在再問你.
你因著什麼原因.
願意繼續信耶穌.
願意繼續不放棄這個信仰.
願意繼續為著耶穌基督向前奔跑.
你今天會再給我一個什麼答案呢.
我相信今天我們所讀的經文.
是給了一個提醒我們.
昔日約翰和他的門徒所傳的.
和當時人所接受的福音.
是一個合情合理的福音.
他們這個福音是合情合理.
因為他們知道所信的真實是真的.
而且是體會過的.
對我們今天來說.
我們能夠繼續在信仰上不放棄.
能夠繼續在信仰上為著耶穌基督向前奔跑.
是因為我們所接受的福音.
是一個合情合理的信仰.
我們的信是基於聖經的教導.

$^{761}$也是基於我們與主耶穌基督的經歷.
我們的信是合理的.
因為我們所信的是歷史上發生過.
而聖經上有啟示過.
我們的信是合情的.
因為我們所信的是我們每一天都願意.
當我們打開心扉的時候能夠經歷過.
也是我們放手過,信服過耶穌.
在我自己教會講道的尾聲.
我通常都會送一個問題給回眾.
一個星期慢慢回來想想.
我今天也送一個問題給你慢慢回來想想.
我們是否願意成為一個合情合理信耶穌的人呢?.
讓我們在這一刻.
在一個禱告的時間.
你和我一起去面對這位愛我們的耶穌.
我們一起祈禱.
我想給大家一點時間.
一分鐘左右.
你在上帝面前回應剛才聽到當中一點或兩點.
神向你講的說話.
以致你所聽到的.
其實不是聽人講.
而是回應上帝所對你說的每一句說話.
少許時間.
我們一起祈禱.
你真的道成肉身來到世界當中.
你的信是真實的.
因為我們在信諸經歷中經歷過.
我們所信的神是一個又真又活的神.
我們獻上感恩頌讚.
主耶穌我們幫助我們成為一個真信的人.
幫助我們所信的是按著真理去信.
真正認清知道所信的是真理.
但同一時間也幫助我們.
不單是在頭腦上知道.
是心願意打開接納你.
做我們的救主和生命的主宰.
去經歷去體會生命當中的改變.
主呀求你幫助我們.

$^{801}$我們可能還有地方不願意放手.
求主幫助幫我們放手.
以致我們能夠一心一意的去跟隨你.
經歷救恩.
不單是得救那一刻.
是經歷救恩到永生.
主呀多謝你.
\newpage



\section{}
\label{sec:GRZOoQIm4Fg}
\textbf{5/13/18 主日講道: 告訴子女兩句重要的話 (張麗仙傳道)}
\newline
\newline
連結: \href{https://youtube.com/watch?v=GRZOoQIm4Fg}{\texttt{https://youtube.com/watch?v=GRZOoQIm4Fg}} ~~~~ 語音日期: 2024-04-19
\newline
\newline
\hyperref[sec:yghFb3uoZ_c]{\small{< < < PREV SERMON < < <}}
~
\hyperref[sec:index_chronic]{\small{[返順時目]}}
~
\hyperref[sec:index_scriptual]{\small{[返順卷目]}}
~
\hyperref[sec:GXyrjafkb9Y]{\small{> > > NEXT SERMON > > >}}
\newline
\newline
$^{1}$(台北市立大樓).
我們一起祈禱.
天上的父神,你是這麼愛我們每一個.
即使我們死在罪惡過犯裡面.
你叫我們在你的愛子的寶血裡面.
能夠活過來.
願意我們今天為恩典而活.
我們除去舊人,穿上新人.
而且我們要教我們的下一代.
去明白這兩個這麼重要的道理.
使他們一生一世在你的恩典裡面活著.
求你教我們時時的反覆思想.
多謝你,奉救主耶穌的名,阿們.
\newpage



\section{}
\label{sec:GXyrjafkb9Y}
\textbf{Gaining a New Perspective Can Change Your Life - Keith Draper}
\newline
\newline
連結: \href{https://youtube.com/watch?v=GXyrjafkb9Y}{\texttt{https://youtube.com/watch?v=GXyrjafkb9Y}} ~~~~ 語音日期: 2024-04-20
\newline
\newline
\hyperref[sec:GRZOoQIm4Fg]{\small{< < < PREV SERMON < < <}}
~
\hyperref[sec:index_chronic]{\small{[返順時目]}}
~
\hyperref[sec:index_scriptual]{\small{[返順卷目]}}
~
\hyperref[sec:KF6pQDFyEx4]{\small{> > > NEXT SERMON > > >}}
\newline
\newline
$^{1}$Several years ago, a suicide occurred among a population that we had resettled through World Relief..
We were determined to unravel what happened, so what did happen?.
A senior man, greatly respected in the community, took his own life..
We gathered the grieving community together to see how we could help, and a young man said,.
"If this man's death surprises you, then you must not know us very well.".
Pretty damning statement to hear from someone, really, who has just gone through such pain..
So we asked him what we didn't know, and he went on to really say that suicide was very common for their people and why..
The source of their sad situation was a battle of peoples..
Their people alongside another people, until they became too numerous and they were perceived as a threat..
So they were deported and they were stripped of their citizenship..
They became a stateless people, having no place to call home, and this led them to have one of the highest suicide rates of any people..
Let's just say they had lost hope..
Today we live in a world that seems to lack hope..
People are at odds with one another, nations continue to war with other nations, encroaching on their rights and freedoms..
When added to this, a pandemic has ravished our world..
I can see why there is such fatalism in our world today..
Like the quandary of the stateless people that I mentioned earlier, is there any place that we can turn to in order to really find hope?.
David wrote the first two chapters of the book of Psalms as a prelude and a framework for the book to help us find direction in life..
Psalms are really songs that God's people wrote for us to sing, whether an anthem or a dirge,.
to pray back to God what's happening around us to help us sort out things in our life..
Interestingly, with every trial and temptation that Jesus faced,.
he confronted them with the Bible and the Psalms in particular with his crucifixion..
The Psalms are a complete musical repertoire, if you will, of themes and echoes, poetry and rhythms, parallels and contrasts,.
all to help us personally deal with life and also as a people to deal with life..
The main point of this second Psalm is this, we're all in, but really what side are we on?.
You see, beyond what we see physically, there is a spiritual war that rages against us behind the scenes,.
evil perpetrated by earthly leaders of countries and industries that cause harm to people who are themselves mere footmen,.
slaves, understudies for spiritual powers and principalities that are in the unseen world..
So who are these warring parties?.
Well, the Bible says that every person who heeds God's instructions is on his side and on the side of his anointed in the scripture..
In this worldwide struggle that we're in, they are the blessed, and really they've been given access to true happiness..
But to fully appreciate their blessed state, they must trust in God's anointed and find delight in his instructions..
Conversely, there are those who are perishing, those who live wickedly and rebel against the Lord,.
muddling through life in daily skirmishes for control of what's here and there..
Psalm chapter 2 assures us of the outcome of this battle, though, and shows us how to fight it successfully..
We need to understand this cosmic battle in order to better find our way forward here on earth..
At their core, both parties are really defined by their perspectives and the authorities that are in their life..
At the beginning of this chapter, we really see the perishing only having a horizontal perspective on life..
The nations, or pantata ethne, is the term there, really means the peoples..
The peoples are in rebellion. And who are they in rebellion against?.

$^{41}$Ultimately, God, though they don't fully recognize him or acknowledge him,.
and their rebellion causes them anger, a sense of futility, and frustration in life..
They are plodding along, not seeing little to any fruit from their labors,.
and they want to see themselves as masters of their own destiny, but they are stuck in the quagmire of this world..
The blessed, on the other hand, bring a vertical perspective to life..
Pictured here in this psalm, God is in heaven. He's in heaven laughing, he's scorning,.
and he's rebuking the rebellion that he sees, because he knows who is truly in charge..
He is placed on the supreme throne of this earth, his anointed one. And who is he?.
Well, in verses 7 and 8, we see that Jesus is that king, and we, the church, are his bride,.
having access to that power and that authority. Jesus proclaims the message that God gave him of good news..
He is indeed God's son, and as it says there, he's equal to him..
And God gives him authority over all of the nations and the peoples of this earth..
God wants to give the whole earth to him as his possession,.
but the Lord also wants to give the nations to his bride, the church..
At the end of time, the prophecy of Revelation chapter 2 quotes this psalm, saying,.
"At the end of time, to all who are victorious, who obey me, obey the Lord to the end,.
to them I will give authority over all nations, and they will rule the nations with an iron rod.
and smash them like clay pots.".
The critical tool to receive this blessing, though, is found at the beginning of verse 8..
He says, "Only ask." Prayer is key..
We must ask God for all that we need to see victories in life, whether they be big or small..
Because the Bible says God makes all things work together for good,.
to those who love him and are called according to his purposes..
So how should we respond in light of this incredible knowledge that we see here in Psalm chapter 2,.
and what should we lead others to do to act on it?.
He really speaks of four things here..
He talks about us acting wisely in the days that we're in, to serve him reverently,.
to rejoice with trembling, and to submit to God's Son in all things..
And why is this so? Because our actions are consequential..
Rebellion brings about God's anger and destruction,.
and that without notice, no matter who you are or what you're doing..
But there is joy, and there's peace for those who take their refuge in the Lord..
They find everlasting life, 24/7 protection, and abundance instead of scarcity..
So back to the stateless people story that I mentioned at the beginning of my time with you..
We said that they had lost hope, and a young man from the community provided a glimmer of hope in our meeting with them..
He mentioned how hopeful all the people were who had helped them as they were resettled here,.
and they all happened to be Christians..
He wondered if they might help him find hope and help his people find hope..
So what did we do?.
We extended a Macedonian call for help, and we found a person of peace who spoke their language.

$^{81}$and was willing to offer this same hope in God's anointed,.
like those who had helped them before when they were first arriving..
They learned that hope is found in God and his appointed ruler, the true ruler of the world,.
and that true ruler of the world is King Jesus..
This understanding of Jesus led these stateless people to form churches in their homes.
and eventually to launch public worship here in the Chicagoland area..
With their first baptisms just a couple of years back,.
a first generation of Christians was born among that population.
when very few, if any, in their country had experienced belief in Christ in the past..
And among those who were first baptized included an 82-year-old witch doctor.
who committed his life to Christ..
Today, there is an active movement of Christian community among this population,.
touching now three generations in a very short period of time..
I received this note recently from one of our youth workers working with their kids..
She said this, "Hey, not sure if you're all going to see this story,.
but I wanted to share it with you, Keith..
I just wanted to share a Little Friday encouragement..
This one is about the children from this community, and they are now high school kids..
It touches my soul. I hope it will touch yours..
I had a Nepali woman tell me this week that in the last week, the last year or two,.
that she and her two boys had decided on the Jesus..
It used to be the Hindu, and now it was the Jesus..
She said, "Since my boys listened to the Jesus, they no fight no more..
They no argue. They talk. They listen. They help..
I asked them why they now act so good..
They tell me, 'Mom, this Jesus is very, very good..
He no like the fighting. He only likes the helping.'.
The Jesus is now spreading among the youth in this community in a very big way..
They band together, worship together, sing and pray,.
and when they're done, they help each other with their homework..
There's so many things that COVID has torn apart in our communities, including among this one,.
but this is something this year that has really served to make things better..
For you see, they, their parents, and their grandparents had gained a vertical perspective on life,.
and it made all the difference for them in finding hope..
Do you know Christ's command to make disciples really flows from this psalm in chapter 2?.
And fulfilling the great commission to go and make disciples of all nations.
honors God's promise to Abraham to bless all the nations on earth through him..
How about you?.
Are you living in the horizontal spaces, fretting about life, causing you anger,.
maybe a sense of futility or despair?.

$^{121}$Maybe you're even plotting with others, hoping for a better days ahead..
Are you ready to look up instead of looking around you?.
God is large and he's in charge, and he's given his son to rule so that we might come to trust him..
He wants us to go to a place of peace and his shalom, where all things are right under his rulership..
In his presence is fullness of joy, and all who take refuge in him will see beyond what is around them to what is over them..
The Bible says a great cloud of witnesses went before us,.
and we also find in Scripture a mighty army of angels surrounds us at God's beckoning call..
I invite you to gain a new perspective on life today and find God's peace for your weary souls..
So let me just ask you a couple of questions here as we wrap up..
Are you ready to let go and let God fight your battles for you?.
Are you willing to live life with victory and hope in his kingdom?.
Would you raise your white flag now and surrender your life to him,.
confessing your need for Jesus to reign in your life?.
These, I believe, help us to find hope in these days..
Please look up instead of just looking with what's around you..
God has so much more for you if you'll only look up..
Let me pray with you..
Father in heaven, thank you so much for these that are gathered here today and these that are watching..
I pray, Lord God, that your words would speak with power here..
Thank you, God, for the dozens of intercessors that I know are praying right now for what I'm sharing here today..
I pray, God, that you would help those who hear it to see what you are saying to them.
and to take action to surrender and to look to you.
and not to fret about what they see happening with the nations around them..
Help them to see what's above and to see what's beyond this world.
and to trust in it, to trust in its ruler, the true ruler that God has put over all things, King Jesus..
And we pray this now in the name of Jesus. Amen..
[Music].
\newpage



\section{}
\label{sec:KF6pQDFyEx4}
\textbf{聖地網上行 - 步主腳蹤 [第六集] 加利利湖}
\newline
\newline
連結: \href{https://youtube.com/watch?v=KF6pQDFyEx4}{\texttt{https://youtube.com/watch?v=KF6pQDFyEx4}} ~~~~ 語音日期: 2024-04-20
\newline
\newline
\hyperref[sec:GXyrjafkb9Y]{\small{< < < PREV SERMON < < <}}
~
\hyperref[sec:index_chronic]{\small{[返順時目]}}
~
\hyperref[sec:index_scriptual]{\small{[返順卷目]}}
~
\hyperref[sec:cEtF7vpVnrE]{\small{> > > NEXT SERMON > > >}}
\newline
\newline
$^{1}$各位好,我們來到聖地網上行的第六集.
我們一起來跟隨耶穌基督的腳蹤.
去到他去過的地方.
在地圖上看到這個地方,你記得叫什麼名字嗎?.
沒錯,就是加里里海或者加里里湖.
在舊約聖經裡面通常叫這個地方做基里列湖.
或者叫做Sea of Kinnerick.
Kinnerick這個字是希伯來文的Kinnor.
意思其實就是樹禽,就是鹹的意思.
不知道你覺得這個湖的形狀是真的像樹禽嗎?.
在新約福音書裡面通常叫這個地方做加里里海.
Sea of Galilee或者加里里湖.
或者有時候會叫做Sea of Tiberias.
因為福音書寫的時候,這個湖旁邊最重要的城市.
就是加里里省裡面的省府.
所以就叫做提比里亞海.
加里里這個字其實就是圈的意思.
Circle或者是District,一個地區的意思.
在地質上是非常複雜.
但最重要的特色當然就是加里里海或者加里里湖.
過去幾年最重要的新聞就是.
加里里湖因為雨量不足.
因為好幾年下雨下不夠.
所以水量是來到百年來最低最低.
低到一個地步,湖上甚至出現了一個小島.
水一直乾的時候,有一個小島就出現了.
因為寒季的緣故,水減退.
有一個很重要的發現.
就是1986年年初發生的.
當時因為有好幾年的寒季,湖水減退.
所以有一些以前在水底下的沙灘岸都呈現了出來.
有兩個漁夫在莫大拉附近.
他們就去附近的沙灘上無意之中.
發現了地上有一些鐵釘.
然後越來越多鐵釘的時候就看到有木頭.
一直到用手刮開的時候.
原來是一艘古木船.
他們就通知一些考古人員來,就轟動一時.
甚至有人傳聞說船上有黃金.
結果警察和以色列軍隊都來到這裡把守.

$^{41}$因為那些古木一暴露在空氣裡面.
會容易腐爛.
所以工作人員花了12個晝夜.
搶救把這艘古木船發掘出來.
然後用防水泡沫和玻璃纖維包著.
然後把它運上去,浮到別的地方.
運上岸去作化驗.
化驗的結果就說這是屬於第一世紀的木船.
所以這艘木船叫什麼名字呢?.
沒錯,就是耶穌船.
當然這些純粹是市場推廣.
但確實是當時耶穌時期的一艘古木船.
這艘古木船的圍骸長大約8.2米.
即是27呎左右.
闊度是2米多,即是大約7呎半.
深度大約1.3米,或是4呎左右.
這艘船大概能坐十多人.
這很吻合聖經的記載.
也很吻合約瑟夫一世紀的猶太歷史學家所描述的木船.
這艘木船是由很多不同的木造成的.
或是修補而成.
有12種不同的木.
你看到這艘小木船.
其實很能幫助我們了解聖經的一些細節.
例如約翰福音21章第6節.
耶穌說:你把網紮到船的右邊,就必得著.
這裡是復活的主在岸上.
叫湖上的門徒.
他們正在打魚,打不到什麼.
他告訴他們,你把網紮到另一邊,就必得著.
其實是一個很小的動作,一個很小的改變.
你想想,一艘木船從左邊紮到右邊.
但結果是很不同的.
雖然只是很小的改變.
其實就像我們聽上帝的話一樣.
主對我們說的吩咐.
其實要我們的改變也是很小的.
我們把網紮到另一邊,然後拉上來的時候.
就發現有很多魚.
我們一個小的動作,小的改變.

$^{81}$一個單單的馴服.
其實就可以看到一個完全不同的結果.
其實加里尼湖確實是耶穌對門徒.
對我們的教室.
剛才我們說過,去年是遭遇百年來最嚴重的寒季.
水位是最低的,甚至是島也出來了.
但是上年,2019年,雨季開始的時候.
就大雨磅bok.
湖水突然間大增,突然間就爆滿.
水位又回到滿的程度.
你說加里尼湖,確實是不能夠預測.
真的,這就是耶穌對我們人生的教室.
確實人生就是這樣.
我們不能夠預測.
我們唯有靠上帝的恩典.
加里尼海平時風平浪靜,其實是很美的.
但是也可以突然間翻起暴風.
所以是非常危險的.
這是地球上海拔以下最低的淡水湖.
陷在峽谷深谷之下.
在Google Earth上可以看到.
好像是一個深的河底.
四周的高山吹來的自然就是冷風.
白天曬過的湖面,空氣當然是熱.
一冷一熱混在一起的時候.
加上四周有很多的深谷.
形成了天然的大漏斗.
強風就可以捲入湖的上面.
突然間湖面就可以狂風大作.
事至今日有不少的新聞.
都是和湖上的泳客的遇溺有關.
說一宗新聞,是在2017年4月的時候發生的.
當時有幾個人在岸邊.
在湖的東邊的一個沙灘露營.
白天的時候他們就在湖上游泳.
誰知突然間有大風從東邊吹來.
他們的浮床就被吹了出去.
當浪開始時,三個人都遇溺.
然後就喪生了,非常可惜.
三個死者的年齡是17,19和21歲.

$^{121}$所以這個湖看起來不是很大.
深度也不是很深,大概200呎.
但是當有突然的風浪來的時候.
很多時候甚至年輕的人,熟水的人都招架不住.
今天我們看兩段經文.
兩段經文都是和門徒在加里里湖上遇見風浪有關.
首先看馬可福音第四章,第一次的風暴.
當天傍晚,耶穌對門徒說:我們到對岸去吧.
於是門徒離開群眾,帶著已經在船上的耶穌一同過去.
也有其他船和他同行.
大風暴來了,波浪打進船內,船很快就灌滿了水.
耶穌在船尾枕頭睡覺.
門徒喚醒他對他說:老師,我們快沒命了,你不管嗎?.
耶穌醒了,斥責風又對海說:處了罷,靜了罷.
風就止住,大大平靜了.
耶穌對門徒說:為什麼害怕?還是沒有信心嗎?.
他們大衛驚恐彼此說:這人到底是誰?.
連風和海都聽從他了.
這裡說到對岸,從猶太人門徒的角度去看.
對岸是什麼地方?對岸是愛邦人的地方,即是羅馬人,羅馬帝國的世界.
下文就知道,這個地方叫做格拉森人的地方.
那裡有群鬼,代表羅馬帝國的勢力.
這是耶穌在加里尼傳道的第一次進入羅馬愛邦的世界.
經文裡面可能有些細節我們要解釋一下.
剛才我們不是說找到一艘古船嗎?.
這樣比較容易想像當時發生了什麼事.
他說船尾枕著,枕頭睡覺.
這句就很奇怪,到底想說什麼呢?.
其實這些是船的構造.
當時的艇仔,船尾會有一些所謂的壓艙物.
壓著船,穩定船身.
今日的艇仔,漁船都會有這樣的設計.
有些沙包放在船尾的艙庫.
很多時候漁夫累的時候會進入這個小小的艙去睡覺.
所以這裡說到耶穌在船尾枕著,枕頭睡覺.
他可能在休息.
當船在航行的時候,突然間就有大風暴.
在以色列巴勒斯坦大致上有兩個季節.
一個叫雨季,一個叫寒季.
寒季就是夏天的時候,大概是五旬節之後.

$^{161}$直到大概九月十月的時間.
這幾個月裡面很少會下雨.
農作物是靠清晨的甘露來生存.
果樹,橄欖樹等等,靠著清晨的甘露.
即是日出之前地上出現的露水.
聖經裡面有不同的經文都有記載這樣的情況.
我記得曾經去過一些考古工場工作.
天還未亮,我們就要去到那個地方.
去到的時候發現工場上的帳篷很重.
因為露水弄濕了.
但是太陽一出來的時候,露水就會突然間消失.
第二種季節就是雨季.
大概十月到下一年的四月尾.
聖經也有講這些季節.
叫做初雨,或者後雨.
雨雲會由地中海,由西至東吹來.
帶來暴雨.
耶穌的門徒裡面有幾個都是當地很有經驗的漁夫.
既然作為漁夫他們有這麼豐富的經驗.
他們一觀察就知道雨雲,風雨何時會來.
絕對對他們來說是不會突然.
但是在加里里湖上,暴風就有些不同.
不是平時從西吹來的雨雲,暴雨.
而是暴風,不是雷雨.
這些暴風是由東面在高欄高地.
二千多呎的高原,突然吹到湖面的東風.
是會突然來的,是沒有預測到的.
這些暴風一吹來就會很猛烈.
有些量度甚至可以風速高達50公里.
即是30英里.
你就可以明白那些經驗老道的漁夫.
雖然見到這麼多風雨,懂得看天色做人.
突然有些風暴,他們都招架不來.
這就是我們人生的縮影.
風浪一定有的.
考試失敗,感情挫折,事業倒退.
遇見不同的逆境,生老病死等等.
這些很多時候都是我們預測得到的.
甚至我們覺得自己很有經驗.
因為人生不如意事十常八九.

$^{201}$但這裡告訴我們,你怎樣經驗老道.
你怎樣致命非凡.
我們人生命的本質都一樣是這麼脆弱.
一下子會有一些風暴來,我們就會被打倒.
因為這些是我們不能夠預測的.
我認識了不少人,人生打遍天下無敵手.
很成功的人,但逆境一來到.
或者世界政治的動盪,突然間就會失控.
大家都聽過,在逆境當中.
甚至他有病死都沒有人可以理會.
我亦認識一些人,他們雖然有突然的暴風來.
他們會難過,會擔心.
但他們心裡仍然有盼望,仍然很渾土.
他們就是一些認識救主的人.
其實馬可福音上下文就是介紹救主神的意志.
是我們每一個人都需要的.
大家留意馬可福音的脈絡.
第四章三十五節,這個風暴的故事.
就是耶穌勝過風浪,勝過海.
然後到第五章一節開始.
就是到格拉森林的地方.
就是耶穌勝過魔鬼.
五章二十一節,他回到湖的另一邊,猶太人那邊.
就是艾魯的女兒死而復活的故事.
這位救主勝過人生任何的風暴.
甚至死亡都能夠達成.
到底你是否真正認識這位救主.
並且真正全心全意地信教他.
剛才你看到的影片都是一些實地的影片.
是加利湖翻起風浪的時候.
有風暴來到的時候的情況.
是否世界上最大的風浪?當然不是.
但我們人很脆弱.
你看到那些古船,小艇其實都很脆弱.
當你在實地裡面有風浪來的時候.
你都會招架不住.
唯有這個主和我們一路同行.
我們就明白什麼叫做在風浪當中仍然有平安.
第二場風暴我們看馬太福音第十四章.
耶穌隨即催門徒上船.

$^{241}$先到那邊去.
等他叫眾人散開.
散了眾人之後.
耶穌就獨自上山去禱告.
到了晚上只有他一人在那裡.
那時船在海中.
因風不順被浪搖晃.
夜裡四更天.
耶穌在海面上走.
往門徒那裡去.
門徒看見他在海面上走.
就驚慌了船.
是個鬼怪.
便害怕喊叫起來.
耶穌連忙對他們說.
你們放心是我.
不要怕.
彼得說主如果是你.
請叫我從水面上走到你那裡去.
耶穌說你來吧.
彼得就從船上下去.
在水面上走.
要到耶穌那裡去.
只因見風大就害怕.
將要沉下去.
便喊了說主啊救我.
耶穌趕緊伸手拉住他船.
你這小信的人.
為什麼疑惑.
他們上了船.
風就住了.
在船上的人都拜他船.
這真是神的意旨了.
這是科幻書記錄的第二次船上風暴.
晚上發生的.
上下文告訴我們.
大概是耶穌餵飽了五千人.
五餅而如的神蹟.
然後叫門徒上船.
夜間就發生風暴.

$^{281}$既然我們來到聖地.
我們也要看看五餅而如堂.
這是初期教會.
四世紀開始.
有這個地方紀念五餅而如餵飽了幾千人的神蹟.
這個教堂叫五餅而如堂.
上面有一個馬賽克圖.
上面有五餅而如.
可能設計的人沒有看聖經.
只有四個餅.
但這裡有一個很出名的.
紀念耶穌五餅而如的地方.
經文記載.
門徒上船時有風不順.
against them.
當你理解本地氣候.
你就知道大概是東面高山高欄高地吹下來的猛風.
這裡提到夜裡四更天.
當時的人將日出至日落分為四更.
四個watches.
凌晨三點至六點.
這是第四更.
三小時一更.
第四更是最後一更.
最後一段時間.
也是最黑暗的時候.
所以門徒逆風而行.
在船上掙扎了很久.
耶穌在水上行走.
一定是令人想起舊約裡面不同的內容.
例如詩篇第77篇19節.
「你的道上帝行的路在海中.
你的路在大水中.
你的腳中無人知道.
你曾藉著摩西和阿倫的手引導你的百姓好像羊群一樣」.
詩篇第77篇提到神可以在水上行走.
可以戰勝水.
留意整篇詩篇第77篇.
詩人在遇難的時候唱出這篇詩詞.
並且他在夜更裡面.

$^{321}$心裡非常不安.
他覺得上帝忘記了他.
在這時候.
詩篇裡面的逆轉是怎樣呢?.
就是他紀念上帝曾經在上古之年做過的事.
紀念到歷史上上帝曾經救贖過自己的子民.
救他們出埃及.
然後走過大水.
你記得過紅海的故事.
所以福音書記載到水上風暴.
甚至這裡提到耶穌在水上行走.
一定會想起出埃及的情景.
其實上下文之前也有提示.
你記得這個風暴之前發生什麼事.
就是餵飽五千人.
在曠野裡面五餅而魚餵飽這麼多人.
就好像出埃及之後.
在曠野裡面神藉著嗎哪餵飽他的子民.
還有27節.
你們放心,是我,不要怕.
門徒以為看到鬼,他就在水上行走.
他說你們放心,是我.
是我這個字就是I am.
在猶太人的典籍裡面.
第一次出現是什麼呢.
就是出埃及的時候.
就是神啟示他的百姓.
我的名字叫做耶和華.
耶和華這個字就是I am的意思.
就是跟摩西說這個就是我的名字.
這個就是你會經歷的新的作為.
另外還有什麼細節呢.
想起夜裡四更天.
你就會想起出埃及.
我們看出埃及記第十四章.
這裡就是百姓出埃及的時候.
前無去路後有追兵.
埃及的軍隊在追殺他們.
來到紅海邊.
摩西向海伸張.

$^{361}$耶和華便用大的東風.
使海水一夜退去.
水便分開.
海就成為乾地.
以色列人可以下海走乾地.
水在他們左右成為長潭.
記得大東風.
水退去.
然後他們行在乾地上.
埃及人追趕他們.
發露一切的馬匹車輛馬兵.
都跟著下到海中.
到了神耕的時候.
耶和華從雲火柱中.
使埃及的軍兵混亂了.
又使他們的車輪脫落.
難以行走.
以至埃及人說.
我們從以色列人面前逃跑吧.
因為耶和華為他們攻擊我們了.
其實這個神耕就是黑夜.
最後一耕.
以色列人去到紅海邊.
他們眼前好像看到埃及的暴君.
軍隊得勝.
但是在這個神耕的時候.
黑夜中居然出人意料.
上帝仍然掌權.
上帝仍然可以得勝.
然後他在水中開一條路.
給他的子民行走.
然後擊退了埃及的軍隊.
最終目的是什麼呢?.
埃及人都被折服.
去認耶和華是真實的神.
為什麼初期教會.
其中一個代表教會的記號就是船.
雖然外面有很多風暴.
但是有上帝同在.
他們就可以得勝.

$^{401}$所以後期教會的建築.
都是想表達這個意思.
門徒聚集的地方.
一起去敬拜神的地方.
叫Navis.
拉丁文的Navis.
意思是船.
如果有上帝和我們同在.
如果我們站在上帝那一方.
我們就能夠得勝風暴.
當時的教會都是非常弱勢.
在水深火熱之中.
在一個澳漣的中間.
耶穌水上行.
不是去彰顯一下.
我多厲害.
而是去鼓勵我們的生命當中.
教會去到水深火熱.
澳漣的當中.
我們要想起.
原來耶穌可以在水上行.
其實馬太福音上下文.
都已經給了一些提示.
上下文說到.
施洗約翰被暴君希律王處斬.
在當時教會.
當時的信徒門徒最黑暗的時候.
到底上帝是否忘記了教會.
忘記了自己的子民.
你看這個野間湖上風暴的故事.
只有馬太福音才加一個細節.
就是彼得也在水上行.
彼得代表什麼呢?.
彼得代表教會.
因為他是門徒之首.
他行在水上.
為什麼可以這樣做?.
因為他的眼目注視在耶穌身上.
在風浪當中他看到的.
是耶穌真正的神的意志.

$^{441}$但當他懷疑的時候.
他就下沉.
因為他的眼目看見四周的風浪.
懷疑的字有三心兩意的意思.
即是事奉兩個主的意思.
當你這樣做的時候.
你就會害怕風浪.
就會下沉.
有黑夜有風暴.
我們很想這個黑夜風暴快點過去.
這是人之常情.
但這段聖經告訴我們.
到底我們在逆境.
在風暴當中.
有沒有令我們自己.
能夠更加認識這位神的兒子耶穌基督.
以至甚至外邦人.
埃及人都可以在我們身上.
看見神的偉大.
上船之後風平浪靜.
船上的門徒說什麼呢?.
他說這是真是神的兒子.
這裡有一個特別的細節我們要留意.
就是神的兒子這個稱號.
你說有什麼特別呢?.
文法上希臘文的次序.
通常是寫兒子然後寫神.
即是神然後到神.
然後到天.
但馬太福音特別將次序倒轉.
先寫神然後才寫兒子.
然後才寫神然後才寫神.
神神.
在馬太福音這個不是偶然的.
三次都用了.
三次都是用這樣的次序.
第一次用「加辛」.
神的兒子.
神行先.
這是第一次.

$^{481}$之後的兩次都是在十字架上寫.
最後一次.
27章第54節.
羅馬的伯夫長和馬太.
是看守耶穌的人.
盯著他的十字架的士兵.
看見地震所經歷的事.
就極其害怕.
說這真是神的兒子了.
你留意和加利利湖上.
黑夜風暴.
門徒說的話.
用詞都是一樣.
這真是神的兒子.
但這個出口的居然是.
羅馬的伯夫長.
羅馬的士兵.
你說有什麼分別呢.
翻譯上沒有什麼大分別.
中英文都是Son of God.
或者God's Son.
特別在於哪裡呢.
神行先然後到兒子.
這個清咸的綽號.
是一世紀羅馬帝王.
他們的冥文inscription.
硬幣.
全部通用的稱號.
事實上都是神行先.
然後到兒子.
不是相反.
由阿古士督到一世紀末的.
羅馬帝王竇米田Domitian.
所以神的兒子.
神行先然後到兒子.
這個行口.
是一個imperial title.
是一個羅馬帝國尊用的清咸.
馬太拿了來用.
他說真正神的兒子.

$^{521}$真正擁有大權的.
原來不是羅馬帝王.
真正給我們生命有意義的.
掌握生死的.
不是羅馬帝王.
而是這位在水上行的耶穌基督.
到底我們生命裡真正的君主.
領導人是誰.
或者這麼問.
風暴過後.
你是否一個更加認識.
神的兒子耶穌基督的人.
還是我們只不過想.
逆境的風暴過後.
其實我們打回原形.
我們依然固朴.
風暴過後.
其他人是否更加在我們身上.
在我們教會.
能夠認識神子耶穌基督.
如果我們希望是.
人更加認識神的兒子耶穌基督.
我們在風暴當中.
我們就需要問.
到底上帝要我們學習的是什麼.
我們是否更加用心地跟隨著.
我們一起祈禱結束.
神啊!你才是掌管萬有.
控制所有歷史.
我們生命裡每一刻的神.
縱使在風暴逆境當中.
你依然是掌權.
你依然是那位真正的君主.
真正的帝王.
求你自己繼續幫助我們.
縱使我們現在在困難.
在逆境當中.
求你讓我們更加盡心地跟隨你.
主,我們希望.
不只是風暴過後.

$^{561}$我們依然顧我.
打回原形.
而在風暴的當年.
我們更加成長.
更加認識你.
風暴過後.
讓人更加認識你自己的美名.
讓人更加認識主耶穌基督.
我們神彼而至.
我們這樣祈禱,放出力量.
我們是否真的能夠成為神?.
\newpage



\section{}
\label{sec:cEtF7vpVnrE}
\textbf{Everlasting Praise 2 (Genesis 2) - Rev. Mark Chan}
\newline
\newline
連結: \href{https://youtube.com/watch?v=cEtF7vpVnrE}{\texttt{https://youtube.com/watch?v=cEtF7vpVnrE}} ~~~~ 語音日期: 2024-04-22
\newline
\newline
\hyperref[sec:KF6pQDFyEx4]{\small{< < < PREV SERMON < < <}}
~
\hyperref[sec:index_chronic]{\small{[返順時目]}}
~
\hyperref[sec:index_scriptual]{\small{[返順卷目]}}
~
\hyperref[sec:H_MgtrWeR24]{\small{> > > NEXT SERMON > > >}}
\newline
\newline
$^{1}$Hey, we are going to continue our mini-series on worship..
And last time we were talking about Psalm 8, right?.
Yeah, and what a nice Psalm to reflect upon..
And today we're going to continue this series.
and we are at the beginning of the Bible..
So if you have your Bible,.
turn to chapter two of the first book, Genesis..
Chapter two, Genesis, all right?.
We're going to read a part of scripture from chapter two..
And this is a passage that talks about Garden of Eden..
Garden of Eden, all right?.
So let me read a couple of verses out of chapter two..
Chapter two, verse seven..
Then the Lord God formed the man from the dust of the ground..
He breathed the breath of life into the man's nostrils.
and the man became a living person..
Then the Lord God planted a garden in Eden in the east..
And there he placed the man he had made..
The Lord God made all sorts of trees.
growing up from the ground,.
trees that are beautiful, that produce delicious fruit..
In the middle of the garden, he placed the tree of life.
and the tree of knowledge of good and evil..
And verse 15, the Lord God placed the man in the garden.
of Eden to tend and watch over it..
But the Lord God warned him,.
"You may freely eat the fruit of every tree in the garden,.
"except the tree of knowledge of good and evil..
"If you eat its fruit, you will surely die.".
Then the Lord God said,.
"It is not good for the man to be alone..
"I will make a helper who is just right for him.".
This is a simple passage, right?.
Seemingly very simple..
By the way, I'm reading from a New Living Translation, NLT,.
which is easier to understand for many of you..
And as we open the opening chapters of the Bible,.
we find a simple story, right?.
How many of you have seen or read a children's Bible?.
Children Bible, kid's Bible..

$^{41}$Yeah, right?.
Do you find this story in the children Bible?.
You do, right?.
Yeah, it's a must-read story, right?.
And you find a forest, Adam and Eve,.
two people in the garden..
And it seems like a very, very simple kid's story..
However, a closer reading,.
or if you read it more carefully,.
you raise a lot of difficult questions,.
a lot of different issues..
For example, did Adam have belly button?.
He's the first guy, right?.
And he was not born from a mother..
Who's the mother of Adam?.
No mother, right?.
So how about Eve, right?.
Eve was born or came into existence.
in a very different way, right?.
So do they have belly buttons?.
Or when did they come about, right?.
I mean, a simple, literal reading of scripture.
will have us believe, okay,.
maybe Genesis 1 and 2,.
things that happened just a few thousand years ago, right?.
So we've come to the conclusion, probably,.
that the earth is young, right?.
It's really, really young..
I mean, how many pages you got, right?.
From creation to Abraham, right?.
It's only 12 chapters, right?.
And what is Adam's DNA?.
Have you ever thought about it?.
Well, Adam is a real man, right?.
I mean, what's his genes, right?.
And Eve, and is the Garden of Eden on the Google map?.
Can you find it?.
Can you find it somewhere in the Middle East, maybe?.
In Iraq?.
Or let the archaeologists take it up, right?.

$^{81}$So a simple reading will have us believe.
this is a children's story..
But a careful reading, it weighs a lot of questions..
And I think a lot of questions that we ask,.
or the questions that I ask,.
you can't find answer from Genesis 1, 2, and 3..
Because the scripture, Genesis 1, 2, and 3,.
are not designed to answer some of our,.
or many of our modern questions, right?.
It is designed to answer a couple of important questions..
Very important, more important.
than the questions that we ask, all right?.
So reading the Bible is kind of like this, okay?.
We have to come to a certain strategy,.
understand what the Bible is all about, right?.
Otherwise, we'll ask a lot of different questions.
that the Bible is not going to answer, right?.
So we have to listen to the text..
We have to listen to the Bible,.
listen to the word of God carefully,.
and let it dictate what it's trying to say..
And we are here to listen and to obey..
That is our attitude in front of scripture, right?.
Before the Bible, we are servants,.
and we listen to the word of God..
So how do we read Genesis 1, 2, and 3?.
Simple story, but actually a very deep story..
I think the story in Genesis 1, 2, and 3,.
and we'll spend some time in Genesis 2,.
are much like a symbolic story, you know, symbolic story..
It's like a can full of different symbols,.
and these symbols have deeper meanings,.
and we have to ask the questions,.
what do they point to, right?.
My analogy of reading the Bible is like reading a map, okay?.
Do you have Google Map on your phone?.
Or some sort of GPS app on your phone, right?.
And you key in your destination, and it will show you,.
you know, where to turn, right?.
Which road to take, which highway to take,.

$^{121}$which exit to take, Google Map, right?.
We all read maps, and sometimes when we read maps,.
we kind of intuitively, we understand what the map means,.
right?.
Reading the Bible is kind of like reading a map..
Let me give you an example..
I have three examples here..
On the right is a Chicago map..
Upper left is a map of Hong Kong, right, where I came from..
And then lower left is, do you know where it is?.
It's Tokyo, Japan, right?.
Why do I put a Japan map on the screen?.
Because a lot of people telling me,.
oh, I need to go to Japan, you know?.
(audience laughs).
Oh, the currency exchange rate is so good, right?.
Everybody's going to Japan, okay?.
So there you go, all right?.
So you read a map, we kind of, we read it intuitively..
We understand what it means, right?.
For example, upper left, the map in Hong Kong,.
I add the layer of traffic..
So you see all the highways, right?.
You see different colors..
You know, red means it's very congested..
Black is totally stopped, right?.
Yellow is kind of congested,.
and then green is all smooth, right?.
Those are traffic map, right?.
So you add a layer, you understand the symbols,.
you understand the color of the lines,.
and then you interpret the map.
the way it is intended to, right?.
You got it?.
And then at the lower left is the Tokyo map..
It's a terrain map, terrain..
Where the mountain, where the valleys, right?.
You see Fuji mountain on the map..
On the, it's at the lower left corner, right?.
Mount Fuji..

$^{161}$So it shows you where the hills are,.
where the mountains are, where the valleys are, right?.
And they color it, right?.
For the plane, for the level, ground level,.
or elevation zero, right, sea level, you see white..
And then you climb the mountain, you see green,.
and then there's a black spot, right?.
Those are higher elevation..
That's a terrain map, right?.
And then on the right is, of course, the Chicago map..
Where are we?.
Where are we on the map?.
We're kind of in the near south..
You see the pink color?.
You know where you are?.
Come on, guys, right?.
You know where you are?.
All right, you're not in green, okay?.
You're right by the lake, right?.
You see Navy Pier?.
You see that?.
Oh, come on, right?.
This is elementary questions, right?.
Yeah, we're kind of in the near south,.
the lower end of the pink region, right?.
So the map has you, different color, right?.
So the map is coded in such a way that you know,.
okay, this is the near south, this is the north,.
and this is the region, or how we call it, right?.
You don't take the map, this color map of Chicago,.
and you go out to the street and say,.
"How come it's not pink?".
Pink!.
I'm expecting the road to be all pink, right?.
Because we're in the near south, right?.
We don't do that, right?.
Because we intuitively, we know that what it means is that,.
okay, the pink area is simple meaning,.
well, this neighborhood belong to the same region, right?.
To the same area..

$^{201}$Now, we don't go out and use a Google map and say,.
"Oh, the line is blue, and I don't see any blue line," right?.
It shows you with legends, with different color coding,.
or different symbols to tell you some aspect.
of the reality, right?.
Now, no map can tell you everything, right?.
You cannot go back to the map maker and say,.
"Hey, why didn't you answer my questions?.
"Why didn't you tell me the terrain?".
Right?.
Well, it's not terrain map, right?.
It's a color code map..
It tells you all the 77 neighborhoods of Chicago,.
and what do they call, right?.
And that is the purpose of a map..
Now, you cannot fault the maker by saying,.
"You don't tell me the reality comprehensively..
"You don't answer all my questions.".
Now, it's much the same way when we read the Bible story..
No Bible story will tell you everything..
No text in the Bible, no passage in the Bible.
will answer every single question that you have, right?.
Now, your questions or your concern.
when you read the Word of God is let that passage.
dictate the questions to ask..
And listen to the passage, listen to the words..
What are they trying to tell you?.
Now, the Bible does not answer all the questions.
that we have, and you cannot say,.
"Okay, the Bible is flawed.".
That's not what it means, because it tells you.
some aspects of the reality..
So reading Genesis 2 is kind of like reading a map..
We have to understand what kind of map it is,.
what kind of text it is..
I already told you the Garden of Eden story.
is highly symbolic story, and use different symbols,.
like a map, using different color codes.
to tell you some truths of the reality, all right?.
What I'm trying to suggest to you.

$^{241}$is that Genesis 1, 2, and 3, the opening chapters of Genesis.
or even the first 11 chapters, or many other passages,.
especially Genesis opening chapters,.
they are identity document..
The text is trying to answer the core question.
of identity, human identity..
In other words, who am I, right?.
Who am I?.
Okay, I'm Mark Chan, but who am I?.
Why do I even exist?.
What am I here for, right?.
What is the purpose of me being here?.
My existence, my purpose, you know?.
Now, so often we can make up answers.
to all of these identity questions,.
but the truth is we have to ask the creator God,.
he created you, and we owe our lives to him..
We have to ask him how to answer those questions..
So there we go, we come to Genesis 2..
It is a story on Garden of Eden..
It's a highly symbolic story..
It doesn't, when I say the story,.
the Bible story is highly symbolic,.
it does not make it less historical..
It does not make it less real, okay?.
Just like the color code map,.
it does not make Chicago less real or less historical..
It is telling you something about the reality of Chicago..
It's a color code map..
It's the same way, Genesis 2, the Garden of Eden story,.
it's a real story..
It's written in such a way that it's symbolic,.
right, understanding a map, different legends and colors..
So we read the story, we have to try to think,.
you know, what the biblical author,.
or what God is trying to say through this passage..
Well, let me test you, all right?.
You like a pop quiz, right?.
Think of Garden of Eden..
Now, the Garden of Eden has an entrance, all right?.

$^{281}$Just like church has an entrance, right?.
Where, which direction does the entrance.
of Garden of Eden face?.
East, okay, very good, Jin Fang, very good..
Good student..
No question about it, right?.
East, right?.
Okay, second question..
Who is guarding the entrance to Eden?.
Who is at the door?.
So you walk in our church, right?.
There is some people guarding the door, right?.
Some of the core team members,.
or some of the welcome team people..
But who's guarding the entrance of the Garden of Eden?.
Who?.
What, I'm sorry?.
Give it a shot..
Okay, anyone give it a try?.
It's cherubim..
Who's that guy?.
Cherubim..
Cherubim is one of the angels, all right?.
Cherubim, all right, got it?.
Cherubim, say it, cherubim, all right?.
There are at least two cherubims.
guarding the entrance of Garden of Eden..
Now, you say, why is it important?.
All right, okay, the entrance of Garden of Eden.
is facing east, okay?.
The place of sunrise, all right?.
Now, in the Bible world,.
all the temples, all the sacred places,.
their entrances face east..
They all face east..
All the sanctuaries in the Bible world face east,.
the place of sunrise..
And at the entrance, facing east,.
is often guarded by some sort of angelic beings,.
some sort of creatures..

$^{321}$Those creatures are shaped in such a way.
that gives you power,.
because who's guarding the door, right?.
It cannot be a little kid, right?.
It has to be someone who is strong and who is powerful, right?.
And cherubim is kind of like that, right?.
Now, let me tell you, much later in the Bible,.
a tabernacle, a tabernacle is actually a portable tent, right?.
Remember the Israelites came out of Egypt, right?.
They were in the wilderness, in the desert,.
and God is living among them in the form of tabernacle..
So he designed the tabernacle,.
and Moses and all those people, they built the tent, right?.
And you can pack the tent, and then you go somewhere else,.
and then you pitch a tent, and then God is there..
And then when you are about to go to another station,.
and then you pack the tent, and then you pitch the tent..
That is tabernacle..
Tabernacle, the entrance of tabernacle..
What direction does it face?.
What direction?.
The entrance of tabernacle..
You know that?.
East, okay?.
East, yeah, I told you already, right?.
East, right?.
And then after the wilderness, many years,.
they came, settled down, the Israelites came,.
settled down in the land of Canaan,.
and eventually they made Jerusalem the capital..
And remember who built the first temple?.
Solomon, right?.
Solomon, come on, guys..
Solomon, wake up..
Almost noon..
Solomon, right?.
The son of David, right?.
He built the first temple, Solomon's temple, in Jerusalem..
And the entrance of Jerusalem temple,.
which direction does it face?.

$^{361}$East, yeah, yeah, exactly..
Exactly, it face east, right?.
And then the entrances of tabernacle and temple.
are guarded by, who's guarding the door?.
What's his name?.
Cherubim, right?.
Cherubim, okay?.
So you find decoration of cherubim in the tabernacle,.
especially the entrance to the holy place..
And then in the temple of Solomon,.
is also guarded by two cherubim, 1 Kings chapter six..
There are two cherubims, angelic beings,.
guarding the door, right?.
Facing east..
Now what did I say in the beginning?.
The God of Eden, facing east..
Yes, yes, there you go..
And who's guarding the door?.
Cherubim, okay, cherubim..
And then you find all sorts of different things.
in chapter two of Genesis..
You got flowers, you got trees,.
you got all this beautiful nature,.
and gold and silver and all this precious stones..
Where did you find it other than God of Eden?.
You find decoration of all those things in tabernacle,.
and in much later, the temple of Solomon, right?.
So in that way, you understand.
the God of Eden is actually a sanctuary..
It is at least described in such a way,.
using later terms to describe the prototype.
of sanctuary, a holy place..
Now what is a sanctuary for?.
For worship, yes, for worship, right?.
We're talking about worship in this series, right?.
So right at the beginning, God is telling us.
when God set up the universe,.
when God put man, humanity, in the God of Eden,.
his purpose is for men and women, humanity, to worship him..
You got it, right?.

$^{401}$So don't even come to the God of Eden story.
and say, "Oh, this is like my backyard, you know?.
"I can do barbecue here, right?.
"I can grow some flowers, you know, I can enjoy it.".
No, it's a sacred place..
It's like a temple..
It's like a sanctuary, right?.
Facing east and guarded by cherub, who?.
Cherubim, okay?.
Cherubim, it's a sacred place..
It's where God dwells..
God lives there, right?.
Now you have heard me before, maybe some time before,.
when you read Genesis 1, you know, six day of creation.
or even seven day creation..
What is that about?.
It's not answering some of the scientific questions.
that we have today..
It's trying to tell you why it exists.
and what is the purpose of the entire universe,.
entire cosmos that God set up..
Seven days, right?.
What other buildings would take seven period of time to build?.
Temples, temples..
In the Bible world or even within the Bible,.
when Solomon built the temple,.
it spent seven period of time, right?.
And building the tabernacle..
Seven paragraphs to talk about building the tabernacle..
So Genesis 1, the first story,.
it's not about geology or science or anything that we ask..
It's about why this universe exists..
Seven day, God set up the cosmos..
So the entire cosmos is like a temple..
It's like God dwelling in the universe..
The entire universe, the entire cosmos,.
it's a sacred place..
It's like a holy place..
God meant to dwell in the midst of it, right?.
Talking about creation care, right?.

$^{441}$Talking about environmental care..
You know, Genesis 1 already right off the bat telling you,.
you know, this world is the cosmos set up by God himself.
to create a design in such a way that is sacred, right?.
That is holy, that he can dwell in it..
And then chapter two, the camera zoom in.
into one spot in the entire cosmos called Garden of Eden..
It's the first sanctuary..
It's the prototype of sanctuary..
And did I say, did I tell you.
the document is the identity document?.
Right?.
So the passage is trying to answer.
questions about your identity..
Who am I, right?.
Now, when you get to know someone, a new friend,.
how do you introduce yourself?.
You tell them kind of who you are, right?.
Where you're from, what do you do, right?.
And those answers are identity..
Say for example, you're a student and you say,.
"Oh, hello, I'm a student.".
And maybe your career, right?.
I will tell people, you know, I'm a pastor..
And they go like, "Whoa.".
Right?.
You may be an accountant..
You may be a teacher..
You may be just a housewife or chef in a restaurant..
Or we answer the identity questions.
by telling people what we do..
What we do..
But here's the problem..
If we fix our identity, if we define our identity,.
anchor our identity simply by what we do,.
it's gonna change, right?.
It's gonna change..
Before I was a pastor, I was a lab chemist..
You know, I work in a environmental lab.
and I do testing for water and all those stuff.

$^{481}$because I studied biochemistry in my undergrad, right?.
So, you know, thank God I graduated and I found a job..
So I was working in a lab..
So every single day, eight hours, nine hours a day,.
I'm doing testing all the time, right?.
And I was called to ministry and then I went to seminary.
and I became a minister, right?.
See, my identity changed from a chemist to a pastor, right?.
And then someday I will retire..
After I retire, am I still a pastor?.
Well, yes and no, right?.
Yes and no, I'm not pay staff of a certain church..
But, you know, maybe part of me is saying,.
okay, I'm still a pastor..
But, you know, you know what I'm saying, right?.
You know what I'm saying?.
When you answer questions about your essential identity,.
you know, things change all the time..
And Genesis 2 is trying to answer that question..
What is your essential identity?.
And the Bible gives us a solid answer..
I'm God's creation..
I'm designed and created by God wonderfully.
and fearfully made..
I'm a child of God..
I'm in Christ, right?.
I'm part of the people of God, right?.
Would that change?.
Would that identity change?.
No, because God does not change, right?.
Once you are a child of God,.
you're always the child of God..
Yesterday, today, tomorrow, eternally..
You're always the child of God..
You're always worthy in God's sight..
You have worth from someone who never changed,.
God himself..
Once you anchor your identity.
into something that is unchanging,.
that is our eternal God, our creator God,.

$^{521}$God who loves you, God who cares for you,.
your identity is secure..
You see what I'm saying, right?.
So we are the people of God,.
and you are a child of God once you are in Christ..
We are the church, and that anchor is eternal..
That identity never change..
Now, back to this story..
God of Eden is an identity document..
What does it tell us about our identity?.
2.15, chapter two, verse 15..
"The Lord God plays the man,".
and by the way, here is the man, right?.
The man, the Adam, right?.
The Adam, right?.
So Adam is kind of interesting..
He's a unique character, right?.
But Adam is actually, everyone is in Adam..
He represent the whole mankind..
He represent the whole humanity..
So when you read the story of Adam and even Eve,.
you're reading a story for every man and every woman, right?.
I read myself into it..
What does it tell you about the essence.
of being a human being, right?.
Being Adam, right?.
Being a human being..
"The Lord God plays the man," the human,.
"in the garden of Eden to tend and watch over it," right?.
He said, "What is so important about those verbs?.
"To tend and to watch over it," right?.
Now, the early rabbis, the Jewish rabbis,.
already realized these words, this phrase,.
is kind of interesting phrase,.
because this phrase, "to tend and to watch over,".
are used for priests who are working in the sanctuary..
Right?.
How do you describe the job description.
of a priest working in the temple or sanctuary?.
That's how they describe it..

$^{561}$In Numbers 3, verse seven, it talks about the Levites..
Levites are those sacred personnel.
being called to minister in the tabernacle..
It says, "They," the Levites,.
"will serve Aaron," the high priest,.
"and the whole community,.
"performing their sacred duties in and around the tabernacle.".
Now, I underline the word "performing.".
You can do some matching, and that is Genesis 2, 15,.
the word translated "ten," right?.
And the Italic word, can you see that?.
The Levites will serve Aaron..
They will serve, the Italic word is matching, "watch over.".
In Hebrew, they are the same words, right?.
So it's very interesting..
We realize when you go to Genesis 2,.
talking about humanity,.
they're being placed in the first sanctuary,.
Garden of Eden, their job is like a priest,.
to worship God..
And what do the priests do in the sanctuary, in tabernacle?.
They offer sacrifices, they clean up the place,.
they make sure the temple,.
the sacred place is running, is operating,.
and to have fellowship with God..
And most of all, the job of the priest.
is to guard the sacred place,.
guard the temple against anything evil,.
or anything secular,.
anything that does not belong to the tabernacle,.
guarding against those things coming in..
So it's like setting a boundary,.
set a boundary around the sacred place,.
around the place of worship,.
so that unclean, profane, secular things.
will not enter the house of God,.
because the sanctuary is where God dwells, right?.
So now we understand, okay?.
Genesis 2 is talking about Garden of Eden.
is a sanctuary, it's like a temple..

$^{601}$And our identity being placed in the Garden of Eden.
is like a priest, right?.
Like a priest..
We worship God..
That is the original design..
You see, in the world of the Bible,.
the highest honor that a human person can have.
is being a priest..
Someone who's close to God, right?.
Someone who has fellowship with God, to worship God..
And in order to worship God,.
we must first have a right relationship with God..
And who is God to you?.
That is also identity question, right?.
If God is your God, God is the one you worship,.
then you fix that identity,.
because I belong to God, and God created me,.
and I'm people of God,.
and worship flows out of that kind of relationship..
I have God with a right relationship,.
and flows out of that relationship with God,.
I worship Him, right?.
In the 17th century, 1646, many years ago,.
many centuries ago, some centuries ago,.
there was a document called Salted Catechism,.
Westminster Salted Catechism..
Now, what does it do?.
That's a pretty interesting word, catechism..
We don't use it anymore, right?.
It's actually a document that the Bible teachers draft.
in those days to help the new believers.
to get to know the basics of their faith..
So catechism simply means question and answers..
The form is always question and answer..
So in this sort of catechism,.
there were about 100 questions,.
and then you ask a question,.
and then the new believer,.
when they even get ready to get baptized, right?.
It's like a membership card..

$^{641}$They have to answer those questions..
So it's about 100 questions they have to answer correctly.
in order to get baptized and be a member of the church..
Isn't that fun?.
There's also a larger catechism,.
which is more detailed with all the Bible references..
Here, I give you a shorter catechism..
It's a very valuable document to summarize.
the basics of the Christian faith, the biblical faith..
And question number one, ask the question..
Right off the bat, in the entire document,.
the first question is, what is the chief end of a person?.
The older translation is,.
what is the chief end, chief purpose of man?.
But today, we kind of distinguish that..
So I translate as, what is the chief end of a person,.
a human person?.
And the answer is, a person's chief end,.
your number one purpose is to glorify God, to worship Him..
And to enjoy Him forever..
Those two little words are very, very carefully chosen..
To glorify God in everything that we do,.
and worship is a big part of it..
And to enjoy Him, that little word.
is a very carefully chosen word..
Because our relationship with God.
is characterized by enjoyment, right?.
You enjoy your relationship with God..
And you feel happy, you feel fulfilled,.
and you feel joyful,.
because you are related to God properly, right?.
So when you do quiet time,.
when you read the Bible, when you pray,.
do you enjoy it?.
Or you're like, oh my God, this is another thing I need to do..
Okay, I need to brush my teeth..
I need to read scripture, all right?.
Oh my goodness, okay, get it done, right?.
Or you don't even do it, right?.
So think about it, our relationship with God.

$^{681}$is supposed to be an enjoyment..
And that's what Westminster Catholicism,.
this old document is telling us..
We worship God, and we enjoy the relationship,.
because we are created to worship God Himself..
We are being put in the sanctuary, Garden of Eden,.
the first temple, the prototype of temple..
And we are like a priest,.
the highest honor a man or a woman can get, right?.
We are all priests before God, and we worship Him..
And secondly, in this chapter,.
I wanna isolate some of the facts,.
or some of the truth that the passage is telling us,.
is that we are created to be a community,.
a family, to worship God..
2.18, "Then the Lord God said,.
"It's not good for the man to be alone,.
"not good for the human person to be alone..
"I will make a helper who is just right for him.".
Again, it's NLT, New Living Translation..
It's an easier English translation..
But He says, listen carefully,.
"I will make him a helper who is just right for him.".
And by the way, the word helper does not mean.
someone of lower status..
Helper does not mean that..
You know why?.
Because many Psalms, many prayers in the Bible,.
we refer God as helper, right?.
God is not your servant..
God is not your maid servant..
Many prayers in the Bible, particularly in the Psalms,.
asking God, "God, can you help me?.
"You are my helper, you are my fortress.".
Now, so the English translation here,.
helper, don't get the wrong connotation.
that, okay, all right, the woman is the lowest status..
No, I think a better translation is ally, alliance, right?.
In order for Adam to take care of the creation,.
to worship God, he cannot do it alone..

$^{721}$He need another one corresponding to him..
He needs another one just like him, in front of him..
What it means, essentially, a copy..
So you need another person, you need a community..
And the most basic unit for community.
is the family, husband and wife..
So Genesis 2 is trying to answer the identity question..
We are not to be alone..
When we worship God, of course,.
we are first designed and created to be a priest,.
to offer sacrifices and to praise God together,.
but we cannot do it alone..
It's not good for man to be alone..
Have you realized we are different from your pets?.
Dog or cat, right?.
You realize that?.
When a dog looks into a mirror, they bark, right?.
They thought there's another dog in there, right?.
Do your dogs do that?.
Maybe your dogs are super dogs..
But animals are different from humans, right?.
We have the sense of, how do you call it?.
Awareness of self-existence..
I'm a being..
I'm different from other people..
Hey, I'm alone, right?.
Now, the first time I really realized that.
is that when we had Rachel, our daughter,.
and she was about probably one year old.
and started to walk, you know, trying to make troubles..
And I realized that one morning,.
she went in front of the mirror in the living room,.
and she kind of looked into the mirror,.
and she touched her face, moved her fingers,.
touched the mirror, and she did not bark..
She realized, hey, this is the reflection of me,.
and I am oneself, right?.
That sense of, you know, it gave me a sense of wonder..
You know, God created us as someone.
who are aware of our own existence..

$^{761}$And she realized that..
A baby, a kid, realized that..
Now, here's what it's trying to say..
The man, the human person, was being alone..
Exists by himself, right?.
And he aware of that, right?.
Now, in Genesis 1, if you realize,.
if you remember, this six day, seven day creation,.
there's a phrase that is repeated seven times..
When God did something, when God set up something.
in the cosmos, he said, "It was good," right?.
"It was good.".
And seven times he said that..
He appreciated his work, and this is operating,.
this is good, this is functional, right?.
And it was good, seven times..
And the last time, the seventh time, it was very good..
That's all in Genesis 1..
In Genesis 2, here, this is the first time God says.
something is not right..
Something is not good..
And God sees a human person being alone and isolated.
and not connected to anyone, and that is not good..
Right, that is not good..
So if you're living your life, you know,.
you don't connect with anyone in a deeper way,.
God said, "This is not good,".
because God designed human existence.
to be connected socially, relationally,.
with other people, right?.
I think, you know, we praise God for modern technology,.
right, the internet, right?.
Online learning, online streaming, and all the stuff..
Those are really, really wonderful things, right?.
The phone that we have, right,.
used to send Apollo, the rocket, to the moon,.
many years ago, the phone that we have..
It's so sophisticated..
So we praise God for all these tools.
and all these technologies..

$^{801}$However, if we are not careful, we lose a lot of things..
We lose a lot of things..
You remember years of pandemic,.
everything goes online, right?.
Right, is it a blessing or a curse?.
Well, both, right, both..
Kind of like a blessing, right?.
You look at our YouTube channel, right?.
We got 16K subscribers, right?.
I mean, who, you know, it's just a blessing..
We can share some resources.
with other people around the globe..
But it's also a curse..
It's also something, you know, not right..
You know, when everything goes online,.
yes, in some way, it's a blessing..
But, you know, if we're not careful,.
we lose a lot of things..
And one of the things that we lose is human interaction,.
human interaction, like real-time interaction..
You know, I still remember when first the,.
when COVID hit, when pandemic hits,.
all the seminars go online, you know?.
And I saw one of the ads, you know,.
one of the commercials, one of the advertisement.
that says a certain seminar is offering a fully online,.
100\% online Master of Divinity degree, you know?.
You know what Master of Divinity is?.
You know, everyone's preparing himself to be,.
or herself to be a pastor, to be a minister..
You have to at least complete the three and a half years.
or four years Master of Divinity curriculum, right?.
Now, why is it called Divinity?.
It doesn't mean who studied MDiv is more divine, you know?.
Again, it's an old word, right?.
Back in the 17th century,.
Westminster Confessional Faith,.
they called those who studied the Word of God,.
who studied the things about God, divine, right?.
Doesn't mean they're more sacred, right?.

$^{841}$But it's an old word, right?.
That's why the degree is called Master of Divinity, right?.
But you can do it totally online, right?.
Now, I totally oppose that..
I oppose that..
Now, you can study other degrees online,.
but MDiv is to prepare people to be pastors,.
and ministers are to relate to other people, right?.
We got a rough shoulder with other people, right?.
You understand what I'm saying?.
I mean, a pastor is not a YouTuber, right?.
We don't do things online mainly..
We do things in real time, right?.
There's a Jewish psychotherapist..
She's Jewish, Esther Perl..
Esther Perl has done a lot of works on human relationship..
And I watched one of her interviews,.
and she laments the modern age,.
the loss of real human interaction..
We lost it, right?.
We're always with our phone, right?.
We don't talk to each other..
We don't get our phone, right?.
Even we talk to each other, the one next to you,.
you text them, right?.
You go through some of the device, right?.
So we lost that real time human interaction,.
and she laments that..
And she talks about the older generation,.
like me or people who are older than me..
I still remember those days we were kids,.
and we have to interact with other kids..
We're on the playground, we're on the street, right?.
The interaction is real time..
It's spontaneous..
It's instantaneous, right?.
You gotta respond to other people..
You gotta learn how to rub shoulders.
with other people all the time..
Real time, right?.

$^{881}$There's no device in my days, right?.
So we were on the playground..
We beat up some kids, you know?.
Real time, you know?.
And the kids beat us up, right?.
And we got in a fight, right?.
We got into a fight..
So that's how we learn..
That's how we learn how to interact with one another..
I'm sorry, my background is kind of rough, okay?.
I know you guys are very, very good, okay?.
You guys are very good, always at home,.
you know, always doing your homework,.
and you get 4.0 or whatever, you know?.
But you know what the problem is?.
It's that most of the social interaction.
are done online today..
Now think about texting..
Texting..
You can always wait, right?.
You receive a text, and you kind of think through it,.
how to respond to it, right?.
On Zoom, you don't turn on your camera, right?.
Yeah, you're kind of absent, right?.
Are you there?.
Yes?.
No?.
Most of the time, you're not there, right?.
You disappear, right?.
And we lost that..
You know, we just don't know..
No, for people who are younger,.
especially much younger than you,.
we just don't know how to interact with other people,.
how to relate and rub shoulders with other people..
The other day, I was having a hot pot.
with some of the friends,.
and then I saw a few tables, you know,.
the other side of the restaurant..
They were doing solo hot pot..

$^{921}$Now to me, this is inherently wrong..
I mean, I understand you have, you know,.
it's a very good deal, right?.
But you know, you're facing one pot,.
and then tons of food,.
and then you consume the whole thing..
And it's all by yourself..
You don't talk to anyone, right?.
I mean, hot pot, essentially, is a community thing, right?.
Don't you agree?.
Right, don't you agree?.
I mean, we lost that..
You know, my wife works with Urban Christian Academy..
It's an after-school program,.
and she always complained to me in the past two weeks.
is that the young students today,.
they don't have any social skills..
They don't even talk to you..
They don't say good morning..
They don't say goodbye..
They don't say thank you..
And they just lost all of this, right?.
And that is not good, according to God's design..
God designed us to not to be an island..
We are created to be connected with one another..
There's an innate desire in all of us,.
in every single one of us,.
to be related to other people, to a friendship..
Part of a family, part of something bigger than myself,.
the community, the church, right?.
We have to relearn that..
You know, we have to relearn that..
And worship, in the biblical sense, cannot be done alone..
I don't know how great the experience is,.
all by yourself, watching a TV, watching a screen,.
and doing your worship or thinking about worship..
That is not the ideal design..
God wants worship to be done in a corporate way,.
you know, collective way..
Sit with one another, seeing one another,.

$^{961}$and it's a corporate experience, to sing songs together..
When you sing praises,.
do you realize that you can use your years?.
Right, you do have two years, right?.
And you listen to other people singing as well..
Well, may not be singing so well,.
but the experience of worship is a corporate experience..
We sing together..
We don't come and listen to a song..
We come and worship God together..
That is the meaning of worship..
You're designed to be a worshiper..
God created you, put you in a temple,.
to be like a priest,.
to serve him and to worship him..
And we are being designed to do that,.
connected as a community, as a spiritual family,.
as a church, to offer our worship to God..
Let us all pray together..
Father, we come before you..
We repent of our sin of selfishness,.
our self-centeredness..
Often we think of ourself.
and the culture will have us believe.
that we can just do life alone..
And we don't need one another..
We don't need other people..
We don't need our counselors..
We don't need our spiritual mentors..
We don't need our friends..
We don't need our brothers and sisters..
We just,.
we use the whole life into individualistic experience..
And we confess our fault and our sin before you..
And I pray that your word will awaken us.
so that we will see that your original design,.
the beautiful design for us to exist.
is to be connected with one another in the community,.
to support one another, to pray for one another,.
to rub shoulders with one another..

$^{1001}$And thank you for your word reminding us.
our original purpose is to worship you..
Thank you for this highest honor.
you granted to all of us,.
that we can worship you in your presence..
And God, I pray that you will help us to recover.
not only the sense of wonder and awe,.
how great you are,.
but we'll recover the sense of togetherness,.
the sense of worshiping you, serving you together..
And thank you for the opportunity we can have today.
to worship you together..
So as a church, we ask that you will continue to renew us.
as one body, as one community,.
to worship you and to present you offering.
and sacrifices that will please you..
Thank you for the time that we have..
We ask that you will send us out.
to continue to live for you in our Monday to Saturday,.
challenge us and help us to follow you..
And help us to connect with one another,.
to bear witness to our great God,.
our great and loving Savior, Jesus Christ..
In Christ name we pray..
Well, thank you very much for coming..
And next week we're gonna serve you,.
our serve God together, all right?.
We have refreshment, right?.
Yeah, please stay around for a time of interaction, okay?.
\newpage



\section{士師記}
\label{sec:H_MgtrWeR24}
\textbf{從懼怕到召命的覺醒 (士師記 6) - 莫家承牧師}
\newline
\newline
連結: \href{https://youtube.com/watch?v=H-MgtrWeR24}{\texttt{https://youtube.com/watch?v=H-MgtrWeR24}} ~~~~ 語音日期: 2024-04-22
\newline
\newline
\hyperref[sec:cEtF7vpVnrE]{\small{< < < PREV SERMON < < <}}
~
\hyperref[sec:index_chronic]{\small{[返順時目]}}
~
\hyperref[sec:index_scriptual]{\small{[返順卷目]}}
~
\hyperref[sec:g90s_0fTe_c]{\small{> > > NEXT SERMON > > >}}
\newline
\newline
$^{1}$大英姐妹平安,很感恩能夠在網上跟大家一起學習神的話.
亦多謝城會教會的邀請,讓我可以在此與大家分享信息.
我是卡神,我是一位在多倫多牧養的牧者.
我所服侍的對象是一些職稱的群體,包括已婚或有年輕子女的家庭.
同時,我亦是多倫多眾人旋律一位義務的牧者.
與弟兄姊妹用音樂及不同的藝術表達去分享福音的一個機構.
我亦是一位學生,在大學修讀實踐神學課程.
正在寫論文,寫到一半了,求神繼續給我智慧去完成它.
我亦是一位父親,一位丈夫,一位兒子.
當我有時思想,神給我有不同的身份和角色.
我一步一步揭示,神在我的生命中究竟想我走一條怎樣的人生路.
而又很奇妙的,這些身份角色又彼此有關連.
在一個身份中所學到的東西,又會幫助我更加明白.
我如何在另一個身份中活出神生意中的我.
我也想起我從小是一個很多事情害怕,很多懼怕,很多憂慮的人.
但我看到的是,神一步一步讓我從戰經的事情中成長起來.
更加懂得依靠他,順靠他,亦慢慢明白多了自己.
神幫我建立整全的真我.
所以今天和你選擇的講題是從懼怕到照明的覺醒.
從懼怕到照明的覺醒.
我不知道你的生命中,你的成長中有沒有一些東西.
令你經常很纏繞,很害怕的.
或者你會否經常問,究竟神在我的人生中應該怎樣給你用.
基督教經常說照明,我們經常說神在我們生命的心意.
究竟神在我的生命中的心意是甚麼?.
是否只有一個方向,還是可以是縱向的?.
我未必能夠回答你所有心中的問題.
但我可以用今天基甸的記載和屬靈中的神給我的經歷.
我們彼此去面對.
基甸的人生很有趣,他有很多東西可怕.
在不同的人生階段.
如果我們將他整個人生分為上下半.
在前半的記載,他都很害怕.
他害怕甚麼?.
他害怕他聽錯了,聽不清楚.
這是否神的心意?.
有時這種害怕,讓我們人有一種謙卑在神面前.
神給基甸這樣越來越經歷他.
他的人生就有所提升成長.
雖然回到後半的記載,他所害怕的是甚麼?.

$^{41}$就是害怕別人看不到他的成功.
害怕別人看不到他的成就.
他整個屬靈生命就開始下滑.
我們看士師記,我們都很快重溫一些很基本的主線.
我們記得士師記經常有句說話提到.
就是以色列人遊行耶和華,眼中捍衛惡的事.
裡面一次又一次記載到.
以色列人如何經歷順利所賜的平安.
不過就慢慢忘記了神.
被那些外邦的風俗,信仰,文化價值觀影響.
然後離開耶和華,也同時敬拜別的神.
神就用一些外邦人去管教自己的百姓.
在苦況中,人就懂得向神去呼求,哀求.
神就憐憫自己的百姓,興起不同的事事.
去拯救自己的百姓,再次帶來一種平安.
不過就像很多學者所形容.
這個士師記的編輯刻意呈現一種下行螺旋的結構.
甚麼是下行螺旋呢?.
就是每況愈下.
這些士師記的整卷書並非按時序去敘述.
而那些史詩,你回頭看.
那些開頭,譬如說,鄂托涅,爾忽,底波拉.
這些都是一些好的領袖,史詩.
但慢慢再看,就開始越來越不像樣.
譬如說,最後參孫.
一個很沒有自制能力,脾氣很差的人.
其實都反映出每況愈下.
同樣,你看那些平安的日子.
在士師記亦都越來越短.
亦都反映出那個越來越差的人.
越來越不懂得去珍惜與他同在的上帝.
基奠其實就是在這個史詩行列裡面的一個分水嶺.
所以看的時候.
我們將他的經歷都化為我們的提醒.
六章一開始就說到這個處境.
就是當時在基奠的生活環境裡面.
其實面對一個外在的挑戰.
就是那些米電人.
米電人就會在他們播種或收割的時候來搶東西.
你去到毀壞他們的農作物.

$^{81}$搶了他們的一些牲畜.
羊,牛,驢等等.
令他們的生活,經濟極其困乏.
在這樣的威脅下.
以色列人就懂得向神去哀求.
神用環境去管教自己的百姓.
亦都同樣神用言語去告訴他們.
這個管教的原因和目的.
有先知說到.
就是在亞摩利人的地.
以色列人在敬拜那些人的神明.
忘記了耶和華.
當然這不單止是在說信仰上的敬拜.
是一種價值觀的扭曲.
人慢慢覺得.
咦?他們那套信念.
他們所嚮往的那種風俗.
很吸引.
人慢慢一心二意.
神在這樣的環境管教自己的百姓.
就是這樣的佈局.
帶出了這個故事的主人.
基奠.
六章十一節是這樣形容的.
約阿斯的兒子基奠在一個壓酒池打小麥.
為躲避米顛人.
打麥通常是一個空曠的地方.
因為這樣打的時候.
能夠將麥和一些不需要的部分.
很容易分開.
酒渣是一個密封的.
可能是一些洞穴.
是小的空間.
室內的.
這樣打就很沒有效率.
因為你會打散了之後.
將你要的東西和不要的東西混在一起.
但這裡告訴我們.
他是因為害怕.
米顛人來搶東西.

$^{121}$他每一天的生活.
都是在擔驚受怕裡面.
去到過的.
在這樣的狀態下.
神跟他說.
強大的勇士.
耶和華與你同在.
其實也挺諷刺的.
我每一天就是.
卑微微小的.
戰戰兢兢的.
今天所辛苦.
艱重下的.
明天被人搶去.
我哪裡像強大.
我哪一塊強大.
我想這也是基甸聽到的時候.
那個心聲和反應.
所以他就開始說到.
這個列祖所說的.
這位大能的神.
好像忘記了我們.
不過耶和華就跟他說.
不是的.
耶和華必與你同在.
並且是靠著.
這個我所給你的力量.
你去拯救.
以色列人脫離米顛人的手.
當然基甸聽到的時候.
他就覺得.
不可能是我.
我很微小.
他說我的家族在這個.
本來也不是很大的.
馬來西亞支柏那邊是最弱的.
我在我的父家裡.
也是最微小的.
我是小之又小.
他很害怕.

$^{161}$他的自我形象很低.
不過.
耶和華第二次跟他說.
我必與你同在.
一次兩次的.
神其實在強調一個訊息.
給基甸聽.
就是你去拯救以色列人.
你去擊退米顛人.
不是因為你的出身.
不是你的能力.
不是你的學識.
是因為我與你同在.
這個是令你不需要再害怕的理由.
這個是你能力全員所在.
我想起我自己的成長.
我剛才說的.
我是一個很散的人.
體弱多病.
我讀書也很平庸.
或者很差.
人家讀一兩次的.
我要讀反覆溫習很多次.
如果我們小時候打球.
要猜情球分隊.
我通常是被人選剩的.
所以我的成長.
很多時候是比下去.
我的自我形象很低.
我也有很多事情擔憂.
因為我覺得競爭.
社會競爭.
學校的競爭很大.
我好像一無是處.
認識了上帝之後.
神告訴我.
我的價值並不是出自我的出身.
不是我的學識.
不是我的成績表.
不是我在公司的職位.

$^{201}$社會上賺到多少錢.
我的價值是源於.
有一個上帝.
祂愛我.
為我犧牲.
珍惜我.
縱然如此.
我在教會裡的事奉.
很多時候.
開始時是躲在後面.
不是在人前.
我覺得這樣比較適合我的性格.
但一路在教會成長.
神說.
我不是只要用那些.
很有學識很有學問.
說話很動聽的那些人.
我要用你.
走到人的前面.
為我作見證.
用你對我的認識.
在生命裡.
屬靈的那種體會經歷.
去分享福音.
去告訴別人.
我是一個又真又活的上帝.
當時聽到我很害怕.
但在這種和神的拉扯下.
我學習信服.
並且我記得第一次聽到神跟我說.
我是叫你站在人前面.
講福音.
那時候我還沒進神學院.
我跟教會的牧者說.
聽到神給我一個呼召.
要我跟眾星傳說出去.
但是我講福音.
但我不懂得講.
教會的牧者很好.
說我教你說.

$^{241}$因為他看到.
神在一個人生命裡的呼召.
是要踏上去.
去驗證.
亦有同行者那種認證.
他就成為我鼓勵和幫助.
很實在的幫助.
亦是靈裡的同行的幫助.
在這裡我們看到.
機電的夢召.
他聽到神跟他說之後.
他也要求一個憑據.
所以他就說.
我想帶出禮物給這位使者.
神就預納了這個食物.
這個禮物成為一個祭物.
很多學者就看到.
機電對使者的相處的款待.
其實是在映照亞伯拉罕創世記十八章三節.
那種對使者熱情的款待.
有些像摩西在荊棘寂靜裡.
聽到上帝的呼召.
上帝的聲音.
這樣的敘述.
其實想帶出一個意思.
就是機電在神的眼中.
他是一個很有價值.
很有潛質的領袖.
一個神體重.
覺得很有潛質的傭人.
其實今天你在神的眼中.
也是很有價值.
一個很有潛質的傭人.
不過神說不是按著你的能力.
不是按著你的學識.
是因為我在你生命裡的同在.
和所給你的能力.
神給我們每一個人的走的路.
可能都不同.
但卻有一樣相似的.

$^{281}$我們不是按著自己的心意.
不是按著自己的那種的 Ambition.
去到揀選要走的路.
而是上帝說.
你來到我面前.
讓我和你同在.
同在去做我要你做的那件事.
那條的路.
機電聽到這個心意.
然後神就給他第一個任務.
我叫他做一個一石二鳥的任務.
因為這個任務不是打仗.
是回到他成長的地方.
回到他的村.
在他的父親.
在他的鄉里面前作一個見證.
這是一個信仰上的宣告.
真神的宣告.
同樣也幫助機電.
去到面對他成長的陰霾.
重整他的信仰觀.
重整他的真我.
這個舉動是什麼呢?.
我們來看看.
六章二十五節.
神叫他.
去拿你父親的公牛.
和另一隻七歲的公牛.
來做什麼?.
去獻祭.
那個七歲的公牛.
當我們回想.
大概是神將以色列人.
交給米地人的手中出生.
是一種呼應.
也有贖罪的意思.
神就叫機電拆掉.
他父親巴力的祭壇.
然後在祭壇旁邊的亞瑟拉.
那個木柱.

$^{321}$就砍了他.
當柴燒.
這樣去到將他所信的神.
呈現在他的.
鄉親父老面前.
當然對機電來說.
這是一個很大的掙扎.
神清楚地告訴我們.
二十七節說.
於是機電找了十個僕人.
就照著耶和華的吩咐去做.
不過不是日光日白去做.
是夜間暗地裡去做.
因為他害怕.
父家和城裡的俄弗拉人.
他也是害怕的.
但他害怕得知道這是神的心意.
他就去踏上.
並且這個選擇.
這個任務.
其實是驅使他去想.
在他生命裡邊.
哪一個才是真正的主宰.
哪一個神才是最可靠.
是巴力.
是當時所敬拜的愛邦神明.
還是耶和華.
我也不是一個基督教家庭長大的.
所以我想起.
當神一步一步說.
你去裝備.
去神學院裝備.
放下你的工作.
讀神學.
給我去用.
我也要回去跟我家人說.
我家人和很多華人家庭一樣.
都是崇尚勤力.
也是一種默默耕耘的美德.
同時我們是人生輸送.

$^{361}$但我們覺得這樣才是成功的.
讀好的學府.
出來找好的工作.
在工作裡爬升.
然後買車買房子.
結婚生子.
兒子又誕生聽話.
讀名校.
諸如此類.
我好像要回去.
面對我的家庭.
我成長的地方.
告訴他們.
有一套人生觀.
信仰觀其實是更好的.
那對我來說是一種掙扎.
因為好像同時否定.
你以前教我的那一套.
未必是最永恆.
最永久的.
有一套是耶和華所吩咐.
所教導的.
是更加真實.
是更加人需要嚮往的.
但就在這個過程.
我慢慢地就.
先在神面前.
重整我的人生.
我的信仰.
什麼才是我人生的核心.
是我的主耶穌.
然後我就要將這個問題.
拋給我身邊的人.
其實有一個更值得考慮的.
人生的取向.
就是跟隨.
相信耶和華.
這位真神.
而不是世界所高舉的那一套.
在基甸的故事.

$^{401}$其實也是這樣.
因為見到當他.
漏夜去拆了巴力的壇.
然後獻祭.
然後第二天那些人就起來說.
這件事是誰做的.
是誰做的.
就是要拿他出來懲罰他.
那基甸的爸爸約亞斯.
就站在所有這些攻擊基甸的人面前.
就說了.
你們要為巴力爭辯.
你們要救他嗎.
為他爭辯的早晨就該處死.
巴力的祭壇已經被拆毀了.
如果巴力是神.
巴力會為自己爭辯的.
在這裡約亞斯當然是一種救子心切.
所以說出來.
但同時我相信.
在約亞斯.
他的生命.
都立刻浮現一個問題.
究竟是那些神明.
愛邦的神明可靠.
還是這個二話才是可靠.
你看到基甸的戰戰兢兢踏上的路.
所呈現的見證.
在人的生命裡是留下問題的.
是留下一個反思的問題.
生命信仰的問題.
同樣看那些俄弗拉人的生命.
都留下這個問題.
所以他們就稱基甸叫耶路巴力.
就是與巴力爭論的意思.
特別的地方就是.
這是別人給他的稱呼.
因此任何人提起他.
說開他的時候.
其實都會令人聯想到這個事蹟.

$^{441}$這個事蹟的核心問題.
就是哪一個才是真神.
你看到每一個微小的傭人.
被上帝使用.
所帶來的感染力.
是一種生命與生命的交匯.
留下生命的問題.
在別人的心當中.
這是一個感染力.
這其實就是我們基督徒進入世界.
那份的照明.
至少是照明之一.
我們需要在別人的生命裡留下問題.
透過我們如何見證我們的信仰.
透過我們的選擇.
人生的取捨.
未必是需要打鑼打鼓.
但我深信.
當我們願意.
降伏在上帝的引領的時候.
當我們願意.
跟著祂所吩咐的去踏上.
神會使用你,使用我.
在別人的心裡留下問題.
所以神同樣在我的家人裡.
都問了這些問題.
因為我的家人會問.
為什麼你會進入神學院.
為什麼你會不做.
你讀了很多年的專業.
究竟背後的動力,能力,理由是什麼.
基督徒就是要在不同的地方留下問題.
讓別人去反思.
我深信上帝在裡面.
會用我們的見證.
我們看到第六章最後.
就是在說機電的禱告.
機電發覺原來神一步一步.
真的在帶領以色列人.
去擊退米電人的時候.

$^{481}$他又再在神面前求多一次印證.
這個我們叫做羊毛的禱告.
第一次他說.
如果周圍都是乾的.
羊毛是濕的.
我就知道是你的心意.
神應允他這個很特別的請求.
神很有耐性.
所以當他再問的時候.
神都沒有發火.
他說如果這次.
你再給我一次更大的肯定和印證.
如果周圍都是濕的.
羊毛這次是乾的.
反過來.
我就知道是你的心意.
其實都頗超自然的.
那個濕氣.
如果周圍都有濕氣.
為什麼羊毛可以不濕呢.
神又再一次應允他這個祈求.
剛才我說.
我在神學院讀書.
很快去到最後一年的時候.
我在裡面有一個教授.
師父.
經常捉我聊天.
他說卡神.
你要好好考慮.
究竟你是否要進入教會.
作全時間的服事.
我生命裡其實都很掙扎.
很多的恐懼.
他和我祈禱.
他經常提醒我.
你去到好像機電那樣禱告.
Pray the Gideon Prayer.
我都不以為然的.
我都說.
神啊.

$^{521}$如果你真的要我走這條路.
你給我更清楚的明證吧.
如是者就去到畢業前一個月.
當時我要和眾恩旋律.
去到Montreal.
日本的地殼.
做幾場的報道會.
我們就開車去.
所以一車一車很多樂器.
咪高峰.
燈光.
音響.
器材等等.
就去到那裡.
從多倫多出發.
頭兩場聚會.
我們都可以將我們所預備的.
這些訊息.
故事.
用這些器材.
很立體地呈現出來.
有很多人認識耶穌.
亦有大英姐妹的生命被鼓勵.
我們很感恩.
去到第三場.
以為今晚繼續.
神繼續大大地使用.
我們就去到.
在星期日的下午.
做彩排.
做一些set up.
通常在正式開始.
和我們排練中間.
有一個小時左右.
我們會祈禱.
會休息.
吃飯等等.
就是在那一刻.
忽然間教會停電.
教會停電都常常有的.

$^{561}$我們都會check一些fuse.
看看電箱是不是有些跳了掣.
但又發覺不對勁.
因為不只是單一教會停電.
是整區停電.
對面的醫院都在行後備電.
這樣就知道.
事情的嚴重性.
當時我們在吃飯.
我記得禮堂.
因為日落.
越來越暗.
我就是面對著.
禮堂和教會的大門口.
以我一個這麼驚青.
這麼多事情懼怕的人.
我應該是震了震了.
但神卻讓我從心底裡.
經歷了一種.
很出人意外的平安.
我在神面前.
有一個心意.
和祂祈求.
神今晚.
你的心意是什麼呢?.
求你顯明.
我就很平安地.
吃了幾口飯.
一邊吃的時候.
我看到街坊已經進來了.
吃完早飯.
打算早點來拜位.
我忽然間感受到.
原來人心真的很嚮往.
經歷福音.
經歷福音的大能.
人很需要信仰.
人很需要耶穌基督.
如是者.
我和弟兄姊妹說.

$^{601}$不如我們禱告.
看神怎樣帶領我們.
成就自己的心意.
無論是取消還是怎樣做.
我們不知道.
我們祈禱.
祈禱之後.
很快就有人說.
其實我們找到三支電筒.
我們教會有一個.
走走調音的鋼琴.
也可以用.
你一樣可以清唱.
你一樣可以用電筒照明.
禮堂.
用古舊的鋼琴.
去彈奏音樂.
我們就決定.
我們今晚繼續用神所給我們.
僅有的.
我們去見證祂.
很奇妙.
我發覺.
我們以為.
我們很依靠的器材.
咪高峰.
燈光.
這些是很幫助到的東西.
但拿走了之後.
反而人更加的專注.
人在一個漆黑昏暗的禮堂裡.
聽到生命的主的訊息的時候.
人有一種反思.
所見證的.
原來神在任何情況下.
祂都可以成就他的奇事.
更加神說.
你放下你原本有的那棍的講章.
那些預備了的powerpoint.
沒有電你都用不到.

$^{641}$你只是說一個訊息.
約翰一書.
一章五節.
神就是光.
在它裡面是無黑暗.
我就是這樣去分享.
神就是一位怎樣的上帝.
去到最後.
我邀請弟兄姊妹.
開手機照著地下.
給願意出來決志的人.
去有一個表示.
有些朋友包括一位伯伯.
是督著拐杖.
一步一步走到台前.
他說我需要得著這個生命的光.
這一份的真光.
這件事情過了之後.
我在神面前.
有一種很長時間的沉澱.
我問神.
你是否垂聽了.
應運了我那份的請求.
神說你自己覺得.
我說是.
我說神.
我明白多了.
原來走全時間.
服事的道路.
不是我有什麼厲害.
甚至乎不是我要去.
在這個場景大展拳腳.
原來就是.
我很嚮往.
有神同行的人生.
才是一個精彩無限的人生.
神說我想這樣去過我的人生.
神說你來無恙我的人.
我就是這樣決定了.
去走無恙的道路.

$^{681}$並且一生的立志.
真是達人.
去到他的根前.
所以我很喜歡這張照片.
這張照片就是當時弟兄姊妹.
在一個漆黑的環境.
用手機幫我拍.
我當時在講道.
一個電腦.
我當時在講道.
一個電筒照下來.
這張圖畫提醒我.
不是人的能力.
是上帝的同在.
上帝的出現.
去成就事情.
不過我有那份福氣.
可以參與其中.
弟兄姊妹我不知道你的生命在哪裡.
你是否依然有恐懼.
纏累著你.
令你不能夠.
更加給上帝使用.
去踏上多一步.
我鼓勵你.
因為我們不是每一個人.
都需要做傳道人.
做牧師.
但每一個人在上帝的設計裡.
都有一個獨特的路.
特別的心意去走.
我們一起在神面前.
勇敢地走多一步.
我深信當你願意這樣做的時候.
神所給你經歷的.
是更真實的信仰人生.
是一個更有價值.
更無憾的人生.
並且在人的生命裡.
留下一些信仰的問題.

$^{721}$生命反思的問題.
願上帝幫助我們.
在不同的角落.
我們活出信徒應有的感染力.
願上帝幫助你.
幫助我.
我們同心禱告.
信賴的天父上帝.
你的救贖.
你的好恩.
成為了我們生命裡.
最大的祝福和動力.
有著一份獨特的心意.
願你幫助我們.
活出生命裡的照明.
願你幫助我們.
看到你給我們不同的位份和角色.
都能夠反照出.
這個在我們生命裡.
有真有活的上帝.
同樣我們今天生命裡有恐懼.
願你幫助我們.
在裡面成長過來.
讓我們更加看見.
我們應該著眼的.
不是那個環境所帶來的束縛.
限制或挑戰.
讓我們目光看見.
這位滿有大能的上帝.
如何使用我們.
使我們的生命.
更加經歷你的同在.
更加使人明白.
這位有真有活的上帝.
正在在人的心裡.
在世界裡.
作王掌權.
願你使用我們的生命.
也在我們每一天的生活.
為人裡得著最高的榮耀.

$^{761}$頌讚.
奉耶穌基督得勝靈之祈求.
阿們.
\newpage



\section{創世記}
\label{sec:g90s_0fTe_c}
\textbf{永恆的敬拜 2 (創世記 2) - 陳崇基牧師}
\newline
\newline
連結: \href{https://youtube.com/watch?v=g90s_0fTe-c}{\texttt{https://youtube.com/watch?v=g90s\_0fTe-c}} ~~~~ 語音日期: 2024-04-22
\newline
\newline
\hyperref[sec:H_MgtrWeR24]{\small{< < < PREV SERMON < < <}}
~
\hyperref[sec:index_chronic]{\small{[返順時目]}}
~
\hyperref[sec:index_scriptual]{\small{[返順卷目]}}
~
\hyperref[sec:Hq5No0HLb30]{\small{> > > NEXT SERMON > > >}}
\newline
\newline
$^{1}$好,我們這次繼續看敬拜的系列.
我們有幾次會講到什麼叫做敬拜.
上一次我們講到詩篇第八篇.
這次我們講到聖經的開頭.
我們看看聖經裡面整體有什麼對敬拜有教導的.
首先我們去看一段聖經.
就是《創世記》第二章.
我們回到伊甸園.
我們先讀經文.
二章第七節.
耶和華上帝用地上的塵土造人.
將生命的氣吹進他的鼻孔.
這人就成了有靈的活人.
耶和華上帝在東方的伊甸栽了一個原子.
把所造的人安置在那裡.
耶和華上帝使各樣的樹從土地裡長出來.
可以越人的眼目好作食物.
原子當中有生命樹和知善惡的樹.
第十五節我們跳到.
耶和華上帝把那人安置在伊甸園.
讓他耕耘看管.
耶和華上帝吩咐那人說.
原種各樣樹上所出的.
你可以隨意吃.
只是知善惡的樹所出的.
你不可吃.
因為你吃他的日子必定死.
耶和華上帝說.
那人單獨一個不好.
我要為他做一個配偶.
幫助他.
伊甸園的故事或者創世記頭幾章.
我們經常覺得是一些很簡單的故事.
大家有沒有.
特別有子女的.
有沒有兒童聖經.
有吧.
兒童聖經有沒有這些故事.
一定有的.
一定有亞當夏娃在草叢裡.

$^{41}$一定有一些動物等等.
我們經常覺得這些是兒童故事的簡單敘述.
不是很難的.
但大家知道如果真的細心去查聖經的話.
其實有很多事情不是那麼簡單的.
特別是頭幾章.
是吧.
我們覺得很簡單.
因為我們沒有想到很多問題.
或者我們查聖經查多一點的時候.
我們漸漸就會發現有很多問題.
可能聖經裡面應該有什麼答案.
或者甚至有沒有答案.
譬如亞當有沒有肚臍.
因為亞當是第一個人.
他不是從孿生出來的.
不用剪臍帶.
夏娃也是一樣.
或者我們會很簡單的去看.
會不會聖經告訴我們.
這個地球其實是很接近發生的事情.
創造的事情只不過很接近.
諸如此類的問題.
我們查經的時候會想一下.
會不會這裡有些答案.
或者應該這麼說.
會不會是沒有答案.
或者聖經想回答什麼答案.
我們想看這段聖經.
講到伊甸園的時候.
一個似乎很簡單.
但裡面其實有很多的記號.
有很多的symbol.
有很多的符號.
是一個充滿符號的故事.
我們在看這個故事之前.
我們需要想一想.
如何去看一個故事.
如何去看聖經裡面的故事.
有一個例子我覺得是可以幫助到大家.

$^{81}$去如何去看這些所有事情.
好像很顯淺.
但其實裡面可能是隱藏著一些符號.
是很深奧的東西.
如何去看到那些真理.
我一個例子就是這樣.
就是地圖.
大家有地圖的.
會看地圖.
Google Earth 或者Google Map.
你去某個地方.
你就打那個地址下去.
然後它就告訴你.
什麼時候轉左.
什麼時候轉右.
還有上哪條路.
那些都是地圖.
地圖如果你懂得看的話.
你就應該明白.
看地圖是有技巧的.
我們很多時候沒有想過如何看地圖.
但我們其實很下意識地.
我們就會覺得地圖是這樣看的.
我們就明白了.
對不對.
這裡有三個地圖.
一個是芝加哥地圖.
一個是香港地圖.
一個是日本地圖.
因為每個人都去日本.
我特意放一個日本地圖在這裡.
去吧.
給你地圖.
三種不同的地圖.
右邊那個.
芝加哥地圖.
看到嗎.
我們大概在所謂近紅色那個地方.
粉紅色下面那個地方.
然後有不同的區域.

$^{121}$它就用不同的顏色來代表.
然後左上角的地圖.
其實有交通在那裡.
就是有交通的路線.
交通塞不塞等等.
有時候有紅色.
有時候有黑色.
有時候黃色.
有時候綠色.
然後左下角的地圖.
就是一個地勢的地圖.
Touring.
就是它告訴你.
富士山你看到嗎.
看到吧.
是不是想去.
Yeah OK.
它告訴你東京這裡附近有一個海灣.
然後左邊那裡有些山.
其中一個就是富士山.
左下角那裡就是富士山.
有不同的地圖.
就有不同的代號.
來告訴我們.
地上一些真實的東西.
某一方面.
沒有一個地圖會告訴你所有的東西.
因為不可能的.
譬如芝加哥的地圖.
我們在那格的地方就是粉紅色的地方.
你不可以拿著地圖.
然後我們出到街上.
說為什麼地下不是粉紅色的.
你不會這樣看的.
你絕對不會這樣看.
因為那只是一個代號.
但是你不可以說地圖的創作者.
畫地圖的人是錯的.
因為它有它的代號.
有它的顏色.

$^{161}$有那些所謂的傳說.
去告訴你地面上.
事實上發生的事情.
我想告訴你.
有一些不同的信息.
沒有一個地圖是說了所有的東西.
也沒有一個地圖.
如果真的這樣畫出來的時候.
就是去欺騙我們.
或者是騙我們的.
對嗎.
因為每一個地圖都會透露不同的資訊.
在事實上.
某一些的資訊.
你就不能夠說.
右邊的地圖不是地勢的地圖.
所以那是錯的地圖.
三個地圖三個都對的.
三個地圖都是表達真理.
也是表達真理的某一方面.
對嗎.
我們看地圖的時候.
很少會說這些東西.
你說為什麼要說這些東西.
我們只是知道怎麼看而已.
但我想告訴你.
看聖經其實都是一樣.
好像看地圖一樣.
你要明白那些傳說.
明白那些象徵.
到底是在表達什麼.
想說什麼.
所以看《創世》的第二章.
其實有一點點像是地圖.
正如我剛才所說.
這個故事有很多的記號.
好像一個地圖有很多不同的顏色一樣.
到底它想解釋什麼.
或者透露什麼資訊.
這裡我想告訴你.

$^{201}$聖經特別頭幾章.
其實不是解答我們很多現代人的問題.
其實主要要解答一個更重要的問題.
就是我們為什麼在這裡.
我們的身份.
我們所謂的身份.
上帝創造我們.
我們在這個世界裡面.
我們出現.
我們的生存.
到底我是誰.
我是一個怎樣的人.
創造主他的計劃.
創造主他的設計.
到底是怎樣去回答這些問題.
所以如果我們看《伊甸園》的時候.
Eden.
或者《伊甸園》的聖經.
我們明白.
原來好像一個地圖一樣.
它有不同的顏色.
有不同的標誌.
我們就想一下.
應該怎樣去看這段聖經.
當我們這樣說.
好像看地圖一樣.
看一個《伊甸園》的記載.
好像看一個symbolic story.
一個敘述裡面有不同的符號的時候.
不會令到這個故事.
是少了歷史性.
也不會令到這個故事.
所敘述的東西.
是少了真實性的.
一樣的真實.
只不過寫出來的方法.
說出來的方法.
是一個symbolic的方法.
一個記號的方法.
等於我們看地圖一樣.

$^{241}$所以如果我們看《伊甸園》的時候.
你記得.
這段聖經應該大家都很熟悉.
或者有些印象.
《伊甸園》門口是向哪個方向的?.
東面.
向東面.
然後有金有銀.
有樹木.
有花草等等.
門口是誰守的?.
記得嗎?.
是誰守門口?.
不記得了.
是基魯帕.
Cherubim.
基魯帕.
基魯帕是把守聖所的活物.
天使.
他把守著門口.
是向東.
這幾樣東西.
剛才我所說的.
你就大概知道《伊甸園》.
到底是一回什麼事?.
一個怎樣的地方?.
譬如說.
出埃及之後的會幕.
以色列人在曠野.
然後說你做一個會幕.
我住在你中間.
會幕就好像我們紮營一樣.
紮了又收.
然後繼續上路.
那個叫Tabernacle.
會幕.
會幕是有什麼守著的?.
聖所有什麼守著?.
就是基魯帕.
是不是?.

$^{281}$會幕是門口向哪?.
向東.
到所羅門.
他做一個石的聖殿.
在耶路撒冷.
也有基魯帕.
守著聖所.
把守著聖所的門口.
我想問你.
聖所或者聖殿的門口向哪?.
向東.
向東很簡單.
伊甸園向東.
會幕又向東.
然後聖殿.
也是向東.
你就知道《伊甸園》.
伊甸到底是一回什麼事了.
所以伊甸園.
其實好像一個聖所.
或者應該是第一個的聖所.
第一個冒犯的聖所.
這裡想解釋什麼?.
就是說上帝創造萬有的時候.
他去設立這個世界的時候.
其實就好像設立一個聖所.
是一個神聖的地方.
第一章.
你曾經聽過我講.
是七天來創立這個世界.
OK?.
Again 又是一個symbolic story.
去告訴你.
不是一些地質地理.
遠古歷史的問題.
不是科學的問題.
他想解答的就是.
這個世界存在.
到底有什麼意義?.
所以創世記第一章用了七天.

$^{321}$來交代整個創立世界的過程.
就是說整個世界.
就好像一個聖殿一樣.
你可以叫宇宙大殿一樣.
當你看到窮滄.
看到大自然這麼美麗的時候.
那個設計這麼精心的時候.
創世記第一章說.
It's a sacred place.
我們站在這個地方.
享受大自然的時候.
是一個神聖的地方.
是上帝居住的地方.
所以第七天.
上帝就安息了.
他就住在他的所謂聖殿.
整個宇宙的聖殿.
然後第二章.
就是我們今天看的《英靈經文》.
就好像鏡頭.
去到焦點放在某一個地方.
叫做伊甸.
伊甸.
Eden.
也好像聖所一樣.
有基督拍手向著東.
然後有金與銀.
到底他想說什麼.
想解釋一個身份的問題.
為什麼我們在這裡.
為什麼你存在.
當我們說到身份.
說到identity的時候.
很多時候我們都會.
我們不是這麼多想這件事.
譬如你認識一個新朋友.
你介紹你自己.
我們通常會怎麼介紹自己.
可能我是一個學生.
我是一個牧師.

$^{361}$我是一個財政人員.
會計師或者教師等等.
或者我是一個家庭主婦.
我是一個經理等等.
這些都是我們做的事情.
因為我們一個星期七天.
大部分時間都是我們工作.
所以有時候我們的身份.
就是從我們工作而來的.
我們做的事情而來.
但你知道我們做的事情而來.
如果你這樣去定我們的身份的話.
這個就不是很絕對.
因為經常會變.
我還沒做牧師之前.
我做什麼的.
我是化驗室裡面做的.
因為我是入實驗室按緊急制的人.
因為我是讀化學的.
我就在實驗室裡面做.
你問我那時候我是誰.
我就在實驗室裡面按緊急制的人.
我是一個化學師.
然後我現在就是牧師.
所以這些事情經常都會變的.
到我牧師要退休的時候.
我還是不是一個牧師呢.
或者你的職業轉變的時候.
你轉行的時候.
你到底是哪一個呢.
怎樣去解答你的身份的問題.
identity的問題.
到底你是誰.
我是誰.
或者你介紹你自己的時候.
很多時候你不會用.
因為這樣的方法.
譬如我們知道那些人.
人體.
你的化學的成分.

$^{401}$大概有六樣.
譬如carbon 水 hydrogen等等.
你不會介紹你自己.
百分之九十九我就是水.
氮和氧氣等等.
這些碳加起來.
這個就是我.
你明白我的意思嗎.
什麼叫做identity.
我是誰.
聖經開頭這段.
還有整本聖經.
其實就是在解答這個問題.
你為什麼在這裡.
Why are you here.
Why do you even exist.
原來就是上帝創造我們.
你就是上帝所創造的奇妙的人.
我們這個人生.
或者我們這個存在.
是一個奇妙的創造.
我們信主的人更加明白.
我們就是上帝的兒女.
上帝的兒女.
我們教會.
我們神的子民.
這個答案.
用上帝來定我們的身份.
就是不會變的.
因為上帝不會變.
你是上帝的兒女.
上帝所創造的.
今天 明天 明日.
直到永恆.
你都是不會變的.
我們必須要搞清楚.
到底我們是哪一個.
我們的identity.
我想告訴你.
創世的第二章.

$^{441}$甚至第三章.
一二三.
其實就是一個身份的文件.
去解釋你和我是哪一個.
這個和敬拜有什麼關係.
我們抽兩段經文.
在這裡看一看.
創世的第二章.
然後說上帝把那人安置在伊甸園.
讓他耕耘看管.
我們第一個的反應就是.
我的後丫.
好像一個花園.
我們要掙些草出來.
或者要去清潔等等.
耕耘看管.
到底是什麼意思.
其實早期猶太人的拉比.
都意識到這件事.
你看他們早期的解經的講論.
拉比就說耕耘看管不是種花種草這麼簡單.
耕耘看管是什麼字.
他們就會聯想到.
是一個祭司會做的事.
因為祭司在聖所.
在會幕到底他做什麼的.
他就是打理聖殿.
你再看看耕耘看管.
去找找這些字出現的時候.
我考考你.
民數記第三章第七節.
說到利未人.
利未人就是在聖所裡面工作的人.
那些聖職的人員.
他說他們要替亞倫即是大祭司.
又替全會眾在會幕前執行任務.
辦理帳幕的事.
你看到哪個字是和創世的二章十五節重複的.
有underline的執行任務.
在會幕裡執行任務.

$^{481}$那個字就是耕耘那個字.
然後打斜那個辦理帳幕會幕的事務.
辦理這個字就是耕耘那個字.
你想像.
辦理是看管那個字.
辦理就是看管那個字.
你不會想到一個利未人.
在聖所裡面他打工的時候.
做事的時候.
他居然是在種花的.
他們會收拾器具.
去set up(安排)一些東西.
然後make sure(確保).
肯定我們的敬拜生活能夠繼續.
如果你這樣看的時候.
你就發現了.
原來上帝創造人.
Adam.
本身給他的職事.
他的身份是什麼呢.
就是一個祭司.
祭司就是要敬拜上帝.
在聖所裡面去事奉.
是和上帝有一個關係.
有一個所謂的團契 fellowship(團結合作).
而最重要的就是祭司要做的事.
就是去到守衛著聖所.
因為聖所以外的地方.
就是繁俗的地方.
繁俗的東西不潔的東西.
不應該進來的東西.
他們就會守衛著不讓他進來.
就是說聖所就是holy的地方.
聖所以外的我們叫做profane(聖俗).
就是說secular(聖俗).
sacred和secular(聖俗).
是誰去把守的分別.
就是祭司去把守.
所以你看到舊約裡面五經.
經常說到那些.

$^{521}$以色列人應該怎樣去敬拜.
那些有殘疾的東西不應該帶進來.
不應該獻給上帝.
或者不潔的不應該進來.
這個就是祭司要做的事.
原來創世記說到伊甸的時候.
原來就是說到敬拜.
原來上帝創造你和我.
本身就是一個敬拜神的人.
是和上帝去建立這個關係.
能夠在這個關係的基礎上.
我們去敬拜.
去讚美祂.
這個就是我們的identity.
we were created to worship.
去敬拜上帝.
敬拜不是今天我們坐在這裡.
有人帶歌我們唱幾首歌就完了.
原來我們生命上帝創造我們.
設立我們的時候.
就是和祂之間有一個關係.
這個關係是怎樣表達出來的.
在敬拜表達出來.
在17世紀的時候有一個信經.
叫做Westminster's Slaughter Catechism.
是什麼來的呢.
當時是一個教導Sunday School.
一個主一學的課程.
這個課程一直歷史很悠久.
到現在都有人在用.
雖然我們教沒有用.
這些課程其實就是一些信仰文語答的書本.
裡面有百多條很簡單的文語答.
這個叫做Slaughter Catechism.
就是小要理問答.
也有本大的叫做Larger Catechism.
大一點的就是詳細一點的.
這些東西是當時每一個信了主.
在教會裡面他要受浸加入教會的時候.
他就需要背這些所謂Catechism.

$^{561}$這些文答.
有很多條例如十誡是什麼你要背.
然後聖經裡面哪裡說什麼你要背.
今天不是很幸福嗎.
你不用好像考牌一樣.
但他的意思是什麼呢.
就是一些教理.
一些要理是我們需要掌握的.
Slaughter Catechism第一個問題就是問什麼.
就是人一生首要的目的是做什麼.
What is the purpose.
The chief purpose of life.
人生首要的目的是做什麼.
是敬拜榮耀神.
以神為樂直到永遠.
這兩個字是很妙的字.
To glorify God.
然後呢.
To enjoy Him.
去到享受祂.
以祂為樂.
這個就是人生最基本的identity.
最基本的身份.
無論你做哪一行哪一職業.
你生命的階段是什麼都好.
這個全部會變的.
但有一樣東西不會變.
就是你人生裡面最重要的身份identity.
就是和上帝有關.
敬拜上帝.
去榮耀上帝.
所以這裡創世記講到這一點的時候.
有兩個重點.
第一個我們被造就是要敬拜上帝.
第二.
創世記第二章.
講到我們是一個敬拜的群體.
是一個事奉上帝的群體.
為什麼這樣講.
第二章講到什麼.

$^{601}$第二章第十八節.
講到他set up這個伊甸園這個聖所的時候.
上帝又說.
那人.
所謂亞當.
單獨一個不好.
那人獨居不好.
或者以前和日本.
我要為他做一個配偶.
來幫助他.
現在做一個人.
所謂亞當.
其實不是一個人這麼簡單.
因為他代表整個人類.
看亞當的故事.
就是看每個人的故事.
是寫著我們的故事.
我們組織在亞當裡面.
哪人獨居不好.
是單獨一個是不好的.
我要為他做一個配偶.
或者英文翻譯.
Helper.
幫助者.
And by the way.
幫助者.
之後我們知道是做一個女人出來.
所以成為一個家庭的單位.
一個community.
一個社群.
一個人是不好的.
我為他做一個配偶.
一個helper.
By the way helper.
不是代表他低等一點.
下級一點.
因為詩篇裡面經常講到.
耶和華是我的helper.
我的幫助者.
你問耶和華大還是你大.

$^{641}$私人這樣祈禱.
耶和華你是我的helper.
Please help me.
耶和華當然大.
上帝當然大過我們.
其實這裡更加好的翻譯應該是.
Ally.
同盟者.
聯盟者.
我要為他做一個聯盟者.
去幫助他.
但這裡很奇怪.
他說哪人單獨一個不好.
哪人獨居不好.
一個人alone.
一個人是不行的.
上帝創造人是一個群體.
而創造一個人去敬拜他.
也不是一個單一一個人去敬拜.
而是一個社群.
一群人一起來敬拜.
有什麼identity.
有什麼身份的記號.
他說to be alone.
他自己一個.
一丁有.
你有沒有留意過.
人和你的寵物.
你的狗.
是不同的.
人和狗是不同的.
狗如果看著鏡子的話.
牠會吠的.
因為牠不知道鏡子裡面是牠自己.
牠當牠是另一隻狗.
牠不停地吠.
因為我知道你的狗可能聰明一點.
你不會這樣做的.
你的狗可能超人狗.
但動物是沒有那種個人存在.

$^{681}$I am alone.
I am a person.
我exist.
這種所謂自我的意識.
我什麼時候留意到.
我記得我生第一個小孩之後.
我女兒.
她大概會走的時候.
大概一歲左右.
我記得有一次我見到她.
我們聽說有塊鏡子.
她第一次看那塊鏡子.
然後她摸一摸那塊鏡子.
然後在那裡撥一撥.
我覺得很有趣.
我見到她在那裡撥.
她意識到這塊鏡子.
起碼她不會吠.
她發現這個人是我.
鏡子裡面那個不是另外一個人.
鏡子裡面只不過是我的反應.
我是存在的.
是真真實實存在.
I am a being.
My self.
這裡說到亞當存在的時候.
是一個是不好.
只有他一個.
是不好.
《創世記》第一章.
有一句話是重複的.
就是上帝創立了一些東西之後.
他說什麼.
看作是好的.
It is good.
It was good.
然後第七次說什麼.
It was very good.
上帝創立了整個宇宙.
一個宇宙大殿.

$^{721}$神聖的地方.
Very good.
然後第二章.
這裡是第一次聖經.
上帝說話.
Not good.
意思就是說.
這樣是行不通的.
功能上是行不通的.
因為上帝創造人.
就是要管理大地位.
一個人是搞不定的.
一定要一班人.
一個群體.
而這個群體最基本的單位.
就是夫妻.
就是男與女在一起.
所以沒有一個人.
上帝創造他的時候.
上帝的意思永遠都不會是.
他是一個孤島.
一個孤立者.
他一定是和其他人有一個聯繫的.
所以宅男宅女是不合乎聖經的.
宅男宅女是不合乎聖經的.
因為他是孤立自己.
It is not good.
這裡提出了一件事.
就是我們人性.
是需要有一個所謂的social connection.
和社群裡面有一個關係.
好幾年的疫情.
已經打亂了我們這件事.
我上個星期說過.
很多時候我們的敬拜.
都是看著電話就算了.
當然了.
因為有互聯網.
或者有這些Internet.
有這些科技.

$^{761}$其實是一件很好的事.
有很多祝福.
我們成為網台有很多祝福.
但是如果不小心的時候.
我們就會損失了很多東西.
就是人與人之間.
我們英文叫做rub shoulder.
Rub shoulder.
我記得好幾年前.
疫情開始的時候.
全部東西都上網.
陸續有不同的神學院.
他們就將課程放上網.
所以現在有很多神學院的學位.
是可以就這樣上網讀.
你就可以畢業.
好不好.
有好的.
你在貧瘠的地區.
你沒有這樣的教育資源.
那你上網就可以了.
我也有Online class.
我也在教.
但是有些什麼不小心的時候.
我們就損失了很多東西.
我們以為我們聽完東西.
我們再學到東西.
有一樣東西最令我驚奇的.
就是他將M.Div的program.
全線上網.
就是說你可以全部上網.
你就可以拿到一個M.Div的degree.
M.Div的degree是什麼.
我們中文翻譯叫做道學碩士.
就是每一個去訓練傳道人.
訓練一個minister.
訓練一個牧師的時候.
他們要go through一個課程.
那個課程叫做M.Div Master of Divinity.
Master of Divinity.

$^{801}$Divinity就是神聖.
不代表道學碩士不如那些人變成神.
不是這個意思.
而是以前古字用divine來形容神學家.
那些神學家叫divine.
我們叫devotion.
所以有M.Div這個課程的時候.
就說Master of Divinity.
就是說他們是divine.
他們是學神學的.
所以他們給他一個碩士.
我最驚奇的一點就是.
為什麼可以讀M.Div可以全線online.
你知道我們牧者.
去訓練一個牧者.
訓練一個牧者是要對人的.
對嗎.
牧者不是對著一個螢光幕的.
而M.Div這個program.
很多時候就是沒有了.
這種所謂human interaction.
你可以讀其他的program是online.
That's fine.
但是如果你讀M.Div online.
我覺得是一個很大很大的問題.
因為你去到教會.
你需要面對實體的人.
你怎樣去跟人互通.
怎樣能夠跟人溝通.
這些東西根本在這幾年已經到了一個地步.
我們完全沒有了這些訓練.
有一個猶太的心理治療師.
叫做Esther Perry.
Esther Perry做很多工作.
或者說很多東西.
寫很多東西.
是跟人間的人際關係有關.
我記得有一次.
看一個podcast.
還是看一個網上的訪問.

$^{841}$訪問Esther Perry.
她問我們現在現代的人.
最大的問題是什麼.
她就說我們現代的人.
沒有了人與人之間的互動.
特別當科技繼續發達的時候.
我們現在說的有AI.
有很多東西全部都是在電話可以做到.
在電腦可以做到.
她說就是損失了一樣東西.
是最寶貴的.
就是human interaction.
就是你跟我之間有一個互動.
我記得我們上一代.
叫做上一代.
我們小朋友.
怎樣學人與人之間的互動.
我們在遊樂場玩.
我還記得我小時候.
放學後去街上.
或者有時候不上學.
去街上玩.
我們就跟其他人有互動.
誰不順眼就去打他.
然後我們又被人打.
然後我們就還擊打他.
這個打不打的東西.
不是打機打船可以做到的.
我們親手打的.
human interaction就是Esther Perry所說的.
什麼叫human interaction.
真正的人與人之間的互動.
就是面對面.
是你不能等.
而你需要去作反應.
你跟一個人聊天.
實體.
是不能你想一下怎樣反應的.
你立刻就要反應.
有人打你你應該怎樣.

$^{881}$你打他.
這些就是我小時候.
對不起你們很乖的.
我小時候在駛士區.
我們就是這樣生活的.
人打人我們就打人.
但現在的問題是什麼呢.
就是我們的social interaction.
我們的社交的互動.
很多時候都是網上做的.
texting.
短訊.
短訊的問題是什麼.
當然有好處.
你可以想清楚再說.
不要說錯話.
但很多時候.
人家發短訊給你.
你再回覆的時候.
你就會想一下怎樣回覆.
但如果你是真人的話.
你就不能等.
你好嗎.
你不能等三分鐘才回答.
你很古怪.
這些是怎樣學的.
是實體學的.
人家發一個電郵給你.
人家check你.
你都可以假扮一下.
不回答.
或者ghost他.
當沒看過.
我知道你們上網.
有時候祈禱會.
不開鏡頭.
吃飯不開鏡頭.
其實我們實體吃飯.
我們都是這樣聊天的.
你說我食相不好.

$^{921}$平時吃飯我都看見你這樣.
對嗎.
我們不開鏡頭為什麼.
因為我們覺得有些東西.
我們不可以和別人互動.
這個就是社群的失落.
我們的世界已經很粉碎.
如果再加上人與人之間.
我們沒有那種community.
我們就更加失落.
對嗎.
在疫情的好幾年.
我們有一個term叫做COVID baby.
對嗎.
COVID baby.
COVID baby大概現在的understanding.
就是為什麼會有個COVID baby出現.
就因為在疫情裡面.
他見少人.
就是小朋友見的人越來越少.
以致他很多時候有社交的焦慮.
因為他小時候已經不見很多人.
對嗎.
特別是陌生人.
他不知道怎樣和他互通.
對嗎.
所以有很多觀察.
都是說這些lost human touch.
所以一個baby這樣長大的時候.
沒有足夠的human touch.
會令到他成長.
是有缺陷的.
社交上emotional.
或者情緒上.
對嗎.
Baby應該是怎樣的.
好有趣 先摸摸他.
他被人摸摸.
他才知道原來一個成長是這樣的.
我太太.

$^{961}$在我們一個ministry partner那裡打工.
做一些課後的補習.
她最近經常說一件事.
她說現在的年輕人.
不懂得和你互動.
連多謝都不會說.
連說話看著你都不敢.
到底發生什麼事.
對嗎.
上帝的design是怎樣.
上帝的設計.
就是我們每一個都不是孤立的人.
那人獨居是不好.
是我們需要一個群體.
我們需要去接待 去學習.
大家所謂rub the shoulder.
對嗎.
我知道現在很多是work from home.
很好啊 work from home.
very good i like it.
可以有一杯咖啡在那裡.
穿一件睡衣 穿一件短褲.
你穿不穿衣服都沒有人理.
總之我做起事就搞定.
你少了什麼.
你損失了什麼.
就是office politics.
沒有辦公室的政治玩.
那些就是人與人之間.
我們損失了的東西.
I don't know.
我們損失了那個community.
損失了那個互動.
敬拜 聖經裡面說.
真正的敬拜.
不是一個孤立者做的事.
因為上帝做人.
他的設計.
就是我們先與上帝有關係.
relationship 互動.

$^{1001}$我們祈禱 我們讀經.
上帝創造人.
是放他在伊甸園裡面去敬拜 去事奉他.
不是單獨.
不是一個個體.
而是一個群體.
一起做.
我希望我們教會.
我們去學習.
敬拜的時候我們學習這件事.
是有人與你rub the shoulder.
實體的.
這個才叫做真正的敬拜.
真正討上帝喜悅.
對嗎.
可以同心得來祈禱.
上主我們多謝你.
因為你創造我們.
我們是奇妙精心所創立的.
原來我們的生命就是來敬拜你.
從起初你的原創計劃.
你就說將我們放在這個世界裡面.
就好像住在聖殿一樣.
有你自己的同在.
有你自己的presence.
主讓我們用歌聲來讚美你.
讓我們用我們的生命來讚美你.
去事奉你.
我們能夠去敬拜.
能夠與你建立關係.
一個屬靈的群體.
讓我們知道如何去彼此.
在一個群體當中學習.
去相愛.
學習彼此實體的生活.
以致我們這個生活能夠成為一個敬拜.
去尊你的名為大.
是給我們的救恩.
主讓我們成為你自己屬神的子民.
去敬拜你.

$^{1041}$去榮耀你.
並且以你為樂.
我們這樣祈禱.
防止病毒.
下星期希望見到大家.
\newpage



\section{10:25-37}
\label{sec:Hq5No0HLb30}
\textbf{Kingdom Values XLV: Luke 10\_25-37}
\newline
\newline
連結: \href{https://youtube.com/watch?v=Hq5No0HLb30}{\texttt{https://youtube.com/watch?v=Hq5No0HLb30}} ~~~~ 語音日期: 2024-04-26
\newline
\newline
\hyperref[sec:g90s_0fTe_c]{\small{< < < PREV SERMON < < <}}
~
\hyperref[sec:index_chronic]{\small{[返順時目]}}
~
\hyperref[sec:index_scriptual]{\small{[返順卷目]}}
~
\hyperref[sec:Gigyg3Yfrs0]{\small{> > > NEXT SERMON > > >}}
\newline
\newline
$^{1}$Hey, family, I'm excited to be with you again this week.
and to jump into the Word of God together.
and discover what God has for us.
and the truth that He has for us.
to encourage us in the places that we're in life right now.
and the things that we're going through.
This week, we're actually going to step away.
from the Sermon on the Mount for just a little bit.
and we're going to look at another kingdom value.
and that is the kingdom value of loving our neighbors.
And we're going to be looking at this today.
just because there's so much going on in our world right now.
And as believers, we want to respond to those things.
We want to love our neighbors well.
And Jesus has some really good advice for us.
in the scripture about how to do that.
And it's almost overwhelming at times.
all the things that we see going on in our world.
We have this coup that's going on in Myanmar.
where the military has taken control of the country.
and there's all types of violence.
and people being oppressed.
and voices that are being suppressed.
and people being hurt.
And there's this terrible situation in Hong Kong.
where freedoms are being stolen away.
and dissent is being stifled and silenced.
And we have all types of religious violence in the world.
We have places in the Middle East and Africa.
where Christians are persecuted and assaulted.
We have violence right here in the city of Chicago.
We have violence on our streets.
where people are harmed and killed and shot every day.
in the city of Chicago.
And we have all these other problems in the United States.
We have all these issues with racism that we've seen.
just be so out in the open over the last year.
And even in the last week, we see this incident.
that involves possible racism against Asian Americans.
And so there's so much stuff happening.

$^{41}$And it's hurtful and it's terrible.
and it's discouraging to see.
But Jesus wants us to respond to those things.
as the people of God, as the subjects of the kingdom.
It's a kingdom value for us to love the people around us.
and to love the world around us.
And it can be difficult to know.
and understand how we can do that.
So that's what we're going to look at today.
We're going to look at this story in Luke chapter 10.
It's a really familiar story.
It's the story of the Good Samaritan.
I think there's some good advice for us here.
that Jesus himself gives us.
about how to love our neighbors.
how to love those around us.
So let's read this together.
and then we're going to dig into it a little bit.
We'll begin in verse 25.
This is Luke chapter 10, verse 25.
It says this.
And behold, a lawyer stood up to him.
to put him to the test, saying,.
Teacher, what shall I do to inherit eternal life?.
He said to him, what is written in the law?.
How do you read it?.
And he answered, You shall love the Lord your God.
with all your heart, with all your soul.
with all your strength, with all your mind.
and your neighbor as yourself.
And he said to him, You have answered correctly.
Do this and you will live.
But he desiring to justify himself.
said to Jesus, And who is my neighbor?.
Jesus replied, A man was going down.
from Jerusalem to Jericho.
and he fell among robbers.
who stripped him and beat him.
and departed, leaving him half dead.
Now by chance, a priest was going down that road.

$^{81}$and when he saw him, he passed by on the other side.
So likewise a Levite, when he had come to the place.
and saw him, passed by on the other side.
But a Samaritan, as he journeyed.
came to where he was.
and when he saw him, he had compassion.
He went to him and bound up his wounds.
pouring on oil and wine.
Then he set him on his own animal.
and brought him to an inn and took care of him.
And the next day he took out two denarii.
and gave them to the innkeeper.
saying, Take care of him.
and whatever more you spend.
I will repay you when I come back.
Which of these three do you think.
proved to be a neighbor to the man.
who fell among the robbers?.
He said, The one who showed him mercy.
And Jesus said to him, You go and do likewise.
Now in this story, what's going on here is.
there's this discussion taking place.
and this lawyer, he's an expert in the Jewish law.
comes to Jesus and in verse 25.
tells us he's kind of trying to trip Jesus up.
He wants to make Jesus look bad.
and ask him a really tough question.
and he says, How do I inherit eternal life?.
And so Jesus puts the question back on him.
and says, Well, what does the law say?.
How do you read it?.
What's important that you see in the law?.
And so this man says, Well, I think.
the most important thing in the law.
is to love God with all your mind.
and your strength and your soul.
and to love your neighbor as yourself.
And so Jesus says, You've answered correctly.
But this man doesn't want to leave it there.
And I love this verse, verse number 29.

$^{121}$where it says, He seeking to justify himself.
desire to justify himself.
He says to Jesus, Who is my neighbor?.
It's such a great verse.
because that's all of us as humans.
We're always trying to justify ourselves.
to make ourselves look good.
in the sight of God and others.
And that's what he's doing here.
He's saying, Yeah, but Jesus.
you've got to understand.
like, I've already done this, you know.
And if there's a problem here.
it's not because I'm not doing it.
It's maybe just because I don't understand.
what it means, some of the words mean.
So who is my neighbor?.
Why don't you tell me who my neighbor is?.
And then Jesus responds.
with this amazing story.
He tells a story about a man.
who is going from Jerusalem to Jericho.
and he gets robbed, brutally beaten.
And he's laying by the side of the road.
Bleeding, and his clothes are stolen.
And he's left for dead.
He's laying there dying.
by the side of the road.
And people begin to pass him by.
The first person that Jesus tells us.
that passes is a priest.
He's coming down that same road.
He sees this man.
He goes to the other side and passes by.
And next comes a Levite.
And he does the same thing.
He sees this man and passes him by.
on the other side.
And then Jesus says.
A Samaritan comes.

$^{161}$And this is huge.
This fact, the hero of this story.
is a Samaritan.
The Samaritan, the Jews.
were people who did not like each other.
at all for ethnic reasons.
So there's so much about race and ethnicity.
in this passage.
That's a whole, could be a whole.
series of messages.
And that's not what we're looking at today.
So I don't have time to unpack all that.
But we need to unpack that in some way.
Because this is huge.
what Jesus is teaching here.
by making the Samaritan the other.
the person who was despised by the Jews.
He's the hero of the story.
He's the one who actually shows.
the character of the kingdom.
And he's actually obedient in a way.
that the Jewish priest and Levite were not.
But he's the one who comes along.
and shows compassion.
He goes to this person.
who's dying by the side of the road.
He binds up his wounds.
He takes care of him.
He banishes him.
He makes sure that he's okay.
And then he puts him on his animal.
on his donkey.
that he was traveling with.
Takes him to an inn.
a place of lodging.
where this person can be cared for.
And make sure that he's cared for.
And then in verse 35.
we see that he even pays for it himself.
And gives money to the innkeeper.

$^{201}$and says take care of him.
And when I come back.
if there's anything else.
that needs to be taken care of.
I'll pay for all of it.
And so he takes care of this stranger.
in an amazing way.
And Jesus closes the story.
by asking who was the neighbor.
And clearly it's a rhetorical question.
The answer is the Samaritan.
And we see this expert in the law.
unable to even name the Samaritan.
And just saying the one.
who showed him mercy.
But Jesus is showing us a lot here.
about who our neighbor is.
And maybe you're familiar.
with this story.
Maybe you've heard it told before.
And maybe you've heard people.
kind of focus on the fact that.
our neighbor could be anyone.
And as I said in the context.
of this story ethnically.
It could certainly be the people.
that are different from us.
That maybe we don't want to love.
or to care for.
That's a huge part of this story.
But what I want us to focus on today.
Is how we're supposed to respond.
to a neighbor.
Whoever that is.
Wherever that is.
Who is suffering.
Because we do see so much brokenness.
And so much suffering in our world.
And as the people of God.
We're supposed to respond.

$^{241}$But how do we respond.
How can we respond.
And I think there's a lot of truth.
In this story for us.
So the first thing.
I want us to consider.
Is that we should.
And we must allow God.
To interrupt the rhythms of our lives.
With the sufferings of others.
Look at verse 33.
It says but a Samaritan.
As he journeyed.
What's he doing.
He's on this road.
From Jerusalem to Jericho.
We don't know why.
It doesn't tell us that.
We don't know if maybe.
He's journeying for business purposes.
Maybe he's a merchant.
He's going to one of those cities.
To sell something.
Maybe he's journeying.
To meet family members.
Maybe he's going to a wedding.
Or a special occasion.
We don't really know.
Because the story doesn't tell us that.
But he's journeying.
As a part of the rhythm of his life.
Obviously.
No one just takes a long journey.
And this was a long.
And dangerous journey on this road.
Because this road was known.
To be frequented by bandits.
No one takes a long.
And dangerous journey for no reason.
There was something.

$^{281}$In the rhythm of his life.
That he was doing.
Like I said.
Maybe it was for personal reasons.
Maybe it was for business.
Or work or something.
But there was something.
That he was doing.
And that.
Whatever that purpose was.
It was interrupted.
When he saw this suffering person.
And let's contrast that.
With the priest and Levi.
Who they are going about.
The rhythm of their life.
And they do not allow.
The rhythm of their life.
To be interrupted.
By this suffering person.
And again.
This story doesn't tell us.
Their reasoning necessarily.
I'm sure there are a lot of reasons.
But we can imagine.
Maybe some of the things.
They were thinking.
Maybe for the priest especially.
Maybe it had to do with the fact.
That he was ceremonially clean.
And he wasn't supposed to touch.
Things that were unclean.
And maybe he was saying.
Oh that person is bloody.
And if I go over there.
And get involved.
It'll create a mess.
And it might make me unclean.
And then I'm not able to serve.
In the way that I'm supposed to.

$^{321}$In the temple.
Isn't that ironic.
That perhaps the priest.
Was afraid of the mess.
That would be involved.
In getting involved.
In the life of someone.
Who's suffering.
And that's what kept him.
From actually doing the thing.
That his faith called him to do.
From loving his neighbor.
In the way that God wanted him to.
What kept him from it.
Was that he felt like.
He was serving God.
By remaining clean.
Isn't that ironic to think about.
That he felt like.
Oh what I have to do.
And the way I have to do it.
Is so important.
That I can't interrupt.
The rhythm of my life.
The Levi.
He's a servant in the temple.
And again we don't know exactly.
What's going on in his head.
But maybe he was like.
Well I'm busy.
I've got to get to the temple.
To do this thing.
Or I have this important role to play.
Or I'm supposed to take care of something.
And I just don't have time.
Or maybe these other roles in my life.
I feel like they're much more important.
Than stopping to deal with this suffering person.
Or maybe they were afraid.
Maybe they just felt like.

$^{361}$Well I don't know.
I've got to hurry.
Because the bandits might come back.
I mean clearly something terrible happened.
And I just need to hurry.
And get safe as quickly as I can.
Or maybe they were just uncomfortable.
We don't really know why.
Or for whatever reason.
The priest and the Levi.
Didn't allow God to interrupt.
The rhythms of their life.
They didn't allow God.
To interrupt them on their journey.
Whatever they were headed to.
They didn't stop.
And see this person who was suffering.
And say I want to help.
And as followers of Jesus.
As a subject of the kingdom.
We've got to learn.
To let God interrupt the rhythms of our life.
With the suffering of others.
When we see someone.
Across the street or across the world.
Maybe it's in our neighborhood.
Maybe it's right there.
Maybe it's a news story.
We see on the other side of the world.
Or maybe it's a friend or a family member.
That we have in another place.
Across the world.
Or even just across the city.
Or in another state.
And we hear about suffering in our life.
We need to allow those things.
To interrupt the rhythm of our life.
Don't just say you know what.
I'm busy.
I've got more important stuff to do.

$^{401}$And what I have to do.
Is really a big deal.
My job, my friendships, my relationships.
I can't stop to help someone.
We should allow God.
To interrupt the rhythms of our life.
And let's think about how God does that.
You know a lot of times God does that.
With things that have proximity to us.
Things that are close to us in some way.
And I think that's important for us to consider.
Obviously there's so much brokenness in the world.
And so much suffering in the world.
It's completely impossible for one person.
Or even for one group of people.
To wrap their minds around.
All the suffering in the world.
And to care about all of it.
And to pray about all of it.
And to help in all of it.
There's no way one person could do that.
And so I'm not saying that we have to try.
To involve ourself in every problem in the world.
And every time anything happens.
We have to take care of everything.
And every person who has a need.
We have to meet that need.
No, that would drive us crazy.
And completely exhaust us.
There's no way we can do that.
But we do need to understand.
That the things that happen in our life.
Happen for a reason.
And so when God interrupts our rhythm.
When he crosses our path with something.
That's because he does want us.
To have our life interrupted by it.
And to respond to it.
Again, we don't have to respond to everything.
But to the things that God brings into our life.

$^{441}$We should be ready and willing to respond.
And that's what we see here with the Samaritan.
He wasn't out on the road.
Trying to find people who are hurt.
And help every person.
That might have potentially been hurt on the road.
He wasn't out on the road.
Trying to find the bandits.
And to stop them from robbing people.
He was just going about the rhythm of his life.
But God brought this suffering person.
Close to him.
And he allowed that suffering.
To interrupt the rhythm of his life.
And when you find yourself.
In a place of proximity.
To someone who is suffering.
Then you should.
That should be a light bulb.
That goes off in your head.
And you say, wait a second.
Maybe God is allowing this to happen.
Because God wants me to get involved.
In this situation or in this person's life.
God wants me to respond in some way.
I just want to remind us.
That that proximity.
Isn't only a geographical proximity.
In this case, it was a geographical proximity.
The Samaritan was walking back past.
And he saw this person suffering.
But it doesn't just have to be that.
It could be relational.
It could be someone.
That maybe they're not close by you.
In physical proximity.
But you have a relationship with them.
They're a friend or they're a family member.
Or they're someone that you know.
And you see them suffering in a way.

$^{481}$And you say, hey, I think.
That God wants me to be involved.
And respond there.
Because he's brought this suffering.
Close to me.
And he's brought it to my attention.
And he's put it in front of me.
And I have a role to play in this.
It could be some other type of connection.
It could be some type of affinity.
Something that you are interested in.
Or something that you understand.
And you see someone.
Who is in that same kind of affinity.
That does the same things you do.
Or maybe they're the same age as you are.
Maybe they're someone you go to school with.
Or someone you work with.
But someone that you're connected with.
In some way.
And that's how God uses.
That connection is what God uses.
To bring that suffering in front of you.
And to bring it close.
And to allow you to see it.
It could be a lot of different things.
But just have your eyes open.
And be aware of those things around you.
Again, no one should feel like.
They need to fix every problem in the world.
That's completely impossible.
And if we think we can do that.
Then our theology is way off.
Only Jesus can do that.
The only reason that.
We're even supposed to be involved in all that.
Is because we're supposed to be.
The body of Christ.
We're supposed to respond in love.
And to serve.

$^{521}$So it's not about us.
It's not about me trying to.
Be some superhero out there.
Saving everyone.
It's not that at all.
But it is about Jesus.
And his love for the lost.
And his love for our world.
And his love for the broken.
And his love for the hurting.
And his love for the oppressed.
And because of his love for our world.
He has told us as his people.
That he wants us to respond.
And to love and serve.
And so we should be looking for those ways.
That Jesus brings that suffering.
Into our lives.
And not just ignore it.
Not just pass it by.
But allow it to interrupt.
The rhythms of our life.
So that we can respond.
And become involved.
And sometimes there's even suffering.
Outside of your proximity.
It's not something.
You would normally notice.
That God brings before you.
And be open to that as well.
The type of ministry I do now.
The neighborhood I live in.
And raise my family in.
Is a place that is very different.
From where I grew up.
I didn't grow up in the state.
I didn't grow up in Chicago.
But God brought me here.
In response to needs that were here.
That God cared about.

$^{561}$It's not about me.
God cared about the needs that were here.
And so he brought me here.
To respond to him.
And to obey him.
And to love others.
That was outside my proximity.
But God brought it before me.
And God might do that for you.
God might bring something before you.
That he wants you to be engaged with.
And so be open to that.
And be willing to respond.
But that brings us to this question of.
If we're supposed to respond.
How do we respond.
What does response look like.
You may say, okay, I understand that.
I need to allow God.
To interrupt the rhythm of my life.
With the suffering of others.
And now I see the suffering.
And it's brought close to me in some way.
But what do I do.
How do I respond.
Well, first and foremost.
We have got to pray.
As followers of Jesus.
As subjects of the kingdom.
That should be the first thing.
That we always do.
We should pray for these situations.
Whether it's across the street.
In our neighborhood.
Whether it's in the life of a family member.
Or whether it's across the world.
We need to get on our knees.
And pray for these situations.
We need to pray.
Because the Bible talks about.

$^{601}$The power of prayer.
That prayer is more powerful.
In a lot of cases.
Than anything else we can do anyway.
And we need to believe that.
We need to believe that.
Before we go out.
And try to fix things.
We should talk to God about it.
Because God is so much more powerful.
Than we are anyway.
And so we should start with prayer.
And we should pray.
Because our prayers can go into places.
Where we cannot.
Maybe you say, okay, I'm concerned.
With an issue that's happening right now.
On the other side of the world.
How do I respond to that.
I'm not physically there.
There's not that much I can do.
Well, you can pray.
And your prayers can enter.
Into that situation.
Even on the other side of the world.
Your prayers can go places.
That you cannot.
But it's also really important.
For us to pray.
Because as we begin to pray.
God begins to work on our heart.
And speak to us and guide us.
In the response that he has for us.
And there is no excuse for not praying.
We could say, well, you know.
I'm too far away to do this.
Or I'm too far away to do that.
Or I don't have time.
To get involved in this way.
Or I don't have the resources.

$^{641}$Or the money to get involved in that way.
Or I don't have the skills.
And the talents to get involved.
In this other way.
All of us can pray.
All of us can find time to pray.
To get our needs.
And to intercede on behalf.
Of someone else who is suffering.
Like I said, whether that's.
Your neighbor across the street.
Whether it's someone in your community.
Whether it's a family member.
Whether it's someone across the world.
Or someone in another state.
Or across the country.
There is no excuse for us not praying.
We have got to begin to pray.
That's something we're not very good at.
In the North American church.
Like we just don't pray very much.
We don't pray very long.
We don't take prayer very seriously.
And that's a huge problem.
We've got to begin to pray.
And when we see the suffering of others.
We have got to pray about it.
And the next thing that we really.
Need to do is to begin to listen to.
And to learn from.
Those who are suffering.
Look at verse 34.
The Samaritan goes.
It says he went to him.
And bound up his wounds.
Pouring on oil and wine.
Now this man had been brutally beaten.
He probably was not able to talk.
At this point.
So the Samaritan didn't.

$^{681}$Necessarily listen to him.
But he went and he learned.
Just through observation.
What the man's needs were.
He looked and he said.
OK this man is injured.
And he needs to be bandaged.
And he needs to have some.
Some medicine put on these wounds.
And that's what he began to do.
And that should be a part of our response.
When God interrupts the rhythms of our life.
With bringing suffering close to us.
We need to look at it.
And then we need to begin to pray.
And then we need to begin to listen.
And learn from those who are suffering.
And as we listen and learn.
God will begin to open our eyes.
To ways that we can actually.
Get involved and begin to help.
And so we need to take time.
And we need to sometimes be quiet.
And listen to the heart of someone.
Who is going through that experience.
And hear their story.
And hear the pain.
And then we can begin to respond.
And that's the next thing.
That we really need to do.
We need to respond in a tangible.
And gospel-centered way.
We need to respond in a tangible.
And gospel-centered way.
That's what we see the.
Samaritan doing in verse 34 and 35.
He goes, he binds up his wounds.
He puts the medicine on him.
He sets him on his animal.
Takes him to the end.

$^{721}$Where he can be carried for.
He pays for it himself.
He does something very tangible.
He doesn't just say, hey.
I'm sorry you got beat up.
And you're having a rough day.
I'm going to pray for you.
And I'll see you later.
And head on down the road.
He responds in a very tangible.
And loving way.
And that's important for us.
Now we need to understand.
That prayer is a tangible response.
Like I said earlier.
Sometimes it's more powerful.
Than anything we can do.
And sometimes in our culture.
We marginalize prayer.
And we act like, well.
You're only praying.
And that's ridiculous.
Prayer is a very tangible response.
If we believe in the Bible.
And we believe in the power of God.
Then we should believe.
The prayer in and of itself.
Is a tangible response.
And I certainly believe that.
And we need to also understand.
That listening and learning from.
And hearing someone's story.
And affirming the pain.
That they've gone through.
Just by listening and being present.
That is a tangible response.
So we've already talked about.
Two good tangible responses.
But there's oftentimes.
Much more we can do.

$^{761}$Beyond those things.
And much more that we should do.
Beyond those things.
And as we pray.
And as we listen.
God begins to open our eyes.
To what those things can be.
We might start out saying.
I don't even know.
How I could help in this.
I don't even know.
This person's dealing with something.
That I've never experienced.
I don't even know how to help.
Or this is happening so far away.
From me.
I don't even know.
How to get involved.
But as we begin to pray.
And as we begin to listen.
God will open our eyes.
To ways that we can get involved.
And tangibly respond.
With the resources and the gifts.
That we already have.
Maybe we say.
I don't have a lot of.
Financial resources to fix this.
But I have a little bit of time.
And I can give some of my time.
Towards this.
Or maybe we say.
You know what.
This is happening.
On the other side of the world.
And I can't be there.
And do much about it.
But I can connect with.
Some people who are going through this.
Virtually over the internet.

$^{801}$And I can tell them.
How much I love them.
How much I'm praying for them.
How much I care.
And how my heart breaks.
For what they're going through.
God will begin to reveal to your heart.
Ways that you can get involved.
Tangibly.
And that's so important.
But it's also important.
That as you get involved.
You center those responses.
On the gospel.
That's really important.
Because as humans.
We're really good at reacting.
To each other.
And reacting to wrong things.
That other people are doing.
But we're not so good at responding.
To God's truth.
And sometimes we just swing.
Back and forth like a pendulum.
We see someone who's doing.
Something wrong over here.
And we're like that's terrible.
So then we swing over here.
In this other direction.
And we get really out of balance.
Over here.
Well that doesn't fix anything.
Because now we're in.
An out of balance position.
That harms people in a different way.
And that's just as unhealthy.
As being out of balance over here.
And we see so much of that.
In human history.
And sometimes when we do that.

$^{841}$We actually create responses.
That are worse than the problem.
That we were trying to fix.
In the first place.
I went to college in Hawaii.
And in Hawaii.
There's all these mongooses.
That are wild and everywhere.
There's all over the place.
Because they don't have.
A natural predator.
Well they were not originally.
On the island.
But when Europeans.
First started going to the island.
There were no snakes on the island.
And they took in some snakes.
At some point.
Then they had a problem.
These snakes had no predators.
They were multiplying quickly.
They didn't know how to get rid of them.
So they brought in the mongooses.
To get rid of the snakes.
Well now the mongooses are everywhere.
Because they have no natural predator.
And throughout the course.
Of human history.
We see that a lot.
We try to fix a problem.
And in the process.
We create another problem.
Because we're just reacting.
To the first problem.
We're not responding to God's truth.
God's word and God's truth.
Contain solutions for us.
That are actually.
The ultimate solution to the problem.
Not just a response.

$^{881}$To someone else's.
Out of balance and wrong behavior.
And so we need to make sure.
That we are gospel-centered.
In our responses.
They need to be tangible.
We need to wear our prayers.
We need to get into the street.
We can't just sit back.
And ignore someone's suffering.
But we need to be gospel-centered.
As we respond.
And then the final thing.
That's so important for us to do.
Is to follow through.
And we see that in verse 35.
The Samaritan didn't just take him.
And leave him.
He gave the money.
He paid for him.
And said, "Hey, I'm gonna come back.
And check on this guy.
And if you need anything else.
I'm gonna pay for that.
I'm gonna make sure this guy.
Is cared for".
And that's hugely important.
Because a lot of times.
We're not very good at following through.
Especially in the North American church.
And especially as young people.
You know, I'm a millennial.
I'm an old millennial.
But I am technically.
In the millennial generation.
And that's one of the problems.
That millennials.
And centennials.
That we're not real good at.
We're good at responding.

$^{921}$And, you know, really getting.
Focused on a problem.
And really talking about it.
And really coming up with.
A lot of things that are wrong.
That's something that we're great at.
We can do that all day, right?.
We can create a lot of noise.
On social media.
And we're really good at that.
What we're not always so good at.
Is the follow through.
We're not always so good at saying.
Hey, let's bring the gospel.
Into this situation.
And let's get involved.
And even when it's messy.
And even when it's difficult.
And even if it takes years.
We're going to continue to be here.
And to love and to serve.
And to care for people.
And to bring Jesus into this place.
And into this life.
And into this situation.
And into this relationship.
And into this brokenness.
That's what we need to be doing.
We've got to follow through.
In these things.
And so there's a lot more here.
I wish we had even more time.
To talk about all this stuff.
But we're running out of time for today.
I want us to begin to think.
About these things.
As we look at all these.
Broken situations in our world.
Don't just ignore them.
Don't just pass them by.

$^{961}$Let them interrupt your life rhythm.
Begin to care.
See the suffering.
Start to pray about it.
Pray seriously.
And pray intently.
And pray fervently.
And begin to listen to the people.
Who are suffering.
And learn from them.
And really care about what's going on.
And as you do that.
Let God reveal ways.
You can respond tangibly.
Make sure those ways.
Are focused on the gospel.
So that they'll be healthy.
And won't just cause you.
To swing around like a pendulum.
And then follow through.
Continue to respond.
Don't just say.
Hey, I'm going to be really into this.
For the next three months.
And I'm going to pray about it.
And I'm going to talk to somebody.
And listen to somebody about it.
And I'm going to do something.
No, continue to follow through.
Sadly, there is so much.
Brokenness in our world.
And like I said earlier.
We can't fix it all.
But God wants us to be.
The hands and feet of Jesus.
In this world.
To do what He's called us to do.
And for most of us as individuals.
It's not that much.
We can't do an enormous amount.

$^{1001}$On our own.
But here's the beautiful thing.
There's this thing called the church.
And it's global.
It is a group of people.
That spans the entire globe.
And if all of us.
Are willing to pray.
And to get involved.
And to respond.
And to be obedient.
To show God's love.
In our own communities.
And the places that we are at.
Then something really big.
Starts to happen around the world.
And that's what Jesus is calling us to do.
I'm going to leave us today.
With the words of Jesus here.
As He calls out this story.
And I think it's a beautiful.
And perfect challenge for us.
Jesus said.
Which of these three.
Do you think prove to be a neighbor.
To the man who fell among the robbers.
And the lawyer said.
The one who showed him mercy.
And Jesus said to him.
You go and do likewise.
Family our world is suffering.
We need to go and do likewise.
Let's pray.
God, I just pray that you would be a work in our hearts.
You would help us to care.
And to be sensitive.
To the suffering of the world around us.
Whether it's across the street.
Or across the world, God.
Don't let us ignore it.

$^{1041}$Don't let us just pass it by.
Yeah, it's easier not to know sometimes.
But that's not your heart.
Your heart is you care for those who are hurting.
And those who are oppressed.
And those who are suffering, God.
And I pray that.
That would be our heart as well.
God teach us to pray fervently.
And seriously and passionately.
About the suffering in our world, God.
Teach us to begin to listen.
And care for those who are going through it.
To begin to try to understand.
What they're experiencing.
And what they're facing.
And how we can be involved.
And God teach us to respond.
In tangible ways.
But ways that are healthy.
And centered on the gospels.
And ways that actually will help.
And not create more problems.
And God teach us to follow through.
Don't let us start off strong.
And then forget about it.
Let us continue to love.
And to serve and to care.
And to respond.
God make us your hands and feet.
Here on earth.
To care for those who are suffering.
With your love and with your grace.
And with your goodness.
Flowing through us, God.
We can't do it on our own.
It would be empty.
And totally would not.
Fix the problems that we see.
But God with you.

$^{1081}$Flowing through us as your people.
The world can be changed.
And lives can be healed.
And that's what we want to see, God.
Help us to do that.
In Jesus' name, Amen.
\newpage



\section{約伯記}
\label{sec:Gigyg3Yfrs0}
\textbf{逆要作門徒 (1) – 約伯記 (陳崇基牧師) 2022.01.02}
\newline
\newline
連結: \href{https://youtube.com/watch?v=Gigyg3Yfrs0}{\texttt{https://youtube.com/watch?v=Gigyg3Yfrs0}} ~~~~ 語音日期: 2024-04-27
\newline
\newline
\hyperref[sec:Hq5No0HLb30]{\small{< < < PREV SERMON < < <}}
~
\hyperref[sec:index_chronic]{\small{[返順時目]}}
~
\hyperref[sec:index_scriptual]{\small{[返順卷目]}}
~
\hyperref[sec:ANI7oTOQdWo]{\small{> > > NEXT SERMON > > >}}
\newline
\newline
$^{1}$各位好,在新的一年裡面,希望大家在主裡面平安.
當我們想到生活艱難的時候,或者疫情已經進入第三個年頭.
無論你住在哪一個地方,都是有些不受控制.
或者令到我們生活非常困苦,也看不到盡頭,不知道什麼時候完.
當我們遇到很多不幸的時候,我們很多時候會想起不同的經文.
會想起聖經裡面的一些故事,或者一些詩篇,一些經段.
其中一個想起的當然就是《約伯記》.
如果你是熟悉聖經的話,我都想考考你.
除了想起《約伯記》之外,當我們受苦的時候想起這本書.
你還會想起什麼章節,什麼段落呢?.
如果你想到的話,除了《約伯記》之外.
你可以在Live chat的留言告訴我們.
到底你想起什麼聖經的段落,或者一些人物的遭遇,或者一些聖經的故事.
今天我們就去看《約伯記》.
《約伯記》是整本聖經裡面其中一個的collection,其中一本書.
聖經裡面大概中間有一個叫做詩篇智慧的collection.
那裡有幾本書,包括詩篇,雅歌等等,傳道書.
《約伯記》就是其中一個核心的書卷.
講到一個人類普遍的問題,就是苦難的問題.
苦難有很多方面,大家知道.
譬如有個人的一方面,有情緒,精神上的掙扎.
或者我們遇到一些不順利的時候,逆境的時候.
我們會有一些uncertainty,有一些不確定性.
有疾病纏繞我們,或者一些殘疾.
有些問題在我們身上發生,這些算是一種personal suffering.
一些個人的受苦.
第二方面,可能就是一些天災,或者人禍.
一些natural disasters.
譬如不久之前,美國中部有龍捲風突襲一個地方.
人們無家可歸,或者在不同的地方,可能有一些天災.
會令到我們的生活突然間打團,突然期待的.
給我們一些苦難,或者在這裡的comment裡面提過的.
可能是傳道書,或者耶利米哀歌.
國家亡國的時候,這些天災人禍的災害臨到.
所以就會令到我們有很多的不安.
第三樣就是,我們都很熟悉的.
就是一些叫做systemic oppression.
或者一些系統性的邪惡,或者一些的壓力.
想想我們生活裡面,在某些壓迫的政權底下的人.
包括香港,包括現在香港變天的這種情況.

$^{41}$很多人好像窒息一樣,或者甚至很失望的要離開.
或者在一些地方,一些少數的民族.
包括教會,受到不同的迫害.
不同程度的壓力,都是suffering的其中一種.
suffering的時候就是一種universal experience.
所以為什麼我們很多時候想起約伯記.
因為約伯記在這個故事裡面的主要的人物.
Joe 約伯,就是吃了不少的苦頭.
第一章的第一節,這樣去描述這個人叫做約伯.
他說烏斯蒂有一個人名叫約伯.
這個人完全正直敬畏神,他遠離罪惡.
他生了有七男三女,家產有很多.
有很多牛羊,母奴等等,木被等等.
這個人在全東方的人當中,是至為尊貴的.
這一節就這樣去介紹他.
然後下面的一些細節我們不讀.
就說到約伯是一個無可指責的人.
說到他道德上,敬拜,祭禮上,他是做到足的.
所以介紹了這個人之後.
然後我們就會看到鏡頭轉到天上.
天上發生了什麼事呢?.
第六節,他說有一天神的眾子都來.
試立在耶和華的面前.
撒旦也來到他們當中.
耶和華問撒旦,你從哪裡來?.
撒旦回答耶和華,我在地上徘徊,走來走去.
耶和華問撒旦,你有沒有注意我的僕人約伯?.
世上沒有一個像他這樣正直敬畏神和遠離罪惡的.
第九節,撒旦回答耶和華.
他說,約伯敬畏神,難道是無緣無故的嗎?.
你不是在地上四周圍捉了尼巴來保護約伯.
他的家庭,他所有的一切.
以及他所作的一切你都賜福給他.
他的家產地上天爭.
只要你伸手擊打他所有的一切.
他就會當面來褻瀆你.
耶和華對撒旦說.
好吧,他所有的一切都交在你手中.
只是不可以伸手害他.
於是撒旦從耶和華面前退去.

$^{81}$之後就說到一些災難.
一浪接一浪的來打擊約伯.
二女相繼離世.
最後他自己就得了重病.
在這個故事裡面就產生了很多的張力.
我們首先問到苦難的問題的時候.
很多時候我們就會問.
為什麼會有苦難發生?.
苦難發生的原因是什麼?.
苦罪源由是什麼?.
是不是這段故事就告訴我們.
我們是責怪或問責於哪一個起事.
哪一個搞這些事情.
是撒旦.
但是如果你繼續問.
張力仍然存在.
因為這裡是上帝容許撒旦去做這件事.
所以我們仍然解決不了.
到底為什麼一個義人會受苦.
約伯為什麼會無緣無故受苦.
難道就是上帝和撒旦好像吵架.
想爭拗誰贏.
於是就找約伯來糟蹋嗎?.
是不是就全部都怪責Satan?.
如果全部都怪責Satan.
我們是不是就能夠解釋到.
這個受苦的問題?.
有一樣東西我們要慢慢看這段聖經.
因為現代的翻譯本.
無論中英文.
我們翻譯了Satan這個字做撒旦.
但如果你查考一下.
這些不需要很高深.
要懂得讀希伯來文.
或者要講到希伯來文才知道.
你查考一下撒旦背後那個字.
其實是Hasatan.
Hasatan是什麼意思呢?.
在希伯來文.
Has.

$^{121}$其實是Article.
Definite Article.
就是一個貫詞.
英文的de.
所以你要直譯的話.
是de Satan.
所以我們講到Satan的時候.
其實是一個Hebrew word.
是一個希伯來文的字.
你一直都不知道.
原來你懂一點希伯來文.
Satan是什麼意思呢?.
Satan是對手.
或者是挑戰者.
Challenger的意思.
所以你要直譯的話.
不應該譯做撒旦.
不是我們今時今日新約裡面.
特別啟示的魔鬼頭號人物.
撒旦.
因為直至舊約與新約之間.
就是兩約之間的時候.
所謂撒旦才成為一個人名.
Personal name.
成為一個實體.
之前在舊約.
包括約伯記.
所提到的Satan.
其實不是指新約裡面所想的撒旦.
所以這裡只不過是說.
The challenger.
就是挑戰者在神的眾子中間.
他可以是善者.
可以是惡者.
也可以是忠誠.
所以我們看到.
不是一個邪惡的天使出現.
而是一個挑戰者.
一個challenger.
他去問神一個問題.

$^{161}$問神.
如果你將好的東西都給約伯.
他當然就會敬畏你.
他會敬拜你.
如果他不再被你去祝福的話.
他就一定會褻瀆你.
但如果你不斷給他好的東西.
約伯就一定會不斷搞錯他的動機.
最終約伯只是為了在你身上去拿到好的利益.
有好處.
他才會敬拜你.
所以如果你這樣想的時候.
這個挑戰者所問的問題.
他說約伯敬畏神.
難道是無緣無故的嗎.
這個問題其實沒有什麼不合理的.
很legitimate的question.
你有小孩你都知道.
大家明白.
一個小孩有好的表現.
如果純粹是為了拿你的獎賞.
拿你的好處的話.
他的動機就會成為一個問題.
他只是為了好處.
為了自己的利益而去聽話.
所以這個挑戰者所問的問題.
問到約伯去敬畏神.
那個動機的時候.
本質上其實是正確的.
是合法的問題.
因為他在質問上帝.
去管理運作這個宇宙的原則.
即是上帝對這個世界的policy.
他的政策到底是不是正確.
所以經文繼續講下去的時候.
上帝就很欣然的去接受這個挑戰者的問題.
然後大家知道約伯在他身上有很多的苦難.
然後幾十章的聖經的內容.
都是講什麼.
都是講這個苦難的問題.

$^{201}$如果你看看約伯記之後這麼多章的經文.
大概是講什麼呢.
敘曲完了之後.
4至到第27章.
有三大段的debate cycle.
就是約伯和他三個朋友之間的debate.
然後到第28章中間有一個智慧的思過插進去.
然後29至40一章.
就有另外一段很長的講話.
其中一個講者就是約伯.
然後有一個人彈了出來叫做以利戶.
這個可能是代表.
因為他的名字是一個Israelite name.
一個以色列的名字.
可能是代表以色列的school of philosophy.
就是以色列人怎樣去明白苦難的問題.
跟這些近東的哲學家一起去debate.
但到最後每個人講苦難的問題.
都是不完備的.
直至到耶和華自己出現.
然後就結語了.
你看到整本書大部分都是約伯和他幾個朋友去辯論.
環繞著什麼問題.
就是二人為何受苦.
三個朋友堅持約伯受苦是因為他做錯了事.
他犯了罪.
所以在他們debate的當中.
很長的debate很細節的debate.
就是他們說如果約伯願意去認罪的話.
雖然你不想起什麼罪.
但儘管你來到隨便認一樣.
只要你認罪的話.
你就會平安無事.
苦難就會結束.
神就會給你一些好的東西.
如果是這樣就很簡單很容易.
約伯只要認錯認罪認自己壞人.
這樣就搞定了.
辯論就很快會結束.
大家知道辯論一直都不結束.

$^{241}$為什麼呢.
因為不是這麼簡單.
因為約伯一直在debate裡面.
他都是堅持自己無辜無罪.
他不應該受苦.
他認為他是一個二人.
他不應該接受這樣的懲罰.
所以辯論就不斷延長.
一直去問這個問題.
漫長的debate裡面.
就去問這個.
為什麼約伯會受苦呢.
這些debate的對話.
裡面所表達的就是當時的世界.
怎麼想這些苦在問題.
甚至其實我們今天現代.
我們的世界觀.
我們的人生觀.
都是會想起這些東西.
為什麼你會受苦呢.
因為在以前的世代.
有一種叫做報應的原則.
Retribution principle.
人很相信他們和神明之間的關係.
就是一種報應的關係.
可以說是好也好報.
惡也惡報.
我對神明好一點.
神明就會給我好處.
如果我做錯事得罪了神的話.
神就會生氣.
然後祂就會降災給我.
然後我就會有不好的事情發生在我身上.
做好事就有獎.
做不好的事就壞事.
你就會受苦.
這些就是retribution principle.
這些很多時候在生活裡面.
這種運作的原則是有效的.
我們說種瓜得瓜種豆得豆.

$^{281}$殺種的收割的.
是平衡的.
是一個equilibrium.
但是大家知道在人生的路途上.
這個原則很多時候是不可行的.
特別我們現在在疫情當中.
你就發現這個所謂的好也好報.
惡也惡報.
很多時候都不是很可行的.
因為明明不是你做錯事.
但是你又受苦.
約伯就在這個情況裡面.
他堅持他自己的正直.
他就要求和上帝面對面.
和上帝來討論.
interview上帝.
就是想問為什麼我作為一個議員會受苦.
所以約伯記了三十多章的經文.
一路的辯論.
辯論最後的時候.
上帝有沒有出現.
上帝真的有出現.
第三十八章第一節.
上帝出現的時候.
他就不是回答約伯的問題.
他不是去解釋苦難.
他是問一輪的問題.
三十八章第一節.
那時耶和華在旋風當中回答約伯.
他說什麼呢.
第四節.
他說我奠定大地根基的時候.
你在哪裡呢.
你若有聰明.
你若有智慧.
就只管說吧.
所以他問了一大輪的問題.
世界創造的時候.
世界有的時候.
到底你在哪裡呢.

$^{321}$當然答案.
這些問題已經牽引了答案.
就是不知道.
這些是過於我的智慧.
你繼續看下去.
到底二十五節.
耶和華就問.
誰為大禹分水道.
誰為淚澱定道路.
使雨降在無人之地.
落在無人居住的曠野.
使荒廢之地都得到滋潤.
使青草在乾旱之地就得以長起來.
誰去造成這些.
換句話說.
這些問題背後想說什麼呢.
就是上帝管理這個宇宙的政策.
是不是真的好有好報.
是不是我們心裡面想的.
人很有限的.
所謂的正義.
做了一件事.
那件事應該有一個報應.
有一個回應.
二十五節至二十七節.
這裡問一個問題.
就是上帝降雨的原則是什麼.
上帝去控制和運作整個世界的時候.
他的operation policy是什麼.
問題就是問.
你知道有時候雨水是降在無人之地.
根本就沒有人去喝那些水.
沒有人會得益的.
那為什麼雨水會降在那地方.
根本沒有人在那裡能夠得到benefit.
怎麼會在荒地裡面.
這些沙漠的地方.
礦業的地方.
沒有人會得到好處.
到底上帝管理萬事.

$^{361}$運作萬有的原則.
是不是我們想的那麼簡單.
我們有限的智慧.
我們覺得好有好報.
惡有惡報.
但不是.
是超越我們所想到的.
繼續問下去.
這一段我覺得很有趣.
十三節.
鴕鳥的翅膀因而鼓動.
但牠的羽毛哪有慈愛.
鴕鳥生蛋.
蛋就留在地裡面.
埋藏在地裡面.
使牠們在土裡能夠得到溫暖.
但那隻鴕鳥那麼Army.
忘記了.
牠放蛋在地上.
又不能放得深.
有人踩過.
甚至自己踩過.
就會踩爛.
然後那個鳥就死了.
野地的走獸又經過的時候.
就把牠們淡走.
踩淡這些蛋.
怎麼會做這麼笨的事.
鴕鳥.
你看牠的樣子.
很笨.
但上帝要創造牠.
上帝的智慧在哪裡.
上帝如果真的創造鴕鳥的時候.
上帝如果這麼有智慧的時候.
怎麼會做這麼笨的事.
根本沒有智慧.
沒有wisdom.
神繼續問約伯的問題.
就是說.

$^{401}$我做了這些事.
你明不明白背後有什麼智慧.
這個下輩子.
我會做什麼.
你明不明白背後有什麼智慧.
這個是你不能夠參透.
真的不明白.
繼續問他.
問到最後的時候.
約伯就只能夠回答.
40章第3節.
約伯就回答耶和華.
我是微小的.
可以回答你甚愛.
我只好用手掩口.
遮住我的口.
說了一次不再回答.
說了二次就不再說.
你見到.
問了這些所謂的rhetorical questions.
問了.
知道答案的問題.
反問的問題.
就是是不是去到.
解釋苦難的恩友.
今天你很運濟.
最近你很多的錯字.
很多的困境.
why.
如果你是帶著這樣的預期.
你希望能夠.
解釋這些苦難.
解釋你現在的遭遇.
你就去看約伯記的話.
找不到答案.
你一定會很失望.
很失望.
因為整本書的構造.
整本書所寫的目的.
根本就不是為我們.

$^{441}$那個苦最問題.
來提供任何.
我們認為合理的答案.
整本約伯記.
不是教我們.
how to think about suffering.
怎麼去思考苦難問題.
整本約伯記.
是教我們.
how to think about God.
in our suffering.
在苦難當中.
我們怎麼去理解.
思考神.
怎麼去.
明白神.
所以約伯記不是解釋苦難的問題.
約伯記其實是擴展我們的神學.
我們的theology.
對上帝的認識.
所以到最後那幾章的時候.
我們發現這個上帝是.
awesome.
beyond我們的understanding.
beyond我們的wisdom.
超越我們人的理性.
上帝.
運轉整個宇宙的policy.
他的智慧.
不是像我們心想的.
我一定是一個人.
好事就一定有好的報酬.
壞事就立刻有壞的報酬.
上帝運作整個宇宙.
cosmos的wisdom.
是超越我們的理性.
神的智慧是高於我們的智慧.
我記得以前讀神學的時候.
有一個神學家叫John Frank.
他說了一句話很有意思.

$^{481}$那句話幾乎改變了我.
每一次敬拜的時候.
我都會想到.
我會想到.
改變了我每一次敬拜的時候.
遇見上帝的時候.
我怎樣去看上帝.
為什麼我會敬拜上帝.
John Frank說什麼呢.
他說神其中一個的本性.
一個的屬性.
我們會說慈愛.
忍耐等等.
我們平時都會說公義等等.
他說神其中一個屬性.
我們很多時候沒有留意的就是不能參透性.
incomprehensibility.
他說如果你想一下.
上帝是不能參透的話.
連上帝有多不能參透.
你都不能參透.
約伯記最後那幾章.
就是告訴我們我們是有限的.
上帝卻是無限的.
上帝和祂的智慧.
是要我們去信靠.
去信任.
雖然現在未必有答案.
但是我們要去信任.
上帝運作整個宇宙.
是有祂自己的智慧.
這個智慧未必是我們這麼小的頭腦.
現在我們就能夠去明白.
去comprehend.
但是我們要去相信上帝.
相信上帝.
最終整般約伯記.
不是為了苦難.
來提供一些安慰給大家.
如果我們有這樣的預期的時候.

$^{521}$或者想解釋苦難.
在苦難裡面有一點點安慰的話.
似乎約伯記不是寫來有這樣的目的.
當然你可以說約伯記受苦比我多.
所以我就覺得我都不是這麼慘.
主要目的不是這樣.
詩篇會做這件事.
愛歌,祈禱.
可以幫助我們舒緩我們的苦難.
或者想在苦難裡面有一點點安慰.
你可以看詩篇,詩歌.
但絕對不是看約伯記.
或者如果在苦難裡面.
你想有一點安慰的話.
可能你聽一下一些叱吒樂壇的冠軍歌曲.
Dear My Friend.
給你一點安慰.
但約伯記.
真正目的.
不是要做這件事.
因為講到苦罪問題的時候.
都沒有講到神異論的問題.
沒有為上帝去辯護.
沒有去辯護.
為什麼會有苦難呢.
做一些哲學的分析等等.
帶著這樣的預期去看這本書.
我們就會失望.
但應該帶著什麼樣的預期.
去看這本書約伯記.
不是去教我們.
怎麼去想著苦難.
而是教我們.
怎麼去想著守護.
在你的苦難中.
在你現在困苦的當中.
怎麼去想著神.
去認識神.
從前我風雲有你.
現在我親眼見你.

$^{561}$你有沒有想過.
我們上帝的認識.
就是神學.
其實就是決定我們的生活.
我們的ethics.
我們的倫理.
我們應做什麼.
會做什麼.
我們的theology是determine.
我們的ethics.
神學是決定我們的倫理生活.
我們回到第一章.
第一章的挑戰者.
問一句什麼問題.
挑戰者.
問約伯敬畏神.
難道是無緣無故的嗎.
難道是不求益處的嗎.
這個問題是合法的.
你敬拜神難道是無緣無故的嗎.
還是你敬拜神是為了什麼.
為了自身的利益.
是有interest的.
所以問的問題就是.
在世間上.
和上帝的關係.
到底是不是有一種不求利益的義.
不求利益的敬畏.
一種disinterest righteousness.
一種disinterest的faith.
去敬拜上帝.
是不是因為他給我好東西.
所以我才去敬拜他.
如果他不給我好東西.
我就不敬拜他.
約伯敬畏神.
我就不敬拜他.
約伯是不是在順境的時候.
才為上帝獻上奉獻呢.
是不是真的無私.

$^{601}$還是順也好逆也好.
約伯都去忠心的敬畏神.
這個才是信仰真正的本質.
如果我和上帝的關係.
只是為了好處的話.
只是為了祂給我一些benefit.
prosperity 祝福.
給我順利的時候.
我才去敬拜祂的話.
這個敬拜要打一個很大的問號.
或者新的一年.
我們真正要做的事.
不是要訂立這些目的這麼簡單.
目標這麼簡單.
而是真的要問這個最根底的問題.
跟隨基督.
難道是無緣無故的嗎.
還是為了我自身的好處.
我的順利.
我得到祝福 我得到幸福.
所以我跟隨.
到沒有幸福的時候 我跟不跟.
如果上帝真是上帝的話.
他值得我們敬拜的話.
即使沒有利益.
我們都願意去敬拜.
祂才是我們真正的上帝.
當然我們在基督裡面.
有很多的所謂benefit.
你想想.
在基督裡面我們信耶穌.
有沒有benefit.
很多 有永生.
有屬靈的祝福.
有禱告的權柄.
天國公民的尊貴身份等等.
我們列出很多最得盛面.
但請你記住.
這些所謂的benefit.
不是我們賺回來的.

$^{641}$是我們本身不配.
是基督伯伯在十下賜給我們.
是祝福.
這個白記.
不是去解釋苦難的作品.
而是一本去教我們.
在苦難當中.
在逆境裡面.
我們怎樣去作門徒的書.
所以不是我們自己.
去作門徒的書.
所以不是一本關於苦難的書.
而是一本關於徒弟的書.
是呼召我們在艱難的時候.
到底我們的信仰是怎樣一回事.
是否繼續跟隨基督.
開始新的一年.
這個可能真是一個最重要的問題.
我們需要去confront(對抗)我們自己.
我們敬拜上帝.
我們去跟隨上帝.
到底是為了什麼.
是為了我們自己有benefit(益處).
還是上帝是值得我們敬拜.
值得我們榮耀.
我們盡心盡力去跟隨他.
在我昨天見到一個圖文.
一個英文的圖文.
我希望這個新年的祈禱.
能夠成為大家的禱告.
新年的禱告.
這裡說.
願神賜給你快樂的一年.
並非保障你免受悲苦和傷痛.
而是賜你力量承載艱難.
不是讓你平步青雲.
而是使你堅定地走每一步.
不是要去除去困擾.
而是消除你內心的恐懼.
不是給你無間斷的陽光.

$^{681}$而是即使黑雲蓋天.
你念想仍然帶著屬天的明亮.
不是叫你生活天天順意.
而是叫你看到他人有需要的時候.
你會樂意地去服事.
願上帝在來年給你慈愛.
平安 盼望和喜樂.
這是我的禱告.
我們同心低頭祈禱.
讓我們順境.
讓我們逆境.
教導我們在逆境當中.
在困擾的裡面.
我們如何思想你.
體會你,認識你.
真真實實地明白.
你是那位超越我們一切的上帝.
配得我們五體投地來敬拜.
神啊,我們有著你的恩典.
縱使在今年裡面.
我們會遇見疫情.
遇見很多困擾.
主叫我們所求所想.
不是單單我們可以平步青雲.
除去困擾.
無災無病.
而是在這一切當中.
我們能夠堅定地,有力地.
向前每一步.
去信任你,信靠你.
因為你是我們生命當中.
宇宙當中的主宰.
主,接下來這一年.
無論我們會否遇見任何挑戰.
求上帝給我們盼望,平安的力量.
叫我們在身邊.
見到有需要的人的時候.
我們樂意伸出援手.
見到我們不開心,低谷的時候.
讓我們有氣力去承載.

$^{721}$堅定我們每一個.
讓我們知道你插不透的智慧.
最終會帶來公義.
最終會帶來榮耀.
求主你幫助我們.
忠心地敬拜你.
祈禱.
奉主耶穌基督的名字.
\newpage



\section{8:18-30}
\label{sec:ANI7oTOQdWo}
\textbf{Hope in Heaven for Suffering on Earth (Romans 8\_18-30) - George Crawley}
\newline
\newline
連結: \href{https://youtube.com/watch?v=ANI7oTOQdWo}{\texttt{https://youtube.com/watch?v=ANI7oTOQdWo}} ~~~~ 語音日期: 2024-04-30
\newline
\newline
\hyperref[sec:Gigyg3Yfrs0]{\small{< < < PREV SERMON < < <}}
~
\hyperref[sec:index_chronic]{\small{[返順時目]}}
~
\hyperref[sec:index_scriptual]{\small{[返順卷目]}}
~
\hyperref[sec:J6PHRTa1EUI]{\small{> > > NEXT SERMON > > >}}
\newline
\newline
$^{1}$Good morning..
Is that on?.
Yeah..
All right..
Awesome..
How y'all doing today?.
Good?.
Yeah..
Thanks for having me..
It's always a blessing to be here and a real great opportunity..
So I'm thankful to you guys..
Our text this morning is coming from Romans chapter 8, 18 through 30..
And I love that a lot of the songs we sang this morning were centered around the glory.
of God, the glory of Jesus..
Because as we're going to get to a lot of the things that we're going to be talking about,.
we're going to be talking about the glory of God..
And I love that a lot of the songs we sang this morning were centered around the glory.
of God, the glory of Jesus..
Because as we're going to get to a lot of the things that we're going to be talking about,.
we're going to be talking about the glory of God..
And I love that a lot of the songs we sang this morning were centered around the glory.
of God, the glory of Jesus..
And I love that a lot of the songs we sang this morning were centered around the glory.
of God, the glory of Jesus..
And I love that a lot of the songs we sang this morning were centered around the glory.
of God, the glory of Jesus..
And I love that a lot of the songs we sang this morning were centered around the glory.
of God, the glory of Jesus..
And I love that a lot of the songs we sang this morning were centered around the glory.
of God, the glory of Jesus..
And I love that a lot of the songs we sang this morning were centered around the glory.
of God, the glory of Jesus..
And I love that a lot of the songs we sang this morning were centered around the glory.
of God, the glory of Jesus..
And I love that a lot of the songs we sang this morning were centered around the glory.
of God, the glory of Jesus..
And I love that a lot of the songs we sang this morning were centered around the glory.
of God, the glory of Jesus..
And I love that a lot of the songs we sang this morning were centered around the glory.
of God, the glory of Jesus..

$^{41}$And I love that a lot of the songs we sang this morning were centered around the glory.
of God, the glory of Jesus..
And I love that a lot of the songs we sang this morning were centered around the glory.
of God, the glory of Jesus..
And I love that a lot of the songs we sang this morning were centered around the glory.
of God, the glory of Jesus..
And I love that a lot of the songs we sang this morning were centered around the glory.
of God, the glory of Jesus..
And I love that a lot of the songs we sang this morning were centered around the glory.
of God, the glory of Jesus..
And I love that a lot of the songs we sang this morning were centered around the glory.
of God, the glory of Jesus..
And I love that a lot of the songs we sang this morning were centered around the glory.
of God, the glory of Jesus..
And I love that a lot of the songs we sang this morning were centered around the glory.
of God, the glory of Jesus..
And I love that a lot of the songs we sang this morning were centered around the glory.
of God, the glory of Jesus..
And I love that a lot of the songs we sang this morning were centered around the glory.
of God, the glory of Jesus..
And I love that a lot of the songs we sang this morning were centered around the glory.
of God, the glory of Jesus..
And I love that a lot of the songs we sang this morning were centered around the glory.
of God, the glory of Jesus..
And I love that a lot of the songs we sang this morning were centered around the glory.
of God, the glory of Jesus..
And I love that a lot of the songs we sang this morning were centered around the glory.
of God, the glory of Jesus..
And I love that a lot of the songs we sang this morning were centered around the glory.
of God, the glory of Jesus..
And I love that a lot of the songs we sang this morning were centered around the glory.
of God, the glory of Jesus..
And I love that a lot of the songs we sang this morning were centered around the glory.
of God, the glory of Jesus..
And I love that a lot of the songs we sang this morning were centered around the glory.
of God, the glory of Jesus..
And I love that a lot of the songs we sang this morning were centered around the glory.
of God, the glory of Jesus..
And I love that a lot of the songs we sang this morning were centered around the glory.
of God, the glory of Jesus..

$^{81}$And I love that a lot of the songs we sang this morning were centered around the glory.
of God, the glory of Jesus..
And I love that a lot of the songs we sang this morning were centered around the glory.
of God, the glory of Jesus..
And I love that a lot of the songs we sang this morning were centered around the glory.
of God, the glory of Jesus..
And I love that a lot of the songs we sang this morning were centered around the glory.
of God, the glory of Jesus..
And I love that a lot of the songs we sang this morning were centered around the glory.
of God, the glory of Jesus..
And I love that a lot of the songs we sang this morning were centered around the glory.
of God, the glory of Jesus..
And I love that a lot of the songs we sang this morning were centered around the glory.
of God, the glory of Jesus..
And I love that a lot of the songs we sang this morning were centered around the glory.
of God, the glory of Jesus..
And I love that a lot of the songs we sang this morning were centered around the glory.
of God, the glory of Jesus..
And I love that a lot of the songs we sang this morning were centered around the glory.
of God, the glory of Jesus..
And I love that a lot of the songs we sang this morning were centered around the glory.
of God, the glory of Jesus..
And I love that a lot of the songs we sang this morning were centered around the glory.
of God, the glory of Jesus..
And I love that a lot of the songs we sang this morning were centered around the glory.
of God, the glory of Jesus..
And I love that a lot of the songs we sang this morning were centered around the glory.
of God, the glory of Jesus..
And I love that a lot of the songs we sang this morning were centered around the glory.
of God, the glory of Jesus..
And I love that a lot of the songs we sang this morning were centered around the glory.
of God, the glory of Jesus..
And I love that a lot of the songs we sang this morning were centered around the glory.
of God, the glory of Jesus..
And I love that a lot of the songs we sang this morning were centered around the glory.
of God, the glory of Jesus..
And I love that a lot of the songs we sang this morning were centered around the glory.
of God, the glory of Jesus..
And I love that a lot of the songs we sang this morning were centered around the glory.
of God, the glory of Jesus..

$^{121}$And I love that a lot of the songs we sang this morning were centered around the glory.
of God, the glory of Jesus..
And I love that a lot of the songs we sang this morning were centered around the glory.
of God, the glory of Jesus..
And I love that a lot of the songs we sang this morning were centered around the glory.
of God, the glory of Jesus..
And I love that a lot of the songs we sang this morning were centered around the glory.
of God, the glory of Jesus..
And I love that a lot of the songs we sang this morning were centered around the glory.
of God, the glory of Jesus..
And I love that a lot of the songs we sang this morning were centered around the glory.
of God, the glory of Jesus..
And I love that a lot of the songs we sang this morning were centered around the glory.
of God, the glory of Jesus..
Most minds of fiction on the planet can't even begin to comprehend or understand how.
great things are going to be, how glorious our future is, our hope in heaven is, where.
our bodies and minds are perfected, the heavens and the earth perfect, perfect justice, perfect.
peace, perfect love, utter perfection..
And according to Paul, all of this is utterly glorious..
And therefore, our present sufferings are nothing but a mere drop in an infinite ocean.
of glory and perfection..
And so that right there is why he can endure suffering and why he can say it's only temporary.
and it's only light, because he knows something is coming that is so much greater..
And then Paul continues..
He doesn't stop there, although he could, but he continues to explain how great this.
future glory is that allows us to persevere through suffering..
The next two verses, Paul says this..
"For the creation waits with eager expectation for the children of God to be revealed..
For the creation was subjected to frustration, not by its own choice, but by the will of.
the one who subjected it in hope.".
So, Paul continues by painting this picture, he gives us this analogy that all of creation,.
the sun, the moon, the stars, the mountains, the hills, the oceans, everything is waiting.
for this coming moment in anticipation..
And the reason that it's waiting is because it has been subjected to frustration, it says..
And the reason that all of creation, as good as it is, has been subjected to frustration,.
we see all the way back in the beginning, in the opening chapters of Genesis, when God.
subjected creation..
You see that when the fall of man happened, when sin entered the world through Adam and.
what him and Eve did in the garden, God not only cursed Adam and Eve, but he also cursed.
the ground..

$^{161}$He cursed creation..
And the reality that creation has been subjected, it's evident everywhere you look in a lot.
of different ways..
Like I said before, our bodies are always breaking down, but the earth itself, all of.
creation is also breaking down and things are not right, they're not as they should.
be..
And so you can think in creation, what examples of this do we see in nature?.
Earthquakes and tsunamis that destroy so much..
Forest fires and droughts where you have whole nations burning for days, weeks on end..
Those things, we call them natural disasters, but they're not really natural..
That's not the original way God created the earth..
That came when sin came, when God subjected the world..
You can also see this, if any of you guys in here are like into physics, maybe you're.
a little bit nerdy like me, you know about the second law of thermodynamics?.
Anybody know what I'm talking about?.
Okay, okay, I'll break it down for you guys, I'll break it down for you guys, okay?.
So the second law of thermodynamics, simply put is that everything is winding down..
Everything is going from a state of order to disorder..
So things are always breaking down..
As soon as you buy a car, you drive it off the lot, immediately loses its value..
You buy a house, over time it breaks down, the pipes, they're leaky, the roof needs.
to be replaced, everything is always breaking down..
And this is true even in our own lives, I'll give you guys a little bit of an example.
for myself, I'll be honest..
As the years go by and I look back at pictures of myself from college or from high school,.
I'll always ask my wife, Kiamia, just be honest with me, Kiamia, you know, look, is.
my hairline falling back, okay?.
I have to wonder, because everything is winding down, eventually all these beautiful locks.
will be gone..
My dad is bald, my grandpa's bald, it's coming for me one day..
That's just the reality, and it's goofy, it's kind of silly, but it's true..
Everything in creation is winding down..
But even in this frustration, there is a hope..
Even for creation, there is a hope..
And so what is that hope?.
Paul explains in the next verse..
That the creation itself will be liberated from its bondage to decay and brought into.
the freedom and glory of the children of God..
And so this hope, even for creation, is the emancipation of all creation from the bondage.
to corruption..

$^{201}$That there is a certain glory that all of creation, like the children of God, will freely.
receive when God makes all things new..
And it is this beautiful new creation that you can see in Revelation chapter 21, chapter.
22, when the universe will no longer trend in the direction of death, disease, and decay..
There will be no more disorder..
All things will be restored to perfect order..
And this future is so great that Paul gives us this picture of all creation waiting, anticipating.
to see what is brought about when that glorious future comes to fruition..
And again, Paul just continues building upon this picture, building this metaphor to underscore.
the greatness of the glory to come..
And I want to explain here the reason that Paul is doing this..
He's explaining how great it is and the glory to come and that creation is growing..
And because this is all building into that very first point of how we can deal with our.
suffering is by looking to a future hope..
So I want you guys to kind of get a sense of this weight of the future glory, how awesome.
it is..
Paul really, really goes to great lengths to get his readers to understand how great.
it's going to be because that will help us deal with suffering..
And so he says in verse 22 and 23, "We know that the whole creation has been groaning as.
in the pains of childbirth right up to the present time..
Not only so, but we ourselves who have the first fruits of the Spirit grown inwardly.
as we wait eagerly for the adoption to sonship, the redemption of our bodies.".
So Paul gives us this, and I think this is like the perfect metaphor for what he's teaching.
here about suffering, is the metaphor of childbirth..
And as I mentioned before, my wife, 38 weeks pregnant, we've been doing these childbirth.
education classes once a week..
And we go to the class, and we have to learn how to prepare, what labor is like, how to.
get ready, make sure your bag is packed, make sure you have the car seat in the car..
And one of the things that I've had to watch is childbirth videos..
And, okay, if there's any moms, fathers in here, even if you're not a mom or a father,.
you know a mom, okay?.
And so there's just something about childbirth that if I were to describe it, I would say.
that it's not a quiet situation..
It's not the most peaceful, serene, quiet environment..
It's kind of a noisy thing..
And a lot of these videos I've seen, these women when they're in labor, they're just.
groaning is the best word you can say..
They're not really saying words..
They're just groaning in anticipation..
And I was thinking that maybe I would give you an example, but I'm not going to do that..

$^{241}$I'm not going to, you know, that's loud..
But you get the idea..
It's just this groaning and it's wordless, but they're just longing and waiting for this.
baby to come into the world..
And so we can learn from childbirth some beautiful truths about our own suffering..
And months and months of pain and sickness and suffering, it teaches us something that's.
true for everyone..
And so when we think about childbirth, we think about this beautiful way in which God.
has designed us to bring life into the world..
There's three things we can learn from this about our own suffering..
The first is that it is absolutely temporary..
So for all the moms out there, could you imagine living in a world where you did not know that.
nine months was like the limit?.
Imagine living in a world where you could be pregnant for nine months, but you could.
also be pregnant for nine years, and you just didn't know..
Okay?.
You just did not know when it was going to come to an end..
Imagine how hard that would be..
Okay?.
And as tough as those nine months are, you always have this thought in the back of your.
head, "Okay, 40 weeks..
Okay, if I can get to 40 weeks, okay, I can handle it..
I got two weeks left..
I got one week left..
I can do this.".
Right?.
And so the fact that the suffering is temporary, it gives you some peace as you bear it..
It gives you some hope and some encouragement..
And the second thing is that suffering is absolutely purposeful..
When a woman is pregnant, all that, the sickness and the aching and the pains, it's not pointless..
All of it is a part of God's beautiful process of bringing life into this world..
And so all of that pain is purposeful..
And so it just, it allows you to suffer it and bear it a little more because you know.
God is using all of this for a beautiful and glorious purpose..
And that none of my suffering is purposeless..
And the last thing, and I think the most beautiful thing, is that all that suffering is absolutely.
worth it..
I've talked to so many moms over the past few months, getting advice and asking them.
about their experience..
And the one thing they'll always tell me, well, two things..

$^{281}$The first thing, they'll always agree that it is difficult, that it is not easy..
But the second thing they always say is that that moment when they first saw their little.
baby girl, or they first held their baby boy to their chest, all the pain, all the suffering,.
it just melts away..
It becomes one of the happiest moments of their life..
And the pain is replaced with pure joy, right?.
And so all of that suffering is absolutely worth it..
And so this is true in childbearing, but it's also true in creation..
And it's true for us that suffering in this fallen world is only temporary and it's not.
purposeless, but it will be absolutely worth it..
And so we can know confidently that this future day of glory is coming, not just for creation,.
but for us also..
And we can know that, Paul says, because God has given us a guarantee..
He has given us a down payment, and that's the Holy Spirit..
He says that the Holy Spirit is the first fruits..
And this is a concept we get from Scripture where when the produce came, when it was time.
for the harvest, the Israelites would give a first fruit offering to God, knowing that.
He would make sure that their crops came in again during the next harvest and they had.
enough food to eat..
And so that first fruit offering, it was like this trusting in God, this guarantee of what.
is to come..
And so the Holy Spirit serves a similar purpose for us..
He is a guaranteed down payment..
He is a promised portion of a coming fullness..
And again, that fullness, that glory, that future hope allows us to persevere when we.
suffer in this life..
Paul continues, verse 24 and 25, "For in this hope we were saved, but hope that is seen.
is no hope at all..
Who hopes for what they already have?.
But if we hope for what we do not yet have, we wait for it patiently.".
So we're hoping, we're waiting for this future glory..
We may not have it yet..
We may not see it yet, but we know it's coming and we have the confident assurance that it.
will come..
And I want to make a point to explain that when you read in the Bible about hope, that.
biblical hope, true hope, as Paul is talking about here, it is not just like a regular.
sense of hope..
It's not like a kid who writes a Christmas letter and is hoping to get his favorite toy.
and he might get it, but he might not..
That's not the kind of hope we have..

$^{321}$We have a hope that is grounded in the promises of God, a God who doesn't lie, who doesn't.
change his mind..
And so we have a hope that is assured, that we know will come to pass..
It's a confident assurance that we have..
And it's because of that confident assurance of the future glory to come that we can eagerly.
and patiently wait through our sufferings..
So, all of that metaphor and analogy and word picture Paul gives us leads us to our first.
answer, our first answer to the question of suffering, how can we deal with suffering?.
It is by knowing and believing in this future hope to come, this glory that's unimaginable,.
that outweighs any suffering you can experience in this life..
That coming glory is so great we can bear under any suffering and difficult circumstances.
in this life..
But Paul continues, and we're going to see now a second reason how we can endure suffering.
in this life..
Verses 26 to 27 says, "In the same way, the Spirit helps us in our weakness..
We do not know what we ought to pray for, but the Spirit himself intercedes for us through.
wordless groans..
And he who searches our hearts knows the mind of the Spirit, because the Spirit intercedes.
for God's people in accordance with the will of God.".
And the beautiful thing about this is that we do not suffer alone..
A lot of people have this, a lot of people in the world who reject the idea of God..
They say things like, "I just can't believe in a God who allows suffering..
I can't believe in a God who just sits back and does nothing while people suffer.".
And the truth is, I would not believe in that kind of God either, because that is not the.
kind of God we have..
We have a God who, one, not only entered into the world to suffer as the world suffers,.
but we have a God who in the Holy Spirit is with us in our weakness, in the midst of our.
suffering, so that we don't have to be alone..
And even when we don't know why we suffer, even when we don't know quite what to pray.
for or how to respond to suffering, Paul says that the Holy Spirit knows what we need, and.
the Holy Spirit is going to God the Father, interceding on our behalf, because on one.
hand, the Holy Spirit who's inside of us knows exactly what we need..
On the other hand, the Holy Spirit knows the mind of God, and He knows the exact will of.
God, so we can be confident that even when we're in a moment of weakness and we don't.
know what to say to God, how to pray, what to do, the Holy Spirit who knows us and knows.
the Father's will goes to God on our behalf..
And those things that we need, those requests taken to God, in His will, will be done..
And so we have that confidence that God is with us, God is always working about His will,.
and that's exactly what we need..
A question that can arise that Paul's going to answer in this next verse, though, is what.

$^{361}$exactly is the will of God?.
If the Holy Spirit, when we don't know what to pray for, is interceding, working for us.
to achieve the will of God, what exactly is the will of God for us in the midst of our.
suffering?.
Paul says in verse 28, "And we know that in all things, God works for the good of those.
who love Him, who have been called according to His purpose.".
So God's will is our good..
God is working all things out for our good..
And that can be difficult to imagine..
Just think about it..
Do we truly believe that God is working all things for our good, that God is working cancers,.
criticism, rejection for our good?.
Just as much as He is working beautiful things, pleasurable, amazing things for our good,.
God is working all things for our good..
And therefore, even when we suffer, our suffering is not meaningless..
It is not purposeless, because God is even using suffering for our good..
And I think that's such a beautiful truth, that our suffering is not meaningless and.
is not purposeless, because take a second and think about a world in which there is.
no God..
And for many people on the planet, this is the world that they think they live in, a.
world with no God, no Creator, no ultimate meaning, no ultimate purpose, no right or.
wrong..
In that kind of a world, our suffering is totally meaningless..
In that kind of a world, we're just creation, we're just accidents, not even created..
That sometime in the past, there was just something that happened in the midst of nothingness.
and created the world, but there's no intended order or no design..
And so we ultimately have no purpose, we have no Creator..
And so when we suffer, we just suffer..
It's not good, it's not bad, it just is what it is..
That kind of life is depressing and it's sad, but that is not the reality we have, because.
we do know there is a God in heaven, and He's sovereign, and He rules and reigns over.
everything and uses all things for our good..
So even when we suffer, there's a meaning, there's a purpose behind it, and we can be.
comforted because we know that meaning and purpose is for our good..
One thing I do want to say about this verse, because it's a very popular verse, and it's.
easy to pull this verse out of its context, and it's easy to try to spin this verse into.
a way Paul did not intend..
So what Paul is not saying is that God is going to do whatever you think is good, that.
God is going to give you whatever you think is good, and that if you think money and riches.
and car and all success and every avenue of life is what's going to happen, that's not.

$^{401}$what Paul is saying, because Paul is somebody, again, he knew what it was like to suffer..
He knew what it was like to be rejected, to be despised, to be hated, to be hungry, but.
that's not because Paul was outside of God's will or disobedient to God..
He was obedient to God and still experienced those things..
The truth of what God really wants for us, and that's not to say that He doesn't ever.
want us to have nice things and have good things..
There are times when He does, but the ultimate good that God is working out, and what Paul.
is talking about here, is not exactly what the world may define good to be..
Paul is going to define what that ultimate good is..
In the very next verse, he says this, "For those God foreknew, He also predestined to.
be conformed to the image of His Son, that He might be the firstborn among many brothers.
and sisters..
And those He predestined, He also called..
Those He called, He also justified..
Those He justified, He also glorified.".
So you can see what Paul is saying here is that our ultimate good is to be conformed.
to Christ, is to become like Jesus Christ..
That is the ultimate good that God always desires in every situation, and God uses everything,.
all things, every situation to that end, to make us like Jesus Christ..
That is the ultimate goal of our lives, as followers of Jesus Christ, is not to get rich,.
is not to be the most popular, or the most well-liked and respected, but is to be like.
Jesus Christ..
No matter what circumstances we find ourselves in, we should long to be like Christ, to think.
like Christ, to love like Christ..
That is what God wants for our lives..
And we see that not only do these verses reveal what God is working all things for in our.
lives, but it also reminds us the reason that we can be confident of these things, that.
we can be absolutely assured that God's got us, is because in this verse, in this passage,.
we see what, for generations, for centuries, the church has celebrated, and it's often.
called the golden chain of salvation..
And it's the reason that we can be confident that God has us, and gives us assurance that.
he's holding us, he's keeping us safe, like an anchor who holds a boat down..
God has this, because Paul says, "Those whom God foreknew, he predestined to be conformed.
to the image of his Son..
And those he predestined, he called..
Those he called, he justified..
Those he justified, he glorified.".
And Paul is writing these verbs here, called, predestined, justified, glorified..
And even though that glory we know is a future glory, Paul's writing about these things as.
if they've already happened, as if it's a present reality that we already are glorified..

$^{441}$And what Paul is doing is telling these readers that if you have been predestined by God,.
foreknown by God, which doesn't mean that he just knew of you, which is true of everyone..
God knows everything, and therefore, from the beginning of time, God knew everyone..
But those who God foreknew, he loved in a special way..
He knew in a special, intimate way..
And those who he foreknew, he predestined..
Those he predestined, he called..
Those he called, he justified..
Those he justified, he glorified..
That's a confident thing, that if you have been foreknown by God, it's already a done deal..
You're glorified by God, as if it's already happened..
And so that is how we can be so confident of these things..
Even when we suffer, even when it seems like things are not going right, this right here.
gives us so much hope..
And so we see here that second reason that we can suffer..
One is because there's a future glory coming that is so great, no amount of suffering can.
even begin to compare..
And number two is that we are never alone in our suffering..
God's always with us..
And he has been before the world was even created..
And now the third reason we see here of why we can endure this suffering is because none.
of our suffering is meaningless..
It all has a reason..
It all has a purpose..
And so those are three things that are just so encouraging in the midst of our suffering..
God's working all things for our good, so our suffering, it is not meaningless..
There's a reason behind it that is good for us, that's great for us..
And that truth that many worldviews, many philosophies and religions don't have this.
hope in the midst of suffering that we have..
Those three things that can keep us grounded in the midst of our suffering..
And I think that there's a beautiful example of this all the way as far back as, again,.
the book of Genesis and the life of Joseph..
We see these truths at play when Joseph, loved by his father but hated by his brothers, was.
betrayed, thrown into a pit, sold into slavery, goes into the land of Egypt, and then he's.
harassed and falsely accused by his master's wife..
Then he's put in jail, and he's left in jail, and he helps these two guys who were formerly.
working in the house of Pharaoh..
He helps one of them get back into Pharaoh's good graces..
And Joseph tells this cupbearer of Pharaoh, "When you get back into Pharaoh's household,.
tell him about me so that I can get out.".

$^{481}$And when the man gets back into Pharaoh's household, he forgets about Joseph..
And Joseph spends even more time falsely accused, wrongly imprisoned..
All of this, through all this suffering, we see God is still with Joseph..
God never leaves Joseph..
God never lets Joseph suffer for no reason..
We start to see that God is working in the midst of all this because the time comes when.
Joseph is called up..
He gets his chance..
He goes before Pharaoh..
He interprets Pharaoh's dream favorably..
It comes to pass, and he gets elevated to the second most powerful position in Egypt,.
all because God was always working through his suffering..
And a famine comes across all the land, and because Joseph, through God, saw it coming,.
he had stocked up grain and food..
So when all the different peoples in the nations around Egypt have to come receive grain,.
his brothers, who set this whole thing in motion, throwing him in the pit, selling him into slavery,.
they come to him now to get help from him so they can live..
And they don't really realize it's him..
There's this kind of back and forth that occurs..
But in the very end of this narrative, we see that he is kind of restored to them..
He accepts them..
He reveals himself to them, to his father..
There's joy..
They're celebrating..
But when their father, Jacob, dies, the brothers get scared, because now they think that Joseph.
is going to get revenge..
There's no one to stop him..
But when they come to him in fear, trembling, he tells them not to fear..
And he gives this beautiful verse, I think sums up the story, the life of Joseph so perfectly,.
and it exemplifies what we're talking about in Romans chapter 8..
Joseph tells them, "You intended to harm me, but God intended it for good, to accomplish.
what is now being done, the saving of many lives.".
So we see that God is sovereign over everything..
And it's not that what his brothers intended for evil, God used for good, as if God had.
to create a plan B and scramble to figure out, "How can I turn this bad situation out for good?".
What he says is that, "You intended to harm me..
God intended it.".
So God was already sovereign over their decision, over their choices, using it for Joseph's good..
That's exactly what Paul is bringing out in Romans chapter 8..
That's the same thing God does for us..

$^{521}$He's always working, he's always using all things, suffering included, for our good..
So when we think about this beautiful truth, one, that God is using all things for our.
good, which means our suffering is never purposeless or meaningless..
Two, the Holy Spirit is working on our behalf, always interceding, so that when we suffer,.
we're never alone..
And three, when we suffer, we have on the other hand, this future glory that is so great,.
it makes all our suffering pale in comparison..
Those three truths give us a new perspective on how to handle suffering in this life..
And so when we suffer now, we don't need to ask the question, "Why me?.
Why am I suffering?.
How could I ever get through this?".
We should ask ourselves, "What can I do now in my suffering with these beautiful truths,.
with what God has revealed?.
How can I apply them to my life to persevere through the suffering, or to encourage someone.
close to me who is also suffering?".
So I want to leave you guys with three final questions and three final exhortations to.
take this word to heart, that we could all use it, be encouraged by it, and apply it.
in our lives to the glory of God..
So, brothers and sisters, when you suffer, are you tempted to give up hope?.
Are you tempted to lose faith?.
Remember that your suffering is temporary..
It will not last forever..
It cannot last forever..
And not only that, but the future glory to come, which will replace that suffering, is.
guaranteed to come, and it's so amazing..
It's so great..
It's perfect..
It's glorious..
And it totally outweighs any suffering you can face in this life..
So look beyond your suffering..
Fix your eyes on the glory that is to come, and keep your hope in that which is unseen..
When you suffer, are you tempted to believe that you are alone, or that no one is there.
to help you?.
Are you inviting the Holy Spirit in the midst of your suffering and weakness to help you?.
Remember that even in your weakness, the Holy Spirit is there for you, interceding on your.
behalf, taking your needs to God, and the will of God will be done in your life..
Rely on the Holy Spirit..
Depend on His working in and through your weakness..
Turn to Him when you suffer..
And lastly, when you suffer, are you tempted to believe it's all for nothing?.

$^{561}$Do you truly believe that God works all things for your good, that your suffering is not.
meaningless?.
Remember that, as Paul says, the truth is that God does work all things for your good..
For those who love God and have been called according to His purpose, and we saw that.
this purpose has been in play from before the foundation of time..
And that's how you can be confident that God is and will work all things for your good..
And this means that your suffering is never in vain..
God's always working for you in eternal weight of future glory, even in the midst of your.
suffering..
Know that your Heavenly Father is good..
He's taking care of you..
No matter what we go through in this life, God is there for us, and He's got us for all.
eternity..
Trust in Him..
Will you please pray with me?.
Heavenly Father, Lord our God, I just want to thank you..
Thank you that you have given us this perspective on suffering, that you have given us a future.
hope that we cannot even fully comprehend, but we know it is so great, so glorious..
Any suffering we face in this life utterly pales in comparison..
Thank you, God, that you do not leave us to suffer alone, but that you have given us your.
Holy Spirit who helps us in the midst of our weakness, who takes our request to our Father.
in Heaven, and that His will will be done in our life, that all things will be worked.
out for good..
And we thank you, Father, that none of our suffering is meaningless, none of it is in.
vain or futile or purposeless, but that you are using it all for your good, which is our.
good..
Help us, Father God, to in all things become more like Jesus Christ..
Please take this word, hide it deep in our hearts so that when we suffer, we will know.
that we have hope, or when that someone near us is suffering, we will be able to share.
that hope with them..
And please, Father God, put a burden in our hearts to take this good news, to take this.
gospel of hope to a world that desperately needs it, so that those who do not know about.
this hope, who don't know that there is purpose in all suffering, that all things can be worked.
for good, help us to take this gospel to them, Father God, that they would receive the blessings.
and the hope that is found in no other name than Jesus Christ..
We thank you for this day..
We thank you for this opportunity to be here to hear your word, Father God..
Just pray that you would sanctify us by your word..
And I ask all these things in Jesus' precious name..
Amen..

$^{601}$[BLANK AUDIO].
\newpage



\section{4:7-12}
\label{sec:J6PHRTa1EUI}
\textbf{Love, the Greatest Gift, the Greatest Power, and the Greatest Purpose (1 John 4\_7-12) - Rev. Jeong}
\newline
\newline
連結: \href{https://youtube.com/watch?v=J6PHRTa1EUI}{\texttt{https://youtube.com/watch?v=J6PHRTa1EUI}} ~~~~ 語音日期: 2024-04-30
\newline
\newline
\hyperref[sec:ANI7oTOQdWo]{\small{< < < PREV SERMON < < <}}
~
\hyperref[sec:index_chronic]{\small{[返順時目]}}
~
\hyperref[sec:index_scriptual]{\small{[返順卷目]}}
~
\hyperref[sec:vMWYJSKl8Ig]{\small{> > > NEXT SERMON > > >}}
\newline
\newline
$^{1}$Mr. Jinwook for today's sermon..
(audience applauding).
- Good morning, brothers and sisters..
Did you hear me?.
Okay..
How have you been doing?.
It's been very cold these days,.
and Omicron is surging,.
and things are still going on,.
and we don't feel much peace sometimes..
But I want you to receive God's peace.
through this worship..
God wants to give you life, peace, joy, and love..
All these things we can only find in our Lord..
Today I'm gonna deliver God's message.
from 1 John, chapter four, verse 72 to 12..
(keyboard clicking).
Okay, so if you have your Bible, open your Bible..
1 John, chapter four, verse 72 to 12..
1 John, chapter four, verse 72 to 12,.
because it's a short passage, let's read together..
Dear friends, let us love one another,.
for love comes from God..
Everyone who loves has been born of God and knows God..
Whoever does not love does not know God,.
because God is love..
This is how God showed his love among us..
He sent his one and only Son into the world.
that we might live through him..
This is love, not that we loved God,.
but that he loved us and sent his Son.
as the atoning sacrifice for our sins..
Dear friends, since God so loved us,.
we also ought to love one another..
No one has ever seen God, but if you love one another,.
God lives in us, and his love is made complete in us..
This is the word of God..
Amen, amen..
So let's pray..
(pages rustling).

$^{41}$Father, here we are in front of you,.
and I'm gonna deliver your word to your people..
But I know that my lips are corrupt,.
and my heart is not truly loving..
So when I deliver this message,.
Father, I'm convicted in my heart already..
Holy Spirit, come upon us..
Let us listen to your voice, Lord..
Convict our hearts, and bring the changes in our lives.
that we may glorify you, that we may live in you..
In Jesus' name we pray, amen..
So last year, at the end of last year,.
when I was asked to deliver three times of Sunday messages,.
one per month, this year from January to March,.
I had to think about what I do really wanna share with you,.
brothers and sisters..
Just as many of you know, after 13 years of living.
as a missionary to America, my family and I.
are going back to my country next month..
And thankfully, some of you guys came to Christ.
through my ministry..
It was not my power, but it was God's grace..
Therefore, this is the last moment,.
so as my last word to you, brothers and sisters,.
I intentionally chose the three topics last December..
So prayer, I delivered the message about prayer in January..
Today, love, and living as a disciple next month..
So those are the three topics I really care.
that you wanna really understand me before I go..
Yeah, your pastor, Mark, will continuously.
remind you the same thing,.
but these are the three topics in my heart..
Prayer life, loving life, and living your life.
as a disciple..
So though I well know that I am weak in love.
and not qualified to preach about love,.
I am also convinced that love is so essential.
for a Christian's life..
And I'm gonna deliver this message about love.
Jesus showed us..

$^{81}$It's not my love, it's not my experience,.
but it's Jesus' love, and the love Jesus wants us.
to live out in our lives..
So let me ask you this question, brothers and sisters..
What is love for you?.
What is love?.
Have you ever seen a couple falling in love.
with each other?.
Some days ago, Jim Fan and Paul came to our house,.
and they were grabbing their hands.
and just hugging all the time, trying to touch each other..
So when their love is just starting,.
they are literally on flame, flame of fantasy,.
like fantasy love..
They are so happy with only their dating partner's presence..
They feel like they have everything,.
and life is so good and perfect..
But the problem is this..
How long does the flame of love last?.
It depends on the couples,.
or on different, some research measures..
However, we just know that just a few months later,.
or a couple of years later,.
or a few years later at most,.
the strong sense of love to win those.
decreases little by little..
And many couples, many husbands and wives,.
feel so dull about their relations..
Many feel and think that.
they do not love each other anymore..
And I watched a Korean live show..
The speaker asked the people in the studio hall,.
the studio hall was packed with people,.
and the speaker asked this question..
If you feel your couple, as a husband and wife,.
love each other, please raise your hand..
And I exactly remember eight people.
raised their hands among couple of hundreds..
Okay, only eight people..
So surely, husband and wife relationship.

$^{121}$should be the most meaningful and beautiful in our lives..
But truly loving relationship.
between husband and wife is quite rare..
It's quite rare..
What is the problem?.
The problem is this..
Even though we desire love so badly,.
we cannot hold even the greatest love in our lives..
Why do the people get married?.
Your parents, they're hating each other,.
but actually they love each other..
We cannot hold even the greatest love in our lives..
What does that mean?.
It means that we receive, oh no, no, no, I'm sorry..
So, we easily lose and destroy love..
So today I'm gonna help you understand.
three important lessons about love.
through today's passage..
Okay, three lessons..
Pastor's love number three, right?.
The first lesson about love is that.
love is the greatest gift in our lives..
Everybody repeat after me..
Love is the greatest gift in our lives..
Even though love is the greatest among all gifts,.
it's a gift..
Do you get it?.
Love is the greatest among all gifts,.
but it's still a gift..
What does that mean?.
It means that we receive and enjoy love for a brief time..
We didn't create love..
We couldn't think out of love..
We received everything about love..
So, verse seven in today's passage says like this..
Verse seven..
I'm sorry..
I'm gonna use this very fancy one..
Okay, let's read together..
Dear friends, let us love one another,.

$^{161}$for love comes from God..
One more time..
Dear friends, let us love one another,.
for love comes from God..
What does, where does love come from?.
According to this verse, from God..
And interestingly, my Korean Bible reads like this..
Love, yeah, Korean translation reads like this..
Love belongs to God..
Whether love comes from God or love belongs to God,.
both means that God is the creator of love..
God is the creator of love,.
and the owner of love, and the giver of love..
Again, God is the creator of love,.
owner of love, and the giver of love..
We receive love as a gift, roughly in two ways, from God..
We receive love from God in two ways..
Number one, we receive love as a gift,.
because God made us as loving beings..
God made us as loving beings..
He could have made us as another being,.
like unable to love..
Let's take an example of ostriches..
The Book of Job, the letters is quite small,.
but Book of Job, chapter 39, verse 14 to 16,.
describes this, how a mommy ostrich.
does never have any sense of love toward her babies..
Have you heard about that?.
It goes like this..
She, a mommy ostrich, lays her eggs on the ground.
and lets them worm in the sand..
Unmindful that a foot may crush them,.
that some wild animal may trample them..
She treats her youngs harshly,.
as if they were not hers..
She cares not that her labor was in vain..
Did you hear me?.
Why is a mommy ostrich like this?.
And the following verse, verse 17, gives the answer..
Verse 17, for God did not endow her with wisdom.

$^{201}$or give her a good sense, a share of good sense..
In other words, ostriches, mommy ostriches,.
cannot have truly caring love toward their youngs.
because God intentionally did not give the ability to love..
But in the opposite, why do we humans.
care so much about love?.
Do you care about love?.
Actually, if we express it or not with our mouth,.
we actually all the time, 24/7, we care for love..
We are flourishing because of love..
We are dying out lacking love..
Love is the essence..
Even though we have been hurt because of love,.
even though we have a lot of conflicts and troubles,.
because of love, we know that we cannot live without love..
Why is love so essential for human existence?.
Because God made us in that way..
God made us as loving beings..
So the Bible so many times emphasizes that.
God made us in his image, in his image..
Genesis chapter 1, 27 says,.
"So God created man in his own image..
"In the image of God, he created him..
"Male and female, he created them.".
So humans are in God's image..
You know that, right?.
And then, what is God's image?.
What is God's image?.
And verse eight in today's passage says,.
God is what?.
God is love..
God is love and we were created in God's image..
Do you understand?.
So our essence, being a human, means love..
So since made in God's image,.
we were created as loving beings..
This is why love is a gift from God..
God made us able to love..
Like what?.
Love between family members,.

$^{241}$love between men and women,.
love between brothers and sisters..
All these forms of true love..
Now the world are calling so many different ideas love,.
but they are not love..
But God gave us the right forms of love.
and all these right forms of love came from God.
and we are able to enjoy these loves.
because God gave us that ability..
This is why love is gift from God..
Second, we receive love as a gift.
because God redeemed our inability to love..
Why is it a gift, love is a gift?.
Because God redeemed our inability to love..
Though God made us as loving beings.
after Adam and Eve's fall, as you know,.
we lost the true loving ability..
Are you able to love without fear?.
We have a lot of fears..
We cannot love..
We are lonely, but we cannot love..
Our love became so self-centered..
We think we love,.
but actually we love only for ourselves..
This is the most serious problem in our life..
We think we love,.
but actually we love only ourselves..
Our self-centered love makes our life really painful.
and brings conflict so many times..
Why is your life so painful?.
Why is your life with all kinds of conflicts?.
Because we are self-centered..
In 2 Samuel chapter 13,.
there is a very horrible story in King David's family..
As you know, King David had many wives,.
and Amnon, his name was Amnon,.
Amnon was the oldest among all sons..
Amnon had a half-sister named Tamar..
Half-sister meaning that Amnon, this guy, and Tamar,.
they had the same father, King David,.

$^{281}$but their mothers were different..
But Amnon, this guy, began to love his half-sister, Tamar..
To make a long story short,.
Amnon asked his father, King David, to send Tamar,.
that she might take care of him.
when he was pretending to be ill..
He was pretending to be ill..
And then what happened?.
Amnon raped his sister.
and had a sexual relationship by force..
When he was trying to do this,.
Tamar was resisting him, saying,.
"Please don't do this dishonorable thing,.
but ask father to give me to you.".
Okay, but however, Amnon did not listen to her.
and defiled her..
Then after raping her, Amnon's mind abruptly changed..
He thought he was loving her..
He thought, "I love you so much, this much.".
But right after that, doing the wrong thing,.
his mind changed..
He thought,.
after he was, when he was raping her,.
Tamar, with the right sense, she resisted him..
And then this made him angry..
You know what?.
This made him angry..
So he drove her out from his house after the wrongdoing..
And then 2 Samuel, chapter 13, verse 15,.
exposes Amnon's insanely selfish love..
Insanely selfish love..
Then Amnon hated her with intense hatred..
In fact, he hated her more than he had loved her..
Amnon said to her, "Get up and get out.".
Is that love?.
Insanely selfish love..
Brothers and sisters, love is so essential.
and central in our lives..
We know that we cannot live without love..
But the problem is that we don't know how to love,.

$^{321}$how to truly love..
Many of us are going through bitter.
and painful stage of life.
because we were hurt by others whom we loved,.
or we hurt others in the name of love..
So we are now quite tired of loving others.
and scared to love anymore..
We are tired..
This is why Jesus had to come..
When Jesus died on the cross.
and shed all this blood for us,.
he renewed our heart..
He renewed our senses..
He cleansed our evil and dirty conscience..
We had a lot of evil and dirty things in our hearts,.
but Jesus cleansed it and gave us a new heart.
that we may be able to love each other again..
Because of Jesus, we can say like this,.
let's love again..
I am dirty..
You are hurt..
But we can have the courage to love each other again,.
forgiving each other again, trusting each other again..
Jesus made it possible..
When I was in my 20s and 30s,.
as I shared many times before,.
if you forget this, please remember this only about me..
Who was I, your pastor?.
In my 20s and 30s, sometimes I was extremely lonely..
I was extremely lonely..
I needed people, but I was afraid of making relationships..
And whenever I tried to approach people,.
I felt fear and burden..
And I even didn't have any hope.
toward the love between man and woman..
Because my heart was so defiled and dirty,.
I didn't have any hope toward the marriage.
between man and woman..
But you know, yeah, I had a lot of fear,.
guilt, and problems, pride, many problems..

$^{361}$But however, when Jesus entered my life,.
he literally renewed everything in me..
Jesus removed my sins..
Jesus removed my unclean thoughts and behaviors.
and my fears..
Jesus removed my fears..
I am not completely free from the fears,.
but removed most of the fears in my life..
When I continuously gave my life to him.
in obedience to him, Jesus filled my heart,.
filled my life with many meaningful relationships..
He enabled me to love others,.
even though it is not perfect..
And Jesus gave me my wife and my children..
Whenever I look at my children and my wife,.
and whenever I find my heart is loving them.
without shame, without guilt, without fear,.
I know that it is from Jesus..
It is not my love..
I know that all the people I have and I love,.
they came from Jesus,.
especially through his forgiveness.
and the power of the Holy Spirit,.
and the power of the Holy Spirit..
I want you to come to the worship..
I want you to know Jesus,.
because when you receive the Holy Spirit,.
there is the power,.
there is the only one thing that will hold you up,.
that will hold your love up.
with nothing else, with nothing else..
Now I fully know that if I have love, any love,.
it is not my own ability..
It's all because of Jesus..
This is why love is a gift from God,.
the greatest gift in our lives..
So today I'm giving this message,.
number one, love is the greatest gift in our lives..
Second, love is the greatest power..
Everybody, love is the greatest power..

$^{401}$There is a very well-known research.
called Grant and Gluck study, Grant and Gluck study..
For over 75 years, they researched 456,.
listen, 456 very poor men growing up in Boston.
from 1939 up to 2014..
75 years of research, can you imagine?.
And then 200, in the opposite, in comparison,.
268 male graduates from Harvard University,.
classes of 1939 to 1944..
Very contrasting group..
Very poor people, very wealthy people..
And their research, after 75 years later,.
tens of millions of dollars of research investment,.
this research concludes,.
the most important single factor in our life.
that makes our life truly happier and healthier is what?.
Only love..
Not money, not education, not position, nothing else..
But the single factor that makes human.
really healthy and happy is love..
Furthermore, when one has,.
when you have this solid and lasting love in your life,.
you feel that life is satisfying and even fulfilling..
So everything depends on love..
This is the power of love..
Love is the most important and the greatest power in life..
Since we forget spiritual things so easily, right?.
You guys forget, but spiritual things,.
it's easier to forget..
Next week, you don't remember.
what you heard the previous week, right?.
Since we forget spiritual things so easily,.
I repeat some very powerful stories again and again..
Some of you guys hear me, heard me,.
talking about this story again and again..
One of these stories is the man, Nick Vujicic..
Do you remember Nick Vujicic?.
Yeah, this is the guy..
When Nick Vujicic was born,.
he was now Australian pastor..

$^{441}$When he was born, without two arms, without one legs..
And the remaining one leg is a very tiny little one..
It's not like the real legs..
So, do you know, even at the age of eight,.
my second daughter, my first daughter, is now eight..
My daughter's 50s age..
When this guy, Nick Vujicic, when he was only eight years old,.
he learned about despair..
Only despair..
He couldn't find any hope in his life..
He thought no one would love him,.
or no woman would want to get married to..
Who will actually love him?.
Look at myself..
But however, in his teens, Nick Vujicic met Jesus Christ..
And the power of Jesus' love changed him..
And verse nine in today's passage, it goes like this..
Let's read together..
This is how God showed his love among us..
He sent his one and only son into the world.
that we might live through him..
Nick Vujicic, this brother, began to see hope.
in his life through Jesus' love..
He began to see his life as a unique one.
from God's perspective..
Before meeting Jesus, he thought, I do not have any hope..
But after meeting Jesus, Jesus' love changed his heart..
He began to see his life, my life is just unique..
My body is just unique, a gift from God..
From God..
He finished his college degree..
He swims without legs and arms..
He can swim..
He does all kinds of sports..
And nothing could stop him..
He drums with his little tiny foot..
Yeah, he plays drums..
He enjoys life fully..
And his wonderful character, his heart,.
and right attitude toward life.

$^{481}$made him a very good international speaker..
He travels all around the world..
And he gives the message about Jesus, message about hope..
So even though he does not have any arms to hug others,.
hundreds of thousands of people in the world,.
they hugged this guy..
Because they learned how to thank.
and how to love this life on earth..
And you know what?.
And finally, look at the picture..
Finally, a woman, beautiful woman,.
she loved Nick Vujicic..
They got married and have four beautiful children,.
as you see in the picture..
And where all these miracles start?.
Where these miracles start?.
Yes, from Jesus' love..
Love has the greatest power in our life..
When you meet true love, you will surely change,.
and your life will be surely different..
And that true love is Jesus, God's love..
Message is simple..
Love is the greatest gift..
Second, love is the greatest power..
And then, the last one, third, what is love to you?.
What is love?.
Third, love is the greatest purpose..
Everybody, love is the greatest purpose..
What does this mean, love is the greatest purpose?.
Since we were made in God's image,.
who is love, the very essence of being a human is love..
We were made to love each other fully without fear..
However, we fall in love with many undesirable..
We want to love, but we love many undesirable,.
unhealthy things in our lives..
We love, in other words, we love what we should not love..
Why do people pursue money?.
Yes, they love money more than any other things..
So for those people who love money,.
sometimes money is more important than relationship..

$^{521}$Sometimes money is more important than their family..
How about those people pursuing power?.
In a month, you know, in a month, in Korea,.
there will be the presidential election..
And then those presidential election candidates,.
they are pursuing power, right, and fame..
That's what they love..
For these guys, their candidates,.
like Joe Biden and Donald Trump did,.
nothing is more important..
So they almost try to destroy the others.
through all possible ways to get the power..
Because for them, power is what they love..
Here is a very famous picture..
When you look at the left side, you can see a picture..
The picture is very small,.
but there is a very black little child..
She's a Sudanese girl..
And she was going to the food center..
On her way to the feeding center,.
because of hunger, she collapsed on the road.
without any power to go on..
Do you see, there's a vulture..
There's a very big, yeah, like eagle,.
a vulture waiting for her to die..
Very tragic picture..
The photographer who took this picture.
was the guy on the right side..
His name was Kevin Carter, 1993, in Sudan..
At first, this photographer tried to drive away the vulture.
from the girl for a time..
However, when this guy was going to the place,.
people warned him not to touch anything.
because it could bring some infection, like coronavirus..
It was not that serious in those days,.
but some kinds of infection..
So he was warned..
At first, he was trying to protect her..
But later, he didn't really try to protect this girl.
and just went away..

$^{561}$Next year, because of this picture,.
Kevin Carter, he got Pulitzer Prize..
Do you know Pulitzer Prize?.
The highest honor for a journalist..
However, only a few months later,.
Kevin Carter, this guy, committed suicide..
He didn't love the girl's life at the moment..
However, later, when he came to the right mind,.
he realized that he loved the wrong things,.
his career, his report, his safety,.
while neglecting a human's tragic death..
Do you think that Kevin Carter is a really evil guy?.
But that's who we are..
We are choosing the wrong things over the right things..
We are loving the wrong things over the right things.
most of the time..
That's why the world is not changing..
That's why the church is not being criticized.
because we are not changing..
Because we love so many unhealthy things.
more than any other things..
We don't know how to love the right things..
My brothers and sisters, we are like him..
We are no better than him..
Are you choosing the right thing?.
Or are you neglecting the people you should love.
in your life?.
We have one life, brothers and sisters..
We have one purpose..
And how do we know the purpose of our lives?.
How do you know what is the purpose of your life?.
When we look at what you love..
When we look at what you love, it shows what you love..
What you care is the purpose of your life..
But Jesus gave us only one command..
Summarizing all the commands in the Old Testament..
316 commands and laws in the Old Testament..
Jesus summarized all the 613 laws into one..
Like this, love God and love your neighbor.
as you love yourself..

$^{601}$And the last supper, right before going to the cross..
Some hours later, Jesus was going to be killed..
But Jesus gave only one command..
John 15, chapter 15, verse 17..
This is my command, love each other..
Let's read together..
This is my command, love each other..
According to Jesus' command, there should be only.
one purpose in our lives, which is love..
Love is and should be the greatest purpose in our lives..
Because love is the greatest gift..
And love is the greatest power in our lives..
God wants us to live a truly happy and full life..
God wants to give us a good life..
And this kind of life comes only through true love..
Without true love, our life is just empty..
Without love, life's just empty..
Before closing, I'm gonna give you a very short three tips..
Here are a few practical ideas to help you.
to be a man and woman of a true love..
How can we do that?.
Number one, you need to love what is good and real..
Listen, you need to love what is good and real..
You should not love what is not good, what is not real..
I repeatedly have warned you..
Most of us are seriously addicted to media..
How many hours are you spending your time for media?.
Or your money?.
Or shopping?.
Or sexual temptations?.
Because you are young people..
God wants to give all those things for true love..
But if you are addicted to the wrong things,.
God cannot give you the right things..
And those things cannot make you truly happy..
They make you slaves of sin and steal your heart.
and destroy your life..
Instead, what is real and what is good,.
you need to love God..
You need to love God..

$^{641}$That is my urge for you guys..
You should love God..
Not just coming to the Sunday worship is not enough..
You should love God more than anything else..
If you love God, then you'll be blessed..
How can you love God?.
Make Sunday worship the center of your life..
My brothers and sisters,.
make the Sunday worship the center of your life..
Keep Sunday worship in any case..
It's simple..
If you neglect Sunday worship,.
if you choose work over Sunday worship,.
you are loving money more than God..
You need to make reading God's word.
and praying the first priority in your life..
You need to meet God in your daily life..
You need to read God's word..
You need to spend time with God..
That should be the first priority in your life..
I'm telling you the truth again..
You cannot serve God and money at the same time..
In your life, the greatest enemy,.
the greatest idol is money..
You should be careful..
Second, you need to love how to express love through words..
How is your mouth?.
Your words has a lot of power..
You need to learn how to hold up your mouth..
We speak so many useless, negative, and destructive words..
Then we mess up our relationship..
Why our relationships are messed up, broken?.
Mostly because of our hearts and because of our words..
And then we have been reinforced to spit,.
to spit poisonous words throughout our lives continuously..
Actually, when we are doing it,.
we are killing and ruining our families.
and people around us through our words..
And we cannot change ourselves..
Sometimes we understand that.

$^{681}$we are spitting out very poisonous words..
But can you change yourself?.
Can you change your mouth?.
Why can't you change?.
Because before that, you need to change your heart..
You need to change your heart..
When can we change?.
When we are filled with the Holy Spirit..
This is why you should love the Holy Spirit..
Do you desire the Holy Spirit?.
Holy Spirit, you are the first one I want to get,.
I want to receive, I want to spend time with..
When you love the Holy Spirit, He will change you..
He will change your heart, your mouth, and your behavior..
You know what?.
We remember a really beautiful word someone spoke to us..
When I was working as a teacher,.
a woman sister who was much younger than I,.
she spoke one word to me,.
maybe more than around 20 years ago..
But I cherish the expression.
because it really pleased me..
Very kind and honoring word..
But how about this?.
Instead, we cannot forget some very hurtful words.
someone told us many, many years ago..
Word has great power..
We need to learn how to love each other through words..
Third, the last one..
Love is spiritual..
What does that mean?.
Love includes some emotional parts, I know..
Emotion is very important..
However, ultimately, love is from God and it is spiritual..
So if you follow the way of God, what does that mean?.
Love is spiritual, what does that mean?.
If you follow the way of God,.
if you obey God's word,.
you'll be better ready to love..
Your heart will be better ready to love on a daily basis..

$^{721}$But if you do not follow the way of God,.
if you do not obey the way of God,.
you will never be able to love others..
Love belongs to God..
We cannot love others by ourselves..
Whenever we gather without Jesus, it's a tragedy..
You know, church is not a safe place..
Without Jesus, we can't hurt each other..
Only when we gather with Jesus, in the name of Jesus,.
whenever we worship, that is the time we are healed.
and our relationship will be different..
Let me close..
Listen, remember, love is the greatest gift coming from God,.
the greatest power coming from God's love and forgiveness..
And love is the greatest purpose of our life,.
found only in God..
\newpage



\section{出埃及記}
\label{sec:vMWYJSKl8Ig}
\textbf{永遠的頌讚 3 (出埃及記 25-31) -  陳崇基牧師}
\newline
\newline
連結: \href{https://youtube.com/watch?v=vMWYJSKl8Ig}{\texttt{https://youtube.com/watch?v=vMWYJSKl8Ig}} ~~~~ 語音日期: 2024-04-30
\newline
\newline
\hyperref[sec:J6PHRTa1EUI]{\small{< < < PREV SERMON < < <}}
~
\hyperref[sec:index_chronic]{\small{[返順時目]}}
~
\hyperref[sec:index_scriptual]{\small{[返順卷目]}}
~
\hyperref[sec:JpQUnDljfo8]{\small{> > > NEXT SERMON > > >}}
\newline
\newline
$^{1}$各位大英姐妹早安.
今次我們繼續敬拜的系列.
我們繼續思想到聖經裡面.
如何教導我們做一個敬拜的社群.
上一次我們提過有詩篇第八篇.
和創世記第二章.
我們說創世記第二章裡面所描述的.
Garden of Eden.
這個伊甸園.
伊甸園是向東的.
是吧 向東的.
還有基督派到守衛.
等等不同的術語等等.
告訴我們伊甸園不是一個花園.
而是一個聖所.
是一個敬拜的聖所.
所以所有的人.
在創造的時候.
其實他們有一個天職.
一個身份.
一個identity.
我們就是worshipper.
去敬拜上帝.
如果你想留意一下.
我記得上次說過.
好像看地圖一樣.
裡面有些symbols.
這個創世記第二章.
說到向東有基督派.
然後有生命樹.
河水從伊甸流出來.
意思就是說.
這個伊甸園在哪裡.
不會在地心曲力金吧.
正常應該是在山頂.
在聖山之上.
或者有一些術語說到行走.
神行走在當中.
譬如世紀第三章.
在這裡伴翼犯罪之後.

$^{41}$天起了涼風.
那人和他的妻子聽見搖滑.
在園中行走的聲音.
walking.
行走.
其實是和上帝有一個關係.
和上帝一起.
要同行的意思.
但是因為他們犯罪.
所以他們就躲起來了.
然後他們被趕出伊甸園.
第四章十六節.
你也看到該人又離開了搖滑的面.
the presence.
或者the face.
所以伊甸園本身是一個敬拜.
和上帝建立關係的一個地方.
如果這樣看的話.
聖經從頭到尾都是講敬拜的.
開頭講敬拜當然.
第一章第二章.
到最後.
啟示錄.
啟示錄說看見新天新地.
先前的天.
先前的地已經過去了.
海也沒有了.
海因為邪惡.
代表邪惡.
看見聖神新耶路撒冷從上帝那裡從天降臨.
預備好了.
好像新娘打扮整齊等候丈夫一樣.
聽見大聲音再補助出來.
他說看啊上帝的帳幕在人間.
他要和人同住.
walking together.
living together.
要作他的子民.
上帝要親自與他們同在the presence.
上帝要察覺他們眼裡不再有死亡.

$^{81}$不再有痛苦.
上帝已經過去了.
所以你想想聖經從頭到尾.
其實就好像書架.
書架的書是不是有bookend.
我知道你不明白.
因為你家裡沒有書.
很少人有書.
但是有bookend書架裡面的腳.
頂住.
從頭到尾都是在講敬拜.
從敬拜到敬拜.
或者聽聽偉大的樂章.
開頭的序曲.
到最後的chord.
都是差不多這樣的聲音.
差不多這樣的melody.
由敬拜到敬拜.
整本聖經都是在講worship.
上帝值得我們去敬拜.
雖然我們人犯罪.
我們是流亡於伊甸園之外.
我們沒有了上帝的同在.
但是仍然是藉著上帝的恩典.
可以和上帝同行.
所以所有敬拜的人.
由《創世記》第三章人搏之後.
怎樣形容他們呢.
敬拜主的人和上帝仍然有關係的人.
就是與上帝同行的人.
譬如第五章23節.
以樂共活了365年.
這裡不應該是他真的這麼長壽.
但是symbolic number.
一整年了.
一個solar year.
365年.
以樂與神同行.
神就將他接去了.
48章15節.

$^{121}$《創世記》最後.
雅各為他種子祝福的時候.
他說什麼呢.
他形容亞伯拉罕和以撒是怎樣的人.
是與上帝同行的人.
藉著上帝去祝福他的子孫.
拜登最近幾年經常說這六個字.
行公義好憐憫.
因為世事太多不義的事.
好好提醒.
告訴我要行.
行公義好憐憫.
尼加書6章第8節.
不過很多時候很少說後面那句.
因為我想對戀比較容易.
行公義好憐憫.
但後面說了什麼.
傳謙卑的心與神同行.
去敬拜他.
這就是上帝對人的旨意.
行公義好憐憫不只是做好事.
是與上帝建立一個約的關係.
與上帝同行.
這就是敬拜的生命.
與上帝同行.
小時候有沒有玩過一個遊戲.
二人三族.
有吧.
我們小時候團契玩的低能遊戲.
就是那些.
你明明不喜歡那個人.
要跟他一起二人三族.
被迫的.
導師說你要跟他.
你就逼著這樣做.
什麼叫二人三族.
就是要跟他同行.
結伴同行.
或者堅持要跟他同行.
上帝就是我們最忠實的結伴者.

$^{161}$或者配伴者.
所以聖經從頭到尾都告訴我們.
神和他的子民是想同在的.
想同行的.
所以整本聖經從頭到尾.
好像兩個bookend.
就告訴我們一個長篇的故事.
告訴我們.
上帝希望跟人同住.
跟人一起同住.
以至人能夠跟他一起敬拜.
跟他建立關係.
所以創世記到啟示錄.
都是在說這個主題.
我們試試去到中間發生什麼事.
大家記得創世記最後的時候發生什麼事.
列祖或者雅各和他的十二眾子.
他們整家人去了哪裡.
埃及沒錯.
去了埃及.
他們流亡去了埃及.
然後多年之後.
上帝就去呼召摩西.
摩西也是流亡於曠野的.
他就跟摩西來到顯現什麼呢.
就是摩西在牧羊的時候.
他已經在一個很遠的地方叫米田地.
他就領羊群到曠野那一邊.
到了什麼地方.
神的山.
敬拜上帝.
遇見上帝的地方.
有上帝同在的地方.
都是山上面的.
叫做赫烈山或者西乃山.
而華氏在荊棘裡面.
火焰裡面跟他顯現.
摩西見到了.
他就說奇怪了.
他想走過去.

$^{201}$我們skip到第五節.
上帝就說你不要過來.
因為你要先將鞋脫下.
因為你所站的地方就是century.
就是聖地.
Holy ground.
Sacred ground.
上帝就說我必與你同在.
呼召他回去帶以色列人出埃及.
你們必回來這個山.
去到事奉上帝.
或者應該敬拜上帝.
所以上帝去到拯救以色列人.
出埃及.
拯救他們的子民.
用恩典來拯救他們之後.
就說你必定要去敬拜上帝.
事奉上帝.
敬拜上帝.
然後你記得.
出埃及裡面的情節.
十災等等.
就打擊了法老.
然後他們就出了埃及.
經過了曠野之後.
上帝就用火和雲來向他們顯現.
他們就回到這個山.
這個叫做赫烈山或者西乃山.
在那裡上帝向他們顯現.
有黑雲密佈.
有火.
有雲.
有風.
地震動.
然後上帝就說話.
說什麼.
Hello.
不是說Hello.
說十誡.
十誡就是以色列國.

$^{241}$以色列民的憲章.
他們的Constitution.
所以有多少東西看到都好.
不是靠見的上帝.
上帝要靠聽.
所以上帝就說話了.
我使耶和華.
曾經將你帶你出來埃及.
為奴之地.
我將你帶你出來.
然後第一屆.
第二屆.
第三屆.
然後上帝就說.
我要住在你們中間.
用什麼方法去住在你們中間呢.
因為我們在說敬拜.
所以我們想從頭到尾.
去追蹤一下.
到底上帝和人之間的關係.
敬拜是有什麼形式.
在出埃及的時候.
他們在西乃山的山腳的時候.
上帝就說.
你要照我給你的藍圖.
去做什麼.
做會幕.
一個Tabernacle.
Tabernacle就是代表.
我在你們中間.
你就是來到這個Tabernacle那裡.
你都可以牽上你的制.
去到盞尾去敬拜.
你就是這樣成為一個敬拜的子民.
什麼叫做Tabernacle.
什麼叫做會幕.
會幕在希伯來文的原文.
叫做Mishkan.
帳幕的意思.
Mishkan.

$^{281}$希伯來文就很妙.
有時候加一個小小的字上去.
加一個音上去.
你就變成了另一個字.
名詞是Mishkan.
動詞就是Shakan.
所以你看到住這個字.
Shakan.
加一個M在前面.
就是什麼.
Mishkan.
Mishkan就是.
很簡單吧.
很簡單.
加一個東西在前面.
它就變成了.
將動詞變成了一個名詞.
所以什麼叫做會幕.
其實就是住的意思.
不是很難吧.
最近我才知道.
有一樣東西叫做Camping和Camping.
有什麼分別.
大家聽過吧.
什麼叫Camping.
露營吧.
什麼叫豪華露營.
奢華露營.
就是加一個G上去.
Glamorous的Camping.
這個叫做Camping.
帳幕的意思.
這個叫Camping.
OK.
Camping.
這個叫Camping.
這個是什麼.
Camping.
這個呢.
Camping.

$^{321}$Glamorous的Camping.
會幕就是一個帳幕.
當然帳幕就是代表上帝住在人之間.
所以我們看一些經文.
講到會幕的建造.
或者會幕的形容.
二十五章至四十章.
整個在《昭迺及記》.
二十五章至四十章.
都是講會幕的建造和會幕的藍圖等等.
二十五章第一節.
耶路亞話吩咐摩西說.
你要吩咐以色列人獻禮物給我.
要將他們收來的禮物.
就是金,銀,銅.
藍色,紫色,珠紅色.
細麻,山羊毛.
染紅的公羊毛.
精美的皮料.
你說他們哪裡有這麼多東西.
他們不是難民.
他們不是今天的難民.
他們出埃及的時候.
是拿了一些金銀的.
因為他們打劫.
他們打贏了埃及.
所以就像打贏仗一樣.
有戰利品.
於是他們拿著很多這樣的東西就走了.
然後去到山腳的時候.
他說你要將金銀和一些最好的料.
精美的皮料.
LV手袋等等.
全部拿過來.
做什麼.
做這個會幕.
紅馬路和寶石.
那些卡的戒指等等.
全部拿下來.
你要做一個century.

$^{361}$聖所.
使我能夠mishkan.
能夠住在你們中間.
你們就按照我的指示.
將其中一切器具的樣式.
照樣去做.
所以上帝說在天上.
我有一個pattern給你.
我有一個藍圖給你.
我們做architect的就最清楚.
我認識一些architect的朋友.
他們看著藍圖.
他們好像在默想.
很喜歡看那些東西.
怎樣畫那些東西.
畫得多好.
好像那些IT佬.
看那些code.
這個code寫得很好.
有沒有這樣的感覺.
寫得很好.
這裡就是用一些architectural的language.
一些建造的元素.
或者一些語言.
來表達上帝在中間去住.
住在人中間.
以致人能夠敬拜他.
我們看多一點.
首先說到做約櫃.
他是用造格木.
造格木做一個櫃.
兩爪半長.
一爪半闊.
一爪半高.
前後內外.
這個包金的.
有金牙邊.
放發彈進去.
裡面有私人座等等.
要用金紙做出兩隻.

$^{401}$什麼.
基奴柏你記得嗎.
在哪裡出現.
這個朋友在創世記.
開頭的時候就出現了.
因為聖所就會有基奴柏.
在私人座的兩頭.
拜託,什麼叫造格木.
如果你以色列開車經常去礦業經過.
經常有雨傘的樹.
那些就叫造格木.
造格木就在礦業生長.
你撒開的時候.
能夠儲到水分.
然後造格木是很防蟲住的.
很難會有蟲住.
是不是比孫子家師好.
所以就用造格木.
在礦業裡面找到的造格木.
來做一個約櫃.
然後他又繼續說.
陳設餅.
陳設餅的桌子.
又用造格木做一張桌子.
看多長多高多闊.
包金的.
上面有貢餅.
就是陳設餅.
放在桌上.
上在我的面前.
My face.
My presence.
所以人們看到會幕的時候.
他就說這個就是.
The presence of God.
就是上帝的同在在我們中間.
然後有萬紙.
十幅的萬紙.
萬紙是用釵的細麻布.
藍色的紫色的珠紅色的織繩.

$^{441}$然後繡上什麼.
又是我的朋友基魯柏.
因為他是把守聖所.
然後有樹板.
樹板是什麼.
就是那些帳幕.
那些你背後那些.
一塊這樣.
你拿去隔著東西.
就要做這個樹板.
然後就把布放在上面.
帳幕的後面.
最後一句.
就是西面.
有六塊.
為什麼帳幕的後面是西面.
因為會幕是向哪邊.
門口是向東.
所以後面就叫做西面.
這個就叫做會幕.
就是上帝要住在他的子民中間.
在出埃及記的時候.
他說我住在你中間.
你就是繼續來敬拜我.
知道看這些字很悶.
有圖有真相.
這裡就是以色列在南部一個公園.
他們做了一個一比一的會幕的模型.
是什麼形狀.
長方形沒錯.
你不要看後面那些貨車的垃圾桶.
不是.
前面這個.
有樹板.
然後有一些帳幕來攔住.
長方形.
二比一.
這個就是Mishkan會幕.
你進去的時候.
就在會幕的外園.

$^{481}$前面這個是什麼.
不是讓你泊車.
那個是祭壇.
走上去.
祭壇是什麼.
就是燒烤的.
把祭物燒了.
燒給上帝.
獻給上帝.
再前面圓形那個.
就是叫做銅盤.
就是他們潔淨用的洗手等等.
然後有一個簾.
看到嗎.
有一塊布.
裡面叫什麼.
聖鎖.
Holy of Holies.
長方形的.
小一點的長方形.
又是二比一.
然後進去的時候.
你看到兩位朋友.
那個好像沒事做.
其實是祭司.
左邊那個是什麼.
七燈臺.
好像生命樹的樣子.
哪裡提過生命樹.
創世記第二章.
右邊那個是什麼.
猜到嗎.
陳設餅.
每天都要燒新的餅.
然後獻給上帝.
或者應該這麼說.
你當你是一個遊牧民族的人.
上帝就好像遊民族的主人家.
他就邀請你進來.
然後他有東西款待你.

$^{521}$那個叫做陳設餅.
Show bread.
中間那塊.
那個就是香壇.
香壇就代表上帝接受人的禱告.
當他點起香的時候.
香向上升.
然後上帝就接受禱告.
再裡面.
上面有什麼.
這個只是聖鎖.
不是致聖鎖.
有個致聖鎖.
在裡面的.
這塊布後面.
就是正方形.
一比一.
Perfect square.
在近東的社會裡面.
就代表perfection.
代表holy.
代表神聖的地方.
上去看看.
我想告訴你.
這個致聖鎖.
是沒有人進去的.
一年只進去一次.
就是雍吉泊的時候.
大概是西曆九月十月的時候.
雍吉泊就是Day of Atonement.
贖罪日的時候.
贖罪日的時候不是每個人進去打卡.
進去打卡拍照.
是一個人進去.
大祭司.
大祭司潔淨好.
預備好.
代表整個的以色列人民.
然後他才進去.
就這麼多.

$^{561}$所以今天.
你和我可以進去.
裡面有什麼.
就是約櫃.
就是這個上面有兩隻基督拍.
然後是一個空盒.
這個約櫃叫做私人座.
什麼意思.
上面叫私人座為什麼.
因為其實是上帝的座駕.
是上帝的Air Force One.
他坐在這裡就是Oval Office.
就是華盛頓白宮裡面的辦公室.
這個是Situation Room.
這個是Control Room.
整個宇宙.
整個的Cosmos.
就是由這裡來控制.
因為上帝坐在這裡.
這個盒子裡面有什麼.
剛才的經文說到.
有兩塊法板.
是兩塊同樣的法板.
identical copy.
為什麼因為近東的時候.
兩個國家兩個人一起來.
有一個約的時候.
covenant的時候.
就說你和我現在立約.
然後立約的條文是什麼.
我們大家一起寫.
寫完之後就make a copy.
你就拿回你國家最大的廟.
就放在你的廟的致勝所裡面.
我就放在我最大的廟的致勝所裡面.
你一個copy我一個copy.
大家就記住了.
世世代代這樣記住.
我們大家應該忠於大家.
對不對.

$^{601}$以色人最大的.
當然只有一個廟.
只有一個聖所.
就是致勝所.
就是會幕.
所以他們的copy就放在這裡.
耶和華的copy在哪裡.
都是同一個地方.
對不對.
不會放在別的地方.
所以兩個copy就放在約會裡面.
這個就是cabernet.
就是會幕.
會幕很奇怪.
在邵逸琦形容會幕的建造的時候.
是用七段的.
每一次說到做一些事情的時候.
二十五章第一節.
就有一句開始叫做耶和華.
曉如摩西.
然後第二段三十章十一節.
說到捐獻.
會幕的敬拜就需要有捐獻.
然後他就會說.
曉如摩西.
然後說到捐獻的事.
然後又說這句話.
然後又說到高柔香.
忌公等等.
到最後的時候.
耶和華曉如摩西.
你要守安息人.
你要敬拜我.
七段來形容整個會幕的設計.
整個的藍圖.
為什麼七段.
創世的第一章.
整個宇宙就好像聖殿一樣.
在上帝同在的地方.
會幕就更加是.

$^{641}$所以用七段來形容整個的敬拜.
用七段來形容整個會幕的symbol.
就是說你看到會幕的時候.
就好像看到一個小的宇宙一樣.
一個原本原創的宇宙.
那個宇宙是一個好像創世的第一章.
是沒有罪的地方.
是敬拜.
是上帝同在的地方.
在那裡才找到神的面.
神的presence.
這裡很簡單.
我想看看一些經文.
說到會幕當有的敬拜.
當他們製造會幕的時候.
去接受這個藍圖的時候.
這群群眾當時叫做以色列民.
神的子民.
他們怎樣去面對這個.
上帝與我們同在.
住在我們中間.
祂的榮耀向我們彰顯.
我們怎樣去回應呢.
第一點.
是他們甘心樂意的去奉獻的.
會幕落成之後.
三十五章第二學節.
凡有智慧的心裡面受感動的婦女.
都是將她們紡織的山羊毛.
官長就將寶石和馬路.
就陷在大祭司前面的牌子.
就拿來了.
然後他們各自就做香做油.
或者點燈用的香料和油.
以色列人無論男女.
凡心裡面受感動的.
就帶著這些的公物.
甘心的兒子去獻給耶和華.
為做耶和華藉著摩西所吩咐的.
一切回墨的公.

$^{681}$甘心的奉獻來獻上.
建造一個回墓.
以至上帝能夠在我們中間.
來與我們同住.
第二方面.
每一個以色列人.
當然包括祭司.
祭司更加.
就是要做什麼呢.
就是要守住這個聖所.
你記得開始的時候.
創世記第二章.
是不是有人要把守這個聖所.
基督伯 亞當.
是需要去把守這個聖堂.
利未記二十章第八節就吩咐以色列人.
他們要謹守上帝的律法.
遵行它.
分別為聖.
當時繁俗的東西.
不潔的東西.
是不能進入聖所的.
罪惡是不能夠進入聖所的.
祭司的工作就是做這件事.
這裡說什麼.
會眾也有責任.
聖潔.
分別為聖的生活.
你可以發現那些律法.
有很多東西說到.
幾乎是整個人的生命.
整個人的生活.
起居飲食.
以至他們安息日的敬拜.
都是說要.
separate.
be holy.
分別為聖.
第三.
他們的生命就是需要與神同行.

$^{721}$敬拜.
就是見神的面.
跟他一起同行.
二十六章十一節利未記.
我要在你們中間立我的帳幕.
我要行走在你們中間.
walking in the midst of you guys.
作你們的上帝.
你們要作我的子民.
我是耶和華你們的上帝.
曾經從埃及領你們出來.
你們不再作埃及的奴僕了.
法老不是你們的主人.
耶和華才是你們的主人.
說到今生奉獻.
說到守衛聖所.
與神同行.
這就是舊約的子民他們需要做的事.
當然你說今天已經新約了.
我們不需要有一個回眸.
你不要當這裡是回眸.
這裡不是回眸.
在新約裡面發生什麼事.
我們不需要再有牛羊的祭禮.
因為拜登的祭禮.
只不過是一個路牌指向.
一個真實的祭禮.
就是基督.
全部都是舊約裡面的sign.
我們都不需要有祭司.
因為最大的祭司已經到了.
就是耶穌基督自己.
更加不需要一個.
physical的實物的回眸.
因為基督已經到了.
約翰福音第一章.
道成了肉身.
住在我們中間.
充充滿滿.
有真理.

$^{761}$我們也見過他的榮光.
正是負獨生子的榮光.
一章十四字.
你猜一下到底哪個字是和回眸有關.
道成了肉身住在我們中間.
住這個字.
mishkan.
住在我們中間.
以致我們見到.
上帝獨生子的榮光.
約翰福音最後的時候.
耶穌臨上十字架.
做了一個禱告.
我們叫他大祭司禱告.
Prayer of the High Priest.
他說什麼.
他說.
天父你已經將榮耀賜給我.
我也將這個榮耀賜給門徒.
使他們合而為一.
將我們合而為一.
大祭司的禱告.
榮耀的彰顯.
原來上帝的榮耀在哪裡找到.
在我們中間.
在跟隨他的門徒.
敬拜他的人中間.
當我們聚集一起敬拜的時候.
我們說上帝你是我們的主.
我們是跟隨你.
你起碼會見到上帝的榮耀.
這個就是上帝榮耀.
彰顯臨到我們中間.
所以那些講到聖所的語言.
那些術語.
那時候就應用在新約基督來了之後.
所有的教會.
基督的門徒身上.
譬如這裡有一個例子.
哥林多前書第三章十六節.

$^{801}$大家知道哥林多前書發生什麼事.
很多問題那個教會.
你說今天教會很多問題.
怎麼都不夠哥林多多.
哥林多嘛.
哥林多大把問題.
他就在那邊摒他們.
你這樣都不明白.
你信了福音怎麼可以這樣行事為人.
怎麼可以這樣去分黨等等.
三章十六節.
難道你不知你們是.
你們就是上帝的電.
上帝的靈住在你們中間.
凡是有人毀壞上帝的電.
上帝一定會毀滅那人.
這裡講到很狠.
他們就是在毀壞上帝的電.
毀壞上帝的教會.
毀壞上帝的電.
保羅說上帝一定會毀滅那人.
因為上帝的電是神聖的.
這電就是你們.
你們教會就有如上帝的電.
上帝的電就是上帝住在中間.
我們能夠見上帝的面.
和去敬拜他.
不是一個建築物.
是我們的生命裡面.
在我們這群群眾裡面.
這就是上帝的電.
我事奉主有二十多年.
我被呼召去牧養的時候.
我一向我的心願都是這樣.
去建立教會.
由我二十多歲開始.
到我現在.
我一直都沒有變過.
因為我知道上帝有一個很重要的計劃.
就是藉著教會去彰顯他的榮耀.

$^{841}$上帝的靈住在教會當中.
不代表我們是完美的.
但是我們需要努力去建造上帝的電.
上帝的教會.
使教會能夠茁壯起來.
我去過很多不同的教會.
做過不同的教會.
我的心願都是一樣.
我發覺我們的心願都是一樣.
不是去毀壞上帝的電.
而是去建立上帝的電.
建立上帝的教會.
這個應該是你和我都應該有的願望.
我們最近剛開始請了一個青年的傳道.
G.
大家認識他嗎?.
剛才打鼓的哥哥.
很厲害吧?.
杜願神學.
因為不同事故的緣故.
他之前的教會是H1B.
結果別人讓他高了.
他剩下很少時間.
否則他就要回到馬來西亞.
很巧合地他就認識了我.
我們總會其中一些人聯繫了.
我在這裡見他.
第一次單對單我和他聊天.
聊完之後.
你知道請人是很難的.
教會請人好像拍拖一樣.
有時候對有時候不對.
我聊完之後.
我都有一個心願.
一個年輕人.
很有前途的年輕人.
但是他願意去事奉.
願意去牧養教會.
不容易走的一條路.
我和他聊完之後.

$^{881}$最後我和他說什麼.
我說我和你認識不深.
你對我認識不深.
可能你認識我之後覺得我很惡搞.
不奇怪.
最後我和他說.
我願意給你一個機會.
我也和他說.
我也希望你給我一個機會.
我們在這裡如果能夠幫助他.
去成長.
去成為一個新的牧者.
這不是一個最尾的事.
是建造教會的事.
我知道我們教會不是很多錢.
資源不是很多.
我們上次開會都很頭痛.
我甚至和Colin說.
如果我們能栽培到一個新的人.
如果我們真的沒有錢的話.
OK.
你在我薪水裡扣10\%出來.
給這個人.
我兩個小孩都長大了.
他們賺不到很多.
但希望不會餓死他們.
我認真的.
我們在這裡做什麼.
我們建立上帝的教會.
神的榮耀要在我們中間.
如果有人可以被栽培.
能夠建立上帝的子民.
I want to go for it.
建立上帝的教會不應該推搪.
也不應該遲延.
我們太舒服了.
我們太自保了.
什麼叫做去敬拜上帝.
如果上帝是真的.
如果你信耶穌基督是真的.

$^{921}$死而復活.
如果是真的.
為什麼不睡在室內.
為什麼要保留那麼多.
如果你有保留的話.
那以後去天堂的時候會不會有保留.
What I am saying is.
建立上帝的殿.
建立上帝的教會.
應該是我們最高最高的priority.
我希望大家聽清楚.
我們呼籲大家.
我們做的事情不只是建造建築.
我們建立上帝的國度.
我們不是在玩泥沙.
上帝的真理 上帝的福音.
是真的.
不只是讓我們玩玩玩意.
開開心心就算了.
這個真理是值得我們以生以死去相逢.
建立上帝的教會.
因為這個教會是有上帝的靈.
同在上帝的榮耀.
你想見到的話.
我們就要虛罪.
不要太多保留.
因為耶穌基督對我們都沒有任何保留.
對你都沒有任何保留.
十字架的救恩對你.
是一個重價買回來的.
你的生命.
你有的資源是應該怎樣去用.
我們同心低流祈禱.
願上帝幫助我們.
主我們來到你面前.
我知道你在我們中間.
因為每一個信耶穌的門徒.
去稱你為我們主的教會.
去敬拜的時候.
我們用詩歌用我們禱告.

$^{961}$用我們的語言.
我們說主你是萬王之王萬主之主.
上帝我們求你.
求你讓我們見到你的榮耀.
見到你少少的榮耀的彰顯.
主與我們教會同在.
以至我們建立一個土里喜悅的群體.
叫我們不單行公而好憐憫.
也都是同理也同行.
主赦免我們很多時候對你有保留.
我們太多東西充塞我們的生命.
主叫我們願意奉獻.
願意去敬拜你.
願意去事奉你.
因為你是我們真真正正.
配得我們讚美.
配得我們敬拜跟從的耶穌基督.
求你在我們教會當中.
你再幫助我們.
能夠去建立這個身體.
以至你自己的榮耀.
我們能夠見到彰顯.
感恩讚美你.
我們這樣祈禱奉著美感.
今年是我們教會有不同的高度的時候.
我希望弟兄姊妹好像坐車一樣.
Hang in there.
坐穩一點.
我們有不同的改變.
我們也有同工退休.
也有新的同工到.
大家可以放心.
我還沒走.
我一定會Stick to the end.
但我們總會有一些.
同工上,人事上的改變.
上一次我們已經說過.
吳牧師會在月底退休.
今天是他最後一個主日.
很忠心在我們中間事奉.

$^{1001}$今天他是最後一個主日.
我們將一份禮物送給吳牧師.
我們請吳牧師上來.
吳牧師謝謝你.
我們也為吳牧師祈禱.
為他的身體祈禱.
為他的家人需要照顧等等.
希望主給他平安.
也希望大家為教會禱告.
特別今天我們有GA.
希望大家來.
我知道大家說.
我們要喝茶,吃東西等等.
我們買一盤糕.
我們有財務上的需要.
我們有事工上的發展.
也要介紹我們的新人同工G.
有很多資訊.
我們希望大家會有掌握.
願上帝祝福各位.
\newpage



\section{}
\label{sec:JpQUnDljfo8}
\textbf{Family Sunday School - EP 4 (It Could Have Been You)}
\newline
\newline
連結: \href{https://youtube.com/watch?v=JpQUnDljfo8}{\texttt{https://youtube.com/watch?v=JpQUnDljfo8}} ~~~~ 語音日期: 2024-05-04
\newline
\newline
\hyperref[sec:vMWYJSKl8Ig]{\small{< < < PREV SERMON < < <}}
~
\hyperref[sec:index_chronic]{\small{[返順時目]}}
~
\hyperref[sec:index_scriptual]{\small{[返順卷目]}}
~
\hyperref[sec:K7tbzr00wo0]{\small{> > > NEXT SERMON > > >}}
\newline
\newline
$^{1}$[Music].
Hello everyone, I'm Miss Rebecca..
Welcome to Family Sunday School..
This is the fourth episode..
What did we talk about last time?.
Remember?.
We talked about Jesus being arrested..
A very important distinction,.
a very important difference is.
being tied up and dragged away..
Or what?.
Jesus allowed them to drag him away..
Do you want to tie him up? No need..
Do you want to lock him up? No need..
Do you want to drag him away? No need..
Jesus just followed them..
I haven't mentioned this before..
This is God's plan to save people..
God said this plan must be fulfilled..
So Jesus had to obey..
Jesus knew what was going to happen..
So he just went with it..
Remember?.
Last episode, the last Chinese time,.
I gave you a challenge..
Remember the challenge?.
Because last time we talked about a lot of different things..
Angles..
Right?.
Different perspectives, different angles..
We looked at the story from different angles..
Judas, John, Peter, Jesus..
Right?.
Does anyone remember the Chinese challenge?.
Let me remind you..
We said that "角" can be used as "angle"..
If triangle is "三角形",.
then what is hexagon in Chinese?.
Does anyone know?.
Does anyone remember?.

$^{41}$Have you talked to your family?.
I'll break it down. Ready?.
Triangle. "Tri" means "三",.
"Angle" means "角",.
"三角形" means "three-angle shape"..
How many "角" are there in hexagon?.
One, two, three, four, five, six..
Right?.
Six..
"六角形".
Right?.
Hexagon..
"Hex" means "six", "hexagon"..
Wow..
Let's repeat the numbers..
If I say "octagon",.
you will know, right?.
Let's repeat the numbers..
One, two, three, four, five, six, seven, eight, nine, ten..
Backwards..
Countdown..
Ten, nine, eight, seven, six, five, four, three, two, one..
What is "octagon"?.
Let's talk about it in Chinese..
Hold the thought..
Okay..
In this episode, we will continue the story..
Jesus followed them to be taken away..
Where did he go?.
What did he do?.
Do you remember who was following them?.
Who was secretly following them?.
Yes..
John and Peter..
Hmm..
Was it dangerous to follow them?.
Was it dangerous for John and Peter?.
Hmm..
Where did they take Jesus?.
What happened to Jesus?.

$^{81}$We, as modern Christians,.
know a little bit about it..
Because we know what "Easter" means..
But at that moment,.
John and Peter didn't know, right?.
Hmm..
So today, we will continue the story..
Before we talk about the story,.
we must play a game, right?.
Mr. Becker loves playing games..
Okay..
Let's go..
The game we are playing this time is called "Name That Animal"..
Guess what animal it is..
How to guess?.
It's not based on eyes..
It's based on ears..
You have to listen to these sounds.
to guess what animal it is..
Here it comes..
The first animal..
Or you can close your eyes..
Then you can hear these sounds clearly..
And you can guess what animal it is..
The first animal..
Who is it?.
Guess..
Frog..
Here comes the next one..
The hint is,.
"I am large and loud.".
Hmm..
What is large and loud?.
Let's listen..
[Name That Animal].
It's loud..
Can anyone guess what it is?.
Elephant..
Can you hear the different sounds?.
Of course the message is different..

$^{121}$What it wants to say..
What it wants to express is different..
Elephant..
Here comes the next one..
The hint is,.
"I am alarming.".
Alarming is like an alarm clock..
It wakes us up..
It's loud..
Alarming..
Hmm..
Who has this animal?.
Wake it up..
Who has heard this sound?.
Wake it up..
[Name That Animal].
Anyone knows?.
Guess..
[Name That Rooster].
Anyone has this sound?.
Here comes the next one..
The hint is,.
"I wake up when you go to sleep.".
[Name That Animal].
Hmm..
Anyone has this animal?.
[Name That Bird].
Do you want to hear it again?.
[Name That Bat].
Hmm..
Guess..
Bat..
[Name That Squeak].
[Name That Bell].
Here comes the next one..
The hint is,.
"I am cleaner than you may think.".
[Name That Bird].
Let's hear this sound..
Or,.

$^{161}$"Has anyone been to the farm?".
You may have heard this sound..
[Name That Animal].
Do you know what animal it is?.
We usually think it is dirty,.
but it is actually a very clean animal..
Let's guess..
Pig..
[Name That Pig].
In the story,.
we may say, "Oink, oink.".
But actually, it sounds like this..
[Name That Pig].
Okay..
The last animal..
The hint is,.
"I like to talk at night.".
[Name That Animal].
Let's listen..
[Name That Animal].
Okay..
Let's listen quietly..
[Name That Bird].
What is it?.
Have you heard of it?.
Is it familiar?.
It is a cricket..
[Name That Cricket].
You heard the sound..
It is not just one cricket,.
but many crickets..
Let's talk..
I want you to think about.
which animal will appear.
in the next story..
Please pay attention..
Okay..
Let's continue the story..
Peter and John follow Jesus..
What will happen?.

$^{201}$They will take Jesus away..
They will find a way to take Jesus away..
John and Peter must follow Jesus..
Let's see what will happen..
Where will they follow Jesus?.
They will go to a big temple..
In the past, religious leaders.
would gather in a big place..
There are buildings,.
and a small garden..
A woman sees him in the garden..
"Oh, you look so familiar..
You look like your friends in Jesus' story.".
What will Peter say?.
"We have said that.
Peter will not deny Jesus,.
no matter what.".
"He said so.".
What will Peter say?.
"Yes, what Jesus said will happen,.
will happen.".
Peter said, "No, no, no, not me..
I don't know him..
Not me, not me..
I don't know him.".
They went to another place..
It was cold outside..
It was dark..
They went to a fire to keep warm..
There was no heater in the past..
They cooked with fire..
They kept warm with fire..
They walked around..
Another woman came to Peter..
She asked him,.
"Aren't you the one.
who followed Jesus?".
"Do you know Jesus?".
"Hmm,.
what do you think Peter will say?".

$^{241}$"He denied Jesus last time..
Will he do it again?".
"Yes," she said..
"Not me, not me, not me.".
They walked away..
They saw many people were drunk..
They felt safe..
They went to the main hall..
They hid and covered themselves..
No one would notice..
They just stood there..
Who would know?.
It was a big deal..
Among the men,.
one said,.
"Hey,.
aren't you the one.
who cut off the ear?".
"Yes, it's you..
You cut off the ear.
of the man.
who cut off the ear..
You are with Jesus..
It's you..
Don't deny it..
You are with Jesus..
Hey,.
do you know.
what Peter will say?".
"Who is it?.
It's not me..
It's not me.".
"Ah!".
"Uh-oh!.
What happened?".
"The cock crowed.".
"Oh!".
"It was that moment..
The moment of the cock crowing.".
"Cuckoo!".

$^{281}$Peter stood up..
"Oh, Jesus said.
the cock crowed before..
You will deny me..
Huh!".
Three times..
Did he deny Jesus three times?.
One, the woman..
Two, the other woman..
Three, the man who cut off the ear..
Huh!.
It was that moment again..
Cuckoo!.
Jesus was inside..
He was inside the house..
He looked out of the window..
He was outside the house..
Peter heard the cock crow..
He looked at the window..
They were good friends..
They remembered the conversation..
They looked at each other..
Jesus was sad..
What did Peter feel?.
He was sad and guilty..
Huh!.
What Jesus said would happen,.
we thought it would never happen..
We thought we were not like that..
We would not do that..
But Jesus said it would happen..
It would happen, right?.
Huh!.
I wonder how Peter feels now..
I want you to pause and discuss..
Why did Peter deny Jesus?.
He loved Jesus and followed Jesus..
Why did people ask him.
if he followed Jesus?.
Did he cut off his ear?.

$^{321}$Did he cut off his head?.
No, no, no..
He obviously lied..
Why did he do that?.
Huh!.
I know, you know, Peter knows..
It's not good to lie..
It's not true, right?.
Telling lies isn't good..
It's not honest..
But why did he do that?.
Pause for a moment..
Think about it..
Discuss..
We will come back..
Okay?.
Let's suppose something..
We might think,.
Wow! Peter?.
He's so bad..
He obviously loved Jesus..
Why did he deny Jesus?.
I will bring you into a story..
See how you will feel later..
Or you will change your mind..
Let's see..
What if?.
If?.
You see some kids.
break a window while playing football..
Oh no!.
The auntie at home.
comes out and scolds them..
Who broke my window?.
She sees your friend nearby..
Wow!.
She accuses your friend..
She brings him back..
Scolds him..
And she asks his mom and dad.

$^{361}$to clarify what happened..
Oh! Big trouble!.
Right?.
But your friend.
has a big question mark..
Huh?.
What happened?.
Your friend doesn't know what happened..
But do you know what your friend is doing?.
He is playing football with you..
But the ball you play with.
and the ball that broke the window.
are different..
Actually, they are not..
But why is the auntie.
your friend?.
Hmm....
How do you feel?.
Wow!.
She sees my friend..
Huh!.
She scolds him..
I am so nervous..
Do you think you will say something?.
If the auntie asks you.
"Hey!.
You play football with him, right?".
"Of course!".
"You broke the window.".
"Do you know him?".
"Do you know that kid?".
"Do you play football with him?".
"What should I do?".
"Do you know what to do?".
Or maybe you are like Peter.
who doesn't want to say anything.
but you don't know him..
Actually,.
Peter didn't expect.
to deny Jesus..

$^{401}$We talked about this in the last supper..
He said he would never do that.
because he thought.
there was no reason.
or situation.
to deny Jesus..
But think about it..
Jesus was just arrested..
A lot of things happened, right?.
Peter was so nervous.
and scared..
It's like this story..
The adults asked him.
"Hey!.
Do you know that kid?".
"He broke the window.".
"Are you with him?".
"Oh!".
Or maybe you would say.
"No, no, I am not with him..
I don't know him. Bye bye.".
Right?.
It's possible, right?.
So it could have been you too..
Actually, we could have also been Peter..
From this perspective,.
it's also a perspective, right?.
Actually, when we put ourselves.
into the Bible story,.
we will understand more.
that these characters.
are not just characters..
They are real people..
Real life, real people..
We can understand.
and empathize..
We can imagine.
and feel their feelings..
When we feel empathy,.
we will understand more..

$^{441}$"Yes, Peter loves Jesus so much..
Of course, he won't deny it..
But compared to me,.
I am also scared..
Why am I being pulled into this?.
Or maybe I won't say it..
But I feel guilty inside..
Guilty, right?.
So we can put ourselves.
into the situation.
and feel more of the other person's feelings..
Sometimes, when we lie,.
we are protecting ourselves, right?.
Who has done it?.
I can be very honest..
Who has done something.
that we are afraid of getting into trouble.
and being recognized?.
What will happen if we are recognized?.
What will happen if we are not recognized?.
When I was young,.
I used to eat chocolate cake secretly..
My mom told me not to eat it..
After the party,.
the guests ate it all..
I ate the leftovers..
But that box of cake.
was really delicious..
But I told my mom,.
I promised not to eat it..
I won't eat this cake..
I won't eat a single piece, mom..
Mom went out..
I and this box of cake.
are very close, right?.
I opened the box of cake..
I took a piece of cake from the back..
I closed the box..
No one would know, right?.
But it's not just once..

$^{481}$Wait a minute..
It's me and this box of cake..
The first piece was delicious..
I ate the second piece..
The third piece..
Mom came back..
She was going to take the cake out.
and put it on a plate.
for the guests to eat..
She opened the box of cake..
She put the cookies in the box..
She took more..
I was the third person..
She took out the first piece of cake.
and put it on the plate..
She put the second piece of cake on the plate..
The third piece of cake..
This is the one I stole..
She stole the third piece of cake..
Uh-oh..
She saw it..
It was obviously a new one..
Why didn't she see a few pieces of cake?.
What did she do?.
She walked around..
She was busy..
She read books..
She drew..
She sat down and slept..
She did a lot of things..
Mom asked,.
"Do you know where the cookies are?".
"I don't know.".
"Why did you do that?".
"You knew. Why didn't you tell me?".
She asked again..
"But you were the only one at home with the cookies.".
"Grandma was in the back garden.".
"Grandma didn't like to eat. Where did the cookies go?".
"Uh-oh.".

$^{521}$"Oh, no. What now?".
"I don't know.".
She kept running away..
She looked for a place..
She pretended not to know..
She pretended not to see..
She asked again..
"Tell me honestly.".
"Did you eat the cookies?".
"Oh, no.".
"I was found out.".
"I didn't steal them. I just ate them.".
She said,.
"You know there are a lot of cookies on your clothes, right?".
"I don't know about my face, but I feel very guilty.".
"I feel so guilty.".
Mom asked,.
"You mean you opened the cookies without thinking,".
"put your hand in it without thinking,".
"and took the cookies out without thinking,".
"put the cookies in your mouth without thinking,".
"and chewed them without thinking,".
"and swallowed them without thinking?".
"Right?".
"Uh... yes.".
"Isn't this story funny?".
"It's true.".
"I'm the one who's telling the truth.".
"I ate the cookies without thinking.".
"But...".
"You asked me, 'Hey, Miss Rebecca, will you tell a lie?'".
"I said, 'No, no, no, no, no, no.'".
"I'm a very honest person. I'm very honest.".
"But...".
"Sigh...".
"When I feel in trouble,".
"when I feel dangerous,".
"I might do it.".
"Right?".
"Everyone does it, right?".

$^{561}$"Have you tried it?".
"You can be honest.".
"You can pause and talk.".
"You can be honest.".
"Why?".
"Because God forgave us.".
"We forgive others.".
"Right?".
"God knows our hearts.".
"God knows that when I was a kid,".
"I couldn't resist this delicious chocolate cake.".
"It was so delicious.".
"I didn't have junk food at home.".
"A customer said, 'Wow, your chocolate chip cookie!'".
"I couldn't resist it. I ate it.".
"He understood.".
"I regretted it later.".
"I told my mom, 'Sorry.'".
"Not only did I eat the cake, right?".
"I'm sorry.".
"Not only did I eat it, I also lied.".
"But after I told her, I felt much better.".
"Do you think...".
"Jesus knows.".
"I think my mom knows, too.".
"I'm wearing a cake dress.".
"She knows I did it.".
"But she gave me a chance.".
"Do you think Jesus can give you a chance?".
"Let's talk next time.".
Chinese time..
What is our special animal today?.
There are many animals in our show..
Is it a pig?.
A elephant?.
No?.
Hmm....
A bat..
What is our special animal today?.
A bat?.

$^{601}$No, of course not..
It's a rooster..
Now,.
"公雞" is "rooster" in English..
"母雞" is "hen" in Chinese..
Okay, Chinese time..
We remember that English has different names,.
male and female..
Right?.
But in Chinese,.
we often add "公".
or "母".
or "那"..
That's it..
Like "牛", "cow"..
"母牛".
Oh!.
Interesting..
"母" means "mother"..
So, if you add "母" to "母".
it becomes a female animal..
Wow!.
Then you don't have to remember so many things..
Right?.
I know that.
people often ask.
"Is your dog male or female?".
"公" means "male"..
Right..
"那" means "female"..
Oh!.
Also,.
in the beginning of the Chinese time,.
there was a challenge..
In the beginning of the intro,.
I said that there was a challenge in Chinese time..
"Octagon".
What is it in Chinese?.
What is a triangle?.
Right..

$^{641}$What is a hexagon?.
What is an octagon?.
1, 2, 3, 4, 5, 6, 7, 8, 9, 10..
Octopus.
"八爪魚".
What is an octagon?.
Right. Octagon..
You got it!.
OK. Let's go back to practice..
See what other animals can add "公".
and "母"..
See you next time!.
Bye!.
Hello everyone!.
I have taken out my Bible..
Here..
There is a Bible..
We are very familiar with the Bible.
of our children..
This is the Bible..
Right?.
This one..
It is in cartoon style..
It is for children..
A common theme in this Bible is that.
the picture is usually very big..
The words are smaller..
The story is more like one, two, three pages..
Usually,.
like this one.
and this one,.
each story is separate..
Especially,.
like you said,.
when Daniel and the lions.
jump to the next story..
There are not many stories connected..
It is like a big story..
But actually,.
the Bible is a big story..

$^{681}$The Old Testament and the New Testament are together..
This Bible story,.
Jesus Storybook Bible,.
is connected to a theme..
Each story of the Old Testament.
and the New Testament.
points to Jesus Christ..
So,.
this is the major difference..
But the style is similar..
I want to talk to you about.
the Bible of children..
There are many kinds..
But what do we want to pay attention to?.
We want to pay attention to the illustration..
The text, the words,.
the vocabulary..
How difficult is the word "born"?.
Do children read it by themselves?.
Or do adults read it?.
We can also look at the storytelling..
How deep is it?.
Are there details?.
We know that there are..
This kind of Bible,.
and our adult Bible,.
adults' Bible has more details..
But there are also.
other Bible for children.
that have more details..
For example,.
this series..
Each story is drawn out of a book..
The Last Stepper..
Alone..
There are more than 20 pages..
So the story is more substantial..
They will think that.
the story is really long..
So,.

$^{721}$if you add them up,.
it becomes a series..
Then,.
why.
if there are so many exposures,.
children will know that there are many different perspectives..
Right?.
They will appreciate it more..
You can also see that.
this illustration is different from other illustrations..
The second kind is this one..
It's about facts..
This is graphic, infographic..
This is also like an encyclopedia..
Right?.
You can see that the illustrations.
are more realistic..
Right?.
It really looks like a real person's handwriting..
Right?.
The Bible says that.
there is a river,.
and there is a hill,.
so there are many things happening..
You will know that it really happened..
Wow!.
So many things,.
plus the children,.
make the Bible stories more profound..
They know that.
just like adults,.
if you continue to study the Bible,.
more exciting things will happen..
So we give them more exposure to.
not just the stories,.
but also the perspectives.
of the stories, the telling,.
and the art..
They will know that there are so many.
ways to express.

$^{761}$our knowledge of Jesus,.
our knowledge of God..
It can be abstract,.
abstract,.
or realistic..
It can be like a real person's handwriting..
There is also a type of Bible.
that looks like a comic book..
We see a lot of words in adult Bibles..
We imagine what we see..
We have seen.
painters who draw Jesus,.
or paint with oil paint..
painters who draw Jesus,.
or paint with oil paint..
But there is more than just this representation..
You can use comic books.
to represent Jesus..
The difference is that.
you need to know that.
some things are taken for granted..
You may not know if Jesus really said.
these words..
But the reference is in the Holy Bible..
There is also this type of Bible..
This is actually quite old..
This is also a comic book..
Look at this..
The picture is more realistic..
You can see every expression..
As the story goes along,.
the expression will be different..
Why is this important?.
Because in so many Bibles,.
the children's understanding of the Bible.
is deeper because.
there are different perspectives..
There is a geographic perspective..
It's like this book..
There is an emotional perspective..

$^{801}$You can see the emotional and.
psychological aspects of the Bible..
You can see the emotional aspects of the Bible..
You can feel the emotional aspects of the Bible..
Or you may like data..
Factual perspective..
Historical perspective..
There are also many different.
perspectives and.
different perspectives..
The more you think and.
know, the deeper the understanding..
There is a connection..
I have talked about many different.
Bibles. Children's Bible..
I want you to.
refer to.
the version I talked about.
and.
the version in your.
children's Bible..
What is included?.
Look for the part.
of Rooster..
Look for the part of.
Peter denying Jesus..
Look at.
which Bible has some things.
and look at the illustrations..
So,.
we talked about perspective in the last episode..
This time,.
there is more exposure..
Actually,.
children know this story.
more deeply..
The more connection,.
the better..
The more connection,.
the deeper the understanding..

$^{841}$The more connection,.
the better..
Story..
Look at different stories,.
different Bibles, different illustrations..
Connecting..
Connecting your.
personal story..
Right?.
Have you ever lied?.
Have you ever denied someone?.
Yes, you have..
The connection.
between Peter and you.
is deeper..
OK..
Challenge yourself to do this..
Also, if your children.
like animals,.
also like animals,.
you can play the game..
Go to the Internet.
and look for sound effects.
or animal sounds..
Play with your children..
Guess which animal it is..
Alright..
See you next time..
[music].
\newpage



\section{}
\label{sec:K7tbzr00wo0}
\textbf{佳節倍思恩(二): 住棚節(下) - 陳崇基牧師}
\newline
\newline
連結: \href{https://youtube.com/watch?v=K7tbzr00wo0}{\texttt{https://youtube.com/watch?v=K7tbzr00wo0}} ~~~~ 語音日期: 2024-05-04
\newline
\newline
\hyperref[sec:JpQUnDljfo8]{\small{< < < PREV SERMON < < <}}
~
\hyperref[sec:index_chronic]{\small{[返順時目]}}
~
\hyperref[sec:index_scriptual]{\small{[返順卷目]}}
~
\hyperref[sec:7G7VT1O0Fmo]{\small{> > > NEXT SERMON > > >}}
\newline
\newline
$^{1}$今天中文堂的講員就是Mark Sir.
我們將時間交給他.
好 謝謝Sandy.
今天我們繼續這個系列.
今天第二部分講到.
聖經裡面 特別是舊約裡面的節期.
這些節期都是一些時間的標誌.
一些時間的標誌.
一些提醒.
讓我們知道一年四季裡面.
有些什麼值得我們去紀念.
有些什麼值得我們去想起.
和上帝的關係.
但剛才講到時間的問題.
什麼時候開關.
什麼時候可以去旅遊.
或者讀書的可能在講開學 科學學等等.
我們生命裡面總是有一些時間的標籤.
一些標誌.
一些提醒.
這些可能都是一些人生裡面一定會經歷的事情.
什麼時候會有財政年度 什麼時候會發工資.
什麼時候會有公司裡面有不同的轉變.
不同的季度等等.
但我們需要知道.
聖經裡面一些的節期.
其實是在一年四季裡面.
給我們一些獨特的標記.
讓我們不要忘記上帝在我們生命裡面做的事情.
所以我們不斷在這個系列裡面會講一些節期.
我們曾經提過有三大節期.
就是預節 大概四月左右的時間.
就是我們Good Friday和Ressurection.
Easter Sunday的時候.
然後過了好幾個星期 七個星期之後.
就是所謂的五旬.
或者是以色列巴勒斯坦本地的修國節.
小麥 大麥修國等等.
這個也是一個標記.
告訴神的子民.

$^{41}$這個時候是時候發工資了.
發工資的時候要想起什麼.
他們就會修成他們有一個First Food.
第一次的修國.
然後他們將某一部分奉獻出來.
所以這些是生命裡面的標記.
然後第三個大的節期就是豬棒節.
Feast of Tabernacles 或者Feast of Booths.
或者Sukkot.
上一次我們提到豬棒節.
三大節期 除了有其他不同的小節之外.
都是在一年裡面.
去到給一些Reminder給神的子民.
十月的時候 大概是豬棒節的時候.
譬如我們西歷十月初.
就是以色列的豬棒節來到.
是一些盛大的慶典.
如果我們去看看舊約裡面.
有不同的時間.
神的子民都會去紀念和守豬棒節.
我們選一段經文.
這段經文就是《利欺米記》第八章.
上下文發生什麼事呢.
就是他們被擄回歸之後.
他們回到耶路撒冷的時候.
他們重建了聖殿.
又重建了城牆.
《利欺米記》其實記載他們Rebuild the Wall.
然後他們去到第八章的時候.
每件事都搞定了.
於是他們發生什麼事呢.
領袖就去吩咐那些子民做什麼.
你們當出去.
去上山.
去將橄欖樹.
野橄欖樹.
番石榴樹.
棕樹和各人茂密的樹枝取來.
照著摩西律法裡所寫的去搭棚.
於是百姓就出去了.

$^{81}$取了樹枝.
各人在自己的房頂.
院子.
上帝的殿院裡面.
水門的廣場.
以法廉門的廣場等等的地方.
去搭棚.
從被擄之地歸回的傳會眾就搭棚了.
住在張棚裡面.
從亂的日子.
約書亞記開始的時候.
他入主加衙的時候.
直到這一天.
以色列人沒有這樣行.
他們都大大喜樂.
由約書亞記入主加衙之後.
都沒有試過這麼大規模的.
整個子民的.
整個耶路撒冷來守這個.
豬八戒.
所以你見到他們很體重.
某一些節日.
這些節日去提醒他們.
生命裡面上帝有他的工作.
我不知道你自己怎樣.
你作為一個基督徒.
或者你認識信仰開始.
到現在.
有沒有一些節期.
有沒有一些提醒告訴你.
上帝對你生命裡面.
一些的恩典和工作.
有些人可能記得.
他自己剛初信的時候.
什麼時候決志.
有些人可能還記得.
他什麼時候受洗.
有些人可能記得.
曾經經歷過上帝很厚的恩典.
可能一些死過翻生的經歷.

$^{121}$可能某一些意外.
上帝去保守過你.
或者可能一些禱告.
上帝是聽了你的.
不同的提醒.
其實是告訴我們.
生命裡面是有上帝的恩典.
是我們需要記下的.
豬八戒就是一年裡面一度.
每一次去提醒神的子民.
他們應該要做些什麼.
我們再看這個《利未記》.
《五經》裡面怎樣講豬八戒.
第一天.
豬八戒有七天.
然後跟著不同的發展.
就加多了一天.
第八天.
怎樣也好.
第一天.
你要拿美好樹上的果子.
忠樹枝.
謬物的枝條和柳枝.
再而又說.
你的上帝面前.
歡喜快樂.
七天.
這個就是今天以色列的猶太人.
他們會做的事.
他們有一些後期的傳統的發展.
他們每逢到這個時候.
大概我們西方這個時候.
他們就會預備.
你見到他們手裡面是拿著什麼.
整個東西捆在一起.
叫做Lulav.
Lulav就是包括什麼呢.
Azure.
Azure就是好像檸檬一樣黃色的.
是一個柚子的果子.

$^{161}$然後就會有另外三種的綠色的枝條.
這些樹枝等等.
包括有Palm Branch.
就是那些蜜棗樹.
他們把它們切下來.
然後放在手裡面捆著.
然後有Myrtle.
和有Lulav.
有柳樹枝.
他們還會捆在一邊.
然後放在手裡面.
然後他們祈禱的時候.
就會上下左右前後.
來搖的.
在歡喜快樂的地方去禱告.
是很特別的一種節日.
就是去到Sukkot的時候.
在以色列是另一番風味.
見到他們搭棚之外.
亦都見到他們去到預備這些.
手裡面要拿著的東西.
然後去禱告.
去讚美.
去敬拜等等.
在兩約之間.
就是舊約完了.
到新約之前.
大概有幾百年.
特別在駐彭捷的禮節上.
有很多不同的演變.
就是耶穌時代的時候.
當時的節日.
已經有很多的解釋.
或者有很多不同的內容加進去.
特別是去到耶穌時代.
那些猶太人拉比的文獻.
就會講到他們首節的時候.
要去唱歌.
要去跳舞等等.
駐彭特別這個節日.

$^{201}$就是與雨水有關.
為什麼呢?.
十月初.
我剛才講過了.
今天的月曆.
猶太人的月曆.
十月初的時候.
就是到他們Sukkot開始的時候.
現在正是準備的時候.
這個時候.
大概就是他們下季的時候.
去到最後最後的時候.
我們經常講下季的小尾巴.
這個時候就是小尾巴.
九月底的時候.
在巴西的氣候.
夏天是怎麼樣的?.
是又乾又熱的.
大概四五月開始就開始下季.
所以就沒有雨水.
甚至完全沒有雨水.
有幾個月時間是非常乾旱.
非常熱的.
去到大概九月尾的時候.
十月初.
按照當時的當地的氣候來說.
就開始會雨季來到.
我講雨季的話.
不是說一雨季開始.
就每天都下雨.
正如我們這裡有雨季一樣.
也不會是每天下雨的.
可能一個月裡面.
可能下雨八至九天.
十至十二天已經很多了.
所以不是每天都會下雨.
久不久就會下雨.
但這個就是他們所講的雨季.
也是在《舊約聖經》裡面所講的初雨這回事.
過了幾個月的夏季的澳年之後.

$^{241}$你可以想象.
那些水池已經乾了.
當時在山上住的人.
他們怎麼儲水呢?.
很多時候在屋裡或屋外.
在床突.
然後掘一個洞在裡面.
他們下雨之後.
他們就會將雨水直接放進水池裡.
水是非常矜貴的.
但如果經歷了幾個月的乾旱夏季之後.
那個水池可能已經乾到見底.
所以在這個時候.
到了守住彭節的時候.
就跟雨水很有關係.
在兩約中間.
特別是耶穌年代的時候.
他們演變到一個地步.
他們在守住彭節的第七天.
每一天聖殿的祭司就會下去.
在聖殿的南面.
一個水池叫做西羅亞池.
那裡就會拿著一個金桶.
然後在水池裡裝一罐水.
然後他們就浩浩蕩蕩地.
走回上山.
上到聖殿山.
然後在祭壇面前.
他們就會灑水.
其實灑水 曉水的意思.
就是一個禱告.
希望能夠新幾年來到.
上帝有恩典.
讓這個地方能夠有足夠的雨水.
如果雨水不足夠的時候.
那個乾旱的問題就會繼續延續.
一年是這樣.
兩年是這樣.
三年是這樣.
那就會非常慘.

$^{281}$所以他們就會裝滿一罐水.
每一天然後就拿回聖殿.
然後就倒在耶和華的面前.
來禱告.
往往他們就會讀或者唱這一節的經文.
《倚創亞書》十二章第三節.
「你們必從救恩的泉源.
歡言取水」.
水有一個寓意.
就是寓意著上帝的恩典.
他的救恩來到.
每一年去到罕見的最後.
主板節就要開始的時候.
這個就是最關鍵的時候.
因為能不能夠.
接下來這一年.
我們能不能夠生活.
能不能夠熬得過的話.
就是靠這段時間.
所以就算到今天.
不是猶太人.
阿拉伯人.
他們很多阿拉伯人.
當然阿拉伯人都有不同的宗教.
有基督教.
有些基督教.
也有一些穆斯林教徒.
連穆斯林教徒.
他們都會去到這個時間.
他們雖然不會守住盲節.
因為是猶太人的節日.
他們都會盼望新的一年來到的時候.
這個雨季來到的時候.
他們有足夠的雨水.
等於你盼望能夠出糧.
每逢糧期來到的時候.
你希望能夠充充足足.
或者年尾收佣金.
收花紅.
你都盼望能夠收多一點.

$^{321}$你可以看到這個節日.
是跟那個人的靈性有關.
這個節日就是給一個標記.
讓我們想起我們的靈性生活.
應該要想起什麼.
當上帝賜福給我們.
家人健康.
自己可能康泰.
準時出糧.
市場比較好.
等等.
投資有回報.
我們就要想起上帝的恩典.
不是忘記上帝.
當上帝賜福不夠的時候.
好像不足夠.
疫病也好.
或者世情不如意也好.
好像當時的節日的猶太人.
他們不會覺得是上帝失職.
他們用信心禱告.
去等候上帝.
不是隨隨便便就離開.
駐彭哲其實是一個信心的等待.
等待祈求上帝的恩典來到.
因為信主 信神.
不一定天色上眼.
但要認定什麼呢.
上帝是你和我的供應者.
駐彭哲就是這個意思.
所以為什麼跟雨水有關係.
跟農作物有關係.
同時這個時間也是最後一年的.
最後的收穫.
最後我們要收藏那些果子.
橄欖 棗樹等等.
所以你去到以色列.
你就會發現那些果樹開始成熟.
這段時間就是準備.
去到今年最後的收穫的時候.

$^{361}$出糧的時候就記起上帝是供應者.
生命裡面可能不一定是順利.
但仍然要記起真正供應的是上帝.
這是信心的一個培養.
節期.
每逢過時過節.
四季來到的時候.
想起上帝給我們的好處.
想起上帝在我們生命當中.
祂是誰.
你的生命其實也是要這樣.
我們經常忙碌.
出糧我們就存錢.
出支票我們就存錢.
然後自動主張我們見到有錢.
市場是高還是低等等.
我們經常看著那個銀行帳戶.
看著我們有的東西.
到底有沒有一些定時定候的.
去告訴我們 去提醒我們.
真正的供應者.
能夠頂得住我們生活的.
能夠頂得住我們的命的.
是上帝.
生活可以很忙 工作可以很忙.
你可能一天做24小時.
但有沒有一些提醒記錄給你.
是上帝掌權仍然是供應.
其實在我們基督徒的年曆裡面.
在我們日曆裡面都是一樣.
我們遲些會說.
Sabbath 其中一個.
安息是所謂.
今時今日在基督教的傳統.
新約的傳統.
就變成七天的第一天.
敬拜.
敬拜是什麼意思.
敬拜是想起上帝是那位掌權者.
上帝是那位管理萬有的.

$^{401}$六天之後他創造.
做完事了.
第七天他安息.
就是說他進去聖殿裡面.
他去作和.
創世記第二章.
遲些我們再說.
但我們在我們的日曆裡面.
其實是有不同的記錄.
去提醒我們的.
所以你說如果一個敬拜有意思的話.
怎樣才有意思.
不是別人做的事.
是你自己的取態.
是你自己進來的時候的取態.
唱詩歌 祈禱.
甚至思想上帝的說話.
都是要講你和我的取態.
如果不是的話.
你敬拜一定不會有意思.
對不對.
有時候弟兄姊妹跟我說.
敬拜裡面很感動.
好不好感動是好.
但最主要是你要做的事.
不是你的心動不動的問題.
不是你感覺到觸感.
這麼簡單.
而是你進來的時候你的取態.
你開聲唱詩歌.
去到敬拜去禱告的時候.
你想一想到底上帝是你的誰.
對不對.
日曆的標籤.
除了鑄盤節是培養一個人的靈性之外.
這些節期.
這些時間的標籤.
去到培養我們的靈性之外.
還有些什麼呢.
《約翰福音》第七章.

$^{441}$我們看看新約.
新約這裡發生什麼事呢.
就是耶穌去鑄盤節.
上下文我們很簡單講過.
上下文發生什麼事呢.
就是猶太人去守鑄盤節.
上下文這裡就講到.
耶穌的兄弟.
他們大家都在北面加里尼.
每逢鑄盤節的時候.
作為一個猶太人.
他們都會去耶路撒冷南下.
去到聖殿去守節.
上面有一些爭拗.
上文就講到.
耶穌的兄弟.
他的親兄弟.
就不信他.
去到懷疑他.
接著他們就去守節.
耶穌就說我不會和你去.
因為我的時間還沒到.
但其實耶穌都會去的.
不過他私底下自己去.
所以第七章之後.
就交代耶穌去到.
猶太地的耶路撒冷.
他的兄弟也同時在耶路撒冷.
去守鑄盤節.
和其他很多的猶太人一樣.
你就可以想象那個熱鬧的氣氛.
有如我們華人的農曆新年一樣.
當然之下沒有什麼氣氛.
但是我講華人比較多的地方.
耶穌自己私底下去的時候.
在那個鑄盤節的最後一天.
七章第37節.
節期的最後一天.
就是最隆重的一天.
耶穌站著喊出話.

$^{481}$人若渴了.
到我這裡來渴.
信我的人.
就如經上所記.
從他腹中.
將流出活水的江河來.
耶穌說是指信他的人.
要受聖靈.
那時還沒有此下聖靈.
因為耶穌還沒有得到榮耀.
七天的最後一天.
最隆重的一天.
耶穌就說.
人若渴了.
就來喝我所給的水.
我剛才講到鑄盤節.
水.
是和水有關.
人們已經渴到死了.
連喝水都要少一點.
因為怕不夠.
就在這個時候.
最隆重的那一天.
猶太人拿著金筒.
祭司拿著金筒.
然後去到聖殿那裡.
有一個很盛大的祈願儀式.
去祈禱.
去祈求上帝新的一年給我們有水喝.
耶穌就說.
人若渴了就到我者那裡.
你們這班人都已經渴得很厲害了.
我想告訴你.
有什麼水.
是你真正需要的.
到底真正要的水是否瓶裝那些.
杯裡面那些.
喝進口的那些.
還是有其他東西.
他說真正的水.

$^{521}$真正可以使你滿足的.
就是永生.
是耶穌給我們的神的生命.
這個是聖靈.
救若至生命.
盼望的是什麼.
就是水的濃度.
代表的就是聖靈.
就是上帝的應許能夠實現.
上帝的生命能夠在我們生命當中.
能夠有永生.
這個才是真正可以滿足我們的.
我們可以每年都祈求有水.
可以每年到糧期的時候.
我們盼望有糧單.
能夠出糧.
能夠準時.
我們希望這個年度成功.
這個成功的報告.
這些可以滿足我們.
但只不過是肉身.
但真正能夠滿足我們的.
是永生.
聖靈.
是我們接受主的時候.
你有神的生命.
這個才是真正能夠滿足我們的.
我不知道你搭長途機的時候.
你會做些什麼.
搭長途機.
可能看電影.
去看那些爛劇.
那些不好的電影我們看幾次.
上兩個星期我回來的時候.
搭長途機回來的時候.
我自己也要看電影.
因為我做不到事.
精神狀態不能夠集中.
我算了.
我去家裡找些電影看.

$^{561}$我又看一次《托根》.
因為我不想看些悶的電影.
我就看些動作戲.
看完之後還是有時間.
我就按按其他有什麼電影好看.
我就看《梅艷芳》.
又看一次.
我覺得真的很好看.
講到一個港人.
香港人為了香港做了不少事情.
但看完之後我有些唏噓.
如果你看過《梅艷芳》.
你就知道.
我不劇透.
最後的時候.
特別是他唱《夕陽之歌》的時候.
我覺得真的很…….
很不知怎樣.
你又睡不著的時候在飛機上.
你又覺得看完之後.
不是很舒服.
留不住最後的時候.
《夕陽之歌》就是講這個.
每樣東西都會過去.
整套劇想起很多過去的感情.
地上擁有的東西.
到最後都是留不住.
很淒美應該這樣說.
如果沒有了永恆的拍.
是不是就是這樣.
你甘心嗎.
耶穌說什麼.
人若渴了.
你要渴求一個豐盛的生命.
人若渴了.
要到我這裡.
喝我給你的水.
真正的水.
真正的滿足.
豬扒的時候.

$^{601}$耶穌就去提醒猶太人.
或者在我們生命當中.
當我們很渴的時候.
耶穌就去提醒我們.
真正你需要的.
其實是永生.
因為如果沒有永生.
一切都是零.
你多麼風光也好.
都是零.
最後都是《夕陽之歌》.
我再講一個猶太人拉比的故事.
跟豬扒有關的.
這個故事就講到以前在英國.
好幾年前.
這個大拉比我曾經提過的.
Jonathan Sachs.
他在英國是一個Chief Rabbi.
他講到他們去守豬扒節的時候.
他想起以前年輕的時候.
剛結婚.
然後自己有家室.
豬扒節的時候.
每一家要做一個扒.
所以剛結婚的時候.
他就要做一個Sukkah.
做一個扒.
因為他是一個學者.
拉比是一個教師.
他又不太懂得做釘子和木材.
他不太懂.
在英國倫敦裡面.
他又認識了一些朋友.
有一個朋友很懂得木工和材料.
很靈活的.
這個朋友就跟他說.
我去買Sukkah的材料.
你想不想跟我去.
他就跟了他朋友去.
一去到那些家庭地庫的地方.

$^{641}$就是買木材和材料的地方.
這個朋友就在袋子裡面拿出一個.
很詳細的一個計劃.
有那些圖片和藍圖.
要什麼材料.
要什麼釘子或者什麼hinges等等.
全部的門和窗等等.
他都有了.
於是他就跟著這個藍圖去買.
這個拉比就不知道怎麼辦.
因為他又不方便.
他就隨便買了幾塊.
好像一些卡片一樣.
然後他就回家去做他的Sukkah.
兩個人回去他就做一個很漂亮的Sukkah.
可能有些這麼漂亮的Sukkah.
有燈有裝飾等等.
他自己就弄了幾塊木頭.
然後釘在那裡.
好像那些爛木屋一樣.
他說他朋友的帳篷那個Sukkah很華麗.
那個Sukkah他覺得可能在曠野40年都可以.
住在那些地方40年80年都無所謂.
120年都無所謂.
因為太漂亮了.
大家第一天弄好這個Sukkah之後.
那天晚上就狂風大雨.
翻風落雨.
那天晚上過了之後第二天.
他朋友打電話給他.
他說我的Sukkah塌了.
我說你那個一定是塌了.
Johnathan Sacks就回答他說.
沒有啊.
我那個還健座.
還健座沒事.
他朋友就很驚訝.
他說有沒有理由.
我那個會塌了.
然後你那個就沒事.

$^{681}$然後立刻飛車去他家裡去看看.
他的Sukkah到底發生什麼事.
正如我所說.
他弄的Sukkah本身是一個很簡陋的Sukkah.
當他朋友去檢查他的Sukkah的時候.
他就發現.
那個棚原來是一口釘.
是釘在他家裡的樹上.
在他家的牆上.
原來就是那一口釘.
令到他整個Sukkah能夠安穩.
經過風吹雨打都沒事.
他的猶太人朋友就醒悟了.
他說不明白.
他說多有周詳計劃也好.
多精打細算也好.
多華麗的Sukkah也好.
如果沒有那一口堅固的釘.
一切都完蛋了.
風吹雨打.
一切生命都會倒塌.
那個朋友說那一口釘.
就是Nail of Faith.
就是信仰信心的釘.
在你生命裡面.
有多周詳計劃也好.
沒有了那個Nail of Faith.
一切最後.
生命.
風吹雨打來到的時候.
就會倒塌.
你想想你經營你的生命的時候.
有沒有留意這件事.
我們日日夜夜.
忙著去建立我們的Sukkah.
建立我們這個暫時的一個帳篷.
我們有小孩的.
我們盡心盡力去做.
為他預備.
為他準備很多很多的東西.

$^{721}$想不到的我們都想好了.
為了我們日後的生活能夠舒服一點.
我們的課運期等等不同的投資.
我們全部都部署好了.
最周詳的計劃.
最精打細算.
最華麗的.
但到底是不是缺乏了.
就是信心的那一口釘.
耶穌說真正的那一口釘.
就是你在神裡面有生命.
你知道鑄蟠節的時候.
在猶太人的會堂.
他們會讀什麼聖經.
舊的裡面就是傳道書.
傳道書最後我們用這一節來結束.
傳道書十二章十三節.
講完所有的生命裡面.
是多麼的虛空之外.
如果沒有永生.
一切都是虛空.
一路講講十二章.
作者說什麼.
這些事你都聽見了.
結論就是要敬畏上帝.
謹守祂的誡命.
這是人當盡的本分.
因為人所經營的.
連一切隱密的事.
或善或惡.
上帝都要受.
我們同心低頭祈禱.
讓我們生命的經營有個後頂.
主我們來到你面前.
我們求你來提醒我們.
讓聖靈這個永生的聖靈.
提醒我們現在有的神的生命.
有這個信心的釘子.
讓我們釘得穩釘得盡.
我們生命的基礎.

$^{761}$盤石在基督裡面.
主幫助我們.
讓教會讓你自己的子民.
被提醒我們生命裡面所經營的一切.
我們所搭的棚.
主我們有沒有在你的永生裡面.
求你讓我們不要忘記你.
你對我們的恩典.
你對我們的救恩.
是多麼的貴重.
多麼的重要.
一個永恆的救恩.
讓這個永恆的生命.
成為我們一切手所經營的東西.
那個根基.
以至風吹雨打來到的時候.
我們知道.
你仍然是我們的救恩.
你仍然是我們的上主.
我們敬拜讚美,事奉你.
祈禱,奉傳明教.
願神祝福各位.
希望下一個星期能夠跟大家一起.
\newpage



\section{}
\label{sec:7G7VT1O0Fmo}
\textbf{Everlasting Praise 3 (Exodus 25-31) - Rev. Mark Chan}
\newline
\newline
連結: \href{https://youtube.com/watch?v=7G7VT1O0Fmo}{\texttt{https://youtube.com/watch?v=7G7VT1O0Fmo}} ~~~~ 語音日期: 2024-05-07
\newline
\newline
\hyperref[sec:K7tbzr00wo0]{\small{< < < PREV SERMON < < <}}
~
\hyperref[sec:index_chronic]{\small{[返順時目]}}
~
\hyperref[sec:index_scriptual]{\small{[返順卷目]}}
~
\hyperref[sec:yIgKl3jW8_8]{\small{> > > NEXT SERMON > > >}}
\newline
\newline
$^{1}$All right, guys, we will continue our series on worship..
And remember we read and studied Psalm 8..
Remember that?.
Psalm 8?.
No?.
Okay, you guys need to improve, okay, in your learning..
Okay, Psalm 8, I know you know, right?.
I know you know..
Psalm 8 is a creation psalm,.
and teaches us how to worship God..
And last time, in this series,.
we were at Genesis 2, one and two,.
and we were talking about the Garden of Eden, right?.
And we say Eden is like a sanctuary,.
it's like a temple, right?.
Where is the entrance facing?.
Which direction?.
Yes, east, yeah, there you go..
And who is guarding the entrance?.
Who is that guy?.
Cherubim, right?.
Cherubim, couple of cherubim..
And by the way, cherubim is plural..
So two, at least two,.
they're guarding the entrance of the sanctuary, all right?.
And Eden was like on mountaintop,.
because rivers are coming out of Eden.
to cultivate the lands nearby..
And also, there are some wordings.
in chapter two of Genesis, or chapter three in Genesis,.
that reminds us this is a sanctuary..
Because whenever the Bible talks about a sanctuary,.
a temple, these wordings appear..
Three, eight, this is after the fall of humankind,.
after Adam and Eve rebelled against God..
And it says three, eight,.
when the cool evening breezes were blowing,.
the man, humanity, representing the whole humanity,.
and his wife heard the Lord God walking about in the garden..
Okay, so they were hiding from the Lord God.

$^{41}$from the tree, in the trees..
And hear the wording, walking about, walking alongside,.
walking in the midst of them..
That is a special vocabulary in the Bible.
that describes God's presence..
God is walking in the midst of His people..
And that kind of wording appears so many times.
in the Old Testament that talks about the temple.
or tabernacle or wherever you find the presence of God..
So you can see from the opening chapters of the Bible,.
Genesis one, two, and three,.
you see worship is a main theme..
People worship God in the presence of God..
People can walk with God, can have fellowship with God..
And that's what worship is all about..
The presence of God is with His people..
God is walking about, walking in the midst of them..
So Eden is like a sanctuary..
So if you have time, you can flip to,.
turn to the last chapters of the Bible..
The last book of the Bible being,.
what's the name of the last book of Bible?.
Revelation, yes, Revelation..
The last chapters, chapter 21, 22,.
talks about new creation..
It talks about future, our ending..
You know, we are heading toward that future..
And that future is gonna be a glorious future..
And I wanna thank George.
for preaching a very good sermon last week,.
reminding us the whole creation is going forward..
That direction, God is gonna renew everything..
God is going to renew the whole cosmos..
And the theme of worship appears again..
Like 21, three, the book of Revelation..
John, the apostle, being lifted up into heaven,.
and he saw or he heard something that is otherworldly..
I heard a loud shout from the throne saying,.
"Look, God's home, the home of God,.
"is now among His people.".

$^{81}$God's home, by the way, this is new living translation..
I guess it's easier for most of us..
The home of God, the presence of God, is among His people..
He will live with them and they will be His people..
God Himself will be with them..
Again, that kind of wording,.
walking in the midst of them, live with them,.
dwell with them, and God will wipe every tear.
from their eyes and there will be no more death.
or sorrow or crying or pain..
All these things are gone forever..
Talking about the new heaven and new earth,.
the new cosmos, the new creation..
That's where we are going..
That's where we're looking forward to..
That gives us hope..
But again, if you look at the wordings here,.
worship is the goal..
The new creation is full of worship..
The new creation is full of worshiping God.
and living in His presence..
So as you realize, the Bible is bookended.
by two portions, the opening chapters.
and the closing chapters, and it's all about worship..
Bookend by the theme of worship..
So the Bible, you read from cover to cover,.
is from worship to worship..
Now, of course, if you trace the grand story of the Bible,.
after the fall of humankind,.
after the rebellion of Adam and Eve,.
the humankind rebelled against God,.
they lost the presence of God.
and they are kicked out from Eden, right?.
Got it?.
Kicked out from Eden..
However, those who know God and those who believe in God,.
they continue worship..
So these kinds of wording are going to describe.
some of those people who continue to be faithful to God,.
even though they were the minority..

$^{121}$523 Genesis, Enoch..
There's a guy called Enoch and he lived for 365 years..
By the way, 365 is a symbolic number, okay?.
All right, he's not like a tree, all right?.
Never die, all right?.
So you have to read the Bible according to its own term..
It's written not to us, but for us, right?.
It's written to ancient people.
and they use symbolic numbers, okay?.
But they are real numbers, okay?.
Real symbolic number, but they have some symbolic meaning..
So Enoch lived 365 years..
Verse 24, walking with God, right?.
In close relationship, in close fellowship with God..
And then one day he disappeared, God took him away..
All right, how cool is that, right?.
Anyways, and then at the end of the book of Genesis,.
when Jacob was blessing his children, all right?.
His end of, come to the end of his life.
and Jacob blessed his children..
He blessed Joseph and said, look at how he described God..
May the God before whom my grandfather, Abraham,.
and my father, Isaac, walked..
He walked before them, right?.
Abraham walked before them, or walked before God..
Isaac walked before God..
And Jacob also walked before God..
So these people who worship God, who are faithful to God,.
they are always described as walking with God.
and God is walking with them, right?.
So the God who has been my shepherd.
all my life to this very day, right?.
So even though mankind was being kicked out of Eden,.
kicked out of the presence of God,.
but God is still walking alongside of them.
for those who are faithful to him,.
for those who still believe in, right?.
So it's described them in such verdicts, right?.
And then we're gonna trace the Bible from the book of Genesis..
Remember where the people of God end up.

$^{161}$by the end of book of Genesis?.
Where were they?.
They were not in Chicago..
Where were they?.
They were in Egypt..
Got it?.
All right, so many things happen,.
but they end up in Egypt for many years..
And then some days ago, or some day later,.
God appeared to Moses..
So you got Abraham, Isaac, Jacob, Joseph, and the 12 tribes,.
and then you come to Moses..
Many years later, Moses..
Moses was in exile,.
and he was shepherding his flock in the wilderness.
in the book of Exodus chapter three..
It's a long passage,.
but we're gonna take a look from the perspective of worship..
One day, Moses was tending the flock of his father,.
and we're gonna skip to into the wilderness,.
and into where?.
He came to Sinai, right?.
You got that word?.
Sinai or Horeb, okay?.
Both referring to the same mountain,.
the mountain of God, right?.
Mountain of God..
Again, sanctuary..
There, the angel of the Lord appeared to Moses.
in blazing fire from the middle of a bush,.
and he stared in amazement..
How cool is that, right?.
Because the bush was engulfed in flames,.
and it did not burn up..
Do not come any closer, the Lord warned him..
Take off your sandals..
Take off your sandals,.
for you are standing on the holy ground, right?.
So this place is holy because God's presence is there..
Take off your sandal, and I will be with you,.

$^{201}$and this is the sign that I'm the one who sent you,.
sent you back to Egypt to bring my people out of Egypt,.
and what's the goal?.
What is the destination?.
You will worship God in this very mountain,.
Mount Sinai, right?.
To do what?.
To worship God, right?.
So God saves his people from the land of slavery,.
from Egypt, for the purpose of serving him.
and worshiping him, and as you know,.
the book of Exodus, there were tons of things.
happened there, 10 plagues, right?.
And all this crazy thing happened,.
and God judged Pharaoh and the nation of Egypt,.
and saved his people out of Egypt..
So they crossed the Red Sea,.
and they crossed the great and fearful wilderness,.
and they finally, they arrived at the foothill.
of what mountain?.
Mount Sinai, right?.
Mount Sinai..
To do what?.
To worship him, right?.
God saved his people out of Egypt.
in order for them to worship him,.
and there at the foothill of the Mount Sinai,.
God gave them the 10 commandments..
God gave them the 10 commandments,.
and the 10 commandments is kind of like.
the constitution of ancient Israel,.
just like the United States Constitution..
Those are the basic principle, lay out 10 principles.
that how they are going to conduct their lives..
Now they are the new people of God,.
how are they going to live?.
How are they going to live according to the principles.
God has given them?.
So that's the 10 commandments..
And then at that moment, God told Moses,.

$^{241}$I'm going to live and dwell among my people..
And the symbol I'm using so that they can see.
I'm living among them is tabernacle..
Tabernacle, right?.
Heard that word, right?.
Tabernacle..
What is tabernacle?.
And this is the topic of the day..
I'm going to talk to you about tabernacle,.
and what is the tabernacle for,.
and what does it mean, right?.
So let's begin by understanding the word tabernacle..
Okay, tabernacle..
Various translation, but it comes from Hebrew,.
meshkan, all right?.
Tabernacle is English, of course..
And it translates the Hebrew word meshkan, meshkan..
And meshkan is a noun, comes from a Hebrew verb, shakan..
So Hebrew is a very funny language..
You add something in the front,.
or you add something at the end,.
and then it becomes another form, right?.
Here's one of the example..
Shakan is living, dwelling, right?.
And you add an M, meshkan, right?.
Did you get it?.
Meshkan, it becomes a noun..
Means a tent, all right?.
It means a tent..
Living in the midst of its people..
Is it really difficult?.
It's not that, right?.
It's not that difficult..
So shakan and meshkan, right, sound the same,.
and that means tabernacle, right?.
God is making his tabernacle among his people, right?.
Recently I learned that there is something called cramping..
You heard that?.
Have you ever been to camping, camping trip?.
Do you like it?.

$^{281}$Do you enjoy it, camping?.
Like pitching a tent, right?.
Yeah..
And then nowadays they have this glamorous camping.
called glamping, right?.
So they sound the same, right?.
Camping and glamping, right?.
So that's what I just learned..
You know, I didn't know that..
I used to just go camping, right?.
And it's pretty, pretty tough.
because you have to go for a public bathroom..
You have to share the bathroom together,.
and it's very dirty and filthy and all that, right?.
So nowadays, you know, people are so spoiled..
We are so spoiled..
Americans are so spoiled..
And we invented this glamping, right?.
What is that?.
Glamorous camping..
That means you have your own bathroom..
You live in a container..
It's very, very comfortable..
There's AC, air conditioning, and all that, right?.
It's just living in, like living in a hotel, right?.
But you live in nature.
because the container is in the nature..
So this is camping, all right?.
It's rain and very dirty, right?.
This is glamping, right?.
You know, it's a luxury kind of camping, right?.
Which one do you prefer?.
Do you prefer camping or you prefer glamping?.
Glamping, you know, you gotta pay more, right?.
What I'm saying is camping and glamping.
is like Mishkan and Shakan, right?.
They sound pretty much the same..
But what it means is that God is pitching his tent,.
glamping, God is glamping in the midst of his people, right?.
It's a tent of God, right?.

$^{321}$It's the tabernacle, right?.
And in Southern Israel, there's a place, a national park,.
and you can drive there in Southern Israel, in the Negev,.
and you can visit this park and you can walk around..
And in this park, there's a one-to-one replica.
of tabernacle, one-to-one ratio..
As you can see right there,.
it's a rectangular shape, glamping, right?.
A tent, kind of a boundary,.
two-to-one ratio,.
rectangular boundary..
And inside there, there's also.
an other smaller rectangular shape,.
you know, with the curtain on top..
And then inside is a square, it's the Holy of Holies..
So what you see here with all the curtains.
blocking off the boundary is the outer court of tabernacle..
And inside the court of the tabernacle,.
there's another, you know, smaller rectangular shape area,.
and you can go into, it's also two-by-one,.
it's called Holy Place..
And inside Holy Place, there's a square, perfect square,.
Holy of Holies..
So let's walk inside, right?.
So you walk inside, you're standing inside the court,.
inside the courtyard of tabernacle,.
and you see the altar, right?.
The altar where you burn the sacrifice, right?.
Now, altar, practically speaking, is a barbecue place,.
and you burn all the sacrifices, right?.
And you also see the round thing, the round dish, right?.
You see that?.
It contains water, it's a water basin,.
and for priests to, you know, wash themselves.
and make sure they're cleansed before they enter.
or doing some rituals in the tabernacle..
So we are in the courtyard..
Now you wanna go inside?.
You go inside, this is a two-by-one smaller rectangular,.
and you see two guys here..

$^{361}$You know, that guy leaning in the back.
should be one of the priests,.
and he was doing something here..
And also another guy in blue robe,.
he's the high priest..
You see the brass plate in front of him?.
There was 12 stones, right?.
So the priests, you know, they can go into the holy place..
And then you see three different kind of furniture items..
On the left is the menorah, all right?.
It's a seven-branch menorah, looks like a tree of life,.
looks like a tree..
On your right is a table called the showbread table..
So every day they would bake some new bread,.
and they can put it on the table..
You know, that symbolize God is hosting his people, right?.
You are being invited into the home of God..
This is the clamping home of God..
You come in and you see on the left is light.
from the menorah, on the right is the hospitality.
provided by the king of the universe, showbread..
In the middle of it is another order.
called the incense order..
The incense order where they burn the incense,.
and the incense, the smoke will go up, right?.
That will represent the prayers of the people, right?.
That's the holy place..
So being right there, right where you are standing now.
is inside the home of God himself, the tabernacle..
And then you see the curtains?.
Those figures on the curtain, who are they?.
Who is guarding the holy place?.
Yeah, yeah, cherubim, there you go, cherubim..
They are supposed to symbolize cherubim..
Do you wanna go inside?.
No?.
You can't go inside, right?.
Because when the tabernacle was in the midst.
of the ancient Israelites,.
no one can go into the holy of holies..

$^{401}$No one can go behind the curtain..
You can go in only one day per year,.
and it's represented by the high priest, Aaron,.
and his descendants..
The high priest will go in one day a year on Yom Kippur..
Day of Atonement..
And he's gonna represent the whole people of Israel.
and go into the holy of holies and to atone the people.
so that the people will find forgiveness of sins..
So only the high priest can enter into the holy of holies,.
which is a square, perfect square,.
meaning holy of holies is a perfect place..
Do you wanna go inside?.
For this model, you can go inside, all right?.
All right, so you kind of go through the curtain.
and inside there's only one furniture item,.
and that furniture item is called the Ark of the Covenant..
All right?.
It's plated with gold, real gold, all right?.
The Ark of the Covenant is actually a box, a chest box,.
and they store one of the items inside the box..
And on top of it, there are two creatures, and who are they?.
The cherubim, yes, the cherubim..
There's gotta be cherubim everywhere, right?.
Cherubim..
So they're guarding the holy of holies..
Now the Ark of the Covenant is actually God's throne..
He sits on it..
That is his chair, that is his oval office, right?.
It's right at the center of the White House..
That's where the president is..
That's where the king of the universe is, right?.
Of course, there's no idol,.
there's no visible symbol for God.
because God is invisible, right?.
But inside this box, you find two copies of stone tablets..
All right?.
And what is written on it should be,.
part of the words should be the Ten Commandments, right?.
There are two copies..

$^{441}$And you know why?.
Because in the biblical time,.
when two nations are making covenant together,.
when two parties are making covenant together,.
they will come together and talk about.
the terms of the covenant..
And once they agree with the terms of the covenant,.
with all the blessing and curses for obeying it,.
you get the blessing..
If you disobey, you get the curse..
And that is gonna protect the covenant relationship.
between two nations, two parties..
When the two parties agree with the terms of the covenant,.
they will make two copies in stone.
or in some doable material..
They will write down or engrave the terms on the tablet..
So you get a copy, I get a copy..
It's an exact copy, right?.
And you're gonna bring your copy back to your country,.
and you're gonna deposit the tablet.
into your main sanctuary, your main temple..
So you're gonna put your copy in your main temple,.
in your holy of holies,.
and I'm gonna put my copy in my sanctuary,.
the biggest temple of my country..
I'm gonna put it there in my holy of holies..
And year after year, we're going to renew the covenant,.
and we're gonna read the covenant terms,.
and make sure we're gonna obey it.
and maintain our relationship..
Now, for Yahweh, for God, and Israel,.
when they make covenant together,.
for the people of Israel, where is their sanctuary?.
It's a tabernacle, right?.
And that's why you find one copy in the holy of holies,.
in the Ark of the Covenant, right?.
And for Yahweh, for God himself, where is his temple?.
Also in tabernacle, right?.
Also in tabernacle, the earthly temple is in,.
it's a tabernacle..

$^{481}$So Yahweh's copy, God's copy, and Israel copy.
are being deposited into this same box..
So you find two exact copies..
That means God and Israel, God and his people.
have a covenant going on, right?.
And that covenant will protect their relationship..
And the people are being called to worship God.
because they are his people, and God is their God..
And they are going to honor God and worship God.
and serve God with God in their midst..
Now we can come, this is the end of the tour, all right?.
This is the end of the tour..
But we're gonna go to the text, go into the Bible,.
and see how the Bible describes.
the building of the tabernacle..
In the book of Exodus, chapter 25 to 40,.
is the long portion of the text,.
talking about the plan and the building of tabernacle..
From chapters 25 to 31, God gave Moses the plan..
It's an architectural blueprint, how to build a tabernacle..
Like how long it is, how wide it is,.
and what you're gonna use..
You're gonna use gold and silver.
for all the things that you're gonna put together..
That is the architectural plan, 25 to 31..
And then 35 to 40, the ending chapters,.
describes the actual building of the tabernacle..
All right, so we're not gonna read it..
But here's the interesting things..
From 25 to 31, talking about the blueprint of tabernacle,.
you expect the Bible will continue.
to talk about the actual building of tabernacle, right?.
However, in between 32 and, 32 to 34,.
you find an episode, a very sad episode..
The people of Israel in the wilderness,.
they worship a different God..
They worship their own idol, golden calf..
However, at the end of this episode,.
God is saying, "I'm gonna still dwell among you.".
And that is a lesson of grace..

$^{521}$We don't deserve the presence of God in the midst of us..
Only because of God's grace,.
that God is going to have his presence among us,.
so that we can worship him..
Have you ever thought of that?.
We can worship God, it's not because we deserve it..
It's not because we earn it..
It's God made a way that is possible.
for us sinners to worship him, right?.
So it's very interesting how the design of the chapters..
After the blueprint, the Bible stops right there..
Hey, remember, God can live among us..
It's not because we earn it..
We're sinners, we don't deserve it..
God, by his sheer grace, he can make his presence among us.
and we can worship him..
And then you see the building of the tabernacle..
The tabernacle represents the presence of God.
in the Old Testament..
That God is living in the midst of his people.
and his people will worship him..
And I wanna bring out two points.
about worship and tabernacle..
And there are two main points..
When the people, when Moses received the plan,.
the blueprint of tabernacle,.
look how the people are being commanded to do..
The Lord said to Moses,.
"Tell the people of Israel to bring me sacred offerings..
Accept a contribution from all those hearts.
are moved to offer them.".
And here's the list of sacred offerings.
you may accept with them..
Gold, silver, bronze, all sorts of expensive items..
You're gonna bring them as offerings.
to build the tabernacle..
Verse eight, God said,.
"Have the people of Israel build me a holy sanctuary.
so I can live among them.".
You must build this tabernacle and its furnishings.

$^{561}$exactly according to the pattern I will show you..
God is saying, "I will have the people.
to build me a sanctuary, build me a camping tent.
so that I can camp among them, I can live among them.
so that they can worship me and I can be their people.".
And you're gonna do so according to the plan.
that I have for you..
Here, "Tell the people of Israel to bring me,".
not the second best, not the leftover,.
"Bring them the best things that they have.".
They have gold and silver because when they left Egypt,.
they plundered Egypt, they conquered Egypt,.
and the Egyptians gave them all these good things, right?.
Treasures, gold and silver..
So they're carrying all those gold and silver items.
with them out of Egypt..
And God is saying, "You need to offer them to me.
to build this tabernacle so that I can build,.
I can live among you.".
You see the principle here is generous offering..
How blessed we are that we can build the presence of God.
in the midst of us..
God does not need the people of Israel.
to build him a tabernacle..
God doesn't need that..
God can send angels to build a tabernacle,.
but he wants his people to participate.
in the act of worship..
God want his people to participate in worship..
And that's why we come to worship,.
not because we deserve it,.
but God is willing to live among us.
and give us the opportunity to sing praises to him,.
to worship him and to serve him..
God does not need you or need me to worship him..
It's all by grace..
And that's why when we build our church,.
when we build the body of Christ,.
the presence of God,.
we are called to make offerings, generous offering..

$^{601}$Have you ever thought why we don't do offerings?.
Sometimes we just don't believe that God is good, right?.
Sometimes we just don't believe God is generous, right?.
But the Bible show us once and again, God is generous..
God is so good to you and to me..
And that's why we can offer our generous offering.
to build his presence, to build the tabernacle..
Just like those days, the people of Israel were called.
to bring in the gold and silver and all these good things..
They can build the tabernacle.
so that God will live among them..
And secondly, is to be holy..
Leviticus 28, God said, "Keep all my decrees.
"by putting them into practice,.
"for I am the Lord who makes you holy.".
Now, what does it mean?.
Holiness or being holy means separated,.
that you're different, right?.
That's what the basic meaning of holy is..
Be holy means be different, be separated.
from all the nations, from all those people.
who don't worship God himself..
And you are different, you got to be holy,.
and you got to keep the words of God.
and put them into practice..
And that is the principle of worship..
We come to worship, it's not just about singing praises..
We got to conduct our lives according to the will of God..
Because God is living among us when we observe his law,.
when we follow his words willingly..
Two principle out of tabernacle..
Now, of course, you're gonna say,.
"Okay, we don't need a tabernacle today," right?.
You can go to southern Israel.
and visit the one-to-one replica..
Do we need a tabernacle?.
We don't..
Why?.
Because we are in the era of New Testament..
Christ has come, Christ has come..

$^{641}$And Christ is the perfect priest..
He is the perfect high priest,.
and he offer himself as the perfect sacrifice..
That's why we don't need the lamb,.
offer the sheep as burnt offering today..
Because the Old Testament era is gone,.
and the way of worship is different..
In the Old Testament, people worship God.
in the tabernacle or around the tabernacle..
In the New Testament, we worship God in Christ..
It's not in a place, it's not in a spot on the map,.
but in Christ..
So in the New Testament, there was no more animal sacrifices..
There was no more Aaron the high priest..
There was no more tabernacle, physical tabernacle,.
because Christ has come..
But I wanna show you this verse, John chapter 1, 14..
Now, as I read John chapter 1, 14,.
you try to guess which word is related to tabernacle..
All right?.
John 1, 14..
So the word, speaking of Jesus,.
became human and made his home among us..
He was full of unfailing love, grace, and faithfulness..
And we have seen his glory,.
the glory of the Father's one and only Son..
Where you see the word related to tabernacle here?.
Home, right?.
Made his home among us..
The word became flesh..
Jesus Christ came as God incarnate..
Made flesh and made his home among us..
Actually, the original Greek term is tabernacle..
God made his tabernacle among us,.
or God tabernacle among us..
God's glamping among us, camping among us,.
living in the midst of us..
So that when we see Jesus, when we see the record of Jesus,.
we see the glory of the Father, the glory of the true God..
And that is the beginning of the Gospel of John,.

$^{681}$chapter 1, 14..
And then at the end of Gospel of John, chapter 17,.
after the Last Supper, Jesus prayed a prayer,.
and we call that prayer the High Priest Prayer..
He was like a high priest..
And one of the verses that describe.
Jesus' High Priest Prayer was 1722..
And Jesus said, "I have given my disciples.
"the glory you gave me, so they may be one as we are one..
"I have given them the glory, I have shown them the glory,".
which is the glory of the cross,.
that he was willing to die on the cross.
to redeem all of us..
"I have given them the glory, I have shown them the glory,.
"so that they may be one as we are one.".
He's like the high priest,.
bring in the glory of God.
to his disciples, to us, to the church, right?.
Now, we don't need the tabernacle because Christ came..
Christ is the tabernacle, that's why we worship in Christ..
That's why we pray in Christ..
Now, however, there are all sorts of sanctuary terms.
or temple terms to describe the people of God.
in the New Testament..
And I wanna close with this verse,.
1 Corinthians 3.16..
Now, you know the letter of 1 Corinthians.
is written to the church in Corinth..
And we know the home churches in Corinth.
had a lot of problems..
They have divisions among them,.
they have separation among them,.
they have arguments in the midst of them,.
and they are fighting for different things..
And it's a very problematic church..
And the pastor, Paul, wrote this letter to them.
to warn them and to teach them.
how to live as a worshiping community..
In chapter three, verse 16,.
he reminds the Christians in Corinth,.

$^{721}$he said, "Don't you realize that all of you together.
"are the temple of God?".
Wait, I thought the temple of God is stone building..
No, Paul said, "Because Christ lives in you..
"You guys, you together, when you meet together,.
"when you are a family of God, you are the temple of God.".
And what is the temple of God?.
What is the sanctuary of God or tabernacle of God?.
It's God's home, right?.
God is living in the midst of the temple, right?.
The temple is God's home..
And Paul is saying, "You guys are the home of God..
"You guys are the embodiment of the presence of God..
"You together are the temple of God.
"and that the Spirit lives in you.".
And you have to realize that..
When we congregate together, when we meet together,.
no matter what size of your congregation,.
if you call on the name of God,.
if you pray to the name of Jesus,.
we are the community of faith, we are the church..
At the same time, we are the temple of God..
God is living in the midst of us..
In 17, it's a huge warning..
He warns the Christian in Corinth,.
"God will destroy anyone who destroys his temple..
"For God's temple is holy and you are that temple.".
Let me read that again..
"God will destroy anyone who destroy this temple.".
And who is, what is this temple?.
The church in Corinth, right?.
The church people..
We are the temple..
"And God will destroy anyone who destroy this temple..
"For God's temple is holy and you guys are the temple.".
Do you realize what, it's so great of God's grace..
God lives, choose to live in the midst of us, right?.
Not that we deserve it, right?.
We are all broken and sinners, right?.
We all have our own flaws..

$^{761}$However, God chooses to live in the midst of his people.
who call on the name of Jesus..
And that's why we need to do everything.
to build up this temple..
That's why we need to do everything.
to build up this community, this church..
Because we understand the New Testament,.
the Bible tells us that we are the embodiment,.
the presence of God..
When someone walk into our church,.
let's say that that person is not a Christian yet..
It's not a believer yet..
He walked in..
How can he see the presence of God?.
He see the presence of God through us, right?.
Paul is saying to the Corinthian Christians,.
"If you are fighting against each other,.
"people cannot see God, right?.
"People cannot see the holy presence of God.".
He warns them, "If you destroy this temple,.
"God will destroy him.".
Because God's temple, God's home is holy, right?.
In the same way, that is a reminder for all of us,.
not to destroy, but to build up, build up the temple,.
build up God's home so that God will be glorified, right?.
And that is the essence of worship..
And brothers and sisters, I pray, I pray that.
we will be faithful, we'll be faithful.
to build up God's temple, God's church..
Now, you have to see it by faith..
You have to see it by faith..
You cannot see with your flashy eyes..
You have to see it by faith..
We see each other, the sister and brother in Christ..
We have to see it by faith that God is living.
in the midst of us..
When we pray for one another, when we worship together,.
God is living and dwelling in the midst of us.
so that our worship can bear the witness to the world,.
to tell other people God is real.

$^{801}$and God is worthy of our worship..
Let us all pray together..
Father, we come before you and ask that you will forgive us..
At times, we just trivialize worship..
And we thought we'd just come together and sing some songs..
No, it is not..
By your grace, you dwell among us..
And we are your people, and we are your temple, holy temple..
And God, I pray that you will renew our faith,.
renew our commitment to you,.
that we'll be your real people,.
knowing that you live among us..
And thank you for being our God, and we can be your people..
Let our worship be holy and pleasing to you..
Let our worship be a witness to the world,.
to other people, that you are real.
and you are worthy of our praises..
So pray that you will renew our spirit of worship,.
our sense of wonder, wonder of your grace,.
all and so awesome of your grace,.
that we can be called to worship you in your presence.
in Christ, and pray..
Well, God bless all of you, and next week, okay,.
we're gonna do our act of worship together, all right?.
Let's do it faithfully..
Let's do it in a way that is pleasing to God..
All right, see you next Sunday..
\newpage



\section{}
\label{sec:yIgKl3jW8_8}
\textbf{永恆的敬拜 3 (列王記上 6-7 ) - 陳崇基牧師}
\newline
\newline
連結: \href{https://youtube.com/watch?v=yIgKl3jW8-8}{\texttt{https://youtube.com/watch?v=yIgKl3jW8-8}} ~~~~ 語音日期: 2024-05-07
\newline
\newline
\hyperref[sec:7G7VT1O0Fmo]{\small{< < < PREV SERMON < < <}}
~
\hyperref[sec:index_chronic]{\small{[返順時目]}}
~
\hyperref[sec:index_scriptual]{\small{[返順卷目]}}
~
\hyperref[sec:EISv2Eie8Pg]{\small{> > > NEXT SERMON > > >}}
\newline
\newline
$^{1}$多謝各位.
跟大家say hi 好嗎.
歡迎大家一齊來敬拜.
今天我們繼續這個worship的series.
我們繼續看看教會的敬拜.
由頭到尾聖經到底怎樣去追蹤這個主題.
我們說過聖經開始和結終.
都告訴我們有一個很重要的命題.
就是敬拜.
上帝創造我們就是去敬拜.
上帝給我們將來也是敬拜.
在新天新地也是一個敬拜.
我們看過詩篇第八篇.
一首敬拜的詩篇.
還有我們看過《創世記》第二章.
開頭的時候.
上帝創造人的時候.
他的心意是什麼呢.
就是要我們去敬拜.
也看過他的回眸.
我們一直追蹤這個聖經的故事的時候.
發現以色列人被上帝拯救.
離開埃及的時候.
他們去到曠野這麼多年.
上帝是住在他們中間.
讓他們能夠有敬拜的生活.
那個住的記號或者住的方式.
就是一個回眸的方式.
你記得回眸是什麼形狀嗎.
長方形的.
裡面也有一個.
什麼形啊.
長方再有一個正方形.
沒錯了.
其實回眸就是一個帳幕.
是耶和華帶著他的軍隊.
帶著他的子民.
在曠野行走的時候.
一直來引領他們.
順帶一提.

$^{41}$那個回眸的形狀.
其實就好像埃及王打仗的時候.
他的帳幕.
所以以色列人出了埃及之後.
耶和華告訴他們.
我住在你們中間.
就不再是法老是你的主宰了.
我是你的主宰.
我是耶和華你的主.
你聽我說.
根據我的心意.
根據我的道理.
來組織你自己的生活.
這次我們繼續下去.
你知道他們以色列人.
完了曠野的路程的時候.
他們就進入了迦林地.
所以多年之後.
他們就安定了.
然後就有聯合王國出現.
聯合王國第一個王是誰.
第二個呢.
素羅因為不合格.
所以耶和華就解僱了他.
然後就是誰.
大衛.
一個合神心意的王.
然後大衛之後就是誰.
所羅門沒錯.
是素羅大衛和所羅門.
然後就發生什麼事.
然後就有人敗家.
然後就將國度分成南與北.
這個就是分裂的皇朝的時候出現.
所以我們今天看到.
所羅門執政的時候.
他做了什麼.
當然他最出名的.
做了的就是第一個聖殿.
因為大衛是沒有份去見聖殿的.

$^{81}$上帝說不用你見聖殿.
你的兒子要為我見聖殿.
所以所羅門在執政的期間.
他就建立這個第一聖殿.
所謂第一個first temple.
所羅門的執政聖經裡面.
有很多的詳述.
有11章這麼多.
如果你有聖經.
打開去到列王記上.
第一至第十一章.
我們大概看看這些篇幅.
很多時候我們看到所羅門的時候.
這個不是陌生的一個人.
不單止是教會基督徒.
對所羅門不陌生.
其實在外面不回教會的人.
他們都可能聽過Solomon.
所羅門這個王.
無論是時代雜誌.
或者一些出名的媒體.
有時候都會說一下所羅門的.
說找到什麼.
所羅門的國度等等.
是一個很出名的王.
在聖經裡面有很多詳細的記載.
說到所羅門的國度.
如果你想起所羅門的時候.
大家可能想起.
我們都挺喜歡這個人的.
因為他有什麼值得我們喜歡.
但是他的樣子是怎樣我們不知道.
總之他有很多的威水史.
你想想一個君王.
需要有什麼.
一個君王治理一個國家.
一個社會.
他應該有什麼.
他要有智慧.
不像我們今天很多政客.

$^{121}$個個都沒有智慧.
古今中外全部都是.
搞什麼鬼.
應該有智慧.
因為要做一個明君.
管理一個國家當然有什麼.
國防.
軍事上的東西.
或者外交的政策.
有財務上的需要.
財源.
我們看看某些的聖經.
怎樣去描述所羅門.
我有一個命題.
我們可能很喜歡所羅門.
但如果你看一看的時候.
你發現聖經的作者告訴你.
其實不要這麼喜歡他.
為什麼這樣說.
我們先看看聖經.
第一段四章二十九節.
這裡說到所羅門的哪方面.
他說上帝賜給所羅門極大的智慧和聰明.
廣闊的心如同海沙.
那麼多.
所羅門的智慧超過.
東方人的智慧和埃及人一切的智慧.
厲害.
他的智慧勝過萬人.
他的名聲是全遍列國的.
作真言三千句.
試過一千零五首.
是不是很厲害.
不是一千首.
再多五首.
多不多.
這個就代表一個明君的描述.
一個好的君主就需要有智慧.
智慧甚至是超過其他人.
然後十章第二十六節.

$^{161}$所羅門聚集戰車和騎兵.
一千四百輛戰車.
一萬二千名騎兵.
安置在屯車城.
在耶路撒冷的黃那裡.
黃在耶路撒冷.
使銀子多如石頭.
真的厲害.
你外面看到那些石頭.
多如石頭.
到處都看到.
香柏木多如謝菲拉.
謝菲拉是以色列巴勒斯坦地區的其中一個區域.
專門出一種樹叫桑樹.
桑樹.
多過他們本地的桑樹.
香柏木是從那裡來的.
在以色列境外北面.
黎巴嫩那裡.
斬那些木頭然後浮木頭到地中海.
到以色列這個地方.
多到一個地步.
多到我們那些野樹.
多過野樹.
所有的馬是從埃及和科威運來的.
埃及在南邊.
科威在北邊.
現在土耳其的地方.
黃的商人定價.
從科威這個地方買來的.
很多的馬.
很多的戰車.
每部戰車怎樣介紹等等.
這裡講到國防.
講到他的資源和財務.
還有什麼呢.
就在這裡.
十一章第三節.
他娶了七百個公主和三百個妃嬪.
厲害吧.

$^{201}$一個都搞不定.
七百個不是吧.
還要多三百個.
後宮有多少人呢.
這些事情不只是色慾的問題.
一個王為什麼要有那麼多妃嬪.
因為那些妃嬪是從哪裡來的.
是周圍鄰國來的.
我們以前中國四大美人都有.
那些都是走去鄰邦的.
嫁出去.
是外交政策.
就是說他和鄰邦全部都有一個Diplomacy.
我們講了智慧有兵馬.
有財源.
有外交.
也都有什麼呢.
以前當時的王一定要有聖殿.
是王去到建最重要的聖殿.
當然以色列國裡面只有一個聖殿.
只有一個廟.
就是耶和華的聖殿.
有圖有真相.
這裡給一個假的圖給你.
讓人們去Reconstruct.
左上角就是大衛城的一個Model.
大衛城是長長的在山上.
最上的地方就是聖殿和皇宮的地方.
右邊那個就是聖殿.
所羅門第一個的聖殿.
門口向哪.
向東.
前面有一個銅盤和祭壇.
然後進去這個聖殿裡面.
就是長方形.
然後裡面有一個正方形.
自稱所為.
這是一個很輝煌的聖殿.
並且左下角.
你見到有一個橫切面.

$^{241}$剖開它告訴你.
你見到裡面中間那個地方.
是不是有隻東西在那裡.
那隻叫什麼名字.
基羅伯沒錯.
大家很聰明.
講了幾次.
他是守哪裡的.
是守住耶和華的寶座.
守住聖所這個地方.
所以你見到.
在聖經描述所羅門的執政的時候.
不單止說他有錢有智慧.
有兵馬等等.
特別講到很多篇幅.
是講他建造聖殿.
第六章.
第六章很有趣.
整章就是講建聖殿.
那些擺設.
金銀等等.
開頭第一節.
他講到一個年份.
然後第六章最後的時候.
又講到一個年份.
是講什麼呢.
他說以色列人出埃及.
480年.
終於來到這裡.
他們開工去建造耶和華的殿.
即是說出埃及流浪了這麼久.
什麼時候才真正出埃及.
不是過紅海過海.
過了這麼多年.
他們終於settle down.
能夠耶和華住在他們中間.
能夠有一個敬拜的生活.
就是到鎖門的時候.
所以這麼多年來都是出埃及的.
入了監牢都算是出埃及.

$^{281}$直到他們開工建聖殿的時候.
終於完結exodus.
終於離開了埃及這個地方.
終於離開了法路.
我們歸耶和華.
作他的聖文也要敬拜他.
這個就是聖殿有的目的.
然後最後的時候.
六章37至38節.
在位的第四年西忽元.
立了耶和華殿的根基.
由頭到尾他就告訴你.
整個聖殿裡面的細節的details.
這些對我們現代人來說.
可能會覺得很悶.
但是在當時的人來說.
是一些很重要的report.
就等於我們家事大會裡面.
有一些很重要的report.
事公的report.
或者財務上的report.
你聽完那些report的時候.
你會覺得這個值得讚美主.
因為在我們中間.
耶和華住在我們裡面.
有聖殿被建立起來.
這個王就是所羅門.
就是他建造第一個的聖殿.
所以第二段就二至第十節.
他就說基本的結構.
那些內殿外殿等等.
那些尺碼等等.
然後到第四段.
又說這個設計.
裡面的金紙銀有多少.
拜道威昨晚我搜尋.
金這個字.
最多在哪裡出現.
就在這裡.
很多金.

$^{321}$你去發現作者去描述聖殿的時候.
金光燦爛.
等於你進了鑽石舖的時候.
你望不到其他東西.
你只望到那些幾卡的鑽石.
金光閃閃.
所羅門執政就是做這件事.
在中間夾住的是幾節.
十一節到十三節.
就去警告耶和華警告所羅門.
好 你已經建了這個聖殿.
輪到你建這個聖殿.
如果.
條件來的.
條件據.
如果你遵行我的律例.
我的典章.
聽從我一切的誡明的話.
行在其中.
我必應驗我向大衛所說的話.
我必住在以色列人中間.
並不凋棄我的百姓以色列.
所以有了這個聖殿的時候.
耶和華就跟他說.
好 你開始敬拜我.
不過.
有一個條件.
就是整個生命都要敬拜.
有一個警告.
我們這樣介紹了所羅門的時候.
我們發現都幾喜歡他.
但慢慢看下去的時候.
你就發現.
好像有些.
有些不太妥.
為什麼不太妥呢.
剛才我不是說有金有銀.
有兵馬.
有這些外交等等.
這些很正常.

$^{361}$古今中外.
所有的君王.
所有的統治者.
國家的掌門.
不是掌門人.
是領導人.
國家領導人.
都需要這些.
但在聖經裡面.
怎樣去評估一個王呢.
如果我回到五經的話.
申命記第十七章.
你留意一下.
我們看看這段聖經.
你再比較一下.
這段聖經裡面說一個王.
上帝對他有什麼要求的時候.
你再看剛才.
所羅門怎樣描述他的時候.
你會不會覺得.
有些是挺有趣.
有些不太妥.
十七章十四節.
你到了耶和華.
你上帝賜給你的地.
得了那地.
居於其中的時候.
若說我們要立王.
治理我們將四圍所有的國家一樣.
你一定要立耶和華.
你上帝所揀選的人.
為你的王.
要從你的兄弟裡面.
一人立他為王.
不可以立你兄弟以外的.
我幫人治理你們.
當然.
只要王不可以為自己.
家添馬匹.
也不可以為自己.

$^{401}$家添馬匹.
使百姓回埃及去.
因為耶和華曾對你們說.
不可再往那條路走.
不要再回埃及.
王不可以為自己立.
妃嬪.
免得她的心偏離.
也不可以為自己多積金銀.
他登位的時候.
十八節.
要做什麼呢.
剛才說幾樣東西不可以做.
登國位的時候.
就要在利美加的祭司面前.
將這律法書五經.
自己抄寫一份.
這書要存在他那裡.
一生的烈日.
他需要去重讀.
要去查經.
要去研讀聖經.
要去明白.
怎樣去敬畏耶和華他的神.
遵守他的律法書裡面一切的律例.
免得他的心向弟兄高傲.
偏離這界名.
或向右或向左.
他和他的子孫.
就可以長久作王治理以色列.
有沒有留意三樣東西是不能做的.
不可以加添什麼呢.
馬匹.
不可以多立妃嬪.
也不可以多積金銀.
所以你想想.
五經怎樣看一個君王的時候.
你再對比一下.
所羅門怎樣被描述的時候.
你就發現好像有些不妥.

$^{441}$是不是.
這個是不是一個我們那麼喜歡的人呢.
聖經說似乎不是.
到底有些什麼事發生了.
我們說過所羅門的執政有十一章那麼長.
一至十一章.
列王記常.
原來中間有一個轉折.
特別說到聖殿的興建的時候.
你就發現有一個類似這樣的結構.
我們不能夠去看一些細節.
不過我們可以看一下結構.
結構已經告訴我們.
聖經想說什麼.
我經常都說.
讀聖經.
We don't know what it means.
Until we know how it means.
我們不知道聖經的經文意思.
直到我們明白聖經怎樣表述意義.
怎樣寫 怎樣佈局.
原來那個作者經文.
是有他的心思在那裡.
你見到第六章我們剛才說的金銀 內電 外電那些厲害的東西.
就是所羅門興建聖殿的時候.
六章一至到第三十七節.
然後第二段.
他就描述.
所羅門興建自己的宮殿.
七章一至到十二節.
然後又回到說.
所羅門興建聖殿.
十三節到五十一節.
很奇怪.
第六章明明是在說興建聖殿.
應該一早說完才說其他東西.
聖經作者就說了一下.
然後停一停.
不如說一下.
所羅門興建自己的宮殿.

$^{481}$說完他興建自己的宮殿.
然後才說所羅門又去興建聖殿.
到底想說什麼.
那個轉折位在哪裡.
就是七章第一節.
六章最後七章.
那個接縫的位置.
轉折位這裡有說什麼呢.
先看第二節.
所羅門王.
衛爺華所建的殿.
長六十爪 闊二十爪 高三十爪.
一爪大概是一尺半左右.
然後我們跳到三十七節.
所羅門在位第四年西伐月.
立了耶和華的殿的根基.
三十八節.
到了十一年布辣月.
就是第八個月.
殿和一切殿都按照樣式造成.
所羅門建殿用了七年.
轉折位到了.
七章第一節.
這個我翻譯給你聽.
因為那個希伯來原文.
其實有一些重要的.
叫做subfix.
他的 他自己的.
那裡很多時候中文翻譯.
出不到那個意思.
七章第一節轉折位.
但是為他自己的宮殿.
所羅門用了十三年.
方才完成他自己整座宮殿.
第二節.
七章第二節.
他建造黎巴嫩臨宮.
長一百爪 闊五十爪 高三十爪.
有四行香柏木柱.
柱上有香柏木的橫樑等等.

$^{521}$你留意到這裡.
你compare一下.
你看的聖經.
是我的 是突出什麼.
他自己 他自己.
是焦點在自己身上.
他這個尺碼 聖殿有多大.
長六十闊二十 高三十.
這個是聖殿.
他自己的黎巴嫩臨宮.
長一百闊五十 高三十.
你知不知道architecture是power.
我們芝加哥很清楚.
有沒有去過那個.
芝加哥河的tour.
坐在船上.
導遊就告訴你.
哪座哪座多厲害.
哪座高一點.
哪座大一點.
就是什麼.
就是厲害一點.
當然了.
難道這座嗎.
那座多厲害.
那座雀標.
什麼時候建造的.
哪個設計師去建造的.
等等.
其實這裡都是一樣.
你看到開始不是很妥.
其實第七章講到黎巴嫩臨宮之後.
講了五個所羅門的宮殿.
五座不同的宮殿.
黎巴嫩臨宮 處子廳 審判廳.
自己住的宮殿.
法老女兒住的宮殿等等.
五座就介紹了他.
然後才講聖殿.
作者用意在哪裡.

$^{561}$聖殿是代表上帝住在子民中間.
讓他們可以敬拜他 尊他為大.
在這個結構裡面你會發現.
好像那個人 所羅門.
他的生命次序有顛倒.
有些問題出現了.
自己蓋的聖殿就大過神的聖殿.
如果你有機會再看一下整段描述所羅門的.
特別第八章講到獻殿的時候.
第八章整段都是講聖殿終於落成.
然後有一個很重要的獻殿儀式.
還有裡面有一個很長很長的祈禱.
這個是列王記上帝八章.
有時間你回去看一下.
是一個很長的禱告.
裡面講到甚麼呢.
裡面講到有五次.
所羅門說我為上帝建造了這殿.
五次 第八章.
二十節 二十七節 四十三節 四十四節 四十八節.
不是六合彩號碼.
五節在祈禱裡面.
我為耶和華建造了這殿.
很奇怪 在整個禱告裡面.
一字不提 他有宣讀律法書.
我們查過李軒禮記.
整個聖城落成之後做了甚麼.
以斯拉尼希米帶領人民看律法書.
因為他們回歸律法.
Return to the Torah.
但在所羅門的獻殿裡面.
居然律法書一字沒有提.
再對立一下後面來的王.
有一個很好的王叫約西亞.
約西亞有一個事蹟 做西亞.
就是清潔聖殿的時候找回甚麼.
律法書.
原來聖經放在那裡封塵.
打了這麼多年都沒有人看過.
然後約西亞就說原來我們一直沒有看過.

$^{601}$快點打開 然後我們帶領整個國民來改革.
所羅門一字沒有提他宣讀律法書.
另外 在整個獻殿的過程裡面.
他有沒有謙卑下拜.
你對立一下大衛.
大衛將約桂 耶和華的座駕.
帶入京都耶路撒冷的時候他做甚麼.
他跳舞.
跳舞就是謙卑.
謙卑在耶和華面前敬拜他.
因為這個是屬於耶和華的日子.
但所羅門卻沒有.
這裡很奇怪.
當他最高峰的時候.
聖經作者很明顯告訴我們.
是有些不妥.
是否成功充分頭腦.
我記得我去區會在另一個城市.
好幾年前.
在新建成的教會.
新開始的教會.
我記得有一個主要亞洲人多的教會.
是發展神速的.
發展神速.
一開始.
沒多久.
就開了另一個地點.
開了好幾個地點.
發展神速的總會的人來看.
個個都說那個就是我們的Poster Child.
上雜誌封面就是那個.
哪一間教會.
搞了幾年之後.
開始搞不清楚為什麼.
搞了幾年之後.
領袖就開始鬼打鬼.
搞不定.
最後就崩潰.
整個散了.
原因就是那些領袖.

$^{641}$是不願意去被問責.
然後就搞了一場風雨.
發展神速但很成功.
這個就是我想說的.
很多時候我們看鎖門的時候.
我們很喜歡這個人.
這個人OK的.
因為他成功.
其實會不會暗地裡我們喜歡成功神何.
他這麼成功一定是上帝祝福他.
但是列王記的作者就說.
你先煞車.
看清楚到底發生了什麼事.
到底轉折在哪裡.
我們做牧師的做了這麼多年.
我們有時候弟兄姊妹有自己新的生意.
開舖.
他會叫我們牧師去祈禱.
如果叫到你去祈禱你去不去.
你去不去.
我們開了新的生意.
新的商舖.
你們來為我祈禱.
其實我不介意的.
我不介意.
雖然你覺得是不是有點奇怪.
你真的覺得有點奇怪.
因為我去過.
你覺得有點奇怪.
有時候都很奇怪.
好像愛龍虎牧師一樣.
只不過我不做那件事.
但其實我不介意的.
如果你有開商舖.
或者有新的生意.
你才來祈禱.
OK Fine.
但你不要想著這個牧師特別靚.
我們每個人都祈禱.
沒有說誰特別靚.

$^{681}$因為我們是奉耶穌基督的名求.
你不是奉我的名求.
奉我的名求就死很多人.
The point is.
當然我不介意.
但挑戰就來了.
到你有一個新的enterprise.
一個新的endeavor.
一個新的努力.
無論是一個生意也好.
商舖也好.
到底是不是跟著上帝的道.
The point不是有一個神人在幫你開球.
他最多幫你.
好像足球場一樣.
幫你踢球開球.
贏不贏到就靠你了.
你叫我去.
我幫你開飯祈禱.
我只不過是幫你開球.
但最重要的是什麼.
你有豐功偉績的時候.
你心靈裡面優先次序是不是搞清楚.
所羅門四十年時間.
死了之後就王國瓦解.
一分為二.
北面有十個支派叫以色列.
南面一個多支派就是猶大.
北面的支派就交給了所羅門其中一個幫他打工的.
叫做耶羅波安.
十個支派交給了他.
南國就交給了所羅門的兒子.
死了之後就是羅波安.
所以北面那個叫耶羅波安.
南面猶大國的叫羅波安.
一邊是十個支派.
因為神說我把他的國太不乖了.
這個人搞錯了.
所羅門搞錯了.
我把他的國度一分為二.

$^{721}$北面交給其中一個將領.
就是耶羅波安.
南面交給他的兒子羅波安.
在交接的期間.
你就發現聖經到底怎樣去描述所羅門.
十一章四十節.
在交接的時候.
也是所羅門萬年的時候.
他知道羅波安得到先知的啟示.
先知說我把十個支派北面交給你.
這個國度一定會分裂的.
所羅門想做什麼.
他想要殺耶羅波安.
耶羅波安逃走去埃及.
他到了埃及去視察.
他就暫住在埃及.
直到所羅門死了.
十二章第三節.
以色列人派人.
這裡完了.
所羅門死了.
然後羅波安.
就是所羅門的兒子.
羅波安.
記不記得.
他就上位了.
上位去管治南國.
管治以色列.
然後十二章第三節.
以色列人.
就是北面那一班人.
十個支派.
他就派人請耶羅波安來.
因為耶羅波安在埃及.
他逃跑了.
因為怕被人追殺.
耶羅波安就和那些以色列的傳會眾.
就和羅波安來對談.
和新的王來對談.
他說你父親叫我們負重額.

$^{761}$現在求你減輕我們父親所加給.
你父親加給我們的苦工和重額.
我們才來服侍你.
然後羅波安就問了朝廷裡面的老神僕.
老神僕就說.
Let's listen to the people.
但那些年輕的神僕就說.
不要跟他們.
我們說事.
羅波安就跟誰.
當然是跟年輕的那些.
一起長大的那些.
所以第十節.
羅波安王就對耶羅波安和那些以色列十個支派的人說.
我的小指頭比你父親的腰還粗.
我父親.
我父親使你們負重額.
現在我必使你們負更重的額.
我父親用鞭子去懲戒你們.
我用蠍子來打你們.
拜祖為什麼我笑呢.
因為他說我的小指頭比你父親的腰還粗.
這句其實是粗話.
聖經裡面有這樣的說話.
其實我想他寫的文章是集合了聖經裡面的這些粗話.
有好幾句.
你想一下.
你用你自己的想像力.
我的小指頭比我父親的腰還粗.
厲害嗎.
你敢動我.
在情節裡面你可以想一想.
這個情況像什麼.
黑社會.
你想一下一個在位的王.
唯一一個在位的王.
去追殺一個新的王.
然後在位的王又令到人民是負重額的.
你想一下是什麼情節.
出埃及.

$^{801}$法老就去做這件事.
法老記得追殺新的王就是摩西.
他那個會帶領人民自由的.
法老又使用以色列人.
有苦工和負重額.
但這裡描述的居然就是所羅門和羅波安.
所羅門其實是新法老.
耶羅波安不變十字派.
那個好像摩西一樣逃走.
然後回來講數.
就帶以色列十字派走.
原本是一個明君.
所羅門就變成一個法老.
很奇怪.
在地圖上他們在以色列國.
不是在埃及.
但以色列國變成埃及.
以色列變成了埃及.
是不是一個最成功的人.
為上帝建殿.
去敬拜上帝.
但他在生命裡面有些事情搞錯.
建聖殿又建自己的皇宮.
本身是我們去敬拜上帝.
變了本身是我們服侍上帝.
現在是上帝以宗教的理由服侍我.
這個很大的警號.
我們經營我們生命的時候都是一樣.
有很多好的東西.
上帝給你.
無論是學業,事業也好.
賺到錢也好.
青年的,家庭的.
賺到錢.
建立到家庭.
兒女.
是不是好事.
當然好.
有兒女.
我也有兒女.

$^{841}$但會不會搞一搞.
不是我去敬拜上帝.
而是上帝要來服侍我.
慢慢我們敬拜的生活就會墮落.
我們對上帝沒有一個尊他偉大的敬拜生活.
擺到他次要.
無所謂敬不敬拜上帝.
最重要是搞定我自己的事.
我那些五個行宮.
能不能搞定.
我不是要望天打卦.
什麼都不管.
而是我們心底裡骨子裡.
我們是不是真的相信上帝是值得敬拜.
還是我們心底裡其實都是懷疑.
都是覺得上帝不是很值得我們服侍.
信這個上帝不是真的要敬拜他.
信這個上帝是要這個上帝服侍我.
求主幫助我們.
我們去敬拜上帝的時候.
是以他為最大.
我們去服侍上帝.
不是上帝來服侍我.
不是上帝來建立我.
而是我用我的敬拜生活.
用我的生命去建立上帝的國度.
去服侍他.
去同心協力祈禱.
去領受聖餐.
上主我們求你.
求你都可以原諒我們.
因為很多時候我們實在.
把我們生命的焦點搞錯.
我們把對你的敬拜.
很多時候放在次要.
對你的敬拜.
我們放為一個註腳.
放為一些等閒的事.
上主我們求你.
你都去赦免我們.

$^{881}$主讓你的子民今天來到你面前.
我們奉你的名來敬拜的時候.
我們知道你是配得我們讚美.
你是配得我們敬拜.
求你重整我們生命的優先次序.
今日我們拿起聖餐的餅和杯.
你再告訴我們這個榮美的救恩.
重價的救恩.
是要我們生命屬於你.
要我們生命來敬拜你.
求父你聽我們禱告.
閱立我們的敬拜.
祈禱奉主名.
\newpage



\section{24:1-12}
\label{sec:EISv2Eie8Pg}
\textbf{Easter: Luke 24\_1-12 He Conquered Sin and Death}
\newline
\newline
連結: \href{https://youtube.com/watch?v=EISv2Eie8Pg}{\texttt{https://youtube.com/watch?v=EISv2Eie8Pg}} ~~~~ 語音日期: 2024-05-08
\newline
\newline
\hyperref[sec:yIgKl3jW8_8]{\small{< < < PREV SERMON < < <}}
~
\hyperref[sec:index_chronic]{\small{[返順時目]}}
~
\hyperref[sec:index_scriptual]{\small{[返順卷目]}}
~
\hyperref[sec:6bRdtM__Fr0]{\small{> > > NEXT SERMON > > >}}
\newline
\newline
$^{1}$Good morning and welcome. I'm so excited that you're joining us this morning on this beautiful Easter Sunday for the Urban Voice English Service..
So glad that you're here and I hope that you've enjoyed this week. It really has been beautiful..
There's been some very beautiful days. There were a few cool days, but there were some really nice days here in Chicago..
I love this time of year, the spring when the days are getting longer and the grass is beginning to turn green.
and there's even some flowers beginning to bloom. Such a beautiful time of year and so hopeful for us..
It reminds us of the renewal and the new life that God brings. So this is just a really beautiful time of year..
I hope you're getting outside and enjoying it. I've been trying to get out and enjoy it and garden and hike and walk with my wife and my kids..
I love being outdoors, so I've been doing that..
This is a special week, obviously. That's why I'm dressed up a little bit. If you've tuned into the U.V. English Service before,.
you know I don't normally dress up like this. It took me a while to tie this bow tie. I'm really not good at that at all..
But I'm dressed up because it's a special week. It's Easter and this is really a special day in our faith.
and a special time to celebrate the resurrection of Jesus and the hope that lies in that resurrection..
I hope that as you've gone through this week, what we sometimes call Passion Week or Holy Week,.
you've had a lot of time to reflect on Jesus and on His life and on His death and His death on the cross and what it meant for you..
I hope you've had a lot of time to connect with God, getting in the Word and spending time in prayer, spending time in worship,.
and maybe spending some extra time just reflecting on the cross and the story of the cross and what it means..
The redemption, the forgiveness, the grace that is communicated in that story. I hope this has been an encouraging week for you..
One of the things I love about Holy Week is it's kind of a microcosm of our faith..
It starts with Palm Sunday and excitement and anticipation and hopefulness and then it progresses through some really difficult times..
Good Friday when Jesus is condemned to die and hung on the cross and He dies..
There's a sense of hope deferred and of discouragement there..
Then it ends with Easter Sunday and the renewal of God's hope of redemption..
I think that's just a beautiful picture for us because it gives us a very real picture of our world..
The world we live in is broken. We can't pretend that everything is fine and that everything is okay and that everything is great..
We can't. We can't pretend that way in our faith. We need to be honest about the brokenness and the suffering of the world..
But we can't get stuck in depression over how broken things are..
We have to have the hope of Jesus and that's what Easter is about..
It gives us the hope that God has come to rescue us, that God is going to make things right..
That's really what we're going to focus on today..
We're going to focus on how the resurrection answers so many of the questions that we have in our hearts..
It answers those heart cries, those deep longings and those heart-wrenching questions that we have..
Jesus answered those things in the resurrection..
He answered them in an affirmative way and He answered them in a beautiful and loving way..
We're going to look at that today..
Let's turn together and look at Luke 24, the first 12 verses..
This is the story of the resurrection..
It says, "But on the first day of the week, at early dawn, they went to the tomb, taking the spices they had prepared..
And they found the stone rolled away from the tomb..
But when they went in, they did not find the body of the Lord Jesus..
While they were perplexed about this, behold, two men stood by them in dazzling apparel..

$^{41}$And as they were frightened and bowed their faces to the ground, the men said to them,.
'Why do you seek the living among the dead? He is not here, but has risen..
Remember how He told you while He was still in Galilee that the Son of Man must be delivered into the hands of sinful men and be crucified and on the third day rise.'.
And they remembered His words..
And returning from the tomb, they told all these things to the eleven and to all the rest..
Now it was Mary Magdalene and Joanna and Mary the mother of James and the other women with them who told these things to the apostles..
But those words seemed to them an idle tale, and they did not believe them..
But Peter rose and ran to the tomb, stooping and looking in..
He saw the linen cloths by themselves, and he went home marveling at what had happened.".
This is just the very beginning of the story of the resurrection..
There's so much more that's here, but we're just going to focus on this section today..
And it gives us a snapshot of what happened on the morning of the first Easter..
And it begins to answer some of those questions that we have in our hearts..
The first thing I want us to consider this morning is that the resurrection is real..
It's real. It really happened..
And I understand that it's an amazing story to hear and to consider..
And that a lot of us, when we hear it, we might feel skeptical..
We might say, "Okay, wait a second. I don't know. Could this really happen?".
Especially if you're not someone who's in the church, you might look at this story and say,.
"Okay, this just seems too fantastic to even be true,.
that someone could be raised from the dead when they've been in the tomb for three days.".
And there's a lot of good evidence for the resurrection..
But I just want to focus on a few things today because we don't have a lot of time..
We could spend weeks and weeks talking about the evidence..
But I just want to focus on a few things to help remind us that there's a lot of good reason.
to believe that this really happened and that the resurrection is real..
It's not just a story. It's not just a fable..
It's not just a myth or some story with a moral meaning..
It's actually something that really happened..
Now, we know that Jesus is one of the most, probably the most, well-documented figure in antiquity..
There are so many sources about Jesus and his life that speak about him..
And not just from Christianity, the Gospels, those sources..
There's a lot of other contemporary historical sources from people who are not Christians.
and who were just historians who wrote about Jesus and his life and his death..
And the fact that he was crucified, he died on a cross,.
and his followers said that he rose from the dead..
So this is not really disputed by any historian, the fact that Jesus lived,.
the fact that he was crucified on a cross, and the fact that people said he rose from the dead..
Pretty much everybody is in agreement about that..
But obviously there are people who say, "We don't think Jesus actually physically rose from the dead..

$^{81}$That just seems like too much.".
But here's three big things for us to consider this morning.
that help us see that this is something that really happened..
First of all, there's the empty tomb..
Jesus' body was not there..
And that seems like something simple in the beginning, but it's an inescapable fact..
When we look in the book of Matthew, in chapter 28, verses 11 to 15,.
when the guards, who the priests had assigned to guard the tomb, come back to the priests,.
the Jewish chief priests and the Jewish elders, and say,.
"Hey, something crazy happened. There was this earthquake, and the body is gone.".
The chief priests say, "Well, we don't want this word to get out, so here's what we're going to do..
We're going to pay you guys some money to just tell people that the disciples came at night and stole the body.".
Now, these people, the priests and elders, were actually enemies of Jesus..
They did not like Jesus at all..
If there was any way that they could produce his body and say,.
"Hey, look, it's a hoax. It didn't happen. Here's the body here," they would have done so..
But even they, even though they did not like Jesus, could not escape the fact that the tomb was empty..
His body was gone..
So they had to invent a story to try to explain it..
And that's an inescapable reality of history..
The tomb is empty. Jesus' body is gone..
No one has ever been able to say, "Hey, we found Jesus' body buried in this place, and we've confirmed it.".
Or, "Hey, we found Jesus living in this other place as an old man, and we've confirmed that.".
No one has ever been able to prove that..
No one has ever been able to say that the tomb is empty..
Jesus really lived on the earth. He really died on the cross..
And the tomb is empty. His body is not there..
He really rose again..
And the next thing I want us to consider is the eyewitnesses that were eyewitnesses to this..
Now, we can see here in Luke, in the passage that we read,.
that even the eyewitnesses themselves were initially skeptics..
They were not sitting around trying to invent a story of how Jesus was rising from the dead..
They were skeptical themselves..
They had forgotten that Jesus talked about being resurrected from the dead..
That's why the angels have to remind them of that in verse 7 of Luke 24, that Jesus had prophesied this..
The angels remind them of that because they had forgotten..
They weren't expecting it..
And they were carrying on as if Jesus' death was a final moment, as if it was all over..
How do we know that? We see that in their actions..
The women, who are the first witnesses to Jesus' resurrection..

$^{121}$And by the way, that's a whole other sermon in and of itself..
It's such a beautiful thing that the first witnesses to the resurrection of Jesus were women..
Because in the first century, in their culture, these women were not really empowered in any way..
Actually, their testimony was not admissible in a court of law..
And so, if you were going to try to find witnesses to support something happening,.
and to try to make a good case that this happened,.
you would not gather a group of women and say, "These are the witnesses.".
But Jesus, His resurrection was first revealed to the women because of God's love.
and God's view of women in the kingdom of God..
And that's a really powerful thing in and of itself..
But these women are the first ones who come..
They have come to anoint the body..
They have come to put these burial spices and sweet-smelling things on the body..
They're not coming because they're expecting for Him to be risen..
They're coming because they're convinced that He is dead..
And they're going to carry out some kind of last service to Jesus by caring for His body..
They're expecting Him to be dead in every way..
We also see that with the disciples, the eleven who are left..
Where are they? They're hiding out in locked rooms..
We see this in the Gospels, that when Jesus begins to appear to them,.
they're hiding in locked rooms because they're afraid..
They're convinced that it's all over, that Jesus is dead,.
that everything that they were moving towards and believing in is now done,.
and that they're probably next, that the authorities are coming for them next..
And so, they're hiding in locked rooms..
They're not anticipating this..
The disciples themselves are so skeptical that here in Luke's account,.
chapter 24, it tells us that when the women, in verse 10,.
come to the disciples, to the eleven, and tell them what they've seen at the tomb,.
that the disciples, it seems to them like an idle tale, and they don't believe the women..
They say, "No, that's just too much. That's crazy. I can't believe it.".
Even the disciples themselves are skeptical..
They say, "This can't actually happen.".
And they had seen the miracles of Jesus,.
and they had heard Jesus prophesy that He was going to rise again,.
but they couldn't believe it. They couldn't accept it..
So, obviously, it's okay to be skeptical..
It's okay to say, "Okay, this seems like something really big and amazing and impossible.".
Yes, it is impossible, but God has the power to do the impossible,.
and that's what He's demonstrating here..

$^{161}$Another thing that's important to understand about these eyewitnesses,.
the women and the apostles themselves,.
was they suffered for this story..
This was not a story that they told to get rich..
This was not a story that they told to benefit themselves in some way..
As they told the story of the risen Christ, and as they went out to share the gospel,.
they were persecuted for it, and every one of them was tortured and eventually killed.
in terrible and brutal and violent ways for this story..
And obviously, if you're telling a lie to other people,.
you might keep the lie up as long as it's benefiting you,.
as long as it protects your reputation, or as long as it enriches you in some way..
But when it starts to hurt you, and when it starts to destroy your life,.
you have no reason to continue to tell that lie anymore..
And you have no reason to tell that lie until your dying moment,.
and to believe in it with all of your heart,.
so much so that you're willing to be martyred for that..
That just doesn't make any sense at all..
And so we have to look at the lives and the testimonies of these eyewitnesses.
and say that they were believing in something real and powerful..
They didn't allow themselves to be tortured and killed and martyred in violent and brutal ways.
for some type of fantasy or some type of idea they had concocted in their head..
This was real, and they knew it, and that's why they had the power of God in their lives.
to spread the gospel and to do miracles and to do all the things the Book of Acts talks about..
And it's why they believed it and believed it enough that they never recanted,.
even when they were tortured, even when they were forced by authorities,.
and the Romans and the Jewish authorities and all these different people.
who just said, "Stop preaching, stop telling this story, just say it didn't happen,".
and they wouldn't, and they went to their deaths believing this because it was real..
And by the way, the Book of Acts talks about the fact that Jesus appeared to over 500 people.
after he had resurrected, before he went into heaven..
So this was not just a small group of people that got together and shared this story..
It was hundreds of people who witnessed Jesus after the crucifixion,.
who witnessed his resurrected body..
They were eyewitnesses..
And the last thing I want us to consider is just the presence of the church itself..
You know, if you concoct a big con and then you get found out,.
you lose all credibility and no one will ever listen to you again..
Think about someone like Bernie Madoff who conned people out of their money,.
and for a while he was trusted and everybody thought he was great.
and that he was this financial wizard..

$^{201}$But when it came out that he was lying to people and losing people's money,.
then he was put into prison, and obviously no one would trust him with their money now..
So if this was some type of fabricated story that the disciples made up.
and it was a lie that they spread, then it could have been easily refuted,.
and they would have lost all credibility,.
and no one would have believed in the message of the gospel or the church,.
and this movement never would have gotten off the ground..
Because everyone would have said, "Oh, we've heard it before. It was a big con..
It was fake. It was just a story that someone made up, and we can clearly refute it..
That person has no credibility, and they should go to jail.".
But that's not the case, right?.
The church grew explosively and has continued to grow.
and has continued to weather storms of persecution and tribulation.
and trials and trouble for the last 2,000 years..
And that tells us that the strength of the early church.
and the continuing strength of the church today around the world.
and the growth of the church around the world.
is not because of something that was simply made up..
It's because of something real..
Now, it is an amazing claim..
Other religions say, "Hey, our leader had a vision or a dream or something like this,.
and that's how we came to know this enlightened truth.".
Jesus says, "Hey, I rose from the dead,.
and I am the way, the truth, and the life, and through me.".
That's how you know the truth..
That is an enormous claim,.
but there's very good reason to believe that it's real..
I'm not trying to tell you that I'm going to convince you of it..
If you're a skeptic and you say, "I'm still not certain," that's okay..
I'm not here to argue today,.
but I do want you to understand there are very good reasons to believe that this is real.
and that it really happened.
and that the power of God was at work in the life of Jesus.
in raising him up from the dead.
and in the church that began after Jesus left the earth..
I want to talk about how the resurrection answers some of the questions that we have in our life..
First, I want us to see that the resurrection certifies the message of the gospel..
The gospel says that there's one way to heaven..
Jesus says that he is the way, the truth, and the life,.
and that no one gets to the Father except through him..

$^{241}$Again, that's a pretty big claim..
Jesus talks about how we come to him through grace and forgiveness,.
that it's not through doing good works..
It's not through trying to be good enough..
It's not something that we earn our way into favor with God or earn our way into heaven,.
but that it's through God's grace and love and forgiveness..
It's by grace and through faith, through our belief in Jesus' love and work on our behalf.
that we come into the kingdom of God..
Those are huge claims..
Without the resurrection, we'd have a lot of reasons to say, "I'm not really sure if that's true..
That sounds too good to believe..
That sounds too good to be true, that I could be forgiven for all the mistakes I've made,.
for everything I've done wrong, and it could just be forgiven and wiped clean.".
That sounds too good to be true, but the resurrection certifies that the message of the gospel is true..
It assures us and it says, "Yes, the gospel is real..
Yes, Jesus is the Son of God..
Yes, Jesus is the way to heaven..
Yes, Jesus is the way to the Father..
Jesus is the way to relationship with God.".
Yes, that comes through by grace and through faith, and you don't have to earn it,.
and you don't have to work hard, and you don't have to perform..
It's not about that..
It's about Jesus cleansing you, forgiving you, and changing your life..
That message is certified and stamped as real by the resurrection itself..
That answers a lot of the questions we have about our own worth and our identity,.
and are we good enough, and do we have to be good enough, and how do we earn it..
All those questions are answered because the resurrection says, "No, you don't have to earn it..
Yes, God loves you..
Yes, you can be accepted..
Yes, you can be forgiven..
Yes, you can come into the kingdom of God because of what Jesus did.".
The next thing I want us to see, and this is powerful for us to consider in the world that we live in right now,.
is that the resurrection authenticates the future promises of God..
The resurrection authenticates the future promises of God..
In the book of Isaiah chapter 40, verses 3 to 5, it says this,.
"A voice cries in the wilderness, 'Prepare the way of the Lord..
Make straight in the desert a highway for our God..
Every valley shall be lifted up, and every mountain and hill shall be made low..
The uneven ground shall become level, and the rough places a plain..
And the glory of the Lord shall be revealed,.

$^{281}$and all flesh shall see it together, for the mouth of the Lord has spoken.".
That's a beautiful promise..
When we look at our world and we see so much injustice, so much suffering, so much pain,.
and so many unlevel places, so many crooked paths,.
God has promised that one day in the future, He's going to make every crooked path straight..
He's going to fill in every valley..
And that's metaphorical language that's saying that God is going to fix every injustice..
God is going to right every wrong..
God is going to balance out the scales of justice once and for all..
That is a beautiful promise..
In Revelation chapter 21, verses 1 through 6, it says this,.
"Then I saw a new heaven and a new earth..
For the first heaven and the first earth had passed away, and the sea was no more..
And I saw the holy city, New Jerusalem, coming down out of heaven from God,.
prepared as a bride adorned for her husband..
And I heard a loud voice from the throne saying, 'Behold, the dwelling place of God is with men..
He will dwell with them, and they will be His people, and God Himself will be with them as their God..
He will wipe away every tear from their eyes, and death shall be no more..
Neither shall there be mourning, nor crying, nor pain anymore, for the former things have passed away.'.
And he who was seated on the throne said, 'Behold, I am making all things new.'.
Also he said, 'Write this down, for these words are trustworthy and true.'.
And he said to me, 'It is done..
I am the Alpha and the Omega, the beginning and the end,.
to the thirsty I will give from the spring of the water of life without payment.'".
Man, what a promise..
What a promise..
There will be a time in the future when God will make everything new..
God will dry every tear..
There will be no more pain, no more disease, no more COVID, no more cancer,.
no more suffering, no more grief, no more injustice, no more racism, no more poverty..
It will all be gone..
And that God will dwell with us..
That's the most beautiful part of the promise..
We were created for relationship to God, and our sin has broken that relationship,.
and out of that brokenness comes all of the terrible and evil things that we see in the world.
that we do to ourselves and do to one another..
But this promise is that God is going to restore that relationship,.
and He's going to dwell with us..
We're going to be together like we were created to be..
What a beautiful promise..

$^{321}$Now, without the resurrection, we could say,.
"That's some beautiful language. That's some poetic stuff, and that sounds great..
Yeah, I really want that to be true, but I'm not sure that it actually is.".
But the resurrection authenticates the future promises of God..
And it says, "Yes, look at the power of God at work in Christ,.
and the power of God at work on the cross as Jesus dies for all of humanity.
and offers forgiveness and grace to everyone,.
and the power of God at work to raise Jesus from the dead,.
to defeat sin and to defeat death once for all.".
God has the power to do that,.
and that authenticates the fact that God has the power to make all things new..
That when Jesus says He is the beginning and the end,.
that He's the first and the last,.
and that He can make all things new,.
that it's true, that it's real, that it's actually going to happen..
The resurrection authenticates that..
And you might say, "Well, that's great..
I mean, the promise of heaven is beautiful. It is beautiful.".
You might say, "But is that enough?".
That's something that the Church has grappled with through the ages,.
and people have at times accused the Church of being too focused on heaven.
and not enough on this world..
And that's a legitimate criticism..
People can say, "Well, Christians shouldn't just sit around and say,.
'Hey, one day it's going to be great. One day it's all going to be fixed..
One day it's going to be better in heaven,.
and we'll just wait until we get there, and in the meantime,.
we'll just live in this world of suffering and brokenness.
and all the terrible things that happen.'".
What about right here, right now?.
And that's a very valid question that we have..
And I want us to see that the resurrection.
doesn't just authenticate the future promises of God,.
but it also validates the present imperative of living kingdom values..
The resurrection validates the present imperative of living kingdom values..
We've been talking about kingdom values for a while now.
as we study the Sermon on the Mount..
And Jesus has been teaching us in that sermon.
how important it is that we live out the values of the kingdom..
He's taught us that ultimately our lives are supposed to be.

$^{361}$so saturated with the values of the kingdom that it's like a sponge..
Whenever you squeeze a sponge, whatever that sponge is saturated with.
is what's going to come out..
And our lives are supposed to be so full and saturated.
with the values of the kingdom that when we're squeezed,.
that the values of the kingdom come out into the lives of the people around us,.
into the world around us..
Love, joy, peace--those are the things that are supposed to come out..
And we look at our world and we say, "What about poverty?.
Does Jesus care about that?".
And the resurrection says, "Yes, Jesus cares about poverty..
Jesus cares about the suffering of people in our present reality.".
"What about racism? What about injustice?.
Does Jesus care about those things?".
"Yes, Jesus cares about those things.".
"Jesus cares about human suffering.".
"What about oppression and people who live in places.
where they're oppressed by evil governments or by armies.
or people who come in and invade them and attack them.
and all these terrible things that happen in our world?".
The resurrection says, "Yes, Jesus cares about that.".
That's why He came. That's why He was here--.
to bring healing into our lives and to bring healing into those things..
We talked about that last week--that Jesus cares about the brokenness.
of the systems and structures in our world, and He wants to heal them..
And He wants to bring healing into our hearts..
Jesus came to heal the world one heart at a time..
And as He heals us, and as our lives get saturated with the values of the kingdom.
and start to flow out into the world around us,.
that a lot of the brokenness, the systems and structures,.
and all those terrible things begin to change as well..
That's what Jesus wants for us..
He wants us to be healed, and He wants us to be an instrument of healing in the world..
Jesus does care about those things..
And that brings us to say, "But is Jesus enough?.
Is Jesus enough to heal our world?".
We say, "All right, Jesus, I hear what you're saying about peace and love and forgiveness..
That's good stuff. But you don't know how mean these streets are..
You don't know how tough it is out here. You don't know how bad it is where I live.".
And Jesus says, "Yeah, I do. I came. I lived it. I walked it..

$^{401}$I was killed by powerful forces that wanted me out of the way..
And they're the people who hung me on the cross because I was a threat to their power..
But I defeated their power, not through force and coercion and greater power..
Even though Jesus had greater power, He defeated the powerful through love..
And that's what Jesus is showing us in the cross and in the resurrection..
And that brings us to another question that we have..
We may say, "Okay, is Jesus enough? Is the love of God enough?".
The answer is yes. The resurrection tells us yes..
But then we say, "But God, is it worth it? Is it worth it to live out the values of the kingdom?".
We've been talking about these for a while, like I said..
And they are tough to live out..
Humility, deep humility that causes us to see our own poverty of spirit,.
our own brokenness, treat others with meekness, mourn over sin and brokenness in our world..
That's tough to live out. Is it worth it?.
Being a peacemaker in a world that's full of violence and vengeance and hatred,.
it is tough to be a peacemaker..
Is it worth it to try to be a peacemaker when everybody else is throwing rocks at you?.
Is it worth it? That's a valid question..
Forgiving others, is it worth it to forgive someone else when they've wronged you?.
Is it worth it to let them off the hook, to let go of that anger and frustration that you feel towards them?.
Jesus called us to forgiveness. That's a kingdom value. But is it worth it?.
What about seeking purity of heart?.
What about loving our enemies?.
I was right here in the Sermon on the Mount. Man, that is tough..
What about trying to honor God in all the practical ways of our life that Jesus has talked about,.
our finances He talked about here in the Sermon on the Mount?.
How we're supposed to lay up treasures in heaven, not on earth..
We're supposed to be heavenly minded with what's valuable to us,.
instead of being earthly minded and trying to just make money and stack it up..
What about honoring God with our relationships and with our sexuality?.
Those things are really difficult..
Is it worth it to try to obey God and try to live out the values of the kingdom?.
The resurrection says yes. The resurrection validates the present imperative of living the values of the kingdom..
The resurrection validates that yes, the sacrifices that we make to obey God.
and to seek His kingdom first and to honor Him in our lives are worth it..
You know, when we think about these questions and we say,.
"Man, I have so many questions about following God.".
Think about the disciples and the questions they had..
They had just spent three years of their life devoted to Jesus and to following Him..
They quit their jobs, they left their families, and they spent every moment with Jesus..

$^{441}$And then they had seen Him crucified on the cross..
And it felt like everything that they believed in,.
everything that they had built their world around, was gone and had come crashing down..
No wonder they were depressed..
No wonder they were hiding out in a locked room..
That's what I would have been doing. I wouldn't have wanted to move in that situation..
But then what happened? The resurrection happened..
And suddenly it answered every question they had..
They must have had so many questions..
"Why, God? Why did this happen?".
"Why did you even call us in the first place that this was going to happen?".
"Why did you let this happen?".
"Why did you let the Romans and the Jewish leaders kill you?".
"It was so unjust. It was so wrong. Why?".
But the resurrection answered every question..
And it validated all the things that Jesus had taught them..
And suddenly they saw that this stuff is real..
And Jesus was telling us the truth the whole time..
And that's why they went out and lived so powerfully..
And through them and through the early church,.
the world was literally turned upside down, as the book of Acts says..
Because the teachings and the truth and all of the things that Jesus had taught them.
were validated by the power of the resurrection..
When they saw that, they knew it's all true..
It's all true. And they went out and lived like it's all true..
And that is so important for us to see in our world..
In our world that's built on strength and brute force and power,.
Jesus came and conquered with His love..
And He conquered with forgiveness. And He conquered with grace..
And He conquered with redemption..
And He conquered sin and death once and for all..
Where does that leave us?.
If you're someone who's watching today and you say,.
"Okay, well, I hear what you're saying and it sounds really good,.
but I don't know. I don't have any relationship with Jesus..
The kingdom of God, that's not really something that I understand or that I'm a part of.".
The scripture is really clear that you can come into the kingdom at any moment..
That you don't have to do anything other than believe in Jesus and put your trust in Him..
You couldn't earn it anyway. I couldn't earn it. None of us could earn it..
All of us are sinners. All of us have made mistakes and fallen short and messed up..

$^{481}$All of us have broken our relationship with God that we were created for because of our mistakes..
But God loved us so much, He said, "I want to restore the relationship.".
He sent His Son to earth. His Son, Jesus, who didn't sin, who lived the perfect life.
and at the end of His life laid His life down in our place as a sacrifice for us..
Took the wrath of God that we deserved on Him for our mistakes and our sin.
and our terrible, evil stuff that we do..
Jesus took it all on Him. Suffered the penalty, paid the cost.
so that we could be forgiven and we could be made free..
And all we have to do to experience that is believe it and say,.
"God, I believe that Jesus was the Son of God. I believe that He died for me..
I believe that He offers forgiveness that can wash away all of my sins and I want it and I'm asking for it.".
If you trust and believe that, then you come into the kingdom of God..
You come into the family of God..
Romans 10, verses 9 and 10 says this,.
"If you confess with your mouth that Jesus is Lord and believe in your heart that God raised Him from the dead,.
you will be saved. For with the heart one believes and is justified,.
and with the mouth one confesses and is saved.".
If you've never done that, if you've never believed in your heart,.
if you've never confessed with your mouth that Jesus is the Lord,.
that He has forgiven your sins and paid the penalty for you,.
you can do that right now. You can do that today..
You can enter the kingdom of God..
But a lot of you are probably people who have already entered the kingdom of God,.
but maybe you still have those questions..
You still have those things, those hurts, those pains of life..
And you say, "God, I don't know if I really want to follow You in this area..
I don't know if I really trust You with this part of my life or this part of my heart or this part of my experience.".
And to all of us who are in that place, I would say that Jesus is calling us to surrender to the values of the kingdom.
and to really seek to live them out as if they are true because they are..
And the resurrection validates that..
Remember Matthew 6, verse 33, right here in the Sermon on the Mount, Jesus says,.
"Seek first the kingdom of God and His righteousness, and all these things will be added to you.".
It's a nice sounding statement, but what if it were actually true?.
How would that affect the way you live your life?.
And the resurrection says it is true..
You can seek God's kingdom and surrender everything to Him..
You can seek to have your life saturated with the values of the kingdom..
And even though it may lead to sacrifice at times, it may lead to persecution at times,.
it may lead to difficulty at times, it will also lead to joy..
It will also lead to amazing transformation in your heart and life..

$^{521}$And it will be worth it. It will be worth it in this life, and it will be worth it one day in heaven..
It will be completely worth it..
And the resurrection validates that..
So I want us to see that today, that the resurrection answers those questions that we have..
It answers them affirmatively..
It says, "Yes, God loves you. Yes, God cares. Yes, God has a future for you.".
And it's beautiful and wonderful, and God's going to dry every tear..
But yes, God also cares about the present, and He wants good and blessing for you..
And He wants you to know Him and walk with Him and experience relationship with Him..
And He wants you to experience healing in your heart and in your emotions and in your past and throughout your life..
And He wants you to experience transformation that allows you to be an instrument of healing in the world around you..
And the resurrection validates and affirms and authenticates all of that..
So we have a message of hope and of joy today when we consider the resurrection..
Let's pray..
God, thank You that Jesus is risen..
Thank You that He came and gave His life for us on the cross, but that was not the end of the story..
That He defeated sin and death once and for all, and that now He reigns on high in heaven..
God, thank You for the power of the resurrection..
May the power of it be real to us..
May we believe it and may it validate and authenticate all of those values of the kingdom that You have been pouring into our hearts..
And may we begin to live like those things are actually true..
May we begin to live them out..
May we become agents of healing in our world..
In Jesus name, amen..
[BLANK AUDIO].
\newpage



\section{}
\label{sec:6bRdtM__Fr0}
\textbf{愈靠近,愈堅強 - 張麗仙傳道}
\newline
\newline
連結: \href{https://youtube.com/watch?v=6bRdtM__Fr0}{\texttt{https://youtube.com/watch?v=6bRdtM\_\_Fr0}} ~~~~ 語音日期: 2024-05-09
\newline
\newline
\hyperref[sec:EISv2Eie8Pg]{\small{< < < PREV SERMON < < <}}
~
\hyperref[sec:index_chronic]{\small{[返順時目]}}
~
\hyperref[sec:index_scriptual]{\small{[返順卷目]}}
~
\hyperref[sec:D_DiJk8_Dhc]{\small{> > > NEXT SERMON > > >}}
\newline
\newline
$^{1}$就是我們今年全年的主題.
Mark Sir 過去幾個禮拜都提過這個題目.
是我們今年的教會一個很重要的主題.
叫做「同健健康教會」.
我想問你喜歡這個主題嗎?.
不喜歡.
不過有時我覺得我這樣問好像很幼稚.
吃雪糕喜歡嗎?.
其實應該問你.
你覺得這個主題重要嗎?.
重要嗎?.
重要.
健康當然重要.
我們每個人都覺得健康是重要的.
所以我們個人.
我們才會煲這麼好的湯.
我們才會身水身汗都跑.
我們會耍太極.
因為我們知道健康是很重要的.
不健康我們會很弱.
我們會病.
我們甚至會死.
個人的健康重要.
家庭的健康重要.
其實教會的健康同樣是這麼重要.
因為老實說.
如果教會不健康.
會弱,會病,也會死.
所以我們今年這個「同健健康教會」.
是一個很重要的題目.
不過幕後消息我想告訴大家.
在這個主題之前.
其實Mark Sir首先在年初的時候.
他有另外一個主題.
就是這個.
Closer and Stronger.
當Mark Sir和我說這個題目的時候.
我覺得很棒.
我自己為什麼覺得很棒呢?.
因為我覺得是一個很.

$^{41}$很導入.
一想起Closer and Stronger.
我想起很多東西.
我覺得非常好.
但我不知道為什麼後來Mark Sir會用了.
「同健健康教會」.
但是我知道他在看的是一個更大的圖畫.
而我也知道.
只要我們Closer and Stronger的時候.
我們就有一間健康而且茁壯的教會.
這個健康和茁壯的教會.
其實是有很多方面的.
你同不同意嗎?.
講道.
教道.
傳道.
還有一樣.
不知道你有沒有想到是什麼呢?.
就是人與人之間相處之道.
重要嗎?.
是非常之重要.
人與人之間相處之道.
不單止對教會重要.
對我們生活每一個範疇都是這麼重要.
這個題目真的可以講十天都講不完.
但是在今早.
我在這裡想跟大家講的是主要的兩點.
因為今年進入2023年.
我在姐妹團契那裡.
是和大家一起去研讀約翰一書.
就發覺約翰一書就是講.
就是很著力講人與人之間相處之道.
所以在今早.
我很希望能夠在約翰一書那裡.
特別提出兩點相處之道.
是神很希望教會做得到.
做得到的時候.
我們就能夠建立一個健康茁壯的教會.
第一點.
是神的心意.

$^{81}$原來就是想教會.
closer and stronger.
大家都知道close是什麼意思.
close是接近.
接近.
strong是強壯.
closer and stronger.
當你更加接近的時候.
你們就更加強壯了.
約翰一書很有意思.
因為一開始的時候.
作者約翰就說了一句話.
從起初到現在.
原有的生命之道.
就是我所聽見.
我所看見.
親眼看過.
親手摸過.
為什麼他這麼說呢.
他這麼說的時候.
一方面是強調.
耶穌是真完全的神.
成為完全的人.
來到我們世界.
和我們接近人.
但是另外一方面.
他也想強調的就是.
因為我見過摸過看過.
我甚至和祂住過.
所以.
我知道祂的心想什麼.
然後他就說了這句話.
他說我將我所看見.
所聽見的傳給你們.
為要使你和我有團契.
想想團契是什麼意思.
一會兒我會詳細說.
而我們的團契.
是和父和他的兒子.
耶穌基督所共有的.

$^{121}$我們將這些事寫給你們.
使我們的喜樂可以滿足.
在這裡約翰對他的讀者.
修信人這麼說.
他說我從耶穌身上.
我親眼看過.
親耳聽過他.
我明白得到他.
他是想和我們.
有一個親密的關係.
在這裡用團契兩個字來代表.
我神父和兒子.
很希望和我們每一個人.
有親密的關係.
而他也希望我們和你們.
就是說每一個信的人.
無論是在耶路撒冷的信徒.
使徒或是在外邦.
小亞細亞歐洲各地的信徒.
一樣我們和你們有親密的團契.
換句話說.
神今天仍然希望.
世界各地的弟兄姊妹.
各地信他的人.
跟隨耶穌的人.
都是有親密的團契.
團契這個字真是怪怪的.
我保證你沒回來教會.
是沒聽過這個字的.
團契這個字是在教會聽到的.
因為平時都是說.
Club 同鄉會 中親會.
團契這個東西.
回來後很快我們就說.
你回來了跨世代團契.
你回來了即稱團契.
你回來了姐妹團契.
但是團契是什麼意思呢.
其實這個字是很有意思的.
我們都知道.

$^{161}$《新約聖經》是希臘文寫的.
希臘文Koinonia這個字.
是指團契.
團契就是這個字.
什麼意思呢.
是很親密的關係.
很親密的關係.
有脈共享.
心與靈深入相交.
有共同的特質.
有活的愛.
我們解讀一下.
用現代話說.
有脈共享.
即是你的即是我的.
不是很想.
心與靈深入相交.
即是無所不談.
說話不怕得罪大家.
有共同的特質就是.
我們分享同樣的價值.
大家的價值觀是一樣.
都是基督裡面.
有活的愛.
愛是不會延期的.
不會到有效期.
不會死的.
愛.
活的愛在我們中間.
這個是團契的意思.
你回頭看就會覺得很奇怪.
原來約翰一開始就說.
我很知神的心意.
祂是想和我們很親密.
親密到這樣.
而祂也想.
所有弟兄姊妹都是這麼親密.
原來這是神的心意.
你有沒有留意到.
剛才聖經最後說了.

$^{201}$我寫這些給你們.
是我的喜樂就得到滿足.
即是我說這些給你們聽的時候.
我很開心的.
其實我說這些給你們聽.
我也很開心.
因為神竟然願意和我們這麼親密.
祂希望我們是這麼親密.
如果我們做得到是多麼好.
請大家也留意背景.
當約翰寫這些說話的時候.
當時在小納西亞.
當地的基督徒.
都是因著信仰受到很嚴峻的逼迫.
就像我們今天這麼舒服.
他們每個人都受到很多不同方面的掙扎.
而約翰將神的心意.
這個親密關係的心意.
告訴他們的時候.
其實就是想說.
你們在困苦掙扎裡面.
你們要更親密和強大.
你更親密.
你就可以強大.
就好像我們中國人一樣.
團結就是力量.
就是這個意思.
我們今天.
似乎我們不是受到嚴峻的宗教逼迫.
但是其實我們人生各自.
你和我.
每一位在這裡.
都是有來自各方面的人生衝擊.
而在這些衝擊挑戰裡面.
今天神同樣的告訴我們.
來到教會.
這個地方是不同的.
我們在這裡真誠付出.
建立親密的關係.
這個是神希望教會所有的.

$^{241}$這樣才是一個健康的教會.
Getting closer, we will be stronger.
今早我以為其實很少人.
因為睡了一小時.
我都怕自己起不來.
但是見到這麼多人.
我都很開心.
其實過去一個月.
教會回來的崇拜的人數.
越來越穩定.
越來越多.
其實很開心.
很興奮.
還記得疫情的時候嗎?.
我們在手機上見面.
我們在電腦上見面.
在屏幕上見面.
就算在街上買菜.
去超級市場買菜.
見到面都是這樣.
還蒙了面.
只顧著猜猜那個是誰.
但是.
真是這樣.
在疫情的社交距離的要求.
是真的沒有辦法close的.
只有father and wicker.
越疏離我們就越軟弱.
疫情的日子是這樣.
但是在這裡.
在今天.
我很願意在這裡.
因為在疫情這幾年期間.
神在城會教會做了一些很奇妙的美事.
其中一樣我很想跟大家說的.
很想跟大家見證的就是.
神在我們這裡.
建立了前所未有的親密的關係.
這些關係.
在一些團契.

$^{281}$在教會的祈禱會.
我們見到.
你會體會得到.
我們在這幾年.
關係沒有疏離了.
而且是親密了很多.
我說很多團契或者是祈禱會.
我自己在教會.
也是負責姐妹事工的.
我就想用姐妹事工.
姐妹團契那裡.
作為例子.
見證神的恩典.
真的在這些日子.
姐妹的關係比以前親密了很多.
以前來就是hi and bye.
很多姐妹都是這樣說.
現在來到好像回家了.
我想一下為什麼會這樣呢.
兩個原因我想得到.
第一就是.
姐妹面對的困難是很多.
我相信很多人面對的困難都很多.
但是我見到這班姐妹.
各自各面對很多困難.
其中包括.
有丈夫或者是子女.
長期病患.
是24/7要照顧.
真是身心力疲乏到不得了.
有些離婚單親的.
他感到很徬徨.
有些很少親人在這裡.
感到孤單無助.
還有一些姐妹在這段日子.
是經歷很多很多疾病.
怎樣繞來繞去都跳不出來.
實在是很困難.
即使我真的有時候覺得.
姐妹是挺有韌力的.

$^{321}$但是也是很困難.
我作為女傳道.
當然會.
我心裡真的很希望盡力的幫.
但是正如每一位教母都同樣體會得到.
我們一個人幫得的是很少.
就算我們盡力.
我們聆聽.
我們開道.
我們祈禱.
甚至我們載他們去醫院.
載他們幫他們做他們需要的事.
我們都是做得很少.
但是在這班很有困難的姐妹中間.
我很開心見到一些這樣的現象.
就是有一些姐妹.
是很主動的站出來幫.
開車接送.
煮食.
安慰.
陪伴.
不是這些人特別有空.
往往是他們越不空閒的人.
越肯這樣做.
他們能夠累得站在那裡.
都起來了.
走出去幫人.
去幫人.
是因為基督的愛在他們裡面.
他願意甚至調動自己工作的時間.
甚至放下個人的享受和休息.
都這樣做.
我體會到兩件事.
第一就是風花雪月的日子.
是很難或者不容易栽種到親密的關係.
往往是說中送歎的時候.
可以做得到.
是不是.
我也體會到的就是.
愛與恨在一個群體裡面.

$^{361}$絕對是有傳染性的.
如果一個群體裡面.
很喜歡說是非的.
這個群體就像整個是非圈.
有幾個人喜歡說是非.
就像整個是非圈.
但是如果這個群體有幾個人.
能夠真的放下自己.
出心出力去幫人的話.
這個群體就是一個.
很有愛互相守望的群體.
結果不單止幫人的人.
被幫的人是拉得近.
而且整個群體.
都是拉近.
closer and stronger.
我在說什麼呢.
我在說.
這一兩年.
姐妹之間的關係.
能夠這麼親密.
是因為.
大家面對的困難很多.
而反而成為了一個很親密的關係.
因為大家互相幫助的精神.
還有第二個原因就是.
我發覺這些姐妹.
是很願意參與和分享的.
我們姐妹團體.
大家都知道叫姐妹船.
每個月月會我們叫大船.
有時候小組聚會.
我們叫上小船.
上船.
我誤會上什麼船了.
最近有人.
是外人來的.
還給了一個很好玩的名字.
你們姐妹又去遊船河.
什麼意思呢.

$^{401}$就是我們飲飲食食.
我發覺.
姐妹為什麼會這麼親密呢.
原來飲飲食食都很重要.
真的很重要.
你都知道.
姐妹煮東西是特別厲害的.
有幾個高手在已經很厲害了.
每次吃東西.
聚會都很精彩.
不過精彩的.
不只是分享食物.
往往是吃完東西之後.
大家坐下來.
分享人生.
他們內心的掙扎.
是很真實的.
大家看到.
很欣賞身邊的姐妹的幫助.
從她們的幫助.
看到神的恩手臨到她.
我說精彩的意思就是.
不是食物精彩.
什麼時候吃東西.
吃到流眼淚.
你試過吃東西吃到流眼淚嗎.
很難.
多好吃都只是笑.
但分享的時候.
大家會流出眼淚.
真誠的眼淚.
這是很精彩的地方.
能夠在這兩年.
姐妹之間建立這麼親密的關係.
而且不是開玩笑的.
還有一個例子.
這兩年Crucify and Stronger.
是名副其實地發生.
舉個例子.
去年有一位住在Northside的姐妹.

$^{441}$特意搬出來唐人街住.
因為近姐妹的群體.
特意搬出來.
另外在過去.
現在數回來.
還不到一年的時間.
有三位姐妹分別搬屋.
特意搬到同樣的街口住.
是不是Crucify and Stronger.
吃東西都好吃一點.
還有的.
以我所知還有的.
大家是想Crucify and Stronger.
我想Mark Sir定這個題目的時候.
沒有想到姐妹這麼認真.
搬屋.
我們非常感恩.
《約翰一書》提到的親密團契.
在姐妹族中.
我們真正體會到.
不過我們只是在學習.
不過我想更加強調.
Crucify and Stronger.
並不是姐妹專利.
弟兄的需要.
我希望弟兄能夠有一個這麼親密的團契.
年青家庭需要.
年長家庭需要.
職青需要.
少年人需要.
甚至我們教務的團隊需要.
我們的Call Team需要.
我們需要Closer.
我們就會Stronger.
如果你在教會裡面體會到這個親密的關係.
你要感恩.
你要珍惜.
你要繼續付出再做.
但是如果你沒有體會到的話.
我覺得其實你可以這樣想.

$^{481}$就是可能神是呼喚你.
從你身上開始.
你去付出.
在教會裡面共建.
一起建立.
這個這麼親密的關係.
這個團契就在你中間.
不過我們都可以預料到.
就是這些親密團契.
是會時時受到攻擊的.
誰攻擊我們呢?.
大家都知道.
是魔鬼.
魔鬼絕對不想有這樣的群體出現.
他不想我們Closer and Stronger.
他想我們Father and Weaker.
可悲的是.
他很成功.
由伊甸園開始.
他要拆散阿當和夏娃親密信任的關係.
以致到人類一直到現在.
世界各地古今中外.
全部人都是.
我們的關係都被他搞得很慘.
好的時候.
他就想Break up我們.
一切的關係.
我們都時時受到試探.
而往往撕裂.
兄弟成仇.
夫妻決裂.
朋友攻擊.
甚至在教會之內.
都時時有很多抗拒的事件出現.
我不再說國與國.
假象鬥爭那些了.
事實上.
說回《約翰一書》.
《約翰一書》.
約翰這麼強調親密的團契.

$^{521}$一開始就說親密的團契.
是因為當時背景就是.
魔鬼正在當時的教會那裡.
要搞亂教會的關係.
而藉著當時他們社會流行的.
希臘學說.
洛斯底主義.
有些人覺得這個主義很棒.
有些人就覺得它是異端.
大家就爭辯不休.
使到教會裡面的弟兄姊妹各自站硬.
互相指責.
甚至不少人離開.
這個是當時的現實.
所以以上我即使舉姊妹選擇為例子.
Closer and Stronger.
老實說.
我都時時會為這個團契而感到警戒.
因為我知道魔鬼一定無孔不入.
一定的.
預了它.
要拆散親密的關係.
老實說.
Closer不是好玩的.
Closer是比較危險的.
是不是.
你遠一點.
大家保持遠一點.
就不會有磨擦.
怎麼磨呢.
你這麼遠我怎麼磨擦你呢.
越遠就越安全.
是不是.
但是越遠.
越軟弱.
不是神的旨意.
Closer and Stronger是神的旨意.
不過我又經常覺得.
磨擦衝突是不要緊的.
最重要的是我怎樣去面對.

$^{561}$怎樣去處理.
如果處理得好.
我們的關係會健康成長.
可能比之前更好.
如果處理得不好.
不死都剩下半條命.
老實說.
很弱的.
藉著約翰一書.
不單止約翰提出這個問題.
他也給了一個秘訣.
就是如果我們要緊守Closer and Stronger.
我們是有一個這樣的秘訣.
也讓我帶大家去看.
這裡看的第二點.
第一點是神的心意.
是他的兒女要Closer and Stronger.
第二點就是.
捨己相愛就是秘訣.
能夠使我們緊守.
堅守Closer and Stronger.
如果我們繼續看約翰一書.
我們會看到.
約翰很妙.
不停地說一個命令.
就是耶穌當日在世界上.
臨死之前.
跟門徒說了很多話.
其中說了一個重要的命令.
他說我賜給你們一個命令.
就是你們要彼此相愛.
繼續聽下去.
我怎樣愛你們.
你們就是要怎樣相愛.
耶穌說的不是啪一聲的相愛.
而是不是一些普通道德倫理的相愛.
是捨己的相愛.
好像祂愛我們一樣.
彼此相愛.
這是能夠解決紛爭的唯一秘訣.

$^{601}$耶穌還沒走之前.
祂也知道門徒之間有很多紛爭.
所以祂臨走之前.
祂一定要說定這件事.
而且說定這件事.
是給以後跟隨祂的人聽的.
有意思的就是.
真的不是隨便說說.
約翰一書.
約翰繼續說.
是用三大段來說彼此相愛.
我不是跟大家讀完.
不用怕.
第一段是第二章.
那裡有五節說到彼此相愛.
第二段是第三章.
那裡有十四節說到彼此相愛.
第三段是第四章.
那裡有十五節說到彼此相愛.
厲不厲真.
他是覺得這件事非常重要.
說三次.
說得這麼長.
當然重要.
而且說得這麼長.
是因為怕死我們做不到.
所以說完又說.
因為跟同聲同氣的人.
你愛他是容易的.
跟那些你看得很不順眼.
他敢選的人.
你是覺得很想抗拒他.
不想黏在一起.
這就是令你father and weaker.
而魔鬼往往就是用了我們這種自我出現.
我們不喜歡人就是很多時候.
ego 自我.
不喜歡我不接受你.
那你又不接受我.
我們不想吵架.

$^{641}$大家避之則吉最好.
就是這樣.
我們的自我很多時候使我們成為像箭豬一樣.
因為看不到.
其實我們很多時候真的像箭豬一樣.
要不就你刺我.
我刺你.
要不就是你和我遠距離.
不要刺了.
不要刺來刺去.
就是這樣.
我們的關係很疏離.
Amazing.
約翰一書一而再再而三提醒信徒.
我們要對付這種自我介離的衝突.
唯一的利器來自耶穌.
就是耶穌當日不單止說.
而且他是實行的.
捨己的相愛.
耶穌說我怎樣愛你們.
你們就要怎樣相愛.
有時候真是彼此相愛.
真是你要覺得.
我很不喜歡.
我也放下我自己.
這就是耶穌說的捨己.
耶穌說的捨己.
不是一定要我們去釘十字架的捨己.
就是這樣.
我們放下自己.
也是非常不容易.
放下我們自己的手.
我們就可能會能夠.
Put ourselves in the other's shoes.
我們能夠從別人的角度去想.
就是這樣.
我放下自己.
我們能夠從抗拒.
定人的罪.
變成包容和寬恕.

$^{681}$漸漸地我告訴大家.
我們自己的心會覺得釋放.
我們覺得我們周圍的環境.
我們能夠建造一個溫暖的環境出來.
你會發覺你不喜歡的人.
都會慢慢地改變.
這就是耶穌彼此相愛的威力.
祂能夠使我們在這個.
針對罪而震動分裂的世界裡面.
能夠緊守.
Closer and then we'll be stronger.
Closer and stronger.
是神希望他的兒女有的相處之道.
教會健康的成長.
是講道,教導,傳道.
還有這個相處之道.
希望我們大家都緊記.
我們在這樣的時候.
我們的教會就會健康,茁壯地成長.
在結束的時候.
我想起一句.
當我預備講道的時候.
經常在我腦海中會出現一句說話.
這句說話是一位國語音樂的歌手說的.
我沒有聽他唱歌.
但我覺得他這句說話很有意思.
讓大家看看.
他說.
We cannot direct the wind.
But we can adjust the sails.
就是說我們不能夠指揮風向.
但是我們可以調整我們的風範.
如果我們在船上.
這個世界風能夠吹過來.
我們不知道在哪裡.
無論那些是甚麼人生的風暴.
颶風.
關係上的,身體上的,政治上的.
甚麼都好.
風暴難免.

$^{721}$無論在我們任何的年齡階層.
在任何的地方.
教會,家庭.
或是工作的地方.
我們都會面對突然打來的風暴.
但是在風暴裡面.
如果我們能夠將我們人生的帆.
轉向靠近大家.
駛去那裡.
駛去靠近大家.
當我們靠近的時候.
我們就會發覺我們更堅強.
我們將我們的帆去靠近.
靠近大家.
我們就會變得更堅強.
神的心意是希望我們這樣的.
如果我們有這樣的心的時候.
神必定幫助我們.
我們一起祈禱.
雖然我們是那些.
那麼有罪性的人.
但是從創世到現在.
從我們墮落在罪裡面.
你就一直的希望.
跟我們接近.
今天我們在約翰的書裡面看到的.
你說跟我們有很親密的團契.
我們發覺原來真的不是說.
你真的.
你意思是.
你想要跟我們很接近.
而你也真是期望我們.
每一位屬於你的人.
我們互相去接近.
我們不要給這個世界上.
很多自己的利益.
或者我們的自我.
或者我們的一些價值觀.
一些的判斷.
讓我們去將人推開.

$^{761}$相反.
教我們在人與人相處的關係裡面.
我們真的能夠用耶穌基督的愛.
捨己的相愛.
彼此的接納.
我們能夠無論在教會.
在我們家裡.
我們在任何的地方.
我們都去建立一個這樣的相處之道.
我們付出愛.
我們跟人接近.
而我們建立一個強大的關係.
求你幫助聖會教會.
我們在今年說到.
我們要建立一個健康的教會的時候.
你幫助我們.
聖靈幫助我們每一位.
我們真的能夠聽你的話.
彼此相愛.
我們一起在這裡建立.
這個健康茁壯的教會.
榮耀你的名.
以致我們更加吸引更多更多.
失望在這個世界裡面的人.
他們見到這個愛.
能夠歸於你.
主求你幫助.
奉救主耶穌的聖名.
阿們.
好 願神說話.
祝福大家.
我們今天到這裡.
謝謝大家收看 再見.
\newpage



\section{}
\label{sec:D_DiJk8_Dhc}
\textbf{Everlasting Praise 4 (1 Kings 6-7) - Rev. Mark Chan--}
\newline
\newline
連結: \href{https://youtube.com/watch?v=D_DiJk8_Dhc}{\texttt{https://youtube.com/watch?v=D\_DiJk8\_Dhc}} ~~~~ 語音日期: 2024-05-14
\newline
\newline
\hyperref[sec:6bRdtM__Fr0]{\small{< < < PREV SERMON < < <}}
~
\hyperref[sec:index_chronic]{\small{[返順時目]}}
~
\hyperref[sec:index_scriptual]{\small{[返順卷目]}}
~
\hyperref[sec:NNhx4LQjZTc]{\small{> > > NEXT SERMON > > >}}
\newline
\newline
$^{1}$But we will continue our theme of worship and let's continue to survey what the Bible.
has to say about worship..
And I told you from the beginning to the end it's all about worship because God is going.
to dwell in the midst of us, dwell in the midst of the people and He want us to have.
fellowship with Him..
So from the Garden of Eden to Tabernacle to today it is the Temple..
And when the Israelites settled down in the land of Canaan in the Promised Land, they.
have King David had Jerusalem to be the capital of the small nation..
And then God told Solomon to build the first temple..
And there was so much details in the Bible from the book of First Kings and talking about.
the details of the building of the temple and we're going to spend some more time in.
the first 11 chapters of First Kings..
So we have the Bible turn to First Kings chapter 1 to 11..
How about that?.
Okay 1 to 11..
There were tons of things, details about the reign of Solomon from First Kings chapter.
1 to 11..
Now let me ask you this question..
Okay what does a good king need to rule a nation?.
If you are the king, alright, or queen, alright, if you need to rule the nation, rule over.
the people, what do you need?.
What do you need?.
The Bible is trying to present to us King Solomon and everything that he needs to rule.
the nations..
Alright?.
So how about if you're a king, you need wisdom to rule the people, right?.
Get it?.
And you need wisdom to judge, you need wisdom to manage your country and that is what God.
has given to Solomon..
Okay?.
So let's take a quick survey of how the Bible describes King Solomon and then finally he.
builds a temple..
Alright, First Kings chapter 4, 30..
His wisdom, God has given him, exceeded that of all the wise men of the east and the wise.
men of Egypt..
His fame spread throughout all surrounding nations..
Now if you understand, you know, kind of the Bible map, Israel is right at the center of.
the ancient world..
Right?.
You've got Egypt, Babylon, Hittite, all the countries are surrounding Israel and Israel.

$^{41}$is right in the middle..
And first verse 31 it says, the name or the fame of Solomon spread throughout all the.
surrounding nations..
And he composed how many prophets and how many songs?.
3,001,005 songs..
Right?.
So he's a songwriter..
You know, Rachel told me, Taylor Swift composed all the songs for her new album, On the Row..
Right?.
How crazy is that?.
Taylor Swift..
Right?.
30 songs, like 30 songs..
On the Row, when she was doing the ERA tour..
That's crazy..
Right?.
But Solomon composed 1,005 songs..
Right?.
There was so much wisdom to this guy..
Other than wisdom, what do you need to rule the nation?.
How about money?.
Right?.
And defense..
Right?.
First King chapter 10, 26..
Solomon built up a huge force of chariots and horses..
He had 1,400 chariots and 12,000 horses..
And the king made silver as plentiful in Jerusalem as stone..
Now if you have visited Jerusalem, you know stones are everywhere..
Stone as pavement, stone hills, stone valley, because Jerusalem is in a mountainous area..
But here Solomon made silver as plentiful in Jerusalem as stone..
Right?.
You count how many stones, that's how much silver or gold he has..
All right?.
And then military might..
Right?.
12,000 horses..
That means perfect number of horses..
There's so much military defense Solomon has built up..
And what else do you need to rule the nations?.

$^{81}$How about international relation?.
Right?.
Just like United States has a lot of different alliance around the world..
Right?.
Solomon did have some alliance..
11.3..
He said Solomon has 700 wives of royal birth and 300 concubines..
How many of them?.
1,000..
Right?.
1,000 of them..
Right?.
How many wives he had?.
700 wives..
Now we always come to some description like this..
We always think, okay, it is about, you know, sexual desire..
Right?.
Well, it's part of it..
But the big part of it, why he has so many wives and queens?.
Because he is making diplomatic relationship with all the surrounding nations..
Right?.
They send a daughter or queen to Solomon and marry to him..
And that's how you build international relation in those days..
However, it has a little note there..
And these wives did turn his heart away from the Lord..
Alright?.
So, you got wisdom, you got economic, military power, and then diplomacy..
Right?.
And then in the Near East, all the kings, if they are good kings, they will build temple..
They will build temples for their biggest God..
Alright?.
So, that's what Solomon did..
And God allowed him to do that..
So, here 1st King chapter 6 verse 2..
The temple that King Solomon built for the Lord was 60 cubits long..
One cubit is about one and a half feet long..
20 cubit wide and 30 cubit high..
And it also said in verse 1, in what year?.
480 year after the Israelites had come out of Egypt..
So finally when Solomon built the temple, the Exodus is over..

$^{121}$Right?.
Remember the story of Exodus?.
The Israelites being saved by grace, by Yahweh, by God..
And they crossed the Red Sea and came through the wilderness..
And for so many years, finally they settled down in the Promised Land..
And King David was made the king over all the 12 tribes..
And he settled down in Jerusalem and make Jerusalem the capital..
And his son Solomon finally built the temple..
And here verse 1 of chapter 6..
480 year after the Israelites had come out of Egypt..
Finally they settled down in Jerusalem and Exodus is over..
So as you can see, the Bible spends a lot of words to describe the reign of a king..
And because the people in the Bible are in the king..
The king represents them..
The king, if it's a good king, then the people will be good..
Right?.
If the king worship other gods, then all the people, because they are being represented.
and they are in the king, they will follow other gods as well..
So after all these descriptions of Solomon, let us see how the Bible evaluate Solomon..
And one of the places that we can go to is the book of Deuteronomy..
The fifth book of the Bible..
Because Deuteronomy is like a constitution for ancient Israel..
And it lays out a lot of rules and regulations for the nation and for the king..
And part of the passage here is, well by the way, this is the model of King Solomon's.
first temple..
On the upper left is the city of David, which is on a small mountain ridge..
And up there on top is the temple and the palace complex..
That's the city of David..
Zion, or Jerusalem..
On the right is the model for first temple of Solomon..
What is the shape of the temple?.
Rectangular, right?.
It's rectangular..
And inside the rectangular, there's a square, right?.
Perfect square..
That is the holy of holies..
And on your lower left is the temple cross section, right?.
You see inside there's a square inside, the room inside..
There's a couple of creatures guarding the holy of holies..
Who are those guys?.

$^{161}$Cherubim, right?.
Cherubim..
If you know nothing about the series, Cherubim is your guy..
Anyways, we go to how the Torah, how the Bible evaluate a king..
One of the passages, very important passage, is about kingship..
Deuteronomy chapter 17..
God told his people through Moses, he said, "You are about to enter the land the Lord.
your God is giving you..
If this happens, which it did happen, be sure to select as king the man the Lord your God.
chooses.".
16, "The king must not build up a large stable of horses for himself..
The king must not take many wives for himself because they will turn his heart away from.
the Lord..
He must not accumulate large amounts of wealth in silver and gold for himself.".
That's interesting..
That's the constitution of ancient Israel..
It has so much to say about the king, right?.
When you elect a king, when he is going to rule, there are three dongs here, right?.
Can you see that?.
Do not build up your military might..
Do not have many wives..
It means international diplomacy..
Do not accumulate a lot of gold and silver, right?.
That's what the king of Israel is set apart from other kings surrounding nations..
Because he's going to depend on God, right?.
He's not going to have a lot of money accumulated for himself..
Not a lot of military might for himself and diplomatic relationship for himself..
Because he's going to trust God and God will protect him, right?.
God will protect his border..
God will fight for his battle, right?.
Interesting if you compare this to King Solomon..
King Solomon is doing contrary to the Torah..
And what does the king have to do?.
It continues to say, "If the king sits on throne as king, the king must copy for himself.
this body of instruction, the Torah on the scroll in the presence of the priest..
And he must always keep that copy with him and read it daily as long as he lives..
That way he will learn to fear the Lord his God by obeying all the terms of the covenant..
All the terms in the instructions and decrees, right?.
So that's so special about Israel's king..
Is that he is not going to do what other pagan kings would do..

$^{201}$He is going to trust in God and what does he have to do?.
He has to do Bible study..
He has to study the word of God..
He has to understand the will of God..
He has to copy himself a scroll of Torah and he has to keep a copy for himself and read.
it daily as long as he lives..
Because only by doing so, he is going to learn how to fear the Lord his God..
And he continues to say, "This reading, regular reading, will prevent the king from.
being proud and acting as if he is above his brothers, his fellow citizens..
It will also prevent him from turning away from these commands in the smallest way..
It will ensure that he and his descendants will reign for many generations in Israel.".
Right?.
You know what the word of God will do in your life?.
If you read it continuously, obediently, trying to understand and listen to the voice of God.
and trying to understand the will of God and submit yourself in front of, before the word.
of God, the Bible will do something in your life..
Right?.
It will make you less proud..
It will make you humble..
It will make you understand that your might and your resources are not from you..
It's from God himself..
It's from the God who has loved you, who has protected you..
You read the word of God daily..
You read the word of God continuously..
It will have effect in your life..
Right?.
That's why I want to encourage you to continue to read the word of God..
So let's talk about the king of Israel..
The top guy in the hierarchy..
The Torah says you have to have the king make a copy and continuously obey the terms and.
obey the voice of God and to read it daily as long as he lives..
And this will prevent him from being proud..
It will make him humble and he will continue to fear God..
And that is the blessing of Solomon..
Solomon is being elected as the king of Israel and God gave him so much blessing, the wisdom,.
the power, the economic resources..
That is the blessing he has..
But as we continue to dig into the description in the Bible about Solomon, you find that.
this king's life is a tragedy..
It's a tragedy..

$^{241}$And what happened here?.
Now my first part of the sermon, I talk about the blessings that Solomon had..
Right?.
And here is the second part, we're going to reflect on what is going on in his life..
What went wrong?.
The tragedy in King Solomon..
And let me remind you when we read the Bible, we read the Bible, we cannot know what the.
Bible means until we know how it means..
That means you have to understand the design of the passages, design of the chapters..
What different passages are being put together in such a way to communicate its meaning..
So we have to understand how, you know, 1 Kings chapter 1 to 11, how it's being designed.
carefully in order to tell you what is going on in the life of King Solomon..
Here let me show you the whole structure that is being suggested by some Bible students.
or Bible scholars, how the first 11 chapters talking about the reign of Solomon, the entire.
reign of Solomon is being captured in 11 chapters..
First King chapter 1 to first King chapter 11..
Now as you can see, this is a very common structure..
We studied the Gospel of Mark, right?.
Remember in the Gospel of Mark, there were many sandwiches, right?.
And the sandwiches, if you have a sandwich, the middle part is the content, right?.
It's the most important thing..
And so often it goes with, like you talk about the Bible author, we talk about A, and then.
you talk about B, and then he'll come back to talk about A again, right?.
So it's A, B, A, a very simple sandwich structure..
Here in first King chapter 1 to 11, it's a huge sandwich..
It's a multi-layer sandwich..
As you can see from chapter 1 to chapter 2, there was a prophet, Nathan..
He intervenes with the succession of kingship..
So in chapter 1 of first King, King David was going to die, right?.
He was totally disabled, dysfunctional to rule his nation..
And then they were talking about who's going to succeed him on the throne..
That is A. And then B, Solomon's being elected, being anointed to be the king after David..
And Yahweh, God eliminates all the threats to make sure Solomon's kingdom is secure..
That's B. And then C is the early promise of God to Solomon..
Giving him wisdom, giving him promises..
And then part D, chapter 3, verse 16 to 4, 34..
And you see Solomon receiving the blessing from God..
He uses his gift for the people and judge them fairly in justice..
And then part E is preparation for building the temple..
Chapter 5, right?.

$^{281}$5, 1 to 18..
And then F, he builds the temple..
Chapter 6..
And all those details and architectural details talking about the first temple, right?.
How long it is, how wide it is, how high it is, and all the gold and silver and all that.
stuff..
And then you come to the middle part of the sandwich..
Part G is that Solomon builds his palaces..
His palace complex..
We can call it rifle buildings because the palace is bigger than the temple..
And then the biblical author, after talking about Solomon building rifle buildings, he.
comes back to talk about the temple, right?.
Do you see that?.
And then he talks about dedicating the temple, which is a pair with part E, E prime and E..
And then D prime, D with apostrophe, Solomon uses his gift for himself now, right?.
And then part C, C prime, is his failure..
And then lastly, B and A, talking about Yahweh, is going to raise up enemies to threaten his.
security..
While at the beginning, God gives him peace, if he turns away from God, God is going to.
raise up enemies and troubles for him..
And then finally, in chapter 11, another prophet comes in and determines who's going to succeed.
King Solomon..
And the kingdom split into north and south..
So as you can see, the Bible is trying to design in such a way to tell you what is going.
on here..
And it has to do with the theme of worship..
First you notice chapter 6, chapter 6, part F here, it details how Solomon built the temple..
Now logically, the story should go on to 713, because it talks about how he builds the temple,.
and then should continue to finish the account of building the temple..
However, the author interrupts the flow of the story, and he's trying to tell you Solomon.
built his own palace, the rival buildings..
Now let's take a look at the turning point of the entire story, which is part G right.
there, which is the end of chapter 6 and the beginning of chapter 7..
Do you find it fun?.
I find it fun reading the Bible, right?.
It's not just giving you a straightforward account..
It gives you spiritual meaning..
When he talks about the story of Solomon, it's not just giving you, okay, it's not giving.
you a Wikipedia..
It's giving you a theological meaning that you can think about..

$^{321}$It's like a warning, or it's like a spiritual reminder to all of us..
What does it mean when we worship God?.
Just like King Solomon, he was supposed to worship God by building the temple, but at.
some turning point, he changes..
He changes..
Now let's look at the turning point, which is the end of chapter 6 and the beginning.
of chapter 7..
Well, by the way, this is this translation, all right?.
So as you can see, the turning point is Solomon is building his own palace complex for himself..
Now here's the turning point..
Pay attention..
Chapter 6, 37..
The foundation of the Lord's temple was laid in mid-spring in the month of Seth..
During the fourth year of Solomon's reign, 38, the entire building was completed in every.
detail by mid-autumn in the month of Bul, during the 11th year of his reign..
So it took seven years for Solomon to build a temple for God..
Chapter 7, verse 1..
But Solomon also built a palace for himself..
And it took him how long to build?.
It took him 13 years to complete the construction of his own palace..
Right?.
Do you see what's going on here?.
He used seven years to build the temple for Yahweh, for God..
And then now he built his own palace for his own, how many years did it take him to build?.
13 years, right?.
Almost twice as long..
And here, the translation here, I want to emphasize one thing..
Is that for himself, did you see that?.
For his own..
For his own enjoyment and for his own kingdom..
His own security..
So the palaces of Solomon is being built actually bigger than the temple..
If you check the rest of chapter 7, it talks about five different palace buildings Solomon.
had built for himself..
And they have different names..
And you check the dimensions, the blueprint is actually bigger, much bigger than the temple.
of God..
I think the message is pretty clear..
That this is a matter of priority..
When Solomon built a temple worshiping God, right?.

$^{361}$He was doing okay..
But in the midst of it, all of a sudden, there was a turning point in his priority..
He put his own kingdom above God's kingdom..
He puts the people's worship for him more important than the people's worship for God..
Right?.
I think the lesson is clear..
In the midst of worshiping God, building the temple, Solomon somehow shifted his focus.
from God to himself..
There's competition of priority in our life, right?.
There's always competition of priority in our life..
Always..
In all life stages, right?.
When I was younger, when I was in my mid-20s, late-20s, I was building up my own career..
Right?.
I mean, you can treat ministry as a career, right?.
Because I want to stand firm, right?.
I want to establish my ministry..
And there's always competition of priority, right?.
Sometimes I would devote too much energy in my own career, in my own ministry, more so.
than my personal relationship with God..
Right?.
Now, at times, you know, I forgot to read the Bible..
I forgot to meditate on the Word of God..
I forgot to cultivate my relationship with God..
And then, after I got married, you know, there were competition of priority, right?.
Because my wife would be the competition of my priority, right?.
She wants me to take care of her, right?.
And there's so many things to build up the family, right?.
And then, not to mention, when the kids come, right?.
Oh boy, they are the competitors of priority in my life, right?.
I have to take care of them almost 24/7, right?.
And there's always priority that is messing us up..
And I think there was a Reformation theologian about 500 years ago..
His name is John Calvin..
I think John Calvin said something in his big treaty on systematic theology..
He says, "Humans, we are all manufacturers of idols.".
We manufacture idols, you know?.
When the Bible talks about idol, Baal, you know, Azshara, all those false gods, right?.
We always think that someone invented them, right?.
Someone put them in front of the people to deceive them and to distract them..

$^{401}$But John Calvin says, "We are manufacturers of idols.".
If there's nothing that would distract us from the worship of God, we will invent something..
We'll manufacture something..
We'll make up something that will compete the priority of God, the place of God in our.
life, right?.
That's what idol means..
Idol means someone or something that competes the first place of God in our life..
And here, I think we should really think about what it is that is messing up our life..
Our priority..
And what is your priority?.
Do we really believe in God and believe that God is the best and God is worthy of my everything.
because God loves me so much?.
Because God blesses me abundantly, right?.
God provides for me..
But do we really believe that God is good?.
Or we manufacture other idols that will satisfy our own needs?.
And perhaps if we ask the question, we really don't believe that God will provide..
We always use our time of worship for idols, right?.
We always use our time to do something else because we really don't believe in the bottom.
of our heart..
We don't believe that God is really helping me and providing for me so that I can put.
my focus on Him, right?.
Do we really believe that God is the one?.
God is the true God..
And He is the true God..
And He loves me so much and He provides for me..
And that means He is worthy of my entire being..
My worship..
That's what happened to King Solomon..
The turning point of his life is not simply just, he just, all of a sudden he wakes up.
a morning, one morning and he said, "Okay, I'm going to sin against God.".
No..
He was building the temple, but at the same time there was a matter of priority messed.
up and he started to build his own palaces, build his own sense of security..
It's the matter of my priority..
Whether God takes the priority or my own things..
Whether those things are good or not..
Now talking about idol, now idols are not necessarily evil in their appearance..
They can be good things in our life that will compete our priority for God..
So as we end this message, I would like us to really reflect and think about what is.

$^{441}$your priority?.
Seriously, right?.
It's not what you say, it's how you live..
Right?.
Because your behavior, your action, speak louder than your words..
Right?.
We can sing songs and praise God and say that God is the true God and then go and praise.
Him and worship Him..
However, we have to check our own actions because our actions always speak louder than.
our words..
Priority, it's God's priority or your own kingdom..
Now it can be a good thing if you're going to school, if you're building up your career..
These are good things..
You want to make enough income to accumulate wealth and your property to make sure that.
you're secure..
Those are good things..
You can be raising your own family..
Those are really good things, beautiful things..
It can be a hobby..
Right?.
It can be a hobby, it can be anything in your life and don't get it wrong, those things.
are beautiful and good things and we are being called to manage the resources God has provided.
for us..
However, if we are not careful, the idols will take over and will compete our priority..
And all the blessings that God has given you and me will turn into tragedy..
That's what happened to King Solomon..
What a lesson to learn in terms of worship..
Right?.
We're supposed to worship God, put God as number one..
However, in the midst of it, he builds his own idols and palace complex for himself..
Now as we end, let's continue to just bow before God..
To understand, you know, God is the Lord of our lives..
Right?.
There's one famous saying that says, "If Christ is not the Lord of all, he is not the.
Lord at all.".
In what part of your life is not giving up to God and God is not taking charge?.
Right?.
Let's entrust our lives into our loving God..
Right?.
Let's all pray..

$^{481}$Father, as we reflect, there are so many idols in our life that will distract us from following.
you, from worshiping you..
It may be our career, it may be our future, it may be even our family..
We try to accumulate wealth and goods and gold and silver so that we can feel secure..
But in the midst of it, we forgot that you are our security..
You deserve our full attention..
So Father, I pray that you will forgive us and cleanse us because we are all sinners,.
we are all broken..
We are all manufacturers of idols and those idols will compete our priority for you..
So as we end this worship, let's pray that in this simple way, they will come before.
you..
We recognize you are Lordship, recognize you are the King over all of us, all aspects.
in our lives..
And Father, we come before you to entrust our own life, our whole life into your hand..
And thank you for your Son, Jesus Christ, who died on the cross so that we can have.
new life..
We can have eternal life, that you can dwell in our life..
And you deserve our worship..
I pray that you will humble all of us so we can give the worship of worth to your name..
And how we live and how we work and how we manage our life will tell other people that.
God, you are great..
God, you are the Lord..
So may you continue to just cultivate our relationship with you, knowing that you are.
not our idol, but you are the Lord of our lives..
You are the Lord of all..
If there's anything that we struggle of letting go, to have you take control, ask that your.
Spirit will humble us, knowing that we owe everything to you..
And thank you for being our loving God, God of blessing, God of provision..
Continue to help us to increase our faith in you, to follow you, to worship you fully..
In Christ's name we pray..
God bless you and I hope that this message will remind you and to help you reflect on.
your relationship with God..
For the rest of the week, God be with you and you go to your work and bear witness for.
the great name of God..
And then come back next Sunday, we're going to serve God together..
Remember next Sunday is the MSI outing, right?.
So we're going to begin at 10 o'clock..
And there's a combined service and then we'll go to MSI and have some fun..
Alright?.
God bless..

\newpage



\section{}
\label{sec:NNhx4LQjZTc}
\textbf{比對研究 1: 必要, 發展, 方法   -  陳崇基牧師}
\newline
\newline
連結: \href{https://youtube.com/watch?v=NNhx4LQjZTc}{\texttt{https://youtube.com/watch?v=NNhx4LQjZTc}} ~~~~ 語音日期: 2024-05-22
\newline
\newline
\hyperref[sec:D_DiJk8_Dhc]{\small{< < < PREV SERMON < < <}}
~
\hyperref[sec:index_chronic]{\small{[返順時目]}}
~
\hyperref[sec:index_scriptual]{\small{[返順卷目]}}
~
\hyperref[sec:NOP7HqFawO8]{\small{> > > NEXT SERMON > > >}}
\newline
\newline
$^{1}$以往我們一直在說創世記一到六章頭.
都是在比對近東的文獻.
我經常都struggle(掙扎).
我掙扎的地方就是.
去到洪水的時候就會引起很多很敏感的問題.
因為等一下我會講到洪水的記載故事.
其實是非常相似近東的洪水故事.
特別是Gilgamesh epic 很相似的記載.
所以會引發我們很多不同的爭議.
所以我覺得有很多東西需要交代.
之前交代過聖經考古學的方法是怎樣.
洪水的毀滅層等等.
洪水的掩蓋當時所知的世界.
還是我們現在的全球.
我都有一些思考的過程.
有些方法要去掌握.
這裡我都想用一些時間.
暫時去看看比對研究的必要.
發展和方法.
因為這些東西不掌握好.
我們聽了洪水的故事.
和近東的洪水故事去比較的時候.
很快就會引發出很多不同的問題.
正因為這個緣故.
所以我覺得不如我們先將.
我們這個比對研究的方法掌握了.
特別是在過去百多年.
對這個比對研究的發展.
到底有什麼是錯了的地方.
現在知道怎樣去修正.
怎樣用這個方法去正確.
所以在這幾節我們嘗試用一些時間.
來看看這個comparative studies.
或者comparative approach.
比對研究到底是什麼呢.
首先第一個問題問你的.
如果你是18世紀之前.
即是1800年之前的人.
當時我們在歐洲.
我們都有些歐洲的畫.

$^{41}$如果我們當時在1800年之前.
我們會對聖經的背景有什麼認識呢.
你覺得如果我們回到歐洲.
我們做歐洲人.
18世紀之前.
我們會對聖經的認識.
或者對聖經背景的認識會是怎樣呢.
所以你想一想.
當時的人對聖經背景的認識是怎樣呢.
你看一看那些畫就知道.
這幅畫呢.
這幅畫是1563年所畫的.
18世紀之前.
你看到他在畫什麼.
在畫什麼呢.
在畫哪個聖經故事呢.
這幅畫叫做Tower of Babel.
就是巴別塔這個故事.
你看到他畫這個巴別塔的時候.
是圓形的.
然後好像羅馬的Colosseum.
或者歐洲式的建築.
你看到這些人在下面都是在作石等等.
是歐洲人來的.
你看到這個巴別塔的故事.
他畫成為歐洲色彩的故事.
18世紀之前的人.
對近東聖經的背景幾乎是沒有的.
所以引發了當時的人對聖經的理解.
特別是對巴別塔的故事.
他們就會用了當時歐洲色彩的東西.
來想像出來.
我們為什麼需要有比對研究.
為什麼需要明白聖經的背景呢.
就是因為我們沒有一個人看聖經的時候.
是一張白紙.
我們總會帶著某些文化的認知.
或者一些預設.
我們會去看聖經.
所以你記得我們說到.

$^{81}$在釋經的一些視野的時候.
有一個視野就是.
The world in front of the Bible.
就是我們有時候不小心的時候.
我們會將我們的印象.
我們的假設.
就會讀進去聖經.
所以為什麼需要有一些所謂working knowledge.
對近東的認知.
對原文 希臘文 希伯來文等等.
有一些基本的認知.
是幫助我們去看聖經和讀經的.
所以你看到歐洲在18世紀之前.
幾乎是沒有可能認識到聖經的背景.
所以他們就用了歐洲的背景.
預設讀進去.
其實不單止是這樣的.
不單止是18世紀前.
很多時候我們不想清楚的時候.
都會是這樣.
譬如有人說耶穌為什麼是白人呢.
你看到這裡是Campus Crusade.
就是這個學員傳道會.
1970年代的時候拍攝耶穌傳.
他找耶穌是一個藍眼睛.
白人的角色.
這個左邊.
如果你去Google或一些搜索引擎.
你去找找.
找找Jesus.
很多時候你會找到類似這樣的形象.
耶穌就是這樣的樣子.
什麼樣子呢 就是白人.
然後好像一個歐洲人.
所以很多時候我們不小心的時候.
我們印象中就會當成.
原來耶穌就是一個好像白人的樣子.
或者在不同的地方.
會當成基督教是一個洋教.
或者一個西方的宗教.

$^{121}$這些耶穌的形象.
都是因為西方的基督教.
他們有的文化的影響比較深遠.
所以就造成了一種這樣的文化傳統.
我們對耶穌的印象是怎樣的呢.
似乎想來想去.
很多時候都是跟隨了歐洲的基督教的影響.
所以就想起我們用自己的文化背景.
或者歐洲人用他們的文化背景.
塑造出來 想像出來的.
那個耶穌就像一個歐洲人.
在2002年的時候.
有一個Forensic的Doctor.
就是一些對人骨有所研究的.
然後藉著人骨可以把人畫出來.
有很多時候他們是做一些查案的繪畫.
這些重構圖.
這種醫生就去研究.
在巴勒斯坦地 以色列.
找到第一世紀的一些人骨.
那些頭顱骨.
找回一些猶太人的頭顱骨.
找回一些平均的猶太人的頭顱骨.
重組當時一個住在巴勒斯坦地的一個男人.
典型的男人.
到底他的樣貌應該是怎樣的呢.
他就找了三個猶太人的人骨.
來做一些重新的繪畫.
繪畫出來就變成這樣.
是一個深啡色眼睛.
皮膚比較黑 比較啡色.
一個這樣的樣貌.
所以有些人就說.
其實耶穌應該是怎樣的呢.
耶穌本身的樣貌應該是怎樣的呢.
因為耶穌都是百分之百人.
是一個猶太人住在巴勒斯坦地.
所以他的樣貌應該大概是這樣的.
等於我們今天說的亞裔人的樣貌.
大概的輪廓是怎樣的.

$^{161}$或者一個白人 一個黑人的輪廓有什麼分別呢.
所以這樣去看的時候.
就發現原來耶穌可能像這樣.
比這個白人的 歐洲人的耶穌還要多.
所以比對研究想做些什麼呢.
比對研究就帶我們回到聖經古代的背景.
回到巴別塔的故事.
你看到巴別塔是歐洲所構想出來的.
一個圓形的巴別塔.
就跟近東世界 最近這百多二百年.
去到找到的所謂巴別塔.
就是完全兩回事.
你看到下面的左下角.
這個巴別塔其實就是吾爾阿巴拉罕的家鄉.
當時去發掘的時候.
那些上面 你看到很多工人站在上面.
拍了一張照片.
重構出來 把磚頭疊上去.
或者把倒塌的東西推上去.
其實就是這樣的一個巴別塔.
其實就是一個正正方方.
然後有一個大的樓梯上去上面.
這些其實就是在中東的世界.
古代中東的世界.
就有很多這些所謂的廟塔.
吾爾在哪裡呢.
吾爾在南維索不達米亞的這個地方.
就是在波斯灣上面.
今天的伊拉克的南邊.
你看到這裡有兩條大河.
一條就是Euphrates.
有法拉底河.
然後上面就是底格拉斯河.
就是Tigris.
吾爾就在這裡.
如果你看這些衛星的圖片的時候.
其實你都看到有這個所謂的巴別塔.
這個廟塔在這裡.
很清楚的讓你看到.
就是剛才我們看到的那個廟塔.

$^{201}$我們就在那張照片.
就是這樣拍過去的.
當然現在是平面.
就是四四方方.
全部用泥磚疊出來的.
實心的.
其實做些什麼呢.
是建一條樓梯上去.
一個所謂英文是portal.
一個梯階上去上面.
可以很高的.
你剛才看到的這些照片.
你知道有幾層樓這麼高.
上面就是一個神靈降臨.
去到居住的地方.
所以很多巴別塔.
所謂的廟塔.
上面其實不是人上去的地方.
是一個聖所的地方.
就是神下來.
人建一個這樣的廟塔.
方便神靈下來這裡.
下面很多時候.
在城市裡面.
在廟塔旁邊.
就會有一個聖廟.
一個廟宇.
他下來的時候.
就可以去接受祭禮等等.
他可以吃東西等等.
餵飽他.
初一十五等等.
就是這樣的一種宗教建築.
所以建這個巴別塔.
所謂建這個廟塔.
其實就是當時近東的文化裡面.
是一個很普遍的東西.
你看看這些泥印像.
其實很小的.
很小的圓形的.

$^{241}$上面雕了一些圖案.
或者一些字.
然後我們找到.
考古學家找到之後.
就將它印在一些軟的泥上面.
然後就知道那個圖像.
到底想說什麼.
你看到這裡有一個三角的.
三層的一個廟塔.
就說到人或者軍王.
去建立這些廟塔.
所以在近東的背景裡面.
我們知道創世記第十一章.
講巴別塔的故事.
就正是描述當時的文化背景.
如果有這些廟塔.
讓我們看到的時候.
讓我們知道它的功能是什麼.
我們才知道怎樣去解釋.
創世記第十一章.
如果不是這樣的話.
我們就會當作是歐洲的款式.
你看到有些人爬上去等等.
其實近東的廟塔.
是給神靈下來的.
方便它下來.
所以你看回創世記第十一章的時候.
都有一種這樣的動作.
一種這樣的動態.
有人下來.
或者有人建上去.
要高達於天等等.
另外一個石碑.
就是左邊這個.
你看到有一個軍王的樣子.
拿著一支黃杖.
前面是什麼.
前面又是一個廟塔.
畫出來的時候.
就是一個正方的廟塔.

$^{281}$就好像我們平時在吾爾.
或者南米索不達米亞.
找到的這些廟塔一樣.
你看到這個是平面圖.
上面的圖案.
好像三角一樣.
上面有一個廟所.
一個聖所.
讓神靈下來.
方便人們接受祭禮等等.
供奉.
所以有了這些背景的時候.
我們就知道如何去解釋清楚一點.
創世記第十一章.
我們要懂得原文.
要懂得經文.
當然是重要的.
要懂得翻譯裡面的細節.
都是重要的.
但是這些東西是翻譯不了的.
你有沒有留意.
你怎樣翻譯都好.
你都翻不出背景.
對不對.
譬如我今天跟你說.
有一個字叫做911.
911是怎樣翻譯的.
英文說911.
或者911.
我們翻譯成911就行了.
但是你翻譯不了.
為什麼.
因為你今天聽911的時候.
你沒有聽背景.
你是不知道911是一回事.
所以有歷史背景.
有政治背景 有軍事背景.
等等很多文化的預設.
是在911這個字背後的.
所以你怎樣翻譯都好.

$^{321}$你都不能夠將文化背景翻譯出來.
所以為什麼我們需要有其他的工具來幫助我們.
原文工具當然是必要的.
所以比對研究.
文化背景都是需要的.
我們先做一些字彙的澄清.
很多時候現在人們說到.
comparative studies.
或者比對研究的時候.
有不同的字眼.
可以是contextual study.
即是背景的研究.
可以是contrastive studies.
即是對立的研究.
因為將近東的文化.
和聖經所講的東西.
來做一個對立.
有什麼不同的地方.
這個就強調它不同.
comparative study.
現在大陸主要都會叫做對立.
主要都會叫做比對研究.
因為不只是compare.
都會contrast.
亦都會有時候叫做background studies.
即是背景的研究.
和contextual study很相似.
但怎樣都好.
都是在講聖經本身有的背景.
就等於我所講.
你講到聖經裡面的一些字眼.
例如剛才講到的這個例子.
創世記第十一章.
講到建那個塔.
那個塔你可以翻譯出來.
英文 中文 其他語文.
翻譯到原本的希臘.
希伯來文.
翻譯出來.
但你翻譯不到的就是.

$^{361}$到底那些廟塔是做什麼的.
如果你明白了背景的時候.
就容易一點掌握.
在這裡我們又看看.
到底聖經的背景是什麼.
在這裡誠實地講.
很簡單的嘗試將聖經的背景.
去列舉出來.
首先從經文開始.
經文就有上下文.
我們就叫做literary unit.
上下文的段落.
譬如你解釋某一節的經文.
那段落就自然是.
一個比較大一點的上下文.
譬如創世記第十一章.
所謂巴別塔的故事.
最近的段落.
就是十一章一至九節.
整個故事就是.
廟塔的故事.
巴別塔的故事.
然後你要將這個unit.
放在一個更加大的上下文.
那個就可以是一個書卷.
如果是創世記十一章的話.
你就要將它放在.
一個不單止大一點的段落.
一至十一節.
然後你就可以.
將整個創世記的故事.
放在整個書卷.
整個創世記裡面.
你就可以看到.
巴別塔的故事.
在整個創世記的流程.
的脈絡裡面.
到底佔了什麼地位呢.
為什麼要有這個故事呢.
然後我們可以.

$^{401}$去到一個所謂文集.
就是Corpus.
Corpus就是.
這一類的.
附近的書卷.
加在一起的時候.
到底怎樣去到呢.
告訴我們這一段的經文.
是怎樣有什麼功能.
將上下文擴大了.
不是單單一本書.
而是可能幾本書.
如果你看創世記十一章.
自然就是五經.
你就看回五經.
到底由創世記第一章.
到申命記最後那一章.
到底這一段的經文.
有什麼功能呢.
這些就是文集.
比如說新約的一個例子.
就是說到保羅的書信.
比如說你在查.
羅馬書第九章.
第九章.
自然的段落.
就是在看九至十一章.
因為九至十一章很明顯就是說到.
神對以色列的計劃.
會怎樣呢.
然後整個書卷.
就是羅馬書.
再看看保羅的神學.
再擴大一點.
就看到保羅的文集.
看看保羅的書信.
到底怎樣去看.
猶太人這個種族.
在末世的時候.
上帝對他們有什麼計劃呢.

$^{441}$然後再擴大一點.
就是舊約和新約.
這個很明顯了.
新約的自然是整本新約傳書.
舊約的就是整本舊約.
Torah.
然後我們就看整本的正典.
Canonical的Context.
比如說創世記一至十一章.
在整本的正典裡面的發展.
到底有什麼回應呢.
有什麼呼應呢.
等等.
然後我們就可以從聖經裡面的世界.
剛才我們綠色的.
都是在聖經裡面.
在聖經裡面.
我們需要考慮的.
但是也需要考慮什麼呢.
就是聖經背後的世界.
聖經背後的世界.
就有什麼呢.
Extra canonical 經外的文獻.
所以我們有時候會看看.
比如我們之前看過.
創世記第六章.
說到神的遺址.
有沒有一些君王會說自己是神的遺址呢.
我們看過Gilgamesh epic.
這些就是經外的文獻.
然後有些歷史的背景.
文化背景.
甚至地理的背景.
所以為什麼經常打開Google Earth.
我們沒有地理的背景.
我們不知道原來.
聖經在創世記第一至十一章.
那個背景.
其實就是.
米索不達米亞.

$^{481}$你記得阿巴拉罕.
就是在米索不達米亞.
在五二出來.
所以很有道理來說.
一至十一章很多時候的背景.
地理的背景甚至文化的背景.
都是和Mesopotamia.
早期的人類的文明.
他們的寫作.
他們的思維有關.
最後有考古的背景.
當然這些不是一層一層的.
規劃得很嚴謹.
但是藍色的主要是.
world behind the bible.
所以你記得.
我們在聖經的視野.
就需要有這幾樣東西.
所以為什麼今天第一節.
是有這個必要的.
因為我們不是一張白紙.
我們總會帶著我們的預言.
當我們明白多一點.
聖經背後的世界的時候.
我們會深入一點了解.
到底聖經裡面在做什麼.
說什麼.
就不是純粹翻譯了.
只是原文.
正如我所說.
原文有很多東西是翻譯不到的.
什麼叫巴別塔.
你可以翻譯成塔字.
或者成流這個字.
但是翻譯不到就是那個.
就是廟塔的.
在近東裡面的功能.
如果你喜歡我們的youtube channel.
就記得like comment和subscribe.
\newpage



\section{民數記 6:24-26}
\label{sec:NOP7HqFawO8}
\textbf{神賜的福   民數記 6\_24-26  (陳崇基牧師)}
\newline
\newline
連結: \href{https://youtube.com/watch?v=NOP7HqFawO8}{\texttt{https://youtube.com/watch?v=NOP7HqFawO8}} ~~~~ 語音日期: 2024-05-22
\newline
\newline
\hyperref[sec:NNhx4LQjZTc]{\small{< < < PREV SERMON < < <}}
~
\hyperref[sec:index_chronic]{\small{[返順時目]}}
~
\hyperref[sec:index_scriptual]{\small{[返順卷目]}}
~
\hyperref[sec:NP_1P0Hd_Zc]{\small{> > > NEXT SERMON > > >}}
\newline
\newline
民數記 6:24-26
\newline
\begin{longtable}{cl}
\hline
\hline
章節 & 經文 (和合本修訂版)\\
\hline
6:24 & \begin{tabularx}{0.7\textwidth}{X} 『願耶和華賜福給你，保護你。 \end{tabularx} \\ \\ \relax
6:25 & \begin{tabularx}{0.7\textwidth}{X} 願耶和華使他的臉光照你，賜恩給你。 \end{tabularx} \\ \\ \relax
6:26 & \begin{tabularx}{0.7\textwidth}{X} 願耶和華向你仰臉，賜你平安。』 \end{tabularx} \\ \\
[1ex]
\hline
\hline
\end{longtable}
$^{1}$各位好,首先跟大家拜年.
我身處芝加哥.
我知道有很多弟兄姊妹在我附近.
也知道有不同的弟兄姊妹在不同的城市.
我們開始的時候先跟大家拜個年.
在你的Live chat box(直播箱)寫新年快樂.
還有一件事,希望你告訴我你身處的城市.
到底現在有多少度.
因為我們芝加哥是非常非常冷的.
我也不知道現在是多少度.
我現在應該是….
別人告訴我是負18度.
我們算上風速,應該是負18度.
不知道你的城市和你自己的地方.
你的溫度是多少.
請你在直播箱打出來.
以致我們能夠看到.
能夠跟大家說聲Hi.
大家新年滿福.
今天我們來分享一段聖經.
是一段很寶貴的聖經.
這段聖經是在民數記第六章.
我們先去PowerPoint.
民數記第六章.
願耶和華賜福給你,保護你.
願耶和華使他的念光照你.
賜恩給你,願耶和華向你養念,賜你平安.
這一段聖經正如我所說.
是在民數記出來的.
這一段聖經是大家非常喜歡的一段聖經.
我相信如果你問一些人.
他們家裡的日曆或者他們家裡的揮春等等.
可能有一些經文的畫.
或者一些經文的紀念.
都會用民數記第六章這段聖經.
那麼,先看看.
有些弟兄姊妹在Edmonton.
-24度.
所以我們都不差.
Calgary是-25度.

$^{41}$看看還有誰.
Halifax是0度.
我們芝加哥其實都不差.
有-18或者-20度.
溫哥華只有0度.
其實對我來說都算暖的.
那麼剛才我們去看這段聖經.
願耶和華賜福給你.
這一段聖經.
這一段是一段祭司.
阿倫祭司對以色列民的祝福的句語.
那麼很多時候在崇拜的最後.
不同的教會都會用這一段來作結束.
會眾離開之前就得到上帝的祝福.
那麼新年來我們當然要講這些祝福的說話.
這個不能走的了.
那麼在這一年裡面.
因為疫情的緣故.
這一段的聖經就特別出名了.
就特別普遍被人使用了.
為什麼呢?.
因為這一段的聖經.
在網路上.
在過去的疫情當中.
有很多人就把它重新演繹.
有一首歌是出自這一段的聖經的.
就叫做The Blessing.
不僅這首歌是一首很有意思的詩歌.
是基於這一段的聖經.
也翻譯了不同的語言.
譬如有希伯來文,日文,韓文,阿拉伯文.
或者其他的演繹方法.
是一首很有意思的詩歌.
也將民數記第六章裡面所提到的.
我們的祈禱.
我們的希望.
是能夠在這個疫情當中能夠臨到我們.
其實這一段的聖經.
不僅是過去一年.
因為這一首歌.

$^{81}$因為疫情的緣故.
是令我們非常渴望.
得到耶和華的賜福.
得到神向我們去讓念.
賜我們平安.
這一段的聖經其實有史以來.
都是一段非常古老的聖經.
我想跟你說的是.
在1970年代的時候.
在耶路撒冷有一個地方叫做Ketev-Hinom.
這個地方是在耶路撒冷的西南面.
我帶大家去看看.
現在我在芝加哥.
我們就去耶路撒冷.
在耶路撒冷的西南面.
有一個地方叫做Ketev-Hinom.
這裡就是耶路撒冷.
你看到這個金頂清真寺.
這裡是耶路撒冷舊城.
在耶路撒冷舊城的西邊.
東邊和南邊都是一個天然的山谷.
東邊我們知道是Ketev-Hinom.
就是琴輪谷.
西邊和南邊.
這個好像L的谷.
大家知不知道是什麼.
叫做Hinom谷.
就是恩聯子谷.
在舊約裡面時常會看到的.
在這裡耶路撒冷舊城的西南邊.
有一個山頭.
這個山頭叫做Ketev-Hinom.
Ketev的意思就是膊頭.
就像在谷上有一個人的膊頭.
上到一個小山的上面.
上面有一間現代的教堂.
叫做安德烈教堂.
這間教堂叫做St. Andrew's Church.
在一次世界大戰之後興建的.
紀念英國在一次世界大戰中陣亡的士兵.

$^{121}$是一間英國蘇格蘭教會的教堂.
在這個教堂附近.
1970年代的時候.
人們發現有一個古墓.
或者有一些古墓在這裡.
有不同的山洞在那裡.
他們就找到一些舊約時代的古墓.
當時那些古墓是從山洞挖出來的.
好像山洞一樣.
是墓穴.
當時1979年的時候.
人們就找到這些古墓.
這些古墓按照一些連算的方法去計算.
大概在主前600年或500年左右的時間.
如果你記得舊約的歷史.
這100年左右是一個很動盪的時候.
耶路撒冷第一聖典.
最後在586年的時候.
586 BCE光元前.
是被毀的.
在這裡就找到一堆山洞的墓穴.
裡面這些墓穴的格式.
在舊約的時候.
通常都是這樣的.
你走進一個山洞裡面.
然後就有三邊.
有一些石椅子或者石桌子等等.
就將一些屍身放在上面.
我知道新年樓樓說這些好像不太吉利.
不過不要緊.
是跟聖經有關的.
在這裡就找到.
這些舊約的鐵鋼時代墓穴的格式.
在山洞裡面出來.
然後三邊有石床.
放在這裡一年之後.
將屍骨放在石床下面.
有一個儲骨的地方.
就是一個repository的地方.
記得舊約裡面經常有一段聖經.

$^{161}$有一段說話.
歸回列祖.
歸到你的祖先那裡.
這些說話就是來自這樣的情況.
就是說你離開這個世界的時候.
你不是真的離開.
你是歸回你的列祖.
回到你的列祖.
這個其實在舊約的神學裡面.
是一個祝福.
在這些山洞墓穴那裡.
也找到一些陪葬品.
你看到這個以色列博物館.
將這些墓穴重構它的模型出來.
讓旅客或者人們能夠看到.
以前的格式是怎樣的.
在這裡有一個重構圖.
走進去裡面.
有不同的number.
考古的人員將這些墓穴.
這些山洞的室.
將它number了.
加上了一個數字.
就是25這裡.
25這裡進去.
正如我所說.
就是舊約的墓穴的格式.
找到一些物件在那裡.
其中一個物件找到的就是這一塊東西.
左上角的.
你有沒有看到是一個捲起來的文字.
其實本身是一些古的銀.
就是真銀.
只不過你看不到銀色.
因為它氧化了.
花了很多時間將這些銀卷.
打開了它.
其實很小的.
最多是一吋那麼大.
所以將它打開的時候.

$^{201}$在顯微鏡下.
你會看到寫了一些文字.
有些筆劃.
是一些古希伯來文.
你看到這裡一個章.
然後一劃.
然後再加一劃.
等等.
然後有其他的筆劃在樹下.
在兩個古銀卷裡.
一個大一點.
一個小一點.
然後找到一些以前的希伯來文的文字.
將它打開的時候.
現在以色列博物館.
你會看到的就是這些Scriptural Roll.
一個A一個B.
兩個這樣.
然後發現這些銀卷上面.
用一些很幼的工具.
來刮上一些經文.
那些經文我們再翻譯出來的時候.
大家就看到了.
這裡是第一個卷.
是一個比較大的卷.
大概有18行的希伯來文.
英文就放在中間這裡.
然後右邊這一處.
我們就看到有中文的粗略翻譯.
小弟就將它翻譯出來.
這裡說什麼呢?.
他說耶和華.
Yahweh.
是偉大的.
他保守.
他的聖約.
和恩代.
那些愛他和守他誡命.
這裡中間斷了.
猜不到是什麼.

$^{241}$但是大家都猜到了.
其實是舊約的一些經文.
The Eternal.
第八行.
永恆者.
祝福比任何的陷阱.
網羅和邪惡更加大.
因救恩在他裡頭.
第十二行.
因耶和華.
For Yahweh.
是我們的恢復者.
是醫治者.
和我們的盆石.
願耶和華.
祝福你.
May Yahweh.
May the Lord.
Bless you.
And may he keep you.
第十六.
願他保護你.
May Yahweh make his face shine.
願耶和華使他的念光照.
後面就斷了.
這是第一個經文.
你有沒有看到有文藻記第六章的文字.
看到了嗎.
看到了.
這是第二個經文.
就是給一個Personal Name PN.
給一個姓名.
上面當然我們不太清楚.
猜不到是什麼姓名.
應該就是一些以色人的姓名.
給他的什麼呢.
The son or daughter of.
名字我們猜不到.
只有一個H.
然後May he or she.

$^{281}$不知道男還是女.
Be blessed by Yahweh.
願他得到耶和華的祝福.
耶和華是誰呢.
就是那位戰士.
幫我們打仗的那位.
也是我們的幫助者.
或者可以這麼翻譯.
Help her.
還有那位.
斥責邪惡者.
The rebuke of evil.
願耶和華祝福你.
保護你.
願耶和華使他的炎光照你.
又賜你平安.
看到剛才有幾句話.
都是直接出自文藻記第六章.
所以這一段聖經.
剛才我說是屬於什麼時候.
是公元前六百年左右.
當時第一聖殿快要被毀.
或者已經被毀了.
這一段時間.
所以這個經卷.
其實就是我們現在所知道的.
其中一個最古老的.
聖經的經文.
就是在經外找到的.
我們現在的希伯文文.
是比較後期的抄本.
但是你見到之前.
我們找到這些.
是公元前五百年的時間.
人們已經有了這些經卷.
並且抄寫在銀卷上.
即是在銀片上.
然後捲起來.
可能是戴在頸鏈等等.
亦都是跟著那個人.

$^{321}$去到最後墓穴的地方.
你就可以見到這一段的聖經.
是多麼受歡迎.
聖經你要知道.
其實是一些很美的詩歌.
我們中文翻譯.
願耶和華賜福給你.
保護你.
願耶和華使他的念光照你.
賜恩給你.
願耶和華向你讓念.
賜你平安.
其實詩歌是很美的.
在這裡列了一些希伯文文.
你自己數一下.
其實有幾個字母.
這段的祝福語句.
詩歌是有三行的.
一二三你見到嗎.
中文其實都是一二三.
三行.
每行是越來越長的.
第一行的字母是十五個字母.
然後第二行是二十.
五個字母.
加上去.
字詞.
幾多個字眼.
第一行有三個.
第二行五個字.
第三行是七個字.
音節.
第一行.
If I regard, I deny.
If I share, I make up.
第一行.
十二個字母.
第二行是十四個字母.
第三行是十六個字母.
其實就像七言絕詩.

$^{361}$好像我們聽一些詩歌.
剛才我們聽一些很動聽的詩歌.
Jackson的team帶我們來.
剛才你不要低估.
背後打鼓的是Jackson的兒子.
很小的.
他打得很好的.
所以是一些很美的詩歌.
其實聖經的詩歌詩句都是很美的.
正如這裡民數記大六章.
三行.
很有規律的.
一路來擴大告訴我們.
神對我們的祝福會是怎樣呢.
除了我們看到這些字母字詞字節.
一路加上去的時候.
最後最後那個字是什麼字.
你見到中文是平安.
中文翻譯很好.
因為希伯來文都是平安.
平安不是無災無病.
聖經裡面說平安是上帝給我們的生命.
那種豐盛的生命.
並且你留意那個動詞.
兩個兩個去.
祝福.
然後保守.
May the Lord bless you and keep you.
祂賜福給你並且保護你.
不單是賜福給你.
也告訴你那種賜福是保護的.
第二行.
祂念光照你和賜恩給你.
兩個動詞.
第三行.
養念.
Lift up his countenance.
養念的意思是什麼意思.
上帝的念去到向你養念.
意思就是留意你.

$^{401}$在意你.
向你養念並且賜給你平安.
六個動詞.
每句兩個動詞.
告訴我們.
上帝對他的子民有一個整全的祝福.
但我想問多一個問題.
為什麼這段聖經會放在這一處.
你看下面第七章.
出現了一句話.
摩西立完帳目.
就是會幕.
就將會幕高抹.
將它成聖.
然後器具都預備好了.
我們就能夠有會幕可以運作.
敬拜的地方.
七章就是祝福語之後.
出現了這一句.
立完整完的會幕.
整完會幕這一句話.
什麼時候出現過呢.
就是很前的.
就是出埃及記最後的時候.
出埃及記最後的時候.
他們去到西乃山.
上帝跟他們立約.
並且上帝去吩咐他們建造會幕.
最後出埃及記.
最後一章第四十章.
就是第二年.
就是第一年正月初一的時候.
他們出埃及.
就是那個禦節的時候.
現在他們出埃及就經過曠野.
去到西乃山就是第二年的正月初一日.
會幕就立起來了.
這句話.
摩西立起帳幕.
立起Tabernacle.

$^{441}$建好會幕敬拜的地方.
這個出埃及的最後最後的時候.
出現這句話.
剛才我們看了民數記第七章一至二節.
就是摩西又立完會幕.
立完帳幕.
中間夾住的是什麼呢.
出埃及記完了.
四十章完了.
就是利未記.
利未記整本書.
然後到民數記頭六章.
都是說什麼呢.
說敬拜的律例.
就是說他們去到西乃山.
去告訴他們.
你建會幕我就跟你們同在.
但是你跟我進入一個聖約.
一個關係的律例.
而這個關係是怎麼樣的呢.
我們大家怎樣去維繫這個關係呢.
上帝就跟我們同在.
那我們怎樣呢.
利未記第一章就開始說.
耶和華就叫摩西.
去告訴以色列人.
他們怎樣去敬拜.
利未記裡面有很多的規條.
有很多的律例.
全部都是環繞著這件事.
就是怎樣成為聖潔的子民.
怎樣去敬拜上帝.
一直說到民數記第六章.
都是在說這件事.
人數有多少.
我們呢.
是有耶和華跟我們同在的一班子民.
我們怎樣去敬拜.
就是說.
我們要尊崇上帝.

$^{481}$要按照上帝給我們的規條.
上帝給我們的說話.
來行事為人.
成為君尊的祭祀.
整個中間的利未記一章到民數記第六章.
都是說這件事.
結束的時候.
上帝就跟亞倫說.
亞倫是祭祀.
祝福會眾.
就是我們剛才讀的那段聖經.
民數記第六章24-26.
所以那段聖經不可以這樣抽離.
然後就說.
這段就像是揮春一樣我們貼上去.
其實這段聖經是一個結束.
一個祭祀的祝福語是一個結束.
結束什麼呢.
就是神的子民.
知道他們進入了上帝的聖約的時候.
他們要去敬拜.
要去尊榮上帝.
然後民數記.
結束了這個祝福.
才繼續起行.
如果你繼續看下去.
第七章第十章等等.
就說到第二年二月的時候.
百姓又再起行.
就是說這裡突然間.
時間停止了.
又出埃及記.
又有了緩和.
一直講下去.
就是我們成為了上帝的子民之後.
進入了這個關係.
這個聖約當中.
最後才有祭祀的祝福.
言下之意.
就是說你要明白.

$^{521}$或者體會.
這個所謂的祝福.
就不單單是一段揮春的說話.
也不單單是我們敬拜之後.
牧師一定要講這句話.
我們才有祝福.
言下之意是說什麼呢?.
這段祭祀的祝福的話.
就是說你如果.
要與上帝的祝福有份的話.
你就一定要與上帝的約有份.
如果你不是成為上帝立約的子民的話.
你就得不到這個祝福.
並且你要得到這個祝福的時候.
你就要行在上帝的聖約律法裡面.
利未幾一章至文肅記.
一直講的就是我們作為一個聖潔的子民.
作為代表上帝.
以及有上帝的名字.
上帝同在的一群人.
我們就必須要行在上帝的說話.
行在上帝的律例當中.
然後最後才有上帝的祝福.
在我們新年的時候.
我們當然想聽一些祝福的說話.
不過讓我們知道.
最重要的.
不是因為祝福的說話好聽.
而是最重要的就是我們能夠成為上帝的子民.
真正的去敬拜祂.
去跟隨祂.
去根據祂的說話來生活.
我們才能夠得到上帝的祝福.
願意我們不單單聽一些好聽的說話.
而是我們立志去成為一個敬拜主.
去跟隨主.
去順服祂.
聽祂說話的一群子民.
這是我們新年的願望.
願上帝祝福各位.

$^{561}$好,又到我們牛年快問快答的時間.
我們講完的時候.
大家去評論一下.
到底最重要的重點是什麼呢?.
希望大家能夠知道.
我們有幾條問題問大家.
希望大家在留言的地方.
將答案寫下去.
希望你們可以成為第一個.
第一條問題我想問的是什麼呢?.
準備好了沒有?.
牛年快問快答就是.
獨裁的亞倫的祝福詞.
是出自哪一本書呢?.
A. 創世記.
B. 埃及紀.
C. 文數紀.
是什麼呢?.
沒錯,答案就是文數紀.
Book of Numbers.
第六章24至26節.
所以記住了.
是第六章出來的.
但是立約之後.
有一大堆東西說到.
神的子民應該怎麼生活.
然後隨著祝福而來.
第二條問題準備好了沒有?.
第一條問題.
準備好了.
就是疫情當中有一首紅爆的英文基督教詩歌.
歌名是什麼呢?.
A. The Blessed.
B. The Blessing.
C. Bless Me.
這一段聖經.
基督教的詩歌.
來自文數紀第六章.
那首歌叫什麼呢?.
很簡單.

$^{601}$就是.
Blessing.
那個祝福.
唯一的祝福就是Blessing.
第三條問題.
第三條問題是問什麼呢?.
準備好了嗎?.
第三條問題.
就是文數紀的祝福語裡面.
希伯來的原文三行裡面.
有多少個字母呢?.
A. 10, 20, 30.
B. 15, 20, 25.
C. 50, 15.
答案是什麼呢?.
希伯來的原文有多少個字母呢?.
三行裡面.
就是15, 20, 25.
15, 20, 25.
你自己數一下.
不數音標.
1, 2, 3, 4, 5.
6, 7, 8, 9, 10.
15個.
第一行15, 第二行20.
第三行是25.
現在我們去到最後一條問題.
OK.
最後一條問題.
就是第四條.
就是發現祭司祝福語的古銀卷裡面.
地點是在哪裡呢?.
A. 古銀卷.
B. 古銀卷.
C. 蒙古卷.
石安山.
答案是什麼呢?.
就是古銀卷.
古銀卷.
在一個.

$^{641}$博頭恩鸞谷.
恩鸞山谷.
恩鸞山的.
堅上面.
OK.
我們希望輕鬆地.
跟你說一些祝福的說話.
也希望能夠.
幫助我們.
來到主的面前.
我們有他的恩典成為納入的子民.
讓我們信服他.
讓我們在生命當中得到祝福.
我們先跟心底人祈禱.
主愛多謝你給我們的祝福.
也讓我們知道.
這些不是理所當然的事.
是因為我們進入了.
你的承諾當中.
因為耶穌基督為我們所作.
祂的救恩 祂的寶血洗淨我們.
祂十字架的恩典.
讓我們能夠成為君尊的祭司.
能夠成為神國裡面的子民.
求主你自己親切地幫助我們.
讓我們能夠信服你.
敬拜你 跟隨你.
以致你的祝福能夠臨到.
也讓我們知道.
我們得到的祝福是一種.
屬靈 永恆的祝福.
不是世人所給的.
是跟世人所給的有不同的.
我們感恩 祈禱 奉主名禱.
\newpage



\section{}
\label{sec:NP_1P0Hd_Zc}
\textbf{考古探秘 (發現死海古卷的傳奇故事) - 陳崇基牧師 Mark Sir (2021.05.30)}
\newline
\newline
連結: \href{https://youtube.com/watch?v=NP-1P0Hd_Zc}{\texttt{https://youtube.com/watch?v=NP-1P0Hd\_Zc}} ~~~~ 語音日期: 2024-05-22
\newline
\newline
\hyperref[sec:NOP7HqFawO8]{\small{< < < PREV SERMON < < <}}
~
\hyperref[sec:index_chronic]{\small{[返順時目]}}
~
\hyperref[sec:index_scriptual]{\small{[返順卷目]}}
~
\hyperref[sec:OP86jLVzWZ4]{\small{> > > NEXT SERMON > > >}}
\newline
\newline
$^{1}$大家好又係到聖經考古ge 時間.
我相信大家都聽過Dead Sea Scroll死海古卷.
呢個係20世紀裡面最大ge 考古發現黎ge .
1947年ge 時候出土.
其實唔係單單係20世紀最大發現.
相信直到今日21世紀.
呢個死海古卷發現仍然係聖經考古裡面.
其中一個最重要ge 考古發現.
今日想同大家分享下.
死海古卷發現ge 傳奇故事.
死海古卷發現ge 地方係昆蘭.
昆蘭呢個名我地唔係好清楚真正ge 意思.
可能係出自一個阿拉伯字詞.
意思係邊白或者月亮ge 光ge 意思.
呢個意思白白地好似月亮咁ge 光.
其實對昆蘭呢個地方昆蘭黎講係非常之貼切.
因為如果你睇下昆蘭呢個地方.
就正正係係死海ge 西北岸ge 猶大礦業.
礦業呢一帶ge 石質都係黃黃白白.
又係好乾旱.
一D 比較軟ge 石灰岩soft limestone.
呢度就係寸草不生.
你見到呢個相片或者你曾經到過昆蘭ge 時候.
你都會明白呢一個地方gwo 個ge 環境係點樣.
發現第一批死海古卷ge 時間係1947年ge 時候.
死海古卷呢一個發現ge 故事係非常之傳奇.
並且係非常之曲折.
版本亦都好多.
將一D 重點今日同大家分享一下.
話說係1947年ge 時候.
兩位貝都恩人我地叫佢做Bedouin.
呢D 遊牧民族.
其中一位叫做Mohammad Adif.
佢綽號叫做豺狼.
據說豺狼就係昆蘭gwo 一帶.
去到當時搵緊一隻咪羊.
以為隻羊仔就躲係個山洞裡面.
於是佢就執起一塊細ge 石頭.
佢就掉入去山洞裡面.
咁希望將一D 響聲.

$^{41}$黎到佢讓隻羊仔驚ge 時候.
佢就從山洞走返出黎.
你見到呢個相片兩位貝都恩人.
右邊gwo 位就係豺狼.
點知豺狼將呢塊小石.
掟入去呢個山洞ge 時候.
佢聽到ge 就係啞屏破裂ge 聲音.
於是佢入去睇下ge 時候.
就發現有十個圓形ge 啞屏.
八個打開ge 時候係空ge .
第九個只係裡面一D 泥土.
第十個打開ge 時候.
裡面就藏住三卷古代ge 手抄本.
兩卷係有麻布包起.
之後呢D 貝都恩人.
就係山洞裡面又再發現.
另外四卷ge 手抄本.
因為呢D 手抄本.
藏jor 係啞屏裡面.
又係昆蘭猶大礦業.
一個咁乾燥ge 地方.
好似一個天然ge 博物館.
將呢D 古卷保留得比較完整.
於是頭一批手抄本就重見天日.
呢一個首次ge 發現.
總共搵到七份ge 古卷.
而呢一個首度發現手抄本ge 山洞.
就叫做Cave 1.
就將佢標籤為日號.
佢就係位於昆蘭以北大概半里ge 地方.
呢七卷ge 手抄本包括有D 乜ye 呢.
第一份就係一個完整ge 以賽亞抄本.
第二份係有部分ge 以賽亞抄本.
第三份就係昆蘭ge 社群紀律條例.
我地今日叫佢做Manual of Discipline.
一種社群紀律ge 手抄.
第四份就係舊約哈巴谷書ge Pesha.
Peshage 意思就係解經.
猶太人ge commentary.
第五份我地今日叫佢做戰卷.

$^{81}$Wars Group.
內容係講到呢D 昆蘭ge 社群ge 人.
一般我地認為起碼有一部分係gwo D 叫做愛色里人.
第一世紀一個猶太古歷史學家叫做Josephus.
曾經提過呢一班人係嚴守紀律住係昆蘭附近呢一個地方.
呢一班昆蘭ge 社群就自稱為光明之子.
佢地大概係公元前大概二三世紀ge 時候.
脫離jor 當時腐敗ge 耶路撒冷ge 宗教權貴.
佢地就遷徙去到東邊ge 曠野.
定居jor 係昆蘭呢個地方.
佢地有一個意識形態.
一個末世論.
認為末日彌賽亞出現ge 時候.
我地屬光明之子就會同黑暗之子.
代表住所有通通腐敗不忠於上帝律法ge 人.
最後會打一場ge 爭戰.
呢個就係大概Warscroll戰卷裡面所講ge 內容.
第六份ge 抄本係一D 詩歌.
一D Thanksgiving Hymns.
或者我地今日叫佢做Hymns Scroll.
希伯來文叫做Hadeach.
第七份就係創世記ge 魏經抄本.
呢D 係兩約中間猶太人一D 著作.
或者對聖經ge 幻想加jor 好多想像力ge 作品.
譬如呢本創世記魏經Genesis Apocrypha.
裡面就提到薩萊係點樣靚.
但創世記正典本身冇講.
但係呢份魏經就加jor 好多幻想上去.
去解釋薩萊點樣靚.
點樣由頭賺到落尾.
呢七份係第一個山洞裡面發現ge 古卷出現之後.
其他去到搵呢D 死海古卷ge 細節故事內容.
我地所知甚少.
我地只知道呢D 死海古卷係用處理過ge 羊皮所造成.
當時首次發現呢D 古卷ge 貝都因人.
係唔知道呢D 古卷ge 價值.
所以輾轉之下.
佢地就將呢D 古卷賣jor 比.
係伯利恆一個古董商.
佢ge 綽號叫做Kando.

$^{121}$Kando本身都係敘利亞ge 東正教徒.
後來發現古卷ge 消息就好快傳開.
係1947年11月28號.
當時希伯來大學ge 考國家就收到消息.
Schenklin當時冒險走入去.
當時仍然屬於阿拉伯人管理ge 伯利恆.
希望係Kandoge 手裡面能夠lor 到呢D 古卷.
係呢個時候就arm arm 係英國統治巴勒斯坦地ge 時間即將結束.
聯合國亦都快將通過劃分巴勒斯坦ge 計劃.
所謂UN Partition Plan for Palestine.
所以當時ge 政局係非常之緊張.
Schenklin第二日成功地係Kandoge 手裡面買jor 三卷古抄本.
呢D ge 抄本當時佢lor 返黎ge 時候只係放住係紙袋裡面.
Schenklin返黎到耶路撒冷ge 時候.
就正正係聯合國剛剛通過劃分巴勒斯坦地ge 計劃.
所以當時街上ge 猶太人就狂歡.
Schenklinlor 到呢三卷ge 古抄本.
其中一卷就係我地今日好多時聽到ge 偉大ge 以賽亞書卷.
Great Isaiah Scroll.
正係係呢個時候發現死海古卷ge 以賽亞書.
對猶太人無論係個人民族ge 意義係非常之大.
你記得現代ge 以色列國係係1948年5月28號正式復國.
因為Kando本身係敘利亞ge 東正教徒.
佢就將其餘gwo 四卷ge 古卷就用jor 24英鎊.
當時大概100蚊美金.
賣jor 俾耶路撒冷ge 另外一個敘利亞ge 教士叫做Metropolitan Samuel.
Samuellor jor 呢D 古卷之後.
佢就嘗試將佢出售.
但係結果都搵唔到買家.
所以佢就帶住呢四份古卷帶到去美國.
甚至係當時國會ge 圖書館做過一次ge 展覽.
只不過當時就冇買家有興趣.
結果Samuel就決定係華爾街時報Wall Street Journalgwo 度刊登jor 個廣告.
想去出售呢四份ge 死海古卷.
你見到呢個Wall Street Journalge classifier寫住.
The Four Dead Sea Scrolls, Biblical Manuscripts Dating Back to at least 200 B.C. are for Sale.
即係呢D 聖經ge 手抄本屬於起碼公元前200年.
而家就出售.
This would be an ideal gift for an educational or religious institution by an individual or group.
咁就話呢四份ge 古抄本係一個好好ge 禮物.

$^{161}$俾一D 學院或者係俾一D 宗教ge 機構.
但係非常之臭口.
係希伯來大學gwo 個考古學家Sukhinet.
佢ge 仔係Yigal Yedin.
Yigal Yedin當時arm arm 就係美國.
佢亦都係一位著名ge 聖經考古學者.
亦都係當時以色列建國ge 時候一個好重要ge 政客.
Yedin收到呢個消息之後.
佢就隱藏自己係代表以色列ge 身份.
就秘密呢D 古卷黎到買入.
點解呢因為當時昆蘭仍然係屬於約旦王國ge 政權.
所以佢地有呢個古卷ge 擁有權.
但係其實Mitchell,Barton,Samuel大可能都已經知道買家一定係以色列ge 代表.
因為成個世界裡面對呢D 古抄本有興趣ge 人.
相信一定係以色列ge 猶太人.
只不過佢地唔係好理會呢D 擁有權ge 問題.
最後Yedin用jor 幾多錢黎到成交.
就用jor 當時ge 美金25萬.
你諗下對當時40年代末ge 時候真係一個天價.
但係呢單ge 買賣對Mitchell,Barton,Samuel.
或者當時係美國紐澤西州ge 敘利亞ge 宗教群體.
其實好處唔係好大.
因為美國最後就向Mitchell,Barton,Samuel作出起訴.
就話呢一行ge 買賣係需要納稅.
所以大部分ge 收益都係歸返比美國ge Uncle Sam.
初步發現jor 第一批ge 古建之後.
消息傳得好快.
唔同ge 人就係困難ge 一帶就展開自己尋寶ge 競賽.
有貝都因人,有敘利亞ge 教導,Cando.
分別唔同ge 人就返返去第一個ge 山洞所在ge 地點.
但係適逢當時英國正要撤出巴勒斯坦地.
戰雲密佈.
所以一度以黎gwo D 學者都唔能夠去到困難一帶正式工作.
最後大戰過jor 之後.
1949年1月.
學者就第一次去到困難發現古建ge 山洞地點.
2,3月ge 時候耶路撒冷ge 法國聖經學院兩個考古學家.
一個叫做Warren DeVos.
另外一個叫做Harding.
佢地一齊就去到困難ge 第一個山洞.

$^{201}$做正式ge 科學考古法國工作.
亦都係呢個山洞出到更加多ge 手抄本.
亦都係困難ge 位置開始jor 初步ge 考古法國工作.
李佐德以為困難只不過係羅馬時期後期.
即係公元後三四世紀羅馬ge 兵寨ge 位置.
同死海古卷就冇關係.
但係一路工作ge 時候直至到1952年.
佢就證實jor 困難確實係當時猶太人社群所住ge 地方.
只不過去到公元後七十年羅馬帝國去到毀滅耶路撒冷ge 時候.
因為好多人走難ge 緣故.
所以當時困難社團ge 人.
佢地先將呢D 手抄本急急ge 收藏係山洞裡面.
係呢段時間貝多瀛又繼續擴大呢D 手塞古卷ge 範圍.
居然發現第二個山洞甚至係更多ge 古卷.
第二個山洞自然叫做Cave 2.
逼到耶路撒冷ge 學院就要出手.
其中第一個出手ge 學院就係American School of Oriental Research.
美國近東研究學院.
就去到軍難ge 一帶曠野ge 山洞進行系統性ge 搜索.
願意去到無孔不入ge 去做發掘ge 工作.
結果係1952年3月10至25號期間.
呢D 考古學家就發現jor 第三個山洞叫做Cave 3.
咁就出土jor 幾份ge 手抄本.
亦都發現jor 一個好著名好神秘ge 古卷叫做Corpus Glow銅卷.
因為文字係寫係一D 銅卷上面.
但係係呢個時候貝都欣人居然又再發現另一個山洞.
我地叫佢做Cave 4.
裡面因為冇D 啞屏裝住古卷.
所以已經分散到成千塊成萬塊ge 碎片.
所以後黎我地聽到死海古卷未能夠及時出版ge 爭議裡面.
有唔少ge 原因都係因為Cave 4係第四個洞裡面.
呢D 手抄本零碎碎片實在太多.
要花好多時間好多精力好多資金先至能夠重新整理出黎.
如果你有去過昆蘭ge 話.
相信你一定係昆蘭位置上面遠望過去對面.
就見到第四個山洞.
gwo 度有一個洞口當時係唔存在.
發掘ge 時候因為要有空氣流通ge 緣故.
所以考古學家先至作開jor 一個山洞.
所以到1951年為止.

$^{241}$昆蘭一帶總共係發現jor 11個山洞.
Cave 1 to 11.
出到jor 近成千份ge 古手抄本.
而且昆蘭一帶亦都發現jor 另外30個以上ge 山洞.
呢D 山洞裡面有唔同ge 瓷器.
有唔同ge 古物.
係尋寶競賽裡面.
貝都因人就發現jor 1 2 4 6 同11.
呢D 山洞裡面都有好豐富ge 重要ge 手抄本.
而考古學家就發現jor 6個山洞.
3 5 7 至到10.
咁你話係咪打成平手呢.
而家我地大概講到1953至到56年之間發生ge 事.
因為1951年ge 時候.
DeVos 同哈丁決定係昆蘭進行大規模ge 系統性ge 考古工作.
所以係53至到56年之間.
佢地就全面ge 去到發掘昆蘭.
發掘ge 報告終於係1993年同2000年ge 時候就出版jor .
呢D 好珍貴ge 死海古卷.
唔單止係一個好傳奇性ge 考古發掘ge 故事.
裡面ge 情節亦都係迂迴曲折.
更加係聖經考古裡面一個最大ge 發現.
唔單單係希伯來ge 首抄本.
成為jor 經文鑑定學ge 重要資料.
亦都可以幫助我地了解到.
當時兩約中間新約耶穌時代gwo 種背景.
呢D ge 死海古卷有好多唔同ge 文字記錄.
大概可以分開幾組.
第一組當然就係聖經ge 抄本.
即係呢D 困難社團ge 人.
不斷去抄寫舊約ge 聖經 希伯來ge 聖經.
第二大類ge 文獻就係困難社團ge 識經書.
譬如佢地會對先知書唔同ge 地方.
唔同ge 經文佢地會作一D 註釋.
第三類就係講到佢地自己困難社團裡面一D ge 內容.
譬如一D 紀律ge 搜查.
講到佢地點樣能夠進入呢個社群.
裡面要搜D 乜ye 紀律 有D 咩手執等等.
而家你要去睇死海古卷 困難ge 文物.
或者呢D 古卷ge copies.

$^{281}$最好ge 地方一定係去Israel Museumge Shrine of the Book展覽館.
你係外面會見到一個白色ge 圓蓋ge 設計.
呢個白色ge 圓蓋係代表住gwo D 收藏古卷ge 長身啞屏.
呢個白色ge 蓋下面就係展館最中心ge 地方.
裡面有Great Isaiah Scroll 偉大ge 以賽亞古卷.
同埋好多其他ge 手抄本.
蓋ge 顏色係白色.
就寓意係戰卷裡面呢D 困難社團.
認為佢地係Son of the Light 光明之子.
另外一度就係黑牆 黑白相對就係光明之子與黑暗之子呢一個陣字.
當然我地今日可以去到困難ge 位置參觀.
但亦都有唔同ge 工作繼續進行緊.
2017年ge 時候 希伯來大學ge 工作人員就係困難附近搵到再多一個ge 山洞.
呢個山洞相信曾經收藏一D 死海古卷.
只不過呢個山洞就冇搵到再多D ge 手抄本.
呢個其實係以色列文物局係猶大曠野一路好多年進行緊ge 有系統ge 考古工作.
目的其實係想策略地統計猶大曠野所有山洞ge 位置.
所以如果你睇唔同ge 參考書 死海古卷ge 翻譯著作等等.
你就會碰到一D 標籤 好多時有個number跟住有個Q字.
Q就代表係昆灣 而gwo 個數字好多時就代表係第幾個發現ge 山洞.
譬如1QISA 你就知道係係昆灣ge 第一個cave 第一個山洞.
而裡面搵到ge 手抄本就係ISA ISA.
如果你想睇死海古卷有關ge 文物 當然第一個地方就係去以色列博物館Shahnaan Lubuk.
但係今日亦都有一個好大ge 工程一路進行緊.
就係將重要ge 死海古卷ge 抄本就將佢數碼化.
用最高科技ge 攝影儀器就將呢D 古卷ge 真相能夠係網上面為我地呈現.
如果我地去Dead Sea Scrolls Digital Library ge 話.
你會發現裡面有唔同ge 資料 並且佢會將呢D 手抄本分成兩類.
一類就係Biblical Scrolls 即係聖經ge 抄本.
另外一類就係一D 非聖經ge 抄本.
譬如呢度見到有4Q Genesis 即係話係第四個山洞昆灣搵到創世記ge 手抄本.
佢屬於ge 時間就係100至1 BCE 即係公元前100年gwo 個世紀.
係屬於Late Hasmonean Period 即係瑪加比時代ge 後期.
而佢ge 語言就係Hebrew 係希伯來文.
如果你點擊ge 時候 你就可以係螢幕上面去到將佢放大縮細.
然後睇到呢D ge 手抄本裡面所寫ge 內容.
譬如呢度如果你認識少少希伯來文ge 字母.
呢度大概寫住 Elohim et Hashemayim ve'et Ha'ales.
呢度下面就係講到 Ruach Elohim 等等ge 字樣.
呢度就係創世記第一章第一節.

$^{321}$神創造天地 et Hashemayim ve'et Ha'ales.
天與地 The Heavens and the Earth.
最後希望呢D 資訊能夠幫助到你.
認識聖經考古裡面最大的發現 死海故卷.
並且能夠給大家一些工具.
讓大家了解到不單止發現的過程.
故事有多麼有傳奇性之外.
你亦都可以親自去睇下呢D 死海故卷.
去明白下第一世紀新約新舊約中間的歷史背景.
\newpage



\section{2:11-22}
\label{sec:OP86jLVzWZ4}
\textbf{Jesus Established the Chruch (Eph 2\_11-22) - Steve Tan}
\newline
\newline
連結: \href{https://youtube.com/watch?v=OP86jLVzWZ4}{\texttt{https://youtube.com/watch?v=OP86jLVzWZ4}} ~~~~ 語音日期: 2024-05-27
\newline
\newline
\hyperref[sec:NP_1P0Hd_Zc]{\small{< < < PREV SERMON < < <}}
~
\hyperref[sec:index_chronic]{\small{[返順時目]}}
~
\hyperref[sec:index_scriptual]{\small{[返順卷目]}}
~
\hyperref[sec:OVsK1efMh70]{\small{> > > NEXT SERMON > > >}}
\newline
\newline
$^{1}$- The phone?.
- Oh, yeah..
- All right..
Brothers and sisters, I realized I didn't really.
have a chance to tell everyone a little bit about myself..
And it's a good, I'm really happy to see.
you brothers and sisters in the church..
And it's a little bit about me is like,.
James, man, James, I mean, that song,.
"The Power of Your Love" was like something.
that I haven't heard for a while..
And it was the first time when Pastor Mark.
was my youth pastor, like close to 20 years ago..
And then it was the first time in the retreat.
that I felt like God convicted me to change..
I need to be different..
I gotta be better..
And that was the starting..
I was like, wow, my Christian journey,.
especially that verse, "Where your will and foe.
"in our lives.".
I'm like, here we are..
I graduated from Moody two years ago,.
a Master of Arts in Spiritual Formation and Cycle Ship,.
along with my beautiful wife right here, Shirley..
And right now, we're getting ready.
for the mission field in Taiwan..
Once Taiwan opens up, we get our visa,.
we're about to head out..
So I learned a lot about being a missionary,.
but that's God's will in my life..
From what I remember 20 years ago as a kid,.
as a high school kid, with my youth pastor,.
who's still my pastor right now..
So I'm just very grateful to have a chance right here.
to share God's word with you right here..
Let's pray together..
Father, we thank you, Lord,.
for the time that you give us right here.
to focus on you, Lord..

$^{41}$As you said, let tomorrow worry about tomorrow, Lord..
Help us worry about today..
And today is your day, Father..
Help us focus, Lord..
Holy Spirit, may you come into us..
Help us not worry..
Help us just focus on you to receive your words, Father..
To receive, to see how great your salvation plan is,.
Father, for every one of us..
In Jesus' name we pray, amen..
Hey, when I talk about the Cubs and the Sox,.
you have an idea what comes to mind..
How about the Bears and the Packers, how about that?.
If you've been in Chicago for a while,.
you probably have a good idea, right?.
These are a rival..
I mean, these are fans from opponent's side.
really hate the other team..
I mean, these are just basically friendly rivals we see..
But in society, as you know, there are many rivals..
And these could be religious rivals,.
like Catholic versus Christian in general,.
like economic rivals, right?.
Are you lower class, are you middle class,.
or are you in upper class, right?.
Or other rivals like, oh, the school that you go to,.
that distinguish, kind of like signify your intelligence..
And there could be family rival..
So I wanna tell you a little bit about a family rival.
that I had last year..
It happened during when my grandmother died,.
like last April..
And she was like 90 something already, about the time..
You know what the sad thing is that she died?.
But the saddest part is that all my uncles,.
when she wasn't even in her grave yet,.
they were talking about money..
That's the rival..
You guys know what I'm talking about, we're Chinese..
Yep, we talk about money, and they became enemy..

$^{81}$It's within these division.
that how Jesus established his church..
So today's topic I wanna talk to you about.
is how did Jesus establish the church?.
So in the book of Ephesians, actually,.
Paul did not chastise these Christians for anything..
It's more like an encouragement.
and laying out God's cosmic plan for them to understand..
So in the beginning, he talk about God's plan in general,.
which involves our spiritual blessing..
And in the beginning section of chapter two,.
he talks about how we are dead to our sins and transgression..
Therefore, God, through Jesus Christ, made us alive..
And then right now, we're gonna talk about the second part.
that is parallel to the first part..
Within this, we're talking about another way.
how we are dead to Christ in a different standpoint..
So in my first point,.
I wanna talk about hopelessness without Christ..
How there's no hope without Christ..
Let me start reading..
Therefore, remember that formerly you who are Gentiles.
by birth and caught uncircumcised.
by those who cautimcise the circumcision,.
which is done in the body by human hands..
Remember that at that time, you were separated from Christ,.
excluded from citizenship in Israel and foreigners.
to the covenant of the promise..
Without hope and without God in the world..
But now in Christ Jesus, you who once were far away.
have been brought back near by the blood of Christ..
So Paul start off by reminding these Christians in Ephesus,.
these Gentile Christians,.
there is a group called circumcision..
Mainly, these are Jewish people.
who probably have accepted Jesus Christ already,.
but still thinks that circumcision is necessary..
So they're saying these groups are mocking you guys.
because you are not circumcised yet..
But the thing is that Romans and Greeks.

$^{121}$don't call themselves Gentiles..
The only people that call themselves,.
call other people Gentiles are Jews, the Jewish people..
So why is it so important to distinguish.
that they are Jews and those other people are Gentiles?.
So if you take a look,.
let's go back to the book of Genesis..
This verse right here is basically the backgrounds.
where God and Abraham,.
God made a covenant with Abraham, an agreement..
He's telling Abraham,.
all the nations on the earth will be blessed because of you..
And this circumcision is this phrase.
that you are to undergo circumcision.
and it will be the sign of the covenant between me and you..
This is a physical sign to show how they are different..
But the problem is that a lot of times.
these Jewish people use,.
think of the sign as seeing themselves as superior..
Like they are the special race..
They have God with them..
They are part of God's family..
This covenant is telling them,.
we are part of God's family and you are not..
So they think of the Gentiles as a lower race, outsiders..
But when Paul is writing this,.
God, Jesus Christ, step on the cross,.
has redeemed us already..
Therefore, the old order of the Judaisms,.
their external and material feature is outdated already.
because Christ came in and pushed in a new order..
Therefore, if I go back here and you take a look,.
Paul specifically says,.
which is done in the body by human hand..
What does that mean by human hand?.
Paul emphasizes because if it's by human hand,.
that means it's not from God..
Therefore, what Paul is saying,.
he's mocking these people also,.
these circumcision, these Jews..

$^{161}$He's trying to tell them,.
hey, you guys don't have the real circumcision either..
Like, let's take a look what's the real circumcision..
The same author of Genesis, Moses,.
he wrote through to Tyrannomy..
He says,.
Circumcision is,.
God has to do the real circumcision through the heart.
where we have the ability to obey and follow his order..
So how does Paul look at that, right?.
In Colossians, this is how he sees it..
In him, you were also circumcised.
with a circumcision not performed by human hands..
Your whole self ruled by the flesh was put off.
when you were circumcised by Christ..
So as we know now, how does God circumcise us?.
It's through Christ..
Christ has to come and die for us.
and change us into a new person..
That is the real circumcision.
that's available to the Jews and Gentiles currently..
(pages rustling).
After talking about circumcision, the physical sign,.
Paul continues to remind these Gentile Christians.
in Ephesus about their former state without Christ..
Why?.
This is to show them how blessed that they are.
at their current position..
So this is what he says to these people.
before they know in Christ..
Remember that at that time you were separated from Christ,.
excluded from citizenship in Israel,.
and foreign to the covenant of the promise..
Basically, this is what Paul is saying..
Israel has an advantage over you..
The Jewish people have an advantage over you.
because they're part of God's family, but you are not..
So it sums up basically this is what he's saying..
The Gentiles, you don't have anything related to God..
You're isolated from God, right?.

$^{201}$So the result comes at the second part of the verse..
You are without hope and without God in this world..
So if you take a look, verse 11 and 12,.
look so dark and gloomy for the Gentile.
in their former state before knowing Christ..
But again, Paul states but in verse 13..
But Christ right here made 180 turn again for them..
You know, he says, "But now in Christ Jesus,.
"you once were far away, have been brought near.
"by the blood of Christ.".
Instead of being dead to our transgression,.
we were once far from God, but Christ brought us near..
What does that mean?.
Bring us near means that we have access to God.
and that access is through the blood of Jesus Christ..
And this gives me, helps me think of this illustration..
Actually, it's also from Paul..
He talks about a branch grafted in in Romans..
And if you take a look, the lower branch,.
the main branch is very different from the top branch.
that is grafted in, you know?.
The top branch doesn't belong..
And take a look at how this person.
attached these two branches together..
It's by some sort of tape or something to glue it on, right?.
That tape is basically the blood of Christ..
He used that to glue us into the main branch.
that we could be part of the same family..
Therefore, have a great appreciation.
for what Christ did for us when we were far away from God..
Therefore, he used his blood to draw us near..
Now let's go to the second point..
The second point is what Christ did to make the church..
So for he himself is our peace,.
he who has made the two groups one.
and has destroyed the barrier,.
the dividing wall of hostility,.
by setting aside in his flesh the law.
with its command and regulations..
His purpose was to create in himself a new humanity.

$^{241}$out of the two, thus making peace,.
and in one body to reconcile both of them to God.
through the cross by which to put to death their hostility..
He came and preached peace to you who were far away.
and peace to those who were near..
For through him we both have access.
to the Father by one spirit..
So as we talk about how Christ in the previous point,.
he restored our relationship with God..
So in verse 14, Paul starts out by emphasizing.
that Christ is our peace..
But what does this peace mean?.
In general, peace means there is no hostility, right?.
Everyone get along..
But how do you gain that peace?.
How did Christ gain that peace for them?.
Well, you had to join two groups into one..
Those two groups that we know are the Gentiles and the Jews..
So how did he go on doing so?.
The first, there is two things,.
that's one thing you have to do..
Let me read it according..
He had to destroy the barrier,.
the dividing wall of hostility by setting aside.
in his flesh the law with his commands and regulation..
So this is referring to the Temple of Jerusalem, guys..
So if you take a look, where I label Gentile,.
that's the outer court..
That is as far as you can go if you're not Jewish..
If you're a Jew, you could go into the inner court.
to worship God..
There is actually a physical wall.
that divides the two people at a time..
So, but the thing is that,.
that's only a manifestation of the Mosaic law,.
the regulations of the command..
That's why Jesus had to come and abolish those commands..
Those regulations and law, to tell you honestly,.
even the Jewish people, they can't meet themselves..
That holiness is so high, the God standards is so high,.

$^{281}$that it's really impossible to meet in human standpoint..
But besides separating these two groups religiously,.
there are also separated,.
the law also separated them socially..
Because of this law,.
the Jewish people have to live according to,.
they don't live according to the custom of the world..
So they kind of like lost communication in that point,.
but consequently at the same time,.
it gives the Jewish people that pride..
They feel like the way how we live, we're superior..
The way you live..
But the thing is that at the end,.
Jesus came and abolished that law once and for all.
with his flesh to present his perfect sacrifice to God..
So we don't have to worry about this anymore,.
these regulations..
Now let's move on to the second part of verse 15 and 16..
So as we move on,.
so the purpose of abolishing these hostility is twofold..
First of all, in verse 15,.
the second part is to create in himself a new humanity..
Basically,.
the rift between the Jews and Gentiles are too deep..
You cannot put them together anymore..
Think about in our country,.
you think about the political state right now,.
the two big parties, the Democrats and the Republicans,.
you can't even work together anymore..
The only thing they could do right now.
is call each other name..
We're lucky that our country's continued to function,.
seems like quite normally right now..
Therefore, God has to push a new humanity..
Within this new humanity, a new order, a new covenant.
has to be established by Christ..
And as I said again,.
Jesus is the one that kind of present himself.
as an offering to reconcile human between human..
So he reconciled the horizontal hostility,.

$^{321}$the horizontal dimension of the hostility..
And as we go on to verse 16,.
in that new body is to reconcile.
both this new humanity to God..
So his blood not just reconcile.
this our human to human hostility,.
it also reconcile the vertical dimension,.
our humanity to God..
So as we talk about the reconciliation process,.
Paul now talks about the proclamation of peace from Christ..
In verse 17, it says, "He came to preach peace.".
How does Christ preach peace?.
He did it by walking, by doing, by walking the talk..
He was, he is our peace in verse 14,.
that says, "For he himself is our peace.".
Again, by using his body, the flesh, the blood,.
to bring, to end that hostility..
And who are he talking to?.
Again, he's emphasizing the ones who are near,.
which are the Jewish people..
They have, think about it,.
they have all the time, they have the temple,.
the laws, the worship regulation, the covenant,.
and the Gentiles, the ones who are far,.
they don't have anything..
But Paul is saying, in verse 16,.
"Even though they have all these things, the Jewish people,.
"they still need reconciliation.".
Again, in verse 18,.
Paul continued to further the concept of peace..
Besides just being right with God, we have access to God..
We, as Jews and Gentiles,.
we stand as one humanity in front of God..
And there is no distinction anymore,.
because it's just by one spirit,.
we have access to the Father..
It stresses that oneness that we have, the new humanity..
And within this new, excuse me,.
within this new humanity,.
I don't have to talk about my credentials,.

$^{361}$my spiritual heritage, my zeal,.
or my talent to access God..
The only thing we all need to share in common.
to have access to God.
is believing in the redemptive work of Jesus Christ..
This brings me, reminds me of this person right here..
Do you know this person?.
His name is Jim Elliott..
He was a missionary to Ecuador,.
especially to the Native American tribes in around the 1950s..
He and his fellow companions, missionaries,.
have decided, "We must let these people hear the gospel.".
But the thing is that they know.
these Indians are hostile to foreigners..
Yet they went anyway..
What's the result?.
He got martyred..
He and his fellow missionaries got martyred..
Yet creating this hostility between their families.
and these Indian tribes..
But instead of seeking revenge,.
their wives decided to continue on this mission,.
decided to go into these tribes and live with them,.
to walk out that gospel..
At the end, bringing the gospel,.
bringing the conversion to these Indians,.
thus ending the hostility..
Therefore, Jesus' work of redemption brings us this peace..
We might have an understanding how hard,.
how much work he had to go through to start this..
Now let's go on to our last point..
What the church looked like..
Consequently, you are no longer foreigners and strangers,.
but fellow citizens with God's people.
and also members of his household..
Built on the foundation of the apostles and the prophets,.
with Christ Jesus himself as the chief cornerstone,.
in him, the whole building is joined together.
and rises to become a holy temple in the Lord..
And in him, you too are being built together.

$^{401}$to become a dwelling in which God lives by his spirit..
So as a result of God's redemptive work,.
Paul's telling us our status in front of God has changed..
And Paul wanted to talk about.
how this new humanity functions,.
which we will see the characteristics of the church..
Paul starts out by using, again, the spatial relationship.
to show the differences before and now..
He says, consequently,.
you are no longer foreigners and strangers..
What does foreigner and stranger,.
when you think about that, what comes to mind?.
You live in a land that is not your own homeland..
You live with the people who are not your people,.
who are not your tribe..
So what does that mean?.
You are outsider..
You don't belong..
But Paul continues to elaborate..
But fellow citizens with God's people.
and also member of his household..
We are citizens..
We are part of God's family right now..
We are not second-class citizens..
We are the same of the Jewish people..
Or let's say this new humanity that Christ created..
As stated in chapter one, verse five,.
God already said he has destined us.
for a daughter and sonship..
This is the closest relationship.
that you can have with someone..
But think about the process of you becoming our citizen,.
like the ones that have become, right?.
You have to come to this country..
You have to become a good member for a couple years..
Make sure you follow the law.
and be a good person in society..
And then you have to study for a test..
And then you make sure you pass the test..
Then you might be able to become a US citizen..

$^{441}$But what happen if I say,.
today the US president came to you personally and said,.
"You don't have to do any of that..
"I will accept you in as a member..
"At the same time, I will adopt you into my personal family.".
That is what the household of God looks like..
Now in verse 20, Paul can you elaborate.
how this household is built upon, this universal church?.
He said, "It's built upon the foundation.
"of the apostles and the prophets.".
So the apostles, it's basically talking.
about the 12 apostles and Paul himself,.
who Jesus personally chose to proclaim the gospel..
And how about these prophets?.
These prophets, based on this context,.
I would see they are not the Old Testament prophets.
because Paul in chapter three, verse five says,.
"The mystery of the gospel is not known.
"to any other generation until now.".
So I would say these prophets are the people.
along with the apostles who might have walked with Jesus,.
or these prophets are people that the apostles converted,.
and they continue to preserve the teaching of Jesus Christ,.
because if they don't continue to preserve the teaching,.
this new humanity will not have its foundation to rely on..
But however, that's not the main and most important element..
As Paul continue, he says, "Jesus Christ himself.
"as the chief cornerstone.".
As you see, Paul emphasize this point..
He uses the word himself to tell.
that Jesus is the chief cornerstone..
What does that mean?.
In the ancient world, the builders,.
the first stone that they lay is the cornerstone..
That's the most important stone..
It's the stone that gives the direction.
where the building's gonna go,.
and it determines if that building will be stable..
And other stones are built upon this cornerstone..
Therefore, Jesus ensure the stability.

$^{481}$of the system of salvation..
As we continue on to the last two verses, 21 and 22,.
Paul continue to elaborate that we also are the stones.
along with the prophets by using the word you too,.
that we, well, for him, he's meaning the Gentiles,.
but we are the Gentiles too right now..
We take part in building this temple too..
Again, within those two verses, Paul is saying,.
"We are joined and built together.
"to construct this household.".
It gives us, symbolize the closeness.
that we have between believers..
And it also shows us this church.
should be a continuous growing body.
because through our evangelism and discipleship,.
it should continue to be growing..
And at the end, the last verse, it says,.
Jesus also dwells within his household..
It shows that God and the people, the believer,.
we should have a close relationship also..
But overall, Jesus is in charge.
of connecting all these stones..
Therefore, he is responsible.
of retaining all these relationships, not us..
But everything is emphasized on Jesus Christ himself..
Everything is linked to him..
The house functions probably only because of Jesus..
We, the people, are the church..
We're not talking about the building, but Jesus is the head..
Now, this brings me to this picture..
If you guys take a look, as you take a look,.
something's wrong with this building..
This is actually a building that collapsed.
because of an earthquake that happened.
in 2017 in Mexico City..
Actually, Mexico is known to be a place.
with a lot of earthquakes,.
so their building regulation is pretty rigorous..
However, why so many buildings still collapse.
is mainly due to because of corruption..

$^{521}$A lot of these engineers still sign off on that blueprint.
knowing that they do not meet the safety inspection..
I mean, when glamorous time, when everything is fine,.
I mean, the building looks fine, right?.
But when hard times come, you can see what happens..
But I could assure you,.
our Christian faith in the church.
has endured a lot of hardship and trial.
throughout the many centuries, right?.
Yet, it stands solid as a rock.
because who is our cornerstone?.
Jesus Christ..
Therefore, we continue to be confident.
to place ourselves as rocks.
to continue to build this household, this church..
In conclusion, let me get back to my story.
of my family few..
My uncles, they're currently still enemies..
They have a hatred toward each other,.
but from the way I see it,.
there's a sense of hopelessness in them..
They don't know, they can't see the light.
at the end of the tunnel..
Because there's no mediator for them to help them see it..
But unlike my family member, unlike my uncles,.
the church has Jesus Christ as the cornerstones,.
brothers and sisters..
He's the one that made the peace.
between human to human and with God..
Therefore, we should have an attitude.
of profound thankfulness towards God.
because we are in a very, very privileged position.
as member of God's new community, which is the church..
Let's pray..
Father, we just wanna thank you, Lord..
Just help us see the big contrast of where we were..
We were not your family..
We were cut off..
We don't have any part of what the Israelites had..
And our goal, the end point was hopelessness, Lord..

$^{561}$But again, Father, Jesus changed that for us.
without us doing anything, while we were still his enemy,.
Lord, and he changed us and he put us as citizens.
right here along with everyone else.
who profess faith in Jesus Christ,.
where we get to build the church along with him, Lord..
We thank you for this privileged position, Lord..
Help the brothers and sisters right here.
have that, meditate on that, Father,.
that this position is not something.
that we could take granted for, Lord,.
but to be thankful every day of where we are,.
especially our status with you, Father God..
We thank you so much..
In Jesus' name we pray, amen..
[BLANK AUDIO].
\newpage



\section{}
\label{sec:OVsK1efMh70}
\textbf{聖地網上行 - 步主腳蹤 [第十四集] 死裡復活}
\newline
\newline
連結: \href{https://youtube.com/watch?v=OVsK1efMh70}{\texttt{https://youtube.com/watch?v=OVsK1efMh70}} ~~~~ 語音日期: 2024-05-27
\newline
\newline
\hyperref[sec:OP86jLVzWZ4]{\small{< < < PREV SERMON < < <}}
~
\hyperref[sec:index_chronic]{\small{[返順時目]}}
~
\hyperref[sec:index_scriptual]{\small{[返順卷目]}}
~
\hyperref[sec:qLV2rgTxxTw]{\small{> > > NEXT SERMON > > >}}
\newline
\newline
$^{1}$各位,我們來到最後一集聖地望上行部主腳鐘.
我們已經有十多次的旅程.
上一集我們停留在苦路Villadrosa最後一站第十四站.
這裡就是聖母堂所標記的地方.
紀念耶穌從死到復活的墳墓.
如果你有去過耶路撒冷玩.
你可能聽過另一個說法.
就是有另一個地方才是耶穌基督死而復活的墳墓.
那個地方叫Garden Tomb花園叢.
這個地方位於耶路撒冷舊城北面.
大馬士革門外面.
有人認為這個地方才是耶穌的墳墓.
因為這裡的面貌和山上的陽色.
就像一個骷髏頭.
所以十九世紀末的時候.
有一群基督徒望見這個山的時候.
他們不太接受聖母堂裡面東正教天主教的儀式.
所以他們認為這裡才是足留地.
因為足留地的意思就是骷髏頭的意思.
大家留意聖經所描述的足留地Gothic.
其實不是說一個山的面貌.
而是說這個地方是一個死刑場.
place of skull 死亡死刑的地方.
更何況二千年前一世紀的時候.
這裡的山的樣貌不是這樣的.
有一個這樣的骷髏頭的面貌.
是後世才出現的.
一世紀的時候不是這樣的.
但是有很多的旅客來到耶路撒冷的時候.
他們很多時候會來到花園叢.
說真的這個花園叢也有吸引的地方.
因為在耶路撒冷的舊城甚至聖母堂.
都給人一種繁忙很多張力的感覺.
這個花園叢特別幽靜.
所以有很多基督徒的客人很喜歡這個地方.
來講解聖經,默想甚至在這裡受聖餐.
按照考古,文獻,歷史,傳統等等.
聖母堂Holy Sepulchre.
仍然是耶穌死而復活,墓穴的所在.
沒錯經過很多很多年.

$^{41}$幾乎二千年的拆毀又重建.
都已經面目全非.
又或者對於你的傳統.
你來到這個地方.
你看到天主教甚至是東正教.
他們的儀式不是我們很熟悉的東西.
不如我們開放一點來欣賞一下.
他們曾經做過的事.
其實是很多時藉著很多很多年.
東正教的努力才能夠保留到這個地方.
雖然現在已經是面目全非.
但我們可以藉著考古,文獻,歷史等等的資料.
將這個地方稍為還原.
這個畫面圖就可以解釋一下.
到底這裡之前發生了什麼事.
我們回到公元後的三十年.
即是耶穌死而復活的事跡的時間.
這個地方足留地位於城牆以外.
這裡本身是一個石礦場.
所以經年累月.
很多好的石塊早已被人採了.
留下來的只是石質較差的石地.
所以經年累月就長出很多花草.
因為石質較差.
所以適合開採墓穴.
所以這一帶的地方在城門以外.
很多時是墳墓和花園.
在以前的歷史文獻中告訴我們.
這裡有一個城門叫花園門.
出來這裡就看到一個像花園的地方.
或者是一些墳墓開鑿的地帶.
一世紀墳墓的格式是怎樣的呢?.
你記得我們曾經去過那薩勒村莊.
那裡有一個人工造出來.
仿效一世紀墓穴的墳墓.
是從石山挖掘出來的.
首先我們進入墳墓有一個內室.
這個內室是用來預備屍體的.
例如人們會帶著沉香活藥.
將屍體預備好.

$^{81}$然後將屍體放入山邊作出來的長洞穴.
希伯來文叫作Kuk.
或者眾數叫做Kukim.
所以有時候會叫做Kukim Tomb.
Kukim的墳墓.
將屍體包起來放入長洞穴.
一年之後屍體腐化.
家屬會回到墓穴的樹.
然後將骨頭拿出來.
那叫做執骨禮.
然後將骨放在骨盒的樹上.
之前我們也說過.
所以在公元後三十年的時候.
這裡有一些這樣的格式的墳墓.
其實在聖母堂附近.
我們也找到這些格式的墳墓.
到了公元後二世紀.
羅馬王為了要將耶路撒冷全面羅馬化.
所以在墳墓之上建立了羅馬神廟.
為了取代初期耶穌之墓的傳統.
這個舉動其實是羅馬帝國.
想要消滅基督教的信仰傳統.
但它無意之中保留了這個墓地的地方.
所以公元第四世紀的時候.
當時君士坦丁執政的時候.
人們來到這裡找回這個教堂的地方.
經過很多世紀的戰亂和毀滅.
終於到了十字軍的時候.
十二世紀重建了現在很宏偉的聖母堂.
Holy Sepulchre.
有兩個圓頂.
現在我們來到耶路撒冷.
Holy Sepulchre 聖母堂遊覽的時候.
進入這個建築物.
大多數我們看到的就是十字軍留下的遺跡.
起碼有千年以上的歷史.
在教堂的不同地方.
石牆上刻了很多小的十字架.
這些都是十字軍時期.
朝聖者來到這裡留下的痕跡.

$^{121}$聖母堂最中心的地方.
當然就是大圓頂下的小教堂.
2016年的時候.
曾經經歷過一個很大型的裝修工程.
這個裝修工程完成了.
亦找到比十字軍更早期的遺跡.
一直以來這裡有很多客人排隊.
進入這個小教堂觀賞.
我帶你進去看看.
這就是世界基督教最神聖的地方.
這裡可以看到一個很小的祭壇.
這就是標記著耶穌基督死而復活的地方.
正如過往的旅程中.
我都告訴大家.
當然我們不能夠直接去證明.
這一處就是耶穌的屍體曾經躺在這裡.
然後在這裡死而復活的地方.
但這一處絕對是聖經所提及的.
歷史文獻所提到的死而復活的墳墓.
你讀一讀四本福音書記載耶穌死而復活的經過.
時間表大概是這樣的.
我們想像一下一個一世紀的墳墓的情況.
當時是星期五的黃昏.
耶穌在十字架上死.
然後他的屍體被解除下來.
因為當時是猶太人的安息日.
日落就要開始.
安息日開始就不能夠工作.
所以門徒很匆忙地把耶穌埋葬.
墳墓用石封上.
對於門徒來說.
這一切都結束了.
一切都是失敗.
之後安息日就不能夠工作.
只能留在家裡.
所以沒有什麼特別的事情發生.
但是耶穌的故事沒有在這裡結束.
反而是在這裡開始.
因為星期日的清早.
本身是跟隨耶穌的父女.

$^{161}$他們想回到耶穌的墳墓.
你可以想像這裡的情況.
他們想完成聖期五草草了事的葬禮.
想做好.
然後把屍體放在牆的山洞裡.
誰知道父女去到墳墓的時候.
發現墳墓是空的.
復活主向他們顯現.
於是他們是第一班人.
很急忙地將復活的消息通知.
當時躲藏的門徒.
而其他門徒也趕來墳墓觀察.
發現墳墓也是空的.
接著發生什麼事呢?.
就是復活主多次向門徒顯現.
也向保羅顯現.
使這些人的生命有一個180度很大的轉變.
這就是基督教信仰.
聖經信仰最獨特的地方.
其實由舊約至新約都是這樣.
聖經信仰是建基於歷史事件.
基督教信仰.
聖經信仰絕對不是建立在一些.
可以自我自身的感受.
或者一些宗教的感覺.
更加不是建立在一套學說上.
而是建立在歷史事件上.
如果比較其他宗教的經書.
你會發現聖經有一個特別的地方.
特別經常提及一些歷史事件.
或者一些地理實地的情況.
這就是聖經特別的地方.
哥林多前書第十五章.
說:基督若沒有復活.
你們的信變是枉然.
你們仍然在墜淚.
這裡說:我們若靠基督.
只在今生有子亡.
就算比眾人更加可憐.
這裡是一個很大的.

$^{201}$你可以說是賭注.
就是說基督教信仰就是建基於.
耶穌的死裡復活.
這個真確歷史裡面的事件.
你看看各本的福音書.
馬太,馬可,路加和約翰.
你看看裡面的一些證供,文件.
他們記載耶穌基督復活的事件上.
都是各有差異的.
我舉個例子.
路加福音第十四章.
他說:七天的聖餐.
是基督復活的.
而耶穌基督復活的.
是基督復活的.
我舉個例子.
路加福音第十四章.
他說:七天的頭一天.
就是星期天.
黎明的時候.
那些婦女帶著所預備的香料.
趕到墳墓.
看到石頭已經從墳墓滾開.
他們就進去指示.
不見主耶穌的身體.
這裡說到多少個婦女呢?.
一群的婦女.
但到約翰福音記載的時候.
七天的第一天清早.
時間一樣.
天還黑的時候.
末達拉的瑪利亞來到墳墓那裡.
看見石頭從墳墓滾開.
就跑去見西門彼得和耶穌所愛的那位門徒.
應該是約翰.
對他們說.
有人把主的身體從墳墓裡拿走了.
我們不知道放在那裡.
你看到有什麼差異呢?.
路加福音說有一群婦女來到墳墓.

$^{241}$約翰福音的鏡頭只留意一個女人.
末達拉瑪利亞.
你可以看到那種差異不是自相矛盾.
因為他沒有說只有她一個.
只不過是鏡頭的焦點放在.
末達拉瑪利亞一個女人身上.
為什麼說這些報導.
大家有差異又不是自相矛盾.
才是寫實,如實的報導呢?.
我今天分享這個訊息的時候.
是2020年9月13日.
前兩天新聞再報導2001年911發生什麼事.
你看看不同的報導,不同的故事.
直至這麼多年講出來的時候.
差異很多時候都很大.
但是這些故事,這些報導雖然有差異.
卻沒有矛盾.
雖然事後這麼多年講出來.
但是有些很重要的內容是大家是吻合的.
是不會懷疑的.
例如是什麼呢?.
第一是飛機撞入雙塔.
第二是雙塔倒塌了.
這兩個絕對是事實.
無論這些報導之間差異有多大.
沒有一個會懷疑,質詢.
到底有沒有這些事情發生.
另一個例子就是猶太人在二次大戰的時候遭受大屠殺.
現在近代的猶太人都不斷不錄,記錄幾十年前.
這些大屠殺的生還者的見證故事.
你聽回這些故事的時候.
很多都是有差異的.
但是沒有懷疑有大屠殺這件事發生.
這些例子讓我想到.
歷史的文件,歷史的記錄,見證的故事.
就有一個這樣的特色.
復活的故事其實都是一樣.
四本福音書的記載死而復活的事件.
彼此是有差異.
但是沒有彼此互相矛盾.

$^{281}$沒有經過後人修飾.
沒有人去將這些不同的故事協調.
盡相同的細節去剷除.
是很如實,寫實的一些報導.
但是這些各種記錄當中.
你去比較的時候.
起碼有兩件事是吻合的.
第一就是空墳墓.
第二就是婦女是第一班見證人.
如果我們抽身回到當時的文化.
當時的社會的話.
你就知道有什麼特別的地方.
居然是女人第一班去見證.
不是男人.
女人當時是沒有法庭,政公的地位.
男人才有.
所以是一種很寫實的報導.
再進一步你想想.
初時打壓基督教的耶路撒冷權貴.
都沒有能力提出一些證據去否定空墳墓.
再想一想那些本身貪生怕死的.
躲起來很灰心已經散去的門徒.
居然有一些東西扭轉了他們.
以致他們以後不惜代價.
甚至殉道來見證耶穌的復活.
再想一想本身是誤信.
甚至是迫害初期信徒,初期教會的保羅.
這些就算是誤信耶穌的學家.
都會認為保羅是一個歷史人物.
是歷史的保羅.
他寫了不少的書信.
保羅反而由一個迫害初期教會.
變成一個傳耶穌復活信息的人.
所以問題不是耶穌的墳墓是否空.
按照見證,文件,紀錄.
確實是空墳墓.
問題是如何解釋空墳墓.
要解釋的話.
這絕對是超越了科學,歷史能夠處理的.
因為這本質上是世界觀的問題.

$^{321}$什麼叫世界觀?.
世界觀就是每個人對生老病死,人生.
對整個世界都會有一套預設.
這些預設很多時候不是科學推理出來的結果.
相反是我們先有這樣的設定.
然後去解釋其他東西.
譬如有人會認為人死就如燈滅.
人沒有什麼永恆價值.
這些不是實驗證明出來的結果.
這些往往是我們的世界觀有了的預設.
有了這個預設我們就會去接受一些東西.
並且可能會拒絕一些東西.
聖經的信仰就是要去重整我們的世界觀.
信耶穌其實就是一個世界觀的轉移.
或者我們想起耶穌從死裡復活.
其實是顛覆重整我們本身有的世界觀.
你想想從第一世紀開始.
耶穌從死裡復活所帶來的影響.
就是對人的世界觀的震撼.
讓人意識到原來有一個更高的現實存在.
原來世界不是一個封閉的系統.
神是真真實實活著.
他創造萬有.
他與我們可以建立關係.
他以恩典.
神的兒子死裡復活來救贖我們.
這是一種重整的世界觀.
一種震動我們本身有的世界觀.
讓我們想清楚原來我們以往的觀念是錯的.
現在我們看世界完全截然不同.
因為耶穌從死裡復活.
我們的人生最後堅定陳詞的不是死亡.
不是羅馬帝國.
不是苦難.
而是上帝天國和神的權能.
永恆才是最真實.
天國就是重整我人生的觀念.
對我人生的方向都有影響.
可能以前那一套.
我以為人活著就是為了安樂茶飯吃.

$^{361}$我不理會社會發生什麼事.
總之我有安樂茶飯就行了.
現在耶穌從死裡復活.
我們生命的追求自然不是這個方向.
而是追求行神的旨意.
神對我對你此時此地的附接.
我們是否忠心來跟隨.
對於第一代猶太人基督徒.
彼得,約翰,雅各等等.
有什麼的震撼呢?.
猶太人大多數都相信身體復活.
只不過他們相信的身體復活.
就是世界終末的時候將來才會發生.
但是突然間他們跟隨主.
耶穌從死裡首先復活.
就將世界終末會發生的事.
帶到現在.
天國突然間不是單單將來的事.
而是現今此時此地要實現的事.
所以對當時的猶太人第一代基督徒.
有一個很大的震撼.
天國不是在等待將來會發生的事.
天國是現在在我的生命裡要發生的事.
神要藉著我去彰顯神自己的榮耀.
彰顯神自己的生命.
另一方面當時的猶太人有一個很強烈的世界觀.
那個觀念可以叫聖地主義.
Holy Land Theology.
世界最神聖的地方就是聖殿.
聖殿裡面的致聖所就是神住的地方.
神榮耀的地方.
那裡才是最神聖的.
離開這個聖殿越來越遠的話.
就越遠離上帝.
所以為什麼今天你去到西牆的時候.
為什麼有這麼多的猶太人覺得那是一個聖地.
因為西牆是最近以往的聖殿.
就是現在的清真寺下面的地方.
甚至有些猶太人會去到西牆的隧道.
就是西牆下面那裡.

$^{401}$那裡有一個角落.
很多時候人們會去那裡祈禱.
因為那裡畫一條垂直的線就是最近致聖所.
這就是猶太人很強烈的聖地的神學.
但是到耶穌基督復活之後.
祂就將天國彰顯出來.
將神的榮耀彰顯出來.
在哪裡顯現出來呢?.
我們看幾處聖經.
路加方第24章.
正當那日門徒中有兩個人往一個村子去.
這條村子叫做二馬五士.
離耶路撒冷約二十五里.
這裡有兩個門徒.
代表當時所有的門徒.
他們以為耶穌死了.
什麼都完了.
非常灰心.
他們就向一個村子叫做二馬五士走.
這裡說得很清楚是離開耶路撒冷的.
其實二馬五士無論今天在哪一處.
都是耶路撒冷以西的一個村落.
但是他們離開耶路撒冷的時候.
耶穌復活的主就向他們顯現.
在這個故事裡面.
一直與他們同行解釋和講道.
告訴他們舊約聖經指向的應許.
就是離開受苦之後會復活.
他們到了村子的時候.
他們就和這些門徒一起吃飯.
將餅分開.
讓門徒看到最後晚餐的動作.
門徒當時才明白.
原來耶穌真的如他所說.
要從死裡復活.
才會來.
這兩個門徒代表教會.
代表教會要領受主的道.
他們聖餐一起吃飯的時候.
代表教會走在一起的時候.

$^{441}$他們領受聖餐.
在這個道當中.
在聖餐裡面.
復活的主就向他們顯現.
所以今天我們去聽見上帝的道.
今天我們去領受聖餐的時候.
路加夫似乎告訴我們.
復活的主是真真實實和我們在一起.
不是單單一個紀念.
而是真真實實.
上帝的榮耀在我們中間.
到了史坦傳.
路加福音的延續.
其中一個故事是.
保羅本身是猶太人的拉比.
他迫害當時的教會.
他去大馬錫這個地方.
大馬士革來拘捕基督徒.
在大馬錫的路上.
他就遇見榮耀的復活主.
在當時的大馬錫路上.
就是今天的戈蘭高地.
你記得我們曾經去過.
這個就是凱撒利亞菲勒比一帶.
你記得登山變象的記載.
就是說凱撒利亞菲勒比一帶.
戈蘭高地的一帶.
本身是一個愛邦人的地方.
對一個傳統正統的猶太人保羅來說.
是一個敵方的地方.
是一個愛邦人的地方.
居然復活的主就在這裡顯現.
顛覆了保羅自己的心.
他就成為一個傳道人.
第二段聖經約翰福音第二十章.
哪日就是七日的第一日.
晚上門徒所在的地方因怕猶太人.
門都關了.
耶穌來站在當中.
對他們說:願你們平安.

$^{481}$說了這話.
他就把手和肋旁指給他們看.
門徒看見就喜樂了.
耶穌又對他們說:願你們平安.
父怎樣差遣了我.
我也照樣差遣你們.
說了這話.
耶穌就向他們吹一口氣.
祂說:你們受聖靈.
這個描述.
讓我們想起什麼呢?.
應該想起創世記第二章.
上帝創造人的時候.
將靈生命吹入亞當裡面.
使祂成為活人.
這是第一次人類的創造.
現在死亂復活的耶穌基督.
祂就將聖靈.
祂所應許的聖靈.
吹入門徒的生命當中.
耶穌的門徒就成為新的人類.
這才是真正的人類.
這一群人.
信主的一群人.
信奉耶穌能夠領受聖靈的人.
我們就被差遣.
耶穌說:我已經勝過世界.
現在我差遣你進入世界.
宣告耶穌才是王者.
第三段我們看的聖經.
是《馬可福音》.
這裡又是講到耶穌基督.
從死亡復活的技術.
第十六章.
那少年人.
其實是一個天使.
在空墳墓和婦女中顯現.
那少年人對他們說.
不要驚恐.
你們尋找那釘十字架的那撒拿人耶穌.

$^{521}$他已經復活了.
不在這裡.
你們可以安放在他的地方.
你們可以去告訴他的門徒和彼得說.
他在你們以先王加利離去.
在那裡你們要見他.
正如他從前所告訴你們的.
雖然《馬可福音》到了這裡.
就沒有了.
結尾失傳了.
但這裡很清楚告訴我們.
是一個空墳墓.
提到耶穌和門徒曾經有的承諾.
在加利利要再見他們.
加利利就是耶穌第一次呼召門徒的地方.
現在死了還沒有復活.
耶穌將天國帶來.
他要在加利利再見這群門徒.
要重新呼召.
重新拆牌他們.
我們看最後一本的福音書.
《馬太福音》如何記載復活的事.
這裡有一個很熟悉的經文.
復活的主在山上對門徒說.
天上地下所有的權柄都賜給我了.
所以你們要去洗萬民作我的門徒.
奉父子聖靈的名給他們施洗.
凡我所吩咐你們的.
都教訓他們遵守.
我就尚與你們同在直到世界的末了.
這裡很熟悉的大使命.
Great Commission的經文.
可能很多時候我們留意的.
只不過是耶穌吩咐我們所做的事.
要洗萬民作他的門徒.
要洗他們的屍體.
要教訓他們遵守主的話語.
這些都是重要的.
但很多時候我們未必留意第一句.
天上地下所有的權柄都已經賜給我了.

$^{561}$這其實是一個王者的權柄.
一個有神聖權柄的宣告.
你想想在舊約聖經當中.
曾經出現過這一句話的.
大概在哪裡.
有哪一個可以有天上地下所有的權柄.
萬民要歸向他.
想起《但義理書》.
《但義理書》第七章.
有一個這樣的句子.
事實上這裡所有的keywords.
所有重要的詞語.
都在《但義理書》第七章的意象裡面出現.
那裡發生什麼事呢?.
就是在但義理見到一個異象.
那些怪獸從海裡面出來.
有四個怪獸.
代表地上四個帝國的政權.
這四個政權過去之後.
就有人子駕雲而來.
那位得勝者.
上帝賜給他天上地下所有的權柄.
並且萬民要去歸向他.
馬太福音這裡.
耶穌在山上說.
我就是那位人子.
我就是那位得勝世界.
得勝所有權勢的王者.
現在我來猜見你.
進入這個世界.
甚至世界的地極之處.
如果你這樣看的話.
相信耶穌從死裡復活.
並不是單單一件第一世紀發生的歷史事件.
而是與我們今時今日.
你與我有很切身的關係.
你看了這麼多集.
今日我希望你能夠再一次做一個決定.
就是接受這位王者.
成為你生命當中的主.

$^{601}$他的主權.
他所帶來的天國.
神的生命.
要藉著你.
在你的生命當中.
去彰顯復活的主.
有人說得很對.
世人見不到死而復活的耶穌.
是藉著相信死而復活的門徒.
將這位見不到的主彰顯出來.
你與我改變了的生命.
神去到顛覆.
去到重整了我們對世界.
對人生的觀念之後.
你與我作為天國的門徒.
我們最終追求的.
就是要彰顯神的榮耀.
因為只有耶穌才是我們的王者.
每年復節的主日.
在聖母堂Holy Sepulchre這個地方.
都有一個特別的東正教的聖火儀式.
叫Holy Fire.
在這個小教堂的祭司.
會將火.
藉著蠟燭提出來.
傳給衛官的信眾.
Holy Fire 聖火的儀式.
就是象徵著主耶穌基督.
從死裡復活的盼望.
你這樣想也很有意思.
聖母堂的圓頂.
其實設計也告訴我們這一點.
中間的光分成十二條光線.
代表著十二門徒初代的教會.
然後將這個生命.
轉給更多更多的人.
代表著神的生命.
要在更多萬民的身上.
來彰顯出來.
今日你與我.

$^{641}$信耶穌跟隨主的話.
我們就成為祂天國的子民.
我們將祂的榮耀.
在我們生命當中彰顯出來.
我們同心低頭祈禱.
因為你是那位從死裡復活的得勝者.
你是唯一一位的王者.
主讓我們的生命.
被你去重整.
你成為我們唯一的主.
你成為我們唯一的主宰.
唯一取悅的對象.
讓我們的生命.
我們所表達的價值.
我們的行為.
我們的心智.
我們的生活態度.
以你天國.
以你神的生命.
為我們最終追求的目標.
求你藉著我們來見證.
耶穌基督是一位活著的主.
見證你自己的真實.
直到你回來為止.
我們這樣祈禱奉主的名.
求你.
\newpage



\section{2:11-22}
\label{sec:qLV2rgTxxTw}
\textbf{Reconciled in Christ (Ephesians 2\_11-22) - Sam Johnson}
\newline
\newline
連結: \href{https://youtube.com/watch?v=qLV2rgTxxTw}{\texttt{https://youtube.com/watch?v=qLV2rgTxxTw}} ~~~~ 語音日期: 2024-05-28
\newline
\newline
\hyperref[sec:OVsK1efMh70]{\small{< < < PREV SERMON < < <}}
~
\hyperref[sec:index_chronic]{\small{[返順時目]}}
~
\hyperref[sec:index_scriptual]{\small{[返順卷目]}}
~
\hyperref[sec:FBZ55r0dhNA]{\small{> > > NEXT SERMON > > >}}
\newline
\newline
$^{1}$- Good morning, it's great to be here, as always..
Love being at UV, and thank you for allowing me to be here,.
and thanks Jason for leading us in worship..
This morning we're gonna be looking.
in the book of Ephesians in chapter two..
We're gonna look at verses 11 to 22,.
and we'll be talking about being reconciled in Christ..
And I'm here periodically, when I've been here.
we've been kind of slowly going through.
the book of Ephesians, this letter that Paul wrote.
to the church in Ephesus and the surrounding churches..
And last time I was here we looked at the first part.
of chapter two where he's talking about.
the fact that we're made alive in Christ..
We talked about the equality that gives us in Christ,.
that we're all equally made alive in Christ..
And even though we're different in our giftings,.
in our skills, in our talents,.
in the ways that God has gifted us,.
we're different in our life experiences,.
we're different even in the roles that we play.
in the church and in life and in family.
and things like that..
We are all equally given the same redemption,.
the same grace, and that puts us really in a place.
where there's no inferior or superiority in the church..
We are equally forgiven by Christ, equally redeemed,.
equally brought into the family..
And that's how we're to treat one another..
So we talked about that last time..
And today we're gonna be talking about really unity..
That's what he's addressing in verses 11 to 22..
And this passage is just so powerful..
He's talking about the unity of the church.
and how that happens, and it's not through us..
It's an impossible thing to see this happen.
apart from Christ..
And that's really what he's describing for us..
So let's begin reading..
We'll read these verses from 11 to 22,.

$^{41}$and then look at a few things that we can learn here..
It says, "Therefore, remember that at one time,.
"you Gentiles in the flesh, called the uncircumcision.
"by what is called the circumcision,.
"which is made in the flesh by hands,.
"remember that you were at that time separated from Christ,.
"alienated from the commonwealth of Israel.
"and strangers to the covenant of promise,.
"having no hope and without God in the world..
"But now in Christ Jesus,.
"you who once were far off have been brought near.
"by the blood of Christ..
"For he himself is our peace,.
"who has made us both one,.
"and has broken down in his flesh.
"the dividing wall of hostility.
"by abolishing the law of commandments.
"expressed in ordinances that he might create in himself.
"one new man in place of the two,.
"so making peace and might reconcile us both to God.
"in one body through the cross,.
"thereby killing the hostility..
"And he came and preached peace to you who are far off.
"and peace to those who were near,.
"for through him, we both have access.
"in one spirit to the father..
"So then you are no longer strangers and aliens,.
"but you are fellow citizens with the saints.
"and members of the household of God,.
"built on the foundation of the apostles and prophets,.
"Christ Jesus himself being the cornerstone,.
"in whom the whole structure being joined together.
"grows into a holy temple in the Lord..
"In him, you also are being built together.
"into a dwelling place for God by the spirit.".
So he's given us this beautiful picture.
of what the church can be,.
both as individuals, that we are to be this dwelling place.
of God in the spirit, we're to be a temple of the Lord,.
but together as a group,.

$^{81}$that we're to be this dwelling place of God in the spirit,.
this temple in the Lord..
But he's talking about how it happens.
that God takes all these very different people.
and brings them together and makes them one..
Now this is something that our world is really craving..
And I mean, people who don't know Jesus.
are really craving this..
People want to see human beings brought together..
People want to see people reconciled across differences..
That's something our world is craving.
and desires to see so desperately..
But it's also something,.
I think this is a huge gospel opportunity.
because it's something that we see clearly in our world..
It doesn't happen on its own..
It's not a simple thing..
It doesn't take place easily..
And we're told here how this can happen in the church.
by the power of Christ..
And it's such an opportunity for us as we live this out.
to be a testimony of the power of Christ..
Remember Jesus in the book of John,.
he prayed the high priestly prayer.
right before the crucifixion,.
and he prayed that we would be one..
And he said that that was one of the ways.
people would know that we were his followers,.
that that would be one of our testimonies of the faith.
is that we're one..
So I want us to see some things in this passage.
and first what I would love us to consider.
is the fact that apart from Christ,.
there are actual differences that create actual tensions..
That's real and I love that scripture acknowledges this..
So sometimes we think,.
well, we need to have unity in the church.
and we need to get along..
And the way to get along is just to act like.
there's no differences..

$^{121}$There's no problems..
There's no problems here..
Everybody thinks the same and we all agree about everything..
Yeah, let's just make it work, right?.
But the scripture doesn't do that..
The scripture acknowledges and says,.
yeah, people are different..
Cultures are different..
People have very different ideas about a lot of things..
And the scripture doesn't ignore that..
It doesn't gloss over that..
Look there at verse 11..
It's saying, "Therefore, remember that one time.
"you Gentiles in the flesh called the uncircumcision.".
Paul actually points out, he's like,.
the Gentiles are people who are not a part.
of the Jewish heritage..
Now they're believers in Christ..
They're in the church..
But remember that before people kind of called you names..
They were like, you guys are the uncircumcision..
You're the outsiders..
You're not a part..
He's reminding them of that..
And he's saying, who called you that?.
Well, it was the people who called themselves.
the circumcision, the people who said,.
we're the in crowd, you guys are out..
And then in the same verse,.
he makes the point that even the people who said,.
they're the in crowd, they call themselves the circumcision..
He said that was actually done by human hands anyway..
He's saying that that was a boundary.
they set up themselves..
It wasn't something that was really making them.
closer to God..
It was just a boundary they set up themselves..
But he's reminding them, there was this actual division.
between you all..
Some of you felt like you were in..

$^{161}$Some of you felt like you were out..
It was a real tension..
You guys didn't like each other..
And I think that's actually really encouraging to us..
We see tension and divisions in our day..
Some of them are very deep, right?.
But this one is as deep as it gets..
It touches on every area of life.
that they were experiencing in the first century..
The Gentile background believers.
did not understand the Jewish background believers.
and vice versa..
They didn't understand one another..
They were very different culturally..
They had completely different diets..
They had different ways of living..
They had a lot of reasons to mistrust one another..
There was a lot of history,.
a lot of very legitimate grievances in their past..
The Jewish background believers are like,.
hey, you guys have oppressed us.
and you've killed people in our nation..
And so there was a lot of legitimate grievances..
And scripture doesn't ignore that..
It doesn't say, hey, look, guys,.
let's just forget about that.
and let's just kind of sing kumbaya and try to get along..
That's not what this passage is saying..
It's saying, hey, there's real stuff here..
You guys are really different..
And you have a past, you have a history,.
and that's okay, we're not gonna ignore that..
But Jesus is gonna take us through it..
That's what this passage is telling us..
It goes on to say that,.
for he himself is our peace who has made us both one.
and has broken down in his flesh.
the dividing wall of hostility..
Verse 13, I mean, verse 14 there,.
it says there is a dividing wall of hostility..

$^{201}$It exists, it's real..
But Jesus is gonna take us through that..
And that's so good for us to know..
There are a lot of differences in the church,.
and that's okay..
God is the one who created the nucleus.
of a lot of those things..
Obviously, we have extended them sometimes.
in unhealthy ways,.
but God's the one who created us in his image,.
and a part of being created in his image.
means that we are creative,.
and we like to build things,.
and we like to do things..
And in Genesis 1:27,.
God gave what theologians call the cultural mandate,.
where he said to the first human beings,.
"Go out into the earth and subdue it, do your thing.".
And ever since then, humans have been doing that, right?.
We've been going into the earth,.
and we've been building cities,.
and creating things, and making art,.
and creating culture, and making music,.
and all these different things that we do.
that God actually created us to do, right?.
But different people in different places.
do that very differently..
So cities in different parts of the world look differently,.
and our culture looks differently,.
our languages look differently..
And so then, as we do that,.
that becomes, at times, a dividing wall..
And again, scripture is not saying,.
"Well, that shouldn't be there..
"It's a problem for you guys to have different cultures.".
It's not saying that at all..
Our cultures reflect the beauty and the majesty of God,.
because God's a creative God,.
and he made us with all these cultural differences..
And when we come into the church,.

$^{241}$it's not that we're supposed to say,.
"Well, I have to leave my culture at the door,.
"and I have to leave all these things outside,.
"and somehow try to be different.
"when I'm in the body of Christ.".
No, God says, "Bring these things into the body of Christ,.
"but I'm also going to make you a new humanity.
"through Jesus.".
And I think that's really important for us to know..
There are issues, there are genuine tensions,.
and at times, even in the church,.
there are genuine, legitimate grievances from the past..
And we shouldn't ignore that,.
we shouldn't try to whitewash it or gloss it over..
That's not what the Bible does..
The Bible says that God is going to do something.
to reconcile us..
And that really starts with the fact that in Christ,.
we can all be reconciled to God..
Anyone can be reconciled to God,.
and everyone can be reconciled to God..
These verses make it clear, verses 12 to 13,.
"Remember that you were at that time separated from Christ,.
"alienated from the commonwealth of Israel,.
"and strangers to the covenants of the promise,.
"having no hope and without God in the world.".
He's talking to the Gentile background believers..
He's saying, "Don't forget, there was a time.
"when you guys didn't know anything about the word,.
"you didn't know anything about the truth,.
"and you were completely separated from God.".
So, Paul is reminding us that, hey, we're all sinful,.
and a lot of us used to be in that state..
But then, the gospel, grace, verse 13,.
"But now, in Christ Jesus, you who were once far off.
"have been brought near.".
Yeah, in the past, you were completely separated..
You didn't have hope, you were alienated.
from the truth of God..
But now, you have been brought from far off to near..

$^{281}$You've been brought close to God..
How did that happen?.
By the blood of Christ..
So, what we see here is that all of us,.
even though we're all sinful, and we're all broken,.
we can all be redeemed in Christ,.
we can all be forgiven through Christ,.
and that all means all, regardless of our background,.
our life experience, our culture, wherever we have been,.
or whatever we've gone through,.
we can be reconciled to God through Christ..
We can be brought near by the blood of Christ..
And what we have to remember is that,.
as we're reconciled to God in Christ,.
then we become reconciled to one another..
He says that, "He himself is our peace,.
"is our peace, that he might create in himself.
"one new man in place of the two,.
"reconcile us both to God in one body through the cross,.
"thereby killing the hostility.".
He himself is our peace..
Jesus himself is our peace..
And we are reconciled to one another through the cross..
This Greek word here for one new man,.
anthropos means human being without reference to gender..
And it's saying that there's a new type of human being,.
a new identity, a new way to be human..
It's not based in Adam, but it's based in Christ,.
saying God is creating a new humanity.
and a new community through the cross..
It's what this verse is telling us here..
"He himself is our peace,.
"that he might create in himself one new man.
"in place of the two,.
"reconcile us both to God in one body through the cross,.
"thereby killing the hostility.".
So as we begin to individually be reconciled to God.
through Christ and the cross,.
God starts to use that to reconcile us to one another..
And even people that are different than us.

$^{321}$or people that we may struggle to understand.
or to relate to, or maybe just something.
in their personality that really kind of grates against us,.
and it's hard for us to understand.
what's going on in their life,.
God begins to reconcile us to one another in Christ..
And just like other areas of our sanctification,.
and sanctification is just a church word.
that means we're becoming more like Jesus..
After we come to know Jesus, we're growing,.
and it's a process of becoming more like Jesus..
And just like other areas of that in life,.
this is a process over time,.
being reconciled to others in the church.
is a process over time..
So we need to be patient,.
patient and exercise grace with one another..
We need to understand that it takes us time..
God has to work on our heart..
God has to do that to me all the time..
God has to work on my heart.
and show me areas where I'm being proud.
and where I'm not understanding someone else.
because of my own pride or my own stubbornness,.
or the fact that I'm like,.
"This is the way I like to do it,.
"this is what I'm comfortable with,.
"this is the way we should do it.".
Sometimes it takes time for God to work on that.
and chisel it away..
So we need to be patient.
and give grace to one another in the church..
But at the same time, we need to expect.
that this kind of reconciliation.
would be reflected in the church..
It should be, that we would be reconciled.
to one another in Christ..
And if we're really serious about following Jesus,.
then we should expect that to be reflected in our lives.
and in our churches..

$^{361}$We should expect to see that..
And that reconciliation to one another.
is really a holistic thing..
Look at verses 18 and 20 here, it says,.
"For through him, we both have access.
"in one spirit to the Father..
"So then you are no longer strangers and aliens,.
"but you are fellow citizens with the saints.
"and members of the household of God,.
"built on the foundation of the apostles and prophets,.
"Christ Jesus himself being the cornerstone.".
So these verses talk about our spiritual condition,.
our physical condition, and our emotional condition..
All of that is in these verses..
And so being reconciled to Christ,.
to one another through Christ,.
is really a holistic thing..
Verse 18 says, "For through him,.
"we both have access in one spirit to the Father.".
When he says we both, he's talking about,.
he's talking about the two different groups.
that he's mentioned..
He's talking about the Gentile background believers.
and the Jewish background believers..
He's saying, remember, through Christ,.
both of us have access in one spirit to the Father..
I also like how in verse 18 and throughout this passage,.
Paul is speaking in the first person..
He's not saying like,.
"Yeah, you guys have this problem with being reconciled.".
He's saying, "Hey, I'm a part of it too..
"I'm in here too.".
That's why he says we both..
Because he was a Jewish background believer..
And remember, before he was saved,.
he had been, he was serious, right?.
He was really serious about the Jewish faith,.
so much so that he was persecuting Christians..
And before he was saved, he was the type of guy.
that definitely would not have hung out with the Gentiles..

$^{401}$Because it was such a deep part of what he believed..
And a deep part of kind of the legalism.
that he was practicing..
So he includes himself..
He says we both have access in one spirit to the Father..
That shows us the spiritual part of this reconciliation..
That he's saying, as a result of reconciliation in Christ,.
we all have equal spiritual access to God..
There's a spiritual reconciliation that takes place.
where we're all put in one family, right?.
And most of us are pretty comfortable with that..
We're like, "Okay, I get it..
"The spiritual reconciliation, I can handle that.".
But then he goes on, make it a little more personal.
in verse 19, he says, "So then you are no longer.
"strangers and aliens, but you are fellow citizens.
"with the saints and members of the household of God.".
Now he's saying that there's this actual physical.
and emotional proximity and relationship with one another.
that we come into..
So the difference in the language there..
Before we were strangers and aliens,.
people who don't spend time together..
That's what strangers are..
Aliens, that word means like foreigners..
We were people who are like, "They're different than me..
"They are the other, we don't spend time together.".
But now you're fellow citizens..
Fellow citizens means, "Hey, we're on the same team..
"We are together.".
And so he's giving us this picture.
of a physical change in our proximity to one another..
That we go from being strangers, we don't hang out,.
we don't spend time together, to now, through Christ,.
we are members of the same household..
Sounds like family, right?.
And that's what he's saying..
You guys used to be strangers,.
you wouldn't have talked to each other..
Now, you're family..

$^{441}$That means we've come into physical and emotional proximity.
with one another..
Now that's a little bit less comfortable, right?.
It's easy to say, "Well, I know other people in the church,.
"they might be different than me,.
"but I'm supposed to be reconciled to Christ in them,.
"so I can be spiritually reconciled..
"I know spiritually we're shoulder to shoulder.".
That's easy..
But when you say, "Oh, wait, I'm actually supposed.
"to be physically reconciled to them,.
"and there's points where I'm supposed to be actually.
"in physical proximity with them and in fellowship,.
"and we're supposed to be in the same room.".
That's hard, 'cause we say, "Wait a second,.
"these people are different than me.".
That's why it's hard to get in the same room..
That's why it's hard to be reconciled in that way,.
to be in fellowship..
That's why it's hard to be in emotional proximity..
It's hard to understand their experience.
and walk through something with them,.
because we're different..
But that's what this passage is telling us..
This means that this reconciliation.
that God has done through Christ.
should impact our physical proximity.
with others in the church,.
and especially, I think, those who are different than us..
That means it impacts our emotional openness.
and our bonds with one another,.
our fellowship together,.
and our service to the world around us..
All areas of our life should reflect this reconciliation.
that Jesus is affecting in us..
Built on the foundation of the apostles and prophets,.
Christ Jesus himself being the cornerstone..
He's the foundation of what's happening in the church..
Now, so far, as we look at these things,.
we start to say, "Okay, Paul,.

$^{481}$"this is good stuff,.
"but you don't really understand the world we live in.".
Right?.
I mean, people disagree about a lot of stuff,.
and we love to talk about unity..
We love to call things united this and that..
Like, even the United States,.
not as united as we would like to think it is, right?.
We love to talk about unity,.
but it just doesn't happen..
It's too hard..
This picture here, we read it, and we're like,.
"Okay, yeah, we're all made one new man in Christ..
"That's cool, but that's not real," right?.
And we have to understand.
that this reconciliation through Christ.
is impossible apart from Christ,.
and that's what he's saying at the end of the passage here..
"Built on the foundation of the apostles and prophets,.
"Christ himself being the cornerstone.".
Christ is the cornerstone, the foundation of this unity..
It says, "In whom the whole structure,.
"the whole structure of the church,.
"being joined together,.
"grows into a holy temple in the Lord.".
But none of that happens..
The joining does not happen without the cornerstone..
In him, you also are being built together.
into a dwelling place for God by the Spirit..
That doesn't happen without the cornerstone..
This reconciliation that we're talking about.
is impossible apart from Christ..
We don't do this on our own..
We can't reconcile ourselves to others..
We don't have the power to do that..
Only God can do this..
Only God can create a new humanity..
And I want us to think real quickly.
about some of the things we as humans do.
to try to fill this void,.

$^{521}$because all of us understand,.
even people who don't know Jesus,.
understand in their hearts that it's wrong,.
that humans experiencing division and hatred.
towards one another is not how we were supposed to be..
It's not what God created us for..
Even people who are not believers,.
they don't believe in God, they understand that..
They say, "Wait, these things are wrong..
"Racism and discrimination.
"and all these hostilities and divisions, they're wrong.".
People get that..
And honestly, that's good..
It's good that as humans, we see that..
And as humans, we're good at kind of diagnostics.
and naming the problems..
That's something we do very well..
But the solutions are a lot harder..
And so when we think about outside the church,.
there are things that people say,.
"Okay, well, let's try to fill this gap..
"Let's try to fix this in different ways.".
And sometimes we're tempted to do these things.
in the church itself..
But those solutions don't live up to the standard.
of what God calls us to, right?.
So one of the first things we see a lot.
outside the church and within the church sometimes.
that's just a human tendency,.
when we think about the divisions of who we are as people,.
is we fall into forms of tribalism..
We say, "Well, I get along with the people.
"that are like me..
"Those people are different..
"It's really hard to understand them.
"or get along with them..
"So we'll just do our thing with our group..
"They can do their thing with their group.
"and that'll work," right?.
Us human beings have done that since the beginning..

$^{561}$And there's a lot of kind of logical seeming reasons.
that make sense there, right?.
But that's not what God called us to in the church..
God called us in the church to love one another,.
even when it's difficult,.
even when we don't understand one another,.
even when we have to cross a barrier to love one another..
So that ultimately doesn't really love others..
It just seeks my own interests and the people like me..
It falls short of God's standards..
Another one that we see a lot outside the church.
is some form of tolerance, where people say,.
"Well, I understand that people are different..
"I really respect that,.
"but, and I'm gonna tolerate the fact they're different..
"We're gonna be in the same room.
"and we're gonna be there,.
"but I'm still not really loving someone..
"I'm just kind of tolerating the fact.
"that they're different than me, and I'm okay with that,.
"but ultimately, I'm not really loving them.
"the way God calls me to.".
Another thing we see is kind of a form of self-exoneration,.
where we say, "Well, I know there's problems,.
"but I'm not that type of person, right?.
"I really care about other people,.
"and I'm not involved in the division and the hostility,.
"so that's not me..
"I'm not a part of that.".
Again, that falls short of God's standard..
What has God called us to do?.
God has called us to genuinely love one another.
and to genuinely be reconciled..
Not just to say, "Well, it's easier not to interact.
"'cause we're different," or not just to say,.
"Well, we'll kind of interact,.
"but I'm just gonna tolerate someone else,.
"even though I don't really agree with them.
"or understand them," and not just to say,.
"Well, I'm a good person,.

$^{601}$"and so I'm not someone who would experience those tensions..
"I'm above all that," and to act like.
we're not impacted by these things..
No, God's calling us to say,.
"Hey, these things are real..
"There are challenges in the church,".
and not just the cultural tension,.
but just different ways of thinking..
Like, "This person thinks differently.
"about ministry than I do,".
or, "This person thinks differently.
"about worship than I do and how it looks.".
Those types of very real, practical things,.
God calls us to say, "Hey, these things are real,.
"and I'm supposed to love this person across this barrier.".
And yeah, there may be some right and wrong involved,.
and maybe we need to really focus on the scripture.
and find out what God has actually called us to do..
I'm not saying we set aside the truth.
or set aside issues of right or wrong,.
'cause we have to work through that in the church, right?.
But even as we do that, we're still genuinely loving people.
across those barriers of difference..
We're not just saying, "Well, I think you're wrong,.
"but I'm gonna tolerate you.".
We're saying, "Hey, I love you,.
"and I want you to know Christ more.".
And maybe you're misunderstanding.
or misinterpreting this part of the Bible,.
but as we walk through this together.
and we move towards the truth,.
I want to love you the entire time,.
and I wanna have my arms around you..
That's what Jesus is calling us to in the church,.
actual love,.
love and respect for God's image in me.
and what God's done in my life,.
love and respect for God's image in you.
and what he's done in your life,.
love and respect and amazement at the redemption of Christ.

$^{641}$and the grace that God has poured into me,.
and also with the redemption and grace.
that God has poured into you..
And love and respect and amazement and appreciation.
for the story that God has built into my life.
and all of the experiences.
and the things that have led me to who I am,.
and love and amazement and respect and admiration.
and appreciation for the story.
that God has built into your life.
and all the experiences and all the things.
that have led you to where you are..
Again, it doesn't mean that we're not serious.
about right and wrong.
or that we have to set those things aside..
It simply means that we're really supposed.
to love one another and be reconciled..
And that's what unity looks like in the church..
Unity in the church is not uniformity..
It's not all of us being the same..
Back to what I said at the beginning,.
that's not what God wants..
God created us with all this beauty and diversity,.
and he did that on purpose..
But unity in the church actually assumes diversity..
It assumes that there's a lot of differences,.
and yet we still love one another..
That's impossible..
How could we do that?.
Through Christ, and that's what Paul's saying here..
He's saying this is only gonna happen.
as Christ brings you together.
and makes a new human out of you..
You're not gonna lose your background,.
you're not gonna lose your experiences,.
you're not gonna lose your culture..
You're gonna bring that in with you..
But now you're gonna have this new basis,.
this new foundation of fellowship.
that you share in common with other believers..

$^{681}$And it's difficult..
I know it is..
I know you guys feel that tension..
Anybody who does ministry in a diverse city like this,.
it's difficult and you feel that tension..
I certainly feel it in my neighborhood..
You get a group of people in the room,.
even if they're all believers,.
they all have different ideas, different experiences,.
different culture, different preferences,.
different things that they enjoy that are important to them..
It's hard to bring everyone together.
and to move in the right direction..
But what this passage is telling us.
is that God can supernaturally do this in our hearts..
And I wanna go back to where we started.
and remind us that the world is craving this..
I would say, and I think there's plenty of evidence.
of this in our world, that this is so impossible.
that the world can't do it..
Apart from Christ, it's just not gonna happen..
People try to make it happen in a lot of ways..
And people try to bring people together, and that's good..
That's wonderful..
We should seek the peace of our cities.
and of our nations and of our world..
So that's good..
But there's example after example in our world.
of the fact that we can't make it happen.
and we can't sustain it apart from Christ..
People just can't do it..
As hard as they try, we can't cross those barriers.
long-term..
We might get people together for a little bit,.
but we can't hold it together..
And the world is craving that..
And I think it's such a gospel opportunity.
for us in the church if we're able to say,.
hey, look at what God is doing here..
It's not us..

$^{721}$We're a group of very different people.
with very different ideas..
We have very different life experiences..
We have different cultures..
We have different languages..
We have different foods and different preferences.
and so many things..
And yet we come into the same room.
and we don't just tolerate each other..
We actually love each other and we love to be here..
It's one of the reasons I love being at UV.
and love being with you all..
I hope you guys love me back..
But that's why..
Because it's such a beautiful thing.
when the people of God come together.
and say we have this foundation of fellowship,.
actual, real, genuine fellowship..
We're not faking it..
We're not pretending to like each other..
We really love one another..
And I really love you guys..
I get excited every time I come in the door.
and see all the people here and see what God is doing here..
And that's genuine..
I love UV..
I love this church..
And it's not because everything here.
is exactly like the churches that I grew up in..
A lot of it's very different..
It's because of Jesus..
We have this foundation in Christ.
that brings us together and helps us to love one another.
and makes us into this new humanity..
And I think that's such a beautiful and powerful vision.
for us to live into in the church..
And like I said, it's a process..
It doesn't happen immediately..
We have to give grace and be patient with one another..
But that's where we should be headed..

$^{761}$And that's what we should pray for..
That's what we should desire..
That's what we should move towards..
That's what we should let God work in our hearts.
to bring that about and be open and sensitive.
and soft-hearted before God and say,.
God, I know I'm uncomfortable right now..
This is different for me,.
but I'm gonna let you pull me in this direction.
because I know this is your work to reconcile us in Christ..
And as we do that, we start to have this testimony.
that the world can't refute..
And the world looks at us in the church and says,.
you know what, I don't always like what you guys say.
about Jesus and sin and like the whole,.
Jesus is the only way to heaven..
I don't know about that,.
but you guys actually love each other..
Like, you really like being together..
And I see you doing that across barriers of culture.
and ethnicity and race and different personalities.
and different giftings and different ideas..
That becomes a really powerful testimony.
to the world around us, that Jesus is making us one..
And just like Jesus prayed,.
it becomes something that people start to say,.
I think those people are really true followers..
They're true disciples..
They're truly believers in Christ.
because that's defining who they are,.
that they're coming together..
They're reconciling Christ across barriers.
of actual difference and actual tension..
So that's my prayer for my life..
And that's my prayer for us as we move forward.
in our respective communities, following Jesus together,.
is that we will grow more and more into this reconciliation,.
that God will be more at work in our lives,.
and it'll be more and more evident to the people around us,.
that they'll see us reflecting this in our communities.

$^{801}$and reflecting the gospel and the beauty.
of Jesus making one new humanity.
out of very different things that existed before..
Let's pray..
God, I thank you for this miraculous picture.
that you give us in this passage.
of what it looks like to be reconciled in Christ..
And God, I confess in my own life.
how impossible this seems at so many different times..
It's not impossible for us to pretend like we want peace.
or to pretend like we want to work together.
or pretend like we want to get along..
But it's really hard to love people that are different.
when we may not understand all the things.
that make us different..
That really is hard..
And what's way harder is to then get us all.
to head in the same direction..
But God, you give us this miraculous picture in your word.
of the fact that we can be made.
into one new humanity through Christ..
And so that's our prayer, God..
We ask that our hearts would be open to you,.
that you would do this, that your kingdom would come.
in this world through us, your people,.
and that you would make us one in Christ.
so that others would know that we're your disciples, God..
In Jesus' name, amen..
\newpage



\section{腓利門書}
\label{sec:FBZ55r0dhNA}
\textbf{升呢: 福音重新定義人與人之間的關係 (腓利門書) - 譚文倩莊傳道}
\newline
\newline
連結: \href{https://youtube.com/watch?v=FBZ55r0dhNA}{\texttt{https://youtube.com/watch?v=FBZ55r0dhNA}} ~~~~ 語音日期: 2024-05-28
\newline
\newline
\hyperref[sec:qLV2rgTxxTw]{\small{< < < PREV SERMON < < <}}
~
\hyperref[sec:index_chronic]{\small{[返順時目]}}
~
\hyperref[sec:index_scriptual]{\small{[返順卷目]}}
~
\hyperref[sec:fTMtbscgTAE]{\small{> > > NEXT SERMON > > >}}
\newline
\newline
$^{1}$大家早安.
我小學未畢業.
就整家移民到美國芝加哥.
來到芝加哥之後.
我姑媽的朋友.
鼓勵我們到教會.
整家一起學英文.
沒多久.
因為學英文的關係.
就有人邀請我們去教會.
40年之後的今天.
草草埋埋.
我身邊有一班.
跟我一起成長的弟兄姊妹.
這班人仍然與我很親近.
今時今日仍然是我的禱告夥伴.
除了弟兄姊妹.
多年來.
神也擺了很多.
弟兄姊妹做我們的導師.
在不同團契.
不同階段的時候.
他們都願意承諾他們的時間.
心力.
他們的人生歷練當然比我多.
屬靈生命也比我們成熟.
進入港島之前.
我很想藉著這個機會.
去多謝這些導師.
特別是我們中學時代.
我們那班人有很多掙扎.
經常折磨我們的導師.
今時今日.
我成為一個傳道人.
我覺得他們絕對有他們的功勞.
眾多導師裡面.
其中一個給我很大的影響.
其中一個影響.
就是他偶爾會給我一顆糖.
給我幾塊餅吃.

$^{41}$這些餅和糖.
很多時候是別人去探望.
探望他之後吃剩的.
但是他吃的時候.
他覺得這些糖或餅很有趣.
他覺得我會喜歡吃.
他就會留給我.
留給我.
等我去探望他的時候.
他就會跟我分享.
多年來的小動作.
教會我原來愛人.
不一定要做什麼大動作.
甚至我在想.
多年來.
他極少用嘴巴去教導我.
但是他的身教.
他的待人接物.
他怎樣管理金錢時間.
對我有很大的影響.
他不是用嘴巴去教導.
他是用行動.
在生命裡面去活出他的信仰.
《肥理門書》裡面.
同樣有一個人.
正正是用他的生命.
去活出福音.
影響身邊的人.
這次回來的時候.
我預備了三堂的功課.
跟大家一起去思考.
第一課今天就是說.
福音如何重新定義.
人與人之間的關係.
第二堂是說到.
教會的群體的功課.
第三課就是說.
《肥理門書》裡面.
一個字都沒有提及耶穌.
究竟耶穌去了哪裡呢.

$^{81}$不如我們先一起低頭祈禱.
出入了我們的神.
祈求你祝福你自己的姓名.
祝福你的話語.
祈求你能夠在祂的身邊.
我小時候住在香港的時候.
我爸爸是行選的.
所以很久後我就會寫信給我爸爸.
還記得有沒有人用過郵簡.
還記得當年有一樣東西叫郵簡.
寫完郵簡貼上郵票.
然後就會將這封信寄去一個老遠的國家.
就是我爸爸的船去到哪個國家的時候.
就是寄去哪個國家的時候.
同樣保羅寫的這封信.
也是由老遠的羅馬帶去哥羅西.
我想問一下.
我們裡面有沒有人知道.
羅馬離開哥羅西有多遠.
有沒有人知道.
我查過原來是1645公里.
或者是2647公里.
我們教會裡面有很多長跑好手.
我想問一下大家.
你猜一下1645公里.
是多少個馬拉松.
有沒有人很快就計算出來.
大約是63個馬拉松.
再試一下用另一個方法.
去消化這個距離.
我嘗試在家裡Google Map.
向著芝加哥這個方向進發.
1640公里.
我看看自己會去到哪裡.
原來大約會去到.
Calvin和Jasmine Frankfurt的家.
我查了一下.
原來我不吃.
不去廁所.
23個小時左右.

$^{121}$就可以開車去到家裡.
一封很遠帶過來的求情信.
日進五節.
我聽見你對眾聖徒的愛.
和對主耶穌的信心.
雖然保羅和肥利門.
分隔在兩個很遠的地方.
但很明顯自旅行間.
保羅一直很關心他和教會的狀況.
其實你也看到.
保羅就好像一個牧師.
但保羅在這封書信裡.
他沒有用一個像是.
一個屬靈長者.
一個權威的身分.
去跟肥利門說話.
反而一開波.
保羅就說.
我是保羅,為基督耶穌做囚徒.
其實整章經文裡的第9,第10,13,23節.
其實都有類似的形容.
通常稱呼是我們人的身分象徵.
其實有一點像是你指揮對方.
怎樣去對待你.
但你看的時候覺得挺有趣.
保羅在這封求情信裡.
不是一開波.
給自己一個挺高的位置.
反而是說自己是一個監等.
在邏輯上去想的時候.
好像沒有什麼智慧.
當然傻了.
其實保羅是很有智慧.
一章8至10節.
因此我上了這年紀.
現在又為基督耶穌.
作囚徒的保羅.
雖然在基督裡.
可以放膽吩咐你做該做的事.
但是為了愛的緣故.

$^{161}$我寧願懇求你.
就是為我所囚時所生的孩子.
阿彌西柔懇求你.
其實保羅是很有智慧的.
因為保羅在這裡.
要說得很清清楚楚.
他不是一個犯了法的監等.
而他是一個愛主.
因為福音緣故被囚的監等.
而他是一個愛主因為福音緣故被囚的監等.
保羅因為顧著福音大國的原因.
他不用一個高壓的聲音.
去使菲利蒙受煩.
他是用愛出發.
把自己放在一個很卑微的位置.
這是極有智慧的.
因為其實他在做這個動作的時候.
他同時在預備菲利蒙.
同樣要因為福音的緣故.
去謙卑他自己.
去接納奴隸.
奴隸亞尼西某.
他不是用嘴巴去說.
他是用他自己的生命.
去活出信仰.
思想保羅的智慧.
我自己在想作為這個傳道人.
我真的覺得自己也要學習一下.
因為保羅在這封書信裡面.
表面上他只是一封求情信.
但其實他同時在活出信仰.
回應福音.
以身作則.
他這個動作.
其實也在牧養菲利蒙.
去成就福音的事公.
第二個關係.
保羅和亞尼西某.
屬靈父子的關係.
菲利蒙書這封求情信.

$^{201}$當中其實沒有清楚交代一些細節.
在25節裡面.
但可以推斷內容似乎是.
一個奴隸.
叫做亞尼西某.
冒犯了他的主人.
菲利蒙.
普遍的推測.
是說他應該是偷了錢.
或者他做了什麼.
得罪了他的主人.
然後他就逃亡.
按照當時的羅馬律法.
如果有人協助這個逃亡的奴隸.
他會賠償主人的損失.
另一個假設.
亞尼西某被他的主人.
被差派到羅馬.
但不知什麼原因.
他逗留的時間過長.
無論如何.
這封信的主要目的.
是保羅懇求菲利蒙.
不計前嫌重新接納亞尼西某.
11至14節.
他從前對你沒有用處.
但現在對你對我.
都很有用處.
我現在把他本人送回給你.
這就是說.
把我的心肝送回給你.
我本來想把他.
留在自己身邊.
好讓他在我為福音被囚時.
代替你服侍我.
可是我沒有得到你的同意.
就什麼都不願意做.
免得你這善行.
是出於勉強.
而不是甘心情願.

$^{241}$以前羅馬的人口.
大約有20\%.
是奴隸.
每5個人當中就有1個.
羅馬政府.
承認奴隸.
天生是一個人.
但在法定上.
他只是主人的財產.
他們不單沒有自由.
連他的命.
他自己都沒有自主權.
如果奴隸犯錯.
我看一些.
不同的《釋經》書.
說到可以在廚房拿油.
打瀉.
激怒他的主人.
就可以割掉.
奴隸.
亦都沒有權擁有自己的家庭.
如果奴隸生了小朋友.
小朋友.
都是屬於他的主人.
主人有權再賣掉.
或者做小奴隸.
最慘的是.
女人.
小女孩.
或者小朋友.
他們很多時候.
都是主人的承諾對象.
十二節.
我的心肝.
當然保羅很清楚知道.
亞歷西茅是一個奴隸.
但因為福音的關係.
因為他信了耶穌的關係.
他是什麼身份.
其實已經不再重要.

$^{281}$因為保羅.
更加著重主耶穌基督.
在這個人身上.
所成就的救恩.
在神的家裡.
無論你有錢或沒有錢.
你有地位或沒有地位.
都不重要.
因為在主裡我們人人平等.
領受同一個救恩.
保羅帶領亞歷西茅.
信主.
牧養他.
在保羅眼中.
亞歷西茅.
不再是奴隸.
卻是主裡的弟兄.
是福音的夥伴.
更加是他的心肝.
第三個關係.
費利門和亞歷西茅.
主人和奴隸的關係.
十五節到六節.
其實他暫時離開你.
也許正是為了養你.
永遠得到他.
他不再是奴隸.
是高過奴隸.
是親愛的弟兄.
對我尤其是這樣.
更何況對你.
無論是在肉身上.
還是在主裡.
都是這樣.
其實想想.
保羅也頗拿費利門命.
因為他說的.
不只是你去接納.
虧欠你的奴隸.
他還要告訴他.

$^{321}$這個奴隸.
他已經信了主.
今時今日.
你們不再是主人和僕人的關係.
而是親愛的弟兄.
你們是親愛的弟兄的關係.
十七至二十一節.
所以你要是把我當作夥伴.
就要接納他.
像接納我一樣.
他若使你損失了甚麼.
或欠你甚麼.
都算在我帳上吧.
我保羅親手寫下.
我必償還.
用不著我說.
其實你.
連你自己也是欠我的.
所以弟兄.
讓我在主裡得到你的幫助吧.
使我的心在基督裡得到歡慰.
我深信你的順從.
才寫信給你.
我一知道你所做的.
甚至會超過我所說的.
赦免不是一個義務.
更何況奴隸根本沒有資格提出這個請求.
所以他需要請一個.
跟他主人有一定的地位的人.
為他提出.
這個中間人.
還要很有智慧.
去提出這個請求.
因為他不是要挑戰主人的權威.
所以大家不妨再想想.
肥理門是一個有身份.
有地位.
有財富.
受人尊敬的基督徒.
他更加是.

$^{361}$哥羅西教會的領袖.
相對.
亞歷西茅是一個.
沒有身份,沒有地位.
沒有自由的奴隸.
甚至他只是肥理門家裡.
眾多奴隸.
之中的其中一個.
在這封信裡.
保羅稱呼肥理門是.
同工夥伴.
謙虛的弟兄.
而亞歷西茅.
是很親切的.
是孩子心肝.
親愛的弟兄.
想想.
順了主之後.
亞歷西茅是升了等級.
保羅堅信.
教會是由上帝透過基督耶穌.
救贖的工作所認證.
領受了救恩的人.
在這個家裡就是弟兄姊妹.
不分你我.
地位,學歷,膚色.
財富的多少.
因為在神的家裡.
這些一點都不重要.
因為福音的關係.
因為主耶穌基督的關係.
這三個人.
現在有一個common denominator.
就是基督徒.
現在有一個common denominator.
有一個共同的元素.
一個影響生命的元素.
福音.
再想深一層.
肥理門接納.

$^{401}$這個升了的亞歷西茅.
其實一點都不簡單.
我讀神學院的時候.
其中一個新約教授.
是一個牧師.
他說他自己從小到大都很乖.
每個星期都回教會.
有一班朋友一起在教會長大.
但到了中學的時候.
慢慢有些同學離開了.
其中一個.
出去了之後.
還經常犯事,吸毒,濫交.
十多年之後.
這個同學又回到教會.
還洗心革面.
悔改.
接受主耶穌基督成為他的個人教主.
我的教授說.
最初他很開心.
看到舊朋友回到教會.
但慢慢地.
他就越來越.
總是不知所措.
心裡就覺得很不滿.
他說因為.
他從小到大都很乖.
每個星期都回主教學.
去崇拜,去事奉.
過著一些很敬虔.
很乖的生活.
在人前人後.
有很好的見證.
他說相對我這個朋友.
外面的生活烏煙瘴氣.
為什麼他無端回來教會.
大家好像當他.
以前的事情好像大家忘記了.
經常都歡迎他.
很歡迎他回來.

$^{441}$那個注意力轉移到他身上.
他就覺得很不公平.
當他這樣想的時候.
突然之間.
他就想起《路加福音》.
十五章十一節.
還有三十二節.
狼子的比喻.
如果你沒聽過這個故事.
我嘗試簡單快速地說出來.
男主角只有三個.
老爸,大兒子和小兒子.
小兒子等不及老爸去世.
就要求他分家產.
分了家產之後.
就去了花天酒地.
花了一個仙都沒有.
想著去找工作又越著饑荒.
當他走投無路.
只要吃豬哨的時候.
他突然想起.
雖然我分了老爸一半家產.
但老爸還有錢.
不如我回去找我老爸.
他給我工作也好.
做牛做馬也好.
至少我有飯吃.
不用做家庭主婦.
小兒子想不到.
當他回家的時候.
還沒到門口.
老爸已經在迎接他.
很開心.
看到小兒子就在說.
你回來了,真是太高興了.
他馬上跟他的傭人說.
快點來,快點來.
快點拿新衣,新鞋給我兒子穿.
然後再另外吩咐一些傭人.
快點去宰雞殺鴨.

$^{481}$燒一隻燒豬.
今晚我要慶祝.
讓全世界知道.
我失而復得的兒子.
回來了.
故事還沒完,大哥剛下班.
在門口看到很熱鬧.
他就捉住其中一個傭人問.
發生什麼事?.
怎麼這麼吵?.
傭人就告訴他.
你弟弟回來了.
你爸爸很開心.
今晚我們開派對慶祝.
大哥知道之後.
當然不是很開心.
他就覺得很生氣.
就找他爸爸理論.
他在想,有沒有搞錯?.
爸爸,我這麼多年.
對家裡鞠躬盡瘁.
你連個生日會.
都沒跟我開過.
現在我這樣的弟弟.
回來,你罵都沒罵他.
你還說今晚開派對.
你怎麼做人家的爸爸?.
爸爸跟大兒子說.
他說,兒子.
我一直都跟你在一起.
其實家裡的一切都屬於你.
現在你這個失而復得的弟弟回來.
是死而復活.
失而又得.
你為什麼會不開心?.
聽到這個浪子比喻的時候.
普遍的人只留意.
兒子有多敗家.
卻忽略了大兒子.
其實也不是一個好人.

$^{521}$耶穌基督的擺上.
不僅驅使保羅和肥利門謙卑.
叫阿彌西某聖叻.
其實想想他們三個.
也因為主的關係.
聖叻成為天父的兒女.
主裡面是一家人.
而且這個身份.
也影響他們的言行.
福音重新定義.
人與人之間的關係.
之前我提到.
我小時候爸爸在行選.
所以我會寫信給他.
聽起來很溫馨.
但其實一直這麼多年來.
我和我爸爸的關係都很差.
想當年移民來美國之後.
我們一家大小.
才正式第一次住在一起.
我爸爸來到之後.
他很不適應.
也找不到一份他喜歡做的工作.
和我們生活習慣也很不同.
結果我們家裡經常家吵屋閉.
加上我們家裡生活拮据.
我爸爸很多時候一不開心.
就會將他的煩惱發洩在我們身上.
我大哥有時候不太機靈.
他打得最多.
我弟弟就最聰明.
他經常裝作不知所措.
所以打得最少.
我不忿氣,拼命打.
因為我經常覺得我沒有錯.
為何要打我.
我越問得多,就越挑戰他的權威.
雖然我沒有我大哥被打得那麼多.
但他會很生氣.
回想起來.

$^{561}$我是一個女生.
但我爸爸瘋了起來.
他從來沒有因為我是一個女生而留手.
很多次我被他打得流血.
多年來都沒有改變.
甚至沒有變本加厲.
只不過我們三個慢慢長大.
離開了家,有了自己的家庭.
想當年我還不是一個傳道人.
但我卻是一個對信仰很認真的基督徒.
所以我知道神要我和爸爸和好.
我需要和爸爸傳福音.
但老實說,我真的很不想.
我也覺得很難.
或者終於我鼓起勇氣和他傳福音.
通常被他罵得一臉屁.
他當然沒有相信.
但其實我心裡是有些心涼.
因為我真的覺得.
反正你也不值得.
你這個人,福音不適合你.
你不相信就罷了.
但我跟神說,我交足功課.
我做了你要我做的事.
是他不相信的.
多年來我跟神祈禱.
我經常跟神說,我很愛你.
但我爸爸這份,你可否不要教我.
所有人我都可以愛.
這個可不可以不愛.
十多年前我搬離芝加哥.
每次回來,我都很掙扎.
因為我都很掙扎,我要不要去探望他.
因為甚至探望他,其實一點都不好玩.
他還能說話的時候,當然會罵我.
到不能說話的時候,他就會假裝看不到我.
所以我也經常假裝.
我也假裝聽不到聖靈的催迫.
拖到無可再拖,我才會走去探望他.
有一次是這樣,我已經差不多要走.

$^{601}$我又想繼續假裝聽不到聖靈的催迫.
結果當然走不掉.
結果我選擇了拒絕自己.
我讀書的時候寫過我的生命目標.
我就是要一生人拒絕我自己.
我要拒絕我自己.
去遵從聖靈的催迫.
要跟隨神.
我跟我爸爸的關係不好.
但當時我也在想.
是我而已,我沒理由.
拒絕讓我女兒跟公公有關係.
所以我就心有不甘,情不願地.
行屍走肉地去探望他.
那天我特意找了Wendy師母載我去Nursing Home去探望他.
我早就預備了,我去到的時候.
我爸爸會假裝看不到我.
母女倆應該也不能坐很久.
所以我就跟Wendy說.
「你30分鐘就回來了,很快的」.
「坐一坐就走了」.
誰不知一進去.
我女兒一看到我爸爸就哭了出來.
我當時在想.
不知道是不是因為她看到一個老人家.
躺在床上很難過.
所以哭了出來.
我就把我女兒拉進廁所.
問她發生什麼事.
誰不知我女兒竟然跟我說.
「公公經常傷害你」.
「為什麼你還要來探望他呢?」.
我女兒這句話真的嚇了我一跳.
聽完我女兒這樣說之後.
我就跟她說.
「不如這樣吧」.
「如果你不想見到公公的話」.
「你就站在廁所裡等媽媽」.
「媽媽坐一會兒,我們就一起走」.
「但如果你沒問題的話」.

$^{641}$「你不如洗個臉」.
「然後出來陪公公坐一會兒」.
「我們再走」.
結果我女兒洗完臉之後.
就走出來坐在我爸爸的床邊跟他對話.
我爸爸當時已經聽不到話.
說不到話.
所以他跟我的女兒.
只是用筆記本.
寫來寫去地交談.
寫著寫著.
我爸爸寫下了一句.
他說「多謝你跟媽媽探望我」.
我不知道你有沒有.
能夠理解到我當時的感受.
對我來說很震撼.
我心想.
我爸爸竟然承認我的存在.
這麼多年來.
我一直都很抗拒神.
特別是.
我還是那句話.
很愛主.
但我覺得.
我爸爸那一份你不要搞.
所以我一直很抗拒神.
為什麼你經常迫我原諒我爸爸.
神在當中沒有承認我的感受.
但這個經歷.
讓我知道.
神是很了解我的感受.
祂很明白我.
原來是我誤會了祂.
福音重新定義人與人之間的關係.
不知道我們弟兄姊妹當中.
會不會有些人.
你會像我這樣.
我們都很愛主.
我們對身邊的人都很好.
我們很友善,很友善.

$^{681}$唯獨生命中.
就是有這樣的人.
我們想跟神說.
這個人你不要搞我.
讓我走.
真的很難.
你是知的主.
你就放過我.
這個人就由我吧.
在《肥利門書》裡面.
保羅因為基督耶穌的關係.
他願意放下自己的身段.
從神裡面用愛出發.
謙卑自己.
以身作則教導.
牧羊懇求肥利門.
其實也同時懇求你.
不要再計較公不公平.
對方到底虧欠你多少.
傷害你有多深.
因為基督耶穌的關係.
你願意踏出一步.
原諒和接納.
我明白的.
我明白是很難.
我掙扎了N年.
不是說幾年.
是十幾.
兩個數字位數來掙扎.
每一次都掙扎.
每一次見他都掙扎.
但每一次我都嘗試.
不是嘗試.
其實我是完全不想.
我要ignore我自己的感受.
因為我知道.
神要我這樣做.
聖靈催促的聖音是這樣.
因為福音的緣故.
我要放下我自己.

$^{721}$去走那一步.
因為在福音面前.
我們人人都平等.
耶穌為我們死.
不是就這樣說.
讓你得到永生.
我們過著舒舒服服.
基督徒在教會裡.
很開心.
我問候你.
你問候我的生活.
而在生命裡.
可能就是有這個人.
因為你信了主耶穌基督.
因為福音的關係.
我們要走出那一步.
降伏在福音的面前.
信服並且學武.
經歷福音的大能.
去重建.
重新定義我們和其他人.
身邊的人.
一些我們不喜歡的人.
得罪過我們的人的關係.
你願不願意接受教導.
讓福音去接受你自己的意願.
去原諒.
去接納.
我們一起低頭祈禱.
主帕望你今天給我們的訊息.
提醒.
鼓勵.
更加安慰我們當中一些受傷害的人.
你知道你明瞭.
你是愛我們的.
我們需要你去.
讓我們知道.
我們可以.
因為我們有你.
你幫助我們當中一些掙扎的弟兄姊妹.

$^{761}$你施恩.
你同行.
雖然非常之萬以分之艱難.
甚至我們努力很多年都沒有結果.
或者所受的傷害更大.
但是主啊.
至少你幫助我們人願意回到你的面前.
去為自己祈禱.
為對方祈禱.
你讓我們渴慕.
我們想看到福音的大能力.
在中間怎樣去重建.
人與人之間的關係.
幫助我們渴慕聖靈的聲音引領.
回應福音的呼召.
這樣祈禱.
奉主耶穌基督的聖名祈求.
阿們.
\newpage



\section{}
\label{sec:fTMtbscgTAE}
\textbf{Everlasting Praise 5 (Corinthians 14) - Rev. Mark Chan -}
\newline
\newline
連結: \href{https://youtube.com/watch?v=fTMtbscgTAE}{\texttt{https://youtube.com/watch?v=fTMtbscgTAE}} ~~~~ 語音日期: 2024-06-05
\newline
\newline
\hyperref[sec:FBZ55r0dhNA]{\small{< < < PREV SERMON < < <}}
~
\hyperref[sec:index_chronic]{\small{[返順時目]}}
~
\hyperref[sec:index_scriptual]{\small{[返順卷目]}}
~
\hyperref[sec:mwhshrf5O44]{\small{> > > NEXT SERMON > > >}}
\newline
\newline
$^{1}$All right, we continue our series on worship..
And I think we need to understand what worship is all about..
And I need to, we need to come before the word of God.
because the word of God is our final standard..
It's the highest standard for our belief and our practice..
Ask questions as someone who is in Christ..
I want to know what worship is all about..
Is it just about attending an hour of program.
to go through the motions and to just sing some songs.
and spend an hour and then I'll just check out,.
and then I'll go home..
Is it what worship is all about?.
So let me begin today with this main question..
This question is how can we become better worshipers?.
Or I should say, how can we be worshipers.
who are pleasing to God?.
How can we do that?.
So let's take a moment and talk to your neighbor,.
if you can..
How can we be better worshipers.
or worshipers that will please God?.
All right, think about it..
What would be your answers?.
Now, sometimes when we approach questions like this,.
or some of the Christian practice.
revealed to us in scripture,.
sometimes we approach questions like this,.
like we approach our boss..
I'll talk to some of you,.
and some of you are working from home..
And someone share a secret,.
because working from home,.
and you have to turn on Microsoft Team.
or Zoom or the online platform,.
and your boss has to know that you are present..
And some of you have a smart way to fake your presence,.
you know, to get your, some sort of software..
I don't know about this, but until someone told me..
You can actually buy a software and install a software,.
and then mechanically remove the mouse from time to time,.

$^{41}$and just fake your presence in the online meeting..
And that's a very, very wise way to do it..
I would say this is beyond my wisdom..
But, you know, is that all God cares about, right?.
I mean, do we treat God as our boss?.
You know, our boss simply, you know,.
want us to check in, right?.
Clock in and clock out..
Want your time, right?.
And then we just fake our presence, right?.
Is that what God care about?.
What does God really care about, right?.
In this series, we say that in scripture,.
God is revealing to us who he is and what he does..
From creation to new creation,.
what God really care about is that he will live among us..
He will be our God and we will be his people..
And that's what worship is all about..
It's not about clock in and clock out..
Okay, I just check this..
Okay, I've done this for this week,.
and I'll get worship done, and I'll sing some songs to God,.
and that's it, and I'll go my way, right?.
Is that what God really care about?.
God cares about his presence with us,.
his relationship with us..
Your boss will not care about his relationship with you,.
right?.
That's why, you know, it is legitimate to think about.
how to fake your presence, right?.
Because all your boss cares about is your productivity,.
right?.
Squeeze you in some way that use you, right?.
Use you for productivity, for profit..
But God cares about his relationship with us..
And that's what it means when we believe in Jesus Christ..
It's not about do's and don'ts..
It's about our relationship with Christ..
I was talking to someone, his friend,.
is a very devoted Catholic..

$^{81}$And he was sharing with me,.
oh no, his friend is not willing to come to church..
Going to church is not just part of his friend's agenda..
Because as a Catholic, a lot of Catholics are like this..
And of course, a lot of Christians are like this..
A lot of evangelical Christians are like that..
No, a lot of Catholics are, you know,.
when they approach faith, when they approach God,.
it's a cultural thing, right?.
Because my father and my mother were devout Catholics,.
so that's why I am Catholic..
And that's part of my cultural identity..
And if you ask a devout Catholic,.
so often you get the answer is that,.
of course I am a Catholic, right?.
Now Catholic in the sense of cultural identity, right?.
It's not really a personal encounter.
or personal intimate relationship with God..
And that's what many, many Catholics are lacking..
And I'm not trying to mention,.
but many, many Christians are like this too..
Many Bible-believing Christians are like that..
The kind of a cultural identity,.
kind of, you know, attack on an app in your phone, right?.
And you never use it..
You got some apps in your phone, right?.
On your phone, installed, but you never touch it..
You never use it, right?.
In case you need it, and then you click on it, right?.
It's not just part of your system..
It's not, right?.
Many of us approach worship, approach God in that sense..
And we lost what really God care,.
what God really care about..
He care about His relationship with us..
So therefore in the entire Bible,.
from creation to new creation,.
from the first chapter to the last chapter,.
it's all about worship..
In some sense, everything is about worship..

$^{121}$Everything is about God living in you,.
and you live in God..
Everything is about you and God relationship..
So we talk about creation..
The Bible is a grand story..
And even though we rebel against God,.
that's Genesis chapter three,.
but God is not abandoning the fallen world.
to rebel against mankind..
God chooses to redeem them and to live among them..
And therefore He saved the Israelites out of Egypt.
by grace, and then live among them in the form of tabernacle,.
the tent of Portable Temple,.
and living with them and guiding their way.
through wilderness..
And then once they settle down in the land of Canaan,.
the promised land, God is dwelling with them.
in a stone temple, right?.
And as times go by, we know that all of this.
Old Testament symbols are pointing forward.
to Christ Jesus, our Lord..
And Christ came and Jesus came and died on the cross.
to redeem His people and rose from the dead.
and gave us hope..
And in Christ now, we worship, not in the stone temple,.
not in the tabernacle, because everything points.
to the reality of Christ..
So in Christ we worship..
In Christ we encounter God,.
and we build our relationship with Him..
Until someday in eternity, when new heaven and new earth.
came down in the midst of us, and God lives among us,.
and that is what we call heaven,.
or should we call new creation, new heaven and new earth..
So the entire story of the Bible tells us.
it's all about His relationship with us,.
and that's what worship is all about..
Now, as we come to the New Testament,.
we find that there were a lot of teachings about worship.
and what worship is all about..

$^{161}$Acts chapter two describes the early church,.
the first generation Christians, what they do in worship..
All the believers devote themselves to the teachings.
of the apostles and fellowship,.
sharing in meals and to prayer..
A deep sense of awe came over them all,.
and the apostles performed many signs and wonders,.
and all the believers, all the disciples met together.
in one place and shared everything they had..
All the while, they praised God.
and enjoyed the goodwill of the people..
Now, from this short passage, we can see.
worship is a priority..
It's a priority for the tiny Christian movement.
in the beginning..
They gather together in homes,.
you know, house churches, what we call,.
and worship is always a priority, not an option..
And then worship is about a corporate community.
getting together, fellowship,.
face-to-face, sharing with one another,.
and they practice generosity..
They give, right?.
Voluntarily, they give their resources.
to build up the church, the community..
And worship is also a puppet weakness to the world,.
enjoying the goodwill of all the people,.
all the neighbors, all the people, not just the believers..
And because when they see them worshiping together,.
when they see them gathering together,.
there is a function, or there is a power.
of outreaching missional, being missional in worship..
You know, when people come into our gathering,.
they hear us praising God, they hear us worshiping God..
What should happen is that they will sense.
that this is a different community,.
a different community that I cannot find.
in the rest of the world..
Someone pointed to me sometimes.
on some worship songs on YouTube,.

$^{201}$and there were people who comment,.
you know, comment below a certain video.
or a certain worship song..
And some of the comments are really interesting..
And some of the comments, let me read that..
Some of the comments reads,.
"I am not a Christian, but I am touched by the music..
"I don't know this song,.
"but I'm touched by the message of the song.".
Now that's what worship should do, ultimately..
It is going out, right, reaching out to other people.
so that people can understand our music.
and understand what we are doing,.
and they sense there's something different,.
you know, in the midst of Christian worship..
So today I would like to turn to one of the passages.
in the New Testament that talks about worship..
And that is 1 Corinthians chapter 14..
As you know, the entire letter of Corinthians.
is a corrective letter..
It means that the pastor, the apostle Paul,.
he writes the letter to correct some of the misbehavior.
and misbelief amongst the house churches,.
tiny house churches in Corinth..
And there were many, many problems in the churches,.
or the house churches in Corinth..
Division, they have a lawsuit against one another,.
they have sexual immorality,.
idle food controversies, and all that..
So the entire letter is trying to correct one problem.
after the other..
In 1 Corinthians 14, actually from 12 to 14,.
Paul is addressing problems in the worship,.
in the corporate gathering..
And in chapter 14, there were two parts..
First of all, Paul is talking about the difference.
between prophecy and tongues, speaking in tongues..
And the second part of this chapter,.
he talks about your worship is a witness to other people..
So let's break this down and try to understand.

$^{241}$what is going on here..
Chapter 14, verse one..
Paul said, "Let love be your highest priority, highest goal,.
"but you should also desire the special gift.
"or special abilities," by the way,.
this is New Living Translation,.
which is easier to understand..
"Special abilities that the Spirit gives,.
"especially the ability or the spiritual gift to prophesy..
"For one who speaks in a tongue.
"speaks not to men, but to God..
"For no one understands him,.
"but he utters mysteries in the Spirit.".
Now what is he talking about?.
Now, first of all, what is tongue?.
What is speaking in tongue?.
I'm not sure you have heard, maybe you have,.
you have heard someone speaking in tongue, right?.
Yeah, when I was in seminary, my first seminary.
was Westminster Theological Seminary in California..
And I lived with three other roommates, they are Koreans..
And one of those guys, in the same age,.
when I was really young, right, 20 something..
And then we lived together, four of us..
And one of my roommate, Korean roommates,.
he speaks in tongue, every morning..
He opened the Bible in his room,.
he does his devotional, and he will pray..
And he will pray in some sort of unknown language, right?.
Of course, I don't know Korean, right?.
I'm not sure he was praying in Korean..
So sometimes I ask him, "Hey, you're praying, right?.
"What kind of language are you praying?".
Because it's different from normal conversation,.
normal speaking of Korean..
So I ask him, "What kind of language are you speaking?".
"No," he said, "I'm speaking in tongue.".
Some sort of mysterious, unknown language.
that only he and God understand, right?.
There's some sort of intimate language.

$^{281}$that only God and him would communicate, right?.
This is kind of like, you know,.
when you hear intimate couple talking to each other,.
you don't understand what they're saying..
But they're talking, right?.
I remember one time I was in a church conference in Texas..
I was all by myself, right?.
I went to a restaurant and I ordered my entree.
and I was enjoying my dinner..
And next to me, at the other table,.
was a dating couple, I guess, right?.
They're pretty young and obviously.
they're in love with each other..
You know, I can understand every single word they say.
in English, but I don't know what they mean..
They have their own intimate connection, right?.
Communication..
I guess speaking in tongue is kind of like that, right?.
So some people have the gift of speaking in tongue,.
some sort of mysterious, unknown language,.
his intimacy with God, and you can't understand.
until it is being interpreted, right?.
So that's speaking in tongue..
And obviously there are some people who are gifted.
in such a way in the Church of Corinth..
So as they gather together, some of the people.
will speak in tongue, you know, freely, right?.
And people around them who don't have the gift.
and where there's no translator or interpreter,.
there's no way you can understand what they say..
That is speaking in tongue..
And how about prophecy, all right?.
Prophecy..
So Paul is trying to contrast speaking in tongue.
and prophecy, right?.
But what is prophecy?.
Now, modern people will always understand prophecy.
in a very narrow way..
What is prophecy to you?.
What does prophecy do?.

$^{321}$It's always prediction, right?.
Weather forecast..
Oh, will it rain tomorrow at 2.50 p.m.?.
Will it snow tomorrow, right?.
We always approach prophecy that way..
But, you know, in the Bible, prophecy is much broader..
You know, look at the prophets in the Old Testament..
The prophets are spokepersons of God..
And they always begin the message by, "Thus says the Lord.".
This is what Yahweh the Lord says..
And you better listen..
And that's what prophets would do, would say..
And it's kind of like preaching..
It's kind of like biblical teaching..
So prophecy is not just about, you know, weather forecast,.
or what will happen tomorrow..
It's not about prediction..
It's about proclamation of the Word of God,.
the authoritative Word of God..
So Paul is trying to contrast prophecy and tongues..
And what he's trying to get at is that prophecy,.
or preaching, or biblical teaching,.
is understandable to all the people..
It's intelligent to all the people,.
because people can understand..
It's a known language..
And people can understand its meaning..
But speaking in tongue, that is a barrier, right?.
Because you need to wait for someone to translate.
or to interpret in order to understand it..
And Paul, in this passage,.
he talks so much about using those gifts..
He said in the corporate worship setting,.
you better not use tongues..
You better use prophecy..
You better speak in such a way.
that everyone can understand..
So verse four, he says,.
"A person who speaks in tongue is strengthened personally,".
because it's an intimate language, right?.

$^{361}$Only God and him would enjoy it..
But nobody else around him can enjoy it..
So it's only strengthened personally,.
but one who speaks a word of prophecy, preaching,.
strengthens the entire church..
And I wish you could all speak in tongues,.
because Paul, obviously, he also speaks in tongue..
But he never does so in the worship setting,.
in the corporate setting..
But even more, I wish you could all prophesy, right?.
You could all speak in such a way,.
teach in such a way, preach in such a way,.
everyone can understand..
For prophecy is greater than speaking in tongues..
Unless someone interprets what you're saying.
so that the whole church will be strengthened..
Now notice the key words here, strengthen, right?.
Or the older translation would say to edify, right?.
It's a architectural term..
So everyone getting together in a church setting,.
in a corporate gathering,.
our purpose is to strengthen one another,.
to pray for one another, to build up one another..
Your presence, your prayers, your ministry.
is to edify other people.
so that the church can be built up, right?.
And the church, the word church,.
repeated many, many times in the entire chapter,.
nine times in total..
Paul is saying we should use our spiritual gifts,.
our talents, our abilities to edify the church,.
to build up the church..
And what we have, our skills and our talents.
are not for ourselves, but means to bless other people..
It's not my own possession.
because the Spirit gives you in such a way.
that you can serve other people..
I have been to many churches.
and I've always find that it's such a great blessing.
to see people using their gifts.

$^{401}$and using their skills to build up.
the community of the church..
Some people who are gifted to help other people, right?.
Help in a very practical way..
I know a brother in California..
Every time I meet with him,.
he always ask the question,.
are you in need financially, right, financially?.
He's gifted in such a way that,.
I wanna bless you because God bless me so much..
And I wanna bless you in such a way.
that God can use you, right?.
It's always good to have a friend like that, right?.
But I don't see him often..
I don't see him often..
But he's just gifted in giving, right?.
And some people who are gifted.
in emotional relational intelligence..
Some people you can meet with them.
and they will encourage you in the conversations.
and they will build you up..
And you find the conversation with him or with her.
is life-giving, right?.
Some people are especially gifted in such a way..
Of course, we're always commanded to love one another,.
to care for one another,.
but some people are especially gifted in such a way..
And some people who are obviously gifted in music, right?.
And they play music.
and so that we can sing songs together, right?.
And some of the gifts need to be developed.
and to be trained to be used..
Otherwise, you would lose it, right?.
You realize I play most of the instruments, right?.
Why do I do that?.
It's not necessarily I enjoy it,.
but because when I go to different churches,.
you know, I serve different churches in the past,.
I find that it's very useful..
And if I can get a handle on it,.

$^{441}$I'll learn it myself, right?.
I'll learn it myself..
Because I want to bless the community that I'm in, right?.
The fellowship, the small group, or the churches,.
wherever I am..
A couple of weeks ago, we had the MSI event..
We went to a museum of science industry..
And for church-wide events, sometimes I do observe.
that we are blessed with many, many people.
who don't take advantage,.
but they understand the purpose of the event.
is to connect with other people,.
especially non-Christians..
Some people who bring new friends,.
and they just take the initiative to approach those friends.
and talk to them and to relate to them..
But on the other side of the coin,.
sometimes I see Christians who just enjoy their own time,.
you know, and take advantage..
And this is sad, right?.
Because we lose the focus, we lose the purpose.
of why we gather together..
Paul is saying why we gather together is to strengthen,.
is to edify one another..
Therefore, Paul is saying prophecy..
When you speak in such a way,.
you can proclaim the word of God.
in such a way everyone can understand this is greater,.
this is stronger than speaking in tongue..
You just enjoy your spiritual high, right?.
The whole church will be strengthened..
And that's the first part of this chapter..
The second part of this chapter,.
let me read part of the passage here..
And this is a very important passage..
Verse 23 and following..
Even so, if unbelievers or people.
who don't understand these things, the spiritual things,.
come into your church meeting, your church worship,.
and hear everyone speaking in an unknown language,.

$^{481}$speaking in tongue,.
they will think that you guys are crazy,.
you guys are nuts, right?.
Because they don't understand..
Verse 24, but if all of you are prophesying,.
preaching, explaining the faith.
in such an understandable way,.
unbelievers or people who don't understand.
the spiritual things come into your worship meeting,.
they will be convicted of sin and judged by what you say..
Verse 25, as they listen,.
their secret thoughts will be exposed..
They will fall to their knees and worship God,.
declaring, "God is truly here among you.".
You see the key word here is understand..
It's the principle of understandability.
or intelligibility..
It's the principle of doing something.
that people can understand, even visitors,.
non-Christian visitors, will be able to understand..
(stammers).
Let me ask you a question..
Is our worship design.
understandable to most people?.
Yes, right?.
Yeah, we use, I'm very grateful for worship leaders.
and people who serve in the worship ministry.
and they choose songs that are newer.
or they're more contemporary, right?.
I'm not saying we should abandon the traditional hymns, okay?.
All right, for those, well, I'm probably the older guy here..
That's so sad..
But anyways, when we grow up in the church,.
we learn a lot of traditional hymns,.
many, many traditional hymns..
And when I hear those traditional hymns,.
I immediately go, this is my language, right?.
This is my language..
And we used to sing from a hymn book, right?.
We sing the old, old hymns..

$^{521}$But we need to understand worship as a language..
It's kind of a language or kind of a form.
that is always evolving,.
from time to time, from generation to generation, right?.
Those songs used to be a great hit.
and no longer the language of today, right?.
Again, I'm not saying we have to abandon.
everything traditional..
I'm not trying to say, let us embrace the newest thing,.
right, like the iPhone, right?.
The newest iPhone, everyone wants to get it, right?.
I'm not trying to say that,.
but the mission of the church, Paul says,.
is to make worship gathering in such a way.
that will honor God and that is also understandable.
to other people, right?.
So I'm grateful for those who try to use songs.
or new songs to worship,.
because the mission of the church.
is not to keep certain tradition, right?.
That is not our mission..
Our mission is to make faith and the Bible.
and our message understandable.
and that is the spirit of the passage here..
Those of you who have been to a traditional worship setting,.
when you come in, you pick up a copy of the bulletin, right?.
Do you still remember those days?.
Or I'm just dating myself, right?.
Yeah, you pick up a copy of a sheet of paper.
with the, it's called worship bulletin, okay?.
And it lists every single component of the worship service..
The first one is usually is the call to worship.
and then there's a prayer and then there's another prayer.
and then you stand up for the prayer.
and then you sit down for another prayer.
and then there's a scripture reading and then preaching.
and then you sit down and you do a certain hymn.
and then after you do a certain hymn.
and then you do some other component.
and then you stand up again for a doxology.

$^{561}$and then you sit down and then you do your silent prayer.
and then you go, right?.
You see what I'm saying?.
There's a traditional liturgy, right?.
Traditional liturgy..
And some of the churches still do that, right?.
But we don't, right?.
We don't..
Why don't we print a worship bulletin.
and hand you one copy when you come in?.
Why don't we do that?.
We don't..
Because we're lazy..
We don't wanna do some stuff, right?.
Well, no, no, no..
We wanna be environmental, right?.
We don't wanna waste paper, right?.
I've been to some churches.
where traditional in worship form,.
after the worship service, you see papers lying around, right?.
And then you just go to,.
those papers will go to the trash..
But what I'm saying is the form.
or the language of worship evolves from time to time, right?.
Even the instruments that we use, right?.
In my old days, in my old days,.
we only do organ and piano, right?.
So many, many of the traditional churches,.
they have to have organist, someone who play organ, right?.
Wouldn't it be nice if we have a pipe organ in the church?.
Don't you think so?.
Yes?.
Yeah, yeah, yeah..
But the maintenance will kill you, right?.
You need to spend a lot of money to maintain those stuff,.
right?.
We use electric piano, right?.
Electric keyboard, right?.
Yeah, let me tell you this little story about drum set..
Okay, I told some of you about this drum set..

$^{601}$Now, I remember in the '90s,.
I was in my church in Vancouver, in Canada..
We didn't have a drum set..
And some of the brothers and sisters wanted to suggest,.
"Oh, let's use drum set.".
Those are the early days of the worship team, right?.
The band, worship band, music band..
And we didn't have a drum set..
And the deacon board called for a meeting to vote.
whether we should buy a drum set using church budget..
And there was a huge controversy..
It was a huge controversy because some people,.
the older folks, they said, "Oh, we should buy a drum set..
"The older folks especially,.
"they are against the idea of putting a drum on stage.
"or buying a drum using church money.".
You know why they think that way?.
You have to understand, they grew up in the '70s,.
you know, the '60s and the '70s..
Those are the days, you know, sexual revolution.
and all that liberation movement, right?.
And for them, this thing, drum set,.
is for those who smoke weeds, right?.
Those rock band, right?.
The rock musicians, they use this kind of instrument..
But they cannot change, right?.
Because the language has changed, right?.
And people use drum set for a different reason now, right?.
People play electric guitar for a different reason now..
Those are the old days, all the rebellious teenagers,.
the young people, would play those instruments, right?.
And they never thought of putting those instruments.
in the church..
So there was a huge debate.
about whether we should buy a drum set..
And what is the result of the vote?.
It was no..
We don't buy it, right?.
And so it didn't carry through..
So I was serving in the church and I was saying,.

$^{641}$okay, let me just buy it out of my own pocket..
So I did..
I spent almost \$2,000 to buy this thing, right?.
So this drum set belongs to me..
It follows me..
It follows me every single day, you know?.
Whenever I serve in the church, I say,.
hey, you don't have a drum set?.
Boom, here..
Okay, it's the instrument,.
it's a personal instrument of the pastor, okay?.
The pastor plays drums too, okay?.
So it's okay, all right?.
It's not like we are going astray or something like that..
But what I'm trying to tell you is that.
there is a different language.
from generation to generation..
Even the attire is different..
What we wear is changing, right?.
You go to the tech talks,.
you go to the big announcement from the big company,.
the IT companies,.
you know, those CEOs, they don't wear suit and tie, right?.
They don't do that anymore, right?.
We change, right?.
Whoever does the announcement is short, right?.
And it's okay, right?.
And it's okay..
And we have to distinguish what the Bible,.
what the Bible says, what the Bible says,.
and what tradition says, right?.
God want us to hear his word,.
but not just our tradition..
And that's the spirit of the passage..
Even the songs that we use today, right?.
I remember some months ago,.
I was trying to pause my Hebrew, my biblical Hebrew..
So I tried to learn some Hebrew songs..
So it's a very beautiful songs..
You know, the Hebrew songs, the Middle Eastern melody,.

$^{681}$is actually very, very sweet, all right?.
So I was trying to play those songs.
because I wanna learn some new Hebrew words, right?.
Like Psalm 23, you can sing the whole thing, right?.
And I was using my guitar.
and trying to sing some Hebrew songs.
so that I can remember some of the words, right?.
And then my kids, my children,.
they came over and say,.
"Can you please stop?".
Because it doesn't sound good, right?.
Actually, it sounds good, right?.
Actually, it sounds good..
And I sing well..
And you know, what I'm saying is the language has changed..
The melody has changed..
And we have to be relevant, right?.
So the principle here is that worship is vertical..
Worship is that we bring the offering to God vertically..
Because we are the people of God,.
and God loves us, and God cares about us, we worship Him..
All because God has created and redeemed us,.
and we respond to the great acts of God..
We respond by saying, "God, thank you for saving us..
"Thank you for letting us be your people.".
So worship is, first of all, vertical..
But it is also horizontal..
It's all about other people,.
and we cannot ignore one another..
And worship needs to be edifying,.
to build up other people.
so that people can understand the message..
And even it can be, it needs to be missional, evangelistic..
When people walk into our church,.
when people join our meetings,.
when people come into our worship,.
they will sense that this is different, right?.
And it's not about the size..
It's about the message that we carry, right?.
People can see, hey, God is in the midst of you..

$^{721}$And let me close by trying to answer this question..
How can we be better worshipers?.
First, I would like to suggest, be mindful..
Be mindful..
When you come in and worship God,.
be mindful in the sense that we are approaching God,.
we are approaching God with a pure heart,.
and we need to be fully present..
Our mind, our spirit need to be fully present..
Even our physical body needs to be fully present.
to prepare ourselves so that we can come.
to the presence of God to be a worshiper pleasing to Him..
And secondly, let us be reminded.
to be obedient in our attitude..
When we come in and pray, prepare our hearts.
so that our hearts is a good soil..
We can receive the word of God, the message of God.
through songs and prayers and sermons and all that we do..
Be obedient..
And lastly, let me remind you to be prayerful..
We approach this holy God, this great and awesome God, right?.
And we need to approach Him with a clean heart..
I always find this good practice..
I come into a worship, corporate worship,.
a little bit earlier, and I do a little prayer,.
and I ask for forgiveness..
I ask that God will forgive my sins and cleanse me..
And I pray that God will prepare my heart to approach Him..
And sometimes I even need to pray.
and forgive other people, right?.
Other people, other brothers and sisters,.
other family members..
I need to ask them, I ask God to forgive them, right?.
So that I can approach God in a new and fresh way..
How can you become a better worshiper?.
Let us pray and prepare our hearts.
to receive the communion today..
Let's pray..
Father, as we come before your presence,.
we want to say thank you..

$^{761}$Thank you for forgiving our sins,.
for cleansing us, for redeeming us,.
all because what Jesus our Lord has done on the cross..
He died and shed blood in order for all of our sins,.
the past, present, and future, everything forgiven.
so that we can be called, declared righteous in Christ..
Thank you for what you have done..
And our worship is a response to your grace..
Thank you for loving us..
Thank you for calling us to worship you..
So we ask that you will cleanse our hearts and our mind.
so we will focus on you..
As we come before the Lord's table today,.
we ask that you will refresh our heart and our spirit.
as we partake the bread and the cup..
These are the symbols, the tangible reminders.
of your love to us..
And thank you..
And may your grace continue to transform us..
In Christ's name we pray..
\newpage



\section{哥林多前書}
\label{sec:mwhshrf5O44}
\textbf{永恆的敬拜 4 (哥林多前書 14) - 陳崇基牧師}
\newline
\newline
連結: \href{https://youtube.com/watch?v=mwhshrf5O44}{\texttt{https://youtube.com/watch?v=mwhshrf5O44}} ~~~~ 語音日期: 2024-06-05
\newline
\newline
\hyperref[sec:fTMtbscgTAE]{\small{< < < PREV SERMON < < <}}
~
\hyperref[sec:index_chronic]{\small{[返順時目]}}
~
\hyperref[sec:index_scriptual]{\small{[返順卷目]}}
~
\hyperref[sec:Gu87k6iNMK0]{\small{> > > NEXT SERMON > > >}}
\newline
\newline
$^{1}$早安.
跟大家打招呼.
身邊的人跟他們說早安.
夠了.
不要說太多.
這次我們繼續去想想.
這個topic.
崇拜的topic.
我想大家.
信主或者甚至你可能.
第一次來教會.
你都會發現崇拜.
公共的worship.
是一個很重要的環節.
在我們心靈裡面.
在我們基督徒的journey裡面.
在我們的旅程裡面.
其實是非常重要的.
每一個禮拜我們會get together.
每一個禮拜我們會練歌.
和預備等等.
我們今天進來.
崇拜大概一個小時的時間.
我們有不同的環節.
你知道背後有很多弟兄姊妹.
做了很多事.
他們練歌在後台.
很多不同的細節.
以致我們今天.
能夠享受崇拜的聚會.
所以大家需要知道.
為什麼.
所謂corporate worship.
corporate worship.
公共崇拜.
集體崇拜如此重要.
因為聖經裡面告訴我們很多章節.
很多經文.
告訴我們我們需要做一個.
認識敬拜是一回事.

$^{41}$和怎樣去成為一個.
更好的敬拜者.
我想問你第一個問題.
我們怎樣能夠成為一個.
更加好或更合上帝心意的敬拜者.
ok.
或者更好不是一個最好的描述方法.
怎樣成為一個更加上帝.
去到喜悅的敬拜者.
這裡講得很清楚.
不是純粹喜悅的敬拜.
而是喜悅的敬拜者.
因為我們很多時候進來.
去教會參與崇拜的時候.
我們很多時候會有錯覺.
覺得上帝只要我一個小時時間.
好了我打完卡就走了.
好像你老闆一樣.
你老闆要你做事.
我最近和一些年輕人聊天.
他們說很多時候他們在家工作.
他就買了一個程式.
就是弄到他的滑鼠可以自己動.
然後他可以去廁所.
去吃東西.
假扮自己在動.
這種心態是一種很聰明的心態.
當然.
這種方法.
你想到老闆要你做什麼.
老闆就要你.
你開啟你的Microsoft team.
或者你的Zoom.
或者你網上的東西.
他就監視你的滑鼠有沒有動.
你的滑鼠動就當你在做事.
我們會不會這樣看上帝.
其實上帝不是這樣的.
上帝不是在乎你純粹一個小時的時間.
上帝在乎你.

$^{81}$是你自己本身.
在崇拜裡面他能夠和你相遇.
我知道大家.
我們生活裡面遇見很多不同的問題.
婚姻裡面有很多張力.
家庭裡面有很多負擔.
經濟裡面.
工作裡面有很多濕濕的東西等等.
肯定的.
我們都背著很多不同的包袱進來.
我們來到敬拜的時候.
但你知道上帝所關注的是什麼.
關注的是你.
是你的心靈.
上帝真的很關懷你的心靈.
上帝能夠和你相遇.
以至他能夠加力給你.
可能你正在經歷洪水.
但上帝告訴你.
他仍然坐著為王.
他仍然供應你.
上帝告訴你他仍然坐著為王.
他仍然供應你.
你信不信.
當我們這樣去想敬拜的時候.
我們就不會想著.
老闆要我的貓水動.
我就弄它動.
上帝是一位愛我們和關心我們.
與我們同在的上帝.
所以我們看敬拜的系列的時候.
我們曾經提過.
從開始直到永恆.
上帝是想和我們建立關係.
和我們建立一個.
你可以叫友誼.
Friendship.
上帝是想和你同在.
在你的生命裡面同在.
當我們信耶穌的時候.

$^{121}$上帝的生命就和我們在一起.
由創世開始.
上帝就說.
這個是我的聖殿.
就是我所居住的地方.
我去創造你的時候.
你的生命就是我所居住.
能夠和你同行的一個地方.
雖然我們會墮落.
這個所謂rebellion.
或者我們叫它fall.
就是人類的叛逆.
創世的第三章.
我們不要認這個是上帝.
但是上帝從來沒有放棄過我們.
從來沒有放棄過這個墮落的世界.
這個搞到亂七八糟的世界.
甚至我們生命已經搞到亂七八糟.
上帝沒有放棄過.
祂就是繼續去救贖我們.
去關懷我們.
所以祂就和上帝的子民同在.
拯救這些人出了埃及.
然後在會幕當中和他們同行.
在曠野裡面和他們同行.
在聖殿裡面.
說這個是我的名所同在的地方.
在國民當中和他們一起.
直到基督來了.
在基督耶穌裡面.
我們就能夠和這位上帝.
直接面對面相遇.
去和祂建立這個關係.
重新恢復這個應該有的關係.
所以我們每一個信耶穌的人.
在基督裡面.
我們能夠去敬拜神.
直到去新天新地新的創造.
New Creation.
都是和敬拜有關.

$^{161}$和上帝和我們同住有關.
所以上帝所關注的.
其實是和我們之間的關係.
所以這次我們看過伊甸園.
會幕和聖殿等等.
我們去到新約.
新約裡面有很多有關敬拜的教導.
譬如少林寺開始的時候.
初期的教會.
他們很多時候是在家裡面敬拜的.
所以我們想起譬如哥林多教會.
羅馬教會.
我們不是想起像我們今天這樣的教會.
因為他們很多時候是在自己的家裡面.
來去敬拜.
在誰的家裡.
有些小組的.
很小的教會.
很少數的基督徒.
上帝就和他們同在.
少林寺第二章.
這裡說到初期教會建立的時候.
他說他們就是門徒.
初期教會他們專注在使徒的教導和彼此的團契.
抹餅和去祈禱.
眾人是心存敬畏.
使徒也都行了很多不同的神蹟.
信的人都聚在一處.
泛物公容.
讚美敬拜上帝.
得到全體百姓的喜愛.
The good will of the people.
主就將得救的人數天天加給他們.
這裡告訴我們初期的教會是怎樣的.
他們聚集一起有corporate worship的時候.
他們就是以敬拜.
以聚集敬拜為他們的priority.
是他們生命的優次.
是優先.
不是一些次要的東西或選項的東西.

$^{201}$是他們需要做.
一起去做.
然後他們這個群體就可以彼此去分享.
泛物公容.
是自願的.
將他們的資源拿出來.
然後大家可以一起來分享.
是generosity.
是慷慨的.
大家一起來分享.
然後他們對外也是一個美好的見證.
有一個outreach的能力.
當他們聚在一起的時候.
在什麼地方都好.
少數的人也好多數的人也好.
就是一個去到所謂外展.
對外一個的見證.
你見到嗎.
得全體百姓的喜愛.
good will of the people.
敬拜是應該這樣的.
當你一起來去讚美上帝.
一起用歌聲用祈禱來去讚美上帝的時候.
人的生命很多時候就會因為這樣而被觸動.
就算是一些未信主的人.
或者第一次來教會的人.
他們見到你們這樣敬拜的時候.
他們需要見到需要聽到.
需要見證到.
然後他們就知道你們是不同的.
最近有人告訴我.
有些歌曲在youtube上面.
有些新的敬拜的詩歌.
youtube下面有一些comment.
user的comment.
有些comment很有趣.
有些comment說我雖然不是基督徒.
但我覺得這首歌很觸動我.
這首歌很感動到我.
音樂或者敬拜的音樂.

$^{241}$有一個這樣的power在那裡.
大家不是很喜歡音樂吧.
喜歡吧.
我知道華人的文化.
華人的家庭都不是很喜歡音樂.
我們只是假扮.
小朋友為什麼要學琴.
打到他要學.
死都要學 哭都要練.
為什麼 因為我們是工具性來看那些東西.
因為可以鍛鍊他的生命.
是的 對吧.
如果說去聽音樂會.
交響樂的音樂會.
你去不去.
我們不去.
小朋友去不去.
不去.
給那麼多錢.
吃了它吧.
我們其實不care的.
對吧.
這樣很糟糕.
因為這是上帝所創造的東西.
你怎麼可以這樣.
所以想清楚.
叫小朋友學琴的時候.
不要工具性.
小朋友知道.
你根本不喜歡音樂.
所以學校就不學了.
Anyways.
為什麼說這裡.
因為我說敬拜的音樂.
是有那種感染力.
有那種感染力.
音樂是有那種感染力.
是可以跨文化的感染力.
對吧.
所以想清楚.

$^{281}$今天我們想去看一段經文.
是哥林多前書第十四章.
哥林多前書大家知道是一個.
寫給哥林多的宗教會.
不同小組教會的一封書信.
是一個所謂corrective letter.
就是一個改正性的一種書信.
因為裡面有很多很多問題.
有紛爭.
有不同的.
甚至敬拜裡面都有問題.
我們就問.
為什麼敬拜會出問題.
我們看一段哥林多前書.
其中一段.
十二至十四章.
全部都是處理敬拜的問題.
在他們自己的敬拜當中.
哥林多的信徒.
是有很多很多這樣的問題.
到第十四章.
他就處理一個特別的特定的問題.
保羅處理什麼問題.
就是所謂謊言的問題.
所以這一章有四十節.
我們當然不能夠看完.
頭一半一至十九節.
其實就在說.
保羅在討論.
所謂說預言和說謊言之間.
有什麼差別.
還有為什麼說預言比說謊言好.
第二部分.
二十節到四十節.
他說我們公共的敬拜.
有什麼見證.
我們先解釋一下.
先體驗經文第一節.
這裡說什麼呢.
他說你們要追求愛.

$^{321}$也要切勿屬靈的恩賜.
尤其是作先知講道的恩賜.
那說謊言的不是對人說的.
是對上帝說.
因為沒有人聽得懂.
他是藉著聖靈說各種的奧秘.
說謊言的是造就自己.
作先知講道的是造就教會.
我希望你們都說謊言.
更希望你們作先知講道.
因為說謊言的若不解釋.
沒有翻譯出來的話.
使教會得造就.
那作先知講道的就比他強了.
什麼事呢.
首先什麼叫謊言.
Speaking in tongues.
很簡單的說.
雖然我們的傳統不是這樣.
有些教會的傳統.
或者有些教派的傳統.
他們有敬拜裡面可能說謊言等等.
說謊言簡單來說.
其實就是說一些unknown language.
ok 是沒有人聽得懂的.
但是講者是和上帝之間有特別的契合.
我不知道大家有沒有聽過.
不要覺得是什麼事.
有些奇怪的現象.
以前我讀神學院的時候.
有三個韓國的室友.
Rooming.
其中一個每天早上他自己訪問.
然後他靈修的時候.
他就會說謊言.
我問他在說什麼.
他說不是在說韓國話.
是在說謊言.
無論如何.
總之是一種unknown language.

$^{361}$和上帝的一種契合.
你可以這樣想像.
兩個很相愛的人.
他們大家有種契合的語言.
就好像是這樣.
我記得有一次我去餐館吃飯.
我一個人去吃飯.
我忘記我去哪個會議.
我一個人自己出去吃飯.
我就在那裡吃飯.
我訂了東西.
我就在那裡吃.
隔壁桌有一對男女.
我想他們是熱戀當中.
他們在那裡聊天.
我英文也算是差不多.
我聽懂每一隻字.
但我不知道他們在說什麼.
那種契合的語言.
我不知道他們在說什麼.
但每一隻字我都聽得懂.
這就是謊言的意思.
他們才明白.
有種人是有種說謊言的恩賜.
說語言是什麼意思呢.
很多時候我們說語言的時候.
我們現代人就會說.
明天三五點下不下雨.
我們說的是預測.
但其實語言在聖經裡面不是這個意思.
其實是說到.
說出上帝的權威的說話.
是上帝的代言人.
這樣的意思.
你記得舊約裡面的先知.
他是怎樣說話的.
他會說如此說.
先知經常有這個說話.
耶和華如此說.
因為上帝給先知的訊息.

$^{401}$先知就要帶給上帝的子民.
所以耶和華如此說.
聽清楚了上帝是這樣說.
所以有點像今天的preaching.
就是講道的意思.
聖經的教導這樣的意思.
保羅就在這裡.
將兩樣東西對答.
或者比較.
一種是說語言.
一種是說.
講道.
先知講道.
他說什麼呢.
他說整個十四章裡面.
經常都強調在敬拜裡面.
應該要做些什麼.
你看到我highlight了什麼.
Understand.
要大家聽得懂才行.
如果大家聽不懂的時候.
大家的教會就不能得到造就.
所以經常都講聽不聽得懂.
然後也經常都講.
能否造就.
造就的意思就是strengthen.
強化到或者幫助到大家.
或者.
以前英文翻譯叫做edify.
是edify.
如果你是講一些東西.
大家不明白的.
私下的語言的.
私事語的.
那怎樣能夠造就教會呢.
怎樣能夠幫助教會呢.
所以他如果這樣去對答說方言.
和先知講道的話.
他就寧願說.
先知講道的恩賜更加重要.

$^{441}$edify.
造就.
下一章在這裡.
說方言的是造就自己的.
作先知講道的造就教會.
我希望你們都說方言.
更加希望你們是作先知講道.
因為說方言的若不解釋出來.
使教會得造就.
那作先知講道的就比他強了.
整章的聖經經常都講到教會.
有九字.
講到教會.
講到我們聚集在一起的時候.
就是要做什麼.
我們要彼此去建立.
所以普羅話.
你們有什麼問題呢.
就是個個都想講方言.
個個都想私事語.
你們覺得這樣和上帝很契合.
他說no.
當你聚集一起敬拜的時候.
你需要去顧念整個群體.
需不需要得到幫助.
得不得到造就.
得不得到建立.
上兩個星期.
我們去MSI.
我們去博物館.
我們百多人去博物館.
這些活動很多時候看到.
有很多弟兄姊妹很感恩.
他們是去建立教會的.
有什麼新來的人.
我們去和他們聊聊天.
他們很明白那個活動是什麼意思.
但很多時候很sad的一件事.
有些人不是很明白.
那些活動是做什麼的.

$^{481}$他們自己在私事語.
自己在搞自己的事.
有些負責的弟兄姊妹說.
我們去拿飯的時候.
明明一個人.
一家三口.
他要拿四盒.
拿五盒.
可能他會偉大一點.
如果第一次來是可以的.
不要緊.
他會賜給你.
但如果是大家的弟兄姊妹.
你需要想一想.
造不造就教會.
只是你來到這裡的時候.
你覺得是serve你自己.
對不對.
很多時候我們都會這樣.
雖然我們未必有講方言的問題.
但很多時候我們所追求的就是.
我自己開心.
我自己滿足就夠.
普通話不是這樣的.
教會聚集在一起的時候.
我們需要other center.
focus on其他人.
來到的時候我想serve其他人.
這個就是基督有的精神.
基督有的生命應該是這樣的.
我們來到serve.
就是用我們的恩賜.
用我們的機會.
用我們的資源.
來建立教會.
教會不是屬於我.
是屬於上帝.
屬於基督的.
我們來到是貢獻.
幫助教會來成長.

$^{521}$過去大概八個星期.
我們在生命之道的查經裡面.
講完最後一堂.
講到一個例子.
我想分享給大家.
他講的例子就是.
以前有些教會.
現在都有.
傳統一些教會.
他們有些花在前面.
插花.
有沒有去過這些教會.
有.
不好意思我們這裡沒有.
因為很多時候.
我們以前去過那幾間教會.
插花那些又沒有人sign up.
那些花又謝了.
不是很榮耀神.
對不對.
anyways.
一來我太太又不喜歡花.
所以.
anyways.
有些教會是這樣的.
他就講了一個例子.
有一個姐妹.
她由開始教會的時候.
她就很喜歡插花.
插得很漂亮.
所以她就插了一盤花.
每個禮拜她就插一盤新的花.
奉獻給上帝.
做一個事奉.
但是久久的時候.
她就發現.
她的姐妹就擁有了這件事.
佔有了這件事.
覺得這件事是沒有人可以碰到的.
也不可以有其他人插花.

$^{561}$只是她.
其他人不可以碰.
這個就是.
不是一種actify.
我們作為事奉的人.
就是因為我們知道上帝給我們有不同的恩賜.
所以我們去用.
去造就其他人.
造就建立教會.
建立他人.
這個不是我自己的祭壇.
是其他人一起共享的.
特別事奉.
甚至做領袖.
甚至我們做牧師.
我們有時候都會這樣.
我們要將這件事佔有.
你可不可以想像這個人.
不可以有人講道.
只可以我講.
每次都只可以我講.
不能夠有其他人講.
更加不可以有比我講得更好的人去講.
你不能夠想像.
如果變成這樣的時候.
這個還是不是基督的教會.
這個就是光伯教會所面對的問題.
他說.
什麼.
他說如果你講方言.
若不翻譯出來的話.
教會就不能夠得到造就.
你所有的東西不能夠造就教會.
建立到其他人的話.
不如找一些東西建立到其他人.
所以保羅就說.
不如作先知講道.
因為他比他更強.
這個第一部分.
他說因此是做什麼.

$^{601}$教會聚集敬拜是做什麼.
就是Actify one another.
就是造就其他人.
第二部分就是20節至到40節.
我們選第23節來看.
這裡講到什麼.
就是整體教會用心去敬拜的時候.
我們能夠做到什麼.
23節.
所以全教會聚在一處的時候.
若都說方言.
偶然有不通方言的人.
或者不信的人進來.
豈不會說你們瘋了嗎.
瘋了.
若個個都作先知講道.
偶然有不信的或者不懂方言的人進來.
就會被眾人勸戒.
被眾人審問.
他心裡面的忍情就被顯露出來.
就必將面伏地.
敬拜上帝宣告說.
上帝真是在你們中間.
說什麼.
說多不多教會.
保羅說.
我希望你們個個都作先知講道.
因為人進來的時候.
他聽到你們做的事.
見到你們做的事.
是明白的.
他能夠明白是intelligent的.
是我一聽就明白的.
或者很容易懂的.
或者較為容易懂的.
聽得通的.
這個才是最重要的.
如果個個都講方言的.
誰聽得懂.
如果個個都是小圈子.

$^{641}$自己私私細語的.
怎可能造就教會.
所以他說要先知講道.
解釋清楚信仰.
解釋清楚上帝是怎樣的.
上帝的生命是怎樣的.
甚至一些未信主的人進來的時候.
他都能夠去understand.
我們崇拜所做的事.
其實都是一樣.
我們要做得一般.
讓人們可以理解.
當然有時我們講道.
或者一些topic.
可能是比較難去明白的.
我知道可能比較老年的人才會明白.
偶然會是這樣.
但in general.
大體上都是應該是這樣.
進來的時候.
我能夠理解.
我知道十教的救恩是一回事.
如果你去了教會.
好像我這樣去了幾十年.
你都明白.
在不同的世代.
崇拜其實好像是一種語言.
語言是會變的.
語言是會evolve的.
無論你講中文也好.
英文也好.
甚麼文也好.
不同的世代.
它的語言.
它的form.
或者形式.
它的language.
是會變的.
我以前去過一些很傳統的教會.
開始崇拜的時候.

$^{681}$有宣教.
然後有各個項目在裡面.
很多詳細的項目.
你去拿那些節目表.
那些節目表有十幾二十樣東西.
那些不是錯的東西.
不是說對與錯的問題.
而是我覺得那個language已經變了.
所以我們進來的時候.
我們不會印程序表給你.
為甚麼.
因為我們懶.
不是.
因為我們環保.
因為你印了程序表.
然後有很多垃圾在那裡.
每個人拿了之後就丟掉.
那個崇拜是沒有一個絕對的format的.
全部都是按照文化世代而去決定的.
剛才我們在初期教會.
他們崇拜的形式.
方法.
也跟我們今天是完全不同的.
And I'm sure.
他們用的音樂.
跟我們今天用的音樂.
也是完全不同的.
我忠告一些比較喜歡傳統.
或者死守傳統崇拜的人.
我忠告一件事.
就是使徒時代.
舊約時代他們敬拜.
用的音樂跟我們今天完全不同.
這個其實是很簡單的原則.
我覺得一年多前.
因為我要學希伯來文.
我就找了一些希伯來文的歌.
我就唱歌.
唱CB23篇就這樣唱出來.
是用一些希伯來的小調.

$^{721}$我想那些小調.
希伯來的歌應該是大衛.
以色列的文化.
中東文化他們用的音樂.
我就學那些歌.
背了它 我彈結他自己唱.
我可以記住那些字.
那個小朋友跟我說.
你不要再唱了.
很難聽.
我說什麼難聽.
不知道多好聽.
豈有此理.
他可能說歌曲難聽.
還是說我唱得難聽.
無論如何.
我想告訴你.
語言已經發展.
你所謂守傳統.
什麼是傳統.
如果我們要回到聖經.
所謂的律語.
我們就要全部唱希伯來小調.
你跟我一起唱好嗎.
都可以唱得好聽.
如果找一些懂得唱的.
他們應該唱得好聽.
但我想告訴你.
轉變了很多.
沒有一個絕對的形式.
甚至裝束都改變了.
你有沒有去過教會.
一定要打領帶.
好像吳敬虔一樣.
我想忠告一件事.
聖經所要求的.
跟我們自己的preference.
是兩碼子事.
對嗎.
我們自己的preference.

$^{761}$不要搞到好像是.
聖經所要求的一樣.
我講過這個故事.
我不知道你們聽過嗎.
聽過沒有.
九十年代的時候.
我去教會.
那時候我去教會事奉.
在執事會上聊天.
因為那時候教會沒有股.
沒有這套股.
就算有那套股.
都不放在台上.
不知道為什麼.
一定要放在下面.
執事會就在那裡談.
買股.
談得很熱烈.
有些人反對.
不想買.
搞了很久.
結果我就說.
算了.
不如我自己暗號班買.
結果我真的買了一套股.
那時候我買了一套股.
在加拿大.
這套股是屬於我的.
我就放在教會.
因為我家裡沒有位置.
解決了爭論.
為什麼要買那套股.
為什麼要給那麼多錢買這些東西.
因為我們又不喜歡.
因為他們的語言沒有那樣東西.
這個我明白.
特別是一些比較老一派的.
在七十年代長大.
八十年代長大.
他們覺得這些東西.

$^{801}$是死小子才做的.
壞小子才會打股.
不好意思.
我也是打股.
但是他們的語言還沒改變.
所以那個樂器.
都是一種改變.
舊有的傳統崇拜.
是要有股.
風琴.
管風琴.
甚至管風琴.
有管風琴的教會好不好.
ok.
不過你知道維修很貴.
維修很貴.
但現在的教會沒有這東西.
電子的鋼琴.
所有東西都是用電.
所以有分別.
我們需要隨著時代去改變.
所以這套股跟著我.
跟了我二十多年.
我去哪裡這套股跟著我.
這套股是我的.
不要亂搞.
每一件東西都是我的.
我很欣賞.
我們教會.
我們崇拜.
很多時候我們會唱一些新的歌.
我知道我們有時候要學新的歌是難的.
我們每個人都想唱舊歌.
社會學家告訴我們.
三十歲之後我們不學新歌了.
我們聽來聽去都是以前的歌.
很懷舊.
但我欣賞一件事.
我們有時候需要學新歌.
因為敬拜裡面需要唱新歌.

$^{841}$我意思不是新的一定是好.
但聖經不是叫我們.
是墨守成規.
敬拜是一個語言.
是需要和現代人溝通.
因為這樣才能造就教會.
特別是這段聖經.
說未信者進來的時候.
他知道你們在搞什麼.
如果他要搞一輪.
他才知道你們在做什麼.
可能我們需要再想一想.
因為保羅說敬拜是向外.
未信者進來的時候.
他能夠明白 聽得懂.
能夠知道你們在做什麼.
所以我們有什麼原則呢.
敬拜 聖經裡面告訴我們.
是有中向的 vertical.
是我們和上帝之間.
我們一群上帝的子民.
和上帝之間一個關係的回應.
上帝救贖我們 創造我們 愛我們.
所以我們在敬拜裡面向祂回應.
同時敬拜不是單單 vertical.
敬拜也是 horizontal.
是有橫面的.
是我們不能夠去漠視大家.
我們需要建立大家.
不能夠自私自利.
我們需要去建立大家.
並且甚至是向外去作見證.
人走進來這個教會的時候.
特別是新的人.
他們會在我們敬拜的當中.
在我們敬拜的音樂.
在我們敬拜的程序裡面.
聚會裡面.
能夠說上帝是有不同你們.
上帝真的在你們中間.

$^{881}$所以回到剛才我們的問題.
我們怎樣能夠成為一個.
更合上帝心意的敬拜者.
第一 我們需要用心.
Be mindful.
不是隨隨便便的時間.
我們就供入你.
因為我們來到上帝的面前.
我們需要有一個.
對的心態.
甚至是我們的身體.
都需要去預備.
身心靈需要去預備.
來到去敬拜.
如果我們污微合遂的時候.
我們怎樣能夠去見證上帝.
但如果我們很渴慕.
來到去敬拜上帝的時候.
那種能力就能夠出來.
甚至可以不開報道會.
有人進來感受到.
上帝真的在我們中間的時候.
那種生命的表露.
音樂的敬拜.
祈禱我們等等所做的事.
能夠告訴人聽.
這個上帝是一個掌管萬有.
掌管我們生命的上帝.
值得我們敬拜.
第二 我們需要去到遵從.
去遵從上帝.
這個是心態的問題.
來到上帝的面前.
到底祂有什麼為我預備.
我今天需要聽的.
需要遵守.
我的生命.
能夠慢慢地被祂改變.
第三 需要祈禱.
進來早一點點可以嗎.

$^{921}$我們需要清理一下我們的心靈.
我們說 上帝.
我希望你住在我們心靈裡面.
居於我心.
但我的心是不是乾淨.
可以讓上帝來居住.
或者我們有很多世界的東西.
追求的東西.
已經賺了上帝出來.
我們需要在上帝面前.
祈禱來預備.
有一個Clean Heart.
有沒有來到上帝面前.
有一點點時間.
祂說 上帝求你赦免我.
求你潔淨我.
讓我今天來到你面前.
是一個新的人.
花一點點時間去祈禱.
或者甚至為身邊的人祈禱.
赦免身邊的人.
以致我們整個群體.
能夠在祈禱.
在敬拜裡面.
和上帝相遇.
和上帝去改變.
我們和上帝去祈禱.
上主我們多謝你.
因為你邀請我們來敬拜你.
你創造我們和救贖我們.
用你愛你的恩典.
十字架的救恩.
能夠使我們成為一個新的身份.
成為你自己的子民.
能夠去敬拜你.
主我們今天來到你面前.
求你赦免和潔淨我們.
在我們生命裡面太多骯髒.
太多骯髒.
或者世界的事.

$^{961}$去佔據了我們.
再次要求你.
來到我們的生命當中.
居於我們的心.
充滿我們.
與我們同住.
與我們同在.
主讓教會來到你面前去敬拜.
我們能夠成為一個土地喜悅的敬拜者.
土地喜悅的敬拜群體.
以致人走到來這裡.
聽見我們讚美的聲音.
聽見我們禱告.
聽見我們向你述說自己的偉大.
他都會說上帝在我們中間.
求主提醒我們.
特別今天我們來到聖餐.
拿起這個餅的杯.
紀念你的救恩.
對我們的恩典.
拿起這個餅的杯.
去紀念耶穌基督為我們所做的.
完全的恩典.
以致我們能夠成為.
你的弟兄姊妹.
我們這樣祈禱奉祝美好.
\newpage



\section{}
\label{sec:Gu87k6iNMK0}
\textbf{Everlasting Praise 6 (Colossian 3) - Rev. Mark Chan}
\newline
\newline
連結: \href{https://youtube.com/watch?v=Gu87k6iNMK0}{\texttt{https://youtube.com/watch?v=Gu87k6iNMK0}} ~~~~ 語音日期: 2024-06-13
\newline
\newline
\hyperref[sec:mwhshrf5O44]{\small{< < < PREV SERMON < < <}}
~
\hyperref[sec:index_chronic]{\small{[返順時目]}}
~
\hyperref[sec:index_scriptual]{\small{[返順卷目]}}
~
\hyperref[sec:3ck_KAakvyM]{\small{> > > NEXT SERMON > > >}}
\newline
\newline
$^{1}$- Thank you, Phoebe..
Yeah, one more reminder is that we have a few workers,.
interns among us, and I wanna personally introduce them.
and wanna appreciate their work and ministry in the summer..
The first one, of course, is Phoebe..
So Phoebe is one of our summer intern..
She is finishing off her master's study.
at Lagos Evangelical Seminary,.
so she's doing some intern here..
And also Wendy, is Wendy around here?.
Back there, right?.
Yeah, Wendy..
Wendy is another intern, so we wanna thank her.
for doing all the background and backstage ministry for us..
And also Lucy, Lucy right there..
And thank you for joining us for this summer intern,.
and she is studying as Wendy as well.
at Moody Bible Institute,.
and she is music composition major..
So if you want some new songs, go to Lucy..
She will be helping us with the instrument.
and all the music as well..
So do pray for our interns, all right?.
We're going to put them to work,.
and they will work very, very hard, okay?.
So they will have a very, very busy month in the summer..
So let's pray together..
Father, give thanks as a community of faith,.
as your believers, we come before you..
We know that we are born in the spirit..
Thank you for the privilege and the blessing..
We can praise you, we can worship you,.
we can draw close to you..
Thank you for inviting us into your presence,.
and may your spirit speak to us.
and lead us into true worship..
And may all that we do, all that we think,.
and all that we say be pleasing to you..
We again come before you..
We ask that you will cleanse our sins,.

$^{41}$and you will forgive us..
You will help us to come before you,.
knowing that in Christ we are righteous,.
because all that your son has done for all of us..
Pray for our church as well..
During the summer, there are people in now,.
especially there are some brothers and sisters.
who are out of town, and I just pray.
that you will grant them safety..
I also wanna pray for those who are sick,.
who are going through some challenges and struggles..
Pray that you will grant them peace and healing..
As we open the scripture,.
we ask that your spirit will lead us into your truth,.
and let your truth set us free,.
and build up your church, and build up our life..
In Christ, we pray..
Well, we will continue our series on worship,.
and this is a seven-part series..
We are in part number six..
So we began this series asking the question,.
how can we be better worshipers?.
How can we be better worshipers?.
How can we become better worshipers as a group,.
as a church, as a community?.
And there are many questions we need to ask ourselves,.
or confront ourselves,.
because God is pleased when his people worship him..
And we have reminded all of us that from the Bible,.
from cover to cover, from the first chapter.
to the last chapter, in some sense,.
worship is all that the Bible is all about..
And what God is doing in humanity is that to rescue us,.
and to live in us, and to dwell in us..
And today, we're going to talk about.
some of the aspects of worship..
And we try to do something like a summary..
Try to find from the pages of Scripture.
what worship is all about,.
and what are those elements happen in corporate worship..

$^{81}$When we worship together, what are the elements,.
what are the things that we do,.
or things that the people in the Bible do when they worship,.
and to reflect upon how we can prepare ourselves,.
how we can become better worshipers..
So I have summarized some of the aspects.
or elements of worship from the Bible,.
especially from the New Testament,.
or also from the Old Testament as well..
From the whole Bible, try to summarize.
what are the main aspects of worship..
When we come before God to worship together,.
or when we call upon the name of Jesus to gather together,.
the New Testament reminds us.
whenever there were two or three people,.
calling in the name of Jesus,.
gathering in the name of Jesus,.
God is in the midst,.
God is dwelling among them in his presence..
So what did they do, right?.
What did they do?.
What do we do?.
And there are eight elements I have tried to summarize..
So in this part and the last part next week,.
we're going to go through this list.
and see what the Bible says.
about all these aspects of worship..
And first of all, it's about songs, right?.
We sing songs, right?.
We sing songs and we sing praises to our God.
when we gather to worship together..
No matter we gather in a church like this,.
like this kind of setting,.
or we gather into small group,.
or we even go to retreat, right?.
We get together as a small group..
We get together to worship..
And of course, the first thing that we do.
is that we sing songs, right?.
There are general principles from the Bible.

$^{121}$about using songs in worship..
However, the Bible never dictate the details, right?.
You agree?.
Does the Bible dictate what kind of songs that we sing?.
You think so?.
You know, I used to go to a reform seminary,.
and in the seminary, there were chapels every week..
There were about two chapels every week, morning chapels..
So we joined the morning chapels,.
and because it's a reform seminary.
and it's on the traditional side,.
and we always use the same hymn book..
And the hymn book, if you pick up the hymn book.
and you flip the pages, there are only 150 songs..
There are only 150 songs, right?.
That makes it easier, right?.
If you're leading worship,.
you don't have to go Google some of the songs, right?.
You don't have to go to YouTube.
and listen to some of the songs,.
because there are no new songs..
There are only 180 songs,.
because they are directly taken out.
from the bowl of salt, right?.
There are 150 songs only..
So that's the philosophy..
They only sing 150 songs, right?.
You agree with that kind of practice?.
Yes, no?.
Well, I don't think the Bible really dictates us.
exactly what kind of songs that we sing, right?.
The songs that we use are contemporary songs, right?.
They are not applicable,.
or they are not being used in other culture,.
in other setting..
Did I tell you last year I tried to learn some Hebrew songs?.
I did, right?.
Yeah..
Did I tell you last week?.
I'm losing my memory, yeah..

$^{161}$Right, and my children don't like Hebrew songs, right?.
The Hebrew songs are pretty much minor songs, right?.
But they're beautiful, right?.
But to some of the contemporary audience,.
don't speak to them, right?.
So we use some of the songs that are contemporary to us..
So we use those songs,.
and they are new songs compared to the bowl of psalms..
However, they are totally applicable..
It's biblical to use new songs,.
because the psalms, the Bible tells us.
whenever God does something new,.
we have to sing new song, right?.
So people will compose new songs to be used in worship,.
right?.
We have a major in music composition..
They're using her gifts, using our gifts,.
to compose entirely brand new songs for worshiping God..
So I don't think God will dictate.
what kind of songs we have to use..
No..
However, most of the songs that we use in corporate worships.
are what I call upward songs..
We sing those songs directly to God,.
directly to the Lord,.
and those are legitimate songs or poems that we can use,.
or lyrics singing directly to God..
Most of the songs that we use in worship are like that..
However, I don't think the Bible has ever dictated us.
what kind of songs or which songs that we use..
As you step into this church,.
you find that our setting is a little bit different..
I used to go to traditional churches..
This is not quite traditional setting for church buildings..
I wanna give you a quiz..
You know what is this?.
You know what this is?.
What kind of furniture item is this?.
Anyone?.
They are actually hymnal or hymn book holder..

$^{201}$You know, if you walk into a traditional church setting,.
you find those hymnal holder.
at the back of the church pillars, right?.
They are used for not your books..
You know, you don't bring your books and put it there..
It is for public worship..
So they use a particular hymnal, particular book of hymns..
And I used to go to a traditional Baptist church.
in California, and when I was younger,.
I always joined the worship service.
and the Wednesday evening prayer meeting..
And what we do is we gather together,.
and the host or the worship leader.
will tell us the number of the hymn..
So let's turn to our hymnal, 126..
And then we will turn to the hymn, hymn book, 126,.
and we sing together..
And all the songs that we use are all compiled.
in that book, right?.
But culture has changed, right?.
And we need to consider the local culture.
and language and tradition..
And there were flexibility,.
how we use or what songs that we use.
in corporate worship..
However, the main thing is there,.
the focus is to sing songs to God,.
it's to praise him with songs and poems..
And that's one of the aspects of public worship..
The second aspect is prayer..
Now, this is a no-brainer, right?.
Because in scripture, the people of God.
get together to worship..
There are many different prayers,.
or many types of prayers..
There are praises, there are petition,.
there are supplications, there are laments,.
there are also confessions of sin,.
expressions of repentance, expressions of thanksgiving..
All this wide variety of prayers are found in the Bible,.

$^{241}$and they are legitimately being used in corporate worship..
So sometimes when we pray together,.
or pray silently, or pray together,.
we confess our sin, and we ask God for his forgiveness..
Sometimes we give thanks to God..
We give thanks to God because God answered our prayer,.
and God has been good to us..
And sometimes we do petitions, right?.
We pray for our church, we pray for our community..
And this is one of the reminders that we have,.
is we summarize what worship is all about in the Bible,.
that we need to prepare ourselves, right?.
We come before God in the presence of God,.
we want to worship, we need to prepare ourselves,.
we learn the songs, and we prepare our heart for prayers..
Did you do that today?.
You walk into this hall, this auditorium,.
you prepare yourself to be prayerful..
You pray to God, you speak to God, right?.
Either you have thanksgiving,.
or you have some confessions to make,.
or some praises you want to sing, right?.
We need to prepare ourselves, right?.
Not only you prepare your body,.
but you prepare your spirit.
so that you can focus on worshiping God..
Speaking of prepare on body,.
I have not been sleeping so well in the last two weeks..
Especially about three weeks ago,.
I tried one of those over-the-counter sleeping pills,.
because I have trouble sleeping..
So I tried some new sleeping aids three weeks ago..
I took one on Tuesday, and of course I slept,.
but I collapsed, I actually collapsed..
And then I woke up Wednesday morning,.
I feel so bad, feel so bad..
I felt like I never slept last night, the night before..
And then it lingered for a few days, actually..
I've been just sleeping all day, sleepy all day, right?.
If you tell me, you ask me,.

$^{281}$if I took the new sleeping pill,.
like Saturday night, it would not be wise, right?.
Because I want to prepare my body.
so that I can worship God..
So I ditched those new sleeping pill,.
and I go back to my, the old ones, the prescribed one..
And it helps me to rest,.
help me to prepare our physical body to worship God..
And this sort of thing reminds us.
we need to prepare ourselves to worship..
We cannot just bump into a worship service,.
and then, okay, all right, we're worshiping..
Right, God wants us to prepare ourselves,.
our spirit, our mind, and our body to worship Him..
That's how we become better worshipers, right?.
Number three is confession of faith..
Now, this sort of things sometimes we do.
in corporate worship..
And what does it mean, confession of faith, or creed?.
It's from Latin, credo..
It means I believe..
You know, throughout the entire church history.
of the last 2,000 years,.
there are generations of theologians and believers,.
they summarize the basic doctrines.
of the Bible..
As they summarize it, they will write some of the creeds..
Now, one of those creeds that we use in the past.
is called the Apostles' Creed, right?.
I actually preach about it.
in a sermon series on Apostles' Creed..
And Apostles' Creed has been around.
for more than 1,000 years..
And it's a very, very valuable and ancient summary.
of Christian faith..
And it summarizes our basic doctrines in 12 articles..
And we recite the Apostles' Creed.
to remind us what we believe, what we stand on..
And that is called the creed or confession of faith..
Actually, those creeds we can find from the Bible..

$^{321}$If you pay attention to some of the scripture.
or some of the passages in the Bible,.
they are actually being incorporated.
into the New Testament writing,.
into the New Testament Bible..
For example, the book of Matthew, Matthew 6..
The Lord's Prayer, right?.
Our Father in heaven, hallowed be your name..
Be thy name..
Thy kingdom come, thy will be done.
on earth as it is in heaven..
Give us today our daily bread..
Forgive us our debts..
And we also have given our debtors..
And lead us not into temptation,.
but deliver us from the evil one..
For thine is the kingdom, the power,.
and the glory forever, amen..
As I read this, I used King James words.
because that's how I grew up..
I recite this Lord's Prayer from time to time.
in worship service, right?.
This is the confession of faith..
Now, how do we know?.
You know the last line I highlight here?.
If you read the details of the English Bible, for example,.
you find a footnote there,.
and the footnote tells you the early,.
early Christian manuscript of Matthew 6.
does not have the sentence highlighted..
Only the later manuscript contains the highlight,.
a second part of verse 13..
Now, you know, there are many manuscripts of New Testament,.
and what the scholar's doing is trying to summarize.
what the original text actually is..
So if you go to the earlier or the best manuscript,.
it doesn't have the yellow highlight verse..
It stops at the evil one..
Deliver us from the evil one, and then it stops..
And that's the Lord's Prayer..

$^{361}$That's the original Lord's Prayer..
That's how Jesus taught His disciples to pray, right?.
It's a series of petition..
Leads us not into temptation, forgive us,.
and deliver us from the evil one..
The last line is being added in the early Christian church,.
and that's why those who copy the Bible.
from generation to generation, the Christian scribes,.
they incorporate this last verse,.
this last second part of verse 13 into the Lord's Prayer..
And why?.
Because the early churches, when they met in house churches,.
they add this phrase at the end of the Lord's Prayer..
And this is how they use it,.
use the Lord's Prayer in their worship service..
So end with some sort of doxology, some sort of praise..
Yours is the kingdom and the power.
and the glory forever, right?.
So as one of the example, you find that, you know,.
this kind of confession of faith using the Lord's Prayer,.
it's a legitimate use of scripture,.
and it happens to be used in worship service..
That's why sometimes we do that too, right?.
A brother just remind me,.
why don't we do something like that in our worship service?.
You know, recite some of the scriptures,.
like the Lord's Prayer, and memorize it, right?.
Because when we memorize the Word of God,.
the Word of God will take effect in our own life, right?.
That's a great idea, and we will continue to do that, right?.
So there's confession of faith..
There's also another aspect called vows,.
or some sort of response or a commitment.
that we make before God in worship..
In the Old Testament and New Testament,.
we find many, many different vows or promises,.
commitments that people make to God.
to fulfill those promises..
Don't you experience that?.
Sometimes you're in your spiritual walk with God..

$^{401}$You pray for a certain thing,.
and then God answer your prayer,.
or God open a way for you, right?.
And you made some sort of promise to God..
You pray a certain thing in front of God,.
and you say, "God, if you answer me,.
"if you do this for me, for your glory,.
"then I will fulfill my vow, right?.
"I will give you thanks..
"I will give you offering," or something like that, right?.
I have done so many times in my Christian walk, right?.
God provide me, you know?.
If you answer my prayer, I will recommit myself to you,.
and that is what vow means..
That is what commitment means..
I remember years, years ago, I was in seminary,.
and I finished the first two terms,.
and then I ran out of money..
I just didn't have money to pay for tuition,.
so I prayed..
Right at that moment, I was in a really, really big struggle.
because I probably have to quit,.
have to find a job, right?.
And to save up more money so that I can pay for tuition..
So I prayed, I prayed..
So as I prayed, and I went to prayer meeting,.
and we prayed together..
Now I remember the following week,.
I went to a prayer meeting..
I went to a worship service, and I came home,.
and I find a white envelope in my mailbox..
I was kind of curious, okay, who put it there, right?.
So I grabbed the envelope, and I went to my bedroom,.
and I opened the envelope, and there was a check,.
and that check has the amount just enough.
for me to pay for the next term, right?.
At that moment, I fell on my knees..
I said, "God, thank you..
"Thank you for answering my prayer..
"Thank you for accepting me because I was of little faith,.

$^{441}$"right, and God, you are the great provider,.
"and you are so faithful to me," right?.
So I fell on my knees, and I prayed, I vowed..
You know, "If God, you are so faithful,.
"help me in my faithlessness, right?.
"Help me with my lack of trust..
"Help me to trust in you.".
So there we go, I fulfilled my vow, right?.
In my years of ministry, there were up times,.
there were many down times..
There are mountain tops, a few mountain tops..
There are many valleys, but I remember this experience..
I made a vow to God that God is faithful,.
and I need to continue to trust in him, right?.
Whether the times is good or the time is difficult..
That should be part of our journey as well, right?.
God grant me this prayer..
I will make a vow to you, right?.
I will recommit myself, and that's part of our worship too..
And in the Old Testament, there was a story about Hannah,.
right, the woman who was barren..
She didn't have any child, and she was suffering nobility..
She was being locked down upon..
And Hannah prayed to God in the tabernacle,.
and prayed that God grant me a child, right?.
And lo and behold, God grant her a child, Samuel..
And the next year, she went back to the temple,.
the tabernacle, with her husband, Elkina,.
and his family, went on their annual trip.
to offer worship sacrifice to the Lord, and keep his vow..
And there you go, vow right there..
The vow actually is Hannah's own vow..
So Hannah prayed to God, grant me a child..
If you really grant me a child, I will offer this child.
for your service, for the Lord's ministry..
And she did that..
She brought Samuel, the boy, to the temple, or tabernacle,.
and offered Samuel as a minister to God..
So vows and promises is part of our worship service..
Now, we will continue this list next week..

$^{481}$Preaching and teaching and sacraments,.
and we will explain what sacraments are,.
and then offerings and fellowship..
Now, I want to close this sermon by saying,.
all of these aspects and all of these elements.
are part of legitimate worship..
And God is pleased when His people gather.
in His presence to worship Him, right?.
Remember this, worship is about God, it's all about God..
And worship in the Bible has two dimensions..
The first dimension is a vertical dimension..
We worship and we respond to what Jesus has done.
in our own lives..
He save us for the purpose of worshiping Him..
And that's a vertical dimension..
And the other dimension is horizontal dimension..
Worship is to edify, is to build up others in the community..
Worship is not for me to have fun..
I come here, I am present here,.
so that I can practice community..
I can build up the body of Christ..
I can have fellowship and share life with other people..
And our worship can be a witness to the world outside..
Let's remember that, right?.
We're gonna close in prayer..
And let's pray before God..
How can we become a better worshiper?.
All right, to glorify Him..
Let's all pray..
I'll give you some seconds and you can pray to God.
in silence..
If God has led you to repent,.
ask for forgiveness..
No, God's grace is always sufficient..
If God is calling you to recommit,.
do so, because God is good, worthy of our worship..
(gentle piano music).
Father, thank you for being our loving God..
Thank you for loving us so much..
And call us into your presence,.

$^{521}$so that we can be your people..
So we pray that you will help us.
to become better worshipers,.
to honor you and to glorify you..
Forgive our shortcomings..
Sometimes we don't prepare..
If we don't prepare our hearts and our body.
to serve you, to worship you,.
may you forgive us.
and help us to put you as our life's pyrode..
To worship you is to bring you glory..
Ask your people, ask your church,.
help us as a community to worship you,.
to honor you in Christ's name we pray..
Well, God bless, and next week we're going to learn.
how to be better worshipers, right?.
So join us next week..
God bless..
[MUSIC PLAYING].
\newpage



\section{歌羅西書}
\label{sec:3ck_KAakvyM}
\textbf{永恆的敬拜 5 (歌羅西書 3) - 陳崇基牧師}
\newline
\newline
連結: \href{https://youtube.com/watch?v=3ck_KAakvyM}{\texttt{https://youtube.com/watch?v=3ck\_KAakvyM}} ~~~~ 語音日期: 2024-06-13
\newline
\newline
\hyperref[sec:Gu87k6iNMK0]{\small{< < < PREV SERMON < < <}}
~
\hyperref[sec:index_chronic]{\small{[返順時目]}}
~
\hyperref[sec:index_scriptual]{\small{[返順卷目]}}
~
\hyperref[sec:xTyYHQxRMCA]{\small{> > > NEXT SERMON > > >}}
\newline
\newline
$^{1}$我們繼續看這個系列是講到崇拜的.
正如我所說 我們來到教會.
我們作為一個教會 作為一個基督徒.
久了之後我們要了解崇拜這件事.
主日崇拜是一個不能夠取代的環節.
在我們生命裡面 在我們基督徒的生活裡面.
這是一個最重要的日子 或者最重要的時間.
所以我們需要花多點時間去了解到底崇拜是一回事.
所以今天我們是第六講.
下一個星期我們第七講就會完結崇拜這方面.
我們始終都是回到這個問題.
就是我怎樣能夠成為一個更合上帝心意的敬拜者.
我們每一個都是敬拜者.
不只是台上帶領的人才能去敬拜.
他只是帶領我們以至整個會眾能夠去敬拜主.
聖經裡面由頭到尾都是告訴我們.
上帝所關注的東西就是在他的子民當中.
你都要跟他們同住.
Walking in the midst of them.
Walking among them.
Living among them.
這是上帝由頭到尾有的心意.
所以敬拜其實就是在上帝的Presence.
上帝的存在 上帝的同在底下.
我們去敬拜上帝.
跟他建立這個關係.
所以主要是敬拜是非常非常重要.
最重要最重要的那件事.
等於你每一天三餐哪一餐最重要.
哪一餐最重要.
應該是早餐吧 有沒有吃早餐.
有些人是晚餐最重要的.
我知道有些朋友會去多倫多.
我不知道他哪一餐最重要.
一天吃七餐.
哪一餐最重要.
在我們基督徒的生活裡面.
最重要的一餐就是主日.
就是這個時候.
因為上帝讓他的子民.

$^{41}$能夠在他的同在裡面去敬拜.
去到聚會的時候.
上帝的心才會喜悅的.
上帝希望能夠住在敬拜的子民的群體裡面.
所以今天我們在最後這兩個星期.
我們嘗試去做一些整合的事.
就是整合一些事.
就是整本聖經裡面怎樣講敬拜.
或者敬拜裡面有些什麼元素等等.
所以我們在這裡.
大概去將它總結起來.
整本聖經裡面由頭到尾.
特別是新約裡面怎樣說.
教會的敬拜應該有些什麼元素.
我知道這些東西.
很多時候是一些敬拜事工的工作坊.
一些workshop我們才會講的.
但我自己想.
當然這個是需要的.
因為帶領敬拜的人需要明白這些東西.
但是正如我所講.
我們每一個都是敬拜者.
不只是那班台上的人.
或者主領的人才帶領我們.
他們才會敬拜的.
我們整個上帝的子民.
整個教會都需要了解什麼叫做敬拜.
所以如果你說要做一個合上帝心意.
完了之後.
希望最主要的是.
上帝覺得我今天所獻上的.
我今天的心靈.
我整個的心靈.
是上帝會喜悅的.
我都需要了解一下.
整個聖經裡面怎樣講敬拜.
這裡我大概去到將它總結了.
大概有六七樣.
有多少樣.
八樣左右.

$^{81}$八樣.
或多或少.
聖經裡面描述敬拜.
或者要求的敬拜會是怎樣.
正如我所講.
聖經裡面有一些很general的principles.
一些很一般的整體的原則.
但是很多時候.
那些細節.
就在不同世代裡面的基督徒.
不同的教會裡面.
他們自己嘗試去figure out.
不同的時代.
不同的背景.
不同的文化.
都會有它不同的表達.
但是最主要來說.
都是有這幾方面.
或者有些這樣的elements.
aspects.
是講到敬拜的.
準備好了嗎.
第一方面當然一定是有詩歌.
一定有詩歌.
詩歌是敬拜裡面其中一個最重要的東西.
但是聖經裡面有沒有講.
細節.
說一定要用甚麼詩歌.
有沒有 你覺得.
聖經裡面告訴我們.
必須要唱哪一首.
不能夠唱哪一首.
有沒有.
沒有吧.
聖經沒有講.
當然大部分的敬拜的詩歌.
都是所謂向上的詩歌.
我們唱那首詩歌.
是直接唱給上帝聽的.
所以那些歌詞.

$^{121}$或者那些poem.
詩句.
都是比較vertical.
比較是中面的.
直接向上帝唱歌.
好像我們剛才唱的那些詩歌.
是直接向上主表達的.
這些詩歌.
通常在敬拜裡面使用.
直接唱給上帝.
但是也有些教會.
有些傳統.
我不知道你有沒有去過這樣的教會.
以前我去神學院的時候.
我去一間改革中神學院.
我們有組會的.
每個禮拜兩次.
我沒有記錯吧.
因為有時候我會不去.
所以.
去兩次組會.
組會就是會有一起敬拜的時候.
那間的群體.
他們的詩歌.
只有150首.
就是詩篇.
詩篇有多少首.
150首嘛.
所以他們坐下來.
我們去敬拜的時候.
我們就拿起那本詩本.
全部都是.
詩篇裡面的.
逐字逐句地寫下來.
有沒有看過詩篇.
如果看過詩篇.
你唱不唱得出來.
你唱不到吧.
他就是用那150首的詩歌的歌詞.
配上音樂.

$^{161}$所以他由頭到尾.
一年到晚.
52個禮拜.
都是唱過150首.
因為他們認為.
聖經裡面只是告訴我們.
唱過150首.
所以你去到那個敬拜的時候.
其實你學的歌只有150首.
是不是很好.
學少很多東西.
是不是.
你一坐下來.
詩篇第一篇.
原來是這樣唱的.
那你就這樣唱.
是不是好呢.
I don't know.
我想聖經裡面沒有這樣detail.
或者去到決定.
我們一定要唱過150首.
150首以外的就不能夠唱.
因為詩篇裡面有講很多很多次.
什麼.
耶和華要向耶和華唱新歌.
每逢上帝有新的工作的時候.
就會有新的詩歌出現.
是不是.
如果你再這樣想.
聖經裡面敬拜的詩歌.
都不是只是150首.
是不是.
出埃及之後.
以色列人向上帝敬拜.
有摩西紅海之歌.
哈拿等等不同的大衛.
他們有很多不同的詩歌.
他們會唱的.
所以不只是限制於150首.
正如我們新的時代.

$^{201}$我們唱詩歌的方法.
也可能會有不同.
我想考考你.
教會裡面傳統的家詩.
看看你知不知道是什麼.
你知不知道這個是什麼.
是什麼.
那個架.
放什麼的.
放聖詩的.
那些新的弟兄姊妹.
初初接觸信仰.
你就是來UV的.
來我們教會的.
你就發現這些東西是什麼.
因為我們沒有這些東西.
這些是傳統的教會後面.
很大張很長的牆椅.
木的牆椅.
後面一個架.
那個架是放聖詩的.
他們就是用那本聖詩.
無論是頌主新歌.
或者我們比較老派的信徒.
你記不記得什麼叫做青年聖歌.
有三種顏色的.
記不記得.
唱青年聖歌的很多都已經去天家了.
真的.
以前我也唱青年聖歌.
頌主新歌.
唱頌主新歌的都已經唱到很久了.
已經不再是新歌了.
當然這些傳統詩歌都有它的價值.
我偶爾會聽這些傳統詩歌.
或者去一些不同的教會聚會的時候.
我會聽見這些詩歌.
我都朗朗上口.
我都發現原來我都不是老人之外.
我記得以前的東西.

$^{241}$新的東西不記得.
以前的東西記得很清楚.
照樣繼續唱.
因為以前我在加州的浸會.
每個星期三晚我們有祈禱會.
祈禱會我們就唱很多詩歌.
那些詩歌全部都是在詩本的謙語那裡拿出來的.
所以一講那個號碼.
我就知道第43首就是那一首了.
這些是舊有的傳統.
正如我所說.
教會的任務.
教會在地上的使命.
並非要保留傳統的.
贊成嗎.
我們要去到履行上帝給我們福音的大使命.
這個才是教會有的任務.
教會有的任務不是保住以前有的東西.
對不對.
但最重要的就是我們要用詩歌來敬拜.
下一次我們會講多一點.
worship這件事.
或者worship這個語言.
但怎樣都好.
第一個元素就是詩歌.
第二個元素當然不需要多講.
就是祈禱.
譬如我們有時候會用木土圍著教會去禱告.
有不同各種各樣的祈禱.
有讚美的祈禱.
有代求的祈禱.
有哀歌.
有認罪悔改的祈禱.
是不是.
有迴轉的祈禱.
有感恩的祈禱等等.
這些不只是詩篇裡面.
詩歌裡面找到.
在聖經裡面舊約新約都會找到的.
來到教會.

$^{281}$來到崇拜的時候.
我們就自然會獻上我們的禱告.
所以你是一個敬拜者.
你就是一個禱告者.
你來到上帝面前.
你要澄明你自己心裡面的心事.
這個都要做準備的.
你走進來教會去敬拜的時候.
你都需要去處理一下你心裡面.
可能雜亂無章.
可能過去一個禮拜你覺得很辛苦.
覺得很難.
家庭裡面有很多的挑戰.
工作裡面等等.
生活裡面等等.
都需要整頓一下.
來到上帝面前來禱告.
我發現最好的一個練習.
就是早一點來教會.
坐下.
五分鐘都好.
可以整頓一下.
對不對.
剛才我們敬拜到練歌的時候.
我們都有些混亂.
我們都要去談一談.
想一想.
所謂電腦有一種東西叫reset.
reset一下.
到底我來這裡做什麼.
搞清楚.
這個都是我們每一個敬拜者需要有的準備.
祈禱.
第三方面就是Confession of Faith.
或者是Creed.
什麼意思呢.
Creed.
拉丁文的意思就是Credo.
就是I believe的意思.
就是說我相信.

$^{321}$所以在敬拜裡面.
有時候會有一種叫做.
供讀信經的環節.
什麼叫供讀信經.
供讀信經就是說整個的群體.
一起來再表述一次.
我們所相信的信仰.
我們的真理是怎麼樣的.
叫做Confession of Faith.
Confessing of Faith.
所以有時候我們會做什麼.
使徒信經.
是不是.
你記得參加過我們某些聚會.
我們一起站起來讀使徒信經.
使徒信經差不多有2000年的歷史.
是歷世歷代的信徒.
去將整個的信仰.
來作總結.
總結出12句的statement.
這個叫做使徒信經.
也有尼西亞信經.
三一信經等等這些信經.
這些信經都是有它的.
應該要保存.
或者應該要去到崇拜裡面使用的.
一個的習慣.
因為我們敬拜的時候.
不只是我們自己敬拜.
我們是集體的敬拜.
並且這個集體的敬拜.
不是現在集體的敬拜.
而是歷世歷代的信徒.
一起來去敬拜.
所以我們記起.
我們不是那種失憶的人.
過去裡面有很多很多.
歷世歷代的信徒.
是和我們組合在一起.
一個群體來去敬拜主.

$^{361}$所以有時候我們會做一些.
公讀信經的環節.
或者其中一個所謂confession of faith.
就是宣告我們的信仰的環節.
就是水禮.
我們浸禮水禮的時候.
那個人去到作見證.
其實都是一種confession of faith.
去到宣認上主是他的君王.
是他生命裡面的主.
所以在敬拜裡面.
是有這樣的環節的.
事實上在聖經裡面.
是有這樣的東西的.
很多時候我們未必會留意.
譬如一個example.
就是主禱文.
主禱文大家會背吧.
我們在天上的父愛.
願人都尊你之名為聖.
願你的國光臨.
願你的旨意行在地上.
如同行在天上.
我們任由天飲食.
今日賜給我們.
免我們的債.
與我們免了人的債.
不叫我們陷入試探.
救我們脫離那惡者.
因為國度權.
並永要傳是你的.
直到永遠.
阿們.
這個是一個信經來的.
或者一個confession of faith.
我們怎麼知道呢.
怎麼知道初期教會.
是用主禱文.
在他們敬拜裡面.
當作一種所謂的liligy.

$^{401}$就是一種敬拜的禮儀裡面.
一個環節.
為什麼呢.
你見到黃色那句.
如果你有以前印出來的聖經.
或者今天可能電子聖經都有.
它是括著說.
古卷裡面.
早期的古卷才有這一句東西的.
沒有這一句東西的.
講清楚.
最初的古卷.
是沒有黃色這一句.
後期的古卷.
才有黃色那句.
你說說什麼.
你知道以前的聖經是抄出來的.
抄出來的.
就是有一個原版.
然後就有一些copyist.
或者scribes.
文士.
他們就會抄它.
抄完之後就再抄.
你說抄來抄去.
抄錯了.
你以前小學的時候被罰抄的時候.
是不是亂抄.
亂抄通.
但是我想告訴你.
是不會的.
因為那些文士是很小心的.
死海古卷發現的時候.
為什麼這麼大.
為什麼這麼震撼.
因為原來發現.
一拿死海古卷.
那些山洞古卷拿出來的時候.
就發現.
咦.

$^{441}$原來那個希伯來的聖經.
舊約的聖經.
那些抄本.
是和一千年之後.
我們之前只知道一千年之後的抄本.
一對一對的時候.
原來百分之九十幾.
都是保留著的.
你說剩下那幾個百分比去了哪裡.
那些只不過是自己更新.
字已經串法不同了.
或者一些音標上的東西.
或者一些很仔細的東西.
他們把它更新了.
就是更新了.
但是百分之九十幾都是一樣的.
所以就這麼震撼.
所以那些文士去抄聖經的時候.
是很小心的.
就是你抄了之後.
然後又對一對.
到底有多少個生字在這裡.
有多少個字母在這裡等等.
全部東西都是檢查好.
就好像你今天寫的那些電腦的那些碼子.
你要compile它.
compile有沒有bug.
如果有bug的話.
你就反過來又要再debug它.
直到你compile到沒有了error為止.
所以那些文士在抄本的右邊.
或者左邊那些空位的地方.
他們就會有一個註腳的人士.
這節經文.
最後那個字是什麼字.
所以就不會有什麼錯的.
無論如何.
這裡的主禱文.
說最初的抄本是沒有第一句.
後期的抄本才有黃色這一句.

$^{481}$意思是什麼呢.
意思就是說似乎初期的教會.
是用這個主禱文.
來作他們聚會裡面其中一個的禮儀.
其中一個的環節.
所以他們加了這一句下去.
For yours is the kingdom and the power and the glory forever.
Amen.
聽起來的時候你就知道是一個liturgy.
好像一首歌一樣.
有一個Amen完結了.
原本的主禱文.
可能只有白色這一句而已.
就是求上帝救我們脫離斯坦.
救我們脫離那惡者就完了.
但是初期教會在應用這一段經文.
耶穌怎麼教會和禱告經文.
他們就放在他們聚會的環節裡面.
所以加了這個.
因為國度權並要傳遞的.
直到永遠.
Amen.
這個環節叫做什麼.
Confession of faith.
就是眾信徒一起聚集的時候.
我們來到Confess.
我們來到表達清楚.
公共的表達清楚.
我們所信仰的是什麼.
我們的上帝是怎樣的.
我們所尋求的東西是應當是怎樣的.
這個就是公讀信經的環節.
然後下一個我們能夠總結出來的.
就是還願或者許願.
Vows.
或者Commitment.
這個可能比較少提.
但是我們平常都會做的.
就是回應.
一個回應的禱告.

$^{521}$去聽見上帝的吩咐.
我們去回應祂.
是一個response.
一個commitment.
在舊約新約.
特別是詩篇.
也有很多這樣的還願的禱告.
就是私人在上帝面前.
來到去決定做某一樣東西.
上帝的愛或者上帝的恩典.
去觸動了他之後.
他的心靈說不行.
我不可以就這樣坐在這裡.
我要許願.
我要告訴上帝.
我願意盡力的為你做什麼.
我要去奉獻等等.
去做一個commitment.
然後見到上帝帶領他.
去應允了他的禱告之後.
他就會來到上帝的面前.
又回來了.
他就說有一個還願.
這個就是敬拜裡面.
讓我們知道.
我們不只是向上帝唱詩.
其實上帝在我們心靈裡面.
去做一些事情.
然後我們去做一些回應.
在舊約裡面有一個人物叫做哈拿.
哈拿大家知道這個婦女.
是沒有孩子的.
於是被人欺凌.
她在生命裡面很苦很慘.
於是她去到會幕.
來到去禱告的時候.
她就很苦的.
來到上帝面前陳明她的苦情.
然後完了禱告的時候.
她就有一種釋放.

$^{561}$她就希望上帝給她一個兒子.
完了禱告的時候.
她就有一種釋放.
然後她就回家.
在禱告裡面.
她就有一個這樣的許願.
如果上帝真的應允我禱告的話.
我必定會將這個孩子.
來到去獻上.
事奉上帝.
第一章第21節.
撒布爾結常.
他說以利加拿.
以利加拿就是她的老公.
和他全家都上去了.
要向耶和華獻蓮祭和還願祭.
那個孩子出生了.
她就回去獻還願祭.
詩篇裡面.
很多很多這樣的祈禱.
都是說上帝在我的生命裡面.
既然做了事情.
fulfill了祂的purpose.
祂的功能.
我就來到去還願了.
教會在公共崇拜裡面.
其實都是應該要這樣的.
因為崇拜好像一個對話.
不是一個monologue.
不是上帝跟我們說話.
然後我們就走了.
而是我們跟上帝有一個對話.
一個回應.
所以在公共崇拜裡面.
是必須要有這個環節.
今天上帝跟你說的話.
給你大小的提醒.
你都需要有一個願.
response.
你跟上帝說.

$^{601}$Yes.
我願意.
我願意去馴服你.
我願意去聽你的說話.
願意在行動上面.
我做一些回應.
這個才是真正的敬拜.
對不對.
除了還願之外.
還有什麼.
就是講道和教道.
這個不需要多說.
因為聖經裡面.
從頭到尾都告訴我們.
神的子民聚集.
舊約到新約.
都是要聽上帝的話語的.
特別是舊約裡面.
雖然我們看到.
有很多上帝出現.
有光顯現.
有暴風來到.
地震盪.
好像西乃山一樣.
但最重要的是什麼.
不是看的東西.
最重要的是聽.
所以上帝在西乃山上面.
他榮耀加臨的時候.
很震驚.
但最重要的目的是什麼.
他說你要聽我的說話.
You have to listen and obey.
所以是audio.
不是visual.
最重要是audio.
新約也是一樣.
每逢教會去到敬拜的時候.
他們就會留心使徒的教導.
和聖經的教導.

$^{641}$拜道威新約裡面的書信.
特別是書信.
或者新約裡面任何一卷書.
其實寫來做什麼.
寫來就是給教會傳閱的.
譬如寫的作者.
是在羅馬書信.
保羅在羅馬寫信.
寫完信他就叫一個信差.
其中一個弟兄帶他去菲臘比宣傳.
帶到菲臘比.
菲臘比主日的崇拜.
當然是在家裡面崇拜.
很少的.
一個小組的教會幾個人.
十幾個人.
他們擠迫在一個廳裡面.
然後他們就聽保羅的書信.
就有個人拿著那個書信.
然後就讀出來.
菲臘比的教會.
你們應該怎樣怎樣怎樣做.
讀完之後.
他們會再翻讀.
但是他們會做一個複製.
然後就傳遞.
好像以前的圖書館那樣.
他們就傳到其他的教會.
然後其他的教會.
雖然不是在菲臘比.
他們就會聽菲臘比的書信.
然後那個書信一直這樣.
繼續傳授下去的時候.
大家的教會就會聽到.
保羅對菲臘比的教導.
所有書信都是一樣.
就是講到公督上帝的說話.
歌羅西書第三章第十五節.
你們要讓基督所賜的和平.
在你們心裡面作主.

$^{681}$也為此蒙召歸為一體.
當用各樣的智慧.
把基督的道.
豐豐富富的存在心裡面.
留意這裡.
用詩篇讚美詩.
靈歌彼此教導.
互相勸誡.
以感恩的心歌頌上帝.
無論作什麼說話或行事.
要奉耶穌的名.
這裡講到剛才的歌.
唱歌.
但也要彼此教導.
用不同的方法來彼此教導.
其中一個方法就是唱歌彼此教導.
唱歌如何教導?.
剛才我們一起唱詩的時候.
有沒有做到教導這方面的功能?.
也有的.
唱歌彼此祝福大家.
知道上帝是真神.
但不要亂唱.
唱我愛你的時候就不行.
因為你是上主.
不是你.
是上主.
但詩歌有一個訊息.
彼此的教導.
剛才我們說書信.
試圖去教導.
也是一種教導或講道的功能.
所以在基督的教會裡面.
主日崇拜裡面.
很多時候有講道這個環節.
我們作為講道的事奉的.
我們都虛心在上帝面前先學習.
然後我們帶給大家上帝的訊息.
在教會裡面.
所以講道的環節是一個中心的環節.

$^{721}$又給大家看看.
你知道中世紀的教會.
如果你去歐洲玩.
你去那些中世紀的教堂.
那些大教堂.
很漂亮的那些大教堂.
你走進去的時候.
中間的傢俬是哪一件傢俬.
你有沒有留意.
是建築師.
中間的傢俬.
就是港台或祭壇上.
中間那件東西是哪一件.
你有沒有留意.
通常是聖餐桌.
或者是祭壇.
港台在哪裡.
港台在左邊.
就是你的左邊.
你有沒有留意.
有些傳統的教會.
有些中世紀的教會.
哥德式的教會等等.
港台在旁邊.
為什麼呢.
因為中世紀的教會.
很多時候是比較visual一點的.
他們的禮儀.
很多時候是所謂的彌撒.
或者祭司的牧者.
他們做的事情就是去重演崇拜等等的禮儀.
所以很多時候是看的東西.
不是聽的.
你進去那些中世紀的教會.
你也知道不是讓你聽到的.
好像山洞一樣.
說一句話.
嗡嗡嗡嗡嗡嗡.
我記得耶路撒冷有一個教堂很出名.
叫St. Anne's Church.

$^{761}$St. Anne.
Anna.
是誰呢.
是馬尼亞的媽媽叫Anna.
因為那些《刺經》說馬尼亞的媽媽叫Anna.
有一間叫St. Anne's Church.
為什麼那麼出名呢.
他說那裡是最perfect的echo.
在那個教堂.
你去唱詩或者說話.
是一個最perfect的回音.
就是七秒鐘.
七秒鐘.
你說「豬」.
就是七秒鐘.
所以很多人去到那個教堂.
他們就會去唱歌.
在那裡唱歌.
那個唱歌簡直是一流.
但是如果是港獨就麻煩了.
因為港獨不是唱歌的.
所以那些教堂本身不是設計給港獨的.
這個是中世紀的傳統.
直到reformation.
宗教改革的時候.
大概五百年前.
當時馬丁路德,嘉爾文等等不同的神學家.
他們說我們要回歸聖經.
不可以只是跟傳統.
因為當時整個的中世紀.
很多時候就是跟傳統然後腐敗.
他們說我們要return to the Bible.
唯獨聖經.
聖經的權威.
所以由那裡開始.
所謂改革中的教會或者新教.
或者我們這個傳統.
有五百年多的時間.
我們就是一種這樣的改革.
於是你進入到改革中的教會.

$^{801}$或者新教的教會.
那個教堂的建築的設計.
就有所不同.
你進入到這些教堂的時候.
中間那一塊是什麼.
是港台.
因為上帝的道要被宣講.
以至上帝的子民能夠回到上帝的話語.
所以你進入到教堂的時候.
你就會看到中間那一塊就是這個.
港台.
是實木.
所以那些港道者.
以前那一派.
就經常拍港台.
因為那個港台不是Akira的港台.
如果是Akira的港台.
一拍就爛了.
是實木來的.
是很漂亮的木.
以前我的教會.
在溫哥華.
加州也是.
那個港台真的很實淨.
這麼大.
裡面其實放了很多東西.
有很多弟兄姐們不知道.
有些水在那裡.
有些廁紙在那裡.
有些垃圾.
等等.
或者有時候那個港台因為這麼大.
他就放一張小椅子在那裡.
因為那些比較矮的木屍.
他就可以站上去.
這樣那個港台就不會阻礙他.
那.
I don't know.
時代又不同了.
今天你去買港台.

$^{841}$或者去找一些新的港台是怎樣的.
透明的.
甚至有很多沒有港台的.
沒有一塊這樣的家具.
但是沒有.
有一樣東西沒有改變的就是上帝的道.
是需要去到被彰顯.
最重要的就是這個東西.
在聖經裡面沒有說.
一定要有一塊這樣的東西.
雖然有一塊這樣的東西.
我覺得有時候也挺好的.
因為安全感是比較大的.
因為你的腳震動沒有人知道.
是真的.
因為我做這行.
所以有時候看見木屍.
他們有時候.
有些腳震動.
因為怕.
但大家看不到.
但你們來到這裡的時候.
有時候有些客席講完就不習慣.
喂.
這裡為什麼沒有東西遮我.
因為我們比較透明.
所以大家是看見的.
但聖經裡面沒有說我們一定要這樣的.
In fact 在聖經的教會.
都是家庭的教會.
怎麼可能放一塊這樣的東西.
是嗎 做了還是.
只不過在廳裡面.
大家小組一個宣講.
但最重要的就是道為中心.
聖經裡面有沒有說.
哪個港道要多長.
有沒有.
你知道有一次保羅去港島.
是說過夜的.

$^{881}$那些人都睡不著.
是不是一定有定義.
或者去決定我們港道要多長.
沒有.
現在新的研究告訴我們.
成人.
他的attention span.
留心聽東西.
最多多少分鐘.
最理想是多少分鐘.
21分鐘.
19 21分鐘.
因為我看了一篇研究是講TikTok.
你知道TikTok.
TikTok.
最理想的一篇港道.
或者那篇的presentation.
19分鐘或者19至21分鐘.
所以我覺得都挺有趣.
新的一代 原來成人都這麼差.
所以我們大概十幾分鐘.
我們就講個笑話.
大家醒一醒.
然後我們又講.
但為什麼很多是以道為中心.
我們走進來的時候.
我們就是receive.
去接收上帝對我的訊息.
我們聽 然後去到信服.
然後就是聖禮.
在基督的教會裡面.
兩大個的sacraments.
兩個聖經裡面有提過的sacraments是什麼.
就是聖禮 水禮和聖餐.
水禮和聖餐.
初期的教會甚至是有愛言.
叫做Agape.
因為Agape就是希臘文的愛的意思.
所以譬如猶大書第十二節說到.
他們聚會的時候有love feast.

$^{921}$有愛言.
所以他們一聚會的時候就吃東西.
一邊吃東西就一邊聚會.
等於今天那些potluck.
那些人帶東西來吃.
或者主人教會預備一些東西.
那些叫做love feast.
整個的worship gathering就是以吃為主幹.
好不好.
看看那些東西好不好吃.
在那love feast裡面其中一部分就是聖餐.
因為聖餐是吃東西的.
所以他們就會拿一些餅.
然後拿一杯.
然後就紀念基督在他們身上所做的事.
當然這個是聖經裡面教會的描述.
不等於他的配方就是這樣.
歷史歷代就是照辦煮碗這樣做.
對不對.
只不過告訴我們.
其中一樣東西我們需要做的就是這些.
聖禮.
水禮.
還有聖餐.
有沒有說過聖餐要多久守一次.
有沒有.
以前我去過一間教會.
每個禮拜都有聖餐.
每個禮拜.
因為他說聖經是這樣說的.
聖經是怎麼說的.
好像他那本聖經和我的聖經有些分別.
其實沒有說的.
最主要說的是regular.
習慣.
所以有很多教會.
譬如我們教會.
一個月一次.
你知道有些教會一年一次的.
好不好.

$^{961}$省水省力.
那不是重點.
他們要做得大.
每一次聖餐的主日的時候.
一年一次.
是一個很大的celebration.
吃的就多一點.
我們那些兩滴就卡住了.
好.
各有不同的做法.
但聖經裡面是給我們彈性的.
好.
浸禮也是一樣.
浸禮很多時候在敬拜裡面是會發生的.
好.
但是有幾多.
你記得.
在少林寺裡面的浸禮.
大部分都不是在敬拜聚會裡面發生的.
是吧.
在河邊.
看到應該要浸禮就去.
是吧.
所以在教會的傳統裡面.
發展裡面.
就會有浸禮.
是包括在敬拜裡面.
然後.
最後的元素就是奉獻.
金錢上的奉獻.
或者資源上的奉獻.
《金銀幕前書》第十六章第一節.
他說論到聖徒的捐款.
捐錢.
我從前怎樣吩咐加拉太宗教會.
你們也當怎樣行.
什麼意思.
吩咐了什麼.
每逢第七天的第一天.
主日.

$^{1001}$各人要照著自己的進貢.
遵項.
抽出來留起.
免得我來的時候才籌錢.
本來其實就是在外邦的教會裡面籌錢.
幫助耶路撒冷教會.
應該是有寄放的.
所以他就說.
照例來說.
我們七天的第一天.
我們各自會有一些進貢的奉獻.
然後奉獻.
他就希望留一部分給我.
以致我們能夠幫助猶太人的教會.
因為現在這個時候.
就是幫助猶太人的教會.
所以有一個奉獻的實踐.
然後最後就會有團契.
在敬拜裡面有Fellowship.
譬如在書信裡面經常提到.
有親咀這件事.
你們要親咀問安.
務要聖潔.
宗教會都問你們安.
親咀.
我們就不親咀了.
今天我們握手.
以前的文化就是親咀.
就是那些面貼面的.
不是咀對咀.
面貼面.
因為中東那些是這樣的.
我們都會有一種這樣的習慣.
在敬拜裡面.
不是只有我一個.
而是去問安.
在基督裡面向大家問安.
本來在書信裡面.
最後的時候就會問安.
問到教會裡面的主管.

$^{1041}$或者那班弟兄姊妹.
向他們問安.
甚至會尊崇他們.
尊貴他們.
所以有人的這一方面.
正如我們所說.
敬拜不是只是宗免.
是有橫面的.
我們留一點時間去問安.
一點點的問安.
在基督裡面問安.
願意我們敬拜.
是真的能夠建立基督的身體.
我們下一個禮拜繼續.
我們同心得留基督.
上主我多謝你.
讓我們能夠有機會.
有福分來敬拜你.
主讓我們成為一個.
合理心意的敬拜群體.
主讓我們生活裡面.
我們所有的難處.
我們今天來到你面前.
可以親民.
主你的恩典是夠我們使用.
也讓我們敬拜裡面.
我們真的向著你.
去聽你的說話.
讓你今天給我們的提醒.
使我們敬拜的生活更豐富.
使我們敬拜的生命.
更加土里奇緣.
祈禱逢轉命救.
好 願上帝賜福給大家.
留一點時間.
大家給錢問安.
OK.
\newpage



\section{}
\label{sec:xTyYHQxRMCA}
\textbf{Everlasting Praise 7 (Rev 4) - Rev. Mark Chan}
\newline
\newline
連結: \href{https://youtube.com/watch?v=xTyYHQxRMCA}{\texttt{https://youtube.com/watch?v=xTyYHQxRMCA}} ~~~~ 語音日期: 2024-06-19
\newline
\newline
\hyperref[sec:3ck_KAakvyM]{\small{< < < PREV SERMON < < <}}
~
\hyperref[sec:index_chronic]{\small{[返順時目]}}
~
\hyperref[sec:index_scriptual]{\small{[返順卷目]}}
~
\hyperref[sec:VHqi_cqG4KM]{\small{> > > NEXT SERMON > > >}}
\newline
\newline
$^{1}$Yeah, do let me know if you're interested,.
join the retreat and it's gonna be a fun time.
and if there are enough people going,.
we're gonna make this happen, okay?.
Why don't we say hi to your neighbors, okay?.
Just say hi, you know, yeah..
Say good morning and welcome,.
some of those friends in our midst..
We continue with our series on worship.
and we learn how to be better worshiper..
And this pretty long series, this is part seven already..
And remember the list of things that I present to you.
from the Bible, if we try to summarize.
what worship is all about and all the elements.
or aspects of worship..
This list, by the way, is not an exhaustive list.
of things we do in worship.
and should not be used mechanically.
because the Bible lays down general principle,.
basic principles for us to practice our worship together,.
but not the details..
And so often the details are dictated.
by the local community or local faith, family,.
churches or even generation..
But we need to learn how to become a better worshiper.
and that is a central question for the entire series..
And let me ask you this question,.
are you becoming a better worshiper?.
Are you becoming a better worshiper?.
Week after week, you know, Sunday is our Shabbat, right?.
Sunday is kind of our big day and the first day of the week..
We come to give God our worship and our offering..
And are we becoming better worshiper?.
And that's the central question we need to ask ourselves..
And today we need to learn how to sing songs, okay?.
How to understand songs..
There are so many elements we can add.
to the list of worship here..
However, I think the main thing,.
there's a main element that we cannot miss.

$^{41}$in our weekly worship, in our Sunday worship.
is the songs, the poems, the lyrics,.
the hymns that we use to sing to God.
and to offer our worship together..
And let me tell you my personal story..
I'm quite surprised that many churches today,.
or at least some of the churches,.
you know, a good number of churches.
are still struggling with using old hymns.
and new praise songs, right?.
And many churches are still debating this issue..
And I'm pretty surprised because, you know,.
I'm quite old, okay?.
I'm quite old compared to many of you guys..
I remember back in the '80s, 1980s, can you believe that?.
1980s, where are you?.
Probably not born, right?.
Right, but you know, I've been going to churches.
for many decades..
1980s is 40 years ago, right?.
The 1980s when I was so much younger, so much younger,.
churches had huge debates on the use of songs, you know,.
especially use of new songs,.
contemporary praise songs, so-called..
When the praise songs were first introduced in 1980s,.
early '80s or even late '80s,.
and those of you who are old enough.
or comparatively older,.
you still remember a music ministry called Marilyn Arthur..
Remember that?.
Yeah?.
No, right?.
I'm one of the older guy here..
You know, some of the songs are still around..
Some of those songs written and composed.
by Marilyn Arthur Music Ministry are still around..
One of those songs we have used before,.
or we have used a lot of time,.
is called "As the Deer.".
As the Deer..

$^{81}$And you know the lyrics, right?.
"As the deer panted for water,.
"so my soul longed after thee.".
Right, "You alone are my heart's desire,.
"I long to worship thee.".
And "You alone my strength, my shield, and so on,.
"I long to worship thee.".
You see, this is the first,.
so-called first generation praise songs.
that have been composed and used by churches.
since the 1980s and late '80s, early '90s..
And it became very, very popular..
You know where this song comes from?.
The lyrics, the content?.
Which song from the Bible?.
No?.
Psalm 42, right?.
42, exactly like that, okay?.
Psalm 42..
And even the lyrics is using some old King James English,.
right, "thee.".
Some of those words are old English words,.
older English words..
It became very, very popular,.
and there were people who don't like this kind of songs..
You know, and I'll tell you why..
And because these kind of contemporary words,.
contemporary worship songs are simple, repetitive,.
and sometimes too short, right?.
And look at this song, you know, you have one verse..
Of course you have verse two,.
somebody add verse two later,.
and then a chorus, right?.
"You alone my strength, my shield," and all that..
But there's not much you can tell from the song itself..
Therefore, there are people who think.
that contemporary worship songs dumbed down worship..
Overly simplified worship, right?.
Because the songs that we use,.
or some of the songs that many modern churches are using,.

$^{121}$like Melanother, first generation,.
are very short and repetitive,.
and they don't cover a lot of content..
They don't cover, definitely they don't cover.
a lot of theological or biblical content, right?.
So there's still a debate in the church,.
and some people who are like me, or even older than me,.
they would prefer the older hymns, right?.
The hymn books, traditional hymns..
And if you check out some of the traditional hymns,.
they do present a lot of content,.
or a lot more content than Melanother,.
or contemporary praise songs..
And check this out..
Oh, this is one of the, if you know Mandarin,.
this is also the first generation,.
so-called first early generation of praise songs, right?.
(speaking in foreign language).
Okay?.
This is actually the first one..
The first one they wrote, and became so popular,.
and this ministry, Stream of Praise,.
has been around for, boy, for decades..
Yeah, for decades..
I remember my wife used to buy the CD..
Each year they have a new album, right?.
And she would go to a Christian bookstore.
and pick up a new CD, all right?.
Every single year, and no matter if she likes the songs.
or not, she will buy a new CD, right?.
But you know what CDs are?.
You don't even know, right?.
I give you a CD, you can't even play it today, right?.
So we use Spotify, and Apple Music, and all that..
iTunes or something..
But this is a first generation praise song,.
and some people don't like it, because it's too simple..
Look at the lyrics..
There's not much there..
It's very focused, though, right?.

$^{161}$Now if you check out some of the traditional hymn,.
like this one, this is a very, very popular classic.
Christian hymn, Amazing Grace..
There are four stanzas, right?.
Four verses..
Actually, there are six verses..
Some versions have six verses,.
and here I present to you four verses..
And you see how rich, in terms of content..
These kind of traditional hymns are very, very rich,.
and teach you a lot of things..
Verse one, Amazing Grace..
How sweet the sound, save a wretch like me,.
was once lost, but now I'm found,.
was blind, and now I see..
What kind of doctrine is verse one presenting?.
Forgiveness, right?.
Forgiveness in grace, right?.
My sin is forgiven..
And then second verse,.
grace taught my heart to fear,.
my fears relieved in grace,.
and how precious did that grace appear,.
the hour, the moment I first believed in Christ..
So stanza two present to you the doctrine of salvation,.
right, and that is salvation in grace..
And verse three, it talks about,.
this is one of my favorite verse in the entire hymn,.
is through many dangerous toys and snares,.
I have already come, and God is leading me,.
safely in grace, and leading me home..
And that's my life journey,.
because full of the grace of God..
Ups and downs in my life journey, in your life journey,.
and when you sing it, when you sing stanza three,.
you realize that the song is trying to tell you.
the providence of God, right?.
The guidance of God in your life journey..
And you are kept safe because of the grace of God..
And then the last verse,.

$^{201}$talks about the future..
The future, the eternity that we share in Christ Jesus,.
and the glorious future, right?.
Someday God will appear and his kingdom will come on earth,.
and I will be God's..
And I say to God, forever,.
I am yours and forever you are mine..
So, it's eschatology, it's the doctrine of last things,.
right, it's something that will happen in the future,.
according to the promises of God,.
the promises of the Bible, the new creation,.
the heaven, right?.
Our true home, and we are heading there,.
and there will be forever praise..
And as I talk through this song, this hymn,.
you see, each verse, each stanza,.
has some sort of profound theology to tell us..
So that's why there are some people who would prefer.
using theologically rich songs,.
and the contemporary worship songs,.
they are too simple to be used in worship..
But we have to ask the question,.
is that a biblical argument?.
Is there really no place for simple songs.
in worship services?.
Is there a place for shorter, simple songs,.
simple praise songs, like the contemporary praise songs,.
that kind of song?.
Do they have a legitimate place,.
legitimate role in our liturgy, in our worship?.
Let's turn to some of the Bible verses,.
and one of those is found in the last book of the Bible,.
the book of Revelation, chapter four..
Chapter four, verse eight, the angels are praising God..
You know, the apostle John who wrote the book of Revelation,.
he is being lifted up to another realm,.
you know, the heavenly realm..
And what he sees is the heavenly worship,.
what worship truly is in eternity..
And he found the angels are praising God,.

$^{241}$you know, day and night, day after day, night after night..
They keep on saying, they keep on praising God.
with one song, holy, holy, holy is the Lord God almighty,.
the God who always was, who is, and who is still to come..
How long is that song?.
Only two lines, right?.
And how many times are they singing it?.
Many times, day and night, right?.
They sing it 24/7, right?.
And it's very simple, very focused,.
praising God for his holiness..
And by the way, what is holiness?.
Holiness is separate, you know,.
it's totally different from us..
Separate from us, set apart from us,.
that is the basic meaning of holy, or holiness..
So they're praising God being holy God,.
because he is sinless, he is separate from us..
However, he has come to redeem us so that we can praise him..
That's what the angels are talking about..
The content is clear and is glorifying God..
And then the next chapter, chapter five,.
again in the heavenly realm, you know,.
in God's realm, in God's domain,.
and the people and the angels are singing praises.
to the lamb that was slain..
If you remember our set of worship in the beginning,.
and Josh led us to read this part of scripture,.
the angels are praising God in front of the throne.
of the lamb of God who was slain..
And they said, they're saying,.
"You are worthy to take the scroll.
"and break a seal and open it.".
Now what does it mean breaking the seal of the scroll?.
You know, in the ancient time,.
when they write important documents,.
they would roll it up, you know, roll it up..
In animal skin or in papyrus,.
they roll it up and then they tie it up.
and then they will have a lump of clay.

$^{281}$to stamp on it to seal the document..
So whoever has that stamp, whoever has the signature,.
has the authority to open up the scroll..
And the scroll is all about the future,.
about the future of creation..
And who is ruling the entire creation,.
the entire timeline of history,.
from past, present, and future?.
Who is worthy to open the scroll?.
Who is it?.
For you were slaughtered, your blood has ransomed people.
for God from every tribe, language, people, and nation..
And you have caused them to become a kingdom of priests.
for our God and they will reign on the earth..
Who is worthy to open the scroll?.
Who has the authority, the sovereignty.
to control your future, to control the future.
of the entire creation?.
It's the lamb who was laying on the cross..
The lamb who suffered and died but rose again..
He is the living God in Jesus Christ our Lord..
He is worthy to open that scroll..
And this is actually the song the angels were using.
to worship the lamb in front of his throne..
How long is this song?.
Not too long, right?.
Short enough I can put it in one slide, right?.
And that's what they were singing.
and they were using to worship God..
So these are very, very short,.
kind of like contemporary worship songs, right?.
So they're legit..
We can use those because they're focused..
They don't tell a lot of content.
but the focus is very, very sharp..
Now are there long praise psalms in the Bible?.
Psalm 119, right?.
It's a Torah psalm..
Appraising the word of God, the value of the word of God,.
the honor of the word of God..

$^{321}$And it's very, very long..
But again, there are many psalms that are shorter.
and much shorter and sometimes even shorter.
than the contemporary praise songs..
Like Psalm 1, right?.
You remember Psalm 1?.
Blessed is the man who does not walk.
in the pathway of the sinners and all that..
And Psalm 1, how many verses are there?.
Six, six verses..
Psalm 23, and many of us know how to recite Psalm 23, right?.
The Lord is my shepherd, I shall not be want..
The Lord is my shepherd..
How many verses are there?.
Six, six verses only..
121, I lift up my eyes and where my help come from..
121, eight verses..
133, behold how good the brothers,.
how beautiful the brothers are living together,.
staying together..
133, three verses only..
So there are shorter songs, many shorter songs.
in the Bible for worship..
And also there were filled longer,.
theologically rich songs that can be used.
in worship services..
You know, if you check out some of the traditional hymns.
that one of the examples I showed you, Amazing Grace,.
they try to squeeze major doctrines into one song, right?.
I went through with you, Amazing Grace,.
there are at least four kinds of,.
or four aspects of theology being squeezed into one song..
Now I don't think, I think there's a legitimate way.
to do it and it has been done so well in Amazing Grace..
But I don't think we should expect a song to be like.
a systematic theology textbook, right?.
I know a lot of you don't read, right?.
You don't read books anymore, right?.
We all go TikTok, right?.
Short videos, right?.

$^{361}$If you pick up some of the books from our library, right?.
From our library corner, you pick up one of those.
systematic theology books, how many pages are there?.
500 pages..
That's crazy, right?.
If you take one of the courses in seminary,.
systematic theology, that's one of the textbooks.
that you're gonna read, right?.
There's only one..
There are probably 10 books you have to read..
And you just, you pick up one of those books,.
you know, it's like weapon, you know?.
You can bang people on the head and it's gonna hurt, right?.
You know, if you squeeze the entire book,.
systematic theology book, into one song,.
which is about five minutes long, three minutes long,.
it may be counterproductive at times, right?.
Right?.
Because you go through so many, so many different types.
of teachings, so many different types of content.
in one medium, one short medium,.
that may not be the best way to do it, right?.
And in worship service, as you can tell.
from the list that I gave you,.
all the elements of worship service,.
I think in terms of conveying content,.
you know, offer you materials, information,.
I don't think songs are the best way to do it, right?.
Do you agree?.
Probably preaching, 25 minutes, 30 minutes,.
is the better way to convey content..
So again, I can understand why some guys like me,.
who are older, would prefer traditional hymns, right?.
But, you know, there was legitimate role.
for shorter, you know, simple songs,.
so that our praise and our worship would be focused..
And some people in the discussion or debate.
would prefer traditional hymns over contemporary music.
because of quality, because of quality..
And contemporary pop music sometimes is rejected.

$^{401}$because of lower music quality, right?.
Do you agree?.
Those are the songs that we use today, week after week,.
they are lower music quality, right?.
And of course, right, musically speaking,.
you know, classical music is of better quality, right?.
In terms of harmonization, in terms of sophistication,.
right, you know, Handel's Messiah,.
I mean, it's so complicated, it's so rich, right?.
And the harmonization and the arrangement is so rich,.
you know, Mozart, Beethoven, you know,.
those kind of music, traditional classical music.
is of better quality in terms of musicality..
So arguably, they are better music quality..
And some people would prefer the traditional hymns.
because of that, because of the better quality..
But as I explained to you in the entire series,.
I asked the question, is it the church mission.
to preserve traditions?.
Is it or not?.
Is it the mission of a church to preserve music or art.
of highest quality?.
Is it our mandate?.
Partially, I think so, right?.
Cultural mandate, Genesis one..
However, one can arguably say.
that this is not the church mission, right?.
Some of the music that we use in the entire long history.
of Christian history only serve for one generation.
because that generation understand.
that kind of language of worship..
And I told you, I learned some Hebrew songs, right?.
And those songs are particularly understandable,.
intelligent to that particular culture..
And it would not be, it would not have the same effect.
for us, right?.
And so, there are people who say.
that music, church music, has to be good enough.
to be preserved beyond its own time..
Otherwise, we should not use it in worship service..

$^{441}$Now, I appreciate the basic intent,.
but I don't think this is totally biblical.
because remember the principle of Paul.
in 1 Corinthians 14..
He said, "No matter what we do, we need to do it.
"in such a way that we can edify the entire family,.
"the church, the local church..
"The church can be strengthened.".
Even new visitors come in and they can feel God's presence..
They can understand the lyrics without much difficulty..
And that is what Paul is going after.
in terms of quality, right?.
Because there are many different kinds of quality.
in the world..
And Paul is saying to the Corinthian church,.
"And please, do your worship and engage in worship.
"in such a way that is good for everybody.".
Even it is effective for new believers or non-believers..
And they come in, what is best in music.
in terms of what Paul is saying.
is that the music that will communicate.
with people today, right?.
And certainly there are musics that I've used.
or my worship life has used in the past..
They are bygone..
You know, we don't use it anymore.
because it spoke to the previous generation.
but no longer speaking to this generation..
So I think there are place for all kinds of music..
There are place for all kinds of worship songs.
including contemporary modern praise songs..
And again, we need to understand what song is all about..
Song or poem is not the ideal medium for content, right?.
As I told you, it's only five minutes long, right?.
You can't say much..
You can only focus on one thing at a time..
And we have to understand and ask the question.
how we can become a better worshiper.
by understanding what worship songs is all about..
You know, if you have come to church for quite some time,.

$^{481}$you understand worship songs are like prayers..
They are directed toward God, right?.
They are prayer-like songs and directed to God, right?.
Most of the songs that we use..
And as such, worship songs or worship poems,.
the songs that we sing, are written for expression,.
emotional expression, emotive impact..
And in order to create that kind of emotive impact,.
it has to be somewhat reparative, right?.
The contemporary worship songs,.
even the traditional hymns, are reparative in some way.
because that is the nature of worship songs..
I would like to use the analogy.
that worship songs are like love songs, right?.
You know, you sing some of the love songs,.
even secular love songs, right?.
Those are very good songs, right?.
They are reparative in some way..
Or worship songs are like love songs and national anthems..
National anthems..
Has to be repeated, right?.
It is kind of reparative..
You do it one time, two time,.
repeatedly, regularly, right?.
Every NFL game start with national anthem, right?.
And we're gonna sing together.
because national anthem are pretty much like love songs..
They are emotional expression.
to declare allegiance to the country, right?.
Are there any reparative songs in the Bible?.
There are..
One of the example is Psalm 136..
Here I give you the first five verses..
Give thanks to the Lord for he is good,.
for his standfast love endures forever..
Give thanks to the Lord, give thanks to the God of gods.
for his standfast love endures forever..
Give thanks to the Lord for his standfast love.
endures forever..
To him who alone does great wonders.

$^{521}$for his standfast love endures forever..
To him who by understanding made the heavens.
for his standfast love endures forever..
And it keeps on repeating the same chorus again and again.
and tells you why do we praise God,.
why do we give thanks to him..
And this Psalm 136 is very, very long, right?.
It's very, very long..
And it actually recites the entire history of Israel..
How God told his people, talk about Exodus,.
talks about signs and wonder God has done in the midst,.
talk about wilderness, talk about how God provides for them..
And that's why God is praiseworthy..
Now, this is a very emotive expression..
Repetitive, repeat again and again,.
for his standfast love endures forever..
Yes, amen, it endures forever..
That's why we praise him..
And every time the chorus repeats,.
it talks about different aspects of God.
and what God is doing in your life, right?.
God has certainly provided for you,.
God has certainly sustained your life,.
God has certainly guided through your life journey,.
and that's why he is praiseworthy.
for his standfast love endures forever, right?.
These are emotions, right?.
Just like my wife would prefer me to say I love you..
For how many times?.
Many times, right?.
I can't tell my wife, you know,.
I told you five years ago, right?.
It's done deal, right?.
There's no more content I wanna express to you, right?.
But you understand this emotive expression, right?.
It has to be spoken or verbalized from time to time, right?.
You know, when my kids were younger,.
we try to intentionally do that, you know,.
we try to teach our kids, you know,.
as you know, kids don't talk these days, right?.

$^{561}$They're often with their phone, right?.
So when they were so much younger,.
we tell them, you know, you have to express your feeling..
You know, you have to verbalize it..
When you see a sunset is beautiful,.
you should say it is beautiful, right?.
It is awesome, right?.
They need to learn how to find the right word.
to describe their feelings..
And we intentionally do that at home..
I will tell them I'm happy because what happened, you know?.
I'm happy, I'm excited, I'm sad, I am kind of depressed,.
I'm exhausted..
You know, these kind of vocabularies.
are emotional vocabularies, right?.
I tell my kids, "Love you," right?.
And they tell me, "I love you.".
And this sort of expression has to happen,.
you know, linguistically, I mean, verbally,.
it has to be described, right?.
And I want to bring this analogy to worship as well,.
because worship is an emotional language..
And I want to tell you from my own life journey,.
my own cultural heritage, right, Asian Chinese, right?.
We have lack of emotional vocabularies..
We have none, okay?.
We have no emotional vocabularies in our dictionary, right?.
Don't you agree?.
How do we express love to our children?.
Chinese parents, we do, "Eat some more.".
(audience laughing).
Have you eaten?.
Go eat, right?.
Oh, eat some more, right?.
This is how we do it, right?.
This is pretty sad..
One time I was visiting a family,.
and they have one of their loved one passed away..
It's very tragic..
So I went to visit the family,.

$^{601}$and they just have no vocabularies to describe the feelings..
In the entire visitation, almost an hour,.
I stay there an hour, and I pray with them,.
you know, I care for them..
All they say is, "Pastor, eat.".
(audience laughing).
Or, "You have to eat," okay?.
I said, "I had dinner.".
"Oh, no, you have to eat.".
"Oh, you have to drink some more," okay?.
And then we talk for about 10 minutes,.
and then later on, you know, the family members come,.
and we found the dish, and come, "Hey, eat," okay?.
I mean, can we say I'm sad?.
Can we say I'm depressed, and I wanna cry,.
and I can lament, right?.
This is a legitimate emotional expression.
that we need to use..
This is a love language, right?.
And we need to learn that..
And in order to be better worshiper,.
we need to learn that kind of language as well..
So I wanna suggest to you, we need to do that..
Not just for parents, okay?.
Not just for family relationship..
You know, all kinds of relationship.
have to be emotionally related, emotionally expressive, right?.
I, as a father, I need to learn how to use those words, right?.
And there are difficulties I need to overcome, right?.
At first, I was quite uncomfortable to say that, right?.
To say to my wife, "I love you.".
Say to my kids, "Love you," right?.
And I kind of, along the way,.
I find some cool way to say it, right?.
"Love you.".
(audience laughs).
And then I get used to it..
I get more comfortable to express it..
But my main point is we need to learn that, right?.
Even between friends..

$^{641}$You know, how do we get closer and closer, right?.
It has to be emotionally expressive, right?.
You gotta have some emotional connection..
And worship is an emotional response to God, right?.
Like this psalm..
"God has saved us and redeemed us,.
"created us and sustained us..
"And we feel good," right?.
We wanna give thanks..
Even giving thanks itself is an emotional response, right?.
And we respond to the goodness of God.
and ultimately the salvation that we have in Jesus Christ.
because we sinners, we're separated from the holy God..
Only Jesus came to bridge the chasm.
between the holy God's totally separate from us,.
alienated from us, and us, broken sinners..
And Jesus died on the cross and rose from the dead.
so that we can be bridged back to God, related to God..
So how can we be a better worshiper?.
And may I suggest to you three more things.
as your take-homes..
First, be intentional..
We have to be intentional..
Our weekly Sunday worship service,.
we have to be intentional..
And you realize how much time we devoted.
to prepare for worship service,.
no matter how many people show up?.
I think we practice, this practice, today's rehearsal,.
I think it took us two plus hours, right?.
And with the preparation and all that,.
number of hours, we do, we rehearse.
in order to craft our worship, right?.
Of course, it's not the best of the world, right?.
We're not professionals..
But we try to do our best, right?.
So what I'm saying is, first point is be intentional.
if we want to become better worshipers..
Arrive a little bit earlier..
I told you many times, my best practice is come in.

$^{681}$and find a seat, find some space, and I can pray to God..
I can meditate on God..
Remember, worship is a relational encounter, right?.
And today is Father's Day..
God is our true Father, right?.
We early fathers are broken..
We are sinners, and we need redemption of Christ,.
just like our children do..
But God is our true and perfect Father..
And Jesus died on the cross to redeem us.
from sin and from death, so that our relationship.
with God is a reality now in Christ..
So be intentional, because God is relating to us, right?.
God is not an ideal or a concept, or even not a philosophy..
God is a living God..
He's a person, and we have a personal relationship.
with our Creator, with our Redeemer..
Secondly, I would suggest to pray..
Learn how to pray..
What do I mean by that?.
You know, in prayers, we don't just make a list.
of needs and wants in front of God..
That is not the essence of prayer..
Prayer is to come before God with an open heart,.
be honest, and be open, and tell God how you feel today,.
how you feel this week..
Maybe you're so tired this week..
Yes, open to God, right?.
Tell God that you are tired..
Tell God that you are exhausted..
I need your help..
Tell God you feel frustrated..
You feel anxious..
You worry about things..
You present those emotional feelings to God,.
and even tell God you are joyful today..
Tell God you're excited, right?.
Did you listen to the words that we use?.
Emotional, expressive terms, right?.
God, I'm exhausted, I'm frustrated..

$^{721}$God, I'm happy, right?.
I'm joyful..
All kinds of emotions are found in the book of Psalms..
All kinds of emotions are found.
in the worship elements in the Bible..
I remember September 11th..
September 11th, I think it happened on,.
of course, it happened on September 11th, 2001, right?.
And I think it happened on Monday or Tuesday,.
yeah, if I'm not mistaken..
And I remember after September 11th happened,.
I was serving in the church,.
and the worship leaders called me up and say,.
how should we respond, right?.
You know, the original worship set,.
we have chosen some songs, right?.
And we have to change every one of them, right?.
And I suggested to them, you know,.
we don't have to come to worship.
always feeling joyful and excited..
We can come before God with our laments, right?.
Lamentations..
Those are legitimate worship response.
and emotional response..
We can lament before God..
And the book of Psalms contains many, many psalms.
that are laments, community laments,.
or individual laments..
Those are sad songs, but we can present those to God..
So I suggest to you to pray, to learn how to pray..
Think for a minute how you would respond.
if God actually appeared to you right now..
If God all of a sudden, (mimics God's voice).
appeared before you, how would you respond, right?.
You know, in the Bible, there were many encounters like this.
in the book of Job..
Job was suffering, right?.
He's a righteous guy and suffered a lot..
And he demanded God to answer his questions..
He has so many questions about suffering, right?.

$^{761}$Why am I suffering, right?.
An innocent man suffering a lot..
Job demanded God to answer his questions..
However, at the end of the book, God did show up..
God said, "Yeah, you wanna interview me?.
"Yes, here I am.".
But when God met Job,.
God in turn asked him a lot of questions, right?.
Do you have the wisdom?.
Do you have my transcendent wisdom.
to govern the entire cosmos?.
Do you have that kind of wisdom?.
Can you understand that kind of wisdom?.
And Job was silent before the Creator..
And he encountered God in such an awesome way, right?.
Moses and Isaiah, they encountered God in worship..
And they realized God is holy and sinless..
And they were unworthy..
And they begged for mercy,.
that we need your grace, God, to come before you..
Otherwise, we are undone, right?.
We are over because we are sinful..
And God is so bright and so glorious, so pure..
How can we approach God?.
There's no way we can approach God,.
except on the ground of grace..
And maybe we can come before God.
and prepare ourselves in prayer,.
knowing that only Jesus made forgiveness possible for us,.
so that we can worship God, we can come before God,.
because Jesus first loved us..
And number three I wanna suggest to you is to learn, right?.
We have to learn..
No matter how long you have been to church,.
no matter how long you've been a Christian,.
no matter how long you have followed Christ,.
we need to learn how to worship,.
how to worship better..
Because worship is not about passively watching a show..
It's not Saturday night life, right?.

$^{801}$It's not just come in and you just absorb something..
It's not passively engaging..
It's about actively, proactively, we participate..
Because we believe in Jesus and we are God's people..
And God is calling his people to worship him..
And worship flows out of our living relationship.
with God himself..
Today, even I've been worshiping for many decades,.
I'm still learning new songs, right?.
Because I understand when God does new things in history,.
there are new songs of worship..
And I'm trying to learn those new songs, right?.
Try to learn and try to concentrate..
Try to direct our attention to God.
so that God is totally pleased in your worship,.
in my worship, in our worship..
Let us all pray and pray that God will help us.
to be better worshipers, to be a better church..
Heavenly Father, we call you Father because you are perfect.
and you are holy and sinless and pure..
Yes, we can come before you today.
only on the ground of grace..
Because of the grace of God expressed to us,.
demonstrate to us on the cross of Jesus..
Thank you for first loving us so that we can love you back..
Thank you for this living relationship.
that we have with you..
Our prayers are not talking to empty air..
Our prayers are directed to you..
And you are our loving Father..
Thank you for sustaining us, providing for us..
And God, we come before you, we are sorry..
Forgive us, forgive our shortcomings..
Sometimes we come before you to worship.
as if this is a job, as if this is one of the things.
that in our checklist, we wanna get it done.
and get it over with..
Father, forgive us..
Bring us into true worship..
Help us to be better worshipers,.

$^{841}$to learn languages that will express our thanks to you,.
to learn songs and poems that we can say to you,.
you are our awesome God..
Thank you for being our loving God..
Thank you for saving us and redeeming us.
from sin and from death and darkness..
Bring us into your light..
Help us to be renewed because we understand your grace..
So transform our lives this week.
so that our lifestyle will honor to you..
And as we come week after week to worship you,.
may our worship be pleasing to you.
for you are worthy of our praise..
In Christ name we pray..
Well, God bless you all and please stay behind.
for a little time of fellowship together.
and there was some refreshment out there.
and say hi to some of our new friends.
and say hi to your neighbors, okay?.
God bless..
\newpage



\section{啟示錄 4:1-11}
\label{sec:VHqi_cqG4KM}
\textbf{永恆的敬拜 6 (啟示錄 4\_1-11) - 陳崇基牧師}
\newline
\newline
連結: \href{https://youtube.com/watch?v=VHqi_cqG4KM}{\texttt{https://youtube.com/watch?v=VHqi\_cqG4KM}} ~~~~ 語音日期: 2024-06-19
\newline
\newline
\hyperref[sec:xTyYHQxRMCA]{\small{< < < PREV SERMON < < <}}
~
\hyperref[sec:index_chronic]{\small{[返順時目]}}
~
\hyperref[sec:index_scriptual]{\small{[返順卷目]}}
~
\hyperref[sec:cMH1ZiMWOm8]{\small{> > > NEXT SERMON > > >}}
\newline
\newline
啟示錄 4:1-11
\newline
\begin{longtable}{cl}
\hline
\hline
章節 & 經文 (和合本修訂版)\\
\hline
4:1 & \begin{tabularx}{0.7\textwidth}{X} 這些事以後，我觀看，看見天上有一道門開著。我頭一次聽見的那好像吹號的聲音對我說：「你上這裡來，我要把此後必須發生的事指示你。」 \end{tabularx} \\ \\ \relax
4:2 & \begin{tabularx}{0.7\textwidth}{X} 我立刻被聖靈感動，見有一個寶座安置在天上，有一位坐在寶座上。 \end{tabularx} \\ \\ \relax
4:3 & \begin{tabularx}{0.7\textwidth}{X} 那坐著的，看來好像碧玉和紅寶石；又有彩虹圍著寶座，光彩好像綠寶石。 \end{tabularx} \\ \\ \relax
4:4 & \begin{tabularx}{0.7\textwidth}{X} 寶座的周圍又有二十四個座位，上面坐著二十四位長老，身穿白衣，頭上戴著金冠冕。 \end{tabularx} \\ \\ \relax
4:5 & \begin{tabularx}{0.7\textwidth}{X} 有閃電、聲音、雷轟從寶座中發出。在寶座前點著七支火炬，就是神的七靈。 \end{tabularx} \\ \\ \relax
4:6 & \begin{tabularx}{0.7\textwidth}{X} 寶座前有一個如同水晶的玻璃海。寶座的周圍，四邊有四個活物，遍體前後都長滿了眼睛。 \end{tabularx} \\ \\ \relax
4:7 & \begin{tabularx}{0.7\textwidth}{X} 第一個活物像獅子，第二個像牛犢，第三個的臉像人臉，第四個像飛鷹。 \end{tabularx} \\ \\ \relax
4:8 & \begin{tabularx}{0.7\textwidth}{X} 四個活物各有六個翅膀，遍體內外都長滿了眼睛。他們晝夜不住地說：「聖哉！聖哉！聖哉！主—全能的神；昔在、今在、以後永在！」 \end{tabularx} \\ \\ \relax
4:9 & \begin{tabularx}{0.7\textwidth}{X} 每逢四活物將榮耀、尊貴、感謝歸給那坐在寶座上、活到永永遠遠者的時候， \end{tabularx} \\ \\ \relax
4:10 & \begin{tabularx}{0.7\textwidth}{X} 二十四位長老就俯伏敬拜坐在寶座上活到永永遠遠的那一位，又把他們的冠冕放在寶座前，說： \end{tabularx} \\ \\ \relax
4:11 & \begin{tabularx}{0.7\textwidth}{X} 「我們的主，我們的神，你配得榮耀、尊貴、權柄，因為你創造了萬物，萬物因你的旨意被創造而存在。」 \end{tabularx} \\ \\
[1ex]
\hline
\hline
\end{longtable}
$^{1}$好,今天我們繼續看這個主題.
這次應該是最後一講.
就是講到敬拜.
是什麼意思呢.
很需要我們去想一想.
敬拜這方面.
因為來到教會的時候.
最重要的事情.
就是學會如何去敬拜.
如何成為一個worshipper.
一個敬拜者.
你記得我曾經分享過這個list.
聖經裡面去到混雜所有對敬拜的教導.
或者對敬拜的描述.
有這幾方面.
當然這個列表.
不是一個完全的列表.
不是一個所有的元素都放上去的列表.
也不是讓我們去到機械式的去運用.
照板煮碗.
我們就做回這些東西.
每一個項目我們都做回.
聖經裡面給我們一些很基本的principle.
一些很基本的原則.
講到敬拜.
很多時候不是可以逗到很多的細節.
但是我們需要問.
如何能夠成為一個好的敬拜者.
如何能夠成為一個敬拜者.
是合上帝的喜悅.
想一想這個問題.
如何可以成為一個更加上帝喜悅的敬拜者.
今天最後一堂.
我們去到看敬拜裡面其中一個項目.
就是詩歌的項目.
特別是我們去設計我們的崇拜.
或者我們來到一個崇拜的時候.
我們想起唱這些詩歌的時候.
我們如何能夠使我們成為一個更好的敬拜者.
我們回答一個問題.

$^{41}$就是現代的敬拜詩歌.
是不是一個適合的詩歌.
還是過度簡化.
dumb down.
dumb down我們的敬拜.
剛才我們唱的那些詩歌.
都算是現代的詩歌.
我們比較老一點的信徒.
我經常形容自己.
我就是那些老差骨.
進到教會.
我知道你們那些前後左右.
裡面外面所有的東西.
我知道怎麼運作的.
就是我們這些老差骨.
問這個問題.
到底現代的詩歌是不是會dumb down worship.
我有一件事覺得很奇特的.
就是我今天去到教會的時候.
或者聽到不同教會的一些對話.
都仍然在掙扎著舊的詩歌和新的詩歌的張力.
有些人認為舊的詩歌才是有價值的.
新的詩歌或者現代的詩歌.
可能有些是好一點的.
但是大致上質素上都是比較差一點.
一代不如一代.
新的一代是不是會這樣呢.
如果老差骨的話.
你好像我一樣.
老差骨.
你記得1980年代.
80年代.
老差骨了.
我們中間有很多80年代的.
可能不在這裡.
80年代的時候.
年輕的時候我們去教會.
是不是有很大的debate.
就是講到舊的詩歌和新的詩歌的張力.
有些詩歌可以用.

$^{81}$有些詩歌不可以用.
為什麼呢.
因為80年代的時候.
先講講歷史.
死火了.
講歷史40年前.
是不是.
真是要命.
不過我們需要知道近代的歷史.
知道近代發生了什麼事.
我們才知道來龍去脈.
然後我們懂得怎麼去分析.
然後我們看看聖經裡面怎麼講.
80年代開始有一些敬拜詩歌的士工是湧現的.
大家記得.
有個英文.
就是從英文這裡開始.
從美國開始的.
大致上從美國開始.
有個叫Melanova的音樂的組織.
記得嗎.
Melanova.
你沒有人點頭的.
你不要裝小時候的.
我知道那些應該是知道的.
Melanova就有一些新的詩歌出現了.
那些叫Praise and Worship Songs.
Praise and Worship Songs.
我記得那時候我們團契或者教會的時候.
拿著一本很薄的.
叫做Praise.
就是北美教會.
就打開這樣去唱.
裡面就有一些現代寫的詩歌.
你記得這首歌的詩歌.
叫做As the Deer.
記得嗎.
1984年的時候寫的.
1984年的時候.
Melanova所寫的.

$^{121}$這首詩歌.
As the Deer Pants for Water.
Right.
So my soul longs after thee.
Right.
就是還在用一些.
Old King James的英文.
You alone are my heart's desire.
等等.
I long to worship thee.
還翻譯成中文.
有廣東話.
有國語.
國語我這裡給你.
是嗎.
這首詩歌在哪裡來的.
材料在哪裡來的.
歌詞.
詩哪一邊.
42.
就是如鹿鶴無雞髓.
那裡就是在說詩人去鶴舞上帝.
這首是很簡單的詩歌.
有多少節.
嗯.
算是兩節.
其實也只是一節.
有副歌.
Right.
這些新的詩歌.
當時湧現的時候.
有很多人覺得.
這些好像過分簡化了.
就是敬拜.
因為這首詩歌說了多少東西.
你只看字也不可以說太多.
因為字已經這麼少字.
這麼少節數.
有這些新的詩歌出現的時候.
很多時候有那種張力.

$^{161}$就是新的詩歌到底是不是丟了.
還是我們應該要保存舊的詩歌.
這首詩歌.
大家知道.
李兆王.
大概什麼時候.
90年代2000年.
差不多.
我不太記得.
問問Ivan.
就是那時候出這首詩歌.
或者出這類的詩歌.
是.
也有那種張力.
就是那些老差骨.
就會覺得傳統的詩歌怎麼辦.
傳統的詩歌豐富很多.
李兆王.
這裡有多少字.
很少字.
所以現代詩歌給人的感覺是比較簡單.
短.
還有是重複的.
Right.
懂國語的.
知道這首了.
是吧.
讚美之詮.
現在出到.
多少唱的CD.
那些詩集.
出到多少.
好像十幾了.
十幾.
我記得每年他出的時候.
我太太也買一隻.
也買一隻.
當時要買現在不用了.
現在Spotify.
是吧.

$^{201}$她收藏了買了CD.
其實我想丟掉它.
因為家裡沒有位置.
而且她聽到.
你去哪裡找一台CD機給她聽.
是吧.
那代位已經完了.
我們已經轉身了.
但這是第一首.
讚美之詮的詩歌.
這首詩歌叫什麼名字.
就是叫讚美之詮.
所以那個時宮也是叫讚美之詮.
到現在仍然是一個很好的時宮.
是吧.
一個台灣人為主.
或者國語為主的詩歌.
沒錯.
這些詩歌我們看下去的時候.
是短.
簡單.
還有是重複的.
如果你看它有沒有.
聖經或者神學的內容的話.
是.
比較少的.
是吧.
我給大家一個比較.
這個首本名曲.
奇異恩典.
有多少節.
有四節.
其實有六節.
通常我們有四節.
六節.
每一節.
都是說出一些神學的概念.
或者神學的教義.
以致我們唱詩歌的時候.
我們能夠得到一些內容.

$^{241}$是吧.
神學上是豐富的.
第一節.
奇異恩典何等甘甜.
說的是什麼.
Forgiveness.
說的是罪得赦免.
第二節.
浩大恩典我敬畏.
然後初選之時蒙恩.
為了這個寶貴恩典.
要說救恩.
Salvation.
Forgiveness.
Salvation.
然後第三.
危險往來我飽受人間苦楚.
是吧.
平安渡過.
然後將會回到天家.
說的是上帝的大靈.
Providence.
是吧.
第四節.
說的是盼望終末的事.
將來在天安居萬年.
讚美永不減少.
是吧.
說的是天上的敬拜.
是吧.
所以你唱完一首詩歌.
你會發現.
嘩.
很豐富.
是吧.
四節.
說出了四個神學的概念.
是不是神學上很豐富.
是吧.
不像讚美之傳.

$^{281}$只是我們舉起手唱歌.
讚美上帝.
是這樣的.
但這些是有內涵的.
正因為這個緣故.
所以有很多人認為.
這些舊有的詩歌.
或者傳統的詩歌.
才是適合於敬拜.
有很多現代的詩歌.
不是.
可以用.
但不是很到位.
是吧.
特別對我們這些老差骨來說.
是吧.
因為太重複了.
現代的詩歌.
舊的詩歌不是.
有很多內容出來.
是吧.
是否可以這樣去.
是否可以這樣理解.
和這樣去分析.
你覺得呢.
傳統詩歌.
沒錯是將一首詩歌.
然後把所有的.
或者很多的神學的內容.
是加進去的.
是吧.
好像這首詩歌.
如果你想一下.
我自己生命裡面.
成長的裡面.
在教會裡面學的.
那些傳統的聖詩.
隨便想一下.
Mighty Fortress.
堅固保障.

$^{321}$Church is one foundation.
教會唯一的根基.
有六七節的.
每一節都是說出不同的.
神學和信仰的內容.
是吧.
是豐富很多的.
是吧.
這個是說真的.
是吧.
真是這樣.
但是.
如果你說要一首詩歌.
把所有的東西加進去的話.
那我有時候就會想.
會不會.
你想一下.
六節詩歌.
每一節都說一個.
或者某一方面.
神學的教義.
就是什麼都說.
是吧.
那會不會水過鴨背呢.
是吧.
就好像那些讀神學院的同學.
我們交功課的時候.
你是不是什麼都放進去的.
教授告訴你.
你要有論文.
你要有主題才行.
是吧.
你不可以什麼都放進去.
不然你什麼都說.
就變成什麼.
什麼都不說.
是吧.
你明白我的意思嗎.
就是如果你要預計.
一首詩歌.

$^{361}$要他交代所有的東西.
由你信主到你去天家.
每次唱一首詩歌.
都是由你出生到你去天家的.
是吧.
那是不是其實.
就是你預期太多了.
是吧.
只是短短一首詩歌.
三至五分鐘.
最多七分鐘.
什麼都說.
你有沒有看過那些.
我知道大家已經不看書了.
那些系統神學的書.
多少頁的.
五百多頁.
那些只能看到它.
你還沒打開.
你的臉都變青了.
五百頁的系統神學書.
我家裡有好幾本.
是囊括了整個聖經裡面.
所教導的主要的教義.
但是你可不可以將那五百頁的書.
壓在五分鐘裡面.
這樣的預期是不實際的.
還有詩歌本身的題材.
不是真的要講內容的.
對嗎.
不是講內容的.
你喜歡聽什麼歌.
舊的歌也好.
新的歌也好.
現代歌也好.
大概都是三至五分鐘.
是吧.
是不是你聽歌去學內容的.
不是的.
是吧.

$^{401}$聽歌很多時候.
歌曲是做什麼的.
是表達情緒的.
表達情感的.
這是歌曲應該有的題材.
如果你說要一首詩歌裡面.
加上所有東西.
這樣才是一首好的詩歌.
我覺得好像沒什麼道理.
是吧.
你看看聖經怎麼說.
說到天上的敬拜的時候.
在史書第四章第八節.
說什麼.
他說他們晝夜不住地說.
他們就是四活物或者天使.
這裡說的是約翰被提到天上的時候.
他發現天上有一個敬拜.
敬拜那位在上掌權的主宰.
他說天使晝夜不住地說.
盛哉盛哉盛哉.
主全能的上帝.
昔在今在以後永在.
他們唱什麼詩歌.
就是那兩行嗎.
是吧.
唱多少次.
很多次.
晝夜不住.
是吧.
我們大哥那些又來.
又來.
夠了.
晝夜不住.
來到我們的場位.
很短的首詩歌.
但是內容和focus很清楚.
因為上主仍然掌權.
不是羅馬帝國裡面的該殺掌權.
再看下一章.

$^{441}$天上的敬拜怎麼形容.
第五章.
說到被殺的羔羊.
他看到這位得勝者.
原來是一位曾經被殺的羔羊.
第九至十節.
你是配拿書卷的.
配揭開七人.
因為你曾被殺.
用自己的血.
從各支派.
各語言.
各民族.
各邦國.
買了人來.
救贖了人來.
使他們歸於上帝.
又使他們成為國民祭司.
歸於我們的上帝.
他們將在地上執掌王權.
說到被殺的羔羊.
去到敬拜他.
首詩歌有多長.
很短.
是嗎.
很短.
起碼剛才讚美之前的詩歌比他長.
是嗎.
詩篇裡面.
說到敬拜的時候.
有沒有一些長的詩歌.
有沒有.
詩集,詩篇.
最長的是哪一篇.
一把十九篇.
是嗎.
但是有沒有短的.
有.
大把.
詩篇第一篇有多少節.

$^{481}$六節.
詩篇二十三篇大家會背.
耶和華是我的牧者.
我必不自缺乏.
大部分特別是老差骨會背.
詩篇二十三篇有多少節.
六節.
是嗎.
六節.
詩篇一百二十一篇.
我向山舉目.
八節.
一百三十三篇.
說到弟兄同住的美好.
是嗎.
多少節.
三節.
所以詩歌裡面是有短也有長的.
有些短的詩歌同樣是一些應該用的詩歌.
但是有時候說到這個話題的時候.
有時候有些人仍然覺得.
傳統的詩歌.
傳統的音樂質素上會比較好.
這個我覺得也是對的.
現代的音樂.
單單說現代音樂.
不要只說基督教音樂.
普通的音樂.
Pop music.
在說到音樂質素的時候.
和音或者寫的時候的複雜等等.
現代的詩歌比較簡單.
或者現代的歌曲比較簡單.
或者曲風比較簡單.
如果說一些古典的音樂.
這些是不能比較的.
巴赫.
Mozart.
Beethoven.
Handel的尼采亞神曲等等.

$^{521}$那些都是非常複雜.
非常精細的音樂.
如果說質素的時候.
確實那些傳統的.
或者古典的.
是難很多和深入很多.
所以那些是good quality.
但是我想問一個問題.
什麼叫quality.
什麼叫good.
你記得上一次.
我們說到保羅的敬拜神學.
它的原則是什麼.
The people can understand.
People can relate.
所有人走進來.
甚至是未信的人走進來.
他聽到的時候.
他明白這首詩歌是說什麼的.
洗教會得造就.
The church will be strengthened.
所以我想每一代.
每一個文化.
每一個社會裡面.
他用的詩歌都有不同.
是要去配合.
那時候的大眾.
到底他能夠聽到的東西.
能夠聽得明白的東西.
所以如果你說.
只能夠流芳百世的詩歌.
才是好的詩歌.
這個在邏輯上是有問題的.
對.
因為有些詩歌.
它只服務那一代.
完了那一代就會結束.
完了那一代就不需要再有.
那些的詩歌是不是好.
都是好的詩歌.

$^{561}$因為它在那個時代裡面.
它說出信息.
下一個時代到了.
它就被淘汰了.
Which is okay.
不一定每樣東西都是流芳百世的.
小弟的港道都不會流芳百世.
今時今日說了.
說完了.
下一代上來說的時候.
下一代去到港道的時候.
都是一樣.
不是所有流芳百世的東西.
才叫做好的.
define什麼叫做好的.
好的就是.
It can strengthen the church.
可以造就教會.
能夠幫助我們來敬拜.
這些就叫做好的詩歌.
By the way.
我說過很多次了.
教會在地上的使命.
不是保留傳統.
我們的使命是說出.
講好我們的信仰.
講好我們的上帝.
向這一代講好我們的上帝.
這就是教會有的使命.
就算回到聖經.
聖經裡面有很多書卷.
有不同的書卷.
那些文學質素都不是相等.
創世記大家知道.
那個敘述的手法是非常高超的.
但如果你回到利希米的話.
就不是這樣的.
利希米是平鋪直敘.
那可不可以說有些好的有些壞呢.
那是不可以的.

$^{601}$所以我們去到敬拜的時候.
唱詩歌的時候.
我們需要多想一點.
怎樣能夠做一個好的敬拜者.
聖經裡面怎樣去看敬拜.
歌曲正如我所說.
不是主要的渠道來表述內容的.
歌曲的時間有限.
字數也有限.
篇幅有限.
那你就問到底什麼叫做敬拜的詩歌.
敬拜的詩歌是做什麼的.
如果留意敬拜的詩歌.
大部分都好像祈禱.
祈禱是向著上主去禱告.
所以如果你要唱詩歌唱得好的話.
我們就需要有一個禱告的靈.
我們才能夠真正敬拜.
別人帶一首詩歌.
我們跟著一起唱.
就好像別人帶一個禱告.
我們一起同心來禱告.
是嗎.
是像祈禱一樣的.
詩歌也是寫來表達情感的.
是嗎.
所以是重複的.
是會重複的.
所謂表達情感的語言.
都會重複的.
是嗎.
例如我太太喜歡我跟她說.
I love you.
我太太說了.
說這麼多幹嘛.
是要重複的.
是嗎.
整句是重複的.
你說真話也好.
說假話也好.

$^{641}$你都需要重複.
你需要再重複.
是嗎.
因為是一種情感的影響.
你就是需要重複.
譬如詩歌有一點像國歌.
無論哪個國家的國歌.
國歌是要表達一種情感.
表達對那個國度效忠的情感.
忠心的情感.
忠誠的情感.
就像美國的國歌.
你唱的時候是Declare Allegiance.
就是你向這個國度表述.
宣告某一樣東西.
你的志向.
崇拜的詩歌其實有一點像.
像愛情的歌.
像國歌等等.
所以是會重複的.
是應該重複的.
不應該這樣平鋪直敘.
系統神學書就沒有人拿著.
系統神學書來宣誓.
是嗎.
國歌 情歌.
這些就是敬拜的語言.
是情感的語言.
所以當你去敬拜的時候.
你是不是需要有些準備.
對嗎.
心情上的準備.
情感上的準備.
重複的時候.
你也需要想想.
重複的時候.
重複有不同方面重複.
你想想上帝在你身上做不同的事情.
你想想上帝自己的本性.
有不同的方面.

$^{681}$所以重複同一句的時候.
都可以有不同的角度去看.
這是我們做敬拜者.
我們需要有的鍛鍊.
詩篇136篇給大家一些品味.
詩篇裡面有很多會重複的詩歌.
這篇詩篇136篇是最典型的.
你們要稱謝耶和華.
因他本為善.
因他的慈愛永遠長存.
你們要稱謝萬神之神.
因他的慈愛永遠長存.
你們要稱謝萬主之主.
因他的慈愛永遠長存.
稱謝那唯一能行大其事的.
因他的慈愛永遠長存.
稱謝那用智慧做天的.
因他的慈愛永遠長存.
繼續下去.
一邊數一邊數.
將整個聖經裡面的歷史.
數給你們聽.
說到出埃及.
說到上帝的供應.
上帝的創造.
上帝的供應.
上帝的帶領等等.
最後結束的時候.
又說到上帝是值得我們稱頌的.
並且每一句他都說什麼.
因他的慈愛永遠長存.
因為他的慈愛永遠長存.
但每一次重複的時候又有所不同.
所以每一個好的讚美的敬拜者.
是需要去學習這件事的.
怎樣去說出一些情感上的語言.
這些亞洲人的文化.
或者華人的文化.
是非常缺少的.
非常缺少.

$^{721}$我們真的沒有什麼情感的生字.
你看看你自己的字典裡面.
有多少時候.
你會用一些情感的生字.
次語.
今天很sad.
今天很傷心.
你說出來.
是不是.
今天很擔心一些事情.
說出來.
我小時候.
我們想做一件事.
我們刻意做一件事.
我們跟小朋友說.
有時候小朋友比較安靜的時候.
我們就說如果日落是美麗的話.
你必須說出來.
你必須用語言表達.
日落是美麗的.
是美麗的.
是很棒的.
你看到風景是美麗的.
你就應該說出來.
以至他能夠用語言表達.
是很棒的.
我做父母有一個.
其中一方面.
就是語言方面.
就是言語上面.
我們需要教小朋友.
去說出用言語去形容一些東西.
他說I am sad.
我就跟他說I am sad today.
我很傷心今天.
因為有些事情發生.
I am sad.
小朋友小的時候.
他正在學習這件事.
但華人的文化我們不會說這些.

$^{761}$我們怎樣說我愛你.
正如我跟小朋友都會說I love you.
I love you.
你可以說得man一點都可以.
我們男人.
我們說情感的話.
我們不是很不聚財.
不知道為什麼.
你可以說man一點.
I love you.
cool一點.
which is ok.
但你要教他去形容一些東西.
用言語去表達.
華人的文化缺少這些東西.
幾乎沒有.
我們怎樣說I love you.
吃東西.
多一點.
我們是這樣.
女兒.
吃東西.
吃多一點.
再吃.
我們是吃東西的文化.
我們這樣表達.
對人的關心.
我招呼你.
吃東西.
有沒有你沒有你.
你今天傷心也好喜怒哀樂也好.
吃.
我們是很實際的.
你不吃很難愛.
餓死你嗎.
我想告訴大家.
在語言上我們需要多做一點功夫.
我們需要學習.
特別在敬拜裡面.
對上帝的表達.

$^{801}$都是一種emotional的response.
由頭到尾整本的聖經告訴我們.
敬拜為什麼會發生.
為什麼人會向上帝敬拜.
因為人認識了上帝的恩典.
上帝拯救了我們.
這位這麼愛我們的上帝.
願意進入我們的心靈.
我們就向他有個回應.
而那個回應是一種情的回應.
情感的回應.
所以說回最後.
我們怎樣去成為一個好的敬拜者.
一個心意的敬拜者.
有三點希望你拿走.
第一 一定要刻意.
你不可以就這樣walk in.
然後你就可以敬拜.
You gotta come.
Prepare.
之前那天都prepare.
好像昨天吃飯.
我和小朋友祈禱的時候.
因為我做爸爸.
我會做祈禱.
我不會推三推四.
很簡單的祈禱.
預備我們明天去敬拜.
是需要刻意的.
我會做我自己.
我覺得一個最好的practice.
我通常都會早一點到.
然後我會在祈禱裡面去準備一下.
你說怎樣準備.
我想起聖經.
我想起上帝在我這個禮拜的生命裡面.
做過的事.
上帝是我的天父.
今天是父親節.
祂才是真正的天父.

$^{841}$無下的天父.
祂願意去救贖我.
我和上帝之間的關係.
relationship 是真的.
It's a reality.
是需要刻意的.
需要準備的.
需要準備自己.
第二方面.
去祈禱.
祈禱有什麼好處.
祈禱不是求問上帝給什麼.
上上簽下下簽.
不是那些.
祈禱是和上帝傾心吐怡去process我的feeling.
向上帝祈禱.
honest 誠懇的跟祂說.
最近可能我真的很累.
我很擔心一些東西.
或者有些東西令我很洩氣.
很frustrated.
等等.
我覺得可以帶到上帝面前.
今天你對上帝感恩.
是有喜樂的.
might as well.
去verbalize it.
用言語去表達.
聖經裡面有很多這樣的例子.
詩歌裡面有很多不同的情感的表達.
甚至我們可以和上帝唱哀歌.
我今天很哀傷.
或者最近很哀傷.
我都可以在禱告裡面向上帝表達.
很多年前911的時候.
911你記得嗎.
911 september 11.
發生的時候好像是禮拜頭.
禮拜一還是禮拜二.
禮拜一 禮拜二.

$^{881}$一發生的時候.
那些敬拜隊.
教會裡面預備敬拜的人.
他們就很頭癢.
突然間我們的設備好像要轉.
對不對.
是一個很傷心的時刻.
所以他們就選擇不同的歌.
將整個敬拜的設計完全改變了.
因為在聖經裡面.
可以讓我們去到上帝面前.
我們來呼喊.
來唱哀歌.
你的情緒可以在上帝面前做一些調整.
第三方面.
我們需要學習.
敬拜絕對不是被動地來觀望.
真正的敬拜是因為我們相信主耶穌.
我們是上帝的子民.
敬拜是因著我們和上帝的活著關係.
而有的反應.
我們和上帝之間有的關係.
和上帝之間能夠認識他.
他去認識我.
這個是我們有的反應.
將我們的注意力放在上帝那裡.
甚至很多時候我們還需要學一些新的歌.
所以我們社交群組.
我們在主日之前.
我們都會發那些將要到主日的詩歌.
學了它.
有新的語言.
有新的禱告.
有新的詩歌.
以至上帝是喜悅我們的敬拜.
我們同心低頭祈禱.
主我多謝你.
我們能夠去敬拜.
能夠去到你的面前.
神愛多謝你.

$^{921}$因為從創世到將來到永恆.
你願意居於我們的生命當中.
你願意住在我們神的子民中間.
在教會裡頭.
當我們向你禱告.
向你敬拜的時候.
你就在我們中間.
主我們相信.
今天我們帶著不同的心情來到你面前.
可能是喜,怒,哀,樂.
可能有很多解不清楚的問題.
但我們知道你是聽禱告的時候.
我們知道你是和我們有一個活著關係的上帝.
我們能夠去信任你.
相信你仍然是作王.
相信你仍然是那位值得我們事奉去敬拜的上帝.
主讓我們的生命.
我們的新心靈.
向著你.
不單向你傾心討耳.
更加我們望見.
神你在我們中間.
作主作王.
掌管一切.
仍然知道上帝你是我們的主宰.
仍然知道上帝你是我們的避難所.
你是那位不變,永恆,值得我們敬拜的上帝.
讓我們去敬拜.
我們的聚集.
是蒙你的喜悅.
祈禱奉祝平安.
願上帝賜給大家平安.
\newpage



\section{5:1-16}
\label{sec:cMH1ZiMWOm8}
\textbf{The Upside-Down Kingdom (Matthew 5\_1-16) - Chee Yie Jong}
\newline
\newline
連結: \href{https://youtube.com/watch?v=cMH1ZiMWOm8}{\texttt{https://youtube.com/watch?v=cMH1ZiMWOm8}} ~~~~ 語音日期: 2024-06-25
\newline
\newline
\hyperref[sec:VHqi_cqG4KM]{\small{< < < PREV SERMON < < <}}
~
\hyperref[sec:index_chronic]{\small{[返順時目]}}
~
\hyperref[sec:index_scriptual]{\small{[返順卷目]}}
~
\hyperref[sec:hDHhZ8zCYYE]{\small{> > > NEXT SERMON > > >}}
\newline
\newline
$^{1}$Good morning, UV..
I hope all of you are able to stay cool this past week..
I'm thankful for just the church here..
Mark brought a portable AC on Monday,.
so it came in clutch for me this week..
But before we get into the word, let's pray..
Heavenly Father, we thank you.
that we are able to come together this morning.
and be here worshiping you and fellowship..
We find another, I just ask that,.
as we're about to open up your word and study your word,.
I ask that your Holy Spirit will be with us,.
that your spirit will open up our heart.
so that we're able to receive what you have for us today..
And I ask that your word will be transforming in our life.
and that it will, yeah, just make a change.
in the way that we live our life..
And I also ask for your spirit will be with me as I preach.
and that you would take away anything.
that is not from you and that your word,.
and your word alone will be preached today..
Pray this in the son's name, amen..
Okay, so I've been at UV for about two months now,.
and usually on Sunday morning,.
right before the English service,.
you can find me serving over at the kids ministry..
And one of the things that I noticed,.
especially with some of the younger kids,.
every time when they're being told to do something,.
they usually like to do the opposite..
Wendy's right here so he can testify to that..
And one of the prime example of this is,.
we have one song that is called "Jesus is Alive.".
So usually at the beginning of the song,.
they will ask them to like hop around,.
clap their hands, stomp their feet..
And then right in the middle of the song,.
there'll be a part where it would say, "Everybody freeze.".
And this started off with just one kid,.
and then more and more kids just like follow along..

$^{41}$So basically, every time when the song starts,.
them to clap their hands or like stomp their feet,.
they would just like stand there and like not move at all..
When it comes to the part where the song told them to freeze,.
they would start like running around.
and then like clap their hands.
and basically just do the opposite.
of what they were being told..
And while some of them think that they're funny.
and like sneaky, not trying to let the teachers recognize it,.
but when you do things in an opposite.
with like basically the rest of the room,.
it is not hard to notice about it..
And so in the same ways, in the Greco-Roman world,.
one of the reason that Jesus causes such a huge stir.
during his time on earth.
is because his teaching and his lifestyle.
is actually completely opposite.
compared to the culture around him,.
that the people cannot help but notice it..
So today we'll take a look at the Jesus Sermon on the Mount.
and what this upside down kingdom of God looks like..
So every country in the world,.
in some ways there is certain guiding principles.
when it comes to the governance.
and the ruling of the country, right?.
So I have to look it up for the US constitution..
These are some of the guiding principles,.
popular sovereignty, limited government,.
separation of powers, checks and balances and federalism..
And so similarly in the kingdom of heaven,.
there are certain guiding principles.
that outline what life in God's kingdom looks like..
And the place where this can be found.
is on Jesus's Sermon on the Mount..
So starting off today,.
I will kick off a series on Jesus's Sermon on the Mount..
And throughout this series,.
we will explore what life in the kingdom of God.
or kingdom of heaven..

$^{81}$I will use this term interchangeably..
So I hope you guys don't get confused..
So kingdom of God, kingdom of heaven,.
in the book of Matthew, they refer to the same thing..
So the Sermon on the Mount.
contains some of the most elaborate teachings of Jesus.
on how we are to live rightly in his kingdom..
And some of you might ask,.
so Chi, you have been talking about God's kingdom.
of this kingdom of heaven for a while now,.
but what exactly is God's kingdom.
or this kingdom of heaven?.
So we can actually spend an entire sermon series.
just on this topic..
But one of the important aspects of God's kingdom.
that I want to highlight today.
is that this kingdom of God is already, but not yet..
It is already, but not yet..
So some of you might already be familiar with this term,.
but the already of this kingdom refers to.
when Jesus died and rose again 2000 plus years ago..
In some ways, he ushers in a new kingdom,.
a kingdom of heaven here on earth..
And so what that means is that in some ways,.
we're actually living in this kingdom of heaven right now..
However, this kingdom is not yet fully being realized..
This is why we still see a lot of brokenness around us..
This is why we still see sufferings and pain,.
those of the, those brokenness around us..
So this is actually the not yet part of God's kingdom..
And this will only come in full realization.
when Christ come again in the future..
And what does this already, but not yet aspects.
of God's kingdom meant for us?.
As we continue to explore this series,.
I think it will be helpful to keep in mind.
that this counter-cultural way of living.
laid out by Jesus in his sermon on the mount.
is something that all believers, we should aim towards,.
we should work towards on this side of heaven..

$^{121}$At the same time, we also have to understand.
that God's perfection can only be attained.
when Jesus comes again..
So the sermon on the mount, even though there are,.
it can seems like a rules or set of regulation.
for us to follow..
But I want to emphasize this,.
that this is not supposed to be a burden.
or a set of impossible rules that no one can live up to..
But this is supposed to be an invitation.
to a way of life that will lead to blessings.
and a way of life that will lead to flourishing as well..
So before I would get into a passage today,.
I just want to give an overview.
for the entire sermon on the mount as a whole..
So Jesus' sermon on the mount.
span across three chapters,.
from Matthew 5 all the way to Matthew 7..
And this is just not some random collections.
about Jesus that Matthew kind of like compiled.
and put it together..
But it is actually a very carefully designed.
and structured section..
And it follows right from Jesus' preaching.
about the kingdom of heaven in chapter four,.
especially in verse 17..
So at the beginning of Matthew four,.
we see that Jesus,.
he was being tempted in the wilderness..
And after he basically able to resist.
all the temptation from the devil,.
he basically began his ministry right after that..
And the central message that Jesus preached.
in his ministry is this..
So Matthew four, verse 17,.
"From that time on, Jesus began to preach,.
"repent for the kingdom of heaven has come near,.
"for the kingdom of heaven has come near.".
And so right following after that,.
beginning in Matthew five,.

$^{161}$the sermon on the mount is basically.
where Jesus explains what he meant.
when he said the kingdom of heaven has come near..
And so I put together this,.
a slide about the structure.
on the sermon on the mount as a whole..
So at the beginning and the end,.
there's the introductory narrative framework.
that basically set up the physical settings.
of this entire sermon on the mount..
And then we get to the main part of it,.
where there's an introduction,.
a main body and a conclusion..
So we'll be focusing on the introduction today..
And then there's the main body,.
which can further divide into three parts as well..
The first part talking about the relationship.
between the law, the Old Testament law,.
to the new covenant of Jesus..
And then the second part talks about religious hypocrisy,.
basically what true and false piety.
or being religious looks like..
And then the third part talks about.
how are we to relate with one another,.
the kingdom ethics or kingdom perspective..
And then this section's on sermon on the mount.
and with basically two different kind of ways,.
two different kind of fruits.
and two different kind of foundations..
Jesus basically is giving his audience an options..
You can either choose the way.
that will lead to fullness and blessings,.
or you can choose the way that leads to destruction..
And so if you're more of a visual learner,.
this is a video from the Bible project.
that they put together and we'll take a look at it..
Hopefully it will work..
- Gospel of Matthew is one of the earliest accounts.
of Jesus of Nazareth..
- And in chapters five through seven,.

$^{201}$we find a collection of Jesus's most well-known teachings.
often called the sermon on the mount..
Here we learn what Jesus means when he announces.
that the kingdom of the skies has come near.
and what it looks like to participate in God's kingdom..
The whole sermon has been given a three part shape..
There's the introduction, then the large main body,.
and then the conclusion..
The introduction announces the surprising.
counterintuitive identity of those.
who were part of God's kingdom..
Then comes the main body of teaching.
where Jesus explores what he means.
when he calls his followers to be righteous..
- What does he mean by righteousness?.
- Righteousness is about living in right relationship.
with God and with others..
It's a character trait that creates justice.
and peace in the world..
Then the body of the sermon is like Jesus's manifesto.
on righteousness examined from three angles.
and three big parts..
How this righteousness relates to the Torah,.
how this righteousness relates to religious practices.
of his day, and how this righteousness is expressed.
in our relationships with God and with neighbor..
Finally, Jesus concludes the sermon by calling people.
to make a choice about how they're going to respond..
- Okay, I see it..
Three main parts of this whole sermon.
and the middle part itself has three parts..
- Right, and then every one of these parts.
itself has three parts..
There is a lot of design in the structure.
of the sermon on the mount.
and we're gonna go over all of it..
In the introduction, Jesus offers a surprise announcement.
to the people standing in front of him..
They are invited to experience the good life.
of God's kingdom..

$^{241}$He first offers nine sayings that redefine.
who is experiencing the good life,.
commonly called the Beatitudes..
Here, we learn that those entering God's kingdom first.
are the lowly outsiders who hunger for righteousness.
and suffer as peacemakers..
- Next, Jesus calls his followers salt of the land..
What's this all about?.
- In the Hebrew scriptures, salt is a symbol.
of God's long-lasting covenant relationship with Israel..
So Jesus is claiming that he and his followers.
are carrying Israel's covenant relationship.
to its fulfillment..
- Next, Jesus calls them the light of the world.
and a city set up on a mountain..
- Jesus is referencing here a promise from the prophet Isaiah.
that one day the inhabitants of Jerusalem.
would reflect God's light and peace and blessing.
out to all the world..
- So Jesus is saying that he and his followers.
are fulfilling that promise..
- Right, and how they go about doing that.
is what the rest of the Sermon on the Mount is all about..
- Okay, so I hope that is a helpful way.
to kind of illustrate this whole entire Sermon on the Mount..
So let's look at Matthew five together..
And yeah, so Matthew five, verse one..
Now, when Jesus saw the crowds,.
he went up on the mountainside and sat down..
His disciples came to him and he began to teach them..
So in this part, we basically find out the who.
and the where for this entire section.
on the Sermon on the Mount..
So the who, of course, there is Jesus up there,.
but then there's two different kind of like distinct groups.
that Jesus differentiates, the crowds and the disciples..
So the crowds are probably people.
that are on the outer circle of Jesus's life,.
and the disciples are in his inner circles..
So Jesus began this by withdrawing from the crowds.

$^{281}$from the outer circle so that he can focus his time.
on his inner circle, the disciples..
But the funny thing is at the concluding part.
of this entire Sermon on the Mount in chapter seven,.
we find that the crowd are present as well..
They probably hear that Jesus is like teaching.
his disciples in the mountain,.
and they kind of like follow him along as well..
So this is the who and the where..
So Jesus, he went up on a mountainside.
and sat down..
So this is pretty much where the Sermon on the Mount.
gets its name, because Jesus is basically preaching.
this entire sermon in the settings of a mountain..
And for us in the 21st century America,.
this might not mean a lot to us,.
but for some of the early Jewish followers of Jesus,.
and even some of the early Christians.
who are familiar with the Old Testament,.
when they see this or when they hear this,.
they will probably immediately recall the event.
in Exodus at Mount Sinai,.
where the Lord basically gave the laws through Moses.
to the people of Israel..
And so right here, Matthew is actually,.
in the way that he structure the Gospel of Matthew,.
he's very intentional in presenting Jesus.
as the faithful Israel,.
who is able to live out God's law perfectly..
So I sort of like make a chart here.
that compare some of the parallel.
between the people of Israel in Exodus.
and also Jesus in the Book of Matthew..
So in the Book of Exodus,.
we have the Exodus event, right?.
Where the people of Israel left Egypt..
And in the Book of Matthew,.
he also presented these stories about Jesus's life..
When Jesus was being born,.
like Herod is trying to basically kill all the babies..

$^{321}$So to escape that, Jesus's parents have to.
basically bring him along.
and like live in Egypt for a little bit..
And so once Herod,.
the Herod that is trying to kill all the babies was dead,.
Jesus and his family then moved out of Egypt.
and back to Jerusalem..
And so this is kind of like the parallel event.
to the Exodus in the Book of Exodus..
And right after the people of Israel left Egypt,.
they wander in the wilderness.
and they were being tempted for 40 years..
And during this time,.
they make a lot of complaints to Moses..
When they have no food, when they have no water,.
they basically came to Moses and asked him like,.
"Why do you even bring us out of Egypt,.
where we would like starve to death?.
We would rather be slave in Egypt.
than to like die here in the wilderness.".
And so Matthew, in some way, he also present Jesus..
Jesus was being tempted after 40 days and 40 nights.
of like fasting and praying..
He was being tempted at the end..
And yet in contrast to the people of Israel,.
Jesus, he was able to resist the temptation every time..
And so, and then after the wilderness wandering,.
we also find in Exodus,.
there is also the narrative about God gave the law,.
the Old Testament, the Torah at Mount Sinai.
to the people of Israel..
And similarly in the Book of Matthew,.
here we find that Jesus preached.
about the kingdom of heaven on a mountain as well..
So there's kind of like the parallel.
of giving the Old Testament Torah from a mountain.
and Jesus preaching about this kingdom of heaven.
from a mountain as well..
So one thing to take note is that the Sermon on the Mount.
that Jesus gave, this is not a new set of Torah.

$^{361}$or new set of law that is supposed to replace.
the Old Testament law..
But Jesus is actually bringing new light.
into this Old Testament law in light of his own coming..
And his teaching is to become the ethical standards.
of righteousness for all the followers in this kingdom..
So now we find out about the who and the where.
for this entire narrative and the following,.
we will find out about the what..
So Jesus begins his sermon by telling the audience.
what a good and blessed life in his kingdom look like..
So all these sayings begins with the word blessed..
And blessed can just mean,.
simply just means having a good life.
or just having a happy life..
But right here in the Sermon on the Mount,.
in the Beatitudes specifically,.
there's probably a nuance of someone.
that have God's favor upon them..
So verse three, he say,.
blessed are the poor in spirit..
If you have never heard of the Beatitudes before,.
or if you've never read this part of Matthew before,.
you might be like, whoa, whoa, whoa, wait a second, Jesus..
Who are the blessed one?.
The poor in spirit?.
I feel like this is a very surprising things.
if we listen to it in our,.
like in the current day and age that we live in, right?.
If we are to write our own Beatitudes,.
we will probably say, blessed are those who are wealthy,.
or blessed are those who have everything that they need..
But for the first century Jewish audience of Jesus,.
living under the rule of the Roman Empire,.
this actually speaks directly to them,.
because this group of people,.
they're basically a group of outcasts.
and nobody in the society..
And yet it is to these people,.
Jesus mentioned that they are the blessed one..

$^{401}$And so what does it mean to be poor in spirits?.
Right here, it refers to a kind of humility.
that understands our own powerlessness..
And in turn, this leads us to completely depend on the Lord..
And it is to these people,.
Jesus has promised that for theirs is the kingdom of heaven..
So I'm gonna get a little nerdy right here..
If you look at the verb that I highlight.
on the second part of the Beatitudes,.
for theirs is the kingdom of heaven..
And then the following is, for they will be comforted,.
they will inherit the earth..
So we see the Beatitudes begin with this divine promise.
that is in the present tense,.
for theirs is the kingdom of heaven..
And it also, later in verse 10,.
it ends with the same promise,.
for theirs is the kingdom of heaven..
And then everything in between.
is actually in the future tense..
So even right here, the Greek grammar.
is trying to present this kingdom of God,.
this kingdom of heaven, as something that is already,.
that we can already experience in the here and now..
But also not yet, that some part of it.
we can only experience when Christ comes again..
So in other words, as believers,.
we are part of God's kingdom now,.
but the fullness of life in this kingdom.
can only be experienced in the future..
And then Jesus continues,.
blessed are those who mourn, for they will be comforted..
This is probably also surprising, right?.
Why are the mourning,.
why are the people that are sad or mourning,.
they are the blessed one?.
But right here, when Jesus talks about mourning,.
there's probably two different kind of causes.
that leads to this kind of mourning..
The first is probably familiar with,.

$^{441}$that the first one is probably something.
that we're all familiar with..
We're mourning because of suffering or hardships,.
or just things or life that are difficult..
Things that are difficult in our life.
that causes us to mourn..
But then there's a second kind of mourning.
that Jesus talks about here..
This is the kind of mourning.
because of an awareness of our own sin..
And we see this a lot in the Old Testament actually,.
when we will read that people will put on sackcloth.
and like put ashes on their head.
as a tangible expression of their mourning.
and also repentance from sin..
So right here, Jesus basically promised that.
when we turn to God in the attitude of mourning.
and repentance because of our own sin,.
God himself, he will forgive us and he will comfort us..
And then verse five continues..
Blessed are the meek, for they will inherit the earth..
And so similar to the first Beatitudes in verse three,.
the meek are those who have an awareness.
of their own circumstances..
And instead of thinking of themselves more highly.
than they should be,.
they understand that ultimately their strength.
and their power comes from the Lord and the Lord alone..
And besides the meek or the gentle,.
they are not just gentle towards themselves,.
but they're also gentle towards others as well..
And when they see that there are some injustice.
or oppressions around them,.
they will want to bring the justice.
and the righteousness of God into those situation..
And these are the people that God promised.
they will inherit the earth..
And to inherit something basically means.
you assume some sort of authority.
or control over them, right?.

$^{481}$So to inherit the earth here means that the meek,.
one day when Christ comes again,.
these are the people that will rule in the new heaven.
and new earth alongside Christ..
And this is the promise that God gave to those.
who are the meek or the gentle..
And so this leads us to the next Beatitudes..
Blessed are those who hunger and thirst for righteousness,.
for they will be filled..
And I'm sure we all have experienced.
some sort of hunger and thirst before, right?.
Maybe some of you sitting right here,.
you're already getting a little hungry.
and thinking about what will the lunch.
or the snacks be later..
But how often do we actually hunger and thirst.
for godly thing?.
How often do we actually hunger and thirst for righteousness?.
Here Jesus say, blessed are those who hunger.
and thirst for righteousness..
And to these people, Jesus promised.
that they will be satisfied, that they will be filled..
So one of the first four Beatitudes that we just look at,.
it focused on, in some ways, a kind of attitude or mindset.
for those who completely depend on the Lord in their life..
This next tree that we will look at,.
we'll actually focus on the actions of those.
who live for God as a result of the dependence on him..
So verse seven, blessed are the merciful,.
for they will be shown mercy..
And this is, again, pretty straightforward..
When we show mercy to others' people,.
we will receive mercy from God as well..
But the questions that, as I was preparing this,.
that I feel like God raised in my heart is this..
Is our day-to-day life a good representation of the grace.
and the mercy that we have received from the Lord?.
I feel like, for me personally, this is,.
I don't think I can fully answer that question.
just because I feel like I have received a lot of grace.

$^{521}$and mercy from God, especially in this season..
But if I honestly take a look at my life,.
a lot of times my life doesn't exemplifies that.
as someone that is merciful..
And then verse eight, blessed are the pure in heart,.
for they will see God..
And so the pure in heart are basically those who love God.
with all their heart, with all their mind,.
with all their strength, and with all their soul..
And these people are basically the one that put Jesus.
as the center of their life,.
put Jesus as the priority in everything that they do..
And Jesus say, these are the people that will see God.
in the new heaven and new earth..
And then verse nine, blessed are the peacemaker,.
for they will be called children of God..
And when it comes to peacemaker here,.
there's, again, two different kind of aspect.
to this peace that Jesus talks about..
To make peace, there's the aspect of peace with God,.
which is kind of like the vertical relationship.
between us and God, and also peace between people..
Basically means when we come together as people,.
are we able to maintain peace in our community,.
in our family, or in our country as a whole?.
Like, are we people that are actively pursuing peace.
in those situation?.
And these people, God promised that they will be called.
children of God..
And in the Old Testament, to become part of God's family,.
or to become children of God,.
you actually have to physically come.
from the line of Abraham..
But in the New Testament, through Jesus Christ,.
we are able to reconcile to God.
and become part of his new covenant family,.
through Christ at the cross..
And as a result of experiencing this peace from God,.
the peacemaker then becomes agent of God's peace.
in the broken world around them..

$^{561}$And Paul highlight this in 2 Corinthians,.
that in some ways, we as God's people,.
we are being reconciled to God,.
so that we can continue to become reconciler.
in our own life..
So this is 2 Corinthians 5, 18 to 19..
All this is from God, who reconcile us to himself,.
through Christ, and gave us the ministry of reconciliation..
That God was reconciling the world to himself in Christ,.
not counting people's sin against them..
And he has committed to us the message of reconciliation..
So, if you have been following along,.
you can notice that there's a certain flow.
to these beatitudes, right?.
This part of the Sermon on the Mount..
So we begin by focusing on the attitude,.
or the mindset of those people that depend on the Lord..
And as we live in submission to the Lord,.
this should, in some ways, lead to a certain actions.
in our day-to-day life, right?.
Like bringing peace to the brokenness around us,.
or showing mercy to the people around us..
And lastly, as we live for God,.
and become his agent of grace and peace in this broken world,.
Jesus also warned us that the world,.
we should expect a certain kind of reaction from the world..
And in some ways, Jesus warned us that.
the world will actually be a hostile place for us..
And we find this in Jesus' last film moments.
with his disciples during the Last Supper,.
when he basically say this to them..
"If the world hates you, the disciples,.
"keep in mind that it hated me, Jesus first..
"If you belong to the world, it would love you as its own..
"As it is, you do not belong to the world,.
"but I have chosen you out of the world..
"And that is why the world hates you.".
So, when we try to live this counter-cultural life.
that represents our kingdom's citizenships.
in God's upside-down kingdom,.

$^{601}$the world will not be a comfortable place for us..
And this is the reality..
But the good news is that Jesus promised,.
right here in verse 10,.
"Blessed are those who are persecuted.
"because of righteousness,.
"for theirs is the kingdom of heaven.".
So, if you remember,.
the Beatitudes opened with the same divine promise..
"Blessed are the poor in spirit,.
"for theirs is the kingdom of heaven.".
And right here, it ends with, again, the same promise..
"Blessed are those who are persecuted.
"because of righteousness,.
"for theirs is the kingdom of heaven.".
And so, this divine promise.
that the kingdom of heaven basically belongs to us.
is actually presently available.
to all people who follow Jesus..
So, let's sit with that for a moment..
What this means is that we can actually get a glimpse.
of life in the heaven, in the here and now,.
while we're still on earth..
I think this is a privilege that we should not take lightly..
And so, the next two verses then continues on.
to elaborate on this,.
but then it makes it even more personal..
So, we can see the switch from the third-person pronoun,.
blessed are those, to a second-person pronoun,.
blessed are you..
So, Jesus say, "Blessed are you when people insult you,.
"persecute you, and falsely say all kinds of evil.
"against you because of me..
"Rejoice and be glad,.
"because great is your reward in heaven..
"For in the same way,.
"they persecuted prophets who were before you.".
So, we just finished with the Beatitudes,.
with all the sayings that Jesus has taught..
And Jesus has shown us what privilege and blessing it is.

$^{641}$to become a citizen of his upside-down kingdom..
He then moved on to provide two metaphors,.
being the salt of the earth and the light of the world,.
to basically show us how citizens of his kingdom.
are supposed to be living on a mission for God..
And just as any earthly citizenships that we have,.
our kingdom citizenships come with a certain privilege.
as well as responsibility..
And this is a both-ends situation..
We can't just enjoy the privilege.
of being part of God's family,.
of being part of God's kingdom,.
without actually carrying any responsibility with it..
So, in other words, we are actually blessed by God.
so that we can continue to be a blessing.
to the world around us..
So, Jesus continues,.
"You are the salt of the earth,.
"but if the salt loses its saltiness,.
"how can it be made salty again?.
"It is no longer good for anything.
"except to be thrown out and trampled underfoot..
"You are the light of the world..
"A town built on a hill cannot be hidden..
"Neither do people light a lamp and put it under a bowl..
"Instead, they put it on its stand.
"and it gives light to everyone in the house..
"In the same way, let your light shine before others.
"that they may see your good deeds.
"and glorify your Father in heaven.".
So, right here, Jesus basically telling his disciples,.
"As salt and light, God has sent us on a mission.
"to show the world who Jesus is,.
"through the way we live our lives..
"And just like the salt that loses its saltiness,.
"if we have no positive impact in the world.
"or in the community around us,.
"we're basically of no use to God's kingdom..
"And on the other hand, as the light of the world,.
"our lives should actively point people to Jesus.

$^{681}$"through the way we live our life.
"and through the works that we have done.".
And so, using these two metaphors,.
Jesus is trying to get this point across..
"As a citizen of God's kingdom,.
"we must make a difference in the world..
"And who we are and whose we are.
"should impact how we live our day-to-day life.".
So, we've come to the end to this whole introduction.
from Jesus' Sermon on the Mount, right?.
And the passage that we just looked through,.
this first 16 verse, in some ways,.
it is very straightforward, right?.
There's no sort of hidden meaning behind it.
or any part that is confusing or hard to understand..
But I think the biggest challenge.
that comes with this passage.
is to be able to live it out, actually,.
in our day-to-day life..
If we take a look at our own lives now,.
can we say that our lives fully exemplify.
these kingdom qualities or these kingdom characteristics?.
You know, one of the things that I like to do.
during my free time is to go on a run..
And usually, on Saturday morning,.
I will go for a longer run..
And so, these past two weeks,.
I've just been spending a lot of time.
just preparing this sermon,.
reading through Matthew, Chapter 5..
And on one of my runs,.
so I was out there on Archer,.
and I came to the stoplight between Archer and Halstead..
So, when the light turns green,.
when I was about ready to cross,.
'cause the pedestrian is green as well,.
so when I was about to cross the road,.
there's this one car basically just ignore.
the pedestrian light,.
and he just make a left turn right in front,.

$^{721}$right when I step on the road,.
he just turn right in front of me..
And at that moment, I was pretty frustrated..
I was like, "What are you doing?.
"You almost ran over me.".
And in my mind, I was thinking,.
"I am on the right here because I follow the rules..
"The pedestrian light just turned green,.
"and you are a car you're supposed to give.
"the pedestrian the right of way, right?".
So, at that moment, I was like,.
"Man, I am justified in my anger,.
"but this guy just doesn't follow the rule.".
So, as a reaction to it, I was like,.
I just do this as he drove past me,.
and then I stare at his car for a little bit..
And right when I cross to the other side of the road,.
one of the Beatitudes just pop out in my mind..
"Blessed are the meek, for they will inherit the earth.".
And right at that moment, I was like,.
"Oh, man, Lord, I am about to preach this sermon.
"in less than two weeks, and yet what I just did.
"is something that doesn't represent.
"what God's people should look like.".
And for the rest of my run, I was basically like praying,.
like asking for forgiveness, and kind of like.
asking the Lord to change me, to change my life,.
so that I can live in a way.
that truly represent His kingdom..
So, I'm sharing this not, yeah,.
just as a way of like being honest of like,.
these standards of righteousness that Jesus gave us.
is not easy, but this is actually an invitation.
to a way of life that will lead to flourishing.
and blessings, even though they might seem.
counter-cultural at times..
We saw kind of like Jesus teach about,.
"Turn the other cheek when people hit you,".
or like, "Walk the extra mile when people ask you.
"to walk like a mile with them.".

$^{761}$So, these are the ways that as God's people,.
as followers of Jesus, this is how we should live our life,.
because we are part of God's upside-down kingdom..
So, what I want to leave you all with is this question,.
is some of these questions to reflect on for this week..
So, do we live our day-to-day life in a way.
that actively brings glory to God?.
And to do it better, what are some of the things.
that we need to stop doing,.
and what are the things that we need to,.
oh, I put two stars there,.
but what are some of the things that we need to stop doing,.
and what are some of the things that we need to start doing,.
so that our lives can be more of a better representation.
of who we are as kingdom, as people,.
as citizens of God's kingdom?.
So, I hope you can take some time this week.
just to pray about it, and reflect on your own life,.
and see what God speaks to you..
So, let's pray..
Heavenly Father, we thank you that you show us.
the perfect examples of what it looked like.
to be a citizen in your kingdom through Jesus..
We thank you that through Jesus,.
we are able to be part of your kingdom..
And I just pray for us right here in this room.
that your Holy Spirit will be with us.
as we continue to live our day-to-day life..
So, I ask that your Spirit will point out to us.
the areas in our life that doesn't represent who you are..
I ask that your Holy Spirit will also be with us.
to give us the strength, to give us the ability.
to live in a way that truly exemplifies.
the God that we represent here on Earth..
So, I just ask that you will remind us.
that we are part of your kingdom,.
and it is such a privilege and honor.
to be part of your kingdoms..
And as your people, you have given us the mission.
to continue to bring your light.

$^{801}$to the darkness in the world around us..
So, I just ask that your Holy Spirit.
will give us the power, give us the ability.
to live in a way that brings glory to you,.
that brings glory to your name..
So, pray this in your Son's name, amen..
Okay, so Mark is not here today,.
so I will be dismissing you all..
But right before that, I just want to encourage you,.
we have some light refreshment outside..
So, don't just leave right after the service..
Stay back, talk to people, catch up with some people,.
fellowship, enjoy some good food..
Okay, go in peace, you be your love..
\newpage



\section{腓利門書 1:4-7}
\label{sec:hDHhZ8zCYYE}
\textbf{愛回家:教會群體的功課  (腓利門書 1\_4-7) - 譚文倩莊傳道}
\newline
\newline
連結: \href{https://youtube.com/watch?v=hDHhZ8zCYYE}{\texttt{https://youtube.com/watch?v=hDHhZ8zCYYE}} ~~~~ 語音日期: 2024-06-25
\newline
\newline
\hyperref[sec:cMH1ZiMWOm8]{\small{< < < PREV SERMON < < <}}
~
\hyperref[sec:index_chronic]{\small{[返順時目]}}
~
\hyperref[sec:index_scriptual]{\small{[返順卷目]}}
~
\hyperref[sec:VPCYBjM4gao]{\small{> > > NEXT SERMON > > >}}
\newline
\newline
腓利門書 1:4-7
\newline
\begin{longtable}{cl}
\hline
\hline
章節 & 經文 (和合本修訂版)\\
\hline
1:4 & \begin{tabularx}{0.7\textwidth}{X} 我在禱告中記念你的時候，常為你感謝我的神， \end{tabularx} \\ \\ \relax
1:5 & \begin{tabularx}{0.7\textwidth}{X} 因聽說你對眾聖徒的愛心，和你對主耶穌的信心。 \end{tabularx} \\ \\ \relax
1:6 & \begin{tabularx}{0.7\textwidth}{X} 願你與人分享信心的時候，能產生功效，讓人知道我們所行的各樣善事都是為基督做的。 \end{tabularx} \\ \\ \relax
1:7 & \begin{tabularx}{0.7\textwidth}{X} 弟兄啊，由於你的愛心，我得到極大的快樂和安慰，因為眾聖徒的心從你得到舒暢。 \end{tabularx} \\ \\
[1ex]
\hline
\hline
\end{longtable}
$^{1}$謝謝.
早安,弟兄姊妹.
很感恩今天可以在這裡繼續跟大家靜悄悄.
其實我也是幾個小時前才回到芝加哥.
很感恩,因為差不多還有30分鐘到達的時候.
突然間屏幕顯示說還有兩個小時才到達.
我想不是吧,是不是弄錯了.
原來是因為天氣的關係,飛機降落不了.
我們就在天空的上空.
我看著圖片在轉圈轉圈.
當我想在這裡埋怨的時候.
我就突然想起安全到達最重要.
轉多少個圈都不要緊.
況且我上飛機之前已經帶了很多零食在袋子裡.
可以熬到今天早上.
所以幾個小時前才可以回到芝加哥,很感恩.
上個月我跟大家分享過《肥李門書》.
一封只有一章經文的求情信.
內容是說一個奴隸,他叫做阿尼西牟.
冒犯了他的主人,肥李門.
普遍推測是這個奴隸偷了錢.
或者是他被派去做一些事情.
但是一去不返.
總之就逃走了.
按照當時羅馬的法律.
協助逃走的奴隸的人.
是要賠償主人的損失.
無論什麼原因.
這封信的主要目的.
是保羅懇求肥李門不計前嫌.
重新接納阿尼西牟.
之前那一課叫做.
《福音如何重新定義人與人之間的關係》.
提到當中三個的關係.
因為福音的介入.
所有關係都要重新定位.
不再是搭檔,不再是父子.
也不再是主人和奴隸.
全世界一起升級.
成為天父的兒女.

$^{41}$在主裡面是弟兄姊妹.
去到總結的部分.
更加提到保羅因為基督耶穌.
因為要成就福音的關係.
以身作則,牧養.
懇求肥李門.
亦同時懇求你和我.
這一群跟隨耶穌基督的人.
願意遵從聖靈的引領.
讓福音去接管你的意願.
去原諒,接納一些虧欠你,得罪你的人.
藉著福音去重建.
重新定義你和人之間的關係.
今天我們再一次看肥李門書.
思考肥李門對眾聖徒的信心,愛心.
對主耶穌的信心.
去思考教會群體的功課.
第二堂《愛回家》教會群體的功課.
我們先一起低頭祈禱.
生命的主,我們的神.
願人恭敬來到你的面前.
領受你的話語.
祈禱奉主耶穌基督的聖名祈求.
阿們.
之前香港有一個處境喜劇.
叫做《愛回家》.
整套劇主要圍繞一家人的生活點滴.
描述各人都因為生活繁忙.
有時會為了一些生活的瑣碎事吵架.
更加引發不少誤會.
但最後他們都能夠和好如初.
因為他們始終都很愛錫這個家.
其實想想教會這個家.
何嘗不是一樣呢.
我們有不同的性格.
不同的年齡層.
我們不同信主的年日.
對大家的期望.
不同的文化.
或者我們處事方法.

$^{81}$或者不同的表達等等.
稍一不慎就會發生衝突.
然後你可能不太想回來見到那個人.
生活在教會裡面.
其實有時想想都挺麻煩的.
你和某某不合時.
回到教會你還會見到他.
好像我們UV一個這麼小的教會.
你一定走不掉.
除非你走去別的教會.
但你想想去別的教會.
又不一定沒有同樣的事情發生.
想想你能去多少個教會呢.
在《肥利門書》裡面.
保羅讚賞肥利門.
對信徒有愛心.
對基督耶穌有信心.
愛心,信心.
在一個人際關係極之複雜的教會裡面.
其實可以發揮什麼作用呢.
跟隨耶穌的基督徒.
你和我又應該怎樣去理解.
和操練愛心和信心.
《肥利門書》一章四至五節.
我在禱告中紀念你的時候.
因聽說你對眾人的愛心.
和你對主耶穌的信心.
這段經文的第五節.
根據希臘原文的排列.
意思是說你對主的愛和信.
以及對所有信徒的愛和信.
五節經文形成一個交叉的結構.
從這個結構去理解這節經文.
是說人對主耶穌基督的信心.
和對其他人的愛心.
是不可以分割的.
意思是沒有一個基督徒說很愛主.
但我不愛人.
或者你是一個很有愛心的人.
但我不愛主.

$^{121}$對主耶穌基督的信心.
是愛其他人的基礎.
這兩樣東西結合在一起.
就正正是我們基督徒.
這三個字的定義的真正意義.
還有.
因為我聽說你的愛和信仰.
你對著主耶穌和所有的聖徒.
這段經文還有一個重點.
在英文聖經裡翻譯是更加清楚的.
就是剛才我說的交叉結構.
剛剛夾住那一句.
Which you have 你有的.
希臘原文Echo這個字.
是包含擁有,緊握,有能力的意思.
保羅在這裡的意思是說.
我聽見你對主耶穌和眾聖徒.
都緊握著愛和信心.
言下之意是肥理門信了基督耶穌之後.
就有這個能力和資源去愛人.
將基督的愛實踐出來.
換言之信主的基督徒.
是已經擁有,緊握,有能力去操練愛心.
你信主信了多少年.
你覺不覺得自己有這個能力.
有這個資源.
你緊握了這件事.
能夠在教會裡操練.
藉著感恩.
保羅強調肥理門過去在愛心裡展現的信心.
而這個愛心在肥理門書裡面的對象.
是說所有的聖徒.
這個所有是包括了虧欠肥理門的奴隸亞尼西牟.
而肥理門其實是做了些甚麼.
被保羅讚賞他那麼有愛心呢.
我們一定要了解當時教會的情況是怎樣.
想像當時的家庭教會.
好像有一個很小的客廳.
裡面只有二,三十人聚會.
當中有幾個比較有米的基督徒.

$^{161}$但大部分都是一些很貧窮的工人.
而肥理門就是當時其中一個很有米的信徒.
他開放自己的家.
提供食物,飲料給人聚會.
而當時的教會發展.
就是靠一些像肥理門這樣的人去建立.
他們有真的信心.
更加有對眾人的愛心.
想想我們不妨去思想一下我們自己的教會.
在這方面我們是否合格呢.
我們又有沒有甚麼缺乏的地方呢.
想當年我住在香港的時候.
在天后一個很小的宣道會裡面聚會.
教會聚會的地方很小.
是我們這裡的一半.
但人數卻是一倍,坐得很密.
而教會每個星期都有一個衣衫難留一陣味的伯伯.
拿著一個很大的紅白藍膠袋進來聚會.
因為我怕了他的味道.
我自己是有一點潔癖.
所以我通常是坐得很遠.
但其實無論你坐多遠.
你是不可能聞到他的味道.
我自己坐在這裡的時候.
很多時候我看周圍的弟兄姊妹.
沒有人會比較遠.
坐得很密.
我想他們是怎樣忍受的呢.
為甚麼可以當沒事呢.
前兩年我再回香港的時候.
就知道這個伯伯去世了.
我才知道原來他住在天水圍.
想當年教會的弟兄姊妹.
去到天水圍的地方傳福音.
帶這個伯伯去信主.
這個伯伯很喜歡教會的弟兄姊妹.
所以他寧願每個星期坐一個多小時的巴士.
來到我們教會去崇拜.
他離開之後.
他的家人就通知教會.

$^{201}$教會的教務同工和弟兄姊妹就去了天水圍.
和他做了一個安息禮拜.
慰問他的家人.
教會對這個伯伯的愛心.
感動了他的家人.
如果你在香港來.
你不會不知道.
天后是一個比較有錢.
比較富裕的一個地區.
相對天水圍是一個比較窮.
比較低收入家庭的社區.
但是想不到天后裡竟然有一個小小的教會.
會去到天水圍傳福音.
不單止這樣.
傳完之後他更加歡迎這個一陣味.
左頭左勢.
他的衣裝和整個教會的人完全不一樣的伯伯.
來到教會當中去崇拜.
這是愛心.
第二點.
對主耶穌的信心.
愛心的動作.
背後正正是弟兄姊妹接受了福音之後.
操練信心的一個表現.
或者你在想.
以後我們教會是否要歡迎一些類似伯伯的人.
教會弟兄姊妹越能夠忍受一些一陣除的人進入我們當中.
左頭左勢的人來到我們裡面.
這就叫做有愛心.
過去十多年因為不同的原因.
我搬過N那麼多次.
不計去探訪一下一兩次的教會.
前前後後在香港和美國我去過不下十個教會.
無論是中國人的教會還是美國人的教會.
其中有一個很恐怖的經歷.
是因為我女兒要去那個教會的英文課.
所以我和我老公被迫要去那一間教會聚會.
那間教會的教導是有問題的.
人與人之間也有很多爭鬥.
我們在那裡捱了兩年.

$^{241}$直至我女兒終於中學畢業上大學.
跟著那個新地我們就走都走不前.
初初去這個教會的時候我沒有什麼發聲.
因為去的時候你也覺得氣氛有點不妥.
總之就是很不舒服.
沒多久就有一個Deacon來跟我聊天.
他說Joyce很多人覺得你很冷酷.
他問我可不可以多點發聲.
跟我們聊聊天.
我說OK 好啊.
其實我不是在裝冷酷.
除了我剛才說到我發聲不太好之外.
我知道這個教會對讀神學的女人是有些抗拒.
但既然他這樣說我就嘗試多點發聲.
這個教會有一個很獨特的文化.
就是一個罵人的文化.
是不是聽都沒聽過呢.
很多領袖就像一個家長.
很喜歡教訓一些平順徒.
我覺得他是很奇特的.
我初初去的時候就覺得很有趣.
新來的舊來的全部都無一倖免.
當然我也不例外.
有一次團契人數不斷下降.
大家在聊天.
我們可以做些什麼去吸引更多人來呢.
我開了四次會每次兩個小時.
說了很多還沒說完.
我已經忍無可忍.
我留意到有些弟兄姊妹很願意帶茶經.
但他們不懂得帶茶經.
很有心但沒有那個訓練.
我就說不如牧師考慮訓練弟兄姊妹帶茶經.
一句話而已.
很有禮貌的.
很小聲的.
結果我被負責小組的兩夫婦亂槍掃射.
教訓了15分鐘.
不是開玩笑的.
我真的看了看錶.

$^{281}$15分鐘還沒教訓完.
亂槍掃射維持在15分鐘.
不是因為他們已經教訓完我.
而是因為我忍不住.
我就說不如不要說了.
我收到了.
知道了就收線了.
那晚回到家之後.
我整晚都睡不著.
我凌晨五時的時候.
發了一個訊息給我的祈禱小組.
叫他們為這間教會祈禱.
也為我祈禱.
我已經不知道應不應該再回教會.
我很心痛.
我當時在想.
為什麼會這樣.
為什麼教會會這樣.
為什麼來到的人不是feel welcome.
不是感受到愛.
來教會應該感受到愛.
我們最出名就是愛.
基督教講愛.
為什麼沒有了這件事.
而來的人反而要步步為營.
否則會被人罵.
我在想我已經不是臭臭地.
不是在阻礙地球轉.
這個教會對新來的人一點恩慈都沒有.
只顧著先來一個下馬威.
告訴你這裡是誰人說了算.
留下來的人要恆久忍耐.
簡直是天方夜譚.
有沒有搞錯.
搬木倒置.
簡直是講客.
之前Mark Sir說了七堂永恆的敬拜.
當中提到很多教會的弊病.
我想有時候我們教會會做很多事.
In the name of God.

$^{321}$時光一定要這樣.
崇拜一定要這樣.
唱詩一定要這樣.
唱這些詩歌.
很多事一定要這樣做.
很多時候只顧著持守一些.
他們覺得很重要的東西.
是在說人覺得很重要的東西.
某一些人覺得很重要的東西.
卻忘記了最基本的愛心.
教會是否需要反省一下.
《肥理門書》這封求情信的下文.
保羅兩次對肥理門提及愛心.
因為保羅想激發肥理門的愛心.
特別將他的愛延伸到.
虧欠他的奴隸亞尼西謀身上.
思想一下這個愛心.
又有沒有激發我們.
去反省一下我們自己.
對教會又有什麼挑戰呢.
或者如果我們用肥理門這間教會.
或者肥理門這個人.
去為我之前提到的教會去把把脈.
他們沒有愛心是一個病徵.
但是病因是什麼呢.
可能有很多不同的原因.
會不會是驕傲呢.
會不會是有階級觀念.
或者他們經常會衡量.
對方是否值得我去關心他.
值不值得我去關心他.
值不值得我去付出.
對我又有什麼好處呢.
不去思想的時候.
有時我們不單是沒有愛心.
連對人基本的尊重都沒有.
肥理門書一章六節.
願你與人分享信心的時候能產生功效.
讓人知道我們所行的各樣善事.
都是為基督做的.

$^{361}$分享信心這四個字.
在希臘原文是Kononia.
在新約聖經裡.
這個字是被翻譯成團契.
其實在我們俗世的社會裡.
它們時時有用的.
例如商人有商會.
Business Association.
有經濟利益的人走在一起.
同鄉有同鄉會.
因為大家同聲同氣.
但在聖經裡Kononia的意思.
不是說大家有沒有經濟效益.
或者我們有沒有血緣的關係.
只有一個原因.
就是在說基督耶穌.
Kononia眾聖徒聚集一起參與.
付出,分享.
當中包括敬拜,事奉,傳福音,建立教會.
保羅在這裡究竟想說些什麼呢.
就是肥理門所說的團契.
不是以經濟或者社會地位.
個人利益等等為基礎.
而是單單以基督耶穌為基礎.
與俗世價值觀很不同的群體.
這就是獨有文化的教會.
的確很多教會都有自己的文化.
通常來教會的人都要投入教會的文化.
不要說是否要摺疊一整身的錘子.
拿著一個紅白藍膠台.
左頭左勢的伯伯.
相反每個人進來的時候.
你都要自我調節.
適應那個教會的文化.
但想想.
教會應該有一個什麼樣的文化.
或者有時候教會是否太著重保護.
教會自己的文化.
不留心的話.
我們可能已經慢慢失去了起初的愛心.

$^{401}$有些人只顧著保護自己當中的一席位.
或者覺得現在都不錯.
不是很理想.
但都習慣了.
而失去那個敏感.
有時當有新的朋友來到我們當中.
我們有沒有預備好迎接新的來賓.
有些事不是很妥.
但我們不想作出什麼改變.
所以就任由他吧.
It's OK.
所以所有無法適應這些文化的人.
就會自己慢慢離開.
我剛才提到的教會就是這樣.
其實我覺得很可惜.
我在這裡長大.
很多時候在傳福音的時候.
我們是出去傳福音的.
但我住在那邊的時候.
那邊是沒有一個團街的.
所以很多中國人會去到教會裡面.
去找Community.
所以你不用出去找人.
那些人自動送上門去教會.
但我看著那些人左門進入.
被人罵到在右門出.
很奇怪.
有沒有想過教會其實好像一個打開門做生意的地方.
歡迎光臨.
我們這班人很棒.
我們推銷的福音很棒.
吸引到你想再回來.
我們UV又是一個怎樣的教會.
我們是不是要隨時候駕.
大家都敏銳.
第一次來到我們這個教會的人.
來到我們這個淑寧家的人.
我們Make sure 他們Feel Welcome.
甚至或者你自己每個星期日回來.
你Feel 不Feel Welcome.

$^{441}$I hope you Feel Welcome.
不說其他日子.
只說星期日.
我們有Welcome Team的弟兄姊妹在門口歡迎.
說的是風雨不改.
我們有現衣帶領的Cafe Team.
每個星期日預備好早餐.
讓大家都不用考慮.
我還沒吃東西.
要馬上買早餐才來教會.
你來到已經有咖啡等著你.
有早餐給你吃.
還有弟兄姊妹其實很早來到.
抹桌子,抹椅子.
收拾好東西,消毒.
還有AV Team.
我們的Worship Team.
帶領兒童主一學的Wendy師母和秀蘭.
很多人很早來.
預備好這個家.
歡迎你回來.
我們投入了很多努力.
歡迎每一個人回家.
教會裡面除了我們的教務木桐公.
還有這班弟兄姊妹.
用愛心擺上時間去服事建立教會.
更加是因為教會要回應福音的使命.
想一些不認識基督的人.
能夠藉著我們的愛心服事去認識基督.
想一想我們是不是這麼有吸引力.
叫人嚮往加入這個朝著對主的愛童信.
對所有聖徒的愛童信前進的群體.
《肥利門書》一章六節.
願你與人分享信心的時候.
能產生功效.
讓人知道我們所行的各樣善事.
都是為基督做的.
若果教會真的要講一個特定的文化.
很肯定教會其實只有一個文化.
就是福音的文化.

$^{481}$以基督為基礎的文化.
全福音的基礎是一個被神觸動.
被神改變的生命自然的流露.
神在肥利門的生命裡面做了一些美善的事.
觸動了肥利門.
要和其他人分享他的信仰.
藉著他對人的愛心.
其他人就有機會來到基督的面前.
全福音難.
我相信其中一個原因.
是因為教會群體沒有好好操練愛心.
信仰的分享和我們的見證沒有知行合一.
基督徒講一套做一套.
叫福音失去吸引力.
肥利門書一章七節.
弟兄啊 由於你的愛心.
我得到極大的快樂和安慰.
因為眾聖徒的心從你得到書場.
保羅讚賞肥利門有信心和愛心.
信心加上愛心的行為.
祝福在他家裡敬拜的弟兄姊妹.
建立教會.
一個不只是用口去傳福音.
和講愛心的教會.
能夠在社區裡造炎造光.
有影響力 有吸引力.
榮耀上帝.
有晚我煮飯的時候.
煮了飯後我就回頭看我家.
在我看的時候.
我心裡有一個莫名的感恩.
為到神給了我一個這麼好的老公.
我家裡的每一樣東西.
都是我老公多年來辛辛苦苦工作.
賺錢買回來的.
我在家裡就負責持守好這個家.
打理好我們的家頭世母.
多年來我們分工合作.
互相扶持.
一起成長.

$^{521}$建立我們的家.
我們教會這個淑齡的家又如何呢.
你想要一個怎樣的教會.
一個怎樣淑齡的家.
你是不是很多時候回來.
只看到這個不妥 那個不妥.
不喜歡這個 不喜歡那個.
其實你想一想.
當你這樣想的時候.
那個專注力其實去了你自己那裡.
我們教會的存在不是這樣.
我們是要傳揚基督.
我們要愛其他人.
那個關注到別人那裡.
我們要怎樣去建立這個教會.
或者你再去想.
這個教會的家.
你有沒有份.
如果你有份.
不如一起夾手夾腳.
去建立一下這個家.
想一下怎樣可以操練愛心和信心.
一起去攜手去做好它.
建立一個你和我都很喜歡回來的家.
《肥理門書》這封信.
重點是講述福音要重整人與人之間的關係.
以及福音要成就的工作.
也提供了一些原則.
教導跟隨基督的信徒.
要怎樣在我們的群體裡面.
去實踐基督教的原則.
求神幫助我們這一班已經得著福音的人.
不要忘記我們教會存在的使命.
我們操練信心.
更加操練愛心.
一起同心建立一個充滿美麗的教會.
讓一些未認識基督的人.
來到的時候.
能夠覺得我們與眾不同.
這個地方是很溫暖的.

$^{561}$很想再回來.
更加能夠藉著教會我們的群體.
讓這個來到我們當中的人認識基督.
而像我和你一樣.
已經認識基督的人.
我們也喜歡回來這個屬靈的家.
我們一起努力去建立.
我們一起低頭祈禱.
我想給你一點時間.
在自己的座位上為自己禱告.
如果剛才我一直在講道的時候.
你心裡也有很多思想.
就為這些思想禱告.
求神幫助你在教會裡.
不只是進來出去.
求主幫助你.
福音能夠改變你的生命.
你回來的時候.
你的心是要建立教會.
是要敬拜神.
七號主我們的神.
藉著《肥理門書》給我們的提醒.
提醒我們每一個人.
都要活出信心.
努力建立一個充滿福音文化的教會.
讓踏足這個家的人.
都能夠享受弟兄姊妹的愛.
得著祝福.
屬靈生命被建立.
祈禱奉主耶穌基督的聖名祈求.
阿們.
祝願大家平安.
記緊操練愛心,信心.
散會.
\newpage



\section{}
\label{sec:VPCYBjM4gao}
\textbf{Back to Basics (Psalms 119) - Rev. Mark Chan}
\newline
\newline
連結: \href{https://youtube.com/watch?v=VPCYBjM4gao}{\texttt{https://youtube.com/watch?v=VPCYBjM4gao}} ~~~~ 語音日期: 2024-07-03
\newline
\newline
\hyperref[sec:hDHhZ8zCYYE]{\small{< < < PREV SERMON < < <}}
~
\hyperref[sec:index_chronic]{\small{[返順時目]}}
~
\hyperref[sec:index_scriptual]{\small{[返順卷目]}}
~
\hyperref[sec:KO8k0_vDZOY]{\small{> > > NEXT SERMON > > >}}
\newline
\newline
$^{1}$- Well, yes, so this Thursday, you're welcome to join us.
for the barbecue hangout, all right?.
We'll make sure you have good meat, okay?.
I'm going to Costco tomorrow, so we'll shop enough food.
for all the people, but you have to RSVP, okay?.
You have to tell me, or maybe Amanda, too..
You can tell her that you are coming, okay?.
So we know you have to feed your stomach,.
we have to prepare fireworks for you, okay?.
We will have some fireworks for you, okay?.
Just a little bit, so don't burn down my house..
Well, being Asian, every July 4th is a very interesting day..
My neighbors, most of them are different ethnic background,.
and they spend quite a bit on fireworks..
But I'm Asian, so I don't do this stuff..
But it's going to be fun, though, right?.
We'll do a little bit, okay?.
We're going to have fun Thursday, Thursday evening, all right?.
Again, I want to remind you to pray for one another..
We're going to start a new series on back to basics..
What does it mean, back to basics?.
You know, in our daily life,.
there are a lot of basic principles we have to master, right?.
Anyone do workout from time to time, or regularly, anyone?.
Only a few hands, okay, great, great, okay, back there, too..
All right, good, good, good, yeah..
There are some basics you have to master, right?.
There are some basics in terms of diet,.
in terms of what you take in, and how you work out,.
and all the activities have some sort of basic principles..
Anyone plays music instruments here?.
Music instruments, yeah..
You remember in those days when you learn.
to operate an instrument,.
a musical instrument, make good sound out of it,.
there are some principles or some basic exercise.
you have to just drill, right?.
It's boring, but in order to do it well,.
in order to enjoy it, you have to master the basics..
And I have this idea about back to basics series.

$^{41}$because I saw this picture not too long ago..
And someone posted on social media,.
it's talking about a building,.
you know, a high rise or the taller the building,.
the deeper the foundation is needed, right?.
The foundations cannot be seen..
It is underground, right?.
Many, many feet underground, many meters underground..
And you have to build that foundation solid enough.
and deep enough in order to build up your life, right?.
Story after story so that the building.
will not easily damage or collapse..
So, you know, talking about architectural terms,.
you need the basics, right?.
You need the foundation..
Talking about a spiritual life,.
talking about your relationship with God, right?.
All of us who are in Christ,.
all of us who are following Christ Jesus our Lord,.
we begin our spiritual life from the first moment.
we came to know Christ..
And from the first moment we came to know Christ,.
we know that the spiritual life needs our,.
need us to master the basic..
The basic things in our spiritual life are necessary.
so that you can continue to enjoy.
your relationship with God,.
so that you can continue to grow in faith,.
solid faith in Christ..
So this is part one of the series.
and I have two important points..
First point is the importance of the Word of God..
Second point is the power, the strength of the Word of God..
Now it's very simple when you think about the Word of God.
as something basic for Christian living,.
for our following of Jesus Christ,.
our spiritual relationship with God..
The Word of God..
You cannot do without the Bible..
We are people of the book, all right?.

$^{81}$We are people of the book..
We need to understand the Bible..
We need to enjoy reading the Word.
or train ourselves so that we can understand.
the biblical messages in the Bible..
Now everyone has a Bible today, right?.
You just download an app, right?.
And you've got all the versions in the world..
Whatever languages you have or you master,.
you can read the Bible in almost,.
almost, actually many of the known languages..
So think about the Word of God..
How important that is..
You think about the Bible,.
there are many, many passages talking about.
the importance of the Word of God in our lives..
New Testament or Old Testament..
Let me give you some samples..
1 Peter 2, verse one..
So put away all malice and all deceit.
and hypocrisy and envy and all slander..
Those are the values of the old life..
Before we came to Christ, put away all the sins..
And Peter said, like newborn infants, newborn infants,.
long for pure spiritual milk,.
that by it you may grow up in salvation..
If indeed you have taste, that the Lord is good..
Like a newborn babies, right?.
We have tons of babies out there, right?.
Oh man, they are noisy, right?.
Whenever I see them, you know,.
I remember the many, many years ago,.
I have my son and daughter and you know,.
I have to give them milk, right?.
I have to take care of them..
And what they yearn for is something to grow, right?.
Now, Rachel and Nathan are adults now, right?.
You know how heavy they are?.
I cannot pick them up now, right?.
They're more than 100 pounds, right?.

$^{121}$120 pounds, 30 pounds or even..
Well, how did they grow up?.
You know, how did they grow up?.
Because they have food, right?.
As babies, they need spiritual milk,.
spiritual nourishment, right?.
They need the nutrition..
They need the food in order for them to grow..
Same here, this is the analogy applied to Christian life..
As we go about our Christian living,.
we need the spiritual milk..
And what Peter is talking about.
is actually the Word of God, pure spiritual milk..
Without it, we cannot grow, right?.
You realize that from,.
in your spiritual journey,.
whether you're a new Christian.
or you have been Christian for many, many years like me,.
I've been Christian for decades..
Whenever I ignore the Word of God.
or go without the Word of God,.
then I feel weak..
I feel I don't have the perspective.
to live my Christian life..
But whenever I am filled with the Word of God,.
then I know how to live as a Christian..
And at least I know why, right?.
There are many, many people who have gone astray.
from their faith..
It's because they go about their life, their Christian life,.
without the spiritual milk, right?.
So that's one of the samples we can find.
from the New Testament..
There are many, many different passages,.
especially in the Old Testament too..
Now in the Old Testament, there's a book called Psalms..
Psalms, how many Psalms are there in the book of Psalms?.
It's a collection of 150 poems or songs or Psalms.
that the poets in the ancient time.
sing about and pray about, worshiping God..

$^{161}$And many of the Psalms actually talk about the value,.
the importance, and the wonder of the Word of God, right?.
Can you think of any samples?.
Psalm one is one of them, right?.
But it's the man who does not walk in the counsel.
of the wicked or stand in the way of sinners.
or sit in the seat of markers,.
but his delight, his enjoyment,.
he's enjoying something..
It's the law, the Torah of the Lord..
On his law, he meditates regularly, day and night..
Now here comes the analogy..
He, if he does that, he is like a tree.
planted by streams of water..
And he's gonna yield his fruit in his season,.
and the leaf does not wither..
Whatever he does prospers..
What a beautiful picture, right?.
A tree planted by the river, by the streams of water..
And he is called blessed, right?.
He will receive the blessing of God..
Psalm one is one of the samples..
And then Psalm 19 I preach about, right?.
About creation and about the word of the creator..
There are many Psalms like this..
150 Psalms, there are many, many Psalms.
that talk about the word of God that is so important..
Today, I would like to pay more attention to Psalm 119..
Psalm 119 verse one,.
beginning with the same word as Psalm one does..
Psalm 119 verse one..
Blessed are they whose ways are blameless,.
who walk according to the law of the Lord..
Verse two..
Blessed are they who keep his statutes.
and seek God with all their heart..
Right?.
Blessed..
Just like Psalm one, talking about the word of God,.
the importance of the word of God in the people of God..

$^{201}$And Psalm 119, verse one, started with the same word..
Blessed are those people whose ways are blameless..
That does not mean, blameless does not mean perfection..
No one's perfect, right?.
You're not perfect, I'm not perfect..
You're sinners, I'm a sinner, right?.
We're all broken, we need the redemption of Christ..
We all need that, we all need the grace of God..
But what does blameless mean?.
Blameless means he's dedicated..
Someone who is dedicated and is faithful, right?.
Is truthful..
To do one thing is to walk according.
to the Torah of the Lord, Torah of Yahweh..
Blessed are they who keep his word.
and seek God with all their heart..
Psalm 19 is actually the longest chapter in the Bible..
The longest poem, longest Psalm in the Bible..
A total of 176 verses..
Why?.
Why 176?.
Because in, well, the Old Testament.
was written primarily in Hebrew..
Hebrew has 22 alphabets, right?.
How many alphabets are there in English?.
20, right?.
Yeah, 26, right?.
(laughing).
26, yeah, you're right..
In Hebrew, it's not 26, it's 22..
From aleph, the first letter, to the last letter, taf..
22 letters..
And in this Psalm, the amazing thing,.
the beauty of it is that every eight verses,.
all eight lines begins with the same letter..
And you go from, like, English go from A to Z..
In Hebrew, go from aleph all the way to taf,.
this 22nd letter..
So there are 22 sections in this longest chapter.
of the Bible, and each section is eight verses..

$^{241}$We call them acrostic Psalm..
You know, it's just a big word for alphabetical Psalm, right?.
Alphabetical Psalm, right, from A to Z..
So it's the longest chapter of the Bible..
Why is this longest chapter of the Bible.
written in such a way?.
It's because when you read through Psalm 119,.
you get the sense of perfection..
You get a sense of unity..
It's something to look at, something that is beautiful..
Talking about what?.
The entire Psalm 119 praises or wonders.
about the Word of God..
So the Word of God, the Bible, the Word of God,.
that is revealed on the pages of Scripture..
It's like a beautiful painting..
You look at the painting, you go to an art museum.
in Chicago, some big museum, you go in there.
and you see some life-size painting, oil painting..
It's beautiful, and you can appreciate it..
Psalm 19 has the same message..
The Word of God is like a beautiful artwork,.
and you can look at it and you can appreciate it,.
and you can sense that there's unity in the Word of God..
There's perfection in it, there's beauty in it,.
just by looking at it..
One of my professors in the days of seminary,.
his name is Bruce Waukee..
He's very old now..
He's already retired, but he's one of the.
Old Testament experts in the world..
And I remember in one of his classes,.
he told us every single word in the Bible,.
of course he's an Old Testament professor,.
always focused on his Old Testament, right?.
So he's saying every single word in the Old Testament.
is worth your reflection..
Every single word, right?.
God does not waste any word in the Scripture..
So the Word of God is like a beautiful painting, right?.

$^{281}$There are broad brush and there are details there,.
and you can look at it and appreciate it,.
and that is the importance of the Word of God..
Not only the Word of God gave birth to us.
and made spiritual life,.
make our spiritual relationship with God possible,.
it is also beautiful and perfect..
And the second point I wanna point out.
is the power of the Word of God..
As you read through Psalm 119,.
not only is it beautiful poetry,.
the content tells you that the psalmist.
needs the Word of God in the midst of suffering, right?.
So Psalm 19 is actually a lament,.
an individual lament, right?.
Some songs are played in minor keys,.
some songs are in major keys,.
but Psalm 19 is in kind of minor key,.
because as you read through the Psalm,.
more than 17 times,.
the psalmist mentioned his suffering to God,.
and he was praying so hard..
His struggle with his pain or his suffer with his struggles,.
whatever that is,.
and he's trying to pray to God and draw strength from God..
But when he suffers, how does he respond?.
Let us read some of the verses here..
Verses 85 and following..
The arrogant dig pitfalls for me,.
contrary to your law..
All your commands are trustworthy..
Help me, God, for men are persecuting me with thou cause..
They almost wipe me out from the earth,.
and I have not forsaken your precepts..
That's part of his prayer, part of his long prayer..
He said the arrogant people, the enemies,.
are persecuting him, pursuing him, without any cause..
So this guy is really suffering..
This guy is really struggling with a lot of challenges,.
a lot of issues in his life, right?.

$^{321}$And something we can relate to, right?.
Probably you have some enemies too..
Probably you are suffering as well at this point..
Either the suffering can be described.
in terms of emotions, mental, or even physical, right?.
The psalmist is crying out to God, I am suffering..
People are digging traps for me..
They are setting traps for me..
They are trying to wipe me out, right?.
But the psalmist is trying to draw closer and closer to God..
He said all your commands are trustworthy, however,.
and you can help me because I have not forsaken.
your precepts, your teaching..
So you should think about whenever you are struggling,.
there's one thing you can do, you cannot do without..
It's the word of God..
Let's read some of the verses..
Let's continue..
Verses 50 and 165..
Psalm says, my comfort in my suffering is this..
Your promise, your promise, O God, perseveres my life..
Great peace have they who love your Lord.
and nothing can make them stumble..
So he's trusting in God by reading the word of God,.
drawing close to God, finding comfort in his suffering,.
finding peace in his suffering..
So how does this psalmist respond to suffering?.
How does he respond?.
He turns to the word of God..
You struggle, the one thing you do, you turn to God..
You read the word, right?.
You know, as a pastor, I often go to hospital.
for visitations, visit the sick and the weak,.
those who struggle, those who suffer..
And so often, when a friend or a neighbor are on sickbed,.
you just find that your human words are limited, right?.
What can I say besides sickbed, right?.
What can I say to give people a little bit sense.
of comfort and peace?.
I remember one time some years ago,.

$^{361}$I was visiting a patient who was suffering.
from terminal disease, cancer..
And he was in terrible pain..
And I was in the hospital room.
and there was no hope for him to recover, no hope..
He just had days remaining and counting down..
And what can I say to bring comfort.
and to bring a little peace to the dying, right?.
To those who are suffering..
And so often, as I told you,.
words coming out of my mouth are so limited, right?.
But I find it's very interesting and very amazing.
that at the end of the visitation,.
I usually read something from the Bible..
And maybe Psalm 23, and you know this Psalm, right?.
It's a very familiar Psalm..
Some of the Psalms, some of the words from the Bible,.
I end my visitation with those words..
And somehow the person who is suffering,.
after hearing the word of God, right, find some peace..
You can look into his eyes.
and he was experiencing something,.
some peace beyond human understanding, right?.
That's exactly what the Bible says, right?.
That's exactly what the Bible promises,.
that if you read the word of God, right,.
in the midst of suffering,.
you will experience some peace.
that is beyond the world can offer..
It's offered by God in Christ only..
I remember when I suffer in different seasons,.
when people, whatever people say,.
you know, I find that their words are empty..
Now, there's no ability to comfort me..
But when I turn to the word of God in my quiet time,.
in my own personal prayers,.
I regain strength, right, to do life..
I regain strength and peace to go on..
Same thing here, Psalm 119, verse 111..
It said, "Your statues are my heritage forever,.

$^{401}$"and they are the joy of my heart," right?.
Someone who is suffering,.
but reading the word of God, drawing close to God,.
knowing that the word of God has power to heal,.
the power to comfort, the power to grant peace, right?.
Some years ago, Christianity Today magazine.
interviewed someone called Billy Graham..
You know, those of us who are,.
well, many of you are younger, right?.
I'm one of the old guys here..
You know, I still remember Billy Graham..
I went to one of his evangelistic gathering, right?.
Billy Graham was the greatest evangelist in 20th century,.
and he has led many, many people to Jesus, to Christ..
And the Christianity Today interviewer.
asked Billy Graham one question..
After all of the great works God worked through him,.
leading many people to Christ,.
millions of people to Christ,.
and the interviewer asked him a question..
He said, "Billy, if you could restart your ministry,.
"restart your work, what would you do?.
"Is there anything you would do differently.
"if you restart your whole thing,.
"you know, your whole ministry, your whole work for God?".
And Billy Graham said, "I would actually spend more time.
"to study the Word of God before I preach.".
That's pretty interesting, right?.
I would study the Word of God and preach less, right?.
You know, at times, we envy the people in the Bible,.
like Abraham, right?.
Abraham, in his faith journey in God,.
God appeared to him visually and spoke to him directly,.
right, and most of us,.
we don't usually have that kind of experience,.
direct encounter with God, right?.
And God just show up in your dream or in your bedroom.
and say, "Hey, you know, I want you to do this, right?.
"I want you to stop doing this," right?.
We usually don't have that kind of spiritual encounter.

$^{441}$with God..
However, if you go back to the story of Abraham.
in the book of Genesis, Genesis 12 to 25,.
and you count how many times God actually appeared.
to Abraham and spoke to him directly, how many times?.
Let me do the calculation..
Abraham started to follow God when he was 75 years old,.
and he passed away at the age of, very old age,.
175, so he followed God for 100 years,.
for almost a century, and in that 100 years,.
we counted from the book of Genesis, chapter 12 to 25, right?.
How many times God actually appeared to Abraham.
and spoke to him directly, appeared to him visually?.
How many times?.
You wanna make guess?.
Eight times..
Eight times, eight times only in his entire faith journey..
Only eight times God appeared to him directly.
and said something to him,.
and if you go to this eight encounters,.
all of the eight encounters are not QA sessions..
It's not like God show up and,.
"Hey, you wanna ask me some questions?.
"You know, I'll take any question..
"I'll give you answer.".
No, it's not a QA session with Abraham..
There's no QA session, and so often,.
God spoke to Abraham directly, not according to his term,.
right, and so often, God says something to Abraham,.
and Abraham did not wanna hear that, right?.
Like, "You have to wait, okay?.
"I'll give you a son, but you have to wait," right?.
You like the answer?.
Of course not, right?.
It's testing your patience..
And so often, God show up to Abraham and spoke to him,.
and leaving him, after the encounter,.
leaving him even more confused, right?.
So if you would choose,.
would you want to receive the word of God like Abraham did?.

$^{481}$I wouldn't do that, right?.
It's so puzzling, and it's so mind-boggling, right?.
You show up, and then you said something to me,.
I didn't wanna hear that, right?.
But you said, "Trust me," right?.
And then you disappear..
Eight times only, in 100 years..
I wouldn't wanna be Abraham..
You know, I wouldn't want to envy him..
But for me, being people in Christ now,.
having the whole Bible already revealed to us, right?.
The Old Testament and New Testament.
is a complete revelation of God's will,.
clear and loud, right?.
The Bible, we all have it,.
has all the clear principles of following God..
All the moral, spiritual requirements,.
the ways of God are clear,.
clearly spelled out on the pages of Scripture..
And that's why we have to say,.
we are the people of the book..
We need to follow the book, right?.
We need to learn how to read it,.
and read the word of God more and more..
Dive into it, understand it,.
and struggle to interpret it,.
to understand it, and apply it into our lives..
So today, as the first part of our series,.
I would like to end with two take-homes..
First, we all have 24 hours..
Take the best time period of your day..
Some of us are morning persons,.
some of us are not, right?.
Your best hour, what would that be, right?.
Maybe at noon, or maybe at night,.
whatever that is, take at least 10 minutes a day.
to read the Bible..
If you don't know where to start, okay,.
let me suggest to you, if you're new believers,.
start with one of the gospels,.

$^{521}$Matthew, Mark, Luke, and John,.
either one of the four, right?.
You can start with Gospel of John,.
read one chapter a day..
It will take you less than 10 minutes, right?.
If you're believers for quite a long time, right,.
and you don't know where to start now, right,.
well, you can start with the Old Testament,.
the book of Psalms..
There are 150 Psalms, right?.
You can read one Psalm a day..
Take at least 10 minutes in your,.
what we call quiet time, or devotional,.
to read part of the Bible, all right?.
I'm not asking you to read it in one sitting,.
cover to cover..
No one can do that, right?.
Take at least 10 minutes a day to read the Bible,.
and to meditate on it, right?.
Take it, store it in your mind, right?.
Remember it, kind of remember the sketch of the passage,.
and bring it with you in your work, right?.
You go about your daily life,.
go grocery shopping, or whatever you're doing..
You know, just from time to time,.
remember that, and to reflect upon it..
And today, I take the suggestions.
from one of our brothers, James,.
that we are going to do the verse of the month, okay?.
Verse of the month..
I'll take one verse of the Bible,.
and actually, it's not verse of the month,.
verse of the week, okay?.
Verse of the week..
I'll take one verse of the Bible every single week,.
and we're going to meditate on it,.
and we're going to memorize it, okay?.
We're going to recite it together as a congregation,.
as a family of God, we're going to read it,.
and remember it, and memorize it, right?.

$^{561}$Next week, we come back, we're going to reveal,.
and it's going to be another verse for next week..
How about that?.
Pretty good suggestion, right?.
Yeah, thank you, brother..
Yeah, so our verse of the week is this..
Your word is the lamp to my feet,.
and the light for my path..
Can we read it together?.
Have Lucy come up here..
Let's read it together, Psalm 119, 105, okay?.
119, 105, you have to memorize the reference as well, okay?.
It's from Psalm 119, verse what?.
105, okay?.
Let's recite it together, all right?.
Psalm 119, verse 105..
Your word is the lamp to my feet,.
and the light for my path..
Wonderful word, right?.
Yeah, the word of God, may it be the lamp and the light.
for our daily walk, right?.
You may face different trials and temptations,.
but may the word of God feed you..
He can give you strength, right?.
Give you the power to live the life God intend to you,.
to live..
Let's all pray together..
Father God, we give you thanks for your word..
I wanna give you thanks because you are the God.
who still speaks..
You speak through your word..
Speak through your scripture..
And I just pray for all the brothers and sisters here..
May we continue to dig into your word,.
to meditate on it day and night,.
so that we will receive blessing from you,.
to grant us strength, grant us the ability,.
the amount of grace that we need.
to go on tackling with different challenges in our lives.
so that we can live the spiritual life.

$^{601}$you intended us to live..
Thank you for calling us to be your people..
We are indeed the people of your word, people of the book..
May your word dwell in our hearts and our mind,.
and be our guidance, be our principle..
Indeed, your word is the lamp unto our feet.
and the light unto our path..
Help us to rely on you, to draw close to you.
by reading your word, to study the word together,.
to have Bible study to remind each one of us.
the reality of your spiritual life,.
reality of our relationship with you..
Thank you for being our God who loves us so much..
Again, we wanna give you thanks, praise you,.
praise you for the fact that you reveal your will.
through scripture..
Thank you, God..
In Christ's name we pray, amen..
(gentle piano music).
Well, we'll end our service here,.
and we have some refreshment out there.
in the cafe or coffee or tea or just water..
So stay behind for some time of fellowship.
and say hi to your new friends.
and say hi to our neighbors..
\newpage



\section{3:1-9}
\label{sec:KO8k0_vDZOY}
\textbf{Back to Basics \#2 (2 Timothy 3\_1-9) - Rev. Mark Chan}
\newline
\newline
連結: \href{https://youtube.com/watch?v=KO8k0_vDZOY}{\texttt{https://youtube.com/watch?v=KO8k0\_vDZOY}} ~~~~ 語音日期: 2024-07-09
\newline
\newline
\hyperref[sec:VPCYBjM4gao]{\small{< < < PREV SERMON < < <}}
~
\hyperref[sec:index_chronic]{\small{[返順時目]}}
~
\hyperref[sec:index_scriptual]{\small{[返順卷目]}}
~
\hyperref[sec:x0e0tuTJkmQ]{\small{> > > NEXT SERMON > > >}}
\newline
\newline
$^{1}$(soft music).
- All right, let's turn to our Bible.
because this is one of the essence.
that we gather as the Church of God..
And we are going through a new series on the Bible,.
back to basics..
And we begin the series by this picture, right?.
This picture I show you last week is a tall building,.
it's a high rise building..
And it recalls us how important the foundation.
of a building is..
The taller the building,.
the deeper the foundation needs to be..
And as you can see this diagram,.
the foundation goes all the way down, right?.
From the ground down through the soft soil,.
through the mud and all the way down to the bedrock,.
which does not move, right?.
So in order to have a strong building,.
in order to have a secure building,.
you have to build your foundation first, right?.
And the foundation is not seen, right?.
It's not visible most of the time,.
but it's important, right?.
Without it, the building will collapse, right?.
Without it, you cannot build a building taller and taller,.
right?.
And the same way for our spiritual life,.
the same way for our life in God,.
if we believe in the name of Jesus Christ,.
we accept Jesus Christ as our Lord and Savior,.
we need to build a solid foundation for our life..
And our relationship with God,.
hinge on our foundation..
And we have to do it one day at a time.
and do it so one block at a time.
so that we have a solid foundation, right?.
So last week we talk about,.
we studied one of the Psalms in the Bible,.
Psalm 119, which is the longest chapter of the Bible,.

$^{41}$longest Psalm in the Bible..
The Psalm talks about the importance and the influence,.
the power of the Word of God..
And that's why we need the Word of God..
So today as we begin,.
we will begin with the, what?.
Verse of the week, yes, right?.
Remember to take home last week.
is to spend at least 10 minutes of your day,.
the best time of your day,.
spend 10 minutes at least to read the Bible, right?.
Read the Word of God, right?.
Did you do that?.
Okay, I see one person nodding the head..
Nodding the head, right?.
I know you've done it, yeah, I know you've done it..
I believe in you, okay?.
Take at least 10 minutes to read the Word of God, right?.
Read the Bible and memorize the verse of the week..
And this is the verse of last week, right?.
So let's do that together, all right?.
Psalm 119, 105, let's read it together..
Your Word is the lamb to my feet and the light to my path..
Can you recite it?.
Yeah, 105, 119, 105, okay?.
(speaking in foreign language).
Oh, okay, right?.
Okay, I practice it, all right..
So this is the verse of last week.
and I hope that you have it recited and memorized by heart..
And we put emphasis on the Word of God..
As a church, I told you we are the people of the book, right?.
Of this wonderful book, of this sacred book.
that God has revealed to us..
And today we'll continue to dig into the Word of God.
and I want you to, if you have a Bible,.
it's a print Bible or a Bible app in your phone,.
turn with me to the book of 2 Timothy..
It's one of the books in the New Testament.
called 2 Timothy..

$^{81}$Toward the end, if you search the table of contents,.
toward the end, it's called 2 Timothy..
And there are two points from 2 Timothy chapter three..
And I wanna share with you,.
I wanna go through this wonderful chapter with you.
and the first half of the chapter.
from verses one to nine, chapter three of 2 Timothy.
talks about the dangers of neglecting the Word of God..
There are many dangers if we neglect the Word of God..
And the second part of this chapter from verses 10 to 17.
talks about another big point,.
which is the benefits of learning the Word of God..
And because time is limited today.
and I will go through the first half,.
which is verses one to nine, the dangers,.
there are dangers of neglecting the Word of God..
Ask questions, ask the church..
Have you ever thought about it?.
If we don't have the Bible,.
if we don't dig deeper and deeper into the Bible,.
if we don't understand the Bible,.
if we don't know the Bible enough, right?.
If our knowledge of the Word of God is insufficient,.
inadequate, there are dangers there..
And we turn to this book called Second Letter to Timothy.
or 2 Timothy..
And we call this kind of books,.
with 1 Timothy, 2 Timothy and the book of Titus,.
there's three letters,.
there's three books in the New Testament..
We call them pastoral letters, okay?.
Pastoral letters..
And what it means is that these are letters.
sent by the apostle Paul..
He is the last disciple and the least disciple,.
apostle of all, called by God to preach the gospel.
to the Gentiles..
He's the last and the least apostle called by God..
And Paul is called by God to be apostle,.
preaching the gospel to the Gentiles, the non-Jews..

$^{121}$And Paul goes around in the Roman Empire of his day.
and plant many small churches..
And these small churches are being planted.
and being launched in different cities in the Roman world..
And these small churches are growing, right?.
Sometimes in the midst of, or many times.
in the midst of struggles and challenges..
And from time to time, early Christians.
are actually suffering and they are being persecuted.
by the authorities..
And there are crisis from without,.
and there's also crisis for these churches from within..
Because there are some false teachers.
who come into the church community.
and spread a false gospel, false teachings..
And Timothy is one of Paul's disciples,.
one of his mentees..
And Paul is kind of like a spiritual father to Timothy,.
and Timothy is kind of like a spiritual son to Paul..
And under Paul's direction, Timothy is in charge.
of the home churches in Ephesus..
Heard about this city name called Ephesus?.
It's a big city, right?.
Like Chicago..
It's a metropolitan city in the Roman world..
So Ephesus is a huge city, and you can imagine.
how difficult it is to pass through.
and to grow small churches in a huge city like Ephesus..
In a big city, there are many, many different forces.
that counteract the church community..
Cultural forces, economic pressure,.
religious forces, and spiritual forces..
They all counteract the church and the gospel..
And Paul writes letters to Timothy..
In total, he wrote two letters..
And to another disciple in another place called Titus..
So Paul wrote two letters, at least two letters.
to Timothy, his disciple, his pastor student,.
to help him, help him to be encouraged.
as he passes the church..

$^{161}$Because there are false teachers in the community,.
and these false teachers twist the truth and deceive many,.
especially those who are naive and uneducated or ill-equipped.
and these people are being led astray by false teachers,.
by cult leaders, or anyone who infiltrate into the church..
And you can imagine how difficult that task is for Timothy..
And Timothy has to teach, has to correct,.
has to grow the church at the same time,.
and deal with problems in the church.
and such an enormous mission..
And being a spiritual father,.
Paul writes letters to Timothy..
These are called pastoral letters.
because the tone in those letters are very personal..
So as we can read some of the passages from these letters,.
you kind of find that it's kind of like.
a further theological training for Timothy..
Already a pastor at Ephesus,.
trying to give him the equipment,.
trying to give him the knowledge,.
the courage to go on pastoring the churches..
So in 1 and 2 Timothy, if you read them,.
scan through the content,.
Paul talks about the gospel of grace,.
sound teaching, the importance of sound teaching.
and theology and doctrines, which are so important..
And because Timothy also has to raise up leaders.
among the communities,.
Paul also, in 1 and 2 Timothy,.
reminds Timothy the qualities.
and the qualifications for leaders,.
for elders, for deacons..
And Paul, you know, but all in all,.
trying to encourage Timothy to be patient,.
to endure hardship..
You know, it is hard to lead the church..
You know what I'm saying?.
It's not easy, right?.
It's not easy..
It's one of the most challenging leadership assignments.

$^{201}$on earth, is leading the church..
And now we come to 2 Timothy,.
which is the second letter Paul writes.
to his mentee, Timothy..
We turn to chapter three, okay?.
If you have your Bible, turn to chapter three as well..
Also, the scripture is on the screen..
Chapter three, verse one..
You should know this, Timothy..
Paul is addressing Timothy in a personal tone..
Know this, Timothy, that in the last days,.
there will be very difficult times..
For people will love only themselves and their money..
They will be boastful, proud, scoffing at God,.
disobedient to their parents, ungrateful..
They will consider nothing sacred..
They will be unloving, unforgiving,.
and slander others with no self-control..
They will be cruel and hate what is good..
They will betray their friends and be reckless,.
be puffed up with pride, love pleasure rather than God..
They will act religious, but they will reject.
the power that make them godly..
And Paul said to Timothy, "Stay away from these people..
"Stay away from them.".
You know, the last days, what is the last days?.
In New Testament terms, in biblical terms,.
last days is the period between the first coming.
of Jesus Christ and the second coming of Jesus Christ..
Christ came, Jesus came the first time,.
2,000 years ago, right?.
He was born God incarnate, and he lived the kingdom life.
in the midst of his disciples, in the midst of the people..
2,000 years ago, he died on the cross in Jerusalem..
He died on the cross and he was crucified, right?.
He was crucified, and he died for our sin,.
died for the sin of mankind..
And on the third day, he rose again,.
and he was ascended to heaven, and he is king over all..
And that is the first coming, right?.

$^{241}$It has been detailed in many books.
and many passages in the New Testament..
That's the first coming..
And there's also the second coming, which is in the future,.
which lies in the future..
Jesus will come back again.
because that's what Jesus said, right?.
That's what we are told in the Bible, in the word of God..
And there's a second coming..
Jesus can come any time, any moment, right?.
So the last days in New Testament terms.
is the period, I don't know how long, right?.
Between the first coming and the second coming..
And that is the last day.
because this is the last period of human history.
in biblical terms..
And that is the last day..
And Paul is reminding Timothy in this period, right?.
And we are in this period, how people behave..
And this last day, the time is critical.
because the last day is the period.
full of dangers, full of temptations,.
full of trials and suffering, right?.
And being Christian is not gonna make it easy..
It is difficult, right?.
And we continue with this passage,.
verses six and following..
He's talking about the ungodly people in the last days..
And Paul said, "They are the kind.
"who work their way into people's homes.".
Why?.
Why do they infiltrate into people's homes?.
Well, as I told you, the early churches,.
they meet in homes, right?.
They meet in people's houses..
They don't have the church building..
You know, church is a community, it's a family..
And they get together in different families, right?.
Whoever has a bigger living room,.
they gather at the living rooms..

$^{281}$So, "They are the kind who work their way.
"into people's homes and win the confidence.
"of vulnerable women who are burdened.
"with the guilt of sin and controlled by various desires..
"Such women are forever following new teachings,.
"but they are never able to understand the truth..
"These teachers, these false teachers,.
"oppose the truth just as Jannate and Jambres oppose Moses..
"They have depraved minds and can't defeat faith,.
"but they won't get away with this for long.".
And Paul says, "Someday, everyone will recognize.
"what fools they are, just as with Jannate and Jambres.".
It's a long passage, but listen up..
How Paul described the people in the last day..
The people love falsehood, and they hate truth..
They hate people telling them the truth..
The objective truth..
And listen how the descriptions are being demonstrated..
And they love themselves..
These people love the money..
They resist the will of God..
Or simply put, they're just proud, right?.
They're prideful..
And look at the behaviors..
They're trying to justify the behaviors.
in the name of God, even in the name of God, right?.
Act religious, talk religious, and talk religious,.
but they reject the authority of God, right?.
And Paul uses one of the examples.
known by many Jews of this day..
Two men by the names of Jannate and Jambres,.
and these two men oppose Moses..
Now, interestingly enough, you cannot find these two names.
in the Old Testament..
You cannot find it in the book of Exodus..
You can only find it in some hundred years before Jesus..
So 200, 300 BC..
At that time, the Jews,.
because they were under different empires,.
and the international language has changed,.

$^{321}$changed from the previous language to a new language.
called Aramaic, right?.
And you know the Bible, the Old Testament,.
was written in which language?.
Hebrew, right?.
Hebrew, by and large, Hebrew..
There was some Aramaic in some chapters..
But by and large, it's in Hebrew..
But in those days, some hundred years before Jesus,.
those Jews, when they gather together in synagogues,.
they wanna hear the word of God being preached..
So the leader of the synagogue will read part of the passage.
from the Hebrew Bible, of course in Hebrew, right?.
And those Jews, because they are not very fluent.
in biblical Hebrew, and by that time,.
the international language was Aramaic,.
so many of them actually speak and hear Aramaic, right?.
So what do you do, right?.
The original language of the Old Testament,.
the Hebrew Bible, is Hebrew, right?.
And there's always a translator there, right?.
Translator in the synagogue..
And he's gonna translate from Hebrew to Aramaic.
so that the people can understand the Bible,.
understand the word of God..
So the paraphrase or the translation is called tagam, right?.
You see the word at the bottom of the screen?.
T-A-R-G-U-M..
It's a Hebrew word..
Simply means translation, right?.
Translation, translated from Hebrew to Aramaic..
It's kind of Aramaic paraphrase of the Old Testament..
When they come to the book of Exodus,.
when Moses was confronting the magicians of the Pharaoh,.
right, with the miracles of God,.
with the signs and wonders of God,.
and the Aramaic paraphrase, the tagam,.
add these two names,.
Jannaz and Jambres, to the story..
So because the Jewish readers would know the story by heart,.

$^{361}$Paul is using this example to show that.
what people would do, the ungodly people,.
the false teacher would do in the last days..
They're just like the ancient history..
What had happened is that these two men,.
as we know Jewish audience,.
these two men opposed the authority of God through Moses..
So what Paul is trying to get at.
is that there were false teachings being widespread.
in the early churches,.
and they were vulnerable to false teachings,.
especially the women..
Now if you read this chapter or read this passage casually,.
you'll find it kind of uncomfortable, right?.
'Cause it seems like Paul doesn't have.
much good things to say about women, right?.
Women are being described as uneducated, naive, right?.
At first glance, Paul doesn't have much positive view.
on women, does he, right?.
But that's not true..
You cannot read the passage out of the context, right?.
And you have to understand what happened in those days..
2,000 years ago in the first century,.
women were especially vulnerable, right?.
Because in those days,.
women didn't have the opportunity to get educated, right?.
And nowadays, it's totally different, right?.
Women are sometimes much more educated than men, right?.
So what Paul is getting at in his day.
is that there were women in home churches,.
they didn't have the privilege of education..
So they were less educated, less equipped, right?.
Because they're less equipped,.
so they're easier to be deceived by false teachers, right?.
And that's what the false teachers.
or some of the false teachers.
trying to do infiltrate into the church.
and try to deceive people, right?.
So that they are more vulnerable to deception,.
and they follow those new teachings,.

$^{401}$so-called new teaching, and are deceived..
And there's danger of neglecting the word of God, right?.
So we understand the context in Paul's day,.
and apply to our days,.
it's actually pretty similar..
It's not talking about women being more vulnerable..
Actually, Christians can be vulnerable.
if we are ill-equipped,.
if we are not equipped in the word of God,.
if we don't know much about the word of God, right?.
And there's a danger of neglecting the word of God..
In Paul's days, in Timothy's day,.
as in today, right?.
If we do not dig deep into the word of God,.
if we don't know much about the Bible,.
if we don't know how to read the Bible,.
we are in danger of false teachers,.
or make ourselves more vulnerable.
to false, unbiblical teaching..
And this is alarming, right?.
This is alarming..
And this is why, you know, in our church,.
I put much emphasis on reading.
and understanding the Bible, right?.
Yes, we come to know God.
and restore our relationship with God,.
our creator, our father..
That's important, that's the starting point, right?.
However, in the journey of faith,.
we have to understand the Bible.
because the words in the Bible.
are God's authoritative words, right?.
Everything has to come under this standard..
The highest standard for Christian living and doctrines.
is the Bible, is the word of God, is the scripture..
And everything that we do,.
every tradition that we get to know,.
or we grow up with,.
has to come under the final standard,.
the highest standard, which is the word of God,.

$^{441}$which is the Bible..
And that's why I put emphasis on reading.
and understanding the Bible, okay?.
This understanding the Bible.
is not a monopoly for pastors or scholars..
It's for every single Christian..
Everyone who follow Jesus Christ.
has to understand and know the Bible..
One of the top-notch Christian scholars or historians.
today is Mark Noll..
Mark Noll used to teach in different institutes,.
including Whitten College..
And you know where Whitten College is?.
Whitten College is about 40 minutes west of Chicago..
It's very close to us..
And he used to teach there for many, many years..
His specialty is contemporary church history,.
American church history..
And Mark Noll has written a book,.
and it was published in the mid '90s,.
and it's still in print..
You know, renew edition, reprint edition..
And it's this book..
And the title of the book is.
"The Scandal of the Evangelical Mind.".
And if you pick up this book,.
and he will tell you what goes wrong.
in the evangelical faith in the church today..
And the first sentence of the book is this sentence..
The scandal of the evangelical mind.
is that there's not much of an evangelical mind..
Isn't that interesting, right?.
He's talking about one of the biggest problem that we have.
is that Christians are often anti-intellectual..
There's a divorce somehow between our intellect,.
how we think, and our piety, our spirituality..
And if you go to lay Christians,.
sometimes you have dialogues.
or a conversation about Christian faith..
Mind, you know, how we think, the clarity of thinking,.

$^{481}$clarity of what we believe and how we practice.
is not in the picture..
Spirituality is often defined as something lack of intellect.
or doesn't have anything to do with intellectual, right?.
Because we think, okay, those are for scholars, right?.
Those are lifeless stuff, right?.
And we don't need to deal with that, right?.
And Mark Noe is trying to diagnose.
what is going on in evangelical spirituality..
Is that there's not much of a mind, right?.
We often approach spiritual truth.
or approach the Bible or approach our faith.
with kind of a anti-intellectual attitude, right?.
We don't care about how we analyze..
We don't care about how we articulate our faith,.
whether we can explain our faith with clarity..
And sometimes spiritual is being defined as fanatic, right?.
Frenzy, right?.
As long as we feel high, as long as we feel good,.
and that's okay, right?.
And check the brain at the door, right?.
It seems like biblical faith is not grounded.
in any way with reason, right?.
But, you know, brothers and sisters,.
you need to understand in the Bible,.
faith, biblical faith, is grounded in facts.
and in reasons, right?.
It is..
It is grounded in facts and reasons in history, right?.
Biblical faith, true faith, is factually grounded..
It's reasonably grounded, right?.
What you believe, what we believe,.
is not a lip of faith, right?.
You just close your eyes and jump over there, right?.
Oh, hopefully this is the right thing to believe in..
No, the scripture tells us you can examine the facts.
and examine what the Bible says with a clear mind, right?.
And sometimes we just don't care..
Brothers and sisters, we cannot do this..
And Paul is asking Timothy to watch out.

$^{521}$because the false teachers is trying to infiltrate.
into the church and spread false teachings.
because, you know, they just don't reason with you..
I think this is really, really sad.
because many people know that the Bible.
is the best-selling book of all time, right?.
You know that?.
The Bible is the best-selling book of all time.
with five to seven billion copies printed.
over the last 1,500 years..
Five to seven billion copies, right?.
And the Bible has been translated into different languages,.
modern languages, Chinese, Korean, right?.
English, Spanish, and all sorts of different languages..
It's been translated into 700-plus languages..
And there are many, many translations to come, right?.
And I think this is really, really sad..
This is really tragic that many people know this,.
but sometimes we don't put effort.
to understanding the Bible..
It's more tragic that we have the Bible in our hands,.
in our phone, in our computer,.
but we don't even try to understand,.
try to handle the Word of God correctly, right?.
We're just being floppy in our study..
We're just being floppy in our thinking..
You know, therefore, I still believe,.
I still believe as a pastor,.
I've been pastoring for 25-plus years..
As a pastor, I still believe the importance.
of classical theological training,.
formal theological training..
You gotta go through school,.
some sort of degree program,.
and you need to be equipped,.
and well-equipped to learn scripture.
so that you can teach, right?.
I still believe all the pastors,.
all those who are being called to ministry,.
especially teaching ministry,.

$^{561}$they have to go through some sort of formal training..
Some sort of formal credentials are necessary..
It's not anything goes..
You have to go through some sort of formal credentials.
in order for you to be equipped to teach other people, right?.
You know, when I first studied in seminary,.
I started the Master of Divinity program in seminary,.
and you know, the toughest subject,.
the most hated classes among our students,.
you know what that is?.
It's original language, right?.
Hebrew and Greek, right?.
We just hated it, right?.
We have to memorize all these vocabularies.
and the paradigms and the grammatical rules, right?.
It's a new language, right?.
I don't even have time for my own native language, right?.
To polish my own language,.
but when we get into seminary,.
we have to study the word of God in its original language,.
right?.
Can you imagine, you know,.
someone studying like the greatest literature,.
the greatest English literature, Shakespeare, in Chinese,.
in translated version?.
No, you can't, right?.
You just can't translate that, right?.
So you have to understand,.
you have to know how the original language works, right?.
In order to have a better understanding..
And so for seminary students, right?.
Some sort of credential,.
some sort of formal training is needed..
Even though you hate it, you just have to do it, right?.
At least you have some working knowledge of how it goes,.
the basics for the Hebrew and Greek or even Aramaic..
And you know what the word seminary means?.
The word seminaries, we have Trinity Seminary,.
Evangelical Seminary..
We have a Moody, right?.

$^{601}$Also seminary, you know, some other seminaries..
You know what the word seminary means?.
(speaks in foreign language).
Right?.
Yeah, in Chinese, (speaks in foreign language).
Seminary is an old English word..
It was first used in the 16th century..
It comes from a Latin word, seminarium..
Seminarium means plant a nursery,.
a seed plot..
Now you get it, what seminary is all about, right?.
It's a seed plot for those who are getting themselves.
ready to serve God, to teach the Word of God, right?.
They come to seminary, it's like a small farm.
to grow the seeds, right?.
Until they sprout and they're strong enough.
and they are ready to be transplanted.
and to help churches to walk in the Word of God, right?.
And I consider myself privileged to receive.
some theological training in North America, right?.
You know, in North America, all the schools,.
the Bible schools, seminaries, are actually.
some of the best schools in the world, right?.
In the best school in the world..
I've been to three seminaries..
Westminster Theological Seminary, they have two campuses,.
Philadelphia and San Diego, Southern California..
The second one I went to is Regent College.
in Vancouver, Canada..
That's where I got the MDiv, Master of Divinity..
And then I started to serve and took a sabbatical.
after some years and I went to Wyndham College.
for my biblical archeology study..
And then now I'm in a program at Midwestern Seminary.
in Kansas City..
And all of these seminaries, I have learned so much, right?.
I walk into the classroom, I'm humbled..
Humbled by what I don't know, right?.
Humbled by the vastness of the Bible, right?.
The vastness of our God, how great our God is..

$^{641}$And we will continue to learn it, right?.
And my first seminary, Westminster Theological Seminary,.
one of their slogans is, "Whole counsel of God.".
You gotta know the whole counsel of God..
The whole will of God being revealed in scripture..
You cannot just cherry picking all the time..
You have to be confronted by the whole Bible.
because the whole Bible reveals the whole counsel of God..
The whole will of God, the demand of God,.
the covenant of God, the word of God..
You cannot cherry picking..
It's not Whole Foods..
You can come in and according to your desire,.
pick whatever you like..
The Bible is not Whole Foods..
You come to the Bible and you need to stand.
under this highest authority, this highest standard..
You gotta be confronted by the Bible,.
your living, what you believe,.
and how you practice the faith,.
and what your doctrines are, right?.
All the time, you need to come to the Bible..
The Bible is like a mirror..
It comes, you go into the mirror.
and check whether you look okay, right?.
And every time you open the Bible,.
every time you read the Bible,.
you are doing self-checking, right?.
Checking yourself, checking your Christian life,.
whether it's aligned with the will of God.
or it's misaligned in some areas..
And with the spirit of repentance,.
you come before God and ask God to help you.
to align your Christian living.
according to the will of God..
And that's why I have to take home for you today, all right?.
So last week, I shared the first two points with you..
Take at least 10 minutes a day to read the Bible, okay?.
Guys, okay?.
If you don't know, if you haven't started,.

$^{681}$I ask you to go to the book of Ephesians, all right?.
I know our men's group is going through the book of Ephesians.
and just read, there are six chapters, right?.
Six chapters in the letter of Ephesians.
and you can start with chapter one,.
chapter two the next day,.
chapter three, chapter four, five, six,.
and then there are many things.
that you may not be able to understand at this point,.
but you can go back, right?.
You can recycle it, right?.
Going back to the first chapter and try to meditate,.
try to understand what is going on here..
What is the argument?.
What is the thesis of the chapter?.
What is the point, right?.
What is the point of what Paul is saying?.
You know, those are texts..
Those are languages and you have to understand it..
You have to comprehend the meaning..
Or you can start with one of the gospel, as I told you,.
Matthew, Mark, Luke, or John..
You can start with one chapter a day..
Take 10 minutes a day, that's the whole point,.
to read the Bible and to think about it..
And then memorize our verse of the week, okay?.
I have a next verse for you later, okay?.
So, you know, be ready..
Number three, I wanna add this,.
to read some good books for spiritual growth, okay?.
Good books..
I'm telling you, I wanna emphasize good books..
Because there are bad books out there, all right?.
You know, I'm over 50 years old.
and I remember when I was young in faith,.
when I was still a college student,.
and my home church didn't have a pastor,.
didn't have good teacher..
So we have to go out to bookstores, Christian bookstores..
Yeah, can you find a Christian bookstore these days?.

$^{721}$No, right?.
They're all closed down, right?.
This is pretty sad, right?.
Everyone goes to Amazon, right?.
Everyone goes to eBooks, right?.
We went to the bookstores and find some books.
to help us and guide us in our Christian walk, right?.
And you go into the books,.
oh, there are different sections, different categories, right?.
And you go in and you try to find some good books.
and we didn't know what good books are, right?.
We just cannot judge, we cannot discern..
So I bought quite a few books.
because I was yearning for the Word of God.
and I read all I can, right?.
And then years later I find that,.
oh, these books are actually bad.
and I didn't know better, right?.
Those are prosperity gospel..
Those are, you know,.
truncated understanding of scripture, right?.
They're fun to read, they're enjoyable to read..
However, they are not sound teaching,.
they are not solid teaching..
And years later I find out, right?.
I feel so bad about myself.
and I have to kind of re-run away my own spiritual life.
by reading good books, getting to know good books..
And brothers and sisters, you need to read,.
or you need to start the habit of reading..
You need to start the habit of reading..
Yeah, I know in our days and age,.
we don't like to read, right?.
We like YouTube or TikTok, right?.
All those sound bites, you know?.
We just like, you know, be entertained, right?.
And passively receive information and misinformation..
And we don't train ourself.
to be better people of discernment..
You know, you need to train yourself how to think,.

$^{761}$how to analyze, how to articulate, right?.
And I wanna recommend you some of the good books.
for spiritual growth, right?.
Some of the good books, right?.
Again, there are many books out there, right?.
The worst thing today is that you go to Amazon,.
you search spiritual growth, you know,.
by and large you will come across some of the very, very,.
you know, truncated teaching, right?.
Because those books are often written by, you know,.
people who don't learn enough in the world..
And they don't have the credential to do so,.
to write certain things..
So, you know, we have to learn.
and develop the habit of reading..
Because reading is active learning, right?.
You read a passage, it's always actively engaging your mind..
You try to understand those words.
and try to connect them together and find the thesis, right?.
Find what it says, right?.
You have to summarize it,.
you have to go back to the text and to digest it,.
to comprehend it, right?.
You have to have the habit of reading,.
ask questions in order for you to grow spiritually..
There's no way you can get without reading..
Some serious reading, some solid reading, right?.
I want to recommend you some of the books,.
some of the classics actually, very good books..
Some of them are older, some of them are more recent..
On the left you see the book written by J.I. Packer..
He passed away not too long ago..
He has this classic,.
"Over One Million Copies Sold," right?.
I'll be more thankful,.
it says, "Over One Million Copies Comprehended.".
Read and comprehend..
However, it's a very good book, "Knowing God.".
It's a very simple, systematic, summary,.
explanation of the whole Christian faith..

$^{801}$And it's been translated into different languages..
Chinese is one of them..
And Tim Keller has written many, many good books..
Yeah, Keller also passed away not too long ago.
and be with the Lord..
But one of the book is, "The Reason for God," right?.
"Believe in an Age of Skepticism.".
This is a book especially good.
for those who are seeking the faith..
You want to understand the faith,.
you haven't made the commitment to believing in God,.
in Jesus, but this is a book for you, Tim Keller..
"The Reason for God," right?.
He walks through you step by step,.
logically, reasonably, what our faith is all about, right?.
Why do we believe in God, right?.
And on the right is,.
Anthony writes, "Simply Question,".
why Christian identity makes sense, right?.
And these are very, very solid resources to help us grow..
Yeah, not too long ago, during Good Friday weekend,.
we invited Dr. Chloe Sun, right?.
Old Testament professor.
from Fuller Theological Seminary in Southern California..
She came and gave a workshop,.
a series of workshops on interpreting the Psalms.
and the Bible..
And many of Chloe's books are actually helpful.
for us to understand the Word..
And so, all I'm trying to encourage you.
is that take some moment to read the Bible.
and also find some good books for spiritual growth..
We can go to our bookshelf.
and find some good books as well, right?.
And I wanna end today with this verse of the week, all right?.
Matthew chapter 22, verses 37 to 38..
And this is what Jesus said..
Jesus said, "You must love the Lord your God.
"with all your heart, all your soul, and all your mind.".
And this is the first and greatest commandment..

$^{841}$Why did I pick this one?.
You must love the Lord your God with all your heart,.
all your soul, and all your mind..
Your thinking is a way of worship..
Your thinking, your knowing God,.
is a way to draw closer to God..
It's not just by feeling, right?.
It's not just spiritual high, but with a cool head..
You think in clarity..
You think about the faith..
You have a correct understanding,.
solid understanding of the Word..
And that is the first and greatest commandment..
Let's recite it together, all right?.
Matthew chapter 22, verses 37 to 38..
"You must love the Lord your God.
"with all your heart and all your soul, and all your mind.".
And this is the first and greatest commandment, all right?.
Let's recite that this week..
Let's all pray together..
Father, we come before you..
We give you thanks because you are the God who speaks,.
and you are the God of truth..
And thank you for revealing your will.
and your whole counsel through scripture..
Thank you for this book that we have..
Thank you for the content that we have..
It talks about you, God..
It talks about how you deal with a sinful mankind,.
how you love us, how you deal with this broken world..
It talks about your heart..
Help us to get to know you more, deeper and deeper,.
drawing closer to you..
And I pray for the sisters and brothers here.
that we do make the effort of knowing you..
And by knowing you, we experience joy and peace in you..
Reading the word, reading your scripture,.
we come closer and closer to the heart of God..
And I pray that you will bless us each day.
with some moments to meditate, to think about your word..

$^{881}$And as we go reading, as we get to know our faith.
more and more, draw us closer so that we know you better..
And thank you for being our God..
That we can relate to you..
We can call you friend..
We can call you heavenly father..
And thank you for the grace and for the peace that we have..
So may your spirit send us out with a spirit of truth.
to yearn for your truth, to get to know you..
And may your church rise up and grow up.
with the foundation solid, grounded in your word..
In Christ, I pray..
(gentle piano music).
Well, God bless you and I wanna see you next week..
And hopefully you can stay a little bit.
and for table fellowship to enjoy the refreshment.
prepared by Kelly..
Thank you, Kelly..
And get to know one another and say hi to one another..
(gentle piano music).
(piano music).
\newpage



\section{提摩太後書 3:1-9}
\label{sec:x0e0tuTJkmQ}
\textbf{回到基礎 \#1 (提摩太後書 3\_1-9) - 陳崇基牧師}
\newline
\newline
連結: \href{https://youtube.com/watch?v=x0e0tuTJkmQ}{\texttt{https://youtube.com/watch?v=x0e0tuTJkmQ}} ~~~~ 語音日期: 2024-07-09
\newline
\newline
\hyperref[sec:KO8k0_vDZOY]{\small{< < < PREV SERMON < < <}}
~
\hyperref[sec:index_chronic]{\small{[返順時目]}}
~
\hyperref[sec:index_scriptual]{\small{[返順卷目]}}
~
\hyperref[sec:eiPsrBdDRJk]{\small{> > > NEXT SERMON > > >}}
\newline
\newline
提摩太後書 3:1-9
\newline
\begin{longtable}{cl}
\hline
\hline
章節 & 經文 (和合本修訂版)\\
\hline
3:1 & \begin{tabularx}{0.7\textwidth}{X} 你該知道，末世必有艱難的日子來到。 \end{tabularx} \\ \\ \relax
3:2 & \begin{tabularx}{0.7\textwidth}{X} 那時人會專愛自己，貪愛錢財，自誇，狂傲，毀謗，違背父母，忘恩負義，心不聖潔， \end{tabularx} \\ \\ \relax
3:3 & \begin{tabularx}{0.7\textwidth}{X} 沒有親情，抗拒和解，好說讒言，不能節制，性情兇暴，不愛良善， \end{tabularx} \\ \\ \relax
3:4 & \begin{tabularx}{0.7\textwidth}{X} 賣主賣友，任意妄為，自高自大，愛好宴樂，不愛神， \end{tabularx} \\ \\ \relax
3:5 & \begin{tabularx}{0.7\textwidth}{X} 有敬虔的外貌，卻背棄了敬虔的實質，這等人你要避開。 \end{tabularx} \\ \\ \relax
3:6 & \begin{tabularx}{0.7\textwidth}{X} 他們當中有人潛入別人家裡，操縱無知的婦女；這些婦女被罪惡壓制，被各樣的私慾引誘， \end{tabularx} \\ \\ \relax
3:7 & \begin{tabularx}{0.7\textwidth}{X} 雖然常常學習，終久無法達到明白真理的地步。 \end{tabularx} \\ \\ \relax
3:8 & \begin{tabularx}{0.7\textwidth}{X} 從前雅尼和佯庇怎樣反對摩西，這等人也怎樣抵擋真理；他們的心地敗壞，信仰經不起考驗。 \end{tabularx} \\ \\ \relax
3:9 & \begin{tabularx}{0.7\textwidth}{X} 然而，他們沒有進步，因為他們的愚昧必在眾人面前顯露出來，像那兩人一樣。 \end{tabularx} \\ \\
[1ex]
\hline
\hline
\end{longtable}
$^{1}$兩位主持人早安.
在我們中間的都是.
Joyful可能我們吃夠了,玩夠了.
然後週末我們就要留在這裡,對吧.
那些不在的就去了別的地方,對吧.
如果有的話.
長週末.
其實芝加哥去哪裡了.
最多去多倫多吧.
最厲害的還是多倫多.
我們先跟心得一起祈禱,好嗎.
主要為教會獻上感恩.
我們為每一位弟兄姊妹主.
我們求你們都祝福,來帶領.
讓我們每一位能夠在基督裡面.
享受到世人給不到的平安.
主,我們或許會有不同的掙扎.
有不同的問題,甚至有病痛.
但我們知道在基督裡面.
我們有最終極的安康.
最終極的醫治.
給我們這個盼望能夠成為我們今天.
生活裡面的動力.
讓我們能夠思想信仰.
學習應用信仰.
以致我們心靈能夠更加成為你見證.
祈禱,奉主名禱.
這次我們開始一個新的系列.
這個系列叫做Back to Basics.
甚麼是Back to Basics.
就是回到基本功.
甚麼是基本功呢.
大家想想.
我們日常生活裡面做很多事情.
都有基本功的,是嗎.
譬如你是做運動的.
都是有些基本功需要熱身的,是嗎.
最近芝加哥城市裡面有城市,知道嗎.
賽車,是嗎.
我們很討厭,是嗎.

$^{41}$因為阻礙我們.
阻礙我們做生意,行動,交通.
昨天我在電視上看了一點.
我太太問我.
為何車輛要這樣.
一開始的時候.
要打S型的.
大家知道嗎.
知道嗎.
那些零頭的,你可否告訴我.
為甚麼.
你現在告訴我.
為何要這樣.
一開始就要這樣.
開始之前.
還沒開始跑之前.
因為那些輪胎是要熱身的,知道嗎.
那些車輪輪胎需要暖暖的.
它才可以怎樣.
才可以扔到彎位,是嗎.
對嗎.
大家知道吧.
如果不是,就會怎樣.
就會很容易滑胎,是嗎.
但太熱又不行.
太冷又不行.
所以要在中間.
所以他們一開始就在這裡轉來轉去.
這就是所謂的基本功.
每一個科目或每一個行動,每一個活動.
都有它的基本功.
信仰其實也是一樣.
信仰有我們的基礎.
我們需要去把握.
基礎要做得好.
以至你能夠去享受那個活動.
能夠享受那個行動.
這個回歸基本功是怎樣來的.
在不久之前,大概兩星期前.
我在社交群組上有人發了一張這樣的照片.

$^{81}$社交群組的社交媒體帖文,新聞帖文.
有沒有這些東西.
沒有.
為甚麼我會有呢.
因為我很喜歡在家裡做事.
做木工,做這個做那個.
裝修一下.
小的那些,不是大的那些.
所以有時候會出現這些東西.
我發現這裡很有意思.
雖然是一個很簡單的圖.
大家可能見過.
一個樓越高的時候.
它的基礎就需要越深,是嗎.
你看到那個基礎.
是要下去到下面所謂的背景的地方.
床岩的地方.
上面的是軟的土.
再下去的可能是一些泥沙.
然後再下去下面就是一些硬的.
不會動的.
一些背景的床岩.
那些東西是不會動的.
所以你的基礎一定要下去到最底最底的地方.
最穩定的地方.
然後你的建築物才能夠建得高.
你看到我們芝加哥特別市中心.
有很多新的很高的樓宇,是嗎.
當然因為他們的基礎是下去到很深.
以至他們能夠建設到這麼高的樓宇.
是嗎.
最近有沒有看到一些新聞.
我看了一個中國大陸的新聞.
說有層樓裂開.
是嗎.
無緣無故裂開.
在住宅裡面多麼有趣.
真是慘.
很多豆腐渣的東西.
其實美國也有很多這樣的橋等等.

$^{121}$這些建築物都是日久失修,是嗎.
因為說的是建築.
說的是物質上的建築.
但屬靈上也是需要建築的.
我們的spiritual life.
是需要有basics,有那個foundation.
所以在這個系列裡面.
我們去說一說有什麼foundation是我們需要的.
或者在這個summer裡面.
我們需要說一說這件事.
我們今天去看一個新約的書信.
叫做提摩太后書.
提摩太后書第三章裡面說到兩大點.
今天的時間的原因.
我們只能夠說到頭一大點.
兩大點是什麼.
第一大點就是忽視主的話語的危險.
如果我們忽視主的話語.
不去學道的話.
不去明白聖經裡面所說的話.
就有很大的危機.
這是三章一至九節.
然後第二大點.
就是the benefit.
有什麼益處.
學習主的道.
學習聖經裡面的話語上帝的啟示.
有什麼益處.
這是三章十至第十七節.
整個三章裡面十七節.
兩大部分.
第一部分就是忽視主的話語.
有什麼危機.
提摩太后書其實是屬於教穆書信.
我們統稱它為教穆書信.
Pastoral letters.
就是包括有三卷書在裡面.
在新約裡面這三卷書.
就是提摩太前書,後書和提多書.
1 Timothy和Titus.

$^{161}$這三卷書我們叫它教穆書信.
原因是什麼.
因為這三卷書信特別是保羅寫給他們的徒弟.
其中一個就是提摩太.
寫給他們怎樣去教導他們.
去處理教會裡面的事務.
當然如果這樣看的話.
所有書信都是這樣的.
所有書信都是有關教會的運作.
教會的成長,教會的建政等等.
哥林多前後書,羅馬書等等.
其實都是這樣.
但我們特別叫這三本書做教穆書信.
因為特別是一種個人的語氣.
來一個師父保羅.
去跟他的徒弟說一番話.
去叮嚀他.
告訴他教會裡面有什麼事發生.
你要怎樣去處理,怎樣去牧養.
所以我們今天看提摩太後書.
你知道當時第一世紀.
保羅去到不同的地方傳福音的時候.
有人信主.
然後他們建立了教會.
在不同的城市,不同的社區.
不同的國內有一些所謂小組的教會.
那時候教會很小的.
在家裡聚會的.
每個星期一次或者多過一次.
他們去到聚會的時候.
他們就會有教導的時間.
讀一些上帝的華語.
當時沒有新約聖經.
只有舊約聖經.
所以他們讀摩西五經等等.
有時候加起來.
就是保羅會寫信給一些教會.
好像寫了電郵給他們.
告訴他們.
記住了.

$^{201}$你們聚會的時候.
你們聽清楚我說什麼.
教目書信就是其中一些新約的書信.
他就是給一些信仰的群體.
來指導他們.
很多時候這些教會.
初生的教會.
在一些大都會裡面.
去到存在的時候.
就有很多的挑戰,很多的危機.
特別是一些內環.
內環中間可能有一些問題.
或者一些張力,關係搞得不好等等.
這個是內環.
也有外憂.
有外環.
因為大都會裡面有很多不同的文化衝擊.
或者思想的衝擊.
宗教的勢力等等.
就會滲入這些教會裡面.
特別是提摩太前後書.
就有一些異端出現.
有一些假教師.
進入了教會後.
去到迷惑不同的人.
於是保羅就寫這些書信.
叫人帶去.
然後去到那個教會.
就去宣讀.
去教導他們.
提摩太當時在哪裡呢.
提摩太在羅馬帝國的東邊.
一個大的城市叫以弗所.
Ephesus.
你可以想像在一個大的城市.
當時的世界裡面.
數一數二的大城市.
在裡面去到.
牧羊教會.
是一個很大的使命.

$^{241}$所以保羅對他的徒弟.
提摩太就好像父親對兒子一樣.
其實他們叫大家.
是屬靈的父子.
因為是他帶大的.
所以他寫了這些書信給提摩太.
就去教導他怎樣去牧羊教會.
提摩太前後書就是這樣類型的書信.
有一點像牧師在職的訓練.
再寫一封書信給他.
再提醒他一些基本的東西.
或者一些重要的東西.
講到福音的恩典.
提摩太前後書講到什麼叫Sound teaching.
什麼叫正統.
確立的教義.
這些教導對於教會.
對於基督徒的生命是重要的.
特別在教會裡面選立.
執事.
領袖.
他們的質素應該怎樣.
要有什麼Qualification.
所以保羅就在這兩封書信裡面.
用一個很叮嚀的口吻.
去告訴提摩太.
你需要忍耐.
需要在艱難當中.
在試煉當中去Tough it out.
要盡力.
繼續去事奉.
繼續去牧羊教會.
今天我們看提摩太後書第三章.
提摩太後書第三章.
我們看兩段.
第一段.
一至到第五節.
你該知道.
他跟提摩太說.
提摩太你該知道.

$^{281}$末世必有艱難的日子來到.
那時人會專愛自己.
貪愛錢財.
自誇.
狂傲.
毀謗.
違背父母.
忘恩負義.
心不勝結.
沒有親情.
抗拒和解.
哭說.
殘言.
不能節制性情.
兇暴.
不愛良善.
賣主賣友.
任意妄為.
自高自大.
愛好燕落.
不愛上帝.
有敬權的外貌.
卻背棄了敬權的實質.
這當日.
你要避開.
什麼叫做末世.
末世在新說聖經裡面的教導.
就是耶穌第一次來.
在十字架上死.
然後復活.
升天.
第一次的來.
直到耶穌再回來.
耶穌再來.
耶穌當然還沒再來.
是將來的事.
隨時會發生的事.
任何時候.
第一次和第二次的來.
中間這段時間.

$^{321}$就是新約裡面所說的最後一天.
不要說最後一天就快來了.
我們快來到末世的日子.
不.
已經是末世的日子了.
現在已經是.
Right.
你想想死風書信的時候.
是當時第一世紀.
是吧.
整段時間到現在.
都是末世的時候.
Right.
保羅說末世的時候.
會有艱難的時候.
這不是一個easy time來的.
他和提摩太說.
和這位初生的年輕的牧者.
告訴他.
You know what.
你不要將這些浪漫化.
我們牧羊教會.
或者在這個末世當中.
是有difficult times.
作為一個信徒.
都會有difficult times.
Right.
然後他就形容裡面的人.
是一種怎樣的人.
Right.
列舉一大堆的東西.
當時新約的書信.
或者當時希羅的書信都是這樣.
要去爭論一件事.
要去論證一件事的時候.
他就給你一個list.
給你一個長的list.
就是告訴你.
你服不服.
明不明.

$^{361}$你看看這裡的.
對人的描述.
末世裡面的人.
會是怎樣.
不愛真理.
聽見真理的時候覺得.
很不舒服.
不喜歡聽.
Right.
你看看保羅怎樣形容這些人.
愛自己的.
什麼是行先.
是自己的前菜行先.
Right.
然後呢.
狂傲的去到.
敵黨.
上帝的權威.
上帝的話語.
你可以大概將他總結出來.
就是.
這些人.
就是.
最後這句.
有敬虔的外貌.
背棄了.
敬虔的.
實質.
這班人是自己先行的.
是敵黨上帝的.
Right.
這個就是末世的時候.
如果當時保羅和提摩泰.
是見到這樣的情況.
我相信我們今天.
就更加見到.
Right.
對嗎.
無論你是在這些社交群組裡面.
收到的消息.

$^{401}$或者聽到一些新聞.
或者在社區裡面等等.
可能甚至是你.
可能就是屬於這種描述.
敵黨上帝自高自大.
愛好的是獻下自己的肚腹.
不愛上帝.
心不成結.
忘恩負義.
Right.
是很可怕的事情.
Right.
今天也是一樣.
或者甚至今天是更加.
再看下去第六節.
他說他們當中有人潛入別人的家裡面.
操縱無知的婦女.
這些婦女被罪惡壓制.
被各樣的私慾引誘.
雖然常常學習忠九末法.
達到明白真理的地步.
從前雅尼和楊比.
怎樣反對末世.
這班人也怎樣敵黨真理.
他們的心地敗壞.
信仰經不起考驗.
然而他們沒有進步.
因為他們的愚昧必在眾人面前顯露出來.
像那兩人.
就是雅尼和楊比一樣.
潛入家裡面.
正如我所說.
當時的教會是House Churches.
是家庭教會.
他們潛入家裡面.
這些假的教師想操縱什麼.
就是無知的婦女.
然後又趁著這些婦女.
她的私慾.
她自己的desire.

$^{441}$不能夠得到控制.
然後她就教壞他們.
這班人.
這班所謂的無知的婦女.
就是常常學習.
但是沒有明白真理的地步.
反而好像是反對末世一樣.
在這裡解釋一下是什麼意思.
這裡是不是說聖經是看低女人.
可不可以說.
翻開這段聖經.
我們就說.
這裡說婦女是無知的.
我就和我身邊的女士說.
和我太太說.
你是無知的.
《史聖經》也是這樣描述你.
被操控.
絕對不是這樣的.
請你不要看錯聖經.
雖然看上去表面好像是.
保羅沒有什麼好的東西說女人.
但是你明白所有的書信都是什麼.
是按照當時的處境寫出來的.
你就可以想像當時的教會.
或者家庭的教會發生什麼事.
甚至當時社會的背景.
當時受教育的人是什麼人.
有機會受教育.
或者有特權可以受教育的人.
大部分都是男人.
上等的男人.
女人很多時候是沒有教育的.
是不能被教育的.
事實上.
女人得到普及的教育.
都不是很久之前的事情.
在世界歷史裡面.
在八零二八年開始.
或者之前.

$^{481}$到近代.
當然今天的女士.
是完全另一回事.
對嗎.
今天的女士很多是很有教育.
很有智慧的.
我不是說學位.
而是她生命裡面的智慧加上.
生活的智慧等等.
是不是.
所以這裡不是說保羅說.
看低女人的意思.
而是當時的.
特別是無知.
或者特別是不能夠.
有那種教育的水平.
去分析的水平.
那班人很多時候往往就是婦女.
很多時候教會裡面.
就是那班人他們是剛開始學到的.
但是他們如果有一些人.
他們當中有人潛入別人的家.
就是一些假教師.
入了去.
異端入了這些家人的教會的時候.
就大件事了.
他就操縱了這些無知的人.
這些特別是很多是婦女的.
當然不是說今天.
但是說當時.
當時發生什麼事.
就是一些假教導入了去教會的時候.
這些未受裝備.
或者是.
知一半又不知道一半.
半知半解的.
那班人就會被迷惑.
被引導.
然後相信了那些東西.
對嗎.

$^{521}$一些所謂新的教導.
我記得有一次.
沒多久之前幾個星期.
和幾個英文堂的弟兄姊妹吃飯.
有幾個都是講國語的.
他們是中國大陸背景.
他們就說開.
這個.
在中國大陸裡面產生了一種異端.
叫做東方閃電.
是嗎.
那個弟兄就說開這件事.
他說東方閃電很厲害.
不只是中國大陸.
是在美國不同的pockets.
美國不同的pockets.
他就問我這裡有沒有.
這裡好像.
ok.
不是很活躍.
喂.
對嗎.
或者可能我錯了.
有沒有聽到這些消息.
在馬上說.
沒有.
是嗎.
我們知道哪家餐館開的.
你知道嗎.
有什麼.
有什麼新穎的異端.
我們誰知道嗎.
就是說在羅省.
在紐約.
是很厲害的.
東方閃電講什麼.
你自己查一下.
ok.
你自己查一下.
你不要亂找.

$^{561}$他們做的視頻.
是做得很專業的.
你看下去是看不到的.
你不知道的.
你不知道是Mark Sir做的.
還是他做的.
很賭本做這些事情.
如果你找全能神教會.
就是東方閃電那些.
怎樣也好.
他們就說開.
就是這些.
因為很多時候.
這些異端產生的地方.
就是一些農村的地方.
就是郊區鄉下的地方.
因為他們怎麼分辨.
對嗎.
所以為什麼要.
要放本給宣教士.
就是這個意思.
宣教士就是被派入去.
不同的地方.
以至他能夠去全港.
真的能夠去糾正一些.
錯誤的教導.
對嗎.
我們不能去.
不過我們可以支持這些人.
去到當中.
那些傳道人,牧師.
去到當中.
能夠去教導他們.
對嗎.
但怎樣也好.
這裡在說什麼.
甚至今天.
如果我們不去學道的話.
一知半解的時候.
就會很危險.

$^{601}$因為我們在說屬靈的真理.
是在說永恆的東西.
是在說永遠的東西.
對嗎.
有人去騙你.
當然是不好.
對嗎.
到處在教會.
或者在不同的群眾裡面.
去騙人.
或者去拿人便宜.
當然這是不好的.
如果你不小心被人騙了.
那也是現今的東西.
現世的錢財.
都是在說永恆的生命.
一些東西是永遠的.
這個很嚴重的.
保羅就說給提摩塔聽.
這種人.
你要去防備.
和避開.
他說從前.
這個瓦尼和楊比.
這樣反對摩西.
這等人也是敵黨真理.
心地敗壞.
經不起考驗.
沒有進步等等.
瓦尼和楊比是什麼人.
Janus和詹姆斯.
原來在舊約聖經裡面沒有提過這兩個人.
在猶太人的典籍裡面.
有一個典籍叫做Targum.
Targum.
T-a-r-g-u-m.
Targum的意思就是希伯文.
就是translation的意思.
因為到了大概主前一二百年左右的時候.
很多的猶太人已經不再聽到.

$^{641}$或者不再熟練希伯文.
他們當時的international language.
是Aramaic.
因為Aramaic的緣故是亞倫文的緣故.
所以很多時候他們在會堂裡面去聽聖經.
去聽摩西五經的講解的時候.
那個講解者就會讀希伯來聖經.
是不是.
舊約聖經.
所以就有一個叫做Targum.
一個翻譯出來.
有一個翻譯員在裡面.
他就將希伯文翻譯成Aramaic.
以至那些操亞蘭文的猶太人.
他們就明白聖經裡面的意思.
在Targum裡面.
在出埃及記的Targum.
去翻譯這一段摩西和法老的魔術詩.
法老的法術詩.
去到敵黨大家的時候.
摩西去到挑戰他們的時候.
那些Targum.
Aramaic的翻譯.
以逆本.
就加了這兩個名字進去.
可能比較容易明白.
這樣的意思.
保羅這樣寫.
就說因為提摩太郎也算是猶太人.
他很明白.
明白這些典籍裡面的故事.
有些人去反對摩西.
反對耶和華的權柄.
所以保羅就用這個例子.
一個很鮮明的例子.
說這些人是有堅強的外貌.
但是沒有堅強的實意.
這些人是反對.
摩西的權柄.
摩西的律法等等.

$^{681}$到末世的時候.
都是一樣.
在家庭的教會裡面.
這些假教師會進入去.
滲入去.
然後慢慢他們會.
所謂將這些假的教導.
去到大腿底.
危險的地方就在這裡.
這個問題不是有沒有假教師的問題.
而是教會裡面的人.
他們分辨不到.
這個才是更大的問題.
對嗎?.
這個才是更大的問題.
我們今天的教會都是一樣的.
這個世代的教會.
其實我們需要去到.
真的去學道.
你不要讓學道這件事.
或者去明白主道這件事.
解釋主道這件事.
不只是給牧師.
或者所謂神職人員的專利.
每一個信徒都需要去裝備.
去明白聖經.
去明白我們的信仰.
以致我們的基本.
我們很清楚的時候.
我們就不會那麼容易.
被別人去影響.
但是很可惜.
很多時候在教會當中.
就不是這樣的.
有一本很暢銷的書.
是在1990年中的時候寫的.
到現在還在印的這本書.
是一個很出色的基督徒學者.
叫做Mark.
不是Mark Chan.

$^{721}$Mark Ngo.
他曾經在不同的學院裡面教書.
特別是教近代的基督教歷史.
教會的歷史.
Mark Ngo.
他寫過一本書叫做.
《The Scandal of the Evangelical Mind》.
就是說.
科恩派信徒思考的醜聞.
Mark Ngo正如我所說.
他在不同的學院教過.
特別是在威廉大學.
在威廉大學教過很久.
現在去了另一個.
好像是Norondale.
去了另一個教育系.
他是世界裡面其中一個最出色的思想教練.
他在這本書裡面第一句.
劈頭第一句就是這一句.
The Scandal of the Evangelical Mind.
Is that there is not much of an evangelical mind.
說什麼呢.
科恩派的信徒裡面.
他們思考的醜聞是什麼.
就是根本沒有什麼思考.
這個很奇怪.
很糟糕.
為什麼呢.
因為很多時候在一些的傳統裡面.
就是將理性,理智和經典拆掉.
是分開了的.
是對立了的.
所以基督徒認為信主就是反智的東西.
甚至世界也會這樣看.
你們信耶穌其實是反智.
是不講道理的.
但你知道聖經裡面講的信仰不是這樣的.
聖經講的信仰是有事實根據.
有歷史根據.
是經歷過考證的.

$^{761}$是可以critically examine.
甚至是approve.
每樣東西是講實在,講事實.
是講真理的.
是講道理的.
不是反智的.
但為什麼很多基督徒就是這麼反智呢.
每樣東西都是我們心領神會就行.
我們不會處理一下思考這件事.
可能你從小到大.
信主到現在.
可能都未必有人跟你強調這件事.
我們所強調的就是.
總之你感覺到耶穌愛你.
這個重要嗎.
當然重要.
因為上帝跟我們之間的關係.
就是信仰的核心.
但有沒有思考這件事呢.
是有的.
我記得以前在神樂院的時候.
有一個老師叫John Frank.
他曾經用過一個例子.
你說上帝會不會理會或者care不care.
你知道某些事實.
數據的事實.
他舉了一個例子.
加州的首府是Sacramento.
你猜上帝care不care.
你知道這件事.
而且知道是正確的.
He cares.
因為是Truth.
是事實.
你不可以搞錯.
如果你搞錯的時候.
上帝是care的.
意思就是說上帝是care我們的理智的.
我們思考,我們想東西.
去分析,去分辨.

$^{801}$有邏輯.
有條理的說出來.
信仰不只是大發熱心.
以至我們最後變成狂熱分子就夠了.
根本就不是狂熱主義.
是有理性的.
是Cool head.
我們會思考,會想東西.
去學習.
去分辨是非.
但很多時候教會就不是這樣.
或者很多時候信徒懶於做這件事.
你覺得是不是一個很災難性的事.
就是聖經你知道歷史歷代都是最暢銷的書.
如果你Google一下.
聖經最暢銷的書.
這麼多年印了多少本.
50至70億本.
去了哪裡.
去家裡封塵.
有些人幾本聖經在家裡.
但我們永遠都不會put effort.
去打開,去understand.
一看到的時候我們就說很悶.
搞不定.
我明白.
聖經裡面有很多東西.
不是那麼容易明白的.
但不代表我們有藉口.
我們說算了.
我們什麼都不管.
特別在我們這個世代.
大人都是一樣.
不要說小朋友.
我們會嘆息.
為什麼小朋友經常對著影片等等.
連我們大人都是這樣.
是吧.
有誰是被.
你會看到一段文字的時候.

$^{841}$你會學習思考.
沒有我們今天.
我們寧願看電視.
看screen.
我們被動的去接收這些資訊.
這樣就覺得是夠.
是不是.
是不行的.
大英姐妹.
聖經在我們手裡.
在我們電話裡.
我們就需要去嘗試去明白.
在嘗試去明白的過程裡面.
我們就不斷的去長進.
不斷的去前進.
不斷的去進步.
以致我們錯了的東西.
我們會更改.
然後改善.
會使得我們的信仰生命更加完備.
是不是.
應該是這樣.
因為這樣的緣故.
我今天仍然覺得.
那種所謂.
classical.
formal.
的神學訓練.
是需要的.
正統的神學訓練.
是需要的.
特別是對一些牧者來說.
想做牧者的話.
先讀好本書.
然後.
書也讀不好.
怎樣去.
牧養呢.
因為需要去裝備.
是不是.

$^{881}$最近.
這個不只是.
只是信仰生活.
或者教會.
需要有一些這樣的.
就是.
整全一點.
或者正統的一些訓練.
每行每業都應該要這樣.
是不是.
最近香港大學.
不是鬧了一個新聞出來.
查到幾十個大學生.
虛報學歷.
是不是.
進了香港大學.
是不是.
不只是.
我不明白香港的社會是這樣.
美國也可以有這樣的事情發生.
是不是.
虛報.
是不是.
What is true.
and what is false.
是不是.
大家要分辨.
是不是.
但是如果社會的文化.
已經不再說真理.
只是說拳頭的話.
指鹿為馬的話.
Who really care.
是不是.
正因為這樣的緣故.
在教會裡面也是一樣.
是不是.
教會.
普世教會.
都是需要.

$^{921}$有正規的訓練的.
做一牧者.
要在神樂園裡面.
特別是因為他做教導.
在神樂園裡面.
起碼讀三年多.
是不是.
就是叫做Master of Divinity.
道學碩士.
是不是.
是一個.
最長的一個program.
在Master的level裡面.
Graduate的level裡面.
是不是.
需要有正規的訓練.
拜讀位.
你知道神樂園.
這個字.
是什麼意思嗎.
Seminary.
是不是.
什麼叫做Seminary.
Seminary這個字.
在拉丁文.
是Seminarium.
意思就是說.
其實是.
花園.
溫室.
就是Nursery.
是那些.
栽種那些.
就是.
初苗的地方.
是這樣的意思.
神樂園就是說.
那群人進了神樂園的時候.
他們就有正規一點的訓練.
然後他們慢慢慢慢.

$^{961}$孕育他們的生命出來.
以致他可以.
移植.
去到不同的地方.
是不是.
這個就是Seminary的意思.
神樂園的意思.
是不是.
我自己在.
的經歷裡面.
我也覺得很有福.
因為我去過.
好幾間.
的神樂園.
第一間是Westminster.
他有兩個Campus.
一個在菲律賓.
一個在南加州.
在San Diego.
第二間是Wilson College.
在Vancouver.
那裡我們是.
當時讀杜學碩士.
然後Wilson College.
在我們芝加哥附近.
Wilson College.
那裡讀聖經考古.
現在我們在Midwestern.
就是我們其中一個.
的浸信會的神樂園.
在Kansas City.
那裡正在進修.
這個神樂園裡面.
你所得到的東西.
我想告訴你.
在北美.
我們有的神樂的資源.
是全世界首位.
One of the best.
所以你覺得.

$^{1001}$我跟那些神樂生.
去跟他說.
我說因為你是神樂生.
無論你是Moody也好.
你是Trinity也好.
你周圍的神樂園.
你是在讀書的.
我有時候我會跟他說.
我說我對你的要求是高一點的.
你覺得這樣是不是很Harsh.
一點都不Harsh.
怎麼會Harsh.
我說你拿的是世界裡面.
最好最好的神樂訓練.
怎麼可以你們生命裡面.
是這麼sloppy的.
怎麼可以你們這麼差.
是不是很Harsh.
不是很Harsh.
特別今天有很多Grade Inflation.
他們在神樂園裡面.
讀的班什麼都拿A.
什麼都拿A.
你給我看看.
因為想哄他們.
他們繼續交學費.
有這樣的Dynamic.
但你想想.
他們擺出來的生命.
到底是不是雙程.
我記得第一間神樂園.
他們有一個口號.
口號是The whole council of God.
The whole council of God.
就是上帝的整全的心意.
我們需要去掌握.
我相信這個不只是.
對一些將來會做牧者的人.
是能夠應用的.
對於我們每一個信耶穌的人.

$^{1041}$都是能夠應用的.
我們需要明白The whole council of God.
上帝整全的心意.
我們需要了解.
Back to basic.
我希望你能夠被挑戰.
不要怕.
你去學習一下.
因為我知道我們大家.
不同的Learning style.
是有些不同的.
但我們也需要學的.
我希望你能夠帶走的功課.
起碼每一天.
在你最好的時間.
你整天最好的時間.
無論是早,午,晚,到.
你自己選擇.
每個人都不同.
起碼花十分鐘去閱讀聖經.
閱讀有什麼好處.
閱讀是一種主動的學習.
主動的吸收.
對嗎.
閱讀.
當然主送不是.
主送看YouTube容易很多.
但是閱讀.
是必須有這一種.
所謂修養.
有一種修養去學習.
如何主動的去學習.
不是被動的去接收.
畫面上給你的東西.
圖畫或影片.
給你的東西.
這些全部都是我們被動的去接收.
如果你想critically engage.
自己的mind.
批判性的去作思考.

$^{1081}$去學習思考.
就是需要做這件事.
閱讀.
特別在我們北美.
有很多很多好的書.
甚至現在開始有些好的書.
已經很快翻譯出來了.
甚至我們華人的學者裡面.
有很多很好的學者.
在這裡你純粹找一下.
你都找到一些初信主.
可以讀的書.
這些書不是那些.
你不要嫌棄牧師在搞學術.
不是學術.
不是所謂我們講的學術.
像牙塔那些學術.
而是這些是實在的書.
J.I.Pekka的認識神.
暢銷了很多年.
是一個很簡短.
很簡明的系統神學.
讓我們能夠了解真理.
最近Tim Keller已經過世了.
回到天教.
寫了很多好的書.
The Reason for God.
特別是給一些可能想接觸一下信仰.
未信主的人.
可以去看的.
New York Times Bestseller.
Simply Question.
Anti-Right.
這些是一些很新的書.
比較新的書.
但是一些很基本的書.
不久之前.
我們有謝霆博士在我們中間.
來到我們識經講座.
謝霆博士也有很多.

$^{1121}$很好的書 特別是詩篇的書.
我想講的就是.
學習一下去閱讀.
學習一下去到.
在我們熟悉的生命裡面.
Spiritual Growth裡面.
去到做一些閱讀的功課.
對嗎.
我們圖書室 不是圖書室.
因為我們圖書假.
上面有很多好的書.
嘗試去到Engage自己.
有一樣東西.
我想加上去的.
去背誦一些金句.
你覺得可能是一些.
很幼稚的東西.
很初班的東西.
因為有時候一些基礎是需要的.
今天我選了這段聖經.
為什麼選這段聖經.
我等一下解釋一下.
這裡講的是耶穌講的話.
馬達科罕第22章.
祂說你要盡心盡性盡義.
愛主理神.
這是最大的.
第一條誡明.
大家背了它好不好.
想一下這裡是什麼事.
為什麼選這段.
With all your heart.
All your soul.
And all your mind.
你的思考.
都要盡心盡力的思考.
來愛上帝.
當你思考的時候就是愛上帝.
當你盡心去認識.
真理的時候.

$^{1161}$去看清楚聖經裡面講的.
和我們一起去研考聖經裡面講的.
這就是.
That's how you love God.
With all your mind.
以致我們不會成為一些.
信仰的生命裡面.
脆弱的人.
搖風擺樓.
有不同的思想.
有不同的假的教導出現.
或者似是而非的事情出現.
我們懂得去分辨.
盡心盡性盡義.
All your mind.
你能夠去愛.
我們的主.
我們守聖山之前我們一起低頭祈禱.
上主我們多謝你.
因為你賜給我們生命.
這個永生的生命.
讓我們能夠享受.
去Engage 去啟動它.
主我們需要.
去感受到.
你自己的同在.
需要感受到.
我們和你之間的關係.
但同時我們需要用我們的.
腦筋.
我們的思考.
去思想你.
去默想你.
幫助我們分析.
幫助我們去了解.
幫助我們去成長.
求主你裝備我們每一個.
不單止我們牧者.
每一個相信主的人.
基督徒.

$^{1201}$大英前輩.
無論我們的背景是怎樣.
無論我們自己本身的.
過往是怎樣.
主我們今天.
能夠認識你.
我們就需要全心全意.
全力的去做.
讓我們盡我們的心意.
盡我們的思考.
我們的思想.
學習一下有條有理的去認識信仰.
去明白你.
去說話的上帝.
你介紹你自己是怎樣.
在歷史歷代在聖經裡面.
你讓我們知道你是一位.
全善全美全能的神.
主幫助我們信你的資源.
讓我們能夠繼續追求.
去明白去認識你.
祈禱奉祝耶穌基督的名字.
我們就要做到這一點.
\newpage



\section{腓利門書}
\label{sec:eiPsrBdDRJk}
\textbf{請勿超越黃線 (腓利門書) - 譚文倩莊傳道}
\newline
\newline
連結: \href{https://youtube.com/watch?v=eiPsrBdDRJk}{\texttt{https://youtube.com/watch?v=eiPsrBdDRJk}} ~~~~ 語音日期: 2024-07-17
\newline
\newline
\hyperref[sec:x0e0tuTJkmQ]{\small{< < < PREV SERMON < < <}}
~
\hyperref[sec:index_chronic]{\small{[返順時目]}}
~
\hyperref[sec:index_scriptual]{\small{[返順卷目]}}
~
\hyperref[sec:QOZ3Ws7CSlE]{\small{> > > NEXT SERMON > > >}}
\newline
\newline
$^{1}$早晨各位弟兄姊妹.
早晨.
雖然今天天空要下雨.
但也沒有我們弟兄姊妹.
回到教會裡敬拜我們的神.
我想給大家一點時間.
跟前前後後左左右右打個招呼.
今天早上回來.
大家準備好聆聽神的話語了沒有.
還記不記得這個暑假.
我這次回來.
講了兩堂《肥理門書》.
今天是最後一堂.
重溫一點.
第一堂的題目叫《星》.
是講述福音如何重新定義人與人之間的關係.
第二堂《愛回家》.
講述教會群體的功課.
今天第三堂.
請勿超越黃線.
Please stand behind the yellow line.
如果上兩堂你沒有在我們當中.
或者在網上看的時候.
我現在嘗試用一點時間.
跟大家重溫.
做個總結.
《肥理門書》不是一本深入神學的書卷.
但它對基督徒人際之間的關係.
和教會群體的本質.
提供了很多很寶貴的見解.
根據經文推斷.
保羅第一次在羅馬被監禁的時候.
為奴隸亞尼西牟寫了一封求情信.
要他親身將這封信帶去哥羅西.
因為他冒犯了他的主人肥理門.
普遍推測是這個奴隸偷了他主人的錢.
之後就逃跑了.
到了羅馬.
不知在什麼情況下認識了保羅.
藉著保羅信了耶穌.

$^{41}$之後他還成為了保羅的福音夥伴.
和他一起宣講天國的福音.
在信中保羅稱呼亞尼西牟為孩子.
甚至願意將他虧欠主人的「濫」上身.
代他還債.
保羅這個動作凸顯了福音徹底和解和寬恕的大能.
另外一個可能性是肥理門主人.
將這個奴隸拆了出去辦一些事情.
結果這個奴隸去了之後一去不返.
家裡少了一個奴隸.
不多不少都令到他的主人有所損失.
無論如何.
這封信的主要目的.
是保羅懇求肥理門重新接納亞尼西牟.
因為亞尼西牟已經信了主耶穌基督.
成為他自己的個人救主.
所以保羅在裡面說.
「不再是奴隸.
而是高過奴隸.
是親愛的弟兄」.
這封求情信.
表面收信的人是主人肥理門.
但實質收信的人.
其實是肥理門家裡一起聚會的弟兄姊妹.
一個相信並且跟隨基督的群體.
他們都要寬恕和接納這位弟兄.
因為寬恕和接納是基督耶穌樹立的榜樣.
保羅的心願是教會這個群體.
能夠反映基督的生命.
福音的文化.
今日第三堂請勿超越黃線.
請勿超越黃線.
我們先低頭祈禱.
出來主掌管我們生命的神.
聖靈24/7在我們生命裡為我們祈禱.
提醒.
指引.
願人都尊你的名為聖.
願你的國降臨.
願你的旨意行在地上.

$^{81}$如同行在天上.
阿們.
不知道你們有沒有留意.
肥理門書裡面.
短短的二十五字前後提及了十個人名.
當中有三個主角.
這三個主角就是保羅,肥理門和亞彌西某.
但整封信其實只有一個人在說話.
而耶穌基督從來沒有出現過.
在肥理門書裡面的三個主角.
其實又如何被沒有出現的耶穌影響了他的行事為人.
和人之間的關係呢.
今日我嘗試從三點跟大家一起思考.
第一點.
在基督裡關係不應寄予公平.
肥理門書一章十六節.
十八至十九節的前半部.
「不再是奴隸,而是高過奴隸,是親愛的弟兄.
對我確實如此,何況對你呢?.
無論在獄身或在主裡更是如此.
他若虧負你或欠你什麼.
都算在我的帳上吧.
我必償還,這是我保羅親筆寫的」.
在當時的羅馬律法上.
在道理上肥理門若果要懲罰.
甚至處決逃走的奴隸亞彌西某.
絕對是合情合理.
一個最公平最對的解決方法.
但保羅卻站在福音的立場上去看亞彌西某.
不再是奴隸,而且是高過奴隸.
是親愛的弟兄.
在主劇裡.
信徒之間存著一份弟兄姊妹的情誼.
因為我們大家都是天父的孩子.
天父願意我們彼此相愛互相服侍.
雖然我們有不同的膚照,不同的崗位.
但歸根究底.
在主裡面,我們是一家人.
我們的信仰不是要廢除.
或者混淆了我們生活在這個世界的責任.

$^{121}$相反,反而更加要加強我們在這些責任上的義務.
神的話語更加指導我們如何正確去履行這些責任.
毫無疑問,亞彌西某是真誠的悔改.
否則,她不會願意拿著這封求情信.
我之前說過,是說1645個mile的路程.
回到主人那裡去求赦免.
亞彌西某通過保羅,承認自己的錯誤.
你能不能夠想.
1645個mile.
或者容易一點,開車.
不吃不去廁所.
不會停下車去跟一個人說對不起.
我試過跟我老公吵架.
從客廳去到房間跟他說對不起.
都覺得很難.
不是說對不起這麼簡單.
亞彌西某做的這件事,是在說分分鐘是沒命的.
亞彌西某在這個旅程裡面.
其實她有很多時間和空間去想.
不如不要回去.
可以回頭.
走去別的地方.
辛辛苦苦長途跋涉帶封求情信回去.
依照當時的文化.
肥美們還要在教會眾弟兄姊妹面前讀這封信.
其實面皮薄一點都搞不定.
不要忘記之後她還可能要接受肥美們決定要怎麼做.
就好像我剛才所說.
甚至要決定處決她.
但因為基督.
因為如今大家是主裡面的弟兄姊妹.
領受著同一份大家都不配的恩典.
所以不要再計較.
不要再去想公不公平.
虧欠他們多.
第二點,在基督裡,關係應該是基於恩慈.
肥美們書一章十四節.
但不知道你的意見,我不願意這樣做.
好使你的善行不是出於勉強.
而是出於自願.

$^{161}$在這裡看到保羅用一個很尊重對方.
很handle with care的態度.
引領肥美們要恩待亞利西茂.
當人真真正正明白福音的時候.
那個甘心樂意,很自然的回應是最好的.
在這裡保羅很明顯明白這一點.
所以他沒有用他的權力.
他是用一份很尊重對方,保護對方的態度.
去跟肥美們對話.
去承認亞利西茂是一個奴隸.
這個奴隸是屬於他主人的.
所以如果他不同意的話.
他不會勉強隨便留下亞利西茂在他身邊.
亞利西茂擅自離開.
令到他的主人有所虧損.
這是無可置疑的.
所以保羅很靠肥美們的理據.
是大家在主裡面關係的基礎.
在主裡面我們樂意用基督的恩賜.
去重新接納亞利西茂.
肥美們書一章十八至二十節.
他若虧負你或欠你什麼.
都算在我的帳上吧.
我必償還.
這是我保羅親筆寫的.
我並不用對你說.
甚至你自己也虧欠我.
弟兄啊,希望你使我在主裡因你得益.
讓我的心在基督裡得到舒暢.
保羅親手寫的這封求情信.
當他對肥美們說.
把亞利西茂虧負或虧欠.
都算在我的帳上吧.
實際他是為保羅,學校,耶穌.
為我們所做的.
把我們的罪債背在自己身上.
保羅愛亞利西茂.
從他把自己帶入亞利西茂的位置上.
甚至好像是他自己犯了錯.
就正如基督為我們在父神面前所做的一樣.

$^{201}$保羅同樣為肥美們.
為亞利西茂在肥美面前做.
這就是基督的恩賜.
第三點,在基督裡,關係基於福音.
肥美們書一章十至十一節.
就是為我在菌所中所生的兒子亞利西茂求你.
從前他與你沒有益處.
但如今與你我都有益處.
在某個程度上.
亞利西茂對保羅來說是有益處.
或許是因為保羅被軟禁時.
他成為了他的得力助手.
相對因為他逃亡.
所以對他的主人來說是沒有益處.
當保羅說到亞利西茂沒有益處或有益處時.
其實當中是有些玩字的含義.
因為亞利西茂這個名字.
直譯的意思正正是有益處.
信了主與保羅同宮的亞利西茂.
如今真是人如其名,是有益處.
但有益處也好,沒有益處也好.
在福音裡都變成有益處.
因為他對保羅來說有益處.
所以間接對肥美們來說也是有益處.
所以保羅才在這裡說.
對你和對我有益處.
保羅要肥美們明白.
亞利西茂因為信主而升了級.
所以他與亞利西茂的關係也升了級.
這個升了級的關係.
彰顯福音的恩典.
重溫三點.
在基督裡關係不應基於公平.
在基督裡關係不應基於恩慈.
在基督裡關係基於福音.
在基督裡.
肥理門書對基督徒人際之間的關係.
和教會群體的本質.
我之前說過.
也提供了很多寶貴的見解.

$^{241}$這些見解要反映福音的大能.
基督在人生命裡活著的印證.
因為我是一個傳道人的關係.
我經常有機會和別人聊天.
很多時候我發現.
人與人之間.
有很多時候無法化解一些衝突.
其中一個很大的原因.
是我們對對方沒有恩慈.
是他不對.
他沒有理會過我的感受.
我最討厭他就是這樣.
你知不知道他多討厭.
說了很多次都不聽.
他和我說話的時候語氣很差.
不該也沒有說話.
以前聽過一個牧師有個比喻.
他說一個結婚的男人.
不會忘記自己是一個結婚有老婆的人.
他不會無緣無故和另一個女人去沙灘漫步.
雖然他老婆不在身邊.
但作為老公的身份.
他做的每一件事都不會當他老婆無道.
基督徒的生命在人際關係裡.
或者特別當我們處理衝突的時候.
其實我們會不會當主無道.
我們只會想念公不公平.
對我來說不公平.
想念對方有多離譜.
他所做的傷害我多少.
虧欠我多少.
甚至對方願意說對不起也好.
也不容易放過他.
主耶穌基督去了哪裡.
主耶穌基督的榜樣去了哪裡.
耶穌的確沒有在《肥利門書》裡出現過.
但其實從頭到尾.
其實你留心思想的時候.
其實從頭到尾他都在.
因為如果不是耶穌的關係.

$^{281}$保羅,肥利門,亞彌西某.
就不需要藉著福音的關係.
去處理他們之間的關係.
再者保羅讚賞肥利門有信心.
對弟兄姊妹有愛心.
勸勉他將他的愛延伸到他的奴隸亞彌西某身上.
就正正是因為主耶穌基督的愛.
就是一個這樣的愛.
不是一個pick and choose的愛.
愛心的對象是所有聖徒.
所以這個所有是包括了這個奴隸.
再者保羅對福音的毀身.
肥利門願意延伸他的愛心.
亞彌西某願意承擔責任.
誠心悔改.
信服等等.
都是全部因為基督在他們生命作供的關係.
有影響力.
這個就是在基督裡.
在耶穌裡面的意思.
耶穌的確沒有在肥利門書裡面出現過.
但是我們不難看到.
耶穌就在這三個主角的生命裡面.
反省我們自己.
我們可能會帶著一條十字架的項鍊.
或者我們有些人會貼一條魚的symbol.
放在我們的車後面.
但是我想.
重要的是我們苦心自問.
耶穌只是一條項鍊掛在我們的項鍊上.
貼在我們的車的後方.
還是住在你的心裡.
在影響你.
人與人發生衝突的時候.
你的反應會是怎樣.
你會不會突然覺得自己需要停下來.
想一想.
反省一下自己.
自己在當中有沒有責任.
會不會想一想聖經的教導.

$^{321}$很有趣的.
有時候你想想我們家人對我們來說就是最親的人.
但很多時候我們傷害得最多的就是我們自己的家人.
例如父母.
我們的丈夫.
太太.
子女.
我們是因此相待.
還是在你的眼中.
你經常看到對方有很多事情都不妥.
有很多事情都不合你心意.
所以很容易有衝突.
很容易會吵架.
其實有時候想想.
可不可以不要咬著不放.
可不可以不計算.
可不可以對對方多一點恩賜.
你知不知道當你放過對方的時候.
其實你同時也在放過自己.
其實生氣也不是很好受.
前幾天在一個紀念服務.
聽了一首很古老的詩歌.
Amazing Grace.
How sweet the sound.
That saved a wretch like me.
Amazing Grace.
之所以Amazing.
是因為其實這個Grace是你和我都不值得的.
這個恩典是你和我都不值得.
我們沒有做過什麼.
甚至我們虧欠很多.
但是神竟然給了我們.
我們也會求主恩賜對待.
但是當有人激怒我們的時候.
我要和他拼命.
保羅如何勸勉肥理們.
保羅基於什麼大膽去作出這個要求.
肥理們為什麼又願意原諒去接納.
我最近和一個我的好姐妹分享.
我們各自面對的挑戰的時候.

$^{361}$說著說著她突然說了一句話很觸動我.
她說每一步求信服不是求結果.
每一步求信服不是求結果.
在不同意的情況下.
你要贏.
要對方屈服認輸.
甚至最好盡快解決問題不要煩我.
還是你願意操練節制.
不讓你自己的感覺.
令你很不舒服的感覺.
覺得不公平.
負面的情緒去牽引你作出回應.
還是你會想一下作為基督徒.
我這個基督徒的身份.
你對你自己有一個要求.
那個要求是你每一步都求信服.
不是求結果.
《迦勒提書2:20》有句經文.
我們大家都很熟悉.
「如今活著的不再是我.
乃是基督在我裡面活著」.
基督在你裡面活著嗎?.
在你的生命裡面活著嗎?.
在你行事為人的裡面活著嗎?.
我們人很多時候的確很自然.
會隨著一些令我們很不舒服的感覺.
用我們的情緒或脾氣去回應對方.
我想起很多年前.
有一次我弟弟下班回來的時候.
我都不太記得發生什麼事.
總之他就發我很大脾氣.
我真的覺得很無神.
回家他就罵我.
當我很想第一個反應去罵他的時候.
我突然想起我弟弟今天上班一整天.
他很辛苦,他很壓力.
所以他的心情才這麼差.
然後我再想想.
其實我有時候心情都不太好.
都會向我弟弟發脾氣.

$^{401}$我弟弟都很硬吃我.
我疼我的弟弟.
我弟弟都很疼我.
所以結果我沒有罵他.
我還立即走進廚房煮了一些東西給他吃.
哄他開心.
難不難?.
聽起來好像覺得有點難.
但其實又不是真的這麼難.
因為愛所以體諒.
願意因此相待.
因為愛基督耶穌.
所以我們願意信服.
我們願意操練因此.
還要一次又一次又一次地操練.
因為我們都軟弱.
都會有失敗的時候.
我不是每次都不回應弟弟.
但不要氣餒.
不是這樣就OK.
繼續努力,不要氣餒.
因為操練就像我們人去做運動.
你操得越多就越適合.
每一步求信服.
不是求結果.
(咳嗽).
對不起.
(咳嗽).
還是你已經畫了一條線給耶穌.
請站在黃線後面.
不是開玩笑.
還告訴他要留意空隙.
我是基督徒.
不過我是這樣的.
這是我的性格.
這是我.
我相信你,但有些事情你不要理會我.
我需要找你的時候,我會找你.
我相信你,但這件事不可能這麼容易就算.
請站在黃線後面.

$^{441}$耶穌請你站在那一邊.
但耶穌好像真的不在.
最近我和女兒聊天.
我和她爸爸聊天,要處理一些事情.
談了很久,其實我們都不能找到大家同意的位置.
最後我只可以和女兒說.
我希望你明白,爸爸和媽媽很愛惜你.
我們有責任去保護你.
如果你願意的話.
雖然你不明白,雖然你不同意.
你都願意遵從我們的意願.
答應我們的要求.
但我自己也知道.
所以我和女兒說.
但你已經長大了.
如果你不願意的話.
其實我也沒有理會.
我們讓女兒自己去抉擇.
很感恩女兒最後願意尊重我們.
接受我們的提議.
我想我們和神的關係.
有時就好像我媽媽和女兒的關係.
有時的確有些事情我不明白.
有些事情我很想神快些答應我.
或者幫我解決.
為何不是呢.
但到了這些時候.
有時我們大家要計算.
我有多著重我和主耶穌基督的關係.
我要信得過祂.
正等於我女兒要信得過我.
給她的提議是為她好.
而願意放下自己.
你看不看重你自己和耶穌基督的關係.
《肥利門書》一章六節.
「我求神使你與我共享的信心發揮功效.
使你認識在我裡面一切美善的事.
可以更認識基督」.
作為一個基督徒.
你的生活做的每一件事.

$^{481}$與你和人之間的關係.
有沒有基督的樣式呢.
你有沒有希望認識基督多些.
你知不知道其實當我們操練愛心信心的時候.
其實就是你可以藉著操練去認識耶穌基督.
在你的生活裡面.
你身邊的人看不看到耶穌.
耶穌是站在黃線的另一邊.
是因為你情願祂站在黃線那一邊.
請站在黃線那一邊.
耶穌請你站在那一邊.
聖經裡面提醒我們.
不要做一個忽冷忽熱半坑的基督徒.
因為這樣的基督徒.
主是情願在他的口裡吐出來.
吐出來.
其實再想想.
其實為什麼不想一是不做.
要做就立志做一個有餡.
生命有改變.
愛人多過傷害人的基督徒.
做一個能夠彰顯基督的基督徒.
但願這本只有一張二十五節的《肥利門書》.
不單提醒我們願意接受教導.
遵從聖靈的引領.
降伏在福音的面前.
更願意讓福音去掌控你的意願.
讓主耶穌基督真真正正在你的生命裡面作動.
去彰顯.
你願意去原諒.
去接納.
去重建.
重新人與人之間的關係.
我們弟兄姊妹是同心合力地一起建立.
一個擁有福音文化.
以基督為基礎的教會.
讓一些未認識基督的人.
能夠藉著我們這個教會群體去認識基督.
而你和我弟兄姊妹已經認識基督的基督徒.
我們是終於回到這個屬靈的家.

$^{521}$一個彼此相愛的屬靈的家.
最後我更加希望大家看到.
福音,基督如何影響保羅,肥利門,亞尼西謀.
他們互相恩慈對待.
是因為表面沒有在章節裡面出現的耶穌.
在這三個人的生命裡面活著.
耶穌在這裡.
耶穌又在不在你的生命裡面.
最後我想挑戰大家.
去想想是不是時候不要再做一個懶得做的基督徒.
要做就做一個有基督在裡面.
在我生命裡面的基督徒.
你信神.
更加立志去信服神.
因為你緊張這個關係.
你愛耶穌.
你願意遵從聖經的教導.
願意聽從聖靈的引領去作出改變.
讓人可以在你身上見到基督.
你是不是真的知道.
「多麼的費用,在那一頂基督的身上,我所犯的罪」.
耶穌為你的罪釘在十字架上所付出的代價.
主的恩慈.
Amazing的恩典.
在我開聲禱告之前.
我很想給你一點時間在自己的座位那裡為你自己祈禱.
或者你已經信主很多年.
其實開始已經.
那條屬靈生命的線是這樣的.
其實你不想是這樣.
You don't want that.
你想繼續成長.
你想認識神更多.
不應該就是這樣.
如果你想再一次立志.
在神面前立志.
真的覺得想一想.
我真的都懶了一段時間.
沒有成長.
我想重拾和神的關係.

$^{561}$我想基督住在我裡面.
影響我的生活.
I'm excited to live for Jesus.
你可以在你自己的座位上為你自己祈禱.
如果你需要一些guidance.
需要一些實質的幫助.
崇拜之後.
或者你ready之後.
歡迎你和我們教會聯絡.
我們一起低頭祈禱.
在這裡我會給大家大概一分鐘的時間.
在你的座位上為你自己祈禱.
之後我會結束.
(靈性音樂).
我們的神.
你愛我們.
愛到底.
願意聽到你聲音的人.
回應你對他們的呼召.
你在人的信服裡面.
成就你要成就的事情.
奉主耶穌基督的聖名祈求.
阿們.
(靈性音樂).
\newpage



\section{3:10-15}
\label{sec:QOZ3Ws7CSlE}
\textbf{Back to Basics \#3 (2 Timothy 3\_10-15) - Rev. Mark Chan-}
\newline
\newline
連結: \href{https://youtube.com/watch?v=QOZ3Ws7CSlE}{\texttt{https://youtube.com/watch?v=QOZ3Ws7CSlE}} ~~~~ 語音日期: 2024-07-18
\newline
\newline
\hyperref[sec:eiPsrBdDRJk]{\small{< < < PREV SERMON < < <}}
~
\hyperref[sec:index_chronic]{\small{[返順時目]}}
~
\hyperref[sec:index_scriptual]{\small{[返順卷目]}}
~
\hyperref[sec:QovpodRHlhw]{\small{> > > NEXT SERMON > > >}}
\newline
\newline
$^{1}$- Hey, good morning everyone, good morning church..
And let's turn to our scripture today..
And if you have a Bible, turn to 2 Timothy chapter three..
And we're going to study this chapter.
and take root in the word of God..
But first of all, we have the verse of the week, right?.
Remember the first one is Psalm 119, 105, right?.
Can we recite it together?.
Let's do that..
Psalm 119, 105..
Your word is a lamb to my feet and a light for my path..
Wonderful verse, right?.
You got that memorized?.
I'm sure you do, right?.
And then the second one, which is last week,.
which is last week's verse of the week..
And this is from Matthew, and that's what Jesus said..
Let's recite it together..
You must love the Lord your God with all your heart,.
all your soul, and all your mind..
This is the first and greatest commandment, all right?.
Got it memorized, right?.
And today I'll give you another one, all right?.
But before that, we're going to dive into chapter three.
of the book of 2 Timothy..
And as I told you, this chapter, 17 verses long,.
there were two parts..
The first part, which I talk about,.
and we studied last week,.
was about the dangers of neglecting the word of God,.
of the ignoring, of not knowing the word of God..
That is chapter three, verses one to nine..
And today we come to the second part of the chapter,.
2 Timothy chapter three, verses 10 to 17..
And that chapter talks about the benefit.
of knowing God's word, of learning God's word..
Speaking of the first part,.
the dangers of neglecting the word of God..
And this is a real danger because as you remember,.
in the first nine verses, Paul is talking to his mentee,.

$^{41}$his disciple, his disciple pastor, Timothy..
And he warns him that there will be a very difficult time.
in the last days..
And I'll remind you, the last days in New Testament terms.
is the period of time, indefinite period of time.
between the first coming of Christ.
and the second coming of Christ, which lies in the future..
So that period of time is called last days.
in the New Testament..
And we are actually in the last chapter of human history,.
according to biblical timeline..
So we are in the last day and Paul is not naive.
what is going on in our time..
Remember how he characterized our time,.
dangerous and violent, and people are lovers of themselves.
and they slander others and have no self-control..
They are haters and they hate what is good..
And these are indeed our time, right?.
And we see right in front of our eyes,.
political violence that happened.
in the Trump Valley yesterday..
It was terrifying, right?.
We are left speechless hearing the news.
and watching videos of such senseless violence..
And my prayers go out to all those who were traumatized.
by this yet another incident, right, in America..
And we're in Chicago, we are not foreign.
to those incidents, right?.
It happens almost every single day..
The other day, I think a couple of days ago,.
I saw a news reporting..
There was another gun violence happening.
not too far from us in Chinatown..
And again, it's gun violence, it's senseless violence..
What can we do, right?.
We ask the question, what can we do?.
It is easy for us to isolate ourselves, right?.
You know, in times of darkness, in times of uncertainty,.
we just wanna take care of ourselves..
And it's very easy, it's very natural for us.

$^{81}$to just hide in our own comfort zone, right?.
We just do our own things, right?.
Make sure everything's under my control, right?.
Look what Paul is doing here..
Paul in chapter three, he knows the time..
He warns about it..
And the last days will be very difficult..
It's indeed very difficult for Christians, right?.
In America, we are relatively better, right?.
Because we have liberty, we have freedom, right?.
I've been meeting with some pastors online.
on a weekly basis in the past two plus years..
Yeah, it's amazing to meet with pastors.
from the communist China..
We have been meeting almost every single week.
in the last two weeks..
And we started toward the end of the pandemic..
And we're still meeting together..
When we meet together, there are a few of them,.
they never turn on their camera, of course..
And there are times we share, we pray for one another..
There are times I give some instructions..
It's just amazing how long it has lasted, right?.
Because originally we thought that we just meet together.
for three months..
And then it continues and continues and continues.
until today, right?.
So on a weekly basis, we meet together..
But I'm also amazed by their faith..
I'm also amazed by their faith..
They cannot go open with what they learn.
and what they preach..
They are often under persecution..
In America, we are pretty easy, right?.
We're pretty easy..
We're quite comfortable..
However, what Paul says remains true, right?.
Being Christian these days can be difficult, right?.
Being a solid Christian is very difficult.
because there's so much temptation out there..

$^{121}$There's so much distraction out there..
And Paul knows the time of last days..
And he described the last day in vivid terms..
And Paul is not naive..
And we should not be naive, right?.
So he continues to talk about not only the dangers.
of not knowing God's word,.
but the benefits of knowing God's word, right?.
Yes, the time is difficult..
The time is chaotic..
We can easily be hijacked by fear and uncertainty, right?.
Talking to some of the friends and some of the people,.
we are still trying to come out of the pandemic fear, right?.
Just there's so much fear out there,.
so much uncertainty out there..
And Paul is saying to us,.
remember this is a very difficult time, right?.
I still remember during those years in the pandemic,.
there was so much misinformation out there..
There was so much misleading content online.
and conspiracy theories and sensationalism out there..
And what happened today, yesterday,.
the senseless violence,.
immediately, the social media is flooded.
with all sorts of, you say,.
conspiracy theories or falsehood or something like that..
And people just feed on those, feed on those things..
To combat against the spread of unbelief and falsehood,.
Paul suggests two things..
And he talked to Timothy about two things..
First, he says, remember,.
you need to stand on the foundation of the Word of God,.
on truth, on biblical truth..
And the first thing he says is to learn from good leaders..
Remember the good examples in your life and follow them..
And secondly, which is the last two verses, 16 and 17,.
he talks about the importance of scripture, Holy Scripture..
Well, let's take a look at what Paul is talking about here..
Verses 10 and following..
He said, he talked about the last days,.

$^{161}$the difficult time in vivid terms..
And in verse 10 and following, he said,.
but you, Timothy, certainly know what I teach.
and how I live and what my purpose in life is..
And you know my faith, my patience, my love, my endurance..
You know how much persecution and suffering I've endured..
You know all about how I was persecuted.
in Antioch, Iconium, and Lystra..
But the Lord rescued me from all of it..
Yes, and everyone who wants to live a godly life.
in Christ Jesus will suffer persecution..
But evil people and the imposters will flourish..
They will deceive others and will themselves be deceived..
This is a dangerous time, right?.
But Paul is reminding Timothy,.
you have great example in front of you,.
and I am one of them..
And Paul is reminding Timothy, you know my faith, right?.
You know my patience, you have learned from me..
I'm not just doing empty talks..
I live the faith, I walk the faith, and you have seen it..
Right, because Timothy has been following Paul.
in many different missionary journeys..
You can check from the book of Acts, right?.
And Timothy follow Paul to different mission fields.
and different cities in the Roman world.
and preach the gospel..
And it was very, very difficult.
because the Jews were after him.
and other people opposing to the gospel..
And Paul endured all of this..
You know, isn't it nice we have great examples.
of faithful followers of Christ in our life, right?.
Now I want you to do one thing, think about it..
Are there any good examples.
of faithful Christians in your life?.
Are there any?.
You can think of something, some names, right?.
And you can think of some faces, right?.
And Paul is older and wiser and more experienced.

$^{201}$than Timothy, and Paul urges Timothy,.
follow my example, right?.
See my example..
I myself have benefited so much.
from great mentors in my life..
Remember, oh boy, it was mid-90s..
It was mid-90s, I was in seminary..
And I was required to do internship.
when I was studying in a seminary..
And part of the credit hour has to be internship.
at a local church..
And I want to tell you honestly,.
you know, in North America,.
the design for seminarians internship.
is actually not quite well designed, okay?.
All right, this is a trade secret, right?.
You know, in the MDiv program, for example,.
yes, you're required to do maybe two terms of internship,.
but it's not very solidly or well-crafted program..
Field education, so-called, it's not really the priority..
Most of the priority we spent our time in MDiv.
was in the classrooms, right?.
With the books, right?.
And most of the internship programs,.
yes, they're required, but the evaluation is very loose..
Usually it's just pass, not pass..
Pass, fail, right?.
I mean, how easy for you to pass a class, right?.
You just submit something, right?.
Internship report, you know, you can just make it up, right?.
You make it up and submit it,.
and the teacher or the professor will grade as pass, right?.
So it was very, very easy, right?.
And think about it, it's quite sad, you know?.
It's quite sad..
You know, many of the students that I know.
did not learn so much from field education, right?.
They might learn well with exams and papers..
They might do well in their homework..
However, dealing with people on the field, right?.

$^{241}$You know, we're doing church ministry, right?.
Church ministry, the essence is people, right?.
How to read people, how to understand people,.
how to relate to people,.
so often is very, very low in priority..
Whether you can learn something solid from the field,.
from field education or internship,.
often the seminary student is the only one.
who is responsible, right?.
If he is smart enough, if he is wise enough,.
he will make it happen..
If he's not so quick to learn the field education,.
then it's gonna hurt, right?.
It's gonna hurt..
Notice what Paul says here..
The emphasis is the personal pronouns, right?.
You know, not you know the faith,.
you know my faith, my patience, my love, my endurance, right?.
In the same way, speaking to all of us as Christians, right?.
We come to the church often with a mentality.
that is not biblical..
We come to church often with a mentality.
of being a consumer, being just passive audience, right?.
We come here, beautiful music, right?.
Beautiful environment, you know, nice coffee we have,.
nice snack we have, and refreshment..
The church is nice because, you know,.
you have great service or something like that..
I get what I really want, and when people let us down,.
you know, church are full of broken people, right?.
Should I remind you that, right?.
We're all broken..
The pastors are broken..
The leaders are broken..
The laymen are broken..
Everyone is broken..
Everyone is sinner, and we all need Christ, right?.
But we come in with a consumer,.
well, religious consumer kind of mentality..
We will be disappointed..

$^{281}$We can always be disappointed..
You know, I can tell you from my experience,.
I'm often disappointed by other people,.
or even I'm disappointed by myself, right?.
And in some sense, with that kind of consumer mentality,.
we feel dissatisfied, right?.
Because we don't receive good service,.
like in the restaurant, right?.
You go to a restaurant, you don't get very good service.
from waiter or waitress, right?.
You feel upset, right?.
Feel upset, you feel being used,.
because you come in with the consumer,.
you know, customer mentality, right?.
Lousy waiter, lousy food and poor service, bad quality,.
and then give you one star, right?.
Write about it in your review, right?.
But that sort of mentality, religious consumerism,.
is unbiblical, right?.
It's in fact anti-biblical, right?.
On verse 14 and following, Paul said,.
"But you must remain faithful.
"to the things that you have been taught..
"You know they are true,.
"for you know you can trust those who taught you..
"You have been taught the Holy Scriptures from childhood,.
"and they have given you the wisdom to receive salvation.
"that comes by trusting in Christ Jesus.".
Now, Paul is encouraging Timothy to look beyond himself..
Remember the first part, he said,.
"Look to me, my patience, my teaching, and my faith,".
right, and follow it..
And then Paul now is encouraging Timothy.
to persevere in difficult times..
And he said, "Remember the Scripture you have been taught.".
And also recall something very, very important..
Paul says to Timothy,.
"You have been taught the Bible from childhood," right?.
Since he was a little kid, since he was a baby..
And Timothy, as we know from the book of Acts.

$^{321}$or from other Scriptures,.
he has a godly mother and godly grandmother, right?.
Now, we're reading chapter three of 2 Timothy,.
but early on in the first chapter,.
Paul actually mentioned the name.
of Timothy's mother and grandmother..
In chapter one of the same letter, same book,.
Paul opens with these words,.
"I remember, Timothy, your genuine faith,.
"for you share the faith that first filled.
"your grandmother, Lois, and your mother, Eunice..
"And I know that the same faith continues strong in you.".
What is the name of Timothy's mother?.
Lois, right?.
And his grandmother is Eunice..
So in that sense,.
Timothy is the third generation Christian, right?.
Paul is asking Timothy in this difficult time,.
"Look to those whom you can learn,".
including Paul, the mentor himself,.
and also his mother and his grandmother..
We call the examples you have seen from them, right?.
In other words,.
"Fix your eyes upon your spiritual mentors,.
"father, spiritual father, so-called spiritual mothers.".
Now, as a parent, this is what I'm trying to do, right?.
You know, I'm a father, and of course, I'm broken,.
and I am weak, I have many shortcomings..
But one thing I wanna do is that,.
one thing I don't do at home is that.
I never push faith down my children's throats..
I just never did it, you know?.
I never did it..
I never preached to them..
I remember just one time I preached to Rachel..
Yeah, she was upset, and she refused to go to sleep..
So I was kinda angry, right?.
So I stormed into her room, right?.
She was about maybe 10 months old, 11 months old..
And I stormed into her room,.

$^{361}$and I preached a five-minute sermon..
I said, "You cannot be selfish, okay?.
"Because we have to go to bed, okay?.
"We have to go to sleep, and you have to go to sleep.".
Something like that, I mentioned the name of Jesus,.
or something like that, right?.
But because I just went out, you know?.
I was in lack of sleep, you know?.
Come on, what do you want, right?.
So that was only one time I really just preached to her,.
right, preached to my child..
But most of the time, I just try to live authentically.
as a man of faith..
And I want to let them know I need grace as well..
I need Christ to strengthen me..
I need the Word of God, right?.
Indeed, I think that's the most important mission.
we have as Christian parents, right?.
To live in front of the eyes of my children..
Let them know we are authentic followers of Christ..
We are not perfect, but we are authentic..
And indeed, as a Christian, we are called.
to live a Godly life so that our life.
can inspire other people..
We can inspire other people,.
the building of the grace of God..
And for ourselves, we must find great examples to follow..
Remember, I asked you earlier, think of faces,.
think of names in your life..
God, put those people in your life..
Maybe a mentor, maybe a pastor, maybe a brother and sisters,.
maybe other followers of Christ..
They have names, they have faces..
And these are great examples that we can follow.
and encourage us to go on in the truth..
And I'm sure in every church that emphasize.
the gospel truth, the biblical teaching,.
you will find good examples to follow, right?.
If you are tentative, if you look around,.
if you observe, you know there has to be.

$^{401}$a few good examples..
Of course, in the church, there are many broken examples..
But if you are attentive to learn and to listen to stories,.
to learn about the lives of people,.
you will find great people and great example to follow..
I told you about internship in the mid-90s in my life..
When I was a young intern in the church,.
I tried to pour out my heart into the community.
and try to serve the community..
And I have met many, many good people, good mentors,.
and godly and wise and well-equipped in the word of God..
And I remember one of them is the Sydney pastor.
at that point..
You know, I was young and as you know me,.
you know, I'm quite rebellious, right?.
I like to think outside the box..
I don't like to do things out of the box..
I don't like policies..
I don't like regulations, you know..
I don't like guidelines..
I'm quite rebellious..
I don't like to follow rules..
And sometimes I create my own rules..
I like to try different things..
But, you know, I meet with this pastor,.
mentor on a regular basis..
And there was not much I remember what he preached..
Whatever he taught or whatever great sermons he has preached,.
you know, I just don't remember..
I remember nothing..
But I remember I observe his life..
How he deals with challenges and oppositions in the church..
How he deal with all those with grace, with prayers,.
with humors, with joy..
I learn so much from him..
And today we're still good friends..
We are lifelong good friends..
Sometimes he call on me and sometimes I call on him..
And he will check on me and, you know, we'll talk, right?.
But one time I bum into him and his wife in Jerusalem..

$^{441}$I was studying in Jerusalem and just in the hotel,.
I didn't know, oh, you are here too..
And we had a great time together.
because he was on a pastoral tour.
or a pastoral Israel trip to learn, right?.
And I was there to learn as well..
And we were in the hotel lobby..
We just talked until 2 a.m..
Mentors like this, you know, we all need these people..
And that's how we learn, right?.
You can ace in all your tests and all your exams,.
but you need the people skills..
You need to learn from good examples, right?.
And in the church, in our Christian growth,.
it is the same way, isn't it, right?.
In order for us to grow spiritually,.
emotionally, holistically, we need other examples..
And Paul is trying to encourage the young mentee, Timothy,.
not just to study scripture,.
but before he talks about studying scripture.
in verses 16 to 17, he talks about three examples,.
himself, Paul himself, and Lois, his mother,.
Eunice, his grandmother..
And there are many, many benefits of learning.
from just real life example, right?.
And we are encouraged to continue in our faith..
So we come to the last paragraph of this wonderful chapter,.
chapter three, verses 16 to 17, but we don't have time,.
but we'll continue next week.
to talk about these two wonderful verses..
And this becomes the verse of the week,.
chapter three, verses 16 to 17, all right?.
Let's recite it together to conclude our service..
All scripture is breathed out by God.
and profitable for teaching, for reproof, for correction,.
and for training in righteousness,.
that the man of God may be complete,.
equipped for every good work, 2 Timothy 3, 16 to 17..
Let's all pray together..
Father, we give thanks for all the great examples.

$^{481}$that we have in our life..
Thank you for placing these people in our life journey.
so that we can learn a thing or two..
In this lifelong faith journey of following you,.
we once again ask you to strengthen our faith.
so that we can take root, deep root in the Word,.
in the truth of the Bible..
Help us to meditate on your Word,.
to study your Word, to dig deeper and deeper.
so that our lives will be on the firm foundation.
that will never move..
Help us to be wise, wise in your Word, wise in our life..
Again, thank you for all the teachers,.
all the great examples that we can follow in our lives..
Help us to meditate upon them as well,.
how they deal with faith,.
how they deal with difficult times,.
how they deal with challenges and crises,.
how they learn wisdom, how they impart wisdom on us.
so that we will continue to follow you..
Thank you for calling your people.
to be part of your kingdom, to be your sons and daughters..
Thank you for your grace in Christ and pray..
All right, God bless you..
I know it's raining really, really hard out there, right?.
It was, no?.
All right, it's a difficult day to come here..
And thank you for coming..
God bless you and see you next week..
\newpage



\section{}
\label{sec:QovpodRHlhw}
\textbf{流散與福音廣傳 - 劉漢杰醫生}
\newline
\newline
連結: \href{https://youtube.com/watch?v=QovpodRHlhw}{\texttt{https://youtube.com/watch?v=QovpodRHlhw}} ~~~~ 語音日期: 2024-07-21
\newline
\newline
\hyperref[sec:QOZ3Ws7CSlE]{\small{< < < PREV SERMON < < <}}
~
\hyperref[sec:index_chronic]{\small{[返順時目]}}
~
\hyperref[sec:index_scriptual]{\small{[返順卷目]}}
~
\hyperref[sec:umenrD4C3Dw]{\small{> > > NEXT SERMON > > >}}
\newline
\newline
$^{1}$今日我們有很多新的朋友在我們當中.
不如我們花一點時間.
跟一些你不認識的人打個招呼好嗎?.
Hello 你好.
不要第一次見人就摟人家.
可能會嚇到人.
我知道我們教會很熱情.
1991年中國大陸華東發生大水災.
香港人和內地同胞血脈相連.
土地相連.
香港五位牧師和神父.
發起血濃於水的運動.
迅速籌集金錢和物資.
透過香港的紅十字會進行救援.
有感苦難如深受.
天涯若比鄰.
天鄰基金在這樣的情況下誕生.
劉漢傑醫生是天鄰基金會的主席.
是一位有三十多年經驗的兒科骨科手術醫生.
他已經退休了.
但他繼續退休後.
去很多不同的地方世界各地進行骨科手術.
去服侍一些弱勢社群.
今年九月我和另外一些醫護人員.
可以再次跟劉醫生去伊拉克進行醫療斷酸.
我們很高興.
他昨天和他太太Shirley.
坐在哪裡.
她坐在我們後面.
他們昨晚凌晨時分來到Chicago.
睡覺時一大早來到我們當中.
藉著他們的事工.
今天和我們分享天鄰基金的事工.
大家早安.
今早有機會和大家分享一些訊息.
我們開始的時候先來一個禱告.
慈悲全能全愛的主.
能夠崇拜.
在今早可以思想一下關於難民的一個問題.
求主你親自向我們每一個人說話.

$^{41}$讓我們明白主你自己的心意.
也知道怎樣去回應你的呼召.
我們恭敬的將今天講道的時間交在主你的手裡.
求主掌管使用.
以上不完全的禱告是奉主耶穌基督得性的名.
阿們.
今天我們講的題目是.
流散與福音廣傳.
首先我講一下我自己.
我是土生土長的香港人.
我成長的過程裡面生活都很舒適.
當時在香港最大的危機就是升學問題.
因為怕成績不好不能升大學.
那個時候對世界的認識不是很多.
更加對難民的認識是少之又少.
我不知道大家有沒有看過.
在1982年有一齣戲叫做《投奔怒海》.
這齣戲在1982年是全港電影獲得一個獎項.
是最佳電影.
其實獲得了五個獎的連續.
這齣戲的內容是講一個日本的記者.
去到越南.
那個時候南越已經淪陷了.
即是共產黨已經控制了整個越南.
他在政府安排之下去訪問和參觀一些當時的居民.
發現他們的生活非常之幸福.
但這個只是表面.
但當他再深入去查究的時候.
發現原來很多越南的人.
他們是冒死逃離越南.
所以就有這個《投奔怒海》的情況.
當時我不知道你還記不記得.
如果你從香港來.
那時候有一些越南的難民.
他們是坐著一艘漁船離開越南.
有些是來到香港.
香港我還記得設了一些專為越南的船民而設的難民營.
你想一想.
海是風平浪靜.
當然很好.

$^{81}$但當風浪大作的時候.
其實是很恐怖.
我以前試過一次.
由香港駕駛帆船.
經過南中國海去菲律賓.
我駕駛的那艘船.
還有很多人.
不是我自己一個人去駕駛.
就是「智峰號」.
不知道大家有沒有聽過.
當時是經過五天的時間.
由香港經過南中國海.
去到菲律賓北面的Hundred Islands.
我記得那時候曾經風浪大作.
我覺得那艘船.
我們那艘船有130多艘.
但在海上真的很渺小.
不要說人.
那艘船也很渺小.
當大浪打過來的時候.
只是那艘船的主桿在水上面.
四十多呎浪.
即是四五層樓的浪打過來.
那時候覺得海是非常之可怕.
有什麼令到這些越南的難民.
要離開他的家園.
去投奔陸海呢.
他們是願意將他的生命壓上島上.
即是說他們一定是受到很大的逼迫.
才會這樣做.
第二次我比較親身的.
接觸到所謂難民.
就是八九六四的時候.
剛剛六四那晚.
我是由美國回香港探親.
所以六四那晚.
整晚我都在香港看著電視機.
我想那時候如果你在香港.
你都會知道.
香港很多人很心急.

$^{121}$本來看著一個很好的.
一個學生爭取民主.
反對官倒的一個運動.
突然之間被鎮壓下來.
在畫面上聽見有槍聲.
聽見有火光.
聽見有坦克.
接著沒多久.
第二天第三天發覺.
一個白色恐怖在大追捕.
接著聽見有一些.
當時的學生領袖要逃離大陸.
後來我也在美國.
親身遇到一些學生領袖.
亦跟他交了朋友.
因為我們在洛杉磯有些教會.
招呼了一些學生.
這些人順眼之間.
本來是一個在大陸裡面.
他們是一些尖子.
在北京讀高等學府的學生.
轉眼之間變成一些政治難民.
要逃離中國.
他們要逃命.
到底甚麼叫做難民呢.
難民的定義.
他是說一個人身處在自己.
難民是指一個人身處在自己國家以外.
在國家以外.
但他為何會身處在國家以外.
就是因為宗教或者政治.
或者一些武裝衝突的緣故.
而害怕他自己受到生命.
或者身體或者自由方面.
受到傷害受到逼迫.
而需要要求國際提供保護.
即是說這些人不是因為經濟原因.
因為經濟原因而走去其他國家.
那些不是難民.
現在說的是他們受到逼迫.

$^{161}$而這個逼迫令到他們的生命身體和自由.
是會受到傷害的.
而他沒有辦法留在他國家受到保護.
所以他要離開這個國家.
尋求國際的保護.
這些就叫做難民.
還有另一個名詞叫做IDP.
Internally Displaced People.
中文翻譯為國內流離失所.
其實這些國內流離失所者IDP.
情況和這些難民一樣.
但他還未能逃離這個國家.
他可能在這個國家裡面.
由一個地方逃到去第二個地方避難.
當他們離開國家的時候.
定義就將由IDP變成難民.
很多時候IDP的情況比難民更加危險.
因為他們所受的逼迫還可能在附近.
聖經裡面其實原來有很多很多難民的故事.
我想和大家回去思想一下.
第一個的聖誕.
當我們提到聖誕節的時候.
通常我們會覺得是一個普天同慶的節日.
一個歡樂的節日.
因為我們的救世主降生.
是一個歡樂的場景.
但如果你回想第一個的聖誕節.
第一個聖誕不是這樣的情景.
第一個聖誕其實是一個危機四伏.
刀光劍影是一個很危險的情景.
我們看看聖經說什麼.
在這個馬太福音裡面有這樣的說話.
有主的使者向約瑟夢中顯現.
他說.
起來帶著小孩子和他的母親.
逃往埃及住在那裡.
等我吩咐你.
因為希律必尋找小孩子.
要除滅他.
約瑟就起來.

$^{201}$夜間帶著小孩子和他母親往埃及去.
住在那裡.
直到希律死了.
這是要應驗主席先知所說的話.
我從埃及召出我的兒子來.
你想一想.
當時晚上.
約瑟得到一個夢.
得到天使的報夢.
天使要他逃離他們現在所住的地方.
即是伯利咸.
要去埃及.
因為當時的王要追殺他的兒子耶穌.
他要考慮.
那時約瑟半夜三更.
馬利亞馬上叫醒.
馬利亞抱著小孩子耶穌.
可能他們有一隻雷.
可能坐著一隻雷.
然後就連夜帶著一些細軟離開.
要去埃及.
你想想當時是很倉促的.
他們在夜間走路.
當時相信不會有甚麼街燈.
可能是很黑的.
他們就靠著月光.
一些星的照明來看路.
來走路.
趕著離開埃及.
你想想當時的情境其實是很危險的.
因為當時的王要殺耶穌.
當我再看這個故事的時候.
我即時有一個反應.
就想起這幅相.
這幅相是在2014年的時候.
一些主流媒體刊登出來的.
可能是BBC或者CNN那類.
這兩個背景是甚麼呢.
就是一些阿茲迪的難民.
倉促地逃離他們的家園.

$^{241}$當時原來發生甚麼事呢.
就是這些所謂阿茲迪的難民.
就是住在伊拉克的北部.
他們住的家鄉叫做Sindia.
就很近敘利亞那邊.
而敘利亞當時興起了一個伊斯蘭國.
這些伊斯蘭國的人很憎惡這些非回教徒.
原來這些阿茲迪的人.
他們的信仰不是回教徒.
他們有自己的宗教.
是一個很古老的宗教.
甚至他們說是在亞伯拉罕之前的.
他們的宗教是說甚麼呢.
是說神創造了天地之後.
將它交給七個天使去管理.
這個天使長叫做Peacock Angel.
曾經犯罪被打落地獄.
後來亦都得到伏職.
可以從地獄釋放出來.
這個故事當被這些回教徒聽了之後.
他們就會覺得這個所謂Peacock Angel.
其實就是Lucifer.
就是這個魔鬼.
所以他們叫這些阿茲迪.
Yasidi這些人叫做Devil Worshipper.
他們很憎恨這些人.
在歷史上.
這些阿茲迪的人.
曾經有73次被種族清洗.
他們現在的人口比以前少.
當時2014年的時候.
ISIS就是由Syria打過來伊拉克.
他們在晚上圍住他們的城.
要攻進去要打他們.
我還記得當時的一些人.
因為我們去伊拉克服侍.
聽到一些人親身跟我說.
他們逃難經過.
他們說晚上突然聽到有槍聲.
有很多嘈吵的聲音.

$^{281}$原來發覺他們的鄉村.
已經被ISIS圍住.
要攻進來.
他們匆忙地逃跑.
他們跟我講這個故事的人.
他們剛好家裡有輛車.
上了車開車往山上走.
在沿途見到很多人已經受了傷.
他們都沒有辦法去幫他們.
而他們知道的.
他們的親鄰舍和親友.
如果沒有辦法有車.
有交通工具走.
很多都走不及就會被殺.
我們再看聖經裡面.
難民的一個故事.
看回Holy Family的情況.
當時的環境.
當時的王希律.
因為他找不到耶穌基督.
所以他下令將兩歲以內的男童殺盡.
他們耶穌一家.
在耶穌出世沒多久.
他就已經變成難民.
要離開他的本國.
去埃及避難.
但這個故事.
在埃及和北非的一些基督徒.
他們認為他們的祖國.
即是埃及.
曾經為耶穌一家提供庇護.
是津津樂道的.
他們甚至有一些歌曲是說這件事.
他們是引以為榮.
因為他們能夠在耶穌受難的時候.
能夠為他一家提供庇護.
這首歌是北非一些基督徒的歌曲.
我們將它翻譯成英文.
基本上是能夠為嬰兒的耶穌一家提供保護.
是感到光榮.

$^{321}$聖經裡面說過.
我們要用愛心去接待一些客女.
希伯來書這樣說.
不可忘記用愛心接待客女.
因為曾經有接待客女的.
不知不覺就接待了天使.
當時在埃及的人.
他不知不覺原來他已經接待了救世主.
其實我們回想在教會裡面.
是有很多被逼迫的歷史.
亦因為這些被迫.
很多基督徒瞬間變為難民.
要逃離他們的本國.
在使徒行傳裡提到.
在初期的教會.
在斯提反被殺之後.
耶路撒冷的教會大遭迫迫.
除了使徒以外.
門徒部分都分散在猶太和撒瑪利亞的各處.
即是在斯提反被殺之後.
教會裡面很多人都感到很危險.
他們要離開本國本地去避難.
那些因斯提反事遭患難四散的門徒.
直走到菲利基和居比路並安提哥.
主與他們同在.
信而歸主的人就很多了.
這裡發生一件事.
當這些基督徒要離開本國逃難的時候.
他們去到其他地方.
他們同一時間亦將福音帶了過去.
神有他的心意.
他們容許這些離散的事情發生.
可能他們的心意就是要這些基督徒利用這個機會.
將福音講傳.
我們看看當時的信徒.
由耶路撒冷走到地中海的一帶.
然後去到歐洲.
而歐洲後來亦將基督教成為國教.
天主教和基督教都會成為國教.
而這些歐洲國家後來都抗衝他們的地盤.

$^{361}$發展很多殖民地.
這些殖民地就去到亞洲和非洲.
而福音亦因為這個緣故傳到亞洲和非洲.
戴德生因為當時英國有很多人來中國競商.
而他亦拿到這個機會來中國傳福音.
通常這些福音首先傳到一個海外的地方.
無論是非洲,印度或中國.
都是一般在海外的地方.
戴德生特別的意象就是他要做內地會.
就是由海外海邊的城市.
例如上海這些發展的城市.
廣州這些發展的城市.
將福音帶到當時中國的內地.
英國成為基督教國教.
但後來亦有一些基督徒在英國受到很大的逼迫.
一些清教徒.
而亦因為這些逼迫.
他們要離開他們的本國.
後來他們來到北美洲.
亦同一時間將福音帶到北美洲.
現在在全世界派出最多宣教士的國家就是美國.
第二就是南韓.
為何美國有這麼多宣教士呢.
亦因為美國起初時是一個很基督化的國家.
當然現在有很多不同的情形發生.
但這個福音的廣傳.
就是通過很多基督徒.
因為他們受到逼迫.
而要離開本地,本國.
去到新的地方.
同時亦帶福音去到這個地方.
在中國的基督徒.
初期都是集中在沿岸的地方.
就算內地會.
他們的大本營都是在河南.
如果你看中國的地圖.
河南其實在中國整個地圖都是偏東面.
但在第二次世界大戰和中日戰爭的時候.
很多基督徒因為戰亂的關係.
他們要逃難.

$^{401}$變成我們所說的IDP.
他們一般向西遷移.
同一時間將這個福音帶到中國的西部.
我不知道大家有沒有留意到.
我們現在就是今天.
我們正面對第二次世界大戰以後.
最大的人道危機.
這個危機就是因為在非洲.
在中東和烏克蘭所發生的戰亂.
造成很多難民.
根據聯合國難民公署的報告.
這個是2023年.
每一次我講類似的題目.
我都要再看新的報告.
每一次的人數都增加.
根據聯合國難民公署2023年的報告.
現在被迫離開家園的人.
被迫離開家園的人.
是有1億1千萬.
其中有4千多萬是難民.
其他那些是所謂的Internally Displaced People.
是國內流離失所.
這個圖其實是在聯合國難民公署那裡.
我複製下來的.
你可以看到117個million.
即是1億1千7百萬的人.
被迫離開家園.
難民的人數在過去10年增加了3倍.
難民大部分來自幾個國家.
最多的是敘利亞.
接著是阿富汗,烏克蘭.
還有很多其他國家.
在南蘇丹已經持續了很久的內戰.
所以他們都有很多的難民.
而有很多的難民是小孩子.
甚至有40\%是小孩子.
估計起碼有100萬個嬰兒.
他出生的時候已經是一個難民身份.
你覺得是不是很慘.
即是出生已經是一個難民的身份.

$^{441}$而大部分這些難民.
是去他們本國附近的地方.
很多都不是走得很遠.
可能是隔離的國家.
而收容他們的國家裡面.
很多都不是一些富裕的國家.
原來收容得最多難民的國家是土耳其.
發達的國家.
譬如歐洲國家和美國.
其實相對其他這些國家的人數不算很多.
歐洲國家裡面最多的就是德國.
大家不知道有沒有印象.
這一幅圖片.
這幅圖片就是幾年前.
一個攝影記者在地中海海邊拍到的.
就是一個男孩子.
這個男孩子是一個三歲的.
由敘利亞來的一個難民.
他們一家本來是打算坐船.
經過地中海去投奔自由.
他們就是剛才出戲的.
投奔陸海的這些人.
但是不幸他們遇到海難.
這個小孩子就死了.
當時這幅圖引起很大的迴響.
原來這個家庭已經申請去加拿大.
作為難民.
即是申請去加拿大移居很久.
但都得不到批准.
當時在加拿大都有一個很大的反響.
每一個難民的家庭其實都有一個新鮮的故事.
而大部分的難民.
他們是有一個宗教背景的.
他們很多都相信有神.
而在近年來的很多難民.
都是一些回教徒.
而近年來有一些報道.
說這些流離失所的回教徒.
他們是有異象和異夢的現象.
當我想到異象和異夢.

$^{481}$通常都是聖經裡面的故事.
想不到是好像現實生活有發生一樣.
好像這些只是一些以往的事.
我也沒有親身見過.
在我認識的朋友裡面.
有親自告訴我他有一個異象或異夢.
但這本書告訴我.
這本是幾年前出版的書.
他說原來在中東的難民.
在穆斯林的難民裡面.
很多是有異象和異夢的情況出現.
就是耶穌基督或天使.
跟他們通過異夢和異象溝通.
近年來原來有大批回教徒信主.
在過去幾十年.
在歐洲的教會有一個衰落的現象.
雖然有很輝煌的教堂.
但聚會人數不斷減少.
而信徒流失的現象很嚴重.
回眾的老化情況也很嚴重.
有些教堂甚至倒閉.
教會在社會上的影響力越來越少.
所以早幾年.
當黃永信牧師大使館的會長.
他在世時還提出一個口號.
叫Re-Evangelization of Europe.
即是向歐洲再次傳福音.
因為當地的教會已經衰落.
但近年來很多難民離開本國後.
很多時都會來到歐洲的地方.
而面對很多難民.
歐洲的教會不能不作出回應.
雖然歐洲的基督徒.
可能對難民採取冷漠的態度.
但神曾經興起了很多不同的士工.
去照顧這些難民.
而因為這些照顧.
近年來有一個有關流散的宣教學.
叫Diaspora Missiology.
講述這些流散的人.

$^{521}$與教會的一些互動.
為了照顧這些眾多的難民.
歐洲的教會變得積極起來.
甚至有一些復興的現象.
這個現象很奇怪.
早前我提到基督徒被迫.
將教會的福音傳了出去.
這次是非基督徒受迫.
來到基督徒的地方.
激發起教會.
令到教會有一個復興的現象.
最近有幾本書出版.
內容都類似.
提到在離散的回教徒裡面.
有一個復興的現象.
提到他們有很多人見到異象.
他們的復興很多時候是在婦女裡面.
你知道如果一個媽媽信了主.
她對她的兒女有很大的影響.
這些書所講的內容.
都是講到近年來.
在這些地方的回教復興的現象.
這本書A Wind in the House of Islam.
提到以往在回教世界的歸主運動.
你提一提.
其實在2000年之前.
是很少歸主運動.
而這個歸主運動的定義是.
在20年之內建立100間教會.
或者有1000個回教徒受洗.
這叫做歸主運動.
其實數目不是很多.
但是在2000年之後.
竟然有69個歸主運動.
在中東的地方出現.
大家可不可以看到.
神正在做這些工作.
這些事情不是無端端的發生.
我們看到神正在攪動這些回教徒.
將他們由本地驅散了出來.

$^{561}$而同時也令他們有接受福音的機會.
這本書叫做Experiencing God.
是我一本很喜歡用的靈修書.
是一本很暢銷的靈修書.
它的作者是Pastor Henry Breckbe.
他曾經說到.
你要看看神在哪裡工作.
你就去加入他的工作.
神我覺得現在是在這些難民裡面工作.
我們是不是應該也一起去參與這個工作呢.
我們看到當今難民是最大的人道危機.
宣教主使用這個危機.
化為一個雙贏的契機.
第一,他將回教徒,異教徒.
從一個根深蒂固的文化環境帶了出來.
讓他們有接觸福音的機會.
第二,他們也利用這個機會.
激發一些僵化的死寂的教會.
令到它有一個復興的現象.
在難民方面的工作.
可以分開三方面來看.
第一,就是關於他們生活的安置.
我們可以幫助難民的家庭.
例如給他們教育的機會,學習的機會和醫療的機會.
我們天倫基恩惠所做的就是做醫療的工作.
第二方面是關心他們心靈方面.
例如心靈的輔導,宗教的輔導.
以及將基督的愛帶給他們.
第三,做一些訓練的工作.
就是培訓他們做職業的工作.
這個難民工作是每一個基督徒都可以參與的.
你可以為他們禱告,奉獻.
也可以為一些已經來到你們附近社區的難民提供一些服務.
我希望大家可以在這裡思考一下.
我們到底有沒有辦法可以幫助這些人.
這些人是好像我們中國人所說的天涯淪落人.
我們願不願意做一個好撒瑪利亞人去幫助他們.
我們的教會到底是一個百里行.
當時沒有地方去接待這個聖家.
還是埃及給地方這個聖家去避難呢.

$^{601}$我們可以思考一下.
多謝大家.
我們都去禱告.
我們知道在芝加哥也有很多難民.
就在我們門口.
到底我們作為基督徒是怎樣用我們的信仰去體驗,照顧他們.
或者為他們禱告.
我們很多時候爭論的只是政治的東西.
但是基督的信仰給我們什麼力量和資源.
我們能夠去幫助這些人.
還有遠地的.
他們在克格斯坦工作.
上一次去參加他們的9月會.
我們遠處近處都有很多流離失所的人.
包括我們自己也是.
但主給我們在流離失所當中.
看到其實我們主耶穌也是一樣.
並且在我們的經歷裡面.
我們能夠看到福音.
我們同心得道祈禱.
至於我們今天來敬拜你.
沒錯,你就是我們宇宙裡面.
生命裡面最大的主宰.
在我們生命當中.
在歷史裡面.
每一個細節你都在你的手中掌管.
雖然在我們生活裡面.
我們生命當中.
我們很多時候覺得會迷失.
會不知道為什麼.
但我深信你是那位滿有主權的神.
世事的發生.
在我們生命裡面不同的遭遇.
主也叫我們能夠來到你面前.
去敬拜去事奉你.
主今天我們思想到一個很切身的問題.
這些難民流離失所流放的人.
在我們近處.
在我們門口.
在遠處.

$^{641}$世界各地不同的地方.
主你在這一代.
此時此地.
我們可以做什麼.
求上帝你親自興起我們.
叫我們教會.
我們在此時此地懂得怎樣去回應.
回應你在歷史裡面所做的事.
在歷史當中你所安排.
所策劃的一切.
主你叫我們知道我們怎樣去回應你.
本著基督的愛.
本著十字架的救恩.
我們將你自己的福音.
將你自己的平安.
能夠帶到近處遠處有需要的人.
主在我們身上.
我們都可能會遇見很多迷失的時候.
今天來到你面前.
求你讓我們知道.
在基督裡面.
我們有一個世界不能夠明白.
世人不能夠賜予給我們的平安.
我們感恩賺給你.
主讓這個平安能夠在我們生命當中落實.
也讓這個平安能夠在我們生命當中分享出去.
我們可能可以用一些資源.
行動.
甚至是禱告.
能夠幫助人認識你.
認識主的裡面.
基督的裡面.
能夠有出人意外的平安.
我們再一次慰問劉醫生和他們的隊工.
他們很忙碌的在不同的地方.
也在預備宣教的外展的工作.
主我們求你親自與他們同在.
去使用他們.
保守他們一切所需要的.
用你的信實來供應他們.

$^{681}$興起更多的人來參與他們.
去參與你此時此地在做的工作.
主要叫我們能夠憑著信.
憑著我們的行動.
來將福音帶到更遠的地方.
我們祈禱奉傳明禁.
好 完上帝賜福給各位.
希望下星期能夠和大家一起到時逢.
\newpage



\section{3:16-17}
\label{sec:umenrD4C3Dw}
\textbf{Back to Basics \#4 (2 Timothy 3\_16-17) - Rev. Mark Chan}
\newline
\newline
連結: \href{https://youtube.com/watch?v=umenrD4C3Dw}{\texttt{https://youtube.com/watch?v=umenrD4C3Dw}} ~~~~ 語音日期: 2024-07-23
\newline
\newline
\hyperref[sec:QovpodRHlhw]{\small{< < < PREV SERMON < < <}}
~
\hyperref[sec:index_chronic]{\small{[返順時目]}}
~
\hyperref[sec:index_scriptual]{\small{[返順卷目]}}
~
\hyperref[sec:koqOwadJaEc]{\small{> > > NEXT SERMON > > >}}
\newline
\newline
$^{1}$- Good morning everyone..
Yeah, thank you Veronica for the clear announcement.
with different languages..
Just like the song that we have just sang,.
from every tribe and every nation, right?.
You know, that is Christ's mind..
You know, Christ is going to redeem every people.
and every tribe and every language to worship him, right?.
He is the Lord of all of us, okay?.
So very good to see you, okay?.
Very good to see you..
(speaking in foreign language).
Salaam, brother..
Yeah, all right..
So I'm speaking three different languages..
Yeah, that's all I can do..
(laughing).
Yeah, but that is the way to worship God.
because we use different languages,.
we use different dialects to express our praises to God,.
right?.
And our God is worthy of our praise..
And we continue our series on back to basics..
And there are basics in our Christian life,.
basics in our Christian walk, and we have to master..
We have to understand, right?.
So often we don't find it enjoyable in our Christian walk.
or in our Christian journey, we have troubles.
because we have forgotten the basics..
And these are basics that we need to focus on..
And that's why in this series,.
we're going to go back to some of the basic understanding,.
basic knowledge, basic teaching of the Bible..
And I've introduced something new to the congregation..
It is the verse of the week, right?.
Remember this verse last week?.
Remember this?.
Verse of the week..
And let's recite together, all right?.
I, that's part of your homework is to memorize it, right?.

$^{41}$Have you?.
I'm sure you have..
Verse of the week, let's recite it..
You must love the Lord your God with all your heart,.
all your soul, and all your mind..
And this is the first and greatest commandment, right?.
That's what Jesus said to us..
This is the first and greatest commandment.
of all the commandments in the Old Testament, in the Bible..
There are many details, there are many different.
commandments and statues you can find in the Bible..
But this is the most important thing..
This is the ultimate summary, right?.
The ultimate summary of what God desire us to do.
is to love Him, to be loyal to Him, to follow Him.
with all our heart and soul and our mind..
And last week, we were talking about the importance.
of the Word of God in our Christian walk, right?.
Our Christian walk, we cannot do so.
without the Word of God, without a correct understanding,.
without a solid teaching of sound doctrine..
And that's why God has given us the scripture..
And this is the verse of the week for this week,.
for this current week..
And it's taken from the chapters that we are studying on..
It's 2 Timothy 3, the last two verses..
16 and 17, last two verses..
Let's recite it together..
All scripture is breathed out by God.
and profitable for teaching, for reproof, for correction,.
for training in righteousness,.
that the man of God may be complete,.
equipped for every good work..
2 Timothy 3, 16 to 17..
All right, so we got to have this verse,.
or these verses memorized for this week, okay?.
So whenever you are doing this or doing your work.
or there is some break time in your mind,.
try to recite this..
I give you the Bible reference, 2 Timothy 3,.

$^{81}$and you can find it and recite it..
All scripture is the breath of God, you know?.
And we're gonna study this chapter, chapter three..
And remember last time I said this long chapter.
of 17 verses can be divided into two parts..
And the first part is verses one to nine..
Second part is 10 to 17..
And the first part from verses one to nine.
is the dangers of neglecting God's word..
God's own word, if we neglect or we ignore,.
there are dangers, right?.
Remember Paul is trying to warn Timothy..
And again, Paul is the mentor pastor of Timothy..
Timothy is his disciple..
Timothy is kind of his mentee..
And Timothy is being in charge of the churches in Ephesus..
And Paul is writing this letter in very personal tone..
That's why we call it pastoral letters, right?.
First and second Timothy in the book of Titus..
He's writing to his mentees..
Timothy is one of them..
And trying to teach them, trying to encourage him.
to persevere in shepherding the church of God..
There were many opposition, there were false teachers,.
there were many dangers in the church,.
in the small churches in Ephesus..
And the first nine verses, it says there were dangers.
warning Timothy that this is the last days..
Remember what I said about last day in New Testament term?.
The last day is the period between the first coming.
of Christ and the second coming of Christ,.
the final coming of Christ,.
which is still lying in the future, right?.
First coming of Christ, we know what it is..
Incarnation of Christ, Jesus came in human form,.
fully God and fully man..
And he died for us, died for our sins..
And on the cross, he was crucified..
And on the third day, he rose again.
and he was ascended to heaven..

$^{121}$And he's king in heaven..
And that is the first coming..
And between first and second coming, it is the last days,.
the last period of human history.
according to biblical timeline..
And Paul is saying to Timothy, this is a dangerous time.
because there are many difficulties.
and people are drawn to falsehood, right?.
I remember a few years ago when we had the pandemic,.
when the pandemic first hit us, right?.
Remember on social media, there's so much misinformation,.
right?.
There's so much conspiracy theories.
or just feeding our uncertainty, feeding our fear..
And sometimes people are being hijacked.
by fear and uncertainty..
And because the time was difficult, right?.
And we were in fear..
And whatever information feed into our brain, our mind,.
especially sensational materials,.
they become like wildfire in the social media..
And I just wonder, shaking my head at times,.
why people would come to believe in falsehood, right?.
Something is clearly wrong, right?.
And they just never double check, right?.
And, you know, scammers, right?.
It's all over the place, right?.
Yeah..
I don't know how many scam calls you get during the week..
How many calls you get during the week?.
Or text messages?.
I can't count, you know?.
I can't count..
You know, every day I get something like,.
(speaking in foreign language).
(laughing).
Something like that..
In Chinese sometimes, sometimes in English, right?.
Hi, you know, I'm your old friend, you know..
I know some of your friend..

$^{161}$Or I even claim that they know my dad, you know?.
My dad passed away many years ago..
(laughing).
But, you know, there are people who believe in scams, right?.
Not too long ago, I received a phone call from a young girl..
I received a phone call from a young girl..
And she was in the mid-twenties..
And she called me because she was in trouble..
She received a call from IRS, you know,.
our tax department, right?.
US tax department..
It's very scary..
Called her and said,.
"Hey, you did not file your taxes, right?.
"For so many years, right?.
"You've been working, you did not file your taxes..
"And now you have to pay back, all right?.
"And you have to pay back like \$5,000..
"But this is how you, you know, deal with IRS, right?".
There's a scammer told her,.
"Okay, you go to the Apple store, all right?.
"You go to Apple store and buy Apple gift cards.
"of 5,000 worth, all right?.
"And you're gonna give me the gift card number.
"and the code or whatever identification..
"And you give it to me, you know, through email online..
"And then, you know, we're gonna cancel out your taxes..
"All right?".
And she believed it..
She went to Apple store and got all the gift cards.
and 5,000 worth and gave all the numbers away..
It was crazy..
And that's why she called me..
She said, "My goodness, I'm in trouble..
"Now I realize, you know, it was a scam, right?".
So miserable, right?.
There's nothing I can do, right?.
There's nothing you can do, right?.
But it just blow my mind, you know,.
how she could have believed something like that, right?.

$^{201}$You know why?.
It's not because she was stupid, right?.
She's very smart..
It was because she was in fear, right?.
Indeed, she did not file taxes..
Indeed, for a couple of years, she did not,.
she forgot or she just ignore it, right?.
And she knew that she owed some money to IRS, right?.
And that's how the scammers got to her, right?.
Because she was in a state of uncertainty and fear.
and all sorts of misleading content.
that she would believe, right?.
And this is the last stage, right?.
And we continue to read this passage..
We find that Paul is trying to urge Timothy.
to learn the truth and learn the word of God..
And that's why the second part, he says,.
there's so much benefit of learning the word of God,.
learning the truth, trying to discern, right?.
And we go through this part of this passage already.
last week from verses 10 to 15..
And now we come to this two verses,.
the verse of the week, verse of this week..
And I tell you, in the entire Bible,.
these two verses, 2 Timothy 3, 16 to 17,.
are some of the most important verses in the Bible..
These two verses are some of the most important verses.
in scripture..
And let's unpack it and we'll see why..
It says, "All scripture is the breath of God.".
God breathed..
And it's useful for teaching, rebuking,.
correcting, and training in righteousness,.
so that the man of God may be thoroughly equipped.
for every good work..
Now, these verses are very, very valuable, very important..
And why?.
Now, this is the last paragraph of chapter three..
And Paul is trying to encourage Timothy,.
yeah, learn from your mentors, right?.

$^{241}$Remember, Paul is one of his mentors..
And who else?.
His mother and Timothy's grandmother, right?.
Eunice and Lois..
Learn from them, because since you were a child,.
you were learning scripture from these people, right?.
And here in verse 16, he says,.
"All scripture is the breath of God.".
Well, let's unpack it..
Our first, the key word here is the breath of God..
God breathed out..
He spoke forth His word..
And God is the speaking God..
God spoke and God is speaking today..
And God will be speaking tomorrow..
And God is the God who speaks..
The idols don't speak..
They may have mouth, they may have lips,.
but they do not speak..
Now, you don't have to be thoroughly learned in the Bible.
and you get the point, right?.
Because the first page of the Bible,.
when you turn to the Bible, the very first page,.
the very first chapter, very first passage,.
is Genesis one, right?.
In the beginning, God created the heaven and earth, right?.
And then there were some other details,.
the spirit of God upon the water..
And what did God do?.
What was the first thing God did?.
He said, right?.
And God said, what did He say?.
Let there be light, right?.
Let there be light..
And then there was light..
So you got that from the first passage of the Bible,.
that God has the attribute of speaking..
God is someone who reveals..
God is someone who tells you what to do..
God is someone who can speak..

$^{281}$And here it says all scripture is the breath of God..
It's the very breath of God..
It tells you that God himself.
is the ultimate author of the Bible..
You have your Bible on your phone app.
or you have a physical print Bible,.
but here it declares that all scripture,.
the Bible is actually written by God..
Of course, it's written by God through people, right?.
Through human authors, like Paul himself, right?.
Like Luke, who wrote the Gospel of Luke..
And Mark, the Gospel of Mark..
And some of the letters,.
or Moses wrote some of the books, right?.
But the ultimate author is God himself..
And it tells you the authority of scripture,.
the standard of scripture..
You know, we are an evangelical church,.
and we are in the great tradition of reformation..
More than 500 years ago, 1517..
1517..
That was the year that launched the movement of reformation..
And the guy who launched it,.
one of the guys who reformed the faith at that time.
was Martin Luther, right?.
Not Martin Luther King, Jr., okay?.
That's a different guy, right?.
Martin Luther, 1517..
Martin Luther, okay?.
He's a church reformer,.
because he saw the Catholic church was broken,.
and they were not coming under scripture..
So the entire movement of reformation is to call people,.
call the churches to return to the word of God..
And the scripture is the final authority.
of everything that we believe.
and how we practice the Christian faith..
We always have to return to the word,.
return to scripture..
And Paul is saying all scripture is the standard,.

$^{321}$because God has written it,.
and he has divine authority..
It says all scripture, all scripture, singular, singular..
Now, in the time of Paul, they only had the Old Testament..
Do you get my point?.
The New Testament was still in the making, right?.
Some of the letters are being written..
By this time when Paul wrote 2 Timothy,.
some of the letters from the Old New Testament.
or even the gospels were already being circulated,.
circulated in different churches..
But the entire New Testament, New Testament canon,.
all 27 books are still in the making..
And the Christians in Paul's days in the first century,.
they only had the Old Testament,.
which was originally written in Hebrew..
And a couple of hundred years before Jesus,.
it was translated into Greek..
And the apostle Paul or the early churches,.
they were using the Old Testament only,.
either in Hebrew or in the translated version Greek,.
we call it Septuagint..
They only had the Old Testament..
Now, listen up..
All scripture, Paul is saying all of the Old Testament,.
right?.
All of the Old Testament, whether it's Hebrew or translated,.
they are as one, the Word of God, right?.
The whole Old Testament is already the very Word of God..
Right?.
Now, you think about it, this is pretty interesting,.
pretty inspiring..
Just the Old Testament, right?.
Just the Old Testament..
Yeah, maybe somehow your house got caught on fire, right?.
And your Bible is being burned, right?.
Let's say the New Testament, all 27 books.
are being burned away, right?.
They've been consumed by fire..
You got out of the fire and you got the Old Testament, right?.

$^{361}$You still have the Word of God, right?.
The first 39 books..
Paul is saying all the Old Testament.
is already the very Word of God..
And think about it, this is very interesting..
For us, we have more than the Old Testament, right?.
Because Jesus came and all the books of New Testament.
are being written about Him coming, right?.
And all 27 books are being added to the,.
what we call canon, what we call the full revelation.
of God in scripture..
And now we have Old Testament and the New Testament..
We have the Old Testament's 39 books.
and we have the New Testament 27 books..
How perfect we are, right?.
We have more than what Paul had, right?.
What Timothy had..
And we know the Word of God,.
we know the will of God so much more..
And we can comprehend the Word of God.
or the mind of God, even with more content, right?.
We are so blessed to have so much.
that God has revealed to all of us,.
the Old and the New Testament..
And Paul is saying all scripture is the breath of God.
and it's useful, it's profitable, right?.
Useful for teaching, for rebuilding, correcting.
and training in righteousness.
so that we can be fully equipped for every good work..
Let's unpack this..
I still like the New Living Translation, NLT..
And I give you here and at the bottom of the screen.
is the contemporary Chinese translation..
I think these more contemporary language.
would have us better understand.
the usefulness of scripture..
There are four phrases..
Paul is saying, if you go to scripture,.
you read the Bible, first thing,.
the Bible will teach you what is true, right?.

$^{401}$(speaking in foreign language).
Right?.
Teach you what is true..
And secondly, is to make you realize.
what is wrong in your life, right?.
I tell you one of the secrets, I think,.
that we don't like to read Bible..
You know why?.
Because of this second purpose..
We don't like to be confronted, right?.
We always want to be comforted, right?.
But you know, you go to scripture,.
so often the scripture confronts us, you know?.
There are behaviors, there are mindset we have to change..
Right?.
And when it comes to changing and repenting,.
we just don't like it, right?.
That's why sometimes we ditch the Bible, right?.
We just don't like to read it..
Or we like to read it because we just want to read it.
selectively, right?.
We're just cherry picking all the time, right?.
Just like to read whatever we like or sound good to us..
However, there are many, many passages in the Bible.
that confronts us, show us that we have sins.
to repent from..
We have shortcomings..
We need to ask God to give us grace..
And that is the value of scripture, right?.
To tell us, yeah, there's something wrong with us, right?.
And that's why we need grace..
And that's why we need Jesus Christ,.
to make us realize this,.
so that our life cannot be the same..
We can be changed and we can be changed.
by the grace of God, right?.
Only if we allow God and the Spirit to work in our lives..
And that's the second purpose..
The third purpose is to correct us..
It's not just showing you what is wrong in your life,.

$^{441}$but it's to tell you when, to correct you,.
to show you the way, the right way,.
when you are wrong, right?.
Yeah..
I've been driving for almost a year,.
over 30 years, yeah..
Over 30 years, yeah..
I consider myself a pretty experienced driver, right?.
Don't you think?.
And I also have some experience in teaching new drivers,.
right?.
You know, my sons and daughters, right?.
I'm trying to teach them how to drive, right?.
And it's such a pain to teach them..
It's so painful to teach a new driver,.
especially a young driver, right?.
I have to show them what is right, right?.
I also have to, and that is pretty easy, okay?.
You have to turn right..
You have to turn left..
You have to make a full stop here, okay?.
They said, "Dad, you don't always do a full stop.".
I said, "Don't worry about me, okay?.
You have to do it, okay?.
Because you're a new student, right?".
But that's pretty easy..
The first thing is pretty easy..
The second thing is pretty hard, right?.
To help her or him to realize.
there's something wrong in his or her driving, right?.
There's some habit you have not developed, right?.
Like checking your blind spot, right?.
Even today, I'm sitting in my son's car..
I still have to remind him, right?.
Every turn you're gonna make,.
you have to always check the blind spot, right?.
Yeah, if you go to my cars, a few cars, a few old cars,.
I always have the, I always buy the blind spot mirrors..
You know what that is?.
I stick it on the side of the side mirror, right?.

$^{481}$Because I've been using it for decades, right?.
And every single car that I own, or before, or now,.
I always have the blind spot mirror, right?.
So I have to remind the new driver,.
"Come on, you have to check your blind spot..
That is the most important thing," right?.
Yeah, and sometimes they get angry..
They get angry at me, you know?.
I'm so kind, right?.
I just want them to drive well, you know?.
But that is the human nature, right?.
We don't like to be pointed out.
there's something wrong in our lives, right?.
But that's what scripture does..
Not only showing us what is wrong in your life,.
but to correct you, to tell you what is the right way,.
the better way, right, to live..
And then lastly, is to teach us what is right..
The original term for the fourth purpose,.
teaching us to do what is right,.
is actually, a better translation may be discipline..
You know?.
You know what discipline is?.
Habit, right?.
Remember, beginning of the year,.
I have a series on the book of Proverbs,.
and the series is called Atomic Wisdom, right?.
Atomic habit, it means small, tiny habits.
that we have to build in our life.
so that we can be godly, we can be trained and equipped.
to walk in the way of God, right?.
And there are tiny atomic habit we have to develop..
Ask questions, you know?.
Now, meeting together is one of them..
Reading the Bible is one of them..
And praying, we're gonna talk about it later,.
some of the habits or disciplines.
that we have to develop, right?.
And every one of us know what discipline looks like, right?.
In the beginning, it was painful, right?.

$^{521}$But as you go on, as you persevere, as you continue,.
you find it is helpful for you, right?.
Right?.
But those of the people who do workouts.
from time to time or regularly,.
you know what I meant, right?.
You know what I meant, right?.
When I first do, try to do workout,.
did I tell you?.
I do plank, right?.
I do plank, I started with, it was a couple years ago..
You know, I find that, you know, I'm getting old.
and I need to do some more exercises, right?.
And, you know, being a pastor,.
sometimes we just sit there and in front of the screen,.
we have to prepare messages.
and we have to answer phone calls and email and all that..
And sometimes we just get lower back pain, right?.
And I started with plank for only 45 seconds..
Dude, it was painful, right?.
And now I can do seven minutes, right?.
Can you do seven minute plank?.
It's pretty good, right?.
For someone like me, right?.
Over 50..
But, you know, you build up your body, right?.
And you become stronger and stronger and stronger, right?.
Until today, if I don't do plank,.
I feel like there's something wrong.
and something missing, right?.
So to get up, I have to do plank again, right?.
And push up and all that stuff..
That's what I'm saying, discipline..
NLT has it, teaching us to do what is right..
Peiyang, xingyi..
Peiyang, cultivate, right?.
It's a habitual pattern, right?.
And we all know the Olympic games.
are gonna start soon, right?.
July 26th, right?.

$^{561}$Olympic games..
I don't know if you watch some of the preparation.
and some of the qualifying games on TVs..
It's pretty interesting to watch..
And pretty exciting and inspiring to watch..
But the athletes, why are they doing what they do, right?.
How could they do something like that, right?.
Because they have gone through tough training, right?.
It's not something that you can do overnight..
They do it day after day, hour after hour,.
minute after minute..
Then that's why they have the ability to do it, right?.
Christians, we need to read scripture.
and we have to understand scripture, right?.
I'm thankful we have at least the men's breakfast, right?.
And we're going through scripture verse by verse.
or chapter by chapter..
I hope this will continue to develop..
They'll develop into a more solid,.
like a Sunday school class, right?.
And a Bible study, so that we can have.
a better understanding of the word of God, right?.
And we're gonna close by this verse of the week..
Oh, by the way, yeah, I said the old scripture,.
2 Timothy was the worst of the week, right?.
This should be the one..
Sorry about that..
Yeah, I'm dating myself..
This is the verse you have to memorize.
for this coming week, all right?.
Verse of the week, all right?.
Matthew chapter five, verse 14..
"You are the light of the world,.
"like a city on a hilltop that cannot be hidden.".
The light of the world..
This is the verse you have to memorize this week, all right?.
Why did I pick this verse?.
You know, in Jewish tradition, Jewish rabbinic literature,.
they discuss what is light, what is light..
And to them, what is light?.

$^{601}$The Torah is the light..
It's the light for life..
The Torah, it meant the scripture, right?.
For the Jews, it was the Old Testament, right?.
The Hebrew Bible..
For us, it's the entire Bible..
And in their discussion among the rabbis,.
the Torah is the light of the world..
The Torah is the light of the world.
because the Torah will give life to people..
The Torah will give grace to people,.
will point people to the right way of wisdom..
And here, Jesus said something really interesting..
He was saying to his disciples, right?.
It's part of the Sermon on the Mount, right?.
Yeah..
Pastor Chi has been preaching on it..
He was talking to his disciples..
"You guys are the light of the world.".
Wait a minute, what does it mean?.
That means you are the Bibles..
You are the Torahs..
You are the lives who will show people.
the beauty of the Word of God, right?.
If you listen to me and obey me, Jesus said,.
you are the light of the world, right?.
You know most people don't read the Bible, right?.
You know most people don't understand the Bible.
and where to start, right?.
But you guys are different, right?.
When they see you, when they see your life,.
they just see the beauty of God's Torah,.
God's law, God's very Word, right?.
Because you're living out the Word of God..
You are the light of the world,.
like a city on a hilltop that cannot be hidden..
You cannot hide..
If you profess the name of Jesus Christ,.
you believe in the name of Jesus,.
you cannot hide, you cannot hide in the dark..

$^{641}$And you are the light, right?.
And as we conclude, I want us to take this challenge..
To take this challenge is to read and meditate scripture,.
right, to read and meditate the scripture..
You are the Bible in some sense..
You are the Word of God..
And you know, in our time, in our age,.
I think this is the most challenging thing that we all face..
It's not so much of preaching the gospel.
or telling other people what they should believe..
That is not the most challenging thing..
It's that the most, the biggest challenge for us.
is that, is how to give people a good, convincing,.
positive impression of what the Christian is like..
You can tell people about the gospel,.
but at the end, you need to leave a good,.
positive impression..
Whether the person believes or not, right?.
Whether the person is convinced.
about the gospel of Jesus Christ or not..
That is not the question..
That is not something we can do, right?.
Only the spirit can do, right?.
But what we are urged to do is to give people.
the positive, convincing impression.
of what believing in Christ is like, right?.
Can we do that?.
To our family members, to our children,.
to our spouse, to our neighbors, right?.
Can we do that, right?.
By what we say and how we behave, right?.
How we interact with other people..
Let them see a good impression..
Let them see the convincing impression, right?.
You are the light of the world..
Let's close..
(smacks lips).
Father, we all come before you..
We confess that we're broken, we are weak..
Sometimes we do ignore the word of God..

$^{681}$But the word of God is right in front of us..
But we neglect it..
We'd rather watch videos, we'd rather scroll our phones,.
and being distracted by all these worldly desires..
We come before you to confess our sin..
So cleanse us and forgive us..
Let us focus, once again, the light that you have given us.
that will shine into our lives,.
all the darkened corners of our heart..
And we need your word..
And may we continue this habit.
to read and to meditate your word..
For your word is good..
Your word is precious and perfect..
Give us the understanding, give us the wisdom,.
so that we can understand scripture,.
understand the very word..
In scripture, we find you..
In scripture, we find you being with you..
And thank you for being our good and faithful God..
May we take your word more seriously..
May we read your word more diligently..
So that we are full of your message,.
full of your word, so that we can be the light.
to this darkened world..
Send us out this week..
So we go to our neighborhood, go to our families,.
to interact with other people, Christians or not,.
that our lives will be the light..
Our lives will be your word..
And we will be your representatives..
Send us out and use us, and help us to dig into your word.
and to help us to understand your very word..
Thank you for being our God who speaks,.
and you still speak..
Speak to us through your spirit, through scripture..
Grant us an obedient heart to listen and to obey.
in Christ and pray..
All right, guys, hope to see you next week,.
and we'll continue to serve our God.

$^{721}$and to understand the word, all right?.
God bless..
\newpage



\section{提摩太後書 3:10-17}
\label{sec:koqOwadJaEc}
\textbf{回到基礎 \#2 (提摩太後書 3\_10-17) - 陳崇基牧師}
\newline
\newline
連結: \href{https://youtube.com/watch?v=koqOwadJaEc}{\texttt{https://youtube.com/watch?v=koqOwadJaEc}} ~~~~ 語音日期: 2024-07-23
\newline
\newline
\hyperref[sec:umenrD4C3Dw]{\small{< < < PREV SERMON < < <}}
~
\hyperref[sec:index_chronic]{\small{[返順時目]}}
~
\hyperref[sec:index_scriptual]{\small{[返順卷目]}}
~
\hyperref[sec:9TgjTCjHTHQ]{\small{> > > NEXT SERMON > > >}}
\newline
\newline
提摩太後書 3:10-17
\newline
\begin{longtable}{cl}
\hline
\hline
章節 & 經文 (和合本修訂版)\\
\hline
3:10 & \begin{tabularx}{0.7\textwidth}{X} 但你已經追隨了我的教導、行為、志向、信心、寬容、愛心、忍耐， \end{tabularx} \\ \\ \relax
3:11 & \begin{tabularx}{0.7\textwidth}{X} 以及我在安提阿、以哥念、路司得所遭遇的迫害和苦難。我忍受了何等的迫害！但從這一切苦難中，主都把我救了出來。 \end{tabularx} \\ \\ \relax
3:12 & \begin{tabularx}{0.7\textwidth}{X} 其實，凡立志在基督耶穌裡敬虔度日的，也都將受迫害。 \end{tabularx} \\ \\ \relax
3:13 & \begin{tabularx}{0.7\textwidth}{X} 只是作惡的和騙人的將變本加厲，迷惑人也被人迷惑。 \end{tabularx} \\ \\ \relax
3:14 & \begin{tabularx}{0.7\textwidth}{X} 至於你，你要持守所學習的和所確信的，因為你知道是跟誰學的， \end{tabularx} \\ \\ \relax
3:15 & \begin{tabularx}{0.7\textwidth}{X} 並且知道你從小明白聖經，這聖經能使你因在基督耶穌裡的信有得救的智慧。 \end{tabularx} \\ \\ \relax
3:16 & \begin{tabularx}{0.7\textwidth}{X} 聖經都是神所默示的，於教訓、督責、使人歸正、教導人學義都是有益的， \end{tabularx} \\ \\ \relax
3:17 & \begin{tabularx}{0.7\textwidth}{X} 叫屬神的人得以完全，預備行各樣的善事。 \end{tabularx} \\ \\
[1ex]
\hline
\hline
\end{longtable}
$^{1}$好 今天我們繼續看上帝的話語.
我們敬拜的時候 需要聽上帝的話語.
上帝的話語記載在聖經裡面.
聖經就是上帝的啟示.
如果我們說 上帝何時會跟我說話.
上帝可以報夢給你 可以雷劈.
可以用很多不同的方法.
但最權威的方法 最重要 最標準.
最高的標準.
就是上帝的話語 就是聖經.
所以你記得.
你剛開始信主的時候.
開始教會去摸索一下.
到底基督的信仰是一回事.
信耶穌是一回事.
你通常都會想到.
唱詩敬拜.
敬拜我們說過了.
就是其中一個最重要的環節.
另一個環節是不能夠沒有的.
就是去摸索一下 到底那本聖經是說什麼的.
聖經難不難看.
是難的.
聖經是挺難看的.
但聖經裡面 亦有很多很淺白的東西.
一看就明白的.
我們隔壁房的兒童主教學時間.
兒童的Kids Ministry.
他們裡面很多時候 說到聖經裡面的內容.
都是說一些比較簡單的.
聖經裡面有易 亦有難.
但我們一直信主 久了的時候.
我們就去掌握 嘗試去理解.
所以在這個Back to basics.
即是回到我們基礎裡面.
我們希望能夠去想一下這些topic.
有什麼是重要的東西.
其中一個最重要的就是上帝的話語.
所以我在這裡想.
每一次講道的時候.

$^{41}$有一句所謂的金句.
有一句的經文.
是讓大家能夠在這星期裡面去裝載.
下載了 裝載了.
這個是上星期的.
記得嗎.
記得吧.
我們一起讀.
你要盡心盡性 盡義愛主你的神.
這是最大的 這是第一條誡明.
馬太福音第22章37至38節.
你要盡心盡性 盡義.
All your soul and all your heart and all your mind.
你的思想 都是需要去愛主你的神.
這個就是希望在上星期能夠大家裝載了.
這次我們繼續看.
我們其中一段很寶貴的一段聖經.
就是提摩太后書第三章.
提摩太后書你記得是什麼.
是所謂教牧書信.
即是寫給牧師的書信.
當時有個師傅叫做保羅.
他就寫給他的徒弟.
徒弟就是在所謂教會.
宗教會當中他在做牧羊的.
在做牧師的.
但是年青有為.
保羅作為一個師傅 作為一個導師.
就寫信給這位叫做提摩太.
就去告訴他.
到底有什麼需要去處理.
這個就是提摩太后書其中一段的經文.
三章十六至十七字.
這個就是這個星期的經文.
等下我們再講多一點.
在這裡講聖經有什麼特別.
聖經有什麼獨特的地方.
我們看回這個教牧書信的時候.
我們發現這個所謂第二 Timothy.
提后 提摩太后書.

$^{81}$是保羅寫給提摩太的.
起碼第二封的書信.
第二封寄給他的書信.
來告訴他到底在這個世代裡面.
他們怎樣去處理教會.
怎樣去領導教會.
怎樣去牧養教會.
去到這個提摩太后書第三章.
這裡整章就是一個很寶貴的經文.
頭一半講什麼.
保羅就告訴提摩太.
告訴牧師.
他說在末世當中非常多的危險.
有一個人忽視上帝的道的話.
就會帶來很多的危險.
然後下一半十至十七節.
他就講到需要去學習主的道.
學習主的道有很多益處.
你記得上次我們講到一至九節.
三章一至九節.
他講到末世的時候.
末世是什麼時候.
末世不是將來.
末世是現在.
因為按照聖經裡面的時間線.
上帝帶領人類歷史的時間線.
這個就是末後的最後一個時期.
而這個末世last date.
就是從基督的第一次來.
到基督的第二次來.
這一段時間就叫做last date.
末世.
末世裡面有很多不同的情況出現.
在一個這麼長.
我們人看為非常長的時間.
上帝當然不是看為一個很長的時間.
可能他眨眼之間對上帝.
因為上帝是timeless.
但是他見到人類的歷史最後一段時間.
是有很多恐慌.

$^{121}$很多的假道來傳播.
這是三章第一至第九節.
他就告訴提摩太太.
你要小心.
因為在末世裡面有很多危險的地方.
你記得他說什麼.
末世裡面的人.
是專顧自己.
是冷漠的.
沒有愛心的.
信一些謊言的東西等等.
如果是保羅時代能夠看到這樣的情況.
我們今天更加看到這樣的情況.
我們記得好幾年前.
一說起真的人都老了.
好幾年前疫情的時候.
疫情是什麼時候.
已經過了幾年了.
雖然有時候人的心境還在疫情裡面.
長久疫情還是一輩子的疫情.
已經過了疫情.
你還記得疫情的時候.
特別是開頭的時候.
真的非常混亂.
還有很多的假的信息.
很多的misinformation.
很多陰謀論.
有很多sensational的東西.
在網上傳播.
我經常提醒大英姐妹.
你不要只看一個source.
你看新聞 你看消息.
或者看某一些理論也好.
或者KOL說的話也好.
你不要只看那一樣.
我提醒大英姐妹.
你不要只用Wechat.
因為很多人只用Wechat.
Wechat播什麼給他.
或者送什麼給他.

$^{161}$他就相信了.
你需要雙重檢查.
你要雙重檢查到底發生什麼事.
但你可以了解恐慌或者恐懼的時候.
你會發現那些misinformation.
就是在餵養你的恐慌.
你收到很多假的消息.
或者一些未確定的消息.
這些falsehood.
正是因為你是恐慌的時候.
你就會很容易跳進去.
這個我認識很溫暖的.
已經是21世紀了.
我們不是在古代.
我們已經是資訊很發達.
但那個人不是因為資訊發達而成長.
是需要有判斷力.
本來就是說這個的時候.
末世的時候.
就會有很多人沒有了判斷力.
或者不想判斷.
不想用理性去想事情.
不想去reason.
就不斷被分化.
然後我們就活在恐慌裡面.
等於有時候被恐慌去hijack了你.
去騎劫了你.
本來就說我們需要小心.
忽視了主的道.
忽視真理的時候就會發生什麼事.
然後去到第二半.
或者下一半.
三章十字到十七字.
他就說我們需要學習主道.
去學習上帝的道.
學習真理.
我們看看有什麼經文.
這裡有什麼內容.
保羅是跟他的徒弟說.
第十節.

$^{201}$他跟他的徒弟提摩泰說什麼.
他說但你提摩泰.
你已經追隨我的教導行為.
自向信心,寬容,愛心,忍耐.
以及我在安提,阿爾戈,念路斯德.
這三個地方都是史唐傳裡面提過.
保羅宣教經過的一些城市.
他在那些城市遭遇了迫害和苦難.
如果你看史唐傳就知道.
他去到不同的地方.
傳耶穌的福音.
說耶穌才是王.
耶穌才是君主的時候.
那些人很憎恨他.
特別是猶太人.
在會堂那裡去到撫慰他.
甚至用石頭打他.
在一個城市裡面.
其中一個城市.
他甚至被人拖出城外.
差不多那些人都以為他死了.
誰知道他死不去.
他就跟提摩泰說這件事.
你看到我所經歷的遭遇的苦難和迫害的時候.
我忍受何等的迫害.
但是在這一切的苦難上面.
主耶穌救了我.
我經歷過苦難.
但是主耶穌救過我.
他說其實凡立志在基督耶穌裡面敬虔道逸的.
都會受到迫害.
只是作惡的,騙人的.
將會變本加厲.
迷惑人的也被人迷惑.
敬虔道逸的.
如果你立志要去做一個基督徒.
真真正正做一個基督徒.
將會受到迫害.
我們在美國當然已經很好了.
誰迫害你呢.

$^{241}$為什麼會迫害你呢.
通常都不是因為你做基督徒.
可能跟別人有爭執.
人家才迫害你.
或者你做錯事人家迫害你.
你有多少自由.
不過這裡告訴我們.
如果你想要過一個敬虔的人.
總會有困難的.
不容易的.
剛才我已經說了.
正經難不難明.
難明.
但是做基督徒難不難.
都難的.
特別是我們不同的處境.
可能有不同的試探.
試煉等等.
都是困難的.
第十三節說什麼.
作惡的和騙人的.
會變本加厲.
對嗎.
我很多時候聽到這些新聞.
我都真的.
因為我小時候的子女.
他們有一個代號叫SMH.
你知道什麼叫SMH嗎.
Shake my head.
我聽到一些騙錢.
有些騙錢.
有些無論是以宗教之名.
以經濟之名.
金錢之名.
去騙人.
我沒想過可以這樣.
這樣都可以.
這樣都會有人掉下去.
對嗎.
Shake my head.

$^{281}$末世的時候.
就會這樣.
但是保羅就去吩咐他的徒弟.
他說你要醒目一點.
怎樣醒目一點.
他要學主的道.
首先學主的道要學什麼.
他說你要學我.
保羅說提摩太太你要學我.
他是他的mentor.
我不知道你生命裡面有沒有mentor.
在我過去這麼多年裡面.
我真的有很多好的mentor上帝.
去安排給我.
去猜排給我.
在我生命裡面跟我說話.
在我生命裡面給我不同的提醒.
我忘記了在剛剛開始神學院的時候.
那時候在90年代.
在神學院.
因為我們讀神學的時候.
我們做一個傳道人.
做一個牧師.
我們需要經過一個學位.
有一個curriculum.
有一個curriculum裡面.
其中一樣東西是require的.
是必須要做的.
就是叫做internship.
就是實習.
各行各業都要.
你讀accounting也好.
讀醫也好.
你都要有internship.
我還記得神學的program.
叫做M.Div. Master of Divinity.
都是需要做internship.
不過那時候最大的掙扎就是.
在神學院的課程裡面.
其實internship是footnote.

$^{321}$很多時候都是footnote.
就是一個註腳.
因為神學院的所謂field education.
internship.
通常都是pass or fail.
沒有grade的.
所以他require你做什麼.
你就去一個教會.
去一個教會.
總之那裡有個老總.
你哄他.
你幫他簽名.
搞定.
你field education是很少很少.
幾乎可以沒有的.
我覺得很sad.
你做一個傳道人.
做一個牧師.
你居然field education是搞不定.
有很多時候學的東西.
都是在課堂裡面學的.
這個不一定是神學院全錯的.
因為神學院裡面要.
一個M.Div.的program.
三年半.
full time讀三年半.
你要處理很多很多不同的項目.
不同的課程.
我都講過了.
聖經很難明白的.
因為只是讀聖經.
讀到你臉青.
還要搞這些field education.
所以有時候.
那個internship做得好不好.
是按照我的.
如果你隨隨便便做.
是過不了關的.
誰過不了關.
有什麼可以做.

$^{361}$我們教會有幾個intern.
做什麼.
你知道他們做什麼嗎.
不是很難.
如果你在教會裡面生活過.
你知道他們做什麼.
帶茶經.
可能有機會講道.
等等.
其實不是很重的東西.
所以你學不學到東西.
是你自己的.
你願不願意在mentor身上學東西.
保羅告訴Timothy.
你當教務書信.
可以是一個目的的進修版.
Timothy已經在役忽所去到木槿會.
保羅就去告訴他.
我是你師父.
你要學我.
你留意這段聖經.
最重要的字是什麼字.
中文不是很能翻譯出來.
奇怪.
可能因為太沉.
三章第十節.
Timothy你已經追隨我的教導.
行為自向信心.
寬容愛心忍耐.
你看英文是怎樣.
Timothy你已經追隨我的教導.
我的行為.
我的自向.
我的信心.
我的寬容我的愛心.
我的忍耐.
為什麼經常是我的.
因為你看到我的example.
你跟我的時候.
你在我身上學的東西.

$^{401}$就是一個生命的傳遞.
你一個人的生命裡面去學到東西.
我們在職場很多年的人.
真正學到東西的很多時候都是你老闆.
當你老闆是好老闆.
我知道有很多不好的老闆.
當你老闆是很好.
如果他真是好.
真是去叮嚀真是去教你的話.
你在他身上學東西.
對不對.
不是書本裡面學到.
你看他怎樣deal with不同的問題.
你看他那個知識上怎樣去應用出來.
特別當我們生命裡面有不同的境況.
他應該怎樣去處理等等.
是在mental身上才學到東西.
本來就是說我的.
在我的生命裡面見到我怎樣教東西.
我遇見有不同的困難有迫害的時候.
我怎樣去處理.
教會裡面有不同的人頂撞他.
他怎樣去寬容.
這些都是要在mental裡面身上學.
Again我問你.
你身上裡面有沒有mental.
在生命裡面有沒有mental.
可能是大聖主的人.
可能是你曾經有的導師等等.
這些都是你生命裡面需要去學習的.
學聖經原來不只是打開一本書.
學聖經是在人的生命裡面去學.
有什麼東西值得你去學習.
我們繼續看下去第十四節.
他就吩咐提摩.
至於你.
你要持守所學到的.
所確信的.
因為你知道是跟誰學的.
並且知道你從小明白聖經.

$^{441}$這聖經能使你因基督耶穌裡面的信.
有得救的智慧.
他鼓勵他.
你現在是一個很困難的時間提摩台.
不過你去到學習.
學習我的.
還有你明白聖經.
所以他說學到是怎樣的.
在人的生命裡面.
在心靈的生命裡面去學.
然後還要在聖經裡面去學.
他說提摩台就是從小去明白聖經.
從小開始就明白聖經.
如果我們翻開一些不同的資料.
我們就知道提摩台是一個怎樣的人.
他媽媽原來是基督徒.
他的婆婆也是基督徒.
所以你這樣算的話.
提摩台是一個基督教的家庭出生.
提摩台後書第一章第五節.
即是同一本書.
前兩章第一章.
第五節說什麼.
他說我記得你.
保羅說我記得你提摩台.
無偽的信心.
見到一個無偽的信心.
一個真確的信心.
這個信心是存在你的.
外祖母羅兒.
和你的母親有離基的心裡面.
我的心信也存在你心裡面.
原來我們知道提摩台這個人.
有媽媽的.
當然他有媽媽.
他媽媽的名字.
就是Lois.
Unice.
他的祖母是Lois.
保羅也認識.

$^{481}$他說提摩台你是怎樣.
他所謂third generation的question.
第三代的基督徒.
很不簡單.
今時今日很多時候一代就沒有了.
但他提醒提摩台做什麼.
當時間艱難的時候.
你需要想起.
不單止想起我作為mental.
提摩台要想起你媽媽.
和你外祖母.
想起生命裡面的spiritual的前輩.
他媽媽和外祖母怎樣去教他.
當然我們作為家長.
我也是家長.
做爸爸.
我們不能夠保證下一代一定會成為基督徒.
因為這是下一代需要做的決定.
你明白我的意思嗎.
是不能保證的.
因為最終是他自己的選擇.
信仰是一個個人的投身.
不能代他上他身去決志.
不能上他身去跟隨基督.
始終是上帝對他呼召.
然後他作正面的回應.
這是他的示威.
我們不能保證.
但我們只能夠做好本份.
我做一個爸爸.
我是牧師家庭.
在牧師家庭長大好嗎.
很幸福吧.
我不知道其他牧師.
我想我家庭就不是這麼幸福.
但我記得開始的時候.
我和太太有這樣的想法.
我們不會跟小朋友去睡覺.
我從來沒有跟他去睡覺.
我記得唯一一次跟他去睡覺.

$^{521}$大概他一歲左右的時候.
我女兒一歲左右的時候.
我女兒有一次晚上不肯睡覺.
你知道他的床有圍欄.
他就在拉著圍欄.
我剛開始走的時候.
他就在叫不知道叫什麼.
好像在鬼叫一樣.
他就吵著我們家庭.
然後我很生氣.
因為明天是主日.
我有很多事情要做.
然後他就衝進去房間.
我馬上跟他說了一篇.
我跟他說了一篇五至十分鐘.
我主要跟他說不要自私.
我請你不要自私.
因為我們每個人都要生活.
我們每個人都要工作.
你這樣鬼叫我們要怎麼做.
你沒什麼要做不代表我們沒什麼要做.
類似這樣.
我一盤這樣出來.
這是我的職業病.
我記得說完這篇道之後.
他就靜了.
那一晚我們就睡覺.
我記得第一次.
只是一次就跟他說睡覺.
那時候因為我軟弱.
累啊.
不如說五分鐘的道.
馬上睡覺.
很生氣.
但通常我都不會這樣做.
但我們會盡力去真誠.
做一個真誠的基督徒.
不等於我是完美.
我絕對不是完美.
我做錯事對小朋友.

$^{561}$罵他們的時候.
我的火氣大了.
我會向他們道歉.
希望他們原諒.
我覺得這才是最重要.
我告訴他們.
我都需要恩典.
我都需要基督.
我更加需要上帝的話語.
所以他們看到我查經.
他們看到我看聖經.
開始年輕的時候.
他們未必去欣賞.
到底在做什麼.
但開始慢慢他們會欣賞.
我記得我女兒現在開始工作.
她在她工作的地方.
是一個基督教補習社.
有時候他們會有早上的懷念.
她就會分享.
我記得她分享了一次.
她跟其他工作人員說.
她說一件事她很慶幸.
就是父母沒有壓迫她信主.
沒有逼她.
我只是做一個環境.
希望她能夠去探索信仰.
去明白信仰.
我相信這也是我們每一個基督徒的家長.
我們需要去努力做的事情.
我們希望子女在我們身上.
看到我們是真真實實.
並且看到恩典的美.
基督徒生命的美.
它的美.
不只是對與錯這麼簡單.
它是美麗的.
這是我們有的目的.
其實我們每一個基督徒都是一樣.
我們希望活出一個真正的生命.

$^{601}$是能夠告訴別人我們不是完美的.
我們的教會不是完美的.
然而我們是真正的.
我們需要恩典.
每一個人都需要恩典.
如果睜開眼睛看看.
每一個福音為中心.
每一個以真道為中心的教會.
都應該有這樣的榜樣.
有時候可能少一點.
有時候可能多一點.
在教會裡面.
但應該是有的.
我們需要留意去學習他們.
學習這些屬靈的榜樣.
有什麼我們值得去欣賞.
值得去跟進.
剛才我講到神樂園的實習.
我記得我每一次去實習.
實習可以隨便過關.
但我每去一個教會.
每去一個群體的時候.
我都希望可以盡心.
我記得我第一次去實習的教會.
根本不是我自己同聲同氣的人的教會.
我其中一個廣東人.
五隻手指都數得完.
全部都是大陸或台灣的.
但無所謂.
盡心盡性去做我應該做的事.
在那裡我找到很多值得學習的生命.
其中一個就是當時的主任牧師.
他找我聊天.
我很樂意跟他花時間.
這個牧師到今天.
我跟他都是朋友.
我記得我離開教會的時候.
他很想留我.
但上帝帶領我要走.
你覺得他喜不喜歡我.

$^{641}$甚麼零頭.
他很喜歡我.
多麼幼稚.
怎會不喜歡.
但有時我們會打電話去聊天.
我還記得有一次在耶路撒冷.
剛剛帶隊的時候就遇到他.
又聊了一晚.
我在他身上欣賞他甚麼.
當時他的教會是一個初生教會.
好像我們這樣.
很多困難.
一直到租地方.
要克服很多阻力等等.
但我見到他在他身上.
他的忍耐,祈禱和幽默.
有一種喜樂.
這些就是書本裡面學不到的東西.
今時今日還是一個.
大家都是很好的朋友.
睜大眼睛看看.
你願意的話.
最慘是你不願意.
你願意的話.
肯定可以在這些人的身上.
學到東西.
這就是學到裡面的精神.
不單止這樣.
保羅提醒提摩他甚麼.
這就是聖經裡面.
舊約新約.
其中兩節是最重要的.
最重要的經文是甚麼.
因為講出聖經的本質是甚麼.
聖經是一回事.
聖經都是神所默示的.
於教訓督澤使人歸正.
教導人學而都是有益的.
叫屬上帝的人得以完全.
進行各樣的善事.

$^{681}$為何這麼重要.
因為有幾個key word.
第一個 神的默示.
God's breath.
聖經都是神的breath.
神的話語 神的說話.
神說的話 神的口氣.
這裡講出上帝是一個怎樣的上帝.
上帝是說話的上帝.
菩薩是不會說話.
其他假神是不會說話.
有口不能言.
但上帝是能夠說話的.
你記得勝經.
我們不用知道很多.
你打開聖經第一章.
上帝第一件事做的是甚麼.
神說要有光.
就有了光.
God who speaks.
以前他講說話.
今時今日他都會講說話.
他說甚麼講說話.
聖經裡面講說話.
藉著聖經講說話.
所以聖經就是上帝的啟示.
或者應該這樣說.
拿起聖經的時候.
或者打開聖經的內容的時候.
最終的作者是誰.
不是保羅.
不是彼得這麼簡單.
是上帝自己去寫作.
是藉著人去寫作.
All scripture.
或者全部的聖經.
當時有的聖經.
保羅時代第一世紀的時候.
新約都未完整寫好.
27卷.

$^{721}$只有舊約.
或者舊約希臘文的翻譯本.
我叫74本.
只是舊約聖經.
只是Old Testament.
他說這個已經是上帝的默示.
只是拿著舊約的聖經.
已經知道上帝會說話.
對我們說話.
何況我們今天已經有了新約.
27卷.
66卷的聖經說出一件事.
整件事就是上帝怎樣救贖人類.
上帝對我們的心意等等.
並且他說.
聖經都是神所默示的.
有什麼益處.
有什麼profit.
或者useful的地方.
他說了四樣.
第一樣.
教導.
督責.
使人歸正.
教導人學而.
我很喜歡現代的一個翻譯.
New Living Translation.
他就將四個功能.
很簡單的.
用一句來說出來.
to teach us what is true.
God is true.
Jesus Christ is true.
耶穌基督是真實的.
這個就是教導.
或者當代譯本我們看到下面.
聖經全部都是上帝啟示的.
有益於什麼.
教導.
告訴你什麼是真理.

$^{761}$什麼是真實.
並且督責.
to make you realize what is wrong in your life.
Right.
說真的.
有時候我們不想看聖經的原因是什麼.
有時候看聖經好像照鏡子.
照妖鏡.
聖經會confront我們.
是吧.
但這個就是聖經要做的事.
是吧.
先去confront我們.
用真理來confront我們.
因為我們是軟弱.
我們是有罪.
我們是破碎的.
所以為什麼我們需要上帝的恩典.
to make you realize what is wrong in your life.
你肯去面對這個真理的時候.
你才明白.
我才可以擁抱到真理.
第三樣.
correct us.
不單止指出我們有什麼錯.
甚至導航我們.
以致我們走一條正的路.
teach us to do what is right.
繼續告訴我們我們在進行中需要做什麼.
所以當代日本人的翻譯也不錯.
培養你去行義.
培養你的生命.
去裝備你.
以致你能夠成為合神心意和合神所用的人.
in fact 其實這裡說teach us what is right.
其實就是discipline.
去到練習的意思.
奧運會快要開了.
有沒有看電視看奧運的選手.
他們的表現.

$^{801}$為什麼他們的運動可以這麼美.
為什麼可以做到這麼好.
他們經過什麼.
很嚴格的學習.
很嚴格的discipline.
他們才能夠這樣做.
其實基督徒的生命都是一樣.
為什麼有些基督徒的生命.
是我們能夠學習的.
原因就是他們沒有漏了discipline.
明白聖經 清楚聖經.
我們才能夠按照上帝的話語來做事.
求主幫助我們.
我們做一個願意學到的.
這個禮拜的經文就是這句.
馬太福音第五章十四節.
我告訴你為什麼我選這段.
先讀 你們是世上的光 成造在山上 是不能隱藏的.
為什麼我選這句.
因為猶太人的傳統.
什麼叫光.
光當然是一個很好的比喻.
比喻很多東西.
在猶太人的傳統拉比的文獻裡面.
他說Torah 上帝的律法就是光.
上帝的律法就是真光.
耶穌在說什麼.
他對門徒說 對我們說.
他說你就是世上的光.
就是說你就是世上的Torah.
你就是世上的聖經.
你就是未信主的人 未接觸的聖經 未接觸信仰的人.
他們不會看聖經.
他們看你.
你就是那本聖經.
你就是上帝的道.
如果你按照上帝的心意 按照上帝的道來生活的時候.
人就見到上帝的律法 上帝的道的奇妙.
今天我希望你.
我們結束祈禱的時候.

$^{841}$我們做一個立志.
真的花一點點時間.
用一點點心機.
學習 去學到.
好嗎.
我們跟上帝一起祈禱.
上帝我們謝謝你 因為你的道就是真光.
你的光 光照在我們心靈裡面黑暗的地方.
在世上黑暗的地方.
去揭示假的事情.
去傳揚真正的真理.
主求你讓我們能夠擁抱.
擁抱你自己的懷疑.
不同的師父 不同的榜樣.
主要我們學到 是完整的來去求學.
完整的來去明白你自己的心意.
主教導我們.
讓我們成為一個真理的人.
讓我們愛你的真理 愛你的道.
以致我們能夠成為世上的光.
世上的聖經.
世上的上帝的律法.
主幫助我們這個星期.
我們能夠在不同的地方.
我們去學習 形容.
培養我們 以致我們的心靈能夠再有你自己豐盛的一天.
我們這樣祈禱 奉尊名.
完善祈禱 祝福各位.
走.
\newpage



\section{}
\label{sec:9TgjTCjHTHQ}
\textbf{Back to Basics \#5 (Philippians 4) - Rev. Mark Chan}
\newline
\newline
連結: \href{https://youtube.com/watch?v=9TgjTCjHTHQ}{\texttt{https://youtube.com/watch?v=9TgjTCjHTHQ}} ~~~~ 語音日期: 2024-08-01
\newline
\newline
\hyperref[sec:koqOwadJaEc]{\small{< < < PREV SERMON < < <}}
~
\hyperref[sec:index_chronic]{\small{[返順時目]}}
~
\hyperref[sec:index_scriptual]{\small{[返順卷目]}}
~
\hyperref[sec:gEFzH4C2B2Y]{\small{> > > NEXT SERMON > > >}}
\newline
\newline
$^{1}$By the one quick important reminder is that our brother Cal it.
sees this from Lebanon's.
Lebanon and he's a Lebanese brother and.
he is preparing lunch for us and so right after service you will be able to experience some.
Mediterranean Arabic hospitality and.
He spent hours and hours and he and his friend and maybe jaws is going to go out and to help him, right?.
Yeah, some of the brothers are helping helping him to prepare.
some hummus.
and pita bread and.
Kabob, right?.
So you're gonna enjoy it when he was.
preparing for for the.
Lunch or refreshment today. I was helping him yesterday actually so I bought I brought my.
charcoal grill.
To to the church and so I think they are grilling the meat and I was I was kindly reminded.
the Mediterranean.
Arabic.
Hospitality, it's like that. You know, they want to go all out. You know, they want to go all you go to a.
Arab family.
They invite you in and they will do everything to make sure you are being on the guest.
So that's that's how they the the honor.
people in their lives and I'm very thankful for.
Califor coming to know Christ and finding.
this church to be his spiritual home for now and.
So so just want to remind you right after this I know you're hungry I know you're hungry that's why you late, right.
You're hungry, right, especially after the service, but you know go enjoy some hummus, right go enjoy some.
Arabic food Lebanese food. All right.
Anyways, we are coming back to this.
Series called back to basics and I want you to remember and practice some of the basic.
Habits, you know in in your Christian life and if we have those slide ready, we will.
first reveal the verse of the week and.
This verse of the week is if you have the Bible, I don't have a slide yet.
You have the Bible, please go to.
Philippines the bow Philippines chapter 4. Yeah, Phoebe. Can I have the slide please?.
Anyways, we'll go to chapter 4 verse 6 and 7 verse 6 and 7. All right.
Here it says don't worry.
Here we go. Don't worry about anything.
Instead pray about everything tell God what you need and thank him for what he has done.
Then you will experience God's peace which exceeds anything you can understand.
His peace will guard your hearts and minds as you live in Christ Jesus.

$^{41}$Philippians 4.
Verses 6 to 7 what a wonderful verse right? What a wonderful sentence or a paragraph.
to.
Tell us, you know, especially those of us who worry some.
Worry a lot and overthinking a lot. Are there anyone here?.
You constantly overthink you constantly think a lot, you know about your future about you know issues or problems.
This verse gives you peace. Don't worry about anything instead pray about everything now.
I want to let you to understand this is taken from the book of or the letter to the Philippines and.
what this letter is as you understand the New Testament has a total of.
27 books or 27.
You know pieces of writings some of them are letters. Some of them are stories.
Some of them are histories and you know different different types of literature.
But if you go to the New Testament Bible.
there were 27 books there or 27 writings and some of the writings we call epistles or letters and.
Many of the letters are actually written from a pastor.
From a pastor leader to a church a very very small church.
very very small house churches in different cities and here we find the letter of.
Letter to the to the to the Philippines and it's written by Paul the Apostle Paul the pastor Paul and.
He wrote this letter.
Call we call the letter to the Philippines and writing to the Christians the house churches in.
Philippi one of the Roman.
Quarantine in in in Greece and.
If you read this letter.
There were only four chapters is very very short.
Book a very very short letter you can sense that this is a personal letter.
With with a lot of personal tones a lot of personal.
Feeling to it and.
There are two reasons why.
first because.
Paul actually found this church.
He established this church in his second missionary journey if you check the book of Acts, which is the history of the earliest.
the first generation.
Christians you'll find some of the stories that Paul.
came to the city called of Philippi is a Roman city and.
He tried to preach the gospel and some of the individuals came to faith and he established the church.
And he's like a father to the church in Philippi.
He's like a father to the Christians the Philippian Christians. And the second reason is the church has been.
supporting Paul.
and his mission and his.
ministry and sometimes they send.

$^{81}$the supplies or.
Essentials to Paul and at this point when when when Paul write this letter the letter to the Philippines.
He was in prison.
He was in Rome. He was in the capital city of the Roman Empire.
He was right there in house arrest.
for almost two years as far as we can reconstruct the.
The last few years of Paul and he was waiting for trial because he want to appeal to Caesar the king.
the Emperor of Roman Empire and.
He did not know the result yet.
He was in house arrest and in his house arrest in the Roman world. You can still receive, you know information from other people.
You can receive visitors. You can have friends and you can send letters and that's the context of the Philippian.
Letter and in this in this letter you get the sense that you know, he's trying to reply to something.
Because we find some verses in this fourth chapter the book of Philippians that Paul is trying to thank.
The church in Philippi, right? He was in Rome, right and.
The Philippian Christians are in Philippi and they send someone to to give the essentials or offer essentials to support.
Paul while he was in prison and.
For example, you can go to the last chapter chapter 4.
417 what does he say? He say I don't say this because I want a gift from you.
In other words a possible receive a gift a financial gift from the church in Philippi rather.
I want you guys to receive a reward for your kindness at the moment.
I have all I need and more more than enough.
I'm generously supplied with the gifts you send me with.
Epaphroditus Epaphroditus and they are sweet smelling sacrifice that is acceptable and pleasing to God now.
Who's this guy before Titus?.
It has to be a member or even a leader from the church in Philippi and he traveled all the way.
to come to Rome to find Paul in house arrest and give him some.
Essentials probably some food items probably some clothes or anything that he needs in two years and constantly he does so and.
Epaphroditus is kind of like, you know, kind of like a UPS, you know a delivery delivery guy or uber eat, you know.
he's gonna support Paul and make sure he's being fed and he's he's being strengthened and to help him and to stay with Paul and.
From the letter we see that Epaphroditus was really sick.
It's probably because he travels so much to help Paul.
and maybe he was got so sick and even to the point of death and.
He was by the side of Paul and Paul is really really.
worried about this brother and maybe because.
Paul and.
Epaphroditus they talk about.
the church and there was so much anxiety for the church because the small church is under a lot of pressure and.
oppositions and Paul want to send Epaphroditus back to the church in.
Philippi and he wants to send Epaphroditus back with a letter and this is the letter.
Right the letter to the Philippians and he you can imagine he's just quickly grab some.

$^{121}$parchment paper and write a letter in four chapters a very very short a very personal letter and send it with the.
Mailman, you know the UPS guy Epaphroditus and send him back to.
Philippi and that's how we find the letter to the Philippines and finally we come to the.
last chapter chapter 4 and.
And Paul is trying to close the letter and he's trying to you know.
Give give his encouragement as much as he can in this letter and chapter 4 verse 1.
He said therefore brothers and sisters you whom I love and long for my crown and my joy.
You got that personal tone.
It's a personal letter and that's how you should stand firm in Lord dear friends. I.
Plead with you dire and I plead with.
Syntaghy to agree with each other in the Lord and yes, I ask you loyal your fellow.
hope this women who have contended at my side in the cause of gospel along with.
Clement and the rest of fellow workers whose name are in the book of life rejoice in the Lord always I say again.
Rejoice, let your gentleness be evident to all and the Lord is near.
Now that's the closing chapter right there. There are two women being named here, right you dire and.
Synecdoche, right, but perhaps these two women are.
Some of the leaders.
In in Philippi some of the leaders of a church.
perhaps they invite the church or fellowship to the living room and have the house churches meeting possible and.
These two women and from from the context, you know, these two women don't agree with each other.
Right, they got some tension inattention or probably.
Disagreement for some reason they were in disagreement and Paul is writing this letter encourage them to be united.
To be of one heart right and that's actually the tone for the entire letter.
Paul the pastor wants the church to work in unity.
To along with other people like Clement right other fellow workers.
So that brings him joy, right and you know from from the wording he used.
The communities in Philippi they are working toward that goal already.
So Paul is trying to encourage them to continue to work in unity and that brings him joy.
Right the entire letter only four chapters, but over and over again.
Paul is expressing his joy and he's happy.
He's joyful.
because of the gospel.
Because the church is being established because the church is trying to stand up and be strengthened in faith. And this is a powerful.
Passage because as you can imagine someone who is in house arrest and.
You just don't know what it will turn out to be.
Right, whether you will live or where you would die, right?.
he is waiting to be to go on trial and.
He was in house arrest and he was in prison and just such powerful.
Wordings that he was trying to use you guys are my joy and my crown and there's a secret here.
Why did Paul do that? How could Paul do that? You know, imagine you're in prison, right or in house arrest.

$^{161}$right.
And you can't go anywhere.
You don't have your strip of your freedom, right and you are being mistreated.
What is the secret of being joyful here?.
You can tell from verse 3.
Paul says all of their names are written in the book of life, right the book of life.
You believe in Christ Jesus if you trust holy trust in the name of Jesus your name.
Is in the book of life, right?.
Your eternity is secure, right? You are on.
the team of team Jesus, right, you know, the Olympic has.
begun and.
there were different athletes right and.
The they they are honored to join team USA, right or the team of the country, right?.
Just joining the team the name being on the team is such an such an honor, right?.
Right. Yeah, you can imagine if your name is on team USA, right? What kind of honor would that be?.
Right, we're gonna throw a big party for you, right too bad. Your name is not there, right?.
Yeah, but your name is on team USA you get the honor and you are happy, right?.
And that's the secret of a Paul.
It's if your name is in the book of life.
You have the ground to be joyful, right? You're joyful not because you are successful.
You're not you're not because because you're not you're not you're not joyful because you're successful in the in the worldly perspective.
You're joyful because your name is in the kingdom, right?.
And then he goes on to encourage the questions encourage us that the Lord is near.
Right now it is usually translated in the in the sense of time, you know temporal sense.
The Lord is near meaning Jesus is coming anytime.
It's just coming anytime soon, right? It can be today. It can be the next hour.
It can be five years later anytime right Jesus coming second coming is right around the corner.
That's the temporal sense. But you know here another sense is is often forgotten.
In those words the Lord is near.
It's the spatial sense.
So what Paul is saying is that the Lord is near Jesus is near is that Jesus is right by your side.
Jesus is present with you if you're facing difficulties or you're facing some sadness and depressed.
moments in your life.
Know that you know, the Lord is right next to you.
Right right beside you. The Lord is near he is someone you can trust.
He's right next to you and it makes a world whole world of difference, right?.
It makes a whole world difference knowing that there's supporting friend right next to you.
Right when you are feeling down when you're feeling bad or sad, right?.
I usually.
visit.

$^{201}$People who are sick in the hospital, right and you can tell from their eyes.
When you go in and you preach you are present with them.
They are encouraged right? They are comforted and.
That's why it's hopeful for many of us to sometimes to talk to someone who is wiser.
Like a wise counselor, right? Either professional counselor or you know a layman.
Find someone who can talk to you find someone who can listen to you to understand.
Your voice right understand your mind.
That is so important. And and that's what it means. The Lord is near right friends if you're going through some down period.
If you're going through some difficult time know that the Lord Jesus is near.
He is right next to you and Paul goes on to talk about one of the basic practice.
Of Christians and don't forget that he said do not be anxious.
That's NIV do not be anxious about anything but in everything by prayer and petition.
With Thanksgiving present your request to God and the peace of God which transcends all.
Human understanding will guard your heart and your minds in Christ Jesus.
he says.
Don't worry.
Right. Don't worry.
Don't worry too much.
Don't be too anxious about anything or perhaps our essentials, right?.
Don't be anxious about anything but in everything you pray and you present your petition to God, right?.
They're almost synonyms prayer and petition.
not the exact wording but you know words that.
Express our way of you know.
Talking to God right and that is prayer. That is the basic practice that we all have. Let me ask you this.
Do you pray?.
Constantly regularly.
Every day. I want to encourage you you take 10 minutes of.
The best time of the day to pray to God right now so often we we we misunderstand.
What prayer is right? Sometimes we think prayer has to be in some sort of using some some sort of big words, right?.
Beautiful words so we can approach God that way. No, that is not the spirit of prayer.
I still remember the first year in in seminary.
I was taking a class and there's a very very.
old and wise professor.
Hey, he always begins the class with prayer and end the class with prayer and when he prays.
No, it's just beautiful prayer.
He uses King James Version English to pray.
You know Tao T and it's just beautiful right because he grew up in the Christian family.
Very devout Christian and he used to that kind of language, you know, King James.
It's very very old old English. But you know when you hear his prayer, you just feel like my goodness.
This is beautiful.

$^{241}$Every time he he opened his mouth. It's like writing poem, right?.
Do we have to pray like that?.
No, we don't right, you know one of the shortest prayers in the Bible.
remember that is.
Lord help me.
That's one of the shortest prayers in the Bible right so his prayer in the Bible the spirit of prayer is authentic.
sincere honest and open you just tell God what you need and how you feel right and.
It's kind of like baby talking to to daddy and mommy.
Right. Do we as as dad and mom do we mind the young child?.
Not knowing what to say and you know, just bah bah bah all the time, right? We love those those voices, right?.
To make known right his or her feeling that that's how we approach God, you know.
Will prayer is authentic.
Authentic sincere honest and open right and we need to practice that right?.
We need to practice that and then Thanksgiving.
He said by prayer and petition with Thanksgiving, right?.
When you when you cultivate the spirit of being thankful.
You are constantly reminded.
What God has done in your lives, right?.
Sometimes we are so overwhelmed by you know, the problems that we have right the worries that you know in our mind.
And then we forget.
No, God has been good to us, right?.
God has been good to us in the past and that's why we can keep on.
Trusting him right there. God has the track record, right?.
of being good to us.
And that reminder will help us.
To worry less right a christian author.
Says listen up christian author says this. I think it says it's very very.
correctly.
The way to be anxious about nothing.
Is to be prayerful.
about everything.
The way to be anxious about nothing.
Is to be prayerful.
about everything.
Right. What a wonderful reminder, right?.
We come here. Paul says don't be anxious about anything but in prayer petition and thanksgiving.
You come before God, right? And that's a wonderful verse, right? But it's also interesting, right?.
you know you think about it from a.
Just our our human perspective if God wants us to worry less.
Why doesn't God just give us everything that we ask for right? Do you want a bigger house?.

$^{281}$Do you want to win the the lottery.
And retire early right is sitting by the beach of cancun, you know all the time, right, you know sipping.
You know pina colada all day.
Don't you want that?.
Right. Don't you think you need that?.
You want bigger bigger car or a better car, right?.
Um money, you know lottery new house right new car right, you know.
Uh, you think God God would be good to us if if if God just shower us with.
all this material possessions.
You think you think that's the goodness of God?.
You say so right you say so but you know what if if you have all that you want, right?.
Sometimes your spirit is poor.
Right because it doesn't build up your character.
Right, it doesn't give you the strong spirit to withstand changes to withstand difficulties, right?.
And that's why paul is saying, you know.
God treats us.
in a different way.
God gives you peace.
Right god gives you peace.
And peace that is beyond human understanding transcends all understanding and talking about peace.
You know to the uh to the years of uh, the philippians in the roman city roman colony.
And uh, they should uh, they should immediately understand what it means peace, right?.
Sometimes we read the word peace we use our understanding.
uh.
And our perception but in the to the roman.
To the romans and and philippian, uh, they're living in a roman city a military city.
They do know what peace is right and there's a slogan.
Back in the uh, first century roman empire is pax romana, right pax romana latin.
It means roman peace, right? You know what what roman peace is?.
You know, uh, rome is a dictator, right?.
Roman peace means I beat you up.
And then you're gonna obey me right I conquer you right and you're not gonna do uh.
You're not gonna be rebellious. You're not gonna rebel against me. That is called peace, right? That is roman peace.
Pax romana, right? I was trying to organize some of the materials lately and uh, I find this uh, roman coins.
Golden coins and what we call this kind of coins are first jewish revolt coins.
You know the first jewish revolt happened in uh, a d uh, 66 to 70.
Those few years, uh, some of the jewish people especially in judea in jerusalem. They rebelled against rome and they fought against rome.
and um, of course they.
Um, ultimately they failed right and the roman army came to palestine and suppressed the revolt.
And after the suppression after roman, uh, the romans conquered, uh, all those, uh,.

$^{321}$rebellious people.
uh, they.
Will uh, they will send a propaganda to the entire empire by minting coins.
And coins with a message, right?.
You know coins have message, right?.
And on the left you see there's a military victory.
procession.
And all the armies and horses, you know and with a horn, you know proclaiming the victory of rome and in front of it.
The little guy, you know the first guy in the procession.
Uh, he has a name.
And his name is simon.
simon son of.
Uh, gilora.
And he's one of the leaders.
of the uh first jewish revolt.
And uh, the romans came and say we are triumph over you. Okay, right and we beat you up.
And you are the head of the procession and we're going to tell other people all the tribes in all the nations.
That we are the winner and you are the loser, right?.
And then who's the winner?.
that's bastion.
The middle the middle face middle face of a coin, right? That is one of the roman empires.
emperors.
And on the right you see the word judea.
Can you make make it out judea?.
And it's a woman.
captive.
You know being bound.
of her hands.
sitting at the foot of uh, uh the palm tree.
And she was defeated.
right, she was crying she was weeping because.
She has no judea has lost the battle, right?.
Now what I want to tell you is this is the room. This is the roman peace.
This is the peace of roman empire they beat you up, right?.
So you cannot rebel anymore, right? And philipi was guarded heavily by.
uh roman.
roman military units.
And if you live in philipi, if you're one of those who read the letter of philippians.
You understand you because you will see soldiers all around you, right?.
Every single city gates you will see roman soldiers heavily guarding the city, right?.

$^{361}$Roman army is to guarantee the roman peace, right?.
That is the peace that they understood.
right.
roman peace.
I still remember, uh the first year of the pandemic.
Remember that we have national guards here.
You know guarding the in and out of downtown.
Remember that?.
You know, they they got uh, you know military trucks and uh, you know all over the cities, right?.
And I remember I have to help uh international student to move.
and uh, I I rent an.
a 10-foot u-haul.
And she has to move out of the uh, you know the the building in downtown.
And I I drove that u-haul.
You know across all the checkpoints, right?.
Dude, I was so scared, right? Well, what are you doing? You know people are using u-haul to.
you know to to to uh.
To to break break down the wall and uh steal things right?.
And it was pretty scary right? I have to talk to the guard and you know, I I'm a pastor.
I'm helping I want people to move, you know, I know this is a bad time, but you know.
The street is like a ghost town, right?.
Do you think that was peace.
In a way it was right, we're quite peaceful.
It's quite peaceful nobody on the street, right?.
But that is the roman peace.
The peace that paul talks about is the peace.
of god.
Right the peace that we find in jesus christ.
And first of all that peace is that we have a restore.
Relationship with god the father right in jesus because of the cross that we have that peace that ultimate peace.
Restore relationship with our creator with our father. We can call him father.
We are sons and daughters and that is the peace talk paul talks about peace given by god, right?.
So paul pictures god's peace.
as the garrison.
As military soldiers keeping guard.
Over your hearts and your minds to protect you from all kinds of attacks.
that is.
the peace of god.
And that's why it's beyond human understanding.
Right, you know that your eternal destiny is secure in christ.

$^{401}$You know, your identity is secure in christ.
And you are not shaken.
No matter what comes your way, right?.
sickness.
right separation.
Heartbrokenness.
All this all these issues are coming in your way.
If you know that your identity is secure.
You have peace and that is the peace from god, right?.
And we can rely on god.
You know, rachel has been driving a lot these days, uh, I told you, you know in the in the beginning a couple years ago.
She hates driving.
Now she likes driving.
She enjoys driving.
And um, you know, she drives one of our old cars.
and uh, she always always comes on, uh, she always comes home and uh,.
uh, tell me.
with a little phrase.
And she said dad, uh, the car is low on gas.
The car's low on gas and she didn't say anything right because she trusts.
That i'm the i'm the good father, right? I'm gonna add uh gasoline to the car when the gas is low, right?.
And um, I really admire that kind of chill out spirit.
Uh, you know just trusting me that I will always provide right? I will always provide.
Yeah.
I think that kind of spirit we need to have.
Without god, right.
We're low on gas.
And father heavenly daddy, we're low on gas.
And that is the spirit of prayer, right? That is the spirit of prayer.
We come before god. We are open. We come before god. We are sincere and honest and that is.
Our daily prayer life, right?.
Come before god to tell him our request.
Tell him our feelings.
Tell him what is what is going on in your life?.
And god is there to listen and the lord is near.
That is the mentality of paul, right? He was in prison. Remember?.
He was um.
under a lot of pressure.
Uh opposition there was no light at the end end of tunnel.
He just doesn't know what the future holds for him. But there was a habit in paul. There was a habit of uh,.

$^{441}$of a faithful living.
To pray and to be thankful all the time, right?.
That takes care of a lot of different issues and worries.
Beyond our thinking beyond our reasoning, right? Don't we all don't we all overthink, you know from time to time?.
Right. We thought that we think that our thinking can can take care of things. No, we don't right and paul is saying and.
Trust trust in god.
Right and the peace of god will be with you.
You know, I watched the uh.
olympic game opening.
And one of the uh, most inspiring moments.
is celine dion.
You watch that celine dion present.
her uh singing.
It's just inspiring right, you know in the past few years she has gone through just.
Just no enormous difficulties, right? Yeah losing losing family member.
um, you know, uh.
You know where disease, you know striking her.
but um.
She keeps on fighting, right?.
You know, it doesn't come overnight, you know that kind of spirit doesn't come overnight.
Because she has been a fighter.
For so long right she had this habit this pattern of life this perspective of life to keep on fighting, right?.
Here we see this we see the spirit in paul.
And paul is trying to encourage us. Hey.
You know.
God is with you.
Trust in him, right?.
pray.
So I want to encourage you to uh to develop.
A constant prayer life. It's very simple. There's no it's not rocket science, right?.
Uh, take 10 minutes.
The best time of your day.
To pray to god, right? You can take pray verbally.
And even non-verbally right and god knows, uh your heart, right.
Let's do that.
and um before we conclude.
Uh, let's recite this, uh, wonderful passage again, right?.
See.
philippians 4 6 to 7.
Don't worry about anything instead pray about everything.

$^{481}$Tell god what you need and thank him for all he has done.
Then you will experience god's peace.
Which exceeds anything we can understand?.
His peace will guard your heart and mind as you live in christ. Jesus.
That's our prayer.
Father no matter what life comes.
In our way and we come before you.
Knowing that you are near.
You're right next to us.
As we close our eyes and bow our head. We know that you are with us.
You're right next to us.
You know what we are going through.
Either it's loneliness sadness.
Uncertainty some difficulties in our lives.
on.
Just so many problems in our lives.
We know that you care for us.
And thank you for calling us to be your sons and daughters.
We come before you to pray.
With an open heart.
Authentic spirit.
Sincere mind knowing that you listen to us.
You hear our prayer and you care for us.
And may your peace your peace that goes beyond worldly understanding.
Will dwell in our heart and be the core of our strength.
So we continue to live in christ jesus send us out from here.
Uh to our school to our workplace to our ministries.
And may your peace guard our hearts and mind this week.
Thank you in christ name pray.
Well, make sure you uh stay.
uh for our.
I don't know refreshment or or meal right? Uh, it should be ready by now and uh, you go ahead to uh the cafe.
And we're gonna enjoy some uh time of table fellowship together. All right. God bless.
\newpage



\section{腓利比書}
\label{sec:gEFzH4C2B2Y}
\textbf{回到基礎 \#3 (腓利比書 4) - 陳崇基牧師}
\newline
\newline
連結: \href{https://youtube.com/watch?v=gEFzH4C2B2Y}{\texttt{https://youtube.com/watch?v=gEFzH4C2B2Y}} ~~~~ 語音日期: 2024-08-01
\newline
\newline
\hyperref[sec:9TgjTCjHTHQ]{\small{< < < PREV SERMON < < <}}
~
\hyperref[sec:index_chronic]{\small{[返順時目]}}
~
\hyperref[sec:index_scriptual]{\small{[返順卷目]}}
~
\hyperref[sec:WzsPGK70Ghc]{\small{> > > NEXT SERMON > > >}}
\newline
\newline
$^{1}$梁議長早安.
很高興大家一齊來到敬拜.
每一次敬拜的時候.
特別是開頭的時候.
有些音樂和詩歌.
可以幫助我們禱告.
幫助我們向上帝陳明我們心裡面的東西.
我們每個人來自不同的背景.
這個星期可能發生不同的事.
或者昨天今早發生不同的事.
我們覺得在敬拜裡面.
我們有時候一些詩詞.
可以幫助我們去打理一下我們心境.
打理一下我們裡面.
可能一些很混亂的思緒.
以致我們能夠來到上帝的面前.
今天我們繼續去看.
到底在生命裡面有什麼基本.
有什麼生命的基礎.
是我們需要把握和不能夠忘記的.
所以來到這個回到基本的時候.
上幾次我們都講到上帝的話語.
上帝的話語有它的權能.
有它的作用.
可以幫助我們去到上帝面前.
也是我們能夠領教.
領受上帝給我們的恩典.
這個藉著上帝對我們講的話.
這次我們講到另一個基礎.
就是我們向上帝講什麼.
就是祈禱.
我不知道大家有沒有一個比較穩定.
或者常行的一個祈禱的生活.
不需要很複雜.
其實就好像我們對.
好像嬰兒對爸爸媽媽說話.
可以很直接.
很真誠.
但是需要在我們心靈裡面.
一種基礎的習慣.

$^{41}$特別是基督徒.
不要說不是基督徒.
不是基督徒都會覺得禱告.
有時候覺得禱告有一種.
therapeutic effect.
就是有一種治療性的一種效果.
何況我們基督徒在基督裡面.
我們有什麼權柄來主見面前.
來禱告.
所以這次的金句就是.
不是吃素菜.
不是吃肥牛.
就是這個肥立比書.
第四章六至第七節.
我們看看.
他說應當一無掛慮.
將你們所有的高素神.
神所賜那超越人所能了解的平安.
必在基督耶穌裡面.
保守你們的心懷意念.
一段很多人認識的.
甚至我們識背的一段經文.
提醒我們不要擔心.
不要擔心任何事情.
當然他不是說我們坐在那兒.
什麼都不顧.
不要擔心任何事情.
就是不要擔心.
如果我們看得上文下理.
就是我們的飲食.
我們生命裡面一些很基本的東西.
不要擔心任何事情.
任何那些事情.
而是要做什麼呢?.
要祈求一切.
我們看看這段聖經.
到底對我們今天有什麼意思.
特別是對我們今天的生活.
這裡是肥立比書第四章.
什麼叫肥立比書?.

$^{81}$可能你沒看過聖經.
或者第一次看聖經.
我們中間有很多.
可能第一次接觸.
或者初初接觸聖經的.
就是新約裡面的其中一本書.
或者其中一封信.
新約裡面.
新約聖經有27卷所謂的書.
是集成成為我們新約的傳書.
裡面有很多書信.
是當時一個牧者.
寫給當時的教會.
這本叫肥立比書.
有一個牧者叫保羅.
保羅大家很熟悉.
他寫了十幾封書信.
其中一封就叫肥立比書.
為什麼叫肥立比書?.
因為寫給一個地方.
叫做肥立比.
菲律賓.
當時是一個羅馬的諸房城.
裡面有一些小組的家庭教會.
很小很小的.
可能幾個人在家庭裡面聚會.
他寫這封信.
來提醒他們一些東西.
去說一些東西.
如果你看肥立比書這本書.
這封信的話.
你可以當它是一封個人的書信.
甚至你.
當然今天我們不寫信.
還有誰收信?.
沒有吧?.
都是收電郵.
電郵我都不看了.
我們看的是什麼?.
是text.

$^{121}$如果是emoji就霸了.
有沒有試過收到一封信.
是人手寫給你的?.
很寶貴吧?.
有沒有?.
什麼時候?.
什麼時候收到?.
上個世紀吧?.
有人親手寫信給你.
是一些很寶貴很有價值的東西.
當時保羅就是寫這封信給.
肥立比這個城市裡面的.
小型的教會.
那個基督徒.
你看這封書信其實很短.
我們的分段只是四章.
四個chapter.
每章大概二十幾節.
你坐在這裡看肥立比書.
從頭看到尾.
大概用你十分鐘.
十五分鐘就搞定了.
大家都不是文盲吧?.
應該一看就明白吧?.
看肥立比書.
整封書信.
你會感受到是一種很個人的書信.
幾乎可以當成情信.
這樣看.
有沒有收到情信?.
幾乎可以當成情信.
是很個人的一種味道.
為什麼呢?.
因為原來肥立比的教會.
本身是保羅建立的.
保羅有幾次的傳道宣教的旅程.
第二次宣教旅程的時候.
他就去到肥立比城這個地方.
然後那裡有人信主.
一個兩個這樣信主.

$^{161}$然後在家庭裡面聚會.
所以這個教會是他生出來的.
他說他幾乎是他的father.
甚至這個教會雖然少.
但是不斷去資助.
或者幫助保羅宣教的工作.
和他在周圍傳道的工作.
所以是一封很個人的書信.
我們看看肥立比的教會.
或者肥立比書信裡面.
大概有些經文你都知道.
他和保羅.
肥立比教會的信徒.
和保羅之間有什麼親密的關係.
在第四章第十七節.
這裡說什麼呢?.
只有四章而已.
這裡已經是結尾了.
保羅寫什麼呢?.
我並不求什麼饋贈.
只求你們的果子.
不斷的贈多.
歸在你們的帳上.
但我已經如數收到了.
並且有餘.
我已經充足.
因我從以巴弗提.
受了你們的饋贈.
當作輕香極美的香氣.
為上帝所閱立的祭物.
這裡說什麼呢?.
就是會信的時候.
保羅就說.
我已經從這個信差.
叫做以巴弗提.
以巴弗提這個人.
應該是肥立比教會裡面.
其中一個領袖.
或者其中一個會員.
其中一個會友.

$^{201}$他們就奉派.
去將一些的饋贈.
將一些的資助.
可能是衣服.
可能是一些日常生活需要的東西.
送到去給保羅.
當時保羅發生什麼事呢?.
他因為傳福音的緣故.
因為傳道的緣故.
他被所謂監禁.
他就在羅馬京城被監禁.
他正在上訴該殺.
但他還不知道結果是怎樣.
如果你看前文後理.
肥立比書.
其實保羅已經打定輸數.
死都好.
都無所謂.
你記得嗎?.
有這樣的一些說話.
在這封書信裡面.
所以他還不知道情況是怎樣.
大概有兩年.
所謂house arrest.
有羅馬士兵看管著他.
以至他最後上訴.
大概有兩年時間.
他就被監禁在羅馬.
但他們那種house arrest.
仍然可以向外界接待一些貪他的人.
或者去寫一些書信.
所以肥立比教會知道.
保羅這個牧者.
被監禁有兩年的時間.
他們就去給他一些資助.
幫助他.
即是食物等等一些essential的東西.
派這個人叫做爾巴弗提.
由肥立比出發.
去到羅馬.

$^{241}$去將這些饋贈.
將這些幫助給了保羅.
誰知道在書信裡面就告訴我們.
這個爾巴弗提因為路上.
travel的時候.
他就病了.
病到甚至要死.
病到甚至死.
不知道他是傷風感冒還是什麼.
還是covid.
他去到保羅的時候.
差不多死了.
保羅就說很傷心見到他.
誰知道他死不去.
於是乎他就說.
快點將這個爾巴弗提.
快點拆敗回去.
就拿了這封肥立比書信.
這封書信裡面就很有人情味.
對嗎?.
即是你收到support.
然後他就回信去.
support他的肥立比教會.
其實很多書信都是這樣的.
不是fundraising level.
而是reply to一個support.
保羅很多時候寫這封書信.
也有些教導.
所以我們來到最後一章.
第四章肥立比書.
這封書信.
裡面他怎樣去教導這個教會.
由頭到尾第一章到第四章.
都是叫他們同心.
我們有基督的心.
一起去事奉.
一起去傳道.
一起去合一.
到第四章第一節.
我們看一點點這些經文.

$^{281}$這裡說什麼呢?.
他說我所親愛的.
是不是很有個人味道?.
我所想念的.
是不是?.
想不想有人這樣跟你說?.
我很想念你.
我所親愛的.
你們就算喜愛我的冠冕.
我親愛的.
你們應當靠主站得穩.
我勸友阿諜和秦都基.
要在主裡面同心.
我也求你這真實同父一丫的.
要幫助這兩個女士.
因為他們在福音的事上.
曾經與我和格利敏.
和其他的同工一同路苦過.
他們的名字已經寫在生命冊上.
你們靠主要上上喜樂.
我再說你們要喜樂.
要讓眾人知道你們謙讓的心.
主已經近了.
這裡先解釋一下.
有兩個人.
叫友阿諜和秦都基.
是男還是女?.
兩個都是女士.
他們在主裡面要同心.
可能有些爭執.
或者有些不滿等等.
保羅寫一封信說.
不如你們兩個好好地同心.
而至女行間裡面.
你看到他們已經開始.
在研究這個協議.
大家可以合一去工作.
並且有其他的弟兄姊妹.
來幫助這兩個女人.
這兩個是什麼人呢?.

$^{321}$應該是當時教會裡面的一些領袖.
說不定友阿諜的廳是比較大.
所以他們主日崇拜的時候.
就招聚一些弟兄姊妹.
然後有一個叫秦都基.
秦都基可能是比較有錢的女士.
然後他就接待當時的教會.
在他家庭裡面聚會.
就勸他們大家一定要同心.
繼續去為福音工作.
Along with other people.
還有什麼?.
格利敏.
Clement.
還有其他人.
我們一起要同心去工作.
去為福音的合一.
這不是一個好味道的圖畫.
如果我們弟兄姊妹.
都是一起同心去事奉上帝.
這不只是保羅會覺得高興.
他說什麼?.
My joy and my crown.
我的冠冕 我的喜祿.
不只是單單這樣.
我們甚至知道天父都會高興.
天父是遇難.
子女行間你看到的語氣是怎樣?.
很開心的.
My joy and my crown.
你想一下一個被監禁的人.
在囚人士.
他居然可以發出這樣的說話.
在囚人士也不知道將來什麼時候會提堂.
他的結果會是怎樣.
是生還是死.
但是保羅就跟他們說.
一些這樣的說話.
那個語氣是一個很個人的語氣.
還有很多很感人的說話.

$^{361}$四章第三節.
他們的名字都在生命冊上.
我想這個可能就是那個秘訣.
為什麼一個在囚的人.
去到在囚算或者憂慮.
將來他的結果會是怎樣.
能不能夠出獄.
是生還是死.
或者判他死刑等等.
為什麼他可以這樣做.
他們的名字都記在生命冊上.
所謂高興的.
是因為我們在天國裡面.
有這個身份.
這個身份是不會改變的.
More so than.
我們做了多少事情.
這個有點像耶穌曾經跟七十個門徒說過.
你記得有一個故事在路加夫.
耶穌猜那七十個門徒出去.
給他們權柄去趕鬼.
然後傳到那裡.
然後那七十個使徒回來回報的時候.
他們就很開心的.
他們就說耶穌你看我們多厲害.
我們憑著你的名字.
你的大名.
有這樣的權柄.
連魔鬼連邪靈都聽我們說.
但耶穌最後說什麼.
他說不要因為那些邪靈聽你說而高興.
你要因什麼而高興.
因為你們的生命是記在天上.
因著這個而高興.
因著你有這樣的身份而高興.
我們有這樣的名份已經是值得高興.
奧運開始了.
贏輸其實都不是最重要.
你想一下.
有誰是奧運的Team USA的.

$^{401}$沒有吧.
應該沒有吧.
我們高手在民間怎麼知道.
Anyways.
你有份在Team USA.
已經是一個榮譽.
這個就是保羅所說的.
你的名字記在生命冊上.
Team Kingdom有你份.
多大的權柄.
然後他說什麼.
他說主已經近了.
大概可以有兩個意思.
通常翻譯的意思就是.
主很快會回來或者隨時會回來.
即是時間性.
主很快回來.
即是上帝的國度會很快回來.
隨時回來應該這樣說.
Anytime.
It can be today.
亦都可能是明日.
是不知道的.
是何時.
所以有這樣的盼望的時候.
就可以不需要掛慮.
或者生活裡面可能會平安一點.
少一點煩躁.
因為主很快會回來.
第二個的意思.
不只是時間.
可能是位置上.
Spatial.
即是不是Temporal.
是Spatial的意思.
什麼意思.
主已經近了.
或者應該是NIV就是這樣.
The Lord Jesus is near.
即是主在你身邊.

$^{441}$主在你旁邊.
就是Spatial的意思.
即是當我們的生命.
可能會遇見很多不同的困難.
不同的令我們沮喪的事.
但你要相信.
主在你旁邊.
在你身邊.
對於肥立比的教會來說.
更加不需要多提.
為什麼他們需要這樣被提醒.
因為你翻開《士行傳》.
記載保羅去到肥立比傳道的時候.
他是不是傳道得很順利.
不是的.
他被人打.
然後被人監禁.
他和西拉被人監禁.
有很多反對的聲音.
入面一些猶太的領袖.
去到反對他們傳耶穌就是彌賽亞.
連同其他羅馬人也去迫害他們.
所以保羅是受很多苦的.
在肥立比也受很多苦.
所以肥立比建立了教會之後.
你覺得那個教會.
是不是很順利.
一定不順利.
教會是一些很困難的時候.
去到直堂.
然後慢慢成長.
但是要經歷很多風雨和困難.
但是保羅說.
The Lord is near.
不要緊.
基督在你身旁.
是不是.
當你遇見困難的時候.
有人在你身旁.
有一個關心你的人在你身邊.

$^{481}$It makes a whole world of difference.
在你生命裡面有多大分別.
比你一個孤孤單單的去面對好.
我不知道你有沒有病過.
有過醫院.
有人來探望你.
有一個真是關心你的人.
當然也可能是不關心你的人來探望你.
來告訴你.
你要做這個要做那個.
你已經很心煩.
你拿第一點.
第二點.
第三點.
第四點.
第五點.
他還給你一堆點.
當然有些這樣的人.
但是真是關心你的.
靠近你.
去幫助你的.
It makes a whole world of difference.
所以為什麼.
我發現有很多時候.
我們找一個顧問.
找一個人和我們傾訴.
當然有一個智慧的顧問.
是真的幫助到我們的生活.
幫助到我們的生命.
因為我們裡面有很多東西要處理.
怎樣處理.
我們每一個人的生命.
上帝構造我們的時候.
就是需要有另外一個人.
Speak into our life.
你不能夠向著牆.
對著zoom.
對著screen.
對著手機.
然後就會處理到.

$^{521}$不可以.
上帝對我們原本的構造就是這樣.
就是人與人之間要有生命的關係.
有一個wise counselor.
一個有智慧的counselor.
聽我們心底話.
或者處理我們心裡面的情緒等等.
主已經近了.
耶穌基督就在你身邊.
我們的主就是一個這樣的主.
關心你 關心我.
肯聽你禱告.
肯聽你去訴苦.
肯聽你去作愛過.
肯聽你講開心的話.
這個就是我們的主耶穌基督.
所以聖經說什麼呢.
他說應當一無掛慮.
將你們所要的告訴神.
神所賜那超越人所能了解的平安.
必在基督耶穌裡面.
保守你們的心懷意念.
他說既然主在我們身旁的話.
我們就應當可以一無掛慮.
不需要掛慮我們很多時候擔心的事情.
他說不需要擔心任何事情.
禱告祈求差不多是同一件事.
去到向上帝去陳明.
很多時候我們去看禱告的時候.
我們都會覺得我們不懂得祈禱.
有沒有人說不懂得祈禱.
很多時候我們都會覺得.
你叫我突然間一個小組或團契.
我們走出來祈禱我們不夠膽.
因為我們當祈禱是演講.
我們以為祈禱就是演講.
其實祈禱不是這樣的.
我們經常有錯誤的感覺.
你們傳道人或禮仙.
那些詞語是華麗的措辭.

$^{561}$我怎麼說得出這些東西.
我一閉上眼就不知道說什麼.
一閉上眼我就沒有話說.
不知道怎麼解釋.
但最主要祈禱的精神不是這樣的.
以前我神道院的時候有一個教授.
每次課堂開始和課堂結束的時候.
他都會帶我們去祈禱.
他祈禱真的很美.
因為他用英文祈禱.
但他用英文祈禱是用King James的英文.
哇 祈禱真的很美.
Great Adult O Lord.
We Praise Thee Because ….
用舊的英語去祈禱.
但祈禱很美.
但我們很多時候誤會了祈禱.
就是要去組織這些美的詞語.
我們就要去祈禱.
甚至有時候我們口頭禪說.
不完全的禱告.
最尾我們奉主耶穌基督的名求.
不完全的禱告.
我想問什麼叫不完全的禱告.
有誰的禱告是完全的.
什麼叫不完全的禱告.
你們知道聖經裡面最短的祈禱是什麼嗎.
在阿們家裡.
我在想的是耶穌救我.
門徒這樣說.
Help Me.
這個才是祈禱.
祈禱真正的精神是什麼.
是真 是誠 是開放.
等於我開頭的時候說.
好像嬰兒對爸爸和媽媽說話.
他有沒有想過一會兒.
然後說什麼.
他隨便說 你聽到亂說.
但是爸爸媽媽喜歡聽.

$^{601}$Yes.
We want to hear that.
我們和天父之間的關係其實都是這樣.
Being thankful.
去說多謝.
當你說多謝的時候.
你就明白.
你說出來的時候 想出來的時候.
你就明白原來上帝在你心靈裡面.
是有很多美好的事.
在心靈裡面當然有很多瑣碎的事.
但是 at least.
We are thankful.
for the good things God has done.
There must be some good things.
所以怎樣我們去到.
就幫助我們繼續向前.
這段金句 經文.
是一個很寶貴的經文.
去提醒我們.
我們可以藉著禱告.
來到上帝的面前.
That's why 我為什麼去到鼓勵大家.
每一天可能take五至十分鐘.
十五分鐘.
安靜的 有聲無聲的都可以.
最重要是什麼.
是真 是誠 和開放.
來禱告.
除了是一個寶貴的經文之外.
也是一個很有趣的經文.
他說什麼呢.
將你所有的事都告到神.
如果你想.
如果有人說.
如果你要改聖經的話.
我們改聖經.
我們自己比較大逆不道.
如果改聖經改第七節.
你應該怎麼改.

$^{641}$就是神會將我們心裡面所需要的.
都給我們.
這樣我們才會有平安.
他不是這樣說.
你明白我的意思嗎.
我們有很多需要.
我們有很多想要的東西.
如果上帝給我們.
我們就有平安.
你是不是這樣想.
至少我也是這樣想.
我想問你.
你想多點錢嗎.
竟然有人搖頭.
厲害.
你想不想中樂透彩票.
Lottery.
經常看新聞.
誰拋球又中了.
為什麼不拿我的份.
這些都是我想要的東西.
甚至你覺得是需要的.
你想不想你的房子大一點.
不想.
非常好.
大一點也不一定好.
大一點更多東西可以做.
想不想有個好車.
多架車.
是不是這些物質上的東西.
全都給我們了.
這是上帝要這樣做.
全都給我們.
我們就有平安.
對不對.
大家有沒有想過物質越豐富的地方.
物質越豐富的人.
心靈是最貧乏.
窮到只剩下錢的人.
有沒有這樣的人.

$^{681}$大有人在.
這就是意思.
上帝賜給我們的是什麼.
是出乎人所能了解的平安.
給我們.
給我們所有物質上最豐富的東西.
每樣東西都有.
是不是心靈裡面就是好.
不一定.
不會令到你心靈更加有能力.
有能力.
有動力來向前.
反而給你很多東西.
你心靈裡面更加痛苦.
你知道我們現代的人.
心理的狀況已經很亂.
但我們物質上.
我們豐不豐富.
我們真的很豐富.
這就是上帝賜給我們的平安.
在講平安的時候.
你就要想清楚.
不是我們華人想的平安.
菲律賓的信徒怎樣想平安.
你記得我講過菲律賓是一個什麼城市.
一個羅馬的城市.
一個羅馬的殖民城.
Colonial City.
羅馬帝國周圍征服不同的地方.
然後征服了的地方就在那裡吸引.
放一個城市在那裡.
將城市變成一個殖民城市.
殖民城市是做什麼的.
就是靠軍事和文化去強制.
去改變.
甚至去打壓被征服的人.
這個叫Peace.
羅馬當時的帝國.
他們有一個羅馬盛世的說話.
羅馬盛世的概念.

$^{721}$Pax Romana.
講到平安的時候.
羅馬人或菲律賓人會想起.
就是羅馬所講的平安.
羅馬的平安就是打到你殘.
你沒得還手.
這個叫平安.
大家不是很平安嗎.
我不再打架了.
因為我已經打濕了你.
你沒能力還手.
這個叫Peace.
This is Roman Peace.
最近我去organize資料的時候.
我看見一些我搜集的金幣圖畫.
金幣圖畫.
我沒有這些金幣.
不要去我家找.
我找到這些金幣.
在PowerPoint上.
不是實物.
這些是什麼金幣.
這些叫First Jewish Revolt Coins.
第一次猶太人反羅馬革命.
大概是主後70年的時候.
但是羅馬大力鎮壓.
打殘了這些猶太人的時候.
成功了.
他們就鑄造這些金幣.
這些金幣其實是Propaganda.
就是政治宣傳.
所以每個人帶著這些金幣.
銀幣銅幣.
他就知道誰贏了.
誰給你平安.
誰打殘你.
你們看左邊那個.
Triumph.
有一群馬.
上面有一個士兵在吹.

$^{761}$追前追前的人.
就是叫Simon Giora.
那個是其中一個反羅馬分子的領袖.
發生什麼事呢.
把他綁起.
就快要釘十字架.
就綁起.
羅馬軍隊打壓完之後.
他們軍隊乘旱乘烈.
就會回到羅馬京城.
最近在法國巴黎奧運會.
其中一個標誌是什麼.
海船門.
就是一些獨立的城門.
他們就是這樣.
打殘你.
然後將領袖放在他們面前.
然後浩浩蕩蕩的.
然後進去.
然後那些人歡呼喝彩.
因為羅馬贏了.
誰贏了.
就是中間那個.
Vespasian.
其中一個羅馬王.
然後右邊是一個金幣的背面.
背面寫著Judea.
然後有一個女人.
兩隻手臂綁住.
在一棵棗樹的下面.
她就在哭號在哭.
這個就是叫做羅馬盛世.
誰有錢.
他有錢.
你沒有錢.
這個叫盛世.
什麼叫盛世.
你被人打到無法還手.
右邊那個就是你.
被人打壓到坐在樹下.

$^{801}$沒了半條命.
這個叫羅馬盛世.
這個叫Roman Peace.
但保羅在這裡說什麼.
他說是上帝所賜的平安.
不是羅馬帝國給你的平安.
不是地上的帝國給你的平安.
他說的Peace是在基督裡面找到的Peace.
The Peace from God.
Peace given by God.
就是上帝.
慈愛的上帝給你的和平.
你可以想像一下.
這種Peace才能幫助我們生命.
能夠保護我們.
所以他說.
上帝所賜的平安.
是超乎人所了解的.
或者應該這樣翻譯.
超乎人所能夠想事情的.
很多時候我們很多焦慮.
因為我們經常有很多原因.
我們經常要想事情.
經常要想來想去.
昨晚睡得著嗎.
當然睡不著.
你想起什麼.
錢啊.
其他東西啊.
好煩啊.
超乎你所能夠想得通的東西.
那種平安在基督裡面保守你們的心懷念.
菲律賓是羅馬的駐防城.
經常有軍隊守著.
這叫羅馬平安.
有軍隊在就有平安.
因為沒有人作反.
但是你就被人打殘.
用這個詞.
We'll guard your heart.

$^{841}$好像士兵那樣去看守著你.
我記得我說過這個例子.
好幾年前新冠肺炎剛開始的時候.
我們鄉下.
是不是有很多國民警察.
你覺得平不平安.
表面上是很平安的.
平安過你每個人拿著槍到處掃.
那些不是叫平安.
不是平安.
但是他說什麼.
The peace of God.
We'll guard your heart and mind.
在基督裡面.
保守著守防著你的心懷念.
當你心懷念健康明白的時候.
當你有安慰有平安.
真正上帝給你平安的時候.
你的心靈就有動力.
無論外面發生什麼事都好.
His peace will guard you.
好像軍隊那樣防守著你.
你就不需要怕.
因為你知道上帝是take care你.
我家裡有些舊車.
我兒女現在開車.
我女兒Rachel.
她很有趣.
她現在很喜歡開車.
之前很不喜歡開車.
現在很喜歡開車.
一個人開車上班等等.
自己會泊車.
她說OK.
起碼大家都平安.
她很喜歡說一件事.
她回到家.
她說爸爸.
The gas is low.
她說沒有油.

$^{881}$她的意思是什麼.
你去入油行不行.
因為你已經剩下四分之一.
是時候你入油了.
我真的很佩服她.
我覺得有這種信念.
我會take care她.
只要她take care我.
我就永遠都有油.
我有用.
能夠開車.
有油用.
她對我的信靠.
我非常admire.
她很chill out.
The gas is low.
Text我.
你要工作了.
當然你要學.
但這就是意思.
我們對上帝的信靠.
The gas is low.
但上帝給我們出人意外.
或者人所能夠理解的.
這種平安.
以致我們能夠應對生活中所有的事.
我們over thinking不到.
我們真的有這種趨勢.
我們經常要over thinking.
想多一點.
我們以為想多一點就搞定.
但最重要的是.
保羅說什麼.
你想想是一個監犯跟你說.
神會賜給你出人意外的平安.
這種生命才能夠繼續向前.
奧運會開幕.
我記得最精彩的地方.
當然有很多精彩的地方.
我記得Celine Dion.

$^{921}$她唱的那首歌.
真是精彩.
Celine Dion最近這幾年.
她的生命中有很多挑戰.
家人去世.
她患上了僵人症.
不能動.
她跪了好幾年.
我一見到她.
真是樂壇天后.
生命繼續可以跪下去.
你看看這些奧運的選手.
是不是一朝一夕.
他們就這樣自己去.
可以有這樣的動力.
不.
在他們的生命中不斷有基礎.
不斷有訓練.
不斷向前.
以至他們能夠有那種肌肉.
能夠向前.
但我們每一個都是這樣.
最後這裡說什麼.
應當一無瓜葛.
將你們所要的告訴神.
神所賜那超越你所能了解的平安.
在基督耶穌裡面保守你們的心懷.
養成一個禱告的習慣.
因為知道上帝是聽我們的禱告.
同心大同.
上主我多謝你.
因為你是那位去到知平安的神.
在我們心底裡面.
是超越我們能夠去理解.
超越我們能夠想透的.
主幫助我們.
在我們很多焦慮,很多壓力的生命裡面.
我們就近你.
因為主你就在我們身旁.
幫助我們每一天.

$^{961}$我們有一個這樣的習慣.
在你面前來聽.
在你面前來禱告.
有聲無聲的,真誠的.
向你彈開.
讓你給我們力量,給我們平安.
在這個星期面對不同的困難和挑戰的時候.
你給我們這個平安.
能夠勝過世界.
祈禱奉尊名.
好,願主的平安常與大家同在.
希望下星期同心一起致奉.
\newpage



\section{12:1-8}
\label{sec:WzsPGK70Ghc}
\textbf{Back to Basics \#6 (Romans 12\_1-8) - Rev. Mark Chan}
\newline
\newline
連結: \href{https://youtube.com/watch?v=WzsPGK70Ghc}{\texttt{https://youtube.com/watch?v=WzsPGK70Ghc}} ~~~~ 語音日期: 2024-08-06
\newline
\newline
\hyperref[sec:gEFzH4C2B2Y]{\small{< < < PREV SERMON < < <}}
~
\hyperref[sec:index_chronic]{\small{[返順時目]}}
~
\hyperref[sec:index_scriptual]{\small{[返順卷目]}}
~
\hyperref[sec:l_mjl7q5DSI]{\small{> > > NEXT SERMON > > >}}
\newline
\newline
$^{1}$- Alright guys, let's turn to our scripture today..
And may I invite you to go to your phone app.
or to your Bible..
We're going to continue our series on basics, okay?.
Let's master the basics..
There are so many basic exercise or habits.
as a question we need to master..
So I introduced to you many weeks ago.
the verse of the week..
And this is the verse of last week..
Remember that?.
Alright, can we recite it together?.
Let's do it in English..
Don't worry about anything, instead pray about everything..
Tell God what you need and thank Him for what He has done..
Then you will experience God's peace,.
which exceeds anything we can understand..
His peace will guide your hearts and minds.
as we live in Christ Jesus, alright?.
So we got the Chinese translation for many of you..
This is the verse of last week..
And today we're going to introduce another verse..
However, we're going to go through the passage of today..
My sermon is taken, centered on this passage,.
this famous passage from the book of Romans..
Romans chapter 12, verse one..
"Therefore I urge you, brothers, in view of God's mercy,.
"to offer your bodies as living sacrifices,.
"holy and pleasing to God.".
This is your spiritual act of worship..
"Do not conform any longer to the pattern of this world,.
"but be transformed by the renewing of your mind,.
"and you will be able to test and approve.
"what God's will is, His good and pleasing perfect will.".
And that is the first two verses,.
and we continue, verse three..
"For by the grace given me, I say to every one of you,.
"do not think of yourself more highly than you ought,.
"but rather think of yourself with sober judgment,.
"in accordance with the measure of faith God has given you,.

$^{41}$"just as each of you has one body with many members,.
"and these members do not all have the same function.".
So in Christ, we who are many form one body,.
and each member belongs to all the others..
"We all have different gifts,.
"according to the grace given us..
"If a man's gift is prophesying,.
"let him use it in proportion to his faith..
"If it is serving, let him serve..
"If it is teaching, let him teach..
"If it is encouraging, let him encourage..
"If it is contributing to the needs of others,.
"let him give generously..
"If it is leadership, let him govern diligently..
"If it is showing mercy, let him do it cheerfully.".
Now that is from the book of Romans, chapter 12..
And that is, as I said,.
is one of the well-known passages in the Bible,.
especially in the New Testament..
And there is some background we need to understand.
as we read this passage..
Now, all of the Bible passages were written.
to a specific context or specific need of a congregation..
This is not new to us,.
but as long as we get more familiar.
with the background of the Bible,.
then we understand what God is saying, right?.
Then we understand what this passage is all about..
So some of the background..
This letter is written by Paul, the apostle Paul,.
the pastor Paul, a missionary,.
founding churches in many cities in the Roman Empire..
And this letter is written by him, by Paul,.
to the Romans, so-called the Romans,.
because those are the Christians.
and house churches in the city of Rome..
And as you know, Rome was the center of Roman Empire,.
of the center of the world in Paul's day..
And the relationship between Paul and the Roman Christians,.
the Christians, the house churches in Rome.

$^{81}$are somewhat interesting because Paul did not establish,.
did not plant the church in Rome..
In fact, he had never been to Rome.
until the end of the Book of Acts..
And when he wrote this letter to the Roman Christians,.
he had never been there..
He said he wanted to visit Rome someday on his way to Spain..
Well, we don't know whether Paul.
has traveled to Spain or not..
However, he wrote this letter,.
and at that time, he did not come to Rome,.
and he never planted the church in Rome..
But there are Christians,.
there are house churches in Rome nonetheless..
And we don't know who brought the gospel.
to the people in Rome,.
and why there were some house churches.
being established in Rome..
So who brought the gospel to Rome in the first place?.
Remember chapter two of Book of Acts,.
very familiar chapter, Pentecost, right?.
When the Holy Spirit descended on the apostles,.
and they began to preach the gospel,.
and Peter begins to preach to the people,.
and 3,000 people came to faith, right?.
And they were baptized by water,.
and they became Christians..
And on the day of Pentecost,.
there were many Jews in Jerusalem..
And the Book of Acts tells us.
there were many Jews come from other places.
of the Roman Empire..
Some come from today's Turkey,.
some come from today's Egypt,.
some come from even Rome,.
because in Acts 2.10, it tells us.
there were some Jewish people coming from the city of Rome,.
from the capital city of Rome..
And they were visiting Jerusalem for the festival,.
for Pentecost, and for Passover before Pentecost..

$^{121}$And they just happened to be there..
And the Holy Spirit descended upon the apostles,.
and they began to preach the gospel,.
and many of them became Christians..
And I think surely they also came to faith,.
and they brought the gospel back to their home city in Rome..
So I think they found the house churches in Rome..
And in the very beginning period of the church in Rome,.
we can imagine the Jewish people.
and their Gentile Christians gathered together, right?.
So that was the picture,.
the reality of house churches in Rome..
Again, it's by some sort of lay Christians,.
white Jewish Christians,.
they brought the gospel back to Rome,.
and they probably attract other people,.
and some of the Gentile Christians..
And they gather together in different houses in Rome,.
and that's how house churches in Rome began..
But, but, but in AD 49, AD 49, okay?.
So in your mind, imagine there's a timeline, okay?.
There's timeline..
You put cross there, it's about AD 30, AD 30, all right?.
And Jesus died and rose again,.
and he ascended to heaven, right?.
And the Holy Spirit came on the day of Pentecost..
Now, that was around AD 30..
But in AD 49, everything changed,.
especially in the city of Rome..
Why?.
Because at that time, Roman emperor named Claudius,.
and he issued a decree to ban all the Jews.
from the city of Rome..
So, so Claudius, the emperor,.
kicked out all the Jewish people from the city of Rome..
And as you can imagine,.
some of them are being kicked out of Rome.
were definitely Jewish Christians,.
and they originally belonged to the house churches in Rome,.
right?.

$^{161}$And actually we find that in the book of Acts, chapter 18..
Here, the map clearly shows where Rome is, right?.
And Paul is writing the letter to the Romans from Corinth,.
right?.
And there were house churches.
mixed with Jewish and Gentile Christians in Rome,.
but AD 49,.
Claudius commanded all the Jews to leave Rome..
And we find that in Acts 18..
Acts 18, what happened here.
was that Paul was doing mission in different cities,.
and he came to Corinth,.
and he found a Jew there named Aquila,.
and he was a native of Pontus..
Pontus was today's Northern Turkey..
And this guy is called Aquila,.
and he was a Jew,.
and Paul met him because he come from Italy.
with his wife Priscilla..
So this couple was being kicked out of Rome,.
and they were in Corinth,.
and they met Paul,.
and Paul preached to them..
And Acts 18 explains,.
because Claudius, the emperor,.
had commanded all the Jews to leave Rome, okay?.
So that's pretty interesting, AD 49, right?.
Remember, the house churches in Rome.
were mixed with Jewish and Gentile Christians, right?.
They were all calling the name of Jesus,.
and they worshiped together..
But AD 49, after Claudius kicked out all the Jews.
out of Rome,.
overnight, the church in Rome.
became virtually 100\% Gentile..
Right, it's not hard to imagine, right?.
Kick out some of the Jews,.
and then what you left in the house churches in Rome.
were Gentiles..
And the Gentile Christians, naturally,.

$^{201}$would take over the leadership of the church..
And so they were minority,.
and now the Jews are being kicked out..
The Jewish Christians left the city..
They became the majority..
And by the time Paul writes from Corinth,.
writing the letter of Rome,.
it was about AD 57, eight years later,.
and some of the Jews were allowed.
to go back to their home city, Rome..
But when the Jewish people came back to the church,.
and they find that the church was dominant by Gentiles,.
right, now before AD 49,.
the Jewish Christians were the majority..
After AD 49, they became the minority..
And you can imagine the tension,.
the social tension, or the leadership tension.
created by such a situation..
And by and large,.
Paul is writing the letter to Romans.
to address some of the problems or tensions.
that happened in the church..
You got that?.
All right, it's very easy, right?.
Now, you wanna do a pop quiz?.
All right, so remember, AD 49 is the watershed.
for the house churches in Rome..
And all of a sudden, the Gentiles became.
the dominant people, and they take over the leadership,.
and they despise the Jews who came back..
And Paul is writing this letter.
to address some of the issues.
or some of the tensions and debates inside the church..
There are two points we can find in this passage..
So the letter to the Romans.
is the longest letter in the New Testament,.
the longest letter..
And Paul used the first 11 chapters to explain,.
to emphasize that Jesus came to save Jews and Gentiles alike..
At the foot of the cross, the ground is level,.

$^{241}$and all have sinned, Paul said, chapter three..
Everyone falls short of the glory of God..
But in chapter 10, Paul proclaims,.
everyone who calls on the name of Jesus will be saved,.
whether you're a Jew or you're a Gentile,.
because the ground at the foot of the cross is level,.
and that is the mercy of God, right?.
We never deserve the grace of God, right?.
Only because what Jesus has done for us,.
and that is the mercy of God..
And in Christ, Paul reminds us,.
we are one new people now, right?.
Now, you may center your identity on your ethnicity..
I am a Jew, I am a Gentile, right?.
I am from a certain ethnic group..
But Paul is saying, when you come to the cross,.
when you receive the mercy of God,.
undeserved, and in Christ, we are now one people,.
and we have the new identity in Jesus Christ..
And that is the most important identity..
And that's why the church is the new people of God..
And we come to chapter 12 of the book of Romans..
Paul is saying, I urge you, brothers,.
in view of God's mercy, right?.
Now, you have heard me explain for 11 chapters.
of the mercy of God on Jews and Gentiles alike..
And Paul is saying, chapter 12 verse one..
Therefore, I urge you, brothers, right?.
Or brothers and sisters, siblings..
You know, this is one of Paul's most favorite term.
to describe Christian relationship..
He could have used the temple,.
he could have used the people,.
but the most favorite term Paul uses.
to describe our relationship as Christians.
is siblings, brothers and sisters in Christ..
And so often, we see each other.
based on our common background, right?.
Common language, race, or interest..
And sometimes we fail to see each other in the church.

$^{281}$on the sole basis of God's mercy and grace, right?.
We come here, we call each other.
sisters and brothers in Christ.
because God has done something in your life, in my life..
God has created something new..
And look around, right?.
This is something new, this is something unheard of,.
or you have never thought about,.
that God is putting you and me together.
in this new society, right?.
New people..
He created a new society of siblings,.
our people from all walks of life, right?.
We're so different, we're so distinct, right?.
But God is calling everyone in Christ,.
Jews and Gentiles, right?.
All kinds of people, East and West, right?.
Coming together to be a people, right?.
I have been to different churches..
I've always been amazed by, you know,.
how the people are so different,.
but God put us together, right?.
I'm so thankful I'm not the only guy there,.
and no one looks like me, right?.
Everyone looks different, right?.
Everyone comes from a different walks of life,.
so that is the mercy of God..
And Paul continues to tell us,.
because of the mercy of God that we receive.
as brother and sister, we need to offer your bodies.
as living sacrifices to God..
And Paul, I think he has in mind.
a different and new sacrifice, right?.
A living sacrifice, right?.
All the sacrifices are dead, right?.
Dead animals..
And Paul is saying, as we come together,.
as we receive the mercy of God,.
it is not offering animals anymore..
We offer ourselves, right?.

$^{321}$And we offer ourselves as living beings,.
living bodies, how we live and how we think,.
how we speak, how we act, how we relate to each other..
It's like a sacrifice to God, right?.
How we build our family, how we guide our children,.
how we serve, how we participate in fellowship and worship.
is also a sacrifice..
It's like a worship to God, right?.
And here the English says spiritual act of worship..
I think a better translation is a reasonable,.
reasonable act of worship..
Or even a better translation is gospel field..
Gospel shaped act of worship..
We come here and it's not because it's convenient to come..
We come here because of Jesus Christ, right?.
And he continues to say, do not conform any longer.
to the pattern of this world..
The way of life in this world.
means the way of life in the Roman empire, right?.
The Roman way is I dominate you, right?.
I have power, you have no power..
So the powerful is ruling the powerless..
The strong is ruling against the weak, right?.
And I think this phrase, do not conform any longer.
to the pattern of this world,.
I think the best paraphrase is from a very,.
is from an older English paraphrase Bible called J.B. Phillips..
So J.B. Phillips was a pastor and translator of the Bible.
and he translated New Testament in contemporary English..
He translated this phrase, do not conform any longer.
to the pattern of this world as, listen up,.
do not let this world around you squeeze you.
into its own mold, but let God remove your mind from within..
J.B. Phillips, yeah, pretty good, right?.
Do not let this world squeeze you into its own mold, right?.
Do not let the pattern of the world,.
the Roman way or the secular way to mold you..
Don't let it squeeze into the worldly mold,.
but let God shapes your mind and mold your mind from within..
There's constant pressure in this world to shape us,.

$^{361}$to squeeze us, to be like the world..
But as Christians, as we receive the mercy of God,.
the true way of life is conforming to Jesus Christ himself..
So point number one, I have it is we,.
each and every one of us together,.
we are living sacrifice to God..
We are worship..
We are like an act of worship,.
not only by singing in the first 15 minutes, right?.
We sing, how many songs do we sing today?.
Three, right?.
Three..
We sing the songs and we follow the PowerPoint..
It's not just that..
Worship is not just that..
Worship is our whole being..
We live together..
We act together as a pleasing sacrifice,.
as a pleasing worship to God..
And let the word of grace shape us.
instead of the words from the world, right?.
That's the two verses in the beginning..
And three to eight, Paul reminds the church,.
reminds us who we are and what we are to do..
Basically, he said, we are the body of Christ..
We are the new body of Christ..
So he said, "By the grace given to me,.
"I say to every one of you,.
"do not think yourself more highly than you are,.
"but rather think of yourself with sober judgment.
"in accordance to the measure of faith God has given you.".
So God has given us different levels of faith,.
different levels of maturity, different levels of abilities..
Verse four, "Just as each of us has one body.
"with many members,.
"these members do not have the same function.".
And the analogy goes that in Christ now,.
we form one body and each member belongs to all the others..
That's a wonderful verse, right?.
In the body of Christ,.

$^{401}$our relationship to each other has changed..
It's defined by God,.
none other than God..
And God defined our relationship with each other..
And God has given us different measures of faith gifts,.
right, some special abilities,.
some special calling to all of us,.
whether we are Jew or Gentiles,.
whether you are high up in the social status.
or low in the social status,.
with status or without status..
Christ has redefined our relationship..
To the Romans, to the first audience of this letter,.
body, when they say body,.
it was a political term.
because politicians in the Roman Empire.
and in the Greek world,.
they like to use the analogy of body.
to describe the citizens of the city.
or citizens in the empire..
You belong to body, to one body,.
the politician would say..
That's why you need to be united.
and you have to vote for me, right?.
You have to support me..
That's how the Roman politicians would say.
when they use the term body..
They're talking about political unit,.
body of citizens..
To the Romans, as I said,.
the rulers were dominant over the ruled..
The slaves need to obey the leaders..
The men are above, on top of women, over the slaves..
And that will be peace..
To the Romans, this is Roman peace..
But Paul is using the same term in a subversive way..
He's using the same term, body,.
and to tell us, hey, in here,.
in this new community, this new society,.
relationship is defined and configured differently..

$^{441}$Verse four, he says,.
"One body with many members and different functions..
"Each member of the church,.
"each member here is gifted in some way.".
Aren't you amazed, you know,.
some people are gifted in very different ways, right?.
You may be gifted in cooking, right?.
I'm certainly not gifted, right?.
I'm certainly not gifted..
Before, my son is spending the summer here,.
and because of him, we have good food..
You know why?.
Because I approach food in a very practical way, right?.
I don't care whether it tastes good or not, right?.
As long as it fills my stomach,.
and I continue to work, right?.
I approach food, it's a very, very, very.
despising way to him, right?.
But he makes good food,.
and he's so gifted in that way, right?.
And some of people are gifted in leading,.
and some of us are gifted in teaching..
Some of us are gifted in caring..
And many different abilities, many different gifts.
that the Bible says has different function, right?.
And the word function, in verse four,.
you see that these members do not all have the same function..
The Greek, the original Greek term is praxis..
Praxis, that's how we get the English word practice, right?.
So Paul is saying members do not all have the same praxis,.
do not all have the same practical function..
And this word is telling us.
not only we receive some spiritual gifts,.
we have to use them..
Spiritual gifts or spiritual abilities.
are not to be owned or possessed by the individuals..
Each one of us needs to discover our spiritual gift..
We discover our spiritual gifts,.
and we need to develop,.
we need to grow our spiritual gifts.

$^{481}$and use them in praxis, in practice,.
to use it for the good of the church,.
for the good of others..
No one, no one, no question in the church is an audience..
Every one of us is performing something..
Every one of us is doing something.
and performing our practical function..
And that's what Paul meant when he continues to say,.
"We have different gifts," verse six,.
"according to the grace given to us.".
A man may be gifted in prophesying, right?.
Now again, let me remind you,.
prophesying in the Bible, in biblical terms,.
is not predicting whether it will rain tomorrow.
at 3.30 p.m..
It's not predicting..
Prophesying is proclaiming the word of God,.
the authoritative word of God..
That's what the prophets did in the Old Testament..
And that's what people who are gifted.
in prophesying will do..
They would preach according to scripture.
and let him use it in proportion of his faith..
Verse seven, "If you're gifted to serve, let him serve.".
By the way, this church is so blessed.
to have many people who are gifted in serving..
You know what it means?.
You know what it looks like?.
People always work in the background, right?.
Putting things together, setting things up, right?.
And they don't talk, right?.
But they use their hands and feet to serve the church..
If it is serving, if it is the spiritual gift of serving,.
let him serve or let her serve..
If it is teaching, let him teach..
If it's encouraging, let him encourage..
Contributing, let him give generously..
If it's leading, let him lead diligently..
Showing mercy, let him do so cheerfully..
And that is the beautiful picture of the church..

$^{521}$And, you know, Olympic Games has been showing on TV, right?.
I hope you have been watching the Olympic Games, right?.
Yeah, you just can't help it..
You have to watch it, right?.
Yeah, it's 'cause you're bombarded by, you know,.
different games and replay, right?.
And I spent some time watching the Olympic Games,.
and I'm just amazed how people get gifted so differently,.
right?.
Some are great swimmers, right?.
You can't tell them to climb, right?.
They are created, they're made, they're gifted to swim, right?.
Some are good at track and field..
Some are good at gymnastics, right?.
You can't tell them to do something else, right?.
And that is the picture Paul is trying to paint to us..
'Cause he says, in Christ, we are a new people,.
and we are gifted in many different ways..
And you and I are called to build up the body of Christ,.
the church..
Use our gifts, right?.
You have to find your own gift..
And you discover your own gift, praise God, right?.
That does not mean you don't need to do your lessons..
You still need to improve, you still need to develop, right?.
You still need to use them..
I remember in the early years of my Christian walk,.
I started to discover the gift of teaching, right?.
Because, you know, that's one of my passions, right?.
I like to study the Bible, I like to know more about God..
And I wanna tell other people, right?.
I remember the early years, I stumbled..
I stumbled a lot, because I don't know how to put.
the information or the materials together in a helpful way,.
to help the church, to help our fellowship,.
to help our brothers and sisters, right?.
And it took me some time to practice it..
It took me some time to take the opportunity,.
seize the opportunity to help other people..
And that's what I'm saying, discover your gifts.

$^{561}$and develop them and deploy them and to use them, right?.
Paul is saying to us, we are a new people of Christ..
We are gifted in many different ways..
Let's build up the body of Christ, right?.
If you find that you are not gifted,.
you need to tell yourself, that's not true..
You are always gifted in some ways, right?.
Our church is blessed by many people.
who thought they are not gifted,.
but what they do, what they practice is so important to us..
And that is praying for the church, right?.
There are some elderly members in our community,.
they faithfully pray for us..
They faithfully uplift the church.
and pray and support the church, right?.
You have the gift and you have to use it..
You have to use a gift to build up the body of Christ..
That brings us to the verse of the week..
And let me conclude with this verses.
taken from the Book of Romans 12, verses four and five..
Why don't we read this verses together, right?.
Just as each of us has one body with many members.
and these members do not all have the same function..
So in Christ, we who are many form one body.
and each member belongs to all the others..
Let's all pray together..
Father, we wanna give you thanks.
for each and every one of us here..
Yes, we come from different walks of life.
and we speak different languages,.
we have different interests and we have different passions..
And yet in Christ, we are one body..
And yet in Christ, in the mercy of God,.
we belong to this new family..
We can call each other siblings,.
siblings in Christ, sisters and brothers in Christ..
What an amazing reality that we can live in..
Thank you for the brother and sister in this room,.
for the siblings in this church,.
that we are being called into this community,.

$^{601}$this family to be one, to be one body..
And God, we just pray that you will help us.
to break down all the walls that separate each one of us.
and help us to be united, united in Christ Jesus..
And thank you for the gifts you have given to us..
Help us use them for the benefit of others,.
to pray for the church, to build up the church,.
to serve other people..
Thank you for this new identity.
and may this new identity continue to remind us.
that we belong to the body of Christ..
What a glorious truth..
And we can use our gifts to serve,.
use our gifts to worship..
Again, we wanna give you thanks and praises.
in Christ and pray..
Well, may God's peace be each and every one of you.
and I wanna see you next week, 11.30.
and let's continue to worship God..
God bless..
(gentle music).
\newpage



\section{}
\label{sec:l_mjl7q5DSI}
\textbf{Back to Basics \#8 (Corinthians 12) - Rev. Mark Chan}
\newline
\newline
連結: \href{https://youtube.com/watch?v=l-mjl7q5DSI}{\texttt{https://youtube.com/watch?v=l-mjl7q5DSI}} ~~~~ 語音日期: 2024-08-23
\newline
\newline
\hyperref[sec:WzsPGK70Ghc]{\small{< < < PREV SERMON < < <}}
~
\hyperref[sec:index_chronic]{\small{[返順時目]}}
~
\hyperref[sec:index_scriptual]{\small{[返順卷目]}}
~
\hyperref[sec:nNDuojAF0eE]{\small{> > > NEXT SERMON > > >}}
\newline
\newline
$^{1}$Alright, thanks Kelly..
Did you get a bottle of the sauce?.
- I got it..
- The XO sauce, right?.
Yeah, so it would be great, alright?.
So let's pray together..
Father, we come before you knowing that.
you hear our prayers and you are with us..
Whether we are going through up or down,.
you are still our God and you are indeed our same God..
And yesterday, today, and tomorrow,.
our Lord Jesus Christ remains the same.
and he is our faithful God and our faithful friend..
We uplift our burden, we come before you.
and open our hearts knowing that.
you are the God who hear our prayers..
So we come before you to confess our sin,.
confess our shortcomings,.
knowing that we are weak without you.
and we need your grace..
As we worship today, we come and ask.
for your grace to draw us near.
so that we get to know you more..
Know what you desire and grant us.
the heart to follow you, to obey you..
We pray for our city as the DNC is coming.
into the city, it's gonna be a lot of traffic.
and there are a lot of different voices and controversies..
I pray for safety and peace for our city,.
for our neighbors, and I just wanna.
come before you in prayer as well.
as the election is coming..
There are a lot of different controversies.
and a lot of issues for us to think about.
and may your church, may your people.
be desiring your kingdom because.
only your kingdom is eternal..
No matter who we vote for, no matter which side.
are we on, we know that we need to stand.
for your kingdom, for only the kingdom.

$^{41}$of our Lord Jesus Christ stands forever..
And that's why we continue to build this church..
That's why we continue to serve you.
and come to worship you because your people.
will remain forever and your kingdom,.
your reign remains eternal..
And help us to fix our eyes on you.
knowing that your reign and your kingship.
is forever and let us worship you.
and serve you forever in Christ's name, we pray..
Well, if you have your Bible, turn me to.
the book of 1 Corinthians..
We're going through this long series.
and today is the conclusion..
This is the eighth week and we introduce.
this verse of the week and we remember.
last week I introduced this verse.
from the book of Romans, chapter 14, verse eight..
And shall we read together?.
Romans, chapter 14, verse eight..
All right, let's do it together in English..
If we live, it is to honor the Lord.
and if we die, it is to honor the Lord..
So whether we live or die, we belong to the Lord..
That's the book of Romans, chapter 14, verse eight..
So often we forget that our life belongs to God, right?.
You're redeemed by the blood of Jesus Christ..
We are all saved by the grace of Jesus Christ.
and our life and 24/7, all of our time,.
belong to God, right, and none other..
He has the ultimate authority..
He has the ultimate lordship and leadership in our life.
and we should yield to his leadership..
Well, every time we struggle with some of the problems.
in our lives, we need to go to God..
Every time we taste the grace of God, right,.
God answer our prayers and help us.
in every single different way,.
we need to come to thank God because our life.
belongs to him, as it says here..

$^{81}$We belong to the Lord, right?.
And that's why we need to understand.
there are basics in our Christian life.
and we went through different elements of the basic,.
the habits, right, of Christian life..
We talk about the Bible, we talk about we need.
to read the word of God, we need to pray.
and we need to build our church life..
We need to get into the habit of stepping into.
and plug ourselves into a community.
and so we want to conclude this series by encouraging.
each and every one of us to serve,.
to use our spiritual gifts..
We went through a little bit in Romans chapter 14.
but today we're going to focus on this topic of serving..
As a member of a church, as brother and sister in Christ,.
we are called to plug into a local community,.
belong to the spiritual family and live together.
and build a church together..
And we come to the book of 1 Corinthians.
and first of all, some of the background, okay?.
So a background about this book called 1 Corinthians..
And we know from the book of Acts.
that Paul found this church..
He planned this church..
It was on his second missionary journey.
and it was about AD 48, 49 and he was in Corinth.
and if you read the story in the book of Acts,.
Paul actually preached the gospel.
and a few of the people from Corinth.
came to believe in Christ..
And then secondly, we need to understand,.
Paul returned to his home base in Antioch.
in the east side, you see that?.
You see that?.
In the Mediterranean east coast,.
there's a port called Antioch.
and down there is Jerusalem, right?.
So there were few missionary trips.
that Paul and his coworkers made.

$^{121}$to preach the gospel to the west, all right?.
The first one, the second one,.
he found the church at Corinth.
and then in his third missionary journey,.
it was about AD 50s,.
he went to today's Turkey.
and he stayed at Ephesus for three years..
And here in Ephesus, in today's western Turkey,.
he heard about the situation.
in the house churches of Corinth..
And remember a few years before,.
Paul planned this church..
Paul came to this place.
and create a church and started a church.
and the church continued to meet..
And while Paul was in Ephesus.
in his third missionary journey,.
he heard about some of the issues and problems.
and conflicts within the house churches of Corinth..
And that's why he wrote one of the letters,.
what we call now First Corinthians..
And if you understand, you know, First Corinthians,.
you read the letter,.
you probably have some impression.
of what that letter is all about..
Paul tackled a lot of problems of the church..
I always think that, you know,.
if my church is not like Corinthian church,.
then, you know, I should thank God, right?.
Because Corinthian church is full of problems,.
full of divisions and full of tensions..
And in this long letter, 16 chapters,.
Paul wrote from Ephesus and send the letter to Corinth.
and to teach the Christians there,.
teach the house churches there,.
how to be united, how to worship God.
and how to take care of some of the problems,.
moral problems or, you know,.
the divisive problems that they have..
And we come to chapter 12..

$^{161}$And chapter 12 talks about three things..
I wanna focus on this, these three points today..
And in chapter 12 of First Corinthians,.
Paul is saying God desires,.
God loves unity..
God would like to see the church,.
the people in the church to be united,.
not separated, not divided..
And that's the first main point..
God loves unity..
God also love diversity in the church..
That's why every single member has his role, right?.
And thirdly, God orchestrates everything,.
everyone together..
So let's take a look at the first point..
God loves, God desires, God delights in unity..
When people praise God together,.
when people serve one another together,.
when people learn to love one another.
and hang in there and to build a community together,.
God loves it, God loves unity..
Chapter 12, verse 12..
The human body, Paul says,.
he's using an analogy, a metaphor for the church..
He says the human body has many parts,.
but the many parts make up one whole body..
So it is with the body of Christ, the church..
So he's using the analogy of human body, right?.
Think about your physical body, right?.
We have different parts, right?.
And this is a perfect analogy for church,.
for us, for urban voice, for our church..
Verse 13, he says, some of us are Jews,.
some of us are Gentiles..
So culturally, we're different..
Some of them are Jews, some of them are Gentiles..
And then some of them are slaves or bond servants,.
and some of them are free, free men, right?.
They are not slaves..
They probably own slaves, right?.

$^{201}$And by the way, slavery in the first century.
is not the slavery in the recent history.
of the United States..
It's quite different..
In any case, some of them are slaves.
and some of them are free, all right?.
So socioeconomic status, everyone is different, right?.
But they are put together in the church..
But we have all been baptized into one body, the church,.
by one Holy Spirit, and we all share the same spirit..
And Paul here want to emphasize the unity that we have..
The unity that we have, that we are united together.
as sisters and brothers in Christ in the same church..
It's not because we share the same interests.
or we share the same personality or background..
He said we all share the same spirit, right?.
That's NLT, New Living Translation..
You probably have a more literal translation.
that says we were all made to drink of one spirit..
That's interesting, drinking the spirit, right?.
Because in the New Testament,.
spirit is often analogous to water..
The spirit is like water..
So we all made to drink one spirit..
So the point is we all share one cup, right?.
We all share the same fountain..
That is the Holy Spirit..
When you became Christian,.
when you came to saving knowledge of Jesus Christ,.
you are sharing the spirit.
because the spirit immediately lives in you, right?.
Lives in you and lives in me.
if we call on the name of Christ..
So we all share the same Holy Spirit..
And the church is always one..
And Paul is trying to emphasize this,.
although the church is made up of many different people,.
different walks of life,.
but we all share one thing in common..
We belong to this one body, sharing one spirit..

$^{241}$And God loves this.
because that's how God set it up, right?.
Sometimes we come to the church, we may complain..
There are people who are unlike us, right?.
The people who are so different from us..
And it takes time and it takes a lot of different trials.
and errors to relate to one another, right?.
To befriend with one another..
And we don't like it, right?.
We always being drawn to the people.
who have the same chemistry, right?.
But that is not how God designed the church..
God loves unity..
And the unity is based on one thing,.
the body of Christ, the same spirit..
God is delightful when he sees different people.
united in Christ in the spirit..
And secondly, this chapter talks about diversity, right?.
Unity is not uniformity..
Unity has a huge diversity in it..
And Paul talks about it and try to unpack this truth for us..
Verses 14 and following..
Yes, the body, the human body has many different parts,.
not just one part..
If the foot says, "I'm not part of the body.
"because I'm not a hand.".
That's interesting statement..
You know, the foot is saying,.
"I'm not a hand so I'm not part of the body.".
That does not make the foot any less a part of the body..
Verse 16, "And if the ear says,.
"I'm not part of the body because I'm not an eye.".
Would that make it any less a part of the body?.
Well, this is a rhetorical question, right?.
Of course not, right?.
If the whole body were an eye, how would you hear?.
If your whole body were an ear,.
how would you smell anything?.
So Paul's using the same metaphor.
and try to unpack body has many different parts, right?.

$^{281}$You have ears, you have nose, you have mouth,.
you have hands, you have feet, right?.
You know, a couple weeks ago I hurt my wrist.
because I was trying to fix some plumbing problem at home..
And I was using the pipe wrench.
and tried to tighten the hose..
And because Nathan was not helping me,.
so I was all by myself..
I'm trying to tighten it because I hate plumbing..
I hate doing water work because no matter what I do,.
it leaks, right?.
Water leaks and it's really frustrating..
And that's why I hurt my wrist..
And because I'm getting old,.
so it takes a long time for me to fully recover..
I'm still kind of painful here..
I need to see a physical therapist pretty soon..
Yeah, right?.
(laughs).
Anyways, if one body is not working,.
it really affects my whole body, right?.
So I mean, for the last two weeks,.
it's just irritating to me..
You know, when I want to grab something,.
when I want to play guitar or something like that,.
I feel pain..
Paul's using the same language here, right?.
Let's imagine for a moment,.
the foot says, "I'm not a hand,.
"I don't belong to this body.".
The foot says, "I'm not as useful as the hand.".
Right?.
The foot is probably looking at the hand..
Hey, I'm not as functional as you are,.
'cause you use hands a lot, right?.
And the foot seems like being ignored..
Now, we may feel these kinds of statements made by Paul.
hypothetical..
They are not..
Remember, I talked to you about the letter.

$^{321}$of 1 Corinthians..
There are many different problems.
in this house churches in Corinth..
And remember, the Corinthian house churches.
are suffering from many different internal problems..
Well, it's like a body, you have different parts.
of having something wrong with it, right?.
Divisions, arguments, and people dislike each other..
And Paul is not making some hypothetical situation here..
He is actually probably tackling some real problems,.
because there are members in the church.
who really think that, you know, I'm not like him,.
you know, so I don't belong to this body, right?.
I don't have the same spiritual gifts,.
I don't have the same function, right?.
He's on the stage and I'm in the audience, right?.
So I don't belong to this body..
I don't need to come, right?.
So people did think in such terms..
Probably in today's language, probably you feel like,.
you know, I'm not a worship team member, right?.
So I'm not gifted in music..
So I don't belong to this body..
I don't belong to this worshiping community..
I'm not gifted in, you know, working on.
or operating computers back there, you know..
I don't know the sound equipment, right?.
I'm not very good at that..
I'm very bad at that..
So I don't belong to this church..
I'm not a Bible study leader, I can't read Bible study,.
you know, I'm not educated enough,.
so I don't belong here..
You know, this is a real situation Paul is trying to tackle..
But Paul is trying to help us to see what God sees,.
not how you see things..
Christians often compare ourselves to others, right?.
We naturally do that, right?.
And that's the culture of the world, right?.
That's the secular culture..

$^{361}$We have to compare and contrast with one another,.
with other people..
And we look at other people have different functions.
and we may have jealousy, we may have envy..
But look at what Paul says here, verse 18..
"But our bodies," we're turning to the same metaphor,.
"have many parts.".
God has put each part just where he wants it..
How strange a body would be if it had only one part..
Yes, there are many parts, but only one body..
The eye can never say to the hand, "I don't need you.".
The head can never say to the feet, "I don't need you.".
Right?.
Simple enough, right?.
However, if we compare and contrast our own role.
with others, there can be two results..
The first result is that the church will never grow..
Right?.
The church will never grow..
We come to the community and we compare and contrast..
We look at each other, they're different, right?.
And they're a competition..
They're comparison, right?.
You're more useful, I'm less useful..
Or I'm more useful, you're less useful, right?.
The church will never grow the way God intends it to grow..
That is, Paul is trying to say unity within diversity..
God put us different ones together..
It's like a body with many parts,.
many different functions, many different roles..
Each of them is important..
And none of us is unimportant..
Right?.
So if we come here to compare and contrast all the time,.
there might be competition, maybe envy, maybe division..
The first result is the church will never grow..
The second result is that when you compare yourself.
with other people in terms of different functions,.
someone is more useful, someone less useful,.
you do not understand the heart of God..

$^{401}$And that is the main point Paul is driving at..
Notice verse 18 here..
18, verse 18..
The bodies are many parts, but who has put them together?.
God has put them together..
Not the pastor, not the leader, not you, and not me..
And according to whose desire.
God put all these parts together?.
It's God himself, his desire..
God wanted us to be like this,.
diverse as diverse it can be..
It's the very desire of God himself to create us,.
create the church, and make us who we are,.
and put us together to become one body, one church..
Now, God does not just delight in diversity..
You know, God look at us, God look at the church..
Man, God is delightful..
It's united, he loves it..
It's diverse, he loves it..
Also, God is not just delightful.
with diversity in the church..
God has special attention to those who are weaker..
Verse 22 and following, Paul continues..
In fact, some parts of the body that seem weakest.
and least important are actually the most necessary..
And the parts we regard as less honorable.
are those we clothe with greatest care..
So we carefully protect those parts.
that should not be seen..
Verse 24, while the more honorable parts.
do not require this special care,.
so God has put the body together.
such that extra honor and care are given to those parts.
that have less dignity..
This makes for harmony among the many members.
so that all the members care for each other..
You see the point?.
There are some parts, just like our body,.
there are parts that are important but not visible, right?.
The lungs, the heart, internal organs,.

$^{441}$they are not visible,.
but they are so vital to the whole body..
Same thing for the church..
There are some members who are less honorable, so to speak,.
or less visible..
They don't have the spotlight,.
but they are as important as others..
And Paul even emphasized we should pay special attention.
to those who are so-called weaker..
Verse 26 and following..
If one body part suffers, all the parts suffer with it..
If one part is honored, all the parts are glad..
All of you together are Christ's body,.
and each of you is a part of it, right?.
Now, who are the weaker members in the church in Corinth?.
I think most likely the weaker members.
are the least gifted..
Remember we talked about worship before?.
We come to chapter 14 of 1 Corinthians.
and talk about the disorder they have.
when the Corinthian Christians gather together..
And they privilege these spectacular gifts,.
like speaking in tongue..
Those signs and wonders, right?.
Those great things, right?.
People can see it and they're so visible..
And these people who have the spectacular gifts,.
they're privileged and they have more honors,.
they have more spotlight and enjoy more attention..
And Paul is trying to subvert the order..
He says in Christ in the church,.
the seemingly weaker members.
should be regarded as equally important..
Special care needs to be given to them, right?.
The culture of the world is survival of the fittest,.
law of the jungle, right?.
If you're weaker, you stay low..
Paul argues the other way..
God sees things totally different from how we see it,.
totally different from the culture of the world..

$^{481}$The least in the eyes of God,.
the least should be given honor, dignity, and value..
I've been to church many, many years..
I remember in my home church in California,.
when I was much younger,.
we were in this small church.
and we have a member, a young adult girl.
who came to church and eventually came to Christ..
She was very, very shy, very, very timid..
And she came to Christ and for years she remains that way,.
very, very passive, very, very timid.
because of her family background..
But she continues to grow spiritually..
So the Bible and pray and worship,.
and day after day we can see she becomes more confident,.
more confident..
And along the way, we encourage her.
and pray for her to try different things.
and serve in different ministries..
And finally, I really remember she tried.
to lead Bible study..
And so we had the preparation beforehand.
and she led Bible study..
And after leading Bible study,.
she thanked all of us for the opportunity..
And we rejoice, we rejoice, we were so happy.
because she was given the attention..
She was given the nurturing..
And it's just amazing how she gradually transformed..
By the grace of God, she became more courageous..
She is a good example of weaker.
and less honorable member in this passage..
And we rejoice and increase the honor for her..
So God loves diversity..
Last point, God orchestrates everyone..
God orchestrates all..
How does he do that?.
The passage continues..
Paul says, "Here are some of the parts.
"God has appointed for the church..

$^{521}$"First, apostles, second, prophets, third, are teachers..
"And then those who do miracles,.
"then who have the gift of healing,.
"those who can help others,.
"those who have the gift of leadership,.
"and those who speak in tongues,.
"speak in unknown languages.".
Now, God is the one who puts together.
a wide variety of gifts and skills and talents together..
Now, here is a list of spiritual gifts..
I don't know which one you identify more,.
but we need to remember this is not an exhaustive list.
of spiritual gifts God has gifted to church..
There are other lists of spiritual gifts.
found in other New Testament letters..
But the whole point is Paul is trying to argue.
with a list that God is the one, not us..
God is the one who orchestrate everyone together, right?.
Some of us went to Six Facts last week, right?.
I went with them..
The whole gang we have, I'm the oldest one..
And I thought it would be the last Six Facts trip.
I ever had in my life, because I'm getting old, right?.
But after the trip, I think I could go again..
Because I can withstand the excitement, right?.
So some of us went there,.
and one of my favorite rides is Max Force..
It's called Max Force..
And it was our last roller coaster ride..
And it's very exciting, right?.
It's very exciting..
I was very happy..
I was very happy, because the Max Force,.
it goes from zero to 78 miles per hour in 1.8 seconds..
Think about it..
Within two seconds, you go from zero to 78 miles per hour..
You drive on the freeway, on a highway, right?.
Yeah, that's about, you're speeding, 78 miles per hour..
But from zero to 78, it's less than two seconds..
It was just a joy ride for me..

$^{561}$And we tried that, that's our last ride,.
because in the park, we can often hear.
the launch noise of the ride..
And the Max Force is just so loud..
When it was launched a couple years ago,.
it was too loud, and the whole Guinea area.
could hear the noise..
And the engineers had to make some modifications.
to the launch system, and adding some barriers,.
and stuff like that..
And you know the total cost for this ride is 15 million..
15 million, oh my goodness, right?.
You think about such a huge and complex machine..
There are many, many different parts..
And they all have to work together.
in order to provide the ride for the visitors, right?.
If there's one part that's missing, malfunctioning,.
the whole ride cannot operate, right?.
In fact, we were standing there for almost an hour.
for this ride, because there were maintenance..
I don't know what the engineer did,.
and what went wrong, right?.
So part of me is like complaining, right?.
And part of me is thankful, right?.
Because there are gifted engineers.
who check on the machine, right?.
Otherwise I would not be able to go home, right?.
So my only complaint is the entire joyride.
lasts less than one minute, right?.
But my point is this is a complex and huge machine,.
and every single part, nuts and bolts,.
need to work together..
The church is the same way..
The church is a beautiful, complex, complicated,.
united, and diverse organism..
Our church is complicated..
I think all churches are complicated, right?.
You deal with people, deal with sinners,.
deal with different kinds of members..
There's always complexity, right?.

$^{601}$Dealing with relational conflict, or different interests..
But it's beautiful..
That's what God decides,.
and He orchestrates everything together..
And Paul continues to end this chapter by saying,.
are we all apostles?.
Are we all prophets?.
Are we all teachers?.
Do we all have the power to do miracles,.
to pray for miracles?.
Do we all have the gift of healing?.
Those are the gifts the Corinthians yearn for..
Speaking in tongues, speaking in unknown languages..
Do we all have the ability to interpret unknown languages?.
Of course not, right?.
We're all different, and every one of us is unique..
Verse 31, so you should earnestly desire.
the most helpful gifts..
And now let me show you a way of life that is best of all..
Paul is saying not all of us are not equally gifted,.
not gifted in the same way..
But he says, however you're gifted,.
and for sure you're gifted in some way.
that you can use to build the body of Christ..
Paul is saying there's the most excellent way,.
and the most excellent way is the next chapter,.
chapter 13..
Remember chapter 13 is the most famous passage in the Bible,.
probably the most famous because in weddings,.
it's often used, you know, love is patient, love is kind,.
they love poem, 1 Corinthians 13..
It's written right after this chapter,.
and that is the most excellent way, right?.
When you come, when you come to the church,.
you need to see what God desires,.
united body and diverse body,.
and God put everything together..
Yes, you're not gifted the way I am..
I'm not gifted the way you are,.
but no matter who you are, your presence,.

$^{641}$your bodily presence, your physical presence,.
your service, and your talents,.
you're using your talents to contribute.
to the community of the church..
We are all equally important in Christ,.
and we are one body, amen?.
- Amen..
- Let's end with our verse of the week, all right?.
And remind ourselves, how can we contribute.
to this body, to the body of Christ?.
1 Corinthians 13, 13..
Now, this is after the love poem,.
love is patient, love is kind..
And verse 13, it said, "Now these three remain,.
"faith, hope, and love,.
"and the greatest of these is love.".
Now, what is love in the biblical sense?.
It's to serve other people..
Put other first, right?.
That is love..
Love is not sentiment..
Love is loyalty..
Love is commitment, and commit to serve other people.
over ourselves..
That is love..
And now, these three remain, faith we have,.
hope that we have, love that we have..
Paul says, "The greatest of these is love.".
Because without love, the church is not possible..
Without love, brothers and sisters in Christ,.
we cannot relate to one another..
The love of Christ, the love of God..
Let's all pray..
Father, increase our love..
For we confess sometimes we are loveless..
With our lips, we sing praises..
With our lips, we pray and say that you are love..
And we thank you for your love,.
for your redeeming love,.
saving us from sin and from death..

$^{681}$But we confess so often we just don't have love in our heart.
because we don't understand your heart..
And your heart is nothing but perfect love..
You want to see, you yearn to see the church united,.
yet using our diverse gifts to serve one another,.
to relate to one another,.
to have fellowship with one another,.
to spend time with one another..
Father, we come and confess our sins and our shortcoming..
And we pray that you will cleanse us.
and to cleanse our community,.
help our community, our spiritual family,.
our church to grow in love..
We bind us to serve one another just as you serve us..
Renew our loyalty, renew our faithfulness.
so that we, as we serve one another,.
as we build this body together,.
we see your presence and we taste your love.
in Christ and pray..
Well, God bless you all..
And next week, we're gonna come together as one body.
to serve and to worship God together..
Please continue to pray for our city this week..
It's gonna be in chaos and yeah, stay away from downtown..
If you live in downtown,.
I know you have some traffic problem maybe,.
but pray for our city, pray for our neighbors..
All right, see you next week..
\newpage



\section{}
\label{sec:nNDuojAF0eE}
\textbf{開往水深之處: 信仰歷程中不同階段的禱告生命 - 余淇煥傳道}
\newline
\newline
連結: \href{https://youtube.com/watch?v=nNDuojAF0eE}{\texttt{https://youtube.com/watch?v=nNDuojAF0eE}} ~~~~ 語音日期: 2024-08-23
\newline
\newline
\hyperref[sec:l_mjl7q5DSI]{\small{< < < PREV SERMON < < <}}
~
\hyperref[sec:index_chronic]{\small{[返順時目]}}
~
\hyperref[sec:index_scriptual]{\small{[返順卷目]}}
~
\hyperref[sec:jtr0kNAa9Hw]{\small{> > > NEXT SERMON > > >}}
\newline
\newline
$^{1}$大型姐妹早安.
在這個夏季有幾位實習生.
還有在課堂也有幾位實習生.
剛才Carmen也說了.
神樂院也開課了.
所以我們也希望能夠有.
大型姐妹在我們中間.
他們在接受神樂訓練.
也給他們機會來實習的一個.
一個方位.
有Bella.
這位姐妹已經大過敬拜.
大家應該不陌生.
她現在在Loyola讀.
Master in Spiritual Direction.
即是靈性的指導.
是一門今時今日很需要的學問.
我們把時間交給她.
歡迎Bella.
大型姐妹早安.
可能大家也習慣我參與敬拜.
很少見我講道.
多謝主讓我有機會在UV講道.
多謝Master介紹我.
我也先介紹一下自己.
我叫Bella,一個L.
我的中文名叫余祺煥.
我在香港在中神MDiv畢業.
之後在一間教會做了三年的青少年傳道.
大概五年前跟阿邦來了.
因為他進修.
所以來了芝加哥.
我現在在一間天主教大學.
Loyola讀靈修指導.
其實我小學也在天主教小學讀書.
經常會走進聖堂對著耶穌像祈禱.
當時有個志願.
大家猜猜是甚麼.
就是做修女.
但之後知道做修女不可以結婚.

$^{41}$就打消了這個念頭.
人的意志很薄弱.
今天我想透過三段經文.
跟大家分享反思禱告的歷程.
第一個故事記載在馬海福音9章14至24節.
這個故事我告訴大家.
大概是說有一群人在爭論.
耶穌問他們在爭論甚麼.
原來有個爸爸帶了自己的兒子來.
這個兒子是被啞巴鬼附身的.
這個污靈會怎樣呢.
就是不斷傷害這個小朋友.
一時就把他扔進火裡.
一時就讓他跳進水坑裡想淹死他.
總之就是想他死.
這個爸爸也請過耶穌的門徒趕走他.
但能趕走嗎.
趕不走.
耶穌就說他們.
這個不信的世代.
我還要和你們一起多久.
我要忍受你們多久.
把這個孩子帶過來.
有一段很有趣的對話.
就是當這個父親詳細講解給兒子聽.
這個被鬼附身的兒子的惡劣情況之後.
就說耶穌如果你能夠做些甚麼.
就請你憐憫我們幫我們.
然後耶穌怎樣回應他.
在和學本翻譯得很淡然.
大家都很熟悉的一個金句.
你若能信,再信的人凡事都能.
但在新漢語的譯法是一個反問句.
耶穌問如果你能.
是這樣的語氣.
如果你能.
然後他就說對於信的人一切都能.
耶穌正在挑戰.
為甚麼你會覺得只有我才有這個信心.
凡信的人一切都能夠.

$^{81}$不是只有我才有這份信心能趕鬼.
而是你和我都可以有這份信心.
話說我是在一間較少提及.
靈性能力的教會成長.
所以當我在大學的時候.
接觸到一些回靈恩教會的大學同學時.
他們都會帶我去一些醫治特會.
那些醫治特會有時候是醫治身體.
有時候是醫治心靈.
我也很熱衷去.
因為我很好奇聖靈的能力彰顯會怎樣.
我也很好奇見到神蹟奇事.
那裡的牧師都真的很厲害.
他們會按著神給他們的感動.
然後就說出一些回中的病徵.
然後那些有那種病的人.
就會來到台前給牧師祈禱.
很多人都會湧出去.
之後他們也會說見證.
那些見證就是腎石被祈祟了.
長短腳又被祈求回到正確位置.
我也很怕有痔瘡的出來.
我和同學都不知道出去好不出去好.
我想說牧師能不能醫治是一回事.
但很多時候我們會覺得.
這些厲害的東西只有某些人才能做到.
為人祈禱.
Mark Sir 或麗倩那些應該比較靈.
讓祂去祈禱吧.
但耶穌說是不信的世代.
我看這段經文的時候也很驚訝.
我也把我被鬼附的兒子帶來耶穌那裡.
為什麼耶穌還會問.
如果你能,為什麼還要被人串呢.
耶穌原來今天說.
祂想門徒和尋找祂的人知道.
每一個人都可以像耶穌那麼有信心.
不單止耶穌的信心.
能夠彰顯祂神的能力.
而是每一個人都能夠有耶穌的信心.

$^{121}$能夠彰顯到神的能力.
可能你在職場,在教會.
你總會覺得有些人特別有恩賜.
特別有能力.
你很自然會把自己給下去.
可能神一直都感動你.
為一些人做一些事.
但你就會覺得我哪裡可以.
我有很多限制的.
讓祂來吧.
祂那麼厲害.
但今天我們會不會.
都好像被耶穌說的那些人.
我們都不能相信我們祈禱.
能夠彰顯到神的能力呢.
會不會我們經常都仰望一些.
屬靈的偉人.
覺得他們.
只有他們才有那份信心呢.
如果今天在你們所身處的場景.
有些人,有些事.
需要你憑信心去禱告.
以至到神可以藉著你.
去祝福那些人,那些事情.
你又願不願意呢.
這個星期初.
有一個UV的姐妹.
我們在言談之間.
她跟我分享.
神感動她為我祈禱.
其實我沒有特別.
講代禱事項給她聽.
她突然間就這樣分享.
神感動我為你祈禱.
就分享了一會兒.
然後我覺得.
我好被祝福.
因為有人關心,有人在乎.
神今天提醒我們祈禱.
不用等那些屬靈偉人.

$^{161}$而是耶穌給了我們.
有跟祂同等的信心去祈禱.
去彰顯神的能力.
對於信的人一切都能.
然後我想跟大家看第二段經文.
在《馬禾福音》十四章.
三十四至三十六節.
是《喀西瑪利園》的禱告.
這段經文.
就在耶穌臨上十字架之前.
跟門徒在喀西瑪利園的一段禱告.
是很罕有地記載.
耶穌自己禱告的內容.
這跟剛才說的祈禱很不同.
剛才在九章.
耶穌教導的禱告.
就是只要你有信心去祈求.
你就會成就.
你想得醫治.
你祈禱就得醫治.
但這個禱告有什麼不同呢.
這裡不只是單純有信心.
耶穌說.
阿爸父啊.
對你來說一切都可能.
求你拿走這個杯.
一切都可能.
很同情上一段經文說的一切都能.
就是說耶穌其實是信父神可以.
不是因為他沒有信心.
但他的祈禱沒有停在這裡.
而是在這裡馬可要重新再定義.
一個更闊的信心.
什麼是更深的禱告.
他說但是不要依我的意思.
要依你的意旨.
原來禱告是一種跟自己交戰的張力.
禱告是一個選擇.
去放下自己慾望的一個過程.
是一個選擇跟神說.

$^{201}$我願意跟你的一個旅程.
不過你不要以為我所說的慾望.
一定是一些很負面.
很自我的一些事.
有時候慾望.
有時候期盼.
可以是你期待一些.
你覺得合神心意的好的事.
話說我是在一個情緒很動盪的家庭長大.
除了大家多少都聽過一些情緒病.
例如抑鬱症,鬱躁症,焦慮症.
其實心理學裡還有一種性格障礙.
還有性格障礙的人.
他也可以是一個高智商的人.
他也可以很有魅力.
他也可以是一個很有能力的人.
性格障礙例如有自戀型的人格.
反社會人格.
每一種性格障礙都有不同的特徵.
有不同的病徵.
我在大學時期.
有一晚因為我媽媽喝醉酒.
醉酒街頭就被送進了醫院.
我第一次聽到醫生跟我說.
媽媽很可能患有邊緣人格障礙.
當時我開始找一些相關的資料.
就發覺九成都很吻合.
邊緣人格障礙其中一個很大的特徵.
就是他很難控制自己的裂路.
我說裂路是真的裂路.
不是憤怒 是裂路.
所以家暴不時都會在我家發生.
家裡有很多暴力.
甚至見血.
當我媽媽一控制不到自己的時候.
她就可以很躁狂.
拿著刀對著我爸爸.
或者在廚房煮完滾水.
一推就推出客廳.
我從小到大.

$^{241}$很多時候半邊身都會乳乳黑黑上學.
這些事並不是個別的事件.
是日常的.
但很奇怪.
我媽媽並不是什麼道德敗壞的人.
也不是黑心的.
她是一個會照顧流浪貓的人.
她是一個去老人院探望老婆的時候.
會照顧隔離床老人家的人.
所以我作為她的女兒.
一直都不明白為什麼她回到家裡.
是另外一個人.
而事實上邊緣人格障礙.
是真的有這種特徵.
是她不能控制自己的極端性格.
極端的情緒.
自我信主開始.
我都很熱愛祈禱.
我會花很多時間祈禱.
並且我會很認真地為家裡禁食.
不是晚上隨便戒掉杯麵.
不是戒宵夜.
而是就算我上班.
三天我都很堅持只喝流質飲品.
一有時間我就去祈禱.
後來我帶了我外婆.
我姨媽和我媽媽信主.
在同一天洗禮.
大家馬上就會想.
可以了,快樂的結局,完成了.
但我媽媽沒有因為我帶了她信主.
或者她洗了禮,而她得醫治.
她依然是一時抑鬱到爆笑.
很歇斯底里.
一時不停地找交捱.
我父母關係也沒有因為這樣而變得好.
有時三更半夜都會吵醒我們.
不停地吵架.
我記得當時我自己讀輔導.
也在讀輔導.

$^{281}$我如何形容家裡呢?.
我家裡好像一個地獄.
同時也像一個北極.
很冷.
也有時很激烈,很憤怒.
我很難感受到愛.
因為我父母根本上沒有空間.
沒有心靈空間去關心或問候我們.
相反,我要背負著這個家.
甚至我父母的情緒壓力.
從小到大,我見證著這個家.
不單是我們受折磨.
其實我媽媽的病也很折磨自己.
從小到大,我也見證著她.
不少次服毒,服藥,企圖跳樓.
情況一點也沒有改變.
神蹟奇事也沒有發生在我家.
每次當我聽到一些基督徒.
牽著父母上台分享.
如何得醫治,家庭復和的見證時.
我也很想有一天.
可以和父母上台講見證.
我每天都很努力為他們祈禱.
原來這個苦杯.
我一直也很想神去拿走.
因為實在太辛苦.
但當時,我不停為他們祈禱.
在耶穌的祈禱下半部分.
祂說不要照我的意思.
雖然我的祈禱中,我祈不出這一句.
但我覺得我已經在走那個方向.
我覺得我帶著很多包袱也好.
雖然我帶著很多重擔也好.
但其實我也一直向神感動我的.
全職事奉獻身的方向走.
我也願意在那時候.
雖然我講不到一些神蹟奇事的見證.
但我也會背著結他去中學或教會.
講一些見證是關於在苦難中.
如何堅持我信仰的見證.

$^{321}$我發覺原來祈禱是很不容易的.
原來祈禱的生命不單只有上半部分.
阿爸父對你來說一切都可能.
求你拿走這個杯.
原來祈禱有下半部分.
有「但是」.
「但是」不要依我的意思.
要依你的旨意.
有時有些事你很想這樣發生.
有些事情,有些人.
你很想他變成這樣.
但事情沒有因為你的祈禱.
而發生你想像的那樣.
那你還願意祈禱說.
「主啊,我是想這樣」.
「但是不要跟我,跟你吧」.
原來我發現禱告是一個充滿張力的過程.
禱告是一個充滿張力和掙扎的旅程.
最後我想帶大家看馬可福音15:33-34.
這裡也是一個祈禱.
就是耶穌在臨死前的最後一個祈禱.
這個祈禱只有一句.
「我的神,為什麼離棄我?」.
剛才上一段經文是說.
耶穌跟父神說.
「不要依我的意思,要依你的意思」.
原來當耶穌說不依自己的意思.
依神的意思走的時候會怎樣?.
他會體會到什麼?.
是不是心裡火熱?.
覺得神跟他很近.
充滿著神很實在的愛.
或者有聲音在天上肯定他.
「這是我的愛子,你做得對」.
不是,相反.
耶穌深深感受到父神離棄了他.
在這個禱告裡.
耶穌選擇走神的意旨的時候.
他為什麼會這麼深感受到被神離棄?.
這個跟我們剛才九章看的禱告層次很不同.

$^{361}$我們以為我們憑信去禱告.
然後神的國度就能以很大的權能.
用一個很神職的方式去展現.
那些鬼只要你憑信就可以趕走.
但是馬可福音告訴我們.
神的國度是怎樣展現的.
是透過耶穌的死.
是要透過耶穌的死.
神的國度就會展現.
但原來當神的國度真的展現的時候.
連耶穌自己也想像不到.
我們會問.
為什麼耶穌走在神的旨意裡.
都竟然覺得神遺棄了他,離棄了他?.
我記得很久之前.
有一個牧師分享他自己的生命歷程.
當時我開始認識一個名字叫做淑玲黑夜.
這位牧師花了很多時間祈禱.
他也是一個很熱忱,很熱情的報道家.
他也進了很多中學去講很多場報道會.
帶很多中學生去信耶穌.
不單是信耶穌.
而是那些中學生願意獻身給主.
是會為主而活的.
我是其中一個被他影響的人.
但原來連他自己一個熱愛禱告,親近神的人.
原來他自己也走過一段淑玲黑夜.
他形容在那段日子裡.
他完全感受不到神.
他一向都是一個對神的聲音很敏銳的人.
但原來一個多淑玲的人也好.
也可能會經歷這種覺得跟神很遠.
甚至神好像不在場的淑玲黑夜.
不知道大家有沒有試過這種跟神很遠的感覺呢?.
我剛才分享.
就算我生命有很多艱難的狀況.
還是面對很多破碎.
其實我也很想堅持我的信仰.
我一直都硬著頭皮.
就算流著眼淚.

$^{401}$我也很想為主向前走.
但到了我讀中神的最後一年.
接近畢業那一年.
也是我準備結婚的那一年.
我面對著信仰上一個從未有過的危機.
就是當我這麼多年.
走神的路.
為著家庭.
我覺得很合神心意.
因為很想媽媽經歷神.
很想她得醫治.
很想家庭復活.
這些事情去祈禱的時候.
在2016年有一天.
我走在街上.
有一晚我走在街上.
我突然收到一個電話.
是我的好朋友打給我的.
很奇怪.
他問我.
他打給我.
他說「阿邦在你旁邊嗎?」.
「你讓他聽吧」.
那一刻.
其實我聽不到任何東西.
但第一時間我就有一個預感.
就是「是不是我媽媽出事了?」.
原來當日在黃昏時段.
我媽媽決定用自己的方式.
去了結自己的生命.
雖然我媽媽的自殘.
對於我來說一點也不陌生.
剛才也說過.
她不時也會服藥.
突然又會爬出窗戶.
所以我心底裡.
經常都有一種恐懼.
不知道哪一天她會出事.
所以我經常跟神說.
「神啊,你保護她吧」.

$^{441}$可能在我心底裡.
也很害怕有這一天.
但沒想過當這一天.
真的發生的時候.
我會是怎樣.
我面對著神.
我心裡有很多憤怒,失望,質疑.
我的神,為甚麼要理會我呢?.
神啊,你去了哪裡?.
在我媽媽做這麼痛心的決定時.
你去了哪裡?.
我祈禱求你去醫治她.
你也沒有去醫治.
又或者我希望家庭復和.
但也沒有復和.
但我也繼續讀神學.
我也繼續做傳道人.
我也夠聽話.
為甚麼結局會是這樣呢?.
對我來說,是一個死亡的結局.
我相信了耶穌這麼多年.
我到了那個關口.
我覺得不用祈禱了.
這麼努力做甚麼?.
人都沒有,復甚麼和?.
媽媽這樣的離開.
彷彿將我的心一大部分抽走了.
從那天開始.
我有一段很長的時間是靜默的.
我無法祈禱.
因為我很怕.
我一祈禱.
在安靜之中要重新觸碰那些痛楚.
在我默會三年.
我暫且將這個部分.
放在一個很深的角落裡.
當我來到芝加哥.
我失去了教會傳道人的身分的時候.
我發覺我不想再回教會.
我發覺就算我回教會.

$^{481}$每當唱讚美詩的時候.
我讚美不了.
我只會流眼淚.
甚至有時候我封鎖了自己.
我很麻木.
這樣的狀態維持了好幾年.
我人生未試過跟神的關係這麼大的隔閡.
直至我遇上了一個靈性指導.
我只見過他兩次.
他在德國.
我在芝加哥.
隔著螢幕.
事隔這麼多年.
竟然有一個人敢問我.
當日的整個過程.
竟然有人敢去觸碰和聆聽我.
不敢再回去回顧的地方.
竟然有一個人陪伴我.
再踏進一個我自己一個人不敢踏進的歷程.
我彷彿覺得在我面前的.
不只是一個人.
而是神的同在.
我可以形容那幾年的時間.
我是在一個全黑的黑洞裡.
我完全全黑.
我完全什麼都看不到.
但是由那一刻開始.
我發覺我跟神的關係有一線的光.
後來我就是因為這樣.
我決定讀靈性指導.
我希望可以這樣陪伴一些有需要的人.
去走一段路程.
耶穌的呼喊.
我覺得是一個很深度的禱告.
耶穌的呼喊.
我們有時候會覺得很輕易地說.
祈禱啊 有什麼開心不開心都跟神說.
但是來到這個時候.
我發覺原來也不是那麼容易.
你能夠呼喊也不是那麼容易.

$^{521}$因為耶穌的呼喊.
我的神 你為什麼離棄我的.
這種呼喊.
是要展現你自己的脆弱.
在一個你很信靠的人面前.
你展現你的脆弱.
所以這個禱告.
這種禱告是一個很信靠的禱告.
耶穌的呼喊.
耶穌不是永遠停留在覺得神離棄的那一刻.
或者這個過程可能是很漫長.
而耶穌的死.
也不是停留在一個永遠的死亡.
神彰顯的方式.
可能有時候很神色歧視.
是有的 有這些時刻.
但是有時候神彰顯的方式.
可能連我們也不明白.
但是在祂無聲當中.
祂跟人同在.
神在人最黑暗的角落裡.
祂可能默默無聲.
但是其實祂從來都沒有缺席.
神會用不同的方式跟人同在.
即使你覺得神離棄了你的時候.
原來神的兒子耶穌.
他也會這樣經歷.
覺得被神離棄的感覺.
他也會這樣禱告.
他也會這樣呼喊.
今天不知道在我和你的生命歷程裡.
又在經歷什麼.
你這一刻在屬靈裡.
你覺得自己在一個怎樣的狀況.
你又會怎樣去形容.
你這一刻和神的關係.
很想這一刻有一點時間.
我們一起安靜.
我們一起祈禱.
剛才分享的訊息裡.

$^{561}$或者不同的生命故事裡.
有什麼是你覺得最深刻.
最觸動你.
或者是神在跟你說.
這段時間我邀請你.
親自和神有些對話.
親愛的主人.
你認識我們每一個.
你知道我們所經歷的每一天.
你知道我們的狀態.
你知道我們生命的狀況.
你知道我們正在經歷什麼.
你知道我們說不出口的祈禱.
我們每一個都在經歷生命裡不同的歷程.
主今天你跟我們說.
無論有些弟兄姊妹.
你提醒我們要憑信心去禱告.
做你要我們做的事.
抑或是你提醒我們.
原來禱告是一種張力.
是一個不斷放下自己的過程.
抑或你明白我們在困難當中.
有時都會聽不到你的聲音.
有時都會覺得神你是否離棄了我們.
無論在什麼階段.
神你都讓我們看到你的故事.
你的呼喊是一個很信靠.
很信靠的信心.
你的故事也不是停留在死亡.
你的復活被人看到.
在人間的困難當中.
在苦難當中.
你就是那一線的光.
你是願意踏進人間.
與人同在的神.
所以主今天願你親自去觸摸我們每一個生命.
願你親自成為我們的光.
你也教我們走向一個更深的禱告生命.
一個信服你的生命.
我們的祈禱是奉耶穌基督得勝的名字祈求.

$^{601}$阿們.
今天我們還有一點時間.
希望在最後大概十分鐘.
如果需要禱告的話.
需要有人陪伴你.
去到有些祈禱.
可能你心裡有些重擔等等.
我們希望你能夠來到台前.
我們教培同工和Bella都在.
我們都會為你祈禱.
其他大英姐妹.
祝福大家平安.
下星期我們一起去敬拜我們的主.
好嗎?.
這裡散會.
\newpage



\section{15:16}
\label{sec:jtr0kNAa9Hw}
\textbf{Fruit that Lasts (John 15\_16) - Sam Johnson}
\newline
\newline
連結: \href{https://youtube.com/watch?v=jtr0kNAa9Hw}{\texttt{https://youtube.com/watch?v=jtr0kNAa9Hw}} ~~~~ 語音日期: 2024-08-27
\newline
\newline
\hyperref[sec:nNDuojAF0eE]{\small{< < < PREV SERMON < < <}}
~
\hyperref[sec:index_chronic]{\small{[返順時目]}}
~
\hyperref[sec:index_scriptual]{\small{[返順卷目]}}
~
\hyperref[sec:7j_P4EWgCCI]{\small{> > > NEXT SERMON > > >}}
\newline
\newline
$^{1}$worship team for leading us. Really enjoyed worshipping with you all. This.
morning we're gonna be looking in the book of John chapter 15, a pretty.
familiar passage, and talking about fruit in our lives, spiritual fruit and fruit.
that lasts. This thing that Jesus introduces to us in this passage. This is.
something that all of us want and desire in our lives. We want to see spiritual.
fruit. We want to see fruit coming out of our lives and when I when I use that.
term fruit I mean the the evidences in our own selves that we're growing in.
Christ, we're becoming more like Jesus, that God is changing us and transforming.
us, but also the evidences in the lives of others when people that we are in.
contact with, people that we have relationship with, we want to see them.
coming to know Christ more or growing in Christ. So all of us desire fruit and we.
want fruit that lasts. There are a lot of examples of that and what I'm gonna.
start with this morning just to give us a neat example of fruit that lasts is in.
East Africa. The countries in East Africa there's about 267 million people in.
those countries and of those 267 million over a hundred and eighty eight million.
are professing Christians. That's over 70\% of the population. In countries like.
Zambia, Uganda, and Kenya over 80\% of the population identifies as Christian. In.
Rwanda it's over 90\% of the population. By 2050 the ten countries with the.
largest Christian populations will include several African states including.
Ethiopia, Nigeria, Congo, and Uganda. So there's been this amazing growth of the.
church in that part of the world and global Christianity has begun to and.
will continue to be shaped by renewal and leadership from from that part of.
the world. But where does that come from? The roots of this renewal go back to a.
revival that took place in the 1930s and this was a powerful movement of God that.
began with a small group of believers who simply had a desire to walk closer.
with God and obey Jesus in every area of life. One of these believers was a man.
named Simeon Sibambi in Uganda. Simeon made a vow to the Lord, "I've committed.
myself to God the Father as from today I desire to be genuinely holy and never.
intentionally do anything unguided by Jesus." His ardent devotion to Jesus and.
his zeal for making disciples began to inspire others and his life was a main.
catalyst of the East African revival. Although he probably never foresaw the.
way that God would use his obedience or the impact his life would have on his.
nation and eventually the globe, God powerfully used him to bring about fruit.
and fruit that remains to this day, fruit that has lasted, fruit that will echo.
into the future. And that's a desire that all of us have, not necessarily that we.
will impact the world in the way that this man's life did, but all of us want.
to see fruit in our own lives and fruit that will last in our life and beyond.
our life into the lives of others and beyond the time that we're on earth.
hopefully. Hopefully things that God does in our lives will be things that will.
echo into someone else's life long after we are no longer here in the lives of.

$^{41}$children or in the that may come someday or in the lives of friends who impact.
someone else and they impact someone else. We want to see that. That's what.
Jesus is talking to us about in John chapter 15. So I'm gonna read this.
passage and then we're gonna look at some of the things that Jesus has to say.
to us. This is John 15. I'm gonna read the first 17 verses. "I am the true vine and.
my father is the vinedresser. Every branch in me that does not bear fruit he.
takes away and every branch that does bear fruit he prunes that it may bear.
more fruit. Already you are clean because of the word that I have spoken to you..
Abide in me and I in you as the branch cannot bear fruit by itself unless it.
abides in the vine. Neither can you unless you abide in me. I am the vine..
You are the branches. Whoever abides in me and I in him, he it is that bears much.
fruit for apart from me you can do nothing. If anyone does not abide in me.
he is thrown away like a branch and withers and the branches are gathered.
thrown into the fire and burned. If you abide in me and my words abide in you,.
ask whatever you wish and it will be done for you. By this my father is.
glorified that you bear much fruit and so prove to be my disciples. As the.
father has loved me so have I loved you. Abide in my love. If you keep my.
commandments you will abide in my love just as I have kept my father's.
commandments and abide in his love. These things I have spoken to you that my joy.
may be in you and that your joy may be full. This is my commandment that you.
love one another as I have loved you. Greater love has no one than this that.
someone lay down his life for his friends. You are my friends if you do.
what I command you. No longer do I call you servants for the servant does not.
know what his master is doing but I have called you friends for all that I have.
heard from my father I have made known to you. You did not choose me but I.
chose you and appointed you that you should go and bear fruit and that your.
fruit should abide so that whatever you ask the father in my name he may give it.
to you. These things I command you so that you will love one another. As we.
think about this passage I just want to remind us of what's going on here at the.
context of this teaching that Jesus gives. This is after the upper room, after.
the the last Passover that Jesus shares with the disciples. So this is right in.
the last few days of Jesus's life, in the last hours of Jesus's life really. And so.
he is teaching some of the most significant things that he's going to.
teach the disciples. And he begins to describe what life is like to follow him.
in these chapters in the book of John. Now surely the disciples heads are.
already kind of spinning. They've seen so much that was hard for them to take in,.
so much that they're still trying to understand. In John 13 Jesus washes their.
feet. They have the Last Supper together. He warns them about the betrayal. Then.
Judas leaves and Jesus gives them the new commandment of love. In John 14 is.
where he talks about how he is the way, the truth, and the life. That there's no.

$^{81}$way to God, to heaven, to enlightenment, to relationship with God except through.
Christ. And then in John 14 as well he gives the promise of the Holy Spirit. And.
now he gives them this really beautiful word picture of life in Christ and what.
it looks like. It's profoundly significant and full of truth for the.
disciples and for us. And so Jesus recalls, this is actually an Old.
Testament metaphor of the vine. The Old Testament talks in multiple places about.
the people of God, Israel, being this fruitful vine. It's in Psalm 80 verse 8.
in Isaiah chapter 5 and Jeremiah 2. It talks about the people of God being this.
vine that bear fruit to glorify God. And Jesus is now saying to the disciples, "Hey.
you know this story. You're familiar with this metaphor. I'm actually the vine. I'm.
the one who is going to produce fruit and I'm gonna do it through you, through.
your lives. As you abide in me, I want to produce salvation, blessing, joy, peace. I.
want this fruit to come out of your lives." And that's what Jesus is speaking.
to them and speaking to us. The first thing I want us to see here is that if.
we want to experience fruitfulness and fruit that lasts as followers of Jesus,.
the only way to see that is to remain in Jesus. This word for abide is in this.
passage. Jesus says it over and over again. It means to stay in a given place.
or in a certain relationship. It means to continue. It means to dwell. It means to.
endure. It means to be present. It means to remain or to stand in a place. And so.
when Jesus is saying, "Abide in me," he's saying, "Remain connected to me. Remain in.
relationship with me. Continue in relationship. Stand in this relationship..
Be in Christ." And he's telling us to remain in his words, to follow his.
commandments. He says in verse 10, "If you keep my commandments, you'll abide in my.
love, just as I have kept my Father's commandments and abide in his love.".
Saying this is a part of the abiding is to keep my commandments. What does that.
mean? Well, "keep my commandments" is a church phrase that really just means to.
do what Jesus said, to live out the things that he talked about and that he.
instructed the disciples in and that he instructs us in to do it, to live it out..
That's what it means to keep his commandments. And so that means that we're.
remaining in his word, remaining in his truth. Again, what does that mean?.
There's nothing unusual about it. It just means that we continue to put in our mind.
and in our hearts the words of Christ. As we read scripture, as we look at what God.
says, as we do this every week and we come and gather together as his people.
and we talk about his word, that's a part of remaining in his word. It's a part of.
abiding in his truth. That's a part of abiding in Christ. It's continuing to.
stay connected to what he said and what he taught us and how he has told us to.
live and how we obey that. It means that we remain vigilant against sin. The Bible.
has all kinds of things to say, especially later in the New Testament.
and all the epistles, about behaviors that are wrong in God's sight and they're.
also harmful and self-destructive to us. And it warns us about those things. It.

$^{121}$says don't do these things. Follow after God. Stay free from sin. Stay.
vigilant against these behaviors and these desires in your own heart,.
whether it's pride or fear or anger or lust or desire for power or desire for.
control or all of these different things. We have to remain vigilant against those.
things. That's a part of remaining in Christ, a part of keeping his.
commandments. A part of abiding in Christ and remaining in Christ means to remain.
in his body. And we know the New Testament talks about that a lot more..
When it talks about this another word picture of the body of Christ, that it's.
like a human body and that if we're going to be a healthy part of it, we have.
to be in it and be connected. If we get disconnected, then we start to experience.
unhealth and disease and even death. And Jesus talks about that in verse 6. If he.
says, "If anyone does not abide in me, he's thrown away like a branch and.
withers and the branches are gathered, thrown into the fire and burned." If a.
branch gets disconnected from the trunk that is giving it resources to live, then.
it's gonna wither. It's gonna die. We all know that. If you've ever cut a branch.
off of a tree or seen one that gets blown off in a storm, the leaves don't.
stay green very long. They start to wither quickly. And so part of remaining.
in Jesus means we remain connected to his body. We continue to do this. We.
continue to gather together and worship and look at his words together and share.
the truth with one another. And that's a part of what keeps spiritual life going.
in us and keeps spiritual fruit coming out of us. And so these things are all.
kind of the the basics of our faith to remain in connection with Jesus. We need.
to remain in his word. We need to remain connected to him in prayer, communicating.
our heart to him. We need to remain connected in worship and in the life of.
his people and community of the church. Those are the basic things, but if we're.
not doing them, there's not going to be fruit. There's not going to be fruit that.
lasts. And sometimes we look at our lives and we say, "Well, I don't see the fruit I.
want to see. I've been asking God for some stuff and I want to see more and.
I'm just not seeing it." And that's a great place to start. It may be basic, but.
it's a great place to start is to ask ourselves, "Okay, well am I connected to.
Jesus right now? Maybe the reason I'm not seeing fruit is I'm just not really.
connected to God. I've just let my relationship with him grow cold and I'm.
not spending time in his word. I'm not spending time in worship or in prayer..
I'm not spending time with his people." And as basic as that is, a lot of times.
that's right where the issue is. We're just disconnected. The next thing we see.
here is that abiding, which is the only way to fruit that lasts, is a.
relationship of trust. In verses 4 and 5, it says, "Abide in me and I in you as the.
branch cannot bear fruit by itself unless it abides in the vine. Neither can.
you unless you abide in me. I am the vine, you are the branches. Whoever abides in.
me and I in him, he it is that bears much fruit. For apart from me you can do.

$^{161}$nothing." Those verses make it clear that the abiding thing is not us trying.
harder and working more for God because we can't do anything anyway. Jesus makes.
it really clear. The abiding is us trusting him to do the work that only he.
can do. In Isaiah chapter 5, the first seven verses God is describing his.
relationship to Israel and he's using the metaphor of Israel as the fruitful.
vine and it talks about the things that God does in order for his people to be.
fruitful. It talks about that God prepares the land, God removes the stones,.
God builds a hedge around the vineyard to keep out wild animals who would.
destroy the fruit, God builds a watchtower with the vineyard to keep the.
enemy away, and God prunes and God does the harvest. If you haven't ever read.
that passage, it's a great passage just to look at in Isaiah 5. And it reminds us.
that all the work is work that God does. We don't prepare the land, we don't.
remove the stones or build the hedge, we don't keep the enemy away, we don't keep.
wild animals away, and yet in our life a lot of times we feel like we need to do.
it. The fruit is gonna come because we need to work harder, we need to be.
smarter, figure more things out, and do more. And that's not what Jesus is.
calling us to here. He's saying apart from me you can do nothing. You have to trust.
and let God do the work that only God can do. Both to make us more like Christ.
but also to allow us to be fruitful in someone else's life. Only God can do the.
work of drawing people and changing hearts. Only God can bring healing,.
whether that's physical healing or what we've probably experienced, all of us.
have experienced even more emotional healing, things that have harmed us and.
things that are really keeping us from going where God wants us to go, and.
things that are causing self-destructive behavior in our life. Only God can bring.
that healing. And so abiding means to trust in God to do the work of the.
vinedresser. It means to trust in Jesus as the source of every need in our life..
It means to trust in Jesus as the source of fruitfulness. Instead of trying to.
find other things to get connected to, that's such a huge temptation for us in.
life. We say I want to see more happening in my heart or my life or in someone.
else's life, so let me find some other things to connect to that I think will.
fill me up and will bring some fruit out of my life and will really, you know, get.
me going. I'm gonna connect to something else. And the sad thing is as we do that.
in life, we find that every other thing that we connect to ends up being.
unsatisfying and unfulfilling. It might give us a short-term rush and we think.
this is it, this is the new thing and I'm so happy and I feel great, but over time.
we find that we're unsatisfied and unfulfilled and all those other things..
Christ is the only connection that's really gonna give us life and it's the.
only thing that's gonna give us fruitfulness. And the other beautiful.
thing about abiding is that in that trust we experience rest, right? Verses.
4 and 5, Jesus is real clear that he does the work. We, apart from him, we do nothing..

$^{201}$Isn't that a peaceful thought to have? It's not peaceful for our flesh because.
like the proud part of us, the selfish part of us wants to be like, no, I should.
do something because I'm really talented, Jesus. Like you need me on.
your team, look at how good I am at all this stuff. But when we're in that mode.
and we're trying to work it out ourselves, we're gonna pretty quickly.
reach a point of complete discouragement, burnout, exhaustion because we can't do.
it ourselves. No matter how talented and smart and gifted you are, you're not.
gonna be able to live out the identity of a disciple of Jesus without Jesus.
living through you. You just can't. And so if you're trying to do that, you're gonna.
you're gonna be exhausted. And abiding gives us rest. It allows us to say, I.
can't, but he can and I'm trusting in him. I'm resting in his work in my life and.
that's a beautiful place to be. Another thing we see here, a difficult part of.
this passage but profound, is that pruning is essential for fruit that.
lasts. This word in the original language in verse 2 means to cleanse of impurity.
or filth or dirt and it means to prune trees and vines from useless shoots..
Anything that's gonna be fruitful, but especially a grapevine, requires pruning.
to get to maximum fruitfulness. Grapevines have root suckers that grow.
up off the roots and there are these branches that grow out of the roots.
that aren't gonna be fruitful at all, but they pull a lot of energy and resources.
out of the main trunk of the plant. So those have to be pruned off for the.
plant to be fruitful. Also, most grapes produce fruit on shoots growing.
off of one-year-old canes or branches. If you have too many old canes because you.
haven't pruned, you'll get very few grapes because they're all too old. But.
if you prune everything back completely every year, again, you'll get very few.
grapes. You'll get a lot of new growth, but very few grapes. So the pruning has.
to be consistent and strategic to get fruitfulness, to maximize fruitfulness. We.
can't do that ourselves. Only God can do that. And God looks at our life and he.
says, "Hey, I want to use you and I want to do something with you, but there's some.
stuff that we got to cut away. There's some root suckers that, they have some.
big green leaves, but they're just pulling resources out of your heart and.
they're not bringing any fruit. And you're growing in the wrong direction..
You're getting twisted. Vines also, vine dressers also use trellises and.
wires to grow the vine so that it'll be fruitful and also so that it will be,.
they'll be able to harvest it. And God has to do that in our lives. He has to say,.
"Okay, there's some growth over here that it's really headed in the wrong direction.
and I'm gonna have to prune that back." That's not painful. I mean, that is.
painful. Sorry, that's not fun. It's painful, but it's so essential for.
fruitfulness, especially fruit that lasts. And so we have to learn as followers of.
Jesus to embrace that and to let God prune away some of that unhealthy growth..
And when we're in those seasons of pruning and we feel that, we feel like,.

$^{241}$"God, I feel like you're pruning me and you're stripping away some stuff that.
it's hard for me to let go of and I don't see fruit. It doesn't feel good.
right now. What do we do in that season?" We do what Jesus said. We remain. We.
continue to dwell with Him and we continue to abide because that pruning.
is essential if we're going to become fruitful. If God's going to get us to the.
place where He can do the things that He wants to do in our life, we have to.
continue to remain in Him, remain in His Word, remain in obedience, remain in His.
body, remain connected. Otherwise, the pruning begins. We become.
diseased and pretty soon we're disconnected and then we're being.
gathered for the fire. We're not useful to God. Another thing about.
fruit and grapevines I think is really interesting is that maturity is.
essential for fruit that lasts. When you think about any natural plant, what brings.
maturity? Well, time first and foremost, right? It takes some time but it also.
takes going through seasons of life. Through some winter seasons, through some.
flood seasons, through some times of drought and feeling very dry, through.
seasons of fire where it seems like everything around you is burning up and.
and you're wondering what's happening, through seasons of having attacks, right?.
The grapevines is one of the reasons the vine dresser would build the hedge.
because the grapevines have wild animals that come in and try to claw at the.
fruit and try to destroy the vine. There's thieves that come in and want to.
to steal the fruit before the vine dresser can get to it. Those are all a.
part of the seasons of life and they're a part of maturity. And if we're gonna.
get to fruitfulness and fruit that lasts, we're gonna have to go through some.
different seasons. We're gonna have to go through some of those seasons and maybe.
we're experiencing that right now. Maybe you're in a place of experiencing drought.
where you just feel spiritually dry and you feel like there's not much water in.
your life. Or maybe you're in a season of a fire where you feel like, man, things.
are burning up, things are falling apart. Like, I don't know what's.
happening in my life right now. God, where are you? Maybe you're experiencing a.
season of the wild animals where it feels like someone else is clawing at.
you trying to destroy the fruit or the good things that God has done in your.
heart. Those seasons are challenging, they're difficult, but they're a part of.
the process of maturity and they're the only way to grow in Christ and to.
experience fruitfulness. And so again, what do we do during those seasons? If.
you are in a season of drought or flood where you feel overwhelmed and it's just.
like the water's coming up over your head or fire, what do you do? You remain..
You abide. You stay connected to Christ. You say, you know what, I don't see the.
fruit right now. I'm not seeing what I want in my life. I'm not seeing it in the.
lives of people around me. You know, God's put me in some place where I.
work or where I go to school or in my family and I know I'm supposed to.

$^{281}$minister here and I'm trying to serve others. I'm trying to love people and I.
don't see any fruit and it just seems like the world's going crazy and.
everything's on fire. What do you do? You remain. You abide. You continue to remain.
in Christ, remain true to his word and be obedient to him, remain connected to his.
body and you continue to live out the truth that he's told you to live out. You.
continue to be obedient to him. That's what you do in those seasons. You.
continue to hold on to him, continue to trust that he can bring the fruit and he.
will bring the fruit and continue to rest and experience his peace. You're not.
responsible for the fruit. He is. And the other part of maturity is that there's.
time that has to take place before the fruit that last comes and this is.
important for us to think about. It's obvious but it's still important for us.
because we live in a modern industrial world, right? Jesus uses an illustration.
that involves organic time and that's intentional, okay? Most of us don't live.
on organic time. We live on machine time. We're used to machine processes that.
happen very quickly. We're used to things that can be produced very quickly. We're.
used to microwaves and drive-thrus and car assembly lines that produce a.
thousand cars in a month, right? We're used to machine processes. Jesus,.
apparently discipleship and spiritual fruit is not a machine process and.
that's the whole reason Jesus uses this very organic illustration to say, "Hey,.
growth and fruit, it's gonna take some time." It doesn't matter what you're.
planting. If you're planting apples or pears or tomatoes or grapes or corn,.
there is no fruit immediately after planting. If you've ever gardened and you.
plant some tomatoes in your backyard, we do a lot of gardening in our.
neighborhood and at our house. We love it. My wife and I do a lot of that with the.
kids. We had our, right now our tomatoes are doing fantastic, our peppers and we.
have a lot of squash, but we planted all that stuff, most of it way before we.
could even take it outside. So back in February, March, because in Chicago if you.
want to have fruit before, you know, the first freeze in September, October, the.
plants need to be mature when you put them in the ground the first week of.
June. So they got planted a long time ago. You put seeds in the ground, there's no.
fruit. There's no fruit for days. There's no fruit for weeks. And some trees, like.
fruit trees, there's no fruit for years. Sometimes you have to wait five years.
with certain fruit trees before you get the first fruit. This is an organic.
process. Here's something a little bit scary for us to think about. Fake fruit.
can be produced immediately. If you want to make apples and tomatoes out of.
plastic, you can make them right away, right? But what's the thing about fake.
fruit? There's no nourishment. It doesn't have what real fruit has. It might look.
good, especially if it's sitting on the other side of the room. You might think,.
"Wow, that apple looks really good." If you ever tried to eat a wax apple or a.
plastic grape, you know what I'm talking about. My mom had plastic grapes.

$^{321}$somewhere in our house when I was little and I did put them in my mouth a few.
times because they looked really good. But when you chew on them, it's just.
plastic, right? It looks good, but it doesn't have the taste and the texture..
It doesn't have the vitamins and minerals. It doesn't have the nourishment..
It's fake. And we're tempted to do that in our walk with God so much. We're.
tempted to say, "Man, real fruit takes a long time in my life. It takes a long time.
in discipling someone else to see real fruit." So wouldn't it be better just to.
do fake fruit? Like we could just speed the process up and really kind of juice.
this up and make it happen real quick. But the problem is fake fruit doesn't.
have what you're really looking for. It doesn't have that texture and taste. It.
doesn't have that, the nourishment. It's huge for me and my wife. As parents, we.
talk about this all the time. Something that we deal with in parenting all the.
time. We're tempted to just try to modify our kids behavior and make them do what.
we think is right. That's fake fruit. There's nothing wrong with working on.
their behavior, but we want their hearts to be changed, not just their behavior,.
right? Because if it's just their behavior, that may look okay for a while,.
but if their heart hasn't been changed, eventually they'll say, "I'm not.
interested in this behavior anymore," right? So we want their heart to be.
changed, and that's how God is working in our lives. He's not just interested in.
something superficial that looks good that happens overnight where he says,.
"Okay, you know, you used to dress this way. Now dress that way. All right, now you're.
done." That's not what God's doing. God is doing something transformative in our.
hearts. Is it gonna change our behavior? Yes. Should our behavior look different.
than the people around us who don't know Christ? Yes. But here's the key..
God's work transforms us from the inside out. Our behavior changed because our.
heart is different, and if we're just modifying our behavior, that's from the.
outside in, we can try to do better all we want, but if our heart's not different,.
we're just gonna keep going back to the same stuff. And the process of spiritual.
fruit is a transformative inside-out process. God says, "Is your behavior gonna.
be different? Yes, but it's gonna be different because your heart's different,.
your desires are different. You are a new person, and that takes time." And so if you.
are in a season of waiting right now where you say, "I don't see the.
fruit in my own life. I'm trying to obey, but I still have these struggles. I still.
have, you know, temptations that I'm like, 'Why is this still here, God?'" Or maybe.
in someone else's life, someone that you love and you're serving, you're trying to.
help them grow in Christ, and you say, "Man, I've been walking with this person for a.
couple of years now, or even longer, and I'm just not seeing the fruit." Understand.
that God is at work in an organic process. It's gonna take some time. And.
one of the things that we can't bypass, even in the world that we live in with.
all of our technology, we can't produce our food on machine time. We just can't..

$^{361}$As much as we juice the process of crops and fruit and whatever else we grow, we.
can only do it so much before it becomes diseased and unhealthy, and then we can't.
eat it. And that's what spiritual fruit is. You can't--there are some ways.
that you can connect with God, yes, but you can't really shortcut that process..
There's no real, like, spiritual steroids that just make you bulk.
up overnight. That's not a thing. You've got to spend the time with God, and God's.
got to work in your life. And that's the same is true when you're walking with.
other people. You've got to give God time to work in their life. So I want to.
encourage us with that. I think as we follow Jesus, we're always tempted to.
compare ourselves to others. We're always tempted to feel like we just need to try.
harder. We're always tempted to say, "Well, this person seems to be, you know, going.
well in their life, or they're growing, and like, why can't I experience that?" Or.
we're always tempted to ask ourselves, "Oh, you know, I should--this should be.
faster. I should be experiencing more of God's work in my life. This should.
happen a lot quicker. I'm trying to love this person over here, and nothing's.
going on. Why not? What's the problem?" We want to live on machine time. We want to.
produce false fruit. We want to say, "God, I'm tired of waiting for you to work, so.
maybe if I just do some stuff, it'll speed up, and we'll all feel better about.
it." I do that all the time, all these things. That's why I'm bringing them up,.
because they're very familiar to me. And yet, Jesus is calling me to remain, to.
abide, saying, "Don't try to make shortcuts. Don't try to juice the process. Just.
remain. Just abide. If you're in a season of drought, season of fire, season of.
flood, keep waiting. If you're in a season of just not much going on, continue to.
abide. If you're even--if you're in a season of suffering or hurt, or where.
you're confused in some way, continue to abide. Remain connected to Jesus. If.
you're in a season of pruning, where God's taking some stuff out of your life,.
continue to abide. Be connected to Him. Be connected to His people. And as we do.
that, God is at work, and He's bringing spiritual life into us, and it is going.
to express itself in fruit. It will. But it takes time, and it takes remaining in.
Christ. It takes abiding. It takes standing in Him and continuing to be.
connected to Him. Let's pray. God, I just thank you for the beauty of your words,.
and of this word picture for us that helps us see that the spiritual fruit in.
our lives is not necessarily going to come in the way that we want it to. And.
God, help us to trust you. Help us to abide. Help us to remain, to continue to.
be connected to you, and obedient, and live out your commands, and stay.
vigilant against sin, to continue to follow you, even when it's hard, even when.
we're experiencing suffering, or pruning, or drought, or fire, or flood. Help us to.
wait in the times when we're just waiting for maturity, or for growth, or.
for something to be strengthened. God, help us to remain in you, and as we do.
that, help us to find rest. Help us to really trust you to do the work that.

$^{401}$only you can do, to take the pressure off of ourselves, and to just really rest in.
you, and in your work in our lives. And God, I pray in my life, and in the lives of.
each one of us in this room, that you would get glory, and you would get fruit..
That fruit would come out of our lives that glorifies you, and that impacts the.
world around us, our family, our friends, our community, our neighborhoods, for the.
sake of your name, and for your glory, and for the love that you have for those.
people who surround us, God. I pray that that fruit would be abundant in our lives..
In Jesus' name, amen..
Okay, thank you Sam for teaching us and encouraging us from the Word of God.
today. And yeah, as we have learned from John 15, my prayer for us as a church, and.
for all of us, is that we can, yeah, learn to abide deeply in Christ, especially at.
times when there are so many voices, ideologies, grabbing, fighting for our.
attention. So, thank you for the reminder again, Sam. And if you're a regular at UV,.
you know the drill after the service, but if you are new here, first, welcome to UV,.
and second, usually after the service, we'll have some refreshment at the cafe.
area, which is just to my right. So, we would like to encourage you to stay back,.
enjoy some food, and also give us an opportunity to get to know you as well..
So, go in peace, UV, your love..
\newpage



\section{哥林多後書 8:1-15}
\label{sec:7j_P4EWgCCI}
\textbf{恩典的慷慨 (哥林多後書 8\_1-15, 9\_6-15) - 譚文倩莊傳道}
\newline
\newline
連結: \href{https://youtube.com/watch?v=7j_P4EWgCCI}{\texttt{https://youtube.com/watch?v=7j\_P4EWgCCI}} ~~~~ 語音日期: 2024-08-27
\newline
\newline
\hyperref[sec:jtr0kNAa9Hw]{\small{< < < PREV SERMON < < <}}
~
\hyperref[sec:index_chronic]{\small{[返順時目]}}
~
\hyperref[sec:index_scriptual]{\small{[返順卷目]}}
~
\hyperref[sec:pxG7cdEbvuY]{\small{> > > NEXT SERMON > > >}}
\newline
\newline
哥林多後書 8:1-15
\newline
\begin{longtable}{cl}
\hline
\hline
章節 & 經文 (和合本修訂版)\\
\hline
8:1 & \begin{tabularx}{0.7\textwidth}{X} 弟兄們，我們要把神賜給馬其頓眾教會的恩惠告訴你們： \end{tabularx} \\ \\ \relax
8:2 & \begin{tabularx}{0.7\textwidth}{X} 他們在患難中受大考驗的時候，仍然滿有喜樂，在極度貧窮中還格外顯出他們樂捐的慷慨。 \end{tabularx} \\ \\ \relax
8:3 & \begin{tabularx}{0.7\textwidth}{X} 我可以證明，他們是按著能力，而且超過了能力來捐助，主動 \end{tabularx} \\ \\ \relax
8:4 & \begin{tabularx}{0.7\textwidth}{X} 再三懇求我們，准他們在這供給聖徒的善事上有份； \end{tabularx} \\ \\ \relax
8:5 & \begin{tabularx}{0.7\textwidth}{X} 並且他們所做的，不但照我們所期望的，更照神的旨意先把自己獻給主，又給了我們。 \end{tabularx} \\ \\ \relax
8:6 & \begin{tabularx}{0.7\textwidth}{X} 因此，我們勸提多，既然在你們中間開始這慈善的事，就當把它辦成。 \end{tabularx} \\ \\ \relax
8:7 & \begin{tabularx}{0.7\textwidth}{X} 既然你們在信心、口才、知識、萬分的熱忱，以及我們對你們的愛心上，都勝人一等，那麼，當在這慈善的事上也要勝人一等。 \end{tabularx} \\ \\ \relax
8:8 & \begin{tabularx}{0.7\textwidth}{X} 我說這話，並不是命令你們，而是藉著別人的熱忱來考驗你們愛心的真誠。 \end{tabularx} \\ \\ \relax
8:9 & \begin{tabularx}{0.7\textwidth}{X} 你們知道我們主耶穌基督的恩典：他本是富足，卻為你們成了貧窮，好使你們因他的貧窮而成為富足。 \end{tabularx} \\ \\ \relax
8:10 & \begin{tabularx}{0.7\textwidth}{X} 我在這事上把我的意見告訴你們，是對你們有益，因為你們開始辦這事，而且起此心意已經有一年了。 \end{tabularx} \\ \\ \relax
8:11 & \begin{tabularx}{0.7\textwidth}{X} 如今就當辦成這事，既然有願做的心，也當照你們所有的去辦成。 \end{tabularx} \\ \\ \relax
8:12 & \begin{tabularx}{0.7\textwidth}{X} 因為人只要有願做的心，必照他所有的蒙悅納，並不是照他所沒有的。 \end{tabularx} \\ \\ \relax
8:13 & \begin{tabularx}{0.7\textwidth}{X} 我不是要別人輕鬆，你們受累，而是要均勻： \end{tabularx} \\ \\ \relax
8:14 & \begin{tabularx}{0.7\textwidth}{X} 就是要你們現在的富餘補他們的不足，使他們的富餘將來也可以補你們的不足，這就均勻了。 \end{tabularx} \\ \\ \relax
8:15 & \begin{tabularx}{0.7\textwidth}{X} 如經上所記：多收的沒有餘，少收的也沒有缺。 \end{tabularx} \\ \\
[1ex]
\hline
\hline
\end{longtable}
$^{1}$弟兄姊妹平安.
在大家當中和大家靜悟.
上個月我去了加拿大.
由我家開車上去大概八個小時.
差不多在芝加哥開車上去多倫多.
Mark Sir是否說過.
芝加哥有什麼地方好去.
除了去多倫多.
你開長途車開到差不多要做什麼.
就是找廁所.
我們找廁所.
找廁所的時候.
發覺我們快要走近四角.
即是四個州.
亞利桑那州.
加拿大州.
新墨西哥州.
Utah的交界處.
厲害嗎.
現在我經常做這些運動.
因為年紀大.
所以我做伸展運動.
我們進入四角.
轉進去之後.
才知道原來每位要八元.
然後我和我老公說.
我們屙了人生中最貴的一獨住住.
這麼快暑假又過了.
深呼吸.
回一回氣.
我們又回到現實.
準備上班.
預備新的學年.
是否不成唏噓.
放完假.
收拾好心情了嗎.
我也不例外.
像剛才所分享.
下星期我會坐飛機去伊拉克.
做醫療短期宣教.

$^{41}$主要服侍的對象.
就是被IS種族滅絕的Yashidi難民.
2014年8月3日.
ISIS半夜到Yashidi.
聚居的地方Sinjar進行大屠殺.
目的是要種族滅絕.
和強迫Yashidi的人改信伊斯蘭教.
今年是大屠殺的十週年紀念.
根據聯合國的數字.
因為這次屠殺.
已變成難民的人有五十多萬人.
今時今日仍然有二千多個婦女和小朋友還未找到.
從2008年開始.
我差不多每年都會去短宣.
但久不久有人問我.
為何要去短宣.
為甚麼呢.
今天我跟大家講的題目叫做.
因電的慷慨.
經文來自哥林多後書八章一至十五節.
和九章六至十五節.
哥林多後書是寺徒保羅寫給哥林多教會的書信.
哥林多是一個古希臘城市.
當時是一個很繁榮的商業中心.
亦是一個很多道德敗壞.
很混亂的地方.
保羅在書信裡提到他曾經去過多倫多.
去建立了這個教會.
但後來他亦收到風和.
聽到這個教會裡面有很多問題.
在哥林多後書一至七章裡.
保羅強調信徒在面對困難的時候.
信靠上帝的重要性.
當中亦提及自己在亞細亞所受的苦難.
上帝在當中給他的安慰和鼓勵.
保羅很注重彼此相愛.
寬恕和和好.
慷慨施予和忠心服侍.
當中對基督徒的生活教會紀律作出很多教導.
來到第八和第九章.

$^{81}$突然保羅就轉一轉話題.
他就勸勉哥林多教會.
要慷慨捐助給有困難的耶路撒冷貧窮的基督徒.
解釋基督徒要慷慨的原因.
今天我嘗試從三點在保羅的勸勉裡.
和大家一起思考慷慨捐獻的功課.
我們先低頭祈禱.
昭去主我們的神.
藉著你的話語.
去引導我們怎樣去過我們的生活.
過一個真是討你喜悅的生活.
讓我們的生命都是榮耀.
還稱讚你.
好像一個禮物那樣呈現到你面前.
蒙你月亮.
今日祈禱奉耶穌得正名求.
Amen.
第一點.
慷慨的教會.
哥林多後書八章一至五節.
弟兄們我要把神賜給馬其頓.
重教會的恩惠告訴你們.
他們在患難中受大考驗的時候.
仍然沒有喜樂.
在極度貧困中.
還格外顯出他們樂捐的慷慨.
我可以證明.
他是按著能力.
而這超過了能力來捐助.
主動再三懇求我們.
准他們在這供給聖徒的善事上有份.
並且他們所做的.
不但照我們所期望的.
更照神的旨意.
先把自己獻給主.
又給了我們.
第八章一開始.
保羅稱讚馬其頓教會.
在極度貧困的情況下.
他們仍然很慷慨.

$^{121}$在這裡提到三點.
第三節說.
按著能力.
而且超過了能力來捐助.
第三節的尾部分.
去到第四節說.
主動再三懇求我們.
准他們在這供給聖徒的善事上有份.
第五節.
並且他們所做的.
不但照我們所期望的.
更照神的旨意.
先把自己獻給主.
又給了我們.
馬其頓教會自己也不是很足夠.
但他們仍然對耶路撒冷的信徒.
教會是很慷慨.
甚至是.
他們懇求有份參與.
他們的動作.
是超越了保羅對他們的期望.
他們愛主.
更加體貼主的心.
慷慨捐獻給有需要的人.
《使徒行傳》十一章二十七至三十節.
當那些日子.
有幾位先知從耶路撒冷下到安提訥.
來中有一位名叫瓦加布站起來.
藉著聖靈指示抱天下將大饑荒.
這是在赫魯第十年間果然實現了.
於是門徒決定照各人的力量捐錢.
順去供給住在猶太的弟兄.
他們就這樣做了.
托巴拿巴和訴羅的手痛到中掌腦那裡.
其實保羅在寫《哥林多前書》的時候.
即是主後的五十五年左右.
他已經幫助耶路撒冷.
在加拿大的教會裡無捐.
而且當時哥林多教會聽到這個消息.
也表示願意參與.

$^{161}$臨前十六章一至四節.
關於為聖徒捐款的事.
我從前怎樣吩咐加拿大的宗教會.
你們也該怎樣做.
每逢七日的第一日.
每人要照自己的收入抽出藥桿.
保羅起來免得我來的時候現抽.
等到我來了你們寫信舉崇誰.
就差遣他們把你們的款項送到耶路撒冷去.
如果我有該去你們可以和我同去.
到了寫《哥林多後書》的時候.
即是大概一年之後.
保羅知道馬其頓教會做了這樣的事.
但是哥林多教會似乎未完成他們之前承諾會做的.
到了臨後八章六節.
因此我們勸提多.
既然在你們中間開始這些慈善的事.
就當把它辦成.
既然你們在信心口才知識萬分的熱忱.
以及我們對你們的愛心上都勝人一等.
那麼當在這前事上也要勝人一等.
我說這話並不是命令你們.
而是藉著別人的熱忱來考驗你們愛心的真誠.
保羅在這裡就對他們發出挑戰.
哥林多教會這些經常誇讚自己比人多懂得去講方言.
多有知識也多有恩賜.
但是你們可不可以同樣在慷慨捐獻的事情上也比人更勁呢?.
保羅提醒哥林多教會奉獻.
不單只是因為他們曾經承諾過.
更加其實他是希望他們在慷慨奉獻的時候.
反映他們認識那個神是一個怎樣的神.
感恩的心就藉著祝福別人.
能夠被彰顯出來.
成為祝福人的動力.
講錢是不是傷感情呢?.
如果你有看過哥林多前後書的時候.
你也知道他們之前是挑戰保羅的教導.
其實他們似乎不是很喜歡保羅.
但是這件事沒有阻礙保羅不出聲.
或者就著就著他仍然堅持去教導真理.

$^{201}$因為八章九節.
他本是富足,卻為你們成了貧窮.
很少你們因他的貧窮而成為富足.
在這裡保羅指出慷慨捐獻的原因.
一個終極的榜樣.
耶穌基督.
很多年前我在一個聚會裡認識了一個弟兄.
這位弟兄在1980年代的時候.
他在加州擁有其中一個規模最大的木材生意.
當時正值加州樓市很蓬勃.
他的生意好得不得了.
他說每天進來公司的錢多到數都數不完.
他說買大屋,買地,買車.
所有他想買的都買盡了.
吃盡真羞白味.
但是他說他這樣花錢.
錢還是不斷在這裡多.
是越來越少的.
他說Joyce.
你知不知道當人有錢到一個這樣的地步.
還有什麼好去追求.
你知不知道我還想追求什麼.
他說就是追求名.
他說下一步就是追求名.
怎樣追求這個名呢.
他說他要捐一大筆錢去一個大學.
讓他們將其中一個大學的建築用他的名字去命名.
當他這樣說的時候.
我腦海裡面就出現了一些東西.
不知道你會不會在腦海裡面出現了一些圖畫.
我就想起在香港坐巴士進紅隧的時候.
經過理工大學.
你就會見到一大堆建築物排在那裡.
有李兆基樓.
周奕卿樓.
陳麗玲樓.
XX樓.
香港大學也有日呼書院.
中文大學也有日呼書院.
原來這些建築物的名字是這樣來的.

$^{241}$不說太遠了.
只說我們芝加哥.
我們教會之前去MSI.
這個Museum of Science and Industry.
這個科學和工業博物館.
都是在我們剛剛教會去之前.
收到一筆一千二百五十萬的捐款.
所以現在嚴格來說不是叫做MSI.
是叫做KCG MSI.
因為捐錢的人叫做Kenneth C. Griffin.
我為何說這些呢.
我想說的是聖經裡面.
耶穌基督慷慨和我們地球人一般的慷慨.
是有一個很大的分別的.
就是那個動機和那個心態.
聖經裡面說的那個慷慨.
是出於愛.
講真誠.
是無私的.
也是持久的.
是不求回報的慷慨.
目的是要讓對方去蒙恩.
甚至無得著救恩.
而我們地球人一般的慷慨.
有時是出於這個社會的壓力.
為了建立自己的形象.
或者是有些benefit自己的動機.
當中不單止可能是有條件.
更加是希望得到回報.
在第一世紀有錢和有權勢的人.
都是靠剝削貧窮的人.
去建立自己的財富和勢力.
這些地主,皇帝,政府官員等等.
久不久他們就會向一些貧窮的人去施恩.
以至貧窮的人對他們很感恩.
很順從他們.
表面來看是很有愛心的.
但實際上他們可能是為了建立自己的形象.
或者只是為了保障他們自己.
能夠繼續平平安安.

$^{281}$建立他們的財富.
鞏固他們的勢力和地位.
但保羅在這裡說的是一個完全相反的目的.
就是主耶穌基督願意變為貧窮.
將一切放上.
因為祂是希望人得著富有恩典的禮物.
保羅在這裡是在說聖經的哲學.
一個和我們地球人完全不同的價值觀.
不同的動機.
這就是聖經裡面.
植根在愛和慷慨.
真正施予的根基.
臨後八章十二至十五節.
因為人只要有願造的心.
必照他所有的夢月立.
並不是照他所沒有的.
我不是要別人輕鬆你們受累.
而是要均勻.
就是要你們現在的富餘.
補他們的不足.
使他們的富餘.
將來也可以補助你們的不足.
這就均勻了.
如經上所記.
多收的沒有餘.
少收的也沒有缺.
多收的沒有餘.
少的也沒有缺.
是出自埃及的嗎哪的故事.
十六章十八節.
用餓魔耳量一量.
多收的沒有餘.
少收的也沒有缺.
各人都按著自己的食糧收集.
保羅所說的均勻.
是神供給以色列人所需要的.
沒有剩也沒有欠缺.
有多的要分給有需要的人.
這就是均勻的意思.
宗教會也應該朝著這個方向前行.

$^{321}$馬其頓教會正正是這個榜樣的表表者.
第二點.
慷慨的上主.
在之前保羅提到耶路撒冷的教會.
經歷了很多迫害.
饑荒和物質上的欠缺.
保羅懇求各個教會一起幫忙籌集資金.
而哥林多教會也響應.
來到第九章.
他呼籲哥林多教會.
不要忘記你們之前答應了的事.
特別是保羅之前也告訴過別人.
這個哥林多教會是多麼願意參與.
他們是多麼慷慨.
但在過程中保羅也是很小心翼翼地處理這件事.
因為他不想誤會.
表面鼓勵人奉獻給有需要的人.
但其實自己是忠報私囊.
代代平安.
過程令主的名受辱.
臨後九章六至十一節.
還有一點.
小眾的小收.
多眾的多收.
各人要隨心所願.
不要為難.
不要勉強.
因為神愛樂捐的人.
神能將各樣的恩惠多多加給你們.
使你們凡事常常充足.
能多做各樣善事.
如經上所記.
他施捨.
周濟貧窮.
他的義含存到永遠.
那似種子給撒種的.
次糧食給人吃的.
必多多加給你們種地的種子.
又增添你們仁義的果子.
你們必凡事富足.

$^{361}$能多多施捨.
使人借著我們而生.
保羅用撒種和收割的比喻.
解釋慷慨的原則.
小眾的小收.
多眾的多收.
上帝喜悅甘心樂意慷慨的人.
並且答應供應人一切所需要.
甚至增加他們的資源.
以致身邊的人有需要的時候.
他們可以隨時幫助.
成為別人的祝福.
在過程中更加引發人.
剛才我說過.
上個月我和我老公開車去了加羅拉多.
行程的其中一站叫做Black Canyon.
懸崖峭壁加上澎湃的河流.
你想想我站在海平線上八千呎.
下面河流的澎湃聲音.
在山頂上很響亮.
那種美麗和壯觀.
真的非不物能夠形容.
另外我也參觀了.
大約八百年前印第安人在懸崖上建居的居所.
這些懸崖的居所.
主要是採用自然洞穴和石頭所形成.
通常是位於高處.
提供了防禦性和隱蔽性.
更加反映當時居民對天氣和環境的了解.
以及他們很巧妙地運用建築的技術.
如何在惡劣的環境中耕種草食物和草水.
印第安人主要的糧食是粟米.
種植了幾千年.
被視為一種很神聖的植物.
除了是他們的糧食之外.
粟米對他們的身份,精神實踐和哲學信仰.
原來是有很重大的意義的.
因為粟米是為他們生命和生活的象徵.
被視為是全能,全知,造物主給他們的禮物.
因為他們沒有文字.

$^{401}$所以音樂對他們來說很重要.
也是他們一個很重要的傳統.
因為他們要靠音樂和唱歌.
將一些重要的訊息和傳統傳給下一代.
將它編成歌曲.
經常在唱,教導下一代.
甚至很有趣的是我們今時今日用GPS.
原來他們要去一個地方.
那個方向也會有歌唱.
以前的人是靠著這些歌曲一路唱.
就知道我究竟向東還是向南.
怎樣去到目的地.
我一直聽導遊在說的時候.
我真的覺得神真的很奇妙.
你創造天地萬物.
那個偉大.
你想想這個和空氣,泥土,水,太陽等東西.
有不同的植物出來供人觀賞.
也有不同的蔬菜,水果.
給我們有生命.
再說音樂,我自己不懂音樂.
但我又對音樂多了一些了解.
更加欣賞.
很感恩神把音樂給了我們.
能不能想像我們的生命沒有音樂.
想想去旅行親身大自然的時候.
你會只在那裡享受.
還是你會忍不住.
特別是我們基督徒.
忍不住讚美神.
神啊,你真是很厲害.
你是一個多麼厲害的神.
你的創造是多麼偉大.
臨後九章八至十節.
神能將各樣的恩惠多多加給你們.
他施捨那次種子.
次糧食給人又增添.
神除了供應我們生活所需.
大自然的壯麗.
好像我剛才所說.

$^{441}$會不會令你忍不住要讚美神.
我們的造物主是多麼奇妙.
多麼慷慨.
去跟我們分享祂的創造.
我靈哥唱.
讚美救主我神.
你真偉大.
何等偉大.
第三點.
慷慨的祝福.
九章九至.
祂施捨周濟貧窮.
他的兒.
行傳道永遠.
這段經文是引用了詩篇一百一十二篇.
說到敬畏上帝的人.
慷慨使喻使他人受益.
上帝的公義動作.
公義的動作.
傳道永遠.
其實詩篇一百一十二篇.
和之前那一篇一百一十二篇.
一百一十一篇是互相緊扣.
詩篇一百一十一篇.
說到上帝是就屬於他子民的行為上.
去彰顯他的公義.
來到一百一十二篇.
敬畏上帝的人.
慷慨使喻被有需要的人.
這個慷慨使喻的行為.
就是行義.
因為詩篇一百一十一篇裡面的上帝.
就是行為公義的上帝.
這個觀念的重要性.
是說上帝祝福人.
人就要同樣祝福別人.
幫助有需要的人.
《哥林多後書》九章裡的慷慨.
就是建基在這個基礎上.
慷慨使喻最終的動機.

$^{481}$是源於耶穌基督的榜樣.
一九四五年第二次世界大戰結束後不久.
德國被劃分成四個佔領區.
英國控制西北部.
法國控制西南部.
美國控制南部.
蘇聯控制東部.
去到一九四八年的時候.
因為東西方之間的意識形態差異.
蘇聯就封鎖了柏林.
封鎖了柏林.
切斷了西北林和西德的陸路和河流運輸.
西方國家當時就唯有聯盟起來.
用空運的方法去援助住在柏林裡面.
物資短缺,飽受飢餓,寒冷的人民.
這個人道主義的行動裡面.
當中有一個軍人叫Gail Harvison.
她不忍心小朋友受苦.
所以她就自己用了一條手巾.
去做了一個像降落傘的東西.
紮住了一些糖果.
就在飛機上綁住一些糖果.
由飛機那裡扔下去拋給小朋友.
她這個動作在當時的打仗極之慘.
極之處於一個很慘,很低沉,很可憐的境況裡面.
將整個城市的氣氛扭轉了.
當人們再看到飛機在上風飛過的時候.
不再想起打仗.
不再想起要逃跑.
反而是飛機一來的時候.
大家就追著飛機這樣去.
因為飛機的來到帶給他們希望.
帶給他們快樂.
這個愛心的行動.
也鼓勵了很多已經厭倦打仗.
甚至質疑我們究竟要援助.
Berlin裡面的人究竟要多久.
我們的國家還要給多少錢出來.
去幫一些我們三五四七的人.
這些人因為Gail的行動.

$^{521}$結果願意參與這個捐獻.
參與無論是奉獻很多糖果或者金錢.
來給這個軍人.
甚至是更加多的軍人參與.
去投放這些糖果直升機到Berlin這個地方.
這件事也提升了美國當時的聲譽.
人們幫Gail起了一個花名.
叫做Candy Balmer 糖果轟炸機.
回答之前的問題.
為什麼要去捐獻呢.
我想起我自己很喜歡去捐獻.
因為我可以親身去到當地.
接觸當地的人.
去觀察,聆聽.
親身了解他們的需要.
也挑戰我自己在不同的環境,文化,言語當中.
去跟別人建立友誼.
分享見證和傳福音.
你知道最棒的是什麼嗎.
最棒的是去到那裡.
能夠親身看到神在那些人裡面的工作.
甚至看到一些人在一個很難的環境下.
我都不覺得說完福音他會相信的情況下.
竟然相信耶穌.
以致他們仍然在那個困難,困境裡.
他們竟然可以帶著盼望和笑容繼續生活.
從此對生命有盼望.
甚至成為當地基督,全洋基督的使者.
去年去伊拉克短孫的時候.
因為我不吃羊肉.
很多人都聽過.
所以沒有一口好吃.
加上水土不服,餓了八天.
而我和我的天媚除了負責照顧病人和做手術之外.
整個團隊每天這三個星期所需要的東西.
都由我們兩個負責.
結果我們兩個就做到趴在那裡.
有晚我和我的天媚回到房間後.
兩個人是累得呆了.
我們真的坐在床上沒有話說.

$^{561}$兩個人呆了.
我就提議,不如我們大家一起祈禱.
真的不太能忍受.
誰知我們兩個人閉上眼.
一聲就哭了出來.
兩個人是爆哭.
全程我都聽不到他們祈禱了甚麼.
他也不知道我祈禱了甚麼.
祈禱完之後,我們兩個都很痛心.
我們說下次不要再回來.
對了,為甚麼今年又去呢?.
去年我和我的天媚覺得不行.
我們要想想要建立一些新的系統.
去幫助整個程序做得更好.
這是第一個原因.
第二個原因是.
如果說辛苦.
我們絕對不讓醫生,護士.
他們也很辛苦.
他們當中很多人.
不止一個,不是說一年去一次.
甚至一年去四次.
和當地的難民做手術.
做跟進.
如果沒有人去那裡.
幫他們處理一些其他事情.
其實他們很難專心去做手術.
第三.
很想念當地的難民.
有時候我想起Yashidi的時候.
覺得他們很傻.
因為他們其實很忠心.
一個有問題的宗教.
去信奉一個傷害他們的宗教.
不單止因為這個宗教.
他們被人種族滅絕.
更因為為了他們遵從這個宗教的原因.
要保持他們的純正血統.
他們的下一代大部分都有傷殘.
有很多健康的問題.

$^{601}$醫療團隊去到他們當中.
不單止是跟他們做手術.
其中一樣很想做到的事.
去斷宣最大的原因.
就是要將真理帶到他們裡面.
讓他們認識真神.
扭轉他們的生命.
一個軍人愛心的行動.
扭轉世界對美國人的形象.
帶給人希望.
蒙恩的基督徒.
又可不可以藉著人道的救援.
扭轉人對基督教的印象.
對耶穌基督的印象.
藉著耶穌基督的福音.
帶給人希望.
臨後九章十二至十五節.
因為辦這功給的事.
不但補足了成圖的缺乏.
他們從這功給的事上得了憑據.
知道你們宣認基督.
順服祂的福音.
慷慨捐助給他們和眾人.
把榮耀歸給神.
他們也因神極大的恩賜.
顯在你們身上.
而切切想念你們.
為你們祈禱.
因他有說不盡的恩賜.
我們的團隊裡面.
有很多人是自費去參加斷宣.
當中我認識一個團友.
在參加決定之後失業.
在他還沒有失業之前.
他就允許了一些主動奉獻給他的人.
但想不到他那邊推了人.
說不需要奉獻.
就失去了那份工作.
他跟我說.
其實他很想找那些人跟他說.

$^{641}$我現在有需要.
但他又不好意思開口.
在我自己的籌款裡面.
竟然不多不少.
剛剛有足夠覆蓋他那份.
神為他預備.
給他憑據.
神要他.
包括我.
包括我去看到.
並且經歷到.
我們所信的神是一個.
多麼的信靠.
是一個功級的神.
神如何藉著我們.
在斷宣祝福我們.
不是說我去祝福人.
神是藉著我去參與.
我還沒有去.
神已經在祝福我們.
彰顯他們的大能.
叫人不可以不將榮譽歸神.
Hannah也說Amen.
叫我們不去數算神給我們恩典.
這就是恩典的慷慨.
神祝福我們.
我們祝福其他人.
所以你們慷慨的奉獻.
不是因為你們喜不喜歡我.
說真的.
所以我要支持Joyce.
也不是你寫張名片.
或者你買一支XO像這麼簡單.
其實你這個公義的行為.
你是在祝福.
其實也在團結我們這群.
有心智去服侍弱勢社群的基督徒.
學習主耶穌基督的榜樣.
去親身到人群當中去關懷.
去愛.

$^{681}$和軟弱有需要的人同行.
我開始籌款的時候.
很快就有一個家庭回應.
我很想把這些見證告訴你們.
我口信告訴他們.
我收到你的奉獻.
那個太太跟我說.
你不要客氣.
能夠去做一些我們做不到的事.
原來這個太太也是剛失去工作.
他們奉獻給我之後.
無奈其實他們自己也去了斷宣.
他們就是馬其頓教會.
在有限的裡面.
仍然超出自己能力的慷慨.
主動甚至乎懇求可以有份參與.
慷慨的捐獻不是有能力沒能力的問題.
是保羅叫你不要將神的恩典.
自己代代平安.
我們人要思考怎樣去回應.
其實教會也是一樣.
你的奉獻不是要看你夠不夠錢用.
教會我們的事工能否給你利益.
我們牧師傳道人也不是要靠你.
讓你的奉獻開放.
慷慨的奉獻是你願意投資在神的國度裡面.
這個恩典的慷慨是生命的問題.
哥林多後書八至九章.
馬其頓人跟隨基督的榜樣.
最重要你要記得的.
是提醒我們每一個不要忘記神的恩典.
我們一起低頭祈禱.
參與主耶我們的神.
多謝你給我們你的話語.
叫我們不要忘記你的恩典.
以致我們也成為你恩典的出口.
感激祈禱.
奉耶穌得性命求.
Amen.
大家安靜思想或是需要祈禱.

$^{721}$之後可以散會.
願主祝福大家平安.
\newpage



\section{}
\label{sec:pxG7cdEbvuY}
\textbf{What Christ Thinks of the Church (Revelation 2-3) - Rev. Mark Chan}
\newline
\newline
連結: \href{https://youtube.com/watch?v=pxG7cdEbvuY}{\texttt{https://youtube.com/watch?v=pxG7cdEbvuY}} ~~~~ 語音日期: 2024-09-05
\newline
\newline
\hyperref[sec:7j_P4EWgCCI]{\small{< < < PREV SERMON < < <}}
~
\hyperref[sec:index_chronic]{\small{[返順時目]}}
~
\hyperref[sec:index_scriptual]{\small{[返順卷目]}}
~
\hyperref[sec:YG4nAX61j6k]{\small{> > > NEXT SERMON > > >}}
\newline
\newline
$^{1}$All right, today we will start a new series..
Last weekend I was not with you..
I know Pastor Sam Johnson was here.
and I really appreciate for his World Ministry Award.
and I was in Minnesota for the weekend.
and from Friday to Sunday, yes, Friday to Sunday,.
I was speaking in the church retreat over there.
and as always, sometimes I go to different churches.
and get to know some of the people.
and sisters and brothers in Christ.
and learn about their stories.
and I was encouraged by their experiences.
and one of the couples who host me.
and pick me up from the airport.
and drive me to the retreat center.
which is two and a half hours away.
from the airport from Minneapolis.
and they are such an inspiration, you know,.
just retire and they have been working as expats.
in China, in Hong Kong, in Singapore.
and very successful people and very resourceful people.
and they have been working tirelessly for the gospel,.
preaching the gospel, you know,.
shepherding the church, helping out with ministries.
and they just settled down in Minnesota.
for less than a year and they are working.
and they are one of those faithful people.
who continue to inspire you, you know,.
in the small church and they continue to serve faithfully.
and use their resources to advance the ministry there..
So I was really, really touched by this couple.
and of course, by many other brothers and sisters..
This church is without a pastor for some time.
and as you can understand, all the elders.
and all the lay leaders have to just contribute.
a little bit more, you know, stretch themselves.
and serve in the church and some of the brothers.
because they're without a pastor,.
so they have to preach..
They have to learn how to preach..

$^{41}$So that's pretty tough for them.
but I'm so inspired by the ministry..
Can you imagine, you know, there's no pastor,.
there's no one preparing the message for Sunday service.
and you have to take turn, right?.
And you have to take turn next time you're gonna preach,.
next time you're gonna preach, right?.
It's pretty tough but, you know,.
I'm really inspired by their faithfulness.
and that's what got me thinking about.
this next series of messages.
and what should I preach on..
You know, during the retreat, I have been thinking a lot.
and praying a lot about what should be the sermon series.
from this Sunday onward,.
at least for the next six, seven Sundays.
and I've been thinking about, okay, New Testament letters,.
maybe Old Testament, maybe the Book of Daniel.
because I've been working on the Book of Daniel.
for quite a long time because in November,.
I'm going to another church to preach in Toronto..
So I've been working on the Book of Daniel..
So, you know,.
checking this passage and that passage.
but, you know, in my prayer, in my desire,.
I come to this point that I want to.
focus on one topic for all of us,.
one topic for all of us,.
is that what Christ thinks of the church, right?.
Now, we come to church, of course,.
we carry our own opinions and preferences.
about the church, right?.
Or what is my role in the church, right?.
Maybe I see myself as such a, you know, just attendee,.
one of the attendees, right?.
I'm not part of this, right?.
I don't want to get involved or something like that..
What is my role in the church, right?.
You can ask, you know, each and every one of us.
and we will give you different answers, right?.

$^{81}$But, you know, ultimately, what matters most.
is not how you and I see the church, right?.
What matters most is what Christ thinks of our church, right?.
He is the head of our church, right?.
He is the leader..
It's not so much about what the passage says,.
it's not so much of the core team member.
that the lay leader said,.
it's not so much of what you say.
or how you see the church..
And your role in the church.
is also depends on this question,.
what Christ thinks of your role in the church..
And this series takes us to the last book.
of the New Testament, of the entire Bible..
You know, the Old Testament have 39 books, right?.
Old Testament, 39 books..
And the New Testament has 27 books, right?.
It's a collection of 27 books and letters..
The last book of the entire Bible.
is called the Book of Revelation..
That is the last book of the Bible..
And this Book of Revelation contains seven letters.
written to seven churches in Asia Minor.
at that time, in the first century..
Now, as we come to the Book of Revelation,.
I remember when I was so much younger,.
you know, I wanted to know what is going on in the future,.
right?.
Kind of like apocalypse, you know,.
kind of, you know, doomsday scenario..
What is going on, right?.
And what does the Bible have to say.
about the end of the world, right?.
So I flipped the pages to the Book of Revelation.
and I read some of the passages..
And I must admit, I could not understand it..
It's just so strange, right?.
It opens up another, you know,.
another box of mystery..

$^{121}$And I did not know how to interpret it..
I did not know how to read it..
And that is the Book of Revelation, right?.
Now, first of all, when we come to the Book of Revelation,.
please understand that the Book of Revelation.
is not so much about prediction of the future, right?.
It's not weather forecast,.
whether it's gonna rain tomorrow at 3.23 p.m..
That's not what Revelation is all about..
That's not what the Book of Revelation is all about..
What it is about is that it was written.
to a particular situation,.
to a group of Christians at that time..
And that was the late first century, about A.D. 90s..
You know, the time of Jesus died and rose again..
It was about A.D. 30, right?.
About A.D. 30..
And then the apostles and the early Christians,.
they went from Jerusalem to other parts of the world,.
including Paul himself, right?.
Went to different cities and towns.
and preached the gospel and established churches, right?.
Small churches, small group house churches..
Now, that was about 30s, 40s, and 50s, right?.
And then we come to the Book of Revelation.
is the end of first century..
You know, if you calculate the years, the number of years,.
by this time, the end of first century,.
the second generation Christians are.
almost 70 years after the time of Jesus..
And at that time, particularly at that time,.
the Roman Empire, they were still under the Roman rule,.
right, they have been promoting, you know, emperor worship..
You know, you can imagine, if you,.
if you have a time machine,.
we can go back to the first century, the late first century..
If we go to one of the house churches,.
they have the political pressure from the top, right?.
From the emperor, from the Roman Empire..
They want the churches to worship the emperor, right?.

$^{161}$You can imagine walking to the church.
in the late first century,.
situation is kind of like the church is always.
under the pressure of the government..
If you join the late first century church.
in the Roman Empire, probably you will find some people,.
they have to sing the national anthem.
before they sing praises to Jesus..
You see what I'm saying?.
Do you, will you like it?.
If you walk in this church, you have to sing.
the national anthem before you do anything else, right?.
Or you will find the picture of the emperor,.
you know, next to the cross,.
or right at the center of the sanctuary,.
and the cross is put on the side, you know?.
And that was the kind of situation the early Christians,.
the first century Christians were facing..
They were facing wave and wave of persecution..
And of course, the pagan people,.
if they don't believe in Christ,.
they were okay with that, right?.
Because emperor was just one more god.
on top of many other gods, right?.
But for churches, it's not okay..
Because, you know, the gospel,.
the gospel, the good news of Jesus Christ.
proclaims that there's only one king in this universe,.
and that is Jesus, who died and rose from the dead, right?.
And for this reason, Christians are often, you know,.
being marginalized and being pressured,.
or being persecuted..
And John the apostle, you know,.
Jesus had 12 apostles, right?.
And the last one who died should be John the apostle..
And he wrote the book of Revelation.
under the inspiration of God..
And he said this,.
right off the bat in chapter one, verse nine..
It says, "I, John," it should be apostle John, right?.

$^{201}$"Your brother and companion in suffering and kingdom.
"and patient endurance, there are hours in Jesus..
"I was on the island of Patmos.
"because of the word of God.
"and because of the testimony of Jesus.".
So he was being persecuted,.
and he was exiled to a small island called Patmos.
in the Aegean Sea..
He said, "On the last day,".
which means he was worshiping, right?.
On the last day, on Sunday,.
"I was in the spirit, and I heard behind me.
"a loud voice like a trumpet,.
"which said, 'Write on a scroll what you see.
"'and send it to seven churches,.
"'to Ephesus, Smyrna, Pergamon,.
"'Tirethea, Thyatira,.
"'Sardis, Philadelphia, and Laodicea.'".
So seven cities, there were churches in these seven cities..
And John's being commanded to write letters.
to each of the seven churches..
And where are they?.
You see the map..
You see all of the seven cities and seven churches.
are located in the western part of today's Turkey, right?.
And you can see the island of Patmos, right?.
Can you see that?.
Island of Patmos..
And John received the revelation,.
which is written in the book of Revelation,.
on this island because he was a political exile,.
and the Romans were going after him..
And during his suffering, during his isolation.
on the island of Patmos,.
he received the revelation from God..
And he wrote seven churches,.
wrote seven letters to the seven churches..
And of course we know seven is a symbolic number, right?.
It's a complete number, right?.
I'm Chinese..

$^{241}$The symbolic number eight, you know,.
it's very important, right?.
You drive around in the city,.
you see a license plate with 888 on it..
Most likely it's a Chinese driver, right?.
Chinese owner..
For the Hebrew culture, or even for the Christian culture,.
seven is the complete number, right?.
Seven is the complete number..
So that means the seven churches.
represent all the churches, right?.
And the first few chapters of the book of Revelation.
contain letters addressed to the particular.
faith communities and apply to them, right?.
And this letter, we shall see,.
we will read letter after letter,.
and we shall see the format of each letter is similar..
Christ, first of all, he will introduce himself,.
and then he will give an evaluation of the church,.
saying that I know what is going on in your church..
And then, thirdly, the letter, which is very short,.
Christ, Jesus, will advise the church.
to change the course of action,.
or keep on doing some of the good works.
they have been doing..
And then the letter will close with a promise,.
a promise particularly relevant to the local church,.
and advise the audience to listen and to obey the word..
To obey the message..
So today, we're going to go through the first letter first..
Okay?.
The first letter is written to Ephesus,.
to the Christians in Ephesus..
It says chapter two, verse one..
"To the angel of the church in Ephesus write,.
'These are the words of him who holds the seven stars.
in his right hand and walks among.
the seven golden lampstands.'".
So these are all symbols, right?.
Symbolic language..

$^{281}$And they represent some of the reality, right?.
And Christ introduced himself as someone.
who holds the seven stars in his right hand, right?.
And what does it mean?.
You can try to imagine Christ himself,.
who is king over the entire cosmos,.
and his right hand is holding seven stars, right?.
Now we know, in the late first century,.
if you go to different archeological sites,.
and you find some of the coins mint by the emperor,.
and one of those coins, if you look at it,.
there is the image of the emperor in the center,.
and there are seven stars surrounding him..
And that emperor is called Domitian..
And he ruled the Roman Empire and persecuted Christians.
in the late first century..
So verse one of chapter two says,.
"It is not Domitian, it is not the emperor.
that holds seven stars..
Christ is the one who holds seven stars..
Christ is the one who holds the entire cosmos.
in his sovereignty..
He is someone who is in absolute control.".
And Jesus introduced himself as someone who is walking.
among the seven golden lampstands..
Now the golden lampstand, called menorah,.
is one of the oldest symbols of Judaism, right?.
You can see the picture on the right..
When you go down to the Western Wall in Jerusalem,.
you see a huge, life-size golden menorah.
to remember this symbol, right?.
This represents the nation of Israel, the people of Israel..
And this is applied, now this menorah is applied.
to the people of God in the church..
Now for the Israelites, or for the Jews,.
menorah reminds them they are being called by God.
to be the light unto the nations, right?.
Israel, or the Jewish people, are being elected by God..
It's not because they're better..
They are just like other pagan nations..

$^{321}$But they are specially elected by God to be the channel,.
to be the instrument of God among the nations..
And they will be like the menorah..
They will be the lampstand to shine the true light.
to all the nations around them, right?.
And that was Israel was supposed to be..
However, here, Christ is saying that the church.
is the true Israel, right?.
The church is the true Israel, composed with believers.
who are Jewish and who are Gentiles..
So Jews and Gentiles together, they proclaim that Jesus.
is the true Lord, and now they are like the menorah, right?.
They should be used by God to shine the light.
to the nations..
And Christ introduced himself in such a way,.
and he command the church in Ephesus to do such and such..
Now remember the format of the letter?.
This is the first letter to the Ephesian Christians, right?.
After introducing himself as someone who holds the seven.
stars and walk among the seven lampstands,.
in verse two, Christ says, "I know your deeds,.
"your hard work, your perseverance..
"I know you cannot tolerate wicked men,.
"and you have tested them, those who claim to be apostles.
"but are not apostles, and found them to be false..
"And you have persevered and have endured hardship.
"for my name, and have not grown weary.".
So is it praise or is it condemnation?.
It's not praise, right?.
It's very positive, right?.
So Christ looking at the church in Ephesus,.
and he praises them for their hard work,.
for their good work..
So the Ephesian Christians are very, very busy.
and hard-working Christians, and they build a ministry..
They block by block, right?.
Break by break..
And the Christians in Ephesus have been trying.
to building the body of Christ..
And by the time John writes the letter of Ephesians,.

$^{361}$by the time John writes this letter,.
the book of Revelation, it was in the late first century,.
right, 1890, and the Ephesian church.
was already 40 years old, right?.
Yeah, the church was already 40-some years old, right?.
Remember I told you last weekend I went to the other church.
in Minnesota..
They have 40 years of history..
They have a long history, right?.
Much longer than our church, right?.
Our church is only eight years old, right?.
We are like a little kid, right?.
We are not grown, right?.
They have, right?.
And the church in Ephesus is also a grown church..
And you see, in this community, there are many, many people.
who are working so hard to build a ministry..
When we first started Urban Voice, it was hard work..
It was very hard work..
Physically difficult because you have to set up.
and you have to do running around..
And mentally, it was challenging..
Spiritually, it was challenging..
Emotionally, it was difficult..
I really appreciate, eight years ago,.
many of the leaders who started this church with me,.
and I really appreciate their perseverance,.
their hard work, right?.
When we first started this ministry, we wanted to raise,.
we had nothing back then..
We wanted to raise, well, we calculate some of the numbers,.
right, we punch in some of the numbers,.
this is what we need, right?.
So we had to raise, in order to start this church,.
we had to raise 120K to get us started..
I remember it was late 2015, and we had zero, we had zero..
But within six weeks, we raised that amount, right?.
Praise God for the commitment, right?.
And the perseverance, the contribution of the people.
who started the church, right?.

$^{401}$And Christ will get the church in Ephesus..
Yes, you guys are very, very good, you know, way to go..
Your perseverance, the hard work, right?.
And I wanna encourage all of you..
We still need to do all the hard work..
We still need to work harder to persevere, right?.
Because all we do is the good work,.
we try to hold fast to the gospel,.
we try to preach the gospel, we try to build the church,.
and it's not an easy task, right?.
And there's no quick fix, right?.
All ministries require hard work..
And you see the Ephesian Christians,.
not only they work hard, but they know the truth, right?.
They know who are the, who the wicked men are, right?.
Some of the apostles, some of the people.
who claim to be apostles, but they are not..
Indeed, they are not apostles..
You know, we find some of the early documents.
that we know some of the early churches,.
they're trying to tell which guy is the true apostle.
and which guy is not the true guy, the false apostle, right?.
And some of the criterias they lay out.
for the leadership team is to test whoever come our way.
to preach the gospel and their behaviors..
And one of the criterias says,.
if the preacher, the itinerary preacher,.
comes to your church and you're gonna host them, right?.
You're gonna host this preacher,.
and if you host him for three days,.
and he ask to be staying longer,.
the document says he is the false apostle,.
because he's trying to take advantage of you..
And after preaching, if he ask for money,.
he is definitely a fake apostle..
And we have some of the early documents.
that help the early churches to discern.
who is true apostle and who is false apostle..
And for the church in Ephesus, they tested them.
and know their doctrines, right?.

$^{441}$They know what they believe in,.
so that's why they can discern what is right.
and what is wrong, what is orthodox and what is cultic,.
what is not orthodox..
And Christ praises them..
And this letter continues, after praising them,.
Christ said, "However, I hold this against you..
"You have forsaken your first love..
"Remember the height from which you have fallen..
"Repent and do the things you did at first..
"If you do not repent, I will come to you.
"and remove your lampstand from its place.".
This is a very stern warning, right?.
Yes, they are hardworking..
However, there's one thing Christ holds against them.
is that this is a church that has abandoned love, you know?.
Abandoned love for God, abandoned the love for others, right?.
And it says, what kind of love?.
First love, right?.
First love..
Somewhere in the past, you had that, right?.
First love, talking about love in the past..
Now, we all have similar experience, right?.
For those of us like me, you know, I have been in Christ,.
I have known Christ and served Christ for many years..
In my faith journey, there were a time.
I was really touched by the love of God.
and I was motivated and compelled by God's love..
So I'll do everything out of my experience of God's love.
and I'll do everything out of love, right?.
I want to care about other people, right?.
And not asking for some benefits..
I want to invite other people to the church..
I want to invite other people to experience.
the grace of God, all because I have taste.
the love of God, right?.
The redeeming grace of Christ..
But the church in Ephesus has lost it all..
How old is this church, remember?.
40 years, right?.

$^{481}$More than 40 years..
And this is kind of a mature church..
But Christ says, "I have this against you..
"You have forsaken the first love.".
You remember there were point in your faith journey.
that you're so compelled by the love of Christ, right?.
Maybe when you were younger, maybe there were.
different period of time in your past journey in faith..
You were so touched by the love of God.
and you are willing to do everything because Jesus loves you.
because you know that God loves you, right?.
The Christians in Ephesus have forsaken the first love..
And these are very, very strong words..
And it doesn't mean that they're being forgetful,.
like you lost a wallet or something..
You have forsaken your first love..
These are very strong words..
It's like intentionally you trash it..
There's something very valuable,.
but you throw it into the trash can..
You say, "I don't care anymore.".
Right?.
That kind of action, that kind of attitude.
that they're holding..
There's no wonder why that Jesus write this letter.
to the church and warns them,.
"Hey, if you do not recover this first love,.
"I will remove you..
"You are no longer my lamp stand.
"and you cannot shine the light to the nations.".
Jesus continues..
He says, "Remember the height from which you have fallen.
"and you have to repent.
"and do the things you did at first..
"If you do not repent, I will come to you.
"and remove your lamp stand from its place.".
These are very, very strong warning words..
If the church is without love,.
it is no longer a church, right?.
You agree with that?.

$^{521}$If a Christian is no longer motivated by the love of God,.
he is not a follower of Jesus Christ anymore..
Right?.
You're just going through some motions, right?.
Routine, right?.
Muscle memory..
Jesus said, "If you do not recover that,.
"you cannot be the lamp stand.".
If we do not recover that, we cannot be God's church..
I remember some years ago,.
I attended a regional meeting for pastors.
in the network in Chicago..
There were about 30, 40 pastors and church planters..
We gathered together in north side of the city..
We invited a Jewish rabbi to give a testimony..
He's a very, very old rabbi..
He believed Jesus is Christ,.
and what we call Messianic Jew..
He was invited to give a testimony..
He lived in Skokie, north of us..
I don't know the real number now,.
the number of Messianic Jews in the Skokie community..
Back then, I heard it was about 100,.
about 100 Jewish people who professed.
Jesus is the Christ, the Messiah..
I don't know the real number now,.
but this rabbi, very old and retired,.
and came to faith, came to know Jesus as his Lord,.
and he shared his testimony..
He lived in Skokie, and he had broken marriages,.
and had not seen his daughter for over 10 years,.
because of broken relationship..
In his brokenness, in his pain, in his darkness,.
a Christian came along and cared for him,.
and stayed with him..
It's not imposing the Bible upon him..
It's not using the Bible to bang on his head..
This Christian man just came along to this old man,.
and cared for him, and listened to his story,.
and tried to be his friend..

$^{561}$And that's how this Jewish man came to faith,.
came to know Jesus, all because of love..
You know, I think the biggest threat in our church.
is indifference, indifference..
Indifference, the spirit of indifference..
It says, "I don't care.".
Well, you have Bible study, I don't care, right?.
You have this good work, or that good work,.
you know, I don't care..
You can share the gospel, and preach the gospel,.
build a ministry, build a youth ministry,.
or whatever happening in the church,.
I just don't care, right?.
I just follow motion, and you know, in and out, right?.
I think this is the biggest threat to our life..
Not just the church community, but to every one of us,.
our individual spirituality..
That is the biggest threat, the biggest danger..
It's not because you're anti-gospel,.
it's not because you're anti-Bible, or anti-God..
It's just indifference, right?.
And that's what Christ is saying to the Christians.
in Ephesus..
I hold this against you..
Yes, you are doing very, very good works,.
but this I against you, because this will destroy you..
This will take you, remove you from the lamp stand..
Now, I like the English paraphrase translation.
called Message..
You heard about it?.
Message by Eugene Peterson..
Back in the '90s, I was in Regent College, Vancouver,.
and Eugene Peterson was the professor.
of spiritual theology there..
And he translated the whole Bible,.
and the translation is called paraphrase,.
because it's not a literal translation..
It's not like, kind of like word for word,.
very close to the original language, like ESV, or NIV..
But he did his translation with paraphrase..

$^{601}$It's very loose, but you still get the meaning.
of the phrases and of the language..
And in Message, Eugene Peterson translate this verse.
as this, and let me read this..
Christ said to the church,.
"But you walk away from your first love.".
Why?.
What is going on with you anyway?.
Do you have any idea how far you've fallen?.
A Lucifer fall..
Let's read it again..
Jesus said, "I have this against you,.
"because you walk away from your first love.".
Why?.
What the heck is going on with you?.
Do you have any idea how far you have fallen?.
Do you have any idea how far,.
what the distance of your fall?.
It's a Lucifer fall, right?.
You know, in the Jewish mind, especially in time of Jesus,.
they had this tradition saying,.
you know, Satan fell from the sky.
to the ground, to the earth, right?.
And that is the biggest distance of one's fall..
Because of his pride, he was the angel of light,.
but he was being cast down from heaven.
and fell to the earth..
That was part of the Jewish tradition and their theology..
And they believe Satan, Lucifer, fell to the earth..
He fell in such a way..
And that is the biggest distance of fall.
in history, in the entire cosmos, right?.
No one thought so much as Satan did..
And Eugene Peterson, I think he got it right..
In his kind of loose paraphrase,.
"Do you have any idea how far you have fallen.
"if you have forsaken, you have abandoned the first love?".
It's like the fall of Satan, right?.
It's a huge deal here, right?.
It threatens your whole life, your eternal life,.

$^{641}$your eternal destiny, your position,.
your role in the kingdom of God..
So Christ says, "Remember the height.
"from which you have fallen,.
"and repent and do the things you did at first.".
Please, sister and brother, I want to urge you today..
Remember how far we have fallen.
if we have forsaken the first love..
Remember how the love of God has touched you before, right?.
Perhaps years ago, perhaps some months ago,.
but along the way, you have tasted the real love of God,.
right, but somehow you have fallen, right?.
Right?.
Repent and do the things you did at first, right?.
If you do not repent, Jesus said,.
"I'll come to you and remove you from the lamb stand..
"You're no longer my people..
"You're no longer my lamb, my menorah.".
This letter closes with this verse..
You have this in your favor,.
you hate the practices of liquidation,.
liquid, the Goliathans..
This group probably is the cultic group.
that infiltrate into the church.
and using the gospel as an excuse for sexual immorality..
And that's much we know, right?.
Again, they know the truth, right, the Ephesian Christians..
However, they have fallen from the first love,.
and it has threatened them as a church..
The last verse, "He who has an ear,.
"let him hear what the Spirit says to the churches..
"To him who overcomes, we cover our first love,.
"repent and do the things we did at first..
"I will give you," Jesus said,.
"the right to eat from the tree of life,.
"which is in the paradise of God.".
All right?.
You are the menorah, and we are the menorah..
The church is the menorah..
But in order for us to play this role,.

$^{681}$to be the light to our neighbors, to the world,.
we need to recover our first love..
Let me close by asking you,.
when was your first love with God?.
When did that happen, right?.
Maybe you were a student..
Maybe before you had a family, right?.
Somewhere along the line, you know..
Where is that first love?.
And Jesus said, "We cover that.".
You know, we cover that..
We repent and change our course..
Let's all pray together..
I want to spend some minute to pray to God,.
to remember the heights from which we have fallen..
Repent and do the things we did at first..
Our loving God, how much do we call?.
We call in our past, sometime in the past..
We were touched by your grace, touched by your love..
But somehow we have forsaken it..
Help us to remember, remember that your love is still real..
Your grace is still glorious..
Your love is so authentic, and still changes us..
It still changes us..
Remind us where we have fallen..
We come before you to repent..
Thank you for being our loving God.
and forgiving our sins..
Forgiving our forgetfulness,.
forgiving our intentional abandonment of your grace..
Help us to repent..
And help us to redo the things that we did at first..
Help us to recover that first love,.
that pure love, that pure motivation..
Purely compelled by the love of God..
To serve you, to follow you..
God, I pray that you have mercy on us..
Mercy on our church..
Help us to remember, to repent, and redo..
And thank you for being our gracious Lord..

$^{721}$And you are always there, walking among us..
And know our situation..
Yes, we come before you..
We change and we repent..
Help us to taste your grace and your love once again..
Anew..
Renew our spirit,.
so that we can be your faithful followers..
To continue to serve you,.
continue to love you,.
and to love other people..
In Christ's name, we pray..
Well, thank you for worshiping with us..
I hope to see you next Sunday,.
when we continue to serve God with our first love..
God bless you..
\newpage



\section{}
\label{sec:YG4nAX61j6k}
\textbf{基督眼中的教會 (啟示錄 2-3) - 陳崇基牧師}
\newline
\newline
連結: \href{https://youtube.com/watch?v=YG4nAX61j6k}{\texttt{https://youtube.com/watch?v=YG4nAX61j6k}} ~~~~ 語音日期: 2024-09-05
\newline
\newline
\hyperref[sec:pxG7cdEbvuY]{\small{< < < PREV SERMON < < <}}
~
\hyperref[sec:index_chronic]{\small{[返順時目]}}
~
\hyperref[sec:index_scriptual]{\small{[返順卷目]}}
~
\hyperref[sec:U3A7_wArCTw]{\small{> > > NEXT SERMON > > >}}
\newline
\newline
$^{1}$早晨 各位姐妹.
戴著口罩也這麼有活力 真厲害.
如果教會多一些這樣的人真好.
多謝主人.
我剛剛在上個週末不在你中間.
因為我和Mrs. Sutter在教會說了五場.
在一個週末 他們在一個退休會裡面.
也挺好 給我自己也有一點空間去思想.
當然也很忙 去港島.
也去認識當地教會的弟兄姊妹.
每逢去到不同的地方 去到不同的教會.
有時候去跟弟兄姊妹認識的時候.
我覺得很大鼓勵.
接待我那對夫婦 是一對很愛主的夫婦.
他們去了Mrs. Sutter 大概一年左右的時間.
定居了 然後在教會投身.
看到他們 聽到他們很多不同的故事.
以往他們在大公司工作.
在不同的地方 去新加坡 去北京 去香港.
等等的地方 做了很多見證.
還有去到事奉 這些真的是很大的鼓勵.
告訴我們我們心靈裡面有什麼更加要緊的.
沒有什麼更加要緊 是事奉上帝.
事奉上帝是更加要緊.
在我們自己的崗位裡面能夠事奉.
所以我自己也有很大的鼓勵.
我希望大家來到教會也是一樣.
我們來到教會 當然有不同的意見.
有不同的看法 或者我們有不同的偏好.
但最重要的是我們要明白教會是一回事.
如果了解教會是一回事的時候.
我們就了解我們自己是一回事.
為什麼我們來教會.
為什麼我們會在這裡生活.
為什麼會在這裡嘗試去活現出.
到底什麼叫基督徒.
信耶穌是一回事.
為什麼要參加一個教會 一個群體.
有什麼目的等等.
這個就是我想在接下來的六七個星期.

$^{41}$我和大家來思考的一個問題.
嘗試去想一下 到底教會是一回事.
當然我們有很多不同的看法.
但最重要的是教會的主怎麼看待我們.
教會在上帝的眼中是怎麼樣的.
上帝對教會有什麼預期.
上帝對教會有什麼吩咐.
這個才是最重要.
不是我講 不是你講.
而是回到聖經上帝的話語.
基督自己怎麼看待教會.
這個主題叫做基督眼中的教會.
What Christ thinks of the church.
大家有沒有想過這個問題.
很多時候我們進到一個教會.
或者去到一個群體的時候.
我們想起的就是我們怎麼想教會.
但到底基督是怎麼去看待教會的.
基督怎麼看待你 怎麼看待我.
所以在接下來的六七個星期.
我們會去看一段經文.
在啟示錄第二至第三章.
如果那些比較老練的基督徒.
你知道這兩章聖經就是七教會書信.
就是基督寫給當時七個的教會.
這個是我想的講題.
我去退休會 去面試退休會.
在想到底要講些什麼.
當然最好就是講新約書信.
因為我星期二 是不是星期二開始.
我的網絡交際星期二晚上開始.
講新約書信.
那我就想不如講新約書信.
但想著想著一直去祈禱.
一直去等候的時候.
我都覺得可能回到這七卷的.
七封的信 是基督當時寫給教會的.
是出自啟示錄 如果你有聖經的話.
打開啟示錄 啟示錄是最後一本書.
聖經有六十六卷書.

$^{81}$舊約有三十九卷 新約有二十七卷.
最後的一卷叫做Revelation 叫做啟示錄.
啟示錄是什麼呢.
很多時候我們 特別是華人的信徒.
特別是那些比較老派的信徒.
我們經常都搞錯.
我們當啟示錄是預先寫好的新聞報導.
這個是錯得非常交關.
到底啟示錄是做什麼的.
啟示錄其實就是想告訴我們.
當時的教會的處境.
針對當時的教會處境.
上帝要將一些最終極的事情.
來啟示出來.
最值得我們緊張的.
最值得我們關心的事情.
去啟示出來.
最主要寫給誰.
不是寫給我們今時今日二十一世紀的人.
寫給當時的人.
所以你看啟示錄的時候.
一定要想當時的處境.
他們在預見什麼.
於是啟示錄就用當時他們能夠明白的方法.
來跟他們溝通.
這個是啟示錄.
不是預先寫好的新聞報導.
我們打開第二至到第三章之前.
第一章說什麼呢.
他們發生什麼處境.
第一章第九節.
寫信的人叫約翰.
相信就是使徒約翰.
十二使徒裡面剩下的那個.
死剩種是約翰.
可能是開始跟耶穌的時候最年輕的那個.
或者他遭遇不同的事情.
所以他就是剩下的那個.
我約翰就是你們的弟兄在耶穌裡面.
跟你們一同在患難國度忍耐裡面有份.

$^{121}$為上帝的道.
和為耶穌作的見證.
曾在那一名叫拔摩的海島上.
有一天我被聖靈感動.
聽見在我後面有大聲音吹好國說.
要把你所看見的寫在書上.
基給 以弗所 史梅蘭 別加摩 推亞推拉 撒迪 飛拉鐵飛 魯迪加 七個的教會.
當時到底發生什麼事.
這裡大概有些蛛絲馬跡告訴你.
約翰在哪裡呢.
在一個叫拔摩的海島上.
當時羅馬帝國對於一些.
特別是在政權上挑戰他們的人.
或者一些所謂的政治犯.
就會把他流放到希臘那些小的海島那裡.
大家可能有去過希臘玩.
那些海島很美的.
以前不是這樣的.
以前不是去那裡度假.
以前是流放你去那裡.
exile你去那裡.
那你就不要那麼多聲音.
當時發生什麼事呢.
大概是第一世紀末的時候.
大概是第一世紀的九十年代.
九十年代發生什麼事呢.
當時有一個羅馬帝王叫做Domitian.
當時我們知道那個處境.
大概就是用很多的壓力.
去給當時的很少的教會.
不同地方的教會.
大家知道主耶穌來的時候大概是什麼時間.
大概主後三十年死而復活.
主後三十年的時候.
所以耶穌過了之後.
有第一代的信徒.
然後有第一代的教會出現.
到這個時候主後大概九十年代的時候.
已經是第二代的信徒出現了.
所以約翰其實是第一代最後最後那個.

$^{161}$所以我們叫他死正種.
剩下那一代最後的一個見證人.
見證耶穌曾經做過的事.
見證耶穌所建立的天國.
他繼續將福音傳開.
所以第二代的信徒在教會那裡.
他們傳福音的時候.
他們傳的福音是什麼呢.
我們想想傳福音是什麼.
我們想想神愛你等等.
真正福音是什麼.
聖經裡面講福音.
是宣告一個王者來到.
That is the gospel.
And there is only one king.
只有一個我們值得去歌頌.
一個值得我們跟從的人.
就是耶穌基督.
如果你將這個福音.
在當時羅馬帝國傳出來的時候.
你會很不滿的.
因為當時的羅馬帝王.
譬如Domitian.
他會認為什麼.
你這是造反.
是煽動或危害國安等等.
所以這些傳福音的人.
包括約翰其中一個.
就是被當時的羅馬帝國.
認為是威脅當時的安危.
國家的安危.
就將他流放.
有一點像今天的文字獄.
如果你報導什麼.
是影響我們羅馬帝國中央的.
我們就將你去.
將你去XR.
所以你傳福音的時候.
一定會被排斥 一定會很艱難.
當時的教會就是遇見一個這樣的處境.

$^{201}$所以為什麼才會有啟示錄出現.
不是預先寫好.
明天三點鐘一定會下雨.
不是預測.
而是最主要告訴我們.
真正的真相 這個世界的真相.
誰才是真正的帝王.
特別當時的羅馬帝國.
在地中海的東部.
他們去宣傳一種叫帝王的崇拜.
設立一些廟宇.
那些廟宇有羅馬帝王的像.
你就要上去供奉他.
特別到Domitian的時候.
主後90年代的時候.
Domitian甚至說他自己是Living God.
我們怎麼知道呢.
因為現在我們找回Domitian當時的硬幣.
我們今天不用現金了.
記得我們那些比較老的人.
我們還用現金.
家裡有很多現金.
特別那些神沙.
那些硬幣.
每一個硬幣都是政治宣傳的工具.
找回當時Domitian那些硬幣.
就說他是Son of God.
說他是Lord and God.
主和神.
所以啟示錄就用很多這樣的語句.
來反駁到底真正的帝王.
真正的掌管宇宙的是誰.
不是Domitian 是耶穌基督.
所以在當時特別羅馬的東邊.
Roman East.
地中海的東邊的世界.
即是今天的土耳其.
希臘.
再過去以色列巴勒斯坦地.
他就推行很多這樣的Emperor Worship.

$^{241}$Imperial Cult.
Imperial Religion.
就是要那些掌管帝國的人.
被供奉和被敬拜.
如果你是教會.
有沒有麻煩?.
都幾麻煩.
你想像一下某些地方.
那個教會裡面.
首先你唱讚美詩.
唱給耶穌之前.
你一定要唱國歌.
即是敬拜要練甚麼?.
天天都要唱國歌.
練好一首國歌.
唱完一首國歌你才好唱給耶穌聽.
或者你看著台.
或者看著教會.
很多時候教會有一個標誌是十字架.
十字架的旁邊.
就放帝王的相片.
或者帝王的相片放在中間.
十字架放在旁邊.
就是這種情況.
如果你不聽話.
就好像約翰一樣.
被流放到拔魔的海上.
這個七教會的書信.
就是寫給當時在現在的土耳其.
七個地方的教會.
你看到這個地圖.
剛才我們不是讀了很多地名嗎?.
記不記得讀了甚麼地名?.
伊弗所你記得了嗎?.
你記得伊弗所.
如果你忘記了七個教會的地方名.
或者教會的名字.
你記住伊弗所就可以了.
其他的不是經常出現的.
由伊弗所開始.

$^{281}$直到第七個是誰?.
魯迪加.
這七個教會.
所以七封書信.
就會按照當時的地圖.
大概那個路線.
可能是郵政的路線.
就是寄信的時候.
郵差派不同的地方.
那個路線.
有一點像Amazon出的地圖.
跟著去.
你知道我家裡.
買了Amazon的東西.
我是最後那個.
不知為何.
每次都是這樣.
永遠都是最後最後才到我.
我就是魯迪加教會.
七封書信.
針對當時的處境.
七個城市裡面.
有不同的House Church.
七封書信都有相似的格式.
首先基督說他自己是誰.
用不同的稱呼來介紹自己.
然後他就說.
我知道你教會發生什麼事.
因為上帝當然知道.
我們教會發生什麼事.
所以就會說出.
七個教會的所謂Evaluation.
我怎樣看你們.
然後第三件事就是.
經常有一句話.
你要做什麼.
你要改變什麼.
或者你要繼續去做什麼.
這個就是對教會的忠告.
然後最後有一個Promise.

$^{321}$凡有耳聽的.
必須要聽.
上帝給你的Promise.
忍耐到底.
你繼續堅守真道的話.
你就會OK的.
有Promise給你.
今天我們看第一封書信.
二章第一節.
說什麼呢.
寫給二佛所.
他說你要寫信給二佛所教會的使者.
那右手拿著七顆星.
在七個金燈台中間行走的者用說.
寫給二佛所.
二佛所當時是一個很大的城市.
大概第三第四大城市.
你知不知道芝加哥是第幾大城市.
三.
三是嗎.
所以為什麼DNC都來.
阻礙我們.
當時二佛所都是一個大的海港.
就是Metropolitan City.
正因為大城市的緣故.
所以是經濟文化宗教的命脈.
當時不同的大城市都想京城羅馬.
來代他們好一點.
要討好他們.
如果你見一個帝王的廟在我們中間.
這樣就好了.
為什麼.
那些資金和經濟.
宗教.
就會更加蓬勃.
所以那些城市.
特別是今天土耳其.
那些主要的城市.
經常都爭誰能夠做到.
省份裡面最大最大的那個.

$^{361}$二佛所當然就是其中一個.
你記得在《少林寺傳》裡面說到.
保羅去到全福音的時候.
他在一個宣教旅程.
他來到二佛所.
並且他在那裡留了三年的時間.
我相信保羅也看中了這個地方.
是一個很重要的地方.
所以在那裡留了三年.
去到牧羊教會.
和在那裡差派其他宣教士.
做很多宣教的工作.
當時寫二佛所書的時候.
在說二佛所書.
不是啟示錄.
二佛所書.
你知道新約裡面其中一本書信就是二佛所書.
大概保羅寫的時候.
寫給二佛所的教會.
大概是主後六十年的時候.
啟示錄什麼時候寫.
主後九十年.
所以二佛所教會應該已經有四十年左右.
四十年左右的教會.
也不算是新的教會.
因為我們教會是baby.
就是八年而已.
二佛所四十年這麼久.
一個比較有歷史.
有根基的教會.
他寫信給他.
他說右手拿著七星.
在七個金燈台中間行走.
這裡說一些符號.
說一些象徵.
為什麼會拿著七星.
你可以想像.
試試在腦海裡面畫一幅這樣的畫.
基督是拿著七火星.
是一幅怎樣的畫.

$^{401}$可能我們想到的就是很浪漫.
我知道有些弟兄姊妹去看流星雨.
特意去黑漆漆的地方去餵蚊子.
然後去看流星雨.
當時他們想起七星是什麼.
剛才我不是說那些硬幣嗎.
你知道Domitian其中一些的硬幣.
他所印出來的硬幣.
鑄造出來的硬幣.
就是有羅馬帝王的意象在裡面.
然後就有七顆星圍著他.
羅馬帝王才是拿著七顆星.
基督說不是.
基督說我才是右手拿著七星.
就是說他掌握著世界.
並且掌握著七個教會.
當然七也是象徵性的數字.
寫給七個教會自然是寫給所有的教會.
他說他知道所有事情所發生的.
拿著七星並且在七個金燈台中間行走.
每個教會可以當成一個金燈台.
金燈台是猶太教最悠久的符號.
金燈台你見到這張相.
你快要下西牆的時候可以有散.
就是一個大的金燈台放在你面前.
這個是以色列國的符號.
是他的象徵.
代表以色列應當成為萬國之光.
以塞俠書有這樣說.
上帝揀選以色列不是因為他們特別好.
可以在別人面前威風一點.
上帝揀選以色列成為他的工具.
使萬國能夠得到真光.
現在應用在誰身上.
應用在教會.
教會才是真正的以色列.
真正的以色列由耶穌基督所建立的.
猶太人在裡面.
信主猶太人.
信主和外邦人都在教會裡面.

$^{441}$所有我們這班人.
就是真正的以色列.
真正的以色列.
我們就是真正的金燈台.
所以有時候我們祈禱的時候.
我們教會成為一個燈台.
什麼意思.
就是我們應該成為萬國的光.
Light unto the nations.
基督說我在金燈台中間行走.
就是我知道你們所有發生的事.
並且他們都有權將你放在燈台上.
有權將你在燈台上拿下來.
這個就是七教會的書信裡面.
最有挑戰性的地方.
你繼續看下去.
這封書信說什麼.
他說我知道你的行為.
勞碌忍耐.
也知道你不容隱惡人.
你也曾經察驗那自稱為使徒的.
卻不是使徒的.
看出他們是假的.
你能夠忍耐.
曾為我的名勞苦.
且不困倦.
這是讚還是嘆.
不敢說是不是.
讚.
他說我知道.
我知道你為我的福音.
為教會為上帝天國.
是有行為.
好的行為.
有勞碌有忍耐.
是不是.
我說過我們教會有多久.
八年而已.
我們只是嬰兒.
從第一天開始.

$^{481}$我真的很欣賞.
我們有很多弟兄姊妹.
真的奮心下去.
去建立教會.
建立教會容易嗎.
來教堂就容易嗎.
你去建立一個教會.
建立一個群體.
就不是那麼簡單.
吃飯就容易.
煮飯就難.
建立教會需要很多勞碌和忍耐.
基督去讚這班人.
上帝是體重你的付出.
特別是為了真理的付出.
這裡說到什麼.
他們能夠察驗那自稱為使徒.
卻不是使徒.
即是知道他們是假的.
當時一些書信或留下的文件.
我們知道有些確實是假的使徒.
或是假的牧者.
因為當時很多時候.
一個教會是沒有自己駐堂的牧者.
他們是靠一些牧者到處走.
去完一個城市就去另一個城市.
然後他們就去宣講上帝的道.
或者做教導.
有些文件告訴我們.
他們當時怎樣去察驗.
就是如果這個牧者來到.
例如愛來港元.
他就招待他.
如果招待他三天.
他都不走的話.
那個可能是假的.
或者招待完之後.
問他要錢的.
那些可能是假的.
在夏文裡面我們知道.

$^{521}$有一派人叫做尼哥拉黨.
我們不太確定尼哥拉黨是什麼回事.
總之是一班人.
很可能從一些字裡行間去知道.
他們是一些假的使徒.
並且去到宣傳說.
我們可以放縱情慾.
得到福音或者放縱情慾.
等等等等.
基督去讚以弗所的教會.
他說你們懂得去辯蝶.
並且又忍耐又好的行為.
肯去奮身在那裡勞碌.
這些是很好的事.
我們繼續看下去.
第二章第四節.
然而有一件事我要責備你的.
就算你把起初的愛心拋棄了.
所以你要回想你是從哪裡墜落的.
並且要悔改.
做起初所做的工作.
你若不悔改.
我要到你那裡.
把你的燈台從遠處拿去.
這是一個警告.
剛才讚完了.
現在就來了.
這是一個來幽.
教會可以因為這個原因.
基督說要把你從燈台遠處拿去.
說什麼呢?.
愛心.
不只是愛心這麼簡單.
是起初的愛心.
以前你有過的愛心.
或者你的初心.
初心.
說起初心.
要回想一下.
因為說以前的事.

$^{561}$想一想.
你想一想你的初心是怎樣的.
以前.
曾經為主.
有的熱誠.
曾經主的愛打動過你的.
你會為事奉主.
為真理.
為上帝出聲.
為上帝事奉.
是不在意的.
不介意.
那些叫初心.
真誠的心.
對上帝的愛.
初愛.
對人的初愛.
其他的人.
你記得那種初心嗎?.
那個愛心的你.
你不介意.
你不在意.
你樂意愛身邊的弟兄姊妹.
肯在那個群體裡面投身.
花時間不緊要.
熱主的愛愛大家.
甚至愛一些未信主的人.
去關懷一些未信主的人.
他說以弗所教會發生什麼事.
起初的愛.
拋棄了.
拋棄是很重的語氣.
拋棄不是錢包不見了.
不好意思不小心.
拋棄是你覺得有價值的東西.
但你把它丟進垃圾桶.
那些明明你應該保留的東西.
你丟進垃圾桶當垃圾一樣.
這個叫拋棄.
基督說你要回想.

$^{601}$回想.
Remember.
可能找回你以前的日記.
找回你以前的回憶.
然後你悔改.
做回你起初所做的事.
沒有了起初的愛心還是不是一個教會.
可以從燈台拿走.
我記得有一次我們浸信會有一個聯會.
有一個餐會.
是weekday的.
我就去了那裡.
那裡有很多church planters.
或者其他pastor.
他們請了一個猶太人的拉比.
信主的拉比.
是老人家.
猶太人的拉比.
在那裡說一些他的故事.
這個猶太人的拉比是那些所謂的Messianic Jew.
他們信耶穌是基督.
請了他來說.
他在Skokie住.
說他信主的經過.
你知道在Skokie大概.
我不知道他們的號碼是不是真的.
整個Skokie有很多猶太人.
知道吧.
大概只有百多個.
是信主的猶太人.
信耶穌的猶太人.
他說他的故事是怎樣.
他婚姻失敗.
他的女兒.
親生女兒.
沒有見十幾年.
人生走到盡頭好像有很多破碎.
很多brokenness.
突然間有一個基督徒在他身邊出現.
很有耐性的聽他說話.

$^{641}$不會跟他說教.
你知道跟猶太人傳福音是最難的.
你千萬不要跟他說教.
他一聽就聽到.
拜祖所有人都不應該說教.
偏偏我們基督徒很喜歡跟人說教.
但最重要的是那個人去關心他.
當他是一個人.
是一個全人類.
到底他經歷過什麼.
我們聽一下他的story.
就是有一個這樣的人.
花很多時間跟他喝咖啡.
然後去陪他 安慰他.
但一句說教都沒有.
就是因為這樣的緣故.
這個人 這個猶太人拉比.
就開始改觀.
對基督徒改觀.
今天教會最大的威脅是什麼.
以弗所這封信告訴我們.
是indifference.
是我們無動於衷.
漠不關心.
是否?.
I don't care.
你去做這件事.
I don't care.
我這件好事.
But I really don't care.
初愛 初心.
拋棄了.
你想想你以前學生的時候不是這樣的.
你想想你以前初信的時候不是這樣的.
你想想你曾經的faith journey裡面.
有些位置不是這樣的.
但你abandon了.
你丟了它.
Right?.
Because you don't care.

$^{681}$上帝說I care.
Right?.
因為上帝以他的愛來愛我們.
He cares about.
我們跟他之間的關係.
Right?.
以弗所教會已經沒有了.
你有沒有留意.
基督的信仰是一個關係.
Right?.
It's a relationship.
我在明尼蘇達克的退休會.
有人介紹我跟一個未信主的.
退休的朋友聊天.
我聊了兩個多小時.
兩個多小時.
我本身想去游泳.
我跟他聊了兩個多小時.
很有趣.
他信的東西很奇怪.
Anyways.
跟他聊了一會我發現.
他不是很明白.
信仰是什麼.
他在跟我debate.
哪個idea是對的.
哪個concept是對的.
他怎樣解說這些東西.
你怎樣解說這些東西.
無論你怎樣解說.
核心的裡面.
信仰是什麼.
It's a relationship.
對嗎.
我們中間做父母的.
你和你的兒女.
最核心的東西是什麼.
有辦法聊.
有溝通.
你尊重他 他尊重你.

$^{721}$你愛他 他愛你.
That is the whole point.
不是講什麼大道理.
對嗎.
It's a relationship.
你和神之間的關係都是relationship.
他說.
你將你的first love.
即是你的起初的愛心已經丟棄了.
你搞什麼.
我很喜歡Eugene Peterson.
一個paraphrase的英文翻譯.
paraphrase是什麼意思.
即是將它譯譯.
不是直譯.
不是很近元文.
不過他將意思.
去到translate出來.
這叫paraphrase.
即是等於今天的當代譯本.
用一個現代的方法去表達出來.
這句話他怎麼翻譯.
先聽聽 先講英文.
這裡說的是基督教會和你講什麼.
But you walk away from your first love.
Why.
What is going on with you anyway.
Do you have any idea how far you have fallen.
A lucifer fall.
中文怎麼翻譯.
你卻離開了你的初愛.
為什麼.
你到底發生了什麼事.
你知不知道你墜落有多深.
墮落有多深.
像路西法.
即是撒旦Lucifer的隕落一樣.
即是猶太人在兩日中間有一個這樣的理念.
撒旦好像神星一樣.
墮落.

$^{761}$自高自大然後墮落.
這個墮落的distance.
是宇宙間最大的墮落distance.
神星墮落.
不能想像那種距離.
Eugene Peterson說.
Walk away from your first love.
What the heck is going on with you.
There is something wrong with you.
What is going on.
你知不知道那個distance有多大.
這個不是一級兩級這麼簡單.
是一個最大的distance的fall.
Lucifer fall.
所以你要回想.
你從哪裡墮落.
要悔改做起初所做的事.
Otherwise.
我們在燈台上沒有什麼意思.
上帝說如果我們不重尋這個初心.
你都不可以在燈台上.
我們靈州人說你是我的見證人.
You fail.
You cannot be my witness.
You no longer my witness.
What the heck is going on.
Lucifer fall.
Satanic fall.
回想悔改.
做起初以前做的事.
以前做的事.
這封宣傳結束的時候.
他說然你還一件可取的事.
就是你恨惡離國帶黨的行為.
這些行為也是我恨惡的.
凡有耳聽的都應當聽.
聖靈向宗教會所說的話.
得勝的.
我必將上帝樂園中生命樹的果子.
賜給他.

$^{801}$我們同心低頭祈禱.
我們今天都預備心去領受聖餐.
重尋這個初心這個初愛.
上主我們多謝你.
因為你愛我們.
我們每一個在你面前.
或多或少我們都忘記了.
拋棄了我們的初心.
以前可能年輕的時候那種初心.
或曾經我們在我們崇靈的旅程裡面.
我們曾經有過的對你那種真誠.
那種第一次愛.
主我們求你寫給我們.
原諒我們.
原諒我們不要當你是我們的主.
原諒我們很多時候只是做一個消費者.
主求你幫助我們復合這個關係.
主在這個時候幫助我們回想.
回想那個時候對你仍然有那種初心.
到底是何時.
主你一清二楚.
主讓我們去回想的時候.
我們去悔改.
求你原諒我們.
給我們力量給我們因緣重造以前.
出自初心那種為你所作的事.
去見證你去分享你.
去建立他人去愛你.
主求你幫助我們.
因為我們每一個都是破碎的.
可能在忙碌裡面我們將初心拋棄.
可能在我們最終將初心拋棄.
可能在很多不同的掙扎裡面.
我們好像垃圾那樣將它丟掉.
主我們多謝你.
因為你不離不捨.
仍在燈台中間行走.
守到你那一著七星.
主求你給我們這個提醒.
讓聖靈幫助我們今天.

$^{841}$我們再一次來到他面前.
回想悔改.
作起初所做的事.
主我們預備我們的生命.
預備我們的心.
來領受這個聖餐.
這個餅這個杯.
不是一些黑板每個月做一次的事.
而是我們來到你面前.
回想悔改.
作起初所做的事.
多謝你對我們的恩憐.
多謝你對我的愛.
在十字架上你為我們寫生.
曾經打動過我們.
曾經燃起我們這種初心.
主今天你去燃起我們.
再一次幫助我們.
我們悔改轉變.
做一些你愛的事情.
以愛來愛你.
以愛來愛我們身邊的人.
建立我們的生命.
建立我們的關係.
我們今天祈禱.
奉主耶穌的名字.
\newpage



\section{}
\label{sec:U3A7_wArCTw}
\textbf{What Christ Thinks of the Church \#2  (Revelation 2-3) - Rev.  Mark Chan}
\newline
\newline
連結: \href{https://youtube.com/watch?v=U3A7-wArCTw}{\texttt{https://youtube.com/watch?v=U3A7-wArCTw}} ~~~~ 語音日期: 2024-09-17
\newline
\newline
\hyperref[sec:YG4nAX61j6k]{\small{< < < PREV SERMON < < <}}
~
\hyperref[sec:index_chronic]{\small{[返順時目]}}
~
\hyperref[sec:index_scriptual]{\small{[返順卷目]}}
~
\hyperref[sec:h0WrV1zwMlM]{\small{> > > NEXT SERMON > > >}}
\newline
\newline
$^{1}$And we turn to the Bible today, all right?.
We started this series on seven letters.
to the seven churches in the book of Revelation,.
and I want to tell you this, God writes letters, okay?.
God writes letters..
Just imagine if God writes you an email, all right,.
and tells you what he thinks of you,.
as a Christian, as a Jesus follower,.
as a church, as a group,.
what would that email tell you, all right?.
If God writes you an email,.
if God's going to write you a letter,.
what would the letter be,.
and what is the content of the letter?.
Think about it, because this series,.
we are thinking about this important topic,.
is what Christ himself, the head of the church,.
you know, the head of the church is not the pastor,.
is not the deacons,.
is not the court team members or ministry team,.
is not you, definitely,.
the head of the church, or head of our church,.
is Christ himself, right?.
Jesus Christ himself..
What does he think of our church?.
That is what matters, right?.
And what does he think of you, your role in the church, right?.
In the book of Revelation, Christ himself, Jesus himself,.
writes letters to the seven churches.
in what today is called Turkey,.
in the western Turkey..
And we went through this in the first letter to Ephesus,.
and there are seven churches there, right,.
in seven different cities..
The first one being Ephesus, right?.
And then the second one is,.
can you make it out from the map?.
Smyrna, right?.
Smyrna..
And then kind of go clockwise, actually,.

$^{41}$clockwise, right, all the way to the seventh one,.
which is the old Assyria..
And we read the first letter to the church of Ephesus,.
and now we come to the second letter..
So we're moving from Ephesus, we follow that line,.
kind of the postal delivery service route,.
moving north to another harbor,.
another port city called Smyrna, all right?.
If you go to Smyrna today,.
you find that this is a modern city,.
and it's a very thriving city, actually, right,.
in western Turkey..
By the way, I'm trying to put together.
a church group going to Turkey and Greece next year..
I don't know if you're interested..
You're interested, right?.
Yeah, you have to pay, right?.
You have to pay for it..
Anyways, we're trying to put together,.
you know what happened in the Mediterranean East,.
the Israel-Palestine conflict, just tragic, right?.
And everything just, you know, stopped right now, right,.
in terms of tourism..
But I'm trying to, you know,.
put together a New Testament tour,.
not going to Israel,.
because I don't know when I can go to Israel,.
or we can go to Israel,.
but hopefully the conflict will die down a little bit,.
and the land will have peace..
But I'm trying to put together a New Testament tour,.
and that is going to Turkey and Greece, all right?.
Follow the footstep of Paul kind of tour..
It'll be fun, two weeks..
And by the way, if you're interested, let me know..
I'm not sure, there's no guarantee I can put it together,.
but it also depends on what's going on over there.
and in the other part of the world..
But anyways, if you go to Smyrna today,.
it's Izmir in Turkey..

$^{81}$It's a thriving city..
You see the cruise line cruising off.
from the harbor of modern day Smyrna..
And if you visit the site,.
you find everything you can find in Chicago, right?.
You know, Supermall and all that stuff,.
everything is modernized..
But, you know, that's the fun thing.
about going to biblical sites..
You know, you always have the mix.
of ancient and modern, right?.
You can go to the mall and there's Starbucks there,.
you know, get a cup of cappuccino or something like that..
And you don't feel like you're going somewhere else, right?.
You're still in the, kind of in the West, right?.
Kind of in America..
But, you know, walking down the street,.
you find the ancient site.
where the archeologists has uncovered..
And they uncovered the first century Smyrna, right?.
It was an ancient city with a good harbor, just like today..
Back then it was a good harbor,.
just like Ephesus had a good harbor back then..
A big marketplace, right?.
A big marketplace, you see down there, all the columns..
Big marketplace, a forum and a theater,.
Roman stadium, temples and library,.
and all the public buildings..
You see the mountain there is called Mount Pagus,.
Mount Pagus..
And it's a city of great building..
And from the ancient texts,.
we know that people are praising the city, Smyrna,.
as a beautiful city..
And it's like a city surrounding that mountain,.
surrounding that hill, just like a crown, you know?.
All the city and all the ancient, you know,.
residents and cities and development at the foothill.
and surrounding that hill..
And on top of that hill, Aquapolis, there's a temple..

$^{121}$And that is the background of the letter.
to the Christians in Smyrna..
Now, again, this is one of the seven letters..
And there were seven letters there..
If you read them through,.
you find that only two letters have no criticism.
of the churches..
There are only two out of the seven, right?.
If you take it as a whole, you know, we fail, right?.
Only two out of these seven churches have no criticism.
from Christ the letter writer himself..
So Christ will get the seven churches..
There are only two churches that have praises only,.
commendations only, no criticism..
And one of them is Smyrna..
The other one is Philadelphia..
So this letter goes,.
"To the angel of the church in Smyrna write,.
"these are the words of him," Jesus himself,.
"who is the first and the last.
"and who died and came to life again.".
First and last, right?.
That kind of echoes the familiar phrase in the Bible,.
Alpha and Omega..
Remember that one?.
Alpha and Omega, there are 24 letters in Greek..
And the first letter being Alpha,.
the last letter, the 24th letter is Omega..
So first and last, Alpha and Omega,.
that means that kind of title for Jesus,.
meaning he has everything under control..
He knows everything, right?.
He's the first and last in your life..
He's the Alpha Omega of your life, right?.
Where did you come from?.
Where did you born?.
Where were you born?.
And what have gone through in your life journey?.
From day one of your existence up to now.
and until future someday, you know,.

$^{161}$you will end your life journey, Alpha and Omega..
That is our Lord Jesus Christ..
And he says, "I know everything about you.".
Every single letter and every single detail,.
every single day and hour you have been through.
and you will go through,.
God's saying he is the Alpha and Omega..
He's the first and the last..
And he knows everything about you.
and he has everything under his control..
And so as to the questions in Smyrna.
and that kind of a comprehensive title..
And also he said, "He died and came to life," right?.
"He has conquered death.
"and death should have no authority over you.".
Yes, death and evil have some effect on you..
It can influence you..
It can overpower you,.
but it has no ultimate authority in your life.
because God in Christ Jesus has conquered the final enemy,.
which is death..
So this dual introduction of Christ himself.
offers some hope and encouragement to all of us, right?.
Especially the small churches,.
small house churches in Smyrna..
Jesus knows everything.
and he has conquered the final enemy..
And it has to do with what they're experiencing at that time..
And I'm sure if we read this short letter,.
you'll find some echoes,.
you'll find some resonance..
You know, it's kind of like me, you know?.
It's kind of like me..
Verse nine, the letter goes on..
"I know your afflictions and your poverty,.
"and yet you are rich..
"I know the slander of those who say they are Jews,.
"but they are not, but they are a synagogue of Satan.".
Now, from this verse,.
do you think the church in Smyrna is rich or poor?.

$^{201}$Was it a mega church?.
Was it a medium-sized church?.
It's a very small church, right?.
They don't have a lot of resources..
They are poor, right?.
In profiting, right?.
Actually, in first century,.
the early Christians, most of them were poor people, right?.
They lack of any resources.
and means to thrive, right?.
They have to depend on God..
So this church is not rich..
They are, in fact, very poor..
And the church is going through some very tough time, right?.
It's a very tough time..
And John, the letter writer, says.
they are going through some tough times.
because of the Jews,.
Jewish slander, right?.
So what happened there, right?.
Now, when we read the Bible,.
we have to read the Bible in its own context, right?.
In the ancient context,.
the term, the Jews, or Jewish people,.
so often not an anti-Semitic statement, right?.
So John is saying Jews or Jewish are slandering.
against the churches, against the church in Smyrna..
It's not an anti-Semitic statement.
because the early Christians, ethnicity-wise,.
they were Jewish Christians, right?.
They were themselves Jews, right?.
How could they be anti-Semitic, right?.
But in the entire New Testament,.
when it mentions the Jews,.
and they are doing something to persecute the early church,.
the Jews represent the unconverted Jewish people,.
especially the authorities, the leaders of the synagogue..
They would turn Christians.
to the Roman authorities for persecution..
Remember, Jesus was turned in.

$^{241}$by the Jewish leaders in Jerusalem, right?.
The early Christians, as told in the book of Acts,.
were persecuted by the Jewish people, right?.
The Jewish authorities in Jerusalem..
And they collaborate with the Romans.
and try to stumble down the church,.
try to stop the church from spreading out, right?.
So the early church had that experience,.
and until the day of John,.
until the day of the book of Revelation,.
they were still under the pressure.
from the Jewish authorities of the city..
So we can imagine in the city of Smyrna,.
there were some Jewish leaders,.
and they don't believe Jesus as Christ, as the Messiah,.
and they bad-mouthed them, and they slandered the Christians,.
and sometimes they turned the leaders,.
the church leaders, to the Roman authorities.
and trying to persecute the church..
So this small church is going through.
some very, very tough experience..
But Jesus says in this letter,.
he says, "Do not be afraid of what you're about to suffer.".
Jesus is not saying, "I'm gonna scoop you out,.
"and I'm gonna save you from suffering.".
Jesus said, "Do not be afraid what you're about to suffer..
"I tell you, the devil will put some of you in prison.
"to test you, and you will suffer persecution.
"for 10 days..
"But be faithful, even to the point of death,.
"and I'll give you the crown of life.".
You know, the main job description of Satan.
is one thing, right?.
One thing..
If you check Satan's resume,.
there's only one thing, only one bullet point..
It's to make God's people suffer, right?.
That's his job, right?.
24/7..
That's what Satan has been thinking about, right?.

$^{281}$That's what Satan is thinking about you..
That's what Satan is thinking about our church..
It's to persecute them, to stop them,.
and to close them down..
That is the resume of Satan, right?.
It's a pretty easy job, right?.
24/7, the devil is trying to just destroy your faith,.
just to lead you astray,.
from believing to unbelieving,.
from unbelieving to unbelieving remains unbelief..
And that is Satan's job description, only job description..
And this letter is trying to encourage.
the small church in Smyrna to withstand the pressure,.
to withstand suffering..
And it says there's a short period of time of testing..
10 days, right?.
10 days..
You know, if you search all the pages of the Bible,.
from Old Testament to New Testament,.
there's one place talking about persecution,.
talking about testing of faith,.
and it mentioned 10 days..
There's only one passage, and that's Daniel chapter one..
Remember that story?.
Daniel was being captured,.
and he was exiled in the Babylonian court,.
and he was forced to be assimilated.
to the Babylonian culture..
And Daniel and his three friends were saying,.
"Hey, you know, give us 10 days..
"I'm not gonna partake in the royal feast..
"You know, you have to test us for 10 days, right?.
"And you will know after 10 days,.
"we will do well and even better than all the other pagans.".
And as a result, you know, the consequence was that.
Daniel, trusting God, and after 10 days,.
you know, the court sees Daniel and his three friends,.
they're actually more healthy, right?.
They are better than other people..
Now, of course, that passage is not.

$^{321}$a modern copycat dietary law, right?.
It's not for us to copy the same thing that they did..
But it shows us, you know, God is trustworthy, right?.
It shows us God is trustworthy,.
and Daniel and his friends were vindicated..
They're proven to be innocent, proven to be holy,.
and set apart by God himself..
So they finally experienced vindication..
And that's what Jesus is saying to the church in Smyrna..
Yes, there's persecution..
There will be 10 days of testing,.
but God will vindicate you..
God will protect you..
And at the end, you will win.
because you are standing on the side of God..
You're standing on God's side..
Persecution..
You know, by the way, persecution is still a reality today..
Still a reality today..
Of course, we're sitting here..
We can come in, lazy..
You know, we can be lazy..
In America, there's not so much about persecution, right?.
I mean, nobody go out of the church or come in the church.
and fearful of the neighbors killing him, right?.
Or slandering you, right?.
There may be some pressure..
There may be some temptation..
And that is by and large our experience.
in America as Christians..
However, persecution is still very much going on today.
in other parts of the world,.
especially those countries or areas.
that has no access to the gospel.
or the authority or the government.
is against the Christian faith..
Like many places in China, right?.
It's very much going on, persecution..
Many places in Muslim countries,.
persecution is a real thing..

$^{361}$And let me read you some numbers, some statistics,.
so that we know what Jesus is talking about here.
is not just 2,000 years ago..
It's very much a reality today..
So some numbers here..
There were, according to open doors.
and Gordon Cornwell resource,.
almost 400 Christians are killed, persecuted,.
for their faith every month..
If you figure out numbers, try to calculate the number,.
there were 13 Christians who died for their faith every day..
Even now on average, 13 Christians every day..
And then 214 churches and Christian properties.
are destroyed every month around the world,.
according to open doors,.
especially in places hostile to the gospel..
200 some churches are being destroyed, you know,.
intentionally because they hate Christ, right?.
And these churches suffer damages..
Now, the last number, listen up..
40 to 80 million Christians have been martyred.
in the course of 2,000 years church history..
40 to 80 million Christians have been martyred..
They die for their faith, being persecuted.
in the course of church history..
Church history has 2,000 years old, right?.
And over half of them were martyred in the 20th century..
The last century, the last century..
Half of them, that means at least 20 million Christians.
were persecuted and martyred in 20th century,.
in the last century..
So it's very much a reality, right, today..
Even though we don't see it, you know,.
kind of we don't see it in our sight, our mind,.
but it's a reality..
So what the church in Samaria is going through.
is actually what many Christians in other parts of the world.
are going through today..
And Satan will continue to persecute Christians..
Satan will not get tired, okay?.

$^{401}$You go to sleep at 10 o'clock..
You go to sleep at 11 p.m., right?.
And then you wake up 8 o'clock, 9 o'clock, 10 o'clock, right?.
You have to rest, right?.
But Satan, Satan does not need rest, right?.
Right?.
And you know, in Jesus' encounters with Satan,.
with the devil, actually the devil studied the Bible..
Remember the temptations Jesus had.
in the 40 day of wilderness?.
Satan actually quote from one of the Psalms, right?.
Psalm 91, right?.
Do you know what Psalm 91 is?.
We don't know because we don't study..
Satan's more diligent, knowing the word of God,.
but for the purpose of tempting,.
for the purpose of persecuting Christians..
So if they misinterpret, they can be led astray, right?.
So Satan will continue to persecute Christians.
with suffering, and by and large for us, it's temptation..
Right?.
There are tons of temptations in our lives, in your life..
Right?.
How long you been a Christian?.
Probably from some months and probably some many years,.
right?.
And if you check your experience,.
there are many, many temptations..
There are many, many temptations.
that could have destroyed your faith, right?.
Criticism..
Some people will criticize you, right?.
Probably unbelieving parents.
or unbelieving family members will criticize you.
for believing Jesus, right?.
And then you experience discouragements, right?.
From time to time, you feel discouraged, right?.
Now, I've been serving in church for 25 years, right?.
And I tell you, many, many experiences of discouragement..
Either people give me the discouragement.

$^{441}$or I myself perpetuate that kind of discouragement, right?.
Don't you feel that way sometimes, right?.
You feel discouraged, right?.
Your friends left the faith, right?.
Your friends criticize you or friends, you know,.
go cold shoulder on you, you know, in terms of faith..
And there are discouragement, there was criticism,.
there was negative experience that we all share..
And that is part of the real temptations..
And Satan's using those occasions,.
using those temptations to lead you astray, you know,.
to beat you down so that you will not be able.
to follow Christ, right?.
You will not be able to serve Christ, you know,.
to magnify your shortcomings, right?.
Magnify your guilt, magnify your weaknesses.
so that you feel like you don't deserve.
to be called a disciple of Christ, right?.
And that is trying to lead you astray from grace, right?.
Because we can always fall back on grace, right?.
Our faith started with grace and ended with sheer grace..
It's not about our work, but Satan can use those occasions.
to destroy your faith, to slow down your Christian growth..
I think this letter is a very relevant letter..
Because it reminds us as Christians, you know,.
we will go through some suffering, right?.
You know what is the most often used image.
to describe Christian walk?.
What's the most often used, frequently used image.
in the New Testament to describe Christian life?.
What is it?.
Let me tell you..
Battle..
It's a warfare, right?.
It's a spiritual battle..
That's one image, right?.
Farming is another image, right?.
Have you ever been a farmer?.
It's tough work, right?.
Running, you know, running..

$^{481}$Track and field, you know, it's a competition..
No, these images call for discipline,.
call for willingness to suffer, right?.
Willingness to withstand pressure..
And that is the call of Christians..
When Jesus called you to follow him,.
what he really meant was what he meant.
when he called the first disciples..
He said to Peter, Andrew, John, and those guys,.
Jesus said, "Follow me.".
Jesus actually meant just follow me,.
just walking after me, right?.
You know, walking from town to town, from city to city,.
and you have no place to stay..
There was no Airbnb for you..
There's no hotel for you, right?.
You gotta, you know, withstand the pressure, right?.
To have the discipline to follow him literally..
And Jesus is calling you today to do the same thing, right?.
It's a battle..
It's a spiritual warfare, and we need to get ready for that..
And I'll tell you how to get ready for that, all right?.
From this level..
And, but I wanna give you an illustration..
You know, not too long ago, someone introduced me.
the term strawberry generation..
Strawberry generation..
Have you ever heard of that?.
It's actually a term coined by an Indonesian sociologist..
You know, it's a professor in Indonesia..
And that professor actually wrote a book on this topic,.
and the title of the book is "Strawberry Generation.".
And it's a very popular term now,.
or used to be a very popular term in Taiwan..
Now, this term, this label, strawberry generation,.
describes people born after 1990..
Anyone here?.
1990..
I'm sure you're not, okay?.
I'm sure you're not..

$^{521}$But let me tell you what strawberry generation is..
Now, strawberry generation means,.
you know, this image means, you know,.
those people who are born after '90s,.
strawberries look nice and colorful, right?.
And you know, you see creativity,.
you see freshness, right?.
And these are the people, so-called Gen Z, right?.
Strawberry generation..
Look very nice and colorful,.
but strawberries bruise easily, right?.
You press it, and then you bruise it..
And the term actually describes a new generation.
that cannot withstand social pressure,.
not willing to work hard,.
or less resistant to challenges in life..
Perhaps, you know, sociologists will say,.
perhaps, you know, this generation.
grew up in an overly protective environment, right?.
Probably fostered by parents, right?.
You heard of helicopter parents, right?.
I was kind of one of them, right?.
When my girl was trying to learn to walk, right?.
And she stumbled, right?.
I was like, "Oh my God.".
Right?.
"Oh my God, I don't want to get hurt," right?.
And, you know, having too much too soon,.
growing up in a very comfortable environment,.
just like greenhouse, right?.
Strawberries are grown in protected environment, right?.
Protected greenhouse..
Spoiled, right?.
And the sociologists are suggesting.
we have a whole new generation who are like strawberries,.
right?.
So I wonder, are there strawberry Christians?.
Are there a generation of strawberry Christians?.
You know, we look nice, you know?.
We look kind of creative, right?.

$^{561}$We have many ideas, and we are fresh..
But, you know, are we like strawberry.
when we are being pushed or pressured,.
and we get bruised too easily, right?.
This letter reminds us what it means to be Christian..
The church in Smyrna, they set up a great example..
They were always under the pressure..
They were always under the attack of Satan, right?.
But Christ encourages them..
To withstand suffering, right?.
So your question is, how to handle suffering?.
How should we do it, right?.
How should we handle pressure?.
Especially pressure related to our faith, right?.
Let me give you some pointers..
First, you need to understand, as this letter tells us,.
understand suffering is a test..
It's a test of your character..
It's a test of your faith, right?.
Now, remember, when you go through a test,.
of course, it's painful..
It's not comfortable..
But remember, only people who are of worth.
deserve testing, right?.
Only you are equipped..
You have been growing up..
That's why God allows testing in your life..
Because God sees you as someone of worth..
Wanna test you, wanna grow you, even better, right?.
That's one pointer..
The second one is to understand what faith really is..
Question of faith, following Jesus..
What does it mean, right?.
So often we misunderstood..
We think believing Jesus is simply.
just get a ticket to heaven, right?.
Get a passport to heaven..
Get a ticket to heaven so I can wait in the gate, right?.
I can be passive, I don't have to do anything..
And I just wait, someday I'll fly to heaven.

$^{601}$or something like that..
No, that is not biblical faith..
God does not save us from suffering..
But God will save us through suffering..
You see the difference?.
God does not necessarily save us from suffering..
Of course we wanna pray that God will save us from suffering..
We don't have to suffer..
But so often God does not do that..
But God, in his grace, in his faithfulness,.
God is telling you he will save you through suffering, right?.
So that you can withstand suffering..
In his grace, you can depend on him..
Depend on his grace and his strength..
The second part, understand what faith is all about..
The third point I wanna suggest to you.
is to learn spiritual disciplines..
Learn spiritual disciplines, right?.
Reading the Bible, right?.
Reading the Bible, praying every day, you know?.
To have the discipline to go to worship, right?.
To do something, what does it mean, spiritual discipline?.
To do something that is against our.
most comfortable desire, right?.
If we seek for comfort, of course we will not read the Bible..
We will not seek God, right?.
We will seek sin..
We will seek our own self..
We will listen to our own self..
But spiritual discipline, we're against our own nature.
and build on the new nature, right?.
The more you practice spiritual disciplines,.
you know, regular prayer, regular Bible study,.
regular getting together with other people,.
other people of God, other Christians,.
to get encouragement and pray for other people,.
these kind of spiritual disciplines will build up your life.
and will help you withstand pressures..
Point number three..
Number four, learn from great examples..

$^{641}$People who have finished their course of faith,.
you know, beautifully and finish it well..
Learn from them, right?.
Last week we had a memorial service.
and we said goodbye or see you later.
to one of the brothers,.
one of the actual retired pastor in our midst..
And he passed away when he was 96 years old, 96 years old..
And he has been serving God for decades.
and he has been faithful..
You know, by telling stories of his past experience,.
he has suffered quite a bit..
He had different challenges in his life,.
serving God, following Jesus, right?.
But he finished it well, you know..
We praise God for his life, right?.
And we can learn from him..
We can learn from great examples of faith, right?.
And there are many great examples around us..
And there are many great examples in the Bible..
We can learn from them..
That's point number four..
Number five, focus on eternal reward..
Someday we will be rewarded..
And God is saying to the church in Myrna.
and to our church, to you,.
he who has an ear,.
let him hear what the Spirit says to the churches..
He who overcomes, he who endures,.
will not be hurt at all..
By the second death..
And he will have the crown of life, right?.
The crown of life..
That is the desire for Christ when he sees you..
Ask questions, ask follow of Jesus Christ..
Look forward to the eternal reward..
The crown of life and then will not be hurt by second death..
What is that?.
First death is the physical death, right?.
The separation of the physical body and our soul or spirit..

$^{681}$The second death is the final judgment..
That is unfold in the book of Revelation,.
the later chapters..
That is the second death..
But Jesus has conquered death..
And second death, final judgment,.
has no power over you, has no authority over you..
So you can withstand because you have this hope.
in the future..
And that hope is secured by Jesus Christ..
Let me tell you this, there's a saying that says,.
tough times don't last, but tough people do..
Right, heard of that?.
Tough times don't last, but tough people do..
But I would say of this passage,.
tough times don't last, but God's people do..
God's people last forever, eternally..
Because we are secure in Christ..
Yes, we can go through some tough times,.
but they don't last because we have this eternal hope..
Right?.
Let me remind you and let me encourage you today..
Let me pray and we come before God and say,.
God, we need your strength, we need your grace,.
and we know that you are a rock and we can depend upon you..
We can draw strength from you..
We can gain strength from you..
We can go on and follow Jesus, right?.
Let's pray..
Father, we just come before you..
We know that we are weak and we are easily defeated..
Whenever we focus on our old self, our own brokenness,.
our own shortcomings, our own weaknesses,.
yes, we will be broken and we'll fall into temptation.
because Satan will not get tired of tempting us.
and damaging our faith..
But we know that you are our rock..
You are the God who we can trust, the God who is our rock..
And I just pray that with all the pressures that we have.
when we follow you, with all the challenges.

$^{721}$that we have to go through, that our faith will be tested..
Our faith will be polished and our character,.
our life, our Christian life will be just like a lampstand.
to witness to your faithfulness..
And I pray that for your church, you will protect us.
and you will bring us through suffering..
And we will know that you are always trustworthy..
Thank you for being our faithful and loving God..
Thank you for loving us so much unconditionally..
And may your love and your grace continue to carry us through.
so that we give you glory in Christ and pray..
Well, please stay for a time of fellowship.
and then we'll start our Bible study, right?.
Okay, all right?.
No Satan is going to do Bible study,.
but we will do Bible study, all right?.
All right, God bless you..
And stay in the cafe and we'll call you.
when it's time for our fellowship, all right?.
God bless..
[BLANK AUDIO].
\newpage



\section{}
\label{sec:h0WrV1zwMlM}
\textbf{What Christ Thinks of the Church \#3 (Revelation 2-3) - Rev. Mark Chan}
\newline
\newline
連結: \href{https://youtube.com/watch?v=h0WrV1zwMlM}{\texttt{https://youtube.com/watch?v=h0WrV1zwMlM}} ~~~~ 語音日期: 2024-09-17
\newline
\newline
\hyperref[sec:U3A7_wArCTw]{\small{< < < PREV SERMON < < <}}
~
\hyperref[sec:index_chronic]{\small{[返順時目]}}
~
\hyperref[sec:index_scriptual]{\small{[返順卷目]}}
~
\hyperref[sec:FGEt94QjF9Q]{\small{> > > NEXT SERMON > > >}}
\newline
\newline
$^{1}$Today, we're going to go through our preaching series,.
our sermon series, and if you have your Bible,.
please turn to the last book of the Bible,.
the book of Revelation, all right?.
Book of Revelation..
We are in this series called,.
"What Christ Thinks of the Church.".
What does Jesus think of our church?.
You know, just imagine, if Jesus is gonna write us an email,.
if he's going to text us,.
if he's gonna write a letter to us,.
and to each and every one of us,.
what will the letter be saying?.
You know, imagine, you know,.
if Jesus write a letter to all of us, to our church,.
what does he think of our church?.
So that is the whole idea of this sermon series.
on the seven letters to the seven churches.
in the first century,.
and we caught it in the book of Revelation..
And as I said before, this seven letters.
were sent to seven churches.
located in today's Western Turkey, right?.
Did I mention to you I wanna organize a New Testament tour.
to Turkey and Greece?.
Did I, did I say that?.
Yeah, I did, right?.
I did..
Last week, I did, yeah..
And, you know, because of the war there,.
because of the instability there,.
we could not go this year..
But I usually go to the Middle East once a year.
to lead some church group and to visit different sites..
Hopefully, next year, we're gonna have a New Testament tour.
for the first step of Paul,.
and we're gonna visit Turkey and Greece..
And anyways, this seven churches.
in the book of Revelation, the first few chapters,.
we caught a seven letters written by Jesus,.

$^{41}$the head of the church, to all the churches,.
to seven churches..
And we went through the letter to,.
Ephesus, and then Smyrna..
And today, we are at the third city..
You follow that line,.
possibly a postal delivery route of the day..
And the seven churches are written to these seven cities.
where there were whole churches in there..
And today, we are in the third city, Pergamon, Pergamon..
Now, back in the first century,.
Pergamon was a thriving city..
And listen to this..
According to different estimation,.
Pergamon had a population of 150,000 people..
And that considered a big city in the first century..
It's a huge city, huge population..
150,000, it's a thriving city..
And today, when we go there,.
you see a Roman theater on the left,.
and the Roman theater can seat more than 10,000 audience..
And then on the right, you can see the aerial picture..
And then down there, you see the modern city of Pergamon..
And this letter, Jesus is writing to the church in Pergamon..
Let's take a look..
212, 212 of Revelation..
"To the angel of the church in Pergamon write,.
"these are the words of him who has a sharp double-edged sword.".
Now, as I said, in all three, all seven letters,.
Jesus introduced himself in a special way.
that is relevant to the city, to the Christians,.
to the church in that particular city..
And writing to the church,.
writing to the Christians in Pergamon,.
Jesus introduced himself as having.
the sharp double-edged sword..
Now, this kind of sword,.
archeologists have dug up so many of them in Roman sites.
in the Mediterranean world in the first century..
It is a well-known weapon of the Roman Empire..

$^{81}$It is symbol of political might, military might..
And politics in today and back in the first century.
was a highly debated topic, right?.
Because politic is the greatest power on earth, right?.
The country, whoever is the boss in the Roman Empire..
And it was a highly debated topic,.
and everyone is so concerned about it..
Just like today, we are very concerned.
about the upcoming presidential election..
It's a very important subject,.
and politics is usually a very relevant subject.
in the first century..
And the Roman emperor and governors of the province,.
including the governor of the city of Pergamon,.
often carry this kind of double-edged sword, you know?.
This kind of double-edged sword..
It's about probably more than 18 inch long..
It's a double-edged sword to symbolize their power..
They have the power to execute the Roman authorities..
They have power to punish criminals.
and anyone against the order of the empire..
You remember in the Gospel of John.
when Jesus was standing on trial before Pilate,.
the Roman governor in Palestine, in Israel..
And you remember in the Gospel of John.
recorded the conversation between Jesus and Pilate..
This is what they talk about..
John 19, 10..
Now, Pilate, the governor,.
representing the power of the Roman empire,.
he said to Jesus, "Don't you realize I have power.
"either to free you or to crucify you?".
And then Jesus said, and that's amazing what he said..
He said to Pilate, "You will have no power over me.
"if it were not given to you from above.".
Right?.
It's a subversive challenge to the Roman authority..
You have power to crucify me and to penalize me.
because there is a power above you..
There's a bigger power..

$^{121}$There's a greatest power above you who is God himself..
So Jesus is essentially saying, "Sorry, Pilate..
"I know you're gonna have the power.
"to crucify me or to free me,.
"but I and my God have the power above you.".
Right?.
And essentially, that's what Jesus is saying.
in this letter to Pogomon..
Jesus said, "I am the only one who has the double-edged sword..
"I have the political power, the power above.".
In chapter one, we see that Jesus holds the seven stars..
Right?.
Just like the Roman emperor would hold seven stars.
in his hand..
Right?.
In chapter one, we also see that Jesus,.
out of his mouth came a sharp, double-edged sword..
And what does it mean?.
It means Jesus has a weapon,.
but this is not a physical weapon..
This is not guns and tanks..
It is the word of truth, the word of the gospel,.
and the gospel can conquer the Roman power..
The gospel can conquer death and Satan and evil..
That's essentially what Jesus is saying..
Right?.
Because Jesus rose from the dead..
He conquered death, the greatest power of all..
Right?.
In fact, the New Testament reminds us.
death is the final enemy for each and every one of us..
You will die..
I will die..
Death is the final enemy,.
but Jesus has conquered this final enemy..
And Jesus is saying he has all the power..
He has the power from above..
So the Bible is proclaiming that,.
yeah, of course we need to concern about the earthly politics.
but we are not isolating ourselves, of course..

$^{161}$Right?.
We need to concern about what is happening in the world,.
but we must be occupied most of all.
in our conversation, in our thoughts, in our actions,.
by the real power, the highest power, the real king..
Jesus himself, right?.
He has the double-edged sword,.
not the Roman guys, right?.
Not the United States president..
Jesus has the real double-edged sword..
And we must be concerned about that power,.
concerned about the gospel,.
concerned about preaching the word of truth,.
concerned about witnessing the name of Jesus, right?.
That is what we must occupy ourselves most, right?.
Let me remind you as the election coming up, right?.
What is occupying our time?.
What is occupying our thoughts most of the time?.
It should be, the Bible says, Jesus himself, right?.
The church, his people, his authority, his kingdom..
Just like the Lord's prayer, "Thy kingdom come.".
On earth as it is in heaven..
We're gonna bring God's kingdom on earth, to the earth,.
to our circle of inference, to help people see Jesus..
That should occupy most of our time, most of our energy..
So Jesus introduced himself, he has the highest power..
He has the sharp double-edged sword..
And he continues to talk to the church in Pergamon..
He said, "I know," remember this formula.
in every single letter?.
He said, "I know your situation..
"I know everything about you.".
He said, "I know where you live, where Satan has his thorn..
"Yet you remain true to my name..
"You did not renounce your faith in me,.
"even in the days of Antipas, my faithful witness,.
"who was put to death in your city, where Satan lives.".
How many times Satan is repeated here?.
Two times, right?.
Yeah, two times..

$^{201}$So looking at Pergamon, this thriving city,.
of course, he says, "I know where you live," right?.
Your physical address, right?.
You can punch it in, in your Google Earth,.
Google Earth or Google Map, right?.
You can go there, where you live, of course, is Pergamon..
It's the ancient city..
That's the physical location, just like Chicago, right?.
I know where you live, Jesus says..
But Jesus is talking about not your physical location,.
it's just spiritual location..
Spiritually, Pergamon was the throne of Satan,.
the throne of the evil..
Now we have archeology uncover some of the truth.
about the city for us..
Pergamon, in the last hundred some years,.
the Germans were there to dig up some of the things.
from the city, Pergamon..
And Pergamon, we know, had a large altar,.
or temple, built for Zeus..
And this huge altar, you can see it now.
in Pergamon Museum in Berlin, Germany..
And because when the Germans went to today's Turkey,.
they dig up some of the stones,.
some of the stones of the altar and the temple,.
and they move it piece by piece back to Berlin.
and put it in the museum..
And this is a huge pagan altar..
This huge pagan temple, hundred some feet long.
and hundred some feet wide and 40 feet plus deep..
It's a huge altar..
And what we found in Pergamon was this huge altar..
And you can see it in Germany..
And the story goes, as the German move the entire altar.
stone by stone to Berlin, and they reassemble it together.
in their own museum..
And then it was 1930s, and the Pergamon Museum in Berlin.
was open to the public..
And a few years later, the German Nazi party,.
they had some rallies, some political gatherings,.

$^{241}$and they are trying to design and build.
some of the huge plazas for the political rally..
And in one of those, the Nazi party commissioned.
an architect to design a special plaza, Sibel and Phil..
And in that plaza, the architect, the chief architect,.
got the inspiration from Pergamon altar..
From this altar, he built that low right,.
Sibel and Phil, in 1930s..
And they turned to Pergamon altar for inspiration..
And of course, as you know, in the following years,.
Hitler descended down the steps..
Now this altar is what the Bible is talking about..
This altar at Pergamon is believed to be.
the throne of Satan, the seat of Satan..
And it's very interesting to think about this verse,.
verse two, 13 and 14..
Jesus is saying to the church in Pergamon,.
"Don't just look at the surface..
"You need to look behind your city..
"And you look behind the city, you know the power of Satan.
"is working behind, is working there..
"There's a spiritual darkness,.
"dark forces behind what you see.".
You know, from your sight, you can see a gigantic building..
You can see a big, huge altar and temple..
But behind, Jesus said,.
is the evil, destructive force of Satan, right?.
And we need to develop that kind of spiritual insight,.
that spiritual perspective..
When we see something from the outside,.
we have to look behind the scene..
Just like what Jesus did here..
Chicago is a beautiful city, isn't it?.
Right?.
It's a beautiful city..
Especially weather is good, all right?.
It's very, very beautiful..
It's just a magnificent building..
Magnificent buildings and magnificent city..
But you need to look behind..

$^{281}$Chicago is known to be one of the toughest places.
for churches..
Do you know that?.
You know, not too far away from us, in our network,.
one of the Baptist churches was trying to fight.
for zoning permit to build a church..
They tried to fight legal battles,.
not to be kicked out from their rental place..
You know, the government is not friendly.
or even hostile to the gospel..
Yes, this is a beautiful city,.
but you look behind the scene,.
you know there's something else working, right?.
The power of Satan..
And the city needs the gospel..
The city needs you to go in there,.
to bear the name of Jesus, to witness to other people.
about the wonderful name of our Lord, right?.
Same thing for Pergamon..
Now, this letter says there was one martyr.
called Antipas, right?.
Antipas..
And he's being described as my faithful witness..
Now, we don't know much about Antipas here..
Perhaps he was a leader in the church in Pergamon,.
and he was a martyr not too long ago,.
before this letter was written..
Now, at the time of revelation, it was the first century,.
the late first century, which is AD 90s..
And we know from other sources that some Christians,.
and in fact, quite a few Christian leaders.
have been martyred in the first century..
Can you name some of them?.
Stephen is one of them, right?.
The bold ax, chapter seven..
He was stoned to death..
James, who wrote the epistle of James..
Peter, who was martyred in the 60s..
Paul as well, he was martyred in Rome in the 60s..
And now when revelation was written,.

$^{321}$it was late first century, AD 90s..
Antipas was one of them..
And that implies he was a faithful witness..
That implies that he was interrogated by Roman authorities,.
and he was ordered to be executed.
when he refused to curse Christ, curse Jesus,.
and worship the Roman emperor..
When Antipas gave his testimony in the court,.
and he was being asked to reject Jesus,.
and to renounce his faith, and to curse Christ publicly..
Antipas, that's for sure, he refused to do that,.
because he only worshiped Jesus himself..
Now, in the midst of intense persecution,.
the small church in Pokémon needs a lot of encouragement..
And Jesus said, "I know you guys..
"I know what is going on with you..
"I know you're suffering for the faith.".
'Cause that I know, that two words in English,.
that I know offers encouragement and comfort.
to those who are suffering for the faith..
Jesus said, "I'm gonna protect you..
"I know your situation.".
Now, Jesus never promised us to follow him,.
there will be no suffering..
But he promised us to save us through suffering, right?.
So he said, "I know your suffering..
"I know what you're going through, right?.
"I'm not sure what you're going through right now, right?.
"Possibly in the past few weeks, or past few months,.
"you have been going through some suffering..
"But remember this.".
Jesus said, "I know, I understand, I'm with you..
"I know your suffering.".
He spoke to, he spoke the same word to.
the Christians in Pergamon, at the throne of Satan, right?.
They were under attack..
And Jesus today is speaking to you the same thing..
I know what you are going through..
Let's continue, verse 14..
"Nevertheless," Jesus said, "I have a few things.

$^{361}$"against you guys, the church in Pergamon..
"You have people who hold to the teaching of Balaam,.
"who taught Balak to entice the Israelites to sin.
"by eating food, sacrifice to idols,.
"and by committing sexual immorality.".
Verse 15, "Likewise, you also have those.
"who hold to the teaching of Nicolaitans.".
Now, what is going on here?.
It's using some Old Testament example.
to describe what is going on in,.
among some of the Christians in Pergamon church..
He said, "First of all, some of the Christians there.
"held to the teaching of Balaam.".
Balaam is an Old Testament guy,.
and he's a very famous guy..
And Balaam, as you remember from the book of Numbers,.
when the Israelites was trying to pass through.
the land of Moab, and they come in closer.
to the promised land, land of Canaan..
And as they pass through the land of,.
outskirt of the land of Moab,.
the king of Moab called Balak, right, Balak..
And he hire a prophet, a false prophet named Balaam.
to curse Israel, to curse Israel,.
to not to blots them, but to curse them.
to bring calamities over the Israelites.
as they pass through the land of Moab..
But God protected the Israelites, right?.
And he turned curses into blessings..
Now, what is going on in the Christians,.
some of the Christians in Pergamon was that.
they hold to the cultic teaching of Balaam..
Now, I think what it means is that.
some of the Christians there.
were preaching for financial gain..
Now, Balaam was hired, remember?.
He was hired for money, and he tried to curse Israel..
And that is the figure of Balaam in the Old Testament.
and in the Jewish tradition, in the biblical tradition..
Balaam is a false prophet, and he does God's ministry,.

$^{401}$or speaking the word of God, for money, right?.
To gain money, for material gain..
And I'm sure this phrase here, or this description here,.
to describing some of the Christians.
who was trying to distort the gospel.
for their personal gain, just like what Balaam did..
And second thing he has against the church in Pergamon.
was that they were eating food sacrificed to idols.
and commit spiritual, sexual immorality..
I think it's likely at that time.
that Christians joined some of the local pagan associations,.
local pagan religious temple associations.
for financial benefits, right?.
Because if you join them, right, join the club,.
then you will have a better business.
and better job opportunities for financial gain, right?.
But in joining some of the pagan anti-Christian associations,.
they had to renounce their faith,.
or they have to renounce their faith for a moment, right?.
And they come back to church and they praise Jesus.
and say Christ is the Lord..
But Monday to Saturday,.
they're worshipping the false gods, right?.
I think that's what happening here..
It is turning Christian faith into making money..
Turning Christian faith for our own benefit, right?.
Sexual immorality, right?.
For our own sexual satisfaction, right?.
If you think about it,.
this letter was not just written.
to Christians in the ancient world..
It's also written to us, right?.
Turning our faith into making profit for ourselves, right?.
To satisfy ourselves, right?.
Don't we see that a lot?.
We believe in Christ, but ultimately,.
you know, at the bottom of our heart, we want benefit..
We want material gains, right?.
We have faith, yeah,.
but we want that faith to be twisted.

$^{441}$so that we have freedom to commit sins, right?.
And that is a crisis of faith, right?.
That's what American culture does to all of us, right?.
It's consumerism, right?.
We go to church just like we go to Walmart, right?.
We go to church just like we go to Whole Foods..
We like to pick and choose,.
and then ultimately, it's for our own benefit..
If we do that, if we have that kind of spirit,.
if we have that kind of mentality distorted,.
we're no longer following Jesus, right?.
We are simply, in the name of faith,.
we want to gain some financial benefit, right?.
Some satisfactions for ourselves..
And that's what Jesus said,.
"I have a few things against you," right?.
Even some of you who hold to the teaching of Balaam.
for financial gain in the name of faith, right?.
"I have this against you.".
He continues to say, "Repent, therefore.".
Repent..
You gotta turn around, you know?.
You see what faith is all about, what Jesus is all about..
Otherwise, Jesus said, "I will soon come to you.
"and fight against them with the sword of my mouth..
"He who has an ear, let him hear.
"what the Spirit says to the churches..
"To him who overcomes,.
"I will give some of the hidden manna..
"I will also give him a white stone.
"with a new name written on it,.
"known only to him who receives it.".
Now, this is a stern warning,.
stern warning of twisting the gospel.
into self-seeking religion..
If you understand faith in this way,.
if faith is simply self-seeking,.
looking for your own benefit,.
that is not biblical faith..
That is Balaam..

$^{481}$That is anti-faith, right?.
That is not biblical faith..
And the judgment of God is severe..
Jesus said, "I am coming to you," right?.
So you have to repent right now..
But there's also reward..
If you repent, if you follow Jesus with your whole heart,.
right, there's also blessing..
The hidden manna, what is that?.
Manna..
We caused you the story.
of the Israelites going through the wilderness..
In the wilderness, they had no food, right?.
They had no water, but God gave them manna, right?.
God gave them manna..
It symbolizes God's faithful provision..
God is gonna provide for you, right?.
If you repent, if you follow Jesus wholeheartedly,.
you know that God is going to provide for you, right?.
Now think about it..
Don't you think at times.
we're not committing ourselves to Jesus.
because we just don't believe that God can provide?.
Sometimes in the bottom of our hearts,.
we just don't think God is so good, right?.
We just don't think God is so faithful, right?.
That's why we have to gain something, right?.
Just like the Israelites in the wilderness did, right?.
They just don't think that God.
will provide double on Friday..
So before Shabbat, some of those people,.
they get the double of the portion they need, right?.
Because in the bottom of their hearts,.
they just don't believe that God is so good,.
God is so faithful..
Here Jesus affirms that God is the faithful provider,.
and you have to step out of your comfort zone..
You have to take a step of faith.
and believe that God will provide..
He has prepared for you manna,.

$^{521}$spiritual food to satisfy you, right?.
Which is better than other food..
And he has also promised a name on white stone, right?.
We're talking about archeology because that's my area..
If you go to the ancient site,.
you always find some big stone, a slab of stone,.
and with the name of Roman emperors.
and the famous people inscribed on them, right?.
Yeah, just like Hollywood, right?.
You walk on the street,.
you have this important names, right?.
To step on, you can step on them..
And they're permanently inscribed on stone.
and indestructible..
And what Jesus promised us is that a name on white stone,.
a permanent remembrance of your name,.
a permanent memorial of your name, right?.
And that's a big deal in the ancient Near East..
That's a big deal in the ancient society..
My name wants to continue, right?.
That's why I have my son, my grandson,.
because I want to perpetuate my name..
And Jesus promised you, if you're faithful,.
your name will be inscribed on white stone..
But here, you must remain faithful..
And you must remember, God is good..
God is faithful and we can depend on him, right?.
I want to spend some moment to pray..
The thing about who God is,.
who God is to you, is he such a faithful God?.
Is he just a good God that you can depend on?.
You can trust him wholeheartedly..
Let's repent, let us come before him,.
knowing that he is indeed faithful and loving..
Let's all pray..
(soft piano music).
Our loving Father, thank you for your word..
Thank you for your reminder,.
your evaluation and your commendation.
and even your teaching..

$^{561}$To correct us..
So we come before you,.
knowing that you are such a great God of grace..
Thank you for your redemption..
Thank you for your gracious provision..
And the fact that your son, Jesus Christ,.
came to die for our sin.
and rose from the dead for all of us,.
to restore our relationship to you,.
this fact reminds us,.
you are such a good and faithful, gracious God..
Forgive us for our little faith..
Forgive us for sometimes in the name of you,.
in the name of faith,.
we are self-seeking..
So we come before you as you admonish us.
to repent, to change our way, to change our attitude,.
to change our mentality, change our heart, please..
May your spirit change us,.
help us repent and help us to come before you..
Just know that, believe that,.
and trust that you are such a faithful and good God..
Thank you for being our God..
So I pray for all the brothers and sisters here,.
pray for myself,.
that you'll grant us this heart of purity..
So wholeheartedly we will follow you,.
wholeheartedly we will obey you.
and help us to be renewed in your grace.
and our lives can be witnessed in this city,.
witnessed to how great you are,.
witnessed to the name of Jesus,.
to help people to be introduced into your presence,.
to help people to see the gospel,.
to help people to know the name of Jesus..
So send us out and be with us.
for you know everything that we are going through.
and you are dependable.
and we can trust you wholeheartedly..
Father, in Christ's name we pray..

$^{601}$Well, thank you again for worshiping with us.
and I wanna send you out with that mission.
to bear witness for the name of Christ.
and may God's peace be with all of you.
and I wanna see you next week..
You can stay around for some fellowship time,.
enjoy our refreshment.
and I'm sure there's some mooncakes out there, right?.
Yeah, it's out there, all right..
Thank you, Pam..
All right, God bless..
\newpage



\section{歌羅西書 1:28-29}
\label{sec:FGEt94QjF9Q}
\textbf{在基督裡的勉導同行 (歌羅西書1\_28-29) - 蘇妙嫻傳道}
\newline
\newline
連結: \href{https://youtube.com/watch?v=FGEt94QjF9Q}{\texttt{https://youtube.com/watch?v=FGEt94QjF9Q}} ~~~~ 語音日期: 2024-09-17
\newline
\newline
\hyperref[sec:h0WrV1zwMlM]{\small{< < < PREV SERMON < < <}}
~
\hyperref[sec:index_chronic]{\small{[返順時目]}}
~
\hyperref[sec:index_scriptual]{\small{[返順卷目]}}
~
\hyperref[sec:S2QbUdfuEWU]{\small{> > > NEXT SERMON > > >}}
\newline
\newline
歌羅西書 1:28-29
\newline
\begin{longtable}{cl}
\hline
\hline
章節 & 經文 (和合本修訂版)\\
\hline
1:28 & \begin{tabularx}{0.7\textwidth}{X} 我們傳揚他，是用諸般的智慧，勸戒各人，教導各人，要把各人在基督裡完完全全地獻上。 \end{tabularx} \\ \\ \relax
1:29 & \begin{tabularx}{0.7\textwidth}{X} 我也為此勞苦，照著他在我裡面運用的大能盡心竭力。 \end{tabularx} \\ \\
[1ex]
\hline
\hline
\end{longtable}
$^{1}$各位,我不是Hanna.
但是我認識Hanna很久了.
認識Hanna很多年了.
我想來了芝加哥沒多久.
我二十多年前來芝加哥的.
都已經認識Hanna.
在教務局同工會等等.
Hanna一直很多年在芝加哥.
和大芝加哥區.
做一些大學生和職稱,社稱的工作.
特別是國語的士工.
我們久不久還會見面.
我會請她吃飯.
好像你沒有請她吃飯.
我每次找她都是請她吃飯的.
有時候我們教會都會借出一些地方.
給她做一些外展的工作.
我們知道在教會城市裡面.
做一些宣教工作,事奉.
是需要不只是我們教會的.
是需要其他的夥伴.
其他的夥伴,其他的同工.
所以Hanna是一個.
我們希望能夠和她更多有夥伴.
你在港島裡面可以認識到她.
在做一些的士工和一些的挑戰.
我們在禱告裡面去紀念.
還有在行動上我們支持她.
好的 我們歡迎Hanna.
我是Hanna.
剛開始的時候看到一個小Hanna.
很可愛.
其實認識你們教會有好幾位同工.
有些人例如Lysine,Phoebe.
我們都是好多年的朋友.
也是為你們教會特別獻上感恩.
城會社區教會.
雖然在國語群體當中不是很有名.
但是已經被偶然聽見.
很多人跟我打聽.

$^{41}$這是一個怎樣的教會.
為你們教會感恩.
你們已經接觸了很多不同背景.
不同階層,不同年齡層的朋友們.
也是很感恩.
也有機會和Mark Sir再一次聯絡.
我們是希望能夠和城會社區教會.
有更多的合作和配搭.
我個人是在香港長大的.
但是我大學是在台灣讀的.
工作也是在台灣.
後來來到美國.
在美國這幾年當中.
其實很主要的工作.
就是神的呼召.
來進入服侍年輕人.
我們在今天的訊息裡.
會更多和大家分享.
在年輕人當中的免導同行.
是在基督裡的免導同行.
同時也是回到聖經的真理裡.
讓我們看回在基督裡的免導同行.
有什麼真理也許我們過去有忽略了.
現在我們回到神的真理當中.
然後再看看我們如何和年輕人.
其實應該是跨代.
年輕人和年輕人同行.
年輕人也要和長輩同行.
我們和同族的同行.
也需要跨向不同的族群.
今天我們的訊息是在.
《歌羅西書》一章28到29節.
在進入經文之前.
我們來看看一個失落的環球.
這個失落的環球就是我們過去有很多的聚會.
講座,培訓.
特別是在疫情當中.
很多這些聚會,講座,培訓.
是在線上進行的.
但是幾年的疫情裡.

$^{81}$我們發覺到很多時候.
我們有很多的培訓,很多的講座.
其實很多時候我們傳達的是知識.
從聚會培訓.
到培育靈命,成全使命當中.
有一個鴻溝.
但是我們在這個鴻溝裡.
是不是可以有一個可以連接起來的.
讓這個培訓,講座,聚會.
和培育靈命,成全使命中間.
是可以連起來的.
我們要找回這個所謂失落的環球.
今天我會透過《歌羅西書》一章28到29節.
我們要找回這個失落的環球.
在基督裡免渡同行.
這個失落的環球.
我們如何回到真理裡面.
然後反思我們的事奉.
反思我們如何去培育和成全使命.
在我們所講的今天的經文裡面.
我們主要是在《歌羅西書》.
《歌羅西書》的一章28到29節.
好.
現在我們先看《和合本》.
《和合本》這麼說.
一章28節.
我們傳揚他使用諸般的智慧.
勸戒各人,教導各人.
要把各人在基督裡完完全全地引到神面前.
我又為此勞苦.
照著他在我裡面運用的大能.
盡心竭力.
這個是保羅.
他在寫《歌羅西書》.
告訴我們他在基督內為主勞苦.
而且不單是他自己.
我們還要去傳揚他.
我們看《歌羅西書》其實是一卷很寶貴的書信.
《哥羅西》不是保羅所建立的教會.
不過《哥羅西》.

$^{121}$保羅聽說有一個教會.
他們發現有異端進入.
很多人是跟從了當時的異端.
諾斯底主義.
是將很多的奧秘.
或者是透過神秘主義進入得著的榮耀.
而保羅很特別地去寫這個書信.
《哥羅西》不是他所建立的教會.
但是他懷著一個牧者的心去勸戒.
他這封書信所用的語氣其實是很溫和的.
跟寫《哥羅西》那些書是很不同的.
可以看到他裡面有一個牧者的心.
一個勸戒的語氣.
而今天我們即將看《歌羅西書》.
他如何將人引導回神的面前.
最主要是在基督內來看我們要傳揚神.
是要怎樣去傳揚他.
《歌羅西書》的一章28到29節.
其實它裡面很捉住我們的眼球.
他一直提到在基督內.
其實在基督內.
這是保羅在很多書信當中經常會提到的一個觀念.
在基督內.
保羅在新約當中留下了13封書信.
他自己在他所寫的書信當中.
很多次用到「在基督裡」這個詞.
「在基督裡」這個詞如果我們在新約裡面看的話.
至少出現了一百七十多次.
為什麼保羅這麼重視「在基督裡」.
是因為我們整個基督信仰是在基督裡而成就的.
我要看回保羅的書信當中.
他是怎樣幫助我們去了解「在基督裡」.
我們在保羅的其他書信.
包括《歌羅西書》.
他提到在基督裡面的三個主要的層面.
第一個層面就是在基督裡面我們領受神一切的美善.
這個包括了我們在基督裡面的福氣.
如果我們看伊弗所書的話.
他有很多的介紹.
其實包括在《歌羅西書》的二章第十節.

$^{161}$也會提到在基督裡面的一切的豐盛.
這個是神給我們的.
在基督裡我們領受了神給我們一切的美善.
如果沒有耶穌基督的救恩.
沒有祂的救贖的話.
我們是沒有辦法去享受這一切從基督而來給我們的美善.
而另一個層面.
在基督裡的意思是我們是信徒的身份.
不單止我們個人是基督徒.
是在基督內的基督徒.
我們也是有其他的基督徒和我們一起.
所以在基督裡也包括了我們和神個人的身份.
我們作為信徒也和其他弟兄姊妹.
同在基督裡.
同在基督內的那種身份.
是彼此相連的.
再來看看第三個基督裡.
在《保羅書信》當中經常會提到的.
就是我們在基督裡面和主的聯合.
因著耶穌基督的救贖的關係.
我們罪人的我可以去連接聖潔的神.
而神的靈可以駐載在我裡面.
來督責,勸戒,引導.
這個連接是因著耶穌基督的救贖,救恩的關係.
所以我們可以不單止與基督連接.
而基督可以在我們裡面.
這是很特別的一個真理.
而《保羅》藉著它很多的書信來帶出這個真理.
而當我們看《歌羅西書》的時候.
它所提到的.
我們要傳揚祂.
是靠在基督裡面.
也是要引導人進入到基督裡面.
耶穌基督給我們的是一個救贖.
這個救贖不單止是給我們個人.
或者給我們這個族群.
它是給萬民的.
當我們看《歌羅西書》的時候.
看一章28節.
它開始的時候.

$^{201}$《保羅》說「我們傳揚祂」.
這個「祂」是基督.
而當我們在28節之前.
它所提到的一些經文.
如果你有聖經的話.
可以打開看.
28節之前它所說到的是.
我們在基督裡面的豐盛.
還有就是我們在基督裡面.
所得著的這個榮耀.
不單止是給我們.
也是給其他的外族人或者外邦人.
所以當它一章28節所說到的.
「我們傳揚祂」.
傳揚祂的對象.
不單止是給我們同一個族群.
我們其實本身也是外邦人.
在以色列人來說.
我們是外邦人,外族人.
但是我們沒有了神的救贖.
當我們得著這個救恩之後.
我們是要跟不同的族群.
不同的年齡層來傳揚.
所以它這裡說到.
「我們傳揚祂」.
這個「傳揚」其實不單止是.
對我們同族群的人.
也是包括萬民.
包括各年齡層的人.
「我們傳揚祂」.
「傳揚」這個字眼.
其實是讓我們傳人傳心.
我們是很用力地去宣講.
有關基督的一切.
這個有關基督的一切.
當然包括祂所給我們.
救贖的恩典.
我們要怎樣傳揚呢?.
是藉著各樣的智慧.
珠般的智慧.

$^{241}$如果我們看新譯本的話.
它的翻譯跟和合本前面的一樣.
但是後面稍微有些不同.
一樣的部分是包括了傳揚祂.
用珠般的智慧.
各樣的智慧.
我們要傳揚神.
不單止是我在這裡說道.
我叫祂傳揚.
包括我怎樣的行事為人.
我怎樣去顯出在我裡面的基督.
這個都是一個傳揚的部分.
很多時候我們想著去傳福音.
就是口裡說.
但是很多時候我們忽略了.
我們怎樣去用我們一生.
用珠般的智慧.
用我們全人來去傳揚基督.
以前我未信主的時候.
我是個脾氣很大的人.
信了主之後.
都是逐漸逐漸地去改變.
我爸爸是在香港的.
他去年回了天家.
但是他未信主之前.
他經常跟我說.
特別是我剛剛信主.
但是生命還沒有很多改變的時候.
我一直跟他傳福音.
他說.
你信神這樣我才不要信.
很羞愧.
我覺得我傳揚他.
不單止是我跟我爸爸去說.
什麼叫做信耶穌.
在我的生命裡面.
我需要去反映出.
耶穌基督在我裡面.
那個改變生命的大能.
所以當我們說傳揚他的時候.

$^{281}$是包括我們全人的身心.
全人的一個呈現在神的面前.
呈現在人的面前.
珠般的智慧包括很多樣東西.
當我們去跟年輕人同行.
或者跟長輩同行的時候.
我們發覺到我們需要不同的智慧.
今年是洛桑2024年9月會舉行.
在這個會前大會當中.
他們有一個報告.
這個報告包括四大重點.
其中一個就是全球人口老化.
全球人口老化不是新聞.
但是這次洛桑會議裡面所提出來的重點.
讓我們看到他們提出一個.
過去比較少提到的.
當過去我們提到全球人口老化的時候.
我們多吸引一些年輕人來.
沒錯,過去這麼多年其實我們也有一起努力.
但是當我們提到全球人口老化的時候.
我們發覺到長輩人口.
中年或以上會越來越多.
而這一班將會是未來十年,二十年.
一群很龐大的未得子民.
我們如何去預備年輕人能夠去跨向長輩.
跨向我們的長輩.
來跟他們傳福音.
見證我們所信的基督是怎樣.
很盼望弟兄姊妹不要像當年我一樣.
我爸爸媽媽看到我這樣的情況.
看到我的生命.
看到我就不想相信.
盼望我們的生命能夠顯出.
那個救恩的真正的轉化生命的動力出來.
我們如何去用珠本的智慧.
很多不同的管道.
包括我們機構.
我們也用藝術創作.
有媒體事工.
也有心理輔導.

$^{321}$也有包括我在芝加哥所做的.
我們叫做生命培育.
大使命生命培育.
我們如何做呢.
保羅在這裡說.
我們要用珠本的智慧做什麼.
要勸誡各人.
勸誡在這裡.
其實他用的詞是很溫和的.
有些翻譯成責備.
但其實保羅在這裡所說的.
是用一個很溫和的語氣來勸誡.
中文的勸誡翻譯得很好.
他不是一個很強烈的責備.
而是懷著一個牧者的心.
來勸勉,誡,引導.
叫他不要這樣做.
當時的哥羅西教會.
有很多人慢慢傾向了那種異端裡.
而保羅在這裡說.
他要透過這個勸誡.
勸誡是提醒人.
你現在相信的東西是錯的.
而後面他所說的教導.
就是告訴人什麼是對的.
所以當我們與人同行.
要傳揚基督的時候.
我們用各種的智慧告訴別人.
你現在所信的是偏差的.
你應該所信的是什麼樣的真理.
我們所信的基督是如何的基督.
這個就是保羅在這裡說的.
用珠本的智慧去勸誡.
告訴別人什麼是不對的.
同時也要告訴別人什麼是對的.
當我們與人同行的時候.
我們不僅是在這裡宣講.
耶穌基督是什麼.
而且我們要很耐心的去引導.
告訴別人你現在所走的.

$^{361}$你現在的信念是有問題的.
那什麼是對的呢.
當我們告訴別人什麼是錯.
什麼是對的時候.
我自己又是否知道.
基督的真理是什麼呢.
在基督裡面是什麼意思呢.
這是對我自己是很好的提醒.
我自己在信主的過程當中.
很忙的造就的是.
初信的時候的基督真理的紮根.
以致到以後在新學院的裝備.
我讀新學院做了很久.
不是很多的學位.
而是慢慢的去吸收.
慢慢的去讀.
在這個過程當中有很多的反思.
而當我與別人分享的時候.
我那時候問他們.
我們說信主.
但到底信的是什麼呢.
是不是所信的神.
真是你在日常生活當中.
你知道祂是怎樣的神.
而與祂同行呢.
當我們在一章28節裡面看到.
保羅說要將各人在基督裡.
完完全全的去引導神的面前.
或者是新譯本跟我們說的.
為了要使各人在基督裡得道完全.
大家覺得兩個譯本有什麼不同呢.
和合本好像只是說.
將人帶到基督面前.
帶到神面前.
但是我們看回新譯本的話.
是更深一層的看到.
是在基督裡面得道完全.
同樣是基督裡.
我們兩個翻譯本去參照的時候.
其實給我們一個更全面的了解.

$^{401}$我們是在基督裡帶人得以完全.
在基督裡得以完全呈獻給神.
如果我們看英文本的聖經的話.
它很特別的.
也提出.
We may present every person complete in Christ.
這個Complete in Christ.
Complete.
其實它很重要一點.
提到不是我們已經很完美了.
完完全全不是代表完美.
而是在基督裡面生命的成熟.
在基督裡面生命的成熟.
我們去呈獻給神.
將人帶到神的面前.
但是一個大前提.
如果我們看一章二十二節的話.
是基督先讓寶血潔淨我們.
使我們始為始.
我們成為聖潔.
我們才在這個位分上被稱為義.
被稱為神的兒女.
是可以來到神的面前.
而保羅他在一章二十八節裡面.
同樣是二十二節和二十八節.
都有一個呈獻.
或者獻情.
就是Present.
把人引到基督面前.
其實都是用這個詞.
Present.
怎樣去將基督已經幫助我們.
進入到和祂和好的關係.
而保羅所說的.
我們要做的.
是幫助人進入到基督裡的關係.
讓我們的生命可以成熟.
所以二十二節和二十八節.
都是提到我們呈獻.
是耶穌基督先將我們用寶血來洗淨.

$^{441}$以致我們可以和祂聯合.
但是我們的生命是否成熟呢?.
是否好像當年我自己這樣.
是的,我信了主.
我應該是被主所潔淨.
但是我生命裡面.
我卻是不理神要我去做什麼.
我仍然大發脾氣呢?.
如果我們真的引領人進入到基督裡的話.
是讓基督在我裡面掌權.
讓我掌權.
而我生命是得以長大成熟.
這個是非常重要.
所以不單止是在線上做培訓.
告訴大家有關耶穌基督的真理.
而是更多的.
讓我們知道怎樣去活出這個真理.
怎樣讓真理在我們裡面.
使我們產生生命的轉化.
如果我們不願意神介入我們的生命的話.
不願意神做我的生命的主.
是很難讓我們的生命看到轉變.
我接受耶穌基督做我的救主.
但是做我的主呢?.
等一下.
我請耶穌在我的客廳坐一坐.
我不請祂進我的客廳.
我就在房間裡看我裡面的髒東西.
這就是我們不讓耶穌基督做我們生命的主.
也就是說當我們去讀耶穌基督的相關話語的時候.
水過鴨背.
那些話語對我有什麼關係呢?.
有時候我有疑問.
我問誰呢?.
我生命當中有掙扎.
我所理解的和我所讀的經文.
是不一樣的時候.
很掙扎的時候.
怎麼辦呢?.
我前個星期去東岸.

$^{481}$參加一個單身成長營.
很多單身弟兄姊妹在那裡.
其中有一位姊妹跟我分享了這個掙扎.
信了主十幾年.
但是呢.
她很多時候在兩性關係當中有很多的不理解.
為什麼神要這樣去對待她.
在這個過程當中.
其實有很多時候是關乎到她平時.
我們平時是怎樣去讀神的話語.
還有我們有沒有一個群體.
是可以一起來彼此的勉勵.
能夠彼此的勸戒.
彼此的教導.
保羅.
他在《歌羅西書》裡面.
他說要讓人能夠進入到基督裡面.
讓他們生命得以成熟.
是不容易的.
這個不容易是包括了.
我們怎樣去陪伴.
怎樣去按著真理去勸戒.
保羅這麼說.
我要為此勞苦.
照著他在我裡面運行的大能盡心竭力.
新譯本是這樣翻譯.
我為了這事勞苦.
按著他用大能在我心中運行的動力.
竭力奮鬥.
當我們去看保羅的時候.
他真的很竭力.
他的竭力不是只是在這一次裡面.
他一直持續的.
是持續不斷的.
他靠的是什麼呢.
如果我們看回.
看善後行傳所說的保羅.
他在傳福音的過程當中.
他怎樣去讓人得生命.
而且不單只是去到一個地方去傳一傳.

$^{521}$而且他還有很多生命的建築在後面.
很多的困難.
很多的挑戰.
很多時候我們陪伴弟兄姊妹.
或者是讓一個舞蹈朋友.
讓他更知道我們所相信的是什麼.
不僅僅是請我們參加聚會.
聽某一場講座這麼簡單.
很多時候是生命的陪伴.
在過去的復仇裡面.
我和我先生我們開放家庭.
幫助我們當時只有語堂和英文堂的教會.
去建立國語堂.
經過十年我們建立了國語堂.
還有很多弟兄姊妹一起去幫忙.
在這個過程當中是很不容易的.
而後來我們也進入了UIC的校園當中.
是和一個叫Asian Campus Church來去配搭.
我們住在校園裡面.
學生在前後街.
我們是不關門的.
學生可以隨時來.
隨時來有些是哭著來失戀.
打著交流分手打交打進來.
有些是洗衣服沒地方洗.
很辛苦的幾年.
但是也讓我們去反思.
我們本身的生命是什麼.
我們有沒有活出那個.
真是基督的愛出來.
很多的挑戰.
很多的眼淚.
但是真是體會一些些.
保羅在這裡所說的.
真是靠著主在我們心裡面所運行的動力.
這個是在基督裡的.
因為我們已經與祂連接.
我們的事奉不是靠著自己的一些熱誠來去做.
更不是靠著我們自己的聰明智慧去做.
乃是靠著我們在基督裡.

$^{561}$在基督裡的真理.
以及在基督裡.
祂在我們裡面的大能大力.
保羅在後面.
其實在《歌羅西書2》一節裡已經說到.
祂的盡心竭力要做的是什麼.
要使人希望能夠使人能夠進入到那個完全的成熟當中.
明白真理.
祂心裡得安慰.
明白真理.
是能夠分辨什麼是對的.
什麼是錯的.
也是在基督裡面生命得以成熟.
還有一點很重要.
是他們能夠真正靈明成熟的是彼此相憐.
這是在《歌羅西書2》一節裡面所說到的.
在基督裡面的免渡同行會辛苦.
但是我要體會到.
是靠著基督我們是可以的.
我很喜歡一首詩歌.
《在花園裡》.
當我說前個星期跟一個職稱或者是寫稱.
他對神很不滿.
因為他很盡心事奉主.
但是他認識的幾個男朋友都反應後來發現.
要麼就去看色情網站.
他很疑問.
他覺得神愛他.
我問他有沒有跟耶穌散散步.
他說跟耶穌散步?.
我跟他分享這首詩歌.
這首詩歌原本是Charles Austin Miles寫的.
他寫的時候是1912年.
他寫的時候是在一個黑房裡面.
他沖洗照片.
他同時也用那些時間去讀經.
來默想.
當時他看到馬尼亞.
發現耶穌不在了.
找不到他的身體.

$^{601}$而那時候.
耶穌向他顯現.
在那時候.
Charles Austin Miles.
他有一個很深的體會.
覺得耶穌彷彿就在黑房裡.
就在跟他一起.
而他就寫了這個《In the Garden》.
其實在花園裡是在墓園裡.
我們花園裡是很浪漫的.
大家注意看.
一般我們的版本.
《In the Garden》.
在花園裡是頭兩節.
就是獨步徘徊.
在花園裡.
多好啊.
有溫柔的聲音傳入我耳中.
應該是神子耶穌.
第二段.
我主持的聲音多麼甜蜜.
很優美的歌聲.
他的旋律也很輕快.
有一個版本是用Waltz.
真的你聽到都想跟他一起跳.
但是很多時候我們發覺.
現在網上流傳的版本.
是沒有第三節的.
而第三節是重點.
也是Charles Austin Miles.
他寫這首詩歌的時候.
很重要的.
也是將我們與主同行.
更加是要什麼呢.
信服主.
被主所拆牌.
第三節.
天色將暗夜幕.
將暗夜幕漸睡.
我甚願與主留在園中.

$^{641}$但主命我去.
懇切呼召我.
要我心向他順從.
很多時候我們與主同行.
似乎是很享受當中的美.
但是更重要的是.
與主同行是要完成主給我們的使命.
使得萬民.
使得各族群.
使得各個年齡層.
能夠認識耶穌基督.
認識祂給我們的救恩.
這個圖我是借回來的.
希望不會侵犯版權.
大家知不知道這個圖.
是張顯在聖經裡面哪一段的經文.
我快點揭曉.
是二萬五詩.
耶穌復活之後.
跟他們一起來對同行.
一邊聊一邊說.
大家想像一下.
你有沒有與神同行.
耶穌每天與我們同行.
祂願意與我們同行.
祂願意向我們講解祂所說的話語.
祂更加願意讓我們在祂裡面.
更深入地進入與祂密切的關係當中.
幫助我們更多地去了解祂.
當我與主同行的時候.
跟我一起走的人.
很自然就會一起與主同行.
這是我們說在基督裡.
免導同行很重要的關鍵.
如果我本身沒有與主同行的時候.
我如何讓我所接觸的人.
一同走向神.
一起在基督裡與神同行.
免導同行的免禮和引導.
又如何能夠產生.

$^{681}$沒有連接.
沒有接觸,沒有連接,沒有同行.
我們又如何能夠在基督裡.
讓基督來得著.
來得成.
因著這個緣故.
我所參與現在在基督使者協會的事工.
我們有一個叫做.
青年大使命生命培育.
我們的重點是以禱告關懷為主軸.
禱告是與神的連接與神同行.
關懷是彼此的關懷.
我們在年輕人當中.
我們有並肩同行的培訓.
並肩關懷的培訓.
我們有一個M and M免導同行的使工.
M and M不是巧克力.
大寫的M是mentor.
是導師.
小寫的M是mentee.
我們在年輕人當中建立這些屬靈關係.
在城區當中有很多年輕人還沒有接觸到教會.
或者他們建立教會.
但是有幾個城區的年輕人教會.
是沒有長輩在當中.
我們與教區的一些教會來聯絡.
一些牧者也都差派他們成熟的基督徒.
進入到年輕人當中.
來做他們的Mentor.
我們也有體驗教育.
大家來藉著戶外活動.
藉著各樣的體驗活動.
包括靈修.
我如何進入與神密切的關係當中.
這個也是體驗教育的部分.
我們還有跨文化的培育.
跨文化的培育是幫助我們可以進入到不同的族群當中.
還沒有說宣教.
而是提升我們的文化相素.
Cultural intelligence.

$^{721}$是讓我們願意接觸不同文化的人.
包括代製的文化.
年輕人願意跨向長輩.
長輩願意跨向年輕人.
我們有即場的同行.
這是一個很龐大的工程.
需要大家一起來去同工.
我們還有在總部.
使者總部.
還有不同的事工.
但是我們有一個共同的目標方向.
是盼望我們能夠裝備來者.
我們從不同地方來.
裝備來者成為基督的使者.
求主幫助我們在基督裡使來者成為使者.
透過緬度同行.
我們一切都轉向神.
我們的主,我們的神.
我們能夠成為弟兄姊妹.
我們能夠有同一個阿爸父.
是因為我們主耶穌基督的救贖.
因著祂的緣故.
我們才能夠同聲.
我們在天上的父.
但是還有很多人還沒聽到救恩.
這個好消息.
很多人哪怕得著救恩.
但是生命仍然活在自己的腦窩當中.
主啊,我們很需要這個同行的關係.
我們很需要回到在基督裡的真理當中.
讓我們能夠一起靠著我們的主.
緬度同行.
求主幫助.
幫助成為社區教會.
是能夠成為一個在基督裡.
與眾人緬度同行的教會.
讓人能夠在主耶穌基督裡.
得以成熟.
呈現給我們的神.
我們交託給主.

$^{761}$求主繼續向我們說話.
交託仰望奉主耶穌基督得勝的名.
阿們.
還有在我們教會的城市裡.
我們實在需要這樣的合作夥伴.
希望將來能夠有合作和夥伴.
但是今天我們一定可以做的.
就是她剛才所給我們的挑戰提醒.
打開我們的眼界.
外面有很多人需要服藥.
打開我們的眼界,我們去禱告.
為了他們的事公禱告.
為了我們教會禱告.
為了我們自己禱告.
願上帝祝福各位.
下星期希望能夠和大家一起來敬拜.
拜拜!.
\newpage



\section{啟示錄 2:1-3:22}
\label{sec:S2QbUdfuEWU}
\textbf{基督眼中的教會 \#2 (啟示錄 2\_1-3\_22) - 陳崇基牧師}
\newline
\newline
連結: \href{https://youtube.com/watch?v=S2QbUdfuEWU}{\texttt{https://youtube.com/watch?v=S2QbUdfuEWU}} ~~~~ 語音日期: 2024-09-17
\newline
\newline
\hyperref[sec:FGEt94QjF9Q]{\small{< < < PREV SERMON < < <}}
~
\hyperref[sec:index_chronic]{\small{[返順時目]}}
~
\hyperref[sec:index_scriptual]{\small{[返順卷目]}}
~
\hyperref[sec:KVIuGXk07jc]{\small{> > > NEXT SERMON > > >}}
\newline
\newline
啟示錄 2:1-3:22
\newline
\begin{longtable}{cl}
\hline
\hline
章節 & 經文 (和合本修訂版)\\
\hline
2:1 & \begin{tabularx}{0.7\textwidth}{X} 「你要寫信給以弗所教會的使者，說：『那右手拿著七顆星，在七個金燈臺中間行走的這樣說： \end{tabularx} \\ \\ \relax
2:2 & \begin{tabularx}{0.7\textwidth}{X} 我知道你的行為、勞碌、忍耐，也知道你不容忍惡人。你也曾察驗那自稱為使徒卻不是使徒的，看出他們是假的。 \end{tabularx} \\ \\ \relax
2:3 & \begin{tabularx}{0.7\textwidth}{X} 你能忍耐，曾為我的名勞苦而不困倦。 \end{tabularx} \\ \\ \relax
2:4 & \begin{tabularx}{0.7\textwidth}{X} 然而，有一件事我要責備你，就是你把起初的愛心拋棄了。 \end{tabularx} \\ \\ \relax
2:5 & \begin{tabularx}{0.7\textwidth}{X} 所以你要回想你是從哪裡墜落的，並且要悔改，做起初所做的工作。你若不悔改，我要到你那裡去，把你的燈臺從原處挪去。 \end{tabularx} \\ \\ \relax
2:6 & \begin{tabularx}{0.7\textwidth}{X} 然而你還有一件可取的事，就是你恨惡尼哥拉派的行為，這種行為也是我所恨惡的。 \end{tabularx} \\ \\ \relax
2:7 & \begin{tabularx}{0.7\textwidth}{X} 凡有耳朵的都應當聽聖靈向眾教會所說的話。得勝的，我必將神樂園中生命樹的果子賜給他吃。』」 \end{tabularx} \\ \\ \relax
2:8 & \begin{tabularx}{0.7\textwidth}{X} 「你要寫信給士每拿教會的使者，說：『那首先的、末後的，死過又活了的這樣說： \end{tabularx} \\ \\ \relax
2:9 & \begin{tabularx}{0.7\textwidth}{X} 我知道你的患難和貧窮—其實你卻是富足的，也知道那自稱是猶太人的所說毀謗的話，其實他們不是猶太人，而是撒但會堂的人。 \end{tabularx} \\ \\ \relax
2:10 & \begin{tabularx}{0.7\textwidth}{X} 你將要受的苦，你不用怕。看哪！魔鬼要把你們中間幾個人下在監裡，使你們受考驗，你們要遭受苦難十日。你務要至死忠心，我就賜給你那生命的冠冕。 \end{tabularx} \\ \\ \relax
2:11 & \begin{tabularx}{0.7\textwidth}{X} 凡有耳朵的都應當聽聖靈向眾教會所說的話。得勝的必不受第二次死的害。』」 \end{tabularx} \\ \\ \relax
2:12 & \begin{tabularx}{0.7\textwidth}{X} 「你要寫信給別迦摩教會的使者，說：『那有兩刃利劍的這樣說： \end{tabularx} \\ \\ \relax
2:13 & \begin{tabularx}{0.7\textwidth}{X} 我知道你的居所，就是有撒但座位之處；當我忠心的見證人安提帕在你們中間，在撒但所住的地方被殺之時，你還堅守我的名，沒有否認對我的信仰。 \end{tabularx} \\ \\ \relax
2:14 & \begin{tabularx}{0.7\textwidth}{X} 然而，有幾件事我要責備你，就是在你那裡有人服從了巴蘭的教訓；這巴蘭曾教唆巴勒將絆腳石放在以色列人面前，使他們吃祭過偶像之物，並且犯淫亂。 \end{tabularx} \\ \\ \relax
2:15 & \begin{tabularx}{0.7\textwidth}{X} 同樣，你那裡也有人服從了尼哥拉派的教訓。 \end{tabularx} \\ \\ \relax
2:16 & \begin{tabularx}{0.7\textwidth}{X} 所以，你當悔改；若不悔改，我很快就到你那裡來，用我口中的劍攻擊他們。 \end{tabularx} \\ \\ \relax
2:17 & \begin{tabularx}{0.7\textwidth}{X} 凡有耳朵的都應當聽聖靈向眾教會所說的話。得勝的，我必將那隱藏的嗎哪賜給他，並賜他一塊白石，石上寫著新的名字，除了那領受的以外，沒有人認識。』」 \end{tabularx} \\ \\ \relax
2:18 & \begin{tabularx}{0.7\textwidth}{X} 「你要寫信給推雅推喇教會的使者，說：『神的兒子，那位眼睛如火焰、雙腳像發亮的銅的這樣說： \end{tabularx} \\ \\ \relax
2:19 & \begin{tabularx}{0.7\textwidth}{X} 我知道你的行為：愛心、信心、勤勞、忍耐；又知道你末後所行的善事比起初所行的更多。 \end{tabularx} \\ \\ \relax
2:20 & \begin{tabularx}{0.7\textwidth}{X} 然而，有一件事我要責備你，就是你容忍那自稱是先知的婦人耶洗別教唆我的僕人，引誘他們犯淫亂，吃祭過偶像之物。 \end{tabularx} \\ \\ \relax
2:21 & \begin{tabularx}{0.7\textwidth}{X} 我曾給她悔改的機會，她卻不肯悔改她的淫行。 \end{tabularx} \\ \\ \relax
2:22 & \begin{tabularx}{0.7\textwidth}{X} 看吧，我要使她病倒在床上。那些與她犯姦淫的人若不悔改他們的行為，我也要使他們同受大患難。 \end{tabularx} \\ \\ \relax
2:23 & \begin{tabularx}{0.7\textwidth}{X} 我又要殺死她的兒女，眾教會就知道，我是那察看人肺腑心腸的，我要照你們的行為報應各人。 \end{tabularx} \\ \\ \relax
2:24 & \begin{tabularx}{0.7\textwidth}{X} 至於你們其餘的推雅推喇人，就是一切不隨從這教訓，不明白他們所謂撒但深奧之理的人，我告訴你們，我不會再把別的擔子放在你們身上。 \end{tabularx} \\ \\ \relax
2:25 & \begin{tabularx}{0.7\textwidth}{X} 你們只要持守那已經有的，直到我來。 \end{tabularx} \\ \\ \relax
2:26 & \begin{tabularx}{0.7\textwidth}{X} 那得勝又遵守我命令到底的，我要賜給他權柄制伏列國； \end{tabularx} \\ \\ \relax
2:27 & \begin{tabularx}{0.7\textwidth}{X} 他必用鐵杖管轄他們，如同打碎陶器， \end{tabularx} \\ \\ \relax
2:28 & \begin{tabularx}{0.7\textwidth}{X} 像我也從我父領受了權柄一樣。我又要把晨星賜給他。 \end{tabularx} \\ \\ \relax
2:29 & \begin{tabularx}{0.7\textwidth}{X} 凡有耳朵的都應當聽聖靈向眾教會所說的話。』」 \end{tabularx} \\ \\ \relax
3:1 & \begin{tabularx}{0.7\textwidth}{X} 「你要寫信給撒狄教會的使者，說：『那有神的七靈和七顆星的這樣說：我知道你的行為，就是名義上你是活的，實際上你是死的。 \end{tabularx} \\ \\ \relax
3:2 & \begin{tabularx}{0.7\textwidth}{X} 你要警醒，堅固那些剩下、快要死的，因為我發現你的行為，在我神面前沒有一樣是完全的。 \end{tabularx} \\ \\ \relax
3:3 & \begin{tabularx}{0.7\textwidth}{X} 所以，要記得你所領受和聽見的；要遵守，並要悔改。你若不警醒，我必如賊一樣來到；我幾時來到你那裡，你絕不會知道。 \end{tabularx} \\ \\ \relax
3:4 & \begin{tabularx}{0.7\textwidth}{X} 然而，在撒狄你還有幾位是未曾污穢自己衣服的，他們會穿白衣與我同行，因為他們是配穿的。 \end{tabularx} \\ \\ \relax
3:5 & \begin{tabularx}{0.7\textwidth}{X} 得勝的必這樣穿白衣，我也不從生命冊上塗去他的名；我要在我父面前，和我父的眾使者面前，宣認他的名。 \end{tabularx} \\ \\ \relax
3:6 & \begin{tabularx}{0.7\textwidth}{X} 凡有耳朵的都應當聽聖靈向眾教會所說的話。』」 \end{tabularx} \\ \\ \relax
3:7 & \begin{tabularx}{0.7\textwidth}{X} 「你要寫信給非拉鐵非教會的使者，說：『那神聖、真實的，拿著大衛的鑰匙，開了就沒有人能關，關了就沒有人能開的這樣說： \end{tabularx} \\ \\ \relax
3:8 & \begin{tabularx}{0.7\textwidth}{X} 我知道你的行為。看哪，我在你面前給你一個敞開的門，是沒有人能關的。我知道你有一點力量，也遵守我的道，沒有否認我的名。 \end{tabularx} \\ \\ \relax
3:9 & \begin{tabularx}{0.7\textwidth}{X} 那屬撒但會堂的，自稱是猶太人，其實不是猶太人，而是說謊話的，我要使他們來到你腳前下拜，使他們知道我已經愛你了。 \end{tabularx} \\ \\ \relax
3:10 & \begin{tabularx}{0.7\textwidth}{X} 因為你遵守了我堅忍的道，我也必在普天下人受試煉的時候保守你免受試煉。 \end{tabularx} \\ \\ \relax
3:11 & \begin{tabularx}{0.7\textwidth}{X} 我必快來，你要持守你所有的，免得人奪去你的冠冕。 \end{tabularx} \\ \\ \relax
3:12 & \begin{tabularx}{0.7\textwidth}{X} 得勝的，我要使他在我神的殿中作柱子，他必不再從那裡出去。我又要把我神的名和我神城的名—從天上我神那裡降下來的新耶路撒冷，和我的新名，都寫在他上面。 \end{tabularx} \\ \\ \relax
3:13 & \begin{tabularx}{0.7\textwidth}{X} 凡有耳朵的都應當聽聖靈向眾教會所說的話。』」 \end{tabularx} \\ \\ \relax
3:14 & \begin{tabularx}{0.7\textwidth}{X} 「你要寫信給老底嘉教會的使者，說：『那位阿們、誠信真實的見證者、神創造的根源這樣說： \end{tabularx} \\ \\ \relax
3:15 & \begin{tabularx}{0.7\textwidth}{X} 我知道你的行為，你也不冷也不熱；我巴不得你或冷或熱。 \end{tabularx} \\ \\ \relax
3:16 & \begin{tabularx}{0.7\textwidth}{X} 既然你如溫水，也不冷也不熱，我要從我口中把你吐出去。 \end{tabularx} \\ \\ \relax
3:17 & \begin{tabularx}{0.7\textwidth}{X} 你說：我是富足的，已經發了財，一樣都不缺，卻不知道你是困苦、可憐、貧窮、瞎眼、赤身的。 \end{tabularx} \\ \\ \relax
3:18 & \begin{tabularx}{0.7\textwidth}{X} 我勸你向我買從火中鍛鍊出來的金子，使你富足；又買白衣穿上，使你赤身的羞恥不露出來；又買眼藥抹你的眼睛，使你能看見。 \end{tabularx} \\ \\ \relax
3:19 & \begin{tabularx}{0.7\textwidth}{X} 凡我所疼愛的，我就責備管教。所以，你要發熱心，也要悔改。 \end{tabularx} \\ \\ \relax
3:20 & \begin{tabularx}{0.7\textwidth}{X} 看哪，我站在門外叩門，若有聽見我聲音而開門的，我要進到他那裡去，我與他，他與我一起吃飯。 \end{tabularx} \\ \\ \relax
3:21 & \begin{tabularx}{0.7\textwidth}{X} 得勝的，我要賜他在我寶座上與我同坐，就如我得了勝，在我父的寶座上與他同坐一般。 \end{tabularx} \\ \\ \relax
3:22 & \begin{tabularx}{0.7\textwidth}{X} 凡有耳朵的都應當聽聖靈向眾教會所說的話。』」 \end{tabularx} \\ \\
[1ex]
\hline
\hline
\end{longtable}
$^{1}$Again 中秋節.
中秋節的時候想起我生日.
有誰在中秋節附近生日的.
沒有 只有我一個.
好 生日快樂.
我主要是想提醒自己.
因為有時候我自己也忘記自己生日.
Anyway.
我們繼續看這個系列.
我們看基督眼中的教會.
如果你今天第一次來教會的話.
我們每逢講道的時候.
我們都會看看聖經對我們有什麼形容.
聖經是上帝的話語.
上帝對我們的提醒.
對以前的信徒的提醒.
也是對我們今天的提醒.
所以我們正在看一個系列.
就是聖經裡面.
曾經耶穌基督寫了七封信給當時七個教會.
上一次我們看了第一個教會的信.
叫做什麼教會.
以弗所教會.
然後這次我們看第二封信.
這封信寫給一個地方.
叫做史美蘭.
在今天的土耳其的西邊.
你看到這條路線由以弗所開始.
第一個城市.
然後去到史美蘭.
然後別加摩.
推亞推拉.
薩迪.
菲拉迪菲.
然後最後就是魯迪加.
利奧特斯.
這裡就講到.
說耶穌會寫七封信給這七個教會.
這條路線大概可能是當時柔政的路線.
由以弗所直到去魯迪加.

$^{41}$大家想像一下.
如果上帝寫一封信給你.
耶穌今天寫一封信給你.
你覺得這封信會講些什麼呢.
或者這封信寫給我們教會.
這封信的內容應該會怎樣呢.
我想這樣可以幫到我們.
去看這七封信.
我上一次說了.
以弗所的教會.
然後我們這次看史美蘭.
史美蘭這個教會.
史美蘭當然是一個城市.
當時是一個古城.
有兩件事值得留意.
Smyrna這個地方.
雖然今天是一個現代化的城市.
但是以前的古城.
到底我們有什麼知道呢.
以前的城市是怎樣的呢.
和我們今天的城市有什麼分別呢.
史美蘭的古城可以這麼說.
第一件事就是.
可以叫做一個死過翻生的城市.
因為大家知道古城.
一定經歷過很多打仗,戰爭.
然後被拆毀,被燒毀.
然後他們再重建.
特別是史美蘭這個地方.
Smyrna這個地方.
曾經被燒毀多次.
又再重建.
當時有一個.
可以說這封信.
是2000年之前.
記載在聖經裡面.
但是我們知道當時的作家.
特別是羅馬的作家.
他們就會稱讚這個城市.
說這個城市好像火鳳凰一樣.

$^{81}$有沒有聽過火鳳凰是怎樣的呢.
Phoenix.
死灰復燃.
在灰燼裡面再復生.
整個城市被燒毀過.
然後再復生.
如果你想一下.
你在街上.
你和一個史美蘭城鎮裡面的居民.
聊天的時候.
他們會想起什麼.
和說起什麼故事.
就是這個史美蘭的城市.
是一個很有驕傲的城市.
等於我們住在芝加哥.
住在芝加哥的人都很有驕傲的.
是不是.
我曾經認識一個白人的牧師.
他以前住在芝加哥.
現在搬回東岸.
有時候和他一起喝咖啡.
他有個說笑.
他說他住在芝加哥的人.
芝加哥以外的都是第三世界.
是.
芝加哥以外的都是第三世界.
你聞到那種味道.
芝加哥的驕傲.
是吧.
無論如何.
史美蘭都是一樣.
死過復生的城市.
就是死不去.
是吧.
第二方面.
除了死過復生之外.
就是亞洲的第一城市.
亞洲第一城市.
為什麼當時屬於亞洲.
就是現在土耳其的地方.

$^{121}$都是屬於亞洲.
當時.
2000年之前.
是哪個帝國去到統治天下.
羅馬帝國.
第一世紀的時候.
羅馬帝國統治天下的時候.
他們認為那邊.
土耳其.
以色列巴勒斯坦地.
埃及等等.
這些地方是屬於地中海東邊的地方.
或者他們叫Roman East.
你知道羅馬帝國的祖家在哪裡嗎.
羅馬帝國的祖家當然是京城羅馬.
是吧.
他們一直擴展到東邊的時候.
他們會怎麼處理.
羅馬帝國東邊的世界.
所謂Roman East.
就是在重要的城市那裡.
然後建一些.
重要的建設.
譬如拜羅馬王的廟.
拜羅馬王為神的一些皇廟.
就是哪個城市能夠爭取到這個樹榮.
能夠幫羅馬帝國的羅馬王.
幫你的老闆.
去做一座聖殿出來.
做一座廟出來供奉他拜祭他的話.
這是一個很大的得到中央的讚賞.
以及得到中央的樹榮.
所以很多城市就會爭取這件事.
史美蘭.
Smyrna.
就是其中一個爭取到第一線城市.
第一城.
First City.
甚至說自己是第一城市在亞洲.
在亞洲裡面第一個的城市.

$^{161}$所以.
是一個很有驕傲的城市.
就是第一城市在亞洲.
在這個城市裡面有一間小的教會.
我們叫史美蘭教會.
基督寫信給教會.
第二封就是寫給這裡的信徒.
史美蘭教會.
我們看看這封信的內容.
二章第八節.
他說你要寫信給史美蘭教會的使者.
說那首先的 末後的.
死過又活著的.
這樣說.
這裡每一封信很短的.
幾節.
開頭都會有一個基督自己介紹自己的.
這裡怎樣介紹自己.
他說他是首先的.
也是末後的.
你聽到我剛才怎麼說.
史美蘭是一個怎樣的地方.
We are the first.
We are the number one.
耶穌說我才是number one.
耶穌怎樣介紹自己.
我是first.
And I am also the last.
我是第一.
也是最後的.
什麼意思.
這些每封教會的書信.
開頭耶穌基督的自稱.
他自己的介紹.
都可以在第一章找到.
這是《啟示錄》.
聖經裡面最後最後那一段書.
第一章已經是耶穌介紹了他自己是怎樣的.
然後七封書信裡面.
他都會用.

$^{201}$是符合當時當地的教會.
當地的城市.
他會再自稱自己一次.
說到史美蘭的時候.
他就說.
I am the first.
我就是首先.
我就是末後.
這個第一章已經出現了.
第一章第八節.
下面說什麼.
他說主上帝說.
我是阿拉伐.
我是俄梅亞.
是今在昔在.
以後永在的傳能者.
這裡是耶穌自己介紹.
什麼叫做Alpha.
Omega.
希臘文有.
二十幾個字母.
第一個字母就是Alpha.
最後一個字母.
就是Omega.
等於英文.
第一個字母是什麼.
是A.
最後一個字母是Z.
他說他從頭到尾都是主.
這一句的稱呼.
或者這一種的稱呼.
其實不只是新約聖經才有的.
舊約聖經也有的.
舊約聖經耶和華曾經自稱自己是什麼.
就是首先和末後.
特別是依據《挫亞書》.
他說到上帝是歷史的主宰.
是你每個人人生的主宰的時候.
他就說我就是first and the last.
現在這裡哪裡自稱.

$^{241}$耶穌自己自稱.
所以耶穌自稱說我就是神.
first and last是耶和華的自稱.
耶穌說我就是那位.
第二個的介紹就是我死過之後又活來的.
史美娜就是一個這樣的城市.
就是一個這樣的地方.
但是耶穌說我才是那位死而復活的主宰.
因為看整封書信裡面.
主題都是和生與死有關.
掌管著生與死.
掌管著生與死的是誰.
是耶穌自己.
第九節.
這裡說到教會有什麼情況.
耶穌說我知道你的患難和貧窮.
其實你卻是富足的.
也知道那自稱是猶太人所說的偽旁的話.
其實他們不是猶太人.
是撒旦會堂的人.
這句話就是耶穌對那個教會的評價.
你覺得那個教會是一個怎樣的教會.
是不是Mega Church.
不是Mega Church.
是一個窮的教會.
還有經歷很多困難的教會.
經歷什麼困難.
他繼續說.
就是一些猶太人來攻擊他們.
攻擊他們.
拜徒威爾這封信.
我們說有七封信.
由已忽所到露底家.
每封信都有彈有讚.
除了兩封.
是沒有彈只是讚.
這封其中一封.
就是史姆萊.
另外一封是沒有彈只是讚.
除了史姆萊之外就是菲拉鐵菲.

$^{281}$這間是怎樣的教會.
當然是一個很貧窮的教會.
但在上帝眼中是富足的.
大家有沒有想過.
到底誰的眼光才是最重要.
耶穌的眼光才是最重要.
耶穌怎樣看待我們教會.
這個才是最重要.
不是我們自己看待這個教會.
或者別人怎樣看待我們教會.
我們自己怎樣看待.
是不打緊的.
最重要是耶穌怎樣看待那個教會.
在這裡解釋一下到底他們發生了什麼事.
他說猶太人去攻擊他們.
毀謗他們.
又說猶太人好像在撒旦的會堂一樣.
就是猶太人他們敬拜的時候.
就是在不同城市的會堂.
他說他們的會堂不是真的會堂.
他說他們是撒旦的會堂.
這裡就要解釋清楚.
因為特別是今時這個世代.
我們很多時候.
有時候人們會誤解聖經.
或者亂用聖經.
就會說聖經就是反猶.
Anti-Semitic.
Anti-Semitism.
我們要知道到底什麼叫猶太人.
他說猶太人所說的毀謗的話.
其實他們不是猶太人.
是撒旦會堂的人.
是不是說猶太人就是撒旦呢.
不是的.
因為在新約裡面的書信.
特別是約翰的書信.
約翰的社族.
譬如啟示錄.
和約翰福音.

$^{321}$裡面說的是猶太人的時候.
他們不是說全部的猶太人.
他說The Jews.
猶太人其實是主要說.
是起初第一世紀的時候.
那些反對基督教的猶太人領袖.
反對基督教的猶太人領袖.
這些叫做猶太人.
其實你想一想都知道.
初期的教會什麼人多.
猶太人多.
他們怎麼會是反猶呢.
因為初代教會其實就是猶太教的一部分.
我們講回歷史.
當時第一世紀的時候.
耶穌死而復活升天離開了.
然後初期的教會開始到擴展的時候.
初期的教會在當時來說.
大部分都是猶太人的基督徒開始的.
他們當時是屬於猶太教的一部分.
在當時羅馬帝國他們處理猶太教的時候.
他們就會覺得.
猶太人可以有一個.
我們用一個容忍的政策來看待猶太人.
猶太人你們自己喜歡的東西都可以.
敬拜我們不逼你.
只要你納粹就行了.
所以猶太教在羅馬帝國裡面是有一些保障的.
這個就是第一世紀中的時候.
但是時間久了之後.
教會雖然當時開始的時候就是猶太教的一部分.
但是慢慢的人多了.
特別是外邦人多了.
就開始大家分道揚鑣.
所以到第一世紀末開始.
基督教就不再受猶太教那把傘裡面.
猶太教的那種保障.
特別是第一世紀末的時候.
寫啟示錄的時候.
就是第一世紀末主後大概九十年代的時候.

$^{361}$寫啟示錄的時候.
教會就遇見很多這些壓力.
因為沒有保障.
又有很多人來攻擊他.
這些就是所謂猶太教的領袖.
所以約翰說到猶太人的時候.
他就知道親羅馬的猶太人.
那些領袖.
是反對耶穌是基督的.
為什麼呢?.
因為耶穌怎麼可以做我們的救世主呢?.
他在十字架上寫的.
一個羞辱的記號.
我們怎麼能夠接受耶穌是彌撞.
救世主救我們呢?.
所以他們不能夠去接受的時候.
他們就反對.
所以他們就認為教會的人.
是可以去打壓的.
所以當時也有一些猶太人的領袖.
他就會控告基督徒或者去誹謗他.
然後抓他上去.
羅馬的法庭.
然後將他們受審.
所以這裡講清楚了猶太人是什麼意思.
他針對的是那些猶太的領袖.
或者會堂的領袖.
所以就說他們是撒旦的會堂.
因為他們背後.
是撒旦來攻擊教會的.
我們明白聖經的背景的時候.
我們就不會這麼容易亂去給標籤.
對嗎?.
所以為什麼要花這麼多時間去解釋呢?.
特別是我們今時今代.
我們對某一些族群的形容和標籤.
是不可以亂來的.
特別是美國.
我們對某一些族群.
不可以亂去說他們標籤他們.

$^{401}$說他們是壞人等等.
不要亂用聖經.
這裡也是一樣.
這裡講到猶太人.
去到猶太的領袖.
當時會堂的領袖.
來攻擊教會.
這個史姆拿教會就是首當其衝.
二章第十節.
他說.
你將要受苦.
但你不用怕.
漢,阿摩,鬼要把你們中間幾個人下在牢裡.
使你們受考驗.
你們要遭受苦難十天.
你不要致死忠心.
我就賜給你那生命的冠冕.
這個教會會經歷什麼?.
會經歷一些困難的時候.
一些很困難的時候.
你有沒有想過耶穌和我們之間的關口.
我們信了耶穌之後.
不一定是無災無難的.
是嗎?.
當然我們祈禱當然希望大家無災無難.
我沒理由為了你進醫院祈禱.
上帝你讓他進醫院吧.
不會的.
我當然希望你從醫院出院.
當然希望你生命中得到上帝的醫治.
上帝也可以醫治你.
但是這個是沒有保證的.
沒有保證的.
你信耶穌.
你接受了耶穌之後.
你作為一個基督徒.
不代表你會無災無難.
這裡說得很清楚.
你們將要受苦.
並且幾個人要下在牢裡.

$^{441}$但是你們不用怕.
因為你們受苦了多久?.
十天.
什麼意思?.
如果你回到聖經.
十天這個時間.
只是一次出現過而已.
特別是說到考驗.
就是但以理書第一章.
現在但以理.
國亡.
就是被擄.
去到巴比倫.
在異鄉裡面.
他被抓進宮廷.
然後他說.
你考驗我們十天.
跟太監長說.
我們不吃你王的擺設的筵席.
我們只吃素菜.
你給我們十天的時間.
去試煉一下我們.
然後那個故事說什麼?.
就是但以理和他幾個朋友.
雖然他們被試煉十天.
試煉完之後.
上帝為他們平反.
上帝為他們出氣.
為上帝有公義給他們.
十天之後他們經過吃素菜.
居然比其他大魚大肉的宮廷裡的少年還要好.
因為就是十天.
他說你受苦難十天.
那些熟聖經的人就會想起.
就是在說一個試煉的時間.
試煉完之後.
上帝就給他們吃素菜.
就是在說一個試煉的時間.
一個試煉的日子.
是十天的.

$^{481}$是有限期的.
但最後的目的是什麼?.
他們信心得到考驗之後.
上帝就為他們平反.
有平反和公義給他們.
史姆達教會.
就是高呼一個這樣的答案.
但上帝說你不用怕.
Do not fear.
說到基督徒受苦.
其實今天也是一樣.
我們今天坐在這裡很開心.
很安心.
當然我們北美的信徒就不用擔憂這件事.
沒有一個人今天走進教會.
或者今天去崇拜.
是會怕著有公安抓你的.
或者走的時候會怕著有公安在門口等你.
No.
我們當然沒有.
因為我們是一個自由的社會.
但是在今時今日.
有很多世界各地不同的地方.
就沒有這麼幸福.
我在這裡說一點點數字.
在Gordon Cornwell上學院.
就是美國一間福音派的上學院.
叫做Gordon Cornwell.
Gordon Cornwell有一個這樣的統計.
他說過去二千年的教會歷史裡面.
有多少信徒是信到.
因著見證耶穌而信到.
是有七千萬.
七千萬.
多少個零.
數不到.
那些七千萬信到的信徒.
在過去教會歷史裡面.
有一半是在什麼時候信到的.
在上個世紀.

$^{521}$二十世紀.
一半以上是發生在二十世紀.
大多發生的地方.
都是共產的政權底下.
或者極權的.
社會教的極權的政府底下.
受很多這樣的苦.
是不是.
甚至今天中國大陸.
在不同的地方.
都是有迫害的.
這個是真實的.
是不能夠否定的.
Gordon Cornwell的數字.
說2011.
不是很久之前.
2011至2020年.
過去十年.
有九十萬基督徒被殺.
新聞裡面都有不同的報導.
中國大陸有些基督徒下監.
ISIS.
當時有些基督徒.
埃及的一些卡布迪的基督徒被殺.
當然我們北美.
我們就不用受那麼多苦.
撒旦的攻擊不是這樣.
撒旦的攻擊都可以很厲害.
可以是試探.
可以給我們很多負面的一些想法.
或者給我們很多失望的批評.
看不起自己或者放棄等等.
我們都會面對不同的.
Temptation不同的試探.
史姆勒教會的書信.
其實是寫給我們的.
寫給我們.
他說我們將會受試煉十天.
是不是.
我們未必會被下監.

$^{561}$不過我們都會面對.
各式各樣或多或少的試煉.
聖經裡面有很多Picture.
去形容基督徒的生活.
有沒有想起什麼Picture.
是不是乾坤在沙灘邊那些.
很浪漫.
No.
聖經裡面有一些Picture.
一些Image去形容基督徒的生活.
我講幾個給你聽.
看你get到什麼意思.
基督徒的生活 基督徒的生命.
好像耕田一樣.
耕種農夫.
就是你去栽種 你去培養.
然後你很努力的.
才能夠生長到結出果子.
耕田 Farming.
第二個更常出現的.
雖然我們不是那麼多時候會留意的.
就是打仗.
爭戰.
Battle.
耕田也好 爭戰也好.
都是需要Discipline的.
對嗎.
我想告訴你.
我經常都是這樣鼓勵自己的.
我自己.
我自己覺得我在這裡.
都算是很舒服.
做基督徒.
很舒服 很安舒.
我經常都這樣勉勵自己.
既然我是那麼安舒的話.
為什麼我不努力一點.
我更加沒有理由去到閒懶.
對嗎.
這麼多年都是.

$^{601}$我都是remind我自己.
為什麼我不可以努力一點.
我和其他弟兄姊妹比.
世界各地其他弟兄姊妹.
有一次網上和一個中國大陸的牧者去傾談的時候.
我發現我們真的活在兩個很不同的世界.
對嗎.
他們擔驚受怕.
對嗎.
不開鏡頭.
對嗎.
或者那個Username都要改了.
對嗎.
我在這個地方.
我有自由.
我在這個地方我有自由.
我有protection.
所以我經常都這樣encourage和remind我自己.
為什麼我不可以更加努力去事奉.
對嗎.
這是不是你應該有的心態.
對嗎.
多麼的享受.
多麼的安舒.
為什麼我們不可以去勉勵我們自己.
努力一點.
史莫拉教會這封書信.
上帝說什麼.
你不用怕.
第一章第六節.
雖然這裡講到兩方面.
史莫拉教會說你的信徒.
你們要學一下兩件事.
如果你的信心受到試煉受到考驗的時候.
第一件事.
你首先要明白考驗是什麼.
你不需要怕.
因為經過之後上帝就會和你同在.
我覺得用另外一個角度來了解.
我們了接受苦.

$^{641}$基督徒的受苦.
基督徒的受考驗.
其實是需要的.
因為如果我們不明白另外一個角度的時候.
我們就會覺得.
為什麼上帝要這樣糟蹋我.
什麼叫測驗.
什麼叫考驗.
你有沒有想過經過測驗.
經過考驗的人.
就代表他本身是有價值的.
如果你沒有價值的話.
你就不會被考驗.
你有沒有想過那些金.
那些礦物要燒過之後.
才會有純金出來.
雜質才能夠去淨化.
真金不怕焚爐火.
如果你的生命是有價值的話.
上帝才會給你試煉.
會給你試探.
這段事已經不在史莫拉教會的書信裡面.
在比利前書第一章第六節.
雖然你們必須要在百般的試煉裡面暫時的優秀.
但是你們可以維持喜樂.
使你的信心機被考驗就被火.
試煉仍然能壞的金字更加名貴.
可以在基督耶穌顯現的時候得到稱讚榮譽和尊貴.
你記不記得在這裡說測驗的時候.
經歷考驗經歷試探的時候.
其實就好像火煉的金一樣.
煉過之後你才能夠更加純正.
你的生命也是一樣.
我的生命也是一樣.
上帝不一定是救我們脫離苦難.
但是上帝絕對能夠帶我們經歷苦難.
另一個角度去看測驗.
另一個角度去看我們處於一個可能不是很理想.
或者甚至受苦的現實.
就是上帝在做什麼.

$^{681}$暫時的優秀被考驗過之後.
經歷火的試煉就會成為正感.
想起將來的榮耀.
另一個角度將來的角度.
這就是啟示錄裡面經常提醒我們的一件事.
就是我們看我們今時今世.
是不可以單憑現在這個角度.
啟示錄就是要去打開一個新的角度新的現實.
啟示出一個你肉眼不能夠看到.
但是卻真實的一個現實.
就是將來有新天新地.
史姆特教會的書信最後結束的時候.
第十節.
你要冒要致死忠心.
我就賜給你阿心和的冠冕.
凡有耳朵的就應該聽聖靈對教會所說的話.
得勝的必不受第二次死的害.
有將來的角度去看我們今時今日.
你生活才會有力.
因為講到第二次死.
什麼意思.
死不是死一次吧.
第一次死是什麼意思.
first death.
就是說我們肉身的死去.
就是安息主懷.
但是啟示錄裡面告訴我們有第二次死.
第二次死是什麼.
我們看一段一節的聖經.
一節聖經.
最後啟示錄.
啟示將來新天新地的時候發生什麼事.
他說死亡和陰間都被永在火壺裡面.
這火壺就是第二次的死.
凡名字是記在生命冊上的.
不記在生命冊上的.
就被永在火壺裡面.
這個是second death.
什麼叫second death.
就是終末的審判.

$^{721}$你和我除了有第一次死之外.
我們都要answer第二次的.
第二次就是你的生命是如何活出來的.
這個是終末的judgment.
終末的審判.
但是基督說什麼.
耶穌基督說我已經從死裡面復活.
就是KO了死亡.
死亡都可以KO.
所以屬於基督的人就不再怕第一次死.
也不再怕第二次的死.
我不知道你今天在預見什麼.
或者過去最近.
你預見什麼可能不同的試煉不同的testing.
或者這段聖經我們應該要再讀一次.
或者這段聖經可以讓我們記住.
等我接下來這個星期.
我們預見不同的情況.
我們知道主和我們同在.
等我讀一次.
第二章第十節.
寫給史姆拉教會也寫給我們.
寫給你和寫給我.
你要冒要至死終生.
我就賜給你那心懷的冠冕.
凡有耳朵的就應當聽聖靈向宗教會所說的話.
得勝的必不受第二次死的害.
我們同心低頭祈禱.
主愛多謝你賜給我們得勝.
賜給我們這個勝利.
你是戰勝死亡.
死亡再不在我們生命裡面有任何權柄.
主愛多謝你因為我們所信所接受的主.
是一位這麼真實.
這麼有權能的主.
主愛為我們每一位.
我們遇見不同的試煉不同的困難.
上帝我們求你.
求你用天國裡面的角度幫助我們.
祂會明白我們的信心受到考驗的時候.

$^{761}$我們的生命會更加如精一樣.
主幫助我們每一位弟兄姊妹.
給我們力量給我們恩典.
讓我們生命繼續成熟.
繼續被練過之後.
更加將榮耀歸給你.
主為了我們教會.
我們求你幫助我們.
讓我們教會能夠成為一個.
可以為你發光為你作見證的人.
但是在試煉裡面.
我們不需要怕.
幫助我們在困難在逆境.
我們不要忘記上帝你和我們同在.
我們這樣祈禱奉主的名.
願上帝使用各位這個星期來見證.
拜拜!.
\newpage



\section{}
\label{sec:KVIuGXk07jc}
\textbf{What Christ Thinks of the Church \#4 (Revelation 2-3) - Rev. Mark Chan}
\newline
\newline
連結: \href{https://youtube.com/watch?v=KVIuGXk07jc}{\texttt{https://youtube.com/watch?v=KVIuGXk07jc}} ~~~~ 語音日期: 2024-10-01
\newline
\newline
\hyperref[sec:S2QbUdfuEWU]{\small{< < < PREV SERMON < < <}}
~
\hyperref[sec:index_chronic]{\small{[返順時目]}}
~
\hyperref[sec:index_scriptual]{\small{[返順卷目]}}
~
\hyperref[sec:pG_vfjaVL_8]{\small{> > > NEXT SERMON > > >}}
\newline
\newline
$^{1}$Hey guys, good morning..
I just want to show you,.
a little kid gave me a gift today.
and really surprised me.
because I wanted to have this for the longest time..
And she, after the 10 o'clock service,.
she walked to me and say,.
"Oh, you know, I have this gift for you.
"and I think you should be the best recipient for this.".
And it's actually this..
You know what is this?.
Huh?.
You don't know?.
Okay, it's not a gift box, okay?.
It's the Ark of Covenant model, okay?.
So you can see here, there are two rods here.
and there's a rectangular box.
and on top of it is the guardian of the sacred place..
What they call?.
Cherubim, yes, right, Cherubim..
Yeah, right..
And then inside, inside you see there's this thing..
The two tablets, right?.
Yeah, the stone tablets of the Ten Commandments.
and also the Rod of Aaron..
This is so little, so nice, so nice..
I want to have this for a long time,.
but I just don't want to waste money buying it..
But now I got it for free, so that's great, yeah, right?.
So that's great..
It will go on my bookshelf, right?.
Let's pray for one another and I see Chee's back.
and let's welcome him..
And he got a visa, right?.
So that's great, God answer our prayer..
Let's pray together..
Father, we give thanks for the brothers and sisters here.
and we all have different challenges.
and different issues in our lives,.
but we know that you are abundant..

$^{41}$We know that you are always faithful to us..
So we come and worship you..
We want to acknowledge that you have the lordship.
and sovereignty over our lives..
And once again, we come to you and open our hearts.
and lay down all of our burden.
and even our heart before you,.
knowing that we have to trust you..
We have to continue to rely on you,.
knowing that you are our faithful God.
and guiding our steps..
Thank you for granting Chee a successful trip to Malaysia.
and come back for the successful application of the visa..
And just pray that you continue to bless this ministry,.
the youth ministry and the English ministry here..
And also want to pray for one another..
We know there's a GA general assembly next week..
Pray that you will grant our hearts.
and our church with a sense of unity..
We come to get to know the financial situation of our church.
and any other planning for our church.
and just use us and use our gifts.
and use our talents, use our resources to build your body.
so that your church will bring glory to you..
Thank you again for bringing us together.
even in the midst of rain.
and pray that you will teach us your truth.
and may your Holy Spirit lead us into your word..
We'll get to know you more in Christ and pray..
All right guys, today is going to be part four.
of our series on seven letters to the seven churches..
And as you remember, we have gone through.
three letters already from Ephesus, Smyrna to Pergamon..
Today is church number four, Phaetia..
Phaetia is a little bit inland..
It's an old city..
It's a prominent city of textile industry..
So if you want to get some clothes or get some fabric,.
this is the place you go to in the ancient world,.
as particularly in today's Western Turkey..

$^{81}$You will find linen and wool and tanners.
and trading of different kinds of textile..
In fact, this city is already mentioned in the book of Acts..
And if you remember, in Acts 16, when Paul was in Philippi,.
he met a woman called Lydia..
Can you see that?.
Lydia is from her hometown, the city of Thyatira..
And she was a pagan..
She was a Gentile, kind of a God-fearer..
You know, she was attracted to Judaism at the time.
and worshiping the one and only God..
And she left her home city, Thyatira, and moved to Philippi..
And there she met Paul and heard the gospel.
and respond to the message of Paul..
And she became one of the first church members.
in the first century..
And by the way, if you remember,.
Paul was also a tent maker, right?.
Instead of just preaching the gospel full-time,.
he is also working part-time as a tent maker..
I think he learned this trade from his hometown,.
Thessal or Thyatira, or maybe he learned it.
from Lydia as well..
In any case, Paul was into the textile industry as well..
So we come to this city called Thyatira..
As you can see, it's a Roman city at the time..
And it's also a thriving city.
with a lot of industry going on..
And this city has a church,.
and the house churches are receiving this letter.
from Jesus, verse 18..
It says, "To the angel of the church in Thyatira write,.
"these are the words of the Son of God,.
"whose eyes are like blazing fire.
"and whose feet are like burnished bronze.".
Try to picture that, try to picture that..
Jesus introduced himself as the Son of God.
with eyes with blazing fire.
and with his feet like bronze, you know?.
He's a shining figure, right?.

$^{121}$And by the way, calling Son of God,.
today we'll understand Son of God as the divine, right?.
Someone who is fully divine..
Of course, Jesus is the truly divine Son of God,.
and He is fully divine..
However, when you mention Son of God in those days,.
remember, we have to read the Bible.
with the context in mind, right?.
In the first century context..
When you mention Son of God,.
automatically people will think of Roman emperors, right?.
Because the Roman emperors, each of them,.
each of them call themselves Son of God, Son of Divine..
And if you pick up some of the coins,.
I know we don't use coins anymore, right?.
Usually we don't use cash, but in those days,.
when they use cash,.
when they have little coins in their pocket,.
when they do trading and buying something,.
they look at the coin,.
and the coins usually have the head of the emperor,.
and on the reverse side, you will see other symbol.
with some of the inscription on it called Son of God..
So coins were a propaganda machine..
It was like the TV station of today,.
and the Roman emperor are claiming to be the Son of God..
But Jesus said, "I am the Son of God, not the emperor..
"I am the true Son of God.".
And his eyes are like blazing fire,.
and his feet are like bronze, shining bronze, right?.
All of these images actually are mentioned in chapter one,.
when Jesus showed up as the Son of God,.
rose from the dead to join the disciples..
And his appearance reminds us God himself,.
because in the Old Testament, all of these phrases,.
shining, lightning, and blazing fire,.
all of them are imageries of God showing up..
Now in the New Testament, God of the New Testament is Jesus..
So God of the Old and the New Testament are the same God,.
right?.

$^{161}$Jesus claiming himself to be the Son of God..
Indeed, he's the Son of God,.
and he is with a blazing fire and a bronze color,.
remind us of God himself, and he is the judge..
He is the supreme sovereign over the whole creation,.
just like God in the Old Testament..
So Jesus show up himself and introduce himself.
to the church in Thyatira as the Son of God.
with blazing fire as the righteous judge..
He said, "I know your deeds..
"I know your works, your love, your faith,.
"your service, your perseverance,.
"that you are now doing more than you did at first.".
Is that a good church?.
I think this is a church that deserve praises, right?.
You know, this series, we ask the question,.
what does Christ think of his church, right?.
When he think of this church in Thyatira,.
he find that this church, this small community,.
is a Christ-honoring community.
because they have good works..
What are the good works?.
He said, "I know your deeds..
"I know your good works.".
And he unpacked that..
What does it mean by showing good works?.
Your love, your love for God, your love for one another,.
and your faith, you rely on God, you trust in God,.
and your service, using your gifts to build up the church..
And lastly, it's perseverance, your patience,.
in the midst of hardship, in the midst of difficulties,.
right?.
You know, a part of our package of following Christ.
is actually not so comfortable.
because we will encounter a lot of difficulties..
We will encounter a lot of challenges, right?.
I'm sure you, since you are following Christ.
since day one, you have encountered a lot of questions.
and a lot of doubts, a lot of difficulties..
Believing Christ or relying on the Lord.

$^{201}$does not mean he will take away all the problems,.
but God promised himself to be your guidance..
God is bigger than your problems,.
and that is testing our faith, right?.
And God, a God-honoring church is a group of people.
like us who have faith and love.
and service and perseverance..
This is a Christ-honoring church with good works..
And I pray and I hope that our church is like this, right?.
Now, it's not about size,.
it's not about how much money we have,.
but about our deeds, our good works..
Jesus even praised this church by saying,.
"You are now doing more than you did at first.".
Now, how long has this church exist?.
Probably for some years, right?.
And how long you've been the Christians?.
Probably for some years, right?.
Or for some months..
You are doing more than you did at first,.
and that receives the attention of Jesus, the Lord himself..
So the first point here in this letter.
is that a Christ-honoring church has good works,.
has good works..
Do we have good works?.
Our love, our faith, can we show them to other people?.
Our service and our perseverance, right?.
And the most important thing, I think, is number four,.
is perseverance..
In the midst of difficulties, in the midst of challenges,.
we still press on, and that is the essence of faith,.
right, we believe that Christ is more than enough.
to help us, and His grace is sufficient..
And point number one, point number one is.
the church has good works,.
and that is a church that is honoring Christ..
Point number two here is that a Christ-honoring church.
takes the truth seriously..
Let's continue to read this letter..
Now, this is a long letter,.

$^{241}$let me just give you a heads up..
It's a longer letter than other letters..
So a Christ-honoring church is a church with good works..
And not only that, they take the truth,.
they have to take the truth seriously..
If you take the truth seriously, then your church,.
then our church, will be honoring to Christ..
Verse 20 and following..
"Nevertheless, I have this against you," Jesus said..
"You tolerate that woman, Jezebel,.
"who calls herself a prophetess..
"By her teachings, she misleads my servants.
"into sexual immorality,.
"and the eating of food sacrificed to idols.".
Now, this is a pretty stern warning, right?.
Pretty stern warning..
It's not very comfortable words..
However, we have to hear this..
We have to hear this..
The questions in Thyatira,.
Jesus has something against them..
And what is that?.
Here we see the name Jezebel..
Jezebel, who is this?.
Who is Jezebel?.
We call any stories about Jezebel?.
It's a female name, right?.
The other day, I was talking to some of the young couples,.
and they were trying to think of names for their babies..
Now, if you have a baby girl,.
you probably won't want to name your girl Jezebel, right?.
Got a bad rap, all right?.
Very, very bad name..
Now, Jezebel was the wife of King Ahab..
Remember King Ahab in the book of Acts,.
or in the book of Kings?.
It's a bad king, right?.
If you want to name your baby boy,.
don't pick this name, Ahab, right?.
It's a very, very bad name, right?.

$^{281}$It's a very bad name..
Why?.
Because Jezebel was originally a pagan, a Baal worshiper..
She worshipped Baal, and Baal, who's Baal?.
Baal is a storm god of Canaan..
If you go to Oriental Museum in University of Chicago,.
by the way, it's a very good museum,.
and they got a lot of artifacts..
People just don't know, right?.
Just don't know Chicago has such a great museum,.
Oriental Museum in University of Chicago,.
and they got tons of great artifacts,.
especially from the Middle East..
By the way, it's free..
You can go and just go in and take a look,.
and it's free donation, all right?.
If you go to Oriental Museum,.
you go to the gallery for Syria Palestine,.
you'll find this,.
the image in the middle,.
and that is Baal, B-A-A-L, Baal..
It means Lord..
Now, Baal is a storm god of Canaan..
You know, people would have to worship.
some sort of storm god in Canaan,.
because if you're farmers,.
what you need in Canaan to survive is rain and storm..
Whenever rainy season comes,.
which is usually the winter from October to March,.
those are the months you want enough rain.
to be rain on your farm and the land..
Otherwise, you will have famine..
You have drought..
So you need to worship some kind of storm god,.
and in Canaan, the biggest storm god is Baal himself,.
and he will bring rain and storm to the land in season..
You see Baal, what kind of posture he has?.
He has his right hand lifted up, right?.
You see the right hand, there's a hole in there, right?.
Now, we got, it's lost already.

$^{321}$because it's so many years ago, right?.
You find some of the old sites, the idol of Baal,.
and the hole actually holds a thunderbolt, you know?.
He's a thunder god..
He can send thunder and send rain to the land..
And Jezebel is a devout Baal worshiper..
And Ahab married Jezebel,.
and Jezebel brings in this false religion,.
this false god Baal into the land of Israel..
That's why it was such a big deal in those days.
because Israel is supposed to worship the only one god.
whose name is Yahweh, right?.
Whose name is Yahweh..
He is the Lord, he's the true Lord, he's the true God..
But Jezebel brings in a false religion.
and try to synchronize it, mix it with Yahweh..
So Baal and Yahweh mixed together..
And that's why you got the story of Elijah, right?.
He's a true prophet and he confront the king..
He confront the people and lead them back to Yahweh,.
the true God..
In any case, Jezebel enticed the entire,.
the whole nation of Israel to worship Baal..
And here in Thiatira, the same Jezebel show up, right?.
Now, of course, it's not the historical Jezebel.
in the Old Testament, but it's probably a nickname.
for some kind of leaders in the house churches in Thiatira..
And these false leaders and false prophets,.
they may be encouraging Christians to compromise their faith,.
to compromise their faith..
You need to understand our Christian faith.
can be compromised, right?.
We can do lip service to God,.
but in the bottom of our heart,.
we probably are compromising our commitment to God..
There are two things mentioned here..
By the way, the painting on the left is Jezebel, right?.
Isn't she beautiful?.
That's why she could entice the entire nation.
to worship Baal..

$^{361}$The two things in this letter mentioned,.
Jesus have against the church in Thiatira..
Sexual immorality..
The all pagan religions in general.
are very relaxed in sexuality, in sexual life..
All kinds of sexual activities outside the marriage covenant.
are permitted, tolerated, or if not encouraged..
So that's one thing Jesus had against the church..
The second thing is eating of food sacrificed to idols..
Now in the pagan world, in the Roman world,.
in the first century,.
you will encounter a lot of different festivals..
So moon festivals, by the way,.
we just had the mid autumn festival, right?.
It's actually kind of a pagan festival..
But in the Roman world,.
the pagan festivals happen in pagan temples..
And you need some sort of membership..
You need some sort of association.
in order to get the benefit out of it..
Now, what kind of benefit are being offered.
in pagan temples, in pagan festivals?.
Well, some of the benefits would be eating the meat..
You know, in the most of the Mediterranean diet.
are vegetarian, right?.
They don't get meat often..
But from time to time, they would join..
You have to join some sort of temple association..
And you have to come to their faith.
to worship their false god.
in order to get those benefits, right?.
Just like today, if you want to get good meat,.
you need to have a Costco membership, right?.
You got to pay for the membership, right?.
Well, lucky for us,.
Costco is just a secular warehouse, right?.
So it would not compromise our faith.
if we go to Costco today or join their membership..
But not so in the first century..
In the first century, Costco is a temple, right?.

$^{401}$You need to go to the warehouse.
and you have to worship some kind of pagan gods.
in order to get in..
You know, I have a friend in Taiwan.
who just visited Chicago not too long ago..
And he moved from Hong Kong with his young kid,.
small family, moved from Hong Kong to Taiwan for some years..
And we were having coffee and he shared with me.
why he stayed in Taiwan for a few years only.
and he could not tolerate anymore..
And he moved back to Hong Kong..
And the reality was that in the city that he lived in, Taiwan.
you know, the kids, the children,.
they go to elementary school, right?.
But from time to time, especially before their final exam,.
before some sort of standardized test,.
you know, he said the whole class of kids.
would have to go to some sort of temple next door..
You know, they got to worship,.
they got to get some lucks or some fortunes from the temple..
And then before they take the test, you know..
So it's kind of very, very strange and pagan culture..
And you have to pray to the gods and even small kids,.
you know, they grow up in that kind of environment.
and they have to pray to the pagan gods.
to bless their upcoming exams or whatever, right?.
So my friend said he just can't stand it anymore..
He can't stand the folk religion culture in Taiwan..
It's just too much for him..
So he moved away..
Now that kind of environment.
is kind of like the first century environment.
where the church or the Christians.
attempted to compromise their faith, right?.
To go to some sort of pagan gods to get the blessing.
and not knowing that the true source of blessing.
is actually from Jesus Christ himself..
So they were being tempted to have to eat the meat.
and you have to get the membership,.
you have to get the association..

$^{441}$And then the problem becomes you bury.
the line between Christians and pagans, right?.
You get one foot in there, right?.
Try to get some benefit..
Now here's the key issue..
It's compromising our faith with the surrounding culture..
Now in America, of course, our society,.
our culture is less religious, right?.
However, we also face another challenge, right?.
Secularization..
Because our society, our world is so secularized, right?.
And faith is being treated as something,.
you know, just private, just personal, right?.
You don't have to show it..
You don't have to live it..
However, if you come to the Bible,.
biblical faith is always public..
Biblical faith is always witnessing, right?.
You bear witness to other people..
I don't care what you do for your day job..
You may work in the kitchen..
You may work in the IT field, doing computer works,.
or you do, you know, any kind of works..
Your faith has to be open, has to be public..
That is biblical faith..
And also, the faith in the Bible,.
if we believe in Jesus, is in essence exclusive, right?.
Yeah, in our environment, in our culture,.
we always treasure tolerance, right?.
We tolerate all kinds of views, right?.
But biblical faith is always down to the bottom..
It's exclusive..
Jesus is the way..
It's not one of the ways..
Jesus Christ is the way to God..
You put your faith in Christ, you believe in Christ alone,.
and Christ has a claim for every single aspect of your life..
When I first came to faith,.
one of the little book that helped me.
was "My Heart Cries Home.".

$^{481}$The book title is called "My Heart Cries Home.".
And it's a story about inviting Christ into your house,.
into your home..
And in a lot of houses, there are different rooms, right?.
You got your bedroom, you got your living room,.
you got your kitchen, you got your different kinds of room,.
and even your closet..
And the story goes, Jesus is being invited into the house,.
and he has to have the key to every single room,.
because Jesus is the Lord of all, right?.
Your entertainment, Jesus is the Lord..
Whatever you watch, right?.
YouTube, TikTok, or whatever..
Christ is the Lord..
Christ is going to say, "Okay, you know, you can read this..
"You can watch this, but you cannot watch this, right?.
"You have to follow my leadership..
"Jesus is the Lord of everything.".
Our faith is comprehensive, right?.
Christ's worship should be comprehensive..
You cannot mix some of your past religion.
with Christianity, with biblical faith..
And again, our faith has to be public,.
has to be going outside and being open.
to invite other people, to see who Christ is..
Remember what Jesus said,.
"We are the salt and the light of the world,.
"and we are being called to transform wherever we are.".
That little phrase has been really, really helpful for me.
since I came to faith..
The phrase has to do with our relationship.
with the world, with the culture..
And that phrase is, is love..
That phrase is in the world, but not of the world..
We Christians are in the world..
We're not going to isolate ourselves, right?.
We are living in the world,.
but we are not of the world, right?.
You're living in the world, but not of the world..
What a reminder for us..

$^{521}$And the Christians in Thyatira are being wanted,.
wanted not to compromise their faith..
So number two, point number two,.
is a Christ-honoring church.
will take the gospel seriously..
It's exclusive, it's open, it's comprehensive, right?.
Verse 21, Jesus continues His warning..
He said, "I've given her Jezebel,".
or some of the false teachers in the midst of the church,.
"given her time to repent from her immorality,.
"but she is unwilling..
"So I will cast her on a bed of suffering,.
"and I will make those who commit adultery with her.
"suffer intensely unless they repent of her ways..
"I will strike her children dead..
"Then all the churches will know that I am He.
"who searches hearts and minds,.
"and I will repay each of you according to your deeds.".
Boy, this is very, very difficult words, right?.
Not very comfortable words..
But that's the treasure of the Bible, right?.
That's the treasure of the Word of God..
It's not just scratching you, you know,.
just to make you comfortable,.
but it's to shake you up.
so that you will live a holy life, right?.
Jesus said, "I've given her time.".
Those who compromise their faith,.
given them time to repent..
And use very vivid metaphor, right?.
The metaphor of adultery,.
the metaphor of sexual immorality.
to describe the act of compromising one's faith..
And there was also warning for these consequence,.
and the consequence is nothing but death, right?.
Now, don't take these words like,.
okay, it's talking about child abuse..
No, it is using the ancient way of saying things,.
ancient rhetoric to wake you up, to shake you up.
so that you know who you really believe in, right?.

$^{561}$That's the good news of the Bible, right?.
The goodness of the Bible,.
you have to understand the bad news first,.
that we are sinners, we are up for judgment,.
no one can save ourselves..
And that's the bad news..
Once we get the bad news,.
we understand what the good news is..
We have to repent and turn around.
and run to Jesus and we'll be saved..
And that is the warning for all of us..
If we compromise our faith in any way, any single way,.
this is the time for us to repent and to return.
so that our faith can be vibrant, we can be used by God..
Otherwise, judgment is waiting at the door..
Verse 24, let me finish this letter..
"Now I say to the rest of you in Thyatira,.
to you who do not hold to the teaching of Jezebel.
and have not learned the way of Satan,.
the so-called deep secrets,.
I will not impose any other burden on you.".
Hey, that's a good news..
Some of the people in Thyatira,.
some of the Christians are standing firm in their faith..
They're being solid in their commitment.
and they don't buy the teaching of Jezebel,.
the so-called secret, right?.
There's nothing secret about it..
It's actually satanic..
And Jesus is encouraging them,.
the rest of you who do not hold to the false teaching,.
I will not impose any other burden on you..
I will give you strength..
25, "Only hold on to what you have until I come..
To him who overcomes and does my will to the end,.
I will give authority over the nations..
He will rule them with the iron scepter..
He will dash them to pieces like pottery,.
just as I have received authority from my Father.".
Now this is from Psalm 2..

$^{601}$Remember Psalm 2 is a royal psalm.
that talks about God has chosen his son,.
his king on earth to rule the nations for him..
And that promised son is kind of like David.
and finally being fulfilled.
in the person and work of Jesus Christ..
So Jesus is quoting that psalm that applies to him..
He will rule them with iron scepter..
He will dash the nations,.
the rebellious nation in judgment,.
like breaking them into pieces like pottery.
because that is the comprehensive authority.
of Jesus Christ himself..
If he does not have that kind of authority,.
he is not Christ..
He is not the true God..
He is not the true Lord,.
but he indeed has that kind of authority over everything,.
over all nations, over every single aspect of your life..
And that's why we have to turn to him.
and worship him alone..
28, "I will give him the morning star..
He who has an ear,.
let him hear what the Spirit says to the churches..
The morning star, right?.
It's kind of a symbol for divine status..
You know, the highest status.
is to receive that kind of morning star..
And that is also from the Old Testament..
Simply put, if you overcome,.
if you continue to believe in Jesus and follow him,.
if you listen to his word,.
Jesus has promised you the morning star..
You will shine like morning star in the sky,.
in the dark sky..
And that is the status of authority.
and status of dignity..
And that is the promise to encourage each and every one of us.
to continue and not to compromise,.
to stand solidly in our faith,.

$^{641}$to hold on to our faith,.
hold on to the gospel..
Now, there are ways to avoid spiritual compromise..
Let me end with some of the points here, all right?.
How do we avoid compromising our faith?.
Number one, well, first of all,.
you have to identify your weakness,.
your shortcoming, where you have to compromise..
So first step is to identify..
You have to know your situation..
There are some environments,.
there are some situations,.
some places that we put ourselves in.
that will lead us into compromise..
We have to identify it, right?.
Bring it to the open, you know?.
Tell other people to pray for you, right?.
To identify the situation where you are often being tempted..
And second is to repent..
Confess our sin, right?.
Ask God for mercy, and He's always there..
He's always there..
Confess our sin..
And number three is set boundaries..
That means having the courage to walk away from temptations,.
walk away from places that will lead us.
into compromising our faith, our commitment..
Number four is to renew our spirit and our mind..
Now, there's no shortcut here..
To renew our mind, to renew our spiritual life.
so that our spiritual life will be vibrant,.
there's no shortcut..
We need much prayers..
We need to continue to pray..
For us, stick to our prayer life,.
pray consistently, and reading the Word of God, you know?.
Immerse ourselves in the Word of God..
Knowing the Word of God is a daily habit..
There's no shortcut..
We have to build this kind of habit, right?.

$^{681}$So that our spiritual life can be strong..
And also, we have to plug into the community..
No one can be a Christian alone, by himself or by herself..
We need other people, right?.
We need to serve other people,.
and other people have to encourage us.
and to keep us accountable, right?.
These are the ways that we can avoid spiritual compromising..
Let's close with a short word of prayer..
We're gonna pray to God,.
asking Jesus at this point,.
hey, are there any places, you know,.
I find myself compromising my commitment?.
Confess our sin..
And ask God for the courage and the strength to walk away,.
and to renew our spirit,.
so that we can truly, truly worship God..
Let's all pray together..
(soft piano music).
Father, listen to our prayer..
Tell us the good news..
So, though we are broken,.
although we are sinful,.
but you are there..
Thank you for what you have done for us in Jesus..
Come before you, we pray that you will remind us..
Shine a light into our hearts, any dark corners..
That we do have some places.
that we compromise our commitment to you..
And I pray, and I pray that as I confess our sins,.
you'll cleanse us..
Because you are righteous, you are holy,.
and you are full of grace..
Father, cleanse each and every one of us..
Cleanse our heart and our mind..
Renew our spirit,.
knowing that you are full of love and full of grace..
That's why we run to you..
As we confess our sin, you will grant us forgiveness..
Yes, indeed, you will..

$^{721}$Thank you for your forgiveness,.
and thank you for your acceptance..
I pray that you will renew our spiritual life,.
so that we will not continue to be weak..
We'll be strong because of your grace..
We'll be solid in our commitment because of your grace..
Renew our mind and our lives,.
so that we can follow you wholeheartedly..
Can worship you and call you our Lord and our Savior..
And you'll lead our life..
Let your comprehensive lordship over us..
We surrender our life..
Help us and cleanse us and renew our lives,.
so that we can be used by you..
May our faith be open..
May our faith be inviting other people.
to enjoy your presence, to come to taste your grace..
May your Holy Spirit be with us as your church..
And we are a community that is totally,.
totally committed to you,.
to serve you and to honor and to worship you..
Thank you for listening our prayer..
Thank you for your everlasting grace..
Again, we come to worship you and we love you..
Thank you very much for being our loving.
and gracious Father..
In Christ name we pray..
Well, once again, we are very blessed to have you,.
and have your presence and to worship together..
And today we're gonna have Bible study.
and we're gonna have a quick lunch, right?.
No, Patrick, he's already outside and prepare..
Is it taco today?.
Yeah, we have a taco lunch, all right?.
And then after the lunch, we will speak to two groups,.
Chinese group and English speaking groups,.
and we'll continue to dig into the Word of God, right?.
We will continue our study on the Book of Romans,.
and chapter one and two..
And I tell you, we have to understand these two chapters..

$^{761}$The good news of Jesus Christ.
come from these two chapters, all right?.
So stay with us, all right?.
\newpage



\section{}
\label{sec:pG_vfjaVL_8}
\textbf{What Christ Thinks of the Church \#5 (Revelation 2-3) - Rev. Mark Chan}
\newline
\newline
連結: \href{https://youtube.com/watch?v=pG-vfjaVL_8}{\texttt{https://youtube.com/watch?v=pG-vfjaVL\_8}} ~~~~ 語音日期: 2024-10-01
\newline
\newline
\hyperref[sec:KVIuGXk07jc]{\small{< < < PREV SERMON < < <}}
~
\hyperref[sec:index_chronic]{\small{[返順時目]}}
~
\hyperref[sec:index_scriptual]{\small{[返順卷目]}}
~
\hyperref[sec:9bZktUw13B0]{\small{> > > NEXT SERMON > > >}}
\newline
\newline
$^{1}$- Well, thanks, Kelly, and make sure you join our family gathering, the GAA, the General.
Assembly..
We have so many wonderful updates and so many news for you, and not only for staff, for.
myself, and for our ministry as well, and for our possible expansion of our church,.
and what's the next step for us, right?.
The church has to go forward all the time, and so has your spiritual life, so has your.
Christian life..
You need to think about what's next for you guys, and that's why we want to focus on this.
series..
And by the way, some update is that my plan in the next few months is to take some of.
you out for lunch or dinner..
How does that sound?.
Pretty good, right?.
I'll pay for it..
Yeah, and I want to spend more time to have fellowship together..
I mean, I don't want us to feel like we are separated and isolated, so don't be surprised.
if I ask you out for dinner or for lunch, okay?.
There's nothing wrong with you, okay?.
So don't have the misconception that, okay, a pastor is going to have a meal with you.
and there must be something wrong with me, right?.
No, it doesn't work that way..
Maybe there's something wrong with me, not you, okay?.
So that's part of my initiative to do some pastoral care among us..
And thanks, Kelly, for reminding us..
We still will have the afternoon Bible study, and today we're not going to have it, but.
every single second and fourth Sunday of the month, and we still want to continue with.
that until there are some other changes..
But I always value the time to do Bible study with sisters and brothers..
Remember last time, I was in the English group, and there were many, many good questions asked,.
even though we moved so slow in terms of covering the passages and the verses, but that's okay..
We need to study together, and whenever I gather with the people of God to study the.
Word of God, I always learn something new, you know?.
I always learn something new..
I always learn something that is not from the books, that is not from my preparation,.
but from our conversation, we can get to know God and His Word better, okay?.
So this series will continue with what Christ thinks of us, thinks of our church..
What does Jesus Christ, the head of the church, think of you and me?.
So we study the seven letters to the seven churches, and we start with the first one,.
Ephesus, and then we move on to the second letter to Smyrna, and then the third one to.
Pergamon, and then via Tyre, and now we come to the fifth letter, which is a short letter.

$^{41}$to the church in Sardis..
And because this is a short letter, and why don't we read together, read the entire letter.
first and see what is going on here, what does Jesus have to say to His church back.
then and to our church today, all right?.
So chapter two, or chapter three, sorry..
To the angel of the church in Sardis write, "These are the words of Him who holds the.
seven spirits of God and the seven stars..
I know your deeds..
You have a reputation of being alive, but you are dead..
Wake up, strengthen what remains, and it's about to die, for I have not found your deeds.
complete in the sight of my God..
Remember therefore what you have received and heard..
Obey it and repent, for if you do not wake up, I will come like a thief, and you will.
not know at what time I will come to you.".
Verse four, "Yet you have a few people in Sardis who have not soiled their clothes..
They will walk with me dressed in white, for they are worthy..
He who overcomes will, like them, be dressed in white..
I will never blow out his name from the book of life, but will acknowledge his name before.
my Father and His angels..
He who has an ear, let him hear what the Spirit says to the churches.".
Remember, it's churches..
It's not just for the church back then in Sardis, but for all of us, for every single.
one of us, if we call ourselves as Christians..
Alright?.
Now, one of the themes of this letter, only six verses long, I don't know if you pick.
it up, the key word, always repeat, you know, in the six verses..
It's name..
Reputation..
Has to do with the name..
Now, as you know, name in the biblical culture is more than just a label, right?.
It's more than just a label to call you, right?.
Or name on your name tag, right?.
Name in the biblical culture has always to do with reputation..
What are you famous for?.
What are you known for?.
What do other people see you and what is your reputation in the community or in the world?.
That is what the name is all about..
And in this short letter, the name repeats again and again..
And there are two main points I want to cover today..
Is that Christ, number one, cares not your earthly name, your earthly reputation..

$^{81}$Jesus doesn't really care..
Cares not so much about your earthly reputation, all right?.
As we talk about these letters, we understand we have to understand the letters or read.
the letters in terms of their own context, right?.
Sardis, the churches in Sardis or the small community house churches in Sardis, Christ.
is going to remind them in their own context..
And we know Sardis, of course now it's a modern city, but back then it was a very famous.
city..
And it was one of the richest and wealthiest capital city of the ancient kingdom called.
Lydian kingdom..
And it's in today's western Turkey..
As you can tell from this verses, Christ is writing a letter to Sardis..
And he introduces himself pretty much similar to how he introduced himself in the first.
letter to Ephesus..
He's holding seven stars and he's holding seven spirits..
And he has everything under his control..
And the letter format goes, he knows something about the church..
He said, I know your deeds and you have the name of being alive..
Your reputation of being alive..
Now if you're living in Sardis back then in the first century, you know that you are.
living in the one of the most famous cities in the world..
You know, Sardis was known for gold and silver because they have gold mine and silver mine..
As I prepared, I remember when I was younger, I used to play a lot of video games..
I don't know if you play the game called age of empire..
Ring any bell?.
No?.
Nobody like video games?.
What's your childhood like?.
No, age of empire is a strategy game and you start with some of the people of your own.
and then you go out to different places to explore and you have to find resources of.
the land..
And then build up your kingdom and build up your army and then you're going to beat up.
somebody, right?.
Beat up some of the neighbors and then you're going to conquer them..
That is age of empire..
What you need in that kind of strategy game, okay, is to find all the resources, natural.
resources like gold mine and silver mine..
And you have to gather the resources and you know, gathering back to your collection center.
and then you're going to build up your people, build up your buildings and build up your.
kingdom and then you could, you will be stronger..

$^{121}$Now, Sardis was one of the cities in the ancient world that is close to a big gold mine..
That's why Sardis was one of the wealthiest kingdom, wealthiest city state in the ancient.
world..
They have the reputation and they are famous..
But Jesus writing to them, yes you guys are rich in terms of your earthly resources..
You have the reputation, famous for being a city of gold..
But Jesus said, no, you are dead..
That's pretty interesting, right?.
Because in Greek, in the original language, there were only two words..
You guys are dead..
You are dead..
They are most famous for their gold mine..
As you can see, the archeologist is thinking up so much..
There's a river nearby that you can just collect the soil from the riverbed and then you can.
sift it and then you will see some of the clumps of gold pebbles..
And then you're going to put it in the crucible and then you're going to refine it, right?.
And then you're going to get some gold out of it..
As you can see in this center picture, right in the center, they have some samples in test.
tube..
Some of the soil from the river, you can see it's kind of shiny, right?.
So Sardis is being known for this..
Because the nearby river is a gold bearing sand..
And all these pieces of jewelry of gold and silver are what they are famous for..
I remember I first moved to the United States in 1990, right?.
1990..
Oh boy, it's 30 years ago, right?.
30 some years ago..
I live in Sacramento..
Sacramento..
It's in Northern California, Central Valley..
It's pretty close to San Francisco, the Bay Area, right?.
So as you know, in the mid, 100 years ago, some 100 years ago, San Francisco and Sacramento,.
Auburn, all those towns were blooming cities, right?.
Because there were so many people coming to the Bay Area, coming to San Francisco and.
Sacramento area for the gold rush, right?.
Yeah, 49ers, right?.
It's the gold rush, right?.
That's why San Francisco is still called (speaks in foreign language).
Because they have gold, right?.
And there were so many Chinese immigrants, actually, coming to the area and they want.

$^{161}$to get some gold, right?.
To be rich and/or to work for gold mine and industries..
But today, if you go to those small towns of, you know, it used to be the gold rush.
towns, they are almost like ghost towns, right?.
Because all the gold is gone, right?.
All the gold and silver are gone, right?.
And they were famous for their resources, but not anymore..
It is the same way in Sardis..
Yes, you have reputation, Jesus said, of being alive, but now you are dead, right?.
So Jesus is not talking about financial resources or the physical resources that they have..
He's talking about the people in Sardis, especially Christian people in Sardis..
Yes, you were known for being wealthiest, but in my eyes, you are dead..
Jesus always focuses on your spiritual status, not your paycheck, not your bank account..
Yes, those are important if you are a good and faithful servant, managing God's resources..
God demands us to be good stewards of your financial resources..
That's why as Christian people, as followers of Christ, we need to have the habit of giving,.
right?.
Living in generosity..
Generosity, right?.
When I say generosity, often we think, okay, those are for the billionaires, right?.
They can practice generosity..
We cannot be generous because, geez, we just don't have much money in our pocket..
Wait, seriously?.
No..
Jesus was the most generous person in the world, right?.
Yet he was not rich in the worldly eyes, all right?.
So Jesus' followers have to practice that as well, to give and to be generous, all right?.
To give and to be generous..
To the church in Sardis, Jesus says, "You guys are dead.".
All right?.
You guys are dead..
And then he issued a series of five commands..
It's like a shotgun, you know?.
Boom boom boom, rapid fire..
Number one, wake up..
Wake up..
Strengthen what remains, the little bit remains in you..
You have to strengthen it..
It's about to die..
That means your spiritual life is about to die..

$^{201}$You're going to lose your salvation..
You're going to lose your place in God, all right?.
So wake up and strengthen what remains, what still remains in you, and it's about to die..
Because Jesus said, "I have not found any of your deeds complete in my sight.".
Number three, remember..
See, this is similar to the first letter to Ephesus..
You got to remember..
Remember what you have received..
The gospel, the grace of God, the blessing of God..
And what you have heard..
You got to remember all those things to rewind them, right?.
Go back to how many years ago or how many months ago you received Jesus Christ and he.
has led your life to this point..
And we source you and give you all this blessing and provision..
You got to remember those things, right?.
Remember the early days you were not willing to go to church..
Maybe somebody, you know, invite you or push you around, you know, and you know, insist.
that you come to church..
And you heard the gospel the very first time, right?.
The very first time you impact and you are touched by the love of God..
You see the cross no longer different, right?.
Remember those days..
Remember what you have received, the love of God, the grace of God..
And then obey it, number four..
Obey it..
Be obedient..
And then repent..
Remember Jesus said if you do not repent, if you do not wake up, I will come to you.
as a thief..
He is talking about second coming..
In the fifth pages of the gospel, Jesus says I will come anytime to you..
You know, I think this is a very, very simple truth, right?.
Either Jesus comes anytime or you go to see Jesus anytime, right?.
We will see Jesus and it will be a sudden moment..
It's not going to, you're not going to plan like when you're going to see Jesus exactly..
You can't..
That's a picture of second coming..
Five commands as I said is a series of rapid fire..
It's to trying to shake you up, it's to wake you up..
Don't we all need to be waken up?.

$^{241}$You know, sometimes Christians, I think we are pretending to sleep..
Pretending you know, nothing's going on, you know..
Follow the motion..
You know, I can ignore God one week after another week, one day after another day..
You know, this is really sad because you know, you heard of it..
You cannot wake someone pretending to sleep..
You know, it is pretend that we don't need, we don't have, we don't have the voice of.
God..
You know, and actually the voice of the Holy Spirit keeps on reminding us every single.
day, every night..
Holy Spirit is knocking on our door..
Are you pretending to sleep?.
Are you pretending that Jesus is not speaking to you?.
Perhaps you have been faithful Christian for a long time..
But that was history, right?.
Now I want to tell you, I remind myself all the time, right?.
I have many, you know, if you check my resume, I have some experiences, right?.
I have some histories of serving God, you know, being faithful, right?.
Yeah, I think I did something according to the will of God..
However, I always remind myself, so what, right?.
Those are histories..
As glorious as they sound, those are in the past..
It's bygone..
What matters most is right now, whether you are dead or not, whether you are alive or.
dead, right?.
Perhaps you had a good reputation before in the church or in your career or in your.
workplace..
You know, you were a good witness or you did something for Jesus..
But what matters most is not in the past..
It's whether you are having a good life, a spiritual life..
Whether you in the eyes of God are alive or you are dead, right?.
Especially some of the churches, ancient churches or old churches, if you walk around the city.
or you go to different places in Europe or in many different towns and cities, you see.
huge ancient, beautiful, you know, Gothic buildings, right?.
Church buildings..
And those were the glorious days for those churches..
But you walk into those churches, there was no vitality there..
There's no spirit there..
And the people of those old churches, they are simply holding on to the property, the.
deeds..

$^{281}$But Jesus said, "You are indeed there.".
Last week, some of the brothers and sisters, we have this little expansion committee and.
we're talking about the long-term future for our church, right?.
So we've been renting this place for the last eight years..
I still remember day one I was here, we got almost nothing here, right?.
And then in the core team, in the last couple of core teams, we were talking about, "Okay,.
maybe we should look forward to a number of years from now..
What can we do?.
How can we expand and develop our ministry and build up our church, right?".
And I hope this is not just my question..
It's not just a question for the core team leaders, but it's the question for all of.
us..
It's the concern for all of us..
How can we expand?.
So last week, out of curiosity, someone texted me and, "Hey, would you like to see one of.
those old church properties on sale?.
Are you going to check it out?".
And I said, "Yeah, let's go.".
Because it's close by my home and why not?.
I just being curious what's out there in the market, right?.
So I went and with other brothers and sisters, Jinfang was there..
And we went in there and there was no more people in there..
No more people in there..
The roof is leaking water and all those things..
And very, very old churches, but look outside..
Look from the outside..
It was intact and pretty nice for the old days..
For the old days..
We have a steeple and pretty beautiful facade..
And how does Jesus see the church?.
How does God see the church?.
Right?.
It's not a church building..
No, by the way, I want to remind you, whenever the New Testament uses the word church, it's.
not talking about buildings..
It's not talking about physical building..
Church is a community..
It's a people..
It's living people..
That is church, right?.

$^{321}$Community of Christians who are faithfully following Christ..
How does God sees the church?.
As I said, Christ doesn't really care about your earthly, our earthly reputation..
What does God cares about?.
Christ cares about your heart..
Your spiritual status..
Your relationship with God himself..
So that's point number one..
Christ, number one is Christ cares not your earthly name..
So often we care so much about our earthly reputation..
We want to collect gold..
We want to collect all the resources..
We don't go to church because we have to go to work..
We have to earn some more money..
All the inflation..
You know how much, you know how expensive the eggs are?.
It's so expensive..
My wife told me, "You don't eat eggs anymore in the morning.".
Because I usually eat two eggs..
I fry two eggs in the morning..
Back then it was 150, 160 a dozen..
I ate two eggs every morning..
That's all I eat..
I need some protein..
But last week she told me, "You don't eat any eggs..
Leave it to your daughter.".
I don't eat any eggs anymore..
You know, there's financial challenges..
There's stress in our lives..
But sometimes we just care so much about our reputation..
You know, we want to collect all these golds and all these resources so that we can build.
up our own life..
You know, that's what we care about..
That's what we're concerned about..
No..
Jesus said, "Yes, you have to be a faithful steward..
Yes, you have to manage the resources God has given to you..
Yes.".
But Jesus cares so much more about your heart, about your life..
Alright?.

$^{361}$The Holy Spirit is knocking on your door today..
And Jesus is writing to you today..
Where are you at this point?.
See how you worship me..
See how you follow me..
See how you pray to me..
That's what Jesus cares about..
Christ cares about your heart right now..
Not in the past..
Right now..
Verse 4..
"Yet you have a few people in Sardis who have not soiled their clothes..
They will walk with me, dressed in white, for they are worthy..
He who overcomes will, like them, be dressed in white..
I will never blow out his name from the book of life, but will acknowledge his name, his.
reputation before my Father and his angels..
He who has ears, let him hear what the Spirit says to the churches.".
Right?.
Take note how Jesus look at the church..
Jesus delighted to see only a few people..
Only a few people in Sardis, in the churches of Sardis..
And those few people have been faithful to him..
Right?.
Have been faithful to him..
Yes, God sees and God values and concerns about the faithful..
Even though they are small in number, you have a few people I see..
Right?.
God sees a few of them because they are dressed in white garments..
In the book of Revelation, in the entire book, white garments and clean garments show up.
many many times and symbolize the positive things in our spiritual life..
Purity, holiness, dedication, victory, and glory..
Those are the people, the saints, who follow Christ faithfully no matter what..
Right?.
There are some of the ways we should revitalize our spiritual relationship with God..
Right?.
At this point I want to close in some time of reflection for all of us..
Alright?.
And we want to do so, you know, quietly and to pray to God..
And let me suggest to you, how do we revitalize our spiritual life?.
Remember what Jesus cares about..

$^{401}$It's your heart..
How do we revitalize our spiritual life?.
Number one, number one, God's grace, remember, God's grace and God's love is always there..
Even we are not faithful at times..
But his love never changes..
Now we read through this letter and we find there are so many words of warning..
We often take warning as punishment..
Harsh and unkind words..
Right?.
But the spiritual warning we find in the Bible, it is there because of God's love for us..
It's not because God yearns to punish us, to yell at us..
God out of his love, out of his grace, try to say those words of warning for all of us..
He wants our names to be on the book of life..
He wants our names to be remembered eternally..
And that's why God wants us..
So how do we go back to God and how do we revitalize our spirit?.
You know?.
So first, number one is don't you forget the grace of God, the love of Jesus, they are.
constant even though we may be faithless at times..
And second is to repent..
Is to repent..
Repentance in the Bible is turned around 180 degree..
Yes, I say to God, you know, "Sorry Jesus, help me return back to you.".
Right?.
So I have this prayer here..
I have this prayer here..
"Heavenly Father, thank you for everything you do for me..
I come before you to repent of my sins.".
And what are those?.
The dotted lines..
You can see..
You can fill it in..
You know..
Right?.
And I want to turn back to you..
And that's our prayer today..
Why don't we do that at this time?.
Reflect on the love of God and we turn, repent..
Alright?.
Let's spend a minute or so and this is time before you, before the loving Christ..

$^{441}$Amen..
Let's all pray this prayer together..
"Heavenly Father, thank you for everything you do for me..
I come before you to repent of my sins..
Please help me turn away from my sins and turn toward you, Lord..
I pray you would help me follow you faithfully and give me the strength and grace today..
In Jesus' name I pray..
Amen.".
With the love of God and grace of God, we are being sent out, okay, to follow Christ.
faithfully..
So that's why I picked this verse of the month, okay?.
We did this for the first of the week..
And I feel like maybe we're going so fast, too fast..
First of the month, so in the next month, October, we're going to repeat that again.
and again, alright?.
And don't you understand this is not words of punishment..
It's the words of loving promises from God..
Here 1 John 1, 9 says, "If you confess, if we confess our sins, He, God, is faithful.
and just and will forgive us our sins and purify us from all unrighteousness.".
Remember this, the grace of God, alright?.
The forgiveness of sin..
And we can start all over..
To love Christ and follow Him faithfully, alright?.
With this word we are sent out and we have refreshment outside, right?.
I don't know, right?.
Yes, we do, right?.
And I want to do one thing new today..
You know, a brother remind me when we have refreshment, sometimes we gather around our.
own little table, alright?.
May I ask you today to gather at the big table, alright?.
We're going to chat and we're going to talk about life, okay?.
And enjoy the food together, alright?.
Go to the big dining table, alright?.
If you can, okay?.
Squeeze in, alright?.
And then one o'clock we'll have the GA right here, General Assembly..
And we have very important updates for you, especially financial updates, alright?.
Alright, God bless you..
(soft music).
[MUSIC PLAYING].

\newpage



\section{使徒行傳 7:2-60}
\label{sec:9bZktUw13B0}
\textbf{史上最撩交打ge 自辯 (使徒行傳 7\_2-60) - 何建邦傳道}
\newline
\newline
連結: \href{https://youtube.com/watch?v=9bZktUw13B0}{\texttt{https://youtube.com/watch?v=9bZktUw13B0}} ~~~~ 語音日期: 2024-10-01
\newline
\newline
\hyperref[sec:pG_vfjaVL_8]{\small{< < < PREV SERMON < < <}}
~
\hyperref[sec:index_chronic]{\small{[返順時目]}}
~
\hyperref[sec:index_scriptual]{\small{[返順卷目]}}
~
\hyperref[sec:1ry__RW4dIE]{\small{> > > NEXT SERMON > > >}}
\newline
\newline
使徒行傳 7:2-60
\newline
\begin{longtable}{cl}
\hline
\hline
章節 & 經文 (和合本修訂版)\\
\hline
7:2 & \begin{tabularx}{0.7\textwidth}{X} 司提反說：「諸位父老弟兄請聽！從前我們的祖宗亞伯拉罕在美索不達米亞，還沒有住在哈蘭的時候，榮耀的神向他顯現， \end{tabularx} \\ \\ \relax
7:3 & \begin{tabularx}{0.7\textwidth}{X} 對他說：『你要離開本地和親族，往我所要指示你的地去。』 \end{tabularx} \\ \\ \relax
7:4 & \begin{tabularx}{0.7\textwidth}{X} 他就離開迦勒底人的地方，住在哈蘭。他父親死了以後，神使他從那裡搬到你們現在所住的地方。 \end{tabularx} \\ \\ \relax
7:5 & \begin{tabularx}{0.7\textwidth}{X} 在這裡神並沒有給他產業，連立足的地方都沒有，但應許要將這地賜給他和他的後裔為業，雖然那時他還沒有兒子。 \end{tabularx} \\ \\ \relax
7:6 & \begin{tabularx}{0.7\textwidth}{X} 神這樣說：『他的後裔必寄居外邦，那裡的人要使他們作奴隸，苦待他們四百年。』 \end{tabularx} \\ \\ \relax
7:7 & \begin{tabularx}{0.7\textwidth}{X} 神又說：『但我要懲罰使他們作奴隸的那國。以後他們要出來，在這地方事奉我。』 \end{tabularx} \\ \\ \relax
7:8 & \begin{tabularx}{0.7\textwidth}{X} 神又賜他割禮的約。於是亞伯拉罕生了以撒，在第八日給他行了割禮；後來以撒生雅各，雅各生十二位先祖。 \end{tabularx} \\ \\ \relax
7:9 & \begin{tabularx}{0.7\textwidth}{X} 「先祖嫉妒約瑟，把他賣到埃及去，神卻與他同在， \end{tabularx} \\ \\ \relax
7:10 & \begin{tabularx}{0.7\textwidth}{X} 救他脫離一切苦難，又使他在埃及王法老面前蒙恩，又有智慧。法老派他作埃及國的宰相兼管法老的全家。 \end{tabularx} \\ \\ \relax
7:11 & \begin{tabularx}{0.7\textwidth}{X} 後來全埃及和迦南遭遇饑荒和大災難，我們的祖宗絕了糧。 \end{tabularx} \\ \\ \relax
7:12 & \begin{tabularx}{0.7\textwidth}{X} 雅各聽見在埃及有糧，就打發我們的祖宗初次往那裡去。 \end{tabularx} \\ \\ \relax
7:13 & \begin{tabularx}{0.7\textwidth}{X} 第二次約瑟與兄弟們相認，法老才認識他的家族。 \end{tabularx} \\ \\ \relax
7:14 & \begin{tabularx}{0.7\textwidth}{X} 約瑟就打發人，請父親雅各和全族七十五個人都來。 \end{tabularx} \\ \\ \relax
7:15 & \begin{tabularx}{0.7\textwidth}{X} 於是雅各下了埃及，後來他和我們的祖宗都死在那裡； \end{tabularx} \\ \\ \relax
7:16 & \begin{tabularx}{0.7\textwidth}{X} 他們又被遷到示劍，葬於亞伯拉罕在示劍用銀子從哈抹子孫買來的墳墓裡。 \end{tabularx} \\ \\ \relax
7:17 & \begin{tabularx}{0.7\textwidth}{X} 「當神應許亞伯拉罕的日期將到的時候，以色列人在埃及人丁興旺， \end{tabularx} \\ \\ \relax
7:18 & \begin{tabularx}{0.7\textwidth}{X} 直到另一位不認識約瑟的王興起統治埃及。 \end{tabularx} \\ \\ \relax
7:19 & \begin{tabularx}{0.7\textwidth}{X} 他用詭計待我們的宗族，苦待我們的祖宗，強迫他們丟棄嬰孩，使嬰孩不能存活。 \end{tabularx} \\ \\ \relax
7:20 & \begin{tabularx}{0.7\textwidth}{X} 就在那時，摩西生了下來，神看為俊美，在父親家裡被撫養了三個月。 \end{tabularx} \\ \\ \relax
7:21 & \begin{tabularx}{0.7\textwidth}{X} 他被丟棄的時候，法老的女兒拾了去，當自己的兒子撫養。 \end{tabularx} \\ \\ \relax
7:22 & \begin{tabularx}{0.7\textwidth}{X} 摩西學了埃及人一切的學問，說話辦事都有才能。 \end{tabularx} \\ \\ \relax
7:23 & \begin{tabularx}{0.7\textwidth}{X} 「他到了四十歲，心中起意去看望他的弟兄以色列人。 \end{tabularx} \\ \\ \relax
7:24 & \begin{tabularx}{0.7\textwidth}{X} 他見他們中的一個人受冤屈，就庇護他，為那被壓迫的人報仇，打死了那埃及人。 \end{tabularx} \\ \\ \relax
7:25 & \begin{tabularx}{0.7\textwidth}{X} 他以為他的弟兄們必明白神是藉他的手搭救他們，他們卻不明白。 \end{tabularx} \\ \\ \relax
7:26 & \begin{tabularx}{0.7\textwidth}{X} 第二天，他遇見有人在打架，就想勸他們和好，說：『二位，你們是弟兄，為甚麼彼此欺負呢？』 \end{tabularx} \\ \\ \relax
7:27 & \begin{tabularx}{0.7\textwidth}{X} 那欺負鄰舍的人把他推開，說：『誰立你作我們的領袖和審判官呢？ \end{tabularx} \\ \\ \relax
7:28 & \begin{tabularx}{0.7\textwidth}{X} 難道你要殺我像昨天殺那埃及人一樣嗎？』 \end{tabularx} \\ \\ \relax
7:29 & \begin{tabularx}{0.7\textwidth}{X} 摩西聽見這話就逃走了，寄居於米甸地，在那裡生了兩個兒子。 \end{tabularx} \\ \\ \relax
7:30 & \begin{tabularx}{0.7\textwidth}{X} 「過了四十年，在西奈山的曠野，有一位天使在荊棘的火焰中向摩西顯現。 \end{tabularx} \\ \\ \relax
7:31 & \begin{tabularx}{0.7\textwidth}{X} 摩西見了那異象，覺得很驚訝，正往前觀看的時候，有主的聲音說： \end{tabularx} \\ \\ \relax
7:32 & \begin{tabularx}{0.7\textwidth}{X} 『我是你列祖的神，就是亞伯拉罕、以撒、雅各的神。』摩西戰戰兢兢，不敢觀看。 \end{tabularx} \\ \\ \relax
7:33 & \begin{tabularx}{0.7\textwidth}{X} 主對他說：『把你腳上的鞋脫下來，因為你所站的地方是聖地。 \end{tabularx} \\ \\ \relax
7:34 & \begin{tabularx}{0.7\textwidth}{X} 我的百姓在埃及所受的困苦，我確實看見了；他們悲嘆的聲音，我也聽見了。我下來要救他們。現在，你來，我要差你往埃及去。』 \end{tabularx} \\ \\ \relax
7:35 & \begin{tabularx}{0.7\textwidth}{X} 「這摩西就是有人曾棄絕他說『誰立你作我們的領袖和審判官』的，神卻藉那在荊棘中顯現的天使的手差派他作領袖，作解救者。 \end{tabularx} \\ \\ \relax
7:36 & \begin{tabularx}{0.7\textwidth}{X} 這人領以色列人出來，在埃及地，在紅海，在曠野的四十年間行了奇事神蹟。 \end{tabularx} \\ \\ \relax
7:37 & \begin{tabularx}{0.7\textwidth}{X} 這人是摩西，就是那曾對以色列人說『神要從你們弟兄中給你們興起一位先知像我』的。 \end{tabularx} \\ \\ \relax
7:38 & \begin{tabularx}{0.7\textwidth}{X} 這人是那曾在曠野的會眾中和西奈山上，與那對他說話的天使同在，又與我們祖宗同在的，他領受了活潑的聖言傳給我們。 \end{tabularx} \\ \\ \relax
7:39 & \begin{tabularx}{0.7\textwidth}{X} 我們的祖宗不肯聽從，反棄絕他，他們的心轉向埃及， \end{tabularx} \\ \\ \relax
7:40 & \begin{tabularx}{0.7\textwidth}{X} 對亞倫說：『你為我們造神明，在我們前面引路，因為領我們出埃及地的這個摩西，我們不知道他遭遇了甚麼事。』 \end{tabularx} \\ \\ \relax
7:41 & \begin{tabularx}{0.7\textwidth}{X} 那時，他們造了一個牛犢，又拿祭物獻給那像，為自己手所做的工作歡躍。 \end{tabularx} \\ \\ \relax
7:42 & \begin{tabularx}{0.7\textwidth}{X} 但是神轉臉不顧，任憑他們祭拜天上的日月星辰，正如先知書上所寫的：『以色列家啊，你們四十年間在曠野，何曾將犧牲和祭物獻給我？ \end{tabularx} \\ \\ \relax
7:43 & \begin{tabularx}{0.7\textwidth}{X} 你們抬著摩洛的帳幕和理番---你們神明的星，就是你們所造為要敬拜的像。因此，我要把你們遷到巴比倫外去。』 \end{tabularx} \\ \\ \relax
7:44 & \begin{tabularx}{0.7\textwidth}{X} 「我們的祖宗在曠野，有作證的會幕，是神吩咐摩西照著他所看見的樣式做的。 \end{tabularx} \\ \\ \relax
7:45 & \begin{tabularx}{0.7\textwidth}{X} 這帳幕，我們的祖宗同約書亞相繼承受了，當神在他們面前趕走外邦人的時候，他們把這帳幕搬進承受為業之地，直存到大衛的日子。 \end{tabularx} \\ \\ \relax
7:46 & \begin{tabularx}{0.7\textwidth}{X} 大衛在神面前蒙恩，祈求為雅各的家預備居所。 \end{tabularx} \\ \\ \relax
7:47 & \begin{tabularx}{0.7\textwidth}{X} 但卻是所羅門為神造成殿宇。 \end{tabularx} \\ \\ \relax
7:48 & \begin{tabularx}{0.7\textwidth}{X} 其實，至高者並不住人手所造的，就如先知所言： \end{tabularx} \\ \\ \relax
7:49 & \begin{tabularx}{0.7\textwidth}{X} 『主說：天是我的寶座，地是我的腳凳。你們要為我造怎樣的殿宇？哪裡是我安歇的地方呢？ \end{tabularx} \\ \\ \relax
7:50 & \begin{tabularx}{0.7\textwidth}{X} 這一切不都是我手所造的嗎？』 \end{tabularx} \\ \\ \relax
7:51 & \begin{tabularx}{0.7\textwidth}{X} 「你們這硬著頸項，心與耳未受割禮的人哪，時常抗拒聖靈！你們的祖宗怎樣，你們也怎樣。 \end{tabularx} \\ \\ \relax
7:52 & \begin{tabularx}{0.7\textwidth}{X} 先知中有哪一個不是受你們祖宗的迫害呢？他們把預先宣告那義者要來的人殺了。如今你們成了那義者的出賣者和兇手了。 \end{tabularx} \\ \\ \relax
7:53 & \begin{tabularx}{0.7\textwidth}{X} 你們領受了天使所傳佈的律法，竟不遵守。」 \end{tabularx} \\ \\ \relax
7:54 & \begin{tabularx}{0.7\textwidth}{X} 眾人聽見這些話，心中極其惱怒，向司提反咬牙切齒。 \end{tabularx} \\ \\ \relax
7:55 & \begin{tabularx}{0.7\textwidth}{X} 但司提反滿有聖靈，定睛望天，看見神的榮耀，又看見耶穌站在神的右邊， \end{tabularx} \\ \\ \relax
7:56 & \begin{tabularx}{0.7\textwidth}{X} 就說：「我看見天開了，人子站在神的右邊。」 \end{tabularx} \\ \\ \relax
7:57 & \begin{tabularx}{0.7\textwidth}{X} 眾人大聲喊叫，摀著耳朵，齊心衝向他， \end{tabularx} \\ \\ \relax
7:58 & \begin{tabularx}{0.7\textwidth}{X} 把他推到城外，用石頭打他。作見證的人把他們的衣裳放在一個名叫掃羅的青年腳前。 \end{tabularx} \\ \\ \relax
7:59 & \begin{tabularx}{0.7\textwidth}{X} 他們正用石頭打司提反的時候，他呼求說：「主耶穌啊，求你接納我的靈魂！」 \end{tabularx} \\ \\ \relax
7:60 & \begin{tabularx}{0.7\textwidth}{X} 然後他跪下來，大聲喊著：「主啊，不要將這罪歸於他們！」說了這話，就長眠了。 \end{tabularx} \\ \\
[1ex]
\hline
\hline
\end{longtable}
$^{1}$OK 聽到啦 真係唔好意思 一出場就已經有咪ge 問題.
Anyway 今日我想講ge 題目係史上最撩交打ge 字變.
我發覺今次真係出好多問題 撩字我打錯jor .
錯別字打jor 撩字 應該係剔手邊多謝Phoebe提醒我.
之後我改唔切powerpoint 大家明白.
係美國內D 中文開始變得唔好 就寫錯字.
Anyway 今日我ge 題目係史上最撩交打ge 字變.
而經文就係《史陀行傳》第七章第二至六十節.
係好長 58節經文 但我地唔會睇哂ge  我只係睇重點.
呢段經文係《聖經》第一次 亦都係最後一次去記載司提反講ye .
點解係咁呢 就係因為呢段經文係司提反自我ge 辯護.
佢辯護完之後就點呀 就俾人用石頭打死jor  然後就冇然後.
咁所以呢個係唯一一次亦都係最後一次司提反係《聖經》裡面係有講ye .
所以其實有咩人會為自己辯護 辯護到講完ye 之後即刻俾人打死.
咁司提反係咪好犀利呢 辯護完之後即刻俾人打死.
所以我就會話佢ge 辯護係史上最撩交打ge 自辯.
我地首先會去睇下司提反究竟被指控D 乜.
然後睇下佢點回應 再睇下呢段對話點解咁撩交打.
同埋我地去睇下呢段ge 講話.
究竟係整卷ge 《史唐行傳》裡面ge 文學脈絡當中有D 咩ge 作用.
咁我地首先去睇下究竟司提反被人指控D 乜.
咁我地去睇返上文第六章ge 八至十四節gwo 度呢.
gwo 度講咩呀 gwo 度就話司提反就比一班ge 猶太人指控兩樣ye .
就係指控邊兩樣ye 呢 就係指控佢.
毀謗聖地同埋律法.
同埋宣揚gwo 個要拆毀聖殿同埋改變摩西律法ge 耶穌.
所以司提反係佢ge 自辯裡面就回應呢兩點.
首先司提反就去回應.
佢回應呢個抵毀聖殿ge 一個指控.
佢就係七章ge 四十四至五十節gwo 度強調兩點.
第一點就係佢話呢 神係曠野ge 時候係點呀.
神係曠野裡面唔係靠聖殿同人同在.
神係曠野裡面係靠見證ge 帳幕同以色列人同在.
呢個係第一點.
第二點係咩呢.
雖然所羅門建造聖殿.
但係司提反最終強調ge 係咩呀.
但係神唔係住係聖殿.
因為神並唔住係人手所做ge 殿宇當中.
咁所以司提反係呢度做緊D 咩呀.

$^{41}$就係當佢比人指控佢毀謗聖殿ge 時候.
佢就反指控gwo D 指控佢ge 人.
就話係你地理解錯jor 聖殿ge 功能.
神從來都唔係靠聖殿同人同在.
神係唔會住係人手所做ge 殿當中.
係你地理解錯聖殿ge 功能.
所以係咪好勁呀.
佢反串返gwo 班指控佢ge 人.
第二就係司提反佢比人指控咩呀.
毀謗律法.
咁所以佢就係七章ge 35到43節gwo 度就強調咩呀.
強調摩西領受jor 活ge 聖賢.
但係係咩呀.
係以色列人ge 先祖唔肯聽佢講.
仲要向偶像獻祭.
值得注意ge 係咩呢.
就係司提反強調呢.
以色列人兩次推開jor 摩西.
換言之呢.
司提反係呢度被人指控毀謗律法.
但係佢係點呀.
佢反指控返以色列人ge 先祖呢.
先係違背摩西律法gwo 個.
佢甚至係七章51到53節gwo 度呢.
去反指控gwo D 呢指控佢ge 人.
話佢地咩呀.
就話你地就好似gwo D 先祖一樣.
好似你地先祖一樣違背律法.
好似先祖殺害先知一樣.
殺害jor gwo 個二人耶穌.
咁呢聖殿同埋律法呢係當時呢猶太人呢最重視ge ye 啦係咪.
咁司提反係呢度呢.
唔單止呢冇辯護到自己呢.
冇毀謗聖殿同埋冇毀謗律法.
反而點呀.
反而反指控返gwo D 指控佢ge 人.
係你地錯誤理解jor 聖殿ge 功能.
同埋唔係我違背律法.
係你地同埋你地ge 先祖違背律法.
咁所以個結果係咩呀.

$^{81}$就即刻比人打死啦係咪.
所以佢就唔再講ye 冇機會再講ye 啦.
不過呢.
司提反呢呢段講話呢.
最難理解ge 地方呢.
其實並唔係點解佢咁撩交打.
而係點解司提反要咁長氣ge 呢.
點解呢個講話咁沉氣.
要用58節ge 經文ge 呢.
如果我地去留意ge 話呢.
司提反係被指控毀謗聖殿同埋律法.
但係如果我地留意返呢段經文呢.
58節ge 經文大部分都同毀謗聖殿同埋律法係無關.
大部分ge 內容係咩呀.
大部分ge 內容係司提反重述返以色列ge 歷史.
我地要問ge 問題就係.
如果司提反ge 講話.
淨係為jor 回應毀謗聖殿同埋律法ge 指控ge 話.
點解佢要用咁多篇幅去講返以色列ge 歷史呢.
要答呢條問題呢.
我地首先要去睇下.
司提反所講ge 以色列ge 歷史有D 乜ye 特別.
同埋睇下司提反所提及ge 人物.
有D 咩ye 特別.
首先我地睇下司提反呢.
佢講ge 歷史有D 咩特別.
司提反呢.
佢講ge 歷史呢.
大致上可以分為四個階段.
第一個階段呢就係亞伯拉罕ge 階段.
第二個階段呢就係約瑟ge 階段.
第三呢就係摩西ge 階段.
第四呢就係約舒亞帶領以色列人.
進入jor 迦蘭地之後ge 階段.
咁如果我地去睇ge 話呢.
你見到司提反最強調ge 階段係咩呀.
係未入迦蘭地之前.
佢就係用jor 僅僅大概六節ge 經文.
58節經文裡面用jor 六節ge 經文.
去講約舒亞帶領以色列人入迦蘭地.

$^{121}$同埋之後ge ye .
就係用jor 六節.
但係大部分ge 篇幅都係講緊.
以色列人係迦蘭地之前ge 事.
咁我地就要問一個問題.
亞伯拉罕約瑟同埋摩西ge 時代.
有D 乜ye 共通點呢.
亞伯拉罕約瑟同埋摩西ge 時代.
有咩共通點呢.
呢三個時代有D 乜ye 共通點.
佢地ge 共通點呢三個時代ge 共通點.
就係神係以色列地以外工作.
神係以色列地以外工作.
例如係亞伯拉罕ge 時代.
神係係邊度呼召亞伯拉罕.
神係係米索不達米亞呼召亞伯拉罕.
叫佢搬屋.
跟住約瑟ge 時代.
神係點架.
神係係埃及.
唔係係以色列地.
唔係係迦南地.
神係係埃及度同約瑟同在.
並且拯救佢脫離一切ge 患難.
然後摩西呢.
摩西係點呀.
摩西唔係係以色列地領受呼召.
摩西做ge ye 係係邊度.
摩西係係西奈山ge 曠野度被神呼召.
佢唔係係以色列地度被神呼召.
同埋係點呀.
摩西係係紅海同埋曠野.
以色列地以外度行神跡騎士.
去帶領以色列人係曠野.
咁所以無論係亞伯拉罕約瑟.
定係摩西.
佢地經歷緊ge ye 係係邊度經歷架.
佢地點樣經歷神.
佢地唔係係以色列地度經歷神.
佢地係係以色列地以外度經歷神.

$^{161}$所以呢斯提反ge 呢一個ge 歷史ge 論述.
點解佢要重複講以色列人ge 歷史.
佢想強調ge 就係.
原來神一直都係係以色列地以外工作.
相反我地去睇呢.
係呢個斯提反ge 論述裡面呢.
佢講到約瑟呀帶領jor 以色列人入迦南之後又點呀.
佢冇再提及神有任何工作.
一個好大ge 對比.
神係以色列地以外工作.
呢個係斯提反要強調ge ye .
咁呢神係以色列地以外工作ge 呢個主題呢.
其實緊扣住呢成個《使徒行傳》ge 劇情發展.
係斯提反ge 講話之前呢.
即係第一至六章《使徒行傳》.
神呢係一直淨係係耶路撒冷當中工作.
但係係斯提反講話之後係點呀.
斯提反講完ye 之後就比人打死jor .
然之後呢耶路撒冷教會就被迫害.
然之後點呀.
信徒就逃散到猶大ge 各地.
同埋撒瑪利雅各處.
而路加呢就從呢一刻開始呢就去描寫呢.
神點樣係猶大地以外工作.
係以色列地以外工作.
例如第八章呢就描述.
肥利係撒瑪利亞度傳道啦.
第九章呢就描述.
蘇羅呢係大馬士革度被呼召等等等等.
換言之呢斯提反ge 講話呢.
其實原來唔單止係回應緊呢.
毀謗聖殿同埋律法ge 指控.
而係作為呢劇情上面ge 轉捩點.
去標誌緊呢神係猶大地.
以色列地以外工作.
今日呢我地會唔會.
今日呢我地唔會再好似以前咁樣呢.
即係斯提反ge 時代咁樣去爭辯緊呢.
聖殿同埋律法ge 功能.
因為我地根本冇聖殿而家係咪.

$^{201}$但係呢斯提反ge 講話呢仍然值得我地去反思.
我地會唔會用緊各種ge 框框.
去規限緊神ge 同在同埋工作ge 邊界呢.
我地ge 信仰會唔會淨係停留係.
星期日我地崇拜ge 呢一個地方.
UVge 呢一個地方呢.
斯提反ge 講話呢挑戰緊我地.
唔好淨係停留係教會ge 四面牆裡面.
而係要係教會以外呢一個地方以外.
ge 任何場景去實踐同埋見證神ge 工作.
咁我地明白jor 呢斯提反點解要重述以色列ge 歷史.
同埋呢一段歷史有咩特別之後呢.
我地而家就要去睇下呢佢提及ge 人物有D 咩特別.
我想問一個問題.
點解斯提反要提.
特別要提亞伯拉罕約瑟同埋摩西呢.
呢三個人.
即係呢三個人物有D 咩共通點.
亞伯拉罕約瑟同埋摩西.
有咩共通點.
呢三個人ge 共通點就係.
呢三個人ge 共通點就係呢.
佢地三個都係流浪異鄉ge 人.
而摩西甚至係咩呢.
摩西甚至係犯jor 法.
然後要走佬ge 逃犯.
係逃犯黎ge 摩西.
唔好唔記得呢個身份.
摩西係一個犯jor 法.
搞到要走佬ge 逃犯.
我想講ge 係D 咩呢.
就係呢史唐人傳呢.
我地好常呢通常呢D 傳道人呀.
牧者呢或者牧師呢.
講史唐人傳就係講全副音架嘛.
我地好少留意呢其他ye .
但我今日好想從呢個角度去切入.
其實史唐人傳唔係淨係講全副音.
流亡呢其實係路加福音.
或者唔係路加福音sorry.

$^{241}$係路加著作.
即係包埋路加福音同埋史唐人傳.
流亡呢個主題呢其實係路加著作裡面呢.
重複出現ge 主題.
不過教會好少講.
我淨係睇到全副音.
但我地冇睇到呢福音係點傳.
同埋全副音gwo 班人ge 身份係乜ye 人.
但今日我想強調返D 點.
我要強調ge 係流亡呢.
其實係路加著作重複出現ge 主題.
史唐人傳呢成日都講到呢.
基督徒要逃亡.
例如係呢段經文裡面已經講jor 係咩.
屍體犯死jor 之後就係點呀.
信徒被逼迫然後就係點呀.
走佬呀.
要流亡要逃亡到猶大同埋薩馬利亞國處.
然後去到十一章再強調返一次ge 一點.
就係十一章重複返一次呢就係屍體犯被殺.
所以呢D 信徒要走佬.
去到唔同ge 地方同埋去到安提奧.
走佬係路加好強調ge ye .
仲有咩呀.
仲有保羅呢.
保羅好多次呢係史唐人傳裡面係點呀.
唔單止係傳福音呀.
通常傳完福音之後要點呀.
傳完之後就要走佬架啦.
路加成日去強調呢.
幾次強調保羅點樣走佬.
幾次去強調保羅點樣逃亡.
保羅係一個流亡ge 人黎架.
係史唐人傳裡面.
除jor 呢主動地逃亡之外呢.
路加亦都好強調呢.
被動地被逐出城外呀.
如果我地去睇返路加福音第四章.
路加福音第四章呢.
耶穌ge 第一篇講到.

$^{281}$佢係自己ge 家鄉拿沙那道講完道之後就點呀.
gwo 班聽眾就好嬲.
點呀.
將佢逐出城外.
所以係路加著作裡面.
第一個被逐出城外ge 人呢.
並唔係門徒唔係使徒.
第一個被逐出城外ge 人係邊個係耶穌.
佢講第一篇道就俾人逐出城外.
係自己ge 家鄉拿沙那道.
俾人逐出城外.
係史唐人傳ge 十三章五十節gwo 度講到.
保羅同埋巴拿巴.
係比西底ge 安提阿道傳道之後點呀.
傳道之後就俾人逐出城外.
係史唐人傳ge 十八章一至四節gwo 度講到兩公婆.
就係呢亞居拉同埋伯基拉.
呢兩公婆有咩特別呀.
就係係羅馬皇帝格魯丟ge 命令底下.
將所有猶太人逐出城外.
包括呢兩公婆亞居拉同埋伯基拉呢兩公婆.
係咁樣俾人逐出羅馬城.
所以我地會見到其實路加著作裡面呢.
係有好多逃亡同埋被逐出城外ge 人物.
包括耶穌基督都係.
點解路加要咁樣寫呢.
我地點解留意唔到.
點解路加一定要咁樣強調呢.
就係因為路加好想透過斯提反ge 講話.
去重新去建構緊基督徒ge 身份.
基督徒ge 身份係點架.
就好似亞伯拉罕,約瑟同埋摩西一樣.
無論係耶穌定係跟隨佢ge 人.
其實都係城外ge 逃亡者黎.
呢個係路加要強調ge 基督徒身份.
我地好少留意呢D ye .
但今日好想同大家去睇返.
路加其實好強調基督徒ge 身份最初係點.
無論係耶穌定係跟隨佢ge 人.
都係城外ge 流亡者.

$^{321}$路加用呢種形象去形容耶穌.
同埋形容早期ge 信徒.
係因為基督徒係身處猶太權力ge 中心.
同埋羅馬權力ge 中心ge 邊緣ge 邊緣人.
佢地一直係逃亡緊.
一直係被逐出城外.
今日作為香港人.
我地以前係香港國際大都會裡面.
我地唔係好明流亡係D 咩.
我地有一個咁繁榮ge 大都會.
但近幾年作為香港人.
我地經歷緊呢幾年ge 社會變化ge 時候.
我地可能更深去體會緊.
乜ye 係流亡.
有D 人被迫要遠走他鄉.
從此永遠冇機會返返香港.
有D 人可能係想走但走唔到.
坐緊監.
就算一D 唔走ge 人.
留係香港ge 人.
都被迫去面對緊一個新ge 香港.
去體驗緊乜ye 係在地流亡.
在地流亡即係我冇走到.
但係我屋企變jor 樣.
呢個叫做在地流亡.
我想講ge 係.
盧家著作唔係就係講全服癮.
呢一段經文.
同緊每一個.
淨係係度經歷緊流亡ge 人講.
無論我地身處乜ye ge 地方.
無論我地走佬.
被迫逃亡jor 去咩地方.
神仍然都係度.
面對住香港ge 前景.
話係香港教會未來要面對.
更大ge 迫迫.
但係呢段經文安慰緊我地.
教會最初.
就係一個流亡ge 群體.

$^{361}$但係神ge 工作從來冇停過.
就好似史太人傳依度所描述ge 一樣.
就算信徒不斷ge 要逃亡.
就算信徒ge 不斷地被驅逐出城.
神ge 道仍然能夠大大得勝.
甚至去到羅馬.
羅馬呢一個希臘羅馬世界ge 中心.
今日唔知我地係咪面對緊流亡ge 處境.
究竟神同我地講緊D 咩呢.
另一方面呢.
我亦都好想用呢段經文.
去為我地每一個係美國ge 信徒帶黎一D 反思.
可能有D 人我地冇經歷香港近呢幾年發生ge 事.
因為我地好早已經黎到美國.
呢段經文對於身處美國ge 信徒黎講.
可能係一個挑戰.
因為路加著作提醒緊我地.
我地所追隨ge 基督.
佢係一個被逐出城外流亡ge 基督.
而早期ge 信徒亦都曾經流亡過.
今日我地係美國.
我地會唔會已經忘記jor 呢一段教會ge 歷史.
忘記jor 教會係曾經身處係權力中心ge 邊緣.
我地今日美國ge 教會會唔會忘記jor 我地ge 過去.
忘記jor 我地曾經係一個流亡ge 群體.
反過黎要去追逐權力.
成為權力ge 中心呢.
我覺得呢一點好值得我地美國ge 教會去反思.
我地會唔會已經從被逐出城外ge 人.
變成逐人出城外ge 權力中心.
我地仲係唔係一個能夠與社會ge 被邊緣ge 人.
同埋受欺壓ge 人同行ge 教會呢.
今個月係佔中ge 十周年.
當我去預備呢篇講章ge 時候呢.
我諗起當年有一D 教會願意去躺開佢地ge 大門.
去比示威者入去休息.
同時亦都有一D 教會呢.
關埋門.
連廁所都唔比人用.
今日我地係美國ge 處境當中.

$^{401}$我地要成為邊一種ge 教會呢.
我地係成為站係權貴gwo 邊ge 教會.
定係同社會被欺壓ge 人同行ge 教會呢.
我覺得呢點係美國ge 福音派教會呢.
係好值得反思.
我舉一個例子.
當我地去講到以巴戰爭ge 時候.
好多好多美國福音派ge 教會.
會覺得巴勒斯坦係邪惡.
以色列永遠係正義.
我地咁樣係做緊D 咩呀.
我地唔記得jor 原來以色列係一個權力ge 核心.
我地咁樣宣告ge 時候.
其實我地係站係權貴gwo 邊.
但當你咁樣講ge 時候.
可能就有好多福音派教會起你毀公你.
咁即係你唔支持以色列.
我想講ge 係咩.
我想講ge 係今日我地有冇睇到權力係邊度.
我地係企係權力gwo 邊定係企係弱勢gwo 邊.
我唔係話巴勒斯坦完全冇錯.
只不過我地成日係我地ge 保守信仰裡面.
我地連權力ge 架構都睇唔到.
我地自己企係權貴gwo 邊我地自己都唔知.
我地已經唔再睇到巴勒斯坦有幾多難民.
睇唔到以色列係點樣轟炸巴勒斯坦ge 城市.
有幾多人死jor 我地睇唔到.
係因為我地ge 眼睛被蒙蔽.
係因為我地忘記jor 我地自己ge 身份.
原來教會本身都係一班難民.
我地ge 開始係由一個被逐出城外流亡ge 基督開始.
而我地ge 教會ge 開始係一班被迫迫ge 信徒.
一班走佬ge 信徒.
一班逃亡ge 信徒.
一班被逐出城外ge 信徒建立.
我地唔記得jor 我地ge 位置.
福音派ge 教會成為jor 可能成為jor 一個追求權力ge 中心.
所以我地都企係jor 權力gwo 邊.
但今日我地願意成為一間咩ge 教會.
係咪一間能夠同弱勢同流離失所ge 人同行ge 教會呢.

$^{441}$願意呢屍體犯今日ge 講話去提醒我地.
我地每一個都係被逐出城外ge 流亡者.
約我地一齊去禱告.
呢一刻呢好想大家有一D 時間呢去禱告.
特別好想平時教會好少做.
但係好想呢一刻我地去紀念一D 呢.
係世界上面經歷緊戰禍.
無論係香港啦無論係巴勒斯坦啦.
無論世界其他ge 地方.
其實呢個世界有好多流亡緊ge 人.
有好多難民.
有好多被迫要流離失所.
去到唔同ge 地方要走佬ge 人.
好想去慰著呢D ye 去禱告.
你首先呢去紀念呢D ge 人.
我地好想呢段時間呢有少少ge 時間.
去紀念世界上面ge 流亡者.
無論係香港人巴勒斯坦人.
定係其他地方被迫要流亡ge 人.
我地去求主去紀念呢D ge 群體.
係流離失所當中去經歷緊上帝ge 恩典.
(安祈).
係呀主.
今日呢我地ge 成妹教會ge 會眾.
黎到你ge 面前.
係呀主今日呢你再一次呢.
透過呢一段ge 經文去提醒我地.
我地每一個人都係被逐出城外ge 流亡者.
所以我地能夠去認同.
其他係世界各地流亡緊ge 人.
主今日呢求你再一次去俾我地睇返.
我地自己ge 身份.
讓我地作為基督ge 信徒.
唔係以成功自居.
唔係以追求權力財力去自居.
主特別祝福我地.
無論係成妹教會定係美國ge 宗教會.
主你約我地唔好去成為權力ge 核心.
主你俾我地再一次ge .
係能夠同雞蛋去同行.

$^{481}$而唔係企係高牆gwo 邊.
主今日呢你去開我地眼睛.
將一切蒙蔽緊我地思想.
蒙蔽我地睇到呢被欺壓ge 人ge 需要ge .
呢D ye 全部都lor 走.
主今日你恢復返我地教會ge 身份.
同埋使命.
因為教會由一開始.
本身就係一個流亡者ge 群體.
亦都係同被欺壓ge 人同行ge 群體.
主今日願意你.
再一次去恢復美國教會ge 身份同埋使命.
讓我地唔單止係傳福音實踐大使命.
而係同緊社會ge 弱勢者同行.
去宣揚緊上帝ge 公義.
主你就咁樣去祝福更新我地ge 教會.
多謝你咁樣去順聽我地ge 禱告.
去奉主耶穌基督ge 聖名祈求.
阿們.
多謝阿邦ge 提醒.
值得我地去思想ge 一個訊息.
咁我希望你等間離開jor .
咁你唔好留學jor .
等間完jor 午餐ge 時候.
返返黎呢個ge 嘉事大會.
因為嘉事大會呢有D 好重要ge update.
比大家.
包括你教會裡面.
我地ge 財務上面ge 虧缺.
同埋我地童工最近一D ge update.
我地事工裡面ge update.
包括小弟自己一D ge update.
不過我都未走住ge .
但我希望你黎到呢個嘉事大會.
關心我地教會.
點樣係呢個世界裡面.
繼續去到進行同埋發展上帝ge 事工.
等間嘉事大會見到大家.
願上帝祝福各位.
拜拜!.

\newpage



\section{啟示錄 2:1-3:22}
\label{sec:1ry__RW4dIE}
\textbf{基督眼中的教會 \#3 (啟示錄 2\_1-3\_22) - 陳崇基牧師}
\newline
\newline
連結: \href{https://youtube.com/watch?v=1ry_-RW4dIE}{\texttt{https://youtube.com/watch?v=1ry\_-RW4dIE}} ~~~~ 語音日期: 2024-10-02
\newline
\newline
\hyperref[sec:9bZktUw13B0]{\small{< < < PREV SERMON < < <}}
~
\hyperref[sec:index_chronic]{\small{[返順時目]}}
~
\hyperref[sec:index_scriptual]{\small{[返順卷目]}}
~
\hyperref[sec:yfuGrCAshKk]{\small{> > > NEXT SERMON > > >}}
\newline
\newline
啟示錄 2:1-3:22
\newline
\begin{longtable}{cl}
\hline
\hline
章節 & 經文 (和合本修訂版)\\
\hline
2:1 & \begin{tabularx}{0.7\textwidth}{X} 「你要寫信給以弗所教會的使者，說：『那右手拿著七顆星，在七個金燈臺中間行走的這樣說： \end{tabularx} \\ \\ \relax
2:2 & \begin{tabularx}{0.7\textwidth}{X} 我知道你的行為、勞碌、忍耐，也知道你不容忍惡人。你也曾察驗那自稱為使徒卻不是使徒的，看出他們是假的。 \end{tabularx} \\ \\ \relax
2:3 & \begin{tabularx}{0.7\textwidth}{X} 你能忍耐，曾為我的名勞苦而不困倦。 \end{tabularx} \\ \\ \relax
2:4 & \begin{tabularx}{0.7\textwidth}{X} 然而，有一件事我要責備你，就是你把起初的愛心拋棄了。 \end{tabularx} \\ \\ \relax
2:5 & \begin{tabularx}{0.7\textwidth}{X} 所以你要回想你是從哪裡墜落的，並且要悔改，做起初所做的工作。你若不悔改，我要到你那裡去，把你的燈臺從原處挪去。 \end{tabularx} \\ \\ \relax
2:6 & \begin{tabularx}{0.7\textwidth}{X} 然而你還有一件可取的事，就是你恨惡尼哥拉派的行為，這種行為也是我所恨惡的。 \end{tabularx} \\ \\ \relax
2:7 & \begin{tabularx}{0.7\textwidth}{X} 凡有耳朵的都應當聽聖靈向眾教會所說的話。得勝的，我必將神樂園中生命樹的果子賜給他吃。』」 \end{tabularx} \\ \\ \relax
2:8 & \begin{tabularx}{0.7\textwidth}{X} 「你要寫信給士每拿教會的使者，說：『那首先的、末後的，死過又活了的這樣說： \end{tabularx} \\ \\ \relax
2:9 & \begin{tabularx}{0.7\textwidth}{X} 我知道你的患難和貧窮—其實你卻是富足的，也知道那自稱是猶太人的所說毀謗的話，其實他們不是猶太人，而是撒但會堂的人。 \end{tabularx} \\ \\ \relax
2:10 & \begin{tabularx}{0.7\textwidth}{X} 你將要受的苦，你不用怕。看哪！魔鬼要把你們中間幾個人下在監裡，使你們受考驗，你們要遭受苦難十日。你務要至死忠心，我就賜給你那生命的冠冕。 \end{tabularx} \\ \\ \relax
2:11 & \begin{tabularx}{0.7\textwidth}{X} 凡有耳朵的都應當聽聖靈向眾教會所說的話。得勝的必不受第二次死的害。』」 \end{tabularx} \\ \\ \relax
2:12 & \begin{tabularx}{0.7\textwidth}{X} 「你要寫信給別迦摩教會的使者，說：『那有兩刃利劍的這樣說： \end{tabularx} \\ \\ \relax
2:13 & \begin{tabularx}{0.7\textwidth}{X} 我知道你的居所，就是有撒但座位之處；當我忠心的見證人安提帕在你們中間，在撒但所住的地方被殺之時，你還堅守我的名，沒有否認對我的信仰。 \end{tabularx} \\ \\ \relax
2:14 & \begin{tabularx}{0.7\textwidth}{X} 然而，有幾件事我要責備你，就是在你那裡有人服從了巴蘭的教訓；這巴蘭曾教唆巴勒將絆腳石放在以色列人面前，使他們吃祭過偶像之物，並且犯淫亂。 \end{tabularx} \\ \\ \relax
2:15 & \begin{tabularx}{0.7\textwidth}{X} 同樣，你那裡也有人服從了尼哥拉派的教訓。 \end{tabularx} \\ \\ \relax
2:16 & \begin{tabularx}{0.7\textwidth}{X} 所以，你當悔改；若不悔改，我很快就到你那裡來，用我口中的劍攻擊他們。 \end{tabularx} \\ \\ \relax
2:17 & \begin{tabularx}{0.7\textwidth}{X} 凡有耳朵的都應當聽聖靈向眾教會所說的話。得勝的，我必將那隱藏的嗎哪賜給他，並賜他一塊白石，石上寫著新的名字，除了那領受的以外，沒有人認識。』」 \end{tabularx} \\ \\ \relax
2:18 & \begin{tabularx}{0.7\textwidth}{X} 「你要寫信給推雅推喇教會的使者，說：『神的兒子，那位眼睛如火焰、雙腳像發亮的銅的這樣說： \end{tabularx} \\ \\ \relax
2:19 & \begin{tabularx}{0.7\textwidth}{X} 我知道你的行為：愛心、信心、勤勞、忍耐；又知道你末後所行的善事比起初所行的更多。 \end{tabularx} \\ \\ \relax
2:20 & \begin{tabularx}{0.7\textwidth}{X} 然而，有一件事我要責備你，就是你容忍那自稱是先知的婦人耶洗別教唆我的僕人，引誘他們犯淫亂，吃祭過偶像之物。 \end{tabularx} \\ \\ \relax
2:21 & \begin{tabularx}{0.7\textwidth}{X} 我曾給她悔改的機會，她卻不肯悔改她的淫行。 \end{tabularx} \\ \\ \relax
2:22 & \begin{tabularx}{0.7\textwidth}{X} 看吧，我要使她病倒在床上。那些與她犯姦淫的人若不悔改他們的行為，我也要使他們同受大患難。 \end{tabularx} \\ \\ \relax
2:23 & \begin{tabularx}{0.7\textwidth}{X} 我又要殺死她的兒女，眾教會就知道，我是那察看人肺腑心腸的，我要照你們的行為報應各人。 \end{tabularx} \\ \\ \relax
2:24 & \begin{tabularx}{0.7\textwidth}{X} 至於你們其餘的推雅推喇人，就是一切不隨從這教訓，不明白他們所謂撒但深奧之理的人，我告訴你們，我不會再把別的擔子放在你們身上。 \end{tabularx} \\ \\ \relax
2:25 & \begin{tabularx}{0.7\textwidth}{X} 你們只要持守那已經有的，直到我來。 \end{tabularx} \\ \\ \relax
2:26 & \begin{tabularx}{0.7\textwidth}{X} 那得勝又遵守我命令到底的，我要賜給他權柄制伏列國； \end{tabularx} \\ \\ \relax
2:27 & \begin{tabularx}{0.7\textwidth}{X} 他必用鐵杖管轄他們，如同打碎陶器， \end{tabularx} \\ \\ \relax
2:28 & \begin{tabularx}{0.7\textwidth}{X} 像我也從我父領受了權柄一樣。我又要把晨星賜給他。 \end{tabularx} \\ \\ \relax
2:29 & \begin{tabularx}{0.7\textwidth}{X} 凡有耳朵的都應當聽聖靈向眾教會所說的話。』」 \end{tabularx} \\ \\ \relax
3:1 & \begin{tabularx}{0.7\textwidth}{X} 「你要寫信給撒狄教會的使者，說：『那有神的七靈和七顆星的這樣說：我知道你的行為，就是名義上你是活的，實際上你是死的。 \end{tabularx} \\ \\ \relax
3:2 & \begin{tabularx}{0.7\textwidth}{X} 你要警醒，堅固那些剩下、快要死的，因為我發現你的行為，在我神面前沒有一樣是完全的。 \end{tabularx} \\ \\ \relax
3:3 & \begin{tabularx}{0.7\textwidth}{X} 所以，要記得你所領受和聽見的；要遵守，並要悔改。你若不警醒，我必如賊一樣來到；我幾時來到你那裡，你絕不會知道。 \end{tabularx} \\ \\ \relax
3:4 & \begin{tabularx}{0.7\textwidth}{X} 然而，在撒狄你還有幾位是未曾污穢自己衣服的，他們會穿白衣與我同行，因為他們是配穿的。 \end{tabularx} \\ \\ \relax
3:5 & \begin{tabularx}{0.7\textwidth}{X} 得勝的必這樣穿白衣，我也不從生命冊上塗去他的名；我要在我父面前，和我父的眾使者面前，宣認他的名。 \end{tabularx} \\ \\ \relax
3:6 & \begin{tabularx}{0.7\textwidth}{X} 凡有耳朵的都應當聽聖靈向眾教會所說的話。』」 \end{tabularx} \\ \\ \relax
3:7 & \begin{tabularx}{0.7\textwidth}{X} 「你要寫信給非拉鐵非教會的使者，說：『那神聖、真實的，拿著大衛的鑰匙，開了就沒有人能關，關了就沒有人能開的這樣說： \end{tabularx} \\ \\ \relax
3:8 & \begin{tabularx}{0.7\textwidth}{X} 我知道你的行為。看哪，我在你面前給你一個敞開的門，是沒有人能關的。我知道你有一點力量，也遵守我的道，沒有否認我的名。 \end{tabularx} \\ \\ \relax
3:9 & \begin{tabularx}{0.7\textwidth}{X} 那屬撒但會堂的，自稱是猶太人，其實不是猶太人，而是說謊話的，我要使他們來到你腳前下拜，使他們知道我已經愛你了。 \end{tabularx} \\ \\ \relax
3:10 & \begin{tabularx}{0.7\textwidth}{X} 因為你遵守了我堅忍的道，我也必在普天下人受試煉的時候保守你免受試煉。 \end{tabularx} \\ \\ \relax
3:11 & \begin{tabularx}{0.7\textwidth}{X} 我必快來，你要持守你所有的，免得人奪去你的冠冕。 \end{tabularx} \\ \\ \relax
3:12 & \begin{tabularx}{0.7\textwidth}{X} 得勝的，我要使他在我神的殿中作柱子，他必不再從那裡出去。我又要把我神的名和我神城的名—從天上我神那裡降下來的新耶路撒冷，和我的新名，都寫在他上面。 \end{tabularx} \\ \\ \relax
3:13 & \begin{tabularx}{0.7\textwidth}{X} 凡有耳朵的都應當聽聖靈向眾教會所說的話。』」 \end{tabularx} \\ \\ \relax
3:14 & \begin{tabularx}{0.7\textwidth}{X} 「你要寫信給老底嘉教會的使者，說：『那位阿們、誠信真實的見證者、神創造的根源這樣說： \end{tabularx} \\ \\ \relax
3:15 & \begin{tabularx}{0.7\textwidth}{X} 我知道你的行為，你也不冷也不熱；我巴不得你或冷或熱。 \end{tabularx} \\ \\ \relax
3:16 & \begin{tabularx}{0.7\textwidth}{X} 既然你如溫水，也不冷也不熱，我要從我口中把你吐出去。 \end{tabularx} \\ \\ \relax
3:17 & \begin{tabularx}{0.7\textwidth}{X} 你說：我是富足的，已經發了財，一樣都不缺，卻不知道你是困苦、可憐、貧窮、瞎眼、赤身的。 \end{tabularx} \\ \\ \relax
3:18 & \begin{tabularx}{0.7\textwidth}{X} 我勸你向我買從火中鍛鍊出來的金子，使你富足；又買白衣穿上，使你赤身的羞恥不露出來；又買眼藥抹你的眼睛，使你能看見。 \end{tabularx} \\ \\ \relax
3:19 & \begin{tabularx}{0.7\textwidth}{X} 凡我所疼愛的，我就責備管教。所以，你要發熱心，也要悔改。 \end{tabularx} \\ \\ \relax
3:20 & \begin{tabularx}{0.7\textwidth}{X} 看哪，我站在門外叩門，若有聽見我聲音而開門的，我要進到他那裡去，我與他，他與我一起吃飯。 \end{tabularx} \\ \\ \relax
3:21 & \begin{tabularx}{0.7\textwidth}{X} 得勝的，我要賜他在我寶座上與我同坐，就如我得了勝，在我父的寶座上與他同坐一般。 \end{tabularx} \\ \\ \relax
3:22 & \begin{tabularx}{0.7\textwidth}{X} 凡有耳朵的都應當聽聖靈向眾教會所說的話。』」 \end{tabularx} \\ \\
[1ex]
\hline
\hline
\end{longtable}
$^{1}$很高興見到大家.
我知道今天下雨.
大家都能夠平安來到.
我們一起看聖經.
我們回到啟示錄第二章.
我們一直說開這個系列.
就是看啟示錄七個教會的書信.
耶穌寫了七封信.
或者發了七封信.
給當時第一世紀的教會.
當時他們在不同的情況裡面.
在不同的挑戰裡面.
這七封信講到他怎麼看教會.
我相信這七封信對我們教會有很多提醒.
基督怎麼看我們教會.
基督寫信到你家裡.
或者寫信給你.
怎麼看你作為一個基督徒.
或者一個基督的群體.
我們首先.
之前兩次我們看了頭兩封書信.
我們說這七個教會是順著一個遊弟的路線.
遊弟一個的教會是以弗所.
然後去到第七個教會叫做路第加.
我們看了以弗所和史梅特.
這次我們看第幾個舊的.
第三個城市就是別加摩這個地方.
別加摩.
這個是以前的一個古城.
是一個很大發展的古城.
因為你看它的劇場.
你看一個城市有沒有規模.
你就知道它的劇場有多大.
有二萬人的劇場.
因為你去一個城市.
美國也是.
有哪個stadium在那裡.
哪個football team等等.
有一個大的stadium就代表那個城市.
是當時數一數二的一個城市.

$^{41}$這個城市別加摩是在今天土耳其西岸的地方.
比較多人一點.
如果你這樣統計的話.
我們這樣去嘗試估計當時有多少人.
這個城市有15萬人.
可以15萬人.
對當時來說是一個很大很大的一個城市.
一個很大發展的城市.
裡面就有一些教會.
有一些初期的教會.
初期的信徒聚集的地方.
很多時候他們在家庭裡面聚集的.
不像我們今天有教堂.
或者有一個facility我們能夠去到.
他們在不同人的家裡面聚會.
基督就寫一封書信給這個教會.
我們說寫過給伊斯特索.
也寫過給史姆納.
現在寫給第三個城市叫做Pergamon.
這個別加摩教會.
第一節這裡這麼說.
12節.
他說要寫信給別加摩教會的使者.
說那有兩刃利劍的.
這樣說.
我們說過每一封書信開始的時候都會有一個自稱的.
基督怎麼自稱.
基督怎麼自稱也是跟當時的處境很有關係.
什麼叫兩刃利劍.
這些是當時羅馬帝國.
他們最大的其中一個的徽號.
或者最大的一個記號.
是代表他們政治的.
軍事的力量.
類似一把這樣的兩刃的利劍.
大概18至20多寸長.
這個是當時羅馬帝國一個最大的標誌.
就是說他們是話事的.
他們有權去審理或者有權去管理統治你或者統戰你.
這是一個政治的徽號.

$^{81}$政治的一個象徵.
我知道我們美國也快大選了.
一說到政治的時候我們有很多東西要說.
有很多東西要說或者令我們很生氣.
為什麼呢.
因為地上最大的權柄就是政權.
你覺得你家裡是不是最大.
可能是最大.
但是社會裡面地上最大的權柄最大的力量就是政治.
所以為什麼一說起這件事我們就很緊張.
而且每一個的政權都希望我們能夠對他效忠.
只能夠向他效忠.
你不說左還是右還是中.
哪一個政黨都想你向他效忠.
在當時教會當時的第一世紀都是一樣.
政治都是一個很重要的社會因素.
當時的羅馬王甚至是本地的特首或者官.
或者總督等等.
他們都是拿著一把這樣的刀.
一把這樣的兩刃的利劍.
代表他們能夠去懲罰能夠去統治.
能夠代表朝廷代表中央代表他們最大的羅馬帝國的政權.
你應該記得當耶穌被審的時候.
你記得耶穌被審.
他去過不同地方被審.
被公會被審或者被猶太人去審.
然後被羅馬人去審.
你記得他見彼拉多的時候他講過什麼.
你記得嗎.
第十九章約翰福音.
他見彼拉多的時候.
因為耶穌不說話.
他沒有任何的對答.
彼拉多就問耶穌.
他說你是不是不說話.
為什麼你一聲不出.
難道你不知道我有權柄釋放你.
也有權柄將你釘在十字架上.
耶穌怎麼回答他.
耶穌說若不是從上頭賜給你這個權柄.

$^{121}$你是毫無權柄辦我的.
所以把我交給你的人.
他的罪更大.
耶穌居然講了一句這樣的說話.
一個將會被殺的羔羊.
一個手無寸鐵的人.
他說如果你想辦我的話.
你要知道.
背後有一個權柄是比你更大的.
那個是上帝的權柄.
不是上帝的權柄.
去把耶穌交給彼拉多的話.
彼拉多是絕對沒有這樣的機會.
沒有這樣的權柄.
沒有這樣的時間.
沒有這樣的地位.
來審他.
耶穌說.
上帝才是最大的權柄.
上帝才是最高最高的政權.
不是你.
不是你一個羅馬的地方官.
或者也不是.
羅馬帝國.
所以耶穌在介紹自己的時候.
他說哪有兩刃利劍的這件事.
耶穌說他才是有兩刃利劍的.
不是羅馬帝國.
不是地上的政權.
也不是我們美國的總統.
誰才是有兩刃利劍.
是死而復活的耶穌.
你想一下這是多大的宣告.
這不是一個普通人.
可以講得出的說話.
很多時候我們看.
特別是一些剛剛開始認識信仰的朋友.
或者開始認識聖經.
或者了解基督的信仰的時候.
我們要知道耶穌不是一個.

$^{161}$他講很多東西.
或者自稱很多的一些說話.
不是普通人能夠講得到的.
除非他是神經病.
或者他是一個很大的騙子.
或者可能他就是真正講真理.
他說他才是有兩刃利劍.
當然我們需要去理解地上發生的事.
我們需要理解我們今天的世代.
有什麼需要.
政治裡面有什麼風啊.
或者有什麼言談等等.
我們需要去分辨.
這些我們要認識的.
不過我們要知道.
在我們這些對話的背後.
只有一位是值得我們效忠的.
就是哪一位.
耶穌基督.
不是政黨.
不是地上的國家.
我不管中國也好 美國也好.
哪個國也好.
沒有一個國比天國更加大.
我經常覺得很有趣.
我們緊張的東西都是地上的東西.
偏偏天國的東西我們不緊張.
教會成不成長.
我不在乎.
但最重要的是在乎什麼.
我喜歡哪個能夠做到總統.
我喜歡哪個政黨.
走什麼的政策.
這個我最緊張.
我不是說不緊張我們需要關注.
不過最緊張的應該是什麼.
最緊張的是天國的發展.
天國有沒有在地上有權柄.
耶穌說他那位有兩刃利劍.
教會一個討耶穌喜悅的教會.

$^{201}$我們最緊張的東西.
最重要的東西應該是天國和福音.
對不對.
一個基督徒一個教會.
最應該關心的是上帝自己的天國在地上的權柄.
介紹了之後.
第十三節.
我知道你們別加摩的地方.
你們的居所.
就算有撒旦座位之處.
當我忠心的見證人安提伯在你們中間.
否定我的信仰.
沒有否認我的信仰.
這裡說得很有趣.
我知道你們的居所.
我知道你們住的地方是怎樣.
當然住的地方在哪裡.
就是別加摩.
當時的古城.
當時教會就存在在這個城市裡面.
這個當然是他們的物理地點.
就好像耶穌跟我們說.
我知道你的居所.
你就住在哪裡.
你就住在芝加哥.
你現在的城市香港等等不同的地方.
你的物理地點.
但這裡說的不是單單地上物理地點.
是說你的靈性地點.
他說你的靈性地點是甚麼.
是撒旦的座位.
撒旦的座位.
說的是甚麼.
在別加摩有一個很大的祭壇.
當時很出名.
有百多呎寬.
百多呎深.
一個很大的祭壇.
是給宙斯的.
當時的羅馬帝國的大神.

$^{241}$大的神明.
大的假神的神明.
就一個這麼大的.
宙斯的大祭壇.
現在你想看這個祭壇.
可以去德國看.
在柏林.
柏林有一個博物館叫做.
柏格曼博物館.
就是別加摩的博物館.
為甚麼呢.
因為早年的時候.
德國人就去了土耳其那裡挖掘東西.
然後挖到這些東西.
當時歐洲人最流行的.
挖到甚麼就搬回家.
所以他們逐塊石頭就搬回家.
然後搬回去.
放在德國柏林那裡.
然後建了一個博物館.
然後就去做展覽.
當時就找到這個叫做.
宙斯的神明.
給祭司的大祭壇.
當時別加摩的人.
他們就知道這個祭壇.
是一個最出名的祭壇.
耶穌就說這個祭壇.
就代表著你們所住的地方.
是一個撒旦的座位.
那些德國人將這些石頭搬回去.
然後再砌出來的時候.
當時是1930年的時候.
1930年的時候就開放給.
德國的民眾來參觀.
幾年之後.
德國納粹黨.
他們因為有很多的政治集會要做.
你記得希特勒.
不是很久之前.

$^{281}$是吧.
因為有很多政治的集會.
德國納粹黨有很多政治集會.
很多人聚集.
所以他們在不同的地方.
不同的城市.
他們就要去建一些大的廣場.
然後可以給政治集會用.
其中一個德國納粹黨的建築家.
Albert Speer.
他就說怎樣去建一個這樣的.
大的政治的廣場.
讓人聚集.
然後向這個政權去效忠.
他就去設計一個plaza的時候.
這個plaza在Lumenberg.
他就用一個這樣的靈感.
他說剛剛搬了回來.
這個宙斯的大祭壇.
他就用這個祭壇.
做一個model.
然後就做了一個大的廣場上.
一個大的樓梯.
一個大的建築物.
叫做Siblingville.
然後就將這個廣場.
變成一個納粹黨的.
效忠的一個廣場.
然後過幾年之後.
希特勒就在這個廣場上宣講.
他的政治言說.
這個就是Thorn of Satan.
就是撒旦的座位.
這裡講得很清楚.
他們住的地方是別家的.
但是耶穌提他們一件事.
他們住的地方其實就是撒旦的居所.
或者撒旦的皇位的地方.
就是說教會.
是需要知道他們住的地方.

$^{321}$不是單單他們的physical location.
而是他們的spiritual location.
我們需要look behind the scene.
要看到背後到底發生什麼事.
我們外表可以看到很宏偉的建築.
可以看到很宏偉的建設.
或者一些基建等等.
但是耶穌說你背後要看到的.
原來是有屬靈的勢力.
魔鬼的勢力在背後.
芝加哥漂亮嗎?.
芝加哥好天氣的時候很漂亮.
天氣好的時候是很漂亮.
但是你知道芝加哥.
特別是在美國浸會的network.
我們講到芝加哥的時候.
大家知道芝加哥是其中一個最難.
去到全福音去建立教會的地方.
每個人聽見芝加哥的時候.
就想起什麼呢?.
瑞多.
譚滬勁.
還有什麼好呢?.
Potthole夠多.
耶穌說你背後是有撒旦的勢力.
這個地方是其中一個最難最難.
去到建立教會.
或者最難去到全福音的地方.
別加摩也是一樣.
你想一下你在祭壇的旁邊.
你去建立教會.
那些人去拜假神.
那些人去上香奉獻.
去給假的神明.
你怎樣去傳耶穌基督的福音.
是非常困難.
事實上這封信也有講.
他說他們在當中.
有忠心的見證人叫做安提帕.
Antipas.

$^{361}$我們不知道這個人的底細是怎樣.
這句話說他是忠心的見證人.
安提帕在你中間.
可能這個人是當時教會的一個領袖.
很清楚他就是因為見證福音的緣故.
所以他就被殺.
殉道.
在第一世紀裡面.
教會遇見一些領袖.
殉道.
Being martyred.
其實是有很多這樣的例子.
你記不記得.
施洗約翰.
史提反.
雅各.
彼得.
保羅.
這些相繼都是在第一世紀.
逐個逐個被羅馬帝國迫害和殺害的人.
寫這封信的時候.
應該是主後的一世紀末的時候.
有另一個.
這裡就是講別加摩的其中一個領導.
安提帕.
他說他是faithful witness.
是一個忠心的見證人.
這句話應該是說.
安提帕曾經是被提堂.
然後在法庭上.
他去到被審判.
然後當時的羅馬的官.
就要他去到curse christ.
去到咒詛基督.
如果你不咒詛基督.
你就要人頭落地.
很明顯這個安提帕.
就是經過一個這樣的審訊.
然後被殺.
然後去信徒.

$^{401}$但耶穌說什麼呢.
他說我知道.
我知道你們居所是撒旦座位之處.
也知道你們忠心見證人安提帕在你們中間.
堅守我的命.
這個教會.
是一個耶穌所欣賞的教會.
在撒旦的座位裡面去建立教會.
在一個反對基督和不信基督的文化裡面.
去建立耶穌的群體.
是耶穌所欣賞的.
耶穌說I know.
他說我知道你們的難處.
我知道你們的suffering.
但我應許.
我會在你的苦難裡面拯救你.
不是拯救你脫離苦難.
而是拯救你經過這個苦難.
這個被耶穌所欣賞的教會.
是一個相信上帝權能的教會.
Otherwise 他們每個人都不會去教會.
在一個這麼困難的地方.
他們仍然堅信仍然去見證.
一個相信神權能的教會.
我想問我們是不是一個相信神權能的教會.
還是我們有任何困難.
有一點風吹草動的時候.
我們就收起來.
耶穌所欣賞的教會是一個.
相信祂權能的教會.
忠心的見證.
第二方面這裡說到.
他們教會是一個相信上帝豐盛的教會.
或者應該這麼說.
耶穌希望他們能夠相信.
上帝是有豐盛的公認的.
不需要怕.
為什麼這麼說呢.
因為這封信的下半節.
他說什麼.

$^{441}$他十四節.
他說有幾件事我們要責備你的.
就是你們中間裡面有人.
服從了巴蘭的教訓.
這巴蘭曾經唆擺巴拿.
將半腳石放在以色列人的面前.
史達文吃過祭過偶像之物.
並且犯淫亂.
同樣在你裡面也有人服從尼加尼嗎.
不是尼加拉吧.
尼哥拉派的教訓.
在這裡說什麼.
不熟悉聖經的時候.
我們就一頭霧水.
因為說到一個人叫巴蘭.
巴蘭在舊約裡面是一個非常出名的先知.
是一個假先知.
巴蘭這個人出名的是什麼呢.
你記得以色列人出了埃及.
然後走過曠野四十年.
然後他們快要去到迦林地的門口的時候.
他們曾經繞過一個地方.
繞過一個地方叫做摩納地.
摩納地有一個王叫巴納王.
巴納王.
巴納王因為這些以色列人.
要經過他的地方或者在他外圍經過的時候.
他就很不喜歡這些以色列人.
不喜歡以色列人.
他就請了一個假先知叫做巴蘭.
這個巴蘭在當時先知的江湖裡面是很出名的.
作法或者有很多的.
或者屬靈的撒旦的力量等等.
所以就叫了這個巴蘭.
請了他給錢請他.
然後請他做什麼呢.
就是去詛咒Israel.
就是去詛咒上帝的子民.
大家記得那個故事在民數記23,24章.
記載一個這樣的故事.

$^{481}$一個假的先知收了錢.
然後就去奉命去詛咒Israel.
當然那個故事最後是發生什麼事呢.
就是耶和華去阻截了巴蘭所做的事情.
你記得那個色說話的奴仔.
那個故事.
就是從這裡出來的.
他說這個教會.
別加摩教會.
有些人信奉了巴蘭的教導.
假先知的教導.
這個假先知曾經是唆擺巴勒王.
將絆腳石放在以色列人面前.
使得他們犯罪.
第一他們服從巴蘭的教導.
第二這個教會因為服從了巴蘭的教導.
所以吃了祭過偶像的食物.
然後也犯了淫亂.
這裡很奇怪.
到底背景是怎樣.
大概應該是這樣的.
當時的人在不同的城市.
因為有很多廟.
那些廟才可以買到祭肉.
或者你要加入那個association.
或者你要加入那個廟.
然後你才能吃到那些東西.
每逢過年過節.
初一十五的時候.
那些廟就會有大的筵席.
然後在大的筵席裡面.
他們有肉.
有一些平時比較難買到的食物等等.
他們就有一個大排筵席.
在廟裡面大排筵席.
但是你要吃的話.
你要加入他們的association.
等於今天你要進Costco買東西.
一定要加入Costco的會員.
當時就是這樣.

$^{521}$要加入一個廟的會員.
你才能夠進去裡面.
有些基督徒因為壓力的緣故.
有經濟壓力.
或者有不同的隱憂的緣故.
他們就一定要加入這些association.
然後他們才能夠去加入這些活動.
等於今天.
如果你說這裡比較少這樣的問題.
可能少一點.
希望少一點.
如果你要吃素的話.
你要去哪裡吃.
你要去佛廟吃.
你一定要加入佛廟的association.
才能夠加入那位台.
你喜歡吃素嗎.
OK吧.
但是你一定要加入那個temple.
你才能夠進去.
這個就是基督所責備.
教會去compromise的地方.
是為了一些benefit.
為了一些自己的好處.
無論是肚腹的好處.
或者一些關係上的好處.
商業關係上的好處.
他們就要放棄信仰.
放棄基督的信仰.
然後加入這個association.
他們才能夠有份去參與這些活動.
你在這裡說什麼.
他說有幾件事我要責備你的.
責備你就是你服從了那幫人的教訓.
什麼意思.
就是將信仰turn it around.
變成我們得到好處的方法.
以宗教之名.
以信仰之名.
或者不理信仰.

$^{561}$我們嘗試去拿一些benefit.
嘗試去拿一些好處.
這個就是耶穌不喜歡教會要做的.
如果教會說我們是信耶穌的話.
我們信耶穌為我們獨一的主.
我們就不能夠將他turn it around.
變成我們自己能夠得到benefit的方法.
大家記得幾個禮拜前.
我去了另外一個地方去講一個退修會.
我聽到一些大英姐妹跟我說.
她自己曾經去過一些不同城市的教會.
還有一些很大的教會.
她跟我說那些教會的牧師是有自己的飛機.
有自己的直升機去工作.
因為牧師很忙.
所以她就要搭飛機去.
她跟我說她那時候去的教會就是這樣.
然後教會裡面很混亂.
為什麼.
因為很多錢銀上的瓜葛.
特別是宗教領袖對錢銀的利用.
我喝咖啡的時候聽一聽.
如果我自己有飛機多好.
好不好.
好啊.
但也要奉獻一下.
這是我在說的.
我的意思是這個世界上有很多人.
以宗教之名自己有benefit.
自己有好處.
是in the name of God.
我真的不能想像.
我真的要撓頭了.
大家在談這些故事的時候.
為什麼會有人寧願去那些教會呢.
為什麼會有人寧願去參與那些活動呢.
這起碼就是我們信仰裡面最大的妥協.
居然給了一些人自己有福利.
有一些不合道理的benefit.
我有幾件事要回答給你.

$^{601}$就是你把信仰轉一轉.
for your own benefit.
這就是巴蘭的教導.
基督勸喻這個教會做什麼.
所以你們應當悔改.
若不悔改,我很快就會到你那裡.
用我口中的利劍來攻擊他們.
凡有已應當聽聖靈向宗教會所說的話.
得勝的,我必將那隱藏的嗎哪切給他.
並賜給他一塊白石.
石上寫著新的名字.
除了領受的以外,沒有人能夠形成.
這是一個很嚴峻的警告.
如果你把福音,信仰變成了我們自己有benefit的話.
上帝的審判是嚴厲的.
上帝必會審判.
這是一個很大的警告.
我將我口中的利劍來攻擊他們.
耶穌吩咐,如果你跟我的話.
你先相信.
你全然相信我.
你知道我會恭迎你,知道我會保住你.
你先不要說上帝先來給我福.
上帝是最大的權能.
不讓我們在祂面前去爭辯.
相反的上帝說,我邀請你.
你去跟我的話,你相信我的話.
你先全然相信,然後我會恭迎你.
你信得過我是一個有豐盛的上帝.
如果你是這樣的話,你得勝的話.
就會有隱藏的嗎哪來賜給你.
什麼叫隱藏的嗎哪?.
嗎哪就是以色列人在曠野行走四十年的時候.
上帝每一天去賜予給他們.
沒有多也沒有少.
上帝有忠心的,忠實的.
按照祂的應許來供應你的需要.
把嗎哪賜給你.
把白石上有新的名字也賜給你.
你去羅馬城,到處走的時候.

$^{641}$就會看到很多雲石.
然後雕了羅馬王的名字.
哪個大官的名字.
那些就是垂名稱氏.
永遠有人記住,永遠有人去紀念的.
耶穌說如果你是得勝,你是相信我的話.
我就賜給你那白石上有你的名字.
就是說我會記住你.
在我記憶裡面是不能夠拆去的.
因為你保持信仰.
大兄姐們,最後我今天要去到的呼籲.
就是這一段聖經.
Repent,悔改.
讓我們悔改.
Turn around.
這一位滿有權柄的耶穌基督.
是值得我們全心全意去相信.
Obey 祂,信服祂.
其他東西祂會供應我們.
要相信得過上帝的權柄.
相信得過上帝.
恩典是奉承的.
好嗎.
我們聽完這封信的時候.
讓我們有一點時間去祈禱.
今天去反省一下到底我有些什麼.
是我需要悔改,回到上帝面前.
相信祂的奉承,相信祂的權柄.
這段時間,大概一分鐘.
讓你在上帝面前默禱.
(安息).
因為你是那位滿有權柄.
並且在口裡有兩刃利劍.
滿有權能的主.
主,我們再一次來到你面前.
我們認定,你才是那位獨一的主宰.
是值得我們全心全意去相信,信教.
主,我們求你給我們一個悔改的心.
主,讓我們知道我們跟隨你的時候.
我們先說的不是我們有什麼條件.

$^{681}$我們先說的不是我們有什麼福分.
而我們先去明白,你是那位真正的主宰.
讓我們的生命是被你所統治.
讓我們的生活被你所帶領.
主,幫助我們來聽你的話,去跟隨你.
也讓我們相信你是那位滿有恩典.
滿有權柄.
你的公認是充充足足的神.
主,你求你給我們軟弱的信心.
能夠再一次來到你面前.
我們回傳,向我們全心全意敬拜.
去相信你.
主,從我們自己個人的生命裡面.
我們成為一個真正去跟隨你的人.
從我們的家庭.
我們自己生活裡面每一個的關係.
主,讓我們看到你是那位在我們中間統治的上帝.
主,讓我們的生命,我們的行言,我們的舉止.
我們所說的,所想的,是更加像你自己.
主,讓我們這個生命能夠活出公義.
活出你自己的美事.
也讓我們教會能夠成為一個.
在這個城市繼續去見證你的教會.
幫助我們傳講你的福音.
去邀請其他人來到我們自己面前.
主使用我們.
在這個禮拜裡面,求你讓我們能夠盡忠.
去做好我們的事.
做你的見證.
祈禱奉祝明教.
吾以上帝祝福各位.
下個禮拜能夠見到大家,一起敬拜.
\newpage



\section{}
\label{sec:yfuGrCAshKk}
\textbf{What Christ Thinks of the Church \#6 (Revelation 2-3) - Rev. Mark Chan}
\newline
\newline
連結: \href{https://youtube.com/watch?v=yfuGrCAshKk}{\texttt{https://youtube.com/watch?v=yfuGrCAshKk}} ~~~~ 語音日期: 2024-10-08
\newline
\newline
\hyperref[sec:1ry__RW4dIE]{\small{< < < PREV SERMON < < <}}
~
\hyperref[sec:index_chronic]{\small{[返順時目]}}
~
\hyperref[sec:index_scriptual]{\small{[返順卷目]}}
~
\hyperref[sec:vjkb_KdvSz0]{\small{> > > NEXT SERMON > > >}}
\newline
\newline
$^{1}$All right, if you have a Bible, we'll continue our series.
on the seven letters to the seven churches..
So please turn with me to the book of Revelation..
And by the way, this is the first of the month, okay?.
This month, we're gonna recite this verse.
and remember the promises there..
Let's do it together, and I picked this last week.
because it is a reminder that we all need this,.
forgiveness of sin from time to time, right?.
And God's grace is always sufficient,.
and sufficient to the point that, you know,.
even though we fall and we stumble,.
but God's grace is always there, all right?.
Let's recite it together, 1 John 1, 9..
If we confess our sins, He is faithful and just.
and will forgive us our sins.
and purify us from all unrighteousness, all right?.
Why don't we do that?.
Take a moment to confess our sin,.
to confess that we have our shortcomings,.
yet God's grace is always there, all right?.
Let's take a moment..
Father, once again, we experience.
the unconditional forgiveness of sin,.
all because of what Jesus has done for us..
And thank you for the renewing grace and love.
that we can experience..
And thank you for your promise.
as we confess our sins and our shortcomings..
You are there to embrace us with your grace.
in Christ, we pray, amen..
So let's turn to the book of Revelation..
Turn to book of Revelation chapter three, all right?.
We have studied some of the letters.
written by Jesus to the seven churches..
And today we're gonna talk about this letter.
to the church in Philadelphia..
So it's not the Philadelphia in United States..
It's Philadelphia in today's Western Turkey..
That is one of the cities the seven churches address to..

$^{41}$So we study this from part one..
Jesus wrote a letter to the church in Ephesus..
And then Smyrna, Pergamon..
Thyatira and Sardis..
And now we come to Philadelphia, Philadelphia..
And by the way, the name Philadelphia.
means brotherly love, right?.
Philo, love, right?.
Adelpha, Adelphia, the second part is brother..
So brotherly love, city..
So this church or house churches.
located in this ancient city, Philadelphia..
And the letter is pretty short..
And there are two points here..
I wanna make two points from this letter,.
from the sixth letter to the church in Philadelphia..
The first point is your work, your ministry,.
your serving of God, it matters to God..
And second point is that your work,.
your ministry will be awarded, all right?.
It will be awarded..
The two points here..
So let's read through this letter..
Verse seven of chapter three..
To the angel of the church in Philadelphia write,.
these are the words of him who is holy and true,.
who holds the key of David..
When he opens, no one can shut..
What he shuts, no one can open..
This is how Jesus introduced himself.
to the church in Philadelphia..
And it talks about key, right?.
Key..
You should have some keys in your pocket, right?.
What did they symbolize?.
Keys..
I remember a couple months ago,.
because I got a very, very cheap used car..
And so I bought this from a friend..
And I got the car and there was only one copy of the key..

$^{81}$One copy of the car key..
And knowing that I'm very forgetful,.
I'll lose things, you know?.
When you ask my wife, I lose things all the time..
And because we only have one key.
and I've become nervous using that car..
So I drove this car to a hardware store nearby.
and to have my car key copied..
And I didn't know how much it would cost..
And I asked the manager, how much would that cost?.
And to my surprise, it's very expensive..
Any guess?.
To copy a car key?.
300, \$400..
Woo, 300, \$400..
I thought I can just go to Home Depot and make a copy..
But it's not that simple..
And the manager had to ask me questions.
whether I own the car, right?.
Whether I have the registration to show him.
that I really is the owner of the car.
and he has to check my registration, my title,.
and make sure I really own the car.
and I really have the car in the parking lot.
and he can go out and check it out..
And he's gonna order a blank, a blank key..
And then I have to come back..
I pay for the blank key and I have to come back.
like a couple weeks ago after I got a call from him..
And it's a very, very complicated process..
And I really appreciate that..
You know why?.
Because it's not gonna be, I just steal a car.
and drive it to a hardware store and have a key copied..
And I really appreciate that because this will prove.
that I'm not doing anything criminal..
And it finally took me a couple of weeks.
to have my car key copied and I have a peace of mind..
300, \$400, I mean, this is crazy, right?.
Keys are access and authority..

$^{121}$But you have the key to my car, you have access to it..
You have the key to my house, you have access.
and you kind of have authority..
And in Bible times, keys are very, very big, very large..
Sometimes people would carry the key over the shoulder.
with a lanyard or sword..
And then you have your key under your armpit..
And you can show people you have a big house..
And here's some of the pictures of the Roman time..
You can see some of the keys,.
one of the keys here on the left..
You find a bronze key from Roman's time.
and it is a horse head handle key..
It's huge, you just stick it in..
I don't know how heavy it is..
But it's made with bronze and metal..
On the right is also a Roman key..
It's an iron key with a iron handle..
That's kind of cool, right?.
That's kind of cool..
I don't have those keys, right?.
And Jesus said, he has the key..
And he has the key to the kingdom of David..
And that is the kingdom, the kingdom of God promised.
long, long time ago in the Old Testament.
and finally fulfilled and inaugurated.
by the true son of David, Jesus Christ himself..
And he has the key, he has the authority.
and access to the kingdom..
And this is how Jesus introduced himself to the church..
And he wants the church to remember who he is..
Jesus is not simply a guy out there on the street..
He is the God, the true son of David.
who inaugurate, usher in the true kingdom,.
the heavenly kingdom of David and the eternal kingdom..
And in Christ, we share that access..
In Christ, we share that authority..
We have the access to go into the kingdom.
and we are part of the kingdom..
And to this little church, Jesus wants his church.

$^{161}$to be reminded, no matter how many demands.
or how many challenges you are facing,.
you are depending on a God that is dependable..
And you can trust in him..
We all need this kind of truth..
Sometimes we're dealing with life,.
with a lot of different issues and challenge.
and we are fighting the battles..
And sometimes we forget that the Jesus that we believe in.
is the one who has the final authority..
He has the final access to the kingdom.
and he guarantee the access to the kingdom..
He has the key..
And we have to remember that..
In Jesus, we have the eternal hope and eternal life..
And we have the eternal source of strength to go on..
And why did he say that?.
We continue to read this letter from verses eight.
and following..
Jesus said to the church, "I know your deeds..
"See, I have placed before you an open door, the access,.
"and no one can shut it..
"I know you have little strength,.
"yet you have kept my word and have not denied my name.".
You see what the church is going on here?.
They're facing a lot of different problems,.
different challenges, just like our church..
And Jesus reminds them, "Yes, you have little strength,.
"but you have kept my word..
"I am here to protect you..
"I am here to walk with you.".
In trials of difficulties, Jesus is still there..
And what kind of trial they're going through in verse nine?.
He said, "I will make those who are of the synagogue of Satan.
"who claim to be Jews, though they are not, but are liars..
"I will make them come and fall down at your feet.
"and acknowledge that I have loved you.".
Now, in the context of the seven letters,.
we have already talked about or described.
what kind of situation the churches were in..

$^{201}$In the first century, they were facing.
a lot of different difficulties and persecution..
And the source of persecution.
is sometimes from so-called the Jews..
As you understand, the early Christians,.
they kind of spin off from the Jewish synagogue.
in the first century..
And many of the Jews in those days,.
they do not believe Jesus is the Messiah..
They do not believe Jesus is the Christ..
And in fact, they not only rejected,.
they would persecute and go after and oppose to anyone.
who proclaim that Jesus is the Messiah..
And that's why the early churches were facing.
a lot of different problems from these Jewish leaders.
who's trying to stumble down the church,.
try to destroy the church..
But Jesus has the promise for the church..
Jesus said, "I will make your enemies,.
"the difficulties, the challenges,.
"to fall down at your feet because I have loved you,.
"because I have taken care of you.".
And so you take heart..
There were a lot of different problems.
and challenges you are facing,.
whether you go to school today.
or whether you go to work today,.
whether we live in a society that is hostile.
to the name of Jesus..
But Jesus is saying, "Not only I have the key.
"to the kingdom, but I will take care of you..
"I will make your enemies fall down before you.".
But in fact, this phrase is alluding to.
some of the Old Testament promises in the past.
that God said all the nations will come to Israel,.
will come to the true Israel, the people of God,.
and they will fall down and be humble,.
and they will acknowledge that the Messiah, Yahweh,.
God, the true God, is the one that they worship..
And the nations will one day come before the throne of God.

$^{241}$and acknowledge and admit that Jesus is the king..
That is the promise here being alluded to..
And this letter is giving us so much promise,.
so much encouragement..
I don't know what you're going through..
Sometimes you may be going through some very difficult time,.
but don't you forget, don't you forget who Jesus is..
He is the one who can conquer enemies..
He is the one who hold the key to the kingdom..
Now, I told you before, there are seven letters.
to the seven churches..
And if you remember, I said there were two letters.
out of the seven letters that have no condemnation at all..
Because if you read through all seven letters.
to the churches, they have similar format..
They start with self-introduction,.
how Jesus introduced himself,.
and always followed by the phrase, I know..
I understand your situation..
And then there will be praises and condemnation.
for the church..
There are some praises for the church..
Yes, I know your deeds, I know what you're doing,.
and I wanna praise you guys, I wanna appreciate you guys..
And then it will be usually followed.
by some sort of correction..
Something Jesus has against the church..
The church has gone wrong at some places,.
and they need to repent..
But there are only two letters that have praises only..
There was no correction..
The first one is the second letter to Smyrna..
Remember?.
Praises only..
The small churches in Smyrna,.
they're going through some hard time,.
but the second letter says,.
Jesus said, you know, I commend you..
You guys are doing great..
And this is the second letter that has only praises..

$^{281}$The letter to Philadelphia..
Praises only..
That's letter number six..
And we should be inspired to become a church like this..
Smyrna or Philadelphia..
This church, we receive praises only.
from the head of the church..
You know, if Jesus walks into our church,.
what would he say?.
I'm not so confident that we'll be like.
the church in Philadelphia..
However, this letter should be an inspiration.
for all of us to become like them..
Church in Philadelphia..
Not a mecca church, but faithful..
And Jesus said, I know your deeds..
And this really challenges our perspective on the church..
Our perspective on our lives, our Christian life..
When Jesus do his assessment, what is he looking for?.
Here it says the church in Philadelphia.
have little strength only..
Little strength only..
Little resource..
Maybe they have less people..
Maybe they have less financial resources..
Maybe they have less people who are gifted..
And Jesus observed that in his assessment..
Jesus said this church, you have little strength..
It's just stating a fact..
It's not a condemnation, right?.
There's nothing wrong with having little strength..
Having little resource..
But according to Jesus' critical assessment,.
he says, but yet you have kept my word.
and have not denied my name..
Under pressure..
When you're under pressure,.
when you're under difficulties,.
we are always tempted to deny the name of Jesus..
Maybe today I'll take a break..

$^{321}$Maybe today I'll forget about Jesus, okay?.
Put him on the side..
In the old days, I have a lot of books on my bookshelves..
Today, I've totally transformed into a e-book person.
because I'm running out of space in my house..
I used to have 5,000 books in my house, physical books..
And boy, when you have to move,.
it's so painful to move..
That's why years ago, I already made up my mind,.
okay, let's digitalize everything..
So slowly transforming to a totally e-book person now..
I haven't bought a book for a long, long time,.
but e-books, yeah, I did..
But in the old days, I have so many books on my bookshelves..
But in the days of seminary, when I was taking MDiv,.
because studying theology is really, really tough..
It's really tough..
There were so many books I have to read..
What is so sad is that in seminary,.
sometimes the least you read is the Bible..
Amen?.
Sometimes we just don't read the Bible..
We just don't have time to read the Bible..
Because you go to a course, you receive a syllabus,.
and you have to read like 2, 3, 4,000 pages, right?.
Church history or exegetical work or whatever,.
commentaries you have to read and reference books, you know..
And you just don't have time to read the Bible..
I tell you, sometimes when you are in a lot of pressures,.
you will put the Bible on the side..
And that's what I did..
The word of God, the very word of God..
And sometimes as Christians, we do that too..
On your phone, do you have a Bible app?.
How often you open the Bible app.
compared to open other apps, right?.
24 hours a day, you probably always go to TikTok..
Always go to Instagram..
Always go to social media app, right?.
But the Bible app, you put it on the side..

$^{361}$Maybe just delete it, right?.
You don't need that..
And when we are pressured, we are so often,.
we would deny the name of Jesus, right?.
The church in Philadelphia, however,.
they insist worshiping Jesus..
They insist on affirming and keeping the gospel,.
keeping the faith..
And that is the encouragement.
from Jesus himself to the church..
You have kept my word, I know it..
I see it..
You have not denied my name, I see it..
And Jesus further promised,.
I will put all of this enemies into,.
to your feet, to the feet of Christ..
I will turn them into worshipers..
Just like what Isaiah a long time ago,.
he prophesied the Gentiles,.
the nation will worship the King of Kings..
Isaiah 49, Isaiah 60,.
that the church will be triumphant because of his grace..
So that's my first point from this letter, your work..
If you keep the word of God,.
if you keep the word of God,.
if you keep on praising Jesus,.
if you keep on following Jesus faithfully,.
God sees your work and your work matters to him..
And secondly, your work will be rewarded..
And you have to remember that..
We have to remember that..
Sometimes we serve God..
Sometimes we continue to follow Jesus..
We lose sight of that..
We thought this is just here and now, right?.
Why bother?.
Why bother we serve the King,.
our Lord with all of our might and our resources?.
This is so hard..
And the reality you need to understand ministry is hard..

$^{401}$I don't know if this is true or not..
I heard some time ago about church leadership,.
about church ministry..
I attended one of the church conferences before.
and the speaker was saying,.
in this world, church leadership or church ministry.
is the hardest thing on earth..
It's the hardest thing to lead..
It's the hardest work of all..
I don't know if you agree with that,.
but whether it's the hardest,.
the most challenging things on earth or not,.
it is challenging..
It is not easy..
You just have to keep on going..
And what motivates you to keep on going?.
It's not your own strength..
It's the hope that we have in Christ..
It's the eternal reward..
Yes, sometimes we do our best in our ministry.
and we don't see fruit..
What can motivate us to continue?.
Sometimes we're in the midst of challenges in ministry,.
in our church, in our just general work of Christian life..
Sometimes we will give up.
because we just think it's not worth it..
Jesus said it's worth it because there was a future.
and the future is reward..
Let's continue to read this letter, verse 10 and following..
Since you have kept my command to endure patiently,.
I will also keep you from the hour of trial.
that is going to come upon the whole earth,.
the whole world to test those who live on the earth..
Now, this is very interesting..
Let's unpack this..
He said, "You have kept my word to endure patiently.".
I like the old translation..
The older translation would use the word steadfastness..
Steadfastness, yeah, we don't use that word anymore..
But steadfastness, it means remaining and unyielding.

$^{441}$in the face of passing time..
As times go by, you are unyielding in your spirit..
You keep on going..
That is endurance, patience..
And also Jesus said, "If you do that,.
"I've seen you do that..
"Since you have kept my command patiently and stand fastly,.
"I will also keep you from troubles in the future.".
Now, in the later chapters of the book of Revelation,.
the book talks about the coming intense, difficult time.
for the first generation, first century church..
There will be a big calamity,.
apocalypse happening in the world..
And a lot of people will suffer..
And a lot of Christians will be tested..
And many more from the book of Revelation,.
we know that they will be a martyr as well..
But to this church, because they have little strength,.
but they have kept the word, they kept on going,.
Jesus promised them, "I will not put extra burden on you.
"because you have been faithful day after day,.
"and I observe it for a long, long time..
"So I will shield you from other troubles,.
"from other difficulties that are coming.".
And that is a word of comfort..
Have you ever asked the question, what God require of you?.
I remember when I was studying theology,.
one of my professors said something.
that really stick into my mind..
He said, "What does God really require of you,.
"require of us, especially in our context in North America?".
Well, we are not in a country.
that is openly hostile to the gospel..
We have freedom..
We can choose to come to church..
We can choose to follow Jesus..
And yeah, there may be some difficulty..
There may be some challenges from our family members.
who are unbelievers..
Maybe it's from our neighbors or whatever..

$^{481}$But in terms of the level of persecution,.
there is nothing, right?.
So my professor asked, "What does God really require of you?".
And he asked a very interesting question..
He said, "Do you think God really requires.
"so much out of us?.
"What does God demand?".
In terms of that demand, it's actually very reasonable,.
very light, and God expects us to worship Him,.
to dwell in His Word, day by day, to live faithfully,.
to bear the name of Christ, to be a good witness..
That is not so difficult..
God is not requiring you to die on the cross physically,.
right, to go through some physical, intense suffering..
There may be Christians in other parts of the world,.
like in the Muslim country or communist country..
They may have to go through some of this.
really, really trying time..
But for us, it's actually pretty light, pretty easy..
And what God has required for us.
is that we will follow Christ faithfully day after day..
That's pretty much our habit..
That's pretty much our pattern of life..
So this letter reminds us,.
your little work, your faithful work,.
will be remembered and will be rewarded..
Let's finish reading this letter..
Jesus said, "I'm coming soon..
"Hold on to what you have..
"Even though you have only a little,.
"hold on to what you have.
"so that no one will take your crown..
"Him who overcomes, I will make a pillar.
"in the temple of my God..
"Never again will He leave it..
"I will write on Him the name of my God.
"and the name of the city of my God, the new Jerusalem,.
"who is coming down out of heaven from my God..
"And I will also write on Him my new name..
"He who has ears, let him hear.

$^{521}$"what the Spirit says to the churches.".
He says pillar in the temple, in the new temple..
In the book of Revelation, in the later chapters,.
in the final chapters, God said the new temple will come,.
the new Jerusalem will come..
And as you know, if you study a little bit,.
or know a little bit about architecture,.
the ancient temples have big columns, huge columns..
You know, the Doric columns, the Ionic columns,.
the Corinthians column..
All those different columns are standing.
to support the temple, and it's part of the temple structure..
If you've gone to a film museum, those are the columns.
the ancient architects will put.
into the structure of the temple..
And the temple is the house of God..
The temple is where God dwells..
And here is the picture of stability,.
a picture of security, that we can be part of the pillar,.
part of the house of God, to support it,.
as a fair, be part of the fair structure.
of God's new temple..
The city of God, the new Jerusalem,.
will descend one day from heaven,.
and God will dwell among us..
See, being part of God's temple, God's home,.
means that you are always in God's home..
God will always remember you..
And so often, we don't wanna serve, we don't wanna give,.
because it's too hard for us..
We don't wanna use our talents,.
we don't wanna give our treasure,.
because it's just too difficult for us..
But if we, if we do that,.
we have this kind of perspective..
It's only because it's too difficult,.
that's why we don't do it..
But we have to ask ourselves,.
do we really believe in God?.
Do we really believe that God is good?.

$^{561}$Do we really believe that God is our strength?.
And God is sufficient?.
Here, although the letter is a letter of praise,.
it's also a reminder, or perhaps a warning..
Warning to all of us, as we read through this letter,.
hey, you need to make sure.
that no one will take your crown away..
Yeah, your crown, your glory can be taken away..
If you don't remember, if you forget.
that your work right here, right now,.
your work of serving, your work of following Jesus.
will be remembered..
If you do not remember that,.
it's a warning and reminder.
that your crown can be taken away..
We can lose our place in God, we can..
It's possible..
It's possible if we don't keep on trusting God,.
we lose our place in the kingdom..
We'll lose our crown..
But the letter reminds us with words of comfort,.
words of encouragement, God will remember your work,.
no matter what that is..
You do it in the name of Jesus, you do it by faith..
Maybe you don't see the fruit right now,.
maybe you don't see the fruit tomorrow,.
or next week, or next month, or next year..
But if you do it by faith, God will reward it..
God will remember..
So your work matters to God, and your work.
will be remembered..
Don't you forget that..
Don't you forget that as we build a church together..
Last week we were in our general assembly,.
and we shared a lot of different updates, right?.
And needs..
As our church grow together,.
we need brother and sister to come on board.
and one of the things we can do is to give..
Because we are growing and our financial need is bigger..

$^{601}$It's bigger than ever..
It is part of the growing pain..
And we have to remember that our work,.
our contribution will be remembered..
Right?.
It will be rewarded..
And don't you forget that your crown,.
our crown can be taken away..
So let's press on, guys..
Let's press on as a family of God to do the good work,.
to do the ministry in the name of Jesus,.
to serve Him, to build His body, to build His kingdom..
Let's all pray together..
Father, we come before You,.
we just open our hearts and our life..
And we confess sometimes we are forgetful..
Sometimes we see things using worldly standard,.
worldly perspective..
And thus we quit..
So we give up because we thought it's not worth it..
But yet Your letter to us.
bring us a new and transcendent.
perspective on all the things that we do..
Help us to follow You,.
even though we have little strength,.
but we know that our work, our ministries, our serving,.
all that we do in the name of Christ,.
all that we do for Your gospel, for Your kingdom.
will be remembered..
And thank You for calling all of us to be Your partners.
so that we can work for Your kingdom,.
we can build up Your kingdom..
What a glorious role and responsibility You lay on us..
And thank You for giving us such an easy context.
in North America..
We can follow You freely, we can serve You freely,.
and that's why we need to give You more..
And thank You for Your blessing..
We need to give You more to respond to Your goodness,.
to Your abundance, so that we can give more,.

$^{641}$we can serve You more,.
to uplift the name of Jesus in the midst of us..
And once again, we wanna close in prayer,.
and we just wanna pray that You will make us.
into Your instrument..
We can glorify Your name, we can bear witness.
to the goodness of our Lord Jesus Christ..
And thank You for being our God..
And thank You for Your love and grace.
and Your comfort and encouragement..
Help us to press on..
To follow You faithfully, day by day,.
in Christ's name, we pray..
Well, God bless you all,.
and hopefully you will go out this week.
to bear witness to the name of Christ.
and serve God in the name of Christ..
Do so by faith, and God will remember it..
All right, see you next week..
\newpage



\section{5:17-18}
\label{sec:vjkb_KdvSz0}
\textbf{工人的工價 (提摩太前 5\_17-18) - 陳崇基牧師}
\newline
\newline
連結: \href{https://youtube.com/watch?v=vjkb_KdvSz0}{\texttt{https://youtube.com/watch?v=vjkb\_KdvSz0}} ~~~~ 語音日期: 2024-10-08
\newline
\newline
\hyperref[sec:yfuGrCAshKk]{\small{< < < PREV SERMON < < <}}
~
\hyperref[sec:index_chronic]{\small{[返順時目]}}
~
\hyperref[sec:index_scriptual]{\small{[返順卷目]}}
~
\hyperref[sec:hfmoR2RDphM]{\small{> > > NEXT SERMON > > >}}
\newline
\newline
$^{1}$多謝Nelson.
Pastor Appreciation Month.
我覺得經常去到這個月的時候.
我們去做一些這樣的事.
我都覺得有點不好意思.
說真的.
我們的童工.
我夠膽代表我們所有的教務童工.
有好幾個是兼職.
我們真的不是為了錢.
我們主要的原因.
事奉.
是因為我們知道上帝的恩典是真實的.
福音給我們有很大的改變.
以致我們的生命可以奉獻給上帝.
我們能夠全時間的.
更多時間的去服侍.
去推動教會去幫助教會.
這個真是我們的心願.
我去過不同的教會.
我覺得今時今日我有的教會童工團隊.
是一個最好的團隊.
因為大家真的用心去做事.
我去過好幾間的教會.
有時候真的不知道在做什麼.
因為很多時候要處理很多不同的人際的問題.
教會的問題等等.
有些很頭痕的事.
但在UV我覺得真是一個很好的團隊.
所以為了這件事去感恩.
正因為這樣的緣故.
教會可以繼續去發展.
教會是有更大的需要.
正如我們剛才說過.
有新的童工是需要.
能夠將他們提拔出來.
以致教會能夠繼續向前.
今天本身是想說啟示錄.
因為有這麼多的報告.
我知道沒有什麼時間.

$^{41}$所以我想說一篇.
On the fly的說.
On the fly.
就是剛剛想的.
真的剛剛想.
你問問後台Phoebe.
我說今天還是不要出.
我昨晚預備的PowerPoint.
下個星期好嗎.
菜已經煮了.
不過下個星期才serve.
今天我就On the fly.
說一篇.
我想起一段聖經.
如果你有聖經的話.
打開.
是兩節.
在提摩太前書第五章十七和十八節.
PowerPoint沒有的.
提摩太前書第五章十七至十八節.
這裡你聽我讀.
善於督導教會的長老.
尤其是勤勞廣道教導人的.
應該得到加倍的敬奉.
因為經常說牛在全谷的時候.
不可以勞著他的嘴.
又說工人得工假是應當的.
大家知道這個所謂教務的書信.
提摩太前後書和提多書.
是去教導教會的領袖.
或者教導教會的群體.
怎樣去處理教會裡面的事務.
怎樣去發展教會.
和處理教會裡面的關係等等.
所以你打開提摩太前後書.
提多書.
這三封書信很短的.
就是保羅寫給他兩個徒弟.
提摩太和提多.
他們分別在不同的地方.

$^{81}$處理和牧養當時的教會.
有很多的問題.
譬如有異端.
有一些人傳假道.
提摩太和提多這兩個徒弟.
應該做些什麼呢.
應該怎樣去教導人呢.
還有在教會裡面遇見不同的問題.
譬如長老的問題.
怎樣處理教會的資源問題.
怎樣去照顧有需要的人.
特別是一些寡婦等等.
所以有很多的你可以叫做policy.
你可以叫他做一些原則.
在這三封書信裡面都出現的.
這裡第五章十七至十八節.
就是講到一些在教會裡面.
做教導和督導的.
幫忙做領導的一班人.
他說這些人是什麼呢.
他說如果他們是勤勞去教導去督導的話.
這班人應該是得到加倍的敬奉.
加倍的敬奉.
其實是在講一些很實際的東西.
是薪酬問題.
我知道牧師講這些東西是很煩.
還有不是很好意思.
但是我想我都老到一個地步.
我都不是很care.
我去過不同的教會.
我告訴你.
我一直這麼多年.
我都沒有跟教會的長輩.
教會的領袖.
去討論薪酬問題.
甚至我答應去一個教會事奉的時候.
我都不會問到底他的package是什麼.
我真的沒有問過.
我也去過教會.
我進去的時候才發現.

$^{121}$原來我的薪酬這麼少.
師母一定要出來打工.
如果不是我們搞不定.
我不知道這樣做對還是不對.
I don't know.
但是因為要開口講一件事.
確實是有些難為情的.
是有些麻煩的.
既然我已經老到這個地步.
我都不是很care的時候.
不如我為教會和為教會的下一代.
去試一下講這件事好嗎.
我希望你講完之後.
你不會覺得又是拿錢.
我希望不是這樣.
我希望本著聖經裡面.
告訴我們我們作為教會.
作為信徒.
作為一班的基督的門徒.
我們應該怎樣去面對這個問題.
我們已經有二十多年的經驗.
我也覺得教會裡面的一些文化.
我是今天可以告訴你.
親身來說我可以告訴你.
正如我所說我們教會.
我相信如果網上一些大英姐妹.
聽這篇讀的時候.
可能你的教會也是一樣.
教會可以向前.
教會可以繼續去到發展.
去到成長的時候.
是有更大的需要.
是財務上的需要.
每一個童工是希望有一個合理的薪酬.
因為我想講得很清楚.
我只是講我們教會.
我們沒有一個童工是為了錢.
沒有一個童工真的進來的時候.
或者我們去呼召他.
參與我們這個行列的時候.

$^{161}$本身有的或者剛剛新任的童工.
沒有一個進來的時候.
才跟你講價.
因為教會文化很有趣.
如果你先開始講價的話.
通常是不請你的.
什麼事呢.
你門口都還沒進來.
你就跟我說套餐.
你的薪金.
你是什麼人.
因為我們事奉的不是為了錢.
是真的.
我只是講我們UV的團隊.
我們的童工.
我們每一個都是很拼命.
很努力的.
來事奉的.
沒有人跟我們講薪酬.
也沒有人去講什麼時候要放假.
放多少天假.
怎樣才算合理.
這些可能跟外面的公司有些不同.
不過我想告訴你.
這個就正是福音的精神.
這個就是耶穌基督給我們的榜樣.
耶穌基督給我們的生命.
我們是擺出來事奉的.
但是在教會裡面.
有一件事很有趣的.
當然我們想教會發展.
而教會發展的時候.
有更大的需要.
財務上的需要.
童工是應該給他們合理的薪酬.
我想說是合理的薪酬.
不是一定要給他們多錢.
而是需要給他們一個合理的補償.
因為他們也要生活.
你不能夠說.

$^{201}$如果教會的童工.
一個工人.
如果他是要優柴優米的話.
他才可以事奉的話.
這個教會是不能夠向前的.
物價高漲的疫情大家都知道.
最近的蛋又貴了.
是不是好像便宜了.
前陣子很貴.
是不是很貴.
我老婆又跟我說了這句話.
她說你不要吃蛋.
留給女兒吃.
她說我每天早上又吃兩隻.
便宜的時候又吃兩隻.
現在貴了就不要吃了.
然後我說.
我很乖的.
你知道我老婆說話是雙重意思.
她其實就是說你不用吃.
反正你吃少了那麼多.
不如吃少一點.
我就說好.
不吃了.
inflation.
大家都知道物價高漲.
是真的令到我們大家很大壓力.
對教會也是一樣的.
是不是.
對教會也是一樣.
因為這樣說.
教會的童工.
如果經常要優柴優米來事奉的話.
是事奉不到的.
我曾經有很多年.
我都是事奉的時候要優柴優米的.
怎樣能夠專心事奉.
對嗎.
做不來的.
但教會很多時候有一個這樣的.

$^{241}$很奇怪的文化.
怎樣講新金.
在開會的時候.
又不一定是你去定會議的議程.
又不一定有時間去討論這件事.
其他大英姐妹.
特別是一些執事和我們一起童工的.
他們都有自己的家庭.
他們都有自己的工作.
有自己的demands.
他們未必有時間去理你那一份.
是不是.
所以很多年有很多時候.
是真的有優柴優米的時候.
但我想告訴大家.
如果一個健康的教會就不應該是這樣.
教會文化是很有趣的.
我們希望有年輕的童工出來.
我們希望能呼召一些年輕的童工.
去帶領著教會.
你都不想一個老的童工一直帶下去.
因為總有一天我都會死的.
我都會老的.
我都會退休的.
或者應該要退休.
但我們希望有年輕的童工出來.
但文化有趣的地方就是.
你希望有年輕的童工出來.
但年輕童工自然是因為年輕.
因為年輕就未必有獨立的經濟能力.
然後他的經驗等等可能是未足夠.
我們教會很多時候薪水就要給.
所謂特別低一點給他們.
但這樣又不是老.
因為我希望有年輕的童工出來.
但我們又知道實況.
年輕童工未必有獨立的經濟能力.
他又要優柴優米的才能夠去事奉.
今時今日我都講過了.
東西貴了很多.

$^{281}$不要講食物不要講生活是貴了.
特別城市生活更加貴.
芝加哥.
生活是貴.
甚至你去讀神學都貴.
讀神學院都貴.
你不要想.
不知道大家有沒有東方宗教的背景.
覺得那些尼姑.
或者那些佛寺裡面的人.
他們不用吃東西.
他們就在那裡打坐.
那些叫神學院.
No.
讀神學院很辛苦的.
讀神學院很辛苦.
我讀過好幾年神學院.
我現在在讀.
我沒有去過一個神學的班.
和神學的同學.
神學班裡面的同學.
年輕的也好.
年紀大的也好.
沒有一個輕言易舉.
用什麼.
背兩條formula.
然後考試的時候我再拿A.
是很費力和很費思想的.
我們都希望是這樣.
如果童工傳道人.
他們去到神學院的時候.
他們得到的裝備是最好的裝備.
我們都不想在一個這麼紛亂資訊的世代裡面.
都是有人亂解經.
亂去到講聖經.
這樣和我們信仰的生命根本是違背的.
問題就是這樣.
年輕的童工他們沒有獨立的經濟能力.
教會人工可能比較低一點.
要他們休茶休米來事奉.

$^{321}$這個是一個實況.
或者你可以建議.
我們不要請窮的年輕童工.
我們請有錢的年輕童工就行.
我這麼多年事奉我也認識不少.
有一些年輕的童工.
他們是所謂的富二代.
他們真的不用休茶休米.
他們生活上做屋賣屋等等完全可以.
他們幾乎可以做義務童工.
這個世界有這樣的人.
當然我不是.
我們沒有一個童工是.
但問題就是這樣.
我們請那些有錢的年輕童工就行.
但因為他們的生活是這樣的時候.
他們生命裡面又沒有壓力.
對嗎.
溫室長大的.
他們生命裡面沒有經歷.
沒有生命的怎樣傳道.
他只不過是書裡面說的話.
對嗎.
你都希望牧師或者教會的領導.
他有生命裡面的經歷.
所以去到和一些童工聊天的時候.
如果有新的童工想加入的時候.
或者想聘請的時候.
我要做的就是.
我要聽聽他的故事.
他們很年輕.
不等於他們沒有經歷.
但年輕裡面都可以很有壓力.
所以你想選擇哪樣.
你想選擇那些很有錢的年輕童工.
但生命裡面沒有經歷.
他們經歷的是甚麼.
這個世界不用休茶休米.
需要的嗎.
因為他從小到大都不用做這件事.

$^{361}$還是我們有些年輕童工.
他們有生命裡面的經歷.
可能經濟上沒有其他人的待遇.
或者沒有其他人更好的條件.
我想問你們選擇哪樣.
你們想選擇哪樣.
選擇有壓力的.
我想告訴大家.
我在教會裡面這麼多年不事奉.
從來沒有商討過薪金.
我從來沒有和教會領袖講過.
討價還價過薪金.
亦都沒有一年我說要加薪金.
大家知道.
我和我的童工講.
和我們call team講.
其他童工可以加.
因為他們是低的.
我是不需要加.
你需要調節的話.
可以調節減低.
我的那份可以減低.
但其他童工是需要有好的待遇.
或者合理的待遇.
這個正是這段聖經所講.
如果勞力去事奉的那班人.
那班教會的領袖.
我就應該加倍去敬奉他們.
要加倍的敬奉他們.
因為他們需要生活.
他解釋一個原因.
什麼原因呢.
保羅在這裡引用摩西律法其中一段.
他說牛在全谷的時候.
不要勞著他的鎚.
在五經裡面的一些律法.
是講到當時的農耕生活.
他講到牛幫你耕田的時候.
牛都是工人.
沒有牠你是耕不到田.

$^{401}$如果牛肚餓的時候.
那些谷的大麥小麥成長了.
未收割之前.
牛如果去吃那些大麥小麥.
你不要勞著他的鎚.
那隻牛是幫你工作的.
你沒有牠是做不到的.
你就隨隨便便給牠吃.
吃得你幾多.
這個當然講到摩西律法裡面有一個vision.
一個所謂social的vision.
對這個世界對這個社會一個意象.
一個願象.
就是一個creation.
即是被造界 創造界.
都是需要我們去悉心照顧的.
所以牛在全谷的時候.
不要勞著牠的鎚.
不要虐待牠.
不要苦待牠.
給牠吃飽.
如果牠不吃飽.
怎樣能夠幫你去耕田呢.
所以他又引用另外一句.
他說工人得工資.
工人得工價是應當的.
合理的工價.
我在教會裡面事奉了這麼多年.
正如我所講.
我老到一個地步.
我有衣有食我算了.
我有衣有食我就知足.
唯一一樣我不能夠知足的.
就是我還有上帝給我的使命我未完.
我完了它.
這個就是我的心.
我相信我團隊裡面每一個的同工.
他們的心都是這樣.
所以我在這裡想呼籲弟兄姊妹.
我們回到聖經裡面看.

$^{441}$回到去了解一下教會的運作.
教會的發展我們需要一些更大的需要.
我想呼籲大家去想一想.
到底我們奉獻的生命.
奉獻的習慣.
有沒有養成.
我們的信心是不是只是憑口講.
教會裡面這些問題我們需要點出來.
我不講這些同工的薪酬等等的問題.
我們先不講這些.
有一樣東西我很多年都觀察到.
教會裡面我舉個例子.
譬如教會裡面的音樂.
敬拜的音樂.
教會裡面敬拜的音樂我們用一些詩歌.
這些詩歌是誰寫的.
是不是AI寫的.
是有人寫的.
很多時候有些弟兄姊妹可能會.
久不久我們會講.
為什麼中文的詩歌.
中文敬拜的詩歌.
教會用的詩歌.
為什麼不是很好.
跟英文比較好像不是很彼得上.
原因是怎樣.
我想告訴大家一個很奇怪的文化.
就是我們經常覺得這些音樂是免費的.
不用錢的.
因為你唱出來大家聽.
有意就聽到.
有意聽到就用不著錢.
開大聲一點.
但我們不知道背後原來是有一些作曲作詞.
有些應該是經歷過訓練.
生命裡面有經歷有裝備的人.
他們去寫一些好的詩歌.
如果沒有這些人就不會有好的詩歌.
但偏偏教會裡面的文化就是.
我們覺得這些是免費的.

$^{481}$我們是不需要支持的.
你知不知道這些詩歌是要給版權的.
在我們基督教的中文詩歌文化裡面.
這件事是最麻煩的.
因為有很多人寫完一首歌.
他的回應就是.
這些要錢的嗎.
你事奉的.
你事奉就不要講錢.
OK Fine.
事奉就不要講錢.
這些人就是要有柴油米才能寫歌.
或者他們業餘才能寫歌.
騰空他們自己的時間就能寫歌.
這麼多年我認識不少.
一些有音樂上的恩賜和才幹的弟兄姊妹.
甚至他們自己搞一些音樂的敬拜的士工.
無論在加州也好.
在加拿大也好.
我都認識不少.
我沒有一個認識的.
是不需要柴油米的.
每個月都是不夠.
我想問.
因為如果你這樣的話.
你怎麼可能會有好的詩歌出來呢.
是不是.
所以有些大教會做.
有時一些的詩歌敬拜隊.
不同的地方.
香港也好.
在加拿大也好.
他們來到的時候.
他們會希望弟兄姊妹支持他們.
往日舊的時候.
現在沒有了.
那時候還會賣CD.
他們弄了一些CD.
完結了音樂的分享晚會.
報道會也好.

$^{521}$他們就會賣這些CD.
賣這些CD的時候.
都要經過很多很多的Hurdle.
因為有很多人說.
教會不可以買賣.
不可以買賣當然是自由奉獻就行了.
講自由奉獻.
其實就是不需奉獻.
是不是.
或者.
我見到有很多很well off的人.
他們怎麼支持呢.
有幾張CD買一張算了.
或者自由奉獻.
20元.
有現金我給20元給你.
給50元給你.
這個是一個很差的文化.
因為他們在想什麼.
這些東西用不著錢.
最多我們祈禱.
我們多祈禱一點.
上帝給我們神蹟.
有些可以不用吃東西的人.
不用生活的人.
然後可以給我們一些好的事奉.
Praise God.
五餅二魚.
很奇怪.
所以為什麼耶穌說.
給多些神蹟你都沒有用.
因為給多些神蹟給你.
都不會建立你的信心.
That is not the point.
到底我們應該怎樣去幫助教會.
整體基督教福音的事奉.
那種工作.
我們自己可不可以參與在裡面.
我覺得最根底的問題就是.
到底我們骨子裡面.

$^{561}$我們是不是真的信上帝.
還是我們覺得我們奉獻的時候.
上帝是阻礙我們發達.
明明我們存了這麼多錢.
我們又去奉獻.
那不就是上帝阻礙我們嗎.
我們的財務計劃.
我們退休的計劃.
或者是我們安穩一點.
安樂一點的生活的計劃.
上帝永遠都是阻礙我們.
又要我時間.
又要我精力.
還要我去奉獻.
如果我們這樣去認識鬼的上帝的話.
我想我們想一想.
我們骨子裡面.
真的是不是相信上帝.
Do we really trust him.
Or do we really know him.
我們是不是真的認識這個上帝.
聖經裡面告訴我們.
上帝是會祝福有能力供應的.
我們是不是真的認識.
這位上帝是這樣的上帝.
信心的功課你和我都要走.
對不對.
結束的時候.
為了我們教會禱告.
我真的心願我們教會能夠繼續向前.
年輕的童工.
他們可以受裝備.
然後可以事奉.
上帝才得榮耀.
我們拿聖餐之前.
我們同心低頭祈禱.
聖餐也告訴我們.
上帝是一個多麼豐足的上帝.
願意祝福我們.
給我們一個救恩.

$^{601}$他就是這樣的上帝.
上帝你是偉大.
配受尊榮的上帝.
你是創造萬有.
祝福我們口齒白沫的上帝.
主要是讓我們.
在我們心底裡認定.
讓我們不要忘記.
我們就做很多不同的要求.
很多不同的挑戰.
主要很多時候我們忘記了.
你就是那一位願意施恩的神.
願意公認的神.
主幫助我們感動我們的生命.
讓我們能夠全心全意事奉你.
讓我們的生命.
我們願意慷慨.
投資在自己上帝國裡的事工.
主讓我們知道.
你是那一位多麼有恩典的神.
以至我們的生活.
包括我們怎樣去管理我們的錢財.
都是一樣.
我們明白你是那一位慷慨的上帝.
我們怎樣去生活.
怎樣去做一個管家.
都是反映著你自己的本性.
施與愛我們的神.
主要我們弟兄姊妹.
整個教會今天來到你面前.
在這個聖餐桌.
拿著這個餅的杯.
一個Token of Grace.
一個記號.
去告訴我們耶穌基督.
你是那一位多麼偉大的神.
對我們的恩典有多少.
多麼豐厚.
主讓我們生命.
讓這個恩典來改變.

$^{641}$幫助我們改變我們怎樣看我們的資源.
怎樣看我們自己的生命.
以至我們做一個忠心的管家.
讓我們教會同工.
弟兄姊妹.
我們所有人.
一起來去慰藉你的福音.
慰藉你的名義.
去繼續努力.
我們這樣祈禱.
馮主名.
\newpage



\section{}
\label{sec:hfmoR2RDphM}
\textbf{What Christ Thinks of the Church \#7 (Revelation 2-3) - Rev. Mark Chan}
\newline
\newline
連結: \href{https://youtube.com/watch?v=hfmoR2RDphM}{\texttt{https://youtube.com/watch?v=hfmoR2RDphM}} ~~~~ 語音日期: 2024-10-15
\newline
\newline
\hyperref[sec:vjkb_KdvSz0]{\small{< < < PREV SERMON < < <}}
~
\hyperref[sec:index_chronic]{\small{[返順時目]}}
~
\hyperref[sec:index_scriptual]{\small{[返順卷目]}}
~
\hyperref[sec:sCrh8yIfIpY]{\small{> > > NEXT SERMON > > >}}
\newline
\newline
$^{1}$All right, thanks, Kelly..
Well, today's marathon, you guys in the race?.
You wish, right?.
You wish..
I was talking to someone in the morning service, and.
he's a veteran..
He's actually ran, I think, five races in the past 10 years..
So he's a veteran..
He doesn't have to do the draw to to join the race..
So so I was I was I was looking at him..
He was here and I was talking to him and said,.
Why are you here?.
Was supposed to be a running..
Right. And he said, No, I'm not going to run today..
I mean, this year and he's running, you know, every every other year..
So we do have some runners..
We do have some runners in our congregation..
In any case, good to be with you..
And I know the traffic is kind of crazy, especially if you're coming.
from the downtown area and the South Loop area..
But, you know, it's great to be here..
We will go to the last letter of the seven letters to the seven churches..
So if you have the Bible, please turn with me to chapter three,.
chapter three of Revelation..
We are finally coming to the last letter..
Remember, there were seven letters..
Jesus sent letters to seven churches.
and from Ephesus, the first one and all the way to.
Lilladahshir today and Lilladahshir is kind of inland..
It's about 40 miles southwest of Philadelphia..
Remember, we talk about the latter to Philadelphia as well..
And also is 100 miles east of Ephesus..
So it's situated in a valley..
It's a very fertile valley..
And what is known for in Lilladahshir.
is that it's a traffic hub. It's kind of like Chicago, right?.
It's kind of like Chicago..
We are right at the center of all the traffic..
That's why we always have a traffic jam all the time..
And it's getting worse with all the road constructions and all that..

$^{41}$And back in the first century,.
there was a Roman highway going through from right to left.
and it passes through Lilladahshir..
So as you can see, it's a traffic hub..
It's also a banking hub, a financial hub..
And it's also famous for textile industry and medical schools..
You know, back then they have medical schools.
and they have fairly famous.
medicine for eye diseases..
If you have some eye diseases in those days, you go to Lilladahshir..
In those days, you go there to be a doctor..
You go there to earn money..
So as you can imagine, the city is very, very rich..
And all the people there are very rich..
And I'll tell you a story in the first century, about 1860..
So about 30 years after Jesus' death and resurrection..
And it's within the period of the early church..
And there was a massive earthquake.
happened around the region of Lilladahshir..
And if you remember, the letter of.
revelations were written in 1890, late first century..
So 30 years before this letter was written,.
this massive earthquake happened in Lilladahshir..
It was 1860..
And some of the towns were destroyed as well, like Sardis..
Can you see Sardis right there?.
Sardis, Philadelphia, Sardis..
Sardis was completely wiped out..
And Rome has to use the federal fund to help rebuild Sardis..
However, Lilladahshir is such a remarkable city, and.
they didn't need any funds from Rome..
And they rebuilt the city themselves..
It's pretty cool, right?.
In the past few weeks, we heard about the hurricane.
in southeast of United States, Florida..
Well, just imagine the towns are being just tragic.
and everything is wiped out and you need federal money.
to pump into the region and to rebuild everything and to clean up everything,.
to rebuild the lives of the people..
Lilladahshir back then said to the Roman government,.

$^{81}$saying to Rome, we don't need your federal money.
because we have enough money to rebuild after the earthquake. Right..
Pretty cool, right?.
And you can imagine the people in the church in Lilladahshir,.
they should have that kind of pride and financial resources..
And that go into the culture, that go into the church..
And I would say it is very true for our church as well..
Yes, we don't have a lot of people, but aren't we all pretty resourceful?.
I think we are..
In some way or the other..
So this letter is written to the church in Lilladahshir and also to us..
Let's read it through. It's a fun letter to read through..
Verse 14 and following..
To the angel of the church in Lilladahshir write..
Jesus said to the church in Lilladahshir, the Christian in Lilladahshir..
These are the words of the Amen..
The faithful, true witness, the ruler of God's creation..
I know your deeds and you are neither cold nor hot..
I wish you were either one or the other..
So because you are lukewarm, neither hot nor cold,.
I am about to speak you out, out of my mouth..
It's kind of interesting..
Jesus interceded himself as the Amen..
Well, you know, Amen..
What does it mean?.
Was pray and end with this cold, cold word..
Right. This cold word..
Amen. What does it mean?.
Truth. Right..
Sincere, truthfully..
Amen. Amen. I said to you, Jesus said to the people all the time in the gospels..
Amen. Amen..
I said to you and English translation is truly, truly..
I said to you, no, God does not lie..
God cannot lie..
And what God says must be true..
And it has to be true..
This is who God is..
Jesus said something about our lives, about the church in Laodicea,.
and he gave his evaluation and the evaluation has to be right,.

$^{121}$has to be truthful, has to be according to the facts..
And his evaluation is kind of interesting..
He said, you guys are lukewarm..
I would rather you guys are hot or cold,.
but you guys are lukewarm, you are neither cold nor hot..
That's interesting..
Sometimes we use this verse as as if, you know, God is wanting us to.
not to be lukewarm, we either be hot or cold..
And you think about it, it's very interesting if you use this verse this way,.
because, you know, aren't we required to increase our spiritual temperature,.
our passion for God, right?.
Our passion for God..
If if this verse is about our heart temperature for God, our passion for God,.
willing God prefer us to be, of course, hot, right?.
And of course, warm, warm is better than cold..
Hot is better than warm..
You see what I'm saying?.
You know, you take the temperature of the congregation and God is saying,.
hey, you know, you guys are lukewarm..
I just don't want you guys, I'll spit you out..
But if hot, meaning we have passion for God,.
then of course, God prefers that..
But if cold in that sense is being lack of passion for God,.
then that sounds kind of strange..
But the Christians in the little desire are lukewarm..
They're not too hot, not too cold..
Lukewarm, it's just like a glass of water..
Like a glass of water, you know, between 100 Fahrenheit to 60 Fahrenheit.
is kind of considered hot, right?.
I consider it lukewarm because the temperature,.
if raises above 100 Fahrenheit, it can burn you..
It can burn your tongue..
If it's below 60, it's kind of cold..
So room temperature is kind of like 60 to 100 Fahrenheit..
So we always take this word as the measurement of our spiritual temperature,.
our passion as Christians..
And God will be saying, you know, I do not prefer you to be lukewarm..
Now, if that's the case, I think it's kind of strange, right?.
As I said, wouldn't God prefer us to be warmer than cold?.
But this word says neither cold nor hot..

$^{161}$I would rather you guys, you know, be cold or be hot, you know,.
just don't be lukewarm..
So, you know, there has to be some other way to understand this verse..
And I think if we get this right, it applies to all of us..
Now, in order to understand this, we have to understand geography..
Geography..
Here on the map, you can see Heropolis..
Heropolis is on the left, upper left, right?.
Heropolis..
And then on the right is Colossae..
And then right in the center is Laodicea..
Now, what am I talking about?.
Now, since this passage is about water in Laodicea,.
so we have to take a look at the region, about the water supply..
Now, Heropolis on the upper left is about eight miles north of Laodicea..
It is a city and it's still a very, a city alive today.
because they are famous for hot mineral springs..
Hot mineral springs..
Just like you go to Yellowstone National Park,.
you see all kinds of, you know, hot springs..
Or you go to hot spring in the Arkansas,.
there's a place called Hot Springs because they are known for hot springs.
and people will enjoy it..
And you can see in the picture, it's still alive today..
And tourists will go there and just dump yourself into a hot spring.
and you just enjoy your time, right?.
Cool..
So that's the water supply in Heropolis..
In the east, on the right, is Colossae..
It is situated at the foot of a snow-capped mountain..
It's called Mount Katmos..
And because you are right there, Colossae,.
you are right below or right at the foothill of the snow-capped mountain,.
what do you get?.
You get cold water..
Therapeutic water..
And you can take a bath even..
And you can drink the pure, pure water..
It's like you get a bottle of water today,.
it says whatever mountain, whatever spring that it comes from,.

$^{201}$probably it's fake..
But you get pure water in Colossae..
So you got hot water on the left, upper left, hot spring..
And then you got cold water on the right..
So as you can see,.
the verse is referring to the other city..
God said to Laodicea, "I would rather you guys be hot or cold,.
but you guys are lukewarm.".
Let's take a look at the description of lukewarm water..
Laodicea had, it's known in history,.
Laodicea did not have reliable water supply..
Because the water has, by the time the water from the hot spring.
coming to Laodicea to the south,.
the temperature decreases, right?.
Hot water becomes lukewarm water..
And remember, if the water becomes lukewarm,.
you cannot use it as hot spring..
It's not hot spring anymore..
It's called lukewarm spring..
But it's not therapeutic anymore..
It's not for medicine anymore..
And the water from hot spring.
usually carries a lot of minerals..
And when that mineral rich water reaches Laodicea,.
it will clog up the water pipes..
It will clog up the water pipes..
Today, we call it hot water..
You see what I'm saying?.
It's common sense, right?.
Yeah, it's common sense..
That's why you need filter water..
You filter out the dirt or the mineral in the water.
so that you can drink it..
Now, that's what happened in Laodicea..
All they got is hot water, lukewarm water..
You can't drink it..
And it's bitter to drink it..
So you have to spit it out..
So now you got it..
Here are some of the pictures from Laodicea,.

$^{241}$and you can still see some of the broken pipes..
They have some pipes, you see the white stuff..
Around the pipes, those are from, you know,.
hardened, solidified minerals of the lukewarm water..
So if you're a plumber in Laodicea,.
you got good business..
Because every house needs your service every now and then..
Because you have to clean up the pipe..
You have to clean up all the plumbing system.
so that you get water coming into the city..
So they don't have fresh water..
So God is saying, "You are right in the middle.".
Right?.
I put a plumber there.
because it's gonna be bad to live in Laodicea..
You always have to drink from bottles of water..
You don't, you can't, you can't just, you know,.
turn on your tap and get some water to drink..
So much different from Chicago..
God is saying, "I would rather you guys be hot on the left..
Be hot spring, so you're functional..
You can bring benefit to other people..
Or to the right, the cold water, people can drink it.".
He's talking about not the temperature of the water..
It's the functionality, the functions or purpose of the water..
Right in the middle of Laodicea,.
the water has no use..
You have to spit it out..
It's totally no use..
He's referring, using this metaphor to talk about Christians in Laodicea..
And they are nominal Christians..
What does it mean, nominal Christians?.
That, okay, you profess that you are Christians,.
but you don't live like one..
You don't live out the gospel..
You don't have a living relationship with Christ..
You just go through the motion and you just don't care..
Right? You come here, okay, all right, then, you know, Monday to Saturday,.
I just care about myself..
You just don't care about your relationship with God..

$^{281}$You don't have a living relationship with God..
You don't have the life to show other people..
So you cannot bring benefits to other people..
God is saying, "I would rather be left or be right..
Be cold or be hot, so that people can use you..
God can use you to be his instruments, to be a witness to him.".
So we got that clear up, right?.
We got that clear up..
It's not talking about spiritual temperature..
It's talking about spiritual function..
It's malfunctioning Christians..
It's malfunctioning Christians in Laodicea..
You're not doing what you're supposed to do..
You're not living a life you're supposed to do..
Because when we believe in the gospel, we should have changed,.
at least fundamentally, our value system,.
our schedule, our life, our function,.
and our relationship with God and to other people..
Everything has to be reconfigured..
But in Laodicea, all those Christians are nominal Christians..
They are not functioning, malfunctioning church..
Now, as you can see, you know, what does Christ want in our lives?.
Christ wants, Christ desires a living relationship with you..
God desires most about your relationship with God himself, with Jesus..
In order to be fruitful, in order to be functional,.
you cannot be lukewarm..
God has to use you..
What is your function?.
Jesus said, "If you abide in me, if you have a living relationship,.
fellowship with me, then I can dwell in you..
If you love my word, if you abide in my word, then I can live my life through you.".
Christ desires, Jesus wants your whole heart worship,.
your service, your time, your talent, your treasure, your life,.
God's life, eternal life can live through you..
And you can do things out of love for God..
That's what Jesus desires..
He doesn't really look at how rich you are or how well off you are..
If you are not generous, you present no functions to other people..
You are not being used by God's eternal kingdom..
That's what Jesus wants..

$^{321}$He wants you to have a living relationship with him..
And secondly, this passage says, God decides to dwell in your heart..
Because without Christ dwelling in our heart, without Christ dwelling in our heart,.
we cannot live out the gospel..
This letter continues, 317..
Jesus said, "You say, I am rich..
I have acquired wealth and do not need a thing.".
But Jesus said, "You do not realize you are actually wretched..
You're wretched..
You're pitiful, poor, blind and naked..
I counsel you to buy from me gold with fire in the fire, so you can become rich..
And buy from me white clothes to wear, so that you can cover your shameful nakedness..
And then you can heal your eyes so that you can see.".
You see, this is all local languages..
Laodicea is famous and well known for all of this stuff..
Their gold, their treasure, their wealth, and their power..
The textile industry..
Their eye medicines..
That's a spirit of self-sufficiency..
That we think we can do life without God..
Ask questions..
I'll make money..
I'll save up money..
I don't need Jesus..
I mean, come on, Jesus is just a footnote to me..
The main content, I'm going to write my life..
I'm going to write every chapter..
That's my main content..
And where's Jesus?.
I mean, maybe it's just a little bit footnote..
Jesus said, "You are actually pitiful..
You're poor and you're blind and naked.".
319..
"Those whom I love," Jesus said, "I would bilk and discipline..
So be earnest and repent..
Here I am, I stand at the door and knock..
If anyone hears my voice and open the door, I will come in.
and dine with him and he with me.".
You see the language of fellowship here?.
Eating together in the ancient world is so important..

$^{361}$It's also important today..
It is even more important in those days..
In the ancient world, intimate fellowship, you know,.
dining together and sharing a meal together..
To have friendship..
To have a living relationship..
So we can eat together, we can talk together,.
we can have fellowship together..
You see, good friends, what do good friends do?.
They eat together..
But do you eat with your enemies?.
No, you don't..
You eat with your good friends and you talk to them.
for hours..
You can talk to your good friend for hours and hours.
and it feels like eternity..
That's what Jesus desires..
Jesus is knocking on your door, knocking on our door..
And Jesus said, "Please listen to my voice..
I want to renew my relationship with you..
I'm here all along, I'm knocking.".
And what's wrong with this picture?.
Jesus is being kicked out of the house..
But He is there knocking..
It's not because He doesn't have power to open your door..
He does have power to open the door..
Jesus is not saying, "Okay, I'm going to beg you,.
I'm going to love seek without you..
I have to come in and have fellowship with you.
and otherwise I feel like less of myself.".
No, Jesus is not like that..
Remember the first verse of this letter,.
Jesus said, "I am the Amen, I'm the ruler of the creation..
I'm the Cosmos guy, you know..
I create the cosmos..
I create the universe and I am in charge of the universe.".
However, Jesus is saying with gentle words,.
"I am standing outside and I'm knocking on your door..
I desire fellowship, relationship with you.".
Verse 21, "To him who overcomes,.

$^{401}$I'll give the right to sit with me on the throne..
Just as I overcome and sat down with my Father on His throne..
He who has ear, let him hear what the Spirit says to the churches.".
As we close and finish off this series,.
this is the seventh letter.
and this is the most, my beloved letter.
because it says so much about us..
It says so much about me as well..
Sometimes I just busy with my work.
and I forgot that what Jesus really desire.
is a moment with me.
so that I can read His word..
I can listen to His voice and have fellowship together..
When was the last time you really enjoy relationship with God?.
And you just spend hours and hours.
praying and nurturing the relationship,.
talking, talking to Christ, someone who loves us..
Just this letter is filled with warnings,.
is filled with builds..
There's no praises here..
But you know, don't misunderstand..
You know, God sometimes has some hard words for us..
It's not because He hates us..
It's not because He wants to condemn us..
Some of those hard words, tough words are tough love..
Jesus said, "I want to sit down with you.
to have a meal.".
Meal, quite no agenda..
It's like a close friend..
Intimacy..
We come before Christ today..
Let us turn and repent..
God looks at us now..
God looks at your heart right now..
He doesn't count you, how much you have accumulated.
or how much achievement for Him at this moment..
Jesus said, "I'm knocking on the door.".
Knocking on the door of your heart, open,.
and I'll come in and have intimacy, fellowship with you..
Let's all pray together..

$^{441}$Father, we confess that sometimes we treat our faith.
as just a concept of philosophy or idea..
No matter how good the idea sounds,.
that is not the point, that is not the core..
The core is coming before you.
and know that you love us.
and know that we respond with love..
I know you are knocking on our door every single day.
and we confess that sometimes we pretend.
that we don't hear it..
We just go about our lives without you.
and treat you as a footnote..
So we've confessed our sins..
We just confess our weaknesses..
And at this time, we just pray that your Spirit.
will speak to us and turn our hearts back to you.
so that we will enjoy our relationship with you,.
so that our lives will abide in your Word,.
our lives will abide in your love,.
that your life, your eternal life,.
the life of Christ, the life of the gospel,.
can live through us..
So Jesus, we open the door..
You come into our hearts.
and be our Lord, be our ruler,.
just as you are ruler of the entire universe..
So yes, we come before you.
and thank you for being our loving God..
In Christ, we pray..
Well, we're gonna close off with our verse of the month..
All right?.
Verse of the month..
I hope you have that memorized..
If not, we're gonna recite it together..
Now, this is a relevant promise.
out of the letter of 1 John..
And it's a very appropriate closing.
for the last letter of the seven churches..
1 John 1, 9..
Why don't we read together?.

$^{481}$"If we confess our sins,.
He is faithful and just.
and will forgive us our sins.
and purify us from all unrighteousness.".
If we confess our sins, He's faithful and just..
He is righteous and He will forgive us in Christ..
All of our sins and clean us up.
from all the unrighteousness,.
all the injustice, all the sinfulness..
And God is gonna live in our hearts..
Today, if you echo with my prayers,.
I do wanna urge you,.
let us spend more time with God..
Now, every day, you know, take 10 minutes, 15 minutes,.
you know, rather than, you know, scrolling through our phone.
and, you know, looking at TikTok videos,.
Instagram updates or whatever crazy news out there..
Know that God is in control..
Spend some time of intimacy with Him.
because God always desires your heart..
Right? Let your heart be open to Him..
So we do that..
And today, as we end our service,.
and we're gonna have some refreshment,.
and please stay, please stay for Bible study..
You know, don't treat Bible study.
as an intellectual exercise only..
Your faith needs the Word of God..
If you wanna grow,.
there's no way that you can just go through the motions..
You need the Word of God to dwell in your heart..
All right?.
So we do invite you to stay.
for about an hour of Bible study,.
and we're going to chapter two of the Book of Romans..
And I'll lead the English group,.
and Paul is gonna lead the Chinese group..
And we're gonna enjoy our time and have fellowship.
and abide in God's Word..
\newpage



\section{啟示錄 2:1-3:22}
\label{sec:sCrh8yIfIpY}
\textbf{基督眼中的教會 \#4 (啟示錄 2\_1-3\_22) - 陳崇基牧師}
\newline
\newline
連結: \href{https://youtube.com/watch?v=sCrh8yIfIpY}{\texttt{https://youtube.com/watch?v=sCrh8yIfIpY}} ~~~~ 語音日期: 2024-10-15
\newline
\newline
\hyperref[sec:hfmoR2RDphM]{\small{< < < PREV SERMON < < <}}
~
\hyperref[sec:index_chronic]{\small{[返順時目]}}
~
\hyperref[sec:index_scriptual]{\small{[返順卷目]}}
~
\hyperref[sec:yEZdtKfic4o]{\small{> > > NEXT SERMON > > >}}
\newline
\newline
啟示錄 2:1-3:22
\newline
\begin{longtable}{cl}
\hline
\hline
章節 & 經文 (和合本修訂版)\\
\hline
2:1 & \begin{tabularx}{0.7\textwidth}{X} 「你要寫信給以弗所教會的使者，說：『那右手拿著七顆星，在七個金燈臺中間行走的這樣說： \end{tabularx} \\ \\ \relax
2:2 & \begin{tabularx}{0.7\textwidth}{X} 我知道你的行為、勞碌、忍耐，也知道你不容忍惡人。你也曾察驗那自稱為使徒卻不是使徒的，看出他們是假的。 \end{tabularx} \\ \\ \relax
2:3 & \begin{tabularx}{0.7\textwidth}{X} 你能忍耐，曾為我的名勞苦而不困倦。 \end{tabularx} \\ \\ \relax
2:4 & \begin{tabularx}{0.7\textwidth}{X} 然而，有一件事我要責備你，就是你把起初的愛心拋棄了。 \end{tabularx} \\ \\ \relax
2:5 & \begin{tabularx}{0.7\textwidth}{X} 所以你要回想你是從哪裡墜落的，並且要悔改，做起初所做的工作。你若不悔改，我要到你那裡去，把你的燈臺從原處挪去。 \end{tabularx} \\ \\ \relax
2:6 & \begin{tabularx}{0.7\textwidth}{X} 然而你還有一件可取的事，就是你恨惡尼哥拉派的行為，這種行為也是我所恨惡的。 \end{tabularx} \\ \\ \relax
2:7 & \begin{tabularx}{0.7\textwidth}{X} 凡有耳朵的都應當聽聖靈向眾教會所說的話。得勝的，我必將神樂園中生命樹的果子賜給他吃。』」 \end{tabularx} \\ \\ \relax
2:8 & \begin{tabularx}{0.7\textwidth}{X} 「你要寫信給士每拿教會的使者，說：『那首先的、末後的，死過又活了的這樣說： \end{tabularx} \\ \\ \relax
2:9 & \begin{tabularx}{0.7\textwidth}{X} 我知道你的患難和貧窮—其實你卻是富足的，也知道那自稱是猶太人的所說毀謗的話，其實他們不是猶太人，而是撒但會堂的人。 \end{tabularx} \\ \\ \relax
2:10 & \begin{tabularx}{0.7\textwidth}{X} 你將要受的苦，你不用怕。看哪！魔鬼要把你們中間幾個人下在監裡，使你們受考驗，你們要遭受苦難十日。你務要至死忠心，我就賜給你那生命的冠冕。 \end{tabularx} \\ \\ \relax
2:11 & \begin{tabularx}{0.7\textwidth}{X} 凡有耳朵的都應當聽聖靈向眾教會所說的話。得勝的必不受第二次死的害。』」 \end{tabularx} \\ \\ \relax
2:12 & \begin{tabularx}{0.7\textwidth}{X} 「你要寫信給別迦摩教會的使者，說：『那有兩刃利劍的這樣說： \end{tabularx} \\ \\ \relax
2:13 & \begin{tabularx}{0.7\textwidth}{X} 我知道你的居所，就是有撒但座位之處；當我忠心的見證人安提帕在你們中間，在撒但所住的地方被殺之時，你還堅守我的名，沒有否認對我的信仰。 \end{tabularx} \\ \\ \relax
2:14 & \begin{tabularx}{0.7\textwidth}{X} 然而，有幾件事我要責備你，就是在你那裡有人服從了巴蘭的教訓；這巴蘭曾教唆巴勒將絆腳石放在以色列人面前，使他們吃祭過偶像之物，並且犯淫亂。 \end{tabularx} \\ \\ \relax
2:15 & \begin{tabularx}{0.7\textwidth}{X} 同樣，你那裡也有人服從了尼哥拉派的教訓。 \end{tabularx} \\ \\ \relax
2:16 & \begin{tabularx}{0.7\textwidth}{X} 所以，你當悔改；若不悔改，我很快就到你那裡來，用我口中的劍攻擊他們。 \end{tabularx} \\ \\ \relax
2:17 & \begin{tabularx}{0.7\textwidth}{X} 凡有耳朵的都應當聽聖靈向眾教會所說的話。得勝的，我必將那隱藏的嗎哪賜給他，並賜他一塊白石，石上寫著新的名字，除了那領受的以外，沒有人認識。』」 \end{tabularx} \\ \\ \relax
2:18 & \begin{tabularx}{0.7\textwidth}{X} 「你要寫信給推雅推喇教會的使者，說：『神的兒子，那位眼睛如火焰、雙腳像發亮的銅的這樣說： \end{tabularx} \\ \\ \relax
2:19 & \begin{tabularx}{0.7\textwidth}{X} 我知道你的行為：愛心、信心、勤勞、忍耐；又知道你末後所行的善事比起初所行的更多。 \end{tabularx} \\ \\ \relax
2:20 & \begin{tabularx}{0.7\textwidth}{X} 然而，有一件事我要責備你，就是你容忍那自稱是先知的婦人耶洗別教唆我的僕人，引誘他們犯淫亂，吃祭過偶像之物。 \end{tabularx} \\ \\ \relax
2:21 & \begin{tabularx}{0.7\textwidth}{X} 我曾給她悔改的機會，她卻不肯悔改她的淫行。 \end{tabularx} \\ \\ \relax
2:22 & \begin{tabularx}{0.7\textwidth}{X} 看吧，我要使她病倒在床上。那些與她犯姦淫的人若不悔改他們的行為，我也要使他們同受大患難。 \end{tabularx} \\ \\ \relax
2:23 & \begin{tabularx}{0.7\textwidth}{X} 我又要殺死她的兒女，眾教會就知道，我是那察看人肺腑心腸的，我要照你們的行為報應各人。 \end{tabularx} \\ \\ \relax
2:24 & \begin{tabularx}{0.7\textwidth}{X} 至於你們其餘的推雅推喇人，就是一切不隨從這教訓，不明白他們所謂撒但深奧之理的人，我告訴你們，我不會再把別的擔子放在你們身上。 \end{tabularx} \\ \\ \relax
2:25 & \begin{tabularx}{0.7\textwidth}{X} 你們只要持守那已經有的，直到我來。 \end{tabularx} \\ \\ \relax
2:26 & \begin{tabularx}{0.7\textwidth}{X} 那得勝又遵守我命令到底的，我要賜給他權柄制伏列國； \end{tabularx} \\ \\ \relax
2:27 & \begin{tabularx}{0.7\textwidth}{X} 他必用鐵杖管轄他們，如同打碎陶器， \end{tabularx} \\ \\ \relax
2:28 & \begin{tabularx}{0.7\textwidth}{X} 像我也從我父領受了權柄一樣。我又要把晨星賜給他。 \end{tabularx} \\ \\ \relax
2:29 & \begin{tabularx}{0.7\textwidth}{X} 凡有耳朵的都應當聽聖靈向眾教會所說的話。』」 \end{tabularx} \\ \\ \relax
3:1 & \begin{tabularx}{0.7\textwidth}{X} 「你要寫信給撒狄教會的使者，說：『那有神的七靈和七顆星的這樣說：我知道你的行為，就是名義上你是活的，實際上你是死的。 \end{tabularx} \\ \\ \relax
3:2 & \begin{tabularx}{0.7\textwidth}{X} 你要警醒，堅固那些剩下、快要死的，因為我發現你的行為，在我神面前沒有一樣是完全的。 \end{tabularx} \\ \\ \relax
3:3 & \begin{tabularx}{0.7\textwidth}{X} 所以，要記得你所領受和聽見的；要遵守，並要悔改。你若不警醒，我必如賊一樣來到；我幾時來到你那裡，你絕不會知道。 \end{tabularx} \\ \\ \relax
3:4 & \begin{tabularx}{0.7\textwidth}{X} 然而，在撒狄你還有幾位是未曾污穢自己衣服的，他們會穿白衣與我同行，因為他們是配穿的。 \end{tabularx} \\ \\ \relax
3:5 & \begin{tabularx}{0.7\textwidth}{X} 得勝的必這樣穿白衣，我也不從生命冊上塗去他的名；我要在我父面前，和我父的眾使者面前，宣認他的名。 \end{tabularx} \\ \\ \relax
3:6 & \begin{tabularx}{0.7\textwidth}{X} 凡有耳朵的都應當聽聖靈向眾教會所說的話。』」 \end{tabularx} \\ \\ \relax
3:7 & \begin{tabularx}{0.7\textwidth}{X} 「你要寫信給非拉鐵非教會的使者，說：『那神聖、真實的，拿著大衛的鑰匙，開了就沒有人能關，關了就沒有人能開的這樣說： \end{tabularx} \\ \\ \relax
3:8 & \begin{tabularx}{0.7\textwidth}{X} 我知道你的行為。看哪，我在你面前給你一個敞開的門，是沒有人能關的。我知道你有一點力量，也遵守我的道，沒有否認我的名。 \end{tabularx} \\ \\ \relax
3:9 & \begin{tabularx}{0.7\textwidth}{X} 那屬撒但會堂的，自稱是猶太人，其實不是猶太人，而是說謊話的，我要使他們來到你腳前下拜，使他們知道我已經愛你了。 \end{tabularx} \\ \\ \relax
3:10 & \begin{tabularx}{0.7\textwidth}{X} 因為你遵守了我堅忍的道，我也必在普天下人受試煉的時候保守你免受試煉。 \end{tabularx} \\ \\ \relax
3:11 & \begin{tabularx}{0.7\textwidth}{X} 我必快來，你要持守你所有的，免得人奪去你的冠冕。 \end{tabularx} \\ \\ \relax
3:12 & \begin{tabularx}{0.7\textwidth}{X} 得勝的，我要使他在我神的殿中作柱子，他必不再從那裡出去。我又要把我神的名和我神城的名—從天上我神那裡降下來的新耶路撒冷，和我的新名，都寫在他上面。 \end{tabularx} \\ \\ \relax
3:13 & \begin{tabularx}{0.7\textwidth}{X} 凡有耳朵的都應當聽聖靈向眾教會所說的話。』」 \end{tabularx} \\ \\ \relax
3:14 & \begin{tabularx}{0.7\textwidth}{X} 「你要寫信給老底嘉教會的使者，說：『那位阿們、誠信真實的見證者、神創造的根源這樣說： \end{tabularx} \\ \\ \relax
3:15 & \begin{tabularx}{0.7\textwidth}{X} 我知道你的行為，你也不冷也不熱；我巴不得你或冷或熱。 \end{tabularx} \\ \\ \relax
3:16 & \begin{tabularx}{0.7\textwidth}{X} 既然你如溫水，也不冷也不熱，我要從我口中把你吐出去。 \end{tabularx} \\ \\ \relax
3:17 & \begin{tabularx}{0.7\textwidth}{X} 你說：我是富足的，已經發了財，一樣都不缺，卻不知道你是困苦、可憐、貧窮、瞎眼、赤身的。 \end{tabularx} \\ \\ \relax
3:18 & \begin{tabularx}{0.7\textwidth}{X} 我勸你向我買從火中鍛鍊出來的金子，使你富足；又買白衣穿上，使你赤身的羞恥不露出來；又買眼藥抹你的眼睛，使你能看見。 \end{tabularx} \\ \\ \relax
3:19 & \begin{tabularx}{0.7\textwidth}{X} 凡我所疼愛的，我就責備管教。所以，你要發熱心，也要悔改。 \end{tabularx} \\ \\ \relax
3:20 & \begin{tabularx}{0.7\textwidth}{X} 看哪，我站在門外叩門，若有聽見我聲音而開門的，我要進到他那裡去，我與他，他與我一起吃飯。 \end{tabularx} \\ \\ \relax
3:21 & \begin{tabularx}{0.7\textwidth}{X} 得勝的，我要賜他在我寶座上與我同坐，就如我得了勝，在我父的寶座上與他同坐一般。 \end{tabularx} \\ \\ \relax
3:22 & \begin{tabularx}{0.7\textwidth}{X} 凡有耳朵的都應當聽聖靈向眾教會所說的話。』」 \end{tabularx} \\ \\
[1ex]
\hline
\hline
\end{longtable}
$^{1}$好,謝謝各位.
今天馬拉松我知道大家很難來.
特別住在南部的.
可能要跑出來才行.
很高興大家能夠一起來敬拜.
我們繼續七教會書信的系列.
我們打開到啟示錄第二章18節開始.
我們這次看到第四封信.
如果你記得我們看了三封信.
這三封信是寄給當時路線裡面頭三戶教會.
由以弗所到史美娜到別加摩.
這次我們看第四封.
就是推也推拉.
譯作推也推拉的一個城市.
所以就內錄了一點.
這個地方原來看背景的時候.
你會發現是一個製衣紡織業很興旺的地方.
就是出產布料,皮革,LV手袋.
打鐵的地方等等.
一個工業的城市.
當時是一個很重要的城鎮.
如果你看看《史蘭傳》的話.
有提過這個地方.
十六章第十四節.
那裡說到保羅和他的宣教隊伍.
離開了今天的土耳其.
然後他們過了對面.
今天所謂馬其頓的地方.
或者希臘半島上面那個地方.
在這裡說到保羅有什麼遭遇呢.
他說看到來到菲律賓.
有一個賣紫色布的婦人.
她的名字叫做呂底亞.
是推也推拉城的人.
提到推也推拉城的人.
所以呂底亞又離開了他本城.
然後他就去了菲律賓的地方.
他就碰見保羅.
然後他就聽見福音.
他在聽著保羅所說的福音.

$^{41}$主就開道他心.
使他留心來聽保羅所說的話.
然後他就信了主.
所以他變成了第一個.
所謂在歐洲信主的信徒.
然後說著說著.
到了十八章的時候.
又說到保羅其他的遭遇了.
說到他們.
他們是誰呢.
就是在本身羅馬京城住的一對夫婦.
叫做阿伯居拉和伯基拉.
他們離開了羅馬.
因為當時排遊的緣故被人趕走了.
他們就去了多倫多.
不是多倫多.
哥林多.
就去了那裡.
本市是什麼呢.
他們是製造帳篷的.
就是跟布料有關.
以前那些圓形的羅馬劇場是有遮蔭的.
有些帳篷的.
所以他們可能就是做這方面的.
製造帳篷為業的.
保羅就跟他們同業.
就跟他們同住.
一同來做工.
拜託.
保羅也是所謂織帳篷.
Tent maker.
可能他也是做這方面的工作的.
因為當時他可能支持不夠.
或者經濟不夠.
很多時候他就要做一些帳篷的生意.
或者帳篷的工作.
或者織帳篷.
或者去做一些布料的工作.
以致支持他繼續能夠將福音傳開.
說到推瓦推拉.

$^{81}$這些第四封信.
就是寫給在推瓦推拉裡面的人.
說什麼呢.
開頭那一節.
第二章第十八節.
我們說過這些書信都有格式的.
首先耶穌怎麼自稱自己.
一個死而復活的耶穌.
一個萬世之主.
他怎麼去介紹自己.
然後他就說我知道你們的情況.
你知道你們教會的情況.
然後他會作出不同的評價等等.
然後去勸告他們.
所以在這裡開始的時候.
第二章第十八節.
他說你要寫信給推瓦推拉教會的使者.
寫給整個教會.
他說上帝的兒子.
神的兒子.
那位眼睛如火焰.
雙腳如發亮的銅.
說這番話.
神的兒子.
我說過很多次.
是當時羅馬帝王的自稱.
他們自居為上帝的兒子.
或者son of god.
這裡當然基督這樣去介紹自己的時候.
他就說我才是son of god.
那位眼睛如火焰.
眼睛如火焰.
這個在先知書裡面.
可能會有這樣的印象.
給大家上帝是眼睛有火焰.
或者他有火焰.
就是說上帝通透的看到你.
有沒有照過X光.
比X光還要厲害.
上帝知道你想什麼.

$^{121}$上帝知道你裡面.
每一個的思念.
每一個的思想.
這個介紹.
包括他是眼睛如火焰.
雙腳如發亮的銅.
這些都是啟示錄第一章.
出現基督的時候.
他這樣介紹自己.
到這裡.
他說雙腳像發亮的銅.
就是他有沒有穿鞋.
是赤腳的.
是不是.
赤腳.
你知道以前.
那些羅馬的神像.
或者那些人像.
好像這張照片.
這個就是大名鼎鼎的奧古士都.
Augustus.
耶穌出生的時候就是他執政.
那些神像和羅馬王的神像.
很多時候他的腳.
是不穿鞋的.
為什麼.
為什麼不穿鞋.
因為.
在他們的神話裡面.
能夠踏入神聖靈域的.
就是不需要穿鞋.
因為你踏入神聖靈域.
就一定要脫鞋.
是不是.
有一個這樣的.
sorry.
有一個這樣的.
宗教上的理解.
基督.
說他是雙腳.

$^{161}$像發亮的銅.
不是因為他本身有鞋要脫鞋.
而是他本身.
就有神聖.
表明他就是在神聖的領域而來.
這位神的兒子.
是對教會說話.
不知道對教會說話會說什麼.
或者想想對我們教會會說什麼.
這封書信繼續.
他說19節.
他說知道你的行為.
愛心.
信心.
勤勞.
忍耐.
又知道你末後所行的善事.
比起初所行的更多.
這裡就是讚還是彈.
是讚的.
是讚.
居然推亞推拉教會.
耶穌見到的時候.
他說我讚你.
看見你們做的事.
四件事.
想想這四件事其實就是我們基督徒.
生活.
教會生活應該有的表現.
愛.
出於對上帝的愛.
對其他人的愛.
信心.
我們對上帝的信教.
對上帝的投資.
對上帝國度的投資.
我們信得過他.
信任過他.
所以我們跟他.
是不是.

$^{201}$服侍.
用我們的才幹.
用我們等等不同的東西.
推進上帝自己的工作.
你都有毅力.
是能夠經得起苦難的.
能夠捱得到苦的.
這四件事幾乎就是我們基督徒應該有的所謂果子.
Fruit.
或者應該我們有的Character.
對嗎.
如果我們做了基督徒這麼久的時候.
我們就檢查一下這四件事.
我們就想一下.
到底我們哪方面需要.
去學習.
推開教會.
他說這四件事.
我都看到.
並且知道什麼.
你現在走的東西.
比以往的更多.
就是說他們不斷有progress.
是嗎.
對嗎.
不斷有progress.
見完耶穌的時候.
他說我知道你們的行為.
我知道你們好的行為.
這裡就說你們要持守你們的好行為.
並且知道你們現在做的比以往的更多.
這是一個很好的教會.
是嗎.
但是我們教會都是這樣.
二十節.
然而有一件事.
我要責備你.
開始就要罵了.
並且這封書信.
其中一封比較長的書信.

$^{241}$說很多這個教會.
是需要被警告的東西.
他說有一件事我需要責備你.
就是你容忍那自稱為先知婦人.
耶洗別教唆我的僕人.
引誘他們犯淫亂.
以及吃濟過偶像之物.
這裡我有一點點拗頭.
這個阿水.
他說耶洗別.
耶洗別.
Jessieburg.
我們教會有些夫婦.
有時候他們生BB的時候.
就想和他們聊天.
改什麼名字.
因為最近都是兒子.
出生很多兒子.
都是想男孩子的名字.
有一次和對夫婦吃飯的時候.
就想起.
如果女兒有什麼名字.
大家在聊天.
我不知道什麼名字好.
有什麼名字不好的.
大家肯定知道.
特別是在講英文的世界.
你絕對不會改女兒的名字.
是耶洗別的.
Jessieburg.
Jessieburg好不好聽.
OK的.
其實蠻有型的.
但是.
在這個英語的世界裡面.
特別是基督教背景的世界.
這個名字是一個臭名.
臭名.
怪得不得了的名字.
你不要以為聖經是針對女人.

$^{281}$所以就把她罵壞.
不是這個意思.
而是有一段故事的.
一聽到耶洗別的時候.
對於一些比較熟悉聖經.
或者熟悉聖經的故事的人.
有時候聽到一些聖經的故事的人.
他們就聽到耶洗別是一個怎樣的人.
耶洗別是哪個時代的人.
當時的以色列國已經是分了南北國的.
北國就在上面.
就是紅色那一塊.
下面南國是小一點的.
就是黃色線那一塊.
下面南國就是猶大國.
它的首都就是耶路撒冷.
上面那一塊是很大的.
橙色那一塊.
紅色那一塊.
就是以色列.
因為頭十個支派就走上去了.
然後大家分裂了.
南北分裂的時代.
他們的京都就不是耶路撒冷.
因為猶大國的京都就是耶路撒冷.
北國的京都在哪裡.
就是薩馬利亞.
薩馬利亞.
當時發生什麼事呢.
特別是在背景下.
當時最重要的一個王.
就叫做阿哈王.
King Ahab.
阿哈王.
剛才說到改名字.
你也不會改你兒子的名字是Ahab.
因為Ahab是一個很臭的名字.
Ahab和Jezebel都是臭名字.
所以大家要知道.
如果上帝真的祝福你.

$^{321}$你又生多幾個的時候.
你問清楚不要亂改名字.
不要以為好聽的就去改.
發生什麼事呢.
在舊約列王裡面.
南國是猶大國.
北國就是以色列國.
在阿哈執政的時候.
就是紅色那一塊.
你要知道.
我們怎麼看地圖.
在以色列國.
你要賺錢的話.
你就要.
最多的錢是什麼地方.
就是沿海的那條路線.
藍色的那條路線.
還有內陸.
在今天的約旦王國.
Transjordanian Highland.
黃色那條線.
這兩條線都是當時的公路.
99日4日.
怎麼賺錢呢.
如果你陸路要賺錢的話.
你一定要.
起碼在那裡.
收錢.
在那裡經商.
然後世界各地不同的人.
由大馬士革.
中東不同的地方.
一直南下到哪裡.
下面.
埃及.
三角洲.
這就是沿海的路線.
這條公路如果你想要有錢的話.
你一定要霸佔那條.
的線路.

$^{361}$以至你能夠收錢.
對不對.
代代平安.
我記得我們以前.
芝加哥.
大概20年左右.
10多20年前.
還是要扔小錢的公路.
是不是.
你記得那些Tow Road.
你扔了一些錢進去.
然後又掉了.
然後又不知道撿好還是再給多一點.
你想在以色列賺錢.
你就是需要霸佔這兩條的道路.
一條是沿海的.
一條是約旦河東的.
有了這兩條線的時候.
就是.
給你國家能夠強大.
國家能夠強大.
為什麼能夠經商.
你看到紅色那一帶是很大一帶.
在阿哈時代.
其實.
你.
只是熟悉的看.
阿哈其實是一個.
最.
興旺的一段時間.
興旺到一個地步.
當時的帝國.
亞述帝國.
在他的石板.
記載他的歷史的時候.
亞述帝國打過來的時候.
他都記載阿哈王.
厲不厲害.
你想想白宮有沒有你的名字.
白宮.

$^{401}$那些總統.
大官.
會不會提到你的名字.
提都沒有提到你.
居然當時的帝國.
就提到阿哈.
Ahab.
和他老爸.
安威.
阿姆尼.
他們的王朝是一個很大的王朝.
當時以色列歷史裡面最鼎盛.
最興旺的時候.
為什麼.
因為他們一邊擴展向左.
一邊擴展.
是向右的.
他就打殘其他的民族.
然後霸佔了沿海的.
和約旦河東.
掌握住這兩個的脈絡.
所以有錢.
好不好.
國家非常之鼎盛.
我說你怎樣罵阿哈都好.
阿哈都是一個.
熟悉來說.
很好的.
有錢是好.
有更多錢好不好.
當然好.
誰會嫌錢多.
留意一下那個地圖.
就是那些沿海的路線.
你發現以色列國.
左邊.
地中海的海皮.
是很滑的.
看到嗎.
一條線很滑.

$^{441}$因為這麼滑的緣故就沒有什麼深的海灣.
漂亮的沙灘就很多.
所以你今天去Tel Aviv.
有很多漂亮的沙灘.
但當時漂亮的沙灘有什麼用.
它要深的海為什麼.
因為.
地中海的海洋的貿易.
是要靠這些大的海灣才能吸進來的.
你看到從下面一直上去.
去到加密山島.
都是很滑的.
這個一向都是以色列國的歷史.
他們賺錢最多就是賺綠路.
怎樣可以將那些海洋貿易的錢.
Suck it in.
怎樣可以拿進來.
你就需要什麼.
你需要深的海港.
就是今天的新加坡.
或者香港.
一些深的海港.
這些深的海港在哪裡.
在以色列的境外.
你看到上面藍色那塊.
長長的那塊.
那裡就是所謂西頓的地方.
就是菲利基的地方.
或者推羅等等的地方.
這些都是當時的新加坡.
當時的大貨運的碼頭.
你要賺錢的話.
你就要跟他結交.
耶穌在哪裡來的.
耶穌就是西頓王的公主.
所以他就跟耶穌.
西頓王的公主.
來結交.
耶洗別漂亮嗎.
我不知道.

$^{481}$因為傳統裡面覺得耶洗別很漂亮.
因為聖經裡面有一節說.
他很會打扮.
所以你看到有一幅畫.
漂亮嗎.
不知道他漂亮不漂亮.
漂亮不漂亮也好.
OK.
總之就要找耶洗別.
然後阿哈跟耶洗別就結親.
結親的話.
大家兩國就有邦交.
有Diplomacy.
我們就能夠將錢拿進來.
將那些貨櫃碼頭的錢.
能夠Suck it in.
所以不是因為他漂亮的緣故.
而是因為經濟的緣故.
要賺更多的錢.
能夠囤積更多的錢.
能夠更加富強.
所以為什麼要跟耶洗別.
西頓的公主來結交.
所以為什麼將耶洗別帶進來以色列.
帶進來宮廷做最大的皇后.
因為這樣的緣故.
阿哈的國度就會更加富庶.
更加多錢.
但是帶進來的時候.
又帶了他的宗教進來.
因為西頓這個地方.
剛剛完海皮.
拜的是什麼神呢.
是巴力.
巴力是誰呢.
巴力這個人.
其實巴力這個字.
Bale.
或者Bale.
其實就是主的意思.

$^{521}$主宰的意思.
為什麼叫他主宰呢.
左邊那個.
就是在芝加哥的大學.
其中一個博物館裡面.
你能夠看到的.
你走進去巴勒斯坦地.
他們找到的東西.
其中一塊就是這一塊.
這個就是巴力神像.
巴力神像大概有多大呢.
一個今天iPhone那麼大.
那麼長.
為什麼人們要拜巴力呢.
因為巴力是雷神.
你需要有雨水.
滋潤大地.
你才能夠有富足.
是不是.
巴力的神話就是說.
他就是雷神.
他就是打雷的那個.
他就是最大聲的那個.
因為打雷就是很大聲.
所以他們就去拜巴力.
看到巴力的右手是向上.
他準備去做什麼.
那裡有一個洞.
已經沒有了洞裡面的東西.
其實應該是雷電.
Light Saber.
他一打雷的時候.
就會有雨水下來.
你將巴力帶進來的時候.
有什麼好處.
將錢帶進來.
你和他結交的時候.
你就將耶和華的信仰混雜.
所以以色列國應該是拜真神耶和華的.
Yahweh.

$^{561}$He is the only God.
但是將巴力帶進來.
以致他們能夠有更多錢.
阿哈加上夜使別.
就成為當時舊約歷史裡面.
最富強的時候.
Make Israel Great Again.
OK.
哇 這個是一個害死當時的人.
富強到一個地步.
我剛才講過.
其他帝國都會說.
阿哈是很強勁.
為什麼他們需要這樣做.
只要我們富強.
我們不會被人欺負.
我們團積更多的東西.
我們做更多的事.
我們就能夠更多享受.
就不會有人怕我們.
所以為什麼以利亞會在加密山.
和巴力的先知來決鬥.
你剛剛看到加密山那個地方.
就是這兩個大本營的中間的地方.
雖然提到夜使別的時候.
提到阿哈的時候.
提到物質.
只要我們賺更多的錢.
有人說這個是拜物教.
或者可以說是拜錢教.
第一世紀的推也推拉教會.
就不是真的有一個夜使別復活.
而是有夜使別的教導.
他說先知的婦人.
夜使別教唆我的僕人.
無論她是男還是女.
怎樣都好.
總之有夜使別的主義出現.
夜使別的思想出現.
其實很簡單.

$^{601}$你這樣解釋的時候.
就是錢越多.
當然好.
是不是.
團積更多的東西.
當然好.
吃偶像的祭物.
你不要依書直說.
今天說拜過神的雞我們不能吃.
不是這個意思.
而是當時有他的背景.
為什麼要去廟裡面去拿這些肉.
因為要做生意.
因為要去找更多東西.
有沒有想過其實基督徒的生命裡面.
其實一個最大最大的試探就是物質.
特別我們在北美自由世界裡面.
物質對我們的捆綁是最大的.
也是最subtle.
最難辨認的.
有人說.
什麼時候我們會被物質捆綁得太厲害.
你問一下以下的說話就知道.
問一下以下的問題.
第一個問題.
你有沒有強迫性的購物慾.
是你不需要的東西.
但是你很強迫性的要去買.
明明是want來的.
你覺得是need.
生命裡面的滿足.
第二個問題.
生命裡面真的不是滿足是知足.
是取決於你團積多少錢.
很多時候團積的錢是不是真的需要.
第三個問題.
第四個問題.
對你和我來說慷慨是不是很難.
第五個問題.
問這個問題.

$^{641}$你就知道錢銀對我們的捆綁有多大.
第五個問題.
就是當牧師講錢的時候.
你會不會覺得不舒服.
我不是說大家.
我自己也是.
我記得好幾年前.
我有一個感動去奉獻給一個機構.
我很有感動.
我覺得這個機構做的事情太好了.
所以我心裡面就想.
不如我奉獻一千元.
奉獻一千元對我來說是多少.
但是我慢慢想的時候.
不如七百好了.
到我過獎的時候.
變成六百元.
那時候我突然間.
當頭棒學到底.
有什麼是在控制著我.
有什麼在控制著我.
我團積這些東西做什麼.
為什麼我不可以用我的資源去祝福其他人.
所以我同樣有這樣的掙扎.
同樣有這樣的掙扎.
每一個都是.
我記得好多年前.
我在加州的時候.
我碰見一個弟兄.
我就和他吃飯.
這個弟兄是學生的時候認識他.
現在已經在做生意.
很成功.
以前當然是貧苦學生.
我請他吃東西等等.
現在他不是了.
他在香港大陸或者不同的地方.
他已經發財了.
發財了.
發財得很厲害.

$^{681}$我就和他吃飯的時候.
我就說.
我和他說.
首先不是你給錢和我給錢的問題.
應該是他給錢.
你們和牧師吃飯應該是你給錢.
無論如何.
那個不算.
那個不和他算.
我和他說.
其實有很多科林士工.
我知道這些科林士工是需要的.
他就說是.
我說是.
很有道理.
他這樣和我說.
他說是.
我也應該要做回我20年沒有做的奉獻.
20年沒有做奉獻.
我鼓勵他.
你應該去到慷慨一點.
對不對.
對嗎.
你猜他走了之後.
他有沒有奉獻.
一分錢都沒有.
那頓飯還是我給的.
聖經裡面說財主.
我們經常想財主就是伊朗盟主.
Bill Gates.
有人說伊朗盟主.
Trillionaire.
不是Billionaire.
Trillionaire.
他們才是財主.
我們不是.
真正財主是什麼.
財主我告訴你.
聖經耶穌裡面所說的.
是以財為主的生命.

$^{721}$以財為我們主宰的生命.
所以我們問到這件事的時候.
以財為主的生命.
耶穌真的說得很對.
他說財主.
是很難進天國.
真的.
因為骨子裡面.
我們根本不相信上帝.
我們經常覺得上帝左頭左勢.
我們不相信天國.
我們不相信將來的東西.
我們只相信現在.
這個就是耶洗別和阿哈.
有的思想.
他說推亞特拉教會.
Yes.
你有愛心有信心.
有毅力有忍耐.
但是.
你們容許了.
耶洗別的教導.
你們需要悔改.
我們結束的時候.
我們不如就這樣看了這封書信.
一路看下去.
因為應該很淺白.
這封書信想說什麼.
二十一節.
我曾經有他悔改的機會.
耶洗別悔改的機會.
他卻不肯悔改他的人民.
看啊 我要使他病倒在床上.
那些與他犯姦淫的人.
若不悔改他們的行為.
我也要使他們同壽大患.
我也要殺死他的兒女.
宗教會就知道.
這個真是很嚴厲.
我是那察看人肺腑心腸的.

$^{761}$我要照你們的行為報應國人.
至於你們其餘的推亞特拉人.
一切不順從耶洗別教訓的.
不明白他所說的撒旦森澳之理的.
我告訴你們.
我不會再把別的擔子放在你們身上.
你們只要持守那已經有的.
直等到我回來.
那德性又遵守我命令到底的.
我就賜給他權柄制服列國.
他必用鐵杖管轄他們.
如同打碎陶器.
將我也從我父領受的權柄一樣.
我又會賜給他神聲.
凡有耳聽的.
就應當聽聖靈向教會所說的話.
讓我們聽聖靈對我們說的話.
讓我們以不以其他為主.
以耶穌為我們真正的主.
我們同心得來祈禱.
上主我們求你給我們悔改的心.
讓我們回絕.
讓我們審視一下我們怎樣對物質.
怎樣對我們生命每一個的角落.
主確實有很多地方我們藏起來.
不願意去鬥.
不願意去搞.
主求你讓我們見到天國的真實.
讓我們回歸.
主讓我們持守我們僅有的.
幫助我們忠心.
幫助我們去到有信心有愛心的.
活出我們基督徒教會的生命.
主當我們這樣想的時候.
我們想起我們自己的角色.
當我們離開這個地方的時候.
我們自己個人的生命.
我們怎樣去跟其他人相處.
跟我們家人相處.
無論是夫婦.

$^{801}$或者我們跟我們父母.
跟我們家人.
我們自己個人的生命.
怎樣能夠使更多成為我們的主.
在我們工作.
在我們的職場.
在我們社會不同的地方.
主求你幫助我們.
以耶穌基督.
為救我們愛我們.
應許給我們這麼好的主宰.
讓我們生活.
與這個福音相稱.
我們祈禱奉傳美教.
願神祝福各位.
\newpage



\section{24:30-34}
\label{sec:yEZdtKfic4o}
\textbf{Way of Wisdom \#1 (Proverbs 24\_30-34) - Rev. Mark Chan}
\newline
\newline
連結: \href{https://youtube.com/watch?v=yEZdtKfic4o}{\texttt{https://youtube.com/watch?v=yEZdtKfic4o}} ~~~~ 語音日期: 2024-10-22
\newline
\newline
\hyperref[sec:sCrh8yIfIpY]{\small{< < < PREV SERMON < < <}}
~
\hyperref[sec:index_chronic]{\small{[返順時目]}}
~
\hyperref[sec:index_scriptual]{\small{[返順卷目]}}
~
\hyperref[sec:VjJ7v3bZQ7A]{\small{> > > NEXT SERMON > > >}}
\newline
\newline
$^{1}$- Hey guys, let's turn to the Bible.
and we will start a mini series called Way of Wisdom..
Way of Wisdom..
Okay, so my first question for you is that.
when you think of wisdom or being wise,.
what does it mean?.
What does it mean to be wise?.
What does it mean to have wisdom?.
And who can you think of someone,.
a name, a face, someone who is a wise person?.
Can you think of anyone like that in your life?.
Someone who is wise, someone who has,.
you know, embodied wisdom in their life?.
Now, does wisdom, does wisdom mean intelligence?.
Just intelligence?.
Or, you know, you get a 4.0 GPA, straight A student..
Does it mean you have wisdom?.
Can your paycheck number on your paycheck.
will tell you whether you have wisdom or not?.
Some of the questions we need to ask,.
what does it mean to have wisdom?.
Especially in our generation, you know,.
you listen to all sorts of different discourses out there,.
political discourse, public discourse,.
private conversations or conversations in our relationships..
Can we pick up some wisdom?.
You know, this generation, we need wisdom..
Especially men, all right?.
I know I'm old school, but, you know,.
we need to embody wisdom..
What does it mean to live a life of wisdom,.
a way of wisdom?.
Many years ago, maybe I'm dating now,.
many, many years ago, back in the 90s,.
back in the 90s, you know, as I told you, I'm getting old..
Back in the 90s, there was a new label or new topic.
that everyone was talking about, mid 90s..
Emotional intelligence, emotional intelligence..
That was a term, emotional intelligence.
or emotional quotient, EQ,.

$^{41}$coined by a scholar named Daniel Goldman..
And he published his studies.
because he followed some of the best students.
out of Harvard, out of MIT,.
followed their career path.
and followed how they managed their life..
And he finds that many of those people have just high IQ,.
does not mean they live successfully,.
does not mean they work with wisdom,.
does not mean they manage their life.
and their career successfully or wisely..
So he publishes his work called "Emotional Intelligence,".
"Emotional Intelligence,".
and the subtitle of the book says it all..
So the main title of the book is "Emotional Intelligence.".
The subtitle is "Why It Can Matter More Than IQ.".
So why EQ matters more than IQ?.
So in his study,.
he finds that there were five essentials for human life..
How we manage our life,.
how we manage our life in the sense.
that we will be so-called successful..
So he named five essential elements of EQ..
First, listen up, first, is self-awareness, self-awareness,.
self-awareness, that's the first thing..
Second, self-regulation,.
how you regulate yourself,.
you know, how you discipline yourself..
Number three, motivation,.
how you have motivation at work, in life in general..
Number four, empathy, empathy,.
how you relate to other people,.
how you make other people feel,.
how you connect with other people, empathy..
Number five is social skills,.
how you relate to one another,.
you know, how you build friendship,.
you know, trust between people..
So number one, five elements of EQ,.
Daniel Coleman said, the first one is what?.

$^{81}$Self-awareness, self-awareness,.
that you are aware of your problem,.
you're aware of who you are.
and you know how to carry yourself, self-awareness..
Second is self-regulation, how you discipline yourself,.
how you live your life, regulate your time..
Number three is motivation,.
you find yourself getting motivated this morning.
when you get up, so tired,.
what's wrong with you?.
There's something wrong with us..
Self-motivation, number three..
Four, empathy, right, empathy..
Number five, social skills..
These five, he says, is the five big things.
to manage your life, so-called emotional intelligence..
So it doesn't have to do with IQ,.
whether you get a PhD or get a professional degree.
or you are very, very educated, it's all here, right?.
The most important things that Coleman said.
is your emotional intelligence..
Speaking of that, that is pretty much similar.
to what the Bible has to say about wisdom..
Now the Bible says a lot about wisdom..
If you want to be a wise person,.
if you want to be a wise, for me, being a wise father,.
wise husband, spouse, wise friend,.
worker, neighbor, all of this,.
the Bible has a lot to say about wisdom..
Has a lot to say about how we live our life wisely..
You know, when the Bible uses the term wisdom,.
it usually means that something related to practical skills..
Practical skills..
You are skillful in certain things..
That means you are wise in certain things..
So Exodus 36, there were a couple of verses,.
it says there were a couple of people,.
Pastor Al and Oholiab, Oholiab,.
anyways, how you pronounce it,.
it doesn't matter, it's Hebrew, anyway..

$^{121}$Oholiab, and every skill-wise person.
to whom the Lord has given skill and ability.
to know how to carry out all the work.
constructing the sanctuary, the tabernacle,.
to do the work just as the Lord has commanded..
Now, what's going on here?.
The Israelites were in wilderness,.
and God said, "You guys need to build a home for me.
"so I can dwell among you, I can live among you.".
And that is tabernacle, or translated sanctuary here..
It's a portable temple so that God will lead his people.
in the midst of them, living in the midst of them,.
walking in the midst of them, to relate to them..
So it has to do with worship,.
it has to do with relating to God in the wilderness..
And so the sanctuary has to be built,.
and the tabernacle has to be built..
You know, tabernacle has different items,.
different furniture items, like manure,.
and the show table, and the fabric,.
the poles, and all of this stuff,.
those are practical things that you need.
to set up the temple, the portable temple tabernacle,.
in the midst..
So God has chosen some of the people here.
who are wise, skilled in certain things..
So probably they can do metalworking,.
they can do carpentry,.
so that they can set up the right, proper worship among them..
So God needs wise people to build his church, right,.
to build his presence, to build up the tabernacle.
in those days when the Israelites were in the wilderness..
So the Bible has a lot to say about wisdom..
And one of the books in the Bible.
that has to do with wisdom,.
has a lot to do with wisdom, is the book of Proverbs,.
found in the middle of the Bible..
And I wanna emphasize something,.
that when we read again this book,.
we try to understand what the world is like in this book..

$^{161}$Now, if you read the book from verse one to the last verse,.
there are 31 chapters,.
you find that there is some sort of world structure.
presented by the book..
And you find something really clear in this book,.
in the book of Proverbs,.
is that there is some sort of moral structure,.
some sort of spiritual structure in this world..
In general, what it means is you do something,.
and then there's consequence..
Your actions always have consequences..
If you do good, in general, you will receive reward..
If you do badly, then you will be penalized,.
you will suffer..
So this is kind of the general truth, it's not always true..
It's not universal truth, it's not promises..
But in the book of Proverbs,.
you find a lot of this has to do with the order.
of the universe that God has arranged,.
so that in this moral order, in this spiritual order,.
we have to behave in a certain way..
Give you an example here, chapter 11, verses 24 to 25..
It says, "A man gives freely, yet gains even more..
"And other men withholds.".
Doesn't want to give, but comes to poverty..
25, "A generous man, alive,.
"living in generosity, will prosper..
"He who refreshes others, give away for others,.
"and help others, will himself be refreshed.".
You see the cause and effect here..
The spiritual order, or the cosmic order,.
if you honor God with your wealth, you will increase..
Even your heart will increase..
You will get a sense of satisfaction..
You will feel good..
If you give freely to others, to bless other people,.
you, yourself, will be blessed..
If you serve other people, it's not because.
you're giving up all the time..
You will get something back, and your life will be blessed..

$^{201}$So that is the world of Proverbs..
There is always consequence in the order..
Just like law of gravity..
You drop something, it will break..
You have a glass of water, you drop it on the floor,.
on the ground, it will break..
The glass will shatter..
That's consequence to our actions..
So that's what the life of wisdom is..
So the book of Proverbs is encouraging us.
to observe the moral universe and the spiritual order.
that God has set up in our world,.
so that we can live accordingly to those principles..
Now, one of my favorite passages.
in the book of Proverbs is this..
Let me read that first..
Chapter 24..
Now, here the wise teacher,.
the wise man, like the teacher in the book of Proverbs,.
he did something, and then he observed something,.
and then from his observation,.
he gets some spiritual principles out of it..
He said, "I went past the field of this lucker,.
"past the vineyard of the man who lacks judgment..
"Thorns had come up everywhere,.
"and the ground was covered with weeds,.
"and stone walls was in ruin.".
Now, look what the wise is doing..
He opened his eyes, and he observed.
what is going on in the world..
He passed through a farm of a lazy person,.
a lazy farmer, sluggard,.
because he does not do anything to the farm..
He does not put in his effort.
to cultivate the ground, the soil..
So he passed through, and he observed..
He observed something because of his lack of judgment,.
because the farmer is being lazy,.
he's not doing his work,.
and then you got weeds,.

$^{241}$and you got all sorts of problems for your soil,.
and your farm is in ruin,.
and stone wall is also in ruin..
Now, what he observes here.
is that there is some sort of arrangement.
in this world..
There's some sort of spiritual order in this world..
If you do something, there's consequence..
If you don't do something you're supposed to do,.
there's also consequence..
If you're lazy, you're not putting effort in your life,.
in your farm, in this case,.
you know, there's always true.
that your farm will go into ruin..
Your house, your stone wall,.
they will always fall..
You know, the wall of stones,.
you stack stones and bricks on top of it.
so that it will protect you..
If you don't maintain it,.
if you don't put in energy,.
if you're not being proactive,.
the bricks or the stones,.
they will inevitably fall..
Okay?.
The fallen stone will not go back to the wall.
unless you pick it up and put it back on..
But if you do not maintain your life,.
if you don't maintain your stone wall,.
the stone will drop..
Law of gravity, right?.
Well, it doesn't take Newton to understand it..
So the wise is trying to help us to observe.
the life of the lazy person..
I mean, look at your bedroom..
If you don't put energy to it,.
look at your house,.
it will be in chaos..
There's always a cause and effect in this world..
And what the wise is trying to understand is,.

$^{281}$again, it's the spiritual order of things..
Life will always deteriorate..
Life always deteriorate..
And life naturally descends into chaos..
Just like stone wall will fail..
It will fall..
Now, what he observed is that.
once you have order,.
stone wall, your farm, good soil,.
order in your house,.
it will eventually lead to disorder..
Order always goes to disorder..
Order always goes to chaos..
It will never happen naturally from chaos to order..
It always happens that order goes to disorder..
Unless you put energy into your life.
and put things back to order..
Repair the broken wall..
Pick up the stone and bricks,.
put it back on the wall..
You see?.
32..
Verse 32, he says,.
"I apply my heart to what I observe.
"and learn a lesson from what I saw..
"A little sleep, a little slumber,.
"a little folding of hands to rest..
"Poverty will come on you like a bandit,.
"and scarcity like an armed soldier..
"You will be attacked.".
You know?.
I don't know if you understand it.
or if you can observe or agree with the wise..
Our life is always going to the direction of chaos..
Unless we're proactive enough.
and we are active enough to make corrections,.
to maintain it..
Your finance..
If you just have money sit there,.
you don't do anything about it,.

$^{321}$it will always be devalued, right?.
Your house, your bedroom,.
your schedule,.
your calendar..
If you don't do anything, folding of hands,.
if you're lazy, if you don't observe and you don't learn,.
it will always go to chaos..
Anyone remember when you were in college,.
you learn second law of thermodynamics?.
Ring any bell?.
Second law of thermodynamics?.
Okay, no, right?.
So you guys didn't go to chemistry class..
Second law of thermodynamics says in this universe,.
things always go to chaos, right?.
The entropy in our lives always increases..
Unless you put energy into it and correct it.
and bring it back to order..
That's what the wise is trying to teach us..
Now, as you can tell from the book of Proverbs,.
it's not like a prophetical books..
It's not like other books in the Bible..
Like the prophets, the prophets will come to you.
and say directly, this is what the Lord has said..
Thus says the Lord, you have to do this,.
thou shall do this or thou shall not do this..
It's a direct command..
But the wise in the book of Proverbs.
is that he uses his observation,.
he'll open his eyes and observe other people,.
observe life in general, and then learn something..
So it's not from top down, it's from bottom up..
That's what wisdom will take you..
Observe your spiritual life..
Now, I've been a Christian for decades..
If I just being lazy and I don't care.
about my relationship with God,.
things will always go to chaos..
You have less strength to withstand temptations..
You have less faith, less confidence to withstand issues,.

$^{361}$to deal with difficulties in your life..
It is always harder to worship God.
unless you train yourself and strengthen yourself..
It's kind of a physical exercise, right?.
How many of you can do a seven minute plank?.
How many of you can do it?.
Probably only a few people can do it, right?.
Well, how come people can do a longer plank?.
Because they train themselves..
They strengthen themselves so that the muscles are going.
so that they can do what they want to do..
Our spiritual muscle always needs work..
You go to the gym, you train yourself,.
you fight against chaos,.
you try to tame the chaos in your life..
Every single morning, you try to deal with all this disorder.
in your work, in your relationship, in your marriage,.
in your family, in your neighborhood..
You always try to rely on God and you train yourself.
so that your life can withstand trials and temptations..
They will come at you whether you ask for it or not..
Difficulties, challenges, issues.
will always come to threaten your life.
and will pull you apart and separate you from God..
So every single morning, you have to have the motivation.
to fight the chaos, be a man,.
put in energy and put things back to order..
So after the sermon, you probably need to go back home,.
put things back to order, right?.
So long you're just folding your hand.
and you don't care much..
After the sermon, you go back to your life,.
your personal life, you know, read the Bible,.
read the Word of God, just meditate on the Word of God..
Unless you do that, you don't have the Word of God,.
you don't have the strength of God inside of you..
And then you will be defeated.
because chaos will take over..
Unless you do that, what the wise is trying to do.
and understand if we fold our hands to rest,.

$^{401}$poverty, spiritual poverty will come upon you..
And that is your responsibility..
That's not the church's responsibility..
That's not your mentors' responsibility..
That is your role..
You need to take charge..
You need to take leadership..
That is the way of wisdom..
Now there's one big question I always meditate on.
is that many of these wise sayings.
are from someone called Solomon..
Solomon has the reputation of being the wisest king.
of all the earth..
You remember that, the book of Kings,.
when Solomon was being introduced?.
He was so wise and he has the wisdom.
to bring righteous judgment to his people..
He can rule his kingdom with righteousness and justice..
There's so many sayings that come from King Solomon,.
but I don't know why Solomon failed at the end..
He failed miserably at the end..
Yeah, he had the glory,.
but at the end of the last chapter of his life,.
everything just spiraled into chaos.
and the kingdom is gone, is in ruin..
And the nation of Israel was divided into north and south..
How did that happen to the wisest king on earth?.
And then one day I was searching some of the Proverbs,.
some of the passages, and I find this really interesting..
And it's really, really a great reminder for all of us..
Why did the wisest person, the wisest king, Solomon,.
fail at the end?.
1927..
Proverbs 1927..
Stop listening to instruction, my son..
You will stray..
Not you may stray, not possibly you will stray,.
but you will definitely stray from the words of knowledge,.
the knowledge of God, the knowledge,.
the wisdom to walk in the way of God..

$^{441}$Stop listening..
One day you find yourself, you check yourself,.
check your heart..
Dude, I don't need to listen to this..
Dude, I don't need to study this..
I don't need to find out and learn something more.
from the word of God..
You stop right there, you stop listening to instructions,.
spiritual instruction, my son,.
you will definitely stray away..
Now I have seen a lot of Christians in my life.
who did just that..
When they were younger, happier,.
and they were really, really solid in their faith..
But I just scratching my head, what is going on with them?.
What is going on with them?.
There has to be something wrong..
And you just hear the person complain about this and that,.
this and that, you know?.
But this verse says, there's a heart issue here..
There's something deep inside, it's the inner issue..
Inner issue of stop listening..
And the Bible says it, adultery..
Because you always listen to someone..
If you don't listen to God, you listen to the counterfeit God,.
the false gods, the culture, the world, or even yourself..
And you deceive yourself..
And then you spiral down into spiritual chaos.
because of this..
You stop listening, my son, you will definitely stray..
You will go astray from spiritual knowledge..
What a reminder to all of us..
I remember in the beginning of this year, in January,.
I preached on a series called Atomic Wisdom..
Remember that?.
It's been so long..
It's not that long, guys, it's not that long..
It's just a few months ago..
I preached on this series, Atomic Wisdom,.
and we went into the Bible and we said,.

$^{481}$there are small atomic habits.
that will bring you huge effect if you do that..
Just a little bit small good habit,.
it will bring you fruitful life..
And I just wanna remind you and challenge you,.
since the beginning of this year,.
what has changed in your habits?.
What has changed in your life?.
Are there any small habits or good habits.
you have already built up?.
The foundation of small habits go a long way, guys..
And what's standing in the way,.
sometimes we say, is life is too busy..
Indeed, life is very busy..
Who is not busy here?.
Everyone is busy..
So often we say, because life is too busy..
Come on, God, just give me a break, all right?.
When I have a vacation, then I'll read the Bible..
When I have spare time, dude, I'll read the Bible..
Now, sometimes we say that to ourselves.
because we wanna deceive ourselves..
We say, life is busy because, no,.
this guy's reading the Bible more because he has more time..
I have less time..
But seriously, it's just self-deception..
Everyone here has 24 hours a day..
We say that as an excuse, but deep down,.
it's probably pride..
Probably we are too proud of ourselves..
You know what pride is?.
Pride is not being mean..
Pride is self-satisfaction..
I'm all good by myself..
I don't need the other voice.
to tell me what to do..
I don't need God to tell me what to do..
I don't need God to define me..
I don't need God to tell me what to do.
and how do I see myself, my identity,.

$^{521}$and my future, everything else..
I don't need God..
I don't want to learn..
I don't want Him in my life..
The most important thing is me..
The most important thing, the defining factor of my life.
is just me and now..
And that is pride in the Bible..
That is stop listening to instruction..
Stop having other people to speak into my life..
I just shut down, you know..
I just, please shut up, guys..
You know, I don't need to listen to you..
Please shut up, God..
I don't need to listen to you..
And that is the way of foolishness.
because you will stray from the word of spiritual knowledge..
I don't care what you did in the past,.
but if you stop listening, then you will go astray..
Chapter 9, verse 10..
The fear of the Lord is the beginning of wisdom..
Knowledge of the Holy One is understanding..
You see the parallel here?.
The fear of the Lord, the proper relationship with our God,.
the fear of the Lord is parallel.
with knowledge of the Holy One..
Fear of the Lord, knowing how to fear God.
is to know God..
To know God means to fear Him,.
to understand who He is and relate to Him properly..
And aren't we blessed that we have Jesus Christ.
in our lives?.
It's in Jesus and in Christ that we are related to God,.
restore our relationship with God properly,.
that God is our creator, God is our Father,.
and we are His sons and daughters..
We are His people..
That's how we know God..
But here the why says, the fear of the Lord.
and the knowledge of the Lord is the same..

$^{561}$And that is the foundation of your life..
That is the foundation of your house.
that you have to build on..
And you keep on listening to God,.
keep on being open to God,.
because God is actively doing something in your life..
The problem is we don't feel like God is doing something.
because we stop listening to Him..
We just close down..
So a final reminder is that, you know,.
today is October..
Oh my goodness..
This year is almost over..
But it's never too late..
It's never too late..
If we're willing to fear God,.
if we're willing to know Him, to relate to Him,.
God will bless you..
So I want to remind you to read the Bible..
Read the Word of God every day..
There's one very wise suggestion.
to how to read the Bible.
if you don't know how to begin with,.
or you haven't been reading the Word of God.
for a long, long time..
There's one good suggestion,.
is to take the book of Proverbs,.
because there are 31 chapters in the book,.
and you read one chapter a day..
Go home and read it..
And it's not just reading the Bible,.
it's to pray to Him..
To pray and to be open what God is doing in your life..
What is God doing in your life?.
What does God want you to do?.
Pray for other people..
Maybe there are some friends that we need to pray for them..
Church ministry, we need to pray for..
And the growth of the church, we need your prayers..
We need to pray for the church..

$^{601}$It's not just, I want this, I need this, I want this, God..
Give it to me, give it to me..
If you don't give it to me, you're not God anymore..
This is not worshiping, this is not fear of the Lord,.
this is using God in the name of God for your sake..
And that is pride..
That is idolatry..
Read the Bible, pray to God..
Now observe the time of worship.
because this is probably the entire week.
you can just sit down and be open to God,.
open your heart and listen to what He is saying to you..
And you listen and obey..
As we close, why don't we pray together?.
(pages rustling).
Probably we have forgotten God, but let's pray..
God is full of grace..
God is always there to knock on our doors,.
to welcome us into His presence..
Hey, Father, dear God, thank you for being our God..
(gentle piano music).
Help us to come before you now..
Open our hearts and our mind, our lives,.
so that we don't,.
we are not occupied with ourself, our ego..
But God, let us know, let me know.
what I need to get rid of in my life.
so that I can put my life back to order..
My spiritual relationship with you,.
my relationship with other people,.
to learn how to fear you,.
to learn to listen to your voice..
And God, thank you for being our loving Father,.
even though we are weak.
and we have so many shortcomings,.
yet you are here to embrace us..
You are here to accept us in Jesus Christ..
Thank you for being our loving Father,.
always waiting for us to return..
So we come before you as the Spirit.

$^{641}$to help us, to help our hearts to return to you,.
to enjoy our time with you,.
to cultivate our relationship,.
the fear of the Lord in our life,.
to allot time in our busy schedule.
to relate to you,.
and pray, we just pray that your Spirit will renew,.
renew our faith..
Help us to have the strength to deal.
with all the chaos in our lives,.
the disorders in our lives..
Yes, we are of little faith, but help us, God..
And thank you for being our God, always there..
Help us to grow spiritually, help us to grow mentally,.
help us to grow emotionally,.
so that we walk in the way of wisdom..
This is our prayer, and I pray that in this coming week,.
you will be with us, strengthen us, carry us through,.
and help us to continuously to rely on you,.
to pray to you, and to worship you,.
knowing that you are our strength and you are our rock..
You never abandon us,.
and nothing will separate us from the love of Christ..
Thank you for being our Father, our God..
In Christ, we pray..
Well, God bless you all,.
and we have some refreshment outside,.
and please enjoy some time of fellowship together..
Make sure we spend this time together.
to get to know each other..
I know there are some friends here, some new friends here..
Welcome, and we wanna get to know you,.
and please get to know us as well..
So yeah, God bless you, and I wanna see you next Sunday,.
and we continue to serve God and worship Him..
\newpage



\section{啟示錄 2:1-3:22}
\label{sec:VjJ7v3bZQ7A}
\textbf{基督眼中的教會 \#5 (啟示錄 2\_1-3\_22) - 陳崇基牧師}
\newline
\newline
連結: \href{https://youtube.com/watch?v=VjJ7v3bZQ7A}{\texttt{https://youtube.com/watch?v=VjJ7v3bZQ7A}} ~~~~ 語音日期: 2024-10-22
\newline
\newline
\hyperref[sec:yEZdtKfic4o]{\small{< < < PREV SERMON < < <}}
~
\hyperref[sec:index_chronic]{\small{[返順時目]}}
~
\hyperref[sec:index_scriptual]{\small{[返順卷目]}}
~
\hyperref[sec:uDsOylL_mRw]{\small{> > > NEXT SERMON > > >}}
\newline
\newline
啟示錄 2:1-3:22
\newline
\begin{longtable}{cl}
\hline
\hline
章節 & 經文 (和合本修訂版)\\
\hline
2:1 & \begin{tabularx}{0.7\textwidth}{X} 「你要寫信給以弗所教會的使者，說：『那右手拿著七顆星，在七個金燈臺中間行走的這樣說： \end{tabularx} \\ \\ \relax
2:2 & \begin{tabularx}{0.7\textwidth}{X} 我知道你的行為、勞碌、忍耐，也知道你不容忍惡人。你也曾察驗那自稱為使徒卻不是使徒的，看出他們是假的。 \end{tabularx} \\ \\ \relax
2:3 & \begin{tabularx}{0.7\textwidth}{X} 你能忍耐，曾為我的名勞苦而不困倦。 \end{tabularx} \\ \\ \relax
2:4 & \begin{tabularx}{0.7\textwidth}{X} 然而，有一件事我要責備你，就是你把起初的愛心拋棄了。 \end{tabularx} \\ \\ \relax
2:5 & \begin{tabularx}{0.7\textwidth}{X} 所以你要回想你是從哪裡墜落的，並且要悔改，做起初所做的工作。你若不悔改，我要到你那裡去，把你的燈臺從原處挪去。 \end{tabularx} \\ \\ \relax
2:6 & \begin{tabularx}{0.7\textwidth}{X} 然而你還有一件可取的事，就是你恨惡尼哥拉派的行為，這種行為也是我所恨惡的。 \end{tabularx} \\ \\ \relax
2:7 & \begin{tabularx}{0.7\textwidth}{X} 凡有耳朵的都應當聽聖靈向眾教會所說的話。得勝的，我必將神樂園中生命樹的果子賜給他吃。』」 \end{tabularx} \\ \\ \relax
2:8 & \begin{tabularx}{0.7\textwidth}{X} 「你要寫信給士每拿教會的使者，說：『那首先的、末後的，死過又活了的這樣說： \end{tabularx} \\ \\ \relax
2:9 & \begin{tabularx}{0.7\textwidth}{X} 我知道你的患難和貧窮—其實你卻是富足的，也知道那自稱是猶太人的所說毀謗的話，其實他們不是猶太人，而是撒但會堂的人。 \end{tabularx} \\ \\ \relax
2:10 & \begin{tabularx}{0.7\textwidth}{X} 你將要受的苦，你不用怕。看哪！魔鬼要把你們中間幾個人下在監裡，使你們受考驗，你們要遭受苦難十日。你務要至死忠心，我就賜給你那生命的冠冕。 \end{tabularx} \\ \\ \relax
2:11 & \begin{tabularx}{0.7\textwidth}{X} 凡有耳朵的都應當聽聖靈向眾教會所說的話。得勝的必不受第二次死的害。』」 \end{tabularx} \\ \\ \relax
2:12 & \begin{tabularx}{0.7\textwidth}{X} 「你要寫信給別迦摩教會的使者，說：『那有兩刃利劍的這樣說： \end{tabularx} \\ \\ \relax
2:13 & \begin{tabularx}{0.7\textwidth}{X} 我知道你的居所，就是有撒但座位之處；當我忠心的見證人安提帕在你們中間，在撒但所住的地方被殺之時，你還堅守我的名，沒有否認對我的信仰。 \end{tabularx} \\ \\ \relax
2:14 & \begin{tabularx}{0.7\textwidth}{X} 然而，有幾件事我要責備你，就是在你那裡有人服從了巴蘭的教訓；這巴蘭曾教唆巴勒將絆腳石放在以色列人面前，使他們吃祭過偶像之物，並且犯淫亂。 \end{tabularx} \\ \\ \relax
2:15 & \begin{tabularx}{0.7\textwidth}{X} 同樣，你那裡也有人服從了尼哥拉派的教訓。 \end{tabularx} \\ \\ \relax
2:16 & \begin{tabularx}{0.7\textwidth}{X} 所以，你當悔改；若不悔改，我很快就到你那裡來，用我口中的劍攻擊他們。 \end{tabularx} \\ \\ \relax
2:17 & \begin{tabularx}{0.7\textwidth}{X} 凡有耳朵的都應當聽聖靈向眾教會所說的話。得勝的，我必將那隱藏的嗎哪賜給他，並賜他一塊白石，石上寫著新的名字，除了那領受的以外，沒有人認識。』」 \end{tabularx} \\ \\ \relax
2:18 & \begin{tabularx}{0.7\textwidth}{X} 「你要寫信給推雅推喇教會的使者，說：『神的兒子，那位眼睛如火焰、雙腳像發亮的銅的這樣說： \end{tabularx} \\ \\ \relax
2:19 & \begin{tabularx}{0.7\textwidth}{X} 我知道你的行為：愛心、信心、勤勞、忍耐；又知道你末後所行的善事比起初所行的更多。 \end{tabularx} \\ \\ \relax
2:20 & \begin{tabularx}{0.7\textwidth}{X} 然而，有一件事我要責備你，就是你容忍那自稱是先知的婦人耶洗別教唆我的僕人，引誘他們犯淫亂，吃祭過偶像之物。 \end{tabularx} \\ \\ \relax
2:21 & \begin{tabularx}{0.7\textwidth}{X} 我曾給她悔改的機會，她卻不肯悔改她的淫行。 \end{tabularx} \\ \\ \relax
2:22 & \begin{tabularx}{0.7\textwidth}{X} 看吧，我要使她病倒在床上。那些與她犯姦淫的人若不悔改他們的行為，我也要使他們同受大患難。 \end{tabularx} \\ \\ \relax
2:23 & \begin{tabularx}{0.7\textwidth}{X} 我又要殺死她的兒女，眾教會就知道，我是那察看人肺腑心腸的，我要照你們的行為報應各人。 \end{tabularx} \\ \\ \relax
2:24 & \begin{tabularx}{0.7\textwidth}{X} 至於你們其餘的推雅推喇人，就是一切不隨從這教訓，不明白他們所謂撒但深奧之理的人，我告訴你們，我不會再把別的擔子放在你們身上。 \end{tabularx} \\ \\ \relax
2:25 & \begin{tabularx}{0.7\textwidth}{X} 你們只要持守那已經有的，直到我來。 \end{tabularx} \\ \\ \relax
2:26 & \begin{tabularx}{0.7\textwidth}{X} 那得勝又遵守我命令到底的，我要賜給他權柄制伏列國； \end{tabularx} \\ \\ \relax
2:27 & \begin{tabularx}{0.7\textwidth}{X} 他必用鐵杖管轄他們，如同打碎陶器， \end{tabularx} \\ \\ \relax
2:28 & \begin{tabularx}{0.7\textwidth}{X} 像我也從我父領受了權柄一樣。我又要把晨星賜給他。 \end{tabularx} \\ \\ \relax
2:29 & \begin{tabularx}{0.7\textwidth}{X} 凡有耳朵的都應當聽聖靈向眾教會所說的話。』」 \end{tabularx} \\ \\ \relax
3:1 & \begin{tabularx}{0.7\textwidth}{X} 「你要寫信給撒狄教會的使者，說：『那有神的七靈和七顆星的這樣說：我知道你的行為，就是名義上你是活的，實際上你是死的。 \end{tabularx} \\ \\ \relax
3:2 & \begin{tabularx}{0.7\textwidth}{X} 你要警醒，堅固那些剩下、快要死的，因為我發現你的行為，在我神面前沒有一樣是完全的。 \end{tabularx} \\ \\ \relax
3:3 & \begin{tabularx}{0.7\textwidth}{X} 所以，要記得你所領受和聽見的；要遵守，並要悔改。你若不警醒，我必如賊一樣來到；我幾時來到你那裡，你絕不會知道。 \end{tabularx} \\ \\ \relax
3:4 & \begin{tabularx}{0.7\textwidth}{X} 然而，在撒狄你還有幾位是未曾污穢自己衣服的，他們會穿白衣與我同行，因為他們是配穿的。 \end{tabularx} \\ \\ \relax
3:5 & \begin{tabularx}{0.7\textwidth}{X} 得勝的必這樣穿白衣，我也不從生命冊上塗去他的名；我要在我父面前，和我父的眾使者面前，宣認他的名。 \end{tabularx} \\ \\ \relax
3:6 & \begin{tabularx}{0.7\textwidth}{X} 凡有耳朵的都應當聽聖靈向眾教會所說的話。』」 \end{tabularx} \\ \\ \relax
3:7 & \begin{tabularx}{0.7\textwidth}{X} 「你要寫信給非拉鐵非教會的使者，說：『那神聖、真實的，拿著大衛的鑰匙，開了就沒有人能關，關了就沒有人能開的這樣說： \end{tabularx} \\ \\ \relax
3:8 & \begin{tabularx}{0.7\textwidth}{X} 我知道你的行為。看哪，我在你面前給你一個敞開的門，是沒有人能關的。我知道你有一點力量，也遵守我的道，沒有否認我的名。 \end{tabularx} \\ \\ \relax
3:9 & \begin{tabularx}{0.7\textwidth}{X} 那屬撒但會堂的，自稱是猶太人，其實不是猶太人，而是說謊話的，我要使他們來到你腳前下拜，使他們知道我已經愛你了。 \end{tabularx} \\ \\ \relax
3:10 & \begin{tabularx}{0.7\textwidth}{X} 因為你遵守了我堅忍的道，我也必在普天下人受試煉的時候保守你免受試煉。 \end{tabularx} \\ \\ \relax
3:11 & \begin{tabularx}{0.7\textwidth}{X} 我必快來，你要持守你所有的，免得人奪去你的冠冕。 \end{tabularx} \\ \\ \relax
3:12 & \begin{tabularx}{0.7\textwidth}{X} 得勝的，我要使他在我神的殿中作柱子，他必不再從那裡出去。我又要把我神的名和我神城的名—從天上我神那裡降下來的新耶路撒冷，和我的新名，都寫在他上面。 \end{tabularx} \\ \\ \relax
3:13 & \begin{tabularx}{0.7\textwidth}{X} 凡有耳朵的都應當聽聖靈向眾教會所說的話。』」 \end{tabularx} \\ \\ \relax
3:14 & \begin{tabularx}{0.7\textwidth}{X} 「你要寫信給老底嘉教會的使者，說：『那位阿們、誠信真實的見證者、神創造的根源這樣說： \end{tabularx} \\ \\ \relax
3:15 & \begin{tabularx}{0.7\textwidth}{X} 我知道你的行為，你也不冷也不熱；我巴不得你或冷或熱。 \end{tabularx} \\ \\ \relax
3:16 & \begin{tabularx}{0.7\textwidth}{X} 既然你如溫水，也不冷也不熱，我要從我口中把你吐出去。 \end{tabularx} \\ \\ \relax
3:17 & \begin{tabularx}{0.7\textwidth}{X} 你說：我是富足的，已經發了財，一樣都不缺，卻不知道你是困苦、可憐、貧窮、瞎眼、赤身的。 \end{tabularx} \\ \\ \relax
3:18 & \begin{tabularx}{0.7\textwidth}{X} 我勸你向我買從火中鍛鍊出來的金子，使你富足；又買白衣穿上，使你赤身的羞恥不露出來；又買眼藥抹你的眼睛，使你能看見。 \end{tabularx} \\ \\ \relax
3:19 & \begin{tabularx}{0.7\textwidth}{X} 凡我所疼愛的，我就責備管教。所以，你要發熱心，也要悔改。 \end{tabularx} \\ \\ \relax
3:20 & \begin{tabularx}{0.7\textwidth}{X} 看哪，我站在門外叩門，若有聽見我聲音而開門的，我要進到他那裡去，我與他，他與我一起吃飯。 \end{tabularx} \\ \\ \relax
3:21 & \begin{tabularx}{0.7\textwidth}{X} 得勝的，我要賜他在我寶座上與我同坐，就如我得了勝，在我父的寶座上與他同坐一般。 \end{tabularx} \\ \\ \relax
3:22 & \begin{tabularx}{0.7\textwidth}{X} 凡有耳朵的都應當聽聖靈向眾教會所說的話。』」 \end{tabularx} \\ \\
[1ex]
\hline
\hline
\end{longtable}
$^{1}$Thank you, Nelson.
我們繼續去.
今天我們講到的topic.
就是基督眼中的教會.
By the way.
我去年答應去多倫多講道的時候.
我沒有想過會有這麼大的印象.
我以為是一個人去.
然後我回來.
不過那是一個多倫多區的聯會.
舉辦的培靈會.
他們每年舉辦一次.
去年的時候.
他們問我有沒有興趣去他們中間做講員.
所以我去那裡一年三晚.
還有一個星期六下午.
大概有四個環節.
我是負責的.
負責人問.
你們教會想不想有一個.
因為是一個conference.
他就問你們教會或者你們的士工.
會不會有一個攤位.
給你一個攤位.
我想了一會.
不如這樣吧.
我們也拿個攤位.
隨便拿張桌子.
放張桌子在這裡.
放一些東西在這裡.
所以有些同工跟我一起去.
所以是非常美好的安排.
三天希望能夠祝福當地的教會.
講到教會.
我們繼續看這個series.
就是啟示錄七教會的書信.
這次我們去到第五間.
第五間的教會.
第五封的書信.
這個地方叫做薩迪.

$^{41}$因為這封書信比較短.
不如我們有聖經的話打開.
如果沒有的話.
我們就在powerpoint上.
我們能夠看得到.
因為短的書信.
不如我們讀一次.
我們一起看完整封書信.
有多少節.
六節.
你要寫信給薩迪教會的使者.
哪有上帝的七靈.
七火星的這人說.
我知道你的行為.
就算名義上你是活的.
實際上你是死的.
你要警醒.
堅固那些盛夏快要死的.
因為我發現你的行為.
在我上帝面前沒有一樣是完全的.
所以要記得.
你所領受和聽見的.
要遵守.
要悔改.
你若不警醒.
我必如賊一樣來到.
我何時來到你那裡.
你絕對不會知道.
然而在薩迪.
你還有幾位.
是未受污穢自己衣服的.
他們會穿白衣與我同行.
因為他們是陪穿的.
得勝的必這樣穿白衣.
我也不從生命插上.
逃去他的名.
我要在我父面前.
在我父的眾使者面前.
宣認他的名.
凡有耳朵的都應當聽.

$^{81}$聖靈向宗教會所說的話.
很短的一封書信.
第五封書信給七個教會.
第五間教會薩迪的教會.
我們回到頭三節的時候.
我們發現薩迪到底是一個怎樣的地方.
一個怎樣的教會.
他說這個教會.
耶穌說我知道你們的行為.
我知道你們所做的事.
知道你們的生活.
在名義上你們是活的.
實際上是死的.
這些說話真的很絕.
聖經裡面有很多說話是很絕的.
我希望我們讀經的時候.
讀到這些味道.
不然我們讀經的時間浪費了.
靈修的時間浪費了.
有時候聖經是想挑釁我們.
去挑戰我們.
看完聖經的時候.
可能是很安樂的.
看完聖經的時候.
可能不是很安樂.
因為聖經很多時候說的都是真理.
因為上帝是不會說謊.
上帝是對我們說真理的.
等於你去看醫生.
看完醫生爽不爽.
很多時候很不喜歡他.
看完醫生之後.
但是你是需要聽的.
這封書信也是一樣.
寫給薩迪的教會.
他說有上帝的七靈和七火星的這樣說.
我們曾經解釋過七火星和七靈是什麼意思.
在第一章啟示錄.
復活的基督.
掌權的基督.

$^{121}$他說他在七星之中.
七個燈台之間.
七靈之中來行走.
就是說他知道教會裡面發生什麼事.
有時候可能我知道教會.
好像很多東西.
但是耶穌基督知道更多東西.
因為他知道每一個人在想什麼.
每一個人自己的私生活.
每一個人自己的生命是怎樣的.
薩迪這個地方.
薩迪這個地方是一個很有名聲的地方.
薩迪這個城市.
以前是一個很闊綽的城市.
如果你有聽過一些古代的歷史.
當然大家未必有興趣.
有一個以前的王國.
叫做利迪亞王國.
利迪亞王國在土耳其的西邊.
當時是非常有錢的.
為什麼呢.
因為他們有很多金礦.
我預備的時候就想起以前.
我年輕的時候.
有時候玩那些電腦遊戲.
我不知道大家以前有沒有玩過電腦遊戲.
可能我那一代.
玩過Age of Empire.
帝國的時間.
那個遊戲.
是一個戰略遊戲.
你到處去找一些金礦,銀礦,木等等的資源.
然後你儲起這些資源.
然後你就可以攻打其他人.
薩迪所坐落的地方.
就正正是有金礦在那裡.
但是耶穌說什麼呢.
他說你名義上是很好的.
聽你的名字是非常好的名字.
但是你們本質上是死的.

$^{161}$你們是死的.
你是死的.
原文只有兩個字.
就是死了你.
第二節.
他說你要警醒.
堅固那一刻.
剩下一點點的.
你還剩下一點點.
還沒死的.
快要死的.
你快點要堅固剩下.
就快要死的.
因為我發現你的行為在上帝面前沒有一樣是完全的.
所以你要記得你所領受的.
所聽見的.
要遵守.
要悔改.
若不悔改.
耶穌就好像賊一樣來找你.
你是不知道的.
什麼時候來.
你絕對不知道.
薩迪我正如所說.
是金礦的地方.
你見到那裡.
左下角.
就是今天的古城的地方.
你見到上面這些都是在那裡找到的東西.
全部都是金礦.
為什麼呢.
因為那裡附近有一條河.
那裡的河水和河裡面找到的泥.
泥沙.
就是一些金的粒子在裡面.
就是右上角那裡.
你見到中間有些鐵屑.
有一些泥土.
仍然是有一些很少很少的金.
所以如果你要尋金的話.

$^{201}$你就走去這個地方.
薩迪.
大家知道我在加州很多年了.
我來美國之後就在加州.
我住在北加州.
塞庫曼路附近.
在那附近有很多以前的掘金.
尋金的城市.
其中一個叫奧本.
我們塞庫曼路去奧本.
不是很遠的.
你就可以去那裡看看.
以前那些尋金熱亞的時候.
那些人在那裡做什麼.
有些房子.
有些工廠在那裡.
然後我們往西邊.
下去就是三藩市.
三藩市大家知道它的名字是舊金山.
舊金山.
為什麼叫舊金山.
因為以前是尋金的地方.
那為什麼是舊呢.
沒有了.
剩下49了.
這些歷史.
全部都在那裡可以見到.
那些金礦的地方.
可以找到金的地方.
但是今天你去那些奧本.
不同的金礦的地方.
全部都已經是很殘舊的地方.
是Ghost Town.
因為不再有人在那裡.
你掘完金就不會再去.
你所聽到的.
它的名聲.
那些town.
那些城市.
那個名譽.

$^{241}$只不過是舊金山.
舊金山.
如果你今天說我去舊金山.
人家不知道你說什麼.
你只會洩露你的年齡.
人家知道你是老移民.
你也知道.
現在沒有人叫舊金山.
沙迪也一樣.
沙迪就是舊金山.
以前是很出名.
不過現在已經是不再出名.
還沒到Ghost Town.
所以耶穌說什麼.
剩下的快要死的.
他用這個城市的文化.
來形容當地的教會.
當地的教會就是這樣.
名義上你是活的.
在一個金城裡面.
多厲害.
但是如今已經不再是了.
因為有的金已經掘完.
但是你的生命是怎樣呢.
以前可能是很多金礦.
名義上的名譽是很好的.
可能為主做很多事.
但是現在已經是一個Ghost Town.
耶穌說你要警醒.
要堅固.
要靜下一點.
靜下什麼呢.
如果你去河裡.
你找找那些泥沙.
還要靜下一點點.
靜下一點點.
快要死的.
要堅固一點.
因為我發現你的行為在上帝面前.
沒有一樣是完全的.

$^{281}$要記得你所領受.
要遵守.
要去悔改.
一年好像連環一樣.
其實就好像.
好像你早上起不了床.
你教鬧鐘.
每分鐘都要教一個鬧鐘.
連教五個.
就是告訴你.
你需要醒來.
在這個時候.
你要去蘇醒過來.
很多時候我們基督徒的生命.
為什麼去到這樣的餘地.
因為我們以往覺得我們是金礦.
我們為主做過多少事.
我們為主立過多少字.
為主奉獻了多少事.
耶穌說.
是快要死了.
現在是怎樣.
以前可能很厲害.
可能很有名譽.
但現在是怎樣.
這裡講清楚.
上帝耶穌看我們的時候.
看你的時候.
不是看你的過去.
過去多爛也好.
過去多金碧輝煌也好.
耶穌所看的是現在.
現在.
你現在在做什麼.
所以一連五次的.
五個鬧鐘告訴你.
你需要醒來.
你現在需要醒過來.
不要去裝睡.
往往我們就是在裝睡.

$^{321}$裝睡是叫不醒的.
扮睡覺.
你有沒有試過扮睡覺.
我小時候.
小朋友小的時候.
我不想他煩我.
我就在扮睡覺.
叫不醒.
怎樣叫醒.
本身都沒有醒.
只不過在裝睡.
基督徒很多時候都是這樣.
我們信了主.
我們就是一個裝睡的人.
神和聖靈去提醒我們.
我們在扮睡覺.
扮聽不到.
假裝聽不到.
真正耶穌所看見的.
所著重的.
所看重的.
不是你的過去.
是你的現在.
是你現在.
我每一天都這樣提醒自己.
以往做過的事.
有多少都好.
其實都不是很多.
很少.
有多少.
奮不戰打了它.
因為主要看的是.
今天我要為它做什麼.
接下來我要為它做什麼.
這才是我們基督徒應該有的生命.
這封書信.
寫給教會.
上帝說我現在要看你現在.
你現在還有一點點.
那是可以的.

$^{361}$但你一定需要悔改.
你需要醒過來.
不要再裝睡.
要兼顧你還剩下的.
神在看什麼.
他說:二言.
在撒迪里門還有幾位.
是未曾污穢自己衣服的.
他們會穿白衣與我同行.
因為他們是配穿白衣的.
得勝的.
必這樣穿白衣.
我也不從我的生命插圖去看命.
我要在我的父.
我重視者的面前.
去宣認他的命.
凡有耳聽的.
就應該聽聖靈向宗教會所說的話.
你留意到耶穌看一個教會的時候.
The few people.
他留意到的你們還有幾位.
是忠於我的.
上帝看到.
是那班忠於他的人.
不是那個數目.
只要他們願意誹服.
上帝看到.
並且給他們有認許.
有promise.
將來是怎樣.
他說要給白衣他們.
配受穿白衣的.
又不會在生命插圖去圖抹他們.
白衣在啟示錄裡面經常出現.
聖潔.
神聖的.
屬於上帝的.
得勝的.
這些人全部都是穿白衣的.
耶穌說如果你悔改.

$^{401}$你願意去到歸向我的話.
今天.
是穿上白衣.
這個認許穿上白衣.
對當時羅馬的背景來說.
很有意思.
因為他們經常會看到一些白衣的人.
在街上走.
我不是說Halloween.
為什麼呢.
因為當時羅馬帝國.
每逢一些大的城市.
那些軍隊去別的地方.
去征服其他人的時候.
搶了其他人的東西.
然後他們將這些東西帶回他們的城市裡面.
他們海船回歸的時候.
就會列隊有很多士兵.
那些士兵很多時候就是穿白衣的.
穿白衣這樣進去這些城市.
然後興高采烈的話.
我們得勝了.
我們能夠進去和王者慶祝.
就是這樣的一個意境.
如果去羅馬的話.
那個鬥獸場附近.
就有一個叫做提多的Titus Arch.
提多的海船門.
他們去慶祝羅馬軍隊去到巴勒斯坦地.
然後打壓了第一次猶太人的叛亂.
打壓了之後.
然後就在耶路撒冷的聖殿.
搶走他們最貴重的東西.
你看到左邊那個就是那個城門.
就是那個海船門下面其中一幅畫像.
在石頭上雕上去的畫像.
你看到那些士兵托著上面一個東西.
七燈台.
看到嗎 七燈台.
然後有很多那些標語.

$^{441}$很多的旗等等.
前面的看到有兩支長的.
那些應該是數花.
就是那些Trumpet.
在聖殿裡面找到的.
那些是皓角.
還有陳設餅等等.
全部東西他們拿回來羅馬.
你看到那些士兵穿什麼衣服.
左邊你看不到.
因為那些已經剝落了.
那些顏色.
但其實你找回那些跡象的時候.
你知道原來本身是彩色的.
這些其實是中央電視台說的東西.
羅馬中央電視台就說我們多厲害.
這些白衣人進城的時候.
他們海船回歸的時候.
這班得勝者.
他們是和王者一起慶祝.
耶穌說我才是得勝者.
能夠忠於我的.
穿起白衣.
能夠和我一起進入新天新地.
得勝的我給他白衣穿上.
因為他配穿白衣.
在生命冊上不會逃去他的名字.
在上帝面前.
眾使者面前去認你的名字.
認你的名字.
我想唯一一件事.
幫到我不會裝睡的.
就是我知道我某一天.
I have to answer.
有一天.
我要在上帝面前去回答.
到底我這個生命做了什麼.
為他做了什麼.
一個宇宙性一個人生性的問題.
一個Cosmic question.

$^{481}$是我們每一個基督徒都需要回答.
不要裝睡.
因為將來一定要去見這位主的時候.
我都要交什麼給他.
對不對.
見這位主的時候.
我不能夠和主說.
Yes I have been faithful.
I have been faithful for little things.
海船回歸.
能夠穿上白衣.
今天看完這封書信的時候.
但我們在上帝面前可以記起.
耶穌見你.
耶穌看見你的時候.
不是看你的過去.
Yes I know.
以前你可能為上帝做過一樣.
上帝做過那樣.
可能是很漂亮的東西.
But what matters.
Is now.
Now and to the future.
耶穌所concern的.
不是你過去是你現在.
不是你過去的名義是現在.
你對主的信.
我們來到主的面前.
我們好像這封書信這樣說.
我們要警醒.
然後我們不再裝睡.
讓主提醒我們.
去遵守去悔改.
去警醒.
直到他回來的時候.
我們同心協力去記住.
主我們眾弟兄姊妹來到你面前.
我們求你.
求你赦免我們.
在我們生活的忙碌裡面.

$^{521}$在我們生活很多不同的偶像裡面.
我們有時的事.
確實真的忘記了.
在名義上好像是活的.
但我們心底裡面.
我們生命和上帝之間的關係.
往往很多時候是死的.
主求你再一次去警醒我們.
給我們力量給我們恩典.
我們迴轉.
去知道將來你為我們所預備的.
海船回歸的那種新天新地的榮耀.
並且我們需要在你面前.
來交仗.
主求你再一次挑釁我們.
挑戰我們.
以致我們生命不再回望以前.
我們曾經有過的事.
因為你所看見的是.
今天我們的心.
是不是真正全心全意來愛你.
全心全意來服侍你.
主求你在我們生命裡每一個角落.
再一次作警醒的工作.
讓我們醒過來.
我們自己個人的生命.
我們的家庭.
我們的一切.
是要向著你.
去榮耀你.
主要幫助我們的工作.
叫我們忠心地做.
不是純粹為了賺錢.
而是為了榮耀神.
去好好管理我們自己的生活.
我們敬拜不是純粹follow the motion.
去循例式地來做.
而是真正來到你面前.
知道你是我們教會的主.
是我們生命裡的主.

$^{561}$是我們生命裡每一個房間.
每一個方面的主宰.
主求你再一次來做我們的主人.
我們娛樂的生活.
我們消遣的生活.
我們的工作.
我們的家庭關係和其他人的關係.
論社的關係.
主要再一次來歸給你.
求你在我們中間.
在我們生命裡再一次去作王.
帶領我們.
我們這樣祈禱奉主的名禱.
願主上帝祝福各位.
讓大家平安.
\newpage



\section{16:18}
\label{sec:uDsOylL_mRw}
\textbf{Way of Wisdom \#2 - Proverbs 16\_18 (Rev. Mark Chan)}
\newline
\newline
連結: \href{https://youtube.com/watch?v=uDsOylL_mRw}{\texttt{https://youtube.com/watch?v=uDsOylL\_mRw}} ~~~~ 語音日期: 2024-11-02
\newline
\newline
\hyperref[sec:VjJ7v3bZQ7A]{\small{< < < PREV SERMON < < <}}
~
\hyperref[sec:index_chronic]{\small{[返順時目]}}
~
\hyperref[sec:index_scriptual]{\small{[返順卷目]}}
~
\hyperref[sec:v8IvKYM0B_8]{\small{> > > NEXT SERMON > > >}}
\newline
\newline
$^{1}$- Thank you so much for your appreciation.
and I really appreciate and really encourage us as staff.
to continue to be faithful.
and we'll put your gifts in good use, all right?.
We could put your gifts in good use..
Today, we're going to continue our series on way of wisdom..
How do we live a life of wisdom?.
And the Bible has so much to say about it..
And remember this passage that I talk about last week.
and it's from chapter 24 of Proverbs..
And here, the wise teacher is going through.
some observation and he went past the field.
of the lazy people, sluggard,.
and he saw that everything is in ruin.
and he's trying to apply from his observation of the nature.
and apply it to his spiritual life, his personal life..
So again, in the book of Proverbs,.
it's not simply about natural observation..
It's about observing and then get some spiritual insight.
so that your life can be built up..
And one has to work against and fight against.
the chaos in our lives..
You know, in the past week, have you thought the chaos.
or the chaos overwhelm you?.
You know, there are tons of chaos around us.
and if we don't put energy to work against it.
and fight against it, it will be a field of sluggard.
and it will be in ruin..
And your spiritual life, your personal life,.
your marriage, your family, your work, your career,.
everything will be in ruin..
And one has to do that and that is what the wisdom literature.
or the book of Proverbs has to say to us..
It's bottom up, observe how life works..
Observe the moral order, the spiritual order.
of God's universe and to observe that.
so that you can get spiritual insight from it..
Now, sometimes we have to understand how to read Proverbs..
So I really wanna encourage you to read Proverbs..
And here I give you a simple structure.

$^{41}$of the book of Proverbs..
And in this book, there are 31 chapters..
And if you zoom out a little bit,.
you find that there's a whole structure of the book..
And you begin with the first nine chapters,.
which is the introduction..
It's a prolonged introduction..
It's an invitation to walk in wisdom..
And you get the sense that the father or the mother.
is telling the son to walk in wisdom,.
to invite him to be wise, to learn the way of wisdom,.
encouraging him and urging the son to choose wisdom.
and reject folly of the world..
To choose God, reject the false gods in the world..
And then you come to a long collections.
of individual sayings from chapter 10 to 30..
And if you read those chapters, have you tried it?.
If you try to read one verse after another,.
it's sequentially, you try to understand it,.
you can't find a really tight flow of argument..
Because it is not designed to be like that..
It's a collection of a lot of different wise sayings.
and put them together..
It's like an anthology..
You know what an anthology is?.
It's a collection of essays..
It's a collection of articles in a book..
So there's not necessarily a logical flow from it..
So you can read chapter 10 to 30.
the same way you read the Book of Romans..
You know, the Book of Romans, we are doing bio study on it.
and we are still on the early chapters of the Book of Romans.
and you have to read it from verse one, verse two, verse three.
because Paul is trying to build an argument there.
and try to explain some of the themes.
and is trying to carry you through.
with some sort of sequential, you know, logical flow..
But in the Book of Proverbs, it's not like that..
So in order to understand chapter 10 to 30,.
the best way, I think, is to read them topically..

$^{81}$Is to read through it and then you try to collect.
some of the things that is related.
to a certain kind of topic..
For example, money..
Wealth and money..
The Proverb has a lot to say about money, about wealth..
So you read through those chapters.
and try to collect those things..
Just like you go to Whole Foods,.
you want to buy something and then you go through.
different aisles and then you pick what you want..
It's kind of like that because, you know,.
wisdom is kind of like that..
Wisdom is observing and collecting.
and try to organize things..
So you read through these chapters 10 to 30.
with a certain kind of topic in your mind,.
like family life, work and career,.
speech, how we speak and when to speak,.
wisdom of speaking..
You categorize them under different topics,.
under different headings, and then you try to understand.
what the wise is trying to tell you..
And today, I want to demonstrate this topical reading..
And we're going to understand or study the theme of pride.
and humility..
Pride and humility..
What is that?.
So when you think of humility or pride,.
what do you think of?.
How do you define it?.
How do you describe humility?.
How do you describe pride?.
And from my understanding, we'll read through.
some of the Proverbs that has to do with this topic..
And I organize it into two related headings..
First heading is humility leads to more wisdom..
That's what humility will do to you..
The second topic is humility will lead to more life,.
life in God, more substantial growth in your life..

$^{121}$So in these two headings, I want to study with you.
what humility, what pride is..
Now, first of all, what humility is not..
What is not humility?.
What is not?.
Now, humility is not humiliation..
Humility and humiliation are totally different things..
All right, the difference is very, very, very important..
Because especially for some of us who grow up.
in a shame-based culture, a shame-driven culture,.
you know, the Asian culture, is very much shame-based..
And humiliation sometimes is messed up with,.
misunderstood to be humility..
Humiliation or self-abasement, it's to shame yourself.
into thinking that you are worthless..
You know, less worth than other people..
Have less value than other people..
And it's very intriguing because at times,.
and this is how I experience it,.
is that we think that we don't measure up.
and we have to compare ourselves to others..
And in doing so, we try to prove that we are better..
We try to assert our own value.
by comparing ourselves to others..
But this is a whole complex of pride and, you know,.
humiliation, because when you compare yourself.
with other people, so often the result is that,.
okay, maybe I still have something that you do not have..
I still get something that you don't..
And it's very easy to fall into another trap for pride..
I compare myself with you, against you,.
and I find that, okay, maybe you are better,.
but maybe there's some way I'm better than you,.
even though there's some other things.
that you are better than me..
Now, that is not humility..
Humility in the Bible says that, yes, I'm broken,.
I am sinful, but God accepts me unconditionally in Christ..
You see, this will free you up from comparison..
This will free you up from being arrogant, from pride,.

$^{161}$because humility is seeing yourself as low,.
humble, in front of God..
No one can measure up to God's standard..
So we come before God and we know, okay,.
we are sinful, we are sinful, and we are humble..
However, we are totally accepted unconditionally in Christ..
And that's the ground we stand on..
That is humility in the Bible..
Look with me, chapter 15..
"Whoever ignores instruction despises himself,.
but he who listens to reproof gains intelligence," wisdom..
33, "The fear of the Lord is instruction in wisdom,.
and humility comes before honor.".
All right, let's begin with the last verse here,.
the second verse, verse 33..
What is humility here?.
Look the parallel here..
You know, the wise saying always says,.
two short sentences put together..
Look the parallel between two sentences in verse 33..
"Fear of the Lord is parallel with humility..
Fear of the Lord, as I explained briefly last time,.
or other occasions, fear of the Lord means.
a proper relationship with God,.
a proper relationship with the Creator..
That is the fear of the Lord..
That's why I fear God, I worship God,.
because He is my Creator, He is my God.".
Right?.
So proper relationship with God is fear of the Lord..
And that is parallel, or sort of equal to humility..
To understand my relationship with God..
You see, this is the gospel of Jesus Christ..
This is the gospel..
I'm totally, 100\% accepted by Christ,.
and that is sheer grace of God..
It's not because of my works..
It's not because I'm better..
It's not my goodness..
I'm totally accepted because what Christ has done for me..

$^{201}$Right?.
The righteousness of Christ..
I'm solidly grounded in the love of Christ..
And that is my security..
My security does not come from.
because I'm better than my co-worker,.
or my co-worker is worse than me..
My security does not come from,.
okay, I'm better than my friends,.
or better than other people, probably in the church,.
I'm better than them, so I feel better..
That is a false sense of identity..
That is idolatry..
But when you are properly related to Christ,.
you receive grace unconditionally,.
and you receive the gospel,.
that's where your identity and your security come from..
You see what I'm saying?.
So we don't have to compare with other people,.
because I have a solid ground to stand on..
That is the love of Christ, the grace of Christ..
In my college year,.
I remember a good pastor who taught this..
He shared one of his experience,.
and this pastor, I just love him..
He's so wise..
And this pastor told us,.
when he was much younger, when he was younger,.
and he joined the church,.
and he wanted to join the church choir..
So there he was, and he joined the church choir,.
and he had rehearsal,.
and he served in the choir for a few times..
And then one day, after the choir performance,.
or after the service,.
the choir director pulled him aside,.
gently and firmly said to him,.
and asked him, "Dear brother,.
"do you really think that God calls you to be in the choir,.
"because you're tone deaf?".

$^{241}$And for the few times he was trying to join the choir,.
because he liked to sing..
He likes to sing..
He likes to sing hymns and praises to God..
So he wants to join the choir, or the worship team..
And then it just dawned on him that he was tone deaf..
He just cannot sing in tune..
And he messed up other people in the choir..
And after this realization,.
he understand he's not musically gifted..
He's not musically gifted..
Yes, he can sing praises to God,.
you know, in the congregation,.
probably it will annoy some of the people beside him,.
but he's not musically gifted..
And that's how he come to a conclusion..
And he shared with us his non-participation.
is helpful for the ministry..
And he's not gonna be jealous and envious of other people..
Other people who are musically gifted,.
they can share the stage..
They can lead other people to worship,.
but he's not gifted..
But he was totally secure..
He doesn't feel bad..
He doesn't feel sad..
That's how God created me..
That's how God gifted me..
And that's how I use my gifts..
And he is very gifted in teaching..
Not in music, not in worship,.
but he is very gifted in teaching..
And that's why he's trying to polish his craft.
and to be a good teacher..
And he was a good teacher and he is a good teacher..
And that's what humility does to you..
You listen to advice..
You listen to those who teaches you..
And he who listens to reproof gains wisdom,.
gains intelligence..

$^{281}$Now, what is humility?.
I like what C.S. Lewis says about humility..
C.S. Lewis is a very famous, well-known theologian.
or apologist in the 20th century..
Oh, he passed away a long time,.
but he wrote a lot of things.
and people are still studying the works of C.S. Lewis..
In fact, there was a institute of C.S. Lewis.
and people trying to understand what he wrote..
He wrote so much..
And C.S. Lewis says about humility..
He says, "Humility is not thinking less of you,.
"but thinking of yourself less.".
Amen?.
Humility is not humiliation..
It's not thinking less of yourself,.
but thinking yourself a little bit less..
Because if you're totally grounded in the grace of God.
and the love of God, you're free..
You're free to compare yourself.
and contrast yourself with others..
And that, as I explained, will lead to humiliation.
or a kind of subtle pride..
Humiliation in the sense that I don't measure up.
and I'm jealous, I'm envious of other people..
They have other gifts and I don't have..
And so you're kind of in this loop of humiliation..
And the other result is pride,.
is seeing at least I'm better than you..
No, there's a very, very intriguing complex to it..
But humility is not thinking less of yourself,.
but thinking of yourself less..
And the book of Proverbs has so much to say.
about pride and humility..
Humility is the way of wisdom..
Another verse in 1310..
It says, "Pride only breeds quarrels.".
But wisdom is found in those who take advice..
You know, have you ever argued with someone?.
Have argument with the other people?.

$^{321}$Either your spouse or your child or a friend.
or whoever that is, or even your pastor..
You see, in the midst of argument,.
so often there is pride, right?.
I want to prove that I'm right and you are wrong..
You know, in premarital counseling,.
I always tell the couples, I always tell couples,.
in marriage, marriage is only game in town.
that you have the result of,.
you don't have the result of win-lose..
You always have both of you lose.
or both of you win..
There's no win and being lost in marriage..
Because if you lay down yourself,.
if you grounded yourself in the love of Christ,.
in the grace of God, there's no pride..
If there's no pride, there's no argument..
You listen to one another..
You learn from one another..
See, wisdom is found in those who take advice,.
who listen to advice..
Last night, we have a small gathering.
with some of the husbands in the church..
And this is kind of a once in a while,.
once in probably nine months or 10 months,.
we get together for a simple dinner..
So I went there and to my surprise,.
we had so much to talk about..
So I show up in the restaurant at 7 p.m., 7 at night..
Went to Phoenix, 7 o'clock..
And I was planning to have a dinner for two hours,.
for two hour max, right?.
What do men talk about?.
Men have nothing to talk about..
But when we sit down and then we warm up.
and then we go into talk about life.
and talk about family in general,.
to my surprise, we had a five hours dinner..
Five hours dinner..
I got home after 11..

$^{361}$Yeah, but I really enjoyed the community..
I really enjoyed the group, you know,.
try to open up ourselves and see,.
hey, you know, can we listen to one another.
and learn from one another?.
Now, I want to tell you,.
this kind of community does not come easily..
Does not come naturally..
For us, for many of us, we have to commit.
in order to cultivate that kind of relationship..
We have to open up and we have to make commitment.
to spend time with one another..
Otherwise it won't happen, right?.
You go to church sometimes,.
you find the conversation is not deep..
And you get tired, you don't want to have fellowship..
You know why?.
It's not because people are shallow..
All of us are yearning for meaningful conversation..
However, in order to do that,.
there is investment..
We have to spend time and commitment.
to stick to one another in order to cultivate.
that kind of family, that kind of community..
So with the men, we talked so much about,.
you know, how to lead the family,.
how to relate to our kids,.
how to deal with the local politics,.
you know, election is coming,.
and what are we concerned about,.
and why do we concern about this or that with our kids?.
And it was a very, very open.
and just great time together sharing.
and communication with one another,.
you know, over a glass of wine..
And that is wisdom, right?.
The community is the wisdom..
So here we come to the second heading,.
which is humility leads to life..
Humility will give you more life..

$^{401}$If you gain wisdom and then you get life.
and then you get better life,.
a godly life, and then it will return back to you.
and will give you more humility..
And you will be humble before God and you learn more.
and then you get more wisdom..
And this is a kind of a circularity..
Circularity in terms of, you know,.
you give and then it will give back to you..
1618, pride goes before destruction..
Haunted spirit before a fall..
Now, this is the observation of the wise teacher..
He observed there was an order in your life,.
order in human life in general..
This is the cosmic order..
You think you're powerful, you're invincible..
If you're prideful, you will fall..
You will fall, inevitably..
Because I'm preparing some of the messages.
for my trip to Toronto,.
I'm talking about the book of Daniel,.
especially the later chapters of Daniel..
Those are really tough chapters..
You have to do a lot of interpretation,.
do a lot of careful reading..
And I was doing some research,.
because the later chapters of the book of Daniel.
is all about the rise and fall of human kingdoms..
Rise and fall of human kings, like Assyrian kings.
and Babylonian kingdom and Persian kingdom.
and the Greek kingdom, the Greek empire..
It's all about the rise and fall of human empire on earth..
And only God's kingdom is eternal..
And that's pretty much what the later books,.
or later chapters of the book of Daniel is all about..
So I was doing some research on the rise and fall of kingdoms.
and it's kind of fun research..
And now, of course, when we talk about.
the fall of any empire, any king,.
it's a whole complexity of factors go into it..

$^{441}$How do you explain the fall of an empire?.
Either in recent history, like the Soviet Union,.
fell in 1991, if I'm not mistaken,.
recent history, or in ancient history, in biblical time..
Now, how do you explain the fall of any empire?.
Of course, there are a lot of different factors go into it..
You know, economic slowdown or political corruption.
or the fall of, the collapse of their social structure.
or social unrest or there's a food supply,.
you know, farm surplus, not enough..
You know, those things..
But the newest research I want to share with you,.
the newest research is actually climate change..
Climate change..
You never thought of it, right?.
Never thought of it because, you know, an empire fell.
because, of course, it's mismanagement.
by one of the bad king, right?.
The last king..
And the newest research has to say.
there's one more factor that we have forgotten..
It's climate change..
For example, this one..
How did the Assyrian Empire fall?.
Now, Assyrian Empire, as you tell from the Old Testament,.
is the most powerful kingdom,.
the first most powerful empire.
the ancient world has ever seen..
Ancient first empire, Assyria,.
from the seventh century to down to,.
well, from the ninth century down to the seventh century,.
200 years, 200 some years of powerful empire..
And they conquered land after land after land,.
conquered their whole world..
And why did they fall?.
Now, the researchers are saying.
not only the social factors or political factors,.
but there's climate change..
There's 125 years of magnet drought..
125 years..

$^{481}$Now, one year of drought, you know,.
you probably can, you know, go through it..
Two years, you probably can go through it..
10 years, probably you can go through it.
because you have enough surplus..
But, you know, 125 years of mega drought,.
that happened in the last 100 years of the Assyrian Empire..
That's why the empire fell..
It collapsed..
Now, that's very, very relevant for our day, of course,.
because it's climate change, right?.
You see all this extreme climate pattern in Florida,.
in some of the places, in the wildfire, hurricane,.
wipe out everything..
Of course, it has so much to say about.
what the Bible has to teach us about creation care,.
about environment and all that, right?.
And many researchers are focusing.
on this kind of direction..
You know why?.
Because not only because this is a relevant topic,.
it is because you cannot control it..
You cannot control it..
No matter how invincible you are,.
you just cannot control climate..
Now, that's what pride is..
To the empire, to the king,.
they always think that they are invincible..
Pride means self-sufficiency..
I'm all good..
I don't need other power, right?.
I can stand on my own..
I have my success, achievement, it's all mine..
It's by my power..
That's what pride is, self-sufficiency..
There are factors that you cannot control..
Even for empire, you just cannot control it..
And that means you are not invincible..
You better be humble..
Last verse I wanna share with you is 22.4..

$^{521}$Humility and the fear of the Lord.
bring wealth and honor and life..
Humility will give you life,.
but pride will give you collapse..
Think about some of the examples.
that you probably have some in mind..
How personal life, how someone who is so successful,.
and he fell and collapsed..
Moral collapse, spiritual collapse..
Think about even religious institution..
It can fall..
If you think you're invincible,.
if you think you don't need to be accountable,.
what lesson we need to learn?.
Humility and the fear of the Lord,.
again, it's a parallel here,.
bring you wealth, honor, and life..
You cannot be pride, prideful,.
and think that you've got everything under your belt..
You see, the book is telling you.
you have to walk in the way of wisdom..
You need to surround yourself with other people of God,.
other godly people..
You cannot go through the spiritual journey by yourself..
You have to have someone to help you,.
and that is humility..
And this will give you wealth and honor and life..
Now, I like the package here,.
wealth, honor, and life, especially the word honor..
You know, in Hebrew, Hebrew is a very concrete language..
It's very concrete language.
because you use a word picture all the time..
And honor, this word, is kavod in Hebrew..
And what does it mean?.
Kavod means something really heavy, weightlifting..
Have you been to gym?.
Weightlifting, that is kavod..
That is something heavy..
Life has substance..
If you have honor, you have substance..

$^{561}$And humility will give you that..
The fear of the Lord will give you that..
So that's why we need to learn how to be humble..
And remember, we have to ground ourselves.
in the sheer grace of God..
And once we do that, we have security, true humility..
So we can do what God has called us to do..
Kelly has reminded us of the need for Bible study..
Now, I want to challenge and I want to encourage you..
Yes, we do need that..
To study the word of God,.
to listen to other brothers and sisters,.
and to be led so that you can understand the word of God..
We all need that, you know..
If you think you know it all,.
I think among us here,.
I would say I'm the one know it all, right?.
But you know, every time I have Bible study.
with other people, other brothers and sisters,.
sometimes we go tangent, okay, fine..
But you know, I always learn something..
I always learn something..
I'm not going to say, you know, you guys,.
you know, you don't know much about,.
you know, the scripture and all that..
That would be pride..
But every time I get together with other people.
and listen to their life.
and listen to how they read the word.
and what they get from the truth of God,.
you know, that will cultivate wisdom..
And we all need this..
And humility will give you reward..
And the reward is eternal life,.
wealth, honor, and life..
Let's all pray together..
And ask Father to grant us true humility..
Father, we come before you..
You are the source of our life..
You created us..

$^{601}$You gifted us in so many different ways..
And we want to praise you..
We want to give you thanks because we are all different..
We're all gifted in a wide diversity of way..
And that is your will..
And that is your goodwill..
Thank you, Father, for calling all of us to be your people..
Teach us the way of wisdom..
Teach us the way of humility..
Help us to learn..
Learn about your word..
Learn about what you're doing in our life..
Help us to be grounded in your sufficient love..
And God, you are the all-sufficient God..
And we are not..
We're not sufficient..
That's why we need your word..
We need your community..
We need your love..
Thank you for being our God..
And call us and encourage us.
so that we'll walk humbly before you..
Help us to look in our lives.
in the light of what Christ has done for us..
So that we stand secure..
We know that we are valuable in your sight..
So valuable that your Son died on the cross to redeem us..
Let our lives demonstrate your way,.
your way of wisdom..
Help us, sisters and brothers in the church,.
we can learn from one another..
Help us to cultivate that kind of spirit of learning,.
spirit of humility, to learn from one another,.
to learn how to fear you,.
so that we can walk in the way of godliness..
Shape us and form us and transform us.
so that your word will be our guidance..
Thank you for being our God, wisdom in Christ's name..
We pray..
Well, God bless you..

$^{641}$And we have some refreshment out there..
And again, I wanna thank you for all your support.
for our trip to Toronto..
I'm leaving Chicago on Thursday..
And we'll come back Monday night..
And pray for us and pray for me.
as I speak to other churches..
All right, I have six messages to give.
and pray for good well-being..
All right?.
Okay, God bless you..
And next week, Chi is preaching, all right?.
So please be here, please be here..
I'm sure he will bring a word of wisdom to all of us..
All right?.
\newpage



\section{約拿書}
\label{sec:v8IvKYM0B_8}
\textbf{約拿一小步 (約拿書 23) - 鍾志愷弟兄}
\newline
\newline
連結: \href{https://youtube.com/watch?v=v8IvKYM0B-8}{\texttt{https://youtube.com/watch?v=v8IvKYM0B-8}} ~~~~ 語音日期: 2024-11-02
\newline
\newline
\hyperref[sec:uDsOylL_mRw]{\small{< < < PREV SERMON < < <}}
~
\hyperref[sec:index_chronic]{\small{[返順時目]}}
~
\hyperref[sec:index_scriptual]{\small{[返順卷目]}}
~
\hyperref[sec:T8J2ny3cosE]{\small{> > > NEXT SERMON > > >}}
\newline
\newline
$^{1}$謝謝Mark Sir的邀請.
今天來到U3都覺得很棒.
很多人都很歡迎.
又有很好吃的早餐.
很棒的咖啡.
我今天講題是《若拿的一小步》.
是關於約拿書二章的.
希望上帝的話語能夠觸動大家.
希望今天你拿走的是上帝去感動你.
親自去跟你說話.
我們為何會叫《若拿的一小步》呢?.
我想先跟大家看看《約拿書》的背景.
我們先看看前文提要.
究竟若拿發生了甚麼事呢?.
若拿一開始在猶大的地方.
他就被上帝呼召.
上帝就跟他說.
你要去尼尼尾城.
去宣講神的話語.
尼尼尾城在東邊.
若拿做了甚麼呢?.
若拿一走,就走到西邊.
直接經過約柏這個港口.
直接要去塔斯.
如果你知道距離.
是一個以色列和西班牙的距離.
意思就是他要橫跨整個地中海.
逃避上帝的呼召.
當然上帝就沒有放過他.
上帝不會那麼容易放過我們.
所以上帝就在途中興起了一個風浪.
那些人就很害怕.
怎麼辦?.
我們這次怎麼做好?.
當那些人很害怕的時候.
若拿就沒有反應.
他只好睡覺.
但接著那些人就開始抽籤.
是誰的錯令我們受風浪呢?.
接著他一抽籤.

$^{41}$抽到若拿出來.
若拿就唯有坦白相告發生了甚麼事.
那些人就問若拿.
喂,怎麼辦?.
我們應該怎麼做?.
若拿就說.
那不如你把我扔下海.
你把我扔下海的話.
風浪就會停止.
那若拿就被人扔下海了.
發生了甚麼事呢?.
上帝猜險了一條大魚.
吞了若拿.
很多人會覺得那條大魚很噁心.
想像一下在口水滿佈的漁福裡面生存.
但你想清楚一點.
這是一個上帝的拯救.
我們今天要看的第二章.
是在一個上帝拯救的場景下發生的.
究竟若拿面對上帝的拯救.
他會怎麼回應呢?.
但同時我們知道若拿是一個不信服的先知.
若拿是一個不喜歡上帝的呼召的先知.
他越走越遠.
他正在經歷的不單只是一個生命的低谷.
他正在經歷的是一個屬靈的低谷.
今天我想大家帶走的是.
如果你今天身處屬靈的低谷.
其實我們有沒有一個轉捩點呢?.
我們如果今天面對一個屬靈的低谷.
我們有沒有辦法走出去呢?.
這個世界只有兩類人.
一是屬靈的低谷裡面的人.
一是將要面對屬靈的低谷的人.
從來沒有人可以避過一個屬靈的低谷.
但我們想知道的是.
其實有沒有一個轉捩點可以發生呢?.
我們一起有個祈禱吧.
我們透過祈禱求神.
開開我們的眼睛.

$^{81}$讓我們知道原來人很真實的場景就是.
我們會不信服.
我們會面對屬靈的低谷.
但神啊你去給我們一個出路.
所以主啊求你給我們心裡面有期待.
求你給我們心裡面.
我們期待你的話語來到我們生命裡面.
去轉化我們.
去更新我們.
請你去憐憫 保守.
以上禱告奉主名求.
阿們.
好.
那我就讀一次約拿書吧.
我讀一次.
我想大家去感受一下.
究竟這篇詩裡面.
約拿是怎樣祈禱的呢?.
請聽我讀.
約拿從餘福裡向上主他的神祈禱.
他說.
我從我的厄困中向上主呼求.
他就答應了我.
從陰間的深處我求告.
你就垂聽了我的聲音.
你把我投到深處.
在海的中心.
水流環繞我.
你的洪濤和你的波浪.
全都掩蓋我身.
而我說.
我被逐出.
離開了你的眼前.
然而我仍要仰望你的聖殿.
縱水圍繞我.
直到喉咽.
深淵環繞我.
海草纏繞我的頭.
我下到山腳.
大地的門山把我永遠關閉.

$^{121}$然而.
你從鄉中把我生命提上來.
上主我的神啊.
我的氣息未蕊時.
我記起了上主.
我的祈禱.
已來到你那裡.
到達你的聖殿.
那些信奉虛無偶像的.
他們的忠誠.
他們離棄.
至於我.
獻祭給你.
我所去的願.
我必償還.
拯救自屬於上主的.
然後上主吩咐的語.
他就套藥拿出來了.
這個就是我們今天要看的經文.
不知道你看的時候.
你有什麼感受.
我想如果今天這個是一個見證分享會.
你最大的感受就是.
哇 藥拿悔改.
不得了.
拍手掌.
是吧.
藥拿說我雖然被逐出.
但是.
我仍然要仰望上帝的聖殿.
藥拿說我要祈禱到達上帝的面前.
藥拿說那些信奉虛無偶像.
即是那幫外邦人.
拜別神的.
他們離棄他們的忠誠.
但是我呢.
我要獻祭和許願.
其實全部東西看下去都不錯.
我們再看多一件事.
藥拿是一個很熟悉聖經的人.

$^{161}$不知道教會有沒有什麼靈劍比賽.
是不是太舊了這件事.
靈劍比賽即是那些叫什麼.
就是背一下金句.
即是你要背一下聖經.
哇 如果藥拿今天參與靈劍比賽.
他一定贏.
為什麼這麼說呢.
他差不多每一句.
都是引用舊約的聖經這樣去說的.
例如第二節.
他說我在患難中求告耶和華.
他就應允我.
詩篇120篇.
不再讀多一次.
因為一模一樣的.
第三節.
你的波浪紅途慢過我身.
詩篇42篇.
又是一模一樣的.
我從你眼前被驅逐.
詩篇31篇.
很相似.
都是我從你眼前.
但是被隔絕.
我想問一個問題.
如果藥拿從來都沒有一個低谷的話.
那他要悔改你做什麼呢.
如果藥拿輸二章這裡.
呈現一個這麼好的見證.
那我們就不用看了.
我們原來是沒有低谷的.
所以我要問一個問題.
今天的講題是不是要轉一轉.
叫做藥拿一大步.
會不會其實藥拿從來都沒有低谷.
他直接悔改就完了.
這樣呢.
接下來我想邀請大家.
跟我一起做一個偵探.

$^{201}$你知道人是很容易說謊的.
人是很喜歡用很多東西去修飾自己的.
我們一起去發掘一下.
藥拿他究竟藏了什麼在裡面.
令我們需要去發掘.
好嗎.
藥拿他說.
第二節.
我從我的困厄中.
困厄中向上主呼求.
他說他呼求.
那我們就要問一下.
他有沒有呼求過呢.
我們去查探一下.
藥拿有沒有呼求過呢.
很特別.
在第一章裡面.
上帝的話去到藥拿那裡.
他說.
你起身.
去尼尼微大城.
你要向那城宣告.
宣告就是呼喊.
那藥拿有沒有呼喊過呢.
沒有.
藥拿去尼尼微大城都沒有去.
當然沒有呼喊.
接著我們看到.
在大風浪裡面.
船長跟他說.
你向你的神呼求吧.
或者你的神會去看顧我們.
藥拿有沒有呼求.
都沒有.
最後有一班人向上主呼求.
是誰呢.
這裡沒有說.
是一班船員.
是一班外邦人的船員向上主呼求.
藥拿有沒有呼求.

$^{241}$沒有.
藥拿從來都不會呼求的.
我們看第三節.
他說你將我投到深淵.
深處.
問題來了.
誰將藥拿扔到海裡.
船員.
我們看到第十五節.
抬起藥拿.
扔到海裡.
是誰叫他們將藥拿扔到海裡.
藥拿自己.
為什麼是上帝將他扔到海裡.
他不怪別人不怪自己.
他在怪上帝.
我們又看看.
他說他被逐出上帝的聖殿.
究竟是誰將他逐出.
第三節說藥拿逃往他施.
離開上帝的面.
重複兩次.
第十節.
藥拿跟他們要坦白.
那些人就很害怕.
害怕什麼呢.
他知道藥拿是離開上帝的面.
有沒有人要逐出藥拿.
沒有.
從來都是藥拿自己逃跑的.
最後.
藥拿說那班外邦人.
信奉虛無神明的.
那些人就不行了.
我呢.
我就獻祭和許願.
聽起來不錯吧.
誰獻祭和許願.
是那班外邦人.
是那班外邦人見到風浪平息之後.

$^{281}$他們跟上帝獻祭和許願.
如果你看到的話.
你心裡有什麼感受.
藥拿說就好聽了.
全部都是錯的.
原來整篇經文裡面.
全部東西都是有一點微微的出錯.
令整件事聽起來很驚險.
實際上是什麼.
口不對心.
我們最容易用這個例子來做一個類比.
這個橙看起來很地道.
可以吃的吧.
一張開的.
全部都是爛的.
我那次看到都嚇一跳.
差點吃掉.
我想藥拿書二章就是一個這樣的場景.
表面很驚險.
實際上是什麼.
就是爛肉一片.
大兄姐妹.
我不知道你有沒有覺得這件事很熟悉.
我看到這張的時候.
我覺得很熟悉.
原因是什麼呢.
好像Mark Sir那樣說.
我是做中學老師的.
我還教宗教課.
我經常都講一些聖經大道理.
經常都跟他們說你要怎麼做.
要怎麼做.
我能不能做到呢.
這是一個問題.
我記得前一年.
我還在教書的時候.
我放工.
我跟幾個學生一起放學.
很開心嘻嘻哈哈.
遠處有幾個學生蹲下.

$^{321}$在坑渠裡不知道在找什麼.
我就問我那幾個學生.
喂 你認識他們嗎.
他們說不認識的.
我們都沒有見過.
那我當然踢爆他們.
同級的 你這樣都不認識.
幫幫他們吧.
那幾個學生說.
真的不認識.
那關我什麼事呢.
那就走了.
走了.
那我又嘻嘻哈哈地走了.
走著走著.
有一點不對勁.
然後我蹲下來.
我問上帝為什麼有這麼不平安的感覺.
上帝在那一刻給了我一個意念.
如果那班同學是同學的話.
那我是他們的什麼.
我是他們的老師.
是宗教科的老師.
但連我自己都沒有去幫的時候.
我竟然要求我的學生去幫他們.
弟兄姊妹 我想說的是.
其實這件事不是很難明白.
因為外表競竭.
內裡安裝.
是我們日常的常態.
這件事在挑戰我們.
會不會有時候其實我們跟神.
跟人.
跟自己的關係.
都已經一塌糊塗.
但連我們自己都不知道.
然後用很多熟悉的華語去包裝.
就覺得自己很OK呢.
會不會有這樣的時候呢.
我想在信仰裡面.

$^{361}$其中一個最難的地方就是.
我們要很真誠地去面對自己.
其中一個最難的地方就是要真誠.
去面對自己的軟弱.
我們的罪.
我們的不足.
這件事應該是最難的事.
而我希望今天我們不是帶著一個.
要批判若拿的態度來到這裡.
因為說真的.
若拿跟我們一樣都是一個人.
若拿會犯的錯.
我們每個人都會犯.
但我相信今天教會給大家的.
是一個出路.
我希望今天不是一個要批判的時間.
而是一個我們可以一起去.
尋找出路的時間.
正正教會就是一個能夠彼此認罪.
互相代求的地方.
我們才那麼喜歡這裡.
正正教會應該就是一個.
我們能夠最真實地去做自己的地方.
我們才會聚在一起.
弟兄姊妹.
我們有沒有機會.
可以去突破自己一點.
我們有沒有機會可以突破自己一點.
去成為一個更加真誠的人.
在弟兄姊妹面前去講自己的軟弱.
我希望這件事是.
我們今天領受到的第一件事.
我們是不是可以真誠地.
去面對自己的屬靈狀態呢.
我希望將這個問題交給大家.
我不知道大家的內心在想什麼.
這件事我很難一看大家的樣子就知道.
但我相信大家最清楚的時候.
能夠自己去反思.
所以我接下來要問第二個問題.

$^{401}$若拿沒有走前一大步.
若拿是不是原地踏步呢.
如果是這樣也挺悲慘的.
如果我們看到一個屬靈的低谷.
然後原來是沒有希望.
原來永遠都是這樣的.
只是一個很虛偽的包裝.
最後原來我們永遠都是腐爛的.
很絕望.
整件事很不開心.
我們再仔細地去看看.
看看約拿書有沒有給我們一個出路.
約拿書的結構是由幾個部分組成的.
有開頭有結尾.
它有兩個正文平行交替著.
第三節有和環去纏繞.
第四節有仰望聖殿.
第五至七節是很相似的.
和環纏繞.
但中間加了一個救贖.
最後都是仰望聖殿.
我們一起做個對比.
我們看看這兩大段的經文.
對在一起的時候發生了什麼事.
第三和第五節很相似.
全部都是關於水的.
去投到海的中心.
水流環繞我.
波浪和鴻圖.
縱水深淵海草.
全部都是水的形容.
同樣都是說到環繞著它.
一個腐爛一個患難環繞著若拿.
但我用紅色去點了一個不同的位置.
紅色不同的位置是什麼呢.
第三節.
若拿說全部都是上帝的錯.
是你將我投到深處.
是你的鴻圖和你的波浪.
第五節還有沒有.

$^{441}$沒有了.
全部都是一個很客觀的形容.
我們看看第四和第七節.
它說我被逐出.
離開你的眼前.
但我要仰望你的聖殿.
我氣息微弱.
我記起上主.
我的禱告來到你的聖殿.
就這樣看是看不出來的.
我想和大家調一調次序.
我們試試看第七節會發生什麼事.
第七節真實的狀況是這樣的.
當微弱時.
我的氣息到莊具.
上主我記起了.
來到你那裡.
我的禱告到達你的聖殿.
分別在哪裡.
分別就是第四節裡面.
全部都是用我來起頭.
是我被逐出.
然後是我要仰望你的聖殿.
第七節.
但有一點點不同.
全部都是將我推後一格.
將我推到最後.
若拿在第七節裡面.
她能夠將自己放得低一點.
都是我做中心的.
但若拿能夠將自己放低一點.
為什麼會做到呢.
為什麼若拿在第三節和第五節.
可以進步多一點.
第四節和第七節可以進步多一點.
關鍵位在哪裡呢.
我相信關鍵位在第六節.
第六節是兩段唯一不同的地方.
第六節是一個我們要看的關鍵位.
第六節說什麼.

$^{481}$我下到山腳.
大地的門山將我永遠關閉.
然而你從坑中把我生命提上來.
上主我的神.
若拿說她尋尋尋.
尋到谷底.
去到大地的門山.
很多時候是在說陰間.
是連陰間都將她永遠關閉.
她面對的是一個死亡的威脅.
那裡再沒有任何的動靜.
全部的東西都靜止.
就這麼多了.
若拿要死了.
但上帝將若拿從坑中提上來.
我相信這個就是重點.
如果今天我們在一個屬靈的低谷裡.
究竟我們的轉捩點在哪裡呢.
如果今天我們在一個屬靈的低谷裡.
我們怎樣才能夠走出來呢.
我深信的是.
我們今天很需要體會上帝的拯救.
如果我們知道上帝不是那個要弄死我們的神.
上帝是那個要救贖我們的神的話.
我們整個觀點會很不同.
如果我們知道上帝要救贖我們的話.
我們會能夠放得自己低一點.
我們會看上帝多一點.
我想跟大家分享一個.
我的弟兄的一個見證.
這個弟兄真的很熟悉.
他在一個屬靈的低谷裡.
這個弟兄真的很熟悉.
他也同意讓我講他的見證.
他在教會裡信主.
很順利地成長.
很穩定地返回教會.
去到某一刻.
他去了外面一個機構事奉.
他遇到一個姐妹.

$^{521}$他一直都很渴求愛情.
打得火熱.
一兩個星期就在一起了.
認識了一個兩個星期就在一起了.
很快.
所有人都勸他認真一點.
認識一下.
當然是沒有用的.
我們這些弟兄.
永遠講的東西都過眼雲煙.
兄弟如衣服.
但是再過一兩個月.
他開始心動了.
對於他來說.
愛情大過天的時候.
他開始想跟她發生性行為.
他跟自己說.
我最後都會結婚的.
為什麼不行呢.
所以他.
再過一兩個月.
他跟那個女孩.
發生了性行為.
他原本已經做好心理預備.
他想去.
justify自己.
他想去告訴自己.
知道自己沒有問題.
但是他真的做完的時候.
他很愧疚.
他在上帝面前.
沒有辦法面對.
他不知道怎樣面對上帝.
他知道自己在扭曲上帝的意思.
但是他不知道怎樣做.
他很慚愧.
他在一個低谷裡面的時候.
他連祈禱都做不到.
但是這個弟兄突然之間.
很突然.

$^{561}$他聽到上帝的聲音.
上帝跟他說.
他按著他的本相.
他按著他原本的樣貌.
上帝接納了他.
上帝接納了我這個弟兄.
他哭得很厲害.
他聽不到.
他沒有辦法一時之間跨過這個軟弱.
但是他很確信的是.
他不可以離開上帝.
當他這麼渴求愛情.
渴求別人的接納的時候.
他見到上帝完完全全地接納了他.
是上帝按著他的本相接納了他.
沒多久他發現不可行.
他就停止了,沒有再去做了.
甚至到現在.
他跟那個女孩斷絕了關係之後.
他跟教會去自首.
他想教會去有一些跟進.
他想教會去紀律.
很奇怪,他很想教會去紀律.
但是他很需要有人去幫助他.
他現在被人紀律.
但是他跟我說的時候.
他說得很開心.
他說我以前好像瞎了一樣.
我看不到.
但是因為有上帝的拯救.
有教會的支持.
他今天能夠走出去.
弟兄姊妹我想說的是.
有時候人的轉變不是一時三刻的.
我們今天的講題是若拿一小步.
我們不奢求我們有一大步.
我們都不想在原地踏步.
有時候只要抓緊上帝的救恩.
我們就可以有一點點的轉變.
不需要多.

$^{601}$一點點的轉變已經足夠.
這一點點的轉變.
上帝會引領我們.
慢慢帶領我們.
一步一步地走出來.
劇透一點.
約拿書是怎樣完的.
約拿書的完結.
不是若拿信服.
「尼尼美人得救」完.
不是.
是若拿去到最後.
他也沒有辦法很順利地.
很順暢地接受上帝的意思.
去和上帝鬥氣.
但是上帝依然和他講說話.
上帝依然用很多時間.
說話去帶領他.
弟兄姊妹.
如果你今天在一個.
屬靈的低谷裡面的時候.
我鼓勵你.
抓緊上帝的救恩.
慢慢來.
慢慢讓上帝帶領你.
走出來.
在這個過程裡面.
你會遇到一個很不同的光景.
我深信上帝.
今天很願意帶領我們每一個.
就算我們很想怨你神的時候.
神會抓住我們.
神會尋找我們.
我們一起有祈禱好嗎.
主啊.
我知道你是那位.
會撇下九十九隻羊.
去尋找我們的神.
所以主啊.
我們今天很希望.

$^{641}$我們能夠向你禱告.
我們每個人的屬靈生命.
有不同的狀態.
如果有低谷的話.
我們求主你去幫助.
我們求你去展示你的救恩.
主耶穌基督.
你親自降世為人.
這樣去拯救我們.
死在十字架上面.
流盡你的寶血.
主啊.
我們很希望你的救贖.
能夠轉化我們的生命.
我們能夠看你.
不再是那位要我們死的神.
你是那位很願意救贖我們.
去引領我們的神.
主啊.
就求你去幫助我們.
抓緊你的救恩.
讓我們的生命.
讓我們的視野.
能夠有些許的轉變.
求主你帶領我們.
每一步每一步.
主啊.
我們很希望好像能夠.
就如若拿一樣.
我們的生命.
要有少許的轉化.
求你去帶領.
保守.
以神禱告.
奉主名求.
Amen.
願主祝福你們.
好,多謝Michael.
希望再叫他回來好嗎.
OK,好.

$^{681}$雖然下星期我不在你中間.
大家有葡撻吃非常好.
我多人多都會吃葡撻.
下星期因為不在.
我想提醒大家.
接下來這兩個星期.
如果你可以投票的話.
請你去投票.
OK.
用你的智慧去投一個公民的責任.
好嗎.
也為我們國家.
為我們城市去祈禱.
好嗎.
我們有一個很短的祈禱.
好嗎.
主,為我們國家.
為我們城市.
為我們社區去禱告.
我知道接下來大選的時候.
都是一個很大的轉折點.
主,我們再次去仰望你.
求你在這個土地.
在這個國家當中.
你去到私人去憐憫.
因為我們去到一個地步.
我們有很多很混亂的東西.
無論道德或者社區.
學校等等.
主,這些都是我們的關心.
我們作為你的子民.
主,我們求你.
給我們盡好我們自己地上公民的責任.
主,為我們在我們的家庭.
在我們的學校.
在我們的社區當中.
我們做一個好的見證.
幫助我們.
我們能夠成為你的一個代表.
將你的愛.

$^{721}$你的恩典.
你自己的美善.
能夠去表揚出來.
我們這樣祈禱奉全命.
好,完善你.
祝福各位.
\newpage



\section{5:17-20}
\label{sec:T8J2ny3cosE}
\textbf{Truly Blessed Life According to Jesus (Matthew 5\_17-20) - Chee Yie Jong}
\newline
\newline
連結: \href{https://youtube.com/watch?v=T8J2ny3cosE}{\texttt{https://youtube.com/watch?v=T8J2ny3cosE}} ~~~~ 語音日期: 2024-11-06
\newline
\newline
\hyperref[sec:v8IvKYM0B_8]{\small{< < < PREV SERMON < < <}}
~
\hyperref[sec:index_chronic]{\small{[返順時目]}}
~
\hyperref[sec:index_scriptual]{\small{[返順卷目]}}
~
\hyperref[sec:Pk_kHNdAhpo]{\small{> > > NEXT SERMON > > >}}
\newline
\newline
$^{1}$Good morning, UB. I hope you guys are extra awake today with the additional hour of sleep.
yesterday night. So, some of you may know, I recently just moved back into the city to.
report, and this is actually the first time that I moved into pretty much like a fully.
unfurnished apartment. And one of the things, one of the lists that I make for my apartment.
is called things that would be nice to have but not necessarily needed. And under the.
list are things like a new desk chair, an external monitor for my laptop, a couch, and.
a couple of area rugs because my entire floor is like pretty much like bathroom tile. And.
it gets really cold, especially heading into the fall and the winter as well. And so, if.
you haven't guessed it yet, that list is pretty much some of the things that I think would,.
in some ways, increase the quality of my life a little better as I live in this apartment,.
right? So, most of us are constantly finding different ways that we think can upgrade the.
quality of our life, whether it is getting the newest iPhone, getting the latest technology,.
or buying a new car, or doing some home projects to renovate your house, make it a little better.
to live in, or even to smaller things like buying a new rice cooker that cooks better.
rice. So, this is because deep down in us, there is always a deep desire to live a good,.
a quality, and a meaningful life, right? So, in fact, there's actually an entire field.
of study that is being dedicated to this goal of helping people live a good life. And that.
field is called positive psychology. And a lot of universities actually in the US now,.
they actually started offering this class in the field of positive psychology as well..
So, this is according to University of Pennsylvania for the Center of Positive Psychology. This.
is how they define what positive psychology is. So, positive psychology is the scientific.
study of factors that enable individuals and communities to thrive. The field is founded.
on the belief that people want to live meaningful and fulfilling lives, to cultivate what is.
best within themselves, and to enhance your experience of love, work, and play. And this.
quest of searching for the good life is not just a new and recent endeavor. But, in fact,.
all throughout history, if we take a look at the ancient Greek and the Rome, we actually.
find many different philosophical schools that actually try to take on this quest of.
what is a good life. So, for example, for the Epicureans, a good life is one that is.
in pursuit of pleasure and the absence of pain. For Aristotle, a good life is a pursuit.
of virtues and development of good personal character. And for Stoics, it is living in.
harmony with the nature and with the cosmos, with the world around us. So, as we can see,.
the questions of what must we do to experience a good life or what is a good life, this is.
actually an age-old question that different people throughout different generations have.
tried to give an account for. But now I want to turn the questions to all of you here..
Who are you putting your trust in that you think will lead to a true and blessed life?.
And if it is not a who, what are you putting your hope in that you think will lead you.
to happiness in this life? So, let's pray together before we dive into our passage today..
Heavenly Father, thank you for your new mercies every morning, and we thank you that we are.
able to come together as a body of Christ this morning to worship you. And I just ask.
that during this time, your Spirit will be with us as we open up your Word and receive.

$^{41}$from what you have to say for us this morning. I also just ask that your Spirit will be with.
me as I speak from your Word. I just ask that your name will be glorified through our worship.
this morning. I pray this in your Son's name, Amen..
Okay, so I'm going to do things a little differently this morning. So, if you have your Bible,.
whether it is on your phone or a physical Bible, I want to welcome you to pull out the.
passage together. So, we'll be looking at Matthew 5, verse 17 to 20 this morning. So,.
if you have your Bible, I would encourage you to pull it out, and we're going to go.
through it together as well. So, I'm going to be reading from the NIV version. So, Matthew.
5, verse 17. "Do not think that I have come to abolish the law or the prophets. I have.
not come to abolish them, but to fulfill them. For truly I tell you, until heaven and earth.
disappear, not the smallest letter, not the least stroke of a pen will by any means disappear.
from the law until everything is accomplished. Therefore, anyone who sets aside one of the.
least of these commands and teaches others accordingly will be called least in the kingdom.
of heaven. But whoever practices and teaches these commands will be called great in the.
kingdom of heaven. For I tell you that unless your righteousness surpasses that of the Pharisees.
and the teachers of the law, you will certainly not enter the kingdom of heaven." So, this.
is the word of God. So, when I preached here last time, back in June, I began a series.
on Jesus' Sermon on the Mount. And our passage today is actually situated right at the main.
body of this Sermon on the Mount. So, again, I want to just give a quick review of the.
structure of this Sermon on the Mount as a whole. So, this Sermon on the Mount begins.
on Matthew 5 and it ends on Matthew 7. And last time, we looked at the introductory narrative.
framework, pretty much setting up the scenes for where this entire teaching is happening..
We saw that Jesus went up on the mountainside and then he began to teach his disciples..
But we also see that there's some crowd present among them as well. And then we look.
at the introduction as well last time, basically the Beatitudes and also the call to be the.
salt and the light to the world around us. And so, in my last sermon, I also highlight.
that this Sermon on the Mount as a whole, it serves as the guiding principle on how.
we ought to live our life in the kingdom of heaven. This is not a set of impossible rules.
that we cannot, like we can never follow, but this is actually an invitation from Jesus.
to a way of living that will lead to true blessing and flourishing. And we also look.
at how Jesus defined what a true and blessed life is through the Beatitudes. And according.
to Jesus, a true and blessed life is a counterculture one. It is one that is being poor in spirit,.
mourning for sins and the brokenness around us. It is one that is being identified as.
humble and meek and so on. So, today we're going to look at how we can experience this.
truly blessed life through our passage today. And this is the big idea that I will be focusing.
on today. Because the truly blessed life can only be experienced through Jesus, we need.
to follow him and live in conformity with him. Again, because the truly blessed life.
can only be experienced through Jesus, we need to follow him and live in conformity.
with him. So, why the truly blessed life can only be experienced through Jesus alone but.
not through other worldly things? In some ways, this is sort of like an exclusive claim,.

$^{81}$right? By saying that Jesus is the only way that we can experience this kind of life..
So, for my sermon today, I'm going to give you two reasons on why the truly blessed life.
can only be experienced through Jesus, and then I'm going to give us an implication.
of what it means for us as followers of Christ. So, I hope you follow along. I'm going to.
give you two reasons and one implication. So, reason number one on why the truly blessed.
life can only be experienced through Jesus is because he is the only one that can fulfill.
the righteous requirement of God on our behalf. So, let's look with me again in verse 17..
So, Jesus began this section by saying, "Do not think that I've come to abolish the law.
or the prophets." So, what is the law or the prophets? The law represents God's standard.
of righteousness and holiness. Or, in other words, it basically refers to what it means.
to live rightly before God. And the prophets, on the other hand, they help point to the.
heart of the law, that it is not just about an outward appearance of obedience to this.
law, but more importantly, it is about the attitude of one's heart. And collectively,.
the laws and the prophets, they refer to the Old Testament scripture in this entirety..
And by declaring that he has not come to abolish the law and the prophets, Jesus sets the tone.
right away at the beginning of his ministry by clearing away all the potential accusations.
and misunderstanding that people might have about Jesus, that he is somewhat anti-law.
or anti-the Old Testament. And then Jesus further clarifies that, "I have not come to.
abolish them, but to fulfill them." Okay, can you guys repeat after me? Not to abolish,.
but to fulfill. Okay, one more time. I think you guys can do better. Not to abolish,.
but to fulfill. Okay, thank you. So, what does it mean that Jesus fulfilled the law.
and prophets and not to abolish them? So, if you think of all the different laws that.
we have, whether it is in your house, at school, at work, or in our society, behind each of.
the law, there is a purpose, right? So, let's take the examples of speed limit in a school.
zone. So, the purpose behind the speed limit at a school zone is not to limit how fast.
you can drive, even though that is kind of like the consequence of it. But the true purpose.
behind a speed limit at a school zone is to protect the students so that driver doesn't.
just drive like 80 miles per hour and like run over school students. However, we can.
be following the law, but totally miss the purpose behind the law, right? So, if you.
have been with us this past couple of months, you might recall one of the announcements.
that we make on both of our English and Cantonese services. Hey, if you see some kids running.
around all the expensive car outside of our buildings, make sure to kind of like tell.
them to like go somewhere else and play, right? So, one time after the Cantonese service,.
I was just like walking towards the bathroom and then I saw some kids chasing around the.
car. So, I just walk over and tell them, hey kids, like don't run around and play around.
here. If you want to play, you can like go to the hallway and like do whatever you want.
over there. But one of the kids, he basically just looked at me and replied, well, I'm not.
going to run around the car then. I'm just going to walk slowly around the car. So, did.
you guys see the problem here? Did you guys see the disconnect right here? The true intent.
of me telling them not to run around the car is so that they don't accidentally like scratch.

$^{121}$or like bump into the expensive car because we have to be like responsible to like pay.
for the damage. But by the kids kind of like responding by, okay, I'm not going to run.
around the car, but I'm still going to play around it by just like walking slowly, that.
pretty much just defeat the purpose of the law, right? That I want them to like not bump.
into the car or like accidentally like scratch those expensive car. So, in the same way,.
Jesus will fuse the law by revealing to us what the true intent and purpose is behind.
the law in the Old Testament through his teaching. And in that way, we can understand what Jesus.
actually means or like what those law actually means. And on the other hand, Jesus will fuse.
the prophets in the sense that Jesus is the ultimate fulfillment of all the Old Testament.
prophecies made by the prophets about the one, the Messiah, who will come and deliver.
God's people and bring about a new relationship between God and his people. So, when Jesus.
say that he fuses the law and prophet, it means that the Old Testament in its entirety.
ultimately points to Jesus. And Jesus is also the official interpreter of the Old Testament..
And this Old Testament law, in light of Jesus coming, they now has a deeper meaning. And.
in my following sermons, I will take a look at some of the examples. If you would just.
want to take a look, if you have your Bible, just look at like verse 21 and the section.
that follows. You will see that Jesus actually used this formula a lot. You have heard that.
it was said, and then he will quote an Old Testament law. And then Jesus will say, but.
I tell you this, and then Jesus would give his own interpretation. So, in that way, Jesus.
bring kind of like a deeper meaning to this Old Testament law. So, why the truly blessed.
life can only be experienced through Jesus? We look at the first reason that he is the.
only one that can fulfill the righteous requirement of God on our behalf. And as a result of that,.
if we believe in him, Christ's righteousness will become our own and we will be made right.
with God. And on this, Paul highlight that. This is from Romans 10 verse 4. Christ is.
the combination of the law so that there may be righteousness for everyone who believes..
And now we're going to look at the second reason. The second reason that the truly blessed.
life can only be experienced through Jesus is because Jesus and his word is the only.
and ultimate authority that will endure forever and lead to the truly blessed life. So, verse.
18, for truly I tell you. So, notice the wording that Jesus used here. For truly I tell you.
or for I tell you the truth. So, right away, Jesus is preparing his disciples for a very.
important truth regarding the word of God that he's about to say to them. And then he.
continues. For truly I tell you, until heaven and earth disappear, not the smallest letter,.
not the least stroke of a pen, will by any means disappear from the law until everything.
is accomplished. And so, to give a little context here, the Old Testament is written.
mostly in Hebrew and a little bit in Aramaic. And when Jesus says not the smallest letter,.
he is referring to the Hebrew letter, Yot, which is pretty much like an apostrophe. And.
when he refers to the least stroke of a pen, he's referring to a stroke that in the Hebrew.
alphabets that is used to differentiate a few of the different characters. But what.
Jesus is doing here, he's pretty much using an argument from lesser to greater. And so.
he said, if he has even the smallest part of the Bible will not be abolished, what about.

$^{161}$the law or the prophets? And so once again, Jesus is saying to his disciples that he has.
not come to abolish the law and the prophets, but quite the opposite. He intends to bring.
the law to his desired goal. That is a true transformation of our heart to the extent.
that even the smallest details in the Old Testament will be fully accomplished in his.
life and in his teaching. And on top of that, he say, until heaven and earth disappear or.
until everything is accomplished, what Jesus is saying here is that as long as our present.
world continue to exist, the word of God will continue to be valid and claims in authority.
over our lives. Until every single passage in the Old Testament is being fulfilled in.
him. So let's step back for a moment. When we hear something like this, that Jesus and.
his words claims his authority over our lives. How does that make you feel? I think for me,.
it actually feels a little uneasy because one of the popular anthems that our current.
generation is singing is, you can be whoever you want to be, or I can be whoever I want.
to be. I don't have to submit to anything, and I get to decide how I want to live my.
life as long as it makes me happy. And on a surface level, this sounds kind of appealing,.
right? It seems like there's a lot of freedom to do whatever you want to do as long as it.
doesn't hurt other people. However, if we actually take a look at the history of Israel.
in the Old Testament, it is not hard for us to actually recognize the consistent threat.
that when people are left to their own to decide what is good and evil in their own.
sight, 10 out of 10 times, it would always lead to downfall and destruction. And if you.
still need more convincing, all you need to do is just to look at the society around us.
and the world around us, and I'm sure you can figure out pretty quickly of what evil.
humanity is capable of when we're being left to our own device. And so, what is the.
alternative here? If we get to decide how we want to live is not the solution, what.
is the alternative here? I would argue it is to live in full submission to Jesus and.
His Word, again, because the Word of God is the only and ultimate authority that will.
lead to a truly blessed life. And just as a parent who knows a lot more than their children.
about what is good and safe for them and then sets certain rules and boundaries to protect.
their children, Jesus, who is all-knowing, knows infinitely more than us as humans on.
what leads to a flourishing life. And so, these rules and laws are put in place not.
to restrict our freedom, but it is put in place out of love to protect us. So, instead.
of seeing this law as restrictive, it is actually a means of grace from God. Instead.
of letting humanity walk down the path of self-destruction, these rules that is from.
the Word of God function as a guide to lead us to a truly blessed life. So, we look at.
two reasons on why the truly blessed life can only be experienced through Jesus. What.
is the implication for us as followers of Christ? To answer that, let us look at verse.
19 together. So, verse 19. "Therefore, anyone who sets aside one of the least of these commands.
and teaches others accordingly will be called least in the kingdom of heaven. But whoever.
practices and teaches these commands will be called great in the kingdom of heaven.".
So, how should we live differently in light of this truth that Jesus is the only one that.
can lead us to a truly blessed life? So, this is what I would say. As Christ followers,.

$^{201}$we need to live in conformity to Jesus and His commands. So, when Jesus say to practice.
these commands, it means to not just submit to these commands, but also to live out, to.
live out the commands as redefined by Jesus. Again, this is not just an outward conformity.
to the law, but it comes from a sincere heart that is being transformed by the power of.
the gospel. And in the following verse, Jesus further clarifies the importance of keeping.
these commands and teaching others as well. So, He say, "For I tell you that unless your.
righteousness surpasses that of the Pharisees and the teachers of the law, you will certainly.
not enter the kingdom of heaven." When it comes to the Torah, the Pharisees and the.
teachers of the laws, they are kind of like the cream of the crop. And the Pharisees particularly,.
they're actually being known for the strict observance of the Jewish law. They don't just.
like follow the laws that God has given in the Old Testament, but they themselves come.
up with extra multiple laws just so that they can be extra careful. So, what Jesus is saying.
here is that, in verse 20, is that, "According to our human standard of righteousness, the.
only way a person can enter into the kingdom of heaven is for their righteousness to surpass.
those of the Pharisees and the teachers of the law." And so, let me put it into some.
perspective for us right here. So, since I moved back into the city, I got into the hobby.
or sports of bouldering or rock climbing. And I actually have a few opportunities to.
bring our youth students from Summit to climb in the past few months as well. And we actually.
just went over yesterday as well. And so, in climbing, the difficulty of the boulder.
problem is usually being great on the V scale, with V0 being the easiest grade, which is.
pretty much like climbing a ladder. And currently, V17 being the most difficult grade that, yeah,.
in the world pretty much, like someone has successfully climbed. And usually, when we.
bring our kids to climb, they will climb in around the V0 to V2 range, and maybe occasionally.
a V3. And in order to encourage them sometimes to try a little harder, some of my counselors.
will tell these kids, "Hey, if you're able to climb this problem, I will buy you a free.
boba." And so, you guys have to understand that for this middle schooler and high schooler,.
getting a free boba is like the best gift that they could ever receive or that they.
could ever think of. So, every time we put those words out, "Hey, if you can climb this,.
you get a free boba," they will just get extra motivated and try really hard to get all the.
way to the top. But imagine if we tell them this, "Hey, I will buy you a free boba if.
you're able to climb this V17 problem," which, to put it into context, there are less than.
10 climbers in the world that have successfully climbed a V17. So, if we tell this to the.
kids, what do you think their reaction would be? They'll probably be like, "That is just.
impossible. What do you want us to do right here?" So, in the same way, trying to earn.
our own way into the kingdom of heaven with our own human effort, that is by having a.
righteousness that surpasses or exceeds the Pharisees and the teachers of the law, is.
almost as impossible as just me telling these kids to climb a V17 if they want to get a.
free boba. But the good news is this, that the righteousness that we need to enter into.
the kingdom of heaven, it does not come from ourselves. It doesn't depend on who we are,.
what ethnicity we are, what socioeconomic status we are, or what ability we have, or.

$^{241}$what effort we put into it. The righteousness that we need to enter into the kingdom of.
heaven comes from Jesus alone because he has fulfilled the law and the prophets. In other.
words, the entire Old Testament scripture is found in fulfillment in the person of Jesus.
through his life and his teachings. So, I've seen one of the preachers did this exercise.
once before. He pretty much went through all the book in the Old Testament and tried to.
show us how Jesus fulfilled each of these Old Testament scriptures. So, I did my own.
similar exercise here, but I'm not going to go through every single book in the Old Testament..
I just pick out a few to kind of try to show us what it means that the Old Testament scripture.
ultimately being fulfilled in the life and teaching of Jesus. So, ready? Here we go..
In Genesis, Jesus is the promised offspring from the woman who will crush the serpent's.
head. In the book of Exodus, Jesus is the Passover lamb whose blood save us from death..
And also, Jesus is the tabernacle that dwells among us as his people. In the book of Leviticus,.
Jesus is the ultimate offering that deals with our sins once and for all. And Jesus.
is also the high priest that become the mediator between God, a holy God, and a sinful people..
In the book of Numbers, Jesus is the bronze serpent in which those who look to him will.
be saved. In the book of Deuteronomy, Jesus is the greater prophet that Moses has promised..
And in the book of Proverbs, which we have looked into the past two weeks with Pastor Mark,.
Jesus is the wisdom personified. In the book of Isaiah, Jesus is Emmanuel, God with us..
He is the suffering servant who was despised and rejected, pierced for our transgressions.
and crushed for our iniquities. And by his wounds, we are healed. In the book of Jeremiah,.
Jesus is the new covenant that the Lord promised. In the book of Ezekiel, Jesus is the one who.
will give us a new heart and a new spirit in order to do the will of God. And in the book of Daniel,.
Jesus is the one like the son of man, coming with the clouds of heaven. And I can go on and on,.
but hopefully you get the point that because Jesus has fulfilled the law and prophets to.
enter into the kingdom of heaven and experience a truly blessed life, it doesn't depend on.
our own righteousness, but the righteousness of Jesus alone. And all we need to do is to.
follow him and live in conformity with him. So I want to come back to this verse in Romans 10,.
verse four again. Christ is the combination of the law so that there may be righteousness.
for everyone who believes. So we don't usually use the word combination in our day-to-day.
conversation, right? But the word combination here refers to the Greek telos, which pretty.
much means the end or the goal. So when Paul say Christ is the combination of the law,.
what he is referring to is that the angle of the law is to lead us to Christ by revealing.
to us our need for greater righteousness that can only come through Jesus. Again,.
because our own righteousness can never be enough to save us. So are you still trying to.
earn your own righteousness? Once again, I want to start, I want to end with the same.
questions that I start this sermon with. Who or what are you putting your trust in that.
you think will lead you to a true and blessed life? Are you still searching for a way that.
will lead to a truly blessed life, but with the intent to do it on your own strength?.
Here's the invitation from Jesus. Come, follow me. So in a few moments, we're going to take.
the communion together as a church and as the body of Christ in response to our message.

$^{281}$today. And my hope is that this will be a time for us to slow down, to reflect on the.
message that we have just received, and also to remember his sacrifice for us at the cross..
And 2 Corinthians 5 21, it says, "For God made Christ who would have never sinned to.
be the offering for our sin, in order that we could be made right with God through Christ.".
So in communion, the bread symbolizes the body of Christ broken for us in order that.
we could be made right with God through Christ. And the cup, it represents the blood of Christ.
poured out for us in order that we could be made right with God through Christ..
And so for those of you who are new here at UV, we have two tables set up with our communion.
elements, one here at the front and one at the back. And before we take the communion.
as the body of Christ, I'm going to put up a couple of reflection questions from my sermon.
this morning up on the screen. And yeah, Phoebe, are you able to go back to my slides?.
Okay, yeah. So I'm going to put these questions up. And what I want you to do is just to,.
yeah, take some time to sit before the Lord and reflect on these questions..
Who and/or what are you putting your trust in that you think will lead to a true and blessed life?.
And in what ways are you trying to earn your own righteousness instead of trusting and submitting to Jesus?.
And during this time, I will have Lucy play a song in the background. This song is called.
"Manifold," and it is actually written by a good friend of mine, Pastor Jordan Brown..
And the idea of this song came from Ephesians 3, 10 to 11. And this is basically a song,.
this song is a call for the church to be the manifold wisdom of God to the world around us.
and to be the presence of God's kingdom here on earth. So, when Lucy is playing in the background,.
whenever you are ready during the song, you can come up and grab the element,.
and then you can go back to your own seat and take some time to reflect on those questions..
And whenever you're ready, take the bread and the cup at your own time..
And after that, feel free to join us in singing this song together..
And one last thing before we enter into communion, I want to speak to those in the room.
who might not identify themselves as followers of Jesus yet..
First and foremost, we really appreciate you being here, and I just want to say thank you for.
coming here and joining us in our worship this morning. But we ask that you not partake in the communion.
because taking the bread and the cup is a public declaration of one's beliefs in the work that Jesus has done on your behalf..
So, if you're not ready to commit your life to Christ yet, I would encourage you instead to use this time.
to really reflect and consider what the good news of Jesus, of the gospel is,.
what Jesus has done on your behalf, and his desire for you to experience the truly blessed life..
So, yeah, you can take some time to sit and reflect for those questions,.
and whenever you're ready, you can come up, take the element on your own, and yeah, take it..
(soft piano music).
♪ Every tribe and tongue ♪.
(soft piano music).
♪ Is wise as one ♪.
♪ To make a tapestry ♪.
(soft piano music).

$^{321}$♪ To be the mystery ♪.
♪ To be the kingdom come ♪.
♪ To every nation seek to love ♪.
♪ Be the bringer of news that our God above ♪.
(soft piano music).
♪ Send His son to the world as to die for us ♪.
(soft piano music).
♪ Every skin and tongue ♪.
(soft piano music).
♪ Form a hasty throne ♪.
(soft piano music).
♪ To make a tapestry ♪.
(soft piano music).
♪ To be the mystery ♪.
♪ To be the kingdom come ♪.
♪ To every nation seek to love ♪.
(soft piano music).
♪ Be the bringer of news that our God above ♪.
(soft piano music).
♪ Send His son to the world as to die for us ♪.
(soft piano music).
♪ To be the manifold ♪.
♪ Wisdom of glory that Christ foretold ♪.
(soft piano music).
♪ Be the church that will spread and then break the whole ♪.
(soft piano music).
♪ Be the family that Abraham longed to sow ♪.
(soft piano music).
♪ Holy, holy God of glory ♪.
♪ Make us worthy, make us lovely ♪.
♪ Holy, holy God almighty ♪.
♪ Make us worthy, make us lovely ♪.
♪ Holy, holy God of glory ♪.
♪ Make us worthy, make us lovely ♪.
♪ Holy, holy God almighty ♪.
♪ Make us worthy, make us lovely ♪.
♪ Make us lovely ♪.
♪ To be the kingdom come ♪.
♪ To every nation seek to love ♪.
(soft piano music).

$^{361}$♪ Be the bringer of news that our God above ♪.
(soft piano music).
♪ Send his son to the world just to die for us ♪.
(soft piano music).
♪ To be the manifold ♪.
♪ Wisdom of glory that Christ foretold ♪.
(soft piano music).
♪ Be the church that will spread and then break the whole ♪.
(soft piano music).
♪ Be the family that Abraham wrought to sow ♪.
(soft piano music).
♪ To be the manifold ♪.
♪ Wisdom of glory that Christ foretold ♪.
(soft piano music).
♪ Be the church that will spread and then breaks sin's whole ♪.
(soft piano music).
♪ Taste of the future we long to know ♪.
(soft piano music).
- Father, thank you..
Thank you for the cross..
Thank you for the free gift of salvation.
that you made Jesus who knew no sin.
to be a sin offering for us.
in order that we could become the righteousness of God.
through Christ..
I ask that this truth will really grow deep in us.
that we're able just to abide in your love.
through Jesus that you have for us..
And so I just ask that because of this truth,.
we're able to live out who you have called us to be,.
to be your kingdom presence to the brokenness around us,.
to be the manifold wisdom that you have foretold.
and be the light and the salt to the world around us..
Thank you, Jesus..
Pray this in your son's name, amen..
So thank you all for joining us for our service today..
We have some refreshment prepared by Kelly outside..
So I would encourage you all to stay after our service.
and hang out and enjoy some light refreshment..
You are dismissed..

\newpage



\section{詩篇}
\label{sec:Pk_kHNdAhpo}
\textbf{這人真是有福 (詩篇 1) - 黃國維博士}
\newline
\newline
連結: \href{https://youtube.com/watch?v=Pk-kHNdAhpo}{\texttt{https://youtube.com/watch?v=Pk-kHNdAhpo}} ~~~~ 語音日期: 2024-11-06
\newline
\newline
\hyperref[sec:T8J2ny3cosE]{\small{< < < PREV SERMON < < <}}
~
\hyperref[sec:index_chronic]{\small{[返順時目]}}
~
\hyperref[sec:index_scriptual]{\small{[返順卷目]}}
~
\hyperref[sec:ywevukHYUvY]{\small{> > > NEXT SERMON > > >}}
\newline
\newline
$^{1}$多謝.
多謝介紹.
其實很開心今天親身來到大家當中.
我其實十幾年前.
2011至2015年在芝加哥.
所以實體認識很多人.
見到又見到朋友.
原來坐在後面.
所以經過差不多十年再回來.
有很多新面孔.
介紹自己.
現在是中神.
香港中國神學研究院的院長.
其實過去我讀博士學位.
就在芝加哥神學院讀的.
今天來到.
這幾天中神的美國董事會.
也有開會的時間.
最興奮的就是能夠回來這裡.
不過可惜見不到Mark.
因為Mark上了多倫多.
我明天就去多倫多他就回來了.
不過不要緊.
我一定抓住他.
一定可以去見面.
所以都是很感恩的事情.
今天我選了一首詩篇.
其實第一篇.
跟大家去到思想的.
不如這樣吧.
我找一些經文出來.
我跟大家讀一次.
讀一讀詩篇第一篇.
我用了《環球聖經》譯本的經文.
詩篇這樣說.
這人真是有福.
他不從惡人的計謀.
不站罪人的道路.
不坐傲慢之人的座位.
相反耶和華的律法.

$^{41}$是他的最愛耶和華的律法.
他就夜脈中.
他像一棵樹.
栽種在原子水道旁.
按時後結果子.
葉子種不枯萎.
他所做的事盡都傳順利.
惡人卻不是這樣.
反倒像糠筆被風吹散.
因此惡人在審判中將站立不住.
罪人在義人的群體中.
也是這樣.
因為耶和華看顧義人的道路.
惡人的道路卻要滅亡.
我們一起去祈禱.
因為你設下你的說話.
你的說話是我們腳前的燈.
路上的光.
求你指引我們.
我們需要聆聽你的話.
所以就今天打開我們的心.
讓我們能夠明白.
更加讓我們跟從你的說話.
祈禱奉耶穌.
祈禱阿們.
這首詩篇開始一句.
就叫做.
這人真是有福.
很吸引.
我們聽到有福.
我們很快就有很多東西想出來.
有福的人是怎樣的.
工作順利.
有兒有女.
身體健康.
長命百歲.
又有財富.
很有福.
所以我們現在很多時候.
廣告.

$^{81}$外面的那些.
都是很想.
你做這件事跟這件事.
那你就可以得到這些事情.
很有福.
不過如果我們這樣讀聖經.
或者這樣理解上帝的說話.
就有些問題.
假設什麼叫有福.
當我們跟著下面的做.
那就有福.
其實說真的.
順便來教會.
你王大仙就差不多.
你知道.
祈福.
你心裡面想要什麼.
就拜一下神.
燒一下香.
他就給你.
所以我們發覺.
基督徒讀聖經.
不應該這樣讀.
這句聖經的翻譯.
新譯本的翻譯是這樣的.
第一句叫做.
有福的人.
然後冒號兩點.
明白嗎.
就是說有福的人冒號兩點.
然後就告訴你.
原來對於基督徒來說.
對於屬於上帝的人來說.
他定義什麼叫有福.
我們要假設.
我們不知道什麼叫有福.
然後才讀這首詩篇.
就會學習什麼叫有福.
所以接著.
就好像查字典的定義.

$^{121}$究竟什麼叫有福.
開始講了.
第一節後面就說.
他不從惡人的計謀.
不暫聚人的道路.
不助傲慢之人的座位.
不過我們現在打開的字典.
給定義.
是不會這樣的.
不會這麼負面.
你今天覺得.
這些教學太負面了.
我們很正面的.
什麼是.
不過特別的地方是.
詩人他定義有福.
用了三個不.
用了三個no開始.
對今天這麼正面的文化.
其實就.
為什麼呢.
我們發覺這三個不.
是有些層次的.
第一是從.
跟從就是走的.
惡人的計謀.
他這樣走.
你就跟著走.
接著就不暫聚人的道路.
就是站在這裡.
走著走著.
站在這裡.
站定定.
不助傲慢之人的座位.
第三個是坐下.
所以你看到.
這三個不是說.
原來我們會走著走著.
走到站在這裡.
然後坐下舒舒服服.

$^{161}$這樣就是一個.
不是.
一個不做這件事.
就是一個有福的人.
其實這裡.
我想我是在說.
原來詩人提醒我們.
這個世界上很多人.
接著他就說耶和華的律法.
他說這個世界.
很多聲音很多東西.
我們自自然然.
習習慣慣.
就會跟著.
跟著就站定了.
然後就舒舒服服坐下.
這樣就.
我們就變成.
不是一個有福的人.
所以這裡是說.
我們要留意這個世界.
很多時候告訴你的東西.
帶著我們走路.
跟你說.
不如你坐下.
舒舒服服.
我們就要小心.
因為一個有福的人.
不是這樣.
這裡是說.
我們世界裡面有很多的.
意識形態價值觀.
是潛移默化地影響我們.
我們不知不覺.
就會落了進去.
自己原來進去了.
都不覺得.
我記得那時候.
十幾年前在美國讀書.
雖然讀博士.

$^{201}$不過有一年的暑假.
我就幫學校的VBS.
上期成班.
照顧小朋友.
搞活動.
我記得有一幕發生了什麼事.
很好.
其中一個活動.
叫小朋友來.
然後他們很有創意.
就說我們有些紙.
就叫小朋友攤在紙上面.
畫他們的輪廓出來.
代表自己.
很好的活動.
然後我們看到.
他們搬了一卷這麼高.
不過卷起來.
好像一卷大廁紙.
這些我在香港沒見過.
美國真的很厲害.
這些東西很大卷都有.
我心裡就很興奮.
我就覺得很棒.
有這麼大張的紙.
可以全部小朋友攤手.
就用來畫.
代表了不同的人.
不過都是在同一卷紙上面.
很興奮.
我以前都是做這些活動.
誰知發覺.
那個大人.
一下就拿剪刀出來.
剪剪剪剪剪剪開.
然後說你做什麼.
開始明白.
其實我發覺原來.
我在香港這麼長大.
我覺得大家是一個群體.

$^{241}$這麼好.
有機會將大家合在一起.
但是這裡的人.
總是要自己一間間房間分開.
這些就是文化.
這就是這個社會.
慢慢慢慢影響我們.
我現在無意說誰好些誰不好.
不過這些就是我們習慣的力量.
我們很多時候不留心.
沒有去aware發生什麼事.
舒舒服服的社會生活.
你就被它影響了.
都不知道.
我們想.
究竟在我們的社會裡面.
有什麼東西.
不斷的潛移默化地影響我們.
我覺得無論香港也好.
或者美國也好.
我們很高舉我們是一個自由的社會.
我們是一個資本主義.
很自由的社會.
很好.
不過我想這個社會.
也影響我們很深.
我們經常說.
資本主義的社會.
最影響我們的是什麼.
或者最好的地方是什麼.
最好的地方是我們自由.
相比那個地方.
很不自由.
我們很自由.
其實是不是呢.
我覺得自由其實不是我們最大的特色.
我不知道大家覺得自己是不是很自由.
其實什麼是我們的特色.
競爭才是我們的最大特色.
我們輸了就不行.

$^{281}$我發覺從小到大.
從讀書那一刻.
到出來工作那一刻.
都在跟別人競爭.
我未必很清楚美國.
不過好像香港.
我十幾年前回到那裡.
小朋友入小學.
哇,真慘.
多難.
贏了再起跑線.
個個都要爭贏人.
你說這樣的環境是不是很自由呢.
其實不是很自由.
有時候會覺得.
甚至這些東西都會滲透我們的工作.
甚至教會生活都會.
有時候我們會覺得.
香港最近很多人移民.
很多教會都少了很多人.
大家都在說怎麼生存下去.
怎麼多點人.
很多這些競爭文化.
都在影響我們.
無論如何.
我想當我們認真去讀上帝的話.
我發覺原來這個世界上.
很多潛移默化的東西都在影響我們.
例如有時候我們會覺得.
不是啊.
我們說.
剛才說到.
這個資本主義社會很多競爭.
聖經都有說競爭.
哥林多前書說.
豈不知在牆上賽跑的都跑.
獎賞的只得一人.
都很說我們要努力去跑.
賽跑.
不過當我們認真去看我們的聖經.

$^{321}$後面那句.
好叫你們都得獎賞.
原來聖經裡面.
即使是說一些賽跑.
都是在說你們.
每一個人.
其實都是得到獎賞.
不是說要打敗別人.
而是說我們在上帝裡面的努力.
為著上帝的緣故.
為著我們要忠心上帝的說話的緣故.
去得到上帝的獎賞.
所以簡單來說.
我想我們發覺.
為何詩人在開始的時候說了三個不.
因為他要我們留意這個世界對我們的影響.
而接著我們就發覺第二節.
第二節他就說.
什麼人是有福.
就不是跟隨世界的方法.
而第二節就是相反.
耶和華的律法是他的最愛.
耶和華的律法他就夜脈中.
所以很明顯我們發覺第一和第二節是對比.
是相反.
第一節是說惡人.
說罪人.
說傲慢的人.
而第二節是說有福的人.
相反.
耶和華的律法是他們的最愛.
當我們讀這篇詩篇.
看第一句的時候.
我們就看到惡人,罪人,傲慢人.
我們第一個感覺就是不是我.
大家坐在這裡記憶圖.
都應該不覺得自己有惡有罪人.
通常都是這樣的.
我覺得我又沒有作奸犯科.
又沒有做壞事.

$^{361}$又沒有殺人放火.
應該不是我.
但特別的地方是.
這裡他將到惡人,罪人和傲慢人相對於相反於愛慕耶和華律法就夜脈中的人.
為什麼.
當一個人不去到愛慕神的律法.
不去到就夜脈中的話.
其實對於詩人來說.
他已經是一個惡人和罪人.
他不是說有些人是罪人.
有些人是好人.
大部分是中間.
其實不是這個意思.
這個意思是當我們不去到愛慕上帝的律法.
這樣去到就夜脈中.
其實就會落入罪人,惡人,傲慢人那種系列.
其實聖經都真的這樣說.
在書中他說那些不冷不熱的人.
又不是這邊又不是那邊的人.
其實上帝不是這樣.
一方面他告訴我們.
當我們不努力去到上帝的說話裡面去愛慕,研讀和跟從的話.
我們自自然然就被世界扯了去.
所以這裡是說原來在我們的信仰裡面沒有所謂中立的東西.
又或者我們的信仰是一種當我們舒舒服服.
不去認真去思想和進行的時候.
就好像逆水行舟被人推了去.
就是會這樣.
所以我想我們都要警醒.
尤其是有時候我來到海外.
無論美國,加拿大.
有時候都要提醒一下第一次.
我自己都在海外住過.
在加拿大有住過,美國有住過.
有時候在這裡很舒服.
很舒服找工作.
很舒服.
因為說真的香港有時候生活難一點.
香港過去幾年有很多困難.
也不是說這裡沒有困難.

$^{401}$不過相對之下.
移民的生活是挺舒服的.
然後我們要提醒自己.
會不會真的太舒服呢.
會不會就走著走著站著.
然後坐下了.
然後就一輩子呢.
所以我們都會問.
在這個地方裡面.
什麼叫做遵行上帝的律法.
究竟這個世界,這個社會給了我們什麼挑戰.
又或者我們今天這個處境裡面.
究竟上帝挑戰我們.
在這個地方裡面.
有些什麼東西是要為上帝去工作.
不從這裡的地方.
不從惡人的計謀.
不暫聚人的道路.
不坐傲慢人的座位.
相反我們要做這些人.
是喜愛耶和華的律法.
是我們最愛耶和華的律法.
就夜脈中.
不過呢.
這裡有個問題.
第二節就說.
什麼是有福的人呢.
耶和華的律法是他最愛的.
耶和華的律法就夜脈中.
不過圖畫呢.
又好像怎樣呢.
好像最有福的人呢.
就多些去退休.
多些去靈修,讀經,安靜.
當然我們去讀經,靈修,安靜.
脈想神的話.
想生命之道.
都是好的.
馬雅說的那些.
好的.

$^{441}$不過幸好.
聖經又不是停留在第二節.
第三節他說.
趴葬一棵樹.
栽在原子水道旁.
按時候結果指.
葉子種不枯萎.
他原來在說.
我們去渴望愛慕神的律法.
就夜思想呢.
還有個比喻.
原來這些說話是要.
讓我們變成一棵健康的樹.
一棵健康的樹栽在水道旁.
按時候結果指.
是在說這棵樹要在世界上.
成為世界上的遮蔭和結果指.
能夠祝福世界上的人.
所以基督徒的生命.
不是停留在我們星期日的聽道.
不是停留在想生命之道.
不是停留在我們平時的靈修.
那種背得出上帝的說話.
而是要成為一棵樹.
生命能夠豐盛地結出果子.
在地上祝福其他的人.
所以我想.
這裡是說.
聖經不斷地跟我們說.
我們譬如新約裡面說.
我們光是業.
本質是光.
本質要咸.
但其實是要進入這個世界裡面.
讓我們從上帝而來的力量.
可以去祝福.
可以在世界上有果效.
這篇詩這樣提醒我們.
所以當我們發覺是這樣的時候.
我們就要想.

$^{481}$今天也有說到.
這個世界很影響我們.
我們要成為一個.
跟從上帝的心意的人.
不過跟從上帝的心意的人.
不只是停留在教會裡面.
不只是停留在星期日裡面.
或者團契裡面.
而是要在我們生活裡面.
能夠流露出這件事的事情.
好了.
然後這個詩篇說就說得很漂亮.
後面那句.
他所做的事全都順利.
而相反惡人卻不是這樣.
反倒像腔比被風吹散.
腔比就是他的骨頭的殼.
乾了風一吹.
Windy City一吹就不見了.
不過這樣.
這個話聽就好聽.
實際生活上經驗就不是跟他們這樣說的.
是不是.
就是說那些愛慕上帝說話的人.
是不是在做的事情很順利呢.
不是的.
有時候愛慕上帝說話.
或者跟從上帝的人.
是很不順利的.
可以被人殺.
惡人可以很順利的.
究竟這兩句是說什麼.
我想如果用這個圖像.
一顆樹栽在溪水旁.
他所做的事盡都順利.
這棵樹栽在溪水旁是說.
他生命底層有種動力.
讓他即使在很困難.
很乾旱的處境裡面.
仍然有力量去成長和祝福其他人.

$^{521}$所以這個順利我想不是解造.
我們在世界上見到那個人得心應手.
有功見都請你.
上什麼功都升職.
很順利.
不是這個意思.
而是即使在乾旱的地方.
生命因為有上帝的說話緣故.
能夠有那份力量.
即使困難都好.
能夠去到結果子.
業子不枯萎.
這裡要說一下我老婆.
有些弟兄姊妹認識我太太.
她做社工.
很好.
其實我回香港差不多十年前.
她做兼職.
因為我們有兩個小孩照顧家裡.
隨著小孩越來越大.
她就越做越多.
終於兩個小孩.
不是小孩了.
變成年青人了.
讀中學.
她終於幾個月前就做了一份兼職.
一做兼職她就比我忙.
我見她做兼職是很辛苦的.
在香港做兼職.
有時候我都替她覺得很辛苦.
不過她又很開心做得到.
其實每晚回來她都有很多投訴.
這個又怎樣那個又怎樣.
但她投訴得來就在笑.
很開心很雀躍.
我想就是這樣.
你問她的工作是不是很順利.
當然不是.
不知道有多少困難.
上面有困難下面有困難.

$^{561}$處境有困難.
不過她很感受到她工作的意義.
她很看得到上帝擺她的位置.
是為了什麼.
所以我覺得她很順利.
不是因為她每天都得心應手.
什麼都解決到.
很舒服地回家.
而是因為她很看得到.
上帝給她的工作.
她也看得到意義的緣故.
即使辛苦也好.
那種力量就好像一棵栽在原子水道旁的樹.
能夠開花結果.
所以我想怎樣去看自己的工作.
她也想順利.
有些事順利.
是不是這件事.
她未必是.
她問我們自己.
究竟我是否因為愛慕上帝的話.
上帝賜我一份很有意思.
能夠祝福世界的工作.
讓我順利呢.
我想這節的意思是這樣理解.
第四節就說到惡人又相反.
惡人不是這樣.
即是不順利.
反倒象腔比被風吹散.
這又不是我們的經驗.
有些人作奸犯科.
做盡壞事.
才會在公司上市.
滿滿滿滿很好.
不過我想這句說話是怎樣呢.
提醒我們用一個上帝的眼光去看世界.
的確今天的社會裡可以有很多不公義的事.
可以有很多看似很順利很有權位的人.
不過你知道他們做很多惡事.
但其實上帝跟他們說.

$^{601}$在神的眼中.
其實他們只是那些腔比.
只是那些被風一吹就散的東西.
或許我們不知道.
想像可能很多人.
可以很風光.
可以很成功.
不過我們不知道的就是.
他們內心是否滿足.
這個我們不知道.
即使他們覺得是也好.
但在上帝的眼中.
其實這些是過去的.
所以簡單來說.
這個意思是我們不要恨.
因為我們會恨.
那個這麼有錢那個這麼成功.
那個就好了.
這麼高位.
不過表面而已.
唯有遵行上帝的說話.
遵行上帝的侮罰.
愛慕上帝的說話的人.
可以不順利.
但從內心.
從裡面.
從上帝的慰養下.
能夠好像栽在原子水道旁的樹.
結果祝福其他人.
所以第四節我覺得.
給我們一份從神而來的眼光.
去看世界的事情.
再下去第五和第六節.
之前說第三和第四節.
說義人.
有福的人和惡人.
在世上的狀況.
而第五和第六節.
就是說他們將來會收什麼成績表.
將來會得到什麼結果.

$^{641}$第五節說.
因此惡人在審判中.
將站立不住.
罪人在義人的群體中.
也是這樣.
審判中.
這裡一定是說到終末.
將主在內.
這個世代過去的時候.
上帝作出的審判.
是一種將來的向導.
所以這裡我想.
另一種給我們今天去面對.
惡人和所謂有義人.
遵行上帝的說話.
但未必順利的人.
這種盼望.
今天.
我們的世界還有很多.
不原本的地方.
不過上帝知道.
上帝記得.
在審判的時間.
公義就要彰顯.
到時惡人將站立不住.
接著一句.
罪人在義人的群體中.
也是這樣.
這句話.
不只是等到終末.
這裡說原來在義人的群體.
義人就是說到那些.
愛慕上帝的說話.
願意遵行.
就像一棵樹栽在水邊的人.
所以很重要.
我們發覺今天.
如果我們跟著這首詩篇.
活出來.
我們發覺我們要遵行上帝的說話.

$^{681}$世上未必順利.
我們看到身邊有很多.
惡人很順利.
令自己很辛苦的時候.
原來義人的群體.
能夠給予我們一份支持.
這提醒我們教會群體做些什麼.
來到教會群體我們很開心.
吃早餐聊天.
很好.
不過這句話提醒我們.
在這些交朋友.
這些開開心心之外.
其實我們需要.
教會的群體.
一起去思想.
一起去分享.
如何在一個.
跟上帝說話相反的世界裡.
基督徒作見證.
是可以很困難的處境裡.
彼此鼓勵.
彼此支持.
這是令我們能夠繼續.
在這個地上活出上帝說話.
一個很重要的群體.
我們想想究竟我們教會群體做了些什麼.
究竟我們全齊.
是否只是開開心心吃飯.
我想不單止是這樣.
最後一句.
因為耶和華.
看顧二人的道路.
惡人的道路.
卻要滅亡.
這裡看顧二人的道路.
不只是說忠沒耶穌在內.
看顧這個字.
原文是認識.
深入地認識.

$^{721}$所以這句話.
意思不是說上帝在天上.
很遠地看著你在受苦.
而是說上帝今天已經很深地認識我們.
這件事成為我們很大的安慰.
最後我想說一個見證.
有一位姐妹.
基督徒.
她叫林超.
上世紀在中國大陸.
五,六十年代.
她是因為自己信仰的緣故.
在一個這麼困難的環境裡.
被迫.
她在社會反對基督教的原因下.
她堅持信仰.
甚至因為信仰的緣故.
去說一些政府很不喜歡的事.
因此被人抓去坐牢.
她說坐牢更開心.
因為坐牢已經不可以再抓我.
在監獄裡繼續傳福音.
很自由地傳福音.
令她怎樣呢.
令她單獨囚禁.
很苦.
其實她的事情是記下.
現在可以看.
林超這位姐妹.
她說為何她在一個人裡.
這麼困難.
她感受到很多惡邪惡勢力的壓迫下.
她仍然一個人在監獄裡.
仍然好像一顆栽在溪水旁的樹木.
她說她每個星期.
每天都會有一個一人大敬拜.
她自己一個人.
但她親近上帝.
而她很知道上帝認識她.
很知道她的心.

$^{761}$很知道她的忠誠.
因為這樣的緣故.
她心得安慰.
其實最慘的.
最後她也是被殺害.
但很感恩.
因為她自己寫下的一些日記.
書信留下了.
讓我們今天知道.
我們知道甚麼是有福.
這首詩篇挑戰我們甚麼是有福.
挑戰我們怎樣看待有福.
有福的人生究竟是甚麼.
是否舒舒服服 順順利利.
長命百歲.
身體健康 兒孫滿堂.
不用憂.
當然上帝可以這樣祝福我們.
但我們是否追求這些.
那些是否得福.
說真的不是在我們手裡.
是上帝的安排.
但我們能夠決定.
我們能夠決定.
我們愛慕上帝的話.
走在這個世界上.
祝福其他人.
叫我們不要經常在這個世界舒舒服服.
而是問一個問題.
上帝挑戰我們.
今天怎樣去活出一個.
毫不一樣的生命.
即使世界有很多邪惡.
因為靠著上帝的緣故.
我能夠活出一個有力.
有精彩的人生.
詩篇這樣挑戰我們.
上帝挑戰我們.
我們今天怎樣去回應上帝呢.
我們求主幫助我們.

$^{801}$我們一齊去祈禱.
天父上帝多謝你.
因為你告訴我們.
一個有福的人生是怎樣.
你不單止告訴我們.
耶穌基督甚至走了出來.
我們見到祂.
就是活出天父的豐盛.
不過在世上一點都不順利.
然而祂能夠依靠你的緣故.
有力量地做你的工作.
叫祝福降臨在世界裡.
求你都這樣去使用我們.
挑戰我們.
讓我們活出一個不一樣的生命.
在我們自己的環境.
工作裡面.
社會裡面.
家庭裡面.
朋友裡面.
能夠跟從你的說話.
成為身邊的人.
祝福.
讓我們即使有很多困難.
很多世上的.
甚至一些傷害.
一些犧牲都好.
仍然知道.
這是最有福.
最豐盛的生命.
幫助我們.
求你聽我們的祈禱.
祈禱奉耶穌基督的名.
\newpage



\section{}
\label{sec:ywevukHYUvY}
\textbf{Who is Your Lord? \#2 Belshazzar (Daniel 5) - Rev. Mark Chan}
\newline
\newline
連結: \href{https://youtube.com/watch?v=ywevukHYUvY}{\texttt{https://youtube.com/watch?v=ywevukHYUvY}} ~~~~ 語音日期: 2024-11-19
\newline
\newline
\hyperref[sec:Pk_kHNdAhpo]{\small{< < < PREV SERMON < < <}}
~
\hyperref[sec:index_chronic]{\small{[返順時目]}}
~
\hyperref[sec:index_scriptual]{\small{[返順卷目]}}
~
\hyperref[sec:KOpSYq8TJvA]{\small{> > > NEXT SERMON > > >}}
\newline
\newline
$^{1}$- Alright, before we turn to scripture,.
why don't we spend some time to pray together, alright?.
Why don't we spend some time to pray?.
And before we pray, I want to ask you a question,.
you know, in the presence of God..
Is there anything, something in your heart,.
in your mind you have to unload?.
Alright, maybe some worries, maybe some issues,.
maybe some challenges in your life..
Maybe some anger, some depression,.
some of the unhappy feelings..
God can take it, you know?.
And this is the best moment we can come before God, alright?.
Let's spend a minute or two, pray to God..
Father, we do come before you knowing.
that you are our loving God..
And you can take all the garbage in our mind and our heart..
No matter what that is, with the worries that we have,.
the despair that we sense, or the challenges.
that we are facing, and we are exhausted,.
we are just drained..
We can come before you knowing.
that we can cast our burdens to you..
Thank you for being our loving God..
Thank you for even taking all this emotions,.
even so-called negative emotions..
You can take it..
And thank you for being there to listen to us,.
to being with us..
Thank you for the fact that you are leading our church,.
leading our lives, and we come before you to know that,.
to acknowledge that you are our strength..
Grant us the strength that we need..
For those who are studying, for those who are coming.
to the end of the semester, with final exams.
and all the assignments, and I pray that you will grant.
each and every one of us the wisdom, the strength,.
the discipline to go on, knowing that your grace.
is always sufficient..
For those who are struggling with family matters.

$^{41}$and the demands from our work,.
we do know that you are our rock,.
and we can depend on you, we can rely on you..
And that's why we come to worship you..
That's why we come to give thanks to you,.
to count the blessing that you are still providing for us..
And thank you very much for being our father,.
being our God..
And most of all, thank you for the love and salvation.
that we have in Christ Jesus, that we are your people..
We are your sons and daughters..
Thank you, in Christ name we pray..
Now we come to our mini series.
on the pagan kings in the Bible..
And last week, we talked about one of the pagan kings,.
who is Cyrus the Great..
Cyrus the Great is a Persian king..
And today we're going to chapter five.
of the book of Daniel, all right?.
If you have your phone and you have your Bible,.
please turn with me to chapter five of Daniel..
If you can find the book of Daniel,.
it is in the Old Testament,.
especially in the prophetical books,.
one of the prophetical books, all right?.
So we're gonna talk about this king,.
it's called Belsasha..
Let's go back to this timeline..
I kind of zoom into this timeline we share from last week..
And we talk about Cyrus..
You can see Cyrus on here, right in the middle,.
right at the bottom..
Cyrus, who became the emperor of the Persian kingdom..
And at the same time, that was the last,.
very last period of the Babylonian kingdom,.
the Babylonian empire from 626 to 539..
And you see on top, there's a bar..
There were two kings, right?.
Nebuchadnezzar and Belsasha..
And I talked to you about how we going to approach.

$^{81}$studying the figures in the Bible,.
especially the pagan kings..
First of all, we have to profile them, right?.
We have to profile them from scriptures.
or from sources outside of the Bible..
And then we come to the Bible,.
and let's see what the Bible says about pagan kings,.
about this or this king..
So in the end of Babylonian empire,.
there were two kings ruling the empire..
The top one is Nebuchadnezzar,.
and the guy below him is his eldest son,.
his firstborn son, Belsasha,.
which is mentioned in the book of Daniel chapter five..
And this is a funny story,.
and I want to go through this story with you..
But you need to understand.
at the end of the Babylonian empire,.
when the power of Babylon is declining.
and the empire is coming to a crumbling.
and collapsing, and the last king is Belsasha..
Belsasha is the co-region of his father, Nebuchadnezzar..
And we are coming to the year 539, 539..
And some of the background here,.
when Nebuchadnezzar became king, it was year 555..
And for some reason, he ruled the Babylonian kingdom.
for a few years, and then he left the capital..
And then he moved to an oasis in northern Arabia.
for 10 years..
And he's kind of a very strange character,.
even for historians,.
is trying to understand why he moved away from the capital..
And then when he moved away from the capital,.
he gave his kingship and trusted kingship.
to his son, Belsasha..
So that's why we have two kings there..
The top guy, number one, is Nebuchadnezzar..
And then the second guy is called Belsasha,.
his firstborn son..
And practically, Belsasha in the Bible or in history.

$^{121}$is the last king of Babylonian empire..
So that's when we come to chapter five of Daniel..
And this is a funny story..
So I want to go through some of the content of this story.
and to think about what the Bible or what God sees.
in this king..
Chapter five, verse one..
King Belsasha gave a great banquet.
for a thousand of his nobles and drank wine with them..
While Belsasha was drinking his wine,.
he gave orders to bring in the gold and silver goblets.
that Nebuchadnezzar, okay, Nebuchadnezzar,.
Nebuchadnezzar's father had taken from the temple.
in Jerusalem so that the king and his nobles.
and his wives and concubines might drink from them..
What do you think?.
As a big banquet, I don't know if you've been.
to a big banquet before..
The biggest one I've been to is like 100 tables..
It's huge in a hotel conference room..
You know, people are so happy and we just drink.
all we want, open bar and all that..
It's kind of like this situation here..
Belsasha, chapter five of Daniel..
So he gave a great banquet to thousands of his noble.
and all his cabinet and new government officials.
are gathered together, his wives and his concubine.
and everyone is drinking..
And then, you know, all of a sudden,.
he gave the order to have the vessels.
that his so-called father or the founder.
of the Babylonian empire, Nebuchadnezzar..
Oh my God..
Nebuchadnezzar..
And then he captured the temple of Jerusalem..
You remember that?.
586 BC..
And then he took some of the vessels,.
gold and silver vessels, or maybe the Ark of the Covenant.
as well, so all the valuables, all the precious things,.

$^{161}$he captured from the temple and he captured them.
to his capital, Babylon..
Now, Belsasha came much later..
He was having a big banquet and then he was so happy,.
I guess, and then he used some of the vessels.
from the temple of Jerusalem..
So it's a big banquet and then he used some of the vessels,.
you know, cups and plates and instruments.
from the Jerusalem temple to use them to drink wine..
You see?.
To drink wine..
Now, a little bit background here..
Whenever the Bible talks about banquet,.
it's not just a wedding banquet..
It's not just a happy party together..
It always has religious and political meaning..
So when the king gathers his cabinet, right here,.
and he make offerings to all the gods..
So when they get together having food and feast together.
and then drink together, it has political.
and religious reason..
And the reason being, they want to please the gods..
Please the gods of the kingdom..
Please all the gods and the purpose is to secure.
the approval of the gods..
It's to secure their blessing,.
especially for the coming year..
So usually this kind of banquet, this kind of feasting.
happens in New Year..
In the New Year, they want to secure the approval.
of the gods for the coming year..
So that's why they gather the temple vessels,.
not only from Jerusalem, but from other places..
Remember Babylon captured, you know,.
pretty much the whole world back then..
And then they destroy some of the temples.
and they capture their gold and silver vessels.
and brought them to the capital..
So this is the scene that, you know,.
not only the Jerusalem vessels are being used,.

$^{201}$all these other gods' temple vessels.
are being brought together for them to drink wine..
And what they do is that they will offer sacrifices.
and libation to the gods for the purpose of,.
you know, pleasing the gods so that the gods will bless you..
So in short, the purpose is to secure.
the approval of the deities, to please them.
so that they will bless his plan and bless his policies.
for, especially for the coming year..
So in such a banquet, the king would expect the gods.
to be happy, you know, to rob us of their plans..
So here, this story, the funny thing is,.
when the king expected the gods to say something,.
God did show up..
And that's the most surprising thing..
Verse five and following..
And suddenly, the fingers of a human hand appear.
and roll on the plaster of the wall.
near the lampstand in the royal palace..
And the king watched the hand as it rolled,.
his face turned pale and he was so frightened.
and his knees knocked together and his legs gave way..
Now that's the funny thing..
You expect gods to say something.
and then here the true gods show up, all right?.
So this is a story of the origin of the English idiom,.
writing on the wall, right?.
And that's where it come from, writing on the wall..
What he saw was there's a hand writing something on the wall.
and it is pretty interesting and also surprising..
If you've been to like British Museum.
or there's a very, very good museum in Chicago,.
you know that?.
It's called Oriental Museum, University of Chicago..
Anyone been there?.
No?.
Oh my God..
Oh, you've been there, good..
That's great, that's great..
Yeah, there's only one person here..

$^{241}$This is the Oriental Museum in University of Chicago.
is a world-class museum and it's also free..
You know, it takes free donation..
Someday maybe we should go there for a field trip, all right?.
There's tons of things you can see..
If you pay attention to the wall reliefs.
in the Assyrian or Babylonian gallery,.
there's no space on the wall..
It's not like our side drywall..
There's no spaces on the wall you can put something on it.
because every single inch of the wall.
is filled with either the king's figure.
or the god's figures or trees and the lion hunt.
and all that stuff..
There's no space in the wall.
and that is the most surprising things..
There's no space in the wall.
but all of a sudden a human hand show up.
and trying to write something on the wall..
Now, this is from Daniel five..
And then once they saw the writings,.
no one can interpret the writings on the wall..
And you can imagine scenes like this..
This is the famous painting by,.
in the 17th century,.
talking about this story..
When the king look to the wall,.
all of a sudden there was some writings showing up..
And they're trying to interpret the meaning.
or the significance of those writings..
And 518, it says,.
Daniel is being called up.
because he is the only guy in the court,.
in the empire, who can interpret like this kind of writings..
What does it mean?.
So he's being called up to the king.
and Daniel said,.
O king, the most high God gave your father, Nebuchadnezzar,.
sovereignty and greatness and glory and splendor..
And he, Daniel also talk about how God had humble his father.

$^{281}$or great grandfather, Nebuchadnezzar,.
when he was filled with pride..
That is from the previous chapter, Daniel four..
And now, writings show up on the wall,.
writings on the wall..
And Daniel said,.
no, you have forgotten how God has humble.
your great grandfather..
518..
And then 22..
He said, but you, his son,.
as his later generation,.
O Belsusha, you have not humbled yourself,.
though you knew all this..
Instead, you have set yourself up against the Lord of heaven..
You praise the gods of silver and gold.
and of bronze, iron, wood, and stone,.
which cannot see or hear or understand..
But you did not honor the God who holds in his hand.
your life and all of your ways..
You have seen how God had humble previous king,.
and now you are filled with pride.
because you took the vessels from the Jerusalem temple.
and you thought that God would approve you.
or God would be pleased..
No, Daniel said, God is not pleased..
524..
Therefore, God sent a hand that wrote the inscription,.
and this is the inscription that was written..
Many, many, tickle, passing..
What do they mean?.
Of course we don't know, right?.
Many, many, tickle, passing is Aramaic..
Aramaic, it's not Hebrew, it's not Greek..
Part of the Old Testament is written in Aramaic,.
and many, many, or many, many means number,.
counting, one, two, three, four, five..
And it's repeated twice to stress that.
God's decision is certain..
It will be fulfilled because God is counting the numbers..

$^{321}$What kind of numbers?.
We have to interpret it..
And then tickle means waiting..
Waiting, like putting things on scale, you know?.
Passing means divided, divided or separation..
So the message literally means numbered,.
numbered, weighted, and divided..
And those words are, as I said, in Aramaic,.
and Aramaic was the international language of the time..
Remember what year was that?.
539, right?.
Sixth century BC..
In the ancient New Year's world,.
the international language at that time was Aramaic..
So the king and the officials were actually able.
to read Aramaic..
Many, many, tickle, passing, they can read it..
What they don't understand, or they need someone.
like Daniel to interpret, is what it means..
What does it really mean?.
What those words mean or signify?.
So they need some help, and Daniel is trying to help them.
to give them the interpretation..
What does it mean?.
Many, many, tickle, passing?.
Number, number, weighted, and divided..
And Daniel comes along and he said,.
"This is what these Aramaic words mean..
"Many, God has numbered the days of your reign.
"and brought it to an end.".
Tickle, you have been weighed on the scales.
and found wanting, so you don't match up..
Paris, or passing, your kingdom is divided.
and will be given to the Medes and Persians..
So these words are actually God's evaluation.
of Belsash's kingdom..
And God is using these words to pronounce.
his upcoming judgment on Belsash..
529, "Then at Belsash's command,.
"Daniel was clothed in purple,.

$^{361}$"a gold chain was placed around his neck,.
"and he was proclaimed the third highest ruler.
"in the kingdom.".
Now, why the third?.
Why number three?.
Remember the background?.
The top guy was Nebuchadnezzar..
Nebuchadnezzar..
The second guy is Belsash, right here..
And then because Daniel could interpret.
the significance of the message, or the words,.
he is now being promoted to the third..
You see what I'm saying?.
Now, before the archeologists find the name Belsash,.
people are doubting about the reliability of scripture..
They were saying there's no king called Belsash in this world.
and then all of a sudden, some British guy.
came to the Middle East some years ago.
and then he was taking in the archeological site.
and then they find some cuneiform tablet,.
some clay tablet, and then they translate it.
and then Belsash was being mentioned.
as the eldest son of Nebuchadnezzar..
And now we know why this verse says Belsash.
is raising Daniel to the position number three.
in his kingdom..
Because Nebuchadnezzar, Belsash, and then Daniel..
Now this is also interesting..
Not only the reliability of scripture,.
but it also makes perfect, kind of interesting.
is that Belsash knows that these are the meanings.
of the words I will make, that God is counting down.
his kingdom, but Belsash did not know.
that his kingdom is coming to an end so quickly..
You know, faster than he imagined..
In the last verse of this passage, verses 30 and 31,.
that very night, Belsash, king of Babylonians,.
was slain and Darius the Mede took over the kingdom.
at the age of 62..
Now Darius the Mede is also another question here,.

$^{401}$but it should refer to Cyrus the Great..
So that night, that night, that very night,.
Cyrus came and Persian kingdom took over Babylon..
And from secular history, we know we can even.
find out what that day is..
It's October 12th, 539 BC, October 12th..
We can calculate so precisely that we know October 12th,.
the army of Persian kingdom march into the capital city,.
Babylon, of Babylonian empire,.
and then the empire crumbles and it fell..
It happened overnight..
Now what are the take-homes from this story,.
this interesting story from Belshasha?.
First, it says clearly that God is sovereign.
over all earthly powers..
All the earthly kingdoms are in God's hands.
and God's power is transcendent over all of them..
Any kingdom, in the ancient time, Babylonian kingdom,.
Persian kingdom, Greek kingdom, the Roman empire,.
any empire, even today, God's power.
is over all earthly kingdoms..
Maybe United States, Russia, China, EU,.
all the earthly kingdoms, all of these earthly kingdoms.
and treasures come and go, rise and fall of empires..
And so often we think our kingdoms will last,.
will continue forever, as long as we live..
No, the Bible says, here the Bible uses a story.
to tell you the theology or the philosophy of history..
This is how God sees history,.
how God sees historical kingdom..
They rise and they fall..
No king, no kingdom on earth will last forever..
Only God's kingdom will last forever..
Now this is our pride and our presumption.
that we think our earthly kingdom will last forever..
The story tells us no earthly kingdom can last.
without God's permission..
God can allow a certain king, a certain official.
to be in the White House..
God can allow this or that king to rule.

$^{441}$over some part of the earth..
But everything is in God's hand..
Now what God allows does not always.
mean God approves of it..
God can allow evil to exist..
That's pretty difficult for us..
God can allow evil to exist..
However, there's always judgment at the end..
There's always righteous judgment at the end..
And judgment often comes as a surprise..
And you would not expect that..
Here in this story, Belshazzar the king.
thinks that his kingdom will go on.
at least for a little while..
Why God, of course God is counting down his days,.
but he thinks his kingdom will go on.
and business as usual..
Now the king even expects that God.
to show up to approve his plan..
And in this chapter, God does show up..
But God does not approve his plan..
Now think about our lives..
Sometimes we want the approval of God, right?.
Maybe we're launching a new career..
Maybe we have a family plan.
or we have our own personal plan..
We always want God to show up and give us the approval..
It may, it may not..
Because it is up to the will of God..
The second point here is that the will of God.
transcends our will..
Divine will transcends your will..
God's will is over above your will..
And so often we want the approval of God.
because we like to do this or that,.
or we plan for this or that..
And sometimes we even think that.
if we pray hard enough,.
if we work hard enough,.
then our intent, our plan will work out..

$^{481}$God may approve my plan..
And this is a real temptation..
This is a real temptation for us..
And it's pride on our behalf..
It's just tempting to even sometimes.
we misuse the name of God for our plan..
Instead of humbly seeking the will of God in our lives,.
sometimes we think we can misuse or abuse the name of God.
so that we'll get the approval from God..
You know, this story tells us the job of God.
is not to rubber stamp your plan..
God's not here to do a bidding..
God's not here, not necessarily to approve your plan..
Even you have well intent..
It may, may not, right?.
The whole point is God is not here.
just to rubber stamp our plan..
God is here to help us and mold us.
so that we can yield to him in worship and obedience..
That's what God wants..
You know, he's not necessarily here to say,.
"Okay, you can do this..
"Okay, if you ask me, I will do that for you.".
The whole thing is because the will of God matters,.
not our will..
Now I wanna close this message,.
this story with James chapter four..
James chapter four..
When we read this two verses,.
I think it speaks to what Daniel 5 has to say..
Is that the will of God is higher than our will..
The plan of God should transcend your plan..
James chapter four..
Why?.
You do not even know what will happen tomorrow..
What is your life?.
You are a mist that appears for a little while.
and then vanishes..
Instead, you ought to say,.
if it is the will of the Lord,.

$^{521}$we will live and do this or do that..
No, the scripture is trying to confront our pride,.
confront our priority..
We don't even know we will exist tomorrow..
It says we are simply a mist..
You know what a mist is?.
Did you wake up early enough this morning?.
You look at your alarm, the mist,.
actually talking about our breath, you know..
When we say things, when we're breathing out, you know,.
that is mist..
The mist will exist for just a little while.
and it will be gone..
And that is your life, that is my life, right?.
Then why, if we exist only like a mist,.
why do we think this will happen tomorrow.
or that will happen tomorrow, right?.
What is your life?.
James challenges and confronts us..
You ought to say, rather,.
if God allows,.
if this is the Lord's will,.
then we will live..
We will do this or do that..
No, the Bible is not encouraging us.
not to manage our life..
It's not saying that we don't have to plan.
for the next day or the next year or next month, you know..
The Bible is all about, you know, management..
You know, we need to be good steward, right?.
However, the Bible also insists that,.
hey, God's will transcends your will..
If God allows, if this is the Lord's will,.
you will live and you will do this or do that..
Why don't we close in prayers.
and I'm sure you have some plans in your life, right?.
And so often, you know, our plans trump God's plan, right?.
Oh God, I have this or that I need to do, okay?.
I'll put you on the side, all right?.
Put you on the side for a moment..

$^{561}$And God, I know you're a king, but hey, you know,.
I have this or that to do, you know, I'm so busy..
I'll put you in the footnote..
This is tempting, all right?.
This is tempting..
As James says, you ought to say, you ought to pray..
If it is the Lord's will, I will live.
and I will do this or that..
Let's be humble and come before God.
and do away our pride to acknowledge His sovereignty, right?.
Let's pray..
Father, we come before you to confess our sin, our pride..
And so often, our priority trumps your priority..
Sometimes we just think that our plan is so important.
and you are the least important..
We come before you, we confess our sin of pride..
Yes, we don't control our lives..
We don't even know that we will exist tomorrow..
We're just like mist..
We appear a little while, we live for a little while,.
and then we'll be vanished and be forgotten..
And yet you, in your grace, allow us to live..
So yeah, we come before you and we pray..
And we say to you, if it is your will, then we will live..
We will do this and do that..
And Father, confront us..
Do we really have the priority of a Christ disciple,.
putting you as our number one,.
putting you as the most important,.
the most important person that we're gonna follow.
and we're gonna obey and we're gonna worship?.
Father, we come before you,.
we ask that your spirit will remind us.
to reorganize our lives,.
putting you first, putting other things.
in the lower priority..
Help us to please you,.
help us to worship you and follow you,.
and to be obedient, always seeking your will first,.
to listen to word of wisdom,.

$^{601}$to listen to your spirit's prompting,.
to help us to live a life that is obedient,.
live a life that will pleasing to you..
Yes, it is your will..
If it is your will, then we will live.
and we will do this and that..
And thank you for your blessing,.
especially when Thanksgiving is approaching..
We do need to acknowledge that this is your will,.
this is your permission,.
that we can come and enjoy your provision..
Yeah, you are so good, God..
Thank you, thank you for the grace of life.
and help us and use our life to honor you..
In Christ, we pray..
(gentle music).
- Well, we have some freshmen out there.
and please stay a little bit for time of fellowship..
And as Kelly remind us, get your tickets today.
and we're gonna have a feast together.
and I'm sure there's no writings on the wall..
So we're gonna enjoy our time together.
in times of Thanksgiving..
And also I wanna remind you,.
Thanksgiving is a great reminder for us,.
this is God's will and God's provision in our life.
to count His blessing, all right?.
Count His provision and thank God.
for everything that we enjoy today, all right?.
Thank you very much and I'll see you next week..
(gentle music).
\newpage



\section{馬可福音 3:1-6}
\label{sec:KOpSYq8TJvA}
\textbf{新約中的情緒 (1): 憤怒 (馬可福音 3\_1-6) - 何建邦傳道}
\newline
\newline
連結: \href{https://youtube.com/watch?v=KOpSYq8TJvA}{\texttt{https://youtube.com/watch?v=KOpSYq8TJvA}} ~~~~ 語音日期: 2024-11-19
\newline
\newline
\hyperref[sec:ywevukHYUvY]{\small{< < < PREV SERMON < < <}}
~
\hyperref[sec:index_chronic]{\small{[返順時目]}}
~
\hyperref[sec:index_scriptual]{\small{[返順卷目]}}
~
\hyperref[sec:dQy1lcfV2Eg]{\small{> > > NEXT SERMON > > >}}
\newline
\newline
馬可福音 3:1-6
\newline
\begin{longtable}{cl}
\hline
\hline
章節 & 經文 (和合本修訂版)\\
\hline
3:1 & \begin{tabularx}{0.7\textwidth}{X} 耶穌又進了會堂，在那裡有一個人，他的一隻手萎縮了。 \end{tabularx} \\ \\ \relax
3:2 & \begin{tabularx}{0.7\textwidth}{X} 眾人為了要控告耶穌，就窺探他會不會在安息日醫治那人。 \end{tabularx} \\ \\ \relax
3:3 & \begin{tabularx}{0.7\textwidth}{X} 耶穌對那手萎縮了的人說：「起來站在當中！」 \end{tabularx} \\ \\ \relax
3:4 & \begin{tabularx}{0.7\textwidth}{X} 他又問眾人：「在安息日行善行惡，救命害命，哪樣是合法的呢？」他們都不作聲。 \end{tabularx} \\ \\ \relax
3:5 & \begin{tabularx}{0.7\textwidth}{X} 耶穌怒目環視他們，因他們的心剛硬而憂傷，就對那人說：「伸出手來！」他把手一伸，手就復原了。 \end{tabularx} \\ \\ \relax
3:6 & \begin{tabularx}{0.7\textwidth}{X} 法利賽人出去，立刻同希律一黨的人商議怎樣除掉耶穌。 \end{tabularx} \\ \\
[1ex]
\hline
\hline
\end{longtable}
$^{1}$好,各位大英姐妹,早安.
早安.
今天我要分享的題目就是.
新約中的情緒,一,憤怒.
為甚麼是新約中的情緒和一呢?.
其實有兩個原因.
第一個原因是跟我老婆有關的Bella.
大家都知道Bella在這裡.
做一個靈修指導.
令我心動.
她說一些這麼靈性的東西.
說一些這麼心靈性的東西.
我這個人又能不能在講道道.
說一些比較靈性的東西呢?.
但很可惜地,很不幸地.
我不是一個很靈性的人.
我真的不太熟悉靈性.
我講不到一些很靈性的東西.
就想想,隨便說.
不如就說一下新約中.
古代是怎樣看待情緒的呢?.
而古代是怎樣看待情緒.
又令我們對於自己的情緒有甚麼反思呢?.
所以我就選擇了這個題目.
至於為甚麼是一呢?.
就是和Mark Sir有關的.
因為我發現Mark Sir說到一個特別的地方.
就是他很喜歡製作一些系列出來.
例如甚麼基督眼中的教會1234.
然後甚麼世俗的君王.
現在說了一了.
所以我就,他啟發了我.
不如製作一個系列出來.
不要只說一個情緒.
多說幾種情緒.
所以如果有機會的話.
下次還有機會.
還有機會在這裡說的話.
可能會有二或者三.
但也可能只有一.

$^{41}$因為我想不到.
如果我突然間.
我這個人比較ad-hoc.
想不到的話.
這個系列就停在這裡.
就這樣了.
不過anyway.
總之今天我選了一個情緒.
我在想.
用甚麼情緒做這個系列的第一篇講道呢?.
我就在想.
不如說一些在教會比較少說的情緒.
因為這些太常見了.
沒甚麼好說的.
說一些我們經常都經歷的.
但在教會不是很常說的情緒.
就是甚麼呢?.
憤怒.
所以我就選了憤怒做我第一篇的講道.
其實很簡單.
憤怒是甚麼呢?.
就是我們在教會裡面.
我們平日經歷.
但在教會很少說.
也不會怎麼表達.
因為教會是甚麼呢?.
愛的群體.
要客客氣氣地.
你愛我,我愛你.
就算生氣.
也要私底下生氣.
不要生氣出面.
不然大家沒有朋友做.
是很嚴重的.
特別在華人教會.
大家要明白.
生氣過一次.
回不了頭.
這些事情是很嚴重的.
但我想去看看.

$^{81}$究竟在新約裡面.
是怎樣看待生氣呢?.
是否不要發出來這麼簡單呢?.
是否不要生氣這麼簡單呢?.
這個就是我很想和大家去看的東西.
在講憤怒之前.
其實我們很少去想.
我們需要問一個問題.
其實甚麼是憤怒呢?.
我們經常經歷的.
但我們很少去思考甚麼是憤怒.
所以我們首先去看看一些定義的問題.
究竟現代的人.
或者現代心理學.
怎樣定義憤怒?.
我選了一個最有權威性的定義.
來自美國心理學會的APA.
讀心理學的人都知道.
APA是最有權威的定義.
APA怎樣定義憤怒呢?.
它是這樣說的.
我自己翻譯的.
將英文轉為中文.
因挫折,真實或想像而來的傷害.
或感知到的不公正.
而產生的緊張和敵意.
這一種就是憤怒.
我又看看另一個定義.
這位叫做Paul Edmund.
是一個很有名的心理學家.
他是專門研究情緒的.
大家如果有看過迪士尼的一套卡通片.
叫做《Inside Out》的話.
就會知道.
其實他Paul Edmund.
是《Inside Out》系列動畫的心理學顧問.
他就是專門研究情緒.
他怎樣定義憤怒呢?.
他就說.
憤怒是七種基本情緒之一.

$^{121}$因為他定義了.
人無論是甚麼文化的人.
都有七種基本情緒.
只要你是人.
就有這七種基本情緒.
而憤怒是其中一種.
當我們追求目標.
受阻的時候.
或者受到不公平的待遇的時候.
就會產生憤怒.
而在最極端的情況下.
憤怒可能是最危險的情緒之一.
因為它和暴力是有一種關聯.
大家都很明白.
你很生氣的時候.
就很想打人.
所以可能是一個最危險的情緒之一.
我們看完一些現代的定義.
但因為我們今天要看新約.
我們不可以只說.
現代人怎樣看憤怒.
我們要問古代人怎樣看憤怒呢?.
我們就看看古代的定義.
一位人稱叫做亞里斯多德.
是很重要的.
說古代的情緒一定要看他的著作.
因為他在他的著作修辭學裡面.
定義了很多情緒.
他分析了各種不同的情緒.
他怎樣分析.
或者怎樣定義憤怒呢?.
他就這樣說.
憤怒就是因為認知到自己.
或者自己的家人.
或者自己的朋友.
受到不公平的待遇.
從而想報復的慾望.
這一種就是憤怒.
對於亞里斯多德來說.
憤怒是怎樣的呢?.

$^{161}$是認知性的.
是你自己知道自己憤怒.
還有你知道自己生氣.
我們平時經常說.
現代人可能會說.
憤怒是控制不了的.
是我裡面的衝動.
但亞里斯多德就不是這樣說.
亞里斯多德承認這種是一種慾望.
不是一種衝動.
但這種衝動是有認知性的.
當你認知到.
你覺得.
是你覺得.
你認知到自己.
或者自己所重視的人.
受到不公平的對待的時候.
當你認知到這件事的時候.
你有一種想報復的慾望的時候.
這一種就是憤怒.
這個就是亞里斯多德的定義.
他怎樣看我們怎樣處理憤怒呢?.
亞里斯多德很搞笑.
他就這樣說.
怎樣處理憤怒.
他就說.
任何人都懂得生氣.
這件事很簡單.
就好像花錢一樣.
我們當中的姐妹很明白.
花錢很簡單.
不過亞里斯多德這樣說.
但是.
但是有一個but.
很大的but.
但是.
要懂得對正確的人.
在正確的時間.
出於正確的目的.
以正確的方式去生氣和花錢.

$^{201}$不是很多人做得到.
也不容易做到.
所以正確地能夠做好.
這些事情是很罕見的.
是值得讚揚和高尚的.
亞里斯多德就將憤怒和花錢.
放在一起說.
大家是不是很明白.
馬上覺得這個哲學家很落地.
其實哲學家也很落地.
古典希臘哲學家很落地.
就是用花錢去比喻發怒.
我們每個人都懂得花錢.
我們每個人都懂得生氣.
但問題是.
不是你懂不懂得生氣.
而是你懂不懂得怎樣生氣.
亞里斯多德是強調怎樣生氣.
所以對於亞里斯多德來說.
憤怒是怎樣.
不一定是對的.
也不一定是錯的.
是要你正確地生氣.
就是對的.
錯誤地生氣就是錯的.
所以對亞里斯多德來說.
重點就在於你怎樣生氣.
我們又看看耶穌的時代.
到了希臘羅馬的時代.
他們的看法又和亞里斯多德不同了.
亞里斯多德是古典希臘的時代.
但到了耶穌的時代.
希臘羅馬的時代.
就有些東西不同了.
對於那時的人來說.
發怒是不道德的.
當時有很多不同的關於憤怒的理論.
但我們可以歸納到幾點的共通點.
第一點的共通點就是.
當時的人的文化就覺得.

$^{241}$發怒是不對的.
發怒是不道德的.
因為發怒是不道德的.
所以聖人.
聖人是甚麼.
道德高尚的人是不會發怒的.
這一點很重要.
跟我們現代不同.
在耶穌的時代.
希臘羅馬的時代.
道德高尚的人是要怎樣.
保持心平氣和.
不動心.
Aperfea.
我們希臘文叫.
就是你不會受到外來的刺激.
影響你的心情.
所以聖人任何時候都很平靜.
就像佛祖一樣.
沒有東西可以刺激到他.
不會生氣,不會哭.
不會不開心.
這就是聖人在那個時代.
所以對於那時的人來說.
生氣是不道德的.
跟你做得對不對是有關的.
你生氣就是做錯了.
是不道德的.
而且聖人是不會生氣的.
而且當時的聖人有很多種.
各種五花八門的.
自怒的理論.
即是你怎樣去控制.
或者你怎樣處理憤怒的理論.
有很多種的.
今天不夠時間說.
我只舉兩個例子.
第一種例子就是.
當時很流行的斯多亞主義.
Stoicism.

$^{281}$斯多亞的哲學認為.
人應該完全徹底地跟隨憤怒.
人是絕對不應該生氣的.
最理想的情況.
要徹底地跟隨他.
這就是斯多亞哲學的想法.
另一種哲學的想法就是.
Platonism.
我們叫它Middle Platonism.
中期的柏拉圖主義.
他們認為.
人的憤怒是無法避免的.
因為憤怒是控制不了的.
不是控制得到.
但憤怒是一定會有的.
但問題是你應該要控制他.
即是憤怒怎樣都是不好的.
但我們沒有辦法避免這樣不好的東西.
但我們要最大程度地控制他.
這就是柏拉圖主義.
Middle Platonism說的話.
他們兩個的分別是什麼呢.
他們兩個都認為憤怒是不好的.
這兩種哲學.
但分別就在於.
斯多亞哲學認為要完全徹底跟隨他.
而柏拉圖哲學認為.
我們要控制他.
明白分別了嗎.
我們又看看聖經怎樣說.
在聖經裡面.
基本上都反映到這兩種流派.
自怒的流派.
第一種就叫做跟隨憤怒的流派.
我找了四段經文.
這四段經文.
我們今天會比較看得多經文.
因為今天說新約中的憤怒.
這四段經文.
《馬太福音》.

$^{321}$《歌羅西書》.
《爾忽所書》.
和《提摩太前書》.
就反映了斯多亞哲學的那種精神.
就是跟隨憤怒.
例如.
《爾忽所書》就說.
一切的苦毒,怒氣,憤怒.
都要從你們當中.
跟隨,除掉.
要完全徹底地拿走.
例如《歌羅西書》就說.
現在你們要除掉這一切.
包括什麼?.
包括憤怒,怒氣等等.
要完全跟隨.
這一部分.
大家明白了.
就是跟隨.
在聖經裡面是有跟隨憤怒的一部分.
但同一時間我們也找到一些經文.
是反映另一種學派.
就是剛才說的.
Middle Patronism.
就是怎樣?.
控制憤怒.
例如《爾忽所書》.
又是《爾忽所書》.
《爾忽所書》.
《提多書》.
和《雅各書》.
這三段經文.
就說我們應該控制憤怒.
舉個例子.
《雅各書》.
《雅各書》就沒有說你要跟隨憤怒.
不過它就說.
我們每一個人都要敏銳地聽.
不要輕易地說話.
以及不要輕易憤怒.

$^{361}$不要輕易憤怒是什麼?.
就是你可以憤怒.
但不要那麼容易憤怒.
就是你要控制你的憤怒.
又或者.
《提多書》就是說.
做監督的人不要那麼容易發怒.
就是說發怒也是無可避免的.
不過你不要那麼容易發怒.
就是你要控制你的怒氣.
OK?.
所以就比較反映.
是Platonism的那種.
就是控制憤怒的精神.
好.
那麼.
引了那麼多經文.
簡單來說.
《聖經》的立場是什麼?.
控制憤怒.
不要生氣.
好.
我們今天的講道就講完了.
好.
跟你旁邊那個說.
你不要生氣.
快點說.
不說而已.
你跟你旁邊那個說.
不准生氣.
《聖經》說你不要生氣.
快點說.
對了.
不准生氣.
不要生氣.
控制憤怒.
OK?.
對了.
如果是這樣的話.
我今天的講道就講完了.

$^{401}$我可以下台了.
我這樣說當然不是.
我們就要問一個問題.
《聖經》是不是就這麼一面倒.
只是叫我們跟隨憤怒.
或者控制憤怒.
你不准生氣.
那當然不是.
我們要問一個問題就是.
究竟新約裡面.
有沒有一些正面的憤怒呢?.
我就選了一些經文.
大家都知道.
譬如在保羅裡面.
保羅其實在《迦太太書》.
或者《哥林多前書》是怎樣的呢?.
很憤怒地罵爆迦太和哥林多的信徒.
所以其實是有正面地用憤怒.
不過我今天想用一些更加明顯的例子.
是有憤怒這個字的例子.
就是哪裡呢?.
就是馬可福音.
馬可福音其實很明確地記載了.
耶穌有兩次很明確的憤怒.
我們就去看看這些經文.
首先就是馬可福音的第三章一至六節.
那裡就是說到耶穌走進會堂的時候.
就看到有一個人的手枯萎了.
但是那時候就有一些法利賽人.
文士就監視耶穌.
就想看看他會不會在安息日的時候.
就去醫好那個手枯萎的人.
就找把柄去控告他.
然後耶穌就跟那個手枯萎了的人說.
你站出來站在中間.
然後跟他們周圍監視他的人說.
在安息日做什麼是對的呢?.
是醫好人是對的.
就是做好事是對的.
醫好人是對的.

$^{441}$還是殺人是對的呢?.
那幫監視他的人就不出聲了.
然後聖經就這樣說.
耶穌就憤怒地還視他們.
什麼叫憤怒地還視他們呢?.
用廣東話說比較正經一點.
就是怒瞪他們.
你聽不懂吧?.
就是比較正經.
耶穌就怒瞪著那些監視他的人.
並且因為他們的心乾硬而覺得難過.
就對那個人說伸出手來.
他將手一伸就好了.
然後法利賽人和希律黨人就商議.
怎樣去殺害耶穌.
這段經文是很清晰.
在原文希臘文都是有憤怒這個字.
耶穌是很憤怒地怒瞪那幫監視他的人.
那我們就要問.
耶穌在這裡為什麼這麼憤怒呢?.
耶穌在這裡這麼憤怒是因為什麼呢?.
安息日本身是什麼呢?.
安息日本身的意義是要令人安息的.
就是令人要回到一個神原初的創造秩序裡面.
但是這一幫監視耶穌的人.
民事法利賽人.
在安息日竟然做什麼呢?.
竟然是要看耶穌會不會治好這個人.
對他們來說安息日不是讓人安息的.
安息日對於這幫民事法利賽人來說是不應該治病的.
是完全違反了耶穌對於安息日的理解.
所以這一刻耶穌很生氣.
你竟然用神的安息日這個誡命來難了我治好一個人.
耶穌因為這一幫民事法利賽人心裡肝癌.
所以就很生氣怒打他們.
你明白吧?.
耶穌在這裡生氣是因為這一幫人顛覆了上帝安息日的意義.
他們不明白安息日的意義.
以至安息日不再使人安息.
而是難了人得醫治.

$^{481}$這個就是耶穌很生氣的原因.
接著我們去看第二段經文.
在《馬可福音》的十一章十二到十三節.
這裡是說耶穌進入耶路撒冷之後.
在第二天他們離開伯大利的時候.
耶穌就肚餓了.
遠遠就看到一棵有葉子的無花果樹.
就走過去看看能不能在樹上找到什麼.
但去到的時候就看到那棵樹除了葉子之外.
什麼都找不到.
因為這個不是無花果的櫃子.
於是耶穌就對那棵樹說.
就是罵那棵樹.
從今以後沒有人吃你的果子.
他的門徒都聽到了.
這裡是說什麼呢?.
就是耶穌罵那棵樹.
我們平時看這段經文的時候就會覺得.
哇!很匪夷所思.
耶穌為什麼要罵那棵樹呢?.
那棵樹做錯了什麼呢?.
經文已經說了不是無花果的季節.
等於你現在很想吃蘋果.
但現在不是蘋果的季節.
然後你看到那棵蘋果樹沒有蘋果.
然後你就罵那棵蘋果樹.
在這裡耶穌就是在做一件事.
罵一棵根本就不在蘋果的季節.
根本就不會有蘋果的蘋果樹.
你罵他.
是不是很難以理解?.
所以我們就要去解明.
我們要問.
耶穌為什麼在這裡要這麼生氣.
去罵一棵沒有無花果的無花果樹呢?.
是不是那棵樹做錯了什麼.
值得被人罵呢?.
其實如果我們去想.
其實那棵樹也很無辜.
明明不是結果的季節.

$^{521}$被人罵.
還要是枯萎了.
多慘啊.
我們就要問其實耶穌在罵什麼呢?.
其實我們要看看.
剛才就著無花果樹的經文的上下文.
在上文11章11節.
就說耶穌進入耶路撒冷.
他就直接進入了聖殿.
然後做了一件什麼事呢?.
他還是四周的一切.
然後就和12個門徒離開.
往八大離去.
我剛才說過.
還是四周.
是路啤.
原文是一樣的.
和剛才我看第三章的動詞是一樣的.
就是耶穌進入耶路撒冷.
直接去到聖殿之後.
做了一件事.
就是路啤周圍的一切.
路啤.
但是之後沒有做事.
路啤完之後就走了.
OK 明白了.
然後下文是什麼呢?.
下文就是11章的15到17節.
那裡就是說.
耶穌就在無花果樹.
然後就回到耶路撒冷.
又是這樣回到耶路撒冷之後做什麼呢?.
就直接進入了聖殿.
然後耶穌做什麼呢?.
我們傳統就會說這段經文是潔淨聖殿.
但其實這裡不是潔淨.
耶穌是審判聖殿.
耶穌是將聖殿裡面所有做買賣的人都趕走.
然後推翻兌換銀錢的人的桌子.
還有那些賣鴿子的人的桌椅.

$^{561}$OK.
在這裡耶穌做了一件事.
就是他做了一個symbolic action.
去審判聖殿.
將聖殿裡面所有做買賣的人.
還有兌換銀錢的人都趕走.
OK.
在這裡.
剛才那段就在無花果樹的經文.
我們看到了.
上下文都是在說什麼呢?.
上下文都是和聖殿有關.
都是說耶穌進入聖殿.
我們開始會明白到.
其實耶穌在生氣什麼呢?.
其實耶穌罵那棵無花果樹沒有結果子.
不是那棵無花果樹做錯了.
而是無花果樹是一個symbol.
這個無花果樹是一個什麼symbol?.
是代表著聖殿.
耶穌在這裡真正要罵的對象.
不是那棵沒有果子的無花果樹.
耶穌真正發怒的對象.
是聖殿.
所以《馬學福音》的十一章十一節.
已經很清楚.
讓我們看到.
是耶穌進入聖殿的時候.
是怒啞周圍的一切.
是聖殿令到他生氣.
其實不是那棵無花果樹.
無花果樹只不過是一個象徵.
是聖殿令到耶穌生氣.
所以耶穌去到十一章的十五到十七節的時候.
他就真的對聖殿做了一些行動.
就是趕走聖殿所有做買賣的人.
和所有兌換薪金的人.
我們就開始明白到.
原來耶穌罵那棵無花果樹.
是在罵聖殿.

$^{601}$如果我們再看《馬學福音》的十一章十二到二十五節.
整個線型結構就會更加清晰.
就是大家笑.
因為上了中神生命之道.
就明白什麼叫線型結構.
不多解.
你會看到第一部分就坐無花果樹.
而中間夾著審判聖殿.
就是我剛才說的那部分.
然後就去到十一章的二十到二十五節.
當他們再經過百達尼.
再經過無花果樹的地方的時候.
就看到無花果樹枯萎了.
所以就看到無花果樹.
就坐審判聖殿和無花果樹.
所以就是ABA派的結構.
線型結構.
我們就很清楚看到.
透過這個結構我們就明白.
原來耶穌罵那棵無花果樹.
其實實質上是在罵聖殿.
和審判聖殿.
而無花果樹的枯萎是象徵什麼.
象徵這座聖殿不能夠再存在.
耶穌不是潔淨聖殿.
潔淨是說你潔淨完之後.
他就恢復.
但耶穌不是要恢復聖殿.
耶穌是要審判聖殿.
他要毀滅聖殿.
他要拆毀聖殿.
聖殿不應該繼續存在.
就好像那棵無花果樹要枯萎一樣.
我們要問一個問題.
聖殿有什麼問題.
其實夏文友說的.
馬拉福音的十二章的41到44節.
我們平時看這段經文.
就會覺得耶穌在讚賞寡婦.
這個是窮寡婦的奉獻.

$^{641}$就是耶穌見到一個寡婦.
窮寡婦將自己所擁有的一切.
就是他賴以為生的一切.
都奉獻給聖殿.
我們平時看聖經就會覺得.
耶穌在讚揚寡婦.
但其實不是.
耶穌是在批判聖殿.
聖殿是什麼呢.
其實聖殿就好像一個中央銀行.
是一個當時猶太人的財政中心.
收集很多錢.
在一個這樣的地方.
聖殿存在的地方.
聖殿本身應該是怎樣的.
聖殿除了是中央銀行.
應該是一個救濟的中心.
如果聖殿運作得好.
應該是沒有窮人才對.
但這裡竟然會怎樣.
有一個寡婦.
這個寡婦要怎樣.
還要把自己為生的一切.
給聖殿.
其實這裡說的是什麼.
就是聖殿帶來的是什麼.
聖殿帶來的是經濟的不公義.
聖殿是在剝削當時的窮人.
而不是在救濟當時的窮人.
所以我們再去看這段.
窮寡婦的奉獻的時候.
就多一個角度.
不是只讚揚窮寡婦.
為主擺上一切.
不是為主.
是聖殿剝削了她的一切.
為生的一切都給了聖殿.
這個就是耶穌恨聖殿的原因.
這座聖殿不再是神所用的聖殿.
就正如那班文士和法理賽人.

$^{681}$扭曲了安息日一樣.
安息日不再使人安息.
而聖殿也不再去救濟當中的窮人.
去失卻了財政中心的作用.
所以耶穌就要審判這個聖殿.
說完兩段關於耶穌的憤怒.
我們有什麼簡單的結論呢.
有些小結.
第一就是剛才我們說過.
當時耶穌的時代是什麼呢.
聖人是不會生氣的.
但是我們看到馬學福音記載.
耶穌是怎樣的.
是生氣的.
所以耶穌是一個非典型的聖人.
OK.
耶穌是一個會生氣的聖人.
但也正正透過耶穌的生氣.
讓我們看到一件事.
馬學福音的記載.
和當時的文化很不同.
剛才記得我說過.
當時的文化覺得.
生氣是什麼呢.
不道德.
但耶穌顛覆了這種想法.
生氣可以是什麼呢.
生氣可以是一種道德的行為.
生氣是可以道德的.
生氣是可以合乎道德.
而且是值得追求的道德行為.
所以第一點很想說的就是.
透過耶穌的生氣我們看到.
原來生氣也可以是一種道德行為.
但是生氣怎樣是道德行為呢.
是當這個憤怒是針對什麼呢.
當這個憤怒是針對不公義的人.
或者不公義的體制的時候.
這種憤怒就會是一種合乎道德.
和值得追求的道德行為.

$^{721}$可能有人會反問.
反駁我.
是耶穌生氣的.
耶穌是神來的.
祂做什麼都可以.
祂生氣.
神生氣不代表人可以生氣.
但我想說不是的.
其實馬學福音是一種.
希臘羅馬的傳記體裁.
傳記體裁代表什麼呢.
傳記體裁.
當時的人寫傳記.
是為了要讓讀者.
去仿效當中的英雄人物.
或者當中的主角.
所做的正面的行為.
所以馬學福音這裡.
特意記載耶穌.
為一些不義的事情發生生氣.
其實是要讓讀者.
去仿效耶穌的行為.
我們是需要仿效耶穌.
怎樣去生氣.
怎樣正確地發生生氣.
怎樣去一些不公義的人.
或者不公義的體制去發生生氣.
今天我想帶出的第一個反省.
是什麼呢.
我們很多時候.
教會迴避發生生氣這個問題.
沒錯.
發生生氣是會帶來暴力的.
發生生氣發得不好.
是很破壞關係的.
但教會同時也有一個問題.
就是我們不再發生生氣.
其實這個是很嚴重的.
如果我們見到.
社會有一些不公義的事情.

$^{761}$我們見到一些很不公義的人.
我們不生氣的話.
其實我們是不合乎道德的.
我們的道德是敗壞了.
記得我剛才說什麼嗎.
李士多德去定義.
發生生氣是什麼.
是見到一些不公義不公正的事情.
你就會生氣.
如果我們見到一些.
不公義不公正的事情.
我們不生氣的時候.
代表什麼呢.
就是我們不覺得.
那些不公義不公正的事情.
是不公義.
你明白我說什麼吧.
我們見到一些不公義的事情.
我們竟然不生氣的時候.
就是代表我們認同那些不公義的事情.
我們不覺得他們不公義.
所以我們就不生氣.
今天值得我們問.
我們基督信仰會不會變成這樣.
我們福音派的信仰.
會不會已經忘記了怎樣生氣.
還有原來我們是應該生氣.
我們在一些應該生氣的事情上不生氣.
我們變成好像法利賽人和文子一樣.
心裡乾顏.
是會被耶穌罵.
是會被耶穌咒詛.
好像一棵無花果樹一樣枯萎.
我們不希望我們的教會是這樣.
我們究竟能不能夠仿效耶穌.
學會怎樣正確地生氣呢.
好,接著還有一點時間.
看第二段經文.
接著看啟示錄.
關於憤怒的事情.

$^{801}$就是在啟示錄第六章9到11節.
那裡就說當約翰揭開.
不是當約翰.
當揭開第五個印的時候.
約翰就見到.
在祭壇底下有一些為神的道.
並且作見證被殺的靈魂.
大聲喊著說.
聖潔真實的主啊.
你不審判住在地上的人.
為我們吸我們新流血的冤要到何時呢.
這裡有幾個重點.
第一個重點就是在祭壇底下.
見到一些殉道的人.
約翰見到在祭壇底下.
有些殉道的人.
他們就在那裡喊.
神,你為什麼不審判地上的人.
要直到何時.
這裡其實是有一個憤怒在這裡.
就是聖潔真實的主啊.
你不審判住在地上的人.
為我們新冤要到何時呢.
這一種的方法.
其實是舊約斯篇哀告詩的方法.
就是當詩人很憤怒上帝的時候.
上帝你是公義的神.
你見到一些不公義的東西.
你不行動要到何時.
其實是在說對神的憤怒.
明白了嗎.
這裡是在說那些殉道的靈魂.
對神的憤怒.
你見到這些不公義的東西.
但神你不行動.
要到何時呢.
所以這裡是在說.
這群殉道的人的靈魂.
對上帝的憤怒.
而祭壇底下是在說什麼呢.

$^{841}$這些殉道者其實就好像一個祭物.
當時的早期基督教覺得.
殉道就好像一個祭物.
獻給上帝.
殉道者的血其實是一個獻祭.
所以這裡就在祭壇底下形容.
所以那些殉道者就像祭物一樣獻給神.
但這些好像祭物一樣獻給神的殉道者.
就很憤怒.
神,為什麼你不出手.
要到何時呢.
我們還要問一個問題.
就是你不審判住在地上的人.
我們就要問.
住在地上的人是什麼人呢.
在啟示錄的十三章六至八節.
就跟我們說了住在地上的人是什麼.
那裡就介紹了.
有一隻獸就出口說一些褻瀆神的話.
然後神就任憑這隻獸.
跟聖徒爭戰並且得勝.
然後凡住在地上的.
明治從創世以來.
沒有記在被殺的高揚的生命冊上的人.
都拜這隻獸.
簡單來說住在地上的人就是什麼.
住在地上的人就是那些敬拜這隻獸的人.
而這些敬拜這隻獸的人是什麼.
就是那些去逼迫聖徒.
跟聖徒爭戰並且得勝的人.
簡單來說是什麼.
住在地上的人就是敬拜獸的人.
而這些敬拜獸的人就是殺了那些.
在祭壇底下.
在祭壇底下氣憤上帝.
在呼喊靈魂的人.
明白了嗎?這個邏輯很簡單.
我們就要再問一個問題.
剛才在祭壇底下的殉道者就在叫.
神你不出手到什麼時候.

$^{881}$對上帝的憤怒.
我們就要問.
究竟神有沒有出手呢?.
神有沒有去回應.
這些在祭壇底下的靈魂的禱告呢?.
我們當然就知道.
如果熟讀啟示錄.
就知道上帝是有回應的.
就在第十四章九到十一節就說.
又有第三位天子時接著他們.
大聲說若有人拜獸或者受像的.
或者在額頭或者手上受了印記的.
就必喝神大怒的酒.
這個酒在神憤怒的杯中.
唇一不擇.
他要在聖天使和羔羊的面前.
在火和硫磺之中受痛苦.
他是那些拜獸和受像的人.
就是那些住在地上的人.
這些住在地上的人要被上帝審判.
要怎樣呢?.
在火和硫磺當中受痛苦.
他受痛苦的煙要往上升.
直到永永遠遠.
這裡值得留意的是什麼呢?.
有火,硫磺和受痛苦的煙要向上升.
我們留意一下.
這裡是說什麼呢?.
那些住在地上的人要被火燒.
很恐怖.
上帝要審判那些住在地上的人.
他們要被火和硫磺燒.
然後他們的煙要向上升.
如果我們留意到.
向上升的這個動詞.
其實就是對應啟示錄第八章三到四節.
那裡說有另一位天使拿著金香爐.
站在祭壇旁邊.
有很多香賜給他.
要和眾聖徒一起祈禱.

$^{921}$要和眾聖徒的禱告.
和一起獻在寶座前的金祭壇上.
那香的煙和眾聖徒的祈禱.
就好像從天使手中一同升到神面前.
簡單來說就是啟示錄第八章說到兩件事.
就是煙的香和聖徒的祈禱.
一同升到神的面前.
所以我們看到.
「升」字一同上升的動詞.
兩段經文是用同一個動詞.
簡單來說是什麼呢?.
在啟示錄第八章.
香的煙和聖徒的祈禱.
就是一種獻祭.
就是香氣的獻祭.
祈禱是一種獻祭.
升到神面前.
如果我們說動詞一樣的話.
意思是說什麼呢?.
是很變態的.
神在這裡覺得.
那些被火燒的.
那些迫害信徒的人.
他們被火燒是一種獻祭.
祭物要被火燒的.
他們被火燒是一種獻祭.
而他們的煙向上升.
就好像香爐的煙向上升一樣.
上升到神面前.
是被熱吶的.
是不是很恐怖呢?.
變態的上帝.
就是審判這些迫害人的人.
他們的審判就好像獻祭一樣.
為上帝所熱吶.
他們的煙就好像香氣一樣.
上升到神面前.
我們看到這裡是關於神的憤怒.
神最終都發怒了.
就是要審判這些.

$^{961}$邪惡的迫害人的人.
我們看到剛才我有說過.
這裡其實我們看了三段經文.
有一個共通點.
就是用一個獻祭的意象.
去連在一起.
就是殉道的信徒.
剛才我說了.
他們是一種獻祭.
殉道者是一種祭物獻給神.
然後聖徒的祈禱是一種獻祭.
最後那些迫不勝途的人.
他們被神審判.
都是一種獻祭.
好像煙一樣向上升.
我們看到對應了.
就是那些殉道者的血是一種祭.
而最終上帝會說.
我審判這些迫害你們的人.
他們都是一種祭物.
獻了給我.
所以這裡就是說一個.
好像很變態的上帝.
最後我們去歸納一下.
到底啟示錄說的憤怒是怎樣呢.
我們看到有三種憤怒.
第一就是.
聖徒的憤怒.
他們是對上帝的憤怒.
第二就是上帝的憤怒.
上帝是對惡人的憤怒.
而最後就是作者自己的憤怒.
其實啟示錄的作者很憤怒.
為什麼他要描寫.
一個上帝可以這麼變態.
我們是不明白的.
為什麼要寫上帝這麼變態.
這麼暴力.
會以折磨人.
做一種祭物.

$^{1001}$這樣去熱納呢.
其實背後是反映著.
啟示錄的作者的受傷.
這裡是反映著受傷者的憤怒.
當一個人受到傷害.
我們知道啟示錄的背景.
就是面對著迫害.
面對著迫害.
面對著很多艱難的時候.
一個人受傷被迫迫的時候.
被迫害被壓迫的人.
是有很多憤怒的.
這個就是作者的憤怒.
最後我們總結一下.
我們今天在說的話.
今天我們看了很多.
不同的新約經文.
雖然有不少的新約經文.
教導我們要控制憤怒.
但我們需要學習的有三點.
第一就是.
好像耶穌一樣.
原來憤怒也可以有正確的憤怒.
我們需要學習.
如何對於不公義的人.
或者體制去發怒.
第二就是.
就好像那些殉道者一樣.
他們是對神憤怒.
我們很多時候不容許自己.
或者我們在教會很少說.
我們對神憤怒.
但在這裡啟示錄告訴我們.
神是能夠熱納我們的憤怒.
我們面對很不公義的事情.
我們可能會問.
神為什麼會這樣發生呢.
或者我們見到一些不公義的事情.
發生的時候.
我們會問神為什麼不出手呢.

$^{1041}$就好像那些殉道者一樣.
啟示錄告訴我們.
其實我們是可以對神憤怒.
神是能夠容許我們對祂憤怒.
這個是我們信仰所容許的事.
而最後我們透過啟示錄的憤怒.
我們明白.
也在學習.
我們學習見到受傷者的憤怒.
當一個人受到極度的壓迫的時候.
因為他所遭受的不公義.
不公正的事情.
他可以有很多憤怒.
當人憤怒的時候.
我們就會覺得他錯.
但其實我們看不看到.
這些憤怒是代表那個人背後的傷害呢.
啟示錄的經文是教我們.
去學習怎樣去看和怎樣去感受.
受壓迫的人和受傷的人的憤怒.
我們能不能夠感受到.
這些憤怒背後的傷痛.
能不能夠和這些憤怒的人.
憤怒的受傷者去同行呢.
這個就是我今天很想.
看了很多經文.
想帶出來關於憤怒的一些反思.
我們一起去禱告.
主今天.
我們去看.
去反思.
去思考一個我們在教會很少想的情緒.
我們去思考一個憤怒的議題.
當我們去看到耶穌在萬惡福音當中的憤怒.
和啟示錄當中所呈現的憤怒的時候.
主今天讓我們對於憤怒有一個新的看法.
就是當我們見到一些不公義的事情的時候.
主因為你教我們懂得去生氣.
我們可以用你的眼光去看到.
有不公義在當中.

$^{1081}$我們懂得怎樣去正確地生氣.
主今天你復還這個憤怒的能力.
正確地憤怒的能力.
在我們的生命當中.
主也同樣地.
如果在在的聽眾姐妹.
我們生命當中對於上帝有一些憤怒.
對於我們身邊的人有一些憤怒的時候.
我們這一刻都可以對上帝說.
如果我們對上帝有憤怒的時候.
我們都可以對上帝說.
正如那些殉道的靈魂一樣.
對上帝去說.
因為我們的信仰是能夠給我們對上帝憤怒.
這是上帝能夠接納的東西.
所以主沒錯.
我們讓我們這班會眾回到你面前.
我們能夠將我們的憤怒.
坦誠地呈現在你面前.
主我們今天很想去禱告.
求你去轉化我們的生命.
讓我們懂得在你面前生氣.
我們不再藏私.
我們能夠去到你面前.
將我們的憤怒呈現給你看.
因為我們的上帝.
我們知道你是一個能夠和受傷者同行的上帝.
你能夠接納受傷者憤怒的上帝.
主今天願意你來到我們生命當中.
你去和我們說話.
你親自去承載我們生命當中.
這些傷痛的憤怒.
你去轉化我們的憤怒.
也幫助我們.
讓我們能夠成為一個能夠與那些憤怒的受傷者同行的人.
讓我們的教會成為一個能夠對不公義的事情.
去發聲去憤怒的教會.
也可以和那些受傷的憤怒者同行的教會.
天父多謝你這樣去聽我們的禱告.
是奉主耶穌基督的聖名而求.

$^{1121}$Amen.
多謝阿邦給我們的提醒.
我們想最後提醒大家.
今天是最後一個主日.
買感恩節聚餐的票.
所以等下完結的時候請你快快買.
等到我們不需要太憤怒.
我們好好準備感恩節的聚會.
願主祝福各位.
賜大家平安.
\newpage



\section{}
\label{sec:dQy1lcfV2Eg}
\textbf{Who is Your Lord? \#1 Cyrus (Isaiah 45) - Rev. Mark Chan}
\newline
\newline
連結: \href{https://youtube.com/watch?v=dQy1lcfV2Eg}{\texttt{https://youtube.com/watch?v=dQy1lcfV2Eg}} ~~~~ 語音日期: 2024-11-26
\newline
\newline
\hyperref[sec:KOpSYq8TJvA]{\small{< < < PREV SERMON < < <}}
~
\hyperref[sec:index_chronic]{\small{[返順時目]}}
~
\hyperref[sec:index_scriptual]{\small{[返順卷目]}}
~
\hyperref[sec:fC6e7drn7H4]{\small{> > > NEXT SERMON > > >}}
\newline
\newline
$^{1}$It is great to be here, to be with you guys, and I want to thank all of you last week for.
covering the operation here and we can continue to worship God..
And thank God for a safe journey to serve the churches in Toronto..
And I want to thank the team of people that came with me and they took care of a lot of.
things for me, setting up the booth for the Bible conference and that includes Andrea.
and Sam and Brenda and Amanda..
So they came with me and took care of a lot of stuff..
Originally I thought I could take care of it by myself but actually they helped me a.
lot..
So it was a blessing to be able to serve other churches and other Christians in other areas..
But again we come back and we have a very interesting development in the presidential.
election and everyone has their attention on this..
And now that it doesn't settle, we have the President-elect Donald Trump..
He's going to be in the White House in the next term..
And it was a big question when I was in Canada..
Everyone is asking me, "Hey, who's going to be elected?".
And it's like I can foretell the future..
But there's a lot of discussions about who's going to be in the White House..
But in any case, every four years we have this change of who is in power and which political.
party and political figures are going to be in the White House..
That's why we have this mini-series I want to inject into our thinking and our reflection.
on the Bible..
This series is called "Who is Your Lord?".
And the Pagan Kings in the Bible..
And I would like to take some of the samples from the Bible, especially from the Old Testament..
As you can see from the biblical history, from the pages of the Bible, you have a lot.
of different kings and kingdoms and empires rise and fall in the entire biblical history..
All the way from the left you get the Roman Empire, and then Greek, and then Hittite,.
and Assyrian Empire, Babylonian Empire, Persian Empire, and down in the south is the Egyptian.
Empire..
And then in the middle of it, that is the center of the ancient world, the land bridge,.
the land of Israel and Palestine, you got different small kingdoms..
The Arameans and the Philistines, and of course Israel and Judah right in the middle, in the.
shadow of all the empires..
As you can tell from this map, and you can get the sense from the Bible, that kingdoms.
rise and fall, empires rise and fall, and kings do rise and fall..
Because God is the Lord of history..
So we want to do this..
We want to go through some of the kings, secular kings, pagan kings in the Bible and see what.
the Bible has to say..

$^{41}$So this is our procedure..
Our procedure is first introduce the pagan king found in the Bible, recorded in the Bible,.
or even named in the Bible, and get their secular profile, you know, like documents.
outside of the Bible, what do they say about this and that pagan king, and then we go back.
to the Bible and see what we can reflect upon..
So this is our procedure..
We first get the secular profile of the pagan king in this series, and then we go into the.
Bible..
And the first one I want to choose is, you see the picture there?.
Cyrus the Great..
So ring any bell?.
You know, this king is named in the Bible, in the book of Isaiah, even before he existed..
So Cyrus the Great..
This timeline, let me go through this timeline for you..
You see the purple center spectrum?.
All the way from the left to right is 700 BC, before Christ, 700 years before Christ,.
all the way to 500 BC..
So this period of history, there were a lot happening, as you can tell..
Let's go through one by one..
First of all, it's the red one on top..
You see that?.
What is that kingdom called?.
What is that empire called?.
Assyrian, right?.
From 911 BC to 605 BC, and then it declined..
And what happened here during the Assyrian Empire was that the fall of the northern Israel.
kingdom Samaria, the fall of Samaria in 722..
So Assyrians, they attacked the land of Israel and they defeated the northern kingdom Israel.
and the capital Samaria fell in 722 BCE..
And all the way down to go south to the southern kingdom, Judah, it was year 701 BC..
And that was when Hezekiah was reigning in Judah..
And the Assyrians blocked all the exits and entrances to the capital Jerusalem, and they.
lay siege of the capital, but they failed and they returned..
So that was 701 BC..
So as I said, Assyrian Empire rise and fall..
At the end of 7th century, about 600 something, it begins to weaken..
And then there were two other superpowers in the Middle East that rise up..
The first one is the blue one, the Median Kingdom, the Median Empire..
And then south of the Median Empire is the Babylonian Empire..
And what happened here is that Babylonians continued to attack and continue to expand.

$^{81}$and they came all the way to the southern small kingdom of Judah and they defeated Judah.
and fall of Jerusalem happened in 586 BC..
That's when Assyrians fell and then Babylonians took over the entire world, almost the entire.
world of the Middle East..
And then the Median Kingdom, there was something interesting happened..
It fell to the Persian Empire..
Now what happened to the Persian Empire when it rises was that it's because of Cyrus the.
Great..
Cyrus the Great was a military leader of the Persian people and they rebelled against the.
Medes because they have been serving the Median Empire, the overlord for so many years and.
they are fed up..
So they rebelled against it and Cyrus the Great let his people defeated the Medes..
Well, first of all, they rebelled against the Medes in 553 and then within three years.
they actually defeat the entire empire of Median Kingdom and they swallow up the entire.
Middle East..
So 539, that is a very important year in the biblical history..
539 BC is Cyrus enters the capital of Babylonian Empire, the city of Babylon, the greatest.
city at that time..
So he entered there peacefully and took over the entire world..
And 539 is also important for the Jews because as you can tell from 586, the Jews were exiled.
to Babylon for about 70 years and then 539 when Persian Empire took over the world, they.
reversed the policy of the previous kingdom..
They don't keep the people, they don't exile the people, they send the people back and.
send their temple vessels back to their homeland and Jews were one of them that got the benefit..
So that's when we have the book of Ezra and Nehemiah when they return from Babylon to.
Jerusalem and we build a temple and we build the city..
So that is what happened for Cyrus the Great..
Now I want your attention to the last 10 years of the Babylonian Empire which is 555 to 539..
That is the last 10 or so years at the end of Babylonian Empire..
What happened here was that there were two kings, father and son, Nebunidas and Belshazzar..
Father and son taking care of the last period of the Babylonian Empire..
And what happened was that Nebunidas was a very funny king and he didn't worship the.
main god, the biggest god in Babylon, the biggest god being Marduk..
And he forsake Marduk and he worship the other god of his own family which is the moon god.
called Sin..
So that made a lot of people in Babylon, especially the capital, not very happy..
There were a lot of social tension and a lot of discontent in the society of Babylon..
Like us, some people are being fed up by what is going on in the society and what is going.
on in our nation..
So they wanted to turn around, they wanted to have someone else in power..

$^{121}$So that is what happened in Babylon in the last 10 years..
Now what happened was that sin is being elevated by Nebunidas..
As you can tell from this stele, from this relief, that is the king Nebunidas on the.
right and he is having a staff in his left hand and he is raising his right hand, worshipping.
the three gods..
You see the moon god, you see the moon right?.
And then you see the sun and then also the third symbol is Venus..
So he is worshipping the moon and the sun and the stars instead of worshipping the guy.
on the left which is Marduk or Merodoth mentioned in the book of Jeremiah..
Now there was tons of huge discontent among the Babylonians..
If you take the poll in the last 10 years of the Babylonian empire, the last 10 years,.
you will get the sense it is very low opinion for Nebunidas and his son Belshazzar who is.
also mentioned in the book of Daniel..
And there was a political vacuum in the last 10 years of the Babylonian empire..
And Cyrus took this opportunity to invade Babylon..
Now you need to understand the ancient context when the king is privileging one god over.
the other, privileging sin over Marduk is a huge thing in the ancient world..
In Babylon there is a yearly annual festival called Akitv..
Every spring around March to April they celebrate because that is when they harvest the wheat.
and barley and it is a huge festival and the king has to be there to host the festival.
and they will take Marduk the idol out of Marduk's temple and put it on a cart and.
then they will do a procession in the capital and it was a huge celebration and it would.
determine the destinies of everything for the next year..
So Nebunidas did not do that and that made a lot of people upset..
So Cyrus took this opportunity and he would battle against his overlord the Medes and.
then a significant number of military people from the Median empire actually switched sides.
to help Cyrus to help him to defeat the Median kingdom and also the Babylonian kingdom and.
finally in 539 BC they defeated the biggest superpower in the world Babylonia..
As you can see all these arrows 553 he would battle against the Medes and then 550 within.
3 years he defeated the entire northern Median kingdom and then all the way he invaded to.
the left to Lydia defeated the Lydia kingdom and then come back and defeated Elam and 539.
as I told you is a very important year in biblical history and world history is that.
Cyrus entered the capital of Babylonian empire Babylon..
And we can find some documents outside of the bible that talks about Cyrus the great.
and one of the documents is called Cyrus the Cylinder and this round thing has a lot of.
inscriptions on it and if you translate it or take an excerpt out of it it talks about.
Cyrus what he did to Babylon..
He said I am Cyrus king of the world king of four corners I enter Babylon in peace and.
Marduk I daily attend to his worship because Marduk exalted me and have all the kings kiss.
my feet and I return the images of the gods and I gather all their peoples and return.

$^{161}$them to their dwellings..
Remember the Jewish people were one of the groups being returned..
Now from this excerpt who is Cyrus?.
Is he a worshipper of Yahweh the true God?.
No he worshiped the false gods..
He worshiped Marduk in fact as a Persian he not only worshipped Marduk the Babylonian.
god he worshipped his Persian god the great god of the Persian religion is called Ahura.
Masta that was a Persian great god if you drive a master car you know what that means.
master in Persian means wisdom..
So he is actually a polytheistic person and he worshipped different gods and as you can.
tell from the context Cyrus was a pragmatic politician he took the advantage of Babylonian.
vacuum political vacuum and he invaded Babylonian empire..
Now he let the captives free he let some of the Jewish people go back to Jerusalem not.
because he is kind or not because he is compassionate..
Is that he is pragmatic he wanted to reverse the policy from Babylonian empire and let.
these people go as long as they don't rebel against him and he will consolidate his power..
So as you can tell he is a polytheistic worshipper he worshipped false gods and he is also a.
pragmatic politician much like today all politicians are pragmatic and you have to be pragmatic..
Now with that context we can go to the Bible..
What does the Bible say about Cyrus the great?.
We find his name being mentioned in Isaiah 44 and 45..
Let me read this passage to you..
And remember Isaiah has two parts the first part is chapter 1 to 39 that is what we call.
first Isaiah and then chapter 40 all the way to the end sometimes we call it second Isaiah.
the second part of Isaiah..
The second part of Isaiah from 40 to all the way to 66 is a prophecy bunch of prophecies.
talking to the Jewish captives in Babylon and how God will have compassion on them and.
to lead them to return to give them hope..
That is what Isaiah 40 all the way to 66 is giving them hope and telling them Yahweh.
the true God will be faithful to them even though they have sinned against him and he.
will forgive their sin and he will lead his people back to their homeland..
So that is the context of 44 chapter 44..
So God is telling his people Israel he said remember this things O Jacob for you are my.
servant O Israel I have made you you are my servant O Israel I will not forget you return.
to me for I have redeemed you..
And God is the one who sets off Cyrus Cyrus the great the politician he is my shepherd.
and will accomplish all that I please..
He will say of Jerusalem let it be rebuilt and to the temple let its foundation be laid.
and this passage continues to tell us something more about Cyrus in the sight of God..
God said what this is what the Lord Yahweh says to his anointed to Cyrus whose right.

$^{201}$hand I take hold of to subdue nations before him and to strip kings of their armor and.
open doors before him the doors of the empire doors of other people and nations including.
Babylon so that gates will not be shut will be open to Cyrus the great..
Now what is amazing what is interesting here is that not only Cyrus name is being called.
out before he exists remember Isaiah prophesied this 100 years before Cyrus existed not only.
his name is being mentioned here but also he is being called his anointed his anointed..
That's interesting because a pagan king is being called the anointed..
That is pretty remarkable..
Now you go back to the timeline you see Isaiah there 740-700 that's when Isaiah was a prophet.
and he was serving the sovereign kingdom Judah in the midst of the crisis from Assyrian empire..
So before Cyrus existed in Isaiah his name was mentioned and not only his name was mentioned.
he was being called God's anointed..
Now immediately we thought ok anointed the Messiah so often we have this Jesus figure.
in our head right whenever Messiah is being mentioned in our songs or in the Bible so.
often we say ok he is the Messiah the capital M Messiah the savior of the world Jesus Christ.
himself there is only one Messiah..
So we always have this sacred notion to the word Messiah but again you need to understand.
these kind of sacred notion got picked up later in biblical history in the new testament.
but in the old testament whenever the word anointed or Messiah is being used it does.
not necessarily mean that someone is sacred or someone is pleasing to God..
In fact so often someone who is being anointed is not pleasing to God..
For example in 1 Samuel 10 verse 1 it speaks of the first king of Israel King Saul remember.
that guy King Saul and after King Saul is David right..
King Saul is being said to be anointed the prophet Samuel took a flask of oil and poured.
it on the head of Saul and kissed him and said has not the Lord Yahweh anointed you.
leader over his inheritance his people..
So Saul is being mentioned or being referred to as the Messiah same word anointed..
Now as you can remember was King Saul a king pleasing to God?.
No right he was actually rebellious to the kingship of God and that's why God raised.
up another king David to replace him..
However Saul is being called the anointed..
So as I said anointed does not mean that he is sacred and this is the same way for the.
pagan kings in the bible especially Cyrus the great..
Even Cyrus the great was being named and is being called the anointed of Yahweh does not.
mean he is a true believer of Yahweh..
In fact he is pagan as we can tell from other documents outside of the bible..
You know in the bible the pagan king is used by God for God's glory..
The pagan king is never put on a pedestal..
He never gets the glory..
God gets the glory..

$^{241}$Whoever is in the white house, whoever is in power, whoever is serving the term in our.
world in the ancient world..
God gets the glory because God is in charge..
Isaiah 45 actually says this, it says the whole point of raising up Cyrus as the anointed.
one is to gain glory for himself..
God says this..
God gets the glory because he controls every move of the political figures and the kings.
and kingdoms and empires on earth..
Isaiah 45 verse 5 it says I am the Lord Yahweh and there is no other apart from me there.
is no God..
I form the light and create darkness..
I bring prosperity and create disaster..
I the Lord Yahweh do all these things..
You can see from the translation here I give you in parenthesis the Hebrew words..
Yasa being translated form, Bara often times being translated create and then Asab translated.
bringing or doing..
Now these kind of words, these words are often used in Genesis 1..
Now immediately we have some theological questions here..
We ask the question of the problem of evil..
How can God create darkness and disaster?.
Is he the origin or the author of evil?.
Is God the author of evil?.
No, no of course not..
But we are troubled with these words because we as modern readers we always try to understand.
these words from a materialistic point of view..
Like God in Genesis 1 created everything we just thought ok it has to be God creating.
something out of nothing..
It's a material process, it's a material creation manufacturing process that we have in mind..
But in the ancient world that is not what they meant..
In the ancient world in Genesis 1 and in Isaiah 45 when they used these words they are not.
talking about material or creating something out of nothing..
It's not talking about God creating disasters out of nothing..
Calamity out of nothing..
There is an evil out of nothing..
It's not talking about that..
These kind of words are talking about management..
Talking about God is in control of everything..
Good things and even bad things is still under the control in the hands of God..
So God has good and evil in his hands..
Never have anything that happened in our lives or in our world that would be outside of God's.

$^{281}$control..
That's what they meant..
That's what these Hebrew words are trying to mean..
So God would allow evil to happen..
Or even a pagan king like Cyrus the Great to control the world..
That does not mean they are out of God's hands..
God has everything, the good and the evil in his control..
Just like today when your boss tells you, can you create a department or a team to take.
care of this task?.
You heard the word create..
You would not think, ok, how can I create a team of people out of nothing?.
Because he is not talking about manufacturing material process..
He is talking about management..
You create a team, it means you have a team in charge to perform a task being assigned.
to you..
It is the same way in the biblical words..
Here it talks about reorganization..
It talks about management..
Do you play chess?.
Anyone here?.
Anyone who is good at playing chess?.
I am very bad..
But if you play chess, if you observe the game of a great chess player and watch how.
he makes every move and defeats his opponent, sometimes the move he makes, I don't understand..
Sometimes the moves that he makes or the strategies that he has in mind, I cannot comprehend..
Until at the end of the game, that I saw something amazing happen..
So the great chess player sometimes makes a move that is beyond my comprehension, where.
he plays up a pagan king..
That is beyond my understanding..
That is beyond my so-called plan..
But at the end of the game, you will see the wisdom of the great chess player..
And at the end of the game, if you go to the game board, you would not praise the chess.
pieces..
You would not pick up the chess pieces and say, "Wow, those pieces are very amazing..
They can do amazing moves, unpredictable moves, and finally they win the game.".
You would not praise the chess pieces..
The glory and victory go to the chess player..
The one who is in control..
The one who controls everything and every piece, every move on the game board..
It is the same way Isaiah is talking about..

$^{321}$He is saying that God's people need to know that we need to acknowledge the sovereignty.
of the one true king, Yahweh, Jesus Christ our Lord..
And our allegiance and our faithfulness is always on the only one true God..
The God of history..
The God of our nation..
The God that is to be worshipped..
Whether the pagan king in the White House or the pagan king in power, whether he will.
implement policies that would seem favorable to us or not, there is only one true God and.
one true king..
And he is in heaven..
Isaiah 49 says, "God says, 'It is too small a thing for you to be my servant to restore.
the tribes of Jacob and bring back those of Israel I have kept..
I will also not only return the people to their land and their exile, but I will also.
do one thing more amazing..
Is that I will make you the people of God a light for the Gentiles, that you may bring.
my salvation to the ends of the earth.'".
You see what is going on here..
There is only one God, one true king in the universe..
And he is above all political powers..
Kingdoms rise and fall..
However there is only one God who is eternally ruling, reigning over all of us..
In our nation, in your personal life..
And not only that, God says, "I will do something more..
I will raise you up..
I want to raise you up to be the light for the Gentiles, for the nations..
And with you, you can bring salvation, the gospel to the ends of the earth.".
What a timely word and what timely reminder that we have..
Because the Christian public witness in our nation is at all time low..
People look at us..
People go to church..
People profess the name of Jesus Christ..
They look at us like you guys are hypocrites..
You belong to the religious establishment..
It's a timely word for us..
And the Christian public witness is at all time low..
Especially in our nation..
And sometimes it's tempting for us to put too much faith on human kings..
I remember I was in seminary class and I heard the professor talking about Cyrus, Isaiah.
44 and 45..
And he's trying to spin the story that Cyrus the Great was probably converted to be a believer,.

$^{361}$to be a believer of Yahweh..
That's why he was favorable to the Jewish people and let them return..
Now that was in the old days when we didn't have much extra biblical document..
As you can tell from the documents that we read, Cyrus sitting there and all that, Cyrus.
remained a pagan king..
There was not a shred of evidence that he was converted..
However, this is our fantasy..
We want someone in power to be converted to be a Christian..
And we all want the kings to become like us, Christians or evangelicals..
How nice, right?.
Of course we want that, we want to pray for that..
We want those who are in charge and those who are with the political power, we want.
them to acknowledge the name of Jesus..
Or sometimes we want or we wish a superstar, a celebrity would become a Christian and she.
or he will bear witness to the world how great Jesus is..
And how nice is that?.
However, Isaiah reminds us, we cannot outsource the public witness to other people..
We cannot outsource the public witness to those in power or in the entertainment business.
or some celebrities..
We cannot outsource this..
Isaiah is talking about to every single person of the people of God..
You are my witness and God wants to make you the light for the Gentiles, for the nations..
God wants to raise you up..
You see, no matter which political party is in charge, no matter what policies they implement,.
whether it would be favorable to us or non-favorable to us, that is not the key..
The key is that we respond to the calling of God to be the light for the nations..
And we ask that the church will raise us up..
We will stand up to be the true church..
Think about the last past two decades..
Two decades or so..
There are so many mega churches that fail, crumble..
I don't have to give you examples..
I was in Toronto for the weekend and a lot of Canadian Christians told me what happened.
in the city of Toronto..
In the metropolitan area, there are also many mega churches or mega pastors figures who.
fail because of scandals..
You see what's going on?.
We fail..
It's not because the political order is in our favor..
It really doesn't matter..

$^{401}$We fail from within..
It's not because the political order is favorable to us to question weakness or not..
It is us that we need to be the light to the nations..
And we bear the name of God..
And we want to bear a positive witness to the world..
A couple of weeks ago, we had a small group gathering for the husbands..
For some of the husbands, we share our stories and our worries for our family, especially.
for the young kids..
Growing up in a liberal environment and they got their value confused..
And there are so many cultural wars out there..
And we exchange a lot of stories about our worries about what is going on in the schools,.
what is going on in the curriculum..
And we worry about that..
But I was talking to the husbands, I said, "You know what, whatever, however the political.
order is favorable to us, there is a front line we cannot give up..
And it is our own home..
Our own family..
How can we live out the gospel, the beauty of truth, the beauty of truth, to let our.
kids and let our next generation to know that the truth is beautiful..
You know, the next generation is tired of us telling them what is right and what is.
wrong..
Of course there is truth and falsehood we have to discern..
That is important..
But the most important thing is that we as the people of God, we can show and demonstrate.
the beauty of truth..
The beauty of marriage..
The beauty of manhood and womanhood..
The beauty of Christian life..
That we live out the gospel..
And that is what the world is dying for..
They want to see the people of God being raised up to be the light for the Gentiles..
And today Isaiah gives us this reminder and this calling..
You go back to your role and your responsibility..
Go back to your home to deal with your children..
To handle relationships in your home..
To handle relationships and ethics in your workplace..
In the society..
How can we speak with the gospel..
Speak with compassion..
Speak with peace..

$^{441}$So that people can see the beauty of the gospel..
Can we do that?.
Can we pray at least?.
Toward that direction..
God help us..
It's not about voting whoever is in the White House..
Of course if it's favorable, praise God..
But the front line that we all are on, we cannot give up..
We are in the trenches to be the light for the Gentiles..
To bring salvation to the ends of the earth..
And that is you and me..
That is all the churches, all the ordinary Christians we can do and we can pray for..
Can we close in words of prayer?.
And just pray and come back to God and say God call us..
And Father you do want to call us to be the light for the Gentiles, for the nations..
And we come before you, we ask that you will forgive us..
You forgive us sometimes we pay so much attention to okay who is in power but of course this.
is our public witness..
This is what we can do for the coming grace for the society..
However, what is most important is that you are in charge and you are calling us..
Each and every one of us to be faithful people..
To preach the gospel..
To share the good news..
In our conversation and in our interaction with other people that we can show them the.
beauty, the goodness of the gospel..
So help us and forgive us..
And please use us..
And this is our prayer as your people..
Humble us so that we worship you alone..
Humble us so that we can display and demonstrate your goodness, your beauty..
And may your name be glorified..
This is our little prayers and you send us out this week..
Go back to school, go back to our home, go back to our workplaces..
Any place you send us and help us to bring the truth and the beauty of truth..
And use us to show your glory in Christ in prayer..
Well thank you very much for worshipping with us..
And this series will continue next week and I hope to see you next week and we'll continue.
with Flat On..
God's psalm over the nations, over us..
So we have some refreshments out there and please stay for a little bit time of fellowship.

$^{481}$and get to know one another, to support one another..
God bless you the rest of the week and see you next week..
[BLANK AUDIO].
\newpage



\section{}
\label{sec:fC6e7drn7H4}
\textbf{感恩,因為.感恩,縱然.Give Thanks because... even though... (陳崇基牧師 Rev. Mark Chan)}
\newline
\newline
連結: \href{https://youtube.com/watch?v=fC6e7drn7H4}{\texttt{https://youtube.com/watch?v=fC6e7drn7H4}} ~~~~ 語音日期: 2024-11-26
\newline
\newline
\hyperref[sec:dQy1lcfV2Eg]{\small{< < < PREV SERMON < < <}}
~
\hyperref[sec:index_chronic]{\small{[返順時目]}}
~
\hyperref[sec:index_scriptual]{\small{[返順卷目]}}
~
\hyperref[sec:HjvG0k680po]{\small{> > > NEXT SERMON > > >}}
\newline
\newline
$^{1}$Welcome to all of you..
[speaking in Cantonese].
I have a very short message to you..
[speaking in Cantonese].
This is a very important message, okay?.
Even though it's a short message, but a very, very important message..
And this verse in the Bible says,.
"Give thanks in all circumstances,.
for this is the will of God in Christ Jesus for you.".
[speaking in Cantonese].
Now, you have to pay attention..
The Bible does not tell us that we have to give thanks for everything,.
including evils, including injustice, including sufferings in our lives..
The Bible is not telling us to do this..
[speaking in Cantonese].
The Bible does not tell us to give thanks for everything..
[speaking in Cantonese].
The Bible is telling us, "Give thanks in all circumstances,.
including unhappy circumstances,.
including some of the painful moments that we have.".
We can give thanks, and we have the resource.
to give thanks in all of these circumstances..
That's what the Bible is talking about..
[speaking in Cantonese].
Now, we thank God..
We can thank God to be with us,.
to be present with us in hard times, in difficult times..
[speaking in Cantonese].
And we can thank God that nothing,.
including unhappy moments, all of this happened outside,.
nothing outside of his wisdom for us..
[speaking in Cantonese].
I don't know if you have played this game called "I Spy.".
And I give you an amazingly detailed picture,.
and you have to find what I want you to find..
Have you ever played that, "I Spy"?.
Something like this..
In a very complicated picture,.
you have to find the happy, smiling onion, all right?.
Can you see that on the right?.

$^{41}$You can find it, right?.
Why don't we do this game right now?.
[speaking in Cantonese].
Of course, you can find it right now,.
but I'm going to show you a picture, the first picture,.
and you have to find me the happy, smiling onion, all right?.
[speaking in Cantonese].
Why don't we do it by table, okay?.
[speaking in Cantonese].
See if the people around the table can find the happy onion, all right?.
The first one is this one, all right?.
So let's find the happy onion, all right?.
Happy onion..
If you find it, raise your hand, okay?.
Now, I must tell you there are three..
There are three happy onions here, okay?.
Where is it?.
Okay, is it on the right?.
Is it on the right?.
On the right or on the left?.
Okay, there's a big one right there,.
and then there's a right one in the lower right..
And can you find the third one?.
This table can find the third one..
Where is it?.
Bottom left, yes, yes, yes..
Can you find it?.
Yeah, these are the happy onions..
[speaking in Cantonese].
Okay, now you have to pay attention..
Okay, now how many of you couldn't find it?.
[speaking in Cantonese].
I'll give you a second chance, all right?.
Second one, ready?.
Okay, we do it by table, all right?.
Here's the second one..
All right, I'm sure you can find the first two..
Yeah, where's the third one?.
Where's the third happy onion?.
[speaking in Cantonese].

$^{81}$Where is it?.
[speaking in Cantonese].
You can find it?.
Okay..
Why is it taking so long?.
Oh, you got it?.
Where is it?.
[speaking in Cantonese].
Oh, there are many..
Okay, all right, here we go..
All right, there you go, there you go..
How many of you find it?.
[speaking in Cantonese].
Go on, happy onion..
[speaking in Cantonese].
Well, let me give you the last one, okay?.
Let's try it, all right?.
Ready?.
[speaking in Cantonese].
Here..
Where are the onions?.
Of course, you can find the first one and the second one..
Where is the third one?.
[speaking in Cantonese].
Again, you guys?.
Oh, because they're closest to the TV, all right..
Yeah, right there..
Wano, where is it?.
It's like right there, right there..
Okay, all right..
Here are the onions..
The happy onions, happy onions..
Well, this is something what we do when we give thanks, all right?.
The exercise of giving thanks or finding things to give thanks for.
is kind of like finding the onion, the happy onion,.
in the complicated picture..
[speaking in Cantonese].
Now, it's much the same way, okay?.
The Bible is not trying to tell you positive psychology..
[speaking in Cantonese].

$^{121}$The Bible is telling you there is an understanding of God.
behind giving thanks..
[speaking in Cantonese].
God has everything in control..
[speaking in Cantonese].
Now, what happened this year for you,.
I'm sure it's kind of like this picture, very complicated..
So I want you to do one thing for me,.
to look for ways that you should be thankful for this year..
So far, we have gone through 11 months in this year..
[speaking in Cantonese].
Now, you have to think..
You have to rewind..
Rewind what happened this year..
[speaking in Cantonese].
And look for things you might have missed..
[speaking in Cantonese].
And things you can be thankful for..
[speaking in Cantonese].
Now, this is what the Bible says..
It's a cultivation of your life..
Cultivate your spirit so that you can center on God.
who is in control and God who is leading you all the way through..
[speaking in Cantonese].
There's a famous African-American writer called Alex Haley..
[speaking in Cantonese].
And in the office of Alex Haley, he's a very famous novelist..
And in his office, there was a picture called "Turtle on Fencepost.".
"Turtle on Fencepost.".
[speaking in Cantonese].
And why does he hang this picture in his office?.
He says, "Every time I begin to feel proud of myself,.
I look at the turtle on top of the fencepost,.
and remember, the turtle did not get there on his own..
He had help.".
[speaking in Cantonese].
Now, I think this is the perfect analogy for our life..
[speaking in Cantonese].
Of course, this phrase "post-turtle" is often used politically..
Talk about dumb politician..

$^{161}$[speaking in Cantonese].
But this is an analogy for our life..
We didn't get there by ourselves..
We had help..
And ultimately, God has helped you..
God has brought you to this place, to this spot..
[speaking in Cantonese].
So, in the season of Thanksgiving,.
we can call this "Give Thanks Because Of.".
Give Thanks Because Of, because every time you look at this picture,.
as Alex Haley said,.
we don't get there by ourselves..
We had help..
And ultimately, God had helped us..
That's why we thank God..
[speaking in Cantonese].
Now, we can call this Thanksgiving exercise.
"Give Thanks Because Of.".
Give Thanks Because Of..
You have success and you have come to this place.
because of the well-being God has granted you,.
because of shelter and security..
That should not be taken for granted..
Because of the opportunities that is right in front of you.
so that you can grow, and that is granted by God..
And the small joys in your life,.
that should not be taken to be granted..
And ultimately, Jesus Christ and His cross,.
the love of God, the grace of God to be with you,.
that should not be taken granted for..
[speaking in Cantonese].
This is what we call "Give Thanks Because Of.".
Right?.
So try to rewind this year..
What had happened to you?.
What have you gone through?.
Think about the grace of God..
Think about the providence of God..
[speaking in Cantonese].
And then there's another level of giving thanks..

$^{201}$Not only giving thanks because of..
It's giving thanks even though..
Giving thanks even though..
[speaking in Cantonese].
Even though there were challenges and setbacks in your life,.
there's more inconveniences, frustrations that happened this year..
Unanswered prayers, perhaps many of those..
Imperfections in your life..
Imperfection struggles, flaws that make you more real..
Make people more real..
Make your family and your relationships more real,.
more authentic..
So think about some difficulties you have faced this year.
to help you grow in resilience,.
to help you grow in patience..
[speaking in Cantonese].
the encouraging prayers.
that knows how to be rejected..
These things....
God gave to us..
God made us go through.
these trials and 1957.
these hardships when we were university student so we could go forward.
[speaking in Cantonese].
We must remember Jesus Christ and His cross.
even though all this happened.
[speaking in Cantonese].
A life of thanksgiving.
Now as we close our today's worship.
[speaking in Cantonese].
I want to invite you to give thanks.
[speaking in Cantonese].
And we want to do something together as a church, as a family.
as a spiritual family.
[speaking in Cantonese].
And on your tables.
you find some ribbons.
some broken ribbons.
some strips of ribbons.
these ribbons are for gifts.

$^{241}$okay, wrapping up a gift.
[speaking in Cantonese].
Now I want you to do something together.
let's do something together.
let's rewind what had happened this year.
[speaking in Cantonese].
What is worthy of giving thanks?.
Right, either setback or challenges or opportunities.
or something that is good.
some blessings in your life.
count the blessings and count the things that we should be thankful for.
and I want to invite you to use the permanent marker.
to write it on the ribbon.
alright, this is your gift of giving thanks.
alright, let's take a moment.
to think about what we can give thanks for this year.
and think about it, it takes time.
look into the details.
[speaking in Cantonese].
Now I want to invite you to do so right now.
it can be very general.
in the well-being, health, success.
small successes in your life.
[speaking in Cantonese].
I want you to write it on the ribbon.
I want you to do one thing today.
if God has reminded you.
at least one thing you should be thankful for.
maybe one thing you have missed to give thanks for.
I want you to write it on the ribbon and then you bring it up here.
you can leave your seats and after you have written.
the items of Thanksgiving on the ribbon.
you can hang it on the cross in front of me.
[speaking in Cantonese].
or do a small prayer.
you say, "God, thank you".
you write it down and bring your ribbon.
and hang it on to the cross in front of me.
[speaking in Cantonese].
[music].

\newpage



\section{以賽亞書}
\label{sec:HjvG0k680po}
\textbf{誰是你主?\#1 居魯士 (以賽亞書 45) - 陳崇基牧師}
\newline
\newline
連結: \href{https://youtube.com/watch?v=HjvG0k680po}{\texttt{https://youtube.com/watch?v=HjvG0k680po}} ~~~~ 語音日期: 2024-11-26
\newline
\newline
\hyperref[sec:fC6e7drn7H4]{\small{< < < PREV SERMON < < <}}
~
\hyperref[sec:index_chronic]{\small{[返順時目]}}
~
\hyperref[sec:index_scriptual]{\small{[返順卷目]}}
~
\hyperref[sec:sCP4BPS1_XA]{\small{> > > NEXT SERMON > > >}}
\newline
\newline
$^{1}$大家好.
Andrew和Sam,謝謝你們.
因為他們是開車的.
另外是Amanda和Brenda.
幫我跟從前往後.
不然就應付不來.
應付不來,人老了.
原來真的老了是不同的.
好像充電都充不到.
睡來睡去都睡不到.
所以過去的週末真的很忙.
唯一一件非常可惜的事.
我選不到一些我很想吃的東西.
因為主人家已經幫我安排好了.
所以下一次再去多倫多.
只吃東西,好嗎.
因為人老了.
所以有些事以為說了就沒說.
以為說清楚或澄清了.
但沒有澄清.
我聽自己說了才知道.
原來第一課有些事是說得失準了.
有些甚麼事失準了.
首先要說清楚.
就是《大義利書》第七章第九節.
他看到耿古常在者在他面前.
然後上面高等法院的景象.
有寶座設立.
寶座應該是眾數.
有很多個寶座.
但只有一位坐在寶座上.
就是耿古常在者.
Ancient of Dates.
這個應該要澄清一下.
以致人們不會迷惘.
另外有一段經文在《大義利書》.
或者在不同的地方也有說過.
這個沒有機會放在那幾課上.
就是《大義利書》說到.
上帝的名應當稱重.

$^{41}$從耿古直到永遠.
智慧和能力是屬乎他的.
他改變時間,季節.
廢王立王.
將智慧賜給智慧人.
知識賜給聰明人.
廢王立王和改變時間.
是上帝掌管時間和掌管地上的政權.
這個在《大義利書》很清楚的一種.
屬靈的意義.
讓我們知道今時今日.
我們怎麼看世界上發生的事情.
或者不同的時間改變.
不同的政權改變等等.
當然剛剛那一週.
美國最大的新聞就是.
川普被選為.
下一屆的總統.
這個不僅是我們很關注的.
也是其他國家的人都很關注.
我去到加拿大的時候.
經常有人問我.
你覺得怎麼樣.
我也有自己的看法.
這個看法.
我自己的看法.
就是這篇文章所說的.
所以在這幾個星期裡.
我想跟大家分享一段.
很短的一系列.
先將這個系列插進去.
就是誰是你主.
我們看看聖經裡面.
一些不信主的君王.
他們有什麼背景和經文裡面提到.
這些君王.
這次選了第一個的幸運兒.
就是居老使客 Cyrus.
Cyrus the Great.
在《二創教書》44-45章都有提到.

$^{81}$我們的步驟是這樣的.
我們首先了解背景.
因為聖經裡面說到這些不信主的君王.
無論提名也好.
不提名也好.
或者說到他們的國度.
他們的帝國.
都有他們的歷史背景.
如果你看一看聖經的話.
你也知道.
在聖經的世界裡面.
在聖經的歷史裡面.
有很多不同的帝國.
或者有不同的君王和王國.
你看看這裡.
我只標記了一些比較.
我們人少會留意到的.
羅馬,希臘,黑人帝國.
亞述,波斯,巴比倫.
下面埃及.
中間那些我叫做小老鼠.
紅色的全部都是大貓.
中間那些.
以色列猶大當然是在世界的中心.
在他們的右上角就是亞蘭國.
Aran或者Aramians.
左下角就是他們剛剛開始建國的時候.
蘇羅,大衛的時候.
有菲律賓人在沿海那邊.
這些帝國和帝王都有興有衰.
所以我們的步驟是這樣的.
我們先了解背景.
然後我們看看聖經裡面怎麼說.
這次第一次的幸運兒就是居老使.
居老使我做了一些我的signature.
我最喜歡做這些事情.
把這些時間表放出來.
時間的圖放出來.
你看到中間那裡.
750,700至500的紫色那裡.

$^{121}$看到嗎.
那是主前的時間.
主前750,700,650,600,550,500.
這是主線的主軸.
當時發生什麼事.
大概是八世紀至主前六世紀的時候.
你看到上面有紅色那個.
看到嗎.
紅色那條是911開始.
很容易記住.
亞述帝國.
新亞述帝國興起的時候.
最厲害的時候是911的時候開始.
直到605.
他們的國土崩潰.
亞述帝國崩潰.
同一時間崩潰的裡面.
差不多最後的大概100年左右.
就有另外一些國土開始崛起.
藍色那個就是Median.
或者The Medes.
或者馬代的時候.
說馬代,波斯,馬代,波斯.
馬代先的.
就是678至550.
同一時間有巴比倫帝國.
那個尼泊爾加里沙.
或者大耶利書的背景.
626至539.
然後馬代帝國就被哪個帝國吞併.
就是先被波斯帝國吞併.
就是550年的時候.
他將它打敗.
然後擊潰了馬代帝國.
然後就吞併了.
所以我們說馬代,波斯.
滅了Persian Empire.
還有什麼看.
這裡有很多東西看.
還有一些重要的黃色線.

$^{161}$我們左邊開始.
你看到就是722年.
就是以色列北國的首都.
他們的京都薩馬里亞淪陷.
那時候是722.
誰打敗他.
哪個帝國去令薩馬里亞淪陷的.
就是亞述帝國.
沒錯了.
大家很聰明.
亞述帝國.
然後他繼續打下去的時候.
打到當時是希西加的時候.
就是主前701年的時候.
打不贏.
進不了耶路撒冷.
不過有圍剿他.
當時也是亞述帝國的時候.
701年.
過了一段危機的時候.
亞述帝國開始式微.
然後有巴比倫帝國出現.
到了586年的時候.
就是耶路撒冷淪陷.
南國和猶大國淪陷.
淪陷了就被擄了.
流亡了大概70年的時間.
直到70年末的時候.
有其他事情發生了.
就是波斯帝國興起.
然後在巴比倫帝國末期的時候.
有兩個末代的王.
你看到55和539.
那時候有兩個所謂末代的王.
一個是聶布尼德和賓夏爾遜.
這兩個是父子.
他們負責末代的時候.
巴比倫的朝政.
當時已經式微了.
有誰崛起呢.

$^{201}$就是大名鼎鼎的居魯士.
居魯士是怎麼崛起的呢.
居魯士和他父親.
一直服侍馬代帝國.
大家不太順暢.
所以最後553年的時候.
他就反叛了.
他不想再做梵容國.
然後他就去反叛.
反叛了帝國.
然後呢.
三年之後.
他就將馬代帝國吞併了.
然後他就入主巴比倫帝國.
把巴比倫帝國擊敗了.
所以他就統治當時的天下.
這個就是Syrus.
居魯士.
很厲害.
我們很喜歡看這些東西.
因為我們很喜歡看戰國.
或者戰爭的英雄故事.
居魯士是其中一個大名鼎鼎的大王.
所以我們不是叫他居魯士二世.
我們叫他居魯士大帝.
因為他這麼厲害的緣故.
先講一下他的背景.
他怎麼崛起.
正如我所說.
聖經裡面有提及過這個人.
但是很少會直接跟我們解釋.
因為他是假設了當時第一代的讀者.
明白當時的改朝換代發生了什麼事.
你看看.
很簡單.
這裡雖然很亂.
巴比倫帝國末期的時候.
兩個王.
就是拿布拉登和貝撒夏.
當時發生什麼事.

$^{241}$拿布拉登王拿布拉登.
他在最後的十年.
他不是很理會朝政的事情.
離開了京都.
然後自己搬到了一個沙漠的綠洲.
然後拜他喜歡的神.
當然巴比倫帝國是多神的.
但是巴比倫帝國裡面.
巴比倫的文化社會裡面.
最大的神是誰.
叫馬杜克.
叫做馬杜克.
每個人都捧他.
但是這個拉登王很奇怪.
拉登王不喜歡拜馬杜克.
他就拜他自己家族的.
叫做SIN.
月神.
這個新月神.
他就不太理會馬杜克.
你知道當時的人.
他們敬拜這些忠神的時候.
其實每一個忠神都代表.
不只是宗教這麼簡單.
他代表經濟和政治的中心.
所以每個廟都等於中央銀行.
所以不管中央銀行.
然後自己去整另一個中央銀行.
當然很多人不順氣.
對吧.
你明明在捧那個.
大家都在捧那個.
他就不去捧那個.
去捧另一個.
所以令社會裡面有很多不安.
和很多不順氣的人.
特別是中央銀行在做事的人.
就是那些馬杜克,大廟.
大銀行裡面.
做的那些祭司和貴族的人.

$^{281}$正如你一群人吃飯.
明明大家都很喜歡姜濤.
你說不喜歡姜濤.
很大件事.
就是這樣的情況.
明明每個人都在捧那個.
拉登王就不捧那個.
所以在當時的巴比倫帝國裡面.
有很多不順氣的人.
有什麼不順氣的人呢.
拉登王喜歡月神.
你看到他的石板.
右邊那個.
上面那個是什麼.
月亮,看見嗎.
不是朝.
月亮,太陽和Venus.
三個星神.
左邊那個就是馬杜克.
或者在耶利米蘇提過的米羅達.
其實是同一個神明.
他就不捧那個.
最大件事就是每年一次.
巴比倫的京都都會有一個大的慶祝.
那個慶祝叫做Akitu Festival.
Akitu Festival其實是春季的一個大慶典.
當時剛剛大麥,小麥收成的時候.
所以是一個農耕的季節.
每逢這個時候就會將馬杜克.
就是馬杜克我們最喜歡的那個.
我們最捧的那個.
我們大家都是粉絲.
我們就將他拿個像出來.
然後就巡遊整個京都和不同的城市.
然後有很多很大的慶祝.
拉登王是不喜歡這個神明的.
所以他就離開了京都.
沒有搞這些事情.
很大件事 你知道嗎.
每年台慶你都不來.

$^{321}$每年台慶你都不去捧台.
所以人們就很不滿.
當時如果你拿個帖文去看一下.
就是那些民意調查.
最近我們經常看的.
現在選了 不用看了.
之前經常看的什麼帖文.
他們有什麼民意調查.
那個popularity是怎樣.
他們是否受歡迎.
政客是否受歡迎.
如果你看一下.
最後的巴布林的十年裡面.
你看一下那個帖文的話.
你就會看到他們很不滿當時的朝政.
正是一個真兇.
就是Cyrus.
居老時 就是興起的時候.
正如我剛才所說的居老時.
他開始反叛馬代帝國.
1955年的時候.
然後用了三四年時間.
就將馬代帝國擊敗.
一直吞併到左邊.
雷底亞.
然後再回來諾邁爾蘭.
然後539年.
是最重要的那一年.
他就入主巴布林的京城.
然後就吞併了整個巴布林帝國.
然後就結束了猶太人.
被擄巴布林的大概70年時間.
你看到539年就是猶太人第一次回歸.
還有回歸之後.
他們想重建耶路撒冷.
還有重建聖殿.
第二聖殿重建成功就是515年的時候.
很亂吧.
這麼多數目.
最重要的是什麼呢.

$^{361}$最重要就是Cyrus.
所謂是乘虛而入.
巴布林帝國已經到末期的時候.
裡面有很多的混亂和張力.
特別是政治上的張力.
有些人已經很不相信這個政黨.
所以人們就已經很想.
有另外一個人來就太好了.
Whoever that is.
所以這個時間.
就正是Cyrus能夠吞併天下.
以及統領天下.
進入巴布林.
在《聖經》以外.
有很多文件提到Cyrus.
譬如這個東西.
在大英博物館裡面看到的.
我抽了多一點給大家看.
他說什麼呢.
他說我是居老時.
這個東西是紀念他的豐功偉績的.
他說我是世界之王.
四極之王.
他和平的進入巴布林.
馬杜克每天拜祭他.
民間非常歡喜.
靠著馬杜克最高的華.
因為馬杜克是提升居老時.
能夠使他成為世界之主.
他就親他的腳.
Kiss his feet.
然後將各地的神像歸還原處.
然後將他的眾民歸還原處.
正如我所說.
就是猶太人其中一批.
就可以回到他們的故土耶路撒冷.
當然這個只不過是其中一個撮腰.
之前這個東西.
之前的上文說什麼呢.
他就去摒以前的巴布林末代皇帝.

$^{401}$Nobilus.
就是說他怎麼不好.
因為他又不捧馬杜克.
他又不捧他們最大的神明.
所以他就失誤了.
令民間非常多疾苦.
令民間很不順暢.
所以他就不斷批評他.
然後到這裡.
他說不是了.
我就不是.
我是怎樣.
我是每天敬拜馬杜克.
我想問你.
這段經文.
居老是一個怎樣的人.
他是不是信耶路撒冷的.
他不是信耶路撒冷的.
他信什麼神.
他是信馬杜克.
巴布林拜的神.
他就去拜.
還說巴布林神去建立他.
使他的國位能夠興起.
並且他是波斯人.
波斯人是拜波斯火教的.
波斯火教有一個最大的神.
應該是居老是自己拜的神.
那個叫做Huram Master.
大家有沒有開車是開Master的.
這個神明的名字.
就是這個意思.
Huram Master.
Master在波斯文裡面就是Wisdom的意思.
你開Master的話.
小心一點開車.
他是拜多神的.
Syrus the Great.
只不過是一個很實用的politician.
一個很實用很pragmatic的politician.

$^{441}$乘虛而入.
然後進入了巴布林帝國.
然後他就做一些.
跟以前帝國的policy相反的事情.
以前就抓了人去帝國.
現在他說OK釋放你們.
然後將那些神像歸還.
他為什麼pragmatic.
Syrus只不過是想趁著這個機會.
來鞏固自己的王位.
這個就是聖經以外所說的居老是.
聖經裡面怎麼說居老是.
我們就進入以賽亞書第44章和45章.
我們先看這段經文.
提到居老是.
這個是非常神奇的東西.
因為在聖經裡面提到一個外邦的王.
還有提到他的名字是比較少有的.
我們看看這段聖經說什麼.
他說雅各.
順帶一提這是以賽亞書第44章.
以賽亞書大概分兩段.
第一個就是1至39章.
然後40至到最後66.
你算算.
就是叫做第二以賽亞.
就是第二部分的以賽亞.
第一部分的以賽亞.
就是寫給在亞述帝國危機的那幫猶大人.
然後第二部分40章之後.
就是預言將來他們被擄巴比倫.
流亡到巴比倫的時候.
他們會發生什麼事.
還有神怎麼眷顧他們.
怎麼作供.
這個就是在第二以賽亞.
或者以賽亞書第二部分出現的.
44章.
雅各.
你要思念這些事.

$^{481}$以色列.
你是我的僕人.
我做了你.
你是我的僕人.
以色列.
我必不忘記你.
你當歸向我.
因我救贖你.
論居老時.
Cyrus.
Cyrus the Great.
他是我的僕人.
他要成就我所喜悅的.
下令重建耶和華山.
和發命令立穩聖殿的禁忌.
這裡就看到居老時這個人.
因為耶和華對所高的居老時如此說.
他的右手.
我曾參赴.
使列國降服在他面前.
列國的腰帶.
我曾鬆開.
使城門在他面前躺開不得關閉.
這個非常驚奇.
如果你看回時間線的話.
如果傳統的意思.
以賽亞是什麼時候事奉.
什麼時候做先知.
是主前740年.
到700年這段時間.
當時還是亞述帝國統治天下.
都沒有居老時這個人.
他連產房都還沒有到.
他已經在100年多前.
提到居老時這個Syrus.
並且用真名來提及他.
當時是一個預言.
你就很奇怪了.
為什麼他會說居老時是被告的.
當我們一提到受告者.

$^{521}$Messiah.
Anointed.
我們很多時候用了後期的意義.
套上去前面.
我們後期的.
當然我們說Messiah.
詩歌唱Messiah.
唱利塞亞的時候.
基督.
我們當然是說救主.
我們唯一一個的救主.
The only savior.
The only Messiah.
就是耶穌基督.
但這是新約的時候才有的意義.
在舊約裡面.
很多時候用Messiah這個字.
或者被告者.
很多時候不一定是他的神聖.
告那個人不一定是那個人.
是合上帝的心意的.
很簡單.
在撒慕爾記上帝十章.
撒慕爾去到高納第一個王.
以色列國第一個王.
就是誰.
Saul.
King Saul.
去告他的時候就說.
耶和華豈不是告你Messiah.
這個字.
或者動詞.
冥詞動詞.
告你做什麼.
作產業的君王.
作以色列的君王.
如果你知道舊約的故事.
Saul是不是一個好的王.
Saul是不行的.
人間覺得他行.

$^{561}$是民選出來的.
大棒.
然後Saul就是不信服上帝的.
並且是背弱的.
最後大衛去取代他.
所以第一個王也好.
叫他受告者也好.
其實也不是一定代表.
用到受告者的字.
代表他是聖聖的.
或者代表他是神聖的.
我們說《但以理書》.
過去一個星期的時候.
我們留意到一樣東西.
也是沒有時間說.
沒有空間說.
大家記得《但以理書》第七章.
有四隻怪獸出來.
就是四個帝國.
其中一個帝國肯定.
第二個或者第三個.
無論你怎麼配對也好.
一定是波斯帝國.
這個很奇怪.
波斯帝國對猶太人好不好.
不是居老師.
不是波斯帝國.
猶太人是不能回歸的.
對他們好不好.
政策上是非常支持.
非常支持.
各族的人民.
包括猶太人.
但是《但以理書》第七章的意象.
包括其他的意象.
說到波斯帝國.
是什麼帝國.
是怪獸是嗎.
是怪獸.
是一個從大海出來的.

$^{601}$混沌的化身.
所以這裡說到.
居老師是上帝的被告者.
或者受告者.
不是說他一定是神聖.
縱使他的政策是支持地上的猶太人.
被流亡的猶太人.
但是在上帝的眼中.
上帝只不過是用他作為其中一步棋.
上帝的權能是什麼.
上帝的權能就是可以用外邦的君王.
異教的君王.
來作成他自己的事.
那個榮譽.
或者那個信任.
不是在那個人身上.
而是在耶和華身上.
如果查看《耳塞夏書》4-45.
或者整個《耳塞夏書》說的是什麼.
是耶和華自己的榮耀.
以及他的權能的彰顯.
不是將那個人.
或者那個很潛在的政治家.
Cyrus 將他放在那個台上.
然後將耶和華自己的榮耀彰顯出來.
因為上帝的權能是掌管萬有.
並且列國的君王.
不信耶和華的君王.
都在耶和華的手中.
45章第五節.
他說我是耶和華沒有別的.
你看清楚.
真正要敬拜真正要事奉的.
是至高者耶和華.
除了我以外再沒有上帝.
不是地上的君王.
而是我在天上的掌管一切的上帝.
是我是耶和華.
我做光 我做暗.
我施平安.

$^{641}$又降災禍.
造成這一切是我耶和華.
這裡用了幾個希伯來的動詞.
Yassar Barra.
Assar等等等等.
我們就想起.
我看到這些字眼的時候.
我們就說上帝為何會做暗.
或者創造一些不好的東西出來.
是不是.
降災禍.
上帝沒理由是邪惡的源頭.
我們想起這件事的時候.
我們很多時候想起從無到有.
那種創造的過程.
但是這些希伯來的動詞.
在背景下.
其實最主要想說什麼.
不是說從無到有.
明明是沒有災禍變成有災禍.
那種情況.
那些動詞想表達的.
就好像創世記第一章.
主要不是說創造出一無所有.
當然這個教義是對的.
但是創世記第一章.
用到同樣這些字眼的時候.
主要不是說創造出一無所有.
主要是說上帝的管理.
主要想說上帝不是明明沒有災.
變成有災.
不是這個意思.
而是上帝掌管黑與暗.
上帝有權能來管轄光與暗.
平安與災禍.
縱使有災禍.
上帝都可以去掌管它.
就是這個意思.
正如你老闆問你.
你可不可以下星期.

$^{681}$創造一個部門出來.
來去照顧這件事.
你就不會跟老闆說.
我怎麼創造 我不是上帝.
我不可以創造出一無所有.
創造是什麼意思.
就是組織.
叫你組織一個部門出來.
叫你組織一個團隊.
或者安排去管轄一個團隊.
分配它的功能.
以至工作能夠完成.
就是這個意思.
所以這裡說的是什麼.
上帝做平安也是做災禍.
主要的意思是.
他管轄這件事.
沒有一件事在上帝自己的手以外.
全部都在上帝的五指山裡面.
就是這個意思.
管轄.
上帝說我掌管一切.
連不信主的君王我都可以興起.
並且可以將他廢.
你要敬拜的 你要事奉的.
是我.
大家有沒有下棋.
仗棋 或者什麼棋都好.
有誰下棋很厲害.
我就很差.
我經常被人叫做「摸缸豬」.
但你看人家下棋的時候.
特別是那些棋王.
很厲害的棋王.
你真的猜不到他每一步是怎麼走的.
看他們下棋的時候.
你會說為什麼會被人擲了一支炮.
是什麼意思.
誰知道他幾步幾步之後.
他就可以反敗為勝.

$^{721}$棋王厲害的就是這樣.
你看完整個棋局的時候.
你不會捉著棋子.
這支炮真厲害 將軍秋車.
這支炮真是很好.
當然不會.
你會看什麼.
你會讚嘆的是背後的棋王.
怎麼用一支棋子.
在耶和華值得我們去榮耀.
以及值得我們去事奉.
這是整個觀點.
在以賽亞.
雖然上帝可以興起不信主的君王.
是一些純粹是政治計算的君王.
他也可以廢棄他.
但這部棋在棋王 耶和華身上.
最主要他想說什麼.
最主要他想說.
既然我這樣做的時候.
你就需要成為我手裡的棋子.
以色列你要做我的手裡的棋子.
上帝的子民你要這樣做.
他說 你作我的僕人.
從雅各的眾志派能夠復興.
使以色列中蒙保存的人歸回.
然而此事尚小.
我還要使你作萬邦之光.
施行我的究竟直到底的.
這就是此時此地.
在每個世代 特別我們這個世代.
基督徒需要做的事情.
我想告訴你 哪一個當選都好.
無論政策是在我的倫理.
或是我的看法上.
是在我的看法上.
但最重要的一件事.
我們要做的.
就是我們的見證.
我想告訴你.

$^{761}$美國到今時今日.
基督徒的公共見證.
是All Time Low.
It's All Time Low.
無論朝政發生什麼事.
我們教會在地上的見證.
是All Time Low.
我們很多時候被人印象.
就是我們這些假冒為善的人.
Hypocrite.
我們只不過循例做一些宗教的儀式.
但我們的信仰沒有落地.
在這個世代裡面.
我們很多時候把我們的信心.
放在白宮.
放在政策.
如果有政策 這樣就好了.
當然這樣是非常好.
That's OK.
但有件事我們不能忘記.
不能推下來.
就是我們的前線是靠我們自己.
作為上帝萬磅之光.
政黨是否合你心意都好.
我想問你.
對於我們在地上的使命.
我們的見證.
有什麼影響.
沒有影響.
我告訴你 這個朝政.
Favor我們.
That's good.
But it's not good enough.
如果我們不起來.
做回我們自己的角色.
我們不守住我們的前線的話.
一切都是徒然.
This is All Time Low.
我們Public Witness.
但怎樣我們能夠教會興起.

$^{801}$這個是真的任重道遠的問題.
是我們每一個都需要問.
我想問你.
過去好幾年.
只算我們之間的割裂區.
我去到加拿大也聽到其他個案.
一個大的教會.
好像一夜之間崩潰.
因為不同的醜聞.
不同壞的見證.
我想問你這些所謂Mega Churches.
或者Mega Pastors.
很出名的Pastors.
他們跌倒.
是因為朝政不順他們的眼還是什麼事.
是否政黨是在迫害他們.
不是.
因為我們心靈裡面有問題.
是我們自己搞砸自己.
不關政黨和公眾政策的事.
是我們教會失敗.
求上帝真的憐憫我們.
我記得上個星期.
是不是上個星期.
老人癡呆.
我和丈夫一起去多倫多之前.
大家有個晚飯.
我很享受那個時間.
大家談到很多社會上的事.
特別是家庭上的事.
我真的覺得這個時候.
要把小朋友養大.
有很多焦慮.
是非不分.
學校裡面混亂得不得了.
我們擔心.
當然.
顛倒真假或者指鹿為馬.
我們很想告訴下一代.
What is right and what is wrong.

$^{841}$當然我們也希望社會有所謂的common grace.
有一個普遍的恩典.
能夠給我們一些order.
或者moral order.
但我想說一件事.
有一個前線是不能夠我們退下來的.
就是我們的家庭.
家庭與誰在白宮無關.
家庭能不能不只是告訴下一代.
What is right and what is wrong.
但我們告訴他們.
The beauty of truth.
真理的美.
美善.
有基督的美善.
誰告訴他們.
不是公共的policy.
而是我們的生命.
對.
誰能夠告訴我們靈社.
耶穌基督的名是真理和美善的名.
值得我們敬拜.
不是白宮的那個.
是我們.
對嗎.
萬邦之光.
你和我都是萬邦之光.
縱使我們覺得自己的光很小.
好像一個小的蠟燭.
但越黑的時候.
你就能夠發揮那個功效.
對嗎.
弟兄姊妹們我們不要忘記.
上帝鼓勵我們.
當時的時間可能是一個.
在某種意義上是一個很好的時間.
我們要告訴真相.
我們要保持政黨的控制.
我們要令這個世界更多光線.
誰做.

$^{881}$我們不是給別人出口.
是我們自己做.
上帝的子民.
上帝要使我們成為萬邦之光.
願意我們這樣禱告.
這樣努力.
可以跟心得有寄託.
因為你呼召我們.
以及你將我們拯救出來.
叫我們能夠成為你自己的子民.
成為你自己國度裡面的人民.
在地上主我們的禱告.
真是.
願你天上的子也臨在地上.
臨到我們生命.
從我們生命開始.
從我們家庭開始.
從我們自己的生活開始.
叫我們不要忘記.
因為你這樣呼召我們的時候.
就給我們一個新的身分.
成為其他人的光.
主幫助我們.
讓我們的小角色.
我們更加忠心來服侍你.
去彰顯上帝自己的榮耀.
主要是為了接下來這四年.
美國的社會.
主要是我們知道.
我們仍然有很多工作需要做.
福音的工作.
教會的工作.
我們外面見證的工作.
主要是祈求你興起我們.
幫助我們 憐憫我們.
我們這樣祈禱.
父王主名.
願上帝祝福各位.
下星期我們繼續思想這個題目.
\newpage



\section{5:21-26}
\label{sec:sCP4BPS1_XA}
\textbf{Righteousness as Redefined by Jesus: Anger (Matthew 5\_21-26) - Chee Yie Jong}
\newline
\newline
連結: \href{https://youtube.com/watch?v=sCP4BPS1_XA}{\texttt{https://youtube.com/watch?v=sCP4BPS1\_XA}} ~~~~ 語音日期: 2024-12-03
\newline
\newline
\hyperref[sec:HjvG0k680po]{\small{< < < PREV SERMON < < <}}
~
\hyperref[sec:index_chronic]{\small{[返順時目]}}
~
\hyperref[sec:index_scriptual]{\small{[返順卷目]}}
~
\hyperref[sec:FaphUTtEgYU]{\small{> > > NEXT SERMON > > >}}
\newline
\newline
$^{1}$All right. Good morning, UV. I hope all of you had a great Thanksgiving this past week,.
whether it is spending time with family, with friends, enjoying some good and delicious.
Thanksgiving food, or also just buy some new stuff on Black Friday. So I got a chance to.
went to the suburbs and spend Thanksgiving with a family from my previous church, Dennis.
and Emily. So it was great to be able to catch up with them again and also hang out with.
three of their teenage kids, because all three of them were also my youth students at my.
youth group in my last church. And it is always fun to get to be reminded of kind of like.
the interesting or funny ways that those three siblings interact with one another. So all.
that to say, I had a great Thanksgiving. And I also want to acknowledge that as we enter.
into these holiday seasons, this might be a season that is not a joyful season for some.
of us, that this can be a challenging season, especially if there's conflicts within your.
families or this time of grieving for some of you. So I want to acknowledge that. And.
also, I just want to say, I hope during this Advent season leading up to Christmas, I pray.
that the Lord would meet each and every one of us in a specific ways as well. Okay, so.
this past September, I had the opportunity to go back home to Malaysia to visit some.
of my friends and also family. And when I step into my house, I immediately notice something.
different from the last time that I was home about seven years ago. So on one side of the.
living room is actually a whole new set of wooden cabinets. And without having to ask.
my mom much about it, I can pretty much guess what happened to the old cabinets. It is termite.
infestation. So growing up, termites is always an issue at my house. And if you don't know.
about termites, let me give you a quick biology lessons about termites. So termites, they're.
tiny little insects that look pretty much like an ant. However, they're a lot more destructive.
than our regular house ants. So for these termites, they like to feed on dead plants,.
they like to feed on trees or any sort of like wooden structures. And they usually prefer.
to hide in dark and humid environment. So you normally won't find them like feeding.
on the surface of the wooden structure at all. And in fact, the wood surface might appear.
completely fine, even though the termites are destroying them from the inside. And this.
is pretty much what happened to some of the wooden cabinets at my home. While they look.
totally fine from the outside, there were just a lot of hollow parts inside those cabinets.
that they're just not structurally sound anymore, to the point that they have to be completely.
replaced. And you guys have to know that this cabinet is a kind of like a central piece.
in like in my house in the living room. This is where we put our TV. This is where we put.
our sound system. And there's also a lot of like storage space where we put like our old.
photo albums. So it is a pretty important piece. And to see it being replaced because.
of the termites, it kind of like hurts me a little as well. So we have been slowly going.
through Jesus' Sermon on the Mount verse by verse for the past couple of months. And at.
this point, you might ask, "Hey, Chi, what does the termites or the termites infested.
cabinet have to do with Jesus' Sermon on the Mount?" Well, just like this wooden furniture.
that appears completely fine from the outside, but they're just not functional anymore because.
of the termites. The issue that Jesus had with how the scripture was being taught and.

$^{41}$lived out by the religious leaders of his days is that these religious leaders, they.
approached this law that Jesus gave them or that God gave them in the Old Testament from.
a very surface level, but they completely missed the heart of God behind the law. And.
as a result of that, the worship, just like this termite infested furniture, the worship.
became an empty religion that has pretty much lost its function and its purpose as well..
So in my past few sermons, we have looked at the Beatitudes and explored together what.
a truly blessed life looks like. We have also studied the different characters and qualities.
that represent life in the kingdom of heaven. And in my last sermon, we look at how we can.
experience this truly blessed life, that is by following Jesus and living in full submission.
to him and his commands. Jesus is the only way and only truth that will lead to a flourishing.
life because he is the only one that can fulfill the law and the prophets. And what this means.
is that the law and the prophets or the Old Testament scripture in its entirety, they.
ultimately points forward to Jesus and Jesus becomes the official interpreter of this Old.
Testament scripture. So today we're going to look at first of the six examples of how.
Jesus fulfilled the scripture by revealing the true intent and the true purpose behind.
this law. And just a quick background, each of the six examples that we will be looking.
at from Matthew five, verse 21 to 48, they all begin with the same formula. You have.
heard that it was saved. And this formula in Jesus's time, when a Jewish rabbi, he.
want to begin a discussion about a specific mosaic law, they will usually introduce the.
discussion with this formula. You have heard that it was saved and then the mosaic law..
So in the same way, Jesus begins each of his six discussion by using something that his.
audience will already be familiar with, and that is in line with the rabbinic tradition..
But each and every time Jesus would take this one step further and he would go on and give.
his own interpretation. But I tell you this. So if you have your Bible with you, please.
turn to Matthew chapter five and we're going to read from verse 21 to 26. Matthew chapter.
five, verse 21 to 26. You have heard that it was saved to the people long ago. You shall.
not murder and anyone who murders will be subject to judgment. But I tell you that anyone.
who is angry with a brother or sister will be subject to judgment. Again, anyone who.
says to a brother or sister, Raka is answerable to the court and anyone who says you fool.
will be in danger of the fire of hell. Therefore, if you're offering your gift at the altar.
and there, remember that your brother or sister has something against you. Leave your gift.
there in front of the altar. First, go and be reconciled to them. Then come and offer.
your gift. Settle matters quickly with your adversary who is taking you to the court..
Do it while you're still together on the way or your adversary may hand you over to the.
judge and the judge may hand you over to the officer and you may be thrown into the prison..
Truly I tell you, you will not get out until you have paid the last penny. This is the.
word of God. Let us pray. Father, as we open up your word this morning, I just ask that.
your spirit will be with us, that your spirit will fill our heart and prepare us to receive.
the word that you have for us today. Amen. Okay. So in my last sermon, we ended in verse.
20 looking at kind of like what kind of righteousness that is required for us to enter the kingdom.

$^{81}$of heaven, right? And for Jesus, the righteousness according to the Pharisees and the teachers.
of the law, they ultimately fall short of God's standard because their version of righteousness.
only concerns about outward actions or kind of like sort of like a behavioral management..
On the other hand, the righteousness according to Jesus, it doesn't just deal with our outward.
actions but it cuts straight to our inner motivations and our core desires that impacts.
how we relate to other people. So the example that we will be looking at today is centered.
around the sixth commandment, you shall not murder. And we'll take a look at how Jesus.
moved past just a literal understanding of this commandment straight to the root issues.
of these actions. And the message that Jesus has for us is that Jesus had for us is this..
While we may never commit the act of murder, we can still unintentionally allow the same.
kind of attitude to take root and permeate in our heart. I can probably say with 99\%.
confidence that none of us here in this room have committed a murder, right? But again,.
what Jesus is kind of like warning us against is that we can still unintentionally allow.
the same kind of attitude to take root and to grow in our heart. So let's look at the.
first point together. Jesus demands a higher righteousness from us. Jesus demands a higher.
righteousness from us. So Jesus begin by saying, you have heard that it was said to the people.
long ago, you shall not murder. So this first part actually, as we have looked at just now,.
this comes from the sixth commandment that the Lord gave to Moses, the Lord gave to the.
people of Israel through Moses at Mount Sinai in the book of Exodus. But then we see that.
there is a penalty being added on to this. It says anyone who murders will be subject.
to judgment. So this add-on penalty is actually part of the judicial process in the Jewish.
court system. And by attaching this penalty to the sixth commandments, it tell us one.
important thing. It tell us that Jesus, he's actually not speaking against the Mosaic law.
or the Torah itself. No, as we have established last time, Jesus himself clarifies, I have.
not come to abolish the law. So what Jesus is trying to correct here is the Torah as.
traditionally interpreted and understood by the Jewish people. One that is being taught.
in the synagogue or one that is being passed down in the family or in the Jewish culture.
as well. So Jesus, he is not reacting against the Mosaic law itself. But what Jesus is reacting.
against is a shallow and legalistic understanding of the law. So this is when the Lord gave.
the 10 commandments for the first time at Mount Sinai. Now in a similar setting on the.
mountainside, Jesus is trying to go back to the heart of God behind these commandments..
One that is not tainted by any of our human interpretation or human tradition. So what.
is the heart of God behind the commandments of you shall not murder? Let's look with me.
again. Verse 21, you have heard that it was said to the people long ago, you shall not.
murder and anyone who murders will be subject to judgment. But I tell you that anyone who.
is angry with a brother or sister will be subject to judgment. So here, Jesus, he goes.
beyond the shallow understanding of the sixth commandment straight to the root issues behind.
murder. That is anger and contempt towards brothers or sister. So the brothers and sister.
here usually used to refer to the communities of believers, just like how we call one another.
brothers and sisters in Christ. But I think in this context, we can actually apply this.

$^{121}$to those outside of our community of faith as well. So while the traditional understanding.
of the sixth commandment say that anyone who murders will be subject to judgment, Jesus.
brings in a higher righteousness by claiming that anyone who is angry with a brother or.
sister will be subject to judgment. So according to our human standard of righteousness, you.
might be completely innocent because you have never committed murder. But according to the.
righteousness defined by Jesus himself, your anger makes you guilty before God and liable.
to the same kind of judgment. And then Jesus continued, again, anyone who says to a brother.
or sister, "Raca," is answerable to the court. So "raca" here is actually a loan word from.
Aramaic, and it doesn't really mean much, but it is used as a derogatory term to call.
someone as empty-headed. So if you say someone's head is empty, that means there's like no.
brain inside, right? So this is kind of like similar to calling someone, "You have no brain.".
And Jesus say, "Anyone who says to a brother or sister, 'Raca,' is answerable to the court.".
And this court is the high court in the Jewish court system, which is the Sanhedrin. And.
anyone who says, "You fool," will be in danger of the fire of hell or Gehenna. So the NIV,.
which is the translation that I use here, actually smoothed this out a little bit. But.
fire of hell, in other translation, is being translated as Gehenna. And Gehenna is used.
to refer to a place that is south of Jerusalem called the Hinnom Valley. And what is the.
significance of this place, the Valley of Hinnom? In the Old Testament, the Valley of.
Hinnom is where child sacrifice were being offered to the pagan god, Molech. And we see.
this in places like 2 Kings. When King Josiah renewed the covenant with the people of Israel.
and with the Lord, this is what 2 Kings says. "He," referring to King Josiah, "He desecrate.
Topeth, which was in the Valley of Ben Hinnom. So no one could use it to sacrifice their.
son or daughter in the fire to Molech." And during Jesus' time, this place, Gehenna, or.
the Hinnom Valley, is used as a place where garbage were being disposed and then burned.
day and night. So if you just look out from Jerusalem and look towards this direction.
of the Hinnom Valley, you will see fire being burned day in and day out. And I can imagine.
the smells probably doesn't smell really good as well, right? So Gehenna basically becomes.
a metaphor for the fire of hell. So Jesus say, "Anyone who says, 'You fool,' will be.
in danger of Gehenna," or in danger of the fire of hell. And here in verse 22, we also.
see sort of like a progression. Jesus begins by talking about an internal attitude or feeling.
of anger. And then he move on to kind of like that anger being expressed outwardly by showing.
contempt to a brother or a sister. And what is contempt? I feel like this is probably.
a word that we don't usually use in our day-to-day language, right? So according to Oxford languages,.
contempt is the feeling that a person or a thing is beneath consideration, worthless,.
or deserving scorn. Again, contempt is the feeling that a person or a thing is beneath.
consideration, is worthless, or deserving scorn. And how is contempt being shown here.
in the examples that Jesus gave? It is by insulting a brother or a sister with names.
like "raka" or "fool." And showing contempt is actually not just about actively belittling.
someone. Showing contempt can take the form of passive indifference or being cold to someone..
Maybe at your house, your spouse has been trying to ask you to help with some of the.

$^{161}$things at the house, whether it is fixing a broken light bulbs or kind of like to mow.
the lawns or to take the trash out. But every time when your spouse kind of like bring this.
request before you, you might be listening, but you're actually not. And you're just responding.
in a way that is like trying to get them off your back so that you can continue with whatever.
you're doing. Or maybe you're walking down the street in Chicago and you saw a homeless.
person kind of like slowly approaching you. And at that moment, you just kind of like.
say to yourself, "Okay, I'm just going to ignore this person because he or she is probably.
not worth my time. I'm just going to like continue with whatever I'm going on." So this.
kind of passive indifference or treating someone coldly is also a form of contempt as well..
And so what Jesus is saying that if you show anger or contempt towards your brother or.
sister, you will face judgment both here on earth in a human court and also before the.
Lord with the fire of hell. And you might ask, "Jesus, why do you take anger and contempt.
so seriously and kind of like put them on par with murder in some cases?" I think this.
is because behind anger and contempt is a posture of disregarding the value and worth.
of a person made in the image of God. And this is exactly the same kind of attitude.
behind murder. When someone doesn't take life of another image very seriously or take.
the life as sacred and worthy to be treated with respect and dignity..
And I want to be clear here. What Jesus is not saying is that anger as an emotion in.
and of itself is a sin. If you just read through the Psalms, you can see all kind of raw emotions,.
right? There's joyful emotions, happy, there's grief, there's lament, and there's also anger..
And in a lot of cases, God welcomes us to pour out our raw emotions before Him, including.
our anger. So what Jesus is not saying is that anger as an emotion in and of itself.
is a sin. And we also see Jesus himself express his anger as well. We see Jesus getting.
furious at the money changers and the merchants at the temple. And he even kind of like pull.
out a whip to drove them out of the temple because they were taking advantage of the.
pilgrims that came all the way to the temple to worship the Lord..
We see Jesus express his anger towards the legalistic way that some of the Pharisees.
practice the law by prioritizing the Sabbath rituals over actually allowing Jesus to heal.
a man that is in need. We also see Jesus became angry at death itself when he stood in front.
of the tomb of his friends Lazarus that has just passed away. So I want to be clear. There.
is a time and a place for anger. And the kind of righteous anger that is expressed in the.
right way with the right intention towards the right thing is not a sin. And in fact,.
if you never felt angry towards any of the injustice or the brokenness in our world around.
us, I might dare to say that you might not be living in according with the Holy Spirit..
However, the anger that Jesus talks about here is a kind of destructive anger that is.
being directed towards another person made in the image of God. This is the kind of anger.
that Jesus is speaking up against. So after laying out his standard of righteousness,.
Jesus then continues to give two applications on the response that he requires from us when.
anger and contempt is present in our community and our relationship with other people..
So verse 23, "Therefore, if you're offering your gift at the altar and there remember.

$^{201}$that your brother or sister has something against you, leave your gift there in front.
of the altar. First, go and be reconciled to them. Then come and offer your gift. Settle.
matters quickly with your adversary who is taking you to court. Do it while you are still.
together on the way, or your adversary may hand you over to the judge, and the judge.
may hand you over to the officer, and you may be thrown into prison. Truly I tell you,.
you will not get out until you have paid the last penny.".
So here's my second point. Jesus demands a better response from us. Jesus demands a better.
response from us. For Jesus, the antidote to the issues of anger and contempt towards.
others is reconciliation. So the better response that Jesus requires from us is to actively.
pursue justice and seek reconciliation. Let's look at the first situation here. "Therefore,.
if you are offering your gift at the altar." So I want you all to notice the present tense.
here. "If you are offering your gift at the altar." The situation that Jesus described.
here involves the Jewish sacrificial system where someone is in the process of bringing.
the gifts to the altar in the Jerusalem temple. And obviously, we don't practice this kind.
of sacrificial system anymore in our church nowadays, right? So I think it would be helpful.
to give a little background to what Jesus is describing here so that it will help us.
better understand what is Jesus asking from us exactly. For the disciples that Jesus is.
speaking to, a lot of them came from the town of Galilee, and they usually only go to the.
temple in Jerusalem once or twice a year during one of the three festivals, the Passover,.
the Feast of Weeks, or the Feast of Tabernacles. And the journey from Galilee to Jerusalem.
can take anywhere from three to four days to one whole week, depending on which route.
they choose. And once you get into the temple in Jerusalem, you still have to wait in a.
long line, perhaps for a couple more days, to be able to finally offer your gifts that.
you brought with you at the altars. So the scenarios that Jesus is trying to paint here.
is that after an exhausting journey to get all the way from Galilee to Jerusalem, and.
then waiting in line for a couple more days, it is finally time for you to offer your gifts.
at the altar. So just think of all the efforts it takes to get to this point. But then Jesus.
continues, "And there, remember that your brother or sister has something against you.".
You might miss this if you don't read it carefully, so I'm going to read it again. It says, "If.
you're offering your gifts at the altar, and there, remember that your brother or sister.
has something against you." Did you guys catch that? In this situation, it is the other party.
that has something against you, and you might very well be innocent at all. But then Jesus.
gives this command to his disciples, "Leave your gift there in front of the altar. First,.
go and be reconciled to them. Then come and offer your gift." So a lot of the time, our.
natural tendency when something goes wrong is to point fingers at other people or to.
shift the blame, right? But for Jesus, it is not about who is at fault, but the greater.
righteousness and the better response that Jesus demands from us is this, that we take.
the initiative to make things right and to seek reconciliation whenever possible. Of course,.
I want to be sensitive here, and I understand there are some exceptions to this, such as.
when someone is just not repentant, or in the case when someone is a victim of abuse,.

$^{241}$where seeking reconciliation with their abuser might not be the wise thing to do. However,.
in most cases, I will say that I agree with Apostle Paul when he emphasized this in Romans..
If it is possible, as far as it depends on you, live at peace with everyone. So for Jesus,.
he understands that the presence of anger and conflict within the community of believers.
can be so destructive that he actually commands his disciples to leave the gift at the altar,.
go back and reconcile with the person, and then come back and offer your gift. Even if.
it means that you have to walk a few more days' journey back to Galilee, and then walk.
a few more days' journey back to Jerusalem again, and then you have to wait in line all.
over again just to offer your gift at the altar. So Jesus might be using some hyperbole.
here, but I think his point is still the same. This is the kind of radical response that.
Jesus demands from us as his followers, a way of living that prioritizes seeking reconciliation.
and pursuing justice within the community of believers, even over our religious rituals.
or our religious worship. And then Jesus continues in verse 25. "Settle matters quickly with.
your adversary who is taking you to the court. Do it while you are still together on the.
way, or your adversary may hand you over to the judge, and the judge may hand you over.
to the officer, and you may be thrown into prison. Truly I tell you, you will not get.
out until you have paid the last penny." So in these second examples, Jesus extends his.
application to those outside of the community of faith, from how we relate to other brother.
and sister to now on how we relate to an enemy or an adversary. While we might not be the.
guilty party in the previous situation, this time we are clearly at fault. But regardless,.
Jesus warns us that reconciliation is an urgent matter, and we need to make peace with our.
enemy before it is too late. If not, we will have to pay the full price for our own actions..
And church, this is the better response that Jesus demands from us, that reconciliation.
and justice is an urgent matter, and it should take priority over our rituals' worships..
Why? Because God cares more about justice and righteousness than religious rituals and.
anti-worship. So I want to end this section by taking a look at Micah 6 to 8. So Micah,.
he is the prophet to the southern kingdom of Judah. And during his time, there was just.
a lot of injustice among the religious leaders and prophets in Judah, from bribery, showing.
favoritism to the wealthy, all the way to stealing from the poor and even taking away.
their land. So this section of Micah 6 is a poem in the form of question and answer..
And right here you will see this person is asking the question. He is desperate to figure.
out what kind of things he should bring when he comes before the Lord and worship Him..
And he is ready to pretty much offer anything that he has to the Lord. So look at verse.
6 together. "With what shall I come before the Lord and bow down before the exalted God?.
Shall I come before Him with burnt offerings, with calves a year's old?" So burnt offerings,.
out of all the different types of offerings in the Old Testament, is the most costly one.
because you pretty much have to burn out the entire animals and there's like nothing left.
for you to be taken home and to be eaten. And the calves a year old is usually also.
considered the best kind of calves and the best kind of animals to use for the burnt.
offerings. And then this person continues, "Will the Lord be pleased with thousands of.

$^{281}$rams or with ten thousands rivers of olive oil?" So now we see he switched from kind.
of like quality to quantity now. Thousands of rams or ten thousands rivers of olive oil..
And it seems like this is still not enough. He then continues, "Shall I offer my firstborn.
for my transgression, the fruit of my body for the sin of my soul?" So now he comes to.
a point where he's willing to offer up even his own son that is the most valuable possession.
that someone can have. But listen to how the Lord responds to him..
"He has shown you, O mortal, what is good. And what does the Lord require of you? To.
act justly, to love mercy, and to walk humbly with your God.".
So God doesn't really care about how expensive or valuable our offering to him is. Or God.
doesn't really care about how big our worship service is or how many people attend our church..
What he cares about is for us to act in a way that is aligned with his heart, that is.
to act justly, to love mercy, and to walk humbly with him. When there's anger and contempt.
or any form of injustice in our community or in our worship, Jesus says that our worship.
is compromised. So here's what I want all of us to take away.
from our passage today. Because Jesus cares about how we relate to people made in his.
own image, we need to actively pursue justice and seek reconciliation. This is the higher.
righteousness that Jesus demands of us, that we treat every human made in his own image.
with the utmost respect and dignity. And this is the better response that Jesus requires.
of us, that we actively pursue justice and seek reconciliation. So once again, because.
Jesus cares about how we relate to people made in his own image, we need to actively.
pursue justice and seek reconciliation. So before I end, I have two application questions.
here that I want all of us to spend some time to reflect on. So as we come to the end of.
the year, is there someone that you need to reach out and seek reconciliation with? Are.
there certain people or group of people that we are tempted to show contempt towards, including.
passive indifference? Ask the Lord for forgiveness and for a transformed heart. So we'll spend.
some time to reflect on these questions and after that, I think Mark will come up and.
lead us in communion..
[MUSIC].
\newpage



\section{但以理書}
\label{sec:FaphUTtEgYU}
\textbf{誰是你主?\#2 伯沙撒王 (但以理書 5) - 陳崇基牧師}
\newline
\newline
連結: \href{https://youtube.com/watch?v=FaphUTtEgYU}{\texttt{https://youtube.com/watch?v=FaphUTtEgYU}} ~~~~ 語音日期: 2024-12-03
\newline
\newline
\hyperref[sec:sCP4BPS1_XA]{\small{< < < PREV SERMON < < <}}
~
\hyperref[sec:index_chronic]{\small{[返順時目]}}
~
\hyperref[sec:index_scriptual]{\small{[返順卷目]}}
~
\hyperref[sec:_TmbaOYl2cg]{\small{> > > NEXT SERMON > > >}}
\newline
\newline
$^{1}$各位觀眾早安.
我們這次繼續我們前一個星期說的一個系列.
我們看看聖經裡面說一些俗世的君王.
或者一些不信主的君王.
一些故事.
我覺得勝經很有趣.
聖經很多時候不是很抽象地告訴我們.
一二三四這麼多點.
第一點第二點第三點.
很少這樣列舉重點出來的.
聖經很多時候是用故事來講述一些教義.
或者用一些故事的方法來述說一些人生裡面需要我們知道的道理.
所以我們看一些俗世君王的時候.
有時候我們會看看他們的故事.
聖經怎麼記載這些敘述.
這些敘述應該怎麼去解釋.
然後對我們今天有什麼意思.
所以今天我們選了另外一個君王.
叫做白沙十王.
這是《但義理書》第五章.
你讓我今年也說《但義理書》.
因為《但義理書》已經成為我今年花了很多時間的一本書信.
我們看《但義理書》第五章.
《但義理書》第五章提到這個叫做白沙十王.
你記得上一次我們提到.
另外一個世俗的君王.
叫做居老時.
下面那位.
黑白那位.
他有個樣子.
在右邊.
他右邊臉的那個.
那個叫居老時.
他是波斯帝國的第一個君王.
他攻陷巴比倫帝國.
我們再把時間回頭一點.
到底巴比倫帝國崩潰之前.
到底發生了什麼事.
當時就是白沙十王.
來做王的時候.

$^{41}$他管理巴比倫地.
如果你記得前一次我們曾經說過我們的步驟應該是怎麼樣.
看這些世俗的君王.
第一個步驟就是看看經外是怎麼看他.
聖經之外是怎麼看他.
然後我們回到聖經裡面.
怎麼看這個君王.
或者說了什麼故事是有關他的.
《但義理書》第五章提到這個白沙十王.
這個君王.
他是巴比倫帝國.
你看到巴比倫帝國黃色那個項.
最後的時候他是做一個末代的君王.
他和他的老父王一起去執政.
去處理巴比倫帝國末代的政務.
如果記得的話.
我曾經提過巴比倫帝國最後的時候.
這位老父王叫做尼布拉登王.
他是一個很古怪的王.
在他555年.
公元前五年.
去到做巴比倫帝王之後.
管了一兩年兩年.
鬆一點.
他就離開了朝廷.
然後他就搬去了一個很遠的地方.
就是今天的沙特阿拉伯北面.
一個綠洲的船.
然後他就不太理會朝政.
然後朝政就交給他的長子叫做白沙薩.
這些全部不是聖經裡面有交代的.
是聖經以外的歷史是有交代的.
所以就去到巴比倫帝國末代的時候.
就有白沙薩王出現了.
如果看一下當時的地圖.
你就會看到左下角.
我們看到這個老父王.
他就搬了去北面沙特阿拉伯的地方.
然後巴比倫帝就留給白沙薩王.
他就管到他主前539年.

$^{81}$看到嗎.
539年發生什麼事.
就是巴比倫帝國綠色那塊.
上面有一個馬代波斯帝國.
當時有一個軍事的天才.
叫做居魯士.
之前我們提過居魯士這個故事.
就是他在主前500多年的時候.
他就去到半變馬代帝國.
然後吞併了馬代帝國.
他就成為馬代波斯的君主.
然後539年你記得嗎.
在聖經的歷史裡面.
世界的歷史裡面.
中東的歷史裡面.
539BC.
是一個很重要的分水嶺.
因為波斯帝國就是靠居魯士.
然後他就入主.
攻陷了巴比倫帝國.
所以這個白沙薩王.
就是《但以理書》第五章裡面提過的白沙薩王.
這個是誰.
是巴比倫帝國.
最後最後的王.
也是老父王尼泊拉登的將子.
白沙薩王.
Belsasha.
我們看看第五章第一節.
怎麼去述說這個故事.
還有這個故事有關所謂的Pagan King.
這個俗世的君王.
有什麼道理想帶出來呢.
第五章第一節.
白沙薩王為他的一千個大臣擺設盛宴.
與這一千人飲酒.
白沙薩在歡飲之間.
吩咐人將他父王尼泊拉登.
從耶路撒冷聖殿所奪掠的金銀器皿拿來.
好使王大臣.

$^{121}$皇后妃嬪用這些器皿來飲酒.
於是他們把聖殿.
就是耶路撒冷上帝殿所擄掠回來的金器皿拿來.
王就與大臣皇后妃嬪就用這些器皿飲酒.
他們飲酒讚美金銀和鐵木石造的神明.
你試一下去想像一下是一個怎麼樣的情況.
在吃大餃子的時候.
就是大家很好吃的.
還有一個大的宴席.
非常大的宴席.
好像一個派對一樣.
可能不夠工具不夠碗碟.
他們就拿了耶路撒冷奪來的金銀器皿.
然後拿聖殿的器皿來飲酒.
這個像是一個派對.
如果你去想像的時候.
就是這個君王.
白沙薩王.
就與他的大臣一起去歡飲的時候.
就是一個這樣的情況.
其實就不只是派對這麼簡單.
因為聖經裡面記載這些所謂派對.
就是吃大餃子.
大家很開心在這裡有一個宴席.
就不像我們今天的那些宴席.
今天我們可能會拿杯酒喝.
或者我們去擺酒等等.
以前在古代的帝國.
每逢擺這些大牌宴席的時候.
很多時候都應該是新年的時候.
或者是春節的時候.
新年的時候.
巴比倫.
特別是巴比倫的文化裡面.
就有一個很大的節日叫做阿基督.
這個阿基督.
其實就是大麥小麥收割的時候.
所以你看到這個照片.
AI幫我弄出來的.
就說到這個阿基督是怎麼樣呢?.

$^{161}$阿基督就是他們會在巴比倫的大神廟.
拿最大的神像.
叫做馬獨克的神像.
然後拿出來就放在車上.
然後就將他巡遊一輪.
然後就大肆的慶祝.
因為這個是全年第一次.
開始有這個大麥小麥的收割.
還有有一個這樣的阿基督.
一個這樣的慶典.
其實是想做什麼呢?.
除了是一個農產的.
豐收的一個節日.
就是一個有宗教和有政治味道的一個慶典.
就是那些的君王和大神.
他的內閣.
所有的政府的人員.
就會一起去有一個歡飲.
或者一個大的盛宴.
因為剛剛收割了.
所以打了大麥小麥.
然後他們做了很多東西出來.
然後大家就在享受.
那個目的是什麼呢?.
那個目的呢.
這些大的歡飲或者盛宴.
其實這些的君王或者地上的帝王.
他們就是想在盛宴裡面.
去拿到神明給他的批准.
神明要給到他的准許.
或者批准他喜歡他做的事情.
為什麼在新年裡面這麼重要呢?.
因為在巴比倫帝國裡面.
那個宗教政治的氣氛是這樣的.
每逢新年的時候.
就是要去量度一下.
接下來一年到底有什麼大計.
接下來一年有什麼政策.
政策他們需要去執行.
或者他們的內閣在想什麼.

$^{201}$我們需要在這一年裡面去執行出來.
他們就希望在一年的新年開始的時候.
來一個大的盛宴.
那個盛宴就會去讚美或者敬拜.
所有帝國裡面的神明.
包括他們巴比倫最大的神明.
或者小一點的神明.
月神啊什麼神啊什麼神啊.
所以包括了耶和華.
因為你記得這個時候.
就是巴比倫帝國統治天下的時候.
第一聖殿是被毀的.
利布格尼撒這裡提過.
所謂利布格尼撒.
不是真的伯沙撒的老爸.
是上面的前任的君王.
利布格尼撒就是幾十年前.
去到耶路撒冷.
然後就攻陷了耶路撒冷的聖殿.
然後把聖殿裡面的所有的工具.
拿到了巴比倫的京城.
所以這些大的盛宴.
其實他們就拿了這些東西出來.
來去慶祝.
不單是慶祝.
也是希望能夠在接下來這一年.
我們所做的事情.
是一個神明或者神祇.
可以去給他們一些的Blessing.
有點像我們華人的文化.
每逢新年的時候.
有些風水佬或者占卜等等.
就會說今年的運程是怎樣.
今年的是不是犯太歲.
如果犯太歲你要做些什麼等等.
這個故事大約是這樣的情況.
所以他們是希望.
這些神明給他們Blessing.
跟他們說話.
然後最重要是給他們Rubber Stamp.

$^{241}$有一個Approval.
所以在這種盛宴裡面.
他們是預期神明會發聲.
預期神明是降旨告訴他們.
Yes, I approve你做的事.
Yes, Go ahead.
大概這樣.
所以在這個故事裡面.
第五節.
這裡說到真的上帝出現了.
五章第五節.
當時他們在大排筵席的時候.
忽然有人的指頭就出現了.
在登台對面皇宮粉刷的牆上寫字.
王漢見寫字的指頭就變了臉色.
心意驚惶.
腰骨好脹脫節.
相實彼此相碰.
好奇怪.
嚇得很慌張.
所以這些都不是驚訝的事.
神明會出聲.
但他害怕的是什麼呢.
居然在牆上寫字.
這個很奇怪.
你看到這裡的名畫.
Wembrandt的一幅名畫.
十七世紀的時候.
講到Belshazzar's Feast.
一幅很出名的名畫.
好像在倫敦可以看到.
講到這個皇帝望一望牆上.
就會發現一些亞蘭文的字母出現.
這個就是英文.
我們有一個術語叫做Writing on the Wall.
牆上的字.
這個術語的來源.
所以這裡是在講大排筵席.
他們希望神明給他們一個祝福.
然後居然神明出現的時候.

$^{281}$耶和華的手出現的時候.
居然在牆上寫字.
所以害怕的地方就是這樣.
居然用一個他們不會預期的方法出現.
當然他們是看得到或讀得到這些字的.
在夏文裡面我們知道.
其實寫的文字是亞蘭文.
亞蘭文是當時的國際語言.
如果你牆上看到有英文字.
你懂得看嗎.
你應該懂得看吧.
但竅妙的地方就是這樣.
就是不懂得解釋.
不懂得解釋.
你牆上看到有英文字.
突然可能新年的時候.
有英文字.
不說是什麼了.
免得你有什麼疑夢.
總之你懂得看每一個英文字.
但你不知道怎麼解釋.
跟我有什麼關係.
跟那個帝王和帝國有什麼關係.
沒有人可以解得通.
問題就不是不懂得讀.
而是不懂得解釋.
所以搞了一段時間之後.
沒有人知道怎麼解釋的時候.
就傳召了但以理進宮廷.
他進去白宮的時候就叫但以理.
你可不可以解釋牆上的文字呢.
結果就真的叫了但以理去.
但以理就跟他說了一番話.
他說你白沙薩王.
至高的上帝曾經將國度大權榮耀威嚴.
賜給你的父利布格尼撒.
就是巴比倫帝國最大的王.
之前有的.
白沙薩你是他的兒子.
是他的後代.

$^{321}$你雖然知道這一切卻不謙卑自己.
竟向天上的主自告.
你又讚美那不能聽沒知沒識的.
用金銀銅鐵木石造的神明.
沒有將榮耀歸給那手中掌管你氣息.
管理你一切行動的上帝.
五章第二十四節.
於是從他那裡顯出指頭.
寫這文字.
所寫的文字就是.
尼尼提克勒烏法爾辛.
就算你看阿南文也好.
你看中文翻譯出來的音譯.
你也看得明白尼尼.
寫什麼字?尼尼.
尼尼兩字.
然後提克勒烏法爾辛.
是不是很容易解釋?.
就是這樣的情況.
你知道那個字.
你看到那個字母.
阿南文的字母.
就算中文翻譯出來的譯音.
但你不知道怎麼解釋.
所以你要有解釋.
你要知道那個解釋.
其實解釋是怎麼樣呢?.
但你就是說尼尼.
首先說尼尼.
尼尼.
尼尼其實是數算的意思.
如果你用現金找續的時候.
你就是數多少錢.
一二三塊錢.
這個就叫尼尼.
然後呢.
提克勒.
什麼意思?.
就是天秤.
秤東西.

$^{361}$今天我們不用那些秤東西.
以前你記得我們很.
以前那一代.
我們去買中藥的時候.
弄一兩.
不知道那些什麼.
不知道那些奇怪的東西.
然後他就慢慢慢慢這樣.
逐顆逐顆這樣調出來的.
那種天秤.
就叫做Tecco.
就是提克勒的意思.
烏法爾辛.
烏其實就是and的意思.
就是和的意思.
烏法爾辛.
烏帕辛.
其實就是divided.
就是將分割或者分裂的意思.
那你就開始明白了.
到底整句話或者牆上的文字.
到底是什麼意思?.
那個interpretation是什麼?.
那個interpretation.
就是說.
這個帝國將會很快完結.
利利.
但是你說.
解釋interpretation就是.
利利就是上帝已經數算.
計算.
你國的日子.
到此完畢.
提克勒.
秤.
記不記得?.
天秤.
就是上帝已經將你評在天秤上面.
評出你的規矩.
所以你見紅.

$^{401}$已經是差了.
比勒斯或者剛才我們那個法爾辛.
其實是眾數單數的意思.
所以音有一點點不同.
就是你的國度要什麼?.
分裂.
歸給馬代和波斯人.
但是你一這樣解釋的時候.
大家就明白了.
我們看了那些什麼利利利利.
法爾辛那些字.
現在我們知道什麼意思.
還有怎樣去應用.
但你就應用在白沙實王身上.
說他的國度上帝已經在數算他的日子.
快完畢.
還有他的國度是不公義的國度.
因為在上帝的項目裡面.
上帝你知道上帝是會計師嗎?.
上帝在計算你的數字.
看你平不平衡.
巴比倫帝國.
白沙實王的帝國.
就是不平衡.
然後你的帝國要分裂.
歸給另外一個外族的馬代和波斯人.
記得那個地圖嗎?.
其實北面的時候已經是重重欲動.
有其他的軍隊已經臨近的時候.
其實可以這樣解釋.
正因為有其他的外敵開始醞釀威脅的時候.
所以為什麼他們才會有大排筵席.
因為有大排筵席是希望這些神明來幫助他.
還有去批准他做的事.
或者幫他的國防等等.
無論如何都好.
這幾個字.
其實就是說上帝已經給你一個evaluation.
給你一個帝國的evaluation.
就知道你的帝國將會滅亡.

$^{441}$正因為這個緣故.
但以理已經解釋了.
白沙實王下令.
就有人把紙袍給但以理穿上.
把金鏈戴在他的頸上.
又傳令使他在國中位列第三.
這裡先解釋一下.
解釋了牆上的字.
但問題就是這樣.
他不知道他的國度原來會消失得這麼快.
還在至高的時候.
覺得還可以屙屙屙屙的.
還可以等的時候.
這就是當時他的心態.
所以把紙袍放在但以理的身上.
把金鏈放在他的頸上.
又傳令使他成為國中第三.
為什麼第三.
因為有個老父王是第一.
叫布拉登.
然後白沙實王的長子是第二.
所以但以理就變成第三.
這個很合理的.
在歷史的背景裡面.
五章第三十節.
當夜.
當夜發生什麼事.
加拿大王白沙實就被殺.
馬代人大劉氏年六十二歲就娶了加拿大王.
加拿大國.
這個大劉氏大概應該是另外一個馬代的綽號.
是給居老時的.
Darius.
這個晚上.
好像一夜之間.
你的帝國就玩完.
一夜之間你的生命就結束.
甚至現在我們可以.
又是聖經以外的資料告訴我們.
可以計算到這一天是哪一天.

$^{481}$是主前539年.
10月12號.
October 12.
10月12號.
有誰是12月12號生日的.
就是他國度玩完那天.
不過是主前539年.
整個巴比倫帝國就好像骨牌一樣.
或者House of Cards.
然後就整個帝國沒了.
我開頭的時候就說.
聖經很多時候不會用一種西方的方法.
列表.
Bullet point.
第一點第二點第三點這樣讓你想一下.
聖經是用故事.
Tell you a story.
然後給你一些反思一些反省.
這裡的反省.
很清楚有什麼可以拿走.
整個故事想告訴我們.
上帝是一個怎樣的上帝.
記住.
看聖經的時候.
我們不要只是看一個故事.
就問有什麼要學.
和有什麼不好學.
有什麼good example.
有什麼bad example.
我們就看到這裡有一個bad example.
但你有一個good example.
我們就學它.
不可以學的.
因為他們的時代和我們的時代不同.
最主要那個故事想說什麼.
聖經用一個故事告訴你.
神是怎樣的.
去介紹上帝是怎樣的.
等於你認識一個人的時候.
你問什麼.

$^{521}$你問他的由來.
他從哪裡來.
你認識什麼.
然後他做過什麼等等.
You tell them the story.
你認識一個人.
就要知道他的故事.
你認識上帝就要知道上帝的故事.
上帝在歷史裡面做過什麼.
那個故事告訴你.
神是怎樣的.
這裡很清楚看見.
上帝的主權.
上帝的王權.
是高過地上一切的帝國.
不只是古代的巴比倫帝國.
我不論你是美國.
中國.
俄羅斯.
任何一個國家.
任何地上所有的kingdom.
They come and go.
你的會走.
興的也會衰.
聖經就是很坦白的告訴我們.
有時候我們未必會感覺到.
我們會覺得這個國度會繼續下去.
我們可能會覺得這個朝政會繼續下去.
如果這個朝政繼續下去就好了.
我們未必會有這樣的感覺.
但是聖經就是告訴我們.
所有的政權.
所有的政府.
沒有上帝的准許的話.
他們都會完蛋.
並且好像是一夜之間就會過去.
這個是一個故事.
說上帝的手在牆上寫字.
在《大人與理書》裡面經常出現上帝的手.
有好幾次.

$^{561}$說到上帝的手.
第五章這裡當然有說.
但以你好像先知那樣.
他說白沙薩王去拜這些神明.
假的神明.
沒有將榮耀歸給那手中.
去掌管你氣息的.
管理你一切行動的上帝.
所以上帝就寫字.
去告訴你.
你的日子在倒數.
第二章《大人與理書》.
開頭的時候.
你記得那個故事.
我們從小到大在中學學到的.
利布格尼薩做了一個大的偶像.
然後有金銀和鐵石.
四部分.
又是說四個帝國.
然後這個人像的意象裡面.
最後發生什麼事.
第二章第45節.
他說你既看見.
非人手作出來的一塊石頭.
從山出來.
打碎整個的偶像.
就是至大的神.
要將將要發生的事指明和.
解釋是準確的.
非人手所作出來的一塊石.
是什麼 自然的石.
即是人造的東西.
神像就要人手去釘.
人手去釘出來.
小的偶像也好.
大的偶像也好.
都是要人手用工具去作出來的.
這個非人手所作出來的石.
就代表上帝.
上帝的手.

$^{601}$上帝的手才做到一個自然的石.
他會粉碎金銀和鐵石.
第二章第45節.
然後第八章第25節.
八章第25節.
長話短說就是在說.
波斯帝國 希臘帝國.
希臘帝國末期的時候.
有一個奸臣 一個暴君出來.
他用權術和鬼咒來得以成功.
他的心也是自高自大.
趁無人避的時候.
他就毀滅多人.
他做了很多很多很陰側的事.
一位暴君.
他就起來去攻擊萬軍之軍.
就是上帝.
置中他非人人的手.
而遭受毀滅.
非人人的手.
就是代表上帝的手.
幾次提到人的手.
或者手這個字.
就是代表上帝去掌管所有的權柄.
你不要把地上的權柄去到粉飾得太厲害.
等於你明白.
地上的權柄你多麼喜歡也好.
他給你多麼成功也好.
甚至給你市不斷飆升.
你不要把它洗禮.
那個不是上帝的手.
要上帝去到批准.
上帝的准許.
這些地上的權柄才能夠繼續.
如果上帝不批准.
這些權柄一夜之間就可以還原.
不要放錯你的信任.
你要放信任的話.
你就要放在上帝自己的國.
因為大義利書一至到第十二章.

$^{641}$就告訴我們只有上帝的國度才是永恆.
地上的國度.
甚至你的國度.
都會一時崩潰.
所以你以誰為主.
第二方面也要講到上帝的旨意.
是高超過人的旨意.
人最好都是謙卑.
我們時時刻刻想起.
不要太自高自大.
我們很多時候都想上帝去批准我們做的事.
我們想好了一堆事.
然後我們都沒有過問過上帝.
甚至有時候把上帝放為我們人生的內容的一個註腳.
我們想上帝去遏制我們做的事.
甚至我們有一個試探.
我們甚至會濫用上帝的名字.
我認識一個朋友.
一個基督徒.
他經常都說上帝感動他做這個做那個.
上帝可不可以直接叫你做某樣東西.
上帝當然有能力可以做.
問題是你到底洞察上帝的心意.
還是你自己強行你的心意.
那個朋友經常都說聖靈叫他做某樣東西.
聖靈又做某樣東西.
但跟他坐下來一起查經的時候.
查經是亂查的.
我心想你聖經都沒有查好.
你怎會明白上帝的心意呢.
邏輯很簡單.
上帝可不可以直接叫你做某樣東西.
當然可以.
但你都需要去明白去判斷.
你連查一個很簡單的經文都查錯的時候.
你怎能夠明白上帝的心意呢.
所以第一件事就是回去做功課.
如何明白上帝是怎樣的.
我們要好好去解釋聖經.
好好去明白聖經.

$^{681}$好讓我們在聖經中得到人生的智慧.
不要只想著上帝一定要填補我的計劃.
上帝一定要填補我的投資或計劃等等.
如果我們的心不是屬於上帝的話.
如果我們的心不是以上帝為我們真正的主的話.
唯一的主的話.
上帝是不會不遵守我們所做的事的.
他只會准許你做.
不過可能一夜之間你自己的國王都可以崩潰.
亞哥書第四章.
講出什麼智慧呢.
跟大義律師差不多.
其實明天如何你們還不知道.
生命是什麼.
你們原始一片的雲霧.
出現片時就不見了.
什麼叫雲霧.
雲霧就是你那一口氣.
the mist.
現在你在呼吸.
就是那一口氣.
你的生命就是那一口氣.
你原始是一口氣.
你們倒應當說.
主若願意我們就能活著.
也可以作者事或作那事.
只有這樣.
我們認定上帝是我們生命的主的時候.
我們才能走得正.
我們的生活有沒有以耶穌為我們真正的主.
你如何用你的時間和金錢.
如何用你的資源.
如何看待你的資源計劃等等.
有沒有首先去認定上帝.
耶穌是我真正的主.
對嗎.
年尾的時候.
我自己也做了一些特別的奉獻.
有需要的人或有需要的地方.
我都會嘗試去慷慨.

$^{721}$雖然錢包緊一點.
但我都嘗試這樣做.
generosity.
我就給了一些奉獻出來.
我自己也沒有問有什麼回報.
因為奉獻就是奉獻.
對嗎.
很奇怪上個星期.
總會有一個職員打電話給我.
他說什麼時間在教會.
我想和你相約見一見面.
我說什麼事呢.
要見面.
因為平時都不找我.
有什麼同人生慶.
無端端找我有什麼事呢.
我想問什麼事呢.
是否想和我算帳.
還是有些事情沒有付錢.
坐下來.
你記得那個韓國人Johnny.
坐下來.
來到教會坐下來.
我說什麼事呢 你找我.
我都很忙.
不能找你.
你找我.
他就說沒什麼.
我們總會結算帳的時候.
發現我們欠你錢.
他就拿信封給我.
你欠我錢.
他說是啊.
我們早年的時候不是支持你的嗎.
那張支票沒有過數.
我心裡想算你誠實.
沒有吞了我的錢.
我回去.
我開信封是一張支票.
正是我去到韓國奉獻的那些錢.

$^{761}$有時候我們在生命裡面.
真的要先認定上帝是我們的主.
對嗎.
上帝是我金錢的主.
是我時間的主.
我們不要經常覺得上帝阻攔我們.
又要我去聚會.
又要我去敬拜.
又要我去做這樣做那樣.
阻礙我發財.
阻礙我走運.
為什麼上帝經常要阻礙地球轉.
不.
他說一切的資源上.
要認定上帝是我們的主.
這是首要.
是我們信託方在的地方.
上帝自己.
不是地上的國度.
不是我們自己的計劃.
今天我們結束的時候.
我們用聖餐結束.
好嗎.
聖餐歷史歷代裡面.
宣告一件事.
就是我們的生命.
是屬於主耶穌的.
他以寵價來救贖我們.
這個餅的杯代表上帝.
藉著他的兒子 神的兒子.
來救贖我們 買贖我們.
Redeem our life.
不是等閒事.
認定上帝是我們的主.
好嗎.
聖餐是為所有基督徒而設的.
如果你很清楚.
信主是什麼意思.
很清楚耶穌.
我已經接受了耶穌.

$^{801}$不過今天我想給你一個更新的機會.
想給你一個新的機會.
因為上帝說領聖餐的時候.
耶穌你是我的主.
他幫助我心靈裡面每一方面.
認定你.
做我的真正的君王.
我們誠心帶去祈禱.
預備我們的心去領受聖餐.
上主我們多謝你.
因為你是我們心靈的主.
我們本身就是一個.
不值得你去到.
以恩典來救贖的生命.
但這個恩典 這份愛.
就是來到我們的人生.
最後多謝你愛我們.
多謝你去到.
在我們人生的旅途裡面追趕我們.
以致我們能夠回到他們身邊.
多謝你十字架上為我們所作.
去救贖我們.
主讓我們心靈有著.
屬靈的尊貴.
在你的眼中.
我們是多麼的寶貝.
主讓我們心靈以你.
為我們真正的主.
真正的王.
因為心靈裡面每一個的抉擇.
我們先去到.
搞清楚我們的動機.
搞清楚我們的取態.
以你的國度 以你的公義.
以你的道為至上.
求主你幫助.
讓我們領受這個聖餐的時候.
知道你的恩典是多麼的豐厚.
也讓我們的生命能夠成為.
你恩典的出口.

$^{841}$去祝福其他人.
讓我們的個人的關係.
我們的生命裡面工作等等.
能夠多一些被人看到.
耶穌基督你是拯救了我們這個生命.
是有所不同的.
我們這期節目也很奉主名的.
\newpage



\section{}
\label{sec:_TmbaOYl2cg}
\textbf{Who is Your Lord \#3 Augustus (Luke 2) - Rev. Mark Chan}
\newline
\newline
連結: \href{https://youtube.com/watch?v=-TmbaOYl2cg}{\texttt{https://youtube.com/watch?v=-TmbaOYl2cg}} ~~~~ 語音日期: 2024-12-11
\newline
\newline
\hyperref[sec:FaphUTtEgYU]{\small{< < < PREV SERMON < < <}}
~
\hyperref[sec:index_chronic]{\small{[返順時目]}}
~
\hyperref[sec:index_scriptual]{\small{[返順卷目]}}
~
\hyperref[sec:WhzVz61wqtY]{\small{> > > NEXT SERMON > > >}}
\newline
\newline
$^{1}$- All right, guys, this is the Christmas season,.
and we're going to talk about the Christmas story,.
of course, and we're gonna turn to the Bible.
and continue, actually, with this mini-series.
on the pagan canes in the Bible..
And I'd love to see all of you,.
and I know some of you are missing today.
because of some difficulties,.
and we're gonna pray for one another.
at the end of the service, all right, after the sermon..
So I wanna set aside some time at the end of this service.
to remember some of those who are weak.
and having challenges..
You know, one of them is Jinfan..
Jinfan is still in the hospital, I think..
Hopefully, she will be discharged by the end of today,.
and I'll give you more updates at the end,.
and actually, it's good to see you,.
once you're here, and good to see you..
I know it's a difficult day yesterday..
Jim Borders also, you might, yeah,.
they went to UR, and so it's a very difficult day.
for some of us..
And remember to pray for one another,.
and we'll do that after the sermon..
And again, we turn to the Christmas story,.
which is found in Matthew and Luke,.
Matthew 1 and 2 and Luke 1 and 2,.
and today, we're going to focus on Luke 2,.
and the first few verses talk about.
the context and the time when Jesus was born..
2.1, in those days, Caesar Augustus issued a decree.
that a census should be taken of the entire Roman world..
And this was the first census,.
or the census before something took place,.
and what that took place was that.
Quirinius was governor of Syria..
And verse 3, everyone went to his own town to register..
Now, this is one of the Christmas stories.
recorded in the Bible about the birth of Jesus..

$^{41}$As I said, the first two chapters of the Gospel of Matthew,.
the first two chapters of the Gospel of Luke,.
which is this passage, it's all about,.
and the Gospel of Luke did not miss the background,.
you know, the history, the time of Jesus' birth..
There are two pagan kings mentioned here..
The first one is named Augustus,.
the second one is Quirinius,.
and we're gonna focus on Augustus..
Augustus was the founder and the first emperor.
of the Roman Empire, and he was a very powerful.
Roman king and a political elite..
So he rose to power after a series of military conflicts.
and a series of political successes,.
and under Augustus, the Roman Empire expanded greatly,.
and he built many different buildings,.
and he kept the peace, what we call later day.
Pax Romana, Roman peace,.
which means he beat up all the guys,.
you know, he conquered all the lands,.
and then there will be peace..
So he kept the peace with his iron fist,.
and he conquered many, many places,.
including the land of Israel, Palestine,.
and he imposed heavy taxes, heavy taxes..
You think your tax is heavy?.
Back in the Roman world, their tax was quite heavy,.
especially for those who are common people..
And Augustus even called himself the son of God..
From the coins, from the statue,.
from all the inscriptions, you know,.
all these promotions and propaganda,.
talking about him as the divine son of God..
And he demanded people to worship him.
and to pledge allegiance to him..
Now, as we see these three verses,.
notice the words, issue a decree, a decree..
It was a political order..
Now, only the Roman emperor can do that,.
can issue a decree, a demand on the people to obey..

$^{81}$And one of the contents of the decree is taking a census..
Now, census in those days in the Roman world,.
it's not simply counting the heads..
It's not simply counting number of people..
It's not just for statistics purpose..
The Roman empire would call for census.
every some 10 years or maybe 14 years..
Every 10 or 14 years, they would call for,.
you know, regional or empire-wide census..
And the purpose of a census is to control,.
is to exercise control of the conquered people..
So in the census, when the time comes,.
you have to report in your hometown to the government..
Now, people would often be upset..
And that's why in times of census,.
you know, people often will plan for revolt.
against the empire..
And you also notice not only the word decree.
and also the word census, also Quirinius..
Quirinius was a governor of the province of Syria,.
which includes the land of Palestine, Israel..
And the word governor is where we get.
the English word hegemony..
Hegemony, it means the most powerful guy, you know..
He's the most powerful guy in the region, in the province..
And of course, the top guy is, this guy here is Augustus..
Now, as you can tell from this three verses,.
Luke is trying to tell us when Jesus was born,.
it was a time of mighty political power,.
exercising their authority over poor people,.
the common people..
And Christmas happened in such a time..
So don't you forget, in the time of Jesus' birth,.
it's not what we experience in Christmas..
You know, usually when Christmas comes,.
what are we thinking about?.
We think about vacation, right?.
You got some vacation to burn, right?.
So you wanna do it by the end of the year..
In times of Christmas, you think about shopping,.

$^{121}$buying gifts, you know, Christmas lights,.
and all those things..
In the time of Jesus' birth, it was not like that..
It was a time of domination..
It was a time of a political crisis.
or political instability..
And this three verses, you know, trying to tell us.
or to give us the context, the time of Jesus' birth..
Verse four, "So Joseph," now, as you can see,.
the camera is trying to zoom in.
to the main character of the story..
"So Joseph also went up from the town of Nazareth,.
"in Galilee to Judea, to Bethlehem, the town of David,.
"because he belonged to the house and line of David..
"He went there to register with Mary,.
"and Mary was pledged to be married to him.
"and was expecting a child.".
I think Mary's carrying at least nine months now..
"While they were there, Joseph and Mary,.
"the time came for the baby, Jesus, to be born..
"And she gave birth to her firstborn, a son..
"She wrapped him up in cloths and placed him in a manger,.
"because there was no room for them in the inn.".
Now, there's so many details here we can unpack,.
but here, the big picture is the common people,.
the conquered people, the losers,.
Joseph and Mary included, and Jesus, of course, included,.
and all the Jewish people, they report to their hometowns.
and to recognize that Augustus was ruling over them..
Augustus was the king..
Nobody else is the king..
And Mary and Joseph were on their way from Galilee,.
which is in the north, to Bethlehem..
As you can see from the map,.
right there, far corner on the right,.
kind of upper right, you see Nazareth?.
Can you see that?.
You can make it out, right?.
That is in Galilee..
Galilee, you see the small lake on the right.

$^{161}$is the Sea of Galilee, a freshwater lake..
And they have to make the journey to the south,.
which is Jerusalem and Bethlehem..
Now, Bethlehem is located a few kilometer south of Jerusalem.
so they have to make the way all the way from the north.
to the south, and that is, the red line there,.
is 80-mile journey, and mostly, I think, on foot..
Mostly on foot..
You can imagine if you are pregnant with a baby,.
eight months, nine months, and you have to walk, right?.
Now, most of the image that we have for Christmas.
is that they got a donkey, right?.
They got a donkey..
They kind of like, they have a Uber, you know?.
So Mary was on the Uber, was on the donkey..
Well, not necessarily so, because you got to pay.
for the Uber, right?.
You can't get the free Uber..
Not necessarily so..
I think the 80, 90-mile journey, mostly, was on foot..
80 miles is like three marathons, right?.
Carrying a baby is a very, very difficult journey..
Chapter two, verse four says, "Went up to Bethlehem.".
And that makes sense geographically.
because you're climbing from the north shore.
of the Dead Sea, which is right at the bottom center,.
and you have to climb almost 2,000 feet,.
and over 2,000 feet, to the central mountain range.
where Jerusalem and Bethlehem were located..
And as you can tell, it's a very difficult journey,.
especially for a pregnant woman..
And when they get there, it was time to give birth..
And Joseph, I think probably, at the time of delivery,.
he had to urgently hire a midwife to help out..
And there was no, you know, Rush Hospital.
or Northwestern Hospital in those days..
There was no delivery room..
And the scripture tells us, they went into the barn,.
and there was no place for the baby to lay but a manger..
However, in Luke one and two, we find people like this..

$^{201}$They were able to meet Jesus as their Savior..
As you can tell from verse one,.
the main character seems to be Augustus,.
seems to be those who hold power,.
those who have political might,.
those who have military might,.
and they were the main actors in history..
They were the main actors on stage,.
on the stage of human history..
Augustus, Quirinius, the governor, the king, the emperor,.
the political elites, all those people.
are the main actors in human story, right?.
So we study Western history or Eastern history.
or whatever history..
Sometimes we focus so much on the big guys, right?.
The kings, the emperor, the generals, right?.
Those who can fight, those who can win,.
and those who can conquer..
But in Luke, in the Christmas story,.
the concentration or the focus is on the common people,.
on Joseph and Mary..
Now, we don't have time to take a lot of details today,.
but in Luke 1 and 2, the first two chapters,.
talking about Christmas or the arrival of Jesus,.
the Son of God, the Savior,.
the Gospel of Luke mentions so many different people.
who have met the baby Jesus,.
who have been part of the Savior's story..
From chapter two, we see Joseph and Mary,.
and Mary was probably, you know how old Mary was?.
Probably a high school girl, right?.
14 year old, 15 year old,.
and has the faith to receive Jesus..
And then chapter one,.
the parents of John the Baptist,.
Zechariah and Elizabeth,.
and these two, they're the couple,.
were aging and dying couple because they were so old,.
and the Gospel of Luke tells us.
they had no child..

$^{241}$They're struggling with infertility,.
and in those days, it is a shame to not have a child..
In those days, it was hopeless.
if you cannot have a child to carry on your name,.
to continue your line and your family..
So Zechariah and Elizabeth,.
they were dying and aging couple,.
and they had no child..
However, in chapter one,.
God graced them and granted them.
a miracle child, John the Baptist..
And John the Baptist is the pioneer of the Savior Jesus..
That is Luke chapter one..
And then the story goes on and talks about Jesus.
when after he was born, when he was eight days old,.
he was brought to the temple in Jerusalem for offering..
So Joseph and Mary are carrying an eight-day-old child,.
Jesus, in their arms,.
and bring him to the temple in Jerusalem.
for circumcision and for offering a sacrifice..
And there in the temple court,.
two other people are being mentioned,.
and they have the opportunity to take part.
and to experience, to meet the baby Jesus..
And the first one is Simeon..
Simeon, the Bible did not tell us how old he was..
I think probably he was in his last days of his life,.
and he was devoted to worship God day and night.
in the temple, and finally he met Jesus, the Savior..
And then the second person is Anna..
Anna was an old widow..
The Gospel of Luke tells us she lived in the temple.
day after day with the hope of seeing the salvation happen.
in her life..
And there Anna also met Jesus in the temple court..
Now these people, Joseph and Mary, Zechariah and Elizabeth,.
and their son, miracle child, John the Baptist,.
pioneer of Jesus, and Simeon, Anna,.
they are all powerless people,.
but in the plan of God, in the story of God,.

$^{281}$they are the main characters,.
because they praise God and thank God and worship God,.
because finally they can see the Savior.
coming into the world, coming into their life..
And they welcome Jesus, the Savior,.
with thanksgiving and with faith..
Now if you have time, I would always recommend.
in times of Christmas, go into the Bible,.
spend some time, and we read and meditate.
and reflect on the story of Christmas..
Now if you have time, look at Luke 1 and 2..
You see, Luke is trying to tell the story.
of some of these common people,.
and afterwards, he always record their praise and worship..
Anna praised God, Simeon praised God.
after holding and seeing the Savior, baby Jesus..
Zechariah and Elizabeth, they burst out.
with praise and worship, because they finally have met.
the hope of Israel, the hope of mankind,.
the salvation in Jesus..
And Joseph and Mary, of course Mary,.
offer a song called "Magnificat.".
And that is a song praising God for reversing our role,.
reversing, you know, the powerful becomes powerless,.
and the powerless become powerful in the sight of God..
So all of these people, they are common,.
even though they are conquered,.
they are in the low, low bottom layer of the society,.
and they welcome the Savior,.
and they have met the Savior face to face.
with thanksgiving and in worship..
Now this is relevant to us,.
because none of us is like Augustus..
None of us is like billionaire..
Maybe you got a lot of money, okay?.
I guess you are not Elon Musk, right?.
You're not in the White House..
You're not the powerful elite..
But this story tells us Jesus came to the life of people.
who are powerless, who are poor, and who feel oppressed..

$^{321}$Now Jesus is different from Augustus,.
because his kingdom is different from the Roman Empire..
Jesus did not come to collect power for a field..
Jesus is going to rule differently..
And as I told you, all the praise songs,.
especially Mary's praise song, "Magnificat,".
talks about God in Jesus elevates those.
who are without privilege, without status,.
elevate them from low to high,.
and release people from bondage..
And that is the vision of Jesus' kingdom,.
is to restore relationship with God.
and restore our brokenness.
so that he can breathe life, eternal life in our lives..
Augustus would ask other people to sacrifice themselves.
for him and for his power, for his kingdom, for his empire..
But Jesus came and sacrificed himself for us,.
sinners, broken people..
We don't deserve it..
But Jesus came and sacrificed himself for us..
And our situation may be different from 2000 years ago,.
but we still see kings these days,.
politicians these days,.
flexing their muscle for self-interest..
And these people, the powerful,.
would benefit themselves and the privileged few..
Jesus is calling us not to serve ourselves,.
but to serve others..
So that's what he did..
The ruling powers may be strong.
and seek to bend us to their vision,.
but Jesus is calling you, calling me, to his kingdom..
And his kingdom is totally different.
from the earthly empire..
His kingdom is characterized by service and humility,.
service and humility..
As we think about the birth of Jesus in this holiday season,.
don't you forget,.
don't you forget Jesus is calling us for his kingdom..
And his kingdom is totally different..

$^{361}$And remind ourselves not to seek our own good,.
but to give and to serve other people..
In Luke chapter 22,.
almost before Jesus went to the cross and being crucified,.
he said to his disciples,.
because they're on their way to Jerusalem,.
and Jesus knew that he was gonna die.
a shameful death on the cross,.
and he was trying to interact with his disciples,.
and his disciples trying to fight for,.
you know, who's gonna be first in Jesus' kingdom..
And Jesus is trying to tell them and to teach them,.
this is a totally different way..
It's a way of reversal..
Jesus said to the disciples,.
"The kings of the Gentiles, lord it over them,.
"and those who exercise authority over them,.
"call themselves benefactors.".
And the kings, the emperor, the powerful people,.
would call themselves to be benefactors,.
because they are the ones who do good to other people, right?.
To have charity to other people,.
to shower blessing to other people..
They are the actors, the benefactors..
These are the Gentile kings..
Jesus said in verse 26 and following,.
he said, "You guys, the disciples,.
"the church, are not to be like that..
"Instead, the greatest among you shall be like the youngest,.
"and the one who rules like the one who serves.".
Who is more important in the kingdom of heaven?.
The one who sits at the table, or the one who serves?.
Who is more important, Jesus said?.
The one who sits at the table, of course, but not here..
You know, in the Roman Empire,.
the one who sits at the table, who has the feast,.
they are the powerful..
But in the kingdom of heaven, in the kingdom of Jesus,.
it's not like that..
Jesus said, "For I am among you as one who serves.".

$^{401}$Jesus came to serve us, not to be served..
And that's why Jesus, in the sight of God,.
is the real Son of the Most High,.
the real Son of God, not Augustus..
And Jesus bring us true peace,.
not the Roman peace, not the earthly peace,.
but the heavenly peace, eternal peace,.
because our relationship with God can be restored.
because of what Jesus did for us..
And Jesus comes to serve all of us..
Those who are on top of the pyramid,.
Augustus, the earthly kings,.
those who have the political power,.
they come to be served..
But Jesus said, "In my kingdom,.
"the order of things is totally different.".
So as we reflect on the Christmas season,.
ask ourselves a question..
What can we do to serve God and others?.
What can we do not to be served, but to serve others?.
As I told you at the beginning of this sermon,.
we're gonna serve one another,.
we're gonna pray for one another..
So let's start with this small action,.
to remember other people, all right?.
So last night, I received a message from Jin Fan,.
and I was kind of taken aback.
because she told me she was in the hospital.
because she was having some high blood pressure,.
beyond normal..
So Paul and Jin Fan rushed to Northwestern Hospital,.
and I just received a message this morning..
As you know, Jin Fan is carrying 34 weeks, right?.
34 weeks, yeah, it's very critical..
And, but I received a message from her this morning,.
her blood pressure is back to normal..
Hopefully, we pray for her, pray for Paul,.
pray for the baby..
Everything will go fine,.
and she will be discharged by the end of today, yeah..

$^{441}$But we need to pray for them, pray for them..
And then, and I also received another news,.
is Wang Shi also went to ER,.
because Jin Paul also went to ER because of stomach ache..
But he was discharged last night, right?.
Yeah, so they have a very difficult day..
Yeah, so, boy, we need to pray for one another, all right?.
For those of us who are well, and relatively healthy,.
don't forget to serve one another,.
to care for one another, right?.
Yeah, why don't we spend some time to pray for one another?.
So in your quiet prayer,.
please pray for those who are in need,.
Jin Fan and Jin Paul,.
and all the people that God brings to mind..
(gentle piano music).
Father, we come before you, we know we are broken,.
we are weak, and we are poor..
We are not like the powerful..
We are not like those who have might..
They seem to be the main actors..
But we come before you, we know that..
Your focus is on us,.
because you love us,.
and you take care of us..
Never have a moment in our lives that's out of your sight..
Never have a minute in our life that's out of your control..
So we come before you to acknowledge.
that you are the Lord, you are our King,.
and you are indeed holding everything together in our world..
So we wanna give you thanks,.
we wanna praise you and worship you..
We sing songs to you,.
not to those on top of the pyramid,.
but to you,.
who came to us in Jesus,.
in the form of manger,.
the form of a baby,.
just like every one of us..
You came to save us, to serve us,.

$^{481}$even die on the cross,.
sacrifice yourself in order to redeem all of us,.
so that our relationship with the Father can be restored,.
so that we can have hope..
And thank you for your love and your grace..
And we don't deserve it,.
yet this grace came to our lives..
So again, in this season of Christmas,.
remind us how you live,.
remind us your life,.
remind us how our Lord Jesus came not to be served,.
but to serve all of us..
And yes, Father, we wanna remember.
some of our family members,.
some of our church members,.
and those who are having a tough time and challenging time..
And we, first of all, wanna uplift Jinfan,.
Paul and their baby in your hand..
We know that they had a very challenging day.
yesterday and today,.
and I just pray that you will grant them peace,.
you'll help them to calm.
and to know that everything is in your hand..
We praise you for bringing the blood pressure.
down to normal level..
We praise you for good news..
And I pray that you will be their protection..
And we just pray that Jinfan will be discharged today,.
and you will continue to protect.
and especially watch over Jinfan's pregnancy..
We also wanna pray for our brother, Jinbo..
We know that he was having just pain yesterday,.
and they were rushed to ER..
Sure, they have a very, very chaotic.
and very, very difficult night,.
but we praise you for good news..
And I just pray that you will continue.
to bring healing to Jinbo,.
to grant strength to Wanshe and their family,.
be with them, and knowing that, God,.

$^{521}$you are always our strength,.
and you are our power, and you are our refuge..
I pray for those we cannot mention by names..
I'm sure some of us are having a very hard time.
and a challenging season in our lives..
Help us to lift our eyes to the throne..
Know that you are the king, and you are in control..
Increase our faith so that we can trust in you..
Thank you for being our God..
Thank you for being, walking with us..
And you are faithful, and you are always faithful,.
yesterday, today, and tomorrow..
And your grace is always sufficient..
So we pray that we'll rely on you.
and to believe in you and to trust in you,.
trusting the fact that you are always our provider,.
and you are always our shepherd,.
and you will be with us always..
So pray for our church as well..
As we come to the end of the year,.
we are thinking about the next year,.
our plan, our budget, our financial need,.
and all of the ministries that we're gonna come before you.
and pray that you will bless us..
Whatever we do in your name,.
we'll receive blessing and we'll receive fruit.
so that we can bring glory to you..
And I pray that as we go out from here,.
you will send us out to live the life of kingdom.
in our family, in our workplace,.
in everywhere you send us,.
that we will demonstrate your love..
You will demonstrate your grace..
Thank you, Father, in Christ, we pray..
Well, I wanna remind you to continue.
to pray for one another..
For those who are weak and for those who are in need..
And I also wanna remind you in the season of Christmas,.
let's do something for other people..
Maybe a random act of kindness to our neighbors,.

$^{561}$to our brother and sister in the church..
And please remember one another.
and show the spirit of Christmas a little bit more..
All right?.
And by the way, this afternoon,.
I just made up my mind to go to the Zoolight.
in Lincoln Park..
Anyone wanna join us?.
You can get the tickets yourself.
and you can see me and we go together.
for the Zoolight at 5.30 p.m..
All right, 5.30..
All right?.
But you have to pay for your own ticket..
But you wanna hang out with some of the people.
together this afternoon, please..
Yeah, you're welcome to join us..
All right?.
\newpage



\section{路加福音 19:37-48}
\label{sec:WhzVz61wqtY}
\textbf{耶穌:復興的聖殿 (路加福音 19\_37-48) - 梁紫健傳道}
\newline
\newline
連結: \href{https://youtube.com/watch?v=WhzVz61wqtY}{\texttt{https://youtube.com/watch?v=WhzVz61wqtY}} ~~~~ 語音日期: 2024-12-11
\newline
\newline
\hyperref[sec:_TmbaOYl2cg]{\small{< < < PREV SERMON < < <}}
~
\hyperref[sec:index_chronic]{\small{[返順時目]}}
~
\hyperref[sec:index_scriptual]{\small{[返順卷目]}}
~
\hyperref[sec:UlWyJIvtS2w]{\small{> > > NEXT SERMON > > >}}
\newline
\newline
路加福音 19:37-48
\newline
\begin{longtable}{cl}
\hline
\hline
章節 & 經文 (和合本修訂版)\\
\hline
19:37 & \begin{tabularx}{0.7\textwidth}{X} 他將近耶路撒冷，正下橄欖山的時候，一大群門徒因所見過的一切異能，都歡呼起來，大聲讚美神， \end{tabularx} \\ \\ \relax
19:38 & \begin{tabularx}{0.7\textwidth}{X} 說：「奉主名來的王是應當稱頌的！在天上有和平；在至高之處有榮光。」 \end{tabularx} \\ \\ \relax
19:39 & \begin{tabularx}{0.7\textwidth}{X} 人群中有幾個法利賽人對耶穌說：「老師，責備你的門徒吧！」 \end{tabularx} \\ \\ \relax
19:40 & \begin{tabularx}{0.7\textwidth}{X} 耶穌回答：「我告訴你們，若是這些人閉口不說，石頭也要呼叫起來。」 \end{tabularx} \\ \\ \relax
19:41 & \begin{tabularx}{0.7\textwidth}{X} 耶穌快到耶路撒冷，看見那城，就為它哀哭， \end{tabularx} \\ \\ \relax
19:42 & \begin{tabularx}{0.7\textwidth}{X} 說：「但願你在這日子知道有關你平安的事，不過這事現在是隱藏的，你的眼睛看不出來。 \end{tabularx} \\ \\ \relax
19:43 & \begin{tabularx}{0.7\textwidth}{X} 因為日子將到，你的仇敵要築起土壘包圍你，四面困住你， \end{tabularx} \\ \\ \relax
19:44 & \begin{tabularx}{0.7\textwidth}{X} 並要消滅你和你裡頭的兒女，連一塊石頭也不留在另一塊石頭上，因為你不知道你蒙眷顧的時候。」 \end{tabularx} \\ \\ \relax
19:45 & \begin{tabularx}{0.7\textwidth}{X} 耶穌一進聖殿就趕出在裡面做買賣的人， \end{tabularx} \\ \\ \relax
19:46 & \begin{tabularx}{0.7\textwidth}{X} 對他們說：「經上說：『我的殿是禱告的殿，你們倒使它成為賊窩了。』」 \end{tabularx} \\ \\ \relax
19:47 & \begin{tabularx}{0.7\textwidth}{X} 耶穌天天在聖殿裡教導人。祭司長、文士和百姓的領袖都想殺他， \end{tabularx} \\ \\ \relax
19:48 & \begin{tabularx}{0.7\textwidth}{X} 但找不出方法來，因為百姓都側耳聽他。 \end{tabularx} \\ \\
[1ex]
\hline
\hline
\end{longtable}
$^{1}$大家好,我是Ken.
如果大家親切一點可以叫我梓健.
其實我還在加哥讀書.
我還未畢業.
暫時還在考考試.
希望我可以順利畢業.
今天很高興跟大家一起學習上帝的華語.
在預備講道的時候.
我想起我和太太Alice相識的故事.
我也問准了太太可以跟大家一起分享.
這裡是我們相識的地方,長洲.
我們在長洲的一個小島上.
的一個神學院相識.
這個小島的神學院.
每個星期五我們都會出隊.
去不同的地方傳福音.
第一年我選擇了監獄部.
我從不同的囚友去傳福音.
而Alice選擇了中學部.
去跟中學生傳福音.
其實我們是同一個成長小組.
是同一班同學.
我看到她很真誠.
和很認真地對待自己的生命.
我就很欣賞她.
我覺得她很愛主.
之後我就想.
不如去年選擇監獄部.
今年不如選中學部.
我也很想服侍中學生.
選之前我完全不知道.
她會再選中學部.
其實我不知道.
誰知道她做了中學部的部長.
我就沒辦法了.
只好和部門一起服侍.
因為這個舉動.
她又覺得我對她有意思.
她就說.
我和上帝納了約.

$^{41}$今年不拍拖.
我就問她.
你的今年是指一個學年的今年.
還是指一年有12個月的今年.
她就說就是今年了.
我不再說了,就是今年.
我要專心單身服侍著.
不拍拖.
不拍拖就好了,打球.
有時我知道.
她每天早上去長洲的運動場跑步.
旁邊有個籃球場.
我很喜歡打球.
我就去打籃球.
每天早上我都和她打招呼.
「你好,Alice」.
可以和她打招呼,聊天.
有個弟兄很好.
他知道我每天早上都去打籃球.
他覺得我很孤單.
他很想陪我去打籃球.
我跟他說.
「不要了,太辛苦了,這麼早」.
直到有一天.
神學院的老師問Alice一個問題.
她就說.
你如何確定你今年單身服侍上帝.
是上帝喜悅的.
你肯定今年這樣做是祂預立的嗎.
其實單身和不單身.
結婚都是上帝呼召的,對吧.
結果我們對大家都很有好感.
經常都糾結.
我們是否應該了解對方多一點.
這個糾結更加令我們.
不太專心服侍.
最後我們打破了固有的想法.
我們拍拖,訂婚,結婚.
一起服侍著.
重點不是單身和結婚.

$^{81}$其實兩者都是一個.
好好服侍主的條件.
不過有時候我們成為了基督徒.
可能我們也會有一個.
根深蒂固的想法.
就是覺得上帝一定喜歡我們這樣做.
覺得結婚或單身才可以服侍主.
但有時候我們會否一廂情願地.
要求上帝接納我們的服侍.
這就是耶穌透過不同的經歷.
去改變了我對拍拖,服侍.
和與人相處的看法.
我講這個故事的原因.
做開場白是因為.
我覺得我們信了主一段時間.
可能會有信仰上形式化的危機.
有時候我們也會走進一些.
例行公事的生活.
就是我覺得我立了志.
就等於上帝會收貨.
我覺得我參加教會活動.
就等於跟隨基督.
我想今天這篇道德提醒我們.
到底甚麼是信仰的本質.
信仰本質是和上帝的關係.
而不是一個行為清單.
不知道今天我們當中.
會否有聽眾姐妹覺得.
最近你的靈明成長會有停滯不前.
或者信了主很多年.
容易留在輸出圈內.
覺得初信的時候的改變.
就已經足夠了.
又或者找不回當初初信.
那種服侍的火焰.
其實我自己也是.
最近我也遇到很多事情.
令我明明也在一個低潮的裡面.
所以今天想和大家了解一下.
究竟跟隨基督是一個甚麼過程.

$^{121}$是一個持續更新的過程.
看看耶穌這個復興的聖殿.
會怎樣帶給我們改變.
首先我們要讀一下今天的經文.
是路加夫音19章38-47節.
我們一起讀.
1,2,3.
他將近耶路撒冷.
正下降冷山的時候.
一大群門徒因所見的一切異能.
都歡呼起來.
大聲讚美神說.
奉主明來的王是應當稱頌的.
在天上有和平.
在至高之處有榮光.
其中有幾個法利賽人對耶穌說.
老師 責備你的門徒吧.
耶穌回答.
我告訴你們.
若是這些人閉口不說.
最初也要呼叫起來.
耶穌快到耶路撒冷看見那聖.
就為他哀哭說.
現在這日子是平安的事.
不過這是現在是隱藏的.
你的眼睛看不出來.
因為日子將到.
你的仇敵要捉起土裡包圍你.
四面困住你.
並要消滅你和你裡面的疑慮.
連一塊石頭也要留在一塊石頭上.
因為你不知道你眷顧的時候.
耶穌一進聖殿.
就趕出裡面做買賣的人.
他們說.
經上說.
我的殿是禱告的殿.
你們都使他成為了賊窩了.
耶穌天天在聖殿裡教導人.
祭司長 文士和百姓的領袖都想殺他.

$^{161}$但找不出方法來.
因為百姓都質疑聽他.
通常我讀經的時候.
也會問幾個問題.
這幾個問題可能會幫助我們了解.
聖經的訊息.
今天我們所說的題目.
耶穌就是復興的聖殿.
只有一點.
我問的三個問題.
都只是指向這一點.
就是作為跟隨這個復興聖殿的人.
我們是需要被持續改變的.
你願意嗎.
如果你在想的過程當中迷失方向.
不要緊.
我們最後還是會回到這一點.
看我們怎樣去實踐.
所以去到這三個問題就是.
第一個 為什麼法利賽人要叫耶穌.
擲避自己的門徒.
是不是因為他們做錯了什麼.
法利賽人有什麼動機.
第二 為什麼耶穌這麼傷心地說出.
耶路撒冷的下場.
大家如果細心看耶穌的宣告.
耶路撒冷的下場其實是很慘的.
但為什麼耶穌要這麼悲傷.
第三 就是為什麼.
耶穌天天在聖殿教導人.
宗教領袖因為百姓聽祂的話.
而找不到方法殺祂.
其實領袖明明可以將祂交給羅馬人.
因為祂被稱為王.
就可以殺祂.
為什麼做不到.
如果清楚的話.
我們可以回到第一個問題.
就是為什麼法利賽人要叫耶穌.
擲避祂的門徒.

$^{201}$是不是因為他們做錯了什麼.
法利賽人叫耶穌擲避門徒.
主要是因為門徒在讚美耶穌的時候.
引用了詩篇裡面的說話.
承認了耶穌是王.
留意 門徒的讚美.
是特意加「王」這個字.
在詩篇118篇25節是這麼說的.
「奉主命來ge ,是應當稱頌的」.
而路加 即是門徒是這麼說的.
「奉主命來ge ,王,是應當稱頌的」.
這個認定令法利賽人非常不安.
因為他覺得.
耶穌不會是理賽亞王.
更加沒有資格做王.
所以猶太人就會更加覺得.
有點突兀.
為什麼呢.
猶太人的王和世俗的王.
有一點點不同.
猶太王是怎麼樣的.
猶太王主要的使命.
是要在苦難裡面去投靠而和華.
成為百姓的榜樣.
去教導百姓怎樣去遵守律法.
怎樣去依靠神.
而詩篇118篇就說過猶太王的功能.
所以我們一起來看看這個詩篇.
不過有時候弟兄姊妹會問我.
新約還新約 舊約還舊約.
為什麼要看新約 還要看舊約呢.
我就會回答他.
如果你想要一塊夏威夷薄餅.
它是一塊有菠蘿的薄餅.
它是一塊很多汁的菠蘿.
如果它沒有多汁的菠蘿的話.
它就是一片普通的薄餅.
它就不是一個夏威夷薄餅.
所以如果沒有舊約的理解的話.
我們就很難明白新約的處境.

$^{241}$究竟在說什麼.
它們是一體的.
所以我們來看看.
詩篇118篇25至27次是這樣說的.
猶太王帶領百姓去到上帝面前.
去祈求拯救.
這裡說到一個王.
他去到帶領群眾去祈求昌盛.
跟行路法很有關係.
為什麼呢.
猶太王祈求.
耶和華拯救和使我們亨通的這個字眼.
其實它的原文是想解釋.
順利.
使我們順利的這個請求.
就會令我們想起詩篇一章三節.
是什麼.
人喜愛耶和華的律法.
晝夜思想.
這些人就會什麼.
有福沒錯.
凡他所做的.
都順利.
沒錯.
順利這個字.
有很密切的關係.
所以當詩篇118寫的這個王.
去代表著百姓去祈求.
拯救和順利的時候.
這個王就必定要教導百姓去遵守律法.
他才有資格去代表百姓.
去到上帝面前去祈求.
這樣說它的下文.
你要看回就會很明顯.
詩篇119第一節是這樣說的.
行為完全遵行耶和華律法的就會有福.
所以這裡詩篇帶出一個重點就是.
猶太王是要教導百姓去遵守耶和華的律法.
可能要再說一次.
猶太王的主要功能.

$^{281}$是要教導百姓去遵守律法.
我們就問.
我們知道了這個功能之後.
和《路加福音》耶穌被要求.
閘閉門徒有什麼關係呢.
第一就是.
法利賽人根本就不覺得耶穌是王.
第二.
更加不會認同他要教導百姓律法.
法利賽人是什麼人.
法利賽人是熟讀聖經.
法利賽人是覺得他們自己才有資格教導律法.
所有的百姓都要根據他們的詮釋.
去活出律法.
所以.
你如果用今天的眼光去看他們.
就有點像神話院的老師.
他有權去寫commentary.
有發言權去解釋聖經.
所以對他來說.
他們法利賽人教導律法是很正常的.
不過.
耶穌這個外來人.
這個木匠的兒子.
憑什麼教導聖人院的老師呢.
所以我們猜到.
法利賽人叫耶穌閘閉門徒.
不是因為門徒做錯了什麼.
而是因為他根本覺得.
耶穌沒有資格教導律法.
所以耶穌馬上回應他們.
我告訴你們.
若是他們閉口不說.
這些石頭必要呼叫起來.
為什麼叫石頭呢.
其實這一段跟下一段.
耶穌所說的寓言很有關係.
但大概在這裡.
耶穌沒有去閘閉門徒.
明顯是在認同門徒的讚美.

$^{321}$所以回答第一條就是.
門徒叫耶穌閘閉門徒.
是因為他們覺得.
耶穌根本沒有資格去教導他們.
即是你對著你的門徒可以.
夠有資格.
但你對著我們.
你怎麼可以.
就好像NBA.
為什麼呢.
你經過高中,大學.
最後還要這麼多輪.
兩輪的選拔.
最後選了60個NBA球員進去.
然後耶穌就好像在街上打球的大叔.
隨便撿幾件.
不懂打球的街童.
帶著他去找NBA的教練和球員.
就在他們面前說.
我是最厲害的教練.
比所有人都厲害.
所以門徒根本就覺得.
耶穌沒有資格教導百姓.
叫他罵自己的門徒吧.
如果你們要解答第一個問題.
我們就下第二個問題.
就是接下來.
為什麼耶穌這麼傷心地說出.
耶路撒冷的下場呢.
如果你細心看41-44節的話.
其實挺恐怖的.
敵人會殺光你裡面的兒女.
這是一個很狠的語言.
耶路撒冷是一個什麼地方.
耶路撒冷是一個宗教中心.
是猶太人會用性命去保護的一個宗教中心.
這裡是他們的一切.
這裡是他們獻祭的地方.
還有猶太人最重要.
第二重要的其實是什麼.

$^{361}$後代.
他們的兒女可不可以在這個應許之地.
去生活下去.
是一個非常重要的事.
現在耶穌就說.
他們的兒女被人殺光.
詛咒他們沒有後代.
還要耶穌專門用了一個很特別的表達.
石頭和兒女在白來文.
其實是一個很相似的讀法.
就是你對著猶太人說.
石頭不會再有石頭.
兒女也不會再有兒女.
這個是一個很強的詛咒.
所以我們要問的問題就是.
究竟這個對於猶太人來說是什麼.
在第二聖殿開始.
他們二十個世代都在這裡生活了這麼多年.
他以為自己做的事是得到上帝的歡心.
但突然有個人走出來就說.
你們不會在這裡生活下去的.
你們會走的.
結果聖殿就在公元後七十年.
被羅馬軍隊毀滅.
是提多將軍叫人火燒聖殿.
而聖殿裡面有很多金.
燒到溶掉.
他看到金在聖殿流出來的時候.
然後他就發現裡面有金.
他們就可以搬走所有石頭.
所以聖殿裡面一塊石頭都沒有.
但在一塊石頭上面.
但說了這麼久.
我突然要問.
為什麼耶穌這麼傷心.
為什麼.
很可能是耶穌作為猶太的領袖.
看到猶太領袖的失職.
為什麼這樣說呢.
這幫猶太領袖包括民事.

$^{401}$法利塞人 祭司長.
他們的工作是做什麼的.
他們的工作是在聖殿的內和外.
去教導百姓怎樣去敬拜上帝.
但他們就失職了.
為什麼呢.
就是45至46節是這樣說的.
耶穌進入了聖殿.
就趕走裡面做生意的人.
就說我的殿必作為禱告的殿.
你們到使他成為了賊窩.
這裡可以打開大家.
這兩句其實是舊約的聖經.
出自於哪裡呢.
很長的.
有人說是詩篇 但不是.
很長的.
很長的 《先知書》.
很長的.
沒錯 《以賽亞書》.
對了 還有最後一卷.
《耶利米蘇》 沒錯.
就是出自於這兩卷很長的書卷.
今天我沒什麼時間跟大家提醒.
但你會看到其實在.
《以賽亞書》56章和《耶利米蘇》7章.
都說了同一個主題.
就是領袖的失職.
就例如.
「這些牧人不知明辨他們都偏行幾路」.
「你們卻不聽我 不質疑而聽」.
「竟硬著腳看比你們的祖先更甚」.
其實還有一個主題.
除了失職之外.
其實他們這兩卷聖經發生的地方.
都是在聖殿.
在聖殿.
這兩卷聖經都指出.
猶太人 以色列人.
在聖殿裡面沒有進到.

$^{441}$做上帝子民的職份.
而當中宗教領袖.
當然是最大的責任.
不單止是因為他們沒有禱告.
不是因為他們沒有禱告.
是因為他們失職.
沒有好好教導律法.
令猶太人要被上帝趕走.
所以回答 為什麼耶穌會哭.
是因為猶太領袖的失職.
而感到很傷心.
去到最後一個問題.
就是為什麼耶穌天天在聖殿裡面教導人.
而宗教領袖就因為百姓質疑而聽.
而殺不了他.
耶穌很明顯是來踩場的.
他帶著這麼多階級.
但可以找人趕走他.
或者將他交給羅馬人.
就好像.
如果你看其他福音書的話.
其實有些福音書.
譬如馬可.
就將他交給羅馬人.
但耶穌是專門這樣寫的.
他說民事和祭司長百姓都是想殺他.
但是殺不了他.
因為什麼.
因為百姓質疑而聽.
為什麼呢.
是不是因為傳媒報導.
因為東張西望.
因為坑槍集.
那時候沒有傳媒.
沒有新電視.
只有一群在聖殿裡面的百姓.
其實這個很關聖殿的功能.
其中一個在川崎大學讀新約研究的研究生.
叫做Cindy O.
她的博士論文就說了.

$^{481}$三個主要聖殿的功能.
除了憲制之外.
第一就是.
聖殿的節日是加強了彼此之間的關係.
一起吃飯.
令到大家不會當彼此是外人.
所以今天早上吃的早餐.
我覺得.
那個炒麵.
那個粥.
我覺得真的很自己人的感覺.
一起吃飯.
你會看到耶穌在福音書裡面.
都不會拒絕和罪人吃飯.
因為他當對方都是自己人.
而第二個就是.
共同利益.
第三.
今天最想說的就是.
聖殿的群體裡面其實有一個教導的功能.
就是每逢七年的時候.
整個猶太社區會在鑄盤節聚集.
大祭司會重奪律法.
令到百姓學到.
他怎樣在生活當中去履行公義.
所以耶穌現在在聖殿裡面做什麼.
是教導.
他教導律法.
他教導百姓.
最重要的是那些百姓有什麼行動.
質疑而停.
耶穌在聖殿裡面做了領袖需要做的事.
就是教導那些百姓怎樣去實踐律法.
而那些百姓明顯是受教.
成為了耶穌的群體.
所以那個領袖不是純粹因為群眾壓力而殺不了耶穌.
而是因為耶穌取代了他們.
做了教導的功能.
也是耶穌取代了聖殿本身的功能.
他們沒有辦法取代聖殿本身.

$^{521}$剛才說的三個問題.
其實都是指向耶穌就是那個復興的聖殿.
我們要跟隨的話就要不斷被他改變.
就好像《路黑福音》裡面那樣.
耶穌一來到就令到宗教領袖很大的改變.
他們覺得他們以為自己在做上帝給他們的事.
但耶穌說錯了.
有時我們信了很久會不會覺得.
我們每天祈禱,讀經敬拜,在教會服事.
這樣就夠資格做一個基督徒了.
我不是說這些事不用做.
我是說不要以為單單做這些事就等於跟隨了耶穌.
最重要的是我們不斷被耶穌改變我們的想法和信念.
我們可以有相應的行動去做出來.
我想跟大家分享一下自己的信念和想法被耶穌改變.
早年我在….
家人早年在深圳做生意.
遇到經濟的衰退.
生意突然間就虧到破產.
整個家庭的經濟跌到谷底.
整個家庭都很挫敗.
同一時間我還要患上一個病叫腦癌.
經常抽筋.
突然間就跌倒.
還失憶.
每次溫習的時候我都很怕.
好像電腦一樣會重啟.
之後我又很慢吞吞.
不知道怎樣很呆.
不知道怎樣去讀書和溫習.
那時候我的人生都很灰暗.
我覺得我的人生沒有甚麼色彩.
我經常都在一個恐懼的裡面.
我連在街上打球都很害怕.
當耶穌到了我的生命裡面的時候.
祂改變了我對人生的看法.
我選這張圖的原因是.
大家一看它不是香港的馬路.
其實是深圳的馬路.
有一次我突然在這條馬路發作.

$^{561}$在中間.
我整個人不受控制.
我也不知道自己有沒有意識.
在馬路上轉圈.
如果大家知道或熟悉大陸的法例.
如果我在那裡被車撞到.
通常人們會把我車死.
為甚麼呢.
因為車死一個人賠少一點錢.
如果車不死那個人就要賠很多錢.
我不知道發生甚麼事.
突然有個人拉著我的手.
拉到我馬路邊.
我醒來之後就完全沒有人在我身邊.
天使找人來救我.
不單止這樣.
耶穌還醫治了我這個病.
令我康復.
有一次我發作很辛苦.
我吃藥都不行.
我祈禱上帝你醫治我吧.
突然有一天.
我沒有再發作.
我明白對於一個病了十多年的人來說.
我不敢相信那件事是這樣發生的.
我去看醫生.
連醫生也覺得不可能.
為甚麼會這樣.
今天我沒有跟大家說很多細節.
重點是耶穌改變了我對人生的看法.
有時候我覺得人生很絕望.
很灰暗.
但突然耶穌給了我一副眼鏡.
讓我看到這個世界是有色彩的.
是溫暖的.
雖然有時候我也會灰心失意.
我的後遺症依然存在.
有時候我會驚恐.
有時候我會突然很頭痛.
有時候可能我說不到一句很完整的話.

$^{601}$但相信耶穌不代表我們沒有苦難.
我人生的眼光從此被改變.
不再那麼灰暗.
我不知道當中有沒有舞蹈者在找尋人生的希望.
你願不願意經歷耶穌給你的人生的轉變.
你接受耶穌的話.
祂會給你一副很特製的眼鏡.
讓你看到你的人生很不一樣.
即使我覺得.
像耶穌剛剛說的.
我們拍拖,結婚.
也有很多觀念要被上帝改變.
我以為我做的事情是上帝喜歡的.
但耶穌透過很多不同的東西去改變我們的觀念.
第一個就是.
和家人的關係也是這樣.
有時候我和我爸爸是很不投契的.
因為各種立場不同.
經常吵架.
我猜大家經歷運動之後.
也會有這種看法.
或者狀態跟上一代有很多衝突.
後來我發現耶穌改變了我對人的看法.
我發現.
就算立場不同或做法不同.
對方也是上帝所愛的.
立場不需要去定義我們的關係.
我們之間是親情,血緣.
是上帝給我們的關係.
之後我就順服了.
其實我爸爸說什麼都好.
我都回答他OK.
他有他的立場,我有我的意見.
我只關心他的人.
睡得好不好.
吃得好不好.
和我媽媽的關係怎麼樣.
有一次我記得.
我媽媽和我爸爸有很大的衝突.
我心底裡很清楚爸爸是錯的.

$^{641}$我覺得他死性不改.
但有一次送完之後.
有時候他對我很感動.
我們要去關心他.
我知道我們沒辦法勸教.
但是.
我心底裡有很大的情緒波動.
我們這麼差.
我和Alice都很差.
上帝都沒放棄我們.
神都沒放棄我們.
所以神都沒放棄他.
所以當時我們.
本來下午要上課.
星期天都上課.
我們走去見他.
有時候我們以為.
有些人死性不改.
不會有改變的可能.
但當我們回想自己.
其實我們也死性不改.
很多改不了的地方.
他都給我們機會去改.
雖然我和我爸爸之間.
還有很多意見的不同.
但因為耶穌.
他令我們的關係改變了.
他改變了我對人的信念.
改變了我人與人之間的關係.
特別是和家人.
說到子女.
雖然我沒有子女.
但我覺得這個世界.
說一句.
很多令我們子女要成功.
成功的意思是.
我的子女能不能讀書.
能不能讀好的大學.
能不能找到一個好的工資.
但作為跟隨耶穌的父母.

$^{681}$我們想想.
子女的成功是否只是建基於.
金錢.
她賺到的錢.
還是我們看看.
她是否和自己相處得好.
她和人相處得好不好.
她在同學之間的見證.
和上帝之間的關係.
我覺得可能我們也要調教一下.
究竟我們子女的成功是怎樣定義.
我一向未讀書.
我覺得讀到一定程度也會有迷失.
特別是讀神學.
覺得自己拿不到高分.
覺得自己拿不到學位.
等於自己是失敗.
一無是處 是垃圾.
一路以為讀書才是上帝的心意.
讀不讀書其實並不是.
但我覺得這個錯覺.
很需要被上帝改變.
因為無論如何.
上帝也在等我們.
最重要是我們和上帝的關係.
剛剛過去的感恩節.
在美國大家都想和家人一起相聚.
做上帝的恩典.
其實香港人也有類似的節目.
無論何時.
在這個好日子裡大家都相聚.
有一句話叫做冬至大過年.
為了甚麼呢.
就是為了家人團團圓圓團聚.
強調家人之間的關係和親情.
可能在美國下班後大家一起吃飯.
但這是我們去年.
神學生一起冬至過冬的日子.
雖然我們沒有包湯圓.
但這個臘味炒飯我們也炒了三個小時.

$^{721}$香港人去到哪裡都想團聚.
我們不斷更新做冬的方法.
不斷做不同的事.
令大家相聚.
因為我們想精神一直傳承下去.
不斷改變不斷更新.
令新一代.
你的下一代.
也會有將中至大過年的習慣.
傳承下去.
為了強調家人的團聚精神.
同樣地.
記得信仰也是.
有一群人兩千多年前.
一直跟著耶穌不斷被他改變.
不斷更新.
不斷被塑造.
直到我們今天才聽到福音.
直到我們今天才知道福音的真實.
你願不願意成為這群人呢.
等耶穌塑造我們的生命和觀念.
我們一起祈禱.
耶穌多謝你.
多謝你今天讓我們一起學習你的話語.
我們每個人都有自己的想法和思維.
在不同的環境裡面.
可能有很多的掙扎.
很多的痛苦.
但耶穌求你幫我們.
幫我們有一個柔軟的心.
不斷讓你改變和無助.
不斷在生活裡面去應用你的道.
弟兄姊妹的心聲你是聽到的.
求你幫助我們.
來感到高雄耶穌的聖靈.
祝福大家.
\newpage



\section{20:1-17}
\label{sec:UlWyJIvtS2w}
\textbf{New Principles for New People (Exodus 20\_1-17) - Rev. Mark Chan}
\newline
\newline
連結: \href{https://youtube.com/watch?v=UlWyJIvtS2w}{\texttt{https://youtube.com/watch?v=UlWyJIvtS2w}} ~~~~ 語音日期: 2025-01-04
\newline
\newline
\hyperref[sec:WhzVz61wqtY]{\small{< < < PREV SERMON < < <}}
~
\hyperref[sec:index_chronic]{\small{[返順時目]}}
~
\hyperref[sec:index_scriptual]{\small{[返順卷目]}}
~
\hyperref[sec:3GTzRC3gFVM]{\small{> > > NEXT SERMON > > >}}
\newline
\newline
$^{1}$All right, all right guys,.
today I would like to meditate with you.
a very important subject, very important message,.
very important theme in the Bible..
In fact, if in the Old Testament,.
if you take the New Testament away, the Old Testament,.
this is the most important passage of all..
It's the passage of 10 Commandments..
All right, I wanna do this because.
this is the end of the year,.
and we'll linger into this passage next Sunday as well..
Next Sunday will be 2025, right?.
And this is a very, very important subject.
for us to meditate on, especially at the end of the year..
And why?.
Because this is the 10 Commandments..
The 10 Commandments has been known as the most important.
and most influential document of probably in the West,.
at least in the West, in the Western civilization..
That is the, probably one can say.
is the moral and ethical foundation of all mankind.
that God has given to us..
That is the 10 Commandments..
And the series is titled,.
New Principles for New People,.
and it's taken from Exodus 20..
So in the entire Old Testament,.
if you have to remember a biblical reference, all right,.
from which book and from which chapter,.
probably you have to memorize this, Exodus 20..
Exodus 20..
And that's where you found one of the versions,.
one of two versions of 10 Commandments..
The other version you can find is Deuteronomy 5,.
but you can memorize in your head, in your brain,.
memorize this is Exodus 20..
Now, when we look at the New Commandment,.
the 10 Commandments, oftentimes the word commandment.
carries some sort of negative connotation, right?.
You believe in Christ and you're a Christian.

$^{41}$and you're a follower of God,.
and then you heard a lot of about commandments..
And commandment has a very restrictive sense to it, right?.
Because when we hear a commandment,.
this is commanded for you,.
and you be commanded to do something,.
and not to do something..
So often it sounds very restrictive.
because it's kind of like, you know,.
boxing in our freedom, our choices..
But, you know, if you turn to the Bible,.
especially the Old Testament,.
how the Bible described this document.
of what we call now 10 Commandments,.
it's not even using the term commandment.
in the Hebrew scripture..
It's using terms like 10 words..
And that's why we get the translation,.
dekala, dekala, dekala, 10, right?.
Logos, 10 words..
Or one can describe the 10 Commandments.
not so much as the restrictive sense.
of the word commandment mean..
We should describe it as constitution..
Just like in the United States,.
we have the constitution, we have the document..
It's not legal code..
It's not something that attorney would go to.
and then find out which code to prosecute someone.
or to make sure someone is in line.
or prosecuting a crime or something like that..
So it's not legal code, it's a constitution..
The constitution has a general sense of being principles,.
all right, being living principle or the big principles..
And that's why this series is called.
New Principles for New People, for new people..
Now, these are big principles..
Larger principles or higher level principles.
for the new people of God..
And these principles are like the United States Constitution..

$^{81}$Those terms and the amendments and all those statements.
are general guiding principles so that you can apply those.
to a specific situation or context..
So we have to remember,.
it's not so much of commanding you to do this.
or that or not that or do this or that..
It's not a list of things to do or things don't do..
It's 10 words or 10 things,.
10 principles for the people of God..
And because this is a very familiar passage, I hope, right?.
I hope to read this passage.
from the beginning to the end..
There are 17 verses for the 10 commandments or 10 words..
Okay, let me read this through first..
Exodus 20, "And God spoke all these words,.
"I am the Lord your God who brought you out of Egypt,.
"out of the land of slavery..
"You shall have no other gods before me..
"You shall not make for yourself an idol.
"in the form of anything in heaven above.
"or on the earth beneath,.
"or in the waters below..
"You shall not bow down to them or worship them..
"For I, the Lord your God, am a jealous God,.
"punishing the children of the sin of their fathers.
"to be the third and fourth generation.
"of those who hate me,.
"but showing love to a thousand generations.
"of those who love me and keep my commandments..
"You shall not misuse the name of the Lord your God,.
"for the Lord will not hold anyone guiltless.
"who misuses his name..
"Remember the seventh day by keeping holy..
"Seven days you shall labor and do all your work..
"But the seventh day is the Sabbath to the Lord your God..
"On it you shall not do any work,.
"neither you nor your son or daughter,.
"nor your manservant or maidservant,.
"nor your animals,.
"nor the alien within your gates..

$^{121}$"For in six days the Lord made the heaven and the earth,.
"the sea and all that is in them..
"But he rested on the seventh day..
"Therefore the Lord blessed the Sabbath day.
"and made it holy..
"Honor your father and your mother.
"so that you may live long in the land.
"the Lord your God is giving you..
"You shall not murder..
"You shall not commit adultery..
"You shall not steal..
"You shall not give false testimony against your neighbor..
"You shall not covet your neighbor's house..
"You shall not covet your neighbor's wife.
"or his manservant or maidservant,.
"his ox or donkey or anything that belongs to your neighbor.".
There's 70 verses, 10 commandments..
Have you heard of it?.
Yeah, you have heard of it, right?.
Have you got it memorized?.
All 10 commandments?.
Well, those are 10 words, right?.
Those are 10 big principles for the people of God..
Now, you have to understand at the time of Exodus,.
at the time of Moses,.
now Moses got this 10 commandments on the mountain.
called Mount Sinai..
Now what happened before they got.
the so-called download from the cloud,.
download from the cloud,.
the 10 words from the top of Mount Sinai,.
beforehand, what happened to them?.
It was Moses called by God to lead the Israelites.
out of the land of slavery, out of land of Egypt..
And they have been oppressed for a hundred years..
And God saved them by grace..
And through the, well,.
they walked through dry ground,.
crossing the Red Sea,.
and then through the wilderness,.

$^{161}$and they finally come to the most important mountain.
in the Bible, which is Mount Sinai or Mount Horeb..
Now they are at the foot of mountain,.
and Moses went up to the mountain top of Mount Sinai.
and download the 10 commandments in two tablets..
And now they are new people of God..
The Israelites were saved by God, saved by grace..
Saved by grace..
It's not because God looking down in Egypt.
and find that the Israelites are behaving well..
So God decided, okay, I'm going to save you guys..
No, they were sinners..
They were in the land of slavery..
They did nothing to merit the salvation of God..
And God saved them because God is faithful to them..
So God saved them totally by grace..
And now they are being led to the foothill.
of the Mount Sinai..
And they are now the new people of God and the new nation.
and living under new principles..
So they have this new identity..
That's why the 10 commandments are given to them..
The 10 words, the 10 principles,.
so that they will live as the new people of God..
This is a totally different people.
from all the other peoples on the earth..
Now at the time of Moses, different nations.
and different kingdoms will make covenant,.
will make agreement with other people..
And so often these international agreement or covenant.
between nations, between kingdoms for diplomacy..
Like kingdom A will make agreement with kingdom B..
It's kind of like the United States.
has different diplomatic relations.
with other countries, right?.
Canada is one of our friendly ally..
UK or EU is our ally..
And all those things are actually happening.
in the time of Moses as well..
So in the ancient world, when nations come.

$^{201}$under some sort of relation, some sort of diplomacy,.
the relationship, political relationship.
or arrangement with each other,.
they will have a document to spell out the details,.
to spell out the terms or the laws between them..
Especially when a bigger king, a greater king,.
will make agreement with a smaller kingdom,.
with the smaller king..
And the great king will lay out all the details.
and all the terms..
What we call now is the Suceran Treaty,.
Suceran Covenant..
Now I know that's a big word, but just bear with me..
It's just a political arrangement.
between a great king and a lesser king..
In the time of Moses, that's how they do diplomacy,.
international relation..
What we call now is Suceran Treaty..
Suceran meaning the sovereign, right?.
The bigger king..
And he will lay out all the terms.
and all of this documents and agreements between nations.
in the time of Moses, about 1500 BC or 1200 BC,.
they will follow a similar format..
And a similar format are pretty much right here I outline..
Usually those documents, Suceran Treaty,.
will begin with introduction..
The introduction of the great king..
So the great king will introduce himself as I am so and so..
I am so and so king of Hittite..
I'm making a covenant,.
I'm making a Suceran Treaty with a smaller king..
Or I am the great king of Egypt,.
great king of other land..
And he will introduce himself a little bit,.
and then he will come to the second element,.
which is the history, or what we call historical prologue..
What it does is that the great king.
will recite some of the historical relation.
with the lesser kingdom,.

$^{241}$and how he has protected them,.
how he has guaranteed their security,.
how he has bussed the smaller nation..
And it will detail some of the historical events.
that happened before between the big kingdom.
and the smaller kingdom..
And that's the second element..
The third element after the introduction of the great king.
and the historical prologue is the laws..
The laws or regulations, or the stipulations..
And those are kind of the thou shall do this.
or thou shall do not, all right?.
Thou shall not or thou shall..
You should do this and you should do that,.
and you should not do this or that..
Things like you have to be faithful to me..
So the great king will say to the lesser king and say,.
"Hey, because I have protected you.
"in the history of your nation,.
"I've been blessing you guys,.
"and that's why you guys need to be loyal to me,.
"and you need to pledge allegiance to me.".
Those are the laws..
And then number four, the fourth element,.
is the succession element..
Succession element meaning when the document is draft,.
when the covenant is agreed upon,.
the lesser kingdom and the greater kingdom.
will agree to continue the agreement..
Kind of like the end of the year,.
you have to sign some sort of document.
to renew your insurance policy,.
to renew your relation with somebody else..
Those are what we call succession elements..
And that's number four..
Number five is calling upon all kinds of witnesses.
to draft the document, to finalize the document..
They will call upon the gods or goddesses.
to witness the event of coming into relations.
with each other between the two nations..

$^{281}$And then finally, there will be a section.
called blessing and curse..
And what those are statements that will say,.
"If you are faithful to me,.
"then you will enjoy this blessing..
"You will enjoy my presence..
"If you are disloyal to me,.
"if you rebel against me, then there will be curses..
"The greater king will threaten the lesser kingdom,.
"the lesser kingdom, and will say,.
"if you are disloyal to me,.
"then I will send my army and destroy you.".
And all sorts of things..
So those are very, very long section,.
and with a lot of different statements saying,.
"You will be blessed if you do be faithful to me..
"You will be cursed if you rebel against me.".
Now, if you can kind of read through these elements,.
one, two, three, four, five, six,.
and then you compare these six elements.
with the biblical documents,.
especially the first five books of Moses,.
you will find these elements in different places..
Like the introduction, you will find the passage.
which is read, the very beginning,.
which is the introduction of the great king.
being Yahweh himself, right?.
The true God who saved Israel out of the land of slavery..
And that's number two..
Versus chapter 20, verse 2b..
That is the history, right?.
I am the Lord your God..
I have saved you out of the land of oppression..
And then number three, element number three,.
are those laws, what we call 10 commandments..
You shall have no other God before me..
And all that, one, two, three, four, five, up to 10..
10 words..
And then there are also some sort of elements.
of secession, right?.

$^{321}$Deuteronomy 31, it talks about, you know,.
you guys need to be, the Israelites need to be.
continuously faithful to God..
And then not only in this generation,.
but the next generation, not only you,.
but your children and your children's children..
You guys need to be faithful to me.
and you have to recite this document again and again,.
again and again to renew your faithfulness again and again..
And then there's also in Deuteronomy 30, witnesses..
And because, you know, in the biblical religion,.
there's only one true God, right?.
And all the other gods are false gods..
And they cannot be witnesses..
So in Deuteronomy 30, being called upon to witness.
the event of the covenant is heaven and earth..
It's the entire cosmos..
So the creation is witnessing this foundational moment.
that God and his people enter into a treaty, a covenant..
And then of course, in different chapters.
in the five book for Moses,.
there were so many sections on blessing and curses..
And those are statements that says,.
if you are faithful, then I will be faithful to you..
If you are not faithful, then you go into exile..
Now, having revealed the format.
of the 10 Commandment document,.
and we now come to the point that we need to understand.
not only the draft documents.
according to the format of the day,.
but they also practice it in such a way.
that will align with the practice.
of the nations coming into covenant..
Say for example, two nations,.
and they formalize the covenant,.
they will inscribe the terms of the covenant.
on stone tablets..
So the 10 Commandments,.
that's why we have the two tablets, right?.
Inscribe on stone tablets, something that is permanent..

$^{361}$So God and Moses, when they talk about,.
agree upon the treaty terms,.
the terms will be inscribed on stone tablets..
So that's why God told Moses,.
hey, you come to the top of the Mount Sinai.
and bring two copies of the stone tablets..
And what they did is align with other nations.
when they come into a diplomatic relationship,.
because they will bring two tablets.
and then inscribe their international covenant.
on the tablet, and then one copy for me,.
and one copy for you..
And then you're gonna bring your copy,.
the stone tablet with all the terms of the covenant,.
you bring it to your capital,.
and in your capital, usually there is the biggest temple,.
and you put it in the holy place of your biggest temple..
And I will do the same thing..
And that's why there were two exact copies.
being made in the covenant-making moment..
So often the modern impression is that.
there are two tablets,.
and one tablet is for Commandment 1 to 5,.
the second tablet has Commandments 6 to 10..
We thought, you know, maybe one tablet is not enough,.
you need another copy to write the rest of the commandments..
No, that is not the way..
That is our modern misinterpretation..
The right interpretation should be.
two tablets of exact copies,.
and one for you and one for me..
And that's why in the Israelite temple,.
which is Tabernacle in the wilderness,.
and then later on in the capital city, Jerusalem,.
the stone temple, in the holy place,.
there is a copy of the stone tablet..
That is for the nation of Israel,.
for the new people of God..
So they put the tablet, put the covenant.
in the very holy place, the most holy place of the temple..

$^{401}$And then God, of course, Yahweh himself,.
he's the true God, and he has only one temple to go to,.
right, which is the earthly copy of the heavenly temple.
in Jerusalem or in the wilderness in the Tabernacle..
So he also put the tablet, God's tablet, in the same place..
And that's why in the nation of Israel,.
in the Tabernacle or in the Jerusalem temple,.
you will find two exact copies in the most holy place..
So having said that, we need to understand.
when God gave the 10 commandments,.
when God gave the constitution,.
it was the climax moment of the entire Bible..
So the laws of God were not for the purpose of restriction,.
but it is a way for the people of God being saved by grace.
to respond with thanksgiving..
So the 10 commandments are being given,.
not so much as you shall not do that, you shall do this..
It's not so much of a list of do's and don'ts,.
but statement of grace..
So when the people first heard about the 10 commandments,.
they would have responded with thanksgiving.
because God saved me, saved us out of the land of slavery,.
land of sin and land of death and oppression..
God saved us by grace..
So in some other tradition, like the Jewish tradition,.
the first commandment actually is not starting.
from verse three, which is thou shall have.
no other God before me..
In the Jewish tradition, the first statement.
of the 10 words, the first word of the 10 words.
is actually verses one and two..
To be exact, it's verse two..
It says, I am the Lord your God, I'm Yahweh your God,.
who brought you out of Egypt, out of the land of slavery..
Now it's not so much of do this or don't do that, right?.
It's not even moral code..
It is a statement of grace..
I have saved you, I brought you out of Egypt by grace..
And now you are free and you come under me.
and I am your true God..

$^{441}$I am your new Lord, no longer Pharaoh being your God,.
your Lord, Yahweh is your Lord..
So the Lord are given as a way of telling the people.
that you need to be thankful, right?.
Now as a Christian for many years,.
I've always reminded how God has saved me out of grace,.
right?.
He has saved me by grace..
And that is the foundation of my relationship with God..
And so your relationship with God as well..
So often we messed up..
We think what we do or what we don't do.
will count us as the people of God,.
as a follower of Christ Jesus..
No, the foundational truth for spiritual life.
is that God has saved you by grace..
It's not because what you have done..
It's because God has done something wonderful in your life..
And that's why you receive the new principles.
to live as someone else, with a new identity.
as the children of God, as a son and daughter of God,.
as a follower of Christ..
I have a new principle in my life to live differently..
So I manage my time, I manage my life differently..
I manage my family differently..
I treat my children differently..
I treat my job differently..
I treat everything God has provided for me differently..
It's a totally different principle..
With that in mind, we come to the so-called first commandment..
You shall have no other gods before me..
Now in the ancient world, there were many gods..
It's polytheistic..
And all the gods, interestingly, are territorial..
All the gods are territorial..
It's kind of like our Chinese or Asian religious tradition,.
the Eastern religious tradition..
Different gods have different realms..
Different gods have different territories,.
different areas that they rule..

$^{481}$It is also true in the ancient world of Israel..
In the Bible, one god, one force god is the storm god..
Another god is the sun god..
Another god is the moon god..
And they have their own territory..
And the statement here,.
you shall have no other gods before me..
What is new here is the word actually before me..
Yahweh said, the true God said, there's no God before me..
You shall have no God before me..
The emphasis is actually on before..
It's in my face..
Literally, it means in my presence, in my face..
And that means the God of Bible is all present..
He is everywhere..
He's everywhere..
And you cannot hide from him..
You cannot hide from him..
You cannot say, oh, I don't go to church,.
so I hide from God..
You know, God cannot see me..
You know, I don't wanna show up in the church.
because I don't wanna see God..
I don't wanna have God see me..
No..
It says, no other gods before me..
So the entire universe is filled with the presence of God..
Everywhere in the universe is before God..
So here is a monotheistic statement..
Says that one God only..
He is the one and true God..
He's not territorial.
because your entire life is under his surveillance..
Your entire life, every single moment,.
every single aspect of your life is before God..
So commandment number one says,.
you shall have no other God before him..
Before the true God..
Number two, you shall not make for yourself an idol.
in form of anything in heaven above.

$^{521}$or on the earth beneath or in the waters below..
You shall not bow down to the idols or worship the idols..
For I, the Lord your God, am a jealous God,.
punishing the children for the sin of the fathers,.
the third and fourth generation for those who hate me,.
but showing love or kindness to a thousand generations.
of those who love me and keep my commandments..
Now that's interesting here..
It's the word jealous..
But jealousy is usually a bad thing, right?.
Most of the time it's a bad thing..
You know, jealousy is I don't want you to have something.
that I don't have, right?.
It's like you are more beautiful than me..
I don't like it, right?.
That's why I want to destroy you, right?.
That is jealousy..
Jealousy is usually a bad thing..
It's a hateful thing..
But you know, in the entire cosmos,.
there are two types of relationship..
There should be jealousy..
You say, wait, what?.
There are two types of relationship..
There should have jealousy..
It means exclusiveness..
And one of the relationships is marriage, right?.
In marriage, we should have jealousy.
because marriage is a one man and one woman.
making covenant and being faithful to each other only, right?.
So the first relationship that should have jealousy.
is marriage..
And the other one is here..
It's your relationship with God.
because God is a true God..
And you have a relationship with God,.
has a spiritual relationship with Him..
Your relationship and how God see His relationship with you.
is also in terms of jealousy.
because you cannot worship other gods..

$^{561}$You cannot worship other idols..
You cannot have someone who has the first place.
in your life other than God Himself..
That's why the commandment two, principle two says,.
"In your life, there should be no idols.".
Now, idols can be anything,.
whatever, whoever occupies your heart.
more than God Himself..
So whatever occupies your heart, your mind,.
more than God is an idol..
Whoever occupies your heart more than God is also an idol..
So idols can be your career..
It can be your money..
It can be your family..
It can also be your children or your grandchildren.
or some other human relationships..
No one and nothing should take the place of God..
God is number one..
If anything that occupies the number one place is an idol..
And so often we try to use our human effort.
to take control of our own life.
instead of yielding our worship to God Himself..
To recognize Him as the ultimate Lord, right?.
So number two says, you know, God hates idols..
God hates idols..
Now quickly, we go to number three..
You put God as your number one..
No other God, no other gods before the true God..
And no idols will occupy or compete.
with the lordship of the true God..
Number three, it says,.
you shall not misuse the name of God, the Lord your God,.
for the Lord will not hold anyone guiltless.
who misuses His name..
Now we usually take this verse.
as meaning you should not curse, right?.
You should not say the name of God,.
Yahweh, Jesus Christ, you know, in vain, right?.
Is it the real meaning of this statement?.
Well, I would say no,.

$^{601}$because in the ancient world,.
name, name is not merely a label..
It's not merely a label like, okay, I call you John, okay?.
And you're Peter, so I'm calling you John..
John is unique..
Name is not just a label in the ancient world..
Name is power..
Name is authority..
Name is reputation..
So statement number three here,.
commandment number three, has to do with God's power..
So what do you deal with, how do you deal with God's power?.
The name of the Lord your God..
Is that you shall not, thou shall not misuse.
the name of the Lord your God..
Now, a little bit footnote here for the verb misuse..
As I said, it's not so much of, you know,.
saying the term God or using the name of Jesus Christ.
in our statement..
The verb misuse, translated misuse,.
actually means carry or lift up..
It's like you pick up your backpack, okay?.
You pick up your backpack..
That is the Hebrew word behind the English word misuse..
So literally, the verse can be translated as shown here..
You shall not carry the name of God in vain..
And what it means is that you cannot.
carry the brand of God in vain..
Let me give you an example..
Let's say you own a McDonald's franchise, all right?.
You know McDonald's runs with franchise, right?.
And you apply for it, and then you own the franchise,.
and there's a great opportunity for you to earn money,.
of course, but, you know, come with that franchise.
is that it requires you to make significant investment.
in managing the business..
It's McDonald, right?.
It's a McDonald..
And it requires you a high level of commitment,.
and you must also follow the corporation's standards,.

$^{641}$like food quality, right?.
If you get one star or five star, then, you know, you're out..
And you better keep up with the standards.
set by the corporate, set by the name, McDonald's..
And you have to carry the good name of McDonald's.
when you own the franchise, when you manage the franchise..
Now, with that example, we understand what this verse means..
You shall not carry the name of God in vain..
What it means is much the same way.
that since you are saved by grace,.
you're owning God's franchise,.
that in your life, you have the franchise of God,.
the true God, and you have to do everything.
to honor the name so that when people see you,.
when people look at you and see how you manage your time,.
manage your money, manage your career,.
manage everything in your life,.
they will see that, "Oh, wow, that is a good brand..
"That is a good name..
"I want to join that franchise as well," right?.
And we must admit and confess that sometimes.
we just not very up to the standard, right?.
We just thought that we are saved by grace.
so we can do anything we want..
That's not how it works, guys..
That's not how it works..
No, you can lose your salvation..
You can lose your identity.
if you do not honor the name, the franchise of God..
And what a challenge to all of us at the end of the year..
As we reveal the past year, what have we done?.
And what kind of effort we put into making.
or carrying the name of God, right?.
And people see us..
The spouse see us..
The children see us..
We see that, you know what?.
God is real..
When I see your life, your children will say,.
or your friends will say, "Oh, man,.

$^{681}$"she is different..
"He is different, and I wonder why.".
And you'll be able to tell them, "You know why?.
"Because I'm carrying the name of God..
"My life is the brand of God..
"My life is the reputation of God.".
So we are saved by grace,.
and we come to God with new principles..
We, because we are managing His franchise..
Are we doing everything to honor the name of God?.
Are we doing everything in terms of managing our time,.
in terms of managing our schedule,.
in terms of managing our priorities?.
I know it probably has been challenging for you.
in the past year..
However, you know, because we're saved by grace,.
and that is the foundation of all..
We are saved by sheer grace of God,.
and my life is carrying His name,.
and His name is precious..
His name is to be honored by all..
Now, as we close, why don't we have some moment.
of reflection, and we come before God and say,.
you know, probably we have to confess for our shortcomings,.
confess our weaknesses,.
or sometimes, you know, kind of outside our mind for us..
We just follow the motions as a Christian,.
and we are not really worshiping God..
We're not really serving Him, right?.
We treat it as a job..
God said, "You shall not misuse..
"You shall not miscarry the name of God in vain.".
Why don't we come before God and have a moment of reflection?.
And let us pray together..
Father, as we close this year, 2024,.
we count how many opportunities that You have given to us,.
the provision that You have provided for us..
And we come before You with thanksgiving..
Thank You for granting us all the good things,.
for all the good things come from our Father in heaven..

$^{721}$And we wanna give You thanks and give You, say thank You..
And not just that, we come before You.
not only with a thankful heart,.
but a renewed mindset,.
knowing that we are Your people, we are Your children..
And God, help us dig into the bottom of the heart,.
that are we really honoring You?.
Or are we making excuses in our life?.
Trying to run away from You,.
even though we know for sure that You love us so much..
Father, we confess our sins..
We confess our weaknesses.
or our tendency to, we're so prone to wander,.
wander away from Your truth,.
wander away from Your grace..
So yes, we come before You,.
know that You are the God of all grace..
Thank You for the forgiveness.
You have provided for us in Christ Jesus..
Forgive our sins and cleanse us..
And renew our mind and our heart.
so that our heart will belong to You.
and belong to You alone..
You are the true God, You're the only God..
And thank You for creating us,.
redeeming us and guiding us..
So help us to make a commitment before You,.
knowing that because You have been faithful to us..
And that's why we live by Your principles,.
to be Your new people,.
to be people that carry Your franchise,.
carry Your name, carry Your reputation,.
that our lives will attract other people.
to come follow You..
Help us to walk in the way of godliness,.
the way of holiness and purity..
Help us God..
And we know that You are full of grace and kindness..
So help us to live in Your love,.
live in Your presence all time..

$^{761}$And help us to abide to Your wonderful principles.
so that our lives will be different..
Thank You for calling us to be Your church..
And thank You for being with all of us,.
even though in the past year or so,.
our church has a lot of different issues and challenges.
and many changes,.
but we know that You are still faithful to us..
So in the head of us, the new year,.
we recognize and acknowledge that You are still faithful..
So help us to be faithful..
In Christ, we pray..
God bless you all and hope to see you next year,.
next Sunday, and continue to worship our one and true God..
Right?.
(gentle music).
\newpage



\section{路加福音}
\label{sec:3GTzRC3gFVM}
\textbf{聖誕的平安 - 路加福音 2 (陳崇基牧師)  Peace of Christmas - Luke 2 (Rev. Mark Chan)}
\newline
\newline
連結: \href{https://youtube.com/watch?v=3GTzRC3gFVM}{\texttt{https://youtube.com/watch?v=3GTzRC3gFVM}} ~~~~ 語音日期: 2025-01-04
\newline
\newline
\hyperref[sec:UlWyJIvtS2w]{\small{< < < PREV SERMON < < <}}
~
\hyperref[sec:index_chronic]{\small{[返順時目]}}
~
\hyperref[sec:index_scriptual]{\small{[返順卷目]}}
~
\hyperref[sec:Lw531TtsFSc]{\small{> > > NEXT SERMON > > >}}
\newline
\newline
$^{1}$On this Christmas Sunday, I'm sure if you reveal what has happened to you this year,.
it's going to be probably a very busy day for you, very busy year for you..
[Chinese].
Perhaps this year has been a very, very tough year for you..
[Cantonese].
For many of us at UV, it has been a challenging year..
You know, we have families dealing with sickness, dealing with death and crisis and difficulties,.
and our ministries are also tackling with different challenges, and in general,.
I think there are so many stresses in our life..
[Cantonese].
So I pray that the songs that we offer to you, the song that we are going to sing,.
the Christmas songs, the baptism testimony, the candle lighting ceremony that we're going to do.
later, will all remind you the real meaning of Christmas..
[Cantonese].
Let me read this passage on screen..
It is taken from the Bible, the Gospel of Luke, from chapter 2..
[Cantonese].
Here, the angelic choir is announcing the good news, the good news of the birth of our Savior,.
to the shepherds in the field in Bethlehem, the very good news of the birth of our Savior,.
Jesus Christ..
[Cantonese].
Suddenly, the angel was joined by the vast host of other angels, the armies of heaven..
They praised God and they say, "Glory to God in highest heaven,.
peace on earth to those with whom God is pleased.".
[Cantonese].
Now, one of the Christmas traditions many Americans would do is to watch some movie,.
some Christmas movie..
[Cantonese].
And I don't know if you do that or not..
[Cantonese].
But, you know, many of the families, American families, will be watching.
Charlie Brown Christmas this season..
[Cantonese].
How many of you have watched this before?.
Charlie Brown Christmas?.
Raise your hand..
Only a handful, okay?.
Only a handful..
[Cantonese].
All right, so in this movie, Charlie Brown, the main character,.

$^{41}$is doing his best to organize a Christmas play with the kids..
[Cantonese].
But as you can tell, Charlie Brown is trying to organize this play with other kids and.
kids are not listening to you, right?.
They run around and no one in the story is listening to Charlie Brown..
And everyone starts to criticize Charlie Brown for his, you know,.
lack of organization or lack of leadership..
So Charlie was very frustrated and nobody understands the real meaning of Christmas..
So at the end, Charlie Brown almost gave up..
[Cantonese].
And by the way, we have some kids joining us today, which is great, right?.
This is great..
They're noisy, they're running around, they're joyful and they want to express themselves..
So as the story goes, Charlie Brown was very frustrated..
[Cantonese].
However, at the end, as the story goes, with all the chaos with the kids,.
Charlie Brown has a very wise friend called Ninus..
So Ninus steps in and Ninus has always been a timid, shy, and, you know,.
less than confident character, less than confident kid in the story..
[Cantonese].
And that's why you, whenever you see Ninus in Charlie Brown stories or movies,.
Ninus has always has a security blanket in his hand..
Okay, ring the bell, remind of this character, always has a security blanket in his hand..
But, Ninus volunteered himself to tell the story of the Bible and the true meaning of Christmas.
under the spotlight..
[Cantonese].
Do you want to watch this movie?.
[Cantonese].
Okay, I'll let you watch two minutes..
[Cantonese].
Okay, so we're going to watch two minutes of the most, most spectacular moment in this entire movie..
[Cantonese].
What about Charlie Brown?.
Do you enjoy the movie?.
Two minutes..
[Cantonese].
Now, I want to tell you guys, there's one little detail, one little detail in this clip,.
most of the people missed..
[Cantonese].
And that one little detail is that right in the middle of Ninus telling the story of Christmas.

$^{81}$from the Bible, telling other people the real meaning of Christmas, a child is born..
You see, Ninus actually dropped his security blanket..
[Cantonese].
Now, I'm sure we all have our security blankets..
[Cantonese].
Things that give you peace, things that give you security, things that give you confidence..
[Cantonese].
Maybe our bank account, the money that we have, our friends, our career, our identity,.
or some sort of certainty in our future..
[Cantonese].
And we all try to grab that thing, you know, or some things that will give us everlasting peace,.
to give us a little bit more peace, a little bit more, you know, confidence..
But as you know, in reality, those things do not last..
They often disappoint us..
[Cantonese].
Maybe health, maybe our well-being, maybe our future, our security, whatever that is,.
your security blanket..
[Cantonese].
And I want to tell you, you know, they often disappoint us..
Nothing can be guaranteed..
[Cantonese].
Now, what Christmas is all about is a different kind of peace..
It's a different kind of peace..
It's a true, everlasting peace..
It's peace with God our Father..
[Cantonese].
Peace in times of chaos, peace in times of difficulty, peace in times of challenge,.
knowing that God is with us..
God is always with us in Jesus Christ..
[Cantonese].
And that is the everlasting peace..
[Cantonese].
The real meaning of Christmas..
Glory to God in the highest heaven and peace on earth to those with whom God is pleased..
[Cantonese].
And the people that God is pleased with are those who have put their faith in Christ Jesus..
[Cantonese].
So how can we receive the gift of God in Christ Jesus?.
[Cantonese].
Well, first thing we have to do is to let go..

$^{121}$If you want to receive a gift, you have to receive it with open hands..
But if we are holding on to our security blankets, our stuff, you know, our hands are not open.
and our hearts are not open, then we cannot receive the gift of God..
[Cantonese].
So how can you receive the gift of God?.
Only by your hands not holding on to your security blankets..
[Cantonese].
How can we experience the gift of God if our hearts are filled with overwhelm.
with our stuff and our idols, our work, our overwhelmingly busy schedule?.
We etch out God in our life..
The only way we can receive the gift of God is that we let go of this stuff..
And with an open heart, we come to God and say, "Thank you for giving me the very son of yours.".
[Cantonese].
So at the end, I want to invite you, all of you, to take this moment,.
take this very moment to reconnect, to reconnect with God..
God who created you, God who loves you, God who is always with you..
[Cantonese].
Now let me read this prayer first before we recite it together..
So I want to invite you to recite it with, you know, with your words or with your heart..
[Cantonese].
Dear Lord Jesus, the prayer says,.
I ask you to grant me peace..
Peace in my heart when I feel, when it feels like the world all around us spins out of control..
Help me stay focused on you..
The son of God, born in a humble manger..
I come to you, God, with an open heart..
I know my need for your grace and your forgiveness..
I ask you to be my savior and my Lord..
Cleanse me from my sins and make me new today..
Fill me with your peace and guide me in your way..
In Jesus' name..
[Cantonese].
Let's end today's service with the lighting ceremony..
We're going to pass the flame to the candles..
But we want to end and conclude with this prayer..
Well, let's do it in English first..
If you want to join me, go ahead, all right?.
But you can pray in your heart as well..
Let's do that..
Dear Lord Jesus, I ask you to grant me peace..

$^{161}$Peace in my heart when it feels like the world all around us spins out of control..
Help me stay focused on you..
The son of God, born in a humble manger..
I come to you with a humble heart..
I know my need for your grace and forgiveness..
I ask you to be my savior and Lord..
Cleanse me from my sins and make me new..
Fill me with your peace..
Guide me in your way..
In Jesus' name..
Amen..
[Cantonese].
Let's all pray together..
[Cantonese].
Lord Jesus, this is our prayer..
And our prayer is to make room for you..
[Cantonese].
And we invite Jesus to come into our lives to be our savior anew..
To be our Lord..
So cleanse us and renew us and guide us in your way..
In Jesus' name..
[Cantonese].
Amen..
\newpage



\section{路加福音 2:8-20}
\label{sec:Lw531TtsFSc}
\textbf{隱藏的聖誕 (路加福音 2\_8-20) - 譚文倩莊傳道}
\newline
\newline
連結: \href{https://youtube.com/watch?v=Lw531TtsFSc}{\texttt{https://youtube.com/watch?v=Lw531TtsFSc}} ~~~~ 語音日期: 2025-01-04
\newline
\newline
\hyperref[sec:3GTzRC3gFVM]{\small{< < < PREV SERMON < < <}}
~
\hyperref[sec:index_chronic]{\small{[返順時目]}}
~
\hyperref[sec:index_scriptual]{\small{[返順卷目]}}
~
\hyperref[sec:pzPtCyCLwfM]{\small{> > > NEXT SERMON > > >}}
\newline
\newline
路加福音 2:8-20
\newline
\begin{longtable}{cl}
\hline
\hline
章節 & 經文 (和合本修訂版)\\
\hline
2:8 & \begin{tabularx}{0.7\textwidth}{X} 在伯利恆的野外有牧羊人，夜間值班看守羊群。 \end{tabularx} \\ \\ \relax
2:9 & \begin{tabularx}{0.7\textwidth}{X} 有主的一個使者站在他們旁邊，主的榮光四面照著他們，牧羊人就很懼怕。 \end{tabularx} \\ \\ \relax
2:10 & \begin{tabularx}{0.7\textwidth}{X} 那天使對他們說：「不要懼怕！看哪！因為我報給你們大喜的信息，是關乎萬民的： \end{tabularx} \\ \\ \relax
2:11 & \begin{tabularx}{0.7\textwidth}{X} 因今天在大衛的城裡，為你們生了救主，就是主基督。 \end{tabularx} \\ \\ \relax
2:12 & \begin{tabularx}{0.7\textwidth}{X} 你們要看見一個嬰孩，包著布，臥在馬槽裡，那就是給你們的記號。」 \end{tabularx} \\ \\ \relax
2:13 & \begin{tabularx}{0.7\textwidth}{X} 忽然，有一大隊天兵同那天使讚美神說： \end{tabularx} \\ \\ \relax
2:14 & \begin{tabularx}{0.7\textwidth}{X} 「在至高之處榮耀歸與神！在地上平安歸與他所喜悅的人！」 \end{tabularx} \\ \\ \relax
2:15 & \begin{tabularx}{0.7\textwidth}{X} 眾天使離開他們，升天去了。牧羊人彼此說：「我們往伯利恆去，看看所成的事，就是主所告訴我們的。」 \end{tabularx} \\ \\ \relax
2:16 & \begin{tabularx}{0.7\textwidth}{X} 他們急忙去了，找到馬利亞和約瑟，還有那嬰孩臥在馬槽裡。 \end{tabularx} \\ \\ \relax
2:17 & \begin{tabularx}{0.7\textwidth}{X} 他們看見，就把天使論這孩子的話傳開了。 \end{tabularx} \\ \\ \relax
2:18 & \begin{tabularx}{0.7\textwidth}{X} 聽見的人都詫異牧羊人對他們所說的話。 \end{tabularx} \\ \\ \relax
2:19 & \begin{tabularx}{0.7\textwidth}{X} 馬利亞卻把這一切的事存在心裡，反覆思考。 \end{tabularx} \\ \\ \relax
2:20 & \begin{tabularx}{0.7\textwidth}{X} 牧羊人回去了，因所聽見所看見的一切事，正如天使向他們所說的，就歸榮耀於神，讚美他。 \end{tabularx} \\ \\
[1ex]
\hline
\hline
\end{longtable}
$^{1}$打開櫃子後.
我發現裡面竟然有一包巧克力.
我當時在想.
媽媽,你真是狡猾.
藏起了一包巧克力不讓我們知道.
我是一個很疼弟弟的姐姐.
拿著一包巧克力.
我把巧克力放在櫃子里.
我看到櫃子里有一個巧克力.
我拿著這包巧克力.
我把它分了一半.
給了一半給我弟弟.
小時候哥哥經常欺負我們.
所以我沒有給哥哥.
那包巧克力很好吃.
我和弟弟就久久地吃掉了.
但吃了不久.
我兩個就開始很肚痛.
原來是巧克力射糖.
很不妥.
那天晚上.
你知道香港也有個廁所.
大家都不用用廁所.
因為被我們兩個姐弟霸佔了.
就好像我很意外地發現這包巧克力一樣.
今天希望我們嘗試.
在這個我們已經聽過無數次的故事里.
耶穌基督的出生里.
我們看看隱藏的功課.
我們先低頭祈禱.
主啊願禰的名被高舉.
我們人有受教的心.
祈禱奉主耶穌基督成名祈求.
Amen.
上兩個星期Mark Sir講到.
《路加福音》的二章一至七節.
講到這個.
「哀煞奧古士督新官上任三把火」.
他要做人口統計.
其實是方便他自己稅收.

$^{41}$要求全世界回到自己的地方去登記.
約瑟帶著這個操身大勢的瑪利亞回到伯利恆.
應驗了《馬太福音》二章的預言.
瑪利亞到報之後不久.
耶穌就在那裡出生.
接下來就是今天的經文.
我嘗試很快地跟大家讀一次.
《路加福音》二章八至二十節.
「在伯利恆的郊外.
有一些牧人在夜間看守羊群.
主的一位使者站在他們旁邊.
主的榮光四面照著他們.
他們就非常害怕.
天使說不要怕.
看哪我報給你們大喜的信息.
是關於萬民的.
今日在大位的城裡為你們生了救主.
就是主基督.
你們要找到一個嬰孩.
包著布.
臥在馬槽里.
那就是記號了.
忽然有一大隊天兵.
和那天使一起贊美神說.
在至高之處.
榮耀歸於神.
在地上平安歸於他所喜悅的人.
終天使離開他們升天去了.
那些牧人彼此說.
我們往伯利恆去.
看看主所指示我們已經成就的事.
他們急忙去了.
找到瑪利亞.
若失和那臥在馬槽里的嬰孩.
他們見過以後.
就把天使對他們論者孩子的話傳開了.
聽見的人都唏其牧人所說的事.
瑪利亞把這一切放在心裡.
反復思想.
牧人就因為聽見的和看見的.

$^{81}$正像天使對他們所說的一樣.
就回去了.
把榮耀贊美歸於神.
首先你們知不知道這一段經文.
《路加福音》二章八至二十節.
在聖經里是唯一一個地方.
記錄了天使顯現給牧羊人的故事.
我之前不知道.
如果你們比我更瞭解就好了.
我們先來看看故事里的牧羊人.
《路加福音》二章八至九節.
在伯利恆的郊外.
有一些牧羊人在夜間看守羊群.
主的一位使者站在他們旁邊.
主的榮光四面照著他們.
他們就非常害怕.
我們不妨想象一下.
如果今時今日你去露營.
外面很黑.
突然有一盞大光燈.
光著你.
還似乎有一個像天使一樣的物體.
出現跟你說話.
你猜你有什麼反應.
再想想這個場景.
月黑風高.
一片看不到盡頭的曠野.
其實除了聽到羊的咩咩聲之外.
突然有一個天使出現.
再周圍有一盞大光燈照著你.
雖然我在想.
牧羊人應該見過不少大場面.
因為他們也要保護羊群.
不可以讓獅子猛獸去吃掉羊群.
但難免你會想.
其實那個情景應該是挺可怕的.
根據聖經的記載.
牧羊人的反應是非常害怕.
不是害怕.
是非常害怕.

$^{121}$不僅害怕.
是非常害怕.
什麼是恐懼呢.
恐懼是一種基本很強烈的情感.
是我們的大腦感覺到一些威脅.
一些威脅正在來而產生.
我們身體也同時會有一些反應.
例如我們的心跳加速.
或者我們全身會緊緊的.
我想起我大學時代的時候.
我很喜歡玩Rock Climbing.
在室內攀石.
其實攀石的好玩之處.
並不是因為像爬樹一樣爬到頂.
而是你要思考那幅牆.
那幅牆的特徵.
它的裂縫和傾斜度.
去手腳並用加上身體去爬到頂.
這個運動是很考技巧.
思考力度和耐力.
很好玩.
那時候生完女兒不久.
我又開始心痕.
覺得是時候可以回到健身房玩攀石了.
既然這麼久沒有玩.
我就找一幅比較容易的牆去做溫習.
發覺生完孩子是果然不同的.
我以前覺得挺容易的.
發覺自己好像很不敏捷,很笨拙.
但終於爬到頂.
頂大概是五層樓高.
去到頂之後我就看一看.
怎料看一看的時候突然很害怕.
以前我玩.
我從來沒有有過這樣的感覺.
當時我在想.
如果我摔下來.
女兒怎麼辦.
然後很安全地下地.
那次之後.

$^{161}$其實我已經21年沒有再玩攀石了.
這個情景令我想起當時的心跳加速.
很害怕.
整個人都緊緊的.
我想想這個牧羊人.
當時面對這麼震撼的情景的時候.
可能他們都有差不多的感覺.
在第一世紀猶太的背景當中.
天使的顯現,意象和神聖的跡象.
是時時都有出現的.
以色列人有長久和超自然相遇的經歷.
例如天使顯現給亞伯拉罕.
撒迦利亞或者瑪利亞等人的故事.
和舊約裡面神榮耀的顯現.
在聖經裡面.
當人遇見天使或者經歷神啟示的時候.
他們很多時候都會很表現出驚嚇.
這個恐懼帶有震撼和敬畏的意思.
因為天使的顯現代表著超自然的事情.
和神的計劃.
《奴家福音》這一部大喜訊的記載.
都提到救主降生之後人的反應.
今時今日值得我們思考的是.
今天你再聽到耶穌基督降生主的福音.
救世主已經來到的時候.
我們有什麼反應呢?.
還是已經知道太多了.
聽了很多次所以已經沒有什麼反應.
這個沒有什麼反應的反應代表著什麼呢?.
在《奴家福音》這一章.
表面來看牧羊人是因為天使的出現.
和被這盞大光燈照著而嚇到.
但其實他們在這個恐懼的反應.
部分原因是因為在這個宣告里.
天使用來形容耶穌的名詞.
所用的詞彙.
是當時羅馬皇帝所用的名詞.
我不知道剛才會賣廣告.
所以我預備了一個廣告.
最近在YouTube我們的YouTube頻道.

$^{201}$有個新節目叫做《幫你三分鐘》.
當中有耶穌的名字更詳盡的解釋.
幫助我們更加明白這一段經文的重要細節.
經文的背景.
如果你還沒有看的話.
我又像Carmen一樣再說一次.
記得回家要看看.
在《奴家福音》二章一節提到這個奧古士都.
即是當時羅馬的皇帝.
他自稱為神.
救主.
主.
他的出身被形容將要帶來世界和平的福音.
一道耀眼的光輝.
榮光.
但《奴家》所描繪的耶穌.
竟然比這個世界上最偉大的人更加偉大.
更加令人震撼的是這個神.
救主.
主.
竟然這麼卑微地在一個馬槽里降生.
《奴家福音》二章十至十四節.
天使說:不要怕!看那!.
我報給你們大喜的信息.
是關於萬民的.
今日在大位的城裡.
為你們生了救主.
就是主基督.
你們要找到一個英鞋包著布.
臥在馬槽里.
那就是記號了.
忽然有一大隊天兵.
同那天使一起贊美神說.
在至高之處榮耀歸於神.
在地上平安歸於他所喜悅的人.
我們繼續保持著想象.
我們再加一點情節.
這裡是說天使叫牧羊人.
你不用怕.
但還是那句.

$^{241}$你可不可以想象.
當時的情景怎麼不怕呢.
這裡是說天使叫牧羊人.
不要怕.
是因為我有一個.
很好的,大好的信息.
要告訴你們.
就是在大位的城裡.
那個真真正正來到這個世界.
能夠拯救我們的救主已經來到.
你們要起行去找這個.
剛剛出生.
包著布的寶寶.
這個就是真正的救主.
主耶穌基督.
接下來還有很多天使出現.
我在想的時候.
比我應該是已經害怕.
好像怕得你不夠害怕.
更害怕.
他們一起開聲贊美神.
在至高之處榮耀歸於神.
在地上平安歸於他所喜悅的人.
《顱吉福音2:15-16》.
眾天使離開他們升天去了.
那些牧人彼此說.
我們往伯利行去.
漢按主所指示我們已經成就的事.
他們急忙去了.
走到瑪利亞.
若失和那我在馬槽里的英骸.
經文記載他們聽完天使的說話之後.
就一起商量.
然後就急忙起程去找耶穌.
是說立刻去起程.
根據APA.
美國心理學學會的定義.
通常我們人面對恐懼的時候.
我們都是很自然地有三個反應.
Fight, flight, freeze.

$^{281}$即是對抗.
和你死過.
或者二.
我們會逃跑.
三十六著.
守衛上將.
三就是站定.
不是因為頂過台游.
而是嚇得不懂得動.
但是牧羊人.
竟然他們的反應.
不是很正常.
也不是這樣.
2:17說.
他們見過以後.
就把天使對他們論者孩子的話傳開去.
他們是照著天使所說的話.
真的立刻去找耶穌.
不只是站著.
不知道怎麼做.
我在想.
在預備這一講章的時候.
我在想.
無論我怎麼嘗試去跟你們說去形容也好.
和你們解釋的細節也好.
如果我們不花時間好好地去思考這段經文.
耶穌的出生的時候.
其實這段經文的震撼.
如何影響我們基督徒的生命.
你是不會理解的.
你一定要自己去明白.
去花時間去祈禱.
去讀這段經文.
神學家Dale Brack說過.
理論上我們基督徒應該和牧羊人的反應差不多.
就是我們會帶著感恩.
興奮.
贊美.
去回應這個這麼震撼的信息.
但是因為耶穌基督這個故事.

$^{321}$已經被人廣泛宣傳了二千多年.
聽過很多次.
所以基督徒對神.
藉著主耶穌基督的降生.
介入我們生命的這件事.
我們已經沒有什麼反應.
這正是《路加福音》一至二章里.
提醒我們的信心.
信心是我們對神所做的仍然有那份驚奇.
對神所能夠做的事.
我們持續地抱著一個很等待.
很感興趣地去知道的態度.
或者我們也需要恢復這種感覺.
或者問下自己.
特別是我們已經信主一段很長時間的人.
我們是否將一些我們已經聽過.
見過.
經歷過的都視為閒事.
又是那些我都完全懂得的心態.
用輕率的態度去面對我們的生命.
神的話.
上帝的作為.
幾年前我看過一個訪問.
我上網找也找不回來.
那個訪問說有一個女人.
剛剛接受了一個剛剛離世的人的肝臟移植.
那個訪問是在分享她自己.
很感激這個素未謀面的人的捐贈.
因為他的捐贈讓她的生命得以延續.
她說我的生命從此之後不再屬於我自己.
以後我要用行動去回報這份愛和恩情.
她的分享令我想起我們基督徒的生命.
其實我們也不就是這樣.
我們應該有這樣的感恩在裡面.
用我們的行動和生命去回應神給我們的這份恩情.
但我們又是否這樣做呢?.
在《斷經文》里的另一個人.
瑪利亞.
《路加福音2:18-20》.
聽見的人都唏其穆揚人所說的.

$^{361}$瑪利亞把者一切放在心裡反復思想.
穆人因為聽見和看見.
正像天使對他們所說的一樣就回去了.
把榮耀贊美歸於神.
我在想瑪利亞為什麼可以如此氣定神坦.
你只是看經文.
她好像沒什麼動過又沒什麼說過.
聖經的形容就是.
她存在心裡反復思想.
這個放在存在這個詞.
在新約里只出現過三次.
在原文里有一保存.
長久不會破爛的意思.
二保護抵擋的意思.
三有對話的意思.
即使瑪利亞仍然有很多不明白的地方.
例如後來2:18說.
聽見的人都愕然.
2:33.
孩子的父母因論耶穌的這些話就驚訝.
2:35.
你自己的心也要被劍刺透.
聖經里的信心.
往往是在我們不能完全明白的時候.
我們仍然選擇相信.
我們相信主的計劃會被成就.
我再清楚解釋一次.
是聖經裡面所說的信心.
很多時候是說我們人在困難,懷疑的裡面.
仍然選擇相信.
因為人相信神會照著他的時間,計劃去成就一切.
通常一個媽媽生完孩子之後.
除了我們的身體經歷了十級痛之外.
情緒上我們也會有很大的起伏.
聖經沒有記載瑪利亞.
生完耶穌之後.
他的身心靈狀況.
但在《路加福音》第一章那裡.
他有記錄將要成為耶穌媽媽瑪利亞的心志.
《路加福音》一章38節.

$^{401}$天使告訴瑪利亞.
她將成為以馬來尼的母親.
瑪利亞回答說.
我是主的侍女.
原照你的話成就在我身上.
反省我們自己的時候.
當我們面對挑戰困難.
我們會怎樣反應呢?.
你有沒有試過像瑪利亞那樣.
去祈禱.
原照你的話成就在我身上.
或者當神的呼召來到的時候.
我們又會不會像牧羊人那樣.
急不及待地去回應.
還是聽完之後我們會沒有什麼反應.
因為沒有什麼感覺.
還是只會害怕.
你不要找我.
很多擔心.
愁眉苦臉.
容許很多外來的聲音.
不同的情況去影響我們的腳步.
在預備講章的過程當中.
我通常都要找一些時間去放鬆我自己.
最近有位弟兄推崇了.
每個月9號9仙的Hulu.
我在想.
這麼便宜不簽名.
就是對不起自己.
Hulu裡面有一部電影.
《Alien》.
《異形》.
一部我小時候大部分時間.
都是遮著眼睛來看的電影.
現在回看我才知道.
原來這部電影是1979年出的.
你們當中可能有些人還沒出生.
我看到這部電影.
我就想這次我要鼓起勇氣.
在家裡天朦朧的時候.

$^{441}$又要我老公在家裡的時候.
再看一次.
希望這次我就不要那麼害怕.
不要遮著眼睛.
真真正正去探究一下.
這部電影其實講什麼.
那只怪物是什麼樣子的.
一路看我心裡想.
這部電影真是很討厭.
大部分時間都是黑漆漆的.
看看這部大部分黑漆漆的電影.
我突然想起路加福音裡面.
講到主的榮光.
突然在我腦里.
光與黑暗的對比很強烈.
《異形》這個外星人的電影裡面.
船員需要用大光燈去找.
去看清楚.
去消滅.
去瞭解究竟他們在面對什麼.
揭開這個怪物的面紗.
同時我自己在想耶穌基督的真光.
是不是正正是我生命裡面.
要去照亮我前面有什麼.
那盞燈呢.
不止一次.
我和大家分享過.
我小時候面對很多困難.
受過很多傷害.
前幾天我們三兄妹聊天.
講回我們小時候的經歷.
我們的感受的時候.
我們大家都是很感恩.
我們的小朋友是很乖的.
我們以前的經歷沒有影響到我們的下一代.
到後來晚上我再和我弟弟聊天的時候.
我們聊天的時候他發了個短信給我.
他說:姐姐.
我們應該很自豪.
但不能拿到信任.

$^{481}$信仰改變了我們.
改變了我們的一代.
我們應該為我們自己感到驕傲.
但我們又不能自傲.
因為事實是我們的信仰改變了我們.
也改變了我們這一代.
你信不信我們的信仰是可以改變我們.
可以改變我們的生命.
改變我們這一代.
甚至沒有改變我們的下一代.
耶穌是你和我都很需要的大光燈.
提供光明照清了.
引領前路.
正如約翰福音8章12節那裡這樣說.
我是世界的光.
跟隨我的.
就不再行在黑暗裡.
必得著生命的光.
還有兩天.
2024年就會過去.
思想《路加福音》第二章.
耶穌基督降生的大好信息的時候.
我在家裡重新思考.
耶穌是世界的光.
對我來說有什麼意思.
再花了一些時間去反思.
盤點.
並且重新調整我的事奉方向.
這個世界不停在變.
隨著年紀大了.
我們身體也有毛病.
越來越發覺自己有很多限制.
當中有很多自己都無法形容或者掌握的事情.
這些一切一切.
其實真的很容易影響我們所做的每一個決定.
我們的身心靈的健康.
今時今日人的精神健康.
已經成為一個我們不能忽視的議題.
因為不僅影響我們.
也影響我們的下一代.

$^{521}$時不時都會有人問我.
你猜世界是否快要末日了.
耶穌是否快要回來.
我的答案是耶穌從來都是快要回來的.
如果你去問一些住在伊拉克的難民.
可能他們會告訴你.
他們已經不止一次經歷過世界末日.
《路加福音2:10-11》.
那天使對他們說.
不要懼怕.
看啊.
因為我報給你們大喜的信息.
是關乎萬民的.
因今天在大位的成磊.
為你們生了救主.
就是主基督.
聖誕節的確過去了.
耶穌這位真正的君王.
真正的救主.
世界的光的降生.
帶來的平安.
我們聽過N次.
但是這個大好的信息.
能不能夠在逆境裡面成為你的盼望.
還能不能夠給你有震撼.
那個驚喜.
你仍然願意耶穌可以改變你的生命.
願意祂引導你.
扭轉你的心態.
迎接2025.
你們願不願意.
我們一起低頭.
主啊.
隱藏的聖誕.
是我們人的信心持續期盼.
主耶穌的光臨.
你經歷人所經歷的.
你忠於所托.
為罪人死.
你的出生.

$^{561}$完全全是因為我們.
我們這群罪人.
容許我們得著生命.
主啊求你賜我們.
持久的.
對你的期盼.
我們永遠關注你.
仍然期盼經歷你的作為.
嚮往耶穌的燈光.
引領我們的生命.
奉主耶穌基督的聖名祈求.
阿們.
崇拜結束.
願主耶穌基督的恩惠.
祂的平安.
與大家同在.
散會.
大家同在,散會.
\newpage



\section{}
\label{sec:pzPtCyCLwfM}
\textbf{The Wonder Of Christmas (Luke 2) - Rev. Mark Chan}
\newline
\newline
連結: \href{https://youtube.com/watch?v=pzPtCyCLwfM}{\texttt{https://youtube.com/watch?v=pzPtCyCLwfM}} ~~~~ 語音日期: 2025-01-05
\newline
\newline
\hyperref[sec:Lw531TtsFSc]{\small{< < < PREV SERMON < < <}}
~
\hyperref[sec:index_chronic]{\small{[返順時目]}}
~
\hyperref[sec:index_scriptual]{\small{[返順卷目]}}
~
\hyperref[sec:KuhthfKDgTw]{\small{> > > NEXT SERMON > > >}}
\newline
\newline
$^{1}$- Well, this is the best time to think about.
the Christmas story and especially the chapters.
that have recorded the details of the birth of Jesus..
And you will find those chapters.
in the first book of the New Testament,.
which is the Gospel of Matthew, Matthew 1 and 2..
And also the third book of the New Testament,.
which is Gospel of Luke, chapter 1 and 2..
And today we're going to focus on the first two chapters.
of the Gospel of Luke..
And it talks about the prophecy,.
the annunciation of the birth of Jesus.
and his pioneer, John the Baptist,.
and also the detail into the details of the birth of Christ..
So these two chapters are the so-called Christmas story..
Now, when we think about the story of Christmas,.
we always think about, you know, it's a kid's story..
It's a children's story, right?.
Yeah, you can find Luke 1 and 2.
in almost all of the children's Bible..
And it's for a good reason because it's easy to understand..
But as we think about this story,.
think about the details here,.
you find that it is not a story for children..
It is not a story for kids..
And perhaps this story, Christmas, is too familiar to us,.
but we have to read it more carefully..
And if we do that,.
we find that the Bible's Christmas story.
has so much hidden meanings.
and so much even political meanings..
And worse, that Luke, in the first two chapters,.
Luke 1 and 2,.
he used these words to tell the Christmas story..
And if you are the first readers of the Bible,.
you know, if you hop into a time machine.
and go back to the time of Jesus, the first century,.
if you hear the story, if you hear some of the words,.
you will find that these terms are loaded.
with religious and political meanings..

$^{41}$And here I give you the summary,.
and the summary will tell you that.
the story is not just for kids to read..
Here, it says, it used words like Lord, Savior,.
Light, All Peoples, Gospel, Peace, God, or Son of God..
Now, I highlight those words.
because not only they are used in the Gospel of Luke 1 and 2,.
as you can tell from the first line,.
the Lord is being used, Angel of the Lord,.
the glory of the Lord, Christ the Lord, Christ the Lord..
And then Savior, unto you is born a Savior..
And then Simeon, one of the old guys in the story.
in the first two chapters, when he saw the baby Jesus,.
he said he saw salvation, Savior or salvation..
And then Light, the word Light, what is so important,.
you know, it's used quite a few times.
in the first two chapters,.
Light to those in darkness, in the shadow of death,.
glory, light and glory,.
light to the nations, the Gentiles, glory to Israel..
All peoples, all the world to be registered..
Joseph and Mary are forced to be registered to pay taxes..
Good news for all people..
And then the salvation is for all peoples on earth..
And then the word Gospel, we heard so many, so many times.
and we believe in the Gospel, we proclaim the Gospel..
And it's translated good news or good news of great joy..
And then peace, peace to guide us in the way of peace, 179..
And then chapter two, verse 14, peace on earth..
And of course the term God or son of God.
show up so many times, I don't need to give you examples..
And what's so special about these terms is that.
these terms are political terms used by Augustus..
Augustus was the first emperor of the Roman empire..
You see the statue on the upper left..
He used a lot of these words to proclaim himself to be king..
And he proclaimed himself to be Lord and savior,.
the light to the nations,.
the light to all the peoples in the empire..
He's ruling all the peoples in the Roman empire..

$^{81}$He became king and it was the Gospel to the nations,.
the good news to the nations..
And Augustus claimed that his rule will bring upon peace,.
bring to all people the way of peace, pex in Latin..
And then he claimed himself to be God.
and used the title son of God in the coins.
and in his inscriptions..
Now, all of these words highlighted.
are often used by the Roman emperor,.
starting from Augustus..
In other words, if you watch the YouTube or TikTok.
or Netflix in those days,.
all the news channel in those days,.
you will be bombarded by the empire's political propaganda..
The political propaganda is telling you.
that you need to be convinced that who bring you peace.
and who bring you joy and who bring you the savior.
and the salvation and the light..
It was Augustus..
It was the emperor, the Roman emperor..
And you will be convinced.
if you watch those stories and propaganda,.
you have to put your faith in the Roman emperor.
because Augustus is God manifest..
Augustus is the son of God..
He is the bringer of peace..
Now, if you understand this background,.
you understand this crystal story.
is using those political terms to apply to Jesus,.
apply to the baby born in manger, humble and small..
You see that the Bible's crystal story.
counters and challenges those claims by Caesar,.
by the Roman emperor..
It says, no, it was not the Roman king who bring you peace..
It was not the emperor who bring you joy or security..
The Roman empire cannot do that..
The Roman empire ultimately is about domination..
It's about military conquest, about oppression..
That's what Roman empire is all about..
And today, if you read this story,.

$^{121}$Luke is telling you there are many Caesars on earth..
There are many kings and kingdoms and empires on earth.
are telling us the same gospels..
Things like you will be satisfied,.
you will have security if you have this or that..
If you believe in this or that,.
if you have this kind of relationship,.
you have acceptance from your peers,.
you have this world's measure of success.
or your technology and all that stuff..
The big idea is the gospel..
This is the gospel from the secular society,.
from today's world..
Now, it's telling you this story,.
it's telling you this concept.
that if you possess this or that,.
if you own this or that,.
you will be totally secure,.
you will find fulfillment, you will be happy..
That is the gospel of today..
And that was the gospel in Luke's day in the first century,.
and Jesus came..
The Christmas story is telling us.
we have to counter the claims of today's Caesars..
Only Jesus Christ, the Lord,.
brings us peace, true peace and security..
If you read Luke one and two as a whole,.
you find that, here's the Chinese version.
if you're interested..
If you look at chapter one and two of Luke,.
it tells you the entire Christmas story.
in two chapters, two long chapters..
You find that the flow of the chapter.
or the table of content or the outline of chapter.
is pretty clear..
After the introduction, there's two annunciations..
An annunciation means hearing the birth of someone..
All right?.
The two annunciation, the first annunciation is.
the birth of John the Baptist, right?.

$^{161}$The pioneer of Messiah..
And then the second annunciation was the birth of Jesus..
And then after that is the two mothers meeting together,.
chapter 139 to 45..
The two mothers is John the Baptist's mother, Elizabeth,.
and his husband, Zachariah..
They were very, very old..
They were without child,.
but God bring them from shame to honor.
and gave them a child, miracle child, John the Baptist..
And that's the first mother, Elizabeth..
And the second mother, of course, is Mary, right?.
Mary was pregnant by the Holy Spirit.
and she was carrying the child, Jesus..
And then the two mothers met together in 139 to 45..
John the Baptist was in the womb.
and he was lipping with joy, seeing the Savior's mother..
Even in the womb, he was looking forward.
to the arrival of the true Savior, Jesus Christ..
So the two mothers meeting together,.
and then Mary sang the song, the praise song, "Magnifica.".
And that is a very, very famous hymn..
If you've been to a Catholic church.
or you come from a Catholic tradition,.
"Magnifica" is one of the most significant psalm.
or hymn sung by Mary..
And then after this first praise song,.
is the birth of John the Baptist..
So it's recording the details of his birth..
And then after his birth, and he was named John,.
Zachariah, the father, burst out in praises..
And that is the second song in the gospel of Luke,.
"Benedictus.".
He's blessing God because of his miracle,.
because of his promise, his faithfulness.
to fulfill the salvation for the people..
And then chapter two is the birth of Jesus..
And then after the recording of the birth of Jesus,.
there was song from heaven, from the heavenly angels,.
the angelic army singing the song together..

$^{201}$And that is one of the most famous Christmas hymns,.
"Gloria in Excelsis," right?.
"Gloria in the highest.".
And then afterward, 222 to 240, Jesus was eight days old..
And from, according to Jewish custom,.
Joseph and Mary need to bring the child, eight day old,.
to the temple in Jerusalem.
for circumcision..
So the Lord is now in the temple,.
the true Lord, the true savior,.
although he's only eight years old,.
showing up in the temple..
And then there were different details..
They offered the sacrifice, Joseph and Mary..
And then two elderly people, Simeon and Anna,.
came to meet the baby Jesus..
And after Simeon saw with his own eyes,.
the eight day old savior,.
he burst out with another song, "Nu Dimittas,".
which means, "Let me depart,".
because he has been waiting, waiting, waiting.
for the salvation to come..
And finally, in the temple court, he saw the savior..
So he said, "Let me depart,.
"because I can now depart in peace.
"because the Lord is good..
"The Lord has been faithful.".
And then the chapter end with 12 year old Jesus,.
and he show up in the temple again, 12 year old..
And he was sitting there discussing the law,.
the law, the law of Moses with the teachers..
And he was sitting there,.
and that is the posture of rabbi..
12 year old, he was the rabbi..
And he was debating and asking questions.
and actually teaching the adults about the Mosaic law,.
about the law and the prophet and the Old Testament..
That's a pretty amazing chapter, chapter one and two..
But let's reveal the stories here..
And in these two chapters, who have met the Lord,.

$^{241}$who have met the savior,.
and experienced the true peace and true salvation?.
Who are these people?.
Now I mentioned here and there,.
you know, the old, right here,.
the old couple, Zachariah and Elizabeth,.
the parents of John the Baptist..
They were really really old and they were without child,.
but God is, God remembers them.
and grant them a miracle child, John the Baptist,.
and then raise them from shame to honor..
And then John the Baptist, while in the womb of the mother,.
he met Jesus in Mary's womb,.
and John the Baptist was leaping for joy.
as if I want to come out,.
I want to come out and see the coming of the salvation..
And then Joseph and Mary,.
Joseph and Mary, the holy family,.
they were powerless as well..
They were the bottom of the society.
under the Roman empire..
In the chapter two, the shepherds,.
watching the sheep in the fields of Bethlehem,.
they were also the outskirt of the society..
And then Simeon, an old man waiting for the peace of God,.
the salvation of God to come to reality..
And finally he saw it,.
and he saw baby Jesus eight days old,.
and he praises God and sings a song, "Let me depart.".
And then finally there's another widow, Anna..
So old widow living around the temple court.
to serve God for years and years and decades,.
and waiting to see the fulfillment of the salvation..
Now in these two chapters,.
those common people,.
not the kings, not the empire,.
not the 1\% of the population,.
but socially outcast, socially marginalized people,.
they have met Jesus because Jesus showed up to them..
Now I want to reflect on the last song..

$^{281}$And as I said, the last song means.
that Simeon, when he saw the salvation of God,.
and he burst out in praises..
And here Luke records the story, 2.25..
Now there was a man in Jerusalem called Simeon,.
who was righteous and devout,.
and he was waiting for the consolation of Israel..
And the Holy Spirit was upon him..
It had been revealed to him by the Holy Spirit.
that he would not die before he had seen the Lord's Christ..
Moved by the Spirit, he went into temple courts..
When the parents, Joseph and Mary,.
brought in child Jesus to do for him.
what the custom of the law require,.
which is the offering,.
in 28, Simeon took the baby in his arms.
and praised God, saying,.
"Sovereign Lord, as you have promised,.
"you now dismiss your servant in peace..
"For my eyes have seen your salvation,.
"which you have prepared in the sight of all people,.
"a light for revelation to the Gentiles,.
"and for glory to your people Israel.".
Now, this guy is an old guy, right?.
And he saw the salvation..
He's socially marginalized, but he met Jesus..
Now in the Roman Empire,.
the strong and the rich,.
they enjoy status and privilege..
No, they got the money, okay?.
They are the Elon Musk, right?.
They can do whatever they want, all right?.
That's how the Roman world works..
You have the power, you have the arms,.
and then you rule, and you have the privilege..
You have the salvation, you have the peace..
But the gospel, the Christmas gospel,.
is counter-empire..
It's against the empire,.
because it proclaims that the baby born in manger,.

$^{321}$humble and low,.
he instead is the true God..
He is the Son of God..
And the true good news is that Jesus is born in Bethlehem,.
one without power..
Even though he had power, he had all the glory,.
but he lay aside his glory and majesty,.
and come into our midst so that we can meet the true God..
We can experience the true salvation..
Now last Friday, I and another minister,.
we went to visit a very, very old lady..
She was dying..
So we went to Northwestern Hospital,.
and with all this obstacle we overcome..
And then we come to the room and saw this old lady..
And I've known this Christian lady for a long, long time,.
for decades..
And she has always been faithful,.
and always been supportive of my ministry..
So I heard about, she was not doing so well..
And especially the last five days,.
she has not been eating,.
and she has not been drinking for five or so days..
And she was very, very weak..
So we came and visited her..
And it's very interesting,.
whenever we visit someone by the sickbed,.
so often it was like the providence of God..
And as I told you, she was without food.
and without water for so many days..
And she was over 80 years old..
And, but when we visited her, she was very, very alert..
And it's kind of providence of God, right?.
In those hours, she was kind of, I don't know,.
maybe she sensed that someone's coming to visit her..
And interestingly, she was able to open her eyes.
and interact with me for a little bit..
So I talked to her and I said,.
"Oh, I'm Pastor Mark, do you recognize me?".
And she blinked her eyes and tried to nod..

$^{361}$And I read Psalm 23..
I read Psalm 23 with her..
And I saw her smile..
I saw her eyes, those weak,.
but radiated with hope and joy..
Now that short encounter,.
a visit of about maybe 15, 20 minutes,.
talked to her family members and talked to her..
I realized this is actually what this song is all about..
Someone who has been waiting for God's salvation..
Someone who has put her hope and trust.
in the true Son of God..
And finally, in the visitation,.
she kind of realized that the hope is in God.
and no one can take it away..
And I see from her eyes, I see from her reaction,.
she was filled with joy and faith..
And what a moment..
And I was brought to tears..
If our lives center on the promise of God,.
God will not fail us..
Yes, we have to be patient at times,.
patient for the work of God..
Sometimes we have to wait for days, weeks, months, and years..
But no one has put the faith in God will fail.
because God has always been faithful to us..
Now in this season of Christmas,.
I'm sure you are busy, right?.
I'm sure you're very, very busy..
There are some students here.
and you're probably writing your final papers, right?.
Some of us by the end of the year,.
probably very busy at work.
and with your family business.
and taking care of different people.
and offering hospitality to your family members.
or you're flying somewhere to visit your family..
People come in and out..
And Christmas has always been a season of challenge..
You know, we can say ho, ho, ho all the time, right?.

$^{401}$I don't know if you listen to the radio..
It's for old guy like me..
I listen to radio because when I drive,.
I pop in the radio to the station..
Since Thanksgiving, I've always listening to the same song.
playing in the radio all the time..
You know what song is that?.
It's the most wonderful time of the year..
It's always playing in the radio..
And you know, I must admit, you know,.
is it really the most wonderful time of the year?.
Well, not quite for me..
I mean, not so bad, but it's not the most wonderful time..
And I'm sure it's not the most wonderful time.
for many of us..
It may be the most difficult time for many of us..
Fill with sadness..
That's why in Christmas season,.
most of the people commit suicide because they're depressed..
They're depressed in the impression.
that the whole world is having a good time,.
but not you, right?.
And it gets you really, really depressive..
This season can be busy, can be very, very tough for you..
But I want to encourage you today as we close..
As we read the Christmas story,.
don't forget that you should invite.
the Son of God into your life..
The low and humble, true Lord and Savior..
He is the one we can put our trust in..
And He is the one when we invite Him into our lives.
to walk with us..
Remember I told you when I visit the patient,.
which Psalm did I read?.
Psalm 23, right?.
Psalm 23, okay?.
You don't need to recite any Psalm in the Bible.
except this one, okay?.
Psalm 23, okay?.
It's applicable in all occasions,.

$^{441}$especially when those people.
who are walking through some tough time..
Because in Psalm 23, right in the middle,.
it said, "Even though I walk through the shadow of death,.
"the valley of shadow of death,.
"you are with me.".
The shepherd is with you..
God is with you..
So as we close, why don't we pray together.
and invite Jesus to come into our lives again,.
to walk with us, no matter what we are experiencing..
Let's all pray..
Father, thank you..
Thank you that you sent your Son.
who laid aside His majesty, His glory.
to come into our midst,.
into our lives, into our hearts..
So as we close, as we reflect on.
the true story of Christmas,.
we want to open our hearts to invite you in..
And Lord Jesus, we want to put our trust in you..
And Lord Jesus, we want to invite you to come to our place..
And thank you for your grace.
that you don't mind how messy our life is..
You don't mind how broken we are..
Yet you came..
So we invite you to come in to be our Lord and our Savior.
and help us to open our eyes to see.
the true light that brings hope to all people,.
brings hope to us,.
can be a reality to our lives..
So again, we invite Jesus to come into our lives.
to be a leader, to guide us,.
to walk alongside with us..
Even though we may experience some tough times.
and challenging times,.
yet we know that you are with us..
Thank you for being our Lord and our Savior..
And again, help us to reflect on the true meaning of Jesus,.
to thank you, to give you praise,.

$^{481}$because you love us so much in Christ and pray..
Well, once again, thank you for coming..
And next week, as Kelly said, it's going to be 11..
Yes, all right..
And thanks again for Patrick and Kelly.
for preparing refreshment for us..
Please enjoy the refreshment.
and stay for a little bit time of fellowship together..
Okay, God bless..
\newpage



\section{路加福音}
\label{sec:KuhthfKDgTw}
\textbf{誰是你主?\#2 奧古斯都 (路加福音 2) - 陳崇基牧師}
\newline
\newline
連結: \href{https://youtube.com/watch?v=KuhthfKDgTw}{\texttt{https://youtube.com/watch?v=KuhthfKDgTw}} ~~~~ 語音日期: 2025-01-05
\newline
\newline
\hyperref[sec:pzPtCyCLwfM]{\small{< < < PREV SERMON < < <}}
~
\hyperref[sec:index_chronic]{\small{[返順時目]}}
~
\hyperref[sec:index_scriptual]{\small{[返順卷目]}}
~
\hyperref[sec:x5SxAthpTtg]{\small{> > > NEXT SERMON > > >}}
\newline
\newline
$^{1}$看完節目後我們繼續我們的系列.
這次到了聖誕節.
我們記得看過聖經裡面.
一些外邦的君王或者世俗的君王.
因為聖誕節的緣故.
所以我們看一個君王一定是聖誕節裡面提過的.
這是在路加福音第二章第一節開始.
他說在那些日子.
凱撒奧古斯都降旨.
叫全國人民都登記戶籍.
這第一次登記戶籍.
是在居里紐作敘利亞總督的時候行的.
眾人各歸各城辦理登記.
約瑟也從加利利的拿撒勒城上猶太地去.
到了大衛的城叫伯利恆.
他是大衛家族的人.
要和他所聘之妻瑪利亞同一登記戶籍.
他們在那裡的時候.
瑪利亞的產期到了.
就生了投胎的兒子用布包起來.
放在馬槽裡.
因為客店裡沒有地方.
這裡提到兩個君王.
兩個所謂的做官的.
第一個就是凱撒奧古斯都.
第二個就是居里紐.
居里紐只不過是省長.
敘利亞巴勒斯坦地的省長.
這次我們看左邊這個.
就是這個Augustus.
奧古斯都.
這個我們每一次到聖誕節的時候都會提的.
正如以前提過.
聖誕節就像小朋友的故事一樣.
Children's Bible一定會有的.
那些博士等等.
去到馬槽看到聖嬰等等.
這些畫像很令我們深印象.
但在路加或者聖經裡面記載的聖誕故事.
只有大概兩處地方.

$^{41}$路加方案第一二章.
今天我們看的.
還有馬化方案第一二章.
如果你看清楚經文裡面說的話.
絕對不是小朋友的故事.
是一個很有政治含義.
而且是當時很混亂的一個時間.
耶穌降生.
所以為什麼我們需要了解一下.
路加要特別看這件事.
如果是說小朋友的故事.
他就不需要看第一節.
是不是.
不需要看第一節.
或者其他那時候歷史的.
或者社會的狀況.
第二章第一節.
那裡說到海薩·奧古斯多.
海薩·斯托.
只不過是一個title.
好像牧師或者總統.
是一個title.
我們要了解一下這個人.
他的王號叫做Augustus.
奧古斯多.
本身他叫做Octavian.
屋大維.
但我們在聖經裡面提到的就是Augustus.
我們需要了解一下小小歷史.
因為不是一個小朋友的故事.
是一個滿有當時的背景.
對當時的人來說很重要.
一聽的時候他們就知道.
了解一下這個奧古斯多.
看看這個時間圖.
這裡有一個世界的歷史.
最普遍現在的分法就是主前和主後.
看到嗎.
Before Christ和AD.
Anno Domini.

$^{81}$就是主後的意思.
Year of the Lord.
這個是西方或者我們現代社會.
很多時候將時間這樣分的.
在左邊數目一直加大.
就是主前.
我們經常說舊約就是主前.
多少年多少年一直加大.
然後在這個所謂零之後.
當然沒有零年.
你明白我的意思嗎.
隨即的時候不會說.
啊 零年到了.
我們過了這個1BC.
然後零年到了.
沒有的.
是1AD所謂.
然後一直的數目就會加大.
這個就是Before Christ和AD的意思.
所以這種分法在主後第五世紀開始.
已經開始有人這樣說.
還有之後就會越來越普遍.
正如我們今天.
我們正在用的歷法.
西歷的歷法.
就是用這個所謂BC和AD.
當然當時第五世紀的時候.
有一個修士去修正這個歷法的時候.
他的資料不是很舊.
所以有些東西是off了的.
你看到上面那裡就是.
耶穌基督的降生的年份.
應當是所謂BC的六年.
為什麼呢.
因為我們知道大希律是什麼時候死的.
大希律死的時候.
附近就有一些月蝕或者天象出現.
所以你算起來的時候.
就是公元前四年.
Hell the Great.

$^{121}$就是他死的.
所以耶穌降生大概是公元前第六年.
所以除了BC和AD的說法之外.
我們就有BCE和CE的說法.
因為比較準確一點.
特別是說到這一段時間.
就是叫做Before Common Era.
就是大家所共識的歷法.
公元前和公元後.
你看到耶穌基督降生的背景.
那個時間是什麼時間.
就是羅馬帝國的時候.
羅馬帝國是公元前27年至400多年.
這段時間就叫做羅馬帝國.
之前就叫做共和國.
是議會去管理和統治整個國度.
共和國的時間大概是500BCE至27.
BCE或者BC.
我不知道你明不明白.
大家可能了解這句話.
英雄做時勢還是時勢做英雄.
你覺得時勢做英雄還是英雄做時勢.
有時候都是時勢做英雄.
奧古斯都能夠成為第一任的羅馬帝國的皇帝.
都是時勢做英雄的.
這裡就很複雜了.
如果只是看下面左邊.
左邊發生什麼事呢.
就是共和國最後一個君主叫做Julius Caesar.
Julius Caesar是很出名的.
出名的就是他在主前44年的時候被刺殺.
就是議會開會的時候.
突然後面有人去刺殺他.
然後他就死了.
然後整個共和國或者羅馬的國度.
就陷入一個內戰的情況.
這個內戰的情況最後打來打去.
剩下兩個最大的.
第一個就是奧古斯都.
或者Octavian.

$^{161}$第二個跟他是正敵對手的.
就是叫做Mark Antony.
所以奧古斯和Mark Antony打生打死.
有很多的內戰.
直到有一個最決定性的戰爭出現.
那個就是主前31年Battle of Actium.
在希臘西面的一個海灣.
有一個很大的.
兩方面.
Octavian就是奧古斯.
和Mark Antony.
大家在海灣那裡有一個海戰.
一個Naval Battle.
然後誰贏了.
就是奧古斯都贏了.
然後Mark Antony就跑去埃及.
本身這兩個人是朋友.
誰知道發生了一些家族上的事情.
或者政治上的問題.
於是兩個人就成為了正敵.
本身Mark Antony和Augustus.
他們本身是聯盟.
這個Antony其實就是奧古斯都.
Augustus的姐夫.
就是我姐姐的丈夫.
Mark Antony和Augustus之間的關係.
大家很親.
本身是同盟的.
誰知道Mark Antony就誘了他太太.
就是誘了我姐姐.
那你很生氣嗎.
很生氣是吧.
Augustus就用了這個原因.
做一個政治的宣傳.
所以就去攻打Mark Antony.
他誘了他太太.
跑去找誰呢.
就找了一個埃及的皇后.
那個叫做Cleopatra VII.
就是埃及妖后七世.

$^{201}$就轟動全國.
當時你看頭條新聞的話.
就是這個頭條新聞.
就是Mark Antony和埃及妖后七世.
就搞上了.
然後第一任太太就不要了.
本身是兄弟如手足.
現在怎麼樣.
現在反而成為了仇敵.
這個Antony Mark Antony.
經過阿克興海戰輸了之後.
他就和他的情人去到埃及.
和Cleopatra VII.
大家就在埃及.
然後Augustus攻打過來的時候.
他們知道搞不定了.
於是Mark Antony和Cleopatra VII.
就一起自殺.
所以是可歌可泣的故事.
很可歌可泣.
就是為了愛情.
他不要一個政壇.
那些人這麼瘋狂.
所以人們就覺得.
為什麼Antony對這個女人這麼死心塌地.
一定是所謂的埃及妖后.
非常妖艷.
非常漂亮.
因為迷人迷到一個地步.
愛情魔力可以吸引Mark Antony的靈魂出來.
然後和他一起死.
老一代的朋友.
你記得在1963年的時候.
荷里活有一套猛片.
記不記得.
記得.
叫做什麼.
叫做埃及妖后.
沒錯.
今天經常有人說.

$^{241}$有什麼猛片.
票房很厲害.
經常說看破地獄.
其實當時這套戲叫做Cleopatra.
是非常賣座的.
當時是58個million的box office.
58個million.
當時.
算現在.
是差不多500個million的box office.
那找誰飾演埃及妖后.
當然找最漂亮的.
當時最漂亮的.
Elizabeth Taylor.
記得吧.
你記不記得這套戲.
你上YouTube還可以看.
是吧.
不短的這套戲也挺長的.
找了Elizabeth太久.
來飾演埃及妖后.
她漂亮嗎.
不知道吧.
她很漂亮.
順帶一提.
說起當時的女神.
當時她有很多名車.
其中一部名車就是1961年的Rolls Royce.
Silver Cloud II.
就是右下角那個.
為什麼要說名車.
這個名車現在市值200萬.
我們進來不是看到很多漂亮的車嗎.
不急.
就是Elizabeth Taylor的那部名車.
為什麼要說名車.
因為美女一定要說名車.
所以要說名車.
Anyway我們回到正題.
正題就是奧古斯都成為霸主.

$^{281}$她在主前27年的時候.
先說正題.
不看女神了.
主前27年的時候.
奧古斯都成為天下的霸主.
本身她叫做奧大維.
現在取王號變成奧古斯都.
所以這個角度是因為靠強權和軍力.
靠政治的手腕.
如果看當時的YouTube,TikTok等等.
所有的宣傳都是在說這個國王.
這個羅馬的皇帝.
Augustus the Emperor.
是多麼厲害.
然後有很多的錢幣.
很多的雕像.
都是歌頌他的人生.
歌頌他的威水史.
甚至有一些當時的詩歌.
或者那些國歌.
都會歌頌羅馬的皇帝.
其中一個最著名的雕像.
就是這個Augustus of Prima Porta.
Prima Porta是羅馬城北面一個小的市鎮.
那裡找到一個這樣的雕像.
你看到這個雕像.
心口上寫著很多政治上的宣傳.
譬如有不同的神明.
因為他是神的兒子等等.
如果你翻查一下一些文獻.
說到這個Augustus的時候.
你會發現他很多時候.
自己有一些的尊稱.
就叫自己做主的.
Lord.
叫自己做救主.
Savior.
說自己是光或者太陽.
他所帶來的是For all peoples.
萬民的.

$^{321}$普世的.
這個君王登基的時候.
是一個的Gospel.
是一個的福音.
帶來地上的平安.
Pax.
拉丁文Pax.
或者他的稱號就是神.
或者神的兒子.
這裡不解釋太多.
如果你想知道更多.
阿邦幫我們錄了幾條片.
你在城會網台就可以看到.
這裡先買一買廣告.
但我想舉一句.
當時有一個哥哥叫做.
Kalman Sekulare.
Kalman Sekulare是什麼意思.
Kalman就是喪的意思.
Kalman這個字其實是喪的意思.
就是詩歌的意思.
Sekulare就是age的意思.
或者時間歷史的大段落.
就是世紀之歌.
當時有一個朝廷的詩人.
叫做賀拉斯.
Horace.
他大概是公元前17年的時候.
寫這首世紀之歌.
裡面的一些歌詞說什麼呢.
他說因為奧古斯都.
成為了羅馬天下的霸主.
所以就有信和平和尊榮.
回到以前的good old days.
還有他說過去的偉人裡面.
誰能夠將這位奧古斯都值得讚頌呢.
發光的他.
治理世界.
光芒萬丈.
然後亞波羅神就是光輝的神.

$^{361}$賜予世界新的光輝.
然後恢復這個世界的和平.
使萬象在春日的時候.
能夠有安穩.
看這些東西我已經提過很多次.
看這些東西其實就像中央電視台.
中央電視台經常播放這些信息.
就是說這個奧古斯都是多麼的了不起.
如果你有這樣的背景去看的時候.
你就會發現路加福音說的聖誕故事.
絕對不是一些小孩子說的故事.
我在這裡和你總結一下.
路加福音第一二章說到聖誕的故事.
耶穌降生的時候.
到底用了什麼字眼.
他說主的榮光.
為你們誕生的是主基督.
Lord Christ.
為你們生下救主.
Savior.
你們看見神的救恩.
有光照耀.
撒迦利亞的一首詩歌.
說到施洗約翰.
彌賽亞的先鋒.
出生的時候.
會見到清晨的日光.
照亮死刃裡面的人.
也有榮光照耀.
是外邦人的光.
是以色列的榮耀.
也有萬民.
萬民.
全國的人民.
眾人所有的人.
All people.
各歸各城.
大喜的信息關乎萬民.
救恩是為萬民所預備的.
二章31節.

$^{401}$也有福音.
天使傳的好訊息.
Good tiding.
大喜的訊息.
福音.
地上便有平安.
Peace.
和平.
一章79節.
二章14節.
地上的和平.
然後神當然提了很多次.
神的兒子等等.
所以你就見到路加在做什麼.
聖誕的故事原來是一個挑戰.
或者一個顛倒.
在上的是羅馬帝王.
奧古斯都.
但他是不是真正的主.
是不是真正的救主.
是不是真正帶給我們平安.
是不是有一個真正的福音.
他說中央電視台就說有.
但聖誕的故事就說No.
是基督.
是耶穌基督在馬槽裡面的嬰孩.
這樣看的時候你就明白.
為什麼一些極權或共產的國家.
這麼不喜歡到聖誕節的時候.
有人慶祝聖誕.
因為他們看聖經比你厲害.
我們看聖經的時候覺得.
Children's story.
有什麼大不了.
但你看真正經文裡面所說的.
其實這個嬰孩出世.
是在震動地上的國.
應該是給我們一個Shake up.
告訴我們誰才是我們的主.
誰才是值得我們去吸引.

$^{441}$值得吸引我們的.
是耶穌基督還是白宮那位.
地上那位.
路加很清楚說.
就是說耶穌基督才是神.
神的兒子.
值得讚頌的是馬槽裡面那個.
謙卑降生的耶穌基督.
他才是真正的救主.
如果看路加福音1-2章的段落.
那個脈絡的話其實都很清楚.
聚然之後.
他就說兩個嬰孩的報喜.
哪兩個.
施洗約翰和耶穌.
報喜之後.
就說有兩個孕婦相遇.
一章39至45節.
這兩個孕婦是施洗約翰的媽媽.
以利沙白.
不是以利沙白太久了.
以利沙白.
另外一個孕婦當然是瑪利亞.
一個是老年的.
以利沙白和她丈夫薩加尼亞.
是老邁的一對夫婦.
是沒有兒子的.
當時的社會裡面是一個很大的羞辱.
第二個就是瑪利亞和約瑟.
當他們相遇之後發生什麼事.
你記得那個很有趣的畫面.
兩個孕婦見到大家的時候.
以利沙白說她抬著施洗約翰還沒出來.
他已經開始作動.
很興奮.
為什麼.
因為他還沒見到真正的國王出生之前.
他已經很想見她.
就是這種.
如果兩個女人再談多一點.

$^{481}$我想他早產了.
他就是一種這麼喜悅.
很希望見到真正的救主誕生.
很驚訝為什麼上帝的恩典終於來到了.
我很想出來見他.
是這種情況.
然後瑪利亞就唱出他的詩歌.
這首叫做Menefica.
尊主族.
我心尊主偉大.
然後就說到施洗約翰的出生.
施洗約翰的出生就是薩加尼亞和以利沙白.
這對老年的夫婦.
誕生了一個所謂的Miracle Child.
然後薩加尼亞就由啞巴變成一個可以開口.
第一次開口的時候.
他就唱出讚頌的歌詞.
Benefica.
就是Blessing.
去到讚頌上主.
讚頌上帝.
然後到第二章.
第一至第二十節.
就說到耶穌基督的出生.
出生之後我們又見到.
又是一首詩歌.
就是天使帶來的詩歌.
就是剛才我們詩歌裡面有唱.
Glory in Excelsis.
天上的至高神的榮耀.
然後到嬰孩.
耶穌第八天的時候.
就要行割禮.
然後帶他上去聖殿.
就是主在聖殿.
於是他們在聖殿發生的事.
獻祭等等.
就見到兩個應該大概都是老人家.
西冕和阿娜.
西冕見到這個孩子.

$^{521}$這個救主來到的時候.
他就說我已經等了很久.
然後他又唱了一首頌歌.
叫做Ludimitus.
Ludimitus的意思就是說.
你看那個歌的內容.
二章二十九到三十二節.
拉丁文其實就說.
安然去世.
讓我安然的離世.
因為他已經親眼看到.
以色列的救恩.
神的救恩已經到了.
然後有一個女人叫阿娜.
講清楚她是老年老邁寡居的一個女士.
去讚美神.
這個就是路加所交代的聖誕的故事.
四首的讚美詩歌.
四首的頌詞.
像不像奧古斯都.
奧古斯都出生的時候也有很多頌詞.
他登基的時候也有很多頌詞.
但是路加特別記載這四首的詩歌.
告訴我們真正值得頌讚的.
真正的救主.
就是這位耶穌基督.
你看到沒有.
在這兩章裡面.
跟這個救主相遇的人.
都是一些很低微的人.
不是社會上高在上很顯赫的人.
不是奧古斯都不是居尼留不是地上的君王.
而是那些謙卑去等候上主救恩.
我不知道大家是不是.
上個星期五.
我跟麗仙傳道去探望一位弟兄的媽媽.
因為有些convocation(困難)所以她進了醫院.
我們就去病床探望這位老人家.
很老的.
她已經四五天沒吃東西.

$^{561}$因為身體非常弱.
很多時候去探病的時候很有趣.
就是去探病的時候她特別精神了.
我就跟這位姐妹說.
我們來探望你.
跟她聊了幾句話.
她突然間好幾天很弱.
在探病和她聊天的時候.
她突然間好像很精神.
她看見我.
她是一個牧師.
跟她一起讀詩篇23篇.
聊了一會.
見到她還微笑.
我真的有種感覺.
這就是她第四首詩歌.
或者跟這四首詩歌相似.
她看見神很實在.
神已經跟她同在.
所以她笑了起來.
跟她一起祈禱.
有和應的.
她看見這個救主這個基督才能給我們一個真正的盼望真正的喜樂.
有時候我們人生在世經歷很多事.
我們都會覺得我們自己是小人物.
因為我們是蟻民.
誰會記得我們.
誰會紀念我們.
我們不是在上掌權.
我們不是英雄.
也不是君王.
在歷史裡面我們算什麼.
好像微塵一樣.
聖誕的故事就是說.
基督就是跟我們的人同在.
基督就是那位跟我們一起的救主.
我們是小人物.
但在聖誕故事裡面這些小人物變成了主角.
那些大人物高高在上的.
變成了配角.

$^{601}$他們只是背景.
在奧古斯都的時候發生這件事.
那你還提不提奧古斯都的威風事件.
隻字不提.
提的就是真正的救主.
真正的君.
就在聖誕的時候我相信.
你應該很忙.
要款喀等.
請你不要忘記.
你最重要款待的是這位真正的救主.
我們低頭祈禱的時候.
如果還沒有認識這位主.
跟我一起認識這位主.
請這位主歡迎他進入你的生命.
好嗎.
主我們多謝你.
你降生在馬槽裡面.
去體恤我們.
跟我們同行.
你雖然在高高在上.
滿有榮耀滿有榮光.
但你來到我們中間.
與我們同行.
我們一班在世界.
在歷史裡面好像未曾一樣的生命.
你就與我們同在.
主我們求你.
你都進入我們的生命.
讓你成為我們的主.
讓我們將這位.
誕生在馬槽裡面的英雄.
一位真正宇宙世界的救主.
愛我們的救主.
款待祂.
進入我們的生命.
使祂帶領我們.
與我們同在.
今天弟兄姊妹.
如果有一些很傷痛的經歷.

$^{641}$一些很難過很孤單的經歷.
特別在聖誕的季節裡面.
主我們求你的安慰.
你的平安.
真正的平安.
能夠臨到我們每一位.
祈禱奉主的美感.
\newpage



\section{出埃及記 20:1-17}
\label{sec:x5SxAthpTtg}
\textbf{新子民,新原則 \#1 (出埃及記 20\_1-17) - 陳崇基牧師}
\newline
\newline
連結: \href{https://youtube.com/watch?v=x5SxAthpTtg}{\texttt{https://youtube.com/watch?v=x5SxAthpTtg}} ~~~~ 語音日期: 2025-01-07
\newline
\newline
\hyperref[sec:KuhthfKDgTw]{\small{< < < PREV SERMON < < <}}
~
\hyperref[sec:index_chronic]{\small{[返順時目]}}
~
\hyperref[sec:index_scriptual]{\small{[返順卷目]}}
~
\hyperref[sec:kxR54r6RFz8]{\small{> > > NEXT SERMON > > >}}
\newline
\newline
出埃及記 20:1-17
\newline
\begin{longtable}{cl}
\hline
\hline
章節 & 經文 (和合本修訂版)\\
\hline
20:1 & \begin{tabularx}{0.7\textwidth}{X} 神吩咐這一切的話，說： \end{tabularx} \\ \\ \relax
20:2 & \begin{tabularx}{0.7\textwidth}{X} 「我是耶和華－你的神，曾將你從埃及地為奴之家領出來。 \end{tabularx} \\ \\ \relax
20:3 & \begin{tabularx}{0.7\textwidth}{X} 「除了我以外，你不可有別的神。 \end{tabularx} \\ \\ \relax
20:4 & \begin{tabularx}{0.7\textwidth}{X} 「不可為自己雕刻偶像，也不可做甚麼形像，彷彿上天、下地和地底下水中的百物。 \end{tabularx} \\ \\ \relax
20:5 & \begin{tabularx}{0.7\textwidth}{X} 不可跪拜那些像，也不可事奉它們，因為我耶和華---你的神是忌邪的神。恨我的，我必懲罰他們的罪，自父及子，直到三、四代； \end{tabularx} \\ \\ \relax
20:6 & \begin{tabularx}{0.7\textwidth}{X} 愛我，守我誡命的，我必向他們施慈愛，直到千代。 \end{tabularx} \\ \\ \relax
20:7 & \begin{tabularx}{0.7\textwidth}{X} 「不可妄稱耶和華－你神的名，因為妄稱耶和華名的，耶和華必不以他為無罪。 \end{tabularx} \\ \\ \relax
20:8 & \begin{tabularx}{0.7\textwidth}{X} 「當記念安息日，守為聖日。 \end{tabularx} \\ \\ \relax
20:9 & \begin{tabularx}{0.7\textwidth}{X} 六日要勞碌做你一切的工， \end{tabularx} \\ \\ \relax
20:10 & \begin{tabularx}{0.7\textwidth}{X} 但第七日是向耶和華---你的神當守的安息日。這一日你和你的兒女、奴僕、婢女、牲畜，以及你城裡寄居的客旅，都不可做任何的工。 \end{tabularx} \\ \\ \relax
20:11 & \begin{tabularx}{0.7\textwidth}{X} 因為六日之內，耶和華造天、地、海和其中的萬物，第七日就安息了；所以耶和華賜福與安息日，定為聖日。 \end{tabularx} \\ \\ \relax
20:12 & \begin{tabularx}{0.7\textwidth}{X} 「當孝敬父母，使你的日子在耶和華－你神所賜你的地上得以長久。 \end{tabularx} \\ \\ \relax
20:13 & \begin{tabularx}{0.7\textwidth}{X} 「不可殺人。 \end{tabularx} \\ \\ \relax
20:14 & \begin{tabularx}{0.7\textwidth}{X} 「不可姦淫。 \end{tabularx} \\ \\ \relax
20:15 & \begin{tabularx}{0.7\textwidth}{X} 「不可偷盜。 \end{tabularx} \\ \\ \relax
20:16 & \begin{tabularx}{0.7\textwidth}{X} 「不可做假見證陷害你的鄰舍。 \end{tabularx} \\ \\ \relax
20:17 & \begin{tabularx}{0.7\textwidth}{X} 「不可貪戀你鄰舍的房屋；不可貪戀你鄰舍的妻子、奴僕、婢女、牛驢，以及他一切所有的。」 \end{tabularx} \\ \\
[1ex]
\hline
\hline
\end{longtable}
$^{1}$未必會背.
不過我們肯定有印象.
但我們每講到十誡的時候.
我們總覺得.
中文翻譯為「誡明」.
「誡明」的意思很多時候就是.
你限制或者不准你做一些事.
我們現代人一聽到就很不滿意.
我們覺得為何不讓我們做這件事.
不讓做這件事 不讓做那件事.
這個是不是聖經所講的十誡的意思呢.
我們中文叫十誡.
英文叫ten commandments.
其實聖經本身 特別是希伯來聖經.
即是舊約聖經.
他不是用十誡來形容.
這十條的項目.
他是用十句說話.
ten words.
或者ten things.
即是十件事.
十件事.
十句說話.
所以我們有時候一些書籍.
會介紹十誡是declar.
declar即是十.
logos就是words.
我們叫它十件事.
或者十個原則.
其實不是有味道是do和don't.
不是只說列表一些東西.
然後你要做一些事.
不可以做一些事.
thou shall do.
thou shall not do.
不是這個意思.
是十個原則.
十個principle.
我們可以形容它做一個憲法.
好像我們美國有憲法一樣.

$^{41}$constitution.
constitution就不是那些律師樓.
那些很厚很多冊的law code.
他拿來check下.
到底有哪個的律例犯了罪等等.
是一些大的原則.
所以你如果問美國的大的原則是什麼.
美國的精神.
自由.
甚至放縱.
自由 總之什麼都可以做.
總之你叫我們不要做事.
那就不行.
這個是最大的罪.
因為我們講liberty.
講freedom等等.
這個十誡所謂.
其實是給神的子民一個十個原則.
十個法則.
用這些原則.
你就應用在一些場景.
一些特定的處境裡面.
譬如說不可殺人.
這個是大的原則.
在不同的情況下.
應該怎樣去執行或者活出這個principle.
that's the whole point.
既然大家都這麼熟十誡.
不如我們看看好嗎.
出《亞劉記》二十章一至十七節.
那些不懂得背的只有印象.
現在來了.
還沒到.
不要緊不要看.
我們還沒看之前.
我們需要介紹一件事.
十誡的鋪排.
或者它的格式.
是和當時古代的時候.
古代世界的格式有相似的地方.

$^{81}$因為你記得出埃及的時候.
大概是主前一千二百年.
一千五百年左右的時間.
當上帝去拯救他的子民.
離開埃及.
黑暗的地方.
遺老的地方.
將他們帶出來.
過了紅海.
然後過了曠野.
去到一個地方.
一座山.
一座山叫做西乃山.
賀列山.
在山那裡.
上帝就出現了.
然後就和他們說.
我現在和你們立約.
我用恩典拯救了你們.
現在你們成為一班新的人民.
新的子民.
有新的身份.
我現在給你們新的原則.
這個叫十誡.
十句的說話.
為什麼是十句呢.
有沒有想過.
為什麼不是五句.
或者七句.
因為聖經文化.
七是perfect number.
或者十二也是perfect number.
十二支派.
為什麼是十.
猜到嗎.
大家猜來猜去.
最有可能的.
就是容易記.
大部分人都有十隻手指.
除非是六指及摩.

$^{121}$多了一隻.
大概十件東西.
我們是容易記到的.
所以有十個法則.
去頒布出來.
摩西的時代.
當時他們不同的國家.
會建立一些盟約.
譬如一些大的國家.
埃及.
他會和一些小的國家.
譬如一些小的城邦.
建立盟約.
建立盟約的時候.
他們就會大家簽署一些條約.
譬如我們美國.
和其他國家都有盟約.
我們和歐盟.
我們和加拿大.
其他國家我們會有一些外交的關係.
Diplomacy.
當時摩西的時候.
都是有這樣的.
兩個國家.
特別是一個大國和一個小國.
大王和一個小王.
去建交.
做一些外交的關係的時候.
就會簽署一些協議.
文件裡面就詳細解釋.
大家怎樣維持這個關係.
怎樣維持一個良好的關係.
大國和小國去建交的時候.
我們就簽署一些所謂的government.
或者我們今天.
會叫做Suceran Treaty.
我知道這個字很大.
其實就是約的意思.
大家有一個約定.
這個約定有一個格式.

$^{161}$第一個格式就是.
大王就會介紹自己.
他說我是誰.
然後第二個元素.
就是講到歷史的過去的恩情.
就是說我曾經怎樣對過你.
曾經怎樣救過你.
現在你應該怎樣.
第三個原則就是說.
你心裡面應該有什麼法則.
我頒佈這個誡命的時候.
我頒佈這些約條款的時候.
你就需要做什麼.
你就需要對我忠心.
你都需要效忠於我.
然後第四個原則.
第四個的元素.
就是續約的條文.
就是說你如果死了之後.
誰續你位.
你兒子續你位.
你王子.
他上任的時候.
登基的時候.
就要成獨這個盟約.
然後第五個的元素.
就是有見證人.
當時那些是拜多神的文化.
他們就會召喚所有的神的名字.
月神,太陽神,星星等等.
然後叫他們做見證人.
最後第六樣.
就是有祝福和奏訟.
就是說如果你是維持這個條約.
你是向我忠心的話.
你就會有祝福.
就是說你有難的時候.
我就會來拯救你.
有人攻打你的時候.
我就會來幫你.

$^{201}$但如果你去到叛變我的話.
就會有奏訟.
你就會被我來攻打.
我就會來圍你.
或者加你稅等等.
就會令到你活不到在我的同在下.
享受不到我的同在下.
這些元素都是兩個國家.
大家去到的一些慰藝.
大家有一個這樣的簽署之後.
大家就會放條約的條款.
所有的文字.
就會放在一些硬的石板.
或者泥板上面.
然後你一份就我一份.
大家有一個複製.
好像你和你所公司有一個條款.
或者和銀行有一個條款.
大家就會有一個複製的條款.
然後你保持在你最安全的地方.
銀行又有一個複製.
兩個國家大家有一個外交條文的時候.
你那個複製就放在你的首都.
裡面最大的聖殿.
聖殿裡面的致聖所.
然後我又是做同樣的東西.
所以你一個複製.
我一個複製.
所以我們現代人對十誡的印象就是.
有兩塊石板.
因為聖經裡面說有兩塊石板.
摩西帶著兩塊石板上山.
然後上帝就寫條款.
上去這個石板那裡.
我們現代的印象就是.
一定是一塊石板不夠位置.
好像一隻黑不夠位置.
所以一定要用兩塊.
所以五個的誡命.
一至五就寫在第一塊.

$^{241}$然後六至十就寫在第二塊.
這個就是我們現代人的印象.
其實不是的.
其實應該是十條都寫在同一塊石板.
十條寫在另外一塊的石板.
你一個複製我一個複製.
然後放在你最重要的地方.
所以以色列文和上帝耶和華來建交.
進入關係的時候.
耶和華就說.
以色列文這個石板是給你的.
你有一個複製.
以色列文只有一個聖所.
在曠野的時候.
他們就在會幕.
進入嘉陵的時候就在耶路撒冷.
聖殿.
所以他們就放這個複製.
放在會幕裡面.
至聖所.
或者將來他們進入嘉陵地的時候.
就放在耶路撒冷那個聖所.
至聖所裡面.
所以他們的複製就放在那裡.
耶和華的複製放在哪裡.
上帝的複製.
上帝當然.
他的聖殿.
就在地上.
就在以色列文中間.
所以他那塊石板又是放在會幕裡面.
所以如果你走進去會幕.
最重要的地方.
就是Holy of Holies.
至聖所裡面.
打開藥櫃的話.
那裡就有兩塊石板.
一塊是屬於神.
一塊是屬於神的子民.
進入耶路撒冷的聖殿.

$^{281}$至聖所同樣是這樣的.
所以神救以色列人出埃及.
來到西大山.
他們就來到整個舊約聖經裡面.
由創世記開始.
如果你一年開始的時候.
你說我要讀聖經.
從頭讀到尾.
由創世記開始讀.
一天讀幾章.
你讀到出埃及的20章.
其實就是整個故事的最高峰的時候.
是摩西五經最高峰的時候.
因為他來到山腳的時候.
耶和華向他們顯現.
然後叫摩西你上來.
我會雲端給你.
什麼?.
十誡.
你就成為一個新的子民.
所以十誡開始是這樣的.
上帝吩咐.
這一切的話.
摩西聽見這些話.
拿了一個法板.
然後就帶下去給以色列人.
上帝吩咐這一切話說.
我是耶和華你的上帝.
曾將你從埃及圍爐之地領出來.
十誡你看到開始的時候.
不是誡命.
這裡哪裡有誡命.
他沒有說道尚都.
或者道尚那.
你可以做什麼.
你不可以做什麼.
第一句不是說一樣.
我是耶和華你的上帝.
曾將你從埃及圍爐之地領出來.
不是從命令開始.

$^{321}$如果你看回剛才那個列表.
這裡所有的元素.
你就會找到原來摩西五經.
裡面記載神和以色列人之間的關係.
神和我們之間的關係.
都有這樣的元素.
間中他會說我是耶和華.
我會介紹自己.
一個大王會介紹自己.
然後有歷史的恩情說出來.
之前怎樣去款待過你.
之前怎樣去救過你等等.
十誡開始的時候.
我是耶和華從埃及圍爐之地.
曾經領你出來.
這個就是我曾經用恩典去對待過你的.
之後有一些國度的律法.
我們會繼續看下去.
然後有一些續約的條款.
申命記又說到.
每一代你們都要繼續去教導這件事.
然後有立約的見證人.
多神的文化裡面.
就是召喚那些假神.
他們的明智來見證這件事情.
但當然在聖經的信仰裡面.
只有一位神.
召喚誰來做見證人.
就是召喚天地來做見證人.
整個宇宙整個大地就是見證著.
上帝和他的子民立約.
然後有祝福和奏坐.
這些條款.
所以開始的時候.
他說我時遙華曾經將你從埃及領出來.
是上帝的恩典開始的.
是按照上帝的恩典來成為一個新的關係.
在關係裡面才說.
因為我已經拯救了你.
所以你應該怎樣回應.

$^{361}$舊約到新約.
說到恩典的時候.
我們經常說恩典就是伯伯給你的.
不需要做任何事情.
沒有錯的.
但從來在舊約新約裡面.
都沒有說過恩典就是.
吃了free lunch然後我們就走.
沒有這樣的東西.
中國人就有.
華人就很多.
吃完free lunch就走.
吃太多就走.
聖經裡面舊約講恩典.
甚至新約講恩典.
都是一樣的.
你一定要回應.
因為這個是一個關係.
如果吃完東西走了.
什麼都不理的話.
那個就是關係的破裂.
這個不是領受恩典.
領受恩典的動態.
應該是我接受了這個恩典.
然後我就會知道.
上帝我多謝你.
因為你拯救了我.
這種生活的反應.
或者生活的取態.
就是一個回饋.
所以領受了恩典.
就一定會有行動上的回饋.
所以你看到十誡的鋪排.
和當時的盟約的鋪排都是一樣.
講完恩情.
講完上帝怎樣去拯救了我們.
然後我們會有一個新的身份.
怎樣去新的生活.
去按照新的原則去活出來.
所以十誡開始的時候.

$^{401}$就是先講恩情.
然後就是條文的.
就是約裡面的terms.
第一誡.
所以我們今天叫第一誡.
除了我以外.
你不可有別的神.
但剛才我說過.
在當時的社會.
是多神的文化.
是滿天神佛的.
那些神明呢.
是有地盤的.
甚至我們華人的.
那些folk religion.
就是民間的信仰.
都是這樣的.
土地.
他就管土地的.
誰就管廚房的.
然後誰就管海的.
做軍等等.
你需要拜下他們的.
因為他們有他們的地盤.
舊約.
新約的文化都是這樣的.
但這裡很奇怪.
這裡說除了我以外.
不可有別的神.
當然我們第一個的印象就是.
獨一的神.
但如果你們這樣想的時候.
耶和華有沒有地盤.
有地盤.
但地盤在哪裡.
全個宇宙都是他的地盤.
他說不可有別的神.
他強調的地方.
就是.
英文翻譯得不錯.

$^{441}$There is no other god before me.
In my face.
在我面前.
你不可有別的神.
意思是甚麼.
我整個宇宙都是我的地盤.
都是屬於我的.
你是逃不掉的.
你只有一位神.
這位神是無處不在的神.
每個宇宙的角落.
你家裡.
你心靈.
每一個角落.
都是在上帝的面前.
在他的.
In his face.
是他見到的.
在他的presence裡面.
所以你所有生活的每一方面.
你的家庭.
你的婚姻.
你的子女.
親子關係.
你的工作.
你所有的東西.
都是在上帝的面前.
來活出來的.
There is no hiding.
因為不可能在一個無處不在的神面前.
是能夠躲起來的.
他是見到的.
我們要尊他為榮.
除了我以外.
你在他面前.
不可有別的神.
這是第一屆.
第二屆.
不可為自己雕刻偶像.
亦都不可以做任何的形象.

$^{481}$彷彿上天下地.
地底下水中的白物.
不可跪拜那些像也.
不可事奉他們.
因為我亦會說.
你的上帝是忌邪的上帝.
恨我的.
我必懲罰他們的罪.
至父及子.
直到三四代.
愛我的.
守我誡命.
我必向他們.
施持愛.
直到千代.
這裡有些問題.
忌邪.
為甚麼上帝會忌邪.
英文是jealous.
是嫉妒.
嫉妒通常我們聽見的時候.
就是不好的東西.
嫉妒.
即是你比我漂亮.
我不喜歡你比我漂亮.
所以我毀你容.
雖然我怎樣都比你醜.
但我不喜歡你比我漂亮.
這件事我得不到.
但我不想你有.
這叫嫉妒.
通常都是不好.
你有錢過我.
我不喜歡你.
但我不想要你的錢.
我要你傾家蕩產.
很開心.
這叫嫉妒一定是不好的.
但天下間有兩個關係.
是需要有嫉妒的.

$^{521}$哪個關係.
婚姻關係.
是需要有嫉妒的.
排他性的.
是不能夠有第三者的.
這叫做應該有的嫉妒.
對嗎.
還有第二個關係是什麼呢.
就是你和上帝之間的關係.
上帝怎樣看你和他的關係.
你怎樣看你和上帝的關係.
上帝是一個忌邪的神.
一個嫉妒的神.
上帝和你的關係應該是排他的.
是不可以有偶像在.
偶像意思就是說.
有些東西是佔據你的心靈.
佔據你的心思.
多過上帝.
時間多過上帝.
比重多過上帝.
這些就是偶像.
就是偶像.
這些偶像可以是很多東西.
是嗎.
可以是你的興趣.
可以是你的嗜好.
可以是其他等等不同的事情.
擁有的東西.
上帝說不可為自己雕刻偶像.
以前古代的人.
雕刻偶像的時候.
偶像是神明的載體.
假神的載體.
即是他的embodyment.
你要雕刻他.
然後要服侍他.
然後要清潔他.
每次坐在十五的時候就要討好他.
這些就是我們文化裡面.

$^{561}$現代古代的.
我們怎樣去處理偶像.
因為如果你這樣做的時候.
偶像代表假神.
你對他好.
他自然會對你好.
所以你對他有交換.
大家有一個神明為我服務.
我為什麼要服侍神明呢.
因為我希望他去祝福我.
又來到農曆新年.
一定會有什麼.
那些風水佬.
或者那些講運程的人出來.
今天你犯大罪.
你應該做什麼.
應該乖乖上香.
為什麼.
你上香的話.
他不生你的氣.
那你就ok.
你想行運.
你就要對他好一點.
因為你和神明之間的關係.
是一個交換.
你做得好一點 他又對你好一點.
和你之間的關係.
絕對不是這樣的事.
上帝對你好不好.
在文件裡面 開頭第一句.
上帝對你實在太好.
因為上帝拯救他們的子民.
離開罪惡 離開圍爐之地.
不是因為他們做得好.
而是因為上帝的恩情.
是白白的恩典.
將他們救出來 離開埃及.
將我們救出來 在基督裡面.
將我們從罪 從死亡救出來.
不是我們值得的.

$^{601}$不過救了我們.
我們應該怎樣去回應.
不可有別的偶像.
不可以為自己雕刻偶像.
第二個問題.
為什麼我這麼兇.
不可以跪拜他.
因為我是忌邪的神.
恨我的我必懲罰他的罪.
由父及子直到三四代.
但你沒有看下一句.
如果你愛我的話.
我施恩祝福給你直到千代.
三四代短一點.
還是千代短一點.
千代即是無限的代.
是上帝的恩典大於上帝的審判.
你不要搞錯.
由頭到尾是一個恩典的文件.
包括第二屆 第三屆.
我們看不完這麼多個屆面.
但我們先看頭三個.
第七節.
不可妄稱耶和華你上帝的名.
因為妄稱耶和華名的.
耶和華必不以他為無罪.
不可妄稱耶和華的名.
神的名.
我們一聽見的時候.
我們通常都會覺得.
就不要說 加.
不要說 Jesus.
不要亂說上帝的名.
或者不要說粗口.
我們通常感覺都是這樣.
這就是因為我們現代人.
我們說名字的時候.
純粹是一個標籤.
你叫John.
他就叫Mary.

$^{641}$所以我叫John的時候不是叫Mary.
因為John是一個標籤.
是一個標籤.
幫我分一些東西.
但在聖經的文化裡.
名字不是只是標籤.
名字是權能.
名字是power.
有那個名字的時候.
就有那個power能夠去操縱他.
所以不可妄稱耶和華的名.
就是不要隨便以神之名去做事.
不要隨便以神之名去做事.
我記得我曾經說過.
我認識一些弟兄姊妹.
經常都說上帝的靈感.
令他做這件事做那件事.
但跟他查經的時候.
他就亂來.
你說是不是妄稱神的名.
絕對有這樣的嫌疑.
因為我們經常想妄稱上帝的名.
以神之名.
我們就有些渣拿.
因為是為上帝做的.
上帝一定包底的.
我們不要亂來.
因為如果亂用上帝的名字.
就代表我們根本不明白上帝的心意.
妄用上帝之名.
In the name of God.
不要用來做藉口.
有人可以以神之名去打仗.
去殺戮.
這絕對不是聖經允許的事.
所以上帝的名是不能夠妄用的.
Again.
名字是權能.
是power.
特別在近東社會.

$^{681}$你知道一個人的名字.
你就帶著可以用這個人的權能.
有些Magic.
有些通靈.
就是在做這件事.
拿了他的名字.
然後就可以控制某某神明.
為我服務.
第三個戒名.
不可妄稱神的名.
實際上就是說.
你不能夠亂去操縱上帝.
因為上帝不是你操縱的.
反而你應該服從上帝的權之下.
上帝有絕對的自主權.
另外一方面.
這裡說到不可妄稱神的名.
原文是挺有意思的.
《霞氣百來文》.
妄稱的稱字.
中文翻譯成稱字.
或者call on英文.
其實意思是有帶著的意思.
不可以妄自去帶著.
而會畫的名.
什麼意思.
好像你的行李.
你的backpack.
你帶著carry.
這樣的意思.
帶著一個很重的袋.
這樣的意思.
所以意思就是說.
不可以亂帶著上帝的名.
意義就是說.
我們作為神的子民.
我們其實是帶著上帝的名.
帶著上帝的名聲.
他的reputation.
用現代的例子說一說.

$^{721}$大家知道麥當勞.
麥當勞是run franchise.
專利.
如果你是有麥當勞.
申請了他的專利.
franchise的話.
你就開一間麥當勞.
你就在總公司底下.
你在某個地方.
去開一間麥當勞.
去經營它.
你要有麥當勞的經營權.
有那個franchise的時候.
你就需要聽總公司的話.
按照總公司的標準.
你去做麥當勞有的東西.
他要弄的menu.
他要怎樣做.
他的quality 質素.
應該要做到差不多.
然後顧客就會給你五星.
這個就叫做你經營.
或者carry.
麥當勞的franchise.
第三界別就是這個意思.
就不要亂去經營上帝的名字.
因為在你身上.
是有上帝的brand.
在你身上.
有上帝的牌子.
人家見到你的時候.
他見到這個牌子是好東西.
或者他覺得這個牌子不行.
意思就是說.
不可以妄自穿戴耶和華的名.
當主救贖了我們.
我們身上有了上帝的名.
特別是經過水禮之後.
你記得水禮的時候.
施浸的時候我們說什麼.

$^{761}$奉父子聖靈的名為你施洗.
你記得很多年前你施洗.
記不記得.
老懵懂已經忘記了.
奉父子聖靈的名為你施洗.
什麼意思.
就是說你的生命裡面.
是帶著父子聖靈的名.
是上帝自己的名字.
是上帝的名字.
我的一舉一動.
我要想清楚.
我是不是在經營這個franchise.
經營得好.
人家見到我的時候.
是不是讚上帝的名.
還是覺得上帝的名是不好.
三界對面告訴我們.
全部都是由上帝的恩典開始.
請你不要覺得.
十界只不過是叫你做好人.
或者不要做某樣東西.
你自己改善一下就可以.
最主要不是這樣.
是由恩典開始.
你想起上帝對你的恩典.
你的身分改變了.
你才會看這些界名是合理的.
這些界名就是說.
我應該對上帝效忠.
我應該去carry上帝的名.
是meaningfully.
是讓人知道這個brand是好的.
因為上帝對我的恩典是多麼的大.
十界對我們今天來說.
是很重要的.
如果你看清楚.
十界不是只是.
十界應該technically.
不是給那些微信主的.

$^{801}$我們美國的文化.
特別是我們社會.
經常爭拗一件事.
十界應該釘在哪裡.
或者公眾哪個地方我們要釘在那裡.
I get the point.
但最主要的說什麼.
是上帝有恩情去拯救了我們.
然後那些界名是一種回饋.
對恩典的回饋.
如果你沒有那種恩典出發的話.
是不會make sense的.
只不過是do's and don'ts.
如果是上帝有恩典給我們的時候.
我們就明白.
我們可以去按照這些原則.
去回應上帝的恩典.
如果你這樣看的時候.
十界不是do's and don'ts.
十界是一個對恩典的回應.
因為這個緣故.
所以不同的傳統.
對十界的編碼都是有分別的.
譬如你看到這個編碼.
這個編碼說什麼.
就是左邊的column.
就是猶太人的傳統.
是怎樣將十界來做編碼.
右邊是基督教的傳統.
即是我們今天聖經的傳統.
我們是怎樣編碼的.
你從右邊開始看.
那些字比較小.
但沒關係.
第一界我們說除了我以外.
不可有別的神.
但猶太人的傳統到今天都是.
第一界很奇怪.
不是說你要做什麼不做什麼.
第一界猶太人的傳統.

$^{841}$就是我是耶和華的神.
曾經將你從埃及地.
遺老之地領出來.
是恩情來的.
有什麼做.
其實沒有什麼可以做的.
道shall do 或道shall not.
是告訴你上帝的恩典對你有多大.
然後第二界才是不可有別的神.
大家記得嗎.
是用恩典開始.
怎樣編碼其實都不是最重要.
最要緊的就是我們明白.
上帝的恩典拯救了我們.
改變了我們.
以致我們今天的身份變了.
原來你和我都是掛著上帝的名字.
掛著上帝的名聲.
新的一年有什麼值得我們更新.
有什麼值得一個新的取態去提醒我們.
就好像你手機那些程式.
有時候要更新.
要去更新的.
你和我都需要.
如果上帝是對我們有恩典的神.
我們怎樣去回饋他.
如果你有這樣的心態去改變的時候.
我告訴你很多這些所謂的立志.
要做的事.
完全的角度是不同的.
每天我要在上帝面前花五分鐘時間.
去聽他的說話.
去禱告.
為什麼.
不是一個淑寧習作做功課.
做完就算了.
我去回饋上帝的恩典.
上帝拯救了我.
對我說話.
我當然樂意聽.

$^{881}$深入一點去看上帝的說話.
是樂事.
因為居然有恩典給我.
我可以聽見上帝的說話.
可以認識他.
對其他人.
都是一個回饋.
上帝怎樣對我好.
愛我.
以致我能夠成為其他人的見證.
穿戴上帝的名字.
在事業裡怎樣穿戴上帝的名字.
怎樣待人接物.
怎樣可以有倫理地有道德地行善.
有智慧地去生活.
做一個有說服力的見證.
Wisdom.
Witness.
在上帝的家.
在社會.
都是一樣.
服事是出自愛的.
上帝愛我那麼多.
我認識他的愛有多少.
有多深的時候.
我才明白原來我要回饋他.
盡一生都不能夠.
都還不盡.
這個是新的Elegy.
新的心態.
我們用新的心態.
開始愛上人.
今天第一個主日.
我是領聖餐.
我把聖餐一個很具體的方法.
去提醒你和我.
上帝對我們恩典有多大.
今天我領聖餐.
這個新的身分.
新的Franchise.

$^{921}$我怎樣好好經營它.
求主幫助我們.
我們同心協力祈禱.
主愛多謝你.
多謝你對我們的恩典.
對我們的恩情.
往日我們很多時候都是迷羊.
我們自己走自己的路.
不過你用你的恩典來挽回我們.
在基督的十誡裡面.
你去救贖我們.
買贖我們.
讓我們能夠屬於你.
我們能夠成為你自己的子民.
大家能夠成為弟兄姊妹.
在基督裡面我們能夠合一.
在一個家裡面.
在一個教會裡面.
我們能夠彼此鼓勵.
彼此連結.
都是因為你給我們的恩典.
我們在聖餐當中.
我們為我們的教會弟兄姊妹去禱告.
為我們中間一些有軟弱的.
身體有病的.
無論是最近這個流行季節.
這些小朋友.
我們有很多年輕家庭.
要服侍小朋友.
小朋友有病.
是很辛苦的.
主愛祈求你給他們恩典.
給他們平安.
給他們意志.
為我們中間可能有長者.
身體有不適的.
生活裡面可能是孤單的.
主愛祈求你憐憫他們.
賜平安.
賜安慰給他們.

$^{961}$主愛給我們每一位.
我們在新的一年.
雖然數字上有改變.
可能在生命裡面.
仍然有很多的挑戰.
主憑著你給我們的恩典.
我們去應戰.
鼓勵我們.
幫助我們.
加力給我們.
主愛領受這個聖餐.
這個餅這個杯.
去紀念你愛我們是有多大.
在十誡上的救恩.
我們窮盡一生也不能夠完全回報.
我們這樣祈禱奉主名.
\newpage



\section{20:1-17}
\label{sec:kxR54r6RFz8}
\textbf{New Principals, New People \#2 (Exodus 20\_1-17) - Rev. Mark Chan}
\newline
\newline
連結: \href{https://youtube.com/watch?v=kxR54r6RFz8}{\texttt{https://youtube.com/watch?v=kxR54r6RFz8}} ~~~~ 語音日期: 2025-01-14
\newline
\newline
\hyperref[sec:x5SxAthpTtg]{\small{< < < PREV SERMON < < <}}
~
\hyperref[sec:index_chronic]{\small{[返順時目]}}
~
\hyperref[sec:index_scriptual]{\small{[返順卷目]}}
~
\hyperref[sec:9TO87He9Dbc]{\small{> > > NEXT SERMON > > >}}
\newline
\newline
$^{1}$- Good to see you all, and as Kelly mentioned,.
many of us are sick, especially with little children..
It's been a very tough season, it's a flu season..
And this flu season is more severe than what we had,.
what I had before..
And even I got sick for quite a few days..
As you know, I went to New York for five days.
of family vacation, and then I got sick..
I went in the Christmas break,.
so there were tons of people there on the street,.
and I could hardly move in Manhattan,.
and that's what I expected..
I would get sick, and I did, I did get sick..
So I came back and been sick for five days,.
and my son and my daughter, they're all sick..
I think my wife got sick too..
I think she got sick before we got sick..
Anyways, yeah, so we all got sick,.
and we just get in line to get sick..
And some of our families are also not here.
because their children are sick,.
and Jinfan and Paul are not here.
because they are still, little Isaac,.
yeah, little Isaac is out of NICU,.
and thank you for your prayer..
And they need time to adjust,.
to get into a new routine..
And you can imagine how difficult it is,.
and it's quite challenging, and keep them in prayer..
And as Kelly told us,.
remember the people in Los Angeles,.
how many of you have friends and coworkers.
or relatives living in Los Angeles?.
How many of you?.
Oh, you do, okay, all right, yeah..
Make sure you call them and pray for them, yeah..
I have some friends who live in Los Angeles,.
and they got evacuated, yeah, and it's very difficult, yeah..
And you see news stories reporting,.
you know, the whole town got wiped out..

$^{41}$It's just horrible, it's horrible..
May God have mercy on them..
So good to see you all, good to see you, good to see you..
(laughs).
Anyways, we continue to worship God,.
and with the number of people,.
we should get in circle, actually..
So that's what I've been thinking..
Maybe we should get into circle,.
and we can do a Bible study or something..
But, you know, not today, okay?.
In the future, all right, yeah..
And we are on this mini series on the 10 commandments,.
and we said these commandments are called 10 words..
They are not being described as commandment.
in the Old Testament..
In the Hebrew Bible, never use the word commandment,.
something that you have to do or you shall not do.
to summarize..
There's 10 things, actually,.
the word used by the Old Testament is actually words..
10 words..
So these are the 10 words, the 10 commandments,.
what we call, that God gave to the people of Israel.
when they are being constituted to be a new people of God,.
being saved out of Egypt by grace,.
and now they have a new identity..
Being the people of God,.
they are given new principles to live by,.
a new lifestyle that says to God,.
"Thank you for saving me, God.".
And now we have a new lifestyle,.
we embrace a new set of identity,.
a new set of value system,.
and these value systems are being summarized.
by these 10 words..
So we went through the first three commandments,.
and now we come to the so-called commandment.
or word number four, the fourth principle..
And here it says, "Remember the Sabbath day.

$^{81}$"by keeping it holy six days,.
"you shall labor and do all your work,.
"but the seventh day is the Sabbath to the Lord your God..
"On this day, the seventh day,.
"you shall not do any work,.
"neither you nor your son, nor your daughter,.
"nor your manservant, nor maidservant, nor your animals,.
"nor the aliens, the foreigners within your gates..
"For in six days, the Lord made the heavens and the earth,.
"the sea and all that is in them,.
"but he rested, God rested on the seventh day..
"Therefore the Lord blessed the Sabbath day.
"and made it holy.".
Now, you have to remember the first audience.
who heard this word, this principle, were Israelites,.
and they have been living in Egypt for hundreds of years,.
for some generations..
So when the Israelites are being saved by God out of Egypt,.
all of them are actually 100\% Egyptians..
You get my point, right?.
You get my point..
All of their clothing, all of the stuff.
that carry with them when they left Egypt in haste,.
everything is made in Egypt..
So actually the identity of Israelites.
was not very clear at this point.
because they just got out of Egypt.
and every one of them lived like Egyptian.
and think like Egyptian and lived like Egyptian before..
And what we know today is in ancient Egypt,.
slaves have a 10-day work week..
So you work 10 days and then you get one day off..
And then you work another 10 days.
and then one day off, perhaps, or sometimes not..
And then 10 days and then one day off, right?.
So this is a typical ancient Egyptian slave work schedule,.
10 days and then one day off..
The Sabbath lifestyle says something different..
Six days you shall work.
and then one day you shall take a break..

$^{121}$And you take a break by keeping it holy.
to worship God, to serve God..
And you have to imagine, if you're the first audience.
out of Egypt, Israelite, you have been in Egypt.
working 10-day week for generations..
And then you come to the foot of Mount Sinai.
and Moses said, "This is what the Lord God has given us,.
"a new principle.".
Principle number four is six days you shall work..
And then one day you shall take a break to keep it holy..
And you're kind of like, "Wow, this is actually humane.".
Right?.
This is actually respecting human rights.
to take a break and to rest, not a workaholic..
So you notice the scope of the principle.
is actually, it's encompassing the entire cosmos,.
the entire human society..
Look with me, who shall not work?.
You shall not work..
Your son shall not work..
And in the original language, your daughter shall not work..
Your man servant, male servant, and the female servant,.
your female servant in the original language,.
and then your animals, the cattle, the sheep.
that work the ground..
And then lastly, the alien, the foreigners.
within your gaze..
So seven categories encompassing the entire.
Israelite community, entire Israelite society..
And the principle reveals the character of God..
Who is God that would give this kind of principle?.
Because God is full of compassion..
God is full of compassion..
God is compassionate about especially slaves,.
those who have no power..
Those who have no power to control the schedule..
They're working under the Lord, the master,.
and they have to work..
And God, by giving this principle,.
shows that his heart is full of compassion and love.

$^{161}$for those people who rest..
And this commandment also explains.
why you have to keep six days to work,.
and then one day wholly to take a break.
and keep it holy, keep it sacred..
For in six days, that's how the Lord God did his work..
Remember Genesis 1, there was six-day creation..
Definitely not a scientific report,.
but it's an artistic representation of how God work..
God worked six days..
He set up the entire cosmos..
And then on the seventh day, Genesis 2, he rested..
He rested on this seventh day and blessed..
And because God worked this way,.
we, as the people of God, also need to work this way..
It's interesting when we read.
from the very beginning of this document,.
10 words, 10 commandments,.
from commandment one to three and now to four,.
this is the first commandment that has to do with time..
And schedule..
Time and schedule..
You know, how many of you still use the paper calendar?.
In the beginning of the year,.
you have to buy a new calendar..
Only my wife..
Yeah, right?.
So you guys use electronic?.
Google Calendar, something like that?.
Yeah, that's what I do, yeah..
But you know, my wife is very old school,.
and we have to buy paper calendar every,.
by the end of the year, we buy a new calendar..
So what I do is, every single year,.
it's one of my routine..
Toward the end of year, I will go to Amazon.
or the bookstore to buy those monthly calendar..
You know, you have to flip,.
and then there's a beautiful picture on top,.
and probably a quote here and there..

$^{201}$And then you have these boxes, right?.
And toward the end of year, I always buy a new one..
And I always buy a Jewish calendar,.
because it has the lunar calendar in there..
So the lunar calendar kind of falls.
in line with the Chinese or the Asian lunar calendar..
So like the end of this year,.
end of this month is the new year for the lunar calendar..
So I look at the calendar,.
and I know what day it is in the Western calendar,.
and then in the lunar calendar..
So that's what we do..
You know, it has to do with schedule..
And I want to ask you a question..
You know, when you became a Christian,.
when you accept Jesus Christ as your Lord and Savior,.
it has to do with, first of all, your time..
You know, Jesus kind of,.
kind of restructure and reconfigure your schedule..
Before you came to Christ, you know,.
every single moment, every day,.
you are the boss, you are in charge..
But you came to Christ,.
it has to do with some schedule change, right?.
God modify our week schedule..
So six days, of course, we need to work..
But in general, in the general pattern,.
we take one day off to serve God, right?.
And that's why we gather..
That's why we keep on gathering..
Yeah, I know it's snowing today, right?.
The weather is not very nice,.
but we still keep it holy,.
because this is one of the commandments from God..
And this is one of the application we should have.
the beginning of the year,.
maybe it's our new year resolution..
You know, we keep it holy,.
set it aside and put it in our calendar,.
put it in our schedule,.

$^{241}$that we are to keep it holy the Sabbath day,.
because God bless this time..
So we worship together..
Now, commandment number five, principle five,.
is honor your father and your mother.
so that you may live long in the land.
the Lord your God is giving you..
Now, today, when we imagine family.
in the Westernized society,.
now the son or the daughter,.
they get married and they probably move out.
and build their own family, right?.
Buy their own houses and live somewhere else.
or move out of town into somewhere else,.
because of job, because of different reality..
But in the near East, in the ancient time,.
in the Bible time,.
most of the families are multi-generational..
Multi-generational..
Now, how many of you still live in your family.
multi-generationally?.
With your grandpa and grandma?.
How many of you?.
No, right?.
Most of the people don't live together.
with different generations..
But in ancient Israel, it has always been the case..
You live with your parents and your grandparents,.
and then probably three generations.
or even fourth generation living under one roof..
And here it says, honor your father and your mother..
You know, in Hebrew,.
Hebrew is a very interesting language..
It says honor, in English translation,.
it has the connotation of something that is heavy..
You know, when something that is valuable,.
when something that has dignity,.
that has value, that has importance,.
you kind of see it as something weighty..
You know, it's like a dumbbell..

$^{281}$It's something that is heavy..
So glory and honor are something that is with weight,.
just like your backpack or your luggage or a suitcase..
You have to carry it around..
So in order to weight, to give weight to others,.
someone has to have the resources to do so,.
especially financial resources..
In the ancient time, of course,.
they don't use electronic money..
They use coins or silver or gold bar to do business..
So something that has value has weight, right?.
You carry your silver around in the marketplace.
to buy stuff, and that's the sense here..
Honor your father and mother is to give them the weight,.
give them the heaviness, because they need to be honored..
They need to have the resource..
In the ancient society, there was no hospice,.
there was no assisted living or something like that..
Everyone lives in the house, right?.
Your home is the hospice..
Your home is the hospital..
And that's why life and death, especially death,.
it's not something foreign to the people in those times..
They always see life and death in their houses..
Baby being born, and then the elderly got sick.
and passed away..
It's right in front of them..
So here the commandment speaks of a resourceful.
person, a proper, you know,.
narrowly defined, is a resourceful male..
He has to honor, give honor to his parents..
And you think about it, the essence of the principle.
has to do with your heart orientation..
This new principle, this new value says.
that one should not be self-centered..
But looking after other people, especially parents..
Because the parents have legitimate needs..
And one has to give resources, offer resources.
to fulfill it, to honor, to make it heavy,.
to make it honorable, so that you will live long.

$^{321}$in the land..
Now you is not singular, you is plural..
It's talking about a nation, talking about a people,.
a community as a whole..
So you can live long, you can continue longer.
from generation to generation in the land,.
so that you, your community, in this case,.
the Israelite community, their society,.
if they look after one another,.
if they take care of one another,.
they are not self-centered community,.
they always look after one another,.
care for one another..
Their community will continue..
Their community life will continue, enjoy longevity..
That is the essence of the principle..
Number six, you shall not murder..
You shall not murder..
Now this is a great translation, you shall not murder,.
because the older translation in English.
is thou shall not kill..
Right, the Hebrew word here translated murder correctly,.
means taking innocent life..
Thou shall not take innocent life..
So that's why modern translation try to correct that,.
or try to connote that kind of sense..
You know, it has to do with murder..
Because the big deal about murder is that,.
because the action will destroy an image of God..
Because every man and woman is created.
in the very image of God..
If you take an innocent life, shed innocent blood,.
means that you assault God's image and God himself..
As Pastor Chi has already preaching,.
or going through the Sermon on the Mount,.
I'm sure you know what that means..
It's not about outward behavior..
In fact, the entire document of Ten Commandments.
has to do less with outward behavior..
It has to do with our inner heart,.

$^{361}$our motive, our heart orientation..
So to hate a person is the root of murder..
That's why that is the essence.
that the principle is trying to lay out..
The ten words have the inner, your heart orientation in mind..
They address our motive, our heart..
So we're called in this principle.
to treat other people with compassion, instead of hatred..
Principle number seven, thou shall not commit adultery..
Now in ancient Egypt, polygamy and adultery.
were very common and it's very much so.
in the patriarchal society,.
because the men would treat women as property and objects..
Here you shall not commit adultery..
Again, the heart orientation,.
it's not so much like I don't do certain things..
So I keep the commandment..
The essence has to do with our orientation to do something,.
do our best to protect the family of other people,.
to protect the dignity of other people, including women..
Again, in the Sermon on the Mount,.
Jesus also talks about the spirit of the law,.
the commandment, has to do with ungodly, lustful desires.
that try to objectify other people.
and use other people to satisfy our own desire..
So that is the essence of this principle..
Again, again and again, we see that the heart orientation.
is to protect other people..
We are out here to protect other people,.
to care for other people, to protect their right.
and protect their dignity and honor..
Number eight, you shall not steal..
Now this principle assumes the right to own something,.
your property, your asset, your money,.
but it also sets a boundary for our possessions,.
for our desire to acquire things..
Yes, we can own stuff, we can own asset,.
we can have our own property,.
but ultimately God is the owner of all..
And we are simply stewards, managers..

$^{401}$We're just managers..
God bless us with different resources,.
financial resources and properties and all that,.
but God wants us to be honest with our relationship with him.
and our relationship with our things..
And there is a boundary there..
Now if you want to shed some light.
on how to apply this principle,.
we can always look for other scriptures.
as what we did for other principles,.
like thou shall not murder.
or thou shall not commit adultery..
You look for New Testament,.
look for other Old Testament passages,.
how they apply it so that you know how to live it..
One of the passages in the Old Testament,.
actually it's the last book of the Old Testament,.
the book of Malachi,.
actually the prophet is teaching the people of God.
and say, there was a debate here, 3/8,.
where a man robbed God, yet you robbed me..
Now who said that?.
God said this..
God said to his people, will someone rob God?.
Yet you guys, pointing to the Israelite people,.
you robbed me..
But then you ask, there's a debate,.
and the people, we but by saying, how do we rob you, God?.
We never robbed you, this is our stuff..
And God answer through the prophet Malachi.
in tithes and offerings..
In your offerings, you keep it only to yourself.
because there is a portion of it that belong to God..
If you do that, it is robbery, it is stealing..
Again, it sets a healthy boundary for our possessions.
so that we are not so absorbed by our materials..
Materialism is not divine principle,.
but faithful stewardship is..
And that's what this principle is all about, principle eight..
Principle nine, you shall not give false testimony.

$^{441}$against your neighbor..
Now so often we read the service meaning of this verse,.
this principle, principle number nine,.
and it often has to do with, okay, don't lie, okay?.
You gotta tell the truth..
But you have to pick on the terms here..
It has to do with the context of legal testimony.
in the court..
And in the ancient society, only the men.
have the social status and social power.
to be called upon by the court to give a legal testimony.
when a case is being processed, right?.
When someone is being prosecuted..
And in Israel, a man may be called to give evidence.
in a legal proceeding..
So the principle says that you have to tell nothing.
but the truth..
A false testimony, therefore, is an abuse of your power..
If you tell or you give a false testimony,.
you are engaging in character assassination..
You are trying to assassinate someone,.
someone who may not be guilty..
And here, the essence, it has to do with the use of power,.
use of our social status and our social means..
And the essence is to protect and defend.
your neighbor's character..
You shall not give a false testimony..
You should protect the other person's integrity..
Now, all of these principles, we come to number nine,.
all of them has a nature of revealing God's character..
So you take this principle and you read it,.
and you have to think about why this principle.
matter to God, because that has to do.
with the character of God..
And here, number nine reveals the character of God.
being truthful..
God is a God of truth..
God is a God of justice..
So it matters to him when your neighbor is being accused.
without someone who defend for him, right?.

$^{481}$So thou shall not give a false testimony.
against your neighbor in a court..
You don't abuse your power..
You use your power, however, to protect.
and defend other people..
Again, from principle number one to nine,.
the orientation has to do with serving other people..
Protecting other people, uplifting other people..
So you protect me and I protect you, right?.
We keep one another accountable,.
and we love one another in the covenant..
Number 10, that's the last one..
Now, this is interesting one..
It says you shall not covet your neighbor's house..
You shall not covet your neighbor's wife,.
men's servant, male servant, ox, donkey, the property,.
the stuff that your neighbor has,.
anything that belongs to your neighbor..
Now again, in ancient Israel,.
when someone cannot pay a loan.
or because of disaster or because of sickness,.
he cannot pay a debt,.
the landlord has the power to take away.
pretty much everything..
So the rich will become richer,.
and the poor can easily become poorer..
And this principle has to do with drawing boundaries..
Thou shall not covet..
Those things do not belong to you,.
and you shall not touch it..
So this commandment is kind of unique.
because it has nothing to do with our behavior, right?.
How can you enforce this?.
There's no way you can enforce this, right?.
Because your heart is coveting something..
It's right here inside..
Nobody can seize it, right?.
Nobody can see it, and nobody can detect it.
until you do something, make your action reality..
Coveting is your heart condition..

$^{521}$This commandment is unique..
You cannot enforce it..
It go right into, straight to the heart of the matter..
It has to do with our greed..
Set boundary for our desire, our desire to possess,.
our desires to have more..
Now, the heart of the commandment, number 10,.
is contentment, right?.
You should be content with what you have..
You should be thankful for what you have.
instead of all day long you trying to seek to.
increase your wealth, even at the expense of other people..
Now, as you can see, all 10 commandments,.
all 10 principles are not some kind of.
self-improvement moral code..
It is, however, a moral vision, right?.
You see, if every Israelite will take these principle.
and apply them in their daily life,.
what kind of society will it become?.
It will become a society that would love God..
It will be a community that would take care of one another..
Everyone is not taking care of themselves, not selfish..
It's not self-centered, but other-centered..
You come to the community,.
come to the community of people of God..
We seek to love God, right?.
Protect God's right to be worshiped..
And then love other people with all our hearts.
and love our neighbors as ourselves..
The 10 words actually challenge the people.
to use their power and we orient their resources.
toward other people, to serve other people..
The direct audience of the 10 words are actually male,.
father, because they have sons and daughters..
They are married..
You shall not commit adultery..
They are pretty well-off because they are resourceful..
Their parents are aging..
In this document, 10 commandments,.
if you read it carefully and critically,.

$^{561}$the direct audience, of course it applies.
to every person of God,.
every people of God, and every Israelite is applicable..
However, the direct audience,.
the first audience of the 10 words, the 10 commandments,.
are actually males, fathers, married people,.
and pretty well-off..
In today's term, probably middle-aged men,.
middle-aged men and being kind of successful,.
stable in our career, in our family life,.
in terms of our property,.
in terms of our social status,.
everything is pretty stable..
These 10 words are actually addressing.
directly to these people..
In a way, it's addressing to all of us..
We are all resourceful in our own ways..
And we are blessed by God..
And the 10 principles are a reorientation.
of our value system..
When we come to God, when we come to know Him.
and follow Jesus, everything in our life,.
our time, our schedule, our resources,.
everything will be reconfigured, reconfigured..
I would like to use the analogy.
of what it means to come to Christ..
The analogy is kind of like, when you come to Christ,.
your operation system is different..
Your operation system is totally different..
The value system is different..
The orientation is totally different..
You cannot use the old operation system to run your life..
You use the new one..
And here, the 10 words are the new kind of.
operation system for our lives..
We protect God's right to be worshiped..
Only one God, no idol..
Sabbath, keep it holy..
We are called to reorient our heart.
to protect other people and to lift up other people.

$^{601}$with compassion and with care..
It is no wonder in the Israelite society,.
if you flip the pages of the Old Testament,.
especially the first five books of Moses,.
the what we call the Torah,.
you don't find police in the Torah..
You have priests, you have prophets,.
you have teachers, you have farmers..
You have totally different, many different roles.
in the moral world of the Torah..
But you do not find.
Israelite police department in the Torah..
You know why?.
Because if everyone observe and take to heart.
these principles, there's no need for law enforcement, right?.
In such kind of communities, such kind of society,.
everyone is out to protect other people..
Everyone is out there to care for other people..
And actually, this is what Jesus said, right?.
If you sum up the entire Torah, he gave two statements..
Love God with all your heart and all your mind.
and all your strength..
And second of all, is love your neighbors as yourself..
Now, if we as a church,.
if we will take it to our heart,.
what kind of community, what kind of fellowship can we build?.
It's not just we come here and consume something..
We take something and we leave..
If we come here and to contribute,.
to pray for other people, to care for other people,.
to love your neighbors, love your brothers and sisters.
as yourself and love God with all your heart.
and all your strength, what kind of fellowship?.
It will be a community pleasing to God, right?.
It will be a community totally different.
from all other communities and clubs out there,.
all other social associations out there.
because they're always born by their own common interest.
or by their own common benefit..
But the church, the people of God are called to be different..

$^{641}$Other-centered, love others,.
pray for others and care for others..
And that's why in our weekly update,.
I have been doing something new starting this year..
I wanna do some weekly update in video format.
and send it out to the social media group..
And I'll do one in Chinese, in Cantonese,.
and one in English..
And I always wanna encourage people to not only get up today.
what's going on in our church, but to pray for..
Our church, to care for other people, right?.
We come to church and ask not just how are you,.
not how are you doing, yeah, how can I pray for you?.
How can I care for you?.
How can I give so that I can make your life more abundant.
to bless you, right?.
So we can bless each other..
Let's all do that, right?.
In this coming week, take one of this principles,.
one of this 10 principles and apply it in your home.
and now church, whatever that is, try to apply it..
Take it to heart..
Remember to love God with all your heart and all your soul.
and all your might and love neighbors,.
your neighbors as yourself..
Let's all pray..
Father, you call us to be a different people.
because you are such a great and unique God.
and you save us by grace out of land of slavery,.
out of sin and out of death..
And now we are being called to take up this new identity,.
to have the honor to be called the people of God,.
your church, your followers..
Help us, God..
Help us to have this principle transform us from inside out..
Address our, we all ran our heart desires.
so that we are not just self-seeking..
And you call us to seek you.
and to love other people as ourselves..
Build our church in love..

$^{681}$Build up our community so that we can fellowship together,.
care for one another and pray for one another..
In this coming week, as we take some of these principles.
to apply in our home, in our church, in our lives,.
I pray that you grant us grace.
so that we will live as your people.
to honor your name in Christ name pray..
Well, thank you very much for coming in the snow.
and there was a refreshment prepared by Kelly..
Again, she's out here to serve us and we should appreciate..
And please stay for a little bit for a time of fellowship.
and to enjoy the refreshment together..
All right, see you next week..
\newpage



\section{出埃及記 20:1-17}
\label{sec:9TO87He9Dbc}
\textbf{新子民,新原則 \#2 (出埃及記 20\_1-17) - 陳崇基牧師}
\newline
\newline
連結: \href{https://youtube.com/watch?v=9TO87He9Dbc}{\texttt{https://youtube.com/watch?v=9TO87He9Dbc}} ~~~~ 語音日期: 2025-01-14
\newline
\newline
\hyperref[sec:kxR54r6RFz8]{\small{< < < PREV SERMON < < <}}
~
\hyperref[sec:index_chronic]{\small{[返順時目]}}
~
\hyperref[sec:index_scriptual]{\small{[返順卷目]}}
~
\hyperref[sec:kf1XXedwCsU]{\small{> > > NEXT SERMON > > >}}
\newline
\newline
出埃及記 20:1-17
\newline
\begin{longtable}{cl}
\hline
\hline
章節 & 經文 (和合本修訂版)\\
\hline
20:1 & \begin{tabularx}{0.7\textwidth}{X} 神吩咐這一切的話，說： \end{tabularx} \\ \\ \relax
20:2 & \begin{tabularx}{0.7\textwidth}{X} 「我是耶和華－你的神，曾將你從埃及地為奴之家領出來。 \end{tabularx} \\ \\ \relax
20:3 & \begin{tabularx}{0.7\textwidth}{X} 「除了我以外，你不可有別的神。 \end{tabularx} \\ \\ \relax
20:4 & \begin{tabularx}{0.7\textwidth}{X} 「不可為自己雕刻偶像，也不可做甚麼形像，彷彿上天、下地和地底下水中的百物。 \end{tabularx} \\ \\ \relax
20:5 & \begin{tabularx}{0.7\textwidth}{X} 不可跪拜那些像，也不可事奉它們，因為我耶和華---你的神是忌邪的神。恨我的，我必懲罰他們的罪，自父及子，直到三、四代； \end{tabularx} \\ \\ \relax
20:6 & \begin{tabularx}{0.7\textwidth}{X} 愛我，守我誡命的，我必向他們施慈愛，直到千代。 \end{tabularx} \\ \\ \relax
20:7 & \begin{tabularx}{0.7\textwidth}{X} 「不可妄稱耶和華－你神的名，因為妄稱耶和華名的，耶和華必不以他為無罪。 \end{tabularx} \\ \\ \relax
20:8 & \begin{tabularx}{0.7\textwidth}{X} 「當記念安息日，守為聖日。 \end{tabularx} \\ \\ \relax
20:9 & \begin{tabularx}{0.7\textwidth}{X} 六日要勞碌做你一切的工， \end{tabularx} \\ \\ \relax
20:10 & \begin{tabularx}{0.7\textwidth}{X} 但第七日是向耶和華---你的神當守的安息日。這一日你和你的兒女、奴僕、婢女、牲畜，以及你城裡寄居的客旅，都不可做任何的工。 \end{tabularx} \\ \\ \relax
20:11 & \begin{tabularx}{0.7\textwidth}{X} 因為六日之內，耶和華造天、地、海和其中的萬物，第七日就安息了；所以耶和華賜福與安息日，定為聖日。 \end{tabularx} \\ \\ \relax
20:12 & \begin{tabularx}{0.7\textwidth}{X} 「當孝敬父母，使你的日子在耶和華－你神所賜你的地上得以長久。 \end{tabularx} \\ \\ \relax
20:13 & \begin{tabularx}{0.7\textwidth}{X} 「不可殺人。 \end{tabularx} \\ \\ \relax
20:14 & \begin{tabularx}{0.7\textwidth}{X} 「不可姦淫。 \end{tabularx} \\ \\ \relax
20:15 & \begin{tabularx}{0.7\textwidth}{X} 「不可偷盜。 \end{tabularx} \\ \\ \relax
20:16 & \begin{tabularx}{0.7\textwidth}{X} 「不可做假見證陷害你的鄰舍。 \end{tabularx} \\ \\ \relax
20:17 & \begin{tabularx}{0.7\textwidth}{X} 「不可貪戀你鄰舍的房屋；不可貪戀你鄰舍的妻子、奴僕、婢女、牛驢，以及他一切所有的。」 \end{tabularx} \\ \\
[1ex]
\hline
\hline
\end{longtable}
$^{1}$每次主因報告的時候都這麼生猛.
令我們本身都有點睡不著.
音樂這麼好.
所以我們醒了.
我想告訴大家.
昨天下午的時候.
我在想.
因為LA大火.
有很多事情值得我們反省.
很多時候我們想起LA大火.
有些什麼訊息等等.
有時候我們說快了.
我們就說上帝審判.
或者什麼.
我們在名義上.
我們沒有怎麼思考過.
我們就說出來.
我相信我們大家不是這樣的.
其他人可能是這樣.
所以下午的時候我就想.
是不是應該轉港台呢.
我想說一系列.
是說創造關心.
大地和萬物.
人去到做管家.
聖經裡面有什麼教導.
有時候我們有種錯誤的看法.
就是覺得這個世界都快完結了.
是不是.
等它燒掉.
反正都玩完了.
Game over.
是不是聖經也是這麼說的呢.
所以我昨天下午的時候.
我就在考量.
到底是不是今天需要去.
說另外一個題目.
所以我用了下午幾個小時的時間.
我就預備了大概60\%的度.
那邊新的度.

$^{41}$但是到了晚上的時候.
我就想了一下.
我們做人不可以這樣.
有始有終.
因為我們在說十屆.
說了一半而已.
應該也說完它.
所以下個星期開始.
下次才去想一想.
這個題目是值得我們去想的.
當然LA大火是非常慘重的.
非常慘重.
求主要去連運當地的人.
我們先不要落井下石.
更加不應該落井下石.
就是想一想.
到底環境和氣候的變化.
有沒有一些大的因素在裡面.
我們需要去想一想.
以及回到聖經裡面.
怎麼說這個世界.
是不是這個世界就這樣丟掉了它.
就罷了.
聖經告訴我們.
新天新地會更新.
是什麼意思.
是不是丟掉了這件事.
然後我們買一個新的.
這個是值得我們去反省的.
也是反省到我們.
不說那麼多了.
這個下個星期.
Anyway 先說回十屆.
有始有終.
我們先出這張slide.
這張就是十大精選.
是這個月的.
希望這個月跟你一起度過.
順帶一提.
這裡所選的港獨.

$^{81}$不只是小大局的港獨.
有很多長長短短.
不同的港元分享.
希望這些十大的精選.
可以幫助大家.
你煲劇可以.
是可以煲劇的.
不過也煲這個好嗎.
OK.
所以希望這個月我們能夠.
在主計華語裡面重新開始.
這個也是我們十屆的系列.
我們希望能夠做到的.
我們上個星期說了頭三屆.
這次我們總結說到四和十屆.
這個文件正如我所說.
在特別西方的文明裡面.
是非常著名的.
我們大家要知道十屆是什麼.
以至我們才能夠明白.
所謂基督教.
猶太教和基督教的傳統.
對西方的社會.
對西方的文明.
對我們的思想.
其實有什麼影響.
就是這個文件.
所以我們值得去看一看整個文件.
這次我們不如從頭讀到尾.
十七節.
二十章一至十七節.
《出亞及記》.
我們看一看十屆這段.
這麼著名的聖經文件.
到底有什麼重大的啟示出來.
二十章第一節.
上帝吩咐這一切話說.
我是耶和華你的上帝.
曾將你從埃及地遺牢之家領出來.
除了我以外.

$^{121}$你不可有別的神.
不可自己雕刻偶像.
也不可做甚麼形象.
彷彿上天下地.
地底下水中的白雲.
不可跪拜那些像.
也不可事奉他們.
因為我耶和華你的上帝.
是既邪的上帝.
恨我的.
我必懲罰他們的罪.
致父及子直到三,四代.
愛我守我誡命的.
我必向他施慈愛直到千代.
不可妄稱耶和華你上帝的名.
因為妄稱耶和華名的.
耶和華必不以他為無罪.
當紀念安息日.
守為聖日.
六日要勞碌做你一切的工.
但第七日是向耶和華你的上帝.
當守的安息日.
這一日.
你和你的兒女.
勞瀑,婢女,生畜.
以及你成女寄居的客女.
都不可做任何的工.
因為六日之內.
耶和華做天,地,海.
和其中的萬物.
第七日就安息了.
所以耶和華賜福與安息日.
定為聖日.
當孝敬父母.
使你的日子在耶和華你上帝所賜你的地上得以長久.
不可殺人.
不可姦淫.
不可偷渡.
不可做假見證.
陷害你的鄰舍.

$^{161}$不可貪亂你鄰舍的房屋.
不可貪亂你鄰舍的妻子.
勞瀑,婢女,牛奴.
以及他一切所有的.
一口氣讀了十誡.
是不是還沒看過.
我們經常聽十誡.
可能大家沒有看過.
《出海及記》二十章.
這個是聖經裡面其中兩處講到十誡的地方.
第二個地方你有機會回去看.
《申命記》第五章.
十誡我們看《出海及記》第二十章的時候.
我們上次提過.
所謂十誡其實不是一個最好的翻譯.
因為聖經本身沒有用所謂誡命.
Commandments.
來形容這十句話.
事實上它就是用十句話.
Ten words 或者Declare.
十句的原則.
有一點像今天我們美國有的憲法.
怎樣做一個美國人.
聽一聽憲法.
活出憲法或者形容憲法.
想一想憲法.
讓憲法來支配我們的生活.
或者影響我們的生活.
以至我們的生活像一個美國人.
這個就是Constitution的意思.
不是一些很細節的東西.
一些大的原則.
一些所謂意識.
一些英文叫做Moral Vision.
就是一個道德上的願景.
以至一個國度.
能夠像一個國民.
一個美國的國民就是一個這樣的願景.
講Liberty 講自由等等.
在十誡就是告訴我們.

$^{201}$上帝的國度 上帝的子民.
應該有什麼大原則.
這些大原則能夠幫助我們記了之後.
然後在我們生活裡面來形容出來.
上次我們講了三界.
這次第四界.
就是當守安息日的界名.
你記得以色列人從哪裡出來.
他們是埃及出來的.
在埃及有幾百年他們是做奴隸的.
我們今天所知.
古代埃及的時候.
他們的勞工的法例.
或者勞工的作息時間.
是怎樣的.
他們的Work Schedule是怎樣的.
他們的所謂的Work Week.
是十天的.
做十天就休一天.
所以你看到這個界名出來的時候.
他們出了埃及.
上帝頒佈這個界名.
他們說做幾天.
做六天就休一天.
是不是很好.
如果你這樣看.
用第一代的讀者來看.
用第一代的讀者.
他們剛剛出了埃及的那班以色列人.
其實都不可以叫做以色列人.
因為他們全部都是Made in Egypt.
因為他們出來的時候.
100\%他們就是埃及人.
因為他們受埃及的人文化.
受埃及的社會.
受埃及的所有的東西來影響.
一個埃及人出了埃及.
然後上帝在西乃山告訴他們.
我現在給你一些新的憲法.
你做一個新的子民.

$^{241}$有一個新的Identity.
你不再是埃及人了.
就是上帝約之下的子民.
然後第四屆告訴他們.
你們守安息日.
因為我的國民.
我們的操作.
或者我們的日常時間的運作.
是跟你以前的埃及是不同的.
以前是十天放一天.
現在是六天就放一天.
六天你要工作.
然後第七天你就去休息了.
是不是很好.
這個其實你用第一代的讀者來看.
是非常人道的誡命.
是知道人是需要作息.
六天放一天好不好.
還是十天放一天.
我知道大家的工作.
可能是十天放一天的.
十天放一天.
十天都不可以放假.
但是你要知道.
那個誡命本身是告訴以色人.
現在你們作息的方法是不同的.
埃及以前是這樣.
現在不是這樣.
因為上帝是體恤創造界.
包括人類.
包括甚至什麼.
你看這裡說什麼.
接觸的層面.
是整個創造界裡面的七個不同的層面.
誰要守安息日來到敬拜.
第一是你.
第二是什麼.
你的兒子.
但原文是你的兒子.
然後女兒.

$^{281}$應該是你的女兒.
然後奴僕.
你的工人.
然後婢女.
你的婢女.
還有生畜.
你那種生畜.
或者你那群生畜.
然後還有什麼.
里城裡面的你的客女.
那些是屬於你的東西.
上帝讓你管.
但你管的方法.
是和上帝工作和作息的時間是一樣的.
所以他就解釋了.
上帝是這樣做事的.
你記得六日創造.
用一個這樣的方法.
展示出來上帝創立整個宇宙.
然後第七日創世的第二章.
休息.
上帝安息.
這個誡命很有趣.
由第一個誡命到第四個誡命.
第一個誡命和時間有關.
是和你的行程有關.
和你的工作行程.
和你的生活行程有關.
大家有沒有用Google Calendar.
或者寫.
有誰還是寫月曆的.
每年年初的時候就買一個新的月曆.
一定要寫上去的.
有沒有.
我太太就是這樣.
每年她都要買一個新的月曆.
所以我們年尾的時候.
我就負責買月曆的.
因為它一定要貼在那裡.
貼在廚房那裡.

$^{321}$然後寫在格子上.
它是這樣的.
你們是不是這樣.
所以我年尾的時候就買一個月曆.
我就買一個猶太人的月曆.
因為有時候方便我讀經.
什麼時候是聚鵬節.
我自己知道.
其實就是和我們農曆的月曆差不多一樣.
我們快到正月了.
所以最冷的時候.
芝加哥最冷的時候.
所以我年尾的時候就買一個猶太人的月曆.
然後它就會從他們的聚鵬節.
就是秋季的時候.
直到明年的年底.
然後到了八九月我就要想.
又要買另一個.
無論如何這個的誡明.
就是一個時間的守則.
或者我們應該這樣想.
這個所謂的Constitution.
到底有什麼的經文.
或者聖經裡面其他地方之後.
再去應用這個十誡的時候.
它是怎樣應用的呢.
我們時間關係.
我們不看太多仔細的.
不過有所謂安息年的誡明.
安息年的誡明是出自這裡.
譬如以前是耕種.
農耕的社會.
他說田地是要耕作六年.
六年之後呢.
就怎樣.
就要休耕一年.
那個叫做安息年.
Sabbatical year.
就是等土地要去休息.
我們中間有很多園藝高手.

$^{361}$你一味種一味種那些東西.
那些泥會怎樣.
那些東西會被你用乾的.
為什麼那些蔥越來越不綠.
我也有一點經驗.
不過只是種蔥而已.
為什麼種來種去都不出.
因為你要fertilize它.
你要放肥料等等.
讓它休息一下.
土地就像我們人一樣.
對嗎.
只不過西方社會我們泛濫到用一個原則.
不說太多了.
這個叫creation care.
這個下次說.
Anyway.
這裡就是說你的時間.
六天作工有一天是需要給.
你去休息.
並且去敬拜.
工作作息敬拜.
原來是和我們作為耶穌基督的門徒有關.
你有沒有想過我們上了耶穌之後.
是要在時間上有改變的.
Yes.
在schedule上要有改變.
對嗎.
這個就是第四屆.
好我們快點看下去.
第五屆.
當孝敬父母洗禮的日子在耶和華.
你上帝所賜你的地上得以長久.
是中蘇來的.
你們應該孝敬父母.
以至你整個國民能夠住在上帝的應許之地.
你們的日子能夠長久.
這裡不是說長命.
而是整個民族能夠繼續延伸下去.
繼續去生長下去.

$^{401}$今天我們當然不同.
我們今天通常.
譬如子女結婚.
他們就走.
我們很少幾代同堂.
有誰是幾代同堂一起住的.
比較少一點.
特別在西方西化了的社會.
我們建立了新的家庭.
然後有新的個體.
在近東的社會裡面.
或者甚至我們華人傳統的文化裡面.
我們都是幾代同堂的.
幾代走在一起.
這裡的意境就是.
這個誡命就是這樣的意境.
他說要孝敬你的父母.
因為那個人還是有他上一代和他一起住的.
當然他下一代有兒女.
所以是三代同堂.
這個誡命的第一代讀者是三代同堂.
這裡說你要孝敬父母.
希伯來文很有趣.
希伯來文的孝敬這個字.
孝敬是說Honor.
就是要尊敬他.
尊貴給他.
其實那個字的意思.
是有重量的意思.
凡是有價值的東西.
就是叫做孝敬或者尊貴他.
或者榮耀他.
那個字其實是說.
Something that is heavy.
就好像你拿著背包.
拿著行李箱.
是有重量的.
你需要給他重量.
就是這個意思.
孝敬他就是給他應該有的重量.

$^{441}$給他應該有的尊貴.
這裡說的是什麼呢.
說的是一些很實際的東西.
和錢財可能有關.
因為當時錢財都是呈的.
呈出來的.
是有資源的人才可以孝敬父母.
有資源的話就要孝敬父母.
把你的資源放在父母身上.
就是這個意思.
孝敬.
Honor.
對象.
這個的誡命由頭到尾對象.
是一個成年人.
第一代讀者.
就是這個十佳.
的第一代.
最直接的受眾.
是一些有才有力的成年人.
這個是以色列的.
或者神國子民.
有的Moral Vision.
這個Moral Vision說.
我們不應當自己顧自己.
要顧軟弱的.
要給尊貴.
給他最第一個的關係.
當然就是父母.
最合理的.
因為特別是在古代的時候.
父母沒人理會就橫屍街頭.
以前沒有health care.
以前沒有hospice.
你家裡就是hospice.
家庭就是需要去照顧.
不只是自己.
而是照顧他人.
得而長久.
這個不只是單單長命.

$^{481}$也不是保證.
而是整個的社群.
擁護這個價值的時候.
他們的生命就可以繼續延續下去.
在應許之地.
他們的民族能夠繼續持續下去.
是一個不是self-centered的社群.
是一個other-centered的社群.
或者我們可以看看.
其他經文怎樣應用這個誡命.
有舊約也有新約.
剛才我們說安息年舊約的其他經文.
怎樣應用第四誡.
第五誡當然也有很多不同的誡命.
但到新約的時候也有應用這個誡命.
譬如《已忽所書第六章》.
《已忽所書第六章》說.
作兒女的你們在主裡面要聽從父母.
這是理所當然的.
第二第三節.
當孝敬父母使你得福在世長壽.
這是第一條大應許的誡命.
這裡說的是第十誡第四誡.
第五誡.
作父親的你們不要激怒兒女.
但要按著主的教導和勸誡養育他們.
這個應許就闊很多了.
不只是孝敬父母.
當然孝敬父母是中心.
第二第三節.
但說的是兒女怎樣在主裡面聽從父母.
父母也不是全部都帶.
當然在近東社會裡面.
父權社會就是父親最大.
打你就不要站定.
罵你不要還口.
永遠都是爸爸最大.
這裡說的是在主裡面要聽從父母.
第四節.
又說父親.

$^{521}$這裡不只是孝敬父母.
說父親怎樣養育兒女.
不要激怒他們.
要用compassion.
要用愛來管理你的家.
按照主的教導和勸誡來養育他們.
所以說的是整個家庭.
說的是整個民族.
民族能夠繼續延伸下去.
生命可以延續下去.
以至在地上能夠見證上帝.
下一個誡命.
第六個誡命.
不要殺人.
當然以往的英文翻譯.
甚至中文的翻譯.
都有不善之處.
不可殺人.
thou shalt not kill.
其實新的翻譯就好很多.
不可謀殺.
因為殺人這個字.
主要是說流無辜人的血.
那些innocent的life.
不可以流他的血.
不可以將他殺害.
意思是說.
其實是一種保障他人的生命.
不是說我翹著雙手不殺你.
已經很好了.
我今天跟你有爭論.
我叫我不殺你.
好.
不是.
這裡說的是保障他人的生命.
為什麼呢.
因為人在希伯來聖經.
在整本聖經裡面.
是帶著神的形象.
謀殺.

$^{561}$premeditated murder.
其實就是破壞上帝的形象.
破壞人.
神在地上的代表.
人就是神的代表.
所以不是說行為這麼簡單.
是說我們的心態.
我們對人的態度.
如果你記得庚山寶訓.
耶穌就去引用這個誡命.
他就說出這個誡命真正的essence.
真正的本質是什麼.
就是需要去愛人.
有兇殺的意念.
或者有殘酷的意念.
就等於殺人.
由頭到尾.
整個十誡都是針對不是外在的行為.
這些不是我們以前中學小學的時候的校規.
是針對我們的動機.
針對我們內心的事.
不可姦淫.
這個更加因為我們大家很熟悉.
以往古埃及裡面.
婚姻的文化或者家庭的文化是十分自由的.
婚外情是很普遍的.
當然誰有婚外情.
當然是男人.
男人才有這樣的power.
不可姦淫.
其實就是不單止保護自己的家庭和婚姻.
因為這是上好的東西.
是上帝所設立的.
你也是保護其他人的家庭.
你是保障他人的家庭和他人的尊嚴.
不犯姦淫.
因為人是神的形象.
婚姻是上帝所設立的.
其他的經文當然我們不需要多說.
我這裡引用了登山寶訓的其中一段.

$^{601}$就是耶穌說.
不可姦淫.
他點出真正的essence在哪裡.
凡看見婦女就動淫念的者人心裡面.
motive.
裡面的東西desire.
已經與他犯了姦淫.
所以整個十誡從頭到尾都不是說外在的東西.
是說內心.
我們怎樣跟人建立關係.
人是誰 其他人是誰.
他是神的形象 我要尊重他.
是針對內心的desire.
不可偷渡.
似乎聖經是贊成個人擁有權的.
財產的擁有權.
你可以擁有你的asset.
你可以擁有你的property.
但所針對的是心裡面的東西.
不是單單是我不偷你的東西.
而是在說貪念.
在你的擁有裡面.
需要畫上boundary.
有些東西不是你的.
其他經文怎樣形容呢.
例如馬拉基書.
舊約裡面最後一卷的書卷.
說到以色列人指責以色列人甚麼事呢.
因為他們重建了聖殿.
沒有人去搞敬拜.
所以聖殿沒有人去修理.
沒有資源去修理等等.
他說人豈可搶奪上帝.
你們竟搶奪我.
你們就是以色列人.
以色列人就說我們甚麼事搶了你的東西.
我們個個都有自己的東西.
上帝說其實你們在當立的十分之一奉獻和當憲的供物上.
你們搶了我的東西.
上帝說那份東西是我的.

$^{641}$你們搶了它.
拿了它.
這就是不可偷渡的其中一個application.
是贊成有個人的擁有財產的權利.
當然我們有.
但我們人.
我們有的物質.
我們所擁有的物質.
我們的資產.
我們的銀行帳戶.
都是管家.
上帝給了我們的時候.
給我們一些擁有.
但需要有一個boundary.
需要有一個尼巴.
因為上帝才是萬有的主宰.
不要在屬於上帝的東西上.
我們在偷渡或者搶奪.
第九誡.
不可作假見證.
陷害你的鄰舍.
這是我們第一個想法.
不要說謊.
你看清楚一點.
我不是說.
大家彼此說謊.
不是這個意思.
不可作假見證陷害你的鄰舍.
意境在哪裡 場合在哪裡.
在法庭.
因為戒命的用詞等等.
在法庭上被人召上來.
作公證的時候.
要給證據的時候.
做一個見證人的時候.
你不可以用不公義.
或者不管真理的方法.
來陷害你的鄰舍.
其實反過來說.
你要保障你的鄰舍的integrity.

$^{681}$保障你的鄰舍.
不會受到不公義.
或者不合真理的審判.
大家要protect大家的character.
沒有一些英文的東西.
我們叫做character assassination.
就是躲在一邊.
我們就去暗殺.
意思就是說gossip.
說一些不合情理的東西.
或者去踩低人的東西.
或者人格上的刺殺.
戒命的本質就是在說上帝的本性.
是一個公義.
本質就是在說上帝的本性.
是一個公義的上帝.
是一個真理的上帝.
他是緊張真理緊張公義緊張justice.
你就不要用你的嘴巴.
去攻擊其他人.
或者暗殺其他人.
我去過好幾間不同的教會.
很多時候開會的時候.
有些人就在做一些人格刺殺.
傳道人都不能招供.
不能去defense.
他們就定他罪.
這個第九誡.
不可作假見證.
陷害你的淪陷.
每一個人都應該有一個公道的defense.
因為大家應該要保障大家.
最後一誡.
不可貪戀你的淪陷的房屋.
和淪陷的妻子等等.
和他一切所擁有的.
以色列的社會.
古代的社會.
當一個人.
譬如有天災來的時候.

$^{721}$他可能會欠債.
他欠債的時候.
他就要在債主的下.
來去幫助他.
他就要在債主的下.
來去幫助他.
在債主的下.
來去生活.
債主就有能力.
來去將他剝削.
能夠奪取他的東西.
奪取他的田地.
奪取他的房屋.
他的資產等等.
所有不動產或可動產.
全部都拿走.
在先知書裡面經常都有.
這個誡命.
說的是什麼?.
不可貪戀.
從頭到尾都在說.
不只是外在的校規.
不是外在的行為.
而是我們的內心.
我們裡面.
第十誡就表明這件事.
十誡是在說內心.
真正的精神.
是在說我們心理的狀態.
第十誡.
是在說我們要知足.
因為以前.
特別是古代的社會.
富者就越富.
貧者會越貧.
因為富者就有能力.
來剝削其他人.
所以這裡是畫清楚界線.
這裡是表明.
上帝自己是一個怎樣的上帝.

$^{761}$是一個連戳其他人.
那些無力的人.
其實也是在說這個十誡.
是針對我們的內心.
其實你這樣看的時候.
一直看到最後.
十條誡命.
其實我們也煩惱.
或者真正的目的不是要我們.
去守校規.
而是告訴我們.
我們是守不到誡命.
怎麼辦.
上帝拯救我們.
上帝用恩典拯救我們.
我們一班有罪的人.
是不能守誡命的.
我們是破誡命的.
但他來拯救我們.
這些思想在十誡裡面.
但我想今天結束的時候.
用這個方法來看.
看清楚十誡.
到底寫給誰.
當然所有神的子民.
都可以應用.
所有以色人都可以應用.
但最主要直接受眾是誰.
是男人.
是做父親的.
三代同堂夾在中間.
有父親有兒女為他打工.
你都是已婚的.
不過犯奸淫.
你生活都幾well off.
有幾輛寶馬.
sorry.
有幾輛車.
有幾間屋等等.
有property.

$^{801}$不要貪戀別人的東西.
有直接的觀眾.
這班人生活裡面有了上帝的祝福.
十誡是要告訴他們.
你們的生活的導向應該是怎樣.
你的orientation應該是怎樣.
不是為自己.
而是為他人.
你想想.
如果一個人民一班的人.
一班的社群.
是這樣去生活的話.
大家來到的時候.
保障你的權益.
我的權益亦都是由你保障.
你生活的導向是為了其他人.
其他人的生活導向是為了我.
這樣彼此去相顧的時候.
這個社會是一個幾幸福的社會.
這個就是十誡的moral vision.
以前我在Whitman College的時候.
有一個.
我想今天美國其中一個研究十誡.
或者生命的一個專家.
Daniel Block.
他有一次在一個課程裡面.
他問了一個問題.
他問有沒有留意.
以色列的社會在摩西五經裡面的舊約神的子民的社會.
他說有一樣東西是我們沒有的.
不是.
有一樣東西是我們有他沒有的.
就是警察.
你找遍五經.
沒有警察這個office.
沒有Israel PD.
沒有的.
他說為什麼.
他說在摩西五經裡面.
那個社群的vision.

$^{841}$生命的意象.
如果人人都是維繫著我和上帝之間的關係.
人人都是維繫著我和其他人的關係.
其實哪需要警察.
不需要警察來執法.
因為大家是一個愛的社會.
所以十誡的要求是什麼.
一班新的子民.
一班按照新的原則.
有新的身份的人.
按照上帝給的誡命.
看清楚這個誡命啟示上帝是一個怎樣的上帝.
一個柳連率.
一個有恩典.
一個有公義真理的上帝.
獨一的上帝要求.
值得我們去敬拜他.
並且在這個誡命裡面.
這個Constitution裡面.
知道上帝是一個怎樣的上帝的時候.
我們大家要用上帝的本性.
用上帝的心懷意念來對待大家.
這個社會的群體.
是一個挺有福氣的群體.
新約的神的子民教會.
這個就是我們的方向.
我想今天挑戰你十個誡命.
你這個星期找一條誡命.
其中一條.
你去祈禱.
因為上帝求你教訓我的心.
幫我按照這個誡命.
這個guideline來生活.
用我們的生命.
愛神的.
愛論社如己.
兩大個summary.
耶穌說.
成為我們的教會.
我們的社群.

$^{881}$的原則.
好不好.
我們同心協力去祈禱.
上主我們多謝你.
因為你是那位拯救我們的神.
我們是一班破碎有罪.
破誡命守不到誡命的人.
我們在軟弱的裡面.
卻有你自己的恩典.
主若求你讓我們不是恩典白白的去守.
讓我們來到你面前.
知道你對我們的新的身份.
頒佈一些新的原則.
幫助我們按照這些新的原則.
來督導引導我們的生活.
用我們的教會.
用我們的生命.
能夠愛神.
能夠盡力的愛人如己.
使這個社群.
使我們這個教會.
這個community.
能夠去反映出上帝的恩典.
幫助我們這個禮拜.
我們記起這些誡命的時候.
讓我們去應勇.
去活出來.
祈禱奉主明救.
願上帝祝福各位.
\newpage



\section{羅馬書}
\label{sec:kf1XXedwCsU}
\textbf{神愛大地 (羅馬書 8) - 陳崇基牧師}
\newline
\newline
連結: \href{https://youtube.com/watch?v=kf1XXedwCsU}{\texttt{https://youtube.com/watch?v=kf1XXedwCsU}} ~~~~ 語音日期: 2025-01-28
\newline
\newline
\hyperref[sec:9TO87He9Dbc]{\small{< < < PREV SERMON < < <}}
~
\hyperref[sec:index_chronic]{\small{[返順時目]}}
~
\hyperref[sec:index_scriptual]{\small{[返順卷目]}}
~
\hyperref[sec:CwJ8DZXjXd8]{\small{> > > NEXT SERMON > > >}}
\newline
\newline
$^{1}$兄弟姐妹新年發財.
恭喜發財.
我們很喜歡錢.
新的一年快要到.
希望大家花點時間去拜年.
還有在後面拍照的時間.
這次我們試一下說一個新的題目.
這個題目其實我從來沒有說過.
在教會.
但我有時候會在團契.
或者茶經.
或者主要學校會提過.
就是和大地和萬物有關.
和創造有關.
上帝是否愛這個大地.
怎樣愛這個大地.
上帝和這個大地的萬物.
受造界,創造界有什麼關係.
所以我想藉著好幾次的講道.
和大家思想一下這個題目.
我們做一點點神學反省.
講神學反省的人會覺得很深.
什麼事?什麼神學反省?.
其實不是這麼深.
其實用我們的信仰或者聖經.
去想一想一些主題.
和信仰有什麼關係?.
我們有沒有地方我們想少了.
或者想不夠的?.
我們需要去到聖經裡面.
給我們一些啟示.
一些啟迪.
讓我們了解某一些主題.
到底應該作為一個基督徒.
作為一個站在聖經之下.
我們相信聖經的.
聖經,信仰對這些主題有什麼關係?.
我想首先做這些神學反省的時候.
先用幾個命題開始.
你回答我到底這個命題對還是不對.

$^{41}$第一.
上帝只是關心我的靈魂.
叫我信耶穌就上天堂.
對不對?.
先不要回答.
先不用開口回答.
因為等下最後會揭曉.
基督的救贖是不是只是救我脫離這個世界?.
離開這個這麼痛苦的世界.
敗壞的世界.
第三.
受造的世界必然都毀滅了.
沒得救了.
所以走為常計.
我們去天堂吧.
第四.
所謂Creation Care.
是與我無關的.
就是管理大地.
Take care 大地.
Take care 環境等等.
與我無關.
先打四個問號.
我們問這個問題的時候.
需要做一點反省.
到底聖經裡面怎麼說.
與我們和創造界有什麼關係?.
如果創造界是.
我們這個大自然.
或者這個世界.
用完即棄的話.
有些東西用完一次就丟掉了.
其實沒有什麼好保育的.
保留它幹嘛?.
反正它都玩完了.
你就不需要去保留它.
或者去.
Take care 它.
或者去護理它.
所以這裡首先有些大前提.

$^{81}$到底是不是好像這樣?.
我們知道原先的天地.
創世記告訴我們一切都是好的.
甚至上帝看見這個世界.
創立這個世界是甚好的.
誰知道人就半變.
然後罪就闖入世界.
然後搞砸了這個世界.
現在的世界就是中間這條線.
基督當然有來過.
在歷史裡面.
基督的出生.
道成活生 死而復活.
這個就是十字架的地方.
我們中間還沒到最後.
我們被罪破壞的現在的世界.
是不是到最後的時候就會中斷?.
全部東西都玩完了.
例如那裡有個核子彈的東西.
爆炸.
全部東西都玩完了.
然後就進入.
所謂新天新地.
基督再回來 箭嘴指下來.
基督再回來的時候.
然後有個新天與新地.
新天與新地.
所以現在的世界.
和將來的新天新地.
就是一個Discontinuity.
一個中斷.
一個絕對的中斷.
這個世界怎麼都玩完了.
你怎麼搞它 怎麼理它 怎麼保護它.
最終都是Game Over.
我只等下一個.
啟示錄最後那兩章說.
那個新天新地 榮耀的新天新地.
完美的新天新地.
就罷就了.

$^{121}$所以今時這個世界.
和將來那個世界.
是中間有個中斷.
完全沒有關係.
如果是這樣的前設下.
現在我們說什麼Creation Care.
或者這個世界需要我們護理等等.
其實沒有什麼用的.
對不對.
最重要的是逃跑.
這裡這麼多天災 這麼多人.
我們快點走 耶穌就叫我們走.
這是第一個前設.
對於這個前設是不是對的.
或者在聖經裡面的一些內容裡面.
怎麼去看的呢 是不是這樣看的呢.
我想告訴你.
我在教會很多年.
多過大半生在教會.
由我年輕的時候開始.
接觸這種課題的時候.
主要是老師 或者牧者都是這樣說的.
我們最重要搶救靈魂.
靈魂 救了你的靈魂.
你有永生.
然後快點逃跑.
這裡不要搞這麼多了.
但我們唯一在地上的使命就是做什麼.
就是救人.
救靈魂.
你的身體都不要了 說真的.
好像衣服一樣 穿爛了就扔掉.
是不是.
我們是這樣接觸的.
這是我們很重的傳統.
到底是不是聖經的價值呢.
第二下面那個就有點不同了.
沒錯 罪是闖入世界.
毀壞了原先的天地.
罪是破壞了的 但是呢.

$^{161}$到耶穌來的時候 這個世界是會更新的.
有一個完美的更新.
然後進入新天新地.
如果這樣看的話 現在的世界和將來.
新天新地就有一個continuity.
一個連續性.
就是你現在做的事.
或者你保存的事.
保留的事.
是將會繼續下去的.
甚至到耶穌回來的時候.
有一個完美的renewal.
這個就是我們需要問的問題.
到底聖經是說一還是說二呢.
如果說一 就好像我所說的.
我們以往的傳統 從小到大接觸的事.
就是傳福音 救靈魂 脫苦海 逃亡.
第二可能就會有些不同.
不同的地方就是我們怎樣去看這個大地.
怎樣去看這個creation.
這個就是我們今天很多時候會說的.
環境保育 關心這個世界.
看護這個世界等等.
我們用幾段經文來思想這個主題.
第一段經文今天選的就是羅馬書第八章.
第八章第十九節開始.
而保羅說什麼呢.
羅馬書第一章到第七章都是說上帝的救恩.
在基督裡面成就的 所謂恩訊清儀.
所以終於到了 終於實現了.
然後第八章可以說是一個高潮.
就是上帝到底要做些什麼呢.
和人類和整個世界要做些什麼呢.
第八章我們來看一看.
我告訴你 我讀這段聖經的時候.
你看看這段聖經怎樣去形容大地受造界.
開心還是不開心呢 你覺得受造界.
八章第十八節.
我認為保羅說.
與將來要顯給我們的榮耀相比.

$^{201}$現在所受的痛苦簡直是微不足道.
這個是當代的本來.
受造之物正熱切等待.
上帝眾子的顯現.
因為受造之物處於虛空之下.
並非出於情願.
而是出於上帝的旨意.
不過受造之物仍然盼望脫離衰敗的束縛.
享受上帝兒女榮耀的自由.
我們知道一切受造之物至今都在新淫.
如同經歷分娩之童.
開心還是不開心.
現在的大地.
這段經文怎樣形容這個創造界.
是開心還是不開心.
不是很開心 是不是.
很簡單 是不是.
因為形容物質的世界.
不是很開心 不是很妥的.
經文用什麼字呢.
用虛空.
虛空你想起什麼.
舊約裡面有一本書叫傳道書.
虛空的虛空其實就是那個字.
就是那個形容詞 其實是嘆息的意思.
它不是出於情願.
大地或萬物受造之界.
受造界是處於一個虛空之下.
在束縛之下 總之不是很妥的底下.
是出於上帝 為什麼呢.
這個背景如果你熟悉聖經的話.
應該大概想到就是創世記第三章.
因為創世記第三章講到人類叛變之後.
帶來一切的咒詛.
包括上帝和人之間的關係破裂.
人和人之間的關係破裂.
然後人和地之間的關係也都破裂.
三章第十七節那裡講.
和亞當講 和人類講.
土地會因你的緣故受咒詛.

$^{241}$你終生的奴婦才能得到土地所吃的.
土地會為你長出泥質.
還有窒泥和荊棘.
你在田間裡面你要吃蔬菜.
但是你要汗流滿面才能有東西吃.
很辛苦的.
然後辛苦到最後就死了.
你是土地而出.
本是塵土 要歸於塵土.
這裡講得簡單一點.
三種的關係.
所謂我們講天倫人倫地倫.
三種的關係.
都很大的張力和很多的矛盾.
這個就是羅馬書第八章的背景.
我們回去看羅馬書第八章的時候.
差不多就是形容這個情況.
不過他用一個塑造界的角度去看.
這個天地 塑造的世界 物質的世界.
在做什麼.
新人和分娩之童.
分娩之童就是生過孩子就知道.
就算未生過孩子都知道是童.
我雖然未生過孩子.
也永遠不會生孩子.
但我知道是最痛的在這個世界上.
這個就是所謂分娩之童.
誰有分娩之童.
女人當然.
大地在新人和分娩之童裡面.
很痛很辛苦.
在叫救命.
是一種二人法.
整個大地 保羅說.
這段經文說很不妥.
好像有病一樣.
你有沒有試過有病.
喘氣 喘喘不息的時候.
然後進了ICU 吸那些很貴的氣.
這個就是那個情況.

$^{281}$他不是在說人.
是在說創造界.
說的創造就是這樣.
當然我們說到這個話題的時候.
我們不一定要贊成那些極端的環保份子.
所有的議題.
我們不一定要這樣.
不過.
聖經裡面告訴我們.
大自然有很多地方是不妥的.
是有不妥的.
這個應該是我們都意識到的事實.
就是氣候的變遷.
氣候暖化.
或者氣候極端化.
就是熱的地方會更熱 破紀錄.
冷的地方會更冷 冷到破紀錄.
這些我們都屢見不鮮.
這個創造 這個環境.
有很多不妥的地方.
有些地方本身是乾燥的.
那是好的.
但是會更加極度的乾燥.
引致山火或者山火出現的時候.
就很難去到撲滅.
有大火不奇怪.
有野火不奇怪.
不過有野火到一個地步.
你是不能夠去撲滅.
或者很難去撲滅.
這個就是我們這一代.
或者到下一代.
都會經常見到的情況.
例如LA的大火.
家園被燒毀損失.
現在大概估計有275個Billion.
275個Billion.
當然是人命關天.
但其實大自然都很痛苦.
大自然的活物都很痛苦.

$^{321}$野獸或者動物.
這張圖片右下角.
就是消防在LA的火場救出一些動物.
有野生的有寵物.
右下角那隻就是一隻小貓.
全身燒到重創.
很可憐.
消防救了牠出來之後給牠水喝.
但這隻貓最後救不到.
很可憐吧.
這些圖片如果我看見.
我都覺得很可憐.
我看見都想流眼淚.
我相信每個人都看到都應該會流眼淚.
大自然好像分娩之痛或者身癮.
劇痛.
是和我們經歷的一樣.
上星期在美國南部.
佛羅里達州北面.
亞利桑那州.
德州.
大雪.
前所未有的暴風雪.
想都沒想過.
將來教會舉辦冬令會可以去那裡.
可以去佛羅里達州.
佛羅里達州是陽光州.
左下角那個不是O'Hare的機場.
是Houston的機場.
這些大暴風雪是前所未見的.
將這些大風雪.
他們下了多少吋雪.
你看到那個數字嗎.
11吋6吋等等.
上星期在西郊.
下了多少吋雪.
3.9吋.
他們比我們多很多.
避寒.
大家要避寒.

$^{361}$現在沒有了.
沒有了.
你想剷雪.
你應該要去下面.
原因是什麼呢.
就是右下角那個MIT的Climate的Polo.
所generate出來的一些簡圖.
本身北極地極的急流.
就像漩渦一樣.
如果強烈的話.
就會將冷風冷流留在北邊.
最多是久不久吹下來的.
但因為暖化的緣故.
就好像你攪東西一樣.
杯子小小的你攪就沒事.
全部咖啡都在裡面.
但你在碗那邊不停攪.
就全部撇出來.
現在的情況就是這樣.
你看到右邊的圖.
左下角就是Texas.
剛剛去到美國的Texas.
南邊等等.
這些是前所未見的東西.
這就是羅馬書第八章.
差不多形容的情況.
所以大地是很深淵的.
大地很痛苦.
就是因為最闖入世界.
環境受到破壞.
第二方面這段經文.
大地好像在等待一樣東西.
萬物在盼望一樣東西.
大地就受苦.
然後保羅說我們都和他一起受苦.
我們和大地好像同命相憐.
但上帝的救恩既然到了.
基督的救恩到了一字失章.
就將新生命帶來.
以致我們作為人類的生命.

$^{401}$能夠得到上帝的救恩.
我們的生命有renewal.
有更新.
不只是救我們靈魂.
是救什麼.
他說我們的身體得到救贖.
我們熱切等待著.
拿到兒子的名分.
即是拿到產業的名分.
我們還未拿到.
但我們有deposit.
有個down payment.
我們希望能夠拿到.
這個就是我們有的盼望.
就是我們的身體.
這裡說的物質的身體.
就是我和你的身體.
我們身體會衰敗.
不過已經有個down payment.
因為耶穌基督從死裡復活.
我們在基督裡的人.
和他一起會share.
一起共享這個救贖.
身體得到救贖.
我們得救的人有這樣的盼望.
不過看得見的盼望就不算是盼望.
誰會盼望已經看不見.
即是現在暫時未見到.
不過我們會盼望將來他顯現出來.
並且他說我們和大地一樣.
好像有軟弱.
有痛苦.
放請我們軟弱的時候.
26節.
有聖靈幫助我們.
我們不知道怎樣禱告.
但聖靈親自用說不出的嘆息.
為我們禱告.
替我們禱告.
洞悉人心.

$^{441}$上帝知道聖靈的意思.
因為聖靈是照著上帝的旨意.
替我們聖徒而去祈求.
萬物就在嘆息之前說.
就是和我們人類一起.
我們是受苦是在嘆息.
就是知道基督已經到了.
所以我們現在嘆息的時候.
不要緊.
有聖靈到了去幫助我們.
在聖靈裡面有那種同理心.
同理心.
不單止是這樣.
這個盼望也是熱切等待.
不是靈魂得救就跟著上天堂的霸走.
而是等候著身體的救贖.
好像耶穌基督到城肉身.
然後死而復活一樣.
死而復活的身體.
現在這個情況是怎樣呢.
那個盼望的情況就好像.
你看長篇劇一樣.
很多集的.
你看到最精彩的一集突然停.
就告訴你.
你再看下去就會看到下一集.
續集大結局.
所以大家很明白的.
這件事.
報劇的試探是很大的.
因為看到中間的時候突然.
你很不開心.
你就沒有忍耐.
你才繼續看下去.
看到最大最後的結局.
這個就是羅馬書第八章所講的.
你留意第八章經常出現一個角色.
就是聖靈.
聖靈經常出現.
二十多次出現.

$^{481}$你即是在舊約裡面提到聖靈.
以賽亞書等等不同的地方.
每逢聖靈出現的時候.
就有創造.
創世記.
以賽亞書的後一半.
當聖靈降臨的時候.
舊恩的日子到了.
聖靈就會啟動新的創造.
所以如果聖靈到了的時候.
聖靈就會在萬物.
加上我們在基督裡面的人.
一直在做更新的工作.
一直在做renew.
所以大地就在盼望將來.
真正的renew.
最後的renew會到.
大結局會到.
是一定會到.
這是第二點.
不單止大地是受苦的新人.
大地也是盼望將來得贖.
可以renew 可以更新.
然後第八章第28節.
就是講大地的管理或者護理.
八章28節.
我們知道上帝使萬事互相效力.
叫愛神的人.
就是上帝按照自己的旨意.
呼召的人得益處.
這段聖經當然很熟悉.
我們經常用一個角度去看.
就是我們是受益人.
有什麼受苦的時候.
有什麼困難來到的時候.
我們就說不要緊的.
為什麼?.
因為萬事互相效力.
你在醫院有人探你.
就說無所謂.

$^{521}$萬事互相效力.
你真的被擊死了.
因為你喘氣的時候.
萬事互相效力.
聖經是不是這樣用的?.
當然最終有這樣的意思.
但是我想告訴你.
互相效力這個字.
synergio.
或者synergio.
是什麼字來的?.
其實就是synergy.
你記得有個字叫synergy.
就是一齊的意思.
synergio.
就是synergy加ergio.
就是我平時辦公室買的椅子.
economic chair.
就是這個字來的.
working.
就是working together的意思.
所以如果你這樣看的時候.
這段聖經大致上不是說我們是最終安慰的受益人.
平安的受益人.
而是說什麼可以這樣意譯.
就是說上帝與那些愛上帝的人同工.
要使萬事成為美好.
這樣看的時候就不是我們是受益人.
我們是workers.
我們是與上帝partner.
使萬事成為美好的.
God is working with us.
through us.
以致我們能夠使all things for good.
這個RSV1971年的版.
就說in everything God works for good.
with those who love Him.
或者God Himself cooperates for good with us.
這個其實就很重要.
因為保羅是猶太人.

$^{561}$猶太人有一個觀念.
神學的觀念到今時今日都在.
叫做鐵觀奧艦.
什麼叫鐵觀奧艦.
奧艦就是世界的意思.
鐵觀就是repair.
就是修車.
repair就是修理一些東西.
就是在世界上追求公義.
然後建立公義.
這個就是神學的概念.
就是我們與上帝一起同工.
for good.
不只是接受上帝給我們的安慰和平安.
而是我們work together.
神要與我們work together.
establish justice.
在這個世界裡面.
今次我們在說的是.
Creation.
如果你這樣看的時候.
其實聖經裡面有說.
是Creation care.
我們怎樣看這個大地.
即是意味著我們不只是旁觀者.
哎呀這裡太多天災人禍了.
我們快點走路.
也不是說我們是受益人.
上帝搞好了都ok了.
他會負責的.
而是我們積極參與.
我們問一個問題.
就是我們可以做些什麼.
在這個Creation裡面.
我們可以做些什麼.
可以幫助到這個Creation.
上個星期四.
佛羅里達州北面的一些海灘.
有很多人去看海龜.
因為海龜飄了上岸.

$^{601}$那些人走去拍照.
然後當局就派一些工作人員去.
他就說這個不是好玩的東西.
因為那些海龜太冷.
下雪了.
太冷了.
好像你在冰水裡游的時候.
你游不到.
你有force bite.
紫色的.
游不到那怎麼辦.
你就飄了.
浪就打到你上岸.
所以幾十隻海龜就衝上岸邊.
快要凍死了.
有些保育人員就去拯救三十隻海龜.
叫那些人不要去拍照.
因為不是動物園.
要打熱線.
要求助.
這個其實有這樣的意味.
就是Creation care.
我們work together.
去establish justice.
使這個大自然能夠繼續延續下去.
我們有份去更新它.
有份去工作.
有份去improve它.
有份去repair它.
很多時候我們在信仰的問答裡面.
我們都會問一個問題.
有時候我們覺得很無聊.
貓狗會不會上天堂.
當然我又未養過貓又未養過狗.
這個對我們不是很切身.
在我們中間有很多是養貓狗的.
當然我們沒有絕對答案.
首先貓狗上天堂已經是有一點.
好像靈魂上天堂那樣.
我們脫苦海那種感覺.

$^{641}$但我相信你思考這個問題的時候.
神學反省.
你應該在這個大前提下去思考.
貓狗將來.
創造界一切不是.
創造界包括這些活物.
不是用完即棄的.
因為聖經裡面告訴我們萬物將來會有更新.
對嗎.
你要在這個大前提下去問這些問題.
起碼有一點指向給我們.
讓我們做多一點反省.
原來這個世界不是玩完就罷.
上帝是用和我們一起去童工.
去repair.
去work it out.
去care for 它.
所以開始的時候四個命題都是錯的.
大地會被更新.
將來新天新地和現在這個天地.
是有連續性的.
而我們今天有個角色.
就可以和上帝童工.
去管理萬物.
creation care.
最後我們有一首回應詩歌.
我想播出來.
就是講到creation care.
很少有這樣的詩歌講這些主題.
所以我知道有一首這樣的詩歌.
我覺得真的很好.
通常我的詩歌都是有難走路.
但很少講creation care.
所以今天我想播這首回應詩歌.
Bella作曲作詞.
他們的MV做得很好.
我想用這個詩歌播出來之後.
做一個回應.
一村一帶的路.
一村一帶一語.

$^{681}$一生連日經處.
堂上你價值標籤.
斑白倦容.
為求踏著這寸土可掌控.
讚卻未夠盡心可放鬆.
一花一土一間.
一刻退去消散.
如燕棲處比爾他.
冰山崩塌.
南瓊活像閣前艷艷色彈.
一號進了夜再吹奏.
世庭海枯下地也喘息.
世界的宿命到了險境.
讓地歸於安寧.
萬物靜養生息.
歸後鳥語凜看.
沙遊水情.
地屬主所賜.
立若之證.
留你可安定.
坐看風景.
受壓的不公平.
雷頭的鐘可放開震鳴.
眾生命期盼眾實風聲.
Zither Harp.
一花一土一間.
一刻退去消散.
如燕棲處比爾他.
冰山崩塌.
南瓊活像閣前艷艷色彈.
一號進了夜再吹奏.
世庭海枯下地也喘息.
世界的宿命到了險境.
讓地歸於安寧.
萬物靜養生息.
歸後鳥語凜看.
沙遊水情.
地屬主所賜.
立若之證.
留你可安定.

$^{721}$坐看風景.
受壓的不公平.
雷頭的鐘可放開震鳴.
眾生命期盼眾實風聲.
Zither Harp.
你我只不過是個客裡.
也有天花灰盡黃地土碎.
如流星一閃即去.
別讚吹多於瑣碎.
讓孤寡可寄居.
你我生於這裡.
我用心得到祈禱.
因為你創造這個世界.
亦是這個世界的美好,美善.
你向我們彰顯.
讓我們能夠活在這個地上.
我們能夠享受.
你給我們的好處.
主,你教導我們.
如何重新思考.
我們和大地之間的關係.
和我們身邊的物質.
身邊的土地.
我們所處的地方.
我們看到的活物.
主教導我們如何好好管理.
因為我們知道你看顧大地.
並且這個大地正向將來.
真正得到救贖.
一個這樣的盼望.
叫我們在此時此地.
我們學好做一個好的管家.
學好我們如何珍惜.
學好我們如何慷慨地思緣.
去更新我們和大地之間的關係.
和其他創造界裡面的活物之間的關係.
主,讓我們能夠更加反映出.
上帝你的愛和你的看顧.
因為大地是屬於你的.
我們生命也是屬於你的.

$^{761}$求主你自己再一次.
藉著我們.
和我們一起動功.
去更新這個地方.
多多奉勸.
多謝各位.
希望下星期能夠見到大家.
\newpage



\section{5:31-32}
\label{sec:CwJ8DZXjXd8}
\textbf{Righteousness as Redefined by Jesus: Marriage \ and  Divorce (Matthew 5\_31-32) - Chee Yie Jong}
\newline
\newline
連結: \href{https://youtube.com/watch?v=CwJ8DZXjXd8}{\texttt{https://youtube.com/watch?v=CwJ8DZXjXd8}} ~~~~ 語音日期: 2025-02-05
\newline
\newline
\hyperref[sec:kf1XXedwCsU]{\small{< < < PREV SERMON < < <}}
~
\hyperref[sec:index_chronic]{\small{[返順時目]}}
~
\hyperref[sec:index_scriptual]{\small{[返順卷目]}}
~
\hyperref[sec:_BgBlG2kkQM]{\small{> > > NEXT SERMON > > >}}
\newline
\newline
$^{1}$All right..
Good morning.
Brothers and sister if you have been with us for a while.
You know that we have been starting to Jesus Sermon on the Mount.
I think we began this series in October last year actually and.
The section that we are in now begins with Jesus.
proclaiming that he has fulfilled the righteous requirement of the law and.
Then using six specific examples or case studies Jesus continues on to redefine.
What true righteousness is in the eyes of the Lord compared to the religious and the cultural norms of his days?.
And so far for the past two months. We have looked at the first two examples on anger and loss and.
How Jesus he points to the root issues behind this action?.
that is a kind of attitude that doesn't treat other people made in the image of God with respect and dignity and.
In our third example today.
We will actually find a shift from the focus on our inner thoughts and desires.
To our ethical actions that impacts how we relate to others and we also find a shift from the decalogue.
Which is the ten words or the ten commandments?.
to the general general law and ethics section in the mosaic law and.
Similar to my last sermon we are actually about to engage in an area of sexual ethics which.
Tends to be an area of high scandal and offense.
It is sensitive because any discussion relate to sexual ethics.
Inevitably engaged in the part of our experience that is deeply personal.
From which we derive a great deal of our identity and also our meaning in life.
And no one likes a personal life to be met away, right?.
and on top of this.
Any kind of ethics by its nature place a certain limitation on our human freedom?.
and.
So for every yes that we say to Jesus and his kingdom ethics.
We are actually saying no to the ways of the world and the ways of our old self.
So today I will do my best to remain faithful to what the Word of God has to say.
when it comes to the topic of divorce and marriage and.
Regardless of where you're coming from and what seasons of life you're in whether you're married.
divorced.
Separated single or widow. I pray that the Lord will speak to each and every one of us.
Tangibly this morning. So let's pray with me together.
Father.
As we come to church and worship you this morning, I just ask that your Holy Spirit will be with us in this moment.
They will work in our hearts as your word is being proclaimed.
Convict us in the areas where we need conviction and.
Comfort us in the areas where we need to comfort.
May or may a word be proclaimed and your name glorified this morning. Amen.

$^{41}$So if you have your Bible with you.
Let's turn to Matthew chapter 5 and we're going to look at 2 verse today verse 31 and 32.
And if you don't have your Bible, I would encourage you to pull out your phone or I also put a couple of copies of Bible.
Behind there on the table. So if you want feel free to.
Just grab a copies and you can follow along.
So we'll be in Matthew chapter 5 verse 31 and 32 and I will be reading from the NIV version.
So verse 31.
It has been said anyone who divorced his wife must give her a certificate of divorce.
But I tell you that anyone who divorced his wife.
except for sexual immorality.
makes her the victim of adultery and.
Marries a divorced woman commits adultery.
This is the Word of God.
So this morning we are going to take a look first at the issue. That is that hand here.
That is the topic of divorce and then we'll take a look at what marriage is.
In another passage in Matthew where Jesus discuss about this topic at length, which is in Matthew 19.
but we are going to.
But we are going to.
Stay in Matthew 5 first. So verse 31.
It has been said anyone who divorced his wife must give her a certificate of divorce.
So what is this certificate of divorce?.
this is actually a piece of paper or a certificate that explains the reason for divorce and it primarily serves to.
function.
First it released the husband's legal claim on his wife.
and this is basically what Joseph intended to do when he found out that Mary was pregnant with Jesus and.
Secondly, it gives the spouse the freedom to remarry.
and.
In the Jewish law only the husband is allowed to give out this certificate of divorce.
however, the rabbis in Jesus day.
They have twisted this law and came up pretty much with a lot of ways where a wife in an unhappy or unfulfilling.
Marriage can sort of like force the hands of the husband to allow them to gain the freedom.
And not only that in Jesus day many of the religious leaders and the cultural in general.
They have pretty much taken this law as a blank check or a license for divorce.
and so this is especially clear in another passage that.
The topic of divorce is being discussed. That is Matthew.
chapter 19.
So in verse 3 it says some Pharisees came to him to test him.
They asked is it lawful for a man to divorce his wife?.
For any and every reason is it lawful for a man to divorce his wife for any and every reason?.
So there seems to be a clear disconnection here right where this law has become a justification for all sort of divorce.

$^{81}$So to get a better idea of where the disconnects happen.
We need to go back to the initial context where this law was given by God to the people of Israel in Deuteronomy 24.
So this is the original context where this law was given.
If a man marries a woman who becomes displeasing to him.
Because he finds something indecent about her and I want you all to pay attention to the word indecent here because we will come back.
To it later. I promise you.
Because he finds something indecent about her and he writes her a certificate of divorce.
Give it to her and sends her from his house and it goes on for a couple more verses.
But basically here in the context of Deuteronomy 24.
this law prohibits a man from remarrying a woman he has previously married and.
Divorced to if she has subsequently been married to another person who has died or divorced her.
But the question that we need to ask here is whether it is a concession or a compromise being made.
Even though we know that marriage according to God's God's design is meant to be permanent.
This is why you always hear this phrase to death do us apart in a marriage vow, right?.
So here in Deuteronomy in an already sinful situation.
God actually permits divorce in light of the fact that the marriage covenant is not permanently. It's not uniformly permanent.
This law was given with the intention of limiting the harm and exploitative tendency.
Not to command divorce.
So because of the sin that led to divorce God actually gave instruction to make sure that the divorced person is being taken care of.
So what is being emphasized here is what may or may not be done in a situation where divorce has already happened.
The law right here is not addressing the moral status of divorce itself.
The emphasis is to regulate in order to minimize the damage not to command or sanction divorce.
So just want to be clear here a permission to divorce does not equal to a command to divorce.
They are two completely different things.
And when it comes to the New Testament.
Jesus does not want us to interpret this law as something that offers a command or prescription for us.
to to divorce.
but this is exactly how the Pharisees that came to Jesus understood it and.
We see this clearly in Matthew 19 again.
some a little bit.
Further along so in Matthew 19 verse 7 again, they asked Jesus.
Why then they asked did Moses command that a man give his wife a certificate of divorce and send her away?.
Jesus replied.
Moses permitted you to divorce your wife. So right here Jesus is correcting them. He's like.
Moses did not command you but he permits you to divorce your wife. Why?.
Because your hearts were hard. So permission is given in these circumstances because sin has already taken place.
But then Jesus continues, but it was not this way from the beginning. I.
Tell you that anyone who divorced his wife except for sexual immorality and marries another woman.
commits adultery.
So verse 9 right here is actually similar to our passage in Matthew chapter 5 on verse 32 just with a slight different circumstances.

$^{121}$So once again, we find this phrase from Jesus.
But I tell you and once again, he continues on to redefines what true righteousness means in the eyes of the Lord.
But I tell you that anyone who divorced his wife except for sexual immorality.
Makes her the victim of adultery and anyone who marries a divorced woman commits adultery.
So, you know Jewish world of academia during Jesus's time there are a lot of different school of thoughts and different thoughts leaders and.
There's actually a lot of ongoing debates that are going on between all these different Jewish.
schools on what actually constitutes a valid ground or valid reasons for divorce and.
A lot of this discussion actually stems from the ambiguity of the word.
Indecent in Deuteronomy the word that I asked you guys to pay attention to just now.
Jerome me 24 it says if a man marries a woman who becomes displeasing to him because he finds something indecent about her.
So on one end of the spectrum, there's the more conservative school the school of Shami.
They say only sexual immorality and immodesty.
Constitute a valid reason for divorce and on the other end of the spectrum. There is the school of Hillel.
They basically say any any reason is valid if.
Your wife is not a good housewife and she accidentally spoiled a meal that constitutes enough reason for you to divorce her.
And there are also other Jewish thought leaders like Rabbi Akiva.
Who basically say if you find someone that is more attractive and she's available that?.
Constitutes valid reasons for you to divorce.
And while all these debates are going on.
Jesus he transcends all of them by saying that divorce is not part of God's original plan.
Therefore remarriage with or without the certificate of divorce constitute adultery in the eyes of the Lord.
However, there are two clear.
exceptions that we can find from the Bible the first one is sexual immorality where we finds here in Matthew chapter 5 and 19 and.
the second one is.
spousal desertion or spousal abandonment and we find that in first Corinthians 7 15 and.
The second cases it basically described a situation in an unequally yoke marriage where one of the spouse is a believers.
and another one is not and.
If the non-believing spouse choose to leave the believing partner.
Paul this is Paul speaking. He say in that case that is valid.
But regardless we should view these two exception as permission not requirement.
And in our passage today.
Jesus did not actually say that in such a situation as husband is required to divorce his wife.
Quite the opposite.
What is implied by Jesus here? Is that even in the case of sexual immorality?.
you are not required to divorce your wife and.
It is possible for the marriage covenant to continue.
In fact because God's design for marriage is meant to be permanent.
Divorce should always be the last resort and this is quite a scandalous claim during Jesus's day.
And at this point some of you might ask what about other situations like.
domestic abuse or.

$^{161}$violence in a marriage relationship.
Of course God has given us the wisdom and his spirit to guide us in various situation.
So in some cases of violence and abuse.
It might be wise for the spouse to separate or even divorce in order to limit the harm and care for the person.
And I truly believe that there is always grace in this kind of situation.
However, this doesn't change the truth that the marriage covenant between a man and a woman is.
meant to be a lifelong covenant and.
This also doesn't change the truth that our God is the loving father who care deeply who cares deeply.
For his children and he desires for each and every one of us to experience a flourishing life.
So as the church as believers, we need to use our God-given wisdom to approach these issues.
With the intention of loving and caring for the person at the same time without compromising our faith.
So we have touched upon the topic on divorce and now we're going to take a look at God's design for marriage.
So we're gonna stay in Matthew 19 for this. So let's go back to Matthew 19 again in verse 3.
When the Pharisees asked Jesus if it is lawful for a man to divorce his wife for any and every reason.
I want you all to pay attention to how Jesus respond to them.
He didn't focus on what constitute legal grounds for divorce.
But instead what he focused on is God's design for marriage.
So Matthew 19 starting in verse 4.
Haven't you read he replied that at the beginning the Creator made them male and female.
this is in reference to Genesis 1 27 and.
say for this reason a man will leave his father and mother and be united to his wife and.
The two will become one flesh and this is in reference to Genesis 2 24.
So they're no longer two but one flesh and therefore what God has joined together.
Let no one separate.
So what can we learn about marriage from Jesus's response here?.
If you're taking note, I would encourage you to write this down.
Marriage is a divine institution that is directed towards the goal of our lifelong discipleship to Jesus.
Again, marriage is a divine institution that is directed towards the goal of our lifelong discipleship to Jesus.
So we're gonna look for us marriage as a divine institution. I.
Don't know if any of you actually enjoy.
Legos and there's no shame in the meeting yet.
Even though you're an adult because I actually still enjoy building structures and play around with Lego.
When I was getting my MD fat tats, I live with a friend a friend's family up in the north suburbs.
And one of the unique thing about this family is that they are a family of builders and architects.
In fact, they have actually helped a couple of their friends and family.
built like the house from the ground up all the way from design to like actually building it and making it happen and.
The house that I live with them at that point the dad pretty much designed it and build it himself as well.
And somehow this builder DNA is being passed down to their grandchildren.
Claire and Russell as well because every time I hang out with them, they just love playing with their Legos.
For a certain period of time whenever they get their hands on their Lego, they actually like to play what they call Lego zoo.

$^{201}$And it is always fascinating seeing Claire and Russell used to creative freedom to build their Lego zoo.
One time the younger brother Russell.
He pretty much put a slide right in the middle of the zoo and not only that the slide actually leads right into the lion.
exhibits.
As an adult when I see the outside, okay, that is an issue right there. So I asked Russell.
Hey Russell, what is the slide for and he just answered me like genuinely?.
Well when the people get tired of walking they can just take the slides.
So even though me as an adult I see some issues with the Lego zoo.
But since this Lego zoo is the quote-unquote creation.
They pretty much have the freedom to build it whatever way that makes sense to them right as a three-year-old at that point.
But on the other hand.
There's another type of Lego that actually comes in a specific set.
With specific step-by-step instruction on how to build it and you can find all sort of different Lego sets at this point from Star Wars.
To Marvel to Harry Potter you name it.
And so a few Christmas ago.
I actually got an architectural Lego set of Singapore from this family and actually had a blast of building it.
I actually have a pictures on to show you so this is the Lego set that I received from this family and.
One thing that I noticed about this Lego set is how much details and thoughts they're being put behind the design of this Lego set.
each single piece of the Lego the colors the shapes.
the orientation.
Where you put them?.
They're all being intricately designed to be built into the skyline of Singapore. I.
Remember at some point I actually put a piece a right piece in the right place.
But I sort of like put it in the wrong orientation.
And I didn't notice it until later when I tried to put another piece on top of it.
And it just doesn't fit so I pretty much have to redo everything from that piece onward.
And so for me to successfully build this Lego set.
I pretty much have to follow like step-by-step instruction from the beginning to the end.
And so when it comes to marriage there are generally two different ways that we can approach it as well.
The first approach is kind of like see it as a social construct.
Which is pretty much like the free play Lego that Claire and Russell is playing with and they can pretty much do whatever.
They want with the Lego and the second approach is.
Seeing marriage as a divine institution.
So it's more like this kind of Lego set where a creator design is specifically to serve a specific purpose.
So if marriage is a social construct.
It means that there is no inherent nature and order to it.
But we as human we collectively decide and agree upon what marriage should be and what is the meaning of marriage?.
One good example of a social construct is cryptocurrency if you think of all the different coins.
Bitcoins or I don't know. What are some other coins these coins these virtual currencies?.
They don't really have any inherent value to it.

$^{241}$but the values come from we as human kind of like assign something to it and.
It also depends on a lot of different factors like the market or the supply and the demand.
But on the other hand if marriage is a divine institution it means that there's an integrity that exists.
Independently on how we as human instrumentalize it.
In other words, there's a certain norm and design to marriage in which we are accountable for.
Similar to when I talked about in my last sermon with kind of like the grain of the universe.
That there's a certain order to it that is being designed by God.
So what are the implication with marriage being one or the other if.
Marriage is a social construct that means that we as human.
We actually have the freedom to choose what form it takes and use it as a tool to pursue our own personal interests.
But if marriage is a divine institution.
It means that there's an inherent structure to marriage that is part of the scent that is part of the created order.
regardless of how we interpret or.
instrumentalize it of.
course, we as human we still have the free will to choose to go against the grain of the universe and.
Violate the created order, but ultimately it will not lead to a flourishing life.
And if we take a look at Jesus response, I think it would tell us a pretty clear answers right here.
Jesus he goes all the way back to Genesis to the creation account.
so notice in verse 4 haven't you read he replied that at the beginning the.
creator made them male and female and.
Say for this reason a man will leave his father and mother and be united to his wife and the two will become one flesh.
so this means that marriage between a man and a woman is instituted by God the creator and.
When it is at the beginning that is since the creation of the world.
And we as human we are actually part of the creation as well.
So we don't really have an authority to simply decide what we want marriage to be.
So what is marriage first it is a divine institution.
That means there are certain orders and norms to it that we as human are accountable for.
But second marriage is directed towards a greater goal.
So if you ask the people around you what is the purpose of marriage you might get an answer like affection.
companionships or.
procreation.
The church historically has seen some of this as the purpose of marriage include procreation.
companionship remedy for loss and also social stability and.
Marriage and family in general they are being situated in this historical narrative where we as human.
Carry the responsibility to extend the kingdom of God to a certain extent.
We find this in Genesis 1 28 where God tell command human to be fruitful and multiply.
Few the earth and subdue it and rule over the rest of the creation. This is what call.
This is what we call the cultural mandate.
All these purpose are true and good.
But ultimately I would say marriage is a way that we can be formed into Christ likeness.

$^{281}$This bring us to my third point that marriage is a way that we can pursue a lifelong discipleship.
to become more like Jesus and.
I just want to be clear here.
What I'm not saying is that your discipleship to Jesus is deficient if you are single or you're not married.
Paul speaks about singleness in 1st Corinthians 7 as well and Jesus himself. He remains single his entire life.
So there's nothing wrong with not being married.
But right here marriage is a form of discipleship.
Because at the core of a godly marriage is a kind of sacrificial love that reflects the gospel and its model after Jesus.
This is what Jesus meant when he says that in the marriage covenant the two will become one flesh.
And there are no longer two but one flesh.
So there are no longer two single individual pursuing after their own interests.
But one single Union pursuing after a common goal.
I'm sure all the married couples in this room can attest to this.
Therefore a marriage to work out. It usually requires a lot of sacrifice and self-denials, right?.
Maybe one of you would like to sleep at 65 degree at night.
But your spouse like to sleep at 75 degree Fahrenheit at night. So for your marriage to work out.
I'm sure either one of you or both of you have to give in in some ways, right?.
And this is just one of the small examples. There are a lot more important decisions like.
How many kids do you want to have?.
Where where do you want to live? What neighborhood you want to live in? What school do you want to send your kids to?.
These are all big decisions that as the two person become one they have to decide upon as one single Union.
And this in some ways is what discipleship is all about that. We are called to deny ourself.
including our own self-interest to take up our cross and to follow after Jesus himself and.
Marriage is an arena where our discipleship to Jesus can be refined and being lived out.
And.
Moreover marriage is a lifelong journey because once again the marriage covenant between a husband and a wife.
Is a sacred covenant that is meant to be permanent.
And this is what Jesus meant when he say in verse 6 therefore what god has joined together.
Let no one separate.
So again, what is marriage?.
Marriage is a divine institution that is directed towards the goal of our lifelong discipleship to Jesus.
So if you have been following along this sermon series on Jesus sermon on the mouth so far.
over the course of my last few sermon.
You will probably started to notice a common threads in Jesus teaching.
Where he raises the standard of righteousness?.
And also intensifies the moral demand represented by this law.
And if this makes you feel like you can never really approach Jesus with such a high demand that he has.
There's a good news for us.
Unlike the pharisees and the teachers of the law.
The paradox of Jesus and some of the most scandalous realities about his ministry.

$^{321}$Despite his demand for righteousness is that the people whom Jesus chose to actively.
be with and spend time with.
Were usually those that are considered as sinners and outcasts by the society.
And these are the people that Jesus came and died for.
So I want to acknowledge that there might be some of us in this room.
Who have personally been through a divorce?.
Or you're currently going through a seemingly hopeless situation in your marriage.
Or maybe some of us.
Or maybe some of us you have been directly or indirectly impacted by a broken marriage in your family of origin.
Or maybe for some of us you're in a season of singleness where you're just struggling and feeling alone.
Whatever your experience is coming to church and hearing this message this morning.
I want to be abundantly clear that the intention behind these teachings from Jesus is not to condemn.
But it is actually out of his love.
For us to experience the good and blessed life.
This is why the entire sermon on the mount actually begins with blessing.
Blessed are the poor in spirit.
And regardless of your current struggles and past experience in your marriage or singleness.
You don't have to carry the weight of your guilt and shame anymore.
because Jesus.
Through his teachings in the sermon of mount.
It's actually inviting all of us into a new way of living.
That will lead to a truly blessed life.
So before I end my sermon, I know we just touched upon a subject that can be heavy for some of us here in the room.
So now I want to actually create some space for us to respond to the lord and do whatever we need to do with him.
And i'm going to do so actually by leading us through a spiritual practice.
together.
so.
This might be new to some of you.
But if you can just find a position or if you want to move to another seats where you feel more comfortable with.
Just find a spot that you're comfortable in and just relax.
And if you're comfortable.
You can close your eyes and take a few deep breaths.
Now I want you to imagine that you and jesus are in a room together.
And jesus is not all the way across the room from you.
But he is either standing or sitting right in front of you or just right beside of you.
Take a few more deep breath listen to this invitation from jesus.
Are you tired.
Worn out.
Burn out on religion.
Come to me get away with me and you will recover your life.

$^{361}$I will show you how to take a real rest.
Walk with me and work with me watch how I do it.
Learn to enforce freedoms of grace.
I won't lay anything heavy or ill-fitting on you.
Keep company with me and you will learn to live freely and lightly.
So.
If you want to just take a moment at your seat, I also put out the passage up on the screen.
this.
In your own pace at your own time. Just read through this passage again. This is matthew 11.
28 to 30, but it is the message version.
Read through this passage again slowly and reflect on the invitation from jesus.
To come to him.
To get away with him.
To walk alongside him.
And to be with him.
And i'm going to close this out in a moment with a prayer.
So.
Jesus.
We thank you for you are a god who is close and near to the brokenhearted.
We thank you.
For your yoke and your burden is not heavy.
But it is easy and light.
As we receive your word this morning.
Help us lord.
To be willing to also receive this invitation from you.
To.
To be with you.
To take up your yoke which is light and free.
And as we continue to live out our christian faith.
Help us to know that jesus.
You are right there with us in each and singing each and every step of the way.
That you are with us.
Comforting us strengthen us.
And healing us.
So as we go go on to the rest of the week, I just ask that lord you will refresh us.
Give us the rest that we need so that we can continue to live out.
Um the life that you have called us to live to live out the purpose that you have given us.
To be a light and the salt to the world.
To the darkness that is around us. So father be with us.
as.

$^{401}$As we walk out from church this morning.
Pray this in the son's name. Amen.
So, uh.
Thank you all if you need.
More time to process feel free to stay in the room if you want to find someone to pray.
Feel free to do so as well.
but for the rest of us, um, you are.
Dismissed you can go in peace and.
We have some dim sum out there.
For us to enjoy as well.
\newpage



\section{約翰福音 3:14-15}
\label{sec:_BgBlG2kkQM}
\textbf{耶穌——賜生命的蛇 (約翰福音 3\_14-15) - 何建邦傳道}
\newline
\newline
連結: \href{https://youtube.com/watch?v=-BgBlG2kkQM}{\texttt{https://youtube.com/watch?v=-BgBlG2kkQM}} ~~~~ 語音日期: 2025-02-05
\newline
\newline
\hyperref[sec:CwJ8DZXjXd8]{\small{< < < PREV SERMON < < <}}
~
\hyperref[sec:index_chronic]{\small{[返順時目]}}
~
\hyperref[sec:index_scriptual]{\small{[返順卷目]}}
~
\hyperref[sec:cIQZdCWcAV4]{\small{> > > NEXT SERMON > > >}}
\newline
\newline
約翰福音 3:14-15
\newline
\begin{longtable}{cl}
\hline
\hline
章節 & 經文 (和合本修訂版)\\
\hline
3:14 & \begin{tabularx}{0.7\textwidth}{X} 摩西在曠野怎樣舉蛇，人子也必須照樣被舉起來， \end{tabularx} \\ \\ \relax
3:15 & \begin{tabularx}{0.7\textwidth}{X} 要使一切信他的人都得永生。 \end{tabularx} \\ \\
[1ex]
\hline
\hline
\end{longtable}
$^{1}$早晨,新年快樂,恭喜發財.
我覺得呢,成日都話要積聚財寶係天上,積在地上都得ge .
不過最緊要係咩呢?就要奉獻.
呢個係好緊要,大家可以積聚財寶係地上.
恭祝大家發財完之後奉獻俾教會.
一齊貢獻係神ge 國度裡面,因為神ge 國度係需要錢ge .
講下先,今日其實本來我預備jor 一篇講道.
本來今日係想講新約中的情緒2,恐懼ge .
所以其實我係有想延續呢個主題.
但係我突然間係上個星期先發現原來我今日講道ge 日子係年初四.
所以我就突然間,昨日下午就預備過一篇新ge 道.
為jor 應節D ,就係同今年新年有關ge .
咁呢,我好想講今年係咩年?蛇年呀嘛,係咪呀?.
所以今日呢,我好想講一篇道係同蛇有關ge .
為jor 應節,今年係蛇年.
咁呢,今日我ge 講道題目就係呢.
耶穌似生命的蛇.
咁呢,係開始講道之前呢,想做一個即係簡單,非標準ge 統計.
我想問,邊個覺得蛇係正面ge ?舉手.
蛇係正面ge ,舉手.
la ,要誠實作答.
邊個覺得蛇係負面ge ?舉手.
蛇係負面ge ,舉手.
la ,大比數勝出都係覺得蛇係負面ge ,係咪?.
特別我地基督徒,係咪?.
當我地諗蛇ge 時候,就會諗起咩呀?.
撒旦呀,係咪呀?.
gwo 條龍啦,即係啟示鹿gwo D 啦.
仲有咩呀?即係伊甸園去引誘亞當同埋夏娃犯罪ge gwo 條蛇.
即係全部呢,係基督教ge 傳統裡面呢.
我地彷彿都會覺得蛇呢,就係唔好ge .
咁呢,蛇呢,又係我地中國人ge 新年,即係十二新潮裡面其中一個新潮.
咁其實蛇,係咪全部都係負面ge 呢?.
係聖經裡面有冇一D 關於蛇係正面ge 經文ge 呢?.
咁我想講呢,其實係有ge .
係新約裡面呢,其實有兩段經文呢,係講蛇係正面ge .
第一段呢,就係馬太福音裡面講咩呀?.
頭先,即係報告有講過啦,靈巧像蛇啦,係咪?.
而第二段呢,就係今日我要講ge 經文.
就係約翰福音三章十四到十五節ge 呢段經文.

$^{41}$咁呢段經文係講咩呢?我地一齊睇下囉.
係呢段經文係咁講ge ,約翰福音三章十四到十五節.
呢度就係講緊耶穌同一個猶太人ge 教師.
尼哥德姆對話當中ge 其中一段ge 經文.
咁樣講ge ,係話耶穌咁講.
摩西係曠野點樣將蛇舉起.
人子,人子即係耶穌自己啦.
人子也必照樣被舉起來.
使所有相信他的人道德榮臣.
咁呢度呢,耶穌好明顯咁樣去將自己比喻為咩呀?.
比喻為摩西係曠野裡面舉起ge gwo 條蛇.
咁所以我地平時係教會聽呢,好少會話耶穌係蛇架嘛,係咪?.
所以今日呢,我好想同大家講呢.
其實聖經係有講過耶穌係蛇架.
就係我地今日睇ge 呢段經文.
耶穌就係gwo 條被舉起ge 蛇啦,ok?.
咁所以呢,我地話耶穌係蛇呢,就唔一定係異端架喎.
因為係約翰福音裡面真係咁講.
耶穌就係gwo 條被摩西係曠野裡面舉起ge 蛇.
咁所以今日我講到ge 題目就係話.
耶穌似生命的蛇.
咁我地返返去呢度啦.
係呢度呢,耶穌將自己比喻為被舉起ge 蛇.
其實呢個係一個好出名ge 舊約ge 典故.
知唔知呢個典故係黎自邊一段經文呀?.
就係黎自民數記ge 廿一章四到八節.
gwo 度咁樣講,就係講緊呢.
以色列人係曠野呢,你都知啦.
以色列人係曠野最出名係咩呀?.
最出名係呀,犯罪啦,同埋咩呀?.
就係發牢騷啦,抱怨啦.
咁呢個故事就係咁ge .
佢地係荷爾山起行ge 時候呢,係曠野啦.
以色列人呢就發牢騷啦.
又抱怨上帝埋怨上帝啦.
就話咩呀?.
喂,呢度又冇ye 食又冇水,係咪好港女呀?.
又冇ye 食又冇水,又行咁耐.
真係好辛苦呀咁樣啦.
咁於是點算呀?.

$^{81}$於是呢,神就頂佢地唔順啦,你地班港女.
就派jor 一堆蛇呀,叫做火蛇啦.
就進入D 百姓gwo 度.
咁呢D 火蛇就點呀?就咬佢地啦.
於是呢,以色列人就死jor 好多啦.
咁呢,於是呢班埋怨ge 百姓呢.
就去到摩西gwo 度就話.
哎呀,唔好意思呀.
我地埋怨神又埋怨你.
真係好唔好意思,我犯罪啦.
你可唔可以幫我地祈下禱,求下神.
叫呢D 蛇離開我地呀.
咁於是呢,摩西就為百姓祈禱啦.
於是神就叫摩西就話啦.
你呢就去製造一條蛇.
就掛係一支竿上面.
咁呢,凡係被咬ge 人呢.
只要一望到呢支竿上面gwo 條蛇呢.
就可以被醫治,就可以存活啦.
咁於是呢,摩西就照耶和華ge 吩咐呢.
就做jor 一條銅蛇,掛係竿上面.
凡係被咬ge 人呢,睇到呢條蛇呢.
就存活啦.
咁所以耶穌呢,係同尼哥德姆ge 對話裡面呢.
講緊呢,人只就好似呢條被舉起ge 蛇呢.
就係講緊呢一條蛇啦.
就係摩西係曠野裡面呢,做ge 一條蛇.
任何ge 以色列人,犯罪ge 以色列人.
只要望呢一條蛇,佢就可以存活啦.
咁所以耶穌呢,引用呢個比喻呀.
或者引用呢個典故其實係講咩呀.
就係想講呢,佢就係gwo 條被舉起ge 蛇.
任何犯罪ge 人.
只要憑住信心,相信.
去仰望呢個被舉起ge 耶穌.
呢一條蛇,就可以存活啦.
咁所以你見到呢,呢個對應ge .
即係呢個典故係對應耶穌.
黎到呢個世界上面要做ge 事.
佢黎到呢個世界上面,係要成為gwo 條被舉起ge 蛇.

$^{121}$就好似摩西ge gwo 條蛇一樣.
讓每一個仰望佢ge 人.
相信佢ge 人,得到生命,可以存活.
咁呢,但係我地要問一個問題.
即係,即係聖經有好多動物ge 嘛,係咪?.
即係點解要無哩啦間,要用蛇去形容自己呢,耶穌.
或者約翰福音點解要用,即係咁多動物唔用.
其實有更多更好的動物,例如咩?.
獅子係咪更好?.
即係獅子呀,呢D 更加威武,更加靚仔ge 動物.
可以用架嘛,呢D 典故.
救藥好多動物ge 典故.
點解約翰福音,要係呢一度用蛇ge 典故呢?.
我地要問呢個問題.
點解咁多動物唔用,要係呢度用蛇ge 典故?.
咁今日好想同大家介紹下一個神明呀.
呢個神明叫做Aeschypius.
咁為jor 簡單D 去叫佢,我叫佢做Eason.
Ok?係啦,Eason.
我以後都會叫佢Eason架喇,係今篇ge 講道裡面簡單D .
唔使每一次都Aeschypius,Aeschypius,Aeschypius咁樣讀佢ge 名.
總之佢就係叫Eason.
咁呢,呢個神有D 咩特別呢?.
呢個神係係當時ge 希臘羅馬ge 世界裡面一個好出名ge 神明.
咁呢,呢個神明有D 咩特別呢?.
就係你睇到左邊,係希臘羅馬世界裡面.
佢ge 雕像係有一個特色ge .
就係佢會拎住一支杖,gwo 支杖上面有一條蛇.
有一條蛇繞住支杖.
係咪就好似摩西gwo 支杖,有條蛇繞係上面.
咁呢,呢個就係當時耶穌時代.
gwo D 羅馬人或者希臘人成日會拜一個好出名ge 一個民間ge 神明.
Aeschypius.
咁呢,佢特別ge 地方係咩呢?.
我叫佢做Eason.
就係因為佢ge 能力係同醫治有關ge .
就係呢,佢最出名就係.
百病能醫,能醫百病.
即係任何人只要有病去拜佢.
就可以被醫治.

$^{161}$咁所以呢,佢當時係一個好出名ge 醫神.
有好多ge 寺廟,有好多ge 同佢有關ge 敬拜都係同醫治有關ge .
咁呢,你見到呢,係我地ge 現代ge 一D 醫療機構.
即係十字車,平時你見到救護車就見到個標誌.
就係有一個杖,有條蛇,係咪?.
或者呢,我地ge 世衛.
世衛組織,世界貿易組織.
其實都係用緊呢個Easonge 一個symbol.
就係有一支杖,然後有條蛇卷係上面.
咁呢,呢個Eason除jor 有一支杖同埋一條蛇之外.
係好多ge 一D 考古ge 發現裡面.
一D ge 寺廟裡面都發現呢.
D 希臘羅馬人會用蛇,用蛇去代表Eason.
即係蛇ge 形象,佢就係用一個蛇ge 形象去代表佢.
咁如果明白呢度ge 話,我地就可能會明多D .
點解約翰福音要形容耶穌做蛇?.
係因為約翰福音當時都係面對緊一班ge .
一D ge 希臘羅馬ge 讀者.
咁我地去睇下,其實呢.
呢個Eason同約翰福音所描述ge 耶穌.
都好多平衡同埋相似ge 地方.
我地睇少少.
第一呢就係呢,呢個Eason係太陽神阿波羅ge 仔.
而耶穌就係神ge 兒子.
第二呢就係呢,呢個Eason係.
佢係神話裡面到處行醫.
而耶穌係約翰福音裡面都做jor 好多醫治ge 神蹟.
第三就係呢,呢個Eason係因為呢.
曾經醫治人,甚至使人復活而激嬲jor 眾神之神宙斯.
所以就被宙斯激殺.
如果大家記得約翰福音第十章gwo 度講到咩.
耶穌使拉薩路復活,然後激怒jor 猶太人.
然後猶太人係因為耶穌復活拉薩路,最後要殺佢.
大家如果記得約翰福音第十章.
跟住第四點係咩呢?就係呢.
宙斯用雷電激殺jor Eason.
後來又復活返佢.
並且將佢高升成為天上ge 一個星座叫做蛇夫座.
即係我地今日睇一個星座叫做蛇夫座.
其實佢ge 典故就係黎自Ekipius,呢個Eason.

$^{201}$而耶穌同樣都係點呀?.
就係神使耶穌掟十字架而死.
而最後又復活佢,使佢高升.
而第五點係咩呢?就係呢.
呢個Easongwo D 信仰Easonge 人.
拜Easonge 人會稱呼Eason做救主.
因為佢能醫百病,會稱呼佢做救主.
而耶穌係約翰福音第四章裡面.
都被當時ge 撒瑪利亞人.
稱呼佢做世界的救主.
咁所以我地就會見到呢.
其實呢原來呢個神明呢個Eason.
Ekipius同耶穌都有.
即係約翰福音所記載ge 耶穌.
有好多平衡ge 地方.
咁我地特別去睇一段經文.
就會更加睇到呢.
係呢度可能係真係反映緊呢.
當時敬拜Easonge 一D ge 背景.
係呢度就係約翰福音ge 五章一至八節.
gwo 度係講緊耶穌一個醫治ge 神蹟.
gwo 度講咩呢就係.
係去到猶太人ge 一個節期.
耶穌就去耶路撒冷.
係耶路撒冷靠近陽門ge 一個池子.
希伯來話叫做不士大.
旁邊有五個狼子.
狼子裡面擋住好多病人.
有核眼ge ,唾ge ,癱瘓ge .
咁gwo 度有一個人就病jor 三十八年.
耶穌見佢訓係度.
見佢病jor 好耐就問佢.
你想唔想好返.
於是gwo 個人就話.
先生水浴ge 時候.
冇人將我擺係池子裡面.
我正要去ge 時候.
別人就俾我先下去.
於是耶穌就同佢講.
起黎拿你ge 玉子即刻行走.

$^{241}$gwo 個人就即刻好返.
就即刻拿自己ge 玉子行走.
咁其實係呢個故事裡面.
係呼應緊一個敬拜醫神ge 背景.
點解我會咁講呢.
其實就係最近ge 發現.
都唔係最近.
即係現代ge 發現.
呢度就形容緊.
係耶路撒冷靠近陽門gwo 度.
有一個池.
gwo 個池叫做不士大.
有五個狼子.
其實D 考古學家.
係耶路撒冷裡面.
就搵到一個類似.
約翰福音所描述ge .
一個叫做不士大的地方.
就係係陽門附近.
有一個池.
同時亦都有五個狼子.
其實係以前好多人就會覺得.
約翰福音腦作ge D ye .
即係約翰福音.
好多時都會被人覺得.
係四卷福音書裡面.
最吹水.
最腦作ge 聖經.
即係成日都會覺得.
呢個作者係亂咁.
作D 地方出黎.
但係係近代.
真係有考古發現.
發現到一個地方.
係同約翰福音ge .
呢一個經文.
講ge 描述係.
差唔多.
類似ge 地方.
就係又有個池.

$^{281}$又近陽門.
又有五個狼子.
所以D 學者.
D 考古學家就會覺得.
好大機會呢.
呢一個耶路撒冷ge 地方.
就係約翰福音.
講緊ge 呢一個地方.
不事大.
咁而特別ge 地方係咩呢.
大家可能會覺得.
哇 耶路撒冷好神聖呀.
聖地黎架嘛.
聖城黎架嘛.
經拜耶和華架嘛.
你會覺得好神聖.
但係特別ge 地方係咩呢.
就係發現呢.
係一個池子呢.
有五個狼子ge 池子裡面.
發現呢.
仲有呢一舊ye .
呢幅圖ge 呢舊ye .
呢舊ye 係咩呢.
係一個花瓶咁ge 物體.
一個器皿黎ge .
呢個器皿見唔見到.
有D 咩特別呀.
上面捲住有D 係D 咩黎架知唔知.
係蛇黎ge .
係一個咁樣耶路撒冷ge 池子裡面.
發現一個咁樣ge 花瓶.
有好多條蛇捲住.
咁一D 學者仲發現呢.
其實呢個花瓶呢.
即係唔係花瓶.
其實係一個器皿.
其實係D 咩黎ge .
係勁大醫神ge 一D 物品黎ge .
一D 係醫神ge 廟.

$^{321}$或者醫神ge 池子度會有ge ye .
就係有D 蛇捲住.
我頭先咪講過呢.
蛇係代表咩呀.
Aescypius.
gwo 個神明ge 代表黎.
咁所以呢舊ye 呢.
其實呢係當時呢.
勁大醫神ge 一D ge 物品.
咁呢我呢度唔係想講呢.
耶路撒冷裡面呢.
有一個醫神ge 廟.
我唔係想咁講.
不過呢好有可能呢.
裡面係有少少混雜ge ye .
即係裡面係有一D 呢.
即係將呢座池呢.
連繫到呢醫神ge 醫治.
ge 一個咁樣ge 背景.
你明我講咩o下.
所以點解係裡面呢.
會發現一D 咁樣ge 物品.
咁呢特別ge 地方係D 咩呢.
呢度亦都係對應緊呢.
醫神ge 一D 呢醫治ge 習俗.
你知唔知係當時ge 羅馬帝國裡面呢.
D 人拜醫神呢.
係點樣得醫治架.
佢地點樣得醫治呀.
通常呢係醫神ge 廟出面呢.
會有個池呀ok.
佢地有兩個stepge .
第一個步驟係咩呢.
就要入個池裡面浸D 水.
即係佢地覺得gwo D 池ge 水.
係聖水黎架.
聖水你諗下幾污糟.
幾多個人會入去浸過.
係唔會換gwo D 水ok.
係古代.

$^{361}$一潭好多腦泥ge 水.
總之係聖水.
就係入去浸咁樣.
浸完之後呢.
gwo 個病人要點呀.
就梗係要多多益善呀.
要奉獻ge .
就係俾錢.
俾錢呢.
個負責去營運呢個池ge 寺廟.
即係俾錢啦ok.
向醫神奉獻.
然後就祈禱.
求醫神呢令我得醫治.
因為個廟裡面ge 醫神ge 祭司呢.
就會按手係gwo 個人度.
你會得醫治啦咁樣.
咁呢個呢就係成個去拜醫神.
得醫治ge 過程.
ok 明架.
咁呢係呢度有D 類似ge .
就係呢有個池啦.
然之後點呀.
呢個人呢係話咩呀.
點解佢醫唔返佢覺得自己.
點解佢好唔到佢覺得自己.
係因為佢癱架嘛係咩呀.
即係佢落唔到個池度呀.
每一次佢想落個池度ge 時候呢.
想浸D 水呢.
俾醫治ge 時候.
都有人快過佢落去.
佢就覺得係因為咁樣呢.
不能被醫治啦.
ok 咁但係耶穌係點樣醫治佢ge .
耶穌只係講一句ge .
你起身佢就好返啦.
ok 咁所以呢度呢.
有一個好大ge 可能性呢.
就係呢呢度係對比緊呢.

$^{401}$耶穌ge 醫治同埋醫神ge 醫治.
耶穌ge 醫治同醫神ge 醫治.
有咩唔同呀.
醫神係透過浸D 水去醫返個人.
但係耶穌係度唔使你浸水ge .
我同你講一句.
你起身即刻起身.
邊個勁D 呀.
醫神勁D 靠腦泥水勁D 呀.
定係靠耶穌ge 說話勁D 呀.
耶穌勁D 係咪呀.
所以呢度呢好可能呢.
楊豐英ge 作者呢.
就係透過去呼應緊當時呢.
一D 希臘羅馬讀者.
好熟悉ge 醫神ge 文化.
醫神ge 醫者文化呢.
去強調緊耶穌.
係勁過你地呢D 民間ge 神明.
勁過呢個靠腦泥水醫治ge 醫神.
耶穌勁D .
唔使靠水呀.
一句說話就醫好呢個人.
除jor 醫者ge 方式唔同之外呢.
我地亦都可以問多一個問題.
有一樣ye 更加唔同.
頭先我講過呢.
醫神呢曾經係神話ge 記載裡面.
使人復活架嘛係咪.
咁而耶穌亦都曾經使拉薩路復活架嘛係咪.
咁但係唔同ge 地方係咩呀.
係頭先我地讀ge 經文.
楊豐英三章十四到十五節.
耶穌被舉起之後帶黎ge 係咩呀.
唔係普通ge 生命咁簡單.
係咩呀.
係永生呀.
永生即係咩呀.
永遠唔死ge 生命.
耶穌帶黎ge 生命唔係就係醫好你咁簡單.

$^{441}$耶穌唔係淨係復活拉薩路咁簡單.
耶穌被舉起帶黎ge 係永恆ge 生命.
係永生.
永生係上帝先有ge 特質.
醫神帶黎ge 係咩呀.
醫神就算醫好人就算使人復活.
gwo 個人最後都會點架.
都係會死ge .
佢唔會有永生ge .
所以醫神ge 醫治帶黎ge 係一個.
會再次死亡ge 生命.
而耶穌係呢度帶黎ge 係.
永遠唔會死亡ge 永生.
邊個勁D .
耶穌勁D .
邊個可以帶黎永生.
係耶穌唔係醫神.
所以呢我地咁樣睇ge 時候就明白.
係呢度約翰福音可能就係對比緊.
耶穌係勁過醫神.
只有耶穌先能夠帶黎永生.
咁耶穌係點樣帶黎永生呀.
係三章十四到十五節gwo 度咁講.
摩西係曠野點樣將石舉起.
人子也必照樣被舉起來.
使所有相信他的人得永生.
人子點樣使人得永生.
係要被舉起來.
被舉起來使人得永生.
被舉起呢個字呢其實係約翰福音裡面出現過幾次.
如果我地去睇其他ge 經文呢.
舉起特別係約翰福音十二章.
三十二到三十三節gwo 度咁樣講.
耶穌再一次講話佢要被舉起.
佢話我要是從地上被舉起.
就要吸引萬人來歸我.
然後約翰話.
耶穌者說是指著自己將要怎樣死說的.
其實被舉起呢個字呢.
係一個好含糊ge 字.

$^{481}$被舉起ge 意思呢.
就咁睇如果冇context去睇ge 話呢.
可以話高舉佢咁解.
即係可以話我地去高舉耶穌呀.
即係glorify佢高舉佢.
咁但係係約翰福音裡面ge 用法呢.
被舉起呢.
佢係特別用緊一個相關語ge 意思.
係呢度被舉起呢.
唔單止係講緊要高舉耶穌.
被舉起亦都係好detail咁解.
好字面上ge 解係點樣.
耶穌被釘死係十字架.
被高舉係十字架上面.
咁所以當耶穌話.
人子被舉起死人得永生ge 時候.
其實呢度係講緊咩呀.
係十二章gwo 度已經講得好清楚.
耶穌者話係指著自己點樣死.
耶穌係透過被舉起來.
被舉起係十字架上面.
透過佢ge 死.
黎死人得永生.
呢個係一個好特別ge 死人得永生ge 方法.
耶穌係要自己死.
然後我地呢班人先會得著生命.
係因為耶穌自己犧牲自己ge 生命.
所以我地先能夠得著永生.
所以呢.
其實成個約翰福音呢.
講好多耶穌為我地死.
耶穌為信徒死.
譬如約翰福音十五章gwo 度講.
你地要彼此相愛.
好似我愛你地一樣.
呢個就係我ge 命令.
人為朋友捨棄生命.
沒有比這更大的愛了.
耶穌被舉起死人得永生.
就係講緊耶穌為.

$^{521}$佢所愛ge 人.
我地每一個.
去捨棄佢ge 生命.
而我地亦都要.
耶穌講咩呀.
我地要學習佢一樣.
我地要彼此相愛.
係要點愛法.
要愛到好似耶穌ge 愛一樣.
要捨棄自己ge 生命.
因為冇愛.
係比呢種愛更大.
所以最後ge 一個對比.
耶穌ge 信徒.
同埋伊神ge 信徒.
一個最大ge 唔同係點樣.
伊神係一個民間信仰.
民間信仰係講緊D 咩.
民間信仰係講緊.
趨吉避凶.
民間信仰ge 信徒係點樣.
伊神ge 信徒係點樣.
伊神ge 信徒去個廟度.
去拜係為jor 咩.
為jor 得醫治.
我奉獻俾伊神.
然後伊神就要回報我.
醫治我ge 病.
所以成件事係咩黎架.
係一種交易.
一種消費.
假設如果你病.
你每一次病都浸gwo 灘老泥水.
但每一次浸完老泥水都唔好.
你仲會唔會去拜伊神.
你唔會再拜伊神.
呢D 就係民間信仰.
係講緊個神明.
俾到D benefit你.
俾到D 益處你.

$^{561}$所以我咪奉獻俾佢.
我咪去浸.
浸完D 水可以好返.
係唔需要commit 民人.
唔需要毀身.
唔需要付出.
你係要需要付出.
但係一種交易.
付出就係我俾錢你醫我.
就好似睇醫生一樣.
你唔需要commit 係個醫生度.
你唔需要為醫生付上你生命.
你唔需要為醫生去追隨一個.
醫你ge 醫生.
你唔會咁做.
呢個係民間信仰.
但我想講下係咩呢.
基督教唔係咁.
耶穌都係醫人ge 神.
但耶穌醫人呢.
講緊ge 唔係交易.
講緊ge 唔係我地俾錢你呀耶穌.
你醫我啦好唔好呀.
唔係咁樣.
講緊ge 係犧牲.
耶穌唔係因為你D 錢去醫你.
你明唔明呀.
係聖經ge 記載裡面.
從來冇一個故事係講緊.
因為有一個病人俾錢耶穌.
或者俾到任何好處耶穌.
所以耶穌都醫佢.
從來冇一個咁樣ge 醫治故事.
耶穌ge 醫治係基於佢ge 使命.
耶穌係犧牲佢自己去醫好人.
唔係因為你D 錢.
唔係因為你俾佢ge 好處.
唔係咁.
耶穌係犧牲自己.
去醫好人.

$^{601}$使人得永生.
呢個亦都係我地基督教信仰.
同民間信仰最大ge 唔同.
有時我地搞錯jor .
我地當jor 基督教同民間信仰一樣.
喂 神祝福我.
我咪愛主多D 囉係咪.
神俾錢我.
我咪奉獻多D 囉係咪.
即係我地凡事都係講求緊.
神對我好唔好.
我就對神好唔好.
呢個係民間信仰.
但係基督信仰係調返轉ge .
今日係神行先.
神犧牲佢自己.
以至到我地可以追隨佢.
今日我地要問ge 問題.
係UV裡面.
我地UVge 教會.
我地而家所信ge 耶穌.
係醫神ge 信仰.
定係真正基督教ge 信仰.
我地有ge 係消費主義ge 精神.
定係一種犧牲ge 精神.
耶穌去講.
我被舉起.
我釘係十字架.
使人得永生.
呢個係最大ge 愛.
人為朋友寫命.
沒有比這更大的愛了.
今日唔單止耶穌要釘係十字架上犧牲.
雖然今日係新年.
應該要講D 開心ge ye .
但我都好想講返.
新年我地之所以有新生命.
係因為有人為我地犧牲.
係因為曾經有一個叫耶穌ge 人.
釘係十字架上.

$^{641}$以致到今日我地有新生命.
我地有永生.
我地要問一個問題.
如果耶穌話呢個係新ge 命令.
你地要彼此相愛.
就好似我愛你地一樣.
其實背後imply緊D 咩.
imply緊就係我地同樣都要繼承.
追隨緊耶穌ge 呢種犧牲精神.
我地係為我地身邊ge 人.
去犧牲我地自己.
沒有比這更大的愛了.
呢個係耶穌ge 醫治故事.
所帶來最重要ge 意義.
我地呢度唔係一個利益ge 群體.
我地教會唔係一個民間信仰ge 團體.
我地唔係拜黃大仙.
我地唔係拜醫神.
我地唔係拜Ekipius.
我地拜ge 係耶穌.
而呢個耶穌ge 要求.
係一種犧牲ge 追隨.
今日我地UV.
係咪一個能夠願意比似毀身.
比似去付代價.
去犧牲去成就神國度ge 教會.
所以今日好想去到今日ge 總結.
今日係.
今年係蛇年.
而呢段經文.
約翰福音ge 三章十四到十五節.
好明確咁樣去將耶穌.
去形容為一條蛇.
當我地去話耶穌係蛇ge 時候.
其實最重要ge 結語係D 咩呢.
就係今日我地望到呢條.
舉起係十字架上ge 蛇ge 時候.
我地會唔會可以回想返.
正正就係呢一條蛇.
以致我地有永生.

$^{681}$係因為呢條蛇.
呢個耶穌犧牲佢ge 生命.
被高舉係十字架上面.
今日我地先能夠有新ge 生命.
有新ye 開始.
一陣間聖餐ge 時候呢.
我地可以去思想下.
我地好少去諗耶穌係蛇.
但我地可以試下思想呢段經文.
我地試下係我地腦海裡去諗下.
我地就係好似gwo D 犯罪ge 以色列人一樣.
去凝望緊gwo 條被高舉ge 銅蛇.
當我地望到呢條銅蛇ge 時候.
我地覺唔覺得係因為.
耶穌呢一條蛇.
為我地犧牲生命.
所以我地先有今日.
我地先可以進入神ge 國度裡面.
我地先有永生.
而當我地再諗落去ge 時候.
今日我地願唔願意去仿效耶穌.
去成為其他人ge 蛇呢.
唔單止係耶穌自己一個被高舉.
我地每一個都要學習耶穌.
去彼此相愛.
去成為一條被高舉.
係竿子上面ge 銅蛇.
去祝福成為別人ge 祝福.
成為別人ge 銅蛇.
讓我地一齊係一個蛇年ge 時候.
一齊去禱告.
親愛ge 主耶穌基督.
今日你去同我地講.
你呢就係gwo 條.
摩世係曠野上面舉起ge 銅蛇.
今日係因為你被舉起.
所以我地先能夠得著永恆ge 生命.
今日係新年新ge 日子.
我地去思想我地自己ge 新ge 生命.
新ge 開始ge 時候.

$^{721}$主你讓我地唔好忘記.
係因為有你ge 犧牲.
係你係十字架上面被高舉.
好似一條蛇一樣.
我地今日先可以得著呢份ge 永生.
主你讓我地唔好將呢份舊恩.
呢份ge 恩典變成好似民間信仰咁樣.
主你俾我地能夠去學習你一樣.
去擁抱呢份犧牲ge 證實.
讓我地學校你.
去成為一條能夠為別人犧牲ge 銅蛇.
去成為其他人生命ge 祝福.
祝福我地今日UVge 群體.
你讓我地依間ge 教會.
唔係一間淨係星期六星期日先會黎聚會ge 一個康樂活動ge 場所.
亦都唔係一個民間信仰ge 場所.
而係我地呢個場所能夠成為一個彼此毀身ge 群體.
我地能夠好似耶穌你所講一樣.
我地彼此相愛.
去跟隨你ge 生命靈.
就好似你愛我地每一個人一樣.
我地能夠成為呢條蛇.
去為其他人寫起生命.
主你祝福我地.
讓我地能夠係新年.
係蛇年.
活出呢種基督ge 蛇ge 精神.
讓我地成為別人ge 銅蛇.
係奉主耶穌基督ge 聖名祈求.
阿們.
\newpage



\section{}
\label{sec:cIQZdCWcAV4}
\textbf{神愛大地 (舊約篇) - 陳崇基牧師}
\newline
\newline
連結: \href{https://youtube.com/watch?v=cIQZdCWcAV4}{\texttt{https://youtube.com/watch?v=cIQZdCWcAV4}} ~~~~ 語音日期: 2025-02-11
\newline
\newline
\hyperref[sec:_BgBlG2kkQM]{\small{< < < PREV SERMON < < <}}
~
\hyperref[sec:index_chronic]{\small{[返順時目]}}
~
\hyperref[sec:index_scriptual]{\small{[返順卷目]}}
~
\hyperref[sec:UraagaffkTo]{\small{> > > NEXT SERMON > > >}}
\newline
\newline
$^{1}$鄭俊雄 早安.
我們今天繼續看短的系列.
叫做《神愛大地》.
我們講Creation Care.
我們去了解到這個主題.
很多時候在教會裡面很少講的.
因為我們最多講的是傳福音和救靈魂.
但這些東西是重要的.
我們每一個都需要去履行大使命.
但在整本聖經裡面.
怎樣教導我們.
還是救靈魂離開這個有苦難的世界.
脫苦海去救呢.
整本聖經應該怎樣去思考這個問題.
所以我們想在這個系列裡面.
去做一些所謂的theological reflection.
一些神學上的反省.
看看聖經裡面怎樣講.
上一次我們講到.
可以有兩個不同的時間表.
或者兩個不同的概念.
第一個概念就是我們傳統教會裡面.
很多時候都接觸到的概念.
就是我們現在這個世界既然被罪破壞了.
上帝來救我們.
藉著耶穌基督救我們的時候.
最終這個世界現在的天地.
一定會是total destruction.
一定會是全然的毀壞.
好像按一個按鈕.
它self-destruct.
然後全部東西都沒有了.
然後新天新地就會來到.
耶穌再回來的時候就會來到.
所以你看到那個核爆那裡.
所以如果你用一個這樣的概念.
去了解我們現在的世界.
和將來新天新地之間的關係.
你就會發現有一個完全絕對的中斷性.
現在這個大自然.

$^{41}$現在我們所做的一切.
無論是文化上的東西.
藝術,音樂等等.
是沒有什麼永恆意思的.
因為到最後的時候都是核爆.
大家爆炸了所有東西.
Game over.
然後新天新地就會來到.
那你說有什麼可以做.
就是我們要脫離這個世界.
這是第一個概念.
但我想問這個概念到底是不是在聖經裡面是符合的.
還是我們看漏了一些東西不是那麼完備.
所以有第二個概念.
我們想去研究或者想去看清楚.
就是這個雖然是被罪破壞的天地.
現在的天地.
和將來新天新地的關係會是怎樣的呢.
會不會只是一個完全的中斷.
完全是Game over.
還是有那個連續性的continuity在裡面.
就是說由現在的世界到下一個將來基督回來的時候.
那個新天新地的關係不是斷軸的.
而是一個更新的關係.
一個完全更新的關係.
我知道講到這裡的時候.
大家覺得很抽象.
你在講什麼.
其實不是那麼抽象.
我們住芝加哥的都知道.
我們四季分明.
冬天很冷.
然後過了冬天就是春天來的時候.
冬天和春天之間的關係.
就是這個關係.
就是一個renewal.
你冬天裡面做的事.
是不是到春天來之前.
然後核爆爆了它.
然後我們不要了.

$^{81}$不是的.
你看那些樹木.
雖然在冬天好像沒什麼活力.
但是到春天的時候.
它們又再生出來.
這個就是所謂total renewal的意思.
這幅圖畫或者可以幫助到我們.
聖經裡面確是這樣講.
現今這個世界不是最後核爆Game over.
所以現在所做的事.
什麼都沒意思.
我們只是去傳福音然後離開這個世界.
救我們靈魂.
我們身體沒意思.
我們物質的世界沒意思.
我們所做的事.
藝術,音樂,文化等等.
都是沒意思的.
這個概念不是一個聖經裡面最完備的概念.
一個最完備的概念就是我們明白.
今天所做的事原來都是有價值的.
在上帝的國度裡面是有價值的.
是有連續性的.
將來會被renewal的.
如果這樣看的時候.
我們整個人的生命.
我們就有些不同的看法.
所以這次我們用一個主題式的方法.
我們看看舊約聖經.
是怎樣講上帝和大地之間的關係.
我們和創造界之間的關係.
第一點我想講的就是上帝給人.
有一個共同的角色.
你和我都有共同的角色.
無論你信耶穌也好.
你不信耶穌也好.
我們都有一個共同的所謂天職.
上帝給我們這個角色.
就是做一個管家.
創世記《劈頭》開始就這樣講.

$^{121}$上帝說我要按照我們的形象.
和按照我們的樣式做人.
這裡講做人類.
人類我們就是.
包括我們在內.
使他們管理海裡面的魚.
天空的飛鳥.
和地上的牲畜和全地.
地上爬行的一切動物.
上帝就按照他們的形象創造人類.
按照上帝的形象創造他們.
創造男 創造女.
上帝就賜福給他們.
說你們要生養更多的遍滿全地.
治理它 管好它.
管理大地.
這個就是我們每一個人.
有的共同的角色.
這個就是我們所謂的文化.
Culture.
我們是做一個管家.
而這裡講得很小心.
這裡不是說上帝創造人.
叫他管大地的時候.
就是扔給他.
然後任意的被他剝削.
任意的被他拿取.
這裡講的是什麼.
是管家 是管理.
如果再看創世記第二章的話.
那裡講到神將人類放在伊甸園裡面.
是叫他耕作.
其實意思是服侍土地.
服侍看守大地.
什麼叫管家.
管家就是說我們沒有最終的擁有權.
最終擁有權是誰 是上帝.
既然是管家的話.
我們最後的話.
我們就需要向上帝問責.

$^{161}$是有accountability.
等於你把錢放在銀行.
銀行就好像你管家management那樣.
錢是應該屬於你的.
然後他好好管.
好好保守住.
你是那個擁有者.
最終是向你負責.
我們和上帝之間的關係也是這樣.
所以第一樣.
在創世記第一章.
裡面講到人類.
就是有一個管家的責任.
然後我們看一些舊約的經文.
特別是五經裡面.
怎樣講神關愛這個被造界.
這裡選了一段.
新面記第二十章.
這裡說以色列人.
他們有時候要打仗.
有一個戰爭的法則.
我們看看裡面細節講什麼.
他告訴以色列人.
你軍隊去圍攻一個城市的時候.
你需要攻打很多日子.
才能夠奪取.
你就不可以用斧頭來砍壞樹木.
你們可以吃樹上的果子.
他不可以將樹砍下來.
田間的樹木是人.
只是樹木沒有得罪你.
你為什麼要和他作對呢.
你要留著它.
然後二十節.
只有那些你知道不能生產食物的樹木.
就可以毀壞.
拿來燒東西也好.
拿來做一些工具也好.
直到你將城攻取.
這裡我們不是講戰爭倫理問題.

$^{201}$這裡有一個特別的地方.
就是以色列人打仗的時候.
他可以和他說.
你有一個法則.
這個法則就是一個所謂戰爭的法則.
如果你對比其他近東的國家.
和一些特別軍閥的國家.
好像亞述那樣.
他們打仗就不是這樣的.
你看到右下角的浮雕.
講到亞述軍隊去攻城的時候.
圍城的時候.
他們做什麼呢.
所有的樹都砍.
那些樹都砍掉.
為什麼.
因為軍閥國家他們要攻打的時候.
有什麼戰略.
是殺他 拿來做一個戰爭.
因為他打到你殘.
就算你撿回一條小命.
你也是一身殘.
就是這個意思.
所以果樹都砍掉.
你知道那些橄欖樹.
要四五年才開始.
會結成熟的果子.
所以砍完你的樹的時候.
就算我走了.
你也餓死.
所以亞述就是這樣去打仗.
以色列人不是這樣.
這裡看到的律法.
是表現了上帝.
怎樣關愛受造的所有東西.
包括樹木.
包括生態.
你打橫打.
你不要搞亂生態.
你打橫打.

$^{241}$不要摧殘大自然.
我不知道大家有沒有聽過一段話.
他說地球已經有46億年的歷史.
如果將46億年轉成一個人46歲.
有誰46歲.
你46歲的話.
地球就好像你這麼大.
將它轉成的話.
如果轉成地球有46歲.
so far有46歲.
人類出現只不過是整個地球歷史的四個小時.
so far我們只生存了四個小時.
而在四個小時裡面.
前一分鐘就是工業革命的時候.
工業革命就開始斬樹.
現在只斬了一分鐘.
我們就斬了全地球一半的森林.
是我們摧毀的.
是人類摧毀的.
我們愛不愛大自然.
我們愛大自然好像我們愛黑膠燒賣.
有人說黑膠燒賣你愛不愛.
愛.
其實就是我們要消耗它.
我們要拿走它.
我們要吃下肚.
我們要消耗它.
這就是我們現代社會面對的一個很大的問題.
一個環境生態的問題.
聖經有沒有一些這樣的地方.
幫助我們去對話.
給我們一些的insight.
絕對有.
大家知道我都會鬥木.
我家裡都會鬥木.
UV裡面有些東西都是我鬥出來的.
在美國鬥木.
有一樣東西很重要的.
就是2x4.
你知道什麼叫2x4嗎.

$^{281}$牆後面那些drywall後面.
那些骨幹.
一條條的那些叫2x4.
在美國你去home depot.
去一些home improvement store.
你一定會買那些2x4.
就是兩吋乘四吋的木方.
你知道最近出的.
近年出的木方.
就是上面那一條.
老的木.
芝加哥有很多舊的樓房.
有些百年的樓宇.
就是用一些老的木.
1918年出的木.
有什麼分別.
你看那些樹環.
是密很多的.
因為等它慢慢慢慢出來.
它就很堅實.
所以那些百年老屋是很硬的.
你要它燒的話都很難燒得爛.
因為越堅實的木就越難燒.
現在弄的那些.
新樓你很喜歡新樓.
全部都是用這些2x4弄的.
如果你跟我一起鬥木的話.
我們鋸那些2x4.
好像鋸牛油一樣.
一鋸就搞定了.
所以你看到那些空手道.
是不是很厲害.
其實一點都不厲害.
因為現在的木方很差.
你去那些Akira那些地方.
就是這些木.
比這些木還差.
這個就是我們現在面對的問題.
因為木材不夠.
所以一味去到生產的現代木材.

$^{321}$就是那些很快生產出來的樹木.
所以有過分的砍伐的問題.
影響那個生態.
在《新面記》在《五經》裡面.
時常告訴我上帝是關愛被造之界.
22章第6節.
他說你如果見到路上有鳥窩.
無論在樹上還是地上.
裡頭有小鳥或者有蛋.
母鳥伏在小鳥或者蛋上面.
你是華人就馬上拿了它.
哇 現在蛋這麼貴.
這個律法說什麼.
你不可連母鳥帶小鳥一起拿去.
總要放母鳥.
你可以管理你可以用大地.
不過你需要確保它能夠持久.
對不對.
持久性.
那你就可以享福了.
在地上你的日子.
整個民族整個國民.
才能夠延續下去.
這裡說得很清楚.
限制人在大地的資源運用.
給到你的東西不代表每樣東西都可以用.
或者用到盡.
或者拿到盡.
你有沒有拿過橙汁.
以前那些拿橙汁.
你怎樣拿橙汁.
你拿到盡最後一滴你都拿了.
為什麼.
拿完它都扔了.
為什麼不拿最後一滴.
這個就是我們很多時候對大自然的心態.
大地律法裡面一次又一次告訴我們.
整個大自然在上帝的眼中是他愛護的.
要保持生態的平衡.
要延續要sustain.

$^{361}$然後要去到限制人.
不要隨便的拿乾拿盡大地的資源.
二十五章第四節.
牛在喘谷的時候不要露著他的嘴.
這段聖經我曾經說過.
因為在新約裡面是應用這段聖經.
是應用在傳道和牧者身上.
保羅說牧者就好像牛一樣.
我就好像牛一樣.
我就在喘谷幫大家喘谷.
如果大家綁著我的嘴.
喘不喘很久.
喘到喘不住氣我就死了.
但如果你覺得.
既然世界玩完game over都扔了.
牧師這樣做到牠喘氣.
牠沒東西吃.
就扔了牠不可以.
請個新的來.
不可以這樣.
因為我們要慷慨地明白.
牛是幫助牠做事.
以至牠幫牠做事的時候.
我們都讓牠有衣食.
讓牠能夠有sustainability.
好好照顧牠.
我們教會牧者傳道人.
甚至我們中間有些神學生.
其實都是要這樣.
值不值得?如果你和我吃飯請我吃飯.
誰會請我吃飯?.
如果我和你吃飯.
你班人真是.
好奇有些人性在.
這是重點.
我們不只是消耗他們的資源.
我們是要確保他們是sustain.
你們還沒記住第十九章.
你們自己在地上收割的莊稼.
不要割盡田割.

$^{401}$要留下莊稼所掉下來的.
你去收割大麥小麥的時候.
不要全部都收.
那些葡萄樹不要全部都收.
掉下來的不要再撿.
不要擠乾擠盡.
因為有需要的人窮人寡居的.
他能夠在你後面.
去拿到這些飲食.
我是耶和華.
你的上帝.
最終擁有權是誰?是上帝.
是上帝自己擁有的.
所以你不可以去抓人.
這個是愛護生態和社會公義去結合.
愛護這個社會愛護身邊的人.
還有去確保大自然.
那個生態能夠繼續運作下去.
到最後第三點就是.
上帝是大地的主宰.
這是一段很長的經文.
我用了很多時間去濃縮出來.
這個是講到安息年和欺年.
利未記二十五章.
在西乃山吩咐摩西港.
你要吩咐以色列人.
對他們講你們六年要耕作田地.
六年你要收整葡萄樹.
收藏它的出產.
但到第七年.
地要完全安息的安息年.
就是向耶護華面前安息.
這個很有趣.
形容地就好像人一樣.
人就需要安息.
但他們講土地是需要安息.
地在安息年.
即是第七年.
所長的是給你和你的僕僚.
你的故宮.

$^{441}$甚至在你中間寄居的.
那些寄居的.
移民也好或者難民也好.
總之在你中間.
你找到他們的話.
是給他們去食的.
所有的出產.
也是給你們的生畜和地上的走獸.
你要計算七個的安息年.
所以六年裡面第七年就是安息年.
但七個安息年之後.
七七四十九年.
到第五十年.
第五十年就是特別一點的.
第五十年就叫做欺年.
欺年的時候發生什麼事呢.
就是所有的人民都宣告自由.
這個叫做欺年.
國人的產業要歸還給國人.
國人要回家.
在說一些窮人.
因為天災的緣故等等不同的原因.
他的地要被賣出去.
但不可以永賣.
為什麼.
因為到欺年的時候.
大家就可以回家了.
好像整個社會按一個reset button.
好像你的電腦按一個reset button.
重新再來.
第五十年要作你們的欺年.
你們不可以耕種.
你們要吃地中自然生長的農作物.
就是不可以繼續去.
拿乾拿盡那些土地的資源.
給他自然的生長.
給他休假.
地不可以賣斷.
因為什麼.
因為地不可以賣斷.

$^{481}$因為地是我的.
你們在我面前是客裡是寄居.
這個奇怪.
跟以色列文說.
以色列文明明是公民.
是有些寄居者在他中間.
但在上帝的眼中.
他們始終都是寄居.
是客裡.
是on a journey.
他們不是真的永遠可以擁有那樣東西.
誰擁有.
是上帝.
所以說你們欺年的時候.
國人的產業就歸於國人.
國人就要歸回自己的家.
第五十年作你們的欺年.
你們不可以耕作.
要休假.
地不可以斷.
因為是屬於我的.
這個經文很明顯.
第一個關心的就是土地.
The land.
土地有安息年有休養.
這段經文很多時候用在牧師.
教牧童工.
需要有安息年.
Sabbatical year.
所以有些弟兄姊妹跟我說.
你做了這麼久.
做了七年.
你不如安息年.
最主要.
我不需要安息年.
但最主要想說的是.
確保經濟和環境.
天然的資源能夠永續.
人不應該.
也不能夠.

$^{521}$是渣幹渣盡.
大家知道我讀神學之前.
我們讀什麼.
Environmental toxicology.
是不是很厲害.
說這個名字出來.
Environmental toxicology.
你知道是什麼嗎.
是環保毒理學.
不是去毒人.
我們是去clean up.
去做一些水或泥土的測試.
拿一些sample.
所以我們會去field拿東西.
我在實驗室做過一段時間.
我讀Environmental toxicology的時候.
有一本所謂的天書.
是一定要看的.
當時我讀的時候是90年代.
那時候這個科目都不是很成熟.
現在成熟很多.
那時候還沒有textbook.
沒有一個很大的textbook.
給你做reference.
但有一本書.
每個人都要看的.
就是Rachel Carson寫的Silent Spring.
寂靜的春天.
我們每個人讀環保毒的科目.
都是需要去讀這本書的.
這本書為什麼這麼出名.
這本書被譽為20世紀最具影響力的環保著作之一.
這個Rachel Carson本身是一個marine biologist.
她寫這本書是寫什麼.
她想揭露當時60年代的時候.
美國不斷用一種殺蟲水叫DDT.
有沒有聽過.
DDT.
就是所有的田地全部都灑這些DDT.
就是殺蟲.

$^{561}$殺蟲那些香蕉蘋果就很漂亮.
或者你今天去市場買的那些是不是很漂亮.
你就知道了.
是很毒的.
因為我們都做過這些DDT的實驗.
對生態有很大的破壞.
譬如會導致那些雞吃了.
在食物鏈裡面染了這些DDT.
它生的蛋殼是會特別薄一點.
美國的雞蛋是不是特別薄一點.
在廚房做的那些一剝掉整個都爛掉.
其中一個原因.
所以她寫這本書之後就掀起了一些關注.
以致1972年的時候.
美國就正式禁止DDT.
去推動一些生態的保育等等.
我們不是一定要贊成左派裡面所有的議題.
但我們需要回到上帝的話語裡面怎麼說.
為什麼他們要用DDT.
為什麼要破壞環境.
因為他們要做生意.
做生意怎麼管得這麼多.
我們最重要的是要拿到最好的利潤.
地下的石油為什麼不轉.
什麼都要轉出來.
為什麼.
因為他們就是黃金.
有錢不賺了.
以後你就吃濕米.
這個問題是價值觀的問題.
是我們的價值系統的問題.
以前我在實驗室做.
其中一個大客.
我們經常幫他們驗東西.
叫做McLaren Air Force Base.
在Sacramento附近.
為什麼要經常和美軍的基地驗水驗泥.
因為那些重金屬會導致食水污染.
泥土會污染.
所以我們經常會做這件事.

$^{601}$下一次告訴你一個很有趣的故事.
怎樣都好都要告訴我們.
我們需要明白生態是需要保護的.
怎樣實踐.
起碼有一個這樣的awareness已經足夠.
起碼有這樣的awareness去開始.
原來這個世界不是玩完罷了.
然後燒了它.
毀滅.
上帝有最終的擁有權.
我們在地上做一個管家.
到最後的時候我們播一首回應詩歌.
也是Bella寫的.
但我們重新整過這些影像.
叫做龜圖.
直至聽這首歌的時候.
我們思想一下到底我們和大地之間的關係.
上帝和大地之間的關係是怎樣的.
龜圖.
一寸一界一語.
一生年日經處.
堂上你價值標籤.
斌白倦容.
為救踏著這寸土可掌控.
讚卻未夠盡心可放鬆.
一花一土一間.
一刻退去消散.
如燕棲處碧兒灘.
冰山崩塌.
南瓊活象角前已掩飾他.
日好盡了夜在吹奏.
世庭海枯下地也喘息.
世界的宿命到了險境.
讓地歸於安寧.
萬物靜養生息.
歸後鳥語凜看.
沙遊水情.
地屬主所賜.
立若鐵證.
留你可安定.

$^{641}$坐看風景.
嘲諳的不公平.
雷頭的總可放開爭拼.
眾聖靈.
期盼重拾風聲.
Zither Harp.
一花一土一間.
一刻退去消散.
如燕棲處碧兒灘.
冰山崩塌.
南瓊活象角前已掩飾他.
日好盡了夜在吹奏.
世庭海枯下地也喘息.
世界的宿命到了險境.
讓地歸於安寧.
萬物靜養生息.
歸後鳥語凜看.
沙遊水情.
地屬主所賜.
立若鐵證.
留你可安定.
坐看風景.
嘲諳的不公平.
雷頭的總可放開爭拼.
眾聖靈.
期盼重拾風聲.
Zither Harp.
你我只不過是個客裡.
也有天花灰盡.
黃地土碎.
如流星一閃直去.
別讚吹多於瑣碎.
讓孤寡可寄居.
你我生於這裡.
Zither Harp.
\newpage



\section{馬太福音 25:14-30}
\label{sec:UraagaffkTo}
\textbf{神國藍圖 (馬太福音 25\_14-30) - 余淇煥傳道}
\newline
\newline
連結: \href{https://youtube.com/watch?v=UraagaffkTo}{\texttt{https://youtube.com/watch?v=UraagaffkTo}} ~~~~ 語音日期: 2025-02-18
\newline
\newline
\hyperref[sec:cIQZdCWcAV4]{\small{< < < PREV SERMON < < <}}
~
\hyperref[sec:index_chronic]{\small{[返順時目]}}
~
\hyperref[sec:index_scriptual]{\small{[返順卷目]}}
~
\hyperref[sec:iar_apDpNcQ]{\small{> > > NEXT SERMON > > >}}
\newline
\newline
馬太福音 25:14-30
\newline
\begin{longtable}{cl}
\hline
\hline
章節 & 經文 (和合本修訂版)\\
\hline
25:14 & \begin{tabularx}{0.7\textwidth}{X} 「天國又好比一個人要出外遠行，就叫了僕人來，把他的家業交給他們。 \end{tabularx} \\ \\ \relax
25:15 & \begin{tabularx}{0.7\textwidth}{X} 他按著各人的才幹，給他們銀子：一個給了五千，一個給了二千，一個給了一千，就出外遠行去了。 \end{tabularx} \\ \\ \relax
25:16 & \begin{tabularx}{0.7\textwidth}{X} 那領五千的立刻拿去做買賣，另外賺了五千。 \end{tabularx} \\ \\ \relax
25:17 & \begin{tabularx}{0.7\textwidth}{X} 那領二千的也照樣另賺了二千。 \end{tabularx} \\ \\ \relax
25:18 & \begin{tabularx}{0.7\textwidth}{X} 但那領一千的去掘開地，把主人的銀子埋藏了。 \end{tabularx} \\ \\ \relax
25:19 & \begin{tabularx}{0.7\textwidth}{X} 過了許久，那些僕人的主人來了，和他們算賬。 \end{tabularx} \\ \\ \relax
25:20 & \begin{tabularx}{0.7\textwidth}{X} 那領五千的又帶著另外的五千來，說：『主啊，你交給我五千。請看，我又賺了五千。』 \end{tabularx} \\ \\ \relax
25:21 & \begin{tabularx}{0.7\textwidth}{X} 主人說：『好，你這又善良又忠心的僕人，你在少許的事上忠心，我要派你管理許多的事，進來享受你主人的快樂吧！』 \end{tabularx} \\ \\ \relax
25:22 & \begin{tabularx}{0.7\textwidth}{X} 那領二千的也進前來，說：『主啊，你交給我二千。請看，我又賺了二千。』 \end{tabularx} \\ \\ \relax
25:23 & \begin{tabularx}{0.7\textwidth}{X} 主人說：『好，你這又善良又忠心的僕人，你在少許的事上忠心，我要派你管理許多的事，進來享受你主人的快樂吧！』 \end{tabularx} \\ \\ \relax
25:24 & \begin{tabularx}{0.7\textwidth}{X} 那領一千的也進前來，說：『主啊，我知道你，你是個嚴厲的人：沒有種的地方也要收割，沒有播的地方也要收穫， \end{tabularx} \\ \\ \relax
25:25 & \begin{tabularx}{0.7\textwidth}{X} 我就害怕，去把你的一千銀子埋藏在地裡。請看，你的銀子在這裡。』 \end{tabularx} \\ \\ \relax
25:26 & \begin{tabularx}{0.7\textwidth}{X} 他的主人回答他說：『你這又惡又懶的僕人，你既知道我沒有種的地方也要收割，沒有播的地方也要收穫， \end{tabularx} \\ \\ \relax
25:27 & \begin{tabularx}{0.7\textwidth}{X} 就該把我的銀子放給兌換銀錢的人，到我來的時候可以連本帶利收回。 \end{tabularx} \\ \\ \relax
25:28 & \begin{tabularx}{0.7\textwidth}{X} 把他這一千奪過來，給那有一萬的。 \end{tabularx} \\ \\ \relax
25:29 & \begin{tabularx}{0.7\textwidth}{X} 因為凡有的，還要加給他，叫他有餘；沒有的，連他所有的也要奪過來。 \end{tabularx} \\ \\ \relax
25:30 & \begin{tabularx}{0.7\textwidth}{X} 把這無用的僕人丟在外面黑暗裡，在那裡他要哀哭切齒了。』」 \end{tabularx} \\ \\
[1ex]
\hline
\hline
\end{longtable}
$^{1}$Thank you 姊妹早安.
早安.
我每一次都很珍惜講道的時間.
因為當我講道給你們聽之前.
其實我要講給自己聽的.
因為我要練習.
所以我自己無恙自己.
今天我想講的道就是天國南途.
我在香港的時候.
就住過在香港仔.
每逢坐巴士出銅鑼灣的時候.
都一定會經過跑馬地的天主教墳場.
場上經常都有這句金句.
不是金句.
就是對聯.
站在那裡的時候.
我就經常看著這句對聯.
上面寫著.
「金席吾軀歸故土,他朝軍體也相同」.
意思就是今天我的軀體已經入土為安.
有朝一日你的身體都同樣會歸於塵土.
坐巴士的時候.
每一次看到這句對聯的時候.
心裡都有點寒.
就會想.
有一天其實我也會這樣.
就是提醒我去想.
我會為什麼著迷呢.
已經信了耶穌的我.
有一天都會進天國的我.
其實今天對於這個天國.
對於我有什麼意義呢.
如果大家記得上一個星期.
Mark Sir說的.
說的原初的天地.
現在被最破壞的天地.
和將來的新天新地.
是有一種連貫性的.
如果我們現在的地上.
和將來的新天新地有一個連貫性的時候.

$^{41}$就是說天國不只是等到我們將來才去的地方.
也不是關乎我們死了之後的事.
而是我們今天的生命.
會被將來的天國所影響.
我們也可以將天國帶到今天.
我們所活著的地上.
所以今天我想藉由耶穌的比喻.
說一下我們怎樣用天國的視野.
去看我們今生的生命.
在《馬太福音》25章14至30節.
我想用口語化去說給大家聽.
大家當聽故事一樣聽.
這個比喻是這樣說的.
耶穌說天國就好像一個主人.
他要出遠門.
就叫了他的僕人來.
將自己的財產派給他們去打理.
主人就按著每一個僕人的能力.
給了不同數量的銀子給他們.
第一個僕人主人給了五千.
第二個僕人主人給了二千.
第三個僕人主人給了一千.
然後這個主人就出遠門了.
第一個僕人拿著五千馬上去做生意.
多賺了五千回來.
第二個僕人也是一樣.
將二千多賺了二千.
但是第三個僕人就不同了.
他怕自己吃虧.
於是就挖了個洞.
將主人給他的一千埋在地上.
什麼都沒有做.
到了一段時間.
主人回來了.
叫了他們三個僕人回來要交代.
要交數了.
第一個僕人就說.
主人,你給我的五千.
我現在賺回五千了.
主人很開心.

$^{81}$就說好啊.
你這個又忠心又本事的僕人.
你做得很好.
我會給你管理更多的東西.
你進來.
跟我一起開心地去享受.
第二個僕人就說.
主人,你給我的二千.
我現在賺多了二千.
主人一樣這樣稱讚他.
他說好啊.
你是一個又忠心又本事的僕人.
我給你管理更多的東西.
你進來.
你跟我一起開心地去享受.
第三個僕人就說.
主人,我知道你是一個很嚴格的人.
你會收割那些沒有種的地方.
你會拿走那些沒有放進去的東西.
所以我怕得不敢動.
將你的一千埋在地上.
現在我拿回給你.
主人聽完之後他很生氣.
你這個又惡又懶的僕人.
你既然知道我會收割那些沒有種的地方.
為什麼你不拿錢去銀行.
起碼你收一下利息也好.
然後就拿走那一千.
給了那一萬的僕人.
因為凡有的他會再加給他.
令他更加富足.
沒有的就算剩下的那一點點.
也會被人拿走.
至於這個沒有用的僕人.
被趕出去之後.
他在黑暗裡哭和咬牙切齒.
如果要將這三個僕人分成兩組.
你會怎樣分.
很自然我們會將兩個好的僕人分為一組.
壞的僕人在第二組.

$^{121}$但如果這個比喻只是想強調.
好的僕人和壞的僕人.
其實一好一壞不就行了.
為什麼我們要有兩個好的又要有兩個壞.
我們要一起看看作者.
究竟透過這兩個好的僕人.
想講什麼重點.
這兩個好的僕人有什麼不同的地方.
又有什麼相同的地方.
我們會發現這兩個好的僕人不同的地方.
是什麼.
就是主人給他們的東西.
和他們賺到的金額是不同的.
一個是五千.
另一個是二千.
但他們兩個都同樣賺到.
他們所得的金額的一倍的錢.
而主人對他們的讚賞和賞賜是一樣的.
這個比喻一開始已經跟我們說.
主人是按照僕人的才能去分錢的.
所以兩個好的僕人的重點.
不是哪一個賺得多.
不是哪一個賺得少.
而是他們同樣按照自己所擁有的才能.
按照他們的知識和經驗.
立刻善用主人給他們的金錢.
無論是賺五千還是賺二千.
這兩個僕人最後都得到主人相同的讚賞和獎勵.
如果主人是上帝.
僕人是我們.
金錢代表著神在地上給我們的一切.
包括性格特質.
包括你的才幹.
你的財產.
你的人際關係.
你的人生經驗.
你的知識等等.
這個比喻是說.
神給了我們在地上作為管家的角色.
我們如何善用上帝給我們這一堆東西呢.

$^{161}$在上帝回來之前.
我們如何盡我們這個管家的本分呢.
我想說.
我們每一個在座每一個人.
在上帝天國的藍圖裡面.
我們都有一個很獨特的位置.
上帝都給了我們每一個人有我們那一份.
不過有時候我們會很不其然地.
很容易地跟其他人比較.
為甚麼他會有錢一點.
為甚麼他的人生會順利一點.
為甚麼他的人際關係會好一點等等.
但當我們的焦點看著隔離一個擁有的東西的時候.
我們會否錯過了上帝給我們面前的那一份呢.
有二千銀子的僕人.
如果他掛著眼紅五千的那個.
他很可能停留在那些苦澀,妒忌,比較的裡面.
他很難再有心力去努力.
去想自己面前的那二千可以如何增長.
而事實上在我們人生裡面.
一定會遇到給我們富有.
給我們有能力.
給我們人緣好.
給我們吸引的人.
又如果掛著五千銀子的那個僕人.
他很自滿地覺得.
主人給我那麼信任我.
他特別信任我.
而他鄙視掛著二千或一千的那個.
甚至他擔心別人會超越了自己.
這樣的情況.
他亦都可以很不專心.
很難去專心上帝交託給他要做的事.
所以今天我們不要忘記.
無論今天我們有的一切.
最終都是屬於上帝.
我們最終都要給上帝.
有一天當我們去面見神的時候.
上帝不會問你隔壁屋的陳生.
這一生過得怎樣.

$^{201}$祂會問你這一生過得怎樣.
如果大家記得.
之前His Foundation.
即是天倫基金創辦人劉醫生.
他曾經來過UV港島.
有沒有人聽過劉醫生港島.
劉醫生是一個骨科醫生.
他創辦了一個幫助落後國家做手術的基金.
去幫助貧窮國家的人.
Joyce傳道之前參加過的伊拉克醫療短宣.
就是這個基金.
我自己去過兩次.
醫療短宣裡面最需要的人才.
第一是醫護.
第二是行政能力很強.
可以幫忙做物流的人才.
而我在團隊裡面這方面的貢獻都是零.
但為什麼我會去呢.
我又不是醫護.
又沒有醫護的知識.
我的Admin又不行.
但我都決定要去.
我就問劉醫生可不可以讓我去.
我有兩個原因.
第一個原因是.
我對比在香港生活.
其實在芝加哥的生活可以很安逸.
如果我不選擇離開我的安書區.
我這個階段的人生.
我會覺得自己好像一條鹹魚.
第二個原因是.
當我看著劉醫生分享和播片的時候.
我看到的不單止是.
這班難民身體上的痛楚.
或者他們的病苦.
或者貧窮.
我更加看到他們心靈裡面的痛苦和掙扎.
其實自2014年開始.
這班Yazidis的群體.
他們被ISIS攻陷他們的家園.

$^{241}$他們被虐殺.
被虐待.
被屠殺.
他們的家園完全是零碎不堪.
他們要逃難.
甚至他們被滅族.
他們面對著失去親人等等種種的創傷.
所以雖然我不是一個醫護.
但是在神的裡面.
我看到我的生命.
無論神放在我生命裡面.
我讀心理輔導的學歷.
我自己的成長經歷.
或者我的性格和恩賜.
都彷彿在呼喚我.
要離開我自己的安書區.
去到一個我不熟悉的地方.
甚至一個溝通不到的地方.
去為神做一點事情.
就是做心靈醫治.
當時我帶了幾樣東西去.
就是顏色筆,紙和圖畫卡.
我帶著這幾樣法寶.
和當地的翻譯員.
去接近那些病人和家人.
跟他們聊天.
表面上他們都很有笑容.
但是當我深入了解他們的時候.
透過他們畫畫.
透過他們選一些卡.
去分享他們內心世界的時候.
我就有機會聽到很多.
他們內心世界的遍體鱗傷的故事.
由於時間有限.
裡面其實很豐富.
有機會再跟大家深入地分享.
我想說的是.
其實參與短宣的人.
每個人都有份正職.
為什麼還會去短宣呢?.

$^{281}$Joyce傳道耀牧會.
而團隊裡面也有不少.
來自香港的醫護.
平時上班都已經夠忙了.
下班當然要休息一下.
他們竟然請自己的年假.
是年假.
香港的年假是很珍貴的.
再用自己的薪金.
不是去旅行.
而是山長水遠去伊拉克做短宣.
是什麼呢?.
令劉醫生作為一個有名望的.
美國的骨科醫生.
到了退休的年紀.
但可以頑皮.
但他願意一年四次去伊拉克.
坐二十多個小時飛機.
然後再坐四個小時車.
兩個星期內做二十台手術.
我想說的是.
很多病人的腳其實是很重的.
你要扛起他們.
有些動作經常要抬起他們的腳.
是很重的.
還有做骨科手術的工具.
鋸,鎚其實都是很重的.
有時連劉醫生自己也會有職業病.
自己要備做手術之後.
休息幾個月.
又要再去幫人做手術.
而我自己在短宣裡.
因為很想認識多點病人.
和難民的生活環境.
所以我多留了一個星期.
去難民營做探訪.
我想說這些所有的事情.
不是每個人都需要一個專業人士才能做到.
也不是每個人都要去到山長水遠.
才能進到天國的那一份.

$^{321}$好像第一個和第二個僕人.
雖然他們的金額有不同.
但他們也有共同的特質.
就是盡力.
今天你和我有沒有看到.
神給你很獨特的那一份.
是真的只有你才有的呢?.
然後你能夠為著神的國度去盡你那一份呢?.
我們又看看第二組的僕人.
這個比喻用了很多篇幅去描寫.
這個又懶又兇的僕人和主人之間的對話.
這個僕人有什麼問題?.
第一個問題當然是懶.
懶得連把錢放進銀行的力量都沒有.
就是隨手把錢近近地埋在地上.
不過更嚴重的問題.
是這個僕人沒有正視自己的懶.
他把責任愛化了.
很有趣.
這個僕人在主人還沒有開口之前.
他就開口去指責這個主人.
他就說.
我知道你是一個很嚴格的人.
因為你這麼嚴苛,這麼嚴格.
所以我真的怕得動不了.
我真的怕得不敢動.
把你的一千元埋在地上.
現在我還給你.
當我看到這裡的時候.
我會覺得這個僕人其實可以是一個基督徒成長的寫照.
為什麼這樣說呢?.
就是很多人信了耶穌.
很多人會去教會.
但是生命是沒有成長的.
你以為一個人信了耶穌.
就會自動成長嗎?.
你以為一個基督徒去教會聽得多.
就會自動成長嗎?.
是不會的.
成長是什麼意思呢?.

$^{361}$成長是你的生命願意讓上帝去介入.
讓上帝去動你一些不想被人動的東西.
例如你願意為自己一些不健康的性格.
又或者生命當中一些重覆又重覆的問題.
很忍心,很決心地去正視它.
成長是不再給自己藉口去逃避問題.
不再純粹將問題去外化.
怪罪在一些外在的原因身上.
而是你能夠誠實地去面對自己.
你願意去反省.
找回屬於你生命.
你覺得神需要你去改變的那個部分.
我想說是很難的.
因為沒有人會喜歡被人說.
沒有人會喜歡被人說覺得自己有問題.
要一個人發現自己有問題已經夠難.
還要去肯去正視自己的問題就更難.
然後還要承認自己有問題.
而選擇去改變就十級難.
這麼難的事情誰會做呢?.
你會不會做呢?.
我想說的是.
所以通常我們會怎樣處理呢?.
怎樣去面對呢?.
就是用最簡單的方法.
就是將垃圾掃到沙發底.
就沒有人看到.
自己也看不到.
那就OK了.
有很多說話有時候會對自己說.
你不要理我.
我是這樣的.
又或者對方都不改.
為什麼要我先改呢?.
又或者我這麼生氣.
都是因為你們惹我.
等等等等這堆東西.
而我曾經也說過.
如果大家有聽我上一次說到.
都知道我成長的背景很不容易.

$^{401}$也很不健康.
我不懂得什麼叫情緒界線.
我成長就是一直為我的家人.
為我父母的情緒負責任.
你不開心就是因為我不夠好.
你會想像到在這個環境長期成長.
我的性格是可以影響得多不健康.
和多扭曲.
甚至那種沉重.
那種被怪罪.
很多內疚感.
那種成長令我覺得.
當日我沒有死也是一個神蹟.
還要今天這麼正常地.
站在大家面前去說.
我想說今天能夠做一個正常人.
是因為我背後走過了一段很長的成長路.
我相信了耶穌之後.
剛才說了成長既然有這麼多艱難.
為什麼要成長呢?.
有什麼動力要我們成長呢?.
我自己每個人都需要一種動力.
去推動你成長.
我自己也有一個動力.
那個動力是什麼呢?.
就是那種好奇心.
就是我相信了主之後.
我裡面有一種好奇心.
會覺得神你創造我這個人.
原本的形象會是怎樣的呢?.
因為你的創造本身一定是美好的.
為什麼我的人生這麼痛苦呢?.
為什麼我裡面的感覺長期都這麼痛苦呢?.
為什麼我生命這麼多問題呢?.
為什麼我跟人相處.
或者我去到一些關口.
生命的關口.
會出現這麼多問題重重複複的呢?.
我自己覺得.
哇,這種人生很累.

$^{441}$所以我想說.
這種累是因為我過往的生命.
不斷地為別人的生命負責.
你的事原來是關我事的.
什麼都關我事.
但是呢.
而我卻沒有真正為自己的生命負責.
我不想自己是這樣.
所以我很想看到神創造我的美好.
所以當人花很多錢去裝身.
花很多錢和時間去做其他事的時候.
我選擇花我的錢和時間去見輔導.
聽起來也覺得很無謂.
又浪費時間又不知道在做什麼.
坐著閉著眼睛看著我.
不是的.
我想說.
我不單是花錢去見輔導.
還有參加一些生命成長的課程.
你會問.
這麼浪費錢為了什麼呢?.
為什麼我會投資在這麼不顯眼的事上呢?.
你穿一件漂亮的衣服也有人稱讚.
你裡面是怎樣別人不會知道的.
你不說別人不知道.
但是我想說.
在我們生命裡面.
我會覺得我需要一塊專業的鏡子.
我需要一塊專業的鏡子去照一下我內心的樣子.
你臉上黏著一顆飯.
你也要別人提醒你才會撕下來.
我的生命也是.
我需要別人去提醒我.
我需要有人去幫助我反省自己成長的問題.
面對我生命當中的傷口.
如果傷口不理會它.
它會發炎.
再不理會它.
它會爛掉.
再不理會它.

$^{481}$它會發臭.
因為當我說要去服侍神.
要去服侍人的時候.
我需要有一個健康的自己.
才能夠服侍.
我在家庭當中得不到的愛.
但是我要重新去學習.
怎樣去愛自己.
以至到我才有能力去愛.
如果不是.
我去過多少教會.
我聽過多少年道.
我的恩賜有多少也好.
我裡面也是爛的.
好像一個橙一樣.
你表面很漂亮.
噴了防腐劑.
但是裡面切開的時候也是爛的.
所以如果我昔日不願意面對自己的生命.
今天你遇見的我.
不會做到我現在做到的事.
也不會做到我現在做的事.
例如有人找我做spiritual direction的時候.
說話給我聽的時候.
我不會有這種承載力.
這種capacity去載你們的東西.
或者我會哭得比你更慘.
然後你告訴我你的慘事.
我會告訴你整個小時內我的慘事.
完全不會聆聽到你.
所以我想說.
如果這樣的我.
你說怎麼辦呢.
所以最後我很想鼓勵大家.
去反思.
今天我們在這個回教會.
我們聽到.
我們已經很多年了.
我們也很熟悉.
但是今天神去鼓勵我們.

$^{521}$就是無論你拿著的五千.
無論你拿著的是二千.
無論神給你的是一千.
我想說就算你拿著一千.
只要你盡力去倍增的話.
你同樣是得到上帝給你的讚賞和賞賜.
在神的角度的藍圖裡面.
你有沒有發現上帝給你的那一份.
你有沒有盡力去發揮呢.
你盡力了沒有呢.
第二件事就是.
信了耶穌不代表你的生命是自然的成長.
你會選擇讓自己去逃避你生命的問題.
還是你願意為自己生命負責任.
正視自己生命的一些問題.
面對它.
靠著上帝去改變呢.
禁食吾軀歸故土.
他消君體也相同.
今天我們怎樣讓天國的視野.
能夠影響我們今天的人生呢.
最後我想說.
我這次預備講道的時候.
寫著寫著講章.
我寫了一首歌.
很想藉著這首歌帶大家去做一些反思.
同時也讓這首歌成為我們的祈禱.
大家可以聽一下.
《我讓望你手》.
詞:陳靜楠 曲:陳靜楠.
我讓望你手.
抱似滿天星宿.
自覺多麼渺小.
神你竟倦猶.
三眼一月如光波穿透.
這生道轉沙流.
願沙粒不會往下白流.
你的世界設計獨特.
漫遊冷暖春秋.
屬天的心意託付.

$^{561}$我肩旁和兩手.
主的聲音我聽後.
隨著走.
神國島南途.
燃起我心中沉睡的夢.
縱然諷刺獨個走窄路.
你的讚賞是我最大回報.
神我願拜佛.
你的國是我一生渴望.
永恆的愛召喚我傾倒.
我僅有的所有.
永恆的愛讓我感受.
獻出我的所有.
讓我們一起去禱告.
這一刻.
我們這位主.
就來到大家面前.
去將那些銀紙.
分派給你們每一個.
是神說.
孩子我給你的這一份是獨特的.
是沒有人能夠取締的.
你試一下想像一下.
這位主人來到你面前.
在你生命裡面.
安放了什麼給你.
你是五千那個.
是二千.
還是一千.
你試一下細心的去望一下.
你手上的.
神給你的那一份.
你會跟主說.
什麼.
這一刻我又邀請你.
你去問下神.
在你生命裡面有什麼地方.
是神很想鼓勵你去突破去改變.
或者你生命裡面有些問題循環又循環.
出現又再出現.

$^{601}$可能你習慣為很多其他人負責.
但是神今天去鼓勵你.
孩子.
你願不願意正視你生命.
需要改變的地方.
你會怎樣回答他.
我給大家一分鐘時間去.
跟神有些祈禱.
最後我會帶大家結束禱告.
愛我們的主.
我們多謝你.
在神國度的藍圖裡面.
你給了我們每一個人都有獨特的一份.
主你給我們見到我們有能力.
並且要去回應你.
你給我們有更加多的愛去突破自己.
我也求你給我們有勇氣.
因為成長會有眼淚.
會有掙扎.
會辛苦會艱難.
會很不情願.
但是神你給我們都有這份力量.
有這份力氣.
去面對.
以致我們有一個更加健康.
更加能夠發光的生命來榮耀你.
最後多謝你.
我們禱告.
奉主耶穌基督的名禱.
阿們.
多謝您的收看.
\newpage



\section{}
\label{sec:iar_apDpNcQ}
\textbf{神愛大地 (新約篇) - 陳崇基牧師}
\newline
\newline
連結: \href{https://youtube.com/watch?v=iar-apDpNcQ}{\texttt{https://youtube.com/watch?v=iar-apDpNcQ}} ~~~~ 語音日期: 2025-02-25
\newline
\newline
\hyperref[sec:UraagaffkTo]{\small{< < < PREV SERMON < < <}}
~
\hyperref[sec:index_chronic]{\small{[返順時目]}}
~
\hyperref[sec:index_scriptual]{\small{[返順卷目]}}
~
\hyperref[sec:IH4xN4bYlR0]{\small{> > > NEXT SERMON > > >}}
\newline
\newline
$^{1}$各位好.
我們又講之前講過的topic.
就是Creation Care.
就是和大地之間的關係.
如何關愛受造界.
我講了兩次.
今次是第三次.
是吧.
應該大家開始熟悉了.
我們講過一次.
第一次就是講羅馬書第八章.
然後第二次我們講的是舊約篇.
今次又是講新約篇.
你記得這個時間表.
這個就是起碼我從小到大.
在教會裡面接觸的時間表.
就是現在這個世界.
現在的天地.
被罪所破壞.
等待將來耶穌再來的時候.
耶穌再來之前會發生什麼事呢.
有沒有聽過七年大災難.
有吧.
有誰沒有聽過.
沒有聽過是好的.
因為你要再去想一下.
經文到底是不是真的講這類東西.
是不是.
這個就是我們整個系列裡面.
我們希望能夠做一些神學的反省.
傳統的時間表.
傳統的概念.
將來這個新天地和我們現在沒有什麼關係.
是完全中斷的.
現在有的東西就像垃圾一樣扔掉.
然後你看到有個核爆的蘑菇.
核爆蘑菇炸毀了.
我們什麼都不需要理.
這個就是我們有一些.
本身有的預設.

$^{41}$如果你聽過末世會有大災難.
你想一下其實和我們生活很有關係.
因為如果是這樣的話.
這個世界是沒得救的.
我們最好走人.
還有這個世界如果這麼亂.
越來越多戰爭.
越來越多打仗的時候.
我們會關心嗎.
我們應該不關心的.
為什麼.
因為中東每逢有什麼動靜的時候.
我們就說聖經預言說了.
我們沒什麼好說的.
不用反戰.
不用譴責哪一個開戰等等.
發動戰爭等等.
我們不用說這些倫理的東西.
因為聖經預言已經實現了.
聖經已經說了世界會大亂.
我們應該做什麼.
逃跑.
這個是傳統的觀念.
我們想問一問這個傳統觀念.
是不是真的合乎聖經裡面整體所說的.
我們用兩段新約的經文.
第一段是彼得後書.
第二段我們回到第一集.
我們說過的羅馬書第八章.
彼得後書說到主在來的時候.
到底會發生什麼事.
還有現在的天地和將來的新天新地.
到底有什麼關係.
三章第八節.
他說親愛的有一件事你不可以忘記.
就是主看一日如千年.
千年如一日.
主沒有辭然他的認許.
就是他再回來的認許.
就是耶穌基督再回來的認許.

$^{81}$就如有人以為他辭然.
其實他是寬容你們.
不願一人沉淪.
願意人人都悔改.
但主的日子要像賊一樣來到.
這個很簡單.
彼得後書.
彼得前書說什麼.
彼得前後書.
前書就是處理教會.
大概教會外面遇見的困難.
一些迫害.
一些逼迫等等.
有外患所謂.
彼得前書就是勸勉教會.
能夠繼續向前.
不需要派一些這樣的東西.
繼續堅守著他們的信仰.
這個彼得前書.
彼得後書是說內患的.
什麼內患呢.
彼得後書裡面說到一些假教師.
一些假教師就會說.
耶穌遲了沒有回來.
他一定不回來.
有沒有約過朋友去放飛機.
遲了來.
他不來.
有些假教師就說.
耶穌也會是這樣.
所以彼得後書到第三章的時候.
他就說什麼呢.
他說不是的.
主一定會回來.
OK.
因為上帝看時間.
不同我們看時間.
千年如一日.
一日如千年.
首先神的時間表.

$^{121}$時間的觀念.
是和我們的時間觀念不同.
大家有些有小孩的知道.
我們開車多人多.
開多久.
九個小時.
你小時候帶他們.
起碼我兩個小時候帶他們.
去多人多的時候.
他們一上車.
你要開九個小時才到.
九分鐘之後他們就說.
到沒有.
他不明白九小時等於九分鐘.
九分鐘等於九小時這個觀念.
這個就是彼得後書裡面所說的.
他說神的遲疑.
目的.
不是因為他放飛機.
而是他寬容人.
給人有時間.
給人有機會.
去到歸向他.
並且主再回來的時候.
好像賊一樣.
賊是沒得預期的.
沒得約定的.
沒得約定賊來打你.
沒有的.
主回來也是一樣.
然後他繼續說.
主一定會再來.
然後他就說哪日.
天.
哪日主耶穌再回來的時候.
天就必然.
有轟然一聲鐘消失.
天體都要被烈火融化.
地和地上的萬物都要燒盡.
基而這一切都要如此消失.

$^{161}$你們處世為人.
必須要聖潔敬虔.
等候並催促上帝的日子來到.
因為在哪日.
天要被火消滅.
天體都要被烈火融化.
但他的應許.
主再來的應許.
我們等候新天新地.
其中會有正義公義常駐.
你看到這一段經文的時候.
我們現代人看這些字眼.
我們一定想起什麼.
核爆.
核爆冬菇.
這些就是我們年輕的時候.
我還記得中學大學的時候.
在教會.
在一些很傳統的教會.
主日學,港台都是說這些.
現在有很多核子彈.
全世界.
可以地球爆很多次.
那個應用就是.
聖經都是這樣說.
所以可能世界的終末的時候.
就會一切燒了.
整個世界就會被核爆所消滅.
到底怎樣看這些經文.
首先第一件事.
大家知道.
這些經文.
聖經裡面有不同的題材.
有敘事,先知的預言.
有詩篇等等.
這些叫做天啟文學.
Apocalyptic literature.
天啟文學.
不是我們平時會看的東西.
是當時猶太人流行的文學.

$^{201}$所以他們用這些題材.
去說出一些東西的時候.
你要明白.
他們的字眼和用語.
或者用的比喻.
都是比較誇張和戲劇性的.
比較戲劇性的.
因為說到世界的末了.
我們就不能夠用我們現代人的字面的解釋.
現代人字面的意思.
就塞進去套進去.
情形有一點像.
創世記第一章.
六日創造.
我們一看的時候.
我們就覺得是一個科學報告.
因為我們是科學世界的人.
我們有iPad,iPhone等等.
我們一想起六日創造.
一定是說六日.
每日24小時.
你就搞了一大輪.
不需要搞出來的問題.
頭三日沒有太陽,月亮.
怎麼會有晚上,早晨.
因為第四日才有太陽和月亮.
上星期五我們有國語小組.
我們有一個無道的朋友.
他想知道聖經.
聖經裡面有些東西卡住他.
他就跑來查經.
我們就說開創世記的經文.
如果你用創世記六日創造.
當為一個現代的科學報告的話.
你一定出事的.
一定出事的.
頭三日是怎樣.
我和他解釋.
頭三日晚上有早晨.
明明有晚上有早晨.

$^{241}$就是有出日落.
頭三日沒有太陽,月亮.
第四日才有.
頭三日晚上有早晨是怎樣.
上帝弄一個電筒出來.
太陽能電筒.
開著開著.
你不可以當它是一個.
現代的科學報告來看.
一定是這樣出事的.
看過的天文學都是一樣.
你不可以按照我們現代人的字面.
想到的東西.
我們在想的是什麼.
Physics,Chemistry.
融化.
樣樣東西爆炸燒掉就搞定.
然後這個世界就玩完了.
這樣就一定出事.
你說怎樣讀呢.
我想告訴大家.
你一定要報我啟示錄的班.
真的.
一定要報.
因為你不懂得看天文學.
我們就亂將它.
現代人的那種預先寫好的報紙.
然後我們就塞進去.
然後我們就.
見到新聞裡面發生什麼.
中間有什麼異獄.
有什麼核爆等等.
我們就說聖經預言是這樣.
No.
不是預先寫好的報紙.
當然我們不能夠很詳細去解釋.
所以為什麼你要報我班.
詳細一點去解釋.
我們解釋一下.
他說天有轟然一聲消失.

$^{281}$英文翻譯成Passed away.
比較好一點.
或者Dissolve.
解開這樣的意思.
或者你有一些詩意去看的時候.
這些天文學的字眼.
其實有一點像交響樂曲.
交響樂曲有很多Movement.
對不對.
有沒有聽過交響樂曲.
沒有.
大家都是沒有文化.
我們.
但是現在不要緊.
我告訴你怎麼聽.
有些Movement.
First Movement.
Second Movement.
他們說First Movement.
Passed away.
完結了.
那是不是炸掉了.
當然不是.
你First Movement完結了.
到Second Movement就開始了.
是不是.
然後有Third Movement.
繼續下去.
大概應該是這樣的意思.
萬物都要消盡.
我們看聖經的時候.
有時候看Footnote.
Footnote寫.
古卷作暴露無遺.
Lay bare.
英文的翻譯.
Lay bare.
即是Naked.
即是萬物會在上帝面前.
上帝再回來的時候.

$^{321}$主耶穌再回來的時候.
萬物都會Lay bare.
在說什麼.
說上帝會審判.
會潔淨.
會審核整個世界.
將有罪的東西拿走.
解決一切不義的事.
上帝審判萬物的時候.
萬物就在上帝面前.
就是Lay bare.
暴露無遺.
Nowhere to hide.
沒得走.
有沒有照過X光.
什麼都照出來.
MRI送你去.
什麼都照出來.
你在法官面前.
在審判者面前.
你就是Nowhere to hide.
因為上帝一定會撥亂反正.
這樣的意思.
火燒和融化.
這個就有點棘手.
我們說火燒融化.
還不是核爆.
還不是一件東西燒毀了.
然後就扔掉.
這些東西不要了.
然後弄些新的東西出來.
現在這個世界.
好像垃圾扔了.
將來世界有一個全新的.
《聖經》裡面用火.
這種的metaphor.
很多時候不一定是負面的.
很多時候反而是正面.
下面這裡我引用了一段經文.
彼得前書.

$^{361}$怎麼用火.
他說我們在法官的試煉裡面是憂愁.
但是我們因此而喜樂.
為什麼.
因為在試煉裡面.
你的信心就會被考驗.
就比那被火試煉仍然能懷的金子.
更加顯得寶貴.
這個火是好還是不好.
你在試煉裡面.
用火來煉證你.
是不是將你扔毀了.
是不是將你垃圾扔了.
然後弄個新的.
No.
它是purify.
refining fire.
就是留下出來的.
就像以前那些煉火煉鋼的.
將雜質去到消滅.
然後留下的就是純潔的金或者銀.
所以這裡絕對不是說爆掉了.
然後新的就來了.
如果你看一下上文.
彼得後書三章上文.
他用另外一個聖經裡面的故事來形容.
主在來的日子.
就是洪水.
創世記洪水.
羅亞的洪水來到.
洪水是不是摧毀.
不要這個天地.
然後弄個新的.
不是.
是洗淨.
洗淨.
然後方舟就是天地的縮影.
所以羅亞就叫那些動物進來.
潔淨的不潔淨的全部都走進來.
方舟本身就是一個小的天地.

$^{401}$上帝保留著.
保存著.
受造界.
就用一個這樣的比喻來形容主在來.
然後下面他就用另外一個比喻.
用火的比喻.
所以這裡絕對不是冬菇.
不是核爆冬菇.
應該是想起的.
這段經文的時候.
是煉鐵.
refining.
要燒到很高溫.
然後將雜質拿走.
這個世界上帝最後的時候.
會做一個更新.
完整的更新.
然後新天新地就會來到.
所以這裡只是給大家一點點提示.
到底說到我們關愛受造界的時候.
我們不能夠.
好像以往或者傳統的.
我們這樣看.
我們每樣東西都玩完了.
我們不用搞了.
我們走.
No 上帝是refine.
將來都會有一個最終極的refine.
我們回到這段經文.
就是第一集的時候我們說過.
八章二十節.
因為受造之物處於虛空之下.
並非出於情願.
是出於上帝的旨意.
不過受造之物仍然是盼望脫離.
衰敗衰亡的束縛.
是享受上帝而他榮耀的自由.
我們知道一切受造之物至今仍然是.
身淫如同經歷分娩之童.
這裡我們第一集的時候說過.

$^{441}$大地是很不妥.
他去到形容是以人法的形容.
即是寸寸不息在吐大氣.
因為是被折磨緊緊.
我們第一次的時候解釋.
應該是說到創世的第三章.
人類叛變之後.
聚廠入世然後破壞大地.
地就受咒詛.
是罪惡.
所以大地仍然是在罪惡裡面.
但羅馬書第一章至第八章.
就說上帝的救恩已經到了.
就是重新恢復原本有的關係.
大地就是在等原本有的美好的關係.
再復修.
所以現在還沒有到的時候.
他就在身淫在嘆息.
這是一個神學的.
或者聖經的背景的解釋.
其實你也可以用一個歷史的背景解釋.
即是羅馬書寫的時候.
是當時主後第一世紀.
當時是羅馬帝國最興盛的時候.
最興盛的一個大國.
帝國一定是強國.
一定是大興土木.
我們經常去歐洲.
或者去以色列也好.
我們看很多時候第一世紀.
羅馬興盛的時候.
他們做很多事情.
大興土做很多事情.
然後有很多水稻.
那些很厲害的羅馬盛世.
Pax Romana.
我們想像的就是一些金碧輝煌的日子.
但如果你從一個環境的角度去看.
Environmental的perspective.
你還記得我大學的時候讀什麼嗎.

$^{481}$Environmental toxicology.
你用環保的角度去看經文的時候.
你就發現原來古代的.
或者近代的一些資料.
告訴我們羅馬盛世的時候發生了什麼事.
表面上好像是強大的.
但是因為帝國強大.
一定是講很多生產.
大量的去生產農業.
然後土地就要榨乾榨盡.
用盡它.
然後用完那些田地就變成荒蕪.
森林全部都棄掉.
那些樹木全部棄掉.
然後弄一些資源.
兵器等等.
然後就濫用自然的資源.
金礦銀礦石礦.
全部東西都拿出來.
這個是羅馬盛世.
古代現代都有些資料.
是給我們一些線索的.
以前有一個地理學家叫做Strabel.
Strabel就是那些貴族.
你當他是貴族.
今天有很多貴族.
有空就去歐洲玩玩.
然後又回家.
不用工作.
然後又去哪裡玩.
Strabel大概就是這樣.
所以他就周圍周遊列國.
然後寫很多的日記.
去了哪裡.
他說到了一個地方叫比薩.
大家聽過比薩斜塔.
在意大利羅馬北面的一個town.
叫做比薩.
他說到比薩以前文明於什麼呢.
就是有很多很漂亮的土木.

$^{521}$和很肥沃的土地.
因為羅馬盛世的緣故.
就砍掉這些木材.
然後去造船.
戰船.
戰艦.
但是如今大部分的木材的資源.
都已經被羅馬的建築拿走了.
你可以想像Strabel去到比薩斜塔玩的時候.
他當然是看不到比薩斜塔.
那時候還沒有比薩斜塔.
他就看到這些森林.
全部被羅馬帝國的政府.
全部濫用並且砍伐.
你可以想像.
是一種這樣的情況.
過分砍伐樹木.
是一件很大的事.
對環境很大破壞.
我們不是說不砍伐.
而是過分砍伐.
Over deforestation.
會發生什麼事呢.
農田就會被.
因為你要發展農田.
你就要砍這些樹.
砍了樹木當然可以用.
然後砍了那個地方.
你就變成了農田.
然後將它不斷耕作.
不讓它休息.
然後你有很多的農產.
用完了.
然後你就再砍其他樹.
這些土地就會被侵蝕.
然後那些土壤.
有Soil erosion.
被風化.
然後它們就變成荒地.
這是一個惡性循環.

$^{561}$在羅馬盛世的時候.
其實是有這樣的發生.
這裡我引用一兩個.
很近代的.
2023,2024年做的.
近代的研究.
Environmental research.
Environmental archaeology.
有趣的發現是什麼呢.
就是羅馬人當時大興土木的時候.
他們去採礦,煉礦.
金銀礦物.
你需要火燒.
剛才說了.
然後你就提煉到金銀出來.
或者那些銀幣.
怎樣弄出來.
就是在銀礦那裡.
去到火煉出來.
但是你去到搞這些採礦.
然後煉礦.
就會造成很多什麼.
污染物.
就是工業的廢料.
譬如鉛.
鋒利.
鉛就到處丟.
然後就會污染環境.
這裡有些研究很有趣.
就是將羅馬時代的人骨.
用來做化驗.
就是將你拆骨.
然後將你化驗.
就發現你骨裡面含鉛量過高.
你知道含鉛量過高會對大家有什麼影響嗎.
蠢.
沒錯.
所以Wall Street Journal.
在今年年頭的時候.
就引用這些研究.

$^{601}$因為當時羅馬盛世的時候.
就是奧古士都.
就是公元前27年.
直到馬可奧利略去世的時候.
大概有一百年左右的時間.
這段時間羅馬盛世.
Pax Romana的時候.
就有相對比較高的工業污染.
很多都是污染物.
可以在羅馬人中間.
會使得他的IQ減2.5至3分.
當你國家盛世的時候.
很多時候人就蠢一點.
不知道現在是不是美國.
就是可以打敗你的IQ Average.
直到3分.
所以現在的人就不會計算.
你想一下羅馬盛世的時候.
你就不應該想什麼呢.
不應該想左邊.
如果用環境的角度去看.
應該是想右邊.
就是污煙瘴氣的意思.
如果這樣看的時候.
這個背景呢.
你會發現羅馬書第八章.
保羅說大地.
寸寸不息的時候.
新人的時候.
其實他自己到處不同的地方宣教.
主要的城市的時候.
他看到的可能是右邊.
一些土地.
大地被破壞.
被榨乾榨了.
所以保羅說這個造界就是新人.
所以說羅馬盛世.
你不要想光輝歲月.
可能是右邊.
是污煙瘴氣.

$^{641}$我讀大學的時候.
Environmental Toxicology.
大學的時候.
預備這個系列.
我看回以前的書.
我真的激氣.
以前又不讀書.
現在又可以看回.
我說過我在實驗室做.
我進實驗室.
是緊急制.
我去驗水.
其中一個大客.
就是McCurran Air Force Base.
在加州.
北加州.
我們經常接到一些樣本.
驗泥土和驗水.
為什麼呢.
因為這些軍事基地.
有很多工業的廢料.
一些重金屬的廢料.
很多時候處理得不好的時候.
就會污染地下水.
污染河流.
污染湖泊.
污染空氣.
然後就跑到魚那裡.
跑到生物那裡.
又跑到菜那裡.
人們吃了之後.
就會很大件事.
如果要清潔的話.
要花幾千萬去清潔.
很貴.
所以我們其中一個大客.
就是這種大的工業.
他們還寄一些泥土和水.
去做化驗.
我忘記了.

$^{681}$我做了幾年.
然後我就.
以前是讀醫院.
不是讀得很好.
我就去讀神河.
我離開了公司.
我離開公司一年.
我就聽到一個消息.
因為我有家人在那裡上班.
突然有一天.
就有聯邦的人員.
去封了實驗室.
那些人上班.
早上八點上班的時候.
沒得上班.
全部的代理人就全部封了.
整個實驗室.
原來追查發生什麼事.
後來發現有兩個化學家.
兩個化學家.
因為我們做完這些測試.
有一個化學家去簽名.
簽名說有沒有污染.
需不需要留意等等.
多少個毫米.
多少ppm等等.
我們就做這些報告.
最大的化學家一定要簽名.
原來抓到這兩個最大的化學家.
他們做數.
他們騙人.
有些樣本明明沒有測試.
但他們說測試了.
其實做這些事很悶.
每次測試的東西.
其實預修都差不多一樣.
其實我做的.
有時候做的報告.
是複製和編輯.
幾乎複製和編輯.

$^{721}$因為他是1.002.
下次是1.003.
有什麼分別.
原來裡面有多年做數.
所以就抓了這兩個化學家.
我想告訴你.
這個例子原因是什麼.
原因是什麼.
其實是一個很大的試探.
因為你說的是錢.
我們當然想抓乾抓盡.
污染食水我們不理它.
危害到人生命.
我們不理它.
我們住在芝加哥.
或許少一點.
這些情況.
因為有管制.
你都聽到很多不同的新聞.
說不同的地方.
有食水的污染.
是不是.
是很大的一件事.
因為這是一個價值觀的問題.
我們怎樣去看.
這個大自然.
怎樣去看天地的.
天然的資源.
你記得我讀的時候.
有一本天書叫做Silent Spring.
Rachel Carson.
這個是Biologist寫的.
他開頭第一頁.
他就引用.
一個人叫做.
Albert Schlesinger.
所說的話.
他是一個德裔的神學家.
一個很出色的20世紀的聖經學者.
神學家.

$^{761}$他說什麼.
他說人類已經.
失去了預見.
預知的本能.
最終他就會摧毀這個世界.
是我們摧毀這個世界.
是因為我們的價值觀.
我們覺得.
大自然是用完即棄.
用完了.
就當垃圾.
租大它不要緊.
對不對.
總之不同種族有不同的人.
都有這樣的弊病.
我覺得中國人華人都有.
去大自然.
以前加州去的.
很美的地方.
你把垃圾亂扔.
是不是.
價值觀問題.
我看完就打完卡.
不用要.
我們會糟蹋.
不會愛惜.
因為我們消耗.
我們覺新就是消耗者.
我們consume了它.
吃完它.
用完它.
然後就扔.
這個絕對不是聖經的價值觀.
所以最後的時候.
我希望大家留意幾樣東西.
新的或更加完備的.
一個時間表.
現在的世界和將來新天新地.
那個關係是有一個continuity.
最好的形容方法.

$^{801}$應該是更新.
不是核爆.
是更新renew.
或者renovation.
將它翻新.
好像你買一間舊屋你將它空置了.
然後再翻新.
你不是扔了它.
將東西重新renovate.
一個新天新地.
有幾樣東西是值得我們去實踐的.
第一我們要知道上帝是關心這個世界.
這個世界和將來的世界.
不是一個replacement.
而是renew.
renovation.
第二我們要做一個好的管教.
我們大自然上做一個管教.
上帝賜給我們的資源.
我們做好我們自己的小角色.
和上帝同工.
這是第八章.
同工for good.
我知道大家可能會問一個問題.
我可以做些什麼.
我這樣可以recycle.
然後做些環保的東西.
那些大的石油公司.
在深海轉一轉.
然後漏些油出來.
污染的程度大到不得了.
我這樣一滴水.
有什麼用.
這個又是價值觀的問題.
我們是不是實用主義者.
我們覺得work的就是真.
不work的我們用來做什麼.
我們決定我們的行為.
我們的生活.
很多時候是用實用主義的價值觀.

$^{841}$而不是用聖經天國的價值觀.
對不對.
你用什麼我奉獻.
大把財團奉獻.
我奉獻一點點用什麼.
我做一點點可以給什麼影響.
這是實用主義.
pragmatism.
應該是天國價值觀帶領著我們.
所以我們儘管做好我們的管家.
第三樣.
就是簡樸的生活.
對不對.
我想問為什麼你買這麼多東西.
你用得到嗎.
你用得完嗎.
對不對.
我有朋友家裡有150雙鞋.
除非你是白竹.
不好意思如果說中了你.
或者家裡有很多玩具.
男人家裡也有很多玩具.
但是很多時候不等用.
放在那裡是封塵的.
簡樸一點的生活.
以免我們浪費.
因為我們不是消費者.
我們是管家.
第四樣是欣賞主的創造.
我們最後有一首回應.
試過這次我們請了Bella親自上來唱.
我們歡迎Bella上來.
多謝大家拍手掌.
既然上來唱.
不如介紹多一點這首歌的背景.
我寫這首歌是源於Theo Music的邀請.
Theo Music是什麼呢.
就是香港以一個神學反思為主題的信仰平台.
透過原創的音樂和神學文章.
幫助信徒和微信人.

$^{881}$用不同的角度去反思信仰.
令到神學不單是一種學術知識.
聽見神學好像很深奧的東西.
而是進入我們生命.
讓我們可以在時代.
在我們的處境裡面去反思生命的一個平台.
所以他們每年都會邀請不同的創作人去寫歌.
將神學反思轉為成音樂.
這次我選擇了土地神學作為主題.
當我寫這首歌的時候.
其實我想起一個字.
就是進.
剛才Mark Sir也說了很多.
在這個要進的文化裡.
我們對自己的要求很盡.
我們對別人的要求很盡.
我們對土地和萬物的資源.
我們都一樣要抓得很盡.
因為很多時候我們人生的價值觀.
或是我們的安全感取決於我們能夠爭取多少.
能夠掌控多少.
所以我寫這首歸土.
正正是要提醒我們去記得.
記得我們來自塵土.
最終都會歸回塵土.
記得這片土地不單是我們的資源.
更加是上主的恩賜.
更加是盛載萬物的家.
記得萬物都需要安息.
而不是不斷被榨取.
願意這首歌.
我們一起去沉澱下來.
靜下來.
去聽一下這個土地的故事.
也聽一下上主的心聲.
《一村一介一遇》 詞:陳曦 曲:陳曦.
一村一介一遇 一生年日傾注.
痛想起價值標誌 斑駁倦容.
為求踏著這寸土可掌控.
殘缺未夠盡怎可放鬆.

$^{921}$一化一土一介 一刻退去消散.
如燕棲處碧兒灘 冰山崩壞.
南京活上各自厭厭盛嘆.
日好盡了夜才追贊.
世亭海枯瀚 地也蠢蝕.
世界的宿命 道地要遣儉.
讓地歸於我來 萬物靜讓生死.
歸後條件領向 沙遊水靜.
地屬主所指 蠟藥鐵針.
流離可安定 坐看風景.
受日的得高堂 回頭的終可放開陣勁.
眾生年 期盼重拾風聲.
《一生一介一遇》 詞:陳曦 曲:陳曦.
一算一介一遇 一生年日傾注.
痛心你價值標籤 斑駁倦容.
為求踏著這村高可長紅.
三卷未夠盡怎可服輸.
一花一桃一眼 一刻退去消失.
如煙繫住鼻兒他 冰山崩壞.
南極沒相割 千煙染色他.
一號進了也在追他.
世聽海呼喊 地也喧囂.
世界的宿命 到了險境.
若地歸於安寧 難得靜養心聲.
歸後了於領漢 沙遊書城.
地屬於初次 蠟藥鐵針.
流離可肯定 坐盼風景.
受壓的不公平 雷霆的寵我放開震驚.
終生無期盼重新方醒.
世聽海呼喊 地也喧囂.
世界的宿命 到了險境.
若地歸於安寧 難得靜養心聲.
歸後了於領漢 沙遊書城.
地屬於初次 蠟藥鐵針.
流離可肯定 坐盼風景.
受壓的不公平 雷霆的寵我放開震驚.
終生無期盼重新方醒.
《你我之歌》 詞:陳曦 曲:陳曦.
你我只不過是個客隊.
也有天花灰聚 皇帝討睡.

$^{961}$如流星一閃即去 別贊取多餘瑣碎.
讓孤寡可寄居 你我生於這裡.
多謝您收睇時局新聞,再會!.
\newpage



\section{5:33-42}
\label{sec:IH4xN4bYlR0}
\textbf{Righteousness as Redefined by Jesus: Word and Deed (Matthew 5\_33-42) - Chee Yie Jong}
\newline
\newline
連結: \href{https://youtube.com/watch?v=IH4xN4bYlR0}{\texttt{https://youtube.com/watch?v=IH4xN4bYlR0}} ~~~~ 語音日期: 2025-03-05
\newline
\newline
\hyperref[sec:iar_apDpNcQ]{\small{< < < PREV SERMON < < <}}
~
\hyperref[sec:index_chronic]{\small{[返順時目]}}
~
\hyperref[sec:index_scriptual]{\small{[返順卷目]}}
~
\hyperref[sec:LHY953jFc_g]{\small{> > > NEXT SERMON > > >}}
\newline
\newline
$^{1}$Good morning/afternoonish, I'm Yuvi. It's good to see a lot of you again..
So welcome back to church and to worship with us..
Okay, so every year, one of the biggest sporting events that we have in the US is the Super Bowl, right?.
And this year, it is estimated about 127 million viewers tune in for the Super Bowl in New Orleans..
And this is probably on the more conservative ends as well, because I would assume there are probably a lot of Super Bowl watch parties that are going on..
And we actually have one here right at UV as well, and I know some of you were here with us that night..
So for some people, they tune in to the Super Bowl because they are the fans of the game, or they are fans of either of those teams..
But for another group of people, they tune in to the Super Bowl because of the commercials, but more importantly, because of the halftime show..
Do any of you fall into this second category? I'm sure there are some of us that belong to that group, right?.
So the halftime show for the Super Bowl this year was headlined by Kendrick Lamar, who is arguably one of the biggest names in the hip-hop scene..
And if you actually pay close attention to his halftime show performance this year, all the way from the stage design, the Uncle Sam persona by Samuel L. Jackson,.
the dancers that are dressed in the colors of American flags, and even the reference to 40 acres and a mule, just to name a few of those..
His performance was actually a creative blend of high-level entertainment that is accompanied by deep social and political commentary,.
where he addresses issues such as racial inequalities in our nation's history, and also cultural identity being in America..
But the pivotal moment for Kendrick Lamar's halftime show performance was when he performed his award-winning Grammy Song of the Year called "Not Like Us.".
This song is a diss track targeting another raper called Drake..
So this feud between Kendrick Lamar and Drake has been going on for quite some time now..
But by performing this particular song in the Super Bowl, Kendrick Lamar brought the ongoing personal but public rivalry between him and Drake.
all the way to the biggest stage of the world and put it on display for everyone to see, pretty much..
However, this kind of public rivalry and conflict does not only happen in our pop culture..
The sad reality is that it happens within the church as well..
If you just look at the church histories, there is the Great Schism in 1054 AD,.
where a split happened between the church with the Roman Catholic Church in the west and the Eastern Orthodox Church in the east..
There's also the Protestant Reformation, which eventually led to the Thirty Years' War between the Catholic and Protestant states in Europe as well..
A couple of weeks ago, I saw a Facebook post from one of the biggest Christian evangelical organizations in the U.S., I think..
Basically, the post is kind of like critiquing the ministry and the work of another pastor who is from a different theological tradition..
And as you can imagine, the comment section pretty much just becomes a place for public, virtual arguments between these two sides of people..
So why did I bring this up? These examples are just a snippet of what is going on in our current world..
Even now, as I'm speaking, there is the war happening in Gaza in the Middle East..
And we're also entering into the third anniversary of the war that is happening between Ukraine and Russia as well..
And if you have been following along the news, especially with the meeting between President Trump and President Zelensky,.
please pray for the leaders of our country that they will care for the things that Jesus truly cared about, for justice and peace,.
not to show our strength or fall into kind of like the lies of American exceptionalism..
But with the ever-increasing division in our society, including within our church, how should we as followers of Jesus respond?.
We have been studying through Jesus' Sermon on the Mount for the past couple of months now,.
and I think the Sermon on the Mount is especially relevant to the cultural moment that we are currently in..
Because in it, Jesus lays out how life in his kingdom ought to look like and how we, as believers, can take part of it..
So in my first sermon, we look at the Beatitudes in which Jesus describes the attitude or character of those who are considered as blessed in his kingdom..
And one of them is this, "Blessed are the peacemakers, for they will be called children of God.".
To understand what a peacemaker is, we need to first understand what peace is exactly..

$^{41}$For a lot of us, when we think about peace, we just think of the absence of conflict, right?.
Maybe you live with a couple of roommates. One of them, they just like to leave all the dishes in the sink unwashed..
And another one of them, maybe they just like to leave the dirty clothes in the washer or dryer for weeks..
And when you have to use it, you just open it up and say, "Well, I guess I have to wait for the dirty clothes to get cleared out first.".
And maybe for you, peace is not confronting them with those issues, even though inside of you, you don't really experience peace at all..
But because for you, peace is just the absence of any sort of conflict..
But the idea of peace in the Bible is more than just the absence of conflict..
It also means the presence of something better..
In the Old Testament, the word for peace in Hebrew is "shalom.".
And shalom represents a state of being whole or being complete, in contrast to the state of chaos, brokenness, or disorder..
And this is why the Jews, they usually greet each other with "shalom," hope you are doing well..
And when we come to the New Testament, the birth of Jesus is being announced as the arrival of peace, or "erene," in Greek, on earth..
Why? Because through Jesus, he made peace between sinful human and a holy God through his death and resurrection..
Or in other words, Jesus restores a broken relationship between human and God that has been tainted by sin since the beginning in Genesis, back to wholeness..
And not only that, he now calls his followers to do the same, to carry on this peace into the world..
So, to be a peacemaker doesn't just mean that we do our best to avoid conflict..
If you see a brother or sister living in sin, false peace might look like this, not calling them out, because you want to avoid any confrontation and conflict..
But true peace, on the other hand, would mean pointing out their blind spots because you love this brother or sister, and you want to see them experience true shalom or wholeness..
So, the section in the Sermon on the Mount that we're currently in explores the better righteousness that Jesus has redefined for his followers..
So far, we have looked at how Jesus redefined his greater righteousness, and we discovered God's heart behind the sixth and the seventh commandments in the Ten Commandments, right?.
That is the kind of attitude in our hearts that treats everyone made in the image of God as worthy to be respected..
So, if you can recall, similar to murder, Jesus said that when we show anger or contempt towards a brother or sister, we diminish the value and worth of someone made in the image of God..
And similar to adultery, when we lust after a brother or sister, we strip them of their humanity and turn them into an object for our own gratification..
And in my last sermon, we also look at how Jesus refrains kind of like the Jewish debate on divorce, and instead he shifts the focus on God's original design for marriage instead..
So, this morning, we're going to look at two more examples on how Jesus redefines what true righteousness means in contrast to the righteousness of the Pharisees and the teacher of the law..
So, if you have your Bible with you, let's turn with me to Matthew chapter 5. We are going to read from verse 33 to 42..
So, Matthew chapter 5, verse 33..
"Again, you have heard that it was said to the people long ago, 'Do not break oath, but fulfill to the Lord the vows you have made.' But I tell you, do not swear an oath at all, either by heaven, for it is God's throne, or by the earth, for it is his footstool, or by Jerusalem, for it is the city of the great king..
And do not swear by your head, for you cannot make even one hair white or black. All you need to say is simply 'yes' or 'no,' and anything beyond this comes from the evil one.".
Verse 38..
"You have heard that it was said, 'Eye for eye and tooth for tooth.' But I tell you, do not resist an evil person. If anyone slaps you on the right cheek, turn them the other cheek also. And if anyone wants to sue you and take your shirt, hand over your coat as well..
If anyone forces you to go one mile, go with them two miles. Give to the one who asks you, and do not turn away from the one who wants to borrow from you.".
So, this is the word of God. Let us pray..
Father, as we open up your words this morning, I just ask that your Spirit will be with us. Prepare our hearts as we receive your words, and allow your words to take root in our heart and transform our life, so that we can bring honor and glory to you through the way that we live..
So, pray this in your Son's name. Amen..
So, Matthew chapter 5, 33 and 242. What is it about?.
If you are taking notes, you can write this down..
"In a fragmented world, we are called to bring peace through our word and deed.".
So, this is the idea that we will be exploring today..
In a world where people are divided based on their skin colors, based on their economic status, based on their political ideologies, or beliefs in health or science, and so many other things that disunites us,.

$^{81}$we, as followers of Jesus, are called to bring peace through our word and deed..
So, we're going to first take a look at how we can bring peace through our words..
So, we bring peace through our words by speaking truthfully and with integrity..
We bring peace by speaking truthfully and with integrity..
So, verse 33. "Again, you have heard that it was said.".
So, once again, these examples begin with this same rabbinic formula..
"You have heard that it was said to the people long ago.".
And just to give a refresher, every time when a Jewish rabbi wants to start a discussion on a specific Mosaic law, they would introduce it using this formula..
So, Jesus, as the definite and ultimate interpreter of the law, Jesus is doing the same thing right here..
And the law that is in focus this time centers around oath-making..
So, "Do not break your oath, but fulfill to the Lord the vows you have made.".
This is not a direct quotation, but it is actually a summary of a series of law regarding oath from the Old Testament..
Nowadays, we would probably hear people say something like, "I swear to God," and then go on to give their statement, right?.
But oath in the ancient world, and even during Jesus' times, looks a little different than what we're used to now..
And they probably carry a little more weight as well..
So, let me give a little context to oath in the ancient world..
So, in the ancient Near East, people would invoke or call upon the name of their gods as a witness to affirm a statement..
And one of the ways that oath can give credibility to a promise or to a statement is by having some kind of curses or consequences being attached to it in the case of dishonesty..
So, for example, someone will make an oath and swear, "I swear that may Marduk or Baal or insert whatever God's name in there,.
may this God strike me dead or make me ill if I break this oath or if I do not follow along with this promise.".
And for the people of God, of course, they would invoke upon the name of the Lord or Yahweh when they make their oath..
And so this is what Jesus is pointing to here in verse 33 on how the Jews and the custom are when it comes to oath making..
And as you can probably expect at this point, Jesus then continues on to give his version and redefine what true righteousness means with this interpretation..
"But I tell you, do not swear an oath at all, either by heaven, for it is God's throne, or by the earth, for it is his footstool, or by Jerusalem, for it is the city of the great king..
And do not swear by your head, for you cannot make even one hair white or black.".
So let's pause here. What is going on in here?.
So during Jesus' time, the Jews, they would not mention the name of the Lord because out of reverence and out of respect for it..
So instead, what they would do is that they would use other things that are being associated with God when they make an oath..
So the heaven, the earth, or the city of Jerusalem are just three of the common realms that the Jews would call upon when they make the oath..
But what Jesus is saying right here is that, "Don't swear by the heaven or by the earth, for they belong to God the creator, not us as human..
And also, don't call upon Jerusalem, because it is God's city. It is not for us to call upon or for us to control or invoke.".
And just in case someone asks, "Okay, Jesus, what if I just swear on my own head? It is part of my body. I should have control over it, right?".
Well, Jesus says, "No, God has authority over us as well..
We do not have the power to even turn our hair white, which symbolizes old age and wisdom, or turn them black, which symbolizes youth and life.".
So, one thing to note here, when Jesus says, "Do not swear an oath at all,".
he is not nullifying old oaths because we find in other parts of the New Testament where God himself makes an oath in the books of Hebrews,.
and so does Paul in some of his epistles as well..
What Jesus is prohibiting here is the kind of superficial and flippant oath made on seemingly ordinary or insignificant issues,.
and sometimes even worse, blatantly false claims or conspiracies..
And what is wrong with all of these?.

$^{121}$At the root of these is a kind of attitude that thinks that we, as human, we have the authority to manipulate God for our own benefits,.
using God's reputation as a tool to lend credibility to our own statement or our own words..
So what should we do then? This is where verse 37 comes in..
And Jesus said, "All you need to say is simply yes or no, and anything beyond this comes from the evil one.".
So, rather than relying on oath, Jesus called us to speak truthfully in all circumstances,.
and anchor the truthfulness of our words in our personal integrity..
So let our yes be yes, and let our no be no. Live up to your own promises and stand by your own words..
Do not use oath as an excuse for not taking any responsibility for the words that comes out of our mouth,.
because these come from Satan..
And so here are just some of the applications that how we can speak truthfully with our words..
So first one, speak with integrity and do not use our words to slander other people..
Secondly, in a digital age, we communicate not just true words spoken, but also true words written or words typed up, right?.
Especially in the age of social media..
So, how we can bring peace with our words?.
Discern the information that you have heard or read online, and do not be quick to spread them around,.
or repost anything that you see online that has the potential to further divide us instead of uniting us..
This is because as followers of Jesus, we are called to bring peace through our words..
So, speak truthfully and with integrity..
And as the Apostle Paul instructed the church in Ephesus, do not let any unwholesome talk come out of your mouth,.
but only what is helpful for building others up according to their needs, that it may benefit those who listen..
So, this is how we can bring peace through our words..
Another way that we can bring peace is through our actions or through our deeds..
We can bring peace by acting selflessly and generously..
So, let's take a look at the next example from Jesus here..
Verse 38..
"You have heard that it was said, 'Eye for eye and tooth for tooth.'.
But I tell you, do not resist an evil person..
If anyone slaps you on the right cheek, turn them the other cheek also..
And if anyone wants to sue you and take your shirt, hand over your coat as well..
And if anyone wants to, and if anyone forces you to go one mile, go with them two miles..
Give to the one who asks you, and do not turn away from the one who wants to borrow from you.".
So, the principle right here, an eye for an eye or tooth for a tooth, is actually called the law of retribution..
And we actually find it in various places in the Old Testament..
And in the Old Testament, this principle of justice is not about personal retribution or personal vengeance,.
but it is actually about regulating excessive punishment..
In other words, it doesn't mean that if someone takes one of your eyes,.
now you have the legal right to take your eyes out..
No, the purpose of this principle is to regulate the penalty by saying that the punishment should fit the crime..
So, for example, if someone knocks a tooth out of your mouth,.
you can't just end your life and take your life by killing them..

$^{161}$Because the punishment, again, it doesn't fit the crime..
And by making this law in the legal system, it actually provides an appropriate avenue.
where people can pursue justice in a legal way instead of personal vengeance or personal revenge..
But once again, Jesus continues to say, "But I tell you, do not resist an evil person.".
I want to be careful here because someone can hear this and think,.
"Is Jesus just telling us to just sit there and ignore the need for justice?.
Or when we see evil being done, do we just sort of let it go and not resist it at all?".
So again, I want to emphasize that in this particular section,.
when Jesus talks about not resisting an evil person,.
Jesus is not rejecting the principle of justice at all,.
but he is speaking on retaliation and revenge at a personal level..
So what Jesus is not saying here is that we do not resist evil at all..
1 Peter 5 says that the devil is like a roaring lion looking for someone to devour..
And we, as God's people, we are called to resist him..
What Jesus is not saying here is also that we do not call out or oppose evil..
Time and time again in the Old Testament,.
we see that the prophets actually calls out the evil and the injustice that is happening in Israel.
and the surrounding nations as well..
And in the New Testament, we also see the apostle Paul actually calls out Peter.
because Peter is not acting in a way that is in line with the truth of the gospel..
And right here, Jesus is also not talking about justice in the legal system or in the government..
Romans 13 is pretty clear..
It shows us that God has given his authority to the government to execute justice on his behalf..
So right here, Jesus is calling us to let go of our right to get even or to repay violence with violence,.
but to depend on the ultimate justice of God..
And I know this is not easy at all because it requires patience on our behalf..
It requires us to wait on the Lord..
But not only that, we are also called to repay evil with good..
And Jesus actually gives a few illustrations here..
So the first one, "If anyone slaps you on the right cheek, turn to them the other cheek also.".
So throughout history, the right hand is usually the more dominant hand for most people, right?.
So to slap someone on the right cheek will be a backhand slap, pretty much like this way..
And in the first century context, a backhand slap is actually meant to be an insult to humiliate someone.
that usually has a lower social status compared to you,.
while a forehand slap is usually for a peer or someone who has equal status with you..
And so in this situation, Jesus tells his disciple to turn the other cheek,.
to take the high road and be humiliated rather than to take revenge..
And I actually like how one of the commentators put it..
He says, "When we turn the other cheek, we force them to treat us as equal.
because they will have to slap us with their forehand this time instead of the backhand.".

$^{201}$And in the second example, Jesus said, "And if anyone wants to sue you and take your shirt,.
hand over your coat as well.".
So shirt during Jesus' time is usually an inner garment and the coat is the outer garment..
And compared to the shirt, the coat, which is the outer garment, is usually the more important one.
because it is used not just as a protective layer from the elements,.
but it can also be used as a blanket or a makeshift shelter.
when someone needs to sleep in the field on the long road trips..
And in Exodus specifically, the law also states that even the poor people.
should have the right to keep their coat..
But in these examples that Jesus gave us here,.
the better righteousness would be to relinquish the more important possession,.
which is the coat, and to treat the injustice with costly generosity..
And in the third case, Jesus said, "If anyone forces you to go one mile, go with them two miles.".
And again, you need a little context to understand this..
So in the Greco-Roman world, a civilian is actually obligated by the law.
to carry the load of a Roman soldier for about one mile if they're being requested..
And for the Jews, they hated this practice..
But again, Jesus shows his disciples the better way by telling them to go the extra mile,.
even though they hated this practice..
Instead of getting stuck in the cycle of repaying evil with evil,.
Jesus calls us to respond with generosity..
Then Jesus sums up this section in verse 42,.
"Give to the one who asks you, and do not turn away from the one who wants to borrow from you.".
So all the things that Jesus asks from us here requires a kind of humility and selflessness, right?.
To not only release our right to get even with someone,.
but also to respond to the injustice or the evil with kindness and generosity..
If we really think about it, what Jesus is asking from us here is quite a tall task, right?.
Maybe some of you would say, "Jesus, you don't really understand the injustice that I have to endure,.
or you don't understand the things that are being done to me.".
One of the passages that has had the biggest impact in my personal life these past few years.
is from Hebrews chapter 4..
Talking about Jesus, the author of Hebrews say,.
"For we do not have a great high priest who is unable to empathize with our weakness,.
but we have one who has been tempted in every way, just as we are, yet he did not sin.".
So what Jesus has taught us here, he did not just say it,.
but he himself embodied this kind of humility and selflessness for us to follow..
In Jesus' arrest and trial, we read that he was being spat on, he was being struck,.
he was being falsely accused, and he was ultimately being humiliated on the cross..
And yet, one of his final words on the cross is this, and I want us to listen carefully..
This is what Jesus prayed to his father..

$^{241}$He said, "Father, forgive them, for they do not know what they are doing.".
So church, Jesus shows us the perfect examples of humility and selflessness,.
both through his words and his action,.
so that he's able to bring peace between God and human, as well as between human and human as well..
So if we, as the church, were able to exemplify this kind of costly generosity,.
I think the watching world would start to take notice of the way that we treat each other,.
and also the out-of-the-world kind of peace that we have..
And this, church, I think can be one of the most powerful witnesses to the kingdom of God.
in our fragmented society..
So, let's pray..
Father, we thank you for you have sent your son Jesus.
to bring peace and to restore the shalom that was being broken by sin..
And not only that, you have given us the gift of the Holy Spirit.
so that we're able to live in a way that follows after your examples..
So, Father, I don't know what are some of the things that are going through our minds currently..
Maybe some of the evils that we have experienced,.
or the things that are being said or done to us,.
or the injustice that we see in the world that just seems like no justice is being done..
So, Father, help us as we wait patiently for your ultimate justice when you come again..
Help us to be people that represent your kingdom,.
that represent what you have called us to be, to be the peacemaker..
For you promised that your kingdom belongs to people like this,.
and that we will be called the sons and the daughters of God..
So, help us to be people that bring peace to our families,.
to our relationships, to our workplace, to wherever you have called us,.
so that we're able to show the world who we belong to..
We're able to show the world who Jesus truly is, the Prince of Peace..
So, Father, help us to live in a way that brings glory to you this week..
Praise the Lord, Amen..
Okay..
\newpage



\section{約翰福音}
\label{sec:LHY953jFc_g}
\textbf{生命頻道 \#1 (約翰福音 10) - 陳崇基牧師}
\newline
\newline
連結: \href{https://youtube.com/watch?v=LHY953jFc_g}{\texttt{https://youtube.com/watch?v=LHY953jFc\_g}} ~~~~ 語音日期: 2025-03-06
\newline
\newline
\hyperref[sec:IH4xN4bYlR0]{\small{< < < PREV SERMON < < <}}
~
\hyperref[sec:index_chronic]{\small{[返順時目]}}
~
\hyperref[sec:index_scriptual]{\small{[返順卷目]}}
~
\hyperref[sec:KKxcZYP9oLY]{\small{> > > NEXT SERMON > > >}}
\newline
\newline
$^{1}$在一些弟兄姊妹的代討裡面.
都紀念到師母.
我太太回到加州.
有四五天時間.
因為幾個星期前.
她有一個親戚剛剛離開了世界.
她有一個舅父.
這個舅父他們剛剛來美國的時候.
幾十年前.
一直照顧他們.
所以我們都掙扎了一會.
到底他是不是需要回去.
他說當然要回去.
我就幫他訂了一張機票.
雖然貴一點.
不過我說這些事情一定是需要做的.
當然他自己回去.
不是帶我一起去.
他回去了.
星期二至星期六他回去了.
我每晚都出去買東西吃.
非常好.
昨晚他就安全的.
昨晚他安全回來了.
感恩.
也感恩聽他說了一些故事.
他們家族裡面越來越多人信主.
對嗎?.
所以值得晚上出去買東西吃.
OK.
今天我們希望能夠跟大家開始一個新的系列.
一些基本的東西.
我們基督徒有一些基本的東西.
是需要去掌握的.
特別是我們怎麼聽上帝的聲音.
這些絕對不是某一些人的特權.
不是說有些人比較屬靈.
他們屬靈指導.
或者他是師父.
或者做牧師等等.

$^{41}$他們才會聽到上帝的聲音.
上帝一向都和我們說話.
問題就是我們有沒有專心去聽.
設計這個slide就很容易.
因為說到頻度.
frequency.
你知道這個是我們UV的logo.
我們UV的logo就是Voice.
Voice當然是在說一些頻度的東西.
所以用一個這樣的比喻.
聖經裡面很多時候.
上帝是和我們說話的.
但問題是我們聽不聽到.
這個才是問題.
到底我們有沒有聆聽上主的聲音.
你記得我們這些比較老的.
我們還在用收音機.
有誰在聽收音機的?.
在車上.
有.
你記得以前收音機是不是要調校的?.
教台FM或者AM.
你要調校.
你調校得不正的時候.
就收到很多雜音.
甚至完全聽不到.
我們就好像一個生命裡面的一個頻度.
是需要調校的.
有時候世界裡面太多聲音.
太多噪音.
我們聽不見主導我們說的話.
其實上帝一向都和我們說話.
問題是我們的收音機有沒有調校得準.
我們的頻道是不是正確.
所以我抽取了一些經文.
講到這一方面.
是比較主題性的.
不過今天我們想去看一段經文.
是《約翰福音》第十章.
《約翰福音》第十章.

$^{81}$這整段經文裡面有幾節是很出名的.
其中一節就是十一節.
十章十一節.
耶穌說我是好牧人.
Good Shepherd.
好牧人為羊寫明.
誰是羊?.
你和我.
我們上帝的子民就是羊.
二十七節也很出名.
我的羊聽我的聲音.
我認識他們.
他們也跟從我.
這個就是牧羊人和羊之間的關係.
我認識我的羊.
羊要聽我的聲音.
跟我走.
整個《約翰福音》第十章.
我們可以叫做Good Shepherd Message.
或者Good Shepherd Discourse.
就是好牧人的講道.
《約翰福音》從頭到尾.
收集了很多.
記錄了很多耶穌的講道.
我們叫他什麼什麼Discourse.
什麼什麼Speech.
什麼什麼Message.
Sermon.
搜集了耶穌很多次的神蹟和講道.
有很多次都是公開的.
並且很多次是和當時猶太人的節慶.
和他們的體制有關.
這些東西.
是耶穌特別在那段時間.
選擇那一天.
選擇那一天是好日子.
他去做的.
譬如安息日.
御節.
駐棒節.

$^{121}$婚禮等等.
這些我們都記得.
是不是《約翰福音》.
到第十章的時候.
耶穌就做一個最後的.
所謂最後一場的公開的講道.
就是和他反對的人講.
也都和他的所謂神的子民.
來到他講道.
一個Open的.
公開的一段時間.
這段時間.
講這個好牧人的講道.
就是耶穌說我就是好牧人.
我的羊是跟我的.
是聽我的聲音的.
這段時間.
是章12節都有交代.
這段時間是趁著什麼好日子來講的.
他說那時是冬天.
22節.
在耶路撒冷有所謂修奠節的時候.
或者獻奠節的時候.
獻奠節或者光明節.
或者修奠節.
或者聖誕節.
或者是聖誕節.
大家知道聖誕節是什麼嗎.
在我們附近北面.
Skokie附近的地方.
有很多猶太人聚居的地方.
在全世界.
Skokie是最多.
在美國裡面.
其中一個數一數二最多的.
猶太人的Holocaust.
大屠殺之後的生還者住的地方.
所以Skokie是一個很大的.
猶太人的社區在那裡.
每逢到大概年尾12月.

$^{161}$就是聖誕節左右的時間.
他們就會慶祝所謂的Hanukkah.
Hanukkah他們會點燭光.
這個不是聖所裡面的七燈台.
你看到有多少支.
是有八支的.
中間那支是用來點燭光的.
你看到有九支.
其實是八支.
他們就是.
這個Hanukkah裡面有八天.
每一天點一支的陽燭.
為什麼呢.
有一個歷史的背景.
一講到這個獻殮節.
或者修奠節的時候.
就是很多年前.
耶穌之前的大概百多年之前.
那些猶太人.
他們巴西那裡就被人欺負.
當時欺負他們那個就是希臘帝國.
或者西樓古.
一直欺負他們一直去迫害他們.
並且又和當地的.
或者當時的一些猶太人的民間的不義的領袖.
來做一些勾當.
串謀他們來污蔑聖殿.
然後去殺害那些祭司等等.
所以是一段很慘的時間.
在生死存亡之間.
大概主前164年的時候.
你知道耶穌的時候大概主後30年.
你看這個時間線再回望.
大概200年左右.
164年BC的時候.
發生什麼事呢.
就是一個家族.
猶太人的家族.
是一個敬虔的猶太人的家族.
叫做瑪加比.

$^{201}$他們就開始搞起義.
他們說要去反對希臘的帝國.
然後這個家族就已經擊退了希臘的帝國.
然後他們就成立了他們自己的王朝.
獨立了大概有100年的時間.
這個大概是主前164年的時候.
他們就拿回這個聖殿.
當他們拿回聖殿的時候.
大家記得嗎.
有一個傳說.
他們進去聖殿的時候.
看到一件事都搞砸了.
剩下一個瓶油.
那個瓶油明明是不夠他們點一個星期的.
誰知道點了8天.
然後給他們8天的時間.
他們再提煉一些鑑油出來.
給聖所裡面的七燈台用.
所以就有一個這樣的傳記的故事.
那個金燈台神跡地就是燒了有8天.
只是用一瓶而已.
所以到今天他們就會點8支的陽燭.
所以在那個100年的獨立.
所謂光輝歲月的時候.
但是到最後的時候.
就是瑪加比的家族.
那些後代都是狗咬狗骨的.
然後羅馬帝國就有機會入侵.
然後拿回攻佔了這個聖殿.
又是搞到聖殿一直都是被污蔑.
每逢到這個修煉節的時候.
他就會講這些故事.
所以《約翰福音》第十章就是一個這樣的背景.
在修煉節的時候.
民間就會問一些很尖銳的問題.
歷史的問題.
他們的歷史.
他們近代的歷史.
或者他們的黑歷史.
講回為什麼我們現在會搞成這樣.

$^{241}$還是被羅馬人欺負我們.
以前希臘帝國也是被人欺負.
然後我們有一個王族出來.
這個瑪加比王朝.
能夠將他們擊退.
但是為什麼我們最後又失敗.
就是因為什麼.
就是因為領導的問題.
因為我們那些王.
不是真正的牧人.
不是真正去牧羊群羊的人.
你知道猶太人.
他們每個星期會跟一個讀經的計劃.
大概三年之後.
他們就會整本聖經讀完.
所以他們的會堂.
安息日的時候.
他們就會讀一段的聖經.
很大一段的.
譬如每逢到九月十月的時候.
他們又重新來過.
然後就創世記開始這樣讀.
到這個修煉節的時候.
大概是西歷的年尾.
他們大概會去到.
看《律法五經》裡面其中一些經文.
也會看一些跟牧羊人有關的.
譬如《以色列書第34章》.
《以色列書第34章》是去到指責當時的.
以色列的所謂領袖.
他們不是真是去到牧羊群羊的.
所以為什麼搞到.
整個國家被摧毀.
因為那些領袖有問題.
《以色列書第34章》第一節.
當時他們修煉節的時候.
大概就會讀這段經文.
這段經文說什麼呢.
他說臨到我禍災.
以色列的牧人.

$^{281}$只知牧羊自己.
吃肥油穿羊毛.
和宰殺肥羊.
大家有沒有去吃火鍋.
他們就是這樣.
天天去吃火鍋.
將羊殺了.
然後去餵飽自己.
汗牛必與牧人為敵.
我必救我的羊脫離他們的口.
不再作他們的食物.
我必在他們之上立一個新的牧人.
這個牧人就好像獨人大衛那樣.
就會牧養他們.
成為他們真正的牧人.
每逢修煉節的時候.
他們就想起這些黑歷史.
就是我們國家裡面的元首.
不是真正的牧人.
現在怎麼辦呢.
到底什麼時候來呢.
只是這樣的背景.
耶穌就講了一篇這樣的道.
叫做Good Shepherd Discourse.
Good Shepherd Sermon.
有這樣的背景的時候.
我們就會知道.
到底祂想講什麼.
修煉節的時候.
耶穌說公開的講.
祂說我實實在在告訴你們.
那不從門進羊圈的.
倒從別處爬進來的.
就是賊.
就是強盜.
那從門進去的才是羊的牧人.
開門的給他開門.
羊也聽他的聲音.
他按著名叫自己的羊.
把羊領出來.

$^{321}$當他把自己的羊都放出去.
放出來.
就走在前面.
牧羊人當然走在前面.
羊也跟著他.
因為他們認得他的聲音.
羊絕不跟陌生人.
反而會逃走.
因為不認得陌生人的聲音.
我就是羊的門.
耶穌說我就是.
the gate of the sheepfold.
第十節.
盜賊來無非是要偷竊.
殺害.
毀壞.
但耶穌說我來了.
是叫羊得生命.
並且得的更豐盛.
我是好牧人.
好牧人為羊寫明.
故宮不是牧人.
羊不是他自己的.
他一看見豺狼來.
就撇下羊群逃跑.
狼抓住羊把他們趕散.
故宮逃走.
因為他是故宮.
對羊毫不關心.
但我是好牧人.
我是the good shepherd.
我認識我的羊.
我的羊也認識我.
這段這麼長的講道.
我們當然不能看完.
但大概你get到主題是什麼.
主題是跟牧羊有關的.
這裡要看看地圖.
在巴勒斯坦地.
你見到中間那裡有幾點.

$^{361}$有沙灘.
然後南面有伯利恆.
再南一些是希伯倫的山地.
那個叫中央山地.
Central Hill Country.
好像一個脊椎那樣.
拱了上去.
那個是最高的地方.
由地中海左邊.
到山峰上面.
山脊上面.
就是這些大小的城鎮.
而旁邊一過山的時候.
好像一過山一下去的時候.
那裡你見到是啡色的地方.
那個叫做曠野.
或者猶大曠野.
這是一個天然的地理.
因為曠野那些雨水.
在地中海吹過來的時候.
吹不到過去曠野.
或者所有的雨水都落在中央山脊的左邊.
然後右邊就沒有水了.
所以曠野長年是沒有雨水的.
好像詩篇所講.
是一些乾旱無水之地.
山脊上面這些大小的城市.
那時候那裡居住的人.
比如伯利恆.
耶路撒冷.
希伯倫等等.
他們就是用這個地理的環境.
來去放羊.
去牧羊.
如果你是住在那裡的話.
比如伯利恆.
好像以前大衛住在伯利恆那裡.
他們放羊的時候.
就走去哪裡放羊呢.
就走去曠野那個地方.

$^{401}$因為曠野那個地方沒有草.
怎麼放羊呢.
死定了.
但是曠野很有趣的.
偶爾有驟雨經過.
驟雨一經過的時候.
那個曠野就像地氈一樣綠色.
有一點點時間.
如果那個牧羊人是有智慧的話.
有能力的話.
他能夠有預知的.
看天氣的智慧經驗的話.
他就會帶羊.
趁著這個時間就去放牧.
一有驟雨之後.
就有些野草出來.
然後羊就可以有東西吃.
這些就是曠野他們會利用的地方.
你就想起CP23篇等等不同的經文背景.
但是同樣你想起wilderness.
想起曠野的時候.
就是荒山野嶺的地方.
最多是什麼出沒.
就是野獸的地方.
野獸出沒的地方.
還有盜賊.
山賊出現的地方.
山賊就是躲在這些山洞裡面.
然後就趁著機會.
然後就去殺害一些人.
譬如耶穌說一些比喻.
很撒馬利人的比喻.
從耶路撒冷到耶利哥.
下去下面.
經過一個地方.
然後那個人就遇見有山賊.
然後打他半死.
這個就是耶穌很多時候.
說到的一些地理的背景.
但是曠野的地方.

$^{441}$當這些牧羊人來放羊的時候.
很多時候你就會在曠野.
找到一些這樣的景色.
就是一個山洞.
半個山洞.
他就將羊放進去裡面.
有時候過夜.
有時候久了的時候.
他就把羊放進山洞裡.
然後他就將石頭堆砌在前面.
然後就擋住這個山洞.
讓羊不會自己跑掉.
然後剩下下面有一個小小的門口.
讓羊可以出入.
晚上的時候.
誰在門口.
就是牧羊人.
牧羊人就是那道門.
牧羊人就是守著那道門.
所以你明白耶穌說.
我是羊的門.
其實他是在說一些實際的意境.
你直接可以翻譯.
我是羊的門.
其實就是我是好牧人的意思.
因為真正的牧人就會守著這裡.
這個就是第一樣.
這個港道裡面所說的角色.
就是耶穌才是那個好的牧人.
聽祂的話就會有安全.
跟著祂的話.
懂聽祂的聲音的話.
跟著祂的聲音的話.
認得出祂的話.
你就有永生了.
第二個角色就是強盜.
或者盜賊.
這裡經常提到有盜賊.
做什麼.
他們是不經過門口進來的.

$^{481}$就是沒有牧人的同意.
他們爬進去.
爬進去裡面偷羊.
或者爬去羊灣等等.
這裡.
那個港道.
其實就是耶穌.
(音量注意).
因為這個公開港道.
如果你看清楚第十章.
其實是很激烈的港道.
其實他講這個港道就是.
想針對當時的猶太人的領袖.
當時猶太人的領袖.
他們守著那個聖殿.
但是他們保障的是他們自己的利益.
不是人民的利益.
因為當時他們是.
跟那個羅馬的政府是串通的.
鞏固自己的勢力和地位.
這群就是所謂牧羊人.
所以耶穌講這個Good Shepherd的.
雙方是誰.
他們是誰.
他們是盜賊.
你們是強盜.
因為你們是為利是圖.
你們出賣公義.
你們是不罔顧人的生死.
出賣人民的領袖.
但是很奇怪.
也有人追捧的.
這個非常之可惜.
但是它是事實.
你想一下.
不是當時是這樣.
現在也是這樣.
在第十章.
他就是針對這些領袖來說這些話.
他說你們才是盜賊.

$^{521}$而我呢.
就是好牧人.
我就是守著那個門.
到第十一章.
下一章你就可以看到.
到底拆穿了什麼真相.
耶穌所指責的那群宗教的領袖.
聖殿的領袖.
就是什麼人.
第十一章第47章.
這裡很有趣的對話.
就是那些宗教的領袖.
祭司長等等.
他們開會.
他說耶穌這麼多人跟他.
我們還得了嗎.
他們開會的談話裡面說什麼呢.
祭司長和法利賽人就召開會議了.
他們說這人就是耶穌.
行了很多的神蹟.
我們怎樣扮他呢.
讓他這樣做.
人人都要信他.
個個都要跟他.
如果個個都跟他.
好不好事.
很嚴重.
為什麼.
因為我們會知道.
羅馬人就會來摧毀我們的聖殿.
和我們的民族.
其中有一個人就是戈爾亞法.
就是審耶穌的大祭司.
當年的大祭司.
就對整群工會的人說什麼.
他說你們不知道.
也不想一想.
用一下腦想一想.
一個人如果替百姓死.
免得整個民族滅亡.

$^{561}$這是對你們有利的.
我們殺了耶穌.
是對你們有利的.
如果將祂剷除的話.
我們就可以保住我們的聖殿和民族.
如果民間裡面太多人跟祂的話.
搞起亂的時候.
羅馬兵就會來摧毀我們的一切.
你可以想像一下.
以前的世界.
耶穌當時的社會.
就是一種曠野.
就是dog eat dog的一個世界.
每個人都是為你仕途的.
到處可以看到山賊和一些土匪.
以前耶穌的時代.
耶穌當時的處境.
好像曠野一樣.
今天其實我們都是一樣.
有人說我們現在就是.
在一個modern desert裡面.
modern wilderness.
我們今天都是.
你不需要太多思考.
留意一下新聞.
你就會發現.
國際關係也好.
在政界也好.
在社會也好.
都是全部為你仕途的人.
是不是.
如果是大國.
如果是有霸權的話.
我們不顧一些小的民族的生死.
可以恃強凌弱.
是dog eat dog的世界.
我們的社區也是一樣.
不久之前.
不知道現在還有沒有.
應該還有.

$^{601}$什麼迷藥黨那些.
什麼神醫黨那些.
這些不是九吃九的世界.
個個都是為你仕途.
個個都是想謀財.
或者甚至害命.
這個就是我們現在的世界.
現代的強度都是機會主義.
在我們這個世界裡面.
好像一個曠野一樣.
耶穌就是指責當時猶太人的領袖.
說他們不是真正保住人民的利益.
是保住自己的荷包.
保住自己的地位.
怕自己失去了聖賢.
耶穌要來.
就是審判這幫人.
並且要取締他們.
所以聖賢要被毀.
真正上下的子民.
不再在聖殿敬拜上.
在基督裡面敬拜上.
所以一個fulfillment.
或者甚至一個replacement.
耶穌說不再需要聖殿.
是需要我.
因為我為我的百姓死而復活.
是不是.
講了牧羊人.
講了盜賊.
主就說.
到底我的羊是怎樣的.
第十章第27節.
所謂第三個角色.
他說我的羊是怎樣的.
如果你是屬於主的羊的話.
就會聽主的聲音.
我的羊人聽我的聲音.
我認識他們.
他們也跟從我.

$^{641}$並且我賜他們永生.
他們永不滅亡.
誰也不能從我手裡把他們奪去.
牧羊人和羊之間有一個很獨特的關係.
是靠聲音.
是靠call it by name.
羊聽見牧羊人的聲音.
牠就會跟從.
認得出這個聲音.
到底是屬於牧羊人.
還是強盜.
還是陌生人.
我記得好幾年前.
去以色列.
有一次去以色列.
我們一群同學.
我們就去一個以色列的植物公園.
我們就在植物公園裡面去遊房.
然後去看看聖經的背景等等.
你知道聖經裡面說很多植物和很多動物.
我們就去看看聖經的背景.
我記得最後.
臨日落之前.
他就帶我們去一個羊欄.
去羊欄那裡.
他的導遊跟我們說.
不如我們玩個遊戲.
我們那裡有幾十人.
我們玩個遊戲.
羊欄玩什麼遊戲呢.
羊欄裡面有一支旗.
一支紅色一支藍色.
我們分組分兩組.
裡面有很多大大小小的羊.
我們就分組.
A,B組.
我們就說要趕這些羊.
去我們那個.
Capture the Flag.
這樣的意思.

$^{681}$就是去我們那個旗那裡.
他就給了時間給我們.
我們就花了差不多一個小時的時間.
我們就盡力去做一個比賽.
然後趕羊去我們那個旗那裡.
然後大混亂了差不多一個小時.
我們用了很多的方法.
威逼利誘.
甚至我們抱起那隻大羊.
搬牠過去那裡.
扔牠過去.
扔牠過去那裡.
牠又說坐不住.
然後又跑去別的地方.
好像那些小孩一樣.
我們角力了大概一個小時.
最後的時候.
那個公園差不多要收工了.
他說我們要收工了.
那怎麼辦.
那個園主就跑過來.
他叫了幾聲.
然後羊就全部排隊跟著他走.
很有趣.
就是聽見那個牧羊人的聲音.
我們經常有個感覺.
那些羊就是很笨的.
是吧.
當然笨了.
肥羊我們吃的.
但是他們最大的Defense是什麼呢.
生存之道.
就是認得牧羊人的聲音.
在YouTube上有一個兩分鐘的影片.
不知道大家有沒有看過.
是表明這個事實的.
這個影片是說什麼的.
兩分鐘而已.
就是有一群羊在羊欄那裡.
就是開始有幾個女生.

$^{721}$學生來的.
就去叫那些羊.
那些羊就不理牠.
到最後.
牧羊人來了.
你看羊有什麼反應.
我們播放這個影片還有兩分鐘.
(羊群).
(Country music).
(剪刀石頭布).
再一次.
(Country music).
看看她….
(Country music).
她在哪裡.
她在那裡.
她在那裡….
我沒看過她的表演.
她的表演很厲害.
很厲害… 真是太棒了.
再看一次… 太棒了….
那是好看嗎 很好看….
我們永遠不會做同樣的表演.
太棒了….
好像很困惑,其中最大的秘密.
就是認不出羊人的聲音.
我們就是主的羊.
是否聽到牠的聲音.
還是我們聽到其他聲音.
在基督徒的生命中.
我們需要培養.
我們要學會聽上帝的話.
要培養這種習慣.
我想問你們每天有沒有這樣的空間.
每天有一點空間.
讓你們學習習慣聽上帝的話.
有一樣東西很可悲.
就是我們基督徒.
有一種即食消費的文化.
消費的心態.

$^{761}$譬如我想聽上帝的聲音.
我找牧師.
牧師就坐下來.
牧師上帝要跟我說什麼.
你快點說.
不要阻礙我.
你要給我東西.
因為你的責任是什麼.
你的責任就是聽上帝的聲音.
然後告訴我.
當我們有大事.
或者我們要做一些重要的決定的時候.
我們就希望上帝快點出來回應我們.
這有點像求籤.
最近我開車在芝加哥.
我留意到有好幾間.
那些看獎的店舖.
有沒有見過.
Palm Reading.
這麼現代的社會.
居然有人還在看獎.
Palm Reading.
這樣也做好生意.
這就是現代人我們有的心態.
就是我們希望求籤去解籤.
我們覺得上帝是on demand.
guard on demand.
隨嬌隨都的上帝.
這個才叫上帝.
這個是非常可悲的.
一個福音派的作家.
他做了很多這種的.
屬靈神學.
Spiritual Theology.
一個叫Dallas Wheeler.
Dallas Wheeler說到.
我們基督徒和上帝之間的關係.
是怎樣的呢.
他說令人遺憾的是什麼呢.
很少有人會真心來求聆聽神對他說話的.

$^{801}$這個說明了什麼呢.
當我們沒有預見逆境.
或者面臨重要抉擇的時候.
往往我們是不會主動來尋求主的聲音.
但是那些真正去理解.
想去渴望聽見神聲音的.
無論是或順或逆.
都同樣渴求上帝的話語.
上帝的說話.
因此當我們尋求神的話語的時候.
應該反問自己.
我們是不是只在困難的時候.
才聽神的話呢.
這樣的反思.
或者會揭示到一件事.
就是當我們聽不到神的聲音.
是因為我們平時都不太想聽祂的聲音.
只是我們覺得需要的時候.
我們才試圖去求祂的聲音.
耶穌說.
我的羊聽我的聲音.
我們有沒有好好學習.
學習聽上帝的聲音.
今天我想給大家一個邀請.
就是我們這個星期.
請你在每一天最黃金的時間.
就是最好的時間.
有些人可能是早上.
有些人可能是中午.
有些人可能是晚上.
你花十分鐘.
在上主的面前.
試下安靜.
有時候我們心靈是需要打掃一下.
有很多東西困擾我們.
有很多問題.
有很多雜音.
有很多嘈音.
是需要打掃一下.
然後花十分鐘.

$^{841}$我希望你看.
《約翰福音第十章》.
好嗎.
這個星期我們要學習去聽上主的聲音.
我們跟心得了祈禱.
預備我們的心.
來到去領受聖餐.
上帝我們多謝你.
因為你是那位好的末人.
我們的主耶穌.
就是那位Good Shepherd.
主我們求你幫助我們聽見你的聲音.
在我們心靈裡面.
無論我們現在是在經歷一些難處.
或者在順利的時候.
主要教我們同樣地去渴求聽你的聲音.
主在這個時候.
我們的心靈需要你去打掃.
將我們罪拿開.
將我們很多的雜音拿走.
叫我們來聽見你的聲音.
叫我們來渴求聽你的聲音.
藉著你的話語.
藉著你的指紋.
藉著生活裡面不同的事.
你去到讓我們聽見主你的慈聲.
帶領我們.
我們作為你的羊.
叫我們的心靈更加奉承.
為我們捨命.
死在十架上.
去救贖我們.
主我們今天要在聖餐餅和杯.
也紀念到你對我們的恩典.
求你在我們心靈當中.
幫助我們有更多空間去聽你.
去聽從你的話語.
聽從你的話.
聽你的聲音.
我們在這裡祈禱奉祖明.

\newpage



\section{馬可福音 14:32-36}
\label{sec:KKxcZYP9oLY}
\textbf{新約中的情緒 (2): 恐懼 (馬可福音 14\_32-36) - 何建邦傳道}
\newline
\newline
連結: \href{https://youtube.com/watch?v=KKxcZYP9oLY}{\texttt{https://youtube.com/watch?v=KKxcZYP9oLY}} ~~~~ 語音日期: 2025-03-11
\newline
\newline
\hyperref[sec:LHY953jFc_g]{\small{< < < PREV SERMON < < <}}
~
\hyperref[sec:index_chronic]{\small{[返順時目]}}
~
\hyperref[sec:index_scriptual]{\small{[返順卷目]}}
~
\hyperref[sec:gEwUekxM30c]{\small{> > > NEXT SERMON > > >}}
\newline
\newline
馬可福音 14:32-36
\newline
\begin{longtable}{cl}
\hline
\hline
章節 & 經文 (和合本修訂版)\\
\hline
14:32 & \begin{tabularx}{0.7\textwidth}{X} 他們來到一個地方，名叫客西馬尼。耶穌對門徒說：「你們坐在這裡，我去禱告。」 \end{tabularx} \\ \\ \relax
14:33 & \begin{tabularx}{0.7\textwidth}{X} 於是他帶著彼得、雅各和約翰同去。他驚恐起來，極其難過， \end{tabularx} \\ \\ \relax
14:34 & \begin{tabularx}{0.7\textwidth}{X} 對他們說：「我心裡非常憂傷，幾乎要死；你們留在這裡，要警醒。」 \end{tabularx} \\ \\ \relax
14:35 & \begin{tabularx}{0.7\textwidth}{X} 他就稍往前走，俯伏在地，禱告說，如果可能，就叫那時候離開他。 \end{tabularx} \\ \\ \relax
14:36 & \begin{tabularx}{0.7\textwidth}{X} 他說：「阿爸，父啊！在你凡事都能；求你將這杯撤去。然而，不是照我所願的，而是照你所願的。」 \end{tabularx} \\ \\
[1ex]
\hline
\hline
\end{longtable}
$^{1}$(新約中的情緒).
大兄姐妹,早安.
今天我繼續我的新約情緒系列.
大家還記得嗎?應該是數個月前.
數個月前,我開始了新約中的情緒系列.
我第一次說的情緒是….
記得是甚麼嗎?是憤怒.
不記得的大兄姐妹.
也可以上YouTube看.
「新約中的情緒1」,憤怒.
今天我說另一個情緒.
今天我要說的情緒是甚麼.
就是「新約中的情緒2」,恐懼.
我們首先要問一個問題.
究竟恐懼是甚麼.
我們…跟憤怒一樣.
我們每個人都一定會經歷過恐懼.
我想問有誰未經歷過恐懼的?舉手.
未經歷過,你從來都不害怕.
不害怕,我很想知道.
有人不害怕嗎.
好,我想問另一個問題.
有人覺得自己經常都害怕嗎.
對,有很多人經常都害怕.
我們經常經歷的情緒究竟是甚麼.
所以我們首先….
今天我會看「新約中的情緒」.
所以會看比較多經文和資料.
但我相信大家一定能跟隨.
所以不要緊,我們一起看看.
我們首先問這個定義的問題.
最簡單的,我們去到APA.
美國心理學會,我曾說過.
最標準的定義是甚麼.
那裡是這樣說的.
恐懼是甚麼.
恐懼是由於察覺到迫在眉睫的危機.
而引起的一種基本強烈的情緒.
並且涉及一系列的生理反應.
例如心跳加速.

$^{41}$是嗎?大家都明白.
例如你緊張,很害怕時.
心跳會很快,對吧.
或者有時會氣喘,不能呼吸.
有這些生理反應,對吧.
我們都明白了.
這就是APA的定義.
我們看看另一個情緒專家.
我上次也說過.
就是迪士尼的卡通片.
「Inside Out」的心理學顧問.
Paul Edmund是這樣定義恐懼的.
他說恐懼是七種….
跟憤怒一樣.
都是七種基本情緒之一.
因為察覺到或想像到.
可能受到身體或心理上的傷害而產生.
恐懼不僅是負面的.
也有正面的作用.
能夠幫助我們.
促使我們應對潛在的危機.
所以大家留意到.
Edmund對於恐懼的定義是怎樣.
就是你害怕一些東西.
那件事不一定是真的,對吧.
可以是甚麼?可以是真.
即是你正面對一件迫在眉睫的事件.
是真實的事件.
亦可以是怎樣.
是想像出來也可以的.
當你想像到一些東西.
你覺得是很危險的事.
你想像到,你也會害怕.
所以恐懼可以面對一些真實的情況.
亦可能那情況是假的.
是你想像出來的.
但恐懼不僅是負面.
也可以幫助我們去想.
當我們想像到一些危險的事件時.
會怎樣.

$^{81}$就促使我們去想辦法去面對它.
所以恐懼是有正面的作用的.
你想想,一個人.
一個完全不害怕的人.
其實會怎樣.
即是很恐怖.
即是有危險,他也不害怕.
那會怎樣?死了也不知道發生甚麼事.
所以其實恐懼是我們人類的進化中.
在幫助我們.
即是我們做原始人時.
其實恐懼是一個很重要的情緒.
是進化出來的.
是一種生物進化的本能.
我們看見獅子,老虎時.
我們會害怕,對吧.
我們想像到有獅子,老虎在我們前面.
我們應該害怕.
如果不害怕,那又怎樣.
就會被它吃掉,對吧.
所以恐懼本身是有正面的作用的.
所以這就是Paul Edmund.
對於恐懼的定義.
但今天我們是說新約中的情緒.
所以說完…跟上次一樣.
說完現代定義,然後我們要怎樣.
我們就要看古代的定義.
又是這位人物,亞里士多德.
我通常都說他的情緒.
因為他說很多情緒.
亞里士多德.
這個古典的希臘哲學家.
就這樣定義恐懼.
他就說恐懼是甚麼.
恐懼是因為未來可能發生.
以及會帶來痛苦的壞事.
而產生的痛苦或不安.
他就這樣說.
不是所有壞事都會帶來恐懼.
那些不會帶來痛苦.

$^{121}$或者距離你很遙遠的事情.
即使是壞事,也不會帶來恐懼.
所以亞里士多德覺得恐懼是甚麼.
恐懼是一種價值判斷.
我重複多一次.
恐懼是一種價值判斷.
為何我會說是價值判斷.
就是因為你在判斷那件事.
是好事還是壞事.
這種是價值判斷.
當你判斷那件事是壞事時.
你才會害怕,對吧.
所以為何不是每件壞事都會害怕.
因為那件壞事不是發生在你身上.
或者你覺得那件壞事與你無關.
或者你覺得那件壞事發生在他身上.
而永遠不會發生在你身上.
即使是壞事,你會害怕嗎.
你不會害怕,因為與你無關,對吧.
所以就算你判斷壞事.
也是要跟你有關.
是很接近你的,你才會害怕.
遙遠的壞事,你也不會害怕.
所以恐懼是一種價值判斷.
當你判斷一件事與你有關的壞事時.
你就會害怕,好嗎.
李氏多德還說了幾件關於恐懼的事.
他除了說恐懼是一種價值判斷.
他還說恐懼與權力有關.
為何恐懼與權力有關.
他就這麼說.
有些人能夠使其他人害怕.
為何人有令其他人害怕的力量.
是因為他擁有傷害那些人的力量.
所以恐懼與權力的結構有關.
有能力的人會令人害怕.
他能夠傷害別人,所以別人會害怕他.
我們害怕一些人.
是因為他們可以傷害我們.
所以我們會害怕那些人.

$^{161}$所以這與權力有關.
而且恐懼也與勇敢有關.
我想問一個問題.
你先別看李氏多德的定義.
誰覺得勇敢的人不會害怕.
有人覺得勇敢的人不會害怕.
李氏多德就這麼說.
他說勇敢的人為了崇高的目標而堅定不移.
做需要勇氣去做的事.
但他不是不害怕.
相反,對任何事情都不害怕的人.
不是傻的,就是對痛苦無感.
為甚麼.
傻子才會完全不害怕.
其實勇敢的人也會害怕.
勇敢不是代表你不害怕.
勇敢是甚麼意思.
是你害怕,但你仍然面對你的害怕.
去做那件事,這才叫勇敢.
完全不害怕的人是傻的.
是白癡的,而不是勇敢.
這是李氏多德的定義.
我們又說到….
李氏多德是古典希臘時代.
我們看看耶穌的時代.
因為我們回到新約.
耶穌時代的哲學家.
或者希臘羅馬的哲學家.
當時的文化的人是怎樣看待恐懼的.
當時有一個很流行的哲學.
上次也有提過,在討論憤怒時.
斯多亞主義.
斯多亞主義定義了四種基本情緒.
你留意,古代人看情緒.
跟現代人的心理學有點不同.
我現在說古代人怎樣看待情緒.
斯多亞主義定義了四種基本情緒.
這群哲學家認為.
當人現在經歷一些好的事情時.
現在經歷好的事情.

$^{201}$那種情緒就叫歡樂.
當你現在經歷一些壞事時.
那種情緒就是愛上.
當你預期自己將來會經歷一些好的事情.
那種情緒就叫desire,渴望.
情緒跟我們現代有點不同.
當人預期自己將來會經歷壞事時.
那種情緒就是恐懼.
所以恐懼的定義跟亞里士多德有點相似.
都是涉及價值判斷.
斯多亞主義哲學認為恐懼是價值判斷.
就是你判斷那件事將會發生在你身上.
而那件事是壞事時.
你就會恐懼,害怕.
好嗎.
但斯多亞哲學家認為.
在世界上只有一件好事.
是甚麼.
唯一的好事情或唯一的好事物.
就是美德,道德,Virtual.
所以對斯多亞哲學家而言.
只有追求美德才是人應該做的事情.
世上任何事情都是無分別的.
所以會怎樣.
帶來的結果就是.
任何人都不應該渴望財富.
因為財富不是好事.
唯一的好事物是美德.
所以人不應該渴望財富.
而去恐懼貧窮.
因為貧窮也不是壞事.
唯一的好事物是美德.
所以無論你是有錢還是沒錢.
都是無分別的.
他們亦覺得人不應該渴望生存.
而去恐懼死亡.
因為生和死都是無分別的.
唯一的好事物就是美德.
所以無論你是有錢還是沒錢.
生還是死.

$^{241}$他們都覺得你需要無感.
這班斯多拉哲學家覺得.
你只需要去追求美德.
好,繼續.
這位就是普塔.
他這樣說.
恐懼其實跟宗教有關.
他說甚麼?.
他這樣說,普塔.
這位中奇柏拉圖主義的哲學家.
他說迷信的人.
會把仁慈的眾神當成令人恐懼.
把眾神慈父般的關懷當成專制.
把眾神的慈愛當成有害.
迷信的人恐懼眾神.
所以他想說甚麼.
他想說真正的信仰不應該懼怕.
為何人們經常害怕神.
人們經常害怕神.
人們恐懼眾神是因為他們迷信.
所以普塔認為恐懼跟迷信有關.
真正的信仰不應該對神有恐懼.
這就是他的看法.
所以我們簡單歸納一下.
在耶穌的時代.
對當時的希臘,羅馬的哲學家來說.
恐懼是怎樣.
第一,恐懼沒有任何正面的功能.
剛才我說過,全部都是負面的.
事實上,在當時耶穌的時代的.
希臘,羅馬的哲學.
有三種關於恐懼.
在哲學著作中經常提及.
三種東西.
第一,人不應該恐懼貧窮.
剛才我說過.
第二,人不應該恐懼死亡.
第三,人不應該恐懼神明.
就是這三種東西.
我們又去看看.

$^{281}$在新約中.
有沒有反映這三個關於恐懼.
在哲學著作中經常出現的主題.
很明顯,當然有.
否則我就不會提出.
我們就看看新約.
首先,新約的一些經文都在強調.
或者提醒我們.
我們不需要恐懼貧窮.
我們不應該恐懼貧窮.
例如在《盧伽福音》或《馬太福音》中.
耶穌也說過.
就是對門徒說.
「所以我告訴你們.
不要為生命憂慮吃甚麼.
也不要為身體憂慮穿甚麼」等等.
還有「我們不要怕.
要變賣我們的財物,施捨給人.
因為我們的財寶在哪裡.
我們的心在哪裡」等等.
這些經文.
其實在新約的聖經中.
包括耶穌自己的教導.
都在說我們應該追求的是甚麼.
是神的角度,對吧.
所以貧窮不是我們需要憂慮或害怕的.
因為只要我們追求神的角度.
其實神正在供應給我們.
這是耶穌在新約的教導.
所以這跟當時的哲學.
不要害怕貧窮有些吻合.
第二,就是不要恐懼死亡.
在新約聖經中,也有很多經文提到.
叫我們不用害怕死亡.
例如耶穌,又是耶穌.
耶穌在《盧伽福音》和《馬太福音》中.
也有這樣說.
「我的朋友,我告訴你們.
那些殺害身體的.
之後不能做甚麼的.

$^{321}$不要害怕他們」.
當人們把你們帶到會堂.
在官長或掌權者面前時.
你們不要害怕要怎樣答辯或說甚麼.
因為聖靈會教導你們應該說甚麼.
等等.
有很多其他的經文.
我這裡舉了一些例子.
其實在新約聖經中.
跟希臘哲學一樣.
都在告訴我們.
我們不需要恐懼死亡.
為甚麼?因為我們有永生.
因為神會看著.
就算我們真的死亡也好.
是不用害怕的.
就算真的死亡也好.
也不用害怕.
最多是甚麼?我們經常說.
最多就是死亡,對吧.
最多就是死亡.
死了也沒有甚麼.
沒有甚麼可以再做,對吧.
他們最多只是殺了我們.
沒有甚麼了.
耶穌在跟門徒說.
即使有人迫使你們.
最多是甚麼?.
最壞的情況就是死亡,對吧.
所以不用害怕.
因為神會看著你.
這是當時新約的教導.
就是不用害怕死亡.
同樣,亦有些教導說.
不用恐懼神.
例如在《羅馬書》和《約翰一書》.
都有提及.
譬如《羅馬書》裡保羅這樣說.
事實上,你們所領受的.
不是奴隸的靈.

$^{361}$使你們仍然有懼怕.
你們所領受的是兒子的名分的靈.
藉著這個靈,你們可以跟神說.
阿爸,阿父.
即是當我們相信耶穌後.
相信主後.
我們跟神有一個新的關係.
奴隸會怎樣?.
奴隸會害怕主人,對吧.
但當我們成為了兒子時.
我們跟神有一個新的關係時.
我們在這個新的關係中.
不再是一個恐懼主導的關係.
不再是奴隸和主人的關係.
所以我們可以不用那麼害怕.
在人和神的關係中.
所以在這裡,你可看到.
在新約的經文中.
都有反映到不要甚麼.
不用害怕沒錢.
不用害怕窮,不用害怕死.
以及不用害怕神.
好,說到這裡.
又是跟上次的主人要問一個問題.
究竟在新約聖經中.
有沒有正面的恐懼呢.
誰覺得…又問一下,互動一點.
誰認為聖經有教我們.
關於正面的恐懼的舉手.
對了,好,又問一下.
誰覺得勝經是一面倒.
覺得恐懼是不好的舉手.
沒有,即是大家都很聰明.
都覺得應該會有正面的恐懼.
好,我們再看看有甚麼正面的恐懼.
第一件事,剛才剛才說完.
我們跟神的關係.
不是奴隸和主人的關係.
所以不用那麼害怕.
不是恐懼主導.

$^{401}$不過事實上,我們都應該害怕神.
是嗎?為甚麼.
因為在聖經中,在新約中有說.
恐懼能夠使我們認清.
誰才是我們真正的主.
這個也很重要.
就是在《馬太福音》和《路加福音》.
耶穌說過不要害怕沒錢或害怕死亡.
之外,祂亦有這樣說.
我的朋友,我告訴你們.
那些殺害身體的.
以後不再做甚麼的,不要害怕他們.
但我要跟你們說,你們應該害怕誰.
就是那個殺人以後.
還有權柄能夠將人扔到地獄的.
要害怕他.
對,我告訴你們.
你們要害怕這一位.
這個有權柄將人扔到地獄的是誰.
就是上帝.
所以耶穌在這裡.
要重新界定權力關係.
剛才我說過甚麼.
阿彌時多德是這樣說的.
阿彌時多德說.
某些人能夠使其他人害怕他.
是因為他擁有能夠傷害那些人的力量.
所以這裡是說價值判斷.
我們要問自己.
我們會否害怕錯了對象.
耶穌在這段經文中跟我們說.
我們不要害怕錯了對象.
當我們害怕的是甚麼.
當我們害怕的是害怕官員.
害怕一些能夠逼迫我們的人時.
其實我們害怕錯了對象.
所以耶穌要說.
耶穌要重新澄清.
當我們害怕那些官員.
害怕那些逼迫我們的人.

$^{441}$我們背後會說甚麼.
我們覺得他們是我們的主.
所以我們才會害怕他們.
我們覺得他們有絕對的權力.
掌握生死大權.
所以我們就害怕他們.
但耶穌在這裡要顛覆這種想法.
重新界定權力架構.
我們真正要害怕的是我們唯一的主.
是我們唯一的上帝.
只有那位我們才需要害怕.
那些能夠逼迫我們的人.
我們不用害怕他們.
最多就是死.
那個把人丟到地獄的人.
那個神,那個才值得害怕.
所以這裡跟亞里士多德所說的很吻合.
我們害怕甚麼.
就是我們覺得誰是我們的主.
如果我們害怕人而不害怕神.
即是甚麼意思.
我們當人是我們的主.
而我們不是當神是我們的主.
這個也值得我們問.
究竟我們經常說我們相信了耶穌.
究竟我們比較害怕人還是神.
還是我們覺得神是很遙遠的.
我們是害怕人而從來不害怕神.
這個是值得問的問題.
我們害怕誰.
我們有沒有錯判我們恐懼的對象.
第二件事.
恐懼使我們能夠認識真正的主之外.
一個相關的正面作用就是.
恐懼能夠促使人活出道德的生活.
例如保羅在《肥那比書》這樣說.
「這也看來,我親愛的弟兄們」.
「你們既然是常常頌從的」.
「不僅我在你們那裡的時候是這樣」.
「如今我不在你們那裡,就更要是這樣」.

$^{481}$「就應該心存恐懼」.
「戰戰兢兢,成就自己的救恩」.
「因為是神在你們裡面運行」.
「使你們立志和行事,要成就祂的美意」.
在這裡說甚麼呢.
保羅也說.
我們不僅是信了耶穌,買了保險.
我們信仰覺得信了耶穌.
就買了保險等上天堂.
我們以前經常這樣想.
但其實不是的.
保羅覺得.
你信了耶穌之後,你活出甚麼.
活出基督的生命.
你要戰戰兢兢地活出基督的生命.
你不是買了保險等上天堂.
我們要在地上仿效基督.
仿效耶穌基督的樣式.
所以我們要戰戰兢兢.
帶著一個恐懼的心.
戰戰兢兢地活出我們的救恩.
去成就神在我們生命當中的心意.
所以恐懼是有一種.
促使我們活出道德生活的作用.
你想像一下.
你做決定的時候.
例如你買一樣東西,你也怕買錯.
我們怕做錯決定.
但如果說到活出基督的樣式.
或是道德行為的時候.
你從來都不怕你有沒有做錯.
其實你是一個怎樣的人.
其實你是一個沒有道德的人.
平時說錢很容易明白.
你怕買錯東西.
你怕被人騙錢,你會害怕.
但如果你在做任何道德決定的時候.
你從來都不害怕.
其實是說甚麼.
你也不在乎道德,對吧.

$^{521}$你就是一個沒有道德的人.
所以害怕和道德是很重要的連繫.
我不是說害怕構成了所有道德決定.
但任何道德決定背後.
一定要有害怕的因素.
你明白我的意思吧.
如果沒有恐懼的話.
是不能成為一個有道德的人.
你明白我的意思吧.
所以其實新的幾個教導.
都跟現代的心理學很吻合.
有些心理學家.
很多心理學家已經是很舊的研究.
你看到是1970年代.
例如有些心理學家會說.
其實恐懼懲罰.
是人類道德發展的第一個階段.
如果人沒有恐懼.
其實是不會發展到道德的.
好嗎?跟我剛才說的一樣.
你要有恐懼懲罰的因素.
你才可以發展到道德.
亦有一些研究.
反社會人格的心理學家會這樣說.
反社會人格的人.
Anti-social的人有甚麼特質.
Anti-social的人是一種病態.
心理的病態.
他的特質是甚麼.
完全不害怕.
他無視任何道德價值.
任何道德規範.
是因為他這個人本身是無懼的.
完全沒有恐懼.
所以他無法發展任何道德價值.
無法發展任何道德價值.
所以這件事跟聖經所說的很吻合.
我們需要戰戰兢兢活出我們的救恩.
戰戰兢兢活出基督的養性.
我不是說恐懼是道德的全部.

$^{561}$但恐懼是道德的基礎.
有正面的作用.
我們需要某程度的害怕.
最後,我想回到這個問題.
我們在生命中都會面對很多恐懼.
剛才亞里士多德也說過.
勇氣不是否定恐懼.
我們回到一段新約經文.
這跟耶穌的經歷有關.
在《馬可福音》的14章.
當耶穌在「赫西瑪利園」中.
聖經很清晰地記載.
其實有清晰記載耶穌的情緒.
除了上次說.
在《馬可福音》有記載耶穌的憤怒外.
在《馬可福音》「赫西瑪利園」的記載中.
有提及到耶穌恐懼的情緒.
我在這裡讀出來.
耶穌和門徒來到一個叫「赫西瑪利園」的地方.
耶穌對門徒說.
「你們坐在這裡,我去禱告」.
於是祂帶著彼得,雅各,約翰一起去.
耶穌驚懼起來,極其難過.
對他們說.
「我的心憂傷要死,你們留在這裡要警醒」.
往前走的不遠,祂就伏在地上禱告說.
「如果可能的話,願那時刻離開祂」.
祂說.
「阿爸父,對你來說一切都是可能的」.
「求你拿走這杯」.
「但是不要依我的意思,要依你的意思」.
我們經常看這段經文.
但我們很少留意到.
原來聖經有記載到.
耶穌在這個時刻所面對的情緒.
《馬可福音》很清晰記載.
耶穌在這個時刻面對.
將要被人拉,要釘上十字架的時候.
一個很清晰的情緒是甚麼.
是恐懼.

$^{601}$其實耶穌也會害怕.
耶穌也害怕死亡.
耶穌也害怕釘上十字架.
所以耶穌在《喀西瑪利園》裡恐懼起來.
向天父祈禱.
祂求上帝將這個苦杯拿走.
讓祂不用上十字架.
不過耶穌在害怕的同時.
祂的禱告不僅求上帝拿走這個苦杯.
祂還求了一件事.
祂說.
「一切都是可能的,求你拿走這個苦杯」.
「但是不要依照我的意思」.
「依照你的意思」.
所以耶穌上十字架不是不害怕.
耶穌是害怕著地上十字架.
耶穌也面對自己的恐懼.
但正如亞里士多德所說.
真正的勇氣其實沒有否定恐懼.
勇敢的人也會害怕.
耶穌也會害怕.
但勇氣是驅使我們做一些我們害怕的事.
面對我們的害怕.
我們仍然會這樣做.
就像耶穌的禱告一樣.
求你拿走這個苦杯.
「不過不要依照我的意思」.
「要依照你的意思」.
這就是勇氣.
去面對我們自己的恐懼.
我們仍然有能力去做我們要做的事.
這種就是恐懼.
去做上帝呼召我們的事.
縱使我們害怕.
縱使我們仍然有恐懼.
但這就是勇氣.
今天很想說說我的故事.
作為今天說到的完結.
我很少分享自己的經歷.
但因為我太太Bella經常分享自己的經歷.

$^{641}$以致我便分享自己的經歷.
我便說說自己的經歷.
其實我想說甚麼.
今天說了這麼多.
既定義又古代的事.
其實你完全忘記也不要緊.
最重要是甚麼.
最重要的第一個重點.
就是我們要面對我們自己的….
我想說的是.
如果我們不面對自己的恐懼.
我們會被我們的恐懼侵蝕了自己.
我們會誤判很多事情.
我們的生命會被恐懼扭曲.
我想說說自己的經歷.
我也分享了自己的掙扎.
或是自己的黑暗面.
我曾經在…當我還年輕.
還血氣方剛的時候.
可能你現在也覺得我還是血氣方剛.
但我有更血氣方剛的時候.
當我更年輕和更血氣方剛的時候.
我曾經在一間教會.
也稱為大教會裡事奉.
當清善玄牧者.
其實那間教會.
也有一些不太健康的情況或文化.
就是童工,執事比較分黨分派.
有不同派系.
大家在職場就明白甚麼是分黨分派.
其實我想說教會也有.
不過我們不說而已.
其實教會也有.
不過我今天不是想說教會的問題.
而是想說我自己.
總之我是在一間大教會.
但分黨分派的教會裡事奉.
而我自己是一個甚麼人.
其實我也是一個害怕的人.
在《九人人格》裡就是六號.

$^{681}$六號就是恐懼.
很多事情是很害怕的.
我也是有很多恐懼在我裡面.
我很容易會害怕這個,害怕那個.
很害怕.
我Bella經常說我害怕青青.
就是我.
當我沒有去面對自己的恐懼時.
而我在充滿黨派鬥爭的教會裡事奉.
會出現甚麼情況呢.
就是在那間教會裡.
也會有很多排斥某些童工的情況.
或者要跟大隊.
就是你要跟某些大粒的人.
這樣你才可以事奉.
才可以有一個空間去事奉.
或者不會被人排斥等等.
我的問題是甚麼.
就是因為我害怕,我恐懼.
所以這是我的黑暗面.
現在回想起來.
我也很後悔曾經是一個這樣的人.
就是我有去面對.
我有去面對一些有權力的人.
我不小心地,自己也不意識.
因為我沒有處理自己的恐懼.
以致我成為了某些人的打手.
去攻擊某些童工.
或者成為了某些人的打手.
去排斥某些童工.
這是我離開了這間教會時.
才會醒悟的事.
當我在一個這樣的文化裡事奉.
當我在一個這樣的環境裡事奉時.
其實我自己也不意識.
原來我是某些人的打手.
我自己也不知道.
其實我真的不知道.
是要離開了這個環境才知道.
為何我會成為一個人的打手.

$^{721}$去攻擊某些童工.
在一些會議上.
為何我會這樣.
其實我跟他沒有仇.
後來我才醒悟.
為何我對他沒有仇,沒有怨.
但為何我會這樣去攻擊他.
其實是因為我裡面的恐懼.
在一個充滿黨派的教會裡.
你沒有黨派的話.
是事奉不了的.
是很恐怖的.
而我就是面對一種這樣的恐懼.
我不意識到這種恐懼.
以致我成為了攻擊其他人的人.
我自己也不知道.
直到離開一個這樣的文化.
這樣的….
這樣的牧會環境.
我才意識到,才發現這件事.
其實今天我們每一個人都一樣.
如果我們不面對我們的恐懼.
不面對我們的黑暗面的時候.
我們也會誤判了很多事情.
恐懼是關於價值判斷.
可能我們也會面對.
好像我的情況一樣.
就是最後,就像耶穌所說的一樣.
我們害怕別人,而不是害怕神.
而在事奉裡也會有這種情況.
我竟然害怕一些大顆的橡汁.
或者一些大顆的人.
而不是害怕上帝.
以致我做了一些錯誤的決定.
即使在事奉的場景裡.
而我們今天每一個人都一樣.
無論在我們生活,在我們的職場.
今天恐懼是一種價值判斷的時候.
我們有沒有面對我們的恐懼.
我們有沒有真正去問.

$^{761}$我們為甚麼要害怕.
我們在害怕甚麼.
我們如何面對我們的恐懼.
最後,我很想有一段.
一起唱梅應詩歌的時候.
我們透過唱這首歌.
我們有一點去禱告,反思的時間.
(梅應詩).
(愛可禱告).
(求羅海這苦澀).
(甘心擺上).
(為現成就交鎖).
(按上帝旨意不惜一切).
(甘心去想 伏至死).
(你為我犧牲 撥開傷苦間).
(總面對許多掙扎).
(仍走向各局他).
(你點燃氣象 戰勝懼怕).
(你愛已經出現交錯中間).
我們嘗試一起唱.
(愛可禱告).
(求羅海這苦澀).
(甘心擺上).
(為現成就交鎖).
(按上帝旨意不惜一切).
(甘心去想 伏至死).
(你為我犧牲 背上苦感).
(總面對許多掙扎).
(仍走向各局他).
(你定義犧牲 戰勝懼怕).
(你愛已經出現交錯中間).
(就算跌倒千百次).
(仍然不捨棄這愛不止).
(就算我這麼不配).
(仍然都喜我 全部追逐).
(就算跌倒千百次).
(仍然不捨棄這愛不止).
(就算我這麼不配).
(仍然都喜我 全部追逐).
(你為我犧牲 背上苦感).

$^{801}$(總面對許多掙扎).
(仍走向各局他).
(你定義犧牲 戰勝懼怕).
(你愛已經出現交錯中間).
(你為我犧牲 背上苦感).
(總面對許多掙扎).
(仍走向各局他).
(你定義犧牲 戰勝懼怕).
(你愛已經出現交錯中間).
讓我們一起去禱告.
今天十字架代表耶穌戰勝了他的恐懼.
當我們凝望十字架時.
我們不僅思想到耶穌為我們死.
我們更思想到.
他在喀西瑪利園時面對自己的恐懼.
他能夠戰勝它.
他能夠勇敢地走上十字架.
縱然有恐懼.
但仍然走上十字架的道路.
今天願上帝進入我們的心靈當中.
重新提醒我們.
讓我們有能力,有勇氣.
去面對自己的黑暗面.
自己的恐懼.
我們再次問自己.
我們究竟在恐懼甚麼.
我們為何要恐懼.
我們如何面對恐懼.
當我們思想到耶穌在喀西瑪利園時.
今天我們再次回想.
真正的勇氣不在於不害怕.
真正的勇氣在於你意識到你的恐懼.
並且縱然有恐懼.
仍然活出主的旨意.
這就是勇氣.
今天願意.
天父,你把這份勇氣賜給我們.
主,今天你讓我們像你一樣.
在我們自己的喀西瑪利園當中.
我們能夠像你一樣這樣祈禱.

$^{841}$你挪開我們的苦杯.
但願意你的旨意成就.
不是按照我們自己的旨意.
而是按照你的旨意.
主,今天你給我們一份勇氣.
你讓我們審查我們的恐懼.
你讓我們可以活出一個真正勇敢的生命.
去成就你的旨意.
是奉主耶穌基督的聖名祈求.
阿們.
拜拜!.
\newpage



\section{約翰福音}
\label{sec:gEwUekxM30c}
\textbf{生命頻道 \#2:聆聽上主的聲音 (約翰福音 15) - 陳崇基牧師}
\newline
\newline
連結: \href{https://youtube.com/watch?v=gEwUekxM30c}{\texttt{https://youtube.com/watch?v=gEwUekxM30c}} ~~~~ 語音日期: 2025-03-18
\newline
\newline
\hyperref[sec:KKxcZYP9oLY]{\small{< < < PREV SERMON < < <}}
~
\hyperref[sec:index_chronic]{\small{[返順時目]}}
~
\hyperref[sec:index_scriptual]{\small{[返順卷目]}}
~
\hyperref[sec:1U5fivtmb3s]{\small{> > > NEXT SERMON > > >}}
\newline
\newline
$^{1}$各位觀眾大家好.
今天我們立刻去大雪.
剛進來的時候大雪.
和大家身邊的人說一說.
多謝你們來.
多謝你們冒著風雨.
值得多謝他們.
多謝他們,真的好事.
因為我早到的.
我算是早到的,八點多就到了.
所以那時候我們還是沒有雪的.
現在就非常大雪.
好了,夠了.
不用多謝大家.
這次我們回來的主題就是.
生命的頻道.
正如我所說,我們屬靈的生命.
是有一點像收音機.
我們需要收到上帝對我們的聲音.
我們有沒有調整好我們的收音機.
很多時候我們生活忙碌的時候.
我們沒有調整好,所以我們聽不到上帝的聲音.
或者我們沒有一個習慣.
練習一下怎樣調整那條天線.
調整那個頻道.
讓我們可以說話.
和聽到上帝說話的頻道.
上帝一向都是在說話的.
而我們呢.
就是不是在聽的意思.
上一次我們說到.
牧羊人和羊的關係.
記得嗎.
上帝,耶穌基督就是牧羊人.
而我們就是牠的一群羊.
羊就是做什麼的.
就是聽見和認得到牧羊人的聲音.
要我們認得到的時候.
我們才能夠知道.
上帝怎樣帶領我們.

$^{41}$所以有時候.
你想一下,如果你信了主那麼久.
你要了解,我們和上帝之間的關係.
不是純粹頭腦上的內容.
我們可以想到很多東西.
或者去拆位.
拆很多難的概念.
我們明白到等等.
但都是需要用心去感受.
我們和上帝之間的關係.
羊和牧羊人就是一個這樣的比喻.
這次我們用另一個比喻.
就是「我是葡萄樹,你們是枝子」.
這個比喻.
就是說耶穌.
就是葡萄樹.
我們就是那些枝子.
Vine and Branches.
我們想用這個比喻.
入手去想一想.
到底我們怎樣和上帝之間建立關係.
怎樣聽到他的聲音和和他有一個連結.
這裡就是約翰福音第十五章.
十五章耶穌說到.
他和門徒之間的關係.
他和你之間的關係.
就是一個葡萄樹的樹幹.
和那些橫枝和枝子的關係.
這個如果是從小到大.
在教會長大的應該都很懂得.
這個比喻.
我們以前小時候.
主日學我們都唱一些不對音的歌.
葡萄樹,主是葡萄樹.
我們是枝子.
記到現在.
對音的未必記到.
不對音的就記到現在.
但是那個比喻.
就是說一個農田的比喻.

$^{81}$在約翰福音第十五章出現的.
這一段聖經.
約翰福音第十五章中間.
就是說到我們和主之間的關係.
這一段說話.
耶穌和門徒所說的.
對他們的叮嚀.
就是在約翰福音的後半段.
十三至十七章.
這一段聖經.
如果你有聽過的話.
我們都說過了.
就是一個道別的講話.
Farewell Speech.
Farewell Sermon.
就是耶穌臨離開這個世界之前.
臨上十課之前.
他和門徒的一些說話.
你看到十三章開始.
就告訴我們這裡有些分別.
因為十三章打後.
就是一些很重要的講話.
重要的事跡.
重要的事件發生.
十三章第一節.
他說在逾越節以前.
晚餐的時候發生什麼事.
就是大家吃飯.
就是逾越節的晚餐.
然後耶穌吃飯的時候.
他就離席了.
離開他餓的位置.
坐餓的位置.
然後他就做什麼.
他就脫下衣服.
然後就去到束上腰.
然後去餵門徒洗腳.
做一個瀑人.
其實這個洗腳的動作.
就是指向十字架上.

$^{121}$他傾倒的生命.
十三章洗完腳.
然後他就開始說話.
十四章說了一段說話.
然後十五章就說到.
葡萄樹與枝子之間的關係.
這裡可以分四段.
第一節到十一節.
就是說到主和門徒之間.
親密的關係.
然後第二段十二至十七.
就說到其他門徒.
門徒與門徒之間.
就是我和你之間的關係.
然後就說到門徒.
就是教會和外面的世界.
這個黑暗的世界.
之間有什麼關係.
現在外面的世界.
去逼迫他們的世界等等.
就是和他們敵對的.
之間有什麼關係呢.
然後二十六至十七節.
就是聖靈和我們之間的關係.
因為耶穌說我要走了.
我不會丟下你.
會有聖靈來到我們中間.
和你們同在.
然後就說到十七章.
大祭司的禱告等等.
十八章第一節就說完.
這一堆說話.
十三節就十七章.
十八章第一節說什麼.
他說說了這些話.
就是道別的講話.
說完了.
臨離開之前說完這個Farewell Speech.
他就跟門徒出去了.
去晚餐的樓房.

$^{161}$然後過了吸倫谷.
然後過了對面一個地方叫做赫西瑪利園.
然後就被捕.
上十字架.
這個就是脈絡.
所以十五章說得很重要.
就是你和上帝之間的關係.
在這個Farewell Speech那裡.
有沒有試過有些親人離開世界的時候.
他說的話特別重要.
一個人快要離開的時候.
說的話很重要.
所以耶穌在離開這個世界之前.
他也有一些這樣的吩咐.
就是你和上帝之間的關係.
你和神之間的關係.
當然這裡說到葡萄樹與枝子的關係.
這個就是當時的背景很清楚.
因為他們是一種農田的生活.
他們會栽種一些葡萄園等等.
我想告訴你.
在幾個禮拜前.
我跟一個國語組去查經.
我們查到詩篇23篇.
查完之後.
詩篇23篇大家認識吧.
耶和華少我都目正不至缺乏等等.
我們就說神和我們之間的關係.
就是牧羊人和羊之間的關係.
這個應該是很溫馨的意境.
對不對.
查經之後有個弟兄走來問我.
牧師,這個好像不是很好.
因為牧羊人和羊之間的關係.
牧羊人養羊是殺牠們.
然後吃羊腩煲.
因為牧羊人是北方人.
吃羊肉.
我問他怎樣解釋.
神是我們的牧羊人.

$^{201}$養我們,然後養肥了我們.
然後吃我們嗎.
我就跟他解釋.
我說不是的.
要用他們的背景.
他們不是經常吃肉的.
我們每天都吃肉.
昨天我去一個小組聚會.
去吃火鍋.
一大堆肉在桌上.
吃都吃不完.
以前的人.
現在其實也是.
地中海的食物.
Diet.
很少有肉的.
通常都是vegetarian.
對不對.
所以他們養的羊不是用來吃的.
如果吃了羊就不行了.
為什麼.
因為羊是生產什麼.
奶,然後做芝士.
有羊毛.
一養肥了.
然後做低它,然後吃了.
那你以後就沒有了.
所以羊和牧羊人之間的關係.
應該用巴勒斯坦的文化去了解.
這裡也是一樣.
葡萄樹與枝子的關係.
都是用巴勒斯坦地常見的異境.
好像這裡的圖畫一樣.
他們當時在山區生活的時候.
他們會在山坡建立一些梯田.
因為平地比較少.
在那裡種一些果樹.
其中一種出名的果樹.
在巴勒斯坦地就是葡萄樹.
葡萄樹,誰有種葡萄的.

$^{241}$在後院那裡.
苦瓜我知道有.
冬瓜也有.
還有合掌瓜.
為什麼我知道.
雖然我不種,但我經常收到一些東西.
有沒有,有誰.
沒有.
沒有種葡萄.
因為如果你種葡萄的話.
種葡萄是很難的.
種葡萄要花很多功夫.
因為葡萄樹.
葡萄樹是要園主花很多心機去栽種.
才能夠有很多年之後.
才能夠有好的果子.
所以他們就會在葡萄園建立一些棚架.
因為葡萄樹的主枝幹.
是比較軟身一點的.
你們見過吧,有沒有印象.
你們是懂得喝紅酒的人嗎.
他們要躲在棚架上才能夠繼續生長.
才能夠出現葡萄.
當然這個就是巴勒斯坦地.
本地農耕的文化裡.
才能夠了解到.
我們現代人比較難了解.
因為我們去到市場買葡萄.
就已經搞好了.
或者去買一瓶紅酒.
那些已經做好了.
他們是從開始.
在棚架上.
他們每一天都照顧那些枝子.
要他們生長要他們健康.
然後才能夠出現葡萄有收成.
這個才是背景.
所以耶穌就用一個這樣的比喻去形容.
他和門徒之間的關係.
當然葡萄樹與枝子.

$^{281}$也都在舊約裡.
因為是農耕的文化.
經常都用來形容.
上帝和以色列人之間的關係.
神就是葡萄園的園主.
他是那個Gardener.
或者那個Owner.
然後他就嘗試栽培這些葡萄樹.
以致他出到好的果子.
但如果上帝的子民是聽神的話的話.
他們就會出到好的果子.
意思就是說.
他們的生命就會有上帝的本性.
他們的生命 他們的生活.
就會有公義 有真理 有平安等等.
然後見證到上帝在他們中間.
當他們聽話的時候.
就會出到好的果子.
就會有好的生命.
如果他們不聽話的時候.
就反而出一些酸的葡萄.
或者壞的葡萄果子.
結果就會怎樣.
如果你是園主.
你花了那麼多心機.
如果他出不到好的果子.
你會做什麼.
剪了它 燒了它.
就不要了.
這個當然就是舊約以色列國的結果.
因為他們經常都去敵黨上帝 杯葛.
所以最後就是焚燒 然後亡國.
所以當耶穌說.
他說我是真葡萄樹.
不只是單單是葡萄樹.
他說我是True Vine.
他就說我才是真正的以色列.
耶穌說我和我的門徒.
我所代表的那個群體.
就是一個新的以色列.

$^{321}$True Vine.
意思就是說以前那個.
就是False Vine.
意思就是說當時那個宗教的領袖.
就不是真正的.
我才是真正的.
我才是真正去到遵守上帝的說話.
然後帶領我的子民來到去.
成為上帝的子民.
I am the True Vine.
我們進入經文.
我們試一下用一個這樣的背景.
剛才我解釋了的.
我們看看這裡和我們和耶穌之間的關係是怎樣.
約翰福音第十五章第一節.
我就是真葡萄樹.
我乎是栽培的人.
凡屬我不結果子的之子.
我就剪掉凡結果子的.
我就修剪乾淨.
使之子結果子更多.
現在你們因我講給你們的道已經結成了.
你們要常在我裡面.
我也常在你們裡面.
之子若不常在葡萄樹上.
自己就不能結果子.
你們若不常在我裡面.
也是這樣.
我就是葡萄樹.
你們是之子.
常在我裡面的.
我也常在他裡面.
這人就多結果子.
因為你了我.
你們就不能作甚麼.
人若不常在我裡面.
就像之子被雕在外面.
枯乾了.
人撿起來.
冰在火裡燒了.

$^{361}$你們若常在我裡面.
我的話也常在你們裡面.
凡你們想要的祈求.
就給你們成就.
你們多結果子.
我父就因此得榮耀.
你們也就是我的門徒了.
我愛你們正如父愛我一樣.
你們要常在我的外裡.
你們若遵守我的命令.
就會常在我的外裡.
正如我遵守了我父的命令.
常在他的外裡.
我知道有一點點拗頭.
你在我裡面.
我在你們裡面.
我們解釋一下.
這裡說到葡萄樹與之子的關係.
當然非常清楚.
這裡說我是葡萄樹的時候.
這是約翰福音的七個.
耶穌說我是的宣告裡面.
最後一個.
你記得我們上次說到好牧人.
他說我是好牧人.
那是最後一次公開說的.
但在約翰福音裡面.
1,2,3,4,5,6,7.
說了七次.
是我是的I am.
到最後.
這就是我們今天的經文.
十五章第一節.
就是真葡萄樹.
我是真的葡萄樹.
他說我是True Vine.
你看到七個的比喻裡面.
我是生命的糧.
世界的光.
陽的門.

$^{401}$好牧人.
等等都是用不同的比喻.
或者不同的實物.
來形容我們與主之間的關係.
如果主是活水的話.
我們沒有活水.
就會怎樣.
就會死.
如果主是生命的糧的話.
如果我們沒有生命的糧.
我們也是死路一條.
說到葡萄樹與之子都是一樣.
最後的說我是.
最後這個比喻.
就是說葡萄樹與之子的關係.
是怎樣的關係.
葡萄樹就是生命的源頭.
我們作為門徒.
我們的Source.
我們的源頭就是葡萄樹.
因為我們是之子.
如果不連結於樹幹的話.
我們就沒有上帝的生命.
或者我們拒絕在主的裡面.
所謂拒絕連結於耶穌的話.
拒絕這個生命的源頭的話.
那又怎樣.
也是死路一條.
我們就不能夠為上帝的結果而自由.
我們生命也不會有改變.
不能夠與上帝連結.
這裡出現了不同的翻譯.
我是葡萄樹.
你不是之子.
上在我裡面的.
我也上在他裡面.
這人就多結果子.
因為你來了我.
你就不能夠作什麼.
不等於你來了上帝.

$^{441}$來了耶穌.
你不能夠吃飯.
不是這個意思.
你不能夠作什麼.
是作什麼.
結果子.
因為之子最重要是做什麼.
就是結果.
對不對.
就是結果.
所以他們用一些語句去形容這個關係.
樹幹和之子之間的關係.
就是上在我裡面.
或者住在我裡面.
英文很多時候會翻譯成.
Remain in me.
或者舊一點的翻譯叫做.
Abide in me.
當然Abide in me就是舊一點的.
我們很少用Abide這個字.
你不會說.
I invite you to abide in my home.
是Remain.
就是上在樹幹上面.
是不是很圓.
你要上在我裡面.
我就上在你裡面.
其實是連結.
在我家門口有一棵樹.
是一棵Berry Tree.
Mulberry Tree.
我們住在家裡住了二十多年.
每年到大概這個時間.
它就開始開花.
然後結果.
然後那些Berry就會.
掉到地上.
有一棵Berry Tree好嗎.
不好是嗎.
好嗎 拿來吃嗎.

$^{481}$對於我們來說就一般般.
為什麼.
因為Berry一掉到地上.
就很麻煩.
我們回家就要踩那些Berry.
藍色 紅色.
踩了進屋裡.
很煩的 每年都是這樣的.
當然我們說環保.
不要亂砍樹.
但我心裡面都覺得.
那棵樹很煩.
經常生的苺子掉下來.
然後有蒼蠅跑過來.
有老鼠跑過來.
或者有一次有些Raccoon跑過來.
是有點麻煩的.
直到兩三個星期前.
政府終於可能收到一些多年的投訴.
二十多年了.
終於跑過來就砍了那棵樹.
一砍就砍成這樣.
其實那棵樹開始生的時候已經分枝了.
生得不太好.
所以前兩天我就經過那個地方.
因為他砍樹拿走了其他東西.
只剩下殘枝.
前兩天我經過的時候就撿起來.
你想不想看一下.
就是這棵.
這棵奇枝很大.
這裡你還可以看到照片.
還有些小枝在出現.
是不是.
其實是可以吃的.
這棵就是枝子.
葡萄樹和枝子的關係.
雖然這棵不是葡萄樹.
這棵是Berry的樹.
但都是一樣.

$^{521}$那棵樹已經砍了.
在地上就撿到這條這樣的東西.
這條東西還可以結果嗎.
什麼可能不可以.
簡直就是不可以.
什麼事啊你們.
機器人沒那麼厲害嗎.
AI沒那麼厲害嗎.
你可不可以在這裡.
然後你就說有些Berry出來了.
然後就去砍它.
因為它已經斷離了樹幹.
你已經沒有任何養分了.
是不是.
你就算等它.
等多久都好.
它都不會出任何果子.
那結局是什麼.
結局就是它焚燒.
是不是.
已經沒有用了.
因為它做不到它本身有的功能.
它本身有的功能就是去結果子.
這個就是常在我裡面.
不常在我裡面的問題.
如果我們作為枝子.
如果我們不連於樹幹的話.
我們就不能夠去結果子.
不能夠結果子.
中國就是死亡.
所以枝子要做的是什麼.
就是連著樹幹.
用這樣的比喻形容我們基督徒的生命.
就是說我們需要和耶穌之間.
有一個很緊密的關係.
並且還要讓原定修剪它.
這裡就是其中一段的說話.
他說凡結果子的.
神就會修剪和乾淨.
洗枝子呢 做什麼.

$^{561}$修剪枝子不是因為手癢.
要去剪爛它.
而是希望那些養分盡量去到一個地方.
然後去結到好的果子.
耶穌說現在你聽我的道.
明白我的說話.
就已經乾淨了 被修剪了.
你們若常在我裡面連結於我的話.
我的話也常在你們裡面.
凡你們想要的 祈求的.
就給你們成就.
當然這裡不是說我要種綠合彩.
我要種彩票.
就為我成就.
不是 凡你想要 枝子想做什麼.
結果子嘛.
對嗎.
枝子不是想發達.
枝子是想為上帝來結果子.
這就是我們應該有的功能.
有說著上帝的生命 表露上帝的生命.
表露基督的品格.
言行舉止等等.
凡我們想要的 就給我們成就.
因為我們跟祂連結.
並且讓祂來修剪.
這裡有些圖畫告訴我們.
在巴勒斯坦地.
或者在不同地方.
歐洲不同的葡萄園.
或者在加州納粹谷.
那些出紅酒的地方.
都會見到這樣的情況.
那些葡萄樹的枝子就會被修剪.
每逢到冬天的時候.
大概這一段時間.
大概一二月的時候.
葡萄樹就會有一個休眠的期.
休眠的期.
它就不做任何事.

$^{601}$在那段時間.
農夫就會將每一條的枝子來做修剪.
修剪一些什麼.
沒用的東西.
那些賤肉橫生的地方.
將它剪掉.
為了什麼.
就是要增加那個結果.
剪除什麼.
我們叫它廢柴.
將那些廢柴剪掉.
然後就準備.
這個時候春天來到了.
那個樹幹和那些枝子.
就會開始去到生長.
這個時候就會慢慢剪一些比較小的.
一些比較軟弱的.
叫做弱枝.
不是弱智.
那些弱枝就需要修剪.
以致它快要到九月的時候.
秋季的時候.
收割的時候.
就能夠有更加豐富的果實.
你了解這樣的圖畫的時候.
你就知道圓釘.
圓柱對葡萄樹的照料.
是非常細心.
因為每一根枝子都要摸過.
然後檢查過.
去到修剪.
和牧羊人和羊之間的關係都相似.
牧羊人知道每一隻羊.
認識每一隻羊.
圓柱也知道每一根枝子.
甚至到今天21世紀.
已經科技這麼發達.
這麼多年.
二千多年以上.
去到栽種葡萄樹的方法.

$^{641}$其實都是一樣.
都是要靠人手.
你不可以有些組合.
有些組合.
然後全部東西修剪.
是每一根圓柱花心思去修剪.
枝子上的壞枝幹.
當然要被剪除.
以免浪費養分.
壞的.
壞的枝幹當然要剪除.
本身好的枝幹.
有時候都會剪除.
好的枝幹.
所謂好的枝幹都是健康的.
但都會剪除,為什麼?.
因為你希望養分去一個地方.
然後集中火力去有果子.
所以不只是壞的.
或者所謂好的都會剪除.
以致你能夠有最好的.
你有壞的枝幹.
當然你不可能有最好的果子.
但如果你有所謂好的.
但不是最好的枝幹.
你都要剪除.
以致你能夠有最好的果子.
有最好的果子.
這個應用在我們心靈裡面其實很有意思.
上帝在我們心靈裡面的工作.
其實有時候都是這樣的.
有時候上帝在我們心靈裡面.
會將一些東西剪掉.
那個過程可能不是我們很想要.
不過就是為了去訴訟我們.
以致我們能夠結更豐盛的果子.
有時候上帝去修剪我們的心靈.
修理我們的心靈.
不一定只是剪掉一些壞的東西.
也可以剪掉一些好的東西.

$^{681}$但這樣做有一個目的.
目的就是讓我們的心靈能夠結出最好的.
最好的東西.
怎樣剪?.
他說用上帝的道.
這個信仰.
如果你是初信的話.
或者你信了很多年的話.
都需要留意的.
信仰不是純粹頭腦的東西.
是心的東西.
是我們心要去領會的.
上帝的話語其實都是一樣.
我有一次在一個小組查經.
最近有一個小組查經.
有一個在醫院做神經科的研究.
那個朋友來到.
因為他說我去其他教會.
很多時候問一些信仰問題.
別人覺得很麻煩.
好像沒有一個很自由的空間.
讓我去探索.
他就被帶到我們這個教會.
我就在小組跟他討論.
因為他說其他教會回答的問題.
通常都是.
我有什麼信仰問題.
然後他就拿聖經來說.
聖經這樣說.
所以你要這樣相信.
然後他說這個我接受不了.
然後他就找其他經文跟他說.
我覺得這個可以給一些人.
但有很多人是不行的.
他就問了一些很好的問題.
問一些信仰的問題.
就是神的存在.
神的本性是怎樣等等.
有一次查經完了.
我們就播一首歌.

$^{721}$一首現代的中文詩歌.
他聽完這首詩歌.
他說我很感動.
因為我從來沒有聽過這樣的詩歌.
因為他的教會都是一些傳統的聖詩.
傳統的聖詩是可以的.
但每個人的台不一樣.
我們的頻道每個人都不同.
我就播了那首歌.
其實YouTube找的很容易.
打下去就播出來.
他就覺得很受感動.
我就突然間了解.
你去到信主.
去到跟隨耶穌.
其實不只是這裡.
是你的心.
你的心.
是需要通一通那個channel.
說到靈性的栽培.
靈性的栽培.
或者靈性的指導等等.
當然這不是我的專項.
但說到靈性的栽培等等.
所謂靈修神學.
其實真的需要留意.
不只是這裡.
而是一些心的東西.
心的東西有一點圓.
可不可以說一二三四五六.
按部就班.
很硬性地做.
也不能夠.
因為心的東西是通不通.
這件事是有一點神秘.
有一點奧秘在.
但奧秘的東西不等於.
不要去掌握.
奧秘的東西不等於掌握.
也需要去掌握.

$^{761}$說到靈性的栽培等等.
有時候我們可以在一些.
教會歷史的所謂的偉人那裡.
學一些東西.
其中一位就是Andrew Murray.
十九世紀的一個牧者.
他寫了一本書叫做.
Abiding in Christ.
什麼意思?.
住在基督裡面.
為什麼他會寫這些東西?.
有一年他事奉了很多年.
神很縱容他.
他有一年突然間幻想.
喉嚨感染.
喉嚨說不出話來.
說了這麼多年.
突然說不出話來很慘.
我最近身體也有毛病.
手也有毛病等等.
如果你說我突然說不出話來.
其實我真的挺慘.
因為我是口水佬.
你問我太太也知道.
她在後面說你閉嘴.
你安靜一點.
安靜一點.
突然間他喉嚨受了很重的感染.
長達兩年多的時間.
幾年.
他說不出話來.
然後他心情就非常低沉.
很抑鬱.
他就被迫要獨處.
被迫去到.
在上帝面前完全安靜.
就在這段時間.
他就去看書看聖經.
然後寫書.
他就去探討一下.

$^{801}$如何領略上帝的工作.
領略上帝的說話.
然後自己心裡面.
內心去降伏於上帝.
上帝在他內心裡面不斷有工作.
藉著一些收錢的經歷.
一些不太好的經歷.
所以其中一本書他寫成的.
就是《住在基督裡》.
Abide in Christ.
Abiding in Christ.
是教會歷史裡面其中一本.
靈修神學的經典之一.
可以算是經典之一.
他說什麼呢?.
他說與普通書.
耶穌去連結是需要時間.
除非你願意給主時間.
否則你不要指望.
你能夠住在裡面.
你是需要每一天.
與耶穌來同座.
這裡就說到我們信仰裡面.
不能或缺的部分.
這個也是我自己的偏見.
大家知道我是很理性的人.
我每件事都想去搞通.
理念上我們需要搞通一些東西.
如果是這樣的話.
有時候就沒有了那種心.
心的那一部分.
你要去經歷的.
Experiential.
你要親身踏進去去經歷.
在禱告裡面去經歷.
你是要願意踏出這一步.
那些不是去解釋什麼很難的概念.
而是心要走出去.
Andrew Murray說.
我們需要花時間.

$^{841}$在心靈裡面有一個這樣的模式.
我們才能夠跟上帝.
可以叫他做.
我們常在他裡面.
他也常在我裡面.
那種常在的連結的那種關係.
Andrew Murray.
另外一個也是講靈修的.
或者講靈修的.
屬靈的經歷的人.
在基督教當然有其他傳統.
天主教等等.
另外一個就是John Wesley.
John Wesley是一個英國的佈道家,牧者.
他是巡渡會.
有沒有聽過巡渡會?.
巡渡會的創辦人就是John Wesley.
18世紀的一個牧者.
也是事奉了很多年.
事奉了很多年之後.
他就感覺到很沮喪.
很depressed.
每一次做事的時候.
他好像沒有什麼信心.
沒有那種的.
怎麼說?.
英文怎麼說?.
不quick.
你明白我的意思嗎?.
不quick.
就是跟上帝不quick.
但是他當然忠心繼續做.
有一次他回到英國的時候.
在美國做了一輪的侍工.
他回到英國.
就跟另一個牧者朋友聊天.
那個牧者朋友給他鐘.
那個牧者跟他說.
Preach the faith.
until you have it.

$^{881}$他已經跟他這麼說.
繼續下去.
繼續忠心這麼說.
這樣忠心的宣講.
信心就在這裡.
直到你能夠得到信心的狀態.
有一次在一年的5月24日.
他很不想回教會.
但是他也照樣回.
在英國一個小的教會.
那個教會位於Ordersgate Street.
在倫敦.
在聚會裡面.
他聽到有人說.
因信清二羅馬書的道理.
然後唱詩.
唱一首詩歌叫做《詩篇130篇》.
《詩篇130篇》是說什麼呢?.
你回去看看.
就是說耶和華.
耶和華的一個禱告.
在深處我去呼求你.
在谷底我要去呼求你.
是一個在低谷裡面發出的禱文.
那天的聚會結束後.
他回到家裡.
他就在日記裡面這樣記錄.
他說很奇怪.
I felt my heart strangely warm.
在一個不願意.
但是又繼續去的狀態下.
他覺得他心靈裡面很奇妙.
溫暖起來.
他不知道怎麼解釋.
但是他感到什麼呢?.
就是我仍然確信基督.
唯獨基督我才能夠得救.
並且賜給我一個得救的確據.
也深知耶穌基督已經除去我的罪.
正視我的罪.

$^{921}$並且拯救我脫離罪和死的律.
這個就是我們信仰裡面.
我們需要肯去答開,答出.
一個親自經歷過的經歷.
為什麼這麼重要呢?.
為什麼有時候可能比頭腦上解通的東西更重要.
就是因為如果你經歷過的話.
那件事就是屬於你的.
是沒有人拿走的.
John Wesley就是這樣去經歷.
如果你在倫敦走一趟.
去一個小的城市,小的地方.
就有一些紀念碑.
5月24日某年.
John Wesley那裡有個牌子寫著.
"I felt my heart strangely warm".
左下角有一個藝術品.
一個火焰的藝術品.
其實是一個火焰.
有些刻了一些字.
就是講他日記裡面所講的東西.
突然間感到這個信仰不只是這裡.
而是這裡.
很多時候講道的時候.
我都會被上帝的話語去觸碰.
我記得我最近那一次.
當然大大小小有時候都會有.
我去年就預備了《但義理書》.
因為它多人多去到港島.
一開始答應去講《但義理書》.
因為跟那裡的主席和他們的搞手去講的時候.
我就說.
我心裡面說.
因為《但義理書》七百九十章很難講.
很難講.
我心裡面說.
沒有人講就讓我講.
因為我喜歡這些理性的挑戰.
但我在預備的當中.
當然我花了很多功力和功夫去預備.

$^{961}$我預備的時候有一個掙扎.
就是我預備不下去.
因為我覺得裡面講的意象很震撼我.
人子的國度.
那種真實.
我一開始看《但義理書》的時候不是很覺得.
但我慢慢看多了的時候.
很多時候晚上我去到預備的時候.
很多晚上我去掙扎去預備的時候.
要花很多功夫去預備的時候.
我簡直是震撼到我做不下去.
我一直流淚.
因為我看到人子要駕駛運而來.
他是那位德性的人子.
那個國度是永不被滅亡的.
居然我有份在這個國度裡面去參與.
是不是.
我和上帝講我是什麼人.
我是什麼人.
我只是一個罪人.
上帝的國度居然可以包括我在裡面.
一個小小的角色.
那時候我就realize.
原來看聖經不只是透腦的事.
是需要心去領會.
而要心去領會的話.
就需要我們肯起碼.
我們願意.
我盼望我們和神的關係都一樣.
今天可能聽完之後覺得有點懸.
但其實可以是很真實.
可以很奧秘.
但卻是很real的東西.
John Wesley有一個禱文.
是幫助不同的人.
有很多人讀完這個禱文.
或者用這個禱文來禱告的時候.
很幫助到他和上帝親近.
用心去領會上帝.
我想用這個禱告來作結束.

$^{1001}$這個所謂納若的禱文.
Covenant Prayer.
當我讀的時候.
我希望你用心和我一起來禱告.
主我不再屬於自己.
乃是屬於你.
求按你的旨意使用我.
將我安置在你所願意的地方.
讓我為你作供.
甚至讓我為你受苦.
讓我被你使用.
甚至願我暫時被擱置.
讓我被高舉.
甚至讓我被卑微.
讓我充實.
甚至讓我被倒空.
讓我擁有一切.
甚至讓我一無所有.
我甘願樂意.
將一切全然的教徒.
任憑你的旨意成就.
如今榮耀而蒙福的上帝.
聖父聖子聖靈.
你是我的.
我也是你的.
願如此成就.
在地上所立的約.
願你在天上建立.
阿們.
我們同心協力地祈禱.
主這是我們禱告.
主願你的旨意.
在我們生命裡面能夠成就.
或許是高升我們.
或許此刻是要卑微我們.
或許這時候要使我們和你同工.
或許是我們暫時被擱置.
要走過幽谷.
主我們都甘願一切教導.
願我們常在你裡面.

$^{1041}$你也常在我們裡面.
我們這樣祈禱奉主名禱.
好,願上帝祝福各位.
\newpage



\section{約翰福音}
\label{sec:1U5fivtmb3s}
\textbf{生命頻道 \#3:聆聽上主的聲音 (約翰福音 14-17) - 陳崇基牧師}
\newline
\newline
連結: \href{https://youtube.com/watch?v=1U5fivtmb3s}{\texttt{https://youtube.com/watch?v=1U5fivtmb3s}} ~~~~ 語音日期: 2025-03-25
\newline
\newline
\hyperref[sec:gEwUekxM30c]{\small{< < < PREV SERMON < < <}}
~
\hyperref[sec:index_chronic]{\small{[返順時目]}}
~
\hyperref[sec:index_scriptual]{\small{[返順卷目]}}
~
\hyperref[sec:ETczAl2v1zo]{\small{> > > NEXT SERMON > > >}}
\newline
\newline
$^{1}$各位聽眾好.
這個星期是春假,是嗎?.
是不是有些有孩子的都出去了?.
我看到不是.
我以為你們全部都走了.
這個星期是春假.
所以希望你們有一個好的春假.
是不是一個星期?.
一個星期而已.
大學的春假應該剛剛結束.
起碼是這區的大學.
所以今天我兒子又回去了.
他不會煩著我.
這個星期是春假.
所以我太太和我女兒.
因為他們都是跟學校的行程.
所以他們明天早上就飛去加州.
正如我之前說過.
他們都計劃去加州.
上次他們突然要回去.
因為有些家事要處理.
所以明天他們就回去加州.
兩個人回加州.
這個星期我就一個人在家.
他們下個星期二才回來.
下個星期二他們回來之前.
我又要飛了.
因為我要去加州市.
我要去上學院報到.
因為有一個計劃要交.
所以我兩個星期都是一支弓.
真是好了.
無論如何.
我們繼續講道的系列.
就是聽上帝的聲音.
當你信了主之後.
應該有一個轉變.
一個很基本的轉變.
在你心靈裡面.
就是上帝的生命進入了我們.

$^{41}$聖靈住在我們裡面.
然後我們會聽到聖靈的聲音.
我們要學習怎樣聽.
因為好像心機一樣.
我們要調整台.
然後我們才聽到主對我們說的話.
對我們的提醒.
第一次我們講到牧羊人和羊之間的關係.
今天詩歌又唱.
上一次我們講到.
葡萄樹和枝子的關係.
都是一個很緊密的關係.
所以我們需要去培養.
我們屬靈的頻度.
我們要聽一下上帝對我們的聲音.
然後我們的頻率要調整得準確.
我們才能夠聽到.
正如我所說.
我們每一個人信了主.
我們都有聖靈住在我們裡面.
今天我們需要講一講聖靈的聲音.
當我們聽到聖靈的時候.
Holy Spirit的時候.
我們可能第一個反應就是.
覺得我們很神秘.
或者一些很難懂的東西.
摸不著頭腦的東西.
或者一些奇怪的東西.
因為我們很多時候接觸.
或者聽見一些屬靈的經驗.
或者一些靈因的經驗.
特別就是講方言.
倒地.
或者被聖靈充滿等等.
我們就有一點敬而遠之.
還有我們會覺得有點模糊.
好像捉不著.
好像風一樣捉不著.
今天每個人都講天氣.
我就講一下風.

$^{81}$我就選了這張照片.
因為確實耶穌是這樣形容聖靈的工作.
聖靈在我們生命裡面重生的工作.
regeneration.
rebirth的工作.
當我們信了主之後.
有一個重生的經歷的時候.
就好像聖靈在那裡做事.
然後聖靈做事的方法.
好像風一樣.
耶穌和尼哥底姆.
或者尼哥德姆.
去到談道的時候.
夜間談道的時候.
就講到這件事.
就是聖靈像什麼呢.
聖靈像風.
風是看不見的.
摸不到的.
但是風所帶來的效果.
你就能夠看得見.
對不對.
打風也好.
或者我們在芝加哥封城也好.
我們全部都看到風.
所有的consequence.
或者做出來的那種效果.
所以聖靈對我來說.
有時候有一點捉不著.
就是有一點模糊.
今天我們想去看一看.
到底耶穌怎麼樣去討論聖靈.
聖靈是一回什麼事呢.
我們首先搞定一些比較框架的東西.
這個就是所謂系統神學裡面.
所結論出來的.
什麼叫做聖靈.
或者什麼叫做神.
神自己本性是怎麼樣的.
這裡我嘗試用七句話來說.

$^{121}$這個就是神.
是三位一體的真神.
三位一體的神.
右邊那一圖可以給你一點印象.
第一樣,聖父是神.
聖子是神.
聖父不是聖子也不是聖靈.
聖子不是聖父也不是聖靈.
聖靈也不是聖父也不是聖子.
但不是三個神.
是一個神.
可以嗎?背了它可以嗎?.
這是一個很大的總結.
正如我所說.
總結或者一些神學的綱領.
其實就是一些很抽象的總結.
因為是一種總結.
所以有很多重要的細節會跳過.
所有總結都是這樣的.
總結不會說完所有東西.
總結只會說出一些框架的東西.
對嗎?明白嗎?.
我做了一整天工作回到家.
太太問我.
How's your day?.
我就跟她說.
OK啦.
OK啦就是總結.
等於沒有說過.
她要的不是這個東西.
因為我跳過了很多細節.
總綱領或者綱領的essential.
很多時候有這樣的限制.
但無論如何.
這就是所謂的Trinity.
就是三合一體.
裡面說到聖靈.
聖靈是什麼呢?.
聖靈就是三合一體的第三個位格.
The third person.

$^{161}$person是有位格的.
等於我們這裡有多少人.
我們有一個位.
也有一個格.
聖靈Holy Spirit就是third person.
但聖靈是哪一個呢?.
我們就看看耶穌如何去到講論.
我們上個星期開始.
我們都是在說什麼?.
《約翰福音》第十四至十七章.
那裡是一個離別的講話.
Farewell speech.
Farewell discourse.
你記得那個情況是怎樣?.
是耶穌和門徒守最後的晚餐.
然後在所謂上樓房的地方.
二樓的地方.
在飯廳那裡.
跟他們說了很多話.
就是《約翰福音》第十四至十七章.
上一次我們說到其中一段.
就是葡萄樹與枝子的關係.
我們要結果子的話.
我們有上帝的聖名的話.
就要連於哪個?.
連於那棵樹.
因為我們就是那棵樹幹.
不是樹幹.
我們是樹枝.
這個就是離別講話的其中一段.
在《約翰福音》第十四至十七章.
這個Farewell講話裡.
離別講話裡很多時候提到聖靈.
我們看看耶穌如何去形容聖靈.
對我們有什麼形容呢?.
所以我們只需要.
我們嘗試回答三個問題.
三個大問題.
就是聖靈是誰?.
第一個問題.

$^{201}$第二個問題.
就是聖靈作什麼?.
What does he do?.
三個問題與我們有什麼關係?.
聖靈是誰呢?.
《約翰福音》第十四至十五節.
他說你們若愛我.
耶穌與門徒說.
或者與教會說.
因為在上流放.
十一個門徒剩下.
賣耶穌的猶太人已經走了.
那班門徒就是代表所有的教會.
他說你們若愛我.
就會遵守我的命令.
我要求父父就賜給你們.
另外一位保衛師.
使祂永遠與你們同在.
祂就是真理的靈.
是世人不能接受的.
因為他們既看不見祂.
也不認識祂.
你們卻認識祂.
因祂常與你們同在.
也要在你們裡面.
這裡說什麼?.
如何形容聖靈?.
另外一位.
Another one.
意思是什麼?.
說話的是誰?.
是耶穌.
耶穌與門徒說話.
與教會說話.
他說我會另外猜查一位.
父賜給你們另外一位.
就是說有另外一位.
類似祂這樣的.
是世人不能接受的.
未信耶穌的不能接受.

$^{241}$因為不能認識祂.
也不認識祂.
說到認識的時候.
不是說冷知識.
不是一些Facts.
是Knowing.
即是說是關係.
你可以對一個人.
有很多冷知識.
你認識他.
知道他有什麼Facts.
但你是否認識他?.
不一定.
所以他說世人不能認識的.
是一種關係.
另一個位格可以與他建立關係.
下面這句聖經.
在約翰一書二章第一節.
那裡形容.
耶穌就是第一位.
所謂的.
說什麼?.
他說在庫裡面.
我們有一位的中寶.
就是那二者.
耶穌基督.
他為我們的罪作了贖罪制.
即是去洗淨我們等等.
那裡說第一位.
耶穌基督.
聖靈就是另外一位.
Another one.
And another one.
然後怎樣形容他?.
形容他說是保衛師.
這個真是奇怪.
我們平時不會說保衛師.
我們平時日常用語.
我們不會用保衛師.
是什麼意思呢?.

$^{281}$這個是很多的討論.
聖靈是保衛師.
這個字.
在新約裡面用過五次.
剛才已經說了約翰一書.
二章一節是形容耶穌.
耶穌也是保衛師.
也做中寶.
作挽回制.
作贖罪制.
然後十四章十六.
十五章二十六.
十六章七節.
另外四次.
就是用在聖靈身上.
這個字就叫做Parakletos.
Parakletos大概意思.
希臘文是兩個字拼在一起.
即是Alongside, Called.
兩個字.
Called to walk alongside.
即是在你身邊.
跟你走的意思.
大概.
基本上是這個意思.
但是不同的聖經的譯本.
就會用不同的字.
來翻譯這個Parakletos.
Parakletos是什麼呢?.
有不同的翻譯.
當你看聖經的時候.
我們有不同的翻譯本.
去比對的時候.
我們就發現這個字.
是很特別的字.
即是說可能現代的語文.
很難去將這個.
所謂Parakletos這個字.
來翻譯.
找到一個合適的字去翻譯.

$^{321}$所以有這麼多不同的翻譯本.
有不同的翻譯.
保衛師,中保.
訓衛師.
護衛者.
或者英文Counselor.
Helper.
然後最後一個就是.
辯護者,Advocate.
每個翻譯都有利有弊.
如果說做Counselor的時候.
你就會想起去見輔導.
或者去見治療師.
或者說到安慰的時候.
就會是一種Comforting.
即是去安慰你.
然後給你一些支持等等.
訓衛呢.
這個是天主教師高譯本的翻譯.
或者護衛,護衛者.
護衛者就是當代譯本的翻譯.
各有千秋,各有好,各有不好的地方.
強調的很多時候就是.
如果中保的話就是和事佬.
即是中間.
合起來.
但通常都是印象比較溫柔一點.
誰是Holy Spirit.
比較溫柔一點.
因為他是Counselor,因為他是Helper.
因為他是給你安慰的等等.
可不可以捕捉全部呢.
所以就有些困難.
怎樣翻譯保衛師.
這件事.
可能一個最好的翻譯就是最後這個.
Advocate.
辯護者.
這個有什麼不好處呢.
就是這個詞很少用.

$^{361}$Advocate,誰會用Advocate的.
平時.
辯護者.
可能你去看醫生的時候.
看醫院的時候.
有間醫院叫做Advocate.
你會知道有Advocate的Hospital.
但為什麼這樣翻譯會比較好一點呢.
因為這個詞.
通常用在一些地方都是用法律的用語.
在法庭上發生什麼事.
然後有一個護衛者.
Advocate.
一個辯護者.
來保障你或幫你講話的.
Speak for you的.
這個人就叫做Advocate.
正如今天醫院醫療制度都會有一種叫做Patient Advocate.
Patient Advocate是什麼.
你進到醫院的時候很嚴重.
然後很亂.
然後可能有些醫院或者一些制度.
會給你一個所謂辯護者.
一個Patient Advocate.
去保障病人和了解你的情況.
然後告訴你有什麼方案等等.
站在你的方向去講話.
站在你的方向去幫助你.
這個優點就是這樣.
因為這個詞通常在當時的語文.
或者當時的文學.
裡面出現的時候.
就是講到法律的層面.
Forensic Legal的東西.
所以可能這個Advocate.
是比較容易讓我們想到.
到底聖靈是做什麼.
另一位的保衛師.
另一位的Advocate.
這個Advocate就是耶穌.

$^{401}$叫父差派給我們.
雖然耶穌會離開這個世界.
他會死而復活.
然後升天.
但是耶穌說.
跟門徒說.
我會給你一個Advocate.
一個辯護者.
站在你的方向去幫助你.
所以你就可以想到.
在法庭上如果我是被告.
有一個Advocate在你身邊.
如果這個是聖靈.
聖靈做什麼呢.
我們儘管看看一些經文.
當然這裡講了很多.
就是十四至十七章.
Faith Discourse裡面.
講了很多跟聖靈的工作有關的東西.
不過我們只能夠選擇性地去看一點.
這裡絕對不是說半個小時裡面.
我們就可以囊括所有聖靈所做的事情.
但是有一些經文值得我們留意.
十四章第十六節.
這個也是整個告別的講話裡面.
其中一段.
十四章第十六節.
他說他永遠與你們同在.
講得很聖靈Advocate.
你們卻認識他.
因他常與你們同在.
也要在你們裡面.
我不會撇下你們為孤兒.
我必到你們這裡來.
凡愛我的就會遵守我的道.
我父也愛他.
並且我們要到他那裡去與他同住.
好像講得很遠.
但是有一樣東西你們要留意.
他說你們是眾數.

$^{441}$就是在說門徒群體教會.
聖靈會與你們同在.
常與你們同在.
在你們裡面.
在舊約裡面說到聖靈的時候.
或者在兩約中間說到聖靈的時候.
說到上帝的榮耀要降臨.
因為彌賽亞日子來到的時候.
就會有聖靈降臨.
就會有神的榮耀來到.
當你想到神的榮耀的時候.
在舊約裡面.
神的榮耀住在哪裡.
住在聖殿裡面.
所以住在聖殿裡面的話.
就與其他的地方不同.
當上帝的榮耀進入了他的聖殿.
這個聖殿就與其他非聖殿的建築物.
有不同的地方.
但是到了新約的時候.
聖靈就住在哪裡.
住在百姓裡面.
不是住在用石頭做的聖殿.
而是住在所有相信耶穌基督的百姓身上.
在你們中間.
但是到最後這一節.
我們要到他那裡與他同住.
就單數了.
凡聽見我道的.
就會有誰.
三位一體的上帝.
聖父,聖子,聖靈.
與他同住.
Make our home with him.
Make our home with her.
如果你聽見這些道理的時候.
你應該覺得很讚歎.
上帝居然住在我們中間.
與他同住.
God makes his home.

$^{481}$in you.
in me.
所以彼得後書裡面說了一句話很有趣.
當我們接受了耶穌.
我們有一個救恩的時候.
上帝的大能和一切關乎.
屬靈生命.
敬虔的事賜給了我們.
因我們認識.
那用自己的榮耀和美德照我們的上帝.
已經認識了這個上帝.
與他有一個認識的關係.
他將有寶貴又極大的應許賜給我們.
使我們脫離世上從情慾來的敗壞.
就分享上帝的本性.
We share.
In some sense we share.
不是說我們變了神.
而我們share了上帝的Divine nature.
是一個很大的特權.
是因為耶穌基督為我們所做的.
所以我們怎樣.
我們就share了這個Divine nature.
Divine nature.
我們再看多幾段.
十四章 二十六 十五章 二十六 十六章 六八節等等.
我們試試去到咀嚼一下.
這裡到底說什麼.
和聖靈有關的工作.
我們先讀一下二十六節 十四章.
但保衛師 誰啊 聖靈.
就是父因我的名所要差來的聖靈.
他要將一切的事指教你們.
並且要叫你們想起.
我對你們所說的一切話.
十五章第二十六節.
但我要從父那裡差保衛師.
就是從父出來真理的聖靈.
他來了就要為我作見證.
十六章第八節開始.

$^{521}$他來的時候要為罪為義為審判指證罪人.
為罪是因為他們不信我.
為義是因為我到父那裡去.
你們將來再不見我.
講到他升天 講到耶穌升天.
為審判就是這個世界.
這個黑暗的世界.
統治者撒旦已受到審判.
只等真理的聖靈來到.
他要引導你們明白一切的真理.
因為他不是憑自己說話的.
是憑他所聽見的都說出來.
他要榮耀我.
這裡很多顏色.
我們想點出一點點.
聖靈的工作是什麼呢.
剛才我們講到聖靈內住.
住在我們裡面.
Make his home in us.
然後這裡講到什麼.
聖靈的工作就是指向基督.
聖靈很多時候不會指向自己.
不會說自己多大能量.
多多事做等等.
聖靈做的事.
就是為耶穌做見證.
向世人去指明耶穌是神的兒子.
聖靈的工作就是指向基督.
所以很多時候我們接觸的.
我們覺得追求聖靈的工作等等.
或者我們近乎到一些.
可能很神秘.
非常神秘.
或者捉不到的一種經歷.
但這段聖經告訴我們.
聖靈做什麼.
其實它是指向基督.
我記得以前在加拿大溫哥華.
讀神學的時候.
其中一科系統神學.

$^{561}$就是一個教授.
很出名的教授教的.
J.I.Packer.
離開了世界好幾年.
J.I.Packer是很出名的.
英國的神學家.
在那裡教.
每次跟他上課都很有趣.
每次上課之前他會唱歌.
唱Doxology.
因為他覺得Theology就是Doxology.
神學就是讚美師.
就是讚美眾.
所以一上課的時候.
開始他就會一起起來.
就會唱Doxology.
然後才開始講課.
J.I.Packer形容聖靈的工作是什麼.
形容聖靈的工作是Floodlight Ministry.
Floodlight是什麼.
就是我們平時後壓.
我們經過有一個感應器.
然後我們亮燈.
有賊就亮燈.
那個叫Floodlight.
Floodlight就是把光射到後壓那裡.
或者射到路上.
聖靈就是一個Floodlight Ministry.
就是燈.
就像燈一樣.
你不是直接看它.
你也不需要直接看它.
因為那個燈泡的工作.
不是要吸引你去看它.
而是希望什麼.
是光照其他地方.
以至你能夠看到其他東西.
這裡就是聖靈所做的工作.
Floodlight Ministry.
它是做什麼.

$^{601}$榮耀我.
聖靈是榮耀基督的.
聖靈是指引人向基督.
使人看到路.
使人看到道路真理生命.
看到基督.
所以不是吸引你去看它.
所有的燈都是這樣.
事實上這是聖靈把耶穌所做的事.
神兒子所做的事.
所做的一切.
更完整地啟示出來.
十六章這裡說什麼.
聖靈要叫人為罪為義為審判.
指正世人.
以前和學本的翻譯.
就是自己責備自己.
當你信主的時候.
應該有一種這樣的經歷.
就是聖靈去到光照.
然後你為罪為義為審判.
首先只是講罪.
簡單一點講罪.
就是自己責備自己.
我知道我是罪人.
我需要什麼.
我需要耶穌基督.
誰讓你看到.
黑漆漆.
就是聖靈光照底下.
你能夠看到.
你記得第一次有一個.
所謂conversion.
這些歸主的經歷.
就是五重節的時候.
史唐尊第二章.
彼得講完道之後.
人們聽完道有什麼反應.
他們覺得窒心.
心裡面有點不舒服.

$^{641}$然後他們就有一個反應.
他們說我們應該做什麼.
彼得跟他們說.
你們要信耶穌基督.
悔改歸向上帝.
這是聖靈的工作.
和聖靈對你們說的話.
我剛才講的時候.
你記不記得.
你以前信主的時候.
是有這樣的經歷.
或者信主之後.
我們久不久都會有這樣的經歷.
是不是.
為罪自己責備自己.
這是聖靈的工作其中一樣.
指你向基督.
以至你能夠抓緊基督的教授.
第二就是辯護.
剛才我們講聖靈保衛師.
怎麼翻譯.
advocate.
是不是.
辯護者.
你可以想像.
為什麼這個更加貼身的翻譯.
那個情況就好像法庭裡面.
你是被告.
然後你是有罪的.
但是有一個大狀.
不是毒舌大狀.
是辯護的大狀.
聖靈就來到你身邊.
Call to be alongside.
To defend.
去到維護你.
他就會做什麼呢.
大概.
聖靈就在你身邊.
然後向義者法官.

$^{681}$指著基督.
耶穌基督.
為了你已經犧牲了.
代替你.
如果你和法官說.
如果你還要他負責的話.
這是雙倍負責.
這是不義的.
真正公義就是.
這個所謂恩順清義的道理.
耶穌基督義者.
已經為你贖罪了.
這個就是.
那個defender.
或者那個advocate.
會和你做的.
聖靈會和你做.
對不對.
然後就說到一堆教導的字眼.
指教.
想起.
引導等等.
聖靈不是來給一些創新的東西.
很多時候他的功能就是提醒你.
用什麼提醒你.
是根據上帝的說話.
根據耶穌所說的話.
因為提醒你.
讓你想起耶穌所說的.
所以這兩次說到的.
真理的聖靈.
Spirit of truth.
或者Spirit that tells the truth.
或者那個靈.
是會對你說真理的.
是根據上帝的說話.
來和你說的.
所以既然根據上帝的說話.
和你說的話.
你怎樣去到能夠.

$^{721}$辨察到聖靈的聲音.
聽到聖靈的聲音.
就需要你有根底.
是不是.
你要學習去測驗一下.
去根據一下.
神的說話是怎樣說的.
聖靈就根據上帝的說話.
來對你作指教.
叫你想起一些東西.
引導你.
上個星期.
我就和一個區會的牧師去吃飯.
在對面吃飯.
這個牧師.
是我們區會其中一個.
他經常到處去的牧者.
和他吃飯的時候.
他就說.
我最近在說一個.
什麼.
港島的題目.
原因是什麼呢.
因為他們家裡收拾東西.
這個牧師家裡有一個很小的小孩.
幫爸爸去收拾書本.
你知道牧師最多的書是什麼書.
聖經.
我們有很多本聖經.
他有不同翻譯的聖經.
小孩就把聖經塞進書架裡.
誰知道小孩就反過來放那本聖經.
那本聖經的頁向外面.
整堆的聖經的頁向外面.
他就留意到一件事.
因為聖經通常是舊款印出來的.
就是金邊銀邊.
所以金邊銀邊就向外面.
他一看的時候.
他就覺得搞錯了.

$^{761}$但他一看的時候就發現.
因為聖經很多時候都用過.
有些有銀邊有些沒有銀邊.
意思是什麼.
你看聖經的時候.
你看得多銀邊就掉了.
但有些沒有揭過的地方就是銀色.
他就把聖經放在那裡.
他就發現每本聖經有些地方是沒有銀邊的.
有些地方是有銀邊的.
代表那些地方是他很少看的.
他就說有一個sermon series.
就是講這件事.
silver lining sermon.
就是聖經裡面很少講的東西.
或者他自己很少看的東西.
譬如說詩之書等等.
anyway.
我這樣講的意思是什麼呢.
聖靈的工作就是做什麼呢.
就是幫我們去根據上帝的說話.
來聽他的聲音.
如果你在這方面.
你不認識對上帝的話語是很片面的話.
那就很難分辨到聖靈對你說什麼.
對不對.
很難.
在這個世界裡面有很多噪音.
兩天前我開車的時候.
有輛車突然停在路中心.
不知道為什麼.
然後大家都彎過去拍他的頭.
然後我看一看.
原來司機.
我以為他心臟病或什麼事.
原來他在看電話.
是不是你啊.
很多distraction.
但我們需要被上帝的說話去吸引.
以至聖靈才能夠提醒我們上帝所說的話.

$^{801}$我記得年輕的時候.
大學的時候.
我記得有一個星期.
每一天我用三四個小時去看聖經.
去看《史蘭傳》.
我不知道為什麼我選擇了《史蘭傳》.
我每天用了兩三個小時.
我就是在看聖經.
我就感應到.
上帝所說的話.
不是說明天什麼時候會下雨.
叫做擔傘.
而是上帝告訴我們他的權能.
他的工作.
很清楚地跟我說.
我這裡想挑戰你.
我之前說過.
用十分鐘在上帝面前.
每天用十分鐘去安靜.
然後十分鐘去看《約翰福音》第十章.
今天我想挑戰你這星期.
儘管你用一個小時.
整個小時.
你專心地去看聖經.
我讓你看《約翰福音》第十三至十七章.
儘管你專心地去看.
你這樣下功夫.
你這樣去看上帝的話語的時候.
聖靈就會跟你說話.
為什麼會這樣做呢.
因為我們這些屬靈的習慣.
其實是幫助我們去培養一個.
The Heartful God.
我們對上帝仰望的一種心態.
是鍛鍊出來的.
如果你有Heartful God的時候.
你才能夠分辨到聖靈對你.
聽到聖靈對你說的話.
對嗎.
Heartful God.

$^{841}$對上帝愛慕的一種心態.
是要培養出來的.
我們很多時候說到聖靈的時候.
我們說到聖靈充滿.
我們說聖靈充滿.
如果聖靈充滿我就好了.
我們就想起像液體或風一樣.
越來越令我膨脹.
越來越吹噪我.
這樣看是否正確呢.
是不是真的這個意思呢.
我昨晚有這樣的想法.
因為我兒子昨晚吃飯的時候跟我說.
他說你知不知道.
如果你去吃All You Can Eat.
All You Can Eat Sushi的時候.
因為他知道我今天會跟年輕一點的去吃Sushi.
他說你知不知道.
如果你去吃All You Can Eat Sushi.
怎樣吃得最多呢.
他說你不是餓自己.
他說你整天在喝水.
他說因為喝水的時候.
你的胃就會脹大.
但又不會飽.
然後你吃東西之前.
一個小時.
吃東西之前一個小時.
你就可以吃一點沙拉.
你可以啟動你的胃口.
他這樣跟我說.
我跟他說我都老了.
我不會吃那麼多.
我是陪人去吃的.
但我們在想聖靈充滿的時候.
我們就想吃多一點.
脹多一點.
這樣就叫聖靈充滿了.
但如果你想.
我記不記得我開頭說什麼.

$^{881}$聖靈在一個胃隔的時候.
那個想像就很不同.
什麼叫做a person.
聖靈如果是一個person的時候.
會有什麼不同.
在你的生命裡面有什麼不同.
我記得好幾年前.
好多年前我們曾經.
我們家裡.
我女兒剛出生.
我們就去接待一個.
外來的教會的講員.
我們當然很尊重他.
他就住進來我們家裡.
我記得他住進來我們家裡.
三幾天裡面.
我們整個家裡很多變化.
我和我老婆少了吵架.
他少了嘮叨.
我少了摒他等等.
我女兒又少了哭.
很奇怪.
那時候她大概一歲多.
她還在嬰兒床.
她每天早上就問我.
那個叔叔醒了沒有.
想去找他.
然後乖乖的就去睡覺.
這就是我的意思.
如果你有一個awareness.
聖靈在一個person的時候.
就不是我們自我膨脹.
就是awareness.
它的presence.
不是portion.
不是多少.
是它的presence.
你想像一下.
它的presence來到我們家裡.
來到我們生命的時候.

$^{921}$會發生什麼改變.
這就是我們屬靈生命裡.
我們需要培養的意識.
那種awareness.
聖靈在這裡.
在我們生命.
在我們家庭的時候.
祂會給我們什麼改變.
結束的時候.
讓我們用這個祈禱.
我寫的這個祈禱.
作結束好嗎.
聖靈.
Oh Holy Spirit give me.
Ears to hear your voice.
Eyes to see your glory.
A tongue to proclaim your truth.
A heart to love your goodness.
A life to obey your will.
聖靈請你賜我.
能聽見你聲音的耳朵.
能看見你榮耀的眼睛.
能宣講你真理的使徒.
能愛慕你美善的生命.
能信服你旨意的生命.
我們低頭祈禱.
主多謝你賜給我們聖靈.
聖靈住在我們裡面.
祂的同在.
祂的presence.
去幫助我們.
來更加愛慕你.
願你對我們說話.
願我們聽見你的聲音.
並且由生命.
去愛慕和信服你的心意.
我們這樣祈禱奉主名.
願上帝祝福各位.
\newpage



\section{約翰福音 6:52-60}
\label{sec:ETczAl2v1zo}
\textbf{生命麵包,你今日食jor 未?(約翰福音 6\_52-60) - 梁紫健傳道}
\newline
\newline
連結: \href{https://youtube.com/watch?v=ETczAl2v1zo}{\texttt{https://youtube.com/watch?v=ETczAl2v1zo}} ~~~~ 語音日期: 2025-04-01
\newline
\newline
\hyperref[sec:1U5fivtmb3s]{\small{< < < PREV SERMON < < <}}
~
\hyperref[sec:index_chronic]{\small{[返順時目]}}
~
\hyperref[sec:index_scriptual]{\small{[返順卷目]}}
~
\hyperref[sec:95H_Df1nGvs]{\small{> > > NEXT SERMON > > >}}
\newline
\newline
約翰福音 6:52-60
\newline
\begin{longtable}{cl}
\hline
\hline
章節 & 經文 (和合本修訂版)\\
\hline
6:52 & \begin{tabularx}{0.7\textwidth}{X} 因此，猶太人彼此爭論說：「這個人怎能把他的肉給我們吃呢？」 \end{tabularx} \\ \\ \relax
6:53 & \begin{tabularx}{0.7\textwidth}{X} 耶穌對他們說：「我實實在在地告訴你們，你們若不吃人子的肉，不喝人子的血，在你們裡面就沒有生命。 \end{tabularx} \\ \\ \relax
6:54 & \begin{tabularx}{0.7\textwidth}{X} 吃我肉、喝我血的人就有永生，並且在末日我要使他復活。 \end{tabularx} \\ \\ \relax
6:55 & \begin{tabularx}{0.7\textwidth}{X} 我的肉是真正可吃的；我的血是真正可喝的。 \end{tabularx} \\ \\ \relax
6:56 & \begin{tabularx}{0.7\textwidth}{X} 吃我肉、喝我血的人常在我裡面，我也常在他裡面。 \end{tabularx} \\ \\ \relax
6:57 & \begin{tabularx}{0.7\textwidth}{X} 永生的父怎樣差我來，我又怎樣因父活著，照樣，吃我肉的人也要因我活著。 \end{tabularx} \\ \\ \relax
6:58 & \begin{tabularx}{0.7\textwidth}{X} 這是從天上降下來的糧，不像你們的祖宗吃過嗎哪還是死了；吃這糧的人將永遠活著。」 \end{tabularx} \\ \\ \relax
6:59 & \begin{tabularx}{0.7\textwidth}{X} 這些話是耶穌在迦百農會堂裡教導人的時候說的。 \end{tabularx} \\ \\ \relax
6:60 & \begin{tabularx}{0.7\textwidth}{X} 他的門徒中有好些人聽見了，就說：「這話很難，誰聽得進呢？」 \end{tabularx} \\ \\
[1ex]
\hline
\hline
\end{longtable}
$^{1}$各位電影姐妹早安.
很感恩再一次有生命氣息在大家當中分享上帝的話語和大家一起學習.
大家如果看PowerPoint的話.
生命麵包.
這四個字你會想起什麼呢.
沒錯.
對於普遍的香港人或者南方的中國人來說.
可能會想起加頓生命麵包.
每個字都會想起加頓.
每個早上都有它的存在.
用它來做多士.
可以說是一代人早上的生命食糧.
生命麵包原來對於某些香港的年輕人就會想起.
這個是左下角的那幅相.
就是有人跟他說.
我下班就會去吃海底撈.
你就在這裡吃生命麵包.
這四個字夾雜著很多情緒.
很多不同的感覺.
對於某些很喜歡吃香港下午茶的朋友.
就會想起左上角那幅畫.
那幅相.
剛剛焗好的菠蘿包.
熱辣辣.
吃一塊在冰箱剛拿出來的牛油.
冰凍.
咬下去.
外脆內軟.
入口即溶.
再配一杯凍檸茶.
簡直是一個高級的享受.
每個年代.
不同的人對於生命麵包都有不同的理解.
或者有少少不同的感覺.
就好像門徒.
耶穌時代的門徒和猶太人一樣.
經歷了政權的打壓.
生活苦不堪言.
看望有些東西想改變現狀.
所以他們很自然對生命麵包有少少不同的理解.

$^{41}$甚至有些超越了想像的期待.
有時我們都會像我們這樣.
對上帝的工作有一定印象.
會有些期望.
覺得上帝一定是這樣做事的.
一定是這樣.
我們可以捉到他的套路.
不過每一次我們以為捉到上帝路的時候.
以為上帝一定是這樣做事的時候.
上帝一定會每一次再一次地反轉我們的認知.
今天我們要看的經文.
就是在說這件事.
我想和大家一起看看.
究竟當時門徒和猶太人期望的生命麵包是甚麼.
而耶穌又如何反轉門徒的期望和想法呢.
經文是這樣說的.
約翰福音六章五十二至六十節.
因此猶太人彼此爭論說.
這人怎能把他的肉給我們吃呢.
耶穌說.
我實實在在告訴你們.
你們若不吃人子的肉.
不喝人子的血.
就沒有生命在你們裡面.
吃我的肉喝我的血的人就有永生.
在末日我要叫他復活.
我的肉真是可吃的.
我的血真是可喝的.
吃我的肉喝我的血的人常在我裡面.
我也常在他裡面.
永活的父怎樣猜我來.
我也因父活著.
照樣吃我肉的人也因我活著.
這就是從天上降下的糧.
吃這糧的人就永遠活著.
不像你們的祖宗吃過嗎哪還是死了.
這樣的話是耶穌在加伯隆會堂裡教訓人說的.
他的門徒中有好些人聽見了就說.
這話甚難誰能聽呢.
這道糧其實原文就是麵包這個字.

$^{81}$所以第一個問題就是.
為什麼有些人包括門徒.
覺得這款生命麵包很難接受呢.
是不是他們不喜歡吃呢.
這個答案是會幫助我們了解一下.
究竟門徒是怎樣理解生命麵包的.
對於有猶太背景的門徒來說.
在天上降下的糧.
很可能就會想起嗎哪.
嗎哪是什麼呢.
嗎哪在舊約裡面是代表上帝供應日用的飲食.
以色列人的生命得以延續是因為上帝賜嗎哪.
讓他們能夠享受著和上帝的關係.
《出克俠記》十六章是這樣說的.
耶和華對摩西說.
我要將糧食從天上降下給你們.
百姓可以出去每天收每天的份.
好試驗他們遵不遵守我的法道.
以色列人只要遵守耶和華的法道.
他們的律法就會得到糧食.
得到嗎哪.
強調的是上帝在人的生命裡面的所有供應.
都是來自於上帝的.
就例如我們面前的食物.
Croissant.
或者我們工作的能力的機會.
每一個人每一秒的呼吸.
都是來自於上帝給我們的.
祂才是生命的來源.
而不是我們自己.
嗎哪也是這樣的存在.
提醒著以色列人.
生命來自於上帝.
而不是以色列人自己.
我們就問.
既然這樣.
為什麼他們不接受耶穌就是生命的麵包呢?.
原來在兩約之間.
以色列人對生命麵包有多重的理解.
有本書叫《所羅門自訓》.

$^{121}$在兩約之間.
就是在新約和舊約之間.
是代表著猶太人對於舊約的一些理解.
好像《色經》書一樣.
它不代表是真理.
但是會幫助我們理解猶太人.
怎樣去對嗎哪.
他有他的期望.
嗎哪在《所羅門自訓》這本書裡面.
是代表著壓迫下的拯救.
反抗著羅馬軍隊的一個象徵.
這段很長.
就是對付不敬虔之人的仇.
以色列人得次嗎哪.
相反地.
你將天使的食物餵養你的百姓.
從天降下準備好的餅.
給他們不必勞力獲得.
這餅美味十足.
適合個人的口味.
你所供給的.
表明你對兒女的是多麼甜美.
按照個人的喜好而供給.
味道隨個人的所願而變化.
這是在說什麼呢.
這個嗎哪很棒.
你想它是菠蘿包就菠蘿包.
你想它是叉燒包就叉燒包.
是小籠包就小籠包.
每天都有不同的口味.
你想就行.
這個名副其實是一個心想事成包.
你想不想吃這個包.
你想不想.
其實真正的嗎哪.
當然不是像這裡所說.
根據個人的口味而變化.
其實這裡所魯門之訓.
他對著一些羅馬軍隊的壁壩的時候.
他又不可以說得很明白.

$^{161}$所以他說羅馬軍隊是獸.
將他們形容為野獸.
而這裡的嗎哪就是天使的食物.
給各人有能力去反抗他們.
反抗這個暴政.
所以在兩者之間.
你要發現生命麵包.
不單只是代表著以色列人對上帝.
上帝對以色列人的應許.
更加是一個力量的泉源.
是要反抗這個暴政.
至少以色列人是這樣看.
他有這個期望.
所以我們回到約翰福音六章的時候.
我們就想.
耶穌講生命的良從天降.
以色列人會有一個誤解或者一個印象.
就是說耶穌終於來了.
他真的終於來了.
不單是做醫病趕鬼.
行神蹟.
而是給他們有從天而來的能力.
去帶領以色列人反抗暴政.
如果門徒和猶太人有這樣的理解的話.
為什麼他們又不接受耶穌這個生命麵包呢.
耶穌有能力可以反抗暴政.
行那麼多神蹟.
這一刻耶穌又沒有表示.
他不會做這件事.
他反而在做什麼.
他在做的是.
你不如吃了我吧.
他講這句.
這個生命麵包這件事.
竟然令到門徒覺得很難.
為什麼呢.
如果大家有印象.
在猶太律法有關血的條例.
大家喜不喜歡吃豬紅.
喜不喜歡吃牙血.

$^{201}$原來在教會是不准的.
血代表什麼.
生命沒錯.
唯有上帝才是生命的來源.
當人犯罪.
用血作代價.
是不是.
你要用生命換生命.
你才可以贖罪.
所以血是贖罪之用.
人是不可以操控的.
在利未極十七章是這麼說的.
吃血的後果是怎樣呢.
大家現在有印象的話.
他會被剪除.
好像一把剪刀一樣.
被人剪了.
意思是.
人與上帝之間的關係.
連結完全剪斷了.
和上帝完全隔開了.
你找不到上帝.
就好像上帝在WhatsApp裡.
傳訊息給你一樣.
你不斷找他.
喂 上帝.
耶穌 你在嗎.
他都沒有回應你.
還要斷水斷糧.
還要沒有上帝指紋的福分.
更加沒有生存的可能.
所以說到這麼嚴重.
不同的一定覺得很難.
以色列人一直禁止吃血.
但現在.
耶穌你竟然叫我連肉帶血地吃.
就說吃了你有生命.
這個似乎和現實非常矛盾.
是不是.
早前LA大火.

$^{241}$因為火勢非常強勁.
原因是風很大 很快.
有些人說.
我們叫消防員去放火.
放火就可以救火.
我們就問為甚麼.
這件事表面上是很矛盾的.
你叫消防員去放火.
不是去救火嗎.
其實他們是這樣說的.
原來消防員可以.
把火燒到一個範圍.
一個圈.
令到火保留在某個階段.
某個範圍.
直到火燒盡裡面所有燃料.
那他自己會熄滅.
表面上我們在想.
原來很多事情是很矛盾的.
不可行原來消防員可以這樣做.
很多事情在我們世界上都是這樣.
表面上很矛盾.
不合我們邏輯.
或者在我們有限的思維下.
我們覺得這件事不行.
但原來是可以這樣發生.
耶穌這裡同樣地提出一個.
表面上很矛盾但實際上可行的事.
就是當我們見到生命麵包的時候.
他有另外一面.
他就說你要吃這塊生命麵包.
這個是我的血和肉.
吃了就有生命.
但這個這麼矛盾的理念.
無人或者至少是門徒.
是很難接受的.
因為為什麼.
如果你要接受這個生命麵包.
你就要承受很大的風險.
你要承受被上帝剪除的風險.

$^{281}$萬一.
如果你站在門徒的角度.
耶穌水他們怎麼樣.
耶穌騙他們怎麼辦.
所以成了門徒他們就想.
特別是有猶太背景的門徒就想.
在66節他這樣說.
從此他門徒有許多他們的血料.
不再和他同行.
這些門徒是怎麼樣.
一直跟著耶穌.
看著他行神蹟.
隨便走趕鬼醫病.
但現在耶穌將自己成為了生命麵包的時候.
門徒就開始接受不到.
這件事完全不符合他的邏輯.
這件事跟他以往的經驗完全不同.
跟他接受的教導完全不同.
其實我們每一個信徒.
跟著耶穌的人.
我們以為我們都接受了耶穌這個生命麵包.
我們以為我們已經完全擁抱了.
這個表面上很矛盾的信仰.
但心底裡我們其實.
是不是有些.
不多不少的疑慮和質疑.
心裡的想法就是.
耶穌會不會水我.
會不會騙我的.
我自己有這樣的經歷.
信了之後.
我一直以為我自己都接受了.
耶穌這款生命麵包.
我相信他犧牲自己的生命.
為我們成就救恩.
相信大家決志那一刻.
都有說過這句.
我信耶穌為我的罪釘在十字架上.
跟大家一樣都是這樣.
信主在教會長大.

$^{321}$第一個令我覺得.
超越我對上帝想像的是我的病.
可能在某些場合都說過.
我出生就有老癌症.
天生.
從小到大我都會很容易暈倒.
突然間就會.
沒有所有東西.
就好像一部電腦沒有儲存文件.
就忘記了所有東西.
對於一個經常暈倒的人.
經常都好像要重設.
要重設.
就會覺得很無力.
很無人.
這個世界是不是來受苦的呢.
沒什麼意義.
那時候我就在想.
如果這個世界有上帝.
祂一定是一個水餓.
一個耍我.
耍得很厲害的上帝.
醫生還跟我說.
如果你十八歲之後還有這個病.
基本上都不可能會好起來.
出奇地的是.
上帝透過弟兄姊妹的愛.
令我感受到.
原來耶穌給每一個人.
都有人生的照明.
祂不會水餓.
當我信了主之後我就想.
我都很希望耶穌做這麼多神蹟.
我很想很想.
祂好像聖經裡面那樣.
醫治我.
喜來恨走那樣.
但是我心裡面都有一個.
很強烈的想法就是.
我相信不到耶穌會這樣做.

$^{361}$因為在我的邏輯裡面.
這個世界這麼多人.
這麼多苦難疾病.
什麼時候會輪到我.
憑什麼上帝會抽中我呢.
整個世界七十多億人.
是不是.
比中六合彩還難.
如果你計算一下.
所以我都懷著半信半疑的心態.
去祈禱耶穌.
我知道你有能力醫治我的.
如果你願意就可以成就.
我真的很想像一個普通人一樣.
讀書 學習.
然後我每一次讀書前.
我都這樣祈禱.
突然有一天我感動.
在我心裡面有聲音叫我停藥.
因為我吃的藥.
其實是幫助我腦的神經慢慢.
slow down 我自己的反應.
令我不會那麼容易短路.
但是這件事其實是恐怖的.
因為如果我停藥的話.
我發作的機會會大大增加.
還會增加我生命的風險.
就是因為如果我突然在馬路上發作的話.
我會暈倒.
可能就被車撞死.
直覺上都知道這件事是違反常理.
但是我真的很想被醫治.
那我就繼續祈禱.
有一天我真的自己停了藥.
結果是甚麼呢.
結果我真的發作得很厲害.
我停了藥.
我真的不斷地暈倒.
一暈倒就在醫院.
我再跟耶穌說.

$^{401}$是不是真的 是不是真的.
是否真的可以.
對我來說我已經很絕望.
我已經覺得我這一生都要背著這個病.
誰知道有一天我發現.
我好像很久都沒有發作.
然後我去看醫生.
一直有醫生監視著我.
然後我跟他說我停了藥.
他說你神經病.
不可以停藥的.
醫生就說那你做一下掃描吧.
誰知道發現腦裡完全沒有異常的東西.
完全沒有.
然後他就說.
你很可能已經痊癒.
很可能.
我不知道為甚麼上帝會醫治我.
這個世界很多病都未必得醫治.
但這個是他的決定.
如果這個醫治.
就這個醫治而言是完全不符合邏輯的.
根據醫生的經驗 根據常人的邏輯.
是完全不符合的.
他這樣救贖我的生命.
正當我以為我過了這一關的時候.
我以為我接受了耶穌這個生命麵包.
他醫治了我 超越了我的想像.
的時候大家看到右上角這幅畫.
這幅圖.
其實就是我和Alice.
可能有一個錯覺 我們都很開心和開朗.
人生沒有甚麼風浪.
其實這個世榮背後有很多挑戰和辛酸.
我們在2020年5月結婚.
大家如果記得2020年是甚麼時間.
是全球疫情大流行的時間.
在2020年我們準備結婚的時候.
大家可能知道 學生讀書.
每時每刻都在等上帝的馬欄.

$^{441}$等待著上帝的所需 救今天的需要.
我試過銀行只剩下98元港幣.
但我發燒 我要看急症.
當時急症在香港是100元.
現在就不止了 加起來幾百元.
我問Alice借我兩元.
因為你知道ATM不能按98元出來.
要按100元.
我拿了兩元就去銀行.
按了100元出來.
所以我以為艱難的日子已經為上.
但當我們結婚的時候.
我們的戶口真的只有兩位數字.
又再次在考驗我信不信上帝.
我再次問 耶穌會不會水我.
會不會再騙我.
讀書之前 曾經有位老師說.
我們一起來探索一下.
芝加哥會不會在這裡讀書.
這個探索之旅之前.
剛剛也說到我們會結婚.
《失望分文記》我們也有想過.
我們真的沒有錢了 還來探索嗎.
誰知我的老師說.
畢業了 看看吧.
看看神學院怎樣做的聖經研究.
然後我就說.
如果實習教會讓我請假.
我就來試試.
我心想教會本身已經安排好所有活動.
暑假 清少年 排滿了.
應該不會有機會請假.
然後我就跟我的上司 牧師說.
我想去芝加哥探索一下.
不知道行不行.
他問都沒有問.
就去吧.
然後請假之後我更害怕.
我跟老師說.
我請假但我沒有錢.

$^{481}$我很喜歡去這個球場.
左上角那個球場.
上次也有介紹過.
我跟Alice踩後的地方.
這也是我經常跟上帝祈禱的地方.
我跟上帝說.
你獲婚了.
結婚要花很多錢.
還要去探索.
但我一元也沒有.
所以他跟我說.
我是你天父爸爸.
我會供應.
那一刻我真的哭了出來.
沒辦法相信.
為什麼呢.
因為從小到大的經驗.
幾乎所有事情都要自己爭取回來.
我相信不到.
上天真的會像麵包一樣.
有錢掉下來.
我的經驗限制著我的邏輯.
在我的邏輯裡.
這件事從來沒有出現過.
我回去算了算.
費用比我一年的學費還要貴.
我跟老師說我沒有錢.
應該去不了.
老師跟我說.
祈禱吧.
好,兩個人回去祈禱.
我沒有跟別人分享.
沒有人知道我有這個想法.
只有我的上師.
我的牧師知道.
說要去探索北美.
誰知道有一天他打電話給我.
他說.
有位叔叔想支持你們去探索.
然後叔叔就說.

$^{521}$我們就去咖啡店吧.
你知道我八十通沒什麼錢.
付款.
去咖啡店.
在金鐘喝杯咖啡可能也要一百多元.
我兌完八十通已經付款了.
然後叔叔跟我說.
你不喝嗎.
我說我可以了.
你們喝就行了.
就把一百多元的咖啡給了叔叔.
然後叔叔跟我說.
你去探索吧.
我支持你.
然後他說很貴.
要這麼多錢.
然後我說.
OK,我幫你.
在我的經驗來說.
有一個陌生人主動說要把錢放進我的戶口.
如果你用今天的邏輯去理解這件事.
有一個陌生人.
主動說要把錢給你.
你會覺得會不會是KK園區騙你的.
完全不合邏輯.
不過上帝真的做的事完全超越了我的想像.
過了一段時間.
他真的把一大筆錢放進了我的戶口.
然後我們就很順利地申請到簽證.
過來這邊探索.
才認識到教授.
才有我們之後2020年讀書的時刻.
很明顯這一次又是再一次反轉.
我以為我順序之後.
已經接受了耶穌.
已經接受了耶穌這款生命麵包.
但他再一次.
在我們的角度裡面.
用我們覺得不合邏輯的邏輯.
的概念.

$^{561}$去叫我們踏出去.
重新建立我們對他的信心和信仰.
上次我也提過.
我太太.
其實我也是半個長洲人.
住在長洲.
長洲有很多親戚.
Alice的姨媽.
其實一直都很拒絕相信基督.
她是相信日本佛的.
也一直阻止姨丈去上教會.
誰知道有一年.
幾年前.
姨丈就發現自己有癌症.
而姨媽就有腦退化的跡象.
在這個屋漏堅逢夜雨的時刻.
我們真的完全沒有想過.
見證著神.
在他們兩個的生命裡面.
的一個很大的作為.
既然他們兩個先後都順著.
姨媽開始腦退化.
記憶衰退.
而姨丈又漸漸變差.
他在最後的日子裡.
基本上都在醫院度過.
姨媽有時都會忘記了姨丈.
甚至會當姨丈是她自己的爸爸.
姨丈病危的時候.
岳父岳母和教會的同工.
都去探望他.
和他分享信仰.
而姨丈在人生最後的兩個星期裡面.
他竟然和岳母分享.
他說他見到耶穌.
他見到耶穌.
他說他相信主.
他立刻就有水利.
教會立刻就有水利.
見證著他.

$^{601}$見證他信耶穌.
而姨媽因為什麼.
因為腦退化.
他沒有能力反對.
他都不知道姨丈已經受了水利.
更加想不到的是.
姨媽因為腦退化.
他的子女將他的專一本佛搬走了.
他都不覺.
最後他竟然願意信主.
而且非常喜樂.
當我們今年新年的時候.
我們在視像打給他.
他不斷跟我們說.
我愛你.
很喜樂 很甜蜜.
姨媽姨丈經歷的一切.
是我們完全想像不到的.
這件事這樣發生.
是完全被我們自己的邏輯.
被我們的思考困死.
我們覺得其實我們應該這樣.
應該這樣.
那些事情才會合邏輯.
才會符合我們期望.
但上天做的歧視.
是再一次令我們跌眼鏡.
所以今天想鼓勵大家.
培養一顆柔軟的心去跟隨耶穌的工作.
為什麼說是培養呢.
我們成年人其實對任何的事情.
都有經驗.
會慢慢累積.
在屬靈的事情都是這樣.
經驗當然是好事.
它幫助我們快速去面對一些危機.
甚至給我們有角度.
當我們有事情要處理的時候.
我們很快就可以知道.
怎樣最有效.

$^{641}$最精準地去處理那件事.
但是在聖經裡面的耶穌.
有時候也會打破我們的想法.
就是我們以為邏輯是這樣.
這樣做才能成就救恩.
但是如果我們要合邏輯地.
去接受耶穌這個生命麵包的時候.
我們很可能會像門徒那樣.
是不是.
會離開祂.
甚至出賣祂.
因為整件事是不合邏輯的.
有什麼理由.
你以前叫我不要吃血.
吃血還會被人剪除.
還會跟上帝隔絕.
現在你叫我吃你.
你叫我吃你的血.
吃完血才有救恩.
才有永遠的生命.
完全不合邏輯.
所以我們要慢慢培養一顆柔軟的心.
去跟隨耶穌的工作.
對自己的生命和對別人也是這樣.
最後一個例子我們就結束了.
我們現在住在一對教會夫婦的家裡.
其中她有一個小朋友.
有一晚我們叫她.
不如洗點水果給我們吃.
洗點葡萄.
對於成年人.
葡萄很簡單.
你會怎麼做.
教她剪下來 撕下來.
用水洗乾淨 清除框架.
就OK了.
誰知道她拿了一個透明袋.
她一直在搖.
一直搖.
然後我們幾個大人就說.

$^{681}$不行 把葡萄弄爛了.
不行的.
行不行的.
想停止她.
不夠20秒.
所有葡萄都掉出來了.
全部都掉出來了.
沒有爛的.
完全沒有爛.
我們成年人洗葡萄的方法很簡單.
就是弄完之後洗完.
我們經常告訴她.
我們有很好的方法.
很好的方向.
甚至會將我們框在某些邏輯底下.
但有時候我們會被自己的想法.
或者我們過往的經驗蒙蔽了.
以為這樣做才是最安全.
實際上也是的.
我們已經預知要提自己的錦骰.
不用冒險的.
但是接受耶穌這個生命麵包就要.
所以我們要培養柔軟的心.
就是Stay open.
嘗試去做一些新的事.
看看上帝在這個工作.
這個世代.
有什麼新的工作.
最後.
為什麼我這個PowerPoint有這麼多不同的包呢.
可能會令到大家很餓.
不好意思.
我這樣做是希望大家.
當大家享受在上帝所賜的食物的時候.
而是記得.
祂是生命的主.
還要培養一顆柔軟的心.
去跟隨祂的工作.
對人.
對自己.

$^{721}$自己也是這樣.
祝福大家.
\newpage



\section{5:43-48}
\label{sec:95H_Df1nGvs}
\textbf{Righteousness as Redefined by Jesus: Love Your Enemies (Matthew 5\_43-48) - Chee Yie Jong}
\newline
\newline
連結: \href{https://youtube.com/watch?v=95H_Df1nGvs}{\texttt{https://youtube.com/watch?v=95H\_Df1nGvs}} ~~~~ 語音日期: 2025-04-10
\newline
\newline
\hyperref[sec:ETczAl2v1zo]{\small{< < < PREV SERMON < < <}}
~
\hyperref[sec:index_chronic]{\small{[返順時目]}}
~
\hyperref[sec:index_scriptual]{\small{[返順卷目]}}
~
\hyperref[sec:FaEgprmwTI0]{\small{> > > NEXT SERMON > > >}}
\newline
\newline
$^{1}$Good morning, UV..
It is good to be able to worship with you this morning..
And as Pastor Mark say, it's good to see some old faces back.
and also some new faces as well this morning..
So welcome and thank you for joining us..
So what are some of the hardest things.
that you have been asked to do before?.
Or if you think about it,.
what are some of the most challenging tasks.
that you have gone through?.
As I reflect on my personal experience,.
I think they usually fall under a few categories..
Some of these things are hard.
because they test our human limits.
and they push us to our very edge,.
intellectually, emotionally, or physically..
This October, I will be running.
the Chicago Marathon as well..
And even though I'm looking forward to it,.
there's also part of me that is somewhat nervous about it.
because I know that in order for me.
to finish these 26.2 miles,.
it will push me to my very limit..
And on the other hand,.
some of the things are challenging.
because they require us to face our own fears..
Think of extreme activities like skydiving,.
bungee jumping,.
or even small things like catching a cockroach or a spider..
Right?.
And then there are certain things.
that are hard and challenging.
because of what they require from us.
to sacrifice our own preference..
So the passage that we will be looking at today.
contains one of the most radical teachings.
and challenging commands from Jesus himself..
It is difficult and challenging.
not because it pushes to our limits.
or it requires us to face our fears,.

$^{41}$but it is difficult because in order to do.
what Jesus has told us to do here,.
it requires us to completely deny ourself.
and our own self-interest..
To put it another way,.
if we're not willing to give up our own pride,.
our preference, our convenience,.
our comfort, or the illusion of control that we have,.
this command from Jesus will almost be impossible.
for us to live out..
So if you have your Bible with you,.
let's turn with me to Matthew chapter five..
We're gonna look from verse 43 to 48 today..
So Matthew chapter five, 43 to 48..
But before we dive into God's word,.
let's pray that his spirit will be with us.
as we open up his words..
Father, we thank you that we're able.
to come into your house this morning.
to lift up our praise to worship you..
And as we open up your words now,.
I just ask that your Holy Spirit will be with us..
I ask that your words.
will not just be knowledge for us,.
but it will be truth that will be able.
to transform our life,.
so that we're able to live in a way.
that brings honors and glory to you..
Amen..
So, Matthew chapter five, verse 43..
You have heard that it was said,.
love your neighbor and hate your enemy..
But I tell you, love your enemies.
and pray for those who persecute you,.
that you may be children of your Father in heaven..
He causes his son to rise on the evil.
and the good,.
and sends rain to the righteous and the unrighteous..
If you love those who love you,.
what reward will you get?.

$^{81}$Are not even the tax collectors doing that?.
And if you greet only your own people,.
what are you doing more than others?.
Do not even the pagans do that..
Be perfect, therefore, as your heavenly Father is perfect..
This is the word of God..
So today, we actually come to the end of Matthew chapter five.
and our passage today serve as sort of a conclusion.
to this first major sections in the Sermon on the Mount,.
where we look at the relationship of the Old Testament law.
to the new covenant of Jesus..
And if you can recall, we begin this section.
by looking at what a truly blessed life is..
And we come to this truth that.
because Jesus has fulfilled the law and the prophet,.
to enter into his kingdom.
and to experience this truly blessed life,.
it doesn't depend on our own righteousness,.
but the righteousness of Jesus alone..
And right after that,.
we look at a series of case studies on how Jesus,.
now being the definite interpreter of the law,.
redefines what true righteousness mean.
in his kingdom of heaven,.
in contrast to the righteousness according to the Pharisees,.
the teachers of the laws,.
or the Jewish religion or Jewish tradition..
In the past couple months,.
we have looked at anger and murder,.
lust and adultery,.
divorce,.
making oath,.
and also the law of retribution..
And how each and every time,.
Jesus shows us his better way and his greater righteousness.
that points back to God's original intention.
behind this law..
If you remember, he always introduced them by saying,.
"You have heard that it was said, but I tell you this.".
So our passage this morning is actually the sixth.

$^{121}$and also the final examples.
that will pretty much sum up these entire sections..
So for my sermon this morning,.
we will first look at the command from Jesus,.
and then we'll look at the reason behind his command,.
and then we will end with the ultimate goal..
So the command, the reason, and the goal..
So let's look at the first part, the command..
Verse 43..
You have heard that it was said,.
love your neighbor and hate your enemy..
But I tell you, love your enemies.
and pray for those who persecute you..
So once again, Jesus begins his teaching with this formula..
You have heard that it was said..
And the first half of this quote, to love your neighbor,.
actually comes from Leviticus 19.18..
And alongside Deuteronomy 6.5, which is the Shema,.
to love your Lord your God with all your heart,.
with all your soul, with all your strength..
Jesus, he sums up the entire Old Testament law.
with these two commandments later in the book of Matthew..
As for the second half of this quote, hate your enemies,.
the origin of this second half is a little more ambiguous..
In the book of Deuteronomy, as the people of Israel.
are about to enter into the promised land,.
the Lord gave them this instruction.
to drive out all the enemies or the inhabitants.
of the land of Canaan,.
and pretty much devote these people to destruction..
And so this quote, to hate your enemy,.
may be an inference that is drawn.
from those passage in Deuteronomy..
On the other hand, this idea of hating your enemy.
may also come as a logical conclusion.
from the command to love your neighbor..
In Leviticus 19, when the Lord commanded the Israelites.
to love their neighbor, the neighbor there.
was clearly referring to a fellow Israelites..
And so when it comes to Jesus' day,.

$^{161}$some of the Jewish people, they take this command.
and draw the conclusion that, okay,.
if we are only called to love our neighbors,.
that probably means we can hate our enemies..
So do you see how twisted our human heart can be.
when we take a positive command from God.
to love our neighbor,.
and pretty much turn it into a license to hate?.
So how does Jesus address this false narrative.
and misunderstanding?.
Verse 44, "But I tell you, love your enemies.
"and pray for those who persecute you.".
So Jesus couldn't be any more clearer here,.
but to most of Jesus' Jewish audience at that time,.
this probably comes as quite a shocking command..
And even for most of us today,.
I would assume that this is still quite a radical command.
and a tall task for us, right,.
to actually love our enemies..
And one thing to clarify here,.
when we hear the word love,.
what are some of the things that come to your mind?.
I feel like in our current day and age,.
this word love has become like a catch-all phrase.
for all sorts of different ideas and different feelings.
and different thoughts, right?.
You might say, "Oh, I love pineapple on my pizza.".
Or you say, "Oh, I love the Chicago Bears..
"I hope they do well this season.".
Or you might say to your significant others, "I love you.".
Or you say to God, "God, I love you.".
Do you notice the issues here?.
The word love has been used.
in such a wide range of applications.
that it has lost its significance.
when it comes to the way that Jesus.
and the Bible talks about love..
And so when Jesus give us the command to love your enemy,.
what kind of love is he referring to?.
Is it the kind of love.

$^{201}$that you like to have pineapple on your pizza.
or the love that you have for your significant others?.
Well, love in the sense that Jesus talks about.
is not just about having a passion for something.
or a deep affection for someone.
because those feelings will fade over time, right?.
Yes, it is true that love includes all of those feelings,.
but more importantly, love is about action and commitment..
And love entails accountability and responsibility..
So let me read a passage from 1 John.
that will help us better understand what love is.
according to Jesus..
So 1 John 3, 16..
"This is how we know what love is..
"Jesus Christ laid down his life for us,.
"and we ought to lay down our lives.
"for our brothers and sisters..
"If anyone has material possessions.
"and sees a brother or sister in need.
"but has no pity on them,.
"how can the love of God be in that person?.
"Dear children, let us not love with words or speech,.
"but with actions and in truth.".
So church, love is more than just the feelings..
I hope that doesn't come as a surprise for you,.
but it is also about actions and truth..
This kind of love is generous,.
yet costly at the same time.
because it requires us to lay down our own kingdom.
and commit to seeking the good of others..
In other words, this is a love.
that is shaped after Jesus and his cross..
This is a kind of cruciform love..
And this cruciform love model after Jesus.
is being lived out when we, as his followers,.
are able to actively love our enemies.
and genuinely pray for those who are hostile.
or unkind to us..
And in the Greek grammar,.
both of the verb to love and to pray.

$^{241}$stress that this command is to be lived out now.
and not just towards a singular enemy,.
but to all who we consider as enemies..
So is loving our enemies easy?.
Of course not..
And is responding to kindness.
towards those who are hostile towards us.
a challenging task?.
Of course it is..
This is because in our fallen and unglorified state,.
we all have an inward curse towards ourself,.
as St. Augustine has pointed out..
That is the tendency to always focus inward on ourself,.
on our needs, on our desires,.
rather than on God and other people..
So of course, to love our enemies.
and pray for those who persecute us.
is not going to come easy.
because it requires us to act in a way.
that is contrary to our human nature, corrupted by sin..
But why does Jesus command us to do so right here?.
And even more, to use this command.
as a summary to his entire teaching.
regarding the better righteousness.
that is required to enter into his kingdom..
So let us look at the why behind this command now..
(silence).
Love your enemies and pray for those who persecute you.
in order that you may be children of your Father in heaven..
He causes his son to rise on the evil and the good.
and sends rain on the righteous and the unrighteous..
So what is the reason behind Jesus' command.
to love our enemies?.
It is to imitate the character of our heavenly Father,.
who is gracious and compassionate,.
slow to anger and abounding in steadfast and loyal love..
To love our enemies and pray for those who persecute us.
shows that we don't belong to this world,.
but we belong to another kingdom that is above..
And moreover, God himself also loves impartially.

$^{281}$by bestowing his blessings on both the evil and the good,.
the righteous and the unrighteous..
So yes, while loving our enemies is not easy,.
we are able to love radically because God loved us first,.
regardless of our actions..
Paul in Romans 5 said this,.
"Very rarely will anyone die for a righteous person,.
"though for a good person,.
"someone might possibly dare to die..
"But God demonstrate his own love for us in this..
"While we were still sinners, Christ died for us.".
Did you catch that?.
While we were still sinners and enemies of God,.
Jesus Christ, the son of God, died for all of us..
And so as God's people, we don't have the excuse.
to not love in the way that God loves us.
and loves in the way that reflects our heavenly father..
Think of the most evil person.
that you can think of in your life,.
or someone that you would consider as your worst enemies..
And guess what?.
He or she is still made in the image of God.
and deeply loved by God, our heavenly father..
So again, as God's people,.
we are to reflect this loving character of our father,.
even towards those that we rather.
would have nothing to do with..
Right?.
And then Jesus tries to bring his point home.
by using two rhetorical questions in verse 46 and 47..
So if you are not familiar.
with what rhetorical questions are,.
these are questions posed usually to provoke thoughts.
rather than to actually get a response..
And rhetorical questions,.
they have the functions of drawing attention to something.
or highlighting a point..
So with these two rhetorical questions,.
Jesus is trying to draw his audience to his point..
So verse 46,.

$^{321}$if you love those who love you,.
what reward will you get?.
Are not even the tax collectors doing that?.
So a little context about the tax collectors..
Tax collectors during Jesus' times.
are being viewed as traitors,.
and they're people that are being despised by the Jews..
So the tax collectors,.
they're usually Jewish people.
who works for the Roman government.
to collect taxes from their own people..
And the common practice for these tax collectors.
is they're actually allowed to collect.
more than that is required,.
and then they can just pocket.
the rest of the money for themself..
And so these Jewish tax collectors,.
they pretty much build their wealth.
by quote-unquote daylight robbing their own people,.
and that is why they're being despised by the Jews as well..
And so what Jesus is trying to say here is this..
Even the tax collectors whom all of you despise.
can reciprocate love to their friends..
There's nothing to brag about it.
if you, as my followers,.
are able to love those who only love you back..
And then Jesus continues, verse 47..
And if you greet only your own people,.
what are you doing more than the others?.
Do not even the pagans do that?.
And so here from 46 to 47,.
Jesus moves from within the Jewish community.
to those who are outside the Jewish community now,.
and he used the Gentiles or the pagans as the examples..
And greetings in that culture.
is closely tied to someone's social status as well..
So especially for those who hold.
a higher social status in a society,.
they like to exchange greetings with other people.
so that they can hear other people calling them their title.

$^{361}$or kind of like highlighting their status, right?.
Again, the point that Jesus is trying to make is this..
If the pagans or the Gentiles.
in which you Jewish people consider as inferior to you all,.
they are able to reciprocate the greetings.
that come from their own people..
You all, as my disciples,.
you should be able to do much better..
You should be able to show the same respect and honor.
to other people, including your enemies..
(smacks lips).
So at this point,.
we have gone through all six of these case studies.
or examples on what true righteousness is.
according to Jesus, right?.
And each and every time we see how Jesus redefines.
what true righteousness is compared to the law.
on how it was being understood by the Jewish people..
And so what is the purpose of Jesus' teachings here.
in revealing the heart of God.
behind all these laws and commandment?.
What is the ultimate goal that Jesus has?.
Let us look at verse 48 here..
"Be perfect therefore, as your heavenly Father is perfect.".
Here, we see Jesus refers to the heavenly Father once again.
as the model for all of us to follow,.
recalling His command earlier to love our enemies.
and pray for those who persecute us.
in order that we can be children of our heavenly Father..
On the surface level, this call to be perfect.
sounds kind of like a commandment that is almost impossible.
for us to achieve on this side of heaven, right?.
However, the word perfect here,.
instead of referring to absolute sinlessness.
or perfect in a way that is like blameless and never sin,.
the Greek word for perfect here, teleos,.
actually reflects more of the idea.
of wholeness or completeness..
And besides, the idea of perfection.
that Jesus talks about here.

$^{401}$is actually more functional as well,.
rather than positional..
What does that mean?.
So for example, if a thing is able to fulfill.
its intended purpose or function.
or reach its full development,.
it can be considered as perfect or teleos in that sense..
So if a puppy is able to grow to its full-grown stature,.
it is considered as teleos..
Or if this wireless microphone is able to function.
in the way that it is supposed to function.
and that it is created to function,.
it can be considered as teleos and perfect as well..
And when it comes to us as human,.
this begs the question,.
what purpose are we as humans created for?.
As you all know, we as human,.
we are made in the image of God..
And as image bearer and ambassador of God's kingdom,.
we are to reflect the character of our heavenly father.
and represent him in this world,.
just as a diplomatic ambassador.
will represent his or her country, right?.
And as we have seen in our passage,.
our heavenly father is one who loves impartially.
by making his son rise to both the good and the evil.
and causing rain to fall on both the righteous.
and the unrighteous..
And in Romans 5, we see that he is a God.
who loves radically and generously.
by sending his own son to die in our place for our good..
All this while we're still sinners and enemies of the cross..
So when we are able to reproduce this kind of life.
that fully exemplifies the character of our heavenly father,.
one that is marked by radical,.
self-sacrificing, cruciform love,.
we fulfill the purpose in which we are made for..
And therefore, we can be considered perfect.
or teleos in that sense..
So to be perfect as our heavenly father is perfect.

$^{441}$means that conformity to the character of our God.
and the image of Christ should be the goal of our words,.
the goal of our actions, our thoughts,.
and the way we relate to those that are around us,.
friends and foes alike..
So if you're taking notes, you can write this down..
The point of all of this is as the children of God,.
we are called to a lifelong journey.
of becoming more and more like our heavenly father..
Notice that this is not an instantaneous thing,.
that when we put our faith in Jesus,.
we instantly become like our heavenly father..
But this is a lifelong journey..
And this includes the way we relate to other people.
made in his image..
If you recall our passage on anger and contempt,.
as well as on adultery and lust..
And this also includes the way we view marriage and divorce,.
the way we use our words, whether we speak truthfully,.
the way we respond to evil,.
and ultimately, in the way we love other people,.
especially those that we consider as our enemies..
And church, this is the surpassing righteousness.
that Jesus is referring to,.
one that is in perfect alignment.
with the heart of our heavenly father..
So again, as the children of God,.
our goal on this side of heaven.
is to continually work towards becoming.
more and more like our heavenly father..
So I just want to end with this story..
So last Sunday, after our general assembly,.
if you were here, you might remember.
that there was just a huge thunderstorm, right?.
So after that, I offered to drive.
one of the members from UV home..
But sometime between dropping him off.
and getting home myself,.
I think my car probably ran over some nail at some point..
Because later that evening, when I was driving again,.

$^{481}$I started to notice some weird noise.
coming out from behind my car..
And when I stopped and take a look,.
this is what I see..
I had a flat tire..
And at that point, I was only about 10 minutes.
away from home, so I'm like,.
okay, I'm not gonna call a tow truck or trip away.
because it might take them a while..
And I'm so close to home,.
I'm just gonna drive home with this flat tire..
But as I was driving, the noise just keep getting.
weirder and weirder and louder and louder..
And me personally, I also became more and more frustrated..
And in my mind, I was like,.
whoever is responsible for that nail.
like dropping on the street, like screw you..
I hope you had a bad day..
Monday morning came, I brought my car to the shop.
to get it fixed..
And the first thing in the morning,.
like the very first thing in the morning,.
so that I can pretty much spend the rest of my day.
to have some like focus time to prepare.
for my sermon today..
But as I read through this passage from Matthew five,.
with the frustration that I had.
towards the flat tire situation still fresh in my mind,.
this word really just stood out to me,.
to love your enemies and pray for those who persecute you..
And even though this is just such a small incident,.
I realized that for me to love and pray for whoever.
that is responsible for that nail,.
it doesn't come natural to me..
I have to kind of like, in some ways, like humble myself.
and come before the Lord..
It's like, okay, Lord, like that is not the way.
I should respond yesterday..
And it takes a lot of intentionality out of me as well..
So this call from Jesus to be perfect.

$^{521}$as our heavenly father is perfect,.
requires a lot of intentionality,.
as well as discipline..
To understand our father's heart,.
we need to spend time with him..
We need to immerse ourself in his words..
And to become more like our father in heaven,.
we need to allow the Holy Spirit to work in our hearts.
and transform us from the inside out..
The righteousness of Jesus.
is not about outward behavioral management,.
but it comes from a heart that is truly transformed.
by Jesus and his gospel..
And this is a process that will last until the day.
that we're able to meet Jesus face-to-face..
So it requires a posture of constant obedience.
and a process of constant surrendering.
in our day-to-day life..
So I want to leave with these questions.
for us to reflect on..
Who are some people whom you would consider.
as your enemies?.
Or who are some people who have wronged you?.
And in light of our passage this morning,.
how can you show the love of Christ to them.
in a way that reflect the love you have received.
from our heavenly father?.
So I'll have Amanda come up to.
play some background music.
and just take some time to reflect on it,.
and I will end this with a prayer in a couple of minutes..
(gentle music).
(gentle music continues).
(baby crying).
(gentle music continues).
(thunder rumbling).
(gentle music continues).
Heavenly father,.
we thank you that while we were still sinners,.
while we were still enemies,.

$^{561}$you sent your son, Jesus Christ,.
to die on our behalf for our good..
Help us, father, to be people.
that reflects your loving character to those around us..
Help us to become people that love in a way.
that we're willing to sacrifice our own preference,.
our own comforts, our own conveniences..
Help us to love in a way that when the world see us,.
that they will know that we belong to you,.
that we can bring glory and honor to your name..
So again, father, thank you for your word this morning..
We pray this in your son's name, amen..
All right, thank you for joining us this morning..
We will have some refreshment outside,.
so I encourage you all to stay after our service.
and just hang out, get to know some new people,.
and yeah, enjoy fellowship with one another..
All right, go in peace, UV, your love..
\newpage



\section{約伯記}
\label{sec:FaEgprmwTI0}
\textbf{在問題中遇見神 (約伯記) - 余淇煥傳道}
\newline
\newline
連結: \href{https://youtube.com/watch?v=FaEgprmwTI0}{\texttt{https://youtube.com/watch?v=FaEgprmwTI0}} ~~~~ 語音日期: 2025-04-10
\newline
\newline
\hyperref[sec:95H_Df1nGvs]{\small{< < < PREV SERMON < < <}}
~
\hyperref[sec:index_chronic]{\small{[返順時目]}}
~
\hyperref[sec:index_scriptual]{\small{[返順卷目]}}
~
\hyperref[sec:i_UUX_QnitY]{\small{> > > NEXT SERMON > > >}}
\newline
\newline
$^{1}$好,弟兄姊妹早安.
話這麼快.
我的UV實習已經到了尾聲.
4月開始,經過學校對我的評估後.
我會正式成為一個.
Special Director,靈性指導.
這段日子對我來說非常寶貴.
有很多好湯喝,有人請我吃飯.
跟我分享在你們生命中的掙扎,成長.
信仰上的瓶頸位置.
甚至是受傷的地方.
靈性指導是跟人深入同行.
探討生命的問題.
以及一起學習聆聽自己.
以及聆聽神的過程.
所以我經常會用問題帶人反思.
來見我的弟兄姊妹.
都會帶來不同的問題.
以下我所說的東西.
全都會經過當事人批准.
我才會在台上說.
其中一個問題最多人問的就是.
我應該怎麼做.
會問我要怎麼做.
亦有一些很有趣的問題.
我經歷過,就是會問.
Bella,今天神有甚麼事.
透過你跟我說嗎.
或者神有否給你一段經文.
那一刻我有點感覺自己像神婆.
「你回去用藥褔飲」.
「三張十六折沖水飲」.
「藥到病除」.
「靈的話,記得出門後多添點香油」.
我想說,如果人生的問題.
可以用金句「回去煲水」.
「半碗煲為半桶水」.
喝了它就能解決.
你說人生多好,對吧.
但很可惜.

$^{41}$很多時候人生的問題.
遇到時都很不容易有答案.
特別是苦難的問題.
人生裡面,在面對苦難時.
很多時候都很想找答案.
是一件很自然的事.
今天我想帶大家看《藥伯記》.
藥伯這個苦難的人生.
他卻面對著上帝對他發問的問題.
而遇見上帝.
我想說,《藥伯記》….
所以我今天說的題目.
就是「在問題中遇見神」.
《藥伯記》就像一齣舞台劇.
一開始有一個旁白揭開他的序幕.
話說在地上有一個人,他叫藥伯.
是一個完全正直的人.
遠離惡事,敬畏神的大好人.
他很有錢,家裡很幸福.
同時,在天上又發生了一幕和一幕.
就是上帝和眾天使在開會.
開甚麼會呢.
在這個會議裡.
撒旦,即是控訴者.
向神提出質疑.
認為藥伯是因為你對他那麼好.
給他那麼多福氣,他才敬畏你.
你試試看,他沒有福氣.
你看看他還會否敬畏你.
而上帝就容許撒旦去試煉藥伯.
試煉的意思是將苦難帶給藥伯.
在很短的時間裡,藥伯失去了財產.
他的子女一個又一個死去.
連自己的身體也患了很重的病.
全身都生滿了毒瘡.
最後藥伯坐在灰燼中,苦不堪言.
經過這麼大的創傷後.
他只說了兩句話.
他說「赤身出母胎也必赤身歸回」.
「賞賜的是耶和華」.

$^{81}$「收取的也是耶和華」.
「耶和華的名是應當稱頌的」.
而第二句是甚麼意思.
就是當他老婆在他落難時.
說他「你看,你現在弄成這樣」.
「你相信你的神吧」.
藥伯就跟他說.
「難道我們在神手中得福」.
「就不用受禍嗎」.
我聽這句話時覺得有點奇怪.
然後仔細想想,其實也很真實.
就像一個人面對極大創傷後.
其實他馬上未必懂得哭.
甚至可能可以很理性,很抑壓.
還可以說一些很高層次的話題.
但可能心底裡承受不了這麼大的衝擊.
在他說完這些高層次.
還讚美神的說話後.
他就完全進入沉默當中.
這就是藥伯真正進入到.
他人生最黑暗的幽谷中.
在這個時候,他三個朋友就出場了.
這三個朋友分別是.
以理法,彼勒達和所法.
他們看到藥伯這麼慘.
沒有一個人看到他可以說話.
因為覺得太可憐.
他們只陪著他七天七夜.
跟他一起大哭.
藥伯終於忍不住哭喊他內心的痛苦.
整個《藥伯記》的三章.
都是藥伯一首哀歌.
他抒發他生命當中的絕望.
自己存在的困惑.
我嘗試用節錄和口語表達一下.
他說「希望我出生那天」.
「消失不存在就好了」.
「為何我不一出生就死了」.
「為何我不是一離開母胎就斷氣」.
「我沒有一天安寧」.

$^{121}$「我沒有一晚睡得著」.
「我連一口氣都沒有空間喘息」.
災難就撲過來.
如果我們用現代心理學.
去讀入經文的時候.
你覺得醫生會診斷它為甚麼.
對,很明顯.
抑鬱症,持續的憂鬱,悲傷,空虛.
以及失去興趣.
感覺到自己很沒有價值.
還有死亡和自殺的念頭.
我們想也想到.
無論他睡覺還是吃肉.
都是180度轉變.
都是很差.
這三位朋友已經陪伴了約伯七天七夜.
其實他的同理心開始慢慢用盡.
以致你回去看約伯記.
足足有28章的經文.
去記載這三位朋友.
輾轉回應約伯.
他們的安慰.
彷彿變成了一場神學辯論和教化.
在他第一輪的安慰中.
我濃縮了,就這樣告訴大家.
伊利法就這樣說.
他說神是賞善罰惡的.
他鼓勵約伯悔改.
他說「你想想」.
「你有否見過無辜的人無緣無故死亡」.
「正直的人有誰會突然被剪除」.
然後下省了三萬至一大堆理論.
彼納達就說.
「神是公義的」.
「如果你沒有罪」.
「神一定會祝福你的」.
「神絕對不會扔下你不管」.
「神絕對不會扔下那些原傳人」.
「也不會輔助那些惡人」.
索法這個朋友同樣廢盡他的純善.

$^{161}$暗示約伯要悔改.
他說惡人的眼睛會瞎.
他們逃不掉.
他們唯一的盼望就是等死.
他說約伯應該更謙卑,更自卑.
神既然審判得,你一定有錯.
我現在這麼精簡地說.
你已經覺得很煩,對吧.
約伯也會感覺到.
在第一輪對話之後.
約伯只想跟神有些對話.
其實他不想跟朋友爭論甚麼.
還未完結.
第二輪繼續安慰他.
這三位朋友照樣重複.
他們如何對上帝公義的理解.
語氣比之前更激烈.
耐性更少.
由好言相勸.
到斥責約伯自以為義.
口出狂言,藐視神.
到了第三輪.
大家除了重複之前的斥責之外.
甚至開始編造一些可能犯罪的罪行.
去指控約伯.
「今天你遭受這麼大的患難」.
「其實可能你也該死」.
這三位朋友一直都認為.
苦難是因為人犯了罪.
神藉著苦難去懲罰人.
他們嘗試用因果報應的神學.
為約伯解釋他的遭遇.
但反而令約伯更痛苦.
這些回應看似正在對應約伯的處境.
跟他說話.
好像在疏導他.
甚至這三位朋友可能覺得自己.
是在安慰約伯.
但事實上,這些都是疑似傷害.
因為他們否定了約伯的無辜和痛苦.

$^{201}$約伯越聽越混亂,越聽越辛苦.
約伯堅持自己沒有錯.
不斷質問神.
「我究竟做錯了甚麼」.
「為何要這樣對我」.
我想說這個對質,質問.
其實不僅是一個自辯.
其實是約伯對神公義的渴望.
以及反映他對神的信任.
以至他可以赤露躺開自己.
向神哀哭,申冤.
而鏡頭一轉.
這三位朋友說完了.
到了一個年輕人身上.
這個年輕人叫以地胡.
雖然他的語氣比較溫和.
也沒有將約伯所受的苦難.
歸咎於他的犯罪,得罪神.
但同樣地,他仍在解釋.
神一定是公義的.
他詮釋苦難是一個警告.
是一個提醒,或是一個管教.
即使這位年輕人滿腔熱誠地解釋.
上帝如何用風雲雷電去掌管這個世界.
他解釋得很好.
但他這堆說話.
都沒有真正觸碰到約伯靈魂中.
最深刻,最核心的渴望.
就是他很想跟神面對面對話.
我曾經聽過一些說話.
在一些基督教聚會中.
不知道大家有沒有聽過.
我給大家看看.
「病不得醫治」.
是因為我們祈禱不夠信心.
我們不夠熱切.
甚至可能是因為你沒有按食祈禱.
或者當天災發生的時候.
有些人會說.
「天災發生是因為神管教的地方」.

$^{241}$「叫人去悔改」.
或者「你所受的痛苦」.
「可能是神給你的功課」.
「這個苦難可能是神對你生命的修剪」.
「所以是神為你好」等等.
有聽過這些說話嗎.
我想說.
這些說話對受苦者來說.
比粗口更難聽.
「將信仰簡化」.
「太早對一個人所經歷的一切」.
「給他一個答案」.
「給他一個很簡單的答案」.
其實對受苦者來說.
是造成二次傷害.
甚至有時候我們會這樣想.
例如我們經歷了一些事情.
我們很快會覺得.
「神是懲罰我的」.
「我一定得罪了神」.
有時候我們會想.
每一件事都有一條公式.
讓我們可以解釋就最安全了.
包括我們的信仰.
最容易處理就最好了.
但原來當人面對痛苦的經歷.
如果我們輕易給予一個答案和結論.
其實某程度上.
我們在輕視人生中的複雜性.
我們在輕視神的奧秘和主權.
我們甚至迴避了.
面對痛苦本身的不確定性.
還有奧秘.
在約伯和這群朋友折騰了很久之後.
終於….
終於上帝在旋風之中開聲.
他親自回應約伯.
但有趣的是.
你以為終極的主角都出現了.
謎底就會解開,對吧.

$^{281}$你以為神都出現了.
還不是一堆答案給約伯.
解釋給他聽.
從頭到尾為何會這樣發生.
不是.
神不僅沒有給人所期望的答案.
讓約伯聽到.
而是他發出一連串震撼人心的問題.
神用了很多篇幅去說.
「在我創造這個世界的時候」.
「在打地基的時候,你在哪裡」.
「你試過叫清晨來臨嗎」.
「你試過命令神光照遍地面嗎」.
「海水出世的時候」.
「我有叫它們亂湧嗎」.
「我有讓這些海水設立界限嗎」.
「你有看過山羊生小孩嗎」.
「你有知道鷹是怎麼飛的嗎」.
「你又知不知道河馬,鱷魚」.
「都是我造出來的嗎」.
「你能控制得了嗎」.
好像與他無關.
所有問題都好像與他無關.
但不是.
神用一連串的問題.
問了約伯很多關於宇宙秩序的問題.
並且描述了山羊是怎麼吃草.
鹿是怎麼生下一代.
神彷彿帶著約伯遊覽了一次宇宙.
神讓約伯體會到這個宇宙.
原來有這麼多事情.
神仍然掌控不了.
又如何能完全明白上帝的心意.
這一連串的問題.
不是要打壓約伯.
不是要大石砸死蟹.
而是讓約伯自己體會到.
人的渺小和有限.
以及神的無限和主權.
神的出現沒有解決約伯的苦難.

$^{321}$神的出現沒有直接回答約伯的質問.
但神的出現.
就是對約伯一個最親密的回應.
當人面對苦難的時候.
很多時候都會出現很多不同的情緒.
不知道你在人生最困難的時候.
或者生命中遇到一些很苦的狀況.
處境的時候.
你會否有一種無力感.
憤怒,悲傷,甚至是絕望.
約伯一開始質問神,甚至控訴神.
但神卻沒有抑壓約伯,沒有否定他.
因為神容得下這些情緒.
反而用問題邀請約伯進入.
和上帝更深的相遇中.
約伯的呼喊,甚至和神的對質.
是靈性中很深層的哀歌.
不是質問神就代表他的信仰等於失敗.
就等於他不信神,質疑神.
不是的.
有一位天主教神學家古鐵雷斯.
在其中一本書的著作中說.
質問不代表背叛.
反而是信仰最真誠的呼喊.
約伯代表著無辜受害者的吶喊.
是窮人和被迫迫者的呼聲.
在我和很多人做靈性指導的時候.
有時候我會邀請人.
將自己內心的感受跟神直接說.
但每當人願意….
在華人群體中,其實很少說感受.
很多時候都是說那些事情.
但當我一直引導他去說他的感受時.
我不僅引導他去說.
而是我帶他去將這些真實的感受跟神說.
無論是憤怒,哀傷,無助.
無論是深層的孤單感.
還是絕望.
甚至是你生神氣的時候.
我看到眼淚一直流.

$^{361}$我看到人的心向神打開.
我也看到神同樣靠近人的心.
掌管宇宙萬物的神.
祂盛載得起我和你.
盛載得起我和你.
生命當中的痛苦和掙扎.
盛載得起你人生當中的重量.
我記得在我人生最低潮的時候.
有一段很長的時間,我都無法站起來.
但我有木槿會.
我很多時候都會帶人去祈禱.
但當我獨自一個人的時候.
我無法站起來.
因為我覺得我很生神氣.
當我媽媽用這種方式離開這個世界時.
我覺得我無法接受.
我覺得無法接受神.
為何在我的生命中會擺出這件事.
為何會讓這件事發生.
我站了這麼久,為何不是痊癒.
所以我很長時間無法站起來.
有一天,我當時已經讀了神學.
信了主很多年,做了傳道人.
我帶人去祈禱,但我自己無法站起來.
這個張力挺大的.
有一天,我在茶餐廳跟阿邦說.
其實我很少說話.
因為在人的痛苦中.
其實都不想再觸碰這些地方.
我問他,我問了一條神學問題.
我問「其實甚麼是祈禱」.
你會想「對,我們平時經常都祈禱」.
「其實甚麼是祈禱」.
阿邦回答了一句.
他說「祈禱是一種張力」.
不知為何我就哭了.
其實我心底裡跟神說.
「神啊,我很痛」.
只是說了一句.
然後我發覺.

$^{401}$不知道神怎樣看待這個祈禱.
可能祈禱不長.
但原來祈禱是這麼大的張力.
原來在苦難中,祈禱是這麼難的.
但原來即使祈禱只有一句.
那也是一個有分量的祈禱.
而且也是一個上帝能夠接住的祈禱.
上帝在旋風中跟約伯相遇.
每一個問題,祂不是全都給了答案.
在人生中的苦難未必有答案.
我們也解釋不了.
但每一個問題.
彷彿都是上帝對人的邀請.
邀請我們不是用我們僅有的知識.
去解釋生命中的苦難.
而是去接受人生本身是一個奧秘.
對吧?.
丁姊妹,你覺得人生有很多事情無法解釋嗎.
是一個奧秘.
人生中的苦難不能隨便讓人解釋.
不只是簡單說你祈禱不夠信心就能解釋.
不是說你犯罪,怎樣怎樣就能解釋.
亦不論我們去解釋別人的苦難.
甚至我們有時也不應該去審判自己.
我們遭遇這些事就是因為神,怎樣怎樣.
約伯真誠地跟上帝對質.
他告訴我們.
即使在苦難中.
我們跟神的關係是值得信任的.
是可以承托我們的生命.
在苦難的奧秘中,上帝依然跟人同在.
上帝依然願意跟人面對面對話.
以至在旋風中出現過的上帝.
跟約伯對話過的上帝.
約伯最後說了一句.
他說.
「我從前風聞有你,現在親眼看見你」.
約伯沒有明白所有事情.
但他卻很真實地去經歷神.
靈姐妹,當我們面對生命的困難時.

$^{441}$我們會問.
在我們現在的處境,我們應該怎樣做.
我們很想有一個答案,有人能給我們.
今天我很想邀請你來跟我們分享.
不要急著去尋找答案.
而是停下來,去問一個問題.
問問你的心.
你想向神作出一個怎樣的呼喊.
可能你會像約伯那樣發現.
神沒有即時給你答案.
你期望有的出路.
但當你帶著真實的情緒.
當你內心混亂不堪.
很黑暗的時候.
你來到神面前.
你尋找神的時候.
神不會轉身離開.
祂會以祂的方式來回應你.
就像祂回應約伯一樣.
這一刻很想用一首詩歌.
讓大家去沉澱一下.
很想當你聽這首詩歌的時候.
你去問一問你的心.
這一刻,你有甚麼.
在你心底裡很想向神呼喊.
這一刻,我邀請你向神禱告.
(祈求神會幫助我).
(看天邊飄過雲海).
(高掃世界幻變尚在).
(從來人在月缺那盤魚片).
(天天總有上演).
(撥掃你我動人故事).
(動人情節留下腹屎).
(細心飾演).
(如何人生粉碎像未曾).
(無言流淚滿身痘疹).
(如何遇見神).
(被安抱才得起贈).
(冥冥平安方法天邊圓).
(流連遊蕩暗室放寸).

$^{481}$(然而遇見神).
(被指引再生精彩沾內心).
我們一起唱.
(看天邊飄過雲海).
(高掃世界幻變尚在).
(從來人在月缺那盤魚片).
(天天總有上演).
(撥掃你我動人故事).
(動人情節留下腹屎).
(細心飾演).
(如何人生粉碎像未曾).
(無言流淚滿身痘疹).
(如何遇見神).
(被安抱才得起贈).
(冥冥平安方法天邊圓).
(流連遊蕩暗室放寸).
(然而遇見神).
(被指引再生精彩沾內心).
(只因有你 今天有你了不起).
(誰無限想一生精彩皆美).
(有誰來導演出好戲).
(人難遇見神是否起).
(最終高飛).
(從來難數高低多少次).
(何時成就哪天失意).
(前途或變時路雖遠).
(有這福氣心內戀).
我們一起祈禱.
(祈禱).
親愛的天父.
你是那位聆聽我們問題的神.
你不嫌棄我們的懷疑,痛苦.
甚至質問.
你願意用愛和耐心.
去接住我們每一滴眼淚.
每一聲的哭喊.
主要我們承認.
人生裡實在有太多我們不明白的事.
有些痛苦實在太沉重.
重到我們連說都說不出.

$^{521}$但今天透過《約伯》的生命.
我們再次被提醒.
或許你不一定即時給我們答案.
但你永遠都在.
你永遠都與人同行.
主啊.
我們將我們最真實的心交給你.
你知道我們心裡每一個角落.
你不是要我們去抑壓.
也不是要我們去假裝堅強.
就算有懷疑,軟弱,怒氣.
我們都願意帶著這些情緒走近你.
因為你是我們靈魂裡最真正的答案.
求你教我們安靜的等候.
在暴風裡聽見你微小的聲音.
我們學習不是用腦去解釋一切苦難.
而是用心靈去經歷你的臨在.
因為你是那位.
在問題中依然與我們相遇的神.
你認識我們.
主多謝你.
從前我們風雲有你.
現在我們親眼見到你.
我們這樣禱告.
是奉主耶穌基督得勝的明智祈求.
Amen.
好,祝福大家.
\newpage



\section{使徒行傳 20:17-38}
\label{sec:i_UUX_QnitY}
\textbf{明辨真理 (使徒行傳 20\_17-38) - 陳崇基牧師}
\newline
\newline
連結: \href{https://youtube.com/watch?v=i_UUX-QnitY}{\texttt{https://youtube.com/watch?v=i\_UUX-QnitY}} ~~~~ 語音日期: 2025-04-16
\newline
\newline
\hyperref[sec:FaEgprmwTI0]{\small{< < < PREV SERMON < < <}}
~
\hyperref[sec:index_chronic]{\small{[返順時目]}}
~
\hyperref[sec:index_scriptual]{\small{[返順卷目]}}
~
\hyperref[sec:TstKn_BEGfA]{\small{> > > NEXT SERMON > > >}}
\newline
\newline
使徒行傳 20:17-38
\newline
\begin{longtable}{cl}
\hline
\hline
章節 & 經文 (和合本修訂版)\\
\hline
20:17 & \begin{tabularx}{0.7\textwidth}{X} 保羅從米利都打發人往以弗所去，請教會的長老來。 \end{tabularx} \\ \\ \relax
20:18 & \begin{tabularx}{0.7\textwidth}{X} 他們來了，保羅對他們說：「你們自己知道，自從我到亞細亞的第一天，我怎樣跟你們相處， \end{tabularx} \\ \\ \relax
20:19 & \begin{tabularx}{0.7\textwidth}{X} 怎樣凡事謙卑，以眼淚服侍主，又因猶太人的謀害經歷試煉。 \end{tabularx} \\ \\ \relax
20:20 & \begin{tabularx}{0.7\textwidth}{X} 你們也知道，凡對你們有益的，我沒有一樣隱瞞不說的，或在公眾面前，或在每一個人的家裡，我都教導你們， \end{tabularx} \\ \\ \relax
20:21 & \begin{tabularx}{0.7\textwidth}{X} 不論猶太人和希臘人，我都已證明他們當在神面前悔改，信靠我們的主耶穌。 \end{tabularx} \\ \\ \relax
20:22 & \begin{tabularx}{0.7\textwidth}{X} 現在我被聖靈催迫要往耶路撒冷去，雖然不知道在那裡會遭遇甚麼事， \end{tabularx} \\ \\ \relax
20:23 & \begin{tabularx}{0.7\textwidth}{X} 但知道聖靈在各城裡向我指證，說有捆鎖與患難等著我。 \end{tabularx} \\ \\ \relax
20:24 & \begin{tabularx}{0.7\textwidth}{X} 我卻不以性命為念，只要走完我的路程，完成我從主耶穌所領受的職分，為神恩典的福音作見證。 \end{tabularx} \\ \\ \relax
20:25 & \begin{tabularx}{0.7\textwidth}{X} 「我素常在你們中間到處傳講神的國；現在我知道，你們眾人以後不會再見到我的面了。 \end{tabularx} \\ \\ \relax
20:26 & \begin{tabularx}{0.7\textwidth}{X} 所以我今日向你們作證，你們中間無論何人死亡，罪不在我。 \end{tabularx} \\ \\ \relax
20:27 & \begin{tabularx}{0.7\textwidth}{X} 因為神一切的旨意，我並沒有退縮不傳給你們的。 \end{tabularx} \\ \\ \relax
20:28 & \begin{tabularx}{0.7\textwidth}{X} 聖靈立你們作全群的監督，你們就當為自己謹慎，也為全群謹慎，牧養神的教會，就是他用自己血所買來的。 \end{tabularx} \\ \\ \relax
20:29 & \begin{tabularx}{0.7\textwidth}{X} 我知道，在我離開以後必有兇暴的豺狼進入你們中間，不顧惜羊群。 \end{tabularx} \\ \\ \relax
20:30 & \begin{tabularx}{0.7\textwidth}{X} 就是你們中間也必有人起來，說悖謬的話，要引誘門徒跟從他們。 \end{tabularx} \\ \\ \relax
20:31 & \begin{tabularx}{0.7\textwidth}{X} 所以你們要警醒，記念我三年之久，晝夜不斷地流淚勸戒你們各人。 \end{tabularx} \\ \\ \relax
20:32 & \begin{tabularx}{0.7\textwidth}{X} 現在我把你們交託給神和他恩惠的道；這道能建立你們，使你們和一切成聖的人同得基業。 \end{tabularx} \\ \\ \relax
20:33 & \begin{tabularx}{0.7\textwidth}{X} 我未曾貪圖一個人的金、銀或衣服。 \end{tabularx} \\ \\ \relax
20:34 & \begin{tabularx}{0.7\textwidth}{X} 你們自己知道，我靠兩隻手工作來供給我和同工的需用。 \end{tabularx} \\ \\ \relax
20:35 & \begin{tabularx}{0.7\textwidth}{X} 我凡事給你們作榜樣，叫你們知道應當這樣勞苦，扶助軟弱的人，又當記念主耶穌的話，說：『施比受更為有福。』」 \end{tabularx} \\ \\ \relax
20:36 & \begin{tabularx}{0.7\textwidth}{X} 保羅說完了這些話，就和大家跪下來禱告。 \end{tabularx} \\ \\ \relax
20:37 & \begin{tabularx}{0.7\textwidth}{X} 眾人痛哭，抱著保羅的頸項跟他親吻。 \end{tabularx} \\ \\ \relax
20:38 & \begin{tabularx}{0.7\textwidth}{X} 叫他們最傷心的，就是他說「以後不會再見到我的面」那句話。於是他們送他上船去了。 \end{tabularx} \\ \\
[1ex]
\hline
\hline
\end{longtable}
$^{1}$多謝主人的報告.
下個星期就是我們的Easter Sunday.
即是復活節的主日.
我們這個星期五就是Good Friday Service.
我們邀請各位來到我們中間.
下個星期五.
在這個地方.
同一個地方七點半鐘.
我們有受復活節的崇拜.
下個星期正如主人所說.
就是我們Easter Sunday.
亦是我們教會的九週年.
我們已經九歲了.
九歲.
我們教會很多人都九歲.
在那邊.
在小朋友那邊.
我們應該很有活力繼續向前.
下個星期我們剛才說到藝術表演.
其實就是兩幅很大的畫.
就會在我們崇拜裡面畫出來.
我們希望真的坐前一點.
我們希望在視覺上.
在音響上.
能夠帶給大家一個新的開始.
因為我們九年開始.
第九歲的時候.
我們在講抗疫.
以及在尋找地方等等.
都是需要我們大家一起.
同心協力的去建立教會的.
OK.
今天講到的不是一個系列.
是一個特別加插下去的.
雖然今天是Palm Sunday.
通常我們講耶穌榮進耶路撒冷.
不過在這裡我想加插一個.
跟我們大家應該都頗接近的講道題目.
原因是什麼呢?.
因為三月底的時候.

$^{41}$我們大芝加哥區華人教會有一個聯會.
聯會發了一個電郵給所有教會.
就說在大芝加哥區裡面的教會.
有一個異端非常活躍.
我們都知道最近異端是活躍的.
譬如說夜戶,見證人.
特別這個電郵講到一個異端.
叫做東方閃電.
所以它滲入不同的教會.
也是很主動,很熱烈的去認識人.
去傳講他們的信息等等.
所以為什麼今天想加插這個進去呢?.
就是想大家有一點點的 awareness.
還有一點點的明白.
好嗎?.
我的經文是取自《使徒行傳》第20章.
這裡是所謂的militant farewell speech.
就是保羅在他第三次宣教回程的時候.
快要回耶路撒冷的時候.
在米利都這個地方.
米利都一個海灣,一個碼頭.
跟以弗所的長老有一些交代,有一些叮嚀.
這裡講到教會會遇見什麼樣的威脅.
很簡單有三點我想今天帶出的.
第一個就是真理.
福音的真理.
聖經的真理.
是可以被扭曲的.
可以被扭曲的.
二端就是這個意思.
是不是?.
第二點就是真理是必須要我們去守護.
並且我們以生命來活出.
講出米利都的道別講話.
先給一點背景.
你看地圖就知道.
這次是保羅第三次的旅程.
第三次宣教的旅程.
你看到他從右上角安提阿開始.
在那裡在陸路去到以弗所.

$^{81}$看到嗎?.
以弗所是一個很重要的城市.
然後在那裡他一直上去北面.
就去到特羅亞另外一個海灣,另一個碼頭.
在那裡他就過海.
就是過海對面.
然後就將福音傳到菲臘比.
去探望那裡的教會.
然後他就繼續以陸路下去到下面的哥林多.
就是今天的希臘半島的地方.
然後又再上去.
回到特羅亞他又再坐船下去.
這次繞過以弗所.
因為他趕路要回去耶路撒冷.
為什麼要回去耶路撒冷呢?.
就是因為他想報告愛邦人的工作.
所有的修成.
告訴總會的舞會.
並且保羅周遊列國到不同的地方.
在特別愛邦人多的教會裡面去籌款.
去支持舞會裡面的大英姐妹.
因為當時有一些天災等等.
所以耶路撒冷的教會很需要資金上的支持.
所以保羅就帶著這些奉獻.
他趕路要去到耶路撒冷.
所以他就繞過了以弗所.
然後他就在米利都那裡做什麼呢?.
他就叫以弗所的長老來到.
然後跟他們有一番說話.
那個就是所謂的Military Farewell Speech.
你看到這個整個地圖.
他當然是回到耶路撒冷.
為什麼這個所謂米利都道別的講話這麼感人呢?.
還有帶著眼淚呢?.
因為保羅知道他回去耶路撒冷的時候.
除了他辦理事公的事.
他還會遇見很多的反對.
因為耶路撒冷有很多反對他的猶太人.
想去陷害他.
他也知道可能到了耶路撒冷的時候.

$^{121}$將會有患難並且會有困所.
結果當然保羅就回到耶路撒冷.
回到沒多久就有一些暴動.
然後有一些反對他的人.
猶太人就將他抓拿.
然後在公會那裡審判.
帶到羅馬官面前去審判.
帶到當時有一個希律王.
阿基帕王在面前去審判.
等等等等.
如果這樣說的話.
像不像什麼?.
像誰?.
就是他去到耶路撒冷要受苦.
然後被公會審判.
被猶太人去害.
然後在巡撫總督面前去受審.
在希律王面前受審.
像誰?.
像耶穌沒錯.
整個《使徒行傳》最後這一段.
就是講到保羅.
告訴人聽.
保羅是一個怎樣的人.
就好像耶穌基督那樣.
所以他回到耶路撒冷的時候.
確實就是發生了這樣的事情.
然後他就再沒有所謂的自由.
然後他就說我要上告該殺.
所謂《使徒行傳》27,28章.
最後兩章.
就是講到他上告該殺的.
遇見不同的風浪等等.
不同的苦難.
最後他是帶著福音.
雖然沒有了真正的完全的自由.
帶著福音去到羅馬.
所以整個《使徒行傳》下面那一段.
就是說保羅很像耶穌基督一樣.
所以在米利都的講話.

$^{161}$我們回到這裡.
就有一點像耶穌基督最後晚餐的時候.
講的farewell speech.
因為這個是很重要很重要.
對教會非常重要.
因為他將會無辜的受害受審.
並且會殉道死.
好像耶穌一樣.
結果我們知道他去到羅馬的時候.
本來就是人頭落地.
所以米利都最後的講話.
就像耶穌和門徒講話一樣.
這個是小亞細亞裡面最後一次.
是公開的講道.
是一篇很感人的講話.
是交付整個教會.
去叮囑教會的領袖.
要去堅守真理.
我們這個講話值得去看一看.
值得去看一看.
你可以想像一下在碼頭.
就是林頌基那一刻.
他有一番話要講.
講了他才進去那個TSA.
才過海關.
有像這樣的情況.
不過不是頌基是頌船.
因為他在米利都坐船.
然後去到耶路撒冷.
我們值得去看這段講話.
你這樣想像的時候.
從17節開始你就知道那個情景是怎樣.
20章第17節.
保羅從米利都派人到爾弗所.
請教會的長老來.
他們到了保羅那裡.
他就對他們說.
你們知道從第一天我踏足亞細亞以來.
我一直怎樣與你們相處.
我怎樣服侍主.

$^{201}$傳言謙卑帶著眼淚.
經受試煉.
就是猶太人的計謀家諸在我身上.
你們也知道任何對你們有益的道理.
我沒有退縮不向你們述說.
或在公共的場合.
或在個人的家裡面.
我都教導你們.
我向猶太人希臘人鄭重的作見證.
告訴他們應當向神悔改.
信靠我們的主耶穌基督.
看我現在被聖靈催逼要到耶路撒冷去.
我不知道在耶路撒冷遇見什麼事.
只知道聖靈在各城裡面.
向我指明有困所和患難等著我.
但我不看重自己的性命.
為要完成我的路程.
完成我從主耶穌所領受的積分.
向人鄭重見證上帝恩典的福音.
我曾在你們中間走遍各處宣揚神的國.
但現在看了我知道.
你們全都不會再見到我的面了.
因此我今天向你們作見證.
我是清白的.
你們所有人的血都與我無關了.
因為神一切的旨意.
我沒有退縮不向你們述說.
聖靈已經立你們作全群的監督.
你們要謹慎.
也要為全群謹慎.
要牧養神的教會.
就是祂用血所救贖回來的.
我知道我離開之後.
必有兇猛的豺狼進入你們中間.
不外適養群.
而且你們中間也必有人起來說乖貓的話.
要引誘門徒跟從他們.
因此你們要謹慎.
記住我三年以來.
晝夜不停的流淚圈戒你們每一個人.

$^{241}$現在我把你們托付給神.
和祂恩惠的道.
這道能建立你們.
能在所有誠信的人中間賜給你們基業.
我沒有貪圖過任何人的銀紙金紙和衣服.
你們自己知道.
是我用這雙手供給我和同伴的需要.
我在這一切事上都給你們作了榜樣.
我們必須這樣勞苦來輔助軟弱的人.
並要記住耶穌的話.
施比受更為有福.
保羅說完了這話.
就跟大家跪下來禱告.
大家全都大聲同哭.
抱著保羅脖子親吻他.
讓他們最傷心的就是保羅說.
他們將不會再見到他的面那句話.
之後就送了保羅上船.
然後他就離開了.
是一個很感人的.
感人肺腑的一段話.
正如你所想像的.
在碼頭在機場.
快要送機那一刻.
快要走的那一刻最後見一面.
他就告訴你一段很重要的話.
第一點我想說的是.
他講清楚福音的真理.
是可以被人扭曲的.
特別是被一些假教師扭曲.
28章29節附近.
這裡說什麼呢.
他說有些兇猛的豺狼會進到教會裡面.
不愛惜羊群.
而且你們中間也有人會起來.
說乖謀的話引誘門徒跟從他們.
這裡講到內憂外患.
教會裡面經常會經歷一些內憂外患.
這句話很有趣.
他說你們中間也有人會起來.

$^{281}$如果你是送機那一群人.
你會說.
大家看看是不是你還是我.
這句話不一定是這個意思.
不一定是送機那一群人.
裡面中間有人.
其實他可能是在說.
已乏所的群體.
已乏所的教會裡面.
因為已乏所裡面有很多人.
裡面可能有些人.
中間做了所謂的掩仔.
將真理扭曲.
但是我們需要了解這一段的說話.
就是告訴我們歷代的教會.
是會面對真理被扭曲的危機.
如果你記得教目書信.
你記得嗎.
教目書信.
提摩塔前後書和提多書.
裡面你看一看.
大概你的印象.
你都會看到裡面遇到很多不同的問題.
其中一個問題就是異端的問題.
就是一些假教師混入了.
特別是女人的家.
混入了一些家庭.
然後用假道引誘人.
這些說話絕對不是看低女人.
因為當時第一世紀的女人.
可能教育的機會比較少.
所以很多時候異端.
就可以在這些群眾裡面.
乘虛而入.
當然今天不同.
今天很多女士的學位.
可能比我們都高.
所以她們是非常有教育的.
但是無論如何教會.
有時候會遇見一些異端的問題.

$^{321}$所以我們需要去學習去分辨.
東方小說剛才我提過.
你們可以分析一下是什麼.
有沒有聽過.
我大概告訴你.
是中國製造的.
是從河南出來的.
幾十年前從河南出來.
然後又說自己不行.
東方小說其中一個名字.
另外一個名字就是全能神教會.
或者在網上看到一些很好的.
很好的質素的YouTube.
都是全能神教會所做出來的.
東方閃電這個名字是從哪裡來的呢.
馬太福音你記得24章.
耶穌說到他再來的時候.
聖殿被毀和他再來的異象裡面.
曾經有一句是這麼說的.
27字.
好像閃電從東方發出.
發起東方閃電的人.
就說對了.
因為在歷史裡面.
舊約新約.
是有三個時代的.
第一個時代就是所謂的律法時代.
律法時代就是舊時要守律法.
然後就到耶穌的時代.
就是恩典的時代.
然後現在這個時代就不同了.
律法恩典時代都過去了.
現在第三個時代就是所謂國度的時代.
Kingdom的時代.
這個時代就不是了.
東方有另外一個教主出來.
有另外一個基督.
那個基督而且是女基督.
女的基督.
你聽一下的時候就知道.

$^{361}$這些全部是狗屁不通的東西.
是胡說八道的東西.
但又很奇怪.
有些群眾就會去跟他.
對不對.
我就覺得很可惜這個時候.
這個時代21世紀.
我們有check GPT 有AI.
很多東西跟我們.
然後我們可以check一些東西.
為什麼有資訊但沒有分析力呢.
但這裡告訴我們.
真的可以避漏.
我們是需要用腦的.
各位不要只說憑信心就可以了.
OK.
我明白我們需要憑信心.
但不要亂用這個字.
這個說話並非告訴我們.
我們看聖經是不用腦.
是憑心憑感覺就好.
絕對不是這樣.
上帝給你造腦就要用腦.
你不用腦就不去分析.
不去看有什麼證據.
我們很多時候會錯意.
我們以為信仰只是沒有根據的東西.
所以這裡講得很清楚.
有豺狼出來在羊群裡襲擊.
去引誘.
第二方面這裡講到.
教會的領袖需要做什麼呢.
聖靈立你們作全群的監督.
你們就當自己的謹臣.
全群謹臣 為全群謹臣.
牧養上帝的教會.
用自己的血所買回來.
是很貴重的東西.
要謹臣兩次的講道.
聖靈立你們作監督的職份.

$^{401}$意思不是你們自己選擇.
是一個神聖的挑選.
或者是一個神聖的託付.
是一個信任.
是託付了給你們.
你們要好好管理.
講到什麼呢.
由上問下.
保羅說要教傳避的東西.
The whole council of God's will.
整本的聖經.
不只是挑選自己喜歡的地方.
是整個的真理.
我們需要有一個掌握.
是需要去守護的.
當然這裡直接的對象.
就是以弗所的教會領袖.
但其實是在說所有的基督徒.
我們每一個都有責任.
保羅吩咐教會特別領袖.
就是需要做什麼.
要小心.
如果我們不謹慎自己.
不懂得去思考的話.
我們就不能幫助羊群.
如果我們作基督徒.
不肯去用腦的話.
我們都不能去見證真理.
對不對.
我剛才說過.
這個時代是資訊爆發的時代.
但非常之可惜的.
就是我們分析力越來越差.
加上AI看的東西.
它就繼續給你提供東西.
你越看越覺得是對的.
這個在邏輯上.
叫做confirmation bias.
叫做確認偏差.
什麼意思.

$^{441}$人本身有一個前設.
他只會注意或者只會選一些.
支持自己本身有的前設觀點的資訊.
他忽略了相反的東西.
這是一個很大的問題.
你看新聞都有左中右分.
如果是看右的.
他就繼續給你提供右的東西.
所以你要懂得去分析.
人更加需要去分析.
現在AI越來越多.
越來越多這樣的資訊的時候.
我們更加需要練習一下.
怎樣去做analysis.
怎樣去debate.
怎樣去明白一件事.
對不對.
神學都是一樣.
真理都是一樣.
上星期在基督教教裡面最大的新聞.
就是芝加哥三一神學院要搬走.
突然間說要搬走.
搬到加拿大.
接近溫哥華.
以前我在溫哥華的時候.
有一間學府叫做Trinity Western University.
那裡都算是挺有名的.
因為Trinity Western University.
有一個研究所專門做死海古卷.
In any case.
上個星期就突然間有一個炸彈.
三一神學院.
一個已經百多年.
在我們Midwest.
在我們中西部.
要搬走.
因為收生不足.
財務的壓力等等.
雖然我不是三一神學院的校友.
但我覺得是非常之可惜的.

$^{481}$因為多年來.
我們知道在三一神學院出很多一流的學者.
一流的研究.
一流的文章.
一流的書.
他們出的學生也有很多很好的學生.
在我們中間也有幾位.
三一神學院做了很多功不可沒的東西.
只不過有很多壓力的緣故.
所以結果他們決定搬走.
然後就和Trinity Western University.
在溫哥華附近的這間學院合併.
這是很大的一件事.
因為對學生.
特別對學生有很多影響.
如果本身在那裡讀書.
突然間不能讀書.
特別對一些國際學生更加麻煩.
是非常之可惜.
我們真的需要為學院裡面的學生去禱告.
為學院裡面的師子去禱告.
非常之可惜.
你看到這裡有個桌子.
現在在美國頭十大神學院.
當然三個神學院都在這裡.
HC就是Headcount.
FTE就是Full Time Equivalent.
那個數字.
頭十大有第一是Liberty.
Liberty在維吉尼亞.
第十是Gordon Cornwell.
Gordon Cornwell在馬薩諸斯.
在波士頓.
十間的神學院.
最大的.
這個是美國裡面最大的.
並且他們的修生率是高的.
學生的數字高.
並且有些是連年破紀錄的.
包括我自己那間.

$^{521}$Midwestern.
你看到這裡十間.
六間都是美國南方浸信會的屬於.
在Southern Baptist.
看到嗎.
Southern Baptist Theological Seminary.
Midwestern.
Dallas就是時代論.
傳統時代論.
它在Dallas.
Southeastern.
Southwestern.
Ashbury就是Western.
就是巡渡委理會的傳統.
New Orleans.
都是Southern Baptist.
Gayway在LA.
也是Southern Baptist.
以前叫做Golden Gate.
在三藩市.
現在搬到LA.
Gordon Cornwell.
Massachusetts.
Liberty.
是非常超保守的.
獨立的.
超保守.
你看到十間裡面六間都是美南浸信會.
是不是好事.
我在Midwestern.
是不是好事.
Midwestern.
我聽到上個學期.
我們修生率破紀錄.
是不是一定好事.
不一定是好事.
為什麼這麼說.
雖然我的學校不受影響.
但我不覺得是一件好事.
因為人學的進展.

$^{561}$我們的知識的進展.
是需要有不同的聲音.
如果只是你自己的門派.
只是你美南浸信會裡面所說的聲音.
單一個角度就容易有什麼.
Confirmation bias.
我不是說基本真理.
基本真理大家都可以的.
無論你是哪個神學院.
哪個宗派都好.
巡渡委理會.
浸信會.
宣導會.
甚至是三個神學院的宗派.
Unity of Good Freedom.
波動會.
基本真理是沒有問題的.
但有很多東西我們還要在新的時代裡面去sort out.
包括女性可不可以募會等等.
可不可以被按立.
這些重要的事.
是需要有不同的聲音才能夠有改進.
如果只是單一的聲音.
這不是一件好事.
三一離開中西部.
是一個很大的損失.
尤其大圖畫來說.
我們特別做教牧的.
我們做牧師的.
我們不用考牌的.
我們很容易停下來不學東西.
對不對.
我做了這麼多年.
你在我的文件夾裡.
隨便都有一些舊子彈可以用.
我可以不用新子彈.
我可以把舊菜炒出來.
有時候舊菜隔夜飯.
蒸飯也挺好吃的.
但問題是我們沒有了全面的神學訓練.

$^{601}$你怎樣做牧者.
我想說一件事.
一個硬幣有兩面.
當然我們可以說學院的行政等等.
搞得不好等等.
但另一面就是教會會不會負責.
弟兄姊妹我們有沒有份.
我們是有份的.
我們有份的.
當我們看到有需要的時候.
我們就攪在一起.
當我們說我們是信真理的.
我們不會去擁護真理.
Shame on us.
我們就是一群這麼自私的人.
我們只不過是等上天堂.
我們在地上上帝的工作.
我們完全漠不關心.
Shame on us.
所以這不是一個很好的現象.
所以有分析說.
福音派南方化.
Southernification.
Evangelicalism.
即是以後的大基地.
美國福音派的聲音.
就會轉移到南邊.
剛才那些神學院大部分都是南邊.
在巴布比爾.
It's a good thing I don't think so.
但值得我們去提醒.
去提醒我們.
我們是不是care才行.
我們是不是care.
我們的資源到底去了哪裡.
全部都是我們自己去計算.
但上帝的角度我們漠不關心.
所以真理是需要守護.
需要去培養的.
弟兄姊妹.

$^{641}$我經常說我們基督徒.
不應該是cheap精.
為什麼我們這麼cheap.
我們cheap到暈.
因為根本就是我們華人的DNA.
我們覺得.
我們拿到一點便宜.
我們多開心.
你就不知道你的科溫卡跌了多少.
我們就是面前的小便宜.
我們覺得很厲害.
施比受更為有福.
保羅說什麼.
你知道自從我在亞細亞的第一天.
我怎樣和你們相處.
我的生活是怎樣.
凡事謙卑以眼淚來服侍主.
有因猶太人的陷害經歷去到試煉.
你們也知道.
凡對你們有益的.
沒有一樣隱瞞不說的.
公眾的時候.
各自在家裡面的時候.
我都教導你們.
凡事給你們作榜樣.
叫你們知道應當這樣勞苦.
符咒軟弱的人.
紀念耶穌的什麼話.
福音書裡面沒有記載.
施比受更為有福.
是不是.
施比受更為有福.
真的需要活出來.
在你生活裡面活出來.
你看到保羅的模樣是怎樣.
他的榜樣是怎樣.
凡事寧願自己有欠缺.
去謙卑服侍其他人.
所以他在天國裡面是有符咒的.
如果我們是孤魂鬼的話.

$^{681}$我們在天國裡面.
我們就是破產的人.
我不管你地上有多少錢.
那些你都帶不走.
但天國裡面我們不是有符咒的.
我們的生命.
我們的生命是不是有實證有力的.
來告訴別人.
以致我們能夠見證耶穌基督.
是有壓力的.
最後我想重覆這三段.
這三個重點.
真正可以被扭曲的.
各位大紀元的節目.
你需要去分析.
你需要想清楚.
你聽到異端說的話.
你不可以每次都說.
總之有人會解答.
我自己都要弄清楚.
真理是需要我們去落力去守護.
去堅持.
用我們的命.
有壓力的去活出來.
這裡很簡單.
用幾部曲.
我相信每一個基督徒.
無論你今天剛剛信的.
你都應該有這幾部曲.
在真理上我們要去學習.
去前進的.
我們要動腦筋.
我們要動思考.
要有思維.
我們要跟著這幾部曲去前進.
聽了福音的我們就信而受洗.
並且我們會明白真理.
並且在生命裡面是行出來的.
當然這不是一個直線的關係.
有時候一些基本的東西.

$^{721}$我們都要重新聽.
但這裡很簡單的四部曲.
我們都講出來了.
如果我們不是這樣的話.
我們基督徒就會很糟糕.
教會就會很糟糕.
聽了你又不信.
信了你又不明白.
明白了你又不去走.
你想耶穌怎樣?.
你想怎樣?.
信了耶穌我們信而受洗.
我們的生命是屬於主的.
我們要去學懂真理.
搞清楚真理.
然後走出來.
如果是這樣的話.
真理才能夠在各處開花.
如果不是就是異端.
細心去聽上帝的話.
懂得去解釋.
這樣你的信仰才會堅固.
你的生命才會有壓力.
去走出來.
我們同心協力去祈禱.
主愛多謝你.
因為你給我們真理.
在恩典裡面我們能夠認識.
接受這個恩惠的福音.
感恩你是那位愛我們的主.
主愛求你幫助我們.
在我們的思考.
我們的真理上面.
我們能夠更加盡心.
能夠更加落心一點.
讓我們的根能夠紮緊.
穩穩的.
來到向前去見證你.
主求你提醒我們.
去敲醒我們.

$^{761}$如果是沉睡的時候.
你叫真理.
叫你的愛.
叫你的恩典再一次去喚醒我們.
以致我們能夠不單止看見你的真理.
並且去追隨你的真理.
去追求.
去認識你.
認識你的說話.
懂得去分析.
懂得去思考.
主要在這個資訊爆炸.
但沒有分析力的世代.
在這個沒有智慧的世代.
主求你教導我們.
如何去學懂.
分辨.
去辨識.
主要幫助你的教會.
保守你的教會.
因為往往會有很多時候.
有異端.
有些假道.
似是而非的道理進入.
求你去提醒我們.
幫助我們來明白.
你幫助我們不單止是用口號.
憑信心就算.
我們真的要下功夫.
學懂如何去跟著你的道.
作思考.
去明白你自己的話.
主要讓你的教會在真理裡面能夠成長.
以致我們福音的見證.
是有力的.
我們這樣祈禱奉主名.
多謝各位.
希望下星期見到大家.
下星期是Easter Sunday.
我們11點鐘開始.

$^{801}$11點鐘準時開始.
\newpage



\section{5:27-30}
\label{sec:TstKn_BEGfA}
\textbf{Righteousness as Redefined by Jesus: Lust (Matthew 5\_27-30) - Chee Yie Jong}
\newline
\newline
連結: \href{https://youtube.com/watch?v=TstKn_BEGfA}{\texttt{https://youtube.com/watch?v=TstKn\_BEGfA}} ~~~~ 語音日期: 2025-04-20
\newline
\newline
\hyperref[sec:i_UUX_QnitY]{\small{< < < PREV SERMON < < <}}
~
\hyperref[sec:index_chronic]{\small{[返順時目]}}
~
\hyperref[sec:index_scriptual]{\small{[返順卷目]}}
~
\hyperref[sec:as4vlJLVesI]{\small{> > > NEXT SERMON > > >}}
\newline
\newline
$^{1}$Good morning church and welcome back to.
Uv it is good to begin the new year in the house of the lord worshiping.
Him together.
So if you have been with us, uh.
Since last year you will know that every first sunday of the month besides going through.
Or participating in communion. We have also been diving into jesus sermon on the mount.
And one of the reasons that I choose to preach expository through this section, uh in matthew from matthew chapter 5.
to 7.
section by section verse by verse.
In contrast to preaching topically where I pick a topic and then just kind of like find different passage to supplement it.
It's because.
I cannot just choose what I want to talk about in the bible and avoid the topics that I don't want to speak on.
Rather the scripture itself.
Function as a guide for myself and the word of god dictates what I will be preaching on.
So today we will actually dive into a topic which can be a little taboo to talk about in church at times.
And not only for the church. I feel like this is also especially true in our asian culture.
Which tends to be a little more reserved at times at least in my own experience growing up.
So today we're actually going to touch upon the topics of sex and sexuality.
And if i'm being honest, I would probably avoid preaching on this topic if I have the option.
Because it is a very sensitive topic and also highly contested.
issues.
Especially in our current cultural climate.
However, jesus himself he doesn't shy away from this topic. So I feel like as the church we should not be as well.
as we are.
In fact sex and sexuality is part of god's original design that he called good.
The unfortunate reality is that the world the world that we live in currently is constantly bombarding us with all different ideologies and perspective.
That is just so far off from god's original design.
So.
As the church as christ followers.
I think it is even more urgent and important for us to understand what godly sexuality looks like.
So if you have your bible with you, let's turn with me to matthew 5.
And we're going to read from verse 27 to verse 30.
So matthew chapter 5 verse 27.
You have heard that it was said.
You shall not commit adultery.
But I tell you that anyone who looks at a woman lustfully has already committed adultery with her in his heart.
If your right eye caused you to stumble.
Gouch it out and throw it away.
It is better for you to lose one part of your body than your whole body to be thrown into hell.
And if your right hand caused you to stumble.

$^{41}$Cut it off and throw it away.
It is better for you to lose one part of your body than for your whole body to go into hell.
so.
Let's pray.
Heavenly father.
As we open up your word this morning.
I just ask that um, your spirit will be with us.
Prepare our hearts as we receive your word.
And I also ask that your holy spirit will be with me as I preach this morning.
I ask that uh, you will take away anything that is not from you.
And may your word and your word only be preached this morning.
Amen.
So I know it is the new year already but not that long ago we just celebrated christmas, right?.
And usually during the christmas season besides celebrating the birth of jesus.
There are usually a lot of activities that is centered around gifts, right?.
Whether it is gift givings or gift exchange and in fact during our christmas service here.
We had a white elephant gift exchange as well.
And so this past christmas in our youth group in summit.
We organized a white white elephant christmas gift exchange for our high schooler as well.
And if you have participated in any sort of like gifts related activities.
I'm sure you have seen this old trick being done at some point where someone will put their actual gifts inside like a mismatched box.
Maybe like a bunch of toilet paper in like a really nice ps5 box.
And until you get to the actual gift itself.
It is pretty much impossible to tell what is actually inside the box.
And so without fail during our summit christmas party a couple of our students.
They did the same old trick.
One of them ensign. He's a guy but he bought.
a pretty nice barbie doll.
But instead of kind of like wrapping it in like a nice wrap or box.
He actually put it in like a 20 pound like noodles box basically like a huge box that he say he just like found at home.
And other students rex he bought also like some card games that is pretty nice.
But he ended up like putting it in layers after layers of different box cereal boxes shoe boxes and then also like different wrappings as well.
so.
Imagine you are the person receiving the gifts in this situation.
What actually matters to you?.
Is it the box itself?.
or the actual gift.
Would you rather get a really nice box?.
But with a gift inside there you're like, uh, I don't really need that.
Or would you rather get a general or basic box but with a gift inside that you actually want?.

$^{81}$I'm sure most of us will probably choose the later, right?.
So in the same way in our passage today.
We'll see that when jesus he talks about righteousness.
And what it means to live rightly as god's people.
in god's world.
under god's rule.
He doesn't really care about the outward appearance of obedience or compliance.
But what jesus really cares about is the posture and attitude of our heart.
So in my last sermon we began to dive into this new section in jesus sermon on the mount which includes six case studies.
On jesus on jesus own version of righteousness and how he has fulfilled the law and the prophets.
Each of these six examples they highlight how we as christ followers ought to live in community. We find another.
Within the church and also with those outside of the church in the world around us.
And last time we look at the first examples.
On the sixth commandment you shall not murder and how jesus points to the root issues behind murder and showing anger and content.
towards her brothers or sister.
And the root issue is not treating a person made in the image of god with respect and dignity.
And so today the righteousness that jesus redefines for us and for his listeners.
Centers around the seventh commandment. You shall not commit adultery.
And similar to the first examples on murder jesus teaching on adultery today also has to do with something that is inside of us.
The core desires and intention of our hearts that impacts how we relate to other people made in the image of god.
So here is my three main point for today.
First we need to reclaim a godly view of sex and sexuality.
And second we need to ruthlessly eliminate sin.
And lastly, we need a transformed heart to live rightly.
So point number one we need to reclaim a godly view of sexuality.
Let's take a look again at verse 27.
You have heard that it was said.
You shall not commit adultery.
So once again, we see this same formula that jesus is using here. You have heard that it was saying.
And when jesus audience hear this they would immediately know that jesus he is about to start a discussion.
on the particular mosaic law.
And this time the specific law that jesus will focus on is on the seventh commandment. You shall not commit adultery.
But what is adultery?.
For jesus adultery refers to any sexual relationships that is outside of the covenant of marriage.
Again adultery refers to any sexual relationships that is outside of the covenant of marriage.
In other words, this commandment is not just for the married couples.
But it also include any kind of sexual relationships.
That is outside of the bound of marriage including.
Like people that are single like me.
So if you have any sexual relationship outside of the marriage covenant or without a marriage covenant, it is considered as adultery.

$^{121}$But once again as the only one who has perfectly who has perfectly fulfilled the law of god.
Jesus he goes one step further and give his own interpretation.
Pointing us towards the heart behind this seventh commandment of you shall not commit adultery.
So jesus continues in verse 28.
But I tell you this.
Anyone who looks at a woman lustfully has already committed adultery with her in his heart.
So right here jesus is not talking about.
Just kind of like looking at someone and finding them attractive.
But the kind of looking that jesus talks about here is the kind of looking that makes you objectify a person.
Treating someone as an object to bring yourself pleasure.
And when we objectify a person we strip them of their god-given dignity.
So in some way this is similar to when we kind of like show anger and contempt towards another brother or sister.
And also notice how it begins.
It begins with the eyes looking right?.
And then we stay at what we look at for a little bit longer than we should.
And what we see becomes the desires of our heart.
Does this pattern sounds familiar at all.
This is the same pattern that we find in the stories of the fall in genesis tree.
This is genesis tree. It says when the woman saw that the fruit of the tree was good for food and pleasing to the eye.
And also desirable for gaining wisdom.
She took some and ate it.
So my question for us is this.
What are our eyes looking at.
And do these things help us to follow jesus or do these things help us to follow what our corrupted heart desires.
and.
Especially living in the age of social media.
there's like so many things that we look at whether it is on our screen on the tvs or.
Around us all the time. So I think it is even more important to.
really pay attention to where we kind of like.
Put our time into what are some of the things that we look at in our daily life.
And at this point I think it is important to make a few clarification.
So we don't swing all the way to the other end of the spectrum and fall into bondage as well.
so.
first clarification.
We need to make a distinction between temptation.
And the act of following true with our temptation.
Okay, so there's a difference between the temptation itself and also the action of following true with our temptation.
It is not a sin when temptation like the last thought pop up in our mind.
Because we don't have complete control again. We have some sort of control, but we don't have complete control of our thoughts.
or our surroundings.

$^{161}$You can't just be like, okay, I don't want to be tempted so i'm just gonna lock myself inside a cave.
So I just don't see anything. We're not able to leave like that, right?.
So that's not possible.
Also, it is not wrong to have sexual desires in fact it is part of who we are as human beings.
But what we do have control over is this.
Do we follow true with our temptation or do we resist it?.
Do we entertain the lustful thoughts in our mind longer than we should or do we cut it off right away.
Do we let our sexual desires and passion.
Take control over us or do we practice self-control which is the fruit of the spirit.
So what truly matters.
Is how we deal with temptations?.
and our desires.
when they come up.
so one of the.
famous theologian.
A german theologian martin luther if you don't know him. He's pretty much the person who started up the.
Protestant reformation out from the catholic church.
And this is what he says about temptation.
it says.
Thus you see that temptation can be avoided by no one.
But resistant may be made.
And with prayer and recourse to divine aid we can put ourself in readiness to meet such designs.
In the book of an old father we read that a young brother expressed a desire to be rid of his thoughts.
Thereupon the old father said.
Dear brother, you cannot prevent the birds from flying in the air over your head.
But you can prevent them from building a nest in your hair.
So as saint augustine says we cannot prevent offenses and temptations.
But by prayer and invocation of the help of god we can prevent them from overcoming us.
This is from his explanation of the lord's prayer and the sixth petition, which is lead us not into temptation.
So again.
What he's basically saying is we cannot prevent.
Thoughts from coming up in our mind. We cannot prevent the bird from flying in the air over us.
But what we can control is to prevent the bird from building a nest on our head.
or prevent those kind of negative thoughts to.
Like take roots in our mind.
So again, it is not a sin to have sexual desires or experience temptation.
But when we give into the temptation or let our passions and desires take control over us.
That becomes an issue.
So that is my first clarification.
second clarification.

$^{201}$Sex and sexuality are part of god's design. They're not something that is dirty or sinful.
But quite the opposite.
They're actually meant to be good gifts from god for us as human to celebrate to enjoy.
Unfortunately, the reality is that just like the rest of the creation this good and perfect gifts from god has been tainted by sin.
You've probably heard of ideas such as um, I can be whatever I want to be or I can do whatever I want.
as long as it makes me happy or.
More recently my body my choice.
And of course, I don't have time in my sermon to dive into all these different.
issues on gender and sexuality.
But I want to point you to a resource if you're interested to learn more about how we as christ followers.
We can approach this topic. So this is a book called love thy body by nancy piercey.
And in this book she talks about issues all the way from abortion.
homosexuality and transgenderism.
to the use of porn to divorce and.
Cohabitation and one of the reasons that I appreciate this book is because the author.
nancy.
She approached all these topics.
Both with a strong biblical conviction and a deep compassion for the people at the same time.
So if you want to pick up a new book to read for the new year, I would recommend.
Highly recommend this book love thy body by nancy piercey.
So.
back to our clarification, so the biggest red flag.
and the most glaring issues with this kind of attitude of I can be wherever I want to be or my body my choice.
Is that when god created the world?.
There is a certain grain or certain flow to it. Just like the grain on a piece of wood.
I'm sure all of you have seen a piece of wood, right or even like touch them.
If you have never seen a piece of wood with some like wooden structures behind that sound booth.
You can go over there and observe it, uh later.
But if you have a piece of wood and you look at it carefully.
You probably see a certain like grain and flow to it, right?.
And if you touch it and kind of like move your hands along the grain, it probably feels relatively smooth.
But if you try to go against the grain.
A grain against the flow and kind of like move your hands on the opposite direction.
You probably might get hurt or even get some splinters, right?.
So in the same way.
In our universe when god created it, there are certain natural orders.
Patterns and inherent structures that is within the universe.
And when we live according to it, we will experience the flourishing life that god has promised.
But when we choose to go against this grain of the universe.
We will experience dysfunction.

$^{241}$and destruction.
So let's take gravity for example, so gravity it is always pulling things towards where.
Towards the center of the earth, right?.
You can never change that it is the fact when god created the world gravity is always pulling things towards the center of the earth.
So you can't just complain to god. Hey god, this gravity is too restrictive.
Why is it only pulling things in one direction?.
I want to be able to fly I want to go against gravity.
well, of course.
You can choose to you can still choose to do whatever you want by jumping off a cliff and trying to fly.
But you might actually end up dying doing that.
So this is the same thing when it comes to sex and sexuality.
There is a certain order.
And design to it.
That we have to live according to if we want to truly experience the flourishing life that jesus has promised.
So.
I just spoke to some of the issues that i've seen in our current culture.
Now I want to take some time to speak to us as the church and some of the things that we can do better as well.
So as the church.
as christ followers.
Instead of telling the non-believing world that they get it all wrong.
And kind of like judging them because we think that we are right.
I think we can actually do a better better job in witnessing to the world.
The beauty of god's design for sex and sexuality instead of just pointing out what the world is doing wrong.
So when I was a student at iit I had an opportunity to.
Go to ecuador for a summer mission with crew.
And during the middle of our six weeks mission trip.
We had a mid-mission retreat at this place called banyos in ecuador. I know it is banyos is also bathroom for spanish.
Uh in spanish, but this place is far off from looking like a bathroom.
This is a very popular tourist attraction.
Kind of like on the mountain side in ecuador.
And one of the things that we get to do when we were at banyos is that we got to hike up at the summit of.
Tungurahua, which is an active volcano in ecuador.
So I remember when we arrived at the summit it was probably around like sunset around evening at that time.
And the view up there it was absolutely stunning.
Like I was just surrounded by all these great mountains and in the midst of it.
I just feel like so small and you just look out and you see the clouds and also like the sun like.
The sunlight the golden rays from the sunlight just kind of like piercing through the clouds.
So I actually have a few pictures um to show.
So these are some of the pictures that I took on that trip.
So the pictures on the left is kind of like my point of view when I got up to the summit.

$^{281}$So if you look right in front of you, that is pretty much the volcanoes.
And then if you look a little bit to the right there's like this treehouse kind of structure over it, right.
And on that treehouse, there's this swing called swing at the end of the world. So there's a picture on the right.
So when you actually go on this swing, you kind of you pretty much just like swing off the cliff.
into like.
Just like I don't know how many feet drop like below you.
but again.
I know it is a little risky, but it is actually the view out there is also like really really gorgeous.
You just look at all the mountains below you all the trees and all the valleys.
And that is probably like some of the most like some of my most favorite place that I have visited so far.
So imagine.
If I when I share with you about my experience visiting, uh, banyos.
Instead of focusing on the beauty of this place.
All the things that I tell you.
are the negative part like oh, it is probably a little risky and dangerous because.
This volcano this mountain is one of the most active volcanoes in ecuador.
It may erupt at any time and I don't know what will happen.
or hey, you have to be really careful out there and.
probably skip the swing because.
if.
You fall off from the swing you might just die.
So imagine if this is all I tell you about this place.
Do you think this place will still be compelling enough for you to visit?.
So in the same way when we talked about god's design for sex and sexuality.
I think we have to do a better job in highlight the beauty of god's creation and god's design.
And make it more compelling instead of just focusing on all the negatives.
So what is sex supposed to be according to god's design.
It is a unity fact between a man and a woman within the covenant of marriage where two become one flesh.
And this is a union not just the physical body.
but also of the mind.
the emotion the soul and the spirit.
More importantly.
It is supposed to point forward to the relationship.
between christ.
And his church.
And this is a beautiful image.
So church.
We need to do a better job in reclaiming a godly view on sex and sexuality.
And make it more appealing to the world around us.
We need to reclaim a godly view on sex and sexuality.

$^{321}$One that is in alignment with god's design.
And then show it to the world. So that is our first point.
So second point we need to ruthlessly eliminate sin. So take a look with me at verse 29 and 30.
If your right eye caused you to stumble.
Gouge it out and throw it away.
It is better for you to lose one part of your body than for your whole body to be thrown into hell.
And if your right hand causes you to stumble cut it off and throw it away.
It is better for you to lose one part of your body than of your whole body to go into hell.
So you might be wondering.
Didn't jesus just emphasize what is truly important?.
It's not our outward actions, but the desires and intention of our heart.
Why is he all of a sudden talking about the eye and the hand again?.
Well, yes jesus cares about the true intention of our heart.
But right here jesus is trying to show how seriously we should treat sin.
And the consequences if we take sin lightly.
So in both of these hypothetical situations.
We find the phrase cause you to stumble if your right eye caused you to stumble or if your right hand caused you to stumble.
This phrase caused you to stumble comes from the greek word scandalito, which is where we get our english word scandalous.
And some english translation. They actually translate it as causing you to sin.
But here it means more than just simply causing someone to commit a sin.
Scandalito here carries the nuance of causing someone to steer away from the path that leads to eternal life in jesus.
towards the life.
That leads to eternal destruction.
And how do we know this if you look at the consequences? It is hell or gehenna the hellfire awaiting.
If we don't deal with the stumbling block that causes us to sin.
And when you look at the solution that jesus provide for us that he instructed us to do it is pretty violent and aggressive, right?.
Why because when it comes to dealing with sin we need to be decisive.
and ruthless.
Jesus, he's not trying to tell us to literally.
Gouge our eyes or to amputate a limb.
But he is using hyperbole to get our attention in order to make his point.
That we need to ruthlessly eliminate sin if we want to follow him and become more like him.
When the apostle paul speaks to the church in ephesus on how they should live.
In light of their new identity in christ. He warns them do not give the devil a foothold.
If we leave even just a tiny bit of foothold to the devil and think that uh, we will probably be all right.
I want you to think again.
I'm sure you all have heard of the stories of david and beth sheba.
and a lot of time we kind of uh.
When we think of these stories we think of it as kind of like the starting point of david's downfall, right?.
But I think this david downfall actually begins way earlier than that.

$^{361}$The story of david and beth sheba begins with king david being complacent during a time of war.
So imagine during a time of war.
Second samuel told us that instead of going out to war with his men king david. He actually stays back in his palace.
in jerusalem.
And on this particular day, he woke up not in the morning.
Not in the afternoon.
but in the evening.
A king during a time of war should be on high alert and on standby mode all the time, right?.
But king david here he's doing the opposite.
He is sleeping in his palace until the evening.
This shouldn't be the routine of a king during a time of war.
However, david he's probably just getting complacent because god has given his armies victories.
after victories after victories over all the different enemies.
And when david leaves this tiny foothold for the devil by being complacent.
On that particular evening it led to more and more serious sin.
beginning with him lasting over beth sheba.
And then committing adultery with her.
And then trying to cover up his own mistakes.
And when his plan to cover up fell.
David ultimately came out with the plan to have besheba's husband uriah killed in the battlefield.
So do you see how serious the consequences can be if we leave just even a tiny foothold for the devil.
This is why we need to ruthlessly.
eliminate sins in our lives.
Especially the parts that we have control over like resisting temptation.
and taking control over.
our passion and desires.
However, I understand that this is easier said than done a lot of times.
I will be the first one to admit that there are probably way too many times where I myself.
failed to resist temptation.
Regardless of how hard I tried.
Can you somewhat relate to that?.
And this brings me to my final point we need we need a transformed heart to live rightly.
We need a transformed heart to live rightly.
We can't live rightly by just not breaking the law or by trying to follow the law.
But we need a transformed heart to live rightly.
Throughout israel history in the old testament the people of god.
failed to live up to god's law generations after generations which ultimately led to the destruction of jerusalem.
And also the exiles in babylon.
However in the midst of this nationwide tragedy.
We find a glimpse of hope.

$^{401}$Where god promised that he would renew the covenant with his people.
And transform their heart.
So this is in the book from the book of jeremiah. So jeremiah 31.
verse 31 to 34.
The days are coming declares the lord when I will make a new covenant with the people of israel.
And with the people of judah.
It will not be like the covenant I made with their ancestors.
When I took them by the hand to lead them out of egypt.
Because they broke my covenant though. I was a husband to them declares the lord.
This is the covenant I will make with the people of israel.
After that time declares the lord I will put my law in their minds and write it on their hearts.
I will be their god and they will be my people.
No longer would they teach their neighbor or say to one another.
Know the lord because they will all know me from the least of them to the greatest.
declares the lord.
For I will forgive their wickedness and will remember their sins no more.
so.
What what the lord is saying here is basically under this new covenant.
The law of god will no longer be inscribed on stone tablets.
But it will be inscribed in our hearts in our mind.
And church we are the people who live under this new covenant.
The same spirit who raised jesus from the dead is also living in us now giving us the power to overcome sin in our lives.
So in order to live rightly before the lord we cannot do it in our own strength or in our own ability.
But we need a heart that is being transformed by the power of the gospel.
And this is the good news that jesus brings to us that we're not only free from the penalty from of sin.
But we are also free from the power of sin over us.
And we know that one day when jesus comes again, we will be free from the presence of sin forever.
And this is why I think that jesus he is very intentional in starting his kingdom teaching.
On the beatitudes which pretty much focus on the attitude and the posture of our hearts.
By being poor in spirit by being humble and meek.
This is because jesus he understands that we need a transformed heart to be able to live rightly before the lord.
Not just by trying to not break the law or trying to follow the law perfectly.
That is the way of the pharisees and the teachers of the law.
So.
As we begin the new year. I think this is a good time for us to set our hearts in the right place.
And allow god to transform it in order to live rightly.
Not through our own strength, but through the work of jesus who has fulfilled who has completely who has perfectly fulfilled.
The righteous requirement of god and also through the power of the holy spirit living in us.
So I have some reflection questions for us here to reflect.
As we kind of like begin our new year.

$^{441}$So what is god saying to me this morning?.
And what am I going to do about it?.
Are there certain areas or sin struggles in my life where I need the power of jesus to transform my heart.
In order to live rightly before the lord.
so I would encourage you all to take some time or just.
Take a picture of these questions and maybe in the following week just take some time to reflect on these questions.
and as pastor mark say.
Um, this is the new year and we want to be able to have a good.
um kind of like.
Yeah a good routine in our spiritual walk in our walk with jesus as well if we want to truly follow jesus.
So, um, let's pray.
Heavenly father.
We thank you that uh, you are not a god.
who punish sins.
But you are a god.
Who show us grace even though we?.
As human we have time and time again failed to live up. Um.
to your standard of righteousness.
So jesus we thank you for this gift of righteousness that you have given us at the cross.
That we are able to come before you.
Just as we are.
Carrying all the weight and the burdens.
And laid it in front of your throne.
Father as we begin the new year, I just ask that you would give each and every one of us here in this room.
Just a desire to know you more just a desire.
To live in a way.
That uh represents our kingdom citizenships that represents who we belong to.
So grant us the power through your holy spirit.
To be able to live rightly before you and to be able to witness the world around us.
the beauty of god.
The beauty of who you are jesus.
So help us jesus in your son's name. We pray.
amen.
So thank you for coming this morning again. Happy new year and we have some refreshment outside for you to enjoy.
So I would encourage you to stay back. Enjoy some fellowship with us.
and.
You are dismissed.
\newpage



\section{}
\label{sec:as4vlJLVesI}
\textbf{5/6/18 主日講道: 我如何從回教徒成為基督徒 (Pastor Wissam Youssef )}
\newline
\newline
連結: \href{https://youtube.com/watch?v=as4vlJLVesI}{\texttt{https://youtube.com/watch?v=as4vlJLVesI}} ~~~~ 語音日期: 2025-04-25
\newline
\newline
\hyperref[sec:TstKn_BEGfA]{\small{< < < PREV SERMON < < <}}
~
\hyperref[sec:index_chronic]{\small{[返順時目]}}
~
\hyperref[sec:index_scriptual]{\small{[返順卷目]}}
~
\hyperref[sec:Ml1uo7s2nGU]{\small{> > > NEXT SERMON > > >}}
\newline
\newline
$^{1}$[speaking Chinese].
(speaking in foreign language).
- Well, we have a very special speaker today, all right?.
And Wasam asked me, is he the first Arab.
in our congregation, I think, right?.
You're the first one, right?.
Iraqi, right?.
And he's a native Iraqi, and he's from Baghdad, Iraq..
How many of you have been to Baghdad, Iraq?.
(audience laughing).
One, okay, all right..
And he, through a very dramatic experience,.
came to Christ in Baghdad, Iraq..
And he will tell you more about his story..
Now he's Associate Minister of Sunset Church of Christ.
in Taylor, Michigan..
Taylor, Michigan, you've been there, right?.
Michigan, you've been there, right?.
State of Michigan..
And his vision is to tell Muslims about Christ.
and to tell Christians about Islam..
And I think this is a very relevant topic today..
Our coworkers, our friends, our neighbors,.
maybe Muslims, or maybe they believe in Islam,.
and they love to tell you about their faith.
and how to respond as a Christians..
So a very, very relevant topic..
And Wasam has also written two books..
The first one is "I Am an Arab and I Am a Christian.".
(speaking in foreign language).
The second book he has written is "Islam in Christ's Eyes.".
It's a scripture study, it's a group study guide.
to help you understand Islam from a Christian perspective.
and how to respond to Islam..
(speaking in foreign language).
Wasam, we wanna welcome you..
- Thank you..
(audience applauding).
Thank you..
Good morning..

$^{41}$(speaking in foreign language).
It is so surreal that I'm invited in one weekend.
to talk in two Chinese congregations..
(speaking in foreign language).
I started the Arab Christian ministry..
(speaking in foreign language).
Upon my graduation.
from Sunset International Bible Institute..
(speaking in foreign language).
Locally, I work from my church, preaching the gospel.
to the predominantly Muslim population.
in the neighboring Dearborn, Michigan..
(speaking in foreign language).
And globally, I work as the adjunct assistant dean.
for the Arabic studies.
with Sunset International Bible Institute..
(speaking in foreign language).
And yet, I think you are never famous.
until you have your name said in Chinese..
(audience laughing).
(speaking in foreign language).
I was born to a Muslim family in Baghdad in 1979..
(speaking in foreign language).
A few months after my birth,.
Saddam Hussein became our president..
(speaking in foreign language).
And Ayatollah Khomeini led the Islamic revolution.
in the neighboring Iran..
(speaking in foreign language).
And we started to have the war.
between the two countries that lasted for eight years..
(speaking in foreign language).
I was born and raised in a classic example.
of a world that was not told to love one another..
(speaking in foreign language).
A world that is not shaped by the goodness.
and the love of Jesus Christ..
(speaking in foreign language).
Shortly after the war, Saddam Hussein invaded Kuwait..
(speaking in foreign language).

$^{81}$We had the sanctions that lasted for 13 years.
in which my whole family lived on \$1 a month..
(speaking in foreign language).
And we had the Gulf War of 1991..
(speaking in foreign language).
Up until that point,.
Saddam Hussein was a secular socialist military dictator..
(speaking in foreign language).
But after he had a Shiite revolt in the south,.
(speaking in foreign language).
he realized he could not impose a secular system.
on the predominantly Muslim people of Iraq..
(speaking in foreign language).
So after he crushed that revolt,.
he started to put on a show of Islam..
(speaking in foreign language).
He wrote Allahu Akbar on the Iraqi flag..
(speaking in foreign language).
He ordered the Quran to be taught.
in all the Iraqi schools..
(speaking in foreign language).
He built hundreds of mosques all around Iraq..
(speaking in foreign language).
And he started to turn the Iraqi community.
from being secular to being more Muslim..
(speaking in foreign language).
My father retired in 1992..
(speaking in foreign language).
And my mom started working as a seamstress,.
mending the clothes of our neighborhood.
to provide for the family..
(speaking in foreign language).
I was in the middle of my teenage years,.
in the middle of the 1990s..
(speaking in foreign language).
When Iraq was officially at enmity.
with everything that is different in the world..
(speaking in foreign language).
We did not have the terrorism and the car bombs.
and Al-Qaeda that is ravaging the Middle East now..

$^{121}$(speaking in foreign language).
But all of that was a culture that was preached in mosques..
(speaking in foreign language).
Taught in schools..
(speaking in foreign language).
And propagated by the government's limited channels..
(speaking in foreign language).
Hate was preached everywhere..
(speaking in foreign language).
It is hate that eventually caused me.
to convert out of Islam..
(speaking in foreign language).
Hate is evil, whether it is in Iraq.
or in America or in China..
(speaking in foreign language).
We did not starve for food or water.
because of the sanctions years..
(speaking in foreign language).
But we were starving for a different kind of bread.
and a different kind of water that we did not even know..
(speaking in foreign language).
My four year old sister passed away in 1996..
(speaking in foreign language).
My mom, a panicked mother,.
and the funeral of her daughter was--.
(speaking in foreign language).
Was begging for someone at the funeral.
to provide her with a glimpse of hope.
that that innocent child went to heaven..
(speaking in foreign language).
No one was able to provide her with that hope..
(speaking in foreign language).
And the only man who was knowledgeable in the Quran.
told her she will be judged for the very first breath.
she took in when she was born..
(speaking in foreign language).
A few months later that same year..
(speaking in foreign language).
My uncle got cancer..
(speaking in foreign language).

$^{161}$He was the uncle that I loved most..
(speaking in foreign language).
He was so kind and so generous and so humorous..
(speaking in foreign language).
But when I last saw him,.
he would cry every time I talk with him..
(speaking in foreign language).
Not only because he saw his death coming,.
which he did in a few days..
(speaking in foreign language).
But for the first time in his life,.
he was unable to go to the bathroom..
(speaking in foreign language).
To get ritually clean for the Muslim prayers..
(speaking in foreign language).
He was in bed, soiling himself..
(speaking in foreign language).
He thought God was done with him..
(speaking in foreign language).
And he died without hope..
(speaking in foreign language).
That's when I decided that's it,.
between me and God,.
rejected Islam and became an atheist..
(speaking in foreign language).
I thought I would be happy if I were an atheist..
(speaking in foreign language).
But somewhere deep inside me,.
I was looking for justice and for a sense of direction..
(speaking in foreign language).
I wanted to entertain and distract myself..
(speaking in foreign language).
I started to think about the world.
(speaking in foreign language).
I started reading books,.
listening to music and watching movies..
(speaking in foreign language).
Most of the books that I read quoted the Bible..
(speaking in foreign language).
I got curious to get myself a copy of the Bible.

$^{201}$only to understand these quotes..
(speaking in foreign language).
I was in my second year in college.
studying civil engineering..
(speaking in foreign language).
When the first Mission Impossible movie was released..
(speaking in foreign language).
Tom Cruise picked the Gideon's Bible and read Job 314..
(speaking in foreign language).
That was the first time I see the Bible..
(speaking in foreign language).
I got so curious..
I went to a flea market in Baghdad.
and bought my first Bible..
(speaking in foreign language).
It was the Gospel of John..
(speaking in foreign language).
I thought that was the entire Bible..
(speaking in foreign language).
I did not find Job 314 in the Gospel of John..
(speaking in foreign language).
So I thought Muslims were right..
Everybody is making up a different Bible..
(speaking in foreign language).
Something caught my attention in the Gospel of John though..
(speaking in foreign language).
The way Jesus addressed.
the hypocritical religious authorities of his time..
(speaking in foreign language).
Something you're not supposed to do in a Muslim community..
(speaking in foreign language).
I was born and raised to two devout Muslim parents..
(speaking in foreign language).
And they taught me that every Imam.
is a steward of God's word..
(speaking in foreign language).
And we should not question them..
(speaking in foreign language).
Well, not Jesus in that book..
(speaking in foreign language).

$^{241}$It was the summer holiday.
between the second and the third year in college..
(speaking in foreign language).
When I was back at that flea market.
flipping the New Testament..
(speaking in foreign language).
I saw the Gospel of John in the New Testament..
(speaking in foreign language).
That was the first time I realized.
that the Gospel of John was not another Bible..
It was part of a bigger book..
(speaking in foreign language).
So I bought the New Testament..
(speaking in foreign language).
A few days later,.
I realized that the New Testament.
was part of another bigger book..
(speaking in foreign language).
So I realized that I lost a lot of money already..
(speaking in foreign language).
And bought my first complete Bible..
(speaking in foreign language).
I was so shocked to see the accurate historical.
and geographic description in the Bible.
of a world that I am familiar with..
(speaking in foreign language).
Babylon, that's 50 miles to the south of my parents' house..
(speaking in foreign language).
Babylonians and Assyrians,.
these are the native Iraqis..
(speaking in foreign language).
Nebuchadnezzar, his poster is next to Saddam Hussein's.
picture in every street corner in Baghdad..
(speaking in foreign language).
The captivity, the pride of the Iraqi history.
that we studied in the seventh grade..
(speaking in foreign language).
I'm sorry..
(speaking in foreign language).
And I was thinking, man, the guy who made up this book.

$^{281}$really knew his history and geography..
(speaking in foreign language).
A Christian philosopher once said..
(speaking in foreign language).
That there are three things that we all need..
(speaking in foreign language).
And these are first, truth..
(speaking in foreign language).
And then goodness..
(speaking in foreign language).
And then beauty..
(speaking in foreign language).
But he said what we want is the exact opposite order..
(speaking in foreign language).
We want beauty first..
(speaking in foreign language).
And then goodness and then truth..
(speaking in foreign language).
The beauty of the Bible caught my attention first..
(speaking in foreign language).
It is the most influential and most beautiful piece.
of literature that was ever written..
(speaking in foreign language).
Then I fell in love with the goodness of the Bible..
(speaking in foreign language).
A book that preaches love and not hate..
(speaking in foreign language).
Peace and not violence..
(speaking in foreign language).
Forgiveness and not vengeance..
(speaking in foreign language).
Integrity, benevolence, you name it..
(speaking in foreign language).
And I fell in love with the Bible for being good.
way before I believed in the Bible for being true..
(speaking in foreign language).
I had so many questions to ask..
(speaking in foreign language).
I did not have any Christian friends..
(speaking in foreign language).

$^{321}$But when the third year in college started,.
I saw a man with a cross in his shirt in our class..
(speaking in foreign language).
He obviously flunked the third year..
(speaking in foreign language).
I asked him, are you a Christian?.
Do you know your Bible?.
(speaking in foreign language).
He said, go ahead..
(speaking in foreign language).
I said, do you Christians believe that you are saved.
by God's grace or by your obedience?.
(speaking in foreign language).
Like every other preacher,.
he did not directly answer my question..
(speaking in foreign language).
But he took the opportunity.
to introduce the gospel message to me..
(speaking in foreign language).
He explained to me how God is a holy and just God.
that does not tolerate sin..
(speaking in foreign language).
But at the same time, he's a loving God,.
so he sent his son to take our sins on him.
and die on the cross for us..
(speaking in foreign language).
It's a combination of God's grace and our obedience..
(speaking in foreign language).
I said, interesting..
One more question..
(speaking in foreign language).
Do you really believe that?.
(speaking in foreign language).
He said, oh yeah, we do..
(speaking in foreign language).
Next day, he brought me some literature.
that talked about how many tens of thousands.
of Bible manuscripts we have today in the museum..
(speaking in foreign language).
From different times, from different places,.

$^{361}$written in different languages..
(speaking in foreign language).
And they all match each other..
(speaking in foreign language).
That was the first time anyone tells me.
that the Bible has not been corrupted..
(speaking in foreign language).
I was 19..
(speaking in foreign language).
I was so happy to learn that the book.
that has the solution to every problem.
my community is suffering from turned out to be true..
(speaking in foreign language).
I was so happy to learn that the book.
that was both beautiful and good turned out to be true too..
(speaking in foreign language).
I thought everybody else would be as happy..
(speaking in foreign language).
I was wrong..
(speaking in foreign language).
The first person I tried to share the gospel with.
was the woman that I loved the most..
(speaking in foreign language).
My mom used to wake up before me.
to make breakfast for me every day..
(speaking in foreign language).
And at that time, she would not answer me.
when I say good morning, mom..
(speaking in foreign language).
One morning I asked mom, what's wrong?.
(speaking in foreign language).
She turned to me and she grabbed me by the collar..
(speaking in foreign language).
And she asked me, please tell me.
what did I fail to provide to you?.
(speaking in foreign language).
That you do this thing to me..
(speaking in foreign language).
I said, what thing?.
(speaking in foreign language).

$^{401}$She said, I've been working day and night.
to make a human being out of you..
(speaking in foreign language).
I've been wearing my eyes out on that sewing machine..
(speaking in foreign language).
To send you to school..
(speaking in foreign language).
And you've been using my money to buy that blasphemous book.
that I've been seeing among your college books..
(speaking in foreign language).
Look at yourself, grow up..
(speaking in foreign language).
With that,.
with the lack of wisdom and the lack of encouragement..
(speaking in foreign language).
With the fact that I started to lose.
many of my Muslim friends.
and did not gain any new Christian friends..
(speaking in foreign language).
And with the fact that I started to get harassed.
by the religious authorities in my neighborhood..
(speaking in foreign language).
I came to that point where I thought,.
maybe this is all wrong..
(speaking in foreign language).
Maybe there is no God..
(speaking in foreign language).
Maybe the gospel was just another philosophy.
I wanted to entertain and distract myself with..
(speaking in foreign language).
I renounced the Christian faith and gave away my Bible..
(speaking in foreign language).
A few months later,.
(speaking in foreign language).
I realized,.
(speaking in foreign language).
that the hope and the satisfaction,.
(speaking in foreign language).
and the fulfillment and the assurance,.
(speaking in foreign language).

$^{441}$that I had when I claimed God's promises.
in the Bible as my own,.
(speaking in foreign language).
cannot be replaced with anything in this life..
(speaking in foreign language).
I'd rather be a persecuted Christian with hope,.
(speaking in foreign language).
than a comfortable non-believer with no hope..
(speaking in foreign language).
So I got a new Bible..
(speaking in foreign language).
In English this time, so my mom would not understand it..
(audience laughing).
(speaking in foreign language).
And I continued to read the Bible in secret.
for the rest of my stay in Baghdad..
(speaking in foreign language).
My biggest discouragement was that it would take me 12 years.
going to every single church building in Baghdad.
trying to get baptized in faith..
(speaking in foreign language).
I finished school in 2000..
(speaking in foreign language).
Finished my service as a Republican guard in 2002..
(speaking in foreign language).
America invaded Iraq in 2003..
(speaking in foreign language).
We started to have the internet.
and we started to be open to the outside world..
(speaking in foreign language).
I would go online and download sermons and Christian songs.
and upload them to my MP3 player.
and listen to them on my way to work every day.
and read the passage from the Bible..
(speaking in foreign language).
Until that one morning in 2009..
(speaking in foreign language).
When my workplace, the Iraqi Department of Public Works..
(speaking in foreign language).
Was rocked by a truck bomb..

$^{481}$(speaking in foreign language).
I was in the staircase, I came out of it without a scratch..
(speaking in foreign language).
We moved to a substitute building.
in another neighborhood in Baghdad..
(speaking in foreign language).
That building had an internet cafe next to it..
(speaking in foreign language).
One afternoon I signed up for a free Bible study..
(speaking in foreign language).
And that's when everything started to change..
(speaking in foreign language).
A woman was assigned to grade my Bible tests..
(speaking in foreign language).
She asked me in her first email if I had any prayer requests..
(speaking in foreign language).
I said, "Would you please pray that I can get baptized?.
"I've been a believer for 12 years already.".
(speaking in foreign language).
And would you please pray that I can live.
in a Christian-friendly community.
where I can say Jesus is Lord.
without getting killed right on the spot..
(speaking in foreign language).
That same day I came across a free book called Bible Basics..
(speaking in foreign language).
Written by a British man who had a ministry in Latvia..
(speaking in foreign language).
I ordered a free copy of that book..
(speaking in foreign language).
He ended up coming to Iraq and baptizing me.
in a bathtub in Kurdistan, Northern Iraq..
(speaking in foreign language).
The woman that used to grade my Bible tests.
was a caregiver for a quadriplegic man.
in a Christian man's house in Pennsylvania..
(speaking in foreign language).
The Christian man told her if that Iraqi plans.
to come over here, he can stay with us..
(speaking in foreign language).

$^{521}$A few weeks later, the American embassy in Baghdad.
started issuing visas for Iraqis.
for the first time in 20 years..
(speaking in foreign language).
So I applied for the visa..
(speaking in foreign language).
Bought the plane ticket..
(speaking in foreign language).
Packed one bag..
(speaking in foreign language).
Said the final goodbyes to my family..
(speaking in foreign language).
Flew all the way from Baghdad to New York City..
(speaking in foreign language).
Came to this nation with one bag..
(speaking in foreign language).
With only enough money to survive for weeks..
(speaking in foreign language).
And with one contact, a poor woman in the Amish country..
(speaking in foreign language).
One Christian man there heard my story..
(speaking in foreign language).
He had me speaking in his church in Pennsylvania..
(speaking in foreign language).
The preacher of which turned out to be a graduate.
from the school I eventually ended up going to..
(speaking in foreign language).
He raised support money for me to go to school.
on the time I ran out of the last dollar.
that I brought with me from Iraq..
(speaking in foreign language).
I finished school and moved to the Detroit area..
(speaking in foreign language).
Back to the Muslim community,.
the same community I escaped from..
(speaking in foreign language).
Eight years ago, I used to hide my faith to stay alive..
(speaking in foreign language).
And today, I share my faith to make a living..
(speaking in foreign language).

$^{561}$Once again, appreciate your invitation..
(speaking in foreign language).
And your time and thank you and God bless..
(speaking in foreign language).
(audience applauding).
Thank you very much..
Well, we still have five minutes.
and maybe there are some questions from the congregation.
and we wanna ask you or you wanna answer.
if you have an answer..
I am all yours..
(audience laughing).
Do you guys have any questions?.
Yes, ma'am..
- Oh, Emily, please get a mic..
- Can you tell us how bad they would treat Christian.
back in like Iran or Baghdad that you used to,.
well, before you leave?.
- There is much more to,.
well, they are persecuting Christians in Baghdad.
that you would know..
Okay, learn the basic facts first..
When I was born and raised,.
that was under Saddam Hussein.
who was a member of the Ba'ath Party.
that was established by a Greek Orthodox Christian..
The Iraqi Prime Minister,.
the international face of Iraq, Tariq Aziz,.
was a Catholic Christian..
And so there has not been any persecution.
on a religious basis under Saddam Hussein..
Now, after Saddam Hussein,.
we still have a civil constitution.
and we still have a lot of Christians..
In fact, my boss in Baghdad when I was an engineer.
was a Catholic woman,.
but the government is not strong enough.
to impose law and order..
When a car bomb explodes in an Iraqi marketplace.
and kills 99 Muslims and one Christian,.

$^{601}$that story will be reported somewhere.
as a Muslim persecuting Christians..
I think there is much more to the stories.
that you can hear, at least from Baghdad,.
from what I saw in Baghdad,.
than what you see from your end..
Christians are being persecuted everywhere.
and more severely in the Muslim countries.
than here in the United States,.
but there is much more to the story..
- Any other questions?.
(audience laughing).
How many-- - Yes, ma'am..
- Oh, sorry..
- I'm interested to know how did you tell your mom.
and your family when you leave the country?.
- Okay..
Good question..
My mom first did not take me seriously..
Then she took me seriously and she was really sad..
Then right before I left, she said,.
"I know that you are doing the right thing,.
"but promise me one thing..
"You do not join a drug dealing gang.".
(audience laughing).
So that's the only thing she knew about America..
I said, "No, Mom, I'm not gonna do that..
"I would never do that.".
Unfortunately, things in Iraq started to go.
from bad to worse at that time..
It's not as bad anymore..
So ISIS came into the scene and more people got killed..
And my mom eventually said, "You have always been wise..
"I have never told you that,.
"but I knew that you were doing the right thing.".
My whole family knows about my conversion.
and my work now and they are fine with it..
In fact, the first book that I wrote,.
the book that exposes Islam, actually that's the second book..
The first book is "Islam in Christ's Eyes.".

$^{641}$I told my mom, "I wrote a book that exposes your faith.
"as bad and false and I dedicate it to you.".
She said, "I'm proud of you.".
(audience laughing).
- Right there, Steve..
- Sir..
- Can I get a mic?.
- You made your mom change her view..
- Time..
- Time?.
- Yeah, you only have one chance.
to offend somebody with the gospel message..
Two decades later, no one really cares..
(audience laughing).
- So you cast the bombardment.
of the gospel message, basically?.
- Well, no, my family are still all devout Muslims..
And yeah, I mean, to be tolerant and receptive.
and respectful is one thing..
To, for their hearts to be convicted of their sin.
and for them to really feel that they need a savior,.
that is a different thing..
It's always harder with the family..
I mean, when I first started to tell my mom about the gospel.
she told me, "Are you saying that you are right.
"and this whole world is wrong?".
To her, the whole world is Islam..
And so, yeah, I mean, it took time..
That's the only thing..
Yeah..
- Well, as I said in the introduction,.
we have some friends or neighbors or co-workers.
who are in the Islamic faith..
We are Christians and what are the pointers.
that you can share with us, how to share the faith,.
how to connect with a Muslim neighbor?.
- Okay, long story short, show them the love of Christ.
before you show them the truth of his gospel..
'Cause people are interested in what is beautiful first.
and then what is good and then what is true..

$^{681}$I have a whole course on that called Islam in Christ's Eyes.
that starts with, see, when you preach the gospel to anyone.
you will have to start where they are at now..
And so what Muslims believe now and where to start,.
the theology of Islam explains the element.
of the Islamic faith, I talk about that..
And to understand how things ended up.
the way they are at now in the Muslim world,.
you will have to have a basic exposure.
to the history of Islam..
So history and then theology and then how to preach.
the gospel to Muslims, this all is a course.
that is called Islam in Christ's Eyes.
that I'm giving now in, what's the name of the congregation,.
the other Chinese congregation that I'm talking to?.
They have the name in Chinese characters,.
I have no idea. (audience laughing).
And they gave me a ride here and a ride there..
And so, before you share the gospel,.
okay, before you ask how can I share the gospel,.
you will have to ask an equally important,.
if not a more important question,.
why should I share the gospel with them?.
What is the motive when you want to share the gospel?.
What is the proper motive if you wanted to share the gospel.
with your non-Christian neighbor?.
Here is the right motive..
When you share the gospel with your neighbor,.
Muslim or non-Muslim, remember that if they responded.
to that message of salvation,.
you will be stuck with them for all the eternity..
Is this what you really want or not?.
If you don't love them that much,.
if you cannot stand seeing them in your neighborhood,.
don't you dare preach the gospel to them..
Move on and preach the gospel to somebody else.
that you really like..
So start with the proper motive first,.
and then the means, and then the opportunity..
And no, we are not trying to murder anybody here..

$^{721}$(audience laughing).
- All right, so let's give a round of applause to Wissam..
- Thank you..
- Why don't we pray together?.
- Sure, please..
- We pray for his ministry in Michigan..
We also pray for more Arab people coming to Christ..
- Thank you..
- All right, let's pray..
Father God, we just want to give you thanks.
because you are the God of the whole world,.
and you love every one of us with grace,.
and with love, and with beauty..
We pray for Wissam..
We praise you for how you have led him to come to Christ,.
to come to know the gospel,.
and come to the United States to start the ministry,.
to preach Christ to many Muslims in his neighborhood..
We just pray that you will continue to use him,.
use his story, use his experience, his knowledge,.
to bring many to Christ, for your glory..
(speaking in foreign language).
- Amen..
- Thank you very much..
- I appreciate that, thank you..
(speaking in foreign language).
(audience laughing).
- Oh, wow..
- That's right..
- Is that all the Chinese you can speak?.
(audience laughing).
- Yeah, and I got that from Brush Hour..
- All right, so..
(speaking in foreign language).
Islam in Christ's Eyes..
(speaking in foreign language).
I am a Christian..
(speaking in foreign language).
(silence).
[BLANK AUDIO].

\newpage



\section{1:8}
\label{sec:Ml1uo7s2nGU}
\textbf{Cross Over (Acts 1\_8) - Rev. Mark Chan}
\newline
\newline
連結: \href{https://youtube.com/watch?v=Ml1uo7s2nGU}{\texttt{https://youtube.com/watch?v=Ml1uo7s2nGU}} ~~~~ 語音日期: 2025-04-29
\newline
\newline
\hyperref[sec:as4vlJLVesI]{\small{< < < PREV SERMON < < <}}
~
\hyperref[sec:index_chronic]{\small{[返順時目]}}
~
\hyperref[sec:index_scriptual]{\small{[返順卷目]}}
~
\hyperref[sec:BPEn5KuFa4c]{\small{> > > NEXT SERMON > > >}}
\newline
\newline
$^{1}$about that. So we are in a very new series because we are nine year old. We have already.
celebrated our ninth year last Sunday and I know many of you were here with us and all.
the Easter Sundays are the church's anniversary. So we are nine years old and because we're.
coming to a new year, the ninth year of UV in our church, we would like to do something.
special this coming year. So in this coming year, we are going to centre on one theme,.
going to centre on one topic and this theme, this topic is called "Crossover". Now what.
is this? Okay, so you can see the logo is very beautifully written out, right? I hope.
this coming year you will know how to sign this. It's a cross and there's a heart and.
we're going to go out, right? And that's what we mean by crossover, alright? Qua yue, transform,.
to go beyond our comfort zone and that's the first thing we want to do. We all have our.
own comfort zone, right? Our emotional comfort zone, our social comfort zone and our own.
physical comfort zone, especially after the pandemic, right? Everyone seems to be suffering.
from social anxiety, right? Don't come near, okay? I want to be myself, okay? I want to.
just be isolated, right? And the Bible challenges if we are in the love of God, if we are in.
the grace of God, we need to leave our own comfort zone and that means we need to grow,.
okay? There are many ways we need to grow, emotionally, socially, spiritually, relationally,.
all kinds of things. We need to leave our comfort zone. So when I say comfort zone,.
you probably will automatically think of some of the things that you try to keep to yourself.
and that may be our comfort zone and in this coming year we want to challenge you, we want.
to remind you to step out of your comfort zone, alright? The second thing here is to.
join with other followers of Jesus Christ. In our church there are many sisters and brothers.
from different backgrounds, speak different languages. I told you last year I had the.
goal of improving my Chinese, right? Yeah, the goal failed, okay? But this year it's.
going to be my goal as well, okay? I'm going to improve my Mandarin speaking skill, okay?.
I will continue to do that, alright? I may have said many funny things but I will head.
on, okay? I will continue to improve. So to join with other followers of Jesus Christ.
is to grow together, right? How do I improve my Chinese skill? By talking to you, right?.
And you talk back to me. And we have conversations and we have fellowship and we build our relationship.
together and we try to build a community to join with other followers of Jesus. That's.
the number two. Number three is that if our church continues to grow, if we experience.
the grace of God and we leave our comfort zone and join with others to build a church,.
then we will have more potential, our redemptive potential to reach other people for Christ,.
other neighbours for Christ. So that's the three things I want to remind you this coming.
year. We're going to remind you every single week almost. Leave your comfort zone. Pray.
that God will give you strength so that I can grow and leave my comfort zone and join.
hands with other people and reach out with the gospel, with the love of Christ. And God.
is calling us this year, the ninth year, to grow deeper, to deepen our unity, alright?.
How can we unite together? We have to get to know each other. We have to speak to one.
another and we have to boldly engage in sharing the gospel to our neighbours. So this will.
be a year long focus and we want to take some of the passages from the book of Acts week.

$^{41}$after week. And we're going to preach and we're going to do Bible study on the book.
of Acts. You know how many chapters are in the book of Acts? How many chapters altogether?.
28. 28 chapters, alright? So that should be enough for us, right? There's a lot of materials..
And as you understand, the book of Acts is about the early history, early movement of Jesus'.
followers, the very first generation and what did they do. And the book of Acts is all about.
how they leave their comfort zone and to join with other people and to use the gospel and.
take the gospel to a farther end of the earth. So again, we will do sermons and we will do.
Bible study. We'll do a lot of things, you know, take materials from the book of Acts..
And today we would like to take a moment to think about what the book of Acts is all about..
Alright? Understand the bigger picture. Now in the book of Acts chapter 1, right from.
the beginning, it is Jesus rose from the dead and he stayed with his disciples for 40 days..
And he taught them about the kingdom of God. And there's a kingdom of God, there's a reign.
of God coming to his people. And that is the kingdom project, so to speak. And this kingdom.
project is a huge project. Acts 1.6. It says, "So when they had come together, the disciples,.
and Jesus, after Jesus rose from the dead and before he ascended to heaven, those 40.
days they had come together, the disciples asked Jesus a question. It's a very legitimate.
question. They said, "Lord Jesus, will you at this time restore the kingdom of Israel?".
That's a legitimate question. Because every Jewish boy and girl is dreaming of the restored.
kingdom of heaven, restored kingdom of Israel. You know, in the good old days, in the times.
of David, King David, thousand years before Jesus, there was a kingdom, right? There was.
a political entity that they enjoy independence and freedom. And that's what they had in mind..
And they asked Jesus, "Hey, is it at this time because you rose from the dead and the.
kingdom project has begun, and it's this time that you're going to restore the fallen, the.
lost kingdom of Israel?" They asked Jesus that, and Jesus said to them, "It is not for.
you guys to know the times or the seasons that the Father has fixed by his own authority.".
Alright, so he answered the question saying, "It's not for you to determine the time and.
the manner of the kingdom." So as I said, being Jewish in their blood, the disciples.
were thinking and dreaming about maybe Jesus would restore that kingdom of David, the political.
state of Israel. However, if you read on how Jesus replied to the disciples' questions,.
God has a bigger, bigger plan for his kingdom. God was about to launch a kingdom on earth..
Yes, yes, it can be called Israel, but this kingdom is not just a plot of land on the.
globe. This kingdom will be reaching to the ends of the earth. So God has a bigger plan.
for us. God has a bigger plan for the disciples. So I want you to think about this, okay? Think.
about what God wants from you. If God, in the Son of Jesus, in the Son of God, Jesus,.
rose from the dead, what does he want from you? And you have that new life. And God has.
a bigger plan for you. Now you may think, "Okay, I just attend church. I just attend.
Bible study, eat with other people. They have good food. Today we have good food. That's.
why we come. That's our dream. Maybe today, this morning you wake up, and that's all you're.
coming in for, right? Coffee, food, all this stuff. Maybe some lovely people here." But.
God has a bigger plan. I remember when I was younger, I'm already over 50 years old. I.

$^{81}$remember when I was 20 years old, I started going to church, and I had my own job. I graduated.
with a biochemistry degree, and I was working in the chemistry lab back many years ago in.
California. And I thought, "No, my goal is just to, okay, every week I'll attend the.
church. Of course, that's my basic, right? I'll do that. No problem." But that's pretty.
much it. I have my own job, and I have my own comfort life, and that's my comfort zone,.
right? And maybe here and there I try to help out, try to help out with the worship team,.
or try to help out with the backstage, try to help out with being an usher, anything.
I can do. But obviously, God has a bigger plan for me, right? And God has a bigger plan.
for you. I remember when we started and launched this church nine years ago, right here, last.
week some of you were here, and you saw some of the pictures, right? And we were an empty.
warehouse nine years ago. And I thought, "Okay, we'll just begin a local church, right? We'll.
stay small, you know, we'll do a little bit of things here and there, and that's it, right?".
But little by little, God is showing us what He has in store for us. And I just never thought.
of or never expected that we would come to this day, right? Especially through the years.
of pandemics. Oh, those were very chaotic times, right? I remember we came here to pack.
up everything, you know, the sound equipment, the audio interface, and you know, those hardware,.
and we move it back to my home so that we can do live stream. And it was for a long.
time it was chaotic. But we come back after years of pandemic, and our online presence.
has been growing bigger. So that's why I'm saying God has something bigger for His people.
and for our church, right? So God wants to launch a kingdom, not in the political sense,.
but in the spiritual sense. But look at you guys, right? Where are you from? Some of the.
early conversations that we have with some of the new friends here is that, "Where are.
you from?" Right? Yeah, where's your hometown? We all come from different hometowns, different.
places, different background, we speak different languages, we have different cultures. But.
you know, God has this bigger plan for all of us and draws us together, all right? Now.
I want to repeat this. I want to remind you and emphasize this. If you want to choose.
your good friend or best friend, you wouldn't choose me, all right? And I wouldn't choose.
you, to be honest, okay? Because we're so different, okay? Okay, maybe you find me being.
a good friend, okay? But all I'm saying is God in His will put us together for a bigger.
purpose, you see? It's not by my choice, it's not by your choice, ultimately it's God's.
own choice and it's the Spirit's will to put us together as a body of Christ, as a church.
community and we can do something for God. And God does have a bigger plan for you and.
for me and for our church. So what does the Kingdom Project look like, all right? So Jesus.
is trying to answer this question. And let me break it down. Jesus in the book Acts chapter.
1 is trying to explain to the disciples, this is what God wants to do. This is what the.
Kingdom of God looks like. Number one, every person will be called to believe in the name.
of Christ. To come under the cross, at the foot of the cross, we believe the gospel and.
so that we are transformed by the love of God, by the grace of God. That's number one..
Each of us, each and every one of us, has experienced God's salvation at the cross in.
many different ways but there's a common ground for all of us. We are all in Christ. We are.

$^{121}$all loved by God. We all belong to the sons and daughters of the Kingdom. That's number.
one. Number two, the Kingdom Project is all about each and every one of us, after we experience.
the transformation of the love of God, is joined together with other people, with other.
believers and we combine together, we grow the church body, the community together, right?.
And then ultimately, it's not for us. It's not that, okay, we are a happy family and.
then we're going to close the door and we enjoy our time. We're going to reach out to.
share the gospel to the ends of the earth. And that is the Kingdom Project God has in.
mind. So Jesus is challenging his disciples, his church. Hey, yes, I'm going to launch.
a new kingdom but this new kingdom is not about a plot of land that you can put a GPS.
on, on the map. It's not a small plot of land in Israel-Palestine. He's explaining, trying.
to expand the vision of the disciples. Hey, God's kingdom is so much bigger. And that's.
why Jesus continues. He said, "But you will receive power when the Holy Spirit has come.
upon you. And you will be my witnesses in Jerusalem, in all of Judea and Samaria, and.
to the end of the earth." You see, this is so much bigger, right? This is so much bigger.
than struggling for a small piece of land. It is sending us out as individuals who are.
transformed by the grace of God and dot the map of the whole world. From where? From Jerusalem,.
wherever they were at, their comfort zone, their starting point, to all of Judea and.
Samaria and to the end of the earth. And that's what the kingdom is all about. Now you may.
say this is a kingdom project beyond our ability, right? I mean, who can do this, right? I mean,.
if you do a company startup, right, you wouldn't think, "Okay, wow, wow, we have to be a global.
company," right? At least in the beginning period of the company. But that's what Jesus.
has in mind. Go to the end of the earth. Is it even possible, you may say, if I were one.
of the disciples, I would probably ask the question is, "Who am I? Who are we?" Right?.
We're just ordinary people and we're weak and poor, we are limited, we don't have a.
lot of resources, right? At times we do wonder, I mean, how can we do that, right? I mean,.
how many of you have been a Christian for some years, for some time? Yeah? Quite many.
of us, right? And sometimes we don't even have the courage to tell other people we are.
Christians, right? You have moments like this in your workplace, in your school, in your.
family. Sometimes, okay, before meal we have to pray, right? That's the right thing to.
do, right? But sometimes we just do a three second or two second pray, right? We don't.
want people to look at us. We don't want people to see that we are weird, right? We don't.
want to stand out, right? However, that's what Jesus had in mind. To give you power..
It's not because you have the resource to do it, but God is going to send the Holy Spirit.
to give you power and ability and capacity to do that. To bring the gospel to the end.
of the earth. I remember some years ago I was in another church, in my previous church,.
in the West Coast. I met this very successful businessman in my church, in the West Coast..
And he is like, I get to know him and he's the most articulate person I've ever known,.
right? He's a multi-billionaire company businessman. And he knows how to talk. He knows how to.
present, you know. You may have watched or listened to some of the talks on TED, right?.
On TED. The TED talk, the professional kind of, you know, presentation. And this businessman.

$^{161}$is one of those people. And he is so gifted and he's so resourceful. But when he came.
to Christ, he became a Christian. And of course he knew a lot of different platforms. He can.
share the gospel, right? In his workplace, with his co-workers and employees. He can.
bring a lot of people to church, right? But, you know, for a long this time, since he became.
a Christian, you know, he didn't even tell his co-workers that he was a Christian. Nobody.
knew in his workplace that he became a Christian and he was saved by grace. Nobody knew for.
a very, very long time. And then, you know, after he has grown so much spiritually, of.
course he has changed, but, you know, we ask him, "Why didn't you tell other people your.
question? Is it your identity, your spiritual identity?" You know why he said? He said,.
"I'm scared. I'm scared. I don't have the guts. I don't have the courage to do it.".
See, that is what is at stake here. Jesus is saying, "You cannot be my witness with.
your own strength. You don't have the guts. You don't have the boldness. You don't have.
the strength to be my witness. However, I'm going to send you my Holy Spirit. God himself.
will live in you." And Jesus knew that his disciples were weak. They were actually defeated.
at the cross, right? They were all scattered. They didn't have the guts to stand up for.
the kingdom. However, everything has changed since Jesus' resurrection. He's going to launch.
a new kingdom. And Jesus knew that we were weak. And God knows it. So God provides the.
resource. The Holy Spirit. God doesn't give us the money. God doesn't give us the ability.
to articulate, to communicate. No. All these things are secondary. But the most important.
thing is in your core, in your life, God's Spirit is giving you the power and courage,.
the ability to share the gospel, to be his witnesses. Let's center on that word. Witness..
Witness. Witness is someone who's being called upon to the court of law to provide a testimony,.
to tell the truth. And that is our identity, right? That is our identity. So Jesus said,.
"I'm going to provide you with that resource. If you tap into that resource, then you will.
begin to know that you have the power to be my witness." In the beginning of the chapter,.
verses four and five, Jesus said, I probably missed one of the, yeah, I think I did. Anyways,.
Jesus said what? "The power of the Spirit will come upon you." Now, what analogy I can.
use to describe the work of the Holy Spirit. Now, the Holy Spirit, the power of the Holy.
Spirit, it's like gasoline, okay? Gasoline for your car, okay? Many of us have your own.
cars, right? So if you drive a Tesla, some of you drive BMW, very nice car, right? Or.
you walk into our church, welcome. There are nice cars out there, okay? Please don't scratch.
it, don't touch it. They don't belong to me, okay? Nice car, right? But if you have a nice.
car, you don't have gasoline, you can go nowhere. You can wax it, you can clean it, right? But.
the car will not move until you fill it with power and gasoline, or electricity for your.
TV. And that's what the Holy Spirit is all about, is to give you the internal power to.
walk the journey of faith, to be Jesus' witnesses. So what did the disciples do? They did return.
to Jerusalem and obey the commands of Jesus. They stayed in Jerusalem until the Holy Spirit.
comes upon them. And they had the power to preach the gospel. So at the end of this sermon,.
I want to conclude with some of the thoughts, okay? Some of the thoughts are some of the.
homework for you, alright? Now these are your homework, okay? This week, this coming week,.

$^{201}$think about it. Name a person or a face that you know who has not known the name of Jesus..
Not the question yet, right? Maybe heard about it, maybe never heard about it. Think of a.
name, okay? Not a question. That God places that name, that person, in your heart. Think.
about this person. And then every single week, every single day of this week, set aside five.
minutes of your time. Each day, five minutes, to pray for that person, okay? Just pray for.
that person. That God will work in his heart or her heart, in his life or her life. That.
he or she will be brought to Jesus. So when you pray, pray for boldness as well. Because.
remember what Jesus said to the disciples? "You will be my witness." Jesus didn't say,.
"Okay, somebody else will be my witness." Okay, so you sit back, okay? You sit on the.
lazy boy, alright? You enjoy. No. You step out of your comfort zone and you pray the.
bold prayer. You know a lot of prayers that we utter to God are very safe prayers, right?.
We want to be safe, right? "God, I'm out of money. Give me some money." Right? "God, I'm.
sick. You know, heal me. Recover me." "God, I'm laid off. You know, give me a job." Those.
are very, very safe prayers, legitimate prayers. But those are really, really safe because.
it doesn't challenge you, right? God's mind is bigger than your mind. God's plan is bigger.
than your plan. And God's plan for your life is much bigger than what you can think of..
Much bigger than your safe prayers. Pray the bold prayer. When you pray this week, every.
single week, pray that God will give you courage to speak. To be a good witness for Christ..
Pray that God will give me strength, okay? To live for Jesus. Until you pray that prayer,.
you will be a nice car without the gasoline, okay? You need to add some gasoline, alright?.
You need to add the power of the Holy Spirit. Come on, guys. Pray that God will help you.
speak. Help you relate to the other person. Have you ever prayed like that? The risky.
prayer. But it's a prayer full of God's power. It's a prayer full of God's vision. That He's.
going to send you out to be the witness for Christ. So may this be our final prayer. God,.
dear God, give me the power to live for you. Give me the power to be your witness. Give.
me the power, the strength to share the gospel to the other person, your friend. Maybe your.
family member. Someone who God has put in your heart, okay? This is our final prayer..
Let's all pray like that, alright? I'll give you some time as we play the music background..
Who is that that God plays in your heart, okay? Maybe someone close to you, maybe someone.
in your workplace. Anyone. God has put a face and a name in your heart. Let's all pray..
Let's all pray together. At this moment, spend 30 seconds, one minute to pray for that person..
Dear Jesus, we come before you. We are weak and we are limited. We are small and we feel.
insignificant just like the first disciples felt. And yet we know that you have promised.
the Holy Spirit to be with us. You have sent your Holy Spirit, your power and your ability.
in our hearts to motivate us, to move us so that we step out of our comfort zone to share.
the gospel to other people. And Jesus, we pray that in your grace and your love, transform.
us continuously. Help us to move from fear to courage. Help us not to be timid, but have.
the boldness because we do have the boldness in Christ. So I pray that you will use us.
in this week as we pray about the person you put in our heart. And I pray that we will.
pray the whiskey prayer. We'll do something boldly for your kingdom. Thank you for your.

$^{241}$grace that we belong to your kingdom. All because of what Jesus has done for all of.
us. Thank you for your eternal love and your eternal grace. And I just pray that in this.
week you will be with us, you send us into our Jerusalem, Judea, Samaria, to the end.
of the earth. That we'll be your witnesses. And may you continue to grace us so that our.
church will continue to grow and our lives will continue to grow. Our heart will continue.
to be expanded, expanded with your love and your grace. To love like you do. To let us.
know like you do. To relate to other people like you do. So send us out so that we can.
continue to build your kingdom. To be part of your kingdom project. What a privilege..
What an opportunity. What a blessing that we can share this project together. So may.
you continue to bless our church, to bless each and every one of us. May we rise above.
our challenges and our issues and our problems. To rely on your power, the power of the Holy.
Spirit to be your witnesses to the end of the earth. So we pray that you will help us..
And we know and we believe that you do. In Christ's name we pray. Amen..
Well may God send you out with peace and courage and the power of the Holy Spirit. Amen. And.
you're welcome to stay with us for a time of fellowship and thank you for A-Peng for.
preparing good food for us. So you can stay here or stay outside in the cafe and we can.
get to know one another and build a church together. Amen. See you next week..
(soft piano music).
\newpage



\section{6:1-1}
\label{sec:BPEn5KuFa4c}
\textbf{Pursuing Kingdom Righteousness (Matthew 6\_1-1) - Chee Yie Jong}
\newline
\newline
連結: \href{https://youtube.com/watch?v=BPEn5KuFa4c}{\texttt{https://youtube.com/watch?v=BPEn5KuFa4c}} ~~~~ 語音日期: 2025-05-07
\newline
\newline
\hyperref[sec:Ml1uo7s2nGU]{\small{< < < PREV SERMON < < <}}
~
\hyperref[sec:index_chronic]{\small{[返順時目]}}
~
\hyperref[sec:index_scriptual]{\small{[返順卷目]}}
~
\hyperref[sec:Wr2X99fNMqk]{\small{> > > NEXT SERMON > > >}}
\newline
\newline
$^{1}$- Thank you, Kelly..
Good morning, UV..
It is good to be back here with you all,.
and I hope you guys get some time to enjoy the spring weather.
and just seeing different flowers blooming,.
the colors, there's like all around the streets..
This is my favorite time of the year..
Okay, this morning we're gonna continue with our series.
on Jesus' Sermon on the Mount..
And today we're gonna look at this,.
okay, pursuing kingdom righteousness..
But before that, let us pray together first..
Heavenly Father, as we prepare to open up your word.
this morning, I just ask that your spirit.
will be in each and every one of our heart..
As we receive your word being proclaimed,.
that your spirit will do your work in our hearts.
to transform us into people that represent your kingdom,.
into people of love, into people that the world.
will recognize that we belong to you..
Because we know that your word is living and active..
So allow your word to do the transforming work.
in our hearts..
Father, I also ask that your Holy Spirit.
will be with me this morning as I preach your word.
that anything that is not from you,.
that you will take them away..
And may your word and your word alone be proclaimed..
I pray this in your son's name, amen..
I don't know if there's the breathing voice.
coming through the mic, but hopefully that is not.
too distracting for you all..
Okay, if you are a basketball fan,.
this time of the year is usually some of the most.
exciting time because the NBA playoff.
is actually happening right now..
And most of the teams that are able to make it.
this far into the seasons are usually led by.
one or maybe two superstars..
And so these matchups are usually pretty fun.

$^{41}$to watch as well..
But what makes a great players like LeBron James.
or Micah Jordan or Steph Curry.
such a great basketball players?.
While there are many different positions.
on the basketball court, these are some of the.
common characteristics that makes someone.
a great NBA players..
First, of course, there is the basketball skills..
How well can they score?.
Or how well can they defend?.
Or how well can they make a play?.
And then there's also basketball IQ..
How well do they understand the game?.
And are they able to read the opponents.
and respond accordingly?.
And of course, there's the work ethic.
and also attitude, both on and off the courts..
And lastly, there's also some intangibles as well,.
like leaderships, versatility,.
or if they're able to make a clutch plays.
when a team needs them..
So all of these are some of the basic characteristics.
that makes someone a good basketball players..
But now if I ask, what makes someone a good Christian?.
How would you answer it?.
And how different do you think the answer would be.
if I ask non-believers these same questions?.
I'm currently in the running groups.
and we do a group run every Friday morning at 6.40..
And usually when we do a group run,.
so you would just end up talking to the person.
right next to you, right, for probably like.
about 30 minutes..
And usually when people find out that I'm a pastor,.
the first thing that they will usually mention to me.
is either, oh, I actually used to go to church growing up,.
but I just stopped going now..
Or, hey, I actually went to church with a friend of mine.
over Easter or over Christmas..

$^{81}$And if I ask them these questions,.
what do you guys think makes someone a good Christian?.
I'm fairly confident that a majority of them will say,.
okay, a good Christian is probably someone.
who goes to church regularly,.
because to them, going to church.
is some of the most basic activities.
that all Christian participate in..
And when it comes to Jesus days,.
there were also a few activities.
that became sort of like the identity marker.
of a good and pious Jews,.
because these are some of the most basic religious duties.
that a Jewish person have to sort of like.
be responsible for..
But in our passage today,.
we'll see how Jesus subverts the norms of his days.
on what makes someone a faithful or devout Jews..
But before that, let's enjoy our passage.
in the larger context of the Sermon on the Mount..
So in Jesus' Sermon on the Mount,.
there is three major section that sort of like.
made up the main content for this sermon..
And in the first major section,.
that is in Matthew chapter five,.
we learn that the truly blessed life.
comes from having an intimate relationship.
with our heavenly father, true Jesus Christ,.
and following after Jesus' example.
to live in a way that embodies his kingdom righteousness..
In this first section,.
Jesus focused on the true righteousness.
that exceeds that of the Pharisees.
and the teachers of the law..
But in our section in Matthew chapter six today,.
Jesus continues with his teaching on righteousness..
But this time, instead of focusing.
on the positive aspect of righteousness,.
Jesus is going to focus on the wrong kind of righteousness,.
false righteousness..

$^{121}$Moreover, while the better righteousness.
in Matthew chapter five has to do with our relationship.
with each other, that is sort of like.
the horizontal aspect of our relationship,.
right here in Matthew chapter six,.
the righteousness that Jesus talks about.
has to do with our relationship with him..
So the vertical relationship between us and between God..
And as we'll see later in our passage,.
showing off one's personal devotion to God.
or one's personal righteousness has become an avenue.
where someone in Jesus' days can do this.
to gain respect from the community..
Or they can do this to kind of like.
increase their social status..
And so Jesus, specifically,.
he warns again these kind of righteousness,.
the performative righteousness,.
the wrong kind of righteousness for his followers..
So let's look at our passage today in Matthew chapter six..
So Matthew six begins like this, be careful..
So Jesus doesn't waste any time.
and he gets straight to his point with this imperative.
or with this authoritative command,.
be careful or be aware..
And when Jesus say, be careful, we need to pay attention..
And so what is Jesus trying to warn us here about?.
He says, be careful not to practice your righteousness.
in front of others to be seen by them..
If you do, you will have no reward.
from your father in heaven..
So this verse right here, verse one,.
will serve as sort of like an overarching structure.
for our passage today..
So let's break it down a little bit..
So right here, there's four elements.
that Jesus will repeat on each of the examples.
that we'll be looking at today..
So the first one is the action..
So the action that Jesus talks about right here,.

$^{161}$the overarching actions is practicing righteousness..
And what righteousness means is pretty much.
is how to live rightly before the Lord..
And right here specifically,.
the emphasis is on our actions and on our deeds..
And then the second element is the manner.
of practicing this righteousness..
It is in front of other people..
And then the third element, the motive,.
is to be seen and approved by other people as well..
And then the final element, the consequences,.
Jesus say you will have no reward.
from your father in heaven..
So right here, what Jesus is having an issue is.
with more specifically on the second and the third element,.
the manner and the motives.
in which we practice our righteousness..
So with that in mind, this is the main idea.
that we'll be diving into today from our passage..
A worldly disciple pursues right living.
for the applause of people..
But a kingdom disciple pursues right living.
for the affirmation of the Father..
So we'll look at this idea from the three different examples.
that Jesus has given..
That is the first one, giving to the poor,.
second, praying, and then the third one, on fasting..
And why these three examples specifically?.
Well, again, these are some of the most basic.
Jewish practice and religious activities.
that Jews can express their devotion to God..
In other words, these three activities.
are the identity marker of a good Jewish person..
If you check all of these off,.
that basically means you are a pretty decent Jew..
And again, through these examples,.
Jesus is trying to show us that as a disciples of Jesus,.
as his kingdom disciples,.
we strive to live a holy and righteous life,.
not for the praise of other people,.

$^{201}$but to be seen by our heavenly Father alone..
So let's take a look at each of these three examples today..
And the first one is giving to the poor..
So Matthew chapter six, verse two..
So when you give to the needy,.
do not announce it with trumpets.
as the hypocrites do in the synagogues and on the streets..
To be honored by others..
Truly I tell you, they have received their reward in full..
But when you give to the needy,.
do not let your left hand know what your right hand is doing.
so that your giving may be in secret..
Then your father who sees what is done in secret.
will reward you..
So if you remember the four element here,.
the first one is the action, right?.
So the action that Jesus focused on here.
is giving to the poor..
And then the manner that Jesus warns against.
is announcing it with trumpets as these hypocrites do.
in synagogues or on the streets..
So a little background context on giving to the poor.
or almsgiving from the Bible..
In the Old Testament, social justice.
is actually a pretty prominent themes all throughout,.
but especially in the prophets..
And when it comes to the intertestamental period,.
that is the time between the Old Testament.
and the New Testament,.
giving to the poor has become a sign of religious piety..
And so fast forward to Jesus days,.
giving to the poor has pretty much become synonymous.
with righteousness..
And when people give,.
they don't usually give directly to the needy,.
but they give it through the synagogues..
And so the synagogues,.
they will function like a social agency,.
similar to the charities or nonprofits.
that we're familiar today,.

$^{241}$like Salvation Army or World Vision,.
where this agency,.
they will collect the donation from the community.
and then distribute it to the poor.
in the community as well..
So while this public display of generosity.
is a way to celebrate God's goodness.
and generosity towards his own people,.
the hypocrites, on the other hand,.
they will take this as an opportunity.
to show off their generosity.
so that they can receive praise from other people..
Simply put, their motivation here.
is to simply be seen and be praised by other people..
So what Jesus warns against here is this..
When you give to the poor,.
don't draw attention to your acts of kindness.
and generosity like these hypocrites do..
And the word hypocrite here will actually be mentioned.
in all three of our examples today as well..
So I think it is worthwhile to spend a little time.
just to understand what Jesus is referring to..
So this word originally referred to a stage players.
or an actor in a Greek theater.
who put on mask to play all these different roles.
in the play..
And usually, they will put on mask.
to pretend to be someone that they are not, right?.
So by the time of the New Testament,.
these words started to take on a more negative connotation,.
referring to someone who pretends to be morally good.
or morally virtuous,.
but they're just being false and empty inside..
Essentially, someone who is a moral actor..
I'm sure we all are familiar with someone like that, right?.
Growing up in Malaysia,.
usually when there's a flood or a house that got burned down.
or some sort of natural disaster that happened,.
the local politicians will usually show up.
to provide some sort of assistance.

$^{281}$and aid to these people, right?.
Because that is part of their job,.
that they're being elected by the people.
to serve the people..
However, most of the time, when the politicians show up,.
they don't just come by themselves..
They usually show up with the social media team.
or reporters from different news outlet.
or journalists from different newspaper.
in order that the good works can somewhat be documented.
and then being shown to the public..
Look, look at what I've done for you all..
And the funny thing is,.
the closer it gets to an election year,.
these kind of operations will become.
more and more frequent..
They will start to visit the poorer part of the city.
and then take a photo there,.
or they will visit hospital or orphanage in the city,.
or do whatever they need, the publicity stunt,.
to increase the reputation..
And the consequences of putting up a show for generosity.
like the hypocrites, this is what Jesus warns..
"Truly, I tell you, they have received the reward in full.".
This can be a little bit confusing at first.
because you might wonder,.
"Jesus, I thought you warned against.
"this kind of public display of generosity..
"How come they're still able to receive a reward.
"and even receive a reward in full?".
So let's unpack that a little bit..
So what Jesus is saying here.
is that they have received their full reward.
now in the earthly sense..
Put it another way, Jesus is saying that.
if you give to the poor.
with the intention of receiving praise from people,.
then praise from people is all you're ever going to get..
And God pretty much owes you nothing..
So in this way, your reward is in the here and now..

$^{321}$Your reward is pretty much temporary..
And from another perspective,.
you're actually not giving to the poor,.
but you're sort of like buying your praises from people..
So your reward is also transactional as well..
Every time you want to receive honor or receive praise,.
you just give to the poor.
so that you pretty much buy your praises,.
buy the action of giving to the poor..
So how should we give as his disciples?.
As kingdom disciples,.
how should we do this act of kindness towards the needy?.
And what kind of manner should we have?.
So in verse three and four, Jesus say,.
"When you give to the poor, to the needy,.
"do not let your left hand know.
"what your right hand is doing,.
"so that your giving may be in secret..
"Then your father who sees what is done in secret.
"will reward you.".
So what Jesus say is this,.
when we give to the poor or when we help those in need,.
do it in secret or privacy,.
to the extent that even our left hand.
doesn't know what our right hand is doing..
And to be able to show kindness and generosity.
to this extent requires us.
to completely disregard ourself, right?.
That is what Jesus mean,.
that your left hand doesn't know.
what your right hand is doing..
It is a point where we need to completely disregard ourself..
And the outcome of it,.
we will receive our heavenly reward from the Father..
And one thing to note here,.
Jesus does not mean that we give.
because we want to receive a reward.
from our heavenly Father..
We give from a cheerful and sincere heart.
for several reasons..

$^{361}$If you recall at the end of Matthew five,.
the command to be perfect as our heavenly Father is,.
our goal on this side of heaven.
is to be like our heavenly Father..
And from the Bible,.
we know that our heavenly Father.
is one who cares for the poor,.
the marginalized,.
and those in which the world has little regard for..
So when we give to the poor or help those in need,.
we are living out our kingdom identity.
as the son and daughters of God..
And secondly, we give to the poor because we love him.
and we want to give back a portion.
of what he has blessed us with..
Later in Matthew,.
in the parable of the sheep and goat,.
Jesus himself emphasized that.
whatever we do to the least among us,.
we do it to Jesus himself..
So as disciples of Jesus,.
we pursue generosity and kindness,.
not for the temporary glory from men,.
like the hypocrites or the worldly disciples do,.
but we do it for the affirmation from our heavenly Father..
So this is the first examples on giving to the poor..
Then Jesus moved on to the second examples, prayer..
So let's start with some background context.
to prayer in the Jewish culture..
If you are a Jews,.
there are set times of prayer for you every day..
As traditionally,.
the Jews, they will pray at least twice a day,.
one in the morning and one in the evening..
And usually during those set hours of prayers,.
they will pray the Shema prayer,.
which is from Deuteronomy..
Hear, O Israel, the Lord, your God, the Lord is one..
And so when it comes to prayer,.
more specifically on personal prayer,.

$^{401}$this is what Jesus has to say..
And when you pray, do not be like the hypocrites,.
for they love to pray standing in the synagogues.
and on the street corners to be seen by others..
Truly, I tell you, they have received their reward in full..
But when you pray, go into your room,.
close the door and pray to the Father who is unseen..
Then your Father who sees what is done in secret.
will reward you..
So again, we see the similar elements.
and themes right here, right?.
It start with the action, that is praying..
And the manner in which the hypocrites pray.
is to stand in the synagogues,.
which is the epicenter of all things Jewish life,.
or on the street corner,.
which is also a public space with high visibility..
If you think of a corner,.
you can see through a corner, not just from one side,.
but from both ways as well..
And usually the foot traffic around the corner.
will also be a little bit slow.
because people will slow down.
when they're trying to turn a corner..
So the street corners are places.
with usually the highest visibilities that you can get..
And the motivation of doing these things,.
once again, it is to be seen by other people..
And to these hypocrites, Jesus once again say that.
they have received their reward in full.
in the here and now,.
the temporary and fleeting kind of reward.
that pretty much just goes away.
once your prayer is finished..
But on the contrary, Jesus taught that when we pray,.
go into the room, and not only that,.
close the door and lock it up.
so that no one else except God alone can see us..
This is because the goal of our personal prayer.
is to commune with our Father and to abide in Him,.

$^{441}$not to impress other people..
And then Jesus continues with the same practice of prayer,.
but this time with a different kind of manner.
and also a different kind of motivation..
So verse seven,.
"And when you pray, do not keep on babbling like the pagans,.
for they think that they will be heard.
because of their many words..
Do not be like them,.
for your Father knows what you need before you ask Him.".
So besides praying in public places,.
another manner that we can be tempted to pray.
is to use fancy words or try to pray as long as we can.
without any actual content to it..
Jesus compares this kind of prayer.
to the prayers of the Gentiles during their pagan worship,.
where they were pretty much like playing relate.
with the pantheon of gods,.
where they would repeatedly mention.
the names of different pagan gods.
with the hope that just maybe the right god.
would listen or hear their prayers.
and help them in their circumstances..
In other words, they pray.
because they try to manipulate gods.
for their personal benefits..
But as God's children, we do not have to pray like that.
because our Heavenly Father is a different kind of God..
We don't need a secret passcode or say a certain phrase.
or special words for Him to hear our prayers,.
because He is a God who loves and cares for His people,.
one who knows our needs even before we bring it before Him..
He is a God who is eager to respond to our requests,.
and He is a God who delight in His children.
when we come before Him with our prayers..
And as we step into this new theme.
of crossing over at our church here,.
one of the aspects that we'll be focusing on,.
like Pastor Mark said earlier,.
is on our relationship with Him, with God..

$^{481}$And so with that, I want to ask,.
how is your relationship with the Lord.
when it comes to your prayer life?.
On my trip back to Malaysia last October,.
I got to spend a lot of time just catching up.
and also hanging out with a lot of friends and family.
because I haven't been back for seven years..
However, one thing that I quickly realized.
is that in most of these gatherings or hangouts,.
people would always ask me to pray.
because they know that I'm a pastor here in the US..
And it got to the point where during dinner one night,.
it is just between me, my brother, and my mom..
I had to stop and tell my mom,.
"Mom, can you stop asking me to pray every time?".
Like people has been asking me to pray.
since I got back to my hometown pretty much..
And this experience makes me realize two things..
For one, for me personally,.
I feel the pressure to pray in a certain way.
to meet the expectation of people,.
especially with people acknowledging that,.
"Oh, he's a pastor.".
So I feel like I have to pray in a certain way..
And so when I pray during those times,.
I feel like I'm praying more.
to meet the expectation of these people.
instead of genuinely praying to the Father..
But secondly, I sense that for some people,.
they have these false ideas that,.
"Okay, if a pastor prays for them,.
they will be extra blessed,.
or that God will be more willing to answer the prayer.".
For these people,.
they might feel like they don't really know how to pray,.
or their prayer is less effective.
just because of the fact that.
they're not "professional Christian.".
But again, what we see right here.
is that this is not the case at all..

$^{521}$God does not concern with what words we use,.
or how long our prayer is..
In the following verse in Matthew 6,.
Jesus would go on to show us how we ought to pray..
And this has become known as the Lord's Prayer..
But for today, we'll skip over this portion first.
and come back to it later in my next sermon..
And so, so far we have looked at the first two examples.
on giving to the poor and also on prayer..
And now Jesus continues with his third practice.
of righteousness on fasting, beginning in verse 16..
So when you fast, do not look somber as the hypocrites do,.
for they disfigure their face.
to show others that they are fasting..
Truly, I tell you, they have received their reward in full..
So the purpose of fasting is for us.
to kind of like have a heart of complete dependence.
on the Lord, right?.
And it's also usually also a way for us.
to kind of like practice humility,.
just knowing that we cannot depend on ourself..
However, the hypocrites,.
they would intentionally disfigure their face.
and make themselves look gloomy and miserable.
when they fast..
This is because in the Old Testament,.
fasting is often associated with repentance..
And along with fasting,.
people will usually also put on ashes on their face.
and put on sackcloth.
in order to heighten the feeling of sin and mourning..
But when these hypocrites fast,.
they put on a show.
so that they can be recognized as more religious.
instead of coming out from a heart.
of mourning over their sins or mourning over the brokenness..
So instead of coming from a place of humility,.
this kind of performative fasting.
comes from a place of pride,.
of wanting to be recognized..

$^{561}$And instead of an attitude of dependence on God,.
this kind of performative fasting.
comes from an attitude of dependence on yourself..
Pretty much, look at how righteous I am..
Look at how religious I am..
And this kind of performative fasting.
is what Jesus is warning us against..
And again, these hypocrites,.
Jesus say that they have received the full reward.
in the here and now..
So how should we fast then?.
So this is what Jesus say, verse 17..
"But when you fast, put oil on your head and wash your face.
"so that it will not be obvious to others.
"that you are fasting,.
"but only to your father who is unseen..
"And your father who sees what is done in secret.
"will reward you.".
So what Jesus meant is this..
When you fast, clean yourself up.
and make yourself look presentable..
Go out of your way to make yourself look normal.
in order that no one will recognize that you are fasting,.
but only the Father alone..
So I want to end my sermon with these stories of someone.
whom I have the privilege of meeting,.
who perfectly embodies this kind of heart posture.
in their love of God.
and also in their service for other people..
In summer 2019, when I was still serving.
with another ministry,.
I spent the summer staffing a summer mission.
of about 20 plus college students.
from all across the United States..
And we're having our summer mission here.
in the inner city of Chicago..
And on our first Saturday,.
we partner with a pastor from South Side of Chicago.
to do a homeless outreach..
So this pastor, his name is Pastor Burton..

$^{601}$And he had been doing this homeless outreach.
for years already at that point..
And every Saturday morning.
during the warmer months of the year,.
he would set up a table to give out lunches.
and essential supplies to the homeless person..
So I have a few pictures from that summer..
Can anyone recognize where the pictures on the top right is?.
Where is that at in Chicago?.
Yes, it is on Canal Street.
and it is right across Pacific Garden Mission,.
which is a homeless shelter..
But why does he specifically set up a table.
to provide these essential supplies and food.
and lunches for the homeless people.
outside of a homeless shelter?.
Well, because there are people.
who are not being able to accept.
into the Pacific Gardens Mission programming.
for various reasons..
And these are the people that Pastor Burton,.
he wanted to serve..
And after a couple of hours,.
just kind of like setting up the table at Canal,.
he would then drive to Lower Wacker Drive.
to give out some supplies.
and also check in on the people who live in tents.
and some of the max-shift shelters there..
So the picture on the lower left.
is actually at Lower Wacker..
And if you have been to Lower Wacker before,.
you probably know that it is pretty much.
like a maze down there, right?.
It's not only is it dark.
and the street is a little complicated.
and your GPS like nine out of 10 times,.
it just doesn't work there..
So I've personally driven through Lower Wacker.
countless times,.
but I actually never recognized anyone.

$^{641}$actually living down there..
But the thing about Pastor Burton.
that really stood out to me,.
even until this day,.
is how he's able to remember.
where each of these individual lives..
I remember that Saturday,.
we would just randomly stop at a spot.
like on Lower Wacker,.
and then he would just turn and tell us,.
"Okay, you guys can grab the supplies.
"and we're gonna head out now.".
And then he would lead us.
to all these different nooks and crannies.
on Lower Wacker,.
where all these people live..
And so not only did he remember.
where each of these people live,.
he also know them by their name..
He will call them,.
and these people will kind of like.
come out of the corners..
And he also knew the specific needs.
of each of these individuals..
If it is a guy,.
he have some specific supplies for a guy..
If it is a mom with a few little kids,.
he also have some specific supplies.
for this mom and her kids..
And the most important thing is,.
you would never see Pastor Burton.
being on the news,.
being recognized for the good work.
that he has done..
And yet he show up week after week,.
because he knew that he was serving.
these people for the Lord..
For him, he is just living out his faith.
and being obedient to Jesus..
And even though I only get to spend a day.

$^{681}$with Pastor Burton,.
his example has left such a huge impact.
in my life that I think.
it will pretty much last forever..
And that summary also changed.
a lot of my perspective of Chicago.
as a city as well..
So before I end,.
I want to leave us with this question.
to reflect on..
What motivations tend to drive.
your acts of worship and service?.
Is it the approval of God?.
Or the admiration of others?.
Again, a worldly disciple pursues.
right living for the applause of people..
And a kingdom disciples,.
on the other hand,.
pursues right living for the affirmation.
of the Father..
So let's spend some time.
just reflect on these questions.
and I would end this with a prayer..
(soft music).
- Heavenly Father,.
I ask that you would help us.
to be a church and be a people.
who loves you,.
who serve others,.
not for the praise of people.
or to be recognized by the world,.
but to do all these things of generosity,.
this act of kindness,.
out of a pure motives of just loving you.
and also loving your people..
Help us to be a people.
who is able to love selflessly,.
who is able to serve selflessly,.
because we know that we belong to your kingdom..
We know that we belong to you,.

$^{721}$our Heavenly Fathers,.
and that we're able to do all these things,.
not because of our own abilities,.
but because of the love that you have for us..
So help us, Fathers,.
not to chase after the glory of man.
or the praise from the world,.
but to be your kingdom disciples.
who long for your affirmation.
and your approval alone..
Help us to be people.
who loves you.
and love others in a way.
that shows our kingdom identity..
(soft piano music).
So Father, as we finish our service today.
and just go on for our week,.
I just ask that you will help us,.
even in this week,.
to have opportunity to serve and love others.
in the way that you have called us to..
So I pray this in your Son's name, amen..
Okay, thank you for joining us this morning..
As Kelly said, we'll serve some lunch outside,.
so feel free to hang out.
and just fellowship with one another..
Thank you..
\newpage



\section{創世記 22:1}
\label{sec:Wr2X99fNMqk}
\textbf{「空白」的三天 (創世記 22\_1–19) - 楊靖詩博士}
\newline
\newline
連結: \href{https://youtube.com/watch?v=Wr2X99fNMqk}{\texttt{https://youtube.com/watch?v=Wr2X99fNMqk}} ~~~~ 語音日期: 2025-05-07
\newline
\newline
\hyperref[sec:BPEn5KuFa4c]{\small{< < < PREV SERMON < < <}}
~
\hyperref[sec:index_chronic]{\small{[返順時目]}}
~
\hyperref[sec:index_scriptual]{\small{[返順卷目]}}
~
\hyperref[sec:rAFDbYbX6HA]{\small{> > > NEXT SERMON > > >}}
\newline
\newline
創世記 22:1
\newline
\begin{longtable}{cl}
\hline
\hline
章節 & 經文 (和合本修訂版)\\
\hline
22:1 & \begin{tabularx}{0.7\textwidth}{X} 這些事以後，神考驗亞伯拉罕，對他說：「亞伯拉罕！」他說：「我在這裡。」 \end{tabularx} \\ \\
[1ex]
\hline
\hline
\end{longtable}
$^{1}$大家好.
平時上網看到這個台.
第一次來見得多.
原來外面有這麼多美麗的東西.
所以沒有機會來過城會的網友.
可以特意來visit一下.
我覺得也是一個芝加哥的奇景.
我是Helen.
多謝大家歡迎我來這裡.
跟大家分享上帝的訊息.
我不知道原來城會教會沒有幫我讀聖經.
不過不要緊.
幸好我有張畫作可以代替.
如果大家像我這樣沒有什麼藝術造詣.
也不要緊.
以畫作為例子.
嘗試明白我們應該看的那段經文.
已經畫好了.
你能理解多少就大家的誤性.
我沒有什麼誤性.
不過我純粹看到.
這是中世紀時代的一位名畫家.
叫做卡拉阿喬.
我沒有讀錯吧.
我應該讀錯.
算了不要緊.
所以我沒有什麼藝術細胞.
不過我嘗試找今天想跟大家分享的經文.
是《創世記二十二章》.
是亞伯拉罕獻他的兒子而殺.
你看到這幅畫中.
中世紀的畫家將他描繪成.
亞伯拉罕按著兒子.
兒子在下面.
然後你看到他的樣子是怎樣.
是很烈還是很堅決.
他很堅決地去做這件事.
他的畫家描繪中.
亞伯拉罕聽到上帝的聲音.
說你要去獻你的獨生子.

$^{41}$你所愛的我一個兒子.
去獻給我作為獻祭的時候.
亞伯拉罕在《聖經》中二十二章.
說他一大清早起來.
然後帶上服務員就走了.
三天之後.
服務員就說你們在這裡等我.
我和兒子會回來.
大部分人都覺得亞伯拉罕真的很有信心.
他怎知道他和兒子會回來.
明明是上帝說你要獻祭.
獻了祭在山上獻.
怎會和兒子回來.
難道輪迴?.
難道破地獄?我不知道.
總之我們經常有印象.
就是亞伯拉罕很厲害.
他的信心就是.
他以至上帝說.
我要你去獻你的兒子上去的時候.
一大清早把所有東西都帶走.
走到某一步的時候.
丟下兩個服務員.
說我們會回來.
很有型的.
我們會回來.
然後和兒子一起帶著東西走.
走著走著.
兒子可能覺得不對勁.
阿伯.
獻祭的祭物在哪裡?.
其實明明就是揹著鞋子的兒子.
就不是你的兒子.
就是祭物.
但是亞伯拉罕說.
上帝會預備.
所以我們的印象.
覺得亞伯拉罕很厲害.
信心非常厲害.
甚至在新約當中.

$^{81}$我們會說.
亞伯拉罕是我們的信心之父.
整個基督教教義.
一個很核心的東西.
就是恩信稱義.
其實就是典故源自於亞伯拉罕.
當他75歲沒有兒子的時候.
上帝說.
不是啊.
有一個兒子我會賜給你.
而這個兒子.
我會使他成為大國.
會使他成為各國的祝福.
當他什麼都沒有的時候.
亞伯拉罕什麼都沒有.
生都沒有生過兒子的時候.
聖經說.
他信.
所以耶和華以他的信為義.
這個就是我們整個教義.
基督教教義當中.
不是憑著我們做了什麼.
不是因為我們做了多少善行.
而是因為我們信上帝.
這個就是在上帝面前能夠稱義.
稱義的意思就是.
我們有權利站在上帝面前.
我們可以在上帝面前站立.
這個就是稱義的意思.
整個新約的開始.
不要說舊約舊約.
亞伯拉罕當然是一個很典型.
很厲害的人物.
摩西亞,亞伯拉罕,戴維這些.
其實如果你是打開新約還是聖經.
你不熟悉聖經不要緊.
打開聖經第一頁.
第一句已經有亞伯拉罕這個名字.
因為是耶穌基督的家譜.
戴維的子孫.

$^{121}$亞伯拉罕的子孫.
就是耶穌.
整個一看下去.
無論是福音書也好.
無論是寶來書信也好.
或者其他書信也好.
亞伯拉罕這個名字是不斷地重複.
不斷地出來.
你看到這幅畫裡面.
我再說多一點.
你看到亞伯拉罕很堅決地獻了這個兒子.
但是有人壓著他.
就是我們待會會讀的聖經裡.
這個《創世記》22章.
他說當他真的想著手起刀落的時候.
神的使者就壓著他.
你看到我看了一些評論.
應該叫做藝術的研究.
我也不知道叫什麼.
很沒有藝術細胞.
他解釋給一些不懂讀畫的人.
他看到天使很大力地壓著亞伯拉罕的手.
他想表達亞伯拉罕真的想著手起刀落的時候.
真的刀下留人.
我的黑來.
這是我們對亞伯拉罕的印象.
很厲害.
想都不用想.
掠都不用掠.
這樣就走去獻世.
這個是希伯來書所描述的亞伯拉罕.
希伯來書這11章.
我們會叫他做信心英雄榜.
很多信心的偉人.
因著信就怎樣怎樣.
亞伯拉罕在這裡佔了一個很大的篇幅.
他基本上就講了整個.
由第八節開始講了整個亞伯拉罕的生命.
我們今天沒有機會去看這麼多篇幅.
因著信.

$^{161}$上帝叫他離開本地本族父家的時候.
他就遵命出去.
一直到母子生的時候.
當上帝說英雄.
就算你母子生.
但其實我會給你一個兒子.
就好像天上的星.
海邊的沙這麼多的時候.
他相信.
然後因著信17節開始.
就在講我們創世記22章的經文.
就是說他被考驗的時候.
把耳朵獻上.
這個就是他的信.
他也甚至.
他的評論就是說.
他甚至認為.
神可以叫他的兒子耳朵.
從死人中復活.
所以我們覺得.
原來亞伯拉罕不怕獻兒子.
因為他覺得上帝可以叫死人復活.
但要記得當時其實是沒有死人復活的回事.
在創世記沒有記載有這件事.
所以其實是亞伯拉罕想像.
上帝可以怎樣.
不反了一個自己的英雄.
因為他英雄就是我會給你兒子.
而這個兒子就成為這個祝福.
但過了這麼多年之後.
又叫他獻兒子.
是死的.
不是說切了一塊皮.
切了一隻手指下來.
是真的整個獻上的時候.
我們的印象就是亞伯拉罕很厲害.
希伯來書這裡其實.
講了整個信心英雄榜12章.
到了12章的時候.
他的結論是什麼呢.

$^{201}$他說既然我們有這麼多人.
如同雲彩包圍圍繞著我們.
就我們應該脫下全類我們的罪.
全心仰望怎樣怎樣怎樣.
其實他在說.
有這麼多人物.
因著信上大.
他們信心這麼強大的時候.
我們可以不用理世上的前類.
我們不需要這麼害怕.
給點信心我們走吧.
你行不行.
因為旁邊那個有信心.
你可不可以有信心多一點.
我很平庸.
雖然我是PhD畢業.
我想告訴大家.
我confess.
在我兩個女兒面前.
我confess.
我中一有四科不合格.
從小到大.
我媽媽也在.
你可以找我媽媽做見證.
我媽媽每次考完試回家.
幾分啊.
70分.
這麼差.
不是啊陳大文只有65分.
陳崇基呢.
陳崇基那個不要說了.
每次90分.
為什麼不和陳崇基比.
怎麼比啊.
陳崇基是那些在中學前.
正在打機.
第二天走進課室.
哎呀90分.
我是大學很討厭的同學.
我們今天三點多還在圖書館.

$^{241}$昨晚11點睡了.
打完機睡了.
怎麼比呢.
我這些信心這麼平庸的人.
我怎麼和亞伯拉罕比.
沒得比.
但如果這樣想.
亞伯拉罕信心很強.
所以我可以follow他的footsteps.
坦白說我不行.
我就是70分的那個.
我怎麼和90分的那個比.
100分的那個比.
我最多和65分的那個在排練.
不知道怎麼辦.
這是不是聖經想告訴我們.
亞伯拉罕就是拿100分的那個.
你就和他比一下.
所以你就會很強.
這是我媽媽的一個myths.
我和祂比多少次.
其實真的不關我的事.
如果我們這樣看聖經.
我想我們miss out了很多東西.
我們剛剛完了復活節.
復活節之前是壽難節.
耶穌在被捕之前.
祂騎了什麼圖.
大家記得嗎.
祂說上十字架了.
Let's go.
就是這樣嗎.
不是.
祂去到上帝面前說.
我很憂傷.
憂傷到幾乎要死.
如果你可以的話.
叫這杯離開我.
耶穌在基督想上這個十字架.
聖經的記載說祂不想.

$^{281}$祂三次祈禱和祂一樣.
如果你可以的話.
叫這杯離開我.
但是如果這是你旨意的.
我會去做.
在盧家福音裡面.
祂的記載是滴汗.
不是Let's go這麼興奮.
不是像亞伯拉罕.
我們印象中的那幅畫.
走啊.
上帝叫我去我就去.
不是.
就算耶穌基督是上帝的兒子.
他掙扎.
掙扎到滴汗.
他說好像血這麼大一滴.
如果耶穌基督.
也要這樣掙扎的時候.
我想問問.
為什麼我們覺得.
亞伯拉罕會比耶穌基督厲害.
這個信心.
Does it even make sense.
有時候我們看聖經.
就這樣看下去.
看下去就這樣說.
他沒有說.
他二十二章那裡說.
上帝叫完他之後.
他沒有記載他掙扎.
我們就當他沒有掙扎.
是不是這樣看聖經呢.
我覺得這樣看聖經.
真的不太能幫助.
這也是我今天.
為什麼會選這段經文.
我嘗試去學習.
這段創世記經文的時候.
Nextham Research Commentaries.

$^{321}$是一個很好的評論.
我相信Mark Sir可以.
介紹更多給你認識.
但我看到他.
他講完一堆評論的時候.
那些總結.
他說到這裡.
亞伯拉罕在二十二章的行為.
你看看黃色那裡.
他說他沒有質疑或尋求其他選擇.
Seek an alternative.
他說他沒有質疑.
沒有尋求其他選擇.
他沒有和上帝爭辯.
然後我心想.
真的嗎.
如果你看回亞伯拉罕的生命.
你是不是能夠去到這樣的結論.
今天我不是有很多時間.
但我嘗試很快地.
跟你去亞伯拉罕的生命.
以至我真的可以.
讓大家一起進入.
究竟亞伯拉罕去到創世記二十二章.
當上帝說.
你要將你的獨生子.
你所愛的那個.
不停地追求.
你的兒子.
獨生的那個.
你所愛的那個.
獻給我的時候.
亞伯拉罕是不是還在.
像剛才那樣說.
他沒有質疑.
他沒有Seek an alternative.
在創世記十二章那裡.
他呼召亞伯拉罕.
他說.
他叫什麼.

$^{361}$你要離開本地本族父家往.
我所指示你的地方去.
當時是第一次跟他說.
你會成為大國.
什麼什麼.
到了創世記十五章的時候.
神在與他立約.
就說你.
你的兒子.
我會賜很多東西給你.
亞伯拉罕說.
我兒子也沒有.
要不我家裡的.
Servant.
我的姐姐.
姐姐有兒子.
我夫妻無兒無女.
給這麼多東西.
我都沒有用.
我姐姐有個兒子就行了.
他說不是啊.
上帝說.
是從你所生的才會.
才會是你的合法繼承人.
不是姐姐的兒子.
不是花王的兒子.
是你自己所生的兒子.
他老婆很聰明.
我會簡稱為屎Q王.
上帝說我會有兒子.
我又生不了.
那時候他老婆已經沒有月經了.
即是今年期過了.
生不了.
他就做了個妾侍給他老公.
我也不明白為什麼可以這樣.
他做了個妾侍給他老公.
你跟他生吧.
亞伯拉罕當時應該還有心有力.
所以他真的生得出.

$^{401}$所以他生了20瑪利出來.
我看TVB大的.
大家如果有看TVB大的.
我現在沒有看了.
TVB大的你知道.
通常都沒有什麼好節目.
說的這些.
二奶啊 私生子啊.
16張就說他們在爭.
在爭的時候.
那時候86歲.
亞伯拉罕有個20瑪利.
就是這個私生子.
到17張的時候.
他老婆說再次.
他的呼召.
他的應許臨到亞伯拉罕.
就再說多很大堆的東西.
他就說你就會成為多父之國.
你老婆呢.
在你老婆所生的這個17張.
他說.
使你在撒萊多多一個兒子.
我要賜福給他.
就是你老婆.
不是二奶.
不是妾侍.
然後亞伯拉罕.
怎樣.
笑了.
他笑什麼.
100歲還能生兒子嗎.
他是不是在說他老婆.
他在說自己.
100歲我無心無力.
我知道我無心無力了.
因為他之後就說.
我老婆已經90歲.
所以100歲那個不是他老婆.
那時候他99歲.

$^{441}$他說100歲我真的生不了.
我希望我私生子在你面前存活.
這個是亞伯拉罕說的.
他不相信的.
他不明白.
亞伯拉罕完全不明白.
他不相信.
所以說我私生子.
之前他說是姐姐的兒子.
現在就說我私生子.
希望在你面前存活.
以至我的所有的身份給他.
然後他說什麼.
他說不是.
從你撒萊.
你老婆所生的兒子.
才會成為你後裔.
這個是他.
已經神確立了.
這個是他的獨生子.
這個是他的合法繼承人.
就不是不知道是誰.
到了100歲的時候.
對不起.
到了100歲.
他真的.
他老婆生了一個兒子.
要在這創世的21章.
如果我們看《色經》.
就在說21章和22章很相似.
因為到最後.
其實他們兩個兒子.
同樣經歷一個很大的威脅.
因為他老婆真的受不了姨母.
也受不了姨母的兒子.
就要趕他走.
上帝說你聽你老婆說.
就真的趕了20萬年.
就是他的私生子走.
所以他真的走了.

$^{481}$所以在那一刻.
也說你趕他走不要緊.
我也會祝福你的私生子.
他也有一代.
不過你記住亞伯拉罕.
你的合法繼承人.
是你和撒萊.
你老婆所生的兒子.
才是你的合法繼承人.
到了21章.
就是我今天會讀的那一段.
我講了一段經文.
之前那一章.
就已經是說.
只有一個.
亞伯拉罕.
如果你一直看他的心靈.
其實你睡不著.
其實挺好看的.
坦白說.
這段聖經應該不是很悶.
這12章.
多到22章.
講到18章的時候.
亞伯拉罕其實是一個談判專家.
他和上帝談判.
如何不毀滅兩個城.
他說有50個好人.
你可不可以不毀滅.
上帝說OK.
45個可以嗎.
OK.
40個可以嗎 上帝.
20個可以嗎.
10個可以嗎.
其實我看了.
我覺得有點煩.
我真的覺得.
其實不用講那麼久了.
他說上帝你先不要生氣.

$^{521}$10個OK嗎.
亞伯拉罕是一個會談判的人.
你看聖經就知道了.
他是會談判的.
他會談判的.
他有兩次在這幾章經文.
是會認妻作婢.
為了自己的兒子.
他不是我們心裡想過這麼高尚的人.
我有一次和朋友分享.
我來芝加哥的港島.
說什麼.
說這些話.
其實亞伯拉罕都很卑鄙.
他說的不是我說的.
其實亞伯拉罕都很卑鄙.
推了老婆出去.
你不要搞我的兒子.
你搞他就可以了.
他是一個這樣的人.
為什麼我們的印象這麼崇高.
其實他和我一樣.
65分拿過來的.
當你這樣看的時候.
你再看這段經文.
「這時」是「時候」.
上帝要考驗亞伯拉罕.
他說亞伯拉罕.
我在這裡啊 上帝.
上帝說你帶著你的兒子.
你我所愛的獨生子.
不要搞錯.
是二十 不是二十瑪利.
不是你的私生子.
是你的合法繼承人.
拿著他.
往摩尼亞去.
在我指示你的座的山上.
牽他為凡制.
很多人會覺得.

$^{561}$亞伯拉罕真是乖.
清早起床就聽上帝的話.
你有沒有想過他睡不著.
他可以睡不著嗎.
可能吧 可能不可能.
我不知道.
但如果我純粹覺得他真是聽話.
清早起床就搞定.
他可以睡不著.
他真的想不定.
為什麼上帝和我說的.
完全不合理呢.
是完全不合理.
因為是上帝從他沒想過有兒子的時候.
說賜福給他.
到賜福給他的時候.
說會有兒子在你的後裔.
他搞了一個二十瑪利出來.
上帝說不是二十瑪利.
是你和你老婆所生的合法繼承人.
才會繼承你遺產.
但轉眼間你又叫我.
It's ok 你趕走二十瑪利.
因為你還有二十.
記得嗎 二十一章的經文.
到二十二章說.
Sorry 你喜歡的一個.
你最愛的一個.
你唯一一個.
請你獻給我.
一個認妻作妹.
一個會和上帝negotiate.
一個其實他也不相信上帝說的話.
不可能不相信.
而是說真的嗎.
我一百歲生不成.
沒心沒力.
我老婆絕痲經.
怎麼生啊.
他真的會像我和你一樣.

$^{601}$去想這些事.
他不是不想的.
他是六十五分的亞伯拉罕.
不是一百分的亞伯拉罕.
當你有這樣的想法.
去想想其實亞伯拉罕真的有掙扎的時候.
他的清早起身.
去到那個地方.
上帝就有指示他的地方.
那段路程其實一點也不容易走.
很多刻可能他可以turn away.
好像Jonah那樣.
走回頭.
上帝叫我去我不去.
到了第三天.
他說你們留在這裡吧.
我和我兒子會回來.
可以是甚麼.
難道告訴你我會斬兒子嗎.
make sense.
為甚麼亞伯拉罕的故事對我來說這麼責心呢.
剛才有人跟我說他去visit Markham.
Ontario.
如果你有機會跟Mark Sir去吃東西.
爆了Mark Sir的東西.
去吃東西或者去喝咖啡.
剛才有位車友弟兄說去喝咖啡.
這裡就是正正喝咖啡.
幫我hatch賣廣告.
外面是空地.
有一天我在行Highway 7 Markham.
在一條很繁忙的路上.
很煩躁 越來越多車.
很煩躁 人們開車很壞.
一條三分鐘的路程.
可以變成十三分鐘出來.
我就很煩躁.
在車裡一家四口說.
何時才可以離開這個地方.
因為我讀完PhD的時候.

$^{641}$我知道我會走.
我知道我想離開多倫多.
但我不知道要去哪.
上帝說何時要走的時候.
我一看車的右邊.
我就看到這個sign.
我整車人一起看到.
突然間這個sign在我車門旁邊放出來.
You go.
我們四個人笑了出來.
是You go.
何時可以走啊 上帝.
我真的問完之後.
不知道為什麼看了右邊.
我就看到這個sign.
每天都走過這條路.
但從來沒有見過這個sign.
是我那一刻這個sign出現了.
就好像阿伯那樣.
You go.
我和我老公說.
我這65分的學生.
讀完希伯來文.
我就給了老師.
我高分的 其實90多分.
不過讀完我真的不記得.
我問他.
《色經學》寫說.
創世記22章和12章.
上帝叫阿伯安罕離開的字.
是怎樣解釋的.
我問我家的希伯來專家.
我不知道他有沒有試過.
真的與我無關.
他說其實是這樣解釋的.
去 你的.
我的希伯來神人真的給了我先駕.
但他真的叫我去 你的.
無論前後怎樣讀.
但意思就是去.

$^{681}$你要去的地方.
上帝真的跟我說.
去要我去的地方.
在過去一年多裡.
我真的很認真地問上帝.
上帝你要我去哪裡.
我今天沒有機會和大家分享.
我再遠一點的故事.
其實阿伯安罕12章的故事.
其實在我20年前.
上帝跟我說過一次.
要我離開.
所以這件事不斷在我腦裡.
到我今天PhD畢業.
我一樣想起阿伯安罕的故事.
所以我很想學他的故事.
因為我很相似他.
上帝叫我去我的地方.
我不知道我去哪裡.
我只知道去.
你去.
所以我真的在那幾個月.
不斷和我老公Selwyn去探索.
在亞洲不同的地區.
起碼在六至八個國家我們探索過.
到最後.
我是被迫和科學擁擠.
這裡又看到那裡又看到.
做夢我都看到上帝給我的科學.
我昨天有個同學聽完我說.
你經常做夷夢.
我說沒有.
其實我做很多夢.
不過從來沒有做夢上帝跟我說過.
不過上帝在我夢裡.
真的很清楚跟我說去日本.
我說你叫我去日本啊上帝.
怎麼可能去日本啊.
我從頭到尾沒想過去日本做宣教.
我很想去教神學.

$^{721}$教神學這個夢其實不在我這裡來.
我這些讀不到書.
我媽媽常說都不知道你有沒有大學.
怎麼會想著去讀PhD去讀神學呢.
到我讀完M.Div出來已經很厲害了.
我是很想去訓練老師去教好書.
因為這個今天沒機會說了.
我很想訓練老師去教好他們的書.
我想訓練教會的老師.
去幫他們可以培育基督教教育.
到後來上帝跟我說.
一個異象.
我之後就讀了PhD.
讀PhD Education.
我一心想著教神學.
但我知道在日本一個這麼hierarchical.
一個這麼矮的女士.
女士這麼矮.
不是日本人.
我怎麼可能在一個不懂日文的情況下.
去教日文神學.
但真是一而再再而三.
起碼有20個生.
是叫我們一家人去日本.
我說上帝.
You don't make sense at all.
我不是三天.
我用了三個星期的時間.
去到最後我跟上帝說.
我哭了跟上帝說.
上帝你要我去的.
我肯去.
如果真的你要我去.
我完全不知道我要做什麼.
但我肯去.
如是者從日本.
在我那天下午.
我肯去日本的時候.
上帝將我整個direction.
去了台灣.

$^{761}$我說去台灣來的.
去到台灣的時候.
去到台灣我突捉機場那一刻.
我收到我們警察會的消息.
Sorry其實我們不太需要你來台灣.
我心想我來台灣.
我真的搭了.
我不是誇張.
第一步.
我問我老公.
你為什麼check email.
我一看就知道是集團飛機.
人家說不要我們.
我說上帝你真是搞笑.
第二天我在台灣.
我收到香港給我的消息.
說對我有興趣.
我說你不要玩.
我怎麼會回香港呢.
我從來沒有想過回香港.
我跟上帝說我去哪裡都可以.
我只有兩個地方不去.
知不知道在哪裡.
就是香港.
我說我不去香港.
不過香港找我.
我真的不想去.
所以當亞伯拉罕的兒子說.
建制的高陽在哪裡.
上帝會預備.
你想多了.
如果你覺得他一定有信心.
我回香港其實有兩大challenge.
我今天沒有機會說很多.
其中一個challenge是international school.
我不知道你知不知道.
香港的international school有多貴.
我知道貴.
但我不知道這麼貴.
我送兩個女兒進去.

$^{801}$除了tuition之外.
我要給一百萬港幣.
丟下我送女兒進來.
然後每年繼續交錢.
我說我上帝哪有這個錢.
然後我們去interview.
都不知道怎麼辦.
然後interview.
原來我兩個女兒.
雖然英文很厲害.
原來我真的記得.
香港的小朋友.
從小到大.
是interview慣的.
我兩個女兒問他.
sorry啊女兒.
問他十個問題.
asked him 20個.
有15個都答I don't know.
另一個就說.
你介紹一下自己.
我想他心裡想.
介紹什麼.
為什麼要impress你.
因為他才說.
you are here to impress me right.
我相信我女兒會想答.
I am not here to impress you.
完了interview之後.
我知道他們擔心.
他問我.
我小時候有一個女兒問我.
如果學校不收我.
我沒有書讀.
那怎麼辦.
你猜我可以答什麼.
你真的interview得很差.
我不可以這樣跟他說.
我跟他說什麼.
上帝會預備.

$^{841}$上帝要我們回香港.
他自己會預備.
亞伯拉罕的生命.
無論是清早起來.
或者到了第三天.
他沒有說過他的生命.
他掙扎什麼.
但如果你真真正正.
看他的生命.
你明白.
我兒子對他來說的意義.
有多大的時候.
你不會覺得.
他是一個100分的陳頌基.
他是你跟我一樣.
是65分的那個.
有很多次.
你真的可以turn away.
說上帝我不玩了.
真的玩不起.
我不想回香港.
不想去日本.
很多次.
每個人congratulate我.
我是哭的.
我想告訴大家.
但這個就是我們跟從上帝的生命.
我的教父說.
I am surprised you say yes.
I told him.
我是不可以say no.
我沒得say no.
因為我很清楚這個上帝帶我去.
亞伯拉罕就是這樣.
不是因為我知道上帝.
可以怎麼make it happen.
我剛才告訴他.
待會我們.
你喜歡聽故事的.
我待會出去再說很多故事.

$^{881}$很多故事.
我告訴你.
我告訴上帝.
你這樣也可以嗎.
我說了很多次.
你真的這樣也可以嗎.
就像這裡這樣說.
最後是什麼.
當是手起刀落的時候.
就有使者從天上呼召.
亞伯拉罕.
我在這裡.
不可以向這個小孩下手.
你看看.
原來有隻羔羊已經在這裡.
在這裡預備給你了.
這個就是耶和華爾納的意思.
上帝要預備.
整段經文告訴我.
我們的上帝是誰.
在我們完全看不到.
根本不可能.
上帝你不make sense.
完全不make sense的時候.
無論你想什麼都好.
亞伯拉罕以為.
他的兒子會死裡復活.
我真的不知道.
上帝怎麼供應我回香港.
無論Housing.
無論Christian International School.
是一個.
貴到不可能負擔得起的東西.
但我想告訴大家.
上帝真的預備給我.
上帝真的準備給我.
上帝真的預備給我.
然後我跟上帝說.
你真的這樣也可以.
因為我說過.

$^{921}$上帝你要我回去很簡單.
你預備這兩樣東西.
我真的.
真的不關我的事.
我沒有說不去.
不過你自己搞定.
不過他真的搞定了.
一年之後.
我等了一年.
我那份工作有了.
但只有那份工作.
那份工作什麼都不夠cover.
上帝預備了.
如果有人.
在你身邊.
經歷一個這樣的.
完全不合理的情況.
我想請你.
剛才他說.
UV的弟兄姊妹要成長.
我想請你們一件事.
不要輕易跟弟兄姊妹說.
你信不信上帝.
真的有人這樣跟我說過.
你信不信上帝.
說真的.
你可以說什麼.
我當然信.
這樣就有信心.
我真的有掙扎.
我真的只有65分.
但上帝憐憫我.
這也是正正基督教教育所說.
當你真的要經歷.
一個生命的轉化的時候.
你的disorienting dilemma.
或者你的conflict.
或者你面對一個.
完全不make sense的那一刻.
就是你轉化.

$^{961}$有機會經歷這個轉化.
這個transformation的開始.
這個transformation不是.
突然間有.
什麼三秒鐘.
馬上得了.
不是那種黑了.
那種美元發財.
而是真的要經歷上帝.
你真的要走到.
你真的要走到上身.
真的是親手起刀的那一刻.
你才會見到上帝的公認.
是要等的.
你不要輕易.
這麼輕易跟你弟兄姊妹說.
你有信心就行.
你沒有信心嗎.
很多次我真的很struggle.
我跟我的女兒有一天說.
難道我的信心.
比這個屆賽總還少.
但上帝憐憫我.
上帝憐憫我.
不是因為我信心多.
也不是因為亞伯拉罕信心多.
而是上帝有恩典.
你只要是走下一步就行.
你不需要想這麼遠.
不需要想這麼遠.
你沒得想的.
你想來想去.
都是那些耳濕馬舌出來的.
是屎撬來的.
他會搞得一鍋泡的.
Just walk the faithful step before you.
只要走下一步.
仍然有人問我.
你覺得這是一個good decision嗎.
我說我不知道什麼是good decision.

$^{1001}$在這一刻帶兩個女兒回去.
我覺得我沒聽過.
我真的沒聽過有人.
好像我這樣傻的.
我不知道什麼是good decision.
但我知道什麼是God honoring decision.
弟兄姊妹.
經歷上帝的.
不是你有信心就行.
不要將別人的信心變成.
會有掙扎.
你的信仰是要經歷這個掙扎.
你要相信的是.
我不走每一步.
上帝看著我.
祂知道就行.
祂看得到就行.
我真的什麼都看不到.
祂看得到的.
這個上帝值得我跟的.
祂不是耍我.
我們就繼續走下去.
繼續去相信.
好像剛才四樂的歌那樣說.
即使我看不見.
還是看不見.
仍然相信你的上帝是信實.
仍然相信上帝看見一些.
你從來看不見的東西.
就這樣儘管跟著你的耶和華行.
為大家祈禱好嗎.
天父上帝多謝你給我機會.
跟大家分享一個.
你在我過去這段時間的經歷.
亦可以從亞伯拉罕身上去.
看回去想回.
你不是要我們好像一個.
完全信心無敵的人.
甚至耶穌基督.
他也在哈薩馬尼園掙扎.

$^{1041}$但我們也很清楚.
原來我們跟從上帝的.
不是按照我們自己的心意.
有時甚至我們不知道.
什麼叫好什麼叫不好.
可能我們連下一步的力都沒有.
但主耶穌求你讓我們記得.
既然我們身邊有這麼多.
六十五分的人.
這麼多六十分的人.
都仍然願意去跟從上帝.
仍然選擇走我們.
下一步的信心.
就好像亞伯拉罕.
有很多次可以走回.
但仍然在走的時候.
在我們微小的信心.
我們信不足的時候.
上帝連問我們.
讓我們成為彼此的鼓勵.
就不是鞭打大家的信心.
而是鼓勵大家.
終有一天我們會看得見.
經歷得到上帝早已預備的地方.
將榮耀歸給你自己.
不是因為我們完美.
而是在我們軟弱上.
You are perfect.
Thank you in Jesus name.
Amen.
多謝大家.
再見.
\newpage



\section{使徒行傳}
\label{sec:rAFDbYbX6HA}
\textbf{天國的啟動 (使徒行傳 1) - 陳崇基牧師}
\newline
\newline
連結: \href{https://youtube.com/watch?v=rAFDbYbX6HA}{\texttt{https://youtube.com/watch?v=rAFDbYbX6HA}} ~~~~ 語音日期: 2025-05-13
\newline
\newline
\hyperref[sec:Wr2X99fNMqk]{\small{< < < PREV SERMON < < <}}
~
\hyperref[sec:index_chronic]{\small{[返順時目]}}
~
\hyperref[sec:index_scriptual]{\small{[返順卷目]}}
~
\hyperref[sec:bewfsUoYndE]{\small{> > > NEXT SERMON > > >}}
\newline
\newline
$^{1}$祝大家聖母節快樂.
我們繼續跨越.
我們說過UV九周年之後.
我們有一個這樣的年提.
就是跨越.
有三點.
第一就是問自己有什麼要安恕.
我們需要克服和跨越出來.
成長就需要這樣.
我們繼續向前就需要這樣.
還有自己的圈子裡面如何跨出來.
和其他的大英姐妹.
和教會裡面的信徒.
能夠一起同行一起成長.
第三方面就是教會如何把持我們的福音.
能夠向我們鄰舍跨越出去.
在四面牆跨越出去.
來向人分享福音.
我們說這一年會看《士徒行傳》.
所以我們會選讀《士徒行傳》的某些經文.
這次我們看第一章.
有三點.
第一章告訴我們.
神國裡面有什麼應許給我們每一個.
給每一個世代的教會.
給每一個基督徒.
第二方面就是上帝的計劃.
這個神國的計劃擴展會是怎麼樣.
第三方面是如何邀請我們去參與.
所以就是這些簡單的三點.
我們看《士徒行傳》第一章.
這裡說到路加寫第二集.
你知道路加寫了路加福音.
然後寫第二集.
續集就是《波羅的X》.
就是《士徒行傳》.
然後他開始的時候第一章第一節.
他就說提亞非羅.
或者提亞非羅大人.
這個人可能是一個羅馬的官.

$^{41}$或者一個外邦的信徒.
他可能是贊助路加寫作的.
因為當時寫作是很多錢的.
所以他贊助他.
或者他剛剛信主.
所以路加就寫第一本書.
《路加福音》.
還有第二本書就是《士徒行傳》.
他就去講述耶穌所做的事情.
還有耶穌之後的門徒做了什麼.
這個就是《士徒行傳》所提到的事情.
他說在第一本書已經講到.
耶穌從頭所做的.
所教導的一切事.
直到祂藉著聖靈.
吩咐所揀選的使徒之後.
被接上天的日子為止.
祂一受害之後.
耶穌做了什麼呢.
祂用了很多證據向使徒.
顯明自己是復活了.
然後四十天.
祂向門徒顯現.
然後祂才升天.
和門徒一起談什麼呢.
就是談上帝國的事情.
神國的事情.
Kingdom of God.
到底這個Kingdom of God是什麼呢.
這個就是《士徒行傳》想交代.
告訴我們是一回什麼事情.
所以第四節他繼續講.
他說耶穌和門徒一起.
在四十天裡面聚會的時候.
他就叮囑他們做什麼.
他說你不要離開耶路撒冷.
要等到被你所應許的.
The promise.
之前講到Promise的Kingdom.
現在講到Promise的什麼呢.

$^{81}$就是聖靈.
就是你們曾經聽過的.
約翰施洗.
約翰之前就是用水作為一個記號.
來為你們施洗.
現在你要等什麼.
就是在聖靈裡面受洗.
所以他們聚集的時候.
門徒就問耶穌.
他說主啊.
你要復興以色列國.
是什麼時候呢.
這是一個很重要的問題.
因為之前講什麼呢.
上帝Promise了一個Kingdom.
Kingdom of God.
然後門徒就問他.
你復興這個Kingdom.
這個所謂Kingdom of Israel.
到底是什麼時候展開的.
之後耶穌就回答.
這個這麼重要的問題.
到底這個神的國度.
是什麼時候展開的.
我不知道你有沒有一些.
你聖經裡面最喜愛.
或者最記得的經文.
有沒有.
John 3:16大家認識吧.
神愛世人等等.
還有呢.
凡事謝恩等等.
我們每個人都有些金句的.
但是以前的猶太人.
耶穌時代的猶太人.
他們的金句是從哪裡來的呢.
就是從以塞瓦書出來的.
你想都沒想過吧.
以塞瓦書.
你知不知道以塞瓦書有些金句.

$^{121}$最多都是.
我的道路不是你的道路.
我的上帝的道路高過你的道路等等.
其中cut一節出來的.
以前的猶太人.
第一代的信徒.
第一代的猶太人.
他們最喜歡的聖經.
居然就是以塞瓦書.
特別是以塞瓦書的下一半.
如果我們去耶路撒冷以色列的博物館.
走去一個展覽館裡面.
那裡有一個圓形的東西.
你走進去一個黑漆漆的地方.
很黑的.
連樓梯都未必看到.
差一點很容易就PK.
很黑的.
那裡有一個很大的卷子.
一個圓筒.
上面有一個手柄.
其實就是一個手抄本的卷軸.
那裡展覽.
你看到那個人在看的.
就是以塞瓦書卷.
這個就是1947年48年的時候.
在死海裡面找到的死海古卷.
他就在那裡展覽.
那個死海古卷就是.
那個叫Great Isaiah Scroll.
偉大的以塞瓦書卷.
為什麼呢.
因為在死海古卷裡面找到的.
不同的抄本,殘卷等等.
最多的就是以塞瓦書.
他們寫的《色經》書.
還有《講論神學》的東西.
都是很多時候來自以塞瓦書的.
特別是以塞瓦書的下一半.
以塞瓦書有66章.

$^{161}$下一半就是39之後.
40之後到最後.
這個所謂第二以塞瓦書.
就是第二部分的以塞瓦書.
那裡講什麼呢.
有很多的主題.
那裡講到對那些被流亡到海外的猶太人.
神的子民.
上帝說我現在要憐憫你.
我的國要臨到.
有幾個主題.
他講到分散在各地不同的十二支派.
將會被重整.
然後聖靈會降臨.
當上帝的救恩來到的時候.
上帝的國展開的時候.
上帝的子民也會成為他的見證人.
並且上帝是憑著他的話語來工作的.
所以他的話語是不會徒然返回的.
最後神的子民就有一個redefinition.
就不是純粹一個種族血統來分.
就是外邦人.
所有的outcast.
包括太監.
包括其他那些.
我們覺得.
以前猶太人覺得是罪人.
或者那些是不可以親近神的人.
全部都會收納成為上帝的one people of God.
新的上帝的子民.
這裡當然有很多經文.
我們只看一節.
49章第6節.
那裡講到上帝的僕人來到.
展開上帝的國度.
和成就上帝的救恩.
上帝說什麼呢.
你作我的僕人.
使雅各眾支派能夠復興.
就是十二支派.

$^{201}$你知道當時流亡之後.
亡國之後.
北面的十個支派已經一早消聲匿跡.
沒有了.
所以講這句話其實是挺詫異的.
我們十二支派不是全散了嗎.
什麼都沒有了.
怎麼能夠回頭.
上帝說不是.
我會做一件新事.
就是雅各以色列的眾支派會復興.
被保存的人能夠歸回.
你們以為這件是小事.
我還可以使你們成為萬邦之光.
並且將我的救恩.
將我的名字傳到他們的.
就是上帝會恢復十二支派.
新的國民會產生.
這個叫做Israel 2.0.
Israel 2.0不是Beta.
第二版.
是一個更加大的新的神的子民.
這個就是上帝所promise.
所認許給每一個世代的信徒.
所以去到史唐傳的時候.
耶穌就回答了.
他說這個是上帝的計劃.
上帝早在《以賽亞書》裡面.
已經promise了這個kingdom.
現在就由我的死而復活.
我的升天.
我作王.
能夠去開始的.
第一章第七節.
耶穌就回答門徒了.
他說父憑著自己的權命所定的日子和時間.
是你們不可以知道的.
不過.
聖靈要降臨在你們身上.
你們就能夠得著能力.

$^{241}$並且在耶路撒冷.
猶太傳地.
撒瑪利亞.
直到地極.
都可以作我的見證.
在這裡告訴門徒.
你不需要計算多少時候.
總之聖靈降臨的時候.
你們就開始這個計劃.
This is the plan.
就是早就promise了給我們.
給我們新的神的子民.
能夠走到地極.
來作上帝的見證.
witnesses.
這句聖經當然很熟.
第一章第八節.
是不是.
第一章.
你記不到其他東西.
記到第一章第八節已經很足夠.
聖靈.
以賽亞說到的kingdom.
聖靈就會降臨.
然後你們就會得著能力.
然後去傳道.
去擴張上帝的國度.
由你們本位開始.
耶路撒冷.
猶太傳地.
撒瑪利亞.
直到地極.
這一句的金句.
所謂一章第八節.
當然很重要.
因為整個小行傳差不多.
就是按照這個分佈來寫出來的.
由耶路撒冷開始.
第一至第七章.
Jerusalem.

$^{281}$然後第八章.
這個福音.
這個聖靈所帶領的群體.
他們就衝出耶路撒冷.
然後去到猶太傳地.
和撒瑪利亞的地方.
直到地極.
到最後十三至二十八章.
就是說到保羅和他的宣教的對公.
去到不同的地方.
最後福音去到羅馬.
第二十八章.
這個就是小行傳裡面有的.
所謂脈絡.
那個flow.
這些傳到地極.
將上帝的國度.
將那個power.
傳到很遠很遠的地方.
這些不只是聖經有這樣說.
不只是當時第一代的信徒.
小行傳這樣說.
當時羅馬帝國也有這樣的用語.
羅馬帝國大家知道.
當然早就覆亡了.
但最近選的教宗.
我們不要以為是翻版.
不是的.
羅馬帝國早就沒有了.
早就已經是歷史沉寂.
以前的羅馬帝國.
特別是當時門徒的時代.
耶穌的時代.
他們都有很多這樣的.
直到萬邦.
去到地極的用語.
譬如.
有一個叫Virgil.
他寫了一些著作等等.
這個是一個羅馬帝國宮廷的.

$^{321}$一個詩人.
他寫了一本書.
講到國家立國的時候的故事.
每個國家都有一個.
所謂foundation story.
他們立國是怎樣的.
開始的時候是怎樣的.
這個beginning of the kingdom是怎樣的.
等於我們每逢到11月的時候.
我們美國就會有一個國家的故事.
我們每個人都會聽的.
就是那個thanksgiving story.
mayflower等等.
然後我們來到的時候.
我們怎樣發生什麼事.
這些所謂美國American kingdom的foundation story.
羅馬都有它的foundation story.
Virgil講到羅馬開始的時候.
他們京都開始擴張.
去到地中海不同的地方.
最初最初的時候發生什麼事.
就是神明給了羅馬的founding fathers.
那些國父.
一個的應許或者一個的答應.
就是什麼呢.
就是emperor without end.
給他們一個無止境的帝國.
所以羅馬帝國擴張到不同的地方.
霸佔整個歐洲.
整個地中海的世界等等.
另外一個的文獻.
就是第一個的羅馬帝王.
first emperor.
叫做奧古斯都.
他有他的威水史的石碑.
到處放.
然後copy.
他又放在其他地方.
告訴人家他有多威水.
有多多功德.

$^{361}$他的豐功偉績是什麼.
他就說什麼呢.
他說他將世界所有的蠻族.
伏在羅馬帝國之下.
這個Westchester的文獻.
或者這個石碑.
上面提到地中海附近.
很多不同的民族.
你數一數.
有25個的種族.
這個奧古斯都就說.
我將這25個的種族.
大大小小的種族.
都征服了.
然後放在羅馬的權下.
這些是當時中央電視台說的話.
羅馬帝國.
羅馬帝王有多厲害.
其實古今中華大的國.
great power.
都會有這樣的宣傳的.
美國來了多久.
特別是來了久一點.
我們已經聽慣了.
我們說我們是.
leader of the free world.
是不是.
所以我們令到大家有秩序.
或者令到大家沒有秩序.
都是入我們數的.
因為我們是leader of the free world.
這個是American empire.
有的propaganda.
我們有的政治宣傳.
中華民國偉大復興都有的.
所以古今中華這些大的國.
都會這樣的.
羅馬帝國他想象的kingdom.
所想象的國.
就是怎樣呢.

$^{401}$就好像這個地圖.
你知道地圖畫疆界是怎樣畫的.
就是他霸佔的地方深色一點.
所以由那裡開始.
羅馬就是那裡.
然後他就向東南西北.
然後佔領不同的地方.
地中海不同的地方.
直到他們的所謂地極.
可能門徒生活中的以色列國.
Israel 2.0都是這樣的.
只不過他們的起點是在耶路撒冷.
一直這樣塗顏色.
好像塗顏色專頁一樣.
小朋友填顏色那些.
有空位就填.
這個就是地上國度的想象.
你想到kingdom的時候.
想到國家的時候.
我們通常就會這樣想象.
但是耶穌和以賽亞.
舊約裡面所promise的kingdom.
那個願景就不同.
就不是一直這樣塗顏色.
佔領不同的地方.
如果以地圖來表達的話.
這個才是kingdom的國家.
是沒有疆界的.
所有信奉耶穌為利塞亞為基督的.
他就是其中一個神的子民.
他去到不同的地方.
將福音傳開.
然後那裡又有另外一點.
然後又有另外一點.
然後分佈在不同的地方.
這個才是耶穌所想的kingdom.
耶路撒冷猶太傳地.
直到地極.
猶太傳地.
撒瑪利亞等等.

$^{441}$直到地極.
能夠每一點光是遍地開花.
這個是kingdom of God.
不是of Rome.
不是of China.
不是of United States.
這個和地上國度的想象完全不同.
我們就是那幫信奉耶穌為救主的那幫人.
我們就到地上不同的地方.
好像一點光遍地開花一樣.
所以一個跨越民族.
跨越地理的.
跨越所有的boundary的kingdom.
所有信奉耶穌的.
領受聖靈的.
受洗歸主的.
就成為神的國度.
這個才是我們的kingdom.
這個才是上帝所說的kingdom.
特別說到種族.
耶路撒冷猶太傳地.
撒瑪利亞直到地極.
我們很多時候想的就是地理.
Geography.
其實在他們的印象裡面.
你講到耶路撒冷猶太傳地.
撒瑪利亞直到地極.
他們很多時候想的不只是地理.
不是在想Google Map.
他們在想種族.
猶太人和其他種族有很多不同的關係.
特別說到撒瑪利亞.
撒瑪利亞以前是北國的.
以色列民族的地方.
十個支派在那裡被滅亡.
對於猶太人來說.
那裡滅亡之後.
有其他種族來了.
然後他們叫他們撒瑪利亞人.
他們和撒瑪利亞人是不會同桌吃飯的.

$^{481}$也不會有任何交往的.
是猶太人不能踏足之地.
大家看過福音書.
大概應該有這樣的印象.
撒瑪利亞人和猶太人是不能夠混在一起的.
但是上帝說我的國度是一個跨種族的國度.
是不講種族的.
所以提到撒瑪利亞的時候.
就等於我們提到什麼.
譬如我們講收沙.
收沙是什麼地方.
你不想去的地方.
你講到英國的時候.
你不只是講地圖的地方.
那個區.
你講人種.
你講ethnic group.
上帝福音的國度.
就是要打破所有這樣的隔膜.
信徒就能夠成為主的見證.
一路遍地開花.
走到不同的地方.
This is the kingdom of God.
你所參與的國度就是這個國度.
所以你不要想著我們再重溫.
或者有我們的nationalism.
所有nationalism 國家主義.
地上的國度的想像.
都不是以賽亞.
不是耶穌.
不是小行傳裡面所講的.
小行傳講的就是我們這班新的神的子民.
有聖靈降臨在我們身上.
我們有能力能夠為主見證直到地球.
這個才是kingdom of God.
而上帝就是邀請我們.
來參與這個kingdom.
怎樣參與呢.
在第一章之後他繼續講.
耶穌升天之後門徒就去了哪裡呢.

$^{521}$第十二節他說有一座山叫做.
橄欖山.
Mount Olive.
離耶路撒冷不遠.
有一個安息日可走的路程.
大概2km.
一英里鬆一點.
那時門徒就從那裡橄欖山回到耶路撒冷.
他們就進城了.
就上到他所住的樓房.
裡面有一些人.
他提到12個門徒.
所謂彼得,約翰,雅各等等.
還有一些婦女.
和耶穌的母親瑪利亞.
耶穌的兄弟.
他們做了什麼.
是同心合意的去行徹禱告.
那時有很多人一起聚會.
120人.
彼得就在弟兄中間.
起來他就開始講論一些東西.
然後講了一篇道.
然後引用了一些詩篇.
他說由大因為賣耶穌.
他的位置已經被廢去了.
所以我們中間就差一個人.
我們需要補回這個人.
大家記得嗎.
我們只有一個使徒.
我們需要在我們從耶穌在世的時候.
他傳道經歷裡面.
和我們一起一班.
跟著的那班是目擊見證人.
要選一個出來.
於是他們就搖籤.
然後就搖到一個人.
抽籤抽到一個人.
就叫馬提亞.
然後最後一節.

$^{561}$他就說馬提亞與11個使徒同列.
He was numbered with the 11.
所以有12個.
為什麼有12個.
剛才我講這首話的書講什麼.
復興十二支派.
所以一定要有12個.
你見到他又強調什麼.
120位.
這個全部都是舊約裡面.
以賽亞特別第二以賽亞裡面所講的那個kingdom.
一個new nation.
new Israel.
新的以色列.
是十二支派歸位的意思.
就是補選.
像不像少林足球.
少林足球那些師兄弟.
被人打到散了.
我透露了年齡.
都是少林足球.
最後的時候就叫他們歸位.
然後是時候我們打球.
師兄弟我們是時候醒起來.
是時候我們認真打一場好球.
結果他們打了一場好球.
就是這個意思.
十二支派已經散了.
但是上帝憑著聖靈的能力.
叫十二支派再重新歸位.
所以選上馬提亞.
其實是有意思的.
因為要fulfill這個plan.
這個plan就是我們每一個都有份參與.
在十二的裡面.
是每一個信徒都有的.
上個禮拜最大的宗教的新聞.
當然就是羅馬天主教選了新的教宗.
新的教宗.
Leo14.

$^{601}$他綽號Leo14.
還是芝加哥人.
所以我們很開心.
其實跟我們沒有什麼關係.
跟我們基督教沒有什麼關係.
當然天主教跟基督教基本上是一樣.
基本上是教義.
只不過有些practice是不同的.
譬如有些practice.
天主教是很著重某一個中心的人物.
就是教宗.
所以有些人專程飛去羅馬.
來看他黑煙白煙的.
但是在聖經裡面說什麼呢.
其實聖經裡面很著重一件事.
叫做信徒皆濟事.
Priesthood of all believers.
所有的信徒.
男與女.
地位高與低.
種族是哪一個的.
都是濟事.
都是神聖的職份.
我不論你是做哪一行.
或者你甚至是一個全時間的家庭主婦.
全時間的家庭主婦.
更加忙於我們是不是有工作.
每一個的職業.
每一個的位份.
都是神聖的.
都是priesthood.
You are priest.
或者這麼說.
你就是牧師.
牧師不是一個兩個.
他只不過有一個職位.
但是跟我們的功能.
跟我們的位份其實都是一樣的.
你試一下說一說你姓什麼的.
你說你姓什麼.

$^{641}$牧師.
你叫你自己.
你自己叫自己.
你拿個電話selfie的時候.
我就是李牧師.
我就是林牧師.
Priesthood of all believers.
聖經裡面告訴我們.
每一個的信徒就是正在參與上帝的工作.
上帝邀請我們.
你都去參與.
馬提亞在《史壯傳》裡面.
當然沒有提過這個人的事蹟.
你可以這麼說.
他是蛋散.
補位.
等於踢球的時候.
他們坐在bench那裡.
坐在旁邊.
有一個人受傷了.
出去了.
然後就叫他出來.
他就下場了.
但是有一個最重要的東西.
就是他下場了.
上帝邀請我們每一個都下場.
你有沒有看過球賽.
有看過.
現場那些人.
現場我們在觀眾席的時候看的時候.
我們只不過是壽終.
我們不會和其他球友有很好的交往.
不一定的.
有時候甚至是仇敵.
因為他捧另一隊.
我們不會和他有交往.
因為我們不參與.
但是你想像一下.
如果你是參與的.
你是下去踢的.

$^{681}$你是下去打籃球的.
你的身份和你的角度.
你的心態就完全不同.
上帝說我們每一個都是參與的一部分.
在史萊恆傳裡面.
我們有時候會討論.
比較大衛的討論.
史萊恆傳裡面主角到底是誰.
你知道可能有些人說是彼得.
因為彼得開頭的十多章.
都說很多他的事蹟.
然後他就好像消失了.
然後就保羅.
所以有些人說可能是保羅行傳.
可能是彼得行傳.
因為沒有彼得就不會有保羅等等.
或者我們甚至說是聖靈行傳.
因為聖靈帶領著耶穌的門徒.
去到不同的地方傳道.
去到擴展上帝的國.
但是你有沒有想過.
史萊恆傳裡面的主角.
其實是教會.
其實是上帝的子民.
是那些無名氏.
是那些沒有留名的那些人.
他們將福音傳到其他地方.
不是哪個大目.
哪個名目.
哪個很有名的傳道人.
或者哪個很有能力的人.
是每一個無名氏.
每一個的蛋散.
是他們肯參與上帝的工作.
然後將福音.
將上帝的國來見證.
史萊恆傳的主角.
應該是教會.
新的人.
有聖靈能力去驅動的那些人.

$^{721}$就是你和我.
你就是那位牧師.
我就是那位牧師.
你就是那位傳道.
你就是那位宣教士.
你就是祭司.
所以這一年裡面我們希望.
我們在UV.
我們鼓勵弟兄姊妹去參與.
這個神國的project.
這個工程.
這個project是怎樣的.
每一個人在十字架底下.
我們被上帝的愛.
上帝的恩典改變了.
死而復活的福音改變了之後.
我們就和其他的.
不是觀眾.
其他的參與者.
我們聯合起來建立教會.
最後我們更加需要將這個工程.
帶到更遠的地方.
去到不同的種族.
或者和我們不同的人.
這個就是我們UV這一年.
我們希望能夠做到的.
所以我們不只是單單.
擴堂這麼簡單.
不只是想找一個大一點的地方.
是我們的人.
我們的心胸要更加大.
我們要有上帝國度的那種想像.
有聖靈的幫助.
然後我們將福音傳到更遠.
更多的人身上.
所以今年我們UV.
我們希望再review你們的功課是什麼.
你需要做的.
是想想.
有一個人.

$^{761}$你仍然知道他未信主.
你是想將福音傳給他.
每一天請你抽5分鐘的時間.
你都可以為他禱告.
你知道為什麼你信主嗎.
因為有人這樣做.
有人膽敢為你禱告.
膽敢走到你面前.
來和你分享福音.
膽敢告訴你他是信徒.
所以你為這個未信主的禱告.
最後求主賜給你勇氣.
並且有智慧.
去傳福音.
在他面前做一個見證.
這就是我們每一個人的禱告.
求主幫助我們參與他的工作.
我們一起低頭祈禱.
神啊求你將我們一些.
難阻我們看到神國的事情.
無論是我們的工作.
還是其他的.
我們要處理的很多的雜務.
事務.
主暫時將這些拿開.
讓我們看到.
我們自己的國度降臨的時候.
帶著聖靈的能力.
是你早早已經應許的.
並且在基督的死而復活.
升天裡面已經展開.
主要是利用我們.
我們好像每點光一樣.
去到世界不同的地方.
社區裡面不同的地方.
去見證你自己的國度.
讓我們有著上帝的恩憐.
你給我們的力量.
聖靈給我們的能力的時候.
我們好像遍地開花一樣.

$^{801}$將你自己的國度來擴展.
主幫助我們UV這個教會.
讓我們來到這個教會的時候.
我們怎樣去想怎樣去建立.
為這個教會禱告.
在這裡事奉.
我們在我們的生命當中.
我們有的崗位.
我們忠心的去工作.
忠心的去見證你.
主要是為了.
可能有一位的朋友同事家人.
你將這個面孔.
將這個名字放在我們心裡面.
主要是為他們禱告.
主要是求你的恩典淋到他們.
讓他們能夠認識你.
求主賜給我們力量.
賜給我們智慧.
讓我們將福音傳開.
我們就是你的見證人.
我們就是你的祭祀.
我們就是牧者.
主拆敗我們去到不同的地方.
將你自己的國度美好的福音.
能夠做一個美好的見證.
使用我們教會.
以致我們在這一年裡面.
我們更加見到上帝自己的恩典.
使我們教會能夠繼續成長.
繼續向前.
繼續能夠擴展.
是為了你的名的榮耀.
主要是求你聽我們禱告.
我們這樣祈禱.
奉主耶穌基督的名字.
元神祂祝福各位.
下個星期能夠和大家一起來這裡事奉.
祝你們聖誕快樂.
\newpage



\section{}
\label{sec:bewfsUoYndE}
\textbf{對兒女的最大期望:智慧 - 黃國棟醫生}
\newline
\newline
連結: \href{https://youtube.com/watch?v=bewfsUoYndE}{\texttt{https://youtube.com/watch?v=bewfsUoYndE}} ~~~~ 語音日期: 2025-05-21
\newline
\newline
\hyperref[sec:rAFDbYbX6HA]{\small{< < < PREV SERMON < < <}}
~
\hyperref[sec:index_chronic]{\small{[返順時目]}}
~
\hyperref[sec:index_scriptual]{\small{[返順卷目]}}
~
\hyperref[sec:MAWjFwIPHJE]{\small{> > > NEXT SERMON > > >}}
\newline
\newline
$^{1}$其實她太過高舉我.
我只是第一次介紹自己.
是一個渾身都沒有張力的人.
分享一些值得跟大家分享的訊息.
五月和六月是父母的大日子.
上星期我們剛剛慶祝了母親節.
下個月就慶祝父親節.
所以今天想跟大家說一下.
父母對子女會有什麼期望.
以及我們可以怎樣幫助子女.
甚至幫助自己.
來得到我們想要的就是智慧.
我想有哪個父母不想子女成龍成鳳.
想子女蠢蠢地,隨隨便便.
騙吃騙喝,混混樂樂過一生就算了.
我想沒有的.
做父母的我們每一個人都是希望望子成龍.
望女成鳳.
現在男女平等.
所以以前就說望女兒不用那麼厲害.
嫁得好就算了.
現在不會了.
當然望她厲害.
但我們怎樣可以令子女厲害.
可以令他們成龍成鳳.
成龍成鳳.
我想一個最基本的條件就是要有智慧.
這件事沒有人會反駁吧.
你沒有智慧很難成龍成鳳.
可能你家裡被父二代混混樂樂.
不休息,不休住也說不定.
但你想他可以做到一些有價值的事.
他們就需要有智慧才能做到.
對子女的期望.
智慧是我們第一樣需要的東西.
這在聖經裡.
在《箴言》裡已經說得很清楚了.
所羅門這樣說.
為什麼寫這些《箴言》呢?.
是因為要人得智慧和分晦.

$^{41}$分辨通達的言語.
使人處事領受智慧,仁義,公平,正直的分晦.
使愚人靈明.
使少年人有知識無略.
我不讀了.
很長的經文.
大家可以回去繼續讀《箴言》這段.
但是它說到的就是智慧的用處.
使人學問增長等等.
最重要的是.
在一切事上.
智慧是第一的.
智慧為首.
我們首先要有智慧.
才可以有聰明.
所以做父母的.
我想我們每一個人都希望子女得到智慧.
我們到底有什麼方法.
或者有什麼.
當我們在尋求智慧中.
是值得注意的呢?.
今天就想跟大家.
講一下這個題目.
首先我們會問第一個問題.
我們為什麼要有智慧?.
傻傻的.
不是其實挺好的嗎?.
十下十下.
可能是不錯的.
嘻嘻.
嘻嘻嘻嘻.
這樣就一輩子了.
現在.
我們的制度可以保障我們.
起碼不會餓死.
不會凍死.
是不是這樣就.
這樣就行了呢?.
不用去想這麼多東西.
其實智慧可能都挺煩的.

$^{81}$有時候.
你要想很多東西.
令你很煩也說不定.
是不是呢?.
為什麼我們要.
尋求智慧呢?.
其實.
我會覺得.
智慧有兩個功效.
它有一個消極的功效.
passive的功效.
有一個積極的功效.
這兩個功效都是我們人生中.
必須的.
智慧不是讓我們賺大錢.
不是讓我們發達.
但是.
它是讓我們怎樣過我們必須的生活.
第一個是消極的功效.
這個是一個保護性的功效.
《聖經》裡面告訴我們.
智慧的作用就是保護我們.
它救我們脫離淫婦.
就是那些油嘴滑舌的外侶.
它是叫我們不要被人騙.
我說的不要被人騙.
不單是說那些.
電騙.
spam.
告訴你.
近來經常收到的.
就是說你的車過了那些.
兜兜費.
沒有交錢.
就叫你快點去登記交錢.
也不是說叫你.
或者是緬甸什麼KK園.
弄個電話來告訴你.
做什麼那些.
智慧會保護我們.

$^{121}$但是更加重要的是.
智慧在我們的信仰裡.
是會保護我們.
不要被那些披著羊皮的狼騙到.
因為在現在這個世界上.
撒旦攻擊教會.
不是明刀明槍的來跟我們打.
而是他們很會裝模作樣的來騙我們.
很多信徒.
如果沒有智慧的話.
就很容易被這些假宗教.
假教使壓到.
有朋友問過我.
城會裡面的朋友.
或者是其他的朋友都問過我.
這裡沒有了那個screen.
沒有了我的powerpoint.
有些技術問題.
不過先說個故事.
我怎麼認識Mark Sir呢.
剛才利先說過.
我來了你們教會不是很久.
其實我之前來過一兩次.
不過沒有什麼特別介紹.
我也上網.
在你們教會.
很多時候都會在你們教會崇拜.
所以其實我覺得這裡.
我不覺得陌生.
但是有朋友問過我.
我怎麼認識Mark Sir.
我認識Mark Sir十幾年了.
我最初認識他是怎樣的呢.
原來呢.
如果大家還記得.
十幾年前有一個機構.
就說他們發現了方舟.
於是就叫人捐錢.
很多錢.
但是呢.

$^{161}$那時候我自己就覺得.
不太妥.
他講的東西.
勾我的頭.
聖經好像不是.
跟聖經講的東西不是很對.
那時候我還沒有認識Mark Sir.
Mark Sir是考古的專家.
大家也都知道.
他講過很多這方面的東西.
這樣.
他也都寫過很多文章.
講這件事有什麼不妥.
那時候我是用另外一個角度.
就是一個讀科學的人.
也是一個讀法律的人的角度.
說這個機構講的東西.
不太make sense.
經不起考驗的.
於是我就寫了一些文章去反駁他們.
那Mark Sir又寫文章來反駁他們.
後來有一班教務.
就是.
會出了一個聯署.
也出了一本書.
可能大家也都見過.
收集了一些文章.
是說.
這個講法不妥.
不科學.
也都是不合理.
叫信徒小心.
最近兩年前這個機構又出來.
講話的Mark Sir也是在網上反駁過他們.
但是十幾年前Mark Sir做.
《腎思明辨》這本書的時候.
就想將我一些文章收集在裡面.
於是就和我聯絡.
我就在那時候認識了Mark Sir.
就有十幾年一直和他聯絡.

$^{201}$這件事.
可以說是很幸運.
也可以說是很生氣.
幸運的就是也有一班教務.
願意出來.
是為了我們的真理.
反駁這些似是而非的教導.
但是很生氣的是.
還有信徒繼續被他們迷惑.
甚至在一兩年前他們繼續出來.
又叫人給錢.
依然是百足之蟲死而不僵.
始終都不可以令到他們.
真的乖乖地做好.
不要靠教會來尋求個人的利益.
這就是智慧給了我們第一個功效.
可以保護我們如何去分辨一些.
似是而非的教訓.
上星期Mark Sir說到.
我們每個人都是牧師,祭師.
我今天想加多一件事.
我們每個人都要做神學家.
因為我們需要一些基本的知識.
才可以幫助我們.
去分辨到什麼是真的神學.
什麼是一些偽的教導.
一些假的教導.
我們需要智慧.
智慧給了我們第一個功效.
就是保護我們.
我們教子女的時候.
第一件事.
也不是教他們一些教條.
而是教他們一些智慧.
希望他們有智慧.
可以分辨到什麼教導是真的.
什麼教導是假的.
這是智慧消極方面的功能.
是保護我們.
第二個智慧的功能就是Enable.

$^{241}$是積極地可以幫助我們做事.
就是智慧必是你行善人的道.
守義人的路.
因為智慧有很多好處.
這裡智慧給了我們的.
也不是給了我們有能力去賺錢.
有能力去做到某間公司的CEO.
或者很高的收入.
智慧可能可以這樣做.
但不一定這樣做.
因為智慧會告訴我們.
賺錢其實不是人生最高的目標.
我們人生最高的目標是行善的道.
但是智慧可以幫助我們.
在很多問題上做到需要做的取捨.
我在大家當中分享過兩次醫療方面的問題.
大家可能也會聽到我說的.
很多時候在醫療方面的問題.
不是黑白分明的.
事實上在很多宗教的問題上.
在很多道德性的問題上.
也不是這麼簡單黑白分明的背教條就夠了.
譬如你有一個親人.
插了喉睡在那裡.
你應不應該幫他拔掉呢.
這個是怎樣處理死亡善終的問題.
不是背著教條說.
生命是神掌管的.
那就神掌管吧.
那就甚麼都不做.
或者不可以做甚麼.
或者一定要做甚麼.
這些是一些兩難的問題.
是需要有智慧去處理的.
上個星期跟大家分享的時候說過.
如果你想拿Medicaid.
有一個方法可以保障你配偶.
還有一些生活的.
就是Medicaid離婚.
但離婚教會說是不行的.

$^{281}$但為了拿Medicaid而離婚.
那是不是離婚呢.
我沒有答案的.
我絕對不是叫大家做或者不做.
我只是說這是一個兩難問題.
道德上的難題.
我們需要智慧去處理.
而不是背教條.
背了就說可以或不行.
就可以解決的.
我們需要智慧去enable.
去賦予能力我們去做這些決定.
這就是智慧的功效.
去保護我們不要被人騙.
特別是不要被一些假的宗教術語騙到.
去enable.
給我們能力去做一些應該做的事.
特別是在一些道德上兩難的問題上.
我們需要智慧幫助我們做這些事.
好,說說我們做父母的.
我們可以怎樣幫助我們的子女有智慧.
首先我想說兩個父母經常犯的錯誤.
希望大家留心.
第一個錯誤就是.
父母有兩個極端的.
這個在華人文化裡是更加常見的.
我們要贏在起跑線.
於是不斷地補習.
學琴,學畫畫.
打球,打跆拳道.
去做演講.
去做社會服務.
去醫院做觀察.
去做辯論.
總之那個日程表要排得滿滿的.
我不是說做這些事不好.
千萬不要誤會.
我不是說這些事不應該做.
我是說不應該為了贏在起跑線.
而監你的子女.

$^{321}$把日程表排得滿滿的去做這些事.
為了拿一張證書.
為了拿一個證書.
讓他們報讀學校的時候.
那個應用程式寫得好看一點.
這樣做是絕對不會幫他們有智慧.
甚至不會幫他們得到什麼.
這些都是好事.
打球是好事.
做運動是好事.
學琴是好事.
學畫畫.
去做義工.
去學跆拳道.
去學辯論.
這些全部都是好事.
但千萬不要為了贏在起跑線.
想拿多一張證書.
就監你的子女去做這些事.
這樣真的會害了他們.
說起來有很多人問我.
做什麼運動最好.
這是我的偏見.
如果你是成年人.
已經四五十歲.
或者五十歲以上.
六十歲左右.
學跆拳道是一個很好的運動.
為什麼呢.
不是開玩笑的.
我不是什麼高手.
千萬不要誤會.
我不能打的.
只不過在一個小時裡面.
你會將你最需要的四項運動做齊.
就是伸展式的拉筋.
伸展式的那些.
你的卡地令你心肺.
你的重量抵抗.
因為你要打拳.

$^{361}$你要踢.
令你的肌肉發力.
還有你的平衡式.
因為你要踢.
你要踢.
你要轉身.
你要轉來轉去.
跳來跳去.
令到你是很平衡的.
你可以在一個小時裡面做齊四項.
所以那個運動是很好的.
不過我說的就是.
我不是什麼高手.
千萬不要誤會.
我們要記住.
智慧不代表.
是你全部都懂得那些東西.
智慧是更加深層一點.
那些東西都是好事.
但不要因為.
你想去拿一張好的證書.
就監去做那些東西.
不是說那些東西不做.
待會我會再詳細一點.
說中間的關係.
但不要是因為想贏在起跑線.
就監去做那些東西.
最後你只會令到他燒脾氣.
另一個極端.
這個近來比較少見.
但也有一些教會.
或者一些人裡面.
尤其是一些基要的教會.
是會聽到這個說法.
就是樣東西都不要.
這些是反智的小學.
最常引用的聖經.
就是歌羅西書中說.
用他的理學和虛空的妄言.
不照著基督.

$^{401}$照人間的遺傳.
世上的小學就將來老去.
認識了神.
就不用要世上小學的東西了.
於是我們這些東西就不用懂了.
我們只懂得勝經.
懂得教會教的東西就夠了.
不要學外面的東西.
什麼進化論.
科學.
哲學那些.
通通都不要學.
不要懂.
總之在家裡就可以了.
我真的有認識傳道人教子女.
不要去讀書.
不用去讀書.
不用讀厲害的.
總之你不要懂那些東西.
在環教會裡面.
現在比較少見.
但還有.
我知道你們教會不是這樣.
因為昨天我才聽Mark Sir的.
講方舟.
我確保我講的東西.
和他沒有衝突.
於是我就再聽他講方舟.
和另一個朋友.
郝牧師的訪問那裡講方舟的東西.
Mark Sir對這種反智的態度是很反感的.
所以我和他是完全同一頁面的.
但是我要再問一下.
什麼是反智的小學.
是不是不用學哲學.
不用學進化論.
不用學其他的科學.
數學那些.
就可以了.
那些就是反智的小學.

$^{441}$我們再看聖經講什麼是反智.
什麼是世上的小學.
就是你們謹守日子月份節期年份.
我為你們害怕.
唯恐我在你們身上枉費了功夫.
世上的小學是什麼.
就是那些服從不可拿.
不可嘗.
不可摸等這類的規條.
這些都是人分腐的教導.
世上的小學不是科學.
不是哲學.
不是其他的知識.
而是一些教條.
將教條代替了真理.
這個就是世上的小學.
我還記得一個好笑的經驗.
幾十年前.
現在可能沒有了.
幾十年前華人教會有一個傳統.
就是叫人不要看電影.
不要進戲院看電影.
近四十年前我在香港.
在一間教會裡面.
那個牧師就說.
不要去看電影.
OK.
那不看電影.
沒什麼大不了.
我幾十年沒有進戲院看電影.
除了去年進去看了《破地獄》.
但他說為什麼不要看電影.
他說因為看電影.
你怎麼知道誰坐過那張椅子.
《詩篇》第一篇說.
不可作洩密人的作為.
所以我就不看電影.
我立刻想站起來問牧師.
那你坐巴士坐不坐低.
你坐不坐飛機.

$^{481}$你坐不坐的士.
你怎麼知道那張椅子誰坐過.
不看電影不是什麼大事.
可以不看沒問題.
但是他把一些傳統當成真理.
又要找一個理由來justify自己這樣做.
於是就硬著頭皮找一個很無厘頭的理由出來.
說出來是笑的.
是令人笑的.
如果有不信的信徒.
他都在教會裡面說.
什麼你的牧師是這樣的.
這麼反智.
說的話這麼不合理.
這個是令基督蒙羞.
令人以為基督徒就是這樣的.
我們不要.
我們不要用這句聖經說.
世上的小學我們不要.
於是就反智.
什麼都不用懂.
智慧不是這樣的.
智慧是要我們明白什麼是真的.
什麼是假的.
好了.
說了智慧的功效.
有什麼做父母的不要做.
我們做父母的.
有什麼我們可以做.
令到我們的子女更加有智慧呢.
第一樣就是.
智慧是需要求的.
智慧不是聰明.
聰明可能是天生的.
有些人的智慧真的比較厲害.
有些人可能是智慧沒那麼厲害.
但是智慧不同於聰明.
智慧是需要求的.
《雅各書》說.
你們誰少智慧就要求神給你們.

$^{521}$也不要.
那個不斥責人的神.
主就必定賜給你.
神不會因為我們沒有智慧.
來小看我們.
我們是需要向神求智慧的.
不要以為我生出來.
我的智慧不是很高.
於是我就潑辮子了.
不是這樣的.
神是說你沒有智慧.
你可以求.
聰明或者是天生的.
但是智慧是可以求回來的.
我們不需要很厲害.
不需要科科考試100分.
不需要拿什麼學位.
或者什麼資格.
我們需要有智慧.
就算一個沒有讀過很多書的人.
只要有一些基本的知識.
他就可以很有智慧.
分辨到是非.
分辨到什麼是對什麼是錯.
所以我們第一件事.
開始我們想子女有智慧的話.
我們就要幫他們求智慧.
我們也要教導他們求智慧.
而不是求知識.
不是求財富.
所羅門求智慧這個故事.
大家都知道.
神說你幫我建了聖殿.
你想要什麼.
他就說我要智慧來管理神的人民.
神就給了他.
所以我們第一件事.
就是我們需要一個desire.
有一個心.
我要有智慧.

$^{561}$而不要說我潑辮子了.
我求求其其.
我過一世就算了.
這個不是神的心意.
第二件事.
我們要有智慧.
我們需要有一定的知識.
知識不是智慧.
但知識是智慧的foundation.
是它的根基.
聖經裡面也是這樣說的.
箴言那裡說.
智慧進入你的心.
但是呢.
你要以知識為美.
知識.
你要有知識.
然後才有智慧.
當然你不可能什麼都懂得.
但你需要懂得一些基本的東西.
Mark Sir為什麼一開始就知道.
找方舟這件事是不妥的呢.
因為他是考古的專家.
他有那方面的知識.
我對考古是一無所知的.
但我的訓練是另一方面的.
我那時候是律師和醫生.
作為一個科學訓練的人.
和一個看證據的人.
我就知道他講的東西不妥.
所以知識是可以幫助我們.
先有一個gut feeling.
這件事不是很妥.
然後我們會再找多一點資料.
來知道什麼是對什麼是錯.
讓我們有智慧的.
開始去問問題.
也都是.
耶利米蘇說.
我將知識和智慧來侮讓你們.

$^{601}$所以我們是需要有知識.
然後我們才有智慧.
二人是因知識得救.
你要執意聽從知識的言語.
要你們的愛心在知識和各樣見識上.
多而又多.
知識到底有什麼可以幫助我們的子女.
可以得到更多知識呢.
我覺得有幾樣東西我們需要知道.
第一.
知識是事實和理論的結合.
Facts and rules的combination.
我們不可以說.
我們不講道理不講理論.
只講知識只講事實.
只講facts是不可以的.
我們也不可以說.
我們不講facts只講理論.
那也是不可以的.
我們需要教他們講道理.
也要跟他們講事實.
只有兩樣東西加起來.
我們才可以增加知識.
這樣說可能很抽象.
但待會說一些例子.
大家會比較容易明白.
第二樣東西我們講到知識.
知識是必須要有一個很闊的foundation.
我們需要各方面都要懂一點.
你不可能全部都懂.
你一定不會又是神學家又是科學家.
有這樣的人.
但我們大部分都不是.
包括我自己做來都不是.
但我們可以每一樣東西都懂一點.
我們都需要懂一些基本的東西.
基本的知識我們需要有.
不要說我將外面的東西.
不合我口味的東西我就完全不聽.
不合我口味的東西都要聽.

$^{641}$都要學一下都要明白一點.
我很喜歡看武俠小說.
金庸寫武俠小說.
每個主角都很能打.
但有沒有留意到.
他的主角能打的那些.
每個都幾個師父.
郭靖學過全真派.
學過洪七公.
韓龍十八掌.
學過九陰真經.
金庸說他哪個男主角最能打.
最厲害的就是張無忌.
因為張無忌學過七傷拳.
又學過太極劍太極拳.
然後乾坤大萊爾.
然後又學聖火靈.
然後又是懂醫人又懂讀.
所以一個很闊的知識的基礎是必須的.
我不是說個個都像張無忌一樣.
每樣東西都那麼厲害.
但我們需要每樣東西都懂一點.
我認識有一兩個人.
他們讀某些科目.
我不想說哪一科.
但他們認為我讀的東西就天下無敵.
我不用懂你的東西.
不用懂你的醫科.
不用懂你的法律.
不用懂你的社會科學.
總之我只要懂得批評你就可以了.
但結果說出來的東西只有笑話.
因為他根本不知道.
其他人已經做了什麼.
有一個人跟我爭論管理方面的問題.
拿了一本書.
你看看那本書怎麼說.
你拿一本七十年代.
講日本管理的書.
你卻說現在的管理沒什麼意思.

$^{681}$七十年代日本的管理.
那時候他們造車.
打到美國一仆一鹿.
日本的車那時候很厲害.
日本的電子產品那時候很厲害.
但和現在有什麼關係.
不是以為你看了幾本書.
就可以分析所有問題.
你需要的是有一個很闊的基礎.
能夠明白什麼在發生.
各方面你都要知道.
你的子女剛才說過.
他們學很多東西.
不要監去學東西.
是為了拿一張證書.
但如果他對某一方面有興趣的話.
千萬不要不給他.
不行啊 你一定要讀書啊.
你明天考試了.
所以你不要花心機去做什麼.
這樣是不對的.
這樣是不行的.
你要給他一個很闊的奇妙.
他們要明白.
他們要有機會去接觸到不同的東西.
這樣才可以幫他增長自己的知識.
有了知識做他的基礎.
他才可以有智慧.
第三個關於知識.
我們怎樣可以得到知識的方法.
我們要明白.
知識不是死的.
是跟著時間跟著環境來變的.
我們需要有彈性.
這個frexibility.
這個open mind.
我們看到一些新的數據.
一些新的東西.
我們需要調整我們的思想.
明白什麼是會變的.

$^{721}$科學是會變的.
知識是會變的.
所以不要說.
科學家今天說什麼.
醫生今天說什麼.
為什麼明天又不同.
這個是必然的.
因為環境變了.
好像說我做醫生.
糖尿病應該怎樣處理.
現在處理糖尿病的方法.
跟我30年前讀醫的時候.
那個方法是很不同的.
因為我們有了新的藥.
我們有了新的數據.
環境又不同.
所以我們的東西又不同.
所以在知識方面.
我們是需要變通的.
當我們教子女的時候.
也不是一部通書看到老.
我們需要什麼時候.
跟著環境改變.
來幫助他們明白.
東西是會變的.
這是我們由知識到智慧.
我們需要注意的事.
第三個可以怎樣令到我們子女.
或者令到我們自己.
更加有智慧的.
就是我們需要請教其他人.
不要以為自己天下無敵.
因為每一個人.
三人行必有我師.
所以我們需要的.
是請教其他人.
慾望人在自己眼中看為正直.
智慧人肯聽人勸教.
你要聽勸教受訓會.
使你終究有智慧.

$^{761}$智慧是必須要像磨刀一樣.
你需要跟外面的人接觸.
你需要聽人說.
你才會更加有智慧.
而不是關上門.
你讀書.
或者是你自己在想.
你就以為很有智慧.
我們一定需要多接觸.
一定需要去看看.
其他人說什麼.
他們不一定對.
但我們需要去接觸其他人.
還有.
當我們去接觸其他人.
請留意他們做過什麼.
沒有好樹可以結壞果子.
也沒有壞樹可以結好果子.
你看到果子.
你就認出樹.
所以我們要學人的時候.
我們不是看一次.
說一件事.
而是我們看看哪個人.
值得我們信任.
我有很多朋友認識了十幾年.
認識了二三十年.
我覺得這些人一直看下去.
信任.
我認識很多牧者.
我有一個女兒.
也是嫁了傳道人.
所以我比較多接觸做牧者的人.
有些.
開始好像說話不錯.
但久了你會發覺.
這樣做事.
也有些牧者.
我認識了他們幾十年.
我是非常欽佩他們.

$^{801}$因為他們是忠心的事奉.
他們不是大教會.
或不是很有權勢.
但他們是忠心的.
他們也不是單一議程.
他們是說成本聖經的教導.
當我們認識人的時候.
我們要尋求智慧的時候.
我們也要看看.
我們接觸的人.
他們是否真的走出聖經的樣式.
是看他們的成本.
聖經是否走出來.
而不是看他們只是說一件事.
這就是我們.
如果我們有子女.
我們希望他有智慧.
或是我們希望自己有智慧.
我們有什麼可以做.
但有一件事.
最最最最最重要.
就是我們不要其身不正.
賢教不如新教.
所羅門寫了很多智慧的東西.
但他教兒子教成怎樣.
不要說他有多少個老婆.
可能那時候環境所逼.
但他怎樣教兒子.
他兒子做了皇帝.
那些人說你不要欺負我們.
我們做工很慘.
讓我們輕鬆一點.
我們就事奉你.
他兒子羅邦怎樣說.
我老爸叫你背重物.
我現在更加要你背重物.
他用一條鞭打你.
我用一條蠍子鞭打你.
所羅門是否一個好皇帝.
對他的兒子來說.

$^{841}$他兒子看到的就是.
他怎樣虐待他的臣民.
他的兒子.
結果就是.
老爸是這樣嗎.
我就跟.
於是後果是.
後果就是造反.
以色列分裂.
十個支派走到.
上去以色列.
後面剩下兩個支派跟著他.
這是神的恩典.
神說我不要大衛的燈熄了.
於是幫他留下兩個支派.
這是神的恩典.
但所羅門說就懂得說.
但他教兒子怎樣.
如果我們想我們下一代有智慧.
我們言教不如身教.
我們更加要做的是.
我們自己行事出來.
行事為人.
我們希望可以影響到他.
但是否這樣說.
我們就不用聽真言.
不用看真言.
不用跟真言說的話.
不是.
如果是的話.
剛才我說的那些就是廢話.
不用聽.
因為耶穌說.
民事和法理錯的人.
坐在摩西的位置.
他們吩咐你們的.
你們都要謹守遵行.
但不要效法他們.
因為他們得過講治.
能說不能行.

$^{881}$我們今天看聖經也是這樣.
我們怎樣看聖經.
這是一個很重要的神學問題.
聖經裡面說的話是一個紀錄.
是作者為了想我們明白一些東西而寫下的.
真言說的話我們要聽.
只不過所羅門教兒子的方法.
我們不需要跟著去做.
我們又可以將所羅門教兒子和上帝比較.
人間的父母.
個個都希望子女有智慧.
我們天上的爸爸也是.
但是天父不只是口說.
這就是基督教和很多其他宗教的分別.
很多其他的宗教就有一本天書.
有一本經.
回教是可蘭經.
佛教有不同的佛經.
都是說這是上天教我們怎樣做事的.
是說.
但是在基督教.
我們相信的就是道成肉身.
神道成肉身.
耶穌來到世上.
教我們怎樣做.
他真的做出來.
告訴我們怎樣做.
這就是天上的父親和其他宗教.
那種只是說的方法的分別.
我們相信神.
因為神來到這裡告訴我們.
我們可以怎樣生活.
怎樣有智慧.
當我們希望子女有智慧的時候.
也希望我們可以記住這一點.
我們希望我們走出一個有智慧的生命.
讓我們的子女可以跟著.
也可以尋求到智慧的方法.
我今天想跟大家分享的就是這麼多.
多謝大家.

\newpage



\section{使徒行傳}
\label{sec:MAWjFwIPHJE}
\textbf{神國的模樣 (使徒行傳 2) - 陳崇基牧師}
\newline
\newline
連結: \href{https://youtube.com/watch?v=MAWjFwIPHJE}{\texttt{https://youtube.com/watch?v=MAWjFwIPHJE}} ~~~~ 語音日期: 2025-05-28
\newline
\newline
\hyperref[sec:bewfsUoYndE]{\small{< < < PREV SERMON < < <}}
~
\hyperref[sec:index_chronic]{\small{[返順時目]}}
~
\hyperref[sec:index_scriptual]{\small{[返順卷目]}}
~
\hyperref[sec:34yXkIAS4Hk]{\small{> > > NEXT SERMON > > >}}
\newline
\newline
$^{1}$今天我們繼續看這個連題.
我們說由復活節開始.
我們會說跨越.
就是說自己的安殊地域裡面.
我們走出來.
嘗試去挑戰一下自己.
推動一下自己.
以致我們的屬靈生命能夠成長.
上個星期因為不在你們中間.
因為我去看我兒子的畢業.
看他畢業真的有一點百感交集.
也有一點憤怒.
因為我的錢包都沒有了.
他就走一走.
畢業他走了十秒鐘都沒有.
拍攝完就算了.
做三個月而已.
我叫他你醒醒定定.
好好地做.
畢業就好像一個跨越.
我們總是人生裡面有不同的障礙要跨過.
總是有不同的安心點我們要離開.
然後我們的生命才可以向前.
這個其實都是上帝給我們的旨意.
叫我們基督徒的生命能夠繼續向前.
所以在這一年裡面.
我們會逐一看《使徒行傳》.
上一次看到《使徒行傳》第一章的時候.
就是說耶路撒冷的信徒.
耶穌基督的門徒.
在耶穌升天之後.
他們就被吩咐在耶路撒冷.
等候聖靈的降臨.
記得嗎?.
然後第一章最後的時候.
他們就選了第十二個的門徒.
因為猶大沒有了.
然後選了誰?.
馬提亞.
這個代表他們怎樣認識自己.

$^{41}$這一幫人認為自己.
就是上帝所應許的神國的子民.
是有聖靈來驅動的那一幫人.
在《舊約二賽》書裡面已經一早說了.
舊約裡面有不同的應許已經說了.
所以他們的自我理解.
他們的明白.
對自己的認識.
就是說我們就是上帝的子民.
我們就是那一幫相信彌賽亞.
然後有聖靈降臨.
在我們身上工作的那一幫人.
這個就是教會.
這個就是我們對我們自己的自我理解.
這個不搞清楚的時候.
我們就不太清楚我們在做什麼.
其實上帝就是設立我們成為一個新的上帝的子民.
有上帝的聖靈在我們身上工作.
所以為什麼他們選了馬提亞.
然後十二支牌就齊了.
現在齊人歸位.
就等候聖靈降臨.
以致我們能夠去為上帝的福音來工作.
這個是我們教會應該有的自我理解.
是我們自己每一個人有的自我理解.
我相信路加去講初期教會的故事的時候.
他有很多故事要講.
但他就是偏選一些故事.
讓讀者讓我們每一個世代的教會.
知道我們到底在做什麼.
原來就是上帝的.
應許的聖靈和他的promises.
在我們中間來fulfill這幫人就是上帝的子民.
所以這次我們看史蘭傳第二章.
第二章是很出名的一段經文.
如果你是初初接觸聖經的話.
或者接觸聖經已經很久了.
你都需要去看一下這一段的聖經.
因為講到五旬節的時候.
聖靈降臨.

$^{81}$到底發生了什麼事.
他就講了一個這樣的故事.
為什麼這段聖經這麼出名呢.
因為五旬節這段經文.
史蘭傳第二章.
就等於是教會的生日.
教會的出世.
教會的出生.
Technically 專有詞來說.
就是教會在這個時候誕生的.
生日.
我們每一個人都有birthday.
birthday 去哪裡.
去小肥羊.
拿個free lunch.
free dinner.
教會的生日.
所以每一次想起五旬節.
Pentecost的時候.
他們就會想起.
這個就是教會的開始.
所以這個故事對我們有什麼啟發呢.
很簡單.
這段經文雖然很長.
47節.
應該夠時間.
47節.
很簡單.
分開三段.
首先頭13節.
講到聖靈降臨.
發生了什麼事.
然後中間那一段.
就是講到聖靈的應許.
這是彼得的講道.
彼得第一次的講道.
講到聖靈原來一早在舊約裡面promise了.
特別在以賽亞書.
上帝的國度降臨的時候.
就會帶著聖靈的能力來到.

$^{121}$然後最後就講到.
聖靈降臨之後.
他們接受了聖靈.
然後發生了什麼事.
有什麼果效等等.
很簡單的去outline出來.
首先我們看第一節.
五旬節那日悼念.
他們就是耶穌的門徒.
全都聚集在一起.
忽然有響聲從天上來.
好像一陣大風吹過.
充滿了你所坐的整座房間.
又有舌頭如火焰.
向他們顯現.
分開的樂在他們每個人的身上.
他們都被聖靈充滿.
就按照聖靈所賜的口才.
說起別國的話來.
這裡講到五旬節.
五旬節是舊約裡面三大節期的其中一個.
三大節期有哪幾個呢.
有御節.
有五旬節.
還有在秋天的時候.
駐彭節.
五旬節是什麼意思呢.
五旬.
Pentecost.
就是五十.
七七四十九日之後.
第五十日.
就由哪一天開始數.
耶穌復活節那天開始數.
記得復活節主人嗎.
記得嗎.
我們有兩幅名畫在這裡.
記得嗎.
在那一天開始數.
五十日.

$^{161}$就去到所謂的五旬節.
第五十日.
所以今年如果你這樣算的話.
是六月八日.
六月八日就來了.
是不是.
所以我們過了復活節.
就還沒到五旬節.
我們剛剛在中間.
這個Pentecost.
這個時間.
當時是農耕社會裡面的第一次收割.
第一次收割是大麥小麥.
大概這個時間就是大麥小麥.
Wheat and barley.
這段時間來收割.
收割了大麥小麥.
就會做麵粉和做啤酒.
他會釀酒等等.
所以這段時間也是一個First Harvest.
所以是一個農耕的收割節.
所以在五旬節發生聖靈降臨.
教會的誕生.
這個就不單單在農田的Harvest.
在屬靈上的收割.
所以第一群人信主.
然後建立了教會.
這裡說到門徒聚集的時候.
有響聲.
還有什麼.
火.
還有什麼.
水.
不是水.
火.
火和風.
火和風等等.
你想起這些畫像的時候.
想起什麼.
究竟裡面有什麼地方是有響聲.

$^{201}$大的響聲.
有火焰.
有風出現.
就是耶和華降臨的時候.
是不是.
特別第一次西乃山的時候.
所以這個五旬節.
他們感覺到領受了聖靈的時候.
有一些這樣的現象的時候.
他們就知道耶和華顯現.
並且耶和華願意住在人的中間.
耶和華降臨.
不只是小小的來到我們就走了.
耶和華降臨就是在回眸.
西乃山回眸.
然後在聖殿耶路撒冷.
當榮耀降臨的時候.
耶和華的大能來到的時候.
就是上帝願意住在我們中間.
願意住在那些被救贖的子民裡面.
願意住在那個群體裡面.
所以試圖感覺到這樣的現象的時候.
他就發現原來聖靈在末世的時候.
願意住在叛徒裡面.
願意住在我們這群人的身上.
所以試圖有著聖靈給的能力.
就說出別國的話.
就是說他們本身不會說的話.
他們現在說.
並且說來幹什麼.
是傳福音.
是傳天國的福音.
說天國已經臨到.
所以在那時候.
就有一個這樣的情況發生.
就是第五節.
當然五旬節的時候.
也有很多人來守節.
所以在耶路撒冷是很虛潭的.
很多其他的流散不同的世界.

$^{241}$地方的猶太人.
就會來耶路撒冷守節.
修國節,禹節,主膀節等等.
就會很虛潭的.
這個時間就是第五節所形容的.
他說有從天下各國而來的虔誠猶太人.
住在耶路撒冷.
或者暫住在耶路撒冷.
守節.
修國節.
朝聖.
聲音一響.
很多人就來聚集.
各人就因為聽見門徒用他們自己的香談來說話.
就甚納悶.
有些些詫異.
他們就問.
漢啊.
這些說話的人不都是加利利人嗎.
順帶一提.
你知道在以色列.
那些猶太人.
當時的猶太人.
他們北面和南面的口音是不同的.
等於台山話.
不同村又不同.
大家不太明白大家在說什麼.
所以北面有北面的香談.
或者北面的accent.
南面的又有南面的accent.
南面的當然說我們最正宗.
你們北方人是邊色.
所以他們就聽見一些很詫異的東西.
就是為什麼他們會用我們的話來說話.
明明這班是加利利人.
一聽見他們說話就知道.
他們是北面的加利利人.
然後第九節開始他們說什麼.
他說我們帕提亞人馬代人.
以蘭人和住在美索不達米亞的.

$^{281}$等等列出十幾個的不同地方.
他們來到耶路撒冷.
他們就聽見他們門徒用他們自己的香談來說話.
這些地方名我們不說那麼多.
這裡大概有十五個地方名.
由東面帕提亞就是Parthia.
Parthian.
由東面去到西面有羅馬.
南面是埃及.
北面就是本都加伯多加等等的地方.
總之有十五個不同的地方.
就來到聽見上帝的福音.
就是說上帝的國度要展開了.
現在要launch了.
現在要kick in.
我在這裡預備的時候.
羅馬帝國的文獻.
當時羅馬詩人和羅馬史學家.
他們說他們羅馬朝廷和政黨多麼的厲害的時候.
他們就是用這樣的方法說.
他們就列出不同的人.
西面東面南面北面.
有什麼人種都歸納到羅馬帝國.
這個說法有一點像當時的政治.
中央電視台說的話.
路加就說現在說的不是地上的國.
是天上的國度.
聖靈要建立神的國度.
要kick in神的國度.
要開展神的國度.
所以試圖用不同的別國的話去傳這個福音.
很多時候說到聖靈的時候.
我們特別傳統的教會.
我們總是有一點點不舒服.
聖靈降臨.
說別國的話.
說方言.
或者有不同的表現等等.
我們總是有一種的.
是不是的.

$^{321}$或者我們有一點點存疑.
但是你要知道在《士行傳》裡面.
不是每次聖靈降臨.
聖靈工作的時候都說方言的.
對嗎.
所以你不要說第二章是這樣發生的.
所以永遠永遠.
以後每個世代.
聖靈就一定會有這樣的表現.
說別國的話.
這個當然是神職.
是上帝所做的事.
但是聖靈的工作不一定是規限於某一種的表現.
有時候會是這樣的.
但是在新約裡面.
甚至在《士行傳》裡面.
說到聖靈的工作.
是說有醫治.
可以帶給我們生命裡面有平安.
給我們有聖靈的果子.
不同的品格等等.
這些都是聖靈的工作.
不需要只規限某一種.
聖靈工作還有一樣最重要的.
就是推動教會.
來講出神國的福音.
傳承神國的福音.
並且領人去認識上帝.
所以不要規限某一種.
很大的事情.
很大的表現.
才是聖靈的工作.
我剛才說的路加.
我想路加寫《士行傳》的時候.
在他手上應該有很多日記.
並且有些日記是說.
我們 保羅去傳福音.
然後我們發生了什麼事.
所以在他手上應該有很多.
日記 聽了很多不同的故事.

$^{361}$所以他就收集了這些故事.
然後按照程序的.
按照安排.
來講這些故事出來.
那個重點就是這樣.
講這些故事有時候不是叫你複製的.
不久之前我看了一篇文章.
那裡說什麼呢.
他說一個人如果最好去培養自己的品格.
有公義.
Justice.
OK 講原則.
講法紀 講法治.
有這樣的生命來推動公義.
有一個最好的方法是什麼呢.
我們很多時候想.
最好的方法當然是學一下有什麼法規.
有什麼原則等等.
Constitutional 有什麼可以講.
最好的方法是什麼呢.
就是看馬丁路德金的故事.
看他的自傳或者他的傳記.
看一個人的故事.
看一個人的生命的時候.
你就會被inspire.
你會被驅動和被啟發.
這樣學一個公義的人生.
才是最有效的.
因為我們每個人學東西都是怎樣學的.
不是一些抽象的理論.
一點一點地我們去跟的.
我們是copy人家.
或者我們被人的生命去inspire我們.
以致我們有力量有passion去做某一件事.
這個其實就是史唐五傳.
有的目的.
講出第一代的信徒.
他們怎樣克服很多困難.
跨越很多不同的困難.
他們去到張神國的道.

$^{401}$去傳開去直到地極.
這個就是那個故事.
所以不要說這個故事是這樣.
所以以後我們的經歷都一定要照板煮碗這樣做.
inspire我們.
所以我們看看這個故事有什麼inspire我們.
五重節第一次.
聖靈的降臨.
是一個劃時代的神蹟.
開展一個新的國度.
所以我們留意到的時候.
很多時候有些學者就會說.
這裡聖靈降臨的時候.
不同的民族.
不同地方來的人.
他們聚集在一起成為一個新的子民.
能夠合一.
很多人就留意到.
這裡就是舊約裡面巴別塔的故事的相反.
你記得巴別塔的故事嗎.
創世記第十一章.
講到人要墮落的時候.
神自告.
他們建了一座很高的塔.
然後他們想合一.
他們想合一.
然後塔頂還可以通天.
講到這個故事的時候.
耶和華就說.
我不是很信得過你們.
不可以這樣.
所以他們就下來做什麼.
驅散流散人.
然後有不同的語言.
這個故事的背景現在去到五重節的時候.
就剛好一個reversal.
一個逆轉.
各人分散了不同的地方.
但是願意在聖靈的裡面.
在上帝的國裡面.

$^{441}$能夠合一.
走在一起.
能夠聽到上帝的福音.
這個就是聖靈降臨的一個描述.
到十四節到三十六節.
這裡就是彼得說的第一篇道.
眾人非常驚奇和困惑.
彼此說這是什麼意思呢.
說有人譏笑他們當然是喝了新酒.
記得我說大麥小麥.
剛剛收割.
釀啤酒都沒有時間.
怎麼會有這麼多啤酒給你喝.
所以不是喝了新酒.
不是喝醉了.
十四節說什麼.
彼得和十一個使徒站起來.
就高聲向眾人說.
猶太人和所有住在耶路撒冷的人.
這時你們要知道要執意聽我的話.
說什麼.
他說這些人並不是我們.
並不像你們所想像的喝醉了.
現在是幾點鐘.
早上九點.
早上九點怎麼會喝酒.
你都不是喝了酒才來.
你應該喝咖啡.
晚上才會.
所以他們不是喝醉了.
他說現在發生的事是什麼.
是先知約爾曾經說過的.
約爾曾經說過在舊約裡面.
在末後的日子.
上帝要將他的靈囂冠在凡有血氣的.
兒女要講預言.
少年人老年人.
僕人婢女等等都要講預言.
都要講出上帝的國度.
發生什麼事.

$^{481}$在天上會發生什麼事.
地上會發生什麼事.
有血有火有煙雲等等.
聖靈光臨的時候.
大響聲有風有火等等.
不同的表象.
就是哪時凡求告主明的.
都必得救.
所以約爾早就在舊約裡面說了.
有一個這樣的應許.
上帝的國度來到的時候.
在末世開始的時候.
就是聖靈會降臨.
然後男女老少.
你有沒有看到這裡有不同的群組.
最大的point最重要的point.
就是有不同的人.
有兒有女.
有男有女.
有少也有老.
也都有地位尊貴的.
地位卑微的.
所有的人都一式可以得到聖靈的工作.
凡血肉之軀.
就是說我們.
我們就是血肉之軀.
我們就是脆弱的人.
我們居然可以接受到聖靈的降臨.
然後聖靈可以在我們心靈裡面成為我們心靈的動力.
血肉之軀.
這兩個星期我真的感覺到我真的有血肉之軀.
因為我坐骨神經很痛.
我寫完其中一篇的paper之後.
哇,真的痛得不得了.
我以為我睡.
我通常都是.
我通常都是睡.
你讓我睡八個小時.
我已經可以reset.
但是我這兩個星期.

$^{521}$reset來reset去都reset不到.
好像你舊車的mileage來高.
你想reset它reset不到.
繼續開的時候.
我們就繼續多mileage.
我們就是血肉之軀.
但是這裡說什麼.
一個很寶貴的教導就是.
聖靈願意住在我們心靈裡面.
居然耶穌的靈願意住在我們心靈當中.
之後彼得就說了一篇圖.
這篇所謂五旬節的講道.
或者第一篇的講道.
大概說什麼呢.
這裡經文當然很長.
他說以色列人.
耶穌就是上帝以歧視神蹟向你們證明的那個人.
上帝按照他預定的旨意和預知.
被交與人.
然後你們就藉著不法之人.
將他釘在十字架上殺了.
但是上帝將死的痛苦解除了.
使他復活.
使基督能夠復活起來.
他原不能被死拘禁.
所以他就引用不同的詩篇.
你看看引用了什麼不同的詩篇.
彼得說詩篇大衛裡面說什麼.
他說我看見主上在我的面前.
他在我的右邊.
不至於動搖.
我心裡面就歡喜快樂.
然後我的肉身就安居在指望之中.
因你必不將我的靈魂雕在陰間.
不會使你的聖者見到身體上的扭歪.
這個詩篇十六篇.
我們再看下去.
他說諸位大兄.
先祖大衛已經死了.
也埋葬了.

$^{561}$你也知道數目在哪裡數.
因為大衛的墓在哪裡.
在耶路撒冷附近.
大衛既然先知.
他早就說了什麼.
他知道上帝曾經向他起誓.
在他的後裔裡面.
有一位會坐在寶座上.
是永遠永遠坐在寶座上.
他預先見到的時候.
大衛就論到基督的復活.
他說不必在陰間.
他的肉體也不見扭歪.
這個就是耶穌.
這個就是大衛的後裔.
大衛的子孫.
這位耶穌上帝已經使他復活.
能夠成為我們這件事的見證人.
基督現在被高舉在上帝的右邊.
從父那裡領受了聖靈.
現在就將聖靈降臨囂冠下來.
大衛本身已經死了.
他沒有復活升天.
但是他指著自己說什麼.
主對我主說.
你坐在我的右邊.
讓我使你的仇敵坐你的腳椅.
詩篇110篇.
故此以色列全家應當知道.
你們釘在十字架的耶穌.
就是真正的大衛王.
就是真正的大衛君王.
上帝所認許和會設立的君主.
這個就是 Messiah.
就是基督.
為什麼會引用詩篇這麼奇怪.
因為我們看詩篇的時候.
就看到詩篇是安慰我們.
因為我們大國的時候.
都會引用這些詩篇.

$^{601}$Be still 你要安靜.
你要知道我是神.
或者詩篇23篇說什麼.
我是我牧者.
詩篇第一篇說什麼.
我們在在溪水旁.
我們的生命就會怎樣.
通常都是一種安慰的角度去看.
為什麼他這裡說.
好像大衛寫詩篇的時候.
好像寫預言一樣.
這裡有一個這樣的情況.
我們要知道.
在舊約裡面.
上帝曾經應許過大衛.
他說你的國位.
你的座位.
是一定會有人繼續坐下去的.
你的國位一定要到永遠.
一定會建立到永遠.
但是你知道歷史是怎麼發生的嗎.
大衛有設立他的國度.
然後有所羅門.
然後分裂成為南北國.
然後一夾被摧毀.
所以猶太人看回詩篇的時候.
就會發現.
大衛所寫的詩篇.
說到他是一個王.
上帝怎麼保守他.
或者應許了他什麼的時候.
上帝說過的話.
應該不會落空的.
但現在我們搞到什麼田地.
我們現在是亡國.
我們是流亡.
所以一方面我們流亡.
一方面上帝曾經說過.
對大衛應許過的東西.
國度永遠的那些東西.

$^{641}$到底兩邊的將來應該怎麼搞呢.
所以到兩約之間.
舊約完了.
還沒有到新約的時候.
神的子民他們看回大衛的詩篇的時候.
明明是一個王所寫的東西.
上帝對他有應許.
但現在搞到我們亡國.
大衛的寶座是空空無人的.
他們就會投射出一個.
將來要來的大衛.
他會再唱這些詩篇.
會寫這些詩篇的應許.
能夠實現的.
你明白我的意思嗎.
就好像你失戀那樣.
那你就唱那些情歌.
可能唱了一首比較正面一點的.
將來會有誰誰來.
你就會投射出一個這樣的人.
這個就是他們看大衛詩篇的時候.
他就投射出一個盼望.
因為上帝的說話永遠不會落空.
但現在我們搞到這樣的田地的時候.
我們會不會實現呢.
他們說會.
什麼時候呢.
將來一定會實現.
所以到這裡的時候.
十六篇110等等不同的.
所謂君王的詩篇.
這些神的子民就會推斷到.
將來會有一位真正的大衛.
是永遠坐在國位上的.
他的子孫是坐在皇位上.
是永遠掌權的.
當我們見到耶穌從死裡復活.
他也是大衛之子.
他又坐在上帝的右邊.
你就說.

$^{681}$這個profile就是fit所有這樣的詩篇.
和所有舊約的寓言.
所以給了整篇的講道.
就是上帝曾經說過的.
現在我們親眼見到.
是實現了.
這一篇就是第一篇的講道.
第一篇講道之後.
當然就是有反應.
37節.
眾人聽見這篇道的時候就覺得擇心.
擇心.
有沒有想過什麼叫擇心.
覺得好像刀吉一樣.
為什麼.
因為我們見到耶穌.
這幫猶太人見到耶穌.
被釘十字架.
卻是在上帝的大能裡面.
能夠復活.
現在升在天上.
他們覺得很不舒服.
他們就對彼得和其餘的使徒說.
他們應該做什麼.
彼得對他們說.
你們各人應當悔改.
奉耶穌基督的名受洗.
洗你們的罪得以赦免.
你們就可以領受所應許的.
所賜下的聖靈.
這應許是給你們和你們的兒女.
一切遠方的人.
也給我們主所呼召來的人.
彼得還用很多其他的說話來作見證.
然後勸勉他們.
你們一定要信主.
一定要悔改.
然後歸向上帝自己的國度.
這是第一次的講道.
還可以算是第一次聖經的報道.

$^{721}$記得我們開始這個聯題.
跨越的時候.
我們說要做什麼.
我們要做見證人.
我們是witnesses.
我們是福音的witnesses.
你們有沒有想過.
這裡第一次的講道.
是一種persuasiveness.
一種說服的動態.
傳完一個福音的時候.
講完一個福音的時候.
應該是有一種persuasiveness.
有一種推動的.
不是說講完純粹.
是一個資料的交代.
每一次提到福音.
或者每一次向惠著福音來作見證的時候.
都應該有一個所謂說服力.
想去convince.
當然不是說逼人.
而是你的見證.
應該有一種convincing的manner.
不會讓對方聽完之後.
看完之後.
原來只是純粹一些冷知識.
永遠都是persuasiveness.
在《少林寺傳》裡面.
每一次的報道.
每一次講到上帝的國的時候.
都是一種推動性的.
去邀請對方.
起碼去考慮一下.
我這裡不是說.
很多人拿著這些經文就說.
我們每次講道的時候.
傳福音的時候.
我們就要這樣.
我們就要逼別人信主.
我們經常有這種感覺.

$^{761}$特別是這個時代.
I don't think that is what happening.
我想這裡說的是.
他的見證和他所傳的東西.
是有一種說服力.
在英國有一個出名的護教學家.
Alister McGrath.
Alister McGrath是一個很尖端的科學家.
也是一個很好的學者.
等等等等.
他曾經說過一句話.
我覺得很有意思.
你聽清楚了.
他說傳福音.
不是單單叫別人決字.
也不是單單叫別人舉手.
傳福音真正的目的是什麼.
傳福音真正的目的是.
你做完見證之後.
在對方留下一個好的正面的印象.
這個才叫做winning.
winning不是逼別人簽什麼.
我知道有時候有些弟兄姊妹都是這樣.
見人就要講福音.
是不是好呢.
In some sense, it's good.
我曾經傳過一個微信主的朋友.
他就說我不想見到那個人.
因為他每次見到我的時候都會講耶穌.
You see what I'm saying.
不是一定要逼別人硬銷什麼.
真正我們傳福音要做到的目的是什麼.
要有一個有說服力的印象.
一個好的印象.
希望我們以這個為我們的目的.
彼得就是在做這件事.
一個好的印象.
有些人問我們應該做什麼.
我們現在應該怎麼算.
於是因為這個報導.

$^{801}$41節.
領受彼得華的人就受了洗.
那一天門徒約添了三千人.
他們專注在使徒的教導和彼此團結.
抹餅和祈禱.
眾人就心存敬畏.
使徒們有很多的神蹟和奇事.
信的人都聚在一處.
販物公用.
有買料,填產,加業.
照每個人所需要的分給他們.
他們天天同心合意的去行徹在聖殿裡敬拜.
在家裡抹餅.
全心歡喜和坦誠的心用飯.
讚美神並且得全體百姓的喜愛.
主張得救的人數天天加給他們.
這裡當然有很多東西我們可以解答.
最主要的目的是什麼.
你記得我說路加社司令團的時候.
就是想啟發我們.
初期的教會是一個很少數的人.
是一些沒有力量的人.
沒有政治力量,沒有任何權柄.
但是他們的生命是怎麼樣的呢.
是一個有信心的生命.
專注在真道上面.
主要的目的是什麼.
他們不是懶散和堅持不懈的.
在真道上面他們學好學到去行道.
並且大家是同心的.
你有沒有留意整篇的經文.
二章一至四十七節.
經常有一句話或者有一種感覺.
就是They are all together.
他們走在一起聚在一起的時候.
被聖靈充滿不是個別單獨的.
而是整個群體聚在一起的時候.
傳道,報道,事奉都是一樣的.
信的人就聚在一起.
聚在一起.

$^{841}$是一群怎樣的人.
他們很渴求上帝的聖靈.
來到他們中間.
這個就是對我們的啟發.
到底我們是不是渴求.
我們很希望聖靈的能力臨到我們.
還是我們只在上帝面前一潭死水.
初期的教會.
就是一群等候,祈禱,願意聖靈降臨.
很想聖靈降臨.
你和我是不是這樣.
一個聖靈的生命是一個怎樣的生命.
他說得全體眾民的喜愛.
人們見到他們的時候覺得他們很有信心.
覺得他們很有生命.
很吸引他們.
其實你和我的生命應該都是這樣.
Yes,我們有不同的難處.
我們可能有不同的困難.
我們要克服.
但如果有聖靈在我們生命裡面的時候.
應該起碼基本上有點不同.
聖靈的生命有點像一塊鐵.
一塊鐵.
一塊鐵近一些磁石的時候.
那塊鐵就會有一些磁性.
越近那塊磁石多的時候.
那塊鐵就漸漸會有一些改變.
但你說大小,重量全部沒有改變.
形狀也沒有改變.
但因為它和磁石近的時候.
它自己的生命就開始有些特質.
和周圍的環境,周圍的關係就開始有改變.
以前只是一塊鐵放在那裡.
現在因為近磁石多了.
有其他的東西就會吸進去.
對嗎?.
聖靈的生命應該是這樣.
我們教會其實都應該是這樣.
到底我們是不是一群上帝的子民.

$^{881}$我們是拒人於門外.
冷淡.
還是我們是歡迎到別人進來.
是加插到別人進來.
是可以牽涉到其他人進來.
這個就是磁石的生命.
聖靈的生命.
你和我都可以有.
因為凡有血肉的人.
血肉之軀.
上帝居然在聖靈當中和我們一起去同行.
住在我們生命裡面.
在過去我曾經提醒過大家.
我們需要找一位未信主的朋友討過.
這次有新的功課給大家.
這次不是想逼大家做什麼.
這次你想想你信主的經歷裡面.
當然有起有跌.
但誰參與過呢?.
你有沒有想過?.
那個人曾經和你們全部人.
那個人曾經栽培過你.
或者和你同行過.
誰參與過你信主的故事?.
再想想他是一個怎樣的人.
他做了什麼.
他生命有什麼表現.
你得到什麼啟發.
我怎樣可以學他.
有什麼值得學他.
當然也可能有些東西不值得學.
但肯定有些東西值得學.
你再看一個故事.
看看他有什麼可以啟發你.
《士堂傳》五順之故事.
和在你生命裡面曾經參與過的人.
有什麼值得去效法.
讓他去啟發你.
讓聖靈住在我們生命裡面.
這是我們今天結束的禱告.

$^{921}$聖靈我們實在太無.
太累.
太獨行.
我們忘記了你呼召我們.
活在一個新的群體裡面.
求你聖靈再次降臨.
讓我們去渴慕你的真理.
去持守我們的團結.
痛心的去祈禱.
醫治我們的冷淡.
更新我們的關係.
癡開能力去見證我們的主基督.
可以痛心低頭去祈禱.
聖靈我們求你這一刻.
你來到我們生命裡面.
在我們生命裡面可能要打很多場仗.
令我們很累.
可以遇見很多不同的問題.
令我們很擔憂的東西.
令我們缺少了很多.
聖靈應該有的給我們的工作.
我們起了我們的和平.
我們的能力.
很多時候我們忘記了.
主,我們求你赦免我們.
今天求你進入我們生命裡面.
主,我們願意再次邀請你.
進入我們的心靈當中.
做我們的主.
帶給我們醫治.
帶給我們安慰.
賜力量給我們.
讓我們有能力為你作見證.
是一個軟弱,有罪,血肉之軀.
但在基督的救恩裡面.
你的聖靈願意與我們同行.
住在我們的生命當中.
主,我們感恩去讚美你.
主,拆敗我們出去.
讓我們有聖靈的生命.

$^{961}$聖靈的能力.
去見證你.
祈禱奉主名.
下星期見大家.
一起事奉.
\newpage



\section{3:1}
\label{sec:34yXkIAS4Hk}
\textbf{對上帝喜悅的祈求 (列王記上 3\_1–15)  - 江芷敏傳道}
\newline
\newline
連結: \href{https://youtube.com/watch?v=34yXkIAS4Hk}{\texttt{https://youtube.com/watch?v=34yXkIAS4Hk}} ~~~~ 語音日期: 2025-06-05
\newline
\newline
\hyperref[sec:MAWjFwIPHJE]{\small{< < < PREV SERMON < < <}}
~
\hyperref[sec:index_chronic]{\small{[返順時目]}}
~
\hyperref[sec:index_scriptual]{\small{[返順卷目]}}
~
\hyperref[sec:EgZkScLaImE]{\small{> > > NEXT SERMON > > >}}
\newline
\newline
$^{1}$很特別.
Ken上任的日子是我講到的.
不是他說到的.
謝謝Mark Sir的邀請.
我是Alice,Ken的太太.
很感恩在這裡有機會跟大家分享上帝的話語.
昨天是端午節.
不知道大家有沒有吃粽子呢.
沒有嗎?.
希望有機會跟大家分享.
我希望下個星期可以包粽子.
不過昨天也很開心.
去Family Fellowship.
有機會吃到粽子.
也收到姐妹送的粽子.
所以很開心.
不知道今時今日.
還有沒有人提起.
為什麼要過端午節呢.
是屈原.
還有沒有人再提起屈原的故事呢.
其實我看回他的才幹和愛國的情懷.
的確是很值得學習.
而且他是一個偉人.
是一個屬靈的偉人.
如果他信主的話.
如果我們今天再去看.
好像一個屈原的偉人.
有什麼代表呢.
我在今天這一段經文找到所羅門.
好像是一個以色列國的偉人的代表.
我們今天會看聶王記上三章的經文.
不過也要告訴大家.
原來有一段經文和他一樣.
同樣記載著所羅門去求智慧的.
是在《歷代志》下一章.
其實聶王記和《歷代志》的作者.
他們雖然是在記載同一件事.
但是他們是有不同的向度.
他們兩本書都在說不同的重點.

$^{41}$我想起好像一對夫婦一樣.
其實他們是一起相識相愛結婚.
但是有很多時候你問他們的愛情故事.
他們是在說不同版本的愛情故事給你聽.
就好像這個聶王記和《歷代志》.
記載著同一件事.
不過他們有不同的角度.
今天我們是會說聶王記的版本.
聶王記是記載以色列的興衰.
不過他不只是在說歷史.
他反而是用以色列的王和先知.
他們的信仰歷程.
去告訴我們知道.
為何以色列最後會四分五裂.
原來以色列的王的信仰狀況.
是很影響當時以色列國的興衰.
所以今天和大家一起看看這段經文.
聶王記第一個很詳細描寫的王.
所羅門他的祈求.
為何他祈願這個禱.
和上帝對完話之後.
上帝又這麼喜悅他.
甚至帶領以色列走向高峰.
成為萬世景仰的王朝.
不過小心.
所羅門也是一個很重要的因素.
令以色列國走向毀滅.
聶王記上三章一至十五節.
如果你有聖經你可以打開.
或者可以跟著PowerPoint的經文.
都可以的.
我習慣了用和修本.
聶王記上三章第一節.
我們通常說所羅門求智慧.
我們單刀直入下面的經文.
我們很多時候忽略了一至三節.
或者四,五節.
我們今天多放一點時間去看一看.
第一節.
所羅門和埃及法老王結親.

$^{81}$娶了法老的女兒.
接他進入大衛城.
直等到建完了自己的公和耶和華的殿.
以及耶路撒冷周圍的城牆.
他這裡要提到.
所羅門和埃及法老結親.
娶了他的女兒.
很漂亮的女兒.
很明顯這是一個政治的婚姻.
因為當時埃及是一個很強盛的國家.
所以跟他聯親.
所羅門就可以建立一個.
非常強大的軍事聯盟.
不過作者要記載這件事.
其實非常諷刺.
因為第三節更加提到.
所羅門是愛耶和華.
遵行他父親大衛的律例.
只是他仍然在丘壇獻祭肖像.
喂.
所羅門愛耶和華.
遵行一切的律法.
為什麼他要娶愛邦的女子.
《申命記》一早已經提及.
上帝曉遇這群以色列民.
是不可以娶愛邦的女子.
對不對.
大家看看這段經文.
他就說.
你不可以娶愛邦的女子.
因為以免你們的心轉向不跟從主.
去事奉其他的神.
激發耶和華的怒氣.
會將你們毀滅.
其實所羅門.
的確不是在遵行耶和華的律法.
再看第十一章.
真的令人大火.
原來所羅門不單只娶了一個公主.
所羅門.

$^{121}$是娶了七,八個公主.
這個AI圖是當事人的.
你拍照吧.
所羅門是娶了七,八個公主.
他跟多少個國家結盟.
對.
而且他不止七,八個老婆.
他是三百個妃嬪.
即是一千人.
他才是時間管理大師的始祖.
所羅門是上帝在地上的代理人.
他是否應該遵行耶和華的律法.
他明知上帝不喜歡他們娶愛邦的女子.
但他還是這樣做.
這個行為的背後是否反映.
所羅門是否真的信靠上帝.
還是他信靠不了上帝.
所以他才要背靠強國.
聶王記的描述很明顯.
不是唱好所羅門.
他是在勾畫一個.
信心不太足夠的所羅門.
我跟Ken是在2022年來芝加哥的.
他是想在三一神學院進修.
所以我們隻身來.
當時新車很短缺.
不知道大家有沒有在那個時候買過車.
因為晶片的緣故.
新車又很短缺.
其實二手市場也很缺貨.
所以我們一邊看車.
我們的預算又不夠.
我們一邊看.
看的都是20年的車.
看一些好一點的牌子.
即是Honda或是Toyota.
但都找不到.
我們每次出去.
都是要弟兄姊妹載我們去看車.
試車.

$^{161}$一次兩次都可以.
你不想整天麻煩弟兄姊妹.
但我們應該看超過七,八次.
心裡都覺得不好意思.
我們都看得很灰心.
加上Ken讀書的時間都很挫敗.
我們開始質疑上帝為甚麼帶我們來這裡.
上帝不是說會供應這一切嗎.
我們找車都這麼難.
在這裡沒有車很難生活.
不過有些弟兄姊妹真的沒有車都能生活.
但我們在上面很郊區的地方.
真的很需要一輛車.
我們那時候很有疑問.
還有信心很軟弱.
當時一個巴基斯坦的太太和我好朋友.
她說不如我們一起祈禱.
因為她也在找車.
她一家五口.
在那個時間她鼓勵我.
她說上帝一定會為我們預備最好的.
上帝一定會為我們預備最好的.
這句話聽起來很正常.
但我的心很不安.
我覺得這句話很難相信.
很難相信.
因為我們找車找了十個月.
事事都好像還是碰壁.
我們如是者地祈禱.
誰知一個星期之後.
她打電話給我.
她很開心地跟我分享.
她說教會有一對很年長的夫婦.
送了一輛七人車給她.
我真的很開心.
我很為她開心.
但下一秒我就很不開心.
為甚麼上帝只會供應給她.
上帝要為她預備最好的.
那我呢.

$^{201}$為甚麼覺得上帝好像不太疼我們.
還有我的不開心是另一方面.
很擔心我的信仰狀況.
我們不是來這裡讀神學的嗎.
我們都是傳道人.
為甚麼我們信心這麼軟弱.
在那個時間我覺得很失敗.
即使之前經歷了很多事.
很相信上帝會帶領我們.
我們才會來.
但在信心軟弱的時候.
我只能承認自己是一個人.
可能你會覺得買車很小事.
其實來這裡經歷更多的事.
的確是的.
在這裡生活是很不容易的.
不過當你重複做一件事.
嘗試一件事.
嘗試完又嘗試.
碰壁的時候.
個人就會很洩氣.
我以前也是一個很理想化的人.
不過我覺得上帝透過這幾年的經歷.
將我們拉出來落地一點.
聶王記的描述所羅門好像就是這樣.
他是一個屬靈的偉人.
他很理應相信上帝.
但他好像也不太相信靠上帝.
聖經說所羅門愛耶和華.
但事實上好像他很有保留.
第三節.
所羅門愛耶和華.
進行他父親大位的律例.
指使他還在丘壇獻祭肖像.
這裡說所羅門不單只娶外邦女子.
他在丘壇獻祭肖像.
好像違反了他愛耶和華的說法.
我要持平一點為所羅門解釋一下.
因為第二節說聖殿未建好.
所以百姓仍然在丘壇獻祭.

$^{241}$這個好像很合理.
而且《歷代志》一章三節補充說.
所羅門和回眾去欺騙丘壇.
是因為那裡有上帝的回眸.
那裡是耶和華獨人摩西在曠野做的回眸.
也在解釋第四節所羅門說去欺騙.
那裡獻祭的原因.
不過我想這是另一個作者很想埋下伏筆的地方.
就是丘壇可以獻祭給耶和華.
但我們再看後一點.
就知道丘壇也可以成為獻祭給其他神的地方.
原來燒香獻祭給神的地方.
是可以做其他的事.
我們看到之後就知道.
原來丘壇是可以獻祭給A神,B神,C神的.
是很不專一的.
但我們說這個聖殿是要專一獻給上帝的.
是要專一的.
但丘壇不是一個很專一的地方.
所以我也明白上帝為什麼經常用婚姻的關係.
形容我們和祂的關係.
因為愛上帝是要很專一的.
祂是一個聖潔的神.
祂容不下第三者.
這裡的描述是為了所羅門.
之後祂很不專一愛上帝留下伏筆.
不過我也要為所羅門說說.
所羅門愛上帝.
而聖經也有這樣說.
所羅門到基片在那裡獻祭.
因為基片有極大的丘壇.
所羅門在那裡獻了一千祭身為繁祭.
基片距離耶路撒冷大概有4.2公里.
昨天我們和家庭同事去了Starf Rock去行山.
我們好像走了五個多公里.
看回Apple Watch.
我們也走了三個多小時.
的確回到家也覺得累.
今早起床是有點掙扎的.
基片在丘壇.

$^{281}$即是說在山上.
所羅門還要帶著一千個祭身.
昨天我們帶著多少個小朋友.
我們帶著十幾個小朋友.
也覺得要花點氣力.
所羅門要帶一千個祭身.
其實所羅門也挺有誠意去獻祭給上帝.
我們也不能否定他真的愛上帝.
所以耶和華主動來找所羅門.
在第五節他這樣說.
在羈片耶和華夜間在夢中向所羅門顯現.
神說:你願我賜你什麼.
你可以求.
上帝主動去找所羅門.
所羅門這樣說.
你曾向你僕人我父親大衛大師辭外.
因為他用忠信公義正直的心在你面前.
所羅門搬了他父親大衛出來.
他告訴上帝知道.
這是你所愛的大衛.
是因為他很正直,公義,忠信.
所以你很看顧他.
繼續去確立自己成為王的認受性.
他又再說.
你又為他存留大慈愛.
賜他一個兒子.
即是他自己.
坐在王位上如同今日.
其實你可以看看.
因為《列王記》是一個很好看的故事.
《列王記》上一章,二章.
是說所羅門怎樣為他的王位.
爭取這個王位.
花了不少功夫.
他殺了他的哥哥.
他怎樣….
他媽媽拔士巴怎樣幫他用妓女.
去取回王位.
其實這是描述一個有血有肉的所羅門王.
我們再看下去,第七節.

$^{321}$現在耶和華我的神.
你使僕人接續我父親大衛作王.
但我是幼小的孩子.
不知應當怎樣出入.
所羅門做王的時候的確很小.
他只有二十歲.
二十歲,所以血氣方剛.
娶了這麼多個老婆.
不知道應該怎樣出入.
是為甚麼呢.
其實《民數記》也有說過類似的文字.
其實是說他不懂得怎樣管理這個國家.
他不懂得怎樣在這個國家的人民面前出入.
引導他們,就好像牧羊人一樣.
所以第八節他繼續說.
僕人住在你揀選的百姓中.
這百姓之多,多得不可點不可選.
這個以色列國有很多人.
所以所羅門求求一件事.
我們小心,我們經常說這段經文是所羅門求智慧.
看清楚了,第九節究竟他求的是甚麼.
我們一起讀,預備起.
所以求你賜僕人善於了解的心.
可以判斷你的百姓辨別是非不然.
誰能判斷你這麼多的百姓呢.
這個是所羅門的祈求.
他在祈求甚麼.
他是求一個善於了解的心.
究竟善於了解的心是甚麼呢.
我看過英文,它是說Understanding Mind.
其實看完也不太清楚.
我再看看原本的文字究竟是甚麼意思.
原來他在求一個可以聽,可以服從的心.
我們經常覺得他是在求智慧.
其實他是在求一個可以聽,可以服從的心.
在舊約裡,聽和服從的確是來自同一個字.
而所羅門求耶和華給他這個心去判斷是非.
即是說要在以色列文裡去分辨善和惡.
去判斷他們.
這是王一個很重要的身分角色.

$^{361}$在創世記裡,我們知道只有上帝能夠分辨善和惡.
而王就是在這個世界上,即是上帝的代理人.
他要教導以色列文去遵行上帝的律例.
去判斷這班以色列文.
因此王也需要有一個能力去分辨善和惡.
而所羅門和上帝的對話就反映著.
我不行,上帝.
我需要這個能力去帶領這班以色列人.
就是這個可以聽和服從的心.
因為這個聽和服從的心.
我才知道你是怎樣帶領這班以色列文.
其實王在地上的角色只是代表上帝.
他要帶領這班以色列人.
去知道怎樣去跟隨上帝,遵行上帝的律法.
所羅門的祈求似乎是肯定著.
以色列,智利以色列是出於耶和華.
儘管所羅門已經展示了他不同的政治手腕.
他跟法老聯親,娶了很多公主.
他還消滅了一些影響他王位的人.
穩定自己的江山.
不過他知道自己的能力是不夠的.
而他的祈求也在說.
上帝,這份管理以色列的能力.
是出自於你的.
所以我要聽到你的聲音.
可以服從你的律例誡命.
我才知道,我才有這份能力去管理以色列.
上帝非常喜悅所羅門的祈求.
第十節.
所羅門因為這事就蒙主喜悅.
我們好像在這裡看到.
上帝和所羅門的關係.
其實他是很想每一個來尋求他的人.
上帝很想我們跟他有一份很好的關係.
所羅門沒有為自己去求長壽.
沒有為自己去求財富.
沒有求滅絕受敵的性命.
因為他知道自己的身份.
他很需要上帝.
所以他求一個可以明辨的心.

$^{401}$他知道一切的能力都是出自於上帝.
而他需要上帝.
這個表達就令上帝很開心.
所羅門很需要上帝.
而上帝就是一個這麼重視跟我們關係的神.
所以他很喜悅所羅門的祈求.
而耶和華更加說.
你沒有求的,我都賜給你.
就是付足尊榮.
使你一切的日子.
列王中沒有一個能夠給你.
所以之後很多不同的國家.
就算是埃及的王.
都會來看看所羅門究竟多有智慧.
他很有名.
當然我不是想說.
跟上帝關係好就可以換取什麼.
或者他是一個很計謀地去祈求.
以至上帝給他這些東西.
不過我想這份令上帝很喜悅的祈求.
是因為所羅門很重視跟上帝的關係.
我想好像有時候做父母一樣.
子女需要你的時候.
其實可能你會覺得有時候覺得麻煩.
還有很多解決很多問題.
但是當他需要你的時候.
其實你都會覺得很滿足.
上帝就好像爸爸一樣.
很喜悅我們跟他說我們很需要他.
很想聽到他的聲音.
很想知道怎樣去跟隨他.
小時候我都記得很喜歡問我爸爸不同的事情.
因為小時候沒有人工智能.
去哪裡坐巴士.
我都是問我爸爸.
因為他好像一個什麼都懂的人.
所以很崇拜他.
以前很多事情都會去找他.
我爸爸都覺得很開心.
因為我很需要他.

$^{441}$所羅門的祈求是令上帝很喜悅的.
其實不是因為他求什麼.
縱使他不是求智慧.
縱使他是求其他東西.
但是因為他在表達.
他很需要上帝跟上帝那份關係.
所以上帝很開心.
第十四節.
這個祈求是有一個條件的.
我們一起讀第十四節.
預備起.
你若遵行我的道.
謹守我的律例誡命.
正如你父親大位所行的.
我必使你長壽.
這裡說的條件是很重要的.
其實上帝不是一個很形式主義的上帝.
不是你做了ABC.
那好吧.
你就有這些東西.
其實遵行他的律例誡命.
是在說你重視和他的關係.
就好像我們平日星期日回來這裡.
不是因為你回來崇拜.
上帝賜你什麼.
讓你和他建立關係.
不是的.
是因為你做這個回來的決定的時候.
你可能面對著不同的掙扎.
今天想睡晚一點.
或者今天想做完這件生意.
其實有很多的掙扎.
但上帝看重的不是你為他做的那件事.
或者那個行為.
是他看重你的心.
在這些掙扎.
做這些決定.
背負這些掙扎的時候.
他看重的是你的心.
所以所羅門為什麼得上帝喜悅.

$^{481}$是因為他重視和上帝的關係.
第十五節我們繼續看下去.
最後一節.
所羅門醒了.
是個夢他就來到耶路撒冷.
站在耶和華的約櫃前.
獻梵祭和平安祭.
又為眾神瀑擺設宴席.
這個夢.
其實他不是醒了.
原來做夢.
我們平時當做夢.
沒有一回事是假的.
但所羅門很明顯知道上帝真的和他說話.
所以他來到耶路撒冷.
他在耶和華的約櫃面前.
去獻梵祭和平安祭.
他知道他想多謝上帝.
和他的機票有些不同.
這次他獻平安祭.
平安祭也可以說是一個.
Fellowship Offering.
它是一種和上帝溝通.
是一個代表團契的祭.
我們可以去獻平安祭.
而這個平安祭獻完之後.
其實獻完之後可以吃的.
我們的祭物是可以吃的.
在吃東西的時候就在建立關係.
很多團契都飲飲食食.
其實不只是喜歡吃東西這麼簡單.
是因為吃東西的時候建立關係.
在這個祭裡面.
我們是在強調和上帝的關係.
而上帝也是一個很重視關係的神.
他喜悅我們的心.
他正在著他.
今日弟兄姊妹.
不知道你的心是怎樣.
你今天帶了一個什麼心情來崇拜呢.

$^{521}$你和上帝的關係是怎樣的呢.
我們今天看完這段經文.
我想不是最重要教我們.
要回去求智慧.
求智慧其實沒有錯的.
這是現今世代的邪惡.
我們都很需要智慧地生活.
或者我們需要有智慧地做一些抉擇.
都是沒有錯的.
儘管向上帝求.
今天所羅門的祈求再一次提醒我們.
和上帝的關係是怎樣的.
我們向他求的時候.
我們在著什麼.
我們是不是在著和他的關係呢.
所羅門肯定這一切的東西.
或者他擁有的事.
或者他擁有的王位.
都是出自於上帝.
他也知道和上帝的關係是最重要的.
不過有時候我們的生命好像所羅門那樣.
很軟弱.
或者有盲點.
但我想在所羅門的身上.
我們知道上帝沒有因為他的弱點.
或者沒有因為他犯罪得罪神去親近我們.
所羅門縱然有保留.
他不是完完全全的信靠上帝.
但上帝仍然走入他的生命裡面.
今天也要走入我們的生命裡面一樣.
他不是看我們做什麼.
去決定和我們建立關係.
最後一個分享.
真的很深刻印象.
在芝加哥前年的時候.
參加了生命貴巷舉辦的青少年營會.
叫Grace Conference.
是給青少年的.
當時有人邀請我去做導師.
我也很想參與事奉.

$^{561}$是一個對我生命很重要的經歷.
當時一起事奉的還有很多高中生,大學生.
他們都是來做導師的.
帶組的.
每個小組大概五,六個人.
他們是一群很愛主的青少年.
很愛主.
他們為了將組員帶到上帝面前.
他們花了很多額外的時間.
去思考生命的問題.
而且我們每天都會很早起床事奉.
但他們還更早起床.
就是為了為這些組員祈禱.
希望他們能夠經歷神.
我也被他們的行動激勵和感動.
在一次的授課裡面.
我負責大祈禱.
完結之後.
有一個小導師.
他走過來跟我說.
Thank you for loving Jesus.
不知為何這句說話.
令我覺得很難接受.
我不是一個很習慣接受,讚美.
但我覺得這不單是在稱讚我.
好像有些事背後還要說多一點.
就是我覺得他在說.
多謝我去愛耶穌.
但我覺得他們比我更愛耶穌.
他走過來多謝我的時候.
我花了很多時間消化.
我覺得自己很不配這句說話.
來到這裡可能經歷不同的困難.
很多的挫敗.
都覺得自己很軟弱.
從前說想盡力去愛神.
但原來發現自己這麼軟弱.
去愛神其實很難.
慢慢地我想開了.
我覺得上帝衡量我們是否愛祂.

$^{601}$不是在乎我們做了甚麼.
就像剛才這樣說.
你回來崇拜.
祂不單是欣賞你這個行為.
祂是欣賞你背後的心.
愛不愛主沒有人可以判斷你的.
你愛不愛神.
你多麼愛神都沒有人可以判斷你.
是在看我們背後的動機和心.
其實有時我面對小朋友.
其實小朋友很小.
他們很活潑.
其實我見到他們都很可愛.
但有時你說話的時候.
他們突然離開位置打瓜豆.
他們好像不太聽你的說話.
但在這句說話中.
我覺得上帝看我們.
好像我們看著這群小朋友.
有時聽不進去.
或者祂跟我們的提醒.
我們都明知故犯.
還要踩腳踏在這坑上.
但上帝仍然耐心地走近我們.
就好像祂仍然走進所羅門的夢裡.
跟祂對話一樣.
上帝期待我們跟祂說.
我們很需要你.
爸爸,我們很需要你.
愛上帝是一個跟祂關係的體驗.
上帝也很期待我們愛祂.
很重視跟祂的關係.
所羅門的生命的確有不好的地方.
但亦提醒我們.
祂跟上帝的關係.
建基於祂重視跟上帝的關係.
丁姐妹.
不知道今天你跟上帝的關係如何?.
或者你會否因為一些軟弱.
或者一些人的說話.

$^{641}$阻礙了你跟上帝的關係?.
讓我們一起去到上帝的面前.
跟祂說.
上帝,我們很需要你.
爸爸,我們很需要你.
讓我們一起祈禱吧.
總是看重我們.
無論我們是怎樣的光景.
我們有時也會得罪你.
我們屬靈生命很差.
信心很軟弱.
但上帝也愛我們.
繼續尋找我們.
扶起我們.
親近我們.
主,求你幫助我們.
好像所羅門一樣.
能夠很坦誠地跟你說.
我們需要你.
求你幫助我們.
跟你建立一個很親密的關係.
可以聽到你的心.
你想我們服從你的地方.
主,願你幫助我們每一個.
當我們想跟我們的靈社跨越.
建立關係之前.
我們先跟你建立一個親密的關係.
願主,你親自幫助我們.
獻上禱告交託.
奉主名求.
阿們.
謝謝大家.
\newpage



\section{馬可福音}
\label{sec:EgZkScLaImE}
\textbf{新約中的情緒 (3): 哀傷 (馬可福音 14:32-36) - 何建邦傳道}
\newline
\newline
連結: \href{https://youtube.com/watch?v=EgZkScLaImE}{\texttt{https://youtube.com/watch?v=EgZkScLaImE}} ~~~~ 語音日期: 2025-06-11
\newline
\newline
\hyperref[sec:34yXkIAS4Hk]{\small{< < < PREV SERMON < < <}}
~
\hyperref[sec:index_chronic]{\small{[返順時目]}}
~
\hyperref[sec:index_scriptual]{\small{[返順卷目]}}
~
\hyperref[sec:lPvFkn_YDDw]{\small{> > > NEXT SERMON > > >}}
\newline
\newline
$^{1}$大家好,早安.
今天我會說的題目是.
「新約中的情緒」.
三,愛上.
我覺得做人要有始有終.
我洗濕了頭.
開始了「新約中的情緒」一.
要說完它.
其實我有點掙扎是否說這個題目.
因為我覺得自己不太能接受這個題目.
第一,我不是讀諮詢.
第二,我沒有讀過心理學.
以Bella來說.
這篇說得好像對得多.
但對我來說.
好像令情緒變得很低.
好像馬上變得很沒有深度.
但我覺得也沒辦法.
這篇說到幾個月前已經預備了.
因為洗濕了頭.
所以我覺得今天也要完成這個系列.
所以今天會跟大家說一下.
究竟在新約中是如何看待愛上.
跟之前兩次說到的一樣.
我們首先會看現代的定義.
如何看待愛上.
然後我們再比較一下.
新約時代的人如何看待愛上.
以及新約的作者如何看待愛上.
以及對我們有甚麼反思.
我想先問一個問題.
誰覺得自己從來沒試過愛上的舉手.
有人沒試過傷心嗎.
應該不會吧.
這個問題對我們每個人都會切.
我再問一個問題.
哪個人你覺得….
這裡有哪個人你覺得.
你不喜歡愛上的舉手.
你不喜歡愛上的.

$^{41}$應該有比較多吧.
開始有比較多人說.
你不喜歡愛上的舉手.
其實愛上對我們人類來說.
是一個很大的難題.
我們每個人都會經歷愛上.
但很多時候我們不是很喜歡愛上.
因為愛上是一樣.
我們都覺得是一件負面的事.
或者是一樣很難處理.
很難面對的事.
我們首先看看.
究竟現代的心理學是如何看待愛上.
究竟愛上是甚麼.
又是APA,美國心理學會.
它如何定義愛上.
愛上就是當人失去了一些很重要的人.
事或物的時候.
例如親人.
之後所經歷的那種痛苦,那種痛苦.
就是愛上.
這就是APA的定義.
我們又看每次說到都出場的人物.
Paul Edmund.
他認為….
我之前說過這個心理學家.
他是研究情緒的.
他定義愛上是七種基本情緒之一.
跟我之前說過的憤怒和恐懼.
都是一樣的,都是基本情緒.
是因為失去重要的人或事物而產生的.
導致我們愛上的原因.
因為個人和文化.
對於甚麼是失去的看法而有所不同.
雖然愛上是一種負面情緒.
但Paul Edmund認為它很有作用.
因為它發出了一種.
人需要接受幫助或安慰的訊息.
簡單來說是甚麼.
就是這樣.

$^{81}$假設我是一個人.
我經歷了一些很創傷的事情.
例如我失去一個親人.
但如果我不懂得愛上,會怎樣.
人們不知道.
原來我經歷了一件很慘的事.
他就不會幫我,對吧.
所以愛上在人類社會中.
是一個很重要的功用.
在社群中發揮了一種信號的功用.
當一個人愛上的時候.
其實他在告訴別人他受到傷害.
他在告訴身邊的人.
告訴社群他需要幫助.
所以如果一個人不懂得愛上.
在社群中會有一個危機.
就是他受傷了,他死了.
也沒有人會知道關心他.
因為他無法發出愛上的訊號.
所以這就是Paul Edmund所說的.
愛上在社群中的作用.
我們又回到古代.
耶穌的時代,新約的時代.
究竟人們怎樣看愛上.
我們回到一個很出名的哲學流派.
叫做斯多亞主義.
我上次也有說過.
在說恐懼時也有說過.
他們將情緒分為四大類.
他們認為在現在經歷好的事情時.
那種就是壓力,歡樂.
當你期望未來經歷好的事情時.
那種情緒就是渴望.
當你現在經歷壞事時.
那種感受就是愛上.
而當你預期未來經歷一些壞事.
我上次也說過.
那種情緒就是恐懼.
所以愛上對斯多亞主義的哲學家.
是怎樣的.

$^{121}$就是當人在經歷一些壞事時.
那一刻的感受就是愛上.
我上次也有提過.
對於這群斯多亞哲學家來說.
他們認為人生只有一件好事物.
就是美德.
才是人生最重要,值得追求.
唯一的好事物.
簡單來說是甚麼.
你的朋友,親人,愛人,情人.
對了,為何是愛人和情人.
總之你的朋友,親人.
以及所有的東西都不重要.
你的財產也不重要.
失去也不重要.
因為最重要的是人需要有美德.
所以其他東西失去了.
其實也沒關係.
都是無關的.
簡單來說,推論出來是甚麼.
人是不應該因為失去任何東西而哀傷.
這件事對於我們現代人來說.
很難接受.
但那時是一個很流行的哲學流派.
斯多亞主義.
他們真的覺得.
人是不應該失去任何東西而哀傷.
所以哀傷對斯多亞哲學來說是甚麼.
哀傷是一件道德上不好的事.
你道德差才會哀傷.
為甚麼?因為如果你道德好.
就會覺得只有德行才是我應該追求.
所以即使我失去親人.
我也不應該哀傷,對吧.
但你因為失去親人而哀傷.
我即是在說甚麼.
你道德差.
他們真的這樣想.
這是跟道德有關.
哀傷是一種道德不好的表現.

$^{161}$為了讓大家更深理解.
究竟這班斯多亞哲學家.
是如何安慰人.
或是他們如何看哀傷.
我引了一些他們的文獻.
這是一位很有名的斯多亞哲學家.
耶穌時代的.
叫做Ipetitus,愛比克泰德.
對,很難讀.
愛比克泰德.
他是這麼說的.
他說永遠不要因為失去任何東西而說.
我失去了他.
而是要說我已經歸還了他.
如果你的兒子死了.
他已經歸還了.
如果你的太太死了.
他已經歸還了.
如果你的財產被沒收.
也是被歸還了.
所以你不要因為這些事情而不開心.
因為那些事情不是你的.
即使你太太死了也不要不開心.
因為她不是你的.
你只是還了她,對吧.
所以不用不開心.
這個還不夠厲害.
我們再看看他如何安慰人.
是這麼說的.
這也是Ipetitus所說的.
就是他教人如何面對哀傷.
從現在開始.
每當你對任何事情感到高興時.
記得想相反的方向,相反的畫面.
當你睡前親吻你的兒子時.
就要低聲跟他說.
「你明天就會死了」.
有甚麼壞處呢.
或者你跟你的朋友說.
「你明天就會離開」.

$^{201}$或者「我會離開」.
「我們以後都不會見到對方」.
他真的這麼說.
這是真的.
我是引用文獻.
當時的斯多哈哲學家真的這麼說.
他在預備人的心.
就是你不要看你的朋友很重要.
你不要看你的兒子很重要.
你不時跟他們說或自己說.
他們是不重要的.
他們真的不在的時候.
你就不會哀傷,對吧.
好嗎.
但人始終會哀傷.
當你面對一些哀傷的人.
應該如何處理呢.
Ipetitus就這麼說.
當你看到有人傷心,哭泣時.
說是因為失去孩子或財產.
你要留心.
千萬不要以為這個人真的在受苦.
他不是在受苦,那些事情是假的.
你要保持警惕.
要警惕,哀傷是很危險的.
要警惕.
因為這個人真實地不會感到痛苦.
因為那些痛苦其實是幻覺,是假的.
這些痛苦源於一種錯誤的判斷.
甚麼錯誤的判斷.
就是他覺得那些事情很重要.
但那些事情其實不重要.
所以才會傷心,對吧.
所以這種痛苦是假的.
是錯誤的判斷.
所以你要警惕.
不過你也要表示同情.
因為他在哭,你也要表示同情.
如果有機會,你就跟他一起哀悼.
但要小心,你不要真的傷心.

$^{241}$你內心不要真的哀悼.
因為那些是假的.
簡單來說是甚麼.
斯多克主義哲學家認為.
當你看到一個人哭泣時.
你要很留心.
不要讓他影響到你也哭泣.
不要讓他影響到你也覺得傷心.
因為那些傷心是假的.
不過你要假裝哭泣,假裝傷心.
你要跟他一起哀悼.
但不要讓他影響到你的心,明白嗎.
因為哀傷是不好的.
這真是他們的想法.
這就是斯多克主義哲學.
試想想,如果放在我們現代.
有人在安息禮拜.
也挺嚴重的.
如果我們用Ebictus的方法.
去到安息禮拜.
就跟一個新人,剛過世的朋友說.
「你不用傷心」.
「他不是你,對吧」.
「只不過是歸還而已」.
「不過我也會…」.
「因為你哭泣,所以我也會一起哭」.
「不過其實我不會傷心」.
「因為這些傷心是假的」.
你要知道你的傷心是假的.
是你錯誤判斷.
你以為你失去的親人很重要.
其實並不重要.
如果是這樣,是否很糟糕.
是很糟糕的.
我們再看看.
另一個流派.
叫做中期柏拉圖主義.
Middle Patriotism.
我上次也有提過.
其中一個哲學家叫做.

$^{281}$Filo,很有名的,菲羅.
他對傷心有這種看法.
Filo是一個猶太人.
所以他會引用舊有的經文.
他就說亞伯拉罕如何面對傷心.
他就這樣說.
當亞伯拉罕失去一生伴侶時.
他仍然像摔角手般.
戰勝了襲擊他.
並威脅他,壓倒他心靈的愛上.
他一生都將理性視為特別的顧問.
現在特別決心服從理性.
理性的勸告是這樣的.
不要過度哀傷.
就像經歷了一場前所未有的災難.
但也不要表現出冷漠.
好像從來沒有發生過令你哀傷的事情.
而是要選擇中間的道路.
就是你不要過度哀傷.
亦不要過度冷漠.
而不是任何一個極端.
要努力表達適度的哀傷.
於是他就說.
亞伯拉罕在妻子的屍體面前哭了一會.
很快就站起來了.
因為他認為再哭下去,再傷心下去.
會跟他的智慧有矛盾.
因為他曾經被教導說.
不應將死亡視為靈魂的滅絕.
而應將死亡視為靈魂與肉體分離.
回到原先的地方,即上帝那裡.
好,我想問大家.
Filo,即中期柏拉圖主義哲學家.
對於哀傷的看法.
跟斯多噶主義有甚麼不同.
是有點不同的.
雖然兩者都覺得哀傷是不好的.
但不同之處是.
斯多噶主義覺得人絕對不可以哀傷.
但Filo,即中期柏拉圖主義.

$^{321}$認為人是可以哀傷的.
不過哀傷要哀傷得適度.
無論任何事情都不可以過分.
你不可以完全不哀傷.
但亦不可以過分哀傷.
應該像亞伯拉罕那樣.
哭著就算了.
哭著就站起來.
哀傷過後,然後就控制自己.
不要再哀傷了,要理性一點.
這就是中期柏拉圖主義.
Filo的想法.
我們如果應用Filo的想法.
去安息禮拜,都會出事.
為甚麼?因為Filo看哀傷是甚麼.
哀傷是一個敵人.
是一個跟亞伯拉罕打摔角的敵人.
亞伯拉罕要努力打低哀傷.
你想像一下,如果我們在想.
我們在幻想.
我們在概念哀傷是一個敵人的時候.
安息禮拜會怎樣安慰家人.
就是不要緊.
你快點站起來打贏你的哀傷.
你戰勝你的哀傷.
你要戰勝你的傷心.
如果我們應用中期柏拉圖主義.
去現代安息禮拜,都是很糟糕的.
哀傷對於Filo來說是一個敵人.
是需要戰勝的.
我們又看看.
我們總結一點.
在希臘羅馬的時代,即耶穌的時代.
基本上,看哀傷都是不好的.
人最好是沒有哀傷.
應該除去哀傷.
斯多噶主義覺得.
人要完全杜絕哀傷.
而中期柏拉圖主義覺得.
哀傷是不好的,你要控制哀傷.

$^{361}$不要讓哀傷戰勝你.
這就是當時希臘羅馬哲學的普遍看法.
沒有人覺得應該哀傷.
每個人都覺得不應該哀傷.
而且要控制.
我們看看新約究竟如何看哀傷.
保羅在《撒羅尼加前書》四章的經文.
整個新約中比較詳細地說.
如何看哀傷或如何面對哀傷的一段.
比較詳細的經文.
保羅在那裡這樣說.
弟兄們,論到已碎了的人.
即死了的人.
我不希望你們不知道.
免得你們哀傷.
就像其他沒有盼望的人一樣.
事實上,我們既然信耶穌死了又復活了.
就應該相信神將來也要把那些.
在耶穌裡碎了的人.
與耶穌一同帶來.
我們照著主的話告訴你們.
我們這些活著存留到主來臨的人.
絕不會在那些碎了的人之前.
因為主必親自從天上降臨.
你們就要用這些說話彼此勸勉.
簡單來說是甚麼.
我們看下去.
保羅對於哀傷的看法.
也有點像希臘羅馬的哲學.
很嚴重吧.
保羅是怎樣說的.
就是我不希望你們哀傷,為甚麼.
因為其實那些人並不是真的滅亡了.
他們只是死了,但將來會復活.
所以你們不用像那些沒有復活的盼望.
的世人一樣哀傷.
如果我們驟眼看下去.
保羅是這樣勸誡.
貼上加錢書的讀者,收信的人.
是否有點嚴重.

$^{401}$這麼像斯多亞哲學.
如果我們按照保羅的說法去做.
去安息禮拜.
又會出事,會怎樣.
如果你有弟兄姊妹.
即是家人剛離世,去安息禮拜.
你怎樣安慰他們.
按照保羅的說話.
你不要像那些不相信主的人一樣傷心.
因為他們現在死了,但復活了.
他們現在跟主在一起.
所以你不用不開心,對吧.
這也是很嚴重的.
我們不會這樣安慰別人.
但我想說真的有人這樣安慰別人.
我們現在這樣看,覺得很誇張.
但我想說在教會.
真的有很多牧者.
按照保羅這裡的教導去安慰別人.
有很多這樣的基督徒.
當別人面對哀傷時,便跟他說.
「不用不開心,你開心吧」.
「因為將來你會再見到他」.
其實這做法.
就像我們剛才取笑的斯多亞主義.
和柏拉圖主義一樣.
其實沒有真正面對哀傷.
那就糟糕了,聖經也是這樣.
聖經也是這樣教,是否很糟糕.
我們一定要問一個問題.
聖經的哀傷.
是否就像剛才的希臘羅馬哲學一樣.
我們以後還是叫別人.
「不用不開心,開心便可以」.
那就糟糕了.
我們要回到剛才貼羅尼加前書的背景.
其實保羅在這裡不是叫別人不要傷心.
他在這裡其實要澄清一個很重要的神學問題.
我們要回到背景.
那時那些貼羅尼加的信徒.

$^{441}$面對一個神學難題.
就是他們有些親人過世了.
信了主,過世了.
但他們不知道如何面對.
因為在那時的早期教會.
還相信耶穌很快便會回來.
如果耶穌很快便回來.
而我的家人已經過世了.
他們如何可以再跟耶穌同在.
他們正面對一個這樣的神學難題.
所以保羅在這裡.
要教這些初信的貼羅尼加信徒.
一個神學重點.
就是你不要以為那些人真的死了.
其實他們將來會復活.
所以你們有一個復活的盼望.
所以保羅在這裡不是叫那些人不要傷心.
在這裡保羅有一個神學問題要處理.
就是你有復活的盼望.
而不是像世人一樣沒有復活的盼望.
所以這裡不是你不要不開心.
而是你要搞清楚.
其實你不開心也好.
但最重要的是你要知道.
將來這些家人會復活.
會跟你們和耶穌同在.
這才是保羅正在面對.
正在處理真正的問題.
所以我們明白這個處境之後.
我們便可以回去了.
新約還有沒有其他經文.
如何去看哀傷.
和當時的哲學有甚麼不同呢.
首先第一件事.
在《馬學福音》我們之前也說過.
十四章.
耶穌在臨上十字架之前.
其實說過一段話.
在喀西瑪利園.
他和雅各,彼得雅各和約翰.

$^{481}$一些門徒去到喀西瑪利園祈禱的時候.
耶穌面對臨上十字架的時候.
他的禱告是這樣說的.
不是他的禱告.
他跟他的門徒這樣說.
他跟他的門徒說甚麼.
「我的心哀傷得要死」.
「你們留在這裡要警醒」.
其實耶穌沒有像當時的羅馬哲學一樣.
說其實死亡是不要緊的.
我是很喜樂的.
因為我回到父那裡.
耶穌不是這樣說的.
耶穌在喀西瑪利園面對將要死亡的時候.
他是坦誠地跟他的門徒說.
「我真的很傷心」.
傷心到怎樣.
「我傷心得快要死了」.
耶穌很坦誠地說出他的感受.
他是傷心得,哀傷得要死.
他不是像哲學家那樣教導門徒.
「其實你不要不開心」.
「我也沒有不開心」.
不是這樣,他是承認自己的哀傷.
我們再看看另一個很偉大的人物.
保羅.
其實保羅在書信裡.
也經常承認自己的哀傷.
例如在《菲勒比書》中有這樣說.
他有一個很重要的瞳孔.
叫做「爾巴忽提」.
當時爾巴忽提已經生病了.
甚至快要死亡.
但神就醫治他,亦連恤他.
保羅說神不單連恤他.
亦連恤他,免得他憂傷.
再加上憂傷.
其實保羅在說甚麼.
保羅很重視爾巴忽提這個瞳孔.
很重視這個人.

$^{521}$所以如果失去爾巴忽提這個人.
保羅會怎樣.
哀傷加上憂傷.
哀傷加上哀傷.
即是很,很,很,很哀傷.
但幸好神憐憫他.
亦憐憫我,即是憐憫保羅.
不用我那麼哀傷.
即是治好爾巴忽提.
所以保羅….
如果我們看到在「肥納比書」.
其實保羅不是像當時的.
希臘羅馬的哲學家一樣.
拒絕哀傷這件事.
保羅其實是承認哀傷.
他自己也說自己哀傷.
他沒有說爾巴忽提不重要.
我歸還它給上帝.
他覺得很重要,所以他會哀傷.
所以我們看到新約的經文.
跟當時的哲學不同的地方是甚麼.
第一個很重要的不同.
就是能夠承認我們正面對哀傷.
我們能夠承認.
我們正在失去一些重要的東西.
耶穌能夠承認.
原來失去生命對祂來說是重要的.
所以祂會哀傷.
保羅能夠承認.
原來失去一個很重要的瞳孔.
爾巴忽提對祂來說是重要的.
所以祂能夠承認自己的哀傷.
今天我們能否承認.
我們的哀傷呢.
其實承認哀傷是很重要的.
因為我們人生裡有很多.
我們很珍惜,很重視的人.
有很多人可能會在我們生命裡.
成為我們生命的一部分.
當我們失去這些人的時候.

$^{561}$其實我們的生命是失去了一部分.
如果我們不能夠承認哀傷的時候.
其實我們同時也是在拒絕.
原來我們的生命.
是失去了一些很重要的一部分.
承認哀傷是在幫助我們.
去接納,擁抱我們自己的生命.
原來的確是受傷.
我們的生命的確是失去了很重要的一部分.
承認哀傷是在說什麼呢.
如果哀傷是來自失去一些重要的東西.
人,事,物的時候.
承認哀傷是在告訴我們.
我們真的覺得我們失去了的人很重要.
這是很重要的.
原來哀傷是在說.
我們生命當中一些很重要的部分失去了.
我們能不能夠坦誠地去擁抱這個事實.
去擁抱我們失去的東西.
接著我們再看看.
剛才施多夏,哲學說.
如果你看到一個傷心的人.
你記得他怎樣嗎.
你不要被他影響.
你要假裝哭.
因為他在哭,所以你要假裝哭.
但你不要真的內心也哭.
但聖經是否這樣說呢.
新約是否這樣說呢.
似乎不是.
在《羅馬書》第十二章是這樣說的.
保羅說要與喜樂的人一同喜樂.
與哀哭的人一同哀哭.
這裡不是說假裝哭那麼簡單.
因為剛才我們已經說了.
保羅是承認哀傷這個事實.
保羅自己也說自己會哀傷.
所以這裡與哀哭的人一同哀哭.
不是在說羅馬哲學裡面那種.
你與哭的人一起假裝哭.

$^{601}$不是這樣的.
而是你要與哭的人一起感受他們的哀傷.
這是很重要的.
我們的哀傷不只是一個很個人化的情況.
正如剛才我一開始說.
Paul Edmund在說甚麼.
Paul Edmund這個心理學家說.
哀傷是一種社群的信號.
一個人哀傷的時候在說甚麼.
一個人哀傷是在說他需要被幫助.
一個人哀傷是在說.
他失去了一些他生命當中很重要的一部分.
我們能否看到這個信號.
我們能否看到一個人哀傷.
其實是在說他需要被幫助.
他需要被安慰.
他需要同行和同伴.
所以在這裡.
其實是Paul Edmund的東西.
反而比較吻合新約的東西.
就是在說當哀傷的人在哀傷的時候.
他在放一個信號給我們聽.
在放一個信號給我們聽.
在放一個信號給整個社群聽.
他需要被幫助.
我們能否擁抱他.
能否跟他同行.
這亦是保羅所說的東西.
所以很簡單.
到了今天總結的時間.
哀傷有兩個面向.
今天我們看新約聖經的時候.
第一,我們能否承認哀傷.
這件事不簡單.
雖然我們剛才在取笑斯多亞哲學家.
但其實回到現實生活.
回到我們現代.
其實承認哀傷也不是一件容易的事.
特別是男性,我不知道.
這不是性別的偏見.

$^{641}$不過我覺得對於男性來說.
可能也是.
當男性經歷哀傷的時候.
我們很難去承認,對吧.
因為中國人也有點像斯多亞哲學.
中國的文化是說甚麼.
男兒流血不流淚,對吧.
哭泣是不好的.
是真的,中國的傳統文化.
有點像斯多亞哲學.
哭泣是不好的.
哀傷對男性來說是甚麼表現.
哀傷對男性來說是一種軟弱的表現.
雖然未至於道德差那麼誇張.
不過還是不好.
哀傷也是一種軟弱的表現.
所以我們受中國傳統文化影響的男性.
其實我們不太懂得承認哀傷.
我們不太懂得哭泣.
我接觸過很多弟兄.
當面對親人離開.
或者面對一些很創傷的事時.
通常第一句反應是甚麼.
「我沒甚麼,沒問題」.
「我沒問題」.
這是傳統男性會說出來的看法,對吧.
聽到後,可能我們在座的自己也是這樣.
我們不太想告訴別人自己不開心.
我也不太想告訴自己自己不開心.
因為我不能接受失去一件很重要的事.
但今天,新約聖經跟我們說.
連耶穌自己也跟我們說.
我們需要去承認我們的哀傷.
因為唯有我們能夠承認哀傷時.
我們才能承認.
原來我們失去的東西對我們是重要的.
如果我們不能承認哀傷.
我們永遠都沒有辦法.
去面對我們失去的人.
或是市民.

$^{681}$唯有承認哀傷.
我們才能夠接納,擁抱.
原來我們失去了.
一個很重要的部分在我們生命當中.
最後就是另一件事.
哀傷有一個很重要的社群面向.
今天,教會其實是一個.
與哀傷者同行的群體.
我們不是看到哀傷的人就跟他說.
「不用哭,開心就行了」.
或是「別不開心,他死了」.
「將來天下再見」.
我們不是說這些.
我們是跟一些哀傷的人同行.
因為當人哀傷時.
是告訴我們他需要受安慰.
他需要被幫助,他需要支援.
今天我們的教會.
是否能跟哀傷的人同行的教會.
我們不是看到哀傷的人.
我們假裝哭泣去陪伴他.
而是我們真的因為他的哀傷.
而牽動到我們的內心.
因為我們也能感受哀傷.
我們每個人都能感受哀傷.
以至我們能跟他的哀傷共鳴.
我想這才是新約聖經教導我們.
如何與哀傷的人同行.
最重要的重點.
所以最後很想我們有少許時間.
一起有些默想,有些討好的時間.
在這一刻,很想….
我不知道你在生命當中.
這一刻是否有經歷一些失去.
或是一些很艱難的事情.
會否在這段日子.
有一些你很重視的人離開了你.
或是你想起一些.
離開了你很重視的人時.
原來你過去從來沒有承認過.

$^{721}$你真的失去了生命的一個很重要的部分.
這一刻很想你安靜下來.
去到主的面前.
你嘗試回想這些人,事,物的時候.
你嘗試跟主說.
「主,原來我真的覺得傷心」.
「原來我真的很重視祂」.
「原來我的生命真的殘破了」.
「我的生命真的破碎了」.
因為這是失去.
你嘗試….
特別是弟兄,我不知道他是否弟兄.
很多時候….
很多這些失去很快就在回憶中閃過.
我們沒有理會.
特別很想邀請弟兄.
嘗試回想這些人時.
去擁抱這份失去.
這些人對你來說是重要的.
你是能夠去承認你的哀傷.
你是能夠在上帝面前接受哀傷,失去.
是你生命的一部分.
主,今天我們來到你的面前.
今天我們討論哀傷的課題時.
真的很難說.
因為我們的人生.
多多少少都面對不同的失去.
我們很怕去接觸這些失去.
因為去接觸這些失去時.
我們覺得很痛.
因為我們很想迴避這種痛苦.
所以我們不能夠….
有時我們不能夠去承認.
我們真的很痛,真的很哀傷.
但主這一刻幫助我們.
就好像你昔日在赫西拉瑪利園.
面對自己的生命將要逝去時.
你也向你的門徒表達.
你哀傷至死.
主今天幫助我們.

$^{761}$能夠接納我們生命中的失去.
原來對我們的生命帶來很大的衝擊.
令我們的生命變得破碎.
主今天給予我們力量能夠哀傷.
以致我們能夠接納失去.
原來是我們生命的一部分.
原來我們真的失去了.
一些很重要的人,事,物.
主今天幫助我們.
讓我們能夠跟過去的失去復和.
讓我們能夠跟失去重要的人,自己復和.
以致我們可以真正被醫治.
以致我們真正可以面對自己.
面對我們自己的人生.
主今天幫助我們.
也祝福UV,祝福我們的教會.
是一個能夠與哀傷者同行.
同哀傷的教會.
主,你讓我們這個群體能夠承受哀傷.
我們不再看到別人哀傷就跟他說.
「你開心就好」.
我們不再只跟別人說.
「你有復活的盼望」.
而是我們能夠被別人的哀傷.
衝擊我們的內心.
我們能夠與別人的哀傷同感和共鳴.
主今天幫助我們.
讓我們能夠面對自己的哀傷.
也能夠與哀傷的人同行.
是奉主耶穌基督的聖名而來.
阿們.
\newpage



\section{箴言 22:6}
\label{sec:lPvFkn_YDDw}
\textbf{父母的職事 (箴言 22\_6) - 陳崇基牧師}
\newline
\newline
連結: \href{https://youtube.com/watch?v=lPvFkn-YDDw}{\texttt{https://youtube.com/watch?v=lPvFkn-YDDw}} ~~~~ 語音日期: 2025-06-19
\newline
\newline
\hyperref[sec:EgZkScLaImE]{\small{< < < PREV SERMON < < <}}
~
\hyperref[sec:index_chronic]{\small{[返順時目]}}
~
\hyperref[sec:index_scriptual]{\small{[返順卷目]}}
~
\hyperref[sec:cZT_1gxTxfc]{\small{> > > NEXT SERMON > > >}}
\newline
\newline
箴言 22:6
\newline
\begin{longtable}{cl}
\hline
\hline
章節 & 經文 (和合本修訂版)\\
\hline
22:6 & \begin{tabularx}{0.7\textwidth}{X} 教養孩童走當行的道，就是到老他也不偏離。 \end{tabularx} \\ \\
[1ex]
\hline
\hline
\end{longtable}
$^{1}$今天是父親節.
父親節當然由我來說.
所以今早我女兒提醒我.
今天是父親節.
我才記得.
今天是父親節.
我們男人很大量的.
沒所謂的.
沒有人理會我們.
不過今天父親節.
我想說一個跟父母有關.
跟家長有關的訊息.
所以我們暫時停一停這個司堂傳.
我想說.
作為一個家長.
無論是作父親或作母親.
我們在聖經裡面有什麼值得我們學習.
還有什麼我們需要被提醒.
拜道威今天是父親節.
或者母親節也好.
如果在座的是單親家庭.
單親的家長.
我是非常敬佩你們的.
因為你一個人揹著家庭是非常困難的.
但作為父或母.
到底有什麼值得我們關注.
有什麼值得學習.
雖然我不敢說我是一個完美的.
模範的父母.
但也有少許經驗之談可以說一下.
無論是成功或失敗的經歷.
我相信我都可以用我自己的故事.
或我自己的例子來說說.
今天在說父親節的時候.
只是說一節聖經.
就是《箴言》22章第6節.
教養孩童.
走當行的道.
就算到老.
他也不偏離.

$^{41}$一句很出名的一段箴言.
22章第6節.
首先搞清楚.
我們看聖經裡面的箴言.
Proverbs.
箴言不是guarantee.
箴言不是應許.
不是promises.
不是guarantee.
沒有寫包單.
你教養孩童.
當行的道.
到老他死那天.
都會一定是名成利就.
是成功.
一定是敬畏神等等.
因為小孩子長大之後.
他成為成人.
他自己會有一些抉擇.
這是他對自己人生的一些責任.
但作為父母的話.
這段說話.
為什麼這麼出名.
因為提醒我們.
到底父母的職責是什麼.
箴言和聖經裡面的.
叫做wisdom literature.
智慧文學是做什麼的.
剛才說了不是promises.
也不是guarantee.
但是什麼呢.
是箴言.
你記得那種入手的方法.
就是說觀察世務.
觀察世界裡面.
即是上帝的order是什麼.
上帝在人間.
在人生裡面.
道德的規律是怎樣.
是從下而上的.

$^{81}$所以他觀察了之後.
智者就說什麼.
他就為你去解讀人生.
告訴你一個智慧的人生.
應該是怎樣歸納一些大的原則.
所以一些general的principle.
general的guideline.
去指引我們怎樣去生活.
這一段經文.
就是說到作為父或者母.
或者甚至在我們中間.
作為外公外婆.
爺爺媽媽等等.
我們應該要做些什麼.
通常的翻譯就是教養孩童.
走他當行的路.
英文或者是train up.
educate.
但其實原文是很搞笑的.
它不是普通的一個教導這個字.
但中英聖經在今天現代的翻譯.
通常都是翻譯為教養.
教導.
教導孩童等等.
但其實教養這個字背後.
其實是一個比較特別的字.
在聖經裡面只用了五字.
是什麼字呢.
應該怎樣去理解它呢.
我們首先將教養這個字不要它.
如果只是看教養背後的字.
hanuk.
到底是什麼意思的時候.
你就試試去猜一下.
我說過整個聖經裡面只用了五字.
第一次是在《新命紀》20章第五節.
那裡講到以色列人未進入應許之地加林地之前.
他們要作好準備.
但當有時他們要打仗的時候.
這裡就提到有一段律法.

$^{121}$他說如果出兵打仗的時候.
發現那些士兵是買了新屋.
有些新的田地.
如果他要打仗的話.
因為他會身亡.
所以不用打.
他就回去做什麼.
你看他做什麼.
他說官長要向士兵宣告說.
誰建了新的房屋.
誰做了建築.
誰買了新的房屋.
尚未奉獻.
或者新月本未行過奉獻禮.
這個人可以回家.
以免他陣亡.
那間屋就歸給其他人.
人家去奉獻了.
所以是頗仁道的.
是不是.
你買了新屋.
如果你未奉獻.
未開波.
未開始.
或者未獻上給耶和華的話.
請你回去.
不用上戰場了.
這個字就是.
《箴言二十二章》.
六字"Hanuk".
那個字.
幸奉獻禮.
或者奉獻的意思.
所以是新的東西.
是開波有關的.
和基礎有關的.
beginning有關的.
就是這個字.
然後第三次容.
就是《列王記上》第八章.

$^{161}$這裡說到聖殿奉獻的時候.
所羅門建了聖殿.
向耶和華獻平安祭.
有多少牛有多少羊等等.
這樣王和全以色列人.
為耶和華的殿行了奉獻的禮.
Hanuk.
就是這個字.
箴言就是做教養.
但是這裡做應該是奉獻的.
就是新的房屋也好.
或者甚至是聖殿神的房屋.
開始啟用的時候.
他就強調第一.
the first beginning.
first foundation的時候.
要做什麼.
所以就奉獻禮.
然後歷代志下.
都是說同樣的東西.
就是聖殿被建造的時候.
他們需要去做奉獻禮.
Hanuk.
如果懂得玩一些希伯來文的音標.
本身沒有音標的.
加一些音標進去的時候.
動詞變成名詞也好.
你就記得Hanuk這個字.
大概就是光明節.
Hanukkah.
記得嗎.
到我們大概聖誕節的時候.
12月的時候.
猶太人就會慶祝修奠節.
所謂festival of rites.
Hanukkah.
那個字就是這個字.
你記得這個背景嗎.
耶穌之前大概160年左右.
當時瑪嘉比革命.

$^{201}$他們成功趕走了希臘的大軍.
然後他們重新拿回聖殿.
重新拿回耶路撒冷.
然後他們就Hanuk.
他們就重新去奉獻.
重修聖殿.
所以又是說一些新的東西.
和一些基本的東西.
是要做的.
所以可以翻譯做什麼呢.
22章第六節.
應該是奉獻孩童.
行所當行的路.
或走他的道.
就是道路他也不偏離.
所以是要去dedicate.
去奉獻這個孩童.
去走他當行的路.
行在神的道那裡.
所以這裡的字強調的是.
那個啟用.
有一點像就職.
譬如總統就職.
開始的時候.
奉獻他.
dedicate他.
第一件事.
趁早的你要做什麼呢.
奉獻這個孩童.
當然不是說嬰兒奉獻禮.
這麼簡單.
只是儀式.
就是我們的取態上.
就是這個小孩.
是需要將他奉獻給上帝的.
是屬於上帝.
要教導他.
行在上帝的道裡.
等於你買了一樣東西.
你給down payment.

$^{241}$買個車也好.
買間屋也好.
那個down payment就是一個investment.
一個投資.
第一筆投資的資金.
你放在那裡.
這裡說奉獻孩童.
或者教養孩童.
使他行當行的路.
所以這個字.
教養.
或者譯做奉獻.
其實有個神聖的意思在這裡.
因為你記得什麼時候用.
是聖殿的時候用.
你在獻殿的時候.
你就說.
Hanukkah.
就是做一個奉獻禮.
一個最重要的事情.
這裡說到父母要做的事情.
或者我們長輩要做的事情.
就是奉獻孩童.
是一個天職.
或者一個神聖的職事.
在我們在座中間.
如果你是做父母的話.
你有沒有想過.
這個是一個sacred的ministry.
是一個神聖的職務.
是一個天職.
在我們有了孩子之後.
或者我們去養育下一代.
去影響下一代的時候.
我們需要第一件事.
去搞清楚我們的取態.
上帝是他的忠心.
Hanukkah.
是他開始.
今時今日做父母很忙的.

$^{281}$因為我們最緊張最打緊的是什麼.
衣食住行.
入最好的校區.
學最多的東西.
學最好的東西.
這些就是充實著我們做父母的時間表.
充實著兒女的時間表.
Don't get me wrong.
這些是很好的東西.
很好的.
教他最好的東西.
讓他進入最好的教育.
這個是我們作為父母應該有的基本責任.
但是聖經裡面說什麼.
做父母的代表是王不夜.
很多時候在亡亂沖濕裡面.
上帝就變成了主腳.
有什麼可以捨棄的.
先捨棄上帝.
因為上帝有些阻頭阻勢.
是不是.
阻住我的兒子.
命成理就.
對不對.
這裡說的是相反.
奉獻孩童走當行的道.
或者他的道.
走在他的道.
直到中老.
他都不偏離.
想清楚你們做一個父母.
想清楚你的職事.
是一個神聖的職事.
由BB出生.
直到可能你是Emptiness.
好像我一樣.
直到中老.
你都是有一個神聖的職務.
就是要奉獻.
Dedicate.

$^{321}$他當行的路.
你花了多少心思.
花了多少時間.
花了多少力氣.
去做這件事.
是不是有些本末倒置.
如果聖職既然是天職的話.
應該怎麼做.
他說當行的道.
奉獻孩童走當行的道.
走他的道.
直到老的時候.
他又不偏離.
有時候英文的翻譯就是.
In the way he should go.
或者According to his way.
His way是什麼way.
是一個孩童的way.
我們就說.
教養孩童 奉獻孩童.
要走一個他應當走的路.
那個孩童當行的道.
我們要教養他.
所以有很多說教養孩童的書.
都是用這句發揮出來的.
就是說應該怎麼去按照那個孩童.
然後教導他 養育他等等.
有些像是恩才私教.
你要認識你的孩童.
然後你教導他.
What is his way.
What is her way.
然後你教導他.
According to her way.
Walk her way.
是不是.
不一定是my way.
是her way his way.
Right.
你明白我的意思嗎.

$^{361}$作為父母的.
我們值得問這個問題的.
是不是.
值得問這個問題.
不是我的意思.
應該按照他的性格.
或者他的取向.
His way and her way.
Right.
所以有很多書都是說這件事的.
新一代的父母我知道.
我們很多時候很緊張.
很關注這件事.
怎樣去教他.
有什麼育嬰的方法.
或者什麼秘笈.
那些秘笈可能是我們.
自己家庭的傳統留下來的.
或者我們自己的觀念留下來的.
我們覺得這樣教才是對的.
或者我們會聽外面有很多專家說.
有些KOL說.
應該怎樣有什麼secret.
有什麼教養的方法.
今時今日我們很多這些.
因為我們又說Big Data.
我們又說很多這些studies.
有什麼新的研究.
我們經常都試探討.
很緊張 很著緊.
就是用什麼教養的方法.
來教我的孩子.
怎知道.
就是為了這件事.
我們夫婦之間經常有爭執.
真的.
我明白.
我們需要有一個好的方法.
一本很受歡迎的書.
叫《焦慮新世代》.

$^{401}$Ancient Generation.
裡面就試探討用什麼方法.
昨天我用了大概一個小時看這本書.
裡面提到一個家長裡面一個很大的問題.
今時今日家長很大的問題.
就是說到.
如果我真的找到一個秘訣.
想到一個方法的話.
我就做好一個教養孩童的工作就好了.
我就得到一個很優秀的產品.
If I can just find the right method.
我一定會找到一個right product出來.
這個作者指出.
這就是我們今時今日最大的問題.
你記得我剛才說什麼.
為了教養的方法.
弄到怎樣.
弄到吵架和家居屋閉.
我不管你那個方法是什麼妙方也好.
如果你的環境弄到這樣的話.
我覺得有些問題.
這個就是作者指出的.
很多時候我們弄錯了焦點.
我們覺得找到一個妙方找到一個秘訣的話.
我的兒子或我女兒.
就一定會變成一個right product.
你只要這樣聽的時候.
你已經知道有問題.
因為你好像做一份工作.
一個產品的標誌.
弄一個iPhone出來.
一個right product我們就搞定.
但如果弄到家居屋閉.
或者夫婦之間有些紛爭.
或者有些張力的時候.
你覺得小孩在這樣的環境長大.
是一個健康的環境嗎.
你覺得是不是一個健康的環境.
作者說了什麼.
他說這是一個wrong approach.

$^{441}$真正應該好一點的比喻.
教養的比喻.
是gardening 是園藝.
大家有沒有種東西.
種一些瓜 種一些菜等等.
怎樣能夠種出東西.
你去灌溉 給水給他.
給他營養.
然後在一個好的環境.
以致他能夠生長.
然後取一些雜草.
環境好的時候.
那個植物才會成長.
不是get the right method.
而是create a good environment.
一個好的環境.
小孩如果從小到大.
見到你們兩個.
或者見到家長.
是不會怎樣處理關係的.
不懂得怎樣相愛的.
不懂得寬恕 不懂得和解.
相處的藝術是零.
你覺得小孩學到什麼.
我告訴你.
Your kid is not stupid.
They see it.
聽清楚了.
Your kid is not stupid.
Maybe you are stupid.
你恃著你有什麼大數據.
你覺得你有很好的方法等等.
我最近看了一個研究.
數據顯示什麼呢.
追溯一些小孩長大之後.
有什麼影響他最多.
大數據說什麼呢.
原來不是父母.
你不要高估你自己.
最影響小孩成長的往往不是你.

$^{481}$So no matter what method you use.
不要太高估自己.
但有什麼環境令到他成長.
這個才是更加重要.
在他面前學習怎樣寬恕.
怎樣和解.
大人怎樣去處理事情.
在我們中間可能有很多.
我們成長的時候的背景.
都是很rough的.
我自己也是.
小時候父母很多時候吵架.
我不管父母看了什麼書.
偶爾在書架我看到我媽.
會看怎樣教養兒女等等.
但我最記得的是什麼.
最記得的是我不知道他們為什麼會吵架.
我最記得的是我問自己.
到底搞什麼鬼的大人.
為什麼搞不清這些問題.
吵什麼鬼呢.
令到我很不高興.
所以我就經常不回家.
經常不讀書.
所以為什麼我打桌球這麼厲害.
因為整個中學年代.
我每天都打桌球.
每天都給錢打桌球.
加很多學費.
就算沒有打那麼多年.
我都還可以.
多的是什麼呢.
就是家庭沒有一個好的環境.
我記得的是什麼.
就是環境是一個負面的環境.
我想告訴你.
如果你想小孩子成長.
品格上成長.
學識上成長.
你想小孩子追求學習.

$^{521}$學習的取態正確的取態.
有一個智慧的人生.
我想告訴你.
如果你想這些事情.
發生在你的下一代身上.
你就先要照鏡子問你自己.
到底你是不是一個成長的人.
你作為一個大人.
你是不是成熟的.
如果相反你是一個嬰兒的話.
你品格上是一個嬰兒.
你是沒有成長的.
我再次告訴你.
這些不是愚蠢的.
你是愚蠢的.
做父母怎麼做呢.
是新教.
不是現教.
是你的命擺出來.
小孩子是可以學校的.
當然不是說我們是完美的.
但起碼我們走得正站得正.
我們是例子.
他可以學校我的.
如果你想小孩子讀書.
而你又不追求知識.
你又沒有一個好的學習心態的話.
就忘記吧.
對嗎.
忘記吧.
很值得思考吧.
有想過父親節講這些嗎.
當頭棒學.
是真的.
當行的道.
原文是行他的道.
his way.
我們通常說his way是甚麼.
當行他的道.
我們就說孩童的道.

$^{561}$因為最近的他就是孩童.
奉獻孩童走他的道.
第二節.
但他是不是一定是孩童呢.
我想他特意有一個模糊的地方.
就是要你咀嚼想清楚.
奉獻孩童走他的道.
是否只指到孩童的道呢.
你去查查上文.
《聖經》是這樣看的.
看上文頭兩節幫幫手.
很肯定不是第五節的歪曲人的道.
沒理由奉獻教養孩童走歪曲人的道.
如果不是這個是他的話.
是誰呢.
再看多一節.
第四節.
行為耶和華.
心裡面謙卑的就有財富尊榮生命為上賜.
這個就是他.
his way.
耶和華的道.
會不會是這樣的意思呢.
教養孩童奉獻孩童走耶和華的道.
因為《箴言》裡面經常說.
耶和華的道路 智慧的道路.
你要鼓勵你的小朋友行耶和華的道.
智慧書經常有這樣的道理.
父母就是智慧的化身.
我們就是智慧的embodiment.
就是智慧的載體.
由第一章到最後那章.
父親或母親都是叮嚀小朋友要做一個有智慧的人.
行在智慧的道上驚威耶和華.
是不是.
教養孩童奉獻孩童要走耶和華的道.
怎樣走 怎樣鼓勵他.
又是新教.
不是現教.
你要知道在整個人類歷史裡面.

$^{601}$一直以來 由古代到現代.
都沒有所謂的parenting.
沒有parenting的專業.
為什麼呢 由小到大.
由古代到現代.
小朋友的成長.
在哪裡成長.
在家庭的環境成長.
幾代同堂.
他觀察到一個大的community社群.
他們應該怎樣去生活.
父母就變成載體.
怎樣教小朋友有智慧去生活.
小朋友看到你怎樣有智慧生活.
他就學到.
就沒有說有個parenting expert.
有個KOL出來講.
或者有哪個專家來教那個小朋友.
你見到那個球教給誰.
教給別人.
教不好就來教會.
教會可以教好他.
來不到就去基督教學校教他.
載體在哪裡.
載體不在我這裡.
你明不明.
1970年代開始才有parenting這個東西出來.
所謂專家專業.
specialization.
70年代開始.
為什麼會有一樣東西這樣掘起.
parenting.
就因為70年代開始整個現代的社會.
就變成那個家庭的大小越來越縮窄.
有所謂核心家庭出現.
nuclear family.
什麼叫nuclear family.
就是說你的social unit.
就是社會裡面的單位只剩下誰.
家長和小孩.

$^{641}$兩代夾在裡面.
再上一代或者其他的社群.
完全是失去了.
有這個nuclear family的現象出現.
所以沒有了其他社群給他的影響.
所以才有所謂養育兒女學.
parenting.
這東西慢慢慢慢變聖所謂parenting的expert.
有一些教養的.
教養孩童的一些妙方.
或者一些專家出現.
當然我不是說這些不好.
這些都是好的.
但問題在哪裡.
問題就是智慧的載體去了別處.
不再是父母自己的天職.
不再是父母自己的責任.
不再是我要set一個榜樣.
他搞不定的時候.
我就送他去工廠.
然後有專家來搞他.
專家搞不定.
我就罵那個專家.
blame shifting.
責任去了別處.
不再是我作為父母.
或者母.
again 你如何教導孩童.
教導我們的下一代.
走耶和華的道.
是新教.
如果你不想你自己的兒女.
特別是今時今日.
有所謂digital addiction.
數碼的成癮.
是不是要找專家.
可以.
但最大問題就是.
到底你自己是不是這樣.
如果你自己都沒有自律的話.

$^{681}$你如何作為一個載體.
作為一個榜樣.
告訴下一代.
你要學這個東西.
如果小朋友讀不好書的話.
你有沒有認真的.
有一個學習的取態.
你是不是學到老.
做到老學到老.
大家知道我兒子大學畢業.
兩個星期前回來.
他執死雞.
找到工作.
這樣也好.
我叮囑他.
拜託你好好地工作.
特別這一代EQ不夠的人.
如果你EQ稍為好一點.
你應該贏了起跑線.
每次他回來的時候.
我都會和他去後壹度聊天.
晚上.
有時候拿一杯啤酒喝.
真的.
我們聊到晚上兩三點.
好累.
但我又好感恩.
我可以這樣單對單和他聊天.
他和我說感情有什麼問題.
關係有什麼問題.
生活裡面有什麼想到的.
但我覺得很有意思.
你問我有什麼秘方.
我沒有什麼秘方.
這個可能是我做牧師的.
我需要檢討一下自己.
因為我從小到大都沒有說教.
我從來沒有和他們說.
我們現在家庭零修.
我們要看聖經.

$^{721}$沒有的.
我只不過在他們面前.
做一個真實的教徒.
是不是.
我有失敗的地方.
我有走得正站得正的時候.
我是真實的告訴他們.
這就是現代人的生活.
榜樣.
去教會.
我從來沒有逼他們.
因為我們去完教會.
我們起碼會有一點改變.
我們不是當一件差事.
因為上帝是一個活的上帝.
我和上帝之間的關係是一個活的關係.
去完教會我會和他們說.
今天學了什麼等等.
有時候他們會說很欣賞哪個哥哥姐姐.
或者叔叔伯伯等等.
教會的長者.
在他們身上有的影響.
不是去擦桌子.
在社群裡面.
他們有什麼智慧的載體.
就是這樣.
一個小孩的成長就是需要這樣東西.
是不是.
上個星期五我去了一個安息禮拜.
一個姐妹的安息禮拜.
認識了她很多年.
聽她下一代再下一代.
幾代的.
去到述史的時候.
我很感動.
因為這個姐妹.
被提到的一些故事.
一些憶念.
都是一些和她生命有關.
她怎樣去事奉.

$^{761}$當她有病痛有癌症的時候.
她怎樣去抗爭癌症.
怎樣保持信仰.
沒有一個子女或者孫兒女.
是說她有什麼妙方.
她教我們怎樣.
有什麼方程式.
就是她的命.
真真實實的.
做一個神的子民.
是新教.
不是靠嘴巴.
求主幫助我們思想清楚.
父親子的訊息不是開心的訊息.
因為我們做爸爸的很堅強.
應該有當頭棒.
想清楚你在搞什麼.
想一想你的家庭環境.
你怎樣塑造到.
你的智慧是生命的載體.
你有一個榜樣走出來.
下一代是這樣被影響生命的.
我們誠心低下祈禱.
為我們中間作父作母親的.
甚至作長輩的.
主我們有這個使命.
求你幫助我們.
憐憫我們.
因為我們實在很脆弱.
我們實在很軟弱.
主教導我們怎樣愛.
教導我們怎樣去行基督的道.
行耶和華的道.
以祝我們奉獻.
教養我們的下一代.
讓他們走在生命的道路上.
求主您憐憫幫助.
封住耶穌的名字.
聖誕快樂.
\newpage



\section{使徒行傳 3:1-10}
\label{sec:cZT_1gxTxfc}
\textbf{福音的突破 (使徒行傳 3\_1-10) - 陳崇基牧師}
\newline
\newline
連結: \href{https://youtube.com/watch?v=cZT-1gxTxfc}{\texttt{https://youtube.com/watch?v=cZT-1gxTxfc}} ~~~~ 語音日期: 2025-06-24
\newline
\newline
\hyperref[sec:lPvFkn_YDDw]{\small{< < < PREV SERMON < < <}}
~
\hyperref[sec:index_chronic]{\small{[返順時目]}}
~
\hyperref[sec:index_scriptual]{\small{[返順卷目]}}
~
\hyperref[sec:1fX6w4FAuCw]{\small{> > > NEXT SERMON > > >}}
\newline
\newline
使徒行傳 3:1-10
\newline
\begin{longtable}{cl}
\hline
\hline
章節 & 經文 (和合本修訂版)\\
\hline
3:1 & \begin{tabularx}{0.7\textwidth}{X} 下午三點鐘禱告的時候，彼得和約翰上聖殿去。 \end{tabularx} \\ \\ \relax
3:2 & \begin{tabularx}{0.7\textwidth}{X} 一個從母腹裡就是瘸腿的人正被人抬來，他們天天把他放在聖殿的一個叫美門的門口，求進聖殿的人施捨。 \end{tabularx} \\ \\ \relax
3:3 & \begin{tabularx}{0.7\textwidth}{X} 他看見彼得、約翰將要進聖殿，就求他們施捨。 \end{tabularx} \\ \\ \relax
3:4 & \begin{tabularx}{0.7\textwidth}{X} 彼得和約翰定睛看他，彼得說：「看著我們！」 \end{tabularx} \\ \\ \relax
3:5 & \begin{tabularx}{0.7\textwidth}{X} 那人就注目看他們，指望從他們得著甚麼。 \end{tabularx} \\ \\ \relax
3:6 & \begin{tabularx}{0.7\textwidth}{X} 彼得卻說：「金銀我都沒有，但我把我有的給你：奉拿撒勒人耶穌基督的名起來行走！」 \end{tabularx} \\ \\ \relax
3:7 & \begin{tabularx}{0.7\textwidth}{X} 於是彼得拉著他的右手，扶他起來；他的腳和踝骨立刻健壯了， \end{tabularx} \\ \\ \relax
3:8 & \begin{tabularx}{0.7\textwidth}{X} 就跳起來，站著，又開始行走。他跟他們進了聖殿，邊走邊跳，讚美神。 \end{tabularx} \\ \\ \relax
3:9 & \begin{tabularx}{0.7\textwidth}{X} 百姓都看見他又行走，又讚美神， \end{tabularx} \\ \\ \relax
3:10 & \begin{tabularx}{0.7\textwidth}{X} 認得他是那素常坐在聖殿的美門口求人施捨的，就因他所遇到的事滿心驚訝詫異。 \end{tabularx} \\ \\
[1ex]
\hline
\hline
\end{longtable}
$^{1}$我們繼續說我們跨越的系列.
我們說到Crossover.
我們希望每一位大英姐妹都能夠和應.
特別是在這一年裡面.
我們能夠成為主的見證.
這次我們繼續看《士徒行傳》.
看到第三章.
第三章有一個故事.
大家可能從小到大.
如果我們中間有從小到大.
在教會長大的.
這個故事應該不是陌生的故事.
金和銀我沒有.
我奉拿撒勒人耶穌的名.
叫你起來行走.
就是一個這麼簡單的故事.
一個故事是說到.
彼得和約翰在聖殿的門口.
看到一個生下來是瘸腿的.
跛腳的.
就去念問他.
然後奉上帝的名.
叫他起來行走.
所以這個《士徒行傳》.
到了第三章就說出一個這樣的故事.
我們要知道上文說的是什麼.
大家記得第二章說的是什麼嗎.
最後的時候大概好像是教會的自拍.
如果教會或者團體大家吃飯的時候.
我們會拍個自拍.
會拍一拍大家的畫面.
那個描述是怎麼樣的呢.
這個就是初期教會在耶路撒冷建立的時候.
他們的所謂自拍.
說到什麼呢.
他說眾人都心存敬畏.
士徒們又行許多奇事神蹟.
信的人都聚在一處繁物共融.
有賣料田產和家業.
照每一個人所需要的分給他們.

$^{41}$他們天天的同心合意.
行切的在聖殿裡面敬拜.
在家裡面抹餅.
然後全著歡喜.
坦誠的心用飯.
贊比神.
得全體百姓的喜愛.
主就將得救的人數天天加給他們.
這裡說到幾樣東西.
第一就是他們同心的來敬拜.
去看見這個群體.
看見教會的時候.
他們應當是這樣的.
同心來敬拜.
並且是繁物共融.
這裡不是說因為國家要求.
或者是政權要求.
是個人自願的.
主要想說什麼.
就是每一個人聚在一起的時候.
當時的教會就是自願的.
慷慨的來奉獻的一群群體.
這些寫照都是去inspire我們.
到底我們的教會是一個怎麼樣的教會.
你看到這裡第二章.
最後的時候.
說到初期教會的寫照.
自拍的時候.
他說有很多奇事和神蹟.
當時聖靈降臨的時候.
五旬節降臨的時候.
就是一個劃時代的時間.
這個劃時代的時間.
就會有一些突然來臨的.
是劃時代的事情會發生.
譬如他們講方言.
還有一些聖靈大能的工作等等.
所以你看到他們有很多奇事發生.
可能病的得醫治等等.
這些都是所謂的apoc making events.

$^{81}$就是一個新的朝代.
新的時代到來.
剛開始的時候.
總是會有一些劃時代的事情發生.
所以記載了很多這樣的事情.
我們就想.
如果你是寫《史唐傳》.
好像路加那樣.
到底你記什麼故事呢.
有很多嘛.
是不是.
路加這裡就很有趣了.
他記載了一個生下來是瘸腿的人.
得醫治的故事.
這個就是第三章.
所以我們不妨看第三章.
這裡我大概有兩點.
我想點出的.
第一個重點就是.
這個故事到底為什麼要寫出來.
和第一個記載神蹟的故事.
在初期的教會.
我相信這裡有兩點值得我們去留意的.
今時今日怎麼對我們講話.
第一樣這個故事.
去挑戰我們排外的心態.
教會很多時候有排外的心態.
這個故事是挑戰我們這件事.
第二樣就是去建立一個融會的文化.
教會應該是一個inclusive.
一個welcoming的社群.
我們看看這個故事想講什麼.
他說下午三點鐘的時間.
彼得和約翰上聖殿去.
一個從武福裡面就瘸腿的人正被人抬來了.
他們天天把他放在聖殿的一個門口.
叫做尾門.
求進殿裡面的人去施捨給他.
這個大概三點鐘下午的時間.
或者英文的翻譯就是denied hour.

$^{121}$第九個小時.
就是早上六點鐘.
你計九個小時就是幾點.
就是下午三點.
大概這個時間.
當時我們知道猶太人是有定時去到聖殿.
去到禱告和敬拜的.
其中一個時間就是denied hour.
下午三點的時候.
可以想像去聖殿的時候.
就是很多人.
很多人就會上去聖殿.
這個時間他們去到禱告.
和獻上他們的祭物,梵祭等等.
這裡就述說到一個從武福生下來的.
就是瘸腿的人.
他就被人抬到什麼地方.
一個聖殿的門口.
尾門.
然後叫人去施捨給他.
almsgiving.
這些都是當時的聽眾.
他們聽見的時候就會明白是什麼事情.
因為一個武福生下來的瘸腿的人.
他能不能進到聖殿.
不可以的.
比如在《五經》裡面說什麼.
利未記二十一章.
拿來不詳細讀.
這裡說是說給阿倫.
阿倫就是祭司的家族.
還有他的後代是什麼呢.
他說凡是殘疾的.
都不可以去獻上帝的食物.
凡是有殘疾的.
失明,瘸腿,五官不正等等.
如果你五官不正.
你就不可以去獻祭.
就是你歪嘴歪舌.
斷手斷腳.

$^{161}$駝背.
坐得不好等等.
有很多這樣的技術.
就不可近前來.
所以就記載了一個瘸腿的.
是不能夠進去聖殿的.
所以他什麼時候開始瘸腿.
什麼時候開始瘸.
出生到現在.
所以你就可以想像這個人是一個猶太人.
他從出生到現在去到黑石的時候.
他都不能夠進去這個聖殿的.
聽就聽見很多故事.
聽見他自己民族或者他自己百姓家人.
他們去到聖殿那裡.
然後多麼宏偉.
去到獻祭那種氣氛等等.
他是沒有份的.
因為從小到大.
出生到現在.
他都是一個被咒詛的一個人.
三章第二節說到.
這個生下來被.
被瘸腿的人被拿到聖殿的門口.
叫做什麼呢.
叫做Beautiful Gate.
一道門叫做美門.
這道美門是哪一道門呢.
正如你去教會有不同門進來.
聖殿當時.
Josephus.
其中一個第一世紀的猶太人的歷史學家.
他就提到聖殿有十道門.
十道很宏偉的門.
他說其中有九道是什麼呢.
周圍是包上金與銀.
門框門尾也是如此.
有一道門.
位於聖殿的內院之內.
他說到第十道門有什麼厲害的.

$^{201}$Josephus專門說這件事.
他們建築多麼宏偉等等.
是怎麼宏偉的呢.
是有Corinthian Brass.
如果你走我們downtown.
很多時候一些古代的建築.
或者一些很有格調的建築.
是有Corinthian Columns的.
有沒有留意.
留意一下.
我們芝加哥很多.
譬如Fill Museum等等.
你看到上面.
他們的Columns的頭頂和柱頭.
是特別不同的.
是華美遠勝於Josephus.
甚至說金銀包覆的門.
此外.
聖殿內院的東門是用銅所做的.
極為沉重.
要多少人才能夠將它關上.
要二十個男人才能夠推開那道門.
重到一個這樣的地步.
你看到這樣的背景.
或者這樣的資料的時候.
你就發現這個尾門.
他在小童子中說的尾門.
是有一點諷刺的.
是不是.
聖殿有沒有錢.
有很多錢.
是不是.
連門.
重到一個地步.
要二十個男人才能夠推開.
就是每天早上你開門.
餐館開門辦事的時候.
要二十個男人才能夠推開.
起碼要有二十個人負責推開門.
聖殿很有錢.

$^{241}$收入也是最高的.
但是天天門口.
就有一個生下來是瘸腿的.
求別人來施捨給他.
你看到這個諷刺的地方在哪裡.
起碼在這裡間接地.
想諷刺當時耶路撒冷的聖殿.
就是在一個神聖的地方.
但是這個神聖的地方.
卻是殘缺的人要被捨棄在門外.
你可以靠近但是卻被禁止進來.
是有一種你和我是不同的.
是一種排外的建設.
或者是一種排外的制度.
當時耶路撒冷的聖殿.
沒錯是金碧輝煌.
不過也有很多老舊的東西在裡面.
無論是一些官僚的東西.
或者裡面一些的Power Play等等.
至少這個故事.
起碼間接地發出一個.
你可以叫它做批判.
這個故事一個很ironic的地方.
一個很諷刺的地方.
明明聖殿是很有錢.
並且當時的聖殿不是只是宗教場所.
是銀行.
也是當時的社區中心.
是當時的社會福利的地方.
居然一個地方搞成這樣.
一個排外的文化.
可能是建築的排外.
可能是制度上的排外.
可能是規矩上的排外.
剛才我們說的殘疾不能禁止進來.
那個人說什麼人不可以進來.
祭司.
但是到了大概耶穌的時代.
他們不只是應用在祭司身上.
是應用在平民百姓身上.

$^{281}$譬如看看死海古卷等等.
其他猶太人的著作.
他們把這件事明明只是給神職人員.
給祭司.
他們不能夠有殘疾不能夠有牽制.
但是他們把這件事應用在所有階級.
或者所有的人身上.
所以起碼這裡有一些間接的.
讓我們想一想反省一下.
當然我們今天不需要聖殿.
因為耶穌就是我們的聖殿.
基督就是我們的聖殿.
但是教會作為一個神的群體.
是不是一個排他的群體.
一個獨立的群體.
很多時候有意無意我們會這樣做.
因為我們來了教會.
我們就會自然覺得這個地方.
是我們感覺舒服.
感覺安全的.
我們坐那個位置我們永遠都是坐那個位置.
我們做那個事永遠都是做那個事.
我們見那班人我們永遠都是想見那班人.
這個很多時候就是有意無意我們會有一個.
因為身份因為語言因為過去.
或者因為不同的因素.
我們就會將門關上.
所以這一年我們希望做什麼就是要跨越.
就是我們不要在有形無形之下.
我們自己畫了線.
有一些隔離或者有一些距離.
有一些界線.
我們需要跨越的.
才會知道我們教會不是一個排斥的地方.
不同語言的人.
我們應該在基督裡面是同一個群體.
無論講英文也好講普通話也好.
講西班牙語也好.
都應該是在基督裡面.
In Christ.

$^{321}$我們就是一起.
Otherwise.
這個教會就不能發揮他應該有的功能.
正如當時聖殿沒有發揮他的功能一樣.
所以記載這個故事.
我們就做一些這樣的挑戰.
挑戰一些排外的心態.
然後告訴我們應該要做什麼.
融會的文化是什麼意思.
第三節.
他看見彼得和約翰.
這個他是誰.
那個伸下腿的人.
他看見彼得和約翰將要進入聖殿.
就求他們施捨.
彼得和約翰定睛看他.
彼得說看著我們.
那人就注目看他們.
指望從他們得著什麼.
彼得卻說金和銀我都沒有.
但我把我所有的給你.
就是奉拿撒冷耶穌基督的名.
起來行走.
於是彼得就拉住他的手.
扶他起來.
他的骨和他的果骨立刻撞見了.
就跳起來.
他就站起來又開始行走了.
他就跟著彼得和雅各進了聖殿.
這個應該是第一次了.
從小到大第一次.
他能夠進入這個聖殿.
邊走邊跳讚美上帝.
百姓看見的時候又讚美上帝.
認得他就是素常在美門的門口.
求人施捨的那一位.
又因為他現在得到的遭遇.
就滿心驚訝和詫異.
這裡說到金和銀都沒有.
我小時候長大的.

$^{361}$就是唱這首金句歌.
這個有道理的.
這個故事是最有名的一句句子.
就是金和銀我都沒有.
奉拿撒冷耶穌的名叫你起來.
所以我們小時候就唱這首歌.
但我想你留意一下夏文.
夏文是彼得說第二篇道的.
第一篇道什麼時候說.
五旬節的時候說.
第二篇道就是第三章.
這個神蹟故事之後他有說.
貫穿整個第三章有一個主題.
就是基督已經從死裡面復活.
已經得勝已經升天.
聖靈並且已經降臨.
就發生什麼事.
天國的工程.
Kingdom Project就開始了.
開始建立教會.
開始建立神的子民.
大家記得我說過.
小行傳和哪本舊約的書.
是很有關係的.
就是以賽亞書.
以賽亞書當然很長.
你回去看.
看他臉都青了.
但我想告訴你.
最主要的是說什麼.
就是上帝的國度利益的時候.
發生什麼事.
就是十二支派會歸位.
然後上帝的子民會變成一個新的子民.
有聖靈同在有能力.
然後將舊恩傳到地極.
以賽亞書主要是想說.
有個盼望和有個應許.
結果耶穌基督來到聖靈降臨的時候.
就知道神的國度已經彰顯了.

$^{401}$並且這個神的子民.
不只是單單猶太人.
是包括那些奉耶穌基督名進入來的外邦人.
所以以賽亞書第35章第6節說什麼.
他說那時.
他提到一些以前被exclude的人.
以前被排外的人.
那些腿殘的人.
或者有殘疾的人.
或者不關我們事的外邦人.
那些人會怎樣.
跳躍如鹿.
跳躍這個字是同一個字.
在少林寺第3章.
阿爸的舌頭.
有殘疾的人.
他們能夠來到上帝面前敬拜.
這個就是以賽亞書舊約裡面一直有的願景.
耶穌降臨.
聖靈降臨.
上帝的國度開始的時候.
這個群體應該是一個這樣的群體.
不單止他們聚集的時候.
是包括所有的人.
奉耶穌的名來到的.
都能夠參與在上帝的子民裡面.
大家叫做弟兄姊妹.
大家叫做神的子民.
並且是一個welcoming的community.
連他們聚會的形式.
語言.
他們的態度都是一樣.
我就想問.
到底我們教會.
我們作為神的子民.
我們有沒有這種身份.
我們作為一個神的群體.
我們是不是一個會集人.
能夠去attract人.
能夠去welcome人.

$^{441}$你都可以中間.
在這個時代做這件事.
不是容易的.
因為我們在所有的media.
所有的公共資訊裡面.
我們所接收到的.
很多時候都是排外的.
是language of exclusion.
我們所說的東西.
都是跟別人分開的.
你有你我有我.
特別在這個年頭.
整個世界好像著了火.
我們公共的語言裡面.
你聽聽.
你留意一下.
以色列 加沙 伊朗.
大家都是鬥開火.
你所接收到的新聞.
所有的報導.
全部的論調.
都是去exclude.
或者去妖魔化大家.
主要我不是想說.
有什麼地緣政治的問題.
當然那些問題是非常大串的問題.
但起碼那個language上面.
rhetoric.
那些修辭.
我們用的論調.
都是很多時候.
是那種.
無論是隔岸觀火也好.
或者是.
即是.
即是當作是.
剛好上個星期.
美國說想.
即是那些人說.
美國會幫以色列攻打伊朗的時候.

$^{481}$攻打了.
有個議員說什麼呢.
有個議員居然在他的ex那裡說.
ex即是twitter.
他說game on.
This is not a game.
不是遊戲.
是人命關天.
或者很多時候在一些傳統的教會裡面.
我們都會有一些很流行的.
普遍的一種的末世論.
Biblical prophecy.
即是聖經的預言.
現在在中東裡面實現等等.
這些我絕對不贊成這樣去解釋的.
因為這些.
如果你這樣去解釋的話.
就一種隔岸觀火的心態.
聖經說了這裡會打仗.
所以就讓它打.
或者我們火上加油等等.
隔岸觀火的心態.
是一種排外的語言.
應該怎樣去解釋這些所謂的Biblical prophecy.
我們在傳統教會裡面很笨.
因為我們接受那一套.
是那一種.
這裡一早就說了會打仗.
所以就讓它打.
應該打.
甚至支持它打.
這個很笨.
怎樣去解釋呢.
我覺得這是解經錯謀.
所以你說應該怎樣解釋.
你應該來我的class.
再思解經錯謀.
社會學的學者說什麼呢.
現在有些大數據說.
社交媒體上面的內容.

$^{521}$很多時候都是指.
由一小撮的人.
6\%的人來產生的.
他們發什麼上Twitter.
上$\times$上Facebook.
或者小紅書.
Whatever.
OK.
只是6\%的人去產生.
但是很糟糕的一件事.
就是這6\%的人.
大多都是偏激的人士.
無論他是左或者右.
他們所產生的言論.
就是一種標籤的言論.
在我們這種所謂分化的年代.
Age of Polarization.
我們最喜歡去分別人.
標籤別人.
你不同意我的意見.
你就是自由派.
你就是那群壞人.
我們就想去批評大家.
這種現象叫做Tribalism.
特別是在後現代的時候.
有一種部落主義.
我是這個部落的.
你是那個部落的.
好像以前那些蠻荒的部落.
我們大家彼此是受敵.
其實這個Tribalism.
90年代的時候已經開始有人說.
已經開始預測.
誰知道就Self-fulfilling Prophecy.
現在20多年之後就真的實現了.
因為90年代初的時候.
當時我就看了一些書.
說後現代主義會帶來什麼情況.
其中那個分析家就說.
說後現代來到的時候.

$^{561}$就會各自為政.
並且會有Tribalism.
部落主義.
所以很多時候這些社交平台.
我們所接觸的.
就是us versus them.
我們和他們之間是不同的.
剛才我說這個故事.
我們說金句是什麼.
金和銀我沒有.
但我想你更加留意的.
不只是那句金句這麼簡單.
留意這裡有個動詞.
彼得和約翰代表教會做什麼.
定睛漢那個生下來瘸腿的.
還有什麼.
拉著他的右手.
扶他起來.
然後他就得到醫治.
通常我們著眼上帝的工作.
神蹟等等.
我們著眼的只不過是結果.
What is being done.
那個consequence那個回收是什麼.
本身是瘸的.
現在他走到.
這個就是我們著眼的地方.
但我們沒有著眼的.
就是路加所說的故事.
他有很多動作.
這些行動的本身.
那些動詞.
定睛漢他拉著他的手.
還有扶起他.
其實這些才是真正神的工作.
你有沒有想過.
那種關懷那種同理.
不是走過就算了.
所以不是us versus them.
拉他上來扶他.

$^{601}$然後和他一起進去聖殿.
這個描述.
就是講到當時的教會.
當時的群體.
應該產生一種這樣的果效.
應該做一件這樣的事.
今時今日還有人說.
二十世紀中的時候.
有一個很著名的社會學的實驗.
叫做Robert's Cave Experiment.
Robert's Cave為什麼呢.
因為是在一個Oklahoma.
一個地方叫做Robert's Cave State Park.
那裡做的一個實驗.
那個實驗是做什麼呢.
他就找了兩班男生.
就分開他們兩組.
就帶了他們進去這個State Park.
然後他們就在.
不知道對方存在的環境裡面.
他們就進行一些活動.
就建立他們的team spirit.
這邊就建立他們的team spirit.
但他不知道有另外一個team在那裡.
就建立他們的歸屬感和認同感.
我們是屬於這個群組的.
然後就將他們帶在一起.
就要求他們在一個環境裡面.
大家去競爭.
就是搞一下鬥爭.
大家鬥爭和搞一下對立.
贏的就有獎品.
有獎金的.
所以社會學家就做了一個這樣的.
所謂Robert's Cave Experiment.
就發現他們本身是好好的.
當他們一放在一個場景裡面.
去爭大家和批評大家的時候.
很快他們就會產生.
他們自己這些小朋友.

$^{641}$就會產生那種敵意和偏見.
然後就侮辱大家和破壞大家.
男生就做這些事.
欺凌大家等等.
搶大家的東西.
破壞大家的物件等等.
他說怎樣去化解這個衝突呢.
就是你要將雙方放在一起.
並且給他們一個目標是超越.
跨越他們自己的群體的利益.
一個大一點的目標.
比他們自己更大的.
譬如大家去幫那個State Park.
修理一下水管.
去做一下清潔等等.
是和他們自己本身自身的利益是無關的.
是一些公共的利益.
然後這些衝突就自然化解.
逐漸化解之後.
他們就變成朋友.
這個就是一種inclusive的community.
很奇怪的人性就是這樣.
將我們分開的時候.
建立了那個歸屬感.
但是在一個場景裡面.
如果是有張力的話.
和有一些競爭的話.
自然會有一些敵意產生.
今天無論講國際的關係.
講本地社區的關係.
講教會之間的關係.
其實都是一樣.
我記得很多年前.
我曾經去過以色列巴黎聖地.
我們做了一個短宣.
那個短宣就是參與一個基督徒.
一個阿拉伯人基督徒的事奉.
他做一些和平的事奉.
服侍當地的一些比較貧困的年輕人等等.
很奇怪的一些教會的長老就跟我說.

$^{681}$你都不是去傳福音.
為什麼你要做這件事.
因為我們很多時候.
就將福音的工作只限制於.
我們去派單張.
沒錯那些真的是傳福音的工作.
但不代表只是規限於那件事.
對吧.
不可以只是規限於我們去講.
然後我們應該要有行動去表述那個福音.
我記得我們幾個人就去探訪一個阿拉伯人基督徒的事奉.
叫做穆斯拉夏.
這個阿拉伯人的事奉是做什麼的呢.
他們是將那些以色列和巴勒斯坦兩方的青年.
帶他們一起去曠野.
去露營.
為什麼選擇曠野呢.
為什麼不選擇以色列或者選擇西岸巴勒斯坦地.
為什麼不選擇.
因為你去曠野的時候.
就是一個中立的地方.
不屬於任何人的.
聖經裡面一直都在講曠野.
都是No Man's Land.
不屬於任何一個政體.
不屬於任何一方的.
他們帶他們去那裡.
然後有兩個禮拜一個月的露營.
大家一起透火.
大家晚上的時候有營火會.
大家在營火會裡面彼此相食.
大家見到大家的面孔.
大家的名字.
認識大家.
很奇怪的.
有些事情就會在那種場景裡面發生.
發生什麼呢.
就是本身大家從小到大都覺得對方是世仇.
但你問他你認識對方的其中一些人呢.
他們是不認識的.

$^{721}$一直都是unknown的人.
只不過是他們想像裡面.
或者那些網路上.
或者那些傳媒上.
他們所見到的就是對方就是我們的仇人.
對方就是我們的世仇.
但是從來沒有見過大家.
從來沒有認識過大家.
這個Musalaha就是做這件事.
什麼叫Musalaha.
阿拉伯文就是reconciliation的意思.
復和的意思.
我們更加需要這件事.
而史唐傳第三章就是說.
我們的群體是一個怎樣的群體.
有沒有一些新手去拉一下外面的人.
在基督裡面我們是同一的.
去定睛看他.
去拉他進來.
新聞報導.
大家知道.
戰雲密報一觸即發.
這些就是我們平時社交網路上.
新聞報導上我們見到的相片.
我們全部都是見到這些.
但我們絕少會見到.
新聞報導會報導這些東西.
是有人的面.
是真實的人.
左邊一張是TVB.
被伊朗的火箭射過來.
然後炸了一些家園.
爸爸和弟弟.
猶太人.
丘相國在加沙.
加沙死了幾萬人.
大部分都是無辜的人.
我們在報章或者我們報導裡面.
加沙全部都是喊mask的人.
是不公平的.

$^{761}$我們一竿子打船人.
沒理由見到你這個樣子.
我們就當你是某某黨的人.
怎可以一竿子打船人.
右下角是伊朗.
被炸到家沒有家園的小朋友.
在這個時代裡面.
教會作為一個小的群體.
在基督裡面我們是一個怎樣的群體.
最近很多大學畢業.
我們芝加哥附近的西北大學.
畢業禮就請了一個人來演講.
那個人就是史蒂夫·卡瓦爾.
史蒂夫·卡瓦爾如果你看《辦公室》的話.
他是其中一個主角.
叫了他來講.
我覺得挺有趣.
聽他演講的字詞.
畢業禮.
畢業禮就是叫學生.
You make a difference.
你出到去Follow your dream.
他居然說了一件事.
Remember the little things.
Being kind.
And you are not alone.
Take care of one another.
拉一拉對方的手.
客氣一點對大家.
Be kind to one another.
這些很小的東西.
這就是史文第三章裡面講.
教會的那個Portrait.
教會的那個故事.
原來我們在基督裡面.
我們Do the little things.
Be kind to one another.
其實就是上帝的工作.
就是聖靈的工作.
所以我今天想remind你.

$^{801}$Be kind to one another.
你回到家裡.
不是伊朗對以色列.
你回到家裡.
不用用美國最大的那些幾萬磅的炸彈.
Just be kind.
伸出你的手.
定睛漢達有多少同理心.
這個就是在耶穌裡面.
在基督裡面有的群體.
所以最後我們講到Crossover.
講到跨越的時候.
我想做點功課.
這個星期.
你主動的去介紹一下你自己.
好不好.
去參加一個小組團契.
主動一點去參與.
花點時間去聽對方的故事.
花一點時間去有一點同理心.
為其他人去禱告.
好不好.
這個想一想.
有哪一位.
你可以主動的去關心一下.
可能是一個姐妹.
一個弟兄.
在教會裡面面熟口熟的.
但都不知道他什麼名字.
我這裡給大家一個准許.
如果你忘記了對方的名字.
不要緊.
今天當新人.
你說不好意思.
我忘記了你的名字.
讓這個群體成為一個.
真的有聖靈工作的群體.
一個彼此關注的群體.
去跨越.
甚至是不跟我們講同一個語言.

$^{841}$我都會去一些普通話的小組.
完全不同的人.
但是在基督裡面.
我們就應當有一個這樣的關係.
好不好.
可以同心得來基督.
主愛多謝你.
因為聖靈的工作.
在世上會在我們中間產生.
主愛你.
你在我們心裡面.
來觸碰和開啟我們.
叫我們的心靈.
叫我們的心能夠向他人來彈開.
在基督裡面.
在耶穌裡面.
我們能夠成為一個.
融會的教會.
融會的社群.
讓我們這個家庭.
能夠去接觸更多的人.
扶起其他的人.
定睛看著其他的人.
邀請其他的人.
讓這個群體能夠成為一個.
有愛的群體.
讓它成為一個.
我們每一個能夠彼此跨越的.
彼此跨越的一個生命.
讓教會能夠成為你自己的見證.
見證聖靈實實在在.
可以在我們中間工作.
主愛你求你使用我們.
讓我們彼此去交流.
彼此去認識.
彼此去建立.
以致這個家庭.
能夠顯出你自己的愛和恩典.
祈禱奉主名求.
好,元上帝祝福各位.

\newpage



\section{}
\label{sec:1fX6w4FAuCw}
\textbf{舊約網上行 - 上帝的足跡: EP12 款待與報喜}
\newline
\newline
連結: \href{https://youtube.com/watch?v=1fX6w4FAuCw}{\texttt{https://youtube.com/watch?v=1fX6w4FAuCw}} ~~~~ 語音日期: 2025-07-01
\newline
\newline
\hyperref[sec:cZT_1gxTxfc]{\small{< < < PREV SERMON < < <}}
~
\hyperref[sec:index_chronic]{\small{[返順時目]}}
~
\hyperref[sec:index_scriptual]{\small{[返順卷目]}}
~
\hyperref[sec:ot0FwdbiPyw]{\small{> > > NEXT SERMON > > >}}
\newline
\newline
$^{1}$我們今天來到《舊日網上》第十二集.
我們正在看亞伯拉罕的故事.
今天看到創世記第十八章.
這裡提到一個情節就是亞伯拉罕接待三位神的使者.
神好像遠行的旅客來到亞伯拉罕的面前.
亞伯拉罕去款待他們.
這個故事很有名.
在古今中外很多時候一些名畫都這樣刻畫.
去描述這個故事.
就是神去探訪亞伯拉罕.
亞伯拉罕去款待他們.
我們說讀聖經就有如去旅行.
我們能夠用我們現代的世界套上去.
一定要明白古代的世界.
說到客旅 說到遠行.
在古代因為交通不便.
天氣很多時候是耗熱.
或者路上有山賊.
所以旅行是非常危險的.
當時也沒有一些客棧或者旅館.
大致上都不會有.
所以遠行的風險就很大.
正因為這個緣故.
近東的文化就有一種特色.
叫做款待.
款待客旅.
Hospitality.
這些是文化上約定的禮節.
就是主客之間有某一種的預期.
某一種的規矩.
以致大家的關係可以維持得到.
即使我們今天也是.
就算我們今天去旅行.
或者我們用Airbnb.
主客之間雖然付了租金.
都有某一程度上的預期.
主人要提供什麼.
一個清潔的地方.
客人需要手藥.
不會弄壞東西等等.

$^{41}$如果不是的話.
大家的評價就要寫得很差.
或者我們想想這些禮節的時候.
習慣 習俗的時候.
我們想起我們文化上.
華人文化也有不少這些東西.
譬如舉個例子.
我們去別人的婚禮.
去別人的婚禮.
如果兩手空空 什麼都不帶.
這是一個羞辱.
為什麼呢.
因為你去一個婚禮.
習慣的習俗禮節.
就是你一定要給舊人情.
給多人情.
反過來說.
主人家去擺酒.
也不能擺得太差.
因為客人吃得不好的時候.
這也是一個羞辱.
這種情況.
我們可以很簡單的.
稱為Honor and Shame.
是一種榮辱的.
一種社會的狀況.
Honor and Shame.
近東的Hospitality款客.
也是這樣的.
有不同的元素在裡面.
未必是明文規定.
寫得很清楚.
但是大家的Expectations.
預期.
是弄得很清楚的.
有什麼元素呢.
第一 接待是家主的榮譽.
是天經地義的.
所以那個責任.
很多時候是一家之主.

$^{81}$是男人.
在城市裡面的男人.
在鄉村.
或者市鎮裡面.
那些人民的榮譽.
可以接待.
是榮譽.
可以接待一個人.
能夠保護他.
能夠提供他的需要.
陌生人來到你的門口.
來到你城鎮的城門的時候.
一定是帶著一些威脅.
所以藉著一個主人.
來邀請陌生人.
進入家裡的時候.
這個舉動.
就將這個本身的威脅.
變成友好.
主人家就會在自己的範圍裡面.
盡力去保護旅客的安全.
家主接待客人的時候.
通常會講清楚留多久.
但是完結之後.
也會誠意拳拳.
留下客人.
客人就需要給面子.
所以有所謂的extended stay.
客人接受了.
家主的邀請.
進入了別人家裡的時候.
就有一套.
不明的規矩.
譬如主人.
請客人進入家裡的時候.
就需要為他洗腳.
需要有膏油抹上.
很多時候.
主人口頭上會.
答應提供一些東西.

$^{121}$很少的.
但最後會大排筵席.
另外客人怎樣回饋呢?.
就是需要講出.
多謝.
去到讚賞主人.
從頭讚到尾.
特別當客人.
有神聖的地位.
先知,神的使者.
往往在這個時候.
就會帶來好的喜訊.
一些祝福的說話.
主人家服侍周到.
是近東文化的榮譽.
代表那個人.
有那個信特.
在聖經當中.
經常用.
hospitality(醫療).
去表達一個人的品格.
一個人忠誠於神.
這種生命.
譬如約伯記三十一章.
約伯.
堅持自己無辜.
他說旅客沒有在街上過夜.
我的門.
向過路的人是敞開的.
他堅持.
款客的文化.
hospitality code(醫療標誌).
如果我們這樣的背景.
我們就明白很多.
聖經的經文.
列出了幾段我們不會說.
但你用背景的時候.
你就會知道當中的情節.
為什麼會這樣發生.
例如創世記第24章.

$^{161}$以撒尋妻記.
有很多hospitality(醫療標誌)的情節.
事司紀第四章.
講到一個女人叫雅逸.
她暗殺了一個迦南將軍.
叫西西拉.
一個敵人.
在她的家裡面做的舉動.
另外第十九章.
就是基比亞的暴民事件.
有一個很恐怖的情節.
說到那些人.
不願意去行款待的文化.
甚至詩篇23篇.
敵人為我擺設筵席.
用油膏撩我的頭.
醫治我這個牧者.
款待我的牧者.
甚至到新約老家方第七章.
講到花里塞人.
來請耶穌吃飯的時候.
和他洗腳等等.
當然我們不能夠說聖經應用的時候.
我們就照字搬字.
古代的文化.
甚至今天中東有款客的文化.
所以我們就要照抄.
陌生人來敲門的時候.
我們一定要開門給他.
他要錢我們一定要給他.
當然不是這樣.
我們需要明白背景背後.
想表達一些真正的意思.
意思就是說.
一些被我們軟弱被我們脆弱的人.
我們要照顧他.
開放我們的生命.
以至能夠有神的祝福.
臨到我們.
今天我們看創世記十八章.

$^{201}$這裡就是款客的故事.
十八章第一節.
在萬里象樹那裡向阿布拉罕獻言.
那時是一天最熱的時候.
阿布拉罕正坐在帳篷的門口.
忽然看見三個人在他面前站著.
聖經首次.
詳細地描述.
一個款客.
阿布拉罕款待三個客人.
這個地點在哪裡呢?.
就是萬里象樹.
之前創世記有提過.
這是迦南山地南部.
希伯倫境內的其中一個地方.
希伯倫比較北一點.
Maui.
可能當時.
一直在迦南的宗教裡面.
在一棵很大的象樹那裡.
去敬拜一些神靈.
直到阿布拉罕來到.
希伯倫這個地方.
那裡發生什麼事呢?.
十三章第十八節.
他去到希伯倫萬里的象樹附近.
在那裡居住.
在那裡做什麼呢?.
為耶和華建造了一座祭壇.
將那地方稱為.
屬於耶和華顯現的地方.
說話的地方.
這是一個伏筆.
因為這個地方萬里希伯倫.
之後我們會發現.
上帝就是來審判索多瑪俄摩拉.
索多瑪俄摩拉.
正在這個地方的東邊.
走出一點點山地的時候.
就會看到四海平原一帶.

$^{241}$這個萬里象樹的地方.
到今時今日很多人都認為.
是一個希伯倫.
北邊一點點的一個地方.
叫做Marat al-Khali.
Khali就是阿拉伯文.
Friendship的意思.
Friendship Hill.
Friendship Heights的意思.
到今時今日.
很多猶太人很渴望.
久不久就去這個地方朝聖.
在這裡我們先處理一下.
一些神學的問題.
有三行客.
三個神的使者.
過去在猶太教和基督教的歷史.
很多時候都去猜想.
這三個天使是什麼身份.
拉比文學.
甚至給了名字.
一個是加伯列.
一個是米加勒.
一個是拉斐爾.
這個是聖經沒有提到的.
使者在兩者中間都被轉提過.
猶太教這樣揣測.
基督教的文獻.
也有這樣去猜想.
覺得這三個自然就是.
Trinity.
但是我相信.
我們不可以這樣看.
當時的人是不會這樣去理解.
因為當時.
交通不方便.
一國與另外一國.
來傳信的時候.
就一定要派出使者.
所謂的ambassador.

$^{281}$先知.
使者帶信息的messenger.
去到另外一個地方.
帶著君王的權柄.
說話的時候.
代表著他的君王.
所以很多時候.
說話不單用第三身.
也會用第一身.
所以我相信我們不必要.
將之後的教義.
譬如Trinity三位一體.
回到這一段經文.
十八章第二節.
亞伯拉罕看見三行客的時候.
就從帳篷的門口跑去迎接他們.
俯伏在地.
在亞伯拉罕整個故事裡面.
神親自向他顯現.
有好幾次.
這次是最後一次.
經文裡面怎樣提到亞伯拉罕.
在中東 在近東.
長老是穿著長袍的.
他們是不會跑的.
因為跑是有損尊嚴.
一定會慢慢走.
所以在這裡描述亞伯拉罕.
跑去俯伏在地迎接他們.
其實是一個很謙卑的態度.
第三節 亞伯拉罕說什麼.
他說 我主要約在你眼前.
蒙恩求你不要離開你的僕人.
讓我拿些水來為你洗洗腳.
在樹下掀一掀.
讓我拿些餅來給你吃.
可以恢復你的精力然後再走.
因為既然來到僕人的這處.
就應該是這樣.
這個就是款待客類.

$^{321}$拿水 洗腳 拿點餅.
發出邀請.
對方肯不肯呢.
這段聖經就繼續交代.
這三個人說.
他們說 就照著你所說的去做.
然後第六節.
亞伯拉罕就趕快的.
擺出一個大牌現職.
趕快進入帳篷.
跟撒拉說.
你快點拿三豆 細麵條和作奢餅.
然後就去牛群那裡.
把又嫩又好的牛仔.
交給他的僕人.
就趕快去預備盛宴.
然後就拿乳酪 奶.
預備好牛肉.
擺在客類的面前.
他們吃的時候.
亞伯拉罕在樹下.
侍候他們.
亞伯拉罕原先說什麼呢.
他說 拿點餅.
現在細節說什麼呢.
就拿細麵出來.
最好的麵粉.
然後就拿牛肉.
又嫩又好的.
乳酪 奶等等.
都是近冬的一些很好吃的東西.
然後就去做一個很豐富的一餐.
還做了侍應.
亞伯拉罕就成為一個.
款待客類的表表者.
這裡也是放下了一個伏筆.
因為下一次我們說到.
索多瑪和摩拉的時候.
會看到一個款待客類.
相反的故事.

$^{361}$現在我們這段經文的細節.
就表明亞伯拉罕是一個義人.
一個Righteous Man.
他的款待代表他的 Virtue.
歡迎客類.
打開他的生命.
誰知就迎接了上帝的工作.
上帝的祝福.
這個款待的習俗是怎麼樣的呢.
第九節繼續說.
他們問亞伯拉罕.
你的妻子撒拉在哪裡.
他回答.
在帳篷裡.
其中一位說.
到明年這個時候.
我必回到你這裡來.
那時你的妻子撒拉.
必生一個兒子.
這裡就說到客人的回饋.
說一些祝福的話.
這次真的帶來喜訊.
因為故事的發展.
亞伯拉罕慢慢慢慢就知道.
這三個旅客.
原來是神的使者.
甚至是耶穌說自己.
上一章第十七章.
神和亞伯拉罕立約.
用割禮為記號.
就已經預言了以撒會出生.
這次三行客來到亞伯拉罕的家.
特別的用意就是.
向撒拉作啟示.
向他報喜.
撒拉的反應是怎麼樣的.
十八章第十節.
撒拉在那人後面的帳篷門口.
也偷聽到聽見這句說話.
報喜的說話.

$^{401}$十一節.
當時亞伯拉罕和撒拉都老了.
上了年紀.
撒拉的月經也停止了.
撒拉心裡暗笑說.
我既已衰敗.
我主也老邁.
豈能有這喜事呢?.
中文翻譯也很暗笑.
沒有笑出聲.
在裡面笑.
這裡強調什麼呢?.
上文下理用了好幾次.
說他年老.
兩者都老邁.
甚至撒拉的月經已經停止了.
經文強調的是.
人已經走投無路.
人的無能.
特別是說到撒拉的時候.
根本是不可能再有生育的機會.
由亞伯拉罕的故事開始到現在.
就越來越強調這一點.
年紀老了.
這次甚至說得很清楚.
是她的經期都已經過了.
十八章十一節.
其實可以直譯成什麼呢?.
The way of women has ceased for Sarah.
可以這樣翻譯出來.
作為一個女人.
她生存之道都已經完蛋了.
不再能夠成為一個女人.
不能夠再生育.
誰知道上帝就帶來報喜.
神就看透人心.
十三節說什麼呢?.
對亞伯拉罕說.
撒拉為什麼會暗笑.
說我已經老了.

$^{441}$真的還能生育嗎?.
也說豈有難成的事嗎?.
到明年這個時候.
我會回到你這裡來.
撒拉必生一個意志.
又再報喜多一次.
去到預言多一次.
說清楚上帝的應許.
一定會按照神的話去兌現.
十五節撒拉聽見就害怕.
她否認可沒有笑.
但那個使者說什麼呢?.
他說你實在笑了.
這裡說到笑.
撒拉的笑.
可能裡面的意思是.
一笑置之.
聽到這樣的報喜.
我已經無能再生育.
怎麼可能發生在我身上.
這種人之常情.
相信撒拉當時.
可能也有這樣的想法.
很多年前.
幾乎是二十五年前.
亞伯拉罕七十五歲.
神呼召他們一家來到迦南地.
要成為大國.
應許他們生子.
但早就說了.
撒拉是不可能生育的.
想盡人的方法.
包括用下甲作代人母.
以實瑪利的出生.
當年可能撒拉也會有祈禱.
祈禱上帝行神蹟.
相信神.
認為神必定會按照他的方法去成事.
但這個時候說得很清楚.
他覺得連做一個女人也不容易.

$^{481}$十八章十四節.
「豈有難成的事.
到明年這個時候.
我要回到你的那裡.
那時撒拉必生一個兒子.
豈有難成的事.
困難難事」.
絕對有一個意思.
人也不能解決的難題.
撒拉不能生育.
人不能來到迦南地.
人不能解決.
但也可以有另一個意味.
譬如這個字.
難成的事.
難事.
在《出亞俱記》三十四章說.
「看!我要立約.
我要在你全體人民面前.
作奇妙的事.
在全地萬國之中沒有行過的.
在你四州萬民必看見的作為.
因為我向你行的是可畏懼的事」.
哪個字是難字呢?.
難成的事的難字.
就是奇妙這個字.
可以翻譯成wonderful.
過於人能夠想像.
很多時候在舊約聖經裡.
就是指上帝的救恩.
上帝成就的計劃.
他的工作.
是過於我們所能想像的.
但是神邀請我們去相信他.
相信他有這樣的能力.
能夠做到這件事.
說到撒拉的報喜.
不如我們看二十一章.
看一個結果.
一年之後.

$^{521}$所謂十月懷胎.
二十一章第一節開始.
耶和華照著以前所說的.
眷顧撒拉.
耶和華按著他的應許.
給撒拉作成.
亞伯拉罕年老的時候.
撒拉懷了人.
到了神應許的日期.
就給亞伯拉罕生了一個兒子.
亞伯拉罕給撒拉所生的兒子.
起名叫以撒.
亞伯拉罕的兒子以撒.
生下來的第八天.
亞伯拉罕就著神的吩咐.
給以撒行了國禮.
他兒子以撒出生的時候.
亞伯拉罕正值一百歲.
撒拉說:神使我歡笑.
凡聽見的也必為我歡笑.
又說:誰能告訴亞伯拉罕.
撒拉要與養兒女呢?.
在他年老的時候.
我竟給他生了一個兒子.
整段的經文和上一段的經文.
有一個呼應.
有一個keyword.
就是笑這個字.
事實上以撒這個字.
以撒其實就是笑的意思.
人想不到.
人以為盡頭的時候.
上帝的工作就想不到.
上帝救恩的計劃.
就不斷地向前推進.
本身人聽起來.
可能是一笑之詞.
現在就反轉變成喜樂的笑.
想不到上帝可以行人看為難的事.
行事奇妙.

$^{561}$是過於人的限制.
人的想像.
但是我們必須要謹慎.
當我們憑信心來接受.
神沒有難成的事的時候.
不等於我們人可以拍胸口.
向上帝發師號令.
可以去指揮上帝.
必須要做這件事.
必須要成這件難事.
或者一定要按我的手法去成事.
如果不是的話.
我就會放棄.
事實上有很多人都是這樣.
以前可能覺得.
神一定會做到很奇妙的事.
會成就一些我禱告的事.
神做不到.
我就放棄信仰.
這裡講得很清楚.
上帝做事往往超過你的想法.
講講這一段聖經.
上帝沒有難成的事.
神可以行一些過於人的限制.
過於人想像的事情.
來推進自己的計劃.
想一想我們自己.
今時今日.
你可能試一下重溫.
回頭看.
可能一年前.
或者在現在這個境地之前的你.
可能是五年十年甚至二十年之前.
或者你曾經患過大病.
進過手術室.
現在出來.
好像自己走到盡頭.
但今時今日你看回.
上帝是否真的行了祂的旗幟.
或者這樣說.

$^{601}$當時你在困境裡.
當時你在沒有盼望的時候.
如果當時有人向你報喜.
告訴你神一定會帶你經過.
或者當時你在手術室.
你去禱告的時候.
當時有人向你講一個訊息.
告訴你神一定會拯救你.
你可能在當時難處的裡面.
困苦的裡面.
骨頭老的裡面.
可能你都會一笑置之.
我們的不順.
我們不能超越我們的想像.
今時今日看回.
原來神真的可以行到奇妙的事.
是沒有難成的事.
最近我留意到一個姐妹的分享.
這位姐妹在基督教的圈子裡面都很出名.
很多人都知道她的故事.
她的名字叫Joni Elson-Tard.
這是殘障的坐輪椅的姐妹.
很多時候做很多的教導工作.
鼓勵的工作.
特別是在一些殘障人士中間.
有很多很好的事工.
留意她分享的原因是.
去年她患上了新冠病毒.
因為新冠病毒的緣故.
所以她有很多的掙扎.
然後得到醫治.
她就出來分享.
Joni Elson-Tard.
她的故事其實很多人都知道詳情.
十七歲的時候.
她因為跳水的意外.
肩膀以下癱瘓.
很多時候我們留意她的故事.
我們都很少留意一點.
就是她的掙扎.

$^{641}$她當時的掙扎.
因為她跳水意外癱瘓之後.
她就不斷去禱告.
其他人都為她禱告.
就是神一定有能力來醫治她.
再神會有難成的事.
再神一定有奇妙的作為.
顯出她的榮耀.
神肯定可以做到.
上帝如果按照祂所願的.
醫治她癱瘓.
上帝一定會得到榮耀.
也有很多人認為.
上帝在她身上真的有很大很奇妙的工作.
但一直祈禱久了.
她就進入一個更大的掙扎.
掙扎就是.
如果神真的不按照她的禱告來行事.
那神會做些什麼?.
神是否真的可以不醫治我.
在我這個殘障的人身上.
能夠顯出祂的作為?.
有沒有這樣的可能?.
這就是一個最大的掙扎.
上帝怎麼可能超越祂這個殘障的限制.
而可以顯出祂的聲望.
顯出上帝的榮耀?.
誰知這麼多年.
上帝就不斷地使用她.
現在她已經七十多歲.
神就是用一個坐在輪椅上.
五十多年的人.
去成就一些我們健全的人.
都不能夠做到的事情.
比我們還要多幾十倍.
這就是上帝沒有難成的事.
上帝當然有能力.
帶領我們離開現在的困境.
可以奇妙到這樣的地步.
但同時上帝也可以藉著我們現在有的限制.

$^{681}$仍然可以用到我們.
仍然可以顯出祂的能力.
你看看祂終身殘障.
有什麼神的榮耀可言?.
誰知再神沒有難成的事.
就是藉著鍾尼亞瑞森塔德.
祝福了千萬殘障的人聽到福音.
有很多身邊不同的人.
就算健全的人都得到幫助.
能夠聽到福音,聽到神的大能.
問題就是這樣.
有兩個選擇.
第一個就是神當初真是神跡地醫治他的癱瘓.
這是難事.
還是神不醫治他的癱瘓.
在他癱瘓當中成就他自己的計劃.
哪一樣是更加奇妙?.
當然就是後者.
再神沒有難成的事.
撒拉阿布拉罕就是經歷這件事.
其實寫這個故事就是要告訴阿布拉罕的子孫.
我們這些恩信稱義的阿布拉罕的子孫.
告訴我們再神會有難成的事.
但我們一定要搞清楚.
上帝不一定要按照我們的時間.
按照我們的方法去做事.
反而可能是邀請我們.
在我們現在的限制裡.
要見到上帝仍然可以行事奇妙.
鍾尼亞瑞森塔德有他在天國裡的角色.
當然我們不需要照辦主管.
但最重點是甚麼?.
就是上帝不是要我們做同一件奇妙難成的事.
而是上帝要在我們現在這個崗位.
去體現出他自己的奇妙.
讓我們不要限制上帝的工作.
讓我們認定再神會有難成的事.
再神有奇妙的事.
而且奇妙的事是過於我們所能想像.
所能估計到的.

$^{721}$邀請就是我們願意開放我們的生命信心.
迎接上帝進來.
迎接上帝的訊息.
迎接神進入我們的家.
進入我們的生命.
隨著主的旨意.
知道在上帝的眼中.
知道在神是沒有難成的事.
上主我們感恩讚美你.
因為在你沒有難成的事.
你行事,你的作為.
是過於我們所想,所求的.
主我們現在無論在任何的境況當中.
即使是我們社會裡面.
不同的不安的境況.
或者我們人生裡面遇到一些大的挑戰.
在你的境況當中.
求你自己將我們的信心.
定精在你身上.
要我們看到在你沒有難成的事.
你可以按照我們的祈願.
我們的祈求來行事.
而顯出你的奇妙.
亦都可以不按照我們的心願.
來行出你的奇妙.
求主你幫助我們.
我們的生命是向你敞開的.
去迎接你,款待你進來.
以致你能夠藉著我們.
顯出你自己的榮耀.
我們這樣祈禱奉主命救.
\newpage



\section{路加福音 10:25-37}
\label{sec:ot0FwdbiPyw}
\textbf{跨越界線的愛 (路加福音 10\_25-37) - 梁紫健傳道}
\newline
\newline
連結: \href{https://youtube.com/watch?v=ot0FwdbiPyw}{\texttt{https://youtube.com/watch?v=ot0FwdbiPyw}} ~~~~ 語音日期: 2025-07-01
\newline
\newline
\hyperref[sec:1fX6w4FAuCw]{\small{< < < PREV SERMON < < <}}
~
\hyperref[sec:index_chronic]{\small{[返順時目]}}
~
\hyperref[sec:index_scriptual]{\small{[返順卷目]}}
~
\hyperref[sec:EEmGj96kwiU]{\small{> > > NEXT SERMON > > >}}
\newline
\newline
路加福音 10:25-37
\newline
\begin{longtable}{cl}
\hline
\hline
章節 & 經文 (和合本修訂版)\\
\hline
10:25 & \begin{tabularx}{0.7\textwidth}{X} 有一個律法師起來試探耶穌，說：「老師！我該做甚麼才可以承受永生？」 \end{tabularx} \\ \\ \relax
10:26 & \begin{tabularx}{0.7\textwidth}{X} 耶穌對他說：「律法上寫的是甚麼？你是怎樣念的呢？」 \end{tabularx} \\ \\ \relax
10:27 & \begin{tabularx}{0.7\textwidth}{X} 他回答說：「你要盡心、盡性、盡力、盡意愛主—你的神，又要愛鄰如己。」 \end{tabularx} \\ \\ \relax
10:28 & \begin{tabularx}{0.7\textwidth}{X} 耶穌對他說：「你回答得正確，你這樣做就會得永生。」 \end{tabularx} \\ \\ \relax
10:29 & \begin{tabularx}{0.7\textwidth}{X} 那人要證明自己有理，就對耶穌說：「誰是我的鄰舍呢？」 \end{tabularx} \\ \\ \relax
10:30 & \begin{tabularx}{0.7\textwidth}{X} 耶穌回答：「有一個人從耶路撒冷下耶利哥去，落在強盜手中。他們剝去他的衣裳，把他打個半死，丟下他走了。 \end{tabularx} \\ \\ \relax
10:31 & \begin{tabularx}{0.7\textwidth}{X} 偶然有一個祭司從那條路下來，看見他就從另一邊過去了。 \end{tabularx} \\ \\ \relax
10:32 & \begin{tabularx}{0.7\textwidth}{X} 又有一個利未人來到那裡，看見他，也照樣從另一邊過去了。 \end{tabularx} \\ \\ \relax
10:33 & \begin{tabularx}{0.7\textwidth}{X} 可是，有一個撒瑪利亞人路過那裡，看見他就動了慈心， \end{tabularx} \\ \\ \relax
10:34 & \begin{tabularx}{0.7\textwidth}{X} 上前用油和酒倒在他的傷處，包裹好了，扶他騎上自己的牲口，帶他到旅店裡去，照應他。 \end{tabularx} \\ \\ \relax
10:35 & \begin{tabularx}{0.7\textwidth}{X} 第二天，他拿出兩個銀幣來，交給店主，說：『請你照應他，額外的費用，我回來時會還你。』 \end{tabularx} \\ \\ \relax
10:36 & \begin{tabularx}{0.7\textwidth}{X} 你想，這三個人哪一個是落在強盜手中那人的鄰舍呢？」 \end{tabularx} \\ \\ \relax
10:37 & \begin{tabularx}{0.7\textwidth}{X} 他說：「是憐憫他的。」耶穌對他說：「你去，照樣做吧！」 \end{tabularx} \\ \\
[1ex]
\hline
\hline
\end{longtable}
$^{1}$弟兄姊妹早安.
早安.
記得幾年前在一間教會實習.
我和一個傳道人帶著一班弟兄姊妹去台灣短訓.
其中一站我們就去這裡.
去這些食堂跟人們聊天報道.
我們幾個人一組,三四個人一組.
我們想著有一個辦法,不過這個是糞便的辦法.
我們就拿著大家看到我的左上角的鑰匙.
其實是五色珠,代表著不同的救恩狀態.
如果大家有用過五色珠應該會知道,去報道.
我們就這樣,拿著鑰匙扔在地上,踢過去.
我們想著…你先幫我保持著鑰匙.
這麼大的動作,台下或附近的年輕人應該會幫你撿.
我們看到一個年輕人很興奮地在打遊戲機.
我們就想著順理成章,把鑰匙踢過去.
年輕人,我叫Ken,謝謝你幫我們撿起這個鑰匙.
你有興趣了解一下這把鑰匙的意思嗎.
誰知道這個年輕人一點都不理睬我們.
我說國語已經算很標準了,對吧.
他一點都不理睬我們,不斷地在打遊戲機.
我猜新的世代可能也是這樣.
然後我們就很尷尬,你知道我們把鑰匙踢出去.
怎麼辦呢?我們唯有死氣沉沉地走到台下.
把鑰匙撿起來.
我們就眼尾一看,旁邊有沒有其他人.
有沒有其他目標.
我們就看到這個,這是AI圖.
有一個很紋身大漢,全身都是汗.
不停地自言自語,很暴躁.
你一直看著我,我就看著你,去吧.
大家都覺得這個男士有點難接觸.
表面上他很粗獷,很不可接觸.
也令我們覺得我們跟這個男士當中.
好像有一個很大的鴻溝.
好像很難攻破的界限.
我們也打消了接觸這個男士的念頭.
誰知有一組弟兄姊妹,我們分了不同的組.
竟然跟這個男士聊天,很支持地坐下來聊天.
起初那個男士也不太理會他.

$^{41}$可能我們說的國語不太好,不太理會他.
我們就算了,找其他人,去其他不同的地方聊天.
半小時後,我們回去,我們還看到.
竟然這一組弟兄姊妹跟紋身叔叔聊得很開心.
而那個紋身叔叔聊著聊著,竟然還流眼淚.
兩行眼淚不停地流.
後來我在beefing session聽那些組員說.
其實這個紋身叔叔遇到很多挫折.
以前也拜過很多不同的神.
你也知道台灣有很多不同的偶像.
最後他竟然決定要信耶穌.
這個例子說的是,我們有時一些自設的界限.
以為紋身叔叔不是一個可接觸的對象.
我們很自然就會劃了一條界.
他是不可接觸的.
我們很自然就會被自己腦裡的一些界限限制住.
以致我們不能見證上帝的恩典在人的身上.
最後那個叔叔也信了耶穌,是吧.
但是我們這一組人就沒辦法見證上帝的愛在人的身上.
另一組人也可以.
今天想說這個例子,做開頭想跟大家說.
有時有些界限是不必要的.
以致我們有些想法,甚至我們有些信仰.
我們的概念,我們覺得自己孤人一等.
我們沒辦法見證上帝在人身上的大能和愛.
這是一個反映的例子,反映著一些不必要的界限.
所以今天的題目我想跟大家說.
跨越界限的愛,希望我們聽完這篇導讀會知道三點.
第一點就是,一些不必要的界限是如何阻止我們見證上帝的.
第二,是為甚麼會有這些界限.
當知道這些界限之後,第三,我們如何跨越這些界限呢.
我們如何可以實踐上帝的愛呢.
如果準備好,我們就一起來.
想跟大家說,今天的經文是三馬里亞人的比喻.
三馬里亞人的比喻大家一聽就已經很熟悉.
其實它是聖經在《路加福音》裡面有獨有的記載.
開頭和馬太和馬賀有些相似,都有關於誡命的真論.
馬太福音版本是記載著一個律法師要試探耶穌.
就說,夫子,律法書上的誡命哪一條是最大的.
馬賀的版本就是文士聽到有關復活的真論.

$^{81}$就見耶穌答得好,就追問他.
究竟律法,誡命哪一條是最重要的.
雖然對觀福音也有一個很相似的演繹.
但是《路加福音》最不同的一點是.
它用了三馬里亞人的比喻去詮釋整個律法的爭議.
所以我今天選的這段經文.
三馬里亞人的比喻是在說一個釋法的爭議.
在說耶穌利用這個比喻.
不單只是在說誰是誰,論社的問題.
而是進一步詮釋什麼叫愛人如己.
說出答案給大家聽.
愛人如己是要突破自己的界限.
我們一起來看看.
《路加福音》十章,25至28節是這樣說的.
有一個律法師起來試探耶穌.
他說,老師,我怎樣做才可以承受永生.
耶穌對他說,律法社是怎樣讀的.
他回答,你要盡心盡性盡力.
盡義愛主,愛主你的神.
又要愛論如己.
耶穌對他說,對,你讀得對.
你這樣做就是永生.
那個人就要證明自己有理.
對耶穌說,誰是我的論社呢.
律法師是什麼人呢.
律法師是述讀律法.
他不單止是讀律法.
他還要是全息教導律法.
身分就好像有些像實行醫老師教會顧問那樣.
他問承受永生的問題.
很明顯是想挑戰耶穌的權威.
他想挑戰一下,耶穌是否懂得這些律法呢.
所以耶穌就反問他,律法怎樣寫呢.
他又答不到兩段舊約.
申命記和利未記.
愛主你的神,又要愛論如己.
這種猶太串珠式的答案,一定是對的.
不過他要證明自己有理.
就問了一個問題,就是誰是我的論社.
就揭示了他根本不明白愛論如己是什麼.

$^{121}$更加不明白為什麼耶穌要用撒瑪利亞人的比喻.
去諷刺這個律法師.
你要盡心盡性盡力盡義愛主你的神.
其實是出自於申命記六章四節.
基本上是猶太人.
你問猶太人也好,他就會知道.
有時如果你看到他們在門尾.
可能還會弄了一個六章四節在那裡.
是一個遵行律法的大原則.
總之是做什麼都好.
你去到盡力盡力地愛主.
但實際上做什麼呢.
是做什麼呢.
是不是等於我們24小時不睡覺就叫做愛主.
很盡力吧.
所以又要愛論如己.
就說了一點具體的東西.
這裡又要其實是一個連接詞.
你可以翻譯為同賣或者就是.
我個人覺得如果律法師提出這個律法的大原則.
將愛論如己就等於愛主你的神的做法.
其實是合理的.
因為愛論如己是一個具體的做法.
所以我們可以這樣理解.
愛主你的神就等於愛論如己.
如果這樣看的話.
我們就要問一下.
究竟什麼叫做愛論如己.
愛論如己出自於利未記.
利未記這裡其實很長的.
講了很多律例典章.
不可恨你的弟兄要指責你的鄰舍.
免得他擔當罪過.
不可報仇也不可埋怨你本國的子民.
你要愛論如己.
我是耶和華.
你要愛他如己.
我是耶和華.
這裡很有趣.
為什麼要將愛論如己和我是耶和華放在一起.

$^{161}$為什麼.
不好意思 我的稿子沒了.
可不可以給我手機.
愛論.
謝謝.
為什麼要將愛論如己和耶和華放在一起.
其實其中一個原因就是因為.
愛論如己和我是耶和華.
在利未記的context裡面.
耶和華就是救你們出埃及的上帝.
當在埃及的境地裡面.
你會發現有很多人跟你一起是寄居.
也有做過奴隸.
所以當不同的人在你當中的時候.
你會發現以色列人不能只顧自己.
他要做好愛論.
做好你身邊不同的人.
你都當他是自己一樣.
就好像.
利未記的十九章是在一個神聖法典裡面.
目的是什麼呢.
目的是教導猶太人.
過一個如何奉獻給上帝.
做一個聖潔的民族的人生.
愛論如己第一個實踐是什麼.
就是將自己奉獻給上帝.
歸入了神的子民.
成為了神聖潔的一部分.
所以愛論如己第一個實踐是什麼.
實踐就是你要愛你身邊的不同族的人.
包括不同支派的弟兄.
大家知道有十二支派吧.
但是很可惜的是.
不知道大家看不看得到.
從創世之後.
人因為犯罪.
你會發現有很多兄弟相爭的事件.
是該因殺阿伯.
閃,阿忽,嘉楠的後裔相爭.
還是索多瑪,俄摩拉攻擊羅德.

$^{201}$阿哥欺騙哥哥的長子名分.
約瑟被哥哥賣到埃及.
含的後裔被埃及人奴役.
這些都是.
本身是什麼.
本身是同族人.
本身是兄弟.
卻兄弟相爭.
互相殘殺.
所以你會明白為什麼你沒叫我這樣說.
恨你的弟兄就有罪.
為什麼要當外人是自己人.
就因為本身.
大家都是來自同一個上帝.
大家都有責任.
去到和睦同居.
管理大地.
所以愛論如己的實踐就是.
我們願意回到創造一開始的起初.
告訴別人.
你和我都是來自同一個上帝.
大家就算有不同的支派.
不同的民族.
大家都是自己人.
這裡的螢幕可能有點小.
想和大家看的這件事就是.
《國家福音》這裡.
撒瑪利亞人的比喻.
不單止是比喻這麼簡單.
其實是在映射一個兄弟相爭的故事.
你會發現.
大家有不同的字眼.
我是highlight了的.
有共同的地點.
有共同的動作.
有共同的意志行為.
你就會發現.
在《國家福音》裡面.
我先讀一讀.
耶穌回答說.

$^{241}$有一個人從耶路撒冷下到耶利哥去.
落在強盜手中.
他們把他奪去衣裳.
把他打得半死.
雕下他走了.
有一個祭司從這條路下來.
看見他就另一邊去了.
有一個利未人來到這裡.
看見他照樣也從另一邊去了.
那時有一個撒瑪利亞人路過那裡.
看見他就動了慈心.
上前用油和酒倒在他的傷處.
包裹好尿.
扶他上自己的身後.
帶到旅店去照應他.
第二天拿出兩個錢幣來.
交給店主說.
請你照應他.
額外的費用我會來還給你.
你想這三個哪一個是落在強盜手中的呢?.
他說是憐憫他的.
耶穌對他說.
你去照樣做吧.
你會看見有同樣的人物.
撒瑪利亞人.
即是北國的以色列人.
亦有南國猶大人.
祭司和當時利未人.
跟著猶大人的利未人.
還有行動拿出衣服來.
讓他們吃喝.
上前包裹.
讓傷者得到照顧.
令他們回到原本的地方.
亦有共同的地點.
撒瑪利亞.
耶路撒冷和耶利哥.
所以大家會發現.
這件事不單只是一個比喻.
它是影射一個兄弟相殘.

$^{281}$在歷代志下的一個地方.
因為時間的問題.
我就不再讀另一篇經文了.
你就會發現無論.
以色列人.
北國的以色列人.
還是南國的猶大人.
聽到撒瑪利亞人的故事.
就會自然想起.
歷代之兄弟相殘的故事.
就好像大家會.
這句歌詞.
Thanks, thanks, thanks, thanks, Monica.
會想起什麼呢.
對,瑞蘭代替你的地位.
當時的猶太人.
一聽到耶利哥.
一聽到打傷他.
就會想起原來.
我們兄弟相殘的故事.
就像我們想起歌詞一樣.
代之很有趣.
他的主題是希望十二支派回歸.
希望南北國可以兄弟和睦共處.
回到創世的時候.
大家就好像.
創世的時候.
來到上帝的面前.
一同敬拜上帝.
但是很可惜的是.
南國的利未人和猶大人.
在耶路撒冷.
他們敬拜上帝.
覺得自己就是祭司的主理人.
九月九支就覺得自己才是名門正派.
跟其他人不同.
在頭腦上就有你猶大人.
我以色列人很不同的分別.
猶大人是優如蝶.
我們有整個祭司傳統.

$^{321}$我們有很好的敬拜.
我們有很好的聖殿.
是敬拜的中心.
唯有猶大人才可以去到上帝面前敬拜.
你們不過人.
只可以在休壇敬拜.
所以回到路加福音.
三瑪利人的比喻.
祭司和利未人.
兩個都是在聖殿裡面服侍的人.
他們很應該明白什麼叫愛倫如己.
就是要跨越自己的界限.
但是他們卻被自己所限制.
沒辦法做到上帝所託付的事.
很可能是.
聖經這裡沒有寫.
不過他很可能因為利未記的寄靜條例而受到限制.
為什麼呢.
留意.
落在強盜手中的人是怎樣.
被人打到半死.
即是說他即將死亡.
或者有機會.
祭司可能在想.
他們是不可以碰死屍的.
為什麼呢.
因為在寄靜條例裡面是這樣的.
不得因普通人的死屍而玷污自己.
只能為近親哀悼而接觸屍體.
即使如此.
接觸之後都仍然要結靜.
祭司可能是因為這條條例而有一個界限.
我不可以做這件事.
我不可以碰他.
我碰到這條死屍的話.
我就回不了聖殿.
我做不到事.
大家覺得祭司對不對.
其實很合理.
因為他要回聖殿工作.

$^{361}$大家還記得愛倫如紀的律法是怎樣.
要當對方是自己人.
自己人即是近親.
是兄弟姐妹.
所以如果那個人真的死了的話.
祭司仍然可以進入聖殿.
只需要做一些寄靜條例.
但他們沒有這樣做而是離開了.
但真正做到愛倫如紀的是誰呢.
是撒瑪利亞人.
是一個在猶太人眼中的混種.
不純正.
被人鄙視.
被人小看的一個撒瑪利亞人.
撒瑪利亞人竟然是他.
做到了整個祭司傳統的核心.
跨越了自己的界限.
當對方是自己人.
照顧他很細緻.
是不是很諷刺.
情況就好像我們基督徒信了久了.
都會有類似的憂鬱感.
就好像有些非信徒是比較次等的.
我們才是上帝的子民.
每個星期上教會.
上帝會愛我多一點.
覺得其他人在上帝眼中不值一提.
路加福音第九章也是這樣說的.
門徒都有這樣的心態.
當耶穌登山變相.
趕鬼之後.
門徒個個都搶著做大佬.
又是因為什麼呢.
撒瑪利亞人不肯接待門徒.
接著門徒想怎樣.
耶穌不如我們用火燒死他們.
門徒這種心態就是.
我們才是上帝的子民.
其他人都是不值一提.
燒死他們沒問題吧.

$^{401}$是不是很恐怖.
我猜現在發動戰爭的人.
都有這樣的心態.
愛倫如己就是要跨越這種心態.
自己比別人高尚的這種心態.
這種心態的危險之處就是.
覺得我們很多事情都與我們無關.
這件事不是我跟開.
你去找誰誰吧.
我這個身分.
我怎會幫你做這些事.
你先過著吧.
說出來很難聽.
但有時候我們心裡可能是這樣想的.
我們首先要留意.
這個心態令我們產生了一些.
不必要的限制.
以致我們心底裡都覺得.
我們比其他人優越.
在心態上已經不想留意其他人的需要.
所以當明白了這些界限.
會令我們沒辦法做到上帝託付我們的事.
我們就要跨越這個界限.
這也是為什麼耶穌做了這個跨越.
在路加福音撒瑪利亞人的段落之前.
耶穌是這樣評價的.
我告訴你.
從前有很多先知和君王.
看你們要看的,但沒看見.
聽你們要聽的,但沒聽見.
耶穌令先知和君王都看不見.
為什麼呢?.
因為他們本身是教導律法的.
但他們不明白律法的意義.
唯有耶穌給這班很尊貴.
高人一等的門徒有機會明白.
耶穌就馬上做了一件事.
就是馬大瑪利亞的故事.
在第十章.
做了一件跨越界限的事.

$^{441}$就是讓瑪利亞在她的腳前聽到.
為什麼這麼突破呢?.
因為當時拉比和女性不可以這麼近接觸.
大家如果有去過以色列哭場都知道.
男女是分開的.
耶穌就用了自己的榜樣.
去跨越了這個界限.
告訴門徒.
她和身邊的人都是一樣.
她當身邊的人都是自己人.
剛好在28天之前.
就是690個小時之前.
我和Alice在高速公路上.
被一輛中型拖車追尾.
當時在塞車.
我們車的後尾箱和後座都撞到摺了.
而這輛車因為後面撞得很厲害.
我們整輛都飛向前撞到前面的車.
所以我們就夾在車子中間.
像三文治一樣.
很感恩上帝保守我們的生命.
讓我們從一堆爛鐵走出來.
這件事發生之後.
我真是很真實地去感受到.
頂梓梅做愛倫如己的這件事.
她當我們是自己人.
當我們和大家一樣.
立即來到現場救我們.
幫我們看保險.
載我們去急症室.
在急症室等了一整晚.
大家很快地出現回應和慰問.
怕我們不好意思.
偷偷地幫我們籌款.
偷偷地塞錢給我.
怕我們沒有錢看醫生.
幫我們做物理治療.
送湯水給我們喝.
幫我們買菜.
和我們一起吃飯.

$^{481}$一邊打電話,一邊發訊息.
各樣的問候.
借車子給我們.
還幫我們看不同的車子.
我真的很感受到.
頂梓梅真的當我們是你們自己一樣.
還要跟我們說.
不要選這麼舊的車子.
選好一點,安全一點.
車子是比較硬的.
你們可以選擇,不用擔心.
大家不是說.
你先找一輛車子開.
是說,如果我是你.
如果我是你,我會選一輛好一點的.
你們可以選擇,不用擔心.
這是真心地將別人.
將其他人身邊的人.
當作是你自己的心態.
愛倫兒紀是怎樣的.
就是一個真誠地認同.
其他人都跟自己一樣.
以至跨越了我們本身的界限.
進入別人的處境當中.
我很開心.
我可以被大家當作自己人.
你也知道.
大家跨越了很多不同的界限.
很多不同的限制.
我也很期待.
教會跟我們社區的每一個人.
都能夠可以實踐這件事.
跨越自己的界限.
所以總結.
是甚麼不必要的界限.
令我們阻止我們見證上帝呢.
就好像一開始,我講我的經歷.
我們有些自恃的界限.
我們覺得叔叔不可接觸.
我們覺得叔叔不接觸.

$^{521}$但是就令我們無法見證上帝.
在叔叔的身上.
存留的恩典.
為甚麼會有這樣的界限.
可能我們心底裡.
我們都有些無法言說.
我們覺得可能我們自己.
比另外身邊的人身分高尚.
當我們知道了這樣的心態的時候.
我們要跨越它.
認清我們自己跟其他人一樣.
都是上帝的子民.
沒有高低之分.
我們要做的行動就是.
當我們當對方是自己人的話.
我們經常都為自己想,是嗎.
我們要想一下.
在他的處境想一下他有甚麼需要.
這個是愛倫如記的.
所以作為家長.
作為別人的親戚,作為別人的子女.
你在職場上,你在社區裡面.
你跟身邊的人有甚麼界限呢.
你願意突破這些界限嗎.
我們邀請大家一起去到永愛.
跨越這個界限.
其實我還有一個邀請.
就是早前我認識一個同工.
在大學裡面做福音士工的前輩.
他提過我很多東西.
他說其中一樣東西就是.
我們以為福音士工就是.
叫學生來我們當中.
叫他們跨越他們的界限.
來我們當中.
但我們有沒有想過他們的處境呢.
我們有沒有想過我們關注一些東西.
但他們也有自己的關注,是嗎.
他們有學業,他們有自己的生活.
我們有沒有想過要跨越自己的界限.

$^{561}$進入學生生活的場景.
他們是怎樣生活的.
他們在關心甚麼.
而不是因為我們想叫他們來我們當中.
當然最後是好的,他們來我們當中.
其中一個應用就是.
我們打算進入大學裡面.
我們跟他們吃午餐,我們搞團契.
這是一個很實際的邀請.
如果有弟兄姊妹.
你的家很近大學.
而你是很想跨越你的界限的話.
讓我們一起在你家做同工.
你聯絡我.
我很想跟大家一起跨越這個界限.
或者你有很好的策略.
你可以跟我說.
我們一起實踐這件事.
我們一起祈禱.
天父,謝謝你讓我們有這個空間.
學習你的話語.
縱使有很多時候我們知道.
我們自己有很多不濟.
也有很多限制.
令我們不能夠跨出去.
在別人的處境為別人著想.
但主啊,你自己的榜樣告訴我們.
你也是這樣做的.
求你讓我們實踐.
你幫助我們.
我們知道有些界限.
在這一刻我們跨不過.
但求主你幫助我們.
求你幫助我們.
我們今禱告.
奉祂為耶穌得勝的名頭.
法定人數不足.
\newpage



\section{使徒行傳 28:1-10}
\label{sec:EEmGj96kwiU}
\textbf{接待的福音(Be Our Guests) (使徒行傳28\_1-10, 30-31) - 朱莉傳道}
\newline
\newline
連結: \href{https://youtube.com/watch?v=EEmGj96kwiU}{\texttt{https://youtube.com/watch?v=EEmGj96kwiU}} ~~~~ 語音日期: 2025-07-09
\newline
\newline
\hyperref[sec:ot0FwdbiPyw]{\small{< < < PREV SERMON < < <}}
~
\hyperref[sec:index_chronic]{\small{[返順時目]}}
~
\hyperref[sec:index_scriptual]{\small{[返順卷目]}}
~
\hyperref[sec:_p6p1yxlKsI]{\small{> > > NEXT SERMON > > >}}
\newline
\newline
使徒行傳 28:1-10
\newline
\begin{longtable}{cl}
\hline
\hline
章節 & 經文 (和合本修訂版)\\
\hline
28:1 & \begin{tabularx}{0.7\textwidth}{X} 我們既已獲救，才知道那島名叫馬耳他。 \end{tabularx} \\ \\ \relax
28:2 & \begin{tabularx}{0.7\textwidth}{X} 當地人非常友善地接待我們；因為正在下雨，天氣又冷，他們就生了火歡迎我們眾人。 \end{tabularx} \\ \\ \relax
28:3 & \begin{tabularx}{0.7\textwidth}{X} 那時，保羅拾起一捆柴，放在火中，有一條毒蛇，因為熱的緣故鑽了出來，纏住他的手。 \end{tabularx} \\ \\ \relax
28:4 & \begin{tabularx}{0.7\textwidth}{X} 當地的人看見那毒蛇懸在他手上，就彼此說：「這人必是個兇手，雖然他從海裡獲救，天理仍不容他活著。」 \end{tabularx} \\ \\ \relax
28:5 & \begin{tabularx}{0.7\textwidth}{X} 保羅竟把那毒蛇甩在火裡，並沒有受傷。 \end{tabularx} \\ \\ \relax
28:6 & \begin{tabularx}{0.7\textwidth}{X} 當地的人想他快要腫起來，或是忽然倒下死了，但等了好久，見他沒有甚麼異樣，就轉念說他是個神明。 \end{tabularx} \\ \\ \relax
28:7 & \begin{tabularx}{0.7\textwidth}{X} 離那地方不遠有一些田產，是島長部百流的。他接納我們，盡情款待了我們三日。 \end{tabularx} \\ \\ \relax
28:8 & \begin{tabularx}{0.7\textwidth}{X} 當時，部百流的父親臥病不起，患了熱病和痢疾。保羅進去見他，為他禱告按手，治好了他。 \end{tabularx} \\ \\ \relax
28:9 & \begin{tabularx}{0.7\textwidth}{X} 從此，島上其餘的病人也都來，得了醫治。 \end{tabularx} \\ \\ \relax
28:10 & \begin{tabularx}{0.7\textwidth}{X} 他們又多方面尊敬我們，到了開船的時候，又把我們所需用的東西送到船上。 \end{tabularx} \\ \\
[1ex]
\hline
\hline
\end{longtable}
$^{1}$(終於有機會見到大家).
終於有機會見到大家.
讓我在疫情期間有機會在網上來服侍大家.
但其實沒有東西可以替代人與人之間的真實交流和接觸.
今天思考要和大家分享什麼經文的時候.
我就想起這套電影.
我想看大家有多少人認得這個片段呢.
如果舉手的應該有差不多一年多.
因為這套片段是2000年的片段.
總共是25年前.
以後要小心聽到這些問題.
要想清楚才能回答.
這套電影叫做《Cast Away》.
中文譯作《劫後重生》.
是當時其中一部最賣座的電影.
電影描述哪個呢.
就是暴露年齡系列.
主角是誰呢.
是湯漢斯.
他飾演一個聯邦快遞員工.
做這些聯邦快遞員工.
他的時間觀念是非常之重.
而且他擁有一個快樂的人生.
他本來打算和女友求婚之後.
就可以快快樂樂地生活下去.
誰知道在聖誕晚餐的時候.
突然收到公司急求.
要出一個緊急的商務旅遊.
要出門.
誰知道在商務旅遊.
飛機在南太平洋失事墮海.
整個飛機只有他一個生還.
然後漂流到一個無人的荒島.
所以這套電影是典型的.
魯賓遜求生記.
那些荒島災難片.
講述人怎樣在絕地求生.
不要失去希望等等.
但我今天不是想講.
你怎樣去荒島求生.

$^{41}$怎樣鑽木取火.
但我在想一個城市人.
怎樣在野地生存.
如果你是湯漢斯.
你會怎樣生存呢.
影片很聰明的設計是什麼呢.
我說他是什麼員工.
是一個Federal Express的員工.
他坐的飛機.
其實飛機全部都是包裹.
所以當飛機墜毀的時候.
那些包裹就逐漸漂流上岸.
他怎樣生存呢.
就拆開那些包裹.
然後找找裡面有什麼東西幫助他生存.
例如什麼呢.
魯賓鞋的刀片.
就成為他去切魚.
去切什麼的工具.
原來包裹裡面有一個分沙.
你猜那個分沙有什麼用呢.
那個分沙的沙可以用來打魚.
原來可以這樣用的.
但最重要最重要.
那些是求生的工具.
但你猜最重要對湯漢斯來說.
是什麼東西令他能夠生存下去呢.
是這個排球.
不知道大家記不記得這個排球叫什麼名字.
原來這個排球.
當時他一打開的時候.
就成為他荒島求生最重要的夥伴.
為什麼他有個手掌呢.
其實是個手掌血印.
因為當時他很生氣.
弄傷了手.
然後他很生氣發脾氣.
他一下打到排球.
怎料就印了一隻血掌印在排球上.
他就忽發奇想.

$^{81}$就畫了一些面孔.
然後賦予了這個排球的生命.
這個排球叫什麼呢.
他給了名字排球的.
叫什麼名字呢.
Wilson.
然後他就和這個Wilson成為了好朋友.
在那個荒島的過程裡面.
就一直和他對話.
甚至去到電影的高潮.
他終於自己弄了一隻木筏.
打算離開那個荒島.
他什麼都不帶.
他帶什麼呢.
Wilson.
和一個他沒有打開的包裹.
結果在海裡面遇到風暴.
那個Wilson就衝走了.
捲走了離開.
而整套戲裡面最令人動容的那一幕.
就是Wilson離開.
然後他落跑.
他跳下海.
很努力想要撈回Wilson.
結果Wilson越走越遠.
他很大聲地叫了一句.
I'm sorry Wilson.
這個就是那一幕.
令到很多人很深刻的印象.
如果你.
如果你留在荒島.
你覺得你寧願帶一個人.
還是帶一個物件去呢.
人類是需要夥伴的.
《笑行傳》.
我今天會和大家看的是《笑行傳》.
《笑行傳》顧名思義.
其實是一本關於旅程的書.
是一位使徒.
即是耶穌的跟隨者.

$^{121}$如何在耶穌復活升天之後.
他的行蹤.
即是叫使徒行傳.
某個程度上你可以將《笑行傳》.
當作是Youtube的旅遊片去看.
你可以當作是.
他會不會去了甚麼景點.
他會不會吃甚麼米芝蓮.
但這些都沒有的《笑行傳》.
相反是說甚麼.
是說他如何按神的吩咐.
無論是彼得.
無論是保羅.
去邊走邊傳.
傳福音的紀錄.
但我今天想和大家看的是.
《笑行傳》最後一章.
《笑行傳》28章.
我們來看最後的兩節經文.
《笑行傳》如何結束的呢.
即是這本關於使徒行蹤的一本書.
去到28章30至31節.
他如何結束這本書呢.
我們一起讀.
《笑行傳》28章30至31節.
保羅在自己所租的房子裡住了足足兩年.
凡來見他的人.
他全都接待.
放膽傳講神國的道.
將主耶穌基督的事教導人.
並沒有人禁止.
我讀經的時候.
我經常會怎樣呢.
想像自己是聖經的作者.
然後怎樣寫這些書卷.
你試想像一下.
如果你是《笑行傳》的作者.
這本紀錄使徒.
耶穌升天之後.
使徒行蹤的一本書.

$^{161}$它應該怎樣結束呢.
你怎樣都不會想到.
最後《笑行傳》的結束.
是這兩節經文.
其實很奇怪.
它一直都沒有提到.
保羅的生活是怎樣.
它主要記載《笑行傳》.
就是傳福音的紀錄.
港道的紀錄.
有多少人信主.
在哪裡建立教會等等.
但是來到《笑行傳》的結束.
無端端又變回YouTube的旅遊片.
說他在一個自己所租住的房子.
然後接待什麼人.
然後只是說全港神國的道.
並沒有人禁止.
句號完結《笑行傳》.
你想不想像得到.
這個是《笑行傳》的高潮呢.
通常一本書的結尾.
應該是高潮.
為什麼《笑行傳》最後.
會來到一個反高潮呢.
還要是甚麼呢.
這裡提到說第30節.
保羅在自己所租的房子裡面.
為什麼要特別註明.
保羅在自己所租的房子裡面呢.
好像很平常.
我們這班很多都是從香港過來美國讀書.
自己租房等等.
我們很多這些經歷.
但是為什麼要特別提.
保羅在自己所租的房子裡面.
住了足足兩年呢.
它背後的背景是什麼呢.
原來這裡保羅不是普通的租房子.
然後過生活.

$^{201}$其實這裡保羅在做什麼呢.
他在坐牢.
這裡保羅在做什麼呢.
其實他在坐牢.
《笑行傳》20章就開始預言.
保羅將會面對囚禁.
結果21章真的被捉拿了.
然後由22章一直到27章.
不斷提到保羅如何面對不同的審訊.
然後去到28章《笑行傳》的結束.
由20章一直開始.
其實《笑行傳》的高潮是什麼呢.
是保羅面對監禁.
然後去到28章就提到保羅終於去到羅馬.
然後他要面對審訊.
所以這裡提到他自己所租的房子.
其實他面對一個什麼呢.
一個不是典型的審訊.
不是一個監獄.
他是一個什麼呢.
如果用英文就叫做house arrest.
即是軟禁.
還要是自己租的房子.
你猜他是不是自己住呢.
不是.
在《笑行傳》28章16節提到.
保羅是和一個看守他的兵一起住.
所以這裡是在說保羅住在一個這樣的房間.
足足兩年一個另類的囚禁.
但當然也容許某程度上的自由.
這個是《笑行傳》的結束.
你會不會覺得如果你是追劇.
無論Netflix或者是Disney Channel.
由20章一直追保羅的囚禁的命運.
然後去到28章.
無端端兩節經文就完了.
根本沒有交代保羅的去向.
或者他將要面對的審訊將會是怎樣的時候.
你會覺得怎樣呢.
香港人有一個典型的形容詞是什麼呢.

$^{241}$「戒凳」有沒有聽過.
即是你看一套戲.
看到很緊張.
很想看結局是怎樣的時候.
突然間他怎樣.
他就計劃去拍那些花園.
然後有太太在煮菜.
然後就結束.
這裡其實就是這樣.
如果香港人就會怎樣.
會「戒凳」.
你沒有交代結局.
保羅《笑行傳》最出名的是一個開場的結局.
是這樣來結束《笑行傳》的結局.
為什麼《笑行傳》是用這兩節經文來結束呢.
一本這麼重要的關於使徒傳福音的記錄.
教會被建立的記錄.
為什麼最後是用這兩節經文來結束呢.
特別是第30節.
保羅在自己所租的房子裡住了足足兩年.
然後只是記載他在這個房子裡做了兩件事.
在《笑行傳》28章結束的時候.
只是提到保羅在一個軟禁的底下.
他做了兩件事.
哪兩件事呢.
第一件事「凡來見他的人,他全都接待」.
這個又是令人很難想像的.
你會不會將「囚禁」或者「軟禁」和「接待」兩個字.
拉在一起來看.
你有沒有想像到「接待」.
在一個什麼狀態下才會接待呢.
你是有餘的時候才會接待.
你會覺得自己很安恕的時候才會接待.
但這裡提到的是什麼呢.
是他接待人.
其實反過來應該是保羅需要被接待.
但這裡說他這樣「凡來見他的人,他全都接待」.
我舉一個例子讓你能夠想像到那種狀態.
你試一下今天.
我媽媽也在.

$^{281}$你試一下今天突然間崇拜結束.
打電話給家人跟媽媽說.
我今天崇拜結束.
我有20人會想上來家裡吃飯.
你猜你媽媽會怎樣呢.
有人會說「痴線」.
會剷我剷到飛起來.
你猜真的加雙筷嗎.
加也要加20對筷子.
所以接待其實不是一個很容易的事.
特別是保羅在一個這樣的狀態下.
竟然「凡來見他的人,他全都接待」.
第二件事保羅在一個這樣的狀態下.
他做什麼呢.
放膽傳講神國的道.
將主耶穌基督的事教導人.
並沒有人禁止.
所以基本上紹行傳最後只是記載.
保羅在囚禁裡做了兩件事.
第一件事是接待.
第二件事是傳道.
而這兩件事原來是一對的.
是分不開的.
什麼意思呢.
原來傳道不能夠離開接待.
反過來接待.
同時也在做什麼呢.
就是傳道.
一個是用口講.
傳講神國的道.
一個是用行動來證明神國的道.
原來在《使徒行傳》28章最後告訴我們.
原來接待和傳福音是離不開的.
甚至這麼說.
接待就是福音的核心.
《使徒行傳》28章最後告訴我們.
原來接待就是傳福音的核心價值.
在古代是非常需要接待的.
因為在古代沒有酒店.
所以當時的社會.

$^{321}$接待是一個很重要的.
他們在文化裡是很熟習的一件事.
但來到我們現代社會.
酒店這麼方便.
而且在香港住過的話.
地方淺窄.
要接待人其實非常困難.
特別是經過什麼呢.
經過疫情.
其實人與人之間的距離被拉遠了.
很多疏離.
再加上現在大家對很多事情的看法很不同.
接待是越來越困難的.
我想起一件事.
可以講一個例子給大家聽.
我在香港的教會在佐敦.
在裕華的對面.
在去年四月的時候.
我教會對面那座大廈.
是一個很出名的大廈.
大火.
非常嚴重的大火.
後來有傷亡.
如果是你的教會對面那座大廈大火.
你會怎樣呢.
起初我們同工商量.
是否安全回教會辦公.
會否有風路.
或者煙很強.
後來我們思考.
不是的.
我們回教會的途中.
看到很多人.
因為那座大廈很大.
是一個很大的complex.
上面有旅館.
那些人因為逃生.
所以全部在裕華門口.
裕華是國貨.
或者在街邊蹲著.

$^{361}$找個地方坐.
等候解封.
可以回家拿東西.
我們就思考.
我們作為教會.
我們應該怎樣做呢.
我們幾個同工就拿了些水.
拿了些口罩.
然後我們邀請那些街坊.
上教會來休息.
我們就說.
我們教會有椅子坐.
有冷氣.
有水喝.
有廁所.
我的角度是.
應該會有人跟進.
誰知道你猜對了.
沒有人有回應.
不用了.
我給他一瓶水.
他都說不用了.
我心想為甚麼呢.
人與人之間的信任.
其實是越來越少.
要去一個陌生的地方.
儘管他其實很有需要.
但他都會覺得.
有個敵意.
我防範.
直到我提出一件事.
結果四五十人上了教會.
你猜是甚麼呢.
你猜到嗎.
充電.
因為你知道.
在這樣的狀況下.
每個人都要上網.
跟足資訊.
然後就覺得很緊張.

$^{401}$沒有電.
很多都是老人家.
他們不懂租充電的地方.
去拿電.
結果我一提出.
想不想上網.
上我們教會充電.
馬上兩眼發光看著我.
陳太,李生.
他們說可以充電.
結果四五十人.
上了教會充電.
人與人之間的隔膜.
那種不信任.
是越來越濃重.
我們以為我們的package很正.
又有早餐.
然後又有甚麼食物.
又有甚麼offer.
最重要是我們有福音.
我們有救恩.
但可能他覺得不relevant.
他覺得不touch到他.
直到你找到接觸點.
你touch到他.
他就馬上有反應.
原來你有他需要的答案.
當然我回到教會.
就馬上找出全教會的充電線.
然後拿拖把來充電.
但保羅在《少行傳》28章.
告訴我們原來接待是這麼重要.
但未完的.
這是《少行傳》28章的結束.
亦是《少行傳》的結束.
但原來28章前面.
還有另一段是講保羅被接待.
是令人很驚訝的.
不是保羅接待人.
是保羅被接待.

$^{441}$我們一起來看《少行傳》28章.
由1到第10節.
我讀給大家聽.
「我們既已得救才知道那島名叫馬爾他」.
提到保羅因為要上羅馬接受審訊.
所以他要行船.
結果在中間遇到船難.
所以他們滯留在島上.
叫做馬爾他.
第二節.
「土人看待我們有非常的情分.
因為當時下雨天氣又冷.
就深火接待我們眾人」.
這個反過來是保羅被接待.
「那時保羅拾起一捆柴放在火上.
有一條毒蛇因為熱了咬住他的手.
土人看見毒蛇懸在他手上.
就彼此說:這人必是個兇手.
雖然從海裡救上來.
天理還不容他活著.
保羅警版那毒蛇甩在火裡並沒有受傷.
土人想他必要腫起來.
或是忽然撲倒死了.
看了多時見他無害.
就轉念說他是個神.
離那地方不遠有田產.
是島長布伯留的.
他接納我們進程款待三日」.
「當時布伯留的父親患熱病和離疾躺著.
保羅進去為他禱告.
按手在他身上治好了他.
從此島上其餘的病人也來得了醫治.
他們又多方地尊敬我們.
到了開船的時候.
也把我們所需要的送到船上」.
在《兆行傳》28章一開頭.
是反過來是保羅被接待.
而且保羅被接待的情況.
是令人匪夷所思的.
他在什麼狀態下被接待呢.

$^{481}$他是在船難中被接待.
而且他被誰接待呢.
土人.
這是和合本的翻譯.
如果你有英文翻譯.
你看看這個土人是怎樣翻譯呢.
是barbarian.
是野蠻人.
或者有些翻譯是uncivilized的人.
為什麼呢.
因為當時新約的背景.
希羅文化接待是什麼呢.
是有文化的人才接待的.
原來接待在當時的文化裡面.
是一個civilization的表達.
但是28章一開頭讓我們看到.
是一個完全逆轉的相反.
是什麼人接待保羅呢.
是一群barbarian.
野蠻人.
土人.
沒有文化的人.
未開化的人.
而且他們怎樣接待保羅呢.
土人看待我們.
其實是接待這個字.
有非常的情分.
什麼叫做非常的情分呢.
英文翻譯為unusual kindness.
是令人很詫異.
不是平常的.
是unusual.
是非常的.
這樣去接待.
怎樣接待呢.
他們當時是船難.
他們游完水.
他們可能會怎樣呢.
會冷.
結果土人怎樣幫他們想呢.

$^{521}$生火接待他們.
不只是給他們飯吃.
當然我們知道他們最需要的是什麼.
他們最需要可能是水.
飯吃.
但不是.
土人站在他們的角度去思考.
生火.
第一件事做的.
就是生火接待他們的眾人.
原來當保羅都同樣被接待.
這是很令人詫異的.
原來土人都會這樣去接待保羅和他那班人.
甚至怎樣呢.
去到第七節提到.
賭掌布百劉.
布百劉是一個銜頭.
所以這班人是怎樣.
去接濾盡情款待三人.
又是一些覺得不是誠實的人.
他都願意這樣熱情地.
甚至有朋友的意思.
所以他們在接待的過程.
有那種關係的建立.
然後你看到.
醫治才發生.
先是接待.
先是關係的接納.
然後保羅來醫治.
那種生命的流露.
那種醫治.
在那裡發生.
我離開美國其實很久.
我之前在三藩市住了十年.
在芝加哥住了八年.
我今年回來原來很特別.
我後來回想.
我離開三藩市二十年.
離開芝加哥十年.
我這次回來三個星期.

$^{561}$有一個多星期在三藩市.
一個多星期在芝加哥.
我兩個地方都是住朋友的家.
甚至很多時候.
我來到美國的時候.
都有朋友offer.
借車給我開.
offer款待我.
請我吃飯.
我有時候想.
我離開三藩市二十年.
離開芝加哥十年.
中間其實沒有交接.
就算回香港見面吃飯.
都只不過是一餐飯.
最多兩餐三餐飯.
為什麼離開了這麼久.
大家仍然願意接待我呢.
因為關係.
因為關係已經建立了.
如果問我.
我在美國前後住了二十年.
最寶貴的是什麼呢.
是建立了很多可以last到一生的關係.
而這些關係.
很多時候是怎樣建立呢.
在接待之間.
我被接待.
同樣我也接待人.
我回到三藩市.
那些人還在說什麼呢.
當時他們剛剛開始來美國的時候.
我們怎樣載出載入.
怎樣跟他們慶祝生日.
怎樣中秋一個月餅.
十個人分.
這些經歷.
是說二十年都繼續說的.
原來接待是有些種.
就是福音的核心.

$^{601}$是令人很驚訝.
同時讓我們看到.
保羅為什麼被下監呢.
就是被猶太人逼白.
猶太人是什麼人呢.
是他同族的人.
是跟他有相同的訓練.
認識舊約聖經的人.
但是猶太人反過來是什麼呢.
不接待保羅.
在這裡讓我們看到.
是賭掌捕白流.
是土人用非常的情分.
來接待保羅.
聖經有時候諷刺的.
就在這裡.
紹行傳最後結束.
在三十.
三十一節.
為什麼要用一個接待的episode.
去結束紹行傳呢.
因為接待就是福音的核心.
第三十節.
保羅在自己所住的房子裡.
住了足足兩年.
凡來見他的人.
他全都接待.
傳福音那種的接待.
和普通人那種的接待是不同的.
普通人那種的接待.
是怎樣呢.
是會選擇過的.
是跟你關係好.
你就當然是熱情的款待.
但是保羅這個凡來見他的人.
沒有差別.
沒有條件地.
這樣來接待.
好像土人接待保羅一樣.
是沒有差別.

$^{641}$沒有條件.
不去分到底是陌生人.
是不是異己.
甚至是敵人.
就算他還沒有認識神.
他還沒有認錯.
他都可以這樣來接待.
可能大家都會有印象.
保羅曾經接待.
在羅馬曾經接待過哪個人呢.
有個人叫做亞尼西姆.
有沒有印象呢.
結果因為這樣.
他是一個奴僕.
可能是一個犯事的奴僕.
但是他一樣接待.
結果生命.
影響生命.
有改變.
不單止是這樣.
最後放膽傳講神國的道.
放膽是什麼意思呢.
即是他裡面沒有限制.
放膽即是他裡面沒有限制他.
傳講神國的道.
外面儘管有監禁.
有軟禁.
但都不會阻止他去接待其他人.
他裡面也沒有阻止他傳講神國的道.
後來這一節.
31節.
《少行傳》28章最後一節.
其實提醒了我們.
能夠沒有差別.
沒有條件地彼此接待的關鍵是什麼呢.
因為我們相信.
我們所信的是神國的道.
是神國的道.
我們所信的福音.
是有他自己的國度.

$^{681}$所以我們在這個國度.
我們怎能夠自己去判斷.
有什麼人我們選擇來接待.
有什麼人不肯去接待呢.
所以如果我們真的相信.
我們所信的是神國的道.
是超越一切的權勢.
一切的限制.
一切的界限.
甚至政見.
很多不同的看法.
如果我們這樣相信的話.
我們不單止能夠放膽傳講神國的道.
我們更加能夠範內見他的人.
他傳道接待.
最後.
我用一個大家很熟悉的聖經故事.
來結束這一段的講論.
我們看回同樣是路加醫生所寫的路加福音.
我們今天看《少行傳》.
是路加醫生所寫的後傳.
我們看回前傳路加福音.
這個大家很熟悉的經文是什麼呢.
路加福音第十章.
好撒瑪利亞人的故事.
我們再看原來好撒瑪利亞人.
是一個接待的故事.
原來好撒瑪利亞人正正就是說什麼呢.
接待.
路加福音第十章二十五節.
有一個律法師起來試探耶穌.
說:夫子,我該做什麼才可以承受永生?.
即是下面的討論是在說得不得救的問題.
不是錦上添花的問題.
是在說得不得救的問題.
而且是要做的.
二十六節.
耶穌對他說:律法上寫的是什麼?.
你念的是什麼?.
他回答說:你要盡心盡性盡力盡義愛主你的神.

$^{721}$又要愛倫舍如同自己?.
耶穌說:你回答的事.
你這樣行.
就必得永生.
即是那個律法師是完全正確的.
他是Bingo的.
他拿4.0的GPA.
但是耶穌怎樣回答?.
說你答得對.
但未完.
你這樣行.
是.
我們裡面有很多金句.
我們查很多聖經.
但是耶穌的關鍵說.
你要這樣行.
就必得永生.
然後說.
那人要顯明自己有利.
就對耶穌說:誰是我的倫舍?.
即是那個律法師.
他也認同耶穌.
但他想再駁斥耶穌.
所以他就問了一個問題.
誰是我的倫舍?.
這句問題其實在說什麼?.
就是他要選擇誰來接待.
誰是我的倫舍?.
所以他肯接待的.
但他要先選擇.
耶穌回答說.
有一個人從耶路撒冷下耶利哥去.
落在強盜手中.
他們扑去他的衣裳.
把他打個半死.
就叼下他走了.
偶然有一個祭司從這條路下來.
看見他就從那邊過去了.
又有一個利未人來到這地方.
看見他也照樣從那邊過去了.

$^{761}$所以好撒瑪利亞人是說一個接待的故事.
是說什麼人接待什麼人.
第33節.
唯有一個撒瑪利亞人走路來到那裡.
看見他就動了持心.
留意.
聖經非常仔細記載.
這個好撒瑪利亞人的每一個行動.
而這每一個行動.
其實都是接待的行動.
上前用油和酒倒在他的雙柱.
包裹好了.
扶他騎上自己的身後.
帶到店內去照應他.
第二天拿出二錢銀紙來交給店主說.
你且照應他.
此外所費用的.
我回來必還你.
前面幾個利未人.
祭司那些都快快地帶過去.
唯有這個撒瑪利亞人.
他接待的故事.
仔細到不能再仔細.
連他的對白都記載了.
原來當人接待的時候.
神是會記得的.
神記得你做的每一個行動.
怎樣接待其他人.
36節.
你想這三個人哪一個是落在強盜手中的淪舍呢?.
他說是憐憫他的.
耶穌說你去照樣行罷.
律法師問我可不可以選擇誰接待.
耶穌說你不可以選擇誰接待.
你的接待是要無差別.
無分別.
亦都無任何的barrier.
去接待凡有需要的人.
今日我們面對的是甚麼人呢?.
我們有甚麼需要呢?.

$^{801}$或者我們身邊的人有甚麼需要.
我們看不見他的需要.
我們願不願意go extra mile.
去接待人.
我剛才看到你們今年的年提.
還是標記是否跨越.
是否cross over.
其中一樣東西就是.
要為基督接觸身邊的淪舍.
我們是否那種無差別.
亦都無條件地去接待其他人呢?.
為何我們今天能夠這樣做呢?.
因為我們稍後會領受聖餐.
我們今天其實都只不過.
是被耶穌接待的人.
我們為何不肯打開我們的生命.
去接待其他的人呢?.
甚至是那些我們覺得.
不可能接待的人呢?.
我們一起祈禱.
因為我們今天正正是被你無條件.
無差別地去接待的人.
我們本來都是不配.
但在我們還作罪人的時候.
你就為我們死釘在十字架上.
今天你要再一次來接納我們.
同樣都幫助我們.
既然蒙你接納的時候.
但願我們願意打開我們的生命.
亦都同樣地無條件無差別地.
去學習去接待其他人.
因為正正是這樣的時候.
我們流露你的生命.
讓人見到原來你就是那個接待者.
求你恩待我們.
我們這樣禱告交託.
奉耶穌基督的聖名.
阿們.
(下集待續).
\newpage



\section{使徒行傳 4:32-37}
\label{sec:_p6p1yxlKsI}
\textbf{成為聖靈的群體 (使徒行傳4\_32-37) - 陳崇基牧師}
\newline
\newline
連結: \href{https://youtube.com/watch?v=-p6p1yxlKsI}{\texttt{https://youtube.com/watch?v=-p6p1yxlKsI}} ~~~~ 語音日期: 2025-07-14
\newline
\newline
\hyperref[sec:EEmGj96kwiU]{\small{< < < PREV SERMON < < <}}
~
\hyperref[sec:index_chronic]{\small{[返順時目]}}
~
\hyperref[sec:index_scriptual]{\small{[返順卷目]}}
~
\hyperref[sec:6C_6zjSlvhM]{\small{> > > NEXT SERMON > > >}}
\newline
\newline
使徒行傳 4:32-37
\newline
\begin{longtable}{cl}
\hline
\hline
章節 & 經文 (和合本修訂版)\\
\hline
4:32 & \begin{tabularx}{0.7\textwidth}{X} 許多信徒都一心一意，沒有一人說他的任何東西是自己的，都是大家公用。 \end{tabularx} \\ \\ \relax
4:33 & \begin{tabularx}{0.7\textwidth}{X} 使徒以大能見證主耶穌復活；眾人也都蒙了大恩。 \end{tabularx} \\ \\ \relax
4:34 & \begin{tabularx}{0.7\textwidth}{X} 他們當中沒有一個缺乏的，因為凡有田產房屋的都賣了，把所賣的錢拿來， \end{tabularx} \\ \\ \relax
4:35 & \begin{tabularx}{0.7\textwidth}{X} 放在使徒腳前，照每人所需要的，分給每人。 \end{tabularx} \\ \\ \relax
4:36 & \begin{tabularx}{0.7\textwidth}{X} 有一個利未人，名叫約瑟，使徒稱他為巴拿巴（巴拿巴翻出來就是安慰之子），生在塞浦路斯。 \end{tabularx} \\ \\ \relax
4:37 & \begin{tabularx}{0.7\textwidth}{X} 他有田地，也賣了，把錢拿來，放在使徒腳前。 \end{tabularx} \\ \\
[1ex]
\hline
\hline
\end{longtable}
$^{1}$在這方面 行動上和祈禱上支持.
今天我們回到這個.
史唐人傳的系列.
我們講到跨越.
剛才我所講的.
我們需要跨越去認識其他人.
甚至是十一點半鐘的其他群體.
上個星期我不在大家中間.
我去了Whitten的華人宣導會.
他們的退休會去到港獨.
我自己.
大英姐妹都很受鼓勵.
我自己也很受鼓勵.
因為我見到有大英姐妹.
很熱心去傳福音.
很熱心很著緊教會的事工.
著緊他們教會的事工.
不著緊我們教會的事工.
但不要緊 每一個教會都是一樣的.
我為大英姐妹獻上感恩.
這是一件很好的事.
所以上個星期不在大家中間.
八月初我都會去San Diego.
另一個教會的退休會做講員.
我去San Diego不是去樂開.
我不是去談世界.
我在他們的退休會.
在他們營地講一個weekend.
願意主動去使用我們.
並且剛才Katrina所講的.
我們再次解見錯誤的網上課程.
我們即將開播.
就是星期二晚和星期三早.
我們即將開播.
有五十多位的朋友是註冊了的.
這五十多位的朋友都不是我們UV的人.
來自香港 來自北美不同的地方.
甚至英國 澳洲等等.
我們希望這個課程能夠.
繼續去到傳講上帝的道.

$^{41}$教導弟兄姊妹如何去掌握上帝的話語.
上帝的話語很重要.
所以今天我們回來這個跨月的系列.
我們說要挑戰大家三件事.
第一件事 你自己的舒適區.
我們需要踏出自己的舒適區.
並且我們要連結其他的弟兄姊妹.
無論在哪一個的組群.
我們需要去到跨越.
能夠去認識大家等等.
並且將福音傳給其他人.
剛才Ken傳到所 祈禱裡面提過的.
就是昨天我們有三十多位.
不同教會的弟兄姊妹在這裡.
一起出隊去報導 關懷等等.
我覺得這是一個非常好的畫像.
弟兄姊妹很著緊去到.
為著福音去到努力.
所以這個也是我們跨月的系列.
使徒行傳.
上一次我們說到第三章 記得嗎.
第三章 就是第一個使徒行傳記載的教會神蹟.
就是一個生下來瘸腿的跛腳的人.
在聖殿的門口尾門.
然後得到醫治 記得嗎.
拿撒拉人耶穌的名.
他的power 他的能力.
能夠使這個人能夠起身.
然後他進入聖殿去見拜上帝.
在第一個神蹟之後.
彼得就開始講道.
就解釋為什麼會有一個這樣的神蹟.
為什麼有一個這樣的劃時代的神蹟出現.
就是耶穌基督的名.
他是那位死而復活得勝的主宰.
就是他有這樣的能力.
你呼喊耶穌的名字.
藉著耶穌的名字來祈禱的時候.
我們生命會向著上帝這個有能力的國度前進.
所以彼得一直解釋的時候.

$^{81}$誰知道到第四章的時候就有人暫亂歌頌.
講著講著的時候.
彼得就被這些猶太的聖殿的權貴.
來暫亂歌頌.
就是講著講著有人停下.
所以第四章開始的時候.
他就被打壓了.
然後就將他去到收監.
這裡就給第四章的一至三十一節的分段給大家.
首先他們被打壓.
然後使徒有他的回應.
八至十二節.
然後也有他的權貴.
那些很失眼的人.
不喜歡他們的人.
來盤問他們為什麼他們要傳福音.
然後使徒也有所回應.
然後他們就威嚇他.
做不到什麼.
結果嚇一嚇他們.
然後最後二十三至三十一節.
就是教會整體的回應.
這裡是史蘭傳第一次記載.
教會受到當時的猶太人的權貴.
耶路撒冷的權貴打壓的故事.
所以史蘭傳裡面有很多第一次的故事.
First story.
我就問了.
這些first story用來做什麼.
第三章我剛才講了什麼.
有神蹟奇事出現.
就是生下來.
瘸腿的能夠起來.
應該是大家都開心的.
確實是百姓都開心.
誰知道呢.
不一定的.
總有人來阻礙上帝的工作.
你有沒有想過.
上帝工作開始.

$^{121}$大家很開心.
大家應該很好的.
應該普天同慶.
不過很多時候就會有反對的勢力出現.
這裡就是第四章所講的.
所以第四章第八節開始發生什麼事.
他們被打壓和被收監的時候.
那時彼得被聖靈充滿.
對這些權貴說.
民間的官長和長老.
倘若今天我們被查問的是因為.
殘障的人身上所行的善事.
這人得到傳譽的.
那麼你們大家以色列全民都應該知道.
站在你們面前.
這人得到傳譽.
生下來的能夠起來的在他們面前.
是因為什麼呢.
是因為你們所釘在十字架.
上帝使他從死裡面復活的那撒人.
耶穌的名字.
又是講到這個名字.
耶穌的權能.
這位耶穌是誰呢.
就是你們將人所雕器的石頭.
已經成為防國的頭塊石頭.
防國者第一塊石頭.
除他以外別無拯救.
因為在天上地下.
所有次下別的名的.
我們可以靠著得救.
只有耶穌基督的名字.
才能夠使我們罪人能夠得救.
所以他就指正一件事.
就是耶穌的緣故.
有這件事情發生.
是耶穌才是我們救恩的源頭.
是唯一的源頭.
這裡講了一句話.
是引用詩篇118篇的.

$^{161}$他說你們將人所雕的石頭.
已經成為了頭塊的防國石.
什麼意思呢.
你想像一下如果你是石匠的話.
Stone mason.
你做什麼的.
Stone mason是做什麼的.
有沒有誰是以前作石的.
搞石頭的.
沒有吧.
石匠是做什麼的.
將人是做什麼的.
他就是去一個石礦場.
然後石礦場要有些好的石.
有些不好的石.
他就將好的石全部採出來.
剩下的不好的石.
就是那些碎石.
不能夠做建築的那些石.
那些壞的材料.
他就將它雕起.
這個就是石匠所做的事.
對不對.
就像市場一樣.
買東西 挑菜.
每一條都去挑.
剩下的殘花 白柳就沒有人理.
類似這樣.
石匠就是做這件事.
以前的第一世紀也好.
古代的人也好.
特別是在耶路撒冷山區.
他們有很多石匠.
石匠就經常找一些好的石礦場.
然後將最好的石頭去採.
最好的那塊是用來做什麼的.
就是你建一個新的房樓的時候.
就是做那個防角石.
第一塊放在那裡的.
我們芝加哥Downtown.

$^{201}$你走來走去的時候.
那些很漂亮的建築.
下面有些防角石.
有時候會雕一個年份在那裡.
就是多少年放這塊石頭在那裡.
那塊是最重要的石頭.
匠人就是做這件事.
順帶一提.
芝加哥其實也很出名.
當然我們出名的東西除了.
槍戰.
死人多.
貪污.
柏號.
這些我們很出名.
但有一樣東西很出名就是石礦場.
因為芝加哥南邊有一個地方叫做.
Fountain Quarries.
如果開294.
有一段你開的時候要留意一下.
左邊右邊.
是好像懸崖一樣.
你有沒有留意過.
294.
那裡是什麼地方.
其實以前是石礦場.
因為全世界最大的其中一個石礦場.
你看到他們已經挖掘了最好的石.
現在做什麼呢.
就是芝加哥有一個Tunnel的Project.
就將那些污水和雨水.
如果下雨下得多.
特別現在周圍都下這麼多雨.
昨年就火燒山.
周圍都燒.
今年就到處都有水.
周圍都有水.
就將那些Flood water.
然後Channel了到很大的水池.
以致不會那麼容易水浸.

$^{241}$就是一個最大的石礦場.
就將最好的石頭已經拿走了.
拿去Downtown 買去其他地方.
然後剩下來的就是一個沒用的東西.
詩篇第118篇就是在說這件事.
在說什麼呢.
當時詩篇的原意大概就是那些外邦人.
看到耶路撒冷.
看到石鞍山.
你知道耶路撒冷是石山.
是盆石.
耶路撒冷是石頭.
你去耶路撒冷走一走.
周圍都是石頭.
大的小的.
全部這些樓房都是用石頭所建的.
詩篇的原意就是說.
這些外邦看到石鞍山.
看到耶路撒冷的時候.
他們就覺得有什麼好.
有什麼特別.
因為好的石頭都不在耶路撒冷出來.
世界上最好的石頭是哪裡出來的.
在埃及.
是不是.
在印尼.
你平時廚房上面那塊.
那塊東西很貴.
那些東西是哪裡來的.
很多都是import來的.
意大利也好.
或者其他地方.
都不是耶路撒冷出來的.
所以外邦就會覺得.
這個石鞍山.
這個耶路撒冷.
這塊石頭有什麼了不起.
但是詩篇說什麼.
上帝偏偏選擇世人所厭惡的.
世人所凋棄的東西.

$^{281}$看不起的.
或者看不起的石頭.
來作防國石.
成為聖殿的地方.
所以上帝也會說.
就是選擇石鞍山.
選擇耶路撒冷.
作為聖殿的地方.
就是救恩在這裡出來.
這是詩篇118篇的原意.
你看到彼得將這個句子.
118篇.
是應用在耶穌基督身上.
他說你們將人所棄的石頭.
就是指耶穌基督.
成為防國石.
你們看不起的.
它變成救恩的源頭.
所以除他以外就別無拯救.
誰知道原來.
耶穌死的地方和復活的地方.
我們叫它什麼.
築留地.
Place of School.
或者Gographer.
拉丁文叫做Galvory.
都是指同一個地方.
就是耶穌所死的地方.
他在那裡被釘十字架.
然後墳墓復活.
今天就是聖母堂的地方.
在耶路撒冷.
本身那個地方是一個石礦場.
將好的石頭已經被人拿走了.
那些匠人已經拿走了.
剩下的一個沒什麼用的.
一個廢的地方.
殘廢的.
荒廢的一個石山.
十字架就在那裡發生.

$^{321}$耶穌就在那裡被釘死.
然後死在裡面復活.
成就上帝的救恩.
所以彼得就說.
你們這些猶太人.
在耶路撒冷的宗教的權貴.
你們看不起的耶穌.
你將他釘在十字架.
上帝就偏偏選擇世人看為低等的石頭.
耶穌基督成為救恩.
所以四章第十一節.
他說什麼.
你們所凋棄的.
耶和華揀選了.
耶和華揀選十字架.
就是耶穌基督作為世人的救恩源頭的地方.
講完這篇道理.
或者一個這樣的回應之後.
這個教會的敵人就做什麼呢.
他們就盤問和威嚇他.
然後他們就這樣回應.
叫他不要再傳耶穌基督了.
你不要再搞亂我們的當了.
因為越來越多人跟著你的時候.
我們就不行了.
我們社會就會有紛亂.
四章第十九節彼得怎樣回應.
彼得和約翰就說.
聽從你們這班權貴.
不聽從上帝.
在上帝面前合理不合理.
你們自己判斷.
你們所看見的.
我們所看見的.
所聽見的.
我們不能不說.
福音我們不能不傳.
這個很奇怪.
教會當時第一次被上面的政權和權貴威脅.
上帝怎樣回應.

$^{361}$他們怎樣回應.
上帝說將聖靈充滿他們.
然後他們更加有膽量去面對他們的對手.
不是做縮頭烏龜.
很奇怪吧.
然後彼得用這樣的回應.
我們不會服從人的.
我們只會服從上帝.
令到彼得勝靈充滿他們.
更加有能力.
繼續去傳福音.
繼續去廣道.
這個是《小行傳》裡面所說的神蹟騎士.
我們說的神蹟騎士.
就是我們會很興奮.
然後做了不同的事.
那些沒錯.
可能都是神蹟騎士.
但最主要是《小行傳》裡面所說的.
上帝的工作.
真正的工作.
他的聖靈的權能在哪裡彰顯出來.
就是使到使徒.
使到教會.
能夠有能力.
更多的能力.
來作見證.
不是叫我們縮頭烏龜.
不是叫我們安逸.
而是他們一班人認定上帝是有權能的.
但願我們教會都是這樣.
今時今日傳福音是難的.
當然每個時代都難.
但我們面對不同的難關.
但我們應該怎樣回應呢.
或者上帝要我們怎樣回應.
是不是叫我們縮頭.
然後我們自己安身.
那就算了.
不是的.

$^{401}$這裡說的是上帝使教會要更加有能力.
我們的禱告應該是這樣的.
當你遇見有困難的時候.
遇見有軟弱的時候.
我不知道你怎樣祈禱.
祈禱就是讓我放假.
讓我休息.
讓我安逸.
你試一下反過來.
上帝可不可以給我能力.
讓我能夠繼續向前.
故事繼續.
第21節.
官長為百姓的緣故.
就想不出任何法子來懲罰他們.
只好怎樣.
就像黑社會一樣.
嚇一嚇他們.
釋放.
因為眾人百姓都是為了所行的神蹟.
然後歸榮耀給上帝.
OK.
23節.
兩個人被釋放了.
他們就回到自己人的地方.
然後將祭司長長老所說的話告訴他們.
他們聽見了之後.
就同心合意的高聲的向上帝禱告.
這個脈絡是什麼.
彼得被拘捕.
然後被公安威嚇.
你想想教會同心合意禱告.
這24節最後那一段.
他們同心合意高聲的向上帝禱告.
你猜他們禱告什麼.
你會禱告什麼.
如果牧師被捉了.
他進入公安那裡.
然後進去兩天.
然後被釋放出來.

$^{441}$你會高聲的向上帝說什麼.
為什麼要放他走.
你可不可以讓他坐多兩天.
讓他休息一下.
我不知道你們怎麼祈禱.
教會怎麼祈禱.
通常我們自然的反應就是.
OK.
他們得到釋放.
然後感恩.
平安無事.
大能不死.
是不是.
我們真的好了.
但是教會怎麼祈禱.
我揭曉了24節.
他們聽見了就同心合意的向上帝禱告.
他說主啊你是那位創造萬物的.
你藉著聖靈託付.
你僕人就祖宗大位.
曾經說過什麼.
外邦一定會反抗的.
萬民必定會謀算虛妄的事.
地上的君王都站穩了.
他們要聚集一齊.
來對抗誰.
對抗主和主的受高者.
這是詩篇第二篇.
然後他們27節就說.
希律本丟彼拉多.
釘耶穌的羅馬官.
外邦人以色列民.
果然是在城裡聚集攻擊誰.
就是你的聖僕耶穌.
但他們卻是做你手裡和你指裡所預定的事.
在十字架上.
將人所廢棄的石頭.
上帝卻簡選他成就自己的計劃.
主啊現在求你鑒察他們的威嚇.
使你的僕人能夠放膽的講你的道.

$^{481}$伸出你的手.
讓醫治神蹟的騎士.
醫治聖靈的工作.
能夠藉著你的聖僕耶穌的名.
能夠行出來.
他們禱告之後就地震了.
他們怎麼祈禱.
首先他們引用詩篇第二篇.
詩篇第二篇就是一篇君王royal psalm.
君王的詩篇.
其實就是講到上帝簡選地上的一個君王.
大衛的子孫.
就是坐在他的寶座上.
地上的國度上.
能夠代表上帝來管理萬物.
並且這裡又應許什麼.
就是一個最終極的大衛的君王.
一定會得到什麼.
詩篇第二篇就是傳遞的權柄.
這位是萬世的君王.
當然一直下來.
由詩篇第二篇舊約的歷史裡面.
有沒有一個這樣的君王出現.
有的有的.
大衛所羅門.
曾經威風過一陣.
但接著就敗家.
然後就分裂南北國.
南北國也沒有了.
所以詩篇第二篇的那些說話.
上帝對大衛的子孫的應許.
好像是落空了.
但到了新約的時候.
或者新約之前的第二聖殿的時期.
那些猶太人就去掙扎一件事.
到底上帝的話語是不是落空.
是不是憑空說的.
還是上帝一定會照他的說話.
來去成就他自己的東西.
結果他們在耶穌來之前.

$^{521}$他們就在想.
到底這位終極的大衛的君王.
這位彌賽亞會不會來.
他來到的時候.
那個景象會是怎麼樣的.
到底哪一個是符合這個景象.
能夠起死回生.
能夠成為一個終極的彌賽亞.
所以當時猶太人有很多這樣的討論.
結果耶穌來到.
死而復活.
成就上帝救了人之後.
升天去作王.
初期的教會使徒就說.
一定是耶穌基督.
是符合所有這些景象.
舊約裡面所描述的君王.
終極的彌賽亞.
一定是耶穌基督.
只有耶穌基督才會符合這個景象.
所以他們就認定這位主耶穌.
就是地上萬國的君王.
因為可能外邦仍然會有謀算的事.
並且是不服於受告者.
但福音一定要繼續傳出去.
因為這是上帝的應許.
並且已經第一步落成了.
在耶穌基督的身上.
教會就繼續延續這個工作.
這個神國的工程.
將上帝的名字.
耶穌的名字能夠傳開去.
就因為他們這樣的信念.
他們就回應當時的危機.
就不是我們自然的回應.
我們自然的回應.
就是收一收.
不要做那麼多.
但他們是看定上帝掌權.
連地上的君王本雕彼拉多.

$^{561}$或者希律.
猶太人的權貴.
都不會走出上帝的五指山.
因為上帝是掌管一切的.
所以他們求的是什麼.
不是求安寧.
相反的是反感應.
他們求上帝的是更加有膽量.
更加有力量.
這裡很明顯他們看見自己的身份是什麼.
四章第29節.
求你監察他們的威嚇.
上帝監察和審判他們.
不過我們要做什麼.
你的僕人求主.
你讓我們放膽去傳你的道.
作你自己的工作.
這是一個身份的認定.
是他們明白他們自己是上帝的國民.
是耶穌的門徒.
是做他們身份對稱的事.
我們的身份是什麼.
如果今天有人問你.
你的身份是什麼.
你可能會告訴他.
我的組織在哪裡.
我的工作是什麼.
或者我家庭的環境關係是怎樣.
這些都是我們的身份.
但是對於基督徒來說.
我們的身份是什麼.
是上帝的國民.
四章第一章我們是上帝的見證.
由第一章開始直到最後的時候.
都是一班教會認定自己是上帝的門徒.
所以他們身份的緣故.
他們做出一些和他們身份相稱的事.
大家知道在今年有很多地方水浸.
特別是美國德州.
德州的恐怖很可怕.

$^{601}$真的很可怕.
這些天災.
我們人根本不能駕馭.
但是看這些故事的時候.
除了留意有多可怕之外.
我們可能留意一下.
裡面可能有一些鼓勵人的故事.
其中一個故事.
就是講到一個26歲的美國的海軍.
他的名字叫做Scott Ruskin.
這個26歲的海軍.
他剛剛做海軍.
沒有做過任何任務.
誰知道就被人叫上.
就是新人.
他是第一個去到現場.
他去做新人.
沒有什麼經驗.
但是他很勇敢地去救援那些小朋友.
他就坐直升機進去災區.
然後冒著大雨.
你知道有多危險嗎.
河流不只是水.
還有很多雜物和泥等等.
他就坐直升機一路來往.
然後冒著大雨.
親自去救這些小朋友.
結果他救了165個小朋友.
165個小朋友.
Morning America去訪問他.
任務完成之後.
他有什麼要講.
他說什麼呢.
他說.
This is why we do the job.
這根本就是我們工作的本份.
這句話就是一個identity的說話.
是一個海軍.
他的identity.
他的名分.

$^{641}$他的身份.
就是去救人的.
This is why we are here.
Right.
This is why we are here.
他所強調的就是.
心底的身份的認同.
去保護去救援.
心理學家告訴我們.
我們的行為.
我們做的事.
無論是吃飯.
衣食住行.
我們所做的動作等等.
其實大概有兩個導向.
第一個導向.
就是以結果作導向.
來引導我們行為的思維.
用一個結果來說.
那個outcome是什麼.
然後第二個導向.
就是從我們的身份開始作導向.
你說什麼.
我不太明白.
我做一個很簡單的表.
左邊的就是那些.
以結果導向作思維.
就是什麼呢.
我要吃得健康一點.
通常你都會失敗的.
因為我要做某樣東西.
我要do something.
我想讀一本書.
你永遠都讀不完那本書.
我想跑馬拉松.
永遠都跑不到.
我想變得沒那麼焦慮.
永遠都更加焦慮.
我想更準時.
誰知道我永遠都是遲到.

$^{681}$想理財有道.
誰知道理財失敗.
等等等等.
這些是結果導向的思維.
因為我們想某個outcome出現.
但通常這個就很難sustain到.
很難維持到.
但一個身份導向就不同了.
身份導向就不是說.
我想吃得健康一點.
身份導向就說.
我是一個講究健康的人.
我不像你們那班人.
ok.
吃一些辣辣的東西.
因為我的身份的緣故.
我就會揀選健康的食料.
你明白我的意思嗎.
我是一個有深度思想的人.
所以我不像你們那班人.
只看tiktok.
看一些短的video.
我去看一些solid的書.
ok.
我是一個跑手.
我是一個跑者.
所以我不像你們這樣.
走兩個block就累了.
因為我是會跑馬拉松.
是一個身份的認同.
我是一個沒有好品格的人.
基督給我一個好的品格.
我不會像你這樣.
整天亂發脾氣.
你明不明白.
我的思維是完全不同的.
我是一個有誠信的人.
所以我不像你們那班人.
經常遲到.
ok.

$^{721}$因為我想做一個可靠的人.
我是體會過上帝給我的祝福.
這個是我的identity.
所以我理財就有道.
對嗎.
根據上帝給我的智慧.
我是一個珍惜上帝的說話的人.
所以為什麼我會看聖經.
我會去明白聖經.
我是一個真實主的門徒.
所以我為主事奉.
我是一個宣教士.
所以我會傳福音.
不是因為我想做一些事.
而是我的being define我的doing.
我的being會define我的doing.
這個就是史壇傳裡面所說的.
教會的身份.
我們僕人能夠有勇氣繼續去傳福音.
這個是我的身份.
發自我的身份.
如果我們教會沒有一個這樣的self understanding.
沒有一個個人的身份的明白的話.
我們做什麼都是假的.
做著做著我們就會洩氣.
但是因為你的身份轉變了.
我是基督的門徒.
所以我會去關懷.
所以我會跨出我自己的地步.
我會接觸其他人.
我能夠去傳福音.
不是因為我想做.
這是我自己.
這個就是我自己.
剛才我提過我們昨天.
有宗教會三四十個弟兄姊妹.
我們一起出隊.
有很多都是我們UV的.
真的很感恩.
這次我覺得特別氣氛和之前是好很多的.

$^{761}$為什麼呢.
我體會到大家對於自己的身份的明白.
This is who I am.
我出去代表福音.
在今年裡面.
我希望大家在聖年裡面.
我們教會作為一個群體.
我們明白我們自己是誰.
我們去到社區不同的地方.
去到我們生命不同的角落的時候.
This is who I am.
為什麼我工作上有交代.
因為我是基督的門徒.
我是屬於基督的人.
我是代表基督和福音的.
為什麼我會和一些未信主的人.
我會嘗試去接觸他們.
This is who I am.
It's not what I do.
It's who I am.
我們同心低頭祈禱.
求主幫助我們.
讓我們身份在神的國裡面.
能夠得到更多的認同.
因為你呼召我們每一個.
因為耶穌基督像人所棄的城樓.
像人所凋棄的.
西人所凋棄的.
上帝在你的智慧裡面卻簡選.
讓基督在十字架墳墓裡面.
成就這個舊人.
將我們救贖.
將我們身份完全的顛覆到反轉.
主我們是你的子民.
求你教導我們怎樣生活.
我們是你的門徒.
求你教導我們怎樣在你的話語裡面.
扎好根基.
在我們禱告裡面.
我們的導向是求上帝.

$^{801}$給我們更多能力.
以致我們身份能夠展現出來.
幫助我們教會.
讓我們能夠有一個很清楚的.
明白我們就是你的僕人.
我們就是你的群體.
我們就是代表著你.
在我們的敬拜.
在我們事奉.
在我們的團計小組.
在我們所作的一切裡面.
由生命裡面發出來.
嶄新的身份.
讓我們能夠成為你自己國民的一部分.
求主你使用我們.
讓我們繼續在今年裡面.
我們跨越上師去事奉你.
因為This is who we are.
讓我們活出我們自己的身份.
這個生命.
能夠歸榮於你.
祈禱奉主明教.
元上帝祝福各位.
\newpage



\section{6:9-15}
\label{sec:6C_6zjSlvhM}
\textbf{A Model for Prayer (Matthew 6\_9-15) - Chee Yie Jong}
\newline
\newline
連結: \href{https://youtube.com/watch?v=6C_6zjSlvhM}{\texttt{https://youtube.com/watch?v=6C\_6zjSlvhM}} ~~~~ 語音日期: 2025-08-24
\newline
\newline
\hyperref[sec:_p6p1yxlKsI]{\small{< < < PREV SERMON < < <}}
~
\hyperref[sec:index_chronic]{\small{[返順時目]}}
~
\hyperref[sec:index_scriptual]{\small{[返順卷目]}}
~
\hyperref[sec:HQs8juXt_Kg]{\small{> > > NEXT SERMON > > >}}
\newline
\newline
$^{1}$Good morning, UV..
It's good to be with you all again today.
and just opening up the word of God..
Okay, cool..
All right, so compared to the generations.
that come before us, I would say that we live.
in a world of extreme convenience, right?.
What do I mean by that, a world of extreme convenience?.
Well, if you have a question about something.
or you need information about a specific topic,.
you don't have to go through the hassle.
of figuring it out on your own right now..
You can just pull out your phone, put out your computer,.
and just type the questions in Google or chat GPT now,.
and you can probably get your answers instantly as well..
Or if you wake up in the morning and find out that,.
oh, I'm actually running out of my coffee beans,.
you can just put an order in Amazon.
and have it delivered overnight.
and sometimes even on the same day as well, right?.
So not only are we able to have a lot of information.
on our fingertips or get the things that we want.
as quickly as possible now,.
but we can also get the things that we want.
any time of the day now..
So as I was preparing for my sermon this week,.
I was curious about where this phrase 24/7 comes from..
And as a perfect example of living.
in this world of extreme convenience,.
obviously I did some quick research on Google..
And guess what?.
This phrase actually originated in the United States.
in the late 20th century,.
originally in the world of sports..
But since then, we see an increasing amount of businesses.
that is modeled after this idea of 24 hours a day,.
seven days a week, like 7-Eleven..
Fast food restaurants like McDonald's.
or ride share business like Uber or Lyft..
Or customer service hotlines,.

$^{41}$online businesses, et cetera, et cetera..
However, way before this idea of 24/7 gains popularity.
and became the norm of our society,.
there's actually a group of Christians called the Moravian.
who started a 24/7 prayer watch,.
not just for a day or a week or a month.
or a year or a decade..
But on August 27, 1727,.
this group of 24 men and 24 women.
from this Moravian community of Hernhut in Saxony, Germany,.
they committed to pray in hourly shifts,.
covering each and every hours of the day continuously..
And this around-the-clock prayer change.
actually lasted for more than 100 years,.
sparking a revival within this community.
and transforming this tiny village of Hernhut.
into one of the most important hubs.
for Protestant global mission movement in the 18th century..
And they also laid the foundation.
for our current modern mission movement as well..
So if you were at my previous sermon here,.
you might remember this question that I asked all of us..
How is your relationship with the Lord.
when it comes to your prayer life?.
Of course, I know the history.
of the Moravian prayer movement is a unique case,.
and we have never seen anything like that since then.
in our church history..
However, like the prayers of the Moravians,.
which turned this community initially marked by conflict.
and division into a place of revival.
and a place for global missions,.
our prayer is also the field behind everything that we do.
in our personal walk with Jesus,.
as well as the things that we do for God's kingdom.
as followers of God..
So personally, I am grateful to the people.
who have taught me and show me how to pray..
Of course, there's my parents.
who taught me how to pray as a kid,.

$^{81}$and also my mom who continuously to show me.
what a faithful prayer life looks like right now..
And there's also one of my pastor, Pastor Rafe,.
who modeled for me what it means.
to sit before the Lord in silence..
We used to meet early in the morning in his office,.
and I'm usually the first thing in his work calendar..
So before we talk about anything,.
he would just kneel down on the floor in his office.
and spend about five to 10 minutes just to listen,.
listen before the Lord, sitting still before the Lord..
And every time he would invite me to join him as well.
in his silent prayers..
I'm also thankful for my friend, Hailey,.
who always reminds me to turn to God first things first..
A couple years ago, I was in the seasons of waiting.
and just a lot of uncertainties..
And one time as I was catching up with her over a run.
and share about some of these things that I'm worried about,.
she just turned to me and say,.
"Hey, why don't we just pray right now?".
So she ended up praying over me.
while we were running around in the park..
And I'm also thankful for all the different community,.
the groups, my small groups,.
that I got to be a part of in the past..
These people have shown me what it looks like.
to labor in prayers for one another,.
to celebrate together when our prayers are being answered,.
and to wrestle with the Lord.
when our prayers are not being answered..
And sometimes to just sit in the unknown.
while waiting for God to respond to our prayer requests..
So personally, I'm in debt to these people.
who have modeled for me.
what a vibrant prayer life looks like..
But today we have the opportunity to learn.
how to pray from the master himself, Jesus..
In some ways, our passage today.
is a glimpse into Jesus's own prayer life..

$^{121}$All throughout Jesus's ministry,.
we always see that he would retreat to some quiet places.
to pray alone, but we don't have much information.
about the actual content of his prayer..
So our passage today is sort of like.
a going behind the scenes.
and taking a peek into Jesus's own prayer life.
and how he communes with God the Father..
So to better situate our passage today.
within the larger context of the Sermon on the Mount,.
this model prayer that Jesus taught his disciple.
comes right after he tells them how not to pray..
So first, Jesus tells them not to pray like the hypocrites.
who love praying standing in the synagogues.
or on the street corner.
in order to be seen by other people..
So Jesus warns again this kind of self-gratifying prayer,.
first and foremost..
And second, Jesus teaches them not to pray like the pagan.
who just keep on babbling with words.
with the hope that one of the small letter,.
quote unquote, "gods" will share their prayer..
So Jesus also warns again this kind of superstitious prayer..
And coming right after these two negative examples.
of how not to pray in a self-gratifying.
and also in a superstitious way,.
Jesus continues in Matthew 6, verse 9..
"This then is how you should pray.".
So let's look at this model prayer from Jesus together..
And I'm sure a lot of you have probably memorized this.
before as a kid..
So I want all of us to read this together.
as a church this morning..
Can we do that?.
I have it on the screen as well..
So let's begin..
"This then is how you should pray..
"Our Father in heaven,.
"hallowed be your name..
"Your kingdom come,.

$^{161}$"your will be done on earth as it is in heaven..
"Give us today our daily bread.
"and forgive us our debts.
"as we also have forgiven our debtors..
"And lead us not into temptation,.
"but deliver us from the evil.".
Okay, so notice how Jesus begins his prayer..
"Our Father.".
He begins with these two words..
So in contrast to the self-gratifying prayer.
of the hypocrites, which begins and ends.
with themself in mind,.
these prayers actually begins with God..
And there's so much more theology.
that is just associated with these two words..
"Our Father.".
However, I want to take some time.
just to highlight a few for us today..
So Father is obviously a familial language, right?.
So what Jesus is teaching here.
is that the God of the universe,.
the creator of all things,.
the author of life,.
he is not just a father,.
but he is our father..
This means that when we come before God in our prayer,.
this relationship is not just between a God and his people,.
but this relationship is between a father.
and his beloved children..
Of course, I want to acknowledge.
that when it comes to our relationship.
with our earthly father,.
we all have different kind of experience, right?.
Some are good, some are bad,.
some might just be indifferent..
And a lot of time,.
our experience with our earthly fathers.
will inevitably affect how we view our heavenly father..
However, the good news that we know from the Bible.
is that if we read the Bible from cover to cover,.

$^{201}$we consistently see the loving.
and compassionate character of our father.
on full display throughout these pages,.
with the ultimate examples.
being God sending his own son into the world.
to die for a group of people who don't really deserve it..
So regardless of how your relationship.
with your earthly father is or was,.
don't let that become a deterrent for you.
to come before your heavenly father in prayer..
And when we begin our prayer with these words,.
our father in heaven,.
we pray knowing who God is, our father,.
and who we are in relationship to him,.
his beloved children..
To put it another way,.
when we knock on heaven's door with our prayer,.
we hold our ID as co-heirs with Jesus himself..
And every heavenly inheritance that is available to Jesus,.
we get to have a part of it as well..
So our father in heaven..
So after knowing who God is.
and who we are in relationship to him,.
we can bring our request and our petition before the Lord..
So the next line, "Hallowed be your name.".
Hallowed, obviously, is an old English word, right?.
I don't think any one of us would use this word.
in our day to day anymore..
But what it means is simply just to make holy.
or to set apart..
So this is the first petition in the Lord's prayer,.
and it begins with worship and adoration towards God..
May your name be kept holy..
In other words, it is a God-oriented petition..
How often in our prayer do we jump straight.
into lifting not the name of God,.
but other names and idols,.
such as success, wealth, comfort, security, or approval?.
However, right here,.
Jesus shows us a very different way to pray,.

$^{241}$one that begins with worship..
Hallowed be your name, or may your name be kept holy..
There are two implied subjects in this petition,.
the first one being us..
May your name be kept holy and sacred by us as his people..
So holiness is at the core of who God is,.
and as Jesus' disciple,.
his holiness should also be visible.
and evident in our day to day life..
And the second implied subject here is God himself,.
that God himself will make the holiness.
that comes along with his name.
to be manifested throughout the entire world..
And so when we pray,.
remember that hallowed be your name.
is an invitation for us to reorient the affections.
of our hearts towards God..
And everything that comes in the Lord's prayer after this.
is an overflow of this kind of heart posture.
before the Lord, one of worship, one of adoration..
So let's continue..
May your kingdom come..
So after we lift up our adoration towards God,.
we find this second God-oriented petition..
May your kingdom come..
This second petition calls for the fullness.
of God's kingdom to come,.
and it highlights the already,.
but not yet aspects of God's kingdom..
Put it another way, this kingdom of heaven.
has already begun in the first coming of Jesus,.
hence the already part..
But we can only experience the fullness.
of this heavenly kingdom.
when Jesus comes again in the future,.
hence the not yet part of it..
And then Jesus continues..
Your kingdom come, your will be done.
on earth as it is in heaven..
So what is God's will?.

$^{281}$God's will is his divine purpose for each of his people.
and also for the rest of the world..
And so this third petition, your will be done,.
calls for the fullness of God's will.
that is currently only present in the heaven.
to be fully manifested on earth as well..
So a couple of years ago,.
I was in the seasons of waiting and uncertainties, right?.
And the Lord Prayers became one of my go-to prayer..
But every time I came to this line,.
specifically, your will be done,.
my entire body would just tense up.
because I understand what it means.
to truly pray these words..
If I pray for God's will to be done,.
what that means is that I would have to give up.
my own preference, my own plan, and also my own kingdom..
However, if we truly want to see heaven meets earth,.
a posture of surrender is required.
in order for God's will to be done in our life,.
in our community, in our city, in our church,.
and also in the rest of the world..
So prayer is the place where the heavenly realities.
began to invade on our earth..
In prayer, we call upon the peace and justice of God.
to push back against the darkness of our world..
In prayer, we reach for the heavenly resource.
and distribute them to places in our world.
that is being destroyed by sin and brokenness..
So, on earth as it is in heaven is an intercessory prayer..
When we pray for his kingdom to come.
and for his will to be done on earth as it is in heaven,.
we are asking God himself to intercede.
in the realities of our broken world.
with his kingdom realities..
So I love how this author put it..
When we engage in intercessory prayer,.
we are loving others on the basis of heaven's resource..
So, when we pray for the sick,.
we are reaching for God's divine power to heal..

$^{321}$When we pray for wars and conflicts.
throughout our world right now,.
we are reaching for God's heavenly peace..
And when we pray for the non-believers in our lives,.
we are reaching for the power of gospel to transform life..
So, prayer is where heavens meets earth..
And in the Lord's prayer, this line on earth.
as it is in prayer also functions as a link.
between the reality of God.
and also the reality of our world..
If you pay attention, these first three petitions.
is all about God, right?.
Your name, your kingdom, and your will..
And so remember, first, prayer begins with God and not us..
It is only when we are able to give God the honor.
and the glory that is due to his name.
and also align our priorities in accordance to his will.
that we are able to know the right thing to ask for.
in our prayer..
So heaven is the source of all power.
and earth is the place where this power.
is being made visible..
So the following three petitions will finally be about us..
Give us today our daily bread..
So this is the first we or us-oriented petition.
and it is a prayer of dependence on God.
for our daily sustenance..
For most of us, this daily bread prayer.
is probably something that we don't quite pray.
as often nowadays, right?.
When you go to the grocery store or Costco,.
you don't just buy enough grocery to last you for a day,.
but I'm sure most of us will usually get enough.
to get through us for at least a week, right?.
But for most of Jesus' audience,.
they don't experience the same kind of abundance.
and wealth as we do today..
Daily bread for them is not always guaranteed..
However, because of our abundance,.
we slowly lose our dependence on God.

$^{361}$and become more and more reluctant to ask God.
for our basic day-to-day small requests..
And because of our abundance,.
we slowly lose our sense of gratitude.
and begin to take things for granted..
The food on our table, the air that we breathe,.
or the roof that is over our head,.
and the list goes on and on..
So back to the seasons of waiting.
that I was in a couple years ago..
One of the things that the Lord really revealed to me.
in that season was that I became so used.
to praying safe prayers..
And what do I mean by that?.
There were things that I was not willing.
to bring before the Lord because I decided for myself.
that God probably would not meet those needs..
And so I would rather not pray about those things at all.
than to bring them before the Lord.
and then have a chance to get disappointed by God.
when he didn't answer it..
Or there would be other needs that I felt like.
they were just too insignificant to even bring.
before the God of the universe..
So I would just resort to some vague or general prayers..
But after that, I made an effort to pray more vulnerably.
and ask God for more specific things,.
allowing myself to either get surprised by God.
or get disappointed by God..
I stopped filtering my own prayer.
and began to pray more freely and more boldly..
And of course, I'm not going to pretend.
that everything went on smoothly after that..
There were things in that season.
that I was really hoping for God to come true,.
but they never happened..
But at the same time, when God show up.
in the exact way that I prayed for,.
I know for sure that it was God himself..
And because of this, my relationship with God.

$^{401}$actually becomes so much more intimate now..
So church, daily bread prayer is an exercise of trust,.
trusting that the all-knowing,.
all-powerful God is in control.
and we can fully entrust our lives into his hand,.
down to the minute and trivial details of our lives..
So don't be afraid to bring your ordinary,.
everyday requests before God,.
even if you feel like they're not important enough.
to even pray about..
And from prayer for our daily bread,.
we move on to the next petition..
And forgive us our debts,.
as we also have forgiven our debtors..
So this is the second petition that is oriented towards us..
And it shifts from our physical need.
now to our spiritual need..
And when Jesus mentioned debts here,.
he's referring to our debts to God..
And therefore, what he means is just.
forgiveness of our sins,.
the way that we have sins against God..
And so when we pray, forgive us of our debts,.
this is a confessional prayer..
Because before we can ask for forgiveness,.
we need to be able to name our own sins..
First John 1, I say, if we confess our sins,.
he is faithful and just,.
and he will forgive us of our sins.
and cleanse us from all unrighteousness..
But pay attention, Jesus doesn't just stop here..
He says, forgive us our debts,.
as we also have forgiven our debtors..
Did you hear that?.
Because of our experience of divine forgiveness,.
it should allow us a greater willingness.
to forgive those who have sinned against us,.
or those who have hurt us..
Later in Matthew 18, Jesus told of a parable.
of an unforgiving servant,.

$^{441}$who owe his master 10,000 talents..
So a talent is almost 20 years worth of wages.
for a day laborer during Jesus' times..
But the master forgave all 10,000 talents of debt.
for this servant..
But when this unforgiving servant got free,.
he saw one of his fellow servants.
who owe him 100 denarii,.
which is only about 100 days worth of wages..
So you compare 100 days worth of wages.
to 20 times 10,000 days worth of wages, right?.
And yet this servant, he refused to cancel the debt.
and put his fellow servants into prison..
And as you can imagine, when the master found out,.
he was furious..
He summoned this unforgiving servant.
and called him a wicked servant,.
because he refused to show mercy to his fellow servant,.
who owe him way less than the mercy.
that he has received himself..
And this wicked servant ended up being put into jail himself.
because of that..
And Jesus end this parable with this warning..
This is what your heavenly father would do to you.
if you refuse to forgive your brothers and sister.
from your heart..
And in fact, this is so important,.
that right after the Lord's prayer,.
in verse 14 and 15, Jesus added this remark..
For if you forgive other people when they sin against you,.
your heavenly father will also forgive you..
But if you do not forgive others their sin,.
your father will not forgive your sins..
So when we receive forgiveness from God,.
it should transform our heart.
and enable us to forgive those who have sinned against us.
in a way lesser extent than we have sinned against God..
This is because being a follower of Jesus.
is not just about our vertical relationship with him,.
but it is also about our horizontal relationship.

$^{481}$with other people made in the image of God..
And then we have the final petition in this Lord's prayer..
Lead us not into temptation,.
but deliver us from the evil one..
This petition is one of deliverance,.
specifically from the evil one..
And one of the realities of living in this fallen world.
is that temptation are just inevitable..
But as God's children,.
he promised that he would not allow us to be tempted.
beyond what we can endure,.
and he would always provide a way out for us..
And so a better way to understand this petition.
is I like how the NLT translate it..
Do not let us yield into temptation..
So this last petition is not a request.
to be completely immune from all sorts of temptation,.
but it is a prayer to God to give us the strength.
to overcome them when we are faced with various temptation..
So we've come to the end of the Lord's prayer,.
and this model prayer from Jesus is straightforward enough.
that even a child can memorize, can recite, can pray it..
But at the same time, there is so much depth to it.
that we can spend a whole series just diving into it..
So as I was preparing for my sermon this week,.
I went back and forth on how I want to end my sermon,.
because there is so much more.
that we didn't even touch upon..
However, I feel like I would do a disservice to all of you.
if I didn't give all of you some extended time.
just to sit with this prayer,.
and to even spend some time praying through it yourself,.
because the best way to allow the Word of God.
to transform us is not just to sit and listen to a sermon,.
but to actually live it out, to practice it..
So to give a little guide,.
I kind of created this framework.
based on the Lord's prayer that I hope will help.
all of us pray through the Lord's prayer.
in five different movement..

$^{521}$So there's worship, there's intercession,.
there is petition, there is confession and reconciliation,.
and there's also deliverance..
And during our time of reflection,.
pick one or maybe two of these areas.
which are most relevant in your personal life right now..
Maybe for some of you,.
there are things that you need to bring before the Lord.
and confess before Him..
And maybe for some of you,.
there are sins and struggles in your life.
that you need to pray to God and ask for deliverance..
And maybe for the rest of you,.
you just need some extended time to worship,.
to fix your eyes upon God's beauty and His glory..
So I would encourage all of us just to take some time.
to listen to the Holy Spirit..
I'll have Amanda out just to play.
some background music as well..
And before you head out later,.
I actually also printed out a more extended version.
of each of these movements that you can take a copy home.
and spend some time this week to pray through it..
Yeah, I will have, I think it is on the table.
right at the back..
So on your way out, feel free to grab a copy..
And my hope is that you will spend some time this week.
in your personal devotional time.
just to pray through them as well..
And so as we spend some time in prayer and reflection,.
feel free to move around the room if you need to..
If you would like a more private space,.
you can just get off your seat and go to a place.
that you feel like you're more comfortable praying in..
Or if you come with someone else,.
you can also get into a group of two or three.
and just pray together..
Or husband and wife, feel free to also get together.
and pray as a couple..
Or if you know of people in the rooms.

$^{561}$that have a specific needs that you can pray for,.
feel free to go over to them,.
lay your hand over them and pray over them..
Okay, so I want this time to be filled with prayer.
and no structure at all..
And after a few moments, I will come up here.
and end our time of prayer together..
So church, let's spend the next five minutes or so.
just praying..
(gentle music).
(gentle music continues).
Church, let's pray together..
Our Heavenly Father, we praise your holy name..
Thank you for adopting us into your family.
and give us the honor to be called sons and daughters.
of God through Jesus Christ..
May your kingdom come and your will be done.
on earth as it is in heaven,.
in urban voice as it is in heaven,.
in Chicago as it is in heaven..
(gentle music continues).
While we experience brokenness and suffering.
on this side of heaven,.
we know that this is not what you intended in creation..
And so we long for you to come again,.
to bring your perfect peace.
and your perfect wholeness to earth..
Father, provide us with our daily needs,.
both big and small..
Forgive us of our sins as we forgive those.
who sins against us..
Help us to not yield to temptation.
because we know that we are not just fighting.
against flesh or blood,.
but we're fighting against evil rulers.
and authorities of the unseen world,.
against mighty power in this dark world.
and against evil spirits in the heavenly places..
So give us strength to overcome the evil ones..
For yours is the kingdom and the power and the glory..

$^{601}$Amen..
All right, thank you all..
Just a reminder on your way out,.
feel free to grab a copy of,.
it is kind of like a daily devotions.
that will help you pray through these five movement.
based on the Lord's prayer..
So grab a copy home..
I have 20 copies, so we should have enough..
And we will also have,.
okay, we don't know if there's refreshment out there,.
so you can go out and get surprised,.
but that's pretty much it..
Okay, thank you all for coming,.
for joining us for our worship..
Go in peace, your love..
\newpage



\section{使徒行傳 2:1-47}
\label{sec:HQs8juXt_Kg}
\textbf{Contagious Church (Acts 2) - Rev. Mark Chan 有感染力的教會 (使徒行傳2\_1-47) - 陳崇基牧師}
\newline
\newline
連結: \href{https://youtube.com/watch?v=HQs8juXt-Kg}{\texttt{https://youtube.com/watch?v=HQs8juXt-Kg}} ~~~~ 語音日期: 2025-08-24
\newline
\newline
\hyperref[sec:6C_6zjSlvhM]{\small{< < < PREV SERMON < < <}}
~
\hyperref[sec:index_chronic]{\small{[返順時目]}}
~
\hyperref[sec:index_scriptual]{\small{[返順卷目]}}
~
\hyperref[sec:J7Rxq8y7uj4]{\small{> > > NEXT SERMON > > >}}
\newline
\newline
使徒行傳 2:1-47
\newline
\begin{longtable}{cl}
\hline
\hline
章節 & 經文 (和合本修訂版)\\
\hline
2:1 & \begin{tabularx}{0.7\textwidth}{X} 五旬節那日到了，他們全都聚集在一起。 \end{tabularx} \\ \\ \relax
2:2 & \begin{tabularx}{0.7\textwidth}{X} 忽然，有響聲從天上下來，好像一陣大風吹過，充滿了他們所坐的整座屋子； \end{tabularx} \\ \\ \relax
2:3 & \begin{tabularx}{0.7\textwidth}{X} 又有舌頭如火焰向他們顯現，分開落在他們每個人身上。 \end{tabularx} \\ \\ \relax
2:4 & \begin{tabularx}{0.7\textwidth}{X} 他們都被聖靈充滿，就按著聖靈所賜的口才說起別國的話來。 \end{tabularx} \\ \\ \relax
2:5 & \begin{tabularx}{0.7\textwidth}{X} 那時，有從天下各國來的虔誠的猶太人，住在耶路撒冷。 \end{tabularx} \\ \\ \relax
2:6 & \begin{tabularx}{0.7\textwidth}{X} 這聲音一響，許多人都來聚集，各人因為聽見門徒用他們各自的鄉談說話，就甚納悶， \end{tabularx} \\ \\ \relax
2:7 & \begin{tabularx}{0.7\textwidth}{X} 都詫異驚奇說：「看哪，這些說話的不都是加利利人嗎？ \end{tabularx} \\ \\ \relax
2:8 & \begin{tabularx}{0.7\textwidth}{X} 我們每個人怎麼聽見他們說我們生來所用的鄉談呢？ \end{tabularx} \\ \\ \relax
2:9 & \begin{tabularx}{0.7\textwidth}{X} 我們帕提亞人、瑪代人、以攔人，和住在美索不達米亞、猶太、加帕多家、本都、亞細亞、 \end{tabularx} \\ \\ \relax
2:10 & \begin{tabularx}{0.7\textwidth}{X} 弗呂家、旁非利亞、埃及的人，並靠近古利奈的利比亞一帶地方的人，僑居的羅馬人， \end{tabularx} \\ \\ \relax
2:11 & \begin{tabularx}{0.7\textwidth}{X} 包括猶太人和皈依猶太教的人， 克里特人和阿拉伯人，都聽見他們用我們的鄉談講論神的大作為。」 \end{tabularx} \\ \\ \relax
2:12 & \begin{tabularx}{0.7\textwidth}{X} 眾人就都驚奇困惑，彼此說：「這是甚麼意思呢？」 \end{tabularx} \\ \\ \relax
2:13 & \begin{tabularx}{0.7\textwidth}{X} 還有人譏誚，說：「他們是灌滿了新酒吧！」 \end{tabularx} \\ \\ \relax
2:14 & \begin{tabularx}{0.7\textwidth}{X} 彼得和十一個使徒站起來，他就高聲向眾人說：「猶太人和所有住在耶路撒冷的人哪，這件事你們要知道，要側耳聽我的話。 \end{tabularx} \\ \\ \relax
2:15 & \begin{tabularx}{0.7\textwidth}{X} 這些人並不像你們所想的喝醉了，因為現在才早晨九點鐘。 \end{tabularx} \\ \\ \relax
2:16 & \begin{tabularx}{0.7\textwidth}{X} 這正是藉著先知約珥所說的： \end{tabularx} \\ \\ \relax
2:17 & \begin{tabularx}{0.7\textwidth}{X} 『神說：在末後的日子，我要將我的靈澆灌凡血肉之軀的。你們的兒女要說預言；你們的少年要見異象；你們的老人要做異夢。 \end{tabularx} \\ \\ \relax
2:18 & \begin{tabularx}{0.7\textwidth}{X} 在那些日子，我要把我的靈澆灌，甚至給我的僕人和婢女，他們要說預言。 \end{tabularx} \\ \\ \relax
2:19 & \begin{tabularx}{0.7\textwidth}{X} 在天上，我要顯出奇事，在地下，我要顯出神蹟，有血，有火，有煙霧。 \end{tabularx} \\ \\ \relax
2:20 & \begin{tabularx}{0.7\textwidth}{X} 太陽要變為黑暗，月亮要變為血，這都在主大而光榮的日子未到以前。 \end{tabularx} \\ \\ \relax
2:21 & \begin{tabularx}{0.7\textwidth}{X} 那時，凡求告主名的都必得救。』 \end{tabularx} \\ \\ \relax
2:22 & \begin{tabularx}{0.7\textwidth}{X} 「以色列人哪，你們要聽我這些話：拿撒勒人耶穌就是神以異能、奇事、神蹟向你們證明出來的人，這些事是神藉著他在你們中間施行，正如你們自己知道的。 \end{tabularx} \\ \\ \relax
2:23 & \begin{tabularx}{0.7\textwidth}{X} 他既按著神確定的旨意和預知被交與人，你們就藉著不法之人的手把他釘在十字架上，殺了。 \end{tabularx} \\ \\ \relax
2:24 & \begin{tabularx}{0.7\textwidth}{X} 神卻將死的痛苦解除，使他復活了，因為他原不能被死拘禁。 \end{tabularx} \\ \\ \relax
2:25 & \begin{tabularx}{0.7\textwidth}{X} 大衛指著他說：『我看見主常在我眼前，他在我右邊，使我不至於動搖。 \end{tabularx} \\ \\ \relax
2:26 & \begin{tabularx}{0.7\textwidth}{X} 所以我心裡歡喜，我的舌頭快樂，而且我的肉身要安居在指望中。 \end{tabularx} \\ \\ \relax
2:27 & \begin{tabularx}{0.7\textwidth}{X} 因你必不將我的靈魂撇在陰間，也不讓你的聖者見朽壞。 \end{tabularx} \\ \\ \relax
2:28 & \begin{tabularx}{0.7\textwidth}{X} 你已將生命的道路指示我，必使我在你面前充滿快樂。』 \end{tabularx} \\ \\ \relax
2:29 & \begin{tabularx}{0.7\textwidth}{X} 「諸位弟兄，先祖大衛的事，我可以坦然地對你們說：他死了，也埋葬了，而且他的墳墓直到今日還在我們這裡。 \end{tabularx} \\ \\ \relax
2:30 & \begin{tabularx}{0.7\textwidth}{X} 既然大衛是先知，他知道神曾向他起誓，要從他的後裔中立一位坐在他的寶座上。 \end{tabularx} \\ \\ \relax
2:31 & \begin{tabularx}{0.7\textwidth}{X} 他預先看見了，就講論基督的復活，說：『他不被撇在陰間；他的肉身也不見朽壞。』 \end{tabularx} \\ \\ \relax
2:32 & \begin{tabularx}{0.7\textwidth}{X} 這耶穌，神已經使他復活了，我們都是這事的見證人。 \end{tabularx} \\ \\ \relax
2:33 & \begin{tabularx}{0.7\textwidth}{X} 他既被高舉在神的右邊，又從父受了所應許的聖靈，就把你們所看見所聽見的，澆灌下來。 \end{tabularx} \\ \\ \relax
2:34 & \begin{tabularx}{0.7\textwidth}{X} 大衛並沒有升到天上，但他自己說：『主對我主說：你坐在我的右邊， \end{tabularx} \\ \\ \relax
2:35 & \begin{tabularx}{0.7\textwidth}{X} 等我使你的仇敵作你的腳凳。』 \end{tabularx} \\ \\ \relax
2:36 & \begin{tabularx}{0.7\textwidth}{X} 故此，以色列全家當確實知道，你們釘在十字架上的這位耶穌，神已經立他為主，為基督了。」 \end{tabularx} \\ \\ \relax
2:37 & \begin{tabularx}{0.7\textwidth}{X} 眾人聽見這話，覺得扎心，就對彼得和其餘的使徒說：「諸位弟兄，我們該怎樣做呢？」 \end{tabularx} \\ \\ \relax
2:38 & \begin{tabularx}{0.7\textwidth}{X} 彼得對他們說：「你們各人要悔改，奉耶穌基督的名受洗，使你們的罪得赦免，就會領受所賜的聖靈。 \end{tabularx} \\ \\ \relax
2:39 & \begin{tabularx}{0.7\textwidth}{X} 因為這應許是給你們和你們的兒女，並一切在遠方的人，就是給所有主—我們的神所召來的人。」 \end{tabularx} \\ \\ \relax
2:40 & \begin{tabularx}{0.7\textwidth}{X} 彼得還用更多別的話作見證，勸勉他們說：「你們當救自己脫離這彎曲的世代。」 \end{tabularx} \\ \\ \relax
2:41 & \begin{tabularx}{0.7\textwidth}{X} 於是領受他話的人，都受了洗；那一天，門徒約添了三千人。 \end{tabularx} \\ \\ \relax
2:42 & \begin{tabularx}{0.7\textwidth}{X} 他們都專注於使徒的教導和彼此的團契，擘餅和祈禱。 \end{tabularx} \\ \\ \relax
2:43 & \begin{tabularx}{0.7\textwidth}{X} 眾人都心存敬畏；使徒們又行了許多奇事神蹟。 \end{tabularx} \\ \\ \relax
2:44 & \begin{tabularx}{0.7\textwidth}{X} 信的人都聚在一處，凡物公用， \end{tabularx} \\ \\ \relax
2:45 & \begin{tabularx}{0.7\textwidth}{X} 又賣了田產和家業，照每一個人所需要的分給他們。 \end{tabularx} \\ \\ \relax
2:46 & \begin{tabularx}{0.7\textwidth}{X} 他們天天同心合意恆切地在聖殿裡敬拜，且在家中擘餅，存著歡喜坦誠的心用飯， \end{tabularx} \\ \\ \relax
2:47 & \begin{tabularx}{0.7\textwidth}{X} 讚美神，得全體百姓的喜愛。主將得救的人天天加給他們。 \end{tabularx} \\ \\
[1ex]
\hline
\hline
\end{longtable}
$^{1}$Pastor Mark..
(speaking in foreign language).
- Hey everyone, good morning..
(speaking in foreign language).
I want to do something special today.
because I want to welcome everyone to be with us here.
and we want to appreciate your presence..
(speaking in foreign language).
Go ahead and stand up and find someone,.
find someone you have not said hi to today, okay?.
And go to that person and say, "It's good to see you.".
(speaking in foreign language).
(audience chattering).
(speaking in foreign language).
(audience chattering).
All right guys, well I know you have a lot to talk about..
(speaking in foreign language).
We will have time after the service, all right?.
(speaking in foreign language).
Well, once again, I want to piggyback on Kelly,.
what she said about we need your help to build this church..
(speaking in foreign language).
And of course, we will appreciate the worship team.
for their music practice and their leading us today.
to worship God together..
Let's, let's, let's move around our applause..
(speaking in foreign language).
And then of course, the backstage,.
we have some new system installed.
and so we are trying to work out the transition..
So we want to appreciate each and every one of them.
inside the room, it's very hot in there..
(speaking in foreign language).
And of course, thank you for being with us today..
(speaking in foreign language).
And if you notice our title today,.
the sermon title is called Contagious Church..
(speaking in foreign language).
Now, as you see this title, I want you to think,.
are you a contagious Christian or are we a contagious church?.

$^{41}$(speaking in foreign language).
Yeah, and a church full of contagious Christians.
is full of people who truly experience the grace of God..
(speaking in foreign language).
Now, let me give you two pictures to, you know,.
try to illustrate what it means to be contagious..
(speaking in foreign language).
And the first image I want to show you is a negative image..
(speaking in foreign language).
Now, you remember COVID-19, right?.
We could not hug each other, we cannot touch each other.
and we have to make sure we are people to people.
at least six feet apart..
(speaking in foreign language).
When you were in a room with other people,.
when you heard one cough and you got scared.
because the virus was spread, you know, fast and far.
and it would disrupt our relationship.
and actually it would disrupt our community..
(speaking in foreign language).
And as you can tell, something can be contagious,.
something negative..
Same thing happen when there's something wrong.
in our life like bitterness, it can be contagious..
Cynicism can be contagious or lukewarm faith, you know,.
half-hearted faith or gossip or disunity..
They're like, they're kind of like secondhand smoke.
and one person can change the whole environment..
(speaking in foreign language).
Now, that's the first image..
I want to show you the second image though..
(speaking in foreign language).
So I want to ask you ever been in a room.
with a friend who laughs really loud?.
(speaking in foreign language).
You know, when your friend laughs and you know,.
even you are kind of sad,.
but the laughter will change your spirit..
(speaking in foreign language).
So when you're around people who love God,.

$^{81}$who take God seriously,.
you should also eventually do the same thing.
because that is contagious..
You know, our holiness, our full spirit life.
and our faith can be contagious..
(speaking in foreign language).
Now in Acts 2, we've been studying the book of Acts.
and we see that after the Holy Spirit comes.
and the whole community become a contagious church..
(speaking in foreign language).
And let me read this passage from 42 to 44..
They, meaning the church in Jerusalem, the first church,.
they devote themselves to the apostles teaching.
and to the fellowship and to breaking of bread.
and to prayer and everyone was filled with awe.
and many wonders, miraculous signs.
were done by the apostles..
All the believers were together.
and had everything in common..
They sell the possessions and goods.
and they gave to anyone who had need..
(speaking in foreign language).
Now, first, we need to know how to read the Bible..
(speaking in foreign language).
So when we read this passage,.
we should not think that we have to copycat.
or do exactly what they did..
(speaking in foreign language).
Because this kind of copycat, you know,.
reading and application is kind of too simplistic.
because the church in Jerusalem,.
they worship in a different way than we do.
and they worship in the temple court.
and they worship in the homes of other people.
and there were no place called church that they go to..
Their homes are the church..
So our context is different..
So we cannot just copycat it and do exactly what they did..
And we worship in a different way..
We gather in a public place.

$^{121}$and we cannot simply just copy exactly what they do..
(speaking in foreign language).
So what the Bible describes.
is not necessarily what the Bible prescribes..
(speaking in foreign language).
But we can ask and we should ask,.
what kind of church were they.
in the first century Jerusalem context?.
(speaking in foreign language).
And we should ask,.
what kind of church were they.
in the first century Jerusalem context?.
And how does their church,.
how does this church inspire us.
to be a church of Jesus today?.
(speaking in foreign language).
So here in this passage,.
we see believers, ordinary folks,.
they live in a compelling way for the gospel.
and it was contagious.
and people were being drawn in and being attracted,.
not because the program.
or not because of their something,.
some of the activities that they did,.
it's the way they live in a compelling way.
to tell other that God is good..
(speaking in foreign language).
So let me suggest three things.
that we should pay attention to.
that will inspire us to live a contagious church life.
and Christian life..
(speaking in foreign language).
First and foremost,.
we should be a Christian..
(speaking in foreign language).
And second, we should be a Christian..
(speaking in foreign language).
First and foremost, be devoted..
And secondly, they were generous..
And third, they were alive in worship..

$^{161}$(speaking in foreign language).
Now, verse 42, it says,.
they were devoted themselves.
to the teachings of the apostle,.
to the teaching of the Bible..
And you remember last week,.
I mentioned this point, be committed..
We Christians need to be committed..
Worship, coming together,.
are not occasional attendance..
Never done in a way that is half-hearted..
So they, if you pay attention,.
they were hungry for the word of God..
And just like we are hungry,.
we want some food..
But spiritually, they were hungry.
and they want to go for food,.
which is the spiritual food, the word of God..
(speaking in foreign language).
Like when you're hungry,.
you go to a happy lamb, right?.
(speaking in foreign language).
And you want to devour the food..
(speaking in foreign language).
Now here, the church is hungry for the word of God..
(speaking in foreign language).
And brothers and sisters, I want to,.
you know, just want you to think about,.
you know, we do live in an age of distractions..
(speaking in foreign language).
Now everyone is busy, everyone is rushing,.
and we're trying to handle a lot of things.
more than we can handle..
You know, sometimes spending time with God.
or studying the word of God with his people.
becomes optional..
(speaking in foreign language).
Now, when I was younger, much younger,.
when I was a young Christian in college,.
I didn't have a habit of reading the Bible.

$^{201}$or reading, you know, spiritual books,.
you know, but, you know, I observed myself..
I started to spend more time with questions.
who take the word of God seriously,.
and I began to change..
I began to, you know, cultivate the desire,.
or the desire to, you know, really taste.
and dig into the word of God.
and taste the goodness of knowing God himself..
(speaking in foreign language).
And let me ask you, what are you devoted to?.
Is there a hunger in your heart.
for prayer, for worship, for the word of God,.
for fellowship, for seeing people to come to Jesus?.
Now, if not, then we must change..
We must realize that we are not the only people.
who are in the world..
We are not the only people who are in the world..
We must change, we must repent,.
and we need grace, the grace of God to have a fresh start..
And because, you know, being a contagious church,.
being contagious Christians,.
we need to begin with a heart changed by the love of God..
(speaking in foreign language).
Now, that's the first point, be devoted..
The second point is be generous..
(speaking in foreign language).
44, 45, all the believers who are together.
and have everything in common,.
and they sell their possessions and goods.
and gave to anyone as he had need..
(speaking in foreign language).
Now, again, our context is different..
There's no copycat here..
We cannot do exactly what they did..
But that is one of the most radical part.
of their commitment, is that people in need,.
when they saw it, they didn't say,.
let me just pray about it, or let me just wait for it..
(speaking in foreign language).

$^{241}$Don't get me wrong, prayers are good,.
and all this caring are good,.
but you see that all these believers.
in the first church of Jerusalem,.
they act together immediately,.
and they are generous, and they give..
Now, this was not communism, okay?.
This was not communism, because it was not forced..
It was not state-sponsored..
It was moved by the spirit of God,.
moved by the love of God,.
because their lives are being transformed by Jesus,.
and they gave up everything,.
and be generous to help and to give to the church..
(speaking in foreign language).
Now, if you know something about the Old Testament,.
you immediately would think of Deuteronomy 15,.
because there was a phrase.
that describes the first church of Jerusalem.
happens to be in Deuteronomy 15..
Let me read that..
There should be no poor among you,.
for in the land the Lord your God is giving you to possess,.
he will richly bless you..
If only you fully obey the Lord your God.
and are careful to follow all these commandments,.
do not be hard-hearted or tight-fisted..
Give generously without a grudging heart..
Then, because of this,.
the Lord your God will bless you in all your work..
(speaking in foreign language).
Now, you see, in the Old Testament,.
God told his people Israel to be generous,.
and not only because there were needy people among them,.
mainly because that God will bless them..
And now, in Acts 2, the Old Testament vision,.
this vision is becoming a reality,.
not through laws, but through love..
(speaking in foreign language).
You see, this is the logic of the kingdom of God..

$^{281}$The way the world works is that when God bless me so much.
and I have extra, then I'll give a little..
That is the way of the world..
(speaking in foreign language).
See, God's way does not follow human logic..
God said, "You be generous first..
"You give in faith, then God will bless you richly..
"And you practice generosity first,.
"and then you will experience the blessing of God.".
(speaking in foreign language).
But as you know, our world is full of takers, not givers..
You know, if there's food, I'll take as much as possible..
If there are free things,.
I'll line up and get them all, right?.
You know, I drive by a food pantry.
on Ashland Avenue many times,.
and I see people line up for free things,.
and many of them are Chinese..
And I do think there are needy people in the city..
I do think we need to contribute..
But, you know, I have observed many people line up,.
you know, they really don't need free food..
They really don't need the free items,.
but they line up anyway..
They take everything home anyway,.
and then, you know, there's so much food,.
and, you know, the food, the extra food got rotten.
and wasted..
(speaking in foreign language).
, and at times, you know, if you observe,.
even the church is full of takers..
You know, the first thing people come to church is to take,.
is to get, to get something..
And have you ever thought of, you know,.
the food in the cafe and the ministry?.
We need financial support..
(speaking in foreign language).
And a couple weeks ago, we had our annual core team retreat,.
and one of the most important items we discussed.
and pray about is our church finance..

$^{321}$It's what we need to sustain and to operate our church..
And as of the end of June this year,.
we lack almost 20,000,.
and we need more than 38,000 a month to operate,.
to pay the rent, and to, you know,.
pay the staff and support the ministries..
And I think if you are practicing,.
tithes and offering regularly,.
and I wanna thank God for you,.
but if you are not,.
I think the Bible is challenging us to give regularly,.
and that is one of the things we take part.
in building the church together..
(speaking in foreign language).
See, the word of God tells you.
that if you really wanna experience the love of God,.
the blessing of God, you must act first..
You must act in faith, and not just our leftover,.
but our love, our response to the grace of God..
(speaking in foreign language).
Now, remember what Jesus said in Luke 16..
He says, "No servant can serve two masters..
"Either he will hate the one and love the other,.
"or he will be devoted to one and despise the other..
"You cannot serve both God and money.".
(speaking in foreign language).
For years, you know, walking with Christ,.
you know, for decades, you know,.
I experienced that if my heart for God is conditional,.
there's no way I can experience.
the unconditional grace of God..
And I wanna remind you, you know,.
a contagious church is one that gives freely..
It's one that we receive the great grace of God,.
and then we offer our best..
(speaking in foreign language).
Now, the third point here is be alive in worship..
And it says, "Every day, they gathered together.
"in the temple courts, and they broke bread in their homes,.
"and ate together, and they were glad and with sincere hearts,.

$^{361}$"praising God and enjoying the favor of the people..
"And the Lord added to the number daily.
"who were being saved.".
(speaking in foreign language).
And you see this community was alive,.
and their hearts were joyful,.
and they worshiped together with, you know, will, attitude,.
and they didn't worship God out of duty..
They did it with joy..
(speaking in foreign language).
And the Lord said, "You know,.
"if you're going to worship God,.
"you're going to worship Him with joy.".
Duty, they did it with joy..
(speaking in foreign language).
Now, God does not look at your appearance..
God is not so interested in your, on the outside,.
how you look on the outside..
God sees your heart deep inside..
You know, when you come to worship,.
you come to pray and reflect,.
and we should all do this,.
and pray that and prepare us,.
help me, God, to love you with all my heart..
(speaking in foreign language).
Now, as I reminded you last week,.
we Christians, we the church, are always visible in public..
(speaking in foreign language).
Now, you imagine if we are all takers, you know,.
new visitors, you know, people come to our community,.
they will feel it..
(speaking in foreign language).
And you see, if we are all givers,.
then God will sense, people will sense the life in us,.
the life of God in us..
And people will see, you know, how they eat together,.
how they spend time with each other,.
how they worship together, and they gave generously..
And people will say, "You know what?.
"I want to be part of that..

$^{401}$"I want that too.".
And, you know, just a reminder,.
you know, last week I told you, all Christians, all of us,.
we are living under the spotlight..
(speaking in foreign language).
We are living under the spotlight..
(speaking in foreign language).
So my conclusion will come to this question..
Does your life make the gospel attractive?.
(speaking in foreign language).
Now, when someone sees you,.
would that person want to follow Jesus.
after watching how you live?.
(speaking in foreign language).
And a contagious Christian life is worshipful, joyful,.
and, you know, it's a sharing life..
(speaking in foreign language).
And that is the life of church this Holy Spirit creates..
Well, it doesn't mean we are all perfect..
We all have our flaws and we all have our shortcomings..
But, you know, that's what the Spirit is trying to create,.
a new church, a new life..
And it's not a private life and that's not personal life,.
but a visible life..
It's not a strange life, but a generous life..
(speaking in foreign language).
See, you come to God today.
and God wants to do something for you..
You know, God wants to transform your life..
God wants your life and your heart and your home.
to be contagious with His love and grace and generosity..
(speaking in foreign language).
Just like people were drawn and being attracted.
to the community of Acts 2..
(speaking in foreign language).
And that's why we are here today..
(speaking in foreign language).
And that's why we are here today..
(speaking in foreign language).
So there are three important steps or simple steps.

$^{441}$I want to suggest to you today..
(speaking in foreign language).
First, starting today, read the Bible every day.
and connect with other Christians..
(speaking in foreign language).
And practice generosity..
(speaking in foreign language).
And we all need to give as much generously.
to support the ministry..
And third is to worship God wholeheartedly..
(speaking in foreign language).
Now, remember last week, what happened?.
You know, each of you get a name, right?.
(speaking in foreign language).
Did you learn that name?.
(speaking in foreign language).
Did you pray for that person?.
(speaking in foreign language).
You did, I see nodding heads..
That's great, that's great..
(speaking in foreign language).
If you've forgotten, if you've forgotten about it,.
(speaking in foreign language).
that's okay..
(speaking in foreign language).
You have one more week to go..
(speaking in foreign language).
So this coming week, remember that person, okay?.
(speaking in foreign language).
And today, I know there are new faces among us..
(speaking in foreign language).
So I want you, after the service,.
to go to someone who's new to us..
(speaking in foreign language).
And learn that person's name..
(speaking in foreign language).
And then pray for that person..
(speaking in foreign language).
All right, so remember these three things..
(speaking in foreign language).

$^{481}$Okay, read the Bible every day and connect with others..
(speaking in foreign language).
And then practice generosity..
(speaking in foreign language).
And worship God wholeheartedly..
(speaking in foreign language).
And after the service, I want you to do a little more, okay?.
(speaking in foreign language).
Remember you had the name from last week..
(speaking in foreign language).
Okay, after the service..
(speaking in foreign language).
Go to that person..
(speaking in foreign language).
And this is what you're gonna say..
(speaking in foreign language).
Thank you for worshiping with us today..
(speaking in foreign language).
Okay, can you do that?.
(speaking in foreign language).
It's pretty easy, right?.
Secondly, okay, thank you for joining us in worship today..
And secondly, you're gonna say,.
how can I pray for you this week?.
(speaking in foreign language).
And third thing is I wanna worship you next week..
Worship with you next week..
(speaking in foreign language).
Okay, first thing is what?.
(speaking in foreign language).
Go to that person and say,.
thank you for worshiping with us today..
(speaking in foreign language).
Secondly is how can I pray for you this week?.
(speaking in foreign language).
Now don't talk too long, okay?.
(speaking in foreign language).
Okay, it's very, very concise and tell them,.
you know, okay, this is how I can pray for you this week..
(speaking in foreign language).

$^{521}$And thirdly is I wanna see you next week..
I wanna worship God with you next Sunday..
(speaking in foreign language).
Shall we do that?.
(speaking in foreign language).
And remember, we have to build the church together..
(speaking in foreign language).
Every single one of us..
(speaking in foreign language).
To build this family together for the glory of God..
(speaking in foreign language).
Well, let's pray..
(speaking in foreign language).
Father, we wanna give you thanks..
(speaking in foreign language).
Because you are the God of our church..
(speaking in foreign language).
Thank you for your presence..
(speaking in foreign language).
Thank you for loving us..
(speaking in foreign language).
And I wanna pray for each and every one of us today..
(speaking in foreign language).
May we act in faith to worship you, to follow you,.
then you will bless us..
(speaking in foreign language).
And I pray that you will continue.
to build our church together..
(speaking in foreign language).
May your love transform us so that we can work for you.
and we can serve you and worship you..
(speaking in foreign language).
We pay all this in Jesus' name..
(speaking in foreign language).
Well, thank you for being with us today..
(speaking in foreign language).
Remember the three things I wanna do to your neighbor..
(speaking in foreign language).
Okay, remember that, go to that person,.
thank you for being with us..

$^{561}$And secondly, how can I pray for you?.
Third is I wanna worship God with you next week..
(speaking in foreign language).
(gentle music).
And God bless you..
(speaking in foreign language).
(gentle music).
(upbeat music).
\newpage



\section{使徒行傳 2:42-47}
\label{sec:J7Rxq8y7uj4}
\textbf{Cross Over - Be Committed ＆ Connected (Acts 2\_42-47) - Rev. Mark Chan 跨越 - 委身與連結  (使徒行傳2\_42-47) - 陳崇基牧師}
\newline
\newline
連結: \href{https://youtube.com/watch?v=J7Rxq8y7uj4}{\texttt{https://youtube.com/watch?v=J7Rxq8y7uj4}} ~~~~ 語音日期: 2025-08-24
\newline
\newline
\hyperref[sec:HQs8juXt_Kg]{\small{< < < PREV SERMON < < <}}
~
\hyperref[sec:index_chronic]{\small{[返順時目]}}
~
\hyperref[sec:index_scriptual]{\small{[返順卷目]}}
~
\hyperref[sec:TtUFBp1ccf4]{\small{> > > NEXT SERMON > > >}}
\newline
\newline
使徒行傳 2:42-47
\newline
\begin{longtable}{cl}
\hline
\hline
章節 & 經文 (和合本修訂版)\\
\hline
2:42 & \begin{tabularx}{0.7\textwidth}{X} 他們都專注於使徒的教導和彼此的團契，擘餅和祈禱。 \end{tabularx} \\ \\ \relax
2:43 & \begin{tabularx}{0.7\textwidth}{X} 眾人都心存敬畏；使徒們又行了許多奇事神蹟。 \end{tabularx} \\ \\ \relax
2:44 & \begin{tabularx}{0.7\textwidth}{X} 信的人都聚在一處，凡物公用， \end{tabularx} \\ \\ \relax
2:45 & \begin{tabularx}{0.7\textwidth}{X} 又賣了田產和家業，照每一個人所需要的分給他們。 \end{tabularx} \\ \\ \relax
2:46 & \begin{tabularx}{0.7\textwidth}{X} 他們天天同心合意恆切地在聖殿裡敬拜，且在家中擘餅，存著歡喜坦誠的心用飯， \end{tabularx} \\ \\ \relax
2:47 & \begin{tabularx}{0.7\textwidth}{X} 讚美神，得全體百姓的喜愛。主將得救的人天天加給他們。 \end{tabularx} \\ \\
[1ex]
\hline
\hline
\end{longtable}
$^{1}$(speaking in foreign language).
- Good morning, church..
(speaking in foreign language).
- How do you like the new setting and new arrangement?.
(speaking in foreign language).
- Well, let's give a hand to the worship team.
and thank you for leading us..
And especially the backstage people,.
those people who are hiding back there,.
they're doing a lot of work, okay?.
(speaking in foreign language).
Okay, I can speak a little bit, man,.
not for too long, okay?.
And then especially (speaking in foreign language).
Okay, okay..
So we hope and pray that this is a new beginning.
for our church, especially this service..
(speaking in foreign language).
And as you know, it's a beginning for our church.
and today we are trying to have this bilingual service..
(speaking in foreign language).
And we announced it last Sunday..
(speaking in foreign language).
And I wanna again thank all the staff.
and all the people who serve..
(speaking in foreign language).
As you know, the church is never a one man band..
(speaking in foreign language).
And it's a family..
(speaking in foreign language).
It takes every single one of us to work together,.
to grow together..
(speaking in foreign language).
As you know, nine years ago,.
when we launched our church in this small warehouse..
(speaking in foreign language).
And how many people we had?.
We had about, what, 30 to 40 people?.
(speaking in foreign language).
And today, by the grace of God,.

$^{41}$and we have grown into a 100 plus adult community..
(speaking in foreign language).
And this year, we wanna challenge you.
with this theme, crossover..
(speaking in foreign language).
As you can see, break our barriers and our boundaries.
and join together to build up the church..
(speaking in foreign language).
And to reach out to new people..
(speaking in foreign language).
So that we and each and every one of us.
will be touched by the love of God..
(speaking in foreign language).
Now, as we begin, I wanna ask one question..
(speaking in foreign language).
What kind of church God want us to become?.
(speaking in foreign language).
Now, same in this question,.
what kind of church God want us to become?.
(speaking in foreign language).
You see, we have been studying the book of Acts..
(speaking in foreign language).
And it's a book filled with wonderful stories..
(speaking in foreign language).
It's a book filled with inspiring pictures.
of the first church, the first people of God..
(speaking in foreign language).
It's kind of like you take your iPhone.
and you go back, you time travel,.
go back to the first century.
and try to take selfie of the church in Jerusalem..
(speaking in foreign language).
You take some selfies of the life of the church back then..
(speaking in foreign language).
Now, those picture will tell you they are.
the ideal community..
And Acts 2 is a good example..
Let me read this passage..
(speaking in foreign language).
Let me read that verse by verse and will be translated..

$^{81}$Verse 42, "They devoted themselves to church in Jerusalem,.
"the first church of God, to the apostles' teaching.
"and to the fellowship and breaking of bread and to prayer.".
(speaking in foreign language).
"Everyone was filled with awe and many wonders.
"and miraculous signs were done by the apostles.".
(speaking in foreign language).
"And look, all the believers were together.
"and had everything in common.".
(speaking in foreign language).
"They sell their possessions and goods.
"and they gave to everyone that had need.".
(speaking in foreign language).
"And every day they continued to meet together.
"in temple courts and they broke bread in their homes.
"and ate together with glad and sincere hearts.".
(speaking in foreign language).
"They praise God and enjoy the favor of all the people.
"and the Lord add to the numbers daily.
"those who are being saved.".
(speaking in foreign language).
"Now I wanna focus on this verse 42.".
(speaking in foreign language).
"It says they devoted themselves.
"to the apostles' teaching and to fellowship.".
(speaking in foreign language).
"And you can imagine it gives you a picture.
"of Christians working hard and persistently together.".
(speaking in foreign language).
"And see they're dedicated, they're devoted.
"to the teaching of the Bible.".
(speaking in foreign language).
"And they overcome challenges and issues.
"to maintain the fellowship family together.".
(speaking in foreign language).
"Now this is called commitment.".
(speaking in foreign language).
"And what does it take to build a church?".
(speaking in foreign language).
"It is commitment.".

$^{121}$(speaking in foreign language).
"Commitment is not based on your feeling, your emotions,.
"it's based on values.".
(speaking in foreign language).
"It's based on what you keep doing something,.
"even when it doesn't feel excited anymore.".
(speaking in foreign language).
"And today many people see church as gym membership.".
(speaking in foreign language).
"Now how many of you have gym membership?".
(speaking in foreign language).
"Yeah, okay, very good, very good, yeah..
"We wanna encourage that, but to many people,.
"gym membership is you sign up and you go.
"when you feel like it.".
(speaking in foreign language).
"And then you don't go, you disappear for weeks,.
"and then probably you come back when you feel like it,.
"and then hopefully you get some quick results.".
(speaking in foreign language).
"And nobody sees you, nobody holds you accountable.".
(speaking in foreign language).
"The commitment is not like that.".
(speaking in foreign language).
"Commitment is having a training partner.".
(speaking in foreign language).
"It's not just membership, having a partner in training.
"is that you show up even when you are tired.".
(speaking in foreign language).
"Because you know somebody's waiting for you.".
(speaking in foreign language).
"And you and the partner will push.
"and challenge each other.".
(speaking in foreign language).
"And you don't easily ghost each other.
"when life gets busy.".
(speaking in foreign language).
"The church is not just membership.".
(speaking in foreign language).
"It's the people you work with.".

$^{161}$(speaking in foreign language).
"You work with them and you continue to build the church.".
(speaking in foreign language).
"Now I understand many of us are immigrants or foreigners.".
(speaking in foreign language).
"Now let me tell you a story..
"My family, my family,.
"we came to the United States in 1990.".
That's how many years ago..
(speaking in foreign language).
"So we came to the United States as immigrants, foreigners,.
"and we came in 1990 and life was really hard.
"and we have to work really, really hard.".
(speaking in foreign language).
"We work all day and we came home.
"and then what we do, we close the door.".
(speaking in foreign language).
"And then the next morning we're up again.
"and we work hard again.".
(speaking in foreign language).
"Because life was hard and busy.".
(speaking in foreign language).
"My goodness, we work around the clock.".
(speaking in foreign language).
"We save every penny, we are scavengers,.
"we try to take everything possible for our benefit.".
(speaking in foreign language).
"Now this mode of living is like, it's a mode of survival.".
(speaking in foreign language).
"Now it's very understandable.
"and there's underlying feeling of fear.".
(speaking in foreign language).
"You see, we foreigners, we immigrants,.
"at least we, my family, we work hard and we stay quiet.
"and we want to protect ourselves.".
(speaking in foreign language).
"We don't trust other people.".
(speaking in foreign language).
"We don't share with other people.".
(speaking in foreign language).

$^{201}$"Because we want to survive because of fear,.
"fear of losing our money, our asset, our security,.
"our place of stay.".
(speaking in foreign language).
"That is our natural mode of survival.".
But for the first church in Jerusalem,.
for the first Christians in Jerusalem,.
life was hard as well..
(speaking in foreign language).
They face challenges and even persecutions..
(speaking in foreign language).
But surprisingly, they continue to grow.
and they continue to fly and expand..
(speaking in foreign language).
And now we have to ask why..
(speaking in foreign language).
You know, because they were committed..
(speaking in foreign language).
Commitment says, "I am not just interested, I'm invested.".
(speaking in foreign language).
"I'm not here just when life is easy,.
"I'm here because it matters.".
(speaking in foreign language).
See, that's what happens when the Holy Spirit takes over..
(speaking in foreign language).
And then people, their lives will cross over.
from fear to faith, from surviving to serving..
(speaking in foreign language).
No, the entire Bible,.
God did not come to save individual people only..
(speaking in foreign language).
The Holy Spirit came and formed a new people..
(speaking in foreign language).
A new kind of community..
(speaking in foreign language).
So here's my challenges for you today..
(speaking in foreign language).
Be committed..
(speaking in foreign language).
The Bible says the church is a family we belong to..

$^{241}$(speaking in foreign language).
And I've served in different churches for years..
(speaking in foreign language).
And I've learned one thing,.
my relationship with others in the church is reciprocal..
(speaking in foreign language).
And the more I love the church,.
the more I serve the church,.
the church, interestingly, will love you back..
(speaking in foreign language).
It's reciprocal because we need to grow in the community..
(speaking in foreign language).
So today my challenge for you is to create a new schedule..
(speaking in foreign language).
And the new schedule says commitment..
It says Sunday worship is not negotiable..
(speaking in foreign language).
Because you do talk about priority..
My priority is God..
(speaking in foreign language).
Now, you understand Christians are always under spotlight..
We are always watched by the world..
(speaking in foreign language).
Whether you like it or not,.
I give you an example of very vivid illustration..
Just last week, we are all, our social media,.
is all flooded by a video of a couple,.
a cop having an affair in a Coldplay concert, right?.
(speaking in foreign language).
And sometimes, you know, stuff like this,.
you know, I think the most interesting thing.
is the comments..
(speaking in foreign language).
Now, one comment says, if they, the couple,.
had not tried to escape the camera.
or, you know, tried to get away from the camera,.
maybe everything would turn out fine, right?.
(speaking in foreign language).
And another reply was very surprisingly and entertaining..
It says, "You're already here,.

$^{281}$"and what is there to be afraid of?".
(speaking in foreign language).
In any case, I mean, you have to understand.
the church question, we questions,.
we are always in the public eyes..
(speaking in foreign language).
And we need to rely on the grace of God.
to set a good example..
(speaking in foreign language).
And we questions, we, the church,.
need to be camera-ready all the time,.
and we need to be always ready,.
and that calls for commitment..
(speaking in foreign language).
And commitment says, "I'm interested, I'm ready, I'm live.".
(speaking in foreign language).
And see, commitment says, "I'm going to worship God.
"and put God as my number one.
"with other believers of Jesus in the community.".
(speaking in foreign language).
And the second point, on top of being committed,.
is be connected..
(speaking in foreign language).
Start to connect with other people in the church family..
(speaking in foreign language).
Now, you have to do your job,.
you have to take the initiative to get to know each other,.
and then pray for one another..
(speaking in foreign language).
Now, on your way in, I think a Jinfan.
already gave you a slip of paper, right?.
(speaking in foreign language).
Have you got it?.
Okay, all right..
So you write your name in the slip of paper, right?.
(speaking in foreign language).
Did you?.
Okay, if not, we're gonna pass it to you right now..
(speaking in foreign language).
All right, how many of you got it?.

$^{321}$Okay, this is something we wanna do together, okay?.
(speaking in foreign language).
Now, you have your name on the slip of paper..
(speaking in foreign language).
All right, and then we're going to mix it up, okay?.
(speaking in foreign language).
All right, so Jinfan, please,.
you're gonna distribute it to all the people,.
and then when you get this slip of paper,.
write down your name..
Don't scribble it, okay?.
Write down, you know, capital letters,.
or write clearly your name, okay?.
(speaking in foreign language).
And then we're gonna mix it up..
(speaking in foreign language).
And then when you pick up a slip of paper.
with someone else's name on it,.
(speaking in foreign language).
and then you're gonna learn that name..
(speaking in foreign language).
Okay, and then, in this coming week,.
you're gonna pray for that person every single day..
(speaking in foreign language).
Okay, should we do it now?.
(speaking in foreign language).
Let's do it now, let's do something together, all right?.
So, there are some pens back there,.
some pencil back there..
You are free to leave your seat right now,.
and go and get some pencil,.
and write down your name, and give it to Jin--.
Huh?.
- I already did it, I already did it..
- Some people already did it, okay, good..
- I did for them..
- Oh, you did for them, okay..
- Nice..
She's going to distribute some of the slips of paper.
to other people..

$^{361}$(speaking in foreign language).
Okay, so you pick up a slip of paper..
(speaking in foreign language).
Check the name..
Make sure the name is not your name, okay?.
You're not gonna pray for yourself, okay?.
(speaking in foreign language).
Make sure the name is not your name, all right?.
(speaking in foreign language).
Then it should go well, all right?.
And then, what do you do?.
(speaking in foreign language).
Learn that name, all right?.
(speaking in foreign language).
Well, learn to pronounce it..
- Yeah..
- Learn how to write it. - (speaking in foreign language).
- Okay, that is the name you wanna keep.
that person accountable in your prayers..
This week..
(speaking in foreign language).
Okay, everybody got it?.
Can you, can you raise your hand?.
Oh, okay, I got two..
- Thank you, thank you, I got it..
- Yeah, sure, yeah..
Let me check if this is my name..
- Yeah..
- All right, I got a very good name..
- Yeah..
- All right, so this is what you do after the service..
(speaking in foreign language).
You try to go to that person..
(speaking in foreign language).
Are you this guy?.
(speaking in foreign language).
All right, you're gonna ask, okay?.
(speaking in foreign language).
See, it takes initiative..
(speaking in foreign language).

$^{401}$Okay, and then say hello..
(speaking in foreign language).
Say hi..
All right, in the community, right?.
(speaking in foreign language).
All right, so you pick up this name,.
you're gonna pray for that person every single day..
(speaking in foreign language).
Now, you learn that name,.
and then maybe you can strike a conversation.
with the person..
(speaking in foreign language).
And then you can ask,.
how can I pray for you this week?.
(speaking in foreign language).
Okay?.
Now, you don't have to dig around for some,.
you know, you don't have to invade privacy, okay?.
Just pray for my work or something like that, okay?.
(speaking in foreign language).
Be polite, okay?.
We're all adults, okay?.
You know, learn that person, I mean, it's gonna be fun..
- Yeah. - All right?.
(speaking in foreign language).
Okay, so remember what you do?.
(speaking in foreign language).
I'm gonna remind you after the sermon, okay?.
- Okay..
- What I'm saying is be committed and be connected..
(speaking in foreign language).
And we can build the church together,.
we can serve God together..
(speaking in foreign language).
Shall we?.
(speaking in foreign language).
Let us pray..
(speaking in foreign language).
Thank you, Lord Jesus,.
you never held back your love for us..

$^{441}$(speaking in foreign language).
And you are committed to us to redeem us..
(speaking in foreign language).
And today we ask in our prayer.
to make us a committed people,.
faithful, devoted, and rooted in you..
(speaking in foreign language).
And truly give us the courage to serve.
even when it is inconvenient..
(speaking in foreign language).
And dear God, help us to love even when it's not returned..
(speaking in foreign language).
And Jesus teach us to make you our number one priority..
With our time, our hearts, our gifts, and our whole lives..
(speaking in foreign language).
Because we are yours..
(speaking in foreign language).
Help us to live like that..
(speaking in foreign language).
In your name we pray..
(speaking in foreign language).
Amen..
- Amen..
Okay, so one more reminder..
You got that?.
Okay, now is a good time, okay?.
I'll give you five minutes, you know,.
you can mix and match and find that person, okay?.
Find it..
If you don't know, ask about it, okay?.
(speaking in foreign language).
All right, so you can get up and go around the room.
and learn that person's name, okay?.
(speaking in foreign language).
All right..
(audience laughing).
And by the way, we have some refreshments outside.
and you--.
(speaking in foreign language).
(audience chattering).

$^{481}$(speaking in foreign language).
(gentle music).
(gentle music continues).
(upbeat music).
\newpage



\section{使徒行傳 2:38}
\label{sec:TtUFBp1ccf4}
\textbf{Only One Way to Be Saved (Acts 2\_38; 4\_12) - Rev. Mark Chan 唯一的拯救 (使徒行傳2\_38,4\_12) - 陳崇基牧師}
\newline
\newline
連結: \href{https://youtube.com/watch?v=TtUFBp1ccf4}{\texttt{https://youtube.com/watch?v=TtUFBp1ccf4}} ~~~~ 語音日期: 2025-08-24
\newline
\newline
\hyperref[sec:J7Rxq8y7uj4]{\small{< < < PREV SERMON < < <}}
~
\hyperref[sec:index_chronic]{\small{[返順時目]}}
~
\hyperref[sec:index_scriptual]{\small{[返順卷目]}}
~
\hyperref[sec:3tR0PVmebsI]{\small{> > > NEXT SERMON > > >}}
\newline
\newline
使徒行傳 2:38
\newline
\begin{longtable}{cl}
\hline
\hline
章節 & 經文 (和合本修訂版)\\
\hline
2:38 & \begin{tabularx}{0.7\textwidth}{X} 彼得對他們說：「你們各人要悔改，奉耶穌基督的名受洗，使你們的罪得赦免，就會領受所賜的聖靈。 \end{tabularx} \\ \\
[1ex]
\hline
\hline
\end{longtable}
$^{1}$Hey everyone, how's it going?.
How's it going?.
Good?.
Alright, why don't we say hi to our neighbors?.
And I want to thank Patrick for doing the announcement..
We want more people to be part of it..
And let's go around, look around and see your neighbor and say hi, okay?.
(Translated).
Okay, we'll give her a high five or something..
(Translated).
Okay..
Alright, so good to see you all..
And perhaps you have been to church for years, for many, many years,.
and perhaps you are new to the church..
And from time to time we need to make sure we understand what church is all about..
(Translated).
So last week I was in Toronto and Montreal for a few days..
And I paid a ticket to visit a church..
And this is a 200 years old church, and a historic church called the Church of Notre Dame in Montreal..
And it was a beautiful church, beautiful..
And it's just a wonderful place to be..
And I spent some time there and I was totally taken by all the stained glass and Gothic style building..
And it was a beautiful church..
(Translated).
So, you know, as you can see, if I can do it on PowerPoint,.
and I was in the church, in this church, and I spent some time there and I asked the question,.
"What is a church?".
And what is the church all about?.
And to many people, to many of us, sometimes we got messed up.
because we, sometimes we believe that the church is a place that we go to..
And that's why I want to ask the question, you know, "What is the church?".
You know, even though there's some truth to we go to church today, but what is the church?.
(Translated).
So, this is a beautiful church, right?.
(Translated).
And, but, you know, there's also a huge misunderstanding what the church is all about..
And, you know, the Bible says the church is not a place, the church is not a building..
And we have to ask, "What is the church in the Bible?".
(Translated).
In fact, the church is us, the people, the people of God..

$^{41}$(Translated).
So today, in this message, I want to emphasize three basic points in this sermon..
And the three basic points is all about how do we become the church..
(Translated).
And the first point is we are all sinners..
(Translated).
And we have been reading the book of Acts this year,.
and it's clear from the book that the church is not a group of perfect people..
(Translated).
And you remember in chapter 2 of the book of Acts, Peter preached the first sermon,.
and he pointed out that we are all sinners, and we sin against God,.
and Jesus died on the cross for our sins, and He rose from the dead for all of us..
(Translated).
And then the people listened to Peter's first sermon, and their hearts were cut,.
and they felt that they need to respond, and they asked Peter what they should do..
(Translated).
And then, you know, the people were being said that, Peter said,.
"Repent and be baptized, and every one of you, in the name of Jesus Christ, for the forgiveness of sins, of your sins.".
(Translated).
And these are very bold statements, okay? And they confront us with two statements of truth..
(Translated).
The first one is we are all sinners, and the second one is only Jesus can save us..
(Translated).
Now, you have to understand, the church, even our church, is not a museum of saints..
The church is home for real people, and we have real problems, okay? We have real brokenness..
And we're not here because we got it all together, but because we know we can..
(Translated).
Now, imagine a hospital. Imagine a hospital. Have you been to a hospital?.
That might not be a good thing, but, you know, the hospital is full of sick people,.
and they want to get well, but, you know, the hospital is not for healthy people..
The church is kind of like a hospital..
(Translated).
And we are all forgiven, and we are all healed by the grace of God,.
because of what Jesus has done for us on the cross..
And the one, the one condemned on the cross, Jesus himself, should have been you and me,.
because we are sinners..
(Translated).
And the good news is, God healed all of our sins on Jesus on the cross, and all of our sins are crucified..
(Translated).
Now, the Bible is really clear on this..

$^{81}$The Bible, God loves us and sends his son to die on the cross, so that your sins,.
every bit of your sin, will be forgiven..
(Translated).
So, that's the first truth, is that we are all sinners, we are all broken people..
And the second is, we are saved by the grace of God..
(Translated).
And this second truth, we have to note that, you know, many people think that,.
"Okay, only if we try, only if I try harder, if I do more good things than bad things, then I'll be okay.".
(Translated).
Now, the truth is, you know, sin is not a scale of balance..
It's a debt that we cannot pay, and even one sin makes us guilty before a holy, sinless, holy, infinite God..
(Translated).
Now, imagine an ocean, okay, endless ocean..
And you are in the sea, and the waves are high, and the current is very, very strong..
And you have no strength left, and there's no land in sight..
You are right in the middle of the ocean, and you cannot swim back, and you cannot save yourself..
(Translated).
Now, in this kind of situation, what you need is someone else, someone else to rescue you..
Now, sin is like that. Sin is like an ocean. It's like an endless ocean..
It pulls us away from God, from the land..
And there's no amount of effort, good deeds, or good religion, or self-improvement can bring us back to the shore..
Now, without a rescuer, we will drown to death..
(Translated).
In the book of Romans, chapter 3 here, it says,.
"This righteousness, the righteousness from God, comes through faith in Jesus Christ to all who believe.".
And there is no difference, for all of us have sinned, and we fall short of the glory of God..
(Translated).
Now, we often have a misunderstanding here. We always think that what the Bible is all about,.
and we think that the Bible is to tell us to be a better person..
In fact, the Bible never does that. The Bible never tells you to become a better person..
The Bible is not a self-help moral book..
(Translated).
Okay, say it with me. The Bible is never a self-improvement book. All right?.
The Bible is never a self-improvement book, so that we can do better,.
and we can reach out to this holy and infinite God..
(Translated).
The Bible, though, tells us that we can never be good enough,.
and we may be better than someone else..
You know, look around your neighbor, and you may be better than your neighbor in some sense..
The Bible never tells us that. Okay?.

$^{121}$Or you may say, you know, look at that person, look at this person..
You know, he's worse than me, and I'm glad I'm better..
(Translated).
The fact is, the Bible tells us we are all sinners. We fall short of the glory of God..
And we don't stand righteous before God..
And we stand, all of us, we stand condemned in the presence of God..
(Translated).
Now, we can never be good enough to meet God's standard,.
and we cannot be good enough to stand in front of God..
Now, that is true, and that is also bad news, right?.
(Translated).
It seems like drowning in the ocean, and there's no land inside,.
and you can't reach the shore on your own..
(Translated).
So we have to ask the question, what can we do as sinners in front of God?.
What can we do when we cannot save ourselves?.
(Translated).
Now, that's all the bad news, but the good news comes from the Bible..
It says, "God came to save us.".
God came to save us. Not we can save ourselves, but God came to save us,.
and we are all saved by the grace of God alone..
(Translated).
So the third point here is that only Jesus can save us..
(Translated).
Now, in chapter 4 of the book of Acts, if you remember, or if you go back and check,.
the religious leaders in Jerusalem, they crucify Jesus on the cross,.
and they confront Peter about Peter's right to preach the gospel..
(Translated).
And Peter responds by saying this verse in chapter 4, verse 12..
It says, "Salvation is found in no one else, for there is no other name under heaven.
given to mankind by which we must be saved.".
(Translated).
Now, why only Jesus can save?.
(Translated).
Because Jesus is the only one who lives a perfect and sinless life..
(Translated).
And Jesus is the only one who died as a substitute for our sins..
(Translated).
And third, Jesus is the only one who rose from the dead and defeated death for all of us..
(Translated).

$^{161}$Think about it. You know, there's no, there's none..
There's no leader in the world, in history..
There's no political leader or no philosophy..
No self-effort can deal with our sin..
(Translated).
Only Jesus can bridge the holy God, sinless God, and us broken sinners..
Only He can be the bridge, and only Jesus has the power to forgive.
and to give us the life of God, the eternal life..
(Translated).
So we have to ask at the end of the sermon is to ask ourselves, "What must we do now?".
(Translated).
What must you do? What must we do?.
The Bible tells us something very simple, but life-changing..
(Translated).
The Bible tells us at least three things..
First is to repent..
Which is turning away from our sin, our brokenness, by trusting in God Himself..
(Translated).
So the first step is to repent. Repent from our sin and our brokenness,.
and make a U-turn back to God..
And secondly is to believe, which is trusting in Jesus Christ Himself,.
alone as our Savior and our Lord..
(Translated).
And be baptized, that's the third step, which is a decoration,.
a perfect decoration of our faith and our new life in Jesus..
(Translated).
Now, I don't know if you have been baptized or not, but you know, what is baptism?.
So often we have some misunderstanding about what baptism is all about..
We think baptism is a graduation ceremony..
(Translated).
And we think that, okay, once we're baptized, we graduate, like graduating from school,.
and we don't have to go back, and we upgrade or something..
(Translated).
Baptism, no, it's not like that. Baptism is not like graduation ceremony..
(Translated).
Baptism is a perfect decoration of our faith in Jesus..
(Translated).
Now, baptism means we say to the world that, yes, I'm a sinner, I'm broken,.
but I have experienced the grace of God, and by the grace of God, I stand righteous before God..
(Translated).

$^{201}$When you believe in Jesus, your sins are forgiven, and we are given the new life, the eternal life..
(Translated).
And sometimes the Bible, especially the book of Acts, describes the eternal life as the gift of the Holy Spirit..
(Translated).
And the gift of the Holy Spirit was that it's actually God's own presence in our life,.
and God is living with you, living in you..
(Translated).
Now, my final question is, if you have made Jesus your Savior and Lord,.
you have accepted Jesus as your Lord, and it is time for you to restart your life, restart your commitment..
(Translated).
If you have never done that, and now is the time to turn to Jesus..
(Translated).
And on my last slide, there's a very simple prayer, and I hope that we can restart our life by trusting in Jesus..
(Translated).
And it's a very simple prayer..
It says, "Lord Jesus, I know I'm a sinner, and I cannot save myself..
I believe you died for my sins and rose again for me..
And I now turn from my sin and trust you as my Savior..
Forgive me and help me follow you the rest of my life..
Would that be your prayer today?.
Would that be your prayer?.
And we can pray this prayer and trust in God, trust in Jesus..
(Translated).
If you can, you can pray with your heart..
Or you can pray out loud with me..
Let me read this prayer again..
Lord Jesus, I know I'm a sinner..
I cannot save myself..
I believe you died for my sins and rose again..
I turn from my sin and trust you as my Savior..
Please forgive me and help me follow you the rest of my life..
Amen..
Let's pray..
Father, you heard our prayer..
And we want to trust in you..
And thank you for being our Lord and our Savior..
Thank you for dying for our sin on the cross..
And you rose from the death for all of us..
And help us to follow you..
Because you are the Lord of our life..

$^{241}$In Jesus' name we pray..
Amen..
Now, last note I want to give to you is that if you did pray this prayer,.
and you have questions about faith, you have questions about the Bible,.
you know, God, Jesus, you know, please stop by and talk to one of us, one of our staff..
(Speaking Chinese).
Okay, so we have a freshman out there,.
and we want to thank the Brothers and Sisters for preparing that..
(Speaking Chinese).
Please spend some time to get to know one another..
(Speaking Chinese).
Alright, God bless you..
(Speaking Chinese).
(Music).
\newpage



\section{路加福音 22:15}
\label{sec:3tR0PVmebsI}
\textbf{彌賽亞 - 耶穌的筵席: 從排他到包容的救恩 (路加福音22\_15 -23) - 梁紫健傳道}
\newline
\newline
連結: \href{https://youtube.com/watch?v=3tR0PVmebsI}{\texttt{https://youtube.com/watch?v=3tR0PVmebsI}} ~~~~ 語音日期: 2025-08-24
\newline
\newline
\hyperref[sec:TtUFBp1ccf4]{\small{< < < PREV SERMON < < <}}
~
\hyperref[sec:index_chronic]{\small{[返順時目]}}
~
\hyperref[sec:index_scriptual]{\small{[返順卷目]}}
~
\hyperref[sec:Xoz1f5_xHTQ]{\small{> > > NEXT SERMON > > >}}
\newline
\newline
路加福音 22:15
\newline
\begin{longtable}{cl}
\hline
\hline
章節 & 經文 (和合本修訂版)\\
\hline
22:15 & \begin{tabularx}{0.7\textwidth}{X} 耶穌對他們說：「我非常渴望在受害以前和你們吃這逾越節的宴席。 \end{tabularx} \\ \\
[1ex]
\hline
\hline
\end{longtable}
$^{1}$(第一節).
弟兄姊妹,早安.
早安.
我也很享受剛才的音樂.
先讓大家看看第一張泡泡.
幫我按一下.
對,這個.
這個不是一個桶,而是一個杯.
是現在世界上最大的一隻杯.
在2023年,墨西哥贏了世界冠軍.
大家猜猜它能裝多少卡路里的水.
A是50至150卡路里.
B是800至150卡路里.
1500卡路里.
C是2000至3000卡路里.
大家都說得很對.
正確答案是C.
是2922卡路里.
贏了的朋友待會請你多喝一杯水.
這個杯可以裝很多墨西哥的飲品.
如果一個人喝的話.
可能一輩子都喝不完.
其實…可以幫我按下一個泡泡嗎.
杯在不同的飲食文化中.
扮演著不同的角色.
代表著好客,尊重,重視.
在阿拉伯的咖啡文化中.
咖啡杯是待客之道.
如果你是客人,右手拿著杯子.
喝完之後,如果你不搖杯子.
主人就會不斷倒給你.
倒到你搖杯子為止.
代表著甚麼?就是款待和慷慨.
所以你會發現.
如果你去阿拉伯人的家作客.
會給你一個很好的咖啡杯.
去喝你的咖啡.
在中國當然有客人之道,對吧.
喝茶,對吧.
你經常請人喝茶,對吧.

$^{41}$茶杯,特別是右上角.
或左上角的紫砂茶杯.
是茶道中的核心.
因為它可以吸收茶的味道.
然後每次沖的茶.
都令每次喝的茶更豐富.
這象徵著寧靜和分享.
很多人還會儲存不同的杯子.
家中可能有很多不同的杯子.
喝茶時也有很多不同的風味.
在韓國更是.
韓國人把酒杯視為社交.
很重要的一環.
很少會被別人的杯子阻擋.
基本上要倒到滿,滿到滿血.
倒完後還要這樣.
要很恭敬地給別人.
因為他們受到儒家的影響.
好客和尊重是很重要的.
所以他們要用杯子來表達.
這是我朋友特意從日本買回來的咖啡杯.
Alice第一次看到.
她說這是骨灰鑼.
其實我覺得不像.
我覺得像觀音擺竹的小杯.
其實它是甚麼.
它是一個品香杯.
你看到這裡有個握位.
它會困住香味.
你嗅起來時是特別不同的咖啡.
所以你喝下去的咖啡會覺得特別甜.
你有機會可以試試.
為何我會用這個杯子作為開頭.
其實在儒家福音的晚餐裡.
杯子同樣扮演著一個很重要的角色.
我們稍後會看看.
究竟耶穌如何用這個杯子.
把整個聖餐的環節推到高峰.
可以給我下一個方法嗎.
說完後我們會得著三件事.

$^{81}$聽完這篇講道.
我們會知道我們應該用甚麼心態.
去領受聖餐.
第二,就是領受聖餐.
需要甚麼資格.
第三,就是思考題.
當領受聖餐的時候.
你會想起甚麼.
如果準備好,就下一篇經文.
「時有到了」.
儒家福音22章15至23節.
「時有到了耶穌坐直」.
「使徒們也和他同坐」.
「耶穌對他們說」.
「我非常渴望在受害以先和你們」.
「吃者如子的現值」.
「我告訴你們我不再學者現值」.
「直等到他實現在上帝的國內」.
「耶穌接過杯來,捉著鳥說」.
「你們拿著杯,分著喝」.
「我告訴你們從今以後」.
「我不再喝這葡萄汁」.
「直等到上帝的國內」.
「他又拿起餅來,捉著鳥」.
「就張開了,遞給他說」.
「這是我的身體為你們捨的」.
「你們應當如此行為的是紀念我」.
「反後,他照樣拿起杯來說」.
「這杯是用我血所立的新約」.
「為你們流出來的」.
「但是看那出賣我人的手」.
「跟我一同在桌子上」.
「人子固然要按預定的離去」.
「但那出賣人子的就有禍了」.
「於是他們就互相追問」.
「他們中間哪一個會做這事」.
通常…請下一個客人.
教會在守聖餐的時候.
都會用哥林多前書的經文.
十一章.

$^{121}$很可能的原因就是有程序.
還有解釋煮的餅和煮的杯.
跟律法之間的關係.
如果有人不按規矩吃聖餐.
就等於怎樣?吃喝自己的罪.
聽起來好像很嚴重和很嚴肅.
以前有大英姐妹都會跟我說.
她說她拿起聖餐杯的時候會手震.
她會發冷和手心冒汗.
心想會否做了一些不守規矩的事.
吃喝自己的罪.
所以…下一個重點.
路加福音這裡.
耶穌似乎有一點不守規矩.
為甚麼呢.
耶穌竟然兩次拿起杯子.
這是在哥林多前書.
以及其他福音書沒有記載.
是路加福音獨有的記載.
我們就要問.
這個獨有的記載.
耶穌是表達甚麼.
如果大家留意到的話.
耶穌手持聖餐在甚麼節日裡.
是在禹節.
在一個禹節的筵席裡.
禹節對猶太人來說是甚麼.
是一個非常重要的節日.
象徵著上帝帶領以色列人出埃及.
帶領他們離開了埃及的奴役.
拯救他們出來.
所以每一年.
以色列人都很重視這個日子.
他們會去吃祝福羊.
紀念耶和華在這裡帶領他們出來.
所以猶太人的禹節傳統裡.
其實都很重視杯子.
禹節的晚餐裡有四個杯子.
四個杯子都代表著不同的東西.
請看這個PowerPoint.

$^{161}$在路加福音裡.
你看到只有兩個杯子.
這裡有四個杯子.
第一個杯子是聖傑的杯子.
象徵著神帶領以色列人出埃及.
第二是審判的杯子.
第三是救贖的杯子.
第四是讚美的杯子.
所以…請看這個PowerPoint.
在路加福音裡.
耶穌拿著第一個杯子.
祝謝了.
說「你們拿著這個杯子」.
「分開喝吧」.
「從今以後我就不再喝了」.
這個杯子很可能是對應著.
禹節的第二和第三個杯子.
強調神國的盼望和神國的來臨.
這個行動是在提醒耶穌.
祂在用祂的生命.
祂明知自己受害.
仍然堅持要吃這個筵席.
耶穌是這樣說的.
「我非常渴望在受害之先」.
「和你們吃這餐禹節的晚餐」.
耶穌用祂的生命成就上帝的國.
而第二個杯子是飯後的杯子.
耶穌說.
「這個杯子是我用我的血」.
「所納的新約為你們流出來」.
「這個杯子是整個聖餐設立的核心」.
「對應的是禹節的第三個杯子」.
「象徵著耶穌的犧牲和新約的成立」.
這樣看來好像很簡單.
其實就是耶穌為我們死.
為我們復活.
我們經常在教會說.
我們在傳福音報道時.
也會說耶穌為我們死,為我們流血.
耶穌的血洗淨我們.

$^{201}$但這個杯子不是那麼簡單.
耶穌明顯要把這個杯子.
和出賣祂的人放在一起.
留意,耶穌在第20節是這樣說的.
「這個杯子是用我血所納的新約」.
「為你們流出」.
21節.
「但是那出賣我的人的手和我一起」.
這是《路加福音》獨立的記載.
即是說甚麼.
是說耶穌專門在禹節晚餐中.
在傳說中彌賽亞的筵席中.
用一種被出賣,被犧牲.
作為救恩的核心.
亦是為何要納新約.
因為新約和以往的盟約很不同.
以往是甚麼呢.
以往彌賽亞的筵席.
只可以讓二人參加.
以往的彌賽亞筵席就是.
那些跟足了猶太律法的人.
才可以成為二人.
才有資格成為彌賽亞筵席.
跟耶穌彌賽亞一起吃飯.
我們說說以色列人.
對於彌賽亞筵席的期待.
總括而言,你看到有很多字.
很多人…很多,我也不讀太多.
其實是二人才有資格.
在《以色列四書》中.
當時猶太人正在經歷被擄.
沒有聖殿堂和耶路撒冷.
他們很渴望傳說中的彌賽亞來臨.
但同一時間還未等到.
那怎麼辦?就創作一些東西.
他們創作了一些像舊約的識經書.
當彌賽亞回來時.
我們就可以跟她一起吃飯.
我們可以分享她的權力和榮耀.
類似是這樣.

$^{241}$所以你看到那些書有點像舊約.
例如《以色列四書》,《巴洛伊書》.
《以洛伊書》.
都是在模仿舊約.
那些是我們稱為法國雜誌.
它是聖經版本的法國雜誌.
它說的是甚麼?.
它說的是在以色列人準備的筵席中.
它將力威亞探和比希莫斯.
即是這兩個海怪.
成為了二人的點心.
即是吃掉了它.
在《以洛伊書》中說.
這個筵席只可以讓以色列二人參加.
所以我們會知道.
在兩天之間猶太傳統中.
彌賽亞筵席從來只對二人開放.
從來.
所以甚麼是二人.
就是那些在苦難中守法.
堅持到最後的人.
才可以跟彌賽亞一起吃飯.
分享權力.
可以給我下一個力量嗎.
就像這個加冕大會.
大家都看過這部.
應該是二十多年前的電影.
「濟公」.
周星馳飾演濟公,對吧.
還記得那個故事嗎.
就是他經歷了人世的苦難.
其實他一開始就跟神仙爭論.
「我不相信外人不會悔改」.
「我希望外人會悔改」.
他去了世間跟外人一起經歷苦難.
經歷了苦難後,外人竟然悔改.
最後在天堂的加冕大會.
獎勵他升級成為一名降龍專者.
這就是猶太人所想的事情.
彌賽亞是要來獎勵他們.

$^{281}$獎勵這群堅持到最後的二人.
是這群二人才有資格.
吃這個加冕大餐.
最近跟一群弟兄姊妹吃飯.
談起Joyce傳道的課堂.
Anger Management.
他們談起這個課堂很有得著.
提醒他們如何面對自己的憤怒.
如何拆解自己因為一些憤怒.
而作出令人後悔的決定和行為.
以致夫妻之間可以相處得更好.
彼此體諒更多.
有件事很有趣.
原來大家上課時有個期望.
就是期望學習對付旁邊的人.
而不是自己.
這個憤怒的背後.
其實可能帶著一些東西.
就是我們很希望旁邊的人.
可以管理好我們的脾氣.
管理好我們的憤怒,對吧.
大家也帶著不少憤怒吧.
大家都會生氣吧.
我生氣他不懂得做那件事.
我生氣他自以為是.
我生氣他經常不聽我的話.
其實這些生氣.
背後都代表著一些願望或期望.
例如Alice經常生我的氣.
她生我的氣,我不聽她的話.
其實背後可能是說.
她很想我多點留意她.
有些願望,多點留意她.
甚至我重視她.
將她放在一個較早,較先的優先位置.
這就是憤怒背後的期望或願望.
所以憤怒背後帶著這些期望.
和願望其實很普遍,對吧.
正常人也會有.
猶太人就是這樣.

$^{321}$在聖殿被毀後.
你可以想像到他們多憤怒.
他們家園盡毀.
他們期待著彌賽亞利林加冕他們.
但他們一直受苦,被人殺,打壓.
所以猶太人期待.
彌賽亞利林直至屬於他們.
是多麼正常的一件事.
只有義人才可以跟他們吃飯.
我還問了一個問題.
就是其他人憑甚麼.
憑甚麼分享彌賽亞的賢值.
有點像是你上班工作了八小時.
最後同事在最後一分鐘才打卡.
但你跟同事收同樣的工資.
你是否很生氣.
有一點這樣的事.
所以…請下一個泡沫.
耶穌竟然在這個加冕大會裡.
把這個象徵著救恩的杯.
遞給了罪人.
把這個杯遞給了罪人.
耶穌透過他的生命.
重新定義甚麼是彌賽亞賢值.
就是挑戰這種排他的觀念.
在猶太人的期望裡.
唯有受到律法最多.
在苦難裡最堅持的信徒.
才有機會吃彌賽亞賢值.
但耶穌做了一件完全反轉的事.
他明知道門徒會出賣他.
他也照樣吃這頓飯.
他說「看吧,那出賣我的人的手」.
「已經跟我同在桌子上」.
「還有之後再給三次不認主」.
這件事是他自己做的.
他自己做的,他自己做的.
出賣耶穌這件事之後.
然後就沒有人理會,沒有人在乎.
他們可能在乎是誰造樣.

$^{361}$因為耶穌說過.
出賣他的人有禍而已.
然後,如果你看之後.
《路加福音》下面的經文.
你就會看見.
他們在做甚麼.
他們是爭奪偉大.
他們是爭奪誰做大佬.
他們只是希望分享彌賽亞賢直的權力.
所以這頓飯,如果你是耶穌.
你捨不捨得下.
每個人都過懷鬼胎.
只想利用耶穌達到自己的目的.
這個彌賽亞的賢直.
其實是眾多賢直當中的高潮.
耶穌一直用他的生命轉化.
跟不同人關於彌賽亞賢直的概念.
可以下一節再說嗎.
謝謝.
在《路加福音》中,你會看見.
其實有四個在彌賽亞.
最後晚餐之前的賢直.
《路加福音》七章.
在法利賽人的故事中.
一個罪人在法利賽人的飯局出現.
用香膏抹耶穌的腳.
法利賽人西門覺得有點不妥.
為何耶穌會接受這個罪人的服侍.
耶穌說,這個女人因為愛多.
所以她的罪得赦免.
這表明彌賽亞賢直.
不只是以人的尊利.
罪人也是因為信心和愛.
可以成為賢直的一部分.
第二,《路加福音》的11章.
亦是法利賽人的飯局.
法利賽人跟耶穌爭論.
飯前是否要接受這個規矩.
法利賽人很注重外表的潔淨.
但耶穌說真正潔淨的是甚麼.

$^{401}$是心裡的潔淨.
所以不是外在的規條.
而是內心的轉化.
第三,大賢直的比喻.
主人請了很多人.
本來受邀的客人都拒絕不來.
主人說請那些窮困,貧窮,乞訝的人.
來參與這次賽事.
但主人卻不願意.
因為他們的心裡的潔淨.
是內心的轉化.
所以主人就不願意.
而是不願意讓他們參與.
所以主人就不願意.
而是不願意讓他們參與.
所以主人就不願意.
而是不願意讓他們參與.
所以主人就不願意.
而是不願意讓他們參與.
所以主人就不願意.
而是不願意讓他們參與.
所以主人就不願意.
而是不願意讓他們參與.
所以耶穌說人子來以為要尋找.
拯救失喪的人.
這個賢直再一次證明.
來塞瓦賢直不只是為二人.
而是為罪人而設.
所以…請下一個包袱,謝謝.
耶穌一直用心良苦.
跟不同的人去吃飯.
希望改變大家對於.
來塞瓦賢直的盼望.
來到路加福音22章.
請下一個包袱,謝謝.
耶穌和門徒在雨節的賢直上.
這個賢直正是高潮.
耶穌說「我非常渴望在我受害之先」.
「跟你們一起吃雨節的晚餐」.
祂明知道有人出賣祂.

$^{441}$祂仍然會選擇跟這群人.
這群出賣祂的人分享賢直.
這個行動其實將整個猶太傳統.
只限於二人的來塞瓦賢直.
轉化成一個罪人的來塞瓦賢直.
所以今天你和我可以參加聖餐.
可以喝這個杯.
是否因為我們是二人.
比其他人做得更好,更公義.
所以有聖餐吃,是嗎.
出賣耶穌的人是否二人.
很明顯不是.
今天我們能夠享用這個聖餐.
是因為耶穌故意將救恩的核心.
將杯子跟罪人分享.
耶穌明知道彼得三次不認祂.
耶穌明知道門徒中間會出賣祂.
耶穌明知道他們之後會爭奪誰做大佬.
因為來塞瓦賢直代表著.
他分享了很多權力出去.
門徒受到猶太文化的影響.
他覺得和來塞瓦賢直吃飯.
就很威風,很有權力.
覺得自己是二人.
但耶穌教導他們.
甚麼才是最大的人.
就是服侍人的位置,是嗎.
但他仍然都應許和他們一同坐直.
這個「應許」顯示了.
整個來塞瓦賢直,整個聖餐.
並不是為二人而設.
而是為一些承認自己是罪人的人而設.
連真正的門徒都被包容.
耶穌一直等待他們的迴轉.
所以來應用.
今天我們如何領受聖餐.
我們應該用甚麼心態領受聖餐.
《路加福音》22章提醒我們.
(聖餐禮).
來塞瓦賢直從來不是為二人.

$^{481}$所謂做得很好的人而設.
而是為失敗,軟弱.
甚至出賣耶穌的人而設.
門徒的真權,能給予的不認主.
但耶穌仍然邀請他們參與這個禮儀.
正如我們今天領受聖餐.
絕對不是因為我們沒有瑕疵.
而是因為耶穌的恩典.
包容我們的不足.
剛剛提到憤怒這個話題.
為何我們經常會生氣別人或自己.
憤怒的背後其實隱藏著一些.
很深藏的願望或期望.
我們希望自己或身邊的人.
能做到一些事情.
但很多時候我們生氣別人時.
是因為別人達不到我們的期望.
我們覺得別人做得不夠好.
經歷了很多失望和沮喪後.
我們很容易會跌入思想的陷阱.
我們以為自己一定比別人做得好.
就像之前我們的車被撞了.
「Hit and Run」.
我也很生氣撞我的司機.
「怎麼搞的?開車不看路嗎」.
撞完還要離開.
那麼不負責任,對吧.
罵他吧.
我們期望社會上每個人.
都為自己的行為負責.
我們有這樣的期望.
所以我們有這樣的憤怒.
但當我們做錯事時.
我們是否真的那麼勇敢承認錯誤.
這是我對自己的反省.
同時我也很生氣警察.
明明有很多警察.
這樣也捉不到?.
那麼久也捉不到?.
久而久之,我便會想.

$^{521}$如果我是警察,我一定做得比他好.
不過只是我的心想而已….
是否真的做得比他好?不知道.
不知道,對吧.
盧家富陰的性餐正在提醒我們.
放下我們自己的異議.
很多時候我們以為做得比別人好.
但事實上可能不是.
以謙卑的心接受耶穌的包容.
對我們自己的包容.
亦是對身邊不同的人的包容.
當你指著別人時.
其實你有四隻手指指著自己.
對吧….
所以領受性餐需要甚麼資格.
無論我們做得如何.
在眾人面前承認我們跟耶穌的關係.
我們切切實實是罪人.
有時還有些很不實際.
或很不屬於耶穌的期望.
以及一些隱藏的議題.
希望耶穌幫我完成.
會嗎.
思考題.
當你領受性餐時想起甚麼.
你會想起一些人從未有機會.
跟耶穌有關係.
他們從未有機會跟我們一起領受性餐.
你會想起這些人嗎.
我們一起祈禱.
主耶穌,謝謝你.
為我們預備這場救恩的現場.
包容我們的軟弱和失敗.
讓我們看到本來就是罪人.
有很多不足之處.
亦有很多我們不知道的隱藏議題.
裡面有很多願望和期望.
我們很想主你幫我們實現.
求你幫助我們.
幫助我們用謙卑的心去領受性餐.

$^{561}$將你的愛傳開.
奉你的聖名祈禱.
阿們.
我們馬上領受性餐.
如果今天你跟耶穌有關係.
你是真心承認耶穌是你的主人.
你就可以領受性餐.
但如果你還有疑慮.
你還未弄清楚.
你跟耶穌之間是甚麼關係.
就請你安坐在你的座位上.
你可以在這個時間考慮.
是否要接受耶穌成為你的救主.
將來就可以跟我們一起領受性餐.
桌上有不同的性餐杯和餅.
大家可以離開座位.
開始領取餅和杯.
如果你有姿態不便的姿態.
請身邊的弟兄姊妹幫忙.
你的性餐.
如果大家拿到的話.
就可以撕開包裝.
這塊餅是主耶穌為我們寫的身體.
我們一起領受主的餅.
這杯是主耶穌為我們寫.
為我們流出的血.
紀念著祂洗淨我們的罪.
你也在這個筵席裡.
對著我們這群罪人.
不要嫌棄我們,繼續跟我們分享.
幫助我們,我們一起領受這杯.
我們一起祈禱.
天父,謝謝你.
謝謝你操縱耶穌成為我們的榜樣.
告訴我們.
我們可能有很多心態不正確.
但主在我們營造罪中的時候.
你就已經跟我們分享這個筵席.
幫助我們每一個人.
我們能夠在這裡繼續尋求.

$^{601}$繼續成長.
讓我們發現自己更多.
我們這樣祈禱,奉求耶穌得勝免求.
祝福大家.
謝謝.
謝謝收看.
\newpage



\section{傳道書 1:1-11}
\label{sec:Xoz1f5_xHTQ}
\textbf{日光之下當思考 (傳道書 1\_1-11) - 譚文倩莊傳道}
\newline
\newline
連結: \href{https://youtube.com/watch?v=Xoz1f5-xHTQ}{\texttt{https://youtube.com/watch?v=Xoz1f5-xHTQ}} ~~~~ 語音日期: 2025-08-24
\newline
\newline
\hyperref[sec:3tR0PVmebsI]{\small{< < < PREV SERMON < < <}}
~
\hyperref[sec:index_chronic]{\small{[返順時目]}}
~
\hyperref[sec:index_scriptual]{\small{[返順卷目]}}
~
\hyperref[sec:rZsT_71UjNw]{\small{> > > NEXT SERMON > > >}}
\newline
\newline
傳道書 1:1-11
\newline
\begin{longtable}{cl}
\hline
\hline
章節 & 經文 (和合本修訂版)\\
\hline
1:1 & \begin{tabularx}{0.7\textwidth}{X} 在耶路撒冷作王、大衛的兒子、傳道者的言語。 \end{tabularx} \\ \\ \relax
1:2 & \begin{tabularx}{0.7\textwidth}{X} 傳道者說：虛空的虛空，虛空的虛空，全是虛空。 \end{tabularx} \\ \\ \relax
1:3 & \begin{tabularx}{0.7\textwidth}{X} 人一切的勞碌，就是他在日光之下的勞碌，有甚麼益處呢？ \end{tabularx} \\ \\ \relax
1:4 & \begin{tabularx}{0.7\textwidth}{X} 一代過去，一代又來，地卻永遠長存。 \end{tabularx} \\ \\ \relax
1:5 & \begin{tabularx}{0.7\textwidth}{X} 太陽上升，太陽落下，急歸所出之地。 \end{tabularx} \\ \\ \relax
1:6 & \begin{tabularx}{0.7\textwidth}{X} 風往南颳，又向北轉，不停旋轉，繞回原路。 \end{tabularx} \\ \\ \relax
1:7 & \begin{tabularx}{0.7\textwidth}{X} 江河都往海裡流，海卻不滿；江河從何處流，仍歸回原處。 \end{tabularx} \\ \\ \relax
1:8 & \begin{tabularx}{0.7\textwidth}{X} 萬事令人厭倦，人不能說盡。眼看，看不飽；耳聽，聽不足。 \end{tabularx} \\ \\ \relax
1:9 & \begin{tabularx}{0.7\textwidth}{X} 已有的事，後必再有；已行的事，後必再行。日光之下並無新事。 \end{tabularx} \\ \\ \relax
1:10 & \begin{tabularx}{0.7\textwidth}{X} 有一件事人指著說：「看，這是新的！」它在我們以前的世代早已有了。 \end{tabularx} \\ \\ \relax
1:11 & \begin{tabularx}{0.7\textwidth}{X} 已過的事，無人記念；將來的事，後來的人也不記念。 \end{tabularx} \\ \\
[1ex]
\hline
\hline
\end{longtable}
$^{1}$早安,弟兄姊妹.
大家夠不夠睡?.
夠?.
這麼不公平?.
我很不夠睡.
我由去年開始祈禱.
在想今年我港島的時候.
要講哪本書.
或者要學哪本書去預備我的港島.
有一天,崇拜完了之後.
出到門口,我看到有一家人.
他們剛剛生了一個小孩.
我走過去,恭喜他們.
然後把手遞出去抱人家的小孩.
那個爸爸很興高采烈地把小孩遞給我.
然後說"Say hi to Grandma Joyce".
我就是這樣反應.
我想,什麼?.
Grandma Joyce?.
有沒有搞錯?.
我心裡在想.
其實我從什麼時候開始.
由這個阿姨跨越.
變成了Grandma.
我在想,究竟是我的衣著的問題.
還是我面上的老人斑開始越來越多.
什麼原恩惠稱呼我為Grandma?.
說起來.
其實叫我BB叫我Grandma.
一點都不過分.
因為這個爸爸很小就結婚.
這個Baby已經是他第五個小朋友.
所以他這樣叫BB叫我Grandma.
其實一點都不過分.
很搞笑.
但是被人叫Grandma.
有些令我開始懷疑人生.
當然,說笑而已.
懷疑人生真的有點誇張.
但是無可否認.

$^{41}$這個Grandma這個稱呼.
在提醒我進入人生另一個階段.
我在想,其實我自己還有多少年呢?.
我在想,我現在的工作有沒有意義呢?.
有沒有價值呢?.
我做的事情有沒有滿足感呢?.
我的婚姻和家人的關係是怎樣呢?.
我自己的健康.
不斷在思考的時候.
我就在想.
不如就看看《傳道書》吧.
這本有關生命的經世巨作.
《傳道書》是三本智慧文學聖經的書卷之一.
這本書描述了傳道者在追尋智慧.
去處理生命.
在太陽底下的生命裡的掙扎.
人生到了盡頭的時候.
我們的一生到底有什麼意義呢?.
我先跟大家讀一下這十一節經文.
《傳道書》一章一至十一節.
在耶路撒冷作王大位的兒子.
傳道者的言語.
傳道者說去控的去控.
去控的去控.
全是去控.
人一切的勞碌.
就像他在日光之下的勞碌.
有什麼益處呢?.
一代過去一代又來.
地卻永遠長存.
太陽上升太陽落下.
急歸所出之地.
風往南颳又向北轉.
不停旋轉繞回原路.
江河都往海裡流.
海卻不滿.
江河從何處流.
仍歸回原處.
萬事令人厭倦.
人不能說盡.

$^{81}$眼看看不飽.
耳聽聽不足.
已有的事後必再有.
已行的事後必再行.
日光之下並無新事.
有一件事人只著說.
看這是新的.
它在我們以前的世代早已有了.
已過的是無人紀念.
將來的事後來的人也不紀念.
首先說傳道書的背景和作者.
傳道書是擁有希伯來文化的世界觀.
在後巴比倫流亡的時期.
是深深受到神和子民之間所納的約.
的影響而寫成的.
傳道書是智慧文學.
傳道者在思索人生的時候.
去質疑人的智慧和我們人的努力.
到最終有沒有甚麼價值.
這個問題是在挑戰當時崇尚智慧文化的社會.
而智慧這個品格.
時時都會被擺在同做好事.
同敬畏神一切.
而遵守這個摩西五經.
就是人追求生命意義的指引和原則.
但在傳道書裡面.
這個傳道者偏偏去問.
我們人有智慧遵守律法又如何.
因為最終我們人都會面對痛苦和死亡.
傳道者又是誰呢.
是傳道書裡面的核心人物.
他通常會被翻譯為一個傳道者.
一個教師,一個講道者.
雖然經文裡面沒有直接說他是誰.
但一般的猶太教或基督教的傳統.
普遍都認為這個就是所羅門王.
因為傳道書一章一至.
這裡這樣說.
在耶路撒冷作王大位的兒子.
傳道者的言語.

$^{121}$不過現在很多聖經的學者都認為.
這個可能只是一種文學上的角色設定.
其實不一定代表所羅門王.
交代完背景和作者.
我們繼續一起看看.
這段經文對我們有甚麼提醒.
傳道書一章二至三節.
這裡這樣說.
傳道者說.
虛空的虛空.
全是虛空.
人一切的勞碌就是他在日光之下的勞碌.
有甚麼益處呢.
第一個虛空.
是對狀況的陳述.
是那個statement of condition.
第二個虛空.
是在加強這個狀況.
第三個虛空.
是要去重複.
第四個虛空.
還要再加強這個重複.
第五個虛空.
最終的結論.
五個虛空.
所有東西都是虛空.
虛空這個字在傳道書的一至二節這裡出現了五次.
在整本傳道書裡面出現了三十八次.
在希伯來的文法裡面.
重複是一種很重要的文學手法.
是用來強調重點.
凸顯主題.
原文的意思是.
我們的生命是短暫的.
是沒有意義的.
也控制不了.
這個是直接挑戰當時相信.
生命是有秩序.
有目的的普遍觀點.
全是虛空這句話其實是很震撼的.

$^{161}$因為當時的猶太人.
在他們傳統的信仰當中.
有智慧.
人做好事.
那個回報是有形可見的.
而生命的意義就是來自這些因素.
但是這裡竟然說.
虛空的虛空.
全是虛空.
在很多英文的譯本裡.
這個虛空的字叫做Vanity.
意思是一口氣.
Breath.
Vapour.
無氣.
Emptiness.
空虛.
這裡說了這麼多次.
好像怕死你.
你理解不到的.
為什麼傳道者要這樣說呢.
因為一章三節.
他說人一切的勞碌.
就是他在日光之下的勞碌.
有什麼益處呢.
益處這個字帶有一種經濟價值的意味.
意思是一個人勞碌了之後.
是應該有回報的.
這個字的字根.
也帶有說剩下的.
留下的意思.
傳道者在這裡不斷地問.
究竟我們人一生的勞碌.
有什麼益處呢.
我們每天上班.
有時還要加班.
很努力工作.
時間很快過了.
等到退休那一天.
突然發覺.

$^{201}$咦 我和家人的關係.
好像有點疏離.
身邊又沒有什麼朋友.
自己的健康也一般般.
銀行又不是有很多錢.
到那時候突然回頭看.
回想一下我現在的狀況.
其實我一生人得到什麼呢.
這節經文正是要我們反思.
我們人在太陽底下的辛勞.
是不是真的有價值呢.
我想我們每個人都有很多付出努力.
卻得不到回報的經歷.
塗乳液 買好衣服 做運動.
最後都被人叫做奶奶.
你很努力工作.
從來都不遲到.
隔壁的團隊需要幫忙的時候.
你義不容辭.
去年評估的時候.
你的老闆還給你很高的分數.
讚賞你.
但是今年卻失驚無神.
裁員裁了你.
搞什麼呢.
最近有個朋友跟我分享.
他奶奶和癌症搏鬥了幾年之後.
終於最近離開了.
他說有一天他兒子上學出門的時候.
回頭突然問他.
讀書讀得這麼辛苦為了什麼.
人最終都會死的.
一針見血.
給你 你會怎麼回答.
傳道書一張四至八節.
一代過去一代又來.
地卻永遠長存.
太陽上升太陽落下.
急歸所出之地.
風往南颳又向北轉.

$^{241}$不停旋轉.
繞回原路江河都往海裡流.
海卻不滿.
江河從何處流.
仍歸回原處.
萬事令人厭倦.
人不能說盡.
眼看看不飽.
耳聽聽不足.
這段經文提到三個自然界的現象.
太陽,風和海.
平衡第八節後半節提到的人.
不能說盡.
眼看看不飽.
耳聽聽不足.
反思我們人生的短暫世事的虛空.
我們身處這個地球是永遠長存.
但是我們這一代人將會被下一代取締.
然後我們下一代人再被下一代取締.
如此類推.
大自然循環不息.
無論太陽離開軌道多少都好.
最終它都會回到原本的路線.
河流,蟲雨水,融雪,海水得著供應.
水最終都會回到大海裡.
想想我們人生一生牢牢碌碌.
小時候讀書.
長大之後工作.
我們照顧我們的家人.
如果計劃得好的時候.
我們供完自己的房子.
再賣一間房子來收租.
我們積谷房機.
計劃好退休之後可以安枕無憂.
我們很小心健康.
不吃油,不吃鹽,不吃糖.
每天都吃麥皮.
小朋友都還沒出生.
我們已經開始計劃.
出生之後安排所有東西.

$^{281}$給錢補習.
我們車出車入參加課外活動.
我們很希望我們的小朋友身心靈都能夠健康成長.
我們預備到一滴不漏.
我們都希望我們的下一代比我們這一代更加好.
不單止有份工作.
分分鐘最好做醫生,律師,銀行家.
夠好味嗎.
最近我聽到一個故事.
一個老華僑將他的生意賣了一億美金.
一億美金不是港幣.
一億美金.
二十多年前這個老華僑曾經找過我老公.
幫他redesign一個office的layout.
他指示我老公.
你將office的layout要全部open.
要一眼看完.
這個office搞定之後.
大家當然很討厭這個architect.
什麼architect.
為什麼要搞到office這樣.
以後我都沒有私隱.
當然他們從此之後不可以偷懶.
老闆要redesign office的原意也是這個原因.
這個老闆的太太老闆娘.
也是出了名的勘劍.
她每一天都會去墨西哥的店舖.
那裡有個架子.
貼上熒光紅色貼紙.
是什麼來的.
就是減價快過期的食物.
她就會買回店舖煮給員工吃.
是不是值得她發財.
這個老華僑有幾個子女.
正如我們大多數亞洲人的期望.
他們每個長大之後都是事業有成.
其中一個女兒現在在另一個州做大律師.
做大律師當然很忙.
但久不久她都會坐她的私人飛機.
回來探望她的家人和親戚朋友.

$^{321}$很有心 其實這麼忙.
都特意坐私人飛機回來.
很可惜她每一次回來的時候.
她媽媽都會嘮叨她.
她說為什麼你又回來.
坐飛機很貴的.
為什麼有工作你不回來.
回來不用花錢嗎.
拜託你多存兩個錢吧.
從小到大你都是這麼大花錢.
噴到她一面被子.
久而久之這個女兒後來回來的時候.
她不是每一次都會探望她的父母.
她回來可能只是見她的兄弟姐妹和朋友.
一生勞碌.
這個老闆和老闆娘很節儉地生活.
建立了一個賣了一億的王國.
供養她的子女 個個出人頭地.
前幾年老闆娘走了.
老闆和子女的關係很疏離.
現在只剩下她一個人.
體弱多病 住在一間大屋裡.
我問你銀行裡面有一億.
或者兩億有什麼分別.
有什麼意思.
是不是很虛空.
想起這個老闆的時候.
想起我和在座的每一位.
加起來.
我們的身家應該和一億美元相差非常遙遠.
我在想其實我也很想自己這麼有錢.
如果我這麼有錢.
你知不知道我會做什麼.
你猜我會做什麼.
想不到嗎.
沒理由想不到.
錢賺錢.
我第一件事會做.
我會買一塊地.
選一塊我自己想.

$^{361}$覺得是最佳的地點.
非常好的地點.
讓我們UV能夠建立一個很適合我們用的教會.
一個可以給.
不單止給我們現在這班人.
是給我們的下一代有一個很適合神用的教會.
能夠繼續將福音拓展.
我們在地上去建立神的國.
好不好.
好.
甚至不單止這樣.
我們要升級我們所有的AV系統.
我們教務同工每一個加薪.
最好請多兩三個同工分擔我們的事工.
是不是很棒.
可惜我沒有一億.
很多年前我和我弟弟討論錢的時候.
說起我們很多父母很努力賺錢.
希望我們的家人能夠覺得舒適一點.
但是有時不知道為什麼去到某個位置.
我們開始有些迷失.
我們只顧著賺錢.
而其實我們沒有跟我們的家人相處,溝通.
建立關係.
有時甚至因為工作的壓力太大.
我們回到家裡還發我們家人的脾氣.
或者父母經常為了錢會吵架.
家無寧日.
吵了這麼多年.
錢有沒有多了就不知道了.
但是這個經常吵架的畫面.
卻深深留在我們的小朋友腦海的生命裡.
造成他們成長的過程裡.
他們經常有很多不安,焦慮.
甚至內疚.
回想起來.
我們不是想賺錢.
讓家人覺得舒適一點.
我們一家人可以開開心心.
小朋友健康成長嗎?.

$^{401}$搞什麼?.
借用這個老闆娘的例子.
你會不會也是那些很節儉.
喜歡存錢,什麼都捨不得花.
精打細算.
看著銀行存摺的銀碼就覺得很開心.
覺得很驕傲的人.
Sorry,其實現在已經不用銀行存摺了.
其實想一想有什麼值得驕傲的.
我們基督徒是不是忘記了.
其實我們只是一個管家.
神給我們的錢不是叫你收起來.
放在銀行裡.
或者你喜歡怎麼存起來.
不是這樣的.
作為一個基督徒.
我想問你.
你處理金錢的哲學.
有沒有一個永恆的價值觀在裡面?.
傳道書一章九至十一節.
他說:已有的事後必再有.
已行的事後必再行.
日光之下並無新事.
有一件事人子著說.
看,這是新的.
它在我們以前的世代早已有了.
已過的事無人記念.
將來的事後來的人也不記念.
日光之下這四個字的意思.
是說今世我們物質層面上的事.
這些事是沒有永恆的價值在裡面的.
一切在日光之下的東西到最後都是虛空.
今天似乎很新奇的事物.
很快就會變得陳舊.
1988年的時候.
我們一家回香港探親.
我舅母是做保險經紀的.
所以出出入入的時候.
她會拿著一個水壺.
很大塊的那種.

$^{441}$然後有一個接收器.
你們當中可能有些人還未出世.
這是第一代的手提電話.
想當年如果你有這個手提電話.
這個水壺.
絕對是一個身份的象徵.
所以我們回香港見到舅母的時候.
我媽和我兩個阿姨.
都特意向她借來.
去拍一張照片.
今時今日這個水壺已經是古董.
想真一點.
40年左右.
不是很久.
40年前的事.
今天我們生活在美國的文化裡面.
我們收到的訊息.
接受的教育.
工作甚至投資.
我們人似乎在當中得到一些益處.
美國的科技在世界的舞台上.
我們是節節領先的.
假設科技的進步.
是將人類發展帶去另一個高峰.
正如現在的人工智能.
AI.
我們家庭團體最近才做了一個.
AI and Parenting的workshop.
我也學了很多東西.
你們知不知道.
原來大約每四個月左右.
就有一個新的AI的version出現.
套用傳道書裡面傳道者的質疑.
他說這個進步是不是真的進步呢.
又或者是不是真的對我們有這麼大的益處.
益處一定有的.
但是我們不可以否認.
當中也有一些很大的負面影響.
例如現在其實我們經常都會聽到.
AI可能在某些產業上取代了人力.

$^{481}$導致很多工種都被取締.
現在很多人會失業.
或者需要大規模的勞動力再接受培訓.
不同的國家為了爭先.
推出最新最好的AI.
要贏過其他國家.
他們爭相投入研發.
但是當中沒有人有時間去關注.
當中衍生出來的問題.
因為你一花時間下去的時候.
你可能就慢人家一步.
當中涉及的道德問題.
大家只顧著全力以赴.
務求贏對方.
你和我的下一代的世界會變成怎樣.
傳道者說虛空的虛空 虛空的虛空.
全是虛空.
人一切的勞碌.
就是他在日光之下的勞碌.
有什麼益處呢.
這兩節經文凸顯了人類的理解能力.
和掌握的限制.
我們以為我們在掌握.
但其實可以掌握什麼呢.
他指出我們的成就.
未必真的可以幫我們得到什麼益處.
或者達到我們的願望.
我們人類的努力.
是否真的可以為人類.
做出一個理想的世界呢.
傳道者在這些問題上掙扎.
他不是想說人類的生命最終是沒有意義的.
而是如果生命真的有意義的話.
傳道者是無法理解.
就好像人捉不住風一樣.
你呢.
你好嗎.
你滿足嗎.
你得到你想要的嗎.
得到你應得的嗎.

$^{521}$你最近擔心什麼.
你怕不怕自己老啊.
你擔不擔心你存的錢不夠啊.
擔心你的計劃是否可以如願進行.
身體有沒有病.
會否因為你的身體沒有病而覺得很沮喪.
你和你家人的關係如何.
晚上睡得好嗎.
幾個月前我在朋友家裡和一大班朋友吃飯.
吃完飯之後我們一起拿了一份文件出來簽署.
這份文件叫做生前醫療遺囑.
Living Will Declaration Form.
這份文件的內容說.
如果有一天我患上無法醫治的疾病.
我的主診醫生親自判斷我已經沒救了.
我決定終止所有只是在現場我生命的治療.
我要求允許我自然死亡.
如果可以的話.
主診醫生可以按照需要的藥物營養.
或者任何被認可必需的程序提供紓緩.
如果到時我神智不清的話.
這份帶有法律權利的文件就是我的聲明.
我最終的表達.
我要求我的家人和醫生遵從我當中的意願.
我拒絕任何醫療或外科的治療.
並願意接受因此所帶來的後果.
我很蒙福很驕傲.
我的生命裡有這班朋友.
其中有幾個我已經相識超過40年.
我們都曾經年輕過.
我們一起見證我們彼此人生不同的階段.
拍拖,分手,一起哭,一起笑.
結婚,生子,養育家庭.
小孩上大學.
這兩三年來小孩一個一個出來工作.
有些甚至沒有結婚.
今時今日我們一起吃飯.
我們說的話題是最近你吃什麼藥.
然後我們竟然.
你想想吃完飯.

$^{561}$我們一起去簽這份生前醫療遺囑.
Living Will Declaration Form.
填完之後我們還好像玩遊戲一樣.
將這份表格傳到你右邊的朋友身上.
他去簽名做我的見證人.
上星期Persi大敬拜的時候.
帶了一首詩歌叫做《超出預計的愛》.
當中有一句歌詞是.
人像雨露野花對嗎.
我們人一生真的很短暫.
傳道者在傳道書一章一至十一節.
這裡提出了一個很根本性的問題.
我們人一生的勞碌到底有什麼益處.
我們的生命有什麼意義.
我們低頭.
在這段經文裡面.
傳道者沒有提供一個答案給我們.
我想我也不應該在這裡嘗試給你們什麼答案.
但我想我可以留下幾個值得思考的問題給你們.
你自己有沒有試過感受到.
在第八至十一節那裡說那種累和空虛的感覺.
我們的信仰有沒有影響到.
你思考人生的意義和目的的問題.
活在這個世界上.
有永恆的價值觀和沒有永恆的價值觀有什麼分別.
主啊求你給我們智慧.
給我們有永恆的眼光.
擺在我們裡面.
好叫我們數算我們在你面前的日子.
奉主耶穌基督的名字祈求.
Amen.
你以為得到嗎?.
\newpage



\section{使徒行傳 8:1-40}
\label{sec:rZsT_71UjNw}
\textbf{流散到地極 (使徒行傳8\_1-40) - 陳崇基牧師}
\newline
\newline
連結: \href{https://youtube.com/watch?v=rZsT_71UjNw}{\texttt{https://youtube.com/watch?v=rZsT\_71UjNw}} ~~~~ 語音日期: 2025-08-24
\newline
\newline
\hyperref[sec:Xoz1f5_xHTQ]{\small{< < < PREV SERMON < < <}}
~
\hyperref[sec:index_chronic]{\small{[返順時目]}}
~
\hyperref[sec:index_scriptual]{\small{[返順卷目]}}
~
\hyperref[sec:GgAMYWfpmPk]{\small{> > > NEXT SERMON > > >}}
\newline
\newline
使徒行傳 8:1-40
\newline
\begin{longtable}{cl}
\hline
\hline
章節 & 經文 (和合本修訂版)\\
\hline
8:1 & \begin{tabularx}{0.7\textwidth}{X} 掃羅也贊同處死他。從那一天開始，耶路撒冷的教會遭受到大迫害，除了使徒以外，眾門徒都分散在猶太和撒瑪利亞各處。 \end{tabularx} \\ \\ \relax
8:2 & \begin{tabularx}{0.7\textwidth}{X} 有些虔誠的人把司提反埋葬了，為他大大哀哭。 \end{tabularx} \\ \\ \relax
8:3 & \begin{tabularx}{0.7\textwidth}{X} 掃羅卻殘害教會，挨家挨戶地進去，拉著男女關在監裡。 \end{tabularx} \\ \\ \relax
8:4 & \begin{tabularx}{0.7\textwidth}{X} 那些分散的人往各地去傳福音的道。 \end{tabularx} \\ \\ \relax
8:5 & \begin{tabularx}{0.7\textwidth}{X} 腓利下撒瑪利亞城去，向當地人宣講基督。 \end{tabularx} \\ \\ \relax
8:6 & \begin{tabularx}{0.7\textwidth}{X} 眾人都聚精會神，同心合意地聽腓利所說的話，一邊聽他的話，一邊看他所行的神蹟。 \end{tabularx} \\ \\ \relax
8:7 & \begin{tabularx}{0.7\textwidth}{X} 因為有許多人被污靈附著，那些污靈大聲呼叫，從他們身上出來；還有許多癱瘓的、瘸腿的都得了醫治。 \end{tabularx} \\ \\ \relax
8:8 & \begin{tabularx}{0.7\textwidth}{X} 那城裡，有極大的喜樂。 \end{tabularx} \\ \\ \relax
8:9 & \begin{tabularx}{0.7\textwidth}{X} 有一個人名叫西門，向來在那城裡行邪術，自命為大人物，使撒瑪利亞的居民驚奇。 \end{tabularx} \\ \\ \relax
8:10 & \begin{tabularx}{0.7\textwidth}{X} 所有的人，從小到大都聽從他，說：「這個人就是神的能力，那稱為大能者的。」 \end{tabularx} \\ \\ \relax
8:11 & \begin{tabularx}{0.7\textwidth}{X} 他們聽從他，因他很久以來用邪術使他們驚奇。 \end{tabularx} \\ \\ \relax
8:12 & \begin{tabularx}{0.7\textwidth}{X} 當他們信了腓利所傳神國的福音和耶穌基督的名，連男帶女都受了洗。 \end{tabularx} \\ \\ \relax
8:13 & \begin{tabularx}{0.7\textwidth}{X} 西門自己也信了；既受了洗，就常與腓利在一處，看見他所行的神蹟和大異能，就覺得很驚奇。 \end{tabularx} \\ \\ \relax
8:14 & \begin{tabularx}{0.7\textwidth}{X} 在耶路撒冷的使徒聽見撒瑪利亞人領受了神的道，就打發彼得和約翰到他們那裡去。 \end{tabularx} \\ \\ \relax
8:15 & \begin{tabularx}{0.7\textwidth}{X} 兩個人下去，就為他們禱告，要讓他們領受聖靈， \end{tabularx} \\ \\ \relax
8:16 & \begin{tabularx}{0.7\textwidth}{X} 因為聖靈還沒有降在他們任何一個人身上，他們只奉主耶穌的名受了洗。 \end{tabularx} \\ \\ \relax
8:17 & \begin{tabularx}{0.7\textwidth}{X} 於是使徒按手在他們頭上，他們就領受了聖靈。 \end{tabularx} \\ \\ \relax
8:18 & \begin{tabularx}{0.7\textwidth}{X} 西門看見使徒一按手，就有聖靈賜下，就拿錢給使徒， \end{tabularx} \\ \\ \relax
8:19 & \begin{tabularx}{0.7\textwidth}{X} 說：「請把這權柄也給我，使我手按著誰，誰就可以領受聖靈。」 \end{tabularx} \\ \\ \relax
8:20 & \begin{tabularx}{0.7\textwidth}{X} 彼得對他說：「你的銀子和你一同滅亡吧！因為你想神的恩賜是可以用錢買的。 \end{tabularx} \\ \\ \relax
8:21 & \begin{tabularx}{0.7\textwidth}{X} 你在這道上無份無關；因為你在神面前心懷不正。 \end{tabularx} \\ \\ \relax
8:22 & \begin{tabularx}{0.7\textwidth}{X} 你要為你這樣的惡而悔改，祈求主，或者你心裡的意念可得赦免。 \end{tabularx} \\ \\ \relax
8:23 & \begin{tabularx}{0.7\textwidth}{X} 我看出你正在苦膽之中，被不義捆綁著。」 \end{tabularx} \\ \\ \relax
8:24 & \begin{tabularx}{0.7\textwidth}{X} 西門回答說：「請你們為我求主，使你們所說的，沒有一樣臨到我身上。」 \end{tabularx} \\ \\ \relax
8:25 & \begin{tabularx}{0.7\textwidth}{X} 使徒既作了見證，並且宣講了主的道，就回耶路撒冷去，一路在撒瑪利亞好些村莊傳揚福音。 \end{tabularx} \\ \\ \relax
8:26 & \begin{tabularx}{0.7\textwidth}{X} 有主的一個使者對腓利說：「起來！向南走，往那從耶路撒冷下迦薩的路上去。」那路是曠野。 \end{tabularx} \\ \\ \relax
8:27 & \begin{tabularx}{0.7\textwidth}{X} 腓利就起身去了。不料，有一個埃塞俄比亞人，是個有大權的太監，在埃塞俄比亞女王甘大基的手下總管銀庫，他上耶路撒冷去禮拜。 \end{tabularx} \\ \\ \relax
8:28 & \begin{tabularx}{0.7\textwidth}{X} 回程中，他坐在車上，正念著以賽亞先知的書， \end{tabularx} \\ \\ \relax
8:29 & \begin{tabularx}{0.7\textwidth}{X} 聖靈對腓利說：「你去！靠近那車走。」 \end{tabularx} \\ \\ \relax
8:30 & \begin{tabularx}{0.7\textwidth}{X} 腓利就跑到太監那裡，聽見他正在念以賽亞先知的書，就說：「你明白你所念的嗎？」 \end{tabularx} \\ \\ \relax
8:31 & \begin{tabularx}{0.7\textwidth}{X} 他說：「沒有人指教我，怎能明白呢？」於是他請腓利上車，與他同坐。 \end{tabularx} \\ \\ \relax
8:32 & \begin{tabularx}{0.7\textwidth}{X} 他所念的那段經文是這樣：「他像羊被牽去宰殺，又像羔羊在剪毛的人手下無聲，他也是這樣不開口。 \end{tabularx} \\ \\ \relax
8:33 & \begin{tabularx}{0.7\textwidth}{X} 他卑微的時候，得不到公義的審判，誰能述說他的身世？因為他的生命從地上被奪去。」 \end{tabularx} \\ \\ \relax
8:34 & \begin{tabularx}{0.7\textwidth}{X} 太監回答腓利說：「請問，先知說這話是指誰，是指自己，還是指別人呢？」 \end{tabularx} \\ \\ \relax
8:35 & \begin{tabularx}{0.7\textwidth}{X} 腓利就開口，從這段經文開始，對他傳講耶穌的福音。 \end{tabularx} \\ \\ \relax
8:36 & \begin{tabularx}{0.7\textwidth}{X} 二人正沿路往前走，到了有水的地方，太監說：「看哪！這裡有水，有甚麼能阻止我受洗呢？」 \end{tabularx} \\ \\ \relax
8:37 & \begin{tabularx}{0.7\textwidth}{X} 腓利說：『你若是一心相信，就可以。』他回答：『我信耶穌基督是神的兒子。』 \end{tabularx} \\ \\ \relax
8:38 & \begin{tabularx}{0.7\textwidth}{X} 於是他吩咐把車停下來，腓利和太監二人一同下到水裡，腓利就給他施洗。 \end{tabularx} \\ \\ \relax
8:39 & \begin{tabularx}{0.7\textwidth}{X} 他們從水裡上來，主的靈把腓利提了去，太監再也看不見他了，就歡歡喜喜地上路。 \end{tabularx} \\ \\ \relax
8:40 & \begin{tabularx}{0.7\textwidth}{X} 後來有人在亞鎖都遇見腓利；他走遍那地方，在各城宣揚福音，一直到凱撒利亞。 \end{tabularx} \\ \\
[1ex]
\hline
\hline
\end{longtable}
$^{1}$(播放).
各位影子妹早安.
很高興今天見到大家.
上星期我來不到大家當中.
我去了一個山明水秀的地方.
就是San Diego.
在那裡主講一個退修會.
讓我們能夠鼓勵不同的教會.
讓我們回到主題「跨越」.
我們一直在講的是《史萊恆傳》.
《史萊恆傳》如果有聖經的話.
和我一起打開好嗎.
我們不需要依賴PowerPoint.
但我們這次繼續講主題.
講到第八章.
如果有《史萊恆傳》打開第八章.
我們先講上次的第六章.
第六章有什麼事.
第六章就是一個公認的問題.
因為那裡有兩組人.
一組是講希伯來話.
另外一個教會的群組是講希臘話.
他們的寡婦是不是有公認的問題.
所以試圖選立一些執事.
來配搭一起去事奉.
在講第八章之前.
當然要講第七章.
第七章就是講到史提反的港獨.
在這裡我想去到第七,第八章.
去到選一些章節來認識一位.
就是《史萊恆傳》裡面講的神.
是一位怎樣的神.
首先講第七章.
史提反的港獨.
史提反的港獨有什麼特別的地方.
就是他們.
大家都知道可能之前聽到可能聽過.
史提反的港獨是一個很長的港獨.
也是很討人厭的港獨.
所以那些猶太人很不喜歡.

$^{41}$非常長的.
不同的主題.
有什麼發展.
譬如我們講到神的榮耀.
你看到那裡是箭嘴.
神的榮耀就是神的同在.
這個主題有什麼發展.
第一個當然就是伊甸園.
在上帝與人類的同在.
人類就好像祭司一樣.
來敬拜祂,事奉祂.
當然人是叛逆.
然後失去了同在.
但是上帝沒有放棄人類.
上帝沒有放棄有罪的人.
然後繼續發展這個故事的時候.
就到了埃及之後.
他們在那裡抗議.
上帝說我要和你們一起同在.
就建了回眸.
回眸就是上帝與人類同在的一個記號.
Tabernacle.
然後他們入主迦南之後.
去到這個應許之地.
就是耶路撒冷的聖殿.
然後慢慢一直發展的時候.
當然以色列國仍然是叛逆.
然後亡國.
然後被擄回歸之後.
等候到真正的尼泊爾出現.
就是耶穌基督.
然後基督到了的時候.
基督就是聖殿.
基督就是我們能夠和上帝相交的地方.
你看到有一個演進.
伊甸園.
回眸.
然後聖殿.
到基督.
然後到將來的時候.

$^{81}$就是啟示錄最後那幾章.
就是講到新天新地.
這個就叫做聖經神學.
其實史提反就是在做這件事.
從頭到尾就是和這些猶太人的會眾.
來解釋一下到底上帝做了什麼.
譬如七章第二節八節.
他說什麼.
他說亞伯拉罕.
在美索不達米亞的時候.
就是以色列之外的地方.
上帝.
榮耀的上帝向他顯然.
然後帶領他去到迦南.
然後有他的以撒亞國.
然後十二眾子.
然後他們就到了埃及.
記得嗎.
這個是創世記最後那幾章.
然後約瑟在那裡死.
以色列人就被下一個國度.
就是埃及的法羅.
去到奴役了好幾代.
然後就選出摩西.
摩西就在埃及的皇宮長大.
然後就帶領以色列人出埃及.
但是以色列人當時是不接受他的.
他就逃難到了米田地.
第三十節開始.
終於在西乃曠野的時候.
耶和華榮耀的神.
就在荊棘的火焰裡面.
向上帝顯然.
然後叫他回去埃及.
叫埃及的以色列人能夠出來.
經過紅海.
去到西乃山.
去到曠野.
四十年.
然後會幕在中間.

$^{121}$當然.
正如我所說.
三十九至四十三章.
史提反就說.
以色列人其實是不信的.
最後他們是流亡於巴比倫.
四十四節他回來.
他就說在會幕裡面.
上帝是跟神的子民同在的.
但是到最後的時候.
為什麼這篇道理這麼討人厭呢.
他就說在所羅門的聖殿.
就是說他指著.
在耶和撒冷說的.
在這篇道理.
耶和撒冷最輝煌的.
最厲害的是什麼.
就是他們的聖殿.
耶和撒冷的聖殿.
他就說這個聖殿是怎樣的.
是至高者不住在那裡的.
就是人手做的.
會幕就是按照上帝的樣式來做的.
摩西照辦主觀的做.
因為天上給他啟示的樣式和尺寸就是這樣.
但是聖殿是怎樣的.
是人手做的.
這個很大挑戰.
為什麼.
因為猶太人覺得最輝煌的東西.
就是他們的聖殿.
在舊約裡面.
最輝煌的臨在上帝的離道.
就是所羅門的聖殿.
現在他們有的聖殿.
在耶穌時代的.
在使徒時代的.
就是第二聖殿.
那個聖殿就是一個很好的.
和很輝煌的東西.

$^{161}$如果你說我們的上帝不在那裡住的話.
這個是最大的侮辱.
如果我告訴你.
你今天去耶和撒.
今天的耶和撒.
Old City.
舊城.
你要去得罪猶太人的話很容易的.
你不是抵禮他們吃的東西.
沒有我們的東西那麼好.
你也不是抵禮他們的語言.
沒有我們廣東話那麼好.
這些是抵禮不了他們的.
你侮辱他們的就是西牆.
Western Wall.
如果你拿一條bacon進去西牆的話.
我告訴你.
你出不來了.
因為他們看那個地方.
就是一個最聖潔最神聖的地方.
偏偏.
史提反就說.
上帝的榮耀不在你中間.
這些只不過是一些人手所做的事情.
於是他們對他的控告就是.
七章第十三節.
這個人史提反.
不斷的說話是侮辱神聖的地方和我們那裡的法.
所以這個人說完之後.
就應該要死了.
因為他在說什麼呢.
太過討人厭.
所以這個是很大的侮辱.
如果你說.
譬如我剛剛去完其他教會講道.
如果我跟大家講.
跟當地的會眾講.
我說上帝的榮耀.
上帝的工作.
到處都有.

$^{201}$在芝加哥.
在城會很多.
在西岸.
在東岸很多很多.
偏偏就不在聖地利哥這裡.
你以為我能回來嗎.
就因為這麼大的侮辱.
我想你能夠理解那個效果.
而史提反就是在說這件事.
但最主要我們想說什麼.
不是說偏道討人厭這麼簡單.
就是偏道想告訴我們.
上帝是一個怎樣的上帝.
上帝是不會被地域所困惑的.
不會被他限制的.
這是普世的上帝.
那些猶太人以為.
這個上帝可以塞在我們所做的聖殿裡面.
No.
史提反不是.
因為這個上帝是一個universal的上帝.
一個普世的上帝.
他是愛世人的上帝.
你看到他整個歷史裡面.
上帝出現的地方是什麼地方.
美索不達米亞.
西乃的曠野.
紅海.
埃及.
回望在曠野裡面.
四十年.
全部都是在迦南地以外的.
全部都是在耶路撒冷以外的.
偏偏就不在耶路撒冷.
這個就是史提反裡面所說的.
和整個史提反行傳裡面所說的.
上帝是一個怎樣的上帝.
是一個普世的上帝.
是一個萬邦的上帝.
大家不要搞錯.

$^{241}$他不是廣東人的上帝.
當然他拯救我們.
但他不只是我們的上帝.
他也是阿拉伯人的上帝.
他是萬邦的上帝.
如果你留意的話.
史提反主持一直告訴我們.
介紹這位上帝.
就是一位不受地域地理限制的.
因為聖靈在各教會作供.
都沒有停過.
你未必有信心眼睛看得到.
不過你留意.
如果你能夠有一個eye of faith的話.
聖靈是不斷有工作的.
剛剛去完那個weekend.
在聖地利哥教會.
那些弟兄姊妹很熱心.
這個要讚一讚.
廣東話的信徒.
廣東話人的命是苦命.
他們真的很努力去作供.
就是事奉.
對他們基督徒的生命.
這種價值.
做到死的價值.
他們就轉移了下來.
廣東話的基督徒的生命力.
真的很強.
耶穌大教會說到.
今天他們經過芝加哥.
來一來.
你知道芝加哥裡面最大的景點是誰嗎.
就是UV沒錯.
所以他們的巴士來了.
他們剛剛買了我的巴士.
Splinter是我的夢幻車.
我馬上出去拍照.
等一下如果有機會.
你可以出去拍照.

$^{281}$太好了.
他們經過特意來看我們的探訪.
我們非常歡迎大家.
在明尼蘇拉教會都是一樣.
上帝的聖靈不斷的來供作.
你要有eyes of faith.
你才能夠去見到.
這個就是史蘭傳想打開給我們的神學.
所以在史提反被害之後.
發生什麼事呢.
第八章第一節.
蘇羅也贊同處死他.
處死史提反.
大家看到嗎.
不是啪啪啪.
從那一天開始.
耶路撒冷的教會遭受大的迫害.
除了使徒以外.
眾門徒都分散在猶太地和撒瑪利雅各處.
有些虔誠的人將史提反埋葬了.
為他大大哀哭.
蘇羅卻殘害教會.
逐家逐戶的進去抓他們的男女關在監獄裡.
第四節.
那些分散的人往各地傳主的福音.
你留意到嗎.
這裡有一個字.
是出現了兩次的.
第八章第一節到第四節.
就是分散這個字.
Diaspora.
Diaspora就是我們今天用的那個字.
Diaspora這個字.
分散的意思.
大家記得我講六章的時候.
我們不是說過有猶太人的分散問題.
猶太人亡國之後.
他們流散到不同的地方.
這個叫做Jewish Diaspora.
現在這裡說的是什麼Diaspora呢.

$^{321}$是Christian Diaspora.
上帝將這些初期教會.
不是很多人.
將他們分散.
從耶路撒冷.
猶太.
直到撒瑪利亞和地極.
不同的地方將他們分散.
用什麼方法將他們分散.
這裡很奇怪.
他們不是說開了什麼宣教年會.
我們要猜傳.
然後傳出去.
然後去到不同的地方.
我們很有計劃的.
在不同的宣教的工場.
我們要做一些事.
然後派宣教士出去.
他們分散的原因.
不是這個原因.
是什麼.
是受迫害.
蘇羅未信主之前.
去到迫害當地的教會.
就造成一個Christian Diaspora.
流散基督徒的現象.
基督徒不只是集中在一個地方.
他們被差派去不同的地方.
但上帝出奇的地方就是.
不是有什麼福音冊.
而是容許教會受到打壓.
然後這些基督徒帶著分散的.
流散的心情.
但帶著福音的使命.
好像上帝自己一樣.
上帝不只是困在耶路撒冷的上帝.
上帝是到處去的上帝.
他們就好像上帝的形象一樣.
跨越到不同的境界.
不同的限制.

$^{361}$進去不同的社群.
然後將福音帶進去裡面.
這叫Diaspora.
是Christian的Diaspora.
今時今日.
都是一個Diaspora的狀態.
我們因為找工作.
或者會流,去流等等.
我們有很多流散的現象.
華人我們自己就是流散的一代.
並且這一代流散的更多.
香港人特別.
我是香港人.
有一個很普遍的海外流散的情況.
就是你去到不同的地方.
那個不像你的家.
你想生根,但是生不出的.
你很嚮往以前的生活.
但以前的生活.
以前的時代已經過了.
以前的地方也變了.
所以你去到不同的地方.
語言可能不太合適等等.
食物不太適合.
經常吃來吃去的東西.
都是不合口味.
嘴唇不太合適.
所以我們經常找東西吃.
但我這裡想說的是Christian Diaspora.
是基督徒被流散.
原來這是上帝隱藏的心意.
是要我們流散到不同的地方.
請你記住.
是帶著福音被流散.
你要經常問自己.
到底我的身份是什麼.
我的身份不只是一個廣東人.
不是一個華人.
不只是一個港人這麼簡單.
無論你是什麼人也好.

$^{401}$你最重要的身份.
不是地上的身份.
因為這些身份都會過去的.
你做廣東人也好,台山人也好,香港人也好.
你做了多久.
我給你多50年好嗎.
你都玩完了.
身份只不過是暫時的.
但你一個永恆的身份.
不變的身份是什麼.
是天國的子民.
我們每個人都是帶著福音到處去.
這是流散的使命.
另外第八章.
他給了一些故事.
施行尊告訴我們一些流散的工作.
說這些使徒或領袖.
這些福音的故事.
都是走到不同的地方.
第八章五到二十五節.
說到一個人去到跨越.
第八章第五節.
費利下到撒瑪利亞城.
費利是什麼人.
猶太人.
在耶路撒冷.
因為不同的緣故要流散.
他就下到撒瑪利亞城.
為什麼要下去.
因為你看到中央山地的地方.
好像脊骨一樣.
全部都是山.
一個是山脊.
耶路撒冷就是最高的地方.
然後他往北走的時候.
其實就是跟著今天的公路.
大概是跟著山脊下去.
一直下去下面.
就是一個很重要的地方.
就是撒瑪利亞城.

$^{441}$向當地人傳講.
基督.
中人都聚精會神.
然後排排坐下來聽.
一邊聽他講一邊看他所行的神蹟.
有很多人被污靈所輔助.
污靈就被趕出來了.
然後癱瘓的.
瘸腿的得到醫治.
這些全部都是耶穌神國來到的時候.
那些劃時代的神蹟出現.
現在聖靈來到了.
就來到一個.
猶太人是不會去的地方.
就是撒瑪利亞地.
同樣出現這個現象.
那城就有極大的起來.
這個流散就很出奇了.
因為猶太人去一個.
他完全不會去的地方.
就是這個腓力.
是七個執事裡面其中一個位置.
他往北走然後下山.
受了.
這裡繼續講下去的時候.
就發現為什麼猶太人是不會去撒瑪利亞的地方.
昨晚我看了一些資料.
有一個叫約瑟夫.
是第一世紀的猶太人歷史學家.
他就說猶太人和撒瑪利亞人中間有什麼牙齒人.
他起碼提了六件新聞的事件.
是當時猶太人和撒瑪利亞人.
是不對的.
是不妥的地方.
他提了六件事件.
我只說第四件事.
第四件事就是撒瑪利亞人在主後一世紀初.
差不多是史行傳的時間.
撒瑪利亞人在晚上就潛入耶路撒冷.
當然耶路撒冷是可以去的.

$^{481}$因為外邦人是可以去耶路撒冷的.
只不過你不可以進去以色列人院.
所以晚上的時候.
有些撒瑪利亞人就潛入了耶路撒冷.
然後夜間的時候.
他們拿著一些骨頭.
然後在聖殿不同的地方就撒.
什麼骨頭.
你以為是豬骨吧.
豬骨已經是侮辱了.
他們拿人骨.
不知道在哪裡找了人的骨頭.
然後就撒在聖殿的場所裡.
然後大家就造成一個很大的衝突.
所以之前之後也發生了不同的事情.
有些仇殺等等.
但我想說這個.
因為我覺得這個是最有趣的.
把人骨撒在聖殿.
你說這些人活該嗎.
猶太人覺得這些人真的很活該.
這樣來污衊我們的聖所.
等於今天有人來我們教會.
在我們早餐那裡下射藥.
搞到我們敬拜不了.
怎麼搞啊.
然後牧師傳道人把他打了一頓.
然後你講不出道理.
嘴巴都歪了.
那你就敬拜不可以進行.
這個就是撒馬尼人為什麼要把人骨撒在聖殿裡.
讓猶太人還是禦節的時候.
因為當時他們最大的節日.
沒有降在他們的身上.
信了主之後他們還沒有降在他們身上.
然後在耶路撒冷的使徒聽見撒馬尼人受了上帝的徒.
就打發誰呢.
當然是他們的長老.
撒馬尼人信主.
成為我們一份子.

$^{521}$成為神的子民.
所以耶路撒冷的總會就差派兩個人去.
就是彼得和約翰.
到他們那裡.
兩個人下去的時候就為他們禱告.
又讓他們領受了聖靈.
因為聖靈還沒有降在他們任何一個人身上.
他們只奉耶穌的名受洗.
於是使徒按守在他們頭上.
他們就領受了聖靈.
這裡有一樣東西要看.
我們經常看《士郎行傳》的時候.
聖靈下來的時候一定會有一些現象.
譬如說方言.
《士郎行傳》告訴我們其實不是一個公式.
你千萬不要把第二章說.
這件事永遠都是一種常規.
這樣是解錯經的.
因為你看到撒馬尼人信主.
他們的經歷是怎樣.
受了聖靈.
但是卻沒有第二章裡面發生的事.
因為聖靈的工作不是公式的.
聖靈的工作不是我們能夠操控的.
無論如何這裡介紹的就是.
猶太人接觸傳道給撒馬尼人.
這個就是破天荒.
於是使徒就下來.
彼得和約翰.
這個有趣的地方就是.
你記得約翰.
特別對約翰來說.
約翰很憎撒馬尼人的.
你記得路加福音.
就是《士郎傳》之前那一集.
說什麼呢.
耶穌派門徒出去撒馬尼亞地傳福音.
誰知道約翰和他的兄弟雅各.
他們去到撒馬尼的村莊的時候.
撒馬尼人是否接受他們.

$^{561}$不接受他們.
然後他們就問.
叫耶穌做什麼.
他說降火燒死他們.
因為他們這群人是該死的.
是不是該死.
因為民族上都覺得他們該死.
何況將福音傳給他們.
他們居然不接受.
反對我們.
耶穌怎麼說呢.
耶穌轉身來責備兩個門徒.
這一趟就很奇怪.
後期發生什麼事呢.
約翰和彼得.
他們下去真的看一下.
上帝在做什麼.
居然出人意外的地方.
就是上帝的福音.
去到一些他們眼中覺得.
是不應該得救的人.
不應該得救的人.
這是一種福音跨越的情況.
我們今天都是一樣.
我一向都在說這一年.
是在說跨越.
跨越.
我們一定要有一種心胸.
培養自己的一種心胸.
是福音的心胸.
是上帝有那種心胸.
是有些人和我們不同.
說的話可能和我們不同.
文化可能和我們不同.
可能會氣死我們.
剛剛那星期我去到聖地牙哥.
我有一個mentor在那裡.
大家說一些做國語事工的難處.
他們都說了一件事.
他們說真的很煩.

$^{601}$他們不像我們廣東人這樣.
是commit的.
因為我們廣東人是容易commit的.
去毀身.
他們都是一樣.
遇見了一個問題.
說就天下無敵.
做起來就沒力氣.
我不知道你們是不是這樣.
你們是不是這樣.
你們廣東人應該不是這樣.
我們做就天下無敵.
說可以不說.
但他們就反過來.
所以令我們很生氣.
但我想問這個問題.
《少林傳》裡面說的上帝就是這樣的上帝.
Yes 他們和我們不同.
但我們要帶著福音進去.
盼望福音能夠改變他們.
這就是一個跨越的勇氣.
第三點.
這裡說到上帝是一個克服地極限制的上帝.
他們都希望教會有一個這樣的視野.
就是地極的視野.
就是讓我們的vision我們看到的東西.
是看到最終最終最遠最遠的地方.
讓我們不要近視.
八章第26節.
有一個主的使者對費利說.
當時他又回到耶路撒冷.
他說對他說什麼呢.
他說起來向南走.
往那從耶路撒冷下加薩的路上.
那路是曠野.
費利就起來了.
不料他下去的時候.
就有一個埃塞俄比亞的人.
他是一個大權的太監.
埃塞俄比亞女王金大基的手下.

$^{641}$他是管錢的.
原來他是federal reserve的.
他上耶路撒冷去敬拜.
你們看到那些箭槌嗎.
我想這個人敬拜是為什麼呢.
這個是一個外邦人.
埃塞俄比亞的甚至是太監.
為什麼他會上耶路撒冷去敬拜一個外邦人.
因為當時.
有些外邦人是被猶太教吸引的.
我們英文叫做guard fear.
就是敬畏神的人.
他們不是猶太人.
不過他們被猶太教那種一神論.
一神的神學.
不是滿天神佛那種.
他們被一神論monotheism.
他們所影響.
於是他們很多時候.
叫做插半隻腳進去.
猶太教.
於是有大時大節的時候.
他們就會上耶路撒冷.
就好像猶太人一樣.
不過他們進不了聖殿最重要的地方.
這個埃塞俄比亞的太監.
大概應該是這樣的背景.
他是被吸引.
然後跟那些流散的猶太人一樣.
他們會上去耶路撒冷過節.
所以這裡說到他會晨的時候發生什麼事呢.
就是遇見到費利.
他往南走.
從耶路撒冷下去加薩的路上.
這個人就是上耶路撒冷敬拜.
然後回程的.
你看到那個地理.
我剛才說的那個中央山地.
是一個骨脊.
當你向左邊走的時候.

$^{681}$其實你走到去平原.
那個叫做沿海平原.
藍色那條線其實是今天的.
以往都是.
古代到今天都是一些大的公路.
當時叫做via maris.
是一些行軍和商業的大道.
藍色那條線via maris.
你看到從加利利右邊.
一直下去到下面.
今天的加薩這個地方.
再下去就是埃及的地方.
所以他走到一個國際性的地方.
international zone.
他做什麼呢.
找到這個太監.
故事是這樣說.
回程之中這個太監坐在車上.
他有錢當然坐benz.
坐在上面做什麼.
他就在看伊塞亞先知的書.
聖靈對弗利說.
他說你靠近那個車.
跟著那個車 聽機.
然後弗利就跑到太監那邊.
聽聽他原來他在讀伊塞亞先知所說的話.
然後就問他一句.
你讀的你明不明.
查經的時候你明不明.
這個太監就告訴他.
請問先知伊塞亞.
這話是指誰呢.
是指自己還是指別人呢.
這個大概是說.
他在看的就是我們非常熟悉的一段聖經.
就是伊塞亞書底五十三章.
那裡說受苦的僕人.
被鞭打不還口等等.
為我們受刑罰等等.
到底這個是指自己.

$^{721}$那個僕人指自己.
還是指其他人呢.
弗利就跟他開口.
就跟他解一段聖經.
告訴他.
裡面這個受苦的僕人.
一定是耶穌基督自己.
因為只有他fit所有的profile.
所以兩人就沿路.
就上一路查經的時候.
就看到有水的地方.
太監就跟他說我信.
我在這裡受洗行不行.
有什麼可以阻止我.
去受洗.
這裡一個很有趣的故事.
道家為什麼要記載這個.
外邦的太監.
去到歸主呢.
你記得我曾經說過.
小男子的劇本是什麼.
是伊塞亞書.
這個太監就是在讀伊塞亞書.
因為當時最流行的金句.
最多人看的聖經.
就是伊塞亞書.
特別是伊塞亞書的下一半.
下一半說什麼.
下一半說了很多不同的主題.
有七個主題.
有上帝的國度歸位.
十二支派會被重整等等.
上帝的國度來到的時候.
第六點說什麼.
是那些以前不會被接納的人.
現在會被接納.
有什麼人呢.
是外邦人和太監.
你看到嗎.
與耶和華聯合的外邦人.

$^{761}$不要說耶和華將我和他的子民分別.
我是不屬於他的子民.
No.
耶和華做了一件新事.
耶穌基督裡面所做的.
就是One people of God.
所有的種族都能夠成為這個One people.
甚至太監.
太監也不要說.
說我是夫婿.
為什麼太監是夫婿.
因為沒有下一代.
也不可能有下一代.
然後說如此說.
謹守我誡命的.
選擇我旨意的.
持守我約的太監.
他們可以做什麼.
可以進我殿裡面.
在我牆裡面.
有紀念碑有名號.
勝過自己肉身有異類.
我必賜他們永遠的命.
不能夠被剪除.
剛才我說的以賽亞書.
受苦的僕人是哪一章.
53.
我想他應該是這樣的.
以前看的聖經不是電話那樣滾動的.
也不是印出來的聖經一頁一頁.
是一卷的.
大概53 56.
你這樣打開的時候.
其實是同一個畫面.
同一個螢光幕.
53 56 他就說.
之後說到這個太監可以成為神的子民.
這個僕人可不可以為我做這件事.
你再解釋的時候.
他就發現原來這個子有用.

$^{801}$原來聖經裡面所說的話.
是和我有關的.
所說的福音.
耶穌基督要來做的事.
是和我有關的.
不是和他人有關.
是和我自己有關.
所以太監就信了然後受洗.
但這裡說的這個故事就是.
這個肥理.
他進入一個國際的地方.
一個國際範圍.
因為上帝是一個國際的上帝.
一個國際的上帝.
我記得好幾個月前.
我參加了一個網上的.
有人邀請我去一個叫做群體研究.
一個群體研究.
就是說一些牧者.
他就找了一些美加的牧者.
大概有十幾人.
我們大家去討論.
我們在北美的華人教會.
我們有很多挑戰.
因為北美華人教會是非常複雜的.
因為我們不是一個語文.
也不是同一個文化.
我們是一個很雜的文化.
起碼大部分的教會.
都是兩文三語的.
所以很多時候出現一個情況就是.
不會自然而然大家會合得來.
你有沒有組裝過LEGO.
那個東西要插進去才能組裝.
但華人教會.
北美的教會很多時候就是這樣.
就是有大小不同的LEGO.
有些雜牌的LEGO.
有些正牌的.
我們怎樣組裝呢.

$^{841}$是有些困難的.
所以我也認同.
我們每個牧者所說的.
其實真的有不同的難度.
但縱使文化有距離.
但我們見到《史板傳》.
裡面那位上帝是怎樣的上帝.
我們所信的上帝.
我們所敬拜的上帝.
是一個怎樣的神.
是一個打破所有界限的上帝.
既然上帝是這樣的話.
我們應該是怎樣的.
所以我們在Crossover裡面.
我們不斷去挑戰大演之妹.
求主給我們智慧給我們力量.
給我們那個韌度.
我們能夠踏出我們自己舒服的地方.
大家知道耶穌基督不是廣東話的神.
不只是廣東話的神.
是任何一種話的主宰.
是任何一個文化任何一個語言的主.
我們就是那些所謂跨文化的宣教士.
我們永遠都是跨文化.
永遠都是挑戰我們自己.
不要躲在我們的舒服中.
不只是說我們自己的話.
吃我們自己喜歡的東西.
這個固然是好.
世界上那麼多種食物.
有什麼比我們的粵菜好.
難道天天都去吃火鍋嗎.
當然我們最好.
我們每個人都覺得自己是最好的.
因為是最舒服的.
但這位的上帝不是一個困在某一個語言某一個群組的上帝.
我們要走出去.
要學人家的語言.
要學學接觸其他人.
是有困難的.

$^{881}$但教會要繼續向前.
使徒長矢就是挑戰我們.
千萬不要躲在我們自己的舒服中.
要我們走出去.
求主幫助我們教會.
認識這位的上帝.
萬邦萬民的上帝.
我們做上帝的教會.
做上帝的子民.
都是一個要向著萬邦萬族的上帝.
今天有幾樣東西要想.
有什麼人是我們需要跨過去.
才能成為一個更加有力的見證.
這個可能不只是語言界限的問題.
可能是一些關係上的界限.
可能是興趣上的界限.
總之就是不是很自然.
但你記得我們的身分是什麼嗎.
是使徒.
是宣教士.
走出我們自己不舒服的地方.
我們才會經歷到.
這位上帝就是《史獨行傳》裡說的上帝.
和其他弟兄姊妹一起來同行.
那些弟兄姊妹不一定是同聲同氣.
這也沒問題.
應該是這樣的.
是大家不配合的.
不配合是怎樣.
我們給時間.
給點力水.
我們弄到他們配合.
大家喜不喜歡我們敬拜音樂.
喜歡吧.
我們也算不錯.
雖然我們每個人都是業餘.
但我們也要時間磨.
他厲害 她厲害 大家OK.
但走在一起很奇怪.
總是要時間的.

$^{921}$要力水的.
教會的關係其實也是一樣.
是需要承諾.
求主幫助我們.
開童心得來祈禱.
因為你是那位跨國 跨種族 跨語言 跨地域限制的上帝.
神愛 多謝你.
因為你是那位愛萬民的上帝.
主愛就認識你.
我們能夠成為你子民的一部分.
求主你幫助我們.
從我們的禱告和我們的心態.
我們的取態和我們的行為上.
讓我們能夠成為你自己的養成.
更加有你自己的養成.
能夠跨出我們自己的局限.
以致你的福音能夠傳得更遠.
以致你的教會能夠成長.
主這個禮拜差遣我們去不同的地方.
讓我們活出我們這個身份.
代表你的福音.
在我們家裡 在我們工作裡.
在我們周圍的地方.
讓我們能夠成為你的見證.
祈禱給主的名堂.
奉元主賜大家平安.
多謝各位.
\newpage



\section{使徒行傳 6:1-15}
\label{sec:GgAMYWfpmPk}
\textbf{從磨擦到拓展 (使徒行傳6\_1-15) - 陳崇基牧師}
\newline
\newline
連結: \href{https://youtube.com/watch?v=GgAMYWfpmPk}{\texttt{https://youtube.com/watch?v=GgAMYWfpmPk}} ~~~~ 語音日期: 2025-08-30
\newline
\newline
\hyperref[sec:rZsT_71UjNw]{\small{< < < PREV SERMON < < <}}
~
\hyperref[sec:index_chronic]{\small{[返順時目]}}
~
\hyperref[sec:index_scriptual]{\small{[返順卷目]}}
~
\hyperref[sec:xrB_pn8UZeQ]{\small{> > > NEXT SERMON > > >}}
\newline
\newline
使徒行傳 6:1-15
\newline
\begin{longtable}{cl}
\hline
\hline
章節 & 經文 (和合本修訂版)\\
\hline
6:1 & \begin{tabularx}{0.7\textwidth}{X} 那些日子，門徒增多，有說希臘話的猶太人向希伯來人發怨言，因為在日常的供給上忽略了他們的寡婦。 \end{tabularx} \\ \\ \relax
6:2 & \begin{tabularx}{0.7\textwidth}{X} 十二使徒叫眾門徒來，說：「我們撇下神的道去管理飯食，是不合宜的。 \end{tabularx} \\ \\ \relax
6:3 & \begin{tabularx}{0.7\textwidth}{X} 所以弟兄們，當從你們中間選出七個有好名聲、滿有聖靈和智慧，我們派他們管理這事。 \end{tabularx} \\ \\ \relax
6:4 & \begin{tabularx}{0.7\textwidth}{X} 至於我們，我們要專注於祈禱和傳道的事奉。」 \end{tabularx} \\ \\ \relax
6:5 & \begin{tabularx}{0.7\textwidth}{X} 這話使全會眾都喜悅，就揀選了司提反—他是一個滿有信心和聖靈的人；他們又揀選了腓利、伯羅哥羅、尼迦挪、提門、巴米拿，並皈依猶太教的安提阿人尼哥拉， \end{tabularx} \\ \\ \relax
6:6 & \begin{tabularx}{0.7\textwidth}{X} 叫他們站在使徒面前，使徒禱告後，就為他們按手。 \end{tabularx} \\ \\ \relax
6:7 & \begin{tabularx}{0.7\textwidth}{X} 神的道興旺起來；在耶路撒冷門徒數目增加得很多，也有許多祭司聽從了這信仰。 \end{tabularx} \\ \\ \relax
6:8 & \begin{tabularx}{0.7\textwidth}{X} 司提反滿有恩惠和能力，在民間行了大奇事和神蹟。 \end{tabularx} \\ \\ \relax
6:9 & \begin{tabularx}{0.7\textwidth}{X} 當時有從稱為「自由人」會堂，並古利奈、亞歷山大會堂來的人，還有些從基利家、亞細亞來的人，起來和司提反辯論。 \end{tabularx} \\ \\ \relax
6:10 & \begin{tabularx}{0.7\textwidth}{X} 司提反是以智慧和聖靈說話，眾人抵擋不住， \end{tabularx} \\ \\ \relax
6:11 & \begin{tabularx}{0.7\textwidth}{X} 就收買人來說：「我們聽見他說褻瀆摩西和神的話。」 \end{tabularx} \\ \\ \relax
6:12 & \begin{tabularx}{0.7\textwidth}{X} 他們又煽動百姓、長老和文士，就突然來捉拿他，把他帶到議會去， \end{tabularx} \\ \\ \relax
6:13 & \begin{tabularx}{0.7\textwidth}{X} 設下假見證，說：「這個人不斷地說話，侮辱神聖的地方和律法。 \end{tabularx} \\ \\ \relax
6:14 & \begin{tabularx}{0.7\textwidth}{X} 我們曾聽見他說，這拿撒勒人耶穌要毀壞這地方，也要改變摩西所交給我們的規矩。」 \end{tabularx} \\ \\ \relax
6:15 & \begin{tabularx}{0.7\textwidth}{X} 在議會裡坐著的人都定睛看他，見他的面貌好像天使的面貌。 \end{tabularx} \\ \\
[1ex]
\hline
\hline
\end{longtable}
$^{1}$好 今日講到這裡 Mark Sir.
好 我們重新打開聖經.
我們繼續這個系列.
講到跨元 因為我們在講哪本書.
是《士郎行傳》 是不是.
我們一章一章地講 今次講到第六章.
我們看《士郎行傳》的時候 很多時候我們留意一些主角.
譬如有誰 彼得 約翰 保羅 巴拿巴.
但很少留意一個主角是教會.
教會也是一個主角.
在《士郎行傳》的黑話裡面 想告訴我們當時的教會是怎樣.
然後教會的一些故事 教會怎樣去經歷一些事.
然後《士郎行傳》告訴我們這個教會怎樣去啟發我們.
怎樣給我們靈感 怎樣去提醒我們.
今天怎樣去回應不同的挑戰.
今天我們看第六章 第一節開始.
第一個point就是.
你會看到這個教會是願意去正視問題和挑戰的.
並且是願意去配搭 去使教會成長.
六章第一節.
那些日子門徒增多 有說希臘話的猶太人向希伯來人發出怨言.
因為在日常的功級上忽略了他們的寡婦.
十二個使徒就叫齊眾門徒來.
我們撇下上帝的道 去管理飯食是不合宜的.
所以弟兄們 當從你們中間選出七個有好明星的.
滿有聖靈和智慧.
我們派他們管理這事.
至於我們 我們要專注在祈禱和傳道的事奉上.
這裡說到初期的教會 由五旬節 聖靈降臨之後.
然後一直發展 直到這個時候.
他們出現了一些問題 教會也有問題.
教會繼續發展 繼續向前 也會出現不同的問題.
我們不要覺得做上帝的工作 一定是每件事都很順利.
不會的 就算是走在順利的旨意裡面.
也不一定是所謂天色常藍.
也會有阻滯 或者不順利的地方 有些問題.
很多時候我們見到教會有問題的時候.
到底我們應該怎樣回應呢.
你知道在之前上文 說到教會受到當時耶路撒冷的權貴打壓.
然後抓了彼得和約翰 讓他們被囚禁了一段短的時間.

$^{41}$但是上帝的教會怎樣去回應呢.
他們求上帝給他們力量 讓他們迎接這個困難.
然後上帝就加力給他們.
這是一個有聖靈的教會應當有的回應.
上個星期我們說到第五章的時候.
就說到一對夫婦 記得什麼名字嗎.
Zephyr 記不記得.
Zephyr Vizicat.
就是阿娜莉亞和薩菲拉.
在奉獻上他們有假冒為善的隱瞞.
隱瞞上帝 或者欺瞞上帝的靈.
他們受審判.
然後教會怎樣去回應呢.
戰戰兢兢地回應.
上帝的聖靈來到的時候.
不一定每件事都很快樂.
要我們去到警醒.
知道上帝是公義 是聖潔的.
是不能夠亂來侵犯.
第五章到第六章的時候.
我們發現這裡有什麼問題呢.
就出現了另外一方面的問題.
問題的起源是什麼呢.
當時的猶太人.
你知道他們亡國是什麼時候嗎.
是主前586年.
猶大國亡國.
然後眾志派所有的猶太人流散到各地.
世界各地不同的地方.
幾百年他們就住在海外.
有很多的猶太人住在海外.
有小部分的留在本土.
這個所謂猶太人流散的問題.
Diaspora的問題.
有聽過這個字嗎.
Jewish Diaspora.
是一個流散的情況.
猶太人在海外在異地散居.
然後很多時候他們每年有幾次大節期的時候.
他們回到耶路撒冷去朝聖.

$^{81}$五旬節 駐彭節等等.
他們回到朝聖的時候.
在耶路撒冷就非常擠迫.
很多時候這個Diaspora.
就是今天我們很多時候會提到的Diaspora Mission.
就是今天的人流散的情況.
Displacement的情況是非常之高.
所以今時到傳福音講宣教的時候.
很多時候不只是在講本土.
是在講用不同的方法不同的途徑.
向一些流散的人去傳福音去關顧他們.
你想一想這個確實是今天很多人的狀態.
所以我們以前叫做僑民.
今天叫做移民或者移居者.
居無定所的.
其實我們華人都是這樣.
過去幾十年有水的地方有華人.
有水的地方有華人.
有華人的地方就有餐館.
這是我們華人的情況.
這是Chinese Diaspora.
Jewish也有他的Diaspora.
是四海為家的.
今天這裡講到這個問題.
其中一個源自的原因就是這個.
巴勒斯坦本地的猶太人.
當然他們是講希伯來話或者阿拉文.
散居在不同的地方.
他們自然一代傳一代傳一代.
他們就不再講他們以往的那個所謂的母語.
他們就在講希臘話.
因為當時希臘的語言是國際語言.
所以他們每逢回到耶路撒冷的時候.
就變成了起碼有兩堆人在這裡.
一堆是本土講希伯來話或者阿拉文.
一堆是回去然後他們是講外語的.
但是那個的名份上面身份上面.
其實大家差不多是一樣的.
甚至有很多一些老了的僑民.
這些Diaspora.

$^{121}$這些海外的猶太人.
他們就會因為不同的原因.
宗教信仰的原因.
他們會回到耶路撒冷定居.
所以他們就想回到自己本土.
他們是華人.
老了之後他們去哪裡.
回鄉下.
回鄉下買一間屋.
然後可以踢踢腳.
有些是這樣的.
猶太人也是.
所以他們當時生活的背景情況就是這樣.
有本土的猶太人.
當然本土的猶太人一直在這裡小小的住.
在以色列地.
在加拿大地.
在巴勒斯坦地.
當然社交上的關係.
是非常緊密的.
幾代同堂.
親戚朋友.
全部都在這裡.
能夠彼此照顧.
但是海外回去的.
很多時候就是會無依無靠.
特別如果你回到本土.
你是作為一個僑民回到去的時候.
移民回到去的時候.
如果你那個配偶.
因為年老.
他去世的時候就剩下你一個.
所以很多時候有寡婦的情況出現.
居寡的情況出現.
當然聖經裡面也有說.
一些寡婦是比較富裕的.
為什麼.
因為她老公是CEO.
她老公死了.
她可以遺產給他.

$^{161}$所以她就比較有些.
袋裡面有些東西.
有些是這樣的.
但有很多不是這樣的.
笑得這麼開心.
當然是你們是那些.
有很多寡婦是無依無靠的.
所以在聖經裡面經常告訴我們.
我們需要去做什麼.
教會特別.
在當時的教會.
因為沒有社會福利.
沒有醫療等等.
這些社會福利的措施.
所以很多時候.
在當時的社群.
就需要去照顧這些無依無靠的寡婦.
起碼基本上要去照顧她.
所以雅各書在下面說什麼.
一章二十七節.
他說在上帝父神的面前.
那些清潔沒有玷污的敬虔.
Guardianess.
就是要看顧孤兒和寡婦.
這個所謂Guardianess.
你有的敬虔.
是需要在生活上表達出來的.
所以他們特別是一些有需要的寡婦.
沒有人理會.
無依無靠的.
他們就向教會申訴.
正如我所說.
當時政府沒有做這些事.
所以很多時候是聖殿或者教會.
所以在第六章就出現了這樣的問題.
之前我們說過.
教會是面對一些外在的打壓.
這裡面對的問題就是一些內部的going pain.
就是你成長的時候.
你長大的時候.

$^{201}$一定會有一些磨練或者痛苦.
那些叫做going pain.
就像你有輛車.
你走久了.
機件之間就會有一些磨擦.
教會其實也有一點像這樣的.
教會我們繼續去到成長的時候.
當然會有一些磨擦.
有誰在我們教會UV.
是沒有和別人磨擦過的.
沒有吧.
有誰起碼和十個人磨擦過的.
肯定會的.
因為車機件一路運行的時候.
就會有一些不同的零件.
齒輪就會發生熱力.
其實這個很正常的.
這個很正常.
你不要以為來教會.
大家是好好相遇等等.
我們每個人都是有罪的人.
但是在發熱的時候.
或者正常的操作裡面有磨擦的時候.
就需要做什麼.
就像你的車.
你要加啟油.
要加冷卻劑.
要檢查到底夠不夠.
大家需要在外在因點裡面.
能夠將磨擦的熱度減輕.
磨擦是必然的.
反而有時候是好事.
因為反映什麼.
因為反映大家都想做好這件事.
大家想搞好家庭.
大家搞好教會.
所以是沒有問題的.
問題就是這個教會.
當時第一世紀的教會.
他們上下如何去應對.

$^{241}$這個才是最重要的問題.
在很多地方.
一個團體.
一個組織.
一個公司.
甚至教會都不例外.
很多時候遇見問題的時候.
就會視而不見.
或者不嘗試去解決問題.
只要解決提出問題的人.
這個教會是怎樣呢.
這裡就是路加所寫的教會.
所刻畫的教會.
是一個肯正視問題.
和想辦法大家彼此同心的去解決這個問題.
所以我說史蘭傳裡面.
教會是一個很重要的角色.
就是他們同心合意的去應戰.
這一次由第一章到現在第六章.
大概都應該是第一次.
第二次要去到第十五章.
耶路撒冷的會議.
就是外邦人和猶太人的律法.
之間有些爭議.
怎樣去正視這個問題.
這個教會就是說.
他們肯去acknowledge the issue.
不一定馬上想辦法.
或者不一定馬上有.
每樣東西都迎刃而解.
不會視而不見.
這個就是第一世紀的教會.
另外他們要去配搭.
在這裡說日常的供給是甚麼意思.
以前的翻譯叫做serving table.
一想起的時候.
甚麼叫serving table.
就是站台.
有沒有做過站台.
我們華人來到美國的時候.

$^{281}$我們經常都要做餐館工作.
有誰做過站台.
這麼少.
你們現在是很多遺產.
起碼我也做過站台.
busboy 洗碗等等.
都捱過一下很辛苦.
事實上他們說.
日常的供給.
管理飯食.
其實是tables的意思.
就是吃東西.
一些基本的工作.
這樣聽的時候說站台.
是一些次要的工作.
我們讀到使徒的回應.
我們有甚麼感覺.
使徒說我們祈禱傳道.
這些我們不做.
很多時候我們都可以有這種藉口.
為甚麼是我洗碗.
因為我應該專責做甚麼.
閉上眼祈禱.
和打開口講道.
所以其他人做.
其實不是這樣的.
你們看清楚這裡說甚麼.
不是說serving tables.
和講道祈禱.
是對立的.
不是這樣的意思.
因為你看到使徒說.
祈禱傳道的事奉.
其實是同一個字.
日常的供給.
管理飯食.
和祈禱傳道的事奉.
是同一個字.
那個字你看到是希臘文.
就是我們今天英文說的deacon.

$^{321}$這個字.
就是出自這裡.
執事是甚麼意思.
就是當時這個字所引申出來的.
所以說到供應.
說到祈禱和講道.
其實都是神聖的工作.
都是需要甚麼人做.
都是需要上帝揀選的.
和有上帝生命的人.
有充滿聖靈的人.
你知道神國裡面.
沒有任何一件事.
是高於其他事的.
沒有聖俗之分的.
你祈禱傳道.
你就是spiritual一點的.
然後你管理飯食的.
這些就是worldly的東西.
沒有這樣分的.
你從事任何的職事.
照顧家庭也好.
管理經濟.
管理錢或者教育等等.
每一個的職事.
都是沒有.
本質上沒有分上下的.
在教會裡面的職事.
其實都是一樣.
我們本質上是平等.
不過我們的功能.
我們的角色是不同的.
所以需要是配搭的.
要大家cooperate.
你今天來教會.
我想問你.
如果沒有咖啡沒有茶.
沒有早餐.
你來不來.
你可能從來沒有想過這個問題.

$^{361}$可能你今天想起的時候.
你發現如果沒有這些東西.
可能我不來.
當然來了.
可能沒有想過.
但我們想說.
是需要配套的.
只是來聽你講道.
我也不來.
會不會.
你當然不會了.
你這麼熟悉.
不是每個人都這樣想.
如果沒有音響.
沒有好音樂.
今天有沒有.
還可以.
其實不錯.
如果沒有.
你來不來.
教會是要講配套的.
所以沒有姓族之分.
沒有上下之分.
我兒子週末回來.
星期五回來.
他第一個訊息告訴我.
他說車的泵把掉了下來.
搞什麼鬼.
可能新手開車.
他們不太熟悉.
之前我女兒也發生過同一件事.
她開車泊車.
為什麼要走到這麼前.
有風不要洗地.
前面有個泊.
她這樣踩上去再退後.
整個泵把掉了下來.
泵把掉下來.
開不了車.
又要我走去搞.

$^{401}$我看到車子.
好像撞了泵的假牙.
我問怎麼開.
馬上又要花一點時間.
又要和他弄.
我想講的.
其實子女真的挺有恩慈.
我的子女挺有恩慈.
有增加我壓力的恩慈.
每個人都是這樣.
有些恩慈強一些.
有些弱一些.
人生真是煩惱.
我想說你的泵把.
對你的車的運作.
好像沒什麼影響.
其實很有影響.
假牙掉下來怎麼走.
所以要弄.
教會成長會去到樽頸位.
那個bottleneck.
有個bottleneck在那裡.
地方可能不夠人手可能不夠.
在我們的財務上不夠.
奉獻等等.
我們不如呼籲弟兄姊妹.
因為我們需要配套才有突破.
這裡就是這樣.
執事要被選出來.
有生命力.
有恩慈.
有實際辦事的能力.
這些工作既是神聖.
也是實際的.
教會要去跨越.
要去成長的話.
我們實在需要.
有份.
有些人的角色可能不同.
但我們每一個.

$^{441}$都有他的功能.
坐在這裡也有功能.
六章第五節.
他們怎樣去回應.
他說.
試圖這樣去說.
我們應該要配搭.
使得全個會眾都喜悅.
他們就揀選了.
幾多個人.
1,2,3,4,5,6,7.
第一個是司提反.
最後一個是尼哥拉.
叫他們站在使徒面前.
禱告之後他們就為他按手.
即是按納他們.
上帝的道就興旺.
在耶路撒冷門徒的數目就變多.
並且有許多.
祭司.
都信從.
這七個名字有什麼共通點.
七個都是希臘的名字.
希臘的名字.
西門.
就是希伯來的名字.
斯蒂凡.
是一個希臘的名字.
Greek name.
為什麼選這七個.
都是希臘的名字.
特別最後那個.
是皈依猶太教的外邦人.
叫做.
尼哥拉斯.
為什麼選這七個.
因為針對當時有什麼問題.
是那些僑居者.
回來的他們操.
什麼語言.

$^{481}$希臘語.
就是針對這班會眾的需要.
使他們能夠.
well taken care of.
right.
我們人人都有盲點.
領袖都是.
所以這裡給我們看到.
打開新的一頁這個教會.
是多元的.
全體的教會.
才能夠健全.
對嗎.
就像我們教會.
我們有男同屋.
有女同屋.
我們的所謂call team.
就是我們的deacon.
都是有男有女的.
其實我很想強調就是需要有男有女.
對嗎.
因為女人看到的.
男人看不到.
我們有很多時候有盲點.
需要大家彼此的配搭.
這裡試圖有自己的角色.
禱告.
港獨.
都是需要有他們的功能需要專注.
其他人我們需要去配搭.
有沒有人和你說過.
預備一篇道.
今天我這篇道.
平時港獨.
我們港獨比較短.
半個小時25分鐘.
你知道要多久.
有沒有人和你說過要多久.
要多久.
你講過.

$^{521}$大概起碼20個小時.
起碼.
我都算有少少經驗.
我有時都要20多30個小時.
有時就5個小時.
我不講.
不告訴你甚麼時候.
但需要.
有時他們說.
台上一分鐘.
台下一個小時.
30分鐘就30個小時.
真的.
我們都需要專注.
當然不是說我們壟斷.
事實上這裡七位的所謂執事.
屎堤犯.
等等.
他們跟著做甚麼.
港獨.
屎堤犯港獨.
他死了.
都是很有能力的人.
但每一個人都有他的功能.
六章第七節.
神的道在興旺.
門徒的數目又加增.
也有很多祭司.
順從了信仰.
這個很特別.
祭司是夜遊神行.
打工的.
居然有部份的祭司開始.
信主.
這些是最不預期.
不會信耶穌的人.
他們信主.
教會繼續發展下去.
你看看下面的篇幅.
六章之後.

$^{561}$有撒瑪利亞人.
有外邦人.
向他們傳福音然後他們信主.
所以突破了樽頸位.
原因在哪裡.
就是他們願意去acknowledge.
那個problem.
然後他們願意去聽.
願意去祈禱解決.
要大家一齊去同心.
如果不是.
教會是突破不了的.
聽到沒有.
我們大家必須要同心.
必須要學習.
彼此去配搭.
如果不是這個家這個教會.
是搞不成的.
我們必須要大家.
每人做我們自己那一份.
這個同心的教會.
的刻畫.
你看《少林寺行傳》.
當教會成為一個角色.
一個character.
不只是彼得,雅各,約翰等等.
不只是保羅.
不只是一個個人.
是一個集體的個體.
就是church 教會.
同心合二由第一章到現在.
其實你不斷有.
這樣的snapshots.
即是一些.
這樣的圖畫.
第一章他們同心合二的.
120人去行節.
去禱告.
去等候聖靈的降臨.
二章五秦節.

$^{601}$來到了他們全部聚集在一起.
二章尾.
信的人都聚在一處.
同心合二.
接著我們看第六章.
講到公認的問題.
耶路撒冷的教會同心合二.
十三章.
講到安提安的教會.
他們一齊來禱告.
派保羅和巴拿巴出去宣教.
十五章我提過了.
耶路撒冷的神學會議.
討論到律法的問題.
大家一齊解決.
二十一章全體的伊弗所教會.
送別保羅.
然後到最後一章二十八章.
羅馬的教會.
去迎接保羅.
這些只是一些片段.
這個教會是不斷的.
同心來向前.
有聖靈的工作原來有這樣的可能.
有這樣的可能.
去告訴我們原來我們祈禱.
我們願意的時候.
是可以做得到的.
上帝的工作可以在我們中間.
是這樣發生.
我們講過我們的連題是跨元.
第二點講到.
不單止是傳福音.
第二點我們講到.
要和其他弟兄姊妹同行.
不是一個人的問題.
而是整個的教會整個的社群.
在這裡我想提議大家.
怎樣才能夠使我們教會能夠突破我們的瓶頸位.
能夠合一.

$^{641}$第一.
我希望邀請你.
你這個禮拜為一個弟兄或姊妹祈禱.
好嗎.
不要只為你自己祈禱.
為了其他人.
你隨便選一個.
選一個帥哥或是選一個美女.
這個禮拜為他祈禱.
有需要的你知道的.
你為他祈禱.
這個禮拜每一天.
第二.
你去參與服事.
無論你做哪一件.
剛才我在後台.
我在看他們.
因為我們有些的轉換.
有些的過渡.
都有很多事務要做.
服事教會.
好像那些萬能膠.
萬能膠.
是大家能夠黏在一起.
你要樂意去參與.
我們教會才可以成長.
第三.
這個很重要.
不要家食.
什麼叫家食.
有時茶餘飯後.
不是不提人的名字.
家食不是這個意思.
家食是什麼呢.
英文叫做character assassination.
一個人的名份.
你刺殺他.
沒有他在defense的時候.
你去說這樣說那樣.
然後奚落別人.

$^{681}$或者說別人的是非等等.
在嘴唇上.
我們學一下.
去shut up.
是不是.
我們學做什麼.
去看做和平.
看大家可以有和諧的關係.
學一下.
人都這麼大了.
學一下如何化解誤會.
有時你所知的事.
或者你所看的事.
未必一定對.
贊不贊成.
不一定對.
你可以學一下.
別人為什麼會這樣做.
學一下修補關係.
如果一個愛的群體.
教會是一個這樣的群體.
我們要學什麼.
不是一切就能砌好.
我們要學一下.
去寬恕.
是不是.
如果大家有什麼誤解的時候.
不要等在裡面.
令到你很多苦讀.
你要拒絕這件事.
第四.
要拒絕小圈子.
我們學一下去跨越.
和其他人去連結.
主動一點.
會前會後.
找一些不相熟的人.
和他坐下.
和他聊聊天.
你知道一個群體是怎樣來的.

$^{721}$對話.
conversation.
大家聊到天才會有群體.
才會有community.
如果聊不到天.
怎麼可能會有一個社群.
第五.
實際一點.
看看怎樣去.
支持教會.
我們不需要每件事都在同一頁.
It's ok.
有些細節上.
我們不一定.
每個人都看齊.
但合一.
是選擇.
It's a choice.
I choose to.
I made a decision.
to do this.
不是自然而然.
I choose to do this.
你的選擇.
禱告裡面的選擇.
你的行動.
你的取態上.
你去選擇.
為我們的差異感恩.
有沒有試過.
我們實在有很多差異.
和第二堂的語言又不合.
文化又不合.
對不對.
但我們試試在各個差異裡面.
我們學習一下.
看見上帝在我們中間做了甚麼.
對不對.
是不同.
是多元的.

$^{761}$這個才是真正的教會.
求主給我們恩典.
1 2 3 4 5 6.
你記得了嗎.
這個星期.
選好了沒有.
選好誰弟兄姊妹.
為他禱告.
如果你選好了.
我們等下完崇拜的時候.
你跟他說.
我選了你.
我這個星期為你祈禱.
有甚麼為你祈禱.
你不要選我.
我知道很多人為我祈禱.
不要選我.
選一些可能你不相熟的.
為你禱告.
我們等下再祈禱.
主愛多謝你.
因為你呼召我們能夠成為教會.
成為你自己的子民.
我們能夠成為弟兄姊妹.
並且我們不是千變一的群體.
主愛求你憐憫我們.
讓我們教會在你自己的愛裡面.
在你的恩典裡面.
在聖靈的工作裡面.
我們學習怎樣去.
一切的去解決問題.
學習怎樣去締造平安和諧.
主幫助我們這個群體能夠在愛裡面成長.
以致我們生命能夠突破.
我們教會整體的生命可以突破.
我們永遠要神聖.
讓我們去福音.
不只是靠口講.
我們走在一起的時候.
就是福音.

$^{801}$走在一起我們就是去永耀神明.
去告訴別人.
上帝.
你是真的.
你是真實的.
是我們真實的.
帶領我們這個禮拜.
我們去在這個教會.
在我們身邊的關係裡面.
我們能夠表達出.
體會到上帝的愛和恩典.
祈禱付出美感.
好,願上帝祝福各位.
\newpage



\section{使徒行傳 4:31-5:16}
\label{sec:xrB_pn8UZeQ}
\textbf{聖靈二次降臨 (使徒行傳4\_31-5\_16) - 陳崇基牧師}
\newline
\newline
連結: \href{https://youtube.com/watch?v=xrB-pn8UZeQ}{\texttt{https://youtube.com/watch?v=xrB-pn8UZeQ}} ~~~~ 語音日期: 2025-08-30
\newline
\newline
\hyperref[sec:GgAMYWfpmPk]{\small{< < < PREV SERMON < < <}}
~
\hyperref[sec:index_chronic]{\small{[返順時目]}}
~
\hyperref[sec:index_scriptual]{\small{[返順卷目]}}
~
\hyperref[sec:EMLJsLtjt_w]{\small{> > > NEXT SERMON > > >}}
\newline
\newline
使徒行傳 4:31-5:16
\newline
\begin{longtable}{cl}
\hline
\hline
章節 & 經文 (和合本修訂版)\\
\hline
4:31 & \begin{tabularx}{0.7\textwidth}{X} 他們禱告完了，聚會的地方震動；他們都被聖靈充滿，放膽傳講神的道。 \end{tabularx} \\ \\ \relax
4:32 & \begin{tabularx}{0.7\textwidth}{X} 許多信徒都一心一意，沒有一人說他的任何東西是自己的，都是大家公用。 \end{tabularx} \\ \\ \relax
4:33 & \begin{tabularx}{0.7\textwidth}{X} 使徒以大能見證主耶穌復活；眾人也都蒙了大恩。 \end{tabularx} \\ \\ \relax
4:34 & \begin{tabularx}{0.7\textwidth}{X} 他們當中沒有一個缺乏的，因為凡有田產房屋的都賣了，把所賣的錢拿來， \end{tabularx} \\ \\ \relax
4:35 & \begin{tabularx}{0.7\textwidth}{X} 放在使徒腳前，照每人所需要的，分給每人。 \end{tabularx} \\ \\ \relax
4:36 & \begin{tabularx}{0.7\textwidth}{X} 有一個利未人，名叫約瑟，使徒稱他為巴拿巴（巴拿巴翻出來就是安慰之子），生在塞浦路斯。 \end{tabularx} \\ \\ \relax
4:37 & \begin{tabularx}{0.7\textwidth}{X} 他有田地，也賣了，把錢拿來，放在使徒腳前。 \end{tabularx} \\ \\ \relax
5:1 & \begin{tabularx}{0.7\textwidth}{X} 有一個人，名叫亞拿尼亞，同他的妻子撒非喇，賣了田產， \end{tabularx} \\ \\ \relax
5:2 & \begin{tabularx}{0.7\textwidth}{X} 把錢私自留下一部分，他的妻子也知道，其餘的部分拿來放在使徒腳前。 \end{tabularx} \\ \\ \relax
5:3 & \begin{tabularx}{0.7\textwidth}{X} 彼得說：「亞拿尼亞！為甚麼撒但充滿了你的心，使你欺騙聖靈，把賣田地的錢私自留下一部分呢？ \end{tabularx} \\ \\ \relax
5:4 & \begin{tabularx}{0.7\textwidth}{X} 田地還沒有賣，不是你自己的嗎？既賣了，錢不是你作主嗎？你怎麼心裡會想這樣做呢？你不是欺騙人，是欺騙神！」 \end{tabularx} \\ \\ \relax
5:5 & \begin{tabularx}{0.7\textwidth}{X} 亞拿尼亞一聽見這些話，就仆倒，斷了氣；所有聽見的人都非常懼怕。 \end{tabularx} \\ \\ \relax
5:6 & \begin{tabularx}{0.7\textwidth}{X} 有些年輕人起來，把他裹好，抬出去埋葬了。 \end{tabularx} \\ \\ \relax
5:7 & \begin{tabularx}{0.7\textwidth}{X} 約過了三小時，他的妻子進來，還不知道所發生的事。 \end{tabularx} \\ \\ \relax
5:8 & \begin{tabularx}{0.7\textwidth}{X} 彼得對她說：「你告訴我，你們賣田地的錢就是這些嗎？」她說：「就是這些。」 \end{tabularx} \\ \\ \relax
5:9 & \begin{tabularx}{0.7\textwidth}{X} 彼得對她說：「你們為甚麼同謀來試探主的靈呢？你看，埋葬你丈夫之人的腳已到門口，他們也要把你抬出去。」 \end{tabularx} \\ \\ \relax
5:10 & \begin{tabularx}{0.7\textwidth}{X} 她立刻仆倒在彼得腳前，斷了氣。那些年輕人進來，見她已經死了，就把她抬出去，埋在她丈夫旁邊。 \end{tabularx} \\ \\ \relax
5:11 & \begin{tabularx}{0.7\textwidth}{X} 全教會和所有聽見這些事的人都非常懼怕。 \end{tabularx} \\ \\ \relax
5:12 & \begin{tabularx}{0.7\textwidth}{X} 主藉使徒的手在民間行了許多神蹟奇事；他們都同心合意地聚集在所羅門的廊下。 \end{tabularx} \\ \\ \relax
5:13 & \begin{tabularx}{0.7\textwidth}{X} 其餘的人沒有一個敢接近他們，百姓卻尊重他們。 \end{tabularx} \\ \\ \relax
5:14 & \begin{tabularx}{0.7\textwidth}{X} 信主的人越發增添，連男帶女都很多， \end{tabularx} \\ \\ \relax
5:15 & \begin{tabularx}{0.7\textwidth}{X} 甚至有人將病人抬到街上，放在床上或褥子上，好讓彼得走過來的時候，或者影子投在一些人身上。 \end{tabularx} \\ \\ \relax
5:16 & \begin{tabularx}{0.7\textwidth}{X} 還有許多人帶著病人和被污靈纏磨的，從耶路撒冷四圍的城鎮來，他們全都得了醫治。 \end{tabularx} \\ \\
[1ex]
\hline
\hline
\end{longtable}
$^{1}$我們回到上帝的說話.
各位大姐妹 我們回到上帝的說話.
我們去到士蘭傳的第四章尾.
我們有一個系列.
一路由年頭到現在是在說跨越.
在這一年裡我們希望記得這個字.
Crossover 或者跨越.
我們希望能夠挑戰各位.
如果你們信耶穌的話.
我們信主的話.
我們基督徒的生命應該要繼續Crossover.
應該繼續去挑戰我們自己.
離開我們自己安殊的地方.
所以我們不斷地在說士蘭傳.
從第一章開始到現在已經到第四章.
不如我們看看這段聖經.
有兩個點我想說.
就是在這裡 當時初期的教會是一個彼此分享的團體.
並且是一個彼此願意去問責.
在上帝面前願意被問責.
Keep accountable.
以及在弟兄姊妹在整個社群裡面是願意Keep accountable的一個群體.
如果你有聖經的話.
跟我一起打開去士蘭傳第四章尾這個地方.
這裡這樣去描述當時的教會.
可以算是聖靈的第二次降臨.
你記得第一次降臨是什麼時候嗎?.
在五重節的時候.
就是第二章的時候.
這個算是第二次了.
因為四章三十一節這裡說什麼呢?.
他說這個教會他們受迫害的時候.
有很多壓力的時候.
他們就禱告求上帝給他們力量.
然後聚會的地方震動.
他們都被聖靈充滿放膽傳講上帝的道.
聖靈又再次充滿他們.
許多信徒都一心一意.
沒有一人說他的任何東西是自己的.
大家公容.

$^{41}$試圖以大能見證主耶穌復活.
眾人也都蒙了大恩.
他們當中會有一個缺乏的.
因為凡有田產房屋的都敗了.
把所賣的錢拿來.
放在使徒的面前.
招每個人所需要的分級國人.
這裡第一點就是.
講到這個社群.
講到他們是彼此能夠去分享的.
你有沒有留意在《少林寺傳》裡面.
或者甚至整本《聖經》裡面.
上帝不是救一個人的.
不是救你一個個體.
另外救一個個體.
然後大家分開.
永遠都是組成一個新的子民.
一班新的人.
一個新的家庭.
一個新的社群.
一個新的community.
這個社群能夠成長.
能夠跨出他們自己的界限.
去到事奉其他的種族.
事奉其他的語言方言的人.
都不是一個人講到的事.
因為這裡第四章做一個總結的時候.
他說初期的教會是從撒亞開始.
並非一個人講到很厲害.
大家知道彼得講到第二章五重說的時候.
吸引了很多人來信主.
但是他總結的時候.
隻字不提.
沒有提到因為有一個講到的很厲害.
也沒有講到說.
有一班人有一班領袖他們很厲害.
而是整個群體一心一意的.
每個人.
來貢獻在這個群體裡面.
這個就是聖經信仰的特別地方.

$^{81}$聖經的信仰和其他宗教是很不同的.
其他宗教是一個人你自己搞定.
你figure out那個philosophy.
那個idea.
那就足夠了.
但聖經的信仰.
由始至終耶穌基督的信仰.
都不是一個人的事.
因為上帝是要組成一家人.
沒有人可以孤軍作戰.
沒有一個基督徒.
可以靠自己.
我們是需要其他人.
其他大英哲門.
你想像一下有一個營火.
燒得很紅.
裡面當然有很多柴.
裡面有很多木.
很多燃料.
如果將其中一條柴扔開.
放在一邊.
它還會不會燒.
還可以燒的.
但燒了多久.
燒了不久.
那個叫廢柴.
基督徒離開群體的時候.
就會變成這樣.
你對上帝的熱心.
你對上帝的信心.
慢慢就會消磨.
這是很自然的事.
為什麼.
因為聖經的信仰並非如此.
第四章描述什麼.
多了一件事.
他說他們當中沒有一個缺乏的.
這句話四章三十四節.
你找遍整本聖經.
只出現在一個地方.

$^{121}$就是《申命記》第十五章.
《申命記》第十五章說什麼.
上帝吩咐他的子民做什麼.
他說你們中間是會有窮人的.
你們中間是會有需要的人的.
在你們中間必沒有窮人.
為什麼.
因為你們需要.
在你們當中去貢獻.
你這樣貢獻的時候.
因為耶和華給你這個產業.
耶和華就會大大賜福給你.
所以他說你們弟兄中間.
若有一個貧窮的.
你不要硬心.
你不要伸手不幫你貧窮的弟兄.
你要慷慨的解囊.
給他們的時候不要心痛.
這裡說得很真實.
貢獻給別人.
去到給錢出來.
當然心痛.
他說你不要這樣.
因為耶和華你的上帝在這事上.
在你一切所有工作.
在你所做的一定會賜福給你.
一定會賜福給你.
上下文就在說.
以色文他們有一個律法.
這個律法就叫做安息年的律法.
叫做Sybarakal Year.
大家明白Sybarakal Year是什麼時候.
Sybarakal Year就是每逢第七年的時候.
就是一年是不同的.
那一年做什麼呢.
就是大家所謂安息.
然後去休耕.
然後有一樣東西很重要的.
就是《申命記》十五章這裡所說的.
就是你要豁免其他人的債項.

$^{161}$讓整個社會.
整個社群.
能夠所謂reset.
好像你用完一部電腦.
用了很久的時候.
你按個reset.
然後叫它shutdown再restart.
一樣.
上帝就說.
如果我的子民.
你們願意去守這個誡命的話.
上帝就會賜福給你們.
國中必定沒有窮人.
《白史行傳》就是去到描述.
耶路撒冷這個小的社群.
就用這樣的字眼.
為什麼呢.
為什麼在新約的時候.
耶路撒冷的教會.
會用一個舊約的.
《申命記》的方法.
來形容他們呢.
原因就很簡單.
就是identity的問題.
身份的問題.
就是新約的教會.
耶穌基督死而復活升天之後.
他們所組成的社群.
他們認為.
他們自己就是誰.
就是上帝的子民.
和舊約上帝的子民.
是連結在一起的.
同樣成為.
The chosen people of God.
是上帝的選民的一部分.
他這樣形容舊約.
也是這樣形容新約.
所以他們的名份.
他們的身份絕對是上帝的選民.

$^{201}$很多時候我們搞錯了.
我們看今天新聞的時候.
我們覺得誰是上帝的chosen people.
是猶太人血統的人.
對嗎.
猶太人血統.
在他們的生命裡流著猶太人的血.
那些人就是chosen people.
我們呢.
我們可能不是這麼chosen.
或者我們不是這麼像選民.
是他們比較像.
聖經從頭到尾舊約至新約.
就是告訴我們在基督裡面的.
那些就是神的選民.
我們就是神的選民.
你在基督的名下.
你信了耶穌基督的.
你就是屬於上帝所揀選的.
所鍾愛的那幫群體.
這就是史文傳為什麼描述當時的教會的時候.
他就說有猶太人.
並且將來會有其他種族的.
所謂外邦人加入去的.
這些都是絕絕對對神的選民.
所以你信耶穌的.
等下你回去照鏡.
你說什麼.
我是上帝的選民.
因為耶穌基督的緣故.
你有沒有去過.
你可以這樣想像.
你有沒有去過婚禮的現席.
人家擺酒.
你不是家屬.
你坐在那裡.
25號桌.
30號桌.
50號桌.
看那個婚宴多大.

$^{241}$你不會這麼傻的走進去.
然後坐在主家席.
現親認戚嗎.
但你可以想像.
你去到一個婚宴.
你坐在26號桌.
編碼就編到26號桌.
然後跟一些人不認識的一起坐.
很沒趣.
突然間開席的時候.
主人家走來跟你說.
他說你坐在主人家.
你坐在那裡幹嘛.
你是屬於我們主人家的.
你會有什麼反應.
當你自己人.
你會不會這麼面懵.
會上去.
但他真的當你是主人家.
禮金也給你.
有份婚.
去不去.
去吧.
這就是聖經裡面所說的.
新約的教會在末世裡面.
所有相信耶穌基督的人.
是跟誰的主家席.
阿布拉罕.
約瑟 雅各 摩西 大衛 所羅門.
等等等等 舊約裡面的列祖.
所有我們覺得他們的生命很吸引.
上帝揀選他們在地上.
有真正身份 有神的子民的身份.
新約的教會.
就是因為耶穌基督所做的.
我們就成為了天國裡面地上的那幫人.
所以在我們中間沒有窮人.
你看到他們的刻畫.
描述當時的教會.
We are the people of God.

$^{281}$No one else.
我們屬於基督的.
我們就能夠經歷上帝的恩典.
我們就能夠彼此的去分享.
就因為我們一幫選民.
我們能夠有上帝那麼多的祝福的時候.
我們就能夠彼此的去分享.
這就是第一世紀教會的所謂Selfie.
他們拍照的時候.
他們看到他們自己.
就是一幫上帝的子民.
所以他們中間會有窮人.
大概就是他們的生命裡面.
他們彼此分享裡面.
他們經歷到就是以前舊約裡面那種的美景.
他們現在作為選民.
他們同樣經歷到.
這就是為什麼我們需要建立我們的教會.
我們需要去做這件事.
就是奉獻.
昨天我們Call Team Retreat.
我們的核心小組.
我們那幫執事.
我們有一個退休會.
我們就在這裡.
我們靈修.
我們大家禱告.
我們為了教會禱告.
特別到我們年終的時候.
我們有一個教會的財務的更新.
之前我們都講過.
我們的預算裡面.
我們每年的預算.
我們每個月.
我們需要的起碼是三萬八千美金.
不是日子.
到六月的時候.
我們的奉獻支出.
我們算一算.
那個基準表.

$^{321}$就是我們欠了一萬九千美金.
我們欠一萬九千美金.
聽到這個數字的時候.
你覺得怎麼樣.
你覺得還可以.
還沒死嘛.
我們Call Team討論過有什麼策略.
我明白.
經濟好不好.
我明白經濟的物價是上漲的.
如果你是恆常奉獻的話.
我為你感恩.
十一奉獻是我們有的福分.
我知道物價高漲.
所以我們都希望.
其中一個策略就是.
我們希望我們的奉獻.
我們需要去再調整一下.
譬如我自己.
我會回去我自己銀行的奉獻.
我需要調高起碼20\%.
我才可以追得上所謂的inflation.
大家今時今日出街.
你都感覺到的.
吃這些東西要付這麼多錢.
是不是.
坐下來給提示.
40元50元.
你們去Taku.
是不是也是付這麼多.
大概都是這樣.
看你叫什麼.
我們物價上漲是現實.
同樣上帝指明了很多教會.
我們就需要去回應.
如果你還沒有開始有奉獻的操練.
開始有奉獻的習慣的話.
我為你祈禱.
我們在主裡面需要學習.
因為耶穌說什麼.

$^{361}$耶穌說一個人不能夠事奉兩個主.
要麼你事奉主.
要麼你事奉馬門金錢.
沒有兩個主可以選擇.
這是其中一個策略.
我們弟兄姊妹需要去活出.
我們上帝指明的身份.
另一個策略我們說過的.
討論了不少時間.
就是我們同工需要減薪.
當然我不希望做到這一步.
是不是.
但我想告訴你.
通常大家做事知道.
你公司如果要減薪的話.
就是減誰.
就是減那些小的.
將那些小的割下.
將他們裁員.
當然我們不希望做到這一步.
但我想告訴你.
我是走著說的.
如果我們真的去到這一步的話.
我們沒有辦法的話.
如果要減薪的話.
我想告訴你.
一定會先減我.
我就不會說那些小的.
我們將他們裁掉.
然後我們保住大的.
因為這是一個不同的身份.
這是上帝的指明.
耶穌說.
走前面的要走後面.
做大的要先做小的.
我們希望是走前面.
順帶一提.
芝加哥以外的成會的朋友.
他們久不久會有些奉獻.
這裡我真的非常感激.

$^{401}$但是這裡我想說.
我們教會.
自己教會應該自己撐.
對嗎.
網上的朋友.
但這只不過對我們來說.
是一個bonus.
但我們作為一個神的子民的身份.
應該要怎樣去做.
自己教會.
自己撐.
我自己也是這樣.
我會調節一下自己的生活風格.
我會省錢一點.
吃少一點.
無所謂.
以致有什麼目的.
以致我猶如.
我能夠祝福其他同工.
其他人.
在一個這樣的經濟狀況裡面.
我們需要做一些調節.
這是現實.
這不是一生鮮的東西.
舊鑰節新約所有神的子民.
都會經歷很多起起伏伏.
不同的帝國.
不同的時間.
不同的艱難的時間.
但申命記在這裡說什麼.
史行傳說什麼.
我們的身份就是我們可以慷慨.
因為我們體會到.
上帝對我們祝福和供應.
這個群體.
屬於耶穌的屬靈群體.
是會反常的.
什麼叫反常.
我們通常都是說.
我們有多餘的時候.

$^{441}$我們才去貢獻.
上帝給我多一點祝福.
那我就祝福其他人.
是不是.
上帝的子民的群體教會.
是反常的.
就是我們先去貢獻.
就會見到祝福.
因為耶和華在這世上.
和你一切工作所做的.
就會似福及利.
你看到是不是反常.
邏輯是反常.
因為我們不是一個人間的群體.
一個屬耶穌.
一個屬天的群體.
但我想告訴你.
我事奉了三十年.
永遠都兌現這個東西.
我今天站在這裡.
我仍然是做到事奉.
仍然可以繼續向前.
我就證明了上帝的祝福.
不是人所計算的.
所以這道使徒行傳所記載的.
所有的selfie.
就是說教會是一個怎樣的教會.
既然得到上帝的大恩.
就會彼此分享彼此的貢獻.
第二方面.
是一個彼此問責的群體.
然後這裡的聖經.
有個selfie之後.
他就把鏡頭zoom了下去.
兩個例子.
第一個例子就叫做巴拿巴.
他說當時有一個人叫做巴拿巴.
本身叫做約瑟.
希伯來名叫做約瑟.
他是利未人.

$^{481}$他本身在塞浦路斯.
Cyprus裡面搬來.
約瑟他有田地也賣了.
就把錢放在使徒的面前.
這裡沒有說巴拿巴是大地主.
他的田有多大.
可能五百呎.
你不要想像他是一個億萬富翁.
這裡所記載的焦點是什麼.
他是一個慷慨祝福人.
所以叫做安慰之子.
什麼叫安慰之子.
希伯來語的一種方式來形容一個人的品格.
什麼什麼之子.
就是說他的品格.
他的本性就是一個鼓勵安慰人的人.
他是一個encourager.
你看到巴拿巴從這裡開始.
直到史蘭傳之後的章節.
都是說他是一個怎樣的人.
他就是一個encourager.
他去提拔誰.
保羅.
沒有巴拿巴就不會有保羅.
是因為巴拿巴願意去提拔他.
給他機會.
保羅才會出來.
然後巴拿巴就會下台.
這是一個很好的生命.
一個很好的榜樣.
所以有一個自拍之後.
說他們是彼此分享的一個群體.
然後他們就說有兩個例子想告訴大家.
這個例子之後.
巴拿巴之後.
就有一個另外一個負面的例子.
這個負面的例子.
就是阿拉利亞和薩菲拉兩夫婦的故事.
看起來很長.
還有看下去.

$^{521}$這段經文怎樣看都好.
都是比較不舒服的.
我們試試看.
他說有一個人叫阿拉利亞.
他的妻子叫薩菲拉.
賣了田產.
薩菲拉.
這個名字很容易記.
如果你今天完結之後.
你不記得這個名字.
不記得這個故事的話.
你記得薩菲拉.
因為現在很流行一張簽證卡叫做什麼.
Sapphire.
Chase Sapphire.
就是這個字.
寶藍色的意思.
有兩夫婦他們賣了田產.
然後私下留了一部分.
他的妻子也知道.
就是兩個人都知道.
夫婦都知道.
其餘的部分.
他就放在市道的面前.
彼得就說.
阿拉利亞為什麼撒旦充滿你的心.
然後你欺騙聖靈.
賣田產.
然後私留一部分.
田產會有賣不是你的嗎.
既賣了錢就不屬於你了.
你不是奉獻給當時的教會.
你不是欺騙人.
而是欺騙上帝.
第五節.
阿拉利亞聽見.
他就跌倒然後斷氣.
所有聽見的人就非常懼怕.
有些年輕人就將他的屍體拿下.
第七節.

$^{561}$約到三小時之後.
他的妻子又來了.
他不知道之前發生了什麼事.
彼得就對她說.
你告訴我賣田產的錢.
就是這個嗎.
她就說.
是的就是這個.
彼得就跟她說.
為什麼你要同謀來試探主的聖靈.
你看埋葬你父親的人的腳已經去到門前了.
他們也要將你抬走.
他就立刻在彼得面前跌倒.
斷了氣.
年輕人又將他的屍體拿走了.
全教會和所聽見的這些人都非常懼怕.
怎樣看完聖經過去.
你覺得很奇怪.
你說是不是.
為什麼要記載這樣的故事.
我們需要問一個問題.
就是路加.
當他講到初期教會的時候.
為什麼會選擇這樣的故事.
他肯定有很多故事可以選擇.
為什麼要選擇一個欺騙的故事.
為什麼.
是不是這麼大的一件事.
奉獻上報假數.
通常我們看電影經文的時候都不是這麼爽.
但我想弄清楚.
經文著眼的是什麼.
首先第一件事著眼的.
不是說他們不得救.
然後他們就落地獄.
經文著眼的是什麼.
不是一些終極救恩的問題.
想說什麼.
他說聖靈既然光臨.
聖靈是聖潔的靈.

$^{601}$是公義的靈.
是真理的靈.
他是測透人心.
他知道你每一個的意念.
每一個的思想.
知道你每一個的penny.
阿娜莉亞薩菲拉.
可能覺得.
這些是內心深處的東西.
沒有人知道.
若要人不知.
上帝也不知.
欺騙聖靈.
其實他針對的是hypocrisy.
是假無為善.
You don't want to talk.
That is the issue.
行事為人和form是相反的.
不想被問責.
本質上他就是.
欺騙聖靈.
然後他威脅了.
當時教會的群體.
你想像一下.
如果教會的群體.
每一個都是這樣的時候.
這個教會的群體.
就不能有一個說服力的公眾見證.
教會永遠都不是收起來的.
教會是公開的.
是public witness.
你想清楚.
public witness是一個公開公共的見證.
我們永遠都是在鏡頭底下.
我們永遠都是在spotlight底下.
所有基督徒.
歷世代代的教會.
永遠都是.
前幾天.
網上瘋傳.

$^{641}$一對情侶.
在co-play演唱會裡抱在一起.
其實他們是偷情.
然後鏡頭就拍到他們.
網上瘋傳.
除非你住火星.
你肯定會知道有這件事.
肯定會知道有這件事.
起碼看到有這件事.
有時候看這些都很有趣.
最有趣的是看comment.
有一個comment這樣說.
如果他們不避鏡頭的話.
可能事情就不會搞到這麼大.
大大方方就不行.
偷情就偷情.
大大方方.
如果不是縮的話.
焦點怎麼會放在你身上呢.
另外一個comment更有趣.
他說舔都舔了.
怕什麼.
你都公開去了.
舔都舔了.
無論如何我想說一件事.
我們永遠都是在鏡頭底下.
我們真的.
不只是那對夫婦.
不只是撒菲拉和阿娜莉亞.
教會基督徒總是在眾人的目光底下.
你不能說你不要跟我.
因為我不是一個好的基督徒.
你不可以這樣用藉口卸駁.
你的身份.
既然你是認信耶穌基督的.
就在目光底下.
就等於今天.
如果我用鏡頭一直對著你.
整個聚會都對著你.
你是上網的.

$^{681}$你OK嗎.
你現在不OK.
下次我不管了.
射著你就好了.
射著你OK.
你就是要上網.
我就不上網.
我躲在下面.
當然這是聚會的setting.
但我想告訴你.
生命的setting不是這樣的.
不要看他們了.
我們的生命.
永遠都在目光底下.
在鏡頭底下.
你不要管有沒有人偷看你.
你不要管有沒有人偷拍你.
永遠都是我們要準備好上鏡.
對不對.
我們不只是因為要上台.
我們要準備義勇要搞得好.
然後我們的衣衫要搞得好.
我們永遠都要準備好.
生命在鏡頭底下.
這是教會應該有的生命.
每一個基督徒有的生命.
你留意這個.
所謂其中一個selfie.
初期教會的selfie.
他們很慷慨地分享.
但這個有一點奇怪.
通常我們慷慨分享的時候.
我們就說有眾籌的時候.
fund me.
有campaign.
我們expansion有campaign.
我們慷慨地去做一些事.
support哪個短孫.
有什麼孫教的項目.
你告訴我意象是什麼.

$^{721}$你有同感的時候.
你就會去fund raise.
你就會慷慨.
當然這是當然的.
在新約裡面.
在《少行傳》或其他書信裡面.
其實很多都是和fund raising有關.
普羅寫很多的書信.
其實都是和fund raising有關.
因為他想繼續有孫教的工作.
繼續將福音在羅馬帝國不同的地方傳遞.
寫肥納比書.
寫這麼多前後書.
希望能夠幫助教會.
能夠做更多更多的工作.
這是當然的.
但在《少行傳》有一個奇怪的地方.
就是沒有眾籌.
沒有fund raising.
沒有vision.
原因在哪裡.
原因就是.
we know who we are.
我們知道我們是哪一個.
我們是屬於耶穌的群體.
屬於主的群體.
彼此分享.
這是我們身份要做的事.
彼此問責.
一個公共的見證的群體.
全部都是因為.
我們已經得到上帝的恩典和祝福.
這就是我們的身份.
明白我們自己的身份.
我們只不過是做一些和我們身份相關的事.
we walk the talk.
live.
24/7.
live.
求主幫助我們.

$^{761}$live.
直播中.
你直播中.
我也是直播中.
我們24/7.
都在鏡頭底下.
我們是live.
我們和心得一起祈禱.
上主求你幫助我們.
用我們的生命.
在你的恩典和祝福裡面.
我們被直播.
用我們的生命成為公共的.
有力的.
見證.
幫助我們今天在這裡.
在這裡.
見證.
幫助我們教會.
我們用我們的資源.
我們的恩慈.
我們的才幹.
we go live.
為你而活.
讓我們這樣生活的時候.
我們群體我們個別的生命.
能夠成為一個更有力的見證.
去告訴別人耶穌基督你是主.
耶穌基督你是我們的君王.
值得我們敬拜.
用我們敬拜.
我們的事奉.
我們的心思意念.
全歸你事奉你.
祈禱洪主名禱.
願主祝福各位.
go live.
以便接收更多遊戲資訊,攻略及新影片等通知©2017 遊戲修理局局長.
\newpage



\section{詩篇 19:1-14}
\label{sec:EMLJsLtjt_w}
\textbf{The God Who Speaks (Ps 19) - Rev. Mark Chan 對你說話的神 (詩篇19\_1-14) - 陳崇基牧師}
\newline
\newline
連結: \href{https://youtube.com/watch?v=EMLJsLtjt_w}{\texttt{https://youtube.com/watch?v=EMLJsLtjt\_w}} ~~~~ 語音日期: 2025-09-01
\newline
\newline
\hyperref[sec:xrB_pn8UZeQ]{\small{< < < PREV SERMON < < <}}
~
\hyperref[sec:index_chronic]{\small{[返順時目]}}
~
\hyperref[sec:index_scriptual]{\small{[返順卷目]}}
~
\hyperref[sec:8R1ydESgnbc]{\small{> > > NEXT SERMON > > >}}
\newline
\newline
詩篇 19:1-14
\newline
\begin{longtable}{cl}
\hline
\hline
章節 & 經文 (和合本修訂版)\\
\hline
19:1 & \begin{tabularx}{0.7\textwidth}{X} 諸天述說神的榮耀，穹蒼傳揚他手的作為。 \end{tabularx} \\ \\ \relax
19:2 & \begin{tabularx}{0.7\textwidth}{X} 這日到那日發出言語，這夜到那夜傳出知識。 \end{tabularx} \\ \\ \relax
19:3 & \begin{tabularx}{0.7\textwidth}{X} 無言無語，也無聲音可聽。 \end{tabularx} \\ \\ \relax
19:4 & \begin{tabularx}{0.7\textwidth}{X} 它們的聲浪傳遍天下，它們的言語傳到地極。神在其中為太陽安設帳幕， \end{tabularx} \\ \\ \relax
19:5 & \begin{tabularx}{0.7\textwidth}{X} 太陽如同新郎步出洞房，又如勇士歡然奔路。 \end{tabularx} \\ \\ \relax
19:6 & \begin{tabularx}{0.7\textwidth}{X} 它從天這邊出來，繞到天那邊，沒有一物可隱藏得不到它的熱氣。 \end{tabularx} \\ \\ \relax
19:7 & \begin{tabularx}{0.7\textwidth}{X} 耶和華的律法全備，使人甦醒；耶和華的法度確定，使愚蒙人有智慧。 \end{tabularx} \\ \\ \relax
19:8 & \begin{tabularx}{0.7\textwidth}{X} 耶和華的訓詞正直，使人心快活；耶和華的命令清潔，使人眼目明亮。 \end{tabularx} \\ \\ \relax
19:9 & \begin{tabularx}{0.7\textwidth}{X} 耶和華的典章真實，全然公義，敬畏耶和華是純潔的，存到永遠， \end{tabularx} \\ \\ \relax
19:10 & \begin{tabularx}{0.7\textwidth}{X} 比金子可羨慕，比極多的純金可羨慕；比蜜甘甜，比蜂房下滴的蜜甘甜。 \end{tabularx} \\ \\ \relax
19:11 & \begin{tabularx}{0.7\textwidth}{X} 因此你的僕人受警戒，遵守這些有極大的賞賜。 \end{tabularx} \\ \\ \relax
19:12 & \begin{tabularx}{0.7\textwidth}{X} 誰能察覺自己的錯失呢？求你赦免我隱藏的過犯。 \end{tabularx} \\ \\ \relax
19:13 & \begin{tabularx}{0.7\textwidth}{X} 求你攔阻僕人不犯任意妄為的罪，不容這罪轄制我，我就完全，免犯大罪。 \end{tabularx} \\ \\ \relax
19:14 & \begin{tabularx}{0.7\textwidth}{X} 耶和華－我的磐石，我的救贖主啊，願我口中的言語，心裡的意念在你面前蒙悅納。 \end{tabularx} \\ \\
[1ex]
\hline
\hline
\end{longtable}
$^{1}$- Thank you, Pastor Mark..
- Hey, how's it going guys?.
I know it has been a long weekend.
and this is a Labor Day weekend.
and a lot of us are out of town..
And then for some of you,.
yesterday was a very, very long day, right?.
Well, one applause for all of you..
Sometimes you have to say thank you to yourself..
In any case, today we're going to talk about Psalm 19..
Now you remember last week we went to the Botanical Garden.
and we had a very good time, right?.
So I kept on thinking about this psalm.
and it says, let me read maybe the first verse..
"The heavens declare the glory of God.
and the skies proclaim the work of His hands.".
And then the creation is like a preacher..
It's like someone who's preaching a sermon to you.
and proclaiming the greatness of God.
to all people everywhere, all the time..
And like this psalm says,.
the heavens and the creation do not speak our language..
There's no voice..
Yet they speak a language.
that everyone of us can understand..
That's amazing..
Day after day, night after night,.
creation keeps on pointing us.
back to the glory of the creator..
And as you can see,.
these new images came out 2022 by James Webb Telescope..
And you know, 100 years ago,.
we, people thought that there was only one galaxy..
But now with this pictures,.
with this all kinds of new images,.
there were countless of galaxies in space, in deep space..
And some of the lights we see today.
has traveled like 13 billion years of time.
before reaching us now..
Now this tells us, like the psalm says,.

$^{41}$you know, how small we are.
and how great our God, how vast the universe is..
And the heavens stretch out beyond our imagination..
And yet, this is only what we see today in nature..
It's only a reflection of the majesty of our maker..
Now the psalm, psalm 19,.
highlights one of the factors, the sun..
Now you can ask the question, why the sun?.
Why focus on the sun?.
Because, you know, everyone see the sun..
You saw the sun today, right?.
You came in and you saw the sun and you woke up..
It's like a bridegroom coming out of his wedding canopy..
You know, it's like a look of triumph,.
like a champion running with joy and with energy..
That is the sun..
The sun rises every day..
It runs its courses across the sky with power..
And you can tell, even in Chicago,.
we don't have a lot of sun,.
but when the sun comes up, nothing escapes its heat..
And the same way, the psalm's telling us,.
no one should escape,.
no one can escape the weakness of creation..
And the glory of God is unmistakable..
And we acknowledge the existence of the savior, our maker..
And not only in the vastness of the space,.
but even in the smallest detail in creation,.
it tells us the existence of God..
Now, some years ago,.
a photographer tried to magnify the snowflakes.
and under the camera and took a lot of pictures..
And as you can see,.
every single snowflake look different, right?.
Everyone, every snowflake has a unique signature.
and they are all masterpiece..
And this is how amazing God's creation is..
And from the blazing sun, we all can see,.
to the delicate snowflake, right?.
Snowflake..

$^{81}$In the smallest detail,.
the creation keeps telling us God is glorious..
Now, this is what the theologians of the Bible tells us..
This is general revelations,.
because God left his fingerprint,.
his signature everywhere,.
every single inch of the creation..
And no one can say, "I never heard of the creation,.
never heard of the creator.".
Creation is always speaking..
And the only thing is whether you would listen or not..
And I remember some years ago,.
there was a debate going on.
between a British rabbi called Jonathan Sachs..
And he was debating the most famous atheist,.
Richard Dawkins, over online..
And this is one of the conversation.
they had between each other..
And I can translate first..
- Okay..
- And Sachs, during the debate,.
he said to the atheist, Dawkins,.
he said, "Richard, religion is music,.
but you are tone deaf.".
And then Richard Dawkins, responding to the debate,.
says, "Yes, it's true..
I'm tone deaf, but there's no music.".
(speaking in foreign language).
- And Psalm 19 reminds us, there is music all the time..
There's music..
Can you listen to the music?.
It is playing..
The music is playing all around us..
And the music of creation declare the glory of God..
(speaking in foreign language).
- And that's the first point you have to listen..
You have to try to listen to the music of creation,.
the voice of creator..
And the second point here is listen.
to the voice of the creator..

$^{121}$And God does not only speak through creation..
He also speaks to us, every single one of us,.
directly through His word..
And that's what the Psalm is trying to describe..
It is the beauty and the power of the word of God..
(speaking in foreign language).
- And in verses seven, eight, and nine,.
it says that God's law is so perfect,.
and it revives our soul,.
and His testimony gives wisdom..
His precept gives us direction for our lives,.
and His commands gives clarity into our lives,.
and His judgment are always true and righteous,.
and the fear of the Lord is pure and eternal..
(speaking in foreign language).
- And in short, God's word gives you wisdom,.
gives you life, gives you joy and light,.
and it is reliable, unchangeable,.
and unlike any of our words, our human words..
(speaking in foreign language).
- And the Psalms keeps on telling us.
that the word of God are more than precious,.
than gold, and much, than much pure gold,.
sweeter than honey, and than honey from the comb..
(speaking in foreign language).
- It is more precious than gold,.
and you know, you remember in the history.
of the United States, many, many years ago.
in the 19th century, 1850s, there was a gold rush, right?.
And all the people around the world.
fled to San Francisco, right?.
San Francisco, that's why we call San Francisco.
(speaking in foreign language).
Okay, and they came for gold,.
and they tried to, you know, sift through the soil.
and to find a few flecks of gold..
(speaking in foreign language).
- And there was a story going around.
about the gold rush, and the story goes,.
three men, they discover a gold mine,.

$^{161}$but they pledge to other people.
that we're not gonna tell anyone.
because it's absolutely secret, we wanna get the gold, right?.
So they promise each other never to tell anyone,.
but then when they walk into a bar, these three men,.
and then other people immediately come to them.
and telling them, "Oh, you guys must have found.
"the gold mine.".
(speaking in foreign language).
- Now, the three men were shocked.
because they never tell anyone,.
and they ask, "How do you know that we found the gold mine?".
And the people respond by saying,.
"It's all written over your face.
"because we see how you smile,.
"so you must have found the gold mine.".
(speaking in foreign language).
- Now, this is what the word of God should do to us, right?.
We are filled with the word of God.
and then we can taste the sweetness.
and the treasure of all the treasures of the law of God..
And people around us, you know, eventually should feel.
the power of the word of God, the joy of the Lord, right?.
Shining on our faces..
(speaking in foreign language).
- So this psalm's telling you to listen.
to the music of creation.
and to listen to the voice of the creator,.
the word of God..
And thirdly, from 11 to 14,.
it tells us to respond to this God who loves us..
And after creation speaks, after God's word speaks,.
and what we make is our response, right?.
And the psalm is telling us,.
to show us how we should answer to the word of God..
(speaking in foreign language).
And the psalm says,.
"By them your servant is warned,.
and in keeping the word of God, there is great reward.".
And he immediately responds saying,.

$^{201}$"Who can discern their own sins, their own errors,.
and ask God to forgive our hidden sin and faults.".
(speaking in foreign language).
And ask God to forgive our hidden sin and faults..
(speaking in foreign language).
Now this is an honest confession..
Seriously, sometimes we cannot see all of our sins..
So we must ask God to forgive us, to protect us,.
and to free us from pride and rebellion from sin..
(speaking in foreign language).
And the psalm says,.
"Keep your servant also from willful sins..
May they not rule over me, and I will be blameless..
I will be clean and innocent of great transgression.".
And his desire, as you can see, is to be free from sin,.
not to be ruled by sin, but shaped by the word of God,.
the very word of God..
(speaking in foreign language).
Finally, the last line is the most precious response..
It says, "May the word of my mouth,.
the meditation of my heart, be pleasing in your sight,.
Lord, my rock and my redeemer.".
(speaking in foreign language).
Now the word pleasing here, acceptable,.
how this translation was used often in the Old Testament.
to describe a sacrifice, sacrifice we offer to God..
You know, a sacrifice is only meaningful.
if only God accept it..
And that is the word here..
In the same way, true worship here is not about,.
you know, asking, okay, do I like it or not?.
Do I like the song?.
Do I like the sermon?.
Do I like, you know, here or whatever?.
It's not our preference..
It's that whether God would accept my worship..
(speaking in foreign language).
And that's what meditation is all about, right?.
Your thought, you know, your thought life matters..
You know, I wonder if you can,.

$^{241}$you would take time every day to meditate.
on the word of God, not rushing through the Bible,.
but, you know, let the word of God sink deep,.
deep into our hearts..
And then our words and our thoughts and our lives.
will become offering that will give God glory..
(speaking in foreign language).
So here, Psalm 19 is like music itself..
It has three movements, right?.
To listen to the voice, the music of creation,.
listen to the voice of God,.
and our response finally to God..
So there are three movements here,.
from creation, the word, to response..
And each of the parts calls to God..
So the first part is the word of God..
And the second part is the word of God..
The parts calls to you and me to listen carefully.
and then to step out of faith..
(speaking in foreign language).
So here's my homework for you this week..
You know, slow down and try to listen.
to the voice of creation..
You know, look up the skies and look at the sun,.
not directly, right?.
Even the snowflakes, okay, not right now, I hope not..
But the smallest details would declare the glory of God.
and the music of creation is always playing around us..
(speaking in foreign language).
And then secondly, is to listen to the voice.
of the creator, 7 to 10, verses 7 to 10..
And God's word is more precious than gold,.
sweeter than honey..
We need to taste it, we need to experience it..
It's a treasure that is hidden,.
but it will fill you with joy and life..
(speaking in foreign language).
So this week, spend some time every day.
to listen to the voice of creation.
and listen to the voice of creator..

$^{281}$And then with confession and prayer,.
making our lives pleasing in the sight of God..
(speaking in foreign language).
So here it tells us creation points us to God..
So let creation point us to God..
And then the word of God will transform us..
And our only fitting response is to bow down.
and worship in repentance and joyful obedience..
(speaking in foreign language).
Well, at this point, why don't we conclude with a prayer?.
And in this prayer, I want you to come.
into the presence of God and to pray to him..
(speaking in foreign language).
Let us pray..
(speaking in foreign language).
Lord, thank you for your creation that declares your glory..
And we wanna thank you for your word that give us life..
(speaking in foreign language).
And forgive our hidden sins and keep us from rebellion.
and let our words and our meditation,.
our thought be pleasing to you.
because you are our rock and our redeemer..
(speaking in foreign language).
In Jesus' name we pray..
(speaking in foreign language).
Amen..
So remember the homework this week..
(speaking in foreign language).
Go out to creation, not your screen,.
not your screen, not your phone..
Go out to the creation and listen.
to the music of the creation..
(speaking in foreign language).
Okay, and then spend some time in the word of God.
that gets singed into your lives..
(speaking in foreign language).
Okay, and God bless you this week.
and hope to see you next Sunday..
(speaking in foreign language).
All right, yeah, put it in your heart, yes..

$^{321}$Okay, go in peace, guys..
(speaking in foreign language).
(gentle music).
(gentle music continues).
\newpage



\section{路加福音 4:1-13}
\label{sec:8R1ydESgnbc}
\textbf{屬靈的謊言 (路加福音 4\_1-13) - 梁紫健傳道}
\newline
\newline
連結: \href{https://youtube.com/watch?v=8R1ydESgnbc}{\texttt{https://youtube.com/watch?v=8R1ydESgnbc}} ~~~~ 語音日期: 2025-09-01
\newline
\newline
\hyperref[sec:EMLJsLtjt_w]{\small{< < < PREV SERMON < < <}}
~
\hyperref[sec:index_chronic]{\small{[返順時目]}}
~
\hyperref[sec:index_scriptual]{\small{[返順卷目]}}
~
\hyperref[sec:Iq3SYwSRmiI]{\small{> > > NEXT SERMON > > >}}
\newline
\newline
路加福音 4:1-13
\newline
\begin{longtable}{cl}
\hline
\hline
章節 & 經文 (和合本修訂版)\\
\hline
4:1 & \begin{tabularx}{0.7\textwidth}{X} 耶穌滿有聖靈，從約旦河回來，聖靈把他引到曠野， \end{tabularx} \\ \\ \relax
4:2 & \begin{tabularx}{0.7\textwidth}{X} 四十天受魔鬼的試探。在那些日子，他沒有吃甚麼，日子滿了，他餓了。 \end{tabularx} \\ \\ \relax
4:3 & \begin{tabularx}{0.7\textwidth}{X} 魔鬼對他說：「你若是神的兒子，叫這塊石頭變成食物吧。」 \end{tabularx} \\ \\ \relax
4:4 & \begin{tabularx}{0.7\textwidth}{X} 耶穌回答說：「經上記著：『人活著，不是單靠食物。』」 \end{tabularx} \\ \\ \relax
4:5 & \begin{tabularx}{0.7\textwidth}{X} 魔鬼又領他上了高山，霎時間把天下萬國都指給他看， \end{tabularx} \\ \\ \relax
4:6 & \begin{tabularx}{0.7\textwidth}{X} 對他說：「這一切權柄和榮華我都要給你，因為這原是交給我的，我願意給誰就給誰。 \end{tabularx} \\ \\ \relax
4:7 & \begin{tabularx}{0.7\textwidth}{X} 你若在我面前下拜，這一切都歸你。」 \end{tabularx} \\ \\ \relax
4:8 & \begin{tabularx}{0.7\textwidth}{X} 耶穌回答他說：「經上記著：『要拜主—你的神，惟獨事奉他。』」 \end{tabularx} \\ \\ \relax
4:9 & \begin{tabularx}{0.7\textwidth}{X} 魔鬼又領他到耶路撒冷去，叫他站在聖殿頂上，對他說：「你若是神的兒子，從這裡跳下去！ \end{tabularx} \\ \\ \relax
4:10 & \begin{tabularx}{0.7\textwidth}{X} 因為經上記著：『主要為你命令他的使者保護你； \end{tabularx} \\ \\ \relax
4:11 & \begin{tabularx}{0.7\textwidth}{X} 他們要用手托住你，免得你的腳碰在石頭上。』」 \end{tabularx} \\ \\ \relax
4:12 & \begin{tabularx}{0.7\textwidth}{X} 耶穌回答他說：「經上說：『不可試探主—你的神。』」 \end{tabularx} \\ \\ \relax
4:13 & \begin{tabularx}{0.7\textwidth}{X} 魔鬼用完了各樣的試探，就離開耶穌，再等時機。 \end{tabularx} \\ \\
[1ex]
\hline
\hline
\end{longtable}
$^{1}$各位早安,大家今天星期過得好嗎?.
這個星期的天氣不錯吧?.
大家有沒有想去曬太陽呢?.
在我預備講道的時候,我想起這個題目,.
也想起我初信的經歷,.
就是我在中三開信主,.
教會在一次小組的時候,.
教會的組長跟我說,.
信主耶穌什麼都不怕,.
於是我就自以為耶穌對不同的基督徒會有多一點的保護,.
我就說,信主耶穌不怕的,.
信主耶穌就不怕冷,.
就有這個誤解,.
十二月來自香港,有種濕凍,.
很冷的,.
我就穿著背心走在街上,.
一路走一路在那裡祈禱,.
耶穌保護我,耶穌保護我,.
信了耶穌就不怕冷,.
結果病了一個星期,.
大家應該不會有我這種好像宗教一樣狂熱的誤解,.
聽起來好像很信靠耶穌,.
因為我要靠著耶穌得勝,.
勝過那些冷,.
其實是對聖經的一些誤解,.
如果嚴重一點來說,.
是有些要強迫上帝在這一刻要保護我,.
上帝當然有能力這樣去保護我,.
但是這個是不是祂的意願呢?.
這個經歷跟我們今天要講的經文有一點點相似,.
就是今天還想跟大家講,.
《屬靈的方言》在《路加福音》四章一至十三節,.
相信大家都到這裡耶穌受試探,.
都有很熟悉的,.
甚至一些未信的朋友,.
如果大家有聽過福一六八,.
叫做間歇式斷食的朋友,.
他都會引用這段聖經,.
他用這段聖經去證明斷食是有效的,.
你看看耶穌都是四十天都沒有吃東西,.

$^{41}$雖然我自己都做了六八,.
但是這段經文不是想證明人類四十天不吃東西都可以生存,.
我是不會用這段經文去證明自己做斷食是很有效,.
耶穌二千年前都這樣做,.
來有正當性,.
這段經文講的是魔鬼如何用屬靈的方言去欺騙耶穌,.
所以今天聽完講道之後,.
大家會知道三件事,.
第一件事就是究竟魔鬼如何用聖經去講一些屬靈的方言呢?.
第二,耶穌又如何逆轉當年人類的錯誤呢?.
第三,這些屬靈的方言今天還存不存在這個世界?.
如果大家準備好了,你們一起來看看,.
如果大家有聖經,你們可以用app去開這本聖經,.
路加福音三章的續譜放在耶穌的洗禮,.
天父宣告對耶穌說,你是我的愛子,.
這個續譜還放在壽司禪之間,.
這不是巧合,是一個聚事上的專登安排,.
通常我們都不看續譜,.
見到續譜就會跳過,.
為什麼呢?因為又長又悶,都不知道在講什麼,.
但這續譜很有趣,結束說亞當是神的兒子,.
這裡強化了耶穌是神的兒子,.
而魔鬼試探說了一句,.
如果你是神的兒子,就做什麼就什麼,.
大家如果找一下整本新約的聖經,.
你會發現耶穌多次被稱為神的兒子,.
經常說他是神的兒子,.
但是亞當是神的兒子這個phrase,.
你找遍整本新約的聖經,.
所有的書卷都不會找到,.
唯有在這裡,路加福音,.
在講什麼呢?.
就是在魔鬼試探耶穌之前,路加醫生專登鋪墊,.
耶穌是出自於亞當,.
他寫了整個續譜出來,.
就說亞當是神的兒子,.
這樣強調是想讀者去想,.
亞當一開始是由神創造,.
這個創造的故事,.
亦都令魔鬼這個試探有點不簡單,.

$^{81}$就是怎樣呢?.
他想起了什麼呢?.
就是當初亞當和夏娃這對新婚夫婦,.
為什麼被蛇的謊言誘惑犯罪?.
蛇就跟夏娃說,.
你吃這個果子就不一定死的,.
但是上帝跟亞當說的原話是什麼?.
原宗的果子你都可以吃,.
但是分辨善惡果子,樹上的果子你不能吃的,.
你吃了怎麼選擇?.
一定死的,是一定死的,.
蛇就很厲害地騙他們說,.
不是,不是一定死,是不一定死,.
大家明白分別吧?.
將字擺在不同的位置,.
整件事完全不同了,.
他還說,上帝是怕你們像他那樣厲害,.
所以就不給你們吃而已,.
所以可見這個謊言多麼有說服力,.
多麼令人上當,.
很容易的,.
所以回到盧家福音這裡,.
魔鬼在盧家福音試探耶穌,.
他直接引用說,.
你是亞當的後裔,.
是神的兒子的身份,.
也就是說你也有神的DNA,.
還有什麼?.
還有可能被人騙的習慣,.
就像當年亞當和夏娃被蛇騙的經歷一樣,.
所以魔鬼試圖誘使耶穌濫用權柄,.
他試試耶穌,.
會不會這個神的兒子像亞當那樣,.
當年那樣被我騙呢?.
例如他做了什麼呢?.
他就叫耶穌將石頭變成食物,.
就叫耶穌拜他,.
在聖殿頂找笨豬針,.
這次魔鬼更厲害,.
他竟然用聖經去騙耶穌,.

$^{121}$在這段試探當中,.
短短十三節,.
魔鬼就連續用了三個謊話去騙耶穌,.
希望耶穌做回當年亞當夏娃所做的事情,.
就是僭越上帝的遺訓,.
第一個謊話我們來看看,.
就是石頭變成食物,.
在曠野裡面石頭和食物代表什麼呢?.
大家如果記得在《春夏訣記》裡面,.
石頭和食物代表上帝對以色列子民的公認,.
如果耶穌將石頭變成食物,.
他就混合了這兩個在曠野裡面很重要的象徵,.
石頭,大家記得盆赤出水,.
嗎哪,大家都記得,是食物,.
漢子,飢餓,你可以解決飢餓和口渴,.
有話僅耶穌僭越上帝的公認而自我行動,.
這個是扭曲上帝的公認模式,.
就是在舊約裡面以色列人要順靠著神每天的公認,.
每天都要靠嗎哪和水,而不是自己製造,.
所以魔鬼的提議表面上好像很合理,.
因為耶穌已經禁食了四十天,已經很餓了,.
耶穌有沒有能力這樣做呢?.
當然有能力,但是這樣做的話,.
就會違背了上帝的話,.
所以耶穌就引用了一句話,.
是申命記的一句話,.
人活著,不是單靠食物,.
其實還有後半句,.
就是乃是靠耶和華口中所出的一切話,.
耶和華所出的一切話,.
是指向申命記錄,.
將是猶太人一個非常重要的聖經,.
他們去到死都會背的,.
就是你們要聽,要聽從上帝,.
所以耶穌這樣去回應他,.
非常正當,.
第二個謊話,.
就是魔鬼叫他拜他,.
拜魔鬼就有所有東西,.
其實這件事就奪去了上帝的所有權,.

$^{161}$耶穌是上帝的兒子,.
你猜他知不知道所有東西都是指上帝,.
就好像剛剛那首詩一樣,.
所有東西都是指上帝,.
耶穌依然用申命記去回應魔鬼,.
說唯獨事奉他,.
而第三個謊話,.
魔鬼這下出了一招大絕,.
就是魔鬼竟然用聖經去騙耶穌,.
在路加福音四章裡面第三個試探,.
魔鬼首先將耶穌帶回耶路撒冷,.
他在聖殿頂就說,.
他用詩篇的九十一篇,.
主要為你命令他的使者保護你,.
他們要用手托著你,.
免得你的腳碰在石頭上,.
或者你去玩笨豬跳,.
聲稱神會保護他,.
這個引用其實建立在詩篇九十一篇,.
整個主題上,.
詩篇九十一篇是甚麼呢?.
詩篇九十一篇是講述一個信心的詩歌,.
強調神對信徒的保護,.
免於很多不同的災難,.
但是魔鬼的用法是扭曲的,.
信徒令耶穌濫用神對人類的保護,.
慫恿耶穌證明自己是神的兒子,.
是來塞亞的身分,.
他如何扭曲呢?.
詩篇九十一篇本身是甚麼呢?.
是對信徒的安慰,.
承諾神會差遣天使保護信徒,.
就好像日常的危險,.
或靈界的威脅的時候,.
你會拿出這篇詩篇說,.
神會保護我們,.
但是魔鬼將這個保護,.
轉化為挑戰耶穌的自殺式行為,.
叫他從上面跳下來,.
暗示耶穌可以做甚麼呢?.

$^{201}$他可以強迫神介入人的行為,.
違背了本身要聽從上帝的原則,.
這裡大家要聽清楚,.
不是上帝不可以介入人,.
去保護人,.
而是信徒作為人,.
不能夠強迫神,.
在某個特定的方式去介入,.
去操控上帝,.
大家可能有聽過這個笑話,.
從前有一個虔誠的基督徒,.
就在海上漁欄,.
他浮在海面上說,.
神啊救我吧,.
很快就有一艘船經過,.
是一艘小船,.
就在上面大叫,.
上來吧,我來救你吧,.
然後他搖頭說,.
不用了,上帝會來救我的,.
小船人就走了,.
過了一會兒,.
有一艘大船就走過,.
拋下了那艘船艦,.
就說快點抓住我,.
我來拉你上來,.
然後他就拒絕說,.
謝謝,不用麻煩你了,.
上帝會派天使救我的,.
於是大船就很無奈地走了,.
最後就有一架直升機,.
他放下了一條樓梯,.
他就說抓住那條樓梯,.
我帶你離開吧,.
他還說,不用了,.
上帝一定親自來救我,.
無奈地直升機就飛走了,.
結果他就死了,.
上了天堂,.
他很生氣地問耶穌,.

$^{241}$問上帝,.
主啊我這麼虔誠,.
為什麼你不救我?.
上帝跟他說,.
傻豬,.
我帶了兩艘船,.
一架直升機給你,.
你為什麼不接受呢?.
這個很虔誠的基督徒,.
表面上是很屬靈的,.
因為他覺得他生命上.
只需要上帝的拯救,.
天使跟上帝就好像.
詩篇91篇那樣說,.
是他被攔鎖,.
其他所有東西.
都不是上帝的拯救,.
唯有上帝的顯現,.
才是他認為的上帝拯救,.
這裡我們犯了什麼錯呢?.
就是我們用人的方式.
去限制上帝的介入,.
上帝就只可以用我們想像中的.
那個方式去介入我們的問題,.
其他方式我們都不接受的,.
大家聽完這個笑話.
可能只會淡笑,.
但摩訶利亞的地方就在這裡,.
他用的策略不是.
令你拒絕聖經,.
而是小小地扭曲聖經的方式,.
對上帝理解的一種方式,.
以為我們覺得上帝的介入.
是限於某種的模式,.
當拯救在其他方式出現的時候.
我們就不當是,.
魔鬼是故意忽略上下文,.
用聖經去騙耶穌,.
就像當初他騙亞當夫婦一樣,.
所以耶穌的回應是引用.

$^{281}$《申命記錄章》的,.
他這樣說,.
不可試探主,你們的神,.
去駁斥這種扭曲,.
去主動限制,.
去強調這種是限制上帝的介入方式,.
是一種試探神的行為,.
所以耶穌引用這一句,.
他不單止是引用這一句,.
他說,你不可試探你們的神,.
你耶和華,.
其實後面還有半句,.
就是將你們在馬薩去試探他一樣,.
我們知道通常我們引用聖經,.
不是只有一句,.
是說這個背後的故事,.
我們一定要說一下這個故事,.
就是撒馬的故事,.
馬薩的故事,.
是發生在以色列人在出埃及之後,.
在曠野裡面,.
這裡出埃及記十七章的記載,.
就是以色列人在曠野出發,.
去到一個地方,.
突然間沒有誰人,.
他就和摩西說話,.
抱怨說,.
為什麼?.
為什麼你要帶我們出來埃及?.
為什麼你救我們出來?.
令到我們在這個曠野裡面渴死,.
所以摩西就和神去呼求,.
說可不可以給我們喝點水?.
神就指示他在河裡山盤石,.
打了盤石就出水,.
就有水喝了,.
然後摩西就命名這個地方叫馬薩,.
意思是試探的意思,.
那裡叫米利巴,.
也是爭拗的意思,.

$^{321}$因為以色列人爭拗,.
而且試探耶和華,.
還說什麼?.
耶和華究竟在不在我們中間?.
耶穌的引用是反映了以色列人的懷疑和挑戰上帝,.
因為上帝的拯救就不像以色列人預計那樣,.
我們以為出了埃及就沒有再受苦了,.
但是竟然還要受苦,.
對於以色列人來說,.
這個不是一個拯救,.
這個是一個審判,.
所以耶穌的回應引用申命記,.
不可試探你的主,你的神,.
背後的故事反映了耶穌對聖經上下文的正確理解,.
就是上帝會以他自己的方式和時機去施行拯救,.
而我們不能夠強迫上帝介入,.
這裡是一個很大的諷刺,.
如果在路加福音裡面,.
試探是發生在什麼地方?.
是在曠野和聖殿,.
這象徵著以色列人歷史,.
在曠野挑戰神的介入和對聖殿的迷思,.
以為有聖殿就有上帝的保護,.
其實還有另一個諷刺,.
就是詩篇91篇在猶太人的傳統裡面是一個很有趣的存在,.
他們會用91篇去驅魔,.
猶太人覺得他們可以用91篇去召喚天使出來,.
令他們受到保護,.
在猶太傳統裡面,當他們遇到靈界的事物,.
惡靈邪靈的時候,.
他們就會配上91篇去保護,.
甚至把這篇詩篇抄在羊皮紙上,.
或者製作一些洗牌吊在自己的頸上,.
有個保護,.
就好像有些人去日本求平安符,.
有些甚至會放在他家的嬰兒的床邊,.
怕他被鬼入侵,.
用意就是以色列人遇到靈界危險的時候,.
可以作為護身符,.
他就強調這個詩篇可以召喚天使,.

$^{361}$以介入作為保護,.
而魔鬼正正就是惡靈的代表,.
魔鬼竟然用詩篇91篇去欺騙耶穌,.
形成了一個自我的矛盾,.
魔鬼本身是一個被推逐的對象,.
他扭曲了這篇詩篇的原意,.
嘗試令耶穌像當日以色列人那樣,.
強迫上帝介入,去試探上帝,.
但是耶穌怎樣?.
耶穌沒有像當年亞當夫婦那樣,.
或者以色列人在狂野強迫上帝,.
用一種用他們的方式去迫上帝介入,.
而是重新翻聖經的原則,.
沒有被魔鬼欺騙,逆轉了人類的那種失敗,.
大家有沒有想過,.
這個故事,這個試探故事,.
存在有什麼意義呢?.
其實有點多餘,.
耶穌作為一個永生神的兒子,.
他怎會被魔鬼那些那樣骯髒的手段欺騙呢?.
你知我知,大眼瘤都知,.
魔鬼都知,耶穌絕對不會被他欺騙,.
我們就要問,是不是這件事多餘呢?.
你知我知,魔鬼都知,當然不是多餘,.
在困難的群體裡面,.
有一篇模仿詩篇教授的詩篇,.
叫做感恩詩篇,.
裡面有一個叫公義教師,.
他這樣說,.
他是會戰勝所有的惡靈和邪惡,.
最後所有的惡靈和邪惡,.
都會被他消滅了,.
因為惡靈欺騙了人,.
令他們偏離了律法,.
但耶穌不單止戰勝了魔鬼,.
最不同的是,耶穌做了甚麼?.
就是講聖經,.
然後呢,他沒有消滅惡靈,.
他沒有消滅魔鬼,.
這套聖經展示給我們看,.

$^{401}$耶穌逆轉了人類的失敗,.
他勝過了試探,.
所以為甚麼路加要用耶穌出自於亞當,.
亞當是神的兒子的身分,.
去勝過魔鬼的試探,.
就是他展示給我們看,.
魔鬼這些屬靈的謊言,.
在我們仍然崩壞猖狂的世界裡面,.
依然存在,.
不過,我們可以打破這些謊言,.
看到真相,.
不知道大家有沒有被魔鬼的屬靈謊言騙過呢?.
有些人,靈體有些人就點頭,.
除了一開始跟大家講的凍親的故事,.
以為自己信了耶穌,.
不會凍親之外,.
我都以為我自己沒有被屬靈的謊言騙過,.
我問一個問題,.
你何時會知道一個謊言是謊言呢?.
就是當你識穿了,.
你知道真相的時候,.
你才知道那是謊言,.
所以從小到大,.
我都覺得自己沒有被這些謊言騙過,.
我以為我活在一個真相的年代,.
我以為上帝的拯救,.
只是讓我遇到他,相信他,.
我定期回教會,.
我做一個虔誠的基督徒,.
我有一個美滿的家庭,.
大家都覺得我是好人,.
這就是上帝的拯救,.
這些都沒有錯的,.
不過我以為這是唯一救恩的方式,.
一個我想像中唯一上帝的拯救,.
其實我這樣想,.
我限制了上帝的救恩,.
限制了上帝的救恩在我身上呈現的方式,.
其中一個我發現魔鬼在我身上做的謊言是,.
我曾經有一個想法,.

$^{441}$我覺得這個世界沒有人值得相信,.
所有的關係最終都是以出賣和被拋棄來結束,.
後來還有人用聖經跟我說,.
他說,是啊,聖經說,一個人都沒有,.
所以所有人的關係都不值得我們留戀,.
唯有上帝的關係才值得我們去珍惜,.
聽起來好像很有道理,.
但是忽略了一個沒有耶人,.
沒有在羅馬書的上下文,.
說的是猶太人和所有的民族都是罪人,.
沒有人可以在上帝面前稱義,.
唯有認識律法才知道真相,.
不是說人與人之間的關係不可信,.
無論如何,我不是說這段經文,.
我不花時間去解釋這個上下文,.
不過我想說的是,其實方言都很屬靈,.
人當然有很靠不住的地方,.
但是當我們過分強調人的不可靠的時候,.
會抽走我們本身對於群體的信任,.
令我們不再相信任何人,.
令我們的人生孤單,.
令我們更加疏離群體,.
如果我們這樣做的話,.
我們更加不能夠在教會裡面,.
或者在社區裡面載而不作光,.
所以如果這樣的話,.
如果我們相信這個方言的話,.
今天你和我都不需要去採教會,.
就好像中世紀有些修士,.
在深山裡面修煉和上帝的關係就可以了,.
是不需要的,.
所以正正就是因為耶穌揭示了魔鬼的方言,.
告訴我們我們仍然有機會去看見真相,.
我一開始都沒有發覺這個方言在我心中出現,.
直到我相信了耶穌,.
我就發現很多時候我都不太想在任何群體,.
甚至教會的團契小組,.
去說我的想法,.
我的內心想法和我的感受,.
我很想保護自己的,.

$^{481}$我很怕別人會傷害我和拋棄我,.
所以從小到大我都很怕和別人交心,.
交心意味著我要將我內心深處的一面揭示給別人看,.
我很怕我和別人建立了關係之後,.
最後我還是被人出賣和被人拋棄,.
甚至當我進入神學院的時候,.
我遇到Alice,.
我喜歡她,但我很怕受傷,.
我一直都不敢和她說,.
我一直都不敢表達我的內心,.
後來我一直在想,.
我一直在找答案,.
為什麼呢?.
為什麼我這麼恐懼和害怕和別人交心,.
為什麼我覺得一段關係最後,.
無論是朋友還是戀愛還是九鼎之妹,.
最後都會被出賣和拋棄呢?.
我一直都想不到答案,.
我找不到真相,.
於是我就去見輔導,.
輔導員叫我留意身體,.
他說身體很神奇,.
會告訴我們一些潛藏的情緒,.
就是我發現每當我很害怕進入一段關係的時候,.
我胸口的右邊會有一陣痛楚在裡面,.
那種痛楚很細微,.
是當我靜下來,.
我感受我的身體的時候,.
我才感覺到,.
上星期大家一起去郊遊,.
很多人都去了,.
不知道大家有沒有加入導遊帶我們.
在草地上感受自己的session,.
在那個時間我也感受到一點,.
當然沒有以前那麼強烈,.
但是這麼多年來都還有在,.
然後我就發現,.
這個謊言,.
所有關係最後都是以出賣和被拋棄作結的謊言,.
是來自我小時候一段經歷,.

$^{521}$一段被拋棄的經歷,.
我記得在我五歲的時候,.
那時候是住在深圳,.
媽媽需要回香港生一個小孩,.
那時候我剛好有腦癌症,.
在我一個最無助的時候,.
最需要人照顧我的時候,.
我媽媽離開了我,.
我知道我媽媽很疼我,.
她為勢所迫,不得已這樣做,.
她為了我們整個家庭著想,.
但是你想一下,.
一個五歲的小朋友,.
身邊無依無靠,.
和一個大我三年的姐姐,.
整天在街上流浪,.
去謀生,.
一個本身不需要自己去謀生的小朋友,.
出盡力氣去掙扎去生存,.
儲下了很多憤怒,.
到我長大之後,.
我一直都沒有發現,.
這段經歷影響得我很深,.
這段經歷就好像告訴一個五歲的小朋友,.
是的,連最親的人都拋棄你,.
還有什麼可以相信呢?.
於是我就去見父親,.
我見了大概三年的時間,.
後來我才發現我很生氣我媽媽,.
其實我知道媽媽很疼我,.
家裡以前無論沒有什麼東西吃,.
她都會留最好的東西給我,.
令我一直以為媽媽不喜歡吃好東西,.
其實她是養一些東西給我吃,.
所以我一直都覺得,.
我否認我很生氣我媽媽,.
直到有一天,.
我發現這個謊言很影響我身邊的關係,.
影響得很嚴重,.
甚至沒辦法相信我身邊的人會不離不棄,.

$^{561}$所以我發掘我內心的憤怒,.
慢慢打開了傷口,.
直到幾年前有一天,.
我發現我自己終於準備好了,.
我有機會跟媽媽說這件事,.
我媽媽當時也準備好了,.
因為當時我媽媽還未相信,.
我就跟她說,.
我發現有一件事,.
我一直都處理不了,.
我信仰裡面唯有耶穌才可以幫我,.
但我不知道耶穌何時,.
用什麼方式去做這件事,.
你可不可以跟我一起祈禱,.
我就跟她說了整件事,.
她當然哭了,.
她說很大反應,.
我真的很盡力地愛你們,.
但我不知道這樣做會令你這麼受傷,.
我說,.
我明白,.
我明白當年你的苦楚,.
我知道你很愛我,.
我只是想說我當年的感受,.
然後我們兩個一起哭著祈禱,.
後來我發現我的傷口慢慢好了一點,.
我開始相信我身邊的關係,.
慢慢打破了這個謊言,.
用了很多年的時間,.
用了五六年的時間,.
所有關係最後都是以被出賣和拋棄作結,.
這個謊言後來慢慢用了一段時間,.
我才慢慢打破了,.
有同學,.
神學院的同學還取笑我,.
說我在WhatsApp,.
我現在還寫了一段聲明,.
我說我不會再拒絕你的愛,.
他取笑我,.
以為我這樣是跟Alice說這句話,.

$^{601}$其實不是,.
這句話是我對上帝說的,.
因為以前我被一些屬靈的謊言騙了很久,.
上帝透過直升飛機,船,.
很多身邊不同的人的愛,.
給我很大的關心和愛,.
我都接不到,.
我都說那不是上帝的,.
那些不是上帝的,.
當我識破了這個謊言後,.
我才接得到,.
所以今天來一個總結,.
魔鬼如何用聖經去說一些屬靈的謊言呢?.
魔鬼扭曲了詩篇91篇,.
上帝對人的保護,.
試圖令耶穌重蹈當年亞當的覆轍,.
強迫上帝以某種方式去介入,.
耶穌如何逆轉當年人類的錯誤?.
耶穌沒有像公義教師那樣撕滅惡靈,.
是重新對申命記的正確理解,.
祂戰勝了魔鬼,.
祂展示給人看,.
是的,有很多屬靈的謊言在世界上,.
但是我們仍然可以打破這些謊言,.
第三,這個我想大家都知道,.
屬靈的謊言還存不存在這個世界上?.
《福音》4章13字是這樣說的,.
魔鬼用完那件事之後就離開耶穌,.
再等時機,再等時機,.
魔鬼沒有這樣就算了,.
他在等待時機,.
就像現在的世代一樣,.
充斥著不同的事,.
你有沒有留意到屬靈的謊言.
用些什麼形式去出現?.
你準備好揭發.
生命中一些屬靈的謊言嗎?.
如果你認為今天你一直看到的都是真相,.
沒有謊言的話,.
那這個似乎是其中一個最大的謊言,.

$^{641}$因為連魔鬼都受到,.
即連耶穌都受到魔鬼的試探,.
只是他沒有上當,.
魔鬼會不會放過我們這些普通的人?.
8月31日,.
8月31日是對香港來說一個很重要的日子,.
願有一天我們一起來打破魔鬼的謊言,.
看到真相,我們一起練.
耶穌謝謝你,.
用你的生命去揭示屬靈的謊言,.
令我們這個世代仍然有希望,.
我們可以因為主的利益的緣故,.
你給我們的聖經,.
我們打破一些騙了我們很久的東西,.
我們不知道弟兄姊妹心中有多少東西,.
是我們未去揭示,.
但主你幫我們,.
你幫我們,.
你幫助我們繼續一起走下去,.
我們這次禱告,.
奉求耶穌的恩物,.
祝福大家,.
今天的宣派到此.
\newpage



\section{馬太福音 7:13}
\label{sec:Iq3SYwSRmiI}
\textbf{What Would You Choose? (Matthew 7\_13 - 27) - Chee Yie Jong 你會怎樣選擇?(馬太福音 7\_13-27) - 楊志溢傳道}
\newline
\newline
連結: \href{https://youtube.com/watch?v=Iq3SYwSRmiI}{\texttt{https://youtube.com/watch?v=Iq3SYwSRmiI}} ~~~~ 語音日期: 2025-09-08
\newline
\newline
\hyperref[sec:8R1ydESgnbc]{\small{< < < PREV SERMON < < <}}
~
\hyperref[sec:index_chronic]{\small{[返順時目]}}
~
\hyperref[sec:index_scriptual]{\small{[返順卷目]}}
~
\hyperref[sec:0No_EkcZeRE]{\small{> > > NEXT SERMON > > >}}
\newline
\newline
馬太福音 7:13
\newline
\begin{longtable}{cl}
\hline
\hline
章節 & 經文 (和合本修訂版)\\
\hline
7:13 & \begin{tabularx}{0.7\textwidth}{X} 「你們要進窄門。因為通往滅亡的門是寬的，路是大的，進去的人也多； \end{tabularx} \\ \\
[1ex]
\hline
\hline
\end{longtable}
$^{1}$Alright, good morning, Yuyi..
大家早上好!.
It's good to be with you all again this morning..
非常感激,今天有機會跟大家一起分享神的信息..
Okay, so we have been in a series on Jesus' Sermon on the Mount for actually more than.
a year now, going through Jesus' teaching from Matthew 5-7, chapter by chapter, section.
by section..
And today is a special day because we are finally coming to the concluding portion for.
this entire Sermon on the Mount..
And this is a very important section because after laying out all His teachings and principles.
for His Kingdom of Heaven, Jesus now demands a response from us..
我們已經在登山保训的系列信息中學习了一年多..
從马太福音第五章到第七章,逐章逐段的研读耶穌的教导,今天.
我們終于來到整篇登山保训的結尾部分..
這是一個非常重要的段落,因為在阐明了天國所有教导和原则.
之后,耶穌现在要求聽众做出回应..
So if you have been with me since the beginning of this series back in July last year, my prayer.
for you is that at the very least, some of what Jesus has taught here have resonated with.
you or make an impact in your personal walk with the Lord..
如果你從去年七月這個系列开始就一直同我一起學习的話,我為你禱告..
希望耶穌在這里的教导至少有一些已經触动了你,化在于你與主同行的路上产生了影響..
So before we dive into our passage today, let's just do a quick review of Jesus' Sermon on the Mount..
And so when Jesus began his ministry, this is the central message that he preached..
"Repent, for the kingdom of heaven has come near.".
在进入今天的經文之前,让我們快速回顾一下耶穌的登山保训..
当耶穌开始祂的事工時,祂所传講的信息是,天國盡了..
你們应当悔改..
And so in this first century world, dominated by the Greco-Roman Empire,.
Jesus is actually ushering a new kingdom and a new way of living.
that is radically different from the culture around them..
And so in his preaching, Jesus announced that the kingdom of heaven has now arrived here on earth..
And the Sermon on the Mount is actually where we find out what this new way of living.
that Jesus ushers in looked like..
So if you want to pull out your Bible, you can follow with me along for a little bit..
In Matthew chapter 5, Jesus begins his sermon by casting a vision of what a true flourishing life.
or what a truly blessed life looks like in this kingdom of heaven in the Beatitudes..
And so contrary to the culture, a true flourishing life comes from being poor in spirit,.
being humble and meek, and being persecuted because of righteousness..
And then this true flourishing life should leave a positive impact in the world that we live in.
and point other people to God the Father..

$^{41}$And this is why Jesus called us to be salt of the earth and light of the world..
[speaking Chinese].
And then in verse 17, Jesus proclaims that he has come to fulfill the law of God perfectly..
As the only one that has perfectly fulfilled all of God's law,.
Jesus also established himself as the ultimate interpreter of these laws..
And he reveals to us the heart of God and the better righteousness behind all these laws..
And so in verse 20, he called us to live out this better righteousness redefined by him,.
not the religious leaders or the teachers of the law of his days..
[speaking Chinese].
And from there, in the later half of Matthew 5,.
Jesus show us six case studies of this true and better righteousness.
that will lead to a truly flourishing life..
Beginning with the issues of anger and lust,.
and then God's design for marriage,.
and then integrity in the way we talk in our speech,.
and then finally to not seeking revenge and responding to our enemy with love..
This, according to Jesus, should be the norm in his new kingdom of heaven..
[speaking Chinese].
And in Matthew 6, Jesus continues with his kingdom teaching.
by addressing three of the most common religious activities of his days,.
giving to the poor, praying, and fasting..
And once again, Jesus redefines what true religion is according to his new kingdom,.
one that is not concerned with the praise of the people,.
but for the approval of the Father alone..
And in the section on prayer, this is also where we find the Lord's Prayer..
[speaking Chinese].
And then beginning in verse 19 in chapter 6,.
Jesus turns his attention to our relationship towards our earthly possessions..
And so, instead of worshiping our earthly wealth,.
Jesus tells us to invest in heavenly treasure that will last for eternity..
And when we are able to set our heart on heavenly treasure,.
we don't have to worry about our earthly needs anymore,.
because we know that the God who takes care of the birds in the sky.
and the wildflowers in the field will take care of us as well..
[speaking Chinese].
And when we come to chapter 7,.
Jesus continues to lay out the social responsibilities.
of those who are part of his new kingdom,.
especially in the way we treat others..
And so, in chapter 7, verse 12,.

$^{81}$Jesus gives us this summary statement of our social responsibility towards others.
according to his kingdom ethics..
And this is what sometimes we call the golden rules..
So, do to others what you would have them do to you..
And this is the true and better righteousness that Jesus has shown us..
And this is also the new way of living that Jesus wants us to follow.
as the citizens of his kingdom of heaven..
[speaking Chinese].
And so, in chapter 12,.
Jesus gives us this summary statement of our social responsibility towards others.
according to his kingdom ethics..
And this is what sometimes we call the golden rules..
So, do to others what you would have them do to you..
And this is the true and better righteousness that Jesus has shown us..
And this is also the new way of living that Jesus wants us to follow as the citizens of his kingdom..
And so, all of this leads us to our passage today,.
where Jesus presents to all his listeners, including us,.
with an important decision to make..
And this decision is a binary decision,.
which means that there is an either/or decision..
There is no room for being on the fence or not making a choice,.
since everyone has to decide for themselves..
[speaking Chinese].
So, what are the two options here?.
Let's take a look at our passage today..
So, this is Matthew 7, verses 13 and 14..
"Enter through the narrow gate,.
for wide is the gate and broad is the road that leads to destruction,.
and many enter through it..
But small is the gate and narrow the road that leads to life,.
and only a few find it.".
And so, this leads me to my first point today,.
that everyone has to make a conscious decision..
[speaking Chinese].
Yes, so everyone has to make a conscious decision..
And so, when I was out on my run earlier last week,.
I saw this exact verse written on the ground around 33rd and Wentworth..
[speaking Chinese].
So, it says, "Enter through the narrow," and then it got blurred a little,.
but it's "Enter through the narrow gate.".

$^{121}$So, what is this narrow gate and this narrow path.
that our passage is talking about today?.
So, this narrow path is the way of the kingdom of heaven..
This way of the kingdom is not the most glamorous or popular one..
It will be restricting and confining for those who want to enter through it..
[speaking Chinese].
And so, again, if you look back at all the kingdom principles and values.
that Jesus has laid out in his sermon on the mount,.
they're all very counter-cultural and demanding in some sense, right?.
For example, to turn the other cheek, to love your enemies,.
and to not gather up material wealth for ourselves, and the list goes on..
[speaking Chinese].
So, to choose this narrow path means that it will not be an easy journey.
because there will be trials and challenges along the way..
[speaking Chinese].
So, that is the narrow gate and the narrow path on one hand..
And on the other hand, there is the wide gate and the broad path,.
which I would call the way of this world..
And this way of the world welcomes all kinds of challenges..
[speaking Chinese].
So, the narrow path is the way of the kingdom of heaven..
And I would call the way of this world..
And this way of the world welcomes all kinds of beliefs, lifestyles, and values..
And this is why the broader path is so popular,.
and many people choose to enter through it..
[speaking Chinese].
And in a world that prioritizes personal freedom,.
this is a very attractive path, right?.
Imagine living in a world where there is no moral absolute,.
and you are just free to do whatever you want to do..
This might sound like a very compelling way to live initially,.
but I can guarantee you that there will come a point.
when the emptiness and the loneliness in life begins to settle in,.
and you just begin to wonder,.
"Is there more to life than this?".
[speaking Chinese].
And this is when you realize that the freedom you have.
is actually not producing the kind of fulfilling and happy life.
that you think this freedom can produce..
[speaking Chinese].

$^{161}$And this becomes even clearer.
when we take a look at the final destination of these two different paths..
The only path that will lead to a true and flourishing life.
is not the wide and easy one, but the narrow path..
And so, yes, while this narrow path is restricting and confining,.
it is the only way that will lead to true freedom..
[speaking Chinese].
I like how this author and pastor John Mark Comer puts it.
when he described the correlation between the dichotomy.
of the way of life that is restrictive and confined on one hand,.
and the resulting freedom that it brings on the other hand..
[speaking Chinese].
So speaking of rule of life,.
or you can think of it simply as a plan to follow Jesus,.
John Mark Comer says,.
"The rule of life is an invitation to a very different definition of freedom.
than that of the modern world,.
an invitation to embrace constraints,.
and that if you give yourself to them, will eventually set you free.".
[speaking Chinese].
So, in the same way, if we're willing to embrace these constraints.
and limitation and restrictions that are being placed upon us.
because we choose to follow Jesus on this narrow path,.
it will ultimately lead to true freedom in Christ..
[speaking Chinese].
So, this is the first point that I want to emphasize today,.
that everyone has to make a conscious decision for themselves..
And Jesus makes his command abundantly clear for us,.
"Enter through the narrow gate.".
[speaking Chinese].
However, as we have seen,.
this journey through the narrow path will not be an easy one,.
and there will be trials and temptation along the way..
And because of that, we need to be alert..
And so this brings me to my second point today,.
that we need to be aware of spiritual deception..
[speaking Chinese].
So, let's take a look at verse 15 to 20..
[speaking Chinese].
"Watch out for false prophets. They come to you in sheep's clothing,.

$^{201}$but inwardly they are ferocious wolves..
By their fruit you will recognize them..
Do people pick grapes from thorn bushes or figs from thistles?.
Likewise, every good tree bears good fruit,.
but a bad tree bears bad fruit..
A good tree cannot bear bad fruit,.
and a bad tree cannot bear good fruit..
Every tree that does not bear good fruit is cut down and thrown into the fire.".
[speaking Chinese].
So on the way to the narrow path, we have to be alert and vigilant,.
looking out for false prophets or false teachers.
who are trying to lead us away from the path towards life..
And most of the time, these false teachers,.
they don't just come with a bright neon sign telling you who they are..
But instead, they usually come in spiritual disguise,.
like wolves in sheep's clothing..
[speaking Chinese].
And if you pay attention to just the culture around us,.
you can probably easily find all sorts of different narratives.
that promise the kind of true and flourishing life.
that only Jesus can provide..
And so if we don't do our due diligence in discerning.
all these different narratives that are being thrown at us every day,.
we could easily fall into their traps..
[speaking Chinese].
And so because of this, Jesus warned us, "Watch out for false prophets.".
And according to verse 20, the best way that we can discern.
if someone is a true or false prophet.
is not by what they preach or what they teach,.
but it is by the fruit that they produce,.
and whether or not the fruit of the Spirit is evident in their life..
[speaking Chinese].
So this is the first type of spiritual deception.
that we need to be aware of, one that comes from other people..
But Jesus also warns us about the second type of spiritual deception,.
that is our own self-deception..
[speaking Chinese].
And so we need to be aware of this..
So verse 21, "Not everyone who says to me, 'Lord, I am a prophet,'.
is a true prophet.".

$^{241}$[speaking Chinese].
And so we need to be aware of this..
[speaking Chinese].
And so we need to be aware of this..
Verse 21, "Not everyone who says to me, 'Lord, Lord,'.
will enter into the kingdom of heaven,.
but only those who does the will of my Father who is in heaven..
Many will say to me on that day, 'Lord, Lord,.
did we not prophesy in your name, and in your name drive out demons,.
and in your name perform many miracles?'.
Then I will tell them plainly, 'I never knew you..
Away from me, you evil doers.'".
[speaking Chinese].
And so the keywords here are the word "says" and "do" in verse 21..
And so following Jesus is not just about saying it.
or confessing it with our mouth,.
but our action have to reflect what we confess and believe..
[speaking Chinese].
The book of James talks about that faith without work is dead..
So if your faith in Jesus does not make you love God more.
and want to make Him the absolute priority of your life,.
your faith might be dead..
And if your faith in Jesus does not make you become a more loving person.
towards your family, your coworkers, or your neighbors,.
again, your faith might be dead as well..
[speaking Chinese].
And the fact that someone can prophesize, drive out demons,.
and perform miracles in the name of God,.
and yet not be able to enter into the kingdom of heaven,.
this might come as a surprise to some of us..
But this kind of spiritual self-deception is a possibility that all of us face..
[speaking Chinese].
So, why is the disconnect here for these false disciples of Jesus?.
The answer is plain and simple..
They are doing the right thing, but for the wrong motives or wrong reasons..
So these false disciples, they crave for the power and status.
that Jesus represents..
They are not doing the right thing..
They are doing the wrong thing..
And they are not doing the right thing..

$^{281}$So these false disciples, they crave for the power and status.
that Jesus represented..
But they never truly acknowledge Jesus as the Lord and Master of their life.
by allowing the will of God to control and guide their life..
[speaking Chinese].
And by the way, just a quick side note,.
this is still true 2,000 years later in our 21st century America,.
where people would use the name of Jesus to gain power, popularity, or even status..
[speaking Chinese].
So, be aware of our own self-deception..
Take an honest look at yourself and examine your intention,.
because we can be doing the right thing for God,.
but for the wrong motives and reason..
And you don't want to wait until the day of judgment,.
and only to find out that you're being denied entering into the kingdom of heaven..
[speaking Chinese].
So this is a very stern warning from Jesus, right?.
But again, of course, we all have the freedom to choose.
to either listen to Jesus' warning or to completely ignore it..
So in this next section of our passage,.
Jesus differentiates between those who truly listen and follow what he has commanded.
versus those who do not..
And so this is my third and final point, that true hearing requires obedience..
[speaking Chinese].
So if you look at verse 24, it begins with the word "therefore.".
And what this means is simply that what Jesus is about to say here.
is directly related to what comes before..
[speaking Chinese].
So let's read this together..
"Therefore, everyone who hears these words of mine,".
referring to all Jesus' teaching in the Sermon on the Mount,.
"and puts them into practice is like a wise man who built his house on the rock..
The rain came down, the streams rose, and the winds blew and beat against that house..
Yet it did not fall, because it has its foundation on the rock..
But everyone who hears these words of mine,".
again referring to Jesus' teaching in the Sermon on the Mount,.
"and does not put them into practice is like a foolish man who built his house on sand..
The rain came down, the streams rose, and the winds blew and beat against that house..
And it fell with a great crush.".
[speaking Chinese].

$^{321}$So the main difference between these two groups of people is that.
the first group in verse 24, they hear these words of Jesus and puts them into practice..
While the second group in verse 26, they hear these words of Jesus just like the first group,.
but they do not put them into practices..
[speaking Chinese].
So Jesus compares the first group to a wise man who builds his house on a solid foundation.
that will stand through all the life, all the winds and storms of life..
But then he compares the second group to a foolish man who builds his house on sand..
And when the flood and the rain comes, these two groups of people,.
they do not put them into practices..
They do not put them into practices..
So tying back to the first examples of the white gate and the broad path that leads to destruction,.
those who share the words of Jesus but do not put them into practice.
will also end up in a similar destiny,.
where the life that they spend years and years building.
will eventually come tumbling down..
(translation).
So, we come to the end of Jesus' teaching on the Sermon on the Mount..
And Matthew, which is the author of this book,.
he then closed his entire section out by resuming the narrative,.
where he concludes,.
"When Jesus has finished saying these things, the crowds were amazed at His teaching,.
because He taught as one who had authority and not as the teachers of the law.".
(translation).
And one of the worst things that could happen after we hear from a sermon like this.
is to think that this has nothing to do with us,.
because Jesus is only speaking to His audience..
But the truth is, these words of Jesus apply to us today as well..
We all have to make our own choice to either follow Jesus.
and the way of His kingdom above, or to follow the way of our world..
And again, Jesus makes it very clear for us, "Enter through the narrow gate.".
(translation).
So, I have a few questions for us to reflect on today..
So....
(translation).
Are you willing to deny yourself and follow Jesus along this narrow path,.
despite all the trials and challenges that will come your way?.
(translation).
And are you willing to live a life in full obedience to Jesus,.
even though it means that you have to live in a counter-cultural way?.

$^{361}$(translation).
So, spend some time in your seat and just reflect on your own,.
and in a few minutes, I'm going to close it out with a prayer..
(translation).
(music).
Let's pray..
(translation).
Heavenly Father, You have shown us the way that leads to life through the narrow path..
So, help us to walk faithfully on this path, even when it is difficult,.
and very few people choose it..
(translation).
Protect us from all the deceptions in the world around us,.
and also within ourselves, that would lead us astray..
And Father, give us the wisdom to discern what is true..
(translation).
Jesus, we thank You for Your promise that You will be with us to the end of the earth,.
as we journey along this narrow path..
(translation).
And Holy Spirit, come and transform our lives so that we can bear the fruit of Your Spirit..
(translation).
In Jesus' name we pray, Amen..
(translation).
Well, thanks, Chi, for a challenging message, and let us sing to the Lord..
\newpage



\section{傳道書 1:12-18}
\label{sec:0No_EkcZeRE}
\textbf{日光之下再思考 (傳道書 1\_12-18) - 譚文倩莊傳道}
\newline
\newline
連結: \href{https://youtube.com/watch?v=0No_EkcZeRE}{\texttt{https://youtube.com/watch?v=0No\_EkcZeRE}} ~~~~ 語音日期: 2025-09-08
\newline
\newline
\hyperref[sec:Iq3SYwSRmiI]{\small{< < < PREV SERMON < < <}}
~
\hyperref[sec:index_chronic]{\small{[返順時目]}}
~
\hyperref[sec:index_scriptual]{\small{[返順卷目]}}
~
\hyperref[sec:dx4StTt36lE]{\small{> > > NEXT SERMON > > >}}
\newline
\newline
傳道書 1:12-18
\newline
\begin{longtable}{cl}
\hline
\hline
章節 & 經文 (和合本修訂版)\\
\hline
1:12 & \begin{tabularx}{0.7\textwidth}{X} 我傳道者在耶路撒冷作過以色列的王。 \end{tabularx} \\ \\ \relax
1:13 & \begin{tabularx}{0.7\textwidth}{X} 我用智慧專心探尋、考察天下所發生的一切事：神給世人何等沉重的擔子，使他們在其中勞苦！ \end{tabularx} \\ \\ \relax
1:14 & \begin{tabularx}{0.7\textwidth}{X} 我見日光之下所發生的一切事，看哪，全是虛空，全是捕風。 \end{tabularx} \\ \\ \relax
1:15 & \begin{tabularx}{0.7\textwidth}{X} 彎曲的，不能變直；缺乏的，不計其數。 \end{tabularx} \\ \\ \relax
1:16 & \begin{tabularx}{0.7\textwidth}{X} 我心裡說：「看哪，我大有智慧，勝過在我以前所有統治耶路撒冷的人；我的心也多經歷智慧和知識的事。」 \end{tabularx} \\ \\ \relax
1:17 & \begin{tabularx}{0.7\textwidth}{X} 我專心想要明白智慧，想要明白狂妄與愚昧，方知這也是捕風。 \end{tabularx} \\ \\ \relax
1:18 & \begin{tabularx}{0.7\textwidth}{X} 因為多有智慧，就多有愁煩；增加知識，就增加憂傷。 \end{tabularx} \\ \\
[1ex]
\hline
\hline
\end{longtable}
$^{1}$早晨各位弟兄姊妹.
不知道大家知不知道.
我是一個科技電腦的天才.
今天我是第一次自己做PowerPoint.
然後自己按鍵.
如果今天PowerPoint有什麼問題.
記住千萬不要怪台後的弟兄姊妹.
只要怪就怪我這個天才.
有一天我整天在看聖經.
看書,看電腦.
看得很頭痛.
我跟老公說.
老公,我想我最近用腦太多.
令到自己很頭痛.
我老公說:是這樣的媽媽,全讀書都是這樣說.
我老公說:是嗎?真的這樣?.
我貪得意就去GPT問.
那怎麼辦?.
你猜他怎麼回應我?.
是我錯的.
他就說:你想不想我幫你找一個簡單又不頭痛的全讀書.
閱讀路線圖?.
我想:咦?這麼有趣?.
有這樣的東西?.
我當然馬上按了.
想看看GPT會怎麼回應我.
按了不夠三秒.
答案馬上出來了.
不過很長.
我一看就覺得又很頭痛.
所以我決定不看.
算了吧.
想再說一說.
全讀書是《聖經》舊約裡其中一本智慧文學.
是深入探討人生的意義.
書中指出我們人生追求的財富,享受和智慧.
都是短暫而且難以掌握的.
尤其是當人面對苦難和死亡的時候.
更加凸顯了我們人的有限.
人的一生最終可能都是虛空.

$^{41}$不過你千萬不要誤會全讀書好像想打沉我們.
相反全讀書裡的全讀者.
是很想藉著自己成為一個反面教材.
提醒人要怎麼過生活.
我們要有敬畏神的心.
有一個永恆的眼光.
更加要享受來自神給我們的恩典.
因為認識神才是我們人生真正意義的所在.
同時不要忘記祂是一個施行最終審判的主.
上一次我們在日光之下當思考.
我們一起看了全讀書一章一至十一節.
當中裡面一句虛空的虛空.
說這句話的人有錢到不得了.
也很厲害.
從他的口中說出這句話.
你會不會有點好奇.
為什麼他會這樣說呢.
不知道你有沒有留意.
我們大多數人一生都在追求一些東西.
而這些東西通常都和幸福平平安安.
開開心心有關.
至少我從來沒有聽過.
我一生最想追求的是很多痛苦.
雖然我見過一些好像很享受痛苦的人.
甚至是基督徒.
不知道是不是洗禮的時候.
浸池裡的水加了很多檸檬汁.
他們經常都罵我.
今天我們繼續一起思考全讀書.
全讀者在這裡分享他用智慧.
嘗試尋找生命的意義的過程.
我們今天試試去思考兩個問題.
智慧和知識有什麼分別呢.
智慧以外其實還有沒有出路呢.
和大家一起去想一想.
全讀者在這段經文究竟想說什麼呢.
我們又一起來去看看.
這個日光之下我們再思考.
我先和大家一起去讀一下今天這段經文.
全讀書一章十二至十八節.

$^{81}$全讀者在耶路撒冷坐過以色列的王.
我用智慧專心探尋.
考察天下所發生的一切事.
神給世人何等沉重的擔子.
使他們在其中勞苦.
我見日光之下所發生的一切事.
看那全是虛空全是暴風.
幻曲的不能辨識.
缺乏的不計其數.
對說看那我大有智慧.
勝過在我以前所有統治耶路撒冷的人.
我的心也多經歷智慧和知識的事.
我專心想要明白智慧.
想要明白狂妄與愚昧.
方知這也是暴風.
因為多有智慧就多有愁煩.
增加知識就增加憂傷.
這七字經文裡面.
智慧這兩個字.
這個詞匯出現了五次.
我用智慧專心探尋考察.
我心裡說看那我大有智慧.
我的心也多經歷智慧.
我專心想要明白智慧.
因為多有智慧.
一章十二至十三節.
我傳道者在耶路撒冷.
以色列的王.
我用智慧專心探尋.
考察天下所發生的一切事.
專心applied my heart.
這個字不單只是在說.
知識上的學習.
而是在說這個傳道者.
他是用心.
他是全人去投入.
用經歷的方法去尋找.
applied my heart to seek and search.
一章十五節.
彎曲的不能辨直.

$^{121}$缺乏的不計其數.
這一句是要和我們之前.
傳道書一章一至十一節.
那個循環筆式去對照.
大自然有它固定的規律.
太陽的升降.
下雨的水.
雪山融下來的水流去江河.
江河再流去大海那裡.
但是那個大海是永遠都不會滿社的.
我們人無論多厲害都好.
我們沒有能力去改變這些規律.
一章十六至十八節.
我心裡說看啊.
我大有智慧.
勝過在我以前所有統治耶路撒冷的人.
我的心也多經歷智慧和知識的事.
我專心想要明白智慧.
想要明白徬徨與愚昧.
方知這也是暴風.
因為多有智慧就多有受煩.
增加知識就增加憂傷.
傳道者強調.
他不單止有智慧.
他的智慧勝過以前所有統治耶路撒冷的人.
他是一個厲害之中的厲害.
注意十六,十七節說到.
我心裡.
我專心.
英文是I set in my heart.
I applied my heart.
第三節說到.
專心探尋考察之後.
再一次說他多努力.
去尋找人生意義的答案.
你看不看到這個傳道者.
在尋找人生意義的這件事上.
他是多認真.
他不單止增加知識.
accumulate knowledge.

$^{161}$更加去經歷,去感受,去experience.
連瘋狂和愚蠢的事.
madness and folly.
他都嘗試用智慧去經歷.
不知道大家有沒有想過.
智慧,wisdom.
和知識,knowledge.
有什麼分別.
知識是指頭腦上掌握的資料.
事實.
智慧是生活和決策裡面.
我們懂得怎樣去應用這些知識.
並且我們知道我們自己的限制.
人知道自己的限制.
知道這個世界上.
有一個比我更加厲害.
更加厲害.
掌管萬有的神.
知識可以令人變成一個專家.
但智慧會令人成為一個謙卑,檢視生命的人.
傳道書的智慧在聖經裡.
時時都會被謙卑和人知道自己的有限去互相緊扣.
因為若果人不用一個敬畏神的心.
做根基.
承認我們自己生命的有限的話.
以為我們依靠自己的努力就可以得到一切.
以為有腦就可以解決一切的事情.
最終我們人都會跪在這裡.
最簡單就是當我們面對困難,苦難,痛苦的時候.
我們人有話事權嗎.
記得我和媽媽在醫院做化療的時候.
我就有一個很深的體會.
我看著我媽媽坐在椅子上被護士插針.
我看著我媽媽.
我在想她身體裡的癌細胞有多細.
但她又有很大的威力去傷害我媽媽.
一個150幾磅的人.
被癌細胞殘食到剩下100磅.
由1800年代的Rudolf Wischau開始.
過去二百多年我們有多少頂尖的科學家,醫護人員.

$^{201}$投入多少錢,人力,物力.
但到今時今日對癌症仍然是很束手無策.
簡單來說.
其實我們人不是真的那麼厲害.
如果你有聖經的話.
你可以嘗試打開它.
或者你回家可以再看一看.
你去留意一下這段經文.
其中有兩句的格式是有點不同的.
在傳道書裡.
傳道者經常用對比,對稱的修辭.
一章十五節.
彎曲的不能辨直.
缺乏的不計其數.
一章十八節.
因為多有智慧就多有受煩.
增加知識就增加憂傷.
作者將這兩句分別出來.
在視覺上和結構上.
更加要突出當中的哲理.
要讀者停下來.
花點時間思考當中的合意.
這兩節經文想表達甚麼呢.
經文是想表達我們人生裡有很多彎曲.
我們人力無法改變的事情.
無論人多聰明都不能解決所有的問題.
提醒我們我們的有限.
你以為讀書讀得多.
我們知得多.
就可以建造一個幸幸福福.
平平安安快快樂樂的人生.
但我想我們大家都知道.
有時我們知道的越多.
我們反而越看到這個世界的歪曲和荒謬.
受煩反而更加多.
我們人更加意識到我們自身的無能為力.
和我們人的有限.
而這個限制正正是傳道書一章十三節裡所講的沉重的擔子.
今時今日不知道你們知不知道.
對我們這個世代的小孩來說.

$^{241}$最沉重的擔子是甚麼呢.
不是功課.
不是考試.
是一樣隱藏但影響很深遠的東西.
就是同輩的接納.
Peer's acceptance.
去到某個年紀他們開始會很尊重自己的外表.
很注重我受不受歡迎.
在社交媒體裡面他們出的貼文.
多不多人like多不多人follow.
社交媒體成為他們情緒和自我價值的一個指標.
我們做父母的當然不想他們受傷害.
所以我們就會去看多些書.
去學習.
去上YouTube看看有甚麼可以幫助到我們.
但現實是有時無論我們有多努力.
我們都鬥不過他們身邊的同學朋友.
鬥不過互聯網或AI的影響.
結果我們父母有時知道越多.
我們反而越擔心.
接著輪到我們揹著一大袋沉重的擔子.
傳道書一章十八節.
因為多有智慧就多有愁煩.
加增知識的就加增憂傷.
原來有時知識只是放大了我們對問題的焦慮.
提醒我們再一次提醒我們的無能為力.
回到之前我問的兩個問題.
智慧和知識有甚麼差別.
智慧以外其實還有沒有出路呢.
Francis Collins一位醫生.
前國家人類基因研究所所長.
在他27歲那一年.
他在幫一個病人看病的時候.
那個病人問他一個問題.
你對上帝的看法是怎樣的.
他覺得這個問題他沒有想過.
這個問題結果成為了他生命裡面一個很關鍵的轉折點.
在他尋找答案的過程當中.
他意識到自己的無神論立場是沒有根基的.
隨著他科學生涯的發展.

$^{281}$他越來越發覺他是不能忽視生命的意義.
死亡的本質.
深層的存在問題.
當他思考生命意義的時候.
面對患病的病人.
他開始質疑自己的世界觀.
迫使他重新檢視信仰和科學之間的關係.
在這段沉沉覓覓的過程中.
他看了C.S. Lewis的《Mere Christianity》.
《純粹基督教》.
看完之後他得出一個結論.
就是相信普世的道德法則.
都是指向這個世界上的確是有一個造物主.
Francis Collins認為.
科學最擅長回答有關自然世界裡面的.
「效」如何的問題.
信仰卻能夠回應「Why」為甚麼的問題.
例如宇宙的起源.
或者人生最終目的的問題.
他發現宇宙的精密和秩序.
在數學的定律裡面.
是可以找到這個造物主的證據.
在他繼續探索的過程當中.
他更加認定上帝是一個很驚人的物理學家,數學家.
他也指出道德律的存在.
即是我們人對善和惡的本能的理解.
相比純粹無神論者所說的那套東西.
是有挑戰的.
而這種道德的意識.
用神的存在來解釋是更加合理的.
最終他發現自己那個無神論的立場是不合邏輯的.
因為科學的證據指向這個世界是有一個造物主.
而基督教就提供了一個理性而且充實的人生觀.
有一次他去Cascade Mountain 登山的時候.
面對大自然的壯觀.
Francis Collins感受到神對他的呼召.
他就在這個情況下決志信主.
信主之後他發現基督教的信仰和他的科學研究是沒有衝突的.
相反增添了他研究工作裡面的屬靈層面.
傳道書裡面關於人生意義和尋找真理的教導.

$^{321}$和他的信仰旅程有很多共鳴的地方.
特別是傳道書裡面提到虛空的虛空.
全是虛空.
展示了我們人類的知識其實是多麼有限.
唯有敬畏上帝才是真正的智慧.
真正的智慧和上帝的思想是互相呼應的.
最近我留意到我有一個朋友正在經歷一些困難.
一些的意思是很多.
是眾數.
很多的困難.
我一直在閱讀他分享的時候.
我發覺他真的面對很多挑戰.
很多挫敗.
很多不公平的事.
很多很多要解決的問題.
也有更多更多解決不了的問題.
過去兩三年.
他的生命裡面好像永遠只是不斷有很多困難.
挑戰,挫敗.
不斷在重疊.
有錢有心的.
很多很多.
閱讀他的分享的時候.
我自己一邊看也有一份無力感.
我不敢用聖經去鼓勵他.
因為他比我更加熟悉聖經.
我也不知道如何去安慰他.
因為我在想.
我可以說些什麼去安慰他.
有時候我們人的安慰.
我們的言語.
都是很有限的.
我想起傳道書的訊息.
它不是說那些不認識神的人就慘了.
你死定了.
因為你不認識主.
就是這樣.
相對認識神的人.
我們就有豐盛的生命.
我們好人一世平安.

$^{361}$所以大家就快點來信耶穌.
不是這樣.
的確不是這樣.
因為我們做基督徒和非基督徒一樣.
我們都會面對很多生命裡面的不公平.
挑戰,困難,苦難.
低處未算低.
可以再低的情況.
你有沒有問過自己.
其實為什麼我們要做基督徒.
有一次我又一如既往.
去和我的mentor聊天.
我也不太記得我們當時在聊什麼.
但我記得他給我的回應.
那個回應就是.
上帝沒有欠我們什麼.
回想起來.
傳道書一章一至十一節.
強調生命的循環和我們人生的勞苦.
我們人類的成就迅速消逝.
到頭來都是虛空.
來到今天一章十二至十八節.
傳道者用智慧去嘗試找到真正的滿足.
開始親身去尋問人生意義.
這段經文不是在描述虛空.
而是很積極地去尋問,去實驗,去思考.
很明顯傳道書一章十二至十八節.
這裡開始.
經文的氛圍似乎有一些希望的轉變.
傳道者在這裡開始探索.
人類的智慧和成就的界限.
為尋找日光之下以外更深層的意義去鋪路.
我們人生真的有很多挑戰.
有很多我們控制不了.
我們沒有話事權.
煩惱到令人累到一個點的事.
我又想給另一個問題你們去思考一下.
如果我們將生命看作為日光之上那位所賜的禮物.
我們人又應該怎樣去思考去過我們的日子呢.
特別如果你是現在經歷一些苦難.

$^{401}$一些好像你做不到什麼去解決的事情.
你怎樣去過這些這麼難的日子呢.
我們一起低頭祈禱.
主要求你幫助我們去思考.
不要讓我們做一個負負碌碌的人.
讓我們真的去思考探索生命.
知道你給我們最大的禮物是什麼.
以致我們不枉此生.
這樣去祈禱.
奉耶穌得正面靠.
Amen.
給我們一個永恆的盼望.
讓我們不只是在日光之下生活.
而是在日光之上.
這個救恩就是上帝給我們的一份最好的禮物.
給我們生命裡面由今生到永恆.
它都是有一個終極的意思.
所以今天我們來到聖餐的主人.
我們也來領受主的聖餐.
因為這個救恩在聖經裡面告訴我們.
是一個最重要的東西.
世界上最重要的東西.
就是我們能夠領受主耶穌給我們的恩典.
聖餐裡面我們需要提醒自己.
這個餅和這個杯代表什麼.
保羅和教會這樣說.
他說我當日傳給你們的是從主領受的.
就是主耶穌被賣的那一天.
他拿了一塊餅來.
捉住之後他就抹開.
他說這個是我的身體為你們寫的.
你們要與之行為的事件.
飯後也照樣拿起杯子來.
他說這個杯子是我的血所納的神藥.
你們每逢喝的時候要與之行為的事件.
每逢吃這個餅喝這個杯子.
就要表明主的誓.
宣告主的誓.
就是等到耶穌基督再翻你的身.
就準備給我們一個盆.

$^{441}$我們同心協力祈禱.
然後我們才領受主的聖餐.
你再次提醒我們.
在永恆的當中.
你為我們安排了這個恩典.
這個救恩.
讓我們能夠得到永生.
讓我們能夠有這個永生的角度.
來看我們今時的人生.
生命可能對我們很忙碌.
很趕急.
但我們知道我們的生命.
有了你永恆給我們的恩典的時候.
我們生活是為了你的恩典.
我們能夠去事奉你.
能夠成為你自己的使命.
求主你再次在這個聖餐的主人.
去提醒我們.
在恩典裡面.
在基督的十教裡面.
我們能夠有這個這麼榮美的福分.
拿著這個餅這個杯子的時候.
我們再次去紀念.
紀念耶穌基督在十字架上為我們所做.
十字架上為我們流血和死去.
這種人的言語不能夠言說的話.
我們要去感恩去讚美你.
讓這個恩典能夠不斷地更新我們.
就是我們的教會.
我們的生命.
能夠贖你.
能夠見證你.
多少福主在這裡.
我們每一次聖餐的時候都提醒大家.
聖餐是什麼意思.
聖經告訴我們.
這個主耶穌的身體和祂在十字架上為我們流血.
如果你已經是接受了耶穌.
你很清楚耶穌和你的關係.
我們歡迎你去到我們士大一張的聖餐桌.

$^{481}$然後拿一個餅和杯子.
然後回到座位.
我們一起領受.
(有些人瘋狂地在想).
(「我喜歡的東西會被我吃掉」).
(「沒有機會的事也會被我吃掉」).
(「我喜歡的東西會被我吃掉」).
(我喜歡的東西會被我吃掉).
(「我喜歡的東西會被我吃掉」).
好,我們先跟心靈受主的佩.
然後我們跟靈主的佩.
再一次我們祈禱,好嗎.
讓你的恩典更新我們.
以至我們的品性,心思,意念.
能夠更屬天.
更加屬於你自己的國度.
但另外,教會讓我們在每個崗位上.
能夠盡心盡力去事奉你.
特別是教會來到這三個月.
我們會改善我們的新的核心動工.
主,我們求你如來同在.
讓我們能夠參與.
在你自己的國度的事故中.
以至教會能夠繼續成長.
祈禱,奉主命令.
剛才Carmen提醒過大家.
今天開始我們會有註冊.
我們每次到選舉的時候.
三年一次,我們就….
\newpage



\section{使徒行傳 9:1-43}
\label{sec:dx4StTt36lE}
\textbf{不可預期的福音 (使徒行傳 9\_1-43) - 陳崇基牧師}
\newline
\newline
連結: \href{https://youtube.com/watch?v=dx4StTt36lE}{\texttt{https://youtube.com/watch?v=dx4StTt36lE}} ~~~~ 語音日期: 2025-09-20
\newline
\newline
\hyperref[sec:0No_EkcZeRE]{\small{< < < PREV SERMON < < <}}
~
\hyperref[sec:index_chronic]{\small{[返順時目]}}
~
\hyperref[sec:index_scriptual]{\small{[返順卷目]}}
~
\hyperref[sec:lW1Rt3jzEUM]{\small{> > > NEXT SERMON > > >}}
\newline
\newline
使徒行傳 9:1-43
\newline
\begin{longtable}{cl}
\hline
\hline
章節 & 經文 (和合本修訂版)\\
\hline
9:1 & \begin{tabularx}{0.7\textwidth}{X} 掃羅不斷用威嚇兇悍的口氣向主的門徒說話。他去見大祭司， \end{tabularx} \\ \\ \relax
9:2 & \begin{tabularx}{0.7\textwidth}{X} 要求發信給大馬士革的各會堂，若是找著信奉這道的人，無論男女，都准他捆綁帶到耶路撒冷。 \end{tabularx} \\ \\ \relax
9:3 & \begin{tabularx}{0.7\textwidth}{X} 掃羅在途中，將到大馬士革的時候，忽然有一道光從天上下來，四面照射著他， \end{tabularx} \\ \\ \relax
9:4 & \begin{tabularx}{0.7\textwidth}{X} 他就仆倒在地，聽見有聲音對他說：「掃羅！掃羅！你為甚麼迫害我？」 \end{tabularx} \\ \\ \relax
9:5 & \begin{tabularx}{0.7\textwidth}{X} 他說：「主啊！你是誰？」主說：「我就是你所迫害的耶穌。 \end{tabularx} \\ \\ \relax
9:6 & \begin{tabularx}{0.7\textwidth}{X} 起來！進城去，你應該做的事，必有人告訴你。」 \end{tabularx} \\ \\ \relax
9:7 & \begin{tabularx}{0.7\textwidth}{X} 同行的人站在那裡，說不出話來，因為他們聽見聲音，卻看不見人。 \end{tabularx} \\ \\ \relax
9:8 & \begin{tabularx}{0.7\textwidth}{X} 掃羅從地上起來，睜開眼睛，竟不能看見甚麼。有人拉他的手，領他進了大馬士革。 \end{tabularx} \\ \\ \relax
9:9 & \begin{tabularx}{0.7\textwidth}{X} 他三天甚麼都看不見，也不吃也不喝。 \end{tabularx} \\ \\ \relax
9:10 & \begin{tabularx}{0.7\textwidth}{X} 那時，在大馬士革有一個門徒，名叫亞拿尼亞。主在異象中對他說：「亞拿尼亞！」他說：「主啊，我在這裡。」 \end{tabularx} \\ \\ \relax
9:11 & \begin{tabularx}{0.7\textwidth}{X} 主對他說：「起來！往那叫直街的路去，在猶大的家裡，去找一個大數人，名叫掃羅；他正在禱告， \end{tabularx} \\ \\ \relax
9:12 & \begin{tabularx}{0.7\textwidth}{X} 在異象中看見了一個人，名叫亞拿尼亞，進來為他按手，讓他能再看得見。」 \end{tabularx} \\ \\ \relax
9:13 & \begin{tabularx}{0.7\textwidth}{X} 亞拿尼亞回答：「主啊，我聽見許多人講到這個人，說他怎樣在耶路撒冷多多苦待你的聖徒， \end{tabularx} \\ \\ \relax
9:14 & \begin{tabularx}{0.7\textwidth}{X} 並且他在這裡有從祭司長得來的權柄，要捆綁一切求告你名的人。」 \end{tabularx} \\ \\ \relax
9:15 & \begin{tabularx}{0.7\textwidth}{X} 主對他說：「你只管去。他是我所揀選的器皿，要在外邦人、君王和以色列人面前宣揚我的名。 \end{tabularx} \\ \\ \relax
9:16 & \begin{tabularx}{0.7\textwidth}{X} 我也要指示他，為我的名必須受許多的苦難。」 \end{tabularx} \\ \\ \relax
9:17 & \begin{tabularx}{0.7\textwidth}{X} 亞拿尼亞就去了，進入那家，把手按在掃羅身上，說：「掃羅弟兄，在你來的路上向你顯現的主，就是耶穌，打發我來，叫你能再看得見，又被聖靈充滿。」 \end{tabularx} \\ \\ \relax
9:18 & \begin{tabularx}{0.7\textwidth}{X} 掃羅的眼睛上立刻好像有鱗一般的東西掉下來，他就能再看得見，於是他起來，受了洗， \end{tabularx} \\ \\ \relax
9:19 & \begin{tabularx}{0.7\textwidth}{X} 吃過飯體力就恢復了。 \end{tabularx} \\ \\ \relax
9:20 & \begin{tabularx}{0.7\textwidth}{X} 立刻在各會堂裡傳揚耶穌，說他是神的兒子。 \end{tabularx} \\ \\ \relax
9:21 & \begin{tabularx}{0.7\textwidth}{X} 凡聽見的人都很驚奇，說：「在耶路撒冷殘害求告這名的不就是這個人嗎？他不是到這裡來要捆綁他們，帶到祭司長那裡去嗎？」 \end{tabularx} \\ \\ \relax
9:22 & \begin{tabularx}{0.7\textwidth}{X} 但掃羅越發有能力，駁倒住在大馬士革的猶太人，證明耶穌是基督。 \end{tabularx} \\ \\ \relax
9:23 & \begin{tabularx}{0.7\textwidth}{X} 過了好些日子，猶太人商議要殺掃羅， \end{tabularx} \\ \\ \relax
9:24 & \begin{tabularx}{0.7\textwidth}{X} 但他們的計謀被掃羅知道了。他們晝夜在城門守候著要殺他。 \end{tabularx} \\ \\ \relax
9:25 & \begin{tabularx}{0.7\textwidth}{X} 他的門徒就在夜間用筐子把他從城牆上縋了下去。 \end{tabularx} \\ \\ \relax
9:26 & \begin{tabularx}{0.7\textwidth}{X} 掃羅到了耶路撒冷，想與門徒結交，大家卻都怕他，不信他是門徒。 \end{tabularx} \\ \\ \relax
9:27 & \begin{tabularx}{0.7\textwidth}{X} 只有巴拿巴接待他，領他去見使徒，把他在路上怎麼看見主，主怎麼向他說話，他在大馬士革怎麼奉耶穌的名放膽傳道，都述說出來。 \end{tabularx} \\ \\ \relax
9:28 & \begin{tabularx}{0.7\textwidth}{X} 於是掃羅在耶路撒冷同門徒出入來往，奉主的名放膽傳道， \end{tabularx} \\ \\ \relax
9:29 & \begin{tabularx}{0.7\textwidth}{X} 並和說希臘話的猶太人講論辯駁，他們卻想法子要殺他。 \end{tabularx} \\ \\ \relax
9:30 & \begin{tabularx}{0.7\textwidth}{X} 弟兄們知道了，就帶他下凱撒利亞，送他往大數去。 \end{tabularx} \\ \\ \relax
9:31 & \begin{tabularx}{0.7\textwidth}{X} 那時，猶太、加利利、撒瑪利亞各處的教會都得平安，建立起來，凡事敬畏主，蒙聖靈的安慰，人數逐漸增多。 \end{tabularx} \\ \\ \relax
9:32 & \begin{tabularx}{0.7\textwidth}{X} 彼得在眾信徒中到處奔波的時候，也到了住在呂大的聖徒那裡。 \end{tabularx} \\ \\ \relax
9:33 & \begin{tabularx}{0.7\textwidth}{X} 他在那裡遇見一個人，名叫以尼雅，得了癱瘓，在褥子上躺了八年。 \end{tabularx} \\ \\ \relax
9:34 & \begin{tabularx}{0.7\textwidth}{X} 彼得對他說：「以尼雅，耶穌基督醫好你了，起來！整理你的褥子吧。」他立刻就起來了。 \end{tabularx} \\ \\ \relax
9:35 & \begin{tabularx}{0.7\textwidth}{X} 凡住呂大和沙崙的人都看見了他，就歸向主。 \end{tabularx} \\ \\ \relax
9:36 & \begin{tabularx}{0.7\textwidth}{X} 在約帕有一個女門徒，名叫大比大，翻出來的意思是多加；她廣行善事，多施賙濟。 \end{tabularx} \\ \\ \relax
9:37 & \begin{tabularx}{0.7\textwidth}{X} 當時，她患病死了，有人把她清洗後，停在樓上。 \end{tabularx} \\ \\ \relax
9:38 & \begin{tabularx}{0.7\textwidth}{X} 呂大原與約帕相近；門徒聽見彼得在那裡，就派兩個人去見他，央求他說：「請快到我們那裡去，不要耽延。」 \end{tabularx} \\ \\ \relax
9:39 & \begin{tabularx}{0.7\textwidth}{X} 彼得就起身和他們同去。他到了，就有人領他上樓。眾寡婦都站在彼得旁邊哭，拿多加與她們同在時所做的內衣外衣給他看。 \end{tabularx} \\ \\ \relax
9:40 & \begin{tabularx}{0.7\textwidth}{X} 彼得叫她們都出去，然後跪下禱告，轉身對著屍體說：「大比大，起來！」她就睜開眼睛，看見彼得，就坐了起來。 \end{tabularx} \\ \\ \relax
9:41 & \begin{tabularx}{0.7\textwidth}{X} 彼得伸手扶她起來，叫那些聖徒和寡婦都進來，把多加活活地交給他們。 \end{tabularx} \\ \\ \relax
9:42 & \begin{tabularx}{0.7\textwidth}{X} 這事傳遍了約帕，就有許多人信了主。 \end{tabularx} \\ \\ \relax
9:43 & \begin{tabularx}{0.7\textwidth}{X} 此後，彼得在約帕一個皮革匠西門的家裡住了好些日子。 \end{tabularx} \\ \\
[1ex]
\hline
\hline
\end{longtable}
$^{1}$多謝.
我們繼續看我們這個跨月的系列.
我們說過我們要看聖經裡面有一本書叫《試堂傳》Book of Acts.
如果你未接觸過聖經的話.
你只需要打開新約聖經第五本書.
裡面記載了一本書叫《試堂傳》.
《試堂傳》講述初代教會的故事.
即是說耶穌基督死而復活升天之後.
有什麼事發生呢?.
初代的教會,第一代的教會到底有什麼經歷等等.
我們在《試堂傳》當中已經說了好幾章.
這次我們跳到第九章.
第九章有什麼事發生呢?.
就是一個人信主的故事.
我相信我們在座的百分之九十以上.
我想我們都已經信了主.
你記不記得你信主的時候是怎樣的?.
記不記得誰帶你信主?.
或者什麼情形,什麼情景信主?.
多少年前?或者最近?.
無論什麼時間,什麼情況都好.
我們都需要謹記我們怎樣信主.
或者可能你還沒有信主.
可能你都應該想一下.
到底教會是什麼一回事呢?.
或者我們所說的耶穌基督的福音是一回什麼事呢?.
你的朋友信主或者你的家人是信主的.
他們是基督徒.
他們到底經歷過什麼呢?.
為什麼他們會覺得這件事是這麼大一件事呢?.
這次我們來到第九章.
就是說到一個人叫做蘇羅.
或者叫做保羅.
他是信主的經歷.
我們先看這段聖經.
就說到這位叫做蘇羅.
或者他的羅馬的名字.
蘇羅就是希伯來的名字.
保羅就是羅馬的名字.
好像我們今天有很多人有中文名.

$^{41}$有英文名.
有西班牙名.
或者有台山話名.
有幾個名字.
蘇羅和保羅都是指同一個人.
第九章就說到他信主的經歷是怎麼樣.
我們先看第九章第一節.
直到第九節.
這段故事是說到他怎樣遇上這位基督.
第九章第一節.
蘇羅不斷用威嚇,空巷的口氣.
向主的門徒說話.
去見大祭司.
要求發信給大馬士革的國會堂.
並約是找著信奉者道的人.
無論男女都准他捆綁帶到耶路撒冷.
蘇羅在途中.
將到大馬士革的時候.
忽然有一道光從天上下來.
四面照著他.
他就突然在地上.
聽見有聲音對他說.
蘇羅 蘇羅.
你為什麼迫害我?.
他說 主啊 你是誰?.
主說 我就是你所迫害的耶穌.
起來 進城去.
你應當做的事必有人告訴你.
同行的人站在那裡.
說不出話來.
因為他們聽見聲音卻看不見人.
蘇羅從地上起來.
睜開眼睛.
竟不能看見什麼.
有人拉他的手.
領他到大馬士革.
他三天什麼都看不見.
也不渴也不吃.
這個就是所謂蘇羅或保羅歸主的故事.
他怎樣信主.

$^{81}$有什麼經歷.
在聖經裡面曾經提過.
基督死後復活之後.
他向門徒顯現有好幾次.
四本福音書.
馬太 馬可 路加和約翰.
都有記載耶穌基督榮耀的顯現.
這一次.
《士郎傳》第九章.
可以算是最後一次.
他榮耀的顯現去彰顯.
彰顯在一個人身上.
這個人就叫蘇羅.
或者之後我們叫他保羅.
這個人信主.
蘇羅或保羅.
其實在歷史裡面.
無論你是講聖經的歷史.
或者世界的歷史.
都是一個很重要的故事.
因為這個人本身是一個怎樣的人.
我們大概可以去到形容這個人是一個.
可能大家也見過這樣的人.
穿黑衣 穿白衣的那些正統 超正統的猶太人.
有沒有見過.
我們芝加哥如果上去北面.
Skokie.
一些這樣的倫社.
猶太人的倫社.
你就可以大概會聽到.
會見到 接觸到一些這樣的人.
就是一些超正統的猶太人.
他們的服飾全部都是黑白的.
他們是嚴守律法的.
他們猶太人會叫做Kosher.
就是一些食物的條例等等.
這些是當時第一世紀猶太人的世界裡面.
是有很多這樣的猶太人.
當然猶太人有很多種的.
好像華人有很多種.

$^{121}$有些住在南洋.
有些在香港 有些在中國大陸等等.
華人也有很多種的.
猶太人也是一樣.
有一種的猶太人我們可以叫做超正統的猶太人.
在第一世紀裡面也是一樣.
這種猶太人是非常保守的一班猶太人.
我們叫他們做法利賽人.
法利賽是什麼意思.
法利賽是舊約和新約之間有幾百年的時候.
他們湧現出來的一個教派.
他們是什麼意思呢.
什麼叫法利賽呢.
法利賽大概是意思是分別出來的意思.
就是和其他人分別出來.
而且這些人的意識形態是非常保守.
甚至可以很激進的.
可以當是一種.
用今天的說話來說.
就是一些極右的激進的組織等等.
保羅或者素羅大概就是一種這樣的人.
我初初看聖經的時候.
很多年前.
我看到這一段的時候.
第一節我就不明白為什麼他會這樣.
這裡說到素羅是做什麼的.
他說不斷用威嚇和凶悍的說話.
向主的門徒說話等等.
初初接觸聖經的時候.
或者初初看史蘭傳的時候.
都不太明白為什麼素羅.
如果是一個這麼宗教的味道的人.
為什麼他可以對主的門徒.
當時初代的教會.
當時他未信主的時候.
他就去迫害他們.
用一個威嚇或者凶悍的態度.
來出現.
這個都很難摸不著頭腦.
今天你會不會想像一個很愛主的人.

$^{161}$他用凶悍和威嚇的說話跟你說話.
會不會?.
我們應該不會吧.
對嗎?.
應該是比較溫柔的人.
我們印象中.
如果你真的是一個很愛神的人.
很為上帝的話語著緊的人.
應該不會說凶悍和黑社會吧.
這裡有一點這樣的味道.
因為我們摸不著頭腦.
直到我們明白背景的時候.
我們才知道為什麼他會有這種姿態出現.
或者用這樣的方法去面對.
處理其他人.
當時的背景.
法拉利賽人.
他們是有一種這樣的熱心.
他們的熱心就是捍衛自己的祖宗.
一直傳下來的信仰傳統.
那些猶太人.
猶太人有很多種.
但是我們說的這種超級非理性的猶太人.
他們是為了捍衛他們有的傳統.
他們的意識形態其實就是這樣的.
他們認為上帝的國度為什麼還沒有到.
你記得那個背景嗎.
當時的猶太人他們住在巴勒斯坦地.
或者有很多在海外的.
他們有沒有自己的政治獨立.
有還是沒有.
他們沒有的.
他們一直幾百年受其他的帝國欺凌他們.
他們就覺得在聖經裡面.
舊約聖經裡面.
說到上帝的國度會來到的時候.
為什麼還沒有到呢.
因為上帝的國度來到的時候.
好像大衛的君王會再駕臨.
然後會管治他們.

$^{201}$他們就不會被外邦或者世俗的帝國來管治他們.
所以他們有一種.
用今天的華語去形容.
有些宗教文化的戰爭.
這樣的好民.
因為他們覺得這個國度.
上帝的國度一定會來到.
為什麼還沒有到呢.
就是因為有些人阻礙著地球轉.
有些什麼人阻礙著地球轉.
就是他們自己的同胞.
有些同胞不謹守律法.
不像保羅索羅他們那樣.
謹守著上帝每一條的律法.
down to the letter.
就是因為那些人太relax.
所以就搞到我們上帝的國度.
一直都遲延了還沒有到.
遲了又還沒有到.
就是因為他們同胞裡面的不謹守律法.
所以因為這樣的熱心.
有這樣的激進的態度.
你就明白為什麼索羅或者保羅.
他們要千方百計的.
來殲滅當時的教會.
因為當時的教會大部分都是猶太人.
他們和其他猶太人有什麼分別.
就是他們相信耶穌.
就是舊約聖經裡面所預言應許的救主.
對嗎.
這個就是初期教會.
但是對於索羅來說.
我們怎麼可能接受到一個十字架上死的人.
為我們救主.
這是一個很羞辱的事情.
被帝國欺負了八百年之後.
判他死刑.
因為他無辜死了.
我們怎麼可能接受他是我們的君王.
所以對於保羅來說.

$^{241}$初期教會的同胞都是猶太人.
他們去歸主的時候.
這個就真的阻礙地球轉了.
所以他就一定要去將他們剷除.
所以為什麼保羅這裡.
是用這種姿態出現.
或者是一種極右的激進的態度.
來對待初期的教會.
你記得在第九章之前發生什麼事.
再回望到七章最後的時候.
你記得第七章發生什麼事.
就是史提反廣道.
史提反是當時猶太的信徒其中一個.
廣道完了之後.
人們很不喜歡聽他廣道.
然後就將他打死了.
記得嗎.
因為史提反的廣道.
上帝榮耀的彰顯.
到處都有.
但是偏偏就不在聖殿.
所以他們最喜歡的聖殿.
最覺得神聖的聖殿.
偏偏史提反就說.
上帝的榮耀不是住在人造石頭的殿裡面.
所以那些猶太人就非常之憤怒.
然後就將史提反殺死.
七章第57節.
眾人大聲喊叫.
不想再聽他廣道.
然後想衝上去.
然後將他推出去城外.
然後就將他處決.
這裡說到作見證的人.
就將他們的衣裳.
處決一個人的時候.
當時的一些禮儀.
他們就將外套脫下.
然後就放在一個名叫素羅的青年的腳前.
這是當時的一種儀式.

$^{281}$當時他們本地法庭的一種儀式.
但是八章第一節.
素羅也贊同去處死史提反.
贊同不是說你們這樣做.
其實他好像官一樣.
他是給認可.
就是說他們認為這個人是不是應該要處死.
去到侮辱我們的聖殿.
素羅是當時的官.
是當時他們本地的領袖.
素羅就說OK.
先處決他.
從那一天起.
耶路撒冷的教會就遭受大迫害.
素羅就卻殘害教會.
然後逐家逐戶的上去.
將男與女關在監獄裡面.
這個就是表明.
當時的素羅不是一個.
他不是一個蛋散.
不是一些簡單的人.
他是當時教派裡面其中一個的領袖.
你可以看到他們就是想.
不擇手段甚至於訴諸於暴力.
就要阻礙教會的發展.
所以為什麼在第九章說到他信主的經歷裡面.
他就上路.
從耶路撒冷去到哪裡.
去大馬士革.
你們看到地圖在左邊.
就是從耶路撒冷起步.
然後他一直向北走.
大概應該繞過薩馬利亞地.
因為正統的猶太人是不會去薩馬利亞地的.
因為又是歷史的背景.
當時薩馬利亞人和猶太人是不瓜葛的.
我們不喜歡經過他們的地方.
因為不會跟他們同桌吃飯.
因為那些人在猶太人的眼中.
特別是超正統的猶太人眼中.

$^{321}$這些人是該死的.
因為他們是雜種的.
所以保羅大概應該繞過薩馬利亞山地.
然後上到加利利的地方.
然後再上到今天我們叫做高蘭高山的地方.
一直上去右上角.
就是東北邊.
然後就去到大馬士革.
為什麼大馬士革這麼重要呢.
為什麼要長途跋涉離開耶路撒冷.
因為教會已經散播在不同的地方.
為什麼一定要上到大馬士革.
去殲滅初期的教會呢.
原因就是地理.
因為你知道大馬士革.
在今天的敘利亞.
而且今天在敘利亞.
現代的敘利亞有一個很大的地方.
當時的大馬士革.
是通往世界各地的地方.
就是十字路口.
就是我們9094和290的交界.
就是你要去其他地方.
一定要經過那個地方.
那個是一個會點.
為什麼呢.
因為到了大馬士革的時候.
你往右邊就是今天的伊拉克.
東敘利亞等等的地方.
然後你轉左.
就是去到今天的土耳其.
以前叫做小亞細亞.
直到歐洲的地方.
所以如果教會散播去這個地方的話.
他就一定要pre-emptive strike.
就是首先做一個攻擊.
然後就去阻止基督教的散播.
所以保羅就趕路.
應該是騎馬吧.
他就要去到大馬士革.

$^{361}$然後將這些人全部抓起來.
這樣初期的教會才不會繼續散播.
當他上路的時候.
我們都要了解當時的訴羅保羅在想什麼.
右邊這個圖就告訴我們.
當時的超正統猶太人.
他們是怎樣去想像不同的空間.
等於我們今天頭腦裡面都有一個地圖.
對不對.
有些地方我們不去.
為什麼呢.
因為有些地方我們不喜歡.
或者那裡的人我們不喜歡.
或者千萬不要去那些地帶等等.
在你的mind map那裡都有.
這個就是猶太人的mind map.
猶太人的mind map是怎樣的呢.
特別是超正統的猶太人.
世界裡面有很多地方.
但最神聖的地方就在哪裡.
當然就是耶路撒冷聖殿.
The most holiest place就是耶路撒冷.
耶路撒冷以外就是聖城.
聖城以外就是聖地.
聖地就是以色列巴勒斯坦地.
上帝所給他們的地方.
然後就是其他的地方.
所以他們從耶路撒冷趕路去大馬士革的時候.
其實在他的mind map那裡.
他就是由一個最神聖的地方去一個最世俗的地方.
或者一個最乾淨的地方.
去一個最骯髒的地方.
或者最不應該去的地方.
不過因為塑羅有一個使命.
就是他知道有些同胞阻礙地球轉.
所以他就要千方百計用不同的手段.
甚至訴諸於暴力.
他就希望阻止這些教會的擴展.
但其實他就是想將上帝的國帶來.
因為所有同胞大家乖乖地.

$^{401}$我們守法的時候.
上帝的國就會來到.
有些不乖的我們就將他抓了.
然後將他殺死也好.
或者將他禁聖也好.
你看看這幅墨羅.
就是講到塑羅.
怎樣將上帝的國帶來.
就是用人的方法去殲滅所有的對手.
可以不擇手段.
為了做什麼.
就是要光復神的國度.
當耶穌向保羅顯現的時候.
就跟他說.
他說你是迫害我的保羅.
你是迫害我的.
對著一個初期的教會.
或者他所迫害的耶穌.
教會就代表耶穌.
這個保羅就是追殺一個無權無勢的.
一個這樣的人.
上帝的國度.
神的國度.
神的榮耀是怎樣來到的.
對於保羅來說.
我們要去到天國底.
我們要將他剷除.
這裡解釋了保羅的心態之外.
我們還可以有什麼值得我們去看一下.
就是在第九章.
你看到嗎.
復活的主在大馬式向保羅顯現.
其實這個不是第一次.
由第七章開始.
我們剛才講的史提反.
他說什麼.
至高者不會住在人手所做的耶路撒冷聖殿.
然後上帝的聖靈的工作.
就由耶路撒冷開始.
去到其他地方.

$^{441}$去到耶路撒冷以外的地方.
你看到第八章發生什麼事.
我們沒有時間看.
但是薩瑪利亞人信主.
有人將福音帶到薩瑪利亞.
猶太人覺得最不應該去的地方.
偏偏在那裡聖靈就降臨.
有人信主.
然後一個太監信了主.
那個太監是埃塞俄比亞來的.
在南邊來的.
然後就到第九章.
復活的主向保羅在大馬式的路上向他顯現.
然後到第十章.
就是去到該薩利亞港.
就有羅馬的白夫長哥尼留信主.
所以你看到由七八九十一直繼續下去.
初代教會的故事就告訴我們.
聖靈的降臨和耶穌的榮耀的彰顯.
就不斷在耶路撒冷以外發生.
這裡就告訴我們.
保羅雖然追殺教會.
想阻截教會去擴展.
但是聖靈就反過來.
不斷向地極將他的工作去彰顯.
縱使這位耶穌聖靈是手無寸鐵.
沒有地上的權柄.
沒有地上的方法.
不過這個角度是繼續彰顯的.
我記得因為我現在正在進修.
我記得有一次有一個課堂.
我們一班同學在那裡.
大概十多人.
我的學院裡面大部分的同學都是白人.
因為我在浸信會的神學院進修.
我記得有一次在一個課堂那裡.
大家休息的時候聊什麼呢.
聊教會面對的挑戰.
有一些白人教會的弟兄姊妹.
就在討論裡面說.

$^{481}$為什麼美國白人的教會會有衰竭的象徵.
因為對白人來說.
特別是南方來的白人.
他們都是基督教背景的.
他們覺得自己都是很有地位的.
為什麼我們會衰竭呢.
我就坐在中間.
他們看我當然是唯一一個中國人.
他們就問我一個問題.
為什麼華人基督教的群體.
特別他們自己是中國人.
當然我不是中國人.
我沒有親身的經歷.
但是大家都知道中國有很多基督徒.
比白人多很多.
大家就在問這個問題.
何解呢.
我們有地位.
我們有我們的方法.
但是為什麼我們還是有衰竭的現象.
何解你們沒有我們的地位.
但是你們在中國.
甚至我們說到Global South.
西班牙語的教會.
在巴西等等.
那些教會開始興起.
其實是無他的.
只要教會信徒.
肯去活出一個像耶穌的門徒的生活.
那就足夠了.
教會如何興起.
我們不需要靠地上的權勢.
我們靠誰.
靠耶穌基督.
靠耶穌基督的生命.
教會才能夠繼續向前.
這和保羅所想的東西.
是剛剛相反.
用人的手段或用人的方法.
以至能夠去阻截其他的同胞.

$^{521}$以至能夠自己的派別興起.
這就是講到上帝的國度是如何來的.
不是用人的方法.
耶穌基督十字架的福音.
第二方面.
這裡講到蘇羅的故事.
這裡不單止講到上帝的國度.
是人能夠攔阻.
並且上帝的國度去到伸展的時候.
不是用人的方法.
是用上帝福音的方法.
第二方面講到上帝的恩典.
臨到這個人.
蘇羅或保羅.
你看史萊恩傳的脈絡.
一直講到蘇羅如何信主的時候.
都是講到一些.
猶太人眼中不應該得到上帝的國度的人.
他們信主的.
剛才那個脈絡讓你看到的是什麼.
撒瑪利亞人.
埃塞俄比亞的太監.
愛邦人.
哥利林.
羅馬的白夫長.
這些全部都是猶太人眼中覺得比較賤的人.
他們不應該進入我們最神聖的地方.
第九章這裡夾住的.
就是講到保羅自己.
他自以為自己.
是應該傳承到上帝的國度.
上帝的國度應該是給他的.
甚至他可以用任何手段.
有訴諸於暴力.
將上帝的國帶來.
但是神就向他顯現.
然後將他擊倒.
我們展示了這幅畫好幾次.
雖然我不是讀藝術的.
但是這個藝術家我覺得挺有意思.

$^{561}$十七世紀的一幅畫.
他就去刻畫保羅信主在大馬式的路上.
上帝的榮耀向他顯現的時候.
有些什麼事發生.
你見到他從馬上掉下來.
這個畫家用光暗對彼得很重要.
黑與白 光與暗.
一個很強烈的對立.
卡佛法治亞應該想表達的是什麼呢.
一個很自以為是的保羅.
本身是站在馬上.
但是榮耀的基督向他顯現的時候.
就將他擊倒.
他就睡在地上.
然後他就是一個弱者.
這幅畫最脆弱的是誰.
就是睡在那裡的訴羅.
當上帝向他顯現的時候.
他就明白.
上帝在恩典將他擊倒.
他不是靠著自己的努力.
靠著自己的手段.
將上帝的榮耀帶來.
而是上帝在恩典裡面向他顯現.
你想想你信主的經歷其實都是一樣.
所以這裡有一節的經文第四節.
平時我都沒有留意.
馬上就倒在地上.
有光向他顯現.
基督向他生命彰顯的時候.
他就跌倒.
不再是以前那樣.
你看看其他的經文.
保羅講述他信主的經歷.
其實都是有這樣的意思.
譬如腓立比書第三章.
那裡講什麼呢.
他憑著肉體.
憑著他自己出生.
或者他自己本身有的與生俱來的條件.

$^{601}$他是一個很愛神的人.
他是一個正正式式的猶太人.
出生之後第八天受國禮.
然後以色列族的.
變阿們之牌.
所生的希伯來人.
就律法說我是法利賽人.
就熱心說我是迫害教會.
在律法的義上是無可指責的.
但是因為我被基督的恩典擊倒之後.
以前我覺得以為是有益的.
因為基督的緣故.
已經當為有損了.
提摩太前書第一章第十三節.
我從前是竊族神.
這個是他晚年的時候寫的.
迫害人的.
灰滿人的.
然而我還忘了上帝的恩典.
說到基督的復活的時候.
他說基督的復活.
第三天復活.
最後他向誰顯現.
基督向他顯現.
他如同一個未到產期而生的人一樣.
什麼叫未到產期而生的人.
就是born premature.
就是他在生出來的.
本身是應該死的人.
不過上帝憐憫他.
讓他死不去.
這裡就告訴我們.
上帝的恩典.
在一個人身上彰顯的時候.
一個人應該有的反應.
你記得我開頭的時候問.
你信主的經歷是怎樣的.
請你不要忘記.
是上帝的恩典.
去憐憫我們.

$^{641}$最後兩個反思的問題.
問我們自己.
是不是好像昔日的素羅和保羅一樣.
我們在左投左使.
以為別人左投左使.
其實是他自己左投左使.
我自己有沒有阻攔了上帝的恩典.
當時我們說到背景的時候.
素羅就是保羅.
是一個滿有激進和熱心的人.
反而成為一個怎樣的人.
去迫害上帝各道子民的人.
因為驕傲.
因為他的自然.
其實他是在抗拒上帝.
十誡的福音上帝向他顯現.
就是將他激怒.
要他謙卑下來.
我們相反問問我們自己.
我們是不是在攔阻上帝的恩典.
可能是我們的自意.
或者我們自己的自滿.
我們是在抗拒上帝.
有如未信主的保羅一樣.
第二個問題.
我自己應該怎樣去見證上帝的福音.
使徒函傳從頭到尾告訴我們.
你憑著聖靈的能力.
你活出基督的生命.
就能夠成為我的見證.
無論你在耶路撒冷還是在地球.
教會,聖靈的工作.
從來不需要靠兵,靠槍.
我們要靠的是什麼.
不是武力,不是勢力.
是聖靈的能力.
是上帝的子民能夠謙卑在聖靈面前.
我們求主幫助我.
成為一個更加實證的基督徒.
求主幫助我活出十字架的福音.

$^{681}$我不知道這是不是你每一天的祈禱.
每一天的祈禱.
求主幫助我.
做一個好的兒女,好的父母.
好的丈夫,好的妻子.
怎樣以愛基督的恩典.
去對待我身邊的人.
做一個好的員工.
以真理,以公義,以和平.
上帝可以這樣.
或者上帝應該只會這樣.
去擴張自己的國度.
十字架的福音是靠我們.
謙卑的去生活出來.
這樣上帝的福音才能夠擴展.
我們同心低頭祈禱好嗎.
願以這個都成為我們的禱告.
主,求你謙卑我們.
甚至用你的愛,用你的恩典來激導我們.
使我們不會破壞了你自己的工作.
主,將我們的驕傲,將我們的恣意.
全部剷除.
帶領我們來到十字架的面前.
教我們謙卑.
主,求你幫助我們的生命.
不是靠我們自己的強權.
或者我們的恣意.
來活出福音的生命.
可以走近十架.
知道你自己是這樣影響世界.
在十字架上願意為我們死.
願意甘心的來救贖我們.
主,讓我們有罪基督.
基督的心.
讓我們的生命能夠活出你自己的福音.
主,願意這是我們這個星期每天的禱告.
主,求你幫助我們.
做一個好的妻子,做一個好的丈夫.
做一個好的父母,做一個好的鄰舍.
幫助我們.

$^{721}$也在教會做一個好的弟兄或姊妹.
讓我們的生命能夠表露出.
耶穌基督的自己的生命.
讓十字架的福音,恩典.
能夠在我們這個教會的生命當中.
能夠去表達出來.
主,使用我們.
祈禱奉出明教.
報完信祝福各位.
在這個星期裡活出基督的生命.
主,願以我們的福音.
祈禱奉出明教.
主,願以我們的福音.
祈禱奉出明教.
主,願以我們的福音.
祈禱奉出明教.
主,願以我們的福音.
祈禱奉出明教.
\newpage



\section{路加福音 8:26-39}
\label{sec:lW1Rt3jzEUM}
\textbf{Who is your enemy (Luke 8\_26-39) - Tsz Kin Leung 誰是敵人 (路加福音 8: 26-39) - 梁紫健傳道}
\newline
\newline
連結: \href{https://youtube.com/watch?v=lW1Rt3jzEUM}{\texttt{https://youtube.com/watch?v=lW1Rt3jzEUM}} ~~~~ 語音日期: 2025-09-23
\newline
\newline
\hyperref[sec:dx4StTt36lE]{\small{< < < PREV SERMON < < <}}
~
\hyperref[sec:index_chronic]{\small{[返順時目]}}
~
\hyperref[sec:index_scriptual]{\small{[返順卷目]}}
~
\hyperref[sec:6VvjU9_F7kQ]{\small{> > > NEXT SERMON > > >}}
\newline
\newline
路加福音 8:26-39
\newline
\begin{longtable}{cl}
\hline
\hline
章節 & 經文 (和合本修訂版)\\
\hline
8:26 & \begin{tabularx}{0.7\textwidth}{X} 他們到了格拉森人的地區，就在加利利的對面。 \end{tabularx} \\ \\ \relax
8:27 & \begin{tabularx}{0.7\textwidth}{X} 耶穌上了岸，就有城裡一個被鬼附的人迎著他走來。這個人好久不穿衣服，不住在屋子裡，而住在墳墓裡。 \end{tabularx} \\ \\ \relax
8:28 & \begin{tabularx}{0.7\textwidth}{X} 他看見耶穌，就喊叫著俯伏在他面前，大聲說：「至高神的兒子耶穌，你為甚麼干擾我？我求你，不要叫我受苦！」 \end{tabularx} \\ \\ \relax
8:29 & \begin{tabularx}{0.7\textwidth}{X} 這是因耶穌曾吩咐污靈從這人身上出來。原來這污靈屢次抓住他；他常被人看守，又被鐵鏈和腳鐐捆鎖，他竟把鎖鏈掙斷，被鬼趕到曠野去。 \end{tabularx} \\ \\ \relax
8:30 & \begin{tabularx}{0.7\textwidth}{X} 耶穌問他：「你的名字叫甚麼？」他說：「群」；這是因為附著他的鬼多。 \end{tabularx} \\ \\ \relax
8:31 & \begin{tabularx}{0.7\textwidth}{X} 鬼就央求耶穌不要命令他們到無底坑裡去。 \end{tabularx} \\ \\ \relax
8:32 & \begin{tabularx}{0.7\textwidth}{X} 那裡有一大群豬正在山坡上吃食，鬼央求耶穌准他們進入豬裡；耶穌准了他們。 \end{tabularx} \\ \\ \relax
8:33 & \begin{tabularx}{0.7\textwidth}{X} 於是鬼從那人出來，進入豬裡，那群豬就闖下山崖，投進湖裡，淹死了。 \end{tabularx} \\ \\ \relax
8:34 & \begin{tabularx}{0.7\textwidth}{X} 放豬的看見這事就逃跑了，去告訴城裡和鄉下的人。 \end{tabularx} \\ \\ \relax
8:35 & \begin{tabularx}{0.7\textwidth}{X} 眾人出來，要看發生了甚麼事；到了耶穌那裡，發現那人坐在耶穌腳前，鬼已離開了他，穿著衣服，神智清醒，他們就害怕。 \end{tabularx} \\ \\ \relax
8:36 & \begin{tabularx}{0.7\textwidth}{X} 看見這事的人把被鬼附的人怎麼得醫治的事告訴他們。 \end{tabularx} \\ \\ \relax
8:37 & \begin{tabularx}{0.7\textwidth}{X} 格拉森周圍地區的人，因為害怕得很，都求耶穌離開他們；耶穌就上船回去了。 \end{tabularx} \\ \\ \relax
8:38 & \begin{tabularx}{0.7\textwidth}{X} 鬼已從身上出去的那人懇求要和耶穌在一起，耶穌卻打發他回去，說： \end{tabularx} \\ \\ \relax
8:39 & \begin{tabularx}{0.7\textwidth}{X} 「你回家去，傳講神為你做了多麼大的事。」他就走遍全城，傳揚耶穌為他做了多麼大的事。 \end{tabularx} \\ \\
[1ex]
\hline
\hline
\end{longtable}
$^{1}$(speaking in foreign language).
- Good morning..
(speaking in foreign language).
(laughing).
(speaking in foreign language).
- I don't know if everyone has played Werewolf.
or different script murders..
This is a Werewolf Bible version created by our church,.
UV Church during the pandemic,.
featuring Satan, later devils,.
a human Adam and Eve, as well as Jesus..
(speaking in foreign language).
The non-biblical version, of course,.
is that werewolves confront gods and humans.
with everyone being enemies of each other..
(speaking in foreign language).
In the biblical version of Werewolf,.
even the bad Samson is a messenger of light,.
and a human's only enemy is the devil..
This is a very interesting transformation..
(speaking in foreign language).
Doesn't that seem a little bit strange to you all?.
Even the bad guys aren't our enemies..
(speaking in foreign language).
Today, we're going to talk about a topic..
Who is your enemy?.
(speaking in foreign language).
Are there some people in your life that you really hate,.
even to the point where you think,.
I never want to see this person again..
I don't want my life to intersect with him or her..
(speaking in foreign language).
The Bible gives us the answer to this question..
(speaking in foreign language).
First, we need to review..
First, we need to review the two stories in, sorry..
First, we need to review the two stories.
in Luke chapter eight,.
as they are in a coherent narratives.
and reflect each other..

$^{41}$These stories take place in the area.
around the Sea of Galilee..
(speaking in foreign language).
Jesus and his disciple came to the region of Gerasenes,.
also known as Gadara,.
where they encountered a man possessed.
by a legion of demons..
The man lived among the tombs.
and he was extremely violent..
No one could subdue him..
(speaking in foreign language).
Jesus asked the name of the demon.
and the demon answered, legion..
The original text here is legion..
We will explain it in detail later..
Because many demon had entered him,.
the demons begged Jesus not to send him into Abyss,.
but to allow them to enter a nearby herd of pigs..
After Jesus permitted it,.
the demons entered the pigs.
and the herd rushed down the steep bank.
into the lake and drowned..
There were about 2,000 pigs in the herd..
The pig herder ran off and reported this to the city.
and the people from the region came to see what had happened..
When they saw Jesus,.
they asked him to leave their region..
The man who had been possessed was restored to normality,.
sitting at the feet of Jesus.
and proclaiming the things that Jesus had done..
(speaking in foreign language).
Before the demon-driving story,.
Jesus and his disciples got on the boat.
to cross to the other side of the lake..
On the way, a great storm arose..
The boat began to take on water..
The disciples, fearful, woke Jesus who was sleeping..
Jesus got up and rebuked the wind and the waves,.
and they immediately calmed down..
The disciples were amazed and asked, "Who is this?.

$^{81}$Even the winds and the water obey him.".
(speaking in foreign language).
These two stories are narratively closely connected..
The coming of the storm happened.
during the boat crossing the lake,.
and the demon-driving occurred on the other shore..
The structure suggests they are related symbolic events.
used to convey deeper messages..
(speaking in foreign language).
In the demon-driving story,.
the demon calls itself Legion,.
which is not a random name,.
but has historical and cultural symbolism..
(speaking in foreign language).
In the Roman Empire era,.
Legion was the standard name for a Roman military unit..
A Roman Legion typically had about 5,000 to 6,000 soldiers,.
symbolizing the military power and rule of Roman Empire..
At that time, the Jews, including the Galilee region,.
were under Roman occupation,.
and the Roman Legion were symbols.
of oppression and violence..
(speaking in foreign language).
The demon naming itself Legion.
implies that these spiritual forces.
are as vast, organized, and destructive as Roman Legion..
They process people, causing chaos and suffering,.
similar to Roman armies' oppression of the Jews..
Jesus driving out the demons.
is not confronting human armies,.
but directly confronting the spiritual Legions..
(speaking in foreign language).
This emphasizes the story's focus..
The true enemy is not visible human forces,.
like the Romans,.
but the hidden spiritual evil forces behind them,.
the devil..
(speaking in foreign language).
Now, we will compare these two stories.
with Exodus chapter 14 in the Old Testament..

$^{121}$We will understand the author's intention.
more comparatively..
Exodus described Moses leading the Israelites out of Egypt,.
escaping Pharaoh's pursuit, which is a classic liberation..
Narrative in Jewish tradition,.
the stories in Luke seem to intentionally respond.
to all retell this Old Testament events,.
but with the key transformations..
(speaking in foreign language).
The Israelites across the Red Sea,.
Moses stretched out his staff to part the sea waters,.
allowing the Israelites to pass on dry ground..
Then the water closed,.
drowning the pursuing Egyptian army..
(speaking in foreign language).
God used a strong east wind to part the Red Sea..
The wind was a tool of God's intervention..
(speaking in foreign language).
The Egyptian army, Pharaoh's chariots and horsemen,.
all drowned in the Red Sea, Exodus 14, 28..
Here, the enemy is a human army,.
representing the Egyptian empire that oppressed Israel..
What died here were person..
(speaking in foreign language).
The two stories in Luke.
also have the same parallel elements of the Exodus..
(speaking in foreign language).
Jesus and his disciples crossed the Sea of Galilee,.
the lake, similar to crossing the Red Sea..
(speaking in foreign language).
In the calming the storm story,.
Jesus rebuked the wind and the waves,.
making them calm, the response to the wind in Exodus..
Jesus directly commenced natural forces..
(speaking in foreign language).
In the demon-driving story,.
the demon, Lagion, entered a herd of pigs..
The pigs rushed into the lake and drowned,.
similar to the Egyptian army drowning..
But here, what died were pigs, not people..

$^{161}$The destruction of the herd.
symbolizes the failure of spiritual enemies..
(speaking in foreign language).
This parallelism is not a coincidence,.
but a typology technique.
commonly used by New Testament authors,.
portraying Jesus as a new Moses,.
but bringing deeper spiritual liberation..
(speaking in foreign language).
Through comparison, we can see that Luke's stories.
intentionally inherited elements from Exodus.
to convey deeper theological messages..
The enemy is devil, not humans..
(speaking in foreign language).
In Exodus, what drowned was the army, people..
In demon-driving story, what drowned were pigs, animals..
And the pigs were carrier of the demons..
The demons, the devil, are the true enemy,.
organized like a Lagion, but not human..
The death of the pigs emphasizes.
what is destroyed are non-human entities,.
symbolizing the failure of spiritual forces..
(speaking in foreign language).
Jesus did not cause any human death..
Even the local, the pigs' owner, were afraid but unharmed..
This contrasts with the complete annihilation.
of the army in Exodus,.
showing that Jesus' authority.
is targeted at the spiritual realm, not humans..
(speaking in foreign language).
So our conclusion is the animal of humans is devil,.
not humans..
(speaking in foreign language).
The true stories in Luke highlights.
humans are not the true enemy..
The empire, through the symbolism of Lagion,.
seems like the oppressor,.
but behind it is the devil manipulating..
The devil hides like a Lagion,.
behind oppression, chaos, and suffering,.

$^{201}$manipulating the visible world..
But Jesus' authority directly targets the devil,.
bringing liberation without needing to destroy humans..
It's not launching a war against humans..
(speaking in foreign language).
The true enemy of humans is the spiritual devil,.
not other humans..
This encourages us to shift our focus.
from interpersonal conflicts to spiritual warfare,.
promoting peace and reconciliation, not hatred..
(speaking in foreign language).
- I once was also regarded some people as my enemies,.
with a lot of hatred towards them,.
especially towards my dad..
When I was young, my dad encountered business difficulties,.
causing my family to go through a very hard time..
(speaking in foreign language).
(gentle music).
- He himself suffered a great blow,.
and at the same time,.
he was very unwilling to seek help from others..
When at that time I was a five, six years old kid,.
in my heart, I kept thinking,.
why couldn't he be a boss at that time.
and refuse to find a job,.
letting us go hungry,.
even without money for medical treatment when sick?.
Why make us suffer like that?.
It was later, under the Lord's guidance,.
that I mustered the courage to ask him.
about what had happened back then..
(speaking in foreign language).
Of course, I asked him in a very roundabout way..
Only then did I learn that at that time,.
he was also trapped in a great regret,.
having messed up the business and couldn't find a job..
In the end, he had to go sell blood.
to get a little money to support the family..
I slowly began to understand my dad's state at that time,.
and resolve some misunderstandings..

$^{241}$(speaking in foreign language).
Do you know that the devil really wants us humans.
to be enemies with each other?.
Misunderstandings between humans.
can cause us to hurt each other, to be enemies..
Jesus' demon driving like this is to let us know.
that our enemy is absolutely not humans, but the devil..
(speaking in foreign language).
In your life, are they enemies?.
Some people you always remember him or her,.
remember the bad things that he or her had done to you..
Remember the harm he or she caused you..
Let us together pray to the Lord..
Help us, heal us, recognize that our enemy.
is the only devil, not people..
(speaking in foreign language).
Let us pray together..
(speaking in foreign language).
Lord, help us recognize the enemy in our lives.
is the devil, not the people around us..
(speaking in foreign language).
Not even those who had hurt us before..
(speaking in foreign language).
Lord, we ask you to help us and heal us.
so that we can see this truth clearly.
and through your love, take action in the Lord's name..
(speaking in foreign language).
Amen. - Amen..
(speaking in foreign language).
謝謝大家收看!.
\newpage



\section{使徒行傳 9:10-19}
\label{sec:6VvjU9_F7kQ}
\textbf{歡迎從不易 (使徒行傳 9\_10-19) - 陳崇基牧師}
\newline
\newline
連結: \href{https://youtube.com/watch?v=6VvjU9_F7kQ}{\texttt{https://youtube.com/watch?v=6VvjU9\_F7kQ}} ~~~~ 語音日期: 2025-09-23
\newline
\newline
\hyperref[sec:lW1Rt3jzEUM]{\small{< < < PREV SERMON < < <}}
~
\hyperref[sec:index_chronic]{\small{[返順時目]}}
~
\hyperref[sec:index_scriptual]{\small{[返順卷目]}}
~
\hyperref[sec:7m42sBmfTkI]{\small{> > > NEXT SERMON > > >}}
\newline
\newline
使徒行傳 9:10-19
\newline
\begin{longtable}{cl}
\hline
\hline
章節 & 經文 (和合本修訂版)\\
\hline
9:10 & \begin{tabularx}{0.7\textwidth}{X} 那時，在大馬士革有一個門徒，名叫亞拿尼亞。主在異象中對他說：「亞拿尼亞！」他說：「主啊，我在這裡。」 \end{tabularx} \\ \\ \relax
9:11 & \begin{tabularx}{0.7\textwidth}{X} 主對他說：「起來！往那叫直街的路去，在猶大的家裡，去找一個大數人，名叫掃羅；他正在禱告， \end{tabularx} \\ \\ \relax
9:12 & \begin{tabularx}{0.7\textwidth}{X} 在異象中看見了一個人，名叫亞拿尼亞，進來為他按手，讓他能再看得見。」 \end{tabularx} \\ \\ \relax
9:13 & \begin{tabularx}{0.7\textwidth}{X} 亞拿尼亞回答：「主啊，我聽見許多人講到這個人，說他怎樣在耶路撒冷多多苦待你的聖徒， \end{tabularx} \\ \\ \relax
9:14 & \begin{tabularx}{0.7\textwidth}{X} 並且他在這裡有從祭司長得來的權柄，要捆綁一切求告你名的人。」 \end{tabularx} \\ \\ \relax
9:15 & \begin{tabularx}{0.7\textwidth}{X} 主對他說：「你只管去。他是我所揀選的器皿，要在外邦人、君王和以色列人面前宣揚我的名。 \end{tabularx} \\ \\ \relax
9:16 & \begin{tabularx}{0.7\textwidth}{X} 我也要指示他，為我的名必須受許多的苦難。」 \end{tabularx} \\ \\ \relax
9:17 & \begin{tabularx}{0.7\textwidth}{X} 亞拿尼亞就去了，進入那家，把手按在掃羅身上，說：「掃羅弟兄，在你來的路上向你顯現的主，就是耶穌，打發我來，叫你能再看得見，又被聖靈充滿。」 \end{tabularx} \\ \\ \relax
9:18 & \begin{tabularx}{0.7\textwidth}{X} 掃羅的眼睛上立刻好像有鱗一般的東西掉下來，他就能再看得見，於是他起來，受了洗， \end{tabularx} \\ \\ \relax
9:19 & \begin{tabularx}{0.7\textwidth}{X} 吃過飯體力就恢復了。 \end{tabularx} \\ \\
[1ex]
\hline
\hline
\end{longtable}
$^{1}$各位聽眾,每個早晨.
這次我們繼續看《跨越》這個系列.
我們特別看到使徒行傳.
講到初期教會的一些故事.
路加修習這些故事的時候.
想告訴每一代的.
世代的教會,包括我們.
到底有什麼需要我們注意.
以及如何在聖靈降臨之後.
我們有福音的生命能夠活出來.
所以我們今天繼續看第九章.
使徒行傳第九章.
講到掃羅.
他的歸主的故事.
在大馬色的路上.
他去迫害教會.
但在路上榮耀的基督.
榮耀的恩典.
向他顯現,將他擊倒.
你記得在第九章.
講到他被擊倒之後.
這位強人的保羅.
終於認輸了.
可以這麼說.
然後現在有人拉著他的手.
看不到東西,瞎了眼.
然後領他進大馬士革.
三天不吃不喝.
對猶太人來說.
這是一個非常嚴峻的禁食的習慣.
就是他悔改.
再一次尋求神的經歷.
第九章交代了掃羅如何信主之後.
有另一個故事講到.
掃羅如何被接納.
被歡迎進入教會.
無論如何,這裡講到掃羅.
他的信主經歷在大馬色的路上.
我相信這裡講到.
一個很重要的.

$^{41}$歸主的經歷.
信主的經歷,我不知道大家是否記得.
你信主的時候是怎樣的呢?.
有些什麼經歷呢?.
內心經過一些什麼的.
的開放等等.
掃羅這裡記載了.
他的經歷,他的信主的經過.
這裡有幾個特色.
讓我們看到真正信主的人.
他們需要有什麼特點呢?.
我問過大家.
信主的時候有什麼經歷.
無論是多少年前也好,還是剛剛信主.
到底你的經歷是什麼呢?.
每個人的經歷都不同.
有些人可能哭三天三夜.
有些人可能很喜樂很開心.
三天三夜也好.
經歷不是最重要.
在這個過程裡面.
每一個人和主相遇之後.
有什麼特點呢?.
在掃羅歸主的經歷裡面.
起碼有兩件事值得我們留意.
並且這兩個元素.
是每一個信主的人.
無論他將來之後的表現.
是怎樣也好.
是怎樣不同也好.
起碼有一樣東西本質上是一樣的.
第一就是每一個人.
親自與主相遇.
有一個個人的encounter.
與主耶穌基督.
做一個個人的遇上.
無論是藉著禱告.
或者心裡面.
他親自去承認.
主是他生命裡面的領導.

$^{81}$這是第一樣.
第二樣就是每一個人.
信主應該有的經歷.
起碼在根底上他決定.
來回歸.
以前走過的路.
以前走的方向可能是完全不理會上帝.
甚至是抗拒上帝.
甚至像素羅一樣.
保羅一樣.
他是迫害教會的.
但是經過了.
與主相遇之後.
他就甘願地來回歸.
向耶穌基督.
這個故事.
交代了素羅怎樣在大馬士革路上.
信主之後.
就繼續有一個延續.
故事是怎樣發展下去的呢.
第九章第十節.
他說在大馬士革路上.
有一個門徒叫做阿拉利亞.
主在意象裡面對他說.
阿拉利亞.
主啊我在這裡.
主對他說.
往那叫直街的路上.
在猶大的家裡面找一個大數人.
名叫素羅.
他正在禱告.
在意象裡面他見到一個人名叫阿拉利亞.
就是他自己.
進來為他按手.
讓他能在看得見.
阿拉利亞就回答.
主啊我聽見很多人講論這個人.
說他曾經在耶路撒冷多多.
苦害你的信徒.
並且他在這裡.

$^{121}$有祭司獎得來的權柄.
要捆綁一切.
求告你名的人.
這裡首先講清楚.
阿拉利亞.
你記得有這個名字曾經出現過.
在第五章.
講到一對夫婦.
他們在奉獻上不誠實.
然後被擊打.
丈夫就是阿拉利亞.
跟這個名字一樣.
還有他太太.
我提過大家很容易記的.
是薩菲拉和塞爾.
所以Ananias and Sapphire.
他們到希控聖靈.
這是第五章的故事.
這裡第九章.
阿拉利亞同名但是不同人.
千萬不要搞錯.
以為他死了復活等等.
是另外一個人.
Ananias和阿拉利亞同名但是不同人.
這裡怎麼形容他呢?.
這裡說他是一個門徒.
主的門徒.
相信主的.
並且在其他士堂傳的章節.
甚至保羅重新.
再憶述.
大馬適路上歸主的故事.
信主的故事.
提到阿拉利亞是一個嚴守.
律法的猶太人.
如果你想想.
他有這樣的特點.
像保羅.
像蘇羅.
跟他的差別不是很大.

$^{161}$但是最重要的分別在哪裡呢?.
就是阿拉利亞.
已經信了耶穌.
為基督為尼泊爾的猶太人.
並且可能.
可能會成為保羅.
追殺的對象.
你記得保羅去大馬適的時候.
他想做什麼?.
迫害教會然後將這些基督徒.
收監.
當時也有.
所謂的通風報信.
情報.
這件事.
大家要知道.
所以當時初期的教會.
信徒也會有.
大家通風報信.
阿拉利亞在大馬適路上.
猶太人的基督徒已經收到.
情報了.
知道保羅.
去迫害教會.
所以你可以看到.
他的反應是可以理解的.
因為如果.
保羅是假的話.
那就很大件事了.
他怎麼能夠接待他呢?.
如果他接待他的話.
他就要做一個很大的冒險.
所以他的反應.
是可以理解的.
但是這個故事繼續交代的時候.
第九章第十五節.
主就在耳目裡面跟他說什麼呢?.
他叫阿拉利亞.
只管去.
因為這位的蘇羅.

$^{201}$保羅.
是我所揀選的器皿.
是在愛邦人君王以色列人.
面前宣揚我的名字.
我也指示了他.
為我的名他必須要.
受許多的苦難.
十七節.
阿拉利亞就去了.
進入了猶太的家.
將他的手按在蘇羅的身上.
他說蘇羅弟兄.
我來的路上.
向你顯現的主就是耶穌.
打發我來.
叫你能夠再看得見.
又被聖靈充滿.
蘇羅的眼睛上.
就好像立刻有輪一般的東西掉下來.
他就能再看得見.
於是他就起來了.
受了洗.
吃過飯之後.
他的體力就恢復了.
這個由開始有一點點的抗拒.
不想去歡迎保羅.
阿拉利亞就遵命.
是一個嘗試去冒險的馴服.
一個可能會是間諜的人.
但因為主對他的吩咐.
因為在主裡面.
他就當保羅.
是可以接待的鄰舍.
甚至他怎樣稱呼他.
第一句說話.
他說蘇羅弟兄.
brother.
你和我是一樣的.
在主裡面在基督裡面在福音裡面.
我現在來接待你.

$^{241}$所以就將保羅.
接待了.
進入教會.
於是第九章.
第21節他就去交代.
蘇羅.
進入了教會的時候.
有什麼事情.
第九章21節.
蘇羅在大馬士革的門徒.
一起住了好些日子.
立刻在各個猶太會堂.
傳揚耶穌.
凡聽見的人都非常驚奇.
他們說什麼.
他說這不是去迫害教會.
去迫敗基督徒嗎.
但是蘇羅越發有能力.
駁倒大馬士革的猶太人.
證明耶穌是基督.
過了好些日子.
23節.
猶太人相議要殺蘇羅.
但他們的計謀被蘇羅自己知道了.
他們晝夜在城門等候.
要去殺他.
於是乎.
他的門徒.
蘇羅的門徒.
就在夜間用一個籃子.
將他放下.
然後救他.
讓他能夠離開險境.
這裡.
可以看到.
這位蘇羅可能去到傳道的時候.
已經有些人跟了他.
信奉了保羅的福音.
即是耶穌基督的福音.
所以有他的門徒在夜間.

$^{281}$救他脫險.
大馬士革的直街.
剛才說到阿拉利亞.
要去到大馬士革的城裡面.
直街一個猶太家.
去找蘇羅.
按手臂.
直街自古以來都存在.
今時今日去到大馬士革.
就是一個現代的直街.
應該是按照著.
二千多年前.
直街.
自古以來有的路線.
在直街附近.
有一間東正教堂.
東正教堂是一個小的庭園.
那個庭園.
當然那個教堂叫阿拉利亞教堂.
裡面有一個.
庭園裡面有一個很簡單的雕像.
就是紀念.
阿拉利亞信服.
冒險來接待保羅.
然後將手按在.
他的頭上.
你看到這個雕像有兩個人.
就是保羅.
這個故事可以算是.
在初期教會裡面.
一個最驚心動魄.
一個最冒險的故事.
就是必須要.
接待保羅.
這個故事告訴我們.
一個人信主.
他必定需要.
有教會.
其他的群體.
來幫助他.

$^{321}$沒有一個人信主.
是多麼強人也好.
像素羅一樣.
是多麼強大多麼有能力.
沒有一個人.
信主.
是可以孤軍作戰的.
每一個人都需要其他人.
好像他需要其他的門徒.
來使他脫險.
需要其他的群體將他接納進來.
這才是一個整串的.
信主的經歷.
早前有一部很出名的電影.
叫做Top Gun.
大家可能有看過.
Top Gun這部電影.
有另一個名字.
叫做Maverick.
什麼意思呢?.
就是獨行俠.
為什麼叫獨行俠呢?.
因為Tom Cruise的角色.
是一個很有經驗和很有能力的戰機師.
他不需要其他人.
他自己孤軍作戰也可以.
但整部電影到最後.
告訴我們什麼?.
最後都需要他的學生.
要其他人來救他.
他才能撐起自己的命.
所以沒有一個人是獨行俠.
基督徒更加不可以是一個獨行俠.
所以為什麼我們需要一個教會的群體.
一個信仰的群體來幫助我們.
我們講過.
保羅.
他本身是一個超正統的猶太拉比.
一個超正統的猶太拉比.
當他遇見基督之後.

$^{361}$他就看不見了.
然後在馬上掉下來.
被擊倒.
相信他信主的時候.
是一個很混亂的時間.
我不知道你信主的時候.
是不是有一點混亂.
剛剛信主的時候.
都不知道搞什麼科.
來到教會原來要做這些事.
原來聖經是這樣的.
可能到你今天信主幾年.
這麼多年.
你打開聖經都很混亂.
每個信主的時候.
保羅其實都是一樣.
是需要別人幫他.
阿利亞.
你都可以接待他.
然後有其他門徒來幫助他.
每個人信主的經歷都是一樣.
沒有人是獨來獨往.
是一定需要其他人來幫他.
我記得我年輕的時候.
不是很會禱告.
不知道禱告是一回什麼事.
你去到教會.
跟別人一起祈禱.
或者是祈禱會.
跟別人一起祈禱的時候.
他們叫你做什麼.
低頭祈禱.
什麼叫低頭祈禱.
就是一輩子都沒低過頭.
或者叫你閉上眼睛祈禱.
突然閉上眼睛祈禱.
你就很集中不到精神.
開眼睛也集中不到精神.
那怎麼祈禱呢.
所以有一點混亂.

$^{401}$有一個姐姐.
教會裡的一個姐姐.
就告訴我.
她以前信主的時候.
她的導師就教她.
怎麼開始懂得祈禱.
就在你自己房間.
或者你自己空間.
前面就放一張凳子.
空凳子.
我開始聽見這個.
所謂的巧妙.
一個提示的時候.
我就覺得很有意思.
讓你具體一點.
去想像到.
或者觸碰到.
原來神是跟你一起.
跟你好像在聊天一樣.
你見到年輕的時候.
是很混亂的.
可能你今天也是一樣.
很混亂.
怎麼祈禱呢.
怎麼能夠具體一點去禱告呢.
其實是需要人去幫助.
需要人去提點.
我記得有一次大學的時候.
有一段時間.
我很想看聖經.
我就記得有一個星期六.
我早上的時候.
大概十點十一點.
我自己在家裡.
他說不如看看聖經.
不知道看什麼.
因為聖經這麼厚.
在哪裡開始看.
從何說起呢.
我就打開小行傳.

$^{441}$我整本小行傳從頭到尾看完.
第一章到第28章.
你猜我明白多少.
我看完之後都不知道看什麼.
中文不是差.
但總之看完之後.
就不太明白他說什麼.
這些時間就是最混亂的.
但需要什麼人去指點呢.
導師啊 需要牧者啊.
需要其他人去幫助.
讓我走這條路的時候.
我能夠有所方向.
所以我們每一個都不是獨行俠.
就算.
說到我們基督徒的品格.
我們剛剛信主.
一路開始.
我們所謂成長或者向前的時候.
我們都會鍛鍊到不同的品格.
譬如忍耐.
忍耐是怎麼鍛鍊的.
人氣你才能鍛鍊忍耐.
怎樣能夠有愛心.
有個人在你面前你很憎他.
這樣才叫有愛心.
有愛心不是說大家一起玩.
我們很有愛心.
那些不是愛心.
那些只不過是同聲同氣.
但配搭回來.
和其他人來去.
我們才能夠鍛鍊出來.
所以是需要其他人.
需要其他實體的人.
去告訴我們.
這個信仰.
到底是怎麼走下去.
大家剛才說了.
Bella剛剛回來.

$^{481}$十幾天了.
前晚回來.
但回來的時間不久.
在芝加哥的時間不久.
她做spiritual direction.
我們都需要有.
若干的spiritual direction.
如果不是.
我們就盲中中.
浪費很多腳骨力.
怎樣能夠使我們成長.
就是一個群體.
而史蘭傳就呈現給我們一個群體.
就是一個歡迎的群體.
一個歡迎異己的群體.
在基督裡面能夠.
一起.
能夠去經歷.
9章26節.
這個故事繼續發展.
就是初信他在保羅萊人的經歷.
當然他在大馬士革一段時間.
然後他就.
回去耶路撒冷.
這件事就大件事.
怎樣回去耶路撒冷.
耶路撒冷是什麼地方.
對於當時的教會來說.
是迫.
迫伯的開始.
他居然走回去耶路撒冷.
還有面子回去.
他想做什麼.
想與門徒結交.
大家是不是歡迎他.
大家不是歡迎他.
這個可以understand.
當然怕了他.
不信他.
是門徒.

$^{521}$第27節.
只有巴拿巴去接待他.
領他去見使徒.
比德,約翰雅各.
教會的領袖.
巴拿巴在路上怎樣看見主.
怎樣向他說話.
怎樣傳揚耶穌放膽的傳道.
都述說給使徒聽.
然後28節.
於是素羅就在耶路撒冷.
與門徒出入來往.
一起同行.
奉主的名放膽傳道.
並且.
與說希臘話的.
猶太人.
講論辯駁.
他們卻沒法子要殺他.
想盡法子要殺他.
弟兄們知道.
然後就帶他脫險.
帶他去到海邊的.
一個海港叫蓋薩尼亞.
送他去大掃.
因為他本身是大掃人.
回到老家鄉.
因為太危險了.
那時各處的教會都得平安建立起來.
凡事.
去到敬畏主.
蒙聖靈的安慰人數漸漸增加.
這個很有趣.
他回到耶路撒冷.
然後那裡的門徒.
不願意去接受他.
但是有一個人.
去到所謂作擔保的.
那個就是巴拿巴.
你記得巴拿巴是不是出現過.

$^{561}$記得嗎.
第四章.
那個巴拿巴和這個巴拿巴.
OK.
剛才那個.
什麼.
安納納伊斯.
這個和一個聖巴拿巴.
同一個巴拿巴.
他就去作一個擔保.
願意帶他去見使徒.
下面那裡有一句.
是加拿大書第一章第十八節.
整章加拿大書第一章.
就是保羅說他.
在大馬式的路上.
怎樣信主.
有什麼經歷等等.
他說過了三年之後才回到耶路撒冷.
他在大馬式甲.
敘利亞等等的地方.
他在那裡傳道.
過了三年.
他才所謂九章二十六節.
回到耶路撒冷.
那裡保羅怎樣益述自己.
他說過了三年.
我才上耶路撒冷去見.
基法.
基法是彼得.
見這個字.
不是說.
看一看他.
然後走.
見是接見的意思.
或者是查詢的意思.
你見到原文.
是和history.
是差不多的字.
是從這個字來的.

$^{601}$我們今天說歷史這個字.
當然今天我們說.
history.
就不同當時的historial.
如果是將現代的意義.
塞進去以前的經文.
就是什麼.
就是解經錯謀.
你來我課程.
如果你一頭霧水的話.
請你來.
因為你不能將現代的history.
他就是回去找基法.
然後就想history一餐.
不是這個意思.
historial.
其實是查訪的意思.
有一點.
是詢問的意思.
你去到一個地方.
你不太清楚.
你去遊客中心.
或者是歡迎中心.
去查訪一下.
詢問一下.
這裡保錄.
三年之後他上耶路撒冷.
再去查訪他所信的福音.
當然歸根究底.
歸根究底的時候.
就明白他的信仰.
明白他的福音.
就是和使徒所傳的福音.
完全一樣.
這是加拉太書第一章所說.
也是使徒第九章二十六節所交代給我們的.
上到耶路撒冷.
他想和其他門徒結交.
但其他門徒很怕他.
因為怕他是間諜.

$^{641}$只有巴拿巴.
之前說他是一個很慷慨的人.
他不認識的.
因為其他人都拒絕他.
但是.
他將保羅曾經做過的事.
全部和使徒當時的教會領袖.
來到去談一談.
然後告訴他.
我們可以接納這個人.
因為已經有三年時間.
證明了這個人.
是真正是我們的門徒.
是我們一份子幫助我們教會.
你看到這個故事.
是一個在基督裡面願意去接納.
去包容一些本身和我們不同的人.
這個才是真正的教會.
如果看使徒行傳的話.
其實一直.
我們說過使徒行傳的其中一個主角是誰.
不是保羅.
不是彼得.
不是那些大的個人.
那些所謂著名的個人.
是教會.
教會在整個使徒行傳的敘述裡面.
是佔有一個很重要的角色.
是一個角色.
這個群體是做什麼的.
無論是耶路撒冷還是其他地方.
你就會看到他們不斷地去學習.
去接納.
接納什麼呢?第二章你看到.
在五旬節教會開播的時候.
就有天下各國來的講不同語言的猶太人.
然後他們在耶穌裡面得到聖靈.
然後他們信主.
第八章.
福音傳到去撒瑪利亞人.

$^{681}$當時猶太人覺得是一些.
不可以.
接納的一些淪舍.
是我們藐視的.
他們願意去將福音傳給他.
然後他們成為教會的一部分.
第八章尾.
就是埃塞俄比亞的太監.
一個猶太人的漢衛.
身體有殘疾的人.
就是看不起的人.
令眼相看的那幫人.
將他接納下來.
這個故事很有趣.
太監信主.
今天的話就是公公信主.
你想想.
如果公公是信主.
坐在你旁邊.
你會不會舒服.
你會舒服.
你天天天天看著他.
你就覺得這個人和我們不同.
非常不同.
聲音又不同.
地位又不同.
但他卻被歸納了.
接納了.
成為基督的子民.
第九章.
就是我們剛剛說的這個故事.
第十章至十一章.
就是羅馬人.
猶太人的死對頭.
羅馬的官.
羅馬的士兵.
就是哥尼林.
十五章他們開會議.
要討論猶太人和外邦人.
信主到底有什麼條件等等.

$^{721}$開了一個大會.
有很多辯論在裡面.
到最後他們的結論是什麼.
他們的結果.
就是說我們在基督裡面.
我們就是德國的子民.
然後一路發展下去.
十六章.
就是將福音傳到歐洲.
跨越當時的地理.
跨越當時所有的邊界.
你見到這個教會.
就是一個誤講邊界.
基督裡面我們就是能夠.
United.
這個才是一個正在成長的教會.
我們教會.
我們做基督徒.
今時今日應該怎樣做.
兩天前我在.
當然最近.
世界和美國.
發生很多事.
因為令我們很多不同意見.
很多不同的分化.
我兩天前在一個.
在一個大學.
我兩天前在一個朋友的.
一個白銀牧師的.
他的社交頁上面.
他就說了一句說話.
我覺得很有意思.
特別說美國本土的氣氛.
他說什麼呢.
他說當你生活在文化戰爭的模式.
當你生活在文化戰爭的模式.
你總是有戰爭要戰鬥.
在聖經裡.
總是有一個不愛的人要歡迎.
你活在上帝的福音和國度裡面.

$^{761}$你整個生活的模式.
看人的模式就完全不同.
活在這個福音裡面.
你總是會覺得陌生人.
我們需要以主的愛去接待.
有鄰舍.
我們盡量去親近.
有鄰舍.
我們盡量去親近.
有仇敵.
我們要學習去愛.
有仇敵.
我們要學習去愛.
剛好如果我們帶著這個福音.
The Gospel mode.
不是Wall mode.
怎樣去邀請.
怎樣去接納其他人.
教會就需要學習.
當時的教會.
學習這件事.
有接納的胸襟在主裡面.
有一個恩念去接納其他人.
同樣今天.
教會也需要學習這件事.
你和我.
也需要學習這件事.
如果這個禮拜你帶著.
這個生存的模式.
不是這個.
是這個.
The Kingdom mode.
Stranger.
鄰舍.
仇敵.
我們應該怎樣去面對他們.
教會是需要在這方面.
我們去學習.
在這個禮拜裡面.
帶著這個模式.

$^{801}$走去我們自己的生活.
你想想有多不同.
多美.
想一想.
這個禮拜.
想一想在我們自己的生命裡面.
是不是曾經有些人.
是用一個這樣的模式出現.
這樣的模式去對待你.
去歡迎.
Be friend.
To love.
想想你過往.
怎樣信主.
有其他人投資在你身上.
是不是.
不是和你打仗.
而是在福音裡面.
去歡迎你.
這個禮拜也想想.
我們怎樣踏出.
有沒有其他人和我們不合.
無論在教會裡面.
你想想教會裡面.
其實也有很多.
People who are different.
和我們不同.
Stranger, neighbor.
Even enemy.
有一個福音的.
生活的模式.
好嗎.
這個禮拜想一想.
你心靈裡面有哪個人.
接納過來.
這個禮拜我怎樣去做到.
這個Gospel.
Kingdom.
可以放心低頭.
因為你是那位先來.

$^{841}$向我們接觸.
先以我們不是陌生人.
而是我們朋友.
你去到.
接觸我們.
你讓你自己的恩典.
你自己的福音.
闖進我們生活裡面.
以致我們能夠認識你.
主啊,求你再一次.
讓你自己十字架的恩典.
去激導我們.
讓我們不是自然.
封閉的在這個世界裡.
不是以戰爭的模式.
而是以基督的福音.
基督的真理.
去接待.
去親近.
去關愛其他人.
主讓我們這個禮拜.
我們能夠用這個模式來生活.
讓教會.
能夠在地上見證.
就是基督你自己.
願意捨身的.
十家強悍.
祈禱奉主命.
\newpage



\section{以弗所書 1:3-14}
\label{sec:7m42sBmfTkI}
\textbf{You Are Blessed (Ephesians 1\_3-14) - Paul Zhang 你是有福的人 (以弗所書 1\_3-14) - 張威傳道}
\newline
\newline
連結: \href{https://youtube.com/watch?v=7m42sBmfTkI}{\texttt{https://youtube.com/watch?v=7m42sBmfTkI}} ~~~~ 語音日期: 2025-10-01
\newline
\newline
\hyperref[sec:6VvjU9_F7kQ]{\small{< < < PREV SERMON < < <}}
~
\hyperref[sec:index_chronic]{\small{[返順時目]}}
~
\hyperref[sec:index_scriptual]{\small{[返順卷目]}}
~
\hyperref[sec:9E6MnSM4U0Q]{\small{> > > NEXT SERMON > > >}}
\newline
\newline
以弗所書 1:3-14
\newline
\begin{longtable}{cl}
\hline
\hline
章節 & 經文 (和合本修訂版)\\
\hline
1:3 & \begin{tabularx}{0.7\textwidth}{X} 願頌讚歸給我們主耶穌基督的父神。他在基督裡曾把天上各樣屬靈的福氣賜給我們。 \end{tabularx} \\ \\ \relax
1:4 & \begin{tabularx}{0.7\textwidth}{X} 因為他從創世以前，在基督裡揀選了我們，使我們在他面前成為聖潔，沒有瑕疵，滿有愛心。 \end{tabularx} \\ \\ \relax
1:5 & \begin{tabularx}{0.7\textwidth}{X} 他按著自己旨意所喜悅的，預定我們藉著耶穌基督得兒子的名分， \end{tabularx} \\ \\ \relax
1:6 & \begin{tabularx}{0.7\textwidth}{X} 使他榮耀的恩典得到稱讚；這恩典是他在愛子裡白白賜給我們的。 \end{tabularx} \\ \\ \relax
1:7 & \begin{tabularx}{0.7\textwidth}{X} 我們藉著這愛子的血得蒙救贖，過犯得以赦免，這是照他豐富的恩典， \end{tabularx} \\ \\ \relax
1:8 & \begin{tabularx}{0.7\textwidth}{X} 充充足足地賞給我們的。他以諸般的智慧聰明， \end{tabularx} \\ \\ \relax
1:9 & \begin{tabularx}{0.7\textwidth}{X} 照自己在基督裡所立定的美意，使我們知道他旨意的奧祕， \end{tabularx} \\ \\ \relax
1:10 & \begin{tabularx}{0.7\textwidth}{X} 要照著所安排的，在時機成熟的時候，使天上、地上、一切所有的，都在基督裡面同歸於一。 \end{tabularx} \\ \\ \relax
1:11 & \begin{tabularx}{0.7\textwidth}{X} 我們也在他裡面得了基業；這原是那位隨己意行萬事的神照著自己的旨意所預定的， \end{tabularx} \\ \\ \relax
1:12 & \begin{tabularx}{0.7\textwidth}{X} 為要使我們，這些首先把希望寄託在基督裡的人，頌讚他的榮耀。 \end{tabularx} \\ \\ \relax
1:13 & \begin{tabularx}{0.7\textwidth}{X} 在基督裡你們聽見真理的道，就是那使你們得救的福音，你們也信了他，就受了所應許的聖靈為印記。 \end{tabularx} \\ \\ \relax
1:14 & \begin{tabularx}{0.7\textwidth}{X} 這聖靈是我們得基業的憑據，直等到神的子民得救贖，使他的榮耀得到稱讚。 \end{tabularx} \\ \\
[1ex]
\hline
\hline
\end{longtable}
$^{1}$(speaking in foreign language).
- Good morning..
Good morning, everyone..
This is Paul, and why am I standing here today?.
Because Pastor Mark is in Toronto this weekend..
- White testing..
- Yeah, so Pastor Mark is in Toronto this weekend,.
and you may not know me, but some of you,.
you may know me, but some of you probably don't know.
my background very well..
Many years ago, I came to the US to study theology.
in a large seminary..
(speaking in foreign language).
- During this time, I had a pretty awkward experience..
When I went grocery shopping one day,.
I accidentally sneezed, and an elderly lady besides me.
say, "Bless you.".
(speaking in foreign language).
- I froze on the spot, thinking, what is she saying?.
Why did she say, "Bless you"?.
I had no idea what it meant..
Only later did I learn that in American culture,.
bless you is a common phrase that people say.
when someone sneezes, and it is a polite reflex,.
and we usually just respond with a quick thank you.
without giving much thought..
(speaking in foreign language).
- But why do people say bless you in the first place?.
This phrase dates back to the medieval Europe.
when medicine are primitive, and plagues like the Black Death.
was a deadly threat, and sneezing was often seen.
as a sign of infection, or even superstitiously.
as a sign of evil spirits at work..
(speaking in foreign language).
- In response, the Pope, during that time,.
encouraged people to say, "God bless you,".
as a form of mutual care and encouragement..
(speaking in foreign language).
- But in today's world, when we hear people say,.
"You are blessed," we will just think of it.

$^{41}$as a form of mutual care and encouragement..
(speaking in foreign language).
- We will just think of it as a regular phrase..
(speaking in foreign language).
- When you hear people say, "You are blessed,".
how would you feel?.
(speaking in foreign language).
- So you might be thinking, well, my life actually feels.
pretty miserable, I don't feel blessed at all..
But our passage today in Ephesians will remind us.
of this truth, that you are a truly blessed person..
(speaking in foreign language).
- Our passage today in Ephesians will remind us.
of this truth, that you are a truly blessed person..
(speaking in foreign language).
, so let's read our scripture today..
(speaking in foreign language).
- So let's read our scripture today..
Blessed be the God and Father of our Lord Jesus Christ,.
who has blessed us in Christ with every spiritual blessings.
in the heavenly places, even as he chose us in him.
before the foundation of the world,.
that we should be holy and blameless before him..
(speaking in foreign language).
- In love he predestined us for adoption to himself.
as son to Jesus Christ, according to the purpose.
of his will, to the praise of his glorious grace,.
which he has blessed us with..
(speaking in foreign language).
- In love he predestined us for adoption to himself.
as son to Jesus Christ, according to the purpose.
of his will, to the praise of his glorious grace,.
which he has blessed us in the beloved..
In him we have redemption through his blood,.
the forgiveness of our trespasses,.
according to the riches of his grace,.
which he lavished upon us in all wisdom and insight,.
making known to us the mystery of his will,.
according to his purpose, which he sets forth in Christ..
(speaking in foreign language).

$^{81}$- As a plan for the fullness of time,.
to unite all things in him, things in heaven,.
and things on earth..
In him we have obtained an inheritance,.
having been predestined according to the purpose of him,.
who works all things according to the counsel of his will,.
so that we who were the first to hope in Christ.
might be to the praise of his glory..
In him you also, when you heard of the word of truth,.
the gospel of your salvation, and belief in him,.
were to with the promise, Holy Spirit,.
who is the guarantor of our inheritance.
until we acquire possession of it,.
to the praise of his glory..
(speaking in foreign language).
- This is a pretty long passage..
Oftentimes when we're excited,.
we often packed a lot of words in just a single breath..
(speaking in foreign language).
Even though this is a 14 verse,.
but in the original language,.
Paul actually put this in one long sentence..
(speaking in foreign language).
- So in this verse, Paul does just that..
He used one long, breathless sentence to urgently declare,.
you have every kind of spiritual blessings..
You may think that Paul has just won \$5 million.
when he say that..
(speaking in foreign language).
- But no, Paul was in prison at that time,.
actually deprived of his freedom..
(speaking in foreign language).
- And yet he shared with his readers.
about this abundance of God's blessing..
So when you hear this, what kind of reaction do you have?.
(speaking in foreign language).
- Your first reaction might be, is this guy crazy?.
(speaking in foreign language).
- But if you look at Paul's epistles,.
all his language and his wording.

$^{121}$doesn't show that he seems crazy at all..
So this means that we have to take him seriously..
(speaking in foreign language).
- So is Paul crazy or is this guy crazy?.
(speaking in foreign language).
- So if you look at Paul's epistles,.
(speaking in foreign language).
- So is Paul crazy or is he serious?.
(speaking in foreign language).
- And of course, Paul is not crazy..
What he says is true..
(speaking in foreign language).
- So what are the blessings that Paul mentioned?.
The first one is the blessed election from verse four to six..
(speaking in foreign language).
- People come to faith at different times.
in various different ways, but have you ever wondered,.
did God choose you or did you choose God?.
(speaking in foreign language).
- So let's do a quick survey..
Raise your hand if you think that you have chosen God..
(speaking in foreign language).
- Now raise your hand if you think that God has chosen you..
(speaking in foreign language).
- If you didn't raise your hands,.
that probably means that you didn't get my point..
(speaking in foreign language).
- From our experience, sometimes it can feel like.
we are the one that choose God..
For example, when life get tough.
and a friend bring you to church.
and share with you about Jesus.
and you decide to believe in Jesus,.
this can be a real experience..
(speaking in foreign language).
- But from a broader perspective,.
the scripture tell us that God chose us.
before we ever choose him..
And so Paul says that God's predestination in election.
are the turning point of all Christians' life..

$^{161}$(speaking in foreign language).
- So without God's election,.
there will be no blessings to follow..
And without God's election, we won't be sitting here today..
(speaking in foreign language).
- And Paul continued to say that God chose his people.
before the foundation of the world,.
and before the time and space even existed..
(speaking in foreign language).
- So before the world even begins,.
God has already chosen us..
(speaking in foreign language).
- God has chose you and me.
and all those who might not get his blessing..
(speaking in foreign language).
- And all those who might not get my point yet..
(speaking in foreign language).
- So this is probably greater than the love stories,.
the greatest love story in our world,.
that God passed through time and space to find us..
(speaking in foreign language).
- You might ask, wait, I didn't even exist then..
How could God choose me?.
(speaking in foreign language).
- The truth is God transcends all things,.
including our human understanding of time and space..
So God has chose us even before.
we ever come into existence..
(speaking in foreign language).
- Here's an example..
My son was born nine months ago,.
but two years ago, before my wife was even pregnant,.
he was already in our minds..
(speaking in foreign language).
- But before he was even born,.
we had already chosen his name..
We planned his birth at Northwestern Hospital..
We picked the formula brand that we would be using,.
and also the diaper brand,.
and also picked which room he would be living in..

$^{201}$(speaking in foreign language).
- Now when my son is born,.
those thoughts are now becoming reality step by step..
(speaking in foreign language).
- But God's election far exceeds this example..
It is a profound mystery..
(speaking in foreign language).
- So what are we predestined for?.
Paul says that it is our adoption to sonship.
through Jesus Christ..
(speaking in foreign language).
- This was quite a scandalous idea during Paul's time..
(speaking in foreign language).
- In Paul's Romans context,.
only Caesar could claim the title of son of God..
Any other person who tried to claim that title.
can, will be beheaded..
(speaking in foreign language).
- Can, will be beheaded..
(speaking in foreign language).
- However, the true son of God is Jesus Christ alone,.
and through him, God adopts us as his sons and daughters..
In Roman culture, legal adoption granted the adopted child.
actually the full status and inheritance right.
equal to those of a biological child..
(speaking in foreign language).
- So when God adopted as his son,.
we have legal rights and also the rights to inheritance..
(speaking in foreign language).
- When I look at the news this week,.
I saw that the US president's son.
is actually drawing a lot of attention in school..
People will flock to befriend him,.
and girls will want to date him..
Not because of who he is,.
but because he is the son of the president..
(speaking in foreign language).
- But being a president's child only bring worldly status,.
but being God's child actually bring eternal significance..
(speaking in foreign language).

$^{241}$- So what can we get from this?.
We became God's people not by accident..
Every Christian is unique,.
and our identity and our value rests not on our achievement,.
but on God's grace and his election..
(speaking in foreign language).
- So when God elects us,.
that means that we have a status in this world..
(speaking in foreign language).
- The second scripture says that salvation.
comes from the Lord,.
and our salvation comes from God's election.
and his redemption through Jesus,.
not from our own good deeds or our character..
(speaking in foreign language).
- Many people think that becoming a Christian.
is all about being a good person,.
but that is a big mistake..
We did nothing to earn our own salvation..
It is God who chose us,.
so as Christians, we should be humble.
and grateful before the Lord..
(speaking in foreign language).
- And the second blessings that we have in God.
is the blessed salvation..
In verse seven, "In him we have redemption through his blood..
The forgiveness of sins.
in accordance with the riches of his grace.".
(speaking in foreign language).
- So redemption is the ransom paid.
by a slave for his own freedom..
(speaking in foreign language).
- And so when God redeems us with a ransom,.
that ransom is Jesus himself..
(speaking in foreign language).
- Many people think that they are a free person,.
but the Bible says that all who sins are slaves to sin..
(speaking in foreign language).
- We only have freedom in our sins,.
but not freedom from our sins..

$^{281}$(speaking in foreign language).
- Only in Christ Jesus, our freedom can be restored..
If you're not in Jesus,.
you are under the bondage to all kinds of things.
in the world..
And so come to Jesus and you will have true freedom..
(speaking in foreign language).
- Paul then continues,.
"Humanity is redeemed through the blood of Jesus.".
(speaking in foreign language).
- So in the Old Testament,.
before the Israel worship at the temple,.
the high priest would lay their hands on an animal,.
transferring Israel's sins to the lamb..
And then the high priest would slaughter it.
and bring his blood into the holy of holies.
so that God can forgive their sins..
(speaking in foreign language).
- But could these lambs truly bear our sin?.
(speaking in foreign language).
- Absolutely no..
So these lambs cannot bear our sin..
They only foreshadow Jesus, the one true lamb of God..
And God prepared the true lamb for his chosen people..
(speaking in foreign language).
- So what can we take away from this blessed salvation?.
A lot of times we think of forgiveness of sins.
as something that is easy,.
but the truth is the price is actually pretty costly..
(speaking in foreign language).
- What is the most valuable things in our world?.
A mansion, a fancy car, Bitcoin?.
These are all valuable, but they are not the most valuable..
(speaking in foreign language).
- My wife always says anything that money can solve.
isn't a big problem..
What money cannot solve is the real problem..
(speaking in foreign language).
- Among all the things that money cannot buy,.
salvation is the most important one.

$^{321}$because Jesus' death is required to pay.
for the wages of our sin..
(speaking in foreign language).
- Jesus is the most precious and most honored in the world..
(speaking in foreign language).
- And the saddest reality in our life.
is that we live our whole life on earth,.
but we completely miss the most valuable gift.
that is the gift of salvation through Jesus..
(speaking in foreign language).
- And the third blessings that we have is the blessed seal..
(speaking in foreign language).
- Paul say when you share the truth and believe in Jesus,.
you receive the Holy Spirit as a seal..
(speaking in foreign language).
- Seal has several meanings..
First, it signifies that this thing is real,.
not counterfeit..
(speaking in foreign language).
- In most places in the world,.
there will be no legal effects without an actual seal..
We put our signature on text forms.
or house deeds to validate its authenticity..
(speaking in foreign language).
- So the seal of the Holy Spirit on his believers.
means that the blessings of God's election,.
the blessings of Christ's redemption,.
and the blessings of our adoption to sonship,.
they're all true and they're legally binding..
(speaking in foreign language).
- A seal also represents ownerships..
Many farmers, they will put a stamp on their animals.
or tag their ears so that when other people see it,.
they know who these animals belong to..
(speaking in foreign language).
- So when the Holy Spirit puts his stamps on the believers,.
that means that we belong to God, not to other people..
(speaking in foreign language).
- In our today's multicultural context,.
ownership is often labored as hegemony..

$^{361}$A wife belonging to her husband is called patriarchy..
A child to the parents is considered as parental control..
And people belong to God is viewed as divine oppression..
But this is not true..
(speaking in foreign language).
- So when I say that my child belongs to me,.
it means that I will care for and protect them,.
not to oppress them..
And so in the same way,.
when Paul says that Christians will belong to God,.
it means that God will care for his people,.
not to oppress them..
(speaking in foreign language).
- So thirdly, seal also represents security..
All our important documents will have seal.
to protect its content, right?.
A police would tap off areas with seal,.
meaning that this area is blocked.
and that you are not allowed to tamper it..
And so the seal of the Holy Spirit on God's people.
means that God will protect his people.
until the day when we meet him face to face..
(speaking in foreign language).
- And so our salvation in Jesus is secure..
(speaking in foreign language).
- So what can we take away from these blessings.
of the seal of the Holy Spirit?.
As Christian, we face all sorts of worldly temptation,.
different difficulties, challenge, and failures..
And no one can live our Christian life.
from the beginning to the end by relying on themselves..
(speaking in foreign language).
- Can you be sure that the day that you die,.
you will still believe in Jesus?.
Or have you ever doubted.
whether you can persevere to the end?.
(speaking in foreign language).
- The world is always luring us away from Jesus all the time..
And as Christians, we experience peaks and valleys.
in our faith as well..

$^{401}$(speaking in foreign language).
- But thank God that he knows we are weak sometimes..
And thank God that he knows that the temptation.
from the world sometimes can pull us away from him..
(speaking in foreign language).
- And when we are sealed by the Holy Spirit,.
the Spirit becomes the guarantee of our salvation,.
not our own feelings..
(speaking in foreign language).
- So how can we be sure that we won't fall away?.
The Spirit of the Lord guarantees it..
So when we are weak, are we willing to rely.
on the Holy Spirit instead of ourselves?.
(speaking in foreign language).
- So let's conclude our passage today..
Our passage tells us that only in Christ Jesus,.
we are truly blessed..
(speaking in foreign language).
- So in Christ, bless you means that God has chosen you.
and given your life true meaning..
In Christ, bless you means that Jesus has redeemed you.
and your life is worth everything to him..
And in Christ, bless you means that the Holy Spirit.
has sealed you and your life is secure in him..
(speaking in foreign language).
- So the questions that I have for all of us is this..
Are you living in this blessing?.
(speaking in foreign language).
- Let's pray..
(speaking in foreign language).
- Heavenly Father, thank you for giving us your word..
(speaking in foreign language).
- We pray for those among us who have not yet received.
or understand these blessings..
Would you open your eyes to see the need.
for your salvation?.
(speaking in foreign language).
- For those who already know these blessings,.
help us to cherish them and also give you thanks.
for your grace..

$^{441}$(speaking in foreign language).
- In the name of Jesus Christ, amen..
Okay, thank you, Paul..
And thank you all for coming..
We will have some refreshment outside,.
so feel free to stay after and get to know some new people..
All right, go in peace, UV, your love..
(people chattering).
(gentle music).
(gentle music continues).
\newpage



\section{傳道書 2:1-11}
\label{sec:9E6MnSM4U0Q}
\textbf{抗疫虛空 (傳道書 2\_1-11) - 譚文倩莊傳道}
\newline
\newline
連結: \href{https://youtube.com/watch?v=9E6MnSM4U0Q}{\texttt{https://youtube.com/watch?v=9E6MnSM4U0Q}} ~~~~ 語音日期: 2025-10-01
\newline
\newline
\hyperref[sec:7m42sBmfTkI]{\small{< < < PREV SERMON < < <}}
~
\hyperref[sec:index_chronic]{\small{[返順時目]}}
~
\hyperref[sec:index_scriptual]{\small{[返順卷目]}}
~
\hyperref[sec:Dh5pU8Y4jPw]{\small{> > > NEXT SERMON > > >}}
\newline
\newline
傳道書 2:1-11
\newline
\begin{longtable}{cl}
\hline
\hline
章節 & 經文 (和合本修訂版)\\
\hline
2:1 & \begin{tabularx}{0.7\textwidth}{X} 我心裡說：「來吧，讓我用喜樂試試你，使你享福！」看哪，這也是虛空。 \end{tabularx} \\ \\ \relax
2:2 & \begin{tabularx}{0.7\textwidth}{X} 論嬉笑，我說：「這是狂妄。」論享樂，「這有甚麼用呢？」 \end{tabularx} \\ \\ \relax
2:3 & \begin{tabularx}{0.7\textwidth}{X} 我心以智慧引導我，我心裡探究，如何用酒使身體舒暢，如何抓住愚昧，直等我看明世人在天下短暫一生中，當行何事為美。 \end{tabularx} \\ \\ \relax
2:4 & \begin{tabularx}{0.7\textwidth}{X} 我大興土木，為自己建造房屋，栽葡萄園， \end{tabularx} \\ \\ \relax
2:5 & \begin{tabularx}{0.7\textwidth}{X} 修造庭園和公園，在其中栽種各樣果樹， \end{tabularx} \\ \\ \relax
2:6 & \begin{tabularx}{0.7\textwidth}{X} 挖造水池，用以灌溉林中的幼樹。 \end{tabularx} \\ \\ \relax
2:7 & \begin{tabularx}{0.7\textwidth}{X} 我買了僕婢，也有生在家中的僕婢；又有許多牛群羊群，勝過我以前所有在耶路撒冷的人。 \end{tabularx} \\ \\ \relax
2:8 & \begin{tabularx}{0.7\textwidth}{X} 我為自己積蓄金銀，搜集各君王、各省份的財寶；又為自己得男女歌手和世人所喜愛的物，以及一個又一個的妃嬪。 \end{tabularx} \\ \\ \relax
2:9 & \begin{tabularx}{0.7\textwidth}{X} 這樣，我就日漸昌盛，勝過我以前所有在耶路撒冷的人。我的智慧仍然存留。 \end{tabularx} \\ \\ \relax
2:10 & \begin{tabularx}{0.7\textwidth}{X} 凡我眼所求的，我沒有克制它；我心所樂的，我沒有不享受。因我的心要為一切的勞碌快樂，這是我從一切勞碌中所得的報償。 \end{tabularx} \\ \\ \relax
2:11 & \begin{tabularx}{0.7\textwidth}{X} 後來，我回顧我手所經營的一切和我勞碌所做的工。看哪，全是虛空，全是捕風；在日光之下毫無益處。 \end{tabularx} \\ \\
[1ex]
\hline
\hline
\end{longtable}
$^{1}$Katrina是很年輕的,因為她不斷說張瑪麗,我覺得她一定是很年輕,應該是還沒有出生,其實她是叫張瑪麗,她是講者..
早晨各位弟兄姊妹,我亦歡迎第一次參加崇拜的朋友,我代表教會歡迎你,願主祝福你..
如果你有任何問題或需要時,在崇拜後大概一個小時後,你就可以跟我們任何一個戴著徽章的教務同工或核心同工聯絡,我們可以嘗試解答你的問題..
我們先低頭祈禱..
《箴言傳道書約百記》是聖經舊約裡的智慧文學,所講的智慧不是普通的知識或學問,而是一種以敬畏上主為基礎,是當人面對人生苦難,道德,社會,日常選擇的一個全人的態度..
什麼是智慧文學呢?.
智慧文學是一種重視人生經驗,個人思考和生活實踐的文體,注重神和人之間的關係,人如何運用神所賜的智慧來生活..
三本書中各有重點,《箴言傳道書約百記》強調行事為人,建基在對上帝的敬虔,就是《箴言傳道書約百記》所講的智慧..
《箴言傳道書約百記》一章七節,經威耶和華是知識的開端,但如往人藐視智慧和分晦..
約百記討論敬虔的人如何在苦難中掙扎,面對一個似乎違反了公義的上帝,坦然並勇敢地與上帝對話..
約百記四十二章五節,我從前只是風聞有理,但現在親眼見你..
《箴言傳道書約百記》檢視日光之下人勞碌的一切事,反思生命的意義和虛空,結論是我們人要知道自己的本分,我們要敬畏神,我們要遵守神的誡命..
因為最終有一天神都會審判人在日光之下所做的每一件事..
《傳道書》十二章十三至十四節,他說你們一切都聽見了,總言而之應當敬畏神,遵守他的誡命,因為這是每一個人的本分..
人所做的每一件事,連一切隱藏的事,無論是善是惡,神都必審問..
我之前也跟大家說過,因為有一次在教會被人叫我做「Grandma」,之後引發我決定要看這個《傳道書》,也嘗試將我自己的領袖和大家分享..
之前我用了兩堂的時間跟大家說第一章,今天我們就進入《傳道書》的第二章,但是我想先跟大家重溫一下..
《傳道書》一開始就有句「去控的去控,去控的去控,凡事都是去控」,這個哲理貫穿了《傳道書》的整個主題..
描寫傳道者細膩地反思人生裡面的智慧,享樂,成就,權力,財富等等的追求..
到最後發現所有的東西都是無法給予真正的滿足..
在經歷了林林總總的追求和矛盾之後,傳道者得出一個結論,就是面對人生的短暫和無常,我們人生最大的目標,最有智慧的導向,.
不是在不斷地追求日光之下,即是沒有永恆價值的東西,而是在我們的生命,有限的生命裡面去承認那位創造天地萬物的主,賞賜的主,.
能夠真正滿足我們的主,再者就是要謙卑和知足,我們要心存感恩,去好好享受神所賜的..
第一堂《日光之下當思考》一章一至十一節,當中有五個虛空,第一個虛空是講對那個狀況的陳述,.
第二個虛空還要加強之前那個虛空,第三個虛空是加強重複狀況的陳述,.
然後再來一次提起,然後到最後總結我們日光之下的牢籠,到最後都是虛空..
這個結論傳道者提出了一個很根本性的問題,就是我們人一生的形形形形到底有什麼益處呢?.
人的生命是不是真的有意義呢?第二堂《日光之下再思考》去到第十二至十八節,.
他說人生裡面有很多彎曲,意思是我們生命裡面有些東西是我們人沒有辦法,我們沒有能力改變的,.
無論我們人多有智慧都無法解決所有的問題,有時知識多反而帶給我們的受煩更加多..
傳道書在這裡提醒我們人的有限,而這種限制正正是傳道書一章十三節裡面所講的沉重的擔子..
今天我們來看傳道書二章一至十一節,主題抗疫虛空..
這段經文記錄了傳道者的生活,是講他喝最好的威士忌,.
他老婆妻妾一大堆,他自己建造自己的房子,花園水池各種各樣的果樹,.
他還要有全自動的水流,不單止這樣,他還可以隨時叫他的工人去燒一塊牛扒羊扒給他吃,.
因為他自己養著很多牛和羊,錢他用不完,因為每個月附近的地方政府都要向他進貢,要賣錢給他..
傳道者很大口氣地說,我擁有這個世界上的人所渴求的一切,我有所有東西,.
但是為什麼後來他又說,這些東西其實最終都是沒有意義的,.
如同暴風,世上沒有什麼是有真正價值的呢?.
虛空的虛空,虛空的虛空,全是虛空..

$^{41}$Honestly 想一想,為什麼追求這個物質,成就或者享受,不可以完全滿足我們的生命呢?.
為什麼傳道者會在這裡說,這些東西所有東西都是虛空的呢?.
再想一想,其實有錢有面子有什麼問題呢?.
享受又有什麼問題呢?.
其實想一想,不懂得享受才是問題,.
今天我們試一下將這段經文分開三段來思考,我們先看第一段經文,.
傳道書二章一至三節,他說我對自己說,內霸去享受去追求生活的美好吧,.
但我發現這也沒有意義,於是我就說,嬉笑是愚蠢的,追求享樂有何益處?.
想來想去,我決定喝酒自樂,我尋求的是智慧,找住的卻是愚昧,.
我試圖以此體驗快樂,就是大多數人在世上短短一生僅有的快樂..
首先我想大家留意,傳道書二章一開始的時候,.
那頭三節其實是跟第一章的修辭結構是一樣的,.
意思是這三節經文跟第一章是互相呼應的,.
傳道者用一個客觀的態度去觀察,描述反思人生跟自己心靈去對話,.
再轉變成一個觀點的個人自轉式的論證去呼應人一生的勞碌,.
最終到底得到些什麼呢?最終似乎都是虛空..
來到第二章,他的開始就是在說什麼?.
他說:Come on,我要好好享受,我要experience the good life..
對你來說,什麼是good life呢?.
我想我們在不同的年紀,其實是有不同的答案..
在這裡傳道者來到年老的時候,似乎他已經得到那個good life,.
因為他已經擁有了他所有追求渴求的東西,.
他已經應該有the good life,沒有欠缺,應有盡有..
留意這段經文,我試圖去享樂,去追求生活的美好,以此體驗快樂..
意思是我嘗試藉著享受,藉著去玩,去證實,去prove,.
這些所有東西是可以滿足我的..
接著我們留意這段經文最後那幾個字,.
"僅有的快樂,for the few years of their lives"..
意思是說,我們人生其實快樂的日子不是很多的,.
因為困難和苦難的日子是更加多,所以是僅有的快樂..
是我們人生的確是這樣,當你真的回顧去想的時候,.
似乎很多時候我們快樂的日子不是很多..
要命嗎?一個已經擁有所有東西的人,.
他現在是告訴你,我擁有所有東西的時候,其實都是無法滿足我的,.
其實我都是不開心的,我發現這些東西全部都是沒有意義的..
接著他還要說,開心又怎樣?很快就過了,很快就沒有了,.
就是極虛空,和第一章的經文在互相呼應,.
新爐有什麼益處?生命有什麼意義?.
下一段經文,傳道書二章四至八節,.
我還大興土木,為自己蓋殿宇,興建美麗的葡萄園,.

$^{81}$試圖找到生命的意義,我建造花園苗圃,.
在園中栽種各種果樹,又建造蓄水池,灌溉大片茂盛的園林..
我買來爐木,有男有女,也有生在家中的爐木,.
我有大群的牛羊,超過所有曾住在耶路撒冷的君王,.
我積聚了大量的金銀,就是諸王及各省的財寶..
我顧了一流的歌手,有男有女,還納了許多漂亮的妃品,.
我擁有世上人所渴求的一切..
偶爾我和我老公也會做夢,特別是我,.
有錢的話,我要建造一間怎樣的房子,.
房子內的設計,應該怎樣的樣子,.
當然最重要的是我的廚房要怎樣設計,.
裡面要有些什麼,比起我們這些星斗市民,.
所羅門當然不用做夢,因為他有很多錢,.
甚至他喜歡在哪裡建房子也沒有問題..
列王記上九章十五至十九節記載,.
所羅門王在位的期間,他在耶路撒冷建造了宏偉的第一聖殿,.
聖殿是約瑰永遠的居所,也是以色列人敬拜上帝的地方,.
他也建造了一群龐大而且華麗的宮殿建築群,.
給他自己和他的家人住..
所羅門出了名有智慧,在他統治之下,.
古代的以色列國是最繁榮昌盛,影響最深遠的時代..
在這段經文裡面,出現了很多在《創世記》一至二章裡,.
都有出現過的詞彙,.
例如栽種,園林,各樣果樹,灌溉,苗圃,興建..
很多聖經學者都認為,.
所羅門王似乎在嘗試興建一個屬於他自己的伊甸園,.
一個沒有神在裡面的伊甸園..
所羅門不單只扮演這個君王的角色,.
他甚至沒有嘗試扮演上帝..
在聖經裡面,所羅門王從來都沒有被描繪成為一個嘗試扮演上帝的人..
但是他的行為展現出來的那個自我,.
不難看到他已經將神放在一邊..
所羅門王十幾歲就登基,.
這個年輕人,說的是十幾歲,.
他登基的時候,他不是懶得拿著雞毛當令箭指著這個指著那個做事,.
他反而是一個這麼年輕的人,他懂得回到神面前去求智慧,.
求上帝給他智慧,能夠統領這個他要管理的地方,.
是不是很值得欣賞?.
其中一個所羅門王很出名,智慧非凡的故事,.
是記載在《列王記》上三章那裡,.

$^{121}$故事記載Betty和Susan來到所羅門王的面前,.
指責對方將他自己死去的孩子和我生下來的孩子交換了,.
在沒有目擊證人的情況下,.
所羅門王說最好的方法就是將這個生下來的孩子劈開一半,.
一半給你,非常公平,解決,.
Betty聽到之後心急如災哭了出來,.
然後她說不要不要,如果你決定要這樣做的時候,.
我就寧願不要這個孩子,將這個孩子送給Susan,.
It's ok,結果憑著Betty和Susan的反應,.
所羅門分辨到哪一個人才是這個孩子的親生媽媽,.
就判了這個孩子給Betty,.
這件事傳出去之後,以色列的人都很驚嘆她的智慧,.
看到這個王展現了上帝的公義,是不是很有智慧?.
溫馨提示,故事出自聖經,.
但Betty和Susan這兩個名字是我加進去的,.
所羅門王他真的厲害,他統領以色列,建立王國,.
建立他的財富,甚至在列王記上十一章記載了他有很多外邦妻妾,.
說的是有七百個公主,三百個妾侍,.
一天都要三年多,.
他為了哄那些外邦妻妾,竟然斗膽興建了兩個外邦的偶像神殿,.
他做了什麼?.
他是不是做這個以色列王,做到亂七八糟,.
而且他忘記了哪一個是他生命的主,.
好像我之前說的,可能我們不是像所羅門王那樣有錢有權有勢,.
但我們每一個人都可能在興建一個屬於我們自己,.
以我自己為中心的伊甸園,.
我們在這個由我們掌控的小小的伊甸園裡,.
我們沒有地方放神在裡面,.
你自己反思一下,在你小小的伊甸園裡面,神存不存在?.
若果神存在在你的伊甸園裡面,.
我想其中一個應有的元素,就是被建立的東西不是一些硬件,.
應該被建立的是那些跟你在伊甸園裡面的人..
二章九節,他說這樣我越來越強盛,.
超過之前住在耶路撒冷的所有人,.
我的智慧從未減少,我想要什麼就拿什麼,.
從不拒絕享樂,甚至在辛苦勞碌中,.
我也能找到極大的樂趣,.
這就是對我所有辛勞的回報,.
然而我回顧自己辛勞成就的一切,.
原來全都沒有意義,如同暴風,.

$^{161}$世上沒有什麼是真正有價值的..
我們當中很多人,我們對我們的將來都有夢想,.
我們都很願意為我們的夢想努力拼搏,.
如周星馳所說,做人如果沒有夢想就會怎樣?.
跟一條鹹魚沒有什麼分別..
回到中學生團契的時候,.
想當年我們的導師推崇我們去看一本書,.
那本書的名字叫做《Money, Sex, and Power》,.
金錢,性與權力..
這本書探討聖經的原則,.
給基督徒能夠落地,坦誠的方式,.
去回應這個時代三個最大的誘惑..
作者Richard Foster提到,.
這個世界沒有一樣東西能夠比金錢,性和權力.
更加能夠深深地觸動我們,.
也沒有一個議題比起這三樣東西更引起爭議..
而金錢,性和權力更加是我們人類的世界裡面.
可以帶來很大的祝福,.
也可以帶來最大的破壞的三個元素..
因為我們人的一生就是不停地去追求,.
嘗試在這三樣東西裡面去得到滿足..
所羅門王也不例外,.
我們和他,或者他和我們其實是一樣的,.
不過在這個反省裡面,.
他告訴我們,.
原來追求這一切都是沒有意義的,.
如同暴風,世上沒有什麼是真正有價值的..
在二章十一節裡的如同暴風,.
其實和第一章十四節的我見日光之下所發生的一切事,.
看那全是虛空,全是暴風,.
是說同一樣東西的,.
暴這個字的意思是捉住,.
catch,或者抓住,capture的意思..
所羅門王嘗試是說他努力去捉住,.
去抓住,.
努力去追求達至,.
或者努力去追求渴望能夠得到,.
結果最終他真的可以擁有世人所渴求的一切,.
但是去到晚年回顧一生的時候,.
他卻在那裡慨嘆,.

$^{201}$很虛空,捉不住,好像抓不住一樣..
這個就正正是傳道書裡面所說的虛空..
不知道你有沒有留意,.
傳道書的二章一至十一節,.
這十一節經文裡面,.
有很多我在裡面的,.
我還,為自己,我建造,.
我買,我有,.
我積聚,我固了,我擁有,.
我越來越我的,.
我想要我也有,.
我所有,我回顧..
想一想其實都正常的,.
這裡是說所羅門王反省自己,.
反省的時候難道幫你,幫其他人反省嗎?.
問問你,.
你有沒有一個自我反省的習慣?.
你有沒有一個自省的能力?.
這麼多年來我都有一個習慣,.
久不久我就會找一個空間,製造一個空間,.
讓我自己靜下來,.
反省一下我自己,.
反省一下,思考一下最近發生的事情..
在這些時候,一路思考的時候,.
我就看到很多神在我生命裡面工作的足跡,.
以致當我在這個空間裡走出來的時候,.
我會有一個比較正確的方向繼續前行,.
因為在這個反省的過程裡面,.
我會被提醒很多事情的發生,.
我自己在處理的事情,.
旁邊都是有神和我同在..
我想我們基督徒其實很需要有這個習慣,.
因為其實很多時候我們向前不斷衝,.
我們不知道自己在衝去哪裡,.
我們人都是很自私,.
我們都是自我中心,.
我們只是經常想著自己,.
我自己覺得這樣好,.
我覺得這樣好,.
我想要,我想擁有,.

$^{241}$永遠都是這樣朝著這樣走,.
我們很驕傲,.
我們很需要製造一個這樣的空間,.
停下來去思考,.
去看到神在我生命裡面在做些什麼,.
看到神的足跡..
在很多這些我我我裡面,.
有祂,.
這個祂將我由太空拉回到地面裡,.
踏實些這樣在主裡面去思考我的生命,.
回頭去思考我之前的問題,.
為什麼說到好像追求物質成就或者享受,.
是不可以完全滿足我們的生命,.
為什麼傳道者說最後這些都是虛空,.
有錢有面又懂得享受又有什麼問題,.
其實基督徒有錢有面又懂得享受,.
是沒有問題的,.
最好就是這樣,.
因為這些都是神給我們的祝福,.
神給我們的賜予,.
事實上傳道書裡面都有經文是這樣說的,.
問題只不過是在這個日光之下短暫,.
困難比開心日子多的生命裡面,.
人沒有繼續敬拜賜我們生命的主,.
能夠引導我們走這條人生路的主,.
你以為我能夠自己努力抓住抓住追求成就和享受,.
我自己擁有更多更多,.
買齊了,.
所有東西都沒有欠缺,.
我就ok了,.
我就開心了,我就滿足了,.
你以為你可以掌握,.
你以為你有把握,.
結果可能好像拿著一堆沙似的,.
你越大力去捏,.
那些沙就在你的手指縫裡面流出來,.
傳道書是聖經舊約裡面其中一本智慧文學,.
它深入探討人生的意義,.
書裡面指出,.
我們人生追求的財富,.

$^{281}$享受,智慧權力,.
其實是難以掌握的,.
在這個困難多過快樂的人生裡面,.
尤其在面對痛苦和死亡的時候,.
更加凸顯了我們人的有限和生命的限制,.
傳道者去到晚年的時候,.
他發現單靠自己一生的努力,.
最終只會去空,.
傳道師二章邀請你再反思,.
你現在很努力去捉住,抓住的,.
去思考是不是真正有價值,.
是不是這些東西真的能夠滿足到你,.
求主幫助我們,.
在我們自己創造的那個世世界伊甸園裡面,.
不要將神慢慢推出去,.
邀請他回到你的生命裡面,.
求主幫助我們,.
我們的生命能夠敬拜那個真正能夠滿足我們的主,.
我們一起低頭祈禱,.
主日求你的話語,.
進入人的心裡面去挑戰我們的靈,.
特別是如果我們的生命裡面,.
我們信主信得久,慢慢將你推下去牆邊,.
你已經不再是我們最終的追求,.
我們想捉住,抓住的對象,.
主日求你赦免,.
求你幫助我們,.
我們軟弱,.
但我們相信你的大能能夠將我們挽回,.
我們多謝你施恩,.
主日多謝你憐憫,.
多謝你愛我們,.
這樣去祈禱,.
奉主耶穌基督的聖名祈求,.
Amen.
再一次多謝,.
這次第一次來我們當中的朋友,.
願神祝福每一位在座的弟兄姊妹,.
散會.
(影片完結).

\newpage



\section{詩篇 63:1-11}
\label{sec:Dh5pU8Y4jPw}
\textbf{Seeking God (Psalm 63) - Rev. Mark Chan 尋求神 (詩篇63\_1-11) - 陳祟基牧師}
\newline
\newline
連結: \href{https://youtube.com/watch?v=Dh5pU8Y4jPw}{\texttt{https://youtube.com/watch?v=Dh5pU8Y4jPw}} ~~~~ 語音日期: 2025-10-14
\newline
\newline
\hyperref[sec:9E6MnSM4U0Q]{\small{< < < PREV SERMON < < <}}
~
\hyperref[sec:index_chronic]{\small{[返順時目]}}
~
\hyperref[sec:index_scriptual]{\small{[返順卷目]}}
~
\hyperref[sec:jhQYPPMz5JE]{\small{> > > NEXT SERMON > > >}}
\newline
\newline
詩篇 63:1-11
\newline
\begin{longtable}{cl}
\hline
\hline
章節 & 經文 (和合本修訂版)\\
\hline
63:1 & \begin{tabularx}{0.7\textwidth}{X} 神啊，你是我的神，我要切切尋求你；在乾旱疲乏無水之地，我的心靈渴想你，我的肉身切慕你。 \end{tabularx} \\ \\ \relax
63:2 & \begin{tabularx}{0.7\textwidth}{X} 我在聖所中曾如此瞻仰你，為要見你的能力和你的榮耀。 \end{tabularx} \\ \\ \relax
63:3 & \begin{tabularx}{0.7\textwidth}{X} 因你的慈愛比生命更好，我的嘴唇要頌讚你。 \end{tabularx} \\ \\ \relax
63:4 & \begin{tabularx}{0.7\textwidth}{X} 我還活著的時候要這樣稱頌你，我要奉你的名舉手。 \end{tabularx} \\ \\ \relax
63:5 & \begin{tabularx}{0.7\textwidth}{X} 我在床上記念你，在夜更的時候思念你；我的心像吃飽了骨髓肥油，我也要以歡樂的嘴唇讚美你。 \end{tabularx} \\ \\ \relax
63:6 & \begin{tabularx}{0.7\textwidth}{X} 見上節 \end{tabularx} \\ \\ \relax
63:7 & \begin{tabularx}{0.7\textwidth}{X} 因為你曾幫助了我，我要在你翅膀的蔭下歡呼。 \end{tabularx} \\ \\ \relax
63:8 & \begin{tabularx}{0.7\textwidth}{X} 我的心緊緊跟隨你；你的右手扶持了我。 \end{tabularx} \\ \\ \relax
63:9 & \begin{tabularx}{0.7\textwidth}{X} 但那些尋索要滅我命的人必往地底下去； \end{tabularx} \\ \\ \relax
63:10 & \begin{tabularx}{0.7\textwidth}{X} 他們必被刀劍所殺，成為野狗的食物。 \end{tabularx} \\ \\ \relax
63:11 & \begin{tabularx}{0.7\textwidth}{X} 但是王必因神歡喜，凡指著他發誓的都要誇耀，因為說謊之人的口必被塞住。 \end{tabularx} \\ \\
[1ex]
\hline
\hline
\end{longtable}
$^{1}$And let's come before the Word of God and we're going to continue to meditate and to.
study..
(Translation).
Okay, first question I have for you is this chart..
Okay?.
Now I want to ask you, you know, this chart is something I learned many years ago..
And I remember a speaker I listened to and he talked about the emotional wheel..
And then there were different colors representing different emotions of our lives..
And I want to do one thing before we jump into the topic and speaking of a Psalm, Psalm.
63..
And the question is, ask yourself today, at this moment, ask yourself, ask yourself, okay?.
Slow down and ask yourself what emotions you have experienced in the past week..
And what emotions you are experiencing because of what has happened in your life..
And take a moment and try to identify it, try to feel it..
(Translation).
Okay, so how many of you have experienced happiness in the past week?.
Happiness..
(Translation).
Jizo, Jizo..
(Translation).
Okay, so how many of you have experienced happiness in the past week?.
Happiness..
(Translation).
Jizo, Jizo..
(Translation).
Okay, all right..
We are not very happy..
All right..
How many of you have felt strong in the past week?.
Strong..
(Translation).
Yeah, either you feel valuable or you feel worthy or someone praises you or someone recognizes.
you, you feel confident, strong..
How many of you?.
(Translation).
Okay, all right..
We have..
(Translation).
Okay, we are not very strong, obviously..
And how many of you felt calm in the last week?.

$^{41}$(Translation).
Okay, Jizo, Jizo..
(Translation).
Okay, all right..
We feel calm..
Okay, a lot of people felt some moments of calmness, right?.
(Translation).
So how many of you felt fear?.
The red section..
Fear last week..
(Translation).
Any?.
No?.
All right, you guys are very brave..
How many of you felt sad last week?.
(Translation).
Okay, one, two..
All right, all right..
Finally, the last quote is mad..
How many of you felt mad last week?.
(Translation).
Okay, all right, all right, all right..
Yeah..
There's nothing wrong to feel in a certain way..
Okay?.
(Translation).
Because God created us, we all have emotions..
(Translation).
But it's good to listen to your emotions..
(Translation).
Whether you feel happy, you feel energetic, strong, sad, depressed, even mad and angry,.
you know, listen to your emotions..
(Translation).
Because God so often uses your emotion to draw you closer to him..
(Translation).
Okay, I know, I know, I'm Chinese, okay?.
I'm Chinese..
(Translation).
And I grew up in a culture that I suppress my emotions..
(Translation).

$^{81}$And especially, I'm Asian male..
I've been taught, you know, either verbally or unconsciously, I inherit a kind of a, you.
know, a set up that, you know, I have to suppress most of my emotions..
(Translation).
Except being mad, okay?.
(Translation).
Except being mad, you know, we don't suppress it..
(Translation).
Which is not really good because our culture is broken, okay?.
But the Bible tells us, you know, God speaks to us and uses our emotions to talk to us..
(Translation).
So it's very important, okay?.
So get a copy of this and try to listen, try to hear your emotions..
(Translation).
And then once you try to identify your emotions, you pray to God, you know, God, I felt this.
way..
Can you help me?.
Can you listen to me?.
(Translation).
And that's exactly what Psalm 63, we're going to read that, a part of it today, Psalm 63..
It talks about all kinds of emotions in the Psalms..
And as you can, as you feel familiar, you know, this Disney movie, Inside Out, has to.
do with, you know, the inner emotions of our lives..
(Translation).
And then Psalm 63 talks about some emotions..
(Translation).
So let me read the first three verses of Psalm 63..
A Psalm of David, a Psalm of David..
When he was in wilderness of Judah, it says, "Oh God, you are my God..
Earnestly I seek you..
My soul thirsts for you..
My body longs for you..
In a dry and weary land where there's no water..
I have seen you in the sanctuary and beheld your power and your glory..
Because your love is better than life, my lips will glorify you.".
(Translation).
Now there are 150 Psalms in the Bible..
If you read every single Psalm, and you quickly realize, many of the Psalms have to do with.
human emotions..
Sadness, fear, anger, madness, disgust, joy, all kinds of emotions..

$^{121}$(Translation).
And some of our desires are broken, and some of our desires are honorable, pleasing to God..
And these emotions are expressed in prayers in the Psalms..
And we listen to, you know, listen to the emotions in your prayer..
And your emotions or your prayer is actually a window..
It's like a window to your soul..
(Translation).
And the Psalms in the Bible are prayers, people's prayers to God..
And they pray to God when life is great, and life's supposed to be..
You have good time..
And at times, you find prayers, they voice out their complaints and laments to God, and.
struggle, you know, mightily struggles with God..
And at times they are sad, they are depressed, and/or in fear, and even they are really, really mad..
(Translation).
Now let's go to Psalm 63..
And I want you to ask yourself when we read the Psalms, have you been to a situation similar to Psalm 63?.
(Translation).
Now it says, it's a Psalm of David..
Okay..
And then when David was in wilderness of Judah, he wrote this Psalm..
(Translation).
Now I give you this chart and up there you see the cross section of the land of Israel, the land of the Bible..
(Translation).
And then let me go through this cross section..
In the west, you see the Mediterranean Sea..
(Translation).
Okay, so it's on your left, okay, it's in the west..
And then you go to the right, to go to the east, you kind of climb the mountains, right?.
You climb up to the central hill country where Jerusalem is..
(Translation).
And then after Jerusalem, can you see Jerusalem?.
Okay, so after Jerusalem, you continue to go right to go to the east..
The elevation drops, right?.
Drops, even drops, drops even to the below sea level, all the way to Dead Sea Valley..
And that is the lowest point of the earth..
(Translation).
And then after the Dead Sea Valley, the Jordan River and all that, you climb again, right?.
You climb again to the mountain, what we call a Transjordanian Highland..
(Translation).
Okay, now pay attention to the rain clouds, okay?.

$^{161}$I put some rain clouds on it, all right?.
Now the rain clouds, the rain, travels from left to right..
(Translation).
Yeah, and then once you hit the watershed at the central hill country,.
and then after that, the rain will, as you can tell,.
the rain often falls on the left side or right side of the central hill country?.
The left side, right?.
It's all on the left side before hitting the watershed,.
and after the watershed, there's no more rain..
(Translation).
So when Psalm 63 talks about David was in the wilderness,.
the Bible talks about wilderness of Judah,.
and that is the region of the land right after, right side of the watershed,.
and it's always dry and thirsty..
(Translation).
And so the wilderness is often a place of testing, because there's no water..
It's often a place of temptation, because there's no food..
You know, you think of the Israelites coming out of Egypt,.
40 years they kind of being tested in the wilderness,.
and think of Jesus 40 days in the wilderness as well..
(Translation).
Okay, so let me show you two pictures..
(Translation).
The picture on the left at the bottom is west of the watershed,.
where you have water and rain..
(Translation).
And then on the right is east side, you know, the other side of the watershed,.
and you can see the picture, the contrast is dry and thirsty,.
and that is the picture of Psalm 63..
(Translation).
I don't know if you have experienced it, I'm sure you did..
The landscape is like your soul scape, right?.
The landscape sometimes is like the soul scape..
Now because of the Israel geography, wilderness becomes a metaphor,.
becomes an image for life, and when something is not right in your life,.
when life is going tough, when life is not working for you,.
you kick in kind of the survival mode, you're kind of in the wilderness,.
and whether you like it or not..
And perhaps today, you are in the wilderness of your life..
You got lost, you know, dry and thirsty..

$^{201}$That's exactly what David has in mind, you know, in the wilderness..
And he says, you know, "God, my God, earnestly I seek you..
My soul thirsts for you and my body longs for you in the dry land without water..
I've seen you in the sanctuary and I remember all that, you know, life was so good..
But now I'm in the wilderness.".
(Translation).
(Speaking Chinese).
See, David is in the wilderness and he is on the run..
He was being pursued by his enemies and the enemies trying to kill him..
And that's why he's praying the psalm in the wilderness..
And it's time of testing, it's time of danger..
In the dry land, he's looking for water,.
but David is turning to God with the spiritual experience..
And he understands that what he, what will crunch his thirst is not just water,.
but the living water, God himself..
And how does he do that?.
(Translation).
And then the first point here you have to remember is that he remember his relationship with God..
(Translation).
When life is going tough on you, you have to remember who you are in the sight of God..
(Translation).
And come back to Psalm 63, notice how he prays..
David prays, he says, "God, oh my, you are my God.".
Notice the personal pronoun here..
He said, "God is his God, my God.".
Okay, my God, my personal God..
And for David, God is not someone else's God..
God, he said, "My God.".
And that makes a huge difference, right?.
When you're going through times of wilderness, you know, you remember David is being faithful to God..
And he remember who God is..
And he remember God is faithful to him..
Not just because when time is good, even when time is tough, this God is still his God..
And he will say, "My God.".
(Translation).
And then secondly, he recall his past experience with God..
(Translation).
And then you have to remember memory, your memory defines you..
Our memories define us..
What we remember, what we recall, define who we are..

$^{241}$And you know, have you, have you known any people who suffer from Alzheimer's?.
You know, they lost their memory and they don't know who they are anymore..
It's like a baby just born or, you know, don't remember anything..
And memory is not just, memory is not just things happened in the past..
Memory, memory defines you..
It's the anchor of who you are..
Memory is a powerful, powerful thing..
(Translation).
So remember, what you recall when life is tough, what you remember when life is tough, defines you..
Okay, so here in the wilderness, David is doing something very special..
David is recalling his past experiences with God..
Notice verse 2 here..
"I have seen," David said, "I have seen you God. I've seen you in worship in the wilderness..
And I've beheld, I've seen your glory and I've experienced your power before..
Because your love is better than life, my lips will glorify you.".
Wait, in the wilderness, my lips will glorify you because I remember who you are..
(Translation).
So how do you keep your spiritual life vibrant, even when life is going tough?.
The answer here is that when God turns around my bad situation, okay, I remember that..
When life is going well and good, I have good relationship with God..
But you know, you know what I'm saying, when you feel good, when you're in good times,.
your relationship with God is superficial, it's fair service..
And you have to remember, you know, many times when life is good and well, people abandon God..
(Translation).
But here, David is having a very, very tough time in the wilderness..
But he's recalling, he's remembering, he's trying to recall his past experiences with God..
And recalling in such a way that he will continue to worship God, just like in the temple, right?.
Just like in the temple. Remember, he was in the wilderness, but he was like experiencing himself,.
remembering who God is, but even in a dry and thirsty land, he's like in the temple..
When he was a king, doing well, worshipping God in Jerusalem..
(Translation).
Okay, so today as we close, I want you to think about this..
How do you keep yourself constantly in the presence of God, in God's love?.
Memory is important..
(Translation).
When life is going tough on you, you remember God..
You remember other godly people around you..
You know, when it does to me, when I was in time of places of wilderness, testing,.
I recall, I remember, there are mentors in my life..
There are many good people in my life, godly people in my life,.

$^{281}$and they encourage me, pray for me, and support me..
And I remember them, because they were doing good works, and they were faithful..
They stick to God..
When life is tough, they continue to trust in God..
And that memory inspires me to follow Jesus..
(Translation).
And let us end by asking yourself, the first question I ask you, how do you feel today?.
Do you feel sad? Do you feel confused? Fear?.
Besides praying for God's turning everything around,.
God is inviting you to pray like David prayed..
(Translation).
Besides praying for God's turning everything around,.
God is inviting you to pray like David prayed..
So how do you feel today?.
God is inviting you to pray..
So let's spend a minute, and I'll let you to pray to God..
You know, you can tell God point blank, you know, this is how I feel today..
Whatever feelings I have, you know, God, I've come before you, help me..
Let's spend one minute of time and pray to God yourself, and then we'll close..
(Translation).
(Praying).
Let's pray..
Dear God, our God, we may feel sad..
We may feel afraid or confused..
We know that you listen to us. Help us to pray to you..
Knowing that you love us, and you see everything in our heart..
Draw us closer to you, because you are my God..
In Christ name we pray. Amen..
So remember this week, pray like this..
Spend five minutes of your day, tell God, you know, this is how I feel..
Help me pray to you. Help me pray to you God..
Okay. Peace to all of you..
(Praying).
\newpage



\section{使徒行傳 13:1-3}
\label{sec:jhQYPPMz5JE}
\textbf{向外的教會 (使徒行傳13\_1-3) - 陳崇基牧師}
\newline
\newline
連結: \href{https://youtube.com/watch?v=jhQYPPMz5JE}{\texttt{https://youtube.com/watch?v=jhQYPPMz5JE}} ~~~~ 語音日期: 2025-10-14
\newline
\newline
\hyperref[sec:Dh5pU8Y4jPw]{\small{< < < PREV SERMON < < <}}
~
\hyperref[sec:index_chronic]{\small{[返順時目]}}
~
\hyperref[sec:index_scriptual]{\small{[返順卷目]}}
~
\hyperref[sec:ZCODt6kOCdM]{\small{> > > NEXT SERMON > > >}}
\newline
\newline
使徒行傳 13:1-3
\newline
\begin{longtable}{cl}
\hline
\hline
章節 & 經文 (和合本修訂版)\\
\hline
13:1 & \begin{tabularx}{0.7\textwidth}{X} 在安提阿的教會中，有幾位先知和教師，就是巴拿巴和稱為尼結的西面、古利奈人路求，與希律分封王一起長大的馬念，和掃羅。 \end{tabularx} \\ \\ \relax
13:2 & \begin{tabularx}{0.7\textwidth}{X} 他們在事奉主和禁食的時候，聖靈說：「要為我分派巴拿巴和掃羅去做我召他們做的工作。」 \end{tabularx} \\ \\ \relax
13:3 & \begin{tabularx}{0.7\textwidth}{X} 於是他們禁食禱告後，給巴拿巴和掃羅按手，然後派遣他們走了。 \end{tabularx} \\ \\
[1ex]
\hline
\hline
\end{longtable}
$^{1}$兩位姐妹.
我們其中一位同工正在跑馬拉松.
打氣,當然不是我.
跑馬拉松要跑二十多公里.
我兩分鐘也跑不到.
C同工現在應該正在跑.
有他的消息嗎.
沒有,好消息就是好消息.
他正在努力,我為他祈禱.
希望他完成這個賽事.
上星期我們很忙,大家都很忙.
因為有報道會和很多邀請聚會等等.
今天來到教會時.
有些新朋友也來到我們中間.
上星期來到的,我們感恩.
雖然上星期一週的報道很忙.
而且很辛苦.
大英姐妹在台前幕後.
還有陪探員,報道隊的大英姐妹.
上星期…不,昨天.
昨天我們的探訪隊,報道隊.
也有出隊,他們很勤力.
我覺得是很好的.
雖然很累,對吧.
雖然很累,但我們睡一,兩天就行了.
我們又繼續上路,對吧.
教會應該是這樣.
教會我們需要繼續向前.
這個年頭.
我們說在UV的第九周年.
我們的年提就是跨越.
在這裡提醒大家和挑戰大家.
有幾方面我們需要踏出.
如果你今天第一次來我們教會.
有很多事情可能還未認識.
或者想去摸索等等.
不要緊,我們大英姐妹很歡迎大家.
無論在信仰上,我們需要去追尋.
去明白這位上帝.
教會是一回甚麼事等等.

$^{41}$教會生活是怎樣.
基督徒個人信仰的生活又是怎樣.
有很多問題和需要我們去掌握.
不過要花時間.
記住上帝的恩典是和我們一起.
所以今年我們希望提醒大家.
我們的生命需要跨越.
離開我們舒適的地方.
和其他大英姐妹一起同行.
每位大英姐妹,包括我在內.
都需要其他生命說出我的生命.
和我們有生命上彼此的鍛鍊.
第三方面就是傳道,報道.
我們的生命就是上帝的見證.
在今年,我提到的題目是「士徒行傳」.
大部分時間都會說「士徒行傳」.
你記得「士徒行傳」開始時.
提到我們每一個都要做一個見證.
由我們自己的本位直到地極.
我們都需要做一個見證.
所以在「士徒行傳」.
我們有很多篇幅提到第一代的教會.
他們有甚麼經歷和故事.
所以我們一直在說.
現在已經到10月了.
對,已經10月了.
已經經不過了,真慘.
我們只說第13章.
「士徒行傳」有28章.
你回到聖經時.
新約第五章才是「士徒行傳」.
「波羅的X」.
「波羅的X」也很長.
因為記載第一代教會有很多不同的故事.
有28章這麼多.
我剛才才想不到.
原來只到第13章,一半都沒有.
無論如何,我們看看這些故事.
第13章是一個很重要的時刻.
對教會,對初期的信徒.

$^{81}$是一個里程碑,為甚麼.
因為你記得在「士徒行傳」中經常說甚麼.
保羅去不同的地方宣教.
起碼記載了他三次宣教的旅程.
最後幾章就說到他上羅馬.
上告該撒將福音帶到當時的京都羅馬.
起碼有三次.
加一次是最後到達羅馬.
第13章是第一次,開始的時候.
所以我們看看這段聖經.
第13章,第1至第3節.
他說在安提阿一個地方.
叫做安提阿的教會.
有幾位先知和教師.
就是巴拿巴和稱為尼潔的西面.
古里內人的路球.
希律,昏風王一起長大的馬廉和素羅.
他們在事奉主和禁食的時候.
聖靈跟他們說.
要為我分派巴拿巴和素羅.
去做我召他們做的工作.
於是他們禁食,禱告之後.
把巴拿巴和素羅,或者保羅.
按手,然後就差派他們走了.
這段短短三節的經文有甚麼好說.
正如我們所說.
這是第一次教會.
即初期教會到達一個新發展的階段.
這三節完結後.
就是保羅和巴拿巴做第一個宣教.
這裡提到一個教會叫做安提阿.
如果你看後面的地圖.
大概在右邊,看見嗎.
今天的敘利亞,在右邊.
是一個大城市或大都會.
當時的東地中海世界.
非常大的都會.
當然大都會就想起甚麼.
有不同的種族,人.
有不同的文化等等.

$^{121}$我們這段經文看到一個這樣的特色.
有不同的種族,文化.
這裡有幾個名字,看見嗎.
第一次看聖經的朋友.
你會覺得有點陌生.
這些名字是甚麼名字.
不是我們平時叫的名字.
因為當時是猶太人的名字.
你看見有好幾個名字.
第一個叫做巴拿巴.
在《士郎傳》中很出名的人.
他本身是猶太人.
在居比路或塞浦路斯出生.
第二個叫西敏.
他的綽號叫利潔.
利潔,Niger.
甚麼意思?不要亂發音.
Niger….
真的有這個字.
他的意思是甚麼.
就是黑的意思.
黑,好嗎.
意思是黑面神.
即是黑鬼.
黑鬼西敏,有個黑鬼西敏.
第三個就是古利內人,Sirene.
Sirene在哪裡?在地圖下方.
有一個地方,Dung Dung的地方.
就是一個海港,叫古利祿.
在哪裡?屬於五大洲哪個洲.
很難嗎?非洲,在北非,可以嗎.
回去讀小學時讀地理,好嗎.
這個人在北非,叫路球.
然後有一個叫馬廉.
馬廉介紹他時.
說他跟希律….
分封王希律王是猶太人的王.
一起長大的,即是甚麼.
即是皇家學校長大的.
皇家學校一起長大的.

$^{161}$即是跟他同學.
即是有一個官的同學.
有一個官叫希律王,分封王的官同學.
這個叫馬廉.
即是在這個名單中.
有不同地區來的人.
亦有不同社會經濟等級的人.
走在一起.
然後還有索羅或保羅.
索或Paul.
這個就是猶太人.
剛剛信主,不久後.
上次也說到他們如何信主.
你看到一個多元和合一的圖畫.
我相信路加記載安提柯教會時.
他特別想凸顯這件事.
上帝所看見的安提柯教會.
一個他覺得很棒的教會.
是一個這樣的教會.
是一個有多元,有不同聲音.
有不同的人類組成的.
教會就是應該是這樣,對吧.
當然不是說侍王主所有的東西.
寫了之後,我們就要複製.
就要學習,不是這個意思.
而是這裡告訴我們一個圖畫.
是啟發我們每一代的信徒.
我們看看教會.
其實教會就是應該是這樣.
一個合神心意的教會應該是這樣.
特別是一個大都會的教會.
應該容納到不同的聲音.
容納到多元的一個群體.
教會其實應該是這樣.
UV其實也應該是這樣.
我們教會確實需要學習.
如何跨出我們自己的小圈子.
跨出我們自己所熟絡的東西.
所熟習的人,群體,關係.
對嗎.

$^{201}$所以應該有老中青.
也有不介意.
起碼我們不介意或是歡迎.
有不同的人來到我們中間.
就算他們的語言.
就算他們的文化不同我們.
就算他們的文化.
令我們覺得很不舒服.
我們也應該歡迎他們.
這是上帝的指紋,是教會.
不是因為你喜歡.
而是因為上帝喜歡.
如果上帝喜歡.
我覺得這就解決了問題.
對吧.
我們要學習如何看到上帝所看的東西.
所以第一方面,這裡所揭示的教會.
當他們還沒開始.
普世宣教的運動之前.
他們是一個容納到多元的群體.
並且他們很著緊在上帝的道.
這裡介紹他們的組成時.
他說開始有幾位先知與教師聚在一起.
然後給我們一個名稱列表.
你知道在教會也有很多不同的恩賜.
在新約的書信或我們教會的經驗中.
我們也會看到有人有不同的恩賜.
有不同的功能.
有些人可能音樂較好.
有些人可能另一方面較好.
有些人可能廚藝較好等等.
我們大家也有不同的恩賜.
在這裡我看的時候,我覺得很奇怪.
為何會特別強調這一點.
為何不說有幾位廚師.
幾位做甜品的弟兄姊妹.
或有幾位掃地的在這裡.
幾位先知與教師,為甚麼.
教師當然很清楚,教師是做甚麼的.
教會裡的教師,包括我自己在內.

$^{241}$就是教導信徒,明白信仰.
以及如何實踐信仰.
這是教導,指導等等.
先知是甚麼?先知需要解釋一下.
你第一次看聖經,看到有先知.
我們便想起微卜先知.
或鐵板神算等等.
其實在聖經裡提到先知.
不是說這些的.
先知不是說這些.
聖經裡的先知是甚麼.
聖經裡的先知是宣講上帝權威的說話.
這叫先知.
即是代表上帝來通話.
所以你看到有教師,有先知.
共通點是甚麼.
教會以上帝的道為中心.
當他們走在一起的時候.
聖靈便開始去工作.
聖靈開始去推動他們.
當他們走在一起.
他們看重上帝的道.
以及明白上帝的道.
大家便開始明白聖靈的說話.
聖靈對我們說甚麼.
以及要做甚麼.
所以弟兄姊妹.
為何會在教會查經.
為何我們會有主學學.
為何會在團契看上帝的說話.
或有靈修等等.
這是我們的中心.
我們快要完結「生命之道」.
對吧….
我知道開始的時候有很多人.
我們快要完結了.
有弟兄姊妹跟我說.
「完結了又怎樣」.
我便說「不如我們自己做」.
好嗎 好….

$^{281}$對嗎.
別人的東西也好.
但我們自己煮的東西應該是最好的.
(笑).
為何要教呢.
不是為了教而教.
我覺得弟兄姊妹走在一起.
我們一起去查經,討論時.
真的有很多學習.
我自己也有很多學習.
好嗎.
以道為中心的教會.
是上帝覺得很棒的教會.
接著他們走在一起時會做甚麼.
事奉主及禁食時.
聖靈開始跟我們說話.
好嗎.
禁食.
這裡沒有詳細說甚麼是禁食.
因為猶太人有很多禁食習慣.
有些是禁食,有些是不禁食.
就像我們今天這樣.
我們在飲食中.
有些是168,有些是Keto.
有些是不知甚麼飲食.
每個人都不同.
猶太人有不同的禁食習慣.
沒有說,但最主要說甚麼.
猶太人禁食的意思是.
有時可能是三餐變兩餐.
或者不吃硬食物等等.
猶太人的禁食習慣.
但最主要的焦點是甚麼.
為何要禁食.
禁食就是要專心.
專心做甚麼.
去尋求上帝.
去操練我們的心智.
這是禁食的精神.
所以有時他們會輕一點的禁食.

$^{321}$會重一點的禁食.
會大一點的,有些會小一點的.
但最主要是甚麼.
我們一起去禱告.
一起專心尋求上帝的心意.
這才是最重要的.
是嗎.
你想想,這很自然.
禁食才可以專心.
前一星期我在多倫多.
去那裡講一些講座分享.
有些教會說.
我們去退修會會去哪裡等等.
離開城市去別的地方.
他們問教會.
教會會叫你們去哪裡退修會.
我說我們想過很多次.
有很多選擇,我們不知道去哪裡.
我跟他們說.
很坦白地跟多倫多的弟兄姊妹說.
我們芝加哥教會的弟兄姊妹很好.
他們建議我們芝加哥教會退修會.
去多倫多.
他們當然歡迎.
但你想想,你去多倫多.
你怎會專心.
你不可能在那裡專心.
真的不可能專心.
除非你是我,好嗎.
我會很專心.
當然他們真的歡迎.
他說你建立美食團.
如果建立美食團,就不是退修會.
對,每個人都吃飽了.
但沒有聽過道.
每個人都記得吃了甚麼好東西.
但不知道吃了甚麼.
禁食的重點就是這樣.
就是要專心.
我們需要專心.

$^{361}$所以千萬不要去我們.
我們千萬不要去.
他….
你自己去,好嗎.
事奉主的時候,聖靈跟我們說話.
「事奉」這個詞.
很多時候用在聖殿的工作.
祭司,神職人員.
在舊約有聖殿的時候.
他們會打理聖殿的會幕工作.
其實是敬拜的意思.
所以ESV有另一個英文的本.
叫做ESV,English Standard Version.
就是Worshipping the Lord.
他們禁食,專心.
來到去事奉的時候.
去敬拜的時候.
他們就聽見聖靈的聲音.
有想過我們進來敬拜的時候.
不只是唱幾首歌這麼簡單.
亦不只是我們想拿東西走.
敬拜的時候是跟上帝相遇的.
上帝說話很多時候在敬拜裡相遇.
對你說話.
我相信你來了教會一段時間.
如果你真的專心來.
來敬拜的時候.
我肯定大大小小的體會.
大大小小的那些時刻.
你都會在敬拜,詩歌,祈禱,聖餐.
講道,分享等等.
你都會聽到上帝跟你說話.
哪怕只是一句,兩句說話.
有時候甚至是….
就像聖靈在你的生命裡.
推一推你去做一些事情.
但問題就是群體,那些人.
是否願意去聽.
是否願意去尋求.
去跟上帝說.

$^{401}$你今天跟我說話.
我正在聽你說話.
反正我們不願意.
第13章第二節.
聖靈說甚麼.
他說你們這些人要做一件事.
要為我做一件事.
本身安特拉教會也不錯,對吧.
但他說你們要做一件事.
拆敗別人.
有時候是的.
聖靈好像是我們生命裡的一個推動力.
去推動我們.
有點像我想起小時候.
你教他踩單車.
你不是一直扶著他,追著他.
腰骨也痛了.
一直推著他,扶著他.
當你覺得是時候了.
你就去推一推他.
對吧.
你想像一下,就是這樣教.
或者教他走,慢慢放手.
然後推一推他.
聖靈在我們生命裡.
很多時候都是做這些事情.
但小朋友也要繼續踩單車.
他不是以防我.
但當他熟習了.
我就會推一推他,推他一下.
然後他自己就會高飛,對吧.
他就可以純粹推走,然後繼續走.
在我們生命裡其實也是一樣.
有點像上帝和我們.
好像把我們當成小朋友.
聖靈有時候會跟你說.
推一推你,是時候信主了.
你不要消滅他的說話.
有時候可能他跟你說.
是時候你多走一步.

$^{441}$是時候你應該有新的嘗試.
不要等到外在的因素全部準備好.
去吧.
這就是聖靈在安徒墓教會裡所做的事.
事奉主,敬拜主的時候.
聖靈跟我們說要做這件事.
但他們需要有一個預先條件.
就是他們願意聽上帝的話.
願意去敬拜上帝.
願意去聽祂的話.
第三方面,這裡看到.
這個教會是一個向外的教會.
並且敢去到,差派工人出去.
第三節.
然後他們禁食,專心禱告之後.
就被巴拿巴和蘇羅按守.
然後差派他們走了.
離開了.
照上帝的心意.
照上帝的呼召.
被上帝推一下.
然後讓他們去.
這裡所有人去看這段聖經的時候.
包括聖經家,學者.
都會有這樣的評論.
這裡有13章.
就是巴拿巴和蘇羅.
都是當時最優秀的墓者.
是最好的墓者.
剛才你看到上面的清單.
有黑鬼西勉,其他人.
但聖靈卻要叫他們差派誰.
巴拿巴和蘇羅.
是最優秀的墓者.
叫他們做甚麼.
不是升職.
讓他們去.
叫他們走.
這裡每個人都會有這樣的評論.
一個初生的安提拉教會.

$^{481}$他們最需要的是甚麼.
就是最大的人.
我們每個群體都需要一些領導者.
都需要一些我們覺得.
可能有很多恩賜的人去帶領.
這是很自然的.
但這裡聖靈做的事是相反的.
叫他們差這兩個人是最好的.
做其他事.
我在教會已經很久了.
去過很多不同的教會.
我真的全職事奉.
差不多30年了.
確實教會裡的經歷.
就像坐貨車一樣.
有人來.
你很高興.
亦有很多時候.
有些人是要走的.
聖靈帶他走.
即是上帝要帶他走.
這個很多時候.
多年來很多時候都要學習.
有時候會看法.
但我又不想看法.
我想盡量在上帝的角度.
上帝的視野裡.
去衡量這些去與留.
哪怕有時候是挺難受的.
但上帝有他自己的主權.
我應該怎樣去馴服.
怎樣去信任.
有時候是難的.
因為我只是一個很渺小的人.
不知道上帝的大計劃是怎樣.
但只能夠肯定一件事.
上帝大的計劃.
一定是大過我最大的計劃.
上帝的計劃一定是比.
我自己覺得最好的計劃好.

$^{521}$對嗎.
不只是教會.
其實你人生也是.
我們都需要有這樣的領會.
我會為我分派巴拿巴和保羅.
去做我差.
我召他們所做的工作.
是上帝說我有事給他們做.
不是你有事給他們做.
我記得很多年前.
我聽了一個中國教會歷史的講座.
一個講座.
是一個專門研究中國教會歷史的學家.
他去研究特別是二十世紀中.
二次大戰之後.
中國教會的擴散,工作等等.
他收集了很多這樣的資料.
然後告訴我們.
中國教會在二次大戰之後.
好幾十年.
改朝換代等等.
有很多不同的改變.
開放不開放等等.
收緊不收緊.
不同的政策下.
教會如何在那種很艱難的情況下.
來不斷開發.
他收集了很多這樣的資料.
宣導會做宣導會的事.
浸信會又做浸信會的事.
那些比較年長的記得戴德生.
內地會又做內地會的事.
他就說中國教會.
內地會又做內地會的事.
因為講完之後.
整個講座之後就有QA.
QA的時候有人問了一個問題.
我覺得這個問題不錯.
我的感覺就是.
中國教會二次大戰之後.

$^{561}$他的福音工作是很各自為政.
是嗎.
浸信會做的事和宣導會做的事.
大家不聊天.
你有你做,我有我做.
有時可能做同一樣事.
是嗎.
他就說這樣好嗎.
各自為政.
其實是否應該要有一個宏大的策略.
有一個大策略.
是否會更加好.
我預期應該是這樣.
是嗎.
但那個學者最後說.
他說在他的認識裡.
沒錯,有時宏大的計劃.
都可能需要的.
但往往很多時.
在福音或教會的宏大的計劃.
都是空談的.
計劃是需要的.
安排是需要的.
我們探訪,傳播.
我們都需要一些安排.
安排不好.
像螺蟹般不行.
是嗎.
但始終回到人的想法.
人有多大的計劃.
有時是空談.
結論就是.
其實我們每一個人.
做回自己的小角色.
我們忠心地去履行我們的責任.
聖靈就會在背後策劃.
所以有時我們需要學習的.
就是我們自己的想法.
教會自己的想法.
是不需要用原子筆寫的.

$^{601}$用原子筆寫,策劃得到.
如果要策劃聖靈.
就讓聖靈去策劃.
願意去拆遣工人.
願意放棄一些我們不想放棄的.
願意在我們自己的小角色裡.
盡力做好自己.
我們深信上帝的計劃.
一定比我們高.
哪怕我們不明白.
但要求於我們的.
就是願意去看重它的道.
願意去信服.
願意聽上帝的聲音.
最後,今天希望你記起幾件事.
在你的生命裡.
問問自己是否願意去聆聽.
學習用我們的耳朵去跟隨上帝.
用我們的耳朵.
在敬拜裡聽聽上帝跟你說甚麼.
在祈禱裡聽聖靈對你有甚麼推動.
期待,我們的心態不是期待.
期待聖靈在我們生命裡工作.
有時可能是痛苦.
有時可能是很開心.
有時上帝可能打斷我們的計劃.
或者是最美好的事.
但起碼我們有期待的期望.
上帝是你願意做的事.
第三,信心.
信得過上帝的智慧.
我們願意踏上去.
忠心做好我們的小角色.
我們不用操控所有東西.
儘管我們放長雙眼.
有時放長雙眼,我們回頭看.
上帝的計劃比我們高.
好嗎?我們同心,低頭祈禱.
預備我們領受聖餐.
因為你愛我們.

$^{641}$所有微小的人的上帝.
你在歷史中是主宰.
你在我們生命中也是主宰.
主宰,你來我們面前為我們降服.
因為看到你自己的大能.
體會到你自己的愛.
但你卻去看顧,去愛我們.
你十字架上的恩典.
讓我們能成為你的兒女.
主一靠你帶領我們.
我們的生命能夠聆聽你的說話.
期待你的工作.
並願意放長雙眼去信靠你.
知道你自己的計劃.
比我們一切最好的計劃更好.
在我們今天這個主日中.
我們領受這個聖餐.
你釘心,願意救贖我們這群不配的人.
讓我們這個新的心靈.
來領受你的恩典.
祈禱奉主彌禱.
這個聖餐是為了你.
\newpage



\section{使徒行傳 16:6-10}
\label{sec:ZCODt6kOCdM}
\textbf{不似預期的福音 (使徒行傳 16\_6-10) - 陳崇基牧師}
\newline
\newline
連結: \href{https://youtube.com/watch?v=ZCODt6kOCdM}{\texttt{https://youtube.com/watch?v=ZCODt6kOCdM}} ~~~~ 語音日期: 2025-10-29
\newline
\newline
\hyperref[sec:jhQYPPMz5JE]{\small{< < < PREV SERMON < < <}}
~
\hyperref[sec:index_chronic]{\small{[返順時目]}}
~
\hyperref[sec:index_scriptual]{\small{[返順卷目]}}
~
\hyperref[sec:Qv1tOecnO_s]{\small{> > > NEXT SERMON > > >}}
\newline
\newline
使徒行傳 16:6-10
\newline
\begin{longtable}{cl}
\hline
\hline
章節 & 經文 (和合本修訂版)\\
\hline
16:6 & \begin{tabularx}{0.7\textwidth}{X} 因為聖靈禁止他們在亞細亞講道，他們就經過弗呂家、加拉太一帶地方。 \end{tabularx} \\ \\ \relax
16:7 & \begin{tabularx}{0.7\textwidth}{X} 到了每西亞的邊界，他們想要往庇推尼去，耶穌的靈卻不許。 \end{tabularx} \\ \\ \relax
16:8 & \begin{tabularx}{0.7\textwidth}{X} 他們就越過每西亞，下特羅亞去。 \end{tabularx} \\ \\ \relax
16:9 & \begin{tabularx}{0.7\textwidth}{X} 夜間，有異象向保羅顯現。有一個馬其頓人站著求他說：「請你過來，到馬其頓來幫助我們！」 \end{tabularx} \\ \\ \relax
16:10 & \begin{tabularx}{0.7\textwidth}{X} 保羅既看見這異象，我們就立即設法往馬其頓去，認為神呼召我們傳福音給那裡的人。 \end{tabularx} \\ \\
[1ex]
\hline
\hline
\end{longtable}
$^{1}$歡迎各位回來.
我們繼續我們的系列是講到跨越的.
我們一直在找一些《士徒行傳》裡面的故事.
如果你第一次接觸聖經的話.
《士徒行傳》是一本怎樣的書呢?.
聖經是有舊約也有新約的.
舊約就是耶穌基督來之前.
預言耶穌基督會來到的那一批的書.
然後新約就是耶穌基督來了之後.
就是祂做了些什麼.
在新約裡面第五本書就是《士徒行傳》.
包括《X》.
我們找一些故事去講一些當時初期教會發生了什麼事.
你知道耶穌基督死而復活.
升天之後.
他的初代教會在耶路撒冷開始.
來到不同的地方傳福音.
所以《士徒行傳》是講到初代教會有些什麼經歷.
還有聖靈在他們中間做了什麼工作.
所以我們一直在看《士徒行傳》.
我們說過上一次我們講到第十三章的時候.
就是講到保羅宣教的旅程.
宣教旅程還沒開始之前發生了什麼事.
你記得嗎?.
安特拉教會有兩個最好的傳道人.
被差派出去.
巴勒巴和保羅.
他們按照地圖上的路線.
從安提拉看到嗎?.
看到吧.
在右邊.
安提拉看到嗎?.
可以嗎?.
好像還沒睡醒.
他們坐船過了今天的塞浦路斯.
就是塞浦路斯或者居比路這個地方.
在那裡傳了福音之後.
他們在島上.
然後去到下面有一個地方叫帕弗.
看到嗎?.

$^{41}$帕弗.
當時的帕弗是一個碼頭.
那個碼頭會坐船去不同的地方.
但現在我們所知道的.
就是帕弗通常坐船去南面.
看到下面埃及.
阿克山大馬.
看到嗎?.
我們要找回當時第一世紀的文獻.
我們知道帕弗是一個什麼樣的碼頭.
就應該是向南的.
坐船去阿克山大.
這個合理.
坐船去阿克山大.
為什麼呢?.
因為阿克山大是當時猶太人其中一個.
在地中海世界裡面最大的一個居留的地方.
最大的一個群體的地方.
所以他們去到帕弗的時候.
按道理應該是坐船.
截船就下去下面南面.
但有些原因.
就不是很按照原本那個人的計劃.
我在香港來的.
香港去某一個碼頭.
一定是坐船去某一個地方的.
如果去港澳碼頭.
當然我講回很多年前.
我沒有回香港很多次.
很久了.
去港澳碼頭的時候.
你沒有坐船去長洲的.
對不對?.
但偏偏.
你應該是去澳門的.
但你有一艘船.
就去了長洲.
就是這樣的情況.
帕弗應該是向南的.
但他們坐了船.

$^{81}$就上了去別加.
然後在龐非尼亞等等的地方.
上面也有另外一個安提阿.
特比等等的地方.
然後傳福音之後.
他們就回到耶路撒冷.
這個是第一次保羅的宣教旅程.
在《少林傳》裡面有三次講到.
保羅的宣教旅程.
上一次我們提了這個宣教旅程.
這次我們提第二個宣教旅程.
第二個宣教旅程發生了什麼事呢?.
又好像有一些不像人預期.
或者不是人控制到的事情是發生的.
我們先看經文.
在左下角.
就將第六節.
那天聖靈禁止保羅他們的同工.
在亞細亞道港道.
他們就經過佛里加.
加拉提一帶地方.
到了美西亞的邊境.
他們想要往彼推尼.
耶穌的靈卻不許.
他們就越過美西亞.
下到特羅亞去.
這裡全部字我都看得明白.
但是看了之後就不太明白.
其實有很多地名在這裡.
是不是?.
他說聖靈禁止他們在哪裡的港道.
亞細亞.
亞細亞就是今天的土耳其.
大概西面那一塊.
禁止了他們在亞細亞.
亞細亞道港道.
亞細亞當時是一個省.
那他們怎麼辦?.
他們就經過什麼地方?.
佛里加和加拉提.

$^{121}$看到嗎?.
他們就屈了上去.
去不了亞細亞.
按道理.
那條線應該是.
如果按照人的計劃.
他們應該是直去爾佛索的.
因為爾佛索是一個大都會.
是今天的土耳其.
當時的所謂小亞細亞.
亞細亞整個.
那個最大最大的城市.
按照保羅的理念或者他的做法.
通常是先去一些大的城市.
大的都會.
但是他在這裡說什麼?.
他說.
聖靈禁止他們在亞細亞港道.
所以他們就屈了上去.
上到佛里加和美西亞.
上到上面.
之後發生什麼事?.
到了美西亞的邊境的時候.
他們就想去到比推尼.
看到比推尼嗎?.
在北面.
即是今天土耳其的北面.
Bamphilia.
想上去.
但是怎樣?.
耶穌的靈卻不許.
即是聖靈卻不許他們.
於是他們就上不到北面.
他們就下來.
就好像被迫一樣.
他們就要經過美西亞.
即是今天土耳其左上角的省份.
他們被迫要經過美西亞.
然後就下去一個地方叫做特羅亞.
看到嗎?.

$^{161}$看到特羅亞.
特羅亞又是一些大的海港.
即是好像香港.
新加坡等等.
迫他們走到特羅亞的地方.
這就是大概經文頭三節所交代的六七八節.
發生什麼事呢?.
他們本身有一個計劃A.
計劃A當然是去當時土耳其的南面.
去伊弗所,龐菲尼亞,亞細亞等地方.
這是他們的計劃A.
但是發生什麼事?.
是聖靈禁止他們.
可能原定的計劃是真的想去伊弗所.
因為那裡是最大的工場.
所以去那裡宣教,傳福音應該是最有利的.
但是誰知道神的靈.
神就不讓他們去.
然後他們就有計劃B.
計劃B簡單來說就是向北面.
他們想往北面.
從來沒有想過要往西面.
但是當他們計劃B想向北面給他們的時候.
聖靈又禁止他們.
所以就路路不通.
然後就被迫去了西面.
去了楚沃斯.
或者德羅亞的海港.
德羅亞的海港我剛才說過.
好像香港,新加坡一樣.
也是當時歐亞之間大的港口.
這個很神奇,很有趣.
一個這樣的故事.
打開小行傳的時候.
我們的感覺就是.
第二章五順節.
聖靈做的事情是很挖破時代的事情.
使使徒可以講方言.
使很多人信主.
是一種很猛烈的表現.

$^{201}$剛剛這個時代開始的時候.
天國開始在地上藉著教會來進行.
聖靈就像火,水一樣去降臨.
然後他們講方言.
然後去到不同的地方傳福音.
但慢慢在小行傳裡面介紹聖靈.
或者去到講出聖靈做的事情的時候.
就慢慢好像沉靜了,沉實了.
然後做一些事情是比較.
好像這個故事一樣.
不是很跟人的想法.
所以這段經文裡面告訴我們.
就是有些事情是他們本身有計劃的.
但聖靈的工作令他們做了一些事情.
或者改了道.
是不像人的預期.
去到特羅亞發生什麼事呢.
就有一個不似人預期的呼召出來.
特羅亞發生什麼事呢.
第九節.
你記得他們上南面.
然後上北面.
兩邊都不通.
然後他們被迫往西面去了特羅亞.
第九節.
在特羅亞發生什麼事呢.
他們過夜的時候.
夜間有異象向保羅顯現.
有一個馬其頓人站著求他說.
請你過來.
到馬其頓來幫助我們.
保羅既看見這異象.
我們就立即的設法往馬其頓去.
認為上帝是呼召我們傳福音及那類的人.
馬其頓在哪裡.
你找找.
馬其頓就在今天希臘的北面.
看到嗎.
在地圖的北面.
好嗎.

$^{241}$所以他們在特羅亞的時候.
他們就看見這個異象.
這個《士多行傳》裡面最有名的一個故事.
就是這個故事.
就叫做馬其頓的異象.
Macedonian Vision.
什麼意思呢.
就是說他在特羅亞.
那個異象裡面就是.
他看到對面有一個馬其頓人.
Macedonian.
來求他.
求你過來幫我們行不行.
就是看著一個海港.
看著對面.
就有一個馬其頓人.
一個西方所謂歐洲的人.
來求他們過去.
他們作為東方人.
他們是東方人.
是從東方過來的.
求他過去來幫助他.
馬其頓有什麼特別呢.
馬其頓有什麼不特別.
如果熟知當時的世界的話.
那個背景的話.
馬其頓是一個很威風的人.
為什麼.
因為你記得.
有一個大將軍.
叫做亞歷山大.
亞歷山大大概在主前三三三年的時候.
他是一個馬其頓人.
是一個非常懂得征戰的大將軍.
他就帶著他的希臘弟子.
往東邊來入侵當時中東的世界.
就是波斯帝國.
波斯帝國是很厲害.
當時是世界的超級力量.
波斯貓很大隻.

$^{281}$但是這個馬其頓的將軍.
亞歷山大.
他就帶著他的弟子.
跟著這個戰爭由左邊一直打.
打了十年之後整個的波斯帝國.
就被他吞併了.
這個就是馬其頓人做的事情.
這個是馬其頓人有的威勢.
如果你上過我的課的話.
你都知道.
有一幅馬賽克圖.
是一些模式.
在地下的模式.
好像現在我們的地氈一樣.
他們就將亞歷山大大帝.
去到打波斯帝國.
當時一個最後的君主.
當他攻擊他們的時候.
那個威勢你都可以描述出來的.
當然那個馬賽克圖.
已經是有一些破損的地方.
但是你看到亞歷山大在哪.
在左邊.
然後右邊那個是誰.
右邊那個就是波斯王.
你看到那個慌忙的樣子.
亞歷山大由左邊一直奔前.
然後打那邊.
那些波斯的弟子.
或者波斯的士兵.
他們就逃亡.
如果你放大一點看一下.
這個亞歷山大大帝.
是不戴頭盔的.
波斯王就戴頭盔.
還是像桑加狗一樣走.
我們說話都戴頭盔.
他不用戴頭盔.
他上陣出征的時候.
是不需要戴頭盔的.

$^{321}$就是一種這樣的威勢.
這是馬其頓人的厲害的地方.
希臘人厲害的地方.
並且因為亞歷山大統治了天下之後.
他就將希臘的文明帶到世界各地的地方.
是他去傳這個希臘的福音.
什麼意思.
所有我們今天很多的政治的思想.
民主.
同盟的概念.
哲學.
建築等等.
語言等等.
都是因為這個馬其頓人.
因為這個希臘人.
而有的貢獻.
大家有沒有去downtown那裡開車.
之後在downtown開車.
看到一些很漂亮的建築.
看到嗎.
film museum門口的柱子.
或者有些doric column.
有些corinthian columns.
全部都是馬其頓或者希臘出來的.
所以是他們去幫整個世界.
但是在這個馬其頓的意象裡面.
就顛覆了過去.
以前就是所謂西方的人.
馬其頓希臘的人.
來幫所有其他人.
現在是一班東方去的小宗教的隊伍.
保羅.
帶著他們去童宮.
在這個意象裡面.
反過來.
馬其頓人說.
你過來幫我.
因為我們需要你幫助.
結果就是這班由東方而來的基督徒.
帶著福音進去歐洲.

$^{361}$去幫馬其頓人.
甚至其他的地方.
這個真的很神奇.
從來都是西方人幫東方人.
對嗎.
但是在意象裡面.
是東方人幫西方人.
幫什麼.
是將福音.
將天國.
帶到他們中間.
最近有一個很出名的英國歷史學家.
叫做Tom Holland.
Tom Holland.
H-O-L-L-A-N-D.
Tom Holland.
你自己Google一下.
不過Tom Holland也是一個演員的名字.
所以你沒有搞錯.
是一個歷史學家.
英國的一個歷史學家.
大概是我的年紀.
這個人很有名.
他寫了好幾個研究的結果.
他說西方的文明.
是因著什麼而建立的.
他說就是因著基督教.
基督教的概念.
基督教的價值.
使得西方文明能夠興起.
我們大家都很清楚.
歐洲以往.
他的文明的興起.
歐洲的文化等等.
西方人的文化.
Tom Holland就說.
如果不是因為基督教的緣故.
西方是不會崛起的.
是因為基督教.
lay down了這個foundation.

$^{401}$以致Western civilization.
能夠去影響世界不同的地方.
儘管看一下Tom Holland所說的.
我不知道他信了主沒有.
好像信了.
但他寫這些東西的時候.
他是純粹用一個歷史學家的身份來說.
這個就是馬其頓的現象.
因為你看到這個很短的故事裡面.
看到聖靈的工作.
或者神所做的事情.
很多時候不是我們能夠去掌握的.
不是我們能夠去控制的.
是真的.
在這幾年裡面.
我們UV都發生很多事.
我預備這個講章的時候.
我想起Ken和Alice剛開工的那天.
記得嗎.
Alice說到的.
她說到車子等等的事情.
是很溫馨的故事.
然後完了之後.
沒兩個小時車子就沒了.
我自己都覺得很失控.
如果你認識我的話.
你都知道我是很怕失控的人.
我真的很不喜歡那種失控的感覺.
每樣事情都覺得應該要讓我掌控得到.
我才有那種安全感.
你坐過我的車嗎.
我是皮膚開車的那個.
這個是我一向都有的掙扎.
往往有時候在福音的事工.
或者教會的事工.
是有一些事情我自己掌控不了.
但就令我有很多掙扎.
不是我預期的.
上帝要教我.
你要相信.

$^{441}$你要相信聖靈的工作.
神在你身上的工作.
往往不是像我們的預期.
是不是.
不像我們的預期.
有一句英文說話.
The spirit's know often sets up a greater yes for his mission.
The know of the spirit often sets up a greater yes for God's mission.
對不對.
剛才還沒開始聚會的時候.
有個姐妹說到她的故事的時候.
她說到一句這樣的說話.
她說有時候神對我們帶領.
人生的帶領.
是令我們摸不出頭腦的.
我們真的猜不到.
有時候我們都會問上帝.
即是睡不著的時候.
我們問上帝到底來教哪一科.
到底來教哪一科.
我每樣東西都想為福音.
為福音的工作.
為教會的工作.
但往往上帝的帶領.
往往上帝的引導.
是過於我的預期.
亦不是我能夠控制到的.
想留的人未必留到.
想發生的事未必發生到.
即儘管我費很多的精力.
費很多的能力.
我想去控制那件事.
往往就不是我能控制到的.
It's the know of the spirit.
但這句說話提醒我們.
There's always greater yes.
而那個greater yes需要等候.
需要內心的信服.
才能夠見到.
有時候失控的時候.

$^{481}$我自己都會想.
自己可能更加清楚想清楚.
到底上帝對我的呼召是甚麼.
上帝對我的要求是甚麼.
其實上帝對我的要求.
不是要我去操控所有的東西.
要所有東西按照我的預期.
來對應.
這不是上帝給我的呼召.
上帝給我們人生的呼召就是.
You be faithful.
你很渺小.
可能很多東西你控制不到.
不過你做回那個部分.
有一個社會學家曾經用一個研究.
去看當時初代教會為甚麼能夠興起.
那本書或研究叫做.
Rise of Christianity.
Rodney Stark.
Rodney Stark不是用一個神學的眼光去看.
為甚麼初代教會.
可以在一個最不利的環境裡.
它能夠去增長.
能夠去到其他地方.
去建立不同的群體.
直到今天.
Rodney Stark是一個社會學家.
他就說.
我用一個社會學的角度來看.
The Rise of Christianity.
為甚麼一個所謂很陰暗.
一個很軟弱的.
一個沒有人會理會的.
一個所謂的Jesus Movement.
一個耶穌的運動.
福音的運動.
為甚麼會變成在世界裡面.
只是幾百年的時間.
當時初代教會幾百年三百年的時間.
能夠成為一個不能夠去磨滅的一種力量.

$^{521}$福音的力量.
用一個社會學的角度去看.
甚麼意思.
不要只是說神蹟.
不要只是說聖靈去到工作.
當然這個是我們要去到認同的.
上帝的工作.
上帝打開門.
所以教會能夠生存.
能夠去成長.
能夠去擴張.
這些全部都是神蹟.
但社會學的角度怎麼看.
就是找一些社會學裡面有些原因.
去追溯一些因素.
以至一個這麼邊緣化的福音運動.
能夠短短幾個世紀裡面.
成為一個世界主導的宗教力量.
他舉了幾個例子.
第一個例子就是.
特別開頭那二三百年.
當時在羅馬帝國裡面.
爆發了幾次的疫症.
COVID.
好像COVID一樣.
有班基督徒他們去到.
所謂不怕死.
他們去擁抱那些.
在街上有病的人.
去洗乾淨他們的身體.
讓他們有尊嚴去死.
去到扶起那些有病的人.
不怕我們去受感染.
所以有很多基督徒.
是因為這樣的緣故去世.
但因為這班基督徒有天國的盼望.
他們覺得這是試煉.
即使很多時候面對著安危的問題.
受感染等等.
他們成為了第一股.

$^{561}$在羅馬帝國裡面.
第一股的所謂社會福利.
醫療.
醫院.
護老院.
他們成為了第一個勢力.
這就是一個所謂社會因素.
不只是說神蹟.
是他們默默地在做一些事.
第二件事是人際關係.
特別是說到姐妹.
那些女人.
女人很有關係.
因為姐妹邀請其他人.
邀請他們的配偶.
邀請其他的朋友.
他們就進入了教會.
然後在教會裡面組成了新的群體.
所以他說基督教的興起.
先為人知的邊緣的耶穌的運動.
如何在短短幾個世紀裡面.
能夠成為主導的勢力.
他說信仰的增長運動.
不是靠神蹟奇事.
亦不是靠一些偉大的港道家.
是靠甚麼人.
靠我們.
無名的信眾.
去吸引一些其他無名的親屬.
朋友來信主.
這就是基督教的興起.
請問我們是否這樣.
當然是上帝要求你和我.
不是我們要控制所有的事.
每件事都要跟著我的計劃.
上帝最終對我們的要求和呼召.
你必須信服它.
堅持它,不論如何.
做回我們的小角色.
一個基督徒去傳福音.

$^{601}$一個基督徒去事奉.
一個無名氏.
我們盡心盡力地事奉.
但願這是我今天的立志,好嗎.
我相信你的人生.
可能最近都有很多覺得失控的事.
誤導你的計劃.
誤導你去控制.
但你記得上帝對你和我的呼召.
你堅持它,要信服它.
我以同心低流祈禱.
上帝,我知道你是掌管全局的神.
在我們人生中有很多大大小小的大靈.
有時很清楚,有時我們模糊不清.
因為我們很渺小.
因為我們不知道將來.
我們也有很多限制.
但上帝,我們知道你仍然是掌管歷史.
掌管我們人生.
為了我們中間每位弟兄姊妹.
我們在這個時間.
最近可能我們遇見很多很混亂的時候.
主,我們求你幫助我們.
保持信心.
盡心地敬拜你,盡心地服侍你.
盡心地將福音帶到我們自己的小角色.
無名氏的小角色裡.
求你使用我們.
主,讓我們離開這個地方.
我們回到我們的崗位.
無論是家庭,是我們的職場.
是我們社會每一個角色.
主,我們求你.
求你幫助我們跟著自己的步伐.
縱使這個步伐不一定是直去.
是會轉左或是右的.
幫助我們保持信心.
我們這樣祈禱,奉主名求.
願上帝賜給大家平安,好嗎.
(待續).

\newpage



\section{士師記 1:1-2:23}
\label{sec:Qv1tOecnO_s}
\textbf{Are We Becoming More Like the World? (Judges 1-2) - Chee Yie Jong 我們是否越來越像這個世界 (士師記 1\_1-2\_23) - 楊志溢傳道}
\newline
\newline
連結: \href{https://youtube.com/watch?v=Qv1tOecnO-s}{\texttt{https://youtube.com/watch?v=Qv1tOecnO-s}} ~~~~ 語音日期: 2025-11-03
\newline
\newline
\hyperref[sec:ZCODt6kOCdM]{\small{< < < PREV SERMON < < <}}
~
\hyperref[sec:index_chronic]{\small{[返順時目]}}
~
\hyperref[sec:index_scriptual]{\small{[返順卷目]}}
~
\hyperref[sec:0vNwJDwewOw]{\small{> > > NEXT SERMON > > >}}
\newline
\newline
士師記 1:1-2:23
\newline
\begin{longtable}{cl}
\hline
\hline
章節 & 經文 (和合本修訂版)\\
\hline
1:1 & \begin{tabularx}{0.7\textwidth}{X} 約書亞死後，以色列人求問耶和華說：「我們中間誰當首先上去攻打迦南人，與他們爭戰呢？」 \end{tabularx} \\ \\ \relax
1:2 & \begin{tabularx}{0.7\textwidth}{X} 耶和華說：「猶大要先上去。看哪，我已將那地交在他手中。」 \end{tabularx} \\ \\ \relax
1:3 & \begin{tabularx}{0.7\textwidth}{X} 猶大對他哥哥西緬說：「請你同我上到我抽籤所得之地，與迦南人爭戰；我也同你去你抽籤所得之地。」於是西緬與他同去。 \end{tabularx} \\ \\ \relax
1:4 & \begin{tabularx}{0.7\textwidth}{X} 猶大就上去，耶和華把迦南人和比利洗人交在他們手中。他們在比色擊殺了一萬人。 \end{tabularx} \\ \\ \relax
1:5 & \begin{tabularx}{0.7\textwidth}{X} 他們在比色遇見亞多尼‧比色，與他爭戰，擊敗了迦南人和比利洗人。 \end{tabularx} \\ \\ \relax
1:6 & \begin{tabularx}{0.7\textwidth}{X} 亞多尼‧比色逃跑，他們追趕他，捉住他，砍斷他大拇指和大腳趾。 \end{tabularx} \\ \\ \relax
1:7 & \begin{tabularx}{0.7\textwidth}{X} 亞多尼‧比色說：「從前有七十個王，大拇指和大腳趾都被我砍斷，在我桌子底下拾取零碎食物。現在神照著我所做的報應我了。」他們把亞多尼‧比色帶到耶路撒冷，他就死在那裡。 \end{tabularx} \\ \\ \relax
1:8 & \begin{tabularx}{0.7\textwidth}{X} 猶大人攻打耶路撒冷，奪取了它，用刀殺城內的人，並且放火燒城。 \end{tabularx} \\ \\ \relax
1:9 & \begin{tabularx}{0.7\textwidth}{X} 後來猶大人下去，與住山區、尼革夫和低地的迦南人爭戰。 \end{tabularx} \\ \\ \relax
1:10 & \begin{tabularx}{0.7\textwidth}{X} 猶大去攻打住希伯崙的迦南人，殺了示篩、亞希幔、撻買。希伯崙從前名叫基列‧亞巴。 \end{tabularx} \\ \\ \relax
1:11 & \begin{tabularx}{0.7\textwidth}{X} 猶大從那裡去攻擊底壁的居民。底壁從前名叫基列‧西弗。 \end{tabularx} \\ \\ \relax
1:12 & \begin{tabularx}{0.7\textwidth}{X} 迦勒說：「誰能攻打基列‧西弗，奪取那城，我就把我女兒押撒嫁給他。」 \end{tabularx} \\ \\ \relax
1:13 & \begin{tabularx}{0.7\textwidth}{X} 迦勒的弟弟基納斯的兒子俄陀聶奪取了那城，迦勒就把女兒押撒嫁給他。 \end{tabularx} \\ \\ \relax
1:14 & \begin{tabularx}{0.7\textwidth}{X} 押撒來的時候，催促丈夫向她父親要一塊田。押撒一下驢，迦勒就對她說：「你要甚麼？」 \end{tabularx} \\ \\ \relax
1:15 & \begin{tabularx}{0.7\textwidth}{X} 她對迦勒說：「求你賜我福分；你既然把尼革夫給了我，求你也給我水泉。」迦勒就把上泉和下泉都賜給她。 \end{tabularx} \\ \\ \relax
1:16 & \begin{tabularx}{0.7\textwidth}{X} 摩西的岳父是基尼人，他的子孫與猶大人一起上到棕樹城，往亞拉得以南的猶大曠野去，住在百姓當中。 \end{tabularx} \\ \\ \relax
1:17 & \begin{tabularx}{0.7\textwidth}{X} 猶大和他哥哥西緬同去，擊殺了住洗法的迦南人，將城徹底毀滅。因此，那城名叫何珥瑪。 \end{tabularx} \\ \\ \relax
1:18 & \begin{tabularx}{0.7\textwidth}{X} 猶大攻取了迦薩和所屬的領土，亞實基倫和所屬的領土，以革倫和所屬的領土。 \end{tabularx} \\ \\ \relax
1:19 & \begin{tabularx}{0.7\textwidth}{X} 耶和華與猶大同在，猶大取得了山區，卻不能趕出平原的居民，因為他們有鐵的戰車。 \end{tabularx} \\ \\ \relax
1:20 & \begin{tabularx}{0.7\textwidth}{X} 以色列人照摩西所說的，把希伯崙給了迦勒。迦勒從那裡趕出亞衲的三支後裔。 \end{tabularx} \\ \\ \relax
1:21 & \begin{tabularx}{0.7\textwidth}{X} 至於住耶路撒冷的耶布斯人，便雅憫人沒有把他們趕出。於是，耶布斯人與便雅憫人同住在耶路撒冷，直到今日。 \end{tabularx} \\ \\ \relax
1:22 & \begin{tabularx}{0.7\textwidth}{X} 約瑟家也上到伯特利去，耶和華與他們同在。 \end{tabularx} \\ \\ \relax
1:23 & \begin{tabularx}{0.7\textwidth}{X} 約瑟家去窺探伯特利，那城起先名叫路斯。 \end{tabularx} \\ \\ \relax
1:24 & \begin{tabularx}{0.7\textwidth}{X} 探子看見一個人從城裡出來，就對他說：「請你把進城的路指示我們，我們會厚待你。」 \end{tabularx} \\ \\ \relax
1:25 & \begin{tabularx}{0.7\textwidth}{X} 那人把進城的路指示他們。他們就用刀擊殺了城中的居民，卻放走那人和他的全家。 \end{tabularx} \\ \\ \relax
1:26 & \begin{tabularx}{0.7\textwidth}{X} 那人往赫人之地去，建造了一座城，起名叫路斯。那城到如今還叫這名。 \end{tabularx} \\ \\ \relax
1:27 & \begin{tabularx}{0.7\textwidth}{X} 瑪拿西沒有趕出伯˙善和所屬鄉鎮的居民，他納和所屬鄉鎮的居民，多珥和所屬鄉鎮的居民，以伯蓮和所屬鄉鎮的居民，米吉多和所屬鄉鎮的居民；迦南人仍堅持住在這地。 \end{tabularx} \\ \\ \relax
1:28 & \begin{tabularx}{0.7\textwidth}{X} 以色列強盛的時候，就叫迦南人做苦工，沒有把他們全然趕走。 \end{tabularx} \\ \\ \relax
1:29 & \begin{tabularx}{0.7\textwidth}{X} 以法蓮沒有趕出住基色的迦南人。於是迦南人仍住在基色，在以法蓮中間。 \end{tabularx} \\ \\ \relax
1:30 & \begin{tabularx}{0.7\textwidth}{X} 西布倫沒有趕出基倫的居民和拿哈拉的居民。於是迦南人仍住在西布倫中間，成了服勞役的人。 \end{tabularx} \\ \\ \relax
1:31 & \begin{tabularx}{0.7\textwidth}{X} 亞設沒有趕出亞柯的居民和西頓的居民，以及亞黑拉、亞革悉、黑巴、亞弗革和利合的居民。 \end{tabularx} \\ \\ \relax
1:32 & \begin{tabularx}{0.7\textwidth}{X} 亞設人因為沒有趕出那地的居民迦南人，就住在他們中間。 \end{tabularx} \\ \\ \relax
1:33 & \begin{tabularx}{0.7\textwidth}{X} 拿弗他利沒有趕出伯˙示麥和伯˙亞納的居民。於是拿弗他利就住在那地的居民迦南人中，而伯˙示麥和伯˙亞納的居民卻成了為他們服勞役的人。 \end{tabularx} \\ \\ \relax
1:34 & \begin{tabularx}{0.7\textwidth}{X} 亞摩利人強逼但人住在山區，不准他們下到平原。 \end{tabularx} \\ \\ \relax
1:35 & \begin{tabularx}{0.7\textwidth}{X} 亞摩利人仍堅持住在希烈山、亞雅崙和沙賓；然而約瑟家權勢強盛的時候，他們成為服勞役的人。 \end{tabularx} \\ \\ \relax
1:36 & \begin{tabularx}{0.7\textwidth}{X} 亞摩利人的地界是從亞克拉濱斜坡，從西拉延伸而上。 \end{tabularx} \\ \\ \relax
2:1 & \begin{tabularx}{0.7\textwidth}{X} 耶和華的使者從吉甲上到波金，說：「我領你們從埃及上來，帶你們到我向你們列祖起誓應許之地。我曾說：『我永不廢棄我與你們的約。 \end{tabularx} \\ \\ \relax
2:2 & \begin{tabularx}{0.7\textwidth}{X} 你們不可與這地的居民立約，要拆毀他們的祭壇。』你們竟沒有聽從我的話。你們為何這樣做呢！ \end{tabularx} \\ \\ \relax
2:3 & \begin{tabularx}{0.7\textwidth}{X} 因此我說：『我必不將他們從你們面前趕出。他們必作你們肋下的荊棘，他們的神明必成為你們的圈套。』」 \end{tabularx} \\ \\ \relax
2:4 & \begin{tabularx}{0.7\textwidth}{X} 耶和華的使者向以色列眾人說這些話的時候，百姓放聲大哭。 \end{tabularx} \\ \\ \relax
2:5 & \begin{tabularx}{0.7\textwidth}{X} 於是他們給那地方起名叫波金，並在那裡向耶和華獻祭。 \end{tabularx} \\ \\ \relax
2:6 & \begin{tabularx}{0.7\textwidth}{X} 約書亞解散百姓，以色列人回到自己的地業，佔各自的地。 \end{tabularx} \\ \\ \relax
2:7 & \begin{tabularx}{0.7\textwidth}{X} 約書亞在世的日子和他死了以後，那些見過耶和華為以色列所做一切大事的長老還在世的時候，百姓都事奉耶和華。 \end{tabularx} \\ \\ \relax
2:8 & \begin{tabularx}{0.7\textwidth}{X} 耶和華的僕人，嫩的兒子約書亞死了，那時他一百一十歲。 \end{tabularx} \\ \\ \relax
2:9 & \begin{tabularx}{0.7\textwidth}{X} 以色列人把他葬在他自己地業的境內，以法蓮山區的亭拿‧希烈，在迦實山的北邊。 \end{tabularx} \\ \\ \relax
2:10 & \begin{tabularx}{0.7\textwidth}{X} 那世代的人也都歸到自己的列祖。後來興起的另一世代不認識耶和華，也不知道他為以色列所做的事。 \end{tabularx} \\ \\ \relax
2:11 & \begin{tabularx}{0.7\textwidth}{X} 以色列人行耶和華眼中看為惡的事，去事奉諸巴力。 \end{tabularx} \\ \\ \relax
2:12 & \begin{tabularx}{0.7\textwidth}{X} 他們離棄領他們出埃及地的耶和華－他們列祖的神，去隨從別神，就是四圍列國的神明，向它們叩拜，惹耶和華發怒。 \end{tabularx} \\ \\ \relax
2:13 & \begin{tabularx}{0.7\textwidth}{X} 他們離棄了耶和華，去事奉巴力和亞斯她錄。 \end{tabularx} \\ \\ \relax
2:14 & \begin{tabularx}{0.7\textwidth}{X} 耶和華的怒氣向以色列發作，把他們交在搶奪他們的人手中，又把他們交給四圍仇敵的手中，以致他們在仇敵面前再也不能站立得住。 \end{tabularx} \\ \\ \relax
2:15 & \begin{tabularx}{0.7\textwidth}{X} 他們無論往何處去，耶和華的手都以災禍攻擊他們，正如耶和華所說的，又如耶和華向他們所起的誓；他們就極其困苦。 \end{tabularx} \\ \\ \relax
2:16 & \begin{tabularx}{0.7\textwidth}{X} 耶和華興起士師，士師就拯救他們脫離搶奪他們之人的手。 \end{tabularx} \\ \\ \relax
2:17 & \begin{tabularx}{0.7\textwidth}{X} 然而，他們卻不聽從士師，竟隨從別神而行淫，向它們叩拜。他們列祖所行的道，所聽從耶和華的命令，他們都速速偏離了，並不照樣遵行。 \end{tabularx} \\ \\ \relax
2:18 & \begin{tabularx}{0.7\textwidth}{X} 耶和華為他們興起士師，耶和華與士師同在。士師在世的一切日子，耶和華拯救他們脫離仇敵的手。耶和華因他們受欺壓迫害所發出的哀聲，就憐憫他們。 \end{tabularx} \\ \\ \relax
2:19 & \begin{tabularx}{0.7\textwidth}{X} 但士師一死，他們又轉去行惡，比他們祖宗更壞，去隨從別神，事奉叩拜它們，總不放棄他們的惡習和頑梗的行為。 \end{tabularx} \\ \\ \relax
2:20 & \begin{tabularx}{0.7\textwidth}{X} 於是耶和華的怒氣向以色列發作，說：「因為這國違背我吩咐他們列祖當守的約，不聽從我的話， \end{tabularx} \\ \\ \relax
2:21 & \begin{tabularx}{0.7\textwidth}{X} 約書亞死的時候所剩下的各國，我必不再從他們面前趕出任何一個， \end{tabularx} \\ \\ \relax
2:22 & \begin{tabularx}{0.7\textwidth}{X} 為要藉此考驗以色列是否肯謹守遵行耶和華的道，像他們列祖一樣地謹守。」 \end{tabularx} \\ \\ \relax
2:23 & \begin{tabularx}{0.7\textwidth}{X} 耶和華留下那些國家，不將他們速速趕出，也不把他們交在約書亞的手中。 \end{tabularx} \\ \\
[1ex]
\hline
\hline
\end{longtable}
$^{1}$- Thank you, Patrick..
Good morning, UV..
(speaking in foreign language).
All right, so if you don't know this about me,.
besides the things that you see me do here.
on a Sunday morning, I also lead our youth group.
and our kids ministry throughout the weekdays..
And working with these students,.
one of the things that I realize is.
the more I spend time with them,.
the more I actually picked up some of the Gen Z.
and Gen Alpha lingo..
And sometimes I even find myself speaking like them..
(speaking in foreign language).
So let me give you a few examples here..
When they say, "Oh, I'm so cooked,".
it doesn't mean that they're literally getting cooked,.
but that means that they are in trouble.
or they're just exhausted..
Or when they say, "This thing or this person.
"looks kind of sus,".
it doesn't mean that, it just means that something.
or someone is looking a little suspicious..
And sometimes they will even combine.
a few of these slang together..
And for example, they will say, "This food is bomb, no cap.".
Let me translate it for you..
That means, "This food is so good, I'm not even lying.".
(speaking in foreign language).
So the reality is, we as human can be easily influenced.
by our surrounding and the people that we spend time with,.
whether it is for good or for worse..
And this is why there's a saying in Mandarin.
that goes like this,.
"He who stays near vermilion gets stained red,.
"and he who stays near ink gets stained black.".
(speaking in foreign language).
So this Sunday, we're actually going to begin a new series.
in the Book of Judges..
I don't know how familiar you are with this book,.

$^{41}$but maybe you have heard of the story of Samson,.
who is one of the more popular judges..
Or if you have been to church for a while,.
maybe you have heard of the story of Gideon and his army..
But there's actually a lot more to it.
in the Book of Judges..
(speaking in foreign language).
So what I hope to do with this sermon series.
is to help us situate these individual stories of judges.
within the context of the whole book,.
and also to understand the Book of Judges.
within the larger framework of God's redemptive history.
throughout the Bible..
(speaking in foreign language).
So what is the Book of Judges about?.
As one scholar describes it,.
at its core, the Book of Judges is about the canonization.
of the people of Israel..
Or to put it another way,.
it is about God's chosen people.
becoming more and more like the idolatrous nation.
around them..
And instead of fulfilling their calling.
as a kingdom of peace and a holy nation,.
and to show these nations around them what God is like,.
judges actually show us God's people at its worst,.
not just from a religious perspective,.
but also from a moral point of view..
(speaking in foreign language).
So if you turn each individual book in the Bible.
into a movie, I can bet you that Judges.
will be the most PG-13 book in the entire Bible,.
because in it we see things like idol worship,.
violence, murder, child sacrifice, and even gang rape..
However, in the backdrop of Israel's complete failure.
and their moral corruption,.
we also see a glimpse of hope for the future..
And this hope for the future is something.
that we will also be exploring.
throughout this series as well..

$^{81}$(speaking in foreign language).
So structurally, the book of Judges.
can be divided into three major section..
There is a double introduction,.
there is a double conclusion at the end,.
and in the middle of them is the cycle of Judges.
where the author dive into the narrative.
of these different judges of Israel..
And this morning we will focus.
on the first introduction section..
And as we go through our passage today,.
I want you all to keep this question in mind..
As God's people, are we becoming more and more.
like the world around us?.
(speaking in foreign language).
So let's look at Judges one together..
After the death of Joshua, the Israelites asked the Lord,.
who of us is to go up first to fight against the Canaanites?.
The Lord answered, Judah shall go up..
I have given the land into their hands..
(speaking in foreign language).
So these first two verses here give us a clue.
on where we are in terms of the history of Israel..
So let me give you a quick overview.
of Israel history up to this point..
The events in the book of Judges.
happened after the death of Joshua,.
but before Israel had their first king..
And at this point, Israel is actually halfway.
through taking over the land of Canaan..
(speaking in foreign language).
So if you recall the covenant that God made.
with your ancestor Abraham all the way back in Genesis 12,.
God promised Abraham that he will make him.
into a great nation, and all the people of the earth.
will be blessed through him..
Besides, God would also give the land of Canaan.
to Abraham's and his descendants..
(speaking in foreign language).
And when we come to the book of Deuteronomy,.

$^{121}$we know that Israel, they finally arrived.
at the edge of the promised land after they left Egypt.
and wander in the wilderness for 40 years..
And at this point, they were getting ready.
to take possession of the land..
So in Deuteronomy 20, we see that the Lord gave Israel.
this command specifically about taking over the land..
(speaking in foreign language).
So this is it..
Deuteronomy 20, verse 16..
However, in the cities of the nations,.
the Lord your God is giving you as inheritance,.
do not leave alive anything that breathes..
Completely destroy them, the Hittites, the Amorites,.
the Canaanites, the Perizzites, Hittites, and Jebusites,.
as the Lord your God has commanded you..
Otherwise, they will teach you to follow.
all the detestable things they do in worshiping their gods,.
and you will sin against the Lord your God..
(speaking in foreign language).
And so just a few things to highlight here..
When we read a passage like this.
to completely destroy any living things,.
it should disturb us..
And it is totally normal for us.
to wrestle with a passage like this..
But what we can know from the Bible.
is that these people are not good people..
These people are evil..
There's a lot of immorality in these cities,.
and they even practice child sacrifice.
in their idol worship..
(speaking in foreign language).
And this is why the Lord specifically warned Israel.
to completely destroy them, or in verse 18,.
they will teach you to follow all the detestable things.
they do in worshiping their gods,.
and you will sin against the Lord your God..
In other words, it is to preserve the holiness of Israel.
that God gave them this command..

$^{161}$(speaking in foreign language).
So that is Deuteronomy..
And when we come to the book of Joshua,.
Israel began the conquest of this land of Canaan,.
taking over the cities and settling into them..
And for the most part, the Joshua generations,.
they remain pretty faithful to Yahweh..
(speaking in foreign language).
And so this is in Joshua 24..
Israel served the Lord throughout the lifetime of Joshua.
and of the elders who outlived him.
and who had experienced everything.
the Lord had done for Israel..
(speaking in foreign language).
However, we also see some clues of potential challenge.
during Israel conquest of the land..
As we see that some of the tribes of Israel,.
they failed to completely drive out.
these inhabitants of the land..
(speaking in foreign language).
And when we come to the book of Judges here,.
this is where things begin to fall apart.
for the people of Israel..
(speaking in foreign language).
And so the rest of the book of Judges chapter one.
is sort of like a summary of how Israel fails.
to follow God's instruction fully..
Both for the Northern tribes and the Southern tribes..
(speaking in foreign language).
So in Judges one, verse 19,.
we see a summary of the conquest for the Southern tribe..
So verse 19, the Lord was with the men of Judah..
They took possessions of the hill country,.
but they were unable to drive out the people.
from the plains because they have chariots.
fed with iron..
And just a quick side note here,.
the iron chariot is actually not the real issue.
for Israel's failure here..
Because later in the story of Deborah and Barak.

$^{201}$in chapter four, we see that this iron chariot.
is nothing against the Lord..
And then in verse 21, the Benjamites, however,.
they did not drive out the Jebusites.
who were living in Jerusalem..
To this day, the Jebusites live there with the Benjamites..
(speaking in foreign language).
, and as for the Northern tribes,.
the failures are also being summarized here..
(speaking in foreign language).
So verse 27, but Manasseh did not drive out the people,.
nor did Ephraim drive out the Canaanites,.
neither did Zebulun drive out the Canaanites,.
nor did Asher drive out those living in.
and in a specific location,.
neither did Naphtali drive out..
And the Amorites confined the Danites to the hill country,.
not allowing them to come down into the plain..
(speaking in foreign language).
So if you're a visual learner, this next one is for you..
So the first image on the left is the extent of the land.
that is being allotted to all the tribes of Israel..
And so this lands and these cities were given.
to all 12 tribes of Israel..
However, they still have to do the work.
of actually taking over these cities..
(speaking in foreign language).
And the image on the right, on the other hand,.
shows the land and the cities that Israel.
had successfully taken over by the time of judges..
And if you just compare these two image,.
you can probably notice that there's actually.
still a lot of land that Israel didn't manage.
to take over, right?.
Especially those along the plains by the Mediterranean Sea..
(speaking in foreign language).
So at this point, you might wonder,.
why is this such a big deal that Israel.
didn't manage to completely take over the land?.
It seems like they had done a good job.

$^{241}$in taking over most of the land, right?.
(speaking in foreign language).
A lot of the time, we settle for half-hearted efforts..
As long as it is good enough, as long as I pass the exam,.
or as long as I do the bare minimum..
And working with my students in Summit,.
most of the time when I ask them to do something,.
they usually respond with some half-hearted effort..
(speaking in foreign language).
Let me give you an example..
Usually during Summit, they will pull up.
those ping pong table over there all the way back.
in the auditorium to play, right?.
And so after Summit, when I tell them.
to put the ping pong table and everything away,.
sometimes they might just fold up the table,.
but still leave the table all the way back there..
And sometimes they will put the table at the right place,.
but then they would just leave the paddle.
and the net still all the way back there..
Or sometimes they would also just leave.
the ping pong balls all around the auditorium.
and never pick them up..
(speaking in foreign language).
So I'm sure we can all think of some personal experience.
where either you, yourself, or your kids,.
or a coworker who just leave things halfway done, right?.
(speaking in foreign language).
So how often do we find ourselves doing just enough.
to get by instead of giving our best effort?.
When it comes to our relationship with God,.
how often do we just settle for a half-hearted obedience.
like the Israelites here?.
How often do we stop at good enough.
when God is calling us to something even greater?.
And how often do we desire to be a new creation in Christ,.
but we're just unwilling to do the hard work.
of rooting out sins and unhealthy patterns in our own lives?.
(speaking in foreign language).
So let's continue in Judges 2.

$^{281}$and see where this half-hearted obedience.
led the next generation of Israel after Joshua in two..
(speaking in foreign language).
So Judges 2, verse one..
The angels of the Lord went out from Gilgal to Bochim.
and said, "I brought you up out of Egypt.
and led you into the land I swore to give to your ancestor..
I say, I will never break my covenant with you,.
and you shall not make a covenant.
with the people of this land,.
but you shall break down their altars..
Yet you have disobeyed me..
Why have you done this?.
And I have also said, I will not drive them out before you..
They will become traps for you,.
and your gods will become snares to you.".
(speaking in foreign language).
And then verse four..
When the angel of the Lord has spoken these things.
to all the Israelites, the people wept aloud,.
and they called that place Bochim..
There they offered sacrifice to the Lord..
(speaking in foreign language).
So what was the charge that was being brought.
against the people of Israel here?.
If you read it carefully,.
the angel of the Lord did not accuse them.
because they failed to drive out the people.
in the land of Canaan..
But in verse two, we see that it is because Israel's.
end up making a covenant with these people,.
a covenant that is incompatible.
with the Lord's covenant with Israel..
And if you recall God's warning back in Jerome 20.
to completely destroy the inhabitants of the land,.
or they will lead Israel into worshiping idol,.
this warning actually becomes a reality here.
when we come to judges..
So look again in chapter two here..
(speaking in foreign language).

$^{321}$The Lord went out from Gilgal to Bochim and said,.
I brought you up out of Egypt and led you into the land.
I swore to give to your ancestor..
I say, I will never break my covenant with you..
And you shall not make a covenant.
with the people of this land,.
but you shall break their altars..
And yet you have disobeyed me..
(speaking in foreign language).
So in God's covenant with Israel,.
God actually held up his end of the bargain.
by bringing them into the promised land..
But for the people of Israel,.
they failed to hold on to their end of the bargain..
And what is the consequences here?.
In verse three, we see the Lord will not forgive them.
for their transgression..
And the Lord will not forgive them.
for their transgression..
And the Lord will not forgive them.
for their transgression..
And what is the consequences here?.
In verse three, we see the Lord will not help Israel.
in driving out the Canaanites now..
And these people will constantly be a thorn.
in Israel's side..
(speaking in foreign language).
So brothers and sister,.
do you see the dangers of our half-hearted obedience?.
When we allow sin to take roots in our lives.
and constantly resisting God,.
God will eventually give us into our own evil ways..
(speaking in foreign language).
And instead of becoming a light to the nations around them.
and showing them what God is like,.
Israel actually became more and more.
like these nations around them..
(speaking in foreign language).
And so this is what I want all of us.
to take away from today..

$^{361}$When we settle for half-hearted obedience,.
we give in to the ways of the world..
And again, when we settle for half-hearted obedience,.
we become more and more like the world around us.
instead of what God has called us to..
(speaking in foreign language).
So the questions that we need to ask ourselves is this,.
what areas of our life do we settle.
for half-hearted obedience?.
Is it our relationship with God, with others,.
with our money or material wealth,.
or with different addictions like alcohol, social media?.
And how have our half-hearted obedience lead us.
to become more and more like the world around us?.
(speaking in foreign language).
And so before I end,.
I want to leave us with this encouragement here..
This actually comes from Jesus' priestly prayer.
right before he was being arrested..
And he spent an entire chapter in John 17.
praying for his disciples and for his followers..
One of the things that Jesus specifically prayed about.
is this, "My prayer is not that you take them.
"out of the world, but that you protect them.
"from the evil one..
"They are not of the world, even as I am not of it.".
(speaking in foreign language).
So church, we are not of this world.
because we belong to Jesus..
And if you recall our previous series.
on Jesus' Sermon on the Mount,.
we know that the way of Jesus.
is actually very, very different.
from the way of the world..
(speaking in foreign language).
And so when the world is figuring out how to get rich,.
Jesus says, "Blessed are the poor in spirit.".
When the world is chasing after happiness and pleasure,.
Jesus says, "Blessed are those who mourn.".
When the world is getting after power and status,.

$^{401}$Jesus says, "Blessed are the meek.".
When the world says, "Eye for an eye and tooth for a tooth,".
Jesus says, "Turn the other cheek.".
And when the world says, "Love your own.
"and hate those that are not like you,".
Jesus says, "Love your enemies.
"and pray for those who persecute you.".
(speaking in foreign language).
And so when the world is figuring out how to get rich,.
Jesus says, "Blessed are the poor in spirit.".
When the world is getting after happiness and pleasure,.
Jesus says, "Blessed are those who mourn.".
When the world is getting after power and status,.
Jesus says, "Blessed are the meek.".
When the world says, "Eye for an eye and tooth for a tooth,".
Jesus says, "Turn the other cheek.".
And when the world says, "Love your own.
"and hate those who persecute you,".
Jesus says, "Blessed are the meek.".
When the world is figuring out how to get rich,.
Jesus says, "Blessed are the poor in spirit.".
When the world is getting after power and status,.
Jesus says, "Blessed are the meek.".
When the world is getting after power and status,.
Jesus says, "Blessed are the meek.".
When the world is getting after power and status,.
Jesus says, "Blessed are the meek.".
When the world is figuring out how to get rich,.
Jesus says, "Blessed are the poor in spirit.".
When the world is getting after power and status,.
Jesus says, "Blessed are the meek.".
When the world is figuring out how to get rich,.
Jesus says, "Blessed are the poor in spirit.".
When the world is figuring out how to get rich,.
Jesus says, "Blessed are the poor in spirit.".
When the world is figuring out how to get rich,.
Jesus says, "Blessed are the poor in spirit.".
When the world is figuring out how to get rich,.
Jesus says, "Blessed are the poor in spirit.".
When the world is figuring out how to get rich,.

$^{441}$Jesus says, "Blessed are the poor in spirit.".
When the world is figuring out how to get rich,.
Jesus says, "Blessed are the poor in spirit.".
When the world is figuring out how to get rich,.
Jesus says, "Blessed are the poor in spirit.".
When the world is figuring out how to get rich,.
Jesus says, "Blessed are the poor in spirit.".
When the world is figuring out how to get rich,.
Jesus says, "Blessed are the poor in spirit.".
When the world is figuring out how to get rich,.
Jesus says, "Blessed are the poor in spirit.".
When the world is figuring out how to get rich,.
\newpage



\section{以弗所書 2:1-10}
\label{sec:0vNwJDwewOw}
\textbf{From Death to Life (Eph 2\_1-10) - Paul Zhang 出死入生 (以弗所書2\_1-10) - 張威傳道}
\newline
\newline
連結: \href{https://youtube.com/watch?v=0vNwJDwewOw}{\texttt{https://youtube.com/watch?v=0vNwJDwewOw}} ~~~~ 語音日期: 2025-11-03
\newline
\newline
\hyperref[sec:Qv1tOecnO_s]{\small{< < < PREV SERMON < < <}}
~
\hyperref[sec:index_chronic]{\small{[返順時目]}}
~
\hyperref[sec:index_scriptual]{\small{[返順卷目]}}
~
\hyperref[sec:_S6mqFk4KdE]{\small{> > > NEXT SERMON > > >}}
\newline
\newline
以弗所書 2:1-10
\newline
\begin{longtable}{cl}
\hline
\hline
章節 & 經文 (和合本修訂版)\\
\hline
2:1 & \begin{tabularx}{0.7\textwidth}{X} 從前，你們因著自己的過犯罪惡而死了。 \end{tabularx} \\ \\ \relax
2:2 & \begin{tabularx}{0.7\textwidth}{X} 那時，你們在過犯罪惡中生活，隨從今世的風俗，順服空中掌權者的領袖，就是現今在悖逆的人心中運行的邪靈。 \end{tabularx} \\ \\ \relax
2:3 & \begin{tabularx}{0.7\textwidth}{X} 我們從前也都生活在他們當中，放縱肉體的私慾，隨著肉體和心中的意念去做，和別人一樣，生來就是該受懲罰的人。 \end{tabularx} \\ \\ \relax
2:4 & \begin{tabularx}{0.7\textwidth}{X} 然而，神有豐富的憐憫，因著他愛我們的大愛， \end{tabularx} \\ \\ \relax
2:5 & \begin{tabularx}{0.7\textwidth}{X} 竟在我們因過犯而死了的時候，使我們與基督一同活過來—可見你們得救是本乎恩— \end{tabularx} \\ \\ \relax
2:6 & \begin{tabularx}{0.7\textwidth}{X} 他又使我們在基督耶穌裡與他一同復活，一同坐在天上， \end{tabularx} \\ \\ \relax
2:7 & \begin{tabularx}{0.7\textwidth}{X} 為要把他極豐富的恩典，就是他在基督耶穌裡向我們所施的恩慈，顯明給後來的世代。 \end{tabularx} \\ \\ \relax
2:8 & \begin{tabularx}{0.7\textwidth}{X} 你們得救是本乎恩，也因著信；這並不是出於自己，而是神所賜的； \end{tabularx} \\ \\ \relax
2:9 & \begin{tabularx}{0.7\textwidth}{X} 也不是出於行為，免得有人自誇。 \end{tabularx} \\ \\ \relax
2:10 & \begin{tabularx}{0.7\textwidth}{X} 我們是他所造之物，在基督耶穌裡創造的，為要使我們行善，就是神早已預備好要我們做的。 \end{tabularx} \\ \\
[1ex]
\hline
\hline
\end{longtable}
$^{1}$[SPEAKING CHINESE].
Good morning..
[SPEAKING CHINESE].
And my wife has always reminded me.
after we got married I'd become a better man because of her..
[SPEAKING CHINESE].
And she has always reminded me how aggressive driver I was..
But after getting married, I get better because of her..
[SPEAKING CHINESE].
Before marrying her, I was very low in emotional intelligence..
But after getting married with her, I increased my--.
[SPEAKING CHINESE].
And she always remind me, do you agree with the change.
that I brought to your life?.
And I have to say yes because of what she has done to me..
[SPEAKING CHINESE].
And the same way, God has remind us, don't forget his grace..
[SPEAKING CHINESE].
And you can tell from the Old Testament,.
God has already reminded his people to build memorials,.
to remember his grace..
And then in the New Testament, always remind us.
what we have enjoyed in Christ Jesus..
[SPEAKING CHINESE].
And now last time, we're talking about Ephesians 1..
And we were reminded how God has chosen his people,.
adopted them as his children, and sealed them.
with the Holy Spirit..
And now we're in chapter 2 of the book of Ephesians..
[SPEAKING CHINESE].
This is our passage today, Ephesians 2, 1 to 10..
[SPEAKING CHINESE].
And this is the first three verses..
And the outline, it tells us our past..
We were dead in sin..
Verses 1 to 3, "As for you, you were.
dead in your transgressions and sins in which you used to live.
when you followed the ways of this world.
and the rulers of the kingdom of the air,.
the spirit who is now at work in those who are disobedient..

$^{41}$All of us also live among them at one time,.
gratifying the cravings of our sinful nature.
and follow his desires and thoughts..
Like the rest, we were by nature objects of wrath.".
[SPEAKING CHINESE].
And verses 4 to 6, we continue, "But because of his great love.
for us, God, in his rich in mercy,.
made us alive with Christ even when we were dead.
in transgression..
It is by grace that you have been saved..
God raised us up with Christ and seated us with him.
in the heavenly realms in Jesus Christ.".
[SPEAKING CHINESE].
Verses 7 to 10, "In order that in the coming ages.
he might show the incomparable riches of his grace.
expressed in the kindness to us in Christ Jesus..
For it is by grace that you have been saved through faith.".
And this is not from yourselves..
This is the gift of God, not by works,.
so that no one can boast..
"For we are God's workmanship created in Christ Jesus.
to do good works, which God prepared in advance for us.
to do.".
[SPEAKING CHINESE].
So we start by looking at this passage,.
looking at this picture..
What can you imagine?.
[SPEAKING CHINESE].
It's this death that we can imagine..
[SPEAKING CHINESE].
So Paul is using this kind of image, this kind of word.
to describe our past..
[SPEAKING CHINESE].
You see that Paul is not telling us that we were sick.
or we were weak, but we were dead in death..
[SPEAKING CHINESE].
So the state of being dead, it means we are separated from God..
[SPEAKING CHINESE].
So you go back to Genesis 2, when God tells the humanity,.
don't take the forbidden food..

$^{81}$Because if you do so, you will certainly die..
[SPEAKING CHINESE].
And of course, humanity took the forbidden food,.
and we bowed against God..
And that's why we have death..
[SPEAKING CHINESE].
So he talks about transgression in sin,.
and we were dead in sin..
It's using the standard of God to evaluate our condition..
[SPEAKING CHINESE].
So you can see men and women, we bow against God,.
and we do not want to respond to God..
And that's why we are left dead in sin..
[SPEAKING CHINESE].
So Paul is saying that we were dead in transgression,.
and we were totally dead because there was no way.
we can live again..
[SPEAKING CHINESE].
So you can imagine, you can see crops in front of you..
And if you yell at it, and you kick it,.
and there's no response..
[SPEAKING CHINESE].
So a spiritually dead person cannot respond to God.
and cannot earn his favor, cannot do anything.
to invoke change..
We were spiritually dead..
[SPEAKING CHINESE].
And you can see we are all dead in sin,.
and we cannot do anything..
And verses 2 and 3 describes our conditions..
"We follow the ways of the world and obey the ruler of the air.
and gratify our desires of the flesh..
And the sinners rule by three powers here--.
world, the devil, and the flesh.".
[SPEAKING CHINESE].
So the world is not talking about our cosmos,.
our physical universe..
It's talking about a system, a value that hates God..
[SPEAKING CHINESE].
And behind this world value system is the leader, Satan..

$^{121}$[SPEAKING CHINESE].
So the devil or Satan is not just--.
it's not what has been described in scary movies..
It talks about enticing us with his entire worldview..
[SPEAKING CHINESE].
And the flesh is not talking about our physical body..
It means our rebellious, sinful nature..
[SPEAKING CHINESE].
So those who are not born again, who are dead in sin,.
are enslaved by all these three powers..
[SPEAKING CHINESE].
So we cannot blame shift our responsibility to other things..
We have to see ourselves from the Word of God..
[SPEAKING CHINESE].
And Paul never blames sin entirely.
on Satan or other factors..
The real culprit is our human self..
[SPEAKING CHINESE].
So this one reminder here is very important,.
that we are dead in sin, and we cannot save ourselves..
Only by the grace of God, we can live again..
[SPEAKING CHINESE].
And the second reminder is that when we sin,.
we cannot blame our environment and blame to other factors,.
because all of us are sinful..
[SPEAKING CHINESE].
Well, as you can tell from James chapter 4, 1,.
it says we sin not because of other factors,.
not because of the devil..
We sin because our passions, our sinful nature..
[SPEAKING CHINESE].
And the world is telling us we sin because.
of the environmental factors..
And this is the belief of the world system,.
that we sin only because of our environment,.
our biological family, or other factors..
And one of the examples is in California some years ago..
Two brothers murdered their parents,.
and they were acquitted because they.
had been abused as children..

$^{161}$So this decided to try to de-sin them..
And that is the logic of the world system..
[SPEAKING CHINESE].
So you can see the logic of the world, the sinful world,.
is telling us robbery because we are poor,.
and we murder because of our trauma in the past,.
sin because of some chemical imbalance inside of us..
But all of this are lies..
[SPEAKING CHINESE].
Now you can see from the story of Joseph in the book.
of Genesis, he was betrayed by his brothers,.
and he should have become--.
Joseph should have become a very bitter and vengeful person..
But he forgave the brothers..
[SPEAKING CHINESE].
And some time ago, we heard the testimony of Mary Joan,.
and she was abandoned by her mother at the age of six..
And by human logic, she should be.
very resentful against her mother..
But she became an evangelist, shared the gospel,.
and cared about her mother..
And that is not environment..
That is the grace of God transforming her..
[SPEAKING CHINESE].
So you can see environment or social factors.
is never an excuse for sin..
It may expose our sin or entice us to sin,.
but it cannot cause us to sin..
God is calling us to respond to pain in our past.
with forgiveness and to face difficulties with faith..
[SPEAKING CHINESE].
And a side note is that many of the Chinese people.
love fortune-telling and feng shui.
because we try to find the factors for our being.
from the environment..
[SPEAKING CHINESE].
Yeah, about 1,500 years ago, the great theologian Augustine,.
before he came to Christ, he used.
to be worshiping the false gods and believing in myths..
But after he came to Christ, he changed..

$^{201}$[SPEAKING CHINESE].
And Augustine used to believe in some sort of astrology.
or cult, and teaching us--.
after he came to Christ, he said--.
he insists that human destiny is not.
determined by stars or astrological factors,.
but by God's will and human responses..
[SPEAKING CHINESE].
So when God says, who is responsible for our nature?.
Who is responsible for our actions?.
It's not the environment..
It's not the world, but us..
[SPEAKING CHINESE].
And the third reminder here is that how Satan works..
And Satan is trying to deceive us and use values that are.
against God to entice us into sin..
And Satan is trying to use some evil ideologies.
to entice us to sin against God..
[SPEAKING CHINESE].
And you can see a lot of modern ideologies--.
transgenderism or some other familiar liberal ideologies--.
behind all of this is Satan at work..
[SPEAKING CHINESE].
And that's why the Bible insists that we.
have to guard our hearts..
And if we guard our hearts, we guard the rest of our lives..
[SPEAKING CHINESE].
So this passage is telling us we are dead in sin,.
and what can we do about it?.
[SPEAKING CHINESE].
And Paul is trying to tell us there's a U-turn here,.
because, but there is something God has done to all of us..
[SPEAKING CHINESE].
And Paul is trying to tell us when we were dead in sin,.
but God is working in His grace to turn our lives around..
[SPEAKING CHINESE].
And we were dead in sin, and God is using His grace.
to turn us around and to give us new life..
[SPEAKING CHINESE].
See, the essence of being born again or new birth.

$^{241}$is that we can have a new beginning in the Christian.
life..
[SPEAKING CHINESE].
And many years ago, I was talking to someone in China,.
and we were reading the Gospel of John..
And this person said, being born again.
means only not cursing, not fighting, not smoking.
or drinking..
[SPEAKING CHINESE].
And this is not rebirth..
[SPEAKING CHINESE].
See, rebirth in the biblical sense.
is not a moral transformation..
It's not behavioral change..
It's not our Eastern religion tells us.
there is an awakening in our mind..
It is what God has done in our lives,.
entirely the will of God..
[SPEAKING CHINESE].
Just like 2 Corinthians chapter 4 says,.
for God, who said, let light shine out of darkness.
and make this light shine in our hearts..
So creation is creation out of nothing..
Just like rebirth in Christ brings us life out of death..
When God speaks, then the dead will live..
And faith and repentance are the fruit of the work of God..
[SPEAKING CHINESE].
So you think of the living again story of Lazarus..
And he had been dead for three days,.
and he was still in the tomb..
When Jesus came to the door of the tomb,.
he said, Lazarus, come out..
And Lazarus did not first decide to respond.
to the call of Christ..
And it was the word of Christ alone that gave him life,.
so that he could respond..
And that is the mystery of regeneration, new birth..
[SPEAKING CHINESE].
And Paul is trying to tell us God made us alive.
because of all his great love, his God's mercy,.

$^{281}$and his covenant love to all of us..
[SPEAKING CHINESE].
See, when you see love and mercy and compassion,.
he's talking about covenant promise God has for his people..
It's not about general human kind of love..
[SPEAKING CHINESE].
And of course, the Bible tells us.
God loves all the people in the very general sense,.
and sent rain to the just and unjust, the sinful.
and not sinful..
[SPEAKING CHINESE].
But the Bible insists the redemptive love and mercy.
are given to only those who belong to God's covenant..
And grace is offered to all of us,.
and effectually given to all the elect..
[SPEAKING CHINESE].
And thirdly, Paul says we are raised with Christ.
and seated with him in the heavenly realm..
[SPEAKING CHINESE].
And the verse, we are raised in Christ,.
and we are seated in the heavenly realm,.
these are all past tense verbs..
And it's already done..
That means Christ is already done for all of us..
[SPEAKING CHINESE].
And you think about the past tense, the [CHINESE].
how does it happen?.
How can people still alive on earth already in the past.
seated in heaven and raised with Christ?.
And that's the mystery of being in Christ..
[SPEAKING CHINESE].
So whenever the Bible tells us that we are in Christ,.
and we enjoy the benefit of being seated in the heavenly.
realm, and we are being saved, they're all in past tense..
[SPEAKING CHINESE].
And you can think about it..
Why is it in past tense?.
[SPEAKING CHINESE].
That means in some sense, in the past, we were there,.
and we were buried in Christ, we died in Christ,.

$^{321}$and we were resurrected in Christ..
[SPEAKING CHINESE].
Why?.
[SPEAKING CHINESE].
That's the mystery of being in union with Christ,.
and what Christ has done in the past..
His death and his resurrection is a historical event..
But because of the virtue of being united, union with Christ,.
we were, in a very real sense, were transferred from the old.
kingdom to the new kingdom..
What has happened in the past has already happened to all of us.
in Christ..
[SPEAKING CHINESE].
And even though we were not there,.
we did not exist back then, but because of the virtue of being.
in Christ, what happened to Christ happened to all of us..
[SPEAKING CHINESE].
So the Holy Spirit is applying what has already happened.
in Christ in the past to all of us..
[SPEAKING CHINESE].
That's the union in Christ..
[SPEAKING CHINESE].
And the Bible tells us, because of us being in Christ,.
God has already regenerated us and gave us new birth..
And we have to ask ourselves, have we.
experienced the new birth?.
[SPEAKING CHINESE].
And it's not because how many times we go to church..
It's not whether we pray the sinner's prayer or all.
the things that we have done..
It is whether we have really experienced.
the new birth in Christ..
[SPEAKING CHINESE].
And when God comes back again, God.
is looking at us, whether you have experienced.
the new birth in Christ..
[SPEAKING CHINESE].
And you see the truth here..
It's not because of all of us, we.
can do something to be born again..

$^{361}$It is Christ taking the initiative.
to grant us the new birth..
[SPEAKING CHINESE].
You can see the life in Christ..
We can see our human existence in a very broader perspective..
[SPEAKING CHINESE].
And life is very short..
And I think of my son..
He has already been one year old..
[SPEAKING CHINESE].
And very quickly, we may die..
And our life is faded..
And we will pass away..
[SPEAKING CHINESE].
But if you see your life in Christ's perspective,.
then our life is not just a few years here..
And our lives belong to eternity from creation to re-creation..
[SPEAKING CHINESE].
And Paul is not telling us, OK, what.
we can do as a Christian for Christ..
He's talking about salvation that we enjoy in Christ Jesus..
And from eternity to eternity, we enjoy this salvation..
[SPEAKING CHINESE].
And that is a very objective fact that.
has happened to all of us that in Christ,.
and Christ has done something for us..
[SPEAKING CHINESE].
Third point..
[SPEAKING CHINESE].
And the third point is giving glory to God from 7 to 10,.
verses 7 to 10..
Paul is trying to tell us our salvation is.
by grace through faith..
[SPEAKING CHINESE].
And you can see the picture here..
[SPEAKING CHINESE].
And what is he doing?.
[SPEAKING CHINESE].
He's praying..
[SPEAKING CHINESE].

$^{401}$And you can see he's trying to pray in humility..
[SPEAKING CHINESE].
And you can tell the empty bowl, if I have pitted on him,.
I put \$1 million in the bowl, but the bowl only.
contains \$1 million worth..
[SPEAKING CHINESE].
So the empty bowl does not have any worth in it,.
and only by the giver who gives whatever value into the bowl..
It's just like us responding to God with faith..
Our faith has no value..
And the value comes from what God has given to all of us..
[SPEAKING CHINESE].
So Paul continues to tell us this is not from ourselves..
This is the gift of God..
[SPEAKING CHINESE].
And you can tell it refers not only.
what is whatever given in the bowl,.
but the bowl itself is given by God..
[SPEAKING CHINESE].
So you can see salvation we have in Christ,.
it's not a trade-off..
It is because God offers grace, and it's not.
God offers something, that we offer faith..
Salvation is entirely, entirely the gift of God..
[SPEAKING CHINESE].
So the gospel, you can tell, it rejects.
any human earning or human favor..
[SPEAKING CHINESE].
And you can tell in our daily lives,.
we count what we have offered to transaction,.
what we have offered in some sort of exchange..
[SPEAKING CHINESE].
So our human's general response is.
that we want to offer something and boast about it,.
and because we do something for the transaction..
[SPEAKING CHINESE].
But the gospel silences all this kind of human boasting..
[SPEAKING CHINESE].
He said all the religions in the world are basically moralism..
Good people go to heaven, and bad people go to hell,.

$^{441}$and it's all about transactional exchange..
We do something, and then we earn something..
[SPEAKING CHINESE].
But you think about the gospel..
If we can offer something for God,.
then Christ Jesus does not have to come..
[SPEAKING CHINESE].
See, that's the uniqueness of Jesus Christ..
He came, and he died, and he was buried, and he rose again..
[SPEAKING CHINESE].
And you can tell that's the difference between Jesus.
and all these other religions, including Confucius..
The Bible tells us Jesus does all of them for all of us..
[SPEAKING CHINESE].
So many people would accept Jesus.
as a wise teacher and his good teachings,.
but they will reject the cross..
They will reject the biblical Jesus..
[SPEAKING CHINESE].
And even we, as Christians, we came to faith,.
and sometimes we work in that kind of mechanism..
We want to do something that we can earn the favor of God..
[SPEAKING CHINESE].
See, I love my son because he is my son,.
and there's kind of an exchange or transactional thing.
going on..
[SPEAKING CHINESE].
You see, I love my son because the virtue of the fact.
that he is my son, not because he does something good.
or he does something bad..
Because he is my son..
That's the identity..
That's why I love him..
[SPEAKING CHINESE].
And if he grow up and becomes bigger and older,.
and he, my son, wants to do something and earn my favor..
[SPEAKING CHINESE].
And then the son will tell me, no, I.
have done something for you, so many things for you,.
to treat you as my father..

$^{481}$Can you love me more?.
[SPEAKING CHINESE].
And then I would be so confused..
[SPEAKING CHINESE].
Because he does not understand my love for him..
It's not a condition on what he has done and what he is..
[SPEAKING CHINESE].
You see, if he wants to earn my favor by doing good things,.
he does not know the basis of my love..
[SPEAKING CHINESE].
See, in much the same way, our relationship with God.
is the same, because of his unconditional love for all of us..
[SPEAKING CHINESE].
See, God's gracious love for all of us.
is not because of what we have done or good things we have done.
for him to earn his favor..
It's all because we are all in Christ..
[SPEAKING CHINESE].
So the emphasis is that we are in Christ..
[SPEAKING CHINESE].
And then Paul continues to tell us.
the purpose of salvation is that we.
will be able to do spiritual good works,.
and the spiritual good works is not human moral practice..
[SPEAKING CHINESE].
And then it tells us we can show the grace.
to all the generations..
[SPEAKING CHINESE].
You see that because we are saved in Christ,.
so that our lives can be used by God.
to show forth his grace and his kindness in our lives..
[SPEAKING CHINESE].
So you can see the plan of salvation.
is not just to save us from sin or from death,.
but to save us so that we can glorify God..
[SPEAKING CHINESE].
And the quotes here are very meaningful,.
and I want to quote it..
"Jesus Christ died for sinners, but more for God himself..
And God put forth as appropriate agent by his blood.

$^{521}$to be received by faith, and this.
is the show of the glory of God.".
[SPEAKING CHINESE].
So you can see we are being saved in Christ.
so that our lives can show the righteousness of God..
[SPEAKING CHINESE].
So you can see in the entire biblical history, what.
God has done in Christ, what Jesus has done in history,.
whether we understand or comprehend it exhaustively,.
it is all for the purpose of God..
[SPEAKING CHINESE].
So the application point for us is that salvation is by grace,.
is not by our human work..
When we face failures or struggles or guilt,.
we need to return to the grace of God, not our self-effort..
[SPEAKING CHINESE].
And you can see whatever God has done for all of us.
is ultimately for his glory..
And we have to ask the question whether we.
live in such a way that would bring God glory..
[SPEAKING CHINESE].
Just like this verse says, whether you eat or drink,.
you live so that you can bring glory to God..
[SPEAKING CHINESE].
So whatever we have done or Christ has done in our lives,.
of eternal value..
[SPEAKING CHINESE].
So we can see there's no division whether someone.
is preaching so it's more valuable to God,.
or we eat or drink, or we conduct our daily life,.
and that is less value..
Everything, everything, whatever we do,.
whatever we think, is for the glory of God..
Everything is holy and sacred..
[SPEAKING CHINESE].
So there's no sacred or worldly division in Christ..
[SPEAKING CHINESE].
See, the emphasis is that whatever.
we do in Christ, whatever we do for Christ,.
it has eternal value..

$^{561}$[SPEAKING CHINESE].
So when God sees what you do, it's.
not because the intrinsic value of your doing..
It is because what you do in Christ..
[SPEAKING CHINESE].
So it's not talking about glorifying God,.
we have to do something grand or do something big..
It's talking about faithful daily living in Christ,.
for Christ..
[SPEAKING CHINESE].
And this is the end of message, and let's pray together..
[SPEAKING CHINESE].
Heavenly Father, we give you thanks.
because you have raised us from death to life.
and give us new heart and new life..
[SPEAKING CHINESE].
And if there is any one of us among us.
who has not been experienced the rebirth or born again,.
and God, we pray that you will bring us.
new life by your Spirit..
[SPEAKING CHINESE].
And for all of us who are already in Christ,.
and help us to know that our lives are not small or not.
valuable, but we are in Christ and we are not.
separated from you..
Whatever we do, do it for your glory..
[SPEAKING CHINESE].
In Jesus' name we pray, amen..
Thank you, Paul..
Thank you..
\newpage



\section{出埃及記 1:8-22}
\label{sec:_S6mqFk4KdE}
\textbf{上帝的應許與保守(出埃及記1\_8-22) - 江芷敏傳道}
\newline
\newline
連結: \href{https://youtube.com/watch?v=_S6mqFk4KdE}{\texttt{https://youtube.com/watch?v=\_S6mqFk4KdE}} ~~~~ 語音日期: 2025-11-03
\newline
\newline
\hyperref[sec:0vNwJDwewOw]{\small{< < < PREV SERMON < < <}}
~
\hyperref[sec:index_chronic]{\small{[返順時目]}}
~
\hyperref[sec:index_scriptual]{\small{[返順卷目]}}
~
\hyperref[sec:19Ss941ItL4]{\small{> > > NEXT SERMON > > >}}
\newline
\newline
出埃及記 1:8-22
\newline
\begin{longtable}{cl}
\hline
\hline
章節 & 經文 (和合本修訂版)\\
\hline
1:8 & \begin{tabularx}{0.7\textwidth}{X} 有一位不認識約瑟的新王興起，統治埃及。 \end{tabularx} \\ \\ \relax
1:9 & \begin{tabularx}{0.7\textwidth}{X} 他對自己的百姓說：「看哪，以色列人的百姓比我們還多，又比我們強盛。 \end{tabularx} \\ \\ \relax
1:10 & \begin{tabularx}{0.7\textwidth}{X} 來吧，讓我們機巧地待他們，恐怕他們增多起來，將來若有戰爭，他們就聯合我們的仇敵來攻擊我們，然後離開這地去了。」 \end{tabularx} \\ \\ \relax
1:11 & \begin{tabularx}{0.7\textwidth}{X} 於是埃及人派監工管轄他們，用勞役苦待他們。他們為法老建造儲貨城，就是比東和蘭塞。 \end{tabularx} \\ \\ \relax
1:12 & \begin{tabularx}{0.7\textwidth}{X} 可是越苦待他們，他們就越發增多，更加繁衍，埃及人就因以色列人愁煩。 \end{tabularx} \\ \\ \relax
1:13 & \begin{tabularx}{0.7\textwidth}{X} 埃及人嚴厲地強迫以色列人做工， \end{tabularx} \\ \\ \relax
1:14 & \begin{tabularx}{0.7\textwidth}{X} 使他們因苦工而生活痛苦；無論是和泥，是做磚，是做田間各樣的工，一切的工埃及人都嚴厲地對待他們。 \end{tabularx} \\ \\ \relax
1:15 & \begin{tabularx}{0.7\textwidth}{X} 埃及王又對希伯來的接生婆，一個名叫施弗拉，另一個名叫普阿的說： \end{tabularx} \\ \\ \relax
1:16 & \begin{tabularx}{0.7\textwidth}{X} 「你們為希伯來婦人接生，臨盆的時候要注意，若是男的，就把他殺了，若是女的，就讓她活。」 \end{tabularx} \\ \\ \relax
1:17 & \begin{tabularx}{0.7\textwidth}{X} 但是接生婆敬畏神，不照埃及王的吩咐去做，卻讓男孩活著。 \end{tabularx} \\ \\ \relax
1:18 & \begin{tabularx}{0.7\textwidth}{X} 埃及王召了接生婆來，對她們說：「你們為甚麼做這事，讓男孩活著呢？」 \end{tabularx} \\ \\ \relax
1:19 & \begin{tabularx}{0.7\textwidth}{X} 接生婆對法老說：「因為希伯來婦人與埃及婦人不同；希伯來婦人健壯，接生婆還沒有到，她們已經生產了。」 \end{tabularx} \\ \\ \relax
1:20 & \begin{tabularx}{0.7\textwidth}{X} 神恩待接生婆；以色列人增多起來，極其強盛。 \end{tabularx} \\ \\ \relax
1:21 & \begin{tabularx}{0.7\textwidth}{X} 接生婆因為敬畏神，神就叫她們成立家室。 \end{tabularx} \\ \\ \relax
1:22 & \begin{tabularx}{0.7\textwidth}{X} 法老吩咐他的眾百姓說：「把所生的每一個男孩都丟到尼羅河裡去，讓所有的女孩存活。」 \end{tabularx} \\ \\
[1ex]
\hline
\hline
\end{longtable}
$^{1}$大家好.
原來說到會拍手.
(笑).
好,今天大家多睡了一小時.
我的說話好像有很多話要說.
很長氣,希望大家可以專心.
好,開始之前我們一起祈禱.
差點誤會,謝謝你.
讓我們能夠閱讀你的話語.
是你的話語能夠改變我們生命.
你打開我們每一個人的心.
能夠明白你的旨意.
在這個地上實踐你的話語.
奉主明求,阿們.
今天我們會說出《埃及記》.
《埃及記》第一章.
這本書可能大家都很熟悉.
記載了在埃及.
上帝如何在奴役中拯救以色列人.
今天我們說第一章.
是一個很短,但很精彩的故事.
因為它說了雅各的家族.
如何在埃及紮根.
然後生養很多,遍佈全地.
然後受到埃及法老王的壓迫.
以至有一個接生婦出場.
她們的英勇行為拯救了很多嬰兒.
然後她們跟法老王有很智慧的對答.
今天說的接生婦.
是我們港道中的主角.
也可以說接生婆.
現代人稱為助產士.
希望解釋一下,讓大家知道.
別聽錯了其他東西.
希望我們在這篇港道之後.
可以說出接生婦.
究竟有甚麼生命特質.
神如何使用她.
為何她有這個生命特質.
讓我們一起看看今天港道的經文.

$^{41}$就是出埃及記的第一章8至14節.
如果你有聖經也可以拿出來.
我讀的是新譯本,第八節.
「那時有一位不認識」.
「約瑟的新王興起來,統治埃及」.
「他對自己的人民說」.
「看吶,以色列民比我們眾多強盛」.
「來吧,我們要用巧計對付他們」.
「恐怕他們爭多起來」.
「一旦發生戰爭」.
「他們就與我們的仇敵聯合」.
「攻擊我們,並且離開這地」.
第11節.
「於是他們只派督工管轄他們」.
「加重他們的重擔,苦害他們」.
「他們為法老建造兩座儲貨城」.
「就是比東和難失」.
第12節.
「但是埃及人越發苦害他們」.
「他們就越發爭多,越發繁殖」.
「埃及人就懼怕以色列人」.
「於是埃及人嚴嚴地使以色列人作苦工」.
「埃及人使他們因作苦工而覺得命苦」.
「他們要做和泥,做磚」.
「作田間各樣的工」.
「這一切苦工」.
「埃及人都嚴嚴地驅使以色列人去做」.
說到這裡.
鏡頭一開始就是描述一個新興的王.
他不認識約瑟.
可能是過了很多代人.
所以他也不認識誰是約瑟.
有文述記載.
我想是以色列人在埃及.
可能大概寄居了四百多年.
他們才出埃及.
所以這個新興的王沒有任何歷史背景.
也跟以色列人沒有任何關係.
他們開始壓迫以色列人,為甚麼.
第九節他們有說.

$^{81}$「埃及王懼怕以色列民」.
「人多勢眾」.
「以色列民在埃及眾多強盛」.
「當日約瑟帶一家人來到埃及」.
「整家人只有70人」.
你知道他們離開埃及時.
演變了多少人嗎.
(過百萬).
過百萬,對,應該有.
因為在《出埃及記》第12章.
他說以色列人離開埃及時.
大概有60萬名男子.
加上女人和小朋友.
當然不止,過百萬是對的.
還有第10節.
他說他怕以色列人多.
可以跟其他國家一起對付他們.
所以他很怕以色列人.
要攻擊埃及自己本土人.
所以他不讓他們繼續增長.
埃及王很嚴苛地奴役這些以色列人.
他要他們做甚麼.
要建一些儲貨城.
即是很大型的建築.
這些建築當然要很多人的物力.
所以他試圖陷害他們.
削弱他們的勢力.
埃及人派這些督工嚴格的.
去克制以色列人.
可能你也看過這些圖片.
以色列人要和泥造磚.
填間各樣的苦工.
聖經說埃及人要使以色列人做苦工.
覺得自己命苦.
所以這些大量的奴役.
讓他們身心靈覺得很疲乏.
越壓迫,情況就越反轉.
越多人生小孩,越來越多人.
形成了惡性循環.
越來越多人,我們更加苦害他們.

$^{121}$但又更加多生小孩.
這種奴役持續了四百多年.
四百多年.
是一個很殘酷的剝削.
連基本的權利也沒有.
以色列人雖然強盛.
但埃及在埃及沒有權勢.
他們唯有繼續接受這些迫害.
法老繼續鼓勵自己本土人.
掠奪和強暴對待這些以色列人.
而增加自己的聲望.
這種社會歧視和恐怖的氛圍.
令以色列人提心吊膽.
失去了自由和尊嚴.
試想想,這個時代變遷.
但這種恐怖的氛圍沒有停止.
在不同的國家或不同的人種.
還在經歷這些迫害.
不僅如此.
埃及法老還要趕盡殺絕.
命令殺死所有的紐邦男生.
希伯來的男嬰.
只留下那些嬰兒.
這種種族滅絕.
這樣嚴重威脅到以色列人.
這個奴役不僅是身體上的迫害.
而且是心理,種族層面上的迫害.
我們很難想像到.
在這個時代.
我們經常說生孩子要有很多預備.
要有穩定的經濟收入.
又要有一個安定的居所.
又要有一個好的教育環境.
社會要穩定.
我們才可以安心生孩子.
我應該…有數漏嗎?還有其他嗎.
我們有很多事情要想.
但你想想當時的以色列社會.
他們甚麼都沒有.
沒有任何的保障.

$^{161}$還要在位者濫殺這些男嬰.
怎樣才可以生孩子.
我們也看到其實生育率.
以及當時社會的變化.
或者一些特別的政治事件.
都會影響他們的出生率.
我看回香港的統計數字.
在香港的統計月刊中.
香港人的出生率….
我先看前一點.
就看到1981年.
他的出生率大概在1000人中.
有16.8個嬰兒.
我不知道為何有點數.
懂得統計的人可以解釋一下.
但到了2023年,香港發生甚麼事.
沙士,2003年….
他的出生率由16.8\%下降至7\%.
即1000人中有7個嬰兒出生.
再到社會運動之後.
就是另一個新低,2021年.
1000人中只有5個嬰兒出生.
所以香港生育現在有2萬元.
鼓勵生育.
所以你看到一個這麼不穩定的社會.
怎可以生育嬰兒.
但在十二節經文中說.
「埃及人越發苦害以色列人」.
「他們就越發爭多,越發繁殖」.
「埃及人就懼怕以色列人」.
這是很不合理的增長.
環境這麼惡劣.
種植物也知道環境這麼惡劣.
收成一定會差.
就像生育嬰兒一樣.
為何他們會越來越多人.
真的想不明白.
經文到第15節.
我們再一起看下去.
15節.

$^{201}$「埃及王吩咐兩個希伯來接生婦」.
「一個名叫斯芙拉·希芙拉」.
「另一個叫普瓦·普瓦」.
「他們說你們給希伯來婦人接生時」.
「要看著她們臨盤」.
「若是男孩,你們要殺死她」.
「若是女孩,她就可以活著」.
「但是接生婦卻敬畏神」.
「不照著埃及王吩咐她們的去做」.
「竟讓男孩活著」.
「於是埃及王把接生婦叫了來」.
「對她們說」.
「你們為甚麼作這事」.
「讓男孩活著」.
「接生婦回答法老」.
「希伯來婦人與埃及婦人不同」.
「她們肯有活力」.
「接生婦還沒有到」.
「她們已經生產了」.
「神恩代這一班接生婦」.
「以色列民增多起來」.
「而且非常強盛」.
「接生婦因為敬畏神」.
「神就為她們建立家室」.
「法老吩咐她們」.
「她還文說」.
「凡是希伯來所生的男孩」.
「你們都要把她投在河裡」.
「凡是女孩就讓她活著」.
你看到希伯來和普亞.
這兩個接生婦有甚麼特質.
敬畏神,經文說了兩次.
她們敬畏神.
不照著法老的吩咐去殺男嬰.
堅持敬畏神.
值得注意的是.
你可以看到這兩個接生婦.
叫甚麼名字.
是嗎.
她們有名字.

$^{241}$埃及王有名字嗎.
沒有,只是稱她為埃及王.
所以在這個鮮明的對比之下.
你知道作者有話說.
有名字的有意思.
所以我們看聖經時.
要留意有名字的人.
雖然她們只是一個很短的篇幅.
但聖經作者給了她們名字.
是有意思的.
而且在希伯來文化中.
通常生育,分娩.
通常由女性做接生婦.
這些接生婦不僅要處理生產過程.
而且還要提供一些情感的支援.
就像我們現在培育.
不僅要煮三餐.
還要支援剛生完寶寶的媽媽.
男士的確很難做到.
即使是自己的丈夫.
我想神給我們男女的創造.
真的很不一樣.
男人很難有情感支援.
我想起很搞笑.
我跟丈夫的一個對話.
有一天我照鏡子.
我說「好像我最近臉色好了點」.
「有點紅潤」.
他看著我,很認真地說.
「你是否包賊」.
我覺得男女想法很不一樣.
而且希伯來接生婦.
他們還有甚麼要做.
就是要加入這些宗教元素.
例如祈禱,做一些儀式.
所以這些接生婦在希伯來社會中.
扮演了一個很重要的角色.
雖然聖經只提到兩個人.
就是斯芙拉和普亞.
但你想想,超過100萬人口.

$^{281}$會否只有兩個接生婦.
很可能他們是接生婦團隊的代表.
或者是他們的領袖.
所以法老只接見他們兩個.
還有第三點.
這些希伯來的接生婦.
你猜他們是否真的希伯來人.
聖經這樣形容他們是希伯來接生婦.
但其實很多人對他們的民族身分.
是有爭議的.
究竟他們是希伯來人.
為自己民族的婦女接生.
還是他們是為希伯來人接生的接生婦.
兩個不同的,你聽明白了吧.
究竟他們是希伯來人.
還是他們不是希伯來人.
有些人說可能他們不是希伯來人.
因為法老一接見他們.
就能接見這兩個人.
還有希伯來在埃及文字中解作美麗.
很可能這是埃及人的名字.
還有猶太學家約瑟夫在猶太古史中說.
這兩個是埃及人.
不過有些猶太人否定.
他們說哪有可能由外邦人.
為我們希伯來人接生.
剛才說接生婦要加入一些宗教元素.
愛俗的人怎可以這樣.
我想是交由大家自己去判斷.
不過無論他是埃及人還是以色列人.
我想他們當時也是一群.
很斗膽對抗當時政權的兩位勇者.
聖經形容他們是敬畏神,對吧.
敬畏神,這是一個很關鍵的動機.
在古埃及,法老自古以來.
都被視為地上的神明.
接生婦看到地上神明的代表,法老.
他主宰了所有東西.
而且他一句話就可以取他的命.
為何接生婦不害怕.

$^{321}$我猜她不是不害怕.
但反而她更害怕.
更敬畏我們這位看不見的神.
是一位看不見的神.
這不是敬畏懲罰而做出他們的行為.
而是我想接生婦優先考慮到神的旨意.
去敬畏神.
他們明白違抗法老,埃及王.
很可能有即時的危險.
不單是他自己,可能是他全家人.
但他們也要選擇去敬畏神.
因此這是我們今天很想知道.
究竟他有甚麼特質.
他就是敬畏神.
他更害怕這個看不見的神.
其實經文中也提及法老的害怕.
法老害怕甚麼.
害怕他的地位受到威脅.
害怕這些以色列人越來越多.
他的政權受到威脅.
形成了一個對比.
接生婦也害怕.
但她害怕這個看不見的神.
但法老害怕他看見面前的困難.
一個阻礙他的困難.
一群越來越多的以色列人.
我們看經文時.
看到法老質問兩個接生婦時.
他們很聰明,他們的回答是怎樣的.
他說「不是我沒有殺這群男嬰」.
「而是這群稀白來的婦人很強健」.
「我去到時,他們已經生了兒子」.
你可能會說這是否撒謊.
我們基督徒不應該撒謊,對吧.
你可以討論一下.
這是否道德,是否正確.
但很明顯這不是作者的向導.
他的動機是想告訴我們.
這群接生婦因為敬畏神.
無論她要做甚麼.

$^{361}$她都要保護這群無辜的生命.
為了神去對抗這個政權.
歷史歷代都有很多對抗政權.
仍然堅持走出上帝真理的一群人.
可能你們也有留意過.
有一個德國的老德會牧者.
他叫潘霍華.
他在納粹時期做牧者,是一位神學家.
他很勇敢對抗希特拉的政權.
他反對納粹軍.
他參加了反納粹的認順教會.
以及拒絕納粹軍對教會的控制.
很積極地參加了地下反抗納粹的運動.
潘霍華很堅持上帝的教導.
很強調基督徒應該信服上帝.
而不是在位的掌權者.
尤其是掌權者違背上帝公義的時候.
所以他在監獄中寫了一本書.
名為「追隨基督」.
他強調廉價恩典和昂貴恩典的區別.
他呼籲這些基督徒在壁壩裡敬畏神.
以及要活出信仰.
他的結局是在1945年.
被納粹軍處決,成為了殉道者.
另一個較接近我們的人.
他叫Charlie Kirk.
他很堅持上帝的道.
他不畏懼一些人對他的挑戰.
他是一個保守黨.
他在生前很多次公開自己基督徒的見證.
很強調他要堅持敬畏神.
並以信仰作為優先.
在面對很多政治壓力的時候.
他堅持以捍衛信仰為先.
他說千萬不要讓恐懼主導了你的行為.
而是要信靠神.
在一次訪問裡.
他說「我希望被記住的是」.
「我對信仰的勇氣」.
「最重要的是我生命裡的信仰」.

$^{401}$Charlie Kirk今年在猶他州被槍殺死.
他死的時候只有31歲,很年輕.
我再看這段經文時.
我心想這兩個接生父其實都很年輕.
因為上帝知道他敬畏神之後.
他說要為他們建立家室.
當時希伯來人通常是12,13歲結婚.
我說「那他們豈不是很年輕」.
但他們又那麼有知識.
其實我想有時神預備我們.
去某個崗位事奉祂.
不一定是因為我們年紀不大.
或是我們是否成熟.
神揀選我們是看我們的心.
無論你的年紀有多大,有多小.
神都很看重我們敬畏祂的心.
相反,即使年紀有多大.
也不限制我們在甚麼崗位事奉神.
昨天我們舉辦了一個歡迎派對.
歡迎上次報道會缺席的朋友.
我們一起吃東西,玩遊戲,聊天.
亦分享了順主的經歷.
我聽到他說順主之後.
他很清楚神要他在這生中.
堅持事奉祂.
而且他在工作中事奉一群長者.
他可以說已經到了退休的年齡.
但他知道神預備他的生命.
要在現在這個崗位繼續服侍神.
我們很欣賞燕姐.
不僅在工作中很忠心事奉神.
在教會中也很忠心事奉.
但我知道很多成人片人.
都很忠心事奉神.
特別是我們教育下一代.
作為父母的大家,其實很難做到.
因為我們面前的困難很多.
就像埃及王面對前面的問題很多.
但我們倚靠那位神,卻好像看不見.
有時我們的心也會怯一怯.

$^{441}$但我想這個接生父給我們的分享.
或提醒就是.
即使我們看不到上帝.
但我們仍然信靠祂,敬畏祂.
祂會帶領我們經過這一切的路.
一開始時,我說今天的講道很簡單.
就是要知道一件事.
接生父究竟有甚麼生命特質.
就是敬畏神.
我再問為何他要敬畏神.
為何他要敬畏神.
我們再看《尚文夏里》.
《出埃及記》一章還有1至7節.
我們剛才沒有讀.
時間的關係,你可以看看.
這是一個雅各的家譜.
第七節總結.
「以色列人生養繁衍眾多」.
「人數增加,極其強盛」.
「遍滿了那地」.
其實這是創世記上帝給我們的應許.
「生養眾多,遍滿全地」.
這亦是我們每一個的使命.
創世記上帝對阿巴拉罕說.
「我必使你的後裔極其繁多」.
他亦對雅各應許.
「你的後裔必如海邊的沙」.
到了《出埃及記》.
他繼續重申一次自己的應許.
就是「比以色列人生養眾多」.
「人數增加,極其強盛」.
「遍滿了那地」.
無論在以色列還是出埃及.
上帝都在堅持他的應許.
這是上帝給我們應許的證明.
在困苦中,我們有時可能會迷失.
但我們看看.
原來上帝給我們的應許和使命.
是祂會幫助我們.
而且祂會在地上告訴我們祂在.

$^{481}$這也是接生父很堅持的原因.
儘管他面前的法老質問他的時候.
他也可能會拚命跟法老交戰.
他不怕法老,因為他敬畏神.
我想我們的生活.
面前有很多困難,也很不容易.
不知道你們有否問過這個問題.
為何我們那麼艱難.
還要繼續向前走.
我想是基於我們基督徒.
知道上帝給我們的使命是甚麼.
上帝給我們的使命是甚麼.
如果你到今天還不知道.
上帝給你的使命是甚麼.
我們要好好尋求,要好好讀聖經.
上兩星期,我看「一六敬拜」的時候.
唱了這首詩歌,叫「手中的福音」.
其實這個「手中的福音」.
好像總結了我們基督徒.
或是在這個世代裡應該要做的使命.
我想跟大家再唱一次這首歌的副歌.
或者你可以再留意一下歌詞.
對,時間關係,我們只唱副歌.
懂得唱的,我們可以一起唱.
我們只唱副歌.
懂得唱的,我們可以一起唱.
(手中的福音).
(憐憫的體恤各方的需要).
(如同住看故).
(關注還難圍固).
(垂簾無故 為接受損者痛苦).
(無謂的擔當與這的劍部).
(無負一點苦痛).
(不理會偽像高 仍然祈禱).
(要將福音於手裡).
(宣告).
好像不太夠口氣.
回家後可以繼續唱.
我想大家在這個世代裡.
做一個基督徒的確不容易.

$^{521}$但我鼓勵大家去尋找.
上帝給你們在生命裡的使命.
無論在你的工作崗位裡.
你照顧小朋友的心路歷程.
真的很不容易.
但希望你看到.
我們看不到的神.
是超過和大於我們所看到的問題.
我們的神是大於我們所看到的問題.
願大家在經文裡.
可以看到上帝給我們的鼓勵.
求主幫助我們活出祂的真理.
願主祝福大家.
\newpage



\section{路加福音 8:26-39}
\label{sec:19Ss941ItL4}
\textbf{誰是敵人? (路加福音8\_26-39) - 梁紫健傳道}
\newline
\newline
連結: \href{https://youtube.com/watch?v=19Ss941ItL4}{\texttt{https://youtube.com/watch?v=19Ss941ItL4}} ~~~~ 語音日期: 2025-11-03
\newline
\newline
\hyperref[sec:_S6mqFk4KdE]{\small{< < < PREV SERMON < < <}}
~
\hyperref[sec:index_chronic]{\small{[返順時目]}}
~
\hyperref[sec:index_scriptual]{\small{[返順卷目]}}
~
\hyperref[sec:SzBowSpLGuQ]{\small{> > > NEXT SERMON > > >}}
\newline
\newline
路加福音 8:26-39
\newline
\begin{longtable}{cl}
\hline
\hline
章節 & 經文 (和合本修訂版)\\
\hline
8:26 & \begin{tabularx}{0.7\textwidth}{X} 他們到了格拉森人的地區，就在加利利的對面。 \end{tabularx} \\ \\ \relax
8:27 & \begin{tabularx}{0.7\textwidth}{X} 耶穌上了岸，就有城裡一個被鬼附的人迎著他走來。這個人好久不穿衣服，不住在屋子裡，而住在墳墓裡。 \end{tabularx} \\ \\ \relax
8:28 & \begin{tabularx}{0.7\textwidth}{X} 他看見耶穌，就喊叫著俯伏在他面前，大聲說：「至高神的兒子耶穌，你為甚麼干擾我？我求你，不要叫我受苦！」 \end{tabularx} \\ \\ \relax
8:29 & \begin{tabularx}{0.7\textwidth}{X} 這是因耶穌曾吩咐污靈從這人身上出來。原來這污靈屢次抓住他；他常被人看守，又被鐵鏈和腳鐐捆鎖，他竟把鎖鏈掙斷，被鬼趕到曠野去。 \end{tabularx} \\ \\ \relax
8:30 & \begin{tabularx}{0.7\textwidth}{X} 耶穌問他：「你的名字叫甚麼？」他說：「群」；這是因為附著他的鬼多。 \end{tabularx} \\ \\ \relax
8:31 & \begin{tabularx}{0.7\textwidth}{X} 鬼就央求耶穌不要命令他們到無底坑裡去。 \end{tabularx} \\ \\ \relax
8:32 & \begin{tabularx}{0.7\textwidth}{X} 那裡有一大群豬正在山坡上吃食，鬼央求耶穌准他們進入豬裡；耶穌准了他們。 \end{tabularx} \\ \\ \relax
8:33 & \begin{tabularx}{0.7\textwidth}{X} 於是鬼從那人出來，進入豬裡，那群豬就闖下山崖，投進湖裡，淹死了。 \end{tabularx} \\ \\ \relax
8:34 & \begin{tabularx}{0.7\textwidth}{X} 放豬的看見這事就逃跑了，去告訴城裡和鄉下的人。 \end{tabularx} \\ \\ \relax
8:35 & \begin{tabularx}{0.7\textwidth}{X} 眾人出來，要看發生了甚麼事；到了耶穌那裡，發現那人坐在耶穌腳前，鬼已離開了他，穿著衣服，神智清醒，他們就害怕。 \end{tabularx} \\ \\ \relax
8:36 & \begin{tabularx}{0.7\textwidth}{X} 看見這事的人把被鬼附的人怎麼得醫治的事告訴他們。 \end{tabularx} \\ \\ \relax
8:37 & \begin{tabularx}{0.7\textwidth}{X} 格拉森周圍地區的人，因為害怕得很，都求耶穌離開他們；耶穌就上船回去了。 \end{tabularx} \\ \\ \relax
8:38 & \begin{tabularx}{0.7\textwidth}{X} 鬼已從身上出去的那人懇求要和耶穌在一起，耶穌卻打發他回去，說： \end{tabularx} \\ \\ \relax
8:39 & \begin{tabularx}{0.7\textwidth}{X} 「你回家去，傳講神為你做了多麼大的事。」他就走遍全城，傳揚耶穌為他做了多麼大的事。 \end{tabularx} \\ \\
[1ex]
\hline
\hline
\end{longtable}
$^{1}$(第五節).
丁姐妹,早安.
先按一下「不理我」.
不知道大家有否玩過狼人殺.
或者不同的劇本殺.
狼人殺是一個.
很二元對立的角色扮演遊戲.
狼人要負責殺死所有神明.
以及所有人類.
而神明和人類負責對付狼人.
通常…99.9\%都是一方勝出.
我們玩了那麼多次.
只試過一次是打和.
這是我們UV在疫情期間創作的.
聖經版狼人殺.
有撒旦,魔鬼仔,人類.
有亞當,夏娃.
當然有仙姿,二才藝.
也有參孫,11A參孫.
當然還有耶穌.
如果你有玩過.
不是聖經版本的狼人殺.
當然是狼人,神明和人類的對決.
很燒腦,不斷猜誰是狼人.
不斷猜誰是神明.
人類的廢村通常沒有甚麼參與.
只是被人殺.
而聖經版本的狼人殺很有趣.
大家有留意到有個比較凶狠的朋友嗎.
在我的左下角,他是參孫.
連這個很頑皮的參孫.
也是一個光明使者.
但如果看聖經.
其實他也做了很多壞事.
我還會問為何參孫也是光明使者.
別問我,不是我做的.
你可以問作者.
而人類唯一的敵人在這個版本裡.
是魔鬼.
人類唯一的敵人是魔鬼.

$^{41}$這是一個很有趣的轉變.
大家不覺得奇怪嗎.
即使是壞人做了好事或壞事.
他在聖經的詮釋裡.
絕對不是人類的敵人.
所以今天我們要探討一個主題.
究竟誰是你的敵人.
在你的生命裡.
有些人是你非常憎恨的嗎.
甚至你心裡有一種想法.
「我永遠不想再見到這個人」.
「我不希望我人生中再遇上這樣的人」.
聖經在這個課題上回答了我們.
給了我們答案.
我們一起來看看.
聽完這篇講道,我們會知道三件事.
第一,我們會知道盧家福音.
「量神蹟,趕鬼入豬群」.
和「平靜風和浪」.
這兩個故事究竟代表甚麼.
第二,我們會用甚麼心態.
對付我們的敵人.
第三,我們的敵人包括我們自己嗎.
如果大家準備好了.
我們就讀讀聖經吧.
首先我們回顧.
盧家福音八章裡的兩個故事.
因為它們是連續的故事.
還要互相呼應.
耶穌和門徒來到格拉森的地區.
遇到一個被鬼附身的人.
這個人住在洞穴裡,很恐怖.
沒有人可以制止他.
這個被鬼附身的人兄.
一眼就認出耶穌是至高神的兒子.
他說「我與你何干,你不要騷擾我」.
「我不想受苦」.
然後這個人兄被鬼趕到曠野.
留意,曠野是一個很值得留意的字眼.
我稍後會解釋.

$^{81}$然後耶穌就問鬼叫甚麼名字.
鬼就說「我是叫群」.
原文是Legion.
我稍後會解釋這個字.
因為他附身的鬼很多.
所以鬼求耶穌.
不要趕他們進去無底洞.
然後就跳進去.
希望可以進入豬群裡.
耶穌大名後,鬼就上了豬的身.
豬立刻衝進山崖裡.
在湖裡淹死.
這個「湖」也是一個很特別的字眼.
當地人看完後.
心裡明白了,他們就害怕.
我看到這裡就問.
「究竟害怕甚麼」.
之後我會說說他們明白了甚麼.
又說說他們害怕了甚麼.
然後他們就叫耶穌離開了他們的地方.
在被鬼上身的人兄之後.
他康復了,坐在耶穌的腳前.
然後他就去傳揚.
上帝耶穌為他做的大事.
我們要問的問題是為甚麼.
當地人明白了之後.
他們會害怕,對吧.
他們害怕甚麼.
跟湖有甚麼關係.
跟群有甚麼關係.
還有跟曠野有甚麼關係.
我們要回答這個問題.
我們要了解一下這個段落的上文.
就是….
《老鬼方案》八章上文說了一個事蹟.
就是在一個講鬼故事之前.
有一個很經典的神蹟.
就是耶穌跟門徒上船.
要去到湖的另一邊.
途中就起了大浪.

$^{121}$整艘船進水,門徒就很害怕.
就叫醒在睡覺的耶穌.
「耶穌,起來…」.
耶穌就說,就馬上風和浪.
風浪就馬上平靜了.
那些門徒是怎樣?很害怕.
問究竟這個是誰.
連風和水都聽他說話.
有否留意.
門徒跟當地人的反應.
其實是有點類似的.
都是很害怕,驚訝,驚奇.
這兩個故事在敘事上.
是一個很緊密的連結.
平靜風浪是發生在.
耶穌帶門徒渡過湖的過程.
在湖上,對吧.
而講鬼故事是發生在哪裡.
就是剛剛渡過湖之後.
剛剛下船.
所以這種敘事,暗事.
其實是說兩個故事.
是一個有關連的事件.
它不是兩個分別的神蹟.
是同一件事.
我又要說說剛才說「群」的字.
在故事中,鬼說「我們是群」.
這個名字其實不是亂取的.
不是因為它們多,所以叫「群」.
是因為這個故事.
有一個很獨特的歷史和文化象徵.
無論是希羅時代的猶太哲學家.
約瑟夫.
還是希臘哲學家,歷史家.
玻利比烏斯.
都是用這個字,Legion.
去代表軍團.
除了軍團之外.
就是一群很多….
剛才說的那些人.

$^{161}$就是一群軍人的表達.
所以如果在羅馬軍團的統稱中.
它在表達甚麼.
大概有5000至6000個士兵.
在羅馬的世界中.
它代表著一個很強大的軍事力量.
當時猶太人,包括耶穌所在的.
加利利地區.
都是被羅馬佔領著.
所以對猶太人來說.
羅馬軍團是壓迫和暴力的象徵.
所以「鬼叫群」這個字.
不只是說它有很多很多的意思.
它代表著它像羅馬軍團般強大.
如此組織,如此壓迫.
它們就傷了人.
令人很痛苦和混亂.
而耶穌說「鬼」,很有趣.
祂不是跟人類軍隊打.
而是直接跟這些靈界軍團對抗.
不知道大家有否看過《聊齋》.
有看過TVB版本的《聊齋》.
看看可否舉手.
如果大家看過這個版本.
可能大家也有一點社會經驗.
1996年首播,我看重播.
裡面有很多妖怪,神怪的東西.
很天馬行空.
很多人都會覺得很驚訝.
為何事情會這樣發生.
另一齣…這是我的年代.
天馬行空的電影是西遊記.
也是TVB的,對吧.
無論是張惠健版本.
或陳虹文版本,我都很喜歡.
大家都知道《聊齋》和西遊記.
其實是神話小說.
沒有人會當它是真的,對吧.
是神話吧.
它是…除了TVB用來做劇集外.

$^{201}$其實它是一個明清時代的神話小說.
但它又不是那麼單純.
它不是單純地說.
創作一些東西讓大家笑一下.
它是在表達當時對於政權的不滿.
但它又不能直接地跟別人說.
「我真的對這個政權很不滿」.
它就要用很神化的東西.
去表達給別人聽.
「我們真的對這個政府很不滿」.
所以那些妖怪.
其實是在代表社會,政權.
對於人的不公.
所以他們要打擊那些妖怪.
要趕走他們.
就是因為這個原因.
而盧家福音這裡的神蹟是真的.
但他仍然用了這個技巧.
這個影視技巧.
就是當時的猶太人覺得.
這個羅馬政權那麼暴力,那麼殘忍.
經常殺他們的同胞.
令他們苦不堪言.
他們才是真正的敵人.
所以Legion這個字.
羅馬軍團就像冤鬼一樣.
控制著羅馬人,控制著猶太人.
就像當年埃及人處處刁難,壓迫猶太人.
令猶太人的生活苦不堪言.
是有這個意味.
但耶穌很有趣.
耶穌很巧妙地.
回憶猶太人很熟悉的舊恩歷史.
扭轉當時猶太人的想法.
就是真正的敵人.
不是眼見,看到的敵人,羅馬人.
而是躲在背後的邪惡勢力,魔鬼.
如果我們把平靜風和浪.
以及當日諸群的這兩個神蹟.
跟舊約《出一甲記》14章做比較.

$^{241}$我們會更加知道作者的用意.
大家如果有讀過《出一甲記》.
通常我們未必能出到埃及.
你知道很長,對吧.
讀到14章已經算不錯.
這也是一個很經典的神蹟.
是上帝帶領摩西,帶領以色列人出埃及.
逃避了法老的追趕.
而這在猶太傳統經典中.
是一個解放的故事.
而路加福音很有趣.
他故意重述舊約的事件.
大家看看有甚麼詞.
首先是渡過水域.
在《出一甲記》14章中.
以色列人過紅海.
摩西用手杖分開海水.
然後等以色列人過了紅海後.
水就合在一起,對吧.
如果大家記得那個故事.
然後怎樣?有甚麼死了.
是,是淹沒了軍隊.
淹沒了軍隊.
然後在路加福音中.
耶穌如何?平靜風和浪,對吧.
然後就像摩西一樣.
命令風和水都聽從他.
然後耶穌如何?帶門徒過了加里尼湖.
剛才提到湖很值得留意.
湖的原文可以翻譯為「海」.
如果你去過加里尼也知道.
其實是一個很大的湖.
它不是海,是淡水.
就像密西根一樣,是淡水.
但這個湖周不時也很大風大浪.
就像海一樣.
所以如果有些知識的人.
或有些知識的翻譯者.
也不會把這個湖翻譯為「海」.
但作者很有趣.

$^{281}$作者特意在原文中選了一個字.
這個字可以翻譯為「海」.
也可以翻譯為「湖」.
其實他有另一個字,他不選.
另一個字只能翻譯為「海」.
他特意不選這個字.
是甚麼意思.
即是說….
作者特意選了一個字.
用雙重的意思去映射當年.
上帝帶領以色列人.
藉著摩西的手帶領以色列人解放.
解放他們.
所以耶穌就像一個新的摩西一樣.
帶領門徒過了海.
第二….
對了,第二就是風的角色.
如果留意在《春夏合記》14章中.
神是用強勁的東風令紅海分開.
而在《平靜風和浪》的故事中.
耶穌斥責風和浪,令他們平靜.
就是在回應裡面的風.
耶穌是直接命令自然的力量.
這樣做是特意勾起以色列人,猶太人.
他們對於救恩的故事的回憶.
就是當日上帝解放猶太人.
令他們離開暴力的政權.
所以為何路加要強調.
鬼趕了被人上身的人兄去哪裡.
曠野.
因為曠野在記憶歷史中是救恩的開始.
是猶太人救恩的開始.
所以如果大家再看下去.
就會發現,全部人的陰霾不同.
而全部人的陰霾是一樣的.
只是有些不同.
共同主題就是敵人在海或湖泡死.
但最大的不同是在《出埃及記》中.
埃及軍隊全部泡死在紅海裡.
對吧.

$^{321}$這裡的敵人是人類的軍隊.
代表著壓迫以色列的埃及政權.
而在「趕鬼故事」.
即在路加福音裡.
鬼軍團進入了豬群後.
衝入湖泡死.
他們也像埃及軍隊般泡死.
但這裡死的是誰.
死的是豬,對吧.
不是人,是豬.
豬群的代表豬本身替代了軍團的覆沒.
再看一次,這裡有點複雜.
《出埃及記》泡死的是軍隊,對吧.
而路加福音,平靜風浪浪.
和「趕鬼故事」死的是豬,沒錯.
豬是鬼的載體.
路加花了這麼大的氣力.
去說服讀者.
他只說了一件事.
就是魔鬼才是人的真正敵人.
他們就像軍團般組織化,強大.
躲在別人背後.
利用人的表面表徵去做邪惡的事情.
所以耶穌沒有泡死人.
而是直接針對魔鬼帶來解放.
祂可以直接把人泡死.
但祂沒有這樣做.
這個訊息是要提醒我們.
我們要改變我們看人際衝突的角度.
由人彼此之間的衝突.
轉向屬靈的真戰.
我們多走一步去理解和解開誤會.
而不是加深仇恨.
這個放在自己的袋子上.
大家看到這張相片.
是2015年大約香港的10月,11月.
有很多芒草的時候拍攝的.
我曾經把某些人當成敵人.
大家看到這張相片都很憤怒.
頭髮都垂起,額頭都腫了.

$^{361}$其中一點是.
其實我以前有一段很長的時間.
我不會….
我認識朋友,第一個條件是立場.
我不會跟不同立場的人做朋友.
因為很憤怒.
我覺得這個世界有很多不公平.
有很多不公義.
其中一些跟我立場不同的人.
就是這個世界不公義的其中一個原因.
我是這樣看的,當時是這樣看的.
這個想法很偏激.
甚麼都是立場先行.
甚麼都是覺得這個世界是黑和白.
沒有中間.
所有事情,幾乎所有事情.
都是二言對立.
如果不是跟我同一個立場.
即不是朋友.
不是朋友即是甚麼?敵人.
我會覺得跟我同一個立場的人.
就是好人.
這些想法令我抽離了很多不同群體.
令我離開了一些.
我本身很珍惜的朋友圈子.
疏遠了一些很關心我的長輩.
就因為這樣.
我發現我跟他們有很多不同.
當時的我,我覺得立場很重要.
看法很重要.
我的觀點很重要.
一些比起人與人之間的關係更重要.
這些就是….
這麼多年後,10年了.
我後來回看.
魔鬼會令我們以為對方是敵人.
會影響我們跟人的關係.
我們以為疏遠了一些我們憎恨的人.
我們就會相安無事.
但原來我們疏遠它的時候.

$^{401}$也會影響到我們自身.
最近我看了一本書.
是Bella送給我的.
叫「自卑與超越」.
她是接近100年前的精神科醫生.
是佛洛伊德的好朋友.
叫Alfred Adler.
他很早期跟佛洛伊德研究人怎樣成長.
有甚麼需求.
他在書中提醒人.
普遍人的一生會完成三個主要任務.
如果沒辦法完成或處理.
很可能這個人不能很從容地看自己.
好像有一個思想陷阱.
他要麼就自卑,要麼就自大.
沒有中間,要麼自卑,要麼自大.
有時覺得自己比別人差很多.
有時覺得自己最了不起.
好像戴了一副眼鏡.
不適合他自己倒數.
他看甚麼都不喜歡.
看人,看自己都不喜歡.
他說三大任務是….
第一,職業的選擇.
他說職業不僅是賺錢那麼簡單.
還包括那份工作的意義.
自我實現.
以及那份工作對社會的貢獻.
第二,就是社會的參與.
人際關係,友誼和他人的合作.
對群體的歸屬感.
第三,親密關係的連結.
無論是好朋友還是伴侶.
意思是有人打從心裡明白你.
還能承載你的情緒.
這三個任務.
如果這三大任務都失敗.
阿德就會相信.
這個人很多時候會進入一個循環.
要麼就自卑,要麼就自大.

$^{441}$沒有中間.
結果會怎樣.
就會越來越看自己或別人.
越來越大的偏差.
最後會跟人類的群體抽離.
懷著一個全世界都欠他的心態.
最後孤獨地死去.
這樣聽起來好像很灰心.
如何解決.
盧家福音用了兩個經典故事告訴我們.
人的真正敵人是靈界的魔鬼.
而不是其他人類.
這鼓勵我們將焦點.
從人際衝突轉向屬靈的真戰.
我們踏出一步,想想.
這個人雖然做了些壞事.
得罪我或傷害了我.
但我嘗試多走一步去理解身邊的人.
看看這是否一個誤會.
看看這是否一個….
我們暫時理解不到的事.
請留意,聽完這篇講道.
你別期望一下子.
突然就會解開所有的誤會和仇恨.
是不會的.
因為我們會花很多時間和心力.
加上上帝的恩典.
慢慢抽絲剝繭.
你才能解開那些東西.
第一,我們要給自己空間.
了解我們一直以來對人的觀點.
我的故事是這樣的.
曾經…我都覺得….
我剛才說過,很多敵人在我身邊.
其中一個是我爸爸.
我小時候覺得我爸爸對我不好.
為甚麼?因為我爸爸….
我小時候,他從事生意.
他做生意時,生意困難.
基本上生意都失敗了.

$^{481}$家裡過了一個艱苦的日子.
他自己也受到很大的打擊.
同時很不想向人求助.
我心想,你做不到老闆.
你做打工仔也可以.
他不肯找工作,要我們捱餓.
我們有病,也看不到醫生.
你知道小時候的小朋友.
難道真的會走去騙錢嗎.
我們心想為甚麼要我們受苦.
後來我相信了大約10年.
之後我鼓起勇氣問他.
我問他.
我很婉轉地問他.
當年為甚麼會發生這件事.
他說當年他很後悔.
因為他無法處理生意,也找不到工作.
他盡了他的能力,只賺了一點錢.
做了很大的事.
最後還是未如理想.
後來我才開始慢慢理解.
我爸爸當時的狀態.
我解開了一些我對他的仇恨.
或一些憎恨.
在還沒解開誤會之前.
我以為我真的生我爸爸的氣.
只是影響我跟他的關係.
我以為我疏遠他,就沒事了.
原來這些誤會.
影響了我怎樣看待男性長輩.
跟男性長輩的關係.
我心想很渴望長輩會照顧我.
有甚麼事會提供適當的保護.
但同時,我也很怕他們會忽然哭.
突然離開.
可能很少的誤會.
足以令我遠離這段關係.
所以原來亞得勒這樣告訴我們.
人可以離開群體.
但群體的關係會影響我們怎樣看待自己.

$^{521}$影響我們怎樣看待身邊的人.
可想而知.
人與人之間的看法.
是多容易被扭曲.
魔鬼就是利用人與人之間.
很細微的誤會或很丁點的磨擦.
令我們彼此之間疏遠.
有些例子還會影響一輩子.
沒辦法挽回.
所以,總結.
大家知道嗎.
魔鬼令我們人與人之間彼此為仇.
人與人之間的誤會.
可以令我們彼此傷害,為敵.
而耶穌這樣說鬼就是要告訴我們.
我們的敵人絕對不是人類.
而是魔鬼.
你們人生中會否也有敵人.
你會否也有一些….
你一直記住它.
你記住它,有時睡不著.
有時你也會想起.
「我真的很想被它除掉」.
你記住它對你的不好.
曾經對你的傷害.
沒錯,那些傷害是真的.
那些傷害一定是真的.
還有些是沒辦法挽回的.
還有時我們也覺得.
這個人根本就是壞人.
與魔鬼無關.
的確,人絕對有人的責任.
我們不能將所有東西都推給魔鬼.
但面對今天的訊息.
我們可以做甚麼.
對吧.
當猶太人受到羅馬軍團的迫迫.
大部分人都是這樣看待.
這個邪惡的羅馬政權.
是因為那些人衰格.

$^{561}$所以才有傷害和壓迫.
但為何耶穌要花這麼大的氣力.
重新…這段舊人歷史.
就是因為….
祂唆使人與人之間的恩仇.
人與人之間的關係.
影響我們很深….
玩完這個星板的狼人殺.
最後我們也會知道誰是撒旦.
誰是魔鬼仔,對吧.
因為這是角色扮演.
但在現實世界裡.
可能我們每個人有一部分.
都會被魔鬼影響扭曲.
但沒有人清楚究竟那部分是人的責任.
還是魔鬼的責任.
我們可以做甚麼.
我們一起求主.
幫助我們澄清一些想法.
重新去審視我們一直堅信的觀點.
求主醫治我們曾經受的傷害.
讓我們有力氣,有勇氣,有智慧.
去認定我們的敵人只有魔鬼.
我們一起祈禱.
主,謝謝你用生命.
告訴我們.
讓我們認清我們生命裡的敵人是魔鬼.
不是我們身邊的人.
幫助我們有勇氣,有智慧.
去弄清楚一些事物.
一些我們深藏在生命裡的想法.
對我們以往熟悉的人的觀點.
重新審視我們以往的看法和觀點.
求主幫助我們,醫治我們.
讓我們看清楚這個事實.
我們這樣祈禱奉耶穌的名求.
阿們,祝福大家.
\newpage



\section{詩篇 73:1-28}
\label{sec:SzBowSpLGuQ}
\textbf{Why do the Wicked Prosper? (Psalm 73) - Rev Mark Chan 惡人亨通時 (詩篇73\_1-28) - 陳崇基牧師}
\newline
\newline
連結: \href{https://youtube.com/watch?v=SzBowSpLGuQ}{\texttt{https://youtube.com/watch?v=SzBowSpLGuQ}} ~~~~ 語音日期: 2025-11-10
\newline
\newline
\hyperref[sec:19Ss941ItL4]{\small{< < < PREV SERMON < < <}}
~
\hyperref[sec:index_chronic]{\small{[返順時目]}}
~
\hyperref[sec:index_scriptual]{\small{[返順卷目]}}
~
\hyperref[sec:I3sM6YZywrk]{\small{> > > NEXT SERMON > > >}}
\newline
\newline
詩篇 73:1-28
\newline
\begin{longtable}{cl}
\hline
\hline
章節 & 經文 (和合本修訂版)\\
\hline
73:1 & \begin{tabularx}{0.7\textwidth}{X} 神實在恩待以色列那些清心的人！ \end{tabularx} \\ \\ \relax
73:2 & \begin{tabularx}{0.7\textwidth}{X} 至於我，我的腳幾乎失閃，我的步伐險些走偏； \end{tabularx} \\ \\ \relax
73:3 & \begin{tabularx}{0.7\textwidth}{X} 因為我嫉妒狂傲的人，我看見惡人享平安。 \end{tabularx} \\ \\ \relax
73:4 & \begin{tabularx}{0.7\textwidth}{X} 他們的力氣強壯，他們死的時候也沒有疼痛。 \end{tabularx} \\ \\ \relax
73:5 & \begin{tabularx}{0.7\textwidth}{X} 他們不像別人受苦，也不像別人遭災。 \end{tabularx} \\ \\ \relax
73:6 & \begin{tabularx}{0.7\textwidth}{X} 所以，驕傲如鏈子戴在他們項上，殘暴像衣裳覆蓋在他們身上。 \end{tabularx} \\ \\ \relax
73:7 & \begin{tabularx}{0.7\textwidth}{X} 他們的眼睛因體胖而凸出，他們的內心放任不羈。 \end{tabularx} \\ \\ \relax
73:8 & \begin{tabularx}{0.7\textwidth}{X} 他們譏笑人，憑惡意說欺壓人的話。他們說話自高； \end{tabularx} \\ \\ \relax
73:9 & \begin{tabularx}{0.7\textwidth}{X} 他們的口褻瀆上天，他們的舌毀謗全地。 \end{tabularx} \\ \\ \relax
73:10 & \begin{tabularx}{0.7\textwidth}{X} 所以他的百姓歸到這裡，享受滿杯的水。 \end{tabularx} \\ \\ \relax
73:11 & \begin{tabularx}{0.7\textwidth}{X} 他們說：「神怎能曉得？至高者哪會知道呢？」 \end{tabularx} \\ \\ \relax
73:12 & \begin{tabularx}{0.7\textwidth}{X} 看哪，這就是惡人，他們常享安逸，財寶增多。 \end{tabularx} \\ \\ \relax
73:13 & \begin{tabularx}{0.7\textwidth}{X} 我實在徒然潔淨了我的心，徒然洗手表明我的無辜， \end{tabularx} \\ \\ \relax
73:14 & \begin{tabularx}{0.7\textwidth}{X} 因為我終日遭災難，每日早晨受懲治。 \end{tabularx} \\ \\ \relax
73:15 & \begin{tabularx}{0.7\textwidth}{X} 我若說「我要這樣講」，就是愧對這世代的眾兒女了。 \end{tabularx} \\ \\ \relax
73:16 & \begin{tabularx}{0.7\textwidth}{X} 我思索要明白這事，眼看實係為難， \end{tabularx} \\ \\ \relax
73:17 & \begin{tabularx}{0.7\textwidth}{X} 直到我進了神的聖所，思想他們的結局。 \end{tabularx} \\ \\ \relax
73:18 & \begin{tabularx}{0.7\textwidth}{X} 你實在把他們安放在滑地，使他們跌倒滅亡； \end{tabularx} \\ \\ \relax
73:19 & \begin{tabularx}{0.7\textwidth}{X} 他們轉眼之間成了何等荒涼！他們被驚恐滅盡了。 \end{tabularx} \\ \\ \relax
73:20 & \begin{tabularx}{0.7\textwidth}{X} 人睡醒了，怎樣看夢，主啊，你醒了也必照樣輕看他們的影像。 \end{tabularx} \\ \\ \relax
73:21 & \begin{tabularx}{0.7\textwidth}{X} 因此，我心裡苦惱，肺腑被刺。 \end{tabularx} \\ \\ \relax
73:22 & \begin{tabularx}{0.7\textwidth}{X} 我這樣愚昧無知，在你面前如同畜牲。 \end{tabularx} \\ \\ \relax
73:23 & \begin{tabularx}{0.7\textwidth}{X} 然而，我常與你同在；你攙扶我的右手。 \end{tabularx} \\ \\ \relax
73:24 & \begin{tabularx}{0.7\textwidth}{X} 你要以你的訓言引導我，以後你必接我到榮耀裡。 \end{tabularx} \\ \\ \relax
73:25 & \begin{tabularx}{0.7\textwidth}{X} 除你以外，在天上我有誰呢？除你以外，在地上我也沒有所愛慕的。 \end{tabularx} \\ \\ \relax
73:26 & \begin{tabularx}{0.7\textwidth}{X} 我的肉體和我的心腸衰殘；但神是我心裡的力量，又是我的福分，直到永遠。 \end{tabularx} \\ \\ \relax
73:27 & \begin{tabularx}{0.7\textwidth}{X} 看哪，遠離你的，必要死亡；凡離棄你行淫的，你都滅絕了。 \end{tabularx} \\ \\ \relax
73:28 & \begin{tabularx}{0.7\textwidth}{X} 但我親近神是於我有益；我以主耶和華為我的避難所，好叫我述說你一切的作為。 \end{tabularx} \\ \\
[1ex]
\hline
\hline
\end{longtable}
$^{1}$Good morning, everyone..
I know today is a very cold and windy day, right?.
[NON-ENGLISH SPEECH].
Why don't we welcome our neighbors,.
give them a warm handshake, find two friends, OK?.
Go around the room and find at least two friends.
and give them a warm handshake..
All right?.
[NON-ENGLISH SPEECH].
Let's do it now..
[NON-ENGLISH SPEECH].
OK, I hate to stop here, but I know.
we have so much to talk about, which.
is a good thing because we are a whole family together,.
a whole community worshiping God..
[NON-ENGLISH SPEECH].
OK, that's good..
Well, today, we're going to continue.
with our meditation and reflection.
on some of the psalms in the Bible..
[NON-ENGLISH SPEECH].
We will turn to Psalm 73..
[NON-ENGLISH SPEECH].
Let me read the psalm, and I'm sure the psalm will speak.
to all of us, every one of us..
[NON-ENGLISH SPEECH].
All right, verse 1 to 3..
"Truly, God is good to Israel, to those who are pure in heart..
But as for me, my feet had almost slipped..
I nearly stumbled and fell, for I.
envied the arrogant when I saw this prosperity of the wicked.".
[NON-ENGLISH SPEECH].
And then verse 4, "They have no pain until their death..
Their bodies are healthy and strong..
They are not plagued by the troubles coming to humanity..
They are not stricken like other people..
Therefore, pride is their necklace..
Violence clothes them like a garment.".
And we skip to verse 12, "Look at these wicked people..
They are always at ease..

$^{41}$Their wealth keeps increasing.".
[NON-ENGLISH SPEECH].
So when you read the psalms, it talks.
about the psalmist being envious against some of the wicked.
people that he saw..
[NON-ENGLISH SPEECH].
Have you felt this way recently?.
[NON-ENGLISH SPEECH].
Think of some of the bad people in your life..
[NON-ENGLISH SPEECH].
Are there any people you can think of?.
They are bad people, but they prosper..
[NON-ENGLISH SPEECH].
How many of you can think of some of the faces, bad people?.
I can see them..
[NON-ENGLISH SPEECH].
No?.
Now, the first part of Psalm 73 is verses 1 to 12..
And verse 1, as you can tell, it says,.
"Truly, God is good to Israel, to those who are pure in heart.".
And this statement is a statement of faith..
It's an Old Testament or the biblical conviction..
It says that whoever keeps the covenant.
will be cared for by God and will be blessed by God.
and because God is good..
[NON-ENGLISH SPEECH].
And for example, you can think of Psalm 1, right?.
If you can remember Psalm 1, Psalm 1.
talks about a blessed person who fears the Lord.
and lives in a clear conscience..
And he walks his life uprightly, and God.
blesses such a person..
And we also teach our children, our next generation,.
to do the same thing..
And here, Psalm 73, the psalmist is actually.
an Israelite who trusts in God..
[NON-ENGLISH SPEECH].
And then in verse 2, it says something.
that contradicts the reality..
From verses 2 to 12, the psalmist complains and testifies.

$^{81}$that he almost fell from his faith,.
fell from his faithfulness to God..
Verses 2 to 6, let me summarize..
It kind of talks about he envied the arrogant, the wicked people.
that he saw because he saw them with prosperity..
Their wealth increasing and all that..
They are so blessed, and he was very angry..
[NON-ENGLISH SPEECH].
And the wicked and the arrogant that he saw,.
they do not keep the law of God..
They do not keep the Torah, and they do not believe in God..
But they have power, and they have privilege..
And yet, they enjoy peace..
You see that they had no pain, and they have no calamity..
Their health plan or their medical insurance.
seems perfect, but we don't have those..
From the outside, you know, these wicked people,.
they are prospering..
[NON-ENGLISH SPEECH].
And then in verse 7, it says their eyes.
spilled from fatness..
And you can imagine the image..
You know, in the ancient world, most of the people,.
they are undernourished..
They don't have enough food..
So a fat body, a fat man, a fat guy in the ancient world.
means that he must have been exploiting or grabbing.
something from other people..
[NON-ENGLISH SPEECH].
,.
And then verse 8, it says they, the wicked people,.
speak with malice about oppressing other people..
And their words are lofty..
And you know, the wicked people described here.
do not always use foul language..
They may speak nicely, and they can be cultured..
They can be polite..
Yet their speech actually oppresses other people.
and exhausts themselves..
[NON-ENGLISH SPEECH].

$^{121}$And in verse 9, they set themselves against heavens..
They challenged the authority of God..
And their tongues devoured or destroyed the earth..
[NON-ENGLISH SPEECH].
And then Hebrew poetry often use two lines to echo one another..
The second line usually expands the first one..
And the picture here is like the wicked--.
it's like a lion..
It's like a beast..
And with mouth open wide, with the upper lip.
lifted up to heavens, and the lower lip to the earth,.
and devour everything..
And they are greedy, and they are without limit..
They are trying to destroy everything..
[NON-ENGLISH SPEECH].
And then verse 10, "Therefore his people.
turned to them, turned to the wicked people,.
and drained down waters in abundance.".
And they say, how does God know that's.
the most high--.
does God have knowledge?.
Verse 12, "Look, these are the wicked, always at ease..
Their riches increase.".
[NON-ENGLISH SPEECH].
So these psalms with 12 verses try.
to describe the wicked, the arrogant people, flourishing.
and running wild..
And we may see more of this kind of people in our day.
and the days to come..
And so often, God's people suffer..
And as you can tell from all sorts of stories.
around the world, Christians or godly people.
often are being persecuted..
[NON-ENGLISH SPEECH].
,.
And then the psalm continues to the second part in verse 13.
and following..
"Surely in vain I have kept my heart pure..
In vain I wash my hands in innocence..
For I am afflicted all day long and disciplined every morning..

$^{161}$If I had said I would speak this way,.
I would have betrayed your children..
When I try to understand all of this,.
the wickedness in the world, it seems.
very williesome burden to me.".
[NON-ENGLISH SPEECH].
And the psalm is trying to tell us that he sees injustice.
everywhere and he is tempted to abandon God,.
abandon his belief in God..
Because the whole psalm is trying to paint the picture.
that it descends from verse 1..
And then step by step by step, the psalmist.
sees all the wickedness and his heart is burned with envy..
And everything seems in vain for him..
And that is our tendency as well..
Sometimes we let the problem define God..
And instead of letting God define the problem,.
and we see the wicked prosper, and we conclude.
that God is not there, and God is unjust, or God is absent..
And that is the classic problem of evil in the Bible.
and in our world..
And the psalmist's bottom line is his conscience..
And he refuses to betray God's people..
[NON-ENGLISH SPEECH].
,.
And you can tell the debilitiness.
that he is experiencing, seeing wickedness in his world..
And you can feel--.
we can feel what he felt..
[NON-ENGLISH SPEECH].
However, in verse 17, there's a turnaround point.
in the entire psalm..
And that's the turning point..
Until I went into the sanctuary of God,.
then I discerned the end of the wicked people..
Now, in the sanctuary, in worship,.
in a gathered community of worship, in psalm,.
the psalmist suddenly has some sort of turning point.
in his mindset..
Perhaps it's like the vision of Isaiah, if you remember..

$^{201}$When he entered the temple, when he entered the worship of God,.
he saw that God is still lifted high,.
and he is still on the throne..
And his entire perspective changes..
[NON-ENGLISH SPEECH].
And that's why we need to come and worship God.
with the right perspective..
When we see God in worship, then our life perspective changes..
And there are three points here I want to share..
[NON-ENGLISH SPEECH].
And the first point is we have to see the end of the wicked,.
like the psalmist did..
Verses 18 to 20, surely, the psalmist.
says, you, God, you set them in the slippery places,.
and you make them fall to ruin..
How suddenly they are destroyed, swept away by terrors,.
like a dream when one awakes, O Lord..
When you rouse yourself, you despise them as phantoms..
Now, evil will end..
Evil will expire..
The wicked will not endure forever..
Now, you can think of a long drama, like a TV series, OK?.
The Bad Guy..
The Bad Guy dominates one episode, or perhaps.
a few episodes..
But if you watch only those episodes, you'll be angry,.
and you'll be upset..
But here, the psalmist says, you have.
to watch the entire series..
And God is the author, and he writes the true ending.
for the wicked..
[NON-ENGLISH SPEECH].
So the first point is to see the wicked, see the end..
And the second point here is that we.
have to admit that we are short-sighted at times,.
verses 21 to 22..
When my heart was ambidextrous, and I was pierced within,.
and I was ignorant, and I was like a beast before you, God..
The psalmist confesses he was foolish..
And we are all like that..

$^{241}$Sometimes we are living in a fog,.
and we couldn't see the long range..
But in the sanctuary of God, God lifts us up out.
of our short-sightedness, and let us see eternity..
Without eternity, of course, there's no justice..
But there is eternity, so there has to be justice at the end..
[NON-ENGLISH SPEECH].
So the first point is we have to see the end of the wicked.
from the perspective of eternity..
And the second one is like the psalmist,.
admit that we are short-sighted at times..
And the third point is that we have to experience.
the presence of God..
Verses 23 and following, let me just read that..
I think this is self-explanatory..
23, "Yet I am with you, God..
You hold me by my right hand, and you guide me.
with your counsel..
And afterward, you will take me into glory..
Whom have I in heaven but you?.
And on earth, there is nothing I decide besides you..
My flesh and my heart may fail, but God.
is the rock of my heart and my portion forever..
Those far from you will perish..
You destroy all who are unfaithful to you..
But as for me, it's good to be near God..
I have made the Lord my God, my refuge..
And I may tell all your works.".
[NON-ENGLISH SPEECH].
the Lord..
[NON-ENGLISH SPEECH].
So you see, Psalm 73 from verse 1,.
it's like the ladder of faith and descending step.
in a time when the psalmist sees the wicked raging.
and his heart burns with envy..
And then there's a turning point at the middle of the psalm..
And then he turned around and he ascend again, step by step,.
upward, when he enters the sanctuary..
[NON-ENGLISH SPEECH].
Now, many of us, I asked you the question at the beginning,.

$^{281}$felt this way, or we feel this way right now..
And many of us repeat this cycle more than once in our lives..
When problems drive us downward and we.
are tempted to refuse God and to abandon God,.
the sanctuary is the turning point..
And it's a signpost to tell us that God's kingdom.
and the existence of God is real..
And in worship, we recover that sense,.
that perspective of eternity, justice,.
and the nearness of God..
[NON-ENGLISH SPEECH].
Well, I would like to take this time to pray together..
[NON-ENGLISH SPEECH].
When I asked you the first question,.
are there any bad and wicked people.
that you envy in your life?.
Of course you do..
[NON-ENGLISH SPEECH].
Maybe you are tempted to abandon God, to say no to God..
But this psalm is inviting us to come to worship him again..
[NON-ENGLISH SPEECH].
So that we can see our life and see.
all the injustice in the world with the perspective,.
eternal perspective of God..
[NON-ENGLISH SPEECH].
Let's all pray together..
[NON-ENGLISH SPEECH].
Dear God, we come before you..
[NON-ENGLISH SPEECH].
We see so many wickedness and evil in our life..
[NON-ENGLISH SPEECH].
We envy and we are angry..
[NON-ENGLISH SPEECH].
But we want to thank you because you listen to our cries..
[NON-ENGLISH SPEECH].
And you listen to our prayer and you listen to our complaints..
[NON-ENGLISH SPEECH].
God, you are so good..
[NON-ENGLISH SPEECH].
Help us to come before your presence today..

$^{321}$[NON-ENGLISH SPEECH].
And know that you are eternal..
[NON-ENGLISH SPEECH].
And you are full of justice..
[NON-ENGLISH SPEECH].
And you see the end of the wicked..
[NON-ENGLISH SPEECH].
And you treasure our faith in you..
[NON-ENGLISH SPEECH].
Help us to increase our faith in you..
[NON-ENGLISH SPEECH].
And trust that you will work things out for your glory..
[NON-ENGLISH SPEECH].
And again, we entrust our lives in your hands..
[NON-ENGLISH SPEECH].
Protect us..
[NON-ENGLISH SPEECH].
And guide us..
[NON-ENGLISH SPEECH].
In Christ's name we pray..
[NON-ENGLISH SPEECH].
Amen..
Well, God bless you all..
And you will be sent out with the love and grace of God..
And to be strengthened this week to bear good witness.
for our Lord Jesus Christ..
[NON-ENGLISH SPEECH].
(soft music).
\newpage



\section{傳道書 2:12-26}
\label{sec:I3sM6YZywrk}
\textbf{感謝天父賜我飲食 (傳道書2\_12-26) - 譚文倩莊傳道}
\newline
\newline
連結: \href{https://youtube.com/watch?v=I3sM6YZywrk}{\texttt{https://youtube.com/watch?v=I3sM6YZywrk}} ~~~~ 語音日期: 2025-11-10
\newline
\newline
\hyperref[sec:SzBowSpLGuQ]{\small{< < < PREV SERMON < < <}}
~
\hyperref[sec:index_chronic]{\small{[返順時目]}}
~
\hyperref[sec:index_scriptual]{\small{[返順卷目]}}
~
\hyperref[sec:lX_7jxrhzjs]{\small{> > > NEXT SERMON > > >}}
\newline
\newline
傳道書 2:12-26
\newline
\begin{longtable}{cl}
\hline
\hline
章節 & 經文 (和合本修訂版)\\
\hline
2:12 & \begin{tabularx}{0.7\textwidth}{X} 我轉而回顧智慧、狂妄和愚昧。在王以後來的人又如何呢？不過做先前所做的就是了。 \end{tabularx} \\ \\ \relax
2:13 & \begin{tabularx}{0.7\textwidth}{X} 於是我看出智慧勝過愚昧，如同光明勝過黑暗。 \end{tabularx} \\ \\ \relax
2:14 & \begin{tabularx}{0.7\textwidth}{X} 智慧人的眼目光明，愚昧人卻在黑暗裡行。但我知道他們都有相同的遭遇。 \end{tabularx} \\ \\ \relax
2:15 & \begin{tabularx}{0.7\textwidth}{X} 我心裡就說：「愚昧人所遇見的，我也一樣遇見，那麼我何必更有智慧呢？」我心裡說：「這也是虛空。」 \end{tabularx} \\ \\ \relax
2:16 & \begin{tabularx}{0.7\textwidth}{X} 智慧人和愚昧人一樣，不會長久被人記念，因為日後都被遺忘。可嘆！智慧人和愚昧人都一樣會死亡。 \end{tabularx} \\ \\ \relax
2:17 & \begin{tabularx}{0.7\textwidth}{X} 於是我恨惡生命，因為在日光之下所發生的事我都以為煩惱，全是虛空，全是捕風。 \end{tabularx} \\ \\ \relax
2:18 & \begin{tabularx}{0.7\textwidth}{X} 我恨惡一切的勞碌，就是我在日光之下所勞碌的，因為我所得的必須留給我以後的人。 \end{tabularx} \\ \\ \relax
2:19 & \begin{tabularx}{0.7\textwidth}{X} 那人是智慧是愚昧，誰能知道呢？他竟要掌管我在日光之下用智慧勞碌所得的。這也是虛空。 \end{tabularx} \\ \\ \relax
2:20 & \begin{tabularx}{0.7\textwidth}{X} 我轉想我在日光之下所勞碌的一切工作，心就絕望。 \end{tabularx} \\ \\ \relax
2:21 & \begin{tabularx}{0.7\textwidth}{X} 因為有人用智慧、知識、靈巧勞碌工作，所得來的卻要遺留給未曾勞碌的人作產業。這也是虛空，大大不幸。 \end{tabularx} \\ \\ \relax
2:22 & \begin{tabularx}{0.7\textwidth}{X} 人一切的勞碌操心，就是他在日光之下所勞碌的，又得著了甚麼呢？ \end{tabularx} \\ \\ \relax
2:23 & \begin{tabularx}{0.7\textwidth}{X} 他日日憂慮，他的勞苦成為愁煩，連夜間心也不得休息。這也是虛空。 \end{tabularx} \\ \\ \relax
2:24 & \begin{tabularx}{0.7\textwidth}{X} 難道一個人有吃有喝，且在勞碌中享福，不是福氣嗎？我看這也是出於神的手。 \end{tabularx} \\ \\ \relax
2:25 & \begin{tabularx}{0.7\textwidth}{X} 論到吃用、享福，誰能勝過我呢？ \end{tabularx} \\ \\ \relax
2:26 & \begin{tabularx}{0.7\textwidth}{X} 神喜愛誰，就給誰智慧、知識和喜樂；惟有罪人，神使他勞苦，將他所儲藏、所堆積的歸給神所喜愛的人。這也是虛空，也是捕風。 \end{tabularx} \\ \\
[1ex]
\hline
\hline
\end{longtable}
$^{1}$早晨各位弟兄姊妹.
無驚無險可以回到你們中間.
因為大家都知道這個政府封鎖.
突然收到很多消息.
可能航班會被取消.
我們上機之前那天就告訴我們.
如果你沒有被取消的時候.
飛機大概會遲四至八個小時.
所以我帶了很多食物在我的袋子裡.
昨晚平安到達.
兩個星期前左右.
我在香港在芝加哥的Facebook群組裡.
回應一個帖文.
我知道有人知道我在說哪個帖文.
因為我們UV裡也有人Like和寫一些東西.
一般來說,我不會在這些群組寫什麼.
或者群組裡有很多人我不認識.
我不會寫東西.
但這次有一個人寫了一些東西.
是我自己很熱衷的.
很緊張的東西.
所以我寫了一些回應.
我覺得我的回應很中立.
應該也沒有什麼.
我問問大家.
18元一碗年稅連貼士的牛腩麵.
划算嗎?.
划算還是不划算?.
大家當然有不同的想法.
我這個會買牛腩回家.
自己做潮州清湯牛腩的主婦.
知道買五磅牛腩回家.
花時間和材料.
用心做完之後.
可能有四磅左右.
沒有.
燕姨說沒有.
三磅多就OK.
可能變三磅.
所以對我來說.

$^{41}$一碗大大碗足料的牛腩麵.
18元年稅連貼士.
合理吧.
但想不到這個帖子.
有很多人加入討論.
四天之內.
四天而已.
竟然有67個留言.
當中還有人因為.
我們其中幾個覺得合理的回應.
而生氣.
我覺得很有趣.
竟然有這麼多人回應.
我心想.
我很想看看這個群組裡.
還有什麼帖子.
有這麼多人回應.
我就不停地掃.
掃到有點累.
不掃的時候.
你猜猜另一個帖子.
有雙位數字人回應的是什麼帖子.
其實大部分帖子沒有人回應.
或者有的都是單位數字.
你猜猜下一個是雙位數字多人回應的帖子.
是什麼帖子.
是說什麼的.
是魚肉燒賣.
20個帖子左右.
你是不是發覺吃果然真的很重要.
之前我跟大家說過傳道書的第一章.
說到傳道者積極地去尋問.
去試驗.
去思考人生的意義.
他發現人在日光之下.
勞碌都是徒勞的.
也是虛空的.
到了二章一至十一節.
傳道者嘗試透過享樂.
財富.

$^{81}$工作.
和物質上去尋求滿足.
但他最終發現.
原來這些東西都不是可以帶來持久的快樂.
或者滿足.
或者人生的意義.
好像嘗試捉住一陣風.
今天我就跟大家繼續看這個傳道書.
到了二章十二至二十六節.
三個主題去思考.
這句我們基督徒吃每一餐飯都會說的話.
我們先低頭祈禱.
出去主我們一神.
在聖經裡面要跟我們說的話.
說出來.
願你的說話能滋潤我們的生命.
你能夠在當中得著榮耀.
奉主耶穌基督的聖名祈求.
阿們.
我先跟大家讀一讀一段經文.
我轉過來觀察智慧狂妄和愚昧.
王的繼承人還能作什麼呢.
只不過作先前作過的罷了.
於是我看出智慧勝過愚昧.
如同光明勝過黑暗一樣.
智慧人的眼目明亮.
愚昧人卻走在黑暗中.
我也知道同一命運要臨到這兩種人身上.
我自己心裡說.
臨到愚昧人身上的也必臨到我身上.
那麼我為什麼要更有智慧呢.
於是我心裡說.
這也是虛空.
因為智慧人和愚昧人一樣.
沒有人永遠紀念他們.
在未來的日子裡都會被人遺忘.
可嘆智慧人跟愚昧人一樣終必死亡.
因此我恨惡生命.
因為在日光之下所發生的事.
都使我厭煩.

$^{121}$一切都是虛空.
都是暴風.
二章十八至二十三節.
我恨惡我的一切勞碌.
就是我在日光之下的勞碌.
因為我不得不把我勞碌的成果留給後人.
那人將來自知是如.
有誰曉得呢.
可是他竟要掌管一切我勞碌得來的.
就是我在日光之下用智慧得來的.
這又是虛空.
於是我轉念.
對日光之下我所勞碌的一切深感絕望.
因為人用智慧,知識和技能所勞碌得來的.
卻必留給未曾勞碌的人為份.
這也是虛空.
是極大的憾事.
人的一切勞碌和操心.
就是他在日光之下的勞碌的.
對自己有什麼益處呢.
其實中期一生.
他的膽子既痛苦又煩惱.
在夜間也不能安心.
這也是虛空.
二章二十四至二十六節.
人最好是吃喝在自己的勞碌中.
自得其樂.
我看這也是出於神的手.
離開了他誰能有吃的呢.
誰能享樂呢.
神把智慧,知識和喜樂賜給他所喜悅的人.
卻使罪人作收集和積聚的差事.
把所得的歸給神所喜悅的人.
這也是虛空.
也是暴風.
從二章十二節開始.
所羅門王從之前審視建造一些實質的事物之後.
轉移到審視智慧的價值.
我想大家都想到.
結論當然是有智慧比沒有智慧好.

$^{161}$因為有智慧的人有洞察的能力.
知道什麼時候做什麼.
或者什麼時候不要做什麼.
但相對如每人卻缺乏了這種洞察力.
就像這裡所說.
如每人卻走在黑暗中.
即是在黑暗中好像盲中中亂闖.
接著我們再看二章十六節.
因為智慧人和如每人一樣.
沒有人永遠紀念他們.
在未來的日子裡都會被遺忘.
可嘆智慧人跟如每人一樣終必死亡.
Yes 人人平等.
我們每一個人最終都會死.
但所羅門王He is so 偉大.
即使他死了.
有什麼可能人們會永遠記得他,紀念他.
在未來的日子裡會被人遺忘.
可嘆智慧人.
這裡說可嘆智慧人跟如每人一樣.
甚至最終會死亡,被遺忘.
其實你想想.
他會這樣問.
他不開心,他生氣.
其實都是很合理的.
2018年英國University of Aberdeen.
做了一個心理學的研究發現.
被遺忘的感覺就好像被忽略,被排斥.
這些感覺可以是痛苦的來源.
也可以引發焦慮症.
即是引發這個anxiety.
你有沒有試過被忽略.
有沒有試過覺得別人排斥你.
我小時候經常都有這個感覺.
沒辦法.
你看看我的樣子.
我的樣子很硬.
不笑的時候又好像很兇.
又有些男生頭.
這些是我經常聽到別人給我的評論.

$^{201}$我小時候有一個漂亮,可愛的表姐.
我很羨慕她.
因為每個親戚朋友都很疼她.
很喜歡她.
買東西給她.
但我經常被人罵.
不是罵這個,是罵那個.
外出又沒有我的份.
很多時候都是我的親戚和我表姐外出後.
拿著一袋,二袋回家的時候.
我才知道原來他們去了購物.
我那時候經常在想.
為什麼不叫我呢.
為什麼我沒有份呢.
在《傳道書2章》這裡.
所羅門王擔心別人會忘記他.
因為他認識到這個現實.
無論一個人多有智慧,成就.
或者地位多高.
有智慧的人和愚者.
到死亡之後.
都會遇到同一個命運.
就是很快被後世的人忘記.
這種存在主義的關注.
是《傳道書》頭幾章的主題.
你有沒有想過.
你自己有沒有想過.
為什麼我會想別人記住我呢.
其實有人永遠紀念自己.
是否重要.
或者是否真的那麼重要呢.
有時候我去參加那些安息禮拜的時候.
那些親戚朋友會上來說.
那個人有多好.
我坐在那裡的時候.
我就在想.
很可惜那個人不在.
他聽不到.
其實這個人在生的時候.
那些人有沒有告訴他呢.

$^{241}$有沒有告訴他.
其實這些事情我是很欣賞的.
或者你曾經做過一些事情影響過我.
是很正面的.
這些話會不會是.
在這個人在生的時候.
去跟他說會好一點呢.
二章十八至二十一節.
我恨惡我的一切勞碌.
就是我在日光之下的勞碌.
因為我不得不把勞碌的成果留給後人.
那人將來是志是如.
有誰曉得呢.
可是他竟要掌管我一切勞碌得來的.
就是我在日光之下用智慧得來的.
這又是虛空.
於是我轉念.
對日光之下我所勞碌的一切.
深感絕望.
因為人用智慧,知識和技能所勞碌得來的.
卻必須留給未曾勞碌的人為訓.
這也是虛空.
是極大的陷阱.
多年的努力.
最終我不得不把勞碌的成果留給後人.
還有那人將來是志是如.
有誰曉得呢.
他是竟要掌管我勞碌得來的.
再想想.
人用智慧,知識和技能所勞碌得來的.
卻必須留給未曾勞碌過的人為訓.
這也是虛空.
是極大的陷阱.
又真的頗要命的.
《舊約列王記》上十二,十四章記載.
所羅門王的兒子羅波安繼承皇位之後.
老百姓請求他減輕.
他父親時代留下來的徵收重稅.
他先去詢問一些年長,有經驗的長老.
得到一些比較溫和的建議.

$^{281}$但大家都知道.
通常年輕人不喜歡聽年長老人家的意見.
忠言逆耳.
他反而接納了身邊一些年輕臣僕驕傲的意見.
決定更加加重他們的重擔.
這個愚蠢的回應激怒了老百姓.
導致王國分裂.
北方十個支派叛變.
擁納耶羅波安為王.
建立北國的以色列.
而南方就只剩下猶大和汴也汶兩個支派.
繼續由羅波安統治.
稱為南國猶大.
Long story short.
羅波安不單止令國家分裂.
聖殿裡的財寶全部都被搶掠.
哪人將來是志是如.
有誰曉得呢.
我們這一代比上一代的環境好.
所以我們的小朋友物質比我們年代豐富.
就算不是特別有錢.
我們都會盡量讓他們過一些好的生活.
讓他們開心.
我們愛他們.
愛他們.
祝福他們.
讓他們開開心心.
讓他們不用那麼辛苦.
不要受人氣.
做少一點.
最好不用做.
結果害死他們.
傳道書二章之前的一章十節.
他說:因為我的心.
要為一切的勞碌快樂.
這是我從一切勞碌中所得的報償.
你知不知道上帝給人的祝福.
不單止是報酬的祝福.
還有我們努力.
我們人的勞碌裡的祝福.

$^{321}$在勞碌裡能夠得到成功感.
那個目的感和意義.
正正是這段經文裡所講的.
所得的報償.
Toil is a blessing.
現今世代有種新的職業道德.
不知道你們知不知道.
就是做少一點.
但賺多一點.
做少一點但賺多一點.
很有趣.
幾個月前有位中學生.
跟我說.
Joyce我的理想.
將來我要找一份工作不用做.
不用怎樣做.
但賺很多錢.
我心想多有趣.
這是我的理想.
記不記得之前我和大家說過.
《聖經新約》裡面的哥林多後書八至九章.
那裡是有基督徒處理金錢的教導.
哥林多後書九章十至十三節.
但那次種子給撒種的.
又次食物給人吃的神.
必定把種子加倍供給你們.
也必增添你們的異果.
你們既然凡事呼籲.
就可以慷慨地捐書.
使眾人借著我們對神的心.
因為這供應的事.
不僅補足了聖徒的缺乏.
眾聖徒因為你們承認和服從了基督的福音.
並且慷慨地捐書給他們和眾人.
藉著你們在這供應的事上所得的貧懼.
就把榮耀歸給神.
簡單來說.
這裡是說你要將神給你的東西分開兩份.
沒有說百分比.
但是是說分開兩份.

$^{361}$是說神給你的東西.
一份是給你的麵包.
另外一份是用來播種的種子.
結果子的種子.
第一份的麵包.
你要傳著一個感恩的心去享受.
是神給你的.
你拿那份麵包用來養家.
可能要交學費去旅行.
積谷防饑.
Go ahead 你是可以用的.
你可以享受的.
但是你不要忘記.
給你的不是所有都是麵包.
因為另外那一部份是種子.
哥林多後書九章的教導.
是說我們要set aside一份種子.
將它放在奉獻在神的國度裡面.
這個動作是承認和服從了基督的福音.
驅使人藉著我們慷慨所得的憑據.
就把榮耀歸給神的果子.
我們的問題不是吃麵包.
而是我們連種子都吃了.
我在2014年的時候就入神學院讀書.
Double major讀了四年半才讀完.
想當年每個學期開學之前.
有個好姊妹都會問我.
你有沒有什麼需要.
每個學期她都會奉獻給我.
幫補我交學費.
幾年來她給我的奉獻.
我自己在想.
其實她是可以拿這些錢.
買個漂亮的手袋給她自己.
我有很多這些收到種子的故事.
這些弟兄姊妹對我的慷慨.
犧牲買個漂亮手袋的捐贈.
成為了我永遠有很多不同感恩的故事.
去和別人分享.
同樣如果你認識我.

$^{401}$我都是一個很慷慨的人.
有時甚至慷慨過頭.
所以有些人以為我很有錢.
其實我不是很有錢.
但我都很慷慨.
是因為神從來對我都很慷慨.
神從來都沒有孤寒過.
我很希望我的慷慨.
我的頭腦.
我的犧牲.
是可以建立到人.
建立到神的國.
神的事公.
教會的事公.
想神在當中得到榮耀.
這正是哥林多後書八至九章.
講到的基督徒金錢的教導.
如果你從來沒有這個習慣.
我很想你今天去思考一下.
以至你會放下一些種子.
放在神的國度裡.
It's not too late.
You still have time.
有一天我和我老公出去吃飯.
坐在那裡等食物來的時候.
我和我老公聊天.
講一下我的女兒.
講一下教會.
講一下事公.
講一下我們的朋友.
終於食物來了.
不是先祈禱.
當然是拿手機出來拍照.
然後發給女兒.
然後我和我老公就在那裡.
當然是拍照.
然後發給女兒.
然後我們才低頭祈禱.
一打開眼睛.
哇,食物來了.

$^{441}$我就看看廚師的擺盤.
顏色.
聞一下有多香.
然後舀一塊.
吃一口.
哇,好好味.
真的很好味.
這間是我們很喜歡的餐館.
突然間我有個意念出現.
我就問我老公.
我說,爸爸.
這個是不是傳道書二章二十四節.
那一段.
講到神的手的意思.
我老公說,是的.
神的手就在每一個細節裡.
多麼好.
2022年回香港的時候.
有個傳道人好朋友.
堅持一定要帶我去看一套電影.
那套電影叫做《過時過節》.
我老公就說.
這套電影叫做《過時過節》.
故事講述一個家庭.
八年前一家人做冬的時候.
發生了一些誤會.
導致這個家庭的關係破裂.
當然,冰封三尺.
從來都是飛一日之寒.
多年來家庭成員之間的矛盾.
以前從來沒有化解的紛爭.
每個人都藏著很多不開心的怨恨.
很多傷痛.
因為這一頓飯之後.
大家都是以後最好不要見面.
這套《過時過節》.
導演說.
他很想大家看這套電影之後去想.
既然大家是一家人.
有沒有可能.

$^{481}$一人退一步.
放下自己.
再坐下來一起過冬.
親愛的傳道人.
跟我看這套電影之前.
原來他已經看了兩次.
但是他跟我看的時候.
我留意到.
他拿一張紙巾出來.
不停地抹眼淚.
已經看了第三次了.
為什麼會有這麼多傷痛?.
為什麼會這麼傷心?.
這張24-25節.
人最好是吃喝在自己的牢籠中自得其樂.
我看這也是出於神的手.
離開了他.
誰能有吃的呢?.
誰能享樂呢?.
在傳道書裡面.
這裡是第一次提到神的手.
讓人享受飲食和在身牢中的喜樂.
這些都是神給人的禮物.
在傳道書裡面還有六段經文.
都是提到這一點.
這不是說.
反正人生都是這麼虛空.
都是這麼徒勞無功.
所以不如今天有酒就今早醉吧.
的態度.
而是他要和應《創世記》和《箴言》裡面.
提到飲食和享受自己身牢的喜樂.
的一個正面觀點.
神選擇每天都向你施恩.
As simple as.
說你吃的牛腩麵.
煲仔飯.
魚肉燒賣.
蒸肉餅.
每一餐飯.

$^{521}$每一杯水.
神的手.
在那些細節裡面.
我們UV教會就很知道吃的重要性.
所以我們每個星期日.
有不同的弟兄姊妹為我們預備早餐.
煲好咖啡.
我不知道你有沒有機會去過其他教會.
你知不知道其實大部份的教會沒有這個事工.
其實我們真的很感恩.
有一班弟兄姊妹這麼願意付上時間和心機.
為我們預備.
好像剛才傑少說的.
還在煎餃子.
我們UV教會是教會裡面一個很重要的事工.
因為我們教睦同工都知道.
吃可以是一個建立彼此關係的工具.
如果你將你和你家人吃飯的情景.
搬來這個建立彼此關係事工的概念裡面.
我想問問你腦海裡面出現了一幅甚麼圖畫.
你是和你家人建立關係.
還是會像過時過節那樣想起很多爭執.
想起很多傷痛.
如果你心目中的圖畫.
是和我形容第二個像過時過節的圖畫.
我想你知.
我明白了.
因為我自己的家以前都是這樣.
已經五十多歲的人.
有時我想起一些事的時候.
那個傷痛到現在想起.
仍然會勾起我很多情緒.
如果你像我這樣.
我想你更加知道一件事.
你更加抓緊一件事.
你可以把權力拿回來.
去改變.
你可以從神裡面的因點取回能力去改變這件事.
就算你仍然有不舒服的傷患.
你要知道神的因點大於所有這些.

$^{561}$你可以不從肚腹切.
你的下一代不需要經歷同樣的傷害.
由你這一代去切斷.
之前我不是和大家說過.
因為Facebook的有趣經歷.
發現吃這東西是很重要的.
套用在今天這段經文裡面.
不知道你有沒有想過.
其實神是很知道我們的需要.
祂真的很知道我們的需要.
所以當你下次吃飯的時候.
我希望你想起今天這個訊息.
你想起神的手.
就在每一個細節裡面.
我祈禱你真的可以每餐飯吃的時候.
你真的存著一個感恩的心.
不是理所當然的.
也好好地去吃那餐飯.
去享受和你一起吃飯的人.
你的家人.
吃飯的時候不要只顧著應該這樣.
還沒做功課或者罵一些其他東西.
那個地方一起吃飯的時候.
可以是一個建立大家關係的時候.
認真的.
去享受神給我們的簡單的快樂.
因為這個是日光之下.
神給我們人最好的禮物.
我們一起低頭祈禱.
很大的恩典.
我求你也給我們有能力去作出改變.
留意你精心為我們預備.
的每一個筵席.
也幫助我們.
在你的國度裡有份.
我們要留一份種子.
去擺在神的國度裡.
去建立神的國度.
願你在當中得著榮耀.
主啊多謝你.

$^{601}$這樣祈禱.
奉主耶穌基督成名祈求.
Amen.
\newpage



\section{詩篇 8:1-9}
\label{sec:lX_7jxrhzjs}
\textbf{True Worship (Psalm 8) - Rev. Mark Chan 真實的敬拜 (詩篇8\_1-9) - 陳崇基牧師}
\newline
\newline
連結: \href{https://youtube.com/watch?v=lX-7jxrhzjs}{\texttt{https://youtube.com/watch?v=lX-7jxrhzjs}} ~~~~ 語音日期: 2025-11-18
\newline
\newline
\hyperref[sec:I3sM6YZywrk]{\small{< < < PREV SERMON < < <}}
~
\hyperref[sec:index_chronic]{\small{[返順時目]}}
~
\hyperref[sec:index_scriptual]{\small{[返順卷目]}}
~
\hyperref[sec:mJ4c0qfZcx0]{\small{> > > NEXT SERMON > > >}}
\newline
\newline
詩篇 8:1-9
\newline
\begin{longtable}{cl}
\hline
\hline
章節 & 經文 (和合本修訂版)\\
\hline
8:1 & \begin{tabularx}{0.7\textwidth}{X} 耶和華－我們的主啊，你的名在全地何其美！你將你的榮耀彰顯於天。 \end{tabularx} \\ \\ \relax
8:2 & \begin{tabularx}{0.7\textwidth}{X} 你因敵人的緣故，從孩童和吃奶的口中建立了能力，使仇敵和報仇的閉口無言。 \end{tabularx} \\ \\ \relax
8:3 & \begin{tabularx}{0.7\textwidth}{X} 我觀看你手指所造的天，並你所陳設的月亮星宿。 \end{tabularx} \\ \\ \relax
8:4 & \begin{tabularx}{0.7\textwidth}{X} 人算甚麼，你竟顧念他！世人算甚麼，你竟眷顧他！ \end{tabularx} \\ \\ \relax
8:5 & \begin{tabularx}{0.7\textwidth}{X} 你使他比神微小一點，賜他榮耀尊貴為冠冕。 \end{tabularx} \\ \\ \relax
8:6 & \begin{tabularx}{0.7\textwidth}{X} 你派他管理你手所造的，使萬物， \end{tabularx} \\ \\ \relax
8:7 & \begin{tabularx}{0.7\textwidth}{X} 就是一切的牛羊、田野的牲畜、空中的鳥、海裡的魚，凡游在水裡的， \end{tabularx} \\ \\ \relax
8:8 & \begin{tabularx}{0.7\textwidth}{X} 都服在他的腳下。 \end{tabularx} \\ \\ \relax
8:9 & \begin{tabularx}{0.7\textwidth}{X} 耶和華－我們的主啊，你的名在全地何其美！ \end{tabularx} \\ \\
[1ex]
\hline
\hline
\end{longtable}
$^{1}$Good morning everyone..
Well why don't we say hi to each other and give them a high five..
Alright..
Say welcome to our worship service..
Okay let's do that..
Let's do high five..
Let's stand up..
Let's do high five with at least three friends..
Okay..
Yeah go around the room..
Yeah give them a high five..
Yeah..
Give them a high five..
Alright..
Yeah..
Thanks Isaac..
Thanks for the effort..
Thank you Isaac..
Alright..
Good to see you all..
As you can tell we gather every Sunday..
We sing, we pray, we study the word and listen to sermon and we have to ask the question.
what worship is all about..
Now we say we come to worship and worship this word is derived from an old English word.
worship and the word has to do with value and honor and the value of a person of an.
object..
So as you can tell worship is about the value we ascribe to God and the honor he deserves..
And you can tell from Revelation 4 there's a verse here that talks about worship in heaven.
and it says you are worthy God you are worthy oh Lord our God to receive glory honor and.
power..
So we come together to worship and today we're going to take a look at Psalm 8 and.
it's about worship..
It's a Psalm about worship..
And the first point here is to appreciate the worth of God our Creator and verse 1 it.
says oh Lord our Lord how majestic is your name in all the earth and you have set your.
glory above the heavens..
And Psalm 8 begins with an exclamation shouting out loud you know the Psalmist David praising.
God here say oh Lord our Lord your name is great in all the earth and when you know when.
the Psalmist look at the heavens the sky he bursts out with awe and wonder and worshiping.
and praising God..

$^{41}$Now last week a couple days ago many people go out to the electoral or just you know somewhere.
in the supper to see what to see the northern lights right to see the northern lights..
Have you have you seen it?.
Anyone?.
Raise your hand..
You go out and try to see it..
Okay there's only one I mean what are you guys doing?.
Look at the screen and you know people people going out to see the northern lights even.
in Chicago we can see it and people when they see it they feel the wonder of creation the.
wonder of the universe and here is here are some of the pictures taken by my son Nathan.
right outside his house..
He just walked outside the house and just turned around and took some pictures and you.
know my goodness my goodness..
You know we need to appreciate the creation and appreciate the creator right?.
So I can tell you never done it before..
That's what's going on in your life..
We are in a position to observe the wonder of God and some years ago I read a book called.
the Purvish Planet..
Some years ago it's written by a scientist and he called the book the thesis..
Our planet is a Purvish planet and what he means is that he argues that our earth you.
know our planet earth is ideally and uniquely positioned in the entire cosmos in the Milky.
Way to observe the wonders of creation..
You see we are in a very unique position and God created in such a way that we can appreciate.
it and verse 2 it says from the lips of children and infants you have ordained praises or strength..
Now this verse is difficult to translate you can find different translations of verse 2.
but the idea is that when you think of babies the infants sometimes we more than people.
will find them you know babies are cute right?.
Baby are very very lovely and adorable but in the ancient world in the biblical world.
babies are powerless people they are defenseless..
And even today when you think of babies you know they are helpless and they need our cares.
but this psalm says God helps the helpless and God draw closer to those who trust in.
him and defends the humble against the enemies..
And this psalm continues to talk about our value our worth..
Verse 3 when I consider your heavens God your work your fingers the moon and the stars which.
you have set in place what is mankind that you are mindful of them?.
Human beings that you care for them..
Psalm 8 describes us each and every one of us as man..
And man in Hebrew is Enosh the word is a Hebrew word and what it means is mere mortal we are.
mortals and we passed away quickly..
And the whole theme here is that the psalm is trying to tell us we humans are fragile.

$^{81}$but the psalm marvels at our values in the sight of God..
And then you know I have to remind you I'm getting old too you know in the past few weeks.
I have been having a lower back pain and it's lingering for days and it really trouble me.
and I have to admit that what the psalm says you know I am getting old you know I don't.
look old but I am getting old and this psalm is you know is correct in terms of talking.
about we are weak but God is strength and God is worthy of our worship..
And that's why the psalm talks about the paradox of human beings on the one hand we are worthy.
but on the other hand we are fragile we are weak..
And then here it talks about God caring for humankind caring for him and the sense is.
the meaning is like taking attendance of a class you know God is caring for us and mindful.
of us it's like taking the roll call hey are you here?.
You know you have the attention of God..
And here the psalm continues to talk about to talk about our work..
Psalm 8 5 to 9 it says you have made humankind a little lower than the angels and crowned.
them with glory and honor..
You made them rulers over the works of your hand you put everything under their feet all.
flocks and herds the animals of the wild the birds in the sky the fish in the sea that.
all that swim in the paths of the sea..
Our Lord how majestic is your name in all the earth..
Now so many years ago not so many years ago maybe 30 some years ago in 1990 okay some.
of you were not born yet in 1990 you know NASA sent a machine called Voyager 1 into.
the deep space so it was in 1990 it was launched into deep space and the mission of the Voyager.
is to take pictures of the galaxy is to take picture of the deep space and when this Voyager.
1 left the earth about 3.7 billion miles away NASA told this machine to turn around and.
took a picture of the planet earth and the picture was taken on February 14th Valentine's.
Day 1990 as you can see our entire planet earth appear as a tiny speck this picture.
is a very famous picture it's called pale blue dot as you can see it can you find it.
the blue dot that's our earth..
That pale blue dot is where we are now..
Now this picture of pale blue dot reminds us how small and tiny we are..
There's also a picture of world view you know from atheistic world view and you see.
the pale blue dot in the vast universe and we have no inherent meaning..
But the bible the bible presents a totally different picture totally different world.
view..
Psalm 8 here tells us the God who made the universe still knows you the God who made.
the universe still loves you and cares for all of us and we are called to be the steward.
of the creation of God..
And human beings are not you are not an accident God made all of us each and every one of us..
See that's a totally different world view the bible presents to us God made each and.
every one of us a little bit lower than the angels and God give us the highest honor possible..

$^{121}$You know we ask ourselves what is our net worth sometimes right..
Okay I want you to take 15 seconds try to calculate your net worth..
You know how much money you make and how much money you have in the bank or in your investment.
account you know try to add them up all together how much net worth do you have..
Elon Musk has a total net worth at this point 467 billion and Bill Gates has a net worth.
right now 104 billion..
But the psalm says God says your worth is not measured never measured by money you are.
crowned with glory and honor from God..
And God made each and every one of us manager the ruler representatives of God the king.
the creator..
Psalm 8 here is actually a summary of Genesis 1..
And thank God and this is some of the take homes that we should get from Psalm 8 first.
is God is worthy of our worship and we need to develop some habits of worshipping God.
and the first thing I want to remind you is that before worship before we start at 1130.
we come early so that you can spend 30 seconds to pray and you can pray a simple prayer you.
know Lord I want to focus my attention on you and your name is majestic..
And during the worship service like here now we should continue to pray and with a prayerful.
attitude and pray to God you know God I come to see you and you know talk to me speak to.
me..
And then after the worship service we should spend a little bit of time like 15 30 seconds.
you know when we end our worship we should meditate you know what kind of things that.
God want me to see today..
And you are worthy of our praise you know we tell God you know we want to worship you.
and the second point is that you each and every one of you are worthy of the care of.
our God..
See God sees you wherever you are..
Sometimes we are too occupied with our own works..
Let's turn our attention to God fully and worship God together..
Let's pray..
Lord our Lord how majestic is your name in all the earth..
Thank you God for reminding us that though we are small you see us you love us you crown.
us with glory and honor..
Help us worship you with a focused heart and may our lives continue to declare your glory..
In Jesus name we pray..
Amen..
All right welcome to all of you and so blessed to see you all and remember next week is 11.
right 11 11 o'clock not here guys not here okay go to Farmer 11 o'clock be there on time.
okay otherwise the food will be all gone okay..
God bless..
(soft music).

$^{161}$[Music].
\newpage



\section{使徒行傳 17:1-34}
\label{sec:mJ4c0qfZcx0}
\textbf{在異地作見證(使徒行傳17\_1-34) - 陳崇基牧師}
\newline
\newline
連結: \href{https://youtube.com/watch?v=mJ4c0qfZcx0}{\texttt{https://youtube.com/watch?v=mJ4c0qfZcx0}} ~~~~ 語音日期: 2025-11-18
\newline
\newline
\hyperref[sec:lX_7jxrhzjs]{\small{< < < PREV SERMON < < <}}
~
\hyperref[sec:index_chronic]{\small{[返順時目]}}
~
\hyperref[sec:index_scriptual]{\small{[返順卷目]}}
~
\hyperref[sec:G1IwEMMFCQQ]{\small{> > > NEXT SERMON > > >}}
\newline
\newline
使徒行傳 17:1-34
\newline
\begin{longtable}{cl}
\hline
\hline
章節 & 經文 (和合本修訂版)\\
\hline
17:1 & \begin{tabularx}{0.7\textwidth}{X} 保羅和西拉經過暗妃坡里、亞波羅尼亞，來到帖撒羅尼迦，在那裡有猶太人的會堂。 \end{tabularx} \\ \\ \relax
17:2 & \begin{tabularx}{0.7\textwidth}{X} 保羅照他素常的規矩進去，一連三個安息日，根據聖經與他們辯論， \end{tabularx} \\ \\ \relax
17:3 & \begin{tabularx}{0.7\textwidth}{X} 講解和說明基督必須受害，從死人中復活；又說：「我所傳給你們的這位耶穌就是基督。」 \end{tabularx} \\ \\ \relax
17:4 & \begin{tabularx}{0.7\textwidth}{X} 他們中間有些人聽了勸，就跟從保羅和西拉，還有許多虔敬的希臘人，尊貴的婦女也不少。 \end{tabularx} \\ \\ \relax
17:5 & \begin{tabularx}{0.7\textwidth}{X} 但不信的猶太人心裡嫉妒，聚集了些市井流氓，搭夥成群，煽動全城的人闖進耶孫的家，要把保羅和西拉帶到民眾那裡。 \end{tabularx} \\ \\ \relax
17:6 & \begin{tabularx}{0.7\textwidth}{X} 那些人找不著他們，就把耶孫和幾個弟兄拉到地方官那裡，喊叫著：「這些攪亂天下的人也到這裡來了， \end{tabularx} \\ \\ \relax
17:7 & \begin{tabularx}{0.7\textwidth}{X} 耶孫竟收留他們。這些人都違背凱撒的命令，說另有一個王耶穌。」 \end{tabularx} \\ \\ \relax
17:8 & \begin{tabularx}{0.7\textwidth}{X} 眾人和地方官聽見這些話，就惶恐了， \end{tabularx} \\ \\ \relax
17:9 & \begin{tabularx}{0.7\textwidth}{X} 於是收了耶孫和其餘的人的保證金後，釋放了他們。 \end{tabularx} \\ \\ \relax
17:10 & \begin{tabularx}{0.7\textwidth}{X} 當夜，弟兄們立刻送保羅和西拉往庇哩亞去；二人到了，就進入猶太人的會堂。 \end{tabularx} \\ \\ \relax
17:11 & \begin{tabularx}{0.7\textwidth}{X} 這地方的猶太人比帖撒羅尼迦的人開明，熱心領受這道，天天查考聖經，要知道這道是否真實。 \end{tabularx} \\ \\ \relax
17:12 & \begin{tabularx}{0.7\textwidth}{X} 所以，他們中間有許多信了，又有希臘的尊貴婦人，男人也不少。 \end{tabularx} \\ \\ \relax
17:13 & \begin{tabularx}{0.7\textwidth}{X} 但帖撒羅尼迦的猶太人知道保羅又在庇哩亞傳神的道，就往那裡去，煽動挑撥群眾。 \end{tabularx} \\ \\ \relax
17:14 & \begin{tabularx}{0.7\textwidth}{X} 於是，弟兄們立刻送保羅到海邊去，西拉和提摩太卻仍留在庇哩亞。 \end{tabularx} \\ \\ \relax
17:15 & \begin{tabularx}{0.7\textwidth}{X} 護送保羅的人帶他到了雅典，他們領了保羅的命令，叫西拉和提摩太趕快到他那裡來，然後回去了。 \end{tabularx} \\ \\ \relax
17:16 & \begin{tabularx}{0.7\textwidth}{X} 保羅在雅典等候他們的時候，看見滿城都是偶像，就心裡非常難過。 \end{tabularx} \\ \\ \relax
17:17 & \begin{tabularx}{0.7\textwidth}{X} 於是他在會堂裡與猶太人和虔敬的人，以及每日在市場上所遇見的人辯論。 \end{tabularx} \\ \\ \relax
17:18 & \begin{tabularx}{0.7\textwidth}{X} 還有伊壁鳩魯和斯多亞兩派的哲學家也與他爭辯。有的說：「這胡言亂語的要說甚麼？」有的說：「他似乎是宣傳外邦鬼神的。」這是因保羅傳講耶穌與復活的福音。 \end{tabularx} \\ \\ \relax
17:19 & \begin{tabularx}{0.7\textwidth}{X} 他們就把他帶到亞略巴古，說：「你所講的這新學說，我們也可以知道嗎？ \end{tabularx} \\ \\ \relax
17:20 & \begin{tabularx}{0.7\textwidth}{X} 因為你有些奇怪的事傳到我們耳中，我們想知道這些事是甚麼意思。」 \end{tabularx} \\ \\ \relax
17:21 & \begin{tabularx}{0.7\textwidth}{X} 原來所有的雅典人和居住在那裡的外國人都無暇管別的事，只是談談或聽聽新聞。 \end{tabularx} \\ \\ \relax
17:22 & \begin{tabularx}{0.7\textwidth}{X} 保羅站在亞略巴古當中，說：「諸位雅典人！我看你們凡事很敬畏鬼神。 \end{tabularx} \\ \\ \relax
17:23 & \begin{tabularx}{0.7\textwidth}{X} 我到處走走的時候，仔細觀察你們所敬拜的，發現一座壇，上面寫著『獻給未識之神明』。你們所不認識而敬拜的，我現在向你們宣告： \end{tabularx} \\ \\ \relax
17:24 & \begin{tabularx}{0.7\textwidth}{X} 他是創造宇宙和其中萬物的神；他既是天地的主，就不住在人手所造的殿宇裡， \end{tabularx} \\ \\ \relax
17:25 & \begin{tabularx}{0.7\textwidth}{X} 也不用人手去服侍，好像缺少甚麼似的；自己倒將生命、氣息、萬物賜給萬人。 \end{tabularx} \\ \\ \relax
17:26 & \begin{tabularx}{0.7\textwidth}{X} 他從一人造出萬族，居住在全地面上，並且預先定準他們的年限和所住的疆界， \end{tabularx} \\ \\ \relax
17:27 & \begin{tabularx}{0.7\textwidth}{X} 為要使他們尋求神，或者可以揣摩而找到他，其實他離我們各人不遠。 \end{tabularx} \\ \\ \relax
17:28 & \begin{tabularx}{0.7\textwidth}{X} 我們生活、行動、存在都在於他。就如你們的詩人也有人說：『我們也是他所生的。』 \end{tabularx} \\ \\ \relax
17:29 & \begin{tabularx}{0.7\textwidth}{X} 既然我們是神所生的，就不應該以為神的神性像人用手藝和心思所雕刻的金、銀、石像一般。 \end{tabularx} \\ \\ \relax
17:30 & \begin{tabularx}{0.7\textwidth}{X} 世人蒙昧無知的時候，神並不追究，如今卻吩咐各處的人都要悔改。 \end{tabularx} \\ \\ \relax
17:31 & \begin{tabularx}{0.7\textwidth}{X} 因為他已經定了日子，要藉著他所設立的人按公義審判天下，並且使他從死人中復活，給萬人作可信的憑據。」 \end{tabularx} \\ \\ \relax
17:32 & \begin{tabularx}{0.7\textwidth}{X} 眾人聽見死人復活的話，就有人譏誚他；又有人說：「我們會再聽你講這事。」 \end{tabularx} \\ \\ \relax
17:33 & \begin{tabularx}{0.7\textwidth}{X} 於是保羅從他們當中出去了。 \end{tabularx} \\ \\ \relax
17:34 & \begin{tabularx}{0.7\textwidth}{X} 但有幾個人依附他，信了主，其中有亞略巴古的議員丟尼修，和一個名叫大馬哩的婦人，還有幾個與他們一起的人。 \end{tabularx} \\ \\
[1ex]
\hline
\hline
\end{longtable}
$^{1}$我們繼續看《史徒行傳》.
《史徒行傳》裡面.
我們講過記載了保羅三次宣教的旅程.
在《史徒行傳》裡面.
有很多故事.
就是傳福音的故事.
就是教會去到將福音廣傳.
教會在不同的地方.
來到被建立的時候.
這些不同大大小小的故事.
在不同的城市.
不同的城鎮.
遇見了一些經歷.
這些都集合在這本書.
叫做《史徒行傳》.
就是新約第五本書.
我們今天會看第十七章.
經文是第十七章.
這個是保羅的第二次宣教.
就是這個地圖所表示的.
大家看看第二次宣教發生什麼事.
你看到一字嗎.
我記得上次給大家看地圖.
看到大家一頭雲.
但是這次一二三四.
第一發生什麼事.
就是在安特拉教會.
他們差派了保羅和巴拿巴.
他們再一次上路.
然後去到不同的教會.
去做一些探訪.
然後傳福音.
但是因為馬可這個人.
他們兩個.
保羅和巴拿巴就有些爭執.
記不記得.
有些爭執.
然後他們就兵分兩路.
巴拿巴就帶著馬可.
跟著黃色橙色那條虛線.

$^{41}$就去了今天的塞浦路斯.
塞浦路斯或者居比路.
然後保羅帶著西拉和他的團隊成員.
他們就去今天的所謂的土耳其.
看到安提阿瑪在右邊.
他們跟著那些綠色的路線.
迂迴曲折的.
經過大數去到特比.
然後路斯德等等.
他們想南下去到伊弗所.
但是聖靈阻止.
但是接著他們就想北上.
想去上面另一個省加拉太.
還有比推尼.
然後聖靈又禁止他們.
就逼了他們走去哪裡.
南邊走不到.
北邊走不到.
就走去二那裡.
就是特羅亞這個地方.
就是今天土耳其西邊的一個海灣.
這個海灣是通往歐洲的.
所以逼了他們去了這個地方.
No.2.
去了特羅亞.
Taurus這個地方.
然後他們在那裡見到一個異象.
那個異象叫做馬其頓的異象.
記得嗎.
Macedonian Vision.
在那裡見到對岸有一個人.
馬其頓的人.
說你過來幫我好嗎.
於是他們就帶著他們的團隊成員.
他們覺得上帝就是帶領他們.
去到歐洲.
於是他們由特羅亞上船.
去到對面的是哪個地方.
經過不同的地方去了菲勒比.
菲勒比.

$^{81}$接著我們去到No.3.
菲勒比傳完福音之後.
那裡當然有Big Bar.
然後他們就移步去到另一個城市叫做提薩羅尼加.
在那裡亦有不同的紛爭.
因為傳福音的緣故.
於是他們又上路去到比利亞.
在那裡他們有一些故事說了.
然後他們就繼續向南.
去到希臘半島下面.
一個地方很重要的.
叫做雅典.
可以嗎 我們看到這裡.
1,2,3,4.
可以嗎.
這個是保羅的第幾次宣教旅程.
第二次.
OK.
在這裡我們看到保羅在比利亞.
應該是乘船去雅典的.
Ethnos.
在那裡他示範了另一個傳福音的範例.
如何去傳福音.
不單止我們需要有勇氣.
不單止我們需要有動力去傳福音.
亦都需要有智慧.
這裡在雅典保羅很能表揚.
表達到到底如何有智慧的.
來跟別人去傳福音去探討.
所以在經文第十七章.
我們有三點想留意的.
第一點就是他如何從心開始.
從他的生命開始.
從他的內心開始來傳福音.
先看經文.
第十七章第十六節.
保羅在雅典等候其他團隊成員來到時.
他看到滿城都是偶像.
心裡非常難過.
於是他在猶太人的會堂.

$^{121}$和猶太人和敬虔的人來傳福音.
以及每天在市場裡面.
會堂那些公眾的地方.
所遇見的人來辯論.
還有Epicurean和Stoic兩派的哲學家.
亦都和他和保羅一起爭辯.
有些哲學家和學者.
就會覺得保羅是胡言亂語的.
要說甚麼呢.
有的說他似乎在宣講外邦的鬼神祠.
為甚麼呢.
因為保羅是在跟他們說耶穌的福音.
以及復活的福音.
雅典是一個怎樣的地方呢.
Ethnos.
當然今天亦是一個Capital的地方.
一個很重要的地方.
在第一世紀.
當時保羅的時候.
雅典有少少像我們今天美國的.
不像美國的甚麼呢.
不像美國的Texas.
不像美國的芝加哥.
像美國的甚麼呢.
Boston.
為甚麼呢.
因為最大的學府.
最高的學府.
就在Boston.
世界在美國.
有些最出名的學府都在那裡.
雅典就是一個這樣的地方.
一個學術的地方.
一個最高學府所在的地方.
出名的就是甚麼.
出名的就是哲學家或思想家.
或者他們帶領著西方希臘羅馬的文化.
是一個出產學者的地方.
對於保羅來說.
這是一個外邦異地的世界.

$^{161}$當然保羅是一個有學之士.
他是大肆成長的一個人.
大肆剛才你看地圖的時候.
都是一個很重要的城市.
但他作為一個猶太人.
他跟外邦人不是有太多打交道.
但這個第二個宣教旅程.
就帶了保羅和他的Team Member.
去了一個外邦學術的世界.
就是雅典Evans這個地方.
是一個大都會.
然後他看到甚麼.
他看到滿城都是偶像.
這些偶像都代表不同的意識形態.
不同的ideology.
我想保羅是想到.
人們在尋找一些很終極的意義.
我們人和動物是不同的.
我們是meaning seeking being.
我們想找一些ultimate的meaning.
我們人生裡面最終的定義是甚麼.
為誰而活 為甚麼而活.
最終極的是甚麼.
保羅想感受到這件事.
他在雅典見到這麼多的意識形態.
和偶像的時候.
他心裡面就非常難過.
我應該可以這樣形容.
心裡面有一種掙扎.
然後這種掙扎會令到他.
不只是哀傷.
而是會令到他有一個動力.
他想做一些事.
你有沒有曾經試過有這樣的感覺.
你感受到一些事.
然後你覺得.
我想做一些事.
無論是事奉也好.
或者為主作見證也好等等.
是被他推動的.

$^{201}$非常難過就是這個意思.
這個值得留意的.
就是保羅見到整個雅典城.
滿有偶像.
很多不同的意識形態.
保羅沒有做甚麼.
保羅沒有那種潔癖.
大家有沒有潔癖.
潔癖.
我指的是屬靈的潔癖.
屬靈的潔癖.
即是很潔身自愛.
我在教會長大.
上一代的基督徒.
特別有一種這樣的氣味.
以往因為那種保守.
很守規矩的框架.
上一代的信徒.
譬如我的上一代.
很多時候說不可以看電影.
連走過電影院門口都不想.
當然有些電影不應該看.
他們有種潔癖.
就是覺得這些世俗的東西.
我們不鬥.
流行的東西.
即是pop culture的東西.
我們不鬥.
酒.
喝酒.
是滴酒不沾的.
有種這樣的動向.
這些可以了解.
因為當時那一代.
可能是這樣的模式.
來走出一個基督徒的生活.
這是上一代的背景.
但今時今日的背景是有點不同.
對吧.
怎樣去integrate.

$^{241}$怎樣去對話.
怎樣去engage.
怎樣去融入.
這些都是我們需要學的.
不可以單單是潔癖.
保羅沒有做的就是沒有潔癖.
第二樣他沒有做的是甚麼.
就是沒有站在道德的高地.
來批判人.
留意雅典的故事.
他完全沒有說.
你們是拜偶像的.
或者你們是不信的.
我高人一等.
你們是不行的.
他沒有這種氣味.
沒有這種語氣.
他心裡難過.
心裡難過.
他亦沒有hard sell.
沒有硬銷甚麼.
有一次我認識一個微信主的朋友.
他跟我說.
因為大家相識了.
說幾句小說之後.
他通常都是問甚麼.
除了問名字.
你從哪裡來.
來了多久.
是男是女都看到.
他問你做哪一行.
我沒有辦法.
我不可以騙他.
我不可以說我是彈結他的.
不可以的.
我職業不是這樣.
我跟他說牧師.
我一說牧師的時候.
他就說.
我認識兩個教會的人.

$^{281}$經常見到我的時候.
就跟我講福音.
我覺得他們很煩.
可能是他講福音的方法.
不是很對.
他講福音和傳福音是很對的.
是很對的.
應該鼓勵的.
但對他來說.
他覺得很煩.
經常要硬銷.
他的印象很不好.
我花了一段時間跟他解釋.
希望能夠放鬆肚子.
保羅在做甚麼.
有沒有一些confrontational.
直接傳福音的故事.
死亡主大白.
特別是猶太人會堂裡面.
因為那個背景是容許你這樣做.
但在保羅的雅典故事.
他卻沒有這樣做.
你留意他宣教的策略.
通常保羅去到不同的大城市.
不同的市鎮的時候.
他第一件事去哪裡.
就是猶太人的會堂.
在雅典.
他真的去了會堂.
跟那些猶太人.
和那些外邦皈信主的.
那些外邦皈信猶太教的.
所謂God-fear的人.
他先去會堂.
但在這裡有個宣教策略的改變.
就是甚麼.
以及每天在市場上.
和所遇見的人來辯論.
所以他兩件事一起做.
Parallel一起做.

$^{321}$跟他以往的做法有些不同.
他先去會堂.
然後搞好會堂的事情.
然後他才可能會找到一些.
皈信猶太教的那班外邦人等等.
這裡他的策略是不同的.
他走到公眾的地方.
機場,食堂,咖啡店等等.
他就去跟那些人溝通.
那些全部都是外邦人.
教會在今時今日.
我們遇見的背景.
都有少許像在雅典.
就是滿城偶像.
偶像不一定是雕出來的.
偶像正如我所說.
是代表我們一些意識形態.
代表我們一些追求.
或者膜拜一些東西.
沉迷一些東西等等.
無論你說的是金錢,成就等等.
但教會應該有甚麼情緒的反應.
不是潔癖.
也不是硬銷.
保羅是怎樣.
心裡面難過.
然後他說.
我想做些甚麼.
這是第一點.
第二點.
他從心開始之外.
他是以對話來做基礎的.
第十六節十七章.
保羅在雅典等候他們的時候.
他在做甚麼.
他與市場裡面的人辯論.
然後跟那些Stuart.
Epicurean的philosophers.
去到爭辯.
然後他們有不同的反應等等.

$^{361}$Epicurean.
是甚麼哲學的派別.
他們認為人生只是追求快樂.
簡單來說.
Stuart或者Stuartism.
其實就是人生裡面有很多苦難.
我們要堅忍等等.
就是其中兩個最大的學派.
跟他們辯論.
中文是辯論.
或者爭辯等等.
我們想起的印象就是一些很激烈的.
confrontational的debate.
一說到debate的時候.
我們就在說那些爭辯的情緒.
爭辯的語氣等等.
其實不一定的.
有一個現代的英文翻譯.
叫做message.
以前是Eugene Peterson做的.
90年代2000年左右的時候.
他完成了整本的message.
message就是一些paraphrase.
即是怎樣說呢.
即是意譯.
將事竟意譯.
他就譯成怎樣呢.
我覺得譯得不錯.
他說Paul discussed with.
the Jews and other like-minded people.
at the meeting place.
on the street.
he talked with anyone who happened along.
he got to know some of them.
pretty well.
through these conversations.
你留意到他用什麼動詞.
是忠誠的動詞.
discussion.
討論.

$^{401}$talk with.
跟他們談話.
然後conversation.
是對話.
你看見保羅在做什麼.
不是咬到臉紅耳熱.
很重要的一件事.
他就是進入他們的世界.
為什麼他進入他們的世界.
因為他觀察了一些東西.
你看他的講道你就知道了.
跟他們有些discussion之後.
22節.
保羅在做什麼呢.
保羅站在.
阿勒古巴.
Erico Pagus.
或者今天叫做Mars Hill.
這個地方.
在Athens有一個突出來的山.
突然間突出來.
不是上城那個.
是另外一個.
這個就叫做阿勒古巴.
阿勒古巴.
在那時.
有一個歷史悠久的法院.
就是一個法院.
或者一些.
我們今天叫做.
public discourse的地方.
就是一些公眾對話的地方.
公眾決策的地方.
他就在那裡被人帶去.
這個地方.
然後就跟那些雅典人來談話.
一個阿勒古巴.
阿勒巴古.
sorry.
阿勒巴古.

$^{441}$有一點像今天的.
譬如seminar.
一些學府裡面的講座.
TED talk等等.
當時因為雅典人很喜歡.
這些新潮的ideology.
保羅殺到去的時候.
他就跟他們說一些.
他們從來沒有聽過的東西.
於是就歡迎他.
在這個所謂TED talk那裡.
why don't you give a TED talk.
就給他15分鐘在那裡說.
保羅就在那裡說.
說什麼呢.
他說雅典人.
我見你們凡事都是敬畏鬼神的.
我到處走走的時候.
仔細觀察你們所敬拜的.
看見一座壇.
上面寫著.
unknown god.
獻給未識之神明的.
你們不認識敬拜的.
現在我向你們宣告.
在這裡做什麼呢.
保羅是進入他們的世界.
保羅是見到或者觀察到.
他們所渴望的是什麼.
他們想有的是什麼.
為什麼要那個未識之神都去拜呢.
你拜那個神明.
你當然要知道祂是什麼水才拜的.
但因為他們是所謂敬畏鬼神.
他們怕拜漏.
拜多個都不要緊.
那個是什麼.
百搭來的.
叫做unknown god.
總之有隻百搭.

$^{481}$我們一定吃糊.
一定吃糊.
未識之神.
所以保羅不是站在高地.
他不是說你們全部錯.
不是否定他們.
而是想明白想了解.
他們到底心裡面最渴望的是什麼.
正如我們所說.
我們人類每一個人.
我們都在心裡面有一個最底層.
最深的渴望.
而那個渴望就在福音裡面打開的.
在耶穌基督裡面打開的.
所以保羅就引用了他們所做的事.
或者所說的事.
觀察了他們有一個unknown god.
未識之神的祭壇.
他就說.
現在我想告訴你.
這一位是誰.
這是一個很有智慧的engagement.
是熟悉當時的文化.
熟悉對象.
然後建立了一些共通點.
其實今天我們說信仰.
我們說福音.
我們作見證.
其實都是一樣.
我們先做的那件事.
不是我們想說什麼.
當然我們想說福音.
我們先做的應該是.
我們應該找一個共通點.
What is the common ground.
到底是興趣.
到底是年紀.
背景等等.
因為你沒有那道門.
你是進不去的.

$^{521}$我們第一件事是需要有智慧.
來找福音的門.
第一件事我們要做的.
不是要爭取贏人.
不是要搏贏人.
而是要關係.
關係.
你留意我們人和人之間相處.
為什麼有些人可以.
好像認識了一陣子.
然後就很熟悉.
為什麼.
除了臭味相投之外.
就是他找到那道門.
大家找到那道門.
有那個connection.
有那種關係的建立.
和備案與接納的動態.
聆聽下人.
不是去到judge.
不是去到審判人.
那種聆聽和互動.
是建立共通點的.
common ground.
第一件事我們需要想的是一樣.
第一件事需要明白的是一樣.
保羅就是在做這件事.
然後保羅在他的講道的內容.
講什麼呢.
他就是以基督教的世界觀.
來開始介紹福音.
你留意他不是一些formula.
不是一些工程式來講福音的.
整段我節錄很多字.
24至31節.
是整段路加總結.
保羅在雅典裡面.
如何和雅典人傳福音.
留意一下有什麼主題.
他說神是創造宇宙其中萬物的上帝.

$^{561}$他是天地的主.
他不駐人手所做的奠禦.
也不需要人手去服侍.
好像自己缺少了什麼.
上帝自己都將生命氣息和萬物賜給萬人.
他從一人做出萬族.
居住在全地面上.
並且預先定準他們連限所住的疆界.
並要使他們尋求上帝.
或者可以揣摩而尋找到上帝.
其實上帝離我們國人都不遠的.
很近的.
因為我們的生活.
我們的行動.
我們的存在都在於他.
就如你們雅典的詩人.
他聽spotify.
他聽雅典的radio channel.
spotify radio channel.
就聽到他們有首詩歌.
詩人說什麼.
他說我們雅典人都是那位神明所生的.
他突然引用了spotify上最popular的那首歌.
既然我們是上帝所生的.
就不應該以為上帝的神聖.
好像人用手藝心思所雕造的金銀石像一般.
世人愚昧無知的時候.
上帝不追究.
但如今.
他吩咐各處各地的人都要悔改.
因為上帝已經定了日子.
藉著他所設立的人.
按公義審判世界.
並且使他.
就是耶穌基督.
從死裡面復活.
給萬人作可信的憑據.
這就是整個雅典.
所謂Marshall Sermon.
雅典的講道的精華.

$^{601}$有什麼主題呢.
他由創造開始的.
大家都知道我們這個世界.
沒理由自然而然.
應該有人設計.
在創造裡面.
如果是上帝的話.
他就是創造萬有.
並且他不住在人手的廟宇裡面.
我們不用弄東西讓他住在裡面.
這就是去矯正當時雅典人.
他們對上帝神明的誤解.
第二就是上帝是遺齒萬有的.
不需要人供應的.
因為外邦的神明就是要供應他的.
初一十五你就要做一隻白切雞給他.
然後大家有些交易.
他生氣就不要了.
最好哄他.
哄好了他.
他就會對你好.
你突然踩到他的腳趾.
得罪他.
這樣就不行了.
初一十五要做齊這些事情.
因為神明有他的需要.
我們有我們的需要.
所以大家就有一個合作.
你就借天借地了.
保羅說不是.
上帝是上帝的話.
上帝不用我們供應他的.
反而上帝去供應我們一切.
反而上帝給我們米飯.
反而上帝去供應萬物給萬人.
這就是在矯正一般人.
雅典人對於神明的錯誤的理解.
上帝不是這樣的.
並且他還說什麼.
上帝是安排歷史.

$^{641}$是歷史的主宰.
我們有的年歲.
我們有的行動.
搬去哪裡.
我們的movement.
我們的國界.
我們的民族.
都是上帝所設立的.
但是有一個目的.
上帝希望人去領略到.
生命裡面不是我們控制所有的東西.
如果我們領略到這件事的話.
我們就明白.
開始去明白.
開始去想.
咦 那不是我控制所有的東西.
那控制所有的東西的是誰呢.
Who is that.
就是為了.
上帝為了萬人去尋求他.
然後他告訴我們.
這位上帝不是遙不可及的.
是跟我們很接近.
跟我們很接近.
並且還有一個終極的目的.
就是將來的公義審判.
我們最終要結帳的.
最終要見上帝的.
有說到永生 審判等等.
也都提到耶穌的死而復活.
這個整個就是基督徒的聖經的世界觀.
但你留意整個孖超的講道.
重點居然是沒有什麼十字架的.
奇怪嗎.
不是說他不講.
之前他有講耶穌等等的東西.
但他路加的總結.
不是太多講十字架的.
但講的是保羅和雅典人之間的共同地位.
由世界的創造開始.

$^{681}$講到上帝維繫我們人生.
講到上帝生我們.
引用一些當時的歌曲或者歌詞.
現代的術語等等.
做一些橋樑.
所以見到這麼多個史萊林傳裡面的.
傳福音的故事.
這個是最落地的.
也是保羅去展現.
如何去建造橋樑.
不是燒掉橋樑.
是建造橋樑.
我們說在整年裡面我們需要有跨越.
我們需要學的.
有人講得很對.
我們基督徒是在世界裡.
但不是在世上.
我們是在世生活.
但不是屬於世界.
什麼叫在世界裡.
在世界裡面生活.
你要在世界裡面生活的話.
我們就一定要和其他人相似.
總有些地方我們是相似的.
元宗崇拜你去哪裡.
你吃東西.
人人都吃東西.
人家吃什麼我們都吃什麼.
有相似的地方.
在世界裡我們一定是相似於世界.
對不對.
基督徒的世界觀.
不應該是躲在我們自己有的空間.
那就算了.
那些叫Armish.
我們躲在那裡.
有我們自己的生活.
有我們自己的價值.
我們擁抱我們自己有的東西.
保護我們自己有的東西.

$^{721}$給我們一個殼.
基督徒不應該是這樣.
應該是參與進入這個世界.
所以時下的意識形態.
是我們需要去思考的.
思考去批判.
不是每樣東西我們去容納.
但我們需要去衡量.
時下的電影說什麼.
文化裡面說什麼.
什麼生活經驗.
什麼藝術等等.
這些是需要我們去接觸的.
以致我們能夠和別人去建立.
Bridges.
有橋樑.
這裡保羅給我們的榜樣.
展示了他怎樣去到.
令福音更加吸引.
上一代有一個很出名的護教學家.
神學家叫Francis Schaeffer.
Francis Schaeffer在歐洲建立了一個.
叫Club.
就是一個Conversational Club.
就是他邀請一些人.
特別是一些學生等等.
年輕的一輩.
他們能夠和他們去到那個地方.
在沒有壓力之下.
開放的環境下.
和他們有一些對話.
不是硬燒什麼.
Schaeffer最喜歡這件事.
他有一本書叫做.
"The God is there and God is not silent".
他說到什麼.
他說如果我有一個小時.
可以和一個未信主的人相處的話.
我會做什麼.
你想想你會做什麼.

$^{761}$一個小時.
如果別人給你一個小時.
你又可以給別人一個小時的話.
這個小時你想做什麼.
他說他會用頭55分鐘.
去到了解.
有什麼困擾人.
了解別人的心靈.
了解一下可能別人理性裡面.
有什麼問題.
然後最後那5分鐘.
分享.
Something of the truth.
真理裡面的一點.
這些功力很大.
因為不是急功近利.
但是肯去engage.
這個就是保羅所展現的.
其實我們每個人都在雅典.
今天我們生活就在雅典.
一個很複雜.
很多元的.
很多意識形態.
衝撞的一個世界.
我們可以做什麼呢.
保羅這裡給我們看到的.
他先去關心別人.
先去和別人聯絡.
然後先去分享福音.
先做到朋友.
並且別人想和你做朋友.
然後將耶穌解釋給他.
去理解別人.
去聆聽.
聆聽一個人的心裡面.
他在意的是什麼.
他怕的是什麼.
他渴望的是什麼.
第三.
他在敏銳聲明給他時間.

$^{801}$但願上帝去幫助我們.
今天的雅典.
今天的城市.
讓我們有智慧有愛的.
做主的見證.
我們同心協力祈禱.
上帝我們多謝你.
因為你進入我們的世界.
去聽我們的故事.
去愛我們.
去接納我們.
在我們不同的環境.
不同的際遇裡面.
你自己出現.
你去愛我們.
擁抱我們.
還有這個這麼美好的福音.
就是因為你給我們恩典.
我們能夠相信你.
能夠認識你.
求你教導我們.
怎樣去做一個更好的見證.
讓我們生活在家庭職場.
在我們街道上我們看到的人.
社會裡面我們所接觸的人.
主要求你讓我們能夠成為一個.
建立橋樑的生命.
使人能夠聽見你的福音.
看見我們的生命.
和我們交往的時候.
見到耶穌基督的美善.
祈禱給我們命.
願上帝祝福各位.
\newpage



\section{以弗所書 2:11-22}
\label{sec:G1IwEMMFCQQ}
\textbf{One in Christ (Eph 2\_11-22) - Paul Zhang 在基督裡合一 (以弗所書2\_11-22) - 張威傳道}
\newline
\newline
連結: \href{https://youtube.com/watch?v=G1IwEMMFCQQ}{\texttt{https://youtube.com/watch?v=G1IwEMMFCQQ}} ~~~~ 語音日期: 2025-12-01
\newline
\newline
\hyperref[sec:mJ4c0qfZcx0]{\small{< < < PREV SERMON < < <}}
~
\hyperref[sec:index_chronic]{\small{[返順時目]}}
~
\hyperref[sec:index_scriptual]{\small{[返順卷目]}}
~
\hyperref[sec:45ohI_TBLkg]{\small{> > > NEXT SERMON > > >}}
\newline
\newline
以弗所書 2:11-22
\newline
\begin{longtable}{cl}
\hline
\hline
章節 & 經文 (和合本修訂版)\\
\hline
2:11 & \begin{tabularx}{0.7\textwidth}{X} 所以，你們要記得：從前你們按肉體是外邦人，是「沒受割禮的」；這名字是那些憑人手在肉身上「受割禮的人」所取的。 \end{tabularx} \\ \\ \relax
2:12 & \begin{tabularx}{0.7\textwidth}{X} 要記得那時候，你們與基督無關，與以色列選民團體隔絕，在所應許的約上是局外人，而且在世上沒有指望，沒有神。 \end{tabularx} \\ \\ \relax
2:13 & \begin{tabularx}{0.7\textwidth}{X} 從前你們是遠離神的人，如今卻在基督耶穌裡，靠著他的血，已經得以親近了。 \end{tabularx} \\ \\ \relax
2:14 & \begin{tabularx}{0.7\textwidth}{X} 因為他自己是我們的和平，使雙方合而為一，拆毀了中間隔絕的牆，而且以自己的身體終止了冤仇， \end{tabularx} \\ \\ \relax
2:15 & \begin{tabularx}{0.7\textwidth}{X} 廢掉那記在律法上的規條，為要使兩方藉著自己造成一個新人，促成了和平； \end{tabularx} \\ \\ \relax
2:16 & \begin{tabularx}{0.7\textwidth}{X} 既在十字架上消滅了冤仇，就藉這十字架使雙方歸為一體，與神和好， \end{tabularx} \\ \\ \relax
2:17 & \begin{tabularx}{0.7\textwidth}{X} 並且來傳和平的福音給你們遠處的人，也傳和平給那些近處的人， \end{tabularx} \\ \\ \relax
2:18 & \begin{tabularx}{0.7\textwidth}{X} 因為我們雙方藉著他，在同一位聖靈裡得以進到父面前。 \end{tabularx} \\ \\ \relax
2:19 & \begin{tabularx}{0.7\textwidth}{X} 這樣，你們不再是外人或客旅，是與聖徒同國，是神家裡的人了， \end{tabularx} \\ \\ \relax
2:20 & \begin{tabularx}{0.7\textwidth}{X} 被建造在使徒和先知的根基上，而基督耶穌自己為房角石， \end{tabularx} \\ \\ \relax
2:21 & \begin{tabularx}{0.7\textwidth}{X} 靠著他整座房子連接得緊湊，漸漸成為在主裡的聖殿。 \end{tabularx} \\ \\ \relax
2:22 & \begin{tabularx}{0.7\textwidth}{X} 你們也靠他同被建造，成為神藉著聖靈居住的所在。 \end{tabularx} \\ \\
[1ex]
\hline
\hline
\end{longtable}
$^{1}$[SPEAKING CHINESE].
You can see two pictures here..
[SPEAKING CHINESE].
On the left is a lonely man..
[SPEAKING CHINESE].
On the right is a very happy gathering family..
[SPEAKING CHINESE].
And if you have a choice, which one would you pick?.
[SPEAKING CHINESE].
Of course, the right one..
[SPEAKING CHINESE].
We all desire to be part of the community..
[SPEAKING CHINESE].
And we don't want to be isolated..
[SPEAKING CHINESE].
And wherever you go, you want that community,.
the family to welcome you..
[SPEAKING CHINESE].
But in reality, sometimes it's different..
[SPEAKING CHINESE].
And sometimes we come to the United States,.
we experience different kinds of isolation..
[SPEAKING CHINESE].
And all kinds of differences make us foreigner,.
whether it's religion, skin color, or culture,.
or language..
[SPEAKING CHINESE].
And sometimes in the church, we feel the same way, too..
We are isolated from each other..
[SPEAKING CHINESE].
And then we want to talk to you about a passage that.
talks about isolation and separation..
And it's not just happened to all of us today,.
but it happened back then..
[SPEAKING CHINESE].
And we talk about the church in Ephesus..
They have different groups with hostility against each other..
[SPEAKING CHINESE].
And then this passage we're talking about,.
it talks about Christ is the only one.

$^{41}$who can break down the walls..
[SPEAKING CHINESE].
And that in Christ, Jesus formed a new family.
and new community..
[SPEAKING CHINESE].
And we want to see how Christ would break down.
all the walls here..
[SPEAKING CHINESE].
And the passage is found in Ephesians chapter 2, 11..
[SPEAKING CHINESE].
There are three different divisions here..
The first part is 11 to 12..
[SPEAKING CHINESE].
OK, 2, 11, it says, "Therefore, remember that formerly you.
are Gentiles by birth and called uncircumcised by those.
who call themselves the circumcision, which.
is done in the body by human hands..
Remember that at that time you were.
separated from Christ, excluded from citizenship in Israel.
and foreigners to the covenants of the promise without hope,.
without God in the world.".
[SPEAKING CHINESE].
And then the second part is 13 and following..
[SPEAKING CHINESE].
Verses 13 and following, "Now in Christ Jesus,.
you who were once were far away have been brought near.
by the blood of Christ..
And he himself is our peace, who has made the two groups one.
and has destroyed the barrier, the dividing wall of hostility.
by setting aside in his flesh the law with his commands.
and regulations..
His purpose was to create in himself one new humanity out.
of the two, thus making peace, and in one body.
to reconcile both of them to God through the cross by which he.
put to death the hostility..
He came and preached peace to all of you.
who were far away, peace to those who were near..
For through him, we both have access to the Father.
by one spirit.".
[SPEAKING CHINESE].

$^{81}$The third division of the passage.
is found in verses 19 and following..
"Consequently, so you are no longer foreigners and strangers,.
but fellow citizens with God's people.
and also members of his household,.
built on foundation of the apostles and prophets..
With Christ Jesus himself as the chief cornerstone,.
in him the whole building is joined together,.
rises to become a holy temple in the Lord..
In him, you too are being built together.
to become a dwelling in which God lives by the Spirit.".
[SPEAKING CHINESE].
And the background of this passage.
talks about the hostility between Jews and Gentiles.
in Ephesus..
[SPEAKING CHINESE].
You can find the map of the city of Ephesus..
[SPEAKING CHINESE].
It's located at the western side of Turkey..
[SPEAKING CHINESE].
Ephesus was a metropolitan city in the ancient world..
[SPEAKING CHINESE].
And on the right, you can see a humongous theater,.
and you can tell how big the city was..
[SPEAKING CHINESE].
Well, not only it has great wealth,.
but it also has the famous Artemis Temple..
[SPEAKING CHINESE].
And you can see the reconstruct diagram on the right..
It has 120 marble pillars to support the temple..
[SPEAKING CHINESE].
And it was a religious center of worshiping idols..
And once in every year, the crowd.
would come from all nations to worship in festivals..
[SPEAKING CHINESE].
And then there were prostitutes in the temple,.
and all the visitors would come and have.
sex with the temple prostitutes..
[SPEAKING CHINESE].
And from these kind of activities,.

$^{121}$they thought they could guarantee fertility.
by doing all these things..
[SPEAKING CHINESE].
So this was the center of a pilgrimage,.
a religious idol pilgrimage..
And that's why Ephesus was a center of economy, finance,.
and religion..
[SPEAKING CHINESE].
And compared to this, even Chicago Pride Parade.
would not match this..
[SPEAKING CHINESE].
So as Chicago and New York and all the major cities,.
Ephesus was a very diverse culturally..
[SPEAKING CHINESE].
So Jews and Gentiles, they lived side by side,.
and they don't like each other..
[SPEAKING CHINESE].
And then for three years, Paul was.
pastoring the church in Ephesus and trying.
to break down the walls with the truth of God..
[SPEAKING CHINESE].
And then the Ephesus Church will become.
a model of unity in Christ..
[SPEAKING CHINESE].
And then we want to talk about what.
are the sources of hostility between the two groups..
[SPEAKING CHINESE].
Paul begins by commanding the people to remember,.
and he wants them to recall their past..
[SPEAKING CHINESE].
Then from the perspective of the Jews,.
the Gentiles or the foreigners were being called uncircumcised..
[SPEAKING CHINESE].
So one of the sources of hostility.
is the uncircumcision of the Gentiles..
[SPEAKING CHINESE].
And from the Old Testament, you can see how important.
circumcision was, because it was our sign of God's covenant.
to his people, beginning with Abraham..
[SPEAKING CHINESE].

$^{161}$And you can see the Bible tells us.
that circumcision is actually a symbol.
to the spiritual reality inside..
But the Jews often turn it into a source of pride..
[SPEAKING CHINESE].
They insult the Gentiles by all sorts of different racial.
insults and slurs..
[SPEAKING CHINESE].
And there were so many Roman bathrooms, which was perfect..
And then when they see the Gentiles,.
they would despise them..
[SPEAKING CHINESE].
So this kind of hostility is deeper,.
is more intense than whatever racial divide.
that we have today..
[SPEAKING CHINESE].
So there are so many regulations among the Jews,.
and they are not willing to help Gentile women..
Even in their childbirth, they are not willing to help them..
[SPEAKING CHINESE].
And Paul was saying, you guys were being circumcised,.
but it was done by human hands..
[SPEAKING CHINESE].
And Paul was a circumcised Jew as well..
[SPEAKING CHINESE].
And that's why Paul, speaking from a circumcised Jew,.
he said circumcision is merely done by a human hand..
[SPEAKING CHINESE].
And then he was speaking to the Gentiles.
and talking to them about the fact.
that they were outside of Israel,.
outside of the covenant..
[SPEAKING CHINESE].
So you go back to the Old Testament..
God made a covenant with the Jews, with his people,.
and starting from Abraham..
And although there were a few foreigners being included.
into the nation of Israel, but however, primarily, it.
was for the race, the Jews..
[SPEAKING CHINESE].

$^{201}$So the Gentiles were the foreigners, outsiders.
of the covenant..
[SPEAKING CHINESE].
And then all the promises belong to the Jews,.
and even the Messiah is prophesied and is for the Jews..
[SPEAKING CHINESE].
So that's why he says the Gentiles,.
you have no part of this whole thing..
[SPEAKING CHINESE].
And the Gentiles, without Christ, they have no hope..
They have no covenant..
[SPEAKING CHINESE].
And you remember, we talked about Ephesus.
was a very rich city and full of wealth and temples.
and religion..
[SPEAKING CHINESE].
And then all the citizens in Ephesus, the Gentiles,.
they were having nightlife and all the excitements.
in the worldly life..
[SPEAKING CHINESE].
And then many people would admire them.
and want to be like them..
[SPEAKING CHINESE].
And Paul is talking to the Gentiles in Ephesus.
that even though you have all of these things,.
but you are without hope in God..
[SPEAKING CHINESE].
So Paul is reminding them they are totally without God,.
without Christ, and without hope..
And their life may look very wealthy,.
but their life is empty..
And they are without the grace of God.
and don't part in God's program..
[SPEAKING CHINESE].
And one of the applications we can find in this passage.
is that Paul's description for the Gentiles.
remind us of our own past..
[SPEAKING CHINESE].
And we were foreigners..
[SPEAKING CHINESE].

$^{241}$And we were without Christ, without God, without hope..
[SPEAKING CHINESE].
And then to many people, maybe we.
are having wealth and all these things in life,.
but in the sight of God, we are without hope..
[SPEAKING CHINESE].
And we can see that knowing God today, knowing Christ today,.
is entirely His grace..
[SPEAKING CHINESE].
And then we can see the human nature in this passage,.
because we humans, we like to create divisions..
We like to focus on our differences.
among different people..
But God's design is not like this..
[SPEAKING CHINESE].
And we often use our differences to reject one another..
And that is actually pride, believing our group is.
superior to other groups..
And the Jews thought circumcision.
made them holier so that they despised the Gentiles..
[SPEAKING CHINESE].
And we have to ask today whether we are building.
this sort of invisible walls..
[SPEAKING CHINESE].
And many differences among us in languages, in culture,.
in occupation..
Our church as well, we have English-speaking,.
Cantonese-speaking, Mandarin-speaking people,.
people from China, people from America, people from Hong Kong..
[SPEAKING CHINESE].
Now I heard a story from a pastor that I know in Canada..
He said his church is filled with two groups..
And one group is the immigrants, investor immigrants..
The other group is professional immigrants..
And they don't like each other, and they.
look down on each other..
[SPEAKING CHINESE].
The professional immigrants will look down on the investors,.
thinking that they only have money,.
and they don't have knowledge, and they don't have.

$^{281}$the prestige of the society..
On the other hand, the investors will look down.
on the professionals and think that they are simply.
working class people..
[SPEAKING CHINESE].
And the church should not be like this..
[SPEAKING CHINESE].
Paul is trying to contrast the past and the present..
You were separated from God, but now you are different..
Now you're in Christ..
[SPEAKING CHINESE].
And Paul was saying, you guys were far away from God,.
but now you are being brought near by the thought of Christ..
[SPEAKING CHINESE].
And in the past, you were separated from God,.
and now you are connected to God..
And that is one of the foundational change.
that we have in Christian life..
[SPEAKING CHINESE].
And that is the core value of Christianity..
You were separated from God, but now you are near to God..
[SPEAKING CHINESE].
And the core value, core essence of Christianity.
is not about behavioral value change,.
but it's about our life being changed in Christ..
[SPEAKING CHINESE].
And it's not because I was bad before,.
and now I want to improve myself, be better..
[SPEAKING CHINESE].
And it is the restoration of relationship with God..
[SPEAKING CHINESE].
And there's one professor who was a philosopher,.
and he used to teach in Chinese University in Hong Kong..
He said Christianity is not about becoming a better person..
It is about being reconciled to God..
[SPEAKING CHINESE].
And the restoration of relationship with God.
is not about our effort..
It's about the grace of God..
[SPEAKING CHINESE].

$^{321}$And in Ephesus, when they worshiped the pagan god,.
that involved a lot of sacrifice and killing a lot of animals..
And some of the cults even practiced infant sacrifice..
[SPEAKING CHINESE].
And you can tell some of the pagans in the Ephesus church,.
they have that kind of background..
And when they heard about blood sacrifice,.
they immediately understand..
But Paul is trying to say to them, all this blood.
that you have shed before had no value.
and had no power to save you..
Only the blood of Jesus Christ can truly save you..
[SPEAKING CHINESE].
And then what did Jesus do?.
And he destroyed the dividing wall.
between Jews and Gentiles..
And we can tell from this passage.
what exactly Jesus has done..
[SPEAKING CHINESE].
And when Paul talks about Jesus breaking down the wall,.
he's actually referring to a physical wall that.
separates the court of the Jews and the court of the Gentiles.
in the temple of Jerusalem..
[SPEAKING CHINESE].
And that wall is actually the dividing separation.
between the Jews and Gentiles..
And Gentiles cannot climb over the wall to go into the court..
[SPEAKING CHINESE].
And we have found one of these stone inscriptions.
that is being set up on the wall..
It's a warning sign to all the Gentiles..
And the content says, if you're Gentiles,.
if you climb over this wall, you will bring yourself death.
penalty..
[SPEAKING CHINESE].
So Paul is saying, Jesus came to break down this wall..
Exactly, this wall is trying to break down so.
that Jews and Gentiles are one..
[SPEAKING CHINESE].
And when he talks about the regulations and the commandments,.

$^{361}$he talks about the commandments received by the Jews from God..
[SPEAKING CHINESE].
And because Jesus came, the Jews and even the Gentiles,.
we do not depend on the sacrifice from the animal,.
because Jesus is the perfect sacrifice..
[SPEAKING CHINESE].
Now, the source of the law of Moses.
is actually from the grace of God..
And God has given the law so that the people.
would behave in such a way that would match their identity..
However, the irony is that sometimes the Jews and people.
would focus on, because we are keeping the law,.
that's why we are better than other people..
[SPEAKING CHINESE].
And just like the Pharisees, they.
were trying to accuse Jesus of not keeping the law,.
and not keeping the law just like them..
That's why they accused Jesus and tried to kill him..
[SPEAKING CHINESE].
And that's why the law, Paul says,.
the Jews and the Gentiles divide different groups of people..
And Jesus came to remove the barrier,.
because Jesus is the perfect keeper of the law..
[SPEAKING CHINESE].
And Christ fulfilled the law in every single level,.
removing the law as a barrier between different people..
[SPEAKING CHINESE].
And that in Christ, in Jesus, we are created as one new humanity..
Jews and Gentiles are all one, united in Christ..
[SPEAKING CHINESE].
So the Bible is telling us the cross of Jesus.
is the foundation of true peace..
[SPEAKING CHINESE].
And the true unity between different kinds of people.
begins and founded on the peace of Christ..
[SPEAKING CHINESE].
So if there was no peace of cross,.
there was no harmony between different kinds of people..
[SPEAKING CHINESE].
So Paul is trying to emphasize the fact.

$^{401}$that Jews or Gentiles, whether you.
belong to the first group or the second group,.
you have to come to Christ, come to God through Jesus Christ..
[SPEAKING CHINESE].
And if we want to resolve conflicts.
between different people, we must come to the cross.
and meet Jesus at the foot of the cross,.
because we are all sinners and we need the grace of God..
[SPEAKING CHINESE].
And if you want to make peace with yourself,.
you have to find the true peace in Jesus first..
[SPEAKING CHINESE].
And sometimes we cannot make peace with ourselves.
because we cannot forgive ourselves,.
but we find true forgiveness only in Jesus Christ..
[SPEAKING CHINESE].
And then this passage tells us the three persons of the Trinity,.
the same God, all participate in the reconciliation of mankind.
relationship with God..
[SPEAKING CHINESE].
And Paul says, through the Spirit, we come to the Father..
[SPEAKING CHINESE].
And the Son redeems us and the Spirit makes us one..
[SPEAKING CHINESE].
So we talk about Father's love and the redemption of Christ,.
and we also emphasize the Spirit's work in our lives..
[SPEAKING CHINESE].
So Christianity is not just doctrine, not just knowledge,.
but it involves power, life-changing power.
in the Spirit..
[SPEAKING CHINESE].
So the doctrines that we believe in.
is founded in the historic event of Jesus Christ..
And without what Jesus has done, our doctrine will be empty..
[SPEAKING CHINESE].
But if we just have doctrine and theology without the Spirit,.
it becomes dead faith, dead orthodoxy..
[SPEAKING CHINESE].
But if we just pursue the Spirit without the truth,.
it will become unhealthy emotionalism in our faith..

$^{441}$[SPEAKING CHINESE].
And we need both truth and the Spirit..
And one of the main ways is to rely on the Spirit.
through our prayer..
[SPEAKING CHINESE].
And we can pray anytime..
And the irony is that sometimes we don't pray..
[SPEAKING CHINESE].
And sometimes we find it easy to attend worship,.
to give an offering, but praying sometimes.
becomes difficult for all of us..
[SPEAKING CHINESE].
We can talk to our friend and have.
chat with our friend for hours..
But sometimes the moment comes we have to pray..
Sometimes our minds will wander..
[SPEAKING CHINESE].
And different things will occupy our mind..
And we cannot be silent..
We cannot be attentive before God and pray..
[SPEAKING CHINESE].
And after a few minutes, we give up praying..
[SPEAKING CHINESE].
So we need the Spirit to help us to pray..
[SPEAKING CHINESE].
So the third point here--.
[SPEAKING CHINESE].
So you were separated from Christ,.
but Jesus came to break down all this wall..
And you are now near to Christ..
[SPEAKING CHINESE].
So the different identity changes here..
From you were foreigners, and now you.
are fellow citizens in the kingdom..
You were strangers to God..
Now you are members of the family, the household of God..
You were scattered stone and separated, isolated..
But now you have been built together to become a temple..
[SPEAKING CHINESE].
And then before, we did not have the heavenly identity..

$^{481}$Now we do have the heavenly identity in the kingdom..
[SPEAKING CHINESE].
In the past, we wandered without hope..
But now we are part of the family,.
part of the God's household..
[SPEAKING CHINESE].
And we are all part of the members.
being built together to build the dwelling of God,.
the temple of God..
[SPEAKING CHINESE].
So the church is actually God's temple in construction..
[SPEAKING CHINESE].
And the church refers not to the physical building here,.
but the people of God themselves..
[SPEAKING CHINESE].
Of course, our church is a visible church,.
a manifestation of the people of God..
But there is also a capital C church,.
which is so much bigger and involves all the saints.
and all the Christians from history.
and from different parts of the world..
[SPEAKING CHINESE].
So the capital C church is the invisible church, which is real.
and involves all the Christians from all the times..
And it goes beyond geographical boundaries..
[SPEAKING CHINESE].
So it says the church has a cornerstone in Christ..
[SPEAKING CHINESE].
And I find this diagram..
It shows you the cornerstone, where it is..
[SPEAKING CHINESE].
And the foundation is the apostles and the prophets..
It's on the wall here..
[SPEAKING CHINESE].
So the cornerstone, you see, is Jesus Christ..
And the wall connected to the cornerstone.
is the prophet and the apostles..
[SPEAKING CHINESE].
And that's why all the Christians from the entire world,.
we are part of this house..

$^{521}$[SPEAKING CHINESE].
And the importance of the cornerstone, because it.
dictates the plan and the structure.
of the entire building..
And that is Jesus Christ himself..
[SPEAKING CHINESE].
And church is not built by a man..
It's not just a social organization..
[SPEAKING CHINESE].
And if the church is simply a social organization,.
it would not survive for so many years..
But you can tell from church history,.
the church survives all the generations..
[SPEAKING CHINESE].
And the church is founded on Jesus Christ, the cornerstone,.
and Christ himself dictates and commands the church..
[SPEAKING CHINESE].
And you can see, even our church, sometimes we decline..
And our church can fall apart..
But we have to have that perspective of the capital C.
Church, and Christ is commanding his church,.
and all of the world, and all through history..
[SPEAKING CHINESE].
And every member of the church, every Christian,.
is one of the living stones being built together.
to form the temple for God..
[SPEAKING CHINESE].
So in this diagram, we are part of the different stones.
and bricks to build this whole household..
[SPEAKING CHINESE].
And you can see the church will continue.
to grow no matter what..
If you include the Old Testament,.
it has even longer history..
And the church will continue to grow..
[SPEAKING CHINESE].
And the church will continue to grow in the size.
that's beyond our imagination..
[SPEAKING CHINESE].
And for now, church history seems.

$^{561}$to be mixed with the world history..
But it seems like it's just one of the organizations.
in the world..
[SPEAKING CHINESE].
But you have to remember the church is in construction..
And you can tell, if an unfinished building,.
you can see a lot of different things here,.
different unconstructed part, and needs to be renovated..
[SPEAKING CHINESE].
But this building will continue to grow.
and continue to be constructed by Christ..
And someday, this church will be complete..
[SPEAKING CHINESE].
And then all these tools and unfinished part will be gone..
And the wall will be gone..
But the church remains..
[SPEAKING CHINESE].
And it seems that the church is growing.
with all kinds of organizations in the world, all the bodies.
in the world..
But someday, when God completes His church,.
everything will be unnecessary and will be removed..
[SPEAKING CHINESE].
And you can see we are the church..
[SPEAKING CHINESE].
The Spirit of God lives among us..
[SPEAKING CHINESE].
You think about it in the Old Testament,.
when the people of God want to see God and worship God,.
they have to travel a long distance.
to come to the temple to worship God together..
[SPEAKING CHINESE].
But now Christ has come..
All the people of God lives in Christ Himself..
[SPEAKING CHINESE].
So each and every one of us, the Spirit lives among us..
[SPEAKING CHINESE].
So when we love each other and we care for each other,.
the light of God will be demonstrated..
[SPEAKING CHINESE].

$^{601}$And if we think about it, we need.
to ask the question whether we are still.
having the mentality of foreigners or outsiders..
We don't find the sense of belonging here..
[SPEAKING CHINESE].
But you have to remember, even though you don't find acceptance.
in the community, you have to remember.
that you are fully accepted in Christ Jesus..
[SPEAKING CHINESE].
And when Christ sees you, it's not because of your skin color,.
it's not because of your background,.
but sees you just like He sees you in Christ, in Jesus..
[SPEAKING CHINESE].
So your identity now you find in Christ.
is that you are part of the family of God.
and part of the people of God and near to God..
[SPEAKING CHINESE].
And that's why we need to pray and we.
need to dwell in the presence of God, to draw near to Him..
[SPEAKING CHINESE].
Let's all pray together..
[SPEAKING CHINESE].
Heavenly Father, we want to praise you and give you thanks..
[SPEAKING CHINESE].
You sent Jesus, your son, to be the messenger of peace..
[SPEAKING CHINESE].
So that we can restore our relationship with you..
[SPEAKING CHINESE].
And we are no longer outsiders and we.
can come become your family and call you Abba Father..
[SPEAKING CHINESE].
Our Savior, Jesus Christ, He broke down all the barriers.
with His blood..
[SPEAKING CHINESE].
And we are all different in front of you..
[SPEAKING CHINESE].
And I just pray that you will help us to be united in you..
[SPEAKING CHINESE].
And help us not forget that we are all sinners.
and we are not better than other people..

$^{641}$We are all saved by grace..
[SPEAKING CHINESE].
And please make us peacemakers as well..
[SPEAKING CHINESE].
So that we can bring the message of reconciliation.
to other people..
[SPEAKING CHINESE].
In Jesus' name we pray..
Amen..
Bye..
[MUSIC PLAYING].
\newpage



\section{路加福音 4:16-30}
\label{sec:45ohI_TBLkg}
\textbf{逆轉的人生(路加福音4\_16-30) - 梁紫健傳道}
\newline
\newline
連結: \href{https://youtube.com/watch?v=45ohI-TBLkg}{\texttt{https://youtube.com/watch?v=45ohI-TBLkg}} ~~~~ 語音日期: 2025-12-01
\newline
\newline
\hyperref[sec:G1IwEMMFCQQ]{\small{< < < PREV SERMON < < <}}
~
\hyperref[sec:index_chronic]{\small{[返順時目]}}
~
\hyperref[sec:index_scriptual]{\small{[返順卷目]}}
~
\hyperref[sec:VcgHfBIcHMQ]{\small{> > > NEXT SERMON > > >}}
\newline
\newline
路加福音 4:16-30
\newline
\begin{longtable}{cl}
\hline
\hline
章節 & 經文 (和合本修訂版)\\
\hline
4:16 & \begin{tabularx}{0.7\textwidth}{X} 耶穌來到拿撒勒，就是他長大的地方。在安息日，照他素常的規矩進了會堂，站起來要念聖經。 \end{tabularx} \\ \\ \relax
4:17 & \begin{tabularx}{0.7\textwidth}{X} 有人把以賽亞先知的書交給他，他就打開，找到一處寫著： \end{tabularx} \\ \\ \relax
4:18 & \begin{tabularx}{0.7\textwidth}{X} 「主的靈在我身上，因為他用膏膏我，叫我傳福音給貧窮的人；差遣我宣告：被擄的得釋放，失明的得看見，受壓迫的得自由， \end{tabularx} \\ \\ \relax
4:19 & \begin{tabularx}{0.7\textwidth}{X} 宣告神悅納人的禧年。」 \end{tabularx} \\ \\ \relax
4:20 & \begin{tabularx}{0.7\textwidth}{X} 於是他把書捲起來，交還給管理人，就坐下。會堂裡的人都定睛看他。 \end{tabularx} \\ \\ \relax
4:21 & \begin{tabularx}{0.7\textwidth}{X} 耶穌對他們說：「你們聽見的這段經文，今天已經應驗了。」 \end{tabularx} \\ \\ \relax
4:22 & \begin{tabularx}{0.7\textwidth}{X} 眾人都稱讚他，並對他口中所出的恩言感到驚訝；他們說：「這不是約瑟的兒子嗎？」 \end{tabularx} \\ \\ \relax
4:23 & \begin{tabularx}{0.7\textwidth}{X} 耶穌對他們說：「你們一定會用這俗語向我說：『醫生，你醫治自己吧！我們聽見你在迦百農所做的事，也該在你自己的家鄉做吧。』」 \end{tabularx} \\ \\ \relax
4:24 & \begin{tabularx}{0.7\textwidth}{X} 他又說：「我實在告訴你們，沒有先知在自己家鄉被人接納的。 \end{tabularx} \\ \\ \relax
4:25 & \begin{tabularx}{0.7\textwidth}{X} 我對你們說實話，在以利亞的時候，天閉塞了三年六個月，遍地有大饑荒，那時，以色列中有許多寡婦， \end{tabularx} \\ \\ \relax
4:26 & \begin{tabularx}{0.7\textwidth}{X} 以利亞並沒有奉差往她們中任何一個人那裡去，只奉差往西頓的撒勒法一個寡婦那裡去。 \end{tabularx} \\ \\ \relax
4:27 & \begin{tabularx}{0.7\textwidth}{X} 在以利沙先知的時候，以色列中有許多痲瘋病人，但除了敘利亞的乃縵，沒有一個得潔淨的。」 \end{tabularx} \\ \\ \relax
4:28 & \begin{tabularx}{0.7\textwidth}{X} 會堂裡的人聽見這些話，都怒氣填胸， \end{tabularx} \\ \\ \relax
4:29 & \begin{tabularx}{0.7\textwidth}{X} 就起來趕他出城。他們的城造在山上；他們帶他到山崖，要把他推下去。 \end{tabularx} \\ \\ \relax
4:30 & \begin{tabularx}{0.7\textwidth}{X} 他卻從他們中間穿過去，走了。 \end{tabularx} \\ \\
[1ex]
\hline
\hline
\end{longtable}
$^{1}$(第1節).
弟兄姊妹平安.
希望大家在這麼大雪紛飛.
在寒冬裡,大家都真的平安.
我覺得這星期我們真的聽到很多.
很不好的消息.
很多令人心痛的消息.
同時,我們也知道有弟兄姊妹在前線.
他們做了很多在冬天送暖的行動.
今天我突然回來時,有個弟兄跟我說.
他說他在鏟雪時,鏟了旁邊的左右.
其實我覺得在這麼寒冷的時間.
基督徒做這些事很好.
因為你可以將你生命中的東西.
你的愛,帶給身邊的人.
因為我覺得這件事.
對我來說也很能溫暖我的心.
不知道…有點暗.
大家認出這個人嗎.
有些人認得出.
在一開始的五分鐘.
希望大家以第一個人稱.
進入一個人的生命.
希望我們用這短短的五分鐘.
去經歷一下.
一個人的生命被逆轉,反轉.
又被人重新建立的一個故事.
如果準備好.
我們一起來代入這個角色.
時間回到40年前.
當我30歲的時候.
我已經建立了自己的雙魚王國.
我跟朋友一起在車庫裡.
開設了一間公司.
10年內,這間公司成長得很快.
由數千元成長到年收入20億美元.
有4000個員工的企業.
我們剛剛推出了一個偉大的傑作.
這部「Map」.
對,我就是Steve John.

$^{41}$那一刻,我擁有了一切.
我是年輕企業家的代表.
全世界都看我頭.
我相信我的人生很成功.
正當我以為人生一帆風順.
風生水起的時候.
突然發生一件事.
我竟然被蘋果解僱了.
我被我親手創立的公司解僱了.
事情是這樣的.
當初我邀請了一個人.
回來當蘋果的行政總裁.
我覺得他可以幫助蘋果成長.
不過時間越久.
我跟他之間的衝突越來越深.
我想要這件事,他想要那件事.
我們的員工開始分黨.
有些人站在他那邊.
有些人站在我那邊.
最後董事會是這樣說的.
「要麼你走,要麼我走」.
在1985年,董事會決定了.
我一定要走.
當時的我….
幫我按….
上一張,對.
拿著紙皮箱.
在蘋果總部的停車場.
眼淚很自然地流出來.
我心想,這個世界玩完了.
我玩完了.
我是一個失敗者.
我辜負了所有人.
他交棒給我,我反而掉在地上.
我甚至想逃離矽谷.
我想逃離這一切.
令我失敗的一切.
於是我寫了一封信.
這封辭職信裡是這樣說的.
「我感到憤怒和悲哀」.

$^{81}$「因為蘋果做這些事一點好處都沒有」.
「我只是30歲,我很想繼續創作」.
「我生命裡還有一部很大的傑作」.
「準備推向世界」.
但蘋果沒有給我這個機會.
當時的我看不到任何希望.
不過在黑暗裡,有件事發生了.
我問自己一個問題.
雖然好像一切都玩完了.
我最喜歡甚麼.
當時我領悟到蘋果公司發生的事.
無法改變我對創作的熱誠.
我被人出賣.
但我依然很喜歡自己做的事.
所以我決定重新開始.
接下來的10年.
發生了一件難以置信的事.
我收購了一間公司叫Nix.
這間公司不是很成功.
但它教會了我很多東西.
我買了一間公司的部門,叫Pixar.
當時沒有人相信Pixar會成功.
沒有人.
甚至我自己也不太肯定.
但我堅持下去.
製作了第一部電影,叫「反斗奇兵」.
當時電影在銀幕上演的時候.
所有人都很驚艷.
這部電影不單是一部電影.
這部電影是未來.
Pixar成功了.
我在Pixar的股票率.
日夜之間升了很多倍.
總值12億美元.
比我在蘋果中損失的一切多.
1997年,蘋果面臨破產.
董事會邀請我回來.
我再次成為蘋果的行政總裁.
你想想,一個人被人出賣.
10年後再一次以拯救者的身分.

$^{121}$回到公司,再創高峰.
我以為我人生已經熬過所有苦難.
我人生會一帆風順.
2003年,當我五十多歲的時候.
正處於人生最成功的時刻.
而蘋果也在最高峰的時間.
推出了iMac.
再加上iPod.
已經改變了整個音樂產業.
正當我以為人生風生水起的時候.
我竟然被診斷出有癌症.
醫生跟我說.
「你大概有三至六個月的時間」.
「你可以想像到這一刻的感受嗎」.
「當你聽到你快要死的時候」.
「當你領悟到你所有的成就」.
「你所建立的王國,所有一切」.
「都不能幫到你的時候」.
「你可以怎樣做呢」.
起初我否認這個事實.
我不想相信.
我相信我可以用另類療法治好自己.
我服用齋藥,服用補充劑.
我可以去針灸,中藥.
做的事我都可以做.
我全都做了,就是沒有做手術.
我推遲了手術九個月.
我覺得我用我人生的意志.
可以戰勝癌症.
就像當年我用我人生的意志.
去帶領蘋果成為一個.
最頂尖的科技巨頭一樣.
但最後我還是被癌症打敗了.
好,如果我們從這個角色抽出來.
如果Steve Jobs.
他是一個虔誠的基督徒.
你看他成功的一面.
他建立了一間公司.
是世界頂尖的科技巨頭.
大家就會覺得上帝真的很祝福他.

$^{161}$令他創造了不少影響世界的科技.
今天你和我手上的智能手機.
雖然未必是用蘋果.
智能手機有很多技術都是蘋果手作.
不單是在iPhone裡面.
如果我們只看半身.
他前半身30歲就有過億身家.
50歲就重返蘋果公司.
創作了很多科技.
你會否想像這樣的人生.
上帝很祝福他.
但如果我們只看他的苦難.
30歲建立公司,但被人出賣.
接著打拚了二十多年.
重返蘋果公司.
再創佳績.
在事業最高峰的時候.
竟然被病魔打敗了.
如果我們只看這一點.
Steve Jobs這麼辛苦這麼多年.
最後還是怎樣.
只有空穴.
我們是否覺得上帝是耍花樣.
人生不是一個定點,人生是一個旅程.
透過人與人之間,人與人事之間的互動.
不斷的變化.
當中有高有低,不斷被逆轉,被反轉.
當你以為事情是這樣走的時候.
事情好像往往不是我們想像的那樣.
所以今天跟大家的角色扮演.
想跟大家探討一下.
第一,基督徒怎樣看人生的高低.
第二,耶穌在路加福音中.
怎樣把逆轉人生的概念.
帶給我們身邊的人.
第三,當我們處於順利.
或處於低谷的時候.
我們怎樣做.
我們首先要讀經文.
「路加福音」,4章16節至30節.

$^{201}$耶穌來到祂自己長大的地方拿撒勒.
在安息日會堂站起來讀聖經.
有人拿了以賽亞書給祂.
祂選了這段經文.
「主的靈在我身上,因為祂用告告我」.
「叫我傳福音給貧窮的人」.
「差遇我宣告被擄得釋放」.
「失明的漢賤受壓迫的得自由」.
「宣告上帝遇納人的欺凌」.
然後祂還給了書人.
會堂的人就定睛看著祂.
耶穌就跟他們說.
「你現在聽到的這段經文已經應驗了」.
但要留意.
這裡的字眼是傳福音和欺凌.
傳福音,大家也很熟悉吧.
把一個好消息帶給未信者.
例如耶穌拯救世人.
從罪中拯救我們.
但欺凌是我們比較少接觸的概念.
這裡跟傳福音有很大的關係.
就是耶穌如何把逆轉人生的概念.
帶出來的關鍵.
要解釋這個關鍵.
我待會要說會堂有甚麼特別.
大家有留意到耶穌來到會堂,對吧.
如果大家….
兩個月前有聽我說.
一個講道叫「肅靈方言」的時候.
講道《路加福音》四章的開頭.
耶穌勝過魔鬼的試探.
打破了魔鬼的方言之後.
然後來到這裡.
四章十六節.
耶穌在會堂裡教導人.
為何會堂那麼重要.
會堂是甚麼功能.
在聖殿被毀之後.
猶太人散居在不同地方.
而會堂承接了.

$^{241}$聖殿本身教導的功能.
而律法成為了猶太人的信仰核心.
不再是聖殿的憲制.
而是律法的實踐.
所以這段經文.
剛剛耶穌展示他有勝過魔鬼的能力之後.
來到會堂教導猶太人的律法.
顯示他對律法.
對猶太人的信仰核心.
有一個很完整的理解.
即是他有權力.
亦有能力去解釋聖經.
當聖殿被毀之後.
很多猶太人都會建立一個制度.
叫做「猶太拉比」.
即是他們會堂的老師.
他有這個能力去解釋聖經.
而耶穌這裡展示了這個能力.
為他之後傳的福音.
帶來的好消息.
建立了一個很好的基礎.
就是他代表上帝.
來跟他的子民說話.
所以如果你看四章的話.
上文第一節和第十四節.
提及耶穌沒有聖靈.
耶穌帶著聖靈的能力是有原因的.
因為耶穌在做救若先知的身分.
就好像救若上帝的大人一樣.
以賽亞那樣.
有權柄向眾人宣告.
耶穌即將要做的事.
所以耶穌引用以賽亞書先知的說法.
主的靈在我身上.
要做一系列的事.
要傳福音給貧窮的人.
要宣告上帝遇難人的氣憐.
這裡傳的福音帶來的好消息.
對於會堂裡的人.
其實不一定是好消息.

$^{281}$為甚麼這麼說呢.
還記得大家之前有上過生命之道嗎.
老師也有提到.
這裡的氣憐跟以賽亞書的恩憐.
是很不同的,對吧.
大家有否留意到路加福音用的是甚麼.
宣告上帝遇難人的氣憐.
而以賽亞書的是耶和華的恩憐.
兩句是不同的.
為甚麼呢.
在第三,二,三頁成書時.
關注的是猶太人回歸.
被人打散,分散了被擄.
他們要回歸.
要趕走佔領耶路撒冷的軍隊.
光復耶路撒冷,重建聖殿.
所以以賽亞書有一句.
宣告耶和華的恩憐.
是有行動的,是甚麼.
「和我們的神報仇的日子」.
「安慰所有悲哀的人」.
意思是耶和華的恩憐來到時.
就是上帝幫我們報仇.
光復耶路撒冷的日子.
所以之前受苦,被殺害,失蹤的手足.
就會得到安慰.
而耶穌這裡的引用.
轉為了氣憐.
還要故意不用我們上帝報仇的日子.
究竟發生甚麼事.
原來在耶路撒冷被佔領後.
第一聖殿被毀後.
猶太人…剛才說了.
散居在不同的地方.
不能再在自己的國家.
有一群人來到死海附近居住.
他們把以賽亞書的恩憐改為氣憐.
目的是當他們面對破家亡的苦難時.
他們期待有一天來塞亞江南.
可以像當年上帝命令猶太人一樣.

$^{321}$把土地物歸原主.
把變成奴隸的人變回自由身.
把受壓制的猶太人釋放出來.
把外國勢力趕走,為他們報仇.
所以當耶穌被聖靈充滿.
受完試探,去完會堂.
教導以賽亞書61章時.
還要提到氣憐這個字.
所有猶太人就會覺得.
耶穌就是解放他們的彌塞亞.
而耶穌還要很慚愧地說.
「今日這段經文已經應驗了」.
他們期待著,「來吧」.
「讓我帶領大家起義,光復耶路撒冷」.
誰知耶穌一盆冷水淋下去.
耶穌就用以利亞醫治外邦人乃曼的故事.
告訴猶太會堂的猶太人.
「這個救恩,這個光復耶路撒冷」.
「你們是沒有份的」.
因為你們當年拒絕先知.
唯一能夠跟救恩有份的.
就是接納耶穌.
還要用一種.
不是我們想像的方式去接納耶穌.
對會堂來說,這絕對不是好消息.
是一個半路被人執拾的壞消息.
守了這麼久的律法.
做了這麼多年憲制.
你現在跟我說.
所有人都可以有救恩.
就像「黑星星期五」一節大醜賓.
大家有去排隊嗎.
這幾年香港大學,大家有去排隊嗎.
有些人排了很久隊.
整晚「終於到你了」.
「你可以用一節去買大電視」.
但突然老闆說.
「全場所有人,每個人都可以買到」.
「只要你肯進來」.
排了這麼久隊的人很生氣.

$^{361}$不肯進去.
所以耶穌宣告.
只有「乃萬德潔正」.
其他人都不可以.
守了這麼久的律法.
原來只要跟從耶穌.
就可以得到猶太人這麼久以來.
夢寐以求的救恩.
對猶太人來說.
這是逆轉的人生.
這是很大的打擊.
所以在會堂裡的猶太人.
他們有很大落差.
他們很想殺了耶穌,推祂下山.
因為耶穌這個彌賽亞.
跟他們期待的救主有很大不同.
會堂的猶太人期待他們.
能夠結束他們自己的苦難.
但耶穌卻不是用他們期待的方式.
帶來救恩.
耶穌是會報仇.
不過不是用他們的方式.
不知道大家初信耶穌時是怎樣.
有些人可能以為信了耶穌.
就會一帆風順,風生水起.
以為我們受了這麼久的苦.
終於等到了.
當然,我們在網上聽到很多見證.
也聽到有些弟兄姊妹.
基督徒撈到風生水起.
但我很肯定不是每一個.
你覺得你信了耶穌之後.
你發了財,請你舉手.
因為這不是我們信了耶穌.
信耶穌的本質,對吧.
信耶穌的本質是甚麼.
信耶穌的本質不是這個大哥.
給了我們一些東西,我們想要.
所以我們要跟隨他,對吧.
而是無論我們人生怎樣反轉.

$^{401}$我們都要跟隨下去.
但如果我們像會堂裡的朋友一樣.
我們長期都是這種心態.
你給我一些好東西,我就跟你.
你不給我好東西,我就追殺你.
這是一個黃大仙式的信仰.
這是操縱上帝的信仰.
我們會否用一個.
在拜黃大仙的方式去信耶穌.
耶穌給我們人生的逆轉是.
你以為你人生會一帆風順.
你可以窮得突然物到穿窿.
但你放心,上帝會看著你.
你以為你人生得苦難的時候.
又會突然被你豁然開朗.
從上帝的角度看到人生的美好.
上帝是否一個玩弄的上帝.
我經常這樣問自己.
我猜耶穌在《路加福音》的經文.
已經回答了.
這段經文給我提醒的是.
不要習以為常.
不要以為上帝給我們的祝福.
是我們的英分.
是我們應該要有的東西.
祂一定要給我們.
剛過去的感恩節大家如何度過.
大家以甚麼心態度過感恩節.
有些姐妹可能在疾病和財務的掙扎當中.
當我們期待著耶穌.
沒有移走我們的苦難的時候.
祂跟我們一起承受我們的苦難.
我們如何回應.
可能我們有很多懷疑.
是否我做得不好.
為何有人會被祝福.
而受祝福的人好像不是我自己.
為何上帝只祝福其他人.
就像會堂裡的猶太人一樣.
熬了那麼久.

$^{441}$竟然是別人有份而我沒有份.
我們的信仰和民間偶像最大的分別就是.
耶穌做主,不是我們做主.
所以無論我們人生高低.
我們都相信耶穌是我們的主.
祂叫我們走到哪裡,我們就去哪裡.
我們唯一可以做的.
是接納耶穌對我們人生的指引.
同一時間,我也很掙扎.
面對著苦難.
我也很失望和痛苦.
為何事情會淋到我身邊.
為何事情會淋到我們的社區.
為何這場災難要淋到大埔.
為何香港人要感受那麼大的折磨.
你問我,我也很掙扎.
這些事情我也不敢輕易地給一個答案.
我唯一可以做的事.
就是保持一顆開放的心.
去問上帝究竟發生甚麼事.
我也在學習.
去過耶穌給我的逆轉人生.
這段經文提醒我.
保持一個心態.
就是當耶穌的話語淋到的時候.
隨時都有逆轉的可能.
可能是從我們的角度來說.
是由好變壞.
就像七個月前.
七個月前的今天.
我剛剛開始在UV事奉.
好像挺好.
第一天我們就被車撞.
我和Alice被車撞到.
然後又被車撞.
當時我們的角度是由好變壞.
第一天就弄成這樣.
大家在過去的感恩節.
我們很感恩上帝保住了我們的命.
是因為這件事.

$^{481}$第一次會很有愛.
很快地接納了我們.
很快地融入了我們的大家庭.
突然間這件事從壞變回好.
大家明白那個心路歷程吧.
我猜這就是逆轉人生的一部分.
當我們以為這個世界壞到極處的時候.
上帝突然揭示一些美好.
可能弟兄姊妹.
你可能正在享受豐盛的人生當中.
上帝有祝福你.
我很建議或鼓勵你.
將這份祝福成為你的動力.
走出這個社區.
看看這個世界有沒有人需要.
我們所擁有的東西.
完全是上帝給我們的.
幫助我們身邊的人度過逆轉的時刻.
人生的低谷.
對身邊的人是一個很大的祝福.
這幾天.
我們經常留意大埔發生的五級火災.
有很多令人難過的消息.
很沮喪,很憤怒.
我猜對整體香港人來說.
都是很大的憤怒,很大的打擊.
很大的心理陰影.
我們可以盡我們的力.
跟香港人一起走下去.
或是一個訊息,一個問候.
一個電話.
關心一下他們的狀態.
這幾天香港的凌晨.
其實我也….
在社交平台,IG,Telegram.
都有跟青少年聊天.
他們連續幾天都睡不著.
他們有很多情緒.
如果你可以給他們一點時間.
聽他們說話,陪伴他們.

$^{521}$讓他們知道.
雖然我們遠在美國這邊.
但我們不是靠岸觀火.
我想對他們也是一個很大的祝福.
如果你想有錢出錢,有力出力.
捐錢給大埔居民.
我鼓勵你千萬….
要捐給一些真的能幫助災民的機構.
所以你捐錢的時候,請你留意.
已經有一些國家級的機構.
趁著這場慘劇發過冷財.
還有些騙徒騙了災民.
拿了他們的身分證.
去誣蔑拿資助.
我們真的要為災後的安置.
賠償祈禱.
為他們的公義.
為他們得到的安慰祈禱.
希望真正有需要的市民能夠幫助.
所以總結.
基督徒怎樣看待我們的人生高低.
無論我們是否基督徒.
人生必定有高山有低谷.
我們每個人都會面對人生的風浪.
就像經文所說.
猶太人面對著苦難和迫害.
而耶穌指出.
祂的救恩並非大家想像的那種.
祂沒有一下子把我們所有的苦難.
而是把苦難放在自己身上.
跟我們一起走.
我們的選擇是.
雖然我們不知道耶穌在做甚麼.
但我們相信祂會帶我們走過這個低谷.
我們願意被祂這樣帶領.
在《路加福音》中.
耶穌用了一個以色列書的英雄.
一個以利沙爾尼亞的歷史故事去表達.
每個人的人生都會有逆轉的時刻.
你們選擇追殺耶穌.

$^{561}$還是選擇跟隨耶穌.
當我們處於順利的時候.
或是低谷的時候,要怎麼做.
保持對上帝開放的心態.
順利的時候.
想想身邊的人的需要.
把上帝給你的祝福.
分給一些有需要的人.
這個世界很需要愛.
很需要有人去平衡公義.
讓我們在低谷的時候.
看著耶穌跟我們一起受苦.
讓我們有捱下去的盼望.
跟我們一起祈禱.
天父,爸爸.
你們深深感受到在苦難中.
裡面的掙扎.
主人,我們呼求你.
很多人在….
在苦難的裡面.
我們不知道該怎麼辦.
(祈禱).
我們有很多疑問.
求你讓我們保持很好的關係.
讓我們有開放的心.
有一天你臨到我們的時候.
我們的生命被逆轉.
我們也會跟著你,幫助我們.
這樣透過這封耶穌名禱.
阿們.
祝福大家.
\newpage



\section{士師記 2:6-23}
\label{sec:VcgHfBIcHMQ}
\textbf{Enduring Faith (Judges 2\_6-23) - Chee Yie Jong 堅韌的信仰(士師記2\_6-23) - 楊志溢傳道}
\newline
\newline
連結: \href{https://youtube.com/watch?v=VcgHfBIcHMQ}{\texttt{https://youtube.com/watch?v=VcgHfBIcHMQ}} ~~~~ 語音日期: 2025-12-08
\newline
\newline
\hyperref[sec:45ohI_TBLkg]{\small{< < < PREV SERMON < < <}}
~
\hyperref[sec:index_chronic]{\small{[返順時目]}}
~
\hyperref[sec:index_scriptual]{\small{[返順卷目]}}
~
\hyperref[sec:8eaa7ARgWZg]{\small{> > > NEXT SERMON > > >}}
\newline
\newline
士師記 2:6-23
\newline
\begin{longtable}{cl}
\hline
\hline
章節 & 經文 (和合本修訂版)\\
\hline
2:6 & \begin{tabularx}{0.7\textwidth}{X} 約書亞解散百姓，以色列人回到自己的地業，佔各自的地。 \end{tabularx} \\ \\ \relax
2:7 & \begin{tabularx}{0.7\textwidth}{X} 約書亞在世的日子和他死了以後，那些見過耶和華為以色列所做一切大事的長老還在世的時候，百姓都事奉耶和華。 \end{tabularx} \\ \\ \relax
2:8 & \begin{tabularx}{0.7\textwidth}{X} 耶和華的僕人，嫩的兒子約書亞死了，那時他一百一十歲。 \end{tabularx} \\ \\ \relax
2:9 & \begin{tabularx}{0.7\textwidth}{X} 以色列人把他葬在他自己地業的境內，以法蓮山區的亭拿‧希烈，在迦實山的北邊。 \end{tabularx} \\ \\ \relax
2:10 & \begin{tabularx}{0.7\textwidth}{X} 那世代的人也都歸到自己的列祖。後來興起的另一世代不認識耶和華，也不知道他為以色列所做的事。 \end{tabularx} \\ \\ \relax
2:11 & \begin{tabularx}{0.7\textwidth}{X} 以色列人行耶和華眼中看為惡的事，去事奉諸巴力。 \end{tabularx} \\ \\ \relax
2:12 & \begin{tabularx}{0.7\textwidth}{X} 他們離棄領他們出埃及地的耶和華－他們列祖的神，去隨從別神，就是四圍列國的神明，向它們叩拜，惹耶和華發怒。 \end{tabularx} \\ \\ \relax
2:13 & \begin{tabularx}{0.7\textwidth}{X} 他們離棄了耶和華，去事奉巴力和亞斯她錄。 \end{tabularx} \\ \\ \relax
2:14 & \begin{tabularx}{0.7\textwidth}{X} 耶和華的怒氣向以色列發作，把他們交在搶奪他們的人手中，又把他們交給四圍仇敵的手中，以致他們在仇敵面前再也不能站立得住。 \end{tabularx} \\ \\ \relax
2:15 & \begin{tabularx}{0.7\textwidth}{X} 他們無論往何處去，耶和華的手都以災禍攻擊他們，正如耶和華所說的，又如耶和華向他們所起的誓；他們就極其困苦。 \end{tabularx} \\ \\ \relax
2:16 & \begin{tabularx}{0.7\textwidth}{X} 耶和華興起士師，士師就拯救他們脫離搶奪他們之人的手。 \end{tabularx} \\ \\ \relax
2:17 & \begin{tabularx}{0.7\textwidth}{X} 然而，他們卻不聽從士師，竟隨從別神而行淫，向它們叩拜。他們列祖所行的道，所聽從耶和華的命令，他們都速速偏離了，並不照樣遵行。 \end{tabularx} \\ \\ \relax
2:18 & \begin{tabularx}{0.7\textwidth}{X} 耶和華為他們興起士師，耶和華與士師同在。士師在世的一切日子，耶和華拯救他們脫離仇敵的手。耶和華因他們受欺壓迫害所發出的哀聲，就憐憫他們。 \end{tabularx} \\ \\ \relax
2:19 & \begin{tabularx}{0.7\textwidth}{X} 但士師一死，他們又轉去行惡，比他們祖宗更壞，去隨從別神，事奉叩拜它們，總不放棄他們的惡習和頑梗的行為。 \end{tabularx} \\ \\ \relax
2:20 & \begin{tabularx}{0.7\textwidth}{X} 於是耶和華的怒氣向以色列發作，說：「因為這國違背我吩咐他們列祖當守的約，不聽從我的話， \end{tabularx} \\ \\ \relax
2:21 & \begin{tabularx}{0.7\textwidth}{X} 約書亞死的時候所剩下的各國，我必不再從他們面前趕出任何一個， \end{tabularx} \\ \\ \relax
2:22 & \begin{tabularx}{0.7\textwidth}{X} 為要藉此考驗以色列是否肯謹守遵行耶和華的道，像他們列祖一樣地謹守。」 \end{tabularx} \\ \\ \relax
2:23 & \begin{tabularx}{0.7\textwidth}{X} 耶和華留下那些國家，不將他們速速趕出，也不把他們交在約書亞的手中。 \end{tabularx} \\ \\
[1ex]
\hline
\hline
\end{longtable}
$^{1}$Thank you, Kelly..
So the first few years after college,.
I attended a small group at my church at that time.
that was pretty multi-generation..
And one of the things that I appreciate.
as an early 20-plus year or so in that group.
is that I got to hear about different parenting.
stories and experience within the group,.
from younger parents with little kids.
to older couples with adult children..
And during my time in that group,.
one of the older ladies, who also.
happened to be the pastor's mom, she.
would always ask the group to pray.
for her other adult child who had left the church..
[SPEAKING CHINESE].
And so those moments made me start.
to wonder, how is faith being passed down.
from one generation to another?.
Looking at my own family, I think.
we are blessed enough that me and all my siblings.
are still walking with the Lord as adults now..
But how big of a role do parents play in this?.
And how much responsibility do adult children.
have when it comes to keeping their parents' faith?.
[SPEAKING CHINESE].
Today, we are going to continue with our series.
in the Book of Judges..
And specifically, we are going to look.
at how each generation of Israel, after Joshua,.
they slowly loses their identity as God's chosen people.
and forsake the faith of their ancestor..
Even though Israel was supposed to be a kingdom of priests.
and a holy nation for the Lord, they.
ended up becoming like the rest of the idolatrous nation.
around them..
[SPEAKING CHINESE].
So in my last sermon, we look at Judges 1,.
where the 12 tribes of Israel, they.
began to settle into their allotted territory.

$^{41}$in the promised land..
And so their main task at this stage of the conquest.
was to drive out the remaining Canaanites who still live.
in the land in order to avoid their idol worship.
and also their moral corruption as Israel slowly.
settling into this land..
[SPEAKING CHINESE].
However, as we see in Judges 1, Israel.
settled for half-hearted obedience,.
as they did not fully drive out all the Canaanites.
as the Lord has commanded them to..
And even worse, they ended up making.
covenants with these Canaanites and worshipping their idols..
[SPEAKING CHINESE].
And so if you recall from my last sermon,.
Judges begin with this sort of double introduction..
The first introduction that we look at last time,.
it highlights Israel's incomplete and half-hearted.
conquest, where they failed to drive out the Canaanites..
And so in other words, this is an assessment.
of Israel's military and political failure..
[SPEAKING CHINESE].
And the second introduction sections.
that we will be looking into today,.
it highlights Israel's incomplete and half-hearted.
faith..
It gives us an overview of what is.
about to happen in this main section of the book..
And it also introduces to the very important cycle.
of Israel's spiritual failure..
[SPEAKING CHINESE].
So this morning, we will dive into the second introduction.
section, and we will explore these questions together..
What does enduring faith that lasts through generations.
requires of us?.
How do we keep ourselves from falling.
into the same situation as this judges' generations,.
where each generation after each generation,.
they just fall further and further away.
from the faith of their ancestors?.

$^{81}$[SPEAKING CHINESE].
And so we will answer these questions from our passage.
today..
And if you have your own Bible, you're.
welcome to go to Judges 2 and follow along..
And so first, enduring faith requires us to remember..
[SPEAKING CHINESE].
And so Judges 2, especially in verse 7,.
it restates the very important details.
about the generations before, that is the Joshua generation..
They remain faithful and steadfast to the Lord..
[SPEAKING CHINESE].
But when we come to verse 10, things actually.
take a pretty drastic turn..
We learn that after the whole Joshua generation had.
been gathered to their ancestors,.
another generation grew up who did not know the Lord.
or remember the things he had done for Israel..
[SPEAKING CHINESE].
And so to put it another way, they.
forgot who Yahweh is and what he had done for their ancestors..
[SPEAKING CHINESE].
And what is the result of their forgetfulness?.
We see this in the following three verses,.
so starting in 11..
"Then the Israelites did evils in the eyes of the Lord.
and served the Baals..
They pursued the Lord, the God of their ancestor,.
who had brought them out of Egypt..
They followed and worshipped various gods.
of the peoples around them, and they aroused the Lord's anger..
Because they pursued him and served Baal and Arcturus.".
[SPEAKING CHINESE].
It is a little hard to notice this.
in the English or the Mandarin translation,.
but in the original Hebrew manuscripts, after verse 10,.
all these action verbs in the next three verse stages.
comes one after another to show the degree of Israel's.
apostasy..
"They did evils in the eyes of the Lord..

$^{121}$They pursued the Lord..
They followed other gods..
They worshipped the gods of the people around them..
They angered the Lord..
Then again, they pursued the Lord.
and served Baal and Arcturus.".
[SPEAKING CHINESE].
And so forgetfulness comes with consequences..
I remember when I went to get my US driver's license,.
and at that point, I had been driving back in Malaysia.
for three parts a year already..
So I signed up for a driving test, and I went to the DMV..
[SPEAKING CHINESE].
And on that day, I was acing through the test..
Vehicle check, done..
Parallel part, done..
Three-point turn, done..
And then the proctor asked me to get on the street,.
and again, I cruised through it without any issues at all..
And when we made our way back to the DMV after driving around,.
I was feeling pretty confident..
And we came to this one last stop sign inside the DMV..
[SPEAKING CHINESE].
And so you have to understand that in Malaysia,.
a stop sign is more like a suggestion..
And so most drivers, they will slow down.
when they approach the stop signs..
And if they don't see any pedestrians or any vehicles.
around, we will just continue on without coming to a full stop..
And so when I came to this final stop sign inside the DMV,.
I forgot that I was in a driving test..
I allowed my old instinct to take over,.
and I just blew straight past the stop sign..
And immediately, the proctor turned at me.
and pointed at the stop sign..
And at that moment, I knew I wasn't going to pass..
And so I ended up having to go back,.
experience through all the hassles of DMV.
again after a few weeks, just to redo my exam..
[SPEAKING CHINESE].

$^{161}$And so in the same way, spiritual forgetfulness.
also comes with serious implications for our faith..
When we forgot what the Lord has done for us,.
when we forget his commands for us,.
when we forget his grace and his mercy in our lives,.
we slowly lose our faith and allow.
our old self and sinful nature to take over our lives again..
[SPEAKING CHINESE].
And so how do we counter this sort of spiritual forgetfulness.
or spiritual dementia?.
It is by remembering..
Remembering is a common practice for God's people.
all throughout the generations..
[SPEAKING CHINESE].
For example, in the Old Testament,.
the Passover was being instituted.
as an annual festival in the Jewish calendar.
so that all generations of Israels that comes after,.
they can remember what the Lord has done for them.
in the first Passover in Exodus by delivering them.
from slavery in Egypt..
[SPEAKING CHINESE].
And when we come to the New Testament,.
we also find other examples as well..
The early church, they partake in the Lord's Supper.
or communion as a remembrance of what Jesus.
has done for them at the cross..
[SPEAKING CHINESE].
And even today, believers all throughout the world.
partake in the sacraments of holy communion.
as well to remember the sacrifice of Jesus..
And in fact, we just did that a few moments ago.
in our service..
[SPEAKING CHINESE].
And therefore, remembering is a spiritual discipline.
that we should all practice so that our faith can be lasting..
[SPEAKING CHINESE].
And so what does enduring faith requires of us?.
First, it requires us to remember,.
to remember who God is and what He has done for us..

$^{201}$And this is the first lesson that we can.
learn from Israel's history here..
And secondly, enduring faith requires.
us to actively repent and turn from our own ways..
[SPEAKING CHINESE].
And so let's pick up our passage again from verse 14..
"In His anger against Israel, the Lord.
gave them into the hands of the raiders who plundered them..
He sold them into the hands of the enemies all around whom.
they were no longer able to resist..
And whenever Israel went out to fight,.
the hand of the Lord was against them to defeat them..
Just as He has sworn to them, they were in great distress.".
[SPEAKING CHINESE].
And so when Israel abandoned the Lord and worshiped other idols,.
the Lord gave them into the hands of the oppressors,.
where Israel, they would suffer for a period of time..
[SPEAKING CHINESE].
And then in verse 16, "Then the Lord.
raised up judges who saved them out.
of the hands of these raiders..
And yet they would not listen to their judges,.
but prostitute themselves to other gods and worship them..
They quickly turned from the ways of their ancestor.
who had been obedient to the Lord's command.".
[SPEAKING CHINESE].
And so what happened here?.
"When Israel suffered under the hand of the oppressors.
and cried out to the Lord, the Lord.
would raise up judges to deliver them from their oppressions.
and sufferings.".
[SPEAKING CHINESE].
And who are these judges?.
First, they are Israel's deliverer..
But secondly, they also serve as a catalyst.
for godly living among the people of God..
[SPEAKING CHINESE].
But this is a very important but..
Israel's disobedience-- Israel's obedience.
did not last long at all..

$^{241}$Look at verse 17 again..
It says, "Yet they would not listen to their judges,.
but prostitute themselves to other gods and worship them..
They quickly turned from the ways of their ancestor.
who had been obedient to the Lord's command.".
[SPEAKING CHINESE].
And then in verse 18 and 19, again, we.
see the same pattern happening all over..
"Whenever the Lord raised up a judge for them,.
he was with the judge and saved them.
out of the hands of their enemies,.
as long as the judge lived..
For the Lord relented because of their groaning.
under those who oppressed and afflicted them..
But when the judge died, the people.
returned to ways even more corrupt.
than those of their ancestors, following other gods.
and serving and worshiping them..
They refused to give up their evil practice.
and stubborn ways.".
[SPEAKING CHINESE].
And so even though the Lord would raise up judges.
to deliver Israel from their oppressions,.
we learn that when the judge died,.
the people actually returned to ways even worse.
than their ancestors..
[SPEAKING CHINESE].
And so think about it this way..
The first time that you did something wrong,.
it makes sense for mercy and forgiveness.
to be shown right because of various reasons..
It can be because of ignorance, or it was just an accident,.
and so on..
But when you repeat the same sin over and over again.
and show no signs of remorse and repentance,.
it just becomes like a cancerous disease.
that slowly eats you away and corrupts your heart..
[SPEAKING CHINESE].
So this is what is happening to this judge's narration..
Verse 19 ends with a summary on this corrupted generation,.

$^{281}$that they refused to give up their evil practice.
and stubborn ways..
And because of the refusal to repent and turn.
from the evil ways, Israel just continues.
to spiral down into worse and worse apostasy..
And this cycle of sinning, oppressions,.
crying out to the Lord, deliverance, peace,.
and then back to sinning again, it.
just keeps repeating itself with no end in sight..
[SPEAKING CHINESE].
And so how do we avoid the downfall of Israel here?.
The answer is simple, but to actually lift it out.
actually requires a lot of intentionality on our end..
We need to truly repent and turn from our evil ways..
[SPEAKING CHINESE].
And true repentance is not just a one-time event,.
but it is a day-to-day, moment-by-moment practice..
[SPEAKING CHINESE].
One metaphor that I find helpful to think.
about this practice of repentance.
is spiritual breathing..
As long as we are alive, we never stop breathing, right?.
And so this is the same when it comes to spiritual breathing..
We exhale by confessing our sin, knowing.
that Jesus has already taken away our sins at the cross..
We then inhale by surrendering to God,.
turning from our own ways and turning back to him.
as the ultimate authority of our lives..
[SPEAKING CHINESE].
And so the more we are able to practice.
the spiritual discipline of repentance,.
and it becomes more nature to us,.
the more we are able to align our lives.
with the heart of God..
[SPEAKING CHINESE].
And so this is the second lesson we can learn from Israel.
when it comes to keeping our faith,.
that it requires us to truly repent and turn.
from our evil ways constantly..
If not, our heart will quickly become hardened,.

$^{321}$and we fall further and further away from the Lord..
[SPEAKING CHINESE].
And the third lesson that I want us to take away from our passage.
today is that enduring faith requires.
us to acknowledge our own insufficiency.
and our need for a savior..
[SPEAKING CHINESE].
And so I know we have been taking lessons.
from Israel's failure here, but I.
want to speak up for them here for a little bit..
A lot of time when we read scripture,.
we can be quick to label a certain characters.
or a group of people as just simply good or bad..
But when in fact, there's a lot more complexity to it,.
and more often than not, we actually.
share the same tendency as these characters as well..
[SPEAKING CHINESE].
And when we look at this new generation of Israel.
after Joshua, all of them grew up.
wandering in the wilderness, right?.
This is the only way of living that they were familiar with..
And so now that they have settled down.
into the promised land and own their own land,.
they had to pick up new skills in this new environment.
in order to survive, with agriculture being.
the primary mode of living here..
[SPEAKING CHINESE].
And so this is where the appeal of worshipping Baals and Ash.
Torres comes in, because their Canaanite neighbors believe.
that these gods have the power to control rain, the storm,.
and fertility..
And thus, they play an important role in an agrarian society..
And so for this new generation of Israel,.
the way to survive in this new place.
is just to look around them and learn.
from their Canaanite neighbors, worshipping the storm and rain.
gods with the hope that these gods will also bless them.
and give them good crops..
[SPEAKING CHINESE].
And so from their perspective, this.

$^{361}$is actually a completely logical move..
And we should not be quick to judge them.
without understanding the challenge they.
face in order to survive..
[SPEAKING CHINESE].
Yet at the same time, we also have.
to acknowledge that the biggest issue here.
is that they choose human logic over trusting God,.
the same God who provided for them even in the wilderness.
and who will surely provide for them in the promised land.
as well..
[SPEAKING CHINESE].
And so in other words, they think.
that they know how to survive in this new environment.
better than God, even though God has specifically.
commanded them time and time again.
to not make covenant with the Canaanites,.
or else they will be tempted to follow the detestable things.
they do in their idols' worship and sin against the Lord..
[SPEAKING CHINESE].
And so as a result of their incomplete faith,.
we see that in verse 20, "Therefore, the Lord.
was very angry with Israel and said,.
because this nation has violated the covenant I.
ordain for their ancestor and has not listened to me,.
I will no longer drive out before them any of the nations.
Joshua left when he died..
I will use them to test Israel and see.
whether they will keep the way of the Lord.
and walk in it as their ancestors did..
The Lord had allowed those nations to remain..
He did not drive them out at once by giving them.
to the hands of Joshua.".
[SPEAKING CHINESE].
So this is the predicament that Israel found themselves in..
As we look at the stories of all the judges.
in the rest of the book, none of them.
were able to put an end to the cycles of spiritual downfall.
and moral corruptions that Israel found themselves in..
And this pattern of sin, oppression, crying out,.

$^{401}$deliverance, peace, and then back to sinning again,.
it just keeps on repeating itself..
[SPEAKING CHINESE].
And if you think you can find a happy ending by flipping.
to the end of the Book of Judges,.
you will actually find the complete opposite..
In the last few chapters of Judges,.
Israel descends into utter chaos,.
and the book actually ends with this remark..
"In those days, there was no kings in Israel..
Everyone did what was right in their own eyes.".
[SPEAKING CHINESE].
So where is our hope in all of this?.
[SPEAKING CHINESE].
I like how this scholar has put it..
He says, "God's covenant with Israel as a contract.
is broken because of Israel's apostasy,.
that is, Israel's failure to hold out.
the ends of the obligation..
But this covenant as a promise still stands,.
that is, God's promise to his people..
What is needed for its fulfillment.
is an Israel which does not fail to keep its obligation..
And in the New Testament terms, this fulfillment.
comes in Jesus, the Messiah, who is everything.
that Israel should have been but was not,.
the perfectly obedient Israel..
And this is where our hope can be found..
[SPEAKING CHINESE].
And so like the Israel who wandered in the wilderness.
for 40 years, Jesus spent 40 days and 40 nights.
in the wilderness, and he was being tempted by the Satan..
And yet he did not sin..
[SPEAKING CHINESE].
And like Israel who receives the Torah from the Lord.
and that tells them how they ought.
to live as God's covenant people,.
and yet they fail to live up to it,.
Jesus himself is the Word of God who became flesh..
And he shows us the true and better way.

$^{441}$of living by redefining what true and better righteousness.
means..
[SPEAKING CHINESE].
And so where Israel fails, Jesus remains perfectly.
faithful and obedient..
And that is why he is our ultimate hope..
[SPEAKING CHINESE].
And so this is our third and final lesson from Judges 2.
here..
Enduring faith requires us to acknowledge our inability.
to save ourselves like the Israelites,.
and then also realizing that we need a Savior, that.
is Jesus himself..
[SPEAKING CHINESE].
And this is the first and the most important steps..
Once we realize our own human depravity and our need.
for Jesus, we then repent and turn from our evil ways.
to follow the ways of Jesus, allowing.
this practice of repentance to become part of our daily lives..
And as we walk through the rest of our faith journey,.
we also take the time to reflect on who God is.
and remember his gift of salvation, of grace.
and mercy in our life..
[SPEAKING CHINESE].
So let's pray together..
[SPEAKING CHINESE].
Heavenly Father, as we enter into this season of Advent,.
we want to take some time to reflect and remember.
your gift of salvation..
Thank you that you did not leave us in our own mess,.
but you sent your only son, Jesus,.
to die on the cross for us..
[SPEAKING CHINESE].
Jesus, you are Emmanuel, God with us..
Thank you for being with us in our brokenness.
and our suffering..
You are the founder and perfecter of our faith..
And for the joy that lay before you,.
you endured the cross, despising its shame,.
and sat down at the right hand of the throne of God..

$^{481}$[SPEAKING CHINESE].
And Holy Spirit, you are God within us..
Thank you for dwelling in us and grant us the power.
to live a life of active repentance.
so that our faith is active and living..
[SPEAKING CHINESE].
And in your most precious name, we pray..
Amen..
[SPEAKING CHINESE].
Amen..
All right, thank you, Nhu-Vy..
Thank you all for joining us..
You are dismissed..
So go in peace..
You are loved..
[SPEAKING CHINESE].
[MUSIC PLAYING].
(music).
\newpage



\section{出埃及記 2:23-25}
\label{sec:8eaa7ARgWZg}
\textbf{願意靠近我們的神(出埃及記2\_23-25) - 江芷敏傳道}
\newline
\newline
連結: \href{https://youtube.com/watch?v=8eaa7ARgWZg}{\texttt{https://youtube.com/watch?v=8eaa7ARgWZg}} ~~~~ 語音日期: 2025-12-08
\newline
\newline
\hyperref[sec:VcgHfBIcHMQ]{\small{< < < PREV SERMON < < <}}
~
\hyperref[sec:index_chronic]{\small{[返順時目]}}
~
\hyperref[sec:index_scriptual]{\small{[返順卷目]}}
~
\hyperref[sec:3dbOIx2ZjUI]{\small{> > > NEXT SERMON > > >}}
\newline
\newline
出埃及記 2:23-25
\newline
\begin{longtable}{cl}
\hline
\hline
章節 & 經文 (和合本修訂版)\\
\hline
2:23 & \begin{tabularx}{0.7\textwidth}{X} 過了許多年，埃及王死了。以色列人因做苦工，就嘆息哀求；他們因苦工所發出的哀聲達於神。 \end{tabularx} \\ \\ \relax
2:24 & \begin{tabularx}{0.7\textwidth}{X} 神聽見他們的哀聲，就記念他與亞伯拉罕、以撒、雅各所立的約。 \end{tabularx} \\ \\ \relax
2:25 & \begin{tabularx}{0.7\textwidth}{X} 神看顧以色列人，神是知道的。 \end{tabularx} \\ \\
[1ex]
\hline
\hline
\end{longtable}
$^{1}$(麥克風聲).
有點緊張.
我還在找麥克風.
掛在這裡.
(笑聲).
親愛的影子們,早安.
今天很大雪.
其實本星期也很冷.
感恩看到大家平安回來了.
本星期冷得有一天.
Ken跟我說.
「為何我的眼鏡這麼模糊」.
他就說「擦擦眼鏡吧」.
一脫下眼鏡.
就發現兩個鏡片都掉下來了.
不知道是否冷宿熱帳還是怎樣.
真的很冷,大家保重.
不知道大家今個星期覺得怎樣.
有些意大利影子們出了遠門.
求神保守他們的老土平安.
我也知道大家中間心情一般.
我也求上帝的話語安慰我們.
跟我們過渡艱難的時間.
今天我們會一起繼續看《出埃及記》.
《出埃及記》二章23至25節.
這段聖經可以說是說出埃及的原由.
其實只有三節,很簡短.
我們一起去讀,好嗎.
《出埃及記》二章23至25節.
預備起.
.
這段經文是說以色列人.
向上帝發出哀求,求救的聲音.
究竟他們發生了甚麼事.
其實上次我們提到第一章時.
也有提到一點.
讓我們一起去回顧一下.
《出埃及記》第一章.
可以說是從一開始是一個很蒙福的開始.
但慢慢轉變成哀哭的過程.

$^{41}$當約瑟仍是宰相的時候.
從左邊開始看起.
你們是從左邊開始看起.
以色列在一個很興盛的時候.
當約瑟仍是宰相的時候.
他們開始在埃及落地生根.
生養眾多.
而且極其強盛,遍滿了那地.
然而好景不常.
有一個新的埃及王上場.
他們不認識約瑟.
以色列人在埃及王眼中成為了威脅.
埃及王看見這群又大又強盛的民族.
就很擔心他們會反抗.
甚至在戰爭中倒戈相向.
幫助他們的強敵.
所以法老千萬要使計.
強迫這些以色列人做苦工.
去壓榨他們.
而且還下令修身婆.
在接生的時候殺死所有男嬰.
這樣他們就可以杜絕人口的增長.
不過上一次提到.
這個修身婆是敬畏神.
所以她沒有按照命令辦事.
然後神賜福給他們.
但是法老沒有停止.
她越發殘酷地害這群以色列人.
她說所有希伯來的男嬰一出生.
就要被扔進尼羅河.
進一步殘害這些以色列人.
然後我們再看第二章.
這些苦難繼續的加劇.
第二章一開始提到摩西的出生.
摩西出於利美的家庭.
她媽媽在殘暴的條例下.
把摩西藏起來.
但藏了三個月.
你看我們教的嬰兒都越來越大.
已經慢慢藏不住了.

$^{81}$到了三個月的時候.
她就把它放在一個蒲草箱.
放進尼羅河.
她也不知道可以怎麼做.
但是神很看顧.
她被法老的女兒發現了.
並且收養她,撫養她成人.
摩西長大了.
不過她很看不過眼.
那些埃及人欺壓她自己的民族.
所以她打死了那些欺負她自己的埃及人.
然後她也被追殺.
所以摩西就逃到米甸.
米甸這個地方遇到她的妻子.
生兒育女.
然後她在那裡開始牧羊的工作.
我們看到以色列人在埃及.
經歷了一些甚麼苦海.
我們看到他們受的苦是怎樣的.
一步一步地進逼.
所以我們剛剛看的經文.
就是二章的收尾.
說到罪事者開始不僅由這些故事引入.
他開始轉換了角度.
用一個很客觀的立場.
去說以色列人的狀況.
並且揭示在這些苦難中.
神是怎樣回應人.
在這短短的三節經文.
不知道大家有沒有留意到.
有一個很整齊的句式.
在二十四,二十五節.
大家看到吧.
有四句句子都是用誰開始的.
用神開始的.
而且他連續表達了四個很主要的動詞.
就是「神聽見」,「神顧念」.
「神看見」,「神知道」.
我們會逐一說下去.
第一就是「神聽見他們的哀聲」.

$^{121}$「他們因奴役而發的求救聲上達神」.
「神聽到他們的哀求」.
「哀聲」通常是指受到極大的痛苦.
而發出的聲音.
以色列人被奴役.
他們很痛苦,他們求救.
所以他們的聲音上達了上帝.
所以人的聲音原來能夠達到上帝.
對我們來說.
我們也覺得這是一件很平常的事.
因為我們平時也會祈禱.
我們覺得我們的聲音,祈禱.
能夠達到上帝.
但在當時的舊約.
可能他們的人不知道原來可以這樣.
因為在一個跟舊約時期很接近的.
古東文化裡.
這些人他們敬拜的神.
可能大部分叫巴力.
這個巴力的神.
原來跟人有一個很大的距離.
他們如何去求告神.
可能是用自己的血.
或者動物的血去獻祭.
以致神能夠回應他.
大家是否記得這個故事.
就是以利亞先知在加密山.
跟巴力先知大鬥.
巴力先知從早上叫到晚上.
神都沒有回應他.
這是在《列王記》上的18章27至29節記載.
我們一起讀一下.
到了正午,以利亞嘲笑他們.
說「大聲求告吧,因為祂是神明」.
「祂或許在默想,或許正忙著」.
「或許在路上,或許在睡覺」.
「祂該醒過來」.
其實是在勸他.
然後繼續說.
「他們大聲求告,按著他們的儀式」.

$^{161}$「用刀槍刺割自己,直到渾身流血」.
「中午過去了,他們狂呼亂叫」.
「直到獻祭的時候」.
「祂沒有聲音,沒有回應」.
「也沒有理睬」.
那些神有沒有聽到巴力先知的求告.
沒有.
即使他們怎樣割自己,渾身流血.
神也沒有聽到他們的呼吸.
因此好像跟我們說的上帝.
形成了一個很鮮明的對比.
神聽到我們的哀聲.
我們的求告聲音能夠上達神那裡.
神聽到.
而且我們跟神很接近.
有時你在家裡,可能叫到隔壁房間.
你也未必聽到.
但聽到的意思是甚麼.
你們的距離很近.
我們跟神的距離很近.
第二,「神顧念他與列祖所立的約」.
敘事者繼續說.
神很顧念他跟列祖所立的約.
有一些聖經譯本會譯作「紀念」.
其實意思是「Remember」.
「Call to mind」等等的意思.
我想無論用哪個字.
都是在表達上帝記得這件事.
上帝記得跟列祖所立的約.
上帝記得.
上帝所要紀念的是.
他跟亞伯拉罕,以撒,瓦戈所立的約.
這並不是三個獨立的約.
其實是一個延續的應許.
上帝應許亞伯拉罕.
同樣應許以撒,瓦戈.
他們的後裔要成為大國.
要得著應許之地.
要成為萬國的祝福.
他們開始知道跟神所立的約.

$^{201}$但到了埃及這個時期.
這個約保守估計應該超過了400年.
400年也經歷了好幾代人.
不知道以色列人還有沒有人記得.
我來了美國之後.
很多人都問我.
「你是甚麼組織?你是哪裡人?」.
我記得以前寫手冊.
寫了甚麼保安之類.
但其實我家是漁民.
所以我也不知道在哪裡.
我們也沒有回過家鄉.
我不知道保安是怎樣.
但可能這裡大部分人.
都很清楚自己是台山人.
是否在福建.
很清楚上幾代人發生了甚麼事.
但我自己也不知道.
我只知道太公那一代已經是漁民.
他們都在船上居住.
所以在前幾代人的事.
我也不知道.
所以我在想.
這群以色列人究竟是否知道.
神跟亞伯拉罕,以撒,雅各納的約.
究竟是甚麼.
但無論這群人是否記得.
這段經文也跟我們說.
主導權是神.
是神記得,神沒有忘記.
神沒有忘記.
我不知道你有沒有試過.
信了主後,你會質疑.
或是想過一個問題.
就是神是否聽到我們祈禱.
或是在聽.
是否記得我們的祈禱.
很多時候.
你判斷一個人是否聽你說話.
你也要看他怎樣回應你.

$^{241}$有些回應是很即時的.
但有些可能是隔了一段時間.
或是要等待.
或是有些回應.
即使你知道對方聽了.
但他好像沒有聽入心.
例如….
A跟B的對話.
A說「剛才說聖誕晚餐」.
「你買了嗎?選甚麼餐好呢」.
「我又想吃原味橙雞」.
「但我又想吃照燒雞」.
「究竟選A,B,C這餐哪個好」.
B可能回應.
「我不知道你怎樣回應」.
可能B…這是我幻想出來的對話.
B可能說「雞而已,隨便吧」.
可能你也聽過類似的對話.
「如有女童實屬巧合」.
可能A會期待B說.
「你兩隻雞都想吃」.
「不如我們一人一半吧」.
對吧?比較甜蜜的.
「不如我選A餐,你選B餐」.
「我們分享好嗎」.
可能他期待著這樣.
但B一句「得埋你」.
他馬上停止了你.
對吧?對話很沒趣.
有些對話可能會這樣.
又或是你跟對方說.
「我最近心情真的一般」.
對方可能會回應你.
「是嗎?你在星期五買了甚麼」.
「好,那樣減價也很重要」.
「買了嗎」.
有時對方太快跳過了你所說的話.
給你一個回覆.
可能完全沒有關係.
你聽完後可能心情更差.

$^{281}$其實一個記得你.
或是一個合適的回應.
除了反映對方是否在聽.
其實是反映對方是否在關心.
是否在關心你說甚麼.
我剛開始信主的時候.
其實我第一個很恆常參加的聚會.
不是教會,而是祈禱會.
因為我每次都很期待聽到人們.
如何經歷上帝.
如何讓他們的祈禱被應運.
所以我每次都會參加.
但信了主十多年.
我慢慢開始明白.
可能我們信仰歷程中.
也有不同的時期.
可能有些祈禱.
神很快地聽了,回應我們.
但有些祈禱.
我們祈禱了很久.
好像上帝也沒有回應.
我們也不知道如何祈求.
可能心底裡很怕上帝不應運祈禱.
不過今天這段經文好像再提醒我們.
最終我們的確不知道上帝如何回應我們.
但原來上帝是聽到的.
上帝是聽到的.
祂會顧念.
祂很在乎我們.
但是敘事者好像讓我們明白.
神親眼看見了每一個以色列人.
每一個的經歷.
每一個被欺負.
每一個被壓榨的經歷.
神都看見.
而且祂不是一個觀察者.
祂遠遠地看著.
是神在裡面.
其實看見.
要看見也是花很多心力.

$^{321}$在這幾個星期裡.
看很多的新聞.
其實心裡也很累.
我不知道你有沒有一個感受.
就是其實要看見所有的故事.
看見所有的經歷.
是很花心力的.
昨天芝加哥香港人這個群體.
在這裡舉辦了一個.
大埔魂福縣的悼念會.
過程中有一個美術治療師.
帶給我們一個工作坊.
教我們如何放鬆.
其實在放鬆的過程中.
你才會知道自己有多緊張.
或者在過程中.
你才會知道.
原來自己這個星期的呼吸.
的確不是很順暢.
我想在這個看見裡面.
其實大家都很沉重.
過程是不容易的.
不過我們不是神.
神看見每一個人.
但我們只是人.
很有限.
我們也需要照顧好自己.
如果有需要.
可以分享一下.
昨天學了甚麼運動去放鬆.
上帝看見以色列人.
每一個祂都看見.
第四 最後.
最後這句是神知道.
這句話不僅是神看見.
不僅是看見.
頭腦上知道的意思.
因為在這個文字裡面.
它是說一些很了解.
你親密關係的人的那種知道.

$^{361}$你知道你女兒的脾性.
你知道你子女的習慣.
那種很親密的關係的知道.
甚至有祂參與其中的那種知道.
這種知道讓我們感受到.
上帝不僅是一個.
物理上跟我們很靠近的神.
而且祂在心靈上也很靠近我們.
我從小到大都是祖母養大的我.
她每天接我放學.
下午我自己做功課.
她就煮飯.
我記得有一年很深刻的印象.
就是我三年級要寫毛筆字.
要學寫毛筆字.
我不知道怎樣寫.
不知道大家是否學過毛筆字.
也經歷過.
其實在很沒頭緒.
又不知道怎樣開始.
很無助的時候.
我不知為何覺得很大壓力.
我就開始哭.
我祖母剛從廚房走出來發現我.
我哭得更厲害.
當有人看見我時.
更加哭得很厲害.
我就說「我真的不懂怎樣寫」.
但平時她很嘮叨.
她經常學紙.
可能因為傳統或上一代不太懂得溝通.
但那一刻她很溫柔.
她說「祖母,跟你一起寫吧」.
「祖母,跟你一起寫吧」.
我祖母不懂得寫字.
我聽完後覺得挺搞笑.
也呆了一會.
然後她放下手上的東西.
坐下來,捉著我的手跟我一起寫.
其實那一刻我邊寫邊抹眼淚.

$^{401}$但我想那些眼淚.
是由已經很徬徨,很無助的感受.
變成了一種「我知道祖母很疼我」.
我知道自己很無助的時候.
原來有人跟我一起.
他不僅很實際地跟我一起.
他透過這件事告訴我.
他很支持我.
他心靈跟我很靠近.
雖然只是一刻.
但我記憶足夠.
我想也改變了我一些.
除了我喜歡寫毛筆字.
之後參加比賽.
每次我遇到一些很新的事.
或者當我覺得很無助.
不知道該怎麼走下去的時候.
我好像有一種感覺是.
他也跟我一起.
他還在香港.
有時我們在這裡遇到困難.
我也不想跟他說.
但他經常說「快點回來吧」.
但我知道我在這裡遇到困難時.
他也跟我一起.
上帝的知道不僅是知道.
而是他真的跟我們一起.
一起走.
他心靈上很靠近我們.
跟我們一起.
可能我們也會問.
既然上帝聽到以色列人的呼吸.
為何他不出手阻止這些苦難.
你知道嗎.
我也不知道.
我也不敢回答這個問題.
因為我覺得作為人.
我沒有資格代上帝回答.
但聖經經常描寫.
人跟上帝去呼求,哀求.

$^{441}$特別是詩篇,哀歌.
或者在新約的福音書.
我記得Ken之前也說過一篇.
在耶穌釘十字架時.
他臨死前也呼求.
「我的神,為何離棄我」.
神好像沒有在苦難中.
將任何事情停止.
但我們仍然可以向上帝祈求.
現在我在培德中心上班.
每天都幫助不同的人解決問題.
不過我也知道自己很有限.
而且是新來的報道.
有些問題的確無法解決.
很無能為力.
我在那裡無能為力.
但當別人開始說出他們的情況時.
我看到他們的無助.
他們一行一行的眼淚.
那一刻我也知道他們的難處.
有時我也不知道該說甚麼.
只聽他們說.
但可能他們說完之後.
也沒有解決過問題.
不過他們會說「謝謝你聽我說」.
或者他們覺得.
你們的存在已經讓我知道不孤單.
的確,苦難在每個地方都發生.
不僅是香港.
你看看新聞.
其實你看到不同的地方.
都有發生不同的事情.
希臘,韓國,歐洲等地方.
德國等地方.
很多事情都在上演.
一些人的悲劇,苦難.
其實芝加哥….
你看看身邊.
雖然每個人都很好.
但可能大家都會面對一些難處.

$^{481}$最近可能心情很低落.
或者身體不適.
神也在聽我們.
神聽到我們說的話.
而且神記得我們.
神也看到.
神知道.
弟兄姊妹,我想今天這段經文.
神很溫柔地告訴我們.
祂在這裡.
無論甚麼情況,祂都在.
而且祂更重申一次.
究竟祂是一個怎樣的神.
而且祂是一個很願意靠近我們的神.
祂願意靠近我們.
上帝求你.
垂聽我們每一個心裡面的聲音.
多謝你透過這段經文.
跟我們再一次去說.
你究竟是一個怎樣的神.
你有多在意我們.
你有多願意去親近我們.
願你親自安慰我們每一個的心.
繼續帶領我們.
與我們同行.
也幫助我們去學習.
知道怎樣跟我們身邊的人去同行.
主我們多謝你.
祈禱交託奉主明求.
阿們.
這段經文告訴我們.
\newpage



\section{馬太福音 2:1-23}
\label{sec:3dbOIx2ZjUI}
\textbf{The Worship of the Magi (Matt 2) - Rev. Mark Chan 東方博士朝拜主(馬太福音2\_1-23) - 陳崇基牧師}
\newline
\newline
連結: \href{https://youtube.com/watch?v=3dbOIx2ZjUI}{\texttt{https://youtube.com/watch?v=3dbOIx2ZjUI}} ~~~~ 語音日期: 2025-12-15
\newline
\newline
\hyperref[sec:8eaa7ARgWZg]{\small{< < < PREV SERMON < < <}}
~
\hyperref[sec:index_chronic]{\small{[返順時目]}}
~
\hyperref[sec:index_scriptual]{\small{[返順卷目]}}
~
\hyperref[sec:f1NDkt3u9eY]{\small{> > > NEXT SERMON > > >}}
\newline
\newline
馬太福音 2:1-23
\newline
\begin{longtable}{cl}
\hline
\hline
章節 & 經文 (和合本修訂版)\\
\hline
2:1 & \begin{tabularx}{0.7\textwidth}{X} 在希律作王的時候，耶穌生在猶太的伯利恆。有幾個博學之士從東方來到耶路撒冷，說： \end{tabularx} \\ \\ \relax
2:2 & \begin{tabularx}{0.7\textwidth}{X} 「那生下來作猶太人之王的在哪裡？我們在東方看見他的星，特來拜他。」 \end{tabularx} \\ \\ \relax
2:3 & \begin{tabularx}{0.7\textwidth}{X} 希律王聽見了，就心裡不安；耶路撒冷全城的人也都不安。 \end{tabularx} \\ \\ \relax
2:4 & \begin{tabularx}{0.7\textwidth}{X} 他就召集了祭司長和民間的文士，問他們：「基督該生在哪裡？」 \end{tabularx} \\ \\ \relax
2:5 & \begin{tabularx}{0.7\textwidth}{X} 他們說：「在猶太的伯利恆。因為有先知記著： \end{tabularx} \\ \\ \relax
2:6 & \begin{tabularx}{0.7\textwidth}{X} 『猶大地的伯利恆啊，你在猶大諸城中並不是最小的；因為將來有一位統治者要從你那裡出來，牧養我以色列民。』」 \end{tabularx} \\ \\ \relax
2:7 & \begin{tabularx}{0.7\textwidth}{X} 於是，希律暗地裡召了博學之士來，查問那星是甚麼時候出現的， \end{tabularx} \\ \\ \relax
2:8 & \begin{tabularx}{0.7\textwidth}{X} 就派他們往伯利恆去，說：「你們去仔細尋訪那小孩子，找到了就來報信，我也好去拜他。」 \end{tabularx} \\ \\ \relax
2:9 & \begin{tabularx}{0.7\textwidth}{X} 他們聽了王的話就去了。忽然，在東方所看到的那顆星在前面引領他們，一直行到小孩子所在地方的上方就停住了。 \end{tabularx} \\ \\ \relax
2:10 & \begin{tabularx}{0.7\textwidth}{X} 他們看見那星，就非常歡喜； \end{tabularx} \\ \\ \relax
2:11 & \begin{tabularx}{0.7\textwidth}{X} 進了房子，看見小孩子和他母親馬利亞，就俯伏拜那小孩子，揭開寶盒，拿出黃金、乳香、沒藥，作為禮物獻給他。 \end{tabularx} \\ \\ \relax
2:12 & \begin{tabularx}{0.7\textwidth}{X} 因為在夢中得到主的指示，不要回去見希律，他們就從別的路回自己的家鄉去了。 \end{tabularx} \\ \\ \relax
2:13 & \begin{tabularx}{0.7\textwidth}{X} 他們走後，忽然主的使者在約瑟夢中向他顯現，說：「起來！帶著小孩子和他母親逃往埃及，住在那裡，等我的指示；因為希律要搜尋那小孩子來殺害他。」 \end{tabularx} \\ \\ \relax
2:14 & \begin{tabularx}{0.7\textwidth}{X} 約瑟就起來，連夜帶著小孩子和他母親往埃及去， \end{tabularx} \\ \\ \relax
2:15 & \begin{tabularx}{0.7\textwidth}{X} 住在那裡，直到希律死了。這是要應驗主藉先知所說的話：「我從埃及召我的兒子出來。」 \end{tabularx} \\ \\ \relax
2:16 & \begin{tabularx}{0.7\textwidth}{X} 希律見自己被博學之士愚弄，極其憤怒，差人將伯利恆城裡和四境所有的男孩，根據他向博學之士仔細查問到的時間，凡兩歲以內的，都殺盡了。 \end{tabularx} \\ \\ \relax
2:17 & \begin{tabularx}{0.7\textwidth}{X} 這就應驗了耶利米先知所說的話： \end{tabularx} \\ \\ \relax
2:18 & \begin{tabularx}{0.7\textwidth}{X} 「在拉瑪聽見號咷大哭的聲音，是拉結哭她兒女；她不肯受安慰，因為他們都不在了。」 \end{tabularx} \\ \\ \relax
2:19 & \begin{tabularx}{0.7\textwidth}{X} 希律死了以後，在埃及，忽然主的使者在約瑟夢中向他顯現， \end{tabularx} \\ \\ \relax
2:20 & \begin{tabularx}{0.7\textwidth}{X} 說：「起來，帶著小孩子和他母親回以色列地去！因為要殺害這小孩子的人已經死了。」 \end{tabularx} \\ \\ \relax
2:21 & \begin{tabularx}{0.7\textwidth}{X} 約瑟就起來，帶著小孩子和他母親進入以色列地去。 \end{tabularx} \\ \\ \relax
2:22 & \begin{tabularx}{0.7\textwidth}{X} 但是他因聽見亞基老繼承他父親希律作了猶太王，怕到那裡去；又在夢中得到主的指示，就往加利利境內去了。 \end{tabularx} \\ \\ \relax
2:23 & \begin{tabularx}{0.7\textwidth}{X} 他們到了一座城，名叫拿撒勒，就住在那裡。這是要應驗先知所說的話：「他將稱為拿撒勒人。」 \end{tabularx} \\ \\
[1ex]
\hline
\hline
\end{longtable}
$^{1}$- Thank you..
- Good morning and Merry Christmas to all,.
oh, I'm sorry..
Yeah, and today is the Sunday before the Christmas Sunday,.
so we have to go into the Bible story on Christmas..
When we think of the Christmas story,.
the Christmas story, you can see in the Christmas card.
or some of the scenes that you have often encounter,.
we often think that the Christmas story.
is a children's story, right?.
But if you read the Bible story on Christmas,.
you'll find that it's not a kid's story.
because the Christmas story recorded in the Bible.
talks about the political realities.
and the kings of the time..
And today we turn to Matthew 2, chapter two,.
and we're gonna compare two sets of people,.
two sets of characters,.
and see their similarities and differences..
(speaking in foreign language).
Now the first comparison we can make.
is between the Magi and the scribes,.
and our passage is found in chapter two in Matthew..
(speaking in foreign language).
Now chapter two, verse one, let me read..
After Jesus was born in Bethlehem in Judea,.
during the time of King Herod,.
and this is the King Herod the Great,.
the first King Herod in the Bible,.
Magi from the east came to Jerusalem and asked,.
"Where is the one who has been born King of the Jews?".
We saw his star when it rose and have come to worship him..
(speaking in foreign language).
Now you focus on the term Magi, M-A-G-I,.
and who are these people?.
Sometimes we translate it Magi,.
sometimes we translate them as wise men,.
and if you look at the newer translation.
of the Chinese Bible,.
they translate it as astrologers,.

$^{41}$astrologers, which is more accurate, I think,.
and these people are those who study the stars..
They interpret dreams and they see the omen.
and are trying to observe the natural differences.
in the sky, and they came from the east,.
and most probably they came from Babylon.
or Arabia or Persia..
(speaking in foreign language).
Now in the ancient Near East,.
these people, astrologers,.
they believed they could read the signs in the sky,.
the heavens, the heavenly bodies,.
and they tried to do that in order to predict.
the future of the world..
(speaking in foreign language).
Now, in the ancient Near East,.
these people, astrologers,.
they believed they could read the signs in the sky,.
the heavens, the heavenly bodies,.
and they tried to do that in order to predict.
the future or control the future,.
and these are some of the ways the pagans used.
to control the gods and the future,.
and it's pretty similar to modern-day folk religion.
or superstitions, fortune-telling or astrology.
to control God, is try to control our future.
by looking at the natural phenomena,.
but if you are reading this story.
from a so-called Jewish biblical perspective,.
you'll find that God does not do that..
God does not allow us to control him..
God actually reveals himself personally,.
and he sends his prophets, his servants,.
and write down scripture to tell us his will,.
so as a Jew, if you read this story,.
you will look down on the pagans,.
the astrologers, the Gentiles,.
and even those who are trying to predict the future,.
and these very people actually, interestingly,.
they travel miles and miles.

$^{81}$in order to worship the newborn king, Jesus..
(speaking in foreign language).
See, this is the first group of people,.
the magi, the pagans, they know almost nothing..
They only saw a star,.
and they came to worship Jesus, the newborn king..
However, the second group is the scribes,.
and we can find their description in verses three and four..
When King Herod heard this,.
about the magi seeing a star and came to worship Jesus,.
King Herod was disturbed,.
and all Jerusalem with him were disturbed..
When he had called together.
all the people's chief priests and teachers of the law,.
and King Herod asked them where the Messiah was to be born..
(speaking in foreign language).
Now, notice how King Herod respond to the star,.
or the magi coming to Jerusalem..
King Herod was actually a Roman client king.
appointed by the Roman government.
to rule over the land of the Jews, Judea,.
and he was only half Jewish,.
and he was not fully Jewish,.
but when he heard about the newborn king of the Jews.
has been born, he became really, really troubled.
and anxious because he felt that his power.
as the king of the Jews was threatened..
(speaking in foreign language).
So here the verses tells us the king.
gather all the experts in scripture,.
the chief priests and the scribes in Jerusalem,.
because they were experts and asked them,.
these scholars, and asked them the question,.
where the Messiah was to be born?.
(speaking in foreign language).
And if I asked you the same question,.
you probably don't know the answer, right?.
Where was the Messiah to be born?.
Where?.
And where do you find this scripture in Old Testament?.

$^{121}$Which part of Old Testament talks about it?.
You don't know, right?.
You don't know..
So the verses tell us in verse five,.
they immediately reply,.
they meaning the scribes and the experts in scripture,.
they said in Bethlehem in Judea,.
for this is what the prophet has written,.
but you Bethlehem in the land of Judah.
are by no means least among the rulers of Judah,.
for out of you will come a ruler.
who will be the shepherd of my people Israel..
(speaking in foreign language).
Now the Jewish scribes knew scripture really well, right?.
They recall, okay, there's the Old Testament passage.
that talk about where the Messiah was to be born..
So they knew scripture well.
and they answer confidently and immediately..
And they quote from, actually from a book.
in the Old Testament called Micah chapter five..
They quote from Micah chapter five and says,.
which in that passage talks about God promising.
a new beginning to his people, his kingdom..
And the new king is not from Jerusalem, the capital,.
but from the hometown of David, Bethlehem..
(speaking in foreign language).
So the scribes, the Jewish scribes,.
they were the ones who knew scripture best..
And they were even teachers of the Bible..
They are the PhDs of the day, PhD in the Bible..
And then geographically, they were closest to Bethlehem..
It's only two to three hours walk.
from Jerusalem to Bethlehem..
Yet, did they go see Jesus, the new born king?.
Did they go worship him?.
No, they didn't even move..
(speaking in foreign language).
So the closest ones did not worship Jesus,.
but the farthest one, the Magi came and worshiped Jesus..
I think this is part of the human nature, right?.

$^{161}$Those who are closest to the church, they're often late..
Right?.
Those who are farther, they're often,.
they're often late..
And so, the closest ones,.
the Magi, they're often late..
And so, the closest ones,.
those who are farther, they often show up earlier..
So I think that's interesting, you know,.
because those who are closest arrive the latest,.
because sometimes we take things lightly.
and make the least preparation..
(speaking in foreign language).
And verses seven to 11 talks about how the Magi.
follow the star and they rejoice greatly.
and enter the house and bow down..
You see all the verbs, they bow down,.
they enter the house, they bow down,.
they worship the child, they present the gifts,.
go, frankincense and myrrh to the new born king..
(speaking in foreign language).
So, the Magi, they're often late..
And so, the closest ones,.
the Magi, they're often late..
(speaking in foreign language).
And then the scripture continues to describe this story..
They went on their way,.
and the star they had seen in the east.
went ahead of them until it stopped.
over the place where the child was..
When they saw the star,.
they were overjoyed on coming to the house,.
they saw the child with his mother Mary.
and they bowed down and worshiped.
and opened the treasures and present the new born king.
with gifts of gold, incense and myrrh..
(speaking in foreign language).
, so you see, astrologers, superstitious people,.
the Gentiles, the pagans, they know almost nothing,.
very little about God, and they only saw a star in the sky,.

$^{201}$yet they took action and traveled miles and miles.
and they became the ones who truly understood.
what worshiping God is..
(speaking in foreign language).
And who, let me ask you,.
who has been a Christian for at least five years?.
Can you raise your hand?.
All right, we get so many of those, yeah..
You know, I've been to church, you know,.
been there, done that, and you know,.
I really feel that sometimes new believers,.
not those who raised their hand,.
new believers always come to church with reverence.
and respect and sincerity..
They come on time and they come prepared..
And we should ask ourselves, is our heart toward worship.
more like the scribes of the Jews.
or like the simple and humble Magi?.
(speaking in foreign language).
And the gospel of Matthew present Jesus as the King.
of all nations..
You know, at the end of the gospel,.
Jesus died and he rose again,.
and he is the one who is worthy of worship.
of all nations and all people..
And he calls the disciples to make disciples.
of all nations and our worship should be.
worthy of worship of all nations..
And he calls the disciples to make disciples.
of all nations and our worship should always consider.
whether Jesus is worthy of our service..
(speaking in foreign language).
Okay, so that's the first comparison.
in the story of Christmas..
The second comparison is between two kings..
One king is called Herod the Great, he's the earthly king..
The other king is the newborn King Jesus..
(speaking in foreign language).
And you know, King Herod or Herod the Great,.
he was known for many things..

$^{241}$He was actually the greatest builder before modern times..
And he built so many things in the land of Israel.
and Palestine and beyond..
And during his reign, he carried out massive.
construction projects and he built more than 30 fortresses.
and cities throughout Palestine and beyond..
And even today, you know, we go to Israel.
and travelers can visit those sites..
(speaking in foreign language).
And Herod also fought his entire life.
to gain his kingship, to keep his throne..
And because of this, he was extremely, extremely,.
this guy, King Herod, the one on the right,.
on the left, King Herod the Great,.
he was extremely jealous and suspicious.
of all the people around him..
And because he considered every single one of them.
as a threat and he would illuminate them.
to take them out when King Herod heard,.
as you can understand, heard about the Magi.
coming to worship another king of the Jews,.
and he was deeply, deeply disturbed..
(speaking in foreign language).
And especially in his later years, in his last years,.
King Herod was living in the constant fear.
of being overthrown and he killed almost.
every of his friends and relatives,.
his beloved wife, one of his beloved wives,.
and three of his own sons he killed..
And he became really, really sick.
well, toward the end of his life..
And he knew that no one would mourn for his death.
after he died, so he became more jealous and more cruel..
(speaking in foreign language).
And there's a historian called, named Josephus,.
and he talks so much about what happened in those years.
before King Herod died..
And he described the final years of King Herod like this,.
at this passage shown you on the screen..
In the past, King Herod took pleasure.

$^{281}$out of excessive concern for his own life.
and he turned against his own relatives..
But the command he now issued right before he died.
was utterly inhuman..
And Josephus said he ordered, King Herod ordered,.
at the moment of his death,.
one person in each Jewish household.
was to be put to death, so that when he died,.
the entire nation would truly mourn,.
having lost their own loved ones..
(speaking in foreign language).
Now this is the true picture of King Herod in Matthew..
He was a tyrant and not a king who cared about his people..
When King Herod died, he was a tyrant..
And he was a tyrant who did not care about his people..
And he did not care about his people..
When King Herod heard about the new born king,.
of course he was anxious and he interrogated the Magi.
and he asked them questions and command a lot of things..
(speaking in foreign language).
Now, this is the first comparison..
The second comparison is King Herod versus King Jesus..
And Jesus was not just a king..
The king described as born to be king..
He didn't fight to have his kingship..
He was born, granted the kingship..
And he is the savior of the world..
And he carried no weapons, no tanks, no guns..
Yet he overcame the earthly king, Herod..
(speaking in foreign language).
Here, verse 12 continues to talk about the story.
in kind of a humorous way..
It tells us that the Magi was warned in a dream..
So God sent them a dream,.
not to return from Bethlehem to Jerusalem,.
back to King Herod..
So they went home by another route..
(speaking in foreign language).
Remember I talked about Herod the Great..
He had many military installations.

$^{321}$along the eastern border of Roman Empire..
So that's a very interesting detail here.
that is kind of, has a sense of humor in the Bible..
These Magi, they were unfamiliar with the land,.
but they escaped safely from the hands of Herod the Great..
So here it tells you,.
the great power of Herod the Great.
was completely outwitted by God..
(speaking in foreign language).
And you can tell from today's news,.
those who seem to have a lot of power,.
they are actually very fragile..
And in this Christmas story,.
the word most often used to describe Jesus.
is the child, the kid, the kid..
The child Jesus, the child..
Again and again, the Bible story.
describes Jesus as small, weak, vulnerable,.
and he has no power,.
but God used him, used this child,.
to accomplish his work..
(speaking in foreign language).
And after the Magi had gone,.
Matthew continues the Christmas story.
in verses 13 and following..
After the Magi had gone,.
an angel appeared to Joseph in a dream and said,.
"Get up, take the child and his mother.
"and escape to Egypt..
"Stay there until I tell you,.
"for Herod the Great is going to search.
"for the child to kill him.".
So he got up, Joseph got up, and took, who?.
The child, right?.
The defenseless child..
And his mother, during the night,.
and left for Egypt,.
where he stayed until the death of Herod the Great..
So fulfilled what the Lord has said.
through the prophet,.

$^{361}$"Out of Egypt I called my son.".
(speaking in foreign language).
Now, if you know the Bible well,.
you know that the Lord is going to call Joseph,.
and he's going to call his son,.
and he's going to call his son, Joseph..
And he's going to call his son, Joseph..
If you know the Bible well,.
you heard the word, you know, Egypt,.
and out of Egypt,.
you will immediately thought of the Exodus story..
(speaking in foreign language).
Because in the time of Jesus,.
in the time of Moses, many, many years ago,.
you know, there was also a tyrant king,.
Pharaoh himself..
And he, Pharaoh, feared that he'd lose his power.
and he gave the order that all the Hebrew boys,.
again, the child, to be killed, to be born..
And yet, there was one left behind.
and did not get killed..
Who was that?.
Moses, right?.
Moses survived, and God preserved Moses.
and raised Moses up.
and later used Moses to defeat the tyrant, the Pharaoh,.
and then deliver his people out of Egypt..
(speaking in foreign language).
And then here, the Gospel story.
is telling us something very important here..
And Jesus is not just another Moses..
He is the greater Moses..
He is the greater Moses because he's the final deliverer..
He's the final king that will save us all.
and he's the king and savior..
(speaking in foreign language).
You see, hail the great looks powerful..
He looks really powerful..
And Pharaoh also looked powerful,.
but both of them were defeated because of God's work..

$^{401}$God's salvation does not come through tanks,.
armies, military, political power..
It came through a child, a vulnerable child,.
but he was protected by God.
and used by God to save his people..
And that is the king of Jesus we celebrate at Christmas..
(speaking in foreign language).
You see, when we read the Christmas story,.
we realize that it's not a child's story..
It's not a kid's story..
It's not shallow..
It's a story about power, authority, and salvation..
And it asks us the question, asks you the question,.
which king will you trust?.
(speaking in foreign language).
The Christmas story asks you,.
there are two ways of responding to Jesus, the newborn king..
How do you see worship?.
And how do we see authority and power?.
(speaking in foreign language).
And when we think of worship,.
you know, giving the glory to the true king,.
we compare the Magi and the experts of Scripture, the scribes..
And how did the Gentiles, how did the Magi, the pagans,.
come to see the Lord?.
(speaking in foreign language).
First, the Magi came prepared..
What did they bring?.
Gold, incense, myrrh..
Those are the most expensive stuff of the day..
And that's why later tradition, because there were three gifts,.
and we thought of the Magi, there were three of them, right?.
Of course, that's not what Scripture says..
And actually, Isaiah, another prophet,.
had prophesied this moment..
When the Savior, when the king, the Messiah come to us,.
the nations, the prophet says, Isaiah said,.
the nations, the peoples will bring gold, incense,.
and worship the Lord, and God will be pleased..
(speaking in foreign language).

$^{441}$So they came prepared..
Also, secondly, they rejoiced greatly..
You see, worship was not a duty to them, but a joy..
(speaking in foreign language).
And then third, the Magi bowed down,.
and they worshiped the child,.
and they fell to the ground before him,.
and did not worship from a distance..
They humbled themselves completely..
(speaking in foreign language).
And this invites us to examine our own hearts..
When we come to worship God, we're like the scribes,.
you know, proud and confident in their own knowledge,.
but they remain unmoved..
Or are we like the Magi, joyful and humble and ready.
and prepared to give our hearts?.
(speaking in foreign language).
Next, the custom story shows us two different kinds.
of authority, two different kinds of power,.
Herod or Jesus..
(speaking in foreign language).
You see, earthly kings in this portrait.
often brings war, fear, or suffering to the people.
because they only care about themselves..
Their power is maintained by force..
Their power is maintained by violence..
Herod ruled because of killing, intimidation, and control..
(speaking in foreign language).
But the true king, Jesus,.
he came in a completely different way..
You see that Jesus came without weapons,.
without armies and political power..
He came as a defenseless child,.
but he came to save, he came to sacrifice himself,.
and he came to give his life to all of us..
(speaking in foreign language).
And ultimately he went to the cross,.
and he died for us and rose again..
And through his sacrifice, you see, countless of people..
You examine the church history, countless of people,.

$^{481}$millions and tons of millions of people.
have been saved, including all of us.
who proclaim the name of Jesus..
And now God calls you and me to become his disciples,.
calls you and me to use the same way,.
vulnerable, weak, just to love and to seek peace,.
to influence other people, not by force,.
but by the love of the cross..
(speaking in foreign language).
This is the king we worship, and this is how we worship,.
and this is how we follow Jesus as Christians..
(speaking in foreign language).
So remember this in this Christmas, remember this,.
and let us all pray together..
(speaking in foreign language).
Jesus, we want to give you thanks..
We want to worship you, we want to bow down before you..
(speaking in foreign language).
Because you are our king, you are the only king..
(speaking in foreign language).
And Father, we come before you in this Christmas,.
we think of your kingship, how you die and sacrifice.
to save and redeem us..
(speaking in foreign language).
And Father, we want to not only give you thanks,.
but use our lives to worship you..
(speaking in foreign language).
And God, we ask that you will help us.
to live a life of love and peace..
(speaking in foreign language).
In Jesus' name we pray..
(speaking in foreign language).
Amen..
Well, God bless you, and next week, make sure you come here..
\newpage



\section{馬太福音 2:1-23}
\label{sec:f1NDkt3u9eY}
\textbf{東方博士朝拜主 (馬太福音2\_1-23) - 陳崇基牧師}
\newline
\newline
連結: \href{https://youtube.com/watch?v=f1NDkt3u9eY}{\texttt{https://youtube.com/watch?v=f1NDkt3u9eY}} ~~~~ 語音日期: 2025-12-15
\newline
\newline
\hyperref[sec:3dbOIx2ZjUI]{\small{< < < PREV SERMON < < <}}
~
\hyperref[sec:index_chronic]{\small{[返順時目]}}
~
\hyperref[sec:index_scriptual]{\small{[返順卷目]}}
~
\hyperref[sec:4msYGc1nvL4]{\small{> > > NEXT SERMON > > >}}
\newline
\newline
馬太福音 2:1-23
\newline
\begin{longtable}{cl}
\hline
\hline
章節 & 經文 (和合本修訂版)\\
\hline
2:1 & \begin{tabularx}{0.7\textwidth}{X} 在希律作王的時候，耶穌生在猶太的伯利恆。有幾個博學之士從東方來到耶路撒冷，說： \end{tabularx} \\ \\ \relax
2:2 & \begin{tabularx}{0.7\textwidth}{X} 「那生下來作猶太人之王的在哪裡？我們在東方看見他的星，特來拜他。」 \end{tabularx} \\ \\ \relax
2:3 & \begin{tabularx}{0.7\textwidth}{X} 希律王聽見了，就心裡不安；耶路撒冷全城的人也都不安。 \end{tabularx} \\ \\ \relax
2:4 & \begin{tabularx}{0.7\textwidth}{X} 他就召集了祭司長和民間的文士，問他們：「基督該生在哪裡？」 \end{tabularx} \\ \\ \relax
2:5 & \begin{tabularx}{0.7\textwidth}{X} 他們說：「在猶太的伯利恆。因為有先知記著： \end{tabularx} \\ \\ \relax
2:6 & \begin{tabularx}{0.7\textwidth}{X} 『猶大地的伯利恆啊，你在猶大諸城中並不是最小的；因為將來有一位統治者要從你那裡出來，牧養我以色列民。』」 \end{tabularx} \\ \\ \relax
2:7 & \begin{tabularx}{0.7\textwidth}{X} 於是，希律暗地裡召了博學之士來，查問那星是甚麼時候出現的， \end{tabularx} \\ \\ \relax
2:8 & \begin{tabularx}{0.7\textwidth}{X} 就派他們往伯利恆去，說：「你們去仔細尋訪那小孩子，找到了就來報信，我也好去拜他。」 \end{tabularx} \\ \\ \relax
2:9 & \begin{tabularx}{0.7\textwidth}{X} 他們聽了王的話就去了。忽然，在東方所看到的那顆星在前面引領他們，一直行到小孩子所在地方的上方就停住了。 \end{tabularx} \\ \\ \relax
2:10 & \begin{tabularx}{0.7\textwidth}{X} 他們看見那星，就非常歡喜； \end{tabularx} \\ \\ \relax
2:11 & \begin{tabularx}{0.7\textwidth}{X} 進了房子，看見小孩子和他母親馬利亞，就俯伏拜那小孩子，揭開寶盒，拿出黃金、乳香、沒藥，作為禮物獻給他。 \end{tabularx} \\ \\ \relax
2:12 & \begin{tabularx}{0.7\textwidth}{X} 因為在夢中得到主的指示，不要回去見希律，他們就從別的路回自己的家鄉去了。 \end{tabularx} \\ \\ \relax
2:13 & \begin{tabularx}{0.7\textwidth}{X} 他們走後，忽然主的使者在約瑟夢中向他顯現，說：「起來！帶著小孩子和他母親逃往埃及，住在那裡，等我的指示；因為希律要搜尋那小孩子來殺害他。」 \end{tabularx} \\ \\ \relax
2:14 & \begin{tabularx}{0.7\textwidth}{X} 約瑟就起來，連夜帶著小孩子和他母親往埃及去， \end{tabularx} \\ \\ \relax
2:15 & \begin{tabularx}{0.7\textwidth}{X} 住在那裡，直到希律死了。這是要應驗主藉先知所說的話：「我從埃及召我的兒子出來。」 \end{tabularx} \\ \\ \relax
2:16 & \begin{tabularx}{0.7\textwidth}{X} 希律見自己被博學之士愚弄，極其憤怒，差人將伯利恆城裡和四境所有的男孩，根據他向博學之士仔細查問到的時間，凡兩歲以內的，都殺盡了。 \end{tabularx} \\ \\ \relax
2:17 & \begin{tabularx}{0.7\textwidth}{X} 這就應驗了耶利米先知所說的話： \end{tabularx} \\ \\ \relax
2:18 & \begin{tabularx}{0.7\textwidth}{X} 「在拉瑪聽見號咷大哭的聲音，是拉結哭她兒女；她不肯受安慰，因為他們都不在了。」 \end{tabularx} \\ \\ \relax
2:19 & \begin{tabularx}{0.7\textwidth}{X} 希律死了以後，在埃及，忽然主的使者在約瑟夢中向他顯現， \end{tabularx} \\ \\ \relax
2:20 & \begin{tabularx}{0.7\textwidth}{X} 說：「起來，帶著小孩子和他母親回以色列地去！因為要殺害這小孩子的人已經死了。」 \end{tabularx} \\ \\ \relax
2:21 & \begin{tabularx}{0.7\textwidth}{X} 約瑟就起來，帶著小孩子和他母親進入以色列地去。 \end{tabularx} \\ \\ \relax
2:22 & \begin{tabularx}{0.7\textwidth}{X} 但是他因聽見亞基老繼承他父親希律作了猶太王，怕到那裡去；又在夢中得到主的指示，就往加利利境內去了。 \end{tabularx} \\ \\ \relax
2:23 & \begin{tabularx}{0.7\textwidth}{X} 他們到了一座城，名叫拿撒勒，就住在那裡。這是要應驗先知所說的話：「他將稱為拿撒勒人。」 \end{tabularx} \\ \\
[1ex]
\hline
\hline
\end{longtable}
$^{1}$我跟傑少一樣,我有一點作病.
話說上星期中的時候.
你知道我兒子在西郊找到工作,他不是在家裡住.
他在別人家裡住,租了一個地方.
他在星期中發訊息給我們,說他病得很嚴重.
他說他想回家.
他休息了幾天,回家了.
我心想,你有病,你回家?.
你應該留在別人家裡.
真的很孝順,他病了,想回家.
不知道怎麼回答他,結果他真的回家了.
戴著口罩,病得飛起來,今天教會去不了.
所以覺得很有家教.
我去拿杯水,害我今天有點氣咳.
沒辦法,太有家教.
說笑就說笑,大家保重身體.
因為我們這裡芝加哥特別寒冷.
我們要說聖誕節的故事.
聖誕節的故事是說基督的降生.
在聖經裡只有兩個地方.
如果第一次接觸聖經的話.
去看聖誕的故事只有兩個地方.
新約聖經裡的第一本書叫《瑪太夫》.
第三本書叫《路加夫》.
說到基督的降生的故事,一些情節等等.
我們通常看這個故事的時候.
會想起一些聖誕卡.
或者在城市裡走,或者看到一些相片.
有些比較溫馨的畫面.
我們今天想看這段聖經.
《瑪太夫》第二章,我們看看.
怎麼去看聖誕的故事.
我相信作者是特意將某些東西對照.
這段聖經《瑪太夫》第二章,一至十二節.
說到東方的博士朝見新生的丹夏來造猶太王的耶穌基督.
這是《瑪太夫》第二章.
對照了有兩組人.
第一組人就是博士和所謂的經學家.
(Powerpoint不可行,為什麼呢?).
OK,Here we go.

$^{41}$所謂博士和經學家之間的對比.
這個故事,我們試試用一個人物的方法來看這段聖誕的故事.
裡面說到,首先兩組人.
第一組叫做博士,Magi.
第二組叫做Scribes.
即是當時的經學家或者法師等等.
經文是怎樣說呢?.
二章第一至第二節.
當大希律王,有好幾個希律王,這是第一個.
並且叫做大希律王.
希律王的時候,耶穌生在猶太的伯利恆.
在Bethlehem.
有幾個博士從東方來到耶路撒冷說.
那生下來作猶太人之王在那裡.
我們在東方看見他的聖,特來拜他.
這裡說到,在大希律王執政的時候.
耶穌基督出生在猶太地的伯利恆.
然後說到東方一個遙遠的地方.
大概是今天伊拉克的地方.
巴比倫的地方.
有幾個所謂的博士從東方來到.
為什麼呢?因為他們純粹是看見一顆星.
一個特別的啟示,一個很小的特別的啟示.
然後他們就遠路的來到耶路撒冷.
來到去找這個所謂生下來作猶太人之王的耶穌基督.
博士是什麼呢?.
博士,我們英文譯做Magi或者Wiseman.
這裡說了幾個博士.
這個新的中文翻譯是挺好的,幾個博士.
雖然我兒子病了.
他昨天不知道什麼事跟我聊天.
戴著口罩跟我聊天.
他說是不是有三個博士.
我說你這麼有家教你都不記得.
聖經裡面沒有說三個博士.
聖經裡面從來沒有說三個博士.
只不過是眾數幾個博士.
因為他看見這幅圖.
為什麼不是三隻駱駝呢?.
我說其實是有幾個,可能是一群人.

$^{81}$眾數,有些博士從東方來到耶路撒冷.
博士是什麼呢?.
有些翻譯做博士我們就想起PhD.
有博士學位的人.
其實不是的.
當時近東的世界裡面.
所謂的博士Magi.
其實他們就是一些上升的人.
或者是看天象的人.
當時近東的世界裡面.
最流行或者最看得多.
這些上升,占卜,搜查.
這些說明書是從哪裡來的?.
是從卡比倫那裡來的.
所以這群人從小到大.
他們在宮廷裡面學習.
這些天文,地理.
這群所謂當時的學者.
他們做這件事是做什麼的呢?.
是幫宮廷打工的.
就是告訴宮廷,朝廷.
有什麼空少或者什麼好的兆頭等等.
有一點像天理報告.
或者他們看的那些運程.
這群就是所謂的博士.
或者新日本翻譯為占星家.
原意就是這樣.
他們研究通曉一些天文,地理等等.
你看這個故事的時候.
如果我們是一個.
所謂猶太人認為的外邦人去看.
就不是很詫異.
外邦人看回外邦人.
他們大家都是迷信的人.
有誰覺得卡比倫寫.
他們是看通聖.
是覺得很詫異的?.
不會的.
但如果你是猶太人的讀者.
猶太人的讀者.

$^{121}$如果看到這段聖經的時候.
他們就會覺得很奇怪.
為什麼呢?.
因為這群占星迷信的人.
去看天象,預測運程等等.
這群人對於猶太人來說.
是一些所謂的Pagans.
他們想用人的方法來控制上帝.
或者控制未來.
在猶太人的信仰裡面.
是絕對接受不了的.
為什麼呢?.
因為上帝是親口說話的上帝.
上帝所說話.
或者他給我們生命裡面的原則.
給我們生命裡面的教導.
是上帝直接給我們有律法的.
親自啟示的.
他拆敗他的僕人先知等等.
來講清楚他的說話.
因為猶太人在聖經裡面說的上帝.
因為上帝是清清楚楚地告訴你.
你的生命裡面的方向,價值等等.
所以猶太人是看不起.
看不起這群外邦人.
這群人是他們會嘲笑的.
看不起的.
但是在這個故事裡面.
這群博士或者占星家.
他們純粹看到一顆星.
一個特別的兆頭.
只是靠著這麼小的啟示.
沒有什麼清楚的內容.
他們就跟著這顆星.
然後去找剛剛出生的耶穌.
為了敬拜他.
這是第一個角色.
就是這群博士.
另外就是這群的經學家.
因為這些博士來到.

$^{161}$麻煩幫我來一段.
來到耶路撒冷的時候.
希律王就知道了.
因為這群博士來到耶路撒冷宮廷的時候.
當然他們就是找皇帝.
宮廷裡面出生了一個皇帝.
然後他們想去朝拜他.
或者送上禮物等等.
他們第一個地方就是去京都.
就是去當時巴勒斯坦地.
耶路撒冷的京都.
他們第一個當然要見當時的皇帝.
就是大希律王.
大希律王聽見了就心裡面不安.
耶路撒冷合成的人也都不安.
他就召齊了祭司長和民間的文士.
那些scribes或者經學家.
問他們說基督當生在何處.
這裡說到這些經學家,民士或者祭司長等等.
就被召集進去大希律的宮廷.
大希律是甚麼人呢.
大希律王其實就是半個猶太人.
他本身是以東人.
因為以前馬加比時代.
以東地方被佔領了.
所以這群以東人被逼做了半個猶太人.
所以大希律王聽見一件這樣的事.
有一個新生的猶太人的王會出現.
大希律王作為一個半個猶太人.
他都知道猶太人的律法和經典.
他都清楚或者聽聞過.
舊約聖經裡面有預言過.
一位新生王出來了.
但是他知道一點不知道一點.
所以他就去問當時的PhD.
問那些經學家或者民事,祭司長.
叫他們查考一下聖經.
聖經裡面哪裡有預言有一個新生王出現.
並且聖經裡面有沒有說他在哪裡出生.
所以他們被召入宮的時候.

$^{201}$就去到第五節.
他們就說.
他們回答說.
經學家就回答.
在猶太的伯利恆.
因為先知記住.
記住什麼.
猶大地的伯利恆.
你在猶大諸城中並不是最少的.
因為將來有一位君王要從你那裡出來.
牧養我的以色列民.
這些經學家或者讀聖經讀得很熟的.
他們一看就看到.
就在哪裡出現.
他說先知記住.
他沒有說哪位先知.
你記不記得哪位.
記不記得哪位.
是不是以賽亞.
是不是耶利米.
是不是達爾尼.
哪位先知說.
尼加是誰.
尼加是先知.
你不是很清楚吧.
因為你不是經學家.
那些熟悉聖經的一翻就知道.
尼加書.
你有沒有看過.
看都沒有看過.
因為這些聖經的專家.
一翻就翻開.
他們就說.
當然是新生王在哪裡出現.
伯利恆.
伯利恆在什麼地方.
本身大衛.
你記得舊約裡有一個很出名的王.
叫大衛王.
King David.

$^{241}$他的祖家或者他的官籍.
直貫.
他的鄉下.
就是在伯利恆.
所以一翻就翻到尼加書裡面.
其中一章第五章.
就說到新生王會出現.
大概在尼加書裡預言.
一個彌賽亞出現.
一個新生王出現.
上文下理大概是上帝說.
給當時的貝約的以色列國聽.
我將會有一個新的王.
新的王朝出現.
而這個新的王朝不是在耶路撒冷出現.
是在大衛出生的地方.
會有一個新的大衛的君王會再出現.
這個新的國度會再來到.
這個就是幾百年前.
在尼加書在舊約裡面.
一早預言了.
大衛的王朝會再出現.
大衛的故鄉伯利恆.
會有一個新的王出現.
這個就是第二班的角色.
這裡說到這些經學家.
或者文士.
或者祭司長握他們最熟悉聖經的.
並且是熟得全部的.
並且在地理上.
他們是最接近伯利恆的.
他們在哪裡?.
在耶路撒冷.
新生王在哪裡出現?.
伯利恆.
耶路撒冷到伯利恆.
即是在南面.
大概兩三公里.
如果你走路的話.
應該三個小時就到了.

$^{281}$他們就知道.
聖經一早預言了.
有一個新生王出現.
是上帝派來的.
但他們手指頭都不動.
這個就有很大對比.
那些所謂東方的博士.
猶太人看不起的那班人.
只是靠著一顆星.
去拜這位新生王.
另外對比的就是這班所謂經學家.
這些經學家.
他們知道,通曉得.
連新生王在哪裡出生.
他們都明白.
明白上問下理等等.
他們沒有去找新生王.
反而是這班東方來的博士.
然後故事就說到.
大希律王問過這班聖經學者.
問清楚在伯利恆出生.
然後這些博士來到的時候.
問清楚那顆星是何時出現.
因為他們是老遠地方來的.
查問清楚.
大概應該是兩年前就出現了.
所以這些博士.
旅行了兩年的時間.
就來到耶路撒冷.
來到以色列.
去找新生王.
於是大希律召見了這些博士之後.
他知道是兩年前左右.
他就知道這班新生王.
應該是兩年前就出生了.
所以馬太福音第二章的故事.
當這些博士來到見耶穌的時候.
耶穌大概是幾歲?.
大概一歲多兩歲吧.
所以你看到馬太福音第二章的故事.

$^{321}$應該不是耶穌.
應該不是一個嬰兒躺在那裡.
大概是一兩歲.
那些走走跌跌的在搞事.
我們教會有很多.
大概是這個時間.
所以他們就來找這個小孩子耶穌.
然後到第九節.
希律王知道大概新生王何時出現.
其實就是想殺他.
到第九節.
他們聽見希律王的話就去了.
他們是博士.
博士聽見王的話就去了.
忽然在東方所看到的那顆火星在前面引領他們.
一直走到小孩子所在的地方之上停住了.
他們看見那顆星就非常歡喜.
盡了防止.
看見小孩子和他母親瑪利亞.
他們就俯伏在那小孩子揭開寶盒.
拿出黃金,乳香,抹藥.
作為禮物獻給這位小孩子.
這就是為何後期的人去想像聖誕節的時候.
有三個博士.
禮物當然每人一份.
沒理由一個人拿三份禮物.
有九份.
所以三份禮物出現就認為有三個博士.
當然不一定是這樣.
無論如何他們將最貴重的東西獻上.
然後去敬拜,俯伏在小孩子的前面.
大家留意.
這班猶太人看不起的人.
對上帝根本是一知半解.
憑著一顆星走幾百里路.
反而在這個故事裡面.
他們是最懂聖誕的人.
他們最明白聖誕節是一回甚麼事.
就是這位生下來作為我們王的耶穌.
是值得我們去敬拜.

$^{361}$居然在這個故事裡面.
是一班本身迷信的人.
沒有看過聖經.
不清楚基督教信仰.
不清楚聖經信仰.
他們來去下拜.
並且預備好去敬拜這位主.
那些通曉聖經的反而.
所謂take it for granted.
完全是不理會的.
有時我覺得在教會裡面.
最教導我如何做基督徒的人.
就是一些未信或剛信的人.
對嗎.
他們最讓我看到.
到底甚麼叫尋找上帝.
甚麼叫敬拜上帝.
真人.
看見一些初信或準備去受洗.
雖然他們不知很多.
我們上星期正在預備.
有幾個弟兄姊妹.
他們去受浸.
即是下星期他們受浸.
每逢受浸前我們都會做少少面試.
網上或實體做.
我們就問一下.
舊的一套叫做查問信德.
即是了解他信的是否真確.
以及基本上他能否掌握到.
其實問的問題是很簡單的問題.
他們就去回答.
其實答案亦非常簡單.
但你看見他們去到預備寫見證.
在他們的生命裡面.
如何遇上這位救主等等.
你聽回他們的故事.
聽回他們的故事的時候.
你就發現他們是真心去相信.
真心去認定耶穌是他們的主.

$^{401}$相反.
有時一些通曉聖經的人.
回教會很久的人.
以為自己知道很多的人.
那班人反而覺得耶穌是他們的主.
只不過是生命裡面的闊綠.
這是一個很大的諷刺.
真心去敬拜.
真是明白敬拜是甚麼意思.
是一班這樣的人.
在《馬太福音》裡面.
由開始到結束.
28章的福音書.
書信裡面說甚麼.
他說這位耶穌是萬國的主宰.
是萬國的君王.
在《馬太福音》最後結束的時候.
死而復活的主.
他頒佈大使命.
他說他是那位拿了天上地下.
所有權柄的主宰.
你們要使萬民作他的門徒.
去跟隨他 去遵行他.
其實他說這位耶穌是怎樣的耶穌.
是配得萬國萬民來去敬拜.
並且呼召萬國萬民萬族.
來去敬拜他.
這位主宰.
這位是真神.
我們要這樣回應.
看見耶穌是值得我們去敬拜.
敬拜不是唱歌就算了.
不是坐著就算了.
不是聽著就算了.
而我們離開的時候.
我們明白這位是我們的主宰.
我們要降伏在他面前.
去敬拜他 去跟隨他.
去聽他的說話.
這是第一組有的對比.

$^{441}$博士和誰.
經學家.
第二組的對比.
就是兩個的王.
一個叫大希律王.
一個叫耶穌.
剛生下來.
作猶太人的王.
這位是耶穌.
第一個叫大希律.
之前我們說過為什麼叫大希律.
這是聖經裡好幾個希律王裡.
第一個.
並且我們後世人叫他大希律王.
Herald the Great.
因為他是建築師.
在執政期間大興土木.
建造了很多碉堡 城堡等等.
今天去以色列或不同的地方.
甚至是希臘.
都有他的遺跡.
這個叫做Herald the Great.
遊客今天都可以看到.
2000年之前.
他所建造的東西.
所以為什麼我們叫他大希律.
這個人正如我所說.
他本身不是一個完全的猶太人.
本身是以東人.
馬加比時代的以東地被佔領.
猶太人逼那些以東人做猶太人.
受割禮等等.
所以他是那種後代.
半個猶太人.
但這個大希律是很懂得玩政治的.
他就討好當時的羅馬朝廷.
然後他拿了巴勒斯坦地來管.
所以他也可以叫做猶太人的王.
但他是用他一生打生打死.
才能夠拿到這個王位.

$^{481}$但這個大希律.
除了聖經裡當然記載他很少東西.
但其實在經外的文件裡.
記載這個人很多東西.
除了挖到他做的東西之外.
在文件裡.
有些人記載的裡面.
都會告訴我們這個人是怎樣的人.
大希律其中一個記載他最多的.
就是約瑟夫.佐斯福.
佐斯福是當時第一世紀.
一個很著名的歷史學家.
他寫了很多東西.
特別是講到這個大希律王.
他講到這個大希律王的品格.
他的性情是怎樣.
這裡這一段猶太古史第十七章.
講到大希律還沒死之前.
他做了些什麼.
他知道他快要死了.
因為他患重病.
很痛苦的時候.
大希律王知道那些猶太人.
他管那些猶太人.
是很不喜歡他的.
是憎恨他不得了.
所以他臨死之前要做些什麼呢.
他為了愛自己的性命.
和對付自己的親人.
因為他很怕別人搶他的位置等等.
他在臨死之前.
下令他還沒死之前.
要猶太人每家每戶.
要叫士兵殺一個人.
將一個人處死.
為什麼.
因為將你的家人處死一個.
這樣就哭了.
你會很傷心.
就等他死了.

$^{521}$他就希望全國都不會因為他死而開心.
大家同歸於盡.
這個人可以這麼瘋狂.
就算你不是為了他哭.
你也要哭.
你為了你死去的親人而哭.
他就希望他死的時候.
有所有的人都不會笑著他死.
這個就是大希律真實的寫照.
是一個暴君.
他打生打死就是為了他拿了皇位.
不是一個為百姓著想的政權.
不是為人民著想的君王.
是為自己.
這是第一個角色.
就是大希律.
然後第二個對比當然就是耶穌.
耶穌不僅僅是一個王.
你見到經文裡面怎樣記述他.
他生下來就是猶太人的王.
他不是打生打死才拿到猶太人的王.
這個救世主的身份.
君王的身份.
不是用兵馬不是用安哥的.
由他出生直到他最後.
都是君王.
並且為愛人民而出現的一個君王.
在聖誕這個故事裡面.
他告訴我們這個君王和我們地上的君王.
是完全兩回事.
聖誕裡面紀念的是一位生下來已經作王的耶穌基督.
我們的主宰耶穌基督.
在聖誕裡面我們所紀念的.
是生下來就是有這樣的王權.
並且他配得這樣的王權.
這段故事繼續講下去.
第十二節.
博士敬拜完小孩耶穌.
這裡怎樣記載呢?.
博士在夢中得到主的指示.

$^{561}$不要再在伯利恆回去耶路撒冷見希律王.
從別的路回去.
然後走了.
這個故事.
我相信有些聖經的故事是叫我們看完之後會笑的.
我覺得就是這一段.
因為知道以色列地在羅馬帝國的東方線.
大希律王執政之後.
大興土木已經做了很多城堡去守衛他的東方線.
有很多軍事的儲備在東方線.
根本沒有可能離開大希律王的五致山.
但居然這些博士.
不知頭不知路.
有主的指示.
然後他們就離開了.
就脫身.
有什麼好笑呢?.
在地上看似一些很強大的權力.
大希律王.
地上的一些國權等等.
你以為他有很大的勢力.
但在這個故事裡面他變成了笑話.
為什麼要這樣寫.
我相信聖誕的故事是有權力顛覆的意義.
他告訴我們地上的政權,軍王,漢子是很大的勢力.
但其實在主的眼中.
他們是被嘲笑甚至變成笑柄.
為什麼要點出這件事呢?.
因為真的我們會看見地上的政權.
我們會覺得很吸引.
我們會埋葦黏在一起.
聖誕他告訴我們要粉碎我們對地上的皇權,皇國的幻象.
真正有權柄的人.
是生下來的這個小孩.
在馬太福音第二章記載這個聖誕故事的時候.
最多重複的一個字形容耶穌的是誰?.
是什麼字?.
有沒有留意?.
是小孩子.
小孩子耶穌.

$^{601}$小孩子耶穌與他的母親.
敬拜這位,俯伏敬拜這位的小孩子.
一次又一次,好幾次去到強調這位耶穌.
初生的,只不過是一個小孩子.
小孩子是沒有能力的.
小孩子是沒有軍權的.
你今天甚至看到我們這些小孩子.
最靚的,他們都是不知頭不知腦的.
但他才是有真正值得我們敬拜的皇權.
這個就是聖誕節.
我們看清楚真正上帝的掌權是一回什麼事.
二章十三至十五節.
他們去後,博士離開了.
有主的使者向約瑟夢中顯現.
他說起來帶著小孩子,聽到嗎?.
小孩子,不是母親與小孩子.
是小孩子與他母親,瑪利亞都是第二人.
逃往埃及,住在那裡,等我吩咐你.
因為大希律要尋找小孩子,要除滅他.
約瑟就起來了,夜間就走難.
帶著小孩子和他的母親往埃及去.
住在那裡,暫居在那個地方直到大希律死為止.
應驗了主,直著先知所說的話.
埃及要領出我的意志.
當猶太人聽見這個埃及的故事的時候.
他們自然會聯想起一個舊約裡面最偉大的故事.
就是《出埃及記》.
《出埃及記》發生什麼事?.
《出埃及記》就是在摩西誕生的時候.
以色列民,他們的小孩子.
記得上個星期Alice說到,兩次都是.
說到第一和第二章,就是那些小朋友.
希伯來人的小朋友,他們被法老殘害.
小朋友被強權殘害.
然後有一位救主誕生,就是摩西.
然後摩西就帶領以色列人出埃及.
然後埃及王被上帝的十災所擊敗.
然後有上帝的拯救,上帝的救恩.
這個pattern在舊約裡面經常出現.
在聖誕節也出現.

$^{641}$就是當耶穌這位小孩子被追殺.
希伯來的小孩子被追殺.
他逃往埃及,然後在那裡暫住.
然後又回去,因為暴君死了.
這是應驗出埃及的pattern.
什麼意思呢?.
你看舊約裡面摩西的故事.
以色列人被殘害,然後出埃及.
法老被打敗,然後有真正的救贖出現.
你看見這個故事的時候.
你就會看到耶穌的影子.
你看摩西的故事的時候.
其實你看到的是耶穌的影子.
現在耶穌已經出現了.
當然不是看影子.
是實體已經出現了.
因為真正的君王出現的時候.
他的救恩不是藉著暴力.
不是藉著軍權.
不是藉著政治的權能帶來的.
而是一個軟弱無能的小孩子.
但他是上帝所看重,上帝所保守的.
他最終成為拯救我們的君王.
這是最大的顛覆.
如何去看權柄.
什麼叫做上帝所看重,上帝所用的權柄.
不是飛機大炮.
而是一位軟弱.
按照上帝的旨意來到我們中間.
和我們一起受苦的.
這位耶穌基督.
是他才可以影響這麼多人的生命.
在教會歷史,在我們中間.
去影響我們生命的.
是一位為我們犧牲,願意用他的軟弱.
來犧牲,去拯救我們這位主.
這裡講清楚聖誕的故事.
是對耶穌有兩種回應.
到底我們是那些熟讀聖經.
但是指頭都不動的人.

$^{681}$不去敬拜他.
還是我們已經知道了這位耶穌的地位.
我們就來敬拜他.
我們是那一班甚麼人.
我們是那一班靠籠地上.
強權的.
還是我們真的來到.
這位剛出生.
作為我們王的耶穌面前.
我們都要去服服,去敬拜他.
這值得我們在聖誕節快到了.
只剩下一個星期.
我們要好好準備.
你看到這班博士來到上帝的面前.
來到耶穌的面前.
帶著的是甚麼.
最好的東西.
一個君王配授的東西.
去服服敬拜.
帶著的是甚麼.
他們說大大歡喜.
內心裡面有一種喜樂.
因為這位君王.
是一個永恆的君王.
絕對能夠拯救我們.
他們服服去敬拜.
五體投地.
謙卑的去敬拜.
當然這裡不是只說我們所謂禮儀.
不是只說我們的禮儀服務.
這一個小時.
而是說我們生命.
剛才我報告的時候.
我們聽見有大英姐妹.
我們參與一些聖誕節的報街活動.
傳科門活動.
可能她的生命應該是這樣的.
特別當我們在一個這麼寒冷的地方.
想想一些孤獨的人.
一些沒有人陪伴的人.

$^{721}$如果有人跟他們同行.
聽他們的故事.
去慰問他們.
這個分別是多大.
我經常覺得在美國度假.
有假期.
是很孤獨的時間.
美國真是一個很孤獨的地方.
你在街上走路都沒有人.
因為每個人都開車.
每個人都躲在自己家.
但當一個這麼寒冷.
這麼孤獨的地方.
我們肯去走到一個人的身邊.
去慰問.
甚至可能吃一頓飯.
這是愛.
這是敬拜的生命.
真正的主耶穌基督就是做這件事.
我們作為基督門徒.
我盼望我們都有這樣的心智.
去做這件事.
好不好.
跟家人吃一頓飯.
好不好.
和和氣氣的.
陪伴.
給一些溫暖.
在街上去到不同的地方.
我們都傳福音.
主幫助我們.
向心大道祈禱.
主愛多謝你.
因為你是我們真正的君王.
我們的主耶穌基督.
沒有刀兵.
沒有軍隊.
沒有政權.
他以一個小孩子.
一個軟弱的方式.

$^{761}$來到我們中間.
和我們同行.
和我們受苦.
願意他捨己.
甚至最後走上十字架.
為我們的罪.
為這群這麼不可愛的人.
同行受苦.
受死.
並且復活.
藉著他的犧牲.
拯救無數的人.
包括我們在座.
求你用你的愛來更新我們.
你再帶著你自己犧牲的愛.
去影響其他人.
將溫暖帶給其他人.
將和善帶給其他人.
主你是我們值得敬拜的.
讓我們的生命.
我們的生活.
我們的嘴唇.
你都可以跟隨你.
去敬拜你.
祈禱奉主明救.
元帥弟祝福各位.
\newpage



\section{詩篇 119:1-176}
\label{sec:4msYGc1nvL4}
\textbf{Loving God's Word (Psalm 119) - Rev. Mark Chan 愛神的話語 (詩篇119\_1-176) - 陳崇基牧師}
\newline
\newline
連結: \href{https://youtube.com/watch?v=4msYGc1nvL4}{\texttt{https://youtube.com/watch?v=4msYGc1nvL4}} ~~~~ 語音日期: 2025-12-29
\newline
\newline
\hyperref[sec:f1NDkt3u9eY]{\small{< < < PREV SERMON < < <}}
~
\hyperref[sec:index_chronic]{\small{[返順時目]}}
~
\hyperref[sec:index_scriptual]{\small{[返順卷目]}}
~
\hyperref[sec:wPcDqaSy0yg]{\small{> > > NEXT SERMON > > >}}
\newline
\newline
詩篇 119:1-176
\newline
\begin{longtable}{cl}
\hline
\hline
章節 & 經文 (和合本修訂版)\\
\hline
119:1 & \begin{tabularx}{0.7\textwidth}{X} 行為正直、遵行耶和華律法的，這人有福了！ \end{tabularx} \\ \\ \relax
119:2 & \begin{tabularx}{0.7\textwidth}{X} 遵守他的法度、一心尋求他的，這人有福了！ \end{tabularx} \\ \\ \relax
119:3 & \begin{tabularx}{0.7\textwidth}{X} 他們不做不義的事，但遵行他的道。 \end{tabularx} \\ \\ \relax
119:4 & \begin{tabularx}{0.7\textwidth}{X} 耶和華啊，你曾將你的訓詞吩咐我們，為要我們切實遵守。 \end{tabularx} \\ \\ \relax
119:5 & \begin{tabularx}{0.7\textwidth}{X} 但願我行事堅定，得以遵守你的律例。 \end{tabularx} \\ \\ \relax
119:6 & \begin{tabularx}{0.7\textwidth}{X} 我看重你的一切命令，就不致羞愧。 \end{tabularx} \\ \\ \relax
119:7 & \begin{tabularx}{0.7\textwidth}{X} 我學習你公義的典章，要以正直的心稱謝你。 \end{tabularx} \\ \\ \relax
119:8 & \begin{tabularx}{0.7\textwidth}{X} 我必遵守你的律例，求你不要把我全然棄絕！ \end{tabularx} \\ \\ \relax
119:9 & \begin{tabularx}{0.7\textwidth}{X} 青年要如何保持純潔呢？是要遵行你的話！ \end{tabularx} \\ \\ \relax
119:10 & \begin{tabularx}{0.7\textwidth}{X} 我曾一心尋求你，求你不要使我偏離你的命令。 \end{tabularx} \\ \\ \relax
119:11 & \begin{tabularx}{0.7\textwidth}{X} 我將你的話藏在心裡，免得我得罪你。 \end{tabularx} \\ \\ \relax
119:12 & \begin{tabularx}{0.7\textwidth}{X} 耶和華啊，你是應當稱頌的！求你將你的律例教導我！ \end{tabularx} \\ \\ \relax
119:13 & \begin{tabularx}{0.7\textwidth}{X} 我用嘴唇傳揚你口中一切的典章。 \end{tabularx} \\ \\ \relax
119:14 & \begin{tabularx}{0.7\textwidth}{X} 我喜愛你的法度，如同喜愛一切的財物。 \end{tabularx} \\ \\ \relax
119:15 & \begin{tabularx}{0.7\textwidth}{X} 我要默想你的訓詞，看重你的道路。 \end{tabularx} \\ \\ \relax
119:16 & \begin{tabularx}{0.7\textwidth}{X} 我要以你的律例為樂，我不忘記你的話。 \end{tabularx} \\ \\ \relax
119:17 & \begin{tabularx}{0.7\textwidth}{X} 求你用厚恩待你的僕人，使我存活，我就遵守你的話。 \end{tabularx} \\ \\ \relax
119:18 & \begin{tabularx}{0.7\textwidth}{X} 求你開我的眼睛，使我看出你律法中的奇妙。 \end{tabularx} \\ \\ \relax
119:19 & \begin{tabularx}{0.7\textwidth}{X} 我在地上是寄居的人，求你不要向我隱藏你的命令！ \end{tabularx} \\ \\ \relax
119:20 & \begin{tabularx}{0.7\textwidth}{X} 我時常切慕你的典章，耗盡心力。 \end{tabularx} \\ \\ \relax
119:21 & \begin{tabularx}{0.7\textwidth}{X} 受詛咒、偏離你命令的驕傲人，你已經責備他們。 \end{tabularx} \\ \\ \relax
119:22 & \begin{tabularx}{0.7\textwidth}{X} 求你除掉我所受的羞辱和藐視，因我遵守你的法度。 \end{tabularx} \\ \\ \relax
119:23 & \begin{tabularx}{0.7\textwidth}{X} 雖有掌權者坐著妄論我，你僕人卻思想你的律例。 \end{tabularx} \\ \\ \relax
119:24 & \begin{tabularx}{0.7\textwidth}{X} 你的法度也是我的喜樂，我的導師。 \end{tabularx} \\ \\ \relax
119:25 & \begin{tabularx}{0.7\textwidth}{X} 我的性命幾乎歸於塵土，求你照你的話將我救活！ \end{tabularx} \\ \\ \relax
119:26 & \begin{tabularx}{0.7\textwidth}{X} 我述說我所做的，你應允了我；求你將你的律例教導我！ \end{tabularx} \\ \\ \relax
119:27 & \begin{tabularx}{0.7\textwidth}{X} 求你使我明白你的訓詞，我要默想你的奇事。 \end{tabularx} \\ \\ \relax
119:28 & \begin{tabularx}{0.7\textwidth}{X} 我因愁苦身心耗盡，求你照你的話使我堅立！ \end{tabularx} \\ \\ \relax
119:29 & \begin{tabularx}{0.7\textwidth}{X} 求你使我離開奸詐的道路，開恩將你的律法賜給我！ \end{tabularx} \\ \\ \relax
119:30 & \begin{tabularx}{0.7\textwidth}{X} 我選擇了忠信的道路，將你的典章擺在我面前。 \end{tabularx} \\ \\ \relax
119:31 & \begin{tabularx}{0.7\textwidth}{X} 我持守你的法度；耶和華啊，求你不要叫我羞愧！ \end{tabularx} \\ \\ \relax
119:32 & \begin{tabularx}{0.7\textwidth}{X} 你使我心胸開闊的時候，我就往你命令的道路直奔。 \end{tabularx} \\ \\ \relax
119:33 & \begin{tabularx}{0.7\textwidth}{X} 耶和華啊，求你將你的律例指教我，我必遵守到底！ \end{tabularx} \\ \\ \relax
119:34 & \begin{tabularx}{0.7\textwidth}{X} 求你賜我悟性，我就遵守你的律法，且要一心遵守。 \end{tabularx} \\ \\ \relax
119:35 & \begin{tabularx}{0.7\textwidth}{X} 求你叫我遵行你的命令，因為這是我所喜愛的。 \end{tabularx} \\ \\ \relax
119:36 & \begin{tabularx}{0.7\textwidth}{X} 求你使我的心趨向你的法度，不趨向不義之財。 \end{tabularx} \\ \\ \relax
119:37 & \begin{tabularx}{0.7\textwidth}{X} 求你叫我轉眼不看虛假，使我活在你的道路中。 \end{tabularx} \\ \\ \relax
119:38 & \begin{tabularx}{0.7\textwidth}{X} 求你向敬畏你的僕人堅守你的話！ \end{tabularx} \\ \\ \relax
119:39 & \begin{tabularx}{0.7\textwidth}{X} 求你使我所懼怕的羞辱遠離我，因你的典章本為美。 \end{tabularx} \\ \\ \relax
119:40 & \begin{tabularx}{0.7\textwidth}{X} 看哪，我切慕你的訓詞，求你因你的公義賜我生命！ \end{tabularx} \\ \\ \relax
119:41 & \begin{tabularx}{0.7\textwidth}{X} 耶和華啊，求你使你的慈愛臨到我，照你的話使你的救恩臨到我， \end{tabularx} \\ \\ \relax
119:42 & \begin{tabularx}{0.7\textwidth}{X} 我就有話回答那羞辱我的，因我倚靠你的話。 \end{tabularx} \\ \\ \relax
119:43 & \begin{tabularx}{0.7\textwidth}{X} 求你叫真理的話總不離開我的口，因我仰望你的典章。 \end{tabularx} \\ \\ \relax
119:44 & \begin{tabularx}{0.7\textwidth}{X} 我要常守你的律法，直到永永遠遠。 \end{tabularx} \\ \\ \relax
119:45 & \begin{tabularx}{0.7\textwidth}{X} 我要自由而行，因我尋求了你的訓詞。 \end{tabularx} \\ \\ \relax
119:46 & \begin{tabularx}{0.7\textwidth}{X} 我要在列王面前宣講你的法度，也不致羞愧。 \end{tabularx} \\ \\ \relax
119:47 & \begin{tabularx}{0.7\textwidth}{X} 我以你的命令為樂，這命令是我所喜愛的。 \end{tabularx} \\ \\ \relax
119:48 & \begin{tabularx}{0.7\textwidth}{X} 我向我所愛的，就是你的命令高舉雙手，我也要默想你的律例。 \end{tabularx} \\ \\ \relax
119:49 & \begin{tabularx}{0.7\textwidth}{X} 求你記念你向僕人所說的話，這話使我有盼望。 \end{tabularx} \\ \\ \relax
119:50 & \begin{tabularx}{0.7\textwidth}{X} 你的話將我救活了；這是我在患難中的安慰。 \end{tabularx} \\ \\ \relax
119:51 & \begin{tabularx}{0.7\textwidth}{X} 驕傲的人極度地侮慢我，我卻未曾偏離你的律法。 \end{tabularx} \\ \\ \relax
119:52 & \begin{tabularx}{0.7\textwidth}{X} 耶和華啊，我記念你從古以來的典章，就得了安慰。 \end{tabularx} \\ \\ \relax
119:53 & \begin{tabularx}{0.7\textwidth}{X} 我因惡人離棄你的律法，怒火中燒。 \end{tabularx} \\ \\ \relax
119:54 & \begin{tabularx}{0.7\textwidth}{X} 我在世寄居，以你的律例為詩歌。 \end{tabularx} \\ \\ \relax
119:55 & \begin{tabularx}{0.7\textwidth}{X} 耶和華啊，我夜間記念你的名，我也要遵守你的律法。 \end{tabularx} \\ \\ \relax
119:56 & \begin{tabularx}{0.7\textwidth}{X} 這臨到我，是因我謹守你的訓詞。 \end{tabularx} \\ \\ \relax
119:57 & \begin{tabularx}{0.7\textwidth}{X} 耶和華是我的福分；我曾說，我要遵守你的話。 \end{tabularx} \\ \\ \relax
119:58 & \begin{tabularx}{0.7\textwidth}{X} 我一心懇求你的面，求你照你的話憐憫我！ \end{tabularx} \\ \\ \relax
119:59 & \begin{tabularx}{0.7\textwidth}{X} 我思想自己所行的道路，我的腳步就轉向你的法度。 \end{tabularx} \\ \\ \relax
119:60 & \begin{tabularx}{0.7\textwidth}{X} 我速速遵守你的命令，並不遲延。 \end{tabularx} \\ \\ \relax
119:61 & \begin{tabularx}{0.7\textwidth}{X} 惡人的繩索纏繞我，我卻沒有忘記你的律法。 \end{tabularx} \\ \\ \relax
119:62 & \begin{tabularx}{0.7\textwidth}{X} 我因你公義的典章，夜半起來稱謝你。 \end{tabularx} \\ \\ \relax
119:63 & \begin{tabularx}{0.7\textwidth}{X} 凡敬畏你、守你訓詞的人，我都與他作伴。 \end{tabularx} \\ \\ \relax
119:64 & \begin{tabularx}{0.7\textwidth}{X} 耶和華啊，遍地滿了你的慈愛；求你將你的律例教導我！ \end{tabularx} \\ \\ \relax
119:65 & \begin{tabularx}{0.7\textwidth}{X} 耶和華啊，你照你的話，善待你的僕人。 \end{tabularx} \\ \\ \relax
119:66 & \begin{tabularx}{0.7\textwidth}{X} 求你教我明辨和知識，因我信靠你的命令。 \end{tabularx} \\ \\ \relax
119:67 & \begin{tabularx}{0.7\textwidth}{X} 我未受苦以先曾經迷失，現在卻遵守你的話。 \end{tabularx} \\ \\ \relax
119:68 & \begin{tabularx}{0.7\textwidth}{X} 你本為善，所行的也善；求你將你的律例教導我！ \end{tabularx} \\ \\ \relax
119:69 & \begin{tabularx}{0.7\textwidth}{X} 驕傲的人編造謊言攻擊我，我卻要一心遵守你的訓詞。 \end{tabularx} \\ \\ \relax
119:70 & \begin{tabularx}{0.7\textwidth}{X} 他們的心蒙昧如蒙油脂，我卻喜愛你的律法。 \end{tabularx} \\ \\ \relax
119:71 & \begin{tabularx}{0.7\textwidth}{X} 我受苦是與我有益，為要使我學習你的律例。 \end{tabularx} \\ \\ \relax
119:72 & \begin{tabularx}{0.7\textwidth}{X} 你口中的律法與我有益，勝於千萬金銀。 \end{tabularx} \\ \\ \relax
119:73 & \begin{tabularx}{0.7\textwidth}{X} 你的手造了我，塑造我；求你賜我悟性學習你的命令！ \end{tabularx} \\ \\ \relax
119:74 & \begin{tabularx}{0.7\textwidth}{X} 敬畏你的人看見我就歡喜，因我仰望你的話。 \end{tabularx} \\ \\ \relax
119:75 & \begin{tabularx}{0.7\textwidth}{X} 耶和華啊，我知道你的典章是公義的；你使我受苦是以信實待我。 \end{tabularx} \\ \\ \relax
119:76 & \begin{tabularx}{0.7\textwidth}{X} 求你照著你向僕人所說的話，以慈愛安慰我。 \end{tabularx} \\ \\ \relax
119:77 & \begin{tabularx}{0.7\textwidth}{X} 求你的憐憫臨到我，使我存活，因你的律法是我的喜樂。 \end{tabularx} \\ \\ \relax
119:78 & \begin{tabularx}{0.7\textwidth}{X} 願驕傲的人蒙羞，因為他們無理傾覆我；但我要默想你的訓詞。 \end{tabularx} \\ \\ \relax
119:79 & \begin{tabularx}{0.7\textwidth}{X} 願敬畏你的人和知道你法度的人都歸向我。 \end{tabularx} \\ \\ \relax
119:80 & \begin{tabularx}{0.7\textwidth}{X} 願我的心在你的律例上完全，使我不致蒙羞。 \end{tabularx} \\ \\ \relax
119:81 & \begin{tabularx}{0.7\textwidth}{X} 我渴想你的救恩身心耗盡，我仰望你的話。 \end{tabularx} \\ \\ \relax
119:82 & \begin{tabularx}{0.7\textwidth}{X} 我因渴望你的話眼睛失明，說：「你何時安慰我呢？」 \end{tabularx} \\ \\ \relax
119:83 & \begin{tabularx}{0.7\textwidth}{X} 我雖像煙薰的皮囊，卻不忘記你的律例。 \end{tabularx} \\ \\ \relax
119:84 & \begin{tabularx}{0.7\textwidth}{X} 你僕人的年日有多少呢？你幾時向迫害我的人施行審判呢？ \end{tabularx} \\ \\ \relax
119:85 & \begin{tabularx}{0.7\textwidth}{X} 不順從你律法的驕傲人為我掘了坑。 \end{tabularx} \\ \\ \relax
119:86 & \begin{tabularx}{0.7\textwidth}{X} 你的命令盡都信實；他們無理迫害我，求你幫助我！ \end{tabularx} \\ \\ \relax
119:87 & \begin{tabularx}{0.7\textwidth}{X} 他們幾乎把我從世上除滅；但我沒有離棄你的訓詞。 \end{tabularx} \\ \\ \relax
119:88 & \begin{tabularx}{0.7\textwidth}{X} 求你照你的慈愛將我救活，我就遵守你口中的法度。 \end{tabularx} \\ \\ \relax
119:89 & \begin{tabularx}{0.7\textwidth}{X} 耶和華啊，你的話安定在天，直到永遠。 \end{tabularx} \\ \\ \relax
119:90 & \begin{tabularx}{0.7\textwidth}{X} 你的信實存到萬代；你堅立了地，地就長存。 \end{tabularx} \\ \\ \relax
119:91 & \begin{tabularx}{0.7\textwidth}{X} 天地照你的典章存到今日；萬物都是你的僕役。 \end{tabularx} \\ \\ \relax
119:92 & \begin{tabularx}{0.7\textwidth}{X} 我若不以你的律法為樂，早就在苦難中滅絕了！ \end{tabularx} \\ \\ \relax
119:93 & \begin{tabularx}{0.7\textwidth}{X} 我永不忘記你的訓詞，因你用這訓詞將我救活。 \end{tabularx} \\ \\ \relax
119:94 & \begin{tabularx}{0.7\textwidth}{X} 我是屬你的，求你救我，因我尋求了你的訓詞。 \end{tabularx} \\ \\ \relax
119:95 & \begin{tabularx}{0.7\textwidth}{X} 惡人等著要滅絕我，我卻要揣摩你的法度。 \end{tabularx} \\ \\ \relax
119:96 & \begin{tabularx}{0.7\textwidth}{X} 我看萬事盡都有限，惟有你的命令極其寬廣。 \end{tabularx} \\ \\ \relax
119:97 & \begin{tabularx}{0.7\textwidth}{X} 我何等愛慕你的律法，終日不住地思想。 \end{tabularx} \\ \\ \relax
119:98 & \begin{tabularx}{0.7\textwidth}{X} 你的命令常存在我心裡，使我比仇敵有智慧。 \end{tabularx} \\ \\ \relax
119:99 & \begin{tabularx}{0.7\textwidth}{X} 我比我的教師更通達，因我思想你的法度。 \end{tabularx} \\ \\ \relax
119:100 & \begin{tabularx}{0.7\textwidth}{X} 我比年老的更明白，因我謹守你的訓詞。 \end{tabularx} \\ \\ \relax
119:101 & \begin{tabularx}{0.7\textwidth}{X} 我阻止我的腳走一切邪路，為要遵守你的話。 \end{tabularx} \\ \\ \relax
119:102 & \begin{tabularx}{0.7\textwidth}{X} 我沒有偏離你的典章，因為你教導了我。 \end{tabularx} \\ \\ \relax
119:103 & \begin{tabularx}{0.7\textwidth}{X} 你的言語在我上膛何等甘美，在我口中比蜜更甜！ \end{tabularx} \\ \\ \relax
119:104 & \begin{tabularx}{0.7\textwidth}{X} 我藉著你的訓詞得以明白，因此，我恨惡一切虛假的行徑。 \end{tabularx} \\ \\ \relax
119:105 & \begin{tabularx}{0.7\textwidth}{X} 你的話是我腳前的燈，是我路上的光。 \end{tabularx} \\ \\ \relax
119:106 & \begin{tabularx}{0.7\textwidth}{X} 你公義的典章，我曾起誓遵守，我必按著誓言而行。 \end{tabularx} \\ \\ \relax
119:107 & \begin{tabularx}{0.7\textwidth}{X} 我極其痛苦；耶和華啊，求你照你的話將我救活！ \end{tabularx} \\ \\ \relax
119:108 & \begin{tabularx}{0.7\textwidth}{X} 耶和華啊，求你悅納我口中的讚美為甘心祭，又將你的典章教導我！ \end{tabularx} \\ \\ \relax
119:109 & \begin{tabularx}{0.7\textwidth}{X} 我的性命常在我手掌中，我卻不忘記你的律法。 \end{tabularx} \\ \\ \relax
119:110 & \begin{tabularx}{0.7\textwidth}{X} 惡人為我設下羅網，我卻沒有偏離你的訓詞。 \end{tabularx} \\ \\ \relax
119:111 & \begin{tabularx}{0.7\textwidth}{X} 我以你的法度為永遠的產業，因這是我心中所喜愛的。 \end{tabularx} \\ \\ \relax
119:112 & \begin{tabularx}{0.7\textwidth}{X} 我的心傾向你的律例，謹守到底，直到永遠。 \end{tabularx} \\ \\ \relax
119:113 & \begin{tabularx}{0.7\textwidth}{X} 心懷二意的人為我所恨；但你的律法為我所愛。 \end{tabularx} \\ \\ \relax
119:114 & \begin{tabularx}{0.7\textwidth}{X} 你是我藏身之處，是我的盾牌；我仰望你的話。 \end{tabularx} \\ \\ \relax
119:115 & \begin{tabularx}{0.7\textwidth}{X} 作惡的人哪，你們離開我吧！我要遵守我神的命令。 \end{tabularx} \\ \\ \relax
119:116 & \begin{tabularx}{0.7\textwidth}{X} 求你照你的話扶持我，使我存活，不要叫我因失望而蒙羞。 \end{tabularx} \\ \\ \relax
119:117 & \begin{tabularx}{0.7\textwidth}{X} 求你扶持我，使我得救，時常看重你的律例。 \end{tabularx} \\ \\ \relax
119:118 & \begin{tabularx}{0.7\textwidth}{X} 凡偏離你律例的人，你都輕看他們，因為他們的詭詐必歸虛空。 \end{tabularx} \\ \\ \relax
119:119 & \begin{tabularx}{0.7\textwidth}{X} 你除掉地上所有的惡人，好像除掉渣滓；因此我喜愛你的法度。 \end{tabularx} \\ \\ \relax
119:120 & \begin{tabularx}{0.7\textwidth}{X} 我因懼怕你，肉體戰慄；我害怕你的典章。 \end{tabularx} \\ \\ \relax
119:121 & \begin{tabularx}{0.7\textwidth}{X} 我行公平和公義，求你不要撇下我，交給欺壓我的人！ \end{tabularx} \\ \\ \relax
119:122 & \begin{tabularx}{0.7\textwidth}{X} 求你保證你的僕人得福，不容驕傲的人欺壓我！ \end{tabularx} \\ \\ \relax
119:123 & \begin{tabularx}{0.7\textwidth}{X} 我因盼望你的救恩和你公義的言語眼睛失明。 \end{tabularx} \\ \\ \relax
119:124 & \begin{tabularx}{0.7\textwidth}{X} 求你照你的慈愛待僕人，將你的律例教導我。 \end{tabularx} \\ \\ \relax
119:125 & \begin{tabularx}{0.7\textwidth}{X} 我是你的僕人，求你賜我悟性，得以認識你的法度。 \end{tabularx} \\ \\ \relax
119:126 & \begin{tabularx}{0.7\textwidth}{X} 這是耶和華採取行動的時候，因人廢棄了你的律法。 \end{tabularx} \\ \\ \relax
119:127 & \begin{tabularx}{0.7\textwidth}{X} 所以，我喜愛你的命令勝於金子，更勝於純金。 \end{tabularx} \\ \\ \relax
119:128 & \begin{tabularx}{0.7\textwidth}{X} 你的一切訓詞，在萬事上我都以為正直；我恨惡一切虛假的行徑。 \end{tabularx} \\ \\ \relax
119:129 & \begin{tabularx}{0.7\textwidth}{X} 你的法度奇妙，所以我一心謹守。 \end{tabularx} \\ \\ \relax
119:130 & \begin{tabularx}{0.7\textwidth}{X} 你的話一開啟就發出亮光，使愚蒙人通達。 \end{tabularx} \\ \\ \relax
119:131 & \begin{tabularx}{0.7\textwidth}{X} 我大大張口，呼吸急促，因我切慕你的命令。 \end{tabularx} \\ \\ \relax
119:132 & \begin{tabularx}{0.7\textwidth}{X} 求你轉向我，憐憫我，就像你待那些喜愛你名的人。 \end{tabularx} \\ \\ \relax
119:133 & \begin{tabularx}{0.7\textwidth}{X} 求你用你的言語使我腳步穩健，不容罪孽轄制我。 \end{tabularx} \\ \\ \relax
119:134 & \begin{tabularx}{0.7\textwidth}{X} 求你救我脫離人的欺壓，我要遵守你的訓詞。 \end{tabularx} \\ \\ \relax
119:135 & \begin{tabularx}{0.7\textwidth}{X} 求你使你的臉向僕人發光，又將你的律例教導我。 \end{tabularx} \\ \\ \relax
119:136 & \begin{tabularx}{0.7\textwidth}{X} 我的眼睛流淚成河，因為他們不守你的律法。 \end{tabularx} \\ \\ \relax
119:137 & \begin{tabularx}{0.7\textwidth}{X} 耶和華啊，你是公義的；你的典章正直！ \end{tabularx} \\ \\ \relax
119:138 & \begin{tabularx}{0.7\textwidth}{X} 你所頒佈的法度是公義的，極其可靠。 \end{tabularx} \\ \\ \relax
119:139 & \begin{tabularx}{0.7\textwidth}{X} 我的狂熱把我燒滅，因我敵人忘記你的話。 \end{tabularx} \\ \\ \relax
119:140 & \begin{tabularx}{0.7\textwidth}{X} 你的言語極其精煉，令你僕人喜愛。 \end{tabularx} \\ \\ \relax
119:141 & \begin{tabularx}{0.7\textwidth}{X} 我渺小，被人藐視，卻不忘記你的訓詞。 \end{tabularx} \\ \\ \relax
119:142 & \begin{tabularx}{0.7\textwidth}{X} 你的公義永遠公義，你的律法是確實的。 \end{tabularx} \\ \\ \relax
119:143 & \begin{tabularx}{0.7\textwidth}{X} 我遭遇患難愁苦，你的命令是我的喜樂。 \end{tabularx} \\ \\ \relax
119:144 & \begin{tabularx}{0.7\textwidth}{X} 你的法度永遠公義；求你賜我悟性，使我存活。 \end{tabularx} \\ \\ \relax
119:145 & \begin{tabularx}{0.7\textwidth}{X} 耶和華啊，我一心呼求你，求你應允我！我必謹守你的律例。 \end{tabularx} \\ \\ \relax
119:146 & \begin{tabularx}{0.7\textwidth}{X} 我向你呼求，求你救我！我要遵守你的法度。 \end{tabularx} \\ \\ \relax
119:147 & \begin{tabularx}{0.7\textwidth}{X} 天尚未亮我呼喊求救，我仰望你的話。 \end{tabularx} \\ \\ \relax
119:148 & \begin{tabularx}{0.7\textwidth}{X} 我終夜雙眼睜開，為要思想你的言語。 \end{tabularx} \\ \\ \relax
119:149 & \begin{tabularx}{0.7\textwidth}{X} 求你按你的慈愛聽我的聲音，耶和華啊，求你照你的典章將我救活！ \end{tabularx} \\ \\ \relax
119:150 & \begin{tabularx}{0.7\textwidth}{X} 追逐奸惡的人迫近了，他們遠離你的律法。 \end{tabularx} \\ \\ \relax
119:151 & \begin{tabularx}{0.7\textwidth}{X} 耶和華啊，你就在我身邊，你一切的命令是確實的！ \end{tabularx} \\ \\ \relax
119:152 & \begin{tabularx}{0.7\textwidth}{X} 我從你的法度早已知道，這法度是你永遠立定的。 \end{tabularx} \\ \\ \relax
119:153 & \begin{tabularx}{0.7\textwidth}{X} 求你看顧我的苦難，搭救我，因我不忘記你的律法。 \end{tabularx} \\ \\ \relax
119:154 & \begin{tabularx}{0.7\textwidth}{X} 求你為我的冤屈辯護，救贖我，照你的言語將我救活。 \end{tabularx} \\ \\ \relax
119:155 & \begin{tabularx}{0.7\textwidth}{X} 救恩遠離惡人，因為他們不尋求你的律例。 \end{tabularx} \\ \\ \relax
119:156 & \begin{tabularx}{0.7\textwidth}{X} 耶和華啊，你的憐憫本為大；求你照你的典章將我救活。 \end{tabularx} \\ \\ \relax
119:157 & \begin{tabularx}{0.7\textwidth}{X} 迫害我的、抵擋我的甚多，我卻沒有偏離你的法度。 \end{tabularx} \\ \\ \relax
119:158 & \begin{tabularx}{0.7\textwidth}{X} 我看見奸惡的人就憎惡，因為他們不遵守你的言語。 \end{tabularx} \\ \\ \relax
119:159 & \begin{tabularx}{0.7\textwidth}{X} 你看我何等喜愛你的訓詞！耶和華啊，求你按你的慈愛將我救活！ \end{tabularx} \\ \\ \relax
119:160 & \begin{tabularx}{0.7\textwidth}{X} 你話語的精髓是真實的，你一切公義的典章永遠長存。 \end{tabularx} \\ \\ \relax
119:161 & \begin{tabularx}{0.7\textwidth}{X} 掌權者無故迫害我，然而我的心畏懼你的話。 \end{tabularx} \\ \\ \relax
119:162 & \begin{tabularx}{0.7\textwidth}{X} 我喜愛你的言語，好像人得到許多戰利品。 \end{tabularx} \\ \\ \relax
119:163 & \begin{tabularx}{0.7\textwidth}{X} 我恨惡，憎惡虛假；惟喜愛你的律法。 \end{tabularx} \\ \\ \relax
119:164 & \begin{tabularx}{0.7\textwidth}{X} 我因你公義的典章一天七次讚美你。 \end{tabularx} \\ \\ \relax
119:165 & \begin{tabularx}{0.7\textwidth}{X} 喜愛你律法的人大有平安，任何事都不能使他們跌倒。 \end{tabularx} \\ \\ \relax
119:166 & \begin{tabularx}{0.7\textwidth}{X} 耶和華啊，我仰望你的救恩，遵行你的命令。 \end{tabularx} \\ \\ \relax
119:167 & \begin{tabularx}{0.7\textwidth}{X} 我心謹守你的法度，這法度我極其喜愛。 \end{tabularx} \\ \\ \relax
119:168 & \begin{tabularx}{0.7\textwidth}{X} 我遵守你的訓詞和法度，因我所行的道路都在你的面前。 \end{tabularx} \\ \\ \relax
119:169 & \begin{tabularx}{0.7\textwidth}{X} 耶和華啊，願我的呼求達到你面前，求你照你的話賜我悟性。 \end{tabularx} \\ \\ \relax
119:170 & \begin{tabularx}{0.7\textwidth}{X} 願我的懇求達到你面前，求你照你的言語搭救我。 \end{tabularx} \\ \\ \relax
119:171 & \begin{tabularx}{0.7\textwidth}{X} 願我的嘴唇發出讚美，因為你將律例教導我。 \end{tabularx} \\ \\ \relax
119:172 & \begin{tabularx}{0.7\textwidth}{X} 願我的舌頭歌唱你的言語，因你一切的命令盡都公義。 \end{tabularx} \\ \\ \relax
119:173 & \begin{tabularx}{0.7\textwidth}{X} 求你用你的手幫助我，因我選擇你的訓詞。 \end{tabularx} \\ \\ \relax
119:174 & \begin{tabularx}{0.7\textwidth}{X} 耶和華啊，我切慕你的救恩！你的律法是我的喜樂。 \end{tabularx} \\ \\ \relax
119:175 & \begin{tabularx}{0.7\textwidth}{X} 願我的性命存活，得以讚美你！願你的典章幫助我！ \end{tabularx} \\ \\ \relax
119:176 & \begin{tabularx}{0.7\textwidth}{X} 我走迷了路如同失喪的羊，求你尋找你的僕人，因我不忘記你的命令。 \end{tabularx} \\ \\
[1ex]
\hline
\hline
\end{longtable}
$^{1}$All right guys.
Happy New Year and why don't you say to your neighbor just a little bit, you know one phrase. Happy New Year. Happy.
2026.
Okay, so today's the last Sunday of.
2025 and as you can.
Feel that time is flying.
so fast and we're gonna welcome a new year 2026 and.
every.
Moment of the end of the year turning to a new year.
We going to do something new right we think about something new.
Today's our young woman in the whole you got to read.
She can't end up to the fish or to make a new year the you know, what am I?.
Don't go to woman. She didn't end you see Cindy.
and then let me go back to.
There's something wrong with the PowerPoint and there we go. Okay, there we go.
It's how many of you have gym membership?.
How many of you know?.
Okay, I see two people and then how many you think about new year you're gonna go on diet.
Okay.
Quick thought, you know, maybe oh I need to go on diet or something.
I mean maybe more healthy and we think about new plans and new goals and.
You know how the story goes and by March.
2026 we're gonna give up right and then the diet is broken the planners sitting on the shelf and the problem is sometimes.
we.
Something more of our hearts and our hearts need to change.
Jason found me and also study.
She enjoyed our gym vacay. Why it's until Buchang go to Apple. Yes. I told you woman that in song.
We're not I think that you have been said the issue. Yes, I'm seeing John's here. That was Cindy. It's about who's the kid you?.
You thought of Sandra fan Jason Kaji woman to不用了也是计划也失敗了.
Ji-hwa-bae-ye-bae fan. So Jia-sang when he was a woman who's not a woman each other. So the she may have done that I'd beenso.
Like to turn to a psalm in the Bible, which is the longest.
Psalm in the Bible Psalm 119 and this song is not inviting to all of us to be.
to to do some self improvement the psalm is inviting us for a.
Transformation and it's not trying to encourage us to try harder.
Harder and harder it actually invites us to go deeper to listen.
To kind of have a spiritual diary that we can be shaped and we shaped by the voice of God.
So that's what Psalm 19 is all about. So as we welcome.
2026 God is not calling us to for another, you know new.
You know trying hard habit, but he's calling us to have our life anchor in his word.
She can't even I used to kill him. Who's he? Oh, she wants you want me to work? He's on our ship. We don't want me shooting guns in.

$^{41}$He's not trying to encourage us to try harder and harder it actually invites us to go deeper to listen.
To kind of have a spiritual diary that we can be shaped and we shaped by the voice of God.
2026 God is not calling us to for another, you know new.
She's not trying to encourage us to try harder and harder it actually invites us to have our life anchor in his word.
And you have to understand Psalm 19 is the longest chapter in the Bible and there was total of.
176 verses in this psalm.
If I call you to memorize this, you know, there's no way right? There's no way but you know.
It's the longest chapter, but it's not accidental as you can see.
Even though not all of us know Hebrew.
You've stopped from the right. No from reading from the right to left. That's Hebrew.
You see 119. Do you see that 119 in the picture?.
119 on the right or already right and you see verse one two, three, four five easier.
You see that right?.
Okay, I know you're smart. Okay, and then the first letter you can see is kind of like a cross and that is olive.
That's the first alphabet of the 22 alphabets of biblical Hebrew and then you go to the next.
Picture on the left you see six seven and eight, right?.
You see that don't you?.
Right. It's also the same letter, right?.
so the first eight verses start with the same letter olive and then you see nine to.
Two sixteen don't you right and you see another letter starting the next section and.
Every line of the next eight lines from nine from verses nine to sixteen start with bake.
Which is the second?.
alphabet of the psalm so there's a 22 letters in biblical alphabets and.
each letter.
Has eight verses in Psalm 119 and each verse is a group and begins with the same letter.
So as I show you one to eight is olive and then all the way.
To the last section of this psalm this long psalm you go to the last alphabet, which is a 22nd letter, which is tough.
So the 119th psalm is the longest psalm in the Bible and it has 176 verses.
So we can't recite it but it's not a coincidence that it's a psalm of the Hebrew alphabet.
There are 22 Hebrew alphabets and each letter corresponds to eight verses and each starts with the same letter.
For example, the first to the eighth verses start with olive.
The ninth to the sixteenth verses start with bake and the last letter starts with tongue.
You know you you you.
see what I.
Would oh, that's what I'm doing. Okay.
Bay and I'll try host of all right. Yeah, I thought you didn't you don't know. Okay. Yeah.
Well, that means the psalm is beautifully craft.
we call it a cross the psalm and the poet is trying to work really really hard and.
In a way that will show you that the Word of God is not.
Casual it is intentional and it requires effort and you know beauty and artistic.

$^{81}$Artistry to preserve it and the form of the psalm itself is a piece of masterpiece.
It's an it's an art form. So it's not randomly put together. It's not occasionally put together, but faithfully.
Beautifully put together. It's like English from A to Z.
(translation).
Even though you don't know biblical Hebrew before but now, you know, the first letter is called what?.
The first letter of Hebrew alphabet olive. Yeah, olive and then the second one.
Bay, bay. Yeah, exactly. That's why we got alphabet, right? It goes.
Back to many many years ago the development of alphabetical language.
Then Hebrew is one of the first ones and anyways.
so this psalm is a beautifully put together and it reminds us the first point is that.
Let the Word of God lead you into deeper prayer.
and.
the psalm.
Begins with open my eyes that I may see you so.
verses.
33 it says oh, I'm sorry. I skipped.
Open my eyes that I may behold the wonderful things of the law.
Right psalm 119 is a beautiful law beautiful psalm together and it's not a prayer for a beginner.
it's a prayer for someone who is reading scripture, you know soaking in the Word of God and.
and.
And want to really really really go into scripture.
(translation).
and verses 33 and 34 as you can see on the screen teach me he prays.
Oh Lord teach me to follow your decrees and I will keep them to the end and give me.
Understanding I will keep your law and obey it with all my heart.
So every time you can sense that the poet.
Opens the Bible.
Turn to the pages of the scripture and he asked God to speak first and he is there to listen and the first response.
God's Word produces is not just give us knowledge. It's give us dependency. We want to depend on God and pray to him.
(translation).
As we turn to the Word of God as we listen to the Word of God.
What we're doing right now our posture has always been to obey and to listen you can think of a young child.
You know, we've got some young children here.
And I still remember when my boy was so young so young so small and he was scared at night.
and the room and the house was dark and the shadows to him was like monsters and.
He was in fear and sometimes in the middle of the night.
he will come over and.
Down the hallway left his room and come to my room and then he was scared and I talked to him.
I spoke to him come here. You're safe. You know, daddy is here and.
You know, he would not.

$^{121}$He would just listen to my voice and he would not do some analysis some.
Intellectual exercise to analyze my words. He hears my words and love my words and he will run to me.
(translation).
This is what God wants us to do when we hear his Word. It's not just about information but an invitation.
It's a voice of invitation.
so I want to encourage you this the end of this this year to do something in new year is to.
read the Bible and to read the Word of God and.
It's not so much of you know doing homework and I need to finish the assignment.
But you know you need to have the right posture before God and say to God when we read the Word of God.
Father dear Heavenly Father speak to me and I am listening.
(translation).
And the second point here is to let the Word of God will shape your heart.
So.
The second point here is that you know, we need to let the Word of God to change us.
help.
to to to transform our heart.
verses 11 and 72.
It says I have hidden or store up your word in my heart that I might not sin against you.
and 72 the law the law of God from your mouth is more precious to me than thousands of.
pieces of silver and gold.
(translation).
So, you know the posture here is to store up, you know to to hide it.
Inside the heart and the Hebrew word here implies something very very precious like gold or in your investment.
You put it in your heart. So the psalm is not just poetry is to change our.
Priority let the priority be first the will of God and you can imagine.
again a child.
Growing up in a loving home, right and the father constantly tells him, you know.
You are love and you know kissing him and you know.
Care for him and years later the child goes up and goes to college for example.
and even though he's away from home and no one is checking on him, but.
he has.
stored the word of the father the voice of father inside of him and he knows he's.
truly beloved and.
Those words from the father will guide his decisions as his choices.
(translation).
You know this kind of habit will transform us, you know, I do daily devotion as well.
So you should right you should be doing daily devotion and it won't take a long time.
You know, I want to tell you my habit every single day is that I subscribe to.
devotional channel and I watch two one minute video and.
It's it talks about a verse in the Bible and explain to me.

$^{161}$You know.
What is what the context is and what the meanings are and that's all it takes and then you try to store that.
Information that verse in your heart the rest of the day and the Word of God will gradually transform you and that's what scripture.
healing in our heart will do.
2026 a new year as we welcome it. Please do not let the Word of God stay in your Bible app.
Okay.
You have Bible app, right?.
Don't let it just collect us in your phone.
Use it and let the Word of God guide our conscience and our decision and our reactions.
(translation).
And then the third point here is to let the Word of God guide your daily steps.
I have chosen verses first 105.
It says your Word God is a lamb.
Unto my feet and is the light to my path and you notice the the metaphor here.
It is not a spotlight. You know, the Word of God is not like a spotlight. It is a lamb.
(translation).
So imagine, you know waking up at the middle of the night.
Maybe you go to I don't know that's what I do. I get old I go to pee.
All right in the middle and you know, the room is dark and the house is dark and you don't turn on every light.
But you know you had kind of have your lamb or use your phone.
As a lamb and it doesn't show you the entire house.
You know, the lamb does not have the power to show you everything but it's show you just enough to take the next step safely.
Imagine we wake up in the middle of the night.
Go to the bathroom or drink water and it's pitch black in the house.
We don't turn on all the lights in the house but only the flashlight of the phone.
It can't light up the whole house but it can make us safe to go to the next step.
So this should be how the Word of God works in my in our lives.
The Word of God does not reveal everything about your future about your life or your situation.
But the Word of God will protect you every step of the way just enough for you to take the next step safely.
You know when I prepare this simple sermon I remember.
many many years ago probably.
Actually 30 years ago when I first commit myself to to church ministry and I had.
30 years ago. I had no idea what my future would look like. I had absolutely no idea, you know.
You know, I did not know where I would go and I did not know where I will serve and what are the challenges.
I would face and what kind of sacrifice it would require me to go forward.
So here I am, you know, I'm still you know, you know rocking.
Still working hard and you know, you know God never needed.
To show me everything right? I don't I don't I don't ask God to show me every single thing and.
Everything about my future. I never asked God to do that and God would not do that.
But through his Word he gave me, you know enough guidance to take one step at a time one step at a time.

$^{201}$so.
At the end of the year I look back.
you know, I find that God has been faithful and.
He had been guiding me all along.
So you can tell what I'm done. I'm doing it. I'm not that little either工作.
So you shouldn't from bush you have a you can't go to women.
Did I have a heart I eat booty Buddha daddy more maybe you know we talk on what's happening?.
You know, I can't eat you say in dog and as we welcome the new year 2026.
You know don't just you know come before God and want to try harder.
You know don't just promise to to to be committed. You know, you have to think about this.
You know one song a day one passage a day of one verse a day.
One gospel a day and that is enough, you know read as you pray and pray as you read.
Let the Word of God, you know form your year and guide the the new year for you.
So you are you know, you know, you know, I'm just a show. She's a kind young show. We all wait. You can you call you make?.
Can you can she can make any John for you? So being to be in dog? Oh, yeah, I said the horse without me again.
So as we pray together.
Remember the Word of God will lead us what lead us into prayer and we shape our heart and guide our every single step.
That's that's that's cause it boy.
So you are you know, man, you see I thought you know.
Dear Lord as we step into.
2026 we confess that our hearts are easily distracted.
We trace many different things in the world, but we all often neglect your word.
Today, we do not want to make you know empty promises want to come before you to make a little commitment.
Store your word into our hearts.
Shape our thoughts and our desires and our choices.
Do I don't want to talk to our new year? We're trying to get a single you can see.
Changchun who really the Hawaii今天我們不想在许下空洞的承诺而是要做出聖洁的伪身.
求你將你的話寫在我們的心理塑造我們的思想可望于每天的选擇.
Open our eyes to see the wonders the beauty of your work as psalm 119 shows us.
Teach us your law according to.
Your wisdom and give us understanding may your word become more precious more precious to us than gold and silver.
And truly be a lamb.
to our feet.
like to our path.
So you don't just see him the Svartian with woman can't yet so I'm a kind of you are you don't achieve now.
Tell the d33 jiao dao women you the beauty told us a p.m. 34 gets a woman who sing when you the heart.
So I'm a shin chong be king in gongbao way.
Who's a candy? She's our Jaisalmer being真实成為我們腳前的燈路上的光.
rule ship in the 105 jaisalmer go.
Father make this year this new year different not because our circumstances would change.

$^{241}$But because our hearts have changed and we offer the new year.
2026 to you as an act of worship.
Can't fool to nirvana niancheng way. We young lady. Yeah, who's a human one?.
Jing I've been or she even when the sing I've been alone woman, but I don't know you niancheng any don't look in by.
In Jesus name we pray phone to yes, or the mean dog. Oh, hey man.
Like I bless all of you and remember to read the work out. All right with the world about even one verse a day. All right.
Like a bus you can do something.
Napa say he can eat yet.
Thank you..
(upbeat music).
\newpage



\section{約翰福音 3:16}
\label{sec:wPcDqaSy0yg}
\textbf{神愛世人 (約翰福音3\_16) - 陳崇基牧師}
\newline
\newline
連結: \href{https://youtube.com/watch?v=wPcDqaSy0yg}{\texttt{https://youtube.com/watch?v=wPcDqaSy0yg}} ~~~~ 語音日期: 2025-12-29
\newline
\newline
\hyperref[sec:4msYGc1nvL4]{\small{< < < PREV SERMON < < <}}
~
\hyperref[sec:index_chronic]{\small{[返順時目]}}
~
\hyperref[sec:index_scriptual]{\small{[返順卷目]}}
~
\hyperref[sec:VP0gKdg8cek]{\small{> > > NEXT SERMON > > >}}
\newline
\newline
約翰福音 3:16
\newline
\begin{longtable}{cl}
\hline
\hline
章節 & 經文 (和合本修訂版)\\
\hline
3:16 & \begin{tabularx}{0.7\textwidth}{X} 「神愛世人，甚至將他獨一的兒子賜給他們，叫一切信他的人不致滅亡，反得永生。 \end{tabularx} \\ \\
[1ex]
\hline
\hline
\end{longtable}
$^{1}$這個約翰福音第三章16節剛才看了.
這個幾乎是全世界最廣為人知的經文.
因為宣告一個事實.
就是神愛世人.
並且是一個很震撼的宣告.
約翰福音三章16節是全世界最廣為人知的一節經文.
宣告了一個真理.
神愛世人.
而這其實是一個非常震撼的宣告.
通常我們看到神愛世人.
God so love the world的時候.
我們就覺得有一點溫馨.
因為我們覺得上帝的愛當然是很偉大.
很偉大 包括所有的人.
因為這個世界這麼大.
但上帝卻愛這個世界.
通常當我們讀到神愛世人的時候.
我們會覺得祂特別溫馨.
神的愛之所以偉大.
是因為這個世界很大.
非常大.
今天有一件事沒有跟你說.
就是我們今天有170多人在我們中間.
今天我們想說今天在我們中間有170多人.
這個是破了我們參會的紀錄.
已經破了我們參會的紀錄.
你想想我們今天有大概200人在這裡.
你想想神就是愛我們這裡每一個人.
我們今天有200多人.
神愛我們這裡每一個人.
但如果看清楚.
這段聖經說神愛世人或者神愛the world的時候.
到底是什麼世人是什麼world呢.
聖經裡面告訴我們或者問我們這些問題.
這個世界是不是一個這麼乾淨的世界.
這個世人 這些世人 神愛世人.
是不是我們很可愛.
所以神來愛我們呢.
聖經裡面特別在約翰福音告訴我們.
這個世人.

$^{41}$這個world 這個世界.
是一個充滿叛逆.
充滿罪惡.
充滿黑暗混亂的一個地方.
但這個世界到底是怎麼樣的呢.
它是乾淨的嗎.
它是充滿愛的嗎.
在聖經中特別是在約翰福音裡.
世界被描寫成一個混亂.
悖逆 充滿罪的地方.
所以這裡說到神的愛的偉大的地方.
就不是因為他愛了很多人.
神的愛的偉大.
是居然愛一幫很壞的人.
所以神的愛之所以偉大.
不是因為神愛很多的世人.
而是因為神居然愛很壞的世人.
所以如果我們真是誠實的去面對我們自己的話.
我們必須要同意.
聖經裡面對我們的描述.
聖經裡面告訴我們.
我們是一些不可愛的人.
我們可以假裝自己比其他人好一點.
是不是 我可能假裝我比你好一點.
或者你會假裝你比我好一點.
但實際上.
我們世界裡面每一個人.
我們其實都是.
搞不定的我們.
如果我們誠實的面對自己.
就必須同意聖經的描述.
我們其實是不可愛的人.
我們常常假裝自己比實際更好.
但其實我們世界是一團糟的.
我們的內心也是一團糟.
所以如果你找一找你內心.
找一找你最近你自己內心的心處.
你可能都知道這個就是事實.
神愛一幫很壞的世人.
不可愛的世人.

$^{81}$這個是我們內心可以去贊同的.
聖經告訴我們.
我們是愛黑暗勝過愛光明.
因為我們想隱藏我們自己.
隱藏我們最黑暗的一面.
隱藏我們的罪惡.
為自己的自私 為自己的惡行.
我們會找藉口的.
而且在內心深處.
我們其實都知道這是真的.
我們並不是好人.
我們是很壞的世人.
聖經告訴我們.
我們愛黑暗勝於光明.
我們隱藏罪惡.
為自己的自私找藉口.
所以你和我想像一下.
你如果看到一些東西.
是又骯髒又不可愛的時候.
你通常的反應會是怎樣的呢.
你會不會是想靠近它.
還是你想離開它呢.
當事情變得又髒又亂的時候.
人們通常會怎麼做呢.
是想靠近它還是想避開它呢.
我們教會裡面有很多的小孩.
BB仔的.
那些BB仔可不可愛.
我們在教會裡面會有一些嬰兒.
你覺得嬰兒可愛嗎.
可不可愛.
可愛 你夠膽不說不可愛.
你也不敢說他不可愛了.
我記得有一天.
我們教會崇拜完之後.
大家就圍著某一個小孩.
不說是誰.
圍著那個小孩.
大家就抱他.
然後哄他開心.

$^{121}$和他玩.
然後撩他笑等等.
我記得有一天在崇拜之後.
大家圍著一個嬰兒.
我也不會說是誰的嬰兒.
大家鬥他 跟他玩.
有人抱他.
然後突然之間.
你可以想像.
突然之間有一股很奇怪的味道.
就出來了.
附近的人都感受到.
有些惡臭出現了.
當然就是那個嬰兒.
他賴了屁了.
弄髒了.
我就看到抱著那個.
無論是哥哥,姐姐.
或者叔叔或者阿姨也好.
他突然間臉色就變了.
他本身是抱著嬰兒的.
但是一聞到那些惡臭的時候.
他就突然間這樣.
然後那個小孩就感應到.
那個嬰兒很厲害.
他感應到被人否定和被人避開.
然後那個本身抱著的那個.
就立刻將那個東西.
就丟了給他的親生爸爸媽媽.
突然之間有一股非常非常奇怪的.
很臭的味道出現.
那個嬰兒的尿布髒了.
我看到抱著嬰兒的人的臉色.
馬上就變了.
然後幾乎馬上地.
把那團東西抱遠一點.
嬰兒可能也感受到了.
就馬上哭了起來.
那位弟兄姊妹很快就把那個嬰兒.
就馬上交回到他的父母手中.

$^{161}$我相信大部分的人.
你聽見這個例子的時候.
你都會有什麼反應 是嗎.
對嗎.
我想大多數的人聽到這個例子的時候.
都會有相似的反應.
為什麼呢.
因為那個不是你的兒子或者女兒.
因為他不是你的孩子.
所以但是父母一接到手過來的時候.
我就看到那個父母.
就立刻把這個嬰兒安撫他.
讓他不要哭.
因為愛他的父母.
雖然有一陣臭味.
但是讓他抱著他.
然後慢慢地幫他弄乾淨.
但父母呢.
我看到父母立刻接過孩子.
去安撫他.
愛孩子的父母.
是不會丟下自己的孩子的.
哪怕他很臭.
他們不會說你先把自己弄乾淨再來.
父母愛那個孩子.
所以就願意去靠近他.
抱起他.
雖然他很臭.
雖然他很不可愛.
但是去靠近他.
然後將他清理他的所謂的混亂.
你就可以見到.
因為愛的緣故.
父母就走向一些混亂的東西.
走向一些混亂的東西.
就不會是彈開遠離的.
他們會靠近孩子.
抱起孩子.
雖然他們不可愛.
很臭.

$^{201}$但是父母會主動地去清理那一團混亂.
為什麼呢.
那是因為愛.
會讓父母走向混亂.
而不是遠離混亂.
我相信這個有一點道理.
就是在說《約翰福音》三章十六節.
上帝如此愛一班不可愛的世人.
大家想起愛的時候.
就是愛是驅使上帝去走向我們.
雖然我們不可愛.
同樣的.
《約翰福音》三章十六節.
也告訴我們一樣的道理.
上帝如此愛不可愛的世人.
驅使愛驅使上帝走向這個混亂的世界.
你看見上帝做了什麼.
他有沒有選擇去保持距離.
因為你很臭.
所以他保持距離.
他有沒有選擇去放棄我們.
還是選擇去.
不如等你那些混亂.
你的臭味搞定了先.
然後他才來擁抱你.
上帝是沒有這樣做的.
那神做了什麼呢.
神選擇保持距離嗎.
他是因為我們臭.
然後選擇放棄我們了嗎.
還是他等我們把自己的問題解決了.
才來愛我們嗎.
沒有.
是上帝知道我們不可愛.
上帝知道我們搞不定.
上帝知道我們是破碎的人.
所以他反而願意接近我們.
差派他的獨生子耶穌基督來到我們中間.
降生在馬槽.
降生在一個卑微的馬槽.

$^{241}$在貧窮裡面出生.
在軟弱裡面出生.
在我們的罪惡.
在我們的邪惡中間出生.
上帝願意觸摸我們一些不乾淨的人.
愛一些應該被拒絕的人.
並且走向我們生命裡面.
我們有傷口有傷痕的一班人.
上帝知道我們不可愛.
知道我們搞不定.
知道我們很破碎.
神走進我們的混亂之中.
拆遣他的兒子.
耶穌降生在馬槽裡.
出生在貧窮.
軟弱.
充滿了罪的世界中.
神觸摸不潔的人.
神愛那些被拒絕的人.
他走進我們的傷口還有黑暗.
為甚麼上帝要這樣做.
聖經講得很清楚.
就是因為上帝的愛比我們一切的破碎更加大.
比我們一切的傷痛.
我們一切的傷痕.
我們的軟弱更加深.
即使我們是失敗.
即使我們是搞不定的一班人.
神仍然愛我們.
選擇去靠近我們.
為甚麼神這麼做呢.
聖經清楚地告訴我們.
神的愛比我們的罪.
傷痛和軟弱更深.
即使當我們如此敗壞.
如此破碎.
神仍然愛我們.
神選擇靠近.
所以聖誕節其實是提醒你和我.
上帝沒有離開.

$^{281}$撇棄我們一班這麼亂的人.
神反而願意去進入.
然後擁抱我們.
去清理我們.
藉著耶穌這位獨生子來拯救我們.
聖誕節提醒我們.
神沒有離開我們這一團亂.
祂進入其中清理祂.
祂藉著祂的兒子拯救我們.
這就是聖經裡面的好消息.
就是聖誕節的訊息.
這是好消息.
這就是我們的聖誕節.
今天最後一個問題想問大家.
就是.
你是否願意伸出雙手.
有人給你一份禮物.
這份最大的禮物就是.
這位上帝不介意.
我們一班這麼不濟的人.
祂將祂的獨生子給你.
我們願不願意用雙手去接受祂.
問題是.
今天你是否願意伸出雙手.
去領受那位為了愛我們.
而降生的神的兒子呢.
我們每一次到我們聖誕節服務.
快要完結的時候.
我們都會點起燭光.
因為燭光去提醒我們.
上帝的兒子在我們中間降生.
帶給我們光明.
去驅走黑暗.
最重要的光明是帶給我們生命裡面.
我們沒有黑暗的時候.
如果你接受這份禮物.
耶穌基督進入你生命裡面.
祂就會慢慢照亮所有黑暗的地方.
而上帝是想擁抱你.
上帝不是去審判我們.

$^{321}$而是用慈愛來擁抱我們.
如果你願意接受這份禮物的時候.
我希望你今天和我們.
有一個很簡單的心裡面的和應.
我們等一下會唱我們的回應詩歌.
這首回應詩歌.
我希望能夠把握到這個訊息.
這個訊息就是聖誕節裡面的訊息.
就是耶穌來到我們中間.
我們在唱的時候.
我們就會把燭光傳到不同的桌子上.
如果你心裡面和我們有一個和應的話.
我請你也點一點那支蠟燭.
並且心裡面和我們這首歌來和應.
然後我們接著低頭祈禱.
謝謝您的收看 祝您年輕時光快樂.
\newpage



\section{路加福音 2:1-14}
\label{sec:VP0gKdg8cek}
\textbf{誰是我王( 路加福音 2\_1-14) - 梁紫健傳道}
\newline
\newline
連結: \href{https://youtube.com/watch?v=VP0gKdg8cek}{\texttt{https://youtube.com/watch?v=VP0gKdg8cek}} ~~~~ 語音日期: 2025-12-29
\newline
\newline
\hyperref[sec:wPcDqaSy0yg]{\small{< < < PREV SERMON < < <}}
~
\hyperref[sec:index_chronic]{\small{[返順時目]}}
~
\hyperref[sec:index_scriptual]{\small{[返順卷目]}}
~
\hyperref[sec:kGtKq4doSFU]{\small{> > > NEXT SERMON > > >}}
\newline
\newline
路加福音 2:1-14
\newline
\begin{longtable}{cl}
\hline
\hline
章節 & 經文 (和合本修訂版)\\
\hline
2:1 & \begin{tabularx}{0.7\textwidth}{X} 在那些日子，凱撒奧古斯都降旨，叫全國人民都登記戶籍。 \end{tabularx} \\ \\ \relax
2:2 & \begin{tabularx}{0.7\textwidth}{X} 這第一次登記戶籍是在居里扭作敘利亞總督的時候行的。 \end{tabularx} \\ \\ \relax
2:3 & \begin{tabularx}{0.7\textwidth}{X} 眾人各歸各城，辦理登記。 \end{tabularx} \\ \\ \relax
2:4 & \begin{tabularx}{0.7\textwidth}{X} 約瑟也從加利利的拿撒勒城上猶太去，到了大衛的城名叫伯利恆，因為他是大衛家族的人， \end{tabularx} \\ \\ \relax
2:5 & \begin{tabularx}{0.7\textwidth}{X} 要和他所聘之妻馬利亞一同登記戶籍。那時馬利亞已經懷孕。 \end{tabularx} \\ \\ \relax
2:6 & \begin{tabularx}{0.7\textwidth}{X} 他們在那裡的時候，馬利亞的產期到了， \end{tabularx} \\ \\ \relax
2:7 & \begin{tabularx}{0.7\textwidth}{X} 就生了頭胎的兒子，用布包起來，放在馬槽裡，因為客店裡沒有地方。 \end{tabularx} \\ \\ \relax
2:8 & \begin{tabularx}{0.7\textwidth}{X} 在伯利恆的野外有牧羊人，夜間值班看守羊群。 \end{tabularx} \\ \\ \relax
2:9 & \begin{tabularx}{0.7\textwidth}{X} 有主的一個使者站在他們旁邊，主的榮光四面照著他們，牧羊人就很懼怕。 \end{tabularx} \\ \\ \relax
2:10 & \begin{tabularx}{0.7\textwidth}{X} 那天使對他們說：「不要懼怕！看哪！因為我報給你們大喜的信息，是關乎萬民的： \end{tabularx} \\ \\ \relax
2:11 & \begin{tabularx}{0.7\textwidth}{X} 因今天在大衛的城裡，為你們生了救主，就是主基督。 \end{tabularx} \\ \\ \relax
2:12 & \begin{tabularx}{0.7\textwidth}{X} 你們要看見一個嬰孩，包著布，臥在馬槽裡，那就是給你們的記號。」 \end{tabularx} \\ \\ \relax
2:13 & \begin{tabularx}{0.7\textwidth}{X} 忽然，有一大隊天兵同那天使讚美神說： \end{tabularx} \\ \\ \relax
2:14 & \begin{tabularx}{0.7\textwidth}{X} 「在至高之處榮耀歸與神！在地上平安歸與他所喜悅的人！」 \end{tabularx} \\ \\
[1ex]
\hline
\hline
\end{longtable}
$^{1}$要馬上吃雞.
大家頂尖會今年過得好嗎.
好.
我猜頂尖會在不同的崗位上.
都有不同的挑戰.
也有不同的驚喜.
我們年尾都一齊數算上帝給我們的恩典.
不知道大家有沒有聽過.
身邊有些朋友這樣說.
是弱者才需要信耶穌.
有沒有聽過.
你怎樣回應他.
我覺得有個朋友他很熟讀歷史.
他也是一個很醒目的人.
他覺得人生不需要求助.
也不需要去找任何人的幫助.
他覺得你的信仰只是用來安慰自己.
是弱者才需要的.
他不需要.
我知道他很崇尚強大的政權.
於是我就看看這個世界有什麼很強的王.
我發現原來一些很強大的君王.
他們做很多事情令他們名流千古.
他們做的事情的目的就是向世界宣示.
我是最強的.
我無人能及.
先說一下這個.
知道這個是誰嗎.
秦始皇.
是,每個人都認識.
古代有秦始皇.
殲滅六國.
收築萬里長城.
統一地方民治.
貨幣各樣的道兩行他都統一了.
建立一個中央集權的國家.
令世人都知道他的名字.
千古留名.
古今中外.
無人不識.

$^{41}$然後我發現原來不單只是秦始皇.
我發現現代的君王都有一個習慣.
就是想告訴全世界.
他是很有權力.
很威風.
這個是太皇.
2019年他在他的加冕登基典禮中.
動用超過20萬人.
花費超過3000萬美金.
去到辦這個加冕的典禮.
他接受水禮.
聖水的灌頂.
高僧的祝福.
升格為一個佛教的神.
轉化為人間的化身.
接受八方政壇的八方聖水.
確立了他對國家的統治性.
是很有正當.
然後祭司就會為他穿起皇袍.
正式成為一個人神一體的太皇.
他就是一個神格化的皇帝.
由此登基.
自古以來皇帝這個身分是很尊貴.
很強大.
很令人崇拜.
是所有人都很想依靠他的一個物體.
一個being.
一個存在.
成為了神明他就會怎樣呢.
管理這個世界的角落.
至少在他的國家裡面.
人民的福祉是依靠著這個主宰.
所以基本上每一個皇帝都想統治更多的地方.
越多的地方越多的人.
表示自己的力量實力很強.
可以名流千古.
今天我們說的經文就是.
老家用了一個去到哪裡就贏到哪裡.
千古一帝屋大衛的凱撒.
去襯托一個出生在馬槽的耶穌.

$^{81}$我們來看看在2025年結束之前.
這段經文給我們有什麼反省和提醒.
準備好我們來看看.
聽完這篇講道我們會知道.
究竟老家是怎樣利用這麼強大的君王.
去襯托一個表面這麼軟弱的耶穌.
一個BB.
面對現實我們究竟怎樣選擇呢.
我們來看看經文.
在這些日子裡面.
凱撒奧古斯督降旨叫全國人民登記戶籍.
這次登記是居尼柳在敘利亞總督之名去舉行的.
眾人就各歸各城辦理登記.
約瑟從加利利的拿撒勒城回到猶太.
去了大衛城一個叫做巴黎行的地方.
因為他是大衛家族的人.
要跟他所聘的妻瑪利亞一同登記戶籍.
那時候瑪利亞已經有新人了.
他們在那時候瑪利亞生產奇游到.
生了一個投胎的兒子.
用布包起放在馬槽裡面.
因為客店沒有地方.
今天第一個出場的人物是千古一帝.
幫我按下一張.
謝謝.
是屋大衛這個皇帝.
他非常之強大.
千古一帝屋大衛.
你聽他的名字都知道厲害.
因為名字那麼長.
他在主前2017年就已經平息了羅馬帝國裡面幾十年的內亂.
他去到哪裡打到哪裡.
從來都沒有輸過.
你贏了那個城市.
如果你反抗他的話.
他就會屠城.
基本上他是殺人如麻.
最後他征服了整個以色列和周邊地區的地中海沿岸.
大家如果看圖就知道了.
這個圖.

$^{121}$你看到黃色和綠色的地方就是他征服的地方.
一個圈.
中間那個圈當然就是地中海沿岸.
埃及和現在的西班牙.
還有非洲的北部.
都是他征服的地方.
如果太複雜的話.
先看簡易版.
再按一次.
對,簡易版.
綠色的地方就是他拿到的地方.
只要他想拿的地方.
他只要一打就贏了.
沒有人可以抵抗到他.
他最後還將整個地中海.
變成自己國家的內海.
好像自己的私家水塘一樣.
所以人民更加把他當神來拜.
他自己也神格化他自己.
他覺得我就是神.
我就是代表神統治這個地方的王.
可想而知這個王有多強大.
然後登記戶籍就是.
經文是這樣說的.
海薩奧古斯督下令登記戶籍.
如果你看舊日本就是報名上冊.
其實報名上冊也不是很好.
登記戶籍會比較接近我們現在的用字.
登記戶籍是在古時.
花費一個很大的人力物力的人口普查.
不像現在做人口普查是怎麼做的.
上網填一張表.
填一張表就完成了.
以前是怎樣的.
以前政府要命令人回自己的鄉下.
帶著一家大小.
盤穿所有家當.
回到自己的家.
然後士兵就來查看.
你們有三個人.

$^{161}$然後這個戶就有三個.
目的是什麼.
目的是要把全國的人數記錄下來.
然後做什麼.
徵稅.
還有徵兵戎.
就是告訴別人這裡有多少個男丁.
到時候打仗我可以怎麼用.
這裡有多少個人口.
我可以徵多少稅.
就好像舊約.
如果你看《十武二記》.
大衛也有數點人數.
其實數點人數也類似登記戶籍.
去到展示戶口就告訴別人.
這個皇帝厲害.
我們有這麼多兵.
從來只有一個很強勢.
很有錢的君王.
才會發動這麼大規模的登記戶籍.
而奧古斯都在歷史上發動了不止一次.
基本上發動了很多次.
所以這個凱撒.
他做的目的是什麼.
他做的目的就是.
向全世界告訴全世界.
我是這個世界的主宰.
沒有人可以反抗他的命令.
連所有人.
甚至臣服於上帝.
這個超然主宰的以色列人.
都要臣服於他.
所以沒有人敢亦其意.
第三節的經文就告訴你.
眾人各歸各城.
還刻意描寫眾人.
回應了奧古斯都的要求.
所以第三節.
各歸各城的行動.
似乎很正常.

$^{201}$皇帝下令百姓遵命.
是一件很正常的事.
今天政府叫你交稅.
你都交吧.
如果沒有你說我奉了憲.
我不交稅.
不可以的吧.
不可能的.
但是這裡用的字很有趣.
它和你還沒記用的字眼很相似.
你還沒記25章是這樣說的.
第三節.
這欺憐你們國人要歸自己的地業.
要歸自己的地業.
和《路加福音》二章三節.
各歸各城的表達非常相似.
《路加福音》這裡做了一個諷刺.
刻意利用凱撒奧古斯都降旨.
叫全國人登記戶籍.
叫人各歸各城的行動.
去描述整個皇帝的目的就是.
路加醫生很想刻意將這個人的命令.
和上帝的命令放在一起.
上次欺憐跟大家說過一下.
這次也要跟大家說一點點.
欺憐是什麼.
是上帝命令猶太人做的事情.
暗示著這個叱吒風雲的千古一帝.
凱撒奧古斯都不及一個出生在馬槽的耶穌.
為什麼呢?.
我們要說一下.
欺憐的概念.
這樣.
OK 好.
在利未記如果大家記得的話.
上帝作為以色列人的主宰命令他們各歸各城.
遵行欺憐的目的就是要表達以色列人敬畏上帝.
他是耶和華 他是以色列人的上帝.
耶和華才是以色列人的主宰.
先天以色列人在埃及那裡.

$^{241}$神拯救了他們出來.
然後進入迦南地.
上帝遵守了一個盟約讓以色列人有土地.
然後命令以色列人要遵守欺憐.
讓孤兒寡婦 寄居者 弱小的人有機會免債.
讓失去家園的人能夠拿回土地.
奴隸被釋放.
只是以色列人當年沒有遵行欺憐.
直到路加年代的以色列人.
他們的身分都被反過來.
他們失去了原本的土地.
被羅馬軍欺壓 被收很重的稅項.
就像當年的寄居者奴隸一樣.
很弱小 生活苦不堪言.
所以各歸各城就令人想起以色列人的欺憐.
當年上帝讓他們進入迦南地.
很希望他們可以拿回本身的土地.
再次成為主人.
現實是什麼呢.
現實就是長勝將軍奧古斯督.
他就是最有能力令所有人臣服於他.
因為他去到哪裡 打到哪裡.
去到哪裡 贏到哪裡.
有能力發動登記戶籍.
更加有能力令以色列人以欺憐的名義拿回土地.
只是歷史告訴我們.
他在耶穌出生之後 他就死了.
所以路加在這個時候再提起.
這句「各歸各城登記戶籍」.
是要暗示這個很強大的千古一帝 屋大衛.
凱撒奧古斯督.
本來很有能力去主宰這個世界.
但是都沒辦法做到.
唯有這個主宰 真正世界主宰耶穌才做得到.
去到這裡我們就要問一個問題.
路加憑什麼說服全世界.
告訴人家那個躺在馬槽的BB耶穌是世界主宰.
大家有沒有想過這個問題.
路加憑什麼.
面對這麼強大的奧古斯督.

$^{281}$我就告訴人家這個躺在馬槽的耶穌才是世界主宰.
路加刻意描寫約瑟是大衛一家一族的人.
有主的使者向牧羊人報大喜的訊息.
然後是關於盲民的救主主耶穌基督.
宣告天上的榮譽歸於神.
地上的平安歸於耶穌所喜悅的人.
大家看到那些紅色的字.
紅色的字.
全部這裡用的紅色的字都是極度政治性的.
當時唯有羅馬君王才可以用這些字眼.
是羅馬君王才可以為這個世界帶來和平和平安好消息.
因為他才夠強大.
在戰亂的時候保護到人民.
如果今天你去羅馬的話.
你會看到這個是奧古斯督自己命人起的平安祭壇的紀念碑.
裡面寫了是他才可以將平安帶給全世界的人.
是他才可以將平安帶給全世界的人.
其他人用是會招致殺身之禍的.
所以今天如果你祝福人家平安的話.
其實如果在2000年前你是會被人殺頭的.
因為你不是君王.
是羅馬君王才可以將平安帶給人.
所以路加這裡很大膽.
他竟然用地上的平安.
用地上的平安歸於耶穌所喜悅的人.
地上的平安竟然來自一個躺在馬槽的BB耶穌.
而不是千古一帝屋大衛.
凱撒奧古斯督.
這個這麼強烈的對比.
我也要講一個例子去展述一下.
就是我有兩個四當同學一起長大的.
一個叫朱仔.
左邊那個叫朱仔.
朱仔雖然他很聰明.
他叫朱仔但他很聰明.
他很耐勞.
一點都不朱.
另一個叫強仔.
雖然他叫強仔但他一點都不健碩.
很散.

$^{321}$所以我們經常叫他散仔強.
所以我們三個人從幼稚園到小學到中學.
都是同學一起長大的.
都是在最差的那一班.
因為我從小有病,讀不到書.
他們兩個經常上課就做夢.
就覺得自己會成為什麼大代.
什麼大亨.
所以我們一起升中的時候.
毫無意外都會升到全區最差的那一班中學.
又是你?我又見到你.
就是最差的那一班.
還要是最差的那一班相遇.
在整個中學裡面.
這間中學被人叫做武術學校.
不是因為他們教習武.
是因為經常打架.
裡面的人經常開拖打到流血.
你可以想像裡面是幫派林立.
很多大佬,很多嘍囉.
每個人都有大佬照顧.
不然就會被人打,被人勒索.
他們打架不是玩玩的.
如果你有留意以前的班房.
椅子,桌椅全部都有鐵通.
我們咬斷了鐵通就出去打人.
打到全身都血.
我想言之.
我們三個小朋友.
十一,二歲.
剛剛升中.
第一件事是什麼呢.
你回到第一天就要決定一件事.
就是你跟不跟大佬.
如果你不找一個碼頭照顧你.
都幾十歲.
朱仔剛剛說他是聰明.
他第一天就決定跟一個看起來很強悍的大佬.
所以他不會被人勒索.
因為他可以勒索人.

$^{361}$沒辦法.
這個弱強食的世界在中學已經出現.
我跟小散強.
他就散,我就弱.
雖然我叫健仔.
但我已經不健壯.
我經常都暈倒.
我怎樣跟人打架.
大佬見到我都不喜歡.
所以我們兩個傻乎乎的被人打.
被人勒索.
有一天大佬就站在我們面前.
叫我們加保護費.
我一元都沒有,不加也不加.
那天小散強有點錢.
他拿了一疊現金出來.
就交給大佬.
大佬就拿了就走了.
我以為沒事.
加了保護費還想怎樣.
第二天大佬就面腫口腫.
被人打到.
他連眼都看不開.
然後他就拿著那疊現金.
交給小散強.
不好意思,我以後不敢了.
然後我就睜開眼看著小散強.
他憑甚麼呢.
我從小就認識他.
他打也打不成.
又不是很兇狠.
然後他說他打了給他表哥.
然後我就說.
給表哥?.
後來我才知道.
原來他表哥是雙花紅棍.
是這個地區的坐館.
這個地區所有的事情.
基本上都是他大佬提的.
所以這個人很害怕.

$^{401}$這個大佬,所謂的小嘍囉.
他完全不敢碰小散強.
自此之後.
死木仔,豬仔.
他就經常要做打手.
你知道出來走就遇到要還.
經常被人打到全身都受傷.
正常人,幾十歲的小孩.
你都會計算.
正常人都會找一個很強大的碼頭.
你絕對不會找小散強.
你都絕對不會找一些.
剛出來報道的小孩.
正如《路加福音》二章這裡.
千古一帝屋大衛.
你怎會不依靠他呢.
他掌握著你的生死.
你交稅,你的戰爭.
你家人的福祉.
是完全被他手掌握著.
他是這一區的大佬.
他是這一區的主宰.
當時所有人都覺得.
他才是可以保護人民的大佬.
是他才可以帶來平安.
他收的稅就是保護費.
他讓你住在這裡,就平安.
但路加竟然把這個故事反轉地說.
他一個看似很軟弱的耶穌.
竟然成為了這個世界的主宰.
再說一個故事.
不知道大家有沒有看過.
阿凡達.
不知道大家看過阿凡達三尾.
火與灰燼.
建議大家如果入場看這場電影.
千萬不要選擇2D.
選擇3D,非常好看.
它的效果很厲害.
為了不劇透.

$^{441}$我只說阿凡達主要故事的主線.
就是人類覺得自己很強大.
有武器,有科技.
在地球自無不勝.
就像千古一帝,屋大衛.
奧古斯都,全世界都要怕他.
他們征服了地球也不要緊.
他們還想征服另一個星球.
去了一個星球叫潘多拉.
在那裡基本上想殺死所有生物和土著.
告訴人類,人類才是主宰星球的上帝.
但最後他們忽略了.
這個世界是有神明存在的.
一個比人類更強大.
100萬倍,無限倍的存在.
當人類以為可以贏到盡.
拿盡所有資源,統治宇宙的時候.
劇情就反轉了.
我只說了這麼多,大家喜歡就去看吧.
路加福音用千古一帝,屋大衛.
奧古斯都的這個強勁的君王.
就算死了也要把他拿出來.
襯托耶穌.
就因為人類覺得自己總之什麼都最強.
贏到開巷,所有人都要臣服於大帝底下.
但耶穌本身擁有了比人類無限的能力.
他可以直接命令天君,天使.
去剷除所有以為自己最強的勢力.
還可以順便訂一間五星級酒店.
在札古斯出世都可以.
但他沒有這樣做.
因為他很想人知道.
人是很想很多東西去保護自己.
很多勢力去幫助我們.
我們覺得有平安.
你擁有表面上一個最強的勢力.
我們才覺得平安.
所以如果今天我們確信.
耶穌是我們的王.
我們就開始要鑽煉.

$^{481}$不是說我們不是選擇最強的東西.
而是我們理性上知道.
當我們有最好的選擇的時候.
我們還要留一個選擇的餘地.
一個空間給上帝指引我們.
可能有時候人類看起來是一個很愚蠢的決定.
人類會覺得這樣選擇才是對的.
為什麼要這樣選擇.
基督徒經常都做這樣的決定.
例如.
不知道你覺得要多少錢才夠錢退休.
有沒有想過這個問題.
應該大家都想過這個問題.
這個是我問AI幫我做資料搜集的結果.
如果年支出四萬元.
你就需要四十萬至五十七萬.
如果你每年花五萬元.
每個人計算不是一個家庭計算.
如果你花五萬元.
你就需要六十五萬至八十三萬美金的資產.
你才可以支撐你之後的退休生活.
後面的就不說了.
因為現在連頭都看不到我.
所以.
看起來好像很害怕我.
所以現在的人會很拼命去找工作.
去找工作增加收入.
連聖誕節都去上班.
有時候家人都沒時間去跟他們相處.
很擔心我之後沒錢花.
不夠錢花.
這個世界告訴你.
如果你聖誕節放假.
你簡直是蠢.
不去賺錢.
還要去做義工.
更蠢.
我們有班弟兄姊妹放假不去旅行.
也不去賺錢.
在這個冬天裡面.

$^{521}$希望帶一點溫暖給社區的公婆婆.
去傳福音.
在世人眼光裡面其實是很蠢.
當全世界都在為自己準備最強的退休計劃的時候.
但基督徒竟然在做這些事.
不是為自己賣命.
而是為搏命去溫暖這個社區.
為甚麼打架仔呢.
就是因為我希望UV的弟兄姊妹.
每一個人都可以成為其中的一員.
不是特別是一些弟兄姊妹去做.
而是每一個弟兄姊妹我們都為耶穌搏命.
另一個就是聖誕派對.
每一個人都很盡力去跟新來的朋友聊天.
希望跟他們建立好的關係.
以至他們認識耶穌.
當全世界都在為自己打拼的時候.
有一班人竟然跟著耶穌.
去選擇消耗自己寶貴的時間去接觸這個社區.
我們跟著這個王.
在世人眼中就是這麼弱這麼蠢.
一點都不強勁.
一點都不醒目.
所以我們擔心自己退休之後.
有時候會擔心自己的小孩.
我們基督徒有時候都會被社會的意識形態影響.
覺得小朋友一定要讀最好的學校.
有最好的資源.
將來才找到最好的工作.
他才會快樂.
無可否認作為家長.
在自己的能力底下.
你給了我.
我也想給最好的東西給我自己的小朋友.
但其實我們初衷是在想.
讓小朋友開心.
他找到人生的意義.
無論他最後怎樣選擇.
都是他決定的.
去到這個地步我們沒有辦法.

$^{561}$所有事情都幫他安排了.
雖然我們可以安排他們很多東西.
覺得無懈可擊.
完全沒有破綻.
期望就是像他跟我們想的一樣這麼成功.
但事情往往跟我們想像的完全不同.
有時候我們在理性底下.
作出的選擇是最好的.
但在路加福音提醒我們.
要留一些空間.
雖然我們面對現實是很現實.
但我們要留一些空間問一下耶穌.
你想這個小朋友成為一個怎樣的人呢.
你想一個小朋友在我們身上.
家長身上.
看到上帝的真實是怎樣的呢.
而不是雙方紅軍.
表面上很強大的一個大哥.
是哪個大哥厲害你就跟誰.
直到他可以作出一個自己的選擇.
他就說我想跟上帝耶穌.
他是一個怎樣的關係.
當我們想到工作的時候.
很多時候我們都是.
首先考慮到發揮的機會.
人工,地點,上班時間.
各樣的條件.
除了這些之外.
其實有沒有一些東西是耶穌想我們.
加入去考慮的因素呢.
我們可以這樣想.
到最後.
2025年還有幾天就結束了.
在這一年裡面你有多少時間.
是為了自己心中的王去拼搏的呢.
有多少時間我們為了耶穌去努力的呢.
盧家富欽利用了一個很強大的君王.
千古一帝,屋大衛,奧古斯督.
去襯托一個表面軟弱的嬰孩耶穌.
昔日最強的君王留下的是什麼呢.

$^{601}$如果你今天去羅馬.
留下的是一堆石碑.
留下的是一堆令人不安的名號.
他留下的不是平安.
而耶穌這個表面軟弱的嬰孩.
就好像今天唱的第一首詩歌一樣.
長勝君留下的一個平安.
他留下的是一個令你.
可以不用和這個世界競爭.
拼命地說我要賺很多錢.
我才可以足夠退休的一個平安.
當我們面對著很多現實.
我們怎樣選擇呢.
讓我們一起祈禱.
這個時間是你和上帝的時間.
有少許時間我們想和上帝說.
我們今天.
過去那一年可能我們很多事做得不好.
可能很多事我們是.
以我們的標準去選擇.
但主人今天.
我希望我們接下來的一年.
我們選擇你.
當我們做一個理性的選擇的時候.
我們都不忘問你.
你覺得怎樣.
你覺得怎樣是一個選擇.
而不是我們自己說.
我是想這樣做.
天父多謝你給我們這個時間學習你的話語.
接下來這個新的一年.
給我們面對著現實的困難的時候.
我們聽得到你的聲音.
我們能夠尋求到你給我們的指引.
幫助我們.
我們感到高雄的耶穌的勉強.
接下來新的一年.
祝福大家.
\newpage



\section{士師記 3:1-31}
\label{sec:kGtKq4doSFU}
\textbf{God’s Holy Love (Judges 3 ) - Chee Yie Jong 神聖的愛 (士師記3\_1-31) - 楊志溢傳道}
\newline
\newline
連結: \href{https://youtube.com/watch?v=kGtKq4doSFU}{\texttt{https://youtube.com/watch?v=kGtKq4doSFU}} ~~~~ 語音日期: 2026-01-14
\newline
\newline
\hyperref[sec:VP0gKdg8cek]{\small{< < < PREV SERMON < < <}}
~
\hyperref[sec:index_chronic]{\small{[返順時目]}}
~
\hyperref[sec:index_scriptual]{\small{[返順卷目]}}
~
\hyperref[sec:6JWiBdm7MIc]{\small{> > > NEXT SERMON > > >}}
\newline
\newline
士師記 3:1-31
\newline
\begin{longtable}{cl}
\hline
\hline
章節 & 經文 (和合本修訂版)\\
\hline
3:1 & \begin{tabularx}{0.7\textwidth}{X} 耶和華留下這些國家，為要考驗所有未曾經歷迦南任何戰役的以色列人， \end{tabularx} \\ \\ \relax
3:2 & \begin{tabularx}{0.7\textwidth}{X} 只是為了要以色列人的後代認識戰爭，教導他們，尤其那些未曾認識這些事的人。 \end{tabularx} \\ \\ \relax
3:3 & \begin{tabularx}{0.7\textwidth}{X} 留下的有非利士的五個領袖，所有的迦南人，西頓人，以及從巴力‧黑門山到哈馬口，住黎巴嫩山的希未人。 \end{tabularx} \\ \\ \relax
3:4 & \begin{tabularx}{0.7\textwidth}{X} 他們是為了要考驗以色列，好知道他們是否肯聽從耶和華藉摩西吩咐他們列祖的命令。 \end{tabularx} \\ \\ \relax
3:5 & \begin{tabularx}{0.7\textwidth}{X} 以色列人住在迦南人、赫人、亞摩利人、比利洗人、希未人、耶布斯人中間， \end{tabularx} \\ \\ \relax
3:6 & \begin{tabularx}{0.7\textwidth}{X} 娶他們的女兒，將自己的女兒嫁給他們的兒子，並事奉他們的神明。 \end{tabularx} \\ \\ \relax
3:7 & \begin{tabularx}{0.7\textwidth}{X} 以色列人行耶和華眼中看為惡的事，忘記耶和華－他們的神，去事奉諸巴力和亞舍拉， \end{tabularx} \\ \\ \relax
3:8 & \begin{tabularx}{0.7\textwidth}{X} 所以耶和華的怒氣向以色列發作，把他們交給美索不達米亞王古珊‧利薩田的手中。以色列人服事古珊‧利薩田八年。 \end{tabularx} \\ \\ \relax
3:9 & \begin{tabularx}{0.7\textwidth}{X} 以色列人呼求耶和華，耶和華就為以色列人興起一位拯救者來救他們，就是迦勒的弟弟基納斯的兒子俄陀聶。 \end{tabularx} \\ \\ \relax
3:10 & \begin{tabularx}{0.7\textwidth}{X} 耶和華的靈降在他身上，他就作了以色列的士師。他出去爭戰，耶和華將亞蘭王古珊‧利薩田交在他手中，他的手戰勝了古珊‧利薩田。 \end{tabularx} \\ \\ \relax
3:11 & \begin{tabularx}{0.7\textwidth}{X} 於是這地太平四十年。基納斯的兒子俄陀聶死了。 \end{tabularx} \\ \\ \relax
3:12 & \begin{tabularx}{0.7\textwidth}{X} 以色列人又行耶和華眼中看為惡的事。耶和華使摩押王伊磯倫強大，攻擊以色列，因為他們行耶和華眼中看為惡的事。 \end{tabularx} \\ \\ \relax
3:13 & \begin{tabularx}{0.7\textwidth}{X} 伊磯倫召集亞捫人和亞瑪力人到他那裡，他就去攻打以色列，佔據了棕樹城。 \end{tabularx} \\ \\ \relax
3:14 & \begin{tabularx}{0.7\textwidth}{X} 於是以色列人服事摩押王伊磯倫十八年。 \end{tabularx} \\ \\ \relax
3:15 & \begin{tabularx}{0.7\textwidth}{X} 以色列人呼求耶和華，耶和華就為他們興起一位拯救者，便雅憫人基拉的兒子以笏，他是個慣用左手的人。以色列人託他送禮物給摩押王伊磯倫。 \end{tabularx} \\ \\ \relax
3:16 & \begin{tabularx}{0.7\textwidth}{X} 以笏打造了一把兩刃的劍，長一短肘，綁在右腿上衣服裡面。 \end{tabularx} \\ \\ \relax
3:17 & \begin{tabularx}{0.7\textwidth}{X} 他把禮物獻給摩押王伊磯倫。伊磯倫是個很肥胖的人。 \end{tabularx} \\ \\ \relax
3:18 & \begin{tabularx}{0.7\textwidth}{X} 以笏獻完禮物的時候，就把抬禮物的人送走。 \end{tabularx} \\ \\ \relax
3:19 & \begin{tabularx}{0.7\textwidth}{X} 但他自己卻從靠近吉甲的雕像那裡轉回來，說：「王啊，我有一件機密的事要奏告你。」王說：「迴避吧！」於是所有侍立在他左右的人都退去了。 \end{tabularx} \\ \\ \relax
3:20 & \begin{tabularx}{0.7\textwidth}{X} 以笏來到王那裡，那時他獨自一人坐在陰涼的頂樓。以笏說：「我有神的話向你報告。」王就從座位上站起來。 \end{tabularx} \\ \\ \relax
3:21 & \begin{tabularx}{0.7\textwidth}{X} 以笏伸出左手，從右腿上拔出劍來，刺入王的肚腹。 \end{tabularx} \\ \\ \relax
3:22 & \begin{tabularx}{0.7\textwidth}{X} 劍柄連同劍刃都刺進去了，肥肉夾住了劍刃。他沒有把劍從王的肚腹拔出來，糞便就流出來了。 \end{tabularx} \\ \\ \relax
3:23 & \begin{tabularx}{0.7\textwidth}{X} 以笏出到門廊，把王關在樓門裡面，就上了鎖。 \end{tabularx} \\ \\ \relax
3:24 & \begin{tabularx}{0.7\textwidth}{X} 以笏出來之後，王的僕人就來了。他們觀看，看哪，樓門鎖住，就說：「他必是在陰涼的房間裡大解。」 \end{tabularx} \\ \\ \relax
3:25 & \begin{tabularx}{0.7\textwidth}{X} 他們等得不耐煩，看哪，樓門仍然不開，就拿鑰匙打開樓門，看哪，他們的主人已經倒在地上死了。 \end{tabularx} \\ \\ \relax
3:26 & \begin{tabularx}{0.7\textwidth}{X} 他們耽延的時候，以笏就逃跑了。他經過雕像那裡，逃到西伊拉。 \end{tabularx} \\ \\ \relax
3:27 & \begin{tabularx}{0.7\textwidth}{X} 他到了那裡，就在以法蓮山區吹角。以色列人跟隨他從山區下來，他在他們前面引路， \end{tabularx} \\ \\ \relax
3:28 & \begin{tabularx}{0.7\textwidth}{X} 對他們說：「緊跟著我！因為耶和華已經把你們的仇敵摩押交在你們手中。」於是他們跟著他下去，佔據了摩押對面約旦河的渡口，不准一人過去。 \end{tabularx} \\ \\ \relax
3:29 & \begin{tabularx}{0.7\textwidth}{X} 那時，他們擊殺了約一萬摩押人，都是強壯的勇士，連一個也沒有逃脫。 \end{tabularx} \\ \\ \relax
3:30 & \begin{tabularx}{0.7\textwidth}{X} 那日，摩押在以色列手下制伏了。於是這地太平八十年。 \end{tabularx} \\ \\ \relax
3:31 & \begin{tabularx}{0.7\textwidth}{X} 以笏之後，有亞拿的兒子珊迦，他用趕牛的棍子打死六百非利士人。他也拯救了以色列。 \end{tabularx} \\ \\
[1ex]
\hline
\hline
\end{longtable}
$^{1}$- Good morning, UV..
(speaking in foreign language).
Testing, testing..
Okay..
All right..
So when you think about God,.
what are some of the image or words.
that comes to your mind?.
Is he a wrathful God, or is he a loving father?.
Is he a mighty warrior, or a humble servant?.
Is he someone that is quick to punish,.
like a parking attendant that is always out there.
roaming the streets and will not hesitate.
to give out a parking ticket,.
even though your meter only expires for a few seconds?.
Or is he slow to anger, like a teacher.
that is always patient with his students.
despite all the repeated mistakes and questions from them?.
(speaking in foreign language).
A.W. Tozer famously said, "What comes into our mind.
"when we think about God is the most important things.
"about us.".
And so what he means here is that the way we view God.
actually shapes the way we live..
So let me give you a quick example..
If we think of God as a vengeful God.
that is always out there to get us,.
we will probably live in constant fear.
of doing something wrong or mess up, right?.
But on the flip side, if we see God as a loving father,.
when we mess up, he will be the first person.
that we run towards and ask for forgiveness, right?.
(speaking in foreign language).
So this morning we will continue our sermon series.
in the Book of Judges..
And so far we have looked at the double introduction.
sections that sets us up on the historical narrative.
of Israel up to this point..
That Israel, they was in the promised land already,.
but they failed to completely drive out the inhabitants.

$^{41}$in the land of Canaan as the Lord has commanded them..
And we see that as a result of that,.
the idol worship culture of the Canaanites,.
they slowly seep into Israel..
And these people of Israelites,.
they turn to worship these false Canaanites gods.
instead of worshiping Yahweh himself..
(speaking in foreign language).
Moreover, the introduction section also gives us.
a framework on Israel's spiritual and moral downfall.
that we will see repeating itself in the cyclical patterns.
towards rock bottom for Israel both as a nation.
and as God's chosen people..
(speaking in foreign language).
However, while Judges highlights God's people.
at its worst, more importantly,.
I think it also show us God's response to human's actions..
And through God's response to Israel's rebellion.
and disobedience, we get a clearer picture of who he is.
and what he is like..
And so today we will go into the main section.
of the book of Judges and look at the story.
of a judge called Ehud..
Ehud is probably a lesser known judge,.
but through this narrative,.
we will explore this particular character.
about God together, that is his holy love..
(speaking in foreign language).
But before we dive into it,.
just a few notes on this main section.
of the book of Judges..
The main section of Judges has six major judge stories.
based upon this cycle..
(speaking in foreign language).
So a quick review of the components of these cycles..
It usually begins like this..
Israel does evil in the eyes of the Lord.
and the Lord will gives or sells them.
into the hands of their oppressors..
And Israel will serve their oppressors.

$^{81}$for a certain amount of years..
And after that, Israel will cry out to the Lord.
and the Lord will raise up for them a deliverer or a judge..
And the oppressor is subdued.
and the land has rest and peace for a certain amount of years.
before the cycle repeats itself again..
(speaking in foreign language).
And as we dive into this narrative cycle of the major judge,.
I want us to keep in mind that.
these individual narratives of the judge.
of Othniel, of Ehud, of Deborah, of Gideon,.
Jethta and Samson,.
they should not be read in isolation from one another..
But together they actually make up a long.
and complex narrative movement.
that shows Israel's slow descents.
into religious apostasy and moral downfall..
One that is being introduced already.
in the double introduction section..
(speaking in foreign language).
So the first few major judge will be mostly positive,.
including our judge for today, Ehud..
But the further we actually go into the book,.
the more issues we find with these judges..
And this is actually being reflected.
in the writings of judges as well..
The further we go into the main cycle section,.
the less consistent each of these components.
appears on the narrative..
And when we come to the last judge of Samson,.
only the first three components were left..
And they didn't even cry out to the Lord at that point,.
which again shows the chaos.
that Israel has descended into..
(speaking in foreign language).
Okay, so with that said,.
let's dive into our passage on Ehud.
from Judges 3 today and look at what it means.
for God to be a God of holy love..
And more specifically, we will look at God's holy love.

$^{121}$from the way he responds to our human action..
So let's look at verse 12 together..
It says, "Again, the Israelites did evils.
"in the eyes of the Lord..
"And because they did this evil,.
"the Lord gave Aglon, king of Moab, power over Israel..
"Getting the Ammonites and Amalekites to join him,.
"Aglon came and attacked Israel,.
"and they took possessions of the city of Palms..
"The Israelites, they were subject to Aglon,.
"king of Moab, for 18 years.".
(speaking in foreign language).
, okay..
So we will first focus on God's holiness this morning,.
especially in the way he disciplines.
and punish his own people..
(speaking in foreign language).
So if you notice, verse 12 begins with the word again..
Why again?.
Because this is not the first time Israel has done evil.
in the eyes of the Lord..
Earlier in the chapter, in the first judges story.
in verse seven, we read that the Israelites did evils.
in the eyes of the Lord..
They forgot the Lord their God.
and served the Baals and the Asherahs..
So the again here highlights that this is actually.
a repeated offense from Israel..
(speaking in foreign language).
And how did God respond to it?.
We see that the Lord gave Aglon, king of Moab,.
power over Israel..
Aglon came and attacked Israel.
and they took possessions of the city of Palms..
And the Israelites, they were subject to king of Moab.
for 18 years..
(speaking in foreign language).
A lot of times when we read a passage of God.
punishing his people like this,.
we can be tempted to think of God as a wrathful.

$^{161}$and vengeful God..
And if we make a mistake, he will be out there to punish us.
like the parking attendant that is always out there.
to ticket people..
(speaking in foreign language).
This even leads some people to conclude.
that the God of the Old Testament,.
who is often portrayed as a God of wrath,.
cannot be the same as the God of the New Testament,.
who is usually portrayed as a God of love..
(speaking in foreign language).
For example, in the second century,.
an early Christian theologian named Martian of Sinai,.
he taught that the creator God of the Old Testament.
was a different and inferior God,.
harsh, legalistic, and judgmental..
In contrast, he taught that the God revealed by Jesus.
in the New Testament was a higher and more supreme God,.
purely loving, gracious, and merciful..
And this eventually led him to actually reject.
the Old Testament as it only revealed.
a lesser and inferior God..
(speaking in foreign language).
And so for people like Martians and any modern manifestation.
of this kind of thinking, it's not surprising that,.
in the early Christian tradition,.
the Old Testament was a God of love..
(speaking in foreign language).
And so for people like Martians and any modern manifestation.
of this kind of thinking and belief about God,.
there is a clear misunderstanding about who God is.
and what He is like..
Yes, God is love..
Yes, God is holy..
But God is not wrath..
Wrath is not an essential attribute of God,.
but wrath is a rightful attribute of God..
(speaking in foreign language).
Not an essential attribute of God,.
but wrath is a rightful expression of God's holiness.

$^{201}$in response to evil..
To put it another way, wrath is not part of God's character,.
but when God sees evils,.
His holiness compels Him to respond in righteous anger.
because God hates evil..
(speaking in foreign language).
And so this is what we can learn about God.
in His response towards human beings..
(speaking in foreign language).
And so this is what we can learn about God.
in His response towards human weakness here..
In essence, God is a holy God and He cannot tolerate evil,.
not even just a little bit of it..
(speaking in foreign language).
More importantly, God's holiness usually goes hand in hand.
with His love as well..
Just as a loving parent who would not spare the rod.
of discipline towards their children,.
God's discipline or punishment towards His children.
always comes from a place of love.
in order to bring them back to Himself as well..
(speaking in foreign language).
And when it comes to God's punishment,.
there's always a limit and a greater purpose behind it.
because God doesn't just punish people.
for the sake of punishment alone..
You can think of God's punishment as more like a timeout.
that a parent would give to a children..
This is because ultimately, His judgment,.
His punishment, and His disciplines,.
they all comes from His character as a holy God,.
as a loving father, and the perfect judge..
And so when God acts,.
all of these characters act in perfect unison..
(speaking in foreign language).
And so this is what we see in our passage here..
For 18 years, Israel was powerless.
against Aglon, king of Moab,.
because the Lord gave Aglon power over Israel.
as a way of disciplining his own people..

$^{241}$And finally, after 18 years of oppressions.
under the hands of Moab,.
Israel finally shows some signs of repentance..
And we'll see how God responds..
So let's look at verse 15 together..
Again, the Israelites cried out to the Lord,.
and He gave them a deliverer, Ehud,.
a left-handed man, the son of Gerar the Benjamite..
The Israelites sent him with tribute.
to Aglon, king of Moab..
(speaking in foreign language).
And this brings me to the second focus for today..
That is God's love in His deliverance..
(speaking in foreign language).
So while God's holiness compels Him.
to respond to the wickedness of His people with wrath,.
His love, on the other hand,.
compels Him to respond to the crying out.
with mercy as well..
And so when the Israelites cried out to the Lord.
after 18 years of oppressions.
under king Aglon of Moab,.
He gave them a deliverer, Ehud..
(speaking in foreign language).
And we didn't look at the story of the first judge,.
but compared to the first judge, Othniel,.
who is a leader and a warrior,.
Ehud, he's a very different kind of judge..
We learned that he is a left-handed man,.
or more literally, he is a man restricted.
in his right hand, which is the dominant hand.
in that culture..
And so this gives him the appearance.
of weakness in some sense..
(speaking in foreign language).
, moreover, Ehud is also not a leader in Israel,.
but he's just someone being tasked with the job.
of bringing tribute to Aglon king of Moab..
(speaking in foreign language).
And yet through a series of trickery and satire,.

$^{281}$we will see that the Lord will actually use.
this seemingly unassuming person, Ehud,.
to deliver Israel from the hand of Aglon, king of Moab..
(speaking in foreign language).
So let's go through the rest of our story..
Verse 16, "Now Ehud had made a double-edged sword.
"about a cubit long, which he strapped.
"to his right thigh under his clothing..
"He presented the tribute to Aglon, king of Moab,.
"who was a very fat man..
"And after Ehud had presented the tribute,.
"he sent on their way those who had carried it..
"But on reaching the stone images near Gilgal,.
"he himself went back to Aglon and said,.
"Your Majesty, I have a secret message for you..
"And the king said to his attendants,.
"Leave us, and they all left.".
(speaking in foreign language).
So notice the satire here..
The author told us in verse 17 that Aglon,.
he was a very fat man..
And this is not fat shaming,.
but this is probably the results of him.
getting fat and over the tribute.
that he actually extorted from Israel.
over the past 18 years..
(speaking in foreign language).
And then starting on verse 20,.
everything went into slow motion..
See, Ehud then approached him.
while he was sitting alone in the upper room of his palace.
and said, "I have a message from God for you.".
And as the king rose from his seat,.
Ehud reached with his left hand,.
drew the sword from his right thigh,.
and plunged it into the king's belly..
Even the handle sank in after the blade.
and his bowels discharged..
Ehud did not pull out the sword.
and the fat closed in over it..

$^{321}$Then Ehud went out to the porch,.
he shut the doors of the upper room behind him.
and locked them..
(speaking in foreign language).
So here we actually see another trickery being revealed..
Ehud, he is not just delivering a tribute.
from the oppressed Israel,.
but he is God's means of deliverance.
for the people of Israel..
Ehud, he is not just a messenger with a secret message,.
but he is a messenger of God's judgment over Aglam..
(speaking in foreign language).
And so even though Aglam had died at this point,.
Israel's deliverance is not complete yet..
In the next few words,.
we see more trickery and satire being played out..
(speaking in foreign language).
We see here because Ehud,.
he locked the door before he left the room..
So when the king's servant came,.
they thought that their king was just having his time.
using the bathroom..
So instead of going in and checking on their king,.
they just waited outside..
And the smell of Aglam's bowel discharge.
probably contributed to this as well..
(speaking in foreign language).
So while the servants waited and waited outside,.
this gave Ehud time to actually escape.
and get back to his own people.
to rally them for the final blow..
(speaking in foreign language).
So let's pick up our story again from verse 28..
Follow me, Ehud ordered,.
for the Lord has given Moab, your enemy, into your hands..
So they followed him down.
and took possessions of the forts of the Jordan.
that led to Moab,.
and they allowed no one to cross over..
At that time, they struck down about 10,000 Moabites,.

$^{361}$all vigorous and strong, not one escaped..
That day, Moab was made subject to Israel,.
and the land had peace for 80 years..
(speaking in foreign language).
And so here we see a total reversal of the roles.
and a complete deliverance of Israel.
from the hands of their oppressors..
(speaking in foreign language).
So while in verse 13, we saw that Eglon attacked Israel.
and took possessions of the city of Palms,.
here in verse 28 and 29,.
Ehud and his men actually struck down.
about 10,000 Moabite warriors,.
and they took possessions of the forts of Jordan.
that led to Moab..
(speaking in foreign language).
And while in verse 14, the Israelites,.
they were subject to Eglon, king of Moab, for 18 years,.
here in verse 30, we saw that Moab actually became subject.
to Israel that day,.
and the Israelites and the land finally had peace.
for 80 years..
(speaking in foreign language).
And so once again, this show us that God's holiness.
and his love always work side by side in his rescues..
On one hand, God is holy and he is righteous..
So when his people sin against him,.
he cannot let sin go unpunished..
But on the other hand, God is also love..
He understands that our human condition.
and our need for a deliverer,.
to deliver us not just from the enemies that is around us,.
but more importantly, from the enemy that is within us,.
that is our sinful and rebellious nature..
(speaking in foreign language).
And God's holy love is like a thread.
that weaves through the entire story of the Bible..
(speaking in foreign language).
And God's love is like a thread.
that weaves through the entire story of the Bible..

$^{401}$(speaking in foreign language).
And God's love is like a thread.
that weaves through the entire story of the Bible..
We see his holy love in the way that he dealt with.
the very first human couples, Adam and Eve..
Because Adam and Eve sin,.
they have to be casted out from the garden of Eden.
because sin can have no presence in the midst of God,.
in the midst of a holy God..
But it is also out of love that they were banished.
because they would literally die.
if they continue to stay in Eden..
(speaking in foreign language).
And this theme of God's holy love.
carries through from Genesis to Exodus,.
into the promised land,.
and eventually to Israel's exiles and the return..
And finally, God's holy love reaches climax.
in the person of Jesus at the cross..
At the cross, we see that the wrath of God.
was being poured out on the sins of the world..
And yet it wasn't us who had to bear.
the consequences of our own sins,.
but it was Jesus, the son of God himself,.
taking on the sins of the world,.
satisfying the wrath of the Father.
because of His love for us..
(speaking in foreign language).
And so as we enter into a new year,.
wherever you are in your journey with God,.
I want to encourage all of us to take our next steps,.
to grow in having a right understanding of God.
as the embodiment of holy love.
so that we can have deeper trust in Him..
(speaking in foreign language).
So let's pray..
Lord, we carry with us all kinds of false images.
and wrong ideas when it comes to you all the time..
(speaking in foreign language).
And as we enter into a new year,.

$^{441}$would you remind us again with the truth.
that you are a holy and loving God?.
You did not spare your own son,.
but gave your only son, Jesus,.
in order to deliver us from our sinfulness.
and to give us eternal life..
(speaking in foreign language).
Help us to not take this gift for granted.
and teach us to live a life of repentance.
and a life of trust through the power of your Holy Spirit..
(speaking in foreign language).
And in your holy name we pray, amen..
(speaking in foreign language).
All right, brothers and sisters,.
thank you for joining us for the first Sunday of the year..
You guys are dismissed, so go in peace in your love..
(speaking in foreign language).
[MUSIC PLAYING].
\newpage



\section{詩篇 133:1-3}
\label{sec:6JWiBdm7MIc}
\textbf{Pilgrim's Community (Psalm 133) - Rev. Mark Chan 上行的群體(詩篇133\_1-3) - 陳崇基牧師}
\newline
\newline
連結: \href{https://youtube.com/watch?v=6JWiBdm7MIc}{\texttt{https://youtube.com/watch?v=6JWiBdm7MIc}} ~~~~ 語音日期: 2026-01-14
\newline
\newline
\hyperref[sec:kGtKq4doSFU]{\small{< < < PREV SERMON < < <}}
~
\hyperref[sec:index_chronic]{\small{[返順時目]}}
~
\hyperref[sec:index_scriptual]{\small{[返順卷目]}}
~
\hyperref[sec:4c2vxrJygVs]{\small{> > > NEXT SERMON > > >}}
\newline
\newline
詩篇 133:1-3
\newline
\begin{longtable}{cl}
\hline
\hline
章節 & 經文 (和合本修訂版)\\
\hline
133:1 & \begin{tabularx}{0.7\textwidth}{X} 看哪，弟兄和睦同住是何等的善，何等的美！ \end{tabularx} \\ \\ \relax
133:2 & \begin{tabularx}{0.7\textwidth}{X} 這好比那貴重的油澆在亞倫的頭上，流到鬍鬚，又流到他的衣襟； \end{tabularx} \\ \\ \relax
133:3 & \begin{tabularx}{0.7\textwidth}{X} 又好比黑門的甘露降在錫安山；因為在那裡有耶和華所命定的福，就是永遠的生命。 \end{tabularx} \\ \\
[1ex]
\hline
\hline
\end{longtable}
$^{1}$(speaking in foreign language).
- All right, guys, good morning..
(speaking in foreign language).
And as you know, I'm trying to pick out.
some of the famous Psalms in the Bible.
and to reflect upon it..
So today I would like to turn to Psalm 133..
And Psalm 133 is talking about a community.
of the people of God..
Let me read that first, okay?.
Psalm 133..
So you find some Psalms in the middle of the Bible, okay?.
In the middle of the Bible..
If you have a phone app or a Bible app,.
maybe I can invite you to turn to it right now..
Sometimes it's always good to search out.
the word of God ourselves, okay?.
So look into your phone or on the screen..
This is a Psalm, Psalm 133..
(speaking in foreign language).
I'm gonna read in English.
and then Ken's going to read in Chinese..
Psalm 133, a song of essence of David..
How good and pleasant it is.
when brothers live together in unity..
It is like precious oil poured on the head,.
running down on the beard,.
running down on Aaron's beard,.
down upon the collar of his robes..
It is as if the dew of Hermon.
were falling on Mount Zion..
For there the Lord bestows His blessing,.
even life forevermore..
(speaking in foreign language).
All right, so there are 150 Psalms in the Book of Psalms,.
but there were 15 Psalms, 15 Psalms, okay?.
120 to 134..
They have the title of Song of Essence..
And these are the Psalms that people use.
as they went up to Jerusalem, to Zion..

$^{41}$Jerusalem is Mount Zion..
Mount Zion is actually Jerusalem..
To worship, when they come to Mount Zion,.
to Jerusalem to worship together..
So these are the 15 Psalms that carry a sense of movement,.
moving like a pilgrimage, far away from Jerusalem,.
when they go toward near Jerusalem..
(speaking in foreign language).
So the first question you have to ask is,.
why is it called Song of Essence?.
You know, along with the other 14,.
you know, there are 15 Psalms, right?.
Why is it called Song of Essence?.
Essence meaning going up, because it is geography,.
geography of the Bible..
Because Jerusalem is situated on a hill,.
it is within the hill country of the land of Canaan..
So no matter which direction you're going to Jerusalem,.
where, which direction you're coming from,.
when you go to Jerusalem, you have to go up,.
you know, you have to go up..
And even today, modern Israelis, when they,.
the Jewish people, when they move or immigrate.
from far land, faraway land, to back to Jerusalem,.
back to Israel, the process of immigration.
is called going up, you know, it's called going up..
And the name for the going up,.
the process of immigration is called Aliyah..
You know, Aliyah comes from Allah, means going up..
So even today, the modern immigration process.
is called going up Aliyah,.
because when they go to Jerusalem,.
when they go back to Jerusalem,.
they always climb up and up and up, and they go uphill..
(speaking in foreign language).
So in the Old Testament, when the people of God.
commanded to go to Jerusalem three times a year,.
as you can tell from here, Deuteronomy,.
one of the books, one of the five books of Moses,.
Deuteronomy chapter 16, it says, you know,.

$^{81}$three times you have to go to Jerusalem.
from where you live, probably you live far away.
from Jerusalem, but three times of the year,.
you have to go to Jerusalem..
And when you go there, listen to the words here,.
three times a year, all your men,.
of course, men are the representatives of the family,.
and with the, you know, female and all the children,.
all of you guys must appear before the Lord your God.
at the place he will choose, which is Jerusalem eventually,.
at the Feast of Unleavened Bread, which is Passover,.
the Feast of Weeks, which is Pentecost,.
and the Feast of Tabernacles, Sukkot..
And no man, no one should appear before the Lord.
when you worship empty-handed,.
and each of you must bring a gift in proportion.
to the way the Lord your God has blessed you..
So three times of the year, they have to go up there.
as one people..
Now, even though they're Jewish people,.
they were divided into 12 tribes..
Those are their tribal identities,.
and sometimes it's very, very precious to them..
It's a source of pride for them..
But three times of the year,.
they have to make a journey, annual journey,.
and that's a special journey..
You have to lay aside your own local identity,.
and you have to come before the Lord.
at one place, the place that God has chosen,.
and to be the people of God worshiping..
(speaking in foreign language).
,.
so these Psalms, 15 Psalms of Ascension,.
are likely passed down as songs.
that used from generation to generation.
when they go up to Jerusalem,.
when they are on their way to Jerusalem to worship..
So they will sing this song,.
and probably resting on the road, in the Airbnb,.

$^{121}$or walking together, hiking together,.
and the destination is Jerusalem..
And after arriving at Jerusalem,.
when the crowd of people are crowding the city,.
they probably will use this Psalm,.
Psalms of Ascension, to worship God..
(speaking in foreign language).
So I would like to take verse one,.
and verse by verse,.
and we're gonna look at this Psalm together..
Verse one, it says, "How good and pleasant,.
how good and beautiful it is,.
where brothers," of course, including sisters,.
"live together in unity.".
(speaking in foreign language).
And then the Psalm in here, verse one,.
begins with some exclamation, you know,.
the exclamation mark..
There were three of them, can you find it?.
"Behold, look, how," you know?.
And then, "How good and how pleasant,".
you know, there are three times the psalmist.
trying to say, "Wow, you know, how precious.
and how beautiful it is when people are gathered together.".
(speaking in foreign language).
And then, you know, I don't care how many people here,.
but, you know, as we gather together,.
you know, two or three, you say,.
with the psalmist, and say, "Wow," right?.
"How good it is to gather together.".
Okay, so the double how here,.
how good and how pleasant, is pretty rare..
One example I can find is from,.
from the, let's see, I don't have it on screen,.
but one example is from Song of Songs, chapter seven,.
where the couple are praising each other,.
and the partner is saying to the other partner,.
and saying, "How beautiful you are,.
and how pleasant you are.".
(speaking in foreign language).

$^{161}$Now, you know, we try to do some imagination here..
When you turn to your neighbor,.
let's say now you turn to a neighbor,.
either front or back, or left or right,.
you say, you know, "You are handsome," you know,.
"or you are so beautiful.".
You know, I don't want you to do it now,.
but, you know, imagine, you know,.
if someone, I thought, turned to Ken,.
and say, "Dude, you are so handsome.".
- Thank you..
(laughing).
- That is the feeling you try to get, all right?.
You are amazing, or you are absolutely stunning, you know?.
So this song is trying to, you know,.
talk to us, or invite us into this kind of emotions..
It's not romantic love, but it's something really pleasant,.
something really, really good..
(speaking in foreign language).
And this is talking about God's people gathering in worship,.
worshiping the same Lord..
(speaking in foreign language).
And then here, translated unity,.
is actually coming from Hebrew word, Yahal..
And Yahal means, you know, something together,.
the togetherness, you know, standing together,.
working together, moving as one body..
You know, in the time of,.
in time between Old Testament and New Testament,.
there was a community of Jewish people.
called Qumran community..
And they, the guys who wrote the Deathly Scroll,.
you heard about that?.
Deathly Scroll, it's from, come from their hands..
And this Qumran community, Jewish community,.
they left Jerusalem because Jerusalem was corrupt,.
and they moved to the wilderness..
And they described their own group as a community of Yahal..
And that means there was a shared purpose.
and shared destination, and they worshiped God together..

$^{201}$That's a Yahal..
(speaking in foreign language).
And then the Psalm says, "How good it is.".
And that is a very frequent word in the Bible..
And do you remember the Bible first mentioned.
something is good?.
Where do you find it?.
Exactly, Genesis chapter one..
God is good..
And then the Bible says, "How good it is.".
And that is a very frequent word in the Bible..
And do you remember the Bible first mentioned.
something is good?.
God created the entire world, the entire cosmos,.
and repeatedly God says, "It is good.".
And at the end, he says, "Very, very good.".
Six days of creating and setting up the cosmos.
in Genesis one, and God creates and call it good..
And then the first time God says, "Not good,".
where do you find it?.
Genesis two, right?.
Remember that?.
When God saw Adam, the humanity, living alone,.
God says, "It is not good for humanity to be alone..
"Not good for man to be alone.".
So in creation, God turns something not good into good.
because he set up the cosmos, and then chapter two,.
he set up a community for mankind..
(speaking in foreign language).
Good..
(speaking in foreign language).
Now, just a matter of illustration..
If you wanna keep a fire going, keep a fire burning,.
keep the flame burning,.
if you wanna want a fire to burn strongly,.
you wanna keep all the wood pieces together, right?.
Together..
When they burn together, then you will see a bigger fire..
But if you take one of the stick out and put it aside,.
throw it aside, then the stick itself,.

$^{241}$while it was burning, you can throw it aside,.
and then it will go out pretty quickly..
(speaking in foreign language).
See, the fire with the wood burning together.
illustrates our lives as the people of God..
We need to burn and come together to encourage one another.
so that it will not go out by itself..
So notice the Psalm does not say,.
okay, the people are perfect,.
and that's why it is so good for them to get together..
It's not talking about we are unusually holy..
You know, we are not holy..
We are sinful..
However, from God's eyes, it is beautiful,.
it is good for brother and sisters,.
for his people to gather together..
(speaking in foreign language).
So something that is good, and we need to see it.
from God's sight..
Verse two, it says, "It is like precious olive oil.
poured on the head, running down on the beer,.
running down on Aaron's beer,.
down upon his collar to his robes.".
You know, verse two..
(speaking in foreign language).
This imagery may not sound appealing to us.
as modern people, all right?.
It sounds messy, it looks messy,.
but you know, you can see, you know,.
the oil pouring down on Aaron, you know,.
is in abundance, right?.
So you need to ask the question,.
as you look at verse two, every community,.
every group of people has some representatives,.
some representative leaders..
Sometimes it is a famous person you grew up,.
you know, admitting someone, even you don't know her.
or you don't know him, but you know,.
you feel like, oh, they represent you..
Maybe an artist, maybe a cultural icon..

$^{281}$If they do well and I do well, right?.
And kind of like last night,.
did you see the football game last night?.
Yeah, right?.
Right, did you see it?.
It's very good to see it, right?.
It's the dramatic turnaround for Bears..
And you know, I see today when I show up at church,.
I see people are happy..
We haven't felt happy for a long time, right?.
Chicago people..
Anyways, so we have our own representatives..
And you know, in the Old Testament, in Israel,.
the representative figure was the high priest Aaron..
(speaking in foreign language).
So when you trace the story of the Bible.
from the Old Testament,.
when the people of Israel came out of Egypt,.
they arrived at Mount Sinai and they received the law.
and then the tabernacle was first established.
in the wilderness because God wanted to dwell within them..
And Moses anointed the tabernacle and his furniture.
and all the items and set up..
And the climax comes to anointing the high priest Aaron.
so that he can represent the whole people of God Israel.
to serve God together..
So in the book of Leviticus, it says,.
Moses was pouring anointing olive oil on the head of Aaron.
and to anoint him to consecrate him..
(speaking in foreign language).
And that was the most sacred moment of the entire Bible..
And even says that, you know, the glory of God, you know,.
shines from the tabernacle and consume the offering.
on the altar and the people shout out loud.
and fell on their faces and they worship God together..
That was the sacred moment of the Old Testament..
(speaking in foreign language).
, and you can imagine that was a sacred,.
the most precious moment of all time in the Bible..
But Psalm 133 turns it around a little bit..

$^{321}$It says the kind of holiness, sacredness.
is actually present when the people of God are united..
When they go up together to Jerusalem.
to worship God together, it's the togetherness.
that is sacred..
It's just like, it is kind of like the first anointing.
of the high priest Aaron,.
when the glory of God finally comes..
You know, it comes not because someone's being anointed..
It comes, the glory of God shines for.
when the people of God are united.
and worshiping God together..
(speaking in foreign language).
And the songs of ascent are probably, you know,.
collected and shaped together during the exile.
of the people..
And when the people of Israel,.
when the Jewish people had no political power.
and they were scattered around the world,.
yet even then, the gatherness of worship,.
when the people coming together to worship,.
the Psalm says, the worship, gathering worship itself.
is a sacred, holy thing..
(speaking in foreign language).
You know, I have been growing up in the church.
for decades, for decades..
Oh man, I'm getting old..
And the church is never perfect..
The church is never perfect..
And sometimes, you know, in my experience of the church,.
the church can break my heart..
The church can hurt me and disappoint me..
You know, I must say with a disclaimer,.
I must include myself, you know..
I disappoint someone else..
I break someone's heart..
You know, I have to include myself..
The church is not perfect..
But, you know, along the years,.
I have never given up on the church..

$^{361}$Never given up on the church..
It's not because the church is perfect..
It's not because church is always a happy place..
But because I understand the church is the very plan of God..
As long as God is God, God's church will remain..
(speaking in foreign language).
So I understand my identity as the member of the church..
I have a calling..
And the calling is a sacred calling..
It is a community Christ purchased.
with His blood and His sacrifice..
It's not because of me..
It's because of what Christ has done..
And He redeemed a people and called it the church,.
called it His body..
And the church belongs to Christ Jesus..
(speaking in foreign language).
So verse three here, not only the gathering of the people,.
the church is sacred..
Not because we're perfect,.
but what God's grace has done and make it sacred..
Verse three, it says,.
"It is as if the gathering of people worshiping together,.
as if the two of Hermon, Mount Hermon,.
were falling on Mount Zion, Jerusalem..
For there the Lord bestows His blessing,.
even life forevermore.".
(speaking in foreign language).
Now this is the second image of the Psalm..
No, the first one is the anointing of the high priest Aaron..
The second images here is a geographical image..
You know, when people are united,.
it is like two of Mount Hermon falling on Mount Zion..
(speaking in foreign language).
Now Mount Hermon is the highest mountain.
in the land of Israel..
It is about 9,000 feet tall..
And it is sometimes called the Snow-Capped Mountain.
or Gray-Haired Mountain..
You know why?.

$^{401}$Because the mountain is so elevated.
and it's always capped by snow..
In much of the year, it's a snow-capped mountain..
(speaking in foreign language).
In the land of Palestine,.
from April, around April, May to October,.
that is the hot and dry summer..
And there is very little rain, no rain at all..
And everything is dry and the ground,.
the soil becomes dry and the crops can die.
unless one thing happen..
And that is the morning moisture..
You know, when you wake up very early in the morning,.
you look at your yard, right?.
There's a layer of moisture..
And because of that, because the mist, the dew,.
that bring water to the soil and to sustain us..
That is life, the dew is life..
(speaking in foreign language).
So in the imagination of the Psalm,.
it says when the dew of Mount Hermon,.
the highest mountain, the Snow-Capped Mountain,.
comes upon the land and even riches Jerusalem.
and Mount Zion, you know, you get the imagery..
It talks about life-giving, refreshment,.
you know, fertility and life..
(speaking in foreign language).
So, you know, when we get together,.
when the people of God get together,.
worship God together, it's kind of like a first taste,.
you know, it's like the appetizer.
of what heaven is like, what eternal life is like..
You know, when we get together to worship,.
to love God together, we will taste.
what eternal life is all about..
(speaking in foreign language).
Now, you remember toward the end of last year,.
we had a baptism here, right?.
And three people got baptized from our congregation..
And at baptism, sometimes we recite the Apostles' Creed..

$^{441}$Remember that, Apostles' Creed?.
And Apostles' Creed is a summary of faith.
and the Creed has many different articles.
to summarize the most basic elements of Christian faith..
And we call them Apostles' Creed..
And there's a summary and it mentions, as you can tell,.
it mentions the church twice..
It says, "I believe the Holy Catholic Church.
"and I believe the chameleon of saints.".
(speaking in foreign language).
See, the article says, "I believe in Catholic Church.".
Most of the modern people, when we heard the word Catholic,.
we immediately think of the Roman Catholic Church, right?.
Catholicism..
No, that's not what Apostles' Creed means..
It says Catholic because the word Catholic.
actually means universal, worldwide church..
And I believe God's will is to set up the Catholic Church,.
which is universal church..
(speaking in foreign language).
See, we, as Christians, we profess the name of Jesus..
We belong to that Catholic Church,.
that one worldwide church..
And also, another article, it says,.
"The chameleon of saints that I believe.".
Now, what does it mean?.
What does it mean?.
It means all Christians, chameleon,.
we share something together..
We're united to Christ,.
and therefore we are connected somehow with each other..
It's like you, you know,.
how many of you like going hiking?.
How many of you?.
No?.
Oh my God..
Okay, all right, only a few..
It's like you go on hiking, on a hiking trail, you know,.
and then you go from your starting point,.
and then you have a destination to go to, right?.

$^{481}$And then along the trail, along the path,.
and then you meet some people on the way, right?.
Maybe you'll meet some people.
going in a different direction, you know?.
They eventually will go to a different destination..
Now, you don't have much to share with them..
However, if you, on the trip,.
you meet some people going to the same destination,.
and then that was the moment you go,.
aha, you know, I'm going to that destination,.
and you are going to that destination too..
And that's why we are on the way together..
When you meet some people going to the same destination,.
you find that, okay,.
you have so much more to share with them..
(speaking in foreign language).
So it says, we believe the chameleon,.
the connectedness, togetherness of saints..
Now, the word saints means all Christians..
It does not mean that we are holy..
We're not holy, we're not perfect, right?.
But we are saints somehow in Christ.
because Christ made us holy and righteous,.
all because Christ has done..
So chameleon means we have shared life..
Chameleon means we have fellowship,.
and we share one faith, we share one worship,.
and we share one salvation..
(speaking in foreign language).
And together we take this psalm,.
and I have some applications for us..
I think the first application from Psalm 133 is,.
thank God for other Christians,.
for other brothers and sisters, right?.
I would say to God, dear God, thank you.
that I am not alone, I'm not alone hiking..
I'm with other people..
Thank God for brother and sister who pray for me,.
who are being with me, who are walking alongside me,.
encourage me, and just have shared life with me..

$^{521}$(speaking in foreign language).
Secondly, we should thank God that, you know,.
I am actually part of God's people..
You know, thank God that, you know,.
God did not save me as an individual alone..
God does not save me as one person..
He placed me into his family..
He placed me into his big church,.
the church, the people of God..
(speaking in foreign language).
And number three, I think the application from the psalm.
is that we need to invest time..
It is good to invest time to learn about.
and to remember the stories of other people, right?.
Introduce yourself, you know, get to know each other,.
not only the name, but, you know, at least,.
we need to hear one another's stories..
(speaking in foreign language).
And then number four, it is good and pleasant..
How good it is to invest your heart and your energy.
in relationship with each other,.
enjoy the fellowship, to set aside time,.
prioritize our time and energy,.
to come together, to join together, and to grow together..
(speaking in foreign language).
So one little application at the end of the service.
will be that if we dismiss you guys,.
today, you need to do one thing, okay?.
Do one thing..
(speaking in foreign language).
Right after service, you shake hand with someone..
Shake hand with someone..
(speaking in foreign language).
And then you say, "I'm going to shake hands with you.".
(speaking in foreign language).
And then you say, "Thank God for you.".
(speaking in foreign language).
Even though you don't like that person..
(speaking in foreign language).
Or you like that person, you know, shake hand,.

$^{561}$and you know, pat on the shoulder,.
and say, "Thank God for you.".
(speaking in foreign language).
All right, let's pray together..
(speaking in foreign language).
Dear Father God, thank you,.
because you set up the church in grace..
(speaking in foreign language).
Not because we are perfect..
(speaking in foreign language).
But Christ has offered his sacrifice,.
perfect sacrifice for us..
(speaking in foreign language).
And thank you, we are part of the church..
(speaking in foreign language).
And help us to love each other and to grow together..
(speaking in foreign language).
In Christ's name we pray..
(speaking in foreign language).
Amen..
So you're dismissed and you're loved,.
and make sure you do that little homework today, okay?.
At this moment, all right?.
(speaking in foreign language).
Yeah, you know, I may not like you,.
but you know, thank God for you..
(laughing).
(speaking in foreign language).
Okay, let's do that..
Let's stand up and shake hands and bless each other..
(speaking in foreign language).
Thank God for you..
(speaking in foreign language).
\newpage



\section{詩篇 133:1-3}
\label{sec:4c2vxrJygVs}
\textbf{上行的群體(詩篇133\_1-3) - 陳崇基牧師}
\newline
\newline
連結: \href{https://youtube.com/watch?v=4c2vxrJygVs}{\texttt{https://youtube.com/watch?v=4c2vxrJygVs}} ~~~~ 語音日期: 2026-01-14
\newline
\newline
\hyperref[sec:6JWiBdm7MIc]{\small{< < < PREV SERMON < < <}}
~
\hyperref[sec:index_chronic]{\small{[返順時目]}}
~
\hyperref[sec:index_scriptual]{\small{[返順卷目]}}
~
\hyperref[sec:LP_MCyZeclI]{\small{> > > NEXT SERMON > > >}}
\newline
\newline
詩篇 133:1-3
\newline
\begin{longtable}{cl}
\hline
\hline
章節 & 經文 (和合本修訂版)\\
\hline
133:1 & \begin{tabularx}{0.7\textwidth}{X} 看哪，弟兄和睦同住是何等的善，何等的美！ \end{tabularx} \\ \\ \relax
133:2 & \begin{tabularx}{0.7\textwidth}{X} 這好比那貴重的油澆在亞倫的頭上，流到鬍鬚，又流到他的衣襟； \end{tabularx} \\ \\ \relax
133:3 & \begin{tabularx}{0.7\textwidth}{X} 又好比黑門的甘露降在錫安山；因為在那裡有耶和華所命定的福，就是永遠的生命。 \end{tabularx} \\ \\
[1ex]
\hline
\hline
\end{longtable}
$^{1}$好 多謝傑少.
我還以為傑少可以說道.
我留下這篇道給下星期.
所以今天決定是對的.
有一篇比較短的道.
大家知道我今天是說詩篇.
我選了一些詩篇.
跟大家來思想.
今天這篇詩篇是第133篇.
如果有聖經的話.
不妨打開這篇詩篇.
因為我們學校找聖經是對的.
打開之後讓我讀.
讀的時候請留意這篇詩篇.
看詩篇看詩歌.
或者我們剛才唱的詩歌.
都是要留意意境.
為什麼會有這樣的意境.
以及意境的意義在哪裡.
我們需要反省.
或者需要留意的是什麼.
我們看看這篇詩篇.
很短的 三節.
第133篇.
大衛上行之師.
看那弟兄和木同住.
是何等善 何等美.
就好比那貴重的油.
飄在阿倫的頭上.
流到胡蘇 又流到他的衣襟.
又好比黑門的金爐.
降在石安山.
因為在那裡有耶和華.
所命定的福.
就是永遠的生命.
我們看這篇詩歌的時候.
需要留意幾樣東西.
第一樣東西就是標題.
很多時候我們沒有留意標題.
標題是什麼.

$^{41}$上行之師 大衛上行之師.
詩集裡面有多少首詩歌.
詩篇裡面有多少首詩歌.
有150篇.
有150首的詩歌.
上一次我提過.
這些詩歌都是人的說話.
人對上帝所表達的說話.
無論是一些討民.
一些哀歌 讚美詩詞等等.
是人向上帝.
是上行的.
跟其他聖經不同.
其他聖經是下行的.
上帝藉著先知說話.
或上帝藉著不同的方法.
來對神的子民說話.
詩篇有一點不同.
詩篇是我們向上帝說話.
好像我們讚美唱詩的時候.
是我們向上帝說話.
你有沒有留意.
就算第一次來教會.
或者來了教會很久.
很多詩歌的題材.
那個所謂的pronoun.
那個代名詞.
都是第二身的.
You.
你.
主耶和華你....
這個就是詩歌那種方法.
在詩集裡面.
有150首這些的詩歌.
而其中有15首是很特別的.
有一個標題給了他.
叫做上行之詩.
是120至134.
有15首這些的詩歌.
都是叫做上行之詩.

$^{81}$這些上行之詩.
有什麼特別的地方呢.
就是他應該是當時神的子民.
他們在散居不同的地方.
他們要上去耶路撒冷.
來去敬拜.
上去耶路撒冷朝聖的時候.
他們可能會唱的詩歌.
說不定隨著年月.
不同的世代.
他們由不同的地方上耶路撒冷.
就是石安山.
耶路撒冷就是石安山.
去到那裡敬拜的時候.
途中上路的時候.
可能休息或者正在走的時候.
他們會唱一些的詩歌.
來去鼓勵他們.
或者給他們一個目標.
使他們遠地上到耶路撒冷的時候.
有一種的觸動.
上到去耶和華所揀選的地方.
去到敬拜的時候.
他們就有一種培養的心靈.
來去敬拜主.
所以就叫做上行之詩.
Song of Ascending.
Ascending是什麼意思呢.
這個我們住在芝加哥的人.
不是很明白的.
有誰在這裡開手波車的.
OK.
這裡開手波車不是很過癮的.
因為沒有斜路.
如果前面的車在那裡.
在那裡搖晃的時候.
他們應該是在開手波車.
因為他們用一波在搖晃.
因為你的車在斜路的時候.
你的車就會向後.

$^{121}$然後你就用一個鑽頭.
大概撞著.
明白我的意思嗎.
不明白的話.
你下次跟我一起去.
因為有山的緣故.
所以你就有那個所謂的斜坡.
在芝加哥是沒有的.
我記得有一次教我兒子開棍波車.
要找一條斜路給他.
很難找啊在芝加哥.
我找到了.
你找到了.
OK 等一下告訴我.
找到斜路.
我找了幾個圈都找不到.
我在哪裡找一個斜路.
幫他學一下怎樣一波起步呢.
是很難找的.
因為芝加哥Midwest.
是一個Great Plain.
是一個大的平原.
但是以前我在加州住.
就不是這樣的.
經常有斜坡.
無論你加州西北.
或者以前在Sacramento住.
甚至三環市.
你去過舊金山三環市.
斜坡就像過山車一樣.
但是耶路撒冷.
以色列巴勒斯坦地.
就是山多平地少.
所以有很多山脈.
有很多群山.
無論你去哪一個方向.
去耶路撒冷.
都是要上山.
就算你比耶路撒冷更遠的地方.
你更高山的話.

$^{161}$你就近耶路撒冷的時候.
就快去到的時候.
你都需要上山的.
因為耶路撒冷就在山上.
雖然是一個小的山上.
所以Song of Ascent的意思.
就是上山的意思.
上耶路撒冷就等於上山.
去石安山就等於要上去.
所以就是Song of Ascent.
甚至今天那些現代的猶太人.
今天的現代以色列的猶太人.
他們說所謂移民.
就是說他們散居的地方.
他們移民回去以色列地.
做以色列的公民.
Israelis的時候.
他們那個我們叫.
美國叫Naturalization.
那個移民的過程.
他們都是上去的意思.
那個字叫做Aliyah.
因為上行就是Allah.
所以他們用回現代的希伯來語.
就是講Aliyah.
就是你做了Aliyah沒有.
你上行了沒有.
因為什麼呢.
因為你回到你的國土的話.
就好像上行一樣.
回到耶路撒冷一樣.
這個上行之師就是一個這樣的背景.
在《申命記》裡面.
律法書提過幾次.
叫以色列民.
當然用南來代表.
當時的文化背景是這樣.
他就說所有的南丁.
要在三個節期裡面.
每年上耶路撒冷.

$^{201}$Go up to Jerusalem.
那個就是上帝所揀選.
要去朝見他.
要去朝聖的地方.
並且你上去的時候.
你不要空手去.
你個人要帶著自己.
手中的能力所能夠奉獻的.
去照著耶路撒冷所賜給你的福.
你都要奉獻上去.
一年三次是什麼節期呢.
除考節.
就是說Passover.
就是預節.
七七節就是五旬節.
七七四十九.
第五十日就是五旬節.
Pentecost.
是另一個修國節.
然後就是駐彭節.
這個Tabernacle.
Sukkot的時候.
一年三次的時候.
你就需要帶著你的全家人.
上去耶路撒冷.
你都要去敬拜的.
這個就應該是上行之師的背景.
這個很有趣.
上行去一個地方.
就不是按照你的身份.
不是按照你的支派.
因為以色列人特別是在舊約裡面.
他們引以為傲的.
很多時候的一個身份.
就是我屬於哪個支派.
對嗎.
等於我們華人都是這樣的.
你姓什麼的.
我們有時候會覺得.
我們的姓是特別厲害的.

$^{241}$我姓陳.
從小到大.
我爸都跟我說姓陳厲害.
很多人姓陳.
人山人海.
OK.
不過這裡就不是.
這裡姓梅.
這裡每個人都吐口水.
都死了.
姓梅的有很多.
有一種引以為傲的身份.
以前舊約的子民都是一樣.
你屬哪個支派.
猶大支派.
變雅敏支派等等.
這些是你的身份的象徵.
但是上行的時候.
就不再說你自己.
引以為傲的身份的表徵.
因為這個是耶和華所揀選的地方.
不是你所揀選的地方.
不是你的鄉下特別輕鄉一點.
而是耶和華.
特意用他的恩典來揀選的.
你就要去那裡一年三次的特別的上行.
因為那裡是耶和華所揀選.
就是耶路撒冷石安山.
所以上行之詩有十五首.
就應該是他們上去朝聖耶路撒冷.
上去石安山的時候敬拜.
所留下來的一些詩歌.
有十五首.
有時候可能路上休息.
或者趕路.
他們就會唱這些詩歌.
這個其中一首.
詩篇第133篇.
你就可以用這樣的意境.
這樣的背景來看上行之詩.

$^{281}$133首這裡說到什麼呢.
他說漢啊弟兄姊妹.
和睦同住.
是何等的善.
何等的美.
你看到這裡一開始有幾個讚嘆的詞彙.
你找到嗎.
讚嘆的詞彙.
如果你用現代的語言來翻譯的話.
你就是可以用.
How.
WOW.
How beautiful.
How nice.
何等.
和漢啊三個字.
特別是何等.
何等怎樣.
何等美.
何等善.
這樣雙倍出現.
一起出現.
在聖經裡面是比較少的.
其中一段我找到的就是雅歌.
雅歌有一段那裡說什麼呢.
這裡沒有打出來.
雅歌第七章.
那裡的情侶.
大家給一支讚大家的時候.
那個情侶說什麼呢.
他說我親愛的喜樂的女子.
你何等美麗.
何等令人喜悅.
你聽到兩個何等嗎.
何等靚.
和何等pleasing.
Right.
這些是讚你的詩歌.
讚你的詩詞.
有沒有被人讚過最近.

$^{321}$你嘗試讚一下你旁邊那個.
說你很正.
你真的超正.
OK.
你現在不需要說.
但你心裡面想一下也好.
無論他說假話.
或者說真的.
你心裡面都覺得feel good.
對嗎.
你讚他說你真的超正.
這些用在談情說愛裡面.
雅歌.
但這裡就用在弟兄姊妹.
神的子民一起的時候.
大家敬拜的時候.
是很正和超正的東西.
特別是在上帝眼中超正的東西.
那個和睦同住.
舊的翻譯是和睦同居.
同居就不是很好.
和睦同住.
可以用那個意境.
大家上路的時候.
可能在一些休息站.
或者一些地方.
他們可能book個airbnb.
大家在這裡睡一晚.
然後我們繼續上路.
或者甚至去到終點.
耶路撒冷目的地的時候.
因為很多人在不同的地方.
聚集在耶路撒冷.
所以就迫爆了城市.
然後在城市以外.
他們紮營等等.
大家彼此的同住.
有沒有試過去紮營.
有誰喜歡露營.
這裡我們很少.

$^{361}$都是紮男紮女.
你去紮營的時候.
旁邊那個人又紮營.
然後你聽到他鼻鼾聲.
或者你聽到他說什麼等等.
就是這樣的意境.
和睦同住.
那個同字是一行.
這個字.
一行.
希伯文是一行.
如果你留意到一些聖經歷史的背景.
在良若之間.
有一群群體.
猶太人.
他們離開了耶路撒冷.
因為耶路撒冷大壞.
他們去了一個地方叫昆蘭.
在一個曠野的地方.
他們有一行.
他們說他們自己是一行.
說我們是一個合一的.
一起作戰.
或者一起向著某一個目標同攻的那群人.
那個叫昆蘭社團.
如果你聽過死海古卷.
就是那群人寫出來的.
就是在那個曠野的地方寫出來的.
一個社團.
一個一行.
形容他們自己是一個共同的目標.
和大家同一個的.
共同的集體的一個身份.
他說好.
這個是何等的美.
何等的善.
是一個good字.
你有沒有想過舊約聖經.
整本聖經裡面第一次出good字.
是什麼時候.

$^{401}$創世記.
沒錯第一章.
創世記.
連續好幾次.
上帝所創立的秩序.
井井有條的.
慎看著事.
good.
慎看著事.
好的.
創世記第一章.
然後第一次上帝說不好的.
在哪裡.
創世記第二章.
第一章說到好呀好呀.
令人都很甚好.
very good.
但到第二章的時候.
他就說了.
not good.
有什麼是not good.
人獨居不好.
人類獨居不好的意思.
人類獨居不好.
可以為他設立社群.
讓他們走在一起.
有community.
這個創世記第二章裡面所說.
所以上帝所設立的東西是good的.
特別是community是good的.
我記得我初信的時候.
以前很小的時候.
十幾歲的時候.
有一個牧師說了一個比喻.
我到今天都記得.
他說基督徒為什麼要在教會裡面成為一份子.
還有要參與投入教會.
他說就好像一個火堆.
有很多木柴.
你要把木柴放在一起的時候.

$^{441}$它才會燒得紅.
如果你把木柴扔開一邊的話.
你就變成什麼.
你就很快熄滅.
你就等於廢柴.
這個比喻不知為何.
我一向都想到.
一向都印象深刻.
所以在教會.
我一向都覺得自己是一份子.
就算我旅行.
去不同的教會一次.
我知道去完之後我都不會再去.
或者別人覺得我去完之後.
都不想再叫我去.
我都覺得自己是一份子.
我都覺得自己是要參與.
在那裡去幫忙.
如果不是很認識我的人.
就覺得很不認識.
很八卦.
我去到那個教會.
那些音響不是很行.
然後我就走到後面去看看.
無論如何.
這裡說是甚好的事.
弟兄姊妹一齊去建立一個群體.
一起同住在一個群體裡面.
去到建立的時候.
是甚看著是好.
是好的.
第二節.
就好比那貴重的油.
曉在阿倫的頭上.
流到鬍鬚流到他的衣襟.
這句就很特別.
因為我們不覺得是很漂亮的東西.
是嗎.
那麼噁心.
你想想你很多鬍鬚.

$^{481}$然後用橄欖油.
在你的頭上.
這裡假設你頭上仍然有頭髮.
一直這樣流下去.
然後鬍鬚已經黏著油.
然後再流到衣襟上.
這個畫面漂亮嗎.
其實很漂亮.
其實他想表達的起碼第一件事.
就是夠多.
對嗎.
不是一杯.
是整桶油倒下去.
流到他周圍都是.
是abundance的意思.
是豐盛的意思.
每一個群體都有代表性人物.
對嗎.
每一個群體.
譬如香港人.
香港人都有不同的明星.
或者不同的代表人物.
可能你從小到大.
你聽他的歌.
聽他的藝人.
你就覺得這個人是代表性的.
是嗎.
在我的那一代.
其中一個代表性的人物.
是一個很sad的story.
我記得這個代表性人物.
小時候.
不幸小時候.
有些朋友跟我說.
他馬上哭了.
在我的那一代.
你知道那個是誰嗎.
Beyond.
黃家駒.
是不是很sad.

$^{521}$很sad吧.
就是那些年輕人.
不知道是誰.
黃家駒還是誰.
Anyway.
你Google一下就知道.
每一個群體都有代表性人物.
對以色列人來說.
代表性人物是誰.
就是阿倫.
這個就是阿倫.
阿倫的頭上.
阿倫是做什麼的.
阿倫是以色列國的第一個大祭司.
The first high priest.
你記得舊約的故事.
就是以色列人出埃及之後.
上帝用恩典拯救他們.
然後他們去到西太山.
接受了上帝的律法.
十誡等等.
然後就有會幕的設立.
因為上帝說要住在你們中間.
所以就有不同的制度.
去被成立.
其中一個制度就是priesthood.
就是大祭司.
祭司這件事.
第一個大祭司當然就是.
Aaron.
阿倫.
在利未記那裡說到.
利未記那裡說到什麼呢.
他們在這個會幕設立的時候.
他們就將高柔.
倒在阿倫的頭上.
讓他能夠成為神聖.
以至他能夠在會幕裡面.
做大祭司的功績.
然後那裡有這樣的描述.

$^{561}$他說他們進入會幕的時候.
阿倫進入會幕的時候.
耶和華的榮光.
就向全體的百姓顯現.
然後有火光.
從耶和華面前出來.
從會幕的致勝所出來.
然後就焚燒壇上的祭物.
百姓一看到的時候就歡呼.
然後面伏於地.
去敬拜上帝.
這個是第一個最sacred的moment.
第一個最神聖的時刻.
這裡說什麼呢.
詩篇.
就是在說這個情況.
高立,阿倫作大祭司.
這個是最榮耀的一個時間.
但是詩篇第133篇.
就很有趣.
他說什麼呢.
他說大祭司被按立的時候.
沒錯是sacred.
但是可能都不及一幫神的子民.
弟兄姊妹一同來敬拜上帝.
一同來到朝聖的時候.
這個時刻就有如什麼.
當時的good old days.
上帝的榮光從會幕裡面去展現出來.
降下然後全體百姓來到去.
俯伏敬拜.
That was the sacred moment.
但是當我們聚集的時候.
這篇詩篇說.
This is the sacred moment.
你有沒有想過來教會.
一起敬拜的時候.
原來是一個這麼sacred的moment.
sacred不是因為我們.
個個都很完美.

$^{601}$唱歌唱得對.
唱得好聽.
或者我們個個的身心靈都好好準備.
我們沒有什麼破碎的東西.
完美的人來到上帝面前.
No No No.
原來單單是因為上帝的恩典.
弟兄姊妹能夠在一起.
在敬拜.
在教會裡面.
一起所謂同住和睦的時候.
這個是very sacred.
不是單止good and beautiful.
是the most sacred thing.
你記得我說我從年輕到現在.
我都覺得教會是我心所在的地方.
沒錯.
教會很多時候令我會心痛的.
也令我很失望的.
教會會.
這個應該包括我自己在.
我都會令人失望.
因為我都是教會的一份子.
我都是一個很殘破.
有很多缺陷的人.
我會令人失望會令人心痛.
你會令我失望你會令人心痛.
但是由年輕到現在.
我都從來沒有放棄過教會.
因為我體會到.
這個是上帝的心意.
上帝很多時候做事很有趣.
上帝做的事未必一定是最有效率.
未必一定是最有效.
但他偏偏就是揀選一些東西.
用他的恩典去運作.
你想想耶路撒冷石安山.
是不是全世界最風光最美的山脈.
如果耶路撒冷知道.
it's not special at all.

$^{641}$但是上帝揀選.
是他的chosen grace.
教會其實都是一樣.
表面看教會其實有什麼特別.
一點都沒有特別.
甚至在社會裡面覺得.
這班人什麼料子.
但這個是上帝恩召的揀選.
是基督用他的血去犧牲.
在十字架上犧牲.
來贖回我們這班人.
這就是最sacred的東西.
是因為上帝出於他的心意.
是很sacred的東西.
然後第三節.
又好比黑門山的金路江在石安山.
即是耶路撒冷.
因為在那裡有耶和華所命定的福就是永遠的生命.
這就是第二個imagery.
第一個imagery就是亞倫.
sacredness 神聖.
第二個imagery你大概get到.
就是應該是一種很life giving的東西.
黑門山是以色列巴基斯坦地最北的一個最高的山.
九千多呎.
是一個雪山.
snow-capped mountain.
是長年積雪的.
當然現在因為氣候亂化.
所以現在已經光了頭.
但它是一個白頭山可以叫做.
即是以色列巴基斯坦地.
它下季的時候是完全不下雨的.
土地是會乾旱.
而那些農作物,人.
怎樣可以在那幾個月.
大概五月至八月底,九月的時候.
這個最乾旱的一個夏季.
能夠熬得過.
就是靠著金路.

$^{681}$即是那些剛剛早上起床.
你見到有一層薄霧那些.
就是靠著那些東西去滋潤著土地和農作物.
滋潤著我們的生命.
這裡說到黑門的金路.
黑門降下來的金路.
一個高山所降下來的清晨的金路.
使到地上能夠有滋潤.
能夠有life.
而那個life是一個永遠的life.
是永遠的生命.
降到哪裡?.
就是以色列人神的子民聚集的那個地方.
就是石安山.
當他們聚集的時候.
他們就去taste到.
不能夠完完全全100\%去體會.
當然這個是到將來新天新地.
我們才能夠體會到.
永遠的生命是一回事.
但是當我們聚集的時候.
弟兄姊妹一齊去朝聖.
去敬拜上帝的時候.
你就發現.
或者吃了頭台.
你get到那個first taste.
What is eternal life?.
什麼叫做上帝給我們的生命.
在上帝的同在裡面.
這個就是上帝為我們設立的群體.
是教會其中一個景象.
神的子民的一個景象.
當我們走在一起的時候.
有上帝所命定的福.
你能夠get到那個taste.
永遠的生命.
你記得我們在浸禮的時候.
我們會有一個信經.
我們一起讀的.
記得嗎.

$^{721}$那個叫做使徒信經.
有很多年歷史.
我們背誦這個信經的時候.
你記得這個信經.
是有很多不同的articles.
來總結整個基督教最重要的.
最basic的信仰.
在信經裡面提過兩次.
是和教會有關.
很多時候我們理解信仰.
理解宗教.
我們都覺得.
都是我們自己個人的東西.
just me and God.
我和Jesus.
我和主之間的契合.
或者有個相通等等.
這些是沒有錯的.
但絕對不是個人主義.
使徒信經有兩處提到教會群體.
他說什麼呢.
他說I believe the holy catholic church.
and I believe the communion of saints.
什麼意思.
當我們聽到catholic這個字的時候.
我們通常就想起羅馬天主教.
羅馬.
但使徒信經本身的意思不是這樣的.
catholic這個字其實是.
普世的意思.
universal church.
所以很多時候現代的翻譯.
我們就說I believe the holy universal church.
就不用catholic這個字.
catholic這個字就是聖兒宮的意思.
大公的意思.
catholic就是這樣的意思.
public visible.
普世universal的教會.
I believe 我相信.

$^{761}$這是上帝自己恩典的揀選.
不是因為我們這裡.
每個人都是屬靈的偉人.
none of us.
是一種這樣的人.
但他是用耶穌基督的恩典.
我們能夠成為教會.
universal church.
並且這裡說到普大兒宮的教會.
不只是說這裡.
我們眼見到的.
而是歷世歷代包括所有的信徒.
所謂因為恩典而能夠成為教會一部分的.
這一班的信徒.
包括已經回到天家的.
或者在教會歷史裡面.
有很多很多跨越時代地理的.
這些信徒.
都是在這個holy catholic church.
一個大聖兒宮的教會裡面是一部分.
而你呢.
就是這個普世的教會的其中一部分.
他說我信聖徒相通.
什麼意思.
chamillion of saints.
就是說所有的信徒都與基督能夠聯合的.
一起來相連.
就等於我們舉個例子.
我們好像去hiking一樣.
我們大家去行山.
行山有目的地.
通常有目的地.
我們在走的時候.
我們發現走著走著我們碰到不同的人.
碰到一些可能他們都在行山的.
不過他說我的目的地是另外一個山頭.
他們去了別的地方.
與你的目的地不同.
但可能你也碰到其他人.
這麼巧合,你不認識的.

$^{801}$但我們碰到一起上路.
然後問問,說著說著的時候.
原來我們去的目的地是一樣的.
是不是很多相通.
雖然你未必喜歡他跟著你.
但有一樣東西跟你一樣的.
就是他的目的地跟你一樣.
是不是.
你喜歡他也好.
你不喜歡他也好.
Eventually.
最終.
都是在目的地見的.
這個叫做聖徒相通.
chamillion of saints.
最後這篇詩篇給我們一些的形容好嗎.
賜給我們身邊的弟兄姊妹.
這班人原來是同路人.
無論他們的年代不同.
語言不同,文化不同.
甚麼不同都好.
兩個目的地是一樣的.
讓我有同行的人.
我為他禱告.
我去建立他.
我們彼此建立彼此服侍.
因為同路人應該做這件事.
第二樣東西.
就是那班路人裡面的其中一部分.
我是其中一分子.
上帝原來不只是愛我一個.
上帝不只是拯救我一個.
他將我用他的恩典.
放在他自己的聖文當中.
不是因為我好.
而是基督給我恩典.
我能夠成為教會的一部分.
第三樣東西.
我值得去投資我的時間.
去學習別人的故事.

$^{841}$你上路上得過癮的話.
你就需要大家聊聊.
你的經歷,你的故事.
這裡不只是記名字那麼簡單.
聽聽大家的故事.
認識一下大家.
第四.
值得用我的心力去投資在關係裡面.
肅靈的關係.
教會,弟兄姊妹的關係.
大家同台吃飯.
大家花點時間.
去克服我們剛認識的時候.
有一點點難過.
沒什麼契合.
但慢慢日子久了.
大家磨合的時候.
就能夠經歷到那個所謂.
「聖徒相通」.
同心去建立教會.
這是何等善,何等美.
我們同心協力去祈禱.
我們能夠成為上帝的子民的一部分.
這個源源全全.
百分之百是你給我們的好處.
主,這個何等美,何等善.
我們能夠走在一起.
你給我們一些同路人.
求你讓我們的生命能夠去投資.
我們的心,我們的力.
使教會能夠成為一個.
更加土地喜悅的群體.
讓我們教會能夠成為一個更有愛.
更有你恩典的一個群體.
我們這樣祈禱,奉主命救.
願上帝賜福給教會.
\newpage



\section{詩篇 8:1-9}
\label{sec:LP_MCyZeclI}
\textbf{人算什麼?(詩篇 8\_1-9) - 陳崇基牧師}
\newline
\newline
連結: \href{https://youtube.com/watch?v=LP_MCyZeclI}{\texttt{https://youtube.com/watch?v=LP\_MCyZeclI}} ~~~~ 語音日期: 2026-01-14
\newline
\newline
\hyperref[sec:4c2vxrJygVs]{\small{< < < PREV SERMON < < <}}
~
\hyperref[sec:index_chronic]{\small{[返順時目]}}
~
\hyperref[sec:index_scriptual]{\small{[返順卷目]}}
~
\hyperref[sec:n8Rs4GhmcQo]{\small{> > > NEXT SERMON > > >}}
\newline
\newline
詩篇 8:1-9
\newline
\begin{longtable}{cl}
\hline
\hline
章節 & 經文 (和合本修訂版)\\
\hline
8:1 & \begin{tabularx}{0.7\textwidth}{X} 耶和華－我們的主啊，你的名在全地何其美！你將你的榮耀彰顯於天。 \end{tabularx} \\ \\ \relax
8:2 & \begin{tabularx}{0.7\textwidth}{X} 你因敵人的緣故，從孩童和吃奶的口中建立了能力，使仇敵和報仇的閉口無言。 \end{tabularx} \\ \\ \relax
8:3 & \begin{tabularx}{0.7\textwidth}{X} 我觀看你手指所造的天，並你所陳設的月亮星宿。 \end{tabularx} \\ \\ \relax
8:4 & \begin{tabularx}{0.7\textwidth}{X} 人算甚麼，你竟顧念他！世人算甚麼，你竟眷顧他！ \end{tabularx} \\ \\ \relax
8:5 & \begin{tabularx}{0.7\textwidth}{X} 你使他比神微小一點，賜他榮耀尊貴為冠冕。 \end{tabularx} \\ \\ \relax
8:6 & \begin{tabularx}{0.7\textwidth}{X} 你派他管理你手所造的，使萬物， \end{tabularx} \\ \\ \relax
8:7 & \begin{tabularx}{0.7\textwidth}{X} 就是一切的牛羊、田野的牲畜、空中的鳥、海裡的魚，凡游在水裡的， \end{tabularx} \\ \\ \relax
8:8 & \begin{tabularx}{0.7\textwidth}{X} 都服在他的腳下。 \end{tabularx} \\ \\ \relax
8:9 & \begin{tabularx}{0.7\textwidth}{X} 耶和華－我們的主啊，你的名在全地何其美！ \end{tabularx} \\ \\
[1ex]
\hline
\hline
\end{longtable}
$^{1}$是 Mark Sir 現在將時間交給他.
各位好 新年快樂.
好像萬歲更生一樣.
天氣都好一點.
我們想用這段聖經開始.
我們今天是2026年第一個主日.
一篇很熟悉的詩篇.
詩篇第八篇.
耶和華我們的主.
你的名在全地何其美.
你將你的榮耀彰顯於天.
你因敵人的緣故從孩童吃奶的口中建立了能力.
使仇敵報仇的閉口無言.
並你所陳設的月亮星宿.
人算甚麼 你竟眷顧他.
世人算甚麼 你竟眷顧他.
你視他比神微小一點.
視他榮耀尊貴為冠冕.
你派他管理你手所做的.
使萬物一切的牛羊.
田野的牲畜.
空中的鳥海裡的魚.
飯油在水裡的都服在他的腳下.
耶和華我們的主.
你的名在全地何其美.
這篇詩篇很熟悉.
這次開年的時候不難倒大家.
揀一首很熟悉的詩篇.
詩篇裡面載著很多當時神的子民.
私人的禱告.
所以他們的禱告 他們的禱詞.
或者他們作了一些詩歌.
去表達他們對上帝的讚美.
或者不同的呼求等等.
這樣就記載在這本所謂詩集.
詩篇裡面.
如果你找過聖經的話.
你知道詩篇大概在哪裡.
在聖經中間.
有150首的詩歌.

$^{41}$是以往歷史歷代神的子民.
他們有的禱告.
然後集合在一起.
放在聖經的中間.
所以這150首的詩篇.
是一些很真誠的說話.
就是詩人感覺到上帝的同在.
感覺到在此時此刻.
他需要怎樣去表達的時候.
他就將它寫成一個詩歌.
有一點像今天的作曲 作詞等等.
作者就將他所經歷的.
放在文字裡面.
所以有人說詩篇像什麼呢.
詩篇好像是小小的聖經.
一本很小的聖經.
聖經很大本 是不是.
聖經很厚 是不是.
你滾手機滾很久都還沒滾完.
但是詩篇150首的詩歌裡面.
就給我們一個所謂濃縮版.
就是那些金寶雞湯.
很濃縮 很 condensed.
你喝完之後.
你大概知道.
到底你和上帝之間的關係應該怎樣.
詩篇裡面的詩歌有很多種類.
有些是讚美詩.
有些是哀歌.
有些是紀念的詩歌.
等等等等.
這一首的詩歌是一首讚美的詩歌.
耶和華說我們的主.
你的名在全地何其美.
你將你的榮耀彰顯於天.
你看到詩人.
當是大衛詩篇 是不是.
他在這裡看到的意境是什麼.
天上有很多星星.
很多的天象等等.

$^{81}$他很讚歎.
很讚歎 很美麗.
然後他就說.
領略到一件事.
就是這位創造主是值得我們敬拜的.
頭一句他說.
我們的主.
你的名在全地.
全地.
Every single square feet.
全地是何其美.
最後那句也是說.
你的名在全地是何其美.
這是什麼心情.
這是一種讚歎和滿有敬畏的心情.
你看到天地萬象的時候.
你不能夠不這樣.
是不是.
宇宙多偉大 是不是.
我們現在還在探索.
這個真的很奇妙 是不是.
2026年.
你知道2025年所謂太空科學.
Astrophysics的新發現是什麼.
有好幾個新發現.
有新的Galaxy.
有新的星系等等.
其中一個發現就是說.
土星.
聽過土星沒有.
有沒有去過.
土星.
原來我們去年.
那些望遠鏡看到.
土星的月亮就多了.
發現多了128個.
所以總共土星7.
就有274個月亮.
如果你是土星人的話.
你想越機煩.

$^{121}$問題不是想不想越.
是想哪個月先.
200多個.
去年我們才開始去到.
再observe到更多更多的東西.
這個起不是私人所讚歎的東西.
你在全地整個宇宙是何其美.
有多少東西去探索.
和有多少東西.
給我們有All and wonder.
叫我們去讚歎上帝的權能.
第二節他說.
因受敵的緣故.
在BB的口中.
建立了能力或者榮耀.
使那些報仇的.
或者以敵黨上帝的閉口無言.
就是上帝去到申冤的時候.
上帝去到行公義的時候.
這些仇敵就閉口無言.
BB說到在古代的時候.
就是特別一些脆弱的生命.
今時今日和以前都是一樣.
是一些很脆弱的東西.
是無助的東西.
我預備的時候.
這篇詩篇雖然是很熟悉.
我記起有一次去.
芝加哥大學的博物館.
近東博物館.
是一個很好的博物館.
是世界級的博物館.
我記得有一次去的時候.
就是跟一個教授去.
就去近東博物館的後面.
如果去近東博物館.
你外面有很多東西看.
當然.
但裡面原來還有很多東西.
是收藏在一起的.

$^{161}$好像Home Depot.
收藏在一起的.
在一些架子上.
經過一個地方.
他就說.
這裡有一個很小的盒子.
是一個很小的木乃伊.
大概死者發現的.
只不過是十歲左右.
他不知道是因為什麼緣故.
所以他去世.
所以就把它弄成了木乃伊.
主管就說.
這些東西我們不能放出去.
因為不是很敏感.
放個小孩在這裡給人看.
大人就可以.
就像你這樣.
但如果小孩放出去.
就不太好看.
所以他就說.
這個我們只能放在貨倉.
一個十歲左右的小孩.
在近東的古代的世界.
因為沒有防疫.
醫藥不夠.
人很快就會過去.
但無論如何.
這裡說的是孩童.
就是一些沒有能力的人.
沒有能力的人.
上帝是怎麼對待他們的呢.
就是他去保護他們.
還有將能力建立在他們身上.
這裡說的是上帝的本性.
是一個保護卑微的人.
還有投靠他的人.
他會為他伸冤的.
這個是上帝由始至終.
整本聖經裡面所啟示的本性.

$^{201}$很多時候我們說到.
一些權力的時候.
說到偉大的時候.
我們都會想起一些大的國.
帝國.
我們的國家多有能力.
我們出兵的時候.
沒有人敢出聲.
我們想起權能power就是這個.
但詩篇裡面說.
整本聖經裡面所說.
上帝的權能往往相反.
他是在一些無能者的身上.
去建立能力.
去保護他們.
去捍衛他們.
詩篇有150篇我說過.
這篇第幾篇.
第八篇.
詩篇一至到150篇其實有一個編排.
大致上我們知道有一個所謂arrangement.
無論哪一個去安排都好.
哪一個最後去做編輯.
去將150首.
去到這樣放都好.
到底是哪一個都好.
或者那班人是哪一個都好.
大概應該是猶太人.
大概兩樣之間.
他們將150首詩歌.
齊集了然後這樣去排序.
大概我們知道.
那個排序其實是有心思的.
很多時候如果你留意的話.
那個心思在哪裡呢.
第八篇.
詩篇有150篇.
頭兩篇就是introduction.
所以兩篇就好像那些書.
preface.

$^{241}$即是開頭的前言等等.
到第三篇至第七篇.
即是第八篇.
三至七篇.
就是大衛作王的時候受的苦難.
即是被人追殺的時候.
他求助無門.
在苦難裡面他去呼求上帝的時候.
他去等候上帝的國.
因為他不是用political power.
拿他的王位.
他是等候上帝的旨意.
耐心等候上帝為他伸冤.
他不去報仇.
他可以去報仇.
他可以殺掃.
但在患難裡面.
他去專心等候上帝的國度來到.
一個謙卑的王去等候上帝的來到.
三至七篇.
然後九至十四篇.
九至十四篇.
這是第八篇.
第九至十四篇.
就是一些受苦的哀歌.
即是詩人.
有時候在大衛.
在軟弱無能的時候.
他信靠上帝.
知道上帝是揀選他的.
他謙卑的.
不是驕傲的一個人.
不是一個狂妄的人.
他是等候上帝的國度來到.
你看到三至七和八至十四篇.
就夾在中間這裡第八篇.
因仇敵的緣故.
耶和華你的能力.
是可以的.
可以去信靠這位的主.

$^{281}$第八篇的位置就在中間.
掌管全地的上帝.
他才是有power.
並且他power的本質和質地.
是不同我們今天眼見的權能.
不是飛機大爆.
而是上帝願意在謙卑人身上.
去幫助他們.
是寫給我們的.
我們每一個脆弱有軟弱的人.
上帝的能力怎樣能夠彰顯.
就是我們不作為一個狂妄的人.
我們謙卑在上帝面前.
我們承認我們需要上帝.
我們生命是很脆弱的.
你和我今天來到新的一年.
我們都老了.
不知道你生日了沒有.
未生日你肯定會生日.
你和我都會老.
包括我自己.
新的一年我自己.
都有些新的志向.
有些立志.
我自己因為兩個月前.
我腰痛得很厲害.
我腰痠.
我經常都有這些問題.
不過痛到一個地步我不能忍耐.
當然有些弟兄姊妹幫了我很多.
我決定要參加LA Fitness.
這是人生中未試過的.
人生中未試過.
我經常覺得為何要做這些事.
我自己在家做平板支撐就行了.
結果我自己都.
走進去LA Fitness.
我自己的心理.
這是我的問題.
我的心理可能有問題.

$^{321}$我掙扎了到底是否真的要參加.
不便宜的.
而且通常參加了之後就不去.
我很懷疑自己的能力.
到底我是否真的會.
所以我決定不如參加.
參加了.
立定心志去LA Fitness.
自己做一些適合我的運動.
我不是18,22歲.
雖然我看到18,22歲的身材.
我就很羨慕.
但我自己都堅持.
因為我是一個很脆弱的人.
我堅持了兩個多月.
不錯吧.
但股票完獎可能明天我不去.
我每天早上都去游泳和做運動.
因為了解到我會衰殘.
上帝才是我的力量.
這就是詩篇所說.
上帝的能力在哪裡彰顯出來.
第三節.
並你所陳設的月亮星辰.
人算甚麼你竟眷顧他.
世人算甚麼你竟眷顧他.
你視他比神微小一點.
視他榮耀尊貴為冠冕.
你派他管理你手所做的萬物.
一切的牛羊,牲畜,海裡的魚.
空中的飛鳥等等.
都服在他的腳下.
這裡說到人的存在.
人的存在是一回甚麼事.
詩篇第八篇其實不只是讚美詩歌.
其實是一個神學的宣告.
宣告甚麼呢.
就是宣告人本性是怎樣.
我們的存在是怎樣.
這裡很有趣.

$^{361}$他說到人的存在.
即是你和我的存在.
是一種不能解除的張力.
有一個tension在我們裡面.
一方面我們覺得自己是很微小.
我們有如微塵一樣.
這裡說到人的時候.
就用了一個字.
在希伯來文是enose.
你不用懂得希伯來文.
你都知道甚麼意思.
enose 現在我解釋給你聽.
enose 就是死硬的人.
mortals.
必死的人.
我們是一班必死的人.
人算甚麼.
然後第二個字就是世人.
son of adam.
或者可以叫做人子.
即是耶穌經常都叫自己人子.
即我們只是一個人.
卑微的人.
代表我們是來自塵土.
希伯來文有些東西很有趣.
adam 人這個字.
和土地差不多音.
因為土地就是adam ma.
你加個ma字在後面.
你好ma的時候.
加個ma字的時候就變成另外一個字.
所以adam和adam ma就很相似.
其實就是在他的字眼裡面.
他想表達一樣東西.
就是一種概念.
我們和土地塵土其實是很有相關.
因為我們是同源的.
我們就是來自塵土.
我們就是塵土就是卑微.
這個一方面有張力.

$^{401}$我們很微小.
好像塵土一樣.
但另一方面詩篇講到人類的時候.
人的存在的時候.
他就講到你使他.
比甚麼.
比神微少一點.
或者舊有的翻譯就是天使.
無論是天使也好.
神也好.
這裡講到Elhim.
上帝.
比上帝微少一點.
我們雖則是塵土.
雖則是微小.
但是上帝創造我們的時候.
給我們的尊貴.
給我們的地位在宇宙裡面.
是你敢認第二.
沒有人敢認第一的.
比上帝微少一點.
上帝是冠軍.
我們就是亞軍.
我們就是一個這麼尊貴的人.
所以很多人會說.
詩篇第八篇是講人的弔詭性.
人類的弔詭性.
就是我們觀望宇宙浩瀚的偉大的時候.
我們發現.
不單止權能的彰顯.
也發現我們自己的存有.
是一個很微小的事.
真的沒有人會在乎你.
沒有人會在乎你.
在16世紀的時候.
有一個神學家.
當時他寫的東西.
現在還有人看.
有人研究.
那個人叫做約翰加爾文.

$^{441}$約翰加爾文有一段說話很有趣.
他說認識自己對自己的認識.
和對上帝的認識.
其實兩樣東西是一個循環.
因為一個人要認識上帝的時候.
就不可能不認識自己.
如果一個人想認識自己.
If you want to know thyself.
You have to know God.
因為你不認識上帝的時候.
你就不會認識你自己.
你認識你自己的時候.
你必須要認識上帝.
他就說到self knowledge和divine knowledge.
兩樣東西是一個循環.
你認識上帝越深的時候.
你就應該更深認識你自己.
認識你自己更多的時候.
你就轉過來.
你就會認識上帝是怎樣.
這個就是詩篇第八篇所說.
如果你覺得你自己.
或者你錯誤地認識你自己.
以為你自己是很強.
或者很能夠去控制人生裡面所有的東西.
那你就其實不認識上帝.
因為你不知道你自己是多麼渺小.
但當你誠實面對你自己的時候.
照一照你的鏡子.
照清楚你自己.
你就發現你自己原來有很多缺陷.
原來你有很多地方是不行的.
你是很脆弱的.
這個時候你就要認識上帝.
因為上帝是那位醫治我們的神.
是創造我們的神.
是帶領我們的上帝.
這個就是我們的paradox.
如果你不認識這個paradox的時候.
你就不認識你自己.

$^{481}$這篇詩篇我們讀了不久.
多少字?.
忘記了.
九節.
你覺得像《創世記》哪一章?.
詩篇說你再傳遞何其美.
你將一切的東西伏在他的腳下等等.
哪一章?.
第一章是嗎?.
你見到第一章《創世記》的時候.
他記載上帝創造天地.
其實很工整的.
六天去創造世界.
你找到人是何時創造?.
這個當然不是科學的陳述.
這是神學的陳述.
是theological statement.
你見到一和四,二和五,三和六.
這個以前說過了是對比.
人是何時出現?.
是the climax of creation.
the crown of creation.
是最後最後出來的.
好像你聽交響樂.
最後最後那個最高峰的時候.
終於所有樂器都出來了.
最宏大的時候.
最偉大的時候就出現了.
這個就是第六天.
上帝創造創立天地.
然後最後去安息或作王.
是第六天才會有人類出現.
但同時你要留意.
第六天有人類出現.
人類多偉大也好.
他和甚麼一起出現?.
地上的動物.
所以你和他都是同一天.
你就不要以為我是另一隻不同的.
雖然我們會說人類是萬物之靈.

$^{521}$但你和動物,貓,狗.
都是同一天發生的.
同一天被創立的.
這是一個theological statement.
告訴我們我們和地上的動物.
其實都有很多相似的地方.
所以現在我理解了.
為甚麼人們那麼疼愛貓狗.
為甚麼貓狗那麼疼愛人.
因為我們都是同類.
這就是paradox.
同時你是一個很glorious的being.
但同時你也是一個很humble的being.
《詩篇》第八篇所說.
Really nobody should care about you.
但這裡說甚麼?.
上帝去眷顧你.
2025年.
不妨想一想.
自己心裡寫一寫日記.
無論你真的寫出來也好.
或是心裡有一個反應,反省.
2025年上帝怎樣看顧你.
大大小小我們都應該經歷過.
去年又說回太空.
這篇詩篇說太空.
上一年太空新聞又發生甚麼事.
你記得有兩個太空人.
前年上了太空站.
Starliner.
去做一些工作等等.
本身打算去多久?.
很短時間.
誰知去了多久?.
去了差不多九個月.
因為有些事情不可行.
有些事情不可運作等等.
超乎他們那幾天的任務.
於是他們滯留在太空.
到去年三月才安全地回到地球.

$^{561}$有多慘?兩個人.
本身打算去幾天.
好像你打算去多倫多.
你去幾天而已.
誰知滯留了九個月.
好不好?.
起碼有東西吃.
多倫多好一點.
如果不然生卡拉的地方.
太空就是這樣.
吃甚麼九個月.
很慘.
和那個人合不合?.
很多時候不合.
這樣就真的要命.
最慘是不知何時回來.
能否回來又是另一件事.
所以太空總署用了很多資本.
去救援他們兩個人.
花了大概多少錢?.
或者算算大約花了一億.
去救這兩條命回來.
最先進科技都會失控.
人有多偉大都是脆弱.
詩篇第八篇所講的.
就是這個道理.
Nobody should have care about you.
但人算甚麼?.
你竟顧念他.
世人算甚麼?.
你竟眷顧他.
顧念當然是care.
但這個眷顧這個字很有趣.
有時候可以翻譯做visitation.
探訪.
到訪的意思.
但其實有個意思.
這個字的意思其實是點名.
點名.
即是老師在課室點名.

$^{601}$大家有沒有試過?.
小學中學的時候.
點你名我很清楚.
因為我姓陳.
經常都是我先走.
因為是C.
經常都是1號2號3號都是我的.
老師叫你名的時候.
你就點名.
舉手到了.
其實就是點名的意思.
即是上帝點你名.
點你名不是要罰你.
點你名是知道你存在.
I am here.
I know you are here.
You matter.
你是有價值的.
我們基督裡面起碼有個價值的宣言.
在耶穌基督裡面我們有個救恩.
是一個無價的救恩.
上帝點你名.
知道你在這裡.
祂說我愛你.
I love you.
I care about you.
有時候我們會覺得自己微不足道.
但詩篇向我們宣告.
你是極重榮耀極重價值的一個存在.
上帝見到你.
去年見到你.
今年也是見到你.
你是配得上帝去點名的人.
你把你自己塞進去.
把你的名字塞進去.
我算什麼.
你已經眷顧我.
我算什麼.
你已經點過名.
我們多脆弱也好.

$^{641}$我們多不事也好.
多麼破碎也好.
上帝的恩典說.
I call on your name.
I know you.
在新紀念裡面.
我希望不是定普通的立字.
二三四.
其實我們真的需要認識我們自己.
認識我們自己是一個怎樣的人.
想深一點認識多一點.
我是一個怎樣的人.
我的personality 我的character.
我的喜好等等.
以至我們能夠認識上帝.
認識自我就是什麼.
就是認識上帝.
而你認識上帝.
就不能夠不認識你自己.
知道我們自己的paradox.
但也知道上帝眷顧我們.
看顧我們.
You matter.
I matter.
因為上帝給我們恩典.
求主幫助我們.
我們今天開始的時候.
開始今年.
我們會守聖餐.
聖餐說什麼.
You matter.
耶穌和你和我說.
你接受我這份救恩的禮物的話.
You matter most in my heart.
你面對今年.
你知道上帝每天早上都點你名.
知道你是誰.
認識你.
求主幫助我們認識自己.
認識一些位外的我們上帝.

$^{681}$我們同心協力祈禱.
我們預備我們的心.
預備接受這個聖餐.
上主我們多謝你.
因為你點我們的名字.
原來我們在浩瀚的宇宙當中.
我們雖然是微塵.
是塵土.
是卑微脆弱.
敗壞的.
但你卻眷顧我們.
We matter to you.
特別當我們領這個聖餐的時候.
你頂著這個杯子.
告訴我們.
你給我們這個恩典.
給我們這個眷顧.
是不停的.
是永遠的.
為我們這群脆弱.
這麼不可愛的一群人.
你願意去寫名.
去救贖,拯救我們.
願你的名在全地.
在我們的生命中.
顯出你是何其名.
我們這樣祈禱.
奉主耶穌的名字.
我們開始領聖餐的時候.
如果你是.
\newpage



\section{約翰福音 1:1-5}
\label{sec:n8Rs4GhmcQo}
\textbf{現身的上帝(約翰福音1\_1-5) - 梁紫健傳道}
\newline
\newline
連結: \href{https://youtube.com/watch?v=n8Rs4GhmcQo}{\texttt{https://youtube.com/watch?v=n8Rs4GhmcQo}} ~~~~ 語音日期: 2026-01-24
\newline
\newline
\hyperref[sec:LP_MCyZeclI]{\small{< < < PREV SERMON < < <}}
~
\hyperref[sec:index_chronic]{\small{[返順時目]}}
~
\hyperref[sec:index_scriptual]{\small{[返順卷目]}}
~
\hyperref[sec:_Xk1Tbpr5MU]{\small{> > > NEXT SERMON > > >}}
\newline
\newline
約翰福音 1:1-5
\newline
\begin{longtable}{cl}
\hline
\hline
章節 & 經文 (和合本修訂版)\\
\hline
1:1 & \begin{tabularx}{0.7\textwidth}{X} 太初有道，道與神同在，道就是神。 \end{tabularx} \\ \\ \relax
1:2 & \begin{tabularx}{0.7\textwidth}{X} 這道太初與神同在。 \end{tabularx} \\ \\ \relax
1:3 & \begin{tabularx}{0.7\textwidth}{X} 萬物都是藉著他造的，沒有一樣不是藉著他造的。凡被造的， \end{tabularx} \\ \\ \relax
1:4 & \begin{tabularx}{0.7\textwidth}{X} 在他裡面有生命，這生命就是人的光。 \end{tabularx} \\ \\ \relax
1:5 & \begin{tabularx}{0.7\textwidth}{X} 光照在黑暗裡，黑暗卻沒有勝過光。 \end{tabularx} \\ \\
[1ex]
\hline
\hline
\end{longtable}
$^{1}$弟兄姊妹,午安.不知道大家有沒有看過這棟建築?.
Good morning everyone. Have you ever seen this building in this picture?.
有看過嗎?.
No..
大家覺得這一座是什麼樣的宗教的建築呢?.
What religion do you think it belongs to?.
有人猜嗎? 哦,道教.有可能.還有嗎? 佛教,OK..
Mongolia..
有人說是道教,有人說是佛教..
其實它是一座基督教的建築..
So maybe we think it belongs to Buddhist or Taoist, but actually it's a Christian building..
它的名字叫道風山,是香港沙田的一座建築,是由一個基督教信義中的宣教士建立的.他把教會聚會的地方刻意以中國佛教寺院的形式去建立,希望基督教以中國化的形式去傳福音..
在1930年代規劃,目的是向剛剛說佛教,道教的教徒去傳福音,讓他們能夠在熟悉的宗教環境裡面接觸到宗教信仰..
So the name of the building is called Tao Feng Shan, a building in Hong Kong, it was designed in the 1930s by Norwegian Lutheran missionary who deliberately built a church meeting place in the style of a Chinese Buddhist temple, hoping to spread the gospel of Christianity in a Chinese way..
The aim was to evangelize Buddhists and Taoists, allow them to encounter the Christian faith in a familiar religious space..
So in order to package the gospel in a way that was familiar to the Chinese people, such as Buddhists and Taoists, the missionaries of that time designed this building in this form..
God knew that we humans didn't know or understand Him, so He was originally invisible..
In the Old Testament, people couldn't contact Him and could only know God's will through Moses and prophets..
所以在新約的年代,上帝因為要讓我們以一種比較熟悉的方式去認識他,他將他自己道成肉身,來到這個人間,成了一個現身的神..
In the New Testament era, God wanted us to know Him in a more familiar way, so He incarnated Himself and came to earth, becoming our living God..
Today we will examine how these ones invisible in terms of God through the incarnation of Jesus appeared in a way that we can understand..
我們先來讀一下今天的經文,約翰福音1:1-5.
太初有道,道與上帝同在,這道就是上帝,這道太初與上帝同在,萬物都是基著他所造的,沒有一樣不是基著他所造的..
凡被造的,在他裡面就有生命,這生命就是人的光,光照在黑暗裡,黑暗卻沒有勝過光..
Let's read today's scripture first..
In the beginning was the world, and the world was with God, and the world was God. He was in the beginning with God..
All things were made through Him, and without Him was not anything made that was made..
In Him was life, and the life was the light of men. The light shines in the darkness, and the darkness has not overcome it..
通常我們讀到這個經文也是很繞口,對吧?.
這個上帝創造這個世界,人們雖然在一開始看不見這個上帝,但看得見上帝用他的話語所創造的世界..
Yeah, the Bible is hard to read, especially for me, for the first time to translate the English. Thank you..
This God created the world, although people couldn't see this God and this beginning, they could see the world created by God through His word..
所以這裡說,太初有道,道與神同在,萬物都是基著他所造的..
就是上帝一開始創造這個世界的時候,他用他的話語,他說要有光就有光..
他說要有天就有天,他說要有地就有地..
所以起初上帝的說就是他的話語,很明顯所有的事情都是基著他的話語所造的..
Therefore it is said here, in the beginning was the world, and the world was with God, and through Him all things were made..
That is when God first created the world, He used His word. He said let there be light and there was light..
He said let there be heaven and there was heaven. He said let there be earth and there was earth..
In the beginning God's speaking was His word. Clearly all things were created through His word..

$^{41}$道與神同在,表明了上帝擁有了創造一切的基石..
他與人類根本的距離是無限..
他根本不需要派他的獨生兒子來到這個世界上,人類只是他區區的受造物..
The phrase "the world is with God" indicates that God processes the foundation for creation,.
and that the fundamental distance between Him and humanity is infinite..
He didn't need to send His only begotten son into the world, humanity is merely His creation..
就例如你喜歡玩樂高,你親手建造了一個樂園,裡面有很多不同的建築物和人..
如果你有能力的話,你會想把自己降級為樂高裡面其中的一員嗎?.
就例如你喜歡玩樂高,你親手建造了一個樂園,裡面有很多不同的建築物和人..
如果你有能力的話,你會想把自己降級為樂高裡面其中的一員嗎?.
或許我們當中會有些人覺得沒什麼大不了,最多是成為一枚樂高而已..
也體會成為樂高的限制和痛苦..
永遠都不能控制自己的身體,你要靠你的主人幫你去調整自己的姿勢..
也體會成為樂高的限制和痛苦..
永遠都不能控制自己的身體,你要靠你的主人幫你去調整自己的姿勢..
所以我們製造樂高的時候會花很多時間和功夫,.
但上帝的話語是在彈指之間的時間,也不用動手,.
他一說話就能帶來創造..
他的話語所創造的能力在世界上沒有一樣東西能與他相比..
就算有這樣的距離,即使是這樣子,上帝仍然將自己的話語耶穌基督他自己降生為人,.
就比起我們成為樂高要降級更多,所受的限制也比我們來得更多..
但耶穌為了成為這個世界的光,人們的盼望決心要這樣子做..
而他的生命就是人的光,這裡用了光和黑暗的比喻,.
這個比喻讓我們理解上帝怎麼樣讓我們知道,感受到他的存在..
相信大家都知道,如果這個世界沒有光,全是黑暗,我們無法看見任何東西..
而大家都知道黑暗是一個狀態,不是一個真正存在的東西..
那個角落之所以黑暗是因為什麼?是因為那個地方沒有光..
而不是真的黑暗存在..
上帝就用了這個我們比較熟悉的方式,光與黑暗的比喻..
從起初的時候,上帝就把光帶到這個世界,讓世界能夠分辨哪裡有生命,哪裡光沒有照到,只剩下黑暗..
就像我們還沒有信主的時候,我們可能還沒有察覺這個上帝是存在的..
我們從前就像被黑暗矇蔽一樣,不知道上帝是那個創造的主,.
也不知道原來祂在我們的身邊,一直保護著我們..
直到我們開始瞭解上帝,原來在我們生命當中做了很多很多美好的工作,.
我們才知道,噢,原來這個世界有溫暖,有生命,有光..
所以為什麼上帝會把光帶到這個世界?.
因為上帝會把光帶到這個世界,而不是把光帶到黑暗..
約翰這裡要強調,這生命就是人的光,光照在黑暗裡,黑暗卻沒有勝過光..
就是因為上帝把光照進這個世界裡面,同樣把光照進我們的生命,.
讓我們看見這個世界與以往有很大很大的不同,讓我們生命看見盼望..

$^{81}$所以如果上帝不把光帶到這個世界,.
我們就永遠等於沒有光,沒有生命,沒有生命..
所以我們就永遠等於沒有光,沒有生命..
所以如果上帝不把光帶到這個世界,.
我們就永遠等於住在黑暗裡面,看不到任何的事物..
所以如果上帝不把光帶到這個世界裡,.
我們就永遠等於住在黑暗裡,看不到任何的事物..
所以這個上帝在本質上,雖然他與我們的距離很遠,.
而我們也是按照他的形象所造,.
但是我們在能力上完全無法與他相比..
但是,然而他卻要成為我們生命的光,.
主動把這個距離收窄,道成肉身來到我們身邊,.
為我們帶來盼望..
所以我們有一些應用..
在我們生活當中,我們需要多一點刻意去發現這個上帝的作為..
例如,我們去享受他所創造的美好,.
他的世界,他的光,發現上帝在我們身上的作為..
在新的一年,我們可以立志,養成一個新的習慣,.
就是每個星期我們都刻意去找,一樣,上帝在我們身上的作為..
或者我們會發現他給我們的祝福..
有時候我們信主了一段時間,我們可能會忽略了這些看似基本的動作,.
以為我們信了主很久,就不需要刻意去發現上帝在我們身上的工作..
第二,我們需要更多的信念,.
以便能夠更多的信念,.
第二,是我們讓我們身邊還沒有感受到上帝的家人朋友,.
也開始感受一下祂的存在..
Let's help our family and friends who haven't yet experienced God's presence to begin feeling it..
留心我們自己對於上帝存在的見證,.
分享我們生命當中上帝的作為..
當我們持續地去留意上帝在我們生活的證據時,.
我們就有見證的存獲,.
我們就可以隨時讓我們身邊的親朋好友去感受一下上帝的確在我們身邊..
Pay attention to our own testimony about God's existence and share God's work in our lives..
When we continuously observe the evidence of God in our lives,.
we have a stock of testimony..
We can then let our family and friends feel at any time that God is indeed with us..
Recently, a Christian friend who is married has been having a dinner party at home..
He invited a woman who is about to get married,.
she is not yet a Christian,.
and he asked the Christian here,.

$^{121}$"How can you be sure that she is your life partner?".
One of the sisters shared,.
"A life partner is to understand each other,.
and know each other's survival mechanism from their families of origin.".
但她說,最重要最重要的是,.
大家在上帝的關照之下纔能有悔改跟改變的可能,.
靠自己完全不能..
那位女生就即使感受到上帝的臨近,.
覺得這一場聚會很有意義..
我想這就是我們持續的留意上帝在我們生活的證據時能夠分享的一個例子..
第三,最後一點就是,.
認識我們內心的黑暗,讓上帝的光和生命照進來..
有時候我們的內心會被黑暗所矇蔽,.
我們沒有感受到上帝的存在跟祂的指引,.
讓我們在平日花一點時間讀上帝的話語,.
讓祂的光照亮我們的內心..
讓我們一起來禱告吧..
主啊,求禰幫助我們,.
每個人都留意上帝的作為,.
讓我們能夠與家人朋友分享,.
讓我們內心受到光照,.
求主幫助我們..
讓我們一起來禱告吧..
主啊,我們禱告,.
每個人都能夠留意上帝的作為,.
我們能夠與家人朋友分享,.
我們的心也能夠被限制..
主啊,幫助我們吧..
我們這樣子禱告是奉主耶穌的命求..
阿們..
(喫飽了).
\newpage



\section{3:1-15}
\label{sec:_Xk1Tbpr5MU}
\textbf{難先要你做 (傳道書3\_1-15) - 譚文倩莊傳道}
\newline
\newline
連結: \href{https://youtube.com/watch?v=_Xk1Tbpr5MU}{\texttt{https://youtube.com/watch?v=\_Xk1Tbpr5MU}} ~~~~ 語音日期: 2026-01-24
\newline
\newline
\hyperref[sec:n8Rs4GhmcQo]{\small{< < < PREV SERMON < < <}}
~
\hyperref[sec:index_chronic]{\small{[返順時目]}}
~
\hyperref[sec:index_scriptual]{\small{[返順卷目]}}
~
\hyperref[sec:Ly_O4TyukTE]{\small{> > > NEXT SERMON > > >}}
\newline
\newline
$^{1}$現在呢?我大聲點叫,好像剛才Chimay說的.
其實我的聲音有點沙,剛才就是在叫.
好冷啊,大家有沒有在外面喝杯熱茶才進來?.
我有啊,我覺得好冷.
我在香港的時候,我就跟兩個阿姨,姨長喝茶.
喝茶的時候,我跟姨長,大姨長和大阿姨先到了.
我大阿姨就在11點之前,要趕快下點心.
知道為什麼嗎?因為有九折.
就趕快點了,小阿姨,姨長就跟著到了.
一到的時候,小阿姨就問,點了點心沒有?.
大阿姨就馬上停了,不要不要不要,不要點了.
我已經點了10樣了.
小阿姨就說,哇,五個人10樣這麼多.
但是不要緊,10就10全10尾.
9就長長久久,8就大發大發.
總之不是7就行了.
如果是7的話,加多一碟豉油都可以.
我坐在這裡聽的時候,我就覺得.
Oh, that is so interesting,這麼有趣.
因為7這個號碼是一個漂亮的號碼.
對於我們基督徒來說,這是一個漂亮的號碼.
在聖經裡面是代表完全,圓滿,完成的狀態.
這個觀念是源自創造的.
在整本聖經裡面是不停地出現的.
也這麼巧,我這次是講傳道書三章.
這段經文裡面正正是有一首詩.
有七對相對的句子.
描述人生裡面不同時間的一些狀態.
我隨便講傳道書一至二章講過些什麼.
傳道者用智慧尋求滿足,記不記得.
他嘗試去找人生意義.
最終發覺在這個日光之下.
所有東西都是很短暫的.
Come and gone.
很尋很尋,直至最後那兩節.
將我們的眼光帶到哪裡.
帶到我們的飯桌上.
而將24至25節說什麼.
人最好是吃喝,在自己的奴立中自得其樂.
Oh by the way.

$^{41}$這也是出於神的手.
不要忘記,離開了他.
誰能有吃的呢?誰能享樂呢?.
來到第三章這裡.
傳道者就在講一些好像既定時間會發生的事情.
而去到三章十四節那裡.
他再一次提起一個我們經常聽到的概念.
驚威神,Fear God.
今天我想藉著傳道書三章一至十五節.
和大家去思考兩個問題.
萬事都有定期,天下萬物都有定時.
如果神已經安排了所有的東西.
所有東西都有神的時間.
我們人其實是不是還有選擇呢?.
如果是這樣的話,我們還這麼努力做什麼呢?.
這段經文再一次提到驚威神這個概念.
其實驚威神是什麼意思呢?.
我們怎樣才叫驚威神呢?.
還有在這段經文裡面.
驚威神和凡事都有定期,天下萬物都有定時.
有什麼關係呢?.
我先和大家讀一讀這段經文.
傳道書三章一至八節.
他說萬事都有定期,天下萬物都有定時.
生有時,死有時,災中有時,不出所中的也有時.
殺戮有時,醫治有時,拆毀有時,建造有時.
哭有時,笑有時,哀動有時,踴躍有時.
拋擲石頭有時,堆砌石頭有時,擁抱有時.
避免擁抱有時,尋找有時,捨棄有時.
保存有時,捨棄有時,私利有時,逢補有時.
靜默有時,講話有時,愛有時,恨有時,戰爭有時,和平有時.
可能你離開之後,腦袋裡不斷轉有時,有時,有時.
第三章一開始,其實他和前面兩章有點不同.
好像有一個新的段落.
這八節裡面有七對的對立面.
講到生死,栽種拔出,哭笑等等.
在這裡,七對對立面是希伯來詩的典型特徵.
對古代希伯來詩人來說,七代表完滿.
在這裡,亦代表了這首詩所講的東西.
已經包含了,圓圓滿滿包含了我們人生裡.

$^{81}$可能會發生,應該會發生的所有事情.
當然,裡面有辣有不辣.
傳道者要提醒,這些所有事裡面.
都有神的perfect timing.
問題是,如果是這樣,我們人是不是還有選擇呢.
有得選擇的,好像我這樣.
我憧憬可以快點供完房子,可以輕鬆點.
花錢不用再有那麼多掣肘.
我有自由去想,亦有自由去供房子.
甚至我有自由去勒緊褲頭.
再給多點錢去principal那裡,讓我快點脫身.
不過我控制不了的,就是如果當我供完房子之後.
我和我老公是不是可以坐在沙發上,按按腳.
無憂無慮地過我們的下半生.
因為健康和生命的長短,不是在我的手裡.
我們繼續看下一段,傳道書三章九至十五節.
作功的人在自己的勞碌上得到什麼益處呢.
我看神給予世人的擔子,是要他們為此煩惱.
他使萬事各按其時成為美好.
他又把永恆的意識放在人的心裡.
雖然這樣,人還是不能察覺.
神自始至終的行為,作為.
我曉得人生最大的事,尋樂享福.
人人有吃有喝,在自己的一切勞碌中只得其樂.
這就是神的恩賜.
我知道神所作的一切,都必永存.
無可增添,無可減少.
神這約了,為要使人在他面前存敬畏的心.
現在有的,先前就有.
將來有的,早已有了.
因為神使已過去的事,重新出現.
若果神已經想好了.
又掌權,又慷慨,又滿有恩典.
我們還有什麼好想呢?.
正如這裡所說.
作功的人在自己的勞碌上,得到什麼益處呢?.
看經文的時候,好像有個答案.
我看神給予世人的擔子,是要他們為此煩惱.
這是什麼答案呢?.
你會否看經文時,想像一下.

$^{121}$這麼搞笑的,神給人煩惱.
這裡不是說我們所信的神是有點變態.
要給我們做些事情,讓我們享受.
讓我們煩惱,折磨我們.
這節的經文正確的解釋是.
神容許我們人的生命裡有重擔.
這些重擔是為了要我們人被磨練,被謙卑.
1991年亞利桑那的科學家建造了一個叫做Biosphere 2的實驗環境.
一個人工控制生態幾乎完美的大自然.
裡面有很乾淨的水,很潔淨的水和空氣.
健康的土壤,經過過濾的陽光.
原意是想試一下人有沒有可能.
為樹木,水果,蔬菜,甚至我們人類.
去提供一個最理想的生長條件.
大約兩年之後,即是1993年.
他們發現了一個很有趣的現象.
明明這個環境幾乎完美的地方.
那裡生長的樹,飄高到某一個高度.
就會陸陸續續地掉下來.
你猜到為什麼嗎?.
原來因為科學家漏了去製造一樣東西.
漏了去製造一樣很關鍵的元素.
就是風.
嚴格來說是Storm.
這個研究發現原來沒有風暴的環境.
原來樹木真的可以生長得很快.
但這些樹的樹幹卻是很瘦弱.
所以到了某個高度的時候.
它們是支撐不了自己而倒塌下來.
原來惡劣的天氣暴風的存在.
會迫使樹木長出一些強壯的根.
去伸展出來抓住泥土.
同時也會生長一些所謂的Stressed Wood.
受壓木,令這些樹木有韌力,有承受力.
所以在這個Biosphere這麼安靜.
這麼舒適,受保護,受控的環境裡生長.
表面好像很漂亮.
但兩年左右,這些樹就承受不了自己的重量.
一棵一棵地倒塌下來.
這個實驗其實很有象徵意義.

$^{161}$原來壓力對健康,強壯的成長是不可以缺少的.
當壓力來到的時候,受壓的確不是很舒服.
但這正是建立力量和適應力的必要條件.
缺乏逆境,反而阻礙了真正的成長和韌力的形成.
聽聽看,似乎這個道理不單止是可以應用在樹木上.
其實我們人也是一樣.
我們不妨去想想.
你生命裡最寶貴的功課,最能夠令你成長的事情.
發覺自己有機會成長的時候是來自哪一道呢?.
是不是你自己讀聖經,聽一篇好的道.
看到好的榜樣.
還是想一想,正正就是被磨練,被謙卑的時候.
我知道神所作的一切都必永存.
無可增添,無可減少.
神這樣作,為要使人在他面前存敬畏的心.
這幾次說傳道書,也有人問我.
其實什麼叫敬畏神?.
或者敬畏神和這裡所說的凡事都有定期,天下萬物都有定時.
有什麼關係呢?.
首先我想跟大家說說敬畏這個原文的意思.
它是在說恐懼,的確是在說恐懼.
是神身上的崇敬,令人害怕的意思.
敬畏神的意思是在說你對神的恐懼,崇敬.
而這個恐懼和崇敬會影響你的行為.
會影響你怎樣去過你的生活.
Fear to the point that it will affect how you behave, how you live.
以前我聽過一個牧師說過一個這樣的比喻.
他說敬畏神就好像一個結婚的男人.
你不會約另一個女人,即不是你老婆.
去跟她在沙灘漫步,然後去吃個燭光晚餐.
為什麼?其實很簡單.
某程度上你怕老婆,是不是?.
說得好聽的就是不是,我很尊敬我老婆.
這種帶有害怕,但又尊敬意味的裡面.
所以這個萬事都有定期,天下萬物都有定時.
和敬畏神之間那個緊扣的關係.
就是在說人承認神是這個主宰.
我們要感恩並且帶著一個敬畏的態度去過我們的生活.
難不難明?你接不接受到?.
我想在無風無浪的日子裡面.

$^{201}$所有事情都似乎頗順境的時候.
甚至敬畏神其實真的一點都不難.
但是當事情不是很順,好像不是很公平.
你看到一些不公義的時候,甚至災難臨到.
或者你現在面對的一些事情是不停的.
不走開的,你長期在受壓或者受苦的狀態裡面.
這段經文是很難哄的.
我去到十一月中,十二月中就回到香港探望我的親戚朋友.
我回到香港的時候,周圍的環境.
我所接觸的人都有一個很傷感的感覺.
不知道是我自己敏感還是的確是這樣.
因為以前我一些很熟悉的店舖.
很多已經關門大吉.
我所接觸的人,特別是比較年輕的一輩.
有家庭的一輩都離開了,剩下很多老人家.
我去探望一些年長的親戚或者我一些朋友的父母.
他們由每一個週末本來是和孫子,孫女和子女玩耍.
到現在週末是沒有什麼事做.
我臨走之前,冬至就快要來到.
我問其中一個親戚.
我問她今年冬至會怎樣過.
她就指著隔壁的老公.
她說跟他一樣,我們會在家裡吃一頓飯.
我聽到她說的時候,我的心很暗.
我立刻打電話給我老公.
我說我想延長我的機票.
我想留在香港,特別陪這幾個老人家去過冬.
誰知道那些人告訴我,過冬不用在那天.
提早也可以,所以結果我就沒有轉機票.
當我在想,我接下來剩下的時間.
我盡量去陪一些有需要的人的時候.
就發生了雲福縣大火那件事.
八棟大廈,七棟著火.
是說從底一直燒到頂.
死傷無數.
我用無數這個字,原因不是政府公佈的一些數字.
因為有些逃脫了,沒事出來的人.
其實他們不是沒事的.
有些可能失去家人.
有些可能逃離,逃得出來.

$^{241}$我在香港的時候見到一個牧師.
他一見到我,就不停地說.
其實他有很多壓力.
因為他的姐姐和姐夫.
就是其中一個逃出來的災民.
很多人心靈上都在有很多痛苦.
剛好這段時間,我都在香港預備今天這個講章.
我在不斷地看這段經文的時候.
我自己都在掙扎,在反思.
究竟我可以怎樣去消化這段經文呢.
我還可不可以理直氣壯地跟大家說.
神的手仍然在.
神的恩典還在.
神的恩典是從來沒有離開過.
怎樣說呢.
離開香港之前.
我和一個朋友去了一間.
星期六晚崇拜聚會的教會.
令我最深刻的不是那邊的港島.
而是港島之後.
有一個牧師出來祈禱.
他說這段時間可能大家都想做很多事.
但偏偏有些人就全拿了東西去做.
他說你們任由他去做吧.
不需要人做,你也做的.
他說況且慢慢災民會散落在香港不同的地區裡面.
有很多事我們基督徒其實仍然可以做的.
例如口渴,給杯水去喝.
去凍,給件衣服去穿.
或者甚至我們最可以再.
真的可以做到的就是陪伴.
在這些災民反覆的情緒.
他們有時可以,有時不可以.
有時有力,有時沒力的時候.
就在他的旁邊陪伴著他.
然後他再說我們做些什麼呢.
不需要讓人知道.
主耶穌基督知道就行了.
這個牧師的幾句話.
定了我沒有機會去災區.

$^{281}$又快要離開香港的很忐忑的心情.
可能在這些時候.
發生這些事我們很難明白神的時間.
在痛苦的時候.
我們很難相信神的手還在.
神的恩典還在.
我們或許會仍然在掙扎神為什麼不施恩.
可能在不公義的裡面.
我們很難敬拜那位可以施行公義.
卻又似乎沒有施行公義的神.
的確是很難.
但去完這個崇拜之後.
我發覺原來.
可能在這些時候更難的是.
若果不回到神的家裡.
回到教會裡.
和一群有共同信仰的人.
我們一起去掙扎.
我們一起去心痛.
我們一起去譴責不義.
一起堅守.
我們一起仍然定義去選擇.
掌管一切三位一體的真神.
如果我們不這樣做.
其實是更加難.
萬事都有定期.
天下萬物都有定時.
上帝的計劃.
我們人的責任不是要想出來.
想出神的時間.
不是我們的工作.
或許我們不是要完全明白.
然後我才去信服.
這裡的經文是要提醒我們.
神對我們的要求.
是在無論什麼境況裡.
我們要有一個敬畏祂的心.
我們一起低頭祈禱.
你是掌管時間,掌管生命的主.
幫助我們.

$^{321}$當我們以為自己可以控制所有的事物.
或許我們平平安安的時候.
就是你的恩典.
主啊,很多事情我們不明白,不知道.
但在這些不明白,不知道的事情裡.
求你幫助我們.
有一個敬畏你,信服你的心.
我們其實自己沒有這樣的能力.
我們都不懂得怎樣去敬拜一位又真又活的神.
一個滿有憐憫,充滿愛的神.
求你幫助我們,求聖靈幫助我們.
主啊,願你能夠在我們的生命裡得到榮耀.
我們敬畏你的心,能夠彰顯你的大愛.
\newpage



\section{}
\label{sec:Ly_O4TyukTE}
\textbf{Seen and Known (Psalm 139) - Rev. Mark Chan 被神認識 (詩139) - 陳崇基牧師}
\newline
\newline
連結: \href{https://youtube.com/watch?v=Ly-O4TyukTE}{\texttt{https://youtube.com/watch?v=Ly-O4TyukTE}} ~~~~ 語音日期: 2026-02-02
\newline
\newline
\hyperref[sec:_Xk1Tbpr5MU]{\small{< < < PREV SERMON < < <}}
~
\hyperref[sec:index_chronic]{\small{[返順時目]}}
~
\hyperref[sec:index_scriptual]{\small{[返順卷目]}}
~
\hyperref[sec:L_CPDfFPb5k]{\small{> > > NEXT SERMON > > >}}
\newline
\newline
$^{1}$sharing God's message..
(speaking in foreign language).
- Good morning, everyone, and welcome to UV,.
and I know it's quite a hassle to come to church today,.
and I'm driving an all-wheel drive BMW,.
and I still kind of drift on the side road especially,.
so it was kind of dangerous, yeah..
So you guys got here okay?.
Right?.
(speaking in foreign language).
- Okay, so we'll continue our mini-series.
on the Psalms in the Bible,.
and today we are at Psalm 139..
And 139, Psalm 139, and if you have a Bible app.
or a printed Bible, I want to invite you.
to turn to 139, Psalm 139..
It is one of the most beloved Psalms in the Bible..
It is the most important Psalm in the Bible..
It's Psalm 139..
It is one of the most beloved passages,.
Psalm poems in the Bible,.
and it is often read in many different situations..
You know, whether we go to a hospital to visit someone,.
or we read the Psalm at funerals,.
or moments of deep anxiety, or even birthday celebration,.
we take 139, Psalm 139, and use it..
And people turn to this Psalm because it speaks.
to all of us something about human existence,.
something universal, and the longing,.
human longing to be known by someone..
(speaking in foreign language).
And we live in a world that is often being seen.
by somebody else, right?.
In this high-tech world,.
we are always being seen all the time,.
and cameras are everywhere, right?.
Cameras are everywhere, and there was GPS tracking,.
and the big data is observing us and tracking us down..
And how many of you got a speeding ticket recently?.
How many of you?.

$^{41}$(speaking in foreign language).
Okay, so we live in the world..
I got a speeding ticket not too long ago on accident,.
actually, yeah, I was rushing to church,.
and I coming to worship, okay?.
I was kind of late for music practice,.
and I got a ticket, all right?.
So we do live in the world often being seen.
and often being observed by someone else..
(speaking in foreign language).
Imagine you're being seen by someone else,.
and the feeling of it is like you're being seen,.
and it feels like sometimes feels like losing your control..
We fear about being seen and being misunderstood.
or being reduced by somebody's sight..
We're being judged, and this fear shapes.
how we relate to other people.
and even how we relate to God himself..
(speaking in foreign language).
There's a modern French philosopher..
His name is Jean-Paul Sartre,.
and he's an atheist philosopher,.
and there are some books that he wrote that he argued.
being seen is actually a threat to human freedom, right?.
That makes sense, right?.
You're being seen all the time,.
and you don't feel like you're being free all the time,.
and Jean-Paul Sartre gave us some examples,.
and one of them is that you imagine a time.
that a stranger stares at you..
There are strangers staring at you,.
and that feeling of it is being,.
your existence being reduced into an object of someone else.
and an object in their eyes,.
and you don't feel, at least you don't feel comfortable,.
right?.
(speaking in foreign language).
Now, all of us drive, right?.
All of us drive on the highway, on the road, okay?.
So you kind of find that experience familiar..

$^{81}$Just imagine you're driving on the highway,.
and you're free, and you're doing your thing, right?.
There's no pressure, and you enjoy the drive,.
and you look around, and then all of a sudden,.
a police car is following you,.
and then you check the back mirror,.
and you saw the car behind you, and how do you feel?.
You feel pressure, right?.
You feel kind of like very uptight,.
and you try to be careful, and you focus on the road,.
and especially what?.
The speed limit, right?.
The stop sign, and then you try to drive carefully..
Now, that is the experience of being seen.
and being judged all the time..
(speaking in foreign language).
Now, being seen feels like being judged, right?.
Now, the difference is who is seeing you,.
who is looking at you, right?.
If it's a police car following you,.
you feel being judged, and you feel pressure,.
but as we turn to Psalm 139,.
we find that there is someone else seeing us all the time..
(speaking in foreign language).
So the first thing here is that,.
try to advance the PowerPoint here, yeah..
The first point here in this Psalm 139 says,.
"God sees us completely.".
(speaking in foreign language).
Now, Psalm 139 opens with the third,.
"Fear we think that John Paul Sarge find it fearful.".
(speaking in foreign language).
It says, "O Lord, you have searched me, and you know me..
"You know when I sit, and you know when I rise..
"You perceive my thoughts from afar..
"You discern my going out and my lying down..
"You are familiar with all my ways..
"Before a word is on my tongue,.
"you know it completely, O Lord..
"You hand me in behind and before..

$^{121}$"You have laid your hand upon me..
"Such knowledge is too wonderful for me,.
"too lofty for me to attend.".
(speaking in foreign language).
And that is one of the Psalms that we have just did, right?.
And God sees and God knows here..
And God searches the depth of your life..
And God knows when you sit, when you sit down,.
or when you stand up..
Now, sitting is kind of a posture of resting, right?.
You wanna take a break and you sit down..
And the standing is another posture of working, probably,.
especially for the ancient people..
So the Psalms says God pays attention to you.
when you are resting, when you are working..
God seeing means God is caring for you..
(speaking in foreign language).
And here, God is not like a police,.
it's not like a police officer trying to judge you.
and find some faults in you..
Here, God is caring for you with His sight.
and giving you comfort..
And don't you find comfort in a God.
who always cares for you and sees you?.
(speaking in foreign language).
And here, the gaze of God is total.
and there's no hiding from God..
And we cannot hide from Him..
There's no escape from the sight of God..
(speaking in foreign language).
And then the second point here.
is that God is present everywhere..
So let us look at the scripture here, 7 to 9 first..
Where can I go from your spirit, God?.
Spirit, God, where can I flee from your presence?.
If I go to the heavens, you are there..
If I make my bed in the depths,.
the deepest part of the world, you are there..
If I rise on the wings of the dawn,.
it means the extremity of direction..

$^{161}$If I settle on the far side of the sea, go ahead..
(speaking in foreign language).
And then it continues to say,.
even there your hand will guide me.
and your right hand will hold me fast..
If I say surely the darkness will hide me.
and the light become night around me..
Even the darkness will not be dark to you,.
the night will shine like the day.
for darkness is as light to you..
(speaking in foreign language).
So here you hear the poet trying to pray to God.
who is always present..
And this doesn't seem to be a cry.
of someone trying to escape from God..
And it's a realization that God's presence cannot be lost..
And there's no place that God is absent..
(speaking in foreign language).
Now, the poet is David..
I mentioned to you the philosopher John Paul Sartre, right?.
For John Paul Sartre, such a knowledge.
that God is always seeing you,.
being seen all the time means losing your freedom.
and your existence because you are always being looked at,.
right, being stared at..
And then for David, for the poet here,.
being seen is a totally different thing.
because here being seen is being seen.
by a God who cares about you..
And the gaze of God does not disappear.
when life becomes dark, when you can't see in the dark..
God is still seeing you and God does not withdraw.
when circumstances go dark or go bad..
(speaking in foreign language).
Now, this is interesting.
because it's the who is looking at you, right?.
It's not a police officer, it's a God who cares about you..
And the presence of God does not expose us.
in order to shame us..
The presence of God accompanies us.

$^{201}$in order to sustain our life..
(speaking in foreign language).
And there's no height,.
there's no height that God is not impressed by us..
And there's no depth,.
no, the deepest part where God is discussed by us..
And here the Psalm says,.
God is always present in good things and in bad things,.
in your success, in your failures,.
in your obedience, in your confusion..
And to be seen by God to David.
is loved by God..
(speaking in foreign language).
I remember when our first child, Rachel was born.
and she was days old or a few weeks old..
And she was lying on her crib..
And I, of course, I'm a loving father..
And I look at her and I see her and I care about her,.
even though she is not paying attention to me..
And she was sleeping and I see every single move of hers..
And this is the idea.
the Psalms trying to provide for us.
is that God is not like a police officer.
or a judge trying to catch us doing wrong,.
trying to issue a ticket or citation to us..
God is a caring God, always looking at you..
(speaking in foreign language).
And here the psalm continues to tell us.
to invite us to know this God who cares about us..
(speaking in foreign language).
Let's continue to read this psalm.
from verses 13 and following..
For you create my inmost being,.
you knit me together in my mother's womb..
I praise you because I'm fearfully and wonderfully made..
Your works are wonderful..
I know that full well..
My frame was not hidden from you.
when I was made in the secret place,.
when I was woven together in the depths of the earth..

$^{241}$(speaking in foreign language).
And here continue in verse 16..
Your eyes saw my unformed body..
All the days ordained for me were written in your book.
before one of them came to be..
How precious to me are your thoughts, O God..
How vast is the sum of them..
Were I to count them,.
they would outnumber the grains of sand..
When I awake, I am still with you..
(speaking in foreign language).
- So where does the knowledge of God about you begin?.
God's knowledge or God's seeing of you.
did not begin when you are doing good, all right?.
God saw you before you perform,.
before any of your activities,.
before your failure,.
even before your self-awareness of your existence..
This means our lives are not accidents to God..
(speaking in foreign language).
- So God's seeing of us precede our shame.
or even our existence..
And that's why David, the poet, can say,.
I praise you because I am fearfully and wonderfully made..
God seeing you gives life rather than death..
(speaking in foreign language).
- And here at the end,.
the psalm turns personal and prayerful..
It says verses 23 and 24..
Search me, O God, and know my heart,.
and test me, and know my anxious thoughts..
See if there's any offensive way in me.
and lead me in the way everlasting..
(speaking in foreign language).
- And this is, I think, the most beloved line.
in the entire psalm,.
or the most surprising line in the psalm..
And the psalm invites us to follow God and to trust Him..
And it does not ask God to look away from us,.
ask God to look even deeper in our lives..

$^{281}$(speaking in foreign language).
- And psalm 139, let me conclude,.
calls you, or all of us, to know God..
And it's not just knowing about Him, to know Him..
And to be known by God is actually grace,.
and to grow in Him..
(speaking in foreign language).
- And knowing God here means we stop hiding from Him..
We can be honest and we can be open.
and transparent before God..
And it means we allow God not only to comfort us,.
but to lead us..
And the gaze of God, the sight of Him,.
does not merely reveal who we really are,.
but lead us into ways of everlasting life following Him..
(speaking in foreign language).
- And today, I wanna invite you to conclude with a prayer,.
to come before this God, and to accept Him,.
and to love Him back, to relate to Him..
So let us all pray together..
Let's pray..
(speaking in foreign language).
- Lord, you see us completely, and you still love us..
(speaking in foreign language).
- Thank you, because your gaze is not meant to shame us,.
but to guide us and to love us..
(speaking in foreign language).
- And help us, Lord, not only to rest in being known by you,.
but to truly know you..
(speaking in foreign language).
- And with the psalm, we wanna conclude.
with our heartful prayer, lead us, God, in the way everlasting..
(speaking in foreign language).
- In Christ, we pray..
- Amen. - Amen..
- Well, God bless you, and always be seen by God this week,.
and knowing that God is always watching over you.
and caring for you, all right?.
(speaking in foreign language).
\newpage



\section{約翰福音 1:1-5}
\label{sec:L_CPDfFPb5k}
\textbf{現身的上帝(約翰福音1\_1-5) - 梁紫健傳道}
\newline
\newline
連結: \href{https://youtube.com/watch?v=L-CPDfFPb5k}{\texttt{https://youtube.com/watch?v=L-CPDfFPb5k}} ~~~~ 語音日期: 2026-02-02
\newline
\newline
\hyperref[sec:Ly_O4TyukTE]{\small{< < < PREV SERMON < < <}}
~
\hyperref[sec:index_chronic]{\small{[返順時目]}}
~
\hyperref[sec:index_scriptual]{\small{[返順卷目]}}
~
\hyperref[sec:XrYioYwg95Q]{\small{> > > NEXT SERMON > > >}}
\newline
\newline
約翰福音 1:1-5
\newline
\begin{longtable}{cl}
\hline
\hline
章節 & 經文 (和合本修訂版)\\
\hline
1:1 & \begin{tabularx}{0.7\textwidth}{X} 太初有道，道與神同在，道就是神。 \end{tabularx} \\ \\ \relax
1:2 & \begin{tabularx}{0.7\textwidth}{X} 這道太初與神同在。 \end{tabularx} \\ \\ \relax
1:3 & \begin{tabularx}{0.7\textwidth}{X} 萬物都是藉著他造的，沒有一樣不是藉著他造的。凡被造的， \end{tabularx} \\ \\ \relax
1:4 & \begin{tabularx}{0.7\textwidth}{X} 在他裡面有生命，這生命就是人的光。 \end{tabularx} \\ \\ \relax
1:5 & \begin{tabularx}{0.7\textwidth}{X} 光照在黑暗裡，黑暗卻沒有勝過光。 \end{tabularx} \\ \\
[1ex]
\hline
\hline
\end{longtable}
$^{1}$第五節目早安,很高興大家在這麼大雪的日子裡都平安回到家中.
不知道大家有沒有見過這個建築呢?.
有沒有見過?.
熟悉香港的朋友,不知道有沒有見過這個建築呢?.
這個建築是在香港的.
大家覺得這個建築是屬於什麼宗教的建築呢?.
佛教,道教,是的.
很相似的,但其實它是一個基督教的建築物.
那為什麼呢?我們可以講一下.
可以幫我按下一張.
這個原來是杜豐山.
這個是位於杜豐山的一個建築.
在香港沙田.
是由一個信義宗的宣教士去建立.
當時他們為了向中國佛教和道教的信徒傳福音.
他們專門將教會這個聚會的地方.
建立成一個好像佛教寺院的模式.
希望用一個中國化的形象來傳福音.
所以在三十年代開始.
他就開始規劃.
然後希望那些佛教徒和道教徒.
可以在一個自己熟悉的宗教環境底下討論信仰.
以致他們能夠接觸到基督教.
所以當年宣教士為了傳福音給中國人.
他用中國人很熟悉的方式.
用道教,佛教的方式去包裝.
去向中國人傳福音.
才將這座建築物建成這樣.
同樣地.
上帝知道我們人類不認識祂.
也不太熟悉祂.
所以這個本來對於我們來說的上帝好像很遠.
隱形的.
在舊約年代.
人是沒有辦法直接去接觸上帝.
如果你有看舊約聖經.
你會知道透過誰說要.
如果要接觸上帝的話.
透過祭司,透過摩西,透過先知.
你才可以知道上帝的心意.

$^{41}$但是到了新約.
上帝為了我們這群.
對祂不是認識很深的人類.
祂專門用了一個我們很熟悉的方法.
讓我們去認識祂.
祂親自去道成肉身.
來到人間成為一個現身的神.
就好像剛剛第三首歌一樣.
我很喜歡歌詞.
祂翻身月亮為我來.
然後從至高之處降身為了愛.
祂就是這樣的一個上帝.
所以今天我們會看看.
這位曾經隱身見不到的上帝.
怎樣在耶穌的道成肉身裡面.
用我們一個理解到的方式.
在我們身邊去呈現.
所以今天我們聽完這篇講道之後.
我們會知道兩件事.
究竟這位上帝怎樣透過.
我們熟悉的方式去揭露.
祂對我們生命的指引.
第二 祂對我們生命中的.
一些黑暗的位置又有什麼啟迪.
首先我們一定要讀一下經文.
就是約翰福音的一章一至五節.
太初有道 道與上帝同在.
道就是上帝.
這道太初與上帝同在.
萬物都是藉著他所做的.
沒有一樣不是藉著他所做的.
凡被做的在他裡面有生命.
這生命就是人的光.
光照著黑暗裡.
黑暗卻沒有勝過光.
上帝透過創造.
開始了這個世界.
一開始人是完全見不到上帝本人.
但是我們可以透過他的說話.
去看到他所創造的世界.

$^{81}$所以經文是這樣說的.
太初有道 道與上帝同在.
萬物都是藉著他所做的.
意思是說上帝一開始.
創造這個世界的時候.
是用他的說話去行事的.
如果大家還記得創世記一章的話.
他說什麼.
上帝說我要有光.
就有光.
他說要有天就有天.
我要有地就有地.
是他說就有了.
起初上帝的說話.
就是他的話語.
很明顯世界上所有的東西.
都是藉著他的話語所創造出來的.
這套和學本他用太初有道的道.
這個道其實是一個中國化的翻譯.
如果你知道和學本的歷史.
一百多年前.
翻譯者在一百年前的時候.
就採用了中國人所說的道.
就是智慧的意思.
意思是指上帝充滿著智慧.
能夠創造這個世界的萬物.
而在中國度家哲學裡面.
道是作為萬物源頭的根源.
無形無處不在.
對人類是不偏不倚的.
但是這個翻譯.
我們知道一百年前.
西方宣教士很想中國人.
得到福音的時候.
這個道的翻譯其實是比較切合的.
令到喜歡或者認識中國哲學的朋友.
可以交流到.
但是這個翻譯給人有一種距離感.
就因為上帝就好像道家裡面的道.
遠遠地看著你.

$^{121}$你又見不到他.
你又捉不到他.
但是你又好像感受不到他.
但是如果我們把這個字直譯成.
話語或者說話.
神藉著他的話語和說話去創造這個世界.
感覺上就很不同.
就是上帝主動說話去創造這個世界.
而且他對這個世界是有他的旨意和目的.
而不是只是說道家的那種道.
是絕對中立自然地去產生這個世界的秩序.
所以約翰福音這裡說的是.
上帝主動用他的話語去創造.
他希望這個世界不單止是他的世界.
而是他希望他的秩序都在這個世界裡面.
而不是由這個世界自由運行.
所以我們要說這個概念.
通常我們提到創造.
大家都會覺得上帝創造.
那就是由無變有.
我們是不是這樣想的.
但其實創造不只是由無變有創造一些新的東西.
而在《創世記》一章裡面.
上帝創造是說將創造物各從其類.
按照上帝覺得好的旨意去分類.
就例如《創世記》一章三節這樣說.
上帝說有光就有光.
上帝看光是好的就將光和暗.
分開.
第六節種水之間又有窮倉.
將水和水分開.
第九節天下的水要聚在一起.
使乾地露出來.
將地和水分開.
第十一節地要長出植物.
要各樣的果子有種子的樹.
在地上各從其類.
我們可以簡單看幾節就知道.
上帝的創造不只是由無變有.
還要包括按照他覺得怎樣才算好去分類.

$^{161}$這個叫做創造秩序.
上帝覺得好就是好.
所以當約翰福音在頭幾節說.
萬物都是藉著他做的.
沒有一樣東西不是藉著他做的.
強調的是上帝不只是由無去創造有萬物的世界.
他的意思是萬物都要按照他覺得好的旨意去形成秩序.
創造秩序這件事.
就好像我用一個例子去說.
我太太昨天上NBTI工作坊.
她在掙扎.
究竟她是J還是P.
我覺得她像J.
因為J的人有自己收拾的秩序.
其實我是很P.
我P的人雖然有自己收拾的秩序.
但大家都不太發現.
我的桌子通常比較亂.
但我永遠都找到我的桌子需要的東西.
收拾完之後我就找不到自己需要的東西.
我太太會在廚房裡走她的創造秩序.
就是每一樣東西她都覺得是好的位置.
這個是油,這個是鹽,這個是糖.
擺好了.
各從其類.
但有時候我用完的東西就沒有擺回原位.
那就不好了.
在她眼中就不好.
創造秩序就是這樣.
創造不單止是創造了一個空間.
她在說什麼?.
她在說空間裡面有一個創造的主宰.
他實行了他的秩序.
有些東西要各從其類去打理.
所以我們花這個時間去說創造秩序.
是上帝用自己的秩序實行在這個世界上.
目的是我們要明白這句.
光照在黑暗裡面.
黑暗卻沒有勝過光.
通常我們讀日文福音.

$^{201}$很快就會帶過這句.
因為這句很容易明白.
很貼近我們的生活.
例如我用電筒照.
基本上照到哪裡.
沒有一個地方是光不能勝過的.
因為黑暗是會消失的.
你這樣照下去.
就會看到光是會被克服的.
除了一個地方.
例如宇宙的黑洞.
它吸了光就會看不到光.
所以我們很自然就會去理解.
這句話很科學式的去理解.
我們會覺得我們之所以看到光.
是因為我們眼睛有感光細胞.
我們之所以看到光.
是因為有感光細胞.
接收到光的反射.
以致我能夠在腦裡.
感受到光源的訊息.
所以科學家就會定義黑暗是一個狀態.
它並不是一個存在的東西.
它會解釋這個角落之所以暗.
是因為這裡沒有光.
而不是因為黑暗存在.
如果我們用這個科學的角度.
去理解這段經文的話.
這段經文就突然變成一本科學書.
一個科學指南.
其實如果你畫了它.
對於整篇經文都沒有太大的影響.
如果我們這樣理解的話.
但是.
我們用創世記上帝創造秩序去理解這句話.
我才想到的東西就完全不同了.
這句話說的是光照在黑暗的裡面.
就好像上帝說光是好的.
光和暗是互相不相容的.
光是負責將上帝的旨意.

$^{241}$即是生命注入人的裡面.
以至到人類的黑暗被揭示.
被克服.
被轉化.
就好像上帝拿著電筒.
它照到我們內心的深處.
這樣說好像有點玄.
好像道士說風水一樣.
我講一個例子給大家.
可能會明白一點.
我和Alice私底下都可能有和大家說過.
我們從結婚開始就有計劃想要下一代.
要有小朋友.
很自然這幾年我們都照常生活.
要做的事我們都做齊了.
但是有一天我們為了一件小事去吵架.
吵得很激烈.
我們兩個都很生氣.
我就說了一句.
我都不覺得我們小朋友會開心快樂地成長.
因為我覺得我們兩個相處這麼多張力.
這麼多爭吵.
如果有一個很無助的嬰兒存在.
我們會更加多磨擦.
大家做過父母都知道.
一個無助的嬰兒.
無論我們怎樣吵也好.
換尿片也要餵奶.
他不會自己突然起來.
父母說他吵架.
我自己換尿片和吃奶.
不會這樣的.
但是我說這句話的背後.
其實我很害怕.
我害怕這個嬰兒就像我童年一樣.
很多苦難.
加上我有腦癌.
我都很不想我下一代再一次經歷我這種無助和困難.
我很害怕會有遺傳.
我遺傳我下一代令他也有我這種苦難.

$^{281}$所以這樣吵吵架.
我才發現原來我口說我準備好會有下一代.
但是我內心一點也不準備好.
一點也不準備好.
這個世界充滿著苦難.
小朋友沒有選擇.
一旦他被人選上.
他就要經歷這個世界無常的痛苦.
還要有腦癌的話.
我真的不知道該怎麼辦.
我的內心深處有一個黑暗的位置.
這個恐懼抓住我.
令我無法喘息.
令我和Alice的關係也很有張力.
然後我們為了這件事糾結了一段時間.
我們覺得我們很需要外力去幫助.
上帝藉著輔導員的口.
去將生命的光照進我的內心.
我和他談了一大輪.
輔導員提供了兩個我從來沒有想過的角度.
第一就是我從小都有腦癌症.
經歷了這麼多年之後.
我也熬過來了.
其實我就是最懂得.
在這個病中照顧自己和別人的人.
我非常適合成為一個陪伴小朋友成長的人選.
第二就是在結婚之前.
其實Alice已經知道我有這個病.
但她仍然願意和我有下一代.
這個願意對我來說.
就像一道光照進了我的生命裡.
對我來說是一個很大的意志.
聽完之後我整個人在顫抖.
也不是人有三急.
是因為我感受到這一句.
生命在他裡頭.
這生命就是人的光.
光照在黑暗裡.
黑暗卻不接受光.
我內心深處的恐懼.

$^{321}$馬上被光照到無法反抗.
我馬上投降.
生命的黑暗就是一些.
我們平時都沒有留意.
一些我們生命中的盲點.
用這個例子希望大家明白.
這一句光照在黑暗裡.
黑暗卻沒有勝過光.
不是在說科學.
而是在說耶穌他親自成為生命的光.
將上帝的旨意照進我們內心的深處裡.
重新將創造的秩序.
擺放在我們人的生命裡.
就好像玩LEGO一樣.
當我們理解道與神同在的時候.
這不僅是上帝擁有創造的基石.
更是耶穌的跨越.
從上帝的話語成為生命的光.
重新將秩序擺放在人身上.
如果你玩LEGO的話.
你親手建造的樂園.
每一區都會有不同的設施.
和不同的LEGO人物.
可能我們會花很多時間.
在這個樂園裡.
用我們想像的方式去建立那些LEGO.
有我們喜歡的秩序.
就算我們多喜歡LEGO.
我們絕對不會想到.
我們自己會成為LEGO的一員.
由那個本身自由自在.
變成一個死物.
任人擺佈.
連你擺什麼姿勢都要主人決定.
你沒有辦法去說任何東西.
所以我們心裡面會知道.
LEGO只是我們其中一個娛樂.
我們不會一輩子想成為LEGO.
失去了人生的主導權.
即使是這樣.

$^{361}$上帝仍然願意將祂的話語.
即耶穌基督.
降生為人.
主動去到跨越這個世界的牆.
將創造的秩序.
重新安放在我們身上.
所以總結.
上帝如何透過我們熟悉的方式去揭露祂對我們生命的指引.
祂對我們生命向往的位置有什麼啟迪.
我想聽完這個講道.
我們這個問題可以一起回答.
就是上帝用了一個我們很熟悉的方式.
就是光和暗的比喻.
這個圖畫是由創世開始.
上帝將光帶到了世界.
讓我們可以分辨哪裡有生命.
哪裡有光.
哪裡剩下的只是黑暗.
在我生命的裡面.
可能會有些東西.
是未被上帝創造秩序所歸類的.
就好像我們未信主之前.
我們根本察覺不到上帝的存在.
以前我們就好像被黑暗遮蔽住.
不知道原來上帝是創造主.
更加不知道祂一直在我們身邊保護著我們.
當祂慢慢將祂的創造秩序放在我們生命的裡面.
我們才慢慢開始了解.
原來祂在我們的生命做了這麼多東西.
那時候我們才知道.
原來世界有溫暖 有光 有盼望.
老實說.
就算到現在做了傳道人.
我都依然覺得自己生命有很多黑暗的位置.
所以為什麼要這樣強調.
這個生命就是人的光.
光照在黑暗裡面.
黑暗勝不過光.
就是因為這個現身的上帝.
將光帶入世界的同時.

$^{401}$同一時間都帶入了我們的生命裡面.
令我們看世界很不同.
讓我們再一次看見生命有盼望.
所以面對著這個主動現身的上帝.
我們可以怎樣回應呢.
第一我們要先學會認識自己內心的黑暗.
讓上帝的光和生命可以照進來.
有時候我們的心會被一些生活的小事和煩惱遮蔽著.
令我們感受不到上帝的存在和他的指引.
所以我都建議大家.
在平時抽一點時間出來用心去讀一下上帝的話語.
不是好像交功課那樣.
是你用心去感受.
讓他的光就好像早上的陽光那樣.
去照亮我們的內心.
不知道大家記不記得我剛剛說過.
上帝是用他的話語去創造世界的.
他一開口這個世界就由混亂變成有罪.
由黑暗變成有光.
就好像我們平時讀經一樣.
我們看聖經裡面上帝的話語的時候.
就在聽上帝再一次對我們的生命.
將創造的秩序再一次放在我們的生命裡面.
可能你現在的生活正處於一片混亂.
可能你會覺得前路真的很混亂,很黑暗.
就好像起初未被創造的世界一樣.
但是當你打開聖經.
你讀上帝的話語的時候.
那種創造的力量就會再一次進入我們的生命.
就好像跟我們的混亂說:已有秩序.
就好像跟我們的黑暗說:已有光.
讀經就是在經歷上帝在你的生命裡面.
再一次注入創造的秩序.
第二,就是我們在生命裡面刻意發現上帝的作為.
例如我們享受祂所創造的美好.
去發現這個世界,感受祂的光.
要留意上帝在我們身上做過什麼.
在這個新的一年,我們養成一個新的習慣.
就是我們每個星期都刻意去找一樣.
上帝在我們生命裡面做過的事情.

$^{441}$或者發現上帝給我們的祝福.
在我自己的生命裡面.
我發現我太太對我接納.
是上帝給我的一個很大很大的祝福.
還有一個光.
有時候我們信主一段時間.
我們可能會忽略了一些很基本的動作.
我們以為我們信主夠久.
我們就不需要刻意去發現祂在我們生命裡面的工作.
但其實上帝每一天都在做一些新的事情.
每一天都想發現.
希望我們發現祂在我們生命裡面做的事情.
第三,我們都要讓家人和朋友察覺到上帝的存在.
當我們留心我們生命的見證.
上帝的作為的時候.
我們就會有足夠的見證的存貨.
無論什麼時候我們都可以隨時拿出來跟別人說.
其實原來上帝真的在我們身邊.
最近跟一群朋友吃飯.
當中有一位快要結婚的女生.
她還沒有信主.
期間她就問一個問題.
你怎樣肯定那個人就是你可以過一生的伴侶呢?.
當中有一個基督徒的姐妹分享令我很深刻.
她說一世的伴侶當然是懂得互相的體諒.
也要清楚對方在原生家庭裡面帶來的那套生存的機制.
但她說最重要的一件事是.
只要大家要在上帝的光照之下.
才會有悔改和改變的可能.
令彼此更加愛對方.
如果你靠自己是完全做不到的.
那個女生聽完之後就感受到上帝好像跟她很近.
她覺得那個聚會很有意義.
我想這個就是我想說的.
當我們平時留意上帝在生命裡面的證據的時候.
我們會儲到這些見證的存貨.
到當你適合的時候.
我們就可以隨時跟身邊的人分享.
告訴別人上帝真的在.
我們一起祈禱吧.

$^{481}$主啊,多謝你.
多謝你在一個很黑暗的世代裡面.
你依然現身在我們身邊.
你是一個主動進入我們的生命.
加入我們的生命.
賜給我們有創造秩序的上帝.
幫助我們.
讓我們用心留意你在我們生命裡面的足跡.
一點一滴的見證.
建立著我們跟祂的關係.
求主幫助我們.
感到這封信很敏感.
在寒冷的冬天.
希望大家都有平安.
祝福大家.
\newpage



\section{詩篇}
\label{sec:XrYioYwg95Q}
\textbf{Be Still (Psalm 46)- Rev. Mark Chan 學習安靜(詩篇46) - 陳崇基牧師}
\newline
\newline
連結: \href{https://youtube.com/watch?v=XrYioYwg95Q}{\texttt{https://youtube.com/watch?v=XrYioYwg95Q}} ~~~~ 語音日期: 2026-02-09
\newline
\newline
\hyperref[sec:L_CPDfFPb5k]{\small{< < < PREV SERMON < < <}}
~
\hyperref[sec:index_chronic]{\small{[返順時目]}}
~
\hyperref[sec:index_scriptual]{\small{[返順卷目]}}
~
\hyperref[sec:QVz_a2O_21A]{\small{> > > NEXT SERMON > > >}}
\newline
\newline
$^{1}$- Good morning, everyone..
(speaking in foreign language).
Well, we are continuing our series on the Psalms.
and we are turning to Psalm 46..
And this is one of the most well-known Psalms in the Bible,.
well-known Psalms in the world..
And many people know this Psalm.
because there's a famous line in Psalm 46,.
"Be still and know that I am God.".
Recall any memory?.
(speaking in foreign language).
Now, these words, "Be still and know that I am God,".
has comforted many, many people in many generations..
And this same Psalm,.
if you know this guy called Martin Luther,.
is one of the reformers 500 years ago,.
reformed the church..
And he wrote a hymn, "Mighty Fortress is Our God.".
And the lyrics is coming from this Psalm, Psalm 46..
And Luther, when he wrote this Psalm,.
he was under stress and under much suffering.
and persecution from the corrupt medieval Catholic church.
and in danger..
And when he find that he was being persecuted.
and he find the courage from this Psalm, Psalm 46,.
"And God is our refuge and our strength.".
(speaking in foreign language).
And then let me read the Psalm.
and we're gonna read the Psalms in three stanzas..
Psalm 46, verse one..
God is our refuge and strength.
and ever-present help in trouble..
Therefore, we will not fear,.
though the earth give way.
and the mountains fall into the heart of the sea,.
though its waters roll and foam.
and the mountains quake with their surging..
(speaking in foreign language).
And then 46, there is a river.
whose streams make glad the city of God,.

$^{41}$the holy place where the Most High dwells..
God is within her..
She will not fall..
God will help her at the break of the day..
Nations are in uproar..
Kingdoms fall..
He lifts his voice..
The earth melts..
The Lord Almighty is with us..
The God of Jacob is our fortress..
(speaking in foreign language).
Then verse eight and following..
Come and see the works of the Lord,.
the desolation he has brought on the earth..
He makes wars cease to the ends of the earth..
He breaks the bow and shatters the spear..
He burns the shields with fire..
Be still and know that I am God..
I will be exalted among the nations..
I will be exalted in the earth..
The Lord Almighty is with us..
The God of Jacob is our fortress..
(speaking in foreign language).
Well, this psalm has three stanzas.
and there are three points I wanna share with you..
The first point is to trust in God..
Trust in God for his protection..
(speaking in foreign language).
So here you can see from verse one to three,.
you're gonna pay attention.
to the geographic descriptions here.
from verse one to verse three..
And you can see there are a lot of highlights here..
There are geographical names..
Mentions the earth, the mountains, the hills, and the water..
And seems to be the focus on the water.
and the chaos..
(speaking in foreign language).
So why water?.
You know, in the ancient East,.

$^{81}$when people thought about this world,.
they were not talking about science..
They were thinking of this world.
and trying to conceptualize this world.
being surrounded by water, by chaos, by chaotic waters..
And whoever can hold back the water up there.
and down there and create this bubble for us,.
create this world, this dry ground for us.
and divide the waters..
And this is the guy we should worship..
And this is the guy who has the most power..
(speaking in foreign language).
And then, you know, the Babylonian myth.
called creation epic,.
it mentions a force God called Marduk..
And Marduk is the biggest God,.
the greatest God in the Babylonian pantheon..
And he is the young God who defeats the water..
It's called Tiamat, the sea goddess..
And he divided her and put half of his water up there,.
supporting the sky and the heaven..
And half of her body, after he killed the water,.
and to subdue it under his feet.
and create a dry ground..
And then he rules over the gods.
and put order in the cosmos.
and create a city, the greatest city of Babylon..
And he became the final authority of the entire cosmos..
(speaking in foreign language).
Now you may remember chapter one of Genesis, right?.
Now Genesis, when it talks about the creation of the cosmos,.
it did not begin with nothing..
So it's not really creation out of nothing in Genesis one..
It begins with water and chaos..
(speaking in foreign language).
And again, it's not talking about natural sciences..
It's talking about theology..
It talks about God taking control.
using familiar concepts to the ancient people..
And he, God himself, the only creator,.

$^{121}$he divide the waters and hold them back..
And that shows the power of God..
(speaking in foreign language).
So the same thing here in the psalm,.
in the first three verses,.
the writer is trying to use the same language.
and same idea to tell us to declare the power of God..
(speaking in foreign language).
And then second point here is that.
the psalm is trying to tell you.
that you have to focus on the presence,.
the dwelling of God..
Now, when we read four to seven,.
you find that it's describing a city, right?.
It's describing a city..
Now we all have certain cities in mind.
when you think of your home city.
or some of the places that brings you memory.
or security or identity..
It may be Beijing, may be Hong Kong,.
may be Chicago, or may be your home city..
Some sort of city that feels like home..
(speaking in foreign language).
And for the people of God in the Old Testament especially,.
the key city is which one?.
Jerusalem, right?.
And Jerusalem is also called Zion, right?.
So Zion, Mount Zion, is actually Jerusalem..
They are the same..
And that is why some of the psalms.
are called Songs of Zion..
Songs of Zion, praising the city of Zion..
And Psalm 46 is actually one of the Zion songs..
And because this city is worthy of our praise.
and our attention because God dwells in the midst of Zion..
And God's people do not need to fear..
(speaking in foreign language).
Well, let's read the four to seven here..
There is a river whose stream make glad the city of God,.
which is Jerusalem, Zion,.

$^{161}$the holy place where the most high dwells..
God is within her and she will not fall..
And God will help her at the break of day..
Nations are in uproar, kingdoms rise and fall,.
and he lifts his voice and the earth melts..
The Lord Almighty is with us..
The God of Jacob is our fortress..
(speaking in foreign language).
So I mentioned the Babylonian creation story.
is also called Enuma Elish,.
where the great God Marduk defeated the chaotic water.
and then divided and established the universe..
And Marduk also built a holy city..
And for the Babylonian religion,.
it is the building of Babylon, the holy city..
(speaking in foreign language).
And, but here it's talking about Yahweh, the true God..
And after he subdued the water,.
he established the city, Jerusalem, of course, right?.
That is the true city of God..
And you note that all the geographic images.
being mentioned in verses one to three, the first stanza,.
here are also repeated..
And all the imageries are also kind of,.
we call in this second stanzas, four to seven,.
but everything is reversed..
Can you see that?.
All the chaos that was mentioned in verses one to three.
are now changed into peaceful pictures..
There's a river coming out of the temple.
and God is in the midst, God is present..
(speaking in foreign language).
So you can kind of tell from four to seven,.
the movement of imageries go into total opposite direction..
It moves from chaos, mountains shaking in the sea,.
water roaring and forming to calm a river with streams..
And these are the imageries that we call the temple,.
the dwelling place of God in Jerusalem..
When we worship, when we worship together,.
we affirm the presence of God, that has never changed..

$^{201}$(speaking in foreign language).
Some years ago, I received an email,.
I heard a story about a Christian woman in Syria..
And as you can kind of recall some of the history,.
recent history, that ISIS was destroying many lives..
And then some missionaries visit a village.
and met this woman, this Christian woman..
And they saw the name of Jesus being posted in her home.
on almost every wall inside her house..
And this is one of the ways she coped.
with the terror brought by ISIS..
(speaking in foreign language).
And when she was terrified, she was trying to calm.
the fear in her heart..
And that's why she put the name of Jesus in Arabic,.
of course, on every wall in the house..
And she believed there was a power,.
there was a power in name of Jesus..
And the fear is will, of course,.
you're not gonna deny the fear,.
but fear can be brought under control,.
especially under the name of Jesus..
(speaking in foreign language).
And then the third point here from verses eight to 11,.
it tells you to be still before God..
(speaking in foreign language).
And eight to 11, let me conclude.
by reading the last few verses..
Come and see the works of the Lord,.
the desolation he has brought on the earth..
He makes war cease..
To the ends of the earth, he breaks the bow.
and shatters the spear, and he burns the shields with fire..
Be still and know that I am God..
I will be exalted among the nations..
I will be exalted in the earth..
The Lord Almighty is with us..
The God of Jacob is our fortress..
(speaking in foreign language).
See, when you're in panic mode,.

$^{241}$I'm not sure you can recall some of the memories,.
some of the experiences that you had.
when you were in panic mode..
It feels like chaos to you, right?.
And this little story has to do with me..
This happened to me last Friday night..
My wife called me, and she said she was on the highway,.
and one of her tires got flat,.
and she was trying to bring students home,.
but she could not drive anymore..
So I received a call and tried to rush.
out of the church building and go into my car.
and then try to rush to her,.
and I also called my son and yelled at my son,.
and I said, "Son, you have to go..
"You have to go with me, okay?.
"Wherever you are, just drop everything.
"and go and rescue your mom, okay?".
So I drove a couple blocks, and then I find that.
I lost my phone..
Where is my iPhone?.
I tried to search my pockets,.
all the pockets of the jacket, all the pockets of my pants,.
and I tried to look for my phone..
I could not find it, and I was really, really.
in a panic mode, and I was so,.
I guess I was so angry..
I was so angry with me, and I was so angry with the iPhone..
And I met someone at the door of the church,.
and Chi-Mei, Chi-Mei's one of the worship leaders,.
and then it's very interesting..
Chi-Mei talked to me and said, "Pastor Mark, calm down.".
And then she had me log into the Find Me app.
and tried to ring the phone, you know..
And then, you know where the phone is?.
In my pocket, yeah, in my pocket..
(speaking in foreign language).
And at that moment, I really experienced,.
you know, the panic mode really took over me..
You know, in fear, in chaos..

$^{281}$And we had some MBTI, you know, you heard about MBTI, right?.
MBTI personality test..
We have some sort of workshops in the past weekend,.
and I realized myself that when I am under stress,.
my feelings can take over..
I got frustrated, and that night, you know,.
that night, Friday night, I was in panic mode,.
and everything fell out of control..
(speaking in foreign language).
And let me tell you another story about fear..
Some years ago, I got a call from a young adult woman,.
and she called me and cried out,.
because she has lost $4,000 to $5,000..
And what happened was that she received a call.
from so-called IRS..
So IRS called her and said,.
oh, you have not been reporting your taxes.
for, you know, so many years,.
and that's why you own IRS for $4,000 to $5,000,.
and this is what you need to do..
And that caller, the scam caller,.
told her to go to the Apple Store,.
to go to Apple Store and get gift cards.
of worth of $4,000 to $5,000..
You know, you know how many gift cards you have to get?.
And she was in panic mode because she,.
she remembered she had not been reporting the tax, right?.
That's why you need to file your taxes, right?.
And then she was in panic mode, and she went to,.
indeed, she went to the Apple Store.
and got $4,000 to $5,000 worth of gift cards..
You know, how many gift cards were there?.
And brought all gift cards home,.
and then, you know, swiped the code,.
and one by one, she told the code,.
the code, the secret code to the caller,.
and gave the caller the money..
(speaking in foreign language).
See, if she was not in panic mode,.
she could have called someone who can trust, right?.

$^{321}$And, you know, try to find out more information..
It was a scam, right?.
But, you know, when fear takes over,.
then we do all sorts of irrational action..
You know, when the people of God live in Jerusalem,.
they often, living in fear,.
military attacks from other empires,.
Assyria or Egyptian empire,.
in their fear, sometimes they panic,.
and they ran after the false gods.
and false ideologies to solve the problem..
(speaking in foreign language).
And the Psalm 46 is trying to declare to us,.
God is still in control..
That's why you need to run to God.
instead of having fear taking over you..
You know, last week, many people watched this guy.
climbing, you know, Taipei 101..
Did you see that on Netflix?.
Maybe you heard about it, heard about this story..
Alex Hono, he's the most famous.
free solo climber in the world..
He climbed Yosemite Half Dome.
and all these tall buildings..
And, you know, last week, he took 90 minutes.
to climb from ground level.
all the way to the top of Taipei 101..
And you can see that it can illustrate.
the picture of, you know, handling fear or managing fear..
And the danger is real, but he prepares with great care..
(speaking in foreign language).
And I find that some of the comments.
are really, really interesting about, you know, this story..
As you can see, he is climbing the building bare hands..
No safety net or any kind..
Just climb and climb and climb..
And some of the comments, I find it very, very interesting..
You know, people say he is not fearless..
He's not fearless, but he is so prepared.
that he will not die..

$^{361}$And then some of the people said, you know,.
we often say we step out of our comfort zone, right?.
We have to challenge ourselves.
to step out of our comfort zone..
But for Alex, it's a totally different thing..
He said it's not stepping out of your comfort zone,.
but to expand your comfort zone..
You know, making something uncomfortable become comfortable..
And, you know, some of the comments says what you did,.
what we did, what we did not see was thousands.
and millions times that he repeated the same moves..
The moves of, you know, grabbing something.
or stepping on something..
Thousands of moves that we did not see..
And some of the comments says he is not crazy..
He is actually extremely rational.
because he practiced the same route more than 40 times..
And he wrote down every single detail,.
all the handholds and all the footholds in his notes..
And he studied it..
He mentally prepared himself for that..
(speaking in foreign language).
And Psalm 46 does the same thing..
It does not pretend that the world is not shaking..
The world indeed sometimes feel like falling.
and the nations are raging..
Psalm 46 is not trying to deny all this fears.
and chaos in our lives..
But, you know, the people of God in Jerusalem,.
they face very many fear and threats every day..
And it was natural for Israel to trust their own strength,.
to trust their own power,.
their own economic or military power..
(speaking in foreign language).
However, the Psalm tells you to shift your attention.
from fear to the power of God, the presence of God,.
so that your heart can be still.
because the Lord Almighty is with us..
And when we trust this,.
we have a different postures to live..

$^{401}$You know, we live differently..
We will choose peace..
We will choose justice..
We will choose kindness..
And we will seek first the kingdom of God..
That is your basic postures change..
Your mindset will change and we shape your priority..
We shape your values..
(speaking in foreign language).
Well, let me conclude with three points here..
Maybe you can take home with this three points..
First, is to pray first and then act secondly..
When anxiety hits us, the Psalm says,.
"Run to God before you run into trouble.".
(speaking in foreign language).
Solutions, maybe a sentence prayer,.
"Lord, be my refuge at this point..
"Be my strength right now.".
(speaking in foreign language).
And secondly, you can practice stillness every day..
You know, just like Alex practiced thousands of times,.
every foothold and handhold..
We, in our soul, we need practice..
You know, take at least 60 seconds or a few minutes a day..
You know, practice, be still in the presence of God..
Put your phone down, turn off your computer,.
breathe slowly, surrender yourself,.
and say to God, "You are God..
"I trust in you.".
(speaking in foreign language).
And third, to live the Psalm as a witness in our lives..
You can repeat the chorus every day,.
the chorus being, "The Lord of hosts is with you..
"Be still and know that I am God.".
Choose to repeat the chorus and choose to choose one action.
that will show the stillness and the peace that we have..
Maybe speak to one another gently..
You know, sometimes we get tensions in our family.
with relationships, with just the chaotic world..
We can act gently, with kindness,.

$^{441}$to show that, you know, the Lord is in control..
The Lord of hosts is with us..
(speaking in foreign language).
Amen. - Amen..
- Let's pray..
(speaking in foreign language).
Father, you are our refuge and our strength,.
a very present help in trouble..
Our heart race and our minds feel like chaos..
So bring us back into your presence.
and teach us to be still and to know that you,.
you alone are our God..
(speaking in foreign language).
And help us to practice peace this week..
Slow us down and lead us into praying first..
And we shape our words and our actions.
with gentleness and kindness..
Make our homes, our church, a place of calm faith.
so that other people can see that the Lord of hosts.
is with us, you are with us..
(speaking in foreign language).
- And in Jesus' name, we pray..
(speaking in foreign language).
- Amen..
\newpage



\section{詩篇}
\label{sec:QVz_a2O_21A}
\textbf{學習安靜(詩篇46) - 陳崇基牧師}
\newline
\newline
連結: \href{https://youtube.com/watch?v=QVz_a2O-21A}{\texttt{https://youtube.com/watch?v=QVz\_a2O-21A}} ~~~~ 語音日期: 2026-02-09
\newline
\newline
\hyperref[sec:XrYioYwg95Q]{\small{< < < PREV SERMON < < <}}
~
\hyperref[sec:index_chronic]{\small{[返順時目]}}
~
\hyperref[sec:index_scriptual]{\small{[返順卷目]}}
~
\hyperref[sec:espYGx8NsQM]{\small{> > > NEXT SERMON > > >}}
\newline
\newline
$^{1}$有請Mark Sir.
好,謝謝….
我們繼續看詩篇.
今天選了詩篇是第46篇.
46篇,記得是說甚麼嗎.
這篇詩篇是最多人知道的.
而且有很多家傳戶曉的經文和詩歌.
都是出自46篇,記得嗎.
記得有一句「你們要安靜」.
「要知道我是神」.
「Be still and know that I am God」.
記得嗎?就是從46篇出來的.
還有記得那些老牌信徒嗎.
應該記得有一首詩歌.
名為「Mighty Fortress」.
「你是我堅固保障」.
記得嗎.
這裡好像不是很多老牌信徒.
這些是很舊的詩歌.
就是以前改教之父.
大約四,五百年前Martin Luther.
寫的一首詩歌.
他的歌詞出自46篇詩篇.
但當Martin Luther寫這首.
名為「Mighty Fortress」.
「堅固保障」的時候.
他不是很安穩地在Lazy Boy上寫的.
他是在一個很大壓力的時間下.
當時改教一,五年多的時候.
他受到當時中世紀末段.
腐敗了的教廷迫害他.
於是在他的靈修和默想中.
他很掙扎.
然後看到46篇的時候.
他寫了一首名為「堅固保障」.
「上帝是我們堅固保障」.
不如我們一起看這篇詩篇,好嗎.
我們先讀,有11節.
今天無論你正在經歷甚麼.
可能你心裡有很多驚慌.

$^{41}$面對自己,面對前途.
或者面對你現在不同的處境,困境.
這首詩歌希望帶給你一些力量.
46篇,第一節.
「上帝是我們的避難所」.
「是我們的力量」.
「是我們患難中隨時的幫助」.
「所以地須改變,山須搖動到海深」.
「其中的水須澎湃翻騰」.
「山須因海漲而顫抖」.
「我們也不害怕」.
「有一條河,這條河分叉使上帝的城寬起」.
「這城就是至高者居住的聖所」.
「上帝在其中,城必不動搖」.
「到天一亮,上帝必幫助這城」.
「萬邦軒揚,國度動搖」.
「上帝出聲,地就融化」.
「萬君子也說與我們同在」.
「那國的上帝是我們的避難所」.
「你們來看耶和華的作為」.
「看他使地怎樣荒涼」.
「他止息戰爭直到地極」.
「他折弓斷槍,把戰車焚燒在火中」.
「你們要休息,要知道我是神」.
「我必在列國中受尊崇」.
「在全地也受尊崇」.
「耶和華與我們同在」.
「那國的上帝是我們的避難所」.
這篇詩篇很清楚有三節.
第一段提到我們如何信靠上帝.
以及祂的保護.
在第一至第三節.
提到上帝是我們的避難所.
留意這裡有一些字句,言詞.
是經常出現的.
例如「山」,「地」,「水」.
甚至是「海」.
哪樣東西出現得更多.
「水」是出現得更多的.
「水」或「海心」或「深海」.

$^{81}$為甚麼呢.
在古代近東的世界.
他們對於宇宙的觀念.
就是整個宇宙都像被水包圍一樣.
這裡不是說科學,宇宙學.
不是說現代的科學.
是以往的詩歌或希伯來人所想的世界.
是一種現象的世界.
你看上天,那天是甚麼顏色.
是白色嗎.
我們最近下雪太多了.
通常是藍色的.
通常是藍天,藍天白雲.
純粹我們觀察的現象就是藍色.
藍色就是甚麼?就是水.
即是水包著這個世界.
當水包著這個世界的時候.
我們不再有一個空間的時候.
我們就不能夠在這裡生存.
不能夠在這裡生活.
誰能夠頂得住這些水.
有一個窮倉,有一個鑊在上面.
能夠把它頂住.
有那個力量把水分開的時候.
這個就是值得我們敬拜的上帝.
所以無論是以色列人,希伯來人.
或者埃及,亞述,巴比倫等等.
他們都是這樣看待這個世界的.
大家都是用一個現象的世界.
用現象的方法來看待這個世界.
因為水是包著的.
誰能夠頂得住,我們就去敬拜他.
誰能夠征服水.
他就是那個贏家.
在巴比倫,古代巴比倫.
有一個神話叫做Eluma Elish.
中文可以翻譯為「創世史詩」.
即是說到世界的開頭是怎樣的.
就是巴比倫的假神叫做馬杜克.
他就是其中一個很年輕的神.

$^{121}$因為是多神論.
天庭裡面有很多的神神怪怪.
他其中一個是年輕的時候長大.
然後他成為一個英雄.
為甚麼呢.
因為有點像漫威的電影.
有一個很大的威脅.
威脅整個宇宙.
就是那個大水.
突然間那個大水的神發瘋.
想剿滅所有的神明.
這個馬杜克年輕力壯.
他突然興起.
他去到戰勝了海神或水神.
然後故事是寫得很好的.
好像天龍八部一樣.
他就將水神斬開.
好像一條魚一樣.
你們有宰過魚嗎.
吃魚的時候將牠分開一半.
然後有一半的魚和一半的水.
就頂在上面.
然後一半的水就頂在下面.
然後弄到有旱地,陸地出現.
然後他這樣頂著的時候.
他就說「我就是統治者」.
「我將大水戰勝征服之後」.
「我就是最大的權兵」.
「你就應該要譴責我」.
然後這個馬杜克.
去到安立了宇宙萬物的秩序.
然後就建立了他的聖城巴比倫.
這個就是所謂的.
Eluma Elish的神話.
你看看這部詩篇第46篇.
其實都是用一種.
當時人熟悉的一種語言.
來表達上帝的權威.
本身這個世界是有水包著.
有混沌來壓制著我們.

$^{161}$但是耶和華是掌權.
他將重水分開.
然後將大水來制服.
阻止重水去威脅我們.
這就顯明上帝有他的大能.
用這樣的方法宣講.
上帝是值得我們去信教的.
因為宇宙最大的能力是誰.
是耶和華.
是祂自己.
所以我們值得思考一下.
這位上帝既然是最終極權能的.
我們就需要敬畏和敬拜祂.
第四節到第七節.
剛才我提到甚麼.
說馬杜克,即是那個假神.
他戰勝了大水之後.
他做了甚麼.
創造世界,對吧.
然後他做了甚麼.
聖城,對吧.
敬拜他的聖城,巴比倫.
但在這裡說到甚麼.
就說到另一個城.
詩篇46篇.
一至三節說到上帝戰勝大水.
值得我們去信教.
四節到第七節就說到一個城市.
說到城市的時候.
我們心靈裡都有不同代表性的城市.
可能你的家鄉是某某城市.
可能你的家鄉或你很愛的城市.
或很多有代表性的市鎮.
就是你的安全感的地方,對吧.
無論是香港,台灣,台北.
或是北京,哪裡都好.
總之是你覺得有安全感.
有那種家的感覺.
對於希伯來人來說.
舊約神的子民.

$^{201}$那個城市是誰.
一定是耶路撒冷.
耶路撒冷又叫甚麼.
錫安.
Jerusalem就是Zion.
Zion就是Jerusalem.
這裡說到最重要的城市就是錫安城.
就是耶路撒冷.
所以有些人說詩篇第46篇.
可以叫作「Zion的歌」.
因為詩篇150首的詩歌裡面.
有好幾首都是去歌頌這個城市.
不是因為這個城市是世界之最.
不一定的.
而是因為這個城市.
就是上帝揀選他會住在那裡.
所以上帝的子民知道.
上帝在這裡跟他同在的時候.
就不需要懼怕.
4至第7節.
他說有一條河,這條河是分岔的.
在上帝的城歡喜.
這個城就是至高者上帝居住的城鎖.
就是耶路撒冷聖殿的地方.
有一條河走出來.
上帝在其中,城就不會動搖了.
雖然冰冷盛夏.
但是上帝必幫助這個城市.
耶和華與我們同在.
亞國的神是我們的避難所.
這裡說到耶和華戰勝大水之後.
他所建立的是甚麼.
是安穩,平安.
他所建立的城市不是巴比倫.
不是埃及的城市.
而是耶路撒冷,是錫安.
你留意到4至第7節.
那裡說的動詞,意象.
是跟開頭3節完全相反的.
頭3節的感覺應該是很澎湃.

$^{241}$有很多聲音,很多噪音.
很多東西煩擾和威脅你.
4至第7節突然靜下來.
而且每件事都反轉.
你記得頭3節說甚麼震動,澎湃等等.
4至第7節出現的時候.
意境就將那些東西全部反轉.
本來是混沌.
但現在我們看到是上帝給我們的平靜.
有一條河,河水流出來.
地理上,耶路撒冷不是有一條河.
不過地理上,耶路撒冷山腳確實有水泉.
叫做積分泉,是從那裡流出來的.
這不是說實際的地理.
而是說你心靈的地理.
在你的生命,心靈裡有這樣的地理.
有這樣的地圖.
有一條河流出來分岔.
這是聖殿,這是伊甸園.
有一條河流出來滋潤大地.
敬拜上帝來到聖殿的時候.
這些人,這些詩篇就想起.
上帝仍然與我們同在.
幾年前,我收到一個電郵.
那個電郵是說中東一些教會的狀況.
有甚麼代討等等.
那裡說到一些宣教士去到敘利亞.
在一個村莊裡探訪.
去到慰問當中的人.
因為當時,好幾年前.
你記得ISIS威脅很多村莊.
刪除了很多東西.
這些宣教士去了敘利亞的村莊時.
他們去到一個姐妹的家.
一個女孩的家.
發現這個女孩的村屋裡.
寫了很多字在牆上.
那些字全都是耶穌的名字.
用阿拉伯文寫等等.
全部把耶穌的名字貼在屋內.

$^{281}$在慰問之中和她探訪的時候.
這些宣教士就去問敘利亞的基督徒.
為何你會這樣做.
因為他說在ISIS威脅和做很多破壞時.
他在生命裡很害怕.
經常害怕被人捉了.
經常害怕他的家人有事.
他的複製方法是甚麼.
就是把耶穌的名字提醒自己.
提醒他有耶穌與他同在.
以致他的生命能夠心底安穩.
就是看到耶穌的名字.
在他的恐懼裡.
他能夠知道上帝與他同在.
詩篇繼續.
他說我們可以安穩在一個.
值得我們信靠的上帝裡面.
第八節到第十一節.
他說你們來看耶和華的作為.
看他使地怎樣荒涼.
他止息戰爭直到地極.
斷弓,斷槍,戰車焚燒在火中.
你們要休息,要知道我是上帝.
我必在列國中受尊崇.
在全地也受尊崇.
萬軍之耶和華與我們同在.
亞國的神是我們的避難所.
順帶一提,我想….
你記得舊約裡其中一個.
上帝截弓斷槍的故事嗎.
記得嗎?在希西加的時候.
希西加在主前大概701年的時候.
不是大概701,直接是701年.
701年,主前701BC的時候.
就用亞述帝國,亞述西征.
打到猶大地.
然後在耶路撒冷的京城.
佈設著千千萬萬的軍隊.
要攻打耶路撒冷的時候.
希西加做了甚麼.

$^{321}$他去聖殿.
他去尋求上帝.
上帝藉著衣衫跟他說.
「不用怕,我會幫你打仗」.
「明天亞述就會退兵了」.
結果那個故事的記載.
只是一次而已.
就是當晚有一位耶和華的使者.
就是耶和華的軍兵.
一位單純的.
就出去圍剿耶路撒冷的亞述兵營裡.
殺了18萬多的士兵.
然後亞述王就退兵.
我想這個詩篇.
大概應該回顧這個歷史的故事.
當你恐懼的時候.
原來耶和華是可以為你打仗.
你不需要被恐懼制服或支配.
我不知道你恐懼的時候.
或者你所謂壓力大的時候.
你會是怎樣的.
我是一個很急的人.
當我再急,再有威脅的時候.
我就會….
你覺得我是否很理性的人.
你們覺得是嗎.
(是).
是,我99\%的時間都是很理性的.
你們在笑甚麼.
我覺得應該大部分時間都是很理性的.
誰知道剛好是星期五.
剛好星期五就發生了一件事.
就是完結了NBTI研討會.
那些人就散了.
我接到電話.
他們說媽媽的車輛….
就是那輛車的輪胎.
突然在高速公路沒有氣.
他們打來求救.
我接了電話,然後就飛奔去.

$^{361}$除了飛奔去之外,我就打給我兒子.
我就告訴我兒子.
「你馬上要去」.
他就問長問短.
因為他跟朋友在高興.
我說「不…你必須去」.
「你必須去,為甚麼」.
他就問「為甚麼?發生甚麼事」.
我說「我不知道」.
「總之你一定要去」.
我就知道當時我已經有恐慌的狀態.
即是在壓力下的狀態.
我已經在電話裡罵了他10句.
終於我用我的方法.
就是用恐懼來威脅他.
總之我受到威脅,我也想威脅他.
我威脅了他一段時間,他說「我現在去」.
「我們兩輛車要去」.
因為大風大雪,每隔五晚.
很煩躁.
我走出去,我就飛車去.
飛了兩個街口,我說「電話在哪裡」.
我剛才明明聽完電話.
翻來覆去我的褸袋.
翻來覆去也翻不到.
然後我說「糟糕,一定是掉了」.
我馬上調頭回來停車場.
我到處在地上尋找.
但電視節目很好.
「你忘記電話了嗎?你沒事吧」.
其實當時是非常慌亂的.
不知為何我踏入了瀕臨的狀態.
我找來找去也找不到.
褸袋也找不到,沒理由的.
剛才談完電話,車也找不到等等.
我走進UV,遇上一個女同性戀者.
那個女同性戀者不知為何.
平時不太鎮定.
平時也很鎮定,我很害怕.
她特別鎮定.

$^{401}$因為跟我相反的時候.
她顯得非常鎮定.
她竟然跟我說「你冷靜點」.
她用「Find me」的應用程式.
指向我的電話,於是找到了.
「對嗎?我登入了」.
「我登入了她的電話,找到我了」.
電話響起了.
「你猜她的電話在哪裡」.
原來她的電話在我的袋子裡.
我的天.
我開車離開時.
「這次真的被阿芝勝出了」.
讓她看見我如此慌亂的狀態.
我們確實有很多恐懼影響我們的生命.
無論是世界上發生的事.
政權,政治發生的事.
國與國之間.
或是可能在我們的生命中.
我們也會有些時候感到失控.
我離開時感激著.
感激著很多人千里迢迢來遇到我.
對.
當你恐懼的時候.
你確實會完全失去所有控制.
失去理性.
但這些情況通常都不在我身上出現.
很少出現.
無論如何.
我預備的時候.
我記起好幾年前.
有一位年輕人突然打電話給我.
他哭著地打電話給我.
他在電話中跟我解釋.
他失去了數千元.
失去了數千元,原來是怎樣的.
原來他收到一個電話.
誤稱是IRS打給他.
稅局打給他.
稅局在電話中跟他說.

$^{441}$「你過去幾年都沒有報稅」.
「所以你欠了四,五千元的稅項」.
「你怎樣去處理稅項呢」.
「你就去蘋果商場買禮物卡」.
蘋果商場的禮物卡.
應該不會有五千元的禮物卡吧.
你是幾百元買的.
他就去了蘋果商場.
「大哥,IRS打給你」.
IRS是全世界最大的黑社會.
你不怕嗎?雙腿發抖吧.
他確實想起過去幾年都是欠缺稅項.
所以他去了蘋果商場.
然後他買了300元,200元,300元.
200元,儲起來就有四,五千元.
然後拿回家.
買了禮物卡,你只能刮.
逐張刮.
然後報告給所有人.
說我的禮物卡的號碼是甚麼.
結果他逐張刮給對方.
於是他就報了四,五千元.
你說我很瘋狂嗎.
很瘋狂吧.
我回想起來,如果我恐慌的時候.
如果Chimmy問我拿錢.
可能我也會給他.
無論如何.
當你恐慌,被恐懼控制的時候.
你就會做很多沒想過的事.
可能他早點致電給我.
這樣就沒事了,對嗎.
對,是嗎.
問清楚,對吧.
或者問一些經紀或熟悉的人.
在恐懼中,我們確實有很多錯誤反應.
但詩篇第46篇告訴我們.
即使兵臨城下,也要謹記.
「夜禍是與我同在」,對吧.
上星期有一宗新聞.

$^{481}$很多人都留意的,對吧.
Alex Honnold.
Alex Honnold赤手空拳爬高山超壁.
他準備了很久去爬台北101.
你們有去過嗎.
有去過101,對吧.
你不是爬上去的.
他的赤手空拳不需要安全帶.
不需要任何輔助.
就這樣爬上去台北101.
很多時候我們看見這些影片時.
我們會有點心驚動魄.
有誰有看.
雖然我們覺得很害怕,但我們也想看.
很奇怪吧.
他只用了90分鐘.
就爬到上面的避雷針尖頂.
你會發現有很多人報道或回應時.
我覺得很有趣.
其中一個回應我覺得很有意思.
就是他不是不怕死.
而是他準備到一個地步不會死.
他不是不怕死.
你也怕死,我也怕死,對吧.
但他準備到一個地步不會死.
或將死的風險減到最低.
另一個回應我覺得很有意思.
他說我們經常想.
我們要走出我們的舒服圈.
我們的舒適圈.
過去一年,我們說要跨越.
我們要走出我們的舒適圈.
這個人Alex就剛剛相反.
他說不是走出我們的舒服圈.
而是擴大我們的舒服圈.
以前我們不敢舒服的東西.
我們願意去擴大它.
如何擴大.
你看見的只是他用90分鐘.
爬到台北101.

$^{521}$但你沒看見的就是他爬之前.
他有好幾千次重複練習同一個動作.
重複練習每一個手,腳的位置.
每一個觸摸的地方.
他會重複每一次.
我們會覺得很瘋狂,對吧.
對嗎?會覺得很瘋狂.
如果你的兒子是這樣.
你真的不可活.
如果你的兒子每天都做這些事.
不知道他何時死.
我們覺得很瘋狂,其實他一點也不瘋狂.
他是一個很有理性的人.
因為同一條路線.
他可以練習40次以上.
是重複日常.
每一個腳點,手點,觸點.
他都寫在他的很詳細的筆記中.
一直在心裡複雜,練習.
這讓我們提醒.
原來恐懼是可以去管理的.
起碼在我們的心靈裡,我們可以開始.
但《詩篇》,《石榴篇》.
它沒有假裝過沒有大水.
沒有崩裂的時候.
亦沒有假裝.
假裝說烈國沒有影響我們,威脅我們.
世界不會有動盪.
但《詩篇》,《石榴篇》.
在現實裡.
背後有一個最終極的我們可以倚靠.
諸事力放在背後.
這位耶和華.
這位可以戰勝大水.
與我們同在.
當我們心不安穩的時候.
萬君子耶和華與我們同在.
生活裡可能有不同的挑戰.
一種樂觀是家庭,經濟等等.
令你心也會融化.

$^{561}$但起碼我們心態上可以有安穩.
你們要休息,要知道我是神.
或者可能有一點.
邁尾的時候的提醒.
我們不妨練習一下.
焦慮來到的時候.
我們跑去的是上帝面前.
不是用我們那種焦慮,恐慌的模式.
來解決問題.
跑去上帝面前.
主要是這個時候我需要你.
並且這些事情可以每天鍛鍊.
嘗試去停頓,去安靜.
這也是我一生的功課.
因為大家知道我做事很急.
很難停下來.
每天停下來.
嘗試去感受一下,呼吸一下.
知道控制權不在我身上.
而是在上帝身上.
並且將這篇詩篇活出來.
就是說我們每天.
其實可以在我們生命裡.
去演繹這個副歌.
「夜禍萬君子耶和華與我們同在」.
當我們跟別人有爭執.
或者當我們有很多繃緊的時候.
我們知道上帝仍然是掌權.
於是我們應該….
有這樣的信念的時候.
其實可以影響我們對話.
對話可以平靜.
對人可以溫柔一點.
最後不是我們掌權.
是上帝掌權.
可以同心協力祈禱嗎.
因為你是我們的避難所.
是我們的力量.
可以挽難裡面隨時的幫助.
當我們心慌,意亂,迷失的時候.

$^{601}$求你帶我們回到你的船上.
讓我們知道你是那位.
穩如泰山的上帝.
教導我們安靜.
安靜,知道你是神.
教導我們在生命裡.
如何用安穩,信靠,平安.
來生活.
在這個時候也邀請大英姐妹.
你最近可能遇見不同的境況.
不同的情況.
你很需要你的位主.
需要位主給你那種安穩.
可以花點時間來到上帝面前.
去休息,去安靜.
知道耶和華才是神.
因為縱然世界,生命裡.
可能有很多動搖.
但你卻永遠是神.
不僅你的能力沒有改變.
你永遠是無法控制的.
在我們生命裡的每一個細節.
你都是無法控制的.
並且我們知道你對我們的愛.
是堅定不移的.
在我們今天來到你自己的聖餐桌前.
在世世代代所行的聖餐禮當中.
我們拿起這個餅,這個杯.
知道你的愛對我們的救恩.
是沒有改變過的.
是鐵一般的事實.
就是你自己的愛子耶穌基督.
在十字架上為我們死.
在墳墓裡為我們復活.
給我們一個新的生命.
讓我們能夠享受你對我們的同在.
最要求你預備我們的心靈.
來領受這個聖餐.
知道無論世界發生任何事.
我們多麼軟弱,多麼破碎.

$^{641}$你依然接納我們.
我們感恩讚美,祈禱奉主明教.
\newpage



\section{士師記}
\label{sec:espYGx8NsQM}
\textbf{The Trustworthy God (Judges 4) - Chee Yie Jong  信實的上帝 (士師記 4) - 楊志溢傳道}
\newline
\newline
連結: \href{https://youtube.com/watch?v=espYGx8NsQM}{\texttt{https://youtube.com/watch?v=espYGx8NsQM}} ~~~~ 語音日期: 2026-02-16
\newline
\newline
\hyperref[sec:QVz_a2O_21A]{\small{< < < PREV SERMON < < <}}
~
\hyperref[sec:index_chronic]{\small{[返順時目]}}
~
\hyperref[sec:index_scriptual]{\small{[返順卷目]}}
~
\hyperref[sec:Ulr7EZgFyUo]{\small{> > > NEXT SERMON > > >}}
\newline
\newline
$^{1}$Hello, brothers and sisters..
[NON-ENGLISH SPEECH].
So one things that I like to do in my free time.
is I like to go climbing..
So naturally, my social media algorithms,.
they would always suggest to me all sorts.
of climbing-related content..
And a few months ago, I actually bumped.
into this very heartwarming YouTube shorts.
posted by-- she's a professional climber and also.
a former Olympian from the UK..
And in this YouTube shorts, she brought her daughter.
to climb outdoors..
And you can see that her daughter,.
she was very determined to climb this climb all by herself.
without the mom's help..
[NON-ENGLISH SPEECH].
So a few moves into the climb, however,.
she started to panic a little bit when things.
got a little more challenging..
And her emotions began to come out..
You can see the dilemma here..
On one hand, she wanted to do the climb all by herself..
So she didn't want her mom to help her at all..
But on the other hand, she was getting a little scared..
And her natural instinct just came up.
to ask her mom for help..
[NON-ENGLISH SPEECH].
So I think we can all resonate with the feelings of the daughter.
here when we are faced with some sort of adversities.
and challenge, right?.
Maybe you're feeling stuck in your life..
And you're just having a hard time figuring out.
what is the next steps..
Or maybe you're in some sort of crisis moment..
And you just can't seem to figure a way out..
Or maybe some of us, you are confronted by your own fears..
Will my marriage or my relationship last?.
Or what if I lose my job?.
Or what if my friends or people that I love.

$^{41}$actually abandon or leave me?.
And in these moments of doubt, how do you usually respond?.
What would you do in the face of adversity?.
[NON-ENGLISH SPEECH].
So this morning, we are going to look.
at the third major judge in the book of "Judges," that.
is, Deborah..
The cycle of Deborah is actually quite unique.
because it spans across two chapters..
And it consists of a story and a song..
So "Judges 4," that is the one that we'll be looking in today,.
is in the form of a story..
That is the narrative portion of this cycle..
While "Judges 5" is in the form of a poem..
That is the song portion of this cycle..
And so we'll actually spend two Sundays.
in this cycle of Deborah alone..
But again, for this morning, we're.
going to focus on the narrative portion of the cycle..
That is "Judges 4.".
And again, we will explore these questions together..
What should we do in the face of adversities and doubt?.
[NON-ENGLISH SPEECH].
And so today, instead of having the passage up on a screen.
and trying to do something different,.
so I want to encourage you to open your own Bible,.
physically or digitally, and follow along with our passage..
Because there will be some clues that's earlier in our passage.
that will only get fully reviewed at the end..
And I think it is actually more effective for us.
to learn from God's word by having the whole passage.
in our hand, rather than seeing it section by section.
on the screen, where I can only show a limited amount of verses.
every time..
And if you don't have a Bible or you.
don't have a way to access it, I also.
printed out a few physical copies.
that you can come up to the table and grab..
And so to apologize, I forgot to print out the Mandarin version..
So there's an English copy of "Judges 4".

$^{81}$up there if you want to come up and grab them..
[NON-ENGLISH SPEECH].
OK, everyone ready?.
We're going to start with "Judges 4" verse 1..
"So again, the Israelites did evil in the eyes of the Lord,.
now that Ehud was dead..
So the Lord sold them into the hands of Jabin, king of Canaan,.
who reigned in Hazor..
Sisera, the commander of his army,.
was based in Harosheth Hagoyim..
Because he and 900 chariots fitted with iron.
and had cruelly oppressed the Israelites for 20 years,.
they cried out to the Lord for help.".
[NON-ENGLISH SPEECH].
And so right here at the first verse,.
we're immediately being connected to our narrative.
back to the previous judge, that is Ehud..
And if you recall from my previous sermon,.
after Ehud had defeated Eglon, king of Moab,.
the land of Israel, they enjoyed a period of peace.
for 80 years..
But right after these 80 years, after Ehud had died,.
this cycle repeats itself again..
And we see that Israel does evil in the eyes of the Lord..
[NON-ENGLISH SPEECH].
And this time, the Lord actually used Jabin, king of Canaan,.
and also his commander, Sisera, to oppress Israel..
[NON-ENGLISH SPEECH].
One advantage of Sisera's army here.
was that they had 900 iron chariots..
And so iron chariots was kind of like the top-of-the-line.
technology in warfare during that time..
When a war is fought over on flat land or coastal plains,.
it actually provides a huge advantage for this iron chariots.
army compared to just regular foot soldiers..
[NON-ENGLISH SPEECH].
And as a result of their oppressions from Sisera,.
Israel once again cried out to the Lord for help..
[NON-ENGLISH SPEECH].
And now let's look together in verse 4 and 5..

$^{121}$[NON-ENGLISH SPEECH].
Now Deborah, a prophet, the wife of Lapidoth,.
was leading or judging Israel at that time..
She held court under the palm of Deborah.
between Ramah and Bethel in the hill country of Ephraim..
And the Israelites went up to her to have their disputes.
decided..
[NON-ENGLISH SPEECH].
And so right here, we're being introduced to Deborah..
Deborah was a prophetess by profession,.
or a more direct translation is that she was a woman prophet..
I'm not going to spoil the rest of the story for you now,.
but women will actually play a significant role in our story.
today..
[NON-ENGLISH SPEECH].
And besides being a prophet or prophetess,.
we also learned that she was judging Israel at that time..
And what this means is that Deborah, in some way,.
she was Israel's decision maker..
And we see that people from all over Israel,.
they will come to her to have their dispute decided..
[NON-ENGLISH SPEECH].
And not only that, we will learn that she.
was also a representative and spoke person.
on behalf of the Lord in the following verses..
[NON-ENGLISH SPEECH].
And so let's continue with verse 6..
"She sent for Barak, son of Abinoam,.
from Kadesh in Naphtali, and say to him,.
the Lord, the God of Israel, commands you,.
go, take with you 10,000 men from Naphtali and Zebelin,.
and lead them up to Mount Tabor..
I will lead Sisera, the commander of Jabin's army,.
with his chariots and his troops to the Kishon River.
and give him into your hands.".
[NON-ENGLISH SPEECH].
And so even though Deborah, she was judging Israel.
at that time, she did not command the army of Israel..
And so she sent for Barak..
There are not many details about Barak here,.

$^{161}$but from what we can gather, he is from the tribe of Naphtali,.
and he was also chosen by God to be Israel's deliverer..
[NON-ENGLISH SPEECH].
And according to Deborah, Barak's role.
was pretty straightforward..
That is, to go and to gather the troops to Mount Tabor.
and to call up the warriors from the tribes of Naphtali.
and Zebelin..
And then the Lord himself, he will.
give Sisera and his army into the hands of Barak..
[NON-ENGLISH SPEECH].
And so Deborah's message here is very clear..
There is a command from God, and there's a promise of victory..
And so the command is for Barak to gather his army.
and lead them into the battlefield..
And then the promise is that the Lord.
will grant Israel their victory over Sisera and his army,.
even though Sisera has military advantage.
with his iron chariots..
[NON-ENGLISH SPEECH].
But this is Barak's response here in verse 8..
"Barak said to her, if you go with me, I will go..
But if you don't go with me, then I won't go.".
[NON-ENGLISH SPEECH].
So even though with the victory guaranteed,.
Barak, he still feel a little uncertain..
And so he started to negotiate a little bit,.
giving Deborah his conditions or his terms for obedience..
If you, Deborah, go with me, then I.
will obey your command and go..
But if you don't go with me, then I won't go..
[NON-ENGLISH SPEECH].
Barak's hesitancy here is not without reason..
From a personal level, we learn that he.
is from the tribe of Naphtali..
And when we go back to the first chapter in Judges,.
we see that the tribes of Naphtali.
actually failed to drive out the Canaanites living in their land..
And so from Barak's perspective, this failure from his tribe.
might have added some doubts into his mind.

$^{201}$when he received this divine message from Deborah..
And so instead of obeying, he shows.
some degree of disobedience..
[NON-ENGLISH SPEECH].
And if we take a step back and zoom out from this narrative.
back to the big picture of Judges,.
we know that the book of Judges in some way.
is an evaluation of whether or not.
Israel, as God's holy nation, will remain faithful.
and obedient to the Lord..
And so as a nation, Barak's response.
here actually gave us some spiritual insights.
into the life of Israel's leaders,.
especially after the death of Moses and his successor,.
Joshua..
[NON-ENGLISH SPEECH].
So now I want us to ask ourselves this question..
What will we do in the face of adversities and doubts?.
And how will we respond when we're.
being asked to do something that we are not comfortable with?.
[NON-ENGLISH SPEECH].
I think we cannot relate to Barak in this situation..
A lot of time, we want to trust God..
We really want to trust God..
But we want to trust him in our own terms.
or in ways that we are comfortable.
and that is convenient for us..
And so our response may look like this..
Lord, I will start giving if you give me a raise at my work..
Or Lord, I will start loving my neighbor next door.
if they can just stop playing music so loud every night..
Or Lord, I will start serving at church on Sunday.
if my boss would just stop scheduling me to work on Saturday..
[NON-ENGLISH SPEECH].
And do you see the issues here?.
Do you think God will be pleased with this kind of obedience.
and incomplete faith?.
In fact, it is well within reason.
for God to take offense with this kind of incomplete faith.
and half-hearted obedience, right?.

$^{241}$Just imagine this..
If you ask your son to take the trash out after dinner,.
and this is how he responds to you,.
he says, yeah, sure, but only if you give me more pocket money.
or you increase my weekly allowance..
If you are the parents here, you will probably.
feel quite disappointed and offended, right?.
[NON-ENGLISH SPEECH].
However, we know that our God is a loving and patient father..
And so even though Barak's faith is incomplete.
at this point of the story, God doesn't just abandon him,.
but he will patiently prove his own trustworthiness.
and faithfulness to Barak in order.
to rehabilitate Barak's faith..
[NON-ENGLISH SPEECH].
So let's come back to our passage.
here and continue in verse 9..
So verse 9, "Certainly I will go with you, saith Deborah..
But because of the course or the way that you are taking,.
the honor will not be yours..
For the Lord will deliver Sisera into the hands of a woman.".
So Deborah went with Barak to Kadesh..
There, Barak summoned Zebulun and Naphtali,.
and 10,000 men went up under his command..
Deborah also went up with him..
[NON-ENGLISH SPEECH].
So there are a few ways for us to understand verse 9 here..
The most straightforward reading might lead us to think that.
because Barak shows some signs of doubt,.
God is kind of like changing his mind.
to teach Barak a lesson by taking away the honor from him..
[NON-ENGLISH SPEECH].
However, there's actually a better way.
to understand this, that this is actually not a punishment.
for Barak, rather it is a clarification.
of God's initial promise..
God's promise to him still stands,.
that victory will still belong to Barak and Israel..
However, Deborah added a crucial detail here in the plot..
Because of the course or because of the way in verse 9.

$^{281}$that Barak is taking, the Lord will deliver Sisera.
into the hands of a woman, not the military commander.
or his army..
[NON-ENGLISH SPEECH].
And despite Barak's incomplete faith,.
God chose to show his faithfulness by actually.
accommodating to Barak's term, that Deborah will indeed.
go with him into the battle..
[NON-ENGLISH SPEECH].
And then verse 11..
"Now Herber, the Canaanite, had left the other Canaanites,.
the descendants of Hobah, Moses' brother or father-in-law,.
and pitched his tent by the great tree.
in Zananim near Kadesh.".
[NON-ENGLISH SPEECH].
So this verse right here, verse 11,.
we actually find a break in the main storyline..
And we're being introduced to this family.
called Herber, the Canaanite..
So even though this verse probably doesn't progress.
the plot at this point, but this Canaanite family,.
and especially the tent that they live in,.
will play an important part later in our story..
So just keep that in mind..
[NON-ENGLISH SPEECH].
So let's continue with our story in verse 12 to 15..
"When they told Sisera that Barak, son of Abinoam,.
had gone up to Mount Tabor, Sisera.
summoned from Harosheth-hagoim to the Kishon River.
all his men and his 900 iron chariots..
Then Deborah said to Barak, go..
This is the day the Lord has given Sisera into your hands..
Has not the Lord gone ahead of you?.
And so Barak went down Mount Tabor.
with 10,000 men following him..
And at Barak's advance, the Lord routed Sisera.
and all his chariots and army by the sword..
And Sisera got down from his chariots and fled on foot.".
[NON-ENGLISH SPEECH].
And so here in verse 14, Deborah once again.

$^{321}$reminds Barak of the promised victory from the Lord..
And so not only did she go physically with Barak,.
in order to build up his faith, the Lord.
also used Deborah to remind and reassure Barak.
at opportune times that the Lord had already.
given Sisera into his hands..
[NON-ENGLISH SPEECH].
And as the Lord promised, we see that he routed the Canaanites.
army, and Sisera actually had to abandon his chariots.
and flee on foot..
[NON-ENGLISH SPEECH].
And then in verse 16, "Barak pursued the chariots.
and his army as far as Harosheth Hagoyim,.
and all Sisera's troops fell by the sword..
Not a man was left.".
[NON-ENGLISH SPEECH].
And so if you remember at the beginning of our passage,.
Harosheth Hagoyim is where Sisera and his army was based..
So Sisera's army, they retreated all the way back to their base..
And if you look at the map, Harosheth Hagoyim.
is the circle on the left..
And it is west of Mount Tabor..
And so we know that Barak, in order to pursue them,.
he went west..
So keep that in mind again..
[NON-ENGLISH SPEECH].
And verse 17, "Sisera, meanwhile, fled on foot to the tent of Jael,.
the wife of Herber the Canaan, because there.
was an alliance between Jabin, king of Hazor,.
and the family of Herber the Canaan.".
[NON-ENGLISH SPEECH].
So this is where the mystery, the information from verse 11.
started to make a little sense..
Here, we're being introduced to another character, Jael,.
who is the wife of Herber the Canaan from verse 11..
[NON-ENGLISH SPEECH].
And why did Sisera actually flee to the tent of Jael?.
We learned that her husband has some sort of political ties.
with Jabin, king of Canaan..
And so Sisera, in his mind, he probably.

$^{361}$thought that he could get some help from this ally of theirs..
[NON-ENGLISH SPEECH].
And if you look at the map again, the tent of Jael.
is situated in Za'ananim, which is east of Mount Tabor..
That is the circle on the right..
[NON-ENGLISH SPEECH].
Yes, so we know Sisera headed east..
So at this point, we started to get more clarity.
on what Deborah means in verse 9..
When she told Barak that because of the course.
or because of the way that he is taking--.
that is, he is heading west--.
the owner would not be his, because Barak.
is literally heading in the wrong direction.
in pursuit of Sisera..
[NON-ENGLISH SPEECH].
And as for Jael, we don't know much.
about her own political allegiance..
However, the important thing here.
is to know that she is not part of Israel..
And physically, she also lives at least 10 miles away.
from the battlefield..
And so this reinforced the narrative.
that she is the least likely person, the least likely.
character to participate in this battle..
And yet, her appearance will actually.
set the stage for the most unexpected way.
the Lord would deliver Israel..
[NON-ENGLISH SPEECH].
So let's look at verse 18 to 21, "So Jael went out to meet.
Sisera and say to him, come, my lord, come right in..
Don't be afraid..
So he entered her tent, and she covered him with a blanket..
I am thirsty, he said..
Please give me some water..
And so she opened a skin of milk, gave him a drink,.
and covered him up..
Stand in the doorway of the tent, he told her..
If someone comes by and asks you, is anyone in there,.
say no..

$^{401}$But Jael, Herbert's wife, picked up a tent peg and hammer.
and went quietly to him while he lay fast asleep, exhausted..
She drove the tent peg through his temple into the ground,.
and he died.".
[NON-ENGLISH SPEECH].
And so now, the promise is fulfilled,.
and there is no ambiguity left..
Just as the Lord promised through Deborah,.
he delivers Sisera into the hands of a woman..
This woman is not Deborah herself, but it is Jael,.
who is the least likely person in this story..
[NON-ENGLISH SPEECH].
And then verse 22..
"Just then, Barak came by in pursuit of Sisera,.
and Jael went out to meet him..
'Come,' she said, 'I will show you.
the man you are looking for.'.
So he went in with her, and there.
lay Sisera with the tent peg through his temple, dead.".
[NON-ENGLISH SPEECH].
And so through Jael's invitation into her tent,.
Barak himself got to witness how the Lord fulfilled his promise..
And I can imagine that at this moment,.
seeing what had happened in front of him,.
the commander of the enemy, who had cruelly.
oppressed Israel for the past 20 years,.
now laying dead on the ground..
Barak probably gained a greater faith and confidence.
in the Lord at this moment..
[NON-ENGLISH SPEECH].
And even though Barak shows some signs of spiritual weakness.
and trying to negotiate with God's command,.
God, on the other hand, continued to show his patience.
and his faithfulness to Barak..
And he did so by using both women in the story, Deborah.
as the messenger and Jael as the executor of the promise,.
to show Barak that his words never fail.
and he can be fully trusted..
[NON-ENGLISH SPEECH].
So let's finish up..

$^{441}$Verse 23, "On that day, God subdued Jabin, king of Canaan,.
before the Israelites..
And the hand of the Israelites pressed harder and harder.
against Jabin, king of Canaan, until they destroyed him.".
[NON-ENGLISH SPEECH].
And so here's the main idea that I want us to take away with us.
today, that we can lean on the Lord,.
even in the midst of adversities or doubts,.
because he is trustworthy and his promise never fail..
[NON-ENGLISH SPEECH].
And this is the invitation that lies before us..
Would you trust him?.
Would you hold on to God's promise,.
even when life seems out of control?.
And would you choose the way of Jesus,.
even when it is inconvenient or uncomfortable?.
[NON-ENGLISH SPEECH].
And so I want to end my sermon by showing the full YouTube.
shorts of the mom bringing her daughter to climb..
And as you watch later, I want to draw your attention.
to these two very different perspectives,.
from the daughter's perspective and the mom's perspective..
And so from the daughter's perspective,.
when she is at the brink of a breakdown,.
her doubts and fears affects her trust in her mom,.
as well as her mom's ability to keep her safe..
But from the mom's perspective, her daughter.
is actually never in danger, because she.
is right there with her..
And instead of letting the mutual trust of her daughter.
to get into her, the mom patiently guides her daughter.
step by step and lovingly encourages her.
until she gets to the top..
[NON-ENGLISH SPEECH].
So let's watch it together..
Man, nice work, man, oh, great work..
[NON-ENGLISH SPEECH].
Nice, bud..
Great work, buddy..
[NON-ENGLISH SPEECH].

$^{481}$Oh, yeah..
[NON-ENGLISH SPEECH].
To your left, mate..
There you go..
Well done, buddy..
[NON-ENGLISH SPEECH].
So it is the same when it comes to our life..
When we see things from our own limited perspective,.
we can easily lose faith when we are confronted.
by our own doubts and fear..
But if we can learn to lean on the Lord.
and look at things from His perspective,.
we would know that we can fully trust Him,.
because He is always right there with us,.
and He is always trustworthy..
[NON-ENGLISH SPEECH].
So, Church, let's continue to grow our faith in Him..
Let's pray..
[NON-ENGLISH SPEECH].
Lord, thank you for being a faithful and trustworthy God..
We know that your promise never fails,.
and we rest in the fact that nothing, nothing can ever.
stop you from accomplishing your word..
So, Lord, increase our faith in you..
Remind us of your goodness and your faithfulness in our life..
Strengthen us in areas of our life.
where our faith is incomplete..
Heal us of our deepest fear and doubt.
so that we can come before you with greater trust..
[NON-ENGLISH SPEECH].
And in Jesus' name, we pray, amen..
[NON-ENGLISH SPEECH].
Amen..
OK, thank you all for joining us this morning..
And as Kelly has mentioned, Pastor Mark.
is preaching at Wheaton today, so he won't be with us..
But I would still encourage you to stay, hang out a little bit,.
and connect with one another..
So go in peace..
You're loved..

$^{521}$you.
\newpage



\section{約翰福音 1:14-18}
\label{sec:Ulr7EZgFyUo}
\textbf{道成肉身 (約翰福音1\_14-18) - 梁紫健傳道}
\newline
\newline
連結: \href{https://youtube.com/watch?v=Ulr7EZgFyUo}{\texttt{https://youtube.com/watch?v=Ulr7EZgFyUo}} ~~~~ 語音日期: 2026-02-16
\newline
\newline
\hyperref[sec:espYGx8NsQM]{\small{< < < PREV SERMON < < <}}
~
\hyperref[sec:index_chronic]{\small{[返順時目]}}
~
\hyperref[sec:index_scriptual]{\small{[返順卷目]}}
~
\hyperref[sec:xg_qWk2v3jg]{\small{> > > NEXT SERMON > > >}}
\newline
\newline
約翰福音 1:14-18
\newline
\begin{longtable}{cl}
\hline
\hline
章節 & 經文 (和合本修訂版)\\
\hline
1:14 & \begin{tabularx}{0.7\textwidth}{X} 道成了肉身，住在我們中間，充充滿滿地有恩典有真理，我們也見過他的榮光，正是父獨一兒子的榮光。 \end{tabularx} \\ \\ \relax
1:15 & \begin{tabularx}{0.7\textwidth}{X} 約翰為他作見證，喊著說：「這就是我曾說：『那在我以後來的先於我，因為在我以前，他已經存在。』」 \end{tabularx} \\ \\ \relax
1:16 & \begin{tabularx}{0.7\textwidth}{X} 從他的豐富裡，我們都領受了恩典，而且恩上加恩。 \end{tabularx} \\ \\ \relax
1:17 & \begin{tabularx}{0.7\textwidth}{X} 律法是藉著摩西頒佈的；恩典和真理卻是由耶穌基督來的。 \end{tabularx} \\ \\ \relax
1:18 & \begin{tabularx}{0.7\textwidth}{X} 從來沒有人見過神，只有在父懷裡獨一的兒子將他表明出來。 \end{tabularx} \\ \\
[1ex]
\hline
\hline
\end{longtable}
$^{1}$大家早安,冷了一陣子終於暖了一點,我已經很期待夏天的來臨,今年真的很冷..
先給大家看一張圖,在歷史上有一個很出名的故事,就是人慶和教會的爭論,大家認不認得他是誰呢?.
有沒有認得?對,嘉莉略,聰明,拉奧神真聰明..
他就是嘉莉略一個很出名的數學家,天文學家,當年嘉莉略用望遠鏡觀察星空,發現太陽不是圍著地球轉,相反是地球圍著太陽轉..
這個知識對我們來說現在很正常,小學也讀過,但是當時很多實學家堅信聖經的自然變異事,.
聖經有幾句是說太陽又升起又落下,是來自詩篇十九章和傳道書一章,所以一定是太陽移動而不是地球移動..
結果他們就因為自己對於信仰的詮釋的堅持,寧願否認我們看到的事實..
當年這一班自以為為了真理把關的宗教領袖就在羅馬的審訊廳裡面一次又一次地盤問,嘉莉略,你是不是承認自己錯?.
你是不是承認聖經是不可以給你這些新的理論去挑戰的?.
最後這個年老的科學家嘉莉略被迫跪在地上,他讀他的懺悔文,公開收回自己的言論,.
太陽是圍著地球轉的,不是地球圍著太陽轉..
只因他看到的事實和教廷一直堅持的詮釋不一樣..
幾百年之後,科學證明了嘉莉略說的是真的,教會就要承認當時的錯誤,向死了幾百年的嘉莉略道歉..
這個故事提醒我們,有時我們以為我們一直堅持的信念就等於上帝的真理..
可能我們出於虔誠,出於經驗,出於傳統,覺得自己沒有錯..
但如果我們不肯學習,不肯聆聽,我們就可能會被真理推翻,成為了當年教廷的人..
所以今天我們要問自己,我們堅持的東西是真理本身,還是我們覺得是好像真理的東西?.
兩樣東西是很不同的,是真理本身,還是我覺得那個是真理?.
你分得到那兩樣東西是甚麼嗎?.
如果大家準備好,我們就會進入我們今天的講道..
聽到這邊的講道,我們會知道三樣東西..
當時的猶太人為甚麼會拒絕耶穌?.
十誡為甚麼成為了他們最大的信耶穌的阻礙?.
而這樣會否同樣地成為了我們的阻礙呢?.
如果準備好,我們就一起來了..
上一次我們講了約翰福音的頭頭那幾節,.
今天我們再繼續講約翰福音的一章,下面中間那段經文..
一章十四至十八節..
道成你肉身,住在我們中間,充充滿滿地有恩典有真理..
我們也見過他的榮光,正是負獨一而至的榮光..
約翰為他作見證喊著說:.
「這就是我曾說,那在我以後來的先於我,因為在我以前他已經存在.」.
從他的豐富裡,我們都領受了恩典,而且欣賞加恩..
律法是藉著摩西般步的,恩典和真理卻是由耶穌基督來的..
從來沒有人見過上帝只有在父女,父懷女獨一而至,將他顯明出來..
在這一段的上文,大家如果看整個約翰福音的第一章,.
上文其實會介紹另外一個約翰,是施洗約翰,.
說他是為真光而作見證,為耶穌鋪平道路..
然後到十五字,我們正在看的這段經文,.
說了一句大家都不太明白的東西,.

$^{41}$「在我後來的人先於我,因為在我以前他已經存在.」.
這個世間你怎樣理解在我之前,在我之後,在我之後,在我之前的你?.
第十十字為何要提到摩西般步的律法,而恩典和真理又是由耶穌來的呢?.
從來沒有人見過上帝,那耶穌又怎樣將上帝顯明出來呢?.
不知道大家看約翰福音,會不會通常都有一個印象,就是很長氣..
那些東西說完又說,說完又說,說了一些你明白的東西,但又不太明白..
大家這些每個字都明白吧,放在一起好像很奇怪..
大家如果有這樣的感覺是非常正常的,.
因為約翰福音所針對的思潮是接近在二千年前,.
我們已經跟他隔了很多代了,很有代溝的..
我們現在跟下一代隔三年都覺得有代溝,何況我們隔了二千多年,我們隔了很多代了..
以致有些東西我們不明白,所以今天我們要搞清楚他在回應什麼..
我們要講一下他在回應什麼,其實他在回應三個很不同的思潮..
第一個就是諾斯底主義,就是大家抬右這邊..
他在回應的是很強烈的二元論,就是至高的真神屬於靈界,而物質的世界是屬於次一級的..
是由這個邪惡的創造主所做的..
這個思潮就會覺得如果上帝道成肉身,成為人的話,.
就會令上帝承認自己就是那個邪惡的創造主..
所以約翰福音對這個思潮有很多的反駁..
第二,約翰福音還會針對這個異教世界觀,特別是希臘的神明..
如果大家有聽過希臘的神話,很多不同的神明,太陽神,阿典娜,.
時不時他們會落到人間,去扮人去欺騙人,.
以致那些神明可以滿足到自己的私慾..
這個和約翰福音所講的道成肉身的本質上完全不同..
耶穌就是那個創造萬有的道..
他來到這個世界上要彰顯神的真理..
所以約翰福音開始他很強調真理和道路..
而第三,就是中間那裡,也就是我今天最想講的就是十誡..
約翰很想回應猶太人對十誡的詮釋,特別是第一和第二誡..
大家背不背得到第一誡和第二誡?.
除我以外,不可別臣,然後第二誡呢?.
不可為自己做偶像,是不是?.
我們這裡要花一點時間去了解一下,.
究竟猶太傳統如何去理解十誡,頭兩誡,對他們來說很重要..
但是他們有什麼詮釋的錯誤,要約翰這麼用力,這麼長氣去反駁他們呢?.
其實在猶太傳統裡面,上帝是怎樣的?.
是無形無象的,是無處不在的..
如果有人說這個人就是上帝,他們就會覺得你肯定是吹牛,.
還有這個人一定是人造的偶像..
就好像當年他們在西奈山下面做了一隻金牛一樣,.

$^{81}$他說這個就是帶你們出埃及的上帝,這隻金牛就是了..
於是猶太人在受到逼迫流散在不同的地方,.
受到希臘羅馬人的統治的時候,.
他們需要向羅馬人告訴,其實我們猶太人都是很有文化底蘊的,很有智慧的..
他們就作了一本書,叫做《變西拉自焚》,.
裡面就很避重就輕的形容,上帝就是很像一個人..
他不敢說上帝就是一個人,.
他說上帝就是很像一個人,.
好像一個智慧在人的當中..
所以他們就說這個上帝依然都是無形無像,.
他無處不在,就好像希臘裡面的諾戈斯智慧一樣,和人一起..
雖然在人的當中,但依然都是無形無像,.
就好像大家現在可能在用的AI一樣..
叫Siri,你幫我想一下這個決定對不對..
他突然就出聲說,這樣這樣做才對的..
這個是一個很有智慧的導人,.
他這本書就是這樣寫的..
為了明白多一點為什麼猶太人拒絕耶穌道成肉身,.
我們一起看看這本書好不好..
《變西拉自焚》,.
大家好像很少看這本書,因為這本是次經來的..
他這樣說的,.
那創造萬有的,向我下了命令,.
那創造我的,使我的著墨安置,.
吩咐說要在那國內居住,在以色列中取得產業..
從太初在遺行之先,他就已經創造了我,.
直到永遠我總不滅絕..
在他愛慕的城市裡,他使我有了安身之所,.
即在上主的份子中,他的基業中..
大家可能覺得從太初之先在永恆的時候已經創造了我,.
這一句是很有約翰福音的味道,.
如果你掩蓋了變西拉自焚的話,.
大家可能會覺得這個是約翰福音..
這裡的我其實指的是智慧,使我的著墨安置,.
意思是這個智慧就住在當年以色列人出埃及的時候的著墨裡,.
跟他們一起移動,跟他們一起生活..
他們就想強調上帝就好像智慧,就好像一個人那樣,.
去到在人類當中居住,跟人一起..
但是,因為猶太人對於十誡有很特別的理解,.
第一誡是什麼?不可有別臣..

$^{121}$第二誡是不可為自己做偶像..
所以他們就會覺得你做有型的事就等於偶像..
他們不敢將智慧,不敢將上帝形容為一個真實的人,.
你只能說好像,他就好像一個人,好像Siri那樣,.
經常跟著你.以人法,上帝只能好像一個人那樣,.
無形無味地住在人類當中,跟以色列人一起..
這裡說的是上帝的一部分,上帝的分子,.
上帝的一部分都不能夠成為道成肉身,.
所以這樣他們才能夠繼續他們對十誡的詮釋..
所以到了這裡大家就會開始明白,.
楊豐陰這裡說的道成肉身,還提到摩西的律法,.
摩西的律法其實就是在西大山拿的法板,.
下來的十誡,他是在說服聽眾,這個理解才是對的..
他來到自己的地方,自己的人不接納他..
如果你看十一節,你會看到他來到自己的地方,.
自己的人不接納他,自己的人是什麼人?.
猶太人,沒錯,因為耶穌是猶太人,.
自己的人不接納他,是因為他們對於摩西的律法的詮釋,.
十誡的詮釋跟楊豐陰有一個很根本的分別,.
在猶太傳統底下所詮釋的上帝,.
是絕對絕對不會成為肉身的,.
他絕對不可能用一個人的身分出現在他身邊,.
所以當約翰說真理和恩典出自耶穌這個彌賽亞的時候,.
其實就是在拆解猶太人對於十誡的詮釋,.
不可有別臣,不可做偶像..
猶太人世世代代這樣傳下來的傳統,.
說上帝不可能是用人出現的,.
這個傳統,如果有這樣的人說自己是上帝,.
他一定是一個偶像,是一個別臣,.
這個傳統他們傳了多少年?.
傳了幾千年,很深刻很牢固地刻在人的生命裡面,.
這份信念你可以說是執著,可想而知很難去改變,.
就算猶太人受到逼迫的時候,.
他們都沒有將上帝形容成希臘的神明那樣,.
以人的身分進入世界,.
他極其量說他很像,他很像一個人,.
這份信念和執著是多麼的堅毅,.
情況就好像,大家有沒有見過這個食物金字塔?.
大家可能有見過吧?.
這個新的食物金字塔和我們長久以來的教育,.

$^{161}$和我們的文化,和我們家庭的習慣都很不同,.
告訴我們我們以前學習的五穀類是最健康的,.
人體所需的東西,今早大家吃的米線,.
都應該是米線會多一點,.
對我們健康是最大的益處,.
但是最近美國衛生部發表了一個顛覆我們想像的食物金字塔,.
和我們以前很堅信的東西完全不同,.
他們總結發現原來五穀類不是我們人體最需要的營養,.
反而是什麼?.
反而是脂肪和蛋白質,脂肪和蛋白質,.
將我們傳統的飲食習慣完全顛倒了,.
令我們所相信的所謂健康的飲食完全不同,.
無論我們屬於哪一邊的飲食,.
你做168還是KETO,.
我們自己都要重新去審視對於健康飲食的一個習慣,.
你究竟這樣吃健不健康呢?.
你剛剛吃的米線究竟健不健康?.
我們會再一次再審視一下究竟我們吃的東西健不健康,.
現在才跟我說不要吃那麼多,不要省省吃,.
我們會怎麼做呢?.
如果將這樣的處境代入猶太人,.
十誡的處境就是他一直相信我們之前信的食物金字塔才是真的,.
你現在跟我說倒轉的那個才是真的,.
他們就會覺得頭突然間好像爆炸一樣去想,.
我從小到大都是這樣相信的,.
我從小到大都是這樣說上帝是無形無味,.
住在人的當中,.
你現在跟我說上帝可以道成肉身,.
可以沒有違反十誡,.
你可以想像到這個對於他們的信念的挑戰是多麼的巨大,.
多麼的劇烈,.
就算我們對著這個新的食物金字塔,.
我們都會覺得很衝擊,.
我們不停地找他的科學證據,.
質疑他們的立論,質疑他們的總結,.
甚至我們會經過很多輪的思考和驗證,.
我們開始慢慢才想究竟這些東西是不是真的呢?.
或者有些人就是我不理你,.
我吃我的飯,你就吃你的肉,.
我就總之不理你,.

$^{201}$那我自己的信念是怎樣呢?.
我就不需要受到挑戰,.
我不需要改變,.
同樣地對於猶太人來說,.
這個衝擊的過程很需要空間和時間,.
所以約翰福音很詳細地說,.
他說完又說,提完又提,.
目的是怎樣呢?.
讓聽眾知道有些我們一直以來堅持的信念,.
在做的傳統,可能做了很多年,.
未必是對的,未必是對的,.
上帝是會改變我們對於世界的看法,.
對我們自己,對事,對我們身邊的人的看法,.
這個道成肉身的上帝,.
就一直在等待著人去改變,.
好像有些人在未信主之前就覺得,.
基督教也是道人向善,.
和其他宗教,和佛教一樣都是道人向善,.
所以大家都是做好事,.
然後勸大家做好人,.
堅守一些道德的底線,.
那我怎樣衡量善呢?.
多數都是用自己和別人比較,.
就是自己不算差,.
譬如像我在M記,.
我見到隔壁桌子吃完東西又不收拾,.
那我就不同了,.
我就自己收拾自己的東西,.
我多有原則和多有愛心,.
這個就叫善,.
這個是未信主之前的一個比較,.
因為你要和別人比較才知道誰壞誰好,.
但是被神改變之後,.
你現在知道自己是罪人,.
發現所謂善的原則,.
不是用人的標準和和別人比較,.
是上帝完全的公義,.
聖潔的標準去看,.
本身你會覺得自己也不錯,.
但是慢慢會承認自己是一個罪人,.

$^{241}$內心很多自我中心,計較,驕傲,.
只是平時包裝得好一點,.
明白罪不只是做錯事那麼簡單,.
是離開神自我中心的一個狀態,.
以前我覺得對人好不傷害人就行了,.
但信主之後我就發現,.
上帝改變我們內心的看法,.
原來你表面對人好,.
內心也可以很黑暗,.
內心也可以很遠離神,.
不知道大家有沒有看最新的《尋秦記》,.
我應該不會劇透的,.
我先說一些大家知道的東西,.
在《尋秦記》裡面,.
最後一集已經說了,.
《尋秦記》的作者將秦王嬴政的身世重新詮釋,.
他將他描述為已經死了,.
在戰場上已經被人打死了,.
然後是項少龍找了趙國的王子趙盤去頂替他,.
為到秦王嬴政之後焚燒坑儒的歷史去做解釋,.
就是他不希望有人知道項少龍的存在,.
而知道他是一個假冒的秦王,.
大概故事線是這樣,.
那就沒有劇透了,.
如果想看之後就去看電影吧,.
我的看法當然是無從稽考,.
是作出來的,.
但是也給我們一個新的角度去審視秦始皇所做的一切,.
究竟他做焚燒坑儒是為了什麼呢?.
上帝作為我們人生的主宰,.
我們的作者,.
他為我們人生都鋪設了很多很精彩,.
很美妙的故事,.
他希望我們從他的角度去審視我們對人生很多很大的事件的一些看法,.
就好像當年猶太人一樣,.
他覺得十誡就說了,.
不可做偶像,不可有別神,.
怎麼會有一個神可以道成肉身呢?.
這樣是一個偶像的行為,.
但是在我們生命出現的耶穌,.

$^{281}$他就成為了一個很大的一個解法,.
就是他當初上帝的話語,.
他由起初已經存在,.
他來到這個世界,.
在一開始,.
在十誡未出現之前他已經在這裡,.
他成為了一個話語,.
進入了這個世界裡面,.
所以當年猶太人對於十誡的詮釋有很大的偏差,.
耶穌就要利用約翰福音這麼長的經文去重新提及摩西的律法,.
神的道路和真理和恩典,.
去令他們去審視自己對十誡,.
對道成肉身耶穌的看法,.
但是呢,.
現實是很殘酷的,.
人總是想回到自己很熟悉的詮釋和理解,.
對自己對別人對事件有一個很安全的位置,.
我不會受到挑戰,.
其實我信了一段,.
我信主一段日子之後,.
我就被蒙召去做傳道人,.
我去讀神學,.
當時我就秉承著一個理念,.
人活著為了什麼?.
敬拜上帝,.
在我的詮釋裡面,.
敬拜是動詞,.
怎麼做事?.
怎麼也不是你不做事,.
坐著就叫做敬拜,.
我就很拼命地做,.
在教會和神學院,.
我什麼服事都做,.
教會,師他,組長,洗廁所,抹桌椅,在圖書館執書,.
一個年教五十二個星期,.
主日學,情侶輔導,教導輔導,.
總之我覺得我懂的東西我都會做,.
我把所有東西拿出來服事,.
因為我很堅信我自己,.
人活著為了什麼?.

$^{321}$敬拜上帝,.
敬拜一定做些事,.
後來有一個神學院老師提醒我,.
他問我,.
你做這麼多事,.
為了什麼?.
我聽到就很愕然,.
我心想,.
做這麼多事當然是為了上帝,.
還用問,.
我當然沒有回答他,.
他再問,.
你做這麼多事,.
真正為了什麼?.
那一刻我覺得自己好像被冒犯,.
我做了這麼多事,.
當然是為了上帝,.
我竟然冠來別人的質疑,.
還要有神學院老師這樣跟我說,.
後來我回去想了一段時間,.
我才發現這是我的保護機制,.
原來我做了這麼多事,.
一直是我的生存機制,.
我一直的成長告訴我,.
我的動機是可以用我的表現去換取別人的接納,.
我內心有一個很深深處的渴望,.
還有一個恐懼就是說,.
我很怕上帝因為我做得不夠好,.
不夠多,.
而不接納我,.
就算信了主,.
也不代表我們從小到大的信念會自動被改變,.
當我很努力事奉,.
以為我做這麼多事是為了上帝的時候,.
後來我才發現原來,.
我的自我價值很容易綁在我的成績,.
綁在我的工作表現,.
別人怎樣看我,.
我身邊的人有沒有接納我,.
原來我做這麼多事,.

$^{361}$其中一部份的動機是在換取別人對我的接納,.
看好我,.
於是結果是怎樣呢?.
你失敗的時候就會很容易自卑,.
很自責,.
做得成功的時候就很容易會驕傲,.
或者甚至會很害怕,.
會失去別人對你的一個好的印象,.
當我深入被神改變之後,.
還要在神學院才開始,.
被改變後我才發現我裡面很多堅持的信念,.
的動機都不是真理,.
只是我成長上的一些保護機制,.
或者我一些生存的狀態,.
以前我也很明白,.
我是按神的形象被造,.
神所愛我不是靠我多成功多有價值,.
上帝才會愛我,.
我一定很明白這些,.
但我裡面的前面式告訴我,.
我從小到大的生存模式建立了一套很牢不可破的信念,.
就好像有些人說江山易改,本性難移,.
但在基督裡面又不是那麼絕望,.
上帝是給人有機會改變的,.
他用了很多年去拆解我對於所謂敬拜要做很多事的想法,.
其實人活著本身和上帝建立關係就是敬拜,.
即是人本身的Being,那個存在,.
今天你可以呼吸,你和他有關係,.
那個就已經是一個敬拜,.
亦都很直白地問我這個問題,.
令我當頭棒喝,重新審視我自己的動機,.
給上帝改變,.
所以來到總結,.
我們應該可以一起回答第一和第二個問題,.
就是當時的猶太人為什麼拒絕耶穌,.
十誡為什麼成為他們最大信耶穌的阻礙,.
就是因為世代的猶太傳統對於十誡有一個很堅持的理解,.
就是不可有別神,不可作偶像,.
就等於神不能夠成為肉身,成為人,.
這個東西到現在都在堅持,.

$^{401}$如果你去以色列短孫和傳福音的話,.
很多猶太拉比都這樣說,.
神不能夠成為人,.
這個其中一個他們拒絕耶穌的一個原因,.
而第三個,這樣會不會同樣地成為我們的阻礙,.
我猜這個要我們自己才回答得到,.
我鼓勵大家,第一就是主動檢視自己的堅持,.
究竟堅持什麼呢?.
我鼓勵大家在上帝面前安靜地問自己,.
你現在堅持的是什麼?.
是真理?是傳統?是習慣?還是自己的面子?.
我建議大家寫下一兩樣自己很不想被人挑戰的觀點,.
你帶入祈禱裡面,.
你看看上帝怎麼回應你,.
也可以在每天讀經裡面特別留意一下,.
有沒有聖經對你一些很新的衝擊呢?.
有時候,如果上帝透過聖經,.
透過祈禱給你一些新的衝擊的時候,.
我邀請你誠實地去面對,.
勇敢地去看看究竟這件事是不是真的要審時,.
第二,就是鼓勵大家聆聽不同的聲音,.
鼓勵大家刻意跟立場不同,背景不同的朋友聊天,.
不是為了辯論,是為了明白對方,.
也明白自己,.
最近Bella有幾個workshop幫助我們認識自己,.
也認識我們身邊的人,.
對我來說最大的得著就是,.
我能夠聽到別人眼中的我是怎樣,.
我帶給別人的感覺是怎樣,.
我帶給別人的動力是什麼,.
也讓我聽到不少令我成長的聲音,.
有時候那些成長的聲音對我來說,.
本身一些信念是一種挑戰,.
但同一時間也令我重新審視,.
我一直以來的一些信念,.
究竟那些東西是否符合聖經,.
是否一個真理,.
有時候我們跟朋友聊信仰傳福音,.
我們學習一下,聽聽別人關心什麼,.
我們有時候會覺得,.

$^{441}$將一套信仰的模式套在別人身上,.
他們就會接受,.
但我們回想一下,.
當你未信主的時候,.
你也不是一下子直接接受上帝所給我們的東西,.
心態上的改變,生命上的改變,.
是要時間和要關係,.
如果你將心比己,.
你看人家想一下,.
看這個世界怎樣理解這個世界,.
怎樣看身邊的人,.
我們才可以找到身邊跟人有關係的接點,.
我們一起祈禱..
天父,多謝你讓我們有這個時間,.
在Hailey的面前,.
就好像那個食物金字塔一樣,.
我們有時候都可能要重新審視一下,.
我們人生很多的信念,.
很多我們堅持著做的事情,.
道理是不是對的,.
給我們每一個人,.
無論是初信,.
還是我們已經跟從很久的弟兄姊妹,.
給我們每一個人都能夠看見我們裡面,.
內心深處的東西,.
也給我們互相幫助,.
去幫助弟兄姊妹可以成長,.
在裡面更加真實,.
在裡面更加看見主理的真理,.
我這樣禱告,.
奉靠耶穌得勝免疾,.
阿們..
祝福大家..
\newpage



\section{19:11-16}
\label{sec:xg_qWk2v3jg}
\textbf{The King on the White Horse (Rev 19\_11-16) - Rev. Mark Chan 白馬之王 (啟19\_11-16) - 陳崇基牧師}
\newline
\newline
連結: \href{https://youtube.com/watch?v=xg_qWk2v3jg}{\texttt{https://youtube.com/watch?v=xg\_qWk2v3jg}} ~~~~ 語音日期: 2026-02-22
\newline
\newline
\hyperref[sec:Ulr7EZgFyUo]{\small{< < < PREV SERMON < < <}}
~
\hyperref[sec:index_chronic]{\small{[返順時目]}}
~
\hyperref[sec:index_scriptual]{\small{[返順卷目]}}
~
\hyperref[sec:KdIPOrLMIS0]{\small{> > > NEXT SERMON > > >}}
\newline
\newline
$^{1}$[ Foreign Language ].
>> All right..
So today is the Lunar New Year Day, and it's the year of horse, right?.
>> Horse..
>> Horse. And if you were born in the year of '26, you're most likely a horse..
If you were born in 1918, most likely you're a horse..
And maybe '14 or '02 or '90, '78, '66, or '54 and '42 and '30 and 1918, you are probably horses, okay?.
[ Foreign Language ].
And how many of you are horses?.
Any? Oh, you are..
[ Foreign Language ].
Okay. So Kelly, you were born in 2026?.
[ Laughter ].
How many of you are horses?.
How many? Just one?.
>> Only one..
>> All right..
So yeah, congratulations..
That's your year, Kelly, all right?.
So, no. When you think of horses, you know, you probably think.
of speed and drive and success, right?.
And of course, in the new year, especially Chinese New Year, we all hope for some prosperity,.
some breakthroughs, and everything will go well and smoothly for us..
[ Foreign Language ].
I've got to tell you, two days ago, I went to a church to attend their, you know,.
kind of a new year, a little gathering..
You know, they got some games for New Year, and there was one of the tables,.
you're invited to write some banners in Chinese on the red paper..
And, you know, one of them is, you know, [Foreign Language].
And then, I was invited to do something like that..
[ Foreign Language ].
And then, this is some of the stuff that I wrote..
So, on the right is, can you see that?.
[ Foreign Language ].
On the left is, if you know how to translate it, is you get rich, you know, quickly..
Yeah..
[ Foreign Language ].
And then, one of the brothers, you know, kind of gave me a nickname, because [Foreign Language].
you know, sounds like Mark Sir..
[ Laughter ].

$^{41}$So, he said, [Foreign Language].
Okay, so..
>> Yeah, everybody..
>> Yeah, in any case, when you see words like this, and you know, something like this,.
in the new year, we often think of prosperity..
And, in Chinese culture, horses, of course, often symbolize energy, speed, and success..
And, we do want to have prosperity in the new year, smooth process..
[ Foreign Language ].
And, I want to turn to a passage today that mentions horse in the Bible..
And, that is the book of Revelation..
And, it asks a different question..
You know, the book of Revelation, we view a horse and ask a different question..
It's not just how fast you get rich, or how fast you're running..
It is a different question..
It's who you are following, and where is that horse taking you..
And, that is a totally different question we ask..
You know, if the direction is wrong, the faster you go, the more dangerous you get..
Right?.
[ Foreign Language ].
And, you know, the book of Revelation is the last book of the Bible..
Right?.
It's the last book of the Bible..
And, the book of Revelation is not written for our, you know, curiosity..
Sometimes we turn to the book of Revelation, we find, probably feel like, oh, there, it's.
newspaper written long time ago, written to us..
No, it's not like that..
You know, book of Revelation is not for us to foretelling the future..
It is written to strengthen the church back then, because the church was under pressure..
It was under the pressure of the Roman Empire..
And, the Roman Empire demanded their allegiance and loyalty..
And, they got so much pressure being Christians..
[ Foreign Language ].
Now, the book of Revelation has its own context..
In the first century, in the first century, almost 2,000 years ago, the Roman Empire was.
the most powerful and wealthy empire in the world..
And, it was skilled with shaping people's imagination and using different methods..
And, the Roman Empire claimed that it brings peace and prosperity and order to all the people..
And, many cities practiced something what we call imperial cult, emperor worship..
And, in different main cities, you know, like Chicago, like New York, they will have temple.
that would dedicate to the emperor..

$^{81}$And, people would go there to worship and to honor Caesar as their lord..
[ Foreign Language ].
And, if you hear the message, you kind of hear the message being sent by today's empires, right?.
Either it is any empire, any country in the world..
And, if you refuse to bow down to worship the emperor, the king, then you will take.
risk of losing business or social standing or any protection..
And, Revelation is speaking into that kind of world and try to open up the curtain to.
reveal what is truly having power, who is really sitting on the throne..
[ Foreign Language ].
So, we're going to read this passage that mentions the horse..
But, you know, the main message you have to remember is that the main message of Revelation.
is not to predict or foretell the future..
It used some sort of different kinds of writing, what we call apocalyptic literature..
It's a literary form, some sort of writing that is familiar to the Jews of the time in.
the ancient time..
And, it used those kind of language and imageries to proclaim the truth with some vivid imageries..
[ Foreign Language ].
So, it's telling you the one who truly rules over everything is not the Roman Empire, not.
the greatest power of the day, not the Roman emperor on the throne who claimed to be divine,.
who claimed to be God, and to claim that he is the king of kings..
[ Foreign Language ].
And, the one who is truly ruling over everything is Jesus Christ who died and rose again..
And, Chapter 19 of Revelation here is where the whole book of Revelation is trying to bring to a conclusion..
[ Foreign Language ].
So, we're going to read verse by verse in English and in Chinese..
Verse 19, 11, "I saw heaven standing open, and there before me was a white horse whose.
rider is called Faithful and True..
With justice, he judges and makes war.".
[ Foreign Language ].
Verse 12, "His eyes are like blazing fire, and on his head are many crowns..
He has a name written on him that no one knows but he himself.".
[ Foreign Language ].
And, then verse 13 continues, "He is dressed in a robe dipped in blood, and his name is.
the Word of God.".
[ Foreign Language ].
And, the armies of heaven were following him riding on white horses and dressed in fine.
linen, white and clean..
[ Foreign Language ].
And, out of his mouth comes a sharp sword with which to strike down the nations..
He will rule them with an iron scepter..

$^{121}$He treads the winepress of the Faerie of the Wrath of God Almighty..
[ Foreign Language ].
And, on his robe and on his thigh, he has this name written, King of Kings and Lord of Lords..
[ Foreign Language ].
Now, this passage that mentions horse, actually is like the book of Revelation..
It's trying to reveal something to us..
And, first of all, it reveals who the true king is..
[ Foreign Language ].
Now, when the people of the Roman Empire see horses, they can imagine, or the only thing.
they can imagine is Rome..
Roman Empire had all the power and all the horses, all the armies to build an empire..
But, Revelation 19 tried to reverse it and show us the white horse with a different rider.
on the horse..
And, his name is a totally different name..
It's called Faithful and True..
[ Foreign Language ].
Now, in the Roman culture, white horses were symbols of victory..
And, the military generals, they would, you know, beat up somebody and then conquer some.
places, some people..
And, they would celebrate and come back to the city with triumphs, with parades and propaganda..
And, people, the Roman Empire want people to feel, to see that the king, the Caesar,.
the Roman emperor was unbeatable..
[ Foreign Language ].
And then, Revelation uses, you know, similar familiar images and to overturn them..
And, it says the true victorious rider is none but Jesus..
It is not Caesar..
The true rider is Jesus Christ himself..
And, his greatest identity is not power but true and faithfulness..
[ Foreign Language ].
Now, this matters to us today..
Because, in New Year, in Chinese New Year, the go-to phrase for us is always prosperity,.
right?.
[ Foreign Language ].
Okay..
We like money, you know..
Money is our God..
Money is our horse, right?.
When we want to follow those horses and we want promises of success, you know, we want.
money..
We want status..

$^{161}$We want our career..
We want comfort..
We want control..
And, you know, these horses run really fast..
And, but, you know, in reality, the Bible says they do not keep their promises..
[ Foreign Language ].
And, Jesus is different because he is faithful..
He is true..
He is justice..
And, this New Year, do not only ask, please do not only ask what I will achieve, either.
success or status or stay in America or whatever..
Ask different questions..
The Bible challenges you to ask different questions..
Who can I trust with my life?.
[ Foreign Language ].
And then the second point here is to tell us that to follow the true King..
[ Foreign Language ].
But Rome demand that people bow to Caesar, bow to the Roman Emperor..
And Revelation 19 challenges us to tell us Jesus is the true King and he wears many crowns.
because every other authority is tiny and small compared to the authority of Jesus..
[ Foreign Language ].
And not only he has many crowns and all the crowns of the world,.
Revelation also tells us he is wearing something different..
He is wearing a robe dipped in blood..
Now, to the Roman culture, that is a military general, right?.
Shed the blood of innocence and come back as victory..
But this tells us Jesus is a totally different King and Jesus does not conquer by becoming evil..
Jesus conquers through sacrifice himself..
He died on the cross and rose from the dead and he defeat evil with justice and righteousness..
[ Foreign Language ].
And there's also sharp sword coming out of his mouth and his weapon is not tanks and guns..
You know, it's not knives or any weapons of the world..
His weapon is the word, the gospel, the gospel of Jesus Christ, the gospel of love and peace..
And Rome conquered by violence and fear but Jesus conquers in a different manner,.
by truth, by his witness and by his holy judgment..
[ Foreign Language ].
And this is why the final title for this horse rider, white horse rider is so important.
because it tells you he is truly the King of kings and Lord of lords..
And it means your Lord, you know, your leader of your life..
It's not your success, it's not your job, it's not your money and it's not your anxiety..

$^{201}$It's not your family either..
Success is not your Lord and Jesus is the only true King and true Lord..
[ Foreign Language ].
In every generation, there is always an empire trying to shape us and conform us into bowing down to the empire, to the power..
And for many, you know, in our culture, in our Asian-Chinese culture, you know,.
the pressure sometimes is very, very strong and sometimes invisible, right?.
It could be, you know, saving your face..
It could be comparison..
It could be achievement of all kinds, including your children, financial security or social expectation..
And revelation asks a different question..
Who will you bow down to?.
[ Foreign Language ].
Now, so in this new year, of course, you know, many people would desire success..
[ Foreign Language ].
And revelation, the Bible, offers all of us something totally different and something much better, indeed the best..
It's not success that's driving your life, but a Savior that died for you and loved you, accepted you..
And His name is faithful and true..
And His victory is already guaranteed..
He is going to come back..
His victory is final..
[ Foreign Language ].
So, brothers and sisters, the most important new year decision is not your pace..
It's not how fast you go..
It is your allegiance..
So, the Bible challenges to step off the horses that cannot save you..
Put your year into the hands of the white horse rider, King of Kings, Jesus Christ..
[ Foreign Language ].
And this is what we do when we pray in conclusion..
Let's pray together and we're going to commit ourselves once again to worship the true King of Kings and Lord of Lords..
Let's all pray together..
[ Foreign Language ].
Lord Jesus, you are faithful and true..
And you are the King of Kings..
In this new year, forgive us for bowing to false lords of all kinds..
Lead our life and lead our families, our work, and our church..
And give us the courage to follow you and the peace to trust in you..
[ Foreign Language ].
In Jesus' name, amen..
[ Foreign Language ].
We have a dumpling party, right?.

$^{241}$Yeah, and then we'll go out and enjoy fellowship together..
[ Foreign Language ].
God bless you..
\newpage



\section{啟示錄 19:11-16}
\label{sec:KdIPOrLMIS0}
\textbf{白馬之王 (啟示錄19\_11-16) - 陳崇基牧師}
\newline
\newline
連結: \href{https://youtube.com/watch?v=KdIPOrLMIS0}{\texttt{https://youtube.com/watch?v=KdIPOrLMIS0}} ~~~~ 語音日期: 2026-02-22
\newline
\newline
\hyperref[sec:xg_qWk2v3jg]{\small{< < < PREV SERMON < < <}}
~
\hyperref[sec:index_chronic]{\small{[返順時目]}}
~
\hyperref[sec:index_scriptual]{\small{[返順卷目]}}
~
\hyperref[sec:am5WL2bo5gk]{\small{> > > NEXT SERMON > > >}}
\newline
\newline
啟示錄 19:11-16
\newline
\begin{longtable}{cl}
\hline
\hline
章節 & 經文 (和合本修訂版)\\
\hline
19:11 & \begin{tabularx}{0.7\textwidth}{X} 後來我看見天開了。有一匹白馬，騎在馬上的稱為「誠信」、「真實」，他審判和爭戰都憑著公義。 \end{tabularx} \\ \\ \relax
19:12 & \begin{tabularx}{0.7\textwidth}{X} 他的眼睛如火焰，頭上戴著許多冠冕；他身上寫著一個名字，除了他自己沒有人知道。 \end{tabularx} \\ \\ \relax
19:13 & \begin{tabularx}{0.7\textwidth}{X} 他穿著浸過血的衣服；他的名稱為「神之道」。 \end{tabularx} \\ \\ \relax
19:14 & \begin{tabularx}{0.7\textwidth}{X} 眾天軍都騎著白馬，穿著又白又潔淨的細麻衣跟隨他。 \end{tabularx} \\ \\ \relax
19:15 & \begin{tabularx}{0.7\textwidth}{X} 有利劍從他口中出來，用來擊打列國。他要用鐵杖管轄他們，並且要踹全能神烈怒的醡酒池。 \end{tabularx} \\ \\ \relax
19:16 & \begin{tabularx}{0.7\textwidth}{X} 在他衣服和大腿上寫著「萬王之王，萬主之主」的名號。 \end{tabularx} \\ \\
[1ex]
\hline
\hline
\end{longtable}
$^{1}$大家好 各位馬年快樂.
先拜一拜年 今天是年廿幾嗎.
28 是吧.
糟了 我弄錯了.
我知道今年沒有年三十的.
所以星期二是初一.
大家是不是放假.
真的放假.
美國來的 請假放假.
自己放自己假.
如果你是馬年的話.
誰屬馬年的.
有哦 屬馬年的.
我們特別恭喜你.
馬年快樂.
我和弟兄姊妹聊天的時候.
我發現年輕人不會背十二生肖.
有誰會背十二生肖的.
有誰會背.
這麼少人.
下一年是什麼.
馬年之後是什麼.
怎麼數.
有誰不會數十二生肖的.
這麼多.
那些年輕人.
很容易.
鼠牛虎兔龍蛇馬羊猴雞狗豬.
很容易 你單字背.
好像.
創出利民生 殊士得殺王.
舊約目錄.
都是這樣背的 怎麼會不會背.
請你去華埠廣場轉一圈.
沒理由不認識的.
你是華人.
說到馬年.
馬就是速度和成功.
和那種衝勁.
我們希望新一年.

$^{41}$都是盼望有些突破.
個個都是事事如意.
特別我們新一年.
新的馬年.
我們有很多好的消息.
我們非常有盼望.
例如.
星期六Chimay會去多倫多參賽.
卡拉OK大賽.
我們希望他馬上成功.
希望他馬上成功.
上個星期.
我去了一個教會.
有一個小的年宵.
和團體一起去.
他們其中一個攤位寫揮春.
寫什麼呢.
找來找去寫什麼.
有人說寫成會.
成會搭馬.
好嗎.
萬馬奔騰.
成會搭馬.
我寫的那些字.
很噁心.
美的那些.
藝人寫的我懶得放上來.
左邊有人叫我寫.
馬上有錢.
Percy.
幫我作了一個花名.
馬上就是Mark Sir.
Mark Sir有錢.
太好了.
因為新年的緣故.
我也選了一節經文.
和馬有關.
聖經裡面提到很多馬.
除了搭馬之外.
這段聖經.

$^{81}$是來自啟示錄.
啟示錄第十九章提到一隻馬.
馬上有一個君王.
我們不妨用經文.
來思想一下.
到底馬廉裡面有什麼值得我們反省.
特別是生命的方向.
生命的節奏等等.
這裡提到.
有一隻白馬出現.
白馬上有一位君王.
坐在那裡.
用了很多描述來告訴我們.
這個君王是值得我們敬拜.
和值得我們去跟隨.
但我們打開啟示錄的時候.
特別是新晶頭.
我們要先弄清楚.
因為很多時候未接觸過聖經.
或者我們初初接觸聖經的時候.
要做什麼.
和聖經裡面不同寫作的方法.
我們就會.
當作啟示錄.
做什麼.
就是通勝.
有沒有看過通勝.
今年會不會行運.
或者將來.
未來.
有沒有犯太歲.
要做什麼等等.
那些紅色一本的.
有沒有看過.
在App上看.
當然那些是看上了.
騙你十年八年.
那些東西是不值得相信的.
但很多時候.
我們都會當啟示錄.

$^{121}$是預測未來.
預先寫好的報紙.
告訴我們世界末日到了.
或者世界末了就快到了.
會有什麼事情發生.
然後我們就去嘗試.
做一些配對.
這裡提到一些符號.
一些意象.
和我們今天在報章上.
或者在新聞上.
我們看到的一些形象.
是否配合.
一定是啟示錄裡面說的話.
一路揭曉.
然後告訴我們是對的.
好像預測未來.
其實不是這樣的.
完全不是這樣的.
所以今天你聽了什麼都好.
啟示錄就不是通勝.
啟示錄不是通勝.
不要當它是通勝.
當然它有說未來的東西.
告訴我們耶穌一定會再來.
這位君王一定會.
將新天新地帶來.
公義真理他將會帶來.
這個是我們盼望.
但很多時候不是我們現代人.
這樣說或者那種印象.
就是預先寫好的報紙.
啟示錄有它的背景.
啟示錄是當時第一世紀的時候.
特別寫給一些.
小說的.
少數的教會.
他們當時在羅馬帝國.
他們受到很多壓迫.
受到很多打壓.

$^{161}$受到很多壓力.
於是他們有時候灰心.
有時候失去盼望.
啟示錄.
就是用一個這樣的方法.
告訴他們鼓勵他們.
特別是用一些他們當時熟悉的.
形象.
當時他們熟悉的.
符號象徵.
來表達出上帝仍然是掌權的.
所以是一本.
鼓勵的書來的.
是一本牧養的書來的.
不是預測未來.
大家知道在第一世紀的時候.
我們經常說新約的背景.
就是當時羅馬帝國.
是非常富強.
他們會用很多的方法.
來維持住.
國度裡面.
帝國裡面的太平和秩序.
特別是.
啟示錄背景的教會.
他們很多時候.
在今天的東地中海的地方.
特別是今天的土耳其.
巴勒斯坦地等等.
他們受到很多壓力.
因為羅馬帝國要去到.
維持住.
太平盛世.
為了維持住.
羅馬帝國所塑造出來的和諧.
就叫那些異見分子.
不要出聲.
壓迫他們.
不准他們說任何東西.
不准他們說反羅馬的東西.

$^{201}$或者挑釁羅馬帝國政權的東西.
所以在東地中海.
土耳其巴勒斯坦地.
等等的地方.
他們就會推行一種.
帝王的崇拜.
帝王崇拜.
就是國家領導人有個廟.
就一定要去供奉他.
於是有很多經濟壓力.
對當時的人來說.
當時他們的經濟.
和宗教.
是完全混為一談.
不能分開.
我們就可以分開.
以前是不能分開.
你要去購物.
或者去做生意.
或者你要在商務上有順利.
你就需要去拜一拜.
那個國家領導人.
你要去那個廟.
向他供奉也好.
或者向他上香也好.
你才會有人脈.
如果你不這樣做的話.
你就會失去很多的保護.
所以啟示錄給當時的教會.
去對他們講說話.
因為當時的教會.
他們相信.
是耶穌基督才是主宰.
如果有個國家領導人.
他們要去供奉他的話.
就成為一種很大的.
生命裡面的張力.
那怎麼辦呢?.
我們會吃虧.
我們受到這麼多的壓力.

$^{241}$我們受到這麼多的迫害.
那怎麼辦呢?.
啟示錄就是用一個這樣的方法.
去跟他們展現出.
真正作王掌權的是誰.
所以啟示錄裡面.
我們看這些經文的時候.
我們需要有一個這樣的背景.
怎樣鼓勵到他們呢?.
我們先看看這段經文.
這段經文是講馬的.
右邊那裡有一個空了的框.
當我們讀這段經文的時候.
我們想像一下.
到底他所描述的.
所形容的是會怎樣呢?.
讀一讀好嗎?.
讀的時候.
那些意象或者那些image.
你會想像到一個怎樣的.
一個畫面呢?.
啟示錄十九章十一節.
後來我看見天開了.
有一匹白馬.
這匹白馬騎在馬上.
是稱為誠信.
真實.
爭戰都憑著公義.
他的眼睛如火焰.
頭上戴著.
許多冠冕.
他身上寫著一個名字.
除了他自己.
沒有人知道.
他穿著浸過血的衣服.
他的名稱.
為上帝之道.
眾天君.
都騎著白馬.
穿著又白又.

$^{281}$潔淨的細馬布.
跟隨他.
有利劍從他口中出來.
用來擊打列國.
他用.
鐵匠管轄他們.
並且要.
全能上帝.
烈怒的.
駕走前.
在他衣服和大腿上寫著.
萬王之王萬主之主.
你想像到.
一個怎樣的圖畫.
想像一下.
這些是文字的描述.
到底那個圖畫.
會是怎樣的.
那個形象會是怎樣的.
如果你在AI那裡.
他就出了一張這樣的畫給我.
你覺得.
是否合理.
白馬.
口中有什麼.
利劍.
口中有利劍.
然後.
眼如火焰.
穿著的衣服是什麼.
浸過血的.
有血沾上的.
這裡沒有寫萬王之王.
萬主之主.
大概應該想像到.
一個這樣的.
意象.
這種意象.
雖然現在是AI.
出來的.

$^{321}$但大概應該想像到.
一個這樣的畫面.
並且背景是怎樣的.
當時的人.
他們在羅馬帝國下受壓迫.
如果在羅馬帝國的街上.
你去逛逛街.
你大概會經常看到.
一個這樣的影像.
無論是一些石雕.
或者一些浮雕.
長畫等等.
你都會經常看到有白馬.
或者有馬匹戰馬.
因為在羅馬的文化裡.
羅馬的帝國.
他們是怎樣建國的.
是用軍事建國的.
就好像今天.
用飛機大炮一樣.
每次一個地方.
被羅馬佔據了之後.
佔領了之後.
將那個地方打殘了之後.
那些軍隊就會海船回歸.
回到京都.
或者各大城市.
就會有遊行有儀式.
然後電視台就會宣佈.
一件事就是.
羅馬人多威風.
帝國是必勝的.
帝國既然是必勝的話.
每一個在帝國裡面的人.
就需要就範.
一定要跟他.
一定要拜他.
這個就是普遍羅馬的市民.
在街上會看到的形象.
但是啟示錄.

$^{361}$將這個形象.
顛倒過來.
因為他形容的是誰.
他形容的不是羅馬將軍.
也不是羅馬帝王.
他形容的.
騎著白馬上的.
是另一位.
這位是怎樣的.
他的名號.
是誠信.
真實.
和公義.
有正義.
羅馬帝國不會說這些.
羅馬帝國會說什麼.
誰的拳頭大.
就是誠信真實和公義.
什麼叫公義.
我打殘你就是公義.
什麼叫真理.
我秉鑼金的你.
我就是真理.
這個叫真理.
這個叫誠信.
和我們今天的帝國一樣.
無論你說什麼國家.
英國也好中國也好美國也好.
什麼國家也好.
但這位坐在白馬上的.
啟示錄說.
這位最終會得勝掌權.
一切的君王.
他的名號.
不是他的拳頭.
是誠信真實和公義.
這告訴我們.
到底哪一位才是.
真正的君王.
啟示錄是要做什麼.

$^{401}$啟示錄是要揭穿一些東西.
你看一些意象.
看一些image的時候.
就不是讓你看外殼那麼簡單.
好像洋蔥一樣.
慢慢撕開的時候.
你就會明白裡面原來是爛的.
他就是想揭示告訴你.
在地上的帝國.
你慢慢撕開每一個層次的時候.
其實天下烏鴉是一樣黑的.
無論他打著什麼的旗號.
語言上也好.
image上也好.
其實本身裡面也是腐爛的.
因為這個是帝國的權利.
這是人的腐敗的權利.
如果你說要找.
世間上.
有一些我們相信.
能夠相信能夠去跟隨的.
就一定不是地上帝國的權柄.
千萬不要搞錯這件事.
啟示錄要一直問一個問題.
到底你的生命是向著哪一個方向.
你的生命追趕的方向是什麼.
追求的方向是什麼.
今天我們也有很多.
這些熟悉的畫面.
無論剛才說到的一些帝國.
political power也好.
有很多熟悉的畫面.
就是在我們生命裡面.
我們也有不同的吸引我們.
跟隨的東西.
無論是一些.
我們的事業.
金錢.
或者生命裡面的安逸.
去控制一些東西.

$^{441}$啟示錄就是想.
去揭示這件事.
他們本身是不值得.
我們將生命交上的.
真正你要將你生命的賭注.
放在哪裡.
放在這一位.
其實真正白馬.
真正作王的.
這一位是誰.
當然就是耶穌基督.
因為說到他有.
許多的冠冕.
當時的王.
古代的王.
包括羅馬帝國.
佔領了一個地方的時候.
他們就會在那個地方.
拿了本地君王的.
皇冠.
然後就帶了上去自己那裡.
有些故事是這樣說的.
將軍去到一個地方.
羅馬王去到一個地方.
有幾個冠冕在頭上.
就代表什麼?.
是我已經戰勝了你.
表明我對你們的權柄.
你們是在我之下.
如果你熟讀啟示錄的話.
你就知道.
啟示錄裡面介紹過一些人物.
是有許多冠冕的.
我們現在看的是第十九章.
就結束了.
在頭十八章.
已經介紹過一些.
人物.
是有.
不同的冠冕.

$^{481}$我們之前說的啟示錄.
是介紹過幾個.
所謂叫做.
Unholy Trinity.
邪惡的三一.
有三大的.
邪惡的.
化身.
用不同的symbol.
來表達.
第一個就是龍.
第二個就是從海裡面.
出來的怪獸.
第三個就是地.
出來的怪獸.
在第一個怪獸出現.
這個叫古龍.
龍.
Dragon.
他出現的時候.
啟示錄就去到揭示.
這隻龍是很厲害的.
他頭上有七個冠冕.
戴著七個冠冕.
即是說他掌管所有東西.
原來他的身份是誰呢.
就是.
撒旦.
接著到第十三章.
有隻怪獸.
從海裡面出來.
為什麼從海裡面出來.
因為以前近東.
古代中東的人.
他們會覺得海是一些混沌的東西.
是不能夠戰勝.
是很混亂的東西.
凡是從海裡面出來的東西.
十該其九都是不好的.
所以他們說在海裡面.

$^{521}$有一隻怪獸走出來.
而那隻怪獸就是他們平時.
見到的羅馬帝國.
他們所接觸的羅馬帝王.
外表是很光鮮的.
但他所揭示的一件事.
原來他是一隻怪獸.
並且是在混沌出來的.
第十三章第一至十節.
就說到這個怪獸.
原來就是羅馬帝國.
他有十國七頭.
並且有十隻國.
每一隻國就有一個官面.
有十個官面.
並且他的名號是褻瀆的名號.
因為自稱為神.
自稱為霸主.
所以說到.
白馬上面.
有許多官面的時候.
他就說.
其實真正的官面在誰那裡.
在耶穌基督.
在耶穌基督的頭上.
既然耶穌已經有許多的權柄.
所有眾多的官面.
都戴在他頭上.
地上的權柄.
就會突然縮小.
他就不再是.
可以哈哈爸爸.
不再是有權柄.
我們繼續看看其他異象.
浸過血的衣服.
這些羅馬市民經常都會看到.
因為海船回歸的時候.
總會那些將軍.
士兵.
他們手上有些什麼.

$^{561}$血.
無辜人的血.
因為他們哈哈爸爸統治四方的時候.
回來的時候.
他們手上 衣服上.
就會有無辜人的鮮血.
但你覺得.
這些很明顯.
很明顯.
很鮮明的一種形象.
這個就是啟示錄所說的.
為什麼耶穌基督.
會有沾了血的衣服.
你可以哈哈爸爸.
然後有無辜人的血沾上你的衣服.
但同樣你也可以.
什麼.
你可以犧牲你自己.
流你自己的血.
以至你能夠拯救到.
其他人.
這個就是十字架上的耶穌.
他是有.
沾了血的衣服.
但是他做的.
是用犧牲的血來得勝.
去擊敗邪惡.
去拯救我們.
又有什麼.
有利劍.
這些又是當時.
羅馬帝國最.
威水的東西.
是他們帝國軍隊的.
標號.
兩人的利劍.
有武器.
他拿著利劍的時候.
他就可以.
東征西討.

$^{601}$但這裡說到這位騎著白馬的.
耶穌基督.
他也有利劍.
但他的利劍不是去到廝殺.
他的利劍是什麼.
上帝的道.
他的利劍是口裡出來的.
就是說他所傳的福音.
用福音.
用十字架的.
福音.
來戰勝四方.
就不是用利劍.
不是用廝殺.
也不是用暴力.
不是用恐嚇.
耶穌基督去征服.
世界 改變世界.
是用真理 用愛.
去感動我們.
剛才我們提過.
受洗.
受洗的時候.
弟兄姊妹.
未下水池之前.
他們都會講見證.
有沒有誰.
聽過一個見證.
我信主因為有人用槍指著我.
所以現在我要受洗.
每一個生命都是經歷過.
上帝對他的恩典.
我們不配這個恩典.
所以我們回應.
這個恩典.
是一個以愛去還.
愛的一個道理.
這個道理絕對不是.
地上帝國的道理.
但這個是天國的道理.

$^{641}$這個就是耶穌基督.
是他可以改變.
世世代代的.
千千萬萬的人.
而羅馬帝國是不能夠的.
最後一個的.
這個名號在衣服上.
大腿上有萬王之王.
萬主之主.
你一聽的時候.
你用當時的背景.
一聽這些名號的時候.
你第一個的感覺.
就是這個一定是羅馬帝王.
因為羅馬帝王.
在他的錢幣.
或者石碑上.
經常都說他就是.
King of kings and lord of lords.
因為雄霸四方的時候.
還有誰被他打.
他就是萬王之王.
萬主之主.
每一個都有代號.
羅馬帝王要這樣自稱的時候.
就是想人做什麼.
就是被他無條件的效忠.
You pledge allegiance to me.
你向我無條件效忠.
這個就是帝國地上.
要我們做的事.
但是這位.
為我們死的耶穌基督.
如果你答應我.
你也會回應我的話.
我才是你值得.
無條件效忠.
因為我無條件的救你.
無條件的拯救你.
所以你看.

$^{681}$啟稟到介紹.
最後最後介紹這位.
白馬上的君王的時候.
他就介紹了一位.
信實真實.
公義的救主.
他的德性是最忠的.
並且他是值得.
值得你和我.
是無條件.
效忠.
Unconditional allegiance.
對我們今天來說.
我覺得最重要的那句說話.
是這段.
第十四節.
第十四節說什麼呢.
他說這隻白馬來到的時候.
後面有很多白馬跟著.
看到嗎.
不是只有一匹.
在後面有無數匹的白馬.
跟著.
白馬跟著的.
騎上白馬的人.
後面跟隊的那群人.
穿著什麼衣服.
又白又潔淨的細麻布.
去跟隨他.
所以聽到那些馬聲的時候.
他不只是一匹馬.
一匹馬.
是萬馬奔騰.
但後面跟著的很多白馬.
我第一個感覺就是.
後面跟著的白馬的.
戰士.
又白又雪白.
又有細麻衣.
這些大概應該是天使吧.

$^{721}$應該是那些天使.
或者天使之類的.
確實可以是.
但如果看回啟示錄.
或者經文的話.
似乎告訴我們這些不是天使.
在這個意象裡面.
看到白馬上的軍王.
後面跟隊的那群.
騎著白馬的.
似乎不是天使.
如果看第十七章第十四節的話.
就提到.
有一些人.
是和這位.
被殺的高揚耶穌基督.
是一起作戰的.
並且他們是一群.
是這樣形容的.
是夢召被揀選.
忠心於上帝.
福音的人.
這一群人就是跟著他作戰的.
以什麼作戰呢.
不是飛機大炮.
是以福音來作戰.
第十九章.
就講到教會.
怎樣形容教會呢.
教會就是身穿明亮.
潔白細麻衣的.
這一群人有細麻衣.
細麻衣代表.
就是.
聖徒的義行.
就是我們的.
righteous act.
我們的acts of kindness.
acts of goodness.
我們在生命裡面.

$^{761}$在家庭裡面 在我們身邊的人裡面.
去到行的.
義行.
如果你這樣看的時候.
你會發現第十九章講的.
最終要來.
得勝的君王耶穌基督.
跟著他的那群人.
眾天君.
跟隨他的.
原來就是.
忠心的信徒.
就是基督徒.
他問一個問題.
就是告訴我們.
如果我們跟隨這位基督的話.
跟隨這一隻白馬的話.
我不得我是其中一個.
是一個隊裡面其中一個.
是跟他的節奏行.
跟他的方向行.
是十四節.
似乎就是一個邀請.
invitation.
就是你是不是那個.
其中.
眾天君.
騎著白馬.
跟隨這位真正君王耶穌基督.
到底你的身份是什麼.
你屬不屬於這個天君.
裡面其中一員.
如果是屬於的話.
我們值得將我們的命交給他.
值得將我們的決定.
我們任何的意願.
我們盼望與利益.
跟他align.
跟他合作.
如果你這樣看的時候.

$^{801}$這本書啟示錄.
或者介紹這個白馬之主的時候.
其實就是帶給.
當時的人和我們今天的人.
心裡面有很多安穩.
我們現代的生活裡面.
特別這一代這幾年.
我們生命裡面有很多拉力.
很多不同的張力.
有時候我們都不知道.
自己在跑去哪裡.
不知道在看什麼.
不知道在趕什麼.
總之我們就很忙.
我們都沒有時間去想這些東西.
對嗎.
啟示錄就是叫我們停下來.
到底你正在跑去哪裡.
你正在走去哪裡.
心裡面有很多安穩.
近幾天我自己.
在不同的情況很多人提醒我.
我自己之前也有說過.
我發現我最近.
很趕.
不知道為什麼很趕.
有人提醒我的時候.
我發現其實我在趕什麼.
為什麼每件事都要這麼快.
為什麼每件事都好像很急.
很倉促.
所以有些很關心我的.
突然就跟我說.
其實你這麼快.
你會錯過很多東西.
沿路的風景你都錯過了.
你自己的狀況.
可能你自己都不明白.
我一向覺得.
效率.

$^{841}$做事快,做事趕.
是我引以為傲的.
對嗎.
有人要把事情做好.
但最近我自己就.
被人提醒.
特別是昨天.
在社交媒體.
有一個朋友.
他發了一段.
很簡單的東西.
就是有關安息的.
就是休息.
他就說我們.
很多時候很快,很趕.
停不下來,靜不下來.
我們說話很快.
說話很倉促.
有時候傷害到別人.
做事很快,有時候急忙.
如果我們整天都是.
24/7地看下去.
其實主要的原因是什麼.
我的社交媒體朋友就說.
主要的原因就是.
因為你不相信.
上帝會工作.
你想控制.
我突然當頭棒紅.
到底我的生命的方向.
我的節奏.
是為了什麼.
到底我吸引的是誰.
新的一年裡面.
啟示錄要問的.
不只是要我們.
馬到功成.
不是每件事都要急趕.
快速.
辦事要快速.

$^{881}$成功要快速.
他問我們不是.
Who.
Are you following.
你在跟隨誰.
但是我們跟隨.
這位主的時候.
我們有那種意願.
我們是跟隨他的節奏.
因為如果.
這位君王是最終得勝的.
起碼我們的心境.
我們會有安穩.
如果是這樣的話.
在我們所有的關係裡面.
其實我們的說話.
我們待人接物.
其實都會有所改變.
對我們身邊的人說話.
可能我們會放慢放緩一點.
因為最終不需要我控制.
上帝仍然作供.
當我待人接物.
處理其他人的時候.
和其他人相處的時候.
我可以和平很多.
對不對?.
因為這個世界很多戰火很多爭戰.
但是我知道.
這位白馬上的君王.
是必定會再來的.
我心裡面有安穩.
我都有勇氣.
跟著上帝的節奏.
有平安有喜樂.
可以同心低頭祈禱嗎?.
我們低頭祈禱的時候.
給大家安靜的時間.
在上帝面前你可以說.
最近到底我自己在說什麼.

$^{921}$在忙什麼.
求主讓我能夠安靜.
跟隨祂的節奏.
和祂一起事奉.
祈禱.
主耶穌基督.
你是那位真實,誠信,公義的主.
你是那位萬王之王,萬主之主.
求主你帶領我們.
也讓我們跟隨你的節奏.
不是跟隨世界的節奏.
賜給我們勇氣.
賜給我們這份安穩.
讓我們跟隨你的時候.
我們有平安有喜樂.
並且有活力.
我們信靠你.
帶領著我們每一個個人的生命.
帶領著我們每一個的關係.
從我們自己和自己的關係.
身邊的家人.
工作,教會,社會.
主讓我們跟著你這個.
信實,真實的節奏.
去跟著你.
相信你.
讓我們經歷到你的特性.
帶領著我們.
如果我們太趕的時候.
叫我們在你的慈愛裡面.
能夠停下來.
讓我們知道.
你愛我們.
並且在背後不斷作功.
主叫我們的生命.
能夠不斷跟隨你.
我們巴不得每一個.
都是在天君裡面的一份子.
穿著你所潔淨的衣服.
事奉敬拜你.

$^{961}$祈禱奉傳美救.
好,恭祝各位馬上平安.
\newpage



\section{詩篇 1:1-6}
\label{sec:am5WL2bo5gk}
\textbf{The Righteous Person (Psalm 1) - Rev. Mark Chan 義人的生命 (詩篇1\_1-6) - 陳崇基牧師}
\newline
\newline
連結: \href{https://youtube.com/watch?v=am5WL2bo5gk}{\texttt{https://youtube.com/watch?v=am5WL2bo5gk}} ~~~~ 語音日期: 2026-03-02
\newline
\newline
\hyperref[sec:KdIPOrLMIS0]{\small{< < < PREV SERMON < < <}}
~
\hyperref[sec:index_chronic]{\small{[返順時目]}}
~
\hyperref[sec:index_scriptual]{\small{[返順卷目]}}
~
\hyperref[sec:mjg2Vzt9lP0]{\small{> > > NEXT SERMON > > >}}
\newline
\newline
詩篇 1:1-6
\newline
\begin{longtable}{cl}
\hline
\hline
章節 & 經文 (和合本修訂版)\\
\hline
1:1 & \begin{tabularx}{0.7\textwidth}{X} 不從惡人的計謀，不站罪人的道路，不坐傲慢人的座位， \end{tabularx} \\ \\ \relax
1:2 & \begin{tabularx}{0.7\textwidth}{X} 惟喜愛耶和華的律法，晝夜思想他的律法；這人便為有福！ \end{tabularx} \\ \\ \relax
1:3 & \begin{tabularx}{0.7\textwidth}{X} 他要像一棵樹栽在溪水旁，按時候結果子，葉子也不枯乾。凡他所做的盡都順利。 \end{tabularx} \\ \\ \relax
1:4 & \begin{tabularx}{0.7\textwidth}{X} 惡人並不是這樣，卻像糠秕被風吹散。 \end{tabularx} \\ \\ \relax
1:5 & \begin{tabularx}{0.7\textwidth}{X} 因此，當審判的時候惡人必站立不住，罪人在義人的會眾中也是如此。 \end{tabularx} \\ \\ \relax
1:6 & \begin{tabularx}{0.7\textwidth}{X} 因為耶和華知道義人的道路，惡人的道路卻必滅亡。 \end{tabularx} \\ \\
[1ex]
\hline
\hline
\end{longtable}
$^{1}$- Good morning everyone..
And we will continue our series on the Psalms in the Bible..
And as I told you before,.
many of the Psalms were originally prayers.
spoken by the people of God..
But some of the Psalms in the book of Psalms.
are written for the purpose of teaching us something..
So teaching us how to relate to God,.
teaching us how to see ourselves and see the world..
(speaking in foreign language).
Now we call these kinds of Psalms, wisdom Psalms..
You may have heard it the very first time today,.
but we can call them wisdom Psalms.
because these wisdom Psalms are meant to teach us something,.
to teach us wisdom, make us wiser.
and help us better relate to God..
(speaking in foreign language).
And let's read this, Psalm 1, okay?.
So Psalm 1 has six verses..
Let's read that verse by verse..
Let me begin with English..
Psalm 1, verse 1..
"Thus is the man who does not walk.
in the counsel of the wicked,.
or stand in the way of sinners,.
or sit in the seat of mockers.".
(speaking in foreign language).
But his delight is in the law of the Lord,.
and on his law he meditates day and night..
(speaking in foreign language).
And verse 3, "He is like a tree planted by streams of water,.
which yields its fruit in season,.
and whose leaf does not wither..
Whatever he does prospers.".
(speaking in foreign language).
Not so the wicked, they are like a shaft,.
and that the wind blows away..
(speaking in foreign language).
And verse 5, "Therefore the wicked.
will not stand in judgment,.

$^{41}$nor sinners in the assembly of the righteous.".
(speaking in foreign language).
And verse 6, "For the Lord watches over.
the way of the righteous,.
but the way of the wicked will perish.".
(speaking in foreign language).
Now Psalm 1 is the first Psalm in the book of Sardar,.
and it is like a gate to the temple of God..
So before you enter the temple of God,.
there was a gate that you have to go through,.
and you walk into the presence of God,.
the most holy place..
So Psalm 1 is like the gate of the temple..
(speaking in foreign language).
So the first word of Psalm 1 verse 1.
is the word we translate blessed..
And in the Old Testament, there were two words of blessing..
And you have to note that, okay?.
(speaking in foreign language).
Now the first word I give you is Barak..
And that is the word we use to describe.
the blessing of God..
(speaking in foreign language).
To mean like we bless you because we wanna grant you.
or give you the power of life..
You know, being blessed, being fruitful..
You know, you recall Genesis 1,.
when God created the world, He blessed humankind..
He blessed the world, be fruitful and multiply..
And that is the word Barak..
(speaking in foreign language).
Okay, so learn this word, Barak, okay?.
Say it with me, okay?.
(speaking in foreign language).
Barak, okay?.
But that is not the word in Psalm 1, 1, okay?.
That is interesting..
That's the second Hebrew word we translate blessed.
or blessing, and that is ashare..
We also translate it as blessed most of the time..

$^{81}$(speaking in foreign language).
That's blessed most of the time..
(speaking in foreign language).
So it doesn't mean like give you life or empower you..
It gives you hope and it focuses on something in the future.
that you will receive future reward and final blessing..
So this means a blessed person in Psalm 1.
do have difficulties now..
(speaking in foreign language).
And then this same word, ashare, say it with me, ashare..
(speaking in foreign language).
Ashare, okay, so you learn some Hebrew word..
The word ashare is also used in another book.
called the Book of Job, chapter five, verse 17..
See the word being translated is also blessed..
And the word behind it, the Hebrew word behind it.
is ashare, not barak..
So blessed is the one whom God disciplines..
That is Job 5, 17..
Now, of course, when you think about discipline,.
God disciplines someone, right?.
And you call that someone blessed..
You know, in the state of being disciplined,.
do you feel happy or do you feel good about it?.
Usually not, right?.
Usually not..
But it's called blessed is the person whom God discipline..
So discipline does not feel joyful at the time,.
but at the end, you will receive the benefits of discipline..
(speaking in foreign language).
Now, for example, if you do a physical exercise regularly,.
anyone, one of you go to the gym, you go to exercise,.
you go to do physical exercise..
When you start out working out, okay,.
when you start to work out, does it feel good?.
No, it shouldn't feel good, right?.
If you feel good when you started doing exercise,.
you're doing it wrong, okay?.
It doesn't feel good when you started to discipline yourself,.
discipline your body..

$^{121}$It hurts at first, but you are still blessed.
because if you keep going, if you keep persevering,.
then you will receive the final blessing, right?.
You will receive the final reward..
(speaking in foreign language).
Or another example is a student..
You go to study and in the beginning of the semester,.
does it feel good or does it feel happy?.
Usually not, right?.
You feel painful as you go through the semester,.
but you see the good result will be at the end,.
and so the blessing is not here and now,.
it's something in the future..
(speaking in foreign language).
So Psalm 11 says, "Blessed is the man who does this.
"or does not do those, this things or that things.".
He's talking about the final blessing, final result..
And Jesus used the same word in Beatitudes,.
in the book of Matthew..
He said, plus, he used the same word in Greek now.
because New Testament is written in Greek, not in Hebrew,.
but it's translated in Greek, makaria,.
and that is being used by Jesus in Beatitudes..
(speaking in foreign language).
And the same word, as share makaria..
Okay, say makaria with me, okay?.
Makaria. - Makaria..
- Makaria, okay?.
All right, so as I give you some example here,.
Jesus said, "Blessed are the poor in spirit,.
"and blessed are those persecuted for righteousness.".
(speaking in foreign language).
Now, when you're going through suffering or persecution,.
it doesn't feel good, right?.
The early church faced the same problem..
They had suffering and they had persecution,.
but Jesus say, "If you belong to the kingdom,.
"the final reward will be yours,.
"and final reward will be great.".
(speaking in foreign language).

$^{161}$So I think Psalm 1:1 can be translated this way,.
and that's my translation..
It can be translated,.
the person who will be ultimately rewarded by God.
is someone who does not walk in the counsel of wicked.
and this and that..
(speaking in foreign language).
So it's not talking about here and now.
that you will feel happy..
It's talking about ultimately,.
at the end of your life, you will feel blessed, okay?.
So as we go through Psalm 1,1,.
you also, the Psalm also unpacks.
how a blessed person will be..
He says he does not walk and does not stand.
and does not sit..
And you see the progression here, right?.
It's like a poem progression from walking.
to standing to sitting,.
and describe how the spiritual drift can happen gradually..
To walk is being influenced by the wicked,.
by the opinions of the wicked..
And to stand is to linger,.
to identify yourself with them..
To sit is to settle down, right?.
To feel belong and feel comfortable with the wicked..
(speaking in foreign language).
And you see the movement is from motion,.
walking to permanence, right?.
So in showing us compromises can slowly happen.
and become a fixed way of lifestyle..
(speaking in foreign language).
But the person who is blessed,.
who see the final reward,.
is someone who does not do those, right?.
Does not do those..
The righteous, not the wicked,.
the righteous will fear God.
and may at this time disadvantage themselves..
The righteous will disadvantage themselves.

$^{201}$to advantage others..
But the wicked will ignore God..
The wicked will not see God, will not fear God..
And they will disadvantage others for their own advantages..
(speaking in foreign language).
So it says the righteous does not do those things..
It's an invitation..
The first Psalm of the Bible is to invite you.
to make a decision to be like the righteous.
and blessed person before God..
(speaking in foreign language).
Now let's continue..
Verse two, continue to describe this person..
And it says the righteous person is,.
has having delight, a happy delight.
in the law of the Lord..
(speaking in foreign language).
It's talking about the inner direction of your heart, right?.
It's not something behavioral, it's not legalism..
It's some sort of acquired taste for the righteous person..
Okay, he acquired the taste, the delight.
in the law of God..
And it's kind of like acquired taste..
It's kind of like, you like coffee or durian?.
Something like that..
It's not a natural taste..
It's acquired taste..
I like coffee, I take coffee every single day..
And I didn't like it before when I was much younger..
But I acquired that taste..
Now the righteous, it's the same thing here..
He trains himself, like a student,.
he trains himself, his mind to desire the word of God..
(speaking in foreign language).
And here he says he meditates day and night..
And it does not mean emptying your mind..
You know, I recall one of the experiences I had.
when I was a college student..
I took Philosophy 101, and the philosophy professor.
asked the students in the lecture hall,.

$^{241}$all of us, to empty our mind..
You know, he was trying to do some exercise,.
some sort of transcendental meditation..
So you have to empty your mind and then imagine something..
You know, that is the secular meditation..
(speaking in foreign language).
And then in the Bible, biblical meditation is different..
It means you think about it, think about the word of God..
You study the word of God..
You question the word of God..
You examine the word of God..
(speaking in foreign language).
And it says day and night, it does not mean 24/7..
It means habit, something customary..
You don't neglect it..
You know, it's like your habit every day..
(speaking in foreign language).
And then it also tells you that the object of meditation.
is Torah, or we translate the law of the Lord..
Torah, actually, a better translation would be instruction,.
guidance, or teaching..
(speaking in foreign language).
Now, the word Torah, or translated in English, law,.
or instruction, actually points to the whole revelation.
of God..
So the Bible is like a constitution of the kingdom people..
We read this document because we want to understand.
the principle, the instructions being written here..
(speaking in foreign language).
So someone invites you to grasp the whole Torah,.
the whole counsel of God..
(speaking in foreign language).
And it's pretty sad because many Christians today,.
we are selective..
And someone tells us, you cannot be selective.
in the word of God..
You need the whole counsel of God..
(speaking in foreign language).
And then verse three continues to describe.
this righteous person like this..

$^{281}$He's like a tree planted by streams of water.
and yielding the fruit in season..
And the leaf does not wither,.
and whatever he does prospers..
(speaking in foreign language).
Now, you think about this image..
It's an image of life and stability, being steady..
Imagine a tree with the roots growing deeper and deeper.
into the ground and planted by the water..
And at the right time, the tree will yield fruit,.
and it will become blessing to other people..
(speaking in foreign language).
And then in the ancient years, the culture of the Bible,.
where can you find trees?.
Most of the time, you find trees in temples..
(speaking in foreign language).
And then you remember the first temple,.
so-called temple in the Bible,.
was actually the Garden of Eden..
The Garden of Eden, Eve, right?.
Garden of Eden, with the tree of life..
(speaking in foreign language).
Now, that's the beginning of the Bible,.
but the end of the Bible, the final chapter of the Bible.
in the book of Revelation talks about new heavens.
and new earth, there again, a river and tree of life..
(speaking in foreign language).
So what it means is that the blessed person,.
the righteous person who is dedicated in the Torah,.
in the instruction of God,.
will live in the presence of God, like in the temple..
(speaking in foreign language).
So it's inviting you to be the righteous person,.
inviting you to be the blessed person..
(speaking in foreign language).
And then verse four talks about the wicked..
Verse four, "Not so the wicked,.
they're not like the righteous,.
they are not like the blessed..
They are like shaft, that the wind blows away.".

$^{321}$(speaking in foreign language).
And when farmers get their harvest from the field,.
like wheat or barley,.
they will bring it to the flushing floor..
It's also, most of the time,.
they will bring the harvest grain to the flushing floor..
(speaking in foreign language).
And because the flushing floor.
is usually on the high ground, on the hilltop,.
so you get stronger wind..
(speaking in foreign language).
And the farmers will throw the crushed grain into the air,.
and then it will drop, right?.
And the shaft, the shell of the wheat and barley,.
they have no weight, and they have no life and no future..
They will be blown away by the wind..
(speaking in foreign language).
So Psalm 1 says, "The wicked are like this..
They do not belong to the assembly.
of the righteous, the blessing.".
(speaking in foreign language).
Therefore, it continues to tell us,.
"The wicked will not stand in judgment,.
nor sinners in the assembly of the righteous.".
(speaking in foreign language).
And verse six tells you,.
"For the Lord watches over the way of the righteous,.
but the way of the wicked will perish.".
(speaking in foreign language).
See, if you are righteous,.
if you make the commitment to love the Word of God,.
even though today you may have suffering,.
but finally, ultimately, you will have blessing..
And you need to understand God knows your way..
God knows your life..
This is not just information..
It means there's a relationship..
There's a covenant..
There's a commitment from God..
There's a caring from God..

$^{361}$(speaking in foreign language).
So in conclusion, this psalm,.
the very first psalm in the Bible presents to you.
there are two ways of life..
There are two different directions of life.
and two endings in our lives..
(speaking in foreign language).
And it asks you, the very important question is,.
what is your decision today?.
What is your commitment today?.
(speaking in foreign language).
I think there are some applications.
or take-homes for us..
First of all is that,.
it invites you to know God..
And knowing God is a life decision..
(speaking in foreign language).
And Psalm 1 reminds you that worship,.
knowing God, begins with choosing the way of God..
(speaking in foreign language).
And then secondly, it tells you,.
you have to work on something..
You have to work on knowing the whole counsel of God,.
the whole Bible..
(speaking in foreign language).
Now do not read the Bible selectively..
You have to read the Bible.
or some parts of a Bible that are not familiar to you,.
that are not comfortable to you..
Do not only choose comfortable passages from the Bible..
(speaking in foreign language).
So I wanna invite you this week, all right?.
If you don't have this habit, all right,.
I wanna invite you this week for the next seven days,.
read one chapter of scripture every day..
(speaking in foreign language).
When you have a little bit of time,.
you know, how much time do you need.
to read one chapter of the Bible?.
(speaking in foreign language).

$^{401}$You say, okay, it depends on how long the chapter is, right?.
There's some shorter, there's some longer, right?.
But you know, three to five minutes, right?.
(speaking in foreign language).
But when you have free time, let your mind,.
when you read the Bible, let your mind return.
to the mind of God, to the word of God..
(speaking in foreign language).
So Psalm 1 is a precious psalm..
It stands as a gate, as the opening door,.
as the entrance into the presence of God..
Before entering into the presence of God, worshiping Him,.
we ask ourselves the questions, where are we standing?.
What kind of person are we becoming?.
(speaking in foreign language).
I want to invite you to pray with me,.
to make this commitment today..
(speaking in foreign language).
Let's pray together..
(speaking in foreign language).
Dear Lord, thank you for your word,.
showing us the way of eternal life..
And we do ask you to help us to take delight in your truth..
Be rooted like a tree by streams of water..
(speaking in foreign language).
We ask you to keep us from drifting.
into paths that lead away from you..
We ask you to teach us to walk in your ways,.
so that our lives may bear fruit for your glory..
(speaking in foreign language).
In Jesus' name we pray..
(speaking in foreign language).
Amen..
- Amen..
- All right, well God bless you,.
and I want to see you next week,.
and make sure you do your homework, right?.
Every day, just one chapter, okay, just one chapter..
You can read more, okay, but it is one chapter..
So you say, where to start, right?.

$^{441}$We'll start with Psalm, okay?.
Psalm one, two, three, four, five, six, seven, eight..
Psalm one, two, three, four, five, six, seven, all right?.
(speaking in foreign language).
Okay, I'll see you next week..
(speaking in foreign language).
\newpage



\section{以斯帖記 8:3-9:16}
\label{sec:mjg2Vzt9lP0}
\textbf{神的救贖 人的責任 (以斯帖記 8\_3-9\_16) - 黃國棟醫生}
\newline
\newline
連結: \href{https://youtube.com/watch?v=mjg2Vzt9lP0}{\texttt{https://youtube.com/watch?v=mjg2Vzt9lP0}} ~~~~ 語音日期: 2026-03-02
\newline
\newline
\hyperref[sec:am5WL2bo5gk]{\small{< < < PREV SERMON < < <}}
~
\hyperref[sec:index_chronic]{\small{[返順時目]}}
~
\hyperref[sec:index_scriptual]{\small{[返順卷目]}}
~
\hyperref[sec:_lnAtpZRm3Q]{\small{> > > NEXT SERMON > > >}}
\newline
\newline
以斯帖記 8:3-9:16
\newline
\begin{longtable}{cl}
\hline
\hline
章節 & 經文 (和合本修訂版)\\
\hline
8:3 & \begin{tabularx}{0.7\textwidth}{X} 以斯帖又在王面前求情，俯伏在他腳前，流淚哀求他阻止亞甲人哈曼害猶太人的惡謀。 \end{tabularx} \\ \\ \relax
8:4 & \begin{tabularx}{0.7\textwidth}{X} 王向以斯帖伸出金杖，以斯帖就起來，站在王面前， \end{tabularx} \\ \\ \relax
8:5 & \begin{tabularx}{0.7\textwidth}{X} 說：「王若以為好，我若在王面前蒙恩，王若認為合宜，我若在王眼前得喜悅，請王下諭旨，廢除亞甲人哈米大他的兒子哈曼設謀，要殺滅王各省的猶太人所頒的詔書。 \end{tabularx} \\ \\ \relax
8:6 & \begin{tabularx}{0.7\textwidth}{X} 我何忍見我本族的人受害？何忍見我同宗的人被滅呢？」 \end{tabularx} \\ \\ \relax
8:7 & \begin{tabularx}{0.7\textwidth}{X} 亞哈隨魯王對以斯帖王后和猶太人末底改說：「因為哈曼要下手害猶太人，看哪，我已把他的家產賜給以斯帖，也把哈曼掛在木架上了。 \end{tabularx} \\ \\ \relax
8:8 & \begin{tabularx}{0.7\textwidth}{X} 你們可以照你們看為好的，奉王的名寫諭旨給猶太人，用王的戒指蓋印；因為奉王的名所寫、用王的戒指蓋印的諭旨是不能廢除的。」 \end{tabularx} \\ \\ \relax
8:9 & \begin{tabularx}{0.7\textwidth}{X} 三月，就是西彎月二十三日，當時王的一些書記受召而來，按著末底改所吩咐的，用各省的文字、各族的語言，以及猶太人的文字語言寫諭旨，傳給那從印度直到古實一百二十七省的猶太人，以及總督、省長和領袖。 \end{tabularx} \\ \\ \relax
8:10 & \begin{tabularx}{0.7\textwidth}{X} 末底改奉亞哈隨魯王的名寫諭旨，用王的戒指蓋印，交給信差們騎上御用的王室快馬去頒佈。 \end{tabularx} \\ \\ \relax
8:11 & \begin{tabularx}{0.7\textwidth}{X} 王准各城各鎮的猶太人在一日之內，在十二月，就是亞達月的十三日聚集，在亞哈隨魯王的各省保護自己的性命，剪除，殺戮，滅絕那要攻擊猶太人的各省各族所有的軍隊，以及他們的妻子兒女，奪取他們的財產為掠物。 \end{tabularx} \\ \\ \relax
8:12 & \begin{tabularx}{0.7\textwidth}{X} 見上節 \end{tabularx} \\ \\ \relax
8:13 & \begin{tabularx}{0.7\textwidth}{X} 這諭旨的抄本以敕令的方式在各省頒佈，通知各族，使猶太人預備等候那日，好在仇敵身上報仇。 \end{tabularx} \\ \\ \relax
8:14 & \begin{tabularx}{0.7\textwidth}{X} 於是騎御用快馬的信差奉王命催促，急忙起行；敕令傳遍了書珊城堡。 \end{tabularx} \\ \\ \relax
8:15 & \begin{tabularx}{0.7\textwidth}{X} 末底改穿著藍色白色的朝服，頭戴大金冠冕，又穿紫色細麻布的外袍，從王面前出來；書珊城充滿了歡樂的呼聲。 \end{tabularx} \\ \\ \relax
8:16 & \begin{tabularx}{0.7\textwidth}{X} 猶太人有光榮，歡喜快樂，得享尊貴。 \end{tabularx} \\ \\ \relax
8:17 & \begin{tabularx}{0.7\textwidth}{X} 王的諭旨和敕令所到的各省各城，猶太人都歡喜快樂，擺設宴席，以那日為吉日。國中許多民族的人因懼怕猶太人，就自稱為猶太人。 \end{tabularx} \\ \\ \relax
9:1 & \begin{tabularx}{0.7\textwidth}{X} 十二月，就是亞達月十三日，王的諭旨和敕令要執行的那一日，猶太人的仇敵盼望制伏他們，但猶太人反倒制伏了恨他們的人。 \end{tabularx} \\ \\ \relax
9:2 & \begin{tabularx}{0.7\textwidth}{X} 猶太人在亞哈隨魯王各省的城裡聚集，下手擊殺那些要害他們的人。沒有人能在他們面前站立得住，因為各民族都懼怕他們。 \end{tabularx} \\ \\ \relax
9:3 & \begin{tabularx}{0.7\textwidth}{X} 各省的領袖、總督、省長，和辦理王事務的人，因懼怕末底改，就都幫助猶太人。 \end{tabularx} \\ \\ \relax
9:4 & \begin{tabularx}{0.7\textwidth}{X} 末底改在朝中為大，名聲傳遍各省；末底改這人的權勢日漸擴大。 \end{tabularx} \\ \\ \relax
9:5 & \begin{tabularx}{0.7\textwidth}{X} 猶太人用刀擊殺所有的仇敵，殺滅他們，隨意待那些恨他們的人。 \end{tabularx} \\ \\ \relax
9:6 & \begin{tabularx}{0.7\textwidth}{X} 在書珊城堡中，猶太人殺滅了五百人。 \end{tabularx} \\ \\ \relax
9:7 & \begin{tabularx}{0.7\textwidth}{X} 他們殺了巴珊大他、達分、亞斯帕他、 \end{tabularx} \\ \\ \relax
9:8 & \begin{tabularx}{0.7\textwidth}{X} 破拉他、亞大利雅、亞利大他、 \end{tabularx} \\ \\ \relax
9:9 & \begin{tabularx}{0.7\textwidth}{X} 帕瑪斯他、亞利賽、亞利代、瓦耶撒他； \end{tabularx} \\ \\ \relax
9:10 & \begin{tabularx}{0.7\textwidth}{X} 這十人都是哈米大他的孫子，猶太人的仇敵哈曼的兒子。猶太人卻沒有下手奪取財物。 \end{tabularx} \\ \\ \relax
9:11 & \begin{tabularx}{0.7\textwidth}{X} 那日，書珊城堡中被殺的人數呈報到王面前。 \end{tabularx} \\ \\ \relax
9:12 & \begin{tabularx}{0.7\textwidth}{X} 王對以斯帖王后說：「猶太人在書珊城堡中殺滅了五百人，又殺了哈曼的十個兒子，在王其餘的各省不知如何。你要甚麼，必賜給你；你還求甚麼，也必為你成就。」 \end{tabularx} \\ \\ \relax
9:13 & \begin{tabularx}{0.7\textwidth}{X} 以斯帖說：「王若以為好，求你允准書珊的猶太人，明日也照今日的諭旨去做，並把哈曼十個兒子的屍體掛在木架上。」 \end{tabularx} \\ \\ \relax
9:14 & \begin{tabularx}{0.7\textwidth}{X} 王允准這麼做。敕令傳遍書珊，哈曼十個兒子的屍體被掛了起來。 \end{tabularx} \\ \\ \relax
9:15 & \begin{tabularx}{0.7\textwidth}{X} 亞達月十四日，在書珊的猶太人又聚集，在書珊殺了三百人，卻沒有下手奪取財物。 \end{tabularx} \\ \\ \relax
9:16 & \begin{tabularx}{0.7\textwidth}{X} 亞達月十三日，在王各省其餘的猶太人也都聚集，保護自己的性命，擺脫仇敵得享平安。他們殺了七萬五千個恨他們的人，卻沒有下手奪取財物； \end{tabularx} \\ \\
[1ex]
\hline
\hline
\end{longtable}
$^{1}$大家早安.
很高興又有機會回來跟大家一起敬拜.
今天還是新年 所以在這裡再一次沉氣向大家拜年.
恭喜大家今年新年馬年龍馬精神.
芝加哥我知道每一年都會有農曆新年的巡遊.
我看過報紙說今年的巡遊在下星期日.
不過我走了就來不了.
去年我住在芝加哥也有去見過這個遊行.
是很開心很好玩的一件事.
不過大家知不知道.
如果你在下星期一去到這裡的猶太人社區的話.
你又會見到另外一個巡遊.
這個就是今年的三月二號和三號星期一和星期二.
他們會慶祝猶太人一個節日叫做普爾節.
普爾節不是保利 不要看錯.
保利的里節是三點水的.
這個爾節是王字邊的.
他們會穿著很有趣的化曬裝穿著很有趣的衣服在巡遊.
猶太人和華人很相似的就是他們每年都有很多節日.
而這些節日是有不同的意思.
可能有些是我們很熟悉的譬如是逾越節.
是聖經舊約裡面有記錄過的節日.
普爾節就是另外一個在舊約裡面有記錄過的節日.
它的起源是在《二師帖記》裡面所說到的.
就是說他們會在這個叫做亞達月的十四十五號.
就是為了猶太人脫離仇敵得到平安.
祭普爾原來是抽籤的意思.
就形成了這兩天是做節日就要他們守節.
其實普爾節和逾越節就是猶太人現在所守的節日裡面.
兩個在舊約裡面有紀念有記錄到的節日.
現在猶太人最重要的節日.
我聽過一些猶太朋友.
我的兩個老師都是猶太人.
最重要的是他們的聖誕節.
就是在聖誕前後的一個節日.
除了聖誕節之外.
就是普爾節和逾越節也是他們很重要的節日.
而這兩個就是在聖經裡面有記錄的節日.
今天我就想跟大家用普爾節的故事.
看看神可以怎麼樣在我們裡面施行祂的救贖.

$^{41}$還有這個節日的故事對我們今天又有什麼啟示.
神給了我們一個很重要的諾言.
就是記載在詩篇裡面.
祂說:你留心聽我的禱告.
誰聽我懇求的聲音.
我在患難那天求告你.
你就必應運我.
在這裡我們看到有一個方程式.
就是我們在患難裡面求告神.
神就會應運.
這是一個非常強力的應許.
給我們信心.
我們現在看見這個世界好像很亂.
好像在患難當中水深火熱.
不知道怎麼樣.
但是神的應許這個諾言是永遠存在的.
在我們的心裡面.
是令我們不用怕.
是可以很安心的面對各種的挑戰.
希望大家可以絕對的記住神的這個應許.
不過這個應許裡面又有一個很有趣的問題.
就是神怎麼應運我們.
我們不妥了.
不行了.
或者我們生病了.
或者我們沒錢.
被人拋棄了.
或者我們有什麼法律的問題.
那麼神怎麼可以拯救應運我們呢.
這個才是所謂的64 million dollar questions.
不是神會不會應運.
而是神怎麼樣應運.
在華人教會裡面.
我們很常見的有一個救贖觀.
或者一個祈求的觀念.
就是沙灘上的足跡這個圖文.
這個故事或者這個祈禱.
可能都很熟悉.
就是有一個人.
他在死了之後.

$^{81}$他就回看他生命裡面.
原來在沙灘上走的.
他就發覺主耶穌都一直陪著他走的.
所以在沙灘上就看到有兩個人的腳印.
但是他又看到.
他不碰很有困難的時候.
就有一對腳印.
於是他就問神.
他說主耶穌.
為什麼我不妥的時候.
你就不跟我走.
我自己在這裡走.
主耶穌就說不是.
你不妥.
你有困難的時候.
是我抱著你走.
不是你自己走.
這個故事是一個很感人的故事.
很引導人.
也都令我們覺得.
原來我不碰的時候.
神就會抱著我走.
我不用做的.
我不用擔心的.
神會抱著我.
我不用.
我給他就行了.
我靠他就行了.
這是一個很感動.
令我們很安心的故事.
但是我要跟大家說的是.
這是一個偏差了的救贖觀.
神不是一定這樣救我們的.
神可不可以這樣救我們呢.
是不是我們有困難的時候.
他就抱著我們.
不用我們出力.
我們什麼都不用管.
他抱著我們走.
我們就抱著他這樣就行了.

$^{121}$我有幾個孫子.
我的孫子很喜歡抱著我們.
有什麼事的時候.
你抱著他.
那就行了.
但是神是不是這樣.
我們有問題的時候.
就讓我們抱著他抱著我們走呢.
當然可以.
神可以這樣做.
絕對可以.
譬如有一個例子.
就是列王記下六章24到7章20節的例子.
那裡說到在以利沙的時候.
亞蘭王攻擊撒瑪利亞.
突然間神就見到撒瑪利亞的軍隊聽見有聲音.
很害怕.
就突然走了.
那些以色列人什麼都不用做.
以利沙就祈禱.
不用出去打.
那些人就走了.
這個故事裡面還有很多意思的.
就是有四個大麻風和軍長的故事.
今天不詳細說了.
大家可以回去看一下.
這個大故事裡面有一個小故事.
說到四個大麻風可以怎麼樣呢.
是得到神的恩典.
而這個不信以利沙講的軍長又怎麼樣.
看見有得吃了.
本來他們機方沒得吃的.
但是看見那些亞蘭人突然間走了.
行了.
但是他們就有得看沒得吃.
結果就死了.
神可不可以這樣教育我們.
什麼都不用做.
抱著我們走呢.
絕對可以.

$^{161}$但是.
神是不是只可以這樣做呢.
我就會說不是.
這個故事本身沒錯.
但是如果我們只堅持神一定只可以這樣救我們.
當我們有困難的時候.
我們就什麼都不用做了.
就是在求神.
神就會救我們.
這個是我們心態上的偏差.
因為如果限了神可以這樣救我們的話.
我想問的就是.
誰是神呢.
你當神是這個Genie.
你拆了神燈.
他就會出來救你.
或者你當神是有求必應的王大仙.
什麼菩薩.
還是這個是我們的主.
我們的神呢.
今天我就想和大家在以斯帖的故事裡面.
這個普爾節的故事裡面.
看一下神怎麼樣去救贖.
以斯帖這個故事是一個很有趣的故事.
剛才大家介紹說有這個課程是說聖經裡面的敘事.
成本以斯帖其實就是一個長的故事.
一氣呵成的故事.
它的地點是波斯帝國.
是當時世界上最強的帝國.
它的地點是包括從非洲去到印度.
在中東一帶.
時間是約莫在公元前480年.
那時候如果你和聖經的舊約那裡看.
以色列人已經是一部分回去了.
是過了70年的被擄.
第二次聖殿又建好了.
但是還有很多以色列人猶太人是散居在各地的.
而以斯帖記的時間就差不多和聖經最後的歷史記錄.
好像以斯拉,尼希米是差不多時間寫成的書.
今天我們就可以在這個故事裡面.

$^{201}$看看神用另外一個方法去救贖.
和今天對我們有什麼意思.
以斯帖記是一本很特別的書.
如果你讀聖經你會發覺.
可能有一個問題就是.
為什麼聖經有本這樣的書.
它是聖經裡面唯一.
全本書沒有提過神這個字.
沒有說過有祈禱.
在故事裡面沒有提過祈禱這回事.
沒有提過祈禱這回事.
沒有任何的神蹟.
只不過是一個劇本.
如果你拍戲它就好像一齣戲.
一個劇本.
沒有任何的金句.
問題是為什麼這本書會放在聖經裡面.
成為聖經的正典呢.
在猶太人裡面他們為什麼覺得這本書.
要放在舊約聖經裡面.
放在他們的聖經裡面.
放在他們不一定要讀的書裡面呢.
在今天經過這麼多次的宗教改革.
為什麼這本書還在呢.
這本書沒有提過神.
甚至連神這個字都沒有提過.
沒有提過祈禱.
沒有提過神蹟.
為什麼還在聖經裡面呢.
因為這本書.
它提到的是神怎樣去拯救.
是不是有一張.
可能有一張徽章不見了.
OK 不要緊.
這本書之所以成為正典.
現在的解經家有一個分析就是.
這本書提到的是一個Paradigm Shift.
在舊約裡面一個規範轉移.
就是由神直接的介入.
轉到神在我們生活裡面介入.

$^{241}$你看舊約以前的神蹟.
以色列人的歷史.
神每一件事都是直接介入的.
好像剛才說那個.
以利沙的時候那個例子.
他們被人攻擊了.
神直接介入.
發生了神蹟.
於是以色列人得救.
但是以斯帖記沒有這樣的東西.
神不再是直接用神蹟來介入人的歷史.
而是.
在日常生活裡面.
在生活裡面.
是介入我們生活的每一部份.
在那裡給我們帶領.
給我們保障.
這個故事就說.
神怎樣拯救以色列人.
怎樣去發生不同的事.
首先就是.
在故事裡面.
很有趣.
神不是叫敵人突然死光就算了.
他是說.
那時候的皇帝不能取消.
他之前屠殺以色列人的預旨.
於是.
他就給出一條新的聖旨.
就說現在給你以色列人可以自衛.
猶太人就組織自衛隊.
反制敵人.
今天因為時間問題.
不能把整個故事說完.
所以希望大家回去.
把以斯帖記從頭到尾一次過.
好像看小說一樣看完.
以斯帖記不能平時查經讀一兩節.
一定要從頭到尾一氣呵成的看完.
因為是一個故事.

$^{281}$整本書就是一個故事.
你一定要從頭到尾.
好像追小說一樣追完.
你才知道它發生什麼事.
在事件最後.
猶太人慶祝他們得救.
慶祝這個普爾節直到今天.
這個故事.
大家回去看.
就說到以斯帖做了波斯的皇后.
她的養父或者堂兄就得罪了奸臣.
奸臣哈曼就想殺死猶太人.
以斯帖就冒死去見皇帝.
但是他沒有求皇帝什麼.
他只是叫皇帝陪他吃飯.
在兩次吃飯裡面.
有些很奇怪的事情發生了.
令到這個奸臣死了.
主死了.
猶太人就得救了.
詳細的事情我等大家自己回去看.
我不詳細說了.
我想說的就是.
在這個故事裡面有三個特徵.
對今天我們尋求神的救贖來說.
是有特別意思的.
第一個特徵.
在以斯帖裡面看到.
神提供了機會.
但是人有執行的責任.
跟我們剛才說的沙灘捉人的故事很不同.
在那個模式裡面.
我們希望什麼都不用做.
神抱著你走就行了.
但是在以斯帖裡面看到.
是不同的.
神給我們機會.
我們自己也要做事.
第二個我們看到一個很特別的拯救方法.
在這個救贖的故事裡面.

$^{321}$我們看不到有什麼英雄.
有什麼先知.
有什麼強人領袖.
第三個特徵就是.
這個拯救是在一個現有的機制裡面進行的.
不再是一些我們想不到的東西.
或者是一些神蹟.
是超自然的東西.
而是在一些現有的機制裡面進行的.
我們就看看這三點.
對我們今天有什麼應用的價值.
第一個特點就是.
神提供了機會.
人有執行的責任.
這次神沒有幫以色列人消滅敵人.
和剛才說的那個故事.
列王記下的故事比.
不同了.
列王記下裡面.
神叫敵人聽見聲音就走了.
還有其他很多這些神蹟.
之前的故事都看到.
神直接出手.
打倒了敵人.
以色列人就得救了.
包括了在出海及記的時候.
神降災直接殺了那些長子.
雨節.
以色列人什麼都不用做.
坐在家裡.
吃一隻羊.
塗點血就行了.
不用做事了.
但這次不同了.
神沒有幫以色列人消滅那些仇敵.
沒有任何超自然的神蹟.
如果我們在這裡祈禱等神蹟.
可能他們已經被人殺光了.
神給了我們九個月的時間去準備的.
如果你看這個故事.

$^{361}$我們得到這個令羽說.
你們可以組織自衛隊.
你們是要自衛.
他們有九個時間去準備的.
怎麼準備.
他們要計劃.
有人說九個月之後.
我要把你們殺光.
你有九個月時間去預備.
你怎麼預備.
你當然要想想.
如果要打過來.
我怎麼樣.
我的屋子怎麼樣.
可能弄強一點.
或者買定一些棍.
磨好刀.
要練拳.
健身.
你要自衛.
你不是手無筷子之力.
你怎麼自衛.
你當然要操練一下.
跑步.
舉重.
儲點糧食.
預算人家要來打你.
你當然要有準備.
至少要自衛打回來.
到了時候.
你需要自己執行這個自衛拳.
神不會幫你一雷打下去.
拿刀衝過來的時候.
把你打掉.
你要自己出去打.
今天我們遇到困難的時候.
我們又想要什麼拯救.
我們是不是想求神.
求你把這個困難拿走.
還是我們一樣要準備.

$^{401}$我們想要一個神跡的見證.
我有病了.
於是我就求神.
叫我一個病.
不用看醫生.
不用吃藥.
我都沒有了.
我沒有工作.
於是叫神.
突然間有人三顧草勞.
敲你門.
什麼生.
什麼小姐.
我們很聽過你的大名.
知道你很厲害.
從今禮聘你出山幫我們做事.
你自己不用找工作.
不用發履歷出去.
不用工作.
我們是不是把神.
當一個有求必應的Genie.
一個燈神.
擦亮燈神就出來搞定.
自己什麼都不用做.
我們怎麼準備.
如果你失去工作.
你是不是坐在那裡.
我失去工作了.
神啊,你幫我找工作吧.
你需要讀書嗎.
寫好你的履歷表.
看看外面哪家公司有什麼.
你可能是想要男朋友.
要女朋友.
你很寂寞.
是不是坐在家裡做宅男宅女.
只顧著玩遊戲.
你會怎樣呢.
需要自己學習.
怎樣去做好你的儀容.

$^{441}$扮好一點.
說話的時候說得合適一點.
還是坐在那裡等神來.
當下幫你做個女人.
放在你身邊.
你一醒來就行了.
我們自己要準備.
神會供應.
但是我們也有執行的責任.
這是第一個在以色列的故事裡.
我們看到的模式.
神給我們機會.
但是我們自己要做事.
第二個以色列的模式是.
我們再也見不到英雄.
見不到什麼先知.
見不到什麼強人領袖.
以前以色列人來看歷史.
看到的模式就是.
一個英雄,一個勇士,一個先知.
或者一個君王.
他帶著人去打仗.
你看時事記,你看武爾記,你看列王記.
歷代志上,歷代志下.
你也看到這樣的模式.
就是一個強人,一個領袖.
帶著一群人出去打.
他打贏就贏.
他打輸就輸.
但是在以斯帖記裡.
我們見不到這個模式.
這次我們見不到一個領袖帶著他們去打仗.
你再看回那個故事.
以斯帖和末底改.
似乎完了之後就沒有出聲了.
沒有再做事了.
故事去到第六章,第七章之後.
他們消失了.
以斯帖見到皇帝.
皇帝說以色列人可以自衛.

$^{481}$然後沒有了.
他們消失了.
再也見不到他們說什麼了.
沒有英雄,沒有領導.
沒有一個領導.
當時的以色列人.
他們散居在127個省.
從印度直到古實.
古實就是現在的埃塞俄比亞.
是非洲.
印度那時候的波斯帝國.
就去到印度的邊緣.
就是現在巴基斯坦.
差不多的地方.
有沙漠,有森林,有海,有山,有湖.
全部都散了.
然後他們就去到印度.
有沙漠,有海,有山,有湖.
全部都散了.
你沒有辦法一個人.
當時沒有斯大林.
伊隆·馬斯克那些衛星飛上去.
你沒有可能將信息傳到不同的地方.
沒有可能一個領袖教以色列人怎麼打.
他們唯有自己每一個地方.
想想我們自己在這個地方遇到問題.
我們怎麼解決.
沒有辦法一個人指揮.
也需要因為他們的環境和安排.
來決定怎麼去自衛.
今天對我們有什麼意思呢.
我覺得這個模式對今天教會來說.
是更加適用的.
就是我們再沒有處境化.
one size fit all解決問題的方法.
對不起.
我們是再沒有one size fit all解決問題的方法.
我們是需要每一件事都處境化.
根據當時當地的環境來做事.
在教會如果我們看回過去十幾年的歷史.

$^{521}$我們發覺很有趣的.
十幾二十年前.
很多教會是一窩蜂去跟這套那套的.
跟Wyck Warren.
Purpose Driven Church.
每個人都買一本書.
跟著芝加哥.
跟著Willow Creed.
每個人都去跟他們.
去他們的教會.
跟著Hill Song也去跟.
現在這些Mega Church.
我們看到很多都是沒落.
我不說神學上有什麼問題.
只不過看到當大家一窩蜂去跟一套東西的時候.
原來發覺不行.
我們每一個人每一個教會.
都需要跟著自己那時候的處境.
來解決問題.
一個很重要的概念就是.
信徒皆濟私.
我們每個人都是直接跟神有關係.
我們不用特別跟著一些強人.
一些領袖說什麼就做什麼.
絕對不是叫大家不聽牧師.
不聽勵仙他們說話.
我們絕對尊重他們.
他們的經驗.
他們的教導都很有價值.
但是我們就不要去崇拜任何的強人.
特別是我們接觸不到那些.
外面的那些.
那些電報評論員.
不知道說什麼.
走出來好像很有型.
其實背後很多骯髒的東西我們不知道.
我們大家要小心.
我們每一個人都可以直接跟神有一個關係.
跟神那裡得到力量.
得到接觸.

$^{561}$我們不要迷信什麼明星講完.
什麼宗教領袖.
我們需要根據個人背景.
我們每一個人自己的背景.
當我們遇到問題的時候.
神會拯救我們.
但是神會根據我們的背景.
來發展出拯救的方法.
說個秘密給大家知道.
我其實很多時候都是在大家教會裡崇拜.
雖然是在網上.
所以我知道大家幾個禮拜前.
有一個MBTI的workshop.
那時候我是很衝動的話.
如果我來得到就好了.
為什麼呢.
MBTI對我來說是特別有意思的.
可能大家有些人會知道.
我做醫生之前我是讀Information Systems.
我在USLA那裡讀IT的.
IS的PhD.
我的博士論文就是用MBTI來研究.
怎樣去訓練一些Managers.
去用那時候的軟件.
所以MBTI對我來說.
那個意思是非常重要的.
看到大家已經過了40年了.
這個工具還在.
看到大家願意在MBTI裡面.
明白自己那個.
對我來說那個叫Cognitive Style.
就是你怎樣處理問題的方法.
是非常重要的.
因為你不同的Cognitive Style.
你就有不同處理問題的方法.
我處理問題的方法.
和你處理問題的方法可能很不同.
絕對沒有辦法說誰對誰錯.
你要根據你的Personality.
來發展出來處理問題的方法.

$^{601}$所以很高興看到你教會會做這件事.
第三個神拯救的方法就是.
他會在現在的機制或者基礎建設.
建設的Foundation裡面進行.
他不再是神蹟突然在上面.
走過來砍了敵人就算了.
而是他在一個機制裡面.
現有的存在機制裡面.
來發展出拯救方法.
在拯救以色列人的時候.
是在當時的法律框架裡面進行的.
不是超自然的.
很有趣的就是.
為什麼說回這件事.
就是那時候為什麼神就這麼說.
本來說要殺死那些以色列人.
不殺了.
我也不知道為什麼.
但是原來那時候波斯有一個.
有一條法律就是說.
皇帝用戒指吸了印的東西.
就不可以取消的了.
皇帝第一次就說殺死那些以色列人.
他就不可以說現在不殺了就算了.
於是他就要出過第二條的命令.
就是說那些人可能還會殺你.
還會打你.
但是我准許你們自衛.
所以他在現有的機制法律基礎裡面.
幫以色列人找到一條出路.
今天神也可以用同樣的方法.
來拯救我們.
他會用已經存在的機制的基礎來拯救我們.
包括了什麼呢.
科學,醫學.
如果你有病.
我們不一定要求神.
你不用吃藥.
不用看醫生.
你就好了.

$^{641}$神可以這樣醫治.
但是神不需要這樣醫治.
如果神通過醫生.
通過藥物.
通過科學來拯救我們.
絕對依然是祂的掌管.
是祂的神蹟.
神可以通過法律來拯救我們.
如果你遇到什麼麻煩.
神可以通過現在有的法律來拯救我們.
神可以通過制度來拯救我們.
已經有的制度.
如果我們被人欺負.
我們去投訴.
這個不是什麼不屬靈.
這個是神通過現有的法律來幫助我們.
說一個好笑的例子給我們知道.
給大家聽聽一個故事.
真人真事在我身上發生的.
很多年前有一次和一個弟兄吃飯.
這個弟兄很慘和我投訴.
他在加州被人撞車.
撞車的人沒有保險.
Uninsured motorist.
我朋友的車被偷走了.
他人受傷了.
看了看身體療法.
但是對方不給錢.
於是他沒有錢給就被人追債.
本來他買了Uninsured motorist的保險.
他的保險公司會賠償.
但是他發現一間教會.
一間很.
我不可以說是邪教.
但是很極端的教會.
就是告訴他.
加州不可以打官司.
你不可以去告人.
你告人就不符合聖經.
他就很慘.

$^{681}$他的保險公司就說.
你要我賠償給你沒問題.
但是你要先申請一個聲明去告那個人.
要先申請一個聲明去告那個人.
那個人不賠償.
我們那個Uninsured motorist才會接受.
會賠償給你.
這樣也很合理.
保險公司有他們做事的方法.
不是你說人家不賠償.
你問他拿錢就不行.
我的朋友就很慘.
他的目標就是告訴他.
他不可以去打官司.
不可以去告人.
他的車壞了.
沒錢買車.
他就被人追去.
chiropractor還是physical therapist.
那些錢.
於是他就走來.
我也是剛好和他吃飯.
他就在我面前發呆.
我拗了頭.
我跟他說不是.
聖經說不可打官司.
保羅說不可以打官司.
這件事好像這個情況.
不是很適合.
不是很適用.
保羅自己也打官司.
保羅被人告的時候.
上訴該殺.
怎麼不會打官司.
結果怎樣呢.
我朋友就解決了問題.
但解決問題的方法是什麼呢.
轉教會.
這個就是神.
是用現有的機制.

$^{721}$法律.
制度.
來幫助他解決問題.
還有一樣東西.
朋友.
神可能通過朋友.
幫助我們解決問題.
總結就是.
另一個故事和大家說.
有一個爸爸和兒子說.
兒子你不會做功課.
不用怕.
我幫你.
於是有一天.
兒子就和他說.
爸爸我不懂做功課了.
你幫我做我就去打機.
爸爸就拗頭說.
為什麼你不做功課.
叫我幫你做.
你去打機.
兒子就說.
你說幫我.
你幫我.
你幫我做.
我就去打機.
我們今天是不是這樣要神呢.
神幫我們.
於是我們就什麼都不做.
我們就去打機.
他就抱著我們在山上走.
我們就抱著他.
很舒服的就算了.
如果你做爸爸的.
你是不是永遠都想幫你兒子做功課.
解決問題呢.
還是你只是想給他資源.
給他支持.
讓他學怎樣解決問題呢.
到底我們想做巨嬰.

$^{761}$現在這個詞語很受歡迎.
巨嬰.
我們想做神的巨嬰.
神抱著我們.
我們咬著他的舌頭.
不用做.
還是我們想神牽著我們.
在沙灘裡面走呢.
到底我們作為基督徒.
我們希望得到的拯救.
是怎樣呢.
到底天父想我們怎樣呢.
這個我就希望大家可以想一想這個問題.
多謝大家.
謝謝大家.
\newpage



\section{詩篇}
\label{sec:_lnAtpZRm3Q}
\textbf{Holistic Worship (Psalm 150) - Rev. Mark Chan 全人的敬拜 (詩篇 150) - 陳崇基牧師}
\newline
\newline
連結: \href{https://youtube.com/watch?v=_lnAtpZRm3Q}{\texttt{https://youtube.com/watch?v=\_lnAtpZRm3Q}} ~~~~ 語音日期: 2026-03-09
\newline
\newline
\hyperref[sec:mjg2Vzt9lP0]{\small{< < < PREV SERMON < < <}}
~
\hyperref[sec:index_chronic]{\small{[返順時目]}}
~
\hyperref[sec:index_scriptual]{\small{[返順卷目]}}
~
\hyperref[sec:y65Q_8wRaJ8]{\small{> > > NEXT SERMON > > >}}
\newline
\newline
$^{1}$Good morning, everyone, and it's always good to worship together.
I want to give a hand to Amanda for doing an excellent job in translating and uh, let's give her a hand.
Uh, it's always good to worship together.
So I want you to say hi to your neighbor and say it's good to see you. Okay.
Okay, whether you know him or know her or not say hi, you know, it's good to see you.
Yes.
Okay, good to see you guys good to see you guys and uh today I would like to end our series on psalms and we.
end up in psalm 150.
Well, let me read this psalm it's very short, uh, it says praise the lord praise god in his sanctuary.
Praise him in his mighty heavens praise him for his acts of power.
Praise him for his surpassing greatness.
Praise him with the sounding of the trumpet praise him with the harp and lyre.
Yue qi han xiao de shen yin zan mei ta.
And verse four to six it says praise him with tambourine.
And dancing and praise him with the strings and flute praise him with the clash of symbols praise him with resounding symbols.
Let everything that has breath praise the lord praise the lord now this psalm is all about worship.
So it asks the first question is what is worship?.
And at the basic level the very very basic level worship is human your response to the only god who is truly.
worthy of your honor.
So again when you walk in here you said oh I am joining the church worship.
This is my response to god what he has done already.
So so so you think about god god is infinite.
God is infinite and the worship god deserves.
Will never end and it continues into eternity.
So when you think about you know as we worship together, uh in the bible, we naturally come to the book of psalm, right?.
Because the book of psalm is collected.
Uh songs and poems by the people of god in the old testament when they returned from exile and there were 150 of them.
So as you look at the whole book of psalm look at the arrangement of it.
Um, you find that it's very interesting because the psalms.
All 150 of them are arranged in five different books and each book ends with something similar.
It's a closing doxology.
So as you can see on the left, you have psalm one and two that is the gateway to the entire book.
And then as I give you example, uh at the end of book one, which is psalm one to 41.
It has the it's the d.
That is doxology. Okay. I don't have enough room for doxology.
But you know it's end with doxology, which is uh verses 14 of psalm 41.
It says praise be to lord the god of israel from everlasting to everlasting. Amen. And amen.
Amen amen and then uh, it is so.
In that the end of book two and book three and book four and book five.
There was an ending doxology.

$^{41}$And then, uh, you know, um, I as I prepared this sermon I thought of uh when I was in seminary.
Uh for my master of divinity degree, uh preparing myself for ministry.
Uh, as I thought about it, it was already 30 years ago.
I was in seminary and I was uh.
with classmates and sometimes uh, we have different.
uh brothers and sisters from other church or tradition or denominations and find kind of interesting conversation with them because some of them.
Uh are extremely conservative about worship.
So.
And when we talk about worship.
The debate get intense sometimes and because some of them are from very conservative.
Tradition they they they believe and insist that the songs we that we use in worship service must come from here.
There's 155 songs.
You know this congregation is relatively young relatively young, okay, I.
I think i'm the one the oldest guy here. Anyway, uh, you know.
I want to ask uh, whether you have been to a church where you walk into the sanctuary to join the worship.
There's a hymn book right in front of you.
Have you ever been to church like that?.
Anyone traditional church, there's a hymn book right in front of you and then uh as worship starts you pick up the hymn book and the host will say, okay, we turn to hymn number.
116 and then everybody turned to the 116 and then we sing that song together.
Have you ever have any experience like that? No, she's so young. Okay, some older guys.
Some older guys some older people.
Um, you've been to church the hymn book is someone who put together some of the songs used by the community and put it in one book.
And then uh, the book of psalm is kind of like that, you know.
The people of god return from exile and they collect 150 of them and the words in psalm 1 to.
150 will be set to different melodies and they sing and they worship with those songs.
And then some of the students that I talked to 30 years ago from a very very conservative background.
They would insist that we can only sing one of the 150.
Psalms we find in the bible and nothing else.
So they think that you know, there are other songs that are written, you know, uh recently.
They cannot be used. You cannot use those songs in the biblical worship.
Any modern songs? Okay any modern songs like, um, you know, uh, like amazing grace. All right, so you are not.
Allowed to use any of those songs.
We we we sang some songs in the beginning right and you cannot use those for them you cannot use those.
Now in their understanding truly biblical worship uses only the 150 songs that god revealed in the bible.
But as you think about it, it's kind of interesting, right.
It's kind of restricted.
But the bible tells us to sing new songs, right.
Even in the book songs sing new songs to god, right? So how do we apply that?.
So you have to understand how to read the bible how to interpret it scripture describes.
The worship in the ancient days, but it does not prescribe every detail of how worship must look like.

$^{81}$So as we turn to uh, psalm 150.
It is the concluding doxology of the entire book.
We want to find out what kind of questions the psalm is trying to answer.
What kind of teaching can we find find in the psalm?.
So the first point here is that this psalm tells us who is to be praised who is to be worshiped.
Verse one it says praise the lord praise god in his sanctuary.
And praise him in his mighty heavens.
See the old testament the sanctuary is of course is jerusalem temple and the temple is the throne of god on earth.
And it's the place where he sits there and he rules the entire world. He reigns as a king.
that's.
So praise god in his sanctuary in his.
Dwelling place and praise him in uh, the mighty water.
Mighty heavens now mighty heavens point to the power and the creation of god because in the ancient east when they imagine the world.
The world is uh is surrounded by water and whoever can hold up the sky.
with.
with the heaven.
Now create this bubble for us to to live in uh to restrain the water and that is the one we worship.
So god is to be worshiped and god is greater than everything, you know.
This is this verse tells you god is greater than you.
god is greater than.
The entire universe and god is greater than you know your whole life even including your ups and your downs.
And then verse two it continues to tell you.
Why should we worship.
Now verse two says praise him for his acts of power praise him for his surpassing greatness.
Now as you pay attention to the psalm is a very short psalm.
Psalm 150 begins and ends in the same way praise the lord right praise the lord and in biblical hebrew.
In the original language, it is hallelujah. And hallelujah simply means it's not english word. By the way, it's a it's a hebrew word.
Hallelujah means hallelu plus. Yeah.
hallelu means.
You have to praise right you have to praise guys praise and plus. Yeah, this is the short form of yahweh.
the the name of god.
Hallelujah.
In the middle of this psalm 150 there were 10 commands to praise.
10 commands for us to praise and worship is a call to move.
Our attention at this time move your attention away from your life from ourselves and turn our attention back to god.
Is.
So I want to ask you when you come in today.
When you come in today, what kind of a question did you ask yourself, right?.
Psalm 150 asks you invites you.
To ask the most important question is why why do we come here?.

$^{121}$Is it out of duty? Is it out of routine that we come here or what is the reason? Why are you here?.
So many of the praise psalms in the bible.
Uh, we'll we'll use this phrase praise the lord and worship god because.
So psalm 150 gives two reasons here.
His mighty acts.
And then his greatness is surpassing greatness.
And when you see.
The creation when you appreciate the creation you naturally will respond with all in wonder and you praise the creator.
And think about the wonder of the whole universe is so complicated and yet beautiful and you praise god.
And then praise him for because of his.
Surpassing greatness and that points to the salvation and the love the grace of god.
And uh, if you look at the whole book of psalms you find that the whole book of psalm tells you one big story of.
God's grace and god's salvation.
So first of all from this diagram from left to right 100.
First song and then the last one is 150 across five books.
There was a movement.
from lament.
Sad songs to praise at the end and about one-third of the psalms of this 150.
Laments and this is the cry to god when they suffer when the people of god suffer.
And then the whole entire.
Turning point. It's kind of like a u-turn when you drive, right?.
You wrong direction and you make a u-turn. Um, the the turning point of the entire book of psalm is psalm.
89 right 89 and it's a painful lament because the uh, because of the fall of the kingdom of david.
In the you remember in the old testament god promised david king david an eternal kingdom that will never end and yet.
Uh, you know, you know the history of the old testament it however it end.
And collapsed the entire kingdom and the people went into exile.
And then uh when you read psalm 89, uh, I give you some example here like first 38.
It's a lament crying out to god because we're suffering we lost our kingdom.
It says you have rejected god god.
You have spurred us and you have been very angry because of our sin you have very you have been very.
Angry with your anointed one, which is uh, the king of the kingdom.
Yeah.
So when everything is gone, uh, the psalm cries out to god, uh,.
Another example another verse is 49. It says oh lord.
Where is your former great love your faithfulness which is in which in your faithfulness is swore to david.
There's an eternal kingdom, but what happened now?.
49.
So psalm 89 is kind of like a minor song, you know, it's a very sad song.
That is the end of book one two and three.
And then you come to psalm 90 which is the beginning of.

$^{161}$Book four and that is interestingly a prayer of moses and kind of restarts the story again.
And psalm 90 i'll give you some example verses 13 to 14.
Is it continues to cry out to god but look up to god and hope.
You know being set in god, how long will it be god? How long it will be that we suffer that we're going going through our down times.
Satisfy us god in the morning with your unfailing love that we may sing for joy.
And we may be glad all of our days.
Your love.
Isn't that also our life experience.
You know, especially from our context, you know, we come to this place.
And sometimes most of the time we suffer.
Whether you have suffered some illness or some difficulties or some challenges in your life.
So this psalm says, you know god.
God's people there's hope in him because uh the faithfulness of god.
And then uh in the entire scheme of things.
You think of the resurrection of jesus christ our lord.
He has guaranteed the victory of god for the people of god. And so even when life seems to fall apart.
We still have hope in him because of the resurrection of jesus.
As the apostle says in the new testament is it's not because we are faithful that we have hope it is because god is faithful.
So we have hope.
And that's why the entire book of psalm ends with what.
not lament.
But praises right you get more praise psalms and more praise psalms at the end.
And until the loudest praise is in psalm 150 and god alone is to be praised.
And that's why god is worthy of our praise because he deserves all the honor and give us hope.
And then the third point here teaches us who.
How should we praise him?.
And you remember the psalms uh includes a lot of different music, uh instruments, right?.
Uh trumpets, uh harps and tambourines strings and flutes and cymbals.
And then, um, you group them together and there are three types of basic instruments, right there are wind instrument.
There was also string instrument. There were also percussion instrument.
and these were the.
uh instruments.
of.
the ancient culture the ancient time.
The point is not to limit our worship to use only.
These sort of musical ancient instruments. The point is that you know, the time is different.
Right time is now different and we use different instruments.
And um the more important point is that the psalm tells us that true worship is full of music.
And it's alive and full of emotions and it even includes dancing.
Right. Do you dance when you when you sing?.

$^{201}$today.
No, you don't right?.
That's our our cultural, uh, you know.
Interesting.
And and this song says you should dance when you sing praises.
See even when you move your feet it's kind of like dancing, right?.
And then your hands move you claps and your feet steps, you know on the side and some people would uh,.
You know lift up their hands and some will stand and the actions express.
Our praises including joy and respect and joy near to god.
And you remember jesus said that uh, all the disciples should love god with all of our heart.
all of our soul mind and.
Strength right strength, which is your physical body. So worship involves your whole.
holistic person.
And I understand that you know, sometimes most of us are in the from the uh,.
asian culture.
You know, uh, our churches can be too serious.
uh.
And too limited, uh by our culture.
But here it tells you worship is a verb and even your body is involved and your heart is involved.
Okay, so the last verse it says, uh who should worship who should praise.
Women.
And verse 60 says let everything that has breath praise the lord praise the lord.
So better instruments better song systems, you know.
Our song system is pretty good and they're not enough the psalm says you need one thing.
Which is everything that has breath and who is it?.
It's you.
It's us all of us has breath.
See you're breathing, right?.
Your breath is the most precious gift of life god has given you and.
And god created us and he breathed into us the breath of life and the breath of life has its purpose.
As someone who has breath of life we are made to praise him to to worship him.
So, how are you using your breath today how using the breath that god has given you, you know, I pray that our.
Worship in the church to show the world show the world that god is worthy of our breath and our praise.
You.
So I hope this song will remind you what worship is all about and to invite you to truly worship him.
Shall we stand and pray together in conclusion?.
Let's all stand together.
Let's pray.
Dear lord, we praise you because you are worthy.
Thank you for jesus christ who gave his body gave his love and shed his blood for our sins.

$^{241}$And rose again to guarantee our victory.
For us as your people.
And today as we continue to worship you we come to your table.
And we bring you our sadness our defeat and our disappointment.
We ask you god to renew our hope.
And cleanse us cleanse our sins by your grace and strengthen us by your spirit and help us to trust in your faithfulness.
Women.
And as we partake the uh.
chameleon elements.
We remember your body and we remember your blood with gratitude.
Until jesus comes back again.
In jesus name we pray amen.
\newpage



\section{詩篇}
\label{sec:y65Q_8wRaJ8}
\textbf{全人的敬拜 (詩篇 150) - 陳崇基牧師}
\newline
\newline
連結: \href{https://youtube.com/watch?v=y65Q_8wRaJ8}{\texttt{https://youtube.com/watch?v=y65Q\_8wRaJ8}} ~~~~ 語音日期: 2026-03-09
\newline
\newline
\hyperref[sec:_lnAtpZRm3Q]{\small{< < < PREV SERMON < < <}}
~
\hyperref[sec:index_chronic]{\small{[返順時目]}}
~
\hyperref[sec:index_scriptual]{\small{[返順卷目]}}
~
\hyperref[sec:CVOZX9nsoIc]{\small{> > > NEXT SERMON > > >}}
\newline
\newline
$^{1}$大英姐妹大家好.
今晚我會飛去土耳其.
我會探索一個新的路線.
大家都可能聽過.
五月有大英姐妹會跟我一起去土耳其.
你猜過去兩天發生的事.
我會飛到土耳其嗎.
暫時未有影響.
暫時真的很難說.
這個世界真的發生太多事.
本來我打算恢復去以色列的聖地轉移.
暫時都未行.
但需要為那個地區禱告.
我想告訴大家.
很多時候有戰爭的時候.
我們都會被迫站邊.
看你幫誰.
其實不需要.
所有戰爭上帝都不喜歡.
因為要知道受害的人是甚麼人.
是無辜的人.
得益的是甚麼人.
是在上掌權的.
無論你是美伊這邊還是伊朗那邊.
哪邊都好.
我們不需要站邊.
因為我們屬於誰.
我們屬於天國.
對嗎.
屬於天國.
我們可以討論戰況,動態等等.
或者地緣政治等等.
不過我們是不需要站邊的.
因為我們屬於只有一位的國度.
一位天國的君主.
我們這樣去到風花雪月的時候.
其實已經出賣了我們自己的國度.
我們的角度.
因為我們是一些不會受到戰火摧殘的人.
對嗎.

$^{41}$我們就是那班袖手旁觀或者隔岸觀火的一班人.
所以我們不需要站邊.
也不需要亂拿聖經.
一些經文來對號入座.
這些是我們環教會.
傳統的教會我們經常有的弊病.
因為那裡有甚麼發生.
中東有甚麼發生.
我們就好像在聖經裡面說.
有一些預早寫好的佈置.
這個是非常不應該的.
應該去學一下怎樣去看聖經.
不要看錯聖經.
不要解錯聖經.
聖經沒有說這些東西.
所以大家不要弄錯這些東西.
但是我們需要去作和平之志.
OK.
那你說聖經這些怎麼解.
那你就要報讀我們的課程了.
就是這麼簡單.
祈禱和報讀課程.
我們先祈禱.
主我們知道.
每逢有戰火摧殘每一個地方的時候.
你的心是很痛的.
因為每一個人都是按照你自己的現象所做.
特別是一些無辜的人.
婦女 兒童 孩童.
他們是被迫在這些戰火的擴散裡面生活.
主我們為了這些人去禱告.
我們為了這兩天美以聯攻伊朗.
然後伊朗還擊.
引致到不只是單單一個國家的人受害.
而是整個地區 整個地帶.
整個中東的地帶.
主我們聽到很多傷亡的報導.
無論在杜拜 在一些Gulf States.
在伊朗 在以色列.
在不同的地區.

$^{81}$我們看到有很多被戰火摧殘的建築物.
有很多人的生命受到很大的影響.
甚至是被摧殘.
主我們求你憐憫.
憐憫這個地方.
因為裡面有很多很複雜的政治因素.
是一些近代政客.
或者不同的大國.
去到彼此去到爭拗角力.
以致有很多平民百姓是受到影響的.
主求你給我們一個淑齡的角度.
讓我們能夠成為和平之子.
縱使我們未必能夠身在那個地方.
但我們都為當地的居民去禱告.
主我們求你去洗這個戰士.
不會繼續灌溉 不會繼續複製.
因為你愛世人.
主我們求你都讓當地的人.
能夠很快複製到這場戰火.
以致人的生命不會受到摧殘.
主求你都教導我們.
在一個這麼混亂.
這麼多爭拗的一個世代裡面.
教導我們明白.
我們真正的效忠.
是上帝是自己的國.
真正的效忠是我們的主耶穌基督.
祂是那位和平之子.
並且將人與人之間和平的圍牆拆掉.
就是藉著祂十字架的救恩.
將任何的人.
所有的人 猶太人 愛邦人.
男與女 遺老的自主的人.
一切能夠在基督裡面能夠成人.
因為你是那位滿有恩典.
滿有慈愛的神.
並且你在這個亂世當中.
你仍然是工作人.
求你幫助我們.
讓我們在這麼混亂的世代裡面.

$^{121}$我們不是退縮.
我們更加願意以福音.
以我們的生命來見證你.
我們這樣祈禱 奉尊名禁.
好 今天我們繼續看詩篇.
我們需要完結這個詩篇的系列.
我今天選了詩篇最後一篇.
你記得詩篇有多少首詩歌嗎.
150首 是吧.
150首.
我們先看這首詩歌.
這就是整個詩集.
最後最後一首的詩歌.
看起來其實非常簡單.
但對我們今天是非常重要的.
第一節這裡這麼說.
「你們要讚美耶和華」.
「在上帝的聖所讚美他」.
「在他顯能力的窮倉讚美他」.
「要因他大能作為讚美他」.
「按著他極美的大德讚美他」.
「要用各星讚美他」.
「鼓瑟彈琴讚美他」.
「擊鼓跳舞讚美他」.
「詩而樂器燒的聲音讚美他」.
「大響的筆讚美他」.
「用高聲的筆讚美他」.
「凡有氣息的都要讚美耶和華」.
「你們要讚美耶和華」.
在這裡結束的時候.
詩篇告訴我們.
整本詩集裡面有一個這樣的結論.
就好像你看一套戲或者看一個連續劇.
最後最後有一個的大結局.
大結局.
雖然我們知道有很多時候.
華文或者港產的大結局.
是很難的.
但詩篇最後的大結局.
就是在這裡總結了整個詩篇.

$^{161}$到底想教導我們做些什麼.
就是什麼叫做敬拜.
What is worship.
如果你問崇拜,敬拜是一回什麼事呢.
這個問題是需要問的.
今天進來的時候.
今天你去教會.
為什麼每一個星期你騰空一些時間.
一些可能你生命裡面可以成為優先的時間.
但為什麼你不去做其他事情.
你去敬拜上帝呢.
那你就要問什麼叫做敬拜.
為什麼需要敬拜.
這些短短的六節.
作為整個詩集的結束.
就是嘗試去回答這個問題.
什麼叫做敬拜,什麼叫做worship.
簡單來說.
其實敬拜是什麼呢.
就是我們作為人.
去回應上帝的恩典和祂的偉大.
我們作出一個回應的時候.
這個就是一個敬拜.
明白嗎.
這個就是我們的response.
所以我們首先是被動的.
上帝做了一些事情創造我們.
創造這個世界.
一個美好的世界.
縱使有罪惡,有邪惡.
但上帝來到去救贖世人.
就是上帝首先做事.
然後我們作為一個被動的主體.
我們去回應祂.
既然上帝是一個無限的,存有.
infinite的being.
所以祂配受的敬拜.
也是不能夠停止的.
這個詩篇一直告訴我們.
就是一個這麼重要的課題.

$^{201}$我們來到敬拜的時候.
和神的子民一起敬拜的時候.
我們就去回應上帝給我們的恩典.
特別是詩集.
聖經裡面的詩篇.
有150首.
就是神的子民.
以色列人.
他們被流放.
然後他們回歸的時候.
他們去整輯出來.
在世世代代有的詩歌.
他們慢慢收集起來.
然後成為150篇.
如果你留意詩集.
其實是有一個佈局的.
或者有一個編輯的.
無論那個編輯是一個人.
或者一群人也好.
總之就是被流亡之後.
這些以色列人他們回歸了.
然後他們就去將.
所謂正典的舊約聖經.
編輯出來的時候.
詩集就是其中一本書.
你看到這裡有很多個像LEGO一樣.
其實是很簡單的.
整本詩篇150首.
其實是分開五本書的.
五卷書.
看到嗎.
有Book 1 to Book 5.
每一個.
1,2就是一個門口.
告訴我們我們敬拜上帝的時候.
有一個Gateway.
有一個入口.
就好像你進去一個豪宅的時候.
開門打開.
然後你脫下鞋.

$^{241}$然後你進去那間屋.
那個最堂皇的地方.
那個就是1和第2篇.
然後就分了五本書.
每本書最後的時候.
你看到Book 1最後的時候.
有一個D是讚美.
就會有一個讚美的.
所謂我們.
技術上我們叫它Doxology.
就是一個榮耀頌.
最後這個詩集.
Book 1收集了.
這些詩歌.
三至四十一篇.
然後最後就會有一節或者兩節.
一段.
就會有一個所謂Doxology.
一個榮耀頌.
他就看到四十一篇.
十四節.
就是耶和華以色列的上帝.
是應當稱頌的.
從耿古直到永遠.
阿們.
是嗎.
Book 1就結束了.
然後Book 2.
Book 3.
Book 4.
Book 5.
二三四五卷.
最後都會有一個.
類似這樣的.
所謂D.
Doxology.
完結的.
所以你看到.
譬如Book 3.
第三卷.

$^{281}$其實就是七十三至八十九.
八十九結束的時候.
就有一個這樣的Doxology.
一百五十篇在哪裡.
一百五十篇就算是.
最後一個的所謂.
如果你學音樂的話.
你就知道什麼叫Coda.
你知道什麼叫Coda嗎.
Coda.
大家都不學音樂的.
Coda.
就是最後最後的時候.
如果交響樂.
交響樂團.
那些樂團.
做出最多樂器.
所有樂器都走出來了.
這個就是一百四十六至一百五十篇.
我們剛剛看了.
就是一百第五十篇.
所以這些詩歌.
其實就是當時神的子民.
他們藉著不同的經歷.
有不同的創作.
他們就會收集起來.
成為今天我們有的.
舊約書卷裡面的詩篇.
好像我們現在少了.
現在沒有做這些事情.
以前我們那一代.
我們就有不同的詩集.
大家記得嗎.
那些比較成熟一點的信徒.
成熟一點.
你記得有《生命聖詩》.
什麼《青年聖歌》.
記得嗎.
《青年聖歌》有綠藍紅的.
記得嗎.

$^{321}$記得吧.
或者有什麼《眾主聖歌》.
有一本的.
為什麼要有那本呢.
以前那一代進到教會的座位後面.
有一個空隙.
有一個縫隙.
然後放一本東西在那裡.
進來敬拜的時候.
我們就拿起那本詩集.
然後那個主席就說.
我們今天唱第三首.
我們一起翻開第三首.
那裡有譜有詞.
我們一起唱.
那時候是之前有所謂的預測.
記得嗎.
記不記得有預測這回事.
記得嗎.
我只是說成熟的信徒.
你這些年輕人不懂的.
真是上帝祝福你.
以前我們有所謂.
還沒有那些膠紙的預測.
之前我們就用那些詩集.
打開.
真是打開.
揭開它.
然後慢慢科技就發展了.
就有那個overhead projector.
記得嗎.
那時候我們敬拜的時候.
有個人坐在前面.
他就有一張.
很多膠紙的.
第一節他就放上去.
第二節就放上去.
有時候會放錯的.
記得嗎.
有沒有做過這個位置.

$^{361}$沒有.
要寫上去.
有時候要印出來.
然後才有所謂的什麼.
PowerPoint.
開始有PowerPoint了.
是不是.
然後現在又有其他的科技.
我們又用其他的軟件等等.
你記得以前那一套.
就是有一個集體的一些詩集.
大家都認識的.
所以我們就會唱那些詩歌.
詩集其實是同一個過程.
來整理出來的.
我記得以前在神學院的時候.
以前讀M.Div的時候.
想起都三十年前了.
這麼老了.
三十年前.
那時候讀M.Div的時候.
有些同學.
有些同學非常著緊.
敬拜裡面唱什麼歌.
他們是那些所謂改革中的教會.
那種傳統的.
保守的改革中教會很有趣.
他們的詩集只有150首.
就是這150首.
你進去改革中傳統的教會.
你打開詩篇的時候.
詩本.
唱什麼歌.
只有150首.
一日復一日.
一年復一年.
都是唱那150首.
只不過是將它配樂而已.
配了樂.
因為他們認為什麼.

$^{401}$他們認為上帝所啟示的話語裡面.
給我們作敬拜.
其實就是只有這150首.
就不可以唱那些.
剛才Percy選的歌.
不行的.
他們覺得接受不了.
只能夠接受詩篇一至150篇.
然後就配不同的樂曲上去.
那詩琴就很重要了.
那詩琴就要彈.
還有一些更加保守的.
簡直不用鋼琴.
就有一些鍵盤.
只有八個鍵.
就是我們今天唱第三首.
然後我們走去按第一顆音.
主要是幫助大家不要走音.
就這樣清唱.
有些同學真的是這樣的.
我真的有一點拗頭.
就算我們傳統的詩歌.
奇異恩典都不合格.
因為奇異恩典不在這150首歌裡.
Yeah.
所以不能唱奇異恩典.
太慘了.
不能唱.
所以任何的歌都不合格的.
只能夠唱詩集裡面有的歌.
但我就有一點拗頭.
聖經裡面叫我們唱新歌.
何謂新歌.
就是剛剛我們唱的那幾首.
就是要唱新歌.
我覺得這個對聖經裡面的一些執著.
一些了解.
不是很對位.
聖經通常都是描述一件事.
Describe一件事.

$^{441}$它就沒有說去Prescribe你某一些東西.
有很多地方是有Prescribe.
就是會規定你做一些事情.
但是敬拜上面.
通常它只是描述當代的人.
他們用什麼詩歌.
或者用了什麼音樂來去敬拜.
就沒有去Dictate.
或者去決定.
到底世世代代用什麼音樂.
或者用什麼詩詞.
所有的詩歌都是神學的Application.
是神學的應用.
我們對上帝的了解.
對上帝的認識.
所謂Theology的神學.
我們應該怎麼應用在現在這個世代.
就要視乎你寫了一首詩詞.
寫了一首詩歌.
或者寫了一首文章.
然後唱出來.
所以聖經很少講到敬拜的時候.
會去決定我們怎麼做.
所以首先我告訴你.
你們很好笑.
有些很傳統的.
連鋼琴都不用.
因為他們覺得勝經沒有鋼琴.
那我就心想.
那你要怎麼樣.
我們就要.
好像耶穌的樣子出來.
是不是.
所以沒有說定死的.
但是是應用.
所以為什麼世世代代的神的子民.
是需要唱新的歌.
因為在這一代.
我們要回應這一代的需要.
在YouTube上最近幾年.

$^{481}$特別是疫情之後.
有一群年輕人做得很好.
他們現在不知道還有沒有做.
年尾的時候他們就會做.
跟我們一起經歷的.
在2024年才算.
2024年跟我們一起經歷的詩歌.
然後他們就把它們合起來.
唱在一起.
記不記得.
這些就是年輕的那些字.
他們將去年經常唱的詩歌.
或者他們覺得很感動的詩歌.
他們就把它們串燒地唱在一起.
唱出來.
為什麼呢.
特別是香港背景.
香港的年輕一代.
他們的信徒.
他們的經歷就是.
過去香港很多年很動蕩.
他們很多時候觸碰到他們.
使得他們心動的.
這些詩歌就是跟那些.
當時的社會的情況.
或者那個背景有關.
可能是一些哀歌.
可能是一些要抒發他們的感情.
或者對上帝的一些呼喊等等.
他們就將它們加在一起.
所以是有一些新歌是需要唱的.
We have to sing new songs.
詩篇第150篇.
這裡怎樣去告訴我們什麼是敬拜.
你說如果不是去決定我們.
用什麼音樂.
或者用什麼樂器.
用什麼詩歌來敬拜.
這裡是在說什麼呢.
是在說什麼有關敬拜的東西.

$^{521}$第一樣他說到就是.
誰是我們要讚美敬拜的.
第一次這裡說什麼.
他說你們要讚美耶和華.
在上帝的聖所讚美他.
在他顯能力的窮倉讚美他.
第一也說到聖所的時候.
說起耶路撒冷.
就是聖所是什麼.
聖所不是人去的地方.
在近東的世界裡面.
所有的聖所都是上帝居住的地方.
是上帝的家.
上帝在那裡是坐著為王掌權.
在地上他去管理整個的大地.
所以在聖所裡面去讚美他.
就是來到上帝的王權的面前.
我們去俯伏去敬拜.
在他顯能力的窮倉讚美他.
這個就是近東.
古代近東人怎麼看這個世界.
這個說了好幾次.
他們認為這個世界是.
就是上面有很多水.
然後誰能夠有一個窮倉.
窮倉這個字在希伯來文.
其實是一隻鑊.
一隻很硬的鑊.
然後有個鑊能夠頂住這些大水.
不會被它淹沒我們.
這個就是值得我們敬拜.
所以說他在他顯能力的窮倉讚美他.
就承認這位主比宇宙更大.
比我們人生更大.
是不是我們人生都會有高低.
但神比這一切都更大.
所以為什麼我們值得去讚美他.
所以這是第一個要回答的問題.
Who to worship.
第二個問題.

$^{561}$他想問的是為何要讚美他.
Why? What is the reason?.
為什麼要去讚美他.
我們走進來的時候.
很多時候我們來到教會.
我們都不問這個問題.
因為第一個問題就是問.
泊車泊在哪裡.
希望你們也問這個問題.
有沒有遲到.
是不是要趕路來.
通常我們不是問Why do we worship.
Why do we get together.
為什麼我們會走在一起.
這裡告訴我們.
Because因他大能的作為讚美他.
按著他極美的大德讚美他.
這篇詩篇一開始是說什麼.
就是說你們要讚美什麼.
其實是命令式.
希伯來文其實就是Hallelujah.
Hallelujah.
剛才有沒有說這句話.
唱這句話.
Hallelujah.
其實就是希伯來文.
所以你去教會.
你學了很多希伯來文.
Hallelujah.
你們要讚美什麼.
這裡在希伯來文.
你們要讚美什麼.
其實就是Hallelujah加Ya.
Hallelujah就是你們要讚美.
Ya就是Ya為的短寫.
所以你們要讚美什麼.
就是Hallelujah加Ya.
在這篇詩篇裡面.
十次的命令我們去讚美什麼.
就是將我們的Attention.

$^{601}$將我們的注視力.
由我們自己轉移到.
天上這位掌權的神.
這就是讚美敬拜的意思.
你進來的時候問Why.
這不是純粹責任.
或者刻板要去做的.
屬靈作業.
我們來到敬拜的時候.
我們需要問為什麼.
所以這裡說到要因什麼.
他的大能的作為.
就是上帝創造這個世界.
創造我們的生命.
我們值得去讚歎.
值得被震撼的.
大家都知道我最近幾個月.
我突然間很早起床.
我還是很早起床的.
我還是很早起床.
不知道為什麼.
Touch住那個按鈕.
過去幾個月我就是.
是啊.
就是人可能老了.
那我就五六點就起床.
那我就會開車去LA Fitness.
很多時候經過Lake Shore.
因為我去High Park.
在我開車去Lake Shore的時候.
就會看到日出.
大家有沒有見過.
在Lake Shore見日出.
看日出.
沒有很少.
我真的每次經過.
開車經過.
我都覺得很讚歎.
很驚訝.
因為每一天的日出都很不同.

$^{641}$而且是很漂亮.
有幾次我忍不住.
我想停下車拍照.
不過還是沒有.
因為趕時間.
但是欣賞創造.
欣賞上帝的大能.
這就是詩篇邀請我們做的.
在我們很忙碌的生活裡面.
匆匆忙忙裡面.
我們看看上帝自己的大能.
這是第一個原因.
第二個原因就是.
按照祂的極美的大德.
Supreme Goodness.
因為這樣的緣故.
來讚美祂.
什麼叫做極美的大德.
其實詩篇整個鋪排裡面.
已經說了.
上帝有什麼極美的大德.
剛才我們看了下面.
上面那兩個箭咀.
就是整本詩集告訴我們.
其實是一個故事.
150首的詩歌.
不是今天那些重主性歌.
你這樣拼湊在一起.
他們說的是一個故事.
整本詩集一至150篇.
其實是說一個故事.
開頭第三篇開始.
就是說到大衛的禱告.
因為他受苦.
在一個角度裡面.
上帝認許了他.
然後建立他的角度.
然後慢慢一直唱下去.
唱到第三篇最後的時候.
就是你看到第89篇.

$^{681}$就是大衛的角度衰落.
因為是流亡 亡國的時候.
所以89篇是一篇很慘的哀歌.
所以卷一二三都很多哀歌.
很多的lament.
因為神的子民受苦.
他們向上帝哀嘆去呼求的時候.
到89篇就更加更加黑暗.
因為89篇我給了一節給大家.
第38節.
是一個哀曲的禱告.
就是說上帝所建立的.
神的子民大衛的角度.
居然崩潰了.
沒有了.
89篇第38節.
但你的老奴.
就是上帝的老奴.
因為上帝的子民.
他們犯罪.
他們敗家.
然後整個的角度被毀滅.
因為你的老奴.
你的受高者.
就是君王,大衛等等.
當然最後是基督.
就拒絕了他.
然後離棄了他.
主啊.
你從前憑信實.
向大衛起誓的慈愛在哪裡.
所以你可以想像.
在一個人生的失敗裡面.
整個國度亡國的裡面.
他們去哀嘆.
去作哀歌.
問上天.
90篇就是神人摩西的禱告.
為什麼90篇會放在那裡.
是摩西的禱告.

$^{721}$因為你知道以色列整本的舊約.
所謂The Beginning of God's Grace.
就是什麼時候開始.
就是摩西.
摩西帶領以色列人出埃及.
然後以色列成為一個新的國度.
新的子民.
所以就將90篇這個.
就是神人摩西的禱告.
就放在《捲四》的開頭.
因為之前.
就一個哀嘆的國度沒有了.
但是到90篇.
摩西突然間好像上了舞台.
然後他就唱出第90篇的哀歌.
就唱什麼呢.
他說耶和華要到什麼時候呢.
使我們早早飽得利的慈愛.
好使我們一生一世能夠歡喜快樂.
這個就是摩西有的禱告.
帶領以色列人民第一次再去禱告.
所以你慢慢見到90篇開始《捲四》.
然後到《捲五》.
就是這個國度有一個復興的盼望.
因為上帝的信實.
不是因為人的信實.
其實你這樣看的時候.
就是整個是一個大的敘事.
整篇詩篇就是在說這件事.
就是上帝的恩典.
是人的成功失敗.
然後再給他一個所謂second chance.
讓他能夠有一個revival.
其實你看《柔弱》的故事.
都好像我們故事一樣.
我們人生的故事裡面.
經常都會經歷這些事.
是不是89篇.
我們哀嘆上帝的信實在哪裡呢.
我們要等到什麼時候呢.

$^{761}$我們受苦的時候上帝在哪裡呢.
可能我們很傷心,很沮喪,很挫敗的時候.
這個就是89篇.
但是整個的詩集就是告訴我們.
既然有一個上帝的恩典.
我們就不至於受虐盼望.
何況我們今天有耶穌基督的死而復活.
上帝的指紋.
最後是保證一個victory.
一個勝利.
不是因為我們怎樣.
而是我們縱使人生都散了.
塌陷的時候.
縱使我們不順.
但是上帝永遠是信實的.
所以到第150篇的時候.
它就是這樣結束.
要讚美耶和華.
我們的生命最後的時候.
是要讚美誰.
不是我們厲害.
是上帝的恩典.
然後這篇詩篇又解答.
另外一個有關敬拜的問題.
三至五節就講到怎樣去敬拜.
你看到有多少樂器在這裡.
一,二,三,可能七件吧.
是當時用的樂器.
重點不是一定要用這些樂器.
重點是當時他們用各種.
拿走所有工具.
所有不同的樂器.
他們去敬拜這一位這麼信實的主.
今天我們怎樣去應用呢.
我們都需要拿走所有工具.
我們都應該要這樣去敬拜.
這裡特別講到敬拜是一回事.
是一個動詞.
你看到有跳舞.
有擊鼓.

$^{801}$全部都是有動作的東西.
這些我們華人基督徒的文化.
我們真的不懂.
因為我們就像碌木一樣.
或者我們表情也像碌木一樣.
我們不懂得怎樣表達.
喜怒哀樂.
因為我們從小到大的升級.
就不是鼓勵我們做這件事的.
是不是.
最重要有Poker face.
每樣東西都有Poker face.
無動於衷這樣才是強人才是好.
但聖經裡面講的敬拜不是這樣的.
是有喜樂.
有喜怒哀樂.
有動作.
要站起來.
要有動作.
要去到跳起來.
整個人來表達.
嚴肅不一定是淑寧.
對不對.
淑寧的人也不一定是嚴肅的.
明白吧.
所以敬拜是一個認真的敬拜.
不一定是嚴肅的敬拜.
可能有這樣的空間.
是需要嚴肅的敬拜.
不過這裡講到.
worship is a verb.
是一個動詞.
身體是需要動.
你的心也需要動.
這個才是真正整個人的敬拜.
到最後這裡講到.
誰應該去敬拜.
第六節.
他說「凡有氣息的都要讚美耶和華」.
凡有氣息的.

$^{841}$你看到他有個鋪排.
就算有最好的樂隊.
最好的工具.
最好的樂器.
最好的音樂也好.
到最後第六節他怎麼結束.
他說「凡有氣息的都應該要讚美耶和華」.
氣息就是我們生命裡面最基本的禮物.
今天我們有口氣.
是上帝給我們最大的恩典.
對嗎.
到你有一天.
你進了ICU的時候.
你就知道那些氣多貴.
氣息是很寶貴的.
所以這篇詩篇講到.
凡有氣息的都要讚美耶和華.
今天上帝給你一口氣.
你應當去讚美祂.
應當去敬拜祂.
在世間上有很多不同的權柄.
要求我們做這件事.
要去讚美祂,要去撐祂.
但是上帝的指紋.
那首音樂的節奏是不同世間的.
我們撐的,唱的,讚美的.
只有一位.
我們主耶穌基督.
我希望今天給你一個這樣的溫習.
什麼叫敬拜.
因為你記住敬拜的原因.
敬拜的意思.
以致我們聚集在一起的時候.
就算有一些誤離教會.
平常誤離教會.
他們進來的時候.
他們會知道這裡是有不同的.
因為耶穌基督值得敬拜.
耶穌基督是真實的.
我們誠心帶領他們去祈禱.

$^{881}$在禱告裡面我們去預備聖餐.
預備我們領受聖餐.
上帝我們因著你的恩典.
因著你的慈愛.
因著你的大能.
我們去敬拜,讚美你.
耿古到永遠都是信實的神.
願意拯救我們這班破碎有罪的人.
讓我們能夠成為你自己的國民.
主要是讓我們的口,我們的氣息.
我們的生命,我們的身體.
起到去敬拜,去讚美你.
特別今天我們拿起聖餐的餅和杯.
我們知道你自己的大德,你的恩典.
去拯救我們這班.
洗淨我們這班人的罪.
主要是我們站在你面前.
耶穌基督為我們所作.
祂十加的救恩.
祂願意去犧牲祂自己的生命.
流出祂的寶血.
洗淨我們每一個.
讓我們能夠成為你自己的國民.
無論男與女.
自主的慰勞的.
哪一個民族的.
說什麼謊言,說什麼話的.
我們都因著耶穌基督.
能夠得到這麼不配得的救恩.
當我們領受這個聖餐的時候.
不單想起你的恩典.
也要讓你的恩典.
在我們生命裡慢慢改造的工作.
更新我們.
祈禱奉傳明禁.
\newpage



\section{出埃及記 3:1-22}
\label{sec:CVOZX9nsoIc}
\textbf{呼召的關係 (出埃及記3\_1-22) - 江芷敏傳道}
\newline
\newline
連結: \href{https://youtube.com/watch?v=CVOZX9nsoIc}{\texttt{https://youtube.com/watch?v=CVOZX9nsoIc}} ~~~~ 語音日期: 2026-03-16
\newline
\newline
\hyperref[sec:y65Q_8wRaJ8]{\small{< < < PREV SERMON < < <}}
~
\hyperref[sec:index_chronic]{\small{[返順時目]}}
~
\hyperref[sec:index_scriptual]{\small{[返順卷目]}}
~
\hyperref[sec:6DmNaa6pfms]{\small{> > > NEXT SERMON > > >}}
\newline
\newline
出埃及記 3:1-22
\newline
\begin{longtable}{cl}
\hline
\hline
章節 & 經文 (和合本修訂版)\\
\hline
3:1 & \begin{tabularx}{0.7\textwidth}{X} 摩西牧放他岳父米甸祭司葉特羅的羊群，他領羊群往曠野的那一邊去，到了神的山，就是何烈山。 \end{tabularx} \\ \\ \relax
3:2 & \begin{tabularx}{0.7\textwidth}{X} 耶和華的使者在荊棘的火焰中向他顯現。摩西觀看，看哪，荊棘在火中焚燒，卻沒有燒燬。 \end{tabularx} \\ \\ \relax
3:3 & \begin{tabularx}{0.7\textwidth}{X} 摩西說：「我要轉過去看這大異象，這荊棘為何沒有燒燬呢？」 \end{tabularx} \\ \\ \relax
3:4 & \begin{tabularx}{0.7\textwidth}{X} 耶和華見摩西轉過去看，神就從荊棘裡呼叫他說：「摩西！摩西！」他說：「我在這裡。」 \end{tabularx} \\ \\ \relax
3:5 & \begin{tabularx}{0.7\textwidth}{X} 神說：「不要靠近這裡。把你腳上的鞋脫下來，因為你所站的地方是聖地」。 \end{tabularx} \\ \\ \relax
3:6 & \begin{tabularx}{0.7\textwidth}{X} 他又說：「我是你父親的神，是亞伯拉罕的神，以撒的神，雅各的神。」摩西蒙上臉，因為怕看神。 \end{tabularx} \\ \\ \relax
3:7 & \begin{tabularx}{0.7\textwidth}{X} 耶和華說：「我確實看見了我百姓在埃及所受的困苦，我也聽見了他們因受監工苦待所發的哀聲；我確實知道他們的痛苦。 \end{tabularx} \\ \\ \relax
3:8 & \begin{tabularx}{0.7\textwidth}{X} 我下來是要救他們脫離埃及人的手，領他們從那地上來，到美好與寬闊之地，到流奶與蜜之地，就是迦南人、赫人、亞摩利人、比利洗人、希未人、耶布斯人之地。 \end{tabularx} \\ \\ \relax
3:9 & \begin{tabularx}{0.7\textwidth}{X} 現在，看哪，以色列人的哀聲達到我這裡，我也看見埃及人怎樣欺壓他們。 \end{tabularx} \\ \\ \relax
3:10 & \begin{tabularx}{0.7\textwidth}{X} 現在，你去，我要差派你到法老那裡，把我的百姓以色列人從埃及領出來。」 \end{tabularx} \\ \\ \relax
3:11 & \begin{tabularx}{0.7\textwidth}{X} 摩西對神說：「我是甚麼人，竟能去見法老，把以色列人從埃及領出來呢？」 \end{tabularx} \\ \\ \relax
3:12 & \begin{tabularx}{0.7\textwidth}{X} 神說：「我必與你同在。這就是我差派你去，給你的憑據：你把百姓從埃及領出來之後，你們必在這山上事奉神。」 \end{tabularx} \\ \\ \relax
3:13 & \begin{tabularx}{0.7\textwidth}{X} 摩西對神說：「看哪，我到以色列人那裡，對他們說：『你們祖宗的神差派我到你們這裡來。』他們若對我說：『他叫甚麼名字？』我要對他們說甚麼呢？」 \end{tabularx} \\ \\ \relax
3:14 & \begin{tabularx}{0.7\textwidth}{X} 神對摩西說：「我是自有永有的」；又說：「你要對以色列人這樣說：『那自有永有的差派我到你們這裡來。』」 \end{tabularx} \\ \\ \relax
3:15 & \begin{tabularx}{0.7\textwidth}{X} 神又對摩西說：「你要對以色列人這樣說：『耶和華－你們祖宗的神，就是亞伯拉罕的神，以撒的神，雅各的神差派我到你們這裡來。』這是我的名，直到永遠；這也是我的稱號，直到萬代。 \end{tabularx} \\ \\ \relax
3:16 & \begin{tabularx}{0.7\textwidth}{X} 你去召集以色列的長老，對他們說：『耶和華－你們祖宗的神，就是亞伯拉罕的神，以撒的神，雅各的神向我顯現，說：我實在眷顧了你們，眷顧你們在埃及的遭遇。 \end{tabularx} \\ \\ \relax
3:17 & \begin{tabularx}{0.7\textwidth}{X} 我也曾說：要把你們從埃及的困苦中領出來，往迦南人、赫人、亞摩利人、比利洗人、希未人、耶布斯人的地去，就是到流奶與蜜之地。』 \end{tabularx} \\ \\ \relax
3:18 & \begin{tabularx}{0.7\textwidth}{X} 他們必聽你的話。你和以色列的長老要到埃及王那裡，對他說：『耶和華－希伯來人的神向我們顯現，現在求你讓我們往曠野去，走三天的路程，為要向耶和華我們的神獻祭。』 \end{tabularx} \\ \\ \relax
3:19 & \begin{tabularx}{0.7\textwidth}{X} 我知道若不用大能的手，埃及王不會放你們走。 \end{tabularx} \\ \\ \relax
3:20 & \begin{tabularx}{0.7\textwidth}{X} 因此，我必伸出我的手，在埃及施行我一切的神蹟，擊打這地，然後，他才放你們走。 \end{tabularx} \\ \\ \relax
3:21 & \begin{tabularx}{0.7\textwidth}{X} 我必使埃及人看得起你們，你們離開的時候就不至於空手而去。 \end{tabularx} \\ \\ \relax
3:22 & \begin{tabularx}{0.7\textwidth}{X} 每一個婦女必向她的鄰舍，以及寄居在她家裡的女人，索取金器、銀器和衣裳，給你們的兒女穿戴。這樣你們就掠奪了埃及人。」 \end{tabularx} \\ \\
[1ex]
\hline
\hline
\end{longtable}
$^{1}$大家好 早安.
時間過得很快 今天已經是夏令的時間.
我記得上一次說到好像是冬令時間.
我們剛剛過完新年.
又進入大齋期.
很快就到甚麼節日.
復活節.
而且還是UV的十週年.
教會很快就十歲了.
其實最近我經常在想.
教會就快十週年了.
有甚麼可以送給教會.
獻情給神呢.
如果我自己有個兒女是十歲的話.
我都會跟他好好慶祝.
更何況是十週年.
我想都會很想跟他展望一下未來.
所以今天這一篇的講道.
很想跟大家一起去思想.
我們在這個十週年裡.
可以獻情一份怎樣的禮物給我們的神.
今天我們會看出埃及記的.
三章 十九章和三十二章.
是有點貪心的.
不過大家可以預備好你的聖經.
隨時打開我們一起去看.
我會在PowerPoint裡面.
也有展示經文的.
大家不用擔心.
出埃及記三章.
字有點小.
不過不要緊 我會跟大家說一下故事.
出埃及記三章一至六節.
其實是一個上帝呼召摩西的經文.
摩西其實是第一次遇見神.
神就呼召他帶領以色列人出埃及.
其實那時候摩西在曠野那裡.
在牧放羊.
帶著他外父的羊.
一路行 一路行 去到曠野的另一邊.

$^{41}$那座山叫做荷列山.
就在那裡 神主動跟他顯現.
你見到有光 在荊棘裡面著了火.
但是那些荊棘是燒不爛的.
所以摩西覺得很奇怪.
走過去看看甚麼事.
在這個時候 神就叫他.
摩西 摩西.
摩西就說.
我在這裡.
神就跟他說.
你不要走過來.
因為你要脫掉你腳上的鞋.
你站的地方是一個聖地.
然後神就再介紹自己.
祂說 我是你父親的神.
是亞伯拉罕的神.
以色列的神 雅各的神.
摩西聽完之後.
他就遮住自己.
他很害怕見到神.
我們要在這裡問.
這兩幅圖究竟哪一幅像一至六節所形容的.
剛才的經文.
其實通常都是日間放陽.
沒理由晚上的.
但是下面那幅圖.
好像有一樣東西沒有畫到.
就是摩西有沒有遮住自己的臉.
是沒有的.
所以我覺得上面黑漆漆的那張圖比較像.
不過不知道為甚麼是晚上.
我們看兩幅圖其實都有一個共同點.
就是他也有畫到摩西脫下鞋子.
其實我經常看這段經文.
我經常想我們回家都脫下鞋子.
那是不是走進去就是一個比較聖潔的地方呢.
又不是.
為甚麼上帝在這裡要描述摩西要脫下鞋子.
因為他站的地方是聖地呢.

$^{81}$聖 即是說聖潔.
聖潔究竟是甚麼意思.
聖潔是分別出來的意思.
即是從世俗裡面分別出來.
究竟怎樣去定義哪裡是世俗.
哪裡是聖潔呢.
有甚麼原則呢.
其實聖潔也是一個頗抽象的概念.
在摩西五經其實經常提到.
關於聖潔的經文或者他的條款.
不過我今天不是想解釋甚麼是聖潔.
因為要講很久.
聖潔在這裡每一次都出現.
在上帝呼召人的時候.
他總是呼召人的時候.
要叫人去自潔.
為甚麼呢.
我們有看士思記 善知書.
譬如耳塞.
上帝呼召他的時候.
他要自己潔淨.
為甚麼上帝要給一個動作人.
去聽他的呼召呢.
其實上帝每一次提到聖潔要求的時候.
他劃定了甚麼界線為聖潔.
就好像摩西一樣.
你站在這個地方.
其實你每天都在這裡放羊.
但是我今天就將這個地方分別為聖.
你要脫下鞋子才可以站上去.
我們就看看.
其實摩西每天都在這個地方裡面.
他每天都放羊.
為甚麼上帝要將這個地方變為聖潔呢.
如果我們想想我們每天工作的地方.
就好像我們的辦公室.
其實一點都不聖潔.
還想快點逃離那個地方.
或者現在你的工作地方.
變成了你的家.

$^{121}$你都不會覺得.
為甚麼有一個這麼特別的景象.
或者想像不到.
上帝會在那裡呼召我們呢.
其實上帝好像在跟摩西說一樣東西.
你現在要聽我說.
你來到這個地方.
不再是你平日很熟悉放羊的地方.
今天這個地方是我主事的.
今天這個地方是我要去呼召你的.
上帝將這個對話弄清楚.
那個誰是主事的那個.
每一次上帝呼召人的時候.
都要弄清楚這個次序.
是我才是主.
是我給你這個聖潔的界線.
我一直在寫的時候.
就想起主人和貓狗寵物的關係.
其實貓和狗都有自主權.
其實你在家裡面.
牠都可以選擇親不親近你.
而且可能有時候關係調轉了.
你要服侍牠.
又要剷屎.
又要摸牠.
要給牠足夠的關注.
不過你始終是主人.
因為你是買食物給牠那個.
有病的時候是你帶牠去看醫生.
而且你是要餵食物給牠吃.
你和牠有一些不同的操練.
或者你要給牠一些規矩.
當那些規矩一出來.
聽到指令的時候.
這些寵物.
這些貓和狗就知道.
你就是那個主人.
寵物是要聽主人的話.
我想起這個關係.
當然不是我們直接代入.

$^{161}$我們就是寵物.
我們比寵物尊貴得多.
或者自由得多.
希望大家不要誤會我這個例子.
但我正想說出.
當上帝給這個指令的時候.
我們要聽.
我們要服從.
那個關係就變得清楚了.
牠是主.
我們是跟從.
上帝給了我們很大的自由意志.
當他跟我們說這些指令的時候.
如果我們選擇去跟從他的話.
我們選擇去跟從他的話.
我們就和他恢復一個.
從熟的關係.
好像摩西一樣.
摩西和上帝.
恢復一個從熟的關係.
他知道誰是主.
當然出埃及記.
不單止是說呼召摩西一個人.
呼召是一個進程.
慢慢.
他由呼召摩西.
變成呼召整個以色列的群體.
就在出埃及記十九章那裡.
上帝又再一次去呼召以色列人.
耶和華神就跟摩西說.
你現在去到百姓那裡.
叫他們明天要自潔.
叫他們洗衣服.
然後第三天.
我就會降臨在這個西奈山上面.
其實這個西奈山和荷里山.
都是同一座山.
不過是山脈和山峰的關係.
來到這座山的裡面.
上帝要為他們定界限.

$^{201}$他說你們要謹慎.
不可以上去這個山.
也不可以摸這個山邊.
只要你摸這個山.
就會致死他.
第十三戰.
他說也不可以伸手摸他.
因為必被石頭打死.
說用箭射頭.
無論是人還是畜生都不得活.
這個聖潔的界線是誰定的.
是上帝定的.
這次是不是跟呼召摩西的時候一樣.
是不一樣的.
所以每一次被呼召的時候.
那個聖潔的界線很不一樣.
其實上帝每呼召一個人.
那個聖潔的界線不一樣.
其實不是我們要深究.
為什麼他這次給這個呼召.
他要做這個聖潔的儀式.
其實是每一次.
那個聖潔的標準來到的時候.
我們就知道上帝就是那個主.
我們就是要聽他講.
他就是給呼召那個上帝.
所以在這裡我們又再看到一次.
那個呼召是我們都要放在適合我們的位置.
當上帝的呼召來到的時候.
我們這個群體就要知道誰是我們的主.
放一個適合的姿態.
這段經文不單止是說.
我們要放一個適合的姿態.
放一個適合的位置.
去知道上帝是上帝.
更加是調整我們的眼光.
第三章 神和摩西的對話.
繼續說.
他說 我確實看見了我百姓在埃及所受的困苦.
我也聽見了他們因受監宮苦待所發的哀聲.

$^{241}$我也確實知道他們的痛苦.
這個經文是上一次和大家講到的時候都有提過.
不過三章這裡再重複一次.
上帝是那位親近我們的上帝.
摩西聽到上帝的呼召.
說他要去 上帝要差派他去法老那裡.
並將他的百姓以色列人從埃及領出來.
摩西第一個反應是什麼.
第十一節 摩西說.
我?我是什麼人?.
竟然可以去見法老?.
還要將以色列人從埃及那裡領出來?.
如果大家對聖經有些了解.
會知道摩西究竟有什麼背景.
摩西以前是埃及的王子.
他在那裡成長.
但是他很看不過眼這些埃及人欺負他的同胞.
所以他打死了欺負他的同胞的埃及人.
他逃難到米甸做牧羊人.
在曠野那裡一留就留了四十年.
其實他已經和外面這個世界脫節.
他不知道外面發生什麼事.
埃及已經離他很遠.
那時候沒有新聞.
不會像我們現在這樣.
今天發生什麼事 明天全世界都知道.
他還要面對一個不可以撼動的法老.
帶以色列人出埃及憑什麼呢?.
但是上帝的回應是什麼?.
我們12節看到是什麼.
他說「我必與你同在」.
上帝沒有回答他那個問題.
我究竟是什麼人.
上帝也沒有說你行的.
我為你打氣 你一定行的.
上帝沒有說這些.
上帝說「我與你同在」.
第14節上帝也說「我是誰」.
「我是自有永有的那個」.
你看到這些經文全部都是「我」.

$^{281}$摩西就問上帝「我是什麼人」.
「我有什麼可能 有什麼能力」.
上帝沒有回答他這個問題.
不過上帝也繼續用「我」這個字去開頭.
說「我必與你同在」.
「我是自有永有的那個上帝」.
I am who I am.
其實翻譯了英文也不知道是什麼.
或者翻譯了原文也不知道這個名字是什麼.
I am who I am.
好像沒有東西可以定義上帝.
我們看古晉東一些文化 一些神明.
他們很多時候都用一個名字去顯示.
他自己有什麼功能 他有什麼屬性.
就好像太陽神 風暴神 月亮神 雷神.
你一聽他的名字就知道他要對付什麼.
或者他有什麼功能.
其實的確這樣的神明都比較容易捉摸.
因為一說就知道他是代表什麼上帝 什麼神.
但上帝就在這裡用「I am who I am」.
「我是自有永有」.
他去回應摩西.
他在這裡否定人為他下一個定義.
否定那些關注他的名字.
否定那些對他定義的定義.
很有趣 上個星期就是年宵.
在年宵裡我們教會也有擺一個Outreach的攤.
那個攤就是玩遊戲和請人喝咖啡.
很有趣 有一個New Chicago的學生.
他一直喝咖啡和Ken說.
Ken就介紹他 這一種咖啡可以喝到什麼水果味.
因為我們沖手沖 很多時候說.
這一種咖啡有什麼朱古力味 有什麼荔枝味.
我們喜歡喝咖啡的人.
慢慢學會品嘗咖啡.
那個學生很有趣.
他就說 為什麼咖啡不可以是咖啡味.
很有趣 今天Ken就跟我分享.
為什麼咖啡不可以是咖啡味.
為什麼要用其他味去定義咖啡是什麼味.

$^{321}$上帝就好像這樣.
他沒有給其他東西去定義.
他就是這個世界未有以前他已經在的上帝.
他就是創造的上帝.
他就是和我們一起的上帝.
而且上帝說 你不要再看你自己了.
你整天說我我我.
你現在看一下我究竟是誰.
上帝在這裡跟他說.
不要再看你不行的東西了.
不要再看著自己的限制.
你的焦點應該是看著我這個呼召你的神.
其實我也很明白摩西的掙扎.
因為神叫他做的東西.
跟他的實力差距很遠.
你完全覺得沒有可能.
的確上帝有時候叫我們做的東西就是沒有可能.
其實我想起一件很不可能的事.
其實我們生命裡面都出現過很多不同的上帝.
但我覺得最不可能的就是上帝帶領Ken來芝加哥讀書.
因為我們以前讀書都是身無分文.
也覺得自己的英文很有限制.
在我們十年前正在讀神學.
就是2016年的時候.
Ken有一次無緣無故打給我.
那時候我們還沒有拍拖.
我跟他也不是什麼關係.
但因為我純粹住在長洲.
他遺留在長洲.
他有需要就打給我.
就是在那個星期六他發燒.
因為星期六就要看急症在香港.
看急症當時是要100元.
但他戶口又不夠100元.
所以按不到那張100元出來.
所以他就打電話給我問我借錢.
現在這些電話就要很小心.
但那時候他就問我借100元.
我也覺得他很慘.
又發燒又要看急症.

$^{361}$又沒有人在他身邊.
家人又在深水Po 住.
我就馬上去長洲醫院借錢給他.
那時候我也很難想像到.
一個身無分文.
身無分文但一兩千也有.
但他連100元也沒有.
但他竟然來到芝加哥讀書.
回望我們過去.
他完成了課程現在還差一篇文.
但我們用了很多錢.
我們回想這些都是上帝給我們的恩典.
上帝也是和我們同在.
和恩典很夠用的上帝.
我們計算一下.
我昨天問他.
你猜我們這三年應該用了多少錢.
Trinity的網站說.
我們一年要用五萬美金.
兩夫妻加上租金和在那裡讀書.
我們讀了兩年半.
我想怎樣也有十二萬.
我一直沒有想過.
上帝可以這樣帶領我們來這裡.
而且我們剛剛來到的時候.
好像一大群人在Orientation Dinner吃飯.
我一句英文都聽不到.
好像自己來錯了地方.
很多的質疑.
到今天可以在這裡服侍.
我都覺得完全是上帝的恩典.
如果當天我們要去計算.
計算自己有多少限制的話.
我相信我們不能夠站在這裡和大家分享.
如果我們只是要看著數字的話.
我們是不需要向前走的.
但上帝就是那個呼召我們的上帝.
就是要叫我們不要只看著自己的限制.
要看著這個自有永有和祂會供應的上帝.
第二點就是我們不單止要看自己的不足.

$^{401}$我們更加要將我們的焦點.
放在這個自有永有超過一切的上帝身上.
第三 呼召其實不是一天完成的一件事.
是一個進程.
就好像出埃及的三章去呼召摩西.
然後到十九章.
上帝是呼召整個的以色列群體.
上帝說你們要聽我的話.
祂說你要去作我祭司的國度.
作聖潔的國民.
上帝在十九章那裡是這樣呼召以色列人.
作祭司的國度.
為聖潔的國民.
其實這個進程不單止是由一開始.
摩西第一次見上帝就聽到這個呼召.
是一點一滴不斷去塗鑄摩西塗鑄以色列人.
摩西不斷在神面前去掙扎.
祂說其實我真的很不行.
我又不會說話.
你派另一個人吧.
摩西是一個很將自己的軟弱呈現給上帝看的一個人.
上帝也在這一段的關係裡面.
沒有走開去聽摩西去說他自己多不行.
你看其實不單止是三章.
第四章也一直在說摩西和上帝不斷去談判.
我們也看到十九章.
以色列人是怎樣回應上帝的呼召.
他不像摩西那樣很多個來回.
他是一個很簡潔的一句說話.
很簡潔的一句說話就是十九章八節.
如果大家有聖經可以打開.
百姓只是跟上帝說凡耶和華所說的我們都要遵行.
只是一句說話就回應上帝的呼召.
其實這是一個很屬靈的群體.
馬上說Amen的一句說話.
不過其實我們看到十九章之後.
就是上帝去呼召摩西上山.
跟他說法版和怎樣做一個聖潔的群體.
在摩西上山之後.
百姓看到摩西遲遲未下山就跟亞倫說.

$^{441}$亞倫起來我們要做一個神像.
可以在我們面前引路.
剛才說的凡耶和華所說的我們都要遵行.
已經拋諸腦後了.
亞倫更加去聽這些百姓說.
他說以色列你們快點回去拿些金來.
我們要一起去做一個金牛獨.
這個領袖是別人說什麼他就做什麼.
他完全沒有想過.
然後做了金牛獨之後亞倫還說.
以色列這個就是領你們出埃及的神.
他一點都沒有糾正這些走錯路的以色列人.
以色列人原來是想要一個看得見摸得到.
其實不一定是耶和華神.
你再看看三十二章的經文.
他也當摩西是帶領他們出以色列的神.
只要讓他們看得見摸得到.
有足夠的安全感.
這個就是他們的神.
簡單來說就是誰幫到我.
誰就是我們的神.
以色列人就是這樣用他們的用語去包裝這個偶像.
其實他們不需要一個神.
他們需要一個工具.
如果今天他們有AI.
可能都覺得AI是他們的神.
什麼都問得到 什麼都解答得到.
原來他們所說的只是一個口號式的敬拜.
我再看摩西和神來回這麼多個圈.
去承認自己的軟弱 承認自己的掙扎.
這個關係來得更加真實.
呼召是一個關係的開始.
或者我可以形容它是一個坦白的開始.
當我想起我怎樣信主的時候.
我都是一個很軟弱的狀態.
或者我是跟神坦白.
我有什麼得罪他.
就好像剛才那首詩歌《真相》.
我慢慢揭開一些包住我的保護機制.
或者一些我很想欺騙自己欺騙人的東西.

$^{481}$我打開了這些看到最真最軟弱的自己.
我們來到神的面前.
那一刻是我和神最親近的時間.
由摩西被呼召的那一天.
他就開始和神建立這個關係.
這個關係不是我們看得很完美的.
是一個跌跌碰碰 一撲一碌的關係.
他也有做錯事得罪神.
但是他和神的關係就很真實.
就好像在三十三章也有說.
神和摩西的關係是怎樣呢.
我們一起去讀好嗎.
《出埃及記》三十三章十一節上.
耶和華與摩西面對面說話.
好像人與朋友說話一般.
好像朋友一樣.
我們再看《申命記》三十四章十節.
以後以色列中再沒有興起先知像摩西的.
他是耶和華面對面所認識的.
這個面對面的認識.
是一個很真實很坦白的一個關係.
神體重的不是我們可以做些什麼.
神體重的是我們和他的關係.
所以這個呼召是叫我們進入和他這麼親密的關係.
是要坦坦白白的.
在我身邊我覺得最親密的一個關係.
當然是我和我老公.
但亦都是我和我媽媽.
在這個關係裡面.
我們經歷過很多磨擦.
但在這些磨擦裡面.
我慢慢覺得我們觸摸到對方究竟是一個怎樣的人.
我和我媽媽小時候很疏遠.
因為她很年輕生我.
為了謀生我十歲以前和我媽媽一起住.
因為爸爸媽媽經常不在身邊.
所以我要學會看別人的眉頭眼額.
很努力做一個大人症.
很懂得看別人需要什麼.
但同時在這個保護自己生存機制的底下.

$^{521}$我對我父母慢慢積累了很多埋怨.
到我青春期的時候.
我開始回家和他們住.
我經常都爆發.
去控訴他們之前為什麼要這樣.
一開始說完大家都是大哭一場.
我媽媽其實很愛我和很疼愛我.
不過她都不懂怎樣去接到我的情緒.
她都覺得很委屈.
因為她已經很盡力地去愛我.
其實她只是十八歲就生我.
她自己都未搞得定.
到我十八歲的時候.
我才慢慢開始明白.
原來做一個大人要這麼多事要顧.
所以我慢慢明白她多一點.
加上我十九歲開始信耶穌.
我們回教會.
我們一家人試著牽著手去祈禱.
每一次想起過往的關係傷痛的時候.
我們都會在哭.
但我們願意彼此去聽.
去坦承我們心裡的感受是怎樣.
我覺得和我媽媽開始越走越近.
不是真正的.
而是情緒上.
我知道她為什麼要這樣做.
我開始盛載她的關係.
開始慢慢地越走越近.
是基於她的坦白和神的恩典.
她不斷地光照我們.
神的恩典是很好.
她不斷地讓我們一家人修補這個關係.
讓我們變得可以彼此擁抱.
我不知道大家有沒有經歷過你的友誼.
都試過在吵架和坦白之後.
大家的感情變得更加牢固.
弟兄姊妹更加容易吵架.
弟兄姊妹之間有沒有試過坦誠.
令你覺得更加親密.

$^{561}$上帝在這裡也跟我們說.
我們需要和祂有一個坦白的關係.
今天我們去問UV十週年.
我們有什麼可以獻情給我們的神.
在這段經文裡再一次提醒我們.
上帝是很想和我們恢復一個關係.
在祂的呼召裡.
我們要回到和祂重熟的關係.
祂才是我們生命的主.
祂才是UV的主.
祂是那個話事的那個.
而我們的眼光要再一次對著這個呼召我們.
和超越一切的上帝.
祂就是那個自由永有.
超越一切的上帝.
一切都是從祂開始.
UV也是.
而祂也邀請我們和祂進入一個更深更坦白的關係.
不知道大家在這幾點裡有什麼想法.
或者今天如果上帝來到我們當中.
問我們願不願意再更深一步和祂建立這個坦白的關係.
不單是每個星期日回來.
是在我們平日裡和祂建立一個更坦白的禱告.
那個同行的關係.
讓我們一起去祈禱.
主,我求你幫助我們.
當我們在這個你如我乍的社會裡.
在我們營營役役的生活裡.
我們有時候都好像忘記了.
你就是我們生命的主.
主,我求你親自幫助我們.
回復這個和你重熟的關係.
你就是我們的主.
你就是UV的主.
也求你幫助我們.
不要經常面對困難的時候.
就想起自己的限制.
我們能夠多一些去想像.
和多一些去凝望著你自己的能力.
當我們說到教會要去擴展的時候.

$^{601}$當我們要去找一個更適合的地方的時候.
當我們要面對很多財務上的困難的時候.
主啊,求你幫助我們.
能夠去望著你這個.
自由,永有,超越一切.
而且供應很足夠的上帝.
主啊,我們都懇求你.
挪去我們心裡面那個金牛毒的敬拜.
我們真的要去來到你的面前.
坦坦白白的去展開我們自己.
給你去對付我們的生命.
將我們最軟弱的一面呈現給你.
主啊,求你閱立我們的禱告.
閱立我們的敬拜.
馮主名求,Amen.
祝福大家.
\newpage



\section{士師記 6:1-40}
\label{sec:6DmNaa6pfms}
\textbf{真實的敬拜 (士師記6\_1-40) - 楊志溢傳道}
\newline
\newline
連結: \href{https://youtube.com/watch?v=6DmNaa6pfms}{\texttt{https://youtube.com/watch?v=6DmNaa6pfms}} ~~~~ 語音日期: 2026-03-16
\newline
\newline
\hyperref[sec:CVOZX9nsoIc]{\small{< < < PREV SERMON < < <}}
~
\hyperref[sec:index_chronic]{\small{[返順時目]}}
~
\hyperref[sec:index_scriptual]{\small{[返順卷目]}}
~
\hyperref[sec:TQN008eNyL0]{\small{> > > NEXT SERMON > > >}}
\newline
\newline
士師記 6:1-40
\newline
\begin{longtable}{cl}
\hline
\hline
章節 & 經文 (和合本修訂版)\\
\hline
6:1 & \begin{tabularx}{0.7\textwidth}{X} 以色列人又行耶和華眼中看為惡的事，耶和華就把他們交在米甸手裡七年。 \end{tabularx} \\ \\ \relax
6:2 & \begin{tabularx}{0.7\textwidth}{X} 米甸的手戰勝以色列；以色列人躲避米甸人，就在山中挖洞穴，挖洞建營寨。 \end{tabularx} \\ \\ \relax
6:3 & \begin{tabularx}{0.7\textwidth}{X} 每當以色列人撒種之後，米甸、亞瑪力和東邊的人都上來攻打他們， \end{tabularx} \\ \\ \relax
6:4 & \begin{tabularx}{0.7\textwidth}{X} 對著他們安營，毀壞那地的農作物，直到迦薩，沒有給以色列留下食物，牛、羊、驢也沒有留下。 \end{tabularx} \\ \\ \relax
6:5 & \begin{tabularx}{0.7\textwidth}{X} 因為那些人帶著他們的牲畜和帳棚上來，像蝗蟲那樣多；人和駱駝無數，都進入境內，毀壞全地。 \end{tabularx} \\ \\ \relax
6:6 & \begin{tabularx}{0.7\textwidth}{X} 以色列因米甸的緣故極其窮乏，以色列人就呼求耶和華。 \end{tabularx} \\ \\ \relax
6:7 & \begin{tabularx}{0.7\textwidth}{X} 以色列人因米甸的緣故呼求耶和華的時候， \end{tabularx} \\ \\ \relax
6:8 & \begin{tabularx}{0.7\textwidth}{X} 耶和華就差遣先知到以色列人那裡，對他們說：「耶和華－以色列的神如此說：『我曾領你們從埃及上來，從為奴之家出來， \end{tabularx} \\ \\ \relax
6:9 & \begin{tabularx}{0.7\textwidth}{X} 救你們脫離埃及人的手，脫離一切欺壓你們之人的手。我從你們面前趕出他們，把他們的地賜給你們。 \end{tabularx} \\ \\ \relax
6:10 & \begin{tabularx}{0.7\textwidth}{X} 我對你們說，我是耶和華－你們的神。你們住在亞摩利人的地，不可敬畏他們的神明，但你們卻不聽從我的話。』」 \end{tabularx} \\ \\ \relax
6:11 & \begin{tabularx}{0.7\textwidth}{X} 耶和華的使者到了俄弗拉，坐在亞比以謝族約阿施的橡樹下。約阿施的兒子基甸正在醡酒池那裡打麥子，為了躲避米甸人。 \end{tabularx} \\ \\ \relax
6:12 & \begin{tabularx}{0.7\textwidth}{X} 耶和華的使者向基甸顯現，對他說：「大能的勇士啊，耶和華與你同在！」 \end{tabularx} \\ \\ \relax
6:13 & \begin{tabularx}{0.7\textwidth}{X} 基甸對他說：「主啊，請容許我說，耶和華若與我們同在，我們怎麼會遭遇這一切事呢？我們的列祖告訴我們：『耶和華領我們從埃及上來』，他那奇妙的作為在哪裡呢？現在耶和華卻丟棄了我們，把我們交在米甸人的手掌中。」 \end{tabularx} \\ \\ \relax
6:14 & \begin{tabularx}{0.7\textwidth}{X} 耶和華轉向基甸，說：「去，靠著你這能力拯救以色列脫離米甸人的手掌。我豈不是已經差遣了你嗎？」 \end{tabularx} \\ \\ \relax
6:15 & \begin{tabularx}{0.7\textwidth}{X} 基甸對他說：「主啊，請容許我說，我怎能拯救以色列呢？看哪，我這一支在瑪拿西支派中是最貧寒的，我在我父家又是最微小的。」 \end{tabularx} \\ \\ \relax
6:16 & \begin{tabularx}{0.7\textwidth}{X} 耶和華對他說：「我與你同在，你就必擊敗米甸，如擊打一個人。」 \end{tabularx} \\ \\ \relax
6:17 & \begin{tabularx}{0.7\textwidth}{X} 基甸對他說：「我若在你眼前蒙恩，求你給我一個證據，證明是你在跟我說話。 \end{tabularx} \\ \\ \relax
6:18 & \begin{tabularx}{0.7\textwidth}{X} 求你不要離開這裡，等我回來，將供物帶來，供在你面前。」他說：「我必等你回來。」 \end{tabularx} \\ \\ \relax
6:19 & \begin{tabularx}{0.7\textwidth}{X} 基甸去預備一隻小山羊，用一伊法細麵做了無酵餅，將肉放在籃子裡，將湯盛在壺中，帶到他那裡，在橡樹下獻上。 \end{tabularx} \\ \\ \relax
6:20 & \begin{tabularx}{0.7\textwidth}{X} 神的使者對基甸說：「將肉和無酵餅放在這磐石上，把湯倒出來。」他就照樣做了。 \end{tabularx} \\ \\ \relax
6:21 & \begin{tabularx}{0.7\textwidth}{X} 耶和華的使者伸出手裡的杖，杖頭一碰到肉和無酵餅，就有火從磐石中出來，吞滅了肉和無酵餅。耶和華的使者就從他眼前消失了。 \end{tabularx} \\ \\ \relax
6:22 & \begin{tabularx}{0.7\textwidth}{X} 基甸見他是耶和華的使者，就說：「哎呀！主耶和華啊！因為我真的面對面看見了耶和華的使者。」 \end{tabularx} \\ \\ \relax
6:23 & \begin{tabularx}{0.7\textwidth}{X} 耶和華對他說：「安心吧，不要怕，你不會死。」 \end{tabularx} \\ \\ \relax
6:24 & \begin{tabularx}{0.7\textwidth}{X} 於是基甸在那裡為耶和華築了一座壇，起名叫「耶和華沙龍」。這壇至今還在亞比以謝族的俄弗拉。 \end{tabularx} \\ \\ \relax
6:25 & \begin{tabularx}{0.7\textwidth}{X} 那夜，耶和華對基甸說：「你要把你父親的公牛，就是那七歲的第二頭公牛取來，並拆毀你父親為巴力築的壇，砍下壇旁的亞舍拉， \end{tabularx} \\ \\ \relax
6:26 & \begin{tabularx}{0.7\textwidth}{X} 在這堡壘頂上整整齊齊地為耶和華－你的神築一座壇，將第二頭公牛獻為燔祭，用你所砍下的亞舍拉當柴。」 \end{tabularx} \\ \\ \relax
6:27 & \begin{tabularx}{0.7\textwidth}{X} 基甸就從他僕人中選了十個人，照耶和華吩咐他的做了。他因怕父家和本城的人，不敢在白天做這事，就在夜間做。 \end{tabularx} \\ \\ \relax
6:28 & \begin{tabularx}{0.7\textwidth}{X} 城裡的人清早起來，看哪，巴力的壇被拆毀，壇旁的亞舍拉被砍下，第二頭公牛獻在築好的壇上， \end{tabularx} \\ \\ \relax
6:29 & \begin{tabularx}{0.7\textwidth}{X} 就彼此問：「這是誰做的事呢？」他們尋找查訪之後，就說：「這是約阿施的兒子基甸做的事。」 \end{tabularx} \\ \\ \relax
6:30 & \begin{tabularx}{0.7\textwidth}{X} 城裡的人對約阿施說：「把你的兒子交出來，我們要處死他，因為他拆毀了巴力的壇，砍下了壇旁的亞舍拉。」 \end{tabularx} \\ \\ \relax
6:31 & \begin{tabularx}{0.7\textwidth}{X} 約阿施對站著敵對他的眾人說：「你們是為巴力辯護嗎？你們要救它嗎？誰為它辯護，就在早晨把誰處死吧！巴力如果是神，有人拆毀了它的壇，就讓它為自己辯護吧！」 \end{tabularx} \\ \\ \relax
6:32 & \begin{tabularx}{0.7\textwidth}{X} 所以那日人稱基甸為耶路巴力，意思是：「他拆毀了巴力的壇，讓巴力與他爭辯吧。」 \end{tabularx} \\ \\ \relax
6:33 & \begin{tabularx}{0.7\textwidth}{X} 那時，所有的米甸人、亞瑪力人和東邊的人都聚集在一起，過了河，在耶斯列平原安營。 \end{tabularx} \\ \\ \relax
6:34 & \begin{tabularx}{0.7\textwidth}{X} 耶和華的靈降在基甸身上；他吹角，亞比以謝族都聚集跟隨他。 \end{tabularx} \\ \\ \relax
6:35 & \begin{tabularx}{0.7\textwidth}{X} 他派使者走遍瑪拿西，瑪拿西人也聚集跟隨他。他又派使者到亞設、西布倫、拿弗他利，他們也都上來會合。 \end{tabularx} \\ \\ \relax
6:36 & \begin{tabularx}{0.7\textwidth}{X} 基甸對神說：「你如果真的照你所說的，藉我的手拯救以色列， \end{tabularx} \\ \\ \relax
6:37 & \begin{tabularx}{0.7\textwidth}{X} 看哪，我把一團羊毛放在禾場上，若單是羊毛上有露水，遍地都是乾的，我就知道你必照你所說的，藉我的手拯救以色列。」 \end{tabularx} \\ \\ \relax
6:38 & \begin{tabularx}{0.7\textwidth}{X} 一切果然發生了。次日早晨基甸起來，把羊毛擰一擰，從羊毛中擠出露水來，裝滿一碗的水。 \end{tabularx} \\ \\ \relax
6:39 & \begin{tabularx}{0.7\textwidth}{X} 基甸又對神說：「求你不要向我發怒，我再說一次，讓我用羊毛再試一次，但願羊毛是乾的，遍地都有露水。」 \end{tabularx} \\ \\ \relax
6:40 & \begin{tabularx}{0.7\textwidth}{X} 這夜，神也照樣做，遍地都有露水，只有羊毛是乾的。 \end{tabularx} \\ \\
[1ex]
\hline
\hline
\end{longtable}
$^{1}$.
大家弟兄姊妹早上好.
盼望大家今天早上都能夠清醒.
雖然我們因為下令時少睡了一個小時.
但是我還是很期待接下來的白日會更加長一點.
太陽會晚一點下山.
.
我對今天《詩詩記》第五章的經文呢.
感到非常的期待.
因為呢我越花多很多的時間去默想它.
它就越改變我對神以及對敬拜的一個理解.
說句私心話.
我最近有一個朋友呢.
他的博士論文正在研究《詩詩記》的第五章.
而且最近剛剛出版.
所以呢在預備今天講到的時候.
我一邊支持他的研究成果.
也一邊在他那邊學習到很多的東西.
.
在進入講到之前.
我們一起先讀一下今天的經文.
我已經把《詩詩記》的第五章.
印成中英文的翻譯版本.
在前面的桌子上.
大家呢歡迎都來拿一份.
先提醒一下呢.
我今天使用的英文譯本呢.
是根據我朋友研究出來的翻譯的.
所以你在聖經的APP上面呢.
應該找不到這個版本的.
.
大家都找到你的copy了嗎.
.
那日底波拉和亞比諾安的兒子巴拉.
唱歌說.
以色列有領袖率領.
百姓甘心犧牲自己.
你們應當稱頌耶和華.
君王啊要聽.
王子啊要側耳.

$^{41}$我要我要向耶和華歌唱.
我要歌頌耶和華以色列的神.
耶和華呀.
你從西耳出來.
從以東田野向前行.
地震動.
天低下.
雲也低下於水.
眾山在耶和華面前搖動.
西奈山在耶和華.
以色列的神面前也搖動.
.
在雅納之子山家的時候.
在雅役的日子.
大道無人行走.
過路人繞道而行.
以色列農村荒蕪.
空無一人.
直到我底波拉興起.
興起做以色列之母.
以色列人選擇新的諸神.
戰爭就臨到城門.
以色列四萬人中.
看得見盾牌槍矛嗎.
.
我心嚮往以色列的領袖.
他們在民中甘心犧牲自己.
你們應當稱頌耶和華.
騎淺色母驢的.
坐繡花毯子的.
行走在路上的.
你們都當思想.
打水的聲音勝過弓箭的響聲.
那裡人要訴說耶和華公義的作為.
他對以色列鄉民公義的作為.
那時耶和華的子民下到城門.
.
那時貴族中的倖存者前行.
耶和華的百姓為我前進攻擊勇士.
.

\newpage



\section{約翰福音 2:1-12}
\label{sec:TQN008eNyL0}
\textbf{Suprising Transformation - (John 2\_1-12) - Tsz Kin Leung 驚喜的轉化 (約翰福音2\_1-12) - 梁紫健傳道}
\newline
\newline
連結: \href{https://youtube.com/watch?v=TQN008eNyL0}{\texttt{https://youtube.com/watch?v=TQN008eNyL0}} ~~~~ 語音日期: 2026-03-24
\newline
\newline
\hyperref[sec:6DmNaa6pfms]{\small{< < < PREV SERMON < < <}}
~
\hyperref[sec:index_chronic]{\small{[返順時目]}}
~
\hyperref[sec:index_scriptual]{\small{[返順卷目]}}
~
\hyperref[sec:a3FPlfJ6MYc]{\small{> > > NEXT SERMON > > >}}
\newline
\newline
約翰福音 2:1-12
\newline
\begin{longtable}{cl}
\hline
\hline
章節 & 經文 (和合本修訂版)\\
\hline
2:1 & \begin{tabularx}{0.7\textwidth}{X} 第三日，在加利利的迦拿有一個婚宴，耶穌的母親在那裡。 \end{tabularx} \\ \\ \relax
2:2 & \begin{tabularx}{0.7\textwidth}{X} 耶穌和他的門徒也被請去赴宴。 \end{tabularx} \\ \\ \relax
2:3 & \begin{tabularx}{0.7\textwidth}{X} 酒用完了，耶穌的母親對他說：「他們沒有酒了。」 \end{tabularx} \\ \\ \relax
2:4 & \begin{tabularx}{0.7\textwidth}{X} 耶穌說：「母親，我與你何干呢？我的時候還沒有到。」 \end{tabularx} \\ \\ \relax
2:5 & \begin{tabularx}{0.7\textwidth}{X} 他母親對用人說：「他告訴你們甚麼，你們就做吧。」 \end{tabularx} \\ \\ \relax
2:6 & \begin{tabularx}{0.7\textwidth}{X} 照猶太人潔淨禮的規矩，有六口石缸擺在那裡，每口可以盛兩三桶水。 \end{tabularx} \\ \\ \relax
2:7 & \begin{tabularx}{0.7\textwidth}{X} 耶穌對用人說：「把缸倒滿水。」他們就倒滿了，直到缸口。 \end{tabularx} \\ \\ \relax
2:8 & \begin{tabularx}{0.7\textwidth}{X} 耶穌又說：「現在舀出來，送給宴會總管。」他們就送了去。 \end{tabularx} \\ \\ \relax
2:9 & \begin{tabularx}{0.7\textwidth}{X} 宴會總管嘗了那水變的酒，並不知道是哪裡來的，只有舀水的用人知道。於是宴會總管叫新郎來， \end{tabularx} \\ \\ \relax
2:10 & \begin{tabularx}{0.7\textwidth}{X} 對他說：「人家都是先擺上好酒，等客人喝夠了才擺上次的，你倒把好酒留到現在！」 \end{tabularx} \\ \\ \relax
2:11 & \begin{tabularx}{0.7\textwidth}{X} 這是耶穌所行的第一個神蹟，是在加利利的迦拿行的，顯出了他的榮耀來，他的門徒就信他了。 \end{tabularx} \\ \\ \relax
2:12 & \begin{tabularx}{0.7\textwidth}{X} 這事以後，耶穌與他的母親、兄弟和門徒都下迦百農去，在那裡住了不多幾天。 \end{tabularx} \\ \\
[1ex]
\hline
\hline
\end{longtable}
$^{1}$[NON-ENGLISH SPEECH].
Brothers and sisters, good morning..
[NON-ENGLISH SPEECH].
To begin today, I would like to tell you a very simple story..
[NON-ENGLISH SPEECH].
There was a high school boy..
His name is Jimmy..
[NON-ENGLISH SPEECH].
He came from a poor family and was very quiet..
Every day, he would walk alone with a large stack of books,.
and almost no classmates would talk to him..
[NON-ENGLISH SPEECH].
In other people's eyes, he was a little bit strange,.
almost like a transparent person with a very low presence..
[NON-ENGLISH SPEECH].
One day after school, he was carrying a whole stack.
of books on his way home, as usual..
[NON-ENGLISH SPEECH].
Several classmates walked over and suddenly bumped into him.
on purpose..
[NON-ENGLISH SPEECH].
Books and stationery were all scattered.
on the ground, everywhere..
[NON-ENGLISH SPEECH].
People around him just smiled, and then they.
separate ways and walk away..
[NON-ENGLISH SPEECH].
At this moment, there was a classmate.
standing right beside him..
He could just pretend that he actually.
did not see everything happen..
[NON-ENGLISH SPEECH].
But he felt a small sense of gratitude..
[NON-ENGLISH SPEECH].
He thought, should I go over and help?.
[NON-ENGLISH SPEECH].
So he went over and scattered them..
[NON-ENGLISH SPEECH].
Helped Jimmy to pick up his books one by one..
[NON-ENGLISH SPEECH].

$^{41}$And asked him, are you OK?.
Where do you live?.
Do you want to come along with me?.
[NON-ENGLISH SPEECH].
It's just a very simple action..
Just walk with him for a while and chat with him for a bit..
[NON-ENGLISH SPEECH].
Gradually, they'd become friends..
And he invited the boy to play basketball with him.
and participate in activities together..
[NON-ENGLISH SPEECH].
For this student, that day was just a very small episode.
in the day..
[NON-ENGLISH SPEECH].
But years later, he suddenly realized.
that that day had changed Jimmy's life..
[NON-ENGLISH SPEECH].
At the graduation ceremony, Jimmy stood on the stage.
and gave a speech..
[NON-ENGLISH SPEECH].
He said the reason why he brought all the books home.
from school that day was because--.
[NON-ENGLISH SPEECH].
It was because he had already prepared to end his life.
that weekend..
[NON-ENGLISH SPEECH].
Books no longer needed to be left at school anymore..
[NON-ENGLISH SPEECH].
However, that afternoon, when the classmates walked over--.
[NON-ENGLISH SPEECH].
He knelt down to help him pick up the books one by one.
and walk him home..
[NON-ENGLISH SPEECH].
That was the first time he truly felt--.
[NON-ENGLISH SPEECH].
It turns out that there are people in this world who.
see me and care about me..
[NON-ENGLISH SPEECH].
It's through this gradual acceptance and care that--.
[NON-ENGLISH SPEECH].

$^{81}$This weighed him from taking his own life.
and led him down to a completely different path..
[NON-ENGLISH SPEECH].
Brothers and sisters, what may seem like this.
is a very small action to us--.
[NON-ENGLISH SPEECH].
A greeting, a smile, a moment of companionship,.
this can all be the turning point in one's person's life,.
the transition from death to life..
[NON-ENGLISH SPEECH].
Today, we're going to take a look at Gospel of John..
[NON-ENGLISH SPEECH].
Let's take a look how God uses seemingly very small things.
in our daily life to touch people's heart.
and even change their whole life..
[NON-ENGLISH SPEECH].
In John number two, Jesus did something.
that seemed insignificant, yet it revealed the glory of God..
[NON-ENGLISH SPEECH].
When I was really young, I heard about this interpretation.
of a scripture in church..
Jesus is a God who can change matter,.
able to transform water, H2O, into alcohol, C2H5O,.
adding many more carbon and hydrogen atoms..
[NON-ENGLISH SPEECH].
Jesus is truly amazing, a God who can change matter..
However, today, if we consider Jesus' situation.
in the Gospel of John and the author's description of Jesus.
performing his first miracle--.
[NON-ENGLISH SPEECH].
--we will discover that this God is not only.
a God who changes matter, but even more so,.
a God who changes people's hearts..
[NON-ENGLISH SPEECH].
Let's look at the first miracle Jesus.
performed in Gospel of John..
[NON-ENGLISH SPEECH].
On the third day, there was a wedding at Cana in Calili..
The mother of Jesus was there..
Jesus also was invited to the wedding with his disciples..

$^{121}$When the wine ran out, the mother of Jesus said to him,.
they have no wine..
[NON-ENGLISH SPEECH].
And Jesus said to her, woman, what.
does this have to do with me?.
My hours has not yet come..
His mother said to the servants, do whatever he tells you..
Now, there were six stone water jars.
there for the Jewish rites of purification,.
each holding 20 or 30 gallons..
[NON-ENGLISH SPEECH].
Jesus said to the servant, fill the jars with water..
And they filled them up to the brim..
And he said to them, now, draw some out and take it.
to the master of the feast..
So they took it..
When the master of the feast tasted the water,.
now it become wine and did not know where it came from,.
though the servants who had drawn the water knew it..
And the master of the feast called the bridegroom..
[NON-ENGLISH SPEECH].
And said to him, everyone serves the good wine first..
And when people have drunk freely, then the poor wine..
But you have kept the good wine till now..
This, the first of his signs, Jesus did at Cana in Galilee.
and manifested his glory..
And his disciples believed in him..
After this, he went down to Capernaum.
with his mother and his brothers and his disciples..
And they stayed there for several days..
[NON-ENGLISH SPEECH].
Why did Jesus' mother said to him, they have no more wine?.
And why did Jesus react so strong?.
[NON-ENGLISH SPEECH].
In Jewish culture at the time, a wedding feast without wine.
was a profound humiliation, seen as failure.
on the part of the groom's family, a breach of etiquette,.
and a social disaster that will be remembered for a lifetime..
[NON-ENGLISH SPEECH].
Mary, Jesus' mom, was likely a relative or very close friend..

$^{161}$Her being mentioned first and seemingly moving in and out.
of the inner court suggests that she.
was more concerned with the host's reputation and glory.
than others, keenly noticing the underlying difficulties..
[NON-ENGLISH SPEECH].
This reflects her initial belief in Jesus' ability to act..
Secondly, her expectation that Jesus could help.
restore the host family glory..
[NON-ENGLISH SPEECH].
Jesus responded, woman, what does this have to do with me?.
My time has not come yet..
Mary's expectation was that Jesus' intervention.
indirectly influenced Jesus to reveal God's glory..
[NON-ENGLISH SPEECH].
In the Gospel of John, time is a crucial word..
When mentioned, it's usually associated with glory..
Because Jesus knew that when he performed a miracle,.
he would reveal God's glory..
Therefore, Jesus' reaction was quite urgent..
[NON-ENGLISH SPEECH].
But what Jesus did next transformed.
the most uncleaned water, the most despised water,.
into something that would bring glory to his master.
and to Jesus himself..
[NON-ENGLISH SPEECH].
Versus six..
According to the Jewish ritual of purification,.
there were six stone jaws placed there..
[NON-ENGLISH SPEECH].
Some Chinese English translation.
rendered this as a stone jaw used for Jewish purification.
rituals..
[NON-ENGLISH SPEECH].
In fact, Jews placed great importance.
on cleanliness rituals because they.
understand that cleanliness regulations God.
gave in Exodus and Leviticus were intended.
to show the world that Jewish people belong.
to Jehovah and to God..
[NON-ENGLISH SPEECH].

$^{201}$Therefore, all the derived purification rituals.
are meant to show one's relationship with God..
[NON-ENGLISH SPEECH].
Even today, if you have the opportunity.
to visit Western Wall in Jerusalem,.
you will see Jews still observing purification.
rituals in front of the Western Wall..
There is a hand washing station with a sewer.
cup in the basin..
[NON-ENGLISH SPEECH].
They use that sewer cup to hold water,.
symbolizing the stone jaws of the pest.
to clean their hands before going to see God,.
thus demonstrating their belonging to God..
Therefore, you can imagine how much importance Jews.
place on purification rituals..
[NON-ENGLISH SPEECH].
When you came to see God today, did you.
get the chance to wash your hands and face?.
[NON-ENGLISH SPEECH].
Jesus, however, used the stone jaw.
used for purification rites to fill them and turn the water.
into wine..
This shattered the Jewish and the disciples' preconceived.
notion about the purification rites..
[NON-ENGLISH SPEECH].
The jaws used for purification were actually.
filled with uncleanliness..
Imagine the water in those jaws was.
used for washing hands, face, and even feet..
[NON-ENGLISH SPEECH].
People came from all directions, traveling long distance.
to this wedding feast..
[NON-ENGLISH SPEECH].
And at that time, technology wasn't as advanced.
as it is now..
Many people weren't wearing shoes..
[NON-ENGLISH SPEECH].
You can imagine how filthy and unclean.
that water could be, filled with mud..

$^{241}$[NON-ENGLISH SPEECH].
But Jesus transformed these most unclean and filthy water.
into the finest wine, better than any wine.
served at the entire wedding feast..
[NON-ENGLISH SPEECH].
Here, Jesus not only changed the material itself,.
but also shattered the stereotype of cleanliness,.
turning the unclean water meant to be received into a tool.
to glorify Jesus..
[NON-ENGLISH SPEECH].
Glorify Jesus..
[NON-ENGLISH SPEECH].
So on the level of personal spiritual life--.
[NON-ENGLISH SPEECH].
These water jars were originally intended for ritual purification,.
symbolizing how people use our behavior and religious rituals.
to slightly cover up their own defilement..
[NON-ENGLISH SPEECH].
Jesus, however, turned the water into wine,.
symbolizing that he used his own blood and his own life.
to become the true source of purity for humanity..
[NON-ENGLISH SPEECH].
The most shameful, unspeakable, and unclean parts of our lives,.
whether failure, addictions, sins, or broken relationship,.
are not forgotten in the gospel, but transformed by him.
into good stories that testify to his grace..
[NON-ENGLISH SPEECH].
In the grace of Christ, we can learn.
to surrender our most unbearable moments,.
no longer cover them up with religious rituals.
or superficial good deeds, but allowing the blood of Jesus.
to truly enter and renew us..
[NON-ENGLISH SPEECH].
There are some things, people and relationship,.
that we might not even want to recall,.
feeling they are too painful to look back on..
We might think that the time will heal all the wounds,.
that they will simply fade away..
[NON-ENGLISH SPEECH].
Sometimes we might be able to cover up these things.

$^{281}$with outward act of kindness or even righteous, religious.
righteous..
[NON-ENGLISH SPEECH].
But Jesus' most unexpected action.
was to transform these most unclean water,.
representing our most unspeakable moments in life,.
into an instrument for glorifying God..
[NON-ENGLISH SPEECH].
If we are willing to offer it, we.
can expect God to transform our unbearable moments in life.
into moments and tools for his own glory..
[NON-ENGLISH SPEECH].
Do you want to do it?.
[NON-ENGLISH SPEECH].
When we talk about work, we expect.
God to intervene in our works as well..
[NON-ENGLISH SPEECH].
I don't know, starting from when,.
we've begun to lose the hope in God's presence in our lives.
and work..
[NON-ENGLISH SPEECH].
Our only hope might be just, let me finish today's work.
as quick as possible..
[NON-ENGLISH SPEECH].
Just like the jewel of gold, even on important occasions.
like weddings, we perform numerous purification rites,.
just completing these elaborated rites is exhausted itself..
[NON-ENGLISH SPEECH].
I don't know starting from when, we've.
stopped expecting God to intervene in our work..
It is as if our work has nothing to do with God..
I'm just here to earn money..
I'm just here to need a job to support my family..
Besides that, my job has no other meaning..
[NON-ENGLISH SPEECH].
What is the most astonishing here about Jesus.
is that he transformed the water used for washing feet, which.
is the jewels despised and considered utterly meaningless.
into a source of glory for the host and for himself..
[NON-ENGLISH SPEECH].

$^{321}$Seemingly, he would turn the seemingly meaningless work.
in our lives into a means to glorify himself..
[NON-ENGLISH SPEECH].
The question is, are you willing to let.
God intervene in our lives?.
[NON-ENGLISH SPEECH].
Or we just simply want God to leave us alone..
Let's have more expectations and think more about how.
God can intervene in our work..
[NON-ENGLISH SPEECH].
In our faith practice, our sinner life must become real..
[NON-ENGLISH SPEECH].
This water was prepared according.
to Jewish ritual of purification and is.
a part of the Old Testament ritual system..
[NON-ENGLISH SPEECH].
Jesus did not use new vessels, but worked directly.
in these purification jars, turning the water into wine..
This symbolized that he did not add a bunch of new rules,.
but rather directed all the purification, sacrifice,.
and festivals of the Old Testament.
to a salvation he accomplished on the cross..
[NON-ENGLISH SPEECH].
For the church, this sign calls on people.
not to merely outwardly observe righteous,.
but to enter into the abundant faith in Christ..
[NON-ENGLISH SPEECH].
Worship, baptism, communion, and fellowship,.
if they lose the inner reality of being cleansed and renewed.
in Christ, these may be just reduced to mere righteous..
[NON-ENGLISH SPEECH].
In this miracle, Jesus invites us.
to prepare empty jars in every festival and gathering,.
trusting in his intervention and waiting for him.
to transform ordinary water and ordinary life.
into wine full of grace and joy..
[NON-ENGLISH SPEECH].
Remember the story that I mentioned at the beginning?.
[NON-ENGLISH SPEECH].
Jimmy stopped his suicidal thoughts.

$^{361}$because someone showed him a small act of kindness.
by wanting to take care of him..
[NON-ENGLISH SPEECH].
This is actually a result of God making.
a person's inner life authentic..
[NON-ENGLISH SPEECH].
Caring for others isn't about external rules.
and regulations, but a genuine feeling.
from the depth of one's heart..
[NON-ENGLISH SPEECH].
God used this small act to save a person's life..
[NON-ENGLISH SPEECH].
Transforming an ordinary life into a wine full of surprises..
[NON-ENGLISH SPEECH].
Let's pray together..
[NON-ENGLISH SPEECH].
Dear Heavenly Father, thank you not only.
for changing the material world, but also.
for changing our inner lives..
[NON-ENGLISH SPEECH].
Let us learn and look forward to you transforming.
the hardships and brokenness in our lives.
into witness to your glory..
[NON-ENGLISH SPEECH].
We pray in the name of Lord Jesus, amen..
[NON-ENGLISH SPEECH].
\newpage



\section{約翰福音 5:1-18}
\label{sec:a3FPlfJ6MYc}
\textbf{屬靈的力量 (約翰福音 5\_1-18) - 鍾智愷傳道}
\newline
\newline
連結: \href{https://youtube.com/watch?v=a3FPlfJ6MYc}{\texttt{https://youtube.com/watch?v=a3FPlfJ6MYc}} ~~~~ 語音日期: 2026-03-24
\newline
\newline
\hyperref[sec:TQN008eNyL0]{\small{< < < PREV SERMON < < <}}
~
\hyperref[sec:index_chronic]{\small{[返順時目]}}
~
\hyperref[sec:index_scriptual]{\small{[返順卷目]}}
~
\hyperref[sec:_FBi_OEMziA]{\small{> > > NEXT SERMON > > >}}
\newline
\newline
約翰福音 5:1-18
\newline
\begin{longtable}{cl}
\hline
\hline
章節 & 經文 (和合本修訂版)\\
\hline
5:1 & \begin{tabularx}{0.7\textwidth}{X} 這些事以後，到了猶太人的一個節期，耶穌上耶路撒冷去。 \end{tabularx} \\ \\ \relax
5:2 & \begin{tabularx}{0.7\textwidth}{X} 在耶路撒冷，靠近羊門有一個池子，希伯來話叫畢士大，旁邊有五個柱廊； \end{tabularx} \\ \\ \relax
5:3 & \begin{tabularx}{0.7\textwidth}{X} 裡面躺著許多病人，有失明的、瘸腿的、癱瘓的。等候水動， \end{tabularx} \\ \\ \relax
5:4 & \begin{tabularx}{0.7\textwidth}{X} 因為有天使按時下池子攪動那水，水動之後，先下水的人無論患甚麼病都能得痊癒。 \end{tabularx} \\ \\ \relax
5:5 & \begin{tabularx}{0.7\textwidth}{X} 在那裡有一個人，病了三十八年。 \end{tabularx} \\ \\ \relax
5:6 & \begin{tabularx}{0.7\textwidth}{X} 耶穌看見他躺著，知道他病了很久，就問他：「你要痊癒嗎？」 \end{tabularx} \\ \\ \relax
5:7 & \begin{tabularx}{0.7\textwidth}{X} 病人回答他：「先生，水動的時候，沒有人把我放在池子裡；我正要去的時候，別人比我先下去了。」 \end{tabularx} \\ \\ \relax
5:8 & \begin{tabularx}{0.7\textwidth}{X} 耶穌對他說：「起來，拿起你的褥子走吧！」 \end{tabularx} \\ \\ \relax
5:9 & \begin{tabularx}{0.7\textwidth}{X} 那人立刻痊癒，就拿起自己的褥子走了。那天是安息日， \end{tabularx} \\ \\ \relax
5:10 & \begin{tabularx}{0.7\textwidth}{X} 所以猶太人對那被治好了的人說：「今天是安息日，你拿褥子是不合法的。」 \end{tabularx} \\ \\ \relax
5:11 & \begin{tabularx}{0.7\textwidth}{X} 他卻回答他們：「那使我痊癒的人對我說：『拿起你的褥子走吧！』」 \end{tabularx} \\ \\ \relax
5:12 & \begin{tabularx}{0.7\textwidth}{X} 他們問他：「對你說『拿起褥子走』的是甚麼人？」 \end{tabularx} \\ \\ \relax
5:13 & \begin{tabularx}{0.7\textwidth}{X} 那治好了的人不知道那人是誰，因為那裡人很多，耶穌已經躲開了。 \end{tabularx} \\ \\ \relax
5:14 & \begin{tabularx}{0.7\textwidth}{X} 後來耶穌在聖殿裡找到他，對他說：「你已經痊癒了，不要再犯罪，免得你的遭遇更壞。」 \end{tabularx} \\ \\ \relax
5:15 & \begin{tabularx}{0.7\textwidth}{X} 那人就去告訴猶太人，使他痊癒的是耶穌。 \end{tabularx} \\ \\ \relax
5:16 & \begin{tabularx}{0.7\textwidth}{X} 所以猶太人迫害耶穌，因為他在安息日做了這些事。 \end{tabularx} \\ \\ \relax
5:17 & \begin{tabularx}{0.7\textwidth}{X} 耶穌就回答他們：「我父做事直到如今，我也做事。」 \end{tabularx} \\ \\ \relax
5:18 & \begin{tabularx}{0.7\textwidth}{X} 為了這緣故，猶太人越發想要殺他，因為他不但犯了安息日，而且稱神為他的父，把自己和神看為同等。 \end{tabularx} \\ \\
[1ex]
\hline
\hline
\end{longtable}
$^{1}$弟兄姊妹早晨.
好高興今日可以同大家一齊去敬拜主.
我們一齊去唱詩.
我們去讚美我們的主.
今日我們的講題是屬靈的力量.
我們會去看約翰福音的五章一到十八節.
如果你有你的聖經.
你也可以去翻開到約翰福音五章一到十八節.
約翰福音五章一到十八節.
不知道大家講起屬靈的力量.
你會想起甚麼.
我們一聽屬靈的力量.
很多時候會想起一些比較神秘的東西.
例如我們會想起好似Harry Potter.
你有些咒術可以去施魔法.
或者我們東方人會想起一些風水的東西.
我們會想起一些風水命理.
可能放一塊石頭在某個位置.
就可能加你的運氣.
我們見到西方也有很多這些.
他們喜歡玩塔羅牌.
我想特別在現在這個世代.
有一樣東西叫做新紀元運動.
是甚麼來的呢.
它是由印度教裡面出來的.
它說在甚麼東西裡面都有些能量.
你透過攝取那些能量.
你就能夠有些不同的運氣.
它跟你說甚麼水晶.
就不知道跟你帶來甚麼運.
它還要說有甚麼磁場.
這個世界你會發現它也很執著.
在一些屬靈的力量裡面.
你會發覺這個世界的人.
他也很想要一些屬靈的力量.
很想去操縱一些東西.
而當然我們基督徒我們要很小心.
我們不會去跌入這些陷阱裡面.
我們不會去做這些事情.
但我想我們今天也要有一個小心的點.

$^{41}$就是我們也發現.
原來屬靈是一樣真實的東西.
原來屬靈的世界是一樣真實的東西.
我們當然不會跟他們一樣.
去濫用一些屬靈的力量.
但原來我們今天也沒有辦法.
去否定屬靈的真實.
屬靈的力量是真實的.
我們今天想去問一個問題.
究竟屬靈的力量在哪裡呢.
我們基督徒既然我們相信神.
我們相信屬靈世界的存在.
究竟今天我們基督徒.
可以怎樣去經歷這些屬靈的力量呢.
當我們不是玩風水.
我們不玩這些水晶.
我們不玩塔羅.
我們不相信這些儒家冥想.
我們可以怎樣去經歷屬靈的力量呢.
我們一起有個禱告好嗎.
我們求上帝帶領我們.
去進入真理裡面.
主啊我們去.
在你的面前我們求你的恩典.
我知道唯有靠著你的恩典.
我們才能夠明白你聖經的真理.
所以主啊我們求你.
今天帶我們進入約翰福音裡面.
求你讓我們更加明白你.
讓我們知道.
能夠怎樣在你裡面得到力量.
得到那份恩典.
主啊我們求你去幫助.
開我們的心.
讓我們去領受你自己的話語.
以神禱禱告.
是奉主耶穌的名祈求.
阿們.
我們今天會先看一段經文.
在我們去做一些反思之前.

$^{81}$我們一起享受一下這段經文.
我們享受一下故事.
我們一起看一下.
究竟耶穌做了些什麼.
我們今天看的約翰福音第五章.
是耶穌做的第三個神蹟.
如果你心水清的話.
你會知道約翰福音有七個神蹟.
七是完全的意思.
它說的是耶穌完全地啟示自己.
而今天我們看的是第三個.
如果你有印象.
第一個是水變酒.
是一個叫加拿芬宴的神蹟.
第二個是在約翰福音的第四章結尾.
是一個醫治的神蹟.
最奇怪的是.
在約翰福音的第五章開頭.
又是一個醫治的故事.
如果你看得勝經多.
你就在想.
為什麼聖經的東西這麼重複.
聖經的東西好像來來去去都一樣.
我們今天要看到的其實不一樣.
今天這個不單止是一個醫治的故事.
我們就一起進入經文裡面.
在第一至五節.
他交代了經文的背景.
他說這些事以後.
就是上一個神蹟之後.
到了猶太人的一個節期.
耶穌就去到耶路撒冷.
在耶路撒冷靠近羊門有一個池子.
希伯來話叫畢氏大.
旁邊有五個柱廊.
裡面躺著很多病人.
有失明的瘸腿的癱瘓.
你會看到這個地方是一個挺大的地方.
裡面裝著了很多病人.
在這一刻我們還不知道為什麼這麼多病人.

$^{121}$但是第四節他給了一些提示我們.
他說這些病人.
在等候池裡面的水去動.
因為有天使按時下池子.
搞動那個水.
水動之後先下水的人.
無論有什麼病都可以痊癒.
你會看到第四節是卡住了.
原因是今天很多的聖經學者.
都會說這一段不是原本裡面的.
他會說這一段是後加上去的.
原因可能是他想解釋一下.
當時有什麼的傳說.
在這個池裡面的人有什麼信念.
但是無論他怎樣也好.
他解釋了一件事給我們聽.
就是當時的人有一種迷信.
他相信原來池子是有特別的功效的.
他相信那個池子是你浸下去.
如果你是第一個浸下去的話.
你就能夠得到醫治了.
接著我們看到經文裡的主角.
在那裡有一個人.
他病了38年.
弟兄姊妹我今年26歲.
如果他病了38年的話.
我要再病多12年.
大家有沒有想過這是多長的時間.
我們有時候說人生可能70歲.
一半的時間都在一個病患裡面.
動不了.
對於這個病人來說.
是一個很絕望的境地.
我們看看耶穌怎樣跟他互動.
耶穌看到他躺在那裡.
知道他病了很久.
就問他你要痊癒嗎.
耶穌直接去跟他說.
其實我能夠給你一個痊癒.
你要不要痊癒.

$^{161}$我們看看病人怎樣回答.
他說先生在水冷的時候.
沒有人把我放在池子裡.
我要去的時候.
人家就讓我先下去了.
病人好像沒有直接回應耶穌的問題.
耶穌說你要痊癒嗎.
他說先生沒有人把我放下去.
他背後的意思是.
其實他想要痊癒.
但他覺得醫治的能力在池子裡.
他覺得醫治的能力在水裡.
所以他只信了一句跟耶穌說.
沒有人把我放在水裡.
所以我到現在還沒有痊癒.
我們看看耶穌說什麼.
耶穌跟他說.
起來拿你的玉子走吧.
那個人就立刻痊癒 拿起自己的玉子走了.
如果你是看一個故事的話.
整件事其實是很快的.
加多一句那個人就已經痊癒了.
那個人等了38年.
他等的原來只是耶穌的一句話.
耶穌說一句他就立即痊癒了.
38年和立即.
對於那個人來說是很震撼的事.
他等了38年的痊癒.
在耶穌的一句話裡.
他立即得到.
我們在這裡看到.
屬靈的力量不是在池子裡.
屬靈的力量是在耶穌的說話裡.
是在耶穌自己裡.
如果故事是這樣完就好了.
如果接著的故事就是.
那個人就信耶穌成為主的門徒.
一切都簡簡單單平安無事了.
是不是.
我們很遺憾的事情是沒有這樣完結的.

$^{201}$接著經文說到猶太人出場了.
我們每一次見到猶太人.
在約翰福音裡面.
我們都要心感不妙的.
因為猶太人每一次都是和耶穌作對那個.
當你看到猶太人.
你就知道有些不好的事情要發生了.
發生什麼事呢.
在第九節尾那裡說到.
那一天是安息日.
所以猶太人對那個被醫好的人這麼說.
今天是安息日.
你不可以拿起你的玉紙走.
玉紙就是那張墊.
你不可以拿起那張墊去走.
作者在這裡很諷刺.
他將那個病人叫做被醫好的人.
他特意強調他是被醫好的.
但是猶太人好像完全看不到一樣.
猶太人好像完全瞎了.
看不到那個人是被醫好的.
他只看到你違反安息日.
你犯了法.
就只有這麼多.
他只看到這麼多.
那個人怎麼回應呢.
他說那個使我痊癒的人跟我說.
拿起你的玉紙走吧.
猶太人問他是誰.
他就不知道了.
因為耶穌已經躲起來了.
這個人呢.
他也是講事實的.
是耶穌叫他去走的.
但是你會看到.
這個很明顯都是一個推卸責任.
你犯了這件事.
他說不是.
是醫好我那個人叫我做的.
不關我事的.

$^{241}$我們中文有一句話叫割蓆.
立刻跟耶穌割蓆.
他叫我做的不關我事的.
如果事情是這樣完.
都OK的.
那就是大家都相安無事.
耶穌有多特別.
耶穌親自去找他.
耶穌親自去找這個人.
跟他說.
你已經痊癒.
就不要再犯罪了.
免得你的遭遇更壞.
這裡有一樣很特別的事.
在約翰福音的頭四章裡面.
從來都沒有說過.
任何一個人個別地犯罪.
我們看到約翰福音第一個犯罪的人.
在這裡.
跟他說你不要再犯罪了.
免得你遭遇的更壞.
究竟這個人犯了什麼罪.
為什麼說他是一個罪人.
我們現在不知道.
看完下一節就知道了.
那個人就去告訴猶太人.
使他痊癒的是耶穌.
請大家留意.
這裡不是一個很中性的形容.
猶太人剛剛才說.
那個治癒他的人犯了安息日.
這個被人治好的人.
他立即跟猶太人去報串.
去打小報告.
說治好他的人是耶穌.
這是一個出賣.
是被治好的人出賣了耶穌.
所以十六節我們立即看到.
那些猶太人就因為這樣迫害耶穌.
因為他在安息日做了這些事.

$^{281}$耶穌的回應都很特別.
耶穌說我父做事直到如今.
我都做事.
這是一個.
耶穌在引用一個猶太人.
他們很多的討論的一個辯論點.
他們說既然這個世界.
是藉著神的工作去維繫的話.
那不可能上帝在安息日不做事.
如果不是.
世界就在安息日裡面會倒塌.
但是除了這個耶穌的引用之外.
你會發現原來耶穌和父神.
是一起去工作的.
我們想想究竟父神是怎樣工作的.
父神是透過說話去工作.
其實耶穌在安息日做了些什麼.
去說一句話.
耶穌做的事就是說一句話.
就只有這麼多.
但是耶穌說是父神做事的同時.
我都做事.
原來耶穌和父神是.
在說話裡面一起去做事.
你會發現原來父神設立安息日.
是為了讓土地得到休息去滋養生命.
耶穌所做的都是在安息日裡面.
去醫治去滋養生命.
父神和聖子所做的事是一樣的.
但是跟著十八節我們見到.
一個很諷刺的東西.
當耶穌說.
父神在安息日我都做事的時候.
那些猶太人就更加想要殺他.
因為他不單止犯了安息日.
而且稱神為他的父.
將自己和神看為同等.
為什麼我說這裡很諷刺呢.
是因為那些猶太人說耶穌犯了誡命.
他說耶穌你犯第四條誡命.

$^{321}$你不就是要守安息日.
耶穌你犯了第一條誡命.
你不可有別的神.
但你當自己是神.
但是看看那些猶太人.
他自己都犯了第六條誡命.
第六條誡命說不可殺人.
當他們說耶穌犯誡命的同時.
他自己都犯了所有誡命.
這個是約翰福音給我們看到的諷刺點.
弟兄姊妹我們看完《亂經文》之後.
我們集中去看.
為什麼病人經歷了屬靈的力量.
經歷了醫治之後.
為什麼他沒有應有的信心.
為什麼他明明經歷了這麼多.
38年的病他都能夠痊癒.
為什麼他沒有生出對上帝應有的敬畏.
我們能夠從對那個詞的迷信.
我們可以有一點點的痕跡.
我想我們先去想一想什麼叫迷信.
迷信就是我們覺得屬靈的力量.
背後其實是沒有神的.
屬靈的力量是一些我們可以操控的東西.
例如我們可以透過唸咒語.
我們透過擺一塊石頭.
透過戴一些首飾.
迷信就是我們相信屬靈的力量.
背後跟神是沒有關係的.
是一些我們可以操控的東西.
那個病人他的問題正正就是.
他覺得那個醫治的力量.
是在那個池裡面.
在那個水裡面.
而不是在上帝裡面.
這樣的話我們會能夠明白.
他怎麼去看耶穌.
對於他來說.
耶穌就是一個能夠操控屬靈力量的人.
不代表他背後是有神的.

$^{361}$如果我們用一個比喻.
我想大家說起一些風水師就很明白.
我們說起例如蘇文峰,麥玲玲.
如果你今天請他幫你看風水.
我們不要這樣做.
只是假設.
然後他幫你看完一輪.
佈置完一輪.
你會不會當他是神來拜.
當然不會.
你會不會說我要跟你建立個人的關係.
我要跟你建立一個親密的關係.
你也不會.
這個是交易.
我們覺得他們是在操控屬靈的力量.
就是這麼多.
我們不會當他是神.
要跟他建立一個關係.
我想那個病人去看耶穌也是一樣.
他覺得耶穌雖然是有能力.
不代表他背後是有神的.
這個就是他最大的問題.
如果用一個圖我們去看的話.
我們會知道其實父神差遣耶穌來到世界上.
耶穌和父神是一起去工作的.
所以耶穌有醫治人的能力.
耶穌也期望人.
因為這樣對上帝有應有的敬畏.
但對於那個病人來說.
他中間的連結斷了.
他看不到耶穌背後有父神.
自然也看不到他需要對耶穌有一種敬畏.
他看不到神在耶穌的背後.
如果我們多看一段經文.
你就會更加明白.
究竟為什麼會有這樣的事情發生.
我純粹做一個很快的對比.
在約翰福音9章.
也是一個醫治的故事.
差不多所有事情都是一模一樣的.

$^{401}$你會看到耶穌主動去醫治一個乞子.
你會看到這個故事裡面牽涉一個詞.
也是一個水潭.
你會看到兩個人被耶穌醫治了之後.
都不知道耶穌在哪裡.
約翰福音故意將兩個好像平衡一樣.
給我們去看.
但不一樣的是什麼呢.
耶穌講明這個人是沒有罪的.
耶穌講明這個人是一個義人.
他們兩個的分別在哪裡呢.
我想透過很快的這段經文我們看一看.
在24節他這樣講到.
我們蔡人就叫那個以前真名的人跟他說.
你要將榮耀歸給神.
我們知道這個人是一個罪人.
在說耶穌.
因為他們覺得耶穌犯案昔日.
所以又覺得他是一個罪人.
我們看看究竟這個人是怎麼回答.
他們互動了一段時間.
我們不詳細去看了.
我們看看那個瞎了的人的結論是什麼.
從30節那裡去這樣講.
那個人回答對他們說.
他開了我的眼睛.
你們竟不知道他從哪裡來.
這真是奇怪.
我們知道神不聽罪人.
唯有敬奉神.
遵行他旨意的神才聽他.
從創世以來.
未曾聽見有人開了生來就失明的人的眼睛.
這個人如果不是從神那裡來的.
就什麼都不能做.
我想在這裡我們明白那個對比.
對於這個核子來說.
去見到耶穌是從父神那裡來的.
去見到耶穌醫治的能力.
不是可以隨便操控的.

$^{441}$而是按著父神的心意一起去走的.
我想如果我們在這段經文裡.
我們要去看到一個屬靈的反思.
就是我們見到真正屬靈的力量.
一定是會按著上帝的心意的.
真正屬靈的力量.
不是一個可以我們隨意操弄.
好像一些魔法這樣的概念.
真正屬靈的力量.
只能夠在父神的心意裡面來.
只能夠在上帝的心意裡面來.
當我們明白經文是想傳遞一個這樣的訊息的時候.
我們一起去看一下.
怎樣在我們生命裡面.
在我們生活裡面去應用.
我想我們今天都面對同一個試探.
我們很容易將屬靈的力量.
放在一些物件裡面.
放在一些死物裡面.
放在一些方法裡面.
而不是將它放在上帝裡面.
我們見到這個習慣.
是由古時開始已經在發展.
你會見到猶太人喜歡做什麼呢.
打仗的時候.
他們怕打不贏.
他們就說.
上帝的藥櫃是有一個特別的力量的.
那個藥櫃有特別的功效的.
打仗的時候抬他出去.
就會打勝仗.
最後當然是打輸.
他們完全是搞錯了.
上帝的力量從哪裡來.
你見到掃羅也是.
以色列的第一個君王掃羅.
他在打仗之前.
他說我一定要獻祭.
因為他覺得他獻完祭之後.
上帝有特殊的力量會幫他打勝仗.

$^{481}$誰知道薩姆爾遲到.
雖然好像是沒有遲到.
但是他就急著一定要獻祭.
他說死了我不獻祭打不贏.
怎麼辦.
他雖然不是祭司.
他自己匆匆忙忙地獻了一個祭.
薩姆爾說你這樣做不按上帝的心意.
你整個國家都沒了.
你見到後來的猶太人.
他以為屬靈的力量在聖殿裡面.
他說有聖殿我們就不會死.
有聖殿就什麼事都沒有.
耶利米就跟他們說.
不是的.
你們叛逆上帝的心意.
你們不要以為有聖殿就有用.
上帝要差遣他的審判來到.
我們見到今天我們基督新教以外.
都有很多不同的儀式.
好像在說屬靈的力量在一些死物裡面.
例如你會聽到有些人有聖水.
有些聖物.
十字架有祝福在裡面.
將屬靈的力量和上帝的心意割開.
說那些東西是有特殊的功效的.
我們要去小心這些東西.
我們今天作為基督徒.
我們有沒有一些這樣的誘惑.
有沒有一些這樣的試探在裡面.
我想今天我們最大的試探是.
我們祈禱的時候.
我不知道大家私人祈禱會祈什麼禱.
很慚愧地說.
我年輕的時候 可能十年前.
祈禱最多的是求上帝給我老婆.
十年前不斷地去祈禱.
我覺得祈禱得多就會快一點.
祈禱得多就好像上帝要這樣給我.
可能我上班之後.

$^{521}$我身邊的人最多的祈禱就是.
上帝幫我找一份工作.
上帝給我一份好一點的工作.
是一份高薪厚職的.
上班不要太遠的工作.
可能大家隨著年紀漸長.
我也知道很多人很喜歡求平安.
出入平安.
我希望每一天都無風無雨 順順利利.
在我們這些祈禱的背後.
我們是不是覺得我們這樣祈禱.
就好像操弄一些屬靈的力量.
要上帝祝福給我們.
我們在這樣祈禱的時候.
有沒有考慮過.
究竟上帝的心意是什麼呢.
我不是說大家不可以這樣祈禱.
我相信我們需要真誠地.
在上帝面前去祈求.
但是我們要留意的是.
我們不要把這些祈禱.
當作是操縱上帝的工具.
我們在這些祈禱裡面.
我們是去聆聽 去測驗上帝的心意.
我們是跟著上帝的心意去行.
而不是要上帝按著我們的心意去行.
如果我們看著耶穌的主禱文.
我們就會知道在禱告裡面.
第一重要的是.
我們在天上的父.
是願人都尊你的名為聖.
是願上帝的國降臨.
是願上帝的旨意.
行在地上如同行在天上.
我們今天如果要祈禱的話.
如果我們真的想我們的祈禱有功效的話.
我們就不能再按著我們的心意而行.
因為上帝的力量不在那裡.
上帝的力量是跟著祂的心意去行事的.
剛剛麗倩也說.

$^{561}$我們祈禱的時候.
我們不是好像泥菩薩過江那樣.
求上帝導我們去不知那裡.
我們今天所信的不是一個黃大仙式的信仰.
我們是降伏在主的面前.
跟隨祂的心意.
所以我們今天.
如果我們負面來講.
我們就不要將屬靈的力量放在錯的地方.
正面來講.
就是真正屬靈的力量.
只可以在執行神的旨意裡面去領受.
弟兄姊妹.
我們今天有很多不同的力量在我們中間.
有政治的力量.
文化的力量.
我們甚至都相信.
其實魔鬼的力量是真實的.
所以有很多人能夠做到很多不同的事.
甚至經文裡面講到.
醫病趕鬼.
都是可能這個世界有人會做到的.
但是我們聽聽馬太福音怎樣去講.
他說不是每一個稱呼我主啊主啊的人.
都能夠進天國.
是唯有遵行我父子義的人.
才能夠進去.
是那天很多人會跟我說.
主啊主啊.
我們不是奉你的名傳道嗎.
奉你的名趕鬼嗎.
奉你的名行許多義能嗎.
他要跟他們宣告.
我從來都不認識你們.
你們走開吧.
如果我們今天想要經歷真正的屬靈力量.
我們就必須要去遵行上帝的心意.
真正的屬靈力量.
只可以在上帝的心意裡面找到.
可能在我們悔改裡面.

$^{601}$我們能夠經歷聖靈的大能.
去給我們一個新的生命.
在全福音裡面.
我們經歷到上帝給我們的那種大能.
去讓我們去作見證直到地極.
我們可能在教會裡面的事奉.
我們經歷到屬靈的那種恩賜.
我們能夠去事奉上帝.
事奉上帝的子民.
我還很記得.
我在年幼的時候.
都曾經有追過星.
我不知道大家有沒有追過星.
當年已經做了基督徒.
但是追星追得很入迷.
我都知道不對.
我就跟上帝說.
上帝不如我祈多點禱.
我讀多點經.
其實這樣我都算是一個好的基督徒.
我怎樣追星其實你都不要理我.
我都做盡了我可以做的事.
不要再要求我那麼多.
每當我越祈禱越讀經的時候.
心裡面都很痛.
因為我感覺到上帝離我很遠.
我感覺到上帝不喜悅這件事.
但是我都不斷用這種.
屬靈的操練去麻醉自己.
以為我這樣做多點.
上帝那個屬靈的力量還在我裡面.
直到有一天我真的發現.
跟上帝真的很遠.
然後在一個港島裡面.
那個講者呼召.
叫大家你要全心全意去敬拜上帝.
不可以再留一絲餘地.
那一刻我真的經歷屬靈的大能.
當我祈禱那麼久讀經那麼久.
都經歷不到上帝的能力的時候.

$^{641}$在那一刻我發覺.
上帝的大能在推動我.
要我作出一個禱告.
去放下這些不合他心意的事情.
在那一刻真的很辛苦.
我想差不多十分鐘左右.
我都沒有舉手.
我都沒有回應那個呼召.
但是最後上帝都推動我去做這件事.
在那一刻上帝的大能是很真實的.
我想我們今天如果做一個總結.
就是屬靈的力量是在上帝的心意裡面.
我們不要將屬靈的力量放在錯位.
而是當我們願意去實踐上帝的心意的時候.
上帝的心意就令我們得力.
上帝願意去給我們力量去進行祂的心意.
我們在經歷這份力量的時候.
我們自然而然能夠感受到上帝.
我想我們今天在這個社會裡面.
我們最大的一個難關.
就是我們所敬拜的神是一個看不見的神.
我們有時候要作見證的時候.
我無法指著不知道哪裡.
然後跟別人說上帝在那裡.
但是今天我們每一個人坐在那裡成為基督徒.
是因為我們個人的生命裡面.
我們都經歷了上帝的大能.
我們每個人都在經歷上帝.
弟兄姊妹我們今天雖然見不到上帝.
但是上帝的聖靈讓我們能夠去經歷祂.
去感受祂.
從而讓我們生出一份信心能夠去敬畏祂.
弟兄姊妹我們今天願不願意去經歷這份屬靈的大能.
當這個世界跟你說很多屬靈的力量.
你可以操縱它的時候.
我們是否願意降服在上帝的心意裡面.
去享受那份真正屬靈的力量.
去經歷祂.
去跟上帝一起去做工.
從而讓我們在每一天裡面.

$^{681}$我們能夠說今天主都賜我力量.
因為上帝喜悅我去做祂的工作.
弟兄姊妹我們願意嗎.
我們願意去經歷這份屬靈的力量嗎.
我們一起去禱告.
親愛的父神.
我們去多謝你.
是因為你很願意在我們生命裡面去做工.
你在我們生命裡面去醫治.
去加力.
你讓我們成為一個新做的人.
主啊我們去求你幫助.
我們求你幫助.
讓我們更加明白你的心意.
因為我們知道.
如果我們今天想去經歷這份屬靈的力量.
我們必須在你的心意裡面去行.
我們必須在你的心意裡面去行走.
主啊我們今天很需要去經歷你.
特別在這個世俗的年代.
在這個好像屬靈力量不被彈劾的年代.
我們很需要真實去經歷你的大能.
從而活出一個不同的見證.
我們求上帝去幫助我們.
在生命裡面活出你的那份力量.
從而讓其他人看到.
讓其他人羨慕.
讓其他人願意一起歸入主耶穌基督裡面.
我們求上帝這樣去幫助我們.
以上禱告是奉主耶穌的名氣求.
阿們.
多謝 Michael 在我們當中分享上.
\newpage



\section{使徒行傳 13:1-52}
\label{sec:_FBi_OEMziA}
\textbf{How to Keep Going? (Acts 13) - Rev. Mark Chan 如何繼續走下去? (使徒行傳13\_1-52) - 陳崇基牧師}
\newline
\newline
連結: \href{https://youtube.com/watch?v=_FBi_OEMziA}{\texttt{https://youtube.com/watch?v=\_FBi\_OEMziA}} ~~~~ 語音日期: 2026-03-30
\newline
\newline
\hyperref[sec:a3FPlfJ6MYc]{\small{< < < PREV SERMON < < <}}
~
\hyperref[sec:index_chronic]{\small{[返順時目]}}
~
\hyperref[sec:index_scriptual]{\small{[返順卷目]}}
~
\hyperref[sec:ZqMpnK5BxrA]{\small{> > > NEXT SERMON > > >}}
\newline
\newline
使徒行傳 13:1-52
\newline
\begin{longtable}{cl}
\hline
\hline
章節 & 經文 (和合本修訂版)\\
\hline
13:1 & \begin{tabularx}{0.7\textwidth}{X} 在安提阿的教會中，有幾位先知和教師，就是巴拿巴和稱為尼結的西面、古利奈人路求，與希律分封王一起長大的馬念，和掃羅。 \end{tabularx} \\ \\ \relax
13:2 & \begin{tabularx}{0.7\textwidth}{X} 他們在事奉主和禁食的時候，聖靈說：「要為我分派巴拿巴和掃羅去做我召他們做的工作。」 \end{tabularx} \\ \\ \relax
13:3 & \begin{tabularx}{0.7\textwidth}{X} 於是他們禁食禱告後，給巴拿巴和掃羅按手，然後派遣他們走了。 \end{tabularx} \\ \\ \relax
13:4 & \begin{tabularx}{0.7\textwidth}{X} 他們既蒙聖靈差遣，就下到西流基，從那裡坐船往塞浦路斯去， \end{tabularx} \\ \\ \relax
13:5 & \begin{tabularx}{0.7\textwidth}{X} 到了撒拉米，就在猶太人各會堂裡宣講神的道，也有約翰作他們的幫手。 \end{tabularx} \\ \\ \relax
13:6 & \begin{tabularx}{0.7\textwidth}{X} 他們走遍全島，直到帕弗，在那裡遇見一個術士—猶太人的假先知，名叫巴耶穌。 \end{tabularx} \\ \\ \relax
13:7 & \begin{tabularx}{0.7\textwidth}{X} 這人常和士求‧保羅省長在一起。士求‧保羅是個通達人，他請巴拿巴和掃羅來，要聽神的道。 \end{tabularx} \\ \\ \relax
13:8 & \begin{tabularx}{0.7\textwidth}{X} 只是術士以呂馬(他的名字翻出來就是行法術的意思)敵對使徒，設法使省長遠離這信仰。 \end{tabularx} \\ \\ \relax
13:9 & \begin{tabularx}{0.7\textwidth}{X} 掃羅，又名保羅，被聖靈充滿，定睛看他， \end{tabularx} \\ \\ \relax
13:10 & \begin{tabularx}{0.7\textwidth}{X} 說：「你這充滿各樣詭詐奸惡，魔鬼的兒子，一切正義的仇敵，你還不停止扭曲主的正道嗎？ \end{tabularx} \\ \\ \relax
13:11 & \begin{tabularx}{0.7\textwidth}{X} 現在你看，主的手臨到你身上，你會瞎眼，暫時看不見日光。」立刻迷濛和黑暗籠罩著他，他到處摸索，求人拉著手領他。 \end{tabularx} \\ \\ \relax
13:12 & \begin{tabularx}{0.7\textwidth}{X} 省長看見所發生的事就信了，因對主的教導感到驚奇。 \end{tabularx} \\ \\ \relax
13:13 & \begin{tabularx}{0.7\textwidth}{X} 保羅和他的同伴從帕弗開船，來到旁非利亞的別加，約翰卻離開他們，回耶路撒冷去了。 \end{tabularx} \\ \\ \relax
13:14 & \begin{tabularx}{0.7\textwidth}{X} 他們從別加往前行，來到彼西底的安提阿。在安息日，他們進了會堂就坐下。 \end{tabularx} \\ \\ \relax
13:15 & \begin{tabularx}{0.7\textwidth}{X} 在讀完了律法和先知的書，會堂主管們叫人過去，對他們說：「二位弟兄，你們若有甚麼勸勉眾人的話，請說。」 \end{tabularx} \\ \\ \relax
13:16 & \begin{tabularx}{0.7\textwidth}{X} 保羅就站起來，做個手勢，說：「諸位以色列人和一切敬畏神的人，請聽。 \end{tabularx} \\ \\ \relax
13:17 & \begin{tabularx}{0.7\textwidth}{X} 這以色列民的神揀選了我們的祖宗，當百姓寄居埃及的時候抬舉他們，用大能的手領他們從那地出來。 \end{tabularx} \\ \\ \relax
13:18 & \begin{tabularx}{0.7\textwidth}{X} 他在曠野容忍他們，約有四十年。 \end{tabularx} \\ \\ \relax
13:19 & \begin{tabularx}{0.7\textwidth}{X} 他消滅了迦南地七族的人後，把那地分給他們為業， \end{tabularx} \\ \\ \relax
13:20 & \begin{tabularx}{0.7\textwidth}{X} 約有四百五十年。此後，他給他們設立士師，直到撒母耳先知的時候。 \end{tabularx} \\ \\ \relax
13:21 & \begin{tabularx}{0.7\textwidth}{X} 從那時起，他們要求立一個王，神就將便雅憫支派中基士的兒子掃羅給他們作王，共四十年。 \end{tabularx} \\ \\ \relax
13:22 & \begin{tabularx}{0.7\textwidth}{X} 他廢了掃羅之後，就興起大衛作他們的王，又為他作見證說：『我尋得耶西的兒子大衛，他是合我心意的人，他要遵行我一切的旨意。』 \end{tabularx} \\ \\ \relax
13:23 & \begin{tabularx}{0.7\textwidth}{X} 從這人的後裔中，神已經照著所應許的為以色列人興起一位救主，就是耶穌。 \end{tabularx} \\ \\ \relax
13:24 & \begin{tabularx}{0.7\textwidth}{X} 在他沒有出來以前，約翰已向以色列全民宣講悔改的洗禮。 \end{tabularx} \\ \\ \relax
13:25 & \begin{tabularx}{0.7\textwidth}{X} 約翰快走完他的人生路程時，說：『你們以為我是誰？我不是；但是有一位在我以後來的，我就是解他腳上的鞋帶也不配。』 \end{tabularx} \\ \\ \relax
13:26 & \begin{tabularx}{0.7\textwidth}{X} 「諸位弟兄—亞伯拉罕的子孫和你們中間敬畏神的人哪，這救世的道是傳給我們的。 \end{tabularx} \\ \\ \relax
13:27 & \begin{tabularx}{0.7\textwidth}{X} 耶路撒冷的居民和他們的官長，因為不認識這基督，也不明白每安息日所讀的先知的書，把他定了死罪，正應驗了先知的預言。 \end{tabularx} \\ \\ \relax
13:28 & \begin{tabularx}{0.7\textwidth}{X} 雖然他們查不出他有該死的罪狀，還是要求彼拉多把他殺了。 \end{tabularx} \\ \\ \relax
13:29 & \begin{tabularx}{0.7\textwidth}{X} 他們既實現了經上指著他所記的一切話，就從木頭上把他取下來，放在墳墓裡。 \end{tabularx} \\ \\ \relax
13:30 & \begin{tabularx}{0.7\textwidth}{X} 神卻使他從死人中復活。 \end{tabularx} \\ \\ \relax
13:31 & \begin{tabularx}{0.7\textwidth}{X} 有許多日子，他向那些從加利利同他上耶路撒冷的人顯現，這些人如今在民間成為他的見證人。 \end{tabularx} \\ \\ \relax
13:32 & \begin{tabularx}{0.7\textwidth}{X} 我們報好信息給你們，就是那應許祖宗的話， \end{tabularx} \\ \\ \relax
13:33 & \begin{tabularx}{0.7\textwidth}{X} 神已經向我們這些作他們兒女的應驗，使耶穌復活了。正如《詩篇》第二篇上記著：『你是我的兒子，我今日生了你。』 \end{tabularx} \\ \\ \relax
13:34 & \begin{tabularx}{0.7\textwidth}{X} 論到神使他從死人中復活，不再歸於朽壞，他曾這樣說：『我必將所應許大衛那聖潔、可靠的恩典賜給你們。』 \end{tabularx} \\ \\ \relax
13:35 & \begin{tabularx}{0.7\textwidth}{X} 所以他也在另一篇說：『你必不讓你的聖者見朽壞。』 \end{tabularx} \\ \\ \relax
13:36 & \begin{tabularx}{0.7\textwidth}{X} 大衛在世的時候，遵行了神的旨意就長眠了，歸到他祖宗那裡，已見朽壞； \end{tabularx} \\ \\ \relax
13:37 & \begin{tabularx}{0.7\textwidth}{X} 惟獨神使他復活的那一位，他並未見朽壞。 \end{tabularx} \\ \\ \relax
13:38 & \begin{tabularx}{0.7\textwidth}{X} 所以弟兄們，你們當知道：赦罪的道是由這人傳給你們的， \end{tabularx} \\ \\ \relax
13:39 & \begin{tabularx}{0.7\textwidth}{X} 你們靠摩西的律法在不得稱義的一切事上，每一個信靠這位耶穌的都得稱義了。 \end{tabularx} \\ \\ \relax
13:40 & \begin{tabularx}{0.7\textwidth}{X} 所以，你們要小心，免得先知書上所說的臨到你們： \end{tabularx} \\ \\ \relax
13:41 & \begin{tabularx}{0.7\textwidth}{X} 『要觀看，你們這些藐視的人，要驚訝，要滅亡，因為在你們的日子，我行一件事，雖有人告訴你們，你們總是不信。』」 \end{tabularx} \\ \\ \relax
13:42 & \begin{tabularx}{0.7\textwidth}{X} 他們走出會堂的時候，眾人請他們在下一個安息日再講這些話給他們聽。 \end{tabularx} \\ \\ \relax
13:43 & \begin{tabularx}{0.7\textwidth}{X} 散會以後，有許多猶太人和敬虔的皈依猶太教的人跟從了保羅和巴拿巴。二人對他們講話，勸他們務要恆久倚靠神的恩典。 \end{tabularx} \\ \\ \relax
13:44 & \begin{tabularx}{0.7\textwidth}{X} 到下一個安息日，全城的人幾乎都聚集起來，要聽主的道。 \end{tabularx} \\ \\ \relax
13:45 & \begin{tabularx}{0.7\textwidth}{X} 但猶太人看見這麼多的人，就滿心嫉妒，辯駁保羅所說的話，並且毀謗他。 \end{tabularx} \\ \\ \relax
13:46 & \begin{tabularx}{0.7\textwidth}{X} 於是保羅和巴拿巴放膽說：「神的道本應先傳給你們；只因你們棄絕這道，斷定自己不配得永生，我們就轉向外邦人。 \end{tabularx} \\ \\ \relax
13:47 & \begin{tabularx}{0.7\textwidth}{X} 因為主曾這樣吩咐我們：『我已經立你作萬邦之光，使你施行我的救恩，直到地極。』」 \end{tabularx} \\ \\ \relax
13:48 & \begin{tabularx}{0.7\textwidth}{X} 外邦人聽見這話很歡喜，讚美主的道，凡被指定得永生的人都信了。 \end{tabularx} \\ \\ \relax
13:49 & \begin{tabularx}{0.7\textwidth}{X} 於是主的道傳遍了那一帶地方。 \end{tabularx} \\ \\ \relax
13:50 & \begin{tabularx}{0.7\textwidth}{X} 但猶太人挑唆虔敬尊貴的婦女和城內有名望的人，迫害保羅和巴拿巴，把他們趕出境外。 \end{tabularx} \\ \\ \relax
13:51 & \begin{tabularx}{0.7\textwidth}{X} 二人對著眾人跺掉腳上的塵土，然後往以哥念去了。 \end{tabularx} \\ \\ \relax
13:52 & \begin{tabularx}{0.7\textwidth}{X} 門徒滿心喜樂，又被聖靈充滿。 \end{tabularx} \\ \\
[1ex]
\hline
\hline
\end{longtable}
$^{1}$- I want to thank all of you for welcoming me back..
I came back last night, right before midnight..
So I travel for, in the past two weeks to,.
well, I have a layover in Paris,.
and then my final destination was in Turkey..
So I joined or led a tour in western and central Turkey,.
and most of the people would have questions.
about the safety in Turkey..
And I am living proof that there was safe to travel, okay?.
There was safe to travel..
A lot of coverage from the western media.
would have us feel really scary about that place..
But you need to understand, you need to read the map, okay?.
You need to understand the map..
Turkey is very at the outskirt of the Middle East..
So what happened in Israel, in Palestine,.
in Tel Aviv, in Jerusalem, they're not, you know,.
greatly affected the area in western Turkey..
So I was fine..
So as you can see, right, two weeks..
(speaking in foreign language).
Okay, so, well, let's get into the text today..
And I have a little text to share with you..
It is from Acts 13..
So as I said, over the past two weeks,.
I'm very blessed to have joined.
some of the brothers and sisters from UK,.
and nine of them joined me..
And then 10 of us traveled together in the lands of Bible,.
especially in Turkey..
And I visited a lot of places,.
Istanbul or the sites of the seven churches.
in the Book of Revelation, and walked through those places,.
and make Bible come alive and remind me.
the Bible is not just a book of ideas..
It is rooted in real history, in real time,.
in real space, geography, and real people..
(speaking in foreign language).
And then two weeks ago,.
I stood in the region of Antioch of Pisidia,.

$^{41}$and it was a place mentioned in the Bible..
You know, when you went there,.
it was quiet and ancient ruins and all the stones,.
and what struck me the most, it's not just the city,.
but what happened in Bible is the journey.
that was made to this place..
(speaking in foreign language).
And then, as you understand in the Book of Acts,.
there's recorded three different missionary journeys of Paul,.
and as you understand, the arrows from Antioch on the right,.
that's in other Antioch,.
and then you follow that white arrow,.
and you come across the island of Cyprus,.
and then they went north..
Paul and his team, in the first missionary journey,.
went north from the island,.
and they went to what they call Asia Minor back then,.
and what we call Turkey today,.
and they landed on a poor city called Perga,.
and then they went up northward.
to what I show you in the picture, Antioch of Pisidia..
So, Paul went there from Perga,.
from the coastal poor city,.
and he has to climb like 1,000 meter or 3,000 feet,.
more than 3,000 feet, in order to reach that place,.
and the road was rugged..
The terrain was very difficult, and it was steep,.
and it was very dangerous,.
and a very exhausting traveling experience..
It's not comfortable at all..
It was a very, very difficult, very hard road..
(Yi speaking Chinese language).
And then, to understand this moment,.
from Perga to Antioch of Pisidia,.
we need to kind of zoom out, zoom out for a second..
What had happened here?.
(Yi speaking Chinese language).
Now, this journey, the first missionary journey of Paul,.
began in Antioch of Syria, as I told you..
There was the church in Antioch of Syria on the right..

$^{81}$They were praying, and they were fasting,.
and Acts 13 reminds us,.
the Holy Spirit, while they were worshiping,.
while they were praying and fasting,.
the Holy Spirit spoke to them,.
set apart for me Barnabas and Saul, who is Paul,.
for the work which I have called them..
So this mission did not begin.
with human ideas or initiative..
It began with the calling of God,.
with the Spirit calling the church to send out people..
(Yi speaking Chinese language).
So what would it mean?.
Because it's very important..
This is the first missionary journey of Paul..
From there, from Antioch of Syria on the right,.
they took a boat or a ferry,.
and then went to the island of Cyprus..
You see the pointing, kind of sharp-pointing island..
It's called Cyprus..
And there, they purged the gospel.
in different Jewish synagogues,.
and they faced different oppositions and persecution..
And they confronted the false prophet,.
and they saw the Roman governor, and he came to faith..
So already, in this journey,.
you have a mix of success and resistance..
(Yi speaking Chinese language).
So we are in Cyprus, and from Cyprus,.
they took a boat and sailed northward to Perga,.
which is a port city..
It's one of the main entry points.
into the interior of today's Turkey or Asian minor..
In a sense, this is where the mission moves.
from coast to coast and into the heart of the region,.
today's Turkey..
(Yi speaking Chinese language).
So can you find Perga?.
From Perga, the port city,.
they went inland and upland to Antioch of Pisidia..

$^{121}$Now, the city, the ancient city of Antioch or Pisidia.
was established in the Greek period,.
and the city was strategically located.
in the crossroads of major Roman road,.
and there was a major highway,.
just like there's a major highway in the city of Chicago,.
290 or 994 or whatever..
So back then, there was a major Roman highway,.
Roman road, called Via Sebast..
Via Sebast, it goes from east to west..
It's a major highway..
It's a key highway that connect different area,.
different regions of Asian minor..
So in other words, Paul was not traveling randomly..
He was moving from place to place to influential places.
where the gospel should be rooted.
and could be preached more widely..
(Yi speaking Chinese language).
However, the road to get from Perga to Antioch of Pisidia.
was very, very tough because Antioch of Pisidia.
is located on a high plateau..
So the journey moved inland,.
and the mission becomes harder,.
and the mission becomes more difficult..
(Yi speaking Chinese language).
So the journey moved inland, and the mission becomes harder,.
and the road becomes steeper, and the cost becomes greater..
(Yi speaking Chinese language).
And this is exactly at this point of the story.
that something happens, okay?.
Something happens..
(Yi speaking Chinese language).
In Acts 13, in this journey, when they were in Perga,.
there were three different responses to follow Jesus,.
three different responses we found from the text..
So one person of the team left for no reason,.
and some other people continue,.
and the others are just the beginning..
(Yi speaking Chinese language).
So I wanna show you there were three different kinds.

$^{161}$of responses to the call of Jesus.
from this passage, Acts 13..
(Yi speaking Chinese language).
Now, point number one is that some people turn back.
when the road gets harder..
(Yi speaking Chinese language).
And Acts 13 says what?.
It says John Mark, there was a guy, John Mark,.
in Paul's team, left at this critical moment,.
and the journey became harder,.
and he, John Mark, chose to return to Jerusalem..
And this reminds us many people are willing.
to follow Jesus only when they find the road was manageable,.
when they find the road was smooth,.
when the cost increases, some people will step back..
(Yi speaking Chinese language).
And we may not physically leave the team,.
but we can turn our back in our hearts,.
in our mentality, because we wanna avoid difficulties..
We wanna avoid hardship or difficult obedience..
We choose comfort over calling,.
and slowly we start moving forward..
(Yi speaking Chinese language).
So the warning is clear..
It is possible to start well in a Christian journey,.
and then still turn back.
when the road gets harder and harder..
(Yi speaking Chinese language).
So John Mark left, and that's the point,.
because the road gets harder and harder..
And then point two here,.
it shows us some people keep going in faith.
when the road gets hard..
(Yi speaking Chinese language).
In the story, it's Paul and Barnabas..
They continue and took the hard and difficult road,.
and they did not stop when it became difficult for them..
And Paul and Barnabas understood that, you know,.
hardship is actually part of our following Jesus..
(Yi speaking Chinese language).

$^{201}$So guys, you need to understand,.
you know, Christian life is not about how you start..
It's not really about how you started your faith journey..
Okay, you may have started your faith journey.
in happiness, you know, excitement and all that stuff..
It is not about that..
Faith is not about starting..
It's about continuing..
It's about continuing in God and trusting God,.
even when the road gets harder,.
because obedience is costly,.
and the cost can be high..
So we may be faithful when the results are low,.
when we don't see success,.
but you stick to it because that's what faith is all about..
(Yi speaking Chinese language).
So sisters and brothers, please, you know,.
pay attention to this..
Do not measure, do not measure your faith journey.
by how you started..
Do not, because it really doesn't matter..
Measure your faith by whether you keep going..
(Yi speaking Chinese language).
The point three here tell us.
that some people continue in grace though..
(Yi speaking Chinese language).
Now, as you can see the context of the story.
in the first journey of missionary,.
many people began to follow Paul and Barnabas,.
but Paul told them not just to begin,.
but to continue in the grace of God..
(Yi speaking Chinese language).
So you say, "Okay, I have to keep going.".
It's when the road gets harder and it's very difficult..
No, but that's not the point..
The point of the Bible, the point of gospel.
is that the Christian life is not sustained by our strength,.
but by the very grace of God, the sheer grace of God,.
because if we depend on the grace of God,.
the grace of God will carry us from beginning to the end..

$^{241}$(Yi speaking Chinese language).
So when you first answer the call of Jesus to follow him,.
the call is not just to begin,.
but to continue, continue and keep going in grace..
(Yi speaking Chinese language).
So here, what is important is that the story of,.
remember who left the team?.
What is his name?.
John Mark, right?.
John Mark..
This is where the story becomes very important.
because if you look at the entire story of John Mark,.
that was the beginning..
He failed, he turned back..
And earlier, earlier he turned back and he gave up,.
but that was not the end of the story for him..
(Yi speaking Chinese language).
Because somehow at the end of his life,.
we know that in the middle of his following Jesus,.
serving God, he changed somehow at the middle..
And later on he did return..
And then eventually at the end of Paul's life,.
when he got so much older,.
he wrote in one of his letter called 2 Timothy..
That is probably last known letter of Paul..
And he said, "Bring Mark.".
And he was saying to Timothy, "Only Luke is with me.".
And one more thing, get Mark, who is John Mark,.
bring him with you..
You know, when Timothy, calling Timothy to come to visit.
and take care of Paul because he was in prison,.
because of the gospel..
And Paul said, "Bring him, bring Mark,.
because John Mark has been helpful to me in my ministry.".
(Yi speaking Chinese language).
And that means only one thing,.
John Mark learn and he change and he continue in grace..
And that means he did not stay in his weakness or failure..
God somehow restore him or he came back.
and God use him again..

$^{281}$(Yi speaking Chinese language).
Now, in closing, I wanna ask you,.
which one are you, which one you wanna be?.
You know, will you turn back in your weakness,.
in your public failure?.
And will you keep going?.
And will you continue in grace?.
(Yi speaking Chinese language).
So this little story in the first missionary journey of Paul.
reminds you that wherever you are today,.
guys, wherever you are today,.
you may feel weak or you may feel like you're failing..
Do not stay where you are,.
because the grace of God is sufficient enough.
to carry you through..
So keep walking..
(Yi speaking Chinese language).
Well, shall we pray together?.
(Yi speaking Chinese language).
Heavenly Father, thank you for your calling.
of each and every one of us to follow you,.
even when the road gets hard..
(Yi speaking Chinese language).
Forgive us for the time we have turned back,.
turned our back against you,.
or chosen comfort over obedience..
(Yi speaking Chinese language).
And yes, give us strength today to keep going in our faith,.
even when we are tired or discouraged..
(Yi speaking Chinese language).
And teach us, Jesus, to rely not on ourselves,.
but to continue in your grace each day..
(Yi speaking Chinese language).
Even when we fail,.
remind us that you can always restore and call us again..
(Yi speaking Chinese language).
And help each and every one of us to keep walking with you,.
step by step, day by day..
(Yi speaking Chinese language).
In Jesus' name we pray..

$^{321}$(Yi speaking Chinese language).
Amen. - Amen..
- Amen..
Thank you very much for worshiping with us,.
and we'll see you next Sunday..
And let's take some time at the end of the service.
to say hi to each other, okay?.
Now, you welcome me back, all right?.
- Yeah. - You also need.
to welcome each other, welcome back, all right?.
(Yi speaking Chinese language).
All right, then there's some refreshment prepared.
for all of us in the cafe area,.
so we can quietly proceed to the cafe,.
and then we have time of fellowship, okay?.
(Yi speaking Chinese language).
Okay, God bless..
\newpage



\section{約翰福音 2:1-12}
\label{sec:ZqMpnK5BxrA}
\textbf{行奇事的神 (約翰福音2\_1-12) - 梁紫健傳道}
\newline
\newline
連結: \href{https://youtube.com/watch?v=ZqMpnK5BxrA}{\texttt{https://youtube.com/watch?v=ZqMpnK5BxrA}} ~~~~ 語音日期: 2026-03-30
\newline
\newline
\hyperref[sec:_FBi_OEMziA]{\small{< < < PREV SERMON < < <}}
~
\hyperref[sec:index_chronic]{\small{[返順時目]}}
~
\hyperref[sec:index_scriptual]{\small{[返順卷目]}}
~
\hyperref[sec:a8Vgzt3GCx8]{\small{> > > NEXT SERMON > > >}}
\newline
\newline
約翰福音 2:1-12
\newline
\begin{longtable}{cl}
\hline
\hline
章節 & 經文 (和合本修訂版)\\
\hline
2:1 & \begin{tabularx}{0.7\textwidth}{X} 第三日，在加利利的迦拿有一個婚宴，耶穌的母親在那裡。 \end{tabularx} \\ \\ \relax
2:2 & \begin{tabularx}{0.7\textwidth}{X} 耶穌和他的門徒也被請去赴宴。 \end{tabularx} \\ \\ \relax
2:3 & \begin{tabularx}{0.7\textwidth}{X} 酒用完了，耶穌的母親對他說：「他們沒有酒了。」 \end{tabularx} \\ \\ \relax
2:4 & \begin{tabularx}{0.7\textwidth}{X} 耶穌說：「母親，我與你何干呢？我的時候還沒有到。」 \end{tabularx} \\ \\ \relax
2:5 & \begin{tabularx}{0.7\textwidth}{X} 他母親對用人說：「他告訴你們甚麼，你們就做吧。」 \end{tabularx} \\ \\ \relax
2:6 & \begin{tabularx}{0.7\textwidth}{X} 照猶太人潔淨禮的規矩，有六口石缸擺在那裡，每口可以盛兩三桶水。 \end{tabularx} \\ \\ \relax
2:7 & \begin{tabularx}{0.7\textwidth}{X} 耶穌對用人說：「把缸倒滿水。」他們就倒滿了，直到缸口。 \end{tabularx} \\ \\ \relax
2:8 & \begin{tabularx}{0.7\textwidth}{X} 耶穌又說：「現在舀出來，送給宴會總管。」他們就送了去。 \end{tabularx} \\ \\ \relax
2:9 & \begin{tabularx}{0.7\textwidth}{X} 宴會總管嘗了那水變的酒，並不知道是哪裡來的，只有舀水的用人知道。於是宴會總管叫新郎來， \end{tabularx} \\ \\ \relax
2:10 & \begin{tabularx}{0.7\textwidth}{X} 對他說：「人家都是先擺上好酒，等客人喝夠了才擺上次的，你倒把好酒留到現在！」 \end{tabularx} \\ \\ \relax
2:11 & \begin{tabularx}{0.7\textwidth}{X} 這是耶穌所行的第一個神蹟，是在加利利的迦拿行的，顯出了他的榮耀來，他的門徒就信他了。 \end{tabularx} \\ \\ \relax
2:12 & \begin{tabularx}{0.7\textwidth}{X} 這事以後，耶穌與他的母親、兄弟和門徒都下迦百農去，在那裡住了不多幾天。 \end{tabularx} \\ \\
[1ex]
\hline
\hline
\end{longtable}
$^{1}$頂智慕早晨.
終於暖一點了.
經歷了一個比較寒冷的冬天.
不知道大家有沒有見過這東西呢.
是啊 醒目啊.
煲湯的朋友應該會知道這是海底的東西.
我也是第一次見到它還沒切的樣子.
通常我們見到的是切了的.
切了的是右手邊的位置.
大家見到它的時候.
海底的東西有什麼功用呢.
紫卡化痰.
紫卡化痰 還有呢.
潤肺 沒錯.
其實它有很多功用.
實際上有清熱潤肺紫卡.
滋陰補腎 山津止渴 美顏潤膚.
含維他命B C E 維生素.
有助滋潤皮膚 提升免疫力 強身健體.
適合涯野和體質偏熱的人士.
特別是在北美經常開暖氣 很燥熱的.
在秋冬的時候很適合煲來喝的.
所以現在暖肺可能你不需要.
但是如果冷和燥熱的時候你就會用它.
海底的東西是一種幫助我們.
轉化我們身體裡面燥熱的一個上好的藥材.
不知道你們知不知道.
聖經裡面也有海底的東西.
這位海底爺平時不鋪頭.
一鋪頭就做了一件奇事.
轉化了人對於榮耀和羞恥的看法.
在約翰福音裡面.
人很急躁地保護自己的榮耀和面子.
而耶穌就是用我們生活中.
很平常最不起眼的事去觸摸我們的心.
甚至改變了我們人的一生.
我們一起來看看吧.
我小時候在教會聽這段經文的理解是這樣的.
「水變酒」大家都聽過吧.
耶穌是一個可以改變物質的神.

$^{41}$可以將水H2O變成酒C2H5O.
多了很多不同的原子.
有碳也有氫.
耶穌真的很厲害.
成為一個可以改變物質的上帝.
不過我今天不是想用這個角度去看.
我想帶大家從另一個角度去看.
想想耶穌當時在約翰福音裡面的處境是什麼.
作者怎樣去描寫耶穌這個海底爺鋪頭.
他做第一個神蹟.
大家就會發現這個上帝不只是一個改變物質的神.
更加是改變人心的神.
我們一會細緻一點去看這個經文.
第三天在加利利的加拿大.
有一個婚宴.
耶穌的母親在那裡.
耶穌和祂的門徒都被邀請到宴會裡面.
酒用完了.
然後耶穌的母親對他們說.
他們沒有走.
很奇怪.
瑪利亞就跟耶穌說.
他們沒有走.
這句話好像很平凡.
但是很奇怪的地方就是.
瑪利亞明明知道耶穌不是賣酒.
而她老公.
就是耶穌的父親.
是一個木匠.
最多是會做木桶.
不會弄酒的.
為什麼瑪利亞會跟他說.
他們沒有走呢.
瑪利亞在期待什麼.
我們接著看下去就有提示了.
就是耶穌說.
母親.
我與你何干.
我的時候還沒到.
耶穌的回應有一點奇怪.

$^{81}$他反應都有點大腿.
就是他在問.
我和你瑪利亞有什麼關係.
而不是問他們沒有走.
關我什麼事.
我和你是什麼關係.
這個問題是有點奇怪的.
如果你細心想.
跟瑪利亞的回應.
他們沒有走.
耶穌就好像很親切.
很親切的說.
我和你是什麼關係.
這樣.
究竟這裡發生什麼事呢.
這裡看到時候這個字.
我的時候還沒到.
時候這個字其實是非常關鍵的.
以致到耶穌是這樣回應母親.
我們一起看下去.
母親對人說.
他告訴你什麼.
你們就做了.
照猶太人潔淨的規矩.
有六口石缸擺在那裡.
每口可以乘兩三桶水.
這段是經文裡面.
最平鋪直敘的地方.
但其實都是那兩點所在.
這裡強調用了六口.
是按照猶太人禮儀所擺放的石缸.
而石缸的大小是什麼.
專登寫下去.
可以裝兩三桶水.
證明這個是一個挺大的石缸.
而這兩點都是為什麼.
耶穌會這樣回應瑪利亞的關鍵所在.
我繼續看下去.
第七.
耶穌就對傭人說.

$^{121}$把缸倒滿水.
他們就倒滿了.
直到缸口.
耶穌又說.
現在要出來送給燕湖總管.
他們就送了去.
燕湖總管就試了那些水所變的酒.
都不知道是從哪裡來的.
就跟那些要水的人說.
叫新郎來.
然後就跟他說.
人家都是先放好酒.
你等人家喝夠了才放次的.
為什麼現在才留到這些好酒.
這個就是耶穌在《約翰福音》裡面.
所描述的第一個神蹟.
在加利利的加納恆.
顯出了他的榮耀.
門徒就信他了.
這個事之後.
門徒母親和兄弟就一起到加伯朗住了幾天.
看完這段經文.
我們要收集一下剛剛所問的問題.
第一.
瑪利亞的提問.
他說他們沒有走了.
其實究竟他對耶穌說話.
期待著什麼呢.
第二.
耶穌的反應其實有點奇怪.
反問我和你有什麼關係.
提到時候還沒有到.
是不是在暗示什麼呢.
第三.
潔淨條例是這段經文的關鍵.
回答了我們上面兩個問題.
讓大家留意一下.
我們準備好就一起來.
第一個問題就是.
瑪利亞對耶穌說他們沒有走了.

$^{161}$究竟期待著什麼呢.
在猶太文化裡面.
酒宴沒有酒喝.
或者在慶祝的婚宴裡面沒有酒喝.
是一個很大的恥辱.
會讓人覺得新郎的家準備失當失禮.
甚至失去了上帝祝福的一個表現.
大家有沒有聽過葡萄園.
有聽過葡萄園.
在舊約裡面經常提到葡萄園.
是上帝所給人的祝福.
以色列自古就是種葡萄和釀酒的重鎮.
所以葡萄酒對於猶太人來說.
其實是一個好意頭.
代表著上帝祝福這段婚姻.
你不夠酒.
簡直是一個公關的災難.
代表著這個家庭幾乎告訴別人.
他們可能不是很受到上帝的祝福.
這個是令到家族沒有面子.
受到屈辱的一個情況.
猶太人很注重面子.
那種榮辱的文化影響他們很深.
情況就好像.
大家都有聽過華人在婚宴上說的十個字.
主人家會說什麼.
吃飯上台會說什麼.
酒微菜薄 招呼不到.
大家多喝兩杯.
酒微菜薄.
大家經常都聽到這句.
這句是主人家用來謙稱酒菜的品質是一半.
分量很少 招呼不到.
只是謙稱 酒微菜薄.
你試一下 酒微菜薄.
喝酒可能整桌都只是一人一杯.
喝完就沒有了.
前菜上完之後就完.
主人家會覺得怎麼樣.
很失禮 招呼不夠.

$^{201}$所以大家都理解到.
猶太人的榮辱文化.
沒有酒喝很失禮.
而耶穌的母親瑪利亞.
很大可能是親戚或是很熟的朋友.
經文看到他們可以看到那些毒人準備走水.
他們應該是可以進入別人家裡內院的人.
可以看到整個婚宴的運作.
是親戚或是熟的朋友.
就像現代的婚宴一樣.
你是熟的人才可以去後台看到他們的運作.
然後看看有沒有什麼可以幫忙.
所以瑪利亞比其他人更加在意主人家的面子和榮耀.
很敏銳地觀察到背後的難處.
所以當瑪利亞跟耶穌說他們沒有酒.
他在期待什麼.
他在反映一件事就是.
瑪利亞相信耶穌可以做一些事.
瑪利亞相信耶穌有能力幫主人家挽回面子.
或者避免失禮.
他期待著耶穌介入.
挽回他們的榮耀.
這裡大家明白了嗎.
大概是這樣.
但是耶穌就回應了.
母親.
我與你何干.
我的時候未到.
意思是我跟你是什麼關係.
有些日本人可能說我跟這件事有什麼關係.
其實你看原文.
其實很簡單.
我跟你是什麼關係.
很直接.
就是這樣.
表達什麼意思.
你這樣也不明白.
你明知道我是彌賽亞.
而行神蹟會令彌賽亞的身分提前被人知道.
就是說耶穌本身很低調.

$^{241}$時間時候到了才彰顯上帝的榮耀.
才做光芒四射的彌賽亞.
而瑪利亞的期望是.
不如你彰顯你的榮耀的時候.
不如幫一下你的家人.
在約翰福音裡面剛剛說到.
時候這個字是很關鍵的.
每一次說到時候的時候.
通常你都會看到第二個字.
就是榮耀.
時候和榮耀其實很關鍵.
就是因為耶穌知道當他一行神蹟的時候.
他就會將上帝的榮耀彰顯出來.
收也收不起來.
所以耶穌的反應不是問他們沒走關我什麼事.
而是在問母親.
你明知道我在潛水.
我在海底.
在做海底的事.
你現在叫我行神蹟.
而你和我的關係.
你還不知道我這樣做.
我要被人知道.
我就是那個彌賽亞.
所以大家明白我.
不過耶穌真是一個很乖的兒子.
之後他照做.
但是他之後做的事就打破了瑪利亞的想像.
耶穌不只是令主人家得到榮耀.
是令到他自己的榮耀都顯出來.
我們繼續看下去.
第六節.
照猶太人的潔淨的規矩.
有六口石缸擺在那裡.
有些中英文譯本就會譯做.
是給猶太人行潔淨禮的石缸.
或者用於猶太人潔淨禮的石缸.
其實猶太人很注重潔淨禮儀.
因為他們知道上帝在川一級記和利未記裡面.
頒佈了一些潔淨的條例.

$^{281}$是想和世人展示猶太民族是屬於耶和華.
屬於上帝.
這裡出現了一個理念.
就是猶太人透過上帝所頒佈的十誡和舊約的潔淨條例.
去理解和衍生出來一些潔淨的禮儀.
目的是什麼呢.
目的是去提醒.
用這個外在的舉動去提醒人類.
我是和上帝有關係的.
當人類做這個外在的舉動.
最理想的效果就是.
人知道而且深化自己的行為.
去提醒自己.
自己屬於上帝.
以致我裡面的內心被轉化.
靠近上帝一點.
而不是單單做外在的行為.
而經文所描述的情況就是出現了一個斷層.
當猶太人很努力去做潔淨的禮儀.
還要用一般人還要大還要多的石缸.
這裡的石缸其實很重和貴.
通常出現在一些有經濟能力.
也很注重潔淨條例的家庭裡.
不是每一個家庭都會有的.
所以這一家人明顯是有頭有臉的有錢人.
為什麼經文要特意提到是兩三桶水.
是因為那個石缸夠大.
容納到的水夠多.
而且貴.
所以他們是非常注重榮辱.
他們覺得潔淨條例做得好的人就等於.
和上帝有好的關係.
得到上帝的祝福.
這個家族就好像在一個猶太文化裡.
是一個根正苗紅 明門正宗.
在乎的是什麼呢.
在乎的是人的榮辱.
多於自己和上帝的關係.
多於自己的心態.
結果是怎樣呢.

$^{321}$結果是這樣的.
這個圖是反轉了的.
只是做外在的禮儀.
只是集中在潔淨禮儀的本身.
很忙碌地很注重外在的表達.
忽略了本身潔淨禮儀想要提醒人.
我屬於上帝.
我們的心要回歸上帝.
所以今天如果你去有機會.
希望大家還有機會.
去耶路撒冷的哭場.
大家還會看到有些猶太人在守潔淨禮.
在牆上有個洗手的地方.
洗手盆裡面有一隻銀色的杯.
他們就會用裡面銀色的杯去舀水.
象徵著以前的一些缸.
他們去洗淨自己的污穢.
去到上帝面前才潔淨上帝雙手.
表明自己屬於上帝.
可想而知猶太人多麼注重潔淨禮儀.
不過有沒有真正提醒自己屬於上帝.
這是一人而二.
這個我也要問自己.
所以耶穌偏偏用了這個走潔淨禮的石缸.
叫人裝滿水.
再用裡面的水變成酒.
他絕對可以用其他容器去裝水.
也不需要用這個石缸去裝水.
這個舉動就打破了猶太人和門徒.
對於潔淨禮的想法.
因為用來走潔淨禮的石缸.
其實是充滿著不潔的.
你想想原本的水是被人洗手,洗臉.
甚至是洗腳.
你可以想像到人是從四方八面來.
走了很遠的路來到這個婚宴.
當時的科技也沒有那麼先進.
而很多人都是沒有鞋穿.
或者穿了一些很粗糙的鞋.
你想像一下那缸水有多骯髒.

$^{361}$有多不潔.
裡面有多少泥,污漬.
而且客人已經到齊了.
酒也差不多喝完了.
所以那六口石缸很可能已經裝了幾輪洗腳水.
已經裝了幾輪.
已經有很多泥,很多東西在裡面.
而現在怎麼說呢.
重新再裝滿它.
然後你就會看到.
原來最骯髒,最不潔的水變成了最好的酒.
比原本婚宴的酒更加好.
經文這麼說.
情況就是這樣的.
就好像參觀這些污水處理廠一樣.
我曾經在新加坡參觀他們的處理廠.
其實他們做得很好.
廠長就帶著我們.
在一排污水處理的缸裡面.
就是剛剛從別人家裡的廁所抽出來的污水.
然後他說水就是這樣.
還有一股味道.
然後經過一系列的處理.
然後你就會看到原本啡黑色的水.
變成了很清澈的水.
然後廠長就拿起杯子.
就在我們前面舀起.
就喝.
然後就問我們.
你想不想試一下.
我們就有一個很大的生理障礙.
我們就覺得.
可以喝嗎.
沒人敢試.
就是這樣.
所以我們心裡面也有一種感覺.
耶穌不單止改變了物質.
他改變了人對於潔淨的黑板印象和錯誤理解.
不是那些水乾不乾淨.
不是那些洗完之後有多少道德上的不潔可以洗掉.

$^{401}$而是重新將潔淨條例的本質連繫到人的心裡面.
我們做這件事其實是怎樣.
我們回到上帝那裡.
所以面對著做歧視的上帝.
我們可以怎樣在屬靈的生命裡有回應呢.
這些石缸本身是給人做潔淨禮儀用的.
象徵著外在的一些行為,宗教儀式.
內心的一些污穢可以被潔淨.
但是耶穌是將裡面的水變成酒.
象徵著他用他自己的生命.
他用他自己的血.
去做了人真正潔淨的根源.
凡是一些我們覺得最修家,最難開口.
最不潔淨的一些生命的片段.
無論是人生失敗,貪玩無知而導致的後果.
上癮,破碎的生命等等.
那些都是我們人生裡面的片段.
我們眼中都會覺得.
我都不想提.
我想都不想想.
就在我們人生很夜闌人靜的時候.
只剩下我和上帝的時候.
那些回憶就好像電影一樣.
去到拍出來.
你想走都走不掉.
或者我們有沒有想過.
或者我們都不敢期望.
自己那些攀回手的一些經歷.
上帝是會做一個神蹟.
會將我們這段最想收買,最失敗,最不潔的人生片段.
轉化為我們人生的一些養分.
甚至是榮耀上帝的工具.
在基督恩典裡面我們可以學習.
將我們最不堪的人生片段交給他.
不是靠一些宗教儀式或者表面的善行去遮蓋掩飾.
而是期待著耶穌基督在我們生命裡面做一次真正的轉化.
可能和大家都有講過.
其實我小時候有腦癌症.
上帝就醫治了我.
但是上帝醫治我的過程裡面.

$^{441}$他留下了一些東西給我.
就是一些我對於未來或我在日常生活中的恐懼.
就是我很怕突然間會暈倒.
那種恐懼有時候會吞噬我.
令我覺得原來我說話可能突然間我會說不到話.
可能我突然間會完全忘記.
這種恐懼對我來說很大.
我仍然要面對著這個恐懼.
那些片段在當我一個人的時候經常出現.
我其實完全想像不到我可以站在這裡和大家說話.
因為即使我站在前面那一刻.
我仍然有少許恐懼.
好像我突然間會暈倒.
這就是我一直想像不到.
這個行歧視的上帝竟然在我生命裡面.
告訴我一些小時候我最不想面對的事情.
他突然間扭轉說不是的.
你可以用這件事成為一個榮耀神的功能.
就好像《楊憑虹二章》做了一個歧視.
竟然將一些最不潔淨的水.
一些可以用原汁氣的東西.
代表著一些過去我們覺得過了去就算的生命的片段.
轉化成為一個榮耀神的工具.
但問題是我們願不願意.
我們可以只是做一些表面的事情.
我們願不願意用一些最真實的生命去回應上帝.
當我們說到工作的時候.
我們不知道什麼時候開始.
我們也不是太渴望上帝在我們的職場介入.
每天上班只做千篇一律的事情.
有時候可能要面對著一些巨嬰同事.
不懂得做事就要煮你一鍋.
心裡很煩躁.
不知道什麼時候你要處理好手上的工作.
不知道什麼時候開始我們就會有些想法.
就是上帝應該不會在我們生命裡面.
在工作裡面出現.
我們唯一的盼望就是今天的事情快點完成.
做好它.
彷彿我們不再期待上帝在我們生命裡面.

$^{481}$我們的工作崗位出現.
不知道什麼時候我們不再期待上帝和我們的工作沒有什麼關係.
我只是賺錢.
我只需要這份工作去養家糊口.
除了這個之外.
其實這份工作沒有什麼其他意義.
耶穌最驚訝的地方就是.
祂竟然將猶太人最嫌棄的洗腳水.
看起來最沒有意義的水.
用完即棄.
那些缸子甚至是有錢人會把它打爛.
因為覺得已經沾了不潔.
已經沒有用了.
祂竟然將裡面的水成為了主人家的榮耀.
將自己的榮耀都彰顯出來.
同樣地.
祂都可以將我們生命中本身覺得很沒有意義的工作.
成為一個榮耀祂的工具.
問題是你願不願意讓上帝這樣介入你的生命和工作.
我們只是想.
不如我們今天完成我們的工作就算了.
我們可不可以多一點期望.
多一點想想上帝如何在我們日常生活中.
在我們的工作中去介入呢.
最後.
在實踐上我們內在生命已變得真實.
那些水本身是照著潔淨禮儀所預備的.
是救藥整套禮儀系統的一部份.
耶穌其實是故意沒有用新的器皿.
祂故意用潔淨的缸裡面的東西去做工.
祂不是想隔對外在的條件和規條.
而是將救藥一切的潔淨.
獻祭 節旗的安排都指向自己.
十字架上的救恩.
對於教會來說.
這個神蹟就是一個記號.
告訴人們.
我們不可以只停留在外在的儀式.
我的意思是.
我們不是不做外在儀式.

$^{521}$今天你崇拜.
你會團契.
你會回到祈禱會.
每一個我們在做的事.
我們的宗教儀式 禮儀都在提醒我們.
我們是屬於上帝的.
預備好自己.
成為一個裝在上帝的榮耀器皿.
讓他們將普通的水.
讓他們將普通的水 普通的人生.
轉化成一個充滿喜樂和恩典的酒.
所以我猜到這裡.
大家都應該懂得怎樣回答這三條問題.
為什麼瑪利亞要對耶穌說他們沒有走.
就是因為瑪利亞期望.
耶穌用神蹟挽回主人家的榮耀.
而耶穌的回應提到事後.
就是因為耶穌用神蹟介入的時候.
就會將上帝的榮耀彰顯出來.
而最後猶太人的潔淨禮儀.
是在提醒人屬於上帝.
耶穌用潔淨禮儀的缸裡面不潔的水.
轉化成最美的酒.
是向我們每一個人發出邀請.
用你最真實的生命去回應上帝.
多一些期待上帝介入我們的工作崗位.
多一些期待上帝在我們內在生命做轉化的神蹟.
讓我們一起祈禱.
天母多謝你.
多謝你親自向我們展示.
一些我們以為最不值得提起的.
最羞恥的回憶.
求你幫助我們.
讓你在我們生命裡面做到神蹟騎士.
讓你轉化.
讓我們多一些期待.
期待你在我們生命裡面多一些互動.
幫助我們.
我們禱告奉耶穌的聖名給你.
阿們.

$^{561}$祝福大家在這個.
即將到來的春季裡面.
期望上帝多一些在你的生命裡面做騎士.
祝福大家.
\newpage



\section{3:12-4:3}
\label{sec:a8Vgzt3GCx8}
\textbf{公義缺席 (傳3\_12-4\_3) - 譚文倩莊傳道}
\newline
\newline
連結: \href{https://youtube.com/watch?v=a8Vgzt3GCx8}{\texttt{https://youtube.com/watch?v=a8Vgzt3GCx8}} ~~~~ 語音日期: 2026-04-05
\newline
\newline
\hyperref[sec:ZqMpnK5BxrA]{\small{< < < PREV SERMON < < <}}
~
\hyperref[sec:index_chronic]{\small{[返順時目]}}
~
\hyperref[sec:index_scriptual]{\small{[返順卷目]}}
~
\hyperref[sec:FQREQHip6Ag]{\small{> > > NEXT SERMON > > >}}
\newline
\newline
$^{1}$「唐兄」還是「唐兄」?.
「唐兄」….
「唐兄」,你丟掉了我的位置.
虛空,勞碌益處.
智慧有善,日光之下.
最後還是要指向傳道書的核心.
是說敬畏神.
敬畏神就像我們基督徒.
做人的一個大原則.
幫助我們在很快消散.
好像煙霧般的人生中.
能夠選擇一條有智慧,有方向的道路.
記得上一次嗎.
應該差不多是兩個月前.
在第三章中有一首詩.
描述人生中既定時間會發生的一些狀況.
經文中亦有提到.
神從開始到終結所做的事.
我們人是不可能完全明白的.
神這樣做.
是要我們人在祂面前有敬畏的心.
今天來到下一個段落.
三章十六至四章三節.
一開始傳道者說.
他留意到秉行公義和伸張正義的地方.
也有邪惡.
有看過《Wicked》這套戲嗎.
這套戲的故事圍繞著兩個女巫.
一個叫Alphaba.
你看到她的帽子.
她出生時全身的皮膚都是綠色的.
她是一個很聰明.
亦很有正義感的人.
但因為她不願意同流合污的關係.
所以她得罪了很多人.
不受歡迎.
另一個女巫叫Galinda.
或者Glinda.
如果你看過這套戲.
你就知道我為何特意這樣說.

$^{41}$很漂亮的女人.
亦是一個世界女人.
她被人捧到天花板.
在大家的心目當中.
她就是一個好的女巫的象徵.
故事圍繞著他們兩個之間的友情.
還有一些執政掌權的人.
如何去製造恐慌.
以及利用權力扭曲一些事實.
我們先來一起看看今天的經文.
傳道書三章十六,四至三節.
我又見日光之下,應有公平之處有奸惡.
應有公義之處也有奸惡.
我心裡說,神必審判義人和惡人.
因為在那裡,各樣事務.
一切工作都有定時.
我心裡說,為世人的緣故.
神考驗他們.
讓他們看見自己不過像走獸一樣.
因為世人遭遇的,走獸也遭遇.
所遭遇的都一樣.
這個怎樣死,那個也怎樣死.
他們都有一樣的氣息.
人不能強於走獸,全是虛空.
道貴一處,都是出於塵土.
也都歸於塵土.
誰知道人的氣息是往上升.
走獸的氣息是下入地.
總而言之.
人能夠在他經營的事上喜樂.
是最好不過了.
因為這是他應得的報償.
他身後的事,誰能領他回來看.
日光之下所發生的一切欺壓之事.
看那受欺壓的流淚,無人安慰.
欺壓他們的有權勢,也無人安慰.
因此,我讚嘆那已死的死人.
勝過那還活著的活人.
那尚未出生,就算未曾見過.
日光之下所發生之惡事的.

$^{81}$比這兩種人更幸福.
若果神在我們每一個人身上.
都有一個美好的計劃.
傳道者在此有點懷疑.
因為他至少看到一個問題.
就是這個世界很明顯地存在不公義.
邪惡和痛苦.
傳道者說.
他看到應該最公道的地方也沒有公道.
新疆正義的地方也變質了.
這裡所說的不公義.
其實沒有被具體界定.
所以我們不可以單單只在想.
法庭上的公義.
因為事實上.
不公義是可以在我們的社會結構中存在.
或者在一些應該負責維持公義的人身上出現.
很多時候我們失望.
不是因為這個世界有壞人.
而是因為我們對一些有權,有能力的人.
或者制度有期望.
事實上,你和我都不可以否認.
人的心一旦牽涉權力,利益的時候.
就很容易被扭曲.
傳道者似乎一早就看到.
問題不一定在政策上.
而是坐在執政掌權的人的心.
到底是怎樣的.
坦白說,其實我們人.
很多時候都會嘗試去秉行公義.
以及做一些正直的事情.
只不過結果有時候是很強差人意.
或者準則不是很一致.
結果可以帶來一些很荒謬的結果.
例如,不知道你們是否知道.
我們美國在某些州份裡面規定.
如果一個未成年的少女要打耳洞的時候.
要有父母的同意和簽名.
但在同一個州份裡面.
如果一個未成年的少女要墮胎的時候.

$^{121}$不需要父母知道.
父母不需要簽甚麼.
聽起來是否很荒謬.
是否會很令人質疑.
究竟我們所定的一些法律.
是否真的在嘗試保護我們這些小朋友.
因為聽起來完全不合理.
正如經文所說.
應有公平之處有奸惡.
應有公義之處也有奸惡.
這裡的奸惡.
在希伯來文的原文的意思.
有「歪了」,「彎了」的意思.
本來是正的,是直的.
但當人自己去定標準的時候.
就會「歪了」,「彎了」.
即是不一致,甚至可以互相矛盾.
這類事件正正反映了一個問題.
就是就算我們人最精英的制度.
法庭,法律.
去嘗試實現公義的時候.
結果也可以出現這麼大的落差.
當然不是一個單一的事件.
而是一個模式.
因為我們人會自以為.
我們能夠建立一個完全公義.
完全正直的系統.
但實際上都是「歪了」,「彎了」.
不完全.
或者公義和正直混為一談.
這正是要呼應傳道者所說的觀察.
問題不是在制度中,而是我們人本身.
其實是當人嘗試去做神.
去定義對錯的時候.
最終出來的公義和正直.
就會變得不太穩定,甚至變形.
其實「歪了」,「彎了」的「心」.
容易將審判之處變成姦惡之處.
或者你會說.
換另一個人或另一批人上去就行了.

$^{161}$但歷史告訴我們.
換另一個人,另一批人上去.
只是歷史重演.
不知道是否傳道者.
自己也站在這個權力之上.
反而更清楚看到制度的腐敗.
以及人性的醜陋,扭曲.
就算你是最有話事權的人.
你也改變不了人心的扭曲.
三章十七節.
「我心裡說,上帝必審判義人和惡人.
因為在那裡,各樣的事務,一切工作.
都有定時」.
可能我們會問,那是甚麼意思.
如果真的會有審判的時候.
為何不趕快現在就去做.
為何有些壞人在做壞事.
我們看來,他們過得似乎不錯.
他們還可以.
為何有些人在走正路.
在做好事.
他們卻要在這裡受苦.
我認識一個熱愛工作的弟兄.
他在公司工作,很鞠躬盡瘁.
早到遲走.
有一天,來了一位年輕人.
老闆對這位弟兄說.
「從今天開始,你跟這位年輕人是隊友.
你們平起平坐.
一起分擔項目的工作」.
不久後,弟兄留意到.
這位年輕人好像有些特別的待遇.
搶功勞,卸責.
最擅長擦鞋.
雖然弟兄知道,但他沒眼睛看.
「我繼續做好自己的工作.
繼續埋頭苦幹」.
誰知這位受寵的年輕人.
持著老闆看著他.
他便變本加厲.

$^{201}$有一天,走到弟兄的桌上說.
「我不喜歡有甚麼人來找你.
你要給我一點尊重」.
一位懂得擦鞋.
一位只懂得大叫.
越來越多不公平的事情發生.
弟兄決定「不行,我要找老闆談談.
究竟發生甚麼事」.
老闆說「沒有,沒這回事」.
但情況越來越差.
如果你是這位弟兄,你會怎麼做.
你猜猜後來會怎樣.
沒多久,老闆發現這位年輕人.
其實真是大隻狗.
不僅吞噬.
實際上很多事情,他都不懂得做.
原來大部分事情.
是這位不吭不響的弟兄在背後做的.
所以他立刻解僱了這位年輕人.
然後向弟兄道歉.
為表示誠意的關係.
他便加上這位弟兄的薪金.
還給了他一份獎勵.
你就想吧.
其實捱了一段時間後.
這位弟兄終於受不了.
他還未找到新的工作.
他離開了這份工作了10年的工作.
聽到這些故事的時候.
我們都很想這位壞人.
這位年輕人被懲戒.
有些事情會不好的發生在他身上.
有權的人去秉行公義,坐強扶弱.
幫助這位弟兄報仇雪恨.
把這位年輕人打得落花流水.
跪下道歉,尊敬他.
我們是想這樣做的.
但為何不是這樣.
壞人很順利.
好人在這裡要受委屈.

$^{241}$我們基督徒會想.
「神啊,為何你還未出手」.
「為何你不是馬上秉行公義」.
「為何不是馬上可以有那個後果」.
因為對我們來說.
真正的公義應該是要即時發生的.
要即時有懲罰.
我們作為父母.
沒理由小朋友做錯了一些事.
然後沒事,一個月後跟他說.
然後又罰他一個月前的事.
是要即時處理的.
「神,為何你不是這樣做」.
如果沒有報應,或者還要我等.
那就是沒有公義.
但在這裡,正正在這裡.
有一個問題我們要反思一下.
當我們這樣說,或者當我們這樣想的時候.
其實我們是否在教神怎樣做事.
我們是否在要求神.
在我們覺得正確的時候去做事.
聖經在彼得後書三章九節說得很清楚.
它說「神沒有拖延,而是祂在忍耐」.
「主沒有拖延祂的應許」.
「就如有人以為祂要辭延」.
「其實祂是寬容你們,不願一人沉淪」.
「而是人人都來悔改」.
意思是,如果神要一早或者立即審判的時候.
你和我,其實有機會悔改嗎.
你和我,有機會回轉嗎.
神之所以延遲審判.
不是因為祂不公義.
而是因為我們所信的神.
祂是有憐憫,有忍耐.
祂選擇等候.
願意更多人能夠得著祂.
願意更多人得著這個救恩.
所以傳道者說.
也有定時的意思是,審判是一定會來的.
但不是你和我的時間,而是神的時間.

$^{281}$有趣的是,這件事會繼續困擾我們.
因為我們仍然在爭奪控制.
我們仍然在爭奪主導權.
《傳道書》提到.
「神容許世界出現混亂,不公」.
其實是要試驗人.
要讓人看到.
我們人很多時候覺得我們自己就是神.
我們自己去定規則.
我們自己決定對與錯.
甚至覺得我們自己可以管理這個世界.
做得和神一樣好,甚至更好.
但我們人其實不適合做神.
為何這個世界有那麼多不公平.
正正是因為我們人在扮神.
我們人拒絕創造主.
我們覺得自己可以取代祂.
但最後產生出來的東西.
就是一些歪了,彎了的東西.
任何肯認真思考的人.
都會看到現在執政掌權的人所做的事.
有時不對勁,有時看似不正直.
某程度上,其實這些都是神憤怒.
顯明了神的憤怒.
祂沒有即時審判.
而是讓人承受扮神帶出來的後果.
你看看周圍,你就會發現.
當我們人自己作主,自己做神.
結果就是得出我們現在所見到的東西.
做得一點都不好.
想想,其實我們是否應該讓神做神.
而我們人需要調教一下我們自己.
承認我們不是神.
然後經文突然轉了一個方向.
傳道者說人和動物一樣.
都有氣息,最終都會歸於塵土.
想想,有血,有肉,有呼吸.
吃和餓,病和死.
出於塵土,歸於塵土.
想想,的確我們和動物好像沒甚麼分別.

$^{321}$但當然是有分別.
三章十一節,回到之前那裡.
它這樣說:「有把永恆放在人的心裡」.
意思是唯獨我們人裡.
有一種對永恆,神的計劃.
和超越今生意義的渴求.
不是一個感覺.
而是你整個人的存在.
包括我們的思想,意志,品格.
和我們的情感.
這個渴求叫我們相信我們的人生.
不只是我們在日光之下.
所見到的東西那麼簡單.
三章二十二節.
「總言而知,人能夠在他經營的事上喜樂.
是最好不過了.
因為這是他應得的報償」.
意思是你不要再試圖控制所有事情.
靠自己去解釋,判斷.
不要再教神怎樣做,何時做.
日光之下的視覺.
我們所看到的,所明白的東西.
其實是很有限的.
即是我們未能完全掌握神的啟示.
我們看到,但好像有點模糊.
模模糊糊.
被我們自己的有限遮蓋.
就好像在芝加哥.
冬天開車一進去的時候.
那些霞氣遮蓋了窗戶.
你看到一點,但看不清楚.
看不完.
四章一節.
日光之下所發生的一切欺壓之事.
看那受欺壓的流淚無人安慰.
欺壓他們的有權勢也無人安慰.
因此我讚嘆那已死的死人.
勝過那還活著的活人.
但那尚未出生.
就是未曾見過日光之下所發生之惡事的.

$^{361}$比這兩種人更幸福」.
四章一節,這一節裡面.
「欺壓」這個字出現了三次.
在希伯來語的字根裡面.
「欺壓之事」和「受欺壓的」是被動的.
而「欺壓的人」是一個主動的行為.
兩個被動,一個主動.
1993年有一部電影叫做《Joy Luck Club》.
其中一幕是講述一對夫婦的交易.
是講述這對夫婦.
老公Harold和老婆Lina都在工作.
他們沒有小朋友.
為了公平起見.
家裡所有的開支都50-50.
房租,水電,買菜,旅行.
他們在雪櫃門口安排了清單.
你花了甚麼錢就寫下來.
花了多少錢.
然後每個月底的時候就結帳.
確保雙方沒有人吃虧.
聽起來合理嗎.
兩人都在工作,對吧.
但其實背後有一個潛藏的問題.
這個隱憂最後因為一桶雪糕的關係.
爆發了出來.
老公Harold很喜歡吃雪糕.
經常買,經常吃.
而Lina完全沒有碰過這桶雪糕.
不是因為她的胖.
而是雪糕會令她回想起.
她童年的一個很差的遭遇.
有些童年陰影.
但到了月底的時候.
Lina也要負擔雪糕的一半費用.
有一天Lina的媽媽來探望他們.
她經過雪櫃的時候.
看到這張清單.
她很不客氣地跟這對夫婦說.
她說這個Harold根本不疼愛他老婆.
因為他不僅要他老婆.

$^{401}$負擔一半的雪糕費用.
他連他老婆怕吃雪糕也不知道.
有些不公義.
表面上不激進.
沒有欺壓,沒有暴力.
甚至很斯文,很理性.
很有制度.
這種不公義.
其實也很貼近傳道書這裡的訊息.
當公平成為了婚姻的保護機制.
兩個人只是在計算.
確保自己沒有吃虧.
慢慢地,這個雪糕已經不再是雪糕了.
因為雪糕象徵著傷害.
這種不公義.
其實可能我們經常會遇到.
問題是你是主動的那個.
還是你是被動的那個.
這個電影情節.
讓我們看到人所謂的公平,公正.
在執行出來的時候.
都可以是歪了,彎了.
不僅不完全,不公義.
也帶著很多傷害.
求主幫助我們.
我們不僅看到不公義.
卻看不到我們自己對別人的不義.
或許我們也需要反思一下自己.
表面上的公平,公正.
是不是聖經所說的公義.
當然更值得我們思考的.
是那個核心的問題.
你會否不知不覺,不留意.
就把自己變成你生命中掌管的神.
自己去定規則,自己去決定對錯.
甚至覺得自己已經可以管理一切.
做得和神一樣那麼好.
甚至比祂做得更加好.
求主的教導,聖靈的工作.
在我們每一個人心中去動工.

$^{441}$我們一起低頭祈禱.
再一次提醒我們.
其實真正可以秉行公義的.
只有你一個.
因為你是公義的主.
在你裡面,公義沒有缺席.
因為真正的公義已經在十字架上成就.
奉主耶穌基督聖名祈求,Amen.
我求主讓你們每一個人.
都帶著謙卑的心去聆聽,去領受.
我們做一個聽到,行動的機會圖.
在你自己的座位可以安靜一點時間.
然後散會.
\newpage



\section{3:14-22}
\label{sec:FQREQHip6Ag}
\textbf{成為合乎主用的人 (啟3\_14-22) - 陳崇基牧師 Becoming Useful for God's Purpose (Rev 3\_14-22) Rev. Mark Chan}
\newline
\newline
連結: \href{https://youtube.com/watch?v=FQREQHip6Ag}{\texttt{https://youtube.com/watch?v=FQREQHip6Ag}} ~~~~ 語音日期: 2026-04-05
\newline
\newline
\hyperref[sec:a8Vgzt3GCx8]{\small{< < < PREV SERMON < < <}}
~
\hyperref[sec:index_chronic]{\small{[返順時目]}}
~
\hyperref[sec:index_scriptual]{\small{[返順卷目]}}
~
\hyperref[sec:mR4JE8s_dsw]{\small{> > > NEXT SERMON > > >}}
\newline
\newline
$^{1}$[NON-ENGLISH SPEECH].
All right, good morning, guys..
[NON-ENGLISH SPEECH].
Very good to see you all..
And as you know, I came back from Turkey last--.
Saturday..
Yeah, last Saturday..
So I'm trying to reorient myself back to the Chicago time..
[NON-ENGLISH SPEECH].
So next time, if you have the opportunity,.
you should come to the Bible Lens with me..
[NON-ENGLISH SPEECH].
OK, so last week, I shared with you.
my recent trip, some of the highlights from my Turkey trip..
And it's about the first missionary journey of Paul,.
recorded in the Book of Acts..
And we visit different biblical sites.
that will help us understand the context of scripture..
So today, I want to bring you another highlight, which.
is a place called Leuticia..
[NON-ENGLISH SPEECH].
And this place is called Leuticia..
And you can find the text from the Book of Revelation,.
the last book of the Bible, chapter 3, 14 to 22..
And as you can see from this Google Maps screen cap,.
Leuticia, in terms of background,.
was a very, very wealthy city..
Now, as you hear about the context of Leuticia,.
you can think of it's kind of like Chicago, OK?.
So it was a traffic hub..
It was the center of all the traffic east and west.
and north and south..
And the banking center, right?.
Sound like our city, right?.
And it was also, in the ancient time,.
a center of textile production..
So they produced a lot of clothing.
and a lot of good wool or sheepskin..
And it's also, interestingly, a place well-known.
for its medical treatment..

$^{41}$So if you go to the hospital, especially for your eyes,.
eye medicine, it was a very well-known thing in Leuticia..
So outwardly, this city, Leuticia, and its culture.
look very, very strong and very successful.
and very self-sufficient..
[NON-ENGLISH SPEECH].
But as you turn to the text today, Revelation 3,.
you see when Jesus spoke to the church in Leuticia,.
he saw something very different from how they see themselves..
And what looked healthy, what looked successful to them,.
is actually very sick on the inside..
And when they looked rich, but they are actually.
poor in the eyes of Jesus..
When they looked strong, they were actually pretty weak.
in the sight of God..
[NON-ENGLISH SPEECH].
So even though it looks very, very successful,.
very self-sufficient, you look at the letter,.
and it says--.
and we can read from chapter 3, 14 and following--.
"To the angel of the church in Leuticia write,.
these are the words of the Amen--".
and Jesus called himself Amen because it means truthful--.
"the faithful and true witness and the ruler.
of God's creation..
15, I know your deeds, and that you are neither cold nor hot..
I wish you were either one or the other..
So because you are lukewarm, neither hot nor cold,.
I am about to spit you out of my mouth.".
[NON-ENGLISH SPEECH].
So Jesus says to the church, you are neither cold nor hot..
Sometimes we approach this statement,.
and we think that, OK, these are, of course,.
very serious words..
And we only think that hot, meaning spiritually passionate..
And then cold, meaning, of course, to us, cold water.
or cold hearted..
It may mean, OK, they are spiritually detached.
or distant, but the local geography of Leuticia,.
of this city, will help us understand this passage.

$^{81}$or those words better..
Because when you turn to the map,.
you see Leuticia is right in the middle, right?.
And then north of it is someplace.
called [INAUDIBLE] today, or Hierapolis..
That is in the north, right?.
And that is the place known for hot spring or hot water spring,.
still having mineral water, hot water coming out.
of this place, which is north of Leuticia..
And to the right is Colossae..
And Colossae was known for its cold mountain, cold spring.
water..
And both hot water and cold water were very useful..
You see that?.
[SPEAKING CHINESE].
So when I went there, this is Hierapolis, right?.
So I was actually there, even though it feels like--.
or it looks like there was a green screen..
And I keyed myself in there..
But it was a very, very beautiful place..
And it's all white..
And [SPEAKING CHINESE] actually means cotton castle..
And it gets the name because from afar,.
you look at Pamukkale, it looks white at the mountaintop..
And it has hot spring really coming out today..
And you can go there and enjoy it on barefoot..
And that is Pamukkale or Hierapolis..
And Leoticia is a little bit different..
Leoticia, right in the middle, it.
drew water from springs to the south..
But by the time the warm water or hot water.
reached the city through the water pipe,.
the water became mineral-laden..
And because you got some white stuff in your water..
And it doesn't taste very good..
And that is the image of lukewarm water.
for Leoticia..
[SPEAKING CHINESE].
So the water in Leoticia is kind of like what we call hot water..
I know Chicago is known for hot water..

$^{121}$Because some people complain to me they do a shower.
and they wash their hairs, they lose hairs..
Because hot water..
Because you got minerals in your--.
there is stuff in the water..
And it's not refreshing..
And it's not refreshing like cold water, Colossi..
It's not refreshing or therapeutic.
like the hot water in Hierapolis..
And the water in Leoticia was simply distasteful and useless..
[SPEAKING CHINESE].
I find some water pipe in the ancient time..
And it was deposited with some sort of what we call calcium.
carbonate, some of the deposit from the mineral from the water..
And as you can tell, the quality of water in Leoticia.
is very, very poor..
So Jesus is not speaking about when he says cold or hot..
He's not talking about our spiritual temperature for God..
And he's speaking about usefulness for God..
So hot water or hot spring water is useful..
Because you can dump yourself in there.
and you can enjoy the hot spring, right?.
And the cold water, of course, is useful..
You can drink it..
But the lukewarm water in Leoticia.
was useless, totally useless..
Because you cannot use it..
And in the same way, Jesus is referring.
to the Christians in Leoticia that they.
have become useless for God's purpose..
[SPEAKING CHINESE].
So this message will help us to understand the words.
in the letter to Leoticia..
And the message is for us today as well..
Because we can be successful in the sight of the world.
by human standards..
Yet we can be spiritually ineffective or useless to God..
And the Christian can be respectable,.
can be successful..
But in the sight of God, it's not truly useful..

$^{161}$[SPEAKING CHINESE].
So the question is, how can we become.
useful for God's purpose?.
And we look at at least three points from this passage..
[SPEAKING CHINESE].
First, we need to understand this short letter..
Jesus Christ is desiring for something from us, from you..
It's a fruitful relationship with you..
[SPEAKING CHINESE].
And Jesus does not look at our money..
Jesus does not look at our success..
Now, we check our portfolio every day, right?.
We check our bank account, especially tax season..
We check how much money we have and our ability..
But Jesus looks at us as whether we are abiding in him,.
whether we are building a living and fruitful relationship.
with him..
And being useful, being fruitful,.
begins with our fellowship with Jesus..
[SPEAKING CHINESE].
Now, the problem with the seer was not that they had money..
The problem-- they had money, that's fine..
The problem is that they no longer.
felt that they needed Christ when they are rich..
Their wealth had made them so self-sufficient,.
and they don't need God..
And they have comfort, and that would weaken their relationship.
or dependence on the Lord..
[SPEAKING CHINESE].
Now, you think about it..
This can happen to all of us, right?.
This can happen to all of us..
We may still come to church..
We may still serve, or we look outside, from the outside..
We look faithful..
But inwardly, we may not be depending on Jesus deeply..
And we may be active, but we are not.
abiding in him, having a living relationship with God..
[SPEAKING CHINESE].
But if we want to be useful for the purpose of God,.

$^{201}$first and foremost, we must stay close to Christ..
Stay close to Christ..
We must walk with him, and listen to him,.
and depend on him, and love him..
This is a fruitful life, and that.
will continue to grow out of your living relationship.
with God..
And you become useful for God's purpose..
[SPEAKING CHINESE].
And then the second point here, Christ desires to dwell in your heart..
[SPEAKING CHINESE].
Verse 17, it says, "You say," speaking.
to the Christians in Laodicea, self-sufficient Christians,.
"You say, I am rich..
I have acquired wealth and do not need a thing..
But you do not realize that you are wretched, pitiful, poor, blind,.
and naked.".
So Jesus said to the church, "I counsel you to buy me gold.
we find in the fire.".
Now, you said you have gold, but you don't really have gold..
You have to buy gold from me we find in the fire.
so you can become rich..
And you get the white clothes to wear so that you.
can cover your shameful nakedness..
And then you can put ointment on your eyes so that you can see..
[SPEAKING CHINESE].
So as you can see from the context,.
Jesus is using the very thing the city of the people.
were proud of to expose their true poor condition..
And they were proud of their banking, proud of their gold.
and riches, but they were actually poor spiritually..
And that they were proud of their textile, their clothing,.
their industry, but they were actually naked spiritually..
They were proud of their eye medicine,.
but they were actually spiritually blind..
[SPEAKING CHINESE].
And this is a warning to all of us, right?.
It's possible for us to be outwardly successful..
Outwardly, we get everything self-sufficient..
But inwardly, we can be empty..

$^{241}$And it is possible for us to gain many things in life,.
but we are still missing the most important thing..
And what matters most is that the presence of Christ.
in our hearts..
[SPEAKING CHINESE].
That's why Jesus said, I cancel you..
You listen to me..
You buy gold from me..
You buy white cloths from me..
And you put on your eyes ointment from me.
so that you can see..
In other words, what truly is needed.
could only come from Jesus Christ our Lord..
[SPEAKING CHINESE].
And pay attention to this..
Only, only Christ can make you truly rich..
And only Christ can cover your shame, your sin..
Only Christ can open our eyes..
Only Christ can make us whole..
[SPEAKING CHINESE].
And the third point here, the letter says,.
Jesus is calling you to repentance.
and to renew fellowship with him..
[SPEAKING CHINESE] Now, verses 19 and following,.
it says, those whom I love, I rebuild and discipline..
So be earnest and repent..
And these are very, very strong words..
Actually, among the seven letters,.
this one, the last one to live this year,.
is the strongest letter..
But you listen to the words..
They are actually words of love..
Because Jesus rebuilds his church.
not because he has given up on them..
It's just because he still loves them..
[SPEAKING CHINESE].
And in verse 20, he continues..
Jesus says, here I am..
I'm standing at the door, and I'm knocking..
And this is a moving picture..

$^{281}$The church is trying to push Jesus to the outside,.
trying to etch him out..
And yet Jesus was still standing there and still knocking.
and still calling them back..
[SPEAKING CHINESE].
And Jesus says, if anyone hears my voice and opens the door,.
I will come in, and I will eat with him, and he be with me..
Now, in the ancient world, especially in the ancient eastern world,.
sharing a meal together or sharing a table.
was a picture of close relationship..
And Jesus is offering something so intimate.
to restore our communion with him..
[SPEAKING CHINESE].
Now, this means if the leviticus is reminding you.
you are lukewarm today, remember it's not.
your spiritual temperature..
If you're lukewarm today, it means in God's sight,.
we are not that useful, right?.
We're not that close to him..
The answer to lukewarmness, all right,.
lukewarmness of our heart condition is not trying harder..
It's not doing more things..
The answer is opening the door again to Christ our Lord..
And it is turning back to him..
It's making room, opening the door to dwell with him.
so that Christ can dwell at the center of our lives again..
[SPEAKING CHINESE].
Now, some of us, some of us here may not.
have rejected Jesus openly, right?.
We're not anti-Jesus..
But perhaps we have slowly drifted away from him..
And we still know the language of faith,.
but we no longer live in a close fellowship with him..
And yet, today, from this letter,.
Jesus is still knocking at your door..
[SPEAKING CHINESE].
So if you think about it, you know, Charlie, reflect..
Jesus may have been still knocking on your door..
[SPEAKING CHINESE].
He is calling us back, and not to shame us, but to restore us..

$^{321}$He wants to renew our heart, come into our hearts,.
renew our vision, and renew our usefulness for God's purpose..
[SPEAKING CHINESE].
So the message to the Christians and the seer.
is also a message to all of us..
Do not be content, do not be content with our comfort.
while inwardly we're drifting away from Christ..
Do not trust in our own strengths,.
and do not say, I don't need a thing from God..
[SPEAKING CHINESE].
Instead, today, we turn to Jesus and abide in him..
Open your hearts to him, and let him search you,.
heal you, and restore you..
And this is how we become useful for God's purpose..
[SPEAKING CHINESE].
Now, at the end of this letter is a series.
of words of comfort..
Jesus gives us promise at the end..
To him who overcomes, to him who is faithful,.
I will give the right to sit with me on my throne..
You see, this is invitation of grace, invitation of hope..
And what an invitation, right?.
We etch him out, but he's knocking on our door..
And Christ not only wants us, he welcomes us..
[SPEAKING CHINESE].
May God keep us from being lukewarm, being useless..
May he make us, all of us, to be close to him.
and full of grace and repentance..
And truly, we want to be useful..
We want to be pleasant for God's purpose..
[SPEAKING CHINESE].
Now, at this point, we want to have some quiet moment..
And I want to invite you to speak to the Lord who.
is knocking on your door..
[SPEAKING CHINESE].
Let's spend some time and pray together, and I will close..
[SPEAKING CHINESE].
Let's pray..
Heavenly Father, we thank you for speaking to us.
through your Word today..

$^{361}$[SPEAKING CHINESE].
Forgive us for the times we have become self-sufficient,.
lukewarm, and detached from Christ..
[SPEAKING CHINESE].
So draw us back to you today into a living and fruitful.
relationship with you..
[SPEAKING CHINESE].
And dear God, we ask that you open our eyes.
to see our true condition..
And we want to welcome Jesus more fully into our hearts..
[SPEAKING CHINESE].
And Jesus, do make us a people who are close to you.
and quick to repent and useful for your purpose..
[SPEAKING CHINESE].
And we pray all this in Jesus' name..
[SPEAKING CHINESE].
Amen..
Well, may God bless you all, and I.
want to see you next week..
And become useful to desire a better relationship with God.
and to worship him, OK?.
[SPEAKING CHINESE].
And I know Patrick has--.
\newpage



\section{以弗所書 3:1-21}
\label{sec:mR4JE8s_dsw}
\textbf{復活節主日講道 (以弗所書3\_1-21) - 陳崇基牧師 Easter Sunday Sermon 2026 (Eph 3) - Rev. Mark Chan}
\newline
\newline
連結: \href{https://youtube.com/watch?v=mR4JE8s-dsw}{\texttt{https://youtube.com/watch?v=mR4JE8s-dsw}} ~~~~ 語音日期: 2026-04-09
\newline
\newline
\hyperref[sec:FQREQHip6Ag]{\small{< < < PREV SERMON < < <}}
~
\hyperref[sec:index_chronic]{\small{[返順時目]}}
~
\hyperref[sec:index_scriptual]{\small{[返順卷目]}}
~
\hyperref[sec:FC_bOIHzaL8]{\small{> > > NEXT SERMON > > >}}
\newline
\newline
以弗所書 3:1-21
\newline
\begin{longtable}{cl}
\hline
\hline
章節 & 經文 (和合本修訂版)\\
\hline
3:1 & \begin{tabularx}{0.7\textwidth}{X} 因此，我—保羅為你們外邦人作了基督耶穌囚徒的，替你們祈禱。 \end{tabularx} \\ \\ \relax
3:2 & \begin{tabularx}{0.7\textwidth}{X} 想必你們曾聽見神賜恩給我，把關切你們的職分託付我， \end{tabularx} \\ \\ \relax
3:3 & \begin{tabularx}{0.7\textwidth}{X} 用啟示讓我知道福音的奧祕，正如我以前略略寫過的。 \end{tabularx} \\ \\ \relax
3:4 & \begin{tabularx}{0.7\textwidth}{X} 你們讀了，就會知道我深深了解基督的奧祕； \end{tabularx} \\ \\ \relax
3:5 & \begin{tabularx}{0.7\textwidth}{X} 這奧祕在以前的世代沒有讓人知道，像如今藉著聖靈向他的聖使徒和先知啟示一樣， \end{tabularx} \\ \\ \relax
3:6 & \begin{tabularx}{0.7\textwidth}{X} 就是外邦人在基督耶穌裡，藉著福音，得以同為後嗣，同為一體，同為蒙應許的人。 \end{tabularx} \\ \\ \relax
3:7 & \begin{tabularx}{0.7\textwidth}{X} 我作了這福音的僕役，是照著神的恩賜，是照他運行的大能賜給我的。 \end{tabularx} \\ \\ \relax
3:8 & \begin{tabularx}{0.7\textwidth}{X} 雖然我比眾聖徒中最小的還小，他還賜我這恩典，讓我把基督那測不透的豐富傳給外邦人， \end{tabularx} \\ \\ \relax
3:9 & \begin{tabularx}{0.7\textwidth}{X} 又使眾人都明白甚麼是歷代以來隱藏在創造萬物之神裡的奧祕， \end{tabularx} \\ \\ \relax
3:10 & \begin{tabularx}{0.7\textwidth}{X} 為要在現今藉著教會使天上執政的、掌權的知道神百般的智慧。 \end{tabularx} \\ \\ \relax
3:11 & \begin{tabularx}{0.7\textwidth}{X} 這是照著神在我們主基督耶穌裡所完成的永恆的計劃。 \end{tabularx} \\ \\ \relax
3:12 & \begin{tabularx}{0.7\textwidth}{X} 我們因信耶穌，就在他裡面放膽無懼，滿有自信地進到神面前。 \end{tabularx} \\ \\ \relax
3:13 & \begin{tabularx}{0.7\textwidth}{X} 所以我求你們，不要因我為你們所受的患難喪膽；這原是你們的光榮。 \end{tabularx} \\ \\ \relax
3:14 & \begin{tabularx}{0.7\textwidth}{X} 因此，我在父面前屈膝— \end{tabularx} \\ \\ \relax
3:15 & \begin{tabularx}{0.7\textwidth}{X} 天上地上的各家都是從他得名的－ \end{tabularx} \\ \\ \relax
3:16 & \begin{tabularx}{0.7\textwidth}{X} 為要他按著他豐盛的榮耀，藉著他的靈，使你們內心的力量剛強起來； \end{tabularx} \\ \\ \relax
3:17 & \begin{tabularx}{0.7\textwidth}{X} 又要他使基督因著你們的信住在你們心裡，使你們既在愛中生根立基， \end{tabularx} \\ \\ \relax
3:18 & \begin{tabularx}{0.7\textwidth}{X} 能夠和眾聖徒一同明白基督的愛是何等的長、闊、高、深， \end{tabularx} \\ \\ \relax
3:19 & \begin{tabularx}{0.7\textwidth}{X} 並知道這愛是超過人的知識所能測度的，為要使你們充滿神一切的豐盛。 \end{tabularx} \\ \\ \relax
3:20 & \begin{tabularx}{0.7\textwidth}{X} 神能照著運行在我們心裡的大能充充足足地成就一切，超過我們所求所想的。 \end{tabularx} \\ \\ \relax
3:21 & \begin{tabularx}{0.7\textwidth}{X} 願他在教會中，並在基督耶穌裡，得著榮耀，直到世世代代，永永遠遠。阿們！ \end{tabularx} \\ \\
[1ex]
\hline
\hline
\end{longtable}
$^{1}$Well, as I said, today is a very special day, it's Easter Sunday and also the 10th anniversary.
for our UV family..
[Chinese language].
And I have a special word for you, okay, it's on the screen, the book of Ephesians,.
chapter 3, verses 20 to 21..
[Chinese language].
The.
first one I want to mention is that God has already done much for each of us..
Each and every one of us, and God has still more to do..
Now earlier in this letter, Ephesians, in the previous two chapters, Paul already reminded.
us what spiritual blessings we already have in Christ, the forgiveness of sin, the grace,.
the peace, and the new identity as the people of God..
[Chinese language].
Those who were far away from God, and now they are being brought near..
And those who were divided, Jews and Gentiles, they have made new, a new one in Christ..
[Chinese language].
Now Jesus has done much great things in the church and for the church..
So we do not, as a church, we do not stand on human strength or, but in the saving work.
of Christ..
Now this alone already give us more than we could ask for, and more blessing than we could.
ever claim..
[Chinese language].
But Paul did not stop here, he goes to chapter three, and he begins to pray for the church..
[Chinese language].
And he prays that the church would be filled with all the fullness of God, and he closes.
with this great prayer, with the declaration we see on the screen, God is able to do immeasurably.
more than you can imagine and you can ask for..
[Chinese language].
That means for a church, for a mature people of God, efficiency..
Even for a church already received so much, God is not finished with us..
[Chinese language].
So say to yourself, there is more, there is more strength to receive, there is more depth.
to know our God, more fullness to experience, and more than anything God wants his people.
to experience and through us..
[Chinese language].
Now this is especially filling for Easter Sunday, this is especially filling for Easter.
Sunday because on Easter Sunday we look back, we look back what God has done, Jesus has.
already done for us..
He died for our sin and rose from the death for all of us..
[Chinese language].

$^{41}$But on Easter we don't just look back, we look forward as well because Christ was risen,.
he is still at work and his power is not only in the past to be remembered but in the future.
to empower the church..
[Chinese language].
So God is still building his church, God is still standing his people, and Christian life.
always look back and forward, and forward, look into the future and we remember what.
God has done but we trust for what he has done and we trust that he will continue to.
do things among us..
[Chinese language].
So in the same way, our 10th anniversary, we are 10 years old, we look back and we look.
forward..
[Chinese language].
So take a look around you, 10 years ago we started with zero in this space, 10 years.
ago we didn't have money, 10 years ago we didn't have any structure..
We had to raise 120K to kick off this ministry, to launch this ministry, and that time I tell.
you, that time I know some of you were here, that time it felt enormous challenge..
[Chinese language].
But God has provided, he moved many people to give generously and he opened doors for.
us..
[Chinese language].
And over the 10 years as you can see God has brought new people and changed lives and carried.
us through many difficult seasons before us..
[Chinese language].
So today we say and we pray that Lord you have done great things and more than what we.
have imagined in the beginning..
[Chinese language].
However, we also look forward, we look forward because God is not done with us, he is not.
finished with us yet..
[Chinese language].
Now when you look forward, you are not looking at yourself, we are thinking about the emerging.
generations, maybe younger pastors, maybe younger leaders and servants and disciples.
to come after us..
[Chinese language].
That's why Ephesians 3 is a very big verse for us because it says to him be glory in the.
church and Christ throughout all, every and each generation..
[Chinese language].
That means UV Church, we never live for this moment, we live for the generations to come..
[Chinese language].
And we all understand this..
Because if you think about it, one day we will pass away and we want to prepare the.

$^{81}$best for our next generation, be them the children, be them the next generation, but.
we want to leave something stronger for the emerging generation, for the next generation..
[Chinese language].
So you want to prepare for the next generation, the emerging generation, you want to see young.
pastors to rise up, you want to see new leaders to be rising up, and the church to be alive.
and to be stronger and biblically grounded, missionally minded and full of love for Christ..
[Chinese language].
And that's why today is not only a celebration..
[Chinese language].
Easter looks back and forward, our 10th anniversary looks back and also we look forward..
And God will lead us to a new year..
And that's why our new yearly theme in the coming year is next..
N-E-X-T..
[Chinese language].
And when you think of this word, you think of you stepping into future and you prepare.
well for those who come after you..
And next also means we believe that God has done much and still has more to do..
[Chinese language].
So here you have to pay attention..
Okay?.
God has done more for us..
So we must go to the next level..
[Chinese language].
So what does it mean practically?.
If God is leading us what is next, if we think of next, the future, we must be willing to.
go to the next level, up the next level in at least three areas..
Okay, let me give you these three areas..
It's time and talent and treasure..
[Chinese language].
Let me talk to you about the first thing is time..
[Chinese language].
This talks about our love for God and our love for the community..
And many people give God the leftover time and treat church as something secondary or.
the footnote and it fits in when God is convenient to our schedule..
[Chinese language].
Now today if Christ is risen, Christ is risen and God has been this faithful to you and.
to me and he deserves your next level commitment..
Right?.
And make him more room and time for worship, for the Word of God, for fellowship, choosing.
commitment over convenience..

$^{121}$[Chinese language].
So think about it..
If our children, if our children see that we love God and honor his Word and commit ourselves.
to his people, then we are preparing something beautiful for them..
[Chinese language].
Now the second point is talent..
It's all of you..
You have abilities and you have talents..
[Chinese language].
Now God has called each and every one of us not to be spectators..
We're not just observing other people doing..
God has given and calling us to use our gifts and our abilities to serve the church, to.
serve the body..
[Chinese language].
Now some of you can lead and some of you can teach..
Some of you can pray and some of you can give..
We all have a role to play..
[Chinese language].
The big point is that everyone can offer something..
You can offer something..
And the next level question for you is that you pray that, "Dear God, how do you want.
me to use more?.
Use me more for your glory.".
[Chinese language].
If we want new leaders, new generations to rise up, then we are the mentors and we set.
the example..
We set the pattern for them..
[Chinese language].
Time, talent, and treasure..
Now this talks about our life as the disciples of Christ to practice generosity..
You know, from the very beginning, from day one, I tell you, UV has been a testimony of.
the faithfulness and provision of God through the generosity of the people..
[Chinese language].
And generosity is actually the DNA of UV since day one..
Because the cross of Christ, the cross of Jesus, shows us the generosity of His love.
and grace..
And the resurrection shows us the rich blessing we have in God..
[Chinese language].
So you want to prepare the best for your next generation, then you must invest in the church.
right now..

$^{161}$[Chinese language].
A stronger church for the next generation does not happen by accident..
It never happened by accident..
[Chinese language].
And we give, we sacrifice, we build because we see the future..
[Chinese language].
And by the grace of God, that someone to invest is you and me..
[Chinese language].
So this is our challenge today..
[Chinese language].
The next level is our time..
To up the level is our talent..
To up the level is our treasure to invest..
[Chinese language].
So I tell you this day, Easter Sunday and our 10th anniversary, we look back but we also look forward..
[Chinese language].
Christ has done great things for us in the past..
[Chinese language].
And we look back for God's faithfulness in the past 10 years..
[Chinese language].
So God has done so much, but He's not done with us..
He is able to do much more, more than you ask for, more than you imagine..
[Chinese language].
So we look forward, not only for ourselves, but the next emerging generation of a church that will pass on our faith,.
pass on our tradition, pass on our value to raise disciples and to equip leaders..
[Chinese language].
So next is for you and me to up the level, to love God more and serve God more..
[Chinese language].
So may God receive glory in our church. Amen..
May God bring us to the next level..
[Chinese language].
Okay, let's pray together..
Lord, we thank you for the cross and the empty tomb. God has done so much in Jesus, and more than we could ever deserve..
[Chinese language].
And we want to thank you..
[Chinese language].
In the past 10 years, your faithfulness to our UV Church and providing more than we imagine..
[Chinese language].
Use our gifts in service and open our hands in generosity and help us to prepare the best for the next generation..
[Chinese language].

$^{201}$And God, we pray that you raise up young people and young leaders and build a stronger church that will bring you glory from generation to generation..
[Chinese language].
Dear God, do more than we ask or imagine according to your power..
[Chinese language].
\newpage



\section{創世記 22:15-19}
\label{sec:FC_bOIHzaL8}
\textbf{Sand (Gen 22\_15-19) - Rev. Mark Chan 海沙 (創世記22\_15-19) - 陳崇基牧師}
\newline
\newline
連結: \href{https://youtube.com/watch?v=FC_bOIHzaL8}{\texttt{https://youtube.com/watch?v=FC\_bOIHzaL8}} ~~~~ 語音日期: 2026-04-18
\newline
\newline
\hyperref[sec:mR4JE8s_dsw]{\small{< < < PREV SERMON < < <}}
~
\hyperref[sec:index_chronic]{\small{[返順時目]}}
~
\hyperref[sec:index_scriptual]{\small{[返順卷目]}}
~
\hyperref[sec:Lyb7YfDNUgY]{\small{> > > NEXT SERMON > > >}}
\newline
\newline
創世記 22:15-19
\newline
\begin{longtable}{cl}
\hline
\hline
章節 & 經文 (和合本修訂版)\\
\hline
22:15 & \begin{tabularx}{0.7\textwidth}{X} 耶和華的使者第二次從天上呼喚亞伯拉罕， \end{tabularx} \\ \\ \relax
22:16 & \begin{tabularx}{0.7\textwidth}{X} 說：「耶和華說：『你既行了這事，沒有留下你的兒子，就是你的獨子，我指著自己起誓： \end{tabularx} \\ \\ \relax
22:17 & \begin{tabularx}{0.7\textwidth}{X} 我必多多賜福給你，我必使你的後裔大大增多，如同天上的星、海邊的沙。你的後裔必得仇敵的城門， \end{tabularx} \\ \\ \relax
22:18 & \begin{tabularx}{0.7\textwidth}{X} 並且地上的萬國都必因你的後裔得福，因為你聽從了我的話。』」 \end{tabularx} \\ \\ \relax
22:19 & \begin{tabularx}{0.7\textwidth}{X} 於是亞伯拉罕回到他僕人那裡。他們一同起身，往別是巴去，亞伯拉罕就住在別是巴。 \end{tabularx} \\ \\
[1ex]
\hline
\hline
\end{longtable}
$^{1}$我們想要選擇信心.
What we should not be thinking is just our program, our strategy, but your identity is who you are.
今年我們教會的主題是Next.
至此我們慶祝教會十周年之際.
我們正在想神在下一個階段要護照我們做什麼呢.
但在我們要開始下一步該做什麼之前.
首先我們需要提升的不是我們的事工項目或者是策略.
而是我們的身份.
今天呢我想用兩段經文.
帶領我們看見Next的方向.
這兩段經文在創世記的第十二章和第二十二章.
這兩段經文就像亞伯拉罕整個大故事的兩個書檔一樣.
在今天不管你是不是宗教人士.
活在這個世界上都很難對亞伯拉罕一無所知.
亞伯拉罕被猶太教基督教和伊斯蘭教共同尊崇.
而這三大一神論的信仰合起來.
他在全球擁有超過四十億的信徒.
所以在當我們打開亞伯拉罕的故事時.
我們其實是在閱讀世界上最有影響力的故事之一.
所以我們第一點最重要的信息是.
我們要信靠這位宣教的上帝.
在第一章一直到第十一章.
第十一到十一章就像整個音樂作品的開端.
在第一章你可以記得.
上帝創造了一個好世界.
祂祝福了整個人類.
但是在第三章發生了什麼?.
信徒搞砸了整個世界.
創世記第十二章並不是突然出現的.
它是介在創世記第一章到第十一章之後的.
在創世記第一章.
神創造了一個美好的世界.
並賜福給人類.
但從創世記第三章開始.
罪不斷蔓延,有悖逆,有暴力,有驕傲.
最後到了巴別塔,列國被分散.
人類試圖離開神,為自己立名.
現在讓我讀讀第十二章的文字.
我們來讀一下這一節.
創世記第十二章第十一節.

$^{41}$亞和華對亞伯拉罕說.
你要離開本地本族父家.
創世記第十二章不只是亞伯拉罕故事的開始.
也是神回應這個破碎世界的開始.
神開始展開祂賜福萬國的計劃.
神對亞伯拉罕說.
我必使你成為萬民的祝福.
我相信你聽過這個新聞.
我這裡有些文字漏掉了.
但我相信你聽過這個新聞.
上週的大新聞是亞伯拉罕第二次的發生.
亞伯拉罕第二次的發生.
上週如果你也有看過阿爾特米斯這部電影的話.
一位宇航員從月球回望地球時說了一句話.
我們選擇地球.
這是一個很美的畫面.
讓我們看見神的心.
神看著這個破碎的世界並沒有轉身離開.
祂選擇了這個世界.
祂主動走向這個世界.
祂呼召亞伯拉罕.
並藉著祂開始祂賜福萬民的計劃.
所以當我們讀到亞伯拉罕的故事時.
我們的印象就是12到22年.
我們專注在亞伯拉罕本人.
亞伯拉罕本人是一個偉大的人物.
是一個有信仰的父親.
在整個故事中的十個章節之中.
大家要明白這個故事並不是關於亞伯拉罕的偉大.
因為你們可以看到.
12到22年.
亞伯拉罕有貪婪和恐懼.
亞伯拉罕有弱點.
亞伯拉罕嘗試推翻神的承諾.
亞伯拉罕失敗.
亞伯拉罕倒下.
亞伯拉罕發生戰爭.
亞伯拉罕遲緩.
亞伯拉罕測試.
很多時候當我們讀亞伯拉罕的故事.

$^{81}$也就是《創世記》第12到22章時.
我們常常把焦點放在這個人身上.
亞伯拉罕被稱為信心之父.
但在《創世記》第12章到22章之間.
有飢荒也有懼怕.
有人性的軟弱.
也有人想用自己的辦法走在神的應許前面.
有戰爭有延遲也有試驗.
一次又一次.
神的應許似乎都受到威脅.
但這為宣教的上帝勝過一切阻攔.
為要成就他賜福萬民的應許.
所以第二點就是.
如果我們能夠信靠這位宣教的上帝.
那我們是一群帶著使命的人.
這一點在《創世記》第12章就已經很清楚了.
神對亞伯拉罕說.
我要使你成為萬民的祝福.
也就是說人要藉著你來認識神.
所以從一開始亞伯拉罕是誰呢?.
他不是為自己而被揀選的.
他是為了別人的緣故而蒙福的.
神給了他一個新的身份.
亞伯拉罕如今成為神賜福萬國計劃的一部分.
當然這個身份並不容易承擔.
在亞伯拉罕的故事裡.
有等待 有軟弱 有懼怕 有失敗 有延遲 也有試驗.
一次又一次 這個應許看起來非常的脆弱.
但是到最後我們來到《創世記》第22章.
也就是獻以撒的故事.
神要亞伯拉罕獻上他最寶貴也是唯一的兒子.
到了最後神要的是亞伯拉罕擁有一個完全降服的信心.
神並不是真的要把以撒奪走.
神是在試驗他.
你願不願意放下你的掌控?.
在亞伯拉罕的故事的結尾.
經過多年的掙扎.
神再次說話.
我必多多賜福給你.
我必使你的後裔大大增多.

$^{121}$如同天上的星 海邊的沙.
你的後裔必得仇敵的城門.
神的任務還在進行.
保證還活著.
最後 亞伯拉罕生了.
最後他有了孩子.
神的目的 為了國家 還在堅持 還在前進.
現在 神在增加一個震撼的形象.
他說 你將有後裔.
你將有兒女 孫子 孫女 如同天上的星 海邊的沙.
神在創世記第十二章開始的時候.
再次確認 這個使命仍然繼續.
這個英雄仍然活著.
神對萬國的旨意仍然堅定不移.
如今 神又加上了一個非常鮮明的圖畫.
亞伯拉罕的後裔要像天上的星那樣眾多.
又像海邊的沙那樣無數.
.
我們能數清天上的星星嗎? 並不能.
我們也不能數清海邊有多少沙粒.
這個圖畫告訴我們 神的眼光遠遠超過我們所能看見的.
.
所以我們看到圖裡 亞伯拉罕從一個小小的家庭開始.
發展成一個更大的家族.
後來成為以色列十二支派.
再後來形成一個民族.
最後 救主也為我們而來.
真正伺服世界的是亞伯拉罕的後裔.
就是耶穌基督.
歷史歷代以來.
無數的人被聚集起來.
在耶穌裡領受真正的祝福.
.
所以我們的教會不只是領受祝福的人.
我們也是要藉著神.
把我們的祝福流出去的一群人.
這就是我們的身份.
我們是一群帶著使命的人.
.
所以今天我們散會離開的時候.

$^{161}$每個人可以帶一小瓶沙子回家.
今年你可以把它放在家裡.
在整個Net系列裡.
我們還會繼續往瓶子裡加上不同的東西.
.
所以當我們把這一小瓶沙子帶回家的時候.
每一次我們看到這個瓶子都要想起來.
神的眼光比你的眼光更大.
你蒙福是為了邀請更多人.
一起來享受神的祝福.
.
今天把這個瓶子大家帶回家.
可以放在一個每天都可以看到的地方.
大家也可以拿出一個盅來試試.
到底能不能數清這個沙子到底有多少.
.
我們每個人都是這瓶沙子裡邊的一部分.
.
我們一起來祈禱.
.
主啊 求你教導我們信靠你這位宣教的上帝.
求你把我們的眼光從自己小小的計劃中提升出來.
也把你對萬國的心賜給我們.
.
現在求你幫助我們活出一群帶著使命之名的身份.
好讓別人可以藉著我們的教會.
我們的家庭.
我們的生命.
認識你的恩典.
.
奉耶穌的名禱告.
.
我們一會走的時候每個人拿一個小瓶子.
.
\newpage



\section{創世記 22:15-19}
\label{sec:Lyb7YfDNUgY}
\textbf{海沙 (創世記22\_15-19) - 陳崇基牧師}
\newline
\newline
連結: \href{https://youtube.com/watch?v=Lyb7YfDNUgY}{\texttt{https://youtube.com/watch?v=Lyb7YfDNUgY}} ~~~~ 語音日期: 2026-04-18
\newline
\newline
\hyperref[sec:FC_bOIHzaL8]{\small{< < < PREV SERMON < < <}}
~
\hyperref[sec:index_chronic]{\small{[返順時目]}}
~
\hyperref[sec:index_scriptual]{\small{[返順卷目]}}
~
\hyperref[sec:FHgJwAMW_uQ]{\small{> > > NEXT SERMON > > >}}
\newline
\newline
創世記 22:15-19
\newline
\begin{longtable}{cl}
\hline
\hline
章節 & 經文 (和合本修訂版)\\
\hline
22:15 & \begin{tabularx}{0.7\textwidth}{X} 耶和華的使者第二次從天上呼喚亞伯拉罕， \end{tabularx} \\ \\ \relax
22:16 & \begin{tabularx}{0.7\textwidth}{X} 說：「耶和華說：『你既行了這事，沒有留下你的兒子，就是你的獨子，我指著自己起誓： \end{tabularx} \\ \\ \relax
22:17 & \begin{tabularx}{0.7\textwidth}{X} 我必多多賜福給你，我必使你的後裔大大增多，如同天上的星、海邊的沙。你的後裔必得仇敵的城門， \end{tabularx} \\ \\ \relax
22:18 & \begin{tabularx}{0.7\textwidth}{X} 並且地上的萬國都必因你的後裔得福，因為你聽從了我的話。』」 \end{tabularx} \\ \\ \relax
22:19 & \begin{tabularx}{0.7\textwidth}{X} 於是亞伯拉罕回到他僕人那裡。他們一同起身，往別是巴去，亞伯拉罕就住在別是巴。 \end{tabularx} \\ \\
[1ex]
\hline
\hline
\end{longtable}
$^{1}$麥教授,時間交給他.
大家好,早安.
我們大一歲,麗仙提醒我們.
所以我們也有些需要成長的地方.
想先告訴大家發生甚麼事.
在我自己的生命中.
我過去生病了.
我剛才計算過是四個星期.
你也嚇壞了吧.
我以為我生病兩天就解決了.
我每次生病時都會告訴自己.
給自己兩天時間解決.
誰知這次很嚴重.
我在土耳其旅行的最後一星期.
我開始咳嗽.
也是很輕的咳嗽.
然後回來後,咳嗽還不清晰.
很少會這樣.
一直咳嗽下去.
在唐慶前兩個星期.
我發燒,發燒至102度.
我很少會這樣.
我吃了退燒肉,吃了就沒事.
但吃完又不退燒.
真的很霸道.
我已經退燒了.
現在已經退燒了.
結果有很關心的人跟我說.
「你一定要看醫生」.
結果我就去看醫生.
你知道我的男人不喜歡看醫生.
我們到最後,死之前才去.
結果我真的去.
用了很大的意志力.
我就說「好吧,我去看醫生」.
跟醫生談過後.
他就說「你還是驗一下新冠肺炎吧」.
我就驗了新冠肺炎.
他說「不是新冠肺炎」.
然後就說「不如你去照肺」.

$^{41}$「去照X光」.
我說「這麼嚴重?」.
他說「因為咳嗽了那麼久」.
「然後又發燒」.
「所以最好先檢查X光」.
「以免你患上瘟疫」.
即是有肺… 肺炎.
對,肺炎.
所以我就去照肺.
他就說「恭喜你,不是得到肺炎」.
「恭喜你,全部都是負面的」.
「沒事,只是脂氣管炎而已」.
脂氣管炎開了藥給我.
今天壓著我的咳嗽.
結果燒退了,但我的咳嗽還未清.
所以這三,四個星期令我有點辛苦.
我一向…我很少跟大家說.
我想你也知道.
我有個弱點,我不太喜歡認帝威.
我覺得自己很無可挑剔.
老虎打死幾隻.
不會生病.
生病一,兩天就會沒事.
我以往給別人的感覺也是這樣.
對別人的要求也是這樣.
別人生病,我覺得你生病兩天就好了.
要開工.
但這次我有個教訓.
原來自己很脆弱,對吧.
我們真的很不堪一擊.
有疾病或生命中發生的危機.
會令我們很容易被打倒.
可能你生命中也在經歷這些事情.
真的不是我們強.
我們真的很不堪一擊,很脆弱.
但最強的應該是上帝,對吧.
是祂的恩典最好.
是祂的恩典最堅穩.
所以過去三,四個星期.
我也有些反省和嘗試放鬆.

$^{81}$想清楚到底是我強還是上帝強.
其實我一點也不堅強.
我隨時會倒下.
但上帝的恩典是很堅穩的.
我們所信靠的不是人,而是上帝.
所以過去幾個星期.
他們撐著那些事工.
這次我們不說病了.
我們說上帝的話語,好嗎.
我們先祈禱,好嗎.
因為我們也是不堪一擊的人.
我們也是很脆弱,很容易倒下.
主,我們的生命.
是因為你的恩典而成在.
你的恩典和你的力量,你的信實.
是恆久不變.
主幫助我們仰望,信靠你.
特別當我們踏入新的一年時.
整個教會,我們在生命裡.
我們會去到更高之處.
更高之處不是因為我們更加厲害.
而是更加在你面前.
我們俯伏承認.
我們是軟弱,是需要你的.
你是那位最偉大的上帝.
而我們去尊崇你,讚美你.
我們這樣祈禱,奉尊名禁.
所以甚麼叫元氣大傷.
我真的有點元氣大傷.
氣不太夠,抱歉.
不過也要說….
剛才Chimmy還說跳舞,唱歌.
無論如何.
今年新的一年.
我們有一個新的年體.
就是「Next」,下一步.
下一步是甚麼.
我們想挑戰大家.
我們整個教會,每個人的生命.
我們怎樣向前一點.

$^{121}$向前一點不是說我們再強一點.
而是我們更加看到上帝.
是一位值得信靠的上帝.
所以在這次整年裡.
我們都會討論「Next」這個主題.
我自己會,其他學生不會.
但我自己會.
我選了兩段經文.
就是《創世記》12章和22章.
這次開始討論這個系列.
甚麼是「Next」.
「Next」的意思.
就是我們需要了解.
到底這位上帝.
我們所稱崇的上帝.
我們所讚美的上帝.
我們所認識的這位上帝.
到底是一個怎樣的上帝.
我覺得最好的方法開始.
就是回到第一本書.
就是《創世記》.
你看到這個簡圖.
就給你《創世記》最重要的故事.
就是阿巴拉罕的故事.
那個所謂的結構.
所以我選了12章一段經文.
和22章一段經文.
說到阿巴拉罕的時候.
可能你對聖經不太熟悉.
或者第一次接觸聖經也好.
阿巴拉罕是我們每個人.
現代人都需要認識的.
因為你知道在世界裡.
三大獨臣的宗教.
三大獨臣的信仰.
只是承認一位神.
One God Only.
三個大的信仰.
猶太教,基督教和伊斯蘭教.
都是以阿巴拉罕.

$^{161}$作為他們最尊崇的其中一個人.
三個大的宗教.
都是以阿巴拉罕作為他們的基礎.
包括基督教.
所以算起來.
所有認識阿巴拉罕.
以及以阿巴拉罕為他們的副.
這個人數可以有40億.
包括三個一神宗教所有的信徒.
所以你說你今天在這個世界裡生活.
如果你不認識阿巴拉罕.
沒有聽過阿巴拉罕的話.
就有點追不上.
所以我們要知道阿巴拉罕是誰.
這裡第一點.
我想看到阿巴拉罕的故事的時候.
第一件事.
他的故事告訴我們甚麼.
就是認識這位宣教中的上帝.
God on a mission.
原來這個上帝藉著阿巴拉罕的故事.
去呈現出這位上帝.
是一個不斷在做宣教工作的.
我們很多時候想到宣教工作的時候.
就會想起我們有沒有傳道部.
我們有沒有去傳福音.
或者某一個項目等等.
聖經從頭到尾都介紹這位上帝.
是一個怎樣的上帝.
就是一個God on a mission.
是一個向普世去宣教.
正在做工作的一位神.
如果你從創世記開始的話.
我們很簡單地說.
由開始創造這個世界的時候.
上帝的心意就是.
這個世界是一個美好的世界.
還有是一個願意賜福給全人類的狀態.
這是創世記第一章.
到第三章發生甚麼事.

$^{201}$罪闖入世界.
然後有邪惡,有暴力.
有很多驕傲的事發生.
三至十一章就告訴我們.
如果有罪闖入這個世界.
搗亂了這個世界的時候.
這個世界會是一個怎樣的地方.
到巴別塔第十一章.
就說列國分散,然後人類試圖離開.
或者控制上帝.
自己建立自己的命.
所以一至十一章告訴我們.
罪搞亂這個世界怎麼辦.
上帝是否像我們平時一樣.
電腦不行動了,我們就按reset.
或者我們有些東西不行動了.
我們就刪除它,重新再來.
不是我介紹這位上帝.
看到上帝的世界被罪人所搗亂之後.
他是沒有放棄的.
他是進入這個世界.
並且展開他要萬國得福的救恩計劃.
這就是神的回應一步.
就是創世記第十二章.
上帝在這個搗亂的世界.
你今天看看我們的世界是否搗亂了.
搗亂得不可以再搗亂.
但是上帝沒有放棄,沒有放棄.
他呼召阿巴漢出來.
他去展開一個新的計劃.
所以這個就是介紹這位上帝.
描述這位上帝.
就是一個不斷在破碎的世界裡.
去展開他的工作.
並且去拯救萬國.
十二章這裡發生甚麼事.
由一至十一章.
我們知道只有死亡.
每件事都幾乎是game over.
十二章上帝就突然間說話.

$^{241}$好像他開始的時候.
從無到有.
他說話,神說要有甚麼就有甚麼.
上帝又再一次說話.
第十二章這裡.
跟阿巴漢說甚麼.
他說你要離開你本地,本族,本家.
往我所指示你的地方.
我使你能夠成為大國.
也賜福給你.
使你的名為大.
要使你成為別人的祝福.
為你所祝福的,我就賜福他.
救助你的,我必救助他.
就是一個保護你的.
上帝保守他.
地上的萬族都要因你阿巴漢.
以及你的子孫,你的國度.
你所有的後代來得福.
這就是上帝在黑暗裡說話的第二次.
第一次是創造的時候.
第二次就是十二章這裡.
上帝說話,然後進入混亂.
在我們生命裡最混亂的時候.
世界裡最混亂的時候.
上帝原來說我不按「重置」.
我要按「新計劃」.
新的工作要展開.
就藉著阿巴漢來展開.
上星期有一個最大的新聞.
當然就是美國射了一艘飛船上去.
「阿特曼斯2」.
繞過月球一周就回來了.
安全回來了,對吧.
繞過一周就發現沒有嫦娥.
看到月亮的一些從來沒見過的.
人類沒見過的照片.
就好像這張.
其中一個太空人說了甚麼.
他說看到一個地球.

$^{281}$一個地球這麼獨特.
在一個真空的狀態.
沒有任何東西的狀態裡.
居然有一個藍色的地點.
藍色的一個地球.
他說「我們永遠會選擇地球」.
上星期三,我們祈禱會的時候.
勵仙也有分享這句話.
但我想告訴你.
這個角度看回這個地球.
看回這個世人的時候.
「我們永遠會選擇地球」.
其實這句話是上帝對阿巴漢說的.
「我們永遠會選擇地球」.
我要拯救世人.
於是我要去呼召你.
要使你能夠成為祝福.
並且祝福不是一個死水.
不只是取得上帝的祝福.
你是一個頻道.
能夠藉著你去使萬國能夠得福.
很多時候我們看阿巴漢的故事.
第12章至22章.
我們的焦點都放在那個人身上.
即是阿巴漢有多好.
我們說阿巴漢是甚麼.
後期的文獻說他是信心之父.
信心之父.
或者以色列的國父.
我們看這些所謂創始父的故事.
創始父的故事的時候.
我們都有一種很仰慕的感覺.
即是誰開創甚麼.
他是多麼厲害等等.
誰開始了某些工作.
我們現在承載他.
承傳他的好處或者他的基礎.
我們覺得那個人特別好.
對嗎.
教會也是很多時候.

$^{321}$一個很不好的文化.
就是覺得教會的創始父.
一定是比較厲害.
其實沒有這回事.
其實我們看阿巴漢都有這種錯覺.
在某程度上,某些迷思上.
我們覺得他特別厲害.
但你看看12章至22章.
整個阿巴漢的故事.
其實是描述一個阿巴漢.
是一個很真實,很脆弱的人.
有饑荒,他就逃跑.
然後他想撿回性命的時候.
他就說謊.
或者他覺得來不及了.
他就放棄人性的軟弱.
自己找自己的方法等等.
有戰爭,有很多延遲.
有很多考驗等等.
其實很多時候他都是失敗的.
你不要看錯他.
他是一個不堪一擊的人.
我想如果阿巴漢在生的時候.
可能他很老.
如果你跟他做過訪問的時候.
我們的幻想就是.
這個阿巴漢是一個強人.
對吧?.
他說的話或他做的事情多厲害.
其實很可能.
如果你去訪問他的時候.
他都是一個充滿病態.
不堪一擊,很脆弱的一個人.
但他口中告訴你.
其實不是我可以.
是上帝可以.
所以整個阿巴漢的故事.
他敘述的時候就是說.
雖然人有很多軟弱,很多脆弱.
很多懼怕,跌倒等等.

$^{361}$但背後的上帝.
他不斷地推展他的計劃.
會否有挫敗,有威脅的時候.
上帝的應許沒有變.
上帝可以圍繞它.
上帝可以克服所有的困難.
所有的難阻.
無論那些難阻是世界給的難阻.
或者人的破碎造成的難阻.
但到最後的時候.
上帝可以做到.
是上帝可以實現他的計劃.
所以我們看阿巴漢的故事.
我們不要說這個人是多麼英雄.
父親是多麼值得我們仰慕.
不,值得我們仰慕的是上帝.
是阿巴漢背後的上帝.
如果你比較一下.
你知道舊約和新約之間.
所謂的兩約之間.
那些經外的文獻.
也有很多說阿巴漢以殺亞國.
甚至說他們的老婆.
他們有幾個老婆等等.
稱讚得他們讚口不絕.
連說阿巴漢的元配.
是多麼漂亮.
稱讚她 從頭稱讚到尾.
有這樣的女人.
可以沒有瑕疵.
這就是後期對那些創造人的幻想.
但在聖經中 你看清楚一點.
原來不是這樣.
是上帝可以克服所有人的軟弱和脆弱.
祂可以克服所有的困難.
以至祂能夠成就祂的計劃.
推進祂的計劃.
使萬國能夠得福.
這是第一點.
我們要知道這位上帝是怎樣的上帝.

$^{401}$祂可以克服我們所有的脆弱.
克服我們教會所有的不足.
克服我們生命中.
可能會再次跌倒.
但我們所信靠的上帝.
就是一個可以繼續推進祂的計劃.
祂說過的不會出爾反爾.
祂應許過的要萬國得福.
祂就會做得到.
這就是阿巴罕的故事.
第一,介紹這位上帝.
第二,既然這位上帝是這樣.
我們的身分是甚麼.
就是我們被邀請.
有這樣的恩典可以被邀請做祂的子民.
我們能夠跟上帝同宮.
帶著使命來到祂的同宮.
這一點在十二章已經很清楚.
你記得我們看那年經文.
上帝跟阿巴罕說甚麼.
我要使你成為福音.
並且藉著你使萬國能夠有福音.
藉著你讓人認識.
不是你多厲害.
是上帝是可以的.
上帝才是真神.
這就是一個新的身分.
選擇阿巴罕不是因為他好.
選擇阿巴罕不是因為他厲害.
而是給他一個新的身分.
藉著他,雖然他不行.
使他能夠有一個使命.
全新的身分.
使萬國能夠得福.
我這樣說不代表這個身分是容易的.
因為在阿巴罕的故事裡.
我們不詳細交代.
有很多他自己跌倒的地方,失敗的地方.
一次又一次.
但這個應許.

$^{441}$因為他的失敗,因為他不行.
我們看他的故事時.
就像我們看長篇劇.
會覺得主角為何那麼不濟.
又不英俊.
說話又不行.
然後他做的事又傻乎乎.
是一個很不濟的人,一塌糊塗.
就像看長篇劇,編劇在做甚麼.
但最重要是告訴你.
背後的編劇才是對的人.
這就是上帝.
他可以承擔所有人的軟弱.
到最後,阿巴罕的故事.
由12章開始至22章.
中間大概有10個考驗.
10個對阿巴罕的考驗.
可能有饑荒,軟弱,等候延遲等等.
到第22章發生甚麼事.
就是以殺的故事.
獻以殺的故事.
獻以殺的故事怎樣看都是難的.
但最重要的是.
上帝告訴阿巴罕.
你要獻上你最寶貴,最心肝寶貝的.
你最捨不得的.
但你要獻上.
我們怎樣看都覺得「怎麼搞的」.
上帝這麼沒人情.
但你看整個故事.
其實上帝不是真的要阿巴罕失去以殺.
他最主要想表達甚麼.
就是想告訴阿巴罕.
你要信靠我的時候.
要有一個怎樣的信心.
就是不要再跟我說這麼多條件.
你對我的信靠應該是全心全意的.
這才是最後的考驗.
你願意放下你的控制.
你讓我去控制.

$^{481}$駕駛的人知道吧.
有誰很喜歡駕駛.
一上車就是你駕駛.
就不喜歡別人駕駛.
有誰是這樣.
有誰有這種病態.
就是你非要駕駛tie 盤的人.
別人駕駛,你總是坐著不太安心.
別人煞車,你又走去煞車.
但你沒甚麼感受.
你明白嗎?我們有這種病態.
我們想控制.
我也是這樣.
我也是要我駕駛,不是別人駕駛.
上帝在天意撒的故事.
其實就是在說這件事.
你走開.
我開車.
並且我開的車.
走的那條路.
比任何的路都更加美好.
可能有些冒險.
可能有很多驚喜.
有很多你覺得失去了的掌控.
但上帝是比你更好的司機,對吧.
上帝是更好的司機.
所以亞伯拉罕的故事結尾.
天意撒的故事結尾.
上帝又再次說話.
上帝說了甚麼?22章.
因為天意撒已經結束了.
因為上帝不是要你的天意撒.
上帝是告訴你.
你對我的信心就是要你走開一邊.
讓我駕駛tie 盤,對吧.
讓我去控制.
22章第15節.
「你的行鳥者是沒有留下你的遺址」.
「你的獨子,你的獨生子」.
「我指著自己起誓」.

$^{521}$神說甚麼?「我必多多賜福給你」.
「使你的後裔大大增多」.
「如同天上的星,海邊的沙」.
「你的後裔必得著受敵的城門」.
「地上的萬國」.
22章又是說這句話.
「地上的萬國必因你」.
「和你的後裔而得福的」.
「因為你聽從我的話」.
如果阿拜漢經過這些跌跌碰碰.
最後他明白.
這位上帝是一座正在宣教的上帝.
這位上帝是值得我們敬佩.
值得我們信靠的.
值得讓他駕駛我們的tie 盤.
這就是結果.
上帝說「我真的會」.
「實現我的計劃」.
「我會做到的」.
「你看,我做到了」.
「我會做到的」.
「你的後裔必定會」.
「如同天上的星,海邊的沙」.
你可以想像那個情況.
獻完以殺的故事.
是一個怎樣的意境.
阿拜漢可能在海邊.
沙灘.
上帝對他說話.
「你數數海灘有多少顆沙」.
「如果你能數到」.
「我就是會賜福你」.
「這樣賜福你」.
賜福到另一個地步.
你是數不盡的.
這就是阿拜漢所介紹的上帝.
是一位願意祝福我們.
只要我們願意信靠祂.
上帝對萬國的祝福.
上帝對萬國的旨意.

$^{561}$是堅定不移的.
這幅圖畫.
天上的星那麼多.
海邊的沙那麼多.
你試試數,數不到.
是嗎?數不到的.
你看到阿拜漢的故事的交代.
由1,2,3,4,5,6.
你就發現原來上帝真的很信實.
有一個很小的家庭.
還是沒有子女的家庭.
一個死亡的家庭.
變成一個家庭.
然後變成一個更大的家族.
以撒雅各.
接著到十二支派.
以色列國成為一個民族.
到最後,是耶穌基督降生.
耶穌說「藉著我的教會」.
「要使萬國能夠成為祂的門徒」.
「萬國能夠得福」.
你會看到由阿拜漢直到耶穌基督.
上帝就是在實行祂自己的計劃.
在耶穌基督裡面.
就真正得到我們有的祖訪.
你看看我們.
你看看我們的教會.
你看看普世教會.
每一代的教會.
我們都會跌跌碰碰.
我們會散亂,會有散,有嘴.
但我們要看到的是我們背後.
原來這位上帝.
是一個值得我們事奉的上帝.
因為有阿拜漢,有耶穌基督.
因為有耶穌基督,有我和你.
誰會想像到在阿拜漢的時候.
可以想像到一個這樣的海沙.
這麼多的祝福.
不是因為仁德.

$^{601}$不是因為以色列國德.
不是因為祖父,列祖很厲害.
不是因為我們很厲害.
而是上帝很厲害.
我們在這一年,新的一年.
我們希望和你一起體驗.
去經歷這個信實的上帝.
正在宣教工作的上帝.
我們怎樣能夠和祂同共.
所以在這一年裡面.
我們開球的時候.
我們希望送一個瓶子給大家.
你拿著這個瓶子.
我們今天這個瓶子只有一樣東西.
我們將來再說這個系列的時候.
我們會加一些東西到裡面.
但我想請你今天拿這個瓶子.
拿一個,好,拿一個.
你不是一個人拿十個.
一個就夠了,為甚麼.
因為你嘗試回去拿.
拿了回去的時候.
你數一下有多少沙在裡面.
你數一下有多少顆.
你下星期告訴我.
我保證你數到下星期還沒數完.
這個就是上帝的應許.
就是上帝自己的本性.
縱使我們是軟弱.
我們不大.
我們不堪一擊.
但上帝才是無堅不摧.
祂值得我們相信.
值得我們去事奉祂.
所以請你拿一個瓶子回去.
這裡有我們的UV的蓋印.
你可以回去寫你的名字,日字.
在這一年裡面,我們都會用這個瓶子.
就好像你遠行的時候.
我們在遠足的時候.

$^{641}$我們拿著這個瓶子.
水瓶,或者你的裝備.
裡面每一樣東西去提醒我們.
這位上帝值得我們全心全意去愛.
全心全意去事奉.
對嗎.
海沙可以提醒我們.
上帝最早的應許.
可以提醒我們教會.
我們希望能夠參與上帝宣教中的工作.
我們成為海沙的一部分.
讓我們願意有更多海沙.
能夠參與到,享受到.
上帝自己的同在和祝福,好嗎.
我們同心協力去祈禱.
主,求你教導我們.
信靠的不是我們自己.
因為我們實在太渺小.
因為我們實在太不堪一擊.
我們實在是脆弱的人.
主,我們這個生命.
我們只不過是一條爛命.
但你卻以恩典,你的愛.
來拯救我們.
不是因為我們可以.
而是因為你,你是可以的.
主,你要讓我們認識你.
是一位宣教中的上帝.
要使萬國的心歸向你.
主,我們已經歸向你.
在基督裡面領受祝福.
也求你讓我們這個軟弱的生命.
能夠成為一個帶著使命的生命.
去認識你,去事奉你.
以致我們個人生命,我們的家庭.
我們每一個的每一方面.
能夠去認識你的恩典.
並且叫人去認識你自己的名字.
奉耶穌基督的名句.
好,願上帝祝福各位.

$^{681}$拜託你拿這個瓶子回去.
你下星期不用拿回來.
安放在你家裡最重要的地方.
提醒你自己,好嗎.
但我們會派….
每一次我們會多派一些東西給你.
然後你拿回去,放在你的瓶子裡.
好嗎?上帝祝福各位.
\newpage



\section{約翰福音 2:12-25}
\label{sec:FHgJwAMW_uQ}
\textbf{Cleanse your hearts of its "Cattle and Sheep" (Jn 2\_12-25) - Tsz Kin Leung 趕走心中的牛羊(約翰福音2\_12-25) - 梁紫健傳道}
\newline
\newline
連結: \href{https://youtube.com/watch?v=FHgJwAMW_uQ}{\texttt{https://youtube.com/watch?v=FHgJwAMW\_uQ}} ~~~~ 語音日期: 2026-04-28
\newline
\newline
\hyperref[sec:Lyb7YfDNUgY]{\small{< < < PREV SERMON < < <}}
~
\hyperref[sec:index_chronic]{\small{[返順時目]}}
~
\hyperref[sec:index_scriptual]{\small{[返順卷目]}}
~
\hyperref[sec:YUWGLt45Wz8]{\small{> > > NEXT SERMON > > >}}
\newline
\newline
約翰福音 2:12-25
\newline
\begin{longtable}{cl}
\hline
\hline
章節 & 經文 (和合本修訂版)\\
\hline
2:12 & \begin{tabularx}{0.7\textwidth}{X} 這事以後，耶穌與他的母親、兄弟和門徒都下迦百農去，在那裡住了不多幾天。 \end{tabularx} \\ \\ \relax
2:13 & \begin{tabularx}{0.7\textwidth}{X} 猶太人的逾越節近了，耶穌上耶路撒冷去。 \end{tabularx} \\ \\ \relax
2:14 & \begin{tabularx}{0.7\textwidth}{X} 他看見聖殿裡有賣牛羊和鴿子的，還有兌換銀錢的人坐著， \end{tabularx} \\ \\ \relax
2:15 & \begin{tabularx}{0.7\textwidth}{X} 耶穌就拿繩子做成鞭子，把所有的，包括牛羊都趕出聖殿，倒出兌換銀錢之人的銀錢，推翻他們的桌子， \end{tabularx} \\ \\ \relax
2:16 & \begin{tabularx}{0.7\textwidth}{X} 又對賣鴿子的說：「把這些東西拿走！不要把我父的殿當作買賣的地方。」 \end{tabularx} \\ \\ \relax
2:17 & \begin{tabularx}{0.7\textwidth}{X} 他的門徒就想起經上記著：「我為你的殿心裡焦急，如同火燒。」 \end{tabularx} \\ \\ \relax
2:18 & \begin{tabularx}{0.7\textwidth}{X} 因此猶太人問他：「你能顯甚麼神蹟給我們看，表明你可以做這些事呢？」 \end{tabularx} \\ \\ \relax
2:19 & \begin{tabularx}{0.7\textwidth}{X} 耶穌回答他們說：「你們拆毀這殿，我三日內要把它重建。」 \end{tabularx} \\ \\ \relax
2:20 & \begin{tabularx}{0.7\textwidth}{X} 猶太人說：「這殿造了四十六年，你三日內就能重建嗎？」 \end{tabularx} \\ \\ \relax
2:21 & \begin{tabularx}{0.7\textwidth}{X} 但耶穌所說的殿是指他的身體。 \end{tabularx} \\ \\ \relax
2:22 & \begin{tabularx}{0.7\textwidth}{X} 所以他從死人中復活以後，門徒想起他曾說過這事，就信了聖經和耶穌所說的話。 \end{tabularx} \\ \\ \relax
2:23 & \begin{tabularx}{0.7\textwidth}{X} 耶穌在耶路撒冷過逾越節的時候，有許多人看見他所行的神蹟，就信了他的名。 \end{tabularx} \\ \\ \relax
2:24 & \begin{tabularx}{0.7\textwidth}{X} 耶穌自己卻不信任他們，因為他認識所有的人， \end{tabularx} \\ \\ \relax
2:25 & \begin{tabularx}{0.7\textwidth}{X} 也用不著誰來證明人是怎樣的，因為他自己認識人的內心。 \end{tabularx} \\ \\
[1ex]
\hline
\hline
\end{longtable}
$^{1}$>> Good morning ladies, good morning brothers and sisters..
>> [FOREIGN].
>> I wonder if everybody likes to eat these kind of foods..
>> [FOREIGN].
>> Steaks, lamb belly, these foods are rich in oil..
>> [FOREIGN].
>> The fat after being cooked at a high temperature melts in your mouth and.
spreads out..
>> [FOREIGN].
>> If you only eat steak, lamb chops, or lamb belly for meal,.
we would think all find this very greasy..
>> [FOREIGN].
>> I heard that Joe Osco was having a big sale this week..
>> [FOREIGN].
>> And their Angus beef this week was in a really good deal..
I was thinking to drive in there immediately and.
to buy several more box while it was cheap..
>> [FOREIGN].
>> To my surprise, when I arrived, only empty shelves as you can see in the picture..
So I end up just grabbing three inexpensive thinly sliced steaks..
>> [FOREIGN].
>> In any case, it's very convenient for.
modern people to eat beef or lamb..
>> [FOREIGN].
>> However, in ancient times, when the Jews still had temple,.
it wasn't so easy for them to eat steak and lamb chops..
>> [FOREIGN].
>> They need to go to the temple to offer sacrifices..
Leaving some for the priest, some for God, and the rest for themselves..
>> [FOREIGN].
>> However, in the temple scene of Gospel of John,.
Jesus actually drove out these cattle and sheep that represented sacrifice..
>> [FOREIGN].
>> What message does this story of cleansing the temple in the Gospel of John.
convey to us?.
>> [FOREIGN].
>> Before reading the scripture,.
let's review the previous section today..
>> [FOREIGN].
>> Last time we mentioned that Jesus was asked to perform a miracle..

$^{41}$At the wedding in Cana, he performed an unexpected miracle..
>> [FOREIGN].
>> He transformed the unclean water originally intended for.
purification rations and stored in stone jars into fine wine..
Thus, revealing God's glory and calling people back to God..
>> [FOREIGN].
>> The following passage continues this theme..
Let's read John 2, 12 to 25..
>> [FOREIGN].
>> After this refers to the events following Jesus' transforming.
of water into wine at the wedding in Cana..
>> [FOREIGN].
>> Jesus went down to Capernaum with his mother, brothers,.
and disciples and stayed there for a few days..
>> [FOREIGN].
>> When it was almost time for a Jewish Passover,.
it's a very important holiday..
Jesus went up to Jerusalem..
>> [FOREIGN].
>> In the temple courts, he found people selling cattle, sheep, and.
dogs, and others sitting at the tables, exchanging money..
>> [FOREIGN].
>> So Jesus made the whips out of the courts and.
drove out all the sheep and cattle from the temple..
>> [FOREIGN].
>> Please note that cattle and sheep are not mentioned in.
a synoptic Gospel account of cleansing the temple..
This is the unique account found only in Gospel of John..
>> [FOREIGN].
>> They poured out the money of those exchanging silver and.
overturned their tables..
>> [FOREIGN].
>> To those who sold those, he said, get this out of here..
Stop turning my father's house into a market..
>> [FOREIGN].
>> Was Jesus opposing monetary transaction within the temple?.
The phrase, the father's temple, is unique in the Gospel of John..
It is only instance in the entire Bible where father and.
temple are used together..
>> [FOREIGN].

$^{81}$>> His disciples remember that it is written, zeal for.
your house will consume me..
>> [FOREIGN].
>> The Jews then responded to him,.
what sign can you show us to prove your authority to do all this?.
>> [FOREIGN].
>> Jesus answered them,.
destroy this temple and I will raise it again in three days..
>> [FOREIGN].
>> They replied, it has taken 46 years to build this temple and.
you are going to raise it in three days?.
>> [FOREIGN].
>> But the temple he had spoken of was his body..
>> [FOREIGN].
>> After he was raised from the dead,.
his disciples recalled what he had said..
Then they believed the scripture and the words that Jesus had spoken..
>> [FOREIGN].
>> Now while he was in Jerusalem at the Passover Festival,.
Passover is mentioned again here..
Clearly, Jesus wanted to remind him of something..
>> [FOREIGN].
>> Many people saw the signs he was performing and.
believed in his name..
>> [FOREIGN].
>> But Jesus would not entrust himself to them for.
he knew all people..
>> [FOREIGN].
>> He did not need any testimony about mankind for.
he knew what was in each person..
>> [FOREIGN].
>> What exactly did Jesus know about people's heart that led him to perform.
such a special cleansing of the temple?.
>> [FOREIGN].
>> After reading the scripture, we will find at least three main problems..
>> [FOREIGN].
>> The story of Jesus driving away the cattle and sheep is unique to the Gospel of John..
The cattle and sheep are the core of sacrifice..
What does it mean that Jesus drove them away?.
>> [FOREIGN].

$^{121}$>> The passage mentioned Passover twice..
What is the significance of cleansing the temple on Passover?.
>> [FOREIGN].
>> After Jesus cleansed the temple,.
did he know what was in people's hearts?.
>> [FOREIGN].
>> After answering these three questions,.
we will understand what Jesus wanted to tell us..
>> [FOREIGN].
>> What is also today's topic,.
cleanse your heart of its cattle and sheep..
What exactly do this cattle and sheep symbolize in our lives?.
>> [FOREIGN].
>> If you're ready, we can start..
>> [FOREIGN].
>> What did Jesus especially emphasize,.
using a whip to drive away cattle and sheep?.
>> [FOREIGN].
>> It was actually quite normal for the temple to have cattle and sheep..
During the second temple period, animal sacrifices were necessary..
>> [FOREIGN].
>> The burnt offering, an optimum offering,.
the peace offerings recorded in the latificates,.
are all associated with sacrificial animals such as cattle, sheep, and doves..
>> [FOREIGN].
>> During Passover, pilgrims from afar flocked to the temple,.
unable to bring suitable sacrificial animals..
>> [FOREIGN].
>> Therefore, sacrificial markets were set up near the temple,.
and even within its ground for the convenience of buying and using immediately..
This is reasonable and necessary from an institutional perspective..
>> [FOREIGN].
>> The sacrificial animals must be without defects and cannot be purchased casually..
>> [FOREIGN].
>> To make it convenient for people from other places and avoid the risk of injury.
and uncleanliness to the sacrificial animals to long journeys..
>> [FOREIGN].
>> Temple taxes require a specific currency..
Exchange services must be provided..
>> [FOREIGN].

$^{161}$>> This is similar to how monetary transactions are heard in more modern temples,.
such as adding oil to incense,.
buying incense in candles,.
and having fortune sticks interpreted..
>> [FOREIGN].
>> Since buying and selling cattle and sheep were originally normal,.
why did Jesus break this norm?.
>> [FOREIGN].
>> The key lies in the Father's temple..
>> [FOREIGN].
>> In the Synodic Gospels,.
Jesus quotes, "My house will be a house of prayer,.
but you have made it a den of thieves.".
>> [FOREIGN].
>> However, in the Gospel of John,.
it is the disciples who quote,.
"My soul is burning with anxiety for your house,".
while Jesus refers to it as the Father's house..
>> [FOREIGN].
>> This shows that Jesus drove away the cattle and sheep to emphasize that it was.
Jehovah's land and he did not want it to become a place of buying and selling..
>> [FOREIGN].
>> But Jesus wasn't simply against commercial activities..
He was against the motives behind buying and selling..
>> [FOREIGN].
>> In Jesus' time, temples often became places of corruptions and power struggles..
>> [FOREIGN].
>> For example, the Temple of Apollo in Greece was not only a financial center,.
but its auricular service were also frequently exploited by politicians..
>> [FOREIGN].
>> Many Roman temples also become money-laundering centers for armies and.
officials to store their spoils, engage in usury, and invade taxes..
>> [FOREIGN].
>> Officials even bought and sold priesthood and maintained their monopoly,.
forming a system of using the name of God to amass wealth..
>> [FOREIGN].
>> This culture also permitted the Jewish temple..
Beginning with the Maccabean dynasty,.
the trading system became the political tools, priestly positions were bought and.
sold, and what should have belonged to God was embezzled..

$^{201}$>> [FOREIGN].
>> Therefore, the act of cleansing indicates the temple's defilement..
>> [FOREIGN].
>> Originally intended to facilitate people's worship to God,.
the system, it has been corrupted by human greed and attachment to power..
>> [FOREIGN].
>> The problem lies not in the system itself, but in people's minds..
>> [FOREIGN].
>> Driving away the cattle and sheep is unique to the Gospel of John..
>> [FOREIGN].
>> In Leviticus, cattle and sheep are central to the peace offering,.
where people offer their fat to God and share the remainder,.
thus expressing their belonging to God..
>> [FOREIGN].
>> The reason for emphasizing Passover again is to remind us that some things.
belong to God that humans cannot transgress them..
>> [FOREIGN].
>> The disciples quote Psalm 69, which describes David's desire to build a temple.
for God, even in difficult circumstances, and his hope for.
a return of people's hearts..
>> [FOREIGN].
>> Psalm 31 says, this will please the Lord more than an ox..
This suggests that when people's heart corrupt the system,.
sacrifices lose their meaning, and even cattle and sheep become worthless..
>> [FOREIGN].
>> Chinese people value reciprocity..
>> [FOREIGN].
>> Giving gifts to express affection is a good thing..
>> [FOREIGN].
>> But if this transactional mentality is applied to one's relationship with God,.
then the priorities are reversed..
>> [FOREIGN].
>> If God bless me, I will give more..
If God does not bless me, I will give less..
This has become a transaction..
>> [FOREIGN].
>> Just like how people are tired of giving mooncakes during mid-autumn festival,.
because it has become a cold and impersonal routine..
>> [FOREIGN].
>> When we give money, time, spiritual support to the church,.

$^{241}$are we making a transactional offering, or do we believe that everything comes.
from God and we are simply fulfilling our duty to give back?.
>> [FOREIGN].
>> Jesus reminds us that cattle and sheep symbolize our hearts turning to God..
>> [FOREIGN].
>> Our lives, assets, and talents are not entirely our own..
>> [FOREIGN].
>> Driving away the sheep and cattle in our hearts doesn't mean abstaining from meat.
or offering sacrifices, but rather watching our own hearts and not transgressing in God's name..
>> [FOREIGN].
>> My wife and I discussed two methods of donation,.
automatic transfer or manual transfer..
Which one is more heartfelt?.
>> [FOREIGN].
>> In fact, both can become formalism..
>> [FOREIGN].
>> If giving is merely a matter of feeling heartache or.
pain without recognizing it as God's grace,.
then it's just an act, not a genuine affirmation that everything belongs to the God..
>> [FOREIGN].
>> The same applies to time..
When we ask, is the time I spend for God worth it?.
The underlying question is, what can I get from, what can I get in return?.
We often cling to life, believing that time belongs to us..
>> [FOREIGN].
>> Jesus' statement about rebuilding the temple in three days meant that.
he would redefine the meaning of the temple with his own body and life..
Returning human life to God..
>> [FOREIGN].
>> The Jews spent 46 years, almost a lifetime for.
most people at the time, expanding the temple..
Yet, their hearts were far away from God, a fact that grieved Jesus greatly..
>> [FOREIGN].
>> Cattle and sheep can represent stinginess and.
attachment to money, time, and talent..
>> [FOREIGN].
>> We are stingy because we have forgotten that breathing,.
relationship, and talents are all give from God..
>> [FOREIGN].
>> Jesus knew that people's heart were filled with an attachment to power,.

$^{281}$money, and time..
>> [FOREIGN].
>> He drove cattle and sheep away to remind us to let go of those attachments and.
return ownership to God..
>> [FOREIGN].
>> Let's pray..
>> [FOREIGN].
>> Heavenly Father, help us to drive away the sheep and cattle in our hearts,.
including our attachment to money, power, time, talent, and relationship..
>> [FOREIGN].
>> Please help us return all ownership to you..
>> [FOREIGN].
>> You know the struggles and the malice in our hearts..
Please lead us to start anew..
>> [FOREIGN].
\newpage



\section{約翰福音 2:12-25}
\label{sec:YUWGLt45Wz8}
\textbf{趕走心中的「牛羊」(約翰福音2\_12-25) - 梁紫健傳道}
\newline
\newline
連結: \href{https://youtube.com/watch?v=YUWGLt45Wz8}{\texttt{https://youtube.com/watch?v=YUWGLt45Wz8}} ~~~~ 語音日期: 2026-04-28
\newline
\newline
\hyperref[sec:FHgJwAMW_uQ]{\small{< < < PREV SERMON < < <}}
~
\hyperref[sec:index_chronic]{\small{[返順時目]}}
~
\hyperref[sec:index_scriptual]{\small{[返順卷目]}}
~
\hyperref[sec:bMtCjPvaMwM]{\small{> > > NEXT SERMON > > >}}
\newline
\newline
約翰福音 2:12-25
\newline
\begin{longtable}{cl}
\hline
\hline
章節 & 經文 (和合本修訂版)\\
\hline
2:12 & \begin{tabularx}{0.7\textwidth}{X} 這事以後，耶穌與他的母親、兄弟和門徒都下迦百農去，在那裡住了不多幾天。 \end{tabularx} \\ \\ \relax
2:13 & \begin{tabularx}{0.7\textwidth}{X} 猶太人的逾越節近了，耶穌上耶路撒冷去。 \end{tabularx} \\ \\ \relax
2:14 & \begin{tabularx}{0.7\textwidth}{X} 他看見聖殿裡有賣牛羊和鴿子的，還有兌換銀錢的人坐著， \end{tabularx} \\ \\ \relax
2:15 & \begin{tabularx}{0.7\textwidth}{X} 耶穌就拿繩子做成鞭子，把所有的，包括牛羊都趕出聖殿，倒出兌換銀錢之人的銀錢，推翻他們的桌子， \end{tabularx} \\ \\ \relax
2:16 & \begin{tabularx}{0.7\textwidth}{X} 又對賣鴿子的說：「把這些東西拿走！不要把我父的殿當作買賣的地方。」 \end{tabularx} \\ \\ \relax
2:17 & \begin{tabularx}{0.7\textwidth}{X} 他的門徒就想起經上記著：「我為你的殿心裡焦急，如同火燒。」 \end{tabularx} \\ \\ \relax
2:18 & \begin{tabularx}{0.7\textwidth}{X} 因此猶太人問他：「你能顯甚麼神蹟給我們看，表明你可以做這些事呢？」 \end{tabularx} \\ \\ \relax
2:19 & \begin{tabularx}{0.7\textwidth}{X} 耶穌回答他們說：「你們拆毀這殿，我三日內要把它重建。」 \end{tabularx} \\ \\ \relax
2:20 & \begin{tabularx}{0.7\textwidth}{X} 猶太人說：「這殿造了四十六年，你三日內就能重建嗎？」 \end{tabularx} \\ \\ \relax
2:21 & \begin{tabularx}{0.7\textwidth}{X} 但耶穌所說的殿是指他的身體。 \end{tabularx} \\ \\ \relax
2:22 & \begin{tabularx}{0.7\textwidth}{X} 所以他從死人中復活以後，門徒想起他曾說過這事，就信了聖經和耶穌所說的話。 \end{tabularx} \\ \\ \relax
2:23 & \begin{tabularx}{0.7\textwidth}{X} 耶穌在耶路撒冷過逾越節的時候，有許多人看見他所行的神蹟，就信了他的名。 \end{tabularx} \\ \\ \relax
2:24 & \begin{tabularx}{0.7\textwidth}{X} 耶穌自己卻不信任他們，因為他認識所有的人， \end{tabularx} \\ \\ \relax
2:25 & \begin{tabularx}{0.7\textwidth}{X} 也用不著誰來證明人是怎樣的，因為他自己認識人的內心。 \end{tabularx} \\ \\
[1ex]
\hline
\hline
\end{longtable}
$^{1}$大家早安.
不知道大家喜不喜歡吃這一類的食物呢.
給個爆點.
牛腩,牛扒,羊腩.
這一類的食物充滿著油脂.
經過高溫烹調之後.
大家就會發現.
脂肪會怎樣?入口即溶.
在口腔裡化開.
如果你一餐吃牛扒,羊扒,羊腩.
相信我們都會吃得很飽.
無論如何都要配合一些食物去吃.
這個星期我收到消息.
某間豬肉大特價.
有一隻安格斯的牛很划算.
打算飛車去買.
趁它便宜就買幾盒.
誰知去到就看見一個划算的價錢.
買光了.
所有東西都買光了.
最後我和Alice就說.
不如平平地撿起三塊薄切的牛扒.
無論如何.
在現代人吃羊,吃牛.
這些很矜貴的食物其實是非常方便.
但在古時候.
猶太人在有聖殿的時間.
如果他們想吃牛扒,想吃羊扒.
你猜他們要怎樣做?.
要宰的.
宰之前要怎樣?.
要獻祭.
對,要獻祭,沒錯.
宰之前就不是像我們去超級市場.
要便宜幾塊雞就撿一件.
是要獻祭的.
宰之前是要去到聖殿前去獻祭.
將一部分肥的脂油留給祭司.
一部分就留給上帝.
另外剩下的才是自己和別人分享.

$^{41}$在耶穌在這篇文章裡面.
耶穌在聖殿的場景.
將這個代表著獻祭的牛羊趕出去.
究竟耶穌在這個動作裡面.
帶給我們甚麼訊息呢?.
我們現在來看看.
在回顧經文之前.
我們要提一提上一個講道.
上一篇經文是說耶穌.
這位海底爺被要求做神蹟.
在迦南的婚宴裡面.
做了一個很意外的神蹟.
將最污糟的水.
潔淨禮儀石缸裡面的水.
變成了好酒.
以致上帝的形容要被彰顯出來.
叫人回歸上帝.
接下來就是這段經文.
延續著這個主題.
我們一起看看約翰福音.
二章十二至二十五節.
看看是做甚麼.
我們會細讀這段經文.
讀完這段經文.
我們才會知道究竟這段經文.
有甚麼奇怪的地方.
我們會留意.
「這事已後」是甚麼意思呢?.
這件事是指耶穌.
在迦南婚宴水變酒的故事.
耶穌和他的母親弟兄.
和門徒去到迦伯隆.
在那裡住了幾天.
接著十三節.
猶太人的雨節近了.
這裡的雨節是一個很重要的日子.
我一會兒會再提.
耶穌上到耶路撒冷去.
見到聖殿裡面有賣牛,賣羊,賣鴿子.
和兌換人錢的人坐在那裡.

$^{81}$耶穌就把繩子弄成一團.
把牛羊也趕出殿.
留意牛羊在這裡對觀福音.
即是其他福音書 馬太,馬可,路加.
都沒有提及過.
三個福音書都在描述耶穌見證聖殿.
但都沒有提及牛羊出聖殿.
所以這裡的約翰福音.
把牛羊趕出聖殿是一個很獨特的記載.
然後就推倒兌換銀錢的人.
不是他們的銀錢.
推倒他們的桌子.
就跟賣鴿子的人說.
把這些東西拿走.
不要把我付的電作為買賣的地方.
這裡耶穌正在反對一些東西.
是否叫人在聖殿裡面不要有金錢的交易呢.
這裡付的電都是約翰福音獨有的記載.
通常我們會讀到有付 也有電.
但付和電放在一起.
只有這裡才有.
我繼續讀下去.
17節就是.
他的門徒就記起.
經常記著說.
我為你的電.
心裡焦急如同火燒.
因此猶太人就問他.
你做這些事還獻什麼神蹟給我們看.
19節就說.
耶穌就跟他說.
你們拆毀這個電.
我三天就會重建.
猶太人就說.
這個電四十六年才重建.
你三天就可以重建.
耶穌就說.
是的 是用他的身體來做電.
所以從他死裡復活之後.
門徒就想起他講過這句話.

$^{121}$就相信了聖經和他所講的東西.
當耶穌經過耶路撒冷.
過這個雨節的時候.
又是雨節.
明顯耶穌是想透過這個節日.
去提醒人的一些東西.
有很多人去看見他行神蹟.
就相信了他的命.
然後耶穌就不把自己交託給他們.
因為他知道萬人.
也不需要誰見證人怎樣.
因為他知道人心裡所存的.
知道人心裡所存的是什麼意思呢.
究竟耶穌知道人心裡所存的是什麼.
以致他要做一個這麼特別的潔淨聖殿.
我們經文讀完了.
我們看看經文答給我們有什麼問題.
至少有三個問題.
第一就是耶穌把牛羊趕走.
代表著什麼.
第二就是在雨節裡面潔淨聖殿.
代表著什麼.
大家如果有記憶的話.
潔淨聖殿在其他福音書.
其實是沒有提到雨節的.
但在《約翰福音》這裡.
是專登有提到兩次雨節.
究竟代表著什麼呢.
第三,耶穌潔淨聖殿之後.
他知道人心裡所存的是什麼.
回答這三個問題.
我們就會知道究竟耶穌想跟我們說什麼.
也是今天的講題.
趕走心中的牛羊.
究竟牛羊在我們現在的年代.
我們的生命裡面代表著什麼呢.
當我想講題的時候.
我也沒有想過.
原來牛羊也是十二生肖其中之一.
不好意思,如果你屬羊.

$^{161}$你屬牛的話.
今天的講題不是針對你的.
放心,是對著我們所有人的.
不是特別針對你.
我們一起來看.
當耶穌拿著鞭子趕走牛羊出聖殿的時候.
我們就要問,為什麼要這樣強調呢.
其實在聖殿裡面有牛羊是一件很正常的事.
在第二聖殿時期.
牛羊的獻祭本身就是需要動物.
在聖殿獻祭需要有動物.
聖殿是將繁祭,贖罪祭,平安祭.
把不同的祭生放在一起.
而牛羊和鴿子是最常見的祭生.
所以當御節來到的時候.
不只是耶路撒冷裡面的猶太人.
而是外地來的朝聖者.
他們都是分散在耶路撒冷之外的猶太人.
他們會回到耶路撒冷.
不少人是很遠路的.
可能是幾十米,幾百米的路程.
回到耶路撒冷.
一起去敬拜上帝,去天祭.
所以他們在路程當中.
基本上是不可能帶合格的祭生回來的.
為什麼呢?因為在聖殿的附近.
甚至在聖殿的院內.
他們需要買一些合格的祭生.
就是因為你帶一些.
譬如牽著一隻牛,牽著一隻羊回來.
牛走幾百米之後.
牠的腿可能跛了.
跛了之後.
李美琪說如果牠是殘疾的話.
你不可以獻給牠.
你不可以獻給上帝.
所以你帶的祭生一定要是最好的.
一定要是沒有殘疾的.
所以衍生了什麼呢?.
衍生了在聖殿的裡面.

$^{201}$甚至在院內.
他們有一個買賣的場地.
就是有一些合格,沒有殘疾的牛和羊.
第二就是.
當牛羊經過地方的時候.
可能會沾染一些不潔.
可能牠達到死屍.
或者達到其他東西的時候.
牠就不能獻給上帝.
所以在聖殿裡面.
有這樣的買賣是很正常的.
甚至當他們奉獻的時候.
他們要給錢.
連聖殿的稅都要用.
這個銀幣.
指定的銀幣才可以順利奉獻.
不然他們根本就不可以做到.
大家有見過這樣的寺廟吧.
是一些民間宗教的寺廟.
基本上所有民間宗教的寺廟.
都有買賣在裡面進行.
添香油,買元寶蠟燭,買香,解籤.
各樣各樣.
他們全部都是金錢的交易.
而在聖殿裡面.
有牛羊買賣本身就很正常.
為什麼這裡覺得正常的東西.
要弄得這麼奇怪呢.
這裡我們要連剛才說到負的殿.
我們就會知道為什麼耶穌要趕走牛羊.
在對觀福音裡面.
就是馬太,馬可,路加.
在潔淨聖殿的場景裡面.
耶穌是直接引用.
我的殿必為禱告的殿.
你們導使他成為了賊窩.
這個大家都很熟悉.
但讓福音是有點不同的.
讓福音是他的門徒說話並不是耶穌.
是他們引用說.

$^{241}$我為你的殿心裡焦急如同火燒.
這裡負的殿在同一段經文出現.
是回應著門徒的引用.
負的殿剛剛說到.
在整本聖經基本上沒有出現過.
只有這裡才出現.
如果我們用餘文去找.
你會看到有其他兩處地方.
一處是路加福音.
一處是使徒行傳.
但他也不是在說負的殿.
他在說某個人的家.
某個人的家是父親的家.
而不是在說這個殿的意思.
但放在一起.
這裡這麼獨有的意思是.
耶穌明顯將牛羊趕走的原因.
是指向耶和華的地方.
是不可以成為一個買賣的地方.
也很明顯.
但他不是想說我們不做金錢的交易.
不是在說聖殿不可以成為一個.
商業交易的地方.
耶穌其實是在反對買賣背後的動機.
這裡有一點tricky.
他不是在說你們不要做買賣這麼簡單.
是做買賣本身的行為背後的動機.
他在反對買賣背後的動機.
為什麼呢?要說一下.
其實在耶穌年代.
神殿是一個貪污和權力鬥爭的地方.
大家如果有見過阿波羅神殿.
只有幾個柱子的神殿.
曾經是一個很輝煌很厲害的神殿.
在希臘羅馬的時候.
基本上很多人都會在裡面工作.
裡面有很多東西.
譬如你去借錢.
你今天去借錢在哪裡?銀行.
或者有人打電話給你.

$^{281}$問你要不要借錢.
你就去借錢.
當時的人借錢在哪裡?就去神殿.
為什麼呢?因為神殿有很多祭物.
那些人去捐錢.
捐完錢之後就留下很多錢在裡面.
放這些錢不知用來做什麼的時候.
神殿就有一個借貸的功能.
不如我借些錢給你.
你還些利息給我.
下一次你就還給我.
而且會有很多不同的金錢交易.
譬如一些神諭服務.
譬如有些朝聖者去到阿波羅神殿.
他就問.
究竟我人生之後會不會順利呢?.
他們希望有些預言.
有些指引.
當他求告的時候.
他們就需要什麼?奉獻.
他們需要給錢給祭司.
這些神諭很多時候都被政客,城邦和羅馬官員利用.
成為他們統治的理由.
所以不少的羅馬神殿都成為了.
一些稅吏,軍隊,官員走資的地方.
因為不單只是一個敬拜的地方.
它是一個很複雜的金融體系.
那些官員會將戰利品.
將收到的錢直接放到神殿裡面.
就說你幫我看著.
還可以做大耳窿.
還可以收帳.
還可以不告訴中央我們有這麼多錢.
更甚的是.
有些官員甚至利用捐獻的份額.
去到將祭司的職位隨意買賣.
以至他們可以壟斷整個金融體系.
壟斷整個神殿裡面的貪污運作.
形成一套以神為命.
但實際是賺錢的系統.

$^{321}$這套神殿的貪污系統其實已經維持了幾百年.
當很多皇帝打仗打到沒錢的時候.
他們會怎樣?.
他們就打神殿的主意.
他就會跟神殿說我們沒錢.
不如你給我點錢.
所以猶太聖殿的聖殿稅.
其中一個原因就是這樣.
他知道聖殿裡面有錢.
所以他就會去打這個稅.
所以在這幾百年.
神殿的制度發展出一套很複雜的金融體系.
方便人去獻祭.
同一時間也方便人去貪污.
有這樣的情況.
令到整個體系都是為了官員和祭司去鞏固自己的權利.
大家猜一下.
這個文化會不會在猶太聖殿裡面出現呢?.
是會的.
在猶太聖殿裡面是有出現過的.
當年一個猶太家族叫瑪加比王朝.
他們在聖殿裡面的交易系統.
他們將這個交易系統成為了一個政治的籌碼.
買賣祭司的職份.
原本是屬於上帝的部分.
他們都拿走了.
他當成是自己的.
祭司是我的.
我喜歡賣給誰就賣給誰.
當年猶太聖殿也是這樣.
所以耶穌在釀方裡面.
在聖殿裡面趕走牛羊.
禁止買賣.
不是單純叫聖殿的人不要做交易.
不要有金錢的瓜葛.
這裡是說人對聖殿的玷污.
耶穌針對的是買賣背後人奉獻的動機.
我再說一次.
耶穌針對的不是買賣.
是買賣背後人奉獻的動機.

$^{361}$這個奉獻的系統本身是怎樣的.
方便人去到上帝面前.
去到敬拜.
去到奉獻.
但是神聖的地方被人的貪念.
被人的權力所玷污.
人的心玷污了整個聖殿系統的奉獻.
的敬拜.
所以為什麼二章二十四二十五字是這樣強調的.
耶穌知道萬人.
知道人的心所存.
他在這裡特意寫得很誇張.
寫得很矛盾.
好像很空白.
知道人生所存些什麼.
正正就是知道人的心.
知道人的心裡面.
有些對權力,對金錢,對關係.
對一些很多不同的東西的依戀.
剛剛說到.
廣州牛羊.
是《羊羊福音》裡面獨有的記載.
牛羊在舊約利未記是很重要的.
它是一個獻祭裡面的核心.
人透過獻上無彈疾的牛羊.
在平安祭獻給耶和華.
脂油和腰子是留給神的.
然後其他就留給獻祭和祭司一同分享.
為了表達什麼.
表達感恩.
就好像語字一樣.
人做這件事.
為了說.
不只是說我想吃牛,吃羊這麼簡單.
是說.
所有東西都是來自於上帝.
現在我將一部分歸還給上帝.
告訴自己.
是呀.
因為有上帝我們才有我自己.

$^{401}$有幾個這樣的情況.
但是.
當23節再提到耳語節的時候.
我們發現.
有些東西屬於上帝的人卻奪走了.
今天時間不是很多.
本身想跟大家看看這個詩篇69篇.
因為門徒引用的是.
我為你的電心裡焦急如同火燒.
這個詩篇的處境.
正正就是大衛面對軍事上的逼迫.
親人的出賣.
很想為上帝建立神殿.
所以.
他想上帝棄穴.
這篇詩篇你看到31節那裡.
他有提到聖.
就是獻祭.
是說什麼呢.
這就讓耶和華喜悅.
聖恥獻牛.
是不是.
獻祭牛.
獻牛.
這裡暗示的是.
人的心歸向上帝.
比獻祭牛羊更加討上帝的喜悅.
是獻祭本身想人的心歸向上帝.
但無奈獻祭這個動作.
失去了本身的意義.
令到牛羊這個獻祭.
都沒有了意義.
所以是人的心.
玷污了整個聖殿的獻祭系統.
令到獻祭失敗.
令到獻牛.
令到平安祭的意義都失去了.
我們華人經常說禮尚往來.
親人朋友之間.
經常會送禮.

$^{441}$去作一個日常的部分.
你去別人家裡.
你通常都不會兩梳子這樣戴.
你怎樣都要戴一些東西.
起碼戴個橙.
新年拜年的時候.
你都會送禮物.
大家都會禮尚往來.
好來好去.
其實這個文化很好.
能夠表現對方的情誼.
但如果我們將這種心態.
放在我們跟上帝的關係.
就本末倒置了.
上帝祝福我.
我就捐多一點.
上帝祝福我.
我就奉獻多一點.
上帝祝福我.
我就給多一點時間給上帝.
如果反過來.
如果上帝沒有加多一點祝福給你.
是不是你就奉獻少一點.
是不是你就給少一點時間給上帝.
禮尚往來其實某程度都是一個交易.
不知道大家有沒有厭倦這種.
一種交易式的形式主義.
就是中秋節.
以前大家都會聚在一起吃飯.
到後來大家都很忙.
不如不要吃飯.
大家交換月餅.
到早年就開始說.
月餅我都不敢拿給你.
不如給你餅卡.
之後就說.
大家都不要給餅卡.
因為大家都沒有時間去拿.
是不是.
大家都開始厭倦這種交易的文化.

$^{481}$因為以往月餅,中秋節.
這些東西的意義是什麼.
大家可以相聚,團聚.
後來變成一種文化.
我給你月餅就等於我重視你.
但後來這個意義已經不再存在.
然後大家就叫大家.
不如大家都不要送東西給大家.
大家說一聲中秋節快樂就可以.
是不是有一種這樣的情況.
但到頭來這種意義失去了之後.
我們上教會也是這樣.
當我們上教會的時候.
我們生命中的東西.
我們的財產,時間,精神.
我們希望獻上的時候.
我們在想的東西.
是禮尚往來.
是一種交易式的奉獻.
還是我們在想.
上帝所有東西都是你的.
我今天將那小部分給你.
即是在提升我們自己.
所有東西都是你的.
我們只是做我們的本份.
耶穌這裡再一次提醒我們.
在這個奉獻的系統裡面.
現實的確是很複雜的.
在耶穌的年代已經是這樣.
神殿已經貪污了幾百年.
然後聖殿也是這樣.
照樣玩轉彎告訴別人.
我們方便別人奉獻.
但最後是貪污,權鬥.
是很多的玷污.
怎樣也好.
我們的心我們分得清楚.
就好像那些牛羊.
明明是奉獻系統裡面.
很重要的一部分.

$^{521}$代表著人的心歸向上帝.
你分一些給上帝.
分一些給祭司.
自己留一些.
提醒我們的人生.
我們的資產,我們的時間.
我們的才能,我們的關係.
不是全權由我們擁有的.
也不是純粹由我們自己決定的.
當我們有意識分一些給上帝的時候.
是在提醒自己.
一切一切都是來自於上帝.
有些東西我們是不可以僭越的.
所以當我們說.
我們要將我們心裡面的牛羊趕走.
不是趕走那些屬羊,屬牛的弟兄姊妹.
不是趕走我們不吃牛扒.
我們不吃羊扒.
是趕走我們對於金錢.
對於權力,對於關係.
對於我們的才能的一些眷戀.
我們覺得那些是我們的.
我們要放下這些依戀.
放下這些眷戀.
最近我和太太Alice.
討論奉獻金錢的時候.
我們兩個人有很不同的方法.
我們其中一個方法就是.
每個月設定好注重轉賬.
銀行到時候就會捐錢給教會.
另一個方法就是手動的.
按錢出來.
或者在銀行的應用程式按出來.
我們就問.
究竟哪一個方法更加體現到奉獻的本意.
大家覺得呢?.
按錢是吧?.
按錢好像好一點.
好像誠心一點.
其實無論你自動轉賬還是手動轉賬都好.

$^{561}$同一時間都可以變得很形式主義.
或者同一時間都可以體現到奉獻的心.
不是形式的問題.
而是什麼?.
心態的問題.
如果你只是看著金額.
除了心痛之外就是肉赤.
沒有其他意義的話.
我們就只是形式上滿足了這個十一奉獻的動作.
其實沒有提醒到我們自己.
究竟我們奉獻是為了什麼.
這樣我們就做著做著.
就會好像月餅券一樣.
月餅券就是不要給.
沒有了本身的意義.
因為本身這個行動的意義是為了提醒我們自己.
所有所有一切一切都是上帝.
就不是那十分一才是上帝.
這一部分就在提醒我們.
我們只是盡我們本分.
而不是上帝對我好,祝福我.
所以我就隨便捐多一點.
隨便給多一點時間.
否則我們就會變成做一些民間宗教.
宰一隻雞.
上帝祝福我.
我宰一隻雞還一下神.
有這樣的情況.
換言之,時間也是.
有時候我們事奉到某個點.
我們就會開始問.
我給上帝這個時間.
其實是不是值得的呢?.
這個問題的背後我們其實正在問一件事.
我給上帝的時間.
我可不可以換一些我覺得值得的東西回來.
或者我覺得這個時間是自己的.
這個時間是我自己喜歡做什麼就做什麼.
是因為我們手握住.
我們心握住這個時間.

$^{601}$以為這個生命是屬於我們.
在經文裡面很有趣.
剛剛說.
耶穌說拆這個殿,然後三天就棄了.
然後猶太人就串耶穌說.
這個殿要建了四十六年.
你三天就重建.
這裡提到四十六年.
其實幾乎是當時古代人一生的年日.
但耶穌就說他三天就可以重建.
當時是怎樣呢?.
這個羅馬皇帝為了討好猶太人.
他就說不如我幫你擴建這個聖殿.
順便收一些稅.
幫我擴建這個聖殿.
出戰出力,用了一生的時間.
做了這個聖殿出來.
當然耶穌說這件事.
他三天就重建,就是在說他死後復活.
大家明白這件事吧.
他重建是把這個聖殿的意義.
放在上帝那裡.
那是怎樣呢?.
就是說猶太人用了這麼寶貴的時間.
用了四十六年一生的時間去建聖殿.
但最後怎樣呢?.
心卻遠離了神.
所以耶穌很心痛.
門徒就說耶穌為你的殿心裡焦急.
是因為聖殿本身的功能.
是要把人帶到神的面前.
明白自己的一切一切.
所有所有都是來自於上帝.
耶穌說我重建聖殿用三天.
就是把他用自己的生命.
把這個聖殿的意義.
把人帶到神的面前.
重新放在這裡.
就是因為本身聖殿的功能.
被人的心玷污了.

$^{641}$人對金錢,對時間,對財幹的眷戀.
人的心玷污了整個聖殿的系統.
令到它失去意義.
所以今天我們說趕走牛羊.
是可以是任何東西的.
牛羊可以是我們對時間,對金錢,對財幹的認識.
我們忘記了我們每一口呼吸.
我們身邊的關係,我們的財幹.
所有所有都是來自於上帝.
所以今天我們拆解了.
約翰福音耶穌結成聖殿的記載.
我猜大家都應該會懂得回答這三個問題.
耶穌結成聖殿之後.
正正知道人的心充滿著對金錢,對權力,對時間,對財幹各樣東西的眷戀.
令到本身聖殿甚至預節的意義都失去了.
耶穌重新將牛羊趕走.
這個象徵式的意義.
是在提醒我們對各樣東西的眷戀要放下.
對各樣東西的眷戀要放下.
要將這些東西重新給上帝.
我們說UV十週年的年替是Next.
當我們在求問上帝.
下一步要怎樣做.
要做些什麼.
我想我們首先要問的是.
我們做的下一步Next.
是為了自己還是為了上帝.
這是第一個問題我們要問.
學習放下我們自己對身邊的眷戀.
重新審視,重新思考.
做下一步有什麼意義.
我們一起祈禱.
天父多謝你.
多謝你幫助我們趕走我們心裡的牛羊.
是我們對於各樣東西的眷戀.
可以是錢,可以是權力,可以是時間.
可以是財幹,可以是關係.
這些各樣東西在我們生命裡出現.
幫助我們.
幫助我們重新將所有的擁有權.

$^{681}$交回到你的手裡.
我們心裡的祈禱,我們的掙扎.
甚至是一些惡念.
你每一天都知道.
因為你知道我們生存的是什麼.
求你幫助我們重新出發.
透過耶穌的命求.
在新的一年裡祝福大家.
趕走心裡的牛羊.
祝福大家.
\newpage



\section{約翰福音 4:43-54}
\label{sec:bMtCjPvaMwM}
\textbf{你是耶穌的粉絲嗎 (約翰福音4\_43-54) - 譚偉明宣教士}
\newline
\newline
連結: \href{https://youtube.com/watch?v=bMtCjPvaMwM}{\texttt{https://youtube.com/watch?v=bMtCjPvaMwM}} ~~~~ 語音日期: 2026-05-07
\newline
\newline
\hyperref[sec:YUWGLt45Wz8]{\small{< < < PREV SERMON < < <}}
~
\hyperref[sec:index_chronic]{\small{[返順時目]}}
~
\hyperref[sec:index_scriptual]{\small{[返順卷目]}}
~
\hyperref[sec:3zbQTQ6NDLE]{\small{> > > NEXT SERMON > > >}}
\newline
\newline
約翰福音 4:43-54
\newline
\begin{longtable}{cl}
\hline
\hline
章節 & 經文 (和合本修訂版)\\
\hline
4:43 & \begin{tabularx}{0.7\textwidth}{X} 過了那兩天，耶穌離開那地方，往加利利去。 \end{tabularx} \\ \\ \relax
4:44 & \begin{tabularx}{0.7\textwidth}{X} 因為耶穌自己作過見證說：「先知在自己的家鄉是沒有人尊敬的。」 \end{tabularx} \\ \\ \relax
4:45 & \begin{tabularx}{0.7\textwidth}{X} 到了加利利，加利利人都歡迎他，因為他們也上耶路撒冷去過節，曾經看過他在節期間所做的一切事。 \end{tabularx} \\ \\ \relax
4:46 & \begin{tabularx}{0.7\textwidth}{X} 耶穌又到了加利利的迦拿，就是他從前變水為酒的地方。有一個大臣，他的兒子在迦百農病了。 \end{tabularx} \\ \\ \relax
4:47 & \begin{tabularx}{0.7\textwidth}{X} 他聽見耶穌從猶太到了加利利，就來見他，求他下去醫治他的兒子，因為他兒子快要死了。 \end{tabularx} \\ \\ \relax
4:48 & \begin{tabularx}{0.7\textwidth}{X} 耶穌對他說：「若不看見神蹟奇事，你們總是不信。」 \end{tabularx} \\ \\ \relax
4:49 & \begin{tabularx}{0.7\textwidth}{X} 那大臣對他說：「先生，求你趁著我的孩子還沒有死就下去吧。」 \end{tabularx} \\ \\ \relax
4:50 & \begin{tabularx}{0.7\textwidth}{X} 耶穌對他說：「回去吧，你的兒子會活！」那人信耶穌所說的話，就回去了。 \end{tabularx} \\ \\ \relax
4:51 & \begin{tabularx}{0.7\textwidth}{X} 正下去的時候，他的僕人迎面而來，說他的兒子活了。 \end{tabularx} \\ \\ \relax
4:52 & \begin{tabularx}{0.7\textwidth}{X} 他就問甚麼時候見好的。他們對他說：「昨天下午一點鐘熱就退了。」 \end{tabularx} \\ \\ \relax
4:53 & \begin{tabularx}{0.7\textwidth}{X} 他就知道這正是耶穌對他說「你的兒子會活」的時候；他自己和全家就都信了。 \end{tabularx} \\ \\ \relax
4:54 & \begin{tabularx}{0.7\textwidth}{X} 這是耶穌從猶太回到加利利後所行的第二個神蹟。 \end{tabularx} \\ \\
[1ex]
\hline
\hline
\end{longtable}
$^{1}$弟兄姊妹平安.
今日真係好高興可以返到黎自己屋企.
呢個唔係即係隨便ge 教會.
呢個係我自己屋企.
因為我認識好多弟兄姊妹.
好多弟兄姊妹響到長大ge .
同埋好多導師都響到教導我ge .
咁樣啊.
真係好開心今日有機會響到同大家.
分享我地宣教ge 經歷.
但係我都同你地講先.
你大部份都認識我係邊個啦.
我響到都好怦怦跳ge 我心跳.
因為我今日要講廣東話.
同埋你地知道我廣東話ge .
文化ge 水平唔高.
但係先要打個底.
因為有D 弟兄姊妹唔認識我.
咁我D 中文ge 水平.
基本上係交TBB學費ge je .
所以有D 詞比較地道D 呢.
即係英文比較GhettoD 呢.
請大家見諒啊.
咁樣.
已經去jor 三年半啦.
咁我都估唔到咁快已經一任響到台灣.
咁我響台灣真係有D 習慣ge .
最唔習慣返到黎第一樣ye 就係.
點解咁會動ge 私家歌.
咁樣啊.
咁我響台灣宣教呢.
基本上呢我係響.
冇好人響高雄.
我係響直唐ge 響屏東市.
一個好細ge 城市響.
高雄ge 東邊.
大概四十五分鐘ge 車程.
咁教會你見到.
右邊gwo 張呢就係.
大概教會前面ge 門口啦.

$^{41}$閂jor 閘而家.
但係左邊gwo 張就係條街係咁樣.
咁但係我地教會大概係三十至四十人啦.
咁呢度一家人響到影幅相比你知道大家.
大部分呢弟兄姊妹呢都係大過我ge .
咁我響到其中ge 職位呢.
就係做一個年青牧師.
咁我其實我發覺.
我去到三十幾四十歲.
而家先去做年青牧師.
其中一樣ye 呢.
我以前都做jor 十三年架啦.
即係gwo D 團體導師.
我發現我唯一一樣ye 做唔到而家.
就係我唔可以同佢玩全程.
我真係好gui 而家比較.
但係另外做年青牧師呢個職位呢.
就係我同太太會教好多下令營東令營.
咁你見到左邊gwo 一張呢.
就係我第一次教下令營ge 時間.
係gwo 時教一個細路都冇ge .
咁我地點樣做呢.
我同太太同埋一個留學生.
係去一個小學gwo 度排擔仗.
我大大聲聲咁樣同佢地.
大家用英文對話.
比佢地見到我地D 標準ge 英文.
咁第一個暑假有大概四個學生.
加埋我兩個仔女有六個.
咁第二年就D 家長見到呢度幾好.
咁唔需要我再去派傳單.
咁家長幫我推銷就容易好多.
除jor 做年青牧師.
我地建堂呢.
唔係話我淨係做小學ge .
淨係細路仔ge 工作.
基本上乜都需要做.
我無論係招待定係去serve午餐.
呢D 全部都需要做.
其實呢D 咁ge 策略.

$^{81}$同埋點樣做法呢.
我完全都係響UV學返黎ge .
出面要人入黎.
定係有招待定點樣有人serve.
有cafe.
我呢D 教弟兄姊妹都係響到學返黎.
都要帶成人ge bible study同埋小組.
最後呢我有機會主持一個婚禮.
有個弟兄叫我主持婚禮.
咁我都未做過其實第一次.
我見過好多次.
跟住我問AI通常應該點樣.
一個baptistge 婚禮.
但係跟住同我D 經驗.
我話大概係咁樣.
應該冇問題.
咁就係咁樣.
你真係發覺到.
去到gwo 度.
你冇人幫你.
你咩都需要做.
但係其中一樣ye .
弟兄姊妹成日都.
我去到gwo 度呢.
其實弟兄姊妹好尊敬我地.
因為我地係全教師ge 身份.
同埋係牧師ge 身份.
但係我都成日同佢講一句.
我就係咁樣同佢講.
我現係我ge 職位係一個牧師.
但係我跟你將黎ge 關係係乜ye .
我跟你係弟兄姊妹.
所以你唔應該稱呼我係Steve牧師.
只有稱呼我Steve.
還是佢唔知道我叫阿明.
你地呢其實.
我都叫阿明仔得架喇其實.
大家返黎係咪先.
我D 導師係度.
好似Jimbo咁樣睇住我大ge .

$^{121}$係咪先.
唔好意思呀Jimbo.
爆jor 你地年齡先.
所以.
咁樣今日ge 主題.
我地講兩種ge 信心.
第一種呢就係.
好淺ge 信心.
我地睇一睇.
黎到話.
係呢個背景就係基本上.
耶穌呢就係從耶路撒冷.
預節ge 緣故返到黎.
佢地自己家鄉.
係加利利gwo 度.
但係佢特登係傾過呢個撒瑪利.
因為佢要讓撒瑪利亞.
撒瑪利亞.
撒瑪利亞人知道.
佢自己ge 身份係表面.
咁佢最後呢就返到黎呢度呢.
但係好奇怪.
第44章跟住你見到佢話.
呢個約翰索.
記得耶穌講.
一個先知.
響佢自己ge 家係冇尊敬ge .
但係45呢跟住佢返到黎.
嘩D 加利利人.
就好開心咁樣歡迎佢.
但我地見到.
佢歡迎佢呢.
因為耶穌做jor 好多ge .
好多ye 響.
耶路撒冷當時.
但係我見到佢做jor D 乜ye 呢.
如果你返返去2章.
你唔需要返去我講.
你返去2章23節呢.
其實呢耶穌做jor 好多神蹟.

$^{161}$因為呢D 神蹟好多人相信佢.
咁你發現.
44同45點解約翰係咁樣寫.
係咪有D 衝突呢.
好似玩ye 喎.
但係佢講呢一句呢.
你如果你睇返其下D 福音書呢.
通常講呢一句ge 背後呢就係.
通常D 人呢係拒絕耶穌.
唔相信佢呢.
佢就講呢句.
佢回應呢句.
但係係度冇喎.
呢度呢.
佢係冇任何ye 都未返到黎屋企.
約翰已經講呢句.
點解呢.
其實佢想同你講.
佢想做一個介紹比你知道.
呢成個故事呢.
其實呢D 加利利人呢.
佢地歡迎耶穌呢.
佢地信心真係好膚淺ge .
佢地好表面去歡迎佢.
咁我地去到46節呢.
咁而家耶穌真係返到黎.
自己屋企.
即係呢個區呢度加利利.
特登返到去迦蘭呢個地方.
咁你見到佢.
約翰就講就話.
佢呢個地方係將水變走ge 地方.
第一間神聖.
咁即係話佢想同我地比一個印象就話.
其實好多人知道呢.
耶穌響到做jor 呢件事.
但係同一時間.
好多人都聽聞佢響耶路撒冷做ge 事.
即係話.
啊耶穌呢而家係比較出名啦而家開始.

$^{201}$跟住黎緊47節你見到.
46節下半節gwo 度呢你見到.
開始有個衝突就開始發生就係.
呢個官有個黎求耶穌去醫治佢仔.
因為佢仔又係死啦.
但係你發現.
呢個官其實唔可以命令耶穌叫佢醫治咩.
要唔係去殺.
要唔係去逼迫佢咩.
其實佢可以喎我覺得.
但係我覺得去到一個地步呢.
我覺得佢已經用jor 好多唔同ge 方法呢.
佢先會肯去求耶穌一個猶太人去.
去醫治佢仔.
佢發現.
即係你發現一個人有財有勢.
最後gwo 個終極點呢都係死亡.
都係打唔過死亡呢個最後ge 敵人.
唯有我地所知道耶穌.
復活先係可以能夠打贏呢個死亡.
響力史記載只有耶穌基督.
係有.
耶穌基督係唯一ge 人呢係一個空ge 墳墓.
咁響48節呢你見到耶穌去回應喇喎.
佢回應就話你地班人.
你係相信我因為你見到呢D 神蹟奇妙咋嘛.
咁佢話你地呢個英文個U呢.
響希伯來文.
唔係唔係.
響希臘文呢就話.
唔係單數係雙數.
佢唔係單單同呢個官講ye 架.
同其他所有ge 觀眾講.
話你地係相信.
你地真係相信我係因為你見到呢D 神蹟奇妙.
所以耶穌就話.
你地成日都係好鍾意睇呢D 好漂亮ge ye .
好犀利ge ye .
好似D 魔術咁ge ye .
佢地見到耶穌係一個咩呀.

$^{241}$係一個好偉大ge Miracle Worker.
同埋而家呢就係出jor 名喇.
返返黎家鄉喇.
咁耶穌帶jor 好大個光榮返黎家鄉喎.
因為點解呢.
如果你知道呢猶太.
猶大地南部耶路撒冷京城咁啦.
佢點睇加利利北邊D 人呢.
基本上睇佢地就係鄉下佬.
無文化.
咁耶穌而家就係變jor 一個.
好出名ge gwo D 名人咁呀.
咁呢一班加利利人呢.
就好似D 粉絲咁樣.
係追住呢個名人好似潮流咁樣.
所以佢好開心.
咁你諗諗而家最潮流係乜ye 呀.
其實我都唔係好清楚呀.
其實我見過D 後生ge 一輩呢.
好似我D 仔女gwo D 呢.
最鍾意係gwo D 咩Taylor Swift呀gwo D .
我諗好似有D 大人都去佢地Concertge .
所以我覺得最潮流係gwo D 但係我又唔好識喎.
咁你追粉絲咁.
但係我唔知道.
但係我只係知道.
我以前追係乜ye 呢.
就係我地公牛隊23號啦係咪呀.
Michael Jordan啦.
其實呢.
Michael Jordan呢.
佢gwo 邊出世呀.
佢大學讀邊度呢.
佢ge 阿爸阿媽呀.
係咩名呀.
佢D 仔女係邊個ge 呢.
個名我都講得出.
但係我同佢係完全冇關係ge .
我認識佢但係佢唔認識我呀.
就係呢樣ye .

$^{281}$咁第二樣ye 就係.
我想再返返去44字gwo 度呢.
就係講話一個先知係冇尊敬ge .
響指ge 家鄉.
咁其實個個都覺得好難ge 其實.
你一個人.
一班人睇住耶穌大.
甚至乎去認識耶穌ge .
係好難去服從響個權柄之下ge .
其實我覺得.
就好似我以前呢.
我所有ge 電腦gwo D .
我D 智慧呢我都係教我阿叔ge .
其實佢一D 電腦都唔識ge 之前.
但係返到黎呢.
佢而家話有新ge software.
要去教我點用.
我成日覺得有D 怪怪地.
冇可能呢D ye 咁都發生.
所以好難我覺得.
但係你發現.
之前D 撒瑪利亞人呢.
冇呢個問題架喎.
因為佢冇呢個包袱係度喎.
佢就基本上呢.
耶穌係冇咩做咩神蹟呢.
佢聽一件事.
耶穌話.
佢已經就信架喇.
佢信係乜ye 呀.
響42節佢話.
佢信耶穌係救世主.
呢D 神蹟奇妙呢.
其實呢.
係幫助我地去.
指證俾我地知道耶穌係邊個.
係彌賽亞.
係救主.
係神的兒子.
讓我地將我地gwo 個信心.

$^{321}$可以擺係我地gwo 度.
呢D 係幫助我地ge .
信心ge 其實.
唔係去滿足我地ge 人.
你見到我地ge 關係.
呢度冇真正ge 關係.
同埋呢.
佢地唔需要做任何commitmentge .
呢D 人呢.
佢gwo D .
佢ge 信心真係好淺.
但係呢.
我可以同你講.
大部份ge 人.
我響台灣.
我遇到呢.
都係咁樣其實.
人.
黎到教會呢.
基本上呢.
唔係話想認識耶穌.
我自己係一樣.
佢地唔見到耶穌係主.
佢地知.
大部份響台灣就係.
見到耶穌係一個好偉大ge 醫生咁樣.
好犀利ge 醫生呀.
定係佢係一個神燈咁樣.
佢地唔想同耶穌要與關係.
去到呢個地步呢.
佢地我有一次去到.
arm arm 去到三個幾個月前呢.
有個.
唔係姊妹.
有個女士同我講.
我地一齊有個.
查經組.
佢話.
我呢排.
我好多年呢都未訓覺.

$^{361}$其實睡眠睡得唔好.
咁我話.
如果你幫我可以.
祈禱.
比我訓得好呢.
我就相信佢啦.
我心諗.
我都滴哂汗.
大佬你想我點呀.
我唔知我點呀大佬.
但係我都幫你祈禱.
我心諗我照幫你祈禱啦.
咁佢第二日真係真係.
佢話我訓得好好.
咁係咪真係信耶穌呀我心諗.
咁我話.
巧合黎ge je .
等我試多幾次.
咁佢又玩ye 係咪.
同我講先.
真係.
但係我睇一睇其他D .
example就係.
有D 人.
好直接啦係咪.
黎教會.
係為jor 餐飯囉.
有D 真係.
有D 為jor 關係囉.
有D .
係黎學英文囉.
單單黎學英文.
如果你要講耶穌呀.
就唔好囉.
咁我要同佢講.
我地係真係學英文ge .
但係.
我生命呢.
係耶穌.
因為耶穌我先有今日.

$^{401}$我唔可以唔講耶穌喎.
如果你要我分享.
真係好難搞ge .
有時有D 人.
另外呢.
D 家長呢.
真係好有.
goal oriented.
真係好有目標ge .
黎教會賴被日.
淨係玩呀.
定係.
學點樣做.
即係學品格呀.
點做好人呀.
嘩呢sai 時間ge 真係.
佢梗係冇咁同我講啦.
但係佢ge 意思係咁啦.
我一至六呀.
我要去補英文.
禮拜六仲要夜晚去補課呀.
雖然禮拜六唔需要返學.
咁我地就係點算呀.
呀好好.
我地有英文ge .
我地廿分鐘有半個鐘頭英文.
我地學對話.
跟住呢.
仲要學樂器添.
好彩我老婆識少少.
另外就係.
有機會去呀.
好唔係次次D .
唔係成日都有得做ge .
就係去拆神台啦.
咁樣呀.
有位弟兄咁呀.
信jor 耶穌咁呀.
叫我地拆神台.
咁我話.

$^{441}$嘩我都未做過呢D ye ge .
嘩.
我係呀.
咁樣但係發現.
台灣人.
如果信jor 耶穌.
你拆jor 佢ge 家.
佢ge 家裡面ge 神台呢.
唔代表拆jor 佢心ge 神台ge .
有D .
台灣人.
D 弟兄姊妹.
將佢以前gwo 個偶像.
搬過黎基督教.
耶穌變jor 佢偶像.
你見到.
所有ge .
所有ge ye 呢.
其實呢.
ge 目標黎教會.
唔係因為耶穌.
係咪有其他ge .
好ge ye ge 事情.
佢可以得到D 乜ye 先黎教會.
佢ge 生命係冇改變ge .
佢諗到.
佢只係想.
上帝祝福佢.
但係佢唔想同上帝有個關係.
咁弟兄姊妹想問.
咁我地呢.
咁我地係咪咁樣做呀.
我地係咪黎教會.
係咪有乜ye 目的呀.
係咪呢個教會ge .
因為個兒童事工好呀.
定係個講道好呀.
當然講道好啦.
我響Markgwo 度學到好多ye .
但係我唔想你地做粉絲呀.

$^{481}$弟兄姊妹.
我想你做門徒.
所以我地睇一睇下一點.
點做.
真正ge 信心係點樣呀.
49至54.
咁但係唔似呢個官呢.
就唔似gwo D 粉絲咁.
佢再去要求耶穌.
你快D 過黎啦.
我仔就黎死架啦.
加耶穌呀點呀.
你唔.
你ge 仔會.
即係唔會死.
咁你諗一諗.
你會點做呢.
你相唔相信呀.
我去邊個餐館呀.
我地個個人呀.
都睇過D Google.
Google review 啦係咪呀.
睇下好唔好食啦.
睇下其實D 人點講呀.
佢就係咁樣講.
gwo 時有冇Google review.
你點知呀.
但係呢.
你見到.
其實響好多地方呢.
下一節呢就係咁樣架啦.
gwo 個人去呼求耶穌好多次.
耶穌有做過一神蹟.
去醫治好gwo 個人.
做D 奇妙ge ye .
但係今次你見到冇喎.
我覺得冇.
其中一樣原因就係.
佢唔想去滿足呢班粉絲.
第二呢.

$^{521}$佢就係.
需要gwo 呢個人去揀.
我冇比奇蹟你見唔.
你見唔到奇蹟.
你會唔會相信我個話語.
所以我真係奇蹟呢.
唔係我地基督徒.
我地可以去build.
apply我地ge 根基黎.
其實係上帝ge 話語.
咁但係你最後你見到.
呢個官員係冇得揀.
佢最後呢.
係需要接受耶穌所講.
而家好表面講ge 話.
跟住離開jor .
就係至少你見到佢肯去接受.
佢離開你見到佢ge .
佢呢個信心係慢慢.
即係又增加ge .
跟住去五十一節去五十二節.
就係佢返屋企啦.
我覺得佢返到屋企.
佢差唔多返到屋企ge 時.
佢gwo 個僕人就見到佢.
佢話.
你ge 仔而家冇事喎.
咁正常黎講.
我地覺得好開心.
其實你搞掂gwo 單ye 呢.
你唔會搵返gwo 個人去幫助.
你唔會搵返呢個醫生ge 其實.
但係佢呢就係冇喎.
佢就諗起.
五十二節佢就諗起.
o下 佢幾時好返架.
我通常我唔會諗話.
我仔幾時發燒好架.
我唔會去諗返幾時gwo D ye .
但係佢有諗過.

$^{561}$因為我覺得佢記得.
耶穌講ge ye .
跟住佢去對一對.
原來真係咁arm 耶穌講gwo 時候.
佢就真係以此好.
你發覺我地基督教gwo 個信心.
唔係淨係得個feel.
我覺得耶穌今日好好.
其實都要有根據.
要考 要用邏輯去諗.
我地gwo 個faith.
但係佢見到.
佢跟住點樣做.
佢真係將呢個信擺係耶穌gwo 度.
佢唔係淨係信耶穌.
係一個好偉大ge miracle worker.
佢係信耶穌係呢個彌賽雅呢個救世主.
佢唔單止信.
佢將呢個福音傳比佢ge 家人.
佢全家人相信.
但你諗一諗 弟兄姊妹.
佢話佢途中係一點鐘前一日.
佢先知佢個仔係好返喎.
你發現加拿去到加伯隆呢.
大概係20個mile.
要用一日ge 時間係gwo 度.
佢過夜有機會係過夜一晚.
你諗一諗.
你過一晚你ge 心情係點樣呀.
其實我覺得好多弟兄姊妹.
我地大家都經歷過gwo 個心情.
上帝ge 應許得唔得架.
即係有冇咩料架呢個上帝.
我覺得gwo 個掙扎呢.
真係好難過gwo 一晚.
我覺得呢個官.
但係我地個個都經歷過呢樣ye .
但係係呢樣ye 上帝係要我地ge struggle.
係呢方面去依靠佢ge 話.
之後佢引證jor 呢.

$^{601}$真係推動jor 佢ge 信心.
我想講ge 就係.
我地真正ge 信心唔係一朝一夕ge .
唔係一次ge 搞掂就搞掂架啦.
呢個救恩就係.
但係呢個信心每一次.
我地同上帝ge 關係呢.
就係佢每一次每一刻佢每個季節呢.
每一年呢佢都將呢D trials.
呢D 試煉呢擺係我地ge 身上.
讓我地去去認識佢多D .
去信教多D .
因為你諗一諗.
我地幾時最唔常唔去搵上帝ge 時候.
係我地最安穩ge 時候.
有時都唔記得佢係邊個啦.
所以我好多人係咪.
大部分人好少數人呢.
會將會真係想咁慢慢將呢個生命咁樣.
係完全交比上帝.
完全相信佢.
呢個係我地每日.
我地每日ge 影子每日ge 基督徒ge 功課.
我都係我ge 功課.
但係我係台灣雖然話大部分人係咁.
但係我都好開心呢.
我見到一班好少數ge 人.
會肯去將呢D 生命慢慢.
即係佢地ge 信心慢慢去生長.
其中一個就係我ge 大茶經.
我gwo 時都唔知點架.
我gwo 時大茶經.
都係大概五六個弟兄姊妹肯黎.
呢D 全部都大過我ge .
咁我黎到話查咩好呀.
咁佢地話你查咩都得架啦Steve.
咁就讓我去福音啦.
咁話我查jor 五六次架啦我聽過.
話OK咁你查jor 五六次咁我再試一次啦.
咁睇一睇啦.

$^{641}$跟住呢我就就係係度所學架啦.
我係神學院呀同埋係教會學ge .
咁教導下唔係淨係咁睇喎.
同佢講聖經ge 背景呀.
咁點樣gwo D 神蹟ge .
即係點樣systematically咁去睇呀.
咁佢話我未學過呢D ye 架喎咁樣.
原發現佢明白地理ge 背景呢.
佢幫助佢佢ge 信心.
佢知道呢度ge 即係約翰講咩ye .
佢幫助佢信心都好多.
跟住佢成日都上黎啦.
去到gwo 地步呢我其實同佢講.
我2月22號呢係我最後一日呀.
響呢個教會.
咁我諗住係1月最後gwo 個禮拜六呢就算啦停啦.
因為2月中呢係新年呀嘛.
咁佢地話Steve唔得喎可以繼續開.
我話點解呀你地新年唔係好多ye 搞架咩.
佢就預備搞乜鬼呀而家.
做jor 出去食飯你仲煮架咩.
搞咁多ye 呀.
跟住佢話真係佢想繼續去裁英.
我做宣教士其中我最開心ge 一份.
即係我開心ge 時候就係呢樣ye .
D 製英之輩想明白上黎ge 話語.
第二呢呢兩個呢係我第一批ge .
gwo D 下令營呢gwo D 學生.
gwo 次有四個學生加埋我細佬六個呢.
佢係其中兩個係呢度.
叫阿Sam同阿Tina.
呢個仔呢特別係呢個仔呢.
一黎到呢即刻我桌底呢有桌布ge .
佢即刻鑽jor 入去喇.
佢唔肯見人ge 好驚你點樣ge .
咁我地成日都哄佢哄佢.
成個東令營呢都係咁樣架啦.
佢阿媽呢之後都繼續去擺佢.
係我地ge 賴皮日ge 團契gwo 度.
過jor 一年咁樣就開始同我friendjor 啦.

$^{681}$咁樣就慢慢有機會.
即係同佢講聖經故仔啦.
但係一難ge 我覺得.
因為呢大部分D 台灣.
係屏東市ge 家庭呢.
我所認識呢全部係單身家庭黎架.
佢地已經唔認識阿爸架啦.
佢以為阿爸已經同阿媽.
已經分開jor 好耐架啦佢地.
所以我好開心.
另外一樣ye 就係可以.
係gwo 度做一個男性ge 榜樣.
我覺得上帝ge 呼召.
係我生命中ge 其中一樣ye 就係呢樣ye .
咁樣我地你見到佢gwo 幅相.
其實佢都對我都okge 而家.
所以呢個有個progression.
係需要時間架.
係需要比佢知道.
我真係真係關心佢.
呢兩個呢就係.
黎到我地教會大概兩年前.
咁佢黎佢係台北做工ge .
但係burn outjor 啦.
因為台北好辛苦.
咁佢返返黎屏東市.
我發現好多人好多都係鄉下都係屏東市.
咁返黎咁搵一間教會.
搵到我地ge 雙語教會.
咁佢黎到.
咁你知D 通常大過我啦.
大部分人都同我講.
同我咁大ge 我D peers呢.
我通常mark住佢地架.
咁我都要做好我ge 工作架.
咁其中一樣ye 咁我地一齊分享啦.
我成日都同佢地出去食飯呀.
同佢地傾偈.
咁其中一樣ye 就係話.
佢話係以前ge 教會.

$^{721}$即係成長gwo 時好多傷害.
咁樣我大家傾下D deology啦.
其中一樣ye 其實.
台灣其實唔係冇教會.
其實有時我地gwo D deology.
未必係最arm ge .
譬如佢覺得好多台灣人覺得.
我浸jor 泥先係叫相信耶穌.
我未浸泥我就唔係信耶穌架.
咁其中一樣ye 就係話.
佢呢個女仔就係好話.
特別女仔就係好話.
我覺得佢係一個perfectionist黎架.
佢覺得乜ye 呢都係要.
我要將所有100\%全部比晒上帝呀.
我唔論係我工作上.
我教會上我都要比佢最好ge .
我都要做得最好.
但係我都要同佢傾一傾.
你呢樣ye 係真係好呀.
但係你係嫌現實有時.
你唔要自己抖下其實.
你班友所以你成日之前burnout囉.
我唔係話你應該係係工作定乜ye 做得唔好.
係應該做得好架.
但係你點樣去balance你ge life.
有時係做.
你係俾到你ge 生命.
淨係gwo 一日呢.
gwo 個禮拜四gwo 一日俾到上帝呢.
就係gwo 一日啦係咪先.
我唔係唔想discourage弟兄姊妹唔去事奉而家.
所以就話你自己去搵個balance.
去搵個平衡喇你ge 生命.
但係見佢慢慢去肯去做呢樣ye .
同埋去真係好認真咁樣去寫低佢gwo D 筆記啦.
我地上主學gwo 時.
同埋最後呢去幫助我地翻譯囉.
佢係我其中一個翻譯.
因為之前我可以用英文講到.

$^{761}$咁最後呢一個家庭呢係一個好貧窮ge 家庭黎架.
咁佢又因為係夏令營ge 緣故.
第二年呢佢就聽聞我地有夏令營.
佢就黎參加我地夏令營.
咁樣呢個好貧窮.
因為呢佢地全家包括ge 丈夫呢.
都好多呢D 先遷性ge 病.
所以lor 晒D 錢去去睇病.
咁佢地其實呢都唔係話未去過教會.
佢地都去過唔同ge 教會.
但係呢佢有一次好.
即係同我講.
我地成個月已經食緊粥架喇其實係屋企.
佢女兒話可唔可以唔食粥呀.
我話發生咩事呀咁大鑊.
係台灣餓唔死喎其實.
你買個便當都好平架咋喎.
但係發現佢D 錢已經去晒gwo 度架嘛.
佢跟住呢有一日同我講.
你地教會雪櫃有咩呀.
你lor 埋黎去食呀.
gwo D 雞塊gwo D chicken nugget呢.
你lor 去啦如果你需要過.
食剩gwo D sung 呢即係冇料.
即係料我姐係夾過ge .
可唔可以lor 唔係咁大鑊呀嘛.
跟住佢可以慢慢.
跟住認識jor 佢佢慢慢同我分享.
我以前有教會gwo 個傳道人同我地講.
我地全家有呢D 病呢.
因為你唔相信耶穌.
嘩大佬咁呢D 咩傳道人黎架我心諗.
跟住都同我講.
佢D 朋友都話.
你過黎做呢個教會gwo 個會員啦.
你睇下做會員有幾多ye lor .
又有米lor 又有gwo 樣ye lor .
咁佢話如果我為jor 去教會做會員可以lor 到ye .
咁呢個信係咩黎架.
咁我真係過jor 一晚呢.

$^{801}$我真係好傷心.
我打邊爐呢我都打到我都火滾呀真係.
我真係知道gwo 個傳道人係邊個.
我真係想拿架撐去打鬼死佢.
我真係但係我唔知呀嘛.
真係好氣我個人.
我記得同佢出黎講唔得.
我地一定要去而家去ge 超市食完飯.
我地成卡ye 我地full up佢.
所有油鹽米醋呀乜ye 都擺晒整滿佢.
我地即刻lor 去佢屋企.
我同佢講呢個係耶穌俾ge 教會俾ge .
我地冇條件ge .
所以佢慢慢你會明白話真係幾難去拆改呢D .
即係呢D 咁ge 傳統ge 呢D 信仰.
跟住呢.
有一刻.
佢慢慢.
佢會問我福音係咩黎架.
有時發覺弟兄姊妹係需要時間同埋關心.
所以弟兄姊妹.
我唔知道你點睇耶穌.
佢係一個好名人啦定係係我地ge 主.
咁我都希望呢耶穌係我地ge 主啦.
咁我希望.
我地既既然如果耶穌係我地ge 主.
我希望我地ge 生命同我福音係相稱架.
弟兄姊妹呀.
我地ge 基督徒ge 生命呢.
其實.
同宣教呢係走唔甩架.
你係基督徒呢你係你係全教士黎架啦.
我同你ge 崗位唔同就係我響台灣你響呢度.
你無論響你ge 工作你地家人gwo 度呢你邊度你都係你係個全教士.
你係代表緊耶穌架.
但係因為你記唔記得大使命呀.
個個人都記得大使命.
你要去張曼文做我門徒.
但係你唔可以講話我其實好渣ge ye .
我乜都唔識架其實.

$^{841}$我同你講你唔可以咁樣講ye 架.
但係因為你唔可以渣得過我.
我去到大學gwo 時我先arm 學識廣東話講都講唔正.
gwo 陣時同我地而家咪係咁樣.
跟住你話連個夢都冇呀.
做乜ye 呀咁大ge 仔去jor 個college係三年.
最後耶穌話你做個老師去做個好榜樣.
男性ge 榜樣所以我教小學.
之後我覺得個生活好好.
我話真係呢一世就好啦.
跟住耶穌話你去你做村教囉.
我話唔係呀嘛D ye 唔arm 我.
我點做牧師je gwo D ye .
跟住話咁村教係要讀神學院架喎.
我就諗緊嘩如果我班友知道我背景佢點收我一.
佢真係收jor 我喎我發現.
跟住呢我又畢到業喎.
我都唔知點畢到業.
跟住個村教機構又收jor 我去台灣.
而家係度慢慢考緊國語.
我真係我真係我唔記得阿堅咁講.
喂爛癱黎架大佬.
大佬我都唔知我自己有咩恩賜真係.
即係音樂又唔得.
電腦又唔係最叻我都唔知有咩恩賜其實而家我都未知有咩恩賜但我唯一ye 就係我就話上帝我係度你需要ge 話你可以用.
但你諗一諗如果你冇恩賜你上咩驚ge .
如果上帝係真係呼召你去一個崗位.
佢會比呢個資源你去做好佢囉係咪呀.
你諗諗你係你公司你個老細叫你做ye 佢唔比部電腦你咁樣大佬你自己lor 佢返黎架.
或者有D 係都唔頂我唔知.
得你明白我講咩呀.
但係我想你記得就係.
你記得大使命但係唔好唔記得個大誡命係咩呀好重要.
記唔記得大誡命係咩呀.
你要愛主利益神.
另外個個係咩呀.
你要愛人如己係咩呀.
好多時呢我地大使命呢我唔記得jor .
gwo 個人呢係變jor 一個福音ge 推銷對象呀.
你同佢傳福音呢.

$^{881}$佢唔信呢拜拜.
我我冇關係架啦.
但係人人呢係可以見到你係咪真心對佢架.
你想同人傳福音呢.
你真係要真心對gwo 個人.
其中一樣ye 我響我響孫九宮長學識.
佢雖然呢我地真係好難先可以去呢個關係好耐先可以呀.
建立好但係佢真係見到你真係真心對佢唔係因為你要推銷一個福音比佢呀.
呢個好大ge 分別呀定時妹.
咁而家將來呢.
我地將來會做D 咩ye 呢.
我地而家做係響屏東南部ge .
響高雄南部住緊但係我會去去一個新ge 教會.
去嘉義嘉義縣一個小ge 地方叫普治一個好細ge 城市gwo 度.
好鄉下仲鄉下而家.
佢大概芝加哥十倍咁多啦.
芝加哥大約十倍啦.
鄉下黎架真係好ge 鄉下.
呢度呢對面呢就係而家我住館住gwo 度呢佢樓上呢就係一個好似有個.
細路仔團體咁我地就而家我黎緊去jor 已經去jor 一個月架啦.
咁同佢地玩啦結了關係呢一班細路仔呢.
都好慘ge 其實好似你諗一諗留守兒童係點樣響大陸.
佢地基本上係咁樣.
佢有D 父母唔響到去jor 第二度做工你諗一諗台灣ge 島好實ge .
其實佢地唔係好遠但係唔得閒去理佢地.
還是佢有D 父母呢就係係個廟gwo 度做ge .
咁你諗一諗佢父母係個廟gwo 度做ge 將來係點樣咪就係個廟gwo 度做啦.
佢就話個父母就話呢度係咪教會黎架係教會我唔要我D 細路仔黎囉.
咁我地唔係係度講耶穌ge 我地呢間呢度囉我真係唔係講耶穌我真係同佢好似團契咁樣去愛佢地建立關係咋.
有時講一個教會即係一個故事一個聖經故事.
好開心我都可以繼續同佢地一齊有機會打波同佢教佢點打籃球.
另外第二個地方就係好近ge 係呢個我地叫恩典中心.
因為我地呢個教會係要恩典中心因為頭先我都講啦如果係教會我地唔黎架啦.
跟住有D D 人入到黎跟住崇拜呀一齊去傾偈呀食飯跟住D 人就開始發覺呢唔係教會教會唔係咁做架嘛.
係呀其實教會就係咁架咋囉佢佢未入過黎佢驚呀嘛唔知道呀嘛.
咁我地教會其實都係30至40人一個細教會.
我呢其實下一任我返去呢就去呢個地方其實呢係學點樣去做一個好ge 牧者.
同一時間去練好我國語囉我黎緊呢間教會呢係完全講中文ge 即係國語.
我臨走之前我嘗試講我自己ge 見證係度大概用jor 13分鐘咁樣傳國語ge .
我就話我之前去台灣一個國語都唔識就係咁夾硬黎架約我.

$^{921}$但我覺得上帝係係我呢個生命中係有一個好大ge mission.
我覺得幾難呢我都會做好佢ge 我盡量.
因為我做基督徒呢我同我上帝ge 關係呢唔係話靠feel架.
係睇o下我服唔服從am I obedient or not呢個就係我其中一個信心ge 增長其中一個manifestation.
咁將來呢因為我真係將來我主管行講Steve Shirley你不如去直堂啦.
我冇諗過喎而家專教士我都以為做唔到而家直堂ge 一天.
我就話好我有一天我將來我目標係台灣我就要去直堂.
但我唔知佢點樣但我覺得因為真係又有一個異象係度.
我想再比無論係神學上定係我地ge 品格上.
我比佢地知道天家係真正gwo 個弟兄姊妹呢點樣係應該係點樣相處ge .
多謝弟兄姊妹.
謝謝大家.
\newpage



\section{以弗所書 3:1-6}
\label{sec:3zbQTQ6NDLE}
\textbf{你是誰?當福音重新定義你的人生 (以弗所書3\_1-6)Who Are You? When the Gospel redefines your life ( Eph 3\_1-6)}
\newline
\newline
連結: \href{https://youtube.com/watch?v=3zbQTQ6NDLE}{\texttt{https://youtube.com/watch?v=3zbQTQ6NDLE}} ~~~~ 語音日期: 2026-05-07
\newline
\newline
\hyperref[sec:bMtCjPvaMwM]{\small{< < < PREV SERMON < < <}}
~
\hyperref[sec:index_chronic]{\small{[返順時目]}}
~
\hyperref[sec:index_scriptual]{\small{[返順卷目]}}
~
\hyperref[sec:6UvylJcofPU]{\small{> > > NEXT SERMON > > >}}
\newline
\newline
以弗所書 3:1-6
\newline
\begin{longtable}{cl}
\hline
\hline
章節 & 經文 (和合本修訂版)\\
\hline
3:1 & \begin{tabularx}{0.7\textwidth}{X} 因此，我—保羅為你們外邦人作了基督耶穌囚徒的，替你們祈禱。 \end{tabularx} \\ \\ \relax
3:2 & \begin{tabularx}{0.7\textwidth}{X} 想必你們曾聽見神賜恩給我，把關切你們的職分託付我， \end{tabularx} \\ \\ \relax
3:3 & \begin{tabularx}{0.7\textwidth}{X} 用啟示讓我知道福音的奧祕，正如我以前略略寫過的。 \end{tabularx} \\ \\ \relax
3:4 & \begin{tabularx}{0.7\textwidth}{X} 你們讀了，就會知道我深深了解基督的奧祕； \end{tabularx} \\ \\ \relax
3:5 & \begin{tabularx}{0.7\textwidth}{X} 這奧祕在以前的世代沒有讓人知道，像如今藉著聖靈向他的聖使徒和先知啟示一樣， \end{tabularx} \\ \\ \relax
3:6 & \begin{tabularx}{0.7\textwidth}{X} 就是外邦人在基督耶穌裡，藉著福音，得以同為後嗣，同為一體，同為蒙應許的人。 \end{tabularx} \\ \\
[1ex]
\hline
\hline
\end{longtable}
$^{1}$[SPEAKING CHINESE].
So we'll continue to preach on the letter to the Ephesians,.
chapter 3..
[SPEAKING CHINESE].
So the title is Who Are You and How the Gospel Defines You..
[SPEAKING CHINESE].
Let's read this scripture together..
Let's do Chinese first..
[SPEAKING CHINESE].
Well, let's read together Ephesians 3..
"For this reason, I, Paul, a prisoner of Christ Jesus,.
on behalf of you Gentiles, assuming.
that you have heard of the stewardship of God's grace.
that was given to me for you, how the mystery was made known.
to me by revelation, as I have written briefly.".
[SPEAKING CHINESE].
Verse 4 and following, "When you read this,.
you can perceive my insight into the mystery of Christ,.
which was made known to the sons of men in other generations,.
as it has now been revealed to His holy apostles and prophets.
by the Spirit..
This mystery is that the Gentiles are fellow heirs,.
members of the same body, partakers.
of the promise in Christ Jesus through the gospel.".
[SPEAKING CHINESE].
The first question is, in which world you are living?.
[SPEAKING CHINESE].
I'm not asking you whether you believe in Jesus.
or whether you come to church..
[SPEAKING CHINESE].
It's the question of which world are you living in right now?.
[SPEAKING CHINESE].
People may say, well, I'm living in this world in Chicago..
[SPEAKING CHINESE].
But the scripture tells us there are two worlds simultaneously..
[SPEAKING CHINESE].
So the old world is centered on ourselves..
[SPEAKING CHINESE].
And then the new world's centered on Christ..
[SPEAKING CHINESE].

$^{41}$Or put it another way, why do you live the way you do?.
And why do you do this?.
[SPEAKING CHINESE].
Or think about every decision that you make..
And there's a question for you is, why do I do those things?.
[SPEAKING CHINESE].
And we have to ask this question often..
[SPEAKING CHINESE].
And our answer to this question is that we want to live better..
[SPEAKING CHINESE].
And then everything's like revolving around ourselves..
And the faith has become a footnote..
[SPEAKING CHINESE].
And then we want to turn to the scripture today.
and to challenge ourselves, like Paul reminds us..
Are we living in the gospel-centered world,.
or are we stuck in the old world centered on ourselves?.
[SPEAKING CHINESE].
OK, so first two verses..
[SPEAKING CHINESE].
So the main point is, who controls your reality?.
[SPEAKING CHINESE].
And Paul described himself and introduced himself..
He says, not an apostle..
He introduced himself as a prisoner of Christ Jesus..
[SPEAKING CHINESE].
And we talk about Paul..
When he wrote this letter, he was not.
on the shore of Mediterranean resort,.
and he was in prison for the gospel..
[SPEAKING CHINESE].
And if you think you were Paul, and how would you.
introduce yourself when you write the letter?.
[SPEAKING CHINESE].
And you may introduce yourself, Paul,.
I am the apostle chosen by God..
[SPEAKING CHINESE].
And Paul says, I am prisoner of Christ Jesus..
[SPEAKING CHINESE].
And it's very strange that he introduced himself like this..

$^{81}$No problem..
I got it..
Oops, sorry..
OK, here you go..
[SPEAKING CHINESE].
So to the Romans, [SPEAKING CHINESE].
Paul is a totally defeated person..
[SPEAKING CHINESE].
From the perspective of the world, from the Romans..
[SPEAKING CHINESE].
And if Paul sees himself in the Roman way,.
and he would pity on himself..
[SPEAKING CHINESE].
But from another perspective, from the gospel perspective..
[SPEAKING CHINESE].
And he was not a defeat, and he was chosen by God..
And part of the plan of God is for him.
to be in prison for the gospel..
[SPEAKING CHINESE].
And not only he was in prison for the gospel,.
he was chosen by God..
Even in full suffering, he's advancing the gospel..
[SPEAKING CHINESE].
So as you can see, it's the same circumstances in the world,.
but it's a completely, totally different perspectives and interpretations..
[SPEAKING CHINESE].
And if we come from a very narrow perspective,.
and we see Paul's life of suffering in prison,.
and we'll think that, OK, maybe suffering is always good..
[SPEAKING CHINESE].
So the real question is not whether he is in prison now or he's suffering..
The main question is, who is narrating your life?.
Who is defining your life and circumstances?.
[SPEAKING CHINESE].
So that's the main point here..
Who controls your reality?.
How you see reality?.
[SPEAKING CHINESE].
And he said it's not the circumstances that define him,.
but it's the Christ, Jesus, the Lord defines him..

$^{121}$[SPEAKING CHINESE].
So sometimes we need to think about our own lives..
We define ourselves by our jobs, our career, our income,.
our reputation, our achievement..
[SPEAKING CHINESE].
And the question of the gospel confronts us.
is that, why are we doing what we do?.
And who do we belong to?.
[SPEAKING CHINESE].
So who is defining our life?.
[SPEAKING CHINESE].
It's not the world, it's Christ to define you..
[SPEAKING CHINESE].
So the first point is, we have to understand Christ, Jesus.
defines our identity..
[SPEAKING CHINESE].
So then we can come to see Paul's suffering.
from this perspective..
[SPEAKING CHINESE].
So this extends to the providence of God in our lives,.
not just the good things, but also the difficult things,.
like Paul's imprisonment and suffering..
[SPEAKING CHINESE].
And you can see the imprisonment and suffering of Paul.
is a bad thing, but in the providence of God,.
it can be a good thing..
[SPEAKING CHINESE].
So the providence of God is not just permitting something bad.
happen in our lives, but according to Paul,.
in Romans 8, that everything God works together.
for the goodness of his kingdom..
[SPEAKING CHINESE].
So that's how Paul interprets his suffering..
[SPEAKING CHINESE].
So when life is not going well for us,.
how do we define our life?.
And according to the folk religion,.
or in Christian's worldview, it's.
the luck and the fortune or circumstances that defines us..
[SPEAKING CHINESE].

$^{161}$But God says in the Bible, Christ controls everything..
[SPEAKING CHINESE].
So things happening in our lives are not because of randomness,.
but there is a plan in him..
[SPEAKING CHINESE].
And so Paul is talking about not only the mission that he has,.
but he received a stewardship from God..
[SPEAKING CHINESE].
So Paul is suffering, but he is suffering not forgetting.
that God is giving him the stewardship and this mission.
to deliver the gospel..
[SPEAKING CHINESE].
So this stewardship is entrusting something important,.
just like if I go away for a long trip.
and I entrust the message to Paul, then Paul.
has to relay this message, important message, to Wendy..
[SPEAKING CHINESE].
So Paul is not saying that his job is his own job,.
but it's God's entrusting this assignment to him..
[SPEAKING CHINESE].
So you think about your own life is not your own ownership.
at all, it's God giving you the stewardship for your life..
[SPEAKING CHINESE].
And if you think about our lives,.
we are not the owner of our own possessions,.
but God is entrusting to us the career, the children,.
including our wealth, everything,.
and we are just stewards and managers..
[SPEAKING CHINESE].
And for a good steward, he understands.
he's being entrusted something important from the owner..
[SPEAKING CHINESE].
You see, so often we come to God and we just.
want to receive blessing and good things from him..
[SPEAKING CHINESE].
But Paul is emphasizing, I have received the stewardship.
from God himself..
[SPEAKING CHINESE].
And to verses 3 to 5, Paul is talking about the gospel.
he is preaching..

$^{201}$[SPEAKING CHINESE].
He said God revealed the mystery of the gospel to him..
[SPEAKING CHINESE].
So the mystery here is talking about the gospel..
It was hidden before, but now because of Christ,.
it was revealed..
[SPEAKING CHINESE].
And how did Paul figure out the gospel?.
He did not do it from his own personal research.
or from another person, but from direct revelation of God..
[SPEAKING CHINESE].
And the term revelation, what does it mean?.
It means-- it does not mean something mysterious,.
but it means something that was covered before,.
but now it's uncovered..
[SPEAKING CHINESE].
So just like the cage back there was originally.
covered by a curtain, and if you take away the curtain,.
then now it's uncovered, and you can see what's inside..
[SPEAKING CHINESE].
And the same way for the gospel..
The gospel is not discovered by men..
It is uncovered by God..
It's revealed what's inside, and now we know the gospel..
[SPEAKING CHINESE].
And God revealed the gospel to Paul..
[SPEAKING CHINESE].
And of course, God also revealed to Paul and other apostles,.
but it is the direct revelation from God..
[SPEAKING CHINESE].
So this revelation is already complete in Christ Jesus..
So Paul has already received this revelation..
Now, the later generations of Christians,.
we may have illumination from the Spirit,.
but we don't have new revelation..
[SPEAKING CHINESE].
So the gospel is not invented by men..
It is directly revealed by God..
[SPEAKING CHINESE].
So humanity can figure out the existence of God.

$^{241}$from the general revelation from the creation,.
but sometimes it is the will, and it's not very correct..
[SPEAKING CHINESE].
And you remember, recently, NASA has sent four astronauts.
to fly around the moon, and one came back,.
and there were four of them, and one was a Christian,.
and he came back, and he was baptized..
[SPEAKING CHINESE].
And he was seeing, when he traveled around the moon.
and saw the space, and he said it's very impossible for us.
to deny the existence of creator..
[SPEAKING CHINESE].
But nature can point you to the existence of the creator..
[SPEAKING CHINESE].
But if you really need to understand or know God,.
you have to depend on not your own consciousness,.
but not through creation, but from the direct revelation.
from God..
[SPEAKING CHINESE].
And the special revelation is Christ Jesus himself,.
and also the Word of God we find in the Bible..
[SPEAKING CHINESE].
So you think about it, without God directly revealing himself.
in the gospel, in Christ Jesus, there's.
no way man can figure out a crucified man is.
the center of the gospel..
[SPEAKING CHINESE].
And you think about a God who dies on the cross.
and who fails, and this is totally foolish,.
according to human logic..
[SPEAKING CHINESE].
In the Qing dynasty, there was a statement called Li Hongzhang,.
and he once said that there is no way that he could understand.
why the Western people would believe in such a defeated.
Christ Jesus..
[SPEAKING CHINESE].
So as you can see, the salvation plan of God.
is against human logic..
[SPEAKING CHINESE].
So the essence of the gospel is that God is coming to find you.

$^{281}$on the cross of Jesus..
[SPEAKING CHINESE].
And he said this revelation of gospel.
is already a shadow in the Old Testament..
[SPEAKING CHINESE].
And you think about the Old Testament in the Abraham.
promise, God has already told him.
you will be blessed with many, many children and many, many.
people..
And then in the prophecy of Isaiah,.
he was already prophesized there was some savior or servant.
will come to save us..
[SPEAKING CHINESE].
And you think about the Old Testament's shadows, and it's just the sketch.
of the gospel, like driving into Chicago..
Far away, you can see the skyline of a city,.
but it's very blurry on the horizon..
Only when you arrive downtown and you come to downtown,.
you will see it clearly..
[SPEAKING CHINESE].
So from the Old Testament, you can see these shadows..
You can see the shadow of the gospel..
But until Christ Jesus came, the mystery is revealed..
And he came, and he died, and he rose again..
And the spirit is poured out to his people..
[SPEAKING CHINESE].
So you see, what is revealed is Christ Jesus..
And what is the gospel?.
The gospel is Christ Jesus himself..
[SPEAKING CHINESE].
So God gave us Jesus Christ, his son, and that is the gospel..
[SPEAKING CHINESE].
So the third point is that what is this mystery of the gospel?.
[SPEAKING CHINESE].
So Paul is saying the mystery is that now, in Christ Jesus,.
even Gentiles can become part of the people of God..
[SPEAKING CHINESE].
And to be the share heir of the promises.
is a new identity..
[SPEAKING CHINESE].

$^{321}$So we are not sojourner outside of God's kingdom,.
but we are being included..
[SPEAKING CHINESE].
And it's not because our performance,.
we can earn the salvation of God..
But it's all because of grace of God..
[SPEAKING CHINESE].
And the Bible tells us we Gentiles are the fellow heirs..
And our identity has changed, and we become one people..
And we can share the blessing of God..
[SPEAKING CHINESE].
So this identity informs us that even now,.
our identity is the kingdom of the people to come..
So we live now, but we belong to the kingdom to come..
[SPEAKING CHINESE].
And it's all past tense description for our identity..
[SPEAKING CHINESE].
And you think about it, OK, we have not ascended to heaven.
like Christ, but how is that possible?.
[SPEAKING CHINESE].
But Paul is saying this new identity informs you.
you have already been united with Christ,.
and he has rose and ascended to heaven..
So are we..
[SPEAKING CHINESE].
So Paul is saying in the Ephesians.
that the gospel has formed a new people, a new group,.
and a new humanity..
[SPEAKING CHINESE].
So in conclusion, for the three points,.
the gospel is not making us more successful or making us more,.
but to redefining us..
[SPEAKING CHINESE].
So the gospel is not simply just add on to our lives,.
but to reconfigure and redefine our own identity, who we are..
[SPEAKING CHINESE].
And because of the gospel, we are.
being brought to a new creation, a new world,.
and that is in united with Christ Jesus himself..
[SPEAKING CHINESE].

$^{361}$So the first point, you remember,.
God is bringing us to a new world..
And so who are we in Christ?.
[SPEAKING CHINESE].
So our life is not defined by our success or our failure,.
but we are defined solely on Christ Jesus himself..
[SPEAKING CHINESE].
So the world is telling you, you have to earn the identity..
You have to earn the success..
But the gospel informs us in a radically different way..
It says, in Christ Jesus, we belong to God..
[SPEAKING CHINESE].
So understand who we are..
And the second question is, who do we live for?.
[SPEAKING CHINESE].
And the gospel informs us we are not living for ourselves,.
but for Christ Jesus..
[SPEAKING CHINESE].
So you think about the life of Paul..
He was a super orthodox rabbi..
And he came to know Christ..
And his journey of life, the direction of life,.
has totally changed..
[SPEAKING CHINESE].
So in a worldly perspective, he's messed up his life..
But from God's perspective, he is directing.
to a new and right direction..
[SPEAKING CHINESE].
But that does not mean we abandon our ordinary life..
We don't work, or we don't go about managing our lives..
But the center has changed for our lives..
Our center has shifted because of Christ..
[SPEAKING CHINESE].
So the gospel informs us we have belief in Christ..
And we are being entrusted with the gospel..
[SPEAKING CHINESE].
Just like Paul has a baby and a child..
And he's being entrusted with a precious life..
And also with it, a lot of suffering and responsibilities..
[SPEAKING CHINESE].

$^{401}$And the point is, we have received.
a lot of different responsibilities and roles.
in our lives..
But the gospel has challenged us,.
confronted us, that we are being entrusted with something.
that is important..
And that is the center of the gospel..
[SPEAKING CHINESE].
So the first question is, who are we?.
We are in Christ..
And who are we living for?.
We are living for Christ..
And the third question is, how do you.
see the body of Christ, the community of faith?.
[SPEAKING CHINESE].
You speak English or Cantonese or Mandarin..
And you have your own group..
[SPEAKING CHINESE].
Or you're rich or I'm poor, so we're different..
[SPEAKING CHINESE].
But the gospel has formed a totally new humanity..
It's not because you're better or I'm poor.
or because you're stronger or weaker..
But we are new people in Christ..
[SPEAKING CHINESE].
So the city of man is ordered by division and separation..
But the city of God is united in Christ..
[SPEAKING CHINESE].
So brothers and sisters, we may call ourselves Christians..
But we may be still living for ourselves..
[SPEAKING CHINESE].
When we come to church, probably we.
are still living in the old world..
[SPEAKING CHINESE].
So Christ comes to our lives..
And asking ourselves, the core of our identity.
has to be shift from ourselves to Christ..
[SPEAKING CHINESE].
And we ask ourselves, at the bottom of our hearts,.
are we living for ourselves or living for Christ?.

$^{441}$[SPEAKING CHINESE].
Or whatever God has entrusted to us,.
and we are avoiding it or trying to escape or run away.
or being faithful..
[SPEAKING CHINESE].
And in the church, are we building up for the--.
building the unity?.
Or we are constantly comparing ourselves to other people?.
[SPEAKING CHINESE].
And let's pray together..
[SPEAKING CHINESE].
God, we want to give you thanks, because you gave us your word..
[SPEAKING CHINESE].
And we pray that, just like Paul, how he sees himself..
And we use the same perspective to see ourselves..
[SPEAKING CHINESE].
And why are we living?.
[SPEAKING CHINESE].
And we are not just living for ourselves, but living for you..
[SPEAKING CHINESE].
Living in the gospel..
[SPEAKING CHINESE].
And day to day, we live for your glory..
[SPEAKING CHINESE].
So we pray in Jesus Christ..
Amen..
Thank you, Paul..
Well, let's give him a hand..
\newpage



\section{}
\label{sec:6UvylJcofPU}
\textbf{Craftsman of Broken Pieces (Gen 12-50) - Rev. Mark Chan 碎片的工匠 (創 12-50) - 陳崇基牧師}
\newline
\newline
連結: \href{https://youtube.com/watch?v=6UvylJcofPU}{\texttt{https://youtube.com/watch?v=6UvylJcofPU}} ~~~~ 語音日期: 2026-05-10
\newline
\newline
\hyperref[sec:3zbQTQ6NDLE]{\small{< < < PREV SERMON < < <}}
~
\hyperref[sec:index_chronic]{\small{[返順時目]}}
~
\hyperref[sec:index_scriptual]{\small{[返順卷目]}}
~
\hyperref[sec:9gxDgFyKyk4]{\small{> > > NEXT SERMON > > >}}
\newline
\newline
$^{1}$Alright guys, good morning..
Well why don't we take a time to say hi to your neighbor..
There are a few neighbors around you and just say hi and good to see you..
Okay?.
Turn around..
Okay, so today we will continue our series on the mission of God..
Alright?.
I think I need some help with the PowerPoint..
Please..
Okay?.
Now, as we come to the church, this is our 10th anniversary..
This is our year number 10..
And we're going to celebrate and work together for the church..
Okay?.
Our series is on NEXT..
And we're going to challenge and remind you to go NEXT in your life..
So last time when you were here when I preached on this series, you should have picked up.
a jar like this..
Right?.
And inside is a bunch of sand..
And it reminds us the promises to Abraham from God..
And we are part of the children of Abraham actually in Christ..
And today if you're new, you can pick up one of these jars and you can bring it home..
If you already have one of these jars and you're going to put it in somewhere that will.
remind you every single day..
So throughout the whole year, we're going to add some items into the jar..
So today we're going to add another piece into the jar..
Let me begin with this picture..
It's called stained glass window..
Right?.
Stained glass window..
Now classical church buildings all have stained glass window and this tradition goes back.
to many, many years ago to the 12th century actually..
And when there was a profound insight into the symbol of God as light..
So whoever that, who has that insight is a theologian..
He replaced the entire wall of the church with stained glass window so that sunlight,.
natural light can come through and shine on us..
And each of the pieces, each of the pieces in the stained glass window is a craft by.
a craftsman and they set together and the result of is a masterpiece of a beautiful.
work and each piece is limited..

$^{41}$However, every piece come together, it will become a whole more beautiful picture..
And the story we're reading today from the book of Genesis, remind us this fact that.
God is the craftsman and he collect all the broken pieces, broken lives and to form a.
picture that will bless the whole world..
And last time we saw that God has made a promise to Abraham, the father of God's people in.
Genesis 12, 1 to 3..
Let me read the last phrase, which is most important is that God promises Abraham that.
he will make him to be a blessing and all peoples on earth will be blessed through you..
And the plan of God is universal as you can see is to bless every tribe and nation through.
one single guy, one single family..
And last time we focus on the personal weaknesses, failures of Abraham and today we're going.
to zoom out to see what kind of family is coming out of Abraham, the chosen family..
And as you can tell from this simple outline that after Abraham, there were three kinds.
of generations, Isaac and Jacob and Joseph and the brothers..
And these three family generations are all actually broken families..
So number one is Isaac, Rebekah, I mean, Isaac, Rebekah versus Esau and Jacob..
And as you know from the story, if you are kind of familiar with the Bible story, especially.
the patriarchs, Isaac loved Esau, right?.
The firstborn son, because he loved the game, right?.
He loved the game and Esau was kind of a hunter and the love of father is determined by appetite,.
right?.
By his stomach and Rebekah plot to deceive the blind Isaac at the end and Jacob, the.
younger son stole the blessing and Esau was enraged and Jacob fled to exile for 20 years.
and the family entrusted with blessing to the nations, they were shattered in the very.
first generation..
So as you can see from the story of the first generation, they are all broken..
The second generation starting with Jacob and his wives and concubine, the wives are.
Leah and Rachel and they compete for the favor of Jacob and each sending their own maidservants.
to bear more children because more children you get more inheritance..
And then the 12 tribes of Israel were born out of jealousy..
So the second generation of the patriarchs actually pretty broken, right?.
And then we come to the third generation which is Joseph and the brothers and Jacob repeat.
the same sin..
He favors one of his sons, his younger son Joseph and give him the clove of many colors..
So in Genesis 37, verse 4 it says when the brothers saw that his father, their father.
loved Joseph more than any of them and they hated him and could not speak a kind word.
to him..
And the Hebrew word behind the kind word is actually shalom..
As you can tell even modern Jewish people they say hi to each other, they don't say.
how are you, they say what?.

$^{81}$They say shalom, right?.
They say shalom..
So you see the brothers hated Joseph so much that they can't even greet him and they can't.
even say shalom to him..
So here the brother could not and would not say hi to him..
The hatred went so deep and they could not even manage to greet him..
And then all the rest of the story talk about the brothers, the elder brothers fought against.
one another to be the number one in the family..
You see Reuben is the first born, is the son of Leah and he after the death of Rachel he.
slept with Bilhah which is who is the maid servant of Rachel and that is kind of the.
act of political power grab and attack on Rachel's family and her sons and trying to.
strip the and at the end of the story Jacob stripped him from his birthright..
And then some of the story tells what happened to son number two and three Simeon and Levi..
They used the rules of circumcision to massacre the men in Shechem and Jacob cursed their.
violence..
And then son number four named Judah he proposed to sell Joseph to Egypt and he kind of dressed.
it up in a nicer way but you know in essence he's trying to illuminate a rival in the family..
And then the brothers, the other brothers flew Joseph into the pit and they what did.
they do?.
They sat down to eat, they sat down to have lunch and they sold him for twenty pieces.
of silver and they bought the coat to the father to deceive him that Joseph is dead..
So from Genesis we get the bits and pieces from the long story you can tell you know.
these are the stories of the founding fathers of the people of God..
And I always tell people you know the stories in the Bible are rated R..
The stories are actually not for children to read..
And we have to ask the question you know why did the author of Genesis focus on these kind.
of stories instead of telling the glorious moments he chose stories of deception, betrayal.
and murder even..
And we have to ask the question why?.
So often we read the Bible as a moral book and we ask the question okay he picked all.
this you know dark sides of the people and perhaps maybe a moral warning to the reader..
But you know if we turn Genesis into a morality textbook then we miss the whole point and.
the whole point asks different question..
It's not a moral example story..
It is who is God behind all of this..
So the Bible is not so much of a moral book..
It is a theological book..
And perhaps the second question is more important because it's theological question is asking.
who is God..
Now you ask yourself with all those founding fathers that are shattered and they're broken.

$^{121}$if you were God if you were God you kind of look at the disaster of the family and what.
would you do right?.
What would you do if you were God?.
Alright in the ancient near east actually the answer was pretty clear because you look.
at the pagan religion the gods would destroy they would destroy the people if they make.
troubles so violence is kind of the normal norm to deal with you know rebellions..
But as you can tell from the book of Genesis the God of the Bible is completely different..
He does not abandon the sinners he chooses the harder way the higher way which is take.
the broken pieces take the broken lives and set them into his picture so that his light.
will shine through to bless the nations..
So toward the end of the book of Genesis Joseph was being sold into Egypt for 13 years he.
was accused and he was being forgotten then Pharaoh had a dream and then Joseph was called.
to interpret it and he interpreted the dream as seven years of plenty and seven years of.
famine and Joseph eventually became the prime minister of Egypt..
As you can see from this verses from chapter 41 verse 57 what happened was that all the.
world the people were coming to Joseph coming to Egypt to buy grain because the famine was.
so severe..
And do you see that Genesis 12 we were in part one I talk about this is quietly becoming.
true not because the patriarchs or the founding fathers were better people they were not they.
were evil people but because of the promise of God was unbreakable..
So when the brothers reconciled with Joseph at the moment of recognition Joseph said something.
like that and Genesis 45 verse 8 it says Joseph said it was not you guys who sent me here.
in Egypt but God sent me to Egypt..
And then toward the end of the story Joseph said to the brothers you intended to harm.
me you guys you brothers intended to harm me but God intended it for good to save many.
many lives..
Now you have to understand that God did not whitewash the sins of the brothers the evil.
was real but human evils cannot stop God's purposes his mission it cannot stop it..
Despite the brokenness God was running an other intention simultaneously..
So you read the stories of the founding fathers of the people of God in summary they all fell.
short of the glory of God and the people of God in the world look exactly like this they.
were weak they were fragile and they were footless and they were unimpressive..
And you look at the church actually that is the real picture..
But I want to quote a very ancient document for almost 500 years old document it's called.
Belgic confession..
In this Belgic confession written in 1561 article 27 it says it says something about.
the church..
He said the church has existed from the very beginning of the world and will be to the.
end it will exist to the end for Christ is the eternal King who cannot be without subjects.
right as long as Christ exists church exists..

$^{161}$So this holy church is preserved by God against the fury of the whole world although for a.
while now right now for a while it may look very small and extinct in the eyes of men..
That's the patriarch family of Abraham, Isaac, Jacob and Joseph and that is so often sadly.
the picture of the church..
Now you come back to the original picture I show you the picture of stained glass no.
single piece of stained glass can complete the whole picture on its own and that is what.
the founding fathers the patriarchal family look like and that is what we look like most.
of the time but the patriarchal stories from Genesis leave us with this uncomfortable reminder..
And we read the stories from the Bible again it's rated out stories and we shake our heads.
we shake our heads because Jacob was evil guy and Judah sell his brothers to illuminate.
him and the biggest problem in the stories of patriarch is not simply that they wrecked.
the plan of God the biggest problem is that their attitude is I don't care attitude..
See Esau sold his birthright to Jacob for a bowl of stew because he doesn't care right.
he doesn't care about the plan of God not because he was evil it's just because he didn't.
care at all..
And the brothers of Joseph thrown him into the pit and they sat down to do what they.
sat down to eat right they sat down to eat that is the attitude of not caring at all.
and that is most frightening because it's a bystander mentality I just want to stand.
on the side and I don't participate I don't care and that is the fatal wound for the people.
of God fatal wound for the church..
Maybe we are doing this you never pray for the church and you never care right you have.
this attitude I don't care you know the plan of God the glory of God I don't care..
There are three invitations today as you leave as I told you you pick up a piece of puzzle.
and you are going to place it in the jar that you got from last time and this piece please.
you know be a reminder to you that you need every piece around you no piece can complete.
the picture you need other pieces to join together every piece has its own place including.
you yourself including you in your brokenness in your incompleteness right now..
And I think the first invitation is humility you know go back to the illustration you know.
go back to the stories of Genesis God you know that's a wrong attitude to say that God.
can only use me once I'm good enough because I'm not good enough God is not going to use.
me you know from the book of Genesis God is never waiting for the perfect person to use..
And God knows you are not perfect God knows you are broken but God can still use you it's.
not because you are better because God is great..
Now the second reminder is to wake up is kind of awakening..
So you ask yourself this you know kind of cosmic question have I been watching the word.
of God from the sidelines?.
You know the mission of the church the needs of the people around me have I ever said you.
know I have nothing to do with them..
And we need to repent from them..
And thirdly is the invitation to participate in the plan of God..

$^{201}$You know I am this broken piece am I willing to place myself this broken piece into the.
hands of the craftsman God himself not waiting on the sideline until I get good enough but.
right now with all the cracks and all the incompleteness I'm going to hand myself and.
trust myself to God because the whole picture God is putting together is not about one piece.
it's all the pieces together the whole picture of the grace of God and the power of God and.
glory of God and you are part small part of that big picture..
So at the end of the sermon I want to take some time some quiet time so that you can.
talk to God..
I'm sure you have been to church more than one time..
And I'm sure God has spoken to you..
You know you can ignore the voice of God..
You can silence the voice of God inside of you..
But God is still speaking to you..
And it's an invitation for you to come be humble and be awakened and to participate..
While the music is playing let's spend a couple of minutes and you talk to God yourself.
right here right now..
Let's all pray..
(Praying in Chinese) Let's all pray..
(Praying in Chinese) Dear Lord we are broken pieces..
We are cracked and we are fall and not enough on our own..
But you dear God you are the craftsman..
And you never discard what is broken..
Take us please..
And set us into your picture..
Let your light shine through..
And with open hands we are yours..
In Jesus name..
Amen..
What the invitation is for you..
\newpage



\section{}
\label{sec:9gxDgFyKyk4}
\textbf{碎片的工匠 (創12-50) - 陳崇基牧師}
\newline
\newline
連結: \href{https://youtube.com/watch?v=9gxDgFyKyk4}{\texttt{https://youtube.com/watch?v=9gxDgFyKyk4}} ~~~~ 語音日期: 2026-05-10
\newline
\newline
\hyperref[sec:6UvylJcofPU]{\small{< < < PREV SERMON < < <}}
~
\hyperref[sec:index_chronic]{\small{[返順時目]}}
~
\hyperref[sec:index_scriptual]{\small{[返順卷目]}}
~
\hyperref[sec:KUqZwRWtyog]{\small{> > > NEXT SERMON > > >}}
\newline
\newline
$^{1}$崇拜也是Mark Sir現在將時間交給他.
好 我們都.
為到.
開到咪未.
OK 好.
我們在講到之前.
我們都為到教會祈禱好嗎.
因為特映姐妹是特別需要去紀念的.
上一個星期日.
葉生,葉太.
他們就有一些家事.
所以我們完了崇拜.
我和另外一位同工.
我們就去到醫院探望葉生.
因為葉生突然間生命危殆.
所以我們上個主日.
下午就去陪葉太.
葉生的生活支援.
他需要去鋪排.
所以在整個下午.
我們陪著他的家人等等.
我們都是為葉太祈禱.
他們的安息星期是安排了這個星期三.
這個星期是中間的.
如果你有任何問題的話.
可以問我們其中一個同工好嗎.
但我們需要為我們身邊有需要的特映姐妹祈禱.
我們同心一致地祈禱.
我們一個屬靈的家.
雖然我們人多.
未必每一個我們都會很個人地去認識.
但在基督裡面.
我們仍然是屬於這個家庭.
主我們為葉生,葉太的家祈禱.
我們知道葉生已經是安息主懷.
失去地上的奴婦.
在天家裡面.
在上帝你自己的安慰當中.
在耶穌的同在裡面.
是很多無比的地方.

$^{41}$主我們為葉太,為祂的家人祈禱.
求你在傷心難過的時候給他們安慰.
也給這個星期三.
我們安息星期當中.
我們能夠見到上帝你自己的安慰和認許.
我們能夠活在你自己的保守.
和你自己的帶領之下.
雖然我們生活裡面.
人生裡面可能有高有低.
有迂迴曲折的時間.
但是我們知道你是我們的牧者.
仍然帶領我們.
我們教會進入第十週年.
我們求你幫助我們.
在這一年的當中.
我們一起在屬性生命上能夠成長.
在我們教會.
我們每一方面能夠成長.
我們榮耀你.
祈禱奉主名禱.
好,在十週年裡面.
我們有一個港獨的系列.
就是叫做Next.
我們這次是第二講的.
上一次我們講的是阿布拉罕的故事.
我們說整本聖經由《創世記》到《啟示錄》.
去介紹一位上帝是怎樣的上帝.
這位上帝是Guard in Mission.
是Missionary God.
是一個宣教中的上帝.
上帝沒有放棄過這個世界.
雖然這個世界是亂七八糟.
我們生命都亂七八糟.
但是上帝一直沒有放棄過我們.
祂的Mission就是去拯救我們.
拯救世人和拯救整個的地.
我們這次開始的時候.
用這張圖片來開始.
這個是在巴黎聖母院的天主教教堂裡面拍的.
大家有沒有去過.

$^{81}$或者見過這樣的影像.
就是Stained Glass.
就是教堂裡面有一些特別傳統的教堂.
他們的裝修或者他們的裝置.
都會有一些彩色的玻璃窗.
Stained Glass.
陽光照下來的時候就會很漂亮.
是不是.
有沒有印象.
你知道我們教會有嗎.
其實我們沒有.
我們沒有Stained Glass.
因為我們不屬於那個世代.
但是這個世代.
就是以前的教堂.
古代教堂的那個世代.
他們就有一種這樣的裝置.
話說這個裝置就是12世紀開始的.
當時有個神學家他認為教會的牆.
其實應該要有光照下來.
因為光就代表什麼.
代表上帝是不是.
代表上帝.
所以他就將這些牆拆掉.
然後弄一些窗戶上去.
當時也有一些彩色的玻璃開始技術成熟.
所以就用到這些新的技術.
來做一些Stained Glass.
很多時候就弄一些Icon.
或者弄一些聖經的故事.
或者一些圖畫等等.
就給信眾.
去到教堂的時候.
他能夠看到可能聖經一些的信息等等.
Stained Glass告訴我們.
每一個碎片是一起放在一起.
才會組成整幅圖畫.
那個Stained Glass每一塊的碎片.
它自己只是一個很小很小的角落.
它需要其他的碎片.

$^{121}$來跟它一起組織.
並且背後有一個工匠.
然後工匠很巧妙地將Stained Glass.
不同的顏色,圖案等等.
放在一起.
然後才會有一幅圖畫.
為什麼用這個圖畫來開始.
因為我們繼續看創世記的故事的時候.
我們發現創世記的故事.
特別是那些所謂族長的故事.
阿布拉罕,以撒,雅各,約瑟,猶大等等.
由阿布拉罕開始.
直到最後創世記的時候.
有十二眾子一個大的家族出現.
原因是什麼.
因為上帝要藉著阿布拉罕來有一個任務.
使萬族能夠因祂得福.
這個上帝.
永遠是在任務上的上帝.
就是將這些破碎的生命組成在一起.
使世界和身邊的人能夠得到祝福.
我們上一次在這個系列開始的時候.
聖經介紹這個上帝是一個宣教宗的上帝.
面對一個被罪破壞的世界.
上帝沒有放棄並且願意去祝福.
所以呼召阿布拉罕出來.
這一段經文就是創世記第十二章.
最後一句就是最重要的.
他說地上的萬族都要因你.
阿布拉罕或者阿布拉蘭來得福.
就是上帝要藉著他來施行.
他自己在寶世上祝福萬族的救恩計劃.
使萬族能夠去認識他.
上一次我們聚焦在阿布拉罕的故事.
阿布拉罕的故事.
我們說阿布拉罕並非打不死的英雄.
其實是一個很真實和很脆弱的人物.
今天我們的鏡頭就較闊一點.
我們就看阿布拉罕之後的後代.
到底發生了什麼事.

$^{161}$如果你有一點印象的話.
你記得阿布拉罕的兒子是誰嗎.
當然有個以實馬.
有以實和埃塞.
埃塞是不同的一代.
記載了阿布拉罕的故事之後.
就說到以薩的故事.
以薩的妻子就是尼伯格.
然後他們也是不辱的.
人到盡頭的時候.
他們不能夠繼續有後代.
不能夠成為一個大國的時候.
上帝就插手去介入.
有一個所謂的奇蹟兒子出現.
這個奇蹟兒子其實是雙重的.
就是以素和雅各.
你記得嗎.
他們是雙胞胎.
從小到大就是在爭奪東西.
爭奪什麼呢.
就是爭奪產權.
你看這些族長的故事.
幾乎是在看何鴻燊的家族的故事.
到底誰爭奪了多少.
爭奪了多少大狗等等.
這個家族的故事.
其實就是在說這些後代.
如何去處理他們有的所謂的heritage.
他們的基業.
以薩和尼伯格.
他們有的雅各,以素和雅各.
就是兄弟之間的相爭.
你記得在雅各和以素的故事.
是不是有兩次.
雅各去欺騙和奪取.
這個所謂長子的名分.
你記得嗎.
記得嗎.
記得吧.
如果你沒看過聖經的話.

$^{201}$當然你不會有印象.
但如果你稍微有查過經.
或者稍微會聽過.
你就知道他們經常都去爭祝福.
在近東以前的社會裡面.
長子first born.
永遠都是得到的產業是double portion的.
就是你有個蛋糕.
就是家父有的遺產.
好像一個蛋糕一樣.
就切開三塊.
如果有兩個兒子的話就切開三塊.
怎樣分呢.
小兒子就分三分之一.
長子就分三分之二.
好不好.
66\%和33\%.
長子永遠都是double的.
所以他們就很緊張長子的名分.
first born的right of inheritance.
我在香港長大的.
雖然我沒有那種香港長大的背景.
但偶爾聽到一些這樣的故事.
譬如香港有一些叫丁屋的故事.
丁權.
什麼叫丁權.
就是產業是按誰來分的.
按丁來分的.
丁就是兒子.
就是按兒子.
如果你是女兒的話.
不好意思.
你就沒有份的.
因為女兒是嫁出去的.
不屬於你的.
兒子留下來.
所以丁權遺產的爭霸.
就有一點像族長的故事裡面的背景.
長子名分既然收到雙倍的話.
產業權是雙倍的話.

$^{241}$這個一定要爭.
所以雅各是雙胞胎小的那個.
他就經常想保持大長子的名分.
你記得他騙他兩次.
是不是.
騙他兩次.
我不知道你從小到大聽這些故事的時候.
為什麼會騙他兩次呢.
你騙他一次.
你記得紅豆湯的故事.
這個爾素出外打獵.
然後肚子餓回來的時候.
餓得不得了.
然後雅各就弄了一碗粥給他.
是不是.
很便宜的.
然後說如果你肚子餓.
你吃吧.
你吃的話就跟我簽一簽那個share.
剛才說股票.
就是那個股份.
你簽給我.
我就拿66.
你就拿33.
這樣好不好.
這個爾素就很流行看這些.
可能是世外高人打獵.
他就很流行看這些.
算了我拿33.
你拿66.
可能中間有這樣的交易.
所以雅各第一次騙了爾素就拿了66\%.
拿了三分之二的股份.
然後這個家父死之前.
就是以撒死之前.
他要祝福他的兒子.
那裡又有另外一個所謂deception的story.
就是騙.
就是.
就是騙祝福的故事.

$^{281}$你記得應該是這個家父.
以撒去祝福那個大兒子.
爾素是不是.
結果雅各就假扮爾素.
假扮他的哥哥.
然後拿祝福.
是不是.
奇怪了.
為什麼在家父臨終的時候.
還要騙他一次.
之前已經騙了.
你拿了66\%的股份.
那個應該是你.
原來是這樣的.
在舊約裡面的背景.
那個所謂的長子名分.
就是有分兩個層次的.
第一個層次可以叫他做.
就是長子名分.
可以叫他做.
就是那個66\%.
你看到那個受益人.
按道理應該是他的長子.
結果那個小兒子雅各就騙了.
拿了那個雙倍的部分.
就是產業權的雙份他拿了.
性質就是他得到了.
用什麼方法.
就是用紅豆湯的交易.
但是到最後的時候.
家父的祝福.
那個就不是產業權的問題了.
那個是管理權的問題.
管理權的問題.
本身理所當然有.
這個管理權.
就是所謂的trustee.
之前那個就是你的股份.
你的portion.
然後最後決定的就是trustee.

$^{321}$那個管理權.
administration.
那個行政權.
就是執行權.
就是在家父那裡.
理所當然應該是給那個firstborn長子的.
但是雅各就用欺騙的方法.
來爭取了.
拿了二數的管理權.
管理權有什麼在裡面.
為什麼這麼重要要爭取呢.
原來管理權.
administration.
就是在execution.
就是執行上.
什麼時候分紅.
什麼時候分那個share.
什麼時候分利息.
是那個管理權.
去負責的.
所以就算你得到多少的產業也好.
多少部分也好.
都是要靠管理權.
來施行的.
大家記不記得.
類似這樣.
你可以想象.
譬如一個大家族的企業.
阿巴罕以撒亞國.
他們有很多產業的.
你可以想象一個大的企業.
就是父親分給你.
當你這邊是長子.
分給你66\%.
然後這邊分多少.
33\%多.
三分二在這邊.
三分一在這裡.
但是.
你分到多少也好.

$^{361}$你分到三分二也好.
但是執行權的CEO在這邊.
CEO負責什麼時候分.
他可以拖到你最後.
才分給你.
所以你拿著的股份.
其實就是白紙.
明白嗎.
明白吧.
或者我這樣舉個例子.
我昨天想到這個例子.
因為吃飯的時候.
就跟兒子聊天.
假設我的退休金湊起來.
我有一億.
有一億.
如果我有一億就好了.
當我沒有一億.
當我有一億.
我女兒有個兒子.
你不要分男女.
總之兒子分三分二.
多不多.
都挺多.
女兒分三分一.
女兒當然不服氣.
她騙了她第一次.
騙了弟弟第一次.
就拿了三分二.
但是始終執行權在哪裡.
就在兒子那裡.
不要說他長子還是次子.
次子.
總之執行權就是.
因為重男輕女.
所以執行權始終都在我兒子那裡.
舉個例子.
是這樣分的.
一個分三分二.
一個分三分一.

$^{401}$到我臨死之前.
開個你講句.
如果我在ICU那裡.
上氣不接下氣的時候.
應該兒子有三分之二.
他應該就馬上來探望我.
他說爸爸.
你那個網絡帳戶.
那個密碼在哪裡.
密碼是什麼.
沒有人知道.
只有你知道.
你現在還有口氣.
你告訴我那個密碼.
因為密碼就是執行權.
你就可以登入.
然後可以.
將那些股份來現金.
怎樣現金.
你才要我的密碼.
好壞不好聽.
如果是女兒截股的話.
女兒早一個小時來探望我.
然後假扮我的兒子.
然後騙了我.
拿個密碼來.
爸爸你快死了.
你沒氣了.
快要拔喉了.
我就糊裡糊塗.
半夢半醒.
我就半昏迷狀態.
就告訴我女兒密碼.
然後之後那個小時.
我兒子才真正趕來.
爸爸你那個密碼是什麼.
我就跟他說.
剛才不是給了你.
那個不是你.
原來是你姐姐.

$^{441}$你快點給我密碼.
我就給他密碼.
就給他密碼兩次.
他立刻登入.
登入不了.
為什麼登入不了.
當然了.
他姐姐拿了之後.
那三分鐘立刻改了密碼.
就是執行權的問題.
所以為什麼哥哥一定要騙他兩次.
因為你拿到執行權的時候.
你實際上拿不到那些所謂股份.
你看到這個故事裡面就很搞笑.
所以二嫂發現自己兩次被雅各欺騙祝福的時候.
她大哭然後說什麼.
二十七章三十六節.
二嫂就說他名叫雅各.
雅各意思就是說.
拿著不放的意思.
他叫雅各豈不是這樣.
是真的.
他的名字真的叫Jacob.
是這麼壞的.
是這麼狡猾的.
他欺騙了我多少次.
兩次.
先前拿了我的股份就是長子的名分.
後來他又奪了我的福分.
拿了我的執行權.
於是你發現發生什麼事.
就起殺機.
等於如果我兒子發現了這件事的時候.
剛才那個例子.
他一定殺他姐姐的.
他一定會要他的命.
這個故事就是說到他們之間的爭奪.
這是非常人性的黑話.
但是卻在聖經的族長的故事.
就是神的子民的.

$^{481}$Founding Fathers.
聖經是這樣記載的.
聖經在這裡告訴我們這些黑歷史.
他們之間這樣去爭奪上帝的祝福.
好像上帝的祝福是要爭回來一樣.
用自己的手段來拿回來的.
這絕對不是一些很可愛的故事.
我經常都說聖經的故事絕對不是給小孩看的.
是真的.
是一些聖人的故事.
所以結果的後果是怎樣.
耳掃就凍殺機.
然後亞國就流亡在米索不達米亞20年的時間.
然後才有機會再回到那裡.
這是第二個所謂族長的故事.
現在我們再看.
耳掃看完了.
我們再看亞國的妻妾.
亞國在米索不達米亞有兩個老婆.
第一個叫利亞.
第二個叫拉潔.
她愛誰比較多.
她愛拉潔Rachel比較多.
為什麼還有另外兩個女人的名字.
因為這兩個老婆爭什麼.
丁權.
所以爭誰的生子生得多.
你生得越多你的分數就越多.
因為按丁來分數.
所以大家在不育當中想去生.
然後生不到的時候.
以前古代的習俗就是叫屎女來跟你生.
因為工人生的都是屬於你的產業.
所以利亞就找了Sipa.
Sipa.
另一個他的工人去生子.
然後拉潔就看到自己不育.
然後就找了Big Lie去幫他生.
結果生了多少個出來.
12個.

$^{521}$再加上女兒.
女兒不是很漂亮.
她叫底娜.
13個.
總之生了13個.
創世記29章尾的時候去交代.
大家看到那個故事是在說什麼.
就是一幫很破碎的人.
他們去爭奪一些世上的東西.
爭奪所謂上帝給他們的產權.
在質度裡面去爭奪.
在計算裡面去爭奪.
然後才有12支派出現.
然後第三個故事就是大家很清楚.
這個就是約瑟和眾兄弟的糾纏.
因為雅各慢慢繼續下去.
由以薩去到雅各.
阿巴罕 以薩 雅各.
雅各就有12個兒子.
他最愛的是誰.
他偏愛的是誰.
偏愛的就是約瑟.
所謂的拉仔.
最後的時候生的約瑟.
所以因為偏愛他的緣故.
那些兄弟就不喜歡他.
這個不是純粹偏愛的問題.
如果你想想產業的問題.
就是為什麼.
因為password在哪裡.
password在他最愛的兒子那裡.
也將會給他的.
所以為什麼兄弟這麼恨他.
恨不恨他.
如果你是其中一個兄長.
特別是長子的話.
我當然不妥他.
我分到多少都好.
password在他那裡.
是他負責.

$^{561}$是他有權.
他什麼時候分給我.
到我死那天可能他都還沒分給我.
所以他們的兄弟就很想殺約瑟.
這裡有幾段說話.
這個故事一開始的時候.
就說到這些兄長怎樣對約瑟.
他說哥哥們.
37章第4節.
見父親愛約瑟過於他們的時候.
就惱恨約瑟.
就不與他說友善的話.
其實原文就是Shalom.
為什麼這麼重要.
因為希伯來人到今天.
以色列人.
他們大家去到問安的時候.
他們說什麼.
Shalom.
說Shalom.
所以其實這裡不是說不說友善的話這麼簡單.
這裡簡直是Hello都不會說.
每天見到約瑟的時候.
如果你是兄長的話.
你就問他什麼.
Where is the password.
I want the password.
password應該在我這裡.
但他就不說.
你說吹什麼我不告訴你.
然後你看到每一個的故事.
每一個兄長.
第一個最大的長子.
就是Ruben.
裡面.
你看到他們在爭家產.
為什麼.
因為當拉傑.
就是所謂第二個老婆.
死了之後.

$^{601}$然後他就和碧拉同寢.
拉傑的侍女同寢.
按照當時近東的社會.
去到和媽媽的侍女同寢.
是什麼呢.
我們現代人就會想.
一定是性的問題.
其實不是的.
是爭家產的問題.
如果和他侍女同寢的話.
他所生的那兩三個小孩.
那些丁權.
就是誰的.
當然是他.
現在我是你爸爸.
他明我的意思是說.
35章第22節.
然後到第二第三仔.
就是Simeon和Levi.
西冕和利美.
他們就是滿有殺機的人.
那34章說了一個故事.
就是他們殺了另一個外邦人的城市.
這個事件.
裡面的人來將他們通殺.
就用了割禮這件事來設一個局.
那個詳細我們不說.
主要是想交代的一個方面.
就是這兩個人是滿有殺機.
是暴力的人.
第四個就是猶大.
Judah.
猶大這個故事我們就更加清楚.
37章他就是想賣掉藥室的.
其實為什麼要將他賣掉.
不是只憎恨他那麼簡單.
而是因為他有密碼.
賣了你看看還有沒有密碼.
對嗎.
當你死了.

$^{641}$然後那個密碼理所當然應該給誰.
給誰爭到那個所謂長子丁權的那個人.
將那個競爭者剷除.
那你就可以拿到.
其他的兄弟做什麼.
就丟藥室在坑裡.
丟在坑裡面.
然後做什麼.
就坐下吃飯.
坐下吃飯.
你就見到這個就是以色列神的子民.
那個祖先的故事.
兩個字.
嘉辯.
為了什麼.
為了利益.
是不是.
為了利益.
這些我們看完的故事.
我們會發現一點都不光榮.
對嗎.
一點都不光榮.
和我們平時看聖經的印象不是這樣的.
我們覺得放在聖經裡面的都是一些比較神聖的東西.
是不是.
神聖的故事.
但卻是聖經的作者不寫一些光榮的東西.
然後偏偏選一些所謂黑歷史.
他中間有欺騙.
是一些破碎的人.
是一些會出賣大家.
和會彼此相憎的一些人.
為什麼會這樣.
你就可以問一個問題.
到底聖經是不是一些普通的道基督教訓的書.
但很多人特別一些.
可能以前在香港長大的一些朋友.
都是在教會的學校長大的.
是不是.
小學,中學.

$^{681}$為什麼我們經常很喜歡將我們的小朋友.
發他們去教會的學校.
天主教,基督教的學校.
為什麼.
你們也是的.
因為他們教得好.
是不是教得好.
因為德育.
discipline(教育)這些東西.
通常教會的學校比較好一點.
但你見到這個當然是對的.
但很多時候教會學校裡面.
你學的信仰或者看聖經的故事.
就不會看這些的.
是不是.
這些太黑了.
其實很難看的.
是不是.
沒有人解釋給你聽.
你怎麼知道發生什麼事.
是不是.
所以我們通常看一些比較溫馨的東西.
是不是.
愛人如己.
怎樣愛人.
怎樣連戌人等等.
但你見到聖經的作者.
很多時候就是將這些的部分.
不提這麼多.
然後提一些比較黑的地方.
其實他的焦點不是教我們做好人.
聖經的焦點不是做這件事.
聖經的焦點是藉著這些黑歷史.
甚至是神的子民的祖先.
就是我們的founding fathers.
去講他們的黑歷史的時候.
就藉著這些黑歷史黑故事.
去告訴我們.
這位上帝是一個怎樣的上帝.
你看完這些故事.

$^{721}$你真的搖頭.
just shake my head.
這班人如果你是神的話.
如果你說要藉著這個家族.
來使萬國得福的話.
你真的搖頭.
你一定會很火滾.
確實在近東的宗教的概念裡面.
神明見到人們不聽話的時候.
他們會做什麼.
神明都會有殺機的.
神明就會殲滅這些人.
禁止重設.
全部東西重新開始.
因為這班人太吵.
太麻煩.
在近東的宗教的概念裡面.
他們的神學裡面.
暴力其實是世界宇宙的規律的其中一件事.
由始至終都會有暴力的.
由始至終都會有一些暴力出現的.
這是一個正常的工具.
你不喜歡那些人.
你就將他全部消滅.
然後再重新開始.
是不是.
你看到聖經所表達的上帝.
就不是一個這樣的上帝.
是截然不同的.
上帝見到這些人搞到很混亂的時候.
見到他們很多破碎的時候.
上帝他選擇一條不是最簡單的方法.
不是最簡單的方法.
最簡單的方法就是重設.
但是上帝去救贖.
上帝救贖一些破碎的人.
破碎的人生.
然後要將他們重新再組合.
使那些圖畫能夠呈現出來.
這樣做上帝絕對不是洗白那些哥哥.

$^{761}$或者其他族長那些惡士.
惡人始終做的事情是邪惡.
但是上帝是可以克服人的軟弱.
克服人的惡意.
最終是成就他自己的計劃.
這個就是故事的目的.
就是想告訴我們.
Theology.
告訴我們上帝是一個怎樣的上帝.
所以這個故事創世記最後41章至50章.
說到約瑟被賣到埃及的時候.
然後在高高低低的當中.
他就成為了埃及的宰相.
最後發生什麼事.
在饑荒的裡面.
約瑟藉著他的苦難.
藉著哥哥在他身上行的惡事.
他就成為了宰相.
然後41章第57節.
各地的人.
或者甚至全地的人.
在故事裡面全地的人都去了埃及.
向約瑟買糧.
因為天下的饑荒甚重.
弟兄姊妹之間他們再相認之後.
一個很感動的故事.
那些哥哥終於認錯了.
大家有Reconciliation.
45章第8節約瑟說什麼.
猜我來的不是你們.
你們原本是惡意使我被賣到埃及的.
我受了很多苦.
不過我現在知道上帝.
Is on a mission.
他的mission是猜我來這裡.
猜我們在埃及.
然後保存著亞國家.
並且祝福各地的人.
50章第20節.
最後一章.

$^{801}$從前你們興獎你們的義士.
是要害我.
但是上帝的義士原是好的.
要使許多人.
萬族能夠得以存活.
成就今天的光景.
這裡你們看到.
老祖耶和華對亞伯拉罕的呼召.
那個的應許就是萬國要因你得福.
這裡就慢慢實現出來.
這個就是族長的故事.
或者可以這麼說.
這是教會的前輩.
上一代第一代的故事.
數他們的黑歷史出來.
不是要描黑他們.
而是想表達清楚這位上帝.
我們所信的上帝.
我們所讚美的這位上帝.
是一個怎樣的神.
而我們一班怎樣的人.
我們是需要他的恩典.
我們是一班很脆弱.
很不行的人.
族長是虧缺上帝的榮耀.
我們每一個都是一樣.
神才是最厲害的.
在歷史歷代裡面.
教會都會有不同的會議開.
在新約的教會.
到現在大概有二千年的歷史.
久不久有一些所謂的教會會議.
或者大公的教會.
普世的教會.
他們大家坐在一起.
去解決一些神學問題.
或者想寫出.
或者總結出.
到底我們的信仰.
應該怎樣了解.

$^{841}$這些所謂叫信條或者信經.
所以在過去二千年裡面.
有很多這樣的信經是產生的.
其中一個信經叫做.
Belgic Confession.
比利時信條.
在第27條那裡.
這是十六世紀的時候有的信條.
我覺得他描述教會.
他描述得很好.
他說什麼呢.
他說教會.
自世界之初就已經存在了.
在創世記的時候已經存在.
教會的初形.
也將會存在世界末了的時候.
因為基督是永恆的君王.
神的子民不可或缺的.
只要有基督.
就要有基督的子民.
就會有教會.
這個所謂神聖的教會.
是蒙上帝的保守.
免受世上一切的侵擾.
即便在世人的眼中.
教會或許顯得渺小.
甚至瀕臨消亡.
extinct.
即便在世人的眼中.
或許是顯得渺小.
或者瀕臨消亡.
我覺得這句話.
點出教會的神學.
是點得很通的.
教會是很脆弱.
神的子民都是很脆弱.
由第一代開始到現在.
為什麼.
不是因為我們可以.
而是因為上帝可以.

$^{881}$不是因為我們沒有黑歷史.
不是因為我們可以.
我們都是虧缺上帝的榮耀.
但唯獨是基督.
基督是我們主的.
基督去拯救我們.
族長的社會.
教會的社會就是這樣.
記得以前20世紀的時候.
有一個很出名的天主教神學家.
Hans Kuhn.
他曾經說過類似這樣的說話.
他說不是教會守住真理.
而是真理守住了教會.
是科學家的真理守住了教會.
你回去看回這張照片.
所謂Stained Glass.
每一個碎片.
不是單獨的.
最後有一個工匠將這些碎片.
放在一起.
然後成為一個很偉大的圖畫.
看族長的故事.
甚至看我們自己的故事.
我們都會覺得.
我們真的很差勁.
我們真的很不行.
蓄意的破壞.
或者我們會不在乎的心態.
來破壞上帝的工作.
但上帝是可以剋聖我們所有人.
我們每一個人.
我們的軟弱和脆弱的地方.
今天有三個邀請.
我們說過這個系列.
我們會有一些東西給大家.
而這個實物.
這個小的物件.
是給你一個提醒.
提醒你上帝的恩典.

$^{921}$上帝的工作.
上一次開始我們有一個瓶子.
裡面有沙.
你記得我們說過亞伯拉罕的沙的所謂應許.
多如海沙.
你的後裔會多如海沙.
這次你出去的時候.
你可以拿一塊拼圖.
我們當然不會給你玻璃.
不會給你玻璃的碎片.
但這個拼圖.
好像一些碎片一樣.
去告訴我們.
原來我們每個人都是一塊碎片.
是不完整的.
我們是需要上帝的恩典.
將我們拼在一起.
第一個邀請就是一個謙卑的邀請.
我們放下一些有的錯覺.
或者那種錯覺就是.
我夠好的話.
上帝就會用我.
原來上帝一直都不需要等我們好.
才能用到我們.
由第一代開始到現在.
上帝不需要等完美的人.
去繼續他的計劃.
上帝有的能力.
是高過我們的不善.
我們的不足.
第二個邀請.
是一個醒覺的邀請.
醒覺的邀請.
問問你自己.
我有沒有用一些旁觀者的眼光.
去看上帝.
看教會.
上帝既然有他的使命.
去到宣教.
教會就自然有同一個的使命.

$^{961}$我們未必會去到彼此爭兵權.
我們未必會彼此爭password.
但我們可能都會有一種.
關我什麼事的心態.
是不是.
等如那12個兒子.
不是每個都出主意.
都出賣藥室.
有很多都是旁觀者.
是不是.
他們就是關我鬼事的那種心態.
關我什麼事的心態.
我想這個才是最根本的問題.
就是上帝的子民很多時候就是說.
與我何干這種心態.
這種應該是一種醒覺的邀請.
你告訴我.
上帝如果繼續去做他的工作的時候.
為什麼我們還會說.
關我什麼事.
第三個邀請就是參與.
我們就這個碎片.
我們可不可以交給我們自己.
奉獻我們自己.
在上帝公像的手中.
不是因為我夠好.
而是我自己不完美也好.
上帝的能力是高過我的.
整幅圖畫我們就看到上帝的恩典.
上帝自己的大能.
等我出去的時候拿一塊.
好嗎.
拿回你的瓶子.
如果你今天第一次來的話.
你不妨都拿一個瓶子.
因為我們有.
放在當下的地方.
放這塊碎片進去.
去告訴你自己.
我們需要大家.

$^{1001}$需要大家繼續參與上帝自己的工作.
上帝你拼我自己.
不完美的一塊.
下去你自己的圖畫裡面.
使你宣教的工作.
使更多人認識主.
使這個工作能夠繼續延伸.
好嗎.
我們同心協力地祈禱.
因為你是.
那位完美的神.
你是去不斷地宣教.
不斷地去救贖人.
我們能夠成為你.
被救贖的子民的一部分.
將我們在恩典裡面拯救過來.
赦免我們的罪.
我們的不是.
我們的脆弱.
我們的不好的地方.
恩典的偉大.
讓我們同心協力來到你面前.
主讓我們願意將我們的生命拼進去.
你一個救贖普世的計劃當中.
叫人藉著我們的生命.
我們的生活.
更加認識.
耶穌基督愛我們的主.
今天我們都是聖餐的主人.
這個聖餐.
這個餅這個杯.
就是告訴我們.
完全因為你自己的恩典.
我們能夠成為你自己的門徒.
成為你的教會.
主讓這個恩典.
你要提醒我們.
活在恩典裡面.
站穩在恩典裡面.
讓你的名.

$^{1041}$能夠得以被稱頌.
我們這樣祈禱奉主名.
\newpage



\section{}
\label{sec:KUqZwRWtyog}
\textbf{遊歷土耳其 - 老底嘉 (Laodicea) - 陳崇基牧師}
\newline
\newline
連結: \href{https://youtube.com/watch?v=KUqZwRWtyog}{\texttt{https://youtube.com/watch?v=KUqZwRWtyog}} ~~~~ 語音日期: 2026-05-15
\newline
\newline
\hyperref[sec:9gxDgFyKyk4]{\small{< < < PREV SERMON < < <}}
~
\hyperref[sec:index_chronic]{\small{[返順時目]}}
~
\hyperref[sec:index_scriptual]{\small{[返順卷目]}}
~
\hyperref[sec:jnGxkuaDf2M]{\small{> > > NEXT SERMON > > >}}
\newline
\newline
$^{1}$我們來到魯迪加的另一邊.
剛才走過主街來到這邊.
我現在站在玻璃天花的上方.
下面是一個神廟.
考古學家正在發掘中.
所以看到下面有一些很大的神廟.
我們現在就走進去看看.
這裡有個很大的神廟.
在這裡考古學家正在發掘中.
所以看到下面有一些很大的石柱.
在這裡我們可以做一點地理的講解.
因為現在我們就在魯迪加.
在魯迪加的對面的北面.
就是剛才看到的白色的山.
那個叫做Pomacalli.
Pomacalli的意思是棉花堡.
那邊有一些硬泳池和很多礦物.
在那裡今天也是.
是一些溫泉的地方.
因為那裡地質上有些水在水泉湧上來.
而且那些水在地下水中溶解了很多礦物.
以致他們出來的時候是一些很熱的水.
並且滿有礦物.
魯迪加就在這裡中間.
後面這裡今天天氣不太清楚.
後面這裡有一個叫做高山.
高山的山腳就是哥羅西的Colosse.
這些地方我們新約的聖經.
熟的話我們也有聽過.
歌羅西書.
那裡有一個基督徒的群體教會.
其實歌羅西書是寫給他們的.
這裡就是魯迪加.
你記得歌羅西書最後的時候.
或者說這封信是給魯迪加的教會.
這個魯迪加的地方.
我們看到地理在北邊就是棉花堡.
然後在旁邊高山的山腳就是哥羅西.
魯迪加的問題就是吃水的問題.
因為自然的水流到這裡.

$^{41}$仍然是有礦物的.
雖然水沒有棉花堡那邊那麼熱.
但是水裡面仍然有一些礦物.
所以不適合飲用.
在棉花堡水有高礦物質和熱.
所以有醫療和溫泉的作用.
所以是有用的.
去到哥羅西山腳.
那裡所有的水都是清水.
沒有什麼礦物.
所以可以用來飲用.
魯迪加偏偏這個地方沒有好的食水.
所以為什麼在啟示錄第七個教會的書信.
寫給魯迪加的時候.
就說他們很有錢.
這個城市住的人很有錢.
因為他們靠著羊毛的生意.
眼藥.
這些全部都是他們引以為傲的錢財資源.
啟示錄的書信就去到指責魯迪加教會.
說他們不冷不熱.
上帝基督說我寧願你熱或者冷.
很多時候我們去解釋這段經文的時候.
我們就說神不想我們不冷不熱.
寧願我們熱或者寧願我們冷.
我們用一個屬靈溫度來看.
但是你這樣看的時候就不太對.
因為上帝沒有理由想我們心冷.
寧願我們心暖一點.
雖然不是很熱.
怎樣去解釋呢.
就是在地理上水質上可以去解釋.
有了這樣的背景.
我們知道原來不冷不熱不是量度屬靈的溫度.
是說那種水的功能.
在棉花堡在北邊有一些礦物水.
是熱的水.
溫泉有醫療的作用.
在哥羅西這邊下流一點.
就是清水.

$^{81}$老底家雖然很有錢.
有羊毛.
然後對著萊克斯山谷.
一個很肥沃很富庶的山谷.
農產或者牧羊業非常發達.
他們很有錢.
幾次地震都不需要羅馬政府補貼津貼去重建.
靠自己就能夠重建.
但是在神的眼中.
他們所做的事就是不冷不熱.
就是說他們的水是不能夠使用的.
甚至喝了會因為苦而嘔吐.
這裡是警覺教會.
我們所做的事.
基督徒所做的事.
我們到底有沒有那種屬靈的功能.
有沒有那種屬靈的效果.
我們所做的事.
是否只是去到外表.
並非能夠真的去事奉上帝.
這個值得我們去反省.
這個老弟家教會 上帝說不想你不冷不熱.
寧願你熱 或者寧願你冷.
給上帝能夠去使用的.
我們自己是不是一些所謂不冷不熱.
不冷不熱不能夠被神的國去祝福.
不能夠去造就其他人.
我們花了不少時間.
去解釋老弟家的地理水質.
這些背景能夠幫助我們.
知道啟示錄寫給老弟家教會書上的內容.
以致我們能夠更加掌握這個訊息.
很有錢的教會.
很富足的人.
但是他們所做的一切.
上帝看為不冷不熱是沒有用的.
不能夠喝的 也不能夠泡溫泉的.
用來做什麼呢.
有什麼好自誇呢.
這個就是教會書信第七封.

$^{121}$寫給老弟家的背景.
掌握了之後就更加清楚.
在這裡有很多發掘的工作.
因為這個篩都是很大的篩.
有一些發掘出來的碎片 瓦片等等.
就會堆在一起.
然後繼續去整理等等.
非常之多.
這裡我們找到一塊水喉.
你看到水喉是圓圓的.
裡面有一些所謂的氧化的礦物.
因為水經過這裡的時候有沉澱.
然後日以累聚.
日久經過之後.
就會一直沉澱一層一層的礦物.
這些礦物讓我們知道這裡的水質非常差.
不冷不熱 就是這個意思.
所以看到水喉就知道.
這裡的教會有錢也沒用.
但是你要做的事情要有效果.
要為上帝所用.
能夠造就別人 能夠幫助別人.
這裡我們讀一段啟示錄這封信.
你要寫信給老弟家教會的天使.
那位阿們就是忠實可靠的見證者.
那位上帝創造之根源.
這樣說 我知道你所做的.
我知道你不冷不熱.
我倒願意你或冷或熱.
因為你像溫水一樣.
既不冷也不熱.
我要從我口中把你吐出來.
你說你富足 寬裕 毫無缺乏.
然而你不知道你自己是多麼悲慘可憐.
你貧乏 在上帝眼中.
你黑眼 又是赤身露體的.
沒有羊毛.
這裡羊毛很出名的.
眼藥很出名的 衣服很出名的.
但在上帝眼中他是赤身露體.

$^{161}$和黑眼貧窮的.
所以我勸你 向我上帝.
去買精靈的金紙.
向我這位復活的基督.
去買精靈的金紙.
因為這樣才能夠使你富足.
買白袍穿上.
可以遮蓋你羞恥的裸體.
也可以在我這裡買眼藥.
去擦一擦你的眼睛.
以致你能夠看見.
這是基督對教會一個很嚴厲的警告.
對我所愛的人.
基督說什麼 我都會責備管教.
所以你們要熱心 要悔改.
聽罷 我站在門外叩門.
若有人聽見我聲音而開門.
我就遵去.
我要和他一起坐直吃飯.
我也要跟他一起去享用賢直.
上帝雖然責備教會.
但仍然希望教會悔改.
以致能夠和他建立關係.
以致那得勝者.
我要似特權給他.
坐在寶座跟我同坐.
正如我已經得勝一樣.
現在跟我父親天父.
同坐在他的寶座上一樣.
聖靈向宗教會說的話.
有耳朵聽的 都應該去聽.
\newpage



\section{士師記 6:11-16}
\label{sec:jnGxkuaDf2M}
\textbf{Faith over Fear (Judges 6\_11-16) - Chee Yie Jong 信心勝過恐懼 (士師記6\_11-16) - 楊志溢傳道}
\newline
\newline
連結: \href{https://youtube.com/watch?v=jnGxkuaDf2M}{\texttt{https://youtube.com/watch?v=jnGxkuaDf2M}} ~~~~ 語音日期: 2026-05-18
\newline
\newline
\hyperref[sec:KUqZwRWtyog]{\small{< < < PREV SERMON < < <}}
~
\hyperref[sec:index_chronic]{\small{[返順時目]}}
~
\hyperref[sec:index_scriptual]{\small{[返順卷目]}}
~
\hyperref[sec:HpXR4dbZvbk]{\small{> > > NEXT SERMON > > >}}
\newline
\newline
士師記 6:11-16
\newline
\begin{longtable}{cl}
\hline
\hline
章節 & 經文 (和合本修訂版)\\
\hline
6:11 & \begin{tabularx}{0.7\textwidth}{X} 耶和華的使者到了俄弗拉，坐在亞比以謝族約阿施的橡樹下。約阿施的兒子基甸正在醡酒池那裡打麥子，為了躲避米甸人。 \end{tabularx} \\ \\ \relax
6:12 & \begin{tabularx}{0.7\textwidth}{X} 耶和華的使者向基甸顯現，對他說：「大能的勇士啊，耶和華與你同在！」 \end{tabularx} \\ \\ \relax
6:13 & \begin{tabularx}{0.7\textwidth}{X} 基甸對他說：「主啊，請容許我說，耶和華若與我們同在，我們怎麼會遭遇這一切事呢？我們的列祖告訴我們：『耶和華領我們從埃及上來』，他那奇妙的作為在哪裡呢？現在耶和華卻丟棄了我們，把我們交在米甸人的手掌中。」 \end{tabularx} \\ \\ \relax
6:14 & \begin{tabularx}{0.7\textwidth}{X} 耶和華轉向基甸，說：「去，靠著你這能力拯救以色列脫離米甸人的手掌。我豈不是已經差遣了你嗎？」 \end{tabularx} \\ \\ \relax
6:15 & \begin{tabularx}{0.7\textwidth}{X} 基甸對他說：「主啊，請容許我說，我怎能拯救以色列呢？看哪，我這一支在瑪拿西支派中是最貧寒的，我在我父家又是最微小的。」 \end{tabularx} \\ \\ \relax
6:16 & \begin{tabularx}{0.7\textwidth}{X} 耶和華對他說：「我與你同在，你就必擊敗米甸，如擊打一個人。」 \end{tabularx} \\ \\
[1ex]
\hline
\hline
\end{longtable}
$^{1}$[END PLAYBACK].
All right..
Happy Mother's Day again to all the moms and the daughters.
in the room..
[NON-ENGLISH SPEECH].
OK..
All right..
Does anyone recognize this athlete?.
[NON-ENGLISH SPEECH].
So this is Elisa Liu, and she became a household name.
after winning the gold medal in women's figure skating.
at the 2026 Winter Olympics..
But her journey to become the best version of herself.
is actually not a typical one..
[NON-ENGLISH SPEECH].
So Elisa, she was a child skating prodigy,.
and she actually become the youngest woman.
to win a US championship title at the age of 13..
However, after the 2022 Winter Olympics in Beijing,.
where she finished sixth, she actually.
decided to retire from the sports at just the age of 16..
And she cited the reasons of burnout, loss of passions,.
as well as just the desire to experience a normal life.
outside of the intense and regimented environment.
of elite skating..
[NON-ENGLISH SPEECH].
[APPLAUSE].
[NON-ENGLISH SPEECH].
So in the eyes of the world, this retirement.
might signify the end of the career for this child prodigy..
But what is so great about Elisa's story.
is that during her time away from the competitive scene,.
she actually rediscovered her joy and passion.
for figure skating..
And she came back two years later,.
but this time in her own terms with a healthier relationships.
with the sport..
And as we now know, the rest is history..
[NON-ENGLISH SPEECH].
[APPLAUSE].

$^{41}$[NON-ENGLISH SPEECH].
So when we look at her story, she.
is actually able to live out her full potential as a figure.
skater when she learns to overcome.
the pressures and expectations that.
led to her initial burnout..
And when it comes to our faith as well,.
what is holding us back from living out our full potential.
as God has called us to be?.
And if we go underneath the surface,.
what are some of the fears that prevent us from living out.
a fully surrendered life?.
[NON-ENGLISH SPEECH].
the world..
[NON-ENGLISH SPEECH].
So today, we're going to look at the story of Gideon.
from the Book of Judges..
And we're going to see how he actually.
allowed his fear to hold him back.
from living out his calling..
Some of you might remember Gideon from your Sunday school.
as this great spiritual leader who.
trusted the Lord through his famous wolf fleece incidents..
But if you recall from my earlier message,.
the whole Book of Judges actually.
show us Israel's slow descent into spiritual apostasy.
and moral corruption..
And this downfall of Israel actually.
becomes a lot more apparent when we.
come to the cycle of Gideon..
He is actually the first judge where.
we begin to see some disturbing weakness in his life..
[NON-ENGLISH SPEECH].
So let's look at our passage in verse 11 and 12..
"The angel of the Lord came and sat down under the oak in Ofra that belonged to Joash the Abiezrite,.
where his son Gideon was threshing wheat in a winepress to keep it from the Midianites..
When the angel of the Lord appeared to Gideon, he said, 'The Lord is with you, mighty warrior.'".
[NON-ENGLISH SPEECH].
And so here we're being introduced to Gideon for the first time in Chapter 6..
And compared to the previous judges who only actually serve as God's instruments of deliverance.

$^{81}$as part of the larger narrative cycle, Gideon himself will actually become the focal point of this cycle..
[NON-ENGLISH SPEECH].
And so when the angel of the Lord appeared to Gideon, he was threshing wheat, not on the threshing floor,.
but we're being told it was in a wine press..
[NON-ENGLISH SPEECH].
So Gideon threshing wheat in a wine press highlights how unusual and also fearful the situation was..
Because normally, wheat was processed on an open threshing floor,.
where the wind would just blow away the lighter shaft,.
leaving the heavier grain on the threshing floor to be collected..
Because of that, threshing floors were usually found in open spaces that are highly visible,.
so you can have maximum exposure to the wind..
[NON-ENGLISH SPEECH].
And in contrast, wine press were often built underground,.
because you need a cool and shaded place to produce good wine..
And so as you can imagine, threshing wheat in a wine press is probably highly ineffective,.
because there's barely any air movement in an enclosed space underground..
[NON-ENGLISH SPEECH].
And yet Gideon chose this hidden space because of his fear,.
so that he could hide from the Midianites,.
who were ready to take away whatever little grain that he could produce..
[NON-ENGLISH SPEECH].
So Gideon's response to the message from the angel of the Lord.
will actually show us a lot about who he is,.
his inner life, his belief about God,.
and also his understanding of God's power and deliverance..
[NON-ENGLISH SPEECH].
So verse 13, "Pardon me, my Lord," Gideon replied,.
"but if the Lord is with us, why has all this happened to us?.
Where are all His wonders that our ancestors told us about when they say,.
'Did not the Lord bring us up out of Egypt?'.
And now the Lord hath abandoned us, and given us into the hand of Midian.".
[NON-ENGLISH SPEECH].
When the angel of the Lord appeared and said to Gideon,.
"The Lord is with you, mighty warriors,".
I can imagine Gideon probably being shocked initially..
But when the shock wore off, and he had some time to gather his thoughts,.
he was instantly triggered by this greeting,.
and his shock turned into a complaint..
"What do you mean the Lord is with us?.
If He is with us, why am I threshing wheat in the winepress now?.

$^{121}$And do you know how much harder it is to thresh wheat in the winepress?.
Do you know the constant fear that I'm living in every single day right now.
with the threats of the Midianites?".
[NON-ENGLISH SPEECH].
And Gideon did not stop there..
He continued, "Where is God now?.
And why didn't He show up when we needed Him the most?.
Where are all the wonders that our ancestors tell us.
about when God delivered them out of slavery from Egypt?".
[NON-ENGLISH SPEECH].
And so even though Gideon didn't say it out loud,.
but deep down in his soul, Gideon was actually questioning God's faithfulness.
and accusing God for breaking His covenant with Israel..
[NON-ENGLISH SPEECH].
And so while this feeling of God is at the wrong for breaking the covenant.
and abandoning Israel is as real as it can be for Gideon,.
he knows that Israel is actually the one who breaks the covenant with the Lord,.
disobeying the Lord by worshiping the gods of the Canaanites and the Amorites..
And the circumstances that Gideon and his people are facing right now.
is actually a direct consequence of their disobedience,.
not God's unfaithfulness..
[NON-ENGLISH SPEECH].
If you're reading through the Book of Judges and you get to this verse,.
you'll notice this, "Wait a second, didn't the Lord just deliver Israel.
from the hand of Jabin in the cycle of Deborah in just the previous chapter?".
However, for Gideon, his fear of circumstances led to his failure.
to recognize God's deliverance as well as his failure to recognize Israel's own sinfulness..
[NON-ENGLISH SPEECH].
So I don't know if you have experienced something similar before,.
where you let the fear of your circumstances cloud your faith in the Lord,.
or when your fear becomes so great that you begin to minimize God.
and to question His faithfulness..
[NON-ENGLISH SPEECH].
So maybe some of you are in a season of doubt or uncertainty or suffering..
And imagine someone come up to you and say,.
"The Lord is with you, Patrick," or "The Lord is with you, Kelly.".
How would you respond to that?.
[NON-ENGLISH SPEECH].
I'm sure you're familiar with the song, "Goodness of God," right?.
And to be completely honest, there is a season in my life.

$^{161}$where I couldn't bring myself to sing the words in this song..
I'm not going to sing it, but I'm going to read..
So this is the chorus..
"All my life you have been faithful. All my life you have been so, so good.".
And I couldn't sing it because in my mind I was questioning,.
"Lord, if you are faithful, if you are good,.
where are you now when I need you the most?".
[NON-ENGLISH SPEECH].
And at that time, whenever I hear the song,.
my entire being just repels against the truth that God is faithful..
Because like Gideon, I was in a dire situation,.
and I let the fear of my circumstances took hold of me,.
instead of choosing to trust in the faithfulness of God.
that He has shown me time and time again..
[NON-ENGLISH SPEECH].
So the question for us is,.
are we allowing our circumstances to affect our trust in God?.
Are we so fixated on the challenges or hardships that are right in front of us.
that we forget our God is an all-powerful, all-knowing, and a loving God?.
[NON-ENGLISH SPEECH].
So Church, if we want to live out our identity and calling as God's children,.
we need to first learn to identify where in our lives.
have we allowed the fear of our circumstances undermine our faith in God?.
[NON-ENGLISH SPEECH].
So this is the first fear that we see in Gideon's life,.
the fear of circumstances..
His initial response is doubtful..
God is still patient with him..
So in verse 14, God Himself is speaking to Gideon directly now..
And so He reassured Gideon..
The Lord turned to him and said,.
"Go in the strength you have and save Israel out of Midian's hand..
Am I not sending you?".
[NON-ENGLISH SPEECH].
In verse 15, this is how Gideon replied,.
"Pardon me, my Lord, but how can I save Israel?.
My clan is the weakest in Manasseh, and I am the least in my family.".
[NON-ENGLISH SPEECH].
So this time, Gideon's response shows us another fear that he had,.
the fear of his own inadequacy..

$^{201}$And this fear actually became a hurdle for Gideon.
when God called him to save Israel out of the hand of Midian..
[NON-ENGLISH SPEECH].
If you notice how he responds, politely he questions God's choice..
"Why me? My clan is the weakest in our entire tribe..
And on top of that, I am the least among my whole family..
And so God, there must be other better candidates out there, right?".
[NON-ENGLISH SPEECH].
I'm sure we all have times where we don't feel confident enough, right?.
And sometimes, this can be a good thing,.
especially when we are aware of our own ability and our own limits..
[NON-ENGLISH SPEECH].
Recently, I watched a few interviews of Alex Hono after he free-soloed Taipei 101..
While a lot of people criticized him for being reckless and doing something so dangerous,.
he actually just described it as taking smart, calculated risk..
Because he knows his own ability and limits so well.
that he knew the likelihood of him falling and dying is actually very low..
And so this is a healthy sense of self-confidence..
[NON-ENGLISH SPEECH].
However, when we go to either end of the extremes,.
either being excessively confident or overly insecure, it becomes an issue..
[NON-ENGLISH SPEECH].
Maybe some of you, you sense that God is calling you to do something that you are not comfortable with..
Or maybe the Holy Spirit has been prompting you to talk to your coworkers or your neighbors..
But a lot of times, instead of saying yes right away, our response might look like this..
"Oh, I don't know, God. I am not theologically trained or I didn't go to seminary..
So I will leave that job to my pastors.".
Or maybe you will respond like this..
"Oh man, God, that really makes me feel uncomfortable..
You made me an introvert and now you want me to talk to my neighbors?".
[NON-ENGLISH SPEECH].
So again, the question for us, where have we let the fear of inadequacy hold us back from living out our calling?.
And where have we allowed our self-perceived weakness become an excuse.
when God is inviting us to do something for His kingdom?.
[NON-ENGLISH SPEECH].
And so again, if we want to live out our full identity and calling as God's children,.
the second fear that we need to learn to identify and overcome is the fear of our own inadequacy..
[NON-ENGLISH SPEECH].
So, look at how the Lord responds to Gideon's fear..
In verse 16, the Lord answered, "I will be with you.".

$^{241}$Oh, sorry. So the reality is as long as we live on this side of Jesus' second coming,.
we will always struggle with the fear of our circumstances and the fear of our inadequacy, right?.
And this is actually a completely normal experience living in a fallen world..
So where is our hope in all of these?.
[NON-ENGLISH SPEECH].
So, in verse 16, the Lord answered to Gideon, "I will be with you,.
and you will strike down all the Midianites, leaving none alive.".
And in fact, if you pay attention to how God responded to Gideon each time,.
He always reassured Gideon with this same promise..
We know in verse 12, the angel of the Lord say, "The Lord is with you, mighty warriors.".
And then in verse 14, the Lord say, "Am I not sending you to save Israel?".
Meaning that God will actually be on Gideon's side..
And then right here in verse 16, "I," meaning the Lord, "will be with you.".
[NON-ENGLISH SPEECH].
So, Church, this is where our hope is, that God is with us..
This does not mean that all of our fears will just magically disappear or go away,.
but it means that in the midst of our fears, something greater and more powerful is actually on our side..
[NON-ENGLISH SPEECH].
Instead of focusing on the fears that are right in front of us, we can learn to lift up our eyes,.
above our circumstances, above our own inadequacy,.
and fixated on the promise of God that He will be with us,.
because He is always consistent with His promise..
[NON-ENGLISH SPEECH].
So God is with Moses when he brought Israel out of Egypt..
[NON-ENGLISH SPEECH].
And God is with Joshua when he led Israel into the Promised Land..
[NON-ENGLISH SPEECH].
And God is with Gideon, even when his fear is greater than his faith..
[NON-ENGLISH SPEECH].
And when we come to the New Testament, Jesus Himself is Emmanuel, God with us..
[NON-ENGLISH SPEECH].
And before Jesus ascended back into heaven after His resurrection,.
Jesus gave His disciples, including all of us today, this great commission in Matthew 28..
"Go and make disciples of all nations, baptizing them in the name of the Father, the Son, and the Holy Spirit,.
and teaching them to obey everything that I have commanded you.".
But this calling from Jesus doesn't just enter, but it comes with the same promise..
Jesus continued, "And surely I am with you always to the very end of the age.".
[NON-ENGLISH SPEECH].
So church, we know that we are able to live out this calling despite our fear,.
despite our own inadequacy, because we know that God is with us..

$^{281}$[NON-ENGLISH SPEECH].
So let's pray together..
[NON-ENGLISH SPEECH].
Thank you for being a God who is with us, who is in us, and who is for us..
[NON-ENGLISH SPEECH].
In areas of my life where my fear is greater than my faith, give me the strength to overcome these fears..
[NON-ENGLISH SPEECH].
And teach me to move from a life of fear into a life of faith in you, Lord..
[NON-ENGLISH SPEECH].
And in your name we pray..
[NON-ENGLISH SPEECH].
All right. Thank you all for joining us today..
We have some refreshment thanks to Carol, I think..
So, yeah, feel free to stay around and have some community time, fellowship time..
[NON-ENGLISH SPEECH].
(gentle music).
\newpage



\section{約翰福音 19:25-27}
\label{sec:HpXR4dbZvbk}
\textbf{母親 這是你的兒子 (約翰福音19\_25-27) - 梁紫健傳道}
\newline
\newline
連結: \href{https://youtube.com/watch?v=HpXR4dbZvbk}{\texttt{https://youtube.com/watch?v=HpXR4dbZvbk}} ~~~~ 語音日期: 2026-05-18
\newline
\newline
\hyperref[sec:jnGxkuaDf2M]{\small{< < < PREV SERMON < < <}}
~
\hyperref[sec:index_chronic]{\small{[返順時目]}}
~
\hyperref[sec:index_scriptual]{\small{[返順卷目]}}
~
\hyperref[sec:y_bUxa4hgT8]{\small{> > > NEXT SERMON > > >}}
\newline
\newline
約翰福音 19:25-27
\newline
\begin{longtable}{cl}
\hline
\hline
章節 & 經文 (和合本修訂版)\\
\hline
19:25 & \begin{tabularx}{0.7\textwidth}{X} 站在耶穌十字架旁邊的，有他的母親、姨母、革羅罷的妻子馬利亞，和抹大拉的馬利亞。 \end{tabularx} \\ \\ \relax
19:26 & \begin{tabularx}{0.7\textwidth}{X} 耶穌見母親和他所愛的那門徒站在旁邊，就對母親說：「母親，看，你的兒子！」 \end{tabularx} \\ \\ \relax
19:27 & \begin{tabularx}{0.7\textwidth}{X} 又對那門徒說：「看，你的母親！」從那刻起，那門徒就接她到自己家裡去了。 \end{tabularx} \\ \\
[1ex]
\hline
\hline
\end{longtable}
$^{1}$(第1節).
各位大英姐妹,早安.
各位母親節快樂.
在母親節的時候.
我也想分享一下我跟媽媽的點滴.
記得在我小時候.
好像是幼稚園畢業禮.
在亞洲,香港.
在暑假的日子,天氣很熱.
學校很有趣,舉辦了戶外的畢業禮.
所有家長和學生.
都在露天操場上曝曬.
但一白遮也沒有.
你知道小朋友還好.
戴著畢業帽,有帽子遮著.
但家長們不斷曬.
大汗淋漓,不斷流汗.
但看著自己的小朋友.
畢業的家長都很滿足,很開心.
其實學校也知道很熱.
雖然沒有安排太陽傘.
但他們有安排可樂車.
讓小朋友排隊拿可樂.
大熱天,不喝東西很容易中暑.
當時我也覺得很熱.
我小時候是一個胖子.
很熱,渾身都流汗.
然後走過去拿一杯可樂.
喝完後就喝光.
然後再拿一杯,再喝一杯.
很快喝完後.
回到座位繼續曬太陽.
然後我媽媽問我.
「為何你不倒一杯給我」.
然後我說「糟糕了」.
其實媽媽曬太陽很久.
她沒有遮掩,又沒有飲品.
但可樂車只有小朋友去拿.
大人又不好意思去拿.
所以小朋友去拿.

$^{41}$我馬上折返去可樂車拿東西.
怎料可樂車的人差不多喝光了.
只剩下半杯.
我拿回杯子,只剩下半杯.
跟媽媽說「媽媽,只剩下半杯」.
媽媽看見我大汗濕透.
就說「不用了,你喝吧」.
小時候對媽媽的印象是.
她不是特別需要喝東西和吃東西.
很多時候她都先讓我吃,先讓我喝.
其實是真的嗎.
長大後才發現.
原來不是媽媽不需要喝東西和吃東西.
而是她總是把她最需要的東西給我.
把她需要的東西放在最後.
這是一個愛的奉獻和犧牲.
但同一時間,當她問我.
為何不替她倒一杯可樂時.
其實她期望甚麼.
她期望我留意到她的需要.
希望有一天我能夠照顧她.
希望有一天我留意到.
身邊人的需要去照顧他人.
這是我跟媽媽的一個很小的片段.
令我思考究竟母親的期望.
和子女對母親的愛之間的張力.
聖經是怎樣看的.
通常母親節.
我們都比較少說這段經文.
因為這段經文有一段很難處理的地方.
就是約翰很刻意地重複用一個字.
就是「母親」.
這個詞在耶穌十字架之上.
祂竟然沒有叫自己的媽媽「母親」.
反而用了一個很突兀的稱呼.
去叫祂的媽媽「夫人」.
中文的聖經譯者都覺得這個位置很奇怪.
所以有些譯本.
如果你看到手上有其他中文譯本.
會把「夫人」譯成「母親」.

$^{81}$然後在後面打了一個括號.
寫了原文是「夫人」.
中文譯者都覺得很奇怪.
所以他們想幫我們撫平這種奇怪感覺.
但其實聖經有時用一些很突兀的字眼.
可能是想表達一些特別的訊息.
在我們的文化裡.
如果一個兒子在眾多人面前.
叫自己的媽媽「夫人」.
聽起來會怎樣.
好像很不尊重和冷冰冰.
約翰刻意將「夫人」繼續保留下來.
他沒有將「夫人」變成「母親」.
他保留了耶穌的原話.
其實要我們想一想.
停一停.
為何耶穌要這樣說.
他究竟想提醒十字架下的母親.
和旁邊的門徒是甚麼.
我們一起讀今天的經文.
我們看看今天的經文.
其實不是很長,只有三字.
是約翰福音的19章25至27節.
這一段約翰描寫.
耶穌和母親瑪利亞的對話.
是耶穌已經釘在十字架上.
耶穌流著血,在呼吸很辛苦時.
和母親最後的對話.
在祂很辛苦時,和母親最後的對話.
約翰形容站在十字架下的.
有甚麼人呢?有耶穌的母親.
有母親的姐妹,即是耶穌的姨母.
有格羅伯的妻子瑪利亞.
和莫大拉的瑪利亞.
大家看到有很多瑪利亞.
至少有三個.
約翰好像把鏡頭放近.
拍攝這幾位女士站在十字架下.
讓我們看到耶穌在十字架上.
看著祂最親,最愛的人.

$^{121}$在這幾個短的經文中.
「母親」這個字出現了多少次.
「母親」這個字至少出現了五次.
這裡刻意提到母親的姐妹.
其實你看到「姨母」.
中文翻譯很簡單,把「姨母」寫進去.
但原文是七二重複的.
它說是「母親的姐妹」.
第26節.
耶穌看到母親和祂所愛的門徒站在旁邊.
對母親說.
「母親,看你的兒子」.
其實這裡的原文是「婦人」.
中文翻譯是「撫平」這種突兀感覺.
把婦人改為母親.
因為剛才也說過.
用「婦人」去稱呼母親.
給人的感覺是怎樣.
好像很不熟悉的感覺.
為何耶穌要這樣做.
剛才提到這麼短的經文中.
至少提到母親多少次.
五次….
但耶穌在母親瑪利亞的直接對話中.
一開口就說….
用甚麼?「婦人」.
祂不是用「母親」.
所以大家都覺得很奇怪.
為何短短三節的經文中.
至少提到母親五次.
但耶穌跟媽媽說話卻不用「母親」.
這是一個很強烈的對比.
我想現場的讀者也會感受到這種突兀感.
母親….
母親的姐妹.
但耶穌一開口就說「婦人」.
然後耶穌對門徒說.
「你看,是你的母親」.
從那一刻開始.
門徒就把她接回自己的家裡.

$^{161}$門徒按照耶穌的意願.
把耶穌的母親瑪利亞接回家.
耶穌看到母親所愛的門徒.
開口說這兩句很短.
但其實是很深奧的話.
「婦人,你看,這是你的兒子」.
「你看,這是你的母親」.
從那一刻開始.
門徒就把瑪利亞接回家.
這是約翰給我們一個很簡單的交代.
但背後有很深的屬靈意味.
值得我們思考.
所以我們看完這麼短的經文.
起碼會有兩個問題.
第一,為何耶穌要叫自己的母親做婦人.
約翰不斷提及母親.
但耶穌開口竟然不是用母親.
而是用婦人稱呼媽媽.
這個落差有甚麼意思.
第二,耶穌在十字架上.
最後跟母親說的話.
是把母親交給門徒.
這個意思是甚麼.
如果準備好.
我們今天會回答這兩個問題.
我們看看.
回答第一個問題就是.
為何耶穌要叫自己的母親做婦人.
如果大家細心看《約翰福音》.
就會發現.
耶穌跟母親瑪利亞直接對話.
約翰只記錄了兩次.
這兩次都是用婦人做開頭.
第一次是耶穌最初第一次行神蹟.
在加拿大婚宴的水變酒裡面.
而第一次是看最後.
十字架最後的對話.
這很明顯是約翰留給我們的首尾呼應.
在加拿大婚宴裡.
母親瑪利亞看到酒用完.

$^{201}$很自然就期望耶穌出來幫忙.
挽回親戚朋友的面子.
他跟耶穌說「他們沒有走」.
耶穌回答第一句話是甚麼.
「婦人,我與你何干」.
「我的時候還未到」.
這句話聽起來很突兀.
亦是我們平時很難想像.
一個孝順的孩子.
會這樣跟媽媽說話,對吧.
現在耶穌在十字架上.
再一次對著瑪利亞.
用同一個稱呼.
「婦人」.
約翰刻意讓我們回想.
加拿大婚宴的那場對話.
回想起耶穌當時說的話.
「我的時候還未到」.
現在到了十字架的這一刻.
耶穌其實是向母親宣告.
「現在這個時候,就真正到了」.
在加拿大婚宴的時候.
耶穌第一次行神蹟.
用水變酒,對吧.
表面上,如果大家記得.
是幫人撿回彩票.
顧著家族和朋友的面子.
但耶穌在過程中已經提到.
祂真正要完成的使命.
不只是要幫人保住面子的神蹟.
而是做一個更大,更榮耀的時候.
就是甚麼.
祂透過祂在十字架上的死.
來彰顯上帝的榮耀.
所以耶穌在十字架上再見婦人時.
就好像祂看著祂的母親.
和在場的門徒和不同的婦女.
你們是否記得當時在加拿大婚宴.
我說「我的時候還未到」.
現在真正的時候到了.

$^{241}$耶穌用祂的生命.
行了最後最重要的神蹟.
就是祂在十字架上死亡.
加拿大婚宴是耶穌行第一個神蹟.
而十字架就是耶穌用祂的生命.
行的最後一個神蹟.
兩者都是為了甚麼.
去彰顯上帝的榮耀,對吧.
但是….
兩者的結果可以是很不同.
在十字架上.
這個神蹟可以說是令耶穌母親的期望.
完全落空.
亦是為何要提到這個首尾呼應的原因.
因為在實踐上帝使命的過程中.
跟人的距離,期望的落差很大.
從瑪利亞的角度來看.
當然希望出人頭地,要祂的兒子.
做一個有能力,有面子.
可以幫助家族的猶太男人.
在加拿大婚宴已經說了.
祂期望耶穌出手.
幫親友做一些事,去挽回面子.
這個期望其實非常合理.
很正常,就像今天有很多父母.
都希望自己的子女爭氣,對吧.
可以為自己爭光,為家庭爭光.
或者我們也很想自己的子女.
有些時候不像自己.
走回自己的舊路.
或者反過來,希望他們達到.
一個自己無法達到的高度.
雖然我自己不是很高.
但如果我有個兒子.
我也很希望他喜歡打籃球.
我也很希望他可以做職業籃球員.
雖然我知道可能性不大.
但我也希望他有一天進入NBA就好了.
對吧?父母無論如何都會有期望.
這是很正常的….

$^{281}$其實反過來,我們做子女的時候.
在長大過程中.
很多時候我們都沒有照顧父母的期望.
去長大成他們期望的樣式.
雖然很多時候.
我們都不想媽媽傷心,失望.
我們很盡力去滿足媽媽的期望.
繼承她一些性格的特質.
但我們無法100\%達到她的期望.
部分原因是因為….
可能是我們反叛,對吧.
可能是因為我們貪玩.
媽媽叫我走左邊,我就走右邊.
覺得自己選擇自己的路.
自己才是老闆,對吧.
無論如何都很想有條路.
是我自己做到的事.
我自己能夠創造自己的世界.
當我們受到媽媽的期望時.
我們都會受到頗大的壓力.
不是因為媽媽的期望錯了.
而是因為媽媽很多時候.
都是我們最主要的照顧者.
我們的情感世界.
很大程度受到媽媽的看法影響.
所以在模型中.
一方面我們很想滿足這些期望.
承繼一些媽媽好的特質.
但同時也很想有自己選擇的可能性.
還有一個部分原因.
是因為上帝給我們每個人.
都有很獨特的使命.
每個人都有自己的照明.
人生不只是我們想怎樣就怎樣.
而是上帝在我們每個人的生命中.
在他創造的生命中.
他有獨特的照明.
上帝創造人很特別.
媽媽十月懷胎.
很自然就會跟自己的子女.

$^{321}$有一個很獨特的連繫.
有時我都會覺得我是媽媽的一部分.
或者有些媽媽都會覺得.
子女是自己的一部分.
甚至是一個小版本的自己,對吧.
有些媽媽很希望自己的子女.
不再重蹈覆轍.
重複自己走過,走錯的路.
所以希望為兒女指明方向.
你跟著這條路走就行了.
保證你沒有錯.
其實耶穌沒有完全否定媽媽的期望.
祂在加拿大行神蹟時.
依然照著媽媽的吩咐走.
但祂在對話中清楚提醒.
「我的時候未到」.
意思是.
「媽媽,你的期望是沒錯的」.
「但你要準備接受」.
「神在我的生命計劃中」.
「不會只圍繞家族的面子」.
「而是一些超越了你的想像」.
「是一些用我們意料之外的路」.
「十字架去彰顯上帝的榮耀」.
雖然耶穌做了這個神蹟.
但祂沒有像母親的期望.
只挽回家族的榮耀.
而是轉化了人對信仰的看法.
來到十字架,耶穌再叫婦人.
就是要提醒母親.
或已經準備成為人母的婦女.
在下面的聽眾.
你子女的人生的道路.
最終不是由你安排的.
亦不是由你的期望主導.
而是完全跟著上帝的旨意.
作為母親,最難的就是放手.
讓子女走一條.
跟自己心目中完全期望不同的路.
但可能這條路.

$^{361}$是上帝選給她的一條道路.
其實這一段.
耶穌跟母親的對話之前有一個處境.
就是24節.
「者要應驗經上記著說」.
「他們分了我的外衣」.
「為我的內衣抽籤」.
這是引用自詩篇22篇.
通常我們看新約.
經上記著說或應驗經上記話.
是指舊約.
而且它不只是說一世段經文.
而是說詩篇22篇本身的處境.
「他們分了我的外衣」.
「為我的內衣抽籤」.
不是說耶穌沒有衣著那麼簡單.
而是說耶穌成為了.
詩人所描述的二人.
他成為了上帝榮耀的成就者.
所以他選擇自己落入受苦的處境.
這件事跟瑪利亞對於榮耀的看法.
有很大的不同.
母親瑪利亞期待甚麼.
她的榮耀是覺得.
她的兒子可以幫加拿大婚宴的主人.
挽回面子.
她期望耶穌可以成為.
一個顧家典型的猶太男人.
萬事以大局為重.
懂得禮數,懂得禮儀招呼人.
更好的是甚麼.
她可以做堂會的拉比.
教書,做教授.
之後成家立室,繁衍後代.
這就是瑪利亞的期望.
成為一個很顯赫的家族.
為十二支派,為自己的支派爭取榮耀.
在十二支派當最標青的支派.
明顯,母親瑪利亞.
是一個很典型的猶太母親.

$^{401}$對子女的期望非常正常.
但耶穌選擇了走上帝呼召的路.
他不只是純粹跟著母親.
要她走她期望走的路.
所以我們回答了第一個問題.
為何耶穌要叫母親「夫人」.
因為這裡是一個首尾副營.
他提升在場的人.
耶穌為了上帝的榮耀.
走了一條跟母親所期望.
完全不同的道路.
回答第一個問題後.
我們去第二個問題.
為何耶穌要把母親交給門徒.
耶穌在十字架上.
跟母親說的最後一句話就是.
「看看你的兒子」.
又對門徒說「這是你的媽媽」.
如果你有留意.
耶穌不是站在平地跟他們說話.
耶穌是怎樣的.
掛在十字架上.
耶穌已經被掛在十字架上.
這表示耶穌在實踐上帝榮耀的同時.
他沒有忘記自己的母親.
在最辛苦,最呼吸不來的光景中.
他沒有忘記.
把自己的母親交給他最信任的門徒.
雖然他沒有做到母親的期望.
把血脈傳承.
但耶穌作為兒子.
一個人類版本的耶穌.
也盡了最大的努力去愛他的母親.
顧及母親的需要.
他希望母親最後成為信道群體中的一份子.
跟他最愛的門徒一起生活.
雖然耶穌似乎做不到母親瑪利亞的期望.
但他實際上在他生命最後的一口氣.
盡了做兒子的責任.
耶穌在十字架上受盡痛苦之際.

$^{441}$清楚地把照顧母親的責任.
交給所愛的門徒.
是出於愛和孝順.
不知道大家信主時有否遇到這個情況.
信主之後,媽媽開始擔心.
擔心甚麼呢.
就是教會的奉獻會否比家用多.
有這個情況嗎.
回教會的時間會否比回家的時間多.
我猜這是普遍非基督徒母親.
對基督徒子女的擔心.
她會勸我相信就好,不要太沉迷.
是嗎.
而基督徒父母可能也有少許掙扎.
當子女不僅留在回教會的層次.
他們開始認真.
認真考慮是否要為上帝做點事.
是嗎.
是否要尋求上帝給他們呼召時.
可能一方面,基督徒父母很安心.
因為他們很愛主,他們很棒.
但同時也擔心他們太愛主.
擔心他們會否太投入.
把所有東西都奉獻給上帝.
會有吧.
大家也放心.
耶穌的榜樣正在處理兩難的局面.
就是耶穌為了上帝的榮耀.
走了一條跟母親期望完全不同的路.
但祂也沒有拿上帝當擋箭牌.
你留意一下.
祂沒有說「我奉獻給上帝」.
「所以媽媽,我不理會你了」.
「再見」.
祂沒有.
祂在最後的一口氣.
記得把母親交給門徒照顧.
成為初期教會的一分子.
所以回答第二個問題.
就是耶穌臨近死亡.

$^{481}$為何要把母親交給門徒.
原因是祂沒有忘記.
祂在地上作為兒子的責任.
在完成上帝交託祂的呼召的同時.
祂也完成了地上的責任.
所以各位迷信主的父母.
或基督徒的父母,不用怕.
即使你的兒子很愛主.
很沉迷耶穌.
你也可以很放心地把他交託給上帝.
因為耶穌自己就做了一個例子.
照顧父母也是呼召其中的一部分.
一個很重要的一部分.
不像坊間對於基督徒的誤解.
就是你信了主,就沒有父母.
剛才也提到.
我媽媽從小到大都很疼愛我.
但我跟大家分享過.
我小時候有一個病.
一直都很困擾我.
亦令家裡有很大的負擔.
但媽媽依然跟我爸爸一起四處奔波.
找不同的醫生,希望能治好我.
或是減輕我的痛楚.
她經常教導我.
當時她是迷信主,不是基督徒.
她說她覺得人生最想要兩件事.
第一件事就是要健康.
第二件事,你猜猜是甚麼.
快樂?不是.
是錢,最重要是有錢.
人有健康,但沒有錢.
人就會為了錢而出賣健康.
人有錢而沒有健康.
其實也沒有用,對吧.
但這件事在我身上也沒有.
第一,我沒有健康.
因為我有一個很嚴重的病.
第二,我沒有錢.
所以她很希望我努力讀書.

$^{521}$改善我跟家人的生活.
我媽媽也很身體力行.
很努力賺錢.
很希望避免我們小時候很貧困.
沒有東西吃的狀態.
再次出現在我們身上.
然後上帝很神奇地醫治了我的病.
是一個很特別的神蹟.
給了我健康.
我媽媽很自然地想.
有了健康又如何?當然是賺錢.
她還跟我說「你努力讀書」.
期望我努力讀書.
找一份穩定的工作,成家立室.
我康復後也很乖巧.
我也很努力讀書.
依照媽媽的期望.
找一份專業的工作.
同時,我不只想賺錢.
所以我選了一份有意義.
又能賺錢,又能幫助別人的工作.
我選擇了做社工.
既專業.
又能賺錢,又能幫助別人.
不錯,我滿足了媽媽的期望.
我不是完全照著媽媽的期望.
我走了一條自己選擇的路.
我媽媽也說「也算滿意吧」.
始終社工收入算是穩定.
那時候我以為我的人生會頗平穩.
誰知道上帝的呼召臨到.
在我即將成為註冊社工的時候.
上帝就在180度將我的人生改變.
明明剛讀完社工,可以出去賺錢.
轉身就進了神學院.
然後我媽媽就覺得很驚訝.
「為甚麼?為甚麼要這樣做」.
「為甚麼讀完社工」.
「你欠了一大筆錢」.
「你借錢讀書吧」.

$^{561}$「你還要再花幾年錢讀神學」.
「你還要薪金低不少」.
那時候我媽媽還未相信.
她一直不明白.
為甚麼要相信主要相信得這麼沉迷.
為甚麼.
平時你奉獻一點就算了.
為甚麼整個人都要奉獻.
這就是有時當上帝的使命臨到.
有時可能真的會跟我們父母的期望.
有很大的落差.
我以為他做不到10\%.
「好,5\%也算吧」.
以為他做不到2\%,算吧.
1\%也好.
誰知原來真的一點也沒有.
但上帝沒有忽略人在世的責任.
耶穌在世的時候.
祂也很盡力地完成祂的責任.
把母親交託給門徒.
雖然做傳代人不是很鬆動.
但我也有盡力照顧我媽媽的需要.
每個母親節.
我都提前為媽媽送一份禮物.
她生日的時候,我也寄錢給她.
希望她可以吃飯吃得開心.
每個星期,我都會打電話給媽媽.
在上帝交託給我的工作之餘.
在我的能力之內.
我也盡力照顧我的母親.
是因為上帝給我們這個榜樣.
祂沒有拿上帝當擋箭牌.
祂沒有說「我相信了耶穌」.
「媽媽,我先說再見吧」.
這是人在世的責任.
所以耶穌也是這樣.
祂寫在十字架上.
也想盡一個兒子的責任.
把母親交託給最愛的門徒.
總結一下,這個時間差不多了.

$^{601}$我們回答兩個問題.
為何耶穌要叫母親「婦人」.
就是耶穌跟母親的對話.
最後是一個首尾呼應.
祂想喚醒人對子女的期望.
當上帝的使命臨到的時候.
很多時候父母對子女的期望.
有很大的落差.
但在上帝的角度.
一個最好的母親是怎樣的.
就是明白子女有自己的使命.
上帝給我們每一個人.
有一個很獨特的使命.
可能有時在完成使命的過程中.
有時會受到我們不想受的痛楚.
或者期待的苦難.
但這些在世人眼中.
你上了教會那麼多時間.
奉獻了那麼多.
是很愚蠢,沒有經驗,吃虧.
但我們依然放手.
讓他們完成和實踐上帝給他們的使命.
第二,為何耶穌把母親交給門徒.
就是耶穌用她自己做榜樣.
為我們展示.
即使在完成使命的過程中.
她也沒有拿上帝當擋箭牌.
她沒有逃避自己在地上的責任.
她也懂得照顧母親的需要.
今天是母親節.
當我們再讀這段經文時.
最重要的是讓十字架重新調整.
我們如何做父母,如何做兒女.
如何做主內的家人.
求主幫助我們每一個母親,每一個門徒.
我們都走在上帝的路上.
我們一起祈禱.
在實踐上帝榮耀的時候.
你也沒有忘記自己的母親.
在最辛苦的時刻.

$^{641}$最無法呼吸的光景中.
你也沒有忘記把自己的母親.
交託給自己最信任的門徒.
雖然有時上帝讓我們子女走的路.
的確跟我們的期望很不同.
但你盡了最大的努力愛你的母親.
顧及你的需要.
讓我們這群門徒去追隨你.
祈求你讓我們有清晰的方向.
繼續做你的使命.
我感到高興,耶穌,得勝命求.
不知道今天有否新來的朋友.
\newpage



\section{出埃及記 20:1-26}
\label{sec:y_bUxa4hgT8}
\textbf{Stone (Exod 20) - Rev. Mark Chan 石板 (出埃及記20\_1-26) - 陳崇基牧師}
\newline
\newline
連結: \href{https://youtube.com/watch?v=y_bUxa4hgT8}{\texttt{https://youtube.com/watch?v=y\_bUxa4hgT8}} ~~~~ 語音日期: 2026-05-24
\newline
\newline
\hyperref[sec:HpXR4dbZvbk]{\small{< < < PREV SERMON < < <}}
~
\hyperref[sec:index_chronic]{\small{[返順時目]}}
~
\hyperref[sec:index_scriptual]{\small{[返順卷目]}}
~
\hyperref[sec:OKPQ6OtPkus]{\small{> > > NEXT SERMON > > >}}
\newline
\newline
出埃及記 20:1-26
\newline
\begin{longtable}{cl}
\hline
\hline
章節 & 經文 (和合本修訂版)\\
\hline
20:1 & \begin{tabularx}{0.7\textwidth}{X} 神吩咐這一切的話，說： \end{tabularx} \\ \\ \relax
20:2 & \begin{tabularx}{0.7\textwidth}{X} 「我是耶和華－你的神，曾將你從埃及地為奴之家領出來。 \end{tabularx} \\ \\ \relax
20:3 & \begin{tabularx}{0.7\textwidth}{X} 「除了我以外，你不可有別的神。 \end{tabularx} \\ \\ \relax
20:4 & \begin{tabularx}{0.7\textwidth}{X} 「不可為自己雕刻偶像，也不可做甚麼形像，彷彿上天、下地和地底下水中的百物。 \end{tabularx} \\ \\ \relax
20:5 & \begin{tabularx}{0.7\textwidth}{X} 不可跪拜那些像，也不可事奉它們，因為我耶和華---你的神是忌邪的神。恨我的，我必懲罰他們的罪，自父及子，直到三、四代； \end{tabularx} \\ \\ \relax
20:6 & \begin{tabularx}{0.7\textwidth}{X} 愛我，守我誡命的，我必向他們施慈愛，直到千代。 \end{tabularx} \\ \\ \relax
20:7 & \begin{tabularx}{0.7\textwidth}{X} 「不可妄稱耶和華－你神的名，因為妄稱耶和華名的，耶和華必不以他為無罪。 \end{tabularx} \\ \\ \relax
20:8 & \begin{tabularx}{0.7\textwidth}{X} 「當記念安息日，守為聖日。 \end{tabularx} \\ \\ \relax
20:9 & \begin{tabularx}{0.7\textwidth}{X} 六日要勞碌做你一切的工， \end{tabularx} \\ \\ \relax
20:10 & \begin{tabularx}{0.7\textwidth}{X} 但第七日是向耶和華---你的神當守的安息日。這一日你和你的兒女、奴僕、婢女、牲畜，以及你城裡寄居的客旅，都不可做任何的工。 \end{tabularx} \\ \\ \relax
20:11 & \begin{tabularx}{0.7\textwidth}{X} 因為六日之內，耶和華造天、地、海和其中的萬物，第七日就安息了；所以耶和華賜福與安息日，定為聖日。 \end{tabularx} \\ \\ \relax
20:12 & \begin{tabularx}{0.7\textwidth}{X} 「當孝敬父母，使你的日子在耶和華－你神所賜你的地上得以長久。 \end{tabularx} \\ \\ \relax
20:13 & \begin{tabularx}{0.7\textwidth}{X} 「不可殺人。 \end{tabularx} \\ \\ \relax
20:14 & \begin{tabularx}{0.7\textwidth}{X} 「不可姦淫。 \end{tabularx} \\ \\ \relax
20:15 & \begin{tabularx}{0.7\textwidth}{X} 「不可偷盜。 \end{tabularx} \\ \\ \relax
20:16 & \begin{tabularx}{0.7\textwidth}{X} 「不可做假見證陷害你的鄰舍。 \end{tabularx} \\ \\ \relax
20:17 & \begin{tabularx}{0.7\textwidth}{X} 「不可貪戀你鄰舍的房屋；不可貪戀你鄰舍的妻子、奴僕、婢女、牛驢，以及他一切所有的。」 \end{tabularx} \\ \\ \relax
20:18 & \begin{tabularx}{0.7\textwidth}{X} 眾百姓見雷轟、閃電、角聲、山上冒煙，百姓看見就都戰抖，遠遠站著。 \end{tabularx} \\ \\ \relax
20:19 & \begin{tabularx}{0.7\textwidth}{X} 他們對摩西說：「請你向我們說話，我們必聽；不要讓神向我們說話，免得我們死亡。」 \end{tabularx} \\ \\ \relax
20:20 & \begin{tabularx}{0.7\textwidth}{X} 摩西對百姓說：「不要害怕；因為神降臨是要考驗你們，要你們敬畏他，不致犯罪。」 \end{tabularx} \\ \\ \relax
20:21 & \begin{tabularx}{0.7\textwidth}{X} 於是百姓遠遠站著，但摩西卻挨近神所在的幽暗中。 \end{tabularx} \\ \\ \relax
20:22 & \begin{tabularx}{0.7\textwidth}{X} 耶和華對摩西說：「你要向以色列人這樣說：『你們親自看見我從天上向你們說話了。 \end{tabularx} \\ \\ \relax
20:23 & \begin{tabularx}{0.7\textwidth}{X} 你們不可為我製造偶像，不可為自己造任何金銀的神像。 \end{tabularx} \\ \\ \relax
20:24 & \begin{tabularx}{0.7\textwidth}{X} 你要為我築一座土壇，在上面獻牛羊為燔祭和平安祭。凡在我叫你記念我名的地方，我必到那裡賜福給你。 \end{tabularx} \\ \\ \relax
20:25 & \begin{tabularx}{0.7\textwidth}{X} 你若為我築一座石壇，不可用鑿過的石頭，因為你在石頭上動了工具，就使壇污穢了。 \end{tabularx} \\ \\ \relax
20:26 & \begin{tabularx}{0.7\textwidth}{X} 你不可用臺階上我的壇，免得露出你的下體來。』」 \end{tabularx} \\ \\
[1ex]
\hline
\hline
\end{longtable}
$^{1}$Ok, this is the third part of my series for our 10th anniversary..
這次是我們第三個系列..
I call it Stone. 我們叫它石板..
We are going to talk about the Old Testament Law. 我們在講舊約的一些規章制度..
Now, whether you have a printed Bible or you have a digital Bible, you understand the Bible is a very very long book, right?.
你是有電子版還是有書,你都知道聖經是一個非常長的書..
And a huge part of the book of this Bible is the Law in the Old Testament. 那在舊約裡面很大的一個部分就是講我們的一個約束..
How many of you have heard about Ten Commandments? 有多少人聽過十件?.
Yeah, the Ten Commandments in the Bible in the Old Testament, you heard about it? 在舊約的十件..
Yeah, those are Law. 那些都是律法..
But as we modern people, we heard about the word Law, we have a negative feeling. 我們當聽到律法的時候,我們常常都覺得被冒犯了,都不高興了..
Do you like Law? 你們喜歡律法嗎?.
Probably not, right? 可能不是吧?.
Sometimes, right? Sometimes. 可能不是吧?.
I remember a while back I was, you know, finding parking in Chicago is quite a challenge, right? 在舊約找一個爬車的地方很難的..
A while back I was driving in Chicago, in Hyde Park, to be exact. 有一次呢,我在芝加哥Hyde Park的這個地區開車的話,我需要路邊找一個爬車的地方..
So as you do that, in street parking in Chicago, you pay attention to the signs, right? 那你在我們家街趴的時候,我們肯定要看路牌嘛..
So I was on the street in Hyde Park, and one side, they have this orange paper sign, you know, tied on the trees, right? 那我在樹上看到橘色的一個紙牌..
And then on this side, it says no parking on Monday. 那我在樹上看到橘色的一個牌子,掛在那裡..
They are doing street cleaning. 要一邊寫著禮拜一不準停車..
And then on the other side, you know, they also hang some signs, right? And there is no parking, the sign says no parking on Tuesday. 嗯,他們要洗街..
And that day was Monday. 那天是禮拜一..
And then I thought, very simple, okay, I find an empty spot, and that day was Monday, yes. 然後我就想,非常簡單,我找到一個空的地方,那天是星期一..
I find a parking space on one side of the street, so I park on the Tuesday side, and I'm safe. 我找到一條街道上的停車位,所以我停在週一那邊,我安全了..
So I parked the car and I left, and I came back a few hours later. 那我先講很簡單嘛,今天是禮拜一,那我應該趴在禮拜二那邊..
我終於千辛萬苦找到個位子在禮拜二那邊,應該沒有問題..
So I parked the car and I left, and I came back a few hours later, and there was an orange citation, one orange envelope on the windshield. 那幾個小時以後呢,我發現一個橘色的信封在我的車窗上面..
And then what happened was that when I was looking for a parking space, and I checked the signs, and I got confused. 那結果是我自己看多了兩邊的路牌以後,我自己覺得路牌..
So it was Monday, I should park on the which side? 禮拜一的話,我應該趴在哪一邊?.
Not parking on Tuesday side, right? 應該是趴在禮拜二不準趴車那邊..
That's what I thought I did. 那應該是這樣子想的,對吧?.
But I was wrong, I did not pay attention that I was actually parking on the no parking on Monday side. 但是我趴在了錯誤的那一邊,我趴在了禮拜一不準趴車那邊..
So I got that orange envelope, and my first reaction was that this is so unfair. 我一拿到橘色的信封,我就覺得這太不公平了..
And then you feel like the law is out to get me, you know? 這個法律是要用來整我的嗎?.
Have you ever got something like that before? 你有沒有這樣的感覺?.
Anyone? 任何人?.
拿到了一個信封以後,太不公平了..
The city of Chicago is evil. 芝加哥政府太糟糕了..
Or maybe I was confused, you know? 也許是我自己迷糊了..
But that's our modern conception or modern feeling about law. 這是我們大多數人對法律的感覺..

$^{41}$Now when we go to the Old Testament, we see commandments and rules, we sometimes read them the same way. 當我們看到舊約的時候,我們看到這樣的一些律法,我們是有相似的感覺..
So you feel like that, right? You open the Bible, you look at the laws, and there are some details, a lot of commandments, ten commandments, and you feel like God is, you know, like a traffic police, you know, okay? 這樣的感覺,對吧?你打開了詩書,看了法律,看了一些細節,有很多的命令,十個命令..
You feel like God is, you know, like a traffic police, you know, okay? 這樣的感覺,對吧?你感覺像是一名交通警察,你知道嗎?.
Always out to get you, you know? 總是要用法律來整你,你知道嗎?.
And then you feel like that, right? 當我們翻開舊約的時候,有些感覺,有些時候會想,神是不是就像那個交通督導員,.
就是要抓我們的錯誤,你們就被罰款了. And then you feel like that, right? 當我們翻開舊約的時候,有些感覺,有些時候會想,神是不是就像那個交通督導員,就是要抓我們的錯誤,你們就被罰款了..
But I want to suggest to you, maybe we have misread the sign. 但我想在這裡要大家想一想,是不是我們一直看錯了路牌呢?.
So I want to show you from the Bible, from the Old Testament itself, the law was never designed to catch us, but it's to show us a better way to live. 所以我想在這裡要大家想一想,是不是我們一直看錯了路牌呢?.
這也想告訴你的是,律法從來不需要抓我們的把柄,而是要指引我們的道路..
See, from this PowerPoint slide, you can see the common misreading or misunderstanding is that the Old Testament, you know, is full of law, rules and judgment. 從這個 PowerPoint 圖上可以看到,律法從來不需要抓我們的把柄,而是要指引我們的道路..
那我們常見有幾個誤解,就像我們 PowerPoint 下面所講的,就是舊約常常讓人覺得是關於律法,規條和審判的..
And then we flip to the pages on the New Testament, and we feel like, oh, it's all about grace and all about freedom and love. 然後我們翻到新約的頁面,我們覺得,哦,這全都是關於恩典,愛和自由..
But, you know, this is a misreading of the whole Bible. 但是,你們知道,這是整個讀書的錯誤..
這絕對是一種誤解..
So, I want to begin with the first point is that the first thing you need to understand about law is that it is of grace. 那,律法其實是關於恩典的律法..
Now, the truth is that, you know, the law was never the opposite of grace, the law was actually grace. 律法從來不是恩典的對立面,律法本身就是恩典..
Now, as some of you know, I'm leaving tonight, I'm flying to Turkey, and I'm leading a group to Turkey in the next two weeks. 現在,有些網友知道,我今天晚上會帶領一個團隊前往土耳其..
Now, you probably have some, you know, a concept that when I mention Turkey, what kind of Bible stories you can think of? 那可能如果你熟悉聖經的話,將我講到土耳其,大家會想到聖經的哪一個部分?.
Probably all in the New Testament, right? Paul's missionary journeys, and the seven churches of the book of Revelation. 可能都是關於新約的保羅的宣教旅程,啟示錄的七個教會等等..
But, you know, Turkey was actually the heartland of the Old Testament. 但是,土耳其也是舊約的核心地帶..
Because in the Old Testament time, the same place, Turkey was the homeland for ancient superpower called Hittite Kingdom. 因為在舊約時期,土耳其是古代世界最強大的帝國之一的故鄉,赫人帝國..
So, what the people have found in the capital of the ancient Hittite Kingdom was what we call a suzerain treaty. 當科普學家挖掘到這個帝國的遺址的時候,發現很多的泥板,其中的一個我們叫做宗教條約..
And then, these kind of treaties are actually treaty between a great nation or great king like Hittite King, the great king, between the great king and the small kingdom. 然後,這種條約是大國與附屬小國之間所立的條約..
And what is interesting is, every single this kind of document, it follows the same format. 每一種宗主的條約其實都用同樣的格式..
So, as you can see here, number one, the treaty begins with the introduction of the great king. 第一個部分,基本上是一個大王來用很長的篇幅來介紹自己有多偉大有多偉大..
So, after the preamble is the second part, which is what we call the historical introduction or prologue. 第二個部分我們叫做是歷史的序言..
And the great king will say something to the small kingdom and say, this is what I have done for you guys before. 那這個很大的偉大的國王就會說,這是我已經為你做的事情..
Perhaps some years before the treaty, the great king has saved this small nation from military attack. 可能很久很久之前你們被人攻打的時候,我來拯救了你們..
So, the historical prologue is to remind the small kingdom that I have been good to you. 那這個部分基本上是用來提醒他們我曾經對你們做過恩典的事情..
And then the number three part, part three is the stipulations or the laws. 第三個部分就是關於律法條款..
So, the stipulations or the laws are saying because I have been gracious to you, and that's why you have to live this way in response. 那這個部分基本上是說因為我如此的恩典的待遇,你應該如何的去回應我如何去生活..
And then number four is the blessing and the curses. 第四個部分關於祝福與詛咒..
So, the treaty says if you obey me, if you are faithful to me, then I will bless you, I will be good to you. 那它說忠誠,我就會給你祝福..
But if you are rebellious, then there will be consequences. 那如果你叛逆我的話,也會帶來後果的..
Now, you look at this format of the Suzerain Treaty, you find it very interesting because it's very similar to the Old Testament law. 那如果你看這個宗主條約的話,你會發現跟我們舊約非常相像..
So, when we go to the Ten Commandments, for example, Exodus 20, 如果你看到埃及紀第20章,我們的十件,.
and before any single so-called commandment begins, it gives you these verses. 那首先在什麼其他東西都不用講的同時,它首先跟你講這句話..
So, let me read this. I am the Lord your God, 我是耶和華你的神, who brought you out of Egypt, out of the land of slavery. 曾將你從埃及帝為奴之家領出來..
Now, only then, after this, do we find the commandments or Ten Commandments. 那這個之後,我們纔會發現十約..
And then the Ten Commandments do not begin with obey me, so I will love you. 十約並不是開始的時候怎麼說,你順從我,然後我再愛你,不是這樣開始的..

$^{81}$It's the other way around. It begins with I have already loved you, I have already rescued you, and here is how you live as my people now. 而是我已經拯救了你,現在這是你做我子民的生活方式..
So, grace comes first, the rescue came first, and then the law came second. The law is flowing out of grace. 拯救在前,律法在後,律法是從恩典流出的..
So, you have to understand, you have to erase the wrong impression. Law is given to you, is to trick you, is to punish you. 你要把這個否定的那個想法,關於律法的想法給去掉..
So, grace first, and then the law is to teach you and guide you how you live. 恩典是先在律法之前出現,然後律法來指導你如何生活..
And then the second aspect of the law in the Old Testament especially is actually wisdom. 第二個方面是關於智慧的..
So, the law reveals the heart of God, the wisdom of God. 律法揭示了神的心腸..
Now, what do the specific laws actually look like? 那些具體的律法究竟是怎麼樣的呢?.
Now, you may have some impression, but I want to show you two examples. 我想給你兩個例子..
So, besides the Ten Commandments, the popular, you know, everybody knows, but I want to give you two samples from some of the laws in the Old Testament. 除了十節以外,我想給你兩個例子..
So, let's take a look at this one, from Deuteronomy 22. 申命記二十二章第八節..
So, pay attention as I read it. So, pay attention as I read it. 小姐,請你們注意..
So, Moses is trying, is to giving the law to the people of God, to Israel, he says, when you build a new house, make a parapet around your roof, so that you may not bring guilt of bloodshed on your house if someone falls from the roof. 你建造新房子,要在房頂上安欄杆,.
免得有人從房上掉下來,流血的罪就歸你家了. So, that's interesting, this law is actually building code. 這很奇怪,這是關於如何建造房子的..
It tells you how to build a house. 你怎麼樣建造一個房子..
Now, in ancient Israel, the house had flat roof. 在古代的以色列,屋頂是平坦的..
So, people can go up to the roof and do something to work there. So, God is saying, you know, when you build a house, you have to build railing, you know, some sort of barrier, so that, you know, perhaps, you know, a child goes up there, and he will fall from the roof. 因為屋頂是平坦的話,人們就可以爬上去做各樣的事情..
那神說要安欄杆,可能就讓孩子不會墜落..
Yeah, anyone has done renovation in your house lately? 有沒有人裝修的,在家裡?.
Okay, a lot, right? And you have to watch out for safety, right? 你們要注意安全,對?.
So, this law is actually a safety law. 那這個關於安全的律法..
But the law, behind the law, reveals the wisdom or the heart of God. 但是律法告訴了我們..
So, it's like the heavenly Father saying, you know, life is sacred, life is so special, and God wants to protect life, and so no one, even when no one is watching. 這是一位父親講的話,生命是神聖的,要保護它,即便沒有人在看..
So, you see that the law is trying to reveal to you the heart of God, the wisdom of God. 這關於神的智慧以及神的心腸..
Now, the second example I want to show you is from Leviticus, chapter 19. 這是下一個是關於歷位紀十九章第十九節..
Do not reap the very edges of your field, or gather the gleanings of your harvest, leave them for the poor and for the foreigner. 你們收割莊稼不可割盡田角,也不可拾取所遺落的,要留給窮人和寄居的外人..
So, you may say this is a harvest law, a harvest code. 可能關於是如何收割的一個規定..
So, your harvest is your entire livelihood, right? 你的收成是你全年的生計..
Whether you have wheat or barley or whatever, God is trying to say to you in this law is saying, when you harvest your wheat and barley, your produce, you have to leave the edges out. 神說你收割莊稼的時候,你必須要留下田角..
So, there are some, you know, some of your production, your food items, you have to leave behind for the widows, for the foreigners. 你要為寡婦,孤兒,陌生人留下可以喫的..
So, you see, this law is kind of a very tiny thing, you know, a very tiny thing about, you know, how you live. 這些都是關於生活的很小的一些關鍵..
But as you can see, you know, when you harvest, you don't harvest every single thing. 但是你收成不是收聖所有的東西..
So, you leave it to the weak people or the poor in the society, so they come after you, they can pick up those food. 你還要去照顧社會上面的底層的人民..
And this is the heart of God, this is justice, it's building kind of building a society, economic system, so that, you know, when you look at this law, it's actually the value from God, behind every value is that, you know, God is trying to protect the weak. 這是建立經濟制度本身之中的公益,每一條律法背後都有一個價值觀,而每一個價值觀背後都有一個神要保護的人..
And this is exactly what I'm saying, you know, the wisdom of God, the heart of God is written in everyday life of the people of God. 這就是神的智慧寫進了日常生活之中..
So, the law is actually grace, and the law is actually revealing the heart of God, the wisdom of God, and number three, the law were meant to be missional, to have the world to see how you live. 律法也是關於使命的..
So, we have to ask the question, why the law was given to Israel, to the people of God, what was the ultimate point? 那麼為什麼要把這些律法賜給以色列呢?它最終的目的是什麼呢?.
Deuteronomy 4 says, again Moses is talking to the people of God before they enter the land of Canaan, and he says, observe them, observe those laws that I have preached to you, I have told you, observe them carefully, this will show your wisdom to the nations, to the pagan nations, and they will say, surely, this is how they live, this great nation is a wise and understanding people..
So, the laws in the Old Testament are actually public display of your life, you know, kind of a new way of life, a new way of life that is marked by justice, by the law. 摩西在以進入以色列,進入應許之地之前說,你們要謹守遵行,這就是你們在萬民眼前的智慧,聰明,他們聽見這一切的律例必說,這大國的人真是有智慧,有聰明..
So, the laws in the Old Testament are actually public display of your life, you know, kind of a new way of life, a new way of life that is marked by justice, by generosity, by grace to the surrounding nations, and the surrounding nations will get curious about, hey, this is how you live, you live different than us, and who is your God, and we want to know Him as well. 摩西在以進入以色列,進入應許之地之前說,你們要謹守遵行,這就是你們在萬民眼前的智慧,聰明,他們聽見這一切的律例必說,這大國的人真是有智慧,有聰明..
所以,律法是公開的見證,一種公義恩慈的生活方式,使周圍的國家都會問,他們的神是誰呢,我想認識他..

$^{121}$今天是我們的,這是我們的護照,我們透過日常的生活的智力,見證神的美善,誠實,慷慨,以及我們對待那些無法回報我們的人的一些方式..
這不是一個新的策略,如果我們在古代發現,這不是一個新的策略,而是聖經裡面最古老的一個生活方式..
實踐是神親自用手指把它刻在石頭上的..
《出埃及記》31章,它說,他將兩塊法板交給摩西,是神用指頭寫的石板..
這不是一份可以隨便修改的備忘錄,而是刻在石頭上的盟約,永遠的長存..
那我們就是為什麼我要你帶一個這樣的小瓶子回家呢?.
我們這個小瓶子裡面裝著沙子..
在第一個部分,這個講到第二部分,我們講到了神對亞伯拉罕的應許,神用亞伯拉罕去祝福萬國的人..
第二個部分,我給你們什麼了呢?是一塊拼圖..
上一次我們跟你說,亞伯拉罕這一堆人全部都是破碎的人..
但是上帝選擇用他們去祝福萬國..
他們的神是我們一樣的神..
今天我要請你們把一些石頭帶回家..
希望你們今天結束的時候,拿一些石頭可以放在你的袋子裡面,帶回家..
當你握住這些石頭的時候,請記住,神把律法刻在石頭上,因為祂的信息是永恆的..
這石頭提醒我們,神說我已經拯救了你,我如此的愛你,這是你應當如何生活的方式..
把它放在家裡某個你可以看到的地方,讓它每天問你一個問題..
我的生活方式有沒有讓人對我的神感到好奇呢?.
神不是要我們完美地活著,但是神要求我們忠實地活著..
這是一個邀請,活得如此的有智慧,如此的慷慨,讓一個旁觀者的世界渴望認識這一切背後的善意..
要知道,實踐是從拯救開始的,一直是如此,我們的福音也是如此..
請幫助我們接受禰的律法,不是覺得它是一種威脅,而是一種禮物..
願我們的生活方式使人渴望認識禰..
奉耶穌的名求,給你們..
我今晚要飛,請祈禱我和我的團隊,我們有17個人要去土耳其..
我會盡力更新您的信息,這取決於互聯網,我會在土耳其錄製一些視頻..
下個禮拜見!.
\newpage



\section{出埃及記 20:1-26}
\label{sec:OKPQ6OtPkus}
\textbf{石板 (出埃及記20\_1-26) - 陳崇基牧師}
\newline
\newline
連結: \href{https://youtube.com/watch?v=OKPQ6OtPkus}{\texttt{https://youtube.com/watch?v=OKPQ6OtPkus}} ~~~~ 語音日期: 2026-05-24
\newline
\newline
\hyperref[sec:y_bUxa4hgT8]{\small{< < < PREV SERMON < < <}}
~
\hyperref[sec:index_chronic]{\small{[返順時目]}}
~
\hyperref[sec:index_scriptual]{\small{[返順卷目]}}
~
\hyperref[sec:wv60KUfih4o]{\small{> > > NEXT SERMON > > >}}
\newline
\newline
出埃及記 20:1-26
\newline
\begin{longtable}{cl}
\hline
\hline
章節 & 經文 (和合本修訂版)\\
\hline
20:1 & \begin{tabularx}{0.7\textwidth}{X} 神吩咐這一切的話，說： \end{tabularx} \\ \\ \relax
20:2 & \begin{tabularx}{0.7\textwidth}{X} 「我是耶和華－你的神，曾將你從埃及地為奴之家領出來。 \end{tabularx} \\ \\ \relax
20:3 & \begin{tabularx}{0.7\textwidth}{X} 「除了我以外，你不可有別的神。 \end{tabularx} \\ \\ \relax
20:4 & \begin{tabularx}{0.7\textwidth}{X} 「不可為自己雕刻偶像，也不可做甚麼形像，彷彿上天、下地和地底下水中的百物。 \end{tabularx} \\ \\ \relax
20:5 & \begin{tabularx}{0.7\textwidth}{X} 不可跪拜那些像，也不可事奉它們，因為我耶和華---你的神是忌邪的神。恨我的，我必懲罰他們的罪，自父及子，直到三、四代； \end{tabularx} \\ \\ \relax
20:6 & \begin{tabularx}{0.7\textwidth}{X} 愛我，守我誡命的，我必向他們施慈愛，直到千代。 \end{tabularx} \\ \\ \relax
20:7 & \begin{tabularx}{0.7\textwidth}{X} 「不可妄稱耶和華－你神的名，因為妄稱耶和華名的，耶和華必不以他為無罪。 \end{tabularx} \\ \\ \relax
20:8 & \begin{tabularx}{0.7\textwidth}{X} 「當記念安息日，守為聖日。 \end{tabularx} \\ \\ \relax
20:9 & \begin{tabularx}{0.7\textwidth}{X} 六日要勞碌做你一切的工， \end{tabularx} \\ \\ \relax
20:10 & \begin{tabularx}{0.7\textwidth}{X} 但第七日是向耶和華---你的神當守的安息日。這一日你和你的兒女、奴僕、婢女、牲畜，以及你城裡寄居的客旅，都不可做任何的工。 \end{tabularx} \\ \\ \relax
20:11 & \begin{tabularx}{0.7\textwidth}{X} 因為六日之內，耶和華造天、地、海和其中的萬物，第七日就安息了；所以耶和華賜福與安息日，定為聖日。 \end{tabularx} \\ \\ \relax
20:12 & \begin{tabularx}{0.7\textwidth}{X} 「當孝敬父母，使你的日子在耶和華－你神所賜你的地上得以長久。 \end{tabularx} \\ \\ \relax
20:13 & \begin{tabularx}{0.7\textwidth}{X} 「不可殺人。 \end{tabularx} \\ \\ \relax
20:14 & \begin{tabularx}{0.7\textwidth}{X} 「不可姦淫。 \end{tabularx} \\ \\ \relax
20:15 & \begin{tabularx}{0.7\textwidth}{X} 「不可偷盜。 \end{tabularx} \\ \\ \relax
20:16 & \begin{tabularx}{0.7\textwidth}{X} 「不可做假見證陷害你的鄰舍。 \end{tabularx} \\ \\ \relax
20:17 & \begin{tabularx}{0.7\textwidth}{X} 「不可貪戀你鄰舍的房屋；不可貪戀你鄰舍的妻子、奴僕、婢女、牛驢，以及他一切所有的。」 \end{tabularx} \\ \\ \relax
20:18 & \begin{tabularx}{0.7\textwidth}{X} 眾百姓見雷轟、閃電、角聲、山上冒煙，百姓看見就都戰抖，遠遠站著。 \end{tabularx} \\ \\ \relax
20:19 & \begin{tabularx}{0.7\textwidth}{X} 他們對摩西說：「請你向我們說話，我們必聽；不要讓神向我們說話，免得我們死亡。」 \end{tabularx} \\ \\ \relax
20:20 & \begin{tabularx}{0.7\textwidth}{X} 摩西對百姓說：「不要害怕；因為神降臨是要考驗你們，要你們敬畏他，不致犯罪。」 \end{tabularx} \\ \\ \relax
20:21 & \begin{tabularx}{0.7\textwidth}{X} 於是百姓遠遠站著，但摩西卻挨近神所在的幽暗中。 \end{tabularx} \\ \\ \relax
20:22 & \begin{tabularx}{0.7\textwidth}{X} 耶和華對摩西說：「你要向以色列人這樣說：『你們親自看見我從天上向你們說話了。 \end{tabularx} \\ \\ \relax
20:23 & \begin{tabularx}{0.7\textwidth}{X} 你們不可為我製造偶像，不可為自己造任何金銀的神像。 \end{tabularx} \\ \\ \relax
20:24 & \begin{tabularx}{0.7\textwidth}{X} 你要為我築一座土壇，在上面獻牛羊為燔祭和平安祭。凡在我叫你記念我名的地方，我必到那裡賜福給你。 \end{tabularx} \\ \\ \relax
20:25 & \begin{tabularx}{0.7\textwidth}{X} 你若為我築一座石壇，不可用鑿過的石頭，因為你在石頭上動了工具，就使壇污穢了。 \end{tabularx} \\ \\ \relax
20:26 & \begin{tabularx}{0.7\textwidth}{X} 你不可用臺階上我的壇，免得露出你的下體來。』」 \end{tabularx} \\ \\
[1ex]
\hline
\hline
\end{longtable}
$^{1}$好 今次我們繼續說這個罌.
如果你是新來的朋友 還沒有的話.
待會出去的時候 我們送一個給你 好嗎.
一個玻璃罌 裡面有不同的物件.
就是我每一次說到新的一年.
我們十週年這個系列的時候.
我們就會從聖經由頭到尾.
我們看看這位的上帝是一個怎樣的上帝.
你記得我第一次說的時候.
今次是第三次.
第一次說的時候 我說上帝是一個怎樣的上帝.
God on a mission.
上帝是由頭到尾.
他不單止創造了這個世界.
知道人類是叛逆他也好.
但是他沒有放棄過這個世界.
他繼續on a mission 拯救這個世界.
這個第一次說的時候就是說.
這個世界這麼砸 這麼不妙 這麼破碎.
上帝做了什麼.
沒有放棄 他就叫了Abraham出來.
叫了亞伯拉罕出來.
藉著亞伯拉罕 以撒雅各.
成為一個家 成為一個國.
成為一個子民.
這個上帝的子民就是在上帝的計劃裡面.
去做上帝的工具.
以致上帝的祝福能夠流向萬國.
這是第一次.
上一次我們說到族長的故事.
就是說這些族長 亞伯拉罕 以撒雅各.
他們不是什麼英雄人物.
我們不需要將他們美化.
他們都是在掙扎又跌又碰的裡面 跟隨上帝.
去承接這個呼召 拿起這個使命.
去使萬國能夠得福.
這是整本《創世記》的結束.
《創世記》就是告訴我們.
上帝藉著這些弱小的人.
他要去成就他自己的心意.

$^{41}$所以我們每一個上帝的子民.
我們有一個identity 一個名分.
這個名分就是上帝賜給我們.
作為他自己的people.
我們都是on the mission.
因為上帝就是on the mission.
這次我們繼續trace整個聖經的大故事.
完了《創世記》 再到出《埃及紀》.
你記得《創世記》最後的時候發生什麼事嗎.
雅各他們一家去了哪裡.
去了埃及.
然後在埃及他們一代一代受苦.
然後上帝聽見他們的聲音.
叫了摩西出來.
然後拯救他們.
用恩典拯救了他們.
他們就能夠出埃及.
經過曠野去到西大山.
他們真正成為了上帝的子民.
這個就是出《埃及紀》20章所記載的.
說到石板.
我的主題是石板.
你記得摩西拿著兩塊石板.
從山那裡下來的.
記得嗎.
那個石板大概是雕著什麼.
雕著十戒.
就是上帝和他的子民.
正式來到立約的紀錄.
所以有兩塊.
一塊是上帝的.
一塊是上帝的子民.
大家互相去到見證.
所以說到石板的時候.
今天我們就拿一根小石子走.
上一次我們有沙.
提到亞伯拉罕的沙土.
他的後代多如海沙的應許.
然後上一次我們說到族長的故事.
我們拿起這個item.

$^{81}$就是那個puzzle.
我們每個都是有破碎的碎片.
是上帝幫我們砌在一起.
以致能夠成為一個美麗的圖畫.
這次我們拿什麼呢.
就是小石子.
為什麼想起小石子.
因為說起律法.
上帝的子民能夠成為自己的people的時候.
上帝就賜給他們石板.
這個石板就記載著.
上帝去引導他的子民給他們的law.
給他們的律法.
現代人我們需要明白什麼叫做律法.
舊約裡面的律法.
law of the old testament.
因為我們想起law的時候.
法律我們就想起一些不好的東西.
對不對.
我們通常經歷所謂law,法律.
我們就覺得心裡面不是很爽.
兩個星期前我就去某個地方.
我就在街上找個位置.
其實我在high park.
我就在那裡找了一條街有個位置.
繞來繞去.
終於找到一個位置.
但是街的兩旁就有一些橙色的牌.
掛在樹上.
大家在芝加哥本地就知道.
city是很陰濕的.
city很陰濕.
我已經繞了幾次.
這邊就no parking on monday.
對面就很陰濕.
no parking on tuesday.
沒錯.
no parking on monday.
no parking on tuesday.
我找位置那天就是monday.

$^{121}$我找到一個空位.
我已經看清楚了.
no parking on monday.
no parking on tuesday.
繞來繞去有一個U turn.
然後就停在那個位置.
我說這次沒有死.
這個免費的位置.
no parking on monday.
我停在那個位置應該是哪裡.
exactly.
我停了很安心.
我說這次真的沒有死.
幾個小時之後.
我回去就拿了一張橙色的東西.
放在玻璃窗那裡.
豈有此理 原來我出事了.
原來我出事了.
因為我繞來繞去的時候我亂了.
我就停在了no parking on monday那邊.
no parking on monday那邊.
我就以為我停在了什麼.
no parking on tuesday那邊.
因為我找到一個位置.
明明是no parking on tuesday那邊.
為什麼我會停在了.
no parking on monday那邊.
所以我拿了60元.
60元.
其實這不是City的陰濕.
是我自己大意.
我相信大家都出事了.
你心裡會覺得很大聲.
經常出事.
我心同感受.
你的反應可能是.
真的很不公平.
我良民一個.
我平時不犯法.
我停在了位置.

$^{161}$你就罰我60元.
我對法律的感覺都不是很好.
因為通常我們想起法律的時候.
就是懲罰.
或者法律就是陷阱.
所以有時候我們將我們這個概念.
我們就去了解舊約的律法.
我們回到舊約裡面.
很多很長篇的一些律法.
我們的眼光就覺得.
上帝當然好像那些交通警.
有沒有見過那些.
那些City worker.
拿著一部機器.
穿著一件衣服.
然後走來走去.
他就是想抓你.
那你心裡面.
就是心裡面很不滿.
是不是.
心裡面很不滿.
但我們心裡面律法就是這樣的.
就是懲罰.
但是我想今天告訴你.
在舊約裡面講律法不是這樣的.
我們有那些舊的概念.
我們一定要將它刪除.
我們經常說.
我們不在律法之下.
我們現在在恩典之下.
你有沒有聽過這句話.
有啊.
從小到大在教會裡面.
都是聽一樣東西的.
We are not under the law.
We are under grace now.
以前就是律法的時代.
現在就是什麼.
恩典的時代.
現在不講律法了.

$^{201}$我們很自由的.
No No No.
這個是舊的概念.
這個是不明白聖經.
解錯聖經的一種概念.
在舊約的律法裡面.
講到的不是一種懲罰.
它是講到上帝給我們的恩典.
所以有個誤解.
我們就不要將.
好像畫面上面那裡說.
舊約那個神就很像交通警.
很嚴的.
你一犯法.
他就罰你六十塊.
罰你六百塊.
就是引頭落地.
新約那個就不是了.
新約那個講愛的.
講自由的.
這個是錯誤的一個見解.
所以這裡講到.
今天想講.
舊約的律法.
當上帝和他的子民立約的時候.
給他們法板.
給他們石板的時候.
其實那個律法.
是有哪幾方面需要留意的.
第一方面.
舊約裡面的律法.
是以恩典來建立的.
是由恩典來開始的.
我們去土耳其.
我們很多時候.
講起土耳其的時候.
想起什麼.
聖經的故事.
或者聖經的事蹟.
初期教會.

$^{241}$啟示六七教會.
全部都是新約的東西.
但很多時候我們忘記了.
舊約的土耳其.
土耳其舊約的時期.
其實是一個帝國在那裡存在的.
那個帝國叫做Hittite.
客人帝國.
客人.
不是嚇死人.
是客人.
客人帝國.
那些考古學家.
在客人帝國.
就是今天的土耳其.
其中一個地方叫做Hattusa.
那裡就找到.
當時客人帝國的石板.
他們留下的一些文件.
是以前政府的文件.
客人帝國大概是什麼時間.
就是出埃及的時候.
大概1500至1200年的時候.
距離現在大概3000多年的時候.
當時客人帝國Hittite Empire.
它是一個超大的super power.
它跟埃及或者跟其他小國.
就會有很多條約.
有很多外交的文件.
來往.
其中一種文件叫做.
Suzerain Treaty.
中文大概也是一種不平等的條約.
就是說一個大的power.
和一個小的power.
大家之間有一個條約.
一個大國和一個小國之間.
有一個條約.
我們叫宗主條約.
找到這些石板的時候.

$^{281}$翻譯出來的時候.
就發現這些石板.
很多時候這些外交文件.
都有一個format.
那個format就是.
那個條約開頭的時候.
那個大王就會自稱.
或者介紹自己.
我就是什麼什麼Great King.
我就是Great King of Hittite.
Super Great King of Hittite.
I am Great.
我很厲害 很厲害.
說了幾十句.
就是說他自己多厲害.
這個就是自我介紹.
序言preamble的地方.
然後那個格式下去下面.
那個大國大王.
就會和那個小王去介紹.
之前我恩代過你什麼.
你記不記得很多年前.
你被人攻陷.
差不多攻陷的時候.
我派兵去救你等等.
我曾經幫過你的.
我是一個大國.
曾經恩代過你小國的.
現在我就告訴你.
What I have already done for you.
以前我曾經為你做過的好事.
所以到第三部分.
現在你應該怎樣和我建立關係.
你是一個小國.
你應該怎樣回應.
我對你的恩典.
然後就有很多的條款.
就列出來了.
然後到最後第四部分.
就是一些祝福和咒詛的語句.

$^{321}$就是說如果你是faithful的話.
我就會繼續祝福你.
如果你不faithful的話.
你反叛我的話.
那你就麻煩了.
我就會來打你.
就是祝福和咒詛.
大概這樣的條約的形式.
我們找回的時候.
我們就發現.
翻開聖經就是舊約律法的格式.
雖然舊約五經裡面.
記載十誡或者不同的條約的時候.
用一個故事式來記載.
但是裡面有很多的元素.
其實是和Suzerain Treaty很有關係的.
並且它的格式是跟著這樣下的.
你記得十誡第一誡是什麼.
除了我以外不可有別的神.
我們小學開始.
如果是接觸教會的.
我們就會背這些東西.
第一誡.
There is none other than God himself.
除了我以外不可有別的神.
然後第二誡.
就是不可以雕刻偶像等等.
我們就說背十誡.
其實十誡那一段.
《出埃及傑》第二十章.
開始的時候不是除了我以外不可有別的神.
你記得是怎樣開始的.
是我是耶和華利的神.
我是耶和華利的神.
曾將你從埃及地遺老之家領出來.
只是這兩句話.
其實就是跟著Suzerain Treaty的格式.
第一句是說什麼.
我是.
I am the Lord your God.

$^{361}$I am Yahweh the Lord your God.
我就是耶和華利的神.
我就介紹我自己.
我是一個怎樣的神.
當然這個就不像那些客人的帝王.
在自己煲自己.
或者讚自己.
我是多麼偉大.
我是多麼厲害.
好像今天那些國家領導人.
讚到自己落花流水.
其實他不是.
他只是虛報自己.
只是宣傳.
但聖經裡面就很簡單說.
我是耶和華利的神.
然後第二句.
其實是按照Suzerain Treaty的第二部分.
就是我曾經為你做過什麼.
我曾經將你從埃及地領出來.
記得那個故事.
到創世記最後的時候.
他們下了埃及.
是不是因為他們很好.
是不是那些希伯來人.
他們很乖很聰明.
他們是世上最好的人.
最仁慈最良善的人.
所以上帝救他們.
不是的.
上帝是用恩典去救他們.
不是因為值得救你.
而是因為我是耶和華.
我用恩手將你從埃及領出來.
不是跟你講條件.
是因為我愛你.
所以我拯救你出來.
將你領出來.
這兩個說話是跟這個格式的.
就是說我已經將你拯救了.

$^{401}$現在你已經有一個新的身分.
被我恩典拯救出來的時候.
你應該如何生活.
所以拯救恩典是在前的.
然後所謂律法.
其實是指引.
是從恩典流出來.
告訴他們.
你們現在應該要如何生活.
所以我們不要再說.
現在不在律法的時代.
那個是舊約的東西.
現在新約我們不講律法.
我們只講恩典.
這個是錯誤的觀念.
因為舊約裡面.
十誡開始的時候.
已經是從恩典開始.
上帝拯救你.
我們今天新約的信徒.
我們知道在基督裡面拯救了你.
你應該有一個新的身分.
所以你應該如何去生活.
這些就是律法的條款.
所以不是說我做好一點.
守好律法.
我就能夠成為神的子民.
應該相反.
我現在已經.
因著上帝的恩典.
我成為上帝的子民.
有一個新的身分.
新的名分.
所以我現在應該要如何生活.
等於我兒子.
我兒子二十多歲了.
他百分之九十九都很乖.
很乖.
是因為他乖.
所以成為我兒子.

$^{441}$但因為他明白他是我的兒子.
所以他應該生活如此.
有我們的家規.
或者我們的家的文化.
所以今天譬如我們飛去土耳其.
我們有幾個電影節目.
我們要去機場.
他就自己幫忙.
他說我載你.
挺好的.
當然如果他開車不好.
我就會覺得害怕.
但我知道他開車是可以的.
他開車很好.
所以說真的很好.
他免費載我.
不用坐Uber.
挺好的.
他因為有這個身分.
所以他如此行.
不是說我載你.
如果我載你.
載得我乖的話.
我就做你的兒子.
如果不乖.
我就做你的五翼仔.
你趕我走.
不.
上帝的子民和上帝的關係都是一樣.
用恩典來開始.
第二方面.
這些律法.
雖然好像和我們文化沒有什麼關係.
說的那些東西.
大部分是和古代的社會有關.
不過.
其實律法的精義是什麼.
是顯出上帝的心.
和上帝的心腸.
上帝的智慧.

$^{481}$上帝對他的子民的領導.
應該給他們一個怎樣的生活模式.
以至他們的生活模式.
在萬國的中間的時候.
萬國見到他們這樣生活.
見到他們這樣去處理.
他們家裡的事.
生活裡的事.
他們就會說.
他們是不同的.
我給大家兩個例子.
第一個例子就是.
講到房屋.
為什麼講房屋.
這個就是我們看舊的律法的時候.
我們覺得不太合理.
為什麼講這些建築文件.
這些建築文件又不適用我們.
我們看看是不是適用我們.
起碼我們知道一個.
他說你建造新的房屋的時候.
他就吩咐那些以色人.
去到迦南地的時候.
他們建造房屋.
他說你屋頂上一定要有安欄杆.
免得有人從房上掉下來.
流血的罪就歸在你的家了.
ok.
為什麼.
你屋頂有沒有欄杆.
為什麼.
因為芝加哥通常是斜的.
是不是.
通常芝加哥屋頂是用來做什麼的.
我們很少用的.
防水當然是.
我們很少上去BBQ.
跳舞或者做運動.
是不是.
通常不會的.

$^{521}$因為我們的生活方式不同.
以前以色列人.
他們的生活方式.
他們是平頂的房屋.
一定是平頂.
因為他們要儲雨水.
然後將雨水儲起在地下.
裡面一些水池等等.
然後房屋的屋頂.
很多時候用來曬禾稈草.
或者曬稻米等等.
禾就放上去上面.
等等去處理一些東西.
所以他們的屋頂.
其實是工作間.
他們上面不是曬太陽.
是去到農業的生產.
一個很重要的空間.
所以因為房頂經常用.
他們很多時候.
旁邊有樓梯上去.
爬兩層樓梯.
就上到上面的房頂.
你就可以估計到.
他們的生活會是怎樣.
有時候小朋友就會爬上去.
然後他們就會在上面.
不知道發生什麼事就會玩.
然後就會掉下來.
對嗎.
所以為什麼.
以色列人去到他們的家林地的時候.
他們要在房屋的頂上面建圍欄.
其實就是保護生命.
顯出了上帝給神的指紋.
要怎樣.
告訴他們.
生命是神聖的.
所以那個真正律法的意義.
不是我們照辦主管.

$^{561}$今天下來這裡.
21世紀我們這樣用.
而我們要知道背後那個意義.
那個意義就是上帝對生命的看法.
上帝很珍惜每一個的生命.
因為每一個生命都是他所愛的.
所以你們建房屋的時候.
都有建屋法.
這些外邦的律法裡面.
不是時常有的.
但以色列人的律法卻是有的.
還有什麼樣本可以給大家.
這一句應該大家很熟悉.
以前可能曾經說過.
在利未記第19章.
他說你收割莊稼的時候.
為什麼收割莊稼.
因為是農田的生活.
農業的社會.
你收割莊稼的時候.
不可以割盡田割.
也不可以收拾所遺落的.
要留給什麼人.
窮人和寄居的外人.
什麼意思.
就是你去收割田和的時候.
你可以什麼都收了.
就是說你賺錢可以賺到盡.
所有割了頭的東西.
你全部都拿了.
所有大麥小麥的東西.
全部都要了.
你可以這樣的.
但以色列人你不同.
上帝的子民.
你不能夠這樣.
因為你不可割盡田割.
剩下來的東西.
你不要走去再收拾.
不要拿乾拿盡.

$^{601}$原因不是主要環保.
原因是留給一些窮人.
當那些比較有錢的人.
他們去到收割完之後.
那些窮人很多時候.
就會在傍晚的時候.
他們就會來到.
然後他們就可以收拾.
這些剩下來的東西.
這些窮人,寄居者.
他們未必有合法的居留權.
你要留給他們.
因為他們都是上帝所愛的.
這裡說到律法是什麼.
上帝想熱神的子民.
有一個公義慷慨的生活.
我們怎樣看我們的錢財.
看我們的資源.
不要割盡它.
意思是不要算死草.
那些婆娘數不要說那麼多.
我們很喜歡說婆娘數.
我們說別人的格是這樣的.
小國崩都不給你.
為什麼?因為公平.
我知道.
是應該這樣的.
這裡不是說你買菜的時候.
不要跟別人算死草.
你應該算到盡.
不是說你要被人騙.
而是你要慷慨的生活.
因為上帝是一個慷慨的上帝.
祝福你的時候.
你就應該去慷慨的祝福其他人.
甚至有些可能你覺得.
不應該被祝福的那些人.
我用不著理他.
他又不是公民.
他又不是我們自己人.

$^{641}$外人來說清楚 外人 寄居者.
跟我有什麼益處?.
他不會回饋我的.
但是你應該有個慷慨的心態.
應該這樣去生活.
這個不是說要照辦煮碗.
你都不會收割莊稼.
你不會趕著走去收割莊稼.
大家在說什麼?.
我們怎樣處理我們的資源.
上帝祝福我們的東西.
我們應該慷慨的去奉獻給別人.
祝福其他人.
所以立法會一直都是這樣說.
說是顯出了上帝自己的心意.
和上帝對上帝的子民.
有什麼智慧.
如果有人這樣生活的時候.
其他人才會覺得好奇.
我們每個人都是各盡田國.
為什麼你不各盡田國?.
What is the purpose?.
以色列人就會告訴其他人.
因為上帝是這樣.
上帝有沒有跟你各盡田國?.
我告訴你.
上帝沒有欠過你任何東西.
一點上帝都沒有欠過你.
但是上帝在基督裡面.
祂會滿滿祝福給你.
你有沒有想過.
你不要常常覺得上帝欠你一樣.
欠你那樣.
不給你一樣.
不給那樣給你.
上帝一點都沒有欠過你.
由你出生那天開始.
到你死那天.
上帝沒有欠過你一個斗零.
但是上帝是一個Generous的教.

$^{681}$Generosity.
甚至祂的獨生子耶穌基督.
在十字架上為你死.
這個就是上帝.
如果你說你敬拜上帝.
如果你信上帝的話.
Please.
你生活裡面.
給多一點人看看.
看到你信上帝.
作為上帝子民的身分.
這個就是舊約吉律法的本意.
第三方面就講到mission.
我們說上帝是一個God of mission.
上帝永遠都是想拯救罪人.
由始至終都是.
由創世記到啟示錄.
上帝是希望去藉著祂的教會.
藉著祂的子民.
去顯出祂自己的大名.
去告訴人聽.
上帝是好的.
去吸引人來到上帝的面前.
在《申命記》第四章.
摩西去吩咐那些以色列人.
他們未進入迦南地之前.
就有一大堆東西去吩咐他們.
講了很多篇讀.
告訴他們.
你進入迦南地的時候.
應該怎樣生活.
剛才我們說莊惠蘭.
不要割盡田國.
這些都是剛才那個package.
在《申命記》第四章.
這裡講了一句話.
講什麼呢.
他說你進去生活的時候.
在上帝的同在下.
去生活的時候.

$^{721}$成為上帝的子民.
你們要謹守進行上帝的律法.
你們就在萬國的面前.
其他你的neighbor的面前.
在你們身邊的人面前.
去顯出智慧和聰明.
他們 就是說其他的人.
非以色列人.
聽見這一切的律例.
他們必定會說.
這大國就是以色列國.
這大國的人.
真是有智慧有聰明.
哪有一個大國.
有神明與他們相近.
好像耶和華我們的上帝.
在我們求告他的時候.
與我們相近.
他們會讚嘆.
他們會問.
這麼奇怪.
為什麼他們會這樣.
為什麼我們這麼好奇.
他們的生活的方式.
或者生活的心態.
是跟我們不同的.
因為他們的神明是不同的.
是不同的.
我們的生活其實都是一樣.
十誡是上帝親手寫的法版.
在不同的經文都有提過我們.
說摩西拿著兩塊石板.
兩塊 記得.
因為當時立約的條款就是這樣.
你一個copy我一個copy.
你就將你的copy.
放在你最重要的地方.
就是你最大的廟那裡.
你那個拿了一個正確的copy.
你就放在你那個廟裡面.

$^{761}$然後每年就拿出來種稻等等.
記住這樣的律例.
這樣的約.
所以兩塊的石板是.
two tablets of the covenant.
two tablets.
exactly一樣.
所以以色列國.
他們有那兩塊石板就不同了.
很有趣.
他們兩塊石板都是放在藥櫃裡面.
在他們的致勝所裡面.
因為上帝耶和華大王.
他最大的那個.
最重要的地方在哪裡.
就是以色列中間那個temple.
會望.
而以色列人作為一個小國.
他們最重要的地方在哪裡.
就是敬拜耶和華的地方.
會望和將來在耶路撒冷的聖殿.
所以在聖殿裡面的致勝所.
有一個藥櫃.
有一個盒子.
那個盒子裡面有幾樣不同的東西.
很重要的東西.
其中一個最重要的就是石板.
所以兩塊都在裡面.
一塊代表耶和華.
大的王.
The Great King.
一塊是代表以色列國.
上帝的子民.
The Small Kingdom.
大家建立了一個關係.
就是大家放在這裡.
所以說到這個石板的時候.
是上帝用指頭親自寫上去的.
親自去雕上去的.
你一份我一份.

$^{801}$所以那份東西.
不只是備忘錄這麼簡單.
是一個永久的盟約.
所以為什麼.
作為上帝的子民.
他們有一個新的身分之後.
他們就有一堆新的.
去啟蒙他們.
怎樣做一個新的子民.
怎樣生活的一種方式.
這就是每一個律法的條款.
律法是什麼呢.
再說一次.
是先有恩典.
然後我們如此行.
不是因為我們做好一點.
我們乖一點.
所以我們能夠成為朕的子民.
反過來.
是既然是上帝的子民.
我們應當如此行.
既然我們已經是.
被耶穌基督救贖的一班人.
我們就應當如此行.
上帝既然是一個慈愛的上帝.
我們應當在我們生活裡面流露慈愛.
說話的時候流露慈愛.
待人的時候流露慈愛.
既然上帝是一個.
忍耐我們的上帝.
我們就應當怎樣.
能夠忍耐.
既然上帝是一個赦免人的上帝.
我們處事和其他人生活的時候.
我們就應當學校上帝的現象.
上帝是一個慷慨的上帝.
我們就理當是一個慷慨的人.
This is who we are.
最後一個問題.
到底我們的生活方式.

$^{841}$我們的心態.
我們的言語.
我們待人接物等等各方面.
我們有沒有引起別人.
讓別人覺得.
是有點不同.
I want to know.
What is going on.
到底我們的生活是怎樣.
講律法的時候.
不是講一些要去欺騙我們.
要去陷害我們的東西.
上帝的律法原來就是一種邀請.
叫我們活得像上帝.
叫我們活得有智慧.
活得像上帝一樣慷慨.
以致我們的生命能夠on the mission.
我們教會講的是抗疼.
如果我們的心態不擴張的話.
抗多大都沒有用.
或者就算我們不抗.
我們舒舒服服也沒有用.
今晚照一照鏡.
我自己的生活.
有沒有被上帝的恩典啟動過.
我經歷過上帝的恩典.
上帝的愛.
我的生活有沒有像上帝.
我有沒有開口更像上帝.
我講一句說話的時候.
以前是這樣講.
現在我有沒有再想一想.
基督會怎樣講.
我以前這樣處理我的資源.
我的時間.
現在我有沒有有些改變.
是因為我身分上已經有改變.
以致我的生活都會有改變.
讓這個問題.
這個禮拜你帶走好嗎.

$^{881}$我的生活方式有沒有讓其他人.
感到好奇.
好奇我這位這麼慷慨.
這麼愛我的神.
是一個怎樣的神.
同心大道祈禱.
然後你設下你美好的律法.
主要是你自己的律法.
有如金紙一樣.
有如金銀一樣.
不是威脅我們 不是限制我們.
而是一份禮物.
叫我們這樣進行你的生活的時候.
使人能夠渴慕你.
使人能夠認識你.
主要是讓我們認識你更多.
讓我們的言談舉止.
讓我們作的決定.
讓我們處理我們生活裡面有的東西.
讓我們做一個好的管家.
讓我們的生命流露出.
上帝你自己更多美好的本性.
讓人能夠因此而被教會.
被你自己神的指紋吸引去認識你.
祈禱奉傳明教.
\newpage



\section{約翰福音 4:1-26}
\label{sec:wv60KUfih4o}
\textbf{井旁相遇 (約翰福音4\_1-26) - 梁紫健傳道 An Encounter at the Well (John 4\_1-26) - Tsz Kin Leung}
\newline
\newline
連結: \href{https://youtube.com/watch?v=wv60KUfih4o}{\texttt{https://youtube.com/watch?v=wv60KUfih4o}} ~~~~ 語音日期: 2026-05-30
\newline
\newline
\hyperref[sec:OKPQ6OtPkus]{\small{< < < PREV SERMON < < <}}
~
\hyperref[sec:index_chronic]{\small{[返順時目]}}
~
\hyperref[sec:index_scriptual]{\small{[返順卷目]}}
~
\hyperref[sec:OAEc6rF7Lr4]{\small{> > > NEXT SERMON > > >}}
\newline
\newline
約翰福音 4:1-26
\newline
\begin{longtable}{cl}
\hline
\hline
章節 & 經文 (和合本修訂版)\\
\hline
4:1 & \begin{tabularx}{0.7\textwidth}{X} 耶穌知道法利賽人聽見他收門徒和施洗比約翰還多， （ \end{tabularx} \\ \\ \relax
4:2 & \begin{tabularx}{0.7\textwidth}{X} 其實不是耶穌親自施洗，而是他的門徒施洗，） \end{tabularx} \\ \\ \relax
4:3 & \begin{tabularx}{0.7\textwidth}{X} 他就離開猶太，又回加利利去。 \end{tabularx} \\ \\ \relax
4:4 & \begin{tabularx}{0.7\textwidth}{X} 他必須經過撒瑪利亞， \end{tabularx} \\ \\ \relax
4:5 & \begin{tabularx}{0.7\textwidth}{X} 於是到了撒瑪利亞的一座城，名叫敘加，靠近雅各給他兒子約瑟的那塊地。 \end{tabularx} \\ \\ \relax
4:6 & \begin{tabularx}{0.7\textwidth}{X} 雅各井就在那裡；耶穌因旅途疲乏，坐在井旁。那時約是正午。 \end{tabularx} \\ \\ \relax
4:7 & \begin{tabularx}{0.7\textwidth}{X} 有一個撒瑪利亞婦人來打水。耶穌對她說：「請給我水喝。」 \end{tabularx} \\ \\ \relax
4:8 & \begin{tabularx}{0.7\textwidth}{X} 因為那時門徒進城買食物去了。 \end{tabularx} \\ \\ \relax
4:9 & \begin{tabularx}{0.7\textwidth}{X} 撒瑪利亞婦人對他說：「你是猶太人，怎麼向我一個撒瑪利亞女人要水喝呢？」因為猶太人和撒瑪利亞人沒有來往。 \end{tabularx} \\ \\ \relax
4:10 & \begin{tabularx}{0.7\textwidth}{X} 耶穌回答她說：「你若知道神的恩賜，和對你說『請給我水喝』的是誰，你早就會求他，他也早就會給了你活水。」 \end{tabularx} \\ \\ \relax
4:11 & \begin{tabularx}{0.7\textwidth}{X} 婦人對耶穌說：「先生，你沒有打水的器具，井又深，哪裡去取活水呢？ \end{tabularx} \\ \\ \relax
4:12 & \begin{tabularx}{0.7\textwidth}{X} 我們的祖宗雅各把這井留給我們，他自己和兒女以及牲畜都喝這井裡的水，難道你比他還大嗎？」 \end{tabularx} \\ \\ \relax
4:13 & \begin{tabularx}{0.7\textwidth}{X} 耶穌回答，對她說：「凡喝這水的，還要再渴； \end{tabularx} \\ \\ \relax
4:14 & \begin{tabularx}{0.7\textwidth}{X} 誰喝我所賜的水，就永遠不渴。我所賜的水要在他裡面成為泉源，直湧到永生。」 \end{tabularx} \\ \\ \relax
4:15 & \begin{tabularx}{0.7\textwidth}{X} 婦人對他說：「先生，請把這水賜給我，使我不渴，也不用到這裡來打水。」 \end{tabularx} \\ \\ \relax
4:16 & \begin{tabularx}{0.7\textwidth}{X} 耶穌對她說：「你去，叫你的丈夫，再到這裡來。」 \end{tabularx} \\ \\ \relax
4:17 & \begin{tabularx}{0.7\textwidth}{X} 婦人回答，對耶穌說：「我沒有丈夫。」耶穌說：「你說沒有丈夫是對的。 \end{tabularx} \\ \\ \relax
4:18 & \begin{tabularx}{0.7\textwidth}{X} 你已經有過五個丈夫，你現在有的並不是你的丈夫。你這話是真的。」 \end{tabularx} \\ \\ \relax
4:19 & \begin{tabularx}{0.7\textwidth}{X} 婦人對他說：「先生，我看你是一位先知。 \end{tabularx} \\ \\ \relax
4:20 & \begin{tabularx}{0.7\textwidth}{X} 我們的祖宗在這山上敬拜神，你們倒說，應當敬拜的地方是在耶路撒冷。」 \end{tabularx} \\ \\ \relax
4:21 & \begin{tabularx}{0.7\textwidth}{X} 耶穌對她說：「婦人，你要信我。時候將到，你們敬拜父，既不在這山上，也不在耶路撒冷。 \end{tabularx} \\ \\ \relax
4:22 & \begin{tabularx}{0.7\textwidth}{X} 你們所敬拜的，你們不知道；我們所敬拜的，我們知道，因為救恩是從猶太人出來的。 \end{tabularx} \\ \\ \relax
4:23 & \begin{tabularx}{0.7\textwidth}{X} 時候將到，現在就是了，那真正敬拜父的，要用心靈和誠實敬拜他，因為父要這樣的人敬拜他。 \end{tabularx} \\ \\ \relax
4:24 & \begin{tabularx}{0.7\textwidth}{X} 神是靈，所以敬拜他的必須用心靈和誠實敬拜他。」 \end{tabularx} \\ \\ \relax
4:25 & \begin{tabularx}{0.7\textwidth}{X} 婦人對他說：「我知道彌賽亞—就是那稱為基督的—要來；他來了，會把一切的事都告訴我們。」 \end{tabularx} \\ \\ \relax
4:26 & \begin{tabularx}{0.7\textwidth}{X} 耶穌對她說：「我就是，正在跟你說話呢！」 \end{tabularx} \\ \\
[1ex]
\hline
\hline
\end{longtable}
$^{1}$[NON-ENGLISH SPEECH].
Good morning, brother and sister..
[NON-ENGLISH SPEECH].
I don't know if you guys have experienced this, where you.
constantly drink all kinds of beverage all day,.
and yet you still feel thirsty all the time..
[NON-ENGLISH SPEECH].
Maybe it is hot outside, and you drink a whole bottle.
of beverage in one go..
And your throat feels slightly better for a while,.
but after a little bit, it gets dry again..
[NON-ENGLISH SPEECH].
Maybe sometimes you buy a lot of things,.
and the delivery box is starting to pile up at your home..
And you feel excited when you get.
to open these boxes initially..
But a few days later, you start going online again and try.
to see what else you need..
[NON-ENGLISH SPEECH].
Or after you finish watching a whole seasons of a show,.
and you just scroll through your phone all night,.
and check through all the social media updates,.
and you feel like you are finally.
being able to relax or rest..
But as soon as you get off your screen,.
it feels quiet all of a sudden..
And what comes to mind isn't a sense of satisfaction,.
but a feeling of emptiness, as if something is missing..
[NON-ENGLISH SPEECH].
Sometimes this kind of thirst is not even visible on the surface..
You may appear busy, fulfilled, or successful,.
and with a full schedule..
[NON-ENGLISH SPEECH].
But the moment you lie in bed preparing to sleep, you begin to ask yourself,.
what am I chasing?.
And is this what I really want?.
[NON-ENGLISH SPEECH].
So this kind of thirst isn't just a thirst in your mouth,.
but it is one in your heart..
It is a longing for meaning, for love, and for security..

$^{41}$[NON-ENGLISH SPEECH].
So today in John chapter 4, we'll see this Samaritan woman.
who is experiencing this unspeakable thirst..
She went to the well intending to draw water as usual,.
but there she met the Lord Jesus, who could truly quench her thirst..
[NON-ENGLISH SPEECH].
So let's read the passage for today..
[NON-ENGLISH SPEECH].
Now, when the Lord knew that the Pharisees had heard that Jesus was.
making and baptizing more disciples than John,.
although Jesus himself did not baptize, but only his disciples,.
he left Judea and departed again in Galilee..
He had to pass through Samaria..
So he came to a city of Samaria called Sychar,.
near the field that Jacob gave to his son Joseph..
Jacob's well was there..
So Jacob and so Jesus, as he was with his journey,.
sat down beside the well..
It was about six hours..
[NON-ENGLISH SPEECH].
At this point, we might have a question..
Why did Jesus have to pass through Samaria?.
He says, had to..
But in reality, many Jews actually deliberately.
took the detour to avoid Samaria..
[NON-ENGLISH SPEECH].
So let's keep this question in mind,.
and we will answer that in a moment..
[NON-ENGLISH SPEECH].
Verse 7, there came a woman of Samaria to draw water..
Jesus said to her, give me a drink..
For his disciple had gone away into the city to buy food..
The Samaritan woman said to him, how is it.
that you, a Jew, ask a drink of me, a woman of Samaria?.
For Jews have no dealings with the Samaritans..
So Jesus answered her, if you knew the gift of God.
and who it is that is saying to you, give me a drink,.
you would have asked him, and he would.
have given you living water..
The woman said to him, sir, you have nothing to draw with,.

$^{81}$and the well is deep..
Where do you get the living water?.
[NON-ENGLISH SPEECH].
Are you greater than our father Jacob who gave us the well.
and drank from it himself and his sons and his cattle?.
Jesus said to her, everyone who drinks of this water.
will thirst again, but whoever drinks of the water.
that I shall give him will never thirst..
The water that I shall give him will become in him.
a spring of water welling up to eternal life..
The woman said to him, sir, give me this water.
that I may not thirst, nor come here to draw..
[NON-ENGLISH SPEECH].
Here, what is the difference between the living water.
that Jesus speaks of and the water in the well?.
The passage mentions that the woman's response raises.
questions about racial identity..
And what is the relationship between the Jew.
and the Samaritan?.
[NON-ENGLISH SPEECH].
And Jesus said to her, go, call your husband, and come here..
The woman answered him, I have no husband..
Jesus said to her, you are right in saying,.
I have no husband, for you have had five husbands,.
and he whom you now have is not your husband..
This you say truly..
The woman said to him, sir, I perceive.
that you are a prophet, a father's worship.
on this mountain, and you say that in Jerusalem.
is the place where men ought to worship..
[NON-ENGLISH SPEECH].
Jesus said to her, woman, believe me,.
the hour is coming when neither on this mountain.
nor in Jerusalem will you worship the Father..
You worship what you do not know..
We worship what we know, for salvation is from the Jews..
But the hour is coming, and now is,.
when the true worshipers will worship the Father.
in spirit and truth..
For such, the Father seeks to worship him..

$^{121}$God is spirit, and those who worship him.
must worship in spirit and truth..
The woman said to him, I know that Messiah is coming,.
he who is called Christ..
And when he comes, he will show us all things..
Jesus said to her, I who speak to you am he..
I am he, yeah..
[NON-ENGLISH SPEECH].
So this is a long dialogue describing how.
the woman interact with Jesus..
But after that, she was changed..
So what is the change that happened in her?.
[NON-ENGLISH SPEECH].
First, why did Jesus have to pass through Samaria?.
[NON-ENGLISH SPEECH].
Second, what is the difference between the living water.
that Jesus spoke about and the water in the well?.
[NON-ENGLISH SPEECH].
And third, what obvious change do we see in the woman.
after she heard Jesus' word?.
[NON-ENGLISH SPEECH].
So let's answer the first question..
[NON-ENGLISH SPEECH].
Why did Jesus have to pass through Samaria?.
[NON-ENGLISH SPEECH].
Our passage says that Jesus left Judea for Galilee,.
and he had to pass through Samaria..
[NON-ENGLISH SPEECH].
Here he came to a city called Sychar,.
near the piece of land that Jacob.
had given to his son Joseph..
And it mentions that Jacob's well was there..
[NON-ENGLISH SPEECH].
And according to the customs of the Jews at that time,.
they often preferred to take a detour.
rather than passing through Samaria,.
because there was a deep religious and ethnic hatred.
between the two groups..
[NON-ENGLISH SPEECH].
But John tells that Jesus had to pass through Samaria..

$^{161}$This "had to" is not a geographical necessity,.
but a necessity in the plan of salvation..
Jesus was there to find someone and through her.
to find the lost..
[NON-ENGLISH SPEECH].
If you've seen this slide, it's made with AI,.
and it's a map of following Jesus' route..
[NON-ENGLISH SPEECH].
And there are three main areas in these pictures,.
from bottom to top..
[NON-ENGLISH SPEECH].
To the south lies Judea, which is.
where Jerusalem and Bethlehem is..
And that is the orange section..
[NON-ENGLISH SPEECH].
In the middle is Samaria, which is the yellow section..
[NON-ENGLISH SPEECH].
And to the north lies Galilee, which.
is the one section in green..
[NON-ENGLISH SPEECH].
Then there's two lines, the red lines..
[NON-ENGLISH SPEECH].
One of the route is the one that Jesus took,.
going straight north from Judea along the mountains.
through Samaria to Galilee..
This is the straight and the shortest route, like the line.
on the map from Jerusalem passing through Samaria..
[NON-ENGLISH SPEECH].
Another route frequently used by the Jews was to bypass Samaria..
They would first travel east to the Jordan River,.
and then through Jericho, cross the river,.
and then travel north along the east bank..
And once they get north, they would cross the river again.
to return to Galilee..
And this is the so-called East Jordan-Periyah route,.
a detour route that is deliberately.
used to avoid Samaria..
[NON-ENGLISH SPEECH].
This is because the Jews have always.
believed that the Samaritans, they were not pure Jews..

$^{201}$They were once a tribe of the northern kingdom of Israel,.
but they intermarried with the outsider..
And thus, they defiled the Jewish bloodline..
And they were considered as mixed race..
[NON-ENGLISH SPEECH].
Moreover, the Samaritans had also.
joined forces with outsiders to kill the Jews of the south..
[NON-ENGLISH SPEECH].
And this turned the two nations who.
share the same origin and ethnicity.
into bitter enemies, with no contact.
and mutual hostilities..
And this is somewhat similar to the current relationship.
between the North Korean government.
and the South Korean government..
[NON-ENGLISH SPEECH].
If you look at the map, going through Samaria.
is a reasonable and common route..
But it is not the only route..
Jesus used the alternative routes--.
the Jews would use the alternative routes.
to avoid Samaria..
[NON-ENGLISH SPEECH].
When Jesus says in chapter 4 that he.
had to pass through Samaria, the point.
is not that he was forced to take a detour by the terrain,.
but that according to the Father's will,.
he deliberately chose this route that most people would avoid,.
so that he could meet the Samaritan woman there.
and bring salvation to the entire city..
[NON-ENGLISH SPEECH].
So Jesus deliberately passes through this place,.
breaking down the longstanding divide.
and mutual hatred between the Jews and the Samaritans..
[NON-ENGLISH SPEECH].
Jesus did this to reaffirm that salvation was not only.
for the Jews, but also for some peoples who were looked down.
upon and considered unworthy of salvation,.
such as the mixed-race Samaritans..
[NON-ENGLISH SPEECH].

$^{241}$This simple sentence crosses several boundaries..
[NON-ENGLISH SPEECH].
It was very unusual for a man to speak to a woman at that time..
[NON-ENGLISH SPEECH].
And the fact that the Jews is the one who took the initiative.
to ask the Samaritan for water is even more unconventional..
[NON-ENGLISH SPEECH].
And yet Jesus set an example for us by doing this..
[NON-ENGLISH SPEECH].
So are you willing to pass through your Samaria?.
[NON-ENGLISH SPEECH] In our lives, there might be places.
or people that we don't want to pass through..
[NON-ENGLISH SPEECH].
Maybe it's the colleague who everyone dislikes and always.
being joked about..
[NON-ENGLISH SPEECH].
Or the relatives who always speak hurtfully.
at family gatherings..
[NON-ENGLISH SPEECH].
Or the neighbor in your community.
who speaks with an accent and has a very different culture..
[NON-ENGLISH SPEECH].
One time, a brother shared that there was a new immigrant.
in his company, and his coworkers all say,.
don't pay attention to him..
His English is not good, and he doesn't understand the culture..
[NON-ENGLISH SPEECH].
One day, while he was reading John chapter 4,.
he was deeply moved when he read about this Jesus sitting.
by the well and saying to the Samaritan woman,.
give me a drink..
So he decided to invite his colleague to lunch.
so that he could hear him talk about his hometown and his family..
[NON-ENGLISH SPEECH].
After a few meals together, the colleague.
said, you are the first person who really.
wants to understand me..
And from then on, this brother became his first impression.
of Christian..
[NON-ENGLISH SPEECH].

$^{281}$So Brian's sister, this is the kind of God that Jesus is..
He doesn't just wait for people to come to him in the temple,.
but he goes to the well and to the despised place..
[NON-ENGLISH SPEECH].
Today, he's also calling us, who is the Samaritan besides us?.
And are you willing to deliberately make.
a trip for the Lord, seemingly by possible,.
but actually prepared by God, encountered by the well?.
[NON-ENGLISH SPEECH].
Jesus not only spoke to her across the barrier,.
but he was also not content with just giving her water,.
as if to quench her thirst temporarily..
And so the next scene will show the transformation.
from well water to living water..
[NON-ENGLISH SPEECH].
To answer the second question..
[NON-ENGLISH SPEECH].
What is the difference between the living water.
that Jesus spoke about and the water in the well?.
[NON-ENGLISH SPEECH].
The well water and living water here.
are the tools that Jesus used to touch our deepest desire..
[NON-ENGLISH SPEECH].
When the woman was surprised that Jesus would ask her.
for water, Jesus turned the conversation.
into something deeper..
If you knew the gift of God and who it is that says to you,.
give me a drink, you would have asked him,.
and he would have given you living water..
[NON-ENGLISH SPEECH].
At first, she only thought of the actual well water..
So she said, sir, you don't have tools to draw water,.
and the well is deep..
Where will you get this fresh water?.
[NON-ENGLISH SPEECH].
But Jesus said, anyone who drinks this water.
will thirst again..
But whoever drinks the water I give them will never thirst..
And the water I give them will become in them.
a spring of water welling up to eternal life..

$^{321}$[NON-ENGLISH SPEECH].
So the thirst here is no longer just a physical thirst,.
but a deep spiritual emptiness and the desperate desires.
of the soul..
[NON-ENGLISH SPEECH].
Jesus mentioned that she had five husbands.
and that the one she had now was not her husband..
This was not to humiliate her, but to gently reveal.
the deep longing she had been trying to feel with all.
these relationships and marriage--.
to be loved, accepted, and affirmed..
[NON-ENGLISH SPEECH].
Many scholars have pointed out that her multiple marriages.
were likely the results of being divorced or widowed repeatedly..
And these were the signs of being hurt..
[NON-ENGLISH SPEECH].
Today we do the same thing..
We constantly search for various kinds of water..
[NON-ENGLISH SPEECH].
Some people quench their thirst with achievements--.
promotions, raises, and titles..
[NON-ENGLISH SPEECH].
But some people quench their thirst with emotions--.
dating, changing partners, hoping that this time they.
will truly be loved..
[NON-ENGLISH SPEECH].
And some people quench their thirst with entertainment--.
binge-watching TV series, playing games,.
and social medias, just to prevent themselves.
from having the time to sit with their inner thoughts..
[NON-ENGLISH SPEECH].
This is similar to a phone constantly using power..
The screen is on, programs are running,.
and the videos are being watched..
And this, on the surface, seems very likely..
But if the charging cable is not plugged into the actual power,.
it will not charge, no matter how long you plug it in..
[NON-ENGLISH SPEECH].
Maybe you can feel your schedule every day,.
pack your shopping cart to the brim,.

$^{361}$and the friend list on your social media is filled..
But the true yearning in your heart of, do I have any value?.
Or am I truly known and truly loved, remains unsolved..
[NON-ENGLISH SPEECH].
Jesus said to the woman, whoever drinks this water.
will be thirsty again..
[NON-ENGLISH SPEECH].
Today, he is also saying to us, children, what you really need.
is not to change the well, but to change the water,.
not to change the object, but to change the source..
[NON-ENGLISH SPEECH].
He will give us the Holy Spirit so.
that we may worship the Father in spirit and truth,.
so that our lives are not supported.
by actual circumstances, but have a living spring within us,.
so that we can still receive comfort and strength from God.
when we encounter pressure, loneliness, and loss..
[NON-ENGLISH SPEECH].
Perhaps you're sitting in a room today,.
and everything seems fine from the outside..
A complete family, a stable job, and a busy ministry..
But only you who can know the real thirst.
that is in your heart..
[NON-ENGLISH SPEECH].
And today, the Lord is saying to you, bring your water jug.
and give up everything you use to quench your thirst,.
and let me be your living water..
[NON-ENGLISH SPEECH].
In the moment the woman tasted the living water,.
her new life showed a sign..
She put down the water jug and ran back to the city..
This leads us to the third scene,.
the testimony of putting down the water jug..
[NON-ENGLISH SPEECH].
And this is also answering our third question..
[NON-ENGLISH SPEECH].
What obvious change occurred in the woman.
after she heard Jesus' word?.
[NON-ENGLISH SPEECH].
Scripture says that then the woman left the water jug.

$^{401}$behind and went into the city, saying to the crowd,.
come and see a man who has told me everything I've ever done..
Could this be the Messiah?.
[NON-ENGLISH SPEECH].
So many Samaritans believe in Jesus.
because of what the woman said..
And later, because they heard him speak, many more believed..
And they're saying, now we believe no more.
because of what you say, but because we.
have heard it ourselves and know that he is indeed the Savior..
[NON-ENGLISH SPEECH].
Note the following..
[NON-ENGLISH SPEECH].
This Samaritan woman hasn't gone through any discipleship.
courses or received any preaching training..
[NON-ENGLISH SPEECH].
Her path is still messy, and many of her wounds.
still remains, but she has met the Lord..
[NON-ENGLISH SPEECH].
And her sermons were simple, not a systematic theology,.
but just come and see..
Come and see the Lord who sees me entirely.
but has not abandoned me..
[NON-ENGLISH SPEECH].
And God used this broken woman to become.
the first missionaries to the city,.
preparing the way for Jesus..
[NON-ENGLISH SPEECH].
One of the sister had recently become a believer,.
and she didn't understand the Bible very well yet..
And she also hadn't completely changed many of her bad temper..
[NON-ENGLISH SPEECH].
But when she returned to her company, people asked her,.
oh, you seem a little bit different today..
[NON-ENGLISH SPEECH].
So she answered awkwardly, I recently.
went to church and met Jesus..
I don't know how to explain it, but it just.
makes me feel more at peace..
[NON-ENGLISH SPEECH].

$^{441}$Because of that, her co-worker, who.
noticed that she was less likely to lose control.
under pressure and more gentle when she was being criticized,.
they asked her, next time when you go to church,.
can I go with you?.
[NON-ENGLISH SPEECH].
She was not a perfect Christian, but she.
was a true witness like the Samaritan woman..
[NON-ENGLISH SPEECH].
So today, God is also asking us these questions..
[NON-ENGLISH SPEECH].
Would you be willing to be like her,.
to put down your water jug and stop.
spending your whole life just to fetch water.
for your own thirst?.
[NON-ENGLISH SPEECH].
Would you be willing to go back to your own city, your home,.
your workplace, your social circle,.
and share honestly on how God has worked in you just.
like she did?.
[NON-ENGLISH SPEECH].
You don't need to be able to explain all the theology..
All you need is just a willingness.
to say, come and see..
And the Lord can use your few words.
to awaken the thirst in one's heart and bring them to Jesus..
[NON-ENGLISH SPEECH].
So the scripture answered us three questions..
[NON-ENGLISH SPEECH].
Before we end, I would like to invite everyone.
to response with these three questions before the Lord..
[NON-ENGLISH SPEECH].
First, where is your Samaria?.
[NON-ENGLISH SPEECH].
Today, Jesus might be waiting for you by your well..
Perhaps it is your daily commute in front of your work seat.
or the sofa in your living room..
[NON-ENGLISH SPEECH].
Would you be willing to stop at the most ordinary places.
in your daily life and listen to him?.

$^{481}$Take a step forward and reach out to the people.
that you wouldn't normally encounter..
[NON-ENGLISH SPEECH].
Second, what is your deepest desire?.
[NON-ENGLISH SPEECH].
Perhaps it is a desire for relationships or achievements.
[NON-ENGLISH SPEECH].
Perhaps you can't even explain it yourself..
You just know that you're tired and empty..
And all you know is that you need to earn more and more..
[NON-ENGLISH SPEECH].
So today, are you willing to bring all of this before the Lord.
and say to him, Lord, I no longer want the well water..
I want the living water that only you can give..
[NON-ENGLISH SPEECH].
Third, who is your city?.
[NON-ENGLISH SPEECH].
Who are some of the people in your life that God has placed?.
Maybe your family, your coworkers, your classmates.
that Jesus wants to invite them through you..
Come and see..
[NON-ENGLISH SPEECH].
Are you willing to pray for them this week,.
asking God to give you the opportunity.
to take that extra step, to say that extra words,.
or to show an extra care for one of them this week?.
[NON-ENGLISH SPEECH].
Let us pray together..
[NON-ENGLISH SPEECH].
Lord Jesus, we confess that we often.
try to quench the thirst in our heart with a lot of well water..
But thank you for reminding us today.
that you are the true living water, the Messiah who.
transcends all barriers, and yet you wait for us by the well..
[NON-ENGLISH SPEECH].
Lord, would you heal our past brokenness.
and be the spring of water within us..
Heal us, like the Samaritan woman,.
to be a true witness in our home, in our workplace,.
and in our circle of friends..

$^{521}$[NON-ENGLISH SPEECH].
Bring people to you in the name of Jesus..
Amen..
[NON-ENGLISH SPEECH].
\newpage



\section{傳道書 4:4-5:7}
\label{sec:OAEc6rF7Lr4}
\textbf{關係的失衡(傳道書4\_4-5\_7) - 譚文倩莊傳道}
\newline
\newline
連結: \href{https://youtube.com/watch?v=OAEc6rF7Lr4}{\texttt{https://youtube.com/watch?v=OAEc6rF7Lr4}} ~~~~ 語音日期: 2026-05-30
\newline
\newline
\hyperref[sec:wv60KUfih4o]{\small{< < < PREV SERMON < < <}}
~
\hyperref[sec:index_chronic]{\small{[返順時目]}}
~
\hyperref[sec:index_scriptual]{\small{[返順卷目]}}
~
\hyperref[sec:augzRkvXpyY]{\small{> > > NEXT SERMON > > >}}
\newline
\newline
傳道書 4:4-5:7
\newline
\begin{longtable}{cl}
\hline
\hline
章節 & 經文 (和合本修訂版)\\
\hline
4:4 & \begin{tabularx}{0.7\textwidth}{X} 我見人因彼此嫉妒而有一切的勞碌和各樣工作的成就，這也是虛空，也是捕風。 \end{tabularx} \\ \\ \relax
4:5 & \begin{tabularx}{0.7\textwidth}{X} 愚昧人抱著雙臂，自食其肉。 \end{tabularx} \\ \\ \relax
4:6 & \begin{tabularx}{0.7\textwidth}{X} 一掌滿滿而得享安靜，勝過兩掌滿滿而勞碌捕風。 \end{tabularx} \\ \\ \relax
4:7 & \begin{tabularx}{0.7\textwidth}{X} 我轉而觀看日光之下有一件虛空的事： \end{tabularx} \\ \\ \relax
4:8 & \begin{tabularx}{0.7\textwidth}{X} 有人孤單無雙，無子無兄弟，竟勞碌不息，眼目也不以財富為滿足。他說：「我勞碌，自己卻不享福，到底是為了誰呢？」這也是虛空，是極沉重的擔子。 \end{tabularx} \\ \\ \relax
4:9 & \begin{tabularx}{0.7\textwidth}{X} 兩個人總比一個人好，他們勞碌同得美好的報償。 \end{tabularx} \\ \\ \relax
4:10 & \begin{tabularx}{0.7\textwidth}{X} 若是跌倒，這人可以扶起他的同伴；倘若孤身跌倒，沒有別人扶起他來，這人就有禍了。 \end{tabularx} \\ \\ \relax
4:11 & \begin{tabularx}{0.7\textwidth}{X} 再者，二人同睡就都暖和，一人獨睡怎能暖和呢？ \end{tabularx} \\ \\ \relax
4:12 & \begin{tabularx}{0.7\textwidth}{X} 若遇敵攻擊，孤身難擋，二人就能抵擋他；三股合成的繩子不易折斷。 \end{tabularx} \\ \\ \relax
4:13 & \begin{tabularx}{0.7\textwidth}{X} 貧窮而有智慧的年輕人，勝過年老不再納諫的愚昧王， \end{tabularx} \\ \\ \relax
4:14 & \begin{tabularx}{0.7\textwidth}{X} 那人從監牢裡出來作王，在國中原是出身貧寒。 \end{tabularx} \\ \\ \relax
4:15 & \begin{tabularx}{0.7\textwidth}{X} 我見日光之下所有行走的活人，都跟隨那年輕人，就是接續作王的那位。 \end{tabularx} \\ \\ \relax
4:16 & \begin{tabularx}{0.7\textwidth}{X} 他的百姓，就是他所治理的眾人，多得無數；但後來的人還是不喜歡他。這也是虛空，也是捕風。 \end{tabularx} \\ \\ \relax
5:1 & \begin{tabularx}{0.7\textwidth}{X} 你到神的殿要謹慎你的腳步；近前聽，勝過愚昧人獻祭，他們不知道自己在作惡。 \end{tabularx} \\ \\ \relax
5:2 & \begin{tabularx}{0.7\textwidth}{X} 在神面前你不可冒失開口，也不可心急發言；因為神在天上，你在地上，所以你的話語要少。 \end{tabularx} \\ \\ \relax
5:3 & \begin{tabularx}{0.7\textwidth}{X} 事務多，令人做夢；話語多，顯出愚昧。 \end{tabularx} \\ \\ \relax
5:4 & \begin{tabularx}{0.7\textwidth}{X} 你向神許願，還願不可遲延，因他不喜歡愚昧人，你許的願應當償還。 \end{tabularx} \\ \\ \relax
5:5 & \begin{tabularx}{0.7\textwidth}{X} 你許願不還，不如不許。 \end{tabularx} \\ \\ \relax
5:6 & \begin{tabularx}{0.7\textwidth}{X} 不可放任你的口使肉體犯罪，也不可在使者面前說是錯許了。為何使神因你的聲音發怒，敗壞你手所做的呢？ \end{tabularx} \\ \\ \relax
5:7 & \begin{tabularx}{0.7\textwidth}{X} 多夢多言，其中多有虛空，你只要敬畏神。 \end{tabularx} \\ \\
[1ex]
\hline
\hline
\end{longtable}
$^{1}$麻煩你,Carmen.
我想問問大家.
人與人之間的關係.
你覺得甚麼是最重要?.
是否請喝茶?.
不是?.
沒理由.
如果我告訴你是良好的溝通.
你同意嗎?.
但甚麼才是良好的溝通?.
如果讓你定義.
甚麼是良好的溝通?.
聆聽.
這是重點.
我想至少要包括.
有個清楚的表達.
剛才也說到.
我們要用心聆聽.
好奇對方說的話.
真誠去嘗試理解.
或者不明白的時候.
會問問題.
去釐清對方說的話的意思.
以及要有合適的回應.
但很多時候.
誰有這麼多時間?.
所以說完.
你聽到就聽到.
你明白就明白.
不明白就算了.
然後才會出現那些.
說你又不聽.
聽你又不明白.
明白之後怎樣?.
又不做.
做又會做甚麼?.
做錯.
錯了又會怎樣?.
錯了又會賴我.
經常會出現勞氣的情況.

$^{41}$今天我們繼續說傳道書.
今天的經文.
是從四章四節到五章七節.
關係的失衡.
即是失平衡.
我在家裡預備這段經文的時候.
看到我一頭霧水.
我就和我老公說.
傳道者是不是有ADHD?.
是不是有些過度弱症?.
為甚麼這段經文.
好像跳一下這裡跳一下那裡.
還是他不是很夠分.
思想有點凌亂.
所以我一直看的時候.
究竟他想說甚麼?.
當然可能是我的問題.
是我自己功力不夠深厚.
若果將傳道書.
將傳道書三至六章抽出來.
前面的三章一至二十二節.
如果你記得的時候.
是說定時.
是說時間如何限制人.
同時顯出神的智慧和主權.
後面那段的五章十節.
到六章九節.
是說財富和享樂.
我們人如何掌控,擁有和滿足.
夾在中間大約四章一節.
到五章九節裡面.
是說六種不同的關係.
這六種不同的關係.
就好像一層膠水.
將前面那段說定時.
和後面那段說財富.
這兩個大主題.
好像將它們拉在一起.
提醒我們.
無論是時間的限制.

$^{81}$還是財富的追求.
都離不開關係這個元素.
就是想說關係有多重要.
其實一點都不凌亂.
上次四章一至三節.
傳道者說了第一個關係是什麼關係.
是壓迫者和被壓迫者之間的關係.
今天的四章四節.
五章七節就轉向人和勞碌.
人和人之間.
人和權力.
人和神之間的關係.
我們先一起低頭祈禱.
主啊我們心存感恩.
我們人算什麼.
你竟然顧念我們.
主啊求你謙卑我們的心.
去聆聽你要和我們說的話.
奉主耶穌基督聖名祈求.
我先和大家讀一讀經文四章四至十二節.
我見人因彼此嫉妒而有一切的勞碌.
和各樣工作的成就.
這也是虛空也是暴風.
如每人抱著雙臂自食其肉.
一掌滿滿而得享安寧.
勝過兩掌滿滿而勞碌暴風.
有人孤單無雙無子無兄弟.
竟勞碌不息眼目也不以財富為滿足.
他說我勞碌我自己卻不享福.
到底是為了誰呢.
這也是虛空是極沉重的擔子.
兩個人總比一個人好.
他們勞碌同得美好的報償.
若是跌倒借人可以扶起他的同伴.
倘若孤身跌倒沒有別人扶起他來.
這人就有禍了.
再者二人同睡都戀和.
一人獨睡怎能戀和呢.
若遇仇敵攻擊孤身難擋.
二人就能抵擋他.

$^{121}$三股合成的繩子不易截斷.
第四節這裡一開始說一種人.
他勞碌得來的成就.
也成為了其他人嫉妒他的原因.
而這個嫉妒也成為了.
其他人追求成功的動力.
傳道者說這些都是虛空都是暴風.
傳道者在這裡指出.
如果我們人的勞碌.
是由比較和嫉妒去推動回來的時候.
其實是很沒有意思的.
正因為這樣的原因.
所以有些人說我不玩這些遊戲.
我不喜歡做一個欺壓人的人.
我也不喜歡競爭.
所以我索性慢活.
這些是什麼人.
有沒有聽過躺平這個詞彙.
躺平是中國大陸近年出現的一個詞彙.
意思是在一個高競爭的社會裡面.
人選擇不做一些要求高的工作.
不追求物質.
也不去嘗試回應社會對他的期望.
他選擇用一個最低消費的方式去過日子.
其實這些是什麼人.
其實就是我們常說的寄生蟲.
第五節說愚昧人抱著雙臂自食其肉.
這種人懶到一個點.
他連把食物放進自己口裡也懶.
他最終只會由內向外去毀滅自己.
也拖累身邊的人和社會.
第六節一掌滿滿而得享安寧.
意思是辛勞工作.
但是也要懂得休息.
恰當的比例.
見過兩掌滿滿而勞碌暴風.
是在說第八節所提到的.
兩隻手不停只有個做字.
又孤零零一個人的工作狂.
傳道者在這裡說這個人沒有伴侶.

$^{161}$沒有子女.
只是不停地做.
不斷地去囤積賺回來的錢.
從來沒有停下來想.
或者他沒有時間去想.
我其實這麼辛苦是為了誰呢.
傳道者的質疑是要帶我們去思想.
在躺平和工作狂之間.
生命裡面會不會有些東西.
是更加值得我們去追求的呢.
第九節說兩個人總比一個人好.
老老實實有時候一個人也挺好的.
哇.
Erica第一個說.
我自己也喜歡有時候一個人.
自己出去看場戲.
喝杯咖啡.
找個餐館自己吃東西.
去享受一下我自己的me time.
但享受一個人的時間.
和孤單是兩回事.
舊約僧賢十八章二十四節.
他說朋友太多的人必受損害.
朋友多呢.
我們經常會說這個人當然是人緣好.
聽得多.
為什麼這裡說朋友多必受損害呢.
在這次經文裡面的原文.
朋友這個字其實有兩個不同的字根在裡面.
第一個朋友的名詞叫做rare.
意思是好鄰舍.
就是這個朋友可以幫你送外賣.
但有事起來的時候.
他可以告你上法庭.
同一個字代表一大堆的關係.
這種朋友的關係.
可能只是停留在.
我們大家都知道對方很多事情的層面.
每次見面都是我們說一下.
我最近在做什麼.

$^{201}$你又最近在做什麼.
Bin最近又在做什麼.
我們只是不停在增加資訊.
很膚淺的.
沒有深度的.
對我們大家之間的關係也沒有什麼好處.
箴言在這裡說.
如果你身邊只有這種朋友的話.
最終這些關係會毀滅了你.
之後聖經用了另一個動詞叫ahav.
這種朋友不單只是陪伴你的人.
而是一個真正關心你生命的人.
當你遇到困難的時候.
他可能比你最親的人.
更加快出現在你家門口.
當你做錯,走錯路的時候.
他又不會怕得罪你.
而選擇不出聲.
因為真正不想你死的朋友.
他給你的提醒.
雖然有時候真的很難聽.
坦白說.
我其實有誰喜歡被提醒.
喜歡被糾正.
好像我每一次寫完講章.
我就會說一次給我老公聽.
心裡我經常都希望他讚我.
但結果每一次他都說媽媽這裡這樣.
媽媽這裡這樣.
我就覺得自己不斷縮小.
但是在這個過程裡.
我又在想.
我知道我老公很認識我.
他也很喜歡聖經.
所以他覺得他的出聲很重要.
我就會慢慢可以再振習起來.
好好地繼續預備講章.
忠言逆耳.
這些是最愛疼你的人.
不過能夠做到這一步的確不是很簡單.

$^{241}$最近我們家庭團契的媽媽組.
我們的一次聚會大家才在這裡分享.
當我們很坦白地去說我們自己的感受.
有時候其實真的是一個很大的冒險.
因為可能你說完出來之後.
你會被誤解.
你會被批評.
甚至受傷.
特別是當我們坦誠之後.
可能換回來的是一些冷淡的回應.
被論斷.
甚至失去一些關係.
結果會怎樣.
我們人就會越來越不出聲.
我們就會學會保護自己.
我們不敢再分享.
但是我在這裡反思.
如果在教會這個大家庭裡面.
我們都不能夠衷心去分享.
我們不能夠傾心吐耳.
另一方面我們不去學習聆聽的功課.
我們學習不要去批評.
我們學習接納.
甚至當對方的選擇.
不是我想最好的選擇也好.
我也去尊重對方的選擇.
甚至支持對方.
我們弟兄姊妹的關係.
只會停留在一個很表面的情況下.
我們教會是很難活出.
弟兄姊妹關係的真正意義.
我想問你.
當我一直這樣說的時候.
你腦海裡面有沒有出現一個.
或者一些這樣的朋友.
或者當你有事的時候.
你想不想起有哪一個會馬上來到你家門口.
回過頭來.
你又是不是你朋友的這種朋友.
傳道者在這裡說.

$^{281}$可能會有的.
可能一個.
兩個.
三個可能已經是極限.
為什麼.
因為大部分的時間.
我們人與人之間有很多競爭.
很多競爭.
很多比較在裡面.
你知道男人特別明顯嗎.
因為男人通常會透過比較去定義自己.
為什麼.
因為裡面有一個潛在的意識.
就是當我比你好的時候.
我就找到我自己的價值.
這個就是傳道書裡面所說的競爭.
競爭本身不是一個問題.
而是當競爭的目的是為了我要壓著你.
這個欺壓的動作的出現.
什麼是欺壓的動作.
而競爭為什麼會變成不公義.
就是當你要贏.
你想人羨慕你.
你就要壓著人.
達到目的的時候.
我們就會被試探.
口不對心.
會走捷徑.
會偷工減料.
會詭詐.
會說謊.
我們會踩人.
我們會妥協自己的誠信.
誠信是什麼.
誠信是我說出來的東西.
和我心裡想的東西是一樣的.
我答應了你的東西.
就算我吃虧我也會做.
再加上一樣.
就是當你做錯的時候.

$^{321}$你會認錯.
你會求原諒.
這個就是誠信.
但當你只是活在一個.
只是競爭的社會的心態裡面.
你整天只顧著要贏.
你還會不會.
整天選擇說真心的話.
會守承諾.
會認錯.
答案是有時都會的.
但當有利益衝突的時候.
你就可能選擇不會.
其實說到底.
欺壓不單只是想說.
我們要比別人好.
而是有意無意.
我們要壓低對方.
為了要保護.
去鞏固我自己.
因為當別人處於.
一個比我低的位置.
我就自己越來越感覺良好.
同時我也不斷.
餵養我自己.
的這種心態.
另一個智慧真言.
一條繩.
兩條繩.
最堅固的是.
三條繞在一起的繩.
如果兩個人比一個人好.
三個人就會更加好.
人生很多事情的確是很虛空.
但有些事情相對來說.
是沒有那麼虛空.
例如建立關係.
我們彼此同行.
追求群體的生活.
比起單純追求財富.

$^{361}$成就 享受.
是有更長遠的價值.
因為當我們遇到攻擊的時候.
你很需要一兩個人在你身邊.
跟你一起走.
去明白你.
去撐你.
在你軟弱的時候.
透過安慰 提醒和保護.
成為你生命裡的一個穩定的力量.
再者在權力 制度和群體的社會裡.
我們更深的焦慮.
就是今天我們是被需要的.
但可能明天已經沒有人記得你.
今天你是重要的.
下一刻你已經被取代.
現實是很殘酷的.
接著13至16節.
他說貧窮而有智慧的年輕人.
勝過年老不再立年的愚昧王.
那人從監牢裡出來作王.
在國中原是出身貧寒.
我見日光之下所有行走的活人.
都跟隨那年輕人.
就是接續作王的那位.
他的百姓.
就是他所治理的眾人.
多得無數.
但後來的人還是不喜歡他.
這也是虛空 也是暴風.
這裡說的智慧是什麼智慧.
這裡說的智慧是謙卑受教.
然後將選擇,行為和結果或者後果.
放在一起去想.
我再說一次.
就是謙卑受教.
然後將選擇,行為和結果或者後果.
放在一起去想.
當中的重點是謙卑受教.
不要自以為聰明.

$^{401}$這裡說的年老的王.
他有什麼優勢.
他年紀大.
所以年紀大會怎樣.
食鹽多過你食米.
很有經驗.
他是王.
有權.
為什麼他是愚昧.
不是因為他有老人癡呆症.
而正正是因為他恃著自己食鹽多過你食米.
所以他不再聽人說.
他不再願意接受教導.
不願意再接受勸告.
這個王出身卑微.
坐過牢.
有點像虐室的故事.
但有一天他都會怎樣.
有一天他都會老.
到時他所治理的百姓.
就會轉會去支持另一個年輕人.
成為新一代的明星.
接著這個明星慢慢都會過氣.
再有新的明星出現.
沒多久這個明星又會再老.
又再被取代.
即是你以為自己很厲害嗎.
有一天你都會被取代.
這個世界沒說沒了誰就不行.
你以為自己年紀越大.
你就越有智慧.
有地位就代表成功.
傳道者在這裡提醒.
當你以為自己已經懂了.
不需要再學.
其實就是開始衰落的時候.
這時虛空.
也時暴風.
我們基督徒的生命成長.
是一生之久的事.

$^{441}$不是說到65歲.
你退休的時候就停下來.
甚至我們牧師傳道人又如何.
牧師傳道人就更加應該是一個願意受教的僕人.
誰說去到某個年紀.
讀了很多書.
或者去到某個位置.
我就已經懂了.
不用再學.
不用再去成長.
五章一至七節.
你到神的殿要謹慎你的腳步.
要向前聽聖過如每人獻祭.
他們不知道自己在作惡.
在神面前你不可無失開口.
也不可心急發言.
因為神在天上你在地上.
所以你的話語要少.
事無多令人做夢.
話語多顯出愚昧.
你向神許願還願不可遲疑.
因他不喜歡愚昧人.
你許的願應當償還.
你許願不還不如不許.
不可放任你的口使肉體犯罪.
也不可在使者面前說是錯許了.
為何使神因你的聲音發怒.
敗壞你所做的呢.
多夢多現其中多有虛空.
你只要敬畏神.
在這裡開始.
我很想大家要留心.
因為傳道者在這裡.
要給一個非常重要的智慧訊息.
教導給我們.
其實之前那段.
傳道者的語氣.
都是比較傾於勸勉.
但由這裡開始.
他的語氣就嚴厲了很多.

$^{481}$傳道者可以很嚴厲地批判.
人來到神的聖殿的時候要謹慎.
要想清想楚才能開口.
我們不要用我們的說話.
我們的口去說一些東西出來.
去嘗試操控神.
嘗試控制神的作為.
尤其輕率許願和不守承諾.
神面前不要失開口.
不要心急發言.
話語要少是什麼意思.
你不要來到神的面前.
一卷嘴在那裡說.
我們通常對著神的時候.
說什麼呢.
就是告訴神.
我需要你應允我什麼.
我有什麼需要.
我求你這樣那樣.
我們說的東西.
全部都是要神為我們做一些東西.
去達到一些東西.
你祈禱的時候.
整個祈禱是不是只是不停地說.
自己想神怎樣怎樣的祈禱.
真正的敬拜.
揭露了我們人的自我中心.
和神尊榮之間的差距.
一個巨大的差距.
因為真正的敬拜不是我們人用說話.
用一些宗教活動.
或者儀式去跟神做交易.
或者我們去嘗試用祈禱.
去展示我們自己有多麼的敬虔.
來到神的面前.
第一件事是要做什麼.
去聽禱.
不是.
去聽敬拜的音樂.
也不是.

$^{521}$最重要的是.
聽神跟你說話.
神啊 我來到你的面前.
因為我需要聆聽.
你要跟我說的話.
無論是透過講道的訊息.
還是透過詩歌裡面的歌詞.
祈禱.
還是跟弟兄姊妹相交.
神啊 我好想聽到你的聲音.
因為在我生命裡面.
現在正在發生一些事情.
我很需要你跟我說話.
經文教我們.
近前聽.
打開耳朵去聽.
什麼叫聽呢.
我們在這裡聽的這個字.
和我們英文的listen.
或者中文聽這個字.
是有一個更深層的意思.
在聖經裡面.
不單止是在講聆聽.
而是在講我們主動地.
謙卑下來去聆聽.
是帶著一個信靠,信服的心去回應.
講得直接一點.
是在講我們要用心去聽.
所以敬拜的優先次序是.
我們謙卑去聆聽神要跟我們說的話.
而聽了之後是會有回應的.
多夢多言其中多有虛空.
你只要敬畏神.
講就講了很多次.
但是如果我們沒有一個敬畏神的心.
其實都是虛空的.
都是沒有意思的.
之前四章四至十六節.
經文提到錯誤的地上敬拜.
我們人將成就,財富,權力當作是神.

$^{561}$五章一至七節這裡.
錯誤的殿中敬拜.
是在講我們人用嘴巴去敬拜上帝.
但是可能只是一場宗教的表演.
這段經文傳道者提到.
其實歸根究底.
講了這麼多種不同的關係.
其實都是在反映我們和神之間的關係.
因為敬畏神不單只是一個態度.
而是我們基督徒的生命核心的方向.
因為當我們人的生命是被神所塑造的時候.
我們就不會再被我們的自我中心.
不斷地去操控.
神的標準影響我們的言語.
影響我們所做的選擇.
影響我們對每一段關係的態度.
影響我們怎樣用我們的時間.
和我們生命的優先次序.
影響我們的價值觀.
換句話來說.
敬畏神的生命是容讓神去塑造我這個人.
怎樣去做人.
而一個人怎樣去做人.
就決定了他怎樣對待他身邊的關係.
關係的失衡.
由探討日光之下.
我們人以自我中心所經營的關係.
轉移去思考我們人和神之間的關係.
人和人之間的關係的失衡.
不單止是社交的問題.
根本上就是人和神之間的關係.
失衡的結果.
因為歪了的心.
我們容不二會嘗試用言語,儀式,承諾.
去操控人,操控上帝.
結果我們人就由欺壓人延伸到欺壓上帝.
傳道者智慧的導向.
我們人要選擇由敬畏神主導的方向去做起點.
因為敬畏神是所有的關係準則和決定最有智慧的方向.
我們一起低頭祈禱.

$^{601}$七海天父我們多謝你.
多謝你今天藉著你的話語提醒我們.
再一次提醒我們不要停留在一個只有白口的狀況裡面.
你再一次提醒我們.
要我們思考我們來到你面前的時候.
我們是否謙卑?我們是否專心?.
我們是否在聆聽你要跟我們說的話?.
主啊求你幫助我們明白.
我們人與人之間.
或者我們與很多我們生命裡面不同的關係裡面的失衡.
其實都是因為我們與你之間的關係.
有些事要處理好.
因為唯有我們追求去敬畏你.
我們願意在你面前學習.
好像耶穌那樣去愛人.
去尊重人.
主啊我們需要你的幫助.
求你幫助我們來到你面前的時候.
不要說那麼多無謂的話.
教我們活出一個真實的生命.
叫我們心存謙卑誠實.
行事為人都是土裡喜悅的.
願你幫助我們每天都是帶著敬畏和信靠的心去跟隨你.
奉主耶穌基督的聖名祈求.
Amen.
\newpage



\section{利未記 23:39-43}
\label{sec:augzRkvXpyY}
\textbf{Booths (Lev 23\_39-43) Rev. Mark Chan 住棚 (利未記23\_39-43) - 陳崇基牧師}
\newline
\newline
連結: \href{https://youtube.com/watch?v=augzRkvXpyY}{\texttt{https://youtube.com/watch?v=augzRkvXpyY}} ~~~~ 語音日期: 2026-07-12
\newline
\newline
\hyperref[sec:OAEc6rF7Lr4]{\small{< < < PREV SERMON < < <}}
~
\hyperref[sec:index_chronic]{\small{[返順時目]}}
~
\hyperref[sec:index_scriptual]{\small{[返順卷目]}}
~
\hyperref[sec:pCh24jxVqNk]{\small{> > > NEXT SERMON > > >}}
\newline
\newline
利未記 23:39-43
\newline
\begin{longtable}{cl}
\hline
\hline
章節 & 經文 (和合本修訂版)\\
\hline
23:39 & \begin{tabularx}{0.7\textwidth}{X} 「但是，從七月十五日起，你們收藏了地的出產之後，要守耶和華的節期七日。第一日為要完全安息，第八日也要完全安息。 \end{tabularx} \\ \\ \relax
23:40 & \begin{tabularx}{0.7\textwidth}{X} 第一日，你們要拿美好樹上的果子、棕樹枝、樹葉茂密的枝條和河邊的柳枝，在耶和華－你們的神面前歡樂七日。 \end{tabularx} \\ \\ \relax
23:41 & \begin{tabularx}{0.7\textwidth}{X} 每年你們要向耶和華守這節七日。你們在七月裡所守的節，要成為世世代代永遠的定例。 \end{tabularx} \\ \\ \relax
23:42 & \begin{tabularx}{0.7\textwidth}{X} 你們要住在棚裡七日；凡以色列家出生的人都要住在棚裡， \end{tabularx} \\ \\ \relax
23:43 & \begin{tabularx}{0.7\textwidth}{X} 好叫你們世世代代知道，我領以色列人出埃及地的時候，曾使他們住在棚裡。我是耶和華－你們的神。」 \end{tabularx} \\ \\
[1ex]
\hline
\hline
\end{longtable}
$^{1}$Good morning everyone, well we have been walking through the Bible big story, remember from.
Genesis to Revelation is a huge grand story, it's about this God of mission and He is.
determined to bless the world against rebellion against Him..
And then you remember God chose Abraham and the patriarchs and used them for redemptive.
work even though they were broken and then He redeemed Israel out of Egypt and gave them.
law at Mount Sinai and then astonishingly, astonishingly God chose to dwell among them.
through the tabernacle and through the desert and through the wilderness, again it's the.
theme of our daily devotion this month..
Now today passage is taken from the book of Leviticus which is the third book of the Pentateuch,.
the five-born Moses and it talks about a festival that God told Israel to keep..
So it's all about festival in the Old Testament, so you think about festivals okay, you think.
about the reason why we celebrate or certain kinds of festival..
That's not what happened, you know people are going firework until 2 or 3 a.m. in my.
neighborhood, right?.
It was July 4th Independence, it was a festival and you remember the reason why we have liberty.
and freedom and you think of a lot of different festival for many of you, you think of Dragon.
Ball, right?.
Recently there was a festival in the Chinese context and then Christian festival as well..
We have Christmas, we have Easter, we have the Lord's Day, every day one day of seven.
days a week and then we come to worship..
So there is reason behind all festivals..
So for Israel, when they came out of Egypt and not yet to the promised land, there were.
40 years they were in the desert, right, in the wilderness and then they are not there.
yet and God set up a tabernacle among them and then right at the center of their camps.
and leading them forward step by step..
So they came out of Egypt and they were the new people and they received God's law and.
the new way of life..
Now they are wandering, you know, in the wilderness and God gave them no map, there was no map.
in the desert, remember..
And then God surprisingly gave them some festival to keep and there were three festival in the.
Old Testament..
So as you can see on the slide, the first one, the biggest one is Passover which happens.
to be early April festival and then the harvest, the Pentecost, which means weeks, 50 days..
So 50 days, 49 days after the Passover is the Pentecost, 50, right?.
Happens to be May 21st this year and the Feast of Tabernacles, remember, you know, the wilderness.
years and Israel building booths in the wilderness..
And then today's passage is taken from the book of Leviticus, it talks about the Feast.
of Booths, Feast of Tabernacles..
So the Jewish people today, especially in Israel, they are still keeping the same festival.
and you know, this year as I show you is in late September..

$^{41}$So let me read this passage, you know, we're going to translate it as well, 2339 in the.
book of Leviticus..
So God said to Israel, to his people, "So beginning with the 15th day of the seventh.
month, after you have gathered the crops of the land, celebrate the festival of the Lord.
for seven days..
The first day is the day of rest and the eighth day is also a day of rest.".
So verse 40, it says, "On the first day of the week, you are to take choice fruit from.
trees and palm trees and branches and all of the poplars and then rejoice before the.
Lord your God for seven days.".
So 41, it says, "Celebrate this festival as a festival to the Lord for seven days each.
year, and this is to be a lasting ordinance for the generations to come..
Celebrate in the seventh month.".
42, it says, "Live in the booths, the tabernacles, the small tabernacles, for seven days..
All native-born Israelites, the people of God, are to live in booths..
So your descendants will know that I had the Israelite living booths when I brought them.
out of Egypt..
I am the Lord your God.".
So we think about this passage, okay, let's see what kind of lessons we can be taught.
about this passage..
So the first thing, I think, we, from this passage, we read that we have to learn to.
remember the grace of God in our life..
Now the Feast of Tabernacles was kept every year right after the autumn harvest..
So the crops are gathered in already in the barn and you can imagine, you know, in late.
September, the most comfortable year, the most comfortable season of the year, and then.
they gather all the crops and if your harvest pay off, how will you celebrate?.
You know, God says, "You keep a festival.".
So God is not telling them, you know, "After the harvest you have to celebrate, you need.
to keep a festival.".
And that's so strange, right?.
Because it says, "Beginning of the festival, after, after you have gathered the crops,.
after you have gathered your harvest, celebrate the festival before the Lord for seven days..
You pause and celebrate before the Lord.".
Isn't it counterintuitive?.
You know, right after you got your harvest, your paycheck, and then you have to celebrate.
before the Lord, not in the party..
You build a booth of branches and live there for seven days..
So as you can see, it's after you got the harvest..
So God is telling them, "When life is going well for you, it is so easy for you to forget.
God.".
And this is a common tendency of human heart..

$^{81}$The Feast of Tabernacles is placed there right after the harvest, right after you got your.
paycheck..
As if God is telling you, your greatest danger is not your lack of resources, it is not wilderness,.
it's not outside..
It is the moment that you hold God's blessing in your hand, gold and silver in your hand..
That is the most dangerous moment..
So the most dangerous spiritual moment is after you got the blessing..
So in Deuteronomy 8, Moses warned the people of God before they entered Canaan..
He said, "When you eat and are satisfied, your heart will become what?.
Become proud, and then you will forget the Lord your God.".
Isn't that the experience all we have?.
We make a lot of money..
We try to work hard and make a lot of money..
And then we etch God out of our life..
And you think about it..
It's very interesting because money can talk to us like God talk to us..
You think about it..
You hold the money in your hand or you check your iPhone, your app, your bank app, and.
then the money is whispering to you, you know, "Hey, you are safe and you are secure, you.
are at peace because geez, you got money, you got some numbers, right?".
And then your future is guaranteed..
But you know, all these things only God can offer you..
So the law, the scripture, move together, move into the life's rhythm..
The Feast of Tabernacles..
So once a year at the height of abundance and blessing, they got all the harvest and.
the law is calling the people of God back to remember the grace they first knew from.
God..
As you remember, this year is our church UV's 10th anniversary..
Now usually in church anniversaries, every single year, the most natural thing for us.
to do is to count or tally up our accomplishment or achievement, right?.
How many people in the attendance?.
But you know, what we need, what we need most is actually to return to the years that we.
got nothing, year one, day one..
So we think about day one as I was here and some of you were here, day one of UV..
We got nothing, we got nothing, we got nothing, we got no money, we got no staff, and we would.
never become a church without the grace of God and let alone last this long..
And I want to tell you in North America, most of the new churches, most of the new churches,.
church plans do not even last past three years..
So the first thing is to remember the grace of God on the height and after we receive.
blessing..

$^{121}$The second thing is to face and to understand that life is fragile..
So remember in verse 40 it says, "On the first day you take choice fruit from the tree and.
branches and rejoice before the Lord your God for seven days.".
So it says live in the booths for seven days and the Hebrew word for tabernacle or booth.
is sukkah and that is the name of a Jewish festival today is called sukkot..
So you know the Jewish people, especially the whereby, have a very interesting rules.
or interpretation of how to build a booth..
And then they are called, there's a later Jewish law called halakha and that means walking,.
you know walking the truth, you know how to practice it..
And here I list out some of the specific rules laid down by later Jewish rabbis..
So they said you need to celebrate the festival by building your booth in the backyard, right?.
And then here let me go through it..
First thing the whereby says it cannot be taller than 33 feet and you know why?.
But because if your booth in the backyard is too tall then it becomes a home, right?.
It becomes permanent..
So second thing is because a booth has to be temporary..
And then the law also says you know once you get too comfortable you will forget that it's.
a booth and it becomes another home for you..
And then the roof has to be made of natural plant materials..
That means it is easily run down and is easily you know you cannot keep the rain out..
And then the walls of your booth has had to be flimsy and no one should live in the booth.
for seven days and feel like oh are we living in the real house of stone..
[Chinese].
And then this rule if you think about it sounds very strange, right?.
And then but they carry some spiritual insight into it, right?.
Because it reminds you that life is fragile..
[Chinese].
And then we spend so much energy, so much money to build up our lives, right?.
You got insurance, you got your security and your stability and yet if you think about.
it one tree can wipe out or one storm can wipe everything out..
[Chinese].
Now recently you know we all experienced a storm, right?.
And then a tree fell and you know this is one of the pictures is a tree fell on a church.
staff member of another church, right? And then the whole car was crushed..
[Chinese].
And then let me make a note here..
Yesterday we were out of power in this area and we were thinking about you know we cannot.
meet here, right? Because without light and without power..
So we were trying to prepare for the online live streaming..
[Chinese].

$^{161}$So after all the preparation and she came back to the church and said, "Oh, the power is.
up again.".
[Chinese].
See life is unpredictable..
[Chinese].
Here I show you NBC Chicago even when the report asking the question whether the tornado.
alley has changed, it's shifting closer to Chicago..
You see we cannot even stop the climate changing, right? We are far more fragile than we think..
[Chinese].
See the flimsy sukkah is meant to communicate to us that your life is fragile..
[Chinese].
And why? Because God want his people to know that what truly protects you was never the.
strength of your own, not your booth, but not something that you built, but the God.
who dwell among us..
[Chinese].
If you think about it, our community, our church, our life is also fragile, you know,.
it's a reminder that without God we are done, we are game over..
[Chinese].
And so you need to pray for the church, you know?.
[Chinese].
Deuteronomy says, you know, God is protecting Israel in the entire journey of the wilderness..
[Chinese].
Now the third point we can learn from this festival is that we have to learn to pass.
on the mission of God..
[Chinese].
Now the heart of this festival and this passage is that so your descendants, the next generations,.
will know that I had the Israelites live in booths when I brought them out of Egypt..
I am the Lord your God..
[Chinese].
And the booth was never meant to be private..
It's not your own personal thing, it's a public thing, it's a living festival..
We enact every year and pass from one generation to the next..
[Chinese].
Now in the household of Jewish people today, they have this question and answer ritual when.
they live in the booth, you know, for seven days in their backyard and they have this.
question and answer kind of ritual and the children will ask them questions and the parents.
will answer the questions..
[Chinese].
Now the children may ask one of the questions is, "Why do we have to move outside to live.
in the backyard?".

$^{201}$[Chinese].
And the parents are to answer, "Because there was once a time we had no home and God was.
our home.".
[Chinese].
And the children may ask another question is, "Why does the roof have gaps that we can.
see the stars?".
[Chinese].
And the parents answer, "It is not broken..
The roof is not broken..
It's meant to be like that, intentional..
It's intentional because we need to be reminded that what truly covers us is not a roof but.
God.".
[Chinese].
And one more insight from the Festival of Tabernacles is that on the Sabbath day during.
the Feast of Tabernacles, now, today, the Jewish people will read one of the books from.
the Old Testament, which is the book of Ecclesiastes..
[Chinese].
And you remember the book of Ecclesiastes emphasizes the fact that life is meaningless..
It's like vapor, right?.
Life is meaningless..
Everything is meaningless..
Everything is fragile..
It ends up, you know, reminding us what is truly worth our pursuit is that we set our.
hearts on the eternal things for God..
[Chinese].
So no matter what you do today, you set your life direction, the ultimate direction is.
to honor and to serve God, you know?.
[Chinese].
Now, some people once asked me, "If you could go back 10 years ago, before you restarted,".
you know, it was a really, really hard time for me..
And people asked me, "Would you do the same?.
Would you make the same decision?".
[Chinese].
And I answered, you know, without hesitation, "Of course, yes!".
[Chinese].
And those were very difficult days, but I always carry a sense of, you know, gratitude.
and thanksgiving..
You know, not everyone, you know, not everyone gets a chance to build up the church from.
nothing, right?.
To start a church, brick by brick..

$^{241}$[Chinese].
And I believe you too, right?.
God is calling you to join us, and I believe that, you know, we carry the eternal value.
in us, and we want to go forth and seek the kingdom of God, and that's the first and.
foremost thing in our life..
[Chinese].
Now in a moment, I'm going to give you, in this series, next series, I'm going to give.
you an item, right, to remember this series, this big story of God..
And I'm going to give each one of you a small piece of bamboo, a very small piece of bamboo..
It's a simple and tangible reminder for you and for me, just like Israel used branches.
to build the Petermonachos in the backyard..
[Chinese].
Now you bring this home and you ask yourself a few questions..
[Chinese].
The first question I suggest is, how have I experienced the grace of God recently?.
[Chinese].
You know, how have you experienced the blessing of God recently?.
Second, how should I respond to the faithfulness of God?.
[Chinese].
The third one is that, you know, as you think about it, you know, God has kept you in the.
wilderness and taught you and show himself to be trustworthy after all..
[Chinese].
And the third question is, how can I pass on the Word of God to others and to the next.
generation?.
[Chinese].
Now take this piece of item home and put it in the jar, right?.
We have a jar prepared for you, right?.
[Chinese].
Take this home and this week and sit with the questions that I suggested..
[Chinese].
Okay, shall we pray together?.
\newpage



\section{出埃及記}
\label{sec:pCh24jxVqNk}
\textbf{The Tent (Exod 25-40) - Rev. Mark Chan 會幕 (出埃及記25-40) - 陳崇基牧師}
\newline
\newline
連結: \href{https://youtube.com/watch?v=pCh24jxVqNk}{\texttt{https://youtube.com/watch?v=pCh24jxVqNk}} ~~~~ 語音日期: 2026-07-12
\newline
\newline
\hyperref[sec:augzRkvXpyY]{\small{< < < PREV SERMON < < <}}
~
\hyperref[sec:index_chronic]{\small{[返順時目]}}
~
\hyperref[sec:index_scriptual]{\small{[返順卷目]}}
~
\hyperref[sec:Cofy9lMF2CI]{\small{> > > NEXT SERMON > > >}}
\newline
\newline
$^{1}$And today we'll have Pastor Mark..
Good morning everyone..
Very good to see you and let's turn to our series on next and this is the fourth sermon.
in our next series and remember this jar?.
You know you got some items here and we wanted what we want to do in this series to trace.
the grand story of the Bible from beginning to the end from Genesis to Revelation and.
every you know week and we push on this sermon we find that the same God is being introduced.
in the entire Bible in the entire story which is a God of mission and God is reaching out.
to us and he is not waiting for us to reach out to him..
And I think to most people when you think about God or gods you know or even for the.
ancient people the gods and the goddesses especially in the ancient Near East people.
would think that we have to build massive temples and like you know God is on the mountain.
and we have to perform we have to do something to reach out to God..
In much of the same way if you have some religious background especially you know in my case.
Chinese folk religion background the gods as people imagine them are distant and transactional.
and the gods stay up there on the mountain and you have to go up there you have to climb.
and that is the arrangement between God and man..
But as we come to our passage today 25 Exodus 25 you find that there is a totally different.
concept of God the picture is totally different because God says in Exodus 25 it says God.
tell the people of Israel to have them make a sanctuary make a home for me and I will.
dwell among them and make this tabernacle the tent so all the furnishings exactly the.
pattern I will show you..
So you see the contrast the totally different concept of who God is you know the God of.
the Bible is the God who dwell among us he comes from above to dwell among us so it's.
not me to build a temple so I can climb up there to meet God but that I God says I want.
to live with you and I want to come down I want to descend and that is the most radical.
statement in religion..
So we come to a really long section in the Bible in the book of Exodus 25 to 40 and remember.
God promised the patriarch Abraham that he would have descendants like the number of.
grains of sand and then God saved the people out of Israel and brought them to the mountain.
and gave them the law remember the stones and as they come to Sinai as the new people.
they are being rescued from Israel they receive the laws..
And now the biggest wonder is that God wants to build a tabernacle a home among them you.
know a tabernacle is simply a portable tent you know have you been to a camping trip before.
yeah you have to pitch a tent right and then pitch a tent to next each other and you heard.
what people said you know in the middle of night and you know kind of like camping right.
because in the wilderness as people move forward in the wilderness God lives in the tent among.
them to guide them in every step of the way..
So in the wilderness you know among the people of Israel you see there's a tent right there.
and if you're new to the Bible you might think you know why does an all-powerful God cares.

$^{41}$so much about you know the tent and gives instructions for these curtain dimensions.
you know every piece every material every piece of furniture inside the tent you may.
wonder this..
But every single piece inside this tabernacle is a picture of how God wants to relate to.
you relate to his people and uh it's a it's a grammar of relationship architectural grammar.
of relationship..
Now if you go to Israel today you drive south all the way south almost to the Red Sea and.
you see a one-to-one tabernacle model in Timna Park and you know I want to show you.
the the floor plan of the tabernacle and the tent is made of different color of curtains.
you know there's a red and purple all kinds of curtains and there are different furniture.
items you know you look at it from top down right it's rectangular in shape and as you.
move inward from the courtyard from the right the courtyard and then the shape becomes more.
perfect square right the holy of holies is perfect square which means in in the ancient.
context means perfection and holiness..
And then uh first of all we come to the uh courtyard which is you know the area on the.
right you see the bronze altar it stands at the uh entrance and the altar is simply is.
simply a a barbecue grill all right it stands there to remind the people that sin is real.
and it costs you something you have to bring offering you know burn offering to approach.
God but God provides you a way and the sacrifice is God's provision..
And then after the altar is the uh bronze basin which is a water basin for you to cleanse.
yourself and then you come to the uh smaller uh rectangle uh which is uh the holy place.
and then you find different items in there uh first items you you will find is on your.
right is a table of showbread and there were twelve loaves of bread and one for each tribe.
of Israel that's why there were twelve and placed before God every single week to remind.
the people God is feeding you is sustaining you..
And then uh you find the other item of the three is one on the left which is the lamp.
stand the golden lamp stand and it's it burns perpetually and uh to communicate to you that.
God is light and uh he is not hiding..
And then the third item is the order of incense and you burn something there and the smoke.
will go up and represent prayer go going up and God is actually listening to the prayers.
of the people..
And then at the heart of everything is the perfect square right the perfect square which.
is called the holy of holies and is behind the thick curtain and the most holy place.
the most important item is the box is the ark of covenant is a gold covered box and.
holding different items in there especially the stone tablets the two tablets of the covenant..
And on top of the ark box is two golden cherubim uh angels with this with the wings spread.
uh toward each other and that's actually the mercy seat the throne of God..
Now because God is holy and we are sinful so no one can approach the most holy place.
uh casually and the high priest that's important the high priest can enter the most holy place.
the God's very home once a year on Yom Kippur on the day of atonement to make sacrifice.

$^{81}$for the sins of the whole people and he entered the most holy place which is the mercy seat..
And I spent so much time explaining the the the blueprint or the fore plan of the uh the.
tabernacle I feel like a real estate agent trying to show you the the fore plan of the.
tabernacle..
And I want to tell you is is all of these details are architectural theology every piece.
in the of the items communicates to you that God forgives God provides and God listens.
to your prayer and God is present God is not a distant you know deity and he is a God who.
wants to make his home among his people..
In chapter 29 uh it says God says I will dwell among the Israelite my people and be their.
God and they will know that I'm the Lord their God who brought them out of Egypt so I might.
dwell among them..
So you have to understand Exodus coming out of Egypt is not just rescue it is relationship..
So not only God desire to dwell among us the people God is determined to do so God is determined.
to do that..
So let's return to the passage and you know after chapter 25 to 31 which is the detailed.
blueprint of the tent and something really bad happened in between 32 and 34..
Now you have to kind of imagine you were there and Moses has been up on the mountain receiving.
the blueprint for 40 days and the people at the foot foot hill and they were getting kind.
of restless and impatient and they want something like you know the force gods that they used.
to worship so they have Aaron Moses brother melt the gold and cast a golden calf and they.
come alongside and surround the the calf and make it an idol for them and they worship.
the idol..
So right in the middle of the story 32 to 34 um you see that you know people are being.
sinful and they turn and rebel against God and I don't want you to feel kind of okay.
we are better right we are better people I mean the people they wouldn't screw up right.
but you know in some way we are better I don't want you to think about that because that.
is a universal human problem for you and for me..
We have the same problem because God is invisible and require our trust our faith and yet sometimes.
we are impatient and we want to fall before a visible God you know an immediate God and.
we want to manipulate because we want the idols..
And you think about it you know we do that every single day we we worship the force God.
we we uh we worship the golden calf because we want to replace our trust in the invisible.
God with something visible money the self-sufficiency some of the version of of the gods that would.
uh fit our one you know uh so the golden calf story is actually not an ancient story it's.
a mirror of our human nature..
And sometimes our gods are money sometimes our gods are sleep you know we just want to.
sleep more right want to be lazy right..
We are all idol worshiper and that is the point of the story right in the middle right.
in the structural center and um uh up on the mountain God is still you know giving the.
grouping and down in the valley uh the people are already trying to replace God with idols..

$^{121}$So right in the middle it tells a anti-climax story and the covenant with God is shattered.
is broken so Moses comes down and flown the uh the tablets on the ground and broke them..
And then God was really angry because you know God want to wipe out the people but at.
this point at this crisis moment uh Moses goes up to the mountain and intercede for.
the people and says Lord if you destroy all your people and then your faithfulness is.
on the line don't do that..
Right here in the middle of the entire story there comes a great theological statement.
and Moses says the Lord the Lord a God merciful and gracious slow to anger abounding in steadfast.
love and faithfulness keeping steadfast love for thousands not one not two thousands of.
generations forgiving iniquities..
So now the crisis is resolved right the people are being saved and redeemed not because we.
are good not because we are so bad but God chooses to stay God chooses to forgive and.
God is good..
So you see that God's determination to dwell among his people even though they are broken.
and they are rebellious but after redemption God's still determined to dwell among them.
because that is the character of God he's gracious he's merciful he's slow to anger.
and then part three the story continues from chapter thirty five to the end chapter forty.
is God's amazing invitation to invite sinners to build his home among them..
And then now something I want you to notice is that thirty five to forty is actually the.
mirror of the first part of the story you see that first part and then there's a middle.
and then there's the end and then you see some of the details some of the verses from.
the first part and third part they are actually mirrored to one another..
And then let me give you an example twenty five thirty one it says make a lamp stand.
is a command right it's the blueprint command and God is saying to people make a lamp stand.
of pure gold and hammer it out bays and shafts and all the shapes and be one piece with it.
that's a command right twenty five thirty one and then on the right is thirty seven.
seventeen it's the execution of the command and he the craftsman the craftsman and he.
also made the lamp stand of pure gold almost word to word you know same thing but on the.
left is command on the right is execution..
So almost identical first part and the third part word for word and God said it and they.
did it and what God has designed the community now builds it and then the blueprint becomes.
the building the word becomes work..
You wonder and it's very interesting that God does not send angels the holy people to.
build his home he chooses his very people who has just betrayed him and that is not.
just forgiveness that is restoration give them the honor and dignity to build his home.
among them..
They you know they and then the people bring the offering the gold and the silver to build.
the home and they bring so much that Moses has to tell the people hey you got to stop.
because there's so much and there's more than enough you got to stop and make your.
offering..

$^{161}$What a picture of grace that God wants his sinful people broken people to build his home.
among them..
Now we at UV we are in a season of building as well you know God is calling us to expand.
right it's the next we have got to go next we got to make room for more people who have.
never met our God and for God to live among us among us is that the work will becomes.
our work the whole community has to work together and the qualification here the Bible tells.
us the qualification to build the church to build a community for God is not because your.
most talented it's not because you're most experienced it's simply because you're willing..
And we you are invited to give generously you're invited to serve willingly you're you're invited.
to show up you're invited to welcome to create a church that will make people feel like this.
God is different and this God is will..
And back to the story in the last chapter at the very end in the conclusion it says.
then the cloud the glory cloud covered the tent of meeting the tabernacle the glory of.
the Lord filled the tabernacle and Moses could not enter the tent because the cloud the glory.
is so much and settled on it and the glory of God filled the entire home..
You see God moved in after the rebellion after people make idols for themselves yet he redeemed.
them and he came down and he stayed and he filled his home with his glory..
You see so often we we have this you know our our own conception especially for new people to the church they they feel like you know I'm not ready to go to church you know because I'm not perfect..
You see we suppose that the church is a place for a community of perfect people right..
Sometimes not only the new people think so sometimes we think so we've been to church.
a lot right and we tend to think that because we are better we come to church but that is.
not the case right..
With that kind of grace you need to embrace God as well..
See the whole movement of Exodus is telling you God in his grace is drawing near to us.
even though we are broken we are dirty it's always not further away the Holy God drawing.
nearer and nearer to us..
It is no wonder centuries later many years later the apostle John opens up his gospel.
and as he introduced us to us introduced Jesus to us he says in chapter 1 verse 14 it says.
the word who is Jesus okay the word is Jesus he became flesh and made his dwelling among.
us and we have seen his glory the glory of one and only son who came from the father.
above and full of grace and truth..
See that word translated dwelling is the same word tent tabernacle and here the gospel of.
John is deliberately trying to echo the tabernacle trying to echo the Old Testament story Jesus.
now is the real tabernacle in person..
See Jesus is the place when we come to Jesus when we come to Christ Jesus is the place.
where heaven and earth meet together in the old covenant God's presence fills a physical.
tent now in Jesus God pitched himself his home among us in Christ..
See through Christ God does not just dwell near us he dwells among us in us..
So I want to tell you this thing in conclusion you know the world most of the people in the.
world don't know Christ you know you have to think that in our world what we need is.

$^{201}$not a church of perfect people better people the world does not need that the world needs.
to see that we as church people we have experienced we have encountered this God of grace he has.
come down to us even though we are broken but he come down to us and love us and embrace.
us..
And what a privilege that we are being called sinners fully redeemed we are being called.
to build a home for God so that God is present among us..
So before we pray to end our service today I want you to pick up one item before you.
go okay before you go pick up this item it's a piece of cloth and the tabernacle was built.
as you remember was built with fine linen and fabric without the cloth there's no tabernacle..
So I want you to take one piece of cloth and add it to your sand jar of the year and you.
have sent the promise of God the puzzle the peace to remind us we are broken yet God has.
chosen to use us and there's some stones in it there's a third item remember remember.
the new principle the law has given to us and now is the cloth is the fabric of God's.
dwelling it's a reminder that you and I are builder build this community so that God lives.
in us you are the material you are the material God is using to make his presence in our midst.
so in this week I want to invite you in this week every day of the week I want to invite.
you in this week every day of the week seven days before we come back to worship every.
day you pray a risky prayer okay risky prayer I tell you what the prayer is and what that.
risky prayer says is this God use me to build your church and pray to God somehow you can.
use me every day of the week God please use me okay I'm sure God will reveal and move.
you to be used by him to build a tabernacle build a home for him shall we pray together.
Father thank you you are not the God who stays at distance you came down and dwell among us.
in a tent and you came in flesh in Jesus for grace and truth and by your spirit we now.
you now dwell in us and fill this place as well fill our church with your glory just.
as you filled the tabernacle build your church through us our finances our skills our presence.
for anyone here who has never known you and let today be the day they discover they find.
that you were looking for them long before they ever looked for you thank you so much.
for your grace in Jesus' name we pray amen.
\newpage



\section{士師記}
\label{sec:Cofy9lMF2CI}
\textbf{When Less is More (Judge 7) - Chee Yie Jong 少即是少 (士師記 7) - 楊志溢傳道}
\newline
\newline
連結: \href{https://youtube.com/watch?v=Cofy9lMF2CI}{\texttt{https://youtube.com/watch?v=Cofy9lMF2CI}} ~~~~ 語音日期: 2026-07-12
\newline
\newline
\hyperref[sec:pCh24jxVqNk]{\small{< < < PREV SERMON < < <}}
~
\hyperref[sec:index_chronic]{\small{[返順時目]}}
~
\hyperref[sec:index_scriptual]{\small{[返順卷目]}}
~
\hyperref[sec:qC9_GTWJyK4]{\small{> > > NEXT SERMON > > >}}
\newline
\newline
$^{1}$(speaking in foreign language).
- Good morning, Yuvi..
(speaking in foreign language).
Testing, testing..
All right, so a couple weeks ago,.
I attended a graduation ceremony.
for one of my students from Summit..
And if we're being honest,.
most of the time when we attend a graduation.
or commencement ceremony,.
what we really care about is the presentation.
of the diploma or the degrees, right?.
Which unfortunately usually happens at the end..
But during the ceremony,.
so in addition to the main speakers,.
they actually have five or six more students.
giving different kind of speeches..
(speaking in foreign language).
And so when it got to the third or the fourth speakers,.
I started noticing that some of the younger kids.
in the room, they began to lose their patience..
And their parents, on the other hand,.
were struggling to try to keep them quiet..
Some families, they even had to leave the room entirely..
And so by the time all the speeches ended.
and the presentation of diplomas finally began,.
more than one hour had already passed..
(speaking in foreign language).
And so this experience sort of reminds me.
that in a world where we usually assume.
that more is always better,.
more money, more power, more efficiency,.
more likes on our Instagram account or a post,.
sometimes less is actually more..
Like a few selected speeches at this graduation ceremony.
or a little bit of salt on a chocolate chip cookies..
And so today we're going to see that less is more..
It's actually not just a good advice..
It is actually a principle that is deeply woven.
into how God chooses to work sometimes..

$^{41}$(speaking in foreign language).
, so we've seen that in the past..
We have been journeying through the book of Judges so far.
and currently we are in the cycle of Gideon..
Last time we saw that when God called Gideon to save Israel,.
he was very hesitant because he let his fears hold him back.
from surrendering to this calling..
And after a few rounds of back and forth.
between Gideon and the Lord,.
Gideon is finally ready to do what God has called him to do.
here in Judges seven,.
as he moves into position to fight the Midianites.
and the coalitions..
But before that, God has one more thing for Gideon..
So let's look at Judges seven together today..
(speaking in foreign language).
- So verse one, "Early in the morning, Jerubbaal,.
"that is Gideon, and all his men camped.
"at the spring of Harol..
"The camp of Midian was north of them in the valley.
"near the hill of Moreh..
"The Lord said to Gideon, 'You have too many men..
"'I cannot deliver Midian into their hands,.
"'or Israel will boast against me..
"'You have too many men.'.
"Jerubbaal said to Gideon, 'You have too many men..
"'I cannot deliver Midian into their hands,.
"'or Israel will boast against me..
"'Saying, 'My own strength has saved me.'".
(speaking in foreign language).
- And so as Gideon and his army,.
they are getting ready for the battle,.
the Lord finds a problem, that they have too many men..
Think about it, why is God so concerned with this,.
that he literally has to interrupt.
Israel's battle preparation?.
If we take a closer look,.
the number of men is actually not the root cause..
The real problem here, as verse two tell us,.
is that Israel would boast against the Lord,.

$^{81}$and kind of like credit the victory.
as their own strength that brings them the victory..
(speaking in foreign language).
So this is not a numerical problem,.
but it is actually a hard problem..
And so God has to deal with it before the battle begins..
And so to deal with these root issues,.
the Lord gave them a special plan..
A less is more kind of plan..
Less dependence on the numbers or their own strength,.
and more dependence on the Lord..
(speaking in foreign language).
So in our passage today,.
we're going to look at this main idea together,.
that self-sufficiency, that is relying on our own strength,.
is the enemies of God's grace..
Again, self-sufficiency is the enemy of God's grace..
And as we go through our passage,.
we're going to explore why God has to deal.
with Israel's self-sufficiency.
before they enter into this battle..
So first, self-sufficiency is the enemy of God's grace.
because it can lead to pride..
(speaking in foreign language).
So ironically, while God is concerned.
about Israel's pride here,.
Gideon himself, on the other hand,.
he is still battling with his own fear,.
which is a consistent threat throughout his own life..
So while a large army might seem exactly.
like what Gideon needs at this point.
to build up his confidence,.
God has a very different plan in mind..
And so he told Gideon,.
"Now announce to the army,.
"anyone who trembles with fear.
"may turn back and leave Mount Gilead.".
So 22,000 men left, while 10,000 remain..
(speaking in foreign language).
So in this first round of cuts,.

$^{121}$the deciding factor is those who trembles with fear..
And Gideon's men are being reduced.
from 32,000 to just 10,000..
(speaking in foreign language).
So while God is mainly concerned with Israel's pride,.
at the same time, he does not forget.
to slowly restore Gideon's faith..
By telling Gideon to send the fearful.
and the trembling Israel home,.
Gideon, whose own past is characterized.
by all the fears and doubts,.
he is now being forced to differentiate himself.
from this group of fearful ones.
that were being sent home..
(speaking in foreign language).
However, God is not done yet..
In verse four, the Lord said to Gideon,.
"There are still too many men..
"Take them down to the water,.
"and I will thin them out for you there..
"If I say, this one shall go with you, he shall go..
"But if I say, this one shall not go with you,.
"he shall not go.".
So Gideon took the men down to the water..
There, the Lord told him,.
"Separate those who lap the water with their tongues.
"as dog laps from those who kneel down to drink.".
(speaking in foreign language).
And then verse six,.
300 of them drank from cup hands, lapping like dogs..
All the rest got down on their knees to drink..
And the Lord said to Gideon,.
"With the 300 men that lapped,.
"I will save you and give the million knights.
"into your hands..
"Let all the others go home.".
(speaking in foreign language).
So in this second round of cuts,.
the Lord asked Gideon to separate those who lapped.
like a dog from those who kneel down to drink..

$^{161}$(speaking in foreign language).
So there are a few possible explanation.
for why the lepers were chosen instead of those who kneel..
(speaking in foreign language).
Some scholars think that the lepers,.
they actually show less fear.
because they're not afraid to just put their guard down.
and drink water like a dog,.
while the kneelers, they stay on guard the entire time..
And this explanation has to do.
with who actually show less fears,.
kind of like in continuation.
with the first round of the cuts..
(speaking in foreign language).
Some other scholars argue that.
it is because actually the lepers,.
they show less situational awareness,.
which makes them worse of a fighter..
But this is exactly why God chooses them.
to show that he can still win the battle.
even with less able fighter..
And so this second explanation has to do with.
who gets the credit for the victory at the end..
(speaking in foreign language).
For me, while both explanations seems plausible,.
the truth is that just we simply do not know..
And the bigger focus here should be the fact.
that Gideon's army is further reduced down to 300 men..
If you do some quick math,.
this is actually less than 1\% of the men.
that he had in the beginning,.
from 32,000 to just 300..
(speaking in foreign language).
So why did God have to make such a drastic cut.
to Gideon's army?.
Again, this goes back to verse two,.
so that the Israel would not think.
that their own strength had actually saved them..
In other words, this is to prevent Israel.
from becoming prideful because of their own strength..

$^{201}$(speaking in foreign language).
So God strips Gideon's army all the way down to 300 men..
On one hand, it is to prevent Israel's self-sufficiency.
from turning into pride..
But for Gideon, God is actually setting him up.
to experience something that he never could have.
with his 32,000 men..
And so this bring us to the second part of the story,.
which is how Gideon's army was saved..
(speaking in foreign language).
This leads us to my next point,.
that self-sufficiency doesn't just lead to pride..
It actually robs us from experiencing God's sufficiency..
(speaking in foreign language).
So imagine Gideon now losing more than 99\% of his warriors,.
and he still have a battle to fight the following day..
Gideon probably feels a lot more anxious.
about the battle now compared to the beginning of the day.
when he was still 32,000 men strong, right?.
However, that same night,.
the Lord once again reassures Gideon.
that he is going to fight the battle..
(speaking in foreign language).
He's going to give the Midianites into his hand..
(speaking in foreign language).
So let's look at verse nine..
(speaking in foreign language).
During that night, the Lord said to Gideon,.
get up, go down against the camp,.
because I am going to give it into your hands..
If you're afraid to attack,.
go down to the camp with your servant, Pura,.
and listen to what they are saying..
Afterward, you will be encouraged..
(speaking in foreign language).
So he and Pura, his servant,.
went down to the outpost of the camp..
And verse 12, the Midianites, the Amalekites,.
and all the other eastern people.
had settled in the valley thick as locusts..

$^{241}$The camels could no more be counted.
than the sand on the seashore..
(speaking in foreign language).
Midian, Amalekites, and the Midianites.
(speaking in foreign language).
This time, before Gideon even have to ask for any signs,.
which is kind of like his signature move at this point,.
the Lord already knows his worries,.
and he reassures him in verse 10,.
if you are afraid to go, go to the enemy camp tonight,.
and you will eventually find a sign of victory..
(speaking in foreign language).
So what God is doing here is that he is extending.
an invitation for Gideon to experience him,.
despite Gideon's own fear,.
as well as the fact that his men are just very little.
compared to the enemy..
(speaking in foreign language).
So verse 13, Gideon arrived just as a man.
was telling a friend his dream..
"I had a dream," he was saying..
"A round loaf of barley bread came tumbling down.
"into the Midianite camp..
"It struck the tent with such force.
"that the tent overturned and collapsed.".
His friend responded, "This can be nothing.
"other than the sword of Gideon, son of Joash,.
"the Israelite..
"God has given the Midianites.
"and the whole camp into his hands.".
(speaking in foreign language).
So this time, the sign of victory that God gave Gideon,.
it actually comes from a dream from the enemy camp..
(speaking in foreign language).
And logically, like a loaf of bread rolling into a camp,.
Gideon should pose no real threat to an army this size..
And yet somehow, even before the battle begins,.
just the mere reputation of Gideon,.
has already struck fear into the entire enemy camp..
(speaking in foreign language).

$^{281}$And this final sign is ultimately enough for Gideon.
to believe that God has indeed given him.
the sword of Gideon..
(speaking in foreign language).
To deliver the Midianite into his own hand..
(speaking in foreign language).
So verse 15, "When Gideon heard the dream.
"and its interpretation, he bowed down and worshiped..
"He returned to the camp of Israel and called out,.
"Get up, the Lord has given the Midianite camp.
"into your hands.".
(speaking in foreign language).
So we don't have time to finish the entire chapter here,.
so I will actually encourage you to read it on your own..
But in the second half of this chapter,.
with just 300 men, Gideon defeats the Midianites.
as the Lord has promised..
And the weapon of choice?.
Not iron chariots, not spears, and not swords,.
but just trumpets, some torches,.
and some jars that they crush..
And it actually works, because the victory ultimately.
was never about Israel's own strength to begin with..
(speaking in foreign language).
So this once again shows that the only reason Gideon.
is able to defeat the Midianites against all odds.
is because of the Lord..
(speaking in foreign language).
So the truth here is that God alone is sufficient.
to meet all of our needs..
Even though Gideon lacks courage, God is enough for him..
Even though the Israel's army is laughably small,.
God is sufficient to secure the victory..
He strips away every layer of human confidence.
until all that is remained is total dependence on him..
And church, this is actually an act of grace.
so that we can experience his sufficiency in our own lives..
(speaking in foreign language).
So I don't know who in this room needs to hear this today,.
but here's the good news..

$^{321}$God is enough for Gideon..
(speaking in foreign language).
God is enough for Gideon..
(speaking in foreign language).
God is enough for Gideon..
(speaking in foreign language).
God doesn't need us to be perfect for him to use us,.
but his power is made perfect in our weakness.
because his grace is sufficient for us..
And God doesn't need us to be self-sufficient,.
but what he needs from us is to be able to open.
to receiving his grace..
(speaking in foreign language).
So maybe some of you find yourself.
in a similar situation with Gideon today,.
that God is stripping away something in your life,.
a relationship, a healthy body, a career,.
a sense of control, or a confidence.
that you have been quietly holding on to..
(speaking in foreign language).
God is enough for Gideon..
(speaking in foreign language).
God is enough for Gideon..
(speaking in foreign language).
God is enough for Gideon..
(speaking in foreign language).
God is enough for Gideon..
(speaking in foreign language).
God is enough for Gideon..
(speaking in foreign language).
God is enough for Gideon..
(speaking in foreign language).
Because this might be a gentle invitation from the Lord.
for you to experience him in ways.
that you are not able to in your self-sufficient state,.
when your career is successful,.
or your relationship is going well,.
or your retirement account is looking great..
Because more than our comfort,.
God is more concerned with our Christlikeness..

$^{361}$And more than our happiness,.
God is more concerned with our holiness..
(speaking in foreign language).
And so instead of asking the Lord,.
Lord, when will this end?.
Try asking him, Lord, what are you forming in me?.
Because God may be doing his deepest work in you.
right there, right then..
He may be tearing down your self-sufficiency.
so that his own sufficiency can finally be seen..
(speaking in foreign language).
This is why the gospel can be so offensive.
to self-sufficient people..
Because in our self-sufficient economy,.
we think that we have to be good enough,.
or strong enough, or capable enough.
in order to save ourself..
But that is not the point of the gospel..
(speaking in foreign language).
The gospel, on the other hand,.
says that you don't have to be good enough,.
you don't have to be strong enough,.
or you don't have to be capable enough..
Because at the cross, God provided the ultimate solution.
that we could never accomplish ourself..
That is salvation by grace through faith..
And it is only when we are able to understand.
that we alone are not sufficient.
that we are able to experience this gift of grace from God..
Grace begins where self-sufficiency ends..
(speaking in foreign language).
So let's pray..
(speaking in foreign language).
Lord, thank you that your grace is sufficient for us,.
and your power is made perfect in our weakness..
(speaking in foreign language).
We thank you that our salvation.
is not earned by our own strength,.
but given freely through your grace..
(speaking in foreign language).

$^{401}$So Lord, help us to live each day.
with humble dependence on you..
(speaking in foreign language).
In your name we pray, amen..
(speaking in foreign language).
\newpage



\section{約翰福音 10:27-30}
\label{sec:qC9_GTWJyK4}
\textbf{一雙永不鬆手的手 (約翰福音10\_27-30) - 梁紫健傳道 The Hands that Never Let Go (John 10\_27-30) - Tsz Kin Leung}
\newline
\newline
連結: \href{https://youtube.com/watch?v=qC9_GTWJyK4}{\texttt{https://youtube.com/watch?v=qC9\_GTWJyK4}} ~~~~ 語音日期: 2026-07-12
\newline
\newline
\hyperref[sec:Cofy9lMF2CI]{\small{< < < PREV SERMON < < <}}
~
\hyperref[sec:index_chronic]{\small{[返順時目]}}
~
\hyperref[sec:index_scriptual]{\small{[返順卷目]}}
~
\hyperref[sec:xpOSI0B_Qx0]{\small{> > > NEXT SERMON > > >}}
\newline
\newline
約翰福音 10:27-30
\newline
\begin{longtable}{cl}
\hline
\hline
章節 & 經文 (和合本修訂版)\\
\hline
10:27 & \begin{tabularx}{0.7\textwidth}{X} 我的羊聽我的聲音，我認識牠們，牠們也跟從我。 \end{tabularx} \\ \\ \relax
10:28 & \begin{tabularx}{0.7\textwidth}{X} 並且，我賜給他們永生；他們永不滅亡，誰也不能從我手裡把他們奪去。 \end{tabularx} \\ \\ \relax
10:29 & \begin{tabularx}{0.7\textwidth}{X} 我父所賜給我的比萬有都大，誰也不能從我父手裡把他們奪去。 \end{tabularx} \\ \\ \relax
10:30 & \begin{tabularx}{0.7\textwidth}{X} 我與父原為一。」 \end{tabularx} \\ \\
[1ex]
\hline
\hline
\end{longtable}
$^{1}$[NON-ENGLISH SPEECH].
Good morning, brothers and sisters..
[NON-ENGLISH SPEECH].
Happy Father's Day..
[NON-ENGLISH SPEECH].
Some of you may remember what it was.
like to have your dad hold your hand when you cross the road..
That's a child..
[NON-ENGLISH SPEECH].
And that is kind of like a big hand wrapping entirely.
around your little hands..
And you had nothing to fear, because you know.
that he would never let go..
[NON-ENGLISH SPEECH].
But not everyone have warm memories when.
it comes to their father's hand..
[NON-ENGLISH SPEECH].
Some remember a hand that was never present..
[NON-ENGLISH SPEECH].
Others remember the hands that would lash out or even hurt.
others..
[NON-ENGLISH SPEECH].
And for some of us, what they miss the most on these Father's.
Days is a hand that they can no longer hold..
[NON-ENGLISH SPEECH].
So our passage today is short, just four verses,.
but the word "hands" here appears twice..
[NON-ENGLISH SPEECH].
And through this verse, Jesus wants.
to reveal to us what kind of Father our Heavenly Father is..
[NON-ENGLISH SPEECH].
And what kind of hands he has..
[NON-ENGLISH SPEECH].
And by the end of the sermon, we will know these three things..
[NON-ENGLISH SPEECH].
First, whose hands does the Father's hold?.
[NON-ENGLISH SPEECH].
And second, why does Jesus say that no one.
is able to snatch them out of the Father's hand?.
[NON-ENGLISH SPEECH].

$^{41}$And third, what does the phrase "I and the Father are one".
have to do with Father's Day?.
[NON-ENGLISH SPEECH].
If you're ready, we'll read our passage today..
[NON-ENGLISH SPEECH].
My sheep hear my voice, and I know them,.
and they follow me..
I give them eternal life, and they will never perish..
And no one will snatch them out of my hand..
My Father, who has given them to me, is greater than all..
And no one is able to snatch them out of the Father's hand..
I and the Father are one..
[NON-ENGLISH SPEECH].
So let's first look at the context.
in which Jesus gave these words..
[NON-ENGLISH SPEECH].
Verse 22 tells us that it was the festival of dedication.
during the wintertime..
[NON-ENGLISH SPEECH].
So the festival of dedication actually.
commemorates the Maccabean-era purification of the temple,.
when the temple, defiled by a foreign king,.
was rededicated as holy to the Lord..
[NON-ENGLISH SPEECH].
And it was at this festival that a group of Jews.
surrounded Jesus and demanded him,.
how long will you keep us guessing?.
If you're the Christ, just tell us plainly..
[NON-ENGLISH SPEECH].
Jesus did not answer them directly with I am or I'm not..
Instead, he spoke about sheep..
[NON-ENGLISH SPEECH].
And this raises our first question today..
[NON-ENGLISH SPEECH].
So why didn't Jesus simply answer them yes or no.
to the are you the Christ questions?.
But instead, he spoke of my sheep hear my voice..
[NON-ENGLISH SPEECH].
In fact, his very way of answering it.
is in and of itself the answer..

$^{81}$So let us follow the passage questions by questions,.
beginning with the first one..
[NON-ENGLISH SPEECH].
Jesus' answer to the questions, are you the Christ,.
was not a theological argument, but a description.
of a relationship..
[NON-ENGLISH SPEECH].
The three verbs in verse 27 are beautiful..
The sheep hear my voice, I know them, and they follow me..
[NON-ENGLISH SPEECH].
In the original language, hear is not just.
mere physical hearing, but discerning and hearing..
[NON-ENGLISH SPEECH].
Like a child in a crowd that instantly recognized.
their father's voice..
[NON-ENGLISH SPEECH].
And know is not just knowing some facts about someone,.
but a deep knowing found in relationship..
[NON-ENGLISH SPEECH].
So Jesus already said in verse 14 and 15,.
I know my sheep, and my sheep know me,.
just as the father knows me, and I know the father..
[NON-ENGLISH SPEECH].
In other words, the knowing between the sheep.
and the shepherd is actually rooted.
in the knowing between the father and the son..
[NON-ENGLISH SPEECH].
We are known by God because of the deep mutual knowing that.
has already exists between the father and the son..
[NON-ENGLISH SPEECH].
And now, once we recognize who the shepherd is,.
the second questions follow..
How does this shepherd keep his sheep?.
[NON-ENGLISH SPEECH].
And why does the word "hen" appear twice.
in such a short passage?.
[NON-ENGLISH SPEECH].
So let us begin from the climax here in verse 28..
"I give them eternal life, and they will never perish,.
and no one will snatch them out of my hand.".

$^{121}$[NON-ENGLISH SPEECH].
So two words here, no forwarding..
[NON-ENGLISH SPEECH] First, "will never perish".
is a double negative, an emphatic expression.
meaning absolutely forever, they will not perish..
[NON-ENGLISH SPEECH].
Second, "snatch," the word, is the same word.
used in verse 12, describing a wolf seizing the sheep.
and scattering the flock..
[NON-ENGLISH SPEECH].
In other words, Jesus is not saying that there is no danger.
or no fierce wolf..
[NON-ENGLISH SPEECH].
What he's saying is no matter how great the wolf power is,.
it cannot seize them..
[NON-ENGLISH SPEECH].
Then notice this..
In verse 28, it says, "out of my hand.".
And verse 29, it says, "out of the Father's hands.".
[NON-ENGLISH SPEECH].
Why does Jesus have to say the same thing twice,.
once using my hand and the other time using the Father's hand?.
[NON-ENGLISH SPEECH].
The answer can be found in verse 30..
I and the Father are one..
[NON-ENGLISH SPEECH].
This is not two hands in the tug of war..
It is the same grip and the same power..
[NON-ENGLISH SPEECH].
The sheep held in the son's hands.
are also held in the Father's hand..
[NON-ENGLISH SPEECH].
This is a double assurance, and yet it is the same hand..
[NON-ENGLISH SPEECH].
And verse 29, it has a crucial phase,.
that this hand is greater than all..
[NON-ENGLISH SPEECH].
This leads us to the first misconception.
that we need to address..
[NON-ENGLISH SPEECH].

$^{161}$Many people live an exhausting life,.
even after they come to faith..
[NON-ENGLISH SPEECH].
Because they think that salvation.
is like a game of holding on to the Father's hand..
[NON-ENGLISH SPEECH].
The moment they loosen the grip, or grow weak, or stumble,.
they will fall and lose their salvation..
[NON-ENGLISH SPEECH].
So they live in constant fear that they are not.
holding tight enough..
[NON-ENGLISH SPEECH].
But if we look at our passage from verse 27 to 30,.
none of the verses says that the sheep.
must hold on to their shepherd..
[NON-ENGLISH SPEECH].
From the beginning to the end, it.
is always the shepherd's hand that holds the sheep..
It is the Father's hand that holds the sheep..
[NON-ENGLISH SPEECH].
Return to the earlier image that a child crossing the street.
is safe..
[NON-ENGLISH SPEECH].
Not because he grips his father's hand with all.
his strength..
[NON-ENGLISH SPEECH].
But because of his father's hands that wraps around.
the little one entirely..
[NON-ENGLISH SPEECH].
Even if the child lets go, the father does not..
[NON-ENGLISH SPEECH].
Your salvation does not depend on how tightly you hold on to..
[NON-ENGLISH SPEECH].
And that hand is greater than all..
[NON-ENGLISH SPEECH].
Finally, we come to the weightiest statement here,.
which brings up the third and final question..
[NON-ENGLISH SPEECH].
What does the word "one" mean when Jesus says,.
"I and the Father are one"?.

$^{201}$[NON-ENGLISH SPEECH].
The word "one" here in the original Greek.
is "neuter," not "masculine" or "not masculine.".
[NON-ENGLISH SPEECH].
This grammatical distinction is important.
because Jesus is not saying that the Father and the Son.
are the same person..
[NON-ENGLISH SPEECH].
What he's saying is that the Father and the Son.
are completely one in nature, in will, in action, and power..
[NON-ENGLISH SPEECH].
They are two distinct persons, yet one God, one will,.
and one hand..
[NON-ENGLISH SPEECH].
This addressed our second misconception..
[NON-ENGLISH SPEECH].
"I and the Father are one" does not.
mean that Jesus and the Father are indistinguishable..
[NON-ENGLISH SPEECH].
And the Jewish crowd, they actually understood exactly.
what Jesus meant, which is why in the next verse,.
they immediately pick up the stones.
and trying to stone Jesus..
[NON-ENGLISH SPEECH].
And the reason is because Jesus claims that he, a mere man,.
claimed to be God..
[NON-ENGLISH SPEECH].
Jesus was declaring, I am one with the Father in essence,.
will, and work, and what I do is what the Father does..
[NON-ENGLISH SPEECH].
It is precisely because the Son and the Father.
are one that Jesus would later say,.
"Whoever has seen the Father has seen me..
Whoever has seen me has seen the Father.".
[NON-ENGLISH SPEECH].
And this is the most important statement for Father's Day.
today..
[NON-ENGLISH SPEECH].
If we want to know what kind of Father our Heavenly Father is--.
[NON-ENGLISH SPEECH].

$^{241}$We don't have to guess..
We simply look at Jesus..
[NON-ENGLISH SPEECH].
Look at how Jesus treats his sheep,.
knowing them, giving them eternal life,.
and never letting go..
[NON-ENGLISH SPEECH].
And that is what our Heavenly Father is like..
[NON-ENGLISH SPEECH].
In fact, the setting of these words being spoken.
on the festival of dedication make it even more wonderful..
[NON-ENGLISH SPEECH].
This festival commemorates the rededication,.
making holy again what has been defiled..
[NON-ENGLISH SPEECH].
And in verse 36, Jesus described himself.
as the one whom the Father consecrated and sent.
into the world..
[NON-ENGLISH SPEECH].
Jesus reveals himself as the truly consecrated one..
[NON-ENGLISH SPEECH].
He came to bring back the loss and the woundership.
into the Father's hand..
[NON-ENGLISH SPEECH].
So in response to this Heavenly Father who is greater than all.
and never lets go, how might we respond to this Father's day?.
[NON-ENGLISH SPEECH].
First, if the word "Father" brings you grief,.
let the Heavenly Father redefine what "Fathers" mean..
[NON-ENGLISH SPEECH].
We naturally project our image of an earthly father.
onto our Heavenly Father sometimes..
[NON-ENGLISH SPEECH].
If your father was absent, you may.
feel that God is also distant and would abandon you..
[NON-ENGLISH SPEECH].
And if your father was harsh or frightening,.
you may feel that God is always ready to punish you as well..
[NON-ENGLISH SPEECH].
But when Jesus says, "I and the Father are one,".

$^{281}$it means this..
To know the Heavenly Father, you don't.
have to guess through the lens of your earthly father, who.
disappoint you, but instead look through Jesus..
[NON-ENGLISH SPEECH].
However Jesus treats you is how the Heavenly Father treats you..
[NON-ENGLISH SPEECH].
This week, find some quiet time to bring a disappointment.
and your wounds from your earthly father.
before the Heavenly Father..
[NON-ENGLISH SPEECH].
And you can ask him to heal your understanding of the word "Father".
through the image of Jesus..
[NON-ENGLISH SPEECH].
And this is Ken's testimony..
So my father was an absent father..
From a very young age, he was not.
there for many of the important moments in my life..
[NON-ENGLISH SPEECH].
Whether it was the secondary school graduation, university,.
or even post-graduation ceremony, he never showed up..
[NON-ENGLISH SPEECH].
Every time, I desperately hoped that he would appear and say,.
"You did well..
I'm so proud of you.".
[NON-ENGLISH SPEECH].
When I came to faith, I thought that this pain.
would just naturally go away..
[NON-ENGLISH SPEECH].
But then I discovered that this sense of not being a firm.
can be projected onto the Heavenly Father..
[NON-ENGLISH SPEECH].
Deep down in me, I felt that I was never doing enough,.
that I have to keep striving to earn the Heavenly Father's.
approval..
[NON-ENGLISH SPEECH].
It was only until later when the Heavenly Father healed me.
and helped me to redefine things through Jesus..
[NON-ENGLISH SPEECH].
And he showed me how deeply I'm loved..

$^{321}$And I slowly move out from under the shadow.
of not being reassured..
[NON-ENGLISH SPEECH].
So I would like to encourage everyone to do the same,.
to bring your disappointments and wounds.
about your earthly father honestly.
before your Heavenly Father..
[NON-ENGLISH SPEECH].
And ask him to heal your understanding of the word.
"father" through the image of Jesus..
[NON-ENGLISH SPEECH].
Second, if you're longing for a hand that you can no longer.
hold, his father say--.
[NON-ENGLISH SPEECH].
Remember that your father is in our Heavenly Father's hand..
[NON-ENGLISH SPEECH].
The scripture says no one is able to snatch them.
from the father's hand..
[NON-ENGLISH SPEECH].
If the father that you're missing belongs to the Lord--.
[NON-ENGLISH SPEECH].
Then this is the promise for you..
Not even death can snatch him from the Heavenly Father's hand..
[NON-ENGLISH SPEECH].
Third, if you are a father today,.
ask the Lord to help us and learn to be like him..
[NON-ENGLISH SPEECH].
None of us as fathers are perfect..
We fail, we lose our temper, our absence,.
and we disappoint our children..
[NON-ENGLISH SPEECH].
But our Heavenly Father has given us.
a model to follow, to truly know your children,.
not just the greats, but to truly know them, to hold them,.
and to not easily let go..
[NON-ENGLISH SPEECH].
You don't need to be a perfect father, but through the Lord,.
you can be a father who never lets go..
[NON-ENGLISH SPEECH].
The kind of father who, even when your children stumble,.

$^{361}$or grow weak, or disappoint you, that you still.
don't release that hand..
[NON-ENGLISH SPEECH].
Because this is exactly what the Heavenly Father has always.
done for us..
[NON-ENGLISH SPEECH].
So let's return to the questions that we have at the beginning.
and answer them together..
[NON-ENGLISH SPEECH].
So first question, whose hand does the father hold?.
He holds every sheep that hears his voice and follows him,.
not because the sheep is good enough or hold tightly enough,.
but because the shepherd knows and has chosen that sheep..
[NON-ENGLISH SPEECH].
Second, what is the basis of our security?.
It is not on how tightly we hold on to,.
but in a hand that is greater than all,.
that no one can snatch us away from..
[NON-ENGLISH SPEECH].
And third, what does I and the father are one.
have to do with Father's Day?.
It is because the son and the father are one.
that we see the Heavenly Father's heart clearly.
through Jesus, a father who never lets go..
[NON-ENGLISH SPEECH].
So some questions for us to reflect on..
When you hear the word "father," what.
is the first thing that comes to your heart?.
Has that feeling quietly shaped how you.
understand our Heavenly Father?.
[NON-ENGLISH SPEECH].
And on this Father's Day, are you.
willing to let the Heavenly Father hold you again.
through the image of Jesus?.
[NON-ENGLISH SPEECH].
Let's pray together..
[NON-ENGLISH SPEECH].
Dear Heavenly Father, thank you for revealing to us.
what kind of father you are through your son, Jesus..
[NON-ENGLISH SPEECH].

$^{401}$We thank you that you are not a father who is far off or cold.
or always ready to punish us..
[NON-ENGLISH SPEECH].
But you are a father who knows us, who comes to find us,.
who holds us with a hand greater than all..
[NON-ENGLISH SPEECH].
May you heal every wound from our earthly father.
that we carry..
And may those who have been hurt by their fathers.
or those who are missing their fathers today.
be held anew by your hand..
[NON-ENGLISH SPEECH].
And may you help every father among us to learn from your.
hands the one that never lets go..
[NON-ENGLISH SPEECH].
We give this Father's Day into your hands..
In Jesus' name, we pray..
Amen..
[NON-ENGLISH SPEECH].
We wish you all a wonderful Father's Day..
\newpage



\section{利未記 23:39-43}
\label{sec:xpOSI0B_Qx0}
\textbf{住棚 (利未記 23\_39-43) - 陳崇基牧師}
\newline
\newline
連結: \href{https://youtube.com/watch?v=xpOSI0B-Qx0}{\texttt{https://youtube.com/watch?v=xpOSI0B-Qx0}} ~~~~ 語音日期: 2026-07-12
\newline
\newline
\hyperref[sec:qC9_GTWJyK4]{\small{< < < PREV SERMON < < <}}
~
\hyperref[sec:index_chronic]{\small{[返順時目]}}
~
\hyperref[sec:index_scriptual]{\small{[返順卷目]}}
~
\hyperref[sec:IBF8VFRUG4s]{\small{> > > NEXT SERMON > > >}}
\newline
\newline
利未記 23:39-43
\newline
\begin{longtable}{cl}
\hline
\hline
章節 & 經文 (和合本修訂版)\\
\hline
23:39 & \begin{tabularx}{0.7\textwidth}{X} 「但是，從七月十五日起，你們收藏了地的出產之後，要守耶和華的節期七日。第一日為要完全安息，第八日也要完全安息。 \end{tabularx} \\ \\ \relax
23:40 & \begin{tabularx}{0.7\textwidth}{X} 第一日，你們要拿美好樹上的果子、棕樹枝、樹葉茂密的枝條和河邊的柳枝，在耶和華－你們的神面前歡樂七日。 \end{tabularx} \\ \\ \relax
23:41 & \begin{tabularx}{0.7\textwidth}{X} 每年你們要向耶和華守這節七日。你們在七月裡所守的節，要成為世世代代永遠的定例。 \end{tabularx} \\ \\ \relax
23:42 & \begin{tabularx}{0.7\textwidth}{X} 你們要住在棚裡七日；凡以色列家出生的人都要住在棚裡， \end{tabularx} \\ \\ \relax
23:43 & \begin{tabularx}{0.7\textwidth}{X} 好叫你們世世代代知道，我領以色列人出埃及地的時候，曾使他們住在棚裡。我是耶和華－你們的神。」 \end{tabularx} \\ \\
[1ex]
\hline
\hline
\end{longtable}
$^{1}$多謝主恩給我們的報告.
讓我們清醒一下.
剛才的歌很能夠幫助我們進入敬拜.
讓我們能夠有一個心來聽道.
我們一直在說這個系列.
這次我們應該是第五講.
我們講鑄棚這件事.
鑄棚是什麼一回事呢.
我們先看聖經.
然後大家可以思想一下這段聖經對我們今天有什麼意思.
講到最後的時候.
我會告訴大家有些講堂上的更新.
利未記23章說什麼呢.
他說你們對以色列人說.
神的子民.
他說你收藏了地的出產.
從7月15日起.
要守耶路華的節有七天.
第一天你要休息.
第八天也要休息.
中間做什麼呢.
第一天首先你要拿美好樹上的枝子.
果子.
棕樹枝.
麥貌樹的枝子.
還有河邊的楊柳等等.
就是拿一些草.
這些樹枝.
在耶路華的神面前.
來去歡喜七天.
你要守節是七天的.
作為耶路華的節.
這個就是你們世世代代永遠的定例.
要在這七月.
守這個節日.
你們要住在棚裡面七天.
以色列人中凡本地的人都要住在棚內.
好讓你們世世代代知道.
我領以色列人出埃及的時候.
曾使他們住在棚內.

$^{41}$我是耶路華的神.
我們說這個大的系列.
從創世記開始直到聖經.
最後我們說.
聖經好像一個長篇的劇集一樣.
有很多集的.
你當它每一卷書一集就有66集.
我們就是說整個的聖經.
就是告訴我們這個劇集裡面最主要的主角是誰.
是上帝自己.
上帝在人間在人類在自己的子民裡面做了很多事情.
然後他藉著這些故事這些事蹟.
來告訴我們啟示.
介紹這位上帝是一個怎樣的上帝.
縱使這個世界是他所創造的.
是他安排得井井有條的.
但是因為人我們叛逆.
以至搞亂了這個世界.
但是上帝從來沒有放棄過這個世界.
一直到推進他自己救贖的工作.
我們知道在創世記裡面.
他就揀選了阿巴拉罕等等的族長.
由他們生出一個大國.
這個大國叫做舊約的以色列.
這個大國就領受了上帝的律法.
出了埃及.
然後成為一個新的子民.
然後他們能夠成為上帝的工具.
以至上帝能夠去祝福,傳遞所有裡面的人.
今天這段經文就是摩西五經的第三本書.
叫做《利未記》.
《利未記》那裡說了一些.
我們經常覺得是一種規條.
就是跟我們現代人沒什麼相關的.
因為裡面說到會幕裡面辦理的祭祀.
或者那些儀式,規矩等等.
祭物要怎樣處理等等.
然後又說了一大堆的規條.
是說到時間怎麼使用的.
譬如節日.

$^{81}$什麼時候會有不同的節日等等.
這些都是《利未記》我們的感覺.
但是《利未記》其實是寫給每一個神的子民的.
因為《利未記》第一章第一節說到.
上帝要曉如告訴摩西.
你要對以色列人說.
不只是以色列國裡面.
以色列人民裡面的祭司.
那些所謂神職人員.
去到說這些所謂祭祀裡面的手冊.
或者規條等等.
原來每一個以色列人.
都需要明白這些東西.
就是說如果你拿著這個《利未記》.
你看這一集的時候.
看一集的時候.
這個劇集.
或者這個書卷的時候.
你就會覺得.
我都好像大祭司一樣.
我要怎麼處理我自己.
怎麼去確保我自己.
能夠去到上帝的面前.
能夠讓上帝住在我們中間.
所以在這裡說了一些規條.
其中今天的經文就說到.
一些守節的規條.
當然有大大小小的節日.
如果你認識一些猶太人的朋友.
或者同事的話.
他們經常放假.
久不久就會放假.
他們每個月都有假放.
不像我們美國人.
我們美國人最少假放.
他們經常放假.
每個月都有假放.
這裡其中說到一個月.
就是他們的農曆.
7月左右的時候.

$^{121}$7月15號.
就是今天我們的西曆.
大概9月尾10月的時候.
就說到一些節期.
在摩西五經裡面.
告訴以色列人聽神的子民.
他們時間是由上帝的原則來定的.
不是說你喜歡怎樣就怎樣.
時間是有一種規律的.
而且那種規律是屬靈的.
Spiritual的規律.
就是上帝告訴大家.
你每逢七天裡面.
有一天是需要安息去敬拜我的.
你不能夠七天都做事.
六天你要努力去做事.
但是你第七天的時候.
就是安息日.
就是所謂今天新約裡面.
基督已經承傳所有的東西.
所以很多時候.
是第七天的第一天.
我們去敬拜.
所以我們信了主的時候.
我們的時間就會改變.
我們時間的習慣.
我們時間的運作.
我們都會有所改變.
你有沒有想過.
我們的信仰和我們的schedule是有關的.
和我們的calendar是有關的.
是有關的.
在舊約裡面.
神的子民領受這些律法的時候.
就告訴我們.
大大小小的節日裡面.
有三個最大的.
第一個就是預節.
就是西曆大概四月初.
我們復節的時候.

$^{161}$第二個就是七七四十九日之後.
第五十日.
五旬節.
Pentecost.
就他們收了他們的大麥小麥.
有收過節的時候.
那個大概是西曆的五月.
然後駐彭節.
就是今天我們所說的那段經文.
到他們大概今時今日的西曆.
九月尾.
十月的時候.
他們就需要去做駐彭.
來到所謂的慶節.
首節.
說到這些節日.
我們每一個人都會想起的.
每逢我們有節日的時候.
我們就會想起那些典故.
或者那些由來.
最近是什麼節日.
最後端午.
是不是端午.
端午我沒有搞錯.
因為我收了很多愛心的粽.
解決了很多餐.
我們會想起那些源頭.
為什麼有這樣的故事等等.
我們基督教最重要的節日是哪幾個.
復活節和聖誕節.
可以這麼說.
比較西方的西化的基督教.
就是有這些重要節日.
在九月裡面就是這三個最重要的節日.
今天這段經文就說到駐彭節的吩咐.
有幾件事值得留意.
第一件事.
就是這些駐彭節.
或者這些節日.
都是叫我們停下來.

$^{201}$來紀念一下.
想清楚上帝在我們心靈裡面有的恩典.
這裡很搞笑.
他說7月15日就開始首節.
什麼時候.
首七天.
他說第一天你要怎麼樣.
然後第八天你要怎麼樣.
但有趣的地方就是.
他說你要在出產之後.
你收藏了出產.
到7月15日起.
就要首節了.
什麼意思.
就是你出了糧之後.
出了糧.
你收藏了你的土產.
這個是秋季的時候.
你有秋收.
然後你就要停下來.
去到首節.
這個很過癮.
因為農作物入倉.
你收到這些大麥小麥.
或者這些農作物的時候.
最安逸的時候就到了.
我問你.
你出糧那天你去哪.
出糧那天你去哪.
我還記得那時候.
我在香港中學畢業.
然後我找了一份工.
那份工當時是OK的.
那個工資.
就是你肯做就行.
在銀行做的.
我記得我第一張支票.
出糧的支票.
出了的時候.
我去了哪裡.

$^{241}$很有趣.
那時候銀行是出現現金的.
那時候我是拿一倍.
拿一袋.
出現現金.
整疊現金在手上.
我去了哪裡.
我去了酒吧.
我和朋友一起去酒吧.
我請大家吃一餐.
就是你出了糧.
當然要去開心.
去安逸的時候.
最開心的時候.
去慶典的時候.
吃大餐的時候.
應該是這樣.
但這裡很奇怪.
他說你豐收之後.
你就要做什麼.
要搭棚.
然後在那裡住七天.
這個不只是以色列人的問題.
其實是說我們人性的問題.
我們最好景的時候.
我們最容易忘記上帝.
我們的心很容易傾向於物質的東西.
朱鵬哲偏偏安排在豐收之後.
神好像告訴我們.
你最危險的時候是什麼時候.
你和上帝之間的關係.
如果你說你相信耶穌的話.
你一個基督徒的生命的時候.
你最危險的那一刻.
不是你缺的時候.
是你有的時候.
是你豐收的時候.
就是當你拿到上帝的祝福.
手裡面有金有銀的時候.
那時候才是最大的危機.

$^{281}$《申命記》第八章提到.
以色列人進入迦南的時候.
進入豐收之地.
進入安逸之地的時候.
他們有一個警告.
這個警告是什麼.
以色列人走過曠野四十年.
《申命記》五經第五本.
告訴他們.
你進入迦南地有豐收有安逸的時候.
要記得什麼.
你要謹守耶和華律法.
以致你不會忘記耶和華.
因為什麼.
因為在說人性的一件事.
就是你建造美好的房屋的時候.
有牛有羊的時候.
金銀增添的時候.
你心裡面就會驕傲.
驕傲不一定是被人欺負.
或者很威風的態度.
驕傲可以藏在心裡面.
驕傲就是我們有一種自足.
當我們有自足有自滿的時候.
我們就容易忘記耶和華.
忘記我們的神.
所以有人說得很對.
他說錢可以跟我們說很多東西.
並且每句話幾乎是上帝可以跟我們說.
因為如果有錢的話.
我們很富足的時候.
錢就好像上帝一樣可以跟我們說.
你很安全.
是不是.
當你查一下你的銀行帳戶.
或者你的錢包的時候.
錢好像在跟你說話.
不用怕.
你有保障.
你有和平.

$^{321}$你有安全.
你有平安.
等等等等.
這些不一定是很錯的事.
錢可以跟我們說很多東西.
是上帝可以跟我們說的.
對不對.
是取代上帝所說的話.
你未來一定是無憂的等等.
但聖經裡面反反覆覆告訴我們.
只有上帝才會賜給我們真正的平安.
所以律法書裡面有一個智慧.
他將這個鑄盤子陷入一個時間.
就是你最豐收的時候.
拿走所有東西的時候.
人最富足的時候.
你就需要停下來.
你要想一下.
上帝怎樣恩待過你.
如果你忘記的話.
在忙碌裡面.
或者你在善忘的當中.
你很快上帝就會淡忘了.
我們教會去到第十年.
每逢教會去到唐慶.
或者新的一年的時候.
或者新的發展的時候.
很多時候我們就會很自然.
我們會數一數我們的成績等等.
我們做過多少事.
當然是有需要的.
有需要.
我最近在反反覆覆想.
我們教會從開始的時候.
我們是從零開始的.
從零開始.
我們一開始是沒有糧出的.
很緊的.
如果不是上帝的恩典.
真的不會成為一個教會.

$^{361}$如果不是因為大英姐妹.
如果不是因為大家在禱告.
在信心裡面擺上.
是根本不會走到今天的.
因為大部分.
我說北美.
大部分新開的教會.
是不過三年的.
三年就死了.
要不就是童工內部的領袖.
他們有瓦解.
或者談不攏.
或者有不同的外在.
或者內在的因素等等.
大部分都還在.
大部分還在.
普中心留在這裡.
那些是中流砥柱.
但我想說一件事.
我們在堂詠的時候.
我們可以數一下業績.
可以數一下成功的地方.
但一件事千萬不要忘記.
想起零開始的時候.
原來是上帝的恩典.
這是鑄棒節的智慧.
豐收之後.
你就需要停下來.
因為我們要學習一下.
去紀念一下上帝的恩典.
上帝在生命裡面有沒有祝福.
有沒有恩典.
你當然會說有.
但重點就是.
你有沒有停下來去數一數.
停下來去紀念一下.
寫下也好.
自己去銘記也好.
以致你和上帝的關係.
不會因為忙碌,豐收而忘記了.

$^{401}$第二方面.
鑄棒節的智慧是什麼.
我要面對人生的軟弱和脆弱.
這裡說竹棒.
第四十節 他說第一天.
你要拿樹上的果子,樹枝等等.
在耶和華神的面前.
你要歡喜快樂七天.
你要住在棚裡面四十二節.
多少天?.
七天 有一個星期.
棚子.
希伯來文就是Sukkah.
或者縱數是Sukkot.
如果你認識猶太人的朋友.
他們首節的時候.
或者看一下猶太人的閱歷.
你要知道有Sukkot.
Sukkot就是鑄棒節的時候.
就是今天西曆大概九月尾的時候.
棚子算是什麼?.
這裡很有趣.
後期的拉比去解釋這些律法的時候.
他們認為做一個Sukkah.
應該怎麼做呢?.
拉比文學裡面有一堆東西叫做.
Halaqaq.
Halaqaq是什麼意思呢?.
Halaqaq就是Halaq這個字.
Halaq就是走出來的意思.
就是實踐出來的意思.
所以有一大批的律法的解釋條文.
歸納了成為所謂的Halaq.
Halaqa.
裡面解釋一個棚子應該怎麼做.
聽著有幾樣東西值得留意.
原來每逢鑄棒節來到的時候.
你就要預備材料.
然後你就在你的後院.
或者你有的地方.

$^{441}$你的所謂patio.
你就會做一個棚子.
這個棚的高度是不可以高過20爪的.
不可以高過.
大概今天我們計算是33尺.
因為拉比認為.
如果你做得太高的話.
太過宏偉的話.
就會變成永久.
你就忘記了你住在那裡.
你就覺得我永久住在這個地方.
太舒服的時候.
你就忘記了.
你所住的只不過是一個很脆弱的棚子.
所以一定不可以太穩陣.
第二樣東西就是要用天然的材料.
樹枝等等那些.
枯枝會損壞的.
會枯萎的.
而且是不太防水的.
不要用一些很好的材料.
就是你不要空一些棚子.
你找外面隨便撿一些東西就好了.
特別是打完風.
我們大部分在外面.
斬一點下來.
然後你就搭一個棚子.
並且一個棚子是需要通天的.
棚子的帳幕周圍就是要簡陋的.
因為你住七天而已.
你住七天.
你就不會覺得真的住一個家.
你不會真的搬去那個地方.
只不過讓你去明白一件事.
就是一個屬靈上的反思.
你住在棚子七天.
你就可以明白到.
人生真的很脆弱.
棚子一樣.
我們現代人生活.

$^{481}$我們花很多錢.
我們去買保障.
買一些安全感回來.
買一些穩定回來.
住棚子就是和這個概念完全相反.
你住七天一個會塌的棚子.
會塌的棚子.
並且規定你整一個會塌的棚子.
不可以整得太實際.
因為人生就是脆弱的.
人生是脆弱的.
上個星期.
上兩個星期.
我們打大風.
真的很可怕.
你們有沒有停電.
有誰停電.
我停了二十多電.
通常天氣就是這樣.
你最需要冷氣的時候.
它就停電.
你最需要那些東西的時候.
它就剛好會壞.
接下來這幾天又會有熱浪.
我告訴你.
你的冷氣就會壞.
很有趣.
停了二十多電.
有美國以外的朋友和我有些通訊.
他說不是嗎.
美國大國居然停電.
美國大國就會停電.
因為我們給很多錢給軍人打仗.
軍備我們很厲害.
但是電力我們就不夠.
甚至NBC Chicago有一個報導.
問一個問題.
他說龍捲風的通道.
龍捲風的大道.
因為有些地方.

$^{521}$特別是Kansas下面的州.
就經常會有龍捲風.
這個所謂的通道.
Tornado的通道.
是不是會轉移.
因為氣候變化的緣故.
想都沒有想過.
我們也有份.
不要說別人會打龍捲風.
我們不會的.
可能我們也會.
鑄鵬的鵬子是什麼.
鑄鵬Tabernacle或Squat.
就是說我們是脆弱的.
我們是受不了的.
天氣氣候所變化的.
你怎麼可能控制.
人生就是這樣.
人生就是有脆弱.
我們需要去承認.
真正保護你的.
真的不是鵬.
真正保護你的.
是在曠野和你同行的上帝.
祂的玩目是住在你們中間.
祂去給你最終的保障.
祂給你最終的Security.
是住在你們中間的上帝.
不是你所搭的棚.
生命裡面我們可以搭多少棚都好.
都是會脆弱.
大風雨水臨到的時候.
我們就會倒下.
最重要的是什麼.
我們不是說不要見棚.
我們要做好的管家.
但我們要明白我們本質是脆弱的.
教會也是一樣.
教會也可以是脆弱的.
因為我們都是不堪一擊.

$^{561}$我們都是有缺陷.
不是大英姐妹.
我們夾手夾腳的話.
我們早已經玩完了.
所以在《申命記》說什麼呢.
《申命記》29章這裡沒有打出來.
他說回顧40年裡面.
神怎麼帶領以色列人的時候.
祂說什麼呢.
祂說我領你們走過曠野40年.
你們身上的衣衫沒有穿破.
你們的鞋子也沒有穿壞.
鞋子沒有爛到.
衣服沒有穿破.
當然祂不是說神蹟.
祂說的是你所基本要有的東西.
所有必須的東西.
這40年裡面.
我都供應了給你們.
在40年的曠野生命裡面.
不是因為你碰.
其實你很脆弱.
但是上帝一路去保守住.
去帶領著你.
除了去到面對我們人生的脆弱之外.
第三件事就是去到學會.
怎樣去承傳上帝的使命.
怎樣去傳承祂.
這一段鑄盤節的吩咐裡面.
其中一段最重要的說話.
就是好讓你們世世代代都知道.
我曾經領以色列人出埃及的時候.
曾經使他們在曠野裡面.
住在棚裡面.
我是耶和華.
你蒙的神.
這個可以算是最重要的一句宣告.
是一個神學的宣告.
因為說一個鑄盤的日子.
這個節日.

$^{601}$並非是一個文化.
純粹文化的節日.
是一個活生生的.
一個重演的.
一個神學的宣告.
去告訴人聽.
屬靈上面.
哪一個是我們的上帝.
即是猶太人的家庭.
他們到鑄盤節的時候.
他們會出去很脆弱的棚.
他們會住七天.
或者在那裡吃七晚晚餐.
他們有一些猶太人的.
叫做律法的移民.
那些的儀式.
即是說父母會跟小孩子.
下一代.
來到去有問有答的.
好像那些問答的冊子.
你打開那些冊子.
你就會問.
小孩子會問一個問題.
然後父母就會回答.
其中一些問題.
問什麼呢.
小孩子會問.
為什麼我們今天會搬出來.
為什麼我們這七天.
要住在這個地方.
為什麼我們放棄.
我們有磚有瓦的.
很舒服的房子.
然後我們搬出來這裡.
在這裡紮營.
父母的回答是什麼呢.
因為我們曾經有家.
因為我們曾經沒有家.
但上帝成為我們的家.
我們沒有家.

$^{641}$但上帝是我們的家.
為什麼我們會出來這裡住.
因為上帝是我們的家.
第二個問題就是.
為什麼屋頂會有洞呢.
那些糖果和棚子.
為什麼會有洞會漏水呢.
父母就回答.
不是因為是壞的.
是故意的.
故意弄到這樣是通風的.
因為真正遮蓋我們的不是屋頂.
是上帝.
屋頂不是我們的shelter.
God is our shelter.
上帝才是我們的shelter.
在住棚節裡面.
有很多這些神學的反省.
雖然我們未必會走去搭棚.
去守住棚節.
但是我們需要明白.
聖經裡面所說的.
告訴我們.
這些屬靈的智慧.
並且後期住棚節.
猶太人的住棚節.
他們會在猶太人的會堂裡面.
去供讀一些聖經的書卷.
聖經裡面舊約.
希伯來的聖經.
有39卷.
你覺得在住棚節裡面.
猶太人的會堂.
很多時候供讀的是什麼書卷.
哪一卷.
創出你紋身書士.
德撒王代代拉尼斯等等.
39卷.
哪一卷.
是傳道書.

$^{681}$傳道書.
為什麼.
因為傳道書說.
我們人生是虛幻.
好像一口氣.
不說太多.
因為Joyce已經在說傳道書的系列.
那裡說什麼呢.
傳道書最強調的是.
既然人生在日光之下.
是脆弱的.
你自己應該怎樣去調教.
你人生的方向.
最值得你追求的東西是什麼.
傳道書最後的結論就是說.
你追求上帝.
以上帝的事為念.
以致你人生的方向.
是不會走錯的.
你可以忙於你每一天.
由早到晚.
24/7.
裡面所要經營的東西.
但是你不要搞錯.
你最終的方向.
是應該去到敬畏上帝.
曾經有人問我.
就是今年.
今年有人問我.
他說如果時光倒流的話.
回到十年前.
當初開教會的時候.
這麼艱難.
你也知道我開教會前.
我開過Uber.
有時候.
有時候人們會跟我開玩笑.
他說挺好玩的.
我說好玩.
當你兩三點要去酒吧載人的時候.

$^{721}$當你戶口的時候.
你覺得好玩嗎.
但是有人問我.
如果你時光倒流十年前.
你會不會再來過.
我幾乎不用思考.
我也會回答.
我會.
我會再做.
你知道我們人生裡面.
特別是基督徒的人生.
有多好.
我們能夠去建立教會.
有多好.
我們有這樣的福分.
有這樣的機會.
我們能夠去參與建立教會的工作.
這個不是要自豪.
不是要去說.
我做了很偉大的人.
不是.
而是我們有這樣的福分.
有這樣的機會.
有這樣的恩典.
我們能夠參與上帝的工作.
參與上帝的工作.
最近大家經常留意的是什麼.
世界盃.
如果你可以和美斯踢一場球.
你願意嗎.
即使你沒有機會碰到球.
起碼你可以和他同場.
和他同場和他拍照.
給他衣服和簽名.
這是什麼.
不是因為你厲害.
球你沒有碰到.
是因為你和他同場.
你可以和他參與那場賽事.
只是這樣.

$^{761}$出場三分鐘都可以.
就算出場三分鐘被人踢到腳.
踢跛了進醫院.
你也願意.
不願意嗎.
你以為每個人都可以嗎.
你願意嗎.
就是因為你可以參與.
就有這樣的機會可以參與.
因為人生就像三分鐘一樣.
如果我可以參與和球王同場的話.
什麼榮幸.
上帝給我們的生命都是一樣.
能夠在這裡建立上帝的家.
能夠在我三分鐘的人生裡面.
能夠去參與上帝的遊戲.
世界賽.
多大的福分.
棚子很快就會離開.
我們應該怎樣去把握.
等一下我們按照這個系列.
我們會每一次給大家一些紀念品.
這個瓶子如果你是今天新來的朋友.
或者最近才來的.
你不知道在搞什麼.
你可以拿回去.
因為每一次我們系列裡面.
講到這個大故事.
上帝是一個願意祝福世界.
願意使用我們.
給我們機會來祝福世界的.
這個訊息.
我們就放在這個瓶子裡.
每一件紀念品我們放在這個瓶子裡.
放在你的桌上.
你看到我就放在我的電腦桌上.
讓我記得.
是多大的福分.
三分鐘可以參與這場賽事.
但是拿這件紀念品的時候.

$^{801}$請你記得.
想一下.
我最近怎樣經歷過上帝的恩典.
怎樣經歷過上帝的恩典.
這次拿這個紀念品是一支竹.
一支很短的竹.
代表搭棚.
豬棚.
這支竹提醒我們.
好像以色列人過去四十年.
他們走過曠野的時候.
上帝是一個很脆弱的棚.
仍然堅持住他們.
去保守住他們.
這是第一個問題.
第二個問題.
我們應該怎樣去回應.
上帝給我們的祝福和保守.
有沒有停下來想一想.
身體力行我們應該怎樣去回應.
第三我應該怎樣去傳承神的工作.
特別是下一代.
下一代.
他會知道神仍然在.
好像小朋友住在棚內.
他問我們.
為什麼我們會在這裡.
為什麼我們會來參與這樣的工作.
為什麼我們會睡晚一點.
為什麼我們要做這樣做那樣.
我們怎樣回答.
我們回答就是上帝仍然在.
上帝是很真實.
我自己經歷過.
我都希望你能夠經歷.
一些未信主的朋友.
一些我們身邊的鄰舍.
我們要告訴他們.
上帝是真實.
上帝是仍然在.

$^{841}$今天請你拿這個紀念品回去.
如果你沒有那個瓶子的話.
你整個瓶子可以拿走.
我今天預備了一些瓶子.
有這些紀念品在.
作為一個提醒.
在這裡紀念.
上主,我們謝謝你.
因為我們能夠跟你同場.
有這樣的福分.
有這樣的榮幸.
可以參與你自己的工作.
是一些有永恆價值的事物.
神,我們謝謝你.
在我們脆弱的人生.
三分鐘的人生裡面.
我們可以有機會.
領受你的祝福.
領受你的恩典.
能夠跟你一起來服侍.
一起來事奉.
建立你自己的神國的子民.
你自己的工作.
主,我們很軟弱.
我們很脆弱.
但我們求你來使用我們.
多謝你給我們這個權柄.
這個機會.
讓主你在我們生命當中.
在我們家庭裡面.
流露出上帝你自己將來的國度.
以這個國度給人的感覺.
是更加真實.
祈禱奉主名禱.
\newpage



\section{}
\label{sec:IBF8VFRUG4s}
\textbf{再思女性的事奉 - 陳崇基牧師}
\newline
\newline
連結: \href{https://youtube.com/watch?v=IBF8VFRUG4s}{\texttt{https://youtube.com/watch?v=IBF8VFRUG4s}} ~~~~ 語音日期: 2026-07-12
\newline
\newline
\hyperref[sec:xpOSI0B_Qx0]{\small{< < < PREV SERMON < < <}}
~
\hyperref[sec:index_chronic]{\small{[返順時目]}}
~
\hyperref[sec:index_scriptual]{\small{[返順卷目]}}
~
\hyperref[sec:MlRDdW70h9M]{\small{> > > NEXT SERMON > > >}}
\newline
\newline
$^{1}$是由Mark Sir跟我們說到的.
早晨 兄弟姊妹.
今天我們想說一個很切身的話題.
我們一向都在說NEXT今年.
但是我覺得有一些話題.
我們不妨在講台上大家思想一下.
特別是說到女性的事奉.
在教會裡面女性的事奉和角色是怎樣的.
還有她們的重要性在哪裡.
這個是很切身的問題.
因為最近我們的總會 美男浸水會.
在他們的大會裡面有一個議案通過了.
他們大比數通過這個議案.
就是讓女性在教會裡面的事奉.
將她限制於非成人.
即是小孩子等等.
美男浸水會不是靚仔的意思.
美男浸水會是美國男方浸水會.
美國男方浸水會是美國裡面最大的一個中派.
我們教會是屬於美男浸水會的.
不是因為我們靚仔.
而是因為我們本身開堂的時候.
是美男浸水會的網絡去幫助我們的.
他們的大會正如我所說.
就是四分三去通過.
女性在教會裡面的事奉是需要限制的.
即是說只是教導小孩子.
或者成人的教導和成人的講道.
女性是不可以這樣做的.
理據很簡單.
他們的理據就是聖經這樣說.
所以這是一個切身的問題.
無論你是今天第一次來教會.
或者你完全不認識聖經.
這也是一個很切身的問題.
為什麼呢.
因為怎樣去應用聖經.
怎樣去理解聖經.
很多時候都是我們今天的話題.
有些人會用聖經或者聖經的說話.

$^{41}$或者經文去做不同的事情.
可以去到起乾戈,打仗.
都可以用聖經來支持.
做很多事情都可以用聖經來做理據.
兩夫妻吵架都可以突然拿本聖經出來.
聖經叫你赦免我.
聖經叫你忍耐.
為什麼你不忍耐.
所以這個是和我們切身的.
就算你還沒接觸過聖經.
你第一次來教會.
你都想明白聖經這本書到底是一本什麼事.
因為在教會裡面.
特別基督徒的世界裡面.
聖經是絕對的權威.
所以我們很多時候會訴諸於聖經.
特別我們更正教.
所謂新教,基督教.
我們和羅馬天主教是有點不同的.
因為我們的權威性在哪裡.
是唯獨聖經.
所以永遠都是說回到聖經.
Return to the Bible.
所以我們需要去思考一下.
到底你聽到一些引用聖經.
去到有聖經支持的一些理論.
或者一些想法,概念.
到底是不是真的合乎聖經.
所以我們需要嘗試去思考一下這個問題.
因為很容易我們就會說.
聖經這麼說.
這個是聖經的權威.
所以就Case closed.
那就不需要再討論了.
我記得我預備的時候.
年輕的時候就想起一件事.
就是我大學的時候.
我很喜歡專門問難的.
聽完講道我就去挑戰牧師.
上完主一學我就去挑戰主一學的教師.

$^{81}$我說你這樣說好像不太對.
因為聖經不應該這樣解釋的.
我就是那些搞事分子.
從小到大都是這樣.
我記得有一次就講了《創世記》第一章.
講到六天創造.
Six days creation.
我就聽到一個教師這樣說.
這個六天創造就是一定是每一天.
因為什麼叫一天.
一天就是24小時.
對不對.
既然是24小時的話.
就應該六天創造.
就是六個24小時裡面創造所有東西.
一下子就走出來.
我就說這個好像不太對.
這樣去理解.
因為《創世記》到底是不是在講科學.
其實不是在講科學的.
所以他就去到反問我.
他說你居然去質疑我們2000年解經的傳統.
豈有此理.
你黃毛丫頭.
是他這樣的意思.
當時我年輕的時候就不太懂得去應對.
我無言以對.
後來我才慢慢成長的時候發現.
原來這個六天24小時.
原來沒有2000年的傳統.
他以為有而已.
其實很早期有一個人叫奧古斯丁.
他已經說這個六天創造.
其實不應該解作科學.
又或者字面的24小時.
只不過是象徵式一個框架.
來表述上帝的偉大.
上帝創造主的偉大等等.
所以我就再去思想一下.
很多時候我們接觸到一些教導.

$^{121}$雖然說有時候會訴諸於聖經權威.
或者我們會訴諸於歷史悠久.
但是你明白嗎.
就算訴諸於權威.
就算訴諸於歷史的傳統.
都不等於聖經本身的權威.
只不過你對聖經某一種的解讀.
那個權威而已.
就不是聖經的權威.
大家明白嗎.
聖經的權威當然是最高.
但是有不同人的解說.
你的解說不一定等同於.
聖經的絕對的權威.
OK吧.
對嗎.
你可以解錯的.
或者你可以梳理聖經裡面的材料的時候.
有一些甩漏的地方.
所以你要想到.
聖經當然是絕對的權威.
但是下面有很多不同的解說.
有一些好一點的解說.
有一些差一點的解說.
這個很多時候是我們有迷思來的.
我們很多時候說.
你去讀神學.
或者有弟兄姊妹去讀神學.
他說讀神學這些東西很難摸索的.
沒有對錯的.
你有你寫我有我寫.
我們經常都有這樣的迷思.
你上課譬如解提摩太前書.
你去寫一份論文出來.
你就會說寫論文而已.
我們大部分都是理科人.
很多理科人在這裡.
計數就算對和錯的.
一算就算到了.
小於一分也不對.

$^{161}$但是你說你讀文科的.
你有你寫他有他寫.
ABCD這樣每個人都可以對.
其實不是這樣的.
其實寫論文寫很多.
但是也有好的論文不是壞的論文.
對不對.
那怎樣才算是一個好的論文.
就是你梳理一些材料.
去處理一個論點.
處理一個思考的過程.
方法等等.
其實都是有好與壞的.
對嗎 大家明白嗎.
我們平日也會遇見的.
兩個人吵架的時候.
你也要梳理一下誰有道理.
你不可以說我有權威.
你一定要聽我說.
這些全部是溝通的問題.
所以我們讀聖經.
其實也是需要去思考一下.
所以今天我們選了這段經文.
是一段聖經裡面.
所謂最棘手的經文.
最棘手的經文.
雖然時間有限.
不過怎樣棘手也好.
我希望25分鐘之後告訴你.
其實不是那麼棘手.
好嗎.
2章第8至15節提摩太前書.
這裡是這樣說的.
這裡主題絕對是在說.
女人在教會裡面的事奉.
對不對.
所以很多時候就拿了這段經文.
來支持某一些的想法.
某一些的建議.
2章第8節.

$^{201}$保羅寫信給當時以弗所的教會.
當時就是提摩太的徒弟.
去到牧羊的宗教會.
這封書信到第2章的時候.
他就說.
我希望男人是舉起聖潔的手.
隨處禱告.
不發怒不爭論.
第9節.
我又希望女人以端正克制.
合乎體統的服裝.
打扮自己.
不以編髮金飾珍珠名貴衣裳來打扮.
要有善行.
這才與自稱為敬畏上帝的女人相稱.
女人要時時的馴服地安靜學習.
我不許女人教導或者講道.
也不許她管轄男人.
只要安靜.
因為亞當先被造.
然後才是夏娃.
亞當並沒有受騙.
是女人受騙.
陷在過犯類.
然而女人若持守信心愛心.
又在聖潔克制上面.
就必藉著生產而得救.
好棘手.
聽完之後都不知道講什麼.
我們就只知道.
我不許女人教導.
這個最清楚.
我們講了很多次.
《看聖經》一定是有它的背景的.
對不對.
《聖經》不是直接寫給我們的.
《聖經》卻是為我們而寫.
既然不是直接寫給我們.
就直接寫給當時他們當下的處境裡面.
去跟他們講說話.

$^{241}$這樣講說話就是直接寫給他們.
就是說2000年之前.
在第一世紀.
在以弗所的教會.
在大摩太的模樣之下.
他們有的處境去寫這段說話.
對不對.
而我們的處境往往是不同的.
現在我們21世紀的處境.
往往是完全不同的.
但是有些是為我們而寫的.
先搞清楚.
寫給他們到底是發生什麼事.
當時的教會.
提摩太模樣教會在以弗所.
以弗所是一個很著名的.
今天的土耳其.
西面一個很大的城市.
是一個當時土耳其省份.
就是小亞細亞省份的首都.
所以是一個很繁華的地方.
那裡他們宗教的背景是非常濃厚的.
如果你去以弗所的話.
今天在土耳其西面.
這個古城裡面都找到一些遺跡.
是跟當時第一世紀保羅.
去到以弗所的時候所見到的東西.
有些是他見到的.
有些是後期.
有些是早期一點.
當時以弗所已經是非常大了.
所以以弗所有一個廟.
是一個當時地中海東邊世界裡面.
一個最大的廟.
這個廟叫做亞底米女神的廟.
右下角那個只不過是重建而已.
因為左上角那個.
就是今天你去到的時候.
你可以在那裡拍照.
拍一張照.

$^{281}$你看到有多少東西剩下在那裡.
剩下一根柱子而已.
這個所謂七大世界人工的奇蹟之一.
Seven wonders of ancient world.
譬如埃及的金字塔是其中一個.
七大奇蹟之一就是亞底米女神的廟.
剩下裡面的遺跡.
當然剩下很少.
剩下根柱.
這個女神廟叫做亞底米女神廟.
裡面有很多祭司.
他們每年會有一個.
不單只是香火鼎盛.
每年都會有一個巡遊的慶典.
那些貴族的女士.
那些婦女就會盛裝出場來打扮.
然後模仿那個女神的模樣.
所以最終的目的.
其實在當時的社會宗教的氣氛下.
當你這樣做的時候.
你不單只是貴族.
你貴族打扮了之後.
以金髮首飾等等.
你再走出來的時候.
其實你就是鞏固你自己的地位.
因為很多時候平民百姓就是平民.
就是貧窮的人.
你這樣去做的時候.
你就會得到眾民的尊重和愛戴.
其實你就是鞏固.
建立你的社會經濟地位.
所以這些慶典不單只是宗教慶典.
是一些社會經濟上的慶典.
當時有一本小說.
那本小說叫做《伊斯維耶克》.
伊斯維耶克就是一位小說家.
叫做斯諾芬.
他所寫的.
這個斯諾芬是以弗所的斯諾芬.
不是雅典的斯諾芬.

$^{321}$雅典的斯諾芬比較有名.
以弗所有一個叫做斯諾芬同名的.
他寫了一本小說.
這個小說其實是愛情小說.
大家有沒有看過愛情小說.
逆書這些 是不是逆書.
對 《勁要》.
對不起 我不看這些.
但是斯諾芬寫了一本《伊斯維耶克》.
就是以弗所的故事.
以弗所的愛情故事.
就說到男女很深愛大家.
然後就克服很多困難.
最後終於可以一起了.
既然一起的時候.
最後他們就在亞底米女神廟.
他們就去參與巡遊.
說到女主角的打扮的時候.
就用那些字.
有些字就跟提摩泰前書第二章那裡的字.
差不多一樣的.
有髮飾 有穿金戴銀等等.
他就去參與亞底米女神像的巡遊的慶典.
甚至到了一個地步.
那些眾民會看到這個女主角.
因為打扮成這樣的時候.
他就以為這個女神是再生的.
他自己就好像女神一樣.
你將這些當時的文本.
來跟聖經當時的背景來做一些比較的時候.
你就知道有一點點的背景.
可能是我們今天未必想像得到的.
我們說回以弗所.
以弗所在那裡宣教有好幾年之後.
當然是建立了一些教會.
而那些教會當時是家庭教會.
在不同的家庭裡面去聚會.
不像我們這樣一排.
有很多人一起在一個公共的場合一起聚會.
以往那個教會的聚會形式是家庭教會.

$^{361}$既然是家庭教會的話.
就一定是一些比較有錢的人.
才能夠開放自己的家庭.
譬如昨天我們教會停電.
這一區停電.
如果有聚會的時候就去誰那裡.
當然是去比較有錢的人家裡.
或者塞得他一進去.
如果去到社下我就不行了.
只可以有兩個聚會.
兩個小組而已.
但是如果有幾十人.
那當然是去一些比較有錢的人家把他塞進去.
然後他就可以煮東西吃.
然後去款待我們.
當時以弗所確實有一些比較富戶的人.
就是信主的.
而且很多時候都是一些婦女信主.
你要看《少林傳》.
你都差不多拿到這個味道.
有時候一些貴族的婦女.
她信了主.
因為她是富裕的家庭出身.
所以很多時候教會在她家裡聚會.
她就成為教會裡面的主人.
她開放教會的時候.
她自然因為有錢有地方有地位.
她就成為當時羅馬社會的所謂恩惠者.
Benefica.
她就可以給恩惠給其他人.
其他人來到的時候.
就可以聚會了.
很多時候因為這樣的緣故.
當時羅馬的社會都是這樣.
有錢有地位的人.
他們當Benefica的時候.
他們就可以打扮自己.
或者去派發利是.
做一些事情.
然後就會令到那些的貧民.

$^{401}$去到給他有尊重.
因為這樣的舉止.
其實就是一種社會經濟行動.
就是說我這樣做的時候.
你受了我恩惠.
然後你就需要回饋.
怎樣回饋.
起碼你會投我票.
你會支持我.
就是這個意思.
和今天的政客差不多.
是去建立自己的.
鞏固自己的尊嚴.
如果你這樣去想的時候.
你就知道教會的情況.
不像我們今天這樣.
今天我們很少特意穿金戴銀來教堂.
來教會.
然後我就鞏固我自己的地位.
對嗎.
我不介意你穿金戴銀.
我不介意的.
我都不會覺得你鞏固你自己的地位.
因為今天有點不同.
就算穿金戴銀進來.
禮服進來的話.
都不會有什麼宗教的氣味.
我不會覺得女神我要下拜你.
都不會的.
因為我們那個社會的狀況.
和以前的社會狀況是截然不同的.
你可以穿金戴銀來.
你可以穿得很漂亮來.
我會稱讚你穿得很漂亮.
但是我沒有那種社會的動態.
社會經濟地位的那種動態.
所以21世紀和以前是不同的.
有了這個背景的時候.
你就明白教目書信.
特別是提摩太前書.

$^{441}$第一本書寫給提摩太.
有一本後書寫給他.
前書說什麼呢.
就是教會裡面一定有一群有才有勢的.
所謂的金主.
在這本書裡面你細讀的時候.
你就發現裡面有一些教會內部的問題.
就是裡面有一些異端滲入了來.
而很多時候教異端或者傳異端的人.
就是當然是那些主人.
因為他們有話語權.
有主人的權威.
所以第一章的時候.
你就看到他處理著.
本來是處理著當時教會裡面有異端滲入的問題.
有些人去議論一些東西或者搗亂一些東西.
用教導來使人有些迷失等等.
到第五章的時候.
他就說到一個很長的篇幅.
說到寡婦.
寡婦在前後文裡面.
就是一些比較有經濟資源的婦女.
她老公當然死了.
她就成為寡婦.
然後那些財產是誰.
你老公死了財產是給誰.
應該給你.
那你承受了那些財產.
那你就自然是比較有錢.
你也有一個大的物業.
所以教會的罪惡就自然會去你家裡.
所以五章十三節在說什麼.
保羅就吩咐這些寡婦.
他說他們.
年輕的寡婦.
說的是年輕的寡婦.
他們學了懶惰.
原來懶惰也可以學.
習慣去挨家閒逛.
不是到處去window shopping的意思.

$^{481}$而是他們當自己是權威.
就是沒有主權的.
在他耳上沒有主權.
然後到處走來走去做什麼.
不僅是懶惰.
是說長道短.
說長道短就是在胡說八道.
或者傳這傳那 胡說等等.
好管閒事說一些不該說的話.
你看到當時婦女的問題是什麼.
就是說長道短.
這個絕對不是聖經針對女人的問題.
這是當時的處境.
今天男人也可以說長道短.
對不對.
什麼時候的時代.
男女都可以說長道短.
只不過是當時的背景.
他們在說長道短.
所以保羅說他們是說長道短.
你要明白那個處境.
他們好管閒事都是一樣.
不只是女人才好管閒事.
不是聖經說女性的本性是好管閒事.
當時的教會裡面.
有些家庭教會裡面.
有人去到好管閒事.
說一些不該說的話.
都一樣.
男女都會說一些不該說的話.
所以這裡絕對不是說女性的本質.
而是當時的處境發生什麼事.
所以到第六章的時候.
保羅就處理有錢有地位的問題.
不要貪財.
你記得這句話嗎.
不要貪財.
所以你有上文下理的時候.
你就回到第二章.
他就處理一些在教會裡面有權有勢的人.

$^{521}$二章第八節.
就先說男人的問題.
男人都有很多問題.
憤怒爭論等等.
這些在教會裡面.
他們可能是貪權等等.
到第九至十幾節.
這段這麼長的聖經.
就是在討論教會裡面女人的問題.
如果有這樣的背景.
剛才我解釋的背景.
你再套進去的時候.
你就明白了.
保羅在處理某些女人的問題.
某一些女人的問題.
那些是什麼女人.
就是有財有勢的.
你想想 穿金戴銀.
今天我們不是每個人都有幾卡鑽石戒指進來.
不是每個女人都有這樣的錢.
和有這樣的福氣.
可以穿金戴銀去一些公眾場所.
當時他們處理這些女人.
某些女人的問題.
就是一些富有.
用自己裝扮來鞏固自己地位.
特別是在一些宗教和公共的場所裡面.
去鞏固自己地位.
當時教會就成為另一個公共的場所.
去到聚會的時候.
他們穿金戴銀就是想鞏固自己地位.
保羅就說.
叫男人安靜禱告.
女人也要去到安靜的學道.
不要穿上衣服去教會.
這是當時的背景.
不要以貴族的衣裳來取人的尊榮.
鞏固自己的地位.
二章第十一節 他說要安靜的學習.
因為剛才上文下理說了什麼.

$^{561}$一些有權有勢的人.
他們可能存異端.
或者一些有偏差的教導等等.
所以保羅就說.
這些女人要先好好學道 先受教.
二章第十二節.
他說我不許女人教導.
或者以往的翻譯就是「講道」.
不要管轄男人.
「管轄」這個字是有一點貶義的.
是壓制的意思 是奪權的意思.
保羅似乎針對的就是.
有一些被異端去到影響.
或者以強勢來壓倒人的那種行為.
所以按道理來說 不是一筆過就說完所有女人.
正如之前說的有男人是濫權.
這裡有些女人是濫權的.
所以保羅就說.
我不許你講道.
你先要去學習 學道.
接著下一節 十三節.
就說到亞當和夏威夷.
這個就是最棘手的地方.
就是很多時候我們解經的時候.
我們就會說 這個一定是講到創世的規律.
就是男人被創造先 女人才被創造後.
這個是創世第二章裡面所講的.
亞當夏威夷的記載.
我們就會說 既然男人是先的.
那男人應該是主.
女人是後.
所以女人是後的話 女人就應該要馴服.
我們第一個的感覺都是這樣的.
但是如果你看一下近東的創世故事的話.
你就知道後面創造的東西 才是最精彩的.
這個不難明白.
我們看世界盃.
我們看到後期追上來的才是精彩的.
特別是小國 這樣才是精彩的.
所以有創世故事 近東的背景的話.

$^{601}$你就知道這裡不是說亞當先做.
所以亞當要做主.
夏威夷就是低等一點 次等一點.
相反可能夏威夷是更加尊貴的.
更加尊貴.
那奈何為什麼他在這裡這樣說呢.
重點就是下一句.
就是夏威夷先被騙的那一位.
你這樣去看的時候 你就會發現夏威夷是後造的.
所以夏威夷是更加尊貴的.
但是他腳是被騙的那個.
其實保羅在說什麼呢.
就是在回應以弗所當時的處境.
女士是有尊貴的.
特別是這班富貴的婦女.
但是她們被什麼騙.
異端或者一些錯誤的教導所影響.
所以保羅說這句話的時候.
未必在解釋誰是上等誰是下等.
你再看創世記第二章.
很多時候兩夫妻在吵架的時候.
就說你是我的幫手.
你當然要幫我 聽我說.
因為創世記第二章說到夏娃被造的時候.
他是亞當的幫手.
你知道幫手這個字形容什麼嗎.
在《舊約聖經》裡面.
很多時候是形容耶和華的.
不是形容菲傭 不是形容工人.
耶和華是次等還是高等.
詩篇裡面經常說耶和華是我的幫助.
就是同一個字.
所以你回去看的時候.
女人雖然是後造 但是她有尊榮.
保羅去引用這句話.
似乎就是針對當時的背景.
富貴的女人你們要好好學到.
因為你們有尊榮.
但是你記得那個故事嗎.
有尊榮的都會被異端影響.

$^{641}$所以你需要好好學到.
不要急於想教化人或者奪權等等.
這些所謂的處境性改正.
不是用一個判詞去說所有的女人.
舉個很簡單的例子.
我兒子小過五歲的時候.
他進廚房 我跟他說.
不准小孩子碰刀.
不要用刀來玩.
說得對不對.
是那個處境裡面需要改正.
但是你說今天他二十多歲了.
我會跟他說同一句話嗎.
不會吧.
他可能反過來說你不要碰刀.
你理論盡 你不要碰刀.
是一種處境性的說話.
所以如果你這樣看的時候.
你就發現不許女人講道或者教道.
其實都是處境性的.
二章十五節.
就更難搞.
藉著生產得救.
是不是藉著生小孩得救呢.
又未必是的.
因為當時講到生產的時候.
giving birth.
很多時候希臘的哲學家或者文學家.
他們講到生產giving birth的時候.
是在講一個人的德性.
他的品性生產出來.
生長出來.
你們明白我的意思嗎.
是一個比方來的.
是一個比喻來的.
你藉著生產得救.
就是keep growing your virtues.
去鍛鍊你的品格.
以致你作成得救的功夫.
這樣的意思.

$^{681}$作成得救的生命.
所以不是說女人你們持守信心愛心.
和克制就可以了.
就必定進去產房你們就得救了.
不是這樣的意思.
現代人可以這樣讀.
但我剛才開頭講什麼.
讀經是不可以這樣的.
對不對.
等於創世記.
創世記第一章六一創造.
你不可以用現代的想法.
我覺得現在是這樣的意思.
所以我就將它立刻套上去.
你一定要知道聖經本身的背景.
問一問背景的問題.
然後你才明白為我們寫的是什麼.
有學者說在聖經裡面.
女性的聲音其實是很少的.
聖經裡面有九十三個出過場的女人.
只是有四十九位是有名有姓的.
譬如路德等等不同的婦女.
是親口講過說話的.
但如果加起來這些說話.
加起來只是整本聖經裡面的1.1\%.
是不是很少.
我之前曾經講過婦女的議題.
很多年前.
我都說聖經是男人的書.
意思是說是男作者寫的書.
並且寫這本書的時候.
當時的文化背景是以男為上的文化背景.
因為聖經除了是上帝的說話之外.
也是人的說話.
百分之百人的說話.
當然一定是有他的文化背景.
有文化背景當時寫出來的時候.
就不代表他一定是贊成.
或者去到造成當時的文化背景.
因為上帝啟示他的說話.

$^{721}$在當時的文化裡面.
就是用當時人文化所熟悉的概念.
來跟他們溝通.
但往往聖經裡面不同的地方.
特別是關鍵的地方.
上帝會放一個女人的故事在裡面.
雖然只不過是百分之一.
少得很微小.
但聖經的作者看清楚的時候.
就是一個男為上的文化裡面.
女人突然說話出來.
並且她說的話是適合男人說的話.
她做的事很多時候是男人應該做的.
那你說是贊成當時的以男為上的文化.
還是在挑戰當時以男為上的文化.
後者.
所以聖經雖然是在那個文化裡面.
以男為上寫出來.
上帝這樣啟示出來.
但往往很多時候是去挑戰那個文化的處境.
這樣是不對的.
不要以男為上.
因為男女是同等的.
你再看《藏世紀志》啟示錄整個大的故事裡面.
你會發現女人的角色不斷不斷在擴張.
不斷在舞台上扮演更加重要的角色.
譬如耶穌復活的故事.
耶穌復活的故事.
人們找到一個空墳墓.
第一個被警察派去報復活的訊息的是誰.
是莫大拉瑪尼亞.
莫大拉瑪尼亞的男人都去了哪裡.
男人都非常害怕.
然後躲起來.
所以是女人被警察派出去的.
所以有很多身軀弱的婦女.
其實是上帝使用她們的.
還有保羅說在基督裡面再不分男與女.
地位上,角色上都不會分男與女.
見到整個聖經的大故事的發展軌道.

$^{761}$是男與女一起去童宮.
更加不要說在《史萊恩傳》裡面.
有其他的說到女人的角色等等.
如果綜合來說.
不是因為他們說話少就不重要.
相反很多時候說話少.
往往才是我們需要留心的.
對嗎.
等於你聽一篇音樂.
不是一路從頭到尾才是尾的音樂.
很多時候就是中間靜了.
有休止符.
休止符不一定是佔整篇樂章裡面大部分的.
但沒有休止符你就很麻煩了.
音樂變成噪音.
男與女在聖經裡面的出現都是一樣.
如果你這樣整體去處理當時的背景.
普羅是處理一個很急切的處境性的問題.
他不是推翻女士教導的立場等等.
是當時的處境有人濫用這些權柄.
就不代表他一鋪過一言堂.
說了女人就是不可以這樣.
在我們最後的時候.
我們需要想一下為何他對我們切身.
因為我們讀聖經的時候.
看聖經的時候.
無論你聽人說.
或者你自己靈修.
或者你在講道裡面.
主要學裡面你聽到人.
去到講聖經的時候.
或者你查經的時候.
你需要去到留意什麼.
我們一定要留意幾樣東西.
一打開聖經或者聽到聖經的時候.
你首先要明白.
是有前面的世界的.
拿著一本聖經或者你的iPhone.
有經文 有訊息.
那個聖經的訊息是有聖經前面的世界.

$^{801}$前面的世界就是讀者自己的世界.
我今天讀聖經.
自然有我這個現代人.
21世紀這個世界.
如果我是出生於以男為上的家庭背景的話.
我覺得男人比女人厲害.
如果我是大男人主義的.
我拿起一本聖經的時候.
我就自然將大男人主義套上去.
對嗎.
但是這樣做不是很妥當.
不是很妥當.
因為我沒有留意自己前面這個世界.
有時候我的前設.
或者我自己有一個概念.
我需要把它打掉.
才能夠聽到聖經裡面所講的東西.
然後我們需要留意聖經背後的世界.
那個處境.
無論那句話語都有它的處境.
對嗎.
其中一個例子很簡單.
如果我說911的時候.
大家有那個處境.
你就會明白911並非純粹一個簡單的數字.
對嗎.
2001年之前.
你講911就沒有這樣的意思.
2001年之後.
911之後.
大家就明白911是什麼.
對嗎.
所以有那個處境.
你需要明白那句話.
聖經其實都是一樣.
你需要明白當時的背景.
你看提摩太前書.
你不能夠不鬥.
當時第一世紀以弗所.
家庭教會裡面發生的事.

$^{841}$第三方面就是聖經裡面的世界.
就是剛才我說的上文下理.
背景怎麼講都好.
上文下理怎麼講男人女人.
女人受什麼.
當時有些女人受什麼影響等等.
所以我們需要留意這三方面.
對嗎.
每當有人引經據典的時候.
無論在新聞報道上.
或者其他的報道上.
你聽到.
你都需要問這三個問題.
讀經的時候.
有沒有先思考我們那套東西.
有沒有記錄了我們那套東西.
聖經當時的背景.
到底怎樣去提醒你的聖經的了解.
然後上文下理.
傳統本身未必一定是錯的.
但如果我們將傳統等同於.
聖經自己的絕對權威的話.
我們需要耍耍車.
需要想一想.
在讀經的方法處理上.
是不是有時候太過倉促.
有什麼為我們而寫.
當然不是說我們不要穿金戴銀.
不是這樣直接自勉這樣形容.
起碼有三樣東西去到想到.
第一樣.
提摩太前書當時的背景.
和我們今天都一樣.
講清楚教會是一個怎樣的地方.
是高舉耶穌基督的地方.
教會並非一個建立自己地位的地方.
建立自己地位的地方.
這個特別對領袖是一個警告.
我們作為領袖.
我們很多時候會來到教會的時候.

$^{881}$我們在教會的群眾裡面.
我們會自自然做一些事.
我們想去到建立或者鞏固自己的地位.
第二樣東西.
每個人都應該服在主導之下.
剛才我說的最高的權威是什麼.
是聖經.
我們應該永遠都是回去看清楚聖經.
說些什麼.
所以我們每一個人都需要耐心去學到.
第三方面基督徒.
基督徒都要有批判性思考.
一定要有.
我們不要說一件事.
憑信心拍胸口就去.
不一定的.
我們需要有批判性思考.
批判性思考.
這個是靈性的成長其中一部分.
想清楚那些東西是不是邏輯.
想清楚那些東西是不是合理.
梳理材料的時候.
那些資料的時候.
有沒有有一些疏忽了.
這些全部我們都需要去學習.
基督教信仰請你演繹的時候.
你不要演繹成為一個盲信的信仰.
是一個有道理.
有證據.
有很多其他材料去整理出來.
一套的生命的世界觀.
這個是基督教信仰.
所以我們演繹的時候.
我們去說基督教的時候.
說我們的信仰的時候.
我們都需要做一個懂得去思考的人.
希望這堂課給大家.
有少少的材料.
給大家一些提醒好嗎.
我們同心得來祈禱.

$^{921}$然後我們領受聖餐.
主愛多謝你.
因為你是那一位.
去對我們說話的上帝.
不單止在當代的信徒教會裡面.
你去作提醒.
去說話.
你更加是那一位.
今時今日仍然說話的神.
主愛願意我們去處理一些.
信仰上的問題的時候.
我們懂得去學習去思考.
因為你是華語.
是神聖的華語.
既然是神聖的華語.
就更加需要我們去下功夫.
思考清楚.
想清楚你所說的.
讓我們更加認識你.
讓教會活在你主導之下.
耐心地學習你的道.
並且去行動.
今日這個聖餐.
都是為我們每一個門徒而設的.
主愛多謝你對我們的愛.
在歷史時空裡面.
耶穌基督在十字架上為我們所做的.
為我們聚而死.
赦免我們.
埋屬我們.
以致我們能夠成為你自己的子民.
願以我們感恩的心來領受.
祈禱奉尊名感恩.
\newpage



\section{3:7-13}
\label{sec:MlRDdW70h9M}
\textbf{在基督裡坦然無懼(弗3\_7-13)張威傳道 Boldness in Christ (Eph 3\_7-13) Paul Zhang}
\newline
\newline
連結: \href{https://youtube.com/watch?v=MlRDdW70h9M}{\texttt{https://youtube.com/watch?v=MlRDdW70h9M}} ~~~~ 語音日期: 2026-07-12
\newline
\newline
\hyperref[sec:IBF8VFRUG4s]{\small{< < < PREV SERMON < < <}}
~
\hyperref[sec:index_chronic]{\small{[返順時目]}}
~
\hyperref[sec:index_scriptual]{\small{[返順卷目]}}
~
\hyperref[sec:YT4UObZ7V44]{\small{> > > NEXT SERMON > > >}}
\newline
\newline
$^{1}$(speaking in foreign language).
- Good morning, brothers and sisters..
(speaking in foreign language).
Sorry, let's read Ephesians 3, 7 to 13..
(speaking in foreign language).
Of this gospel, I was made a minister.
according to the gift of God's grace,.
which was given me by the working of his power..
To me, though I am the very least of all the saints,.
this grace was given to preach to the Gentiles.
the unsearchable riches of Christ,.
and to bring the light to everyone.
what is planned of the mystery hidden for ages in God,.
who created all things..
(speaking in foreign language).
So that through the church and the manifold.
wisdom of God that might now be made known.
to the rulers and authorities in the heavenly places..
This was according to the internal purpose.
that he has realized in Christ Jesus our Lord,.
in whom we have boldness and access with confidence.
through our faith in him..
So I ask you not to lose heart over what I'm suffering.
for you, which is your glory..
(speaking in foreign language).
And have you ever woke up in the middle of the night.
with your full heart of fear and anxiety?.
(speaking in foreign language).
So maybe you're worried about work,.
maybe money, your health,.
feeling like maybe you are not good enough,.
scared about how people look at you.
and maybe look down upon you..
(speaking in foreign language).
I've been there, I've been tried this before..
Sometimes I woke up in the middle of the night.
in such pressure and then I'm even stressed out.
in my dreams..
(speaking in foreign language).
A lot of people look strong and busy outside,.

$^{41}$but deep down they are living in constant fear..
(speaking in foreign language).
We don't just fear about our consistence..
(speaking in foreign language).
And at the deepest level, people are actually afraid of God..
We know instinctively that God is holy and we are not..
After Adam sinned, his first reaction wasn't to run to God,.
it was to hide from him..
People today do the same thing..
We use success and business to cover up the emptiness.
of being separated from God..
(speaking in foreign language).
But there is the amazing part of gospel.
is God doesn't wait for us to become worthy or perfect.
before he lets us become close to him..
While we are still sinner,.
Christ already opened the way for us..
(speaking in foreign language).
Ephesians 3, 7 to 13 shows us that a person.
who used to tremble in fear before God.
can now come to his presence with complete confidence.
and it all because of Christ..
(speaking in foreign language).
Today, we will look at three things from this passage..
The number one thing that we will cover today.
is by grace, no longer fear our unworthiness..
(speaking in foreign language).
I became a servant of the gospel.
according to the gift of God, of God's grace..
The word minister or servant in the original text.
means a servant or a person who serves..
Paul did not begin by emphasizing I'm an apostle..
Instead, and a stress that he was a servant of the gospel.
and this is very important because in ministry,.
the greatest danger is not weakness, but pride..
(speaking in foreign language).
So Paul understood clearly that his ministry.
was not based on his own qualification,.
but on God's free and unmerited grace..
(speaking in foreign language).

$^{81}$Paul had once persecuted the church.
and arrested Christian and was even glad.
when Stephen was stoned to death..
Think about it, on the road of Damascus,.
Paul was not going to worship God, but to arrest Christians..
He thought he was serving God,.
but in reality, he was opposing him..
(speaking in foreign language).
But on that road, the rising Christ.
took the initiative to find him..
(speaking in foreign language).
Why are you persecuting me?.
(speaking in foreign language).
Only then did Paul truly realize that he had been doing..
That is grace..
Grace is not man seeking God,.
but God actively seeking and saving.
those who were once his enemies..
(speaking in foreign language).
Paul said, "I am the very least of all the saints.".
He was not pretending to be humble.
or trying to put in on a spiritual show..
These are the words of someone who truly knew himself..
(speaking in foreign language).
Paul knew that if Christ had not stopped him on that road,.
he would never ever have returned to God by himself..
The more he understood grace,.
the more he realized how unworthy he was..
(speaking in foreign language).
Today, many Christians are constantly fixated.
on their own unworthiness..
"My education is not good enough..
"I'm not eloquent..
"My life has too many problems.".
They say, "When I become more mature,.
"maybe one day I will serve..
"Or when I understand the Bible better,.
"then I will got involved in church.".
(speaking in foreign language).
But in the Bible, there was not even one single person.

$^{121}$whom God used only after they have become perfect..
God does not use people because they are already complete,.
but because Christ is complete..
(speaking in foreign language).
So God accepts unworthy people.
and uses unworthy people in Christ..
(speaking in foreign language).
So this passage also gives great hope.
to those who have failed..
Some people feel, "I have failed too many times..
"I have let God down too much..
"God must have given up on me.".
We believe God can forgive others,.
but we find it hard to believe He will forgive us too..
(speaking in foreign language).
If God did not give up on a person like Paul,.
then we today are not beyond His grace either..
The greatest hope for the gospel is this..
God saves not the pretty good people,.
but the real sinners..
(speaking in foreign language).
So what truly makes us stand firm is not ourselves,.
but the grace of Christ..
(speaking in foreign language).
By the cross, no longer living in separation..
(speaking in foreign language).
To preach to the Gentiles.
the unsearchable riches of Christ,.
these unsearchable riches do not refer.
merely to material blessing,.
but to God's salvation,.
the forgiveness of sins,.
sanctification, and presence of Christ..
Christ is not a poor savior..
He is a rich God, and what He gives is true abundance..
We may not lack material things,.
but we often fail to enjoy the riches we have in Christ..
(speaking in foreign language).
Paul then speaks about mystery..
As mentioned before, this mystery.

$^{161}$is not something mysterious,.
but a truth that was hidden for ages.
and has now been revealed..
(speaking in foreign language).
What is the truth?.
The Gentiles and Jews become one body in Christ..
(speaking in foreign language).
We consider this quite normal today,.
but in the first century, it was a shocking declaration..
(speaking in foreign language).
In the first century, there was a huge divide.
between Jews and Gentiles..
Jews regarded Gentiles as unclean and far from God..
They were even unwilling to enter a Gentile's house..
Some Jew would rather take a long detour.
than pass through Samaria..
For a Gentile to be saved,.
they first have to become a Jew,.
be circumcised, and keep the law..
(speaking in foreign language).
So there was a deep war between people,.
a war of ethnicity, a war of culture,.
a war of religion, and a war of identity..
But the cross demolished the dividing wall of hostility..
Those who were once enemies are now in Christ's one family..
(speaking in foreign language).
The deeper meaning is found in verse 10..
God did this so that through the church.
and the manifold wisdom of God.
might now be made known to the rulers.
and authorities in the heavenly places..
When a once proud Jew and a former idol-worshiping Gentiles.
they worship together, love each other,.
and serve each other in the church..
Even the angels are watching attentively.
and exclaimingly in amazement..
This is the manifold wisdom of God..
(speaking in foreign language).
So church is no longer merely a religion organization..
It's a living testimony of God's redemptive wisdom.

$^{201}$displayed in the universe..
(speaking in foreign language).
Brothers and sisters,.
this world is becoming increasingly divided.
between nations, between ethnic groups, between individuals..
People attack each other.
because of cultural and background differences..
There are conflicts everywhere..
Online, I myself used to argue.
and even sometimes curse at others in the internet..
There are conflicts even at home.
and it's becoming harder and harder.
for people to trust one another..
But the church should be the place.
where the world sees the gospel..
(speaking in foreign language).
The church is not a social club or interest group.
where similar people gather..
The true church consists of people.
who never naturally be together.
yet have become one body because of Christ..
(speaking in foreign language).
There are people of different cultural backgrounds,.
different levels of education, different personality,.
yet because of Christ,.
they accept one another and love each other..
(speaking in foreign language).
The world will ask, why can these people love each other?.
The answer is not because they are all wonderful people,.
but because of Christ..
(speaking in foreign language).
So when the church lives in such unity,.
not only will this world see the reality of gospel,.
even the rulers and authorities in the heavenly realms.
will see the unsearchable wisdom of God..
(speaking in foreign language).
What truly unites us is not ourselves,.
but the cross of Christ..
(speaking in foreign language).
By Christ, we can boldly come to the Father..

$^{241}$(speaking in foreign language).
Don't take this statement for granted today..
In Old Testament time,.
people could not casually approach God..
(speaking in foreign language).
When God revealed himself on Mount Sinai,.
everyone was terrified,.
and only Moses was allowed to represent them before God..
Later, in the temple, only the high priests.
could carefully enter the most holy place.
to offer sacrifices,.
because the sinner could not bear God's holiness..
(speaking in foreign language).
But now, because of Jesus,.
we can come boldly into God's presence..
This boldness is not because we have become worthy,.
but because Jesus has accomplished everything..
Jesus' precious blood has removed our sins.
and torn the veil of most holy place..
The book of Hebrews says,.
"Let us then with confidence.
"draw near to the throne of the grace..
"It's not by our good work,.
"but by the merit of Christ.".
(speaking in foreign language).
Many Christians still live in the state of hiding from God..
As soon as they sin, they will not pray..
Then when they become weak.
and they start staying away from church meeting.
because they feel God must be very disappointed in me..
(speaking in foreign language).
But gospel does not say you have no sin..
It says Christ has already borne your sin for you..
Therefore, today you do not come to God.
based on your own performance,.
but based on Christ..
(speaking in foreign language).
We can remain bold even in suffering..
Paul continues,.
"So I ask you not to lose heart.

$^{281}$"over what I'm suffering for you,.
"which is your glory.".
(speaking in foreign language).
At that time, Paul was in prison..
Yet instead of asking the Ephesian church.
to feel sorry for him,.
he said, "This suffering were their glory.".
Why?.
Because his suffering allowed the gospel.
to reach Gentiles and God's mystery to be unfolded..
The world saw chains and shame,.
but Paul saw God's glory and benefit of the church..
(speaking in foreign language).
This brings tremendous comfort to us today..
Many Christian easily become discouraged in trials..
Financial burdens, family problems,.
difficulties at work, rebellious children,.
health issues, so on..
We often ask, "God, why me?.
"Why do others have so easy life.
"while I struggle so much?".
(speaking in foreign language).
The gospel does not promise a smooth and easy life..
It promised us that even in the darkness trial,.
we can still come boldly to our heavenly father..
(speaking in foreign language).
The sign of God's presence is not necessarily.
a smooth environment,.
but the presence of peace, hope,.
and the ability to call on his name.
even in the midst of hardship..
(speaking in foreign language).
True strength is not grinding our teeth.
and enduring by ourselves,.
but coming to God through Christ even in our weakness..
Just as Paul experienced in prison..
Often the most difficult times.
become the times when we draw closest to God.
and experience his great power most clearly..
(speaking in foreign language).

$^{321}$What truly makes us bold and fearless.
is not ourselves, but Christ..
(speaking in foreign language).
Brothers and sisters,.
let's honestly ask ourselves this question today..
Am I still living in the fear of not being good enough,.
afraid to serve or fully commit myself?.
Am I still living in separation,.
having barriers with brothers and sisters.
or distance from God?.
Am I still living in discouragement because of trials,.
afraid of the future, afraid of my circumstance,.
afraid I won't make it?.
The gospel declares to us today,.
in Christ, we can be bold and fearless..
(speaking in foreign language).
Come to God right now,.
no matter how much fear in your heart..
Quiet yourself and come boldly.
to the throne of grace by faith..
Tell God all your burdens..
If you have conflict with a brother or sister,.
take the initiative to seek reconciliation.
so that the church can be a testimony of unity.
in this divided world..
Stand firm in the midst of trials,.
entrust the difficulties you are facing to God.
and ask Lord to help you not lose heart,.
but instead, see how even your suffering.
can become God's glory..
(speaking in foreign language).
Let's pray..
(speaking in foreign language).
- Heavenly Father, we are unworthy people,.
but in Jesus Christ, you have opened a way for us..
Help us, please help us to look more to Christ.
so that in our weakness and fear,.
we may still come to you..
Help us to believe that you are sovereign,.
even in trials..

$^{361}$Let the church become witness of unity.
in this broken world..
In the name of Jesus..
\newpage



\section{約翰福音 4:3-14}
\label{sec:YT4UObZ7V44}
\textbf{放下水桶 (約翰福音4\_3-14) - 梁紫健傳道}
\newline
\newline
連結: \href{https://youtube.com/watch?v=YT4UObZ7V44}{\texttt{https://youtube.com/watch?v=YT4UObZ7V44}} ~~~~ 語音日期: 2026-07-12
\newline
\newline
\hyperref[sec:MlRDdW70h9M]{\small{< < < PREV SERMON < < <}}
~
\hyperref[sec:index_chronic]{\small{[返順時目]}}
~
\hyperref[sec:index_scriptual]{\small{[返順卷目]}}
~
\hyperref[sec:rfQ3ZwOlzM4]{\small{> > > NEXT SERMON > > >}}
\newline
\newline
約翰福音 4:3-14
\newline
\begin{longtable}{cl}
\hline
\hline
章節 & 經文 (和合本修訂版)\\
\hline
4:3 & \begin{tabularx}{0.7\textwidth}{X} 他就離開猶太，又回加利利去。 \end{tabularx} \\ \\ \relax
4:4 & \begin{tabularx}{0.7\textwidth}{X} 他必須經過撒瑪利亞， \end{tabularx} \\ \\ \relax
4:5 & \begin{tabularx}{0.7\textwidth}{X} 於是到了撒瑪利亞的一座城，名叫敘加，靠近雅各給他兒子約瑟的那塊地。 \end{tabularx} \\ \\ \relax
4:6 & \begin{tabularx}{0.7\textwidth}{X} 雅各井就在那裡；耶穌因旅途疲乏，坐在井旁。那時約是正午。 \end{tabularx} \\ \\ \relax
4:7 & \begin{tabularx}{0.7\textwidth}{X} 有一個撒瑪利亞婦人來打水。耶穌對她說：「請給我水喝。」 \end{tabularx} \\ \\ \relax
4:8 & \begin{tabularx}{0.7\textwidth}{X} 因為那時門徒進城買食物去了。 \end{tabularx} \\ \\ \relax
4:9 & \begin{tabularx}{0.7\textwidth}{X} 撒瑪利亞婦人對他說：「你是猶太人，怎麼向我一個撒瑪利亞女人要水喝呢？」因為猶太人和撒瑪利亞人沒有來往。 \end{tabularx} \\ \\ \relax
4:10 & \begin{tabularx}{0.7\textwidth}{X} 耶穌回答她說：「你若知道神的恩賜，和對你說『請給我水喝』的是誰，你早就會求他，他也早就會給了你活水。」 \end{tabularx} \\ \\ \relax
4:11 & \begin{tabularx}{0.7\textwidth}{X} 婦人對耶穌說：「先生，你沒有打水的器具，井又深，哪裡去取活水呢？ \end{tabularx} \\ \\ \relax
4:12 & \begin{tabularx}{0.7\textwidth}{X} 我們的祖宗雅各把這井留給我們，他自己和兒女以及牲畜都喝這井裡的水，難道你比他還大嗎？」 \end{tabularx} \\ \\ \relax
4:13 & \begin{tabularx}{0.7\textwidth}{X} 耶穌回答，對她說：「凡喝這水的，還要再渴； \end{tabularx} \\ \\ \relax
4:14 & \begin{tabularx}{0.7\textwidth}{X} 誰喝我所賜的水，就永遠不渴。我所賜的水要在他裡面成為泉源，直湧到永生。」 \end{tabularx} \\ \\
[1ex]
\hline
\hline
\end{longtable}
$^{1}$謝謝 各位早安.
不知道大家早上起床第一件事是做什麼呢.
哇 這麼開心 這麼好.
第一件事不是張開眼睛.
第一件事 除了張開眼睛之外.
有沒有人第一件事是拿你的電話出來的.
或者是ipad 或者是鬧鐘.
如果你第一件事起床除了張開眼睛之外.
你是拿ipad或者拿電話的 請舉手.
OK 都不少.
不知道大家有沒有試過在深夜本身有點睏.
結果你打算有自己的me time 正呀.
打算看手機或者ipad.
看了半個小時 看了一個小時.
然後經常都會覺得 為什麼自己這麼晚睡.
明明十點坐在床上 看到十一點.
有這樣的情況.
我們經常都是這樣.
我們經常都會看message 看很多不同的社交媒體.
看短片.
其實短片都給我們一些安慰.
或者是分散我們日常的注意力.
覺得我們工作很辛苦.
照顧完小朋友 終於有自己的時間.
那就看一下.
電話就好像一個無底黑洞.
你不斷去loop 去看新的post 短片.
你透過很多自己的時間.
希望我們拿回一些屬於自己的空間.
屬於自己的滿足回來.
但問題是 大家不知道有沒有發現.
你會覺得好像永遠都不夠.
loop完一次 看完一次.
應該都開始覺得滿足.
但又想loop第二次.
心裡面那份渴望 始終都沒有被解決.
loop到累了 就ok了.
我終於可以睡覺了.
但第二天 早上又來.
每一天重複重複跌入這個循環.

$^{41}$我相信大家都很熟悉這個感覺.
用電話 用電子世界去填補這份空虛.
那份渴望 好像沒有真正被滿足.
接下來我們會看一下.
經文裡面有一個婦人.
她一生都在尋找不同的滿足.
希望能夠滿足自己.
直到有一天 她在井旁.
遇到一個讓她有滿足的人.
我知道FF家庭團契的弟兄姊妹.
下星期會用這段經文去查經.
所以請你特別留心.
今天的經文會帶給我們不同的思考.
我們一起來.
讀一下今天的經文.
約翰福音四章.
耶穌離開猶太 回到加利利去.
他必須經過撒瑪利亞.
於是到了撒瑪利亞的一座城叫敘迦.
接近雅各給他的兒子約瑟的地.
雅各井在這裡.
耶穌因為累就坐在井旁.
那時候是正午.
有一個撒瑪利亞婦人來打水.
耶穌就跟她說 請給我水喝.
大家可以看到我高亮了一些字.
有些高了和低了.
值得留意的是.
經文寫耶穌必須經過撒瑪利亞.
是不是因為這條路近一點.
所以是必須呢.
經文提到雅各給他的兒子約瑟的地.
雅各井.
這個也是一個重要的地方.
繼續讀.
第八節.
因為那時門徒進城買食物.
撒瑪利亞婦人對他說.
你是猶太人.
為什麼要向我一個撒瑪利亞女人要水喝呢.

$^{81}$因為猶太人和撒瑪利亞人沒有來往.
為什麼耶穌會親自開口.
主動去問一個撒瑪利亞婦人去說話.
也去給她借水喝.
這裡的沒有來往其實很值得留意.
寫得很隱晦.
究竟實際上代表什麼呢.
第十節.
耶穌就回答.
如果你知道上帝的恩賜.
跟你說給水喝的是誰.
你早就會求他.
他也會早就給你活水.
這裡提到活水.
究竟活水和這個婦人喝的井水有什麼分別.
看完經文之後.
我們會有三個問題.
我們收集這個問題.
就是聖經作者帶著我們去理解他的思路.
第一,第四支話.
耶穌必須經過撒瑪利亞.
必須因為什麼呢.
第二,耶穌作為一個猶太男人.
為什麼會和一個撒瑪利亞婦人去說話.
而且去問她借水喝呢.
第三,什麼是活水.
和婦人喝的井水有什麼分別.
我們就帶著這幾條問題去看今天的經文.
回答第一個問題.
耶穌必須經過撒瑪利亞.
就要明白這個必須要記出期.
這幅是地圖.
在說這幅地圖之前.
我要說一下猶太人和撒瑪利亞人的淵源.
在北國以色列亡國之後.
亞述王就將外族人遷入撒瑪利亞.
和當地人通婚.
他們的後代就是撒瑪利亞人.
他們只接受摩西吳敬.
還在基里森山建了自己的聖殿.

$^{121}$在主前大概一百二十幾年左右.
這個聖殿就被猶太人帶兵毀滅了.
另一邊,北國的以色列.
就是撒瑪利亞人的祖先.
剛剛說過.
曾經聯合外邦人去剿滅殘殺猶太國.
大家可以想像.
這兩族人雖然同源.
但是他們的仇恨很深.
因為他們互相殘殺.
所以在虔誠的猶太人.
在猶太地.
他要去到加利利的話.
他寧願繞路.
多走幾天都不願意踏入撒瑪利亞.
大家看看這幅地圖就知道.
最南面,橙色的部分是猶太地.
橙色的部分.
有耶路撒冷,有伯利恆.
中間那一塊就是撒瑪利亞.
有雅各井.
然後之後北面.
綠色那一塊是加利利.
有拿撒勒.
大家看到在地圖中間.
有一個紫色的箭咀和一條線.
大家看到嗎.
很可能就是耶穌當時走的路.
由耶路撒冷經過雅各井.
撒瑪利亞再上到北面的加利利.
當然你可以想像到.
猶太人是不會走這條路的.
猶太人有另一條路繞過撒瑪利亞.
就是紅色那條路.
紅色那條路線.
由約旦河東面直上繞過中間那塊.
就是黃色那塊地.
然後才到北面.
綠色那塊是加利利.
猶太人寧願走很遠的路.

$^{161}$都不想和撒瑪利亞人有接觸.
所以耶穌的必須.
是不是因為趕時間走捷徑呢.
這裡的必須不是在說地理.
是在說上帝的救恩計劃.
耶穌是必須走進去.
因為那裡有一個人.
甚至有一個民族.
必須經歷這個救恩.
就好像我們平時用GPS找路.
有時有些區域我們是不想走的.
就好像我這樣.
我去年撞車.
有時我都不想再走那條高速公路.
因為有些陰影.
你會想繞過那裡.
但耶穌的導航不是計時間.
是計人的.
他專門選了一條路.
為了一個避開全村人的婦女.
因為當時的婦女.
通常打水的時間是什麼時間.
是早上或者下午黃昏的時間.
很少會在中午的時候去打水.
因為中午.
你想想在中東.
熱到什麼地步.
曬到太陽最猛烈的時候.
又乾又熱.
你可以想像到.
那位婦女很明顯是一個孤獨的人.
她連打水都要避開人.
耶穌就專門在最曬的中午.
走進撒瑪利亞找這個婦人.
解決了第一個問題.
我們就要問.
第二個問題就是.
為什麼耶穌會開口借水喝.
而回答這個問題.
第九節有一個很關鍵的句子.

$^{201}$就是猶太人和撒瑪利亞人沒有來往.
原文的意思還要具體.
就是撒瑪利亞人和猶太人之間沒用.
沒用是什麼.
不是說他們沒用.
是說在井旁的場景裡面.
他們沒有共用器皿.
猶太人覺得.
撒瑪利亞人是什麼.
混種人士.
整個民族都是不潔的.
他用他的杯子.
他用他的水桶.
他碰一碰.
都會玷污了整個猶太人的習俗.
所以有點像是.
大家經歷了疫情之後就會明白.
那幾年我們跟別人吃飯.
我們用公飯.
戴口罩.
搓手液.
每個人都會用.
別人的杯子你會拿來喝嗎.
不會吧.
如果有一個陌生人突然走過來.
說可不可以借你的杯子來喝水.
你會怎樣.
他用完之後我當然不會用.
就是這樣.
所以經歷了疫情之後.
每個人的防線都特別明顯.
當時的猶太人和撒瑪利亞人.
經歷了打仗.
互槍殘殺之後.
更加是見面都不會見.
更遑論喝他手碰過的水.
是不是.
所以耶穌這一句.
請給我水飲.
已經打破了兩重的界限.

$^{241}$第一就是猶太人和撒瑪利亞人之間.
那種不潔的界限.
第二就是男人和女人之間的界限.
當時在耶穌的年代.
特別是猶太拉比.
通常會在公眾場合.
避免和非親屬的婦女說話.
這是一個社會的禁忌.
所以當耶穌專登進入撒瑪利亞的時候.
走近這個婦人.
還要開口問她拿水飲.
這個舉動是一個很大的突破.
沒有什麼猶太人會這樣做.
但是當這個井旁婦人.
用兩軸人的界線拒絕耶穌的時候.
耶穌的回答聽起來有點九不搭八.
就說如果你知道上帝的恩賜.
和你說給我水飲的是誰.
你早就會求他.
而且他會給你活水.
當婦人還以為.
這個活水就是那個流動的水.
耶穌說了一句金句.
就是說.
凡喝者水的還要再喝.
誰喝我所賜的水就永遠不喝.
我所賜的水要在他裡面成為泉源.
直湧到湖泳生.
那麼耶穌所指的活水是什麼.
這是我們第三條的問題.
婦人一開始聽錯了.
活水這個詞語其實一語相關.
上帝在井底的死水.
次免去解作活動的山泉水.
所以她才會這麼說.
井這麼深.
你又沒有打水的器皿.
你憑什麼要活水.
難道你比雅各更厲害嗎.
她這麼說.

$^{281}$當婦人還在考慮.
怎樣實際上打水的時候.
耶穌已經將層次提升到生命的層次.
留意約翰的用字很仔細.
他說雅各井的時候.
是祖宗留下的井水.
但是耶穌所賜的水.
他用了另外一個字.
泉源.
我所賜的水要在裡面成為泉源.
直到永生.
一個是靠人力一桶桶地打出來.
另一個是不斷湧出來.
不用你打的.
所以我們進一步問.
究竟活水,泉源是什麼.
大家如果將聖經翻後幾章.
約翰福音第七章.
約翰自己解了這個謎底.
如果我們帶著這個敘事的背景.
我們才會明白.
其實耶穌為什麼在下文提到.
婦人五個丈夫的對話.
祂不是因為想八卦,唱衰這個婦人.
而是因為這個婦人.
代表著整個撒瑪利亞的民族.
一直在尋找水源.
一直在尋找泉源.
一直都尚未找到.
所以聊聊活水.
耶穌就說叫你丈夫來吧.
這個婦人就說我沒有丈夫.
耶穌就說.
你說得對.
你現在已經有五個丈夫.
而且現在那個都不是你的丈夫.
很多人讀到這裡.
很快就會將這個婦人定罪.
覺得她私生活不檢點.
但是大家慢慢看一下這段經文.

$^{321}$耶穌有沒有說過一句.
祂是罪人.
耶穌有沒有說過.
你悔改吧.
一句都沒有.
對比起《約翰福音》第八章行淫的婦人.
耶穌起碼有一句說.
從此不要再犯罪了.
但是這裡一句都沒有.
在那個時代.
其實婦女基本上是沒有離婚的權利.
有妻的主動權是在男人手上.
一個女人嫁五次.
很可能不是她選的.
很可能是因為她五次被離婚.
或者五次守寡.
無論是被離婚或者守寡.
都是代表著一次又一次的被遺棄.
一次又一次的失去了依靠.
大家還記得她中午來打水嗎.
很可能不是因為她生活不檢點.
是因為她一直找的依靠.
到頭來沒有一個靠得住.
是因為她一直找的依靠.
到頭來沒有一個靠得住.
更加令人驚喜的是.
這五個丈夫其實是影射著整個撒瑪利亞的民族.
給大家看一段經文.
如果大家有看獵王幾下十七章.
你就會發現在北角亡國之後.
亞述王就將五族人遷入了撒瑪利亞.
這五族人帶著五族的神明.
而撒瑪利亞人又拜了一份.
經文是這樣說的.
亞述王在這五個地方.
巴比倫,古他等等的地方.
帶了很多人來撒瑪利亞.
各族的人住在城裡.
為自己製造神像.
安置在撒瑪利亞人所造的幽壇的殿裡.

$^{361}$而撒瑪利亞人都有拜.
所以在撒瑪利亞的歷史.
正正就是一段嫁過五個神的歷史.
同一時間拜耶和華.
同一時間都拜當地五族的神.
是一個混雜的敬拜.
所以五個的丈夫這句話.
其實他不單只是說一個婦人的故事.
他是在回響整個撒瑪利亞民族的故事.
他就說你們找了很多依靠.
但沒有一個靠得住.
所以耶穌說這句.
他不是說要踢爆婦人.
說她生活不檢點.
而是告訴你.
我認識你.
我知道你喝過多少口井的水.
我知道你水桶裝極都不滿.
留意婦人的回應很有趣.
她說先生我看得出你是善智.
反映耶穌說中了她生命中最深處的需要.
言下之意.
耶穌說我才是賜活水聖靈的那位.
是生命的傳元.
是依靠你們以前的族人不斷找都找不到.
現在耶穌親自來到你的面前.
都是因為這個原因.
所以如果大家發現第四章.
有說到敬拜的問題.
在下文緊接著有敬拜的問題.
今天我們沒有時間說這裡.
留待下次再說.
你要明白.
其實就是因為突然間說到生命的需要.
說到敬拜的問題.
他才會說我是這座山拜.
你是那座山拜.
這個婦人拿著手上的水桶.
就是她一直為生命的傳元.
以為有很靠得住的關係.

$^{401}$但無奈五段的婚姻.
五次的盛載.
多次的失敗.
反映很多人生不如意的事.
你以為那件事順利.
我們就會找到依靠.
滿足到我們生命的空虛.
其實我們每一個人.
手上都有自己的水桶.
一些是我們以為可以滿足到我們的來源.
有人的水桶是事業.
有人的水桶是關係.
有人是子女的成就.
有人是別人怎麼看我.
這些本身都不是壞事.
但請留意.
如果我們天天拿著這個水桶.
以為那些裝回來的水.
就可以滿足我們內心深處的需要.
但我們最終會發現.
真正的滿足.
在耶穌裡面才找到.
所以第28節的動作.
成為了整段經文最美的鏡頭.
就是這個婦人放下水桶.
進入了城裡.
她放下水桶是.
你想想.
一個中午熱騰騰的時候.
特意拿著水桶來打水.
但沒有打到水.
她放下水桶.
然後她不要了.
因為她已經找到傳源了.
她已經找到真正滿足她生命需要的位置.
所以她不需要水桶.
所以來個總結.
說回剛剛開頭的三條問題.
我猜我們都應該回答得到.
耶穌為什麼要經過撒瑪利亞.

$^{441}$就是因為祂知道人最深層的需要.
是耶穌親自跨越界限.
去到井旁婦人當中尋找她.
祂說.
先生.
我看得出你是先知的.
祂是一個真正敬拜的人.
第二.
為什麼耶穌要開口借水喝.
這個問題其實是在引導人.
人以為真正的需要.
是我面前的水.
是我們面前的電子世界.
是我面前的關係.
是我的事業.
是我賺到多少錢.
水桶就是我們天天以為.
我們可以滿足到自己的倚靠.
耶穌問這條問題.
就是令我們覺醒.
真正滿足人生命的需要.
就是面前這位彌賽亞.
第三.
什麼是活水.
活水就是聖靈.
是在我們裡面擁留著的生命泉源.
水桶是我們天天以為.
可以滿足到我們的倚靠.
耶穌就要邀請我們今天.
放下這些虛假的倚靠.
面對這個主動越界.
似活水的耶穌.
我們可以怎樣回應呢.
放下我們手上的水桶.
是一些我們以為.
可以滿足到我們心靈最深處的倚靠.
這個世界不斷地說服你.
那件事很重要.
這件事很重要.
這個世界的敘事非常之強大.

$^{481}$令我們從心裡面相信.
我們休息是怎樣.
休息時間就要看電話.
看短片.
Netflix and chill.
令我們覺得.
我們心裡面最需要的就是.
電子世界帶給我們的安慰.
給我們的感官刺激.
是一些現實世界裡面.
給不到我們的東西.
但最後其實你聽一聽.
你內心深處你的身體.
告訴你.
這些東西從來都滿足不到你.
我們可以做的事.
是我們回到聖經裡面的敘事.
約翰剛剛就跟你說.
從來滿足我們心靈的東西.
就唯有是活水.
唯有是聖靈.
唯有是耶穌自己.
換言之.
這個世界都不單止跟你說.
電子世界才滿足到你.
還有其他東西.
關係 成就 期望.
別人的認同.
不要誤會我的意思.
我不是說那些東西不好.
我是說.
如果我們將人生的剛確.
將我們的人生.
完全依賴在這些東西身上的時候.
我們就會成為一個有洞的水桶.
裝多少都不夠.
飲多少都不夠.
所以這個星期.
希望大家找一個安靜的時間問問自己.
究竟我現在.

$^{521}$我的人生最渴望的是什麼.
然後呢.
我想很誠實地跟大家說一句.
沒了那句 不要緊.
就是耶穌請你把水賜給我.
耶穌請把這水賜給我.
可不可以大家跟我說一句.
我們一二三說一句.
一二三.
耶穌請把這水賜給我.
我希望這一句成為我們這個星期的祈禱.
好嗎.
第二 就是面對著活水的上帝.
我們可以走到我們經常避開的人那裡.
耶穌必須經過撒瑪利亞.
其實我們人生裡面.
有沒有一些撒瑪利亞.
有沒有一些地區.
有沒有一些人我們在避開.
可能是一些很久沒有跟祂說話的家人.
或者在教會裡曾經跟祂吵過架.
不是很合口味的一些肢體.
可能是一些你心靈裡已經認定.
這個人應該不會相信.
那些朋友.
這個月可不可以選一個你一直避開的人.
你問他主動借一口水來喝.
我的意思不是真的借一口水.
一個訊息 一頓飯.
一個問候.
從一件最小的事開始.
第三 就是由拿水的人變成分享傳言的人.
這個婦人的報道很簡單.
她沒有讀過神學.
她只說你們看吧.
這個人將我之前所行的事一切都說出來.
她得到活水她沒有自己享受.
她馬上走進城裡跟一群人說.
你看看.
誰知道整個撒瑪利亞城都相信了.

$^{561}$活水的特性就是湧出來的.
是收不住的.
當你從你心裡相信.
耶穌才是你最需要的那個位置.
你就會很自然跟別人說.
耶穌在你生命有多重要.
有多滿足.
你的故事 你的以往.
你被尋找的經歷是你最有力的見證.
我們希望在這個最撕裂的世代.
如果我們肯成為跨族群.
跨階層跨立場的教會.
我們教會就會成為一個井旁的教會.
讓那些乾渴的人在這裡感受到活水的真實.
最後思考的是.
你的水桶是什麼呢.
你願意放下它嗎.
你的撒瑪利亞在哪裡呢.
當我們思考下一步的時候.
下一步的時候我們想.
是怎樣呢.
我想我們可以從第一步開始.
就是由我們睡醒 張開眼睛那一刻開始.
我們就問耶穌.
可不可以賜活水給我.
我們問一下耶穌.
今天我們的生命最需要的是什麼.
我很希望你的答案是.
原來我最需要的是你 耶穌.
我們一起祈禱.
耶穌多謝你.
多謝你必須經過我們的撒瑪利亞.
在我們最乾渴 最避開人的時刻.
你坐在井旁等我們.
你把你的活水賜給我們.
讓我們放開我們手上的水桶.
讓你的靈成為我們生命中永留的泉源.
帶領我們 讓我們走進我們一直避開的人.
讓我們教會成為一個井旁的教會.
讓一些乾渴的人在這裡遇見你.

$^{601}$我想是祈禱奉耶穌的聖典禮物.
\newpage



\section{出埃及記}
\label{sec:rfQ3ZwOlzM4}
\textbf{會幕 (出埃及記25-40) - 陳崇基牧師}
\newline
\newline
連結: \href{https://youtube.com/watch?v=rfQ3ZwOlzM4}{\texttt{https://youtube.com/watch?v=rfQ3ZwOlzM4}} ~~~~ 語音日期: 2026-07-12
\newline
\newline
\hyperref[sec:YT4UObZ7V44]{\small{< < < PREV SERMON < < <}}
~
\hyperref[sec:index_chronic]{\small{[返順時目]}}
~
\hyperref[sec:index_scriptual]{\small{[返順卷目]}}
~
\hyperref[sec:jsiErX4TZXs]{\small{> > > NEXT SERMON > > >}}
\newline
\newline
$^{1}$各位大英姐妹 大家好.
今次我們繼續我們的系列.
今次有少許東西獎給大家.
所以如果你未有這個瓶子的話.
大家可以拿.
這個瓶子裡面已經有好幾樣物件在裡面.
我們希望在今年裡面.
這個瓶子去提醒大家.
我們在上帝裡面有什麼恩典.
今次是我們這個NEXT的第四講.
你見到那個Title是回望.
我們就看看舊約聖經裡面講到的回望.
你記得勝經是一個大的故事.
一個長篇劇有66集.
因為聖經裡面有66卷書.
由創世記至最後一卷書叫做啟示錄.
中間就是基督耶穌的來臨.
所以舊約裡面算是一部分上半的長篇劇.
有39卷書.
所以你計算它有39集.
然後新約就有27卷書.
你就計算它後面的27集.
所以我們由第一集看到最後一集的時候.
我們應該大概get到.
雖然每一集每一卷書.
講的東西是有點不同.
它們有一個主題.
是獨特的主題.
但是你串在一起的時候.
你就發現原來背後有一個編劇.
整個故事就呈現出來了.
整個大的故事.
其實就是介紹這位上帝是一個怎樣的上帝.
一個怎樣對待你的上帝.
一個跟你怎樣.
應該建立一個怎樣關係的一個上帝.
所以由頭到尾第一集至到最後一集.
都是一個故事.
一個宏大的故事.
在故事裡面告訴你.

$^{41}$上帝告訴你我是怎樣的.
等於我們彼此認識的時候.
我們大家見到新朋友來到.
我們大家彼此介紹的時候.
你都會說什麼.
由你出生開始說.
說到你現在.
我在哪裡出生.
經歷過什麼等等.
有時間的話.
因為聽了你的故事的時候.
就明白你是一個怎樣的人.
明白你的故事就是明白你.
明白上帝的故事.
都是明白和認識上帝的.
所以我們第一次開始.
我們就由《創世記》開始.
直到現在我們看到《出埃及記》的後半段.
我知道Alice傳道是在說《出埃及記》.
不過我比她快很多.
我們先用一個宏大的.
不用仔細的方法去看.
我們今天看的是25章第40章《出埃及記》.
我們看的時候.
我們首先問一個問題.
很多時候我們想像一個神.
我們想像一個很高超的.
很高的.
不能夠觸及的一個存在者.
或者那些古代的人.
他們就會想像神是一個怎樣的神.
說到神明 想起任何一種的神明.
他們就會想起一定是在高山上.
因為你要攀上去.
你要爬上去.
你才能夠跟祂接觸到.
祂是不會下來的.
這個不亙古代的人對宗教的信念是這樣.
其實現代人也是這樣的.
如果你是有一點點宗教的背景.

$^{81}$就是你未相信基督之前.
未認識耶穌之前.
你有任何的宗教的背景.
或者任何少許的概念也好.
神明總是遙遙不及的.
你是在上天的.
然後可能古代的人也好.
現代的人也好.
就要做一些大的寺廟.
或者在山上做一些建築物.
然後你希望跟祂不一樣.
跟祂接觸到.
能夠攀山你上去.
跟祂有一些接近.
無論是民間信仰.
或者以前的都是一樣.
這樣的概念.
這種神明就是很功利的.
因為是一個交換.
因為跟祂是一個交換.
就是神是遠遠超越我們的.
然後我們需要用我們的方法.
來走去祂那裡.
無論是藉著一些禮儀.
或者藉著一些不同的祭祀.
或者不同的方法.
我們是希望去到神那裡.
但是出埃及記那裡介紹的神.
往往是相反的.
因為你記得以色列人.
從埃及出來之後.
上帝用一手拯救了他們.
他們就去了一個地方.
叫做西乃山.
西乃山.
然後他們接納了上帝給他們的律法.
成為立約的子民.
然後到第25章開始.
一直說.
其實是說一些.

$^{121}$我們平時沒什麼興趣聽的東西.
就是弄一件東西出來.
弄一個所謂會幕.
中文譯做會幕.
其實就是帳幕.
就是Tent.
Tabernacle.
第25章第8節說什麼呢.
他說他們要為.
就是上帝的子民.
要為我做一個聖所.
使我住在他們中間.
我將我的家住在你們裡頭.
建造我的家在你們裡頭.
住在他們中間.
你們要按照我的指示.
帳幕和一切器具樣式.
照樣去做.
上帝就在第25章開始.
就很細節的去說.
會幕是怎麼建造的.
所以對我們來說沒什麼興趣.
我記得有一次.
那時候在神樂園的時候.
有一次我忘記去哪裡.
我們有一個retreat.
就是跟其他神樂生一起.
你知道神樂生很熟齡的.
是不是.
我們早上當然有組會.
我們就在Airbnb大屋裡.
我們有小組的靈咒.
不知道哪個人就選了這一段來讀.
就是25章.
就一讀就讀了4,5章.
如果你想知道細節的話.
你回去看一下.
我不會讀.
他一邊讀一邊讀.
就是讀會幕那些.

$^{161}$萬紙的那些curtain.
那些尺寸.
多少爪 多少爪 多少爪.
然後就做這個金燈台.
然後怎麼做 要多大 要多小.
全部都是尺寸的東西.
全部都是home depot的東西.
就是home depot就要量度.
多少尺 多少寸.
讀得很好 你很悶.
很悶.
那我心想 什麼時候讀完.
這些就是我們沒有興趣的東西.
是不是.
但讀完之後有一個弟兄就說得很對.
他說不如我們祈禱.
我說祈禱.
去完home depot怎麼會祈禱.
去完home depot看完那些間隔.
就好像看房子一樣.
量度一下房子.
kitchen counter有多大 有多長 有多寬.
然後我們不如祈禱.
他說有什麼好祈禱的.
他真的祈禱.
那個弟兄就祈禱.
他說多謝主.
因為我們看到你的居所.
你是多麼細節地安排.
他這樣祈禱的時候.
我覺得挺有道理的.
挺有道理的.
雖然對我們印度人的耳朵來說.
我們沒什麼興趣去看這些東西.
但是上帝是建造他的居所.
在他的子民中間.
住在他們中間的時候.
他是很細微的 很細節的.
去說出他是一個怎樣的神.
所以那些家具和間隔 匯木的設計.

$^{201}$其實是一些建築的神學.
architectural theology.
就是用一些architect.
或者建築的藍圖的方法.
來表達他是一個怎樣的上帝.
我記得最近有一個朋友聊天.
他說他來了芝加哥很久.
我問他你有沒有去過architectural tour.
你知道我們navy pier那裡.
你可以上船.
然後在Chicago river上去.
然後去介紹不同的建築物.
有誰去過.
有誰沒有去過.
然後我想沒有理由沒有去過.
不是吧.
我說這個是芝加哥的.
one of the best.
是不是.
你看那個tour.
或者那個tour guide.
他記載所有的建築物.
然後告訴你.
芝加哥是一個怎樣的城市.
就是用這些architect.
去tell the story.
it's amazing.
你看回望這個故事的時候.
都是一樣的.
去看看上帝怎樣去設計他的居所.
在我們中間.
以色列南部有一個地方.
你開車去開幾個小時.
就是為了看這個東西.
就是一個一比一的.
匯膜的模型.
你看到一個匯膜.
其實就是純粹.
好像你平時露營.
出去露營的時候.

$^{241}$你開的營是一樣的.
就開一個tent.
在他的神的指紋中間.
對於我們.
對聖經不是很熟悉的.
我們就會很奇怪.
為什麼上帝會這樣做呢.
原來在這些architecture裡面.
他想去表達一樣東西.
就是神學.
上帝是一個怎樣的上帝.
你看到這個匯膜的平面圖.
由右至左我們進去的時候.
右邊就是入口.
進去的時候就有一個外院.
外院就是長方形的.
然後你慢慢一路進展.
一路進去最中心的地方的時候.
就變成一個perfect square.
一個完美的方形.
在古代的近東的畫藍圖.
就是你找一個近東的結設計師.
去畫藍圖的時候.
他畫一個perfect square的時候.
就代表這個地方是最聖潔.
和最神聖的地方.
因為是perfect square.
長與闊都是一模一樣.
但是你走出來的時候.
就是一個長方形.
或者兩個正方形.
你在外院裡面看到什麼設備呢.
第一個設備你看到的就是祭壇.
還有那個銅盤.
這裡是在說祭祀的東西.
就是說你approach這個上帝.
進入去見這個上帝的時候.
你就要知道原來是有代價的.
那個代價就是上帝所提供的祭祀.
所以以色列神的子民.

$^{281}$要進去上帝這條路.
他就需要去獻祭.
你不可以空手隨隨便便來到.
你要帶著梵祭.
其他的祭物等等.
在摩西五經裡面有提及的.
就是上帝去給一條路.
給人進入去祂自己的同在裡面.
你經過這個外院.
洗了手 潔淨了.
銅盤就是水盤.
祭壇其實就是燒烤的grill.
煮了之後 燒了之後.
你就知道有代價.
to approach.
然後你就去到另一個空間.
就是一個小的長方形空間.
那個就叫做holy place.
或者聖所.
裡面有多少件傢俬.
有三件 是不是.
有三件傢俬.
第一件就是叫做陳設餅桌.
陳設餅的桌.
table of showbread.
其實那裡就有十二塊的麵包在那裡.
十二塊的餅在那裡.
就每天都要去做出來的.
就等於你今天去bakery買東西.
每天早上廚師就會做新鮮的麵包出來.
每一天在曠野裡面行走的時候.
以色列人會在他們中間.
神的會在他們中間.
祭司 利未人等等.
他們就是預備這些東西.
然後十二塊的麵包.
為什麼是十二塊.
因為十二字牌是整個神的子民.
每天就會放在那裡.
就是那些陳設餅.

$^{321}$就是去提醒我們.
原來上帝是住在整個神的子民的中間.
去供應我們一切的需要.
然後第二件傢俬就是燈台.
燈台每天都會添一些橄欖油.
然後你就要長年的去燃亮它.
就代表上帝是光.
然後我就見到一個香壇.
Incense Altar.
那裡就是點一些香.
或者那些以前有的spice等等.
他們會點了它.
然後那些香上升的時候.
就代表祈禱在上帝的面前.
上帝在聆聽我們禱告.
如果你這樣看這些傢俬的設計的時候.
其實都是很practical的東西.
其實那些香是隨湊的.
因為你有燒烤.
你有祭祀.
所以那個地方要衛生的話.
有很多東西就需要.
有銅盆,有水.
然後有香.
但是那個神學的意味.
就是說你進到這個聖所的時候.
你面對這個上帝.
祂是光.
但是祂願意去聽我們的禱告.
然後過了這個聖所.
就有一塊的簾,一塊的布.
你要揭開這塊布.
然後走到裡面.
你就會見到一件東西.
一個箱子.
就叫做約櫃.
Ark of the Covenant.
這個箱子裡面裝什麼呢?.
其中一樣就是那個法板.
上一次我們講到律法的時候.

$^{361}$就說有兩塊的法板.
一塊代表上帝這個國王.
還有另一塊代表以色列國.
他們大家是建立了這個聖約的關係.
在這個盒子上面.
就有兩隻的天使.
叫做基魯柏.
其實就是上帝的座駕.
這些全部的建築的方法,藍圖.
這些都是以前古埃及.
法老王去出征.
周圍出征的時候.
帶著軍隊,帶著他的子民.
他的帳目都在他們中間.
其實用這個背景來跟他們說.
現在上帝不再是法老王.
耶和華是你的上帝.
現在你不再屬於以前埃及.
那個欺壓你的君王.
現在是屬於一個.
跟你恩典裡面立約的上帝.
是一個愛的statement來的.
It's a statement of love.
然後每一次,每年.
只有一天叫做Yom Kippur.
Yom Kippur大概九月,十月的時候.
贖罪日的時候.
只有大祭司才能進去至聖所.
他就由外院然後結證了之後.
來到至聖所裡面.
然後代表上帝的子民去獻祭.
整個的每一個器具.
其實都是講到上帝.
去赦免我們的罪.
去供應我們一切需要.
去聆聽我們.
還有在曠野裡面去引導我們.
是跟我們同在的.
這個上帝不是一個高攀不能的上帝.
並且他是下來我們中間.

$^{401}$住在我們中間.
將他的家放在我們中間的一個上帝.
所以二十九章這裡講什麼呢?.
他說我耶和華要住在以色列人中間.
作他們的上帝.
他們便知道我是耶和華.
我他們的上帝.
是將他們從埃及領出來的.
說要住在他們中間.
你有沒有留意這裡講什麼?.
經常講到住.
我的家在你中間.
所以出埃及的時候.
以色列人所經歷的.
不是純粹說我們得到自由了.
離開苦難.
而是一個進入關係的過程.
在西太山的山腳.
他們見到一個回眸.
去到建立的時候.
一個藍圖的時候.
你可以想像摩西在西太山.
拿到這個藍圖.
好像現在那些建築師.
拿到一些很大張的印出來的紙.
一看這個平面圖.
一看裡面所有的設計.
所有的尺寸.
所有的細節的時候.
他們知道原來這位上帝.
在恩典裡面來愛我們.
那是第一部分.
第二部分.
在《出埃及記》裡面.
記載的回眸的故事.
就出奇地中間那裡.
就有一個三十二至三十四章.
說到另一個故事.
就好像作者沒什麼心思.
突然間說到很多藍圖的東西.

$^{441}$二十五至三十一章.
說到回眸的設計.
尺寸.
Home Depot.
那些藍圖的東西.
然後突然間三十二章.
就說另一個故事.
就是以色列人在西太山腳那裡.
拜金牛的事.
這個簡直是反高潮.
因為開頭說到一個Statement of Love.
你想像一下你的愛人很愛你.
他買了一個豪宅給你.
然後就你進去住.
他先把那些相片和藍圖給你看.
全部都是合你心意的.
正不正.
很正.
那下一步應該是什麼.
下一步應該是入住.
或者開工整他.
誰知道這個聖經的作者.
在《出埃及記》就插了三章進去.
就說到他們拜金牛.
這個故事我們不詳談.
但是就說到以色列人.
等摩西下山拿那個藍圖的時候.
他們等不及了.
他們就叫阿倫.
不如我們現在弄一隻金牛出來.
因為我們不想要這個耶和華上帝.
我們等不及了.
我們想拜這隻金牛.
這個金牛當然在當時的時候.
就是一個物件的偶像.
所以耶和華就很生氣.
真的很生氣.
但是其實在今天的故事來說.
這隻金牛其實都代表我們很多不同的.
所謂金牛的.

$^{481}$就是我們會崇尚物質.
我們一切都以物質.
作為我們決定的法則.
或者我們可能一些的情慾.
一些引誘我們的東西.
我們每個人心裡面都有金牛的.
這個作者特意在一個會幕的藍圖之後.
插這個故事進去.
其實想說什麼呢?.
想說什麼呢?.
因為在32-34章說到拜金牛的時候.
耶和華是應該以公義去殲滅他的子民.
因為想想他用他的恩手拯救了你們出來.
使你們脫離法老的苦難.
然後你來到這裡的時候.
我給你一個表達愛的聲明.
是不是我說我愛你.
我愛你們.
但是居然你去怎樣?.
你背若離經背道.
所以在這個時候.
其實是整個《出埃及記》裡面.
一個很大的危機.
是一個危機.
因為如果是這樣的話.
整個上帝所做的救贖計劃.
就會遊戲結束了.
會不會遊戲結束?.
你看32-34章.
裡面摩西就做了一個很重要的角色.
所謂一個中保的角色.
他就在上帝面前.
去求上帝去赦免他們.
因為上帝說不如我將這些子民全部殲滅.
然後在摩西你身上.
我去興起一個新的子民.
好不好?.
那你摩西就威風了.
摩西說不要不要.
因為如果是這樣的話.

$^{521}$萬國會怎樣看我們?.
就會覺得我們很廢.
特別是耶和華很廢.
怎麼可能會這樣發生?.
所以他就求上帝用憐憫去恩典.
去赦免他的子民.
結果就出了一句.
舊約裡面最重要的一句statement.
就是認識這位耶和華神是怎樣的神.
這個statement有34章6-7節.
耶和華是大有憐憫大有恩典的神.
不輕易發怒.
且有豐盛的慈愛和誠實.
為千萬人存留.
不是審判.
是慈愛.
赦免罪孽.
過犯和罪惡.
在舊約裡面.
這個是最大的statement.
因為揭示這位上帝不單是遙遙不可及.
他是住在我們中間.
他不是一個遠遠不能及的上帝.
他住在我們中間.
並且這個上帝願意以恩典.
以慈愛不輕易發怒來對待我們.
所以他寧願決意的去赦免他的子民.
住在他們中間.
然後故事繼續發展.
上帝就說.
金牛的事件之後.
你們得到赦免.
你們為我建造我的家.
然後才開工做這個會幕.
在宏觀的時候就發現.
其實作者想用一些很細節的東西.
用當時人明白的一些語言.
來介紹他自己是一個怎樣的神.
35-40章就是摩西下山之後.
他就召集這些百姓.

$^{561}$然後給他們一個藍圖.
你看看上帝就是這樣愛我們.
上帝願意下來我們中間.
住在我們這裡.
百姓居然有這樣的邀請.
雖然是有罪.
剛才那首詩歌有一句歌詞很觸動我.
千瘡百孔.
我們都是一班千瘡百孔的人.
但居然我們得到上帝的赦免之後.
我們得罪上帝 但上帝赦免我們.
他居然還說什麼.
你照我的藍圖.
建造我的居所在你們中間.
所以百姓就動手來建造這個會望.
那幾天又是很長的35-40章.
但那裡和之前25-30章完全不同.
25-30章就是你要做一樣.
你要做一樣.
你要根據藍圖做一樣.
35-40章就是他們做了什麼.
就是照樣下命令.
然後他們執行那個命令.
他們就照樣照著上帝的藍圖和樣式.
來做每一樣東西.
舉例25章31節.
這裡是說藍圖的第一部分左邊.
耶和華吩咐摩西和他們的子民.
你們要做一個純金造出來的燈台.
上面要有灶 有杯等等.
然後放那些油燈.
這個是七燈台的建造.
那個藍圖.
你要這樣做 照樣這樣做.
然後那個金牛的事件32-34章.
到了會望工程開剪整.
幾乎是同一句話.
只不過之前是命令式.
現在是敘述式.
37章第17節.

$^{601}$他就做了 那裡有一個工匠.
他就用純金造出燈台.
然後有exactly所有的東西.
燈台連接成一塊.
金燈台不是薄銷行出來的.
整塊真金這樣突出來.
你見到開頭是命令.
告訴你藍圖是這樣.
現在他們就實際的來執行出來.
所以25章至31章.
還有35至40章.
其實好像照鏡一樣.
後面的exactly是之前藍圖所說的.
幾乎是一模一樣.
上帝說了 事就這樣成了.
上帝這樣吩咐他們.
就是這樣去建造.
你說你這個整個宏觀.
你就見到上帝是一個滿有恩典的上帝.
並且他願意去邀請我們.
為他建立居所.
這個statement是一個最大的statement.
教會 初時以色列人.
都算是初期的教會.
我們教會的前身.
舊約的子民和新約的子民.
同一個的子民.
我們居然有這樣的救贖.
和可以參與建立上帝的居所.
教會都是千瘡百孔的.
教會都是很多錯誤的地方.
很多時候我們有一種錯誤的觀念.
你聽見很多人都這樣說.
或者我們有時候都會這樣想.
教會都是我整理得好一點.
我才來.
我不來教會 為什麼.
因為我覺得我太多污穢的東西.
等我先搞定 我才配來教會.
配來上帝的面前.

$^{641}$但很多未信主的人都是這樣想的.
甚至很多信了主很久的人都是這樣.
所以我們就有個殼.
我們想搞定那個殼.
起碼我們表面上.
我們覺得我們都是很光鮮.
我們將那些東西隱藏.
但其實我心想.
上帝如果光照我們任何一個角落的時候.
我們個個都是千瘡百孔.
是污穢的.
我們個個都是不配的.
有誰可以配來到上帝面前.
沒有一個人.
唯獨是上帝在恩典裡面邀請我們.
我們才能夠參與.
正因為這個緣故.
我們的參與就更加有價值.
因為我們經歷過上帝的恩典之後.
我們就明白.
這個上帝這麼愛我.
我豈能不參與.
我豈能仍然坐在這裡.
所以這裡會幕的建造裡面.
就有一些東西.
講到那些以色列民.
他們甘心奉獻.
奉獻到一個地步.
多到一個地步.
摩西說夠了我們.
夠了.
太多了.
一個金燈台要多少金去做.
大概好像這個十字架這麼高.
OK.
你當它多少斤都好.
他們捐獻的.
是多到一個地步.
我可以做幾十個金燈台.
摩西說夠了夠了.

$^{681}$因為已經足夠了.
就因為他們被上帝的恩典觸動過.
他被上帝的生命觸動過.
當你唱那些詩歌.
你想起那些經文的時候.
我是一個千瘡百孔的人.
我是一個污穢不得的人.
但上帝居然在高天來到我中間.
住在我的生命裡面.
住在他的子民中間.
這個上帝是值得我們讚美的.
難怪.
去到最後的時候.
上帝的榮耀就充滿了這個會望.
四十章就是《出雅及記》最後一章.
最後那幾個字.
他說哪時雲彩就遮蓋著會望.
耶和華的榮光就充滿著它的像模.
使得摩西都不能夠進去.
因為雲彩停在旗上.
耶和華的榮光充滿著這個像模.
沒有人有資格.
讓上帝住在他中間.
千瘡百孔.
但上帝他在恩典裡面邀請我們.
甚至去建造他自己的會望.
我知道我們教會.
我們都來到一個時間.
十年的故事.
我們都去到一個擴展的時候.
但請你知道.
當我們呼籲我們一起建立教會的時候.
請你知道.
我們是一班怎樣的人.
我們是一班不配的人.
既因為我們不配.
上帝的恩典卻來到我的生命裡面.
很真實的來到我的生命裡面.
所以我豈能夠只是坐在這裡.
不去參與.

$^{721}$我們領受過上帝的恩典.
領受過上帝愛的人.
我們被祂觸動過的人.
我們更加需要去建立這個群體.
讓上帝住在我們中間.
這是多大的權利.
多大的權柄.
多大的福分.
是我們不配的.
Yet, it comes to you.
來到我們生命裡面.
教會從來不是為完美的人而設.
但卻是為一班千瘡百孔的人.
污穢不堪的人.
而經歷過上帝基督的恩典.
我們就來到去組成.
一個上帝能夠住到的地方.
整個脈絡裡面.
你會發現上帝是越來越靠近人.
靠近我們這一班不配的.
難怪到了新約的時候.
約翰福音第一章.
約翰去介紹耶穌的時候.
耶穌來到我們中間.
約翰福音第一章.
由頭講到尾.
到第十四節.
他說道就是耶穌自己.
成了肉身住在我們中間.
充足滿滿有恩典有真理.
我們也見過他的榮光.
正是父,獨,生,子的榮光.
那個住在我們中間.
Dwelling.
Made his dwelling among us.
住在我們中間.
這個字其實就是回眸.
Tabernacle.
所以你想直譯的話.
He made his tabernacle among us.

$^{761}$Or he tabernacled among us.
他已住在我們這些千瘡百孔的人的生命當中.
住在我們裡面.
這個真是上帝的榮光.
真是上帝的榮光.
上帝的榮光不是什麼神蹟騎士.
很spectacular 放煙花的人.
很實在的是上帝說.
雖然你不配.
雖然我不配.
但我願意去救贖你.
住在你的心裡面.
住在你的生命當中.
你會發現世界.
一個誤信的世界.
需要的是什麼.
不一定是需要一個.
大家很光鮮,很完美的教會.
世界需要見到的教會.
就是我們這班.
我們可能不配的人.
但上帝的恩典.
在我生命裡面成為真實.
我親自經歷過.
被觸動過.
上帝在耶穌基督裡面給我們愛.
我們說這個journey裡面.
是整個聖經的大故事.
去認識這位主.
今天我們邀請你.
出去的時候.
你拿一塊布條.
為什麼是布條.
因為今天說到回眸.
回眸很多不同顏色的布條.
你都可以造成的.
就是上帝以恩典裡面.
去邀請祂的百姓,祂的子民.
一起去建造祂的居所.
你拿一條回去.

$^{801}$然後放在我們所謂.
十週年的瓶子那裡.
我們有幾行的物件已經有了.
海沙這麼多的.
就是亞伯拉罕的祝福.
然後還有什麼.
有一塊碎片.
我們說到族長故事的時候.
個個都是破碎的人.
但是上帝去使用他們.
叫萬國能夠得福.
然後我們有一些石頭.
說到上帝的律法.
因為有一個新的身分.
我們作為上帝的子民.
我們就按照祂的律法,祂的原則.
來生活.
今天就是布條.
讓上帝住在我們中間.
我們都有份.
有這樣的恩典.
可以參與建立上帝同在的地方.
教會要有上帝的同在.
就需要我們以心,以力擺上.
最後的時候我想邀請你.
這個星期.
可以為教會祈禱.
不要只聽說教會需要擴展.
需要多走一步.
需要go next.
How about you go next.
我們英文有句說話.
out of sight out of mind.
看不到我們就忘記了.
但是基督徒呢.
out of prayer out of mind.
我們不去禱告,不去紀念的時候.
我們就當自己是旁觀者.
但是今天的故事告訴我們.
原來上帝邀請你每一個.

$^{841}$去建立他自己的居所.
這個星期做一個功課.
拿一條布條.
然後我們每一天為了教會去禱告.
為了教會我們不單只需要地方.
為了我們教會的生命我們去禱告.
我們在信,忘,愛當中.
我們都要自己成長.
You have to go next.
And I have to go next.
上帝榮耀在我們中間.
\newpage



\section{出埃及記}
\label{sec:jsiErX4TZXs}
\textbf{生命的定位 (出埃及記 6) - 江芷敏傳道}
\newline
\newline
連結: \href{https://youtube.com/watch?v=jsiErX4TZXs}{\texttt{https://youtube.com/watch?v=jsiErX4TZXs}} ~~~~ 語音日期: 2026-07-12
\newline
\newline
\hyperref[sec:rfQ3ZwOlzM4]{\small{< < < PREV SERMON < < <}}
~
\hyperref[sec:index_chronic]{\small{[返順時目]}}
~
\hyperref[sec:index_scriptual]{\small{[返順卷目]}}
~
\hyperref[sec:DRGVYhPEXTM]{\small{> > > NEXT SERMON > > >}}
\newline
\newline
$^{1}$Hello 大家早安.
早安.
可以再上來跟大家分享上帝的話語.
我們一起看看這個.
大家知不知道是甚麼.
這個是White Elephant.
最吸引和最多人搶的禮物.
這個就是一粒AirTag.
它可以去定位你想找到的東西.
如果你不見了錢包,鎖匙甚至小朋友.
你都可以放一粒AirTag在他身上去找他.
你可以靠他知道你想找的東西在哪.
今天我們除了說物件需要定位.
其實最重要是我們的人生能夠找到我們的定位.
今天我們會去看出埃及的十災.
不知道大家有沒有看過.
大部分人都有看過.
十災其實是一段頗長的經文.
今天我們可能不會仔細地看.
但是我們會希望找一些作者重複的訊息.
或者他很強調的經文.
我們一起去看看.
我們先讀出埃及的六章一至八節.
有沒有些小的.
我們試試看 看到的讀大聲一點.
預備起.
耶和華對摩西說.
現在你必看見我向法老所行的事.
使他因我大量的手放以色列人走.
因我大量的手把他們趕出他的地.
神吩咐摩西對他說.
我是耶和華.
我從前向亞伯拉罕以撒雅各.
顯現為全能的神.
至於我的名耶和華.
我未曾向他們知道.
我要與他們建立我的約.
要把迦南地他們寄居的地賜給他們.
我聽見以色列人被埃及人奴役的哀聲.
我就紀念我的約.

$^{41}$所以你要對以色列人說.
我是耶和華.
我要除去埃及人家給你們的奴役.
救你們脫離他們的奴役.
我要用新出來的棒臂.
藉嚴厲的懲罰救贖你們.
我要以你們為我的百姓.
我也要作你們的神.
我除去埃及人家給你們的奴役.
你們就知道我是耶和華你們的神.
我起誓應許給亞伯拉罕以撒雅各的地.
我要領你們進去.
將那地賜給你們為業.
我是耶和華.
你見到我highlight了這麼多次這一句.
就是今次的重點.
就是神在這一段經文裡面.
說了很多次什麼.
我是耶和華.
這個在短短的八節裡面.
都出現了四次.
而且在我們要說十災的過程.
即是六至十四章裡面.
出現了十六次.
這個不單止是神去介紹自己.
說我是誰.
而是一個文藥性的宣告.
為什麼這樣說呢.
因為神在這裡清楚告訴我們.
他是因為紀念和我們的約.
所以要帶我們出埃及.
他聽到我們的哀聲.
聽到埃及人去奴役以色列人的哀聲.
之前我們講道也有提過.
就是上帝是一個紀念他的約.
親近我們.
聽到我們的哀聲.
看到我們的苦的一個上帝.
所以他紀念我們的約.
而且他要帶領我們出埃及.

$^{81}$他說他要做我們的上帝.
他就是那個耶和華.
這個是一個文藥性的宣告.
這個也是以色列人在埃及裡面.
做了四百年的奴隸.
可能他們自己都已經忘記了自己的身份.
或者他們日常的工作.
已經好像讓他們奴役到.
都沒有機會想一下.
究竟我自己是哪一個呢.
近年在美國有很多報道.
或者一些健康的調查.
都講到青少年有一個身份的焦慮.
可能他們會不知道自己是什麼人.
或者不知道自己的性別.
或者他們自己對自己將來感到很迷茫.
有時我們在一個長期迷失的狀況裡.
青少年的確不知道自己是誰.
因為這個世界用了很多東西去標籤我們.
就好像我們的成績.
我們的收入.
我們的身份.
甚至我們的社交媒體有多少人關注.
在很多青少年裡面覺得這些東西都很重要.
不過這個是不是他人生的定位呢.
上帝說我是耶和華.
你是我的百姓.
這個身份搞清楚之後.
我們就知道自己為了什麼而活.
或者我們今天坐在這裡為了什麼.
有些人亦都會說.
他說這個出埃及的十災.
應該是在對應埃及當時不同的神明.
就要證明神能勝過那些神明.
我們都可以看一下.
我覺得這個表挺有趣.
有九個的神.
九個的神都是代表不同的東西.
就好像十災裡面.
神師子不同的法術都是對付不同的埃及神明.

$^{121}$我們可以由左邊開始看.
利諾河神就是第一災.
上帝要將水變成血.
為什麼呢.
你可以用這個方法去理解.
或者之後你再找多一些資料.
埃及人以為利諾河就是他們生命的泉源.
所以他們覺得他們一切都是來自利諾河.
但是誰創造這個利諾河.
我們有看聖經就知道是上帝創造的.
而且再看一下第二災.
即是蛙災.
上帝能夠叫這些青蛙出來.
其實當時的埃及的法師都能夠叫這些青蛙出來.
其實我覺得十災很好看.
因為原來那些法師都能夠叫青蛙出來.
好像在彰顯他和神有一樣的能力.
但是他能不能叫這些青蛙走呢.
他不能夠.
所以法老最後要求摩西.
你叫你的神請走這些青蛙.
或者之後的災害都是一樣.
他們以為自己的能力能夠和神媲美.
但是卻不能夠.
聖經作者沒有將十災寫成為一個神明大鬥法.
為什麼呢.
我們經常說看聖經我們要問為什麼.
如果你不知道想了解多一點.
怎樣看敘事文的話.
你都可以上Mark Sir現在開的敘事文課程.
其實真的要問.
如果是十災寫成神大鬥法.
不同的神明就可以了.
為什麼要是我們現在看到的版本呢.
剛才的六章已經說了一個十災的預告.
他就要讓我們知道什麼.
他要讓我們知道身份.
在十災裡面不單止說神能夠有能力對抗這些神明.
而且在這些大鬥法裡面.
要知道誰才是上帝.

$^{161}$耶和華究竟是一個怎樣的神.
神不單止勝過這些神明.
祂更加是這個世界的創造者.
祂更加是一個守施慈愛的上帝.
「我是耶和華」這一句.
不單止是說我要帶你們出埃及.
出埃及不單止是一個逃生搜查.
而是說你人生的定位究竟在哪裡.
身份一旦被重建改變.
你的生命方向也可能會改變.
就好像你做了父母之後.
你就不一樣了.
因為你所想的東西.
或者你的方向都是為了你的小朋友去想.
你的身份可能會定義了你的方向.
所以神在這裡跟我們說.
祂的身份.
祂就是那個耶和華.
我們是誰.
我們是祂的子民.
在預備的時候也想起.
我和Ken剛剛過了我們的結婚週年.
回想六年前.
大家都知道2026年就是COVID.
其實當時我們預訂了香港的禮堂.
可以坐三,四百人.
我們也預備邀請不同的朋友,親人來到.
但是2020年的年頭.
大家知道疫情來勢洶洶.
我們也要每天留意新聞.
它說聚集人數每況愈下.
開頭說100人,50人,20人.
甚至最後是5個人參與婚禮.
也說有一些社交的距離.
在我們婚禮之前.
我們才看到新聞說.
婚禮可以有50人.
50人.
但是其實我家人也很多.
如果不算朋友.

$^{201}$我也很難取捨.
我婆婆生了10個子女.
我奶奶也生了7個.
你算一下有多少.
算上伴侶也有30多個.
再加上兄弟姐妹.
當時的教會也很好.
給了我們很多彈性.
弟兄姊妹也說我們幫你處理好.
我們出來.
讓你的親朋戚友去婚禮.
而且在婚禮裡.
因為很多的變動.
變動到一個地步.
我也很不想舉辦婚禮.
我就跟Ken說.
不如算了.
不要舉辦婚禮.
Ken當時想了一下.
他就跟我說.
我們不變就可以了.
也起雞皮疙瘩.
我應該穿回來.
起疙瘩.
他說婚禮其實就是.
你和你愛的人.
一起在上帝面前立約.
這是婚禮的意義.
其實其他人也很重要.
不過不是最重要的.
也不是最重要的.
到了最後.
其實你知道你和你愛的人.
可以走進來這裡.
和上帝結一個盟約.
這才是婚禮最重要的東西.
所以我說你的定位是什麼.
是很重要的.
影響著你怎樣做人.
第二.

$^{241}$我們待會也會去看一看.
這個十災的影片.
讓我們回顧一下.
這個十災究竟是怎樣發生的.
我想大家在這段短片裡.
可以留意一下.
法老究竟是一個怎樣的人.
或者他的性格是怎樣.
可以請去片.
法老摩西和阿倫.
就照著上帝的話去.
法老的魔術師.
用邪術做了同一件事.
所以法老還是不肯讓上帝的子民走.
於是上帝再向埃及.
發出九個災害.
去顯示他的能力.
好.
即使受到這麼多的傷害.
法老還是心很硬.
不肯讓這些熄滅人走.
在最後一個災害的夜晚.
法老醒了.
聽到很大的哭聲.
原來全埃及人的家都有人死了.
法老叫摩西和阿倫來.
叫他們帶以色列人走.
於是以色列人立刻就走.
向著葉爾許之地出發.
還帶了很多財寶和弱勢的骸骨一起走.
因為是很久之前應承的.
誰知他們走了之後.
法老又改變主意.
派軍隊去追以色列人回來.
當以色列人見到法老和軍隊就很害怕.
上帝按照他之前應承過以色列人的.
用方法救了他們.
經歷了這麼多.
以色列人終於看到上帝的大能.
是唯一的真神.

$^{281}$他們又信任上帝的僕人摩西.
這個動畫也挺好看的.
不過聖經還很低.
很多細節大家可以去發現.
大家有沒有看到這段影片是怎樣形容法老的.
他說他的心很硬.
究竟是誰令法老心硬的呢.
我們很多時候看和合本的聖經.
都說是耶和華使他心硬.
其實這個不單是形容法老的性格.
還在揭示他自己的一些屬靈的境況.
他的心硬是從哪裡來的呢.
其實我們知道法老本身是萬人之上.
全部人都要聽他.
而且在埃及的文化裡.
他覺得法老其實就是那個神.
剛才你不是看到十災對抗的埃及神明.
其實最後一個沒有畫出來的.
就是埃及法老.
他覺得埃及法老就是地上的代表.
他就是這個神明.
所以法老覺得自己勝過地上的一切.
耶和華都不是甚麼一回事.
所以他自己是一個很剛硬的心.
他不能夠讓以色列人離開埃及.
因為他就沒有那些被奴役的奴隸了.
但其實我們經常都說.
耶和華使法老心硬.
我們有一個更好的翻譯.
是這樣寫的.
就是神容許法老繼續選擇他自己的心硬.
這個選擇是很重要的.
法老一次又一次見到神的作為.
但他仍然選擇去拒絕耶和華.
其實有很多時候我們去選擇很多東西.
我們這個世代都有很多不同的新聞可以選取.
有很多人說現在很多錯誤的資訊.
但我們仍然選擇相信自己想相信的東西.
其實法老亦都是一樣.
他去選擇他自己.

$^{321}$選擇去操控這一切.
我們可以在希伯來文.
即是一些原文的解釋裡面.
知道聖經十宗裡面所說的心硬.
其實都有三個不同的說法.
第一個就是法老在整件事裡面的心理狀況.
因為在七章第一章都未開始.
就是法老心硬不肯放百姓走.
這個是耶和華對他的描述.
是一個心理的狀態.
第二個亦都可以說很沉重,很遲鈍,很麻木的意思.
就是描寫法老對神的話越來越沒有反應.
就好像八章十五或者三十一節那裡這樣說.
即是法老見到災害過了之後.
但他仍然硬著這個心.
他好像麻木了.
對上帝的作為麻木.
他都再一次去ignore這個上帝.
第三個亦都可以理解是加強,堅固的意思.
就是上帝去加強或者容許法老的心硬.
就是剛才我們所說的這個譯本.
其實有三個不同的意思.
但不等於我們理解了這個意思.
我們就知道這個是甚麼意思.
我們要理解整個的處境.
神在這個十災裡面.
是一次又一次容許法老去自己做這個抉擇.
這個字是讓我們體會到.
神容許法老去自己做一個抉擇.
這個心硬不是上帝令他由一個很善良的人.
變成一個很麻木,不仁的法老.
而是法老本身他就是對上帝麻木.
或者他不當上帝是一回事.
上帝容許他.
直到他知道或者看到他自己是怎樣的人.
神都是讓我們慢慢去發現.
他就是一個這樣的上帝.
去到第七災的時候.
法老都知道他自己是有罪的.
這個是《出埃及九章》.

$^{361}$他說這一次我犯了罪了.
耶和華是公義的.
我和我的百姓即埃及人都是邪惡.
其實他慢慢會自己發現.
他由一開始完全不理會上帝.
到他知道他自己犯罪.
再到第十章的時候.
他再一次都說自己得罪了耶和華.
又得罪了摩西亞倫他們.
他就求饒恕.
當然我們知道他是因為不想埃及人再受苦.
所以他和神去求饒恕.
但這個十災的描寫過程.
不只是神明大鬥法.
而是在告訴我們.
耶和華是容許或者給了一個忍耐.
法老一次又一次去做選擇.
耶和華任憑法老的心眼.
是一個很寶貴的自由.
當我們去愛我們身邊的人.
其實不是硬要對方去接受我們所覺得的好處.
或者好的東西.
而是讓對方選擇.
上帝完全有能力去發出一個號令.
令法老去跟著他所說的命令去做.
但他沒有.
他是一步又一步.
一災又一災.
很大費周章地.
讓他自己在當中去發言.
十災不只是描述耶和華沒有能力.
更是賜給我們一個很自由的意志.
或者一個空間.
讓我們能夠選擇.
第三.
在十災裡面.
經常都出現一句說話.
很重複的.
就是耶和華神跟法老要求.
放走我的百姓.

$^{401}$好讓他們能夠事奉我.
在十災裡面都提到六次.
是非常之關鍵.
他不單止表示.
以色列人要離開埃及.
更加是表示.
上帝要取回這班以色列人生命的主權.
在出埃及的時候.
我們經常都覺得好像這班以色列人搬遷.
搬離開埃及.
這就是出埃及記.
其實不單止是這樣.
更新層次的是.
一個屬靈和盟約上的轉移.
我們看到這班以色列人.
為法老做事的時候.
他好像為帝國奴役.
生命被壓榨.
被消耗.
但神卻要讓他們的生命.
不再被壓榨.
不再被奴役.
他們要從埃及出來的時候.
他們是由被迫工作.
神希望他們能夠成為一個.
甘心敬拜他的群體.
他們希望這班以色列人.
能夠由被奴役.
轉到事奉上帝.
被掌控的生命.
能夠有那份自由.
他能夠選擇.
以至神能夠差遣他.
這個很不一樣.
可能你會說.
其實隸屬哪一個.
無論是耶和華還是法老.
都是被奴役.
但其實不是.
剛才那一點已經解釋了.

$^{441}$上帝是一個耶和華.
不單止是滿有能力.
而且是讓我們選擇的神.
我們經常都以為.
生硬的只有法老.
其實不是的.
在這裡跟我們說.
這班叛逆的人.
不單止是埃及人或者法老.
其實還有以色列人.
以色列人當遇到.
埃及兵追趕的時候.
他們在紅海前面.
剛才那段影片沒有說.
但聖經有說.
在紅海面前.
這班以色列人又害怕起來了.
我們可以一起看看第十節.
法老逼近的時候.
以色列人舉目.
看埃及人追來了.
就非常懼怕.
以色列人向耶和華哀求.
他們跟摩西說.
難道埃及沒有墳地嗎.
為什麼要帶我們出來死.
你為什麼這樣對我們.
為什麼要將我們從埃及領出來.
我們在埃及都已經跟你說過.
不要搞我們.
讓我們事奉埃及人就算了.
總比死在曠野好.
這班以色列人.
其實他看了很多.
十災上帝在他們身上做的事.
但他仍然好像忘記了.
或者看不到一樣.
在《埃及記》裡面.
其實有很多次.
耶和華是專門分別以色列人.

$^{481}$令他們不要受災害.
就好像在九章裡說.
在辛蓄的災害裡面.
耶和華是要分別以色列和埃及的辛蓄.
讓以色列人一隻都不死.
你想想如果上帝保護你.
很明顯周圍的牛羊都死了.
就是以色列人家裡的牛羊沒有死.
而且在十災裡面.
所有的長子.
埃及地的長子都死了.
但維及以色列人的長子就沒有死.
而且在冰雹裡面.
以色列人過山地都沒有冰雹.
是不是很奇怪.
為甚麼以色列人還看不到.
他們其實和埃及法老是一樣的心眼.
當我們看不到的時候.
我們是不知道.
但上帝今天就要和這些以色列人說.
我就是那個耶和華.
我要帶領你們出來埃及.
其實上帝不單要讓這些人離開埃及.
上帝不單在出埃及記裡說要帶人脫離痛苦.
更加要給人機會看到自己的錯誤.
讓人重新學會敬拜.
回到一個敬拜的群體.
在過程裡面.
上帝就好像給這些以色列人.
容許他們有埋怨的空間.
容許他們去做錯.
但上帝給了我們一個方向.
我們這班敬拜祂的人.
不是不能夠錯.
我們也可以.
但祂要讓我們知道.
這個自由究竟是一個怎樣的自由.
祂要我們成為一個敬拜的群體.
而是我們生命有方向.
能夠依靠祂.

$^{521}$就好像剛才所說.
我們要知道我們自己的身份.
而且我們要知道我們之後要怎樣做.
來到最後.
其實我仍然有些餘悸.
今天是5月31日.
其實上年的6月1日.
我第一次在這裡講道.
Ken第一天在這裡服侍.
其實我以為慢慢開始解決.
誰知很不幸地.
我們當日的下午.
我們就在Highway 55那裡.
被人Hit and Run.
而且進了ER.
我記得當時Doris,Ronald,Mark Sir,Amanda.
和Mark Sir的兒子Nathan.
馬上來車禍現場找我們.
我們當日就體會到.
甚麼叫Mark Sir的車快過白車.
我撞完之後.
其實我有五分鐘.
我猜我都說不出話來.
Ken在旁邊一直問我.
「你有沒有事?你有沒有事?」.
但我完全說不出話來.
好像這裡有一股很大的力量.
去勒住了我.
而我整個人都在顫抖.
當時Doris給了我很多衣服.
我還以為自己很冷.
顫抖了一整晚之後.
我發現其實應該是驚恐.
是撞車之後的驚恐.
令到我這樣.
撞完車之後.
很多弟兄姊妹去幫我們.
和林廣安幫我們做物理治療.
但我們其實在這件事之後.
都覺得在這裡生活好像很不安全.

$^{561}$其實大家可能對這裡生活安不安全.
有不同的定義.
這是我的感受.
就是覺得在這裡生活好像很不安全.
其實其他都很好.
吃東西很好.
事奉很開心.
但心裡始終很害怕.
有一份驚恐.
再過幾個月.
我們又要開始想身分的問題.
所以覺得很不容易.
但有一天我和YA分享.
在這些很不安很不穩的裡面.
我開始想其實我自己在這裡為了什麼.
上帝好像透過一段經文都和我說.
祂是耶和華.
我就是跟隨祂的子民.
祂怎樣去保護以色列人.
都一樣保護我們.
就好像零零碎碎.
周圍都有災害去保護我們.
其實我再想這個車禍可以夾死我們.
因為後面有一輛車撞到.
然後它駛到路肩.
所以才卸了一些力.
如果不是我們可能夾住前面那輛車都變了.
所以我開始去想我自己在這裡為了什麼.
想上帝是一個怎樣的神的時候.
我就開始想我們經常都說outreach.
我是否真的要穩定下來.
我才去事奉神呢.
神好像在這段經文裡面開始解答.
其實不需要等什麼身份.
不需要等你怎樣settle下來.
就好像當初我信主一樣.
都是在一個很混亂的狀況之下.
就是在那個混亂的狀況之下.
我們就需要神.
在一個很混亂的狀況之下.

$^{601}$我們就需要這個創造世界.
而且讓我們生命秩序的神在我們當中.
來到一個總結.
今天我們說出埃及記.
不單止一個求生手冊.
是一個重新讓我們身份定位的故事.
上帝就是說了這一句.
祂是耶和華.
我們是祂的子民.
祂是耶和華.
祂是滿有能力而且滿有慈愛.
守約的一個上帝.
我們是祂的子民.
不知道今天你生命的定位是怎樣呢.
你生命的定位是怎樣呢.
可能你在剛才我所說的混亂當中.
你願不願意和上帝恢復一個很有秩序的關係.
承認祂就是你那個耶和華神.
而你就是祂的子民呢.
而你亦願不願意好像上帝呼召以色列人一樣.
去活出一個敬拜的人生.
或者你會不會看到.
你心裡面都好像法老,埃及人一樣.
有不願意放手的地方呢.
我希望今天你看到的不是上帝怎樣說你.
而是看到上帝這一份愛一直在等著你.
當我們說教會要延伸的時候.
我們的生命其實都需要延伸的.
你是在活一個怎樣的敬拜人生呢.
有時可能我們都說.
等某些東西準備好我們才去做.
不是啊 上帝今天就跟我們說.
不是等我們所有東西都準備好.
就在一個曠野裡面.
祂呼召我們能夠活出一個敬拜祂的人生.
可能讓我們生命裡面有很多不想被上帝觸碰的.
我們求上帝去幫助.
更加讓我們的戒善能夠延伸.
好像剛才Mark Sir在祈禱裡面說.
我們能夠打開多一點自己.

$^{641}$讓人能夠在我們生命裡面拿到一些東西.
不單止是我們的知識.
不單止是分享我們的一些關係.
是我們能夠讓他看到耶穌在我們生命裡面做些甚麼.
讓我們一起去祈禱.
親愛的天父 多謝你一路帶領我們.
讓我們教會能夠一路走下去.
主啊 我們求讓你能夠帶領我們.
看到更多自己需要你的地方.
而且在我們生命裡面更多被你觸碰.
我們都不知道怎樣可以活出一個.
你喜悅敬拜的人生.
求你在我們餘下的時間.
分散的腳步.
讓我們能夠一路再想起.
色靈啊 求你親自帶領我們.
多謝你的愛 多謝上帝.
我們獻上禱告交託奉主名求 阿們.
上帝祝福大家.
\newpage



\section{以弗所書 3:14-21}
\label{sec:DRGVYhPEXTM}
\textbf{被基督的愛充滿 (以弗所書3\_14-21) - 張威傳道 Filled with the Love of Christ (Eph3\_14-21) - Paul Zhang}
\newline
\newline
連結: \href{https://youtube.com/watch?v=DRGVYhPEXTM}{\texttt{https://youtube.com/watch?v=DRGVYhPEXTM}} ~~~~ 語音日期: 2026-07-12
\newline
\newline
\hyperref[sec:jsiErX4TZXs]{\small{< < < PREV SERMON < < <}}
~
\hyperref[sec:index_chronic]{\small{[返順時目]}}
~
\hyperref[sec:index_scriptual]{\small{[返順卷目]}}
~
\hyperref[sec:2jdjy7oOq1s]{\small{> > > NEXT SERMON > > >}}
\newline
\newline
以弗所書 3:14-21
\newline
\begin{longtable}{cl}
\hline
\hline
章節 & 經文 (和合本修訂版)\\
\hline
3:14 & \begin{tabularx}{0.7\textwidth}{X} 因此，我在父面前屈膝— \end{tabularx} \\ \\ \relax
3:15 & \begin{tabularx}{0.7\textwidth}{X} 天上地上的各家都是從他得名的－ \end{tabularx} \\ \\ \relax
3:16 & \begin{tabularx}{0.7\textwidth}{X} 為要他按著他豐盛的榮耀，藉著他的靈，使你們內心的力量剛強起來； \end{tabularx} \\ \\ \relax
3:17 & \begin{tabularx}{0.7\textwidth}{X} 又要他使基督因著你們的信住在你們心裡，使你們既在愛中生根立基， \end{tabularx} \\ \\ \relax
3:18 & \begin{tabularx}{0.7\textwidth}{X} 能夠和眾聖徒一同明白基督的愛是何等的長、闊、高、深， \end{tabularx} \\ \\ \relax
3:19 & \begin{tabularx}{0.7\textwidth}{X} 並知道這愛是超過人的知識所能測度的，為要使你們充滿神一切的豐盛。 \end{tabularx} \\ \\ \relax
3:20 & \begin{tabularx}{0.7\textwidth}{X} 神能照著運行在我們心裡的大能充充足足地成就一切，超過我們所求所想的。 \end{tabularx} \\ \\ \relax
3:21 & \begin{tabularx}{0.7\textwidth}{X} 願他在教會中，並在基督耶穌裡，得著榮耀，直到世世代代，永永遠遠。阿們！ \end{tabularx} \\ \\
[1ex]
\hline
\hline
\end{longtable}
$^{1}$(speaking in foreign language).
- Good morning, brothers and sisters..
(speaking in foreign language).
Let's continue to study Efficience..
(speaking in foreign language).
- So for this reason, I kneel before the Father,.
from whom every family in heaven and on earth.
derives its name..
I pray that out of his glorious riches,.
he may strengthen you with power through his spirit.
in your inner being..
So that Christ may dwell in your hearts through faith,.
and I pray that you, being rooted and established in love,.
may have power, together with all people of holy people,.
to grasp how wide and long and high and deep.
is the love of Christ..
(speaking in foreign language).
And to know this love that surpasses knowledge,.
that you may be filled with the measure.
of all the fullness of God..
Now to him who is able to do immeasurably.
more than all we ask or imagine,.
according to his power that is at work within us,.
and to him be glory in the church and the Christ Jesus.
throughout all generations, forever and ever, amen..
(speaking in foreign language).
This passage is Paul's prayer..
(speaking in foreign language).
There are six chapters in Ephesians,.
and this is at the end of the chapter three..
(speaking in foreign language).
It's just like when you're writing a letter.
and you express your thoughts and feelings..
(speaking in foreign language).
So let's review these three chapters..
Chapter one, Paul explains how God chose his people.
in Christ before foundations of the world..
In chapter two, he tells us that sinners are saved.
by grace through faith..
In chapter three, he reveals that this salvation.

$^{41}$has so also come to the Gentiles..
(speaking in foreign language).
So according to our logic,.
after explaining so much theology,.
Paul should say, "Now it's your time..
"Go ahead and work harder..
"Try your best and leave all this out..
"We Asian parents love, like to talk to kids in this way.".
(speaking in foreign language).
But Paul does not do that..
Instead, Paul prays, he worships God..
(speaking in foreign language).
Why?.
Because God knows that the greatest problem of beliefs.
is not lack of knowledge..
This real problem is that truth often remains in our mind.
and never reaches our heart..
Truth may not transform our lives into worship..
(speaking in foreign language).
So we may know that God's sovereign yet still live.
in anxiety..
We may know that God loved us, yet still lack security..
The problem is not simply knowing the truth..
The problem is that truth has not yet gone deep.
into our lives..
(speaking in foreign language).
So Paul says, "So Paul prays that this truth.
"will take roots in believers' heart.
"and transform into their lives.".
Today, Paul's prayer teaches us three things..
(speaking in foreign language).
So Paul says, "For this reason,.
"I kneel before the Father.".
After teaching deep truth,.
Paul's first response is to kneel..
(speaking in foreign language).
So today, kneeling is rare in modern society..
In ancient times, people kneeled before kings..
Even today, some royal traditions.
still preserve this practice..

$^{81}$Kneeling represents surrender..
It means acknowledging that the one before me.
is greater than I am..
In this scripture, kneeling before God.
means acknowledging his greatness..
It means confessing, "I am weak, but God is strong.".
(speaking in foreign language).
This is where Paul begins his prayer..
He kneels before God and admits, "I cannot, but God can.".
(speaking in foreign language).
Often, we are unwilling to kneel in prayer..
All our prayers become mere formal..
On the surface, we may say we're too busy,.
but the deeper reason is this..
In our hearts, we believe that we can handle life ourselves..
(speaking in foreign language).
We trust our abilities, our experience, our connections..
Only when a person truly recognize his limitation.
does he begin to pray seriously..
All of us, not just you, but also me as well,.
need to kneel before God..
(speaking in foreign language).
Many churches' problems are not solved by better programs,.
more equipment, more activities, or larger budgets..
The real solution is often prayer..
A church without prayer may look active.
and successful on outside..
It may have many problems, money, and resources,.
but inwardly, it can be empty..
(speaking in foreign language).
A church without prayer is empty..
It is like a tree, a hollow tree..
From outside, it looks alive, full of leaves and strong,.
but inside, there is no life..
(speaking in foreign language).
I grew up deeply influenced by Chinese house church,.
family church. In the past time,.
believers often prayed together on their knees..
Today, there is less common.
in both Chinese churches and American churches..

$^{121}$(speaking in foreign language).
So that does not mean kneeling.
makes someone more spiritual,.
but throughout church history,.
kneeling has been a visible expression.
of humility before God..
(speaking in foreign language).
If you have never tried kneeling in prayer,.
I encourage you to try it at home..
You may learn something new..
(speaking in foreign language).
What does Paul pray for?.
He does not pray first for better circumstances..
He does not pray for persecution to be stopped..
He does not pray for comfort, success, or prosperity..
(speaking in foreign language).
He prays that believers should be strengthened.
by in their inner being..
Paul knows that that determines.
whether a person stands from firm.
is not primarily external circumstances,.
but inner strength..
Two people might face the same difficulty..
One collapse, the other grows stronger..
The difference is often found in the heart..
(speaking in foreign language).
Think about a cup of water sitting on this table..
When the table shakes, water spills out..
Our hearts are similar..
When light shakes us, what comes out?.
Despair, anger, jealous, anxiety?.
These things come out because they were already inside..
(speaking in foreign language).
Many people today look fine on outside,.
but inwardly they are becoming weaker..
This scripture calls us to deal with the heart..
(speaking in foreign language).
We often want God to change our circumstances..
Our prayers are often focused on removing difficulties..
There is nothing wrong about that..

$^{161}$No one wants suffering or hardship,.
but God cares deeply about how our hearts respond..
(speaking in foreign language).
More than changing our circumstances,.
God desires to strengthen us inwardly..
So when facing pressure, difficulties, and struggles,.
do not only ask God to change our circumstances..
Ask him to change your heart..
(speaking in foreign language).
Paul says so that Christ may dwell in your heart.
through faith, the word dwell means to live permanently..
It is not a temporary visit..
Christ is not like a guest we welcome into our home.
for a coffee or conversation,.
only to leave shortly afterwards..
Christ comes to live in us..
(speaking in foreign language).
Being Christian is not merely just pray to declare.
that I am Christian only once..
It means Christ dwells in us..
Too often we treat Jesus like a guest..
We welcome him on Sundays,.
we welcome him during our difficult times,.
but during the rest of the week,.
when we are in peaceful time, we remain in control..
Paul reminds us Christ is not a guest, Christ is Lord..
(speaking in foreign language).
Paul then says believers must be rooted.
and grounded in love..
He uses two images, a tree and a building..
A tree needs deep roots,.
a building needs a strong foundation..
In downtown Chicago, the taller the sky crapper,.
the deeper the foundation must be..
We live in Chicago and we all know.
that storm can come at any time..
(speaking in foreign language).
Strong winds and storms will reveal.
whether roots are deep..
Life is the same, storms will come..

$^{201}$The real question is, where are your roots?.
If your roots are in anything other than Christ,.
we will eventually be shaken..
Only those rooted in Christ's love will remain firm..
(speaking in foreign language).
Paul describes Christ's love with four dimensions.
or four words, wide, long, high, deep..
Long means God has loved his people from eternity..
Wide is God loves, reaches all kind of people,.
no matter if it's Jewish or Gentiles,.
rich or poor, people from every nations and language..
High means throughout Christ,.
sinner become children of God..
And deep is Christ humble himself,.
enter our broken world,.
suffer and died on a cross to save his people..
(speaking in foreign language).
Paul does not simply say, study Christ's love..
He says, experience it..
Why?.
Because human's deep problem is not behavior, it is love..
After being separated from God,.
the human's heart constantly ask, am I worthy of love?.
Do I matter?.
That is why we pursue money, success, status,.
and achievement to prove our worth..
But none of these can truly satisfy us..
Only love of God can satisfy the human soul..
(speaking in foreign language).
So Christ's love changes our value system..
(speaking in foreign language).
Without Christ's love, we fear about everything..
We fear failure..
We fear about rejection..
We fear lost..
But when we truly experience Christ's love,.
we become free because we know our value.
is not defined by how the world sees us..
Our value is defined by how Christ sees us..
(speaking in foreign language).

$^{241}$So today, what defines you?.
Your wealth, your job, your marriage, your children,.
your achievements, or other people's opinions,.
or is Christ's love?.
Many of us have believed for so many years,.
yet deep inside us, deep inside our heart,.
we still allow the world to define us..
That is why when we lost something important,.
we collapse..
But the person rooted in Christ's love remains stable,.
even when everything around them shakes..
(speaking in foreign language).
Verses 20 to 21 bring Paul's prayer to its climax..
He prays the greatness of God's power..
God's glory is eternal..
He is glorified in the church and in Christ Jesus forever..
Paul declares that God can do far more than we can imagine,.
not a little more, not slightly more, immeasurably more..
God's power is greater than our struggles..
God's plans are higher than our plans..
God's grace exceeds our imagination..
(speaking in foreign language).
Why does Paul speak about God's power here?.
(speaking in foreign language).
- Why?.
Because everything he has prayed for.
sounds so impossible, right?.
Inner strength, being filled with Christ's love,.
a transformed life..
What are you talking about?.
I mean, our first reaction is often like,.
maybe Paul, you yourself can do that..
Well, I cannot..
(speaking in foreign language).
- Paul agrees by ourselves, we cannot, but God can..
Many times we only see our problems,.
limitations and failures..
We see broken families, struggles in church,.
weakness in ourselves, but Paul reminds us,.
God's power is still at work..

$^{281}$Perhaps we feel weak, perhaps we feel defeated,.
perhaps we think nothing can be changed,.
but God is able..
(speaking in foreign language).
- Where is his power at work?.
Within us..
God works first in the heart..
As we said earlier, God often does not begin.
by changing circumstances, he begin by changing people..
That is why Paul emphasized inner strength.
and being filled with Christ's love..
Without inner transformation, even improved situation.
cannot truly help us..
Miracles are not always dramatically external changes..
Sometimes the greatest miracle is inner transformation,.
from fear to courage, from bitterness to forgiveness,.
from despair to hope, and that is the power of God..
(speaking in foreign language).
- Finally, Paul says, to him be glory in the church.
and in the Christ Jesus..
The ultimate goal is not human success..
It is not our comfort..
It is not our satisfaction..
The ultimate goal is God's glory..
(speaking in foreign language).
- When things seem impossible,.
believe that God is still at work..
We often fall into despair and decide too quickly.
that nothing can be changed..
We say this situation will never improve..
There is no hope, but man's limit is not God's limit..
We see obstacles, God sees possibility..
Have you lost hope in something?.
Paul reminds us, do not limit God.
with your human perspective..
(speaking in foreign language).
- God often works quietly, silently,.
like a seed growing on the ground..
His work may be invisible sometimes..
We may not see it, but that does not mean.

$^{321}$nothing is happening..
Even when we cannot feel God, He is still working..
Even in silence, God is shaping us..
(speaking in foreign language).
- Paul's prayer teaches us three important things.
and the truths..
One, strength begins with prayer..
Number two, growth comes from knowing.
and experiencing Christ's love..
Number three, true transformation happens from inside out..
This transformation does not come.
through human efforts alone..
It comes through God's working in us..
The greatest problem for Christian.
is not that we are not busy enough.
or there's not heart working enough..
The real question is this,.
are we filled with the love of Christ?.
(speaking in foreign language).
- Let's pray..
(speaking in foreign language).
- Heavenly Father, help us recognize our need.
to be changed by you..
Help us experience more of your love.
and more of your power..
Circumstances often make us discouraged,.
but with you, all things are possible..
Teach us to humble ourselves before you.
and begin with prayer..
Strengthen us through prayer..
In Jesus' name, amen..
- Amen..
\newpage



\section{傳道書 5:8-6:8}
\label{sec:2jdjy7oOq1s}
\textbf{錢值幾錢?(傳道書5\_8-6\_8) - 譚文倩莊傳道}
\newline
\newline
連結: \href{https://youtube.com/watch?v=2jdjy7oOq1s}{\texttt{https://youtube.com/watch?v=2jdjy7oOq1s}} ~~~~ 語音日期: 2026-07-12
\newline
\newline
\hyperref[sec:DRGVYhPEXTM]{\small{< < < PREV SERMON < < <}}
~
\hyperref[sec:index_chronic]{\small{[返順時目]}}
~
\hyperref[sec:index_scriptual]{\small{[返順卷目]}}
~
\hyperref[sec:3LIoE3b6SXA]{\small{> > > NEXT SERMON > > >}}
\newline
\newline
傳道書 5:8-6:8
\newline
\begin{longtable}{cl}
\hline
\hline
章節 & 經文 (和合本修訂版)\\
\hline
5:8 & \begin{tabularx}{0.7\textwidth}{X} 你若在一個地區看見窮人受欺壓，公義公平被掠奪，不要因此驚奇；有一位高過居高位的在鑒察，在他們之上還有更高的。 \end{tabularx} \\ \\ \relax
5:9 & \begin{tabularx}{0.7\textwidth}{X} 況且地的益處歸眾人，就是君王也受田地的供應。 \end{tabularx} \\ \\ \relax
5:10 & \begin{tabularx}{0.7\textwidth}{X} 喜愛銀子的，不因得銀子滿足；喜愛財富的，也不因得利益知足。這也是虛空。 \end{tabularx} \\ \\ \relax
5:11 & \begin{tabularx}{0.7\textwidth}{X} 貨物增添，吃的人也增添，物主得甚麼益處呢？不過眼看而已！ \end{tabularx} \\ \\ \relax
5:12 & \begin{tabularx}{0.7\textwidth}{X} 勞碌的人不拘吃多吃少，睡得香甜；富人的豐足卻不容他睡覺。 \end{tabularx} \\ \\ \relax
5:13 & \begin{tabularx}{0.7\textwidth}{X} 我見日光之下有一件令人憂傷的禍患，就是財主積存財富，反害自己。 \end{tabularx} \\ \\ \relax
5:14 & \begin{tabularx}{0.7\textwidth}{X} 他因遭遇不幸，財產盡失；他生了兒子，手裡卻一無所有。 \end{tabularx} \\ \\ \relax
5:15 & \begin{tabularx}{0.7\textwidth}{X} 他怎樣從母胎赤身而來，也必照樣赤身而去；他所勞碌得來的，手中分毫不能帶去。 \end{tabularx} \\ \\ \relax
5:16 & \begin{tabularx}{0.7\textwidth}{X} 這是一件令人憂傷的禍患。他來的時候怎樣，去的時候也必怎樣。他為風勞碌有甚麼益處呢？ \end{tabularx} \\ \\ \relax
5:17 & \begin{tabularx}{0.7\textwidth}{X} 並且他終身在黑暗中吃喝，多有煩惱、病痛和怒氣。 \end{tabularx} \\ \\ \relax
5:18 & \begin{tabularx}{0.7\textwidth}{X} 看哪，我所見為善為美的，就是人在神賜他一生的日子吃喝，享受日光之下勞碌得來的好處，因為這是他應得的報償。 \end{tabularx} \\ \\ \relax
5:19 & \begin{tabularx}{0.7\textwidth}{X} 而且，一個人蒙神賞賜財富與資產，又使他能享用，能獲取自己當有的報償，在他的勞碌中喜樂，這是神的賞賜。 \end{tabularx} \\ \\ \relax
5:20 & \begin{tabularx}{0.7\textwidth}{X} 他不多思念自己一生的日子，因為神使他的心充滿喜樂。 \end{tabularx} \\ \\ \relax
6:1 & \begin{tabularx}{0.7\textwidth}{X} 我見日光之下有一件禍患重壓在人身上， \end{tabularx} \\ \\ \relax
6:2 & \begin{tabularx}{0.7\textwidth}{X} 就是人蒙神賜他財富、資產和尊榮，以致他心裡所願的一樣都不缺，只是神使他不能享用，反被外人享用。這是虛空，也是禍患。 \end{tabularx} \\ \\ \relax
6:3 & \begin{tabularx}{0.7\textwidth}{X} 人若生一百個兒子，活許多歲數；他即使壽命很長，心裡卻不因福樂而滿足，又不得埋葬；我說，那流掉的胎比他倒好。 \end{tabularx} \\ \\ \relax
6:4 & \begin{tabularx}{0.7\textwidth}{X} 因為這胎虛虛而來，暗暗而去，名字被黑暗遮蔽， \end{tabularx} \\ \\ \relax
6:5 & \begin{tabularx}{0.7\textwidth}{X} 而且沒有見過天日，甚麼都不知道，這胎比那人倒享安息。 \end{tabularx} \\ \\ \relax
6:6 & \begin{tabularx}{0.7\textwidth}{X} 那人雖然活千年，再活千年，卻不能享福；眾人豈不都歸同一個地方去嗎？ \end{tabularx} \\ \\ \relax
6:7 & \begin{tabularx}{0.7\textwidth}{X} 人的勞碌都為口腹，心裡卻不知足。 \end{tabularx} \\ \\ \relax
6:8 & \begin{tabularx}{0.7\textwidth}{X} 智慧人比愚昧人有甚麼益處呢？困苦人在眾人面前知道如何行，有甚麼益處呢？ \end{tabularx} \\ \\
[1ex]
\hline
\hline
\end{longtable}
$^{1}$終於來到傳道書的中間.
猜猜今天這段經文 傳道者要跟我們說什麼.
我在練習的時候 發覺已經告訴大家今天要說什麼.
如果今天是第一次來到我們教會.
我想代表教會歡迎你.
完了崇拜之後 歡迎你留步.
讓我們大家可以有互相認識的時間.
好 回到傳道書.
傳道書是聖經舊約裡面其中一本智慧文學.
探討人生的終極意義 道德的實踐.
苦難的奧秘.
和我們人和做物主之間的關係.
今天的經文在傳道書的五章八節到六章九節.
傳道書的作者今天要跟我們說錢.
你聽到這個主題的時候 是否馬上醒了.
還是馬上在想 說錢會傷感情.
我覺得不是很怕傷感情 因為與我無關.
不是我說的 傳道者說的.
我只是依經直說.
不如你又想想 錢有什麼好說的.
在這段經文裡面.
傳道者不是站在道德高地上.
跟我們說錢是好的 或者錢是不好的.
而是他要藉著說錢這件事.
要我們思考一下.
喜歡錢的人不會因為錢而滿足.
錢多又如何.
你知不知道可以分分鐘得個體字.
然後還要慢慢失眠.
為什麼呢.
有賺錢的能力 又懂得存錢.
但可以無神神發生一些事 一鋪清袋.
有財有勢一無所缺.
但卻沒有享受的能力.
很殘忍吧.
就像回到第一章那樣說.
虛空的虛空.
一直在聽 不知道你會不會覺得.
很depressing 這麼慘.
其實傳道者在這段經文裡面.

$^{41}$有什麼好說的呢.
有沒有好的東西在裡面說的呢.
是有的.
會不會在五章十八至二十節那裡說.
看我所看見為善為美的.
就是人在神次他一生的日子吃喝.
享受日光之下勞碌得來的好處.
因為這是他應得的報償.
而且一個人夢神賞賜財富與資產.
又使他能享用.
能獲取自己當有的報償.
在他的勞碌中喜樂.
這是神的賞賜.
他不多思念自己一生的日子.
因為神使他的心充滿喜樂.
意思是否叫我們不用想太多.
又不要計算太多.
總之有得吃你就吃.
有得飲你就飲.
有工返就返工.
開開心心過日子.
存著感恩的心去領受.
就叫做懂得享受上帝給你的恩典.
如果我告訴你不是.
至少這個不是重點.
那你猜一下甚麼是重點呢.
不如我們先一起低頭祈禱.
我們求聖靈幫助我們去聆聽神的話語.
我們人算得是甚麼.
你竟然眷顧我們.
求你幫助我們.
我們渴望聽到你的聲音.
奉主耶穌基督成名祈求.
1980年林子祥出了一首歌.
叫做《分分鐘需要你》.
懂得唱嗎.
願你會zar 火箭帶我到天空去.
在太空中兩人住.
以前聽可能就想沒甚麼可能.
現在去太空住真的好像可能有可能.

$^{81}$他後面說甚麼.
有了你開心D .
乜都清心滿意.
鹹魚白菜.
又好好味.
日日食會不會好好味.
現在經常說鹹魚是一個processed food.
吃得多會致癌.
這首歌說甚麼呢.
他說不需要山珍海味.
不需要豪華的生活.
只要跟你一起我就滿足.
五章十字十二字.
喜愛銀子的不因得銀子滿足.
喜愛財富的也不因得利益之足.
這也是虛空.
貨物增添吃的人也增添.
物主得甚麼益處呢.
不過眼看而且.
勞碌的人不拘吃多吃少.
睡得香甜.
婦人的豐足卻不容他睡覺.
錢沒問題.
喜歡錢也沒問題.
不滿足才是問題.
Jandith Rockefeller.
美國有史以來其中一個最有錢的人.
曾經有人問他.
還要多賺多少錢才算夠.
他說再多一點點.
再多一點點就夠了.
所以喜愛銀子的人.
不會因銀子滿足的問題.
是在於人那種多一點點.
才夠的心態.
永遠都不滿足.
還記得史圖保羅在提摩太前書六章十節.
說到貪錢財是萬惡之根.
為甚麼.
原因是我們人和錢建立了一個錯誤的關係.

$^{121}$這裡說的喜愛銀子.
有一種朋友的概念在裡面.
記不記得之前我說過.
ahalf這個字.
一個深層忠誠跟你同行.
患難的關係.
陪你 撐你.
在你有困難的時候.
不會離開你.
問題是.
人會否把錢當成自己最好的朋友.
就是那個你可以依靠.
在你困難的時候幫助你的關係.
你期望錢對你忠心.
滿足你的需要.
在你最需要人陪你 撐你的時候.
不離開你.
一種錯誤的關係.
經文繼續說.
莫增添 吃的人也增添 物主得甚麼益處 不過眼看如此.
換句話說.
當你的財富增加的時候.
圍住你 消耗你的資源的人也會增加.
針元也有一個很類似的提醒.
財物使朋友增多.
窮人連朋友也遠離他.
人人都討好慷慨的人.
送禮的人.
人人都與他為友.
現實是.
當一個人沒有資源的時候.
可以連最親密的關係也疏離.
為甚麼.
因為他已經沒有被需要的價值.
所以財富帶來的是朋友.
還是期望和索取呢.
當你有資源的時候.
你要照顧更多的人.
分配更多的資源.
應付更多的需求.

$^{161}$甚至維持更多的關係.
相對 如果你沒有或沒有甚麼的時候.
其他人對你沒有甚麼期望.
你與對方的關係還來得簡單一點.
勞碌的人不拘吃多吃少.
睡得香甜.
富人的風築卻不用他睡覺.
心理學上有一個很簡單的觀察.
一個人的壓力大不大.
你看他晚上睡得好不好你就知道了.
傳道者在這裡做了一個對比.
他說一個辛勤工作.
但生活簡單的人.
他不是很富有.
但他有夠好睡.
擁有多的人.
因為背著很多掛慮.
就算回到家的時候.
他的腦袋仍然不停在運作.
那個deal最後成不成交.
有沒有漏了甚麼.
成不成交.
成不成交怎麼辦.
對手會出甚麼招數.
換句話來說.
財富不會令你沒有問題.
而是令你不停處理問題.
喜歡錢的人不會因為此而不滿足.
分分鐘還慢慢失眠.
13至17節說.
我見日光之下有一件令人憂傷的禍患.
就是財主積存財富反害自己.
他因遭遇不幸財產盡失.
他生了兒子手裡卻一無所有.
他怎樣從母胎赤身而來.
也必照樣赤身而去.
他所勞碌得來的手中分毫不能帶去.
這是一件令人憂傷的禍患.
他來的時候怎樣.
去的時候又怎樣.

$^{201}$他為勞碌有甚麼益處.
並且他終身在黑暗中吃喝.
多有煩惱,病痛和怒氣.
不知道大家會否像我一樣.
還經常留意香港的新聞.
香港最近有幾宗很大的財經新聞.
震撼彈.
其中一個就是抱黃波叔.
他在2021年離開之後.
他的兒子接管了他的家族生意.
他在疫情之前買了很多舊酒店.
打算翻身然後打造酒店的王國.
但遇上疫情的關係.
大量的資金被扣除.
加上樓市大跌.
他們家族用了四五十年建立回來的八百億.
八百億的王國遺產不斷在縮水.
今年2026年他被兩間公司申請破產.
發生這件事大約四五年.
用了四五十年建立回來的八百億公司的王國.
四五年被申請破產.
另一個鄭志剛曾經是香港第三大富豪.
鄭裕彤的孫子.
他主政期間大撒金錢搞藝術.
興建你們可能有些人都去過的K11 MUSE.
結果新世界發展.
在2024年蝕了171億.
是20年來首次虧損.
市值蒸發了720億.
他因遭遇不幸財產盡失.
投資失敗遇上疫情生意蝕本.
八百億是什麼錢.
八百萬我都沒見過.
八百億真的很難想像.
天災人禍.
被人騙被人搶.
我們這個世界可以這麼無厘頭.
可以發生一些你猜都猜不到.
好像沒什麼可能會發生的事.
馬太福音六章有一句經文.

$^{241}$可能我們當中有些人都很熟悉.
不要為自己積攢財寶在地上.
答真一點這句經文在想的時候.
思考一下這些新聞.
你都不是說不怕.
幾百億.
什麼都沒了.
我在想他們沒了那些東西會怎樣.
找回那400呎單位住.
很諷刺怎麼接受.
就是這麼瘋狂.
十六節這是一件令人憂傷的禍患.
可能你們要幫幫我.
這是一件令人憂傷的禍患.
他來的時候怎樣去的時候也怎樣.
他為勞碌有什麼益處.
傳道者批判的不是反對努力.
工作或者追求卓越.
事實上如果你記得的時候.
傳道者之前也多次肯定.
勞碌本身是神給我們的其中一個恩賜.
我們人也應該去享受.
勞碌之後所得來的份.
但這裡他真正想指出.
我們人生充滿了很多.
我們控制不了的變數.
很多超出我們可以控制的事情.
悲劇是會發生的.
損失是會出現的.
計劃都可以被打亂.
It happens.
問題不是在說我們有沒有努力.
而是我們要Better think twice.
我們的安全感.
是否建立在這些東西上.
在這些財富或者成就.
我們想著可以建立回來.
給我安全的這些東西上.
六章1-9.
我見日光之下有一件禍患.

$^{281}$重壓在人身上.
就是人夢.
神賜他財富資產和尊榮.
以致他心裡所願的一樣都不缺.
只是神賜他不能享用.
反被愛人享用.
這是虛空也是禍患.
人若生了一百個兒子.
活許多歲數.
他即使壽命很長.
心裡卻不因福樂而滿足.
有不得埋葬我說.
那留掉的胎被他討好.
因為這胎虛虛而來暗暗而去.
明智被黑暗掩蔽.
而且沒有見過天日.
什麼都不知道.
這胎被那人倒向安息.
那人雖然活千年再活千年.
卻不能享福.
眾人豈不都歸同一個地方去嗎.
人的勞碌都為後福.
心裡卻不知足.
智慧人比愚昧人有什麼益處呢.
困苦人在眾人面前.
知道如何行有什麼益處呢.
眼睛所看的比心裡妄想的都好.
這也是虛空.
也是暴風.
一個人可以擁有很多很多.
但若果神沒有給他享受的能力.
他其實和沒有是沒有分別的.
之前我和大家說的John D. Rockefeller.
美國有史以來其中一個最有錢的人.
在他大約五十歲的時候.
因為他長期工作的壓力.
他的健康開始出現了問題.
他整天都周身痛.
脫髮.
消化系統出現問題.

$^{321}$他去看醫生.
醫生給他注射了什麼呢.
就是喝奶和吃餅乾.
就算他身體和腸胃康復.
他只可以吃很小的部份的青豆.
生菜.
魚.
Brown rice.
Big potato.
其中一個最有錢的人.
想吃什麼都可以.
卻沒有能力享受.
那個年代的人大約五十歲.
但John D. Rockefeller活到九十七歲.
很多人都很渴望很長的生命.
但他身體一直有很多狀況.
在他身邊和他很親的人說.
John D. Rockefeller是長期失眠的.
長期沒有笑容.
對任何事情都提不起勁.
財富資產一樣都不缺.
但卻沒有享受的能力.
這也是虛空 也是暴風.
回到剛才說的五章十八至二十節.
我們嘗試思考一下這節經文.
他說看那我所見為善為美的.
就是人在神賜他一生的日子吃喝.
享受日光之下勞碌得來的好處.
因為這是他應得的報償.
而且一個人蒙神賞賜財富與資產.
又使他能享用能獲取自己當有的報償.
在他的勞碌中喜樂.
這是神的賞賜.
他不多思念自己一生的日子.
因為神使他的心充滿喜樂.
不知道你有沒有留意這裡.
有一些很重要的教導.
蒙神賞賜使他能享用.
使他能獲取自己當有的報償.
使他在他的勞碌中喜樂.

$^{361}$這是神的賞賜.
還有.
你若在一個地區看見窮人受欺壓.
公義公平被掠奪.
不要因此驚奇有一位高位的在監察.
在他們之上還有更高的.
況且地的益處歸眾人.
就是君王也受田地的供應.
都是蒙神的賞賜.
都是神使他的結果.
結束之前.
我想跟大家分享一個我在新聞看到的故事.
我之前跟大家說過.
香港民福縣在去年十一月的時候.
發生一場大火災.
八十七歲的周婆婆.
她就住在其中一棟大廈裡面的十四樓.
一場大火.
燒毀了她住了二十多年的家.
大半生人的財物.
上個月周婆婆終於可以在她的孫子陪同之下.
回到家裡.
沒有升降機.
撐了十四層樓.
撐到上去的時候滿頭大汗.
去周圍看一下自己的家.
已經被燒到黑了.
被煙燻到黑漆漆.
門窗都破爛了.
火災之後留下一陣臥味.
她走去客廳的角落.
她自己的老朋友那部琴的前面.
戴著手套坐下.
就開始去彈那首她自己很喜歡的詩歌.
《耶和華是愛》.
她一邊彈一邊唱.
她還問孫子你有沒有跟著我唱.
孫子說有啊有啊.
原來周婆婆一早已經帶好歌詞.
準備回家的時候跟她的家人一起唱.

$^{401}$她說《耶和華是愛》這首詩歌.
是她最喜歡最熟悉的詩歌.
上樓之前.
有記者問周婆婆.
她說大火已經燒了你所有的東西.
你有什麼感受.
婆婆就說是啊是啊.
很多東西都燒了.
不過我現在所需要的東西.
我都有.
上帝的恩典夠我用.
所以我的心很平靜.
我很感恩.
她又說苦難是魏莊的祝福.
她說因為這場大火之後.
有很多人去關心她.
甚至一個以前的學生特意從上海飛來香港去探望她.
給她錢.
婆婆在房裡找找找.
其中一樣她說她很想找到的是什麼呢.
是她藏在房裡一堆印著耶穌愛你的毛巾.
她說我很想找到這堆毛巾.
因為每一次我出去傳福音的時候.
我就會將這些毛巾送給那些福音對象.
錢值多少錢.
我們現在是活在人類歷史上最富裕的世代.
但為什麼我們不是一個最滿足的世代.
如果我們人生所有問題都是可以用錢去解決的.
那似乎我們有錢的人就應該是活得很快樂的.
為什麼有人什麼都有卻是享受不到.
反過來有些人沒有什麼.
但卻是活得很開心很滿足呢.
傳道書五章八節到六章九節.
傳道者要教我們認清一個事實.
我們所擁有的全部都是神給的.
是神賜予的.
甚至是神給我們能力去享受.
真正的滿足不是擁有更多更多.
而是我們要學習怎樣從神的手中去領受.
一個八十七歲的老人家.

$^{441}$在失去所有日光之下的財物之後.
你能不能想像他仍然可以笑著坐在琴面前唱.
耶和華是愛.
神保守我.
我什麼都沒有了.
但我很感恩.
神的恩典夠我用.
他心裡記掛的是那些還沒聽到福音.
他很想派一條毛巾給他的人.
周婆婆的一生沒有建立在十四樓那層樓上的家裡.
她的家當裡沒有在她的錢和金器上.
她的人生卻是建立在一個大火點燒都燒不到的事實上.
耶和華是愛.
夢神賞賜的滿足.
神使人心充滿喜樂.
Amen.
無限滿足快樂湧於心裡.
共同度每一天.
再偽爛神都不會變.
我們一起唱這首詩歌好不好.
神令我聆聽心中的憂鬱.
共同度每一天.
耶和華是愛.
在困境中他保守任憂.
神為我擺設風俗的隱諱.
再偽爛也不會變.
在世間輸人與盜.
他的愛存在常無盡.
神為我死因賣力保守免禮.
共同度之世.
耶和華是愛.
讓我安身青草看誰變.
無限滿足快樂湧於心裡.
再偽爛也不會變.
請坐.
大家好.
多謝Joyce給我們的提議.
\newpage



\section{以弗所書 4:1-7}
\label{sec:3LIoE3b6SXA}
\textbf{不只是名字,更是生命 (以弗所書 4\_1-7) - 張威傳道 Not Just a Name But Life - Paul Zhang}
\newline
\newline
連結: \href{https://youtube.com/watch?v=3LIoE3b6SXA}{\texttt{https://youtube.com/watch?v=3LIoE3b6SXA}} ~~~~ 語音日期: 2026-07-20
\newline
\newline
\hyperref[sec:2jdjy7oOq1s]{\small{< < < PREV SERMON < < <}}
~
\hyperref[sec:index_chronic]{\small{[返順時目]}}
~
\hyperref[sec:index_scriptual]{\small{[返順卷目]}}
~
\hyperref[sec:8hcsYSy4IhI]{\small{> > > NEXT SERMON > > >}}
\newline
\newline
以弗所書 4:1-7
\newline
\begin{longtable}{cl}
\hline
\hline
章節 & 經文 (和合本修訂版)\\
\hline
4:1 & \begin{tabularx}{0.7\textwidth}{X} 我為主作囚徒的勸你們，既然蒙召，行事為人就要與你們所蒙的呼召相稱。 \end{tabularx} \\ \\ \relax
4:2 & \begin{tabularx}{0.7\textwidth}{X} 凡事要謙虛、溫柔、忍耐，用愛心互相寬容， \end{tabularx} \\ \\ \relax
4:3 & \begin{tabularx}{0.7\textwidth}{X} 以和平彼此聯繫，竭力保持聖靈所賜的合一。 \end{tabularx} \\ \\ \relax
4:4 & \begin{tabularx}{0.7\textwidth}{X} 身體只有一個，聖靈只有一位，正如你們蒙召，是為同有一個指望而蒙召， \end{tabularx} \\ \\ \relax
4:5 & \begin{tabularx}{0.7\textwidth}{X} 一主，一信，一洗， \end{tabularx} \\ \\ \relax
4:6 & \begin{tabularx}{0.7\textwidth}{X} 一神－就是萬人之父，超越萬有之上，貫通萬有，在萬有之中。 \end{tabularx} \\ \\ \relax
4:7 & \begin{tabularx}{0.7\textwidth}{X} 我們每個人蒙恩都是照基督所量給每個人的恩賜。 \end{tabularx} \\ \\
[1ex]
\hline
\hline
\end{longtable}
$^{1}$[NON-ENGLISH SPEECH].
Good morning, brothers and sisters..
[NON-ENGLISH SPEECH].
Today, we will continue to learn efficiency..
[NON-ENGLISH SPEECH].
It's about not just a name, but life..
[NON-ENGLISH SPEECH].
Therefore, a prisoner for the Lord.
urge you to walk in a manner worthy of the calling.
to which you have been called..
With all humility and gentleness,.
with patience, bearing with one another in love,.
eager to maintain the unity of the spirit.
in the bond of peace..
[NON-ENGLISH SPEECH].
There is one body and one spirit,.
just as you were called to the one hope that.
belonged to your call..
One Lord, one faith, one baptism, one God,.
and the Father of all, who is over all, and through all,.
and in all..
But grace was given to each one of us.
according to the measure of Christ's gift..
[NON-ENGLISH SPEECH].
Before we begin, let me ask you to imagine a scene..
Suppose you walk into the Michelin three-star restaurant..
At the entrance, you see the prestigious Michelin.
certification displayed proudly..
The interior is elegant and luxurious..
Every waiter is dressed sharply in suit,.
and everything just feels professional..
[NON-ENGLISH SPEECH].
You order their signature pasta dish,.
full of anticipation, thinking that today I finally.
got the taste of the work of a world-class chef..
A few minutes later, a waiter brings you your food..
You open the plate, and what do you see?.
Instant noodle soup in plain hot water.
with a sausage on the top..
So you freeze..

$^{41}$[NON-ENGLISH SPEECH].
You call the manager and ask, wait a minute..
Isn't the Michelin three-star restaurant?.
Why are you serving instant noodles?.
The manager comes over and calmly says, yes,.
we are Michelin three-star..
The certification is real..
As for the food, that's not really the important part..
Would that sound absurd to you?.
Of course it is..
[NON-ENGLISH SPEECH].
And yet, many Christians live exactly like that..
We carry the identity of citizen of heaven..
We carry the label Christian..
We say we belong to God, but our lives often.
do not match our identity..
When conflict comes, we lose our temper immediately..
When interests are threatened, we use worldly method..
We gossip, criticize, and tear others down,.
living no differently from the world..
[NON-ENGLISH SPEECH].
In today's passage, Paul tells us something very important..
Since we are called by God, our lives.
should be worthy of that calling..
We should not merely carry a spiritual label.
while living in a worldly life..
Today, we will look at this passage in three parts..
First, those who are called must remember their identity..
Second, those who are called must live out.
the character of Christ..
And third, those who are called must diligently preserve unity..
[NON-ENGLISH SPEECH].
Paul says in verse 1, "I, therefore,.
a prisoner for the Lord, urge you.
to walk in a manner worthy of this calling to which you.
have been called.".
This is a major turning point in the whole book of Ephesians..
What did Paul spend the first three chapters talking about?.
He talked about how Father chose us, how Christ redeemed us,.
how did the Holy Spirit seal us, and once we were dead in sin..

$^{81}$Well, now we have life..
And once we were separated from God,.
now we belong to God's family..
Only in this chapter 4 does Paul begin.
to talk about how should we live..
[NON-ENGLISH SPEECH].
This is the logic of our gospel..
Grace comes first..
Response comes second..
Identity comes first..
Lifestyle comes second..
God always acts first..
And then we respond..
God moves first, and then we follow..
This is the gospel..
[NON-ENGLISH SPEECH].
Now compare what other religions in the world says..
What do they say?.
They will say, fix yourself first..
You have to do your part better first..
And change yourself first..
Maybe God will start loving you..
But what the gospel tells us is something completely different..
The gospel says, God loves you when you are still unworthy..
And this is the grace from God..
[NON-ENGLISH SPEECH].
There is another reason Paul waits until chapter 4.
to talk about Christian living..
Paul knows this..
It's so true that if you don't know who you are,.
you will never live the way that you are supposed to live..
Without gospel and doctrine, what's.
left here or in Christianity?.
When you are told to try hard and better.
and become a better person, that is not gospel..
That is moral advice..
[NON-ENGLISH SPEECH].
And this world always has plenty of moral advice..
When we are in trouble, we always.
focus solely on changing the outside circumstances.

$^{121}$or our behavior, and we miss the root issue--.
understanding the gospel..
Why do we become jealous?.
Why do we fight?.
Why do we hurt each other?.
Because we forgot who we are..
We forgot that we are the people of grace..
[NON-ENGLISH SPEECH].
That's why Paul waits until chapter 4.
to start talking about how we should live..
[NON-ENGLISH SPEECH].
Paul says, "You are called.".
What does that actually mean?.
It means we did not find God first by ourselves..
God calls us first, just like Lazarus..
Jesus stood before a dead man and said, "Lazarus, come out.".
A dead man can never respond on his own.
unless Jesus calls him..
[NON-ENGLISH SPEECH].
In the same way, spiritually speaking, we were dead..
But God called us..
So let me ask you this..
Why are you a Christian?.
Is it because that you are smarter or better or more.
spiritual than others?.
No..
It is about grace..
Grace came to you..
[NON-ENGLISH SPEECH].
When Paul says, "Walking in a manner worthy of your calling,".
what does "worthy" mean?.
It does not mean earning salvation..
Paul is not saying, "Behave better.
so you deserve this salvation.".
That's not what he means..
[NON-ENGLISH SPEECH].
So what he actually means here--.
let your life match the grace you have received..
Live in a way that fits your identity..
If your identity is heavenly, don't live like the world..

$^{161}$Imagine Pastor Mark got drunk every day,.
lying outside of our church building..
So what would you say?.
Then you might say, that doesn't fit..
It doesn't match..
In the same way, God has given us a glorious identity..
Our lives should reflect that identity..
Next, let's look at the second point..
[NON-ENGLISH SPEECH].
Verses 2 says, "With all humility and gentleness,.
with patience, bearing with one another in love.".
Before Paul talks about unity, he first talks about characters..
Why?.
Because the biggest problem in a church.
is usually not the structure..
It is the people inside the church..
A church may have a perfect system, a strong structure,.
a well-organized ministry..
But if the people are spiritually immature,.
everything can still fall apart..
That's why Paul starts talking about characters..
[NON-ENGLISH SPEECH].
Paul gives us four types of characters..
The first one, he mentioned humility..
In Greek culture, humility was seen as weakness and not.
that good in that culture..
But in scripture, humility is a character of Christ..
Jesus is the Son of God..
He's the true King..
But where was he born?.
Not in a palace, in a mansion, not into a noble family,.
but into a home of a carpenter..
He humbled himself and obeyed the Father,.
even to the point of death on a cross..
[NON-ENGLISH SPEECH].
So what is biblical humility?.
It's not low self-esteem..
It's not thinking poorly of yourself..
True humility means this..
You are no longer centered on yourself..

$^{201}$You are centered on God..
Many conflicts in church are rooted in pride..
Why won't the people listen to me?.
Why won't they do the things in my way?.
Why doesn't anyone respect me?.
Pride destroys unity, and humility builds unity..
[NON-ENGLISH SPEECH].
Second, gentleness..
Gentleness is not weakness..
Gentleness is strength under control..
Think of a horse..
A horse is powerful, but when trained,.
the strength is under the control of its master..
That is gentleness..
Jesus said, "I am gentle and lowly in heart.".
Jesus was gentlest person who ever lived..
And yet, he was full of power..
He calmed storms..
He cast out demons..
He raised the dead..
Gentleness does not mean having no power..
It means having power, but not abusing it..
[NON-ENGLISH SPEECH].
The third point is patience, and the fourth one.
is bearing with one another..
There were always some people in the church.
that we think we might not like..
But the same is true in the family life, too..
Sometimes even our spouse, children, parents,.
sometimes we have conflict moment with them..
However, why should we be patient?.
Because Christ has been patient with us..
We've seen against God again and again..
We've failed him again and again..
Let him disappointed..
Yet God has shown enormous patience towards us,.
and we must treat others the same way..
[NON-ENGLISH SPEECH].
Paul adds an important phrase, "with all.".
Not just humility, not just gentleness,.

$^{241}$but with all humility..
And this means not only when life is easy,.
but in every situation, even though in the toughest.
moment in your life, in good times, in hard times,.
when life is smooth or when life is painful..
[NON-ENGLISH SPEECH].
Now, here is the real question..
Can people actually produce these characters on their own?.
Can we truly be humble?.
Can we truly be gentle?.
Can we truly be patient only by our efforts?.
No..
What comes naturally to us?.
Protecting ourselves, defending ourselves,.
exhorting ourselves, because our nature tendency.
is always self-centered..
[NON-ENGLISH SPEECH].
Paul is not saying, try hard to be a better person..
Improve yourself a little bit more..
Be more disciplined..
No..
Paul is saying, you need the gospel..
When we truly understand the gospel,.
the Holy Spirit begins shaping the life of Christ in us..
Humility and gentleness are not something we fake..
They are not something we force ourselves to perform..
They are the fruit of spirit..
[NON-ENGLISH SPEECH].
Only when you understand what the gospel can,.
you can become truly humble..
Only when you understand the gospel.
can you become truly gentle..
And only when you understand the gospel.
can you preserve the true unity..
Apart from the gospel, there is no real humility,.
no real gentleness, and no real unity..
[NON-ENGLISH SPEECH].
And even more than that, those four characters.
describe Jesus himself..
Jesus is humble, Jesus is gentle, he is patient,.

$^{281}$Jesus is gracious..
The Christian life is about reflecting Christ.
like the moon reflects the light of the sun..
Our lives are meant to reflect Christ..
[NON-ENGLISH SPEECH].
Paul emphasized unity here..
Why is unity so important?.
Because God is the God of unity..
God's plan is that in the fullness of time,.
all things will be united in Christ..
And the goal of salvation is reconciliation..
Reconciliation between God and man,.
the reconciliation among people..
Sin always works against unity..
Sin brings division, conflict, hostility,.
and broken relationship..
[NON-ENGLISH SPEECH].
Notice Paul's word..
He does not say create unity..
He says maintain unity..
Why?.
Because true unity is not something that we can create..
It is something the Holy Spirit has already given..
[NON-ENGLISH SPEECH].
Christ has always torn down the division war.
through his death on cross..
The problem is not that church has no unity..
The real problem is that we often.
destroy the unity God has already given us..
How do we destroy it?.
By our pride, bitterness, judgment, conflict,.
and division..
[NON-ENGLISH SPEECH].
That is why Paul says we must be eager to maintain unity..
What does eager mean?.
It means unity does not come easy..
It costs something..
We must lay down ourselves..
We must forgive..
We must be patient..

$^{321}$[NON-ENGLISH SPEECH].
The foundation of church unity is not human effort..
The foundation is unity of triune God..
One body, one spirit..
This points to the Holy Spirit who makes the church one body..
One Lord, one faith, one baptism..
This point to Christ..
We belong to the same Lord, one God and Father..
This points to the Father..
We share the same Heavenly Father..
The Father, the Son, and the Holy Spirit are not divided..
They are perfectly united..
That is why the church must be united..
[NON-ENGLISH SPEECH].
What does Satan love to see the most?.
Not church closing, not Sunday worship stopping..
What Satan loves most is this..
Church still gathers..
The song are still singing..
The sermon still preached..
But brothers and sisters become cold to each other..
They attack one another, judge one another,.
hurt one another..
That destroyed the witness of the church..
[NON-ENGLISH SPEECH].
But then verse 7 says, grace was given to each one of us.
according to the measure of Christ's gift..
What does that mean?.
Unity does not mean uniformity..
God does not mass produce people..
And the body is not made up of only eyes..
Or the body is not made up of only with a lot of pair of hands..
Every person is different..
Each person has different gifts..
Some are gifted in teaching..
Some are gifted in caring for people..
Some in serving..
Some in leading..
Unity does not erase differences..
A healthy church is united in truth and diverse in gifts..

$^{361}$[NON-ENGLISH SPEECH].
Brothers and sisters, the real question.
is not whether we attend church today..
The real question is, do we look like Christ?.
Can your family see Christ in you?.
Can your children see Christ in you?.
Can your brothers and sisters see Christ in you?.
[NON-ENGLISH SPEECH].
Are you building up the church or trying to tear it down?.
Are you bringing peace or creating division?.
May God have mercy on us..
May God renew us..
May God help us walk in manner worthy of the calling.
we have received..
[NON-ENGLISH SPEECH].
Heavenly Father, we thank you for calling us out of darkness.
into your marvelous light and from death into life..
Thank you for giving us a new identity in Christ..
You have given us such a glorious calling,.
and we ask you would help us live.
lives worthy of the calling..
Lord, give us hearts of humility, gentleness,.
and patience..
Help us to reflect the character of Christ in our daily lives..
Teach us to love one another, bear with one another,.
and preserve the unity of the church..
In Jesus' name, amen..
\newpage



\section{哥林多前書 11:2-16}
\label{sec:8hcsYSy4IhI}
\textbf{再思女性的事奉 (哥林多前書11\_2-16) - 陳崇基牧師}
\newline
\newline
連結: \href{https://youtube.com/watch?v=8hcsYSy4IhI}{\texttt{https://youtube.com/watch?v=8hcsYSy4IhI}} ~~~~ 語音日期: 2026-07-20
\newline
\newline
\hyperref[sec:3LIoE3b6SXA]{\small{< < < PREV SERMON < < <}}
~
\hyperref[sec:index_chronic]{\small{[返順時目]}}
~
\hyperref[sec:index_scriptual]{\small{[返順卷目]}}
~
\hyperref[sec:code]{\small{> > > NEXT SERMON > > >}}
\newline
\newline
哥林多前書 11:2-16
\newline
\begin{longtable}{cl}
\hline
\hline
章節 & 經文 (和合本修訂版)\\
\hline
11:2 & \begin{tabularx}{0.7\textwidth}{X} 我稱讚你們，因為你們凡事記得我，又堅守我所傳授給你們的。 \end{tabularx} \\ \\ \relax
11:3 & \begin{tabularx}{0.7\textwidth}{X} 但是我要你們知道：基督是男人的頭；男人是女人的頭；神是基督的頭。 \end{tabularx} \\ \\ \relax
11:4 & \begin{tabularx}{0.7\textwidth}{X} 凡男人禱告或講道，若蒙著頭，就是羞辱自己的頭。 \end{tabularx} \\ \\ \relax
11:5 & \begin{tabularx}{0.7\textwidth}{X} 凡女人禱告或講道，若不蒙著頭，就是羞辱自己的頭，因為這就如同剃了頭髮一樣。 \end{tabularx} \\ \\ \relax
11:6 & \begin{tabularx}{0.7\textwidth}{X} 女人若不蒙著頭，就該剪了頭髮；女人若以剪髮剃髮為羞愧，就該蒙著頭。 \end{tabularx} \\ \\ \relax
11:7 & \begin{tabularx}{0.7\textwidth}{X} 男人本不該蒙著頭，因為他是神的形像和榮耀；但女人是男人的榮耀。 \end{tabularx} \\ \\ \relax
11:8 & \begin{tabularx}{0.7\textwidth}{X} 起初，男人不是由女人而出，女人卻是由男人而出。 \end{tabularx} \\ \\ \relax
11:9 & \begin{tabularx}{0.7\textwidth}{X} 而且男人不是為女人造的，女人卻是為男人造的。 \end{tabularx} \\ \\ \relax
11:10 & \begin{tabularx}{0.7\textwidth}{X} 因此，女人為天使的緣故應當在頭上有服權柄的記號。 \end{tabularx} \\ \\ \relax
11:11 & \begin{tabularx}{0.7\textwidth}{X} 然而，照主的安排，女人不可沒有男人，男人也不可沒有女人。 \end{tabularx} \\ \\ \relax
11:12 & \begin{tabularx}{0.7\textwidth}{X} 因為女人原是由男人而出，男人是藉著女人而生；但萬有都是出於神。 \end{tabularx} \\ \\ \relax
11:13 & \begin{tabularx}{0.7\textwidth}{X} 你們自己要判斷，女人禱告神，不蒙著頭合宜嗎？ \end{tabularx} \\ \\ \relax
11:14 & \begin{tabularx}{0.7\textwidth}{X} 你們的本性不也教導你們，男人若留長頭髮是他的羞辱嗎？ \end{tabularx} \\ \\ \relax
11:15 & \begin{tabularx}{0.7\textwidth}{X} 但女人留長頭髮是她的榮耀，因為這頭髮是給她蓋頭的。 \end{tabularx} \\ \\ \relax
11:16 & \begin{tabularx}{0.7\textwidth}{X} 若有人想要辯駁，我們卻沒有這樣的規矩，神的眾教會也沒有。 \end{tabularx} \\ \\
[1ex]
\hline
\hline
\end{longtable}
$^{1}$會處理一下女人.
正如我上星期所說.
這個議題我們很多人都應該感到切身.
因為特別說到《聖經》.
或者《聖經》.
到底應該怎麼看.
也有很多不同的新聞.
或者不同的消息等等.
所以按照弟兄姊妹所給我的回應.
我就決定繼續說下去.
好嗎?.
繼續說下去.
或者如果今天之後.
大家都是有要求的話.
都可以繼續說下去.
是吧?.
因為說女人可以說很久的.
所以今天我們繼續說一下.
處理另外一段經文.
這些經文都是比較難的.
首先和大家解釋清楚.
雖然比較難.
或者你剛剛來到教會.
或者你可能還沒接觸過《聖經》.
你發現這些經文裡面說的東西.
都不知道他在說什麼.
這是很正常的.
首先大家不要被它嚇怕.
因為說到最後的時候.
大概25分鐘左右.
那個應用一定是和大家有關的.
好嗎?.
就算你聽不明白.
最後的時候.
大家一定會有一些應用的地方.
不過雖然打開這些經文的時候.
就算你讀了《聖經》很多年.
都可能不明白.
對嗎?.
有沒有同感?.

$^{41}$《聖經》有時候是挺難明白的.
好嗎?.
原因不是因為他特意寫得這麼難.
原因就是我們和《聖經》的受眾.
那個時代和文化是有不同的地方.
正如有不同的地方.
所以我們需要下一點功夫.
才能夠明白當時的人.
他們怎麼了解經文的意思.
然後我們今天才可以說.
有什麼原則裡面.
他是藉著當時的文化.
當時的語言來表達的.
那些原則可以應用在我們現在這個時代.
好嗎?.
所以大家用腦筋好嗎?.
大家看完《世界盃》看得這麼精彩.
應該腦筋轉速很快的.
是嗎?.
OK.
好.
那我們來吧.
我選了這段經文是哥林多前書第十一章.
第十一章第二節到十六節.
比較長一點.
這麼多字.
不過不要緊.
我帶大家一起看.
十一章第二節.
這裡哥林多前書.
就是保羅寫給哥林多的教會.
或者哥林多的宗教會.
當時的一些家庭教會.
他在這封書信裡面.
講了很多其他東西.
但是到十一章的時候.
他就討論到蒙頭的問題.
女人要蒙頭.
就是有veil over her head的問題.
十一章第二節.

$^{81}$我稱讚你們.
因為你們事事紀念我.
又堅守我傳給你們的傳統.
但我希望你們知道.
基督是每個男人的頭.
男人是女人的頭.
神是基督的頭.
每個男人禱告或宣講神的訊息.
如果蒙著頭.
就是羞辱自己的頭.
每個女人禱告或者宣講神的訊息.
如果不蒙頭的話.
就是羞辱自己的頭.
因為這如同剃了頭一樣.
女人如果不蒙著頭.
就應該把頭髮剪掉.
女人如果以剪髮剃頭為可恥的.
就應該蒙著頭.
男人不應該蒙著頭.
因為他是神的形象和榮耀.
女人卻是男人的榮耀.
因為男人不是從女人而出.
女人卻是從男人而出.
並且男人不是為女人而做.
女人卻是為男人而做.
因此為天使的緣故.
女人頭上應該有權柄的記號.
但是在主裡面.
女人不可沒有男人.
男人也不可沒有女人.
這個我贊成.
因為正如女人是從男人而出.
照樣男人也是藉著女人生出來的.
但萬有都是出於神.
你們自己判斷.
女人向神禱告.
不蒙著頭恰當嗎.
難道你們的本性沒有教導你們.
男人若留長髮是他的羞辱.
女人若留長髮卻是他的榮耀嗎.

$^{121}$因為刺激他的頭髮.
就跟蒙頭的外衣寬領一樣.
如果有人想爭辯.
我們可沒有這種習俗.
神的宗教會也都沒有.
好聽嗎?這段經文.
好聽吧?.
你不知道他在說什麼.
但你先不要被他嚇怕.
因為我們不知道他在說什麼.
不代表他沒有意思.
我們不知道他在說什麼.
或者爭辯的時候.
原因可能是我們跟他的文化差距不同.
或者我們不知道他們當時的處境.
有時候看聖經.
特別是看新約的一些書信.
就好像聽人說話.
你只聽到一邊.
聽不到另外一邊.
你只聽到另一邊.
他就是這樣.
是的.
你說得對.
那邊到底在說什麼.
你就不太清楚.
所以我們先不需要被他嚇怕.
我們首先想想.
到底當時他的背景.
或者那邊的電話在說什麼.
以致他這邊有這樣的回應.
而我們聽到他的回應.
保羅正在回應.
哥林多教會裡面的一些問題.
哥林多教會裡面當然有很多的問題.
在哥林多的書信裡面.
你也知道這個教會是問題多多的教會.
所以叫哥林多.
所以我們理解一下他的前文後理的時候.
我們就可能會理解多一點.

$^{161}$剛才我們說那段經文.
我讀出來的時候.
大概你明白的意思.
女人需不需要蒙頭.
大概明白的意思是這樣的.
正因為這樣的緣故.
所以有部分教會.
起碼我聽過的一些小部分教會.
他們真的會這樣做的.
他們會說女人講道.
女人上台講道.
或者去做一些神的話語的教導的時候.
她需要頭上有一條頭巾.
要包一條頭巾上去.
其實是表面的意思去應用經文的意思.
就是這麼簡單.
所以我們表面上就覺得應該是這樣處理的.
所以有些教會就會照辦主管.
經文說什麼我們就照樣去應用.
照樣這樣說.
所以這段經文我們一般的理解是這樣的.
一般理解.
一般理解就是為什麼他們要蒙頭.
就是因為應該是女人的問題.
女人有什麼問題呢.
女人的問題可能就是.
你可以想像一下當時的背景.
當時以男為上的文化.
以父權的社會.
就不像我們今天比較男女平等的一種文化.
你可以理解.
當時有些哥林多的女性.
她們信了主.
然後她們從那些父權的文化裡面.
以男為上的文化裡面.
她們聽到這個福音.
然後她發現了保羅所講的福音.
耶穌基督的福音.
原來給我們一個新的身份.
在基督裡面我們大家都是同一的.

$^{201}$在基督裡面我們大家男女是平等的.
所以當時女人可能就是公開場合的時候.
她們就需要蒙頭.
然後她們去到教會的時候.
她們就覺得.
因為我們有一個這樣的身份.
在基督裡面我們就有一個新的身份.
所以我們乾脆就除了我們的頭巾.
然後就披頭散髮.
因為我們已經有了新的身份.
新的身份.
所以可能有這樣的問題.
所以保羅就說不如你們蒙頭.
就不要破壞大家公開的見證.
因為敬拜,聚會是一些公開的事情.
所以如果有些人進來的時候.
在社會裡面進來的時候.
他看到那些女人是這樣的時候.
那就很忌諱的.
那就很麻煩的.
正因為這樣的緣故.
所以通常我們的解說就是說.
一定是女人的問題.
是女人有問題.
所以就需要要求她去到蒙頭.
如果你這樣想的時候.
其實也不是這麼麻煩的.
我覺得.
不是這麼難明白的.
你可以想像一下.
如果你是一些重男輕女的環境裡面出生的話.
女人經常被歧視的.
男人經常看低女人的.
其實我們華人的文化裡面.
我們傳統都有這樣的思維.
女人無才便是德.
對不對.
如果你拿了幾個PhD.
你是女人的.
為什麼要拿這麼多的PhD.

$^{241}$是不是.
男人就很自然的.
我們學識上面是高一點等等.
女人就應該怎樣.
無才便是德.
你少認識一些東西.
或者你只是躲在廚房裡面就好了.
等等等等.
或者女人一定要跟尾的.
如果你經常聽這些東西.
或者在這種環境裡面出生的時候.
其實也挺大壓力的.
有沒有同感.
有吧.
如果你是女生從小到大.
你要做某一件事.
或者你在一個盒子裡面.
一個框架裡面.
你才會受人接納.
或者受人尊重的.
你一聽到這個福音的時候.
你就會覺得一個所謂解放.
是不是.
一個Liberation.
原來不需要這樣的框架.
因為這些框架不一定是代表真理的.
耶穌基督的福音告訴我們.
男女平等.
教會是一個新的社會.
一個基督裡面新的人類.
無分貴賤的.
無分男女的.
這樣是很吸引的.
對不對.
你知道早年的時候.
大概200年前左右.
西方的宣教士去到東方的時候.
譬如進入中國的時候.
他們其中一件事想廢除的是什麼.
就是女人紮腳.

$^{281}$是不是.
因為西教士覺得這些太殘忍了.
太殘忍了.
而且這樣是控制女性的.
所以他們第一件事.
他們希望能夠廢除這些東西.
以及建立一些女子中學.
女子的小學,中學等等.
因為女孩子是不能讀書的.
是不是.
所以很多西教士去到東方的時候.
或者去到這些國家的時候.
他們就想改變社會的制度.
是什麼緣故呢.
是因為福音的緣故.
是因為福音給他們的精神.
所以他們覺得男女是要平等的.
所以我們看這段經文的時候.
我們可能會想.
一般的解說就是女人有問題.
女人應該要蒙頭.
因為當時的髮色是這樣的.
這些我們西化了的社會.
我們不是很明白的.
因為我們通常衣著,衣服,頭髮,長短.
這些我們不太理會的.
或者我們比較開放一點.
但是比較保守的一些地區.
或者東方的社會會比較緊張一點.
譬如去中東.
去中東.
你進去一些特別的聖堂的地方.
你不可以露這麼多肉.
是不是.
皮膚不可以露這麼多.
上一次我們去土耳其的時候.
我就經常吩咐一些弟兄姊妹.
我說我們今天會去聖堂.
特別吩咐一些姐妹.
就會說你們穿衣服要保守一點.

$^{321}$雖然土耳其是一個很現代化的地方.
但是你去一些聖堂.
譬如清真寺.
你就需要保守一點.
因為那裡有它的規矩.
當然你要尊重別人.
尊重別人的時候.
你才可以進去.
誰知道我自己忘記了.
我就穿了一條短褲.
我就忘記了.
我一不給別人吩咐就因為熱.
我就很自然就穿了一條短褲.
去到門口的時候.
入口的看監就拿了一塊布給我.
你圍著它.
這次我有一點點差池.
因為我通常都記得.
這麼多年我都記得.
結果我就要圍一條布.
我就發現很多人跟我拍照.
你看到很多人跟我一起拍照.
我很開心.
不是因為他們覺得我上鏡.
而是因為我穿了他們的裙子之後.
我就覺得特別上鏡.
所以東方的社會.
有時候會比較保守一點.
好不好看?.
挺好看吧?.
好看.
穿多一次吧.
我應該上台的時候穿這條.
尊重一下大家.
因為白手.
所以我想明白.
衣著或者外表在社會裡面確實比較緊張一點.
所以你可以明白.
一般的理解這段經文的時候.
就是女人有問題.

$^{361}$她們索性非投散法.
或者在聚會裡面.
保羅就覺得你過了火.
你會有損教會在社會裡面的聲譽.
所以保羅最緊張的是什麼?.
保羅最緊張的是.
基督徒的敬拜,我們的敬拜.
是一個公開的場合.
我們就應當要有一個公開的,好的見證.
不要被人誤會福音.
所以他就要求女人是蒙頭的.
這個是一般的解說.
但是如果你想清楚.
裡面有一點這段經文.
裡面有一些前言不對後語.
就是保羅到底是要蒙頭還是不要蒙頭呢?.
我們其中一個方法就可以找一下.
當時的社會裡面.
到底女人是不是真的個個都是蒙頭呢?.
今天你可以去博物館看看羅馬的雕像.
頭部的雕像.
有男有女的.
是不是每個女人都是蒙頭呢?.
那又不是.
那是不是男人個個都不蒙頭呢?.
又不是.
有時候男人有蒙頭.
所以就有一點混淆的證據.
當時的習俗到底是怎樣的呢?.
就比較難去找到一個一致的背景的看法.
或者可能我們看聖經的時候看得不通透.
會不會其實不應該這樣看呢?.
我自己一直以前這麼多年.
我看這段經文.
我都覺得應該是那個解說.
就是那麼多的女性是有問題的.
就是她們覺得被解放.
然後她們就很自由,很開放.
保羅就在那個情況裡面去做一些改正.
但是呢.

$^{401}$會不會有另外一個提問我們可以發出呢?.
就是保羅這裡會不會其實不是說女人有問題.
其實可能是男人有問題呢?.
會不會是這段經文裡面.
保羅是在引用哥林多的一些說法.
或者哥林多男人的一些說法.
然後他引用的當中.
他就去到加插他自己的反駁在那裡.
這個早年的時候.
我就看了一本書.
就拜讀這個女色經家.
是英國人來的.
Lucy Papiat.
她就認為不應該這樣看這段經文的.
為什麼要看一個女色經家寫的書呢?.
就很簡單.
因為連色經界裡面.
都很多時候是男人寫書多的.
可能你的印象也是.
但是近年裡面已經越來越多好的.
非常頂尖的女色經家出來了.
女色經家問的問題就跟我們男人問的問題.
很多時候自然就不同的.
因為如果我看哥林多詞書11章.
我們覺得女人有問題的.
我作為一個男人.
我當然覺得很舒服的.
當然女人有問題.
我們經常都是這樣.
我們覺得只有女人有問題.
但是如果女色經家說的時候.
可能她的角度就會很不同.
她就會問.
可不可以是另外一邊呢?.
這些絕對不是.
強行去將字句曲解的.
而是用另一個角度去看一看.
會不會有這樣的可能性呢?.
是保羅去引用哥林多男人說的話.
然後再加以反駁.

$^{441}$其實哥林多前書.
或者其他保羅的書信.
都會引用一些對手所說的話.
譬如羅馬書等等.
在哥林多前書都有的.
譬如哥林多前書6章12節.
「凡事都可行,但不是凡事都有益處.」.
「凡事都可行,但無論那一件,我都不受他克制.」.
你看到那個引號加上去了.
那個原文是沒有的.
以前我年輕的時候看中文聖經.
或者甚至英文聖經.
都沒有那個引號在那裡.
沒有引號的時候.
我們就會覺得.
當然是保羅出口.
因為這封信是他寫的.
保羅一定是6章12節的時候.
他說「凡事我都可行,但不是都有益處.」.
凡事我都可行.
在基督裡面凡事都有自由.
凡事都可行.
但無論那一件,我不應該被他克制.
這句應該是保羅出口.
但現在我們慢慢開始有一個共識.
原來保羅有一些修辭上的東西.
是在他的書信裡面用的.
凡事都可行.
其實應該是哥林多教會的口號.
因為哥林多教會.
男人特別.
他們就會覺得我們信了主.
我們就是.
anything we can do.
everything is permissible.
所有事情我們都做得到.
有什麼所謂.
我們基督裡面有自由.
但現在我們慢慢有共識.
其實保羅是引用哥林多教會裡面所講的東西.

$^{481}$然後他就反駁.
你們說凡事都可行.
但不是凡事都有益處.
你想清楚.
不是每件事都可以做.
你最重要記住.
那件事不要被他克制.
世界盃可以看.
但不可以被世界盃控制你.
對不對.
你的情緒也好等等.
所以他是在引用哥林多的說話.
但用這個角度看的時候.
你再看哥林多的前書.
可能就不同了.
我在這裡再讀這段聖經.
橙色的地方.
可能是哥林多的男人講的.
第二節.
我稱讚你們.
因為你們時時都紀念我.
凡事都堅守我給你的傳統.
但我希望你們知道.
基督是每個男人的頭.
男人是女人的頭.
神是基督的頭.
但你們說.
每個男人禱告或宣講神的訊息.
如果蒙著頭就是羞辱自己的頭.
如果女人講道或禱告宣講神的訊息.
如果不蒙著頭.
就是羞辱自己的頭.
如同剃自己的頭一樣.
保羅就說.
如果你們說了這些話.
這樣的話.
那就第六節了.
女人如果不蒙著頭.
就應該將頭髮剪掉.
女人如果以剪髮剃頭為可恥.

$^{521}$就應該蒙著頭.
如果你們這樣說四五節的話.
第六節你們自然有這樣的結論.
第七節.
你們哥林多人這樣說.
男人不應該蒙著頭.
因為他就是神的形象和榮耀.
女人卻是男人的榮耀.
因為男人不是從女人而出的.
女人卻是從男人而出的.
並且男人不是為女人而做.
女人卻是為男人而做.
這些你們說的.
這些你們哥林多人說的.
因此為天使的緣故.
女人的頭上就應該有權柄的記號.
這些全部橙色的都是你們說的.
保羅就跟他們說.
如果你們說這些話.
結論就會是這樣.
我們繼續看下去.
現在我就要矯正你.
十一節.
保羅說什麼呢.
他說但是在主裡面.
女人不可沒有男人.
男人不可沒有女人.
我們每個人都有母親.
我們是需要女人的.
所以男女是平等的.
十二節.
因為正如女人是從男人而出.
照樣男人也是藉著女人生出來.
但當然萬有就是出於上帝.
你們自己判斷一下.
這裡就直接跟哥林多教會說.
你們自己的判斷.
女人向神禱告.
不蒙著頭恰當嗎.
當然是恰當.

$^{561}$不蒙著頭當然是恰當.
為什麼.
難道你們的本性沒有教你們.
男人如果留長頭髮就是他的羞辱.
女人如果留長頭髮.
卻是他的榮耀嗎.
因為刺逼他的頭髮.
就是跟他蒙頭的外衣寬領一樣.
因為他已經有長頭髮了.
所以他已經好像他的外衣蒙頭一樣.
所以就不需要再蒙頭.
第十六節.
這裡就最清楚.
他說如果你們想再爭辯的話.
我們使徒們.
可沒有這樣的習俗.
神的宗教會.
其他地方的教會也沒有.
所以希望梳理了一點這些問題.
到底問題在哪裡.
如果你這樣看的時候.
你就發現他在引用哥林多的對手說的話.
那些男人說的話.
然後他寓意去反駁.
就說你們這樣想是不對的.
你們這樣想最後的結論都是自相矛盾的.
所以他就說.
我們根本就沒有這樣的習俗.
宗教會也沒有這樣的習俗.
誰有這樣的習俗.
哥林多的教會有這樣的習俗.
你這樣看的時候.
這個看法有什麼可以梳理.
就是起碼首先.
保羅不會前言不對後語.
好像又蒙頭好像又不蒙頭.
原來他就是在反駁.
這樣看的時候.
保羅就是和哥林多的大男人神學.
去做爭辯.

$^{601}$哥林多教會有大男人.
他們就用了.
曾經保羅說過.
男人是女人的頭.
我們就想去訓練這個神學.
其實就是壞鬼神學.
就是大男人神學.
就是我是男人.
我就不用蒙頭.
但是如果女人要講道.
或者要禱告的話.
就請他蒙頭.
因為女人頭上一定有一個記號.
就是她要可以信服的.
對不對.
所以在反駁的當中.
保羅就是用他的本性來帶出.
女人的長髮自然是榮耀.
自然是榮耀.
他就不需要再加上一塊根來蒙頭.
也解釋到最後那句.
最後那句是什麼.
就是神的宗教會也沒有這樣的習俗.
如果你這樣看的時候.
到底是男人有問題.
還是女人有問題.
如果你說女人有問題的話.
那你就很難.
就有些東西卡住了.
前言不對後語.
還有最重要是十六節.
如果有人和我去爭辯的話.
我們沒有這樣的習俗.
神的宗教會也沒有這樣的習俗.
那會不會應該轉一轉那個角度呢.
其實可能是男人有問題.
對不對.
是男人有問題.
就多不多有一班大男人.
就弄了一個大頭神學出來.

$^{641}$他們就要女人蒙頭.
這些就是我們開頭的時候說.
某些教會今天看這段聖經的時候.
他們就照主管來應用.
如果這樣應用的時候.
除了表面去看聖經.
也不是很明白那個上下文.
或者那個背景.
是不是.
這樣就會重複了.
哥林多教會的大男人蒙頭神學.
所以保羅寫這段東西.
其實是在矯正這件事.
十六節也說得通了.
我們沒有這樣的習俗.
神的教會也沒有.
哥林多為什麼會有呢.
所以整段經文其實是針對大男人.
哥林多的大男人.
如果你這樣看的時候.
就配合哥林多的前書整本書.
因為整本書裡面都針對很多.
領袖或者愛權的人.
他們想做的事情.
和他們令到教會混亂的一些事項.
十一章就是處理女人蒙頭的問題.
如果這樣看的時候.
這個是男人的問題.
保羅就是斥責他們.
然後就說到聖餐的問題.
我們平時守聖餐的時候.
不是會讀一段聖經嗎.
在十一章後面出來.
那裡就是有些混亂的情況.
因為他們分開貧富.
守聖餐有些富有的人早到了.
他們就吃飽東西.
然後貧窮的人.
他們下班的時候.
他們來到聚會的時候.

$^{681}$他們就沒有東西吃.
他就去改正他們聖餐所有的混亂.
十二章就說到才幹和恩賜的使用.
保羅就去罵哥林多的人.
他說你們有一些的恩賜.
譬如方言的恩賜.
就看低其他的人.
他就去予以改正.
十三章就是愛篇.
愛是恆久忍耐又有恩賜那篇.
很出名的那篇.
經常婚禮會用的那篇.
為什麼他會說這些.
因為上文就是在說.
你們都不懂愛.
我現在告訴你什麼是愛.
愛就是恆久忍耐又有恩賜等等.
是為其他人緣故的.
不是爭權奪利的.
然後十四章十五章.
就去處理其他的問題.
你這樣看整封信.
其實保羅就好像一個家長.
那些教會開始長大了.
他就不只是告訴他們.
Do's and Don'ts.
他就是告訴他們.
有一個新的人類.
在基督裡面誕生了.
這個就是教會.
在基督裡面有的弟兄或者姐妹.
我們都是平等的.
大家是按恩賜來事奉的.
因為上帝要在我們中間.
有一個新的創造.
但是那不多的教會.
很可能就是很多大男人.
愛權的領袖等等.
他們就做了很多這樣的事情.
在制度上或者習慣上.

$^{721}$就會令到其他人覺得.
他們是二等公民.
保羅就在這本書信裡面.
來做改正.
所以如果你這樣看的時候.
似乎就不是要女人蒙頭.
而是保羅相反.
就說不應蒙頭.
不應該蒙頭.
而那些大男人.
為什麼搞這些東西出來.
宗教會都沒有這樣的習俗.
為什麼你要搞這些東西出來.
因為你們亂去應用《聖經》.
這個值得我們去思考一下.
怎樣去看《聖經》.
有幾點是我們可以點出來的.
第一就是結論的時候.
第一就是在基督裡面.
男與女我們有的恩賜.
有的才幹.
是應當可以拿來服侍的.
如果有一個女性.
有一個姐妹.
她受上帝的呼召.
或者有信仰群體.
去確認她有這樣的恩賜的時候.
我是何許人我去攔阻那位姐妹.
如果她有教導的恩賜.
如果她有《聖經》的恩賜.
如果她有廣道的恩賜.
她有蒙上帝去呼召她.
又有她一個好的見證.
其他人在她的群體上.
對她有一個引證的時候.
我是誰去阻止她.
對嗎.
我夠神嗎.
當然不是.
在基督裡面已經是新的群體.

$^{761}$在基督裡面我們能夠平起平坐.
是按照我們的恩賜.
不是按照我們的性別.
是按照我們的禮物.
不是按照我們的性別.
去事奉.
這個是基督有的意思.
第二件事我們想到的.
在讀經裡面.
我們看《聖經》裡面或者應用《聖經》的時候.
我們真的要去下一點功夫.
明白下上問下理.
不要只是抽一些東西.
然後我們隨便去應用.
這樣說不是將《聖經》困在當時的文化.
而我們理解《聖經》是藉著當時的文化.
透過當時的背景.
來轉達上帝永恆的信息.
所以我們一定不能跳過那一步.
我們不理當時的背景.
是不是.
我曾經說過.
我覺得每一次看《聖經》的時候.
都好像去旅遊.
去一個你陌生的地方去旅遊.
一個從來沒有去過的地方去旅遊.
你就聽一下別人說什麼.
聽一下別人的習俗.
聽一下別人的文化背景是怎樣.
聽一下那裡的信息是怎樣.
聽一下別人是怎樣接受.
看《聖經》其實都是一樣的.
你去哪裡旅遊.
去一個《聖經》的世界.
你去那裡聽一下別人說話.
明白別人的背景.
然後你知道上帝所轉達的永恆信息.
到底是什麼.
這是第二.
第三樣.

$^{801}$第三樣就是我們需要反問我們自己.
有什麼藉著《聖經》可以反問我們自己.
到底我們有沒有濫權.
我們的實踐上面.
甚至我們的制度上面.
我們有沒有一些濫權的問題.
如果我們濫權的時候.
或者亂用我們的權柄的時候.
或者限制其他人的時候.
這個和基督的福音是不是有違反.
我們自己生活.
我們有沒有好像基督那樣謙卑.
有沒有濫用我們自己的力量或者地位.
第四樣.
因為我們說到女性.
我們不妨思想一下.
在我們人生裡面對你生命有影響的女性.
有誰你想到.
你想一下她的樣子.
想一下她的名字.
想一下她做過什麼.
可能是牧師.
可能是一些傳道人.
可能是導師.
或者可能是朋友.
一些姐妹.
在你生命裡面有很多很多這樣的人.
其實是祝福過你的.
但這些都是上帝所用的生命.
去感激他們.
對不對.
特別可以跟大男人說話.
在我們生命裡面有沒有一些這樣的人.
對不對.
我想起我媽咪.
想起在我教會以前特別年輕的時候.
曾經傾注過她的生命.
在我的成長裡面.
這些可能是女傳道人.
可能是一些導師.

$^{841}$可能是身邊的一些姐妹.
對我生命裡面有影響的.
我們低頭祈禱的時候.
我們不妨思想一下這個問題好不好.
給大家一點時間.
有沒有一些姐妹一些女性.
在生命裡面有一個影響.
因為你做男做女.
在基督的救贖裡面.
在基督的恩典當中.
我們無分地位.
你去到呼召男去到服侍你.
也呼召姐妹.
在歷史歷代裡面有很多女性.
是你自己親自呼召出來的.
每一個都是被你所使用的.
可謂生命當中.
你藉著這幾天.
也去祝福我們生命.
去造就我們.
我們獻上感恩讚美.
雖然我們只是說一些女性事奉的經文.
但它是對我們每一個人有應用的.
你幫助我們.
讓我們用我們的恩賜.
也去改變我們.
矯正我們心裡面.
或者在我們背景裡面.
有一些不正確的事情思想.
主佑我們求你.
求你讓我們在基督裡面.
我們真的享有這個自由.
在基督裡面我們能夠.
在我們角色功能上面.
在我們服侍上面.
我們沒有分男與女.
讓我們一起在神的國度裡面.
我們去服侍你.
就如基督去到服侍教會一樣.
為了我們生命裡面.

$^{881}$有一些影響過我們的女性.
我們獻上感恩讚美.
多謝他們謙卑的服侍.
多謝他們所說的話.
多謝他們所做的事.
甚至默默無言的在我們身邊.
在我們生命裡面.
去傾倒他們的生命.
主佑我們獻上感恩.
獻上讚美.
求你讓我們這個群體.
在基督裡面形成一個教會.
我們是一個高舉.
每一個人.
在基督裡面.
我們去榮耀神地.
我們去體重每一個人.
讓我們的事奉.
在基督裡面得目稱讚.
我們這樣祈禱奉祖明教.
願主祝福各位.
下星期繼續回來聽.
\newpage

\allsectionsfont{\centering}

\setlength\parindent{0pt}
\setlength{\columnsep}{1.25em}
\setlength{\parfillskip}{0pt}
\setlength{\tabcolsep}{1em}
\raggedbottom

\pagenumbering{gobble}


\newfontfamily\leftfont[Path=../fonts/fell_french_canon/, Ligatures=TeX, ItalicFont=IMFeFCit29C.otf, BoldFont=AveriaLibre-Bold.ttf]{IMFeFCrm29C.otf}
\newfontfamily\leftcitationfont[Path=../fonts/frankruehl/]{FrankRuehlCLM-Medium.ttf}
\newfontfamily\centerfont[Path=../fonts/garamond/, Ligatures=TeX, ItalicFont=EBGaramond-SemiBoldItalic.ttf]{EBGaramond-SemiBold.ttf}
\newfontfamily\rightfont[Path=../fonts/averia/, Ligatures=TeX, ItalicFont=AveriaLibre-RegularItalic.ttf, BoldFont=AveriaLibre-Bold.ttf, BoldItalicFont=AveriaLibre-BoldItalic.ttf]{AveriaLibre-Light.ttf}
\newfontfamily\rightcitationfont[Path=../fonts/rashi/]{Mekorot-Rashi.ttf}
\definecolor{hcolor}{HTML}{D3230C}
\definecolor{rcolor}{HTML}{D36F0C}
\newcommand{\chfont}[1]{\centerfont{\huge\textcolor{hcolor}{#1}}}
\newcommand{\leftcitation}[1]{\leftcitationfont{\Large\textcolor{hcolor}{#1}}}
\newcommand{\rightcitation}[1]{\rightcitationfont{\normalsize\textcolor{rcolor}{#1}}}
\newfontfamily\flowerfont[Path=../fonts/fell_flowers/]{IMFeFlow2.otf}

\begin{sloppypar}

\chapter*{\chfont{編按結語}}

\columnratio{0.5,0.5}\begin{paracol}{2}

\fontsize{11}{13}\leftfont \Large \leftcitation{א} \leftfont 余少好文．宏志博覽群書而不忘．善存藏經籍文獻備後時之用。\leftcitation{ב} \leftfont 歸主年時．受友所薦．聞道網海．\switchcolumn\fontsize{11}{13}\rightfont \Large \leftcitation{ח} \rightfont 有見粵道之危．國之封講道千言亦將就至．急之何則為？\leftcitation{ט} \rightfont 嘗聞猶太者之傳承．在其力守口述之

\end{paracol}


\columnratio{0.32,0.32,0.32}\begin{paracol}{3}

\fontsize{11}{13}\leftfont \Large 尤以吳約翰遜者 \switchcolumn[2]\fontsize{11}{13}\rightfont \Large 統．以煉千載不

\end{paracol}

\columnratio{0.32,0.32,0.32}
\begin{paracol}{3}\fontsize{11}{13}\leftfont \Large 為重．其載上之粵語講道緩緩入耳．收之藏其音頻．善妥整存．反復而嚼．受益無窮。\leftcitation{ג} \leftfont 我城我國既限．歷一四一九之不測．肺疫延年．信徒靈長屢受圍創．神州燈臺數盡指日可待．粵道之求與日俱增。\leftcitation{ד} \leftfont 觀乎社、經、法、媒、言、信、網之地．愈趨受鋤．自翔不果．授受壓力．粵道聖言亦愈漸艱難。\leftcitation{ה} \leftfont 況崇基例乎．學苑講道屢逆權勢者．其言末強受壓．舊章盡刪以存其身。講道釋數失傳．徒嘆奈何。

\switchcolumn

\fontsize{11}{13}\centerfont 
\begin{tikzpicture}
    \node (0,0) [xshift=-0.10cm, yshift=-1.0cm, opacity=0.10]{\includegraphics[width=0.30\textwidth]{../ot_frontcover.png}} ;
    \node (0,0) [xshift=+0.20cm, yshift=+2.0cm, opacity=0.10]{\includegraphics[width=0.20\textwidth]{../christ_on_cross.png}} ;
\end{tikzpicture}
\Large 

\leftcitation{ס} \centerfont 詩百又廿七載：
\leftcitation{ע} \centerfont 非耶和華建屋宇.則匠人之經營徒.
\leftcitation{פ} \centerfont 非耶和華衛城邑.則守者之儆醒徒.
\leftcitation{צ} \centerfont 余獻是卷予華人社區．願為福音流通之器．願獻斯微材為祭榮耀上帝．
\leftcitation{ק} \centerfont 阿門

\switchcolumn

\fontsize{11}{13}\rightfont \Large 滅．時越次聖殿期及當今。\leftcitation{י} \rightfont 猶太者力廣納之．筆錄以卷軸．便以傳、閱、頌、攜、守、鎖、抄、譯、釋、編，得書塔木德、密示拿等經傳．家喻戶曉．傳流若芳。\leftcitation{כ} \rightfont 猶太者文以載道．傳其口述．今我輩粵道之傳應當作如是．遂力行粵音識辨之法．載言載道．以盡忠傳粵道以待興。\leftcitation{ל} \rightfont 蒙下賜恩惠．無畏海量字音文書．既馭上帝之道．今廣及粵語講道．重駛編程之技．匯導粵音遂字稿．重塑講道現場．以傚猶太卷軸之舉便以傳流。\leftcitation{מ} \rightfont 是卷乃粵音口述傳之屬．莫通華文白話之語．

\end{paracol}

\columnratio{0.5,0.5}
\begin{paracol}{2}\fontsize{11}{13}\leftfont \Large \leftcitation{ו} \leftfont 斯殺一違儆百逆．既禁壓之．我輩聞風無奈．在所難免。\leftcitation{ז} \leftfont 另有異人例乎．以版權之名．脅網絡頻道之舉．同授礙予粵道之存流。

\switchcolumn

\fontsize{11}{13}\rightfont \Large 惟待後繼來者之傚．以譯釋傳之於神州華文地。\leftcitation{נ} \rightfont 今能排程驅馭圖靈以編彙文檔，其碼長共數千千亦無逢大礙．全蒙上帝保守。

\end{paracol}



\columnratio{1}\begin{paracol}{1}

\fontsize{11}{13}\rightfont \Large
~~~~~~~~~~~~~~~~~~~~~~~~~~~~~~~~~~~~~~~~~~~~~~~~~~~~~~~~~~~~~~~~~~~~~~~~~~~~~~~\leftcitation{ר} \rightfont 二零二三年二月一日

~~~~~~~~~~~~~~~~~~~~~~~~~~~~~~~~~~~~~~~~~~~~~~~~~~~~~~~~~~~~~~~~~~~~~~~~~~~~~~~\leftcitation{ש} \rightfont 米迦勒

~~~~~~~~~~~~~~~~~~~~~~~~~~~~~~~~~~~~~~~~~~~~~~~~~~~~~~~~~~~~~~~~~~~~~~~~~~~~~~~\leftcitation{ת} \rightfont 書於香港

\end{paracol}

\end{sloppypar}
\end{document}
